\chapter{Assembly of the proof}\label{chapter-assembly}

This chapter combines the results of the previous chapters into the
half-plane statement of Chapter~\ref{chapter-main-theorems}. For a primitive
character $\psi$, families of definite-side data are built over the host
fields of Chapter~\ref{chapter-host-fields}, and at every large scale of such a
family the bounds of Chapters~\ref{chapter-hermitian-slopes}
to~\ref{chapter-tate-shafarevich} verify the hypotheses of the comparison
inequality of Chapter~\ref{chapter-inert-prime-balance}. The window bound that
results is turned into a zero-free half-plane by
Chapter~\ref{chapter-windowed-landau}. The last section lists the parameters
and says which bound feeds which inequality.

\section{The order of the argument}\label{section-asm-order}

Read from the top, the proof of the half-plane statement has the following
steps. Each is a statement of this chapter.
\begin{enumerate}
\item The Galois representations of
  Chapter~\ref{chapter-galois-representations}, with local-global
  compatibility relative to every map satisfying the local Langlands laws,
  give the bound on Tate-Shafarevich groups of
  Chapter~\ref{chapter-tate-shafarevich} and, together with the results of
  Chapter~\ref{chapter-literature} and the covolume bound of
  Chapter~\ref{chapter-covolumes}, all remaining hypotheses
  (Theorems~\ref{asm-every-character}
  and~\ref{asm-every-character-literature}).
\item These results produce a collection of definite-side hypotheses
  (Definition~\ref{asm-definite-side-hypotheses},
  Theorem~\ref{asm-definite-side-exists}): a class of families of
  definite-side data together with the inputs that the links of
  Chapters~\ref{chapter-period-divisibility} and~\ref{chapter-integral-structure}
  need.
\item With Green's parametrization of the cuspidal characters of
  $GL_d(\mathbf F_g)$, such a collection gives tied families for every host
  requirement (Theorem~\ref{asm-tied-from-green}). The proof finds designed
  families (Theorem~\ref{asm-design-exists}), establishes the links eventually
  along them, and concludes by Theorem~\ref{asm-tied-from-links}, whose last
  input is the global link of Theorem~\ref{asm-global-link}.
\item For each primitive character $\psi$, the host fields of
  Chapter~\ref{chapter-host-fields} meet the host requirement of $\psi$ on a
  set of scales with windows of exponent $3/2$. A family meeting it that is
  tied at its large scales satisfies the window bound
  $G_{K_\psi}(x)\le x^{1/2+9\varepsilon}$ of
  Chapter~\ref{chapter-inert-prime-balance}, and
  Chapter~\ref{chapter-windowed-landau} turns the window bounds for all small
  $\varepsilon$ into $L(s,\psi)\ne0$ for $\operatorname{Re}s>1/2$
  (Theorem~\ref{asm-zero-free-primitive}).
\item The Euler factors at the primes dividing the modulus pass from primitive
  characters to all characters (Theorem~\ref{asm-every-character-green}).
\end{enumerate}

\section{Ties and designed families}\label{section-asm-families}

Families of definite-side data, their members at each scale, their exponent
$\varepsilon$, classes of families and properties that hold eventually along a
class are those of Definition~\ref{pd-family}, with the size data of
Definition~\ref{pd-size-data}; at a scale $(x,\varepsilon,J)$
(Definition~\ref{hf-windows}) we write $L=x^{\varepsilon}$, so that
$b=e^{L}$ and $a=xb$. A family has \emph{windows of exponent $\lambda$} if its
set of scales does (Definition~\ref{wl-windows}). The comparison of
Chapter~\ref{chapter-inert-prime-balance} takes place at the size $a$, and the
primes that enter it lie below $X''(n)=xnX'$.

\begin{definition}\label{asm-tied}
Let $\mathcal F$ be a family with exponent $\varepsilon$
(Definition~\ref{pd-family}), let $x$ be real, and let
$e_D,e_L,K,c_{HH},e_K,C_L$ be real numbers. The family is \emph{tied at $x$}
with these constants if there are
\begin{itemize}
\item a member $i_0$ at $x$, whose Galois datum (Definition~\ref{hf-frame})
  has group $A$, with ranks $m$ and $f$, and a map $\mathrm{pt}$ from $A$ to the
  members at $x$;
\item integer sequences $(a_t)$ and $(b_t)$;
\item comparison data $\mathcal D$ on $(A,W,\mathcal T,m,f)$
  (Definition~\ref{ib-comparison-data}), where $W(u_0)$ is the set of words of
  $D_{x,\mathrm{pt}(u_0)}$ and $\mathcal T$ is a finite set;
\item a function $\ell$ on the characters of $E_0\times A$, an integer $N$,
  subsets $T,T_2\subset A$, a function $a$ on $A$, weights $w(u_0,j,\rho)$,
  bijections of $\mathcal T$ with $T$ and of $\{1,\dots,f\}$ with the split
  subgroup $\Psi_A$ of $E_0\times A$ (Definition~\ref{ib-host-group}), and real
  numbers $\eta,X',\theta$,
\end{itemize}
such that the following hold.
\begin{enumerate}
\item For every $u_0\in A$, the member $D_{x,\mathrm{pt}(u_0)}$ has the Galois
  datum, $m$, $f$, $q$ and $d$ of $D_{x,i_0}$, base point $u_0$, and exponents
  $a_t,b_t$ on its top classes.
\item For every $u_0\in A$, the parameters $a,b,b',k,L,L_a,B_0,N_c,v$ of
  $\mathcal D$ are those of $D_{x,\mathrm{pt}(u_0)}$ and its size data;
  $\lambda=\log(g_0g_1)$; $\gamma_0$ is Euler's constant;
  $\ell_0=\operatorname{Re}\frac{L'}{L}(1,\chi_d)$ for the quadratic character
  $\chi_d$ modulo $d$; and $X''=X''(n)$, $\eta$ and $X'$ are the corresponding
  quantities of the size data.
\item For every $u_0\in A$ and every word $u$ at $u_0$: $r_u$ is the rank
  $d_u$ of the division lattice $N_u$; $P_u$ is the auxiliary prime $T_u$;
  $\mu_N(u),\mu_N^+(u)$ are the slope and maximal slope of $N_u$; $h_u$ is the
  height of $u$; $\mu_M(u),\mu_M^+(u)$ are the slope and maximal slope of the
  first lattice $M_u$ of the chain of $u$; and $E_{u,j}=\sum_w\mathrm{sz}_w
  (E_{u,j}(w))$. Further, $\mu_C(u_0)$ and $\mu_C^-(u_0)$ are the slope and
  minimal slope of the dual $CC$ of the glued lattice, $F(u_0)=F_o$,
  $\mathrm{Ga}(u_0)=G_a$, and the local sums of $\tau_w$ and of the losses
  $\iota_w$ and the row deficit $\delta(u_0)$ are those of
  $D_{x,\mathrm{pt}(u_0)}$.
\item $T$ is the set of top classes of the members at $x$, $T_2$ is the
  image of the set $T_2$ of the members, $a(h)=a_h$, the weights
  $w(u_0,j,\cdot)$ are the Harish-Chandra weights $w_j(\tau)$ of
  $D_{x,\mathrm{pt}(u_0)}$, and $\ell$ is compatible with the host:
  $\ell(\psi)=\frac{L'}{L}(1,\chi)$ for every character $\psi\ne1$ and every
  primitive Dirichlet character $\chi$ with $\chi(n)=\psi(n)$ for $n\ne0$.
\item $\mathcal D$ has parameters $x$ and $\varepsilon$, with
  $e_{\mathrm{Div}}\le e_D$, $e_{\mathrm{Loss}}\le e_L$ and $c_{HH}(\mathcal
  D)\le c_{HH}$. With $H=A$, the arithmetic data of the host of
  $D_{x,i_0}$ and $\mathrm{tr}$ the identity of $A$, conditions (1) to (4) of
  Lemma~\ref{ib-gap-bound} hold with constants $(e_K,K)$; the data lie in the size windows
  (Definition~\ref{ib-size-windows}); and for the three nontrivial characters
  $\psi$ of $E_0$, $|g_\psi(x)|\le x^{1/2+\varepsilon}$ and
  $|\ell(\psi)|\le C_L\log x$.
\end{enumerate}
A class $\mathcal E$ of families \emph{has tied members} with exponents
$e_D,e_L$ if there is $\varepsilon_1>0$ such that for every
$\varepsilon\in(0,\varepsilon_1)$ some family $\mathcal F\in\mathcal E$ with
exponent $\varepsilon$ has constants $K,c_{HH},e_K,C_L,x_0$ such that
$\mathcal F$ is tied at $x$ with $(e_D,e_L,K,c_{HH},e_K,C_L)$ for every scale
$x$ of $\mathcal F$ with $x\ge x_0$.
\end{definition}
\leanref{Formalism/SharpRestriction/GlobalSchurBounds/TieGlobal.lean}{Formalism.MainNotions.ChainTied}
\leanref{Development/AutomorphicArithmetic/CharacterBoxes/DivinNOfHypV.lean}{UpperTheory.Top.FamilyClassConnected}

A family tied at $x$ carries, at $x$, core data with exponents $(e_D,e_L)$
(Definition~\ref{ib-core-data}), read off from its members. So
Theorem~\ref{ib-window-bound-on-scales} gives the window bound
$G_K(x)\le x^{1/2+(e_D+e_L+7)\varepsilon}$ at the large scales of a family that
is tied there and satisfies the host conditions for $K$.

The lower bounds for the constants in Definition~\ref{pd-family} come in three
stages. Some constants of a family are not required to be large; a second list,
among them the height constant $C_H$, must exceed a function of the first; and
a third list must exceed a function of the first two. A \emph{largeness
threshold} $\kappa$ is a choice of these lower bounds, and a family is
\emph{large for $(\kappa,C_H)$} if its constants satisfy them. If two
properties hold eventually along a class, so does their conjunction: take the
smaller $\varepsilon_1$ and, at each $\varepsilon$, the larger lower bounds.

\begin{definition}\label{asm-box-row-tame}
Let $D$ be a definite-side datum (Definition~\ref{th-definite-datum}) with
ranks $k,m$ and $r=k+m-1$, whose host field has auxiliary primes $q,d$ and
integer $Q_S$ (Definition~\ref{hf-framed-host-field}) and whose host has level
primes $g_0,g_1$ (Definition~\ref{hf-host}). Let $\varepsilon>0$ and $Y>1$, put
$Q_1=\exp(Y^{\varepsilon/5})$ and $Q_2=\exp(Y^{\varepsilon/5}\log^2Y)$, and let
$Q_S'$ be the odd part of $Q_S$. For $Q>0$, a primitive Dirichlet character
$\chi$ modulo $N$ is \emph{good for $Q$} if it takes real values, $N\le Q$, and
$L(s,\chi)\ne0$ for every $s$ with $\operatorname{Re}s>1-1/400$ and
$|\operatorname{Im}s|\le Q$. Write $\chi_{-q},\chi_{-d},\chi_8$ for the
Kronecker symbols $\left(\frac{-q}{\cdot}\right)$, $\left(\frac{-d}{\cdot}\right)$
and $\left(\frac{2}{\cdot}\right)$.
\begin{enumerate}
\item $D$ satisfies the \emph{box condition at $Y$} if every primitive
  character that agrees with $\chi_{-q}$ or with $\chi_8\chi_{-q}$ on the
  integers prime to $8qd$ is good for $8Q_1$; every primitive character that
  agrees with one of $\chi_{-d}$, $\chi_8\chi_{-d}$, $\chi_{-q}\chi_{-d}$,
  $\chi_8\chi_{-q}\chi_{-d}$ on the integers prime to $8qd$ is good for
  $8Q_1Q_2$; and every primitive real character $\chi_\theta$ of the host field
  (Definition~\ref{hf-host-field-of-character}) is good for $8Q_S'Q_1Q_2$ if
  $d$ divides its conductor, and good for $8Q_S'Q_1$ if $q$ divides its
  conductor and $d$ does not.
\item $D$ satisfies the \emph{row condition at $Y$} if $q$ and $d$ are primes
  with $q\equiv3\pmod8$ and $d\equiv7\pmod{16}$, $\sqrt{Q_1}<q\le Q_1$,
  $\sqrt{Q_2}<d\le Q_2$, $\left(\frac{-d}{q}\right)=1$,
  $\left(\frac{-d}{p}\right)=1$ for every prime $p\le Y^{\varepsilon/5}$, and
  $\log(Q_Sqd)\le2Y^{\varepsilon/5}\log^2Y$.
\item At each level prime $g_k$, $k\in\{0,1\}$, the datum has a set of $f$
  \emph{level places} and, at each level place $v$, \emph{tame values}
  $\chi_{k,v,1},\dots,\chi_{k,v,m}\in(\mathbf Z/g_k\mathbf Z)^\times$, one for
  each lift. The datum has \emph{distinct tame data} if for $k\in\{0,1\}$ and
  every level place $v$ over $g_k$:
  \begin{enumerate}
  \item the $2m+1$ elements $1$, $\chi_{k,v,j}$ and $\chi_{k,v,j}^{-1}$
    ($1\le j\le m$) of $(\mathbf Z/g_k\mathbf Z)^\times$ are pairwise distinct;
  \item for every prime $\ell$ outside the set $\mathcal B_r(g_k)$ of
    Definition~\ref{hf-good-host}, the prime-to-$\ell$ parts of these $2m+1$
    elements, taken in $\overline{\mathbf F}_{g_k}^\times$, are pairwise
    distinct;
  \item every $\chi_{k,v,j}$ is a square in $(\mathbf Z/g_k\mathbf Z)^\times$.
  \end{enumerate}
\end{enumerate}
\end{definition}
\leanref{Formalism/SharpRestriction/GlobalSchurBounds/MainIn.lean}{Formalism.MainNotions.Boxes_at}
\leanref{Formalism/SharpRestriction/GlobalSchurBounds/GeneralLinearActions/FiltrationCocycleLiftCriterion.lean}{Formalism.MainNotions.GlobalTheoryTheory.BoxE}
\leanref{Formalism/SharpRestriction/GlobalSchurBounds/MainIn.lean}{Formalism.MainNotions.GoodAllFor}
\leanref{Formalism/SharpRestriction/GlobalSchurBounds/GeneralLinearActions/FiltrationCocycleLiftCriterion.lean}{Formalism.MainNotions.GlobalTheoryTheory.GoodFor}
\leanref{Formalism/SharpRestriction/GlobalSchurBounds/MainIn.lean}{Formalism.MainNotions.QDRows_at}
\leanref{Formalism/SharpRestriction/LevelLetterCocycles/AlignN.lean}{Formalism.MainNotions.TameDistinct}
\leanref{Formalism/SharpRestriction/LevelLetterCocycles/AlignN.lean}{Formalism.MainNotions.tameList}

The box condition at $Y$ is condition (5) of Definition~\ref{hf-good-host}
with $Q_1$ and $Q_2$ taken at $Y$, and the row condition at $Y$ with
$x^{1/J}\le Y\le x$ says that $Y$ is a pair scale for the host in the sense of
that definition; the congruences on $q$ and $d$ hold for every host field.

\begin{definition}\label{asm-twist-condition}
Let $D$ be a definite-side datum with host field $E$, ranks $m$ and $t$, level
primes $g_0,g_1$, coefficient field $L$, tuples $\mu_w$ at the places $w$ of
$L$ and base tuple $\mu_0$ (Definition~\ref{pd-period-data}). A place $w$ is
\emph{twisted} if $\mu_w\ne\mu_0$; we write $p_w$ for its residue
characteristic. Let the twisting data of $D$ (Definition~\ref{th-twisting-data})
have twisting indices $i$, with primes $\ell_i$, rings $\mathcal O_i$ with
fraction fields $K_i$ and $\mathcal O_{i,k}$ with roots of unity
$\zeta_{i,k}$, sets $S_0(i)$, ray class groups $\mathrm{RCG}_i(Q)$ with Artin
maps $\operatorname{art}\colon G_E\to\mathrm{RCG}_i(Q)$, the group $G_0$, the
idempotents $e_{i,b}$, the spaces $V_0(i,j,b)$, theta data $a$ with theta
elements $\theta(a)$ and ramification sets, theta lattices, $\psi$-twisted theta
lattices and twisted lattices, the sets $\Lambda_i$ with signs $b_\lambda$,
theta values with bounds $c_0(\lambda)$, signatures $s$ with signs $b_s$,
Bloch-Kato lengths and bounds $\nu_i$. Write $v$ for the valuation of
$\mathcal O_{i,k}$ normalized by $v(\ell_i)=1$. For a discrete valuation ring
$\mathcal O'$ over $\mathbf Z_{\ell_i}$, with normalized valuation $v'$, and
$\mathcal O'$-lattices $M,M'$ in one vector space put
\[
 [M:M']=\frac{\operatorname{len}_{\mathcal O'}\big(M/(M\cap M')\big)
 -\operatorname{len}_{\mathcal O'}\big(M'/(M\cap M')\big)}{v'(\ell_i)}.
\]
For a continuous character $\varphi\colon G_E\to\mathcal O_{i,k}^\times$ and a
place $P$ of $E$, the \emph{tame valuation} $\operatorname{tv}_P(\varphi)$ is
the infimum of $v(\varphi(\sigma)-1)$ over the elements $\sigma$ of the inertia
group at $P$ with $\varphi(\sigma)\ne1$, and $0$ if there are none. A
\emph{twist choice} at a twisted place $w$ is a twisting index $i$ with
$\ell_i=p_w$, a sign $b\in\{\pm1\}$, a finite set $Q$ of places of $E$, an
integer $k'\ge0$ and a character $\psi$ of $\mathrm{RCG}_i(Q)$ with values in
$\mathcal O_{i,k'}^\times$. Write $\psi^{(b)}=\psi$ if $b=1$ and
$\psi^{(b)}=\psi^{-1}$ if $b=-1$. The \emph{avoided places} of $D$ are the
places of $E$ over $3$, $5$, $g_0$, $g_1$, $p_{de}$,
$p'_{de}$, the auxiliary primes $T_u$ of the words, the primes
dividing the norm of the ideal $\mathfrak c$ of the host, and the prime $q_0$
of Definition~\ref{th-grid-places}: the least prime $\equiv1$ modulo $r_\ast$
times the conductor of $E$ and prime to that norm, where $r_\ast$ is
the least prime dividing none of the conductor of $E$, $p_{de}-1$,
$p'_{de}-1$ and the norm of $\mathfrak c$.
% Names (STYLE rule 13): q_0 is the auxiliary prime and r_* the helper prime of Definition
% th-grid-places (Lean q0Av D and rPrAv D), in this chapter and in Chapter 6; this chapter called
% them q^av and r^av. This chapter writes r = k+m-1 for the rank of D (Definition th-ranks), not r_0,
% which Lemma is-onto-step-numerics uses; and p_de, p'_de for the level primes (Lean pde, pde').

Among the further choices of Definition~\ref{pd-family}(4) are the
\emph{twist choices} of a family $\mathcal F$: at every member $D$ at every
scale $x>x_1$ it chooses a twist choice at every twisted place, and for every
twisting index $i$, lift $j$ and sign $b$ a set $\Sigma_D(i,j,b)$ of theta data
of $D$ whose theta elements satisfy $g\theta(a)=\psi_{U,i,b}(g)\theta(a)$ for
$g\in U$. The family satisfies the \emph{twist condition} if for all large $x$
and every member $D$ at $x$:
\begin{enumerate}
\item for every twisted place $w$ with conjugate place $\bar w$, the choices at
  $w$ and $\bar w$ have the same $i$, $Q$, $k'$ and $\psi^{(b)}$, and opposite
  signs when $\bar w\ne w$;
\item there are integers $N_0(i)\ge0$ and finite sets $S(i,j,b)$ of theta
  data such that at every twisted place $w$, with choice $(i,b,Q,k',\psi)$ and
  the character $\varphi=\psi^{(b)}\circ\operatorname{art}$ of $G_E$:
  \begin{enumerate}
  \item for every lift $1\le j\le m$ and sign $b'$, with
    $\Sigma=\Sigma_D(i,j,b')$: the vectors $|G_0|e_{i,b'}\theta(a)$,
    $a\in\Sigma$, lie in $V_0(i,j,b')$ and span it over $K_i$; and
    $S(i,j,b')\subset\Sigma$, and the vectors $|G_0|e_{i,b'}\theta(a)$,
    $a\in S(i,j,b')$, are linearly independent over $\mathcal O_i$ and span
    the $\mathcal O_i$-module $|G_0|e_{i,b'}\Theta(\Sigma)$;
  \item $k'\ge1$; every place in $Q$ is split in $E/F$, of degree $1$ and
    unramified over $\mathbf Q$, and lies outside $S_0(i)$, the avoided places
    and the ramification sets of the theta data in the sets $S(i,j,b')$,
    $1\le j\le m$; and $\varphi$ is continuous, takes values in the powers of
    $\zeta_{i,k'}$, is ramified at the places of $Q$ and nowhere else, is
    trivial on the decomposition groups at the places over $\ell_i$, and
    satisfies $\varphi(\bar\sigma)=\varphi(\sigma)^{-1}$ for the conjugation
    $\sigma\mapsto\bar\sigma$ by the complex conjugation of $E$;
  \item there are $\epsilon_{j,b'}\ge0$ with
    $\sum_{j=1}^m\sum_{b'}\epsilon_{j,b'}\le\ell_i^{-2}$ such that, for every
    lift $1\le j\le m$ and sign $b'$, with $\Sigma=\Sigma_D(i,j,b')$, the
    $\varphi^{(b')}$-twisted theta lattice
    $\Theta'=e_{i,b'}\Theta_{\varphi^{(b')},Q}(\Sigma)$ is a lattice in
    $V_0(i,j,b')\otimes\mathcal O_{i,k'}$, it is contained in the twisted
    lattice $e_{i,b'}\theta_{\varphi^{(b')},Q}(M_{i,j,b'}(k',Q,\varphi^{(b')}))$,
    and for every discrete valuation ring $\mathcal O'$ finite over
    $\mathcal O_{i,k'}$ and every $\mathcal O'$-lattice $M$ in
    $V_0(i,j,b')\otimes\mathcal O'$,
    \[
     [M:\mathcal O'\Theta']\le[M:\mathcal O'\otimes e_{i,b'}\Theta(\Sigma)]
     +\epsilon_{j,b'};
    \]
  \item there are $\epsilon_s\ge0$ for the signatures $s$ of $i$, with sum at
    most $\ell_i^{-2}$ over the good signatures, such that for every good
    signature $s$ the Bloch-Kato length of $s$ twisted by $\varphi^{(b_s)}$,
    with the set $Q$, is at most $\nu_i+\epsilon_s$;
  \item for every $\lambda\in\Lambda_i$, the theta value of $\lambda$ twisted
    by the character $(\psi^{(b)})^{(b_\lambda)}$ of $\mathrm{RCG}_i(Q)$ is
    nonzero, and its valuation $v$ lies between $-N_0(i)$ and
    $c_0(\lambda)+\ell_i^{-2}$;
  \item $t\sum_{P\in Q}\operatorname{tv}_P(\varphi)\le\ell_i^{-2}$.
  \end{enumerate}
\end{enumerate}
The \emph{strong twist condition} asks, in addition, that (b) to (f) of (2)
hold at every twisted place with $\psi$ replaced by $\psi^a$, and so $\varphi$
by $\varphi^a$, for every integer $a$ prime to $\ell_i$.
\end{definition}
\leanref{Formalism/SharpRestriction/LevelLetterCocycles/TwGood.lean}{Formalism.MainNotions.TwGood}
\leanref{Formalism/SharpRestriction/LevelLetterCocycles/TwGood.lean}{Formalism.MainNotions.TwGoodE}
\leanref{Formalism/SharpRestriction/LevelLetterCocycles/TwGood.lean}{Formalism.MainNotions.TwConjAt}
\leanref{Formalism/SharpRestriction/LevelLetterCocycles/TwGood.lean}{Formalism.MainNotions.TwGoodAt}
\leanref{Formalism/SharpRestriction/LevelLetterCocycles/TwGood.lean}{Formalism.MainNotions.AvRec}
\leanref{Formalism/SharpRestriction/LevelLetterCocycles/IntegralHeckeForms.lean}{Formalism.MainNotions.IntegralHeckeForms.GoodTw}
\leanref{Formalism/SharpRestriction/LevelLetterCocycles/TwAllF.lean}{Formalism.MainNotions.TwGoodAtAll}
\leanref{Formalism/LocalToGlobal/ClassOfRecordOfChoice.lean}{Formalism.L2.Top3.TwAllF}

\begin{definition}\label{asm-designed}
Let $\mathcal F$ be a family with exponent $\varepsilon$, and let $C_0$ and
$C_H$ be real. For a member $D=D_{x,i}$ write $\Gamma=\exp(C_\Gamma m^5fL)$,
which is the number $\Gamma_P$ of Definition~\ref{th-definite-datum}(3) by
Definition~\ref{pd-size-data}(9), let $\dim M_u$ be the dimension of the space
$M_u$ of a word $u$ and $D^{*}$ the normalizing dimension
(Definition~\ref{is-orthogonality-link}),
$|W|$ the number of words,
$N_m=k+m(t+1)$, and let $D_\Gamma(u)$ be the stabilizer denominator of $u$
(Definition~\ref{pd-period-data}(2)).
\begin{enumerate}
\item $\mathcal F$ has \emph{chosen constants} $(C_0,C_H)$ if there are
  $C_5,c_6,C_b>0$ such that for all large $x$ and every member $D$ at $x$:
  the set $T_2$ of the host satisfies
  $|T_2|\ge(n_0/L_w)\exp(-C_{T_2}\sqrt{\log n_0\log L_w})$;
  $\delta_1n/2\le|T|\le2\delta_1n$; the bad tops satisfy the bound with
  $C_b$; for $\chi\in\{\chi_{-q},\chi_{-d},\chi_{-q}\chi_{-d}\}$
  (Definition~\ref{hf-good-host}), $|g_\chi(x)|\le x^{1/2+\varepsilon}$ and
  $|\frac{L'}{L}(1,\chi)|\le C_5\log x$; for every Dirichlet character $\chi$
  of the host field and every $y\ge X'/3$,
  $|g_\chi(y)|\le y\exp(-c_6x^{\varepsilon/8})$; for every word $u$,
  $C_0^{-1}\le\dim M_u\cdot\Gamma^2/D^{*}\le C_0$; and
  $C_0\Gamma^2x^{C_H}\le|W|$.
\item $\mathcal F$ satisfies the \emph{size condition on $T_u$} if for all
  large $x$, every member and every word,
  $b^{m/2}\le T_u$ and $\log T_u\le C_{T_u}mL$.
\item $\mathcal F$ satisfies the \emph{size condition on stabilizer
  denominators} if there is $C$ with
  $\log D_\Gamma(u)\le CfN_m\log(fN_m)$ for all large $x$, every member and
  every word.
\end{enumerate}
Let $\mathcal C$ be a class of families and $C_0$ real.
\begin{enumerate}
\setcounter{enumi}{3}
\item $\mathcal C_{\mathrm{tw}}$ is the class of $\mathcal F\in\mathcal C$
  such that, for all large $x$, every prime factor of every $D_\Gamma(u)$ is at
  most $2fN_m+1$, and that satisfy the twist condition
  (Definition~\ref{asm-twist-condition}).
\item $\mathcal C_N$ is the class of $\mathcal F\in\mathcal C$ with
  $C_{YL}\le9$ such that, for all large $x$ and every member: every
  $D_\Gamma(u)$ divides Minkowski's number
  $M(2fN_m)=\prod_pp^{\sum_{j\ge0}\lfloor2fN_m/(p^j(p-1))\rfloor}$; the member
  has distinct tame data; the box condition holds at some
  $Y\in[x^{1/J},x]$; and the row condition holds at some
  $Y'\in[x^{1/J},x]$ (Definition~\ref{asm-box-row-tame}).
\item The class $\operatorname{Des}(\mathcal C,C_0)$ of \emph{designed
  members} consists of the $\mathcal F\in\mathcal C_{\mathrm{tw}}$ that have
  chosen constants $(C_0,C_H)$ for some $C_H$ and satisfy both size
  conditions.
\item \emph{Designed families exist in $\mathcal C$} with constant $C_0$ if
  there is $\varepsilon_0>0$ such that for every
  $\varepsilon\in(0,\varepsilon_0]$ and every largeness threshold $\kappa$
  there are $\mathcal F\in\mathcal C$ with exponent $\varepsilon$ and
  $C_H\ge0$ such that $\mathcal F$ is large for $(\kappa,C_H)$, has chosen
  constants $(C_0,C_H)$, and satisfies both size conditions.
\item Let $c_T$ be a function of two real variables. \emph{Design choices
  exist in $\mathcal C$} with constant $C_0$ and function $c_T$ if there is
  $\varepsilon_0>0$ such that for every $\varepsilon\in(0,\varepsilon_0]$ and
  every largeness threshold $\kappa$ there are $\mathcal F\in\mathcal C$ with
  exponent $\varepsilon$ and $C_H\ge0$, $C_5,C_b>0$ such that $\mathcal F$ is
  large for $(\kappa,C_H)$; all conditions of chosen constants other than the
  thin-scale bound and the proviso $C_0^{-1}\le\dim M_u\cdot\Gamma^2/D^{*}\le C_0$ hold
  with $C_5,C_b$; $|W|\le C_0\Gamma^2x^{C_H}+1$ for all large $x$; the
  conditions defining $\mathcal C_N$ hold; and
  $c_T(\varepsilon,C_{YL})\le C_{T_u}$ and $C_F=C_H$.
\end{enumerate}
\end{definition}
\leanref{Formalism/SharpRestriction/GlobalSchurBounds/Chosen.lean}{Formalism.MainNotions.ChosenE}
\leanref{Formalism/SharpRestriction/ParamsHeightRankConstant.lean}{Formalism.MainNotions.TuSize}
\leanref{Formalism/SharpRestriction/ParamsHeightRankConstant.lean}{Formalism.MainNotions.DGamSize}
\leanref{Formalism/SharpRestriction/SmoothDesignExistence/DesignExistsDes.lean}{Formalism.MainNotions.Des}
\leanref{Formalism/SharpRestriction/ParamsHeightRankConstant.lean}{Formalism.MainNotions.DesignExists}
\leanref{Formalism/SharpRestriction/DesignExistsOfChoices.lean}{Formalism.MainNotions.DesignChoices}
\leanref{Formalism/LocalToGlobal/ClassOfRecordOfChoice.lean}{Formalism.L2.Top3.WithChN}
\leanref{Formalism/LocalToGlobal/ClassOfRecordOfChoice.lean}{Formalism.L2.Top3.D3}
\leanref{Formalism/SharpRestriction/MuTuOfTuSize.lean}{Formalism.MainNotions.WithTw}
\leanref{Formalism/SharpRestriction/SmoothDesignExistence/DesignExistsDes.lean}{Formalism.MainNotions.WithSm}
\leanref{Formalism/SharpRestriction/SmoothDesignExistence/DGamSmooth.lean}{Formalism.MainNotions.Pk.DGamSmooth}
\leanref{Formalism/SharpRestriction/DesignExistsOfChoices.lean}{Formalism.MainNotions.DGamDvdF}
\leanref{Formalism/SharpRestriction/LevelLetterCocycles/AlignN.lean}{Formalism.MainNotions.TameDistinctF}
\leanref{Formalism/SharpRestriction/GlobalSchurBounds/MainIn.lean}{Formalism.MainNotions.BoxesF}
\leanref{Formalism/SharpRestriction/GlobalSchurBounds/MainIn.lean}{Formalism.MainNotions.QDRowsF}

\begin{definition}\label{asm-duty}
A family $\mathcal F$ with exponent $\varepsilon$ \emph{carries its duty
data} if the following hold.
\begin{enumerate}
\item Some of its constants take fixed values: $C_p=1$, $e_P=12$,
  $c_{TO}=28e_P+41=377$, $C_t=c_{TO}+24=401$, $C_{\eta F}=6$, $C_c=2$,
  $c_w=9/10$ and $C_{YL}=9$.
\item At every member at every scale $x>x_1$: the ideal $\mathfrak c_K$ of
  the host is the unit ideal; the Harish-Chandra parameters at all places have
  the same entries $v$; $g_0,g_1\le e^{9L}$; $p_{de},
  p'_{de}<X'$; $j_0=\lceil\log_5(nN_m/m)\rceil$;
  $L_w=10\lceil(C_t+4)\log x\rceil+2\lceil30/\varepsilon\rceil+22$; and
  $T_O=\lceil x^{c_{TO}}\rceil$.
\item At every member at every scale $x>x_1$, the ranks are adapted to the
  scale: $x^{2+8\varepsilon}\le n\le2x^{2+8\varepsilon}$,
  $\sqrt{x\delta_1}\le m\le2\sqrt{x\delta_1}$ and
  $x^{9/2+10\varepsilon}\le k\le2x^{9/2+10\varepsilon}$.
\item At every member at every scale $x>x_1$, the first modulus of the
  absolutely irreducible span is less than $X'$.
\item For all large $x$: every member has distinct tame data
  (Definition~\ref{asm-box-row-tame}); the weight vectors have small shares, a
  sum over candidate primes being at most $x^{-9}$; every tuple of Hecke
  characters is used; and the family satisfies the strong twist condition
  (Definition~\ref{asm-twist-condition}).
\end{enumerate}
\end{definition}
\leanref{Formalism/LocalToGlobal/ClassOfRecordOfChoice.lean}{Formalism.L2.Top3.Duty}
\leanref{Formalism/SharpRestriction/PageFix.lean}{Formalism.MainNotions.PageFix}
\leanref{Formalism/SharpRestriction/PageFix.lean}{Formalism.MainNotions.PageWin}

The size data of a member include the bound $1\le t\le x^{C_t}$ on the twist
parameter. The duty data fix $C_t=401$, and this is the exponent with which
the bound on Tate-Shafarevich groups of Chapter~\ref{chapter-tate-shafarevich}
is needed.

\section{The links}\label{section-asm-links}

\subsection{From links to ties}\label{subsection-asm-links-to-ties}

\begin{definition}\label{asm-sha-bound}
A family $\mathcal F$ with exponent $\varepsilon$ \emph{satisfies the
Tate-Shafarevich bound} if there is $x_0$ such that for every $x\ge x_0$ and
every member at $x$, with host field $E$ and twist parameter $t$, the field
$E$ satisfies the Tate-Shafarevich bound $(t,e^{x^{\varepsilon}/2})$ of
Definition~\ref{sha-bounded-field}: the groups $\text{Ш}^1_\ell(E,k)$ have
finite length for every prime $\ell$ and $0<k<t$, and
\[
 \sum_{\ell\in S}\log\ell\sum_{1\le k<t}\operatorname{len}\text{Ш}^1_\ell(E,k)
 \le e^{x^{\varepsilon}/2}
\]
for every finite set $S$ of primes.
\end{definition}
\leanref{Development/AutomorphicArithmetic/CharacterBoxes/DivinNOfHypV.lean}{UpperTheory.Top.HypV}

The bound $e^{x^{\varepsilon}/2}$ is $\sqrt b$ for the scale of the member. It
is the bound $\sqrt{\exp(x^{\varepsilon})}$ of
Definition~\ref{sha-vandiver-weakening}, at the scales $x\in[Y,Y^2]$ that the
weakening of Vandiver's conjecture provides.

\begin{definition}\label{asm-twisted-links}
Let $\mathcal F$ be a family with exponent $\varepsilon$. For a member $D$ at
$x$ with size data and a lift $1\le j\le m$, let $\eta_j(w)$, $D_0$, $B_0$ and
$N_c$ be as in Definition~\ref{is-twisted-division-bound}, and put
\[
 \eta_j=\sum_{w\in\operatorname{supp},\ \mu_w\ne\mu_0}\ell_w\,e_w\,\eta_j(w).
\]
Let $d_j$ be the dimension of the space of $j$-th lifts at the base tuple
$\mu_0$, and let $E$ and $t$ be the host field and the twist parameter of $D$.
\begin{enumerate}
\item (\emph{Tate-Shafarevich link}.) There are $C$ and $x_0$ such that for
  every $x\ge x_0$, every member at $x$ whose groups
  $\text{Ш}^1_\ell(E,k)$, $0<k<t$, have finite length, and every
  $1\le j\le m$: if the series $\sum_\ell\log\ell\sum_{1\le k<t}
  \operatorname{len}\text{Ш}^1_\ell(E,k)$ converges, then
  \[
   \eta_j\le d_j\Big(C\big(fmbL^2+(\log x)^2N_c\big)+\sum_\ell\log\ell
   \sum_{1\le k<t}\operatorname{len}\text{Ш}^1_\ell(E,k)\Big).
  \]
\item (\emph{Rank equality at small level}.) For all large $x$, every member
  and every $1\le j\le m$, $d_j=D_0$.
\item (\emph{Rank conditions}.) For all large $x$, every member and every
  word $u$, the division lattice $N_u$ is finitely generated, spans its
  space, and has positive rank.
\end{enumerate}
\end{definition}
\leanref{Formalism/SharpRestriction/GlobalTheoremInputs.lean}{Formalism.MainNotions.LinkN}
\leanref{Formalism/SharpRestriction/GlobalTheoremInputs.lean}{Formalism.MainNotions.SmallLevelRankEq}
\leanref{Formalism/SharpRestriction/GlobalTheoremInputs.lean}{Formalism.MainNotions.NuFullF}
\leanref{Formalism/SharpRestriction/AdmissibleParameterJoins/SumLogPrimeFactorsLe.lean}{Formalism.MainNotions.J.MuRank}

\begin{definition}\label{asm-global-inputs}
Let $\mathcal F$ be a family with exponent $\varepsilon$, and let
$C_{\mathrm{abs}},C_H,C_0,e_D,e_L$ be real. For a member with size data and a
word $u$, let $N^{\circ}_u$ and $s(u,w)$ be as in
Definition~\ref{pd-size-data}, and $\mathrm{sz}_w(y)=2\ell_we_wy$ as in
Definition~\ref{is-division-link}.
\begin{enumerate}
\item (\emph{Slope link}.) There are $C$ and $x_0$ with
  \[
   \Big|\mu(N_u)-\mu(N^{\circ}_u)-\sum_w\mathrm{sz}_w(s(u,w))\Big|\le CX''(n)
  \]
  for every $x\ge x_0$, every member at $x$ and every word $u$, the sum over
  the support, where $N^{\circ}_u$ carries the metrics of $N_u$.
\item The \emph{global inputs} with $(C_{\mathrm{abs}},C_H,C_0)$ are the gain
  link (Definition~\ref{pd-gain-link}), the orthogonality link with constant
  $C_0$ (Definition~\ref{is-orthogonality-link}), the irreducibility link
  (Definition~\ref{is-irreducibility-link}), the division link with constants
  $C_{\mathrm{abs}},C_H$ (Definition~\ref{is-division-link}) and the slope
  link.
\item (\emph{Global link}.) If $\mathcal F$ has chosen constants $(C_0,C_H)$
  (Definition~\ref{asm-designed}) and the global inputs with
  $(C_{\mathrm{abs}},C_H,C_0)$ hold, then there are $K,c_{HH},e_K,C_L,x_0$
  such that $\mathcal F$ is tied at every scale $x\ge x_0$ of $\mathcal F$
  with constants $(e_D,e_L,K,c_{HH},e_K,C_L)$ (Definition~\ref{asm-tied}).
\end{enumerate}
\end{definition}
\leanref{Formalism/SharpRestriction/GlobalTheoremInputs.lean}{Formalism.MainNotions.LinkMu}
\leanref{Formalism/SharpRestriction/GlobalTheoremInputs.lean}{Formalism.MainNotions.GlobalTheoremInputs}
\leanref{Formalism/SharpRestriction/GlobalTheoremInputs.lean}{Formalism.MainNotions.LinkGLOBAL}

\begin{lemma}\label{asm-division-from-sha}
Let $\mathcal F$ be a family with exponent $\varepsilon>0$
(Definition~\ref{pd-family}). Suppose that $\mathcal F$ satisfies the
Tate-Shafarevich bound (Definition~\ref{asm-sha-bound}), the Tate-Shafarevich
link and the rank equality at small level
(Definition~\ref{asm-twisted-links}). Then $\mathcal F$ satisfies the
division bound at twisted places.
\end{lemma}
\begin{proof}
Let $C_1$ and $x_3$ be the constants of the Tate-Shafarevich link. Let $x_2$
be a threshold for the Tate-Shafarevich bound and $x_4$ one for the rank
equality. Since $\log x=o(x^{\varepsilon/2})$, there is $x_6$ with
$(\log x)^2\le x^{\varepsilon}$ for $x\ge x_6$. Put $C=|C_1|+1$, and let
$x\ge\max(x_2,x_3,x_4,x_6,x_1+1)$.

Let $D$ be a member at $x$, with host field $E$ and size data $L$, and let
$1\le j\le m$. Write $\Sigma$ for the series
$\sum_\ell\log\ell\sum_{1\le k<t}\operatorname{len}\text{Ш}^1_\ell(E,k)$. The
Tate-Shafarevich bound says that its terms are finite and that its partial
sums are at most $\sqrt b$. Its terms are nonnegative, so $\Sigma$ converges
and $\Sigma\le\sqrt b\le b$, since $b=e^{x^{\varepsilon}}\ge1$. The
Tate-Shafarevich link and $d_j=D_0$ give
\[
 \eta_j\le D_0\big(C_1(fmbL^2+(\log x)^2N_c)+\Sigma\big).
\]
Here $L=x^{\varepsilon}$, so $fmbL^2=x^{\varepsilon}B_0$, and
$(\log x)^2N_c\le x^{\varepsilon}N_c$. Since $f,m,L\ge1$, we have
$b\le B_0\le x^{\varepsilon}(B_0+N_c)$. Hence
\[
 C_1(fmbL^2+(\log x)^2N_c)+\Sigma\le|C_1|x^{\varepsilon}(B_0+N_c)
 +x^{\varepsilon}(B_0+N_c)=Cx^{\varepsilon}(B_0+N_c).
\]
\end{proof}
\leanref{Development/AutomorphicArithmetic/CharacterBoxes/DivinNOfHypV.lean}{UpperTheory.Top.divinN_of_hypV}

\begin{lemma}\label{asm-slope-link}
Let $\mathcal F$ be a family with exponent $\varepsilon$. Suppose that for
all large $x$ and every member at $x$: the local data of the size data vanish
above $X''(n)$, except at places over $T_u$, $g_0$, $g_1$, $d$ and the prime
factors of $D_\Gamma(u)$; and every division lattice $N_u$ is finitely
generated of positive rank. Suppose also that there are $C_T,C_G$ with
$\log T_u\le C_TX''(n)$ and $\log D_\Gamma(u)\le C_GX''(n)$ for all large $x$,
every member and every word. Then $\mathcal F$ satisfies the slope link
(Definition~\ref{asm-global-inputs}) with constant
$2(\log4+3+C_T+C_G)$.
\end{lemma}
\begin{proof}
Fix a member at a large scale $x$, with coefficient field $L_D$, and a word
$u$; write $s_w=s(u,w)$. For the families of Definition~\ref{pd-family},
$N_u$ is the $\mathcal O_{L_D}$-span of $I_uN^{\circ}_u$, where
$I_u=\prod_w\mathfrak p_w^{-\lfloor e_ws_w\rfloor}$, the product over the
support, and the metrics of $N_u$ are invertible.

\emph{Step 1.} Multiplying a lattice by the fractional ideal $I_u$ raises its
slope by $-2\log\mathrm N(I_u)/[L_D:\mathbf Q]$, which is
$\sum_w2\ell_w\lfloor e_ws_w\rfloor$, with $\ell_w=\log\mathrm N\mathfrak
p_w/[L_D:\mathbf Q]$. Hence
\[
 \mu(N_u)-\mu(N^{\circ}_u)-\sum_w\mathrm{sz}_w(s_w)
 =\sum_w2\ell_w\big(\lfloor e_ws_w\rfloor-e_ws_w\big).
\]

\emph{Step 2.} Each term on the right has absolute value at most $2\ell_w$, and
it vanishes when $s_w=0$. So the left side is at most $2\sum_w\ell_w$ in
absolute value, the sum over the places $w$ of the support with $s_w\ne0$.

\emph{Step 3.} By the vanishing hypothesis, a place $w$ with $s_w\ne0$ lies
over a prime $p\le X''(n)$, or over one of $T_u$, $g_0$, $g_1$, $d$, or over a
prime factor of $D_\Gamma(u)$. Call this set of primes $P$.

\emph{Step 4.} For each prime $p$, the numbers $\ell_w$ of the places $w$ over
$p$ add up to at most $\log p$. Hence $\sum_w\ell_w\le\sum_{p\in P}\log p$.

\emph{Step 5.} Chebyshev's bound gives $\sum_{p\le y}\log p\le y\log4$, and
the logarithms of the distinct prime factors of $D_\Gamma(u)$ add up to at most
$\log D_\Gamma(u)$. Therefore
\[
 \Big|\mu(N_u)-\mu(N^{\circ}_u)-\sum_w\mathrm{sz}_w(s_w)\Big|\le
 2\big(X''(n)\log4+\log T_u+\log g_0+\log g_1+\log d+\log D_\Gamma(u)\big).
\]

\emph{Step 6.} The sizes of the family give $\log g_0,\log g_1,\log d\le
X''(n)$ for large $x$, and the hypotheses give $\log T_u\le C_TX''(n)$ and
$\log D_\Gamma(u)\le C_GX''(n)$. So the right side is at most
$2(\log4+3+C_T+C_G)X''(n)$.
\end{proof}
\leanref{Formalism/SharpRestriction/AdmissibleParameterJoins/SumLogPrimeFactorsLe.lean}{Formalism.MainNotions.J.linkMu}
\leanref{Formalism/SharpRestriction/AdmissibleParameterJoins/SumLogPrimeFactorsLe.lean}{Formalism.MainNotions.J.nprimeMu_at}
\leanref{Formalism/SharpRestriction/AdmissibleParameterJoins/SumLogPrimeFactorsLe.lean}{Formalism.MainNotions.J.linkMu_of_linkDIV}

The hypotheses of Lemma~\ref{asm-slope-link} hold for a designed family that
satisfies the division link and the rank conditions. The vanishing of local
data above $X''(n)$ is part of the division link. The size condition on $T_u$
gives $\log T_u\le C_{T_u}mL$, and the size condition on stabilizer
denominators gives $\log D_\Gamma(u)\le CfN_m\log(fN_m)$; both are at most a
constant times $X''(n)=xnX'$ for large $x$.

\begin{theorem}\label{asm-tied-from-links}
Let $\mathcal C,\mathcal C',\mathcal E$ be classes of families
(Definition~\ref{pd-family}) with $\mathcal C'\subset\mathcal C$ and
$\mathcal C'\subset\mathcal E$, such that every family in $\mathcal E$
satisfies the Tate-Shafarevich bound (Definition~\ref{asm-sha-bound}). Let
$C_{\mathrm{abs}},C_0,e_D,e_L$ be real. Suppose that designed families exist
in $\mathcal C'_{\mathrm{tw}}$ with constant $C_0$
(Definition~\ref{asm-designed}), and that the
following hold eventually along $\operatorname{Des}(\mathcal C,C_0)$, as
properties of $(\mathcal F,C_H)$:
\begin{enumerate}
\item the Tate-Shafarevich link (Definition~\ref{asm-twisted-links});
\item the support condition (Definition~\ref{is-division-link});
\item the irreducibility link;
\item the rank equality at small level;
\item if both size conditions hold, then the conditional division link holds
  with constants $C_{\mathrm{abs}},C_H$ (Definition~\ref{is-division-link});
\item if $\mathcal F$ has chosen constants $(C_0,C_H)$, then the rank
  conditions hold;
\item the gain link;
\item if the size condition on $T_u$ holds, then the orthogonality link holds
  with $C_0$;
\item if both size conditions hold, then the global link holds with
  $(C_{\mathrm{abs}},C_H,C_0,e_D,e_L)$.
\end{enumerate}
Then $\mathcal E$ has tied members with exponents $e_D,e_L$
(Definition~\ref{asm-tied}).
\end{theorem}
\begin{proof}
The conjunction of conditions (1) to (9) holds eventually along
$\operatorname{Des}(\mathcal C,C_0)$; let $\varepsilon_1$ be its threshold
and $\varepsilon_3$ the threshold for the existence of designed families. Put
$\varepsilon'=\min(\varepsilon_1,\varepsilon_3)$, and let
$0<\varepsilon<\varepsilon'$.

Let $\kappa$ be the largeness threshold of the conjunction at $\varepsilon$.
There are a family $\mathcal F\in\mathcal C'_{\mathrm{tw}}$ with exponent
$\varepsilon$ and $C_H\ge0$ such that $\mathcal F$ is large for
$(\kappa,C_H)$, has chosen constants $(C_0,C_H)$, and satisfies both size
conditions. Since $\mathcal C'\subset\mathcal C$, the family $\mathcal F$ lies
in $\mathcal C_{\mathrm{tw}}$ and is a designed member of $\mathcal C$, so
conditions (1) to (9) hold for
$(\mathcal F,C_H)$. Since $\mathcal C'\subset\mathcal E$, the family
$\mathcal F$ lies in $\mathcal E$ and satisfies the Tate-Shafarevich bound.

By Lemma~\ref{asm-division-from-sha}, conditions (1) and (4) and the
Tate-Shafarevich bound give the division bound at twisted places. With the
support condition (2), the size conditions and (5), the division link holds
with constants $C_{\mathrm{abs}},C_H$. Condition (6) gives the rank conditions. The
division link, the rank conditions and the size conditions give the slope link
by Lemma~\ref{asm-slope-link} and the paragraph after it. Conditions (7), (8)
and (3) give the gain, orthogonality and irreducibility links. So the global
inputs hold, and (9) shows that $\mathcal F$ is tied at all its large scales
with exponents $e_D,e_L$. Thus $\mathcal E$ has tied members, with threshold
$\varepsilon'$.
\end{proof}
\leanref{Development/AutomorphicArithmetic/CharacterBoxes/DesWin.lean}{UpperTheory.Top.topDesWin}

\subsection{Bounds at a member}\label{subsection-asm-member-bounds}

\begin{definition}\label{asm-member-bounds}
Let $\mathcal F$ be a family with exponent $\varepsilon$
(Definition~\ref{pd-family}), and let $D$ be a member at a scale $x>x_1$, with
its size data (Definition~\ref{pd-size-data}). Write $k,m,n$ for its ranks and
$f=2n$; $A$ for the group of its Galois datum and $\Psi_A$ for the split
subgroup of $E_0\times A$ (Definition~\ref{ib-host-group}), whose elements code
the $f$ places $\tau$ of Definition~\ref{th-definite-datum}(2); $T$ for its top
classes, with exponents $a_t$, and $T_2$ for the set of its host; $g_0,g_1$ for
its level primes; $a,b,b',X''=X''(n)$, $L=x^{\varepsilon}$ and the width
$\eta$ for the quantities of the size data; $L_a$, $B_0$, $N_c$, $F_o$ and
$\mathrm{sz}_w$ as in Definition~\ref{is-division-link}; $v>0$ for the common
entry of the Harish-Chandra parameters; $w_j(\tau)$ for the Harish-Chandra
weight of the lift $j$ at the place $\tau$, and, for $\rho\in\Psi_A$,
$w_j(\rho)=w_j(\tau)$ if $\rho$ codes the place $\tau$ and $w_j(\rho)=1$ if it
codes none; and $W$ for the set of words. For a word $u$ let $T_u$ be its auxiliary prime, $M_u$ the first
lattice of its chain, and $N_u$ its division lattice, of rank $d_u$ and height
$h_u$. Let $RR$ be the glued lattice, $CC$ its dual and $\delta$ the row
deficit (Definition~\ref{is-irreducibility-link}). Put
\[
 G_a=\sum_{w\in\operatorname{supp}}\mathrm{sz}_w\Big(\sum_{j=1}^m
 \mathrm{gain}^\flat_j(w)-\tau_w\Big),\qquad
 E_u=\sum_{j=1}^m\sum_{w\in\operatorname{supp}}\mathrm{sz}_w(E_{u,j}(w)),\qquad
 I=\sum_{w\in\operatorname{supp}}\ell_we_w\iota_w,
\]
\[
 s=\frac{\sum_ud_u(\mu(N_u)-h_u)}{\sum_ud_u},\qquad
 s^+=\max_u\big(\mu_{\max}(N_u)-h_u\big),
\]
and for real $C$ and a word $u$ put
$\mathrm{Lo}_C(u)=CI+C|T|m(a/k+b')\log T_u+CX''+1$. Let
$C,\eta_B,e_B,e_L,e_K,c_6,C_5$ be real numbers. The member satisfies:
\begin{enumerate}
\item the \emph{count} with $C$ if for all words $u,u'$
  \[
   \frac1f\sum_{j=1}^m\sum_{\tau}\log\frac{w_j(\tau)v}{\pi}
   +m\log(g_0g_1)-C\Big(L_a+\frac{m^2L_a}k\Big)\le\log T_u\le CmL
  \]
  and $|\log T_u-\log T_{u'}|\le C\log k$;
\item the \emph{slope bounds for $M$} with $C$ if
  $-C|T|aL_a(1+m/k)\le\mu(M_u)$ and $\mu_{\max}(M_u)\le C|T|aL_a$ for every
  word $u$;
\item the \emph{slope bounds for $N$} with $C$ if for every word $u$
  \[
   \mu_{\max}(N_u)\le\mu_{\max}(M_u)+G_a+F_o-E_u+CmB_0,\qquad
   \mu(M_u)+G_a+F_o-E_u-\mathrm{Lo}_C(u)\le\mu(N_u);
  \]
\item the \emph{lattice conditions} if $d_u>0$ and
  $\mu(N_u)\le\mu_{\max}(N_u)$ for every word $u$, and
  $\mu_{\min}(CC)=-\mu_{\max}(RR)$;
\item the \emph{frame bound} with $(\eta_B,e_B)$ if there are
  $0\le\eta'\le\eta_B$ and $0\le e\le e_B$ such that for every $\Delta>0$ with
  $s-\mu_{\min}(CC)\le\Delta$,
  \[
   s-\log|W|-\frac{4(s^+-s)\log\big(1+\Delta/(s^+-s)\big)
   +(\eta'+4\sqrt e)\Delta}{1-\delta}\le\mu(CC),
  \]
  the term $(s^+-s)\log(1+\Delta/(s^+-s))$ being read as $0$ when $s^+=s$;
  this is the frame bound of Definition~\ref{hs-frame-bound};
\item the \emph{loss bounds} with $(C,e_L)$ if
  $I\le C|T|mbx^{e_L\varepsilon}$ and, for the arithmetic data of the host
  (Lemma~\ref{ib-host-arithmetic-data}) and every integer $N$ with
  $a_t/w_j(\rho)\le N$ for all $t\in T$, $j$ and $\rho\in\Psi_A$,
  \[
   \sum_{t\in T}\sum_{j=1}^m\mathrm{ASK}_N\big(a_t,t,w_j\big)-G_a
   \le C|T|mbx^{e_L\varepsilon},
  \]
  with the inert-prime sums of Definition~\ref{ib-inert-sums};
\item the \emph{node bounds} if for every lift $j$ and every
  $\rho\in\Psi_A$: $w_j(\rho)\ge1$; $b\le w_j(\rho)\le b+\eta+1$ if $\rho$ is
  thick at the base point of $D$ (Definition~\ref{ib-size-windows}); and
  $w_j(\rho)\le2b'$ otherwise;
\item the \emph{character conditions} with $(c_6,C_5)$ if there is a function
  $\ell$ on the characters of $E_0\times A$, compatible with the host as in
  Definition~\ref{asm-tied}(4), with
  $\operatorname{Re}\ell(\chi_C)=\operatorname{Re}\frac{L'}{L}(1,\chi_d)$ for
  the character $\chi_C$ of Definition~\ref{ib-host-group}, such that: if every
  character $\chi$ of the host field satisfies
  $|g_\chi(y)|\le y\exp(-c_6x^{\varepsilon/8})$ for $y\ge X'/3$, then the
  thin-scale condition (W4) of Definition~\ref{ib-size-windows} holds with $X'$
  and $\theta=\exp(-c_6x^{\varepsilon/8})$; and if
  $|g_\chi(x)|\le x^{1/2+\varepsilon}$ and
  $|\frac{L'}{L}(1,\chi)|\le C_5\log x$ for
  $\chi\in\{\chi_{-q},\chi_{-d},\chi_{-q}\chi_{-d}\}$, then
  $|g_\psi(x)|\le x^{1/2+\varepsilon}$ and $|\ell(\psi)|\le C_5\log x$ for the
  three nontrivial characters $\psi$ of $E_0$;
\item the \emph{window conditions} with $(e_K,c_6)$ if $v>0$, $0\le L_a\le
  2x^{\varepsilon}$, $x^{1/2-\varepsilon}\le m\le x^{1/2+\varepsilon/2}$,
  $xm\le k$, $k\le x^{e_K}$, $k\le a$, $k\le n^2x^{1/2-\varepsilon}$,
  $kb'\le a$, $X''\le a$, $\log|W|\le(k-3)\log2$, $x^{-\varepsilon}n\le|T_2|$,
  $2f(x+1)\exp(-c_6x^{\varepsilon/8})\le1/x$, $6f(x+1)(\eta+1)\le b/x$, and
  $a_t\le2a$ and $|a_t/b-x|\le1$ for every $t\in T$.
\end{enumerate}
The family satisfies (1), (2) or (3) \emph{at all large scales} if there are
$C\ge0$ and $x_0$ such that every member at every scale $x\ge x_0$ satisfies
it with $C$, and (4) or (7) if there is $x_0$ such that every member at every
$x\ge x_0$ satisfies it. It satisfies the \emph{frame bound with $C_H$} if
there is $x_0$ such that for every $x\ge x_0$ there are $\eta_B,e_B\ge0$ with
$e_B\le1/4$ and $(\eta_B+4\sqrt{e_B})x^{C_H}\le6$ such that every member at $x$
satisfies (5) with $(\eta_B,e_B)$. It satisfies the \emph{loss bounds with
$e_L$} if there are $C\ge0$ and $x_0$ such that for every scale $x\ge x_0$
there are a member $i_0$ at $x$ and integers $a_t,b_t$, indexed by the common
image of the top classes in $\mathbf Z$, such that the $a_t$ lie in the window
of Definition~\ref{pd-size-data}(2) for the size data of $i_0$,
$b'\le b_t\le b_{\mathrm{top}}(x)$ for the upper bound $b_{\mathrm{top}}$ of
the family, some member at $x$ has the exponents $(a_t,b_t)$, and every member
at $x$ with these exponents satisfies (6) with $(C,e_L)$. It satisfies the
\emph{character conditions} if for all $c_6,C_5$ there is $x_0$ such that every
member at every $x\ge x_0$ satisfies (8) with $(c_6,C_5)$, and the \emph{window
conditions} if for every $c_6>0$ there are $e_K\ge0$ and $x_0$ such that every
member at every $x\ge x_0$ satisfies (9) with $(e_K,c_6)$.
\end{definition}
\leanref{Formalism/SharpRestriction/GlobalSchurBounds/QuadraticLogDerivCongr.lean}{Formalism.MainNotions.Lnk.Count_at}
\leanref{Formalism/SharpRestriction/GlobalSchurBounds/QuadraticLogDerivCongr.lean}{Formalism.MainNotions.Lnk.SlopesM_at}
\leanref{Formalism/SharpRestriction/GlobalSchurBounds/QuadraticLogDerivCongr.lean}{Formalism.MainNotions.Lnk.SlopesN_at}
\leanref{Formalism/SharpRestriction/GlobalSchurBounds/QuadraticLogDerivCongr.lean}{Formalism.MainNotions.Lnk.Lat_at}
\leanref{Formalism/SharpRestriction/GlobalSchurBounds/QuadraticLogDerivCongr.lean}{Formalism.MainNotions.Lnk.Frame_at}
\leanref{Formalism/SharpRestriction/GlobalSchurBounds/QuadraticLogDerivCongr.lean}{Formalism.MainNotions.Lnk.Loss_at}
\leanref{Formalism/SharpRestriction/GlobalSchurBounds/QuadraticLogDerivCongr.lean}{Formalism.MainNotions.Lnk.Nodes_at}
\leanref{Formalism/SharpRestriction/GlobalSchurBounds/QuadraticLogDerivCongr.lean}{Formalism.MainNotions.Lnk.Char_at}
\leanref{Formalism/SharpRestriction/GlobalSchurBounds/QuadraticLogDerivCongr.lean}{Formalism.MainNotions.Lnk.Win_at}
\leanref{Formalism/SharpRestriction/GlobalSchurBounds/QuadraticLogDerivCongr.lean}{Formalism.MainNotions.Lnk.CountF}
\leanref{Formalism/SharpRestriction/GlobalSchurBounds/QuadraticLogDerivCongr.lean}{Formalism.MainNotions.Lnk.FrameF}
\leanref{Formalism/SharpRestriction/GlobalSchurBounds/LossesSmallShares.lean}{Formalism.MainNotions.Lnk.LossFS}
\leanref{Formalism/SharpRestriction/GlobalSchurBounds/QuadraticLogDerivCongr.lean}{Formalism.MainNotions.Lnk.WinF}

\begin{lemma}\label{asm-tie-from-bounds}
Let $\mathcal F$ be a family with exponent $\varepsilon$, and let
$C_{\mathrm{abs}},C_H,e_L\ge0$, $C_5$ and $c_6>0$ be real. Suppose that
$\mathcal F$ satisfies the count, the slope bounds for $M$ and for $N$, the
lattice conditions, the node bounds, the frame bound with $C_H$, the loss
bounds with $e_L$, the character conditions and the window conditions
(Definition~\ref{asm-member-bounds}); that for all large $x$ every member
satisfies $\delta_1n/2\le|T|\le2\delta_1n$, the bounds
$|g_\chi(x)|\le x^{1/2+\varepsilon}$ and $|\frac{L'}{L}(1,\chi)|\le C_5\log x$
for $\chi\in\{\chi_{-q},\chi_{-d},\chi_{-q}\chi_{-d}\}$, and the thin-scale
bound with $c_6$ of Definition~\ref{asm-designed}(1); and that $\mathcal F$
satisfies the irreducibility link (Definition~\ref{is-irreducibility-link}) and
the division link with constants $C_{\mathrm{abs}},C_H$
(Definition~\ref{is-division-link}). Let $C_d$, $C_c$, $C_s$, $C_1$ and $C_l$
be the constants of the division link, the count, the slope bounds for $M$
and for $N$, and the loss bounds. Then there are $e_K$ and $x_0$ such that
$\mathcal F$ is tied at every scale $x\ge x_0$ with constants
$(C_{\mathrm{abs}},e_L,K,C_H,e_K,C_5)$ (Definition~\ref{asm-tied}), where
$K=\max(C_d,0)+C_c+C_s+C_1+C_l+6$.
\end{lemma}
\begin{proof}
Let $e_K$ be the constant of the window conditions for $c_6$. Let $x_0$ exceed
$x_1+1$ and the thresholds of all the hypotheses, and let $x\ge x_0$ be a
scale. Write $A$ for the common group of the Galois data at $x$; it has order
$n$.

\emph{Members.} The loss bounds give a member $i_0$ and exponents
$(a_t)$, $(b_t)$ that some member at $x$ realizes. The members of a family at
one scale cover the base points: for every member at $x$ and every
$u_0\in A$ there is a member at $x$ with base point $u_0$ and the same
exponents on its top classes. Choose for each $u_0\in A$ a member
$\mathrm{pt}(u_0)$ with base point $u_0$ and exponents $(a_t),(b_t)$. All these
members have the Galois datum, ranks, primes $q,d$, set $T_2$ and parameters
$L_a$, $B_0$, $N_c$ of $i_0$.

\emph{Characters.} The character conditions give a function $\ell$, compatible
with the host, with $\operatorname{Re}\ell(\chi_C)=\operatorname{Re}
\frac{L'}{L}(1,\chi_d)$. The thin-scale bound with $c_6$ gives (W4) with
$\theta=\exp(-c_6x^{\varepsilon/8})$, and the bounds for
$\chi_{-q},\chi_{-d},\chi_{-q}\chi_{-d}$ give $|g_\psi(x)|\le
x^{1/2+\varepsilon}$ and $|\ell(\psi)|\le C_5\log x$ for the three nontrivial
characters $\psi$ of $E_0$.

\emph{Comparison data.} Take comparison data on $(A,W,\mathcal T,m,f)$ with
$\mathcal T$ the top classes, $W(u_0)$ the words of $\mathrm{pt}(u_0)$ and the
identifications of Definition~\ref{asm-tied}(1)--(4), with
$N=\lceil2a\rceil$, with constants
\[
 c_{\mathrm{Div}}=c_E=c_H=\max(C_d,0),\quad e_{\mathrm{Div}}=C_{\mathrm{abs}},\quad
 c_{\mathrm{Ct}}=C_c,\quad c_{\mathrm{Sl}}=C_s,\quad c'_{\mathrm{Sl}}=C_1,\quad
 c_{\mathrm{Loss}}=C_l,\quad e_{\mathrm{Loss}}=e_L,\quad c_{HH}=C_H,
\]
with the frame parameters $\eta(u_0)$ and $e(u_0)$ given by the frame bound at
$\mathrm{pt}(u_0)$, $\delta(u_0)$ the row deficit of $\mathrm{pt}(u_0)$, and
the bounds $\bar\eta=\eta_B$, $\bar e=e_B$ and
$\bar\delta=\exp(-x^{\varepsilon/(7J)})$.

\emph{The comparison hypotheses.} We check
Definition~\ref{ib-comparison-hypotheses} at each $u_0\in A$ and the member
$\mathrm{pt}(u_0)$.
\begin{itemize}
\item (H1): $x,b,a,b',k$ are positive since $x>1$; $v>0$ by the window
  conditions; $f=2n>0$; $a_t\ge a>0$; $w_j(\rho)\ge1$ by the node bounds;
  $r_u=d_u>0$ and $\mu_N(u)\le\mu_N^+(u)$ by the lattice conditions; $P_u=T_u$
  is a prime; $c'_{\mathrm{Sl}}=C_1\ge0$; $\bar\delta<1$; and
  $\bar\eta\ge0$, $0\le\bar e\le1/4$ by the frame bound.
\item (H2), (H3) and (H4) are conditions (7), (8) and (9) of the division link
  at $\mathrm{pt}(u_0)$, whose constant $C_d$ is at most $\max(C_d,0)$; the
  terms it multiplies, $|T|aL_a$, $mB_0$ and $m(B_0+N_c)x^{C_{\mathrm{abs}}
  \varepsilon}$, are nonnegative.
\item (H5) is the count, with $\lambda=\log(g_0g_1)$.
\item (H6) is the slope bounds for $M$.
\item (H7) is the slope bounds for $N$; with $\kappa(u_0)=I$, the quantity
  $\mathrm{Lo}_{C_1}(u)$ is $\mathrm{Lo}(u_0,u)$ of
  Definition~\ref{ib-comparison-quantities}.
\item (H8) is the loss bounds, applied with $N=\lceil2a\rceil$: for $t\in T$
  and $\rho\in\Psi_A$ we have $a_t/w_j(\rho)\le a_t\le2a\le N$, by the node
  bounds and the window conditions.
\item (H9) is the frame bound.
\item (H10) follows from condition (2) of the division link,
  $s+\mu_{\max}(RR)\le x^{C_H}b$, and the equality
  $\mu_{\min}(CC)=-\mu_{\max}(RR)$ of the lattice conditions.
\item (H11): $0\le\delta(u_0)\le\bar\delta$ by the irreducibility link, and the
  frame bound gives $0\le\eta(u_0)\le\eta_B$ and $0\le e(u_0)\le e_B$.
\end{itemize}

\emph{Error windows with constants $(e_K,K)$}
(Definition~\ref{ib-error-windows}). Conditions (1) to (4) are the window
conditions with $x\ge1$, $0<\varepsilon\le1$, $b\ge1$, $a=xb$ and $|T|\ge1$,
which follows from $|T|\ge\delta_1n/2>0$. Condition (5) holds since
$x^{\varepsilon/(7J)}\ge1$ gives $\bar\delta\le e^{-1}<1/2$. Condition (6) holds because each of $C_d$,
$C_c$, $C_s$, $C_1$, $C_l$ and $\max(C_d,0)$ lies in $[0,K]$, and
$C_{\mathrm{abs}},e_L,C_H\ge0$. For condition (7),
\[
 (\eta_B+4\sqrt{e_B})x^{C_H}b\le6b\le6|T|a\le Kx^{\varepsilon}|T|a,
\]
since $|T|\ge1$, $a=xb\ge b$, $K\ge6$ and $x^{\varepsilon}\ge1$.

\emph{Condition (4) of Lemma~\ref{ib-gap-bound}.} By definition
$B_0=fmbx^{\varepsilon}=2|A|mbx^{\varepsilon}$ and
$N_c=|T_2|kbx^{2\varepsilon}/|A|^2$, so both bounds hold with equality.

\emph{Size windows} (Definition~\ref{ib-size-windows}), with $a_E=a$. (W1):
$\eta\ge0$, $X'>0$, $b'>0$, $\theta>0$ and $b'X'=a=xb$. (W2) and the bound
$a_t\le2a$ are window conditions. (W3) is the node bounds, with
$a_t/w_j(\rho)\le N$ as above. (W4) was obtained under \emph{Characters}. (W5)
is the window conditions, since $|\Psi_A|=2|A|=f$. (W6): $a\le2xb$, the bound
$\delta_1n/2\le|T|$ with $\delta_1=x^{-\varepsilon/8}$, and the window
conditions $x^{-\varepsilon}n\le|T_2|$, $x^{1/2-\varepsilon}\le m$ and
$k\le n^2x^{1/2-\varepsilon}$.

These are the conditions of Definition~\ref{asm-tied} at $x$, with
$e_{\mathrm{Div}}=C_{\mathrm{abs}}$, $e_{\mathrm{Loss}}=e_L$,
$c_{HH}(\mathcal D)=C_H$, the constants $(e_K,K)$ in
Lemma~\ref{ib-gap-bound}, and $C_L=C_5$.
\end{proof}
\leanref{Formalism/SharpRestriction/GlobalSchurBounds/LossesSmallShares.lean}{Formalism.MainNotions.Lnk.linkCore}

\begin{definition}\label{asm-facts-volume-tie}
Let $D$ be a definite-side datum (Definition~\ref{th-definite-datum}) with host
fields $F\subset E$, $f=[F:\mathbf Q]$, and $v=v_E$ the half root discriminant
of $E$ (Theorem~\ref{cv-siegel-mass-ratio}). Consider Hermitian spaces $W$ over
$E$ with correct signature (Definition~\ref{th-spaces}), with standard lattice
$L_W$, the unimodular lattice of $W$. The datum attaches to every such $W$ a
number $\operatorname{vol}_D(W)>0$, its \emph{volume}, and to every pair
$(P,Q)$ a Haar measure $c_{P,Q}$ on $U(P,Q)$, of total mass $1$ when $Q=0$. Let
$W$ have signature $(P,Q)$ at every complex embedding of $E$, let
$K_W\subset U(W)(\mathbf A_{F,f})$ be the stabilizer of the closure of $L_W$, and
identify $U(W)(F_\rho)$ with $U(P,Q)$ at each real place $\rho$ of $F$ through
the archimedean embeddings of $D$. A \emph{volume measure} for $W$ is a Haar
measure $\mu$ on $U(W)(\mathbf A_F)$ with
\[
 \mu\{h:\ h_f\in K_W,\ h_\rho\in B_\rho\text{ for all }\rho\}=\prod_\rho c_{P,Q}(B_\rho)
\]
for all measurable $B_\rho\subset U(P,Q)$. Put
$\Lambda_E(M)$ as in Theorem~\ref{cv-siegel-mass-ratio}.
\begin{enumerate}
\item The \emph{volume tie} holds for $D$ if for every such $W$ there is a
  volume measure, for every volume measure $\mu$ there is a measurable
  fundamental domain for $U(W)(F)$ in $U(W)(\mathbf A_F)$, and
  $\mu(\mathcal F)=\operatorname{vol}_D(W)$ for every volume measure $\mu$ and
  every such fundamental domain $\mathcal F$.
\item The \emph{mass ratio under volume} holds for $D$ if the following holds
  whenever the volume tie holds for $D$. For every $W$ as above, of rank
  $M\ge3$ and signature $(P,Q)$ at every complex embedding, and every sign
  $s=\pm1$ with $P\ge1$ if $s=1$ and $Q\ge1$ if $s=-1$, there is $x\in W$ with
  $h(x,x)=s$ such that the orthogonal complement $W'$ of $x$, with the
  restricted form, is a Hermitian space with correct signature $(P-1,Q)$ if
  $s=1$ and $(P,Q-1)$ if $s=-1$ at every complex embedding, and
  \[
   \frac{\operatorname{vol}_D(W')}{\operatorname{vol}_D(W)}
   =\Big(\frac{(2\pi)^M}{(M-1)!}\Big)^f(2v)^{f(1/2-M)}\Lambda_E(M)^{-1}.
  \]
\item The \emph{mass formula for genera of rank $2$} holds for $D$ if
  $\operatorname{mass}(L_W)=\operatorname{mass}(I_2)$
  (Definition~\ref{cv-mass}) for every $W$ as above of rank $2$ that is
  positive definite at every complex embedding.
\end{enumerate}
\end{definition}
\leanref{Formalism/SharpRestriction/GlobalSchurBounds/StandardLattice.lean}{Formalism.MainNotions.StandardLattice.VolTie}
\leanref{Formalism/SharpRestriction/GlobalSchurBounds/StandardLattice.lean}{Formalism.MainNotions.StandardLattice.IsVolMeas}
\leanref{Formalism/SharpRestriction/GlobalSchurBounds/MassRatioUnfoldedOfTheoremTwo.lean}{Formalism.MainNotions.RT.MassRatioUnderVolume}
\leanref{Formalism/SharpRestriction/GlobalSchurBounds/ThickPlacesCorollaryOfAnyTranslate.lean}{Formalism.MainNotions.StandardLattice.VolumeRatio₂}
\leanref{Formalism/LocalToGlobal/GlobalTheory/GL5/DiscreteSeriesOfFormalDegree.lean}{Formalism.L2.GlobalTheoryTheory.GenusMassVolumes.GenusMassAt}

\begin{lemma}\label{asm-facts-base-volume}
Let $D$ be a definite-side datum (Definition~\ref{th-definite-datum}) with
ranks $k,m$ and $r=k+m-1$, host fields $F\subset E$ and base space $Y_0$
(Definition~\ref{th-spaces}). Suppose that the volume tie, the mass ratio under
volume and the mass formula for genera of rank $2$ hold for $D$
(Definition~\ref{asm-facts-volume-tie}). Then
$\operatorname{vol}_D(W)=\operatorname{mass}(I_{P+Q})$ for every Hermitian
space $W$ over $E$ with correct signature $(P,Q)$ at every complex embedding
and $P\ge2$. In particular $\operatorname{vol}_D(Y_0)=\operatorname{mass}(I_r)$.
\end{lemma}
\leanref{Formalism/LocalToGlobal/GlobalTheory/GL5/DiscreteSeriesOfFormalDegree.lean}{Formalism.L2.GlobalTheoryTheory.GenusMassVolumes.volume_eq_mass_rank_two}
\leanref{Formalism/LocalToGlobal/GlobalTheory/GL5/DiscreteSeriesOfFormalDegree.lean}{Formalism.L2.GlobalTheoryTheory.GenusMassVolumes.vol_indef}
\leanref{Formalism/LocalToGlobal/GlobalTheory/GL5/DiscreteSeriesOfFormalDegree.lean}{Formalism.L2.GlobalTheoryTheory.GenusMassVolumes.volume_mass_ge_two}
\leanref{Formalism/LocalToGlobal/GlobalTheory/GL5/DiscreteSeriesOfFormalDegree.lean}{Formalism.L2.GlobalTheoryTheory.GenusMassVolumes.volume_mass_succ}
\leanref{Formalism/LocalToGlobal/GlobalTheory/GL5/PacketStatementOfIndependence.lean}{Formalism.L2.GlobalTheoryTheory.GenusMassVolumes.volume_mass_of_genus_rank}
\leanref{Formalism/LocalToGlobal/GlobalTheory/GL5/DiscreteSeriesOfFormalDegree.lean}{Formalism.L2.GlobalTheoryTheory.GenusMassVolumes.volume_mass_of_class_data}
\leanref{Formalism/SharpRestriction/GlobalSchurBounds/ThickPlacesCorollaryOfAnyTranslate.lean}{Formalism.MainNotions.StandardLattice.volumeMass_of_unfolding}
\leanref{Formalism/GL3MeanSquares/HookFormulaTorsion/FiniteEulerProductLeZeta.lean}{Formalism.HermitianLattices.GenusMassVolumes.global_arithmetic_identity}
\leanref{Formalism/LocalToGlobal/GlobalTheory/GL5/FamilyMassVolumeIdentityOfClassDecomp.lean}{Formalism.L2.GlobalTheoryTheory.GenusMassVolumes.Unfold.measure_fd_toReal}
\leanref{Formalism/LocalToGlobal/GlobalTheory/GL5/FamilyMassVolumeIdentityOfClassDecomp.lean}{Formalism.L2.GlobalTheoryTheory.GenusMassVolumes.measure_KSub}
\leanref{Formalism/LocalToGlobal/GlobalTheory/GL5/PClassItems.lean}{Formalism.L2.GlobalTheoryTheory.GenusMassVolumes.class_number_finite_rank}
\leanref{Formalism/LocalToGlobal/GlobalTheory/GL5/PacketStatementOfIndependence.lean}{Formalism.L2.GlobalTheoryTheory.GenusMassVolumes.classMass_of_D7}
\begin{proof}
Write $R_E(M)$ for the right side of Theorem~\ref{cv-siegel-mass-ratio}, so
that $\operatorname{mass}(I_M)=R_E(M)\operatorname{mass}(I_{M-1})$ for $M\ge3$;
the theorem applies to the host field $E$.

\emph{Step 1: the ratio of volumes.} For $M\ge1$,
\[
 \Big(\frac{(2\pi)^M}{(M-1)!}\Big)^f\Big((M-1)!\Big(\frac v\pi\Big)^M\Big)^f=(2v)^{fM},
\]
so the reciprocal of the ratio in Definition~\ref{asm-facts-volume-tie}(2) is
$((M-1)!(v/\pi)^M)^f(2v)^{-f/2}\Lambda_E(M)=R_E(M)$. Since the volume tie holds,
the mass ratio under volume therefore gives, for $W$, $s$ and $W'$ as there,
\[
 \operatorname{vol}_D(W)=R_E(M)\operatorname{vol}_D(W').
\]

\emph{Step 2: unfolding.} Let $W$ be positive definite of rank $n$ at
every complex embedding. Then $U(P,Q)=U(n)$ is compact, and for a volume
measure $\mu$ the set $K=\{h:h_f\in K_W\}$ has $\mu(K)=c_{n,0}(U(n))^f=1$. By
Borel's finiteness of class numbers \cite{PR94}, $U(W)(\mathbf A_F)$ is
the disjoint union of finitely many double cosets $U(W)(F)g_iK$, and each group
$\Gamma_i=U(W)(F)\cap g_iKg_i^{-1}$ is finite, being discrete in a compact
group. Unfolding a fundamental domain $\mathcal F$ over these double cosets
gives
\[
 \mu(\mathcal F)=\sum_i\frac{\mu(K)}{|\Gamma_i|}=\sum_i\frac1{|\Gamma_i|}.
\]
The double coset of $g_i$ corresponds to the class of the lattice
$W\cap g_i\widehat{L_W}$ in the genus of $L_W$, whose stabilizer is
$\Gamma_i$, so the sum is $\operatorname{mass}(L_W)$
(Definition~\ref{cv-mass}). By the volume tie,
$\operatorname{vol}_D(W)=\operatorname{mass}(L_W)$. For $n=2$ the mass formula
for genera of rank $2$ gives $\operatorname{vol}_D(W)=\operatorname{mass}(I_2)$.

\emph{Step 3: definite spaces.} We show
$\operatorname{vol}_D(W)=\operatorname{mass}(I_n)$ for every positive definite
$W$ of rank $n\ge2$, by induction on $n$. Step 2 is the case $n=2$. Let $W$ have
rank $n+1\ge3$. Take $s=1$ in Step 1: the space $W'$ is positive definite of
rank $n$, so
$\operatorname{vol}_D(W)=R_E(n+1)\operatorname{mass}(I_n)=\operatorname{mass}(I_{n+1})$.

\emph{Step 4: indefinite spaces.} We show
$\operatorname{vol}_D(W)=\operatorname{mass}(I_{P+Q})$ for $W$ of signature
$(P,Q)$ with $P\ge2$, by induction on $Q$. Step 3 is the case $Q=0$. If
$Q\ge1$, take $s=-1$ in Step 1. The space $W'$ has signature $(P,Q-1)$, so by
induction
$\operatorname{vol}_D(W)=R_E(P+Q)\operatorname{mass}(I_{P+Q-1})=\operatorname{mass}(I_{P+Q})$,
since $P+Q\ge3$.

\emph{Step 5.} The space $Y_0$ has rank $r$ and signature $(k+r_S,r_S+1)$ at
every complex embedding, with $k+r_S\ge2$ (Definition~\ref{th-spaces}(5)).
Step 4 gives $\operatorname{vol}_D(Y_0)=\operatorname{mass}(I_r)$.
\end{proof}

\begin{lemma}\label{asm-bounds-glued-rank}
Let $\mathcal F$ be a family with exponent $\varepsilon$
(Definition~\ref{pd-family}), and let $C_0>0$ and $C_{\mathrm{abs}},C_H$ be
real. Suppose that $\mathcal F$ satisfies the orthogonality link with constant
$C_0$ (Definition~\ref{is-orthogonality-link}), the irreducibility link
(Definition~\ref{is-irreducibility-link}) and the division link with constants
$C_{\mathrm{abs}},C_H$ (Definition~\ref{is-division-link}), and that for all
large $x$ every member at $x$ satisfies $C_0^{-1}\le\dim M_u\cdot\Gamma^2/D^{*}\le
C_0$ for every word $u$ (Definition~\ref{asm-designed}). Then for all large
$x$ the glued lattice $RR$ of every member at $x$ has positive rank.
\end{lemma}
\leanref{Formalism/SharpRestriction/GlobalSchurBounds/TopDensityNGeSixteen.lean}{Formalism.MainNotions.Lnk.rrRankF_of}
\leanref{Formalism/SharpRestriction/GlobalSchurBounds/TopDensityNGeSixteen.lean}{Formalism.MainNotions.Lnk.dall_pos}
\begin{proof}
Let $x$ exceed the thresholds of the hypotheses, let $D$ be a member at $x$ and
$u_0$ its base word. The orthogonality link gives $e_0\ge0$ with
$\sqrt{e_0}\le e^{-x^{\varepsilon/8}}<1$, so $e_0<1$, and, since the bounds on
$\dim M_u\cdot\Gamma^2/D^{*}$ hold, it gives
$\|K_{u_0u_0}-1\|_{HS}^2\le e_0\dim M_{u_0}$. These bounds also give $\dim
M_{u_0}\ge D^{*}/(C_0\Gamma^2)>0$.

Suppose that $D_{\mathrm{all}}=0$. Then $\mathcal H=0$, so $P_{u_0}=0$ and
$K_{u_0u_0}=0$, and $\|K_{u_0u_0}-1\|_{HS}^2=\dim\mathcal C_{u_0}=\dim M_{u_0}$.
This exceeds $e_0\dim M_{u_0}$, a contradiction. So $D_{\mathrm{all}}>0$.

The irreducibility link gives $\delta\le\exp(-x^{\varepsilon_z/7})<1$, so
$D_{\mathrm{rows}}=(1-\delta)D_{\mathrm{all}}>0$. By condition (5) of the
division link, $RR$ is finitely generated and spans its $L$-vector space, whose
dimension is the dimension $D_{\mathrm{rows}}$ of the space $\mathcal E$ of top
forms. So $\operatorname{rk}RR=D_{\mathrm{rows}}>0$.
\end{proof}

\begin{lemma}\label{asm-bounds-dual-slopes}
Let $K$ be a totally complex number field with $r$ infinite places, let $V$ be a
Hermitian space over $K$ with dual $V^\vee$, and let $A\subset V$ be a nonzero
Hermitian lattice that spans $V$, with dual $A^\vee$
(Definition~\ref{hs-hermitian-lattice}). For a saturated sublattice $N\subset A$
let $N^\perp\subset V^\vee$ be the space of linear forms that vanish on $N$.
\begin{enumerate}
\item $A^\vee$ is a nonzero Hermitian lattice that spans $V^\vee$. For every
  saturated sublattice $N\subset A$, the sublattice $A^\vee\cap N^\perp$ of
  $A^\vee$ has rank $\operatorname{rk}A-\operatorname{rk}N$ and Arakelov degree
  $\widehat{\deg}(A^\vee\cap N^\perp)=\widehat{\deg}(N)-\widehat{\deg}(A)$
  (Definition~\ref{hs-slopes}). In particular
  $\widehat{\deg}(A^\vee)=-\widehat{\deg}(A)$.
\item Every sublattice $R\subset A^\vee$ lies in $A^\vee\cap N^\perp$ for a
  saturated sublattice $N\subset A$ with
  $\operatorname{rk}R=\operatorname{rk}A-\operatorname{rk}N$, and then
  $\widehat{\deg}(R)\le\widehat{\deg}(A^\vee\cap N^\perp)$.
\item $\mu_{\max}(A^\vee)=-\mu_{\min}(A)$ and $\mu_{\min}(A^\vee)=-\mu_{\max}(A)$.
\end{enumerate}
\end{lemma}
\leanref{Formalism/HarderNarasimhan/HarderNarasimhanProofs/DualPl.lean}{Formalism.HN.dual_polygon_identity}
\leanref{Formalism/HarderNarasimhan/HarderNarasimhanProofs/DualPl.lean}{Formalism.HN.DualPolygon.deg_R}
\leanref{Formalism/HarderNarasimhan/HarderNarasimhanProofs/DualPl.lean}{Formalism.HN.DualPolygon.deg_dual}
\leanref{Formalism/HarderNarasimhan/HarderNarasimhanProofs/DualPl.lean}{Formalism.HN.DualPolygon.sat_of_ann}
\leanref{Formalism/HarderNarasimhan/HarderNarasimhanProofs/DualPl.lean}{Formalism.HN.DualPolygon.corr_R}
\leanref{Formalism/HarderNarasimhan/HarderNarasimhanProofs/DualPl.lean}{Formalism.HN.DualPolygon.mu_dual}
\leanref{Formalism/LocalRingCardinality/RankBoundMonotoneInDimension.lean}{Formalism.BasicNotions.bridge_dual_mu_holds}
\begin{proof}
(1) A basis of $V$ inside $A$ shows that $A^\vee$ lies between two lattices of
full rank in $V^\vee$, so it is finitely generated, spans $V^\vee$ and is
nonzero. The rank of $A^\vee\cap N^\perp$ is $\dim N^\perp=\dim V-\operatorname{rk}N$.
The formula for its degree is a computation of covolumes: in coordinates adapted
to $N$, the covolume of $A^\vee\cap N^\perp$ for the dual norms is the covolume
of $N$ divided by that of $A$, up to the factors $2^{-r}|d_K|^{1/2}$ per unit of
rank that the normalization of $\widehat{\deg}$ removes
(Remark~\ref{hs-normalization}). The case $N=0$ gives
$\widehat{\deg}(A^\vee)=-\widehat{\deg}(A)$.

(2) Let $W\subset V$ be the subspace on which every element of $R$ vanishes, and
$N=A\cap W$. Then $N$ is saturated, $KN=W$ because $A$ spans $V$, and
$R\subset A^\vee\cap N^\perp$, with $\operatorname{rk}R=\dim V-\dim W$. So $R$
has finite index in $A^\vee\cap N^\perp$, its covolume is at least that of
$A^\vee\cap N^\perp$, and its degree is at most that of $A^\vee\cap N^\perp$.

(3) Let $\theta$ be real. By Definition~\ref{hs-slopes},
$\mu_{\max}(A^\vee)\le\theta$ exactly when
$\widehat{\deg}(R)\le\theta r\operatorname{rk}R$ for every nonzero sublattice
$R\subset A^\vee$. By (2) it suffices to take $R=A^\vee\cap N^\perp$ with $N$
saturated and $N\ne A$, and by (1) the condition becomes
$\widehat{\deg}(A)-\widehat{\deg}(N)\ge-\theta r(\operatorname{rk}A-\operatorname{rk}N)$
for these $N$. As $\widehat{\deg}(A/N)=\widehat{\deg}(A)-\widehat{\deg}(N)$, this
says $\mu_{\min}(A)\ge-\theta$. Hence $\mu_{\max}(A^\vee)=-\mu_{\min}(A)$.

Similarly $\mu_{\min}(A^\vee)\ge\theta$ exactly when
$\widehat{\deg}(A^\vee)-\widehat{\deg}(R)\ge\theta r(\operatorname{rk}A^\vee-\operatorname{rk}R)$
for every saturated sublattice $R\ne A^\vee$. By (2), such an $R$ is
$A^\vee\cap N^\perp$ for a saturated $N\ne0$, and by (1) the condition becomes
$-\widehat{\deg}(N)\ge\theta r\operatorname{rk}N$. By
Proposition~\ref{hs-hn-polygon}(1), $\mu_{\max}(A)$ is the slope of a nonzero
saturated sublattice, so the condition for all $N$ says
$\mu_{\max}(A)\le-\theta$. Hence $\mu_{\min}(A^\vee)=-\mu_{\max}(A)$.
\end{proof}

\begin{lemma}\label{asm-bounds-lattice}
Let $\mathcal F$ be a family with exponent $\varepsilon$, and let
$C_{\mathrm{abs}},C_H$ be real. Suppose that $\mathcal F$ satisfies the
division link with constants $C_{\mathrm{abs}},C_H$
(Definition~\ref{is-division-link}), and that for all large $x$, at every
member at $x$, every division lattice $N_u$ is finitely generated, spans its
space and has positive rank, and the glued lattice $RR$ has positive rank.
Then $\mathcal F$ satisfies the lattice conditions at all large scales
(Definition~\ref{asm-member-bounds}(4)).
\end{lemma}
\leanref{Formalism/SharpRestriction/GlobalSchurBounds/QuadraticLogDerivCongr.lean}{Formalism.MainNotions.Lnk.latF}
\leanref{Formalism/SharpRestriction/GlobalSchurBounds/QuadraticLogDerivCongr.lean}{Formalism.MainNotions.Lnk.lat_at}
\leanref{Formalism/LocalRingCardinality/RankBoundMonotoneInDimension.lean}{Formalism.BasicNotions.bridge_dual_mu_holds}
\begin{proof}
Let $D$ be a member at a large scale and $u$ a word. Then $d_u>0$ by
hypothesis, and $N_u$ is a nonzero Hermitian lattice, so
$\mu(N_u)\le\mu_{\max}(N_u)$ by Proposition~\ref{hs-hn-polygon}(1). By
condition (5) of the division link, $RR$ is finitely generated and spans its
space, and it has positive rank. By Lemma~\ref{asm-bounds-dual-slopes}, its
dual $CC=RR^\vee$ is a Hermitian lattice spanning the dual space with the dual
norms, and $\mu_{\min}(CC)=-\mu_{\max}(RR)$.
\end{proof}

\begin{lemma}\label{asm-bounds-frame-realised}
Let $D$ be a definite-side datum (Definition~\ref{th-definite-datum}) with
coefficient field $L$, set of words $\mathcal W$, spaces $M_u$, number $D^{*}>0$,
row deficit $\delta$, numbers of rows $D_{\mathrm{all}}$ and
$D_{\mathrm{rows}}=(1-\delta)D_{\mathrm{all}}$, Bessel-period data and Gram
matrices $K_{uu'}$ (Definition~\ref{is-orthogonality-link}), and division
datum $RR$, $N_u$, $B_u$, $h_u$ (Definition~\ref{hs-frame-realisation}), with
$CC=RR^\vee$ and $d_u=\operatorname{rk}N_u$. Let $\sigma,\delta_e,e_B$ and
$\eta_0,\eta_1,e_0$ be real numbers with $-1\le\eta_0\le\sigma/2$,
$\eta_1\le\sigma/2$, $e_0\ge0$, $\sqrt{e_0}\le\sigma/2$, $e_0<1$ and
$2e_0\le e_B$. Suppose that:
\begin{enumerate}
\item the Bessel identity holds for $D$ (Definition~\ref{th-bessel-identity});
  $RR$ is finitely generated and spans its space; and for every word $u$,
  $B_u(N_u)\subset CC$, and $h(B_u)\le h_u$ if $B_u\ne0$;
\item every $N_u$ is finitely generated, spans its space and has positive
  rank;
\item $0\le\delta<1$;
\item $\dim M_u>0$ for every word $u$, and the orthogonality bounds (1) to (3)
  of Definition~\ref{is-orthogonality-link} hold with $(\eta_0,\eta_1,e_0)$;
\item $\dim M_u-d_u\le\delta_e\dim M_u$ for every word $u$.
\end{enumerate}
Then $D$ is frame-realised with parameters $(\sigma,\delta_e,e_B)$
(Definition~\ref{hs-frame-realisation}).
\end{lemma}
\leanref{Formalism/SharpRestriction/GlobalSchurBounds/TopDensityNGeSixteen.lean}{Formalism.MainNotions.Lnk.frameReal_at_of}
\leanref{Formalism/SharpRestriction/GlobalSchurBounds/TopDensityNGeSixteen.lean}{Formalism.MainNotions.Lnk.matS_at_of}
\leanref{Formalism/SharpRestriction/GlobalSchurBounds/TopDensityNGeSixteen.lean}{Formalism.MainNotions.Lnk.matK_at_of}
\leanref{Formalism/SharpRestriction/GlobalSchurBounds/TopDensityNGeSixteen.lean}{Formalism.MainNotions.Lnk.blockReal_at_of}
\leanref{Formalism/SharpRestriction/GlobalSchurBounds/TopDensityNGeSixteen.lean}{Formalism.MainNotions.Lnk.unfolding_constant_pos_of_diagonal}
\begin{proof}
\emph{Coordinates.} Let $v$ be the infinite place of $L$ given by its embedding
in $\mathbf C$. The Bessel identity gives a unitary map $V$ of $\mathcal E$
and, for each word $u$, a number $a_u\ne0$ and a unitary map $U_u$ of
$\mathcal M_u$ with $B_u(\varphi,F)=a_u\mathcal B_u(U_u\varphi,VF)$, and the
pairing $B_u$ is the image of an $L$-rational pairing. Choose an orthonormal
basis $f_1,\dots,f_{D_{\mathrm{all}}}$ of $\mathcal H$ whose first
$D_{\mathrm{rows}}$ vectors are the images under $s\circ V$ of an orthonormal
basis of $\mathcal E$. Writing $B_u$, $\mathcal B_u$ and $P_u$ in these
coordinates (after complex conjugation if the embedding of $L$ is the
conjugate of the one defining $v$), one gets invertible matrices $Y_u$ and $Y'$
and, for each $u$, a matrix $V_u$ with $\dim M_u$ rows, $d_u$ columns and
$V_u^{*}V_u=1$, such that $Y'^{-1}B_{u,v}Y_u$ consists of the first
$D_{\mathrm{rows}}$ rows of $\hat B_uV_u$, where $\hat B_u$ is the matrix of
$P_u$ in the basis $f$. Here $a_u^{-1}$ is absorbed in $Y_u$, and $V_u$ comes
from $U_u$ and the isometry $\mathcal C_u\cong M_u$ restricted to
$\mathcal M_u\cong M_{u,\mathrm{ram}}$. This is a computation with the
coordinates of the Bessel identity.

\emph{The frame.} Take the inner products on $V_{u,v}$ and
$(V_{RR}^\vee)_v$ for which the columns of $Y_u$ and of $Y'$ are orthonormal,
$\ell=D_{\mathrm{all}}-D_{\mathrm{rows}}$, $F_u=\mathbf C^{\dim M_u}$,
$\nu_u=V_u$, $\hat B_u$ as above and $\gamma_u=c_u^{-1}$. The number $c_u$ is
positive: it is nonnegative, and $c_u=0$ would give $K_{uu}=0$ and
$\|K_{uu}-1\|_{HS}^2=\dim M_u>e_0\dim M_u$. We check the conditions of
Definition~\ref{hs-frame-realisation}.
\begin{enumerate}
\item The dual $CC$ of $RR$ is finitely generated and spans the dual space; the
  $N_u$ are finitely generated and span; $\sum_u\dim V_u=\sum_ud_u>0$; and the
  conditions on $B_u$ are hypothesis (1).
\item $0\le\delta<1$, and
  $(1-\delta)(D_{\mathrm{rows}}+\ell)=(1-\delta)D_{\mathrm{all}}=D_{\mathrm{rows}}$,
  the dimension of $V_{RR}$.
\item This is the choice of $Y_u$, $Y'$ and $V_u$.
\item These are the assumptions on $\eta_0,\eta_1,e_0$.
\item $\dim V_u=d_u$, so this is hypothesis (5) summed over the words.
\item By definition of $K_{uu'}$, $\gamma_u\gamma_{u'}\|\hat B_u^{*}\hat
  B_{u'}\|_{HS}^2=\|K_{uu'}\|_{HS}^2$ and
  $\|\gamma_u\hat B_u^{*}\hat B_u-1\|_{HS}^2=\|K_{uu}-1\|_{HS}^2$; and
  $(V_{RR}^\vee)_v\oplus\mathbf C^\ell$ has dimension $D_{\mathrm{all}}$. So the
  orthogonality bounds (1) to (3) are the conditions of near orthogonality
  (Definition~\ref{hs-near-orthogonal}).
\end{enumerate}
\end{proof}

\begin{lemma}\label{asm-bounds-frame}
Let $\mathcal F$ be a family with exponent $\varepsilon$, and let
$C_{\mathrm{abs}}$ be real, $C_0\ge1$ and $C_H>0$. Suppose that $\mathcal F$
satisfies the orthogonality link with constant $C_0$, the irreducibility link
and the division link with constants $C_{\mathrm{abs}},C_H$, and that for all
large $x$ every member at $x$ satisfies: $C_0^{-1}\le\dim
M_u\cdot\Gamma^2/D^{*}\le C_0$ for every word $u$, and $C_0\Gamma^2x^{C_H}\le|W|$;
every $N_u$ is finitely generated, spans its space and has positive rank $d_u$;
$d_u=\dim M_u\,(1-1/(T_u-1))^f$ and $b^{m/2}\le T_u$ for every word $u$; and the
lattice conditions (Definition~\ref{asm-member-bounds}(4)). Then $\mathcal F$
satisfies the frame bound with $C_H$ (Definition~\ref{asm-member-bounds}).
\end{lemma}
\leanref{Formalism/SharpRestriction/GlobalSchurBounds/TopDensityNGeSixteen.lean}{Formalism.MainNotions.Lnk.frameF_of}
\leanref{Formalism/SharpRestriction/GlobalSchurBounds/TopDensityNGeSixteen.lean}{Formalism.MainNotions.Lnk.frameRealF_of''}
\leanref{Formalism/SharpRestriction/GlobalSchurBounds/TopDensityNGeSixteen.lean}{Formalism.MainNotions.Lnk.frameRealF_of}
\leanref{Formalism/SharpRestriction/GlobalSchurBounds/TopDensityNGeSixteen.lean}{Formalism.MainNotions.Lnk.matRealF_of}
\leanref{Formalism/SharpRestriction/GlobalSchurBounds/TopDensityNGeSixteen.lean}{Formalism.MainNotions.Lnk.matSF_of}
\leanref{Formalism/SharpRestriction/GlobalSchurBounds/TopDensityNGeSixteen.lean}{Formalism.MainNotions.Lnk.matKF_of}
\leanref{Formalism/SharpRestriction/GlobalSchurBounds/TopDensityNGeSixteen.lean}{Formalism.MainNotions.Lnk.ramShareF_of}
\leanref{Formalism/SharpRestriction/GlobalSchurBounds/TopDensityNGeSixteen.lean}{Formalism.MainNotions.Lnk.ramShare_at_of}
\begin{proof}
Write $E=\exp(-x^{\varepsilon/8})$ and $\sigma=x^{-C_H}$. Let $x$ be a large
scale, larger than the thresholds of the hypotheses and such that $x\ge18$,
$x\ge20^{1/C_H}$, $x^{C_H+1}\le\exp(x^{\varepsilon/8})$,
$x^6\le\exp(x^{\varepsilon}/2)$ and $4x^{\varepsilon/8}\le x^{\varepsilon}$; all
these hold for large $x$. Let $D$ be a member at $x$.

\emph{The ramified share.} For every word $u$, $T_u\ge b^{m/2}\ge b\ge3$, since
$m\ge2$, and
\[
 1-\frac{d_u}{\dim M_u}=1-\Big(1-\frac1{T_u-1}\Big)^f\le\frac f{T_u-1}\le\frac{2f}{T_u}
 \le\frac{2f}b,
\]
so $\dim M_u-d_u\le(2f/b)\dim M_u$. The window $k\le x^5$ of the sizes of the
family and $f\le k$ give $2f\le x^6\le\exp(x^{\varepsilon}/2)$, so
$2f/b\le\exp(-x^{\varepsilon}/2)\le E^2$.

\emph{The parameters.} The orthogonality link gives
$\eta_0,\eta_1\le E$ and $e_0\ge0$ with $\sqrt{e_0}\le E$ for which the
orthogonality bounds hold; replacing $\eta_0$ by $\max(\eta_0,0)$ keeps them,
so $-1\le\eta_0$. Since $Ex^{C_H}\cdot x\le1$ and $x\ge18$, we have
$Ex^{C_H}\le1/18$ and $E\le\sigma/2$. Put $\delta_e=E^2$ and $e_B=2E^2$. Then
$\eta_0,\eta_1\le\sigma/2$, $\sqrt{e_0}\le\sigma/2$, $e_0<1$, $2e_0\le e_B$,
$0\le\delta_e,e_B\le1/4$, $\dim M_u-d_u\le\delta_e\dim M_u$, and
\[
 (3\sqrt{\delta_e}+4\sqrt{e_B})x^{C_H}\le11Ex^{C_H}\le1.
\]

\emph{Frame realisation.} The irreducibility link gives $0\le\delta<1$, and the
division link gives its conditions (1), (5) and (6) and, as its condition (4),
the Bessel identity (Definition~\ref{is-division-link}). The bounds on $\dim M_u\cdot\Gamma^2/D^{*}$
give $\dim M_u>0$. So Lemma~\ref{asm-bounds-frame-realised} applies: $D$ is
frame-realised with parameters $(\sigma,\delta_e,e_B)$.

\emph{The frame bound.} Since $x\ge20^{1/C_H}$, $\sigma\le1/20$.
Proposition~\ref{hs-frame-at-glued-lattice}, with $C_0\ge1$, the bounds on $\dim
M_u\cdot\Gamma^2/D^{*}$ and on $|W|$, the $N_u$ and the lattice conditions,
gives the frame bound with parameters $(\eta_B,e_B)$, where
$\eta_B=5\sigma+3\sqrt{\delta_e}\ge0$. Finally $e_B\le1/4$ and
\[
 (\eta_B+4\sqrt{e_B})x^{C_H}=5+(3\sqrt{\delta_e}+4\sqrt{e_B})x^{C_H}\le6.
\]
These $\eta_B$ and $e_B$ depend only on $x$, so $\mathcal F$ satisfies the frame
bound with $C_H$.
\end{proof}

\begin{lemma}\label{asm-bounds-nodes-characters}
Let $\mathcal F$ be a family with exponent $\varepsilon$
(Definition~\ref{pd-family}).
\begin{enumerate}
\item Every member at every scale $x\ge3$ with $x>x_1$ satisfies the node
  bounds (Definition~\ref{asm-member-bounds}(7)).
\item For all real $c_6$ and $C_5$, every member at every scale satisfies the
  character conditions with $(c_6,C_5)$ (Definition~\ref{asm-member-bounds}(8)).
\end{enumerate}
\end{lemma}
\leanref{Formalism/SharpRestriction/GlobalSchurBounds/QuadraticLogDerivCongr.lean}{Formalism.MainNotions.Lnk.nodesF}
\leanref{Formalism/SharpRestriction/GlobalSchurBounds/QuadraticLogDerivCongr.lean}{Formalism.MainNotions.Lnk.nodes_at}
\leanref{Formalism/SharpRestriction/GlobalSchurBounds/QuadraticLogDerivCongr.lean}{Formalism.MainNotions.Lnk.numN_at}
\leanref{Formalism/SharpRestriction/GlobalSchurBounds/TopDensityNGeSixteen.lean}{Formalism.MainNotions.Lnk.charF}
\leanref{Formalism/SharpRestriction/GlobalSchurBounds/TopDensityNGeSixteen.lean}{Formalism.MainNotions.Lnk.char_at}
\begin{proof}
(1) Let $D$ be a member at such a scale $x$, and put $w''=w''(n)=b'/(xn)$. The
size data of $D$ place the Harish-Chandra weights of every lift $j$ in windows:
$b\le w_j(\tau)\le b+\eta+1$ if $\tau$ is a thick place;
$b_t+b'/x\le w_j(\tau)\le b_t+b'/x+\eta$ if $\tau$ is the place $\tau_1(t)$ of a
top class $t$; and $w''\le w_j(\tau)\le w''+\eta$ otherwise. The parameter
choices and the sizes of the family give, for $x\ge3$, the inequalities
$1\le w''\le b'$, $w''+\eta\le2b'$ and $b_t+b'/x+\eta\le2b'$, and the size data
give $b_t\ge b'$; this is a computation with the quantities of
Definition~\ref{pd-size-data}(1). An element $\rho\in\Psi_A$ is thick at the
base point exactly when it codes a thick place. If $\rho$ codes no place, then
$w_j(\rho)=1$ and $\rho$ is not thick. In the three windows $w_j(\rho)\ge
\min(b,b',w'')\ge1$; at a thick place $b\le w_j(\rho)\le b+\eta+1$; and
elsewhere $w_j(\rho)\le2b'$. These are the node bounds.

(2) Let $D$ be a member at $x$. For a character $\psi\ne1$ of $E_0\times A$,
the arithmetic data of the host attach to $\psi$ a Dirichlet character, induced
by a unique primitive character $\chi_\psi$; put
$\ell(\psi)=\frac{L'}L(1,\chi_\psi)$. Two primitive Dirichlet characters that
agree at all integers $n\ne0$ are equal, so $\ell$ is compatible with the host.
The character $\chi_C$ is the quadratic character of $\mathbf Q(\sqrt{-d})$
(Remark~\ref{ib-host-group-fields}), so
$\operatorname{Re}\ell(\chi_C)=\operatorname{Re}\frac{L'}L(1,\chi_d)$. The
corrected Chebyshev sum of $\psi$ for $\ell$ is the sum $g_{\chi_\psi}$ of the
Dirichlet character $\chi_\psi$ of the host field, so the bound for all these
characters at $y\ge X'/3$ gives (W4) with $X'$ and
$\theta=\exp(-c_6x^{\varepsilon/8})$. Finally, the primitive characters
attached to the three nontrivial characters of $E_0$ are those induced by
$\chi_{-q}$, $\chi_{-d}$ and $\chi_{-q}\chi_{-d}$, so the bounds for these give
the bounds for the three characters.
\end{proof}

\begin{lemma}\label{asm-bounds-windows}
Let $\mathcal F$ be a family with exponent $\varepsilon$ and let $c_6>0$. Let
$x>x_1$ be a scale with $x\ge3$, $x^5\le\exp(x^{\varepsilon/3})$,
$x^7\le\exp(c_6x^{\varepsilon/8})$, $x^9\le\exp(x^{\varepsilon})$,
$x^5\le\exp(x^{\varepsilon}/2)$ and $2\le x^{\varepsilon-\varepsilon/5-\varepsilon/100}$,
and let $D$ be a member at $x$ whose set $T_2$ satisfies
$|T_2|\ge x^{-\varepsilon}n$. Then $D$ satisfies the window conditions with
$(5,c_6)$ (Definition~\ref{asm-member-bounds}(9)).
\end{lemma}
\leanref{Formalism/SharpRestriction/GlobalSchurBounds/QuadraticLogDerivCongr.lean}{Formalism.MainNotions.Lnk.size_window_at_five}
\leanref{Formalism/SharpRestriction/GlobalSchurBounds/QuadraticLogDerivCongr.lean}{Formalism.MainNotions.Lnk.winE}
\leanref{Formalism/SharpRestriction/GlobalSchurBounds/QuadraticLogDerivCongr.lean}{Formalism.MainNotions.Lnk.winP}
\leanref{Formalism/SharpRestriction/GlobalSchurBounds/LossesSmallShares.lean}{Formalism.MainNotions.Lnk.size_windows_of_thick_density}
\leanref{Formalism/SharpRestriction/GlobalSchurBounds/LossesSmallShares.lean}{Formalism.MainNotions.Lnk.thick_density_of_construction}

This is a computation from the sizes of the family, which place $n$, $m$ and
$k$ within the powers $x^{\pm\theta}$ of $x^{2+8\varepsilon}$,
$\sqrt{x\delta_1}$ and $x^{9/2+10\varepsilon}$ for a small fixed $\theta$, and
$\log q$, $\log d$ and the root discriminant $r_E$ within fixed powers of $x$,
from the parameter choices, and from the five inequalities between powers of $x$
and exponentials. For instance, $L_a=\log(ar_E)$ with $\log a=\log
x+x^{\varepsilon}$ and $r_E\le2nx^{\varepsilon/100}\sqrt{qd}$, and the
inequality $x^5\le\exp(x^{\varepsilon}/2)$ gives $0\le L_a\le2x^{\varepsilon}$;
and $v>0$ because $r_E\ge1$.

\subsection{The level primes}\label{subsection-asm-level-primes}

\begin{definition}\label{asm-level-setting}
Let $D$ be a definite-side datum (Definition~\ref{th-definite-datum}) with
ranks $k,m,t,n$, $f=2n$ and $N(j)=k+j(t+1)$, coefficient field $L$, and
lattices of forms $\Lambda_{u,j}\subset V_{u,j}$ with injections
$\operatorname{emb}_j$ (Definition~\ref{pd-period-data}). Let $u$ be a word,
with auxiliary prime $T_u$, space $M_u$ and subspace $M_{u,\mathrm{ram}}$
(Definition~\ref{th-definite-datum}(3)), dominant weights $\lambda_{u,j}(\tau)$
(Definition~\ref{th-definite-datum}(7)), and Hecke data at $T_u$ with top
embeddings $\mathcal T_{\mathrm{top}}$, sets $V_{\mathrm{top}}(w)$, bijections
$\tau_w$, numbers $b_{\mathrm{top}}(\tau)$ and $W_\tau(u)$ and local lattices
$N_\lambda(u,w)$ (Definition~\ref{pd-hecke-at-T}). Put
$a_{\max}(u)=\max\big(0,\max_{\tau\in\mathcal T_{\mathrm{top}}}\lambda_{u,0}(\tau)_1\big)$.
Let $V^{\mathrm{ram}}_{u,0}$ be the $L$-subspace of the $y\in V_{u,0}$ with
$\operatorname{emb}_0(y)\in M_{u,\mathrm{ram}}$, and let the \emph{ramified
lattice} $\Lambda^{\mathrm{ram}}_u=\Lambda_{u,0}\cap V^{\mathrm{ram}}_{u,0}$
carry the metrics of $\Lambda_{u,0}$. For a finite place $w$ of $L$, with
valuation ring $\mathcal O_{L,w}\subset L$ and ramification index $e_w$, let
$S_w$ be the $\mathcal O_{L,w}$-span of $\Lambda_{u,0}$ in $V_{u,0}$, and
$S^{\mathrm{ram}}_w=S_w\cap V^{\mathrm{ram}}_{u,0}$.
\begin{enumerate}
\item The \emph{setting of the level primes} holds at $u$ if
  $2fN(m)+1<T_u$; $4a_{\max}(u)<T_u$;
  $2\big(\lambda_{u,j}(\tau)_a-\lambda_{u,j}(\tau)_b+N(j)\big)<T_u$ for all
  $j\le m$, $\tau\in\mathcal T$ and $1\le a,b\le N(j)$; and
  $b_{\mathrm{top}}(\tau)\ge0$ for every $\tau\in\mathcal T_{\mathrm{top}}$.
\item The \emph{division side conditions} hold at $u$ if
  \[
   \dim M_{u,\mathrm{ram}}=\dim M_u\Big(1-\frac1{T_u-1}\Big)^f,\qquad
   \mu(\Lambda^{\mathrm{ram}}_u)=\mu(\Lambda_{u,0})-\frac{f\log(T_u-1)}{T_u-2},
  \]
  and $\mu_{\max}(\Lambda^{\mathrm{ram}}_u)\le\mu_{\max}(\Lambda_{u,0})$
  (Definition~\ref{hs-slopes}).
\item The \emph{congruence length bound} holds at $u$ if for every place $w$
  of $L$ over $T_u$ we have $N_\lambda(u,w)\subset S^{\mathrm{ram}}_w$, and
  the $\mathcal O_{L,w}$-module $S^{\mathrm{ram}}_w/N_\lambda(u,w)$ has finite
  length, at most
  \[
   e_w\,m\dim M_u\sum_{v\in V_{\mathrm{top}}(w)}\Big(\frac{a_{\max}(u)}k
   +2W_{\tau_w(v)}(u)+k+b_{\mathrm{top}}(\tau_w(v))\Big).
  \]
\end{enumerate}
A family $\mathcal F$ (Definition~\ref{pd-family}) has the \emph{level data} if,
for all large $x$, the division side conditions and the congruence length
bound hold at every word of every member at $x$.
\end{definition}
\leanref{Formalism/SharpRestriction/GlobalSchurBounds/ProductFirstZerothParts.lean}{Formalism.MainNotions.product_first_zeroth_parts.HidaSetting}
\leanref{Formalism/SharpRestriction/GlobalSchurBounds/ProductFirstZerothParts.lean}{Formalism.MainNotions.product_first_zeroth_parts.Alcove}
\leanref{Formalism/SharpRestriction/GlobalSchurBounds/ProductFirstZerothParts.lean}{Formalism.MainNotions.product_first_zeroth_parts.amax}
\leanref{Formalism/SharpRestriction/GlobalSchurBounds/JointsF.lean}{Formalism.MainNotions.GlobalTheoryTheory.SettingEventually}
\leanref{Formalism/SharpRestriction/GlobalSchurBounds/LinkMin.lean}{Formalism.MainNotions.product_first_zeroth_parts.family_property_and}
\leanref{Formalism/LocalToGlobal/GlobalTheory/GL6/Statement.lean}{Formalism.L2.GlobalTheoryTheory.LocalLatticeBases.CongruenceLatticeLengthBound}
\leanref{Formalism/LocalRingCardinality/DataLatDefs.lean}{Formalism.BasicNotions.LT.Sram}
\leanref{Formalism/SharpRestriction/GlobalSchurBounds/LinkMin.lean}{Formalism.MainNotions.J.LvlF}

\begin{lemma}\label{asm-level-setting-holds}
Let $\mathcal F$ be a family with exponent $\varepsilon$
(Definition~\ref{pd-family}) that satisfies the size condition on $T_u$
(Definition~\ref{asm-designed}(2)). Then, for all large $x$, the setting of the
level primes (Definition~\ref{asm-level-setting}(1)) holds at every word of
every member at $x$.
\end{lemma}
\leanref{Formalism/LocalToGlobal/GlobalTheory/PlacesP.lean}{Formalism.L2.GlobalTheoryTheory.Places.settingF_of}
\leanref{Formalism/SharpRestriction/AdmissibleParameterJoins/DIVfromGLo.lean}{Formalism.MainNotions.J.levelPrime_setting_of_sizes}
\leanref{Formalism/SharpRestriction/GlobalSchurBounds/AlcoveOf.lean}{Formalism.MainNotions.GlobalTheoryTheory.Alc.alcove_of}
\leanref{Formalism/SharpRestriction/AdmissibleParameterJoins/DIVfromGLo.lean}{Formalism.MainNotions.J.amax_le}
\leanref{Formalism/SharpRestriction/GlobalSchurBounds/GlobalRankLatticeProofs/SevenLogLeLOfPowerBound.lean}{Formalism.MainNotions.flipTwoThree.nums_wts_of_tuSize}
\begin{proof}
Let $x$ be a large scale, $D$ a member at $x$ with ranks $k,m,t$ and the
quantities $b=\exp(x^{\varepsilon})$ and $b'$ of its size data, and $u$ a word.
The parameter choices and the sizes of the family, with the bound $\log
T_u\le C_{T_u}mx^{\varepsilon}$ of the size condition and the bound
$|\frac{L'}{L}(1,\chi_d)|\le x^{\varepsilon}$ of the size data, give the
following for all large $x$; this is a computation with the quantities of
Definition~\ref{pd-size-data}, and (ii) and (iii) read the entries of
$\lambda_{u,j}(\tau)$ off Definition~\ref{th-definite-datum}(7).
\begin{enumerate}
\item[(i)] $m\ge4$, $24x\le b'\le b$ and $3N(m)\le b'$.
\item[(ii)] $-(2b+b')\le\lambda_{u,j}(\tau)_a\le2xb+b'^2/64$ for all $j\le m$,
  $\tau\in\mathcal T$ and $1\le a\le N(j)$.
\item[(iii)] $\lambda_{u,0}(\tau)_1\le2xb$ and $b'\le b_{\mathrm{top}}(\tau)$
  for every $\tau\in\mathcal T_{\mathrm{top}}$.
\end{enumerate}
By (i), $b\ge24$ and $xb\le b^2/24$. The size condition gives $T_u\ge
b^{m/2}\ge b^2$, since $m\ge4$. We check the four conditions.
\begin{enumerate}
\item Since $2n\le k$ (Definition~\ref{th-ranks}), $f\le k\le N(m)$, and
  $9N(m)^2\le b^2$ by (i). So $2fN(m)+1\le2N(m)^2+1<9N(m)^2\le T_u$.
\item By (iii), $0\le a_{\max}(u)\le2xb$, so $4a_{\max}(u)\le8xb<b^2\le T_u$.
\item By (ii) and $N(j)\le N(m)\le b'\le b$,
  \[
   2\big(\lambda_{u,j}(\tau)_a-\lambda_{u,j}(\tau)_b+N(j)\big)
   \le4xb+\frac{b'^2}{32}+4b+4b'
   \le\Big(\frac16+\frac1{32}+\frac16+\frac16\Big)b^2<b^2\le T_u.
  \]
\item By (iii), $b_{\mathrm{top}}(\tau)\ge b'>0$ for
  $\tau\in\mathcal T_{\mathrm{top}}$.
\end{enumerate}
\end{proof}

\begin{definition}\label{asm-level-local-facts}
Let $D$ be a definite-side datum and $u$ a word, with the notation of
Definition~\ref{asm-level-setting}. Let $\mathbf F_v$ ($v\in V_u$) be the $f$
fields with $T_u$ elements and $\rho_v$ the actions of $\mathbf F_v^\times$ of
Definition~\ref{th-definite-datum}(3), and $U_v$ the Hecke operators at $T_u$
(Definition~\ref{pd-hecke-at-T}); for $a\in\mathbf F_v^\times$ let
$a_v\in\prod_{v'\in V_u}\mathbf F_{v'}^\times$ have entry $a$ at $v$ and $1$
elsewhere. The forms of $u$ also carry \emph{down data}: for every level
$i\le m$, every $v\in V_u$ and every class $c\in\mathrm{cls}_i$, an $\mathbf
F_v$-vector space $k_{i,v}(c)$ of dimension $N(i)$ with a vector $t_{i,v,c}$,
and for every linear form $g$ on $k_{i,v}(c)$ a class $c^g\in\mathrm{cls}_i$
and a linear automorphism $\rho^g$ of the functions on $Y_i$; and, for every
finite place $w$ of $L$ and every $v\in V_u$, an $\mathcal O_{L,w}$-submodule
$B_w(v)\subset V_{u,0}$ that spans $V_{u,0}$ over $L$, the \emph{second branch
lattice}. The \emph{down operator} at $v$ on functions $\varphi$ on
$\mathrm{cls}_i\times Y_i$ is
\[
 (D_v\varphi)(c)=\sum_{g:\ g(t_{i,v,c})=1}\rho^g\big(\varphi(c^g)\big).
\]
The \emph{local facts} at $u$ are the following.
\begin{enumerate}
\item A \emph{free basis at} $T_u$: a finite set $X$ with a free action of
  $\prod_{v\in V_u}\mathbf F_v^\times$ and a linear isomorphism
  $\Theta\colon M_u\to\mathbf C^X$ with
  $\Theta(\rho_v(a)\varphi)(x)=\Theta(\varphi)(a_vx)$ for all $v\in V_u$,
  $a\in\mathbf F_v^\times$, $\varphi\in M_u$ and $x\in X$.
\item The \emph{slope line}:
  $\mu(\Lambda^{\mathrm{ram}}_u)=\mu(\Lambda_{u,0})-f\log(T_u-1)/(T_u-2)$.
\item \emph{Hecke integrality at level zero}: for every place $w$ over $T_u$
  that sees $\pi_u$ and every $v\in V_u$, $U_v$ maps
  $\operatorname{emb}_0(B_w(v))$ into $\operatorname{emb}_0(S_w)$, and $D_v$
  maps $\operatorname{emb}_0(S_w)$ into $\operatorname{emb}_0(B_w(v))$.
\item The \emph{ramified level zero}: for every $v\in V_u$, $U_v$ and $D_v$ map
  $M_{u,\mathrm{ram}}$ into itself, and $U_vD_v\varphi=T_u^{k-1}\varphi$ for
  $\varphi\in M_{u,\mathrm{ram}}$.
\item The \emph{second branch inclusion}: for every place $w$ over $T_u$ that
  sees $\pi_u$ and every $v\in V_u$ with $b=b_{\mathrm{top}}(\tau_w(v))\ge0$,
  $T_u^bS_w\subset B_w(v)$, and the $\mathcal O_{L,w}$-module
  $B_w(v)/T_u^bS_w$ has finite length
  \[
   e_w\dim M_u\,\frac{\sum_{a=1}^k\lambda_{u,0}(\tau_w(v))_a+kb}k.
  \]
\end{enumerate}
\end{definition}
\leanref{Formalism/SharpRestriction/GlobalSchurBounds/ProductFirstZerothParts.lean}{Formalism.MainNotions.product_first_zeroth_parts.Dop}
\leanref{Formalism/SharpRestriction/GlobalSchurBounds/FiniteModelOfAdmissibleTransports.lean}{Formalism.MainNotions.product_first_zeroth_parts.FreeT}
\leanref{Formalism/SharpRestriction/GlobalSchurBounds/FiniteModelOfAdmissibleTransports.lean}{Formalism.MainNotions.product_first_zeroth_parts.SlopeLine}
\leanref{Formalism/LocalToGlobal/GlobalTheory/GL6/CongruenceNoKernelOfHeckeRamified.lean}{Formalism.L2.GlobalTheoryTheory.LocalLatticeBases.HeckeIntegralLevelZero}
\leanref{Formalism/LocalToGlobal/GlobalTheory/GL6/CongruenceNoKernelOfHeckeRamified.lean}{Formalism.L2.GlobalTheoryTheory.LocalLatticeBases.HeckeRamifiedLevelZero}
\leanref{Formalism/LocalToGlobal/GlobalTheory/GL6/CongruenceNoKernelOfHeckeRamified.lean}{Formalism.L2.GlobalTheoryTheory.LocalLatticeBases.SecondBranchInclusionLength}

\begin{lemma}\label{asm-level-free-count}
Let $T$ be a prime, $I$ a finite set with $f$ elements, $\mathbf F_v$ ($v\in I$)
fields with $T$ elements, and $G=\prod_{v\in I}\mathbf F_v^\times$; for
$a\in\mathbf F_v^\times$ let $a_v\in G$ have entry $a$ at $v$ and $1$
elsewhere. Let $X$ be a finite set with a free action of $G$, and let $R$ be the
space of functions $\xi\colon X\to\mathbf C$ with
$\sum_{a\in\mathbf F_v^\times}\xi(a_vx)=0$ for every $v\in I$ and $x\in X$. Then
$\dim R\cdot(T-1)^f=|X|\,(T-2)^f$.
\end{lemma}
\leanref{Formalism/LogVandermondeBounds/LogVandermondeProofs/MulSingleMulApplySelf.lean}{Formalism.LogVandermondeBounds.product_first_zeroth_parts.finrankFiberSumZero}
\leanref{Formalism/SharpRestriction/GlobalSchurBounds/FiniteModelOfAdmissibleTransports.lean}{Formalism.MainNotions.product_first_zeroth_parts.rank_nat}
\begin{proof}
The action is free, so there are a finite set $Q$ and a $G$-equivariant
bijection $X\cong G\times Q$, with $G$ acting on $G\times Q$ by left
multiplication on the first factor; then $|X|=(T-1)^f|Q|$. A function on
$G\times Q$ is a family, indexed by $Q$, of functions on $G$, and
$\mathbf C^G=\bigotimes_{v\in I}\mathbf C^{\mathbf F_v^\times}$. The condition at
$v$ says that the function is killed by $1\otimes\cdots\otimes
P_v\otimes\cdots\otimes1$, where $P_v$ sends a function on
$\mathbf F_v^\times$ to the constant function with value the sum of its values.
The kernel of $P_v$ is the space $K_v$ of functions with sum $0$, of dimension
$T-2$. Choosing in each $\mathbf C^{\mathbf F_v^\times}$ a basis of $K_v$
completed by the constant function $1$, on which $P_v$ is injective, one sees
that the functions killed by all these operators form $\bigotimes_vK_v$, of
dimension $(T-2)^f$. Hence $\dim R=(T-2)^f|Q|$.
\end{proof}

\begin{lemma}\label{asm-level-division-side}
Let $D$ be a definite-side datum and $u$ a word, with the notation of
Definitions~\ref{asm-level-setting} and~\ref{asm-level-local-facts}. Suppose
that the setting of the level primes holds at $u$, that $M_u$ has a free basis
at $T_u$, and that the slope line holds at $u$. Then the division side
conditions hold at $u$.
\end{lemma}
\leanref{Formalism/SharpRestriction/GlobalSchurBounds/LinkMin.lean}{Formalism.MainNotions.product_first_zeroth_parts.division_side_conditions_at_sizeData}
\leanref{Formalism/SharpRestriction/GlobalSchurBounds/FiniteModelOfAdmissibleTransports.lean}{Formalism.MainNotions.product_first_zeroth_parts.rank_line}
\leanref{Formalism/SharpRestriction/GlobalSchurBounds/FiniteModelOfAdmissibleTransports.lean}{Formalism.MainNotions.product_first_zeroth_parts.muMax_line'}
\leanref{Formalism/SharpRestriction/GlobalSchurBounds/FiniteModelOfAdmissibleTransports.lean}{Formalism.MainNotions.product_first_zeroth_parts.muMax_line}
\begin{proof}
\emph{Dimension.} Let $(X,\Theta)$ be a free basis at $T_u$. Since
$\Theta\big(\sum_a\rho_v(a)\varphi\big)(x)=\sum_a\Theta(\varphi)(a_vx)$, the
isomorphism $\Theta$ carries
$M_{u,\mathrm{ram}}=M_u\cap\bigcap_v\ker\sum_{a\in\mathbf F_v^\times}\rho_v(a)$
onto the space $R$ of Lemma~\ref{asm-level-free-count} for the $f$ fields
$\mathbf F_v$ with $T_u$ elements, and $|X|=\dim M_u$. That lemma gives
$\dim M_{u,\mathrm{ram}}\,(T_u-1)^f=\dim M_u\,(T_u-2)^f$. Dividing by
$(T_u-1)^f\ne0$ gives the first condition.

\emph{Slope.} The second condition is the slope line.

\emph{Maximal slope.} If $\dim M_u=0$, then $V_{u,0}=0$, since
$\operatorname{emb}_0$ carries a basis of $V_{u,0}$ to a basis of $M_u$, and
$\Lambda^{\mathrm{ram}}_u=\Lambda_{u,0}$. Otherwise $T_u-2>0$ by the setting,
so $\dim M_{u,\mathrm{ram}}>0$ by the first condition. The map $\nu_u$ carries
$V(N_u)$ isomorphically onto $V^{\mathrm{ram}}_{u,0}$, and $j_N$ carries a basis
of $V(N_u)$ to a basis of $\mathcal M_u\cong M_{u,\mathrm{ram}}$
(Definition~\ref{pd-period-data}(7)), so $V^{\mathrm{ram}}_{u,0}\ne0$. A nonzero
vector of $V^{\mathrm{ram}}_{u,0}$ has a nonzero multiple in $\Lambda_{u,0}$,
which spans $V_{u,0}$, so $\Lambda^{\mathrm{ram}}_u\ne0$. Every nonzero
sublattice of $\Lambda^{\mathrm{ram}}_u$ is a nonzero sublattice of
$\Lambda_{u,0}$ with the same metrics, so
$\mu_{\max}(\Lambda^{\mathrm{ram}}_u)\le\mu_{\max}(\Lambda_{u,0})$ by
Definition~\ref{hs-slopes}(2).
\end{proof}

\begin{lemma}\label{asm-level-congruence-dvr}
Let $\mathcal O$ be a discrete valuation ring with uniformizer $\varpi$, $V$ an
$\mathcal O$-module, and $P,S,B\subset V$ submodules such that $S$ is free of
finite rank $r$ and $\varpi y\in P$ implies $y\in P$ for $y\in V$. Let
$U,D\colon V\to V$ be $\mathcal O$-linear and $e,b\ge0$ integers with
$\varpi^bS\subset B$, $U(B)\subset S$, $D(S)\subset B$, $U(P)\subset P$,
$D(P)\subset P$, and $U(D(s))=\varpi^es$ for $s\in S\cap P$. Then
$\varpi^bU(S\cap P)\subset S\cap P$ and, with lengths in
$\{0,1,2,\dots,\infty\}$,
\[
 \operatorname{len}_{\mathcal O}\big((S\cap P)/\varpi^bU(S\cap P)\big)\le
 \operatorname{len}_{\mathcal O}\big(B/\varpi^bS\big)+er.
\]
\end{lemma}
\leanref{Formalism/GL3MeanSquares/HookFormulaTorsion/LemmaCTwistRamified.lean}{Formalism.HermitianLattices.RamifiedHeckeModels.lemma_c_twist_ramified}
\leanref{Formalism/GL3MeanSquares/HookFormulaTorsion/LengthQuotientBoundGeneral.lean}{Formalism.HermitianLattices.RamifiedHeckeModels.lemma_c_under_twist}
\leanref{Formalism/GL3MeanSquares/HookFormulaTorsion/LengthQuotientBoundGeneral.lean}{Formalism.HermitianLattices.RamifiedHeckeModels.length_quotient_bound_general}
\begin{proof}
Put $S'=S\cap P$ and $B'=B\cap P$. The hypotheses give $\varpi^bS'\subset B'$,
$U(B')\subset S'$ and $D(S')\subset B'$. So
$\varpi^bU(S')=U(\varpi^bS')\subset U(B')\subset S'$, and
\[
 \operatorname{len}\big(S'/\varpi^bU(S')\big)=\operatorname{len}\big(U(B')/U(\varpi^bS')\big)
 +\operatorname{len}\big(S'/U(B')\big).
\]
The first quotient is the image of $B'/\varpi^bS'$ under $U$, so its length is
at most that of $B'/\varpi^bS'$. For $s\in S'$ we have
$\varpi^es=U(D(s))\in U(B')$, so $S'/U(B')$ is a quotient of $S'/\varpi^eS'$.
As a submodule of a free module of rank $r$ over a principal ideal domain, $S'$
is free of some rank $r'\le r$, and $S'/\varpi^eS'$ has length $er'\le er$.
Finally, the map $B'/\varpi^bS'\to B/\varpi^bS$ is injective: if $y\in B'$
equals $\varpi^bs$ with $s\in S$, then $\varpi^bs\in P$, so $s\in P$ by the
hypothesis on $P$, and $y\in\varpi^bS'$. Hence
$\operatorname{len}(B'/\varpi^bS')\le\operatorname{len}(B/\varpi^bS)$.
\end{proof}

\begin{lemma}\label{asm-level-congruence}
Let $D$ be a definite-side datum and $u$ a word, with the notation of
Definitions~\ref{asm-level-setting} and~\ref{asm-level-local-facts}. Suppose
that the setting of the level primes holds at $u$, and that Hecke integrality
at level zero, the ramified level zero and the second branch inclusion hold at
$u$. Then the congruence length bound holds at $u$.
\end{lemma}
\leanref{Formalism/LocalToGlobal/GlobalTheory/GL6/CongruenceNoKernelOfHeckeRamified.lean}{Formalism.L2.GlobalTheoryTheory.LocalLatticeBases.congruence_no_kernel_minimal}
\leanref{Formalism/LocalToGlobal/GlobalTheory/GL6/CongruenceNoKernelOfHeckeRamified.lean}{Formalism.L2.GlobalTheoryTheory.LocalLatticeBases.rational_level_zero_of_integral}
\leanref{Formalism/LocalToGlobal/GlobalTheory/GL6/CongruenceNoKernelOfHeckeRamified.lean}{Formalism.L2.GlobalTheoryTheory.LocalLatticeBases.hecke_ramified_lattice_stable_length}
\leanref{Formalism/SharpRestriction/GlobalSchurBounds/FiniteModelOfAdmissibleTransports.lean}{Formalism.MainNotions.product_first_zeroth_parts.LemmaC_core}
\leanref{Formalism/SharpRestriction/GlobalSchurBounds/FiniteModelOfAdmissibleTransports.lean}{Formalism.MainNotions.product_first_zeroth_parts.twist_le}
\leanref{Formalism/SharpRestriction/GlobalSchurBounds/FiniteModelOfAdmissibleTransports.lean}{Formalism.MainNotions.product_first_zeroth_parts.Nlam_eq}
\leanref{Formalism/SharpRestriction/GlobalSchurBounds/FiniteModelOfAdmissibleTransports.lean}{Formalism.MainNotions.product_first_zeroth_parts.len_list_le}
\leanref{Formalism/SharpRestriction/GlobalSchurBounds/FiniteModelOfAdmissibleTransports.lean}{Formalism.MainNotions.product_first_zeroth_parts.len_pow_le}
\begin{proof}
Let $w$ be a place of $L$ over $T_u$. If $w$ does not see $\pi_u$, then
$V_{\mathrm{top}}(w)$ is empty, so $U_w$ is the identity and $e_T(u,w)=0$. Then
$N_\lambda(u,w)=S^{\mathrm{ram}}_w$, the quotient is $0$, and the bound holds.
Let $w$ see $\pi_u$.

\emph{Step 1: the operators over $L$.} Let $v\in V_u$. The lattice $B_w(v)$
spans $V_{u,0}$ over $L$, the map $\operatorname{emb}_0$ is injective and
$\iota$-semilinear, and $U_v$ is $\mathbf C$-linear. So Hecke integrality gives
an $L$-linear map of $V_{u,0}$, again written $U_v$, with
$\operatorname{emb}_0\circ U_v=U_v\circ\operatorname{emb}_0$, and it maps
$B_w(v)$ into $S_w$. In the same way, since $S_w$ spans $V_{u,0}$, the down
operator gives an $L$-linear map $D_v$ of $V_{u,0}$ that maps $S_w$ into
$B_w(v)$. By the ramified level zero, both preserve $V^{\mathrm{ram}}_{u,0}$,
and $U_vD_vy=T_u^{k-1}y$ for $y\in V^{\mathrm{ram}}_{u,0}$.

\emph{Step 2: one place of $V_{\mathrm{top}}(w)$.} Let $v\in
V_{\mathrm{top}}(w)$, so that $\tau=\tau_w(v)\in\mathcal T_{\mathrm{top}}$, and
put $b_v=b_{\mathrm{top}}(\tau)$, which is at least $0$ by the setting, and
$g_v=T_u^{b_v}U_v$. Let $\varpi$ be a uniformizer of $\mathcal O_{L,w}$; then
$T_u=\epsilon\varpi^{e_w}$ with $\epsilon\in\mathcal O_{L,w}^\times$. Apply
Lemma~\ref{asm-level-congruence-dvr} with $\mathcal O=\mathcal O_{L,w}$,
$V=V_{u,0}$, $P=V^{\mathrm{ram}}_{u,0}$, $S=S_w$, $B=B_w(v)$, $U=U_v$,
$D=\epsilon^{1-k}D_v$, $e=e_w(k-1)$ and $b=e_wb_v$. Its hypotheses hold: $S_w$
is finitely generated and torsion free over $\mathcal O_{L,w}$, hence free, of
rank at most $\dim V_{u,0}=\dim M_u$; $P$ is an $L$-subspace;
$\varpi^bS_w=T_u^{b_v}S_w\subset B_w(v)$ by the second branch inclusion; and
Step 1 gives the rest. Since $\varpi^bU_v=\epsilon^{-b_v}g_v$, the lemma and
the length in the second branch inclusion give
$g_v(S^{\mathrm{ram}}_w)\subset S^{\mathrm{ram}}_w$ and
\[
 \operatorname{len}\big(S^{\mathrm{ram}}_w/g_vS^{\mathrm{ram}}_w\big)\le
 e_w\dim M_u\,\frac{\sum_{a=1}^k\lambda_{u,0}(\tau)_a+kb_v}k+e_w(k-1)\dim M_u.
\]

\emph{Step 3: the numerator.} Let $\tau\in\mathcal T_{\mathrm{top}}$. Then
$\lambda_{u,0}(\tau)_1\le a_{\max}(u)$, and $|\lambda_{u,0}(\tau)_a|\le
W_\tau(u)$ for $2\le a\le k$ by the definition of $W_\tau(u)$ with $i=0$; in
particular $b_{\mathrm{top}}(\tau)=-\lambda_{u,0}(\tau)_k\le W_\tau(u)$. Hence
\[
 \frac{\sum_{a=1}^k\lambda_{u,0}(\tau)_a+kb_{\mathrm{top}}(\tau)}k\le
 \frac{a_{\max}(u)+(k-1)W_\tau(u)+kW_\tau(u)}k\le\frac{a_{\max}(u)}k+2W_\tau(u).
\]
With $k-1\le k$ and $b_v\ge0$, the bound of Step 2 is at most
$e_w\dim M_u\big(a_{\max}(u)/k+2W_\tau(u)+k+b_v\big)$.

\emph{Step 4: the product.} Let $G$ be the composite of the $g_v$,
$v\in V_{\mathrm{top}}(w)$, in the order that defines $U_w$. Then
$G=T_u^{\sum_vb_v}U_w$, and $G^m=T_u^{e_T(u,w)}U_w^m$ because
$e_T(u,w)=m\sum_vb_v$. So $N_\lambda(u,w)=G^m(S^{\mathrm{ram}}_w)$ by
Definition~\ref{pd-hecke-at-T}(3). If $g$ and $h$ are endomorphisms of a module
that map a submodule $A$ into itself, then so does $gh$, and
$\operatorname{len}(A/ghA)\le\operatorname{len}(A/gA)+\operatorname{len}(A/hA)$,
because $gA/ghA$ is the image of $A/hA$ under $g$. Applying this repeatedly to
$A=S^{\mathrm{ram}}_w$ gives $N_\lambda(u,w)\subset S^{\mathrm{ram}}_w$ and
\[
 \operatorname{len}\big(S^{\mathrm{ram}}_w/N_\lambda(u,w)\big)\le
 m\sum_{v\in V_{\mathrm{top}}(w)}\operatorname{len}\big(S^{\mathrm{ram}}_w/g_vS^{\mathrm{ram}}_w\big).
\]
With Steps 2 and 3 this is the congruence length bound at $w$.
\end{proof}

\begin{lemma}\label{asm-level-commensurable}
Let $D$ be a definite-side datum with size data (Definition~\ref{pd-size-data})
and $u$ a word, with the notation of Definition~\ref{asm-level-setting}; let
$N_u\subset V(N_u)$, $\nu_u$ and $M^{\mathrm{ram}}_u=\nu_u^{-1}(\Lambda_{u,0})$
be as in Definition~\ref{pd-period-data}(7), and $N^\circ_u\subset V(N_u)$ the
submodule of the size data. For $\mathcal O_L$-submodules $A,B$ of $V(N_u)$ and a
finite place $w$ of $L$, say that $A\subset B$ \emph{locally at} $w$ if for
every $y\in A$ there is $\sigma\in\mathcal O_L\smallsetminus\mathfrak p_w$ with
$\sigma y\in B$. Suppose that:
\begin{enumerate}
\item for every finite place $w$ not over $T_u$, $N^\circ_u\subset
  M^{\mathrm{ram}}_u$ locally at $w$, and also $M^{\mathrm{ram}}_u\subset
  N^\circ_u$ locally at $w$ when $p_w$ is dropped or not inert for $D$;
\item for every place $w$ over $T_u$, a vector $y\in V(N_u)$ has $\sigma y\in
  N^\circ_u$ for some $\sigma\in\mathcal O_L\smallsetminus\mathfrak p_w$
  exactly when $\nu_u(y)\in N_\lambda(u,w)$;
\item $N_\lambda(u,w)\subset S^{\mathrm{ram}}_w$ for every place $w$ over $T_u$;
\item $N_u$ is the $\mathcal O_L$-span of $I_uN^\circ_u$ for a nonzero fractional
  ideal $I_u$ of $L$.
\end{enumerate}
Then $N^\circ_u\subset M^{\mathrm{ram}}_u$, $N^\circ_u$ is finitely generated,
and $cM^{\mathrm{ram}}_u\subset N^\circ_u$ for some nonzero $c\in\mathcal O_L$.
Moreover $N_u$ is finitely generated and spans $V(N_u)$.
\end{lemma}
\leanref{Formalism/SharpRestriction/GlobalSchurBounds/LinkPk.lean}{Formalism.MainNotions.Lnk.nuCommens_of}
\leanref{Formalism/SharpRestriction/GlobalSchurBounds/LinkPk.lean}{Formalism.MainNotions.Lnk.nuAtT_of}
\leanref{Formalism/SharpRestriction/GlobalSchurBounds/LinkPk.lean}{Formalism.MainNotions.Lnk.nu_le_mram}
\leanref{Formalism/SharpRestriction/GlobalSchurBounds/LinkPk.lean}{Formalism.MainNotions.Lnk.mem_of_locally}
\leanref{Formalism/SharpRestriction/GlobalSchurBounds/LinkPk.lean}{Formalism.MainNotions.Lnk.exists_free_prime}
\leanref{Formalism/SharpRestriction/GlobalSchurBounds/LinkPk.lean}{Formalism.MainNotions.Lnk.exists_common_mul}
\leanref{Formalism/SharpRestriction/GlobalSchurBounds/TopDensityNGeSixteen.lean}{Formalism.MainNotions.Lnk.NuCommens}
\leanref{Formalism/SharpRestriction/GlobalSchurBounds/TopDensityNGeSixteen.lean}{Formalism.MainNotions.Lnk.nuFull_of}
\begin{proof}
\emph{The lattice $M^{\mathrm{ram}}_u$.} The map $\nu_u$ is injective and carries
$M^{\mathrm{ram}}_u$ into the finitely generated module $\Lambda_{u,0}$, so
$M^{\mathrm{ram}}_u$ is finitely generated. It spans $V(N_u)$: for $y\in
V(N_u)$, the vector $\nu_u(y)$ lies in $V_{u,0}$, which $\Lambda_{u,0}$ spans,
so $d\nu_u(y)\in\Lambda_{u,0}$ for some nonzero $d\in\mathcal O_L$, and $dy\in
M^{\mathrm{ram}}_u$.

\emph{Places over $T_u$.} Let $w$ lie over $T_u$ and $y\in N^\circ_u$. By (2)
with $\sigma=1$, $\nu_u(y)\in N_\lambda(u,w)$, which lies in $S_w$ by (3). So
$\nu_u(y)$ is a combination of elements of $\Lambda_{u,0}$ with coefficients in
$\mathcal O_{L,w}$, and clearing denominators gives $\sigma\in\mathcal
O_L\smallsetminus\mathfrak p_w$ with $\sigma\nu_u(y)\in\Lambda_{u,0}$, that is
$\sigma y\in M^{\mathrm{ram}}_u$.

\emph{Inclusion.} By (1) and the previous step, every $y\in N^\circ_u$ lies in
$M^{\mathrm{ram}}_u$ locally at every finite place $w$. The ideal of the
$\sigma\in\mathcal O_L$ with $\sigma y\in M^{\mathrm{ram}}_u$ is therefore
contained in no maximal ideal, so it contains $1$ and $y\in M^{\mathrm{ram}}_u$.
As $\mathcal O_L$ is noetherian and $M^{\mathrm{ram}}_u$ is finitely generated,
so is $N^\circ_u$.

\emph{The multiplier.} The auxiliary prime $q$ is not inert for $D$, by
Definition~\ref{pd-size-data}, and neither is $2$, which divides $Q_S$
(Definition~\ref{hf-framed-host-field}). As $q\ne2$, one of them, say $p$,
differs from $T_u$; let $w$ be a place of $L$ over $p$. By (1), every $y\in
M^{\mathrm{ram}}_u$ has a nonzero $\sigma_y\in\mathcal O_L$ with $\sigma_yy\in
N^\circ_u$. If $y_1,\dots,y_s$ generate $M^{\mathrm{ram}}_u$, then
$c=\prod_i\sigma_{y_i}$ is nonzero and $cM^{\mathrm{ram}}_u\subset N^\circ_u$.

\emph{The lattice $N_u$.} There is a nonzero $d\in\mathcal O_L$ with
$dI_u\subset\mathcal O_L$. Then $dN_u\subset N^\circ_u$ by (4), and
multiplication by $d$ is injective on $V(N_u)$, so $N_u$ is isomorphic to a
submodule of the finitely generated module $N^\circ_u$ and is finitely
generated. Let $a_0\in I_u\cap\mathcal O_L$ be nonzero. For $y\in
M^{\mathrm{ram}}_u$ we have $cy\in N^\circ_u$, so $a_0cy\in N_u$ by (4), and
$y=(a_0c)^{-1}(a_0cy)$ lies in the $L$-span of $N_u$. As $M^{\mathrm{ram}}_u$
spans $V(N_u)$, so does $N_u$.
\end{proof}

\subsection{The count and the slope bounds}\label{subsection-asm-count-slope-bounds}

\begin{lemma}\label{asm-global-second-layer}
Let $\mathcal F$ be a family with exponent $\varepsilon$ (Definition~\ref{pd-family}) that
satisfies the \emph{second-layer conditions at the real place}: there is $x_a$ such that every
member $D$ at every scale $x\ge x_a$, with coefficient field $L$, has its theta realization tied at
the real place and has tied columns, as in the proof of
Lemma~\ref{asm-proviso-bottom-dimension}, whose characters in degree $0$ have finite order with
values in $L$. That is, for every word $u$ the character $\chi$ of the level $U$ of the forms of
$u$ in degree $0$ given by the tied columns satisfies $\chi(y)^n=1$ for some $n\ge1$ and all
$y\in U$, and every value $\chi(y)$ is the image of an element of $L$ under its embedding in
$\mathbf C$. Then there is $x_0>x_1$ such that every member $D$ at every scale $x\ge x_0$ has the
following properties.
\begin{enumerate}
\item The volume tie holds for $D$ (Definition~\ref{asm-facts-volume-tie}(1)).
\item $D$ has tied columns whose characters in degree $0$ have finite order with values in $L$.
\item For every word $u$ of $D$, every class $c\in\mathrm{cls}_0$ of the chain of $u$ has
  trivial stabilizer $\Gamma_c$.
\end{enumerate}
\end{lemma}
\leanref{Formalism/SharpRestriction/GlobalSchurBounds/LinkTied.lean}{Formalism.MainNotions.RealPlaceSecondLayerConditions}
\leanref{Formalism/SharpRestriction/GlobalSchurBounds/LinkTied.lean}{Formalism.MainNotions.family_conditions_of_layer_two}
\leanref{Formalism/SharpRestriction/GlobalSchurBounds/MassRatioUnfoldedOfTheoremTwo.lean}{Formalism.MainNotions.RT.volTie_of_realTied}
\leanref{Formalism/SharpRestriction/GlobalSchurBounds/LinkTied.lean}{Formalism.MainNotions.stabF}
\leanref{Formalism/SharpRestriction/GlobalSchurBounds/JointsF.lean}{Formalism.MainNotions.flipTwoThree.stabiliser_family_of_definite}
\leanref{Formalism/SharpRestriction/GlobalSchurBounds/JointsL.lean}{Formalism.MainNotions.GlobalTheoryTheory.JDB.gBig_of}
\begin{proof}
Let $x_a$ be as in the hypothesis and put $x_0=\max(x_a,3,x_1+1)$. Let $D$ be a member at a
scale $x\ge x_0$.

(1) The tie at the real place contains the volume tie.

(2) This is the hypothesis.

(3) Since $x\ge3$ and $x>x_1$, the sizes of the family give $(kb)^8\le g_0$, and as in Step~2 of
the proof of Proposition~\ref{asm-proviso-reduction} this gives $g_0>2fk+1$.
Lemma~\ref{asm-proviso-trivial-stabilizers}, applied with the level data among the further choices
of the family (Definition~\ref{pd-family}(4)), shows that $\Gamma_c=1$ for every word $u$ and
every $c\in\mathrm{cls}_0$.
\end{proof}

\begin{definition}\label{asm-global-column-values}
Let $D$ be a definite-side datum (Definition~\ref{th-definite-datum}) with rank $k=N(0)$, host
field $E$, set $\mathcal T$ of $f$ places and coefficient field $L$, and let $u$ be a word, with
auxiliary prime $T_u$, bottom space $X_0$ of its chain with Hermitian form $h$, lattices
$L_u(0,c)\subset X_0$ for $c\in\mathrm{cls}_0$ (Definition~\ref{th-definite-datum}(7)), and
lattice of forms $\Lambda_{u,0}\subset V_{u,0}$ of weights
$\lambda=(\lambda_{u,0}(\tau))_{\tau\in\mathcal T}$ (Definition~\ref{asm-level-setting}). For a
finite place $w$ of $L$ let $\mathcal O_{L,w}\subset L$ be its valuation ring, and for an
$\mathcal O_L$-module $M$ let $M_w$ be its $\mathcal O_{L,w}$-span. Let $L[x]$ be the polynomial
ring over $L$ in variables $x_{\tau,a,b}$ ($\tau\in\mathcal T$, $1\le a,b\le k$), put
$\nu_{\tau,i}=\lambda_{u,0}(\tau)_i-\lambda_{u,0}(\tau)_k$, and with the Weyl modules
$\Delta_R(\nu)$, the extremal vectors $h_\nu$, the Fock pairing and the action $\rho_\lambda$ of
Definition~\ref{pd-weight-lattices}, let $V_\lambda\subset L[x]$ be the $L$-span of the products
$\prod_\tau f_\tau(x_\tau)$ with $f_\tau\in\Delta_L(\nu_\tau)$, and $\Delta_{\lambda,w}$ the
$\mathcal O_{L,w}$-span of those with $f_\tau\in\Delta_{\mathbf Z}(\nu_\tau)$. For complex
polynomials $p,q$ in the $x_{\tau,a,b}$ put
$\{p,q\}_\lambda=\sum_\kappa\kappa!\,p_\kappa\bar q_\kappa\big/\prod_\tau[h_{\nu_\tau},h_{\nu_\tau}]$.
The \emph{columns of $u$ have values} if there are the following data.
\begin{enumerate}
\item The set $\mathrm{cls}_0$ is nonempty, and there are field embeddings
  $\rho_\tau\colon E\to L$ for $\tau\in\mathcal T$.
\item For $c\in\mathrm{cls}_0$ and $\tau\in\mathcal T$: a Hermitian lattice $Q(c,\tau)$ over $L$,
  finitely generated and spanning its space, with a basis $b_1,\dots,b_k$ of that space in which
  its metric at every infinite place $v$ of $L$ has an invertible matrix $G_{c,\tau,v}$; and an
  additive map $\beta_{c,\tau}\colon X_0\to V(Q(c,\tau))$ to the space of $Q(c,\tau)$, with
  $\beta_{c,\tau}(ay)=\rho_\tau(a)\beta_{c,\tau}(y)$, such that $Q(c,\tau)$ is the
  $\mathcal O_L$-span of $\beta_{c,\tau}(L_u(0,c))$ and
  $\|\beta_{c,\tau}(y)\|_v^2=\operatorname{Re}v(\rho_\tau(h(y,y)))$ for all $y\in X_0$ and all
  infinite places $v$.
\item For every $c\in\mathrm{cls}_0$ the lattice $L_u(0,c)$ lies in its dual
  $\{y\in X_0:h(L_u(0,c),y)\subset\mathcal O_E\}$, with index $T_u^{2f}$.
\item An injective $L$-linear map $\operatorname{co}$ from $V_{u,0}$ to the functions
  $\mathrm{cls}_0\to L[x]$, whose image consists of the functions with values in $V_\lambda$.
\item For every finite place $w$ of $L$ and all matrices $\gamma_{c,\tau}\in GL_k(L)$ such that
  $Q(c,\tau)_w$ is the $\mathcal O_{L,w}$-span of the vectors $\sum_i(\gamma_{c,\tau})_{ia}b_i$
  ($1\le a\le k$): an element $y\in V_{u,0}$ lies in $(\Lambda_{u,0})_w$ if and only if
  $\rho_\lambda\big((\gamma_{c,\tau}^{-1})_\tau\big)\operatorname{co}(y)(c)\in\Delta_{\lambda,w}$
  for every $c\in\mathrm{cls}_0$.
\item For every infinite place $v$ of $L$ and every $y\in V_{u,0}$,
  \[
   \|y\|_v^2=\sum_{c\in\mathrm{cls}_0}\operatorname{Re}\{P_c,P_c\}_\lambda,\qquad
   P_c=\rho_\lambda\big((G_{c,\tau,v}^{-1})_\tau\big)\,v\big(\operatorname{co}(y)(c)\big),
  \]
  where $\|y\|_v$ is the norm of the Hermitian lattice $\Lambda_{u,0}$ and $v(\cdot)$ applies the
  embedding $v$ to the coefficients.
\end{enumerate}
\end{definition}
\leanref{Formalism/SharpRestriction/GlobalSchurBounds/FlipTwoThree.lean}{Formalism.MainNotions.flipTwoThree.ColVal}
\leanref{Formalism/SharpRestriction/GlobalSchurBounds/FlipTwoThree.lean}{Formalism.MainNotions.flipTwoThree.LatT}
\leanref{Formalism/SharpRestriction/GlobalSchurBounds/FlipTwoThree.lean}{Formalism.MainNotions.flipTwoThree.TSpan}
\leanref{Formalism/SharpRestriction/GlobalSchurBounds/FlipTwoThree.lean}{Formalism.MainNotions.flipTwoThree.IntSpanAt}
\leanref{Formalism/SharpRestriction/GlobalSchurBounds/FlipTwoThree.lean}{Formalism.MainNotions.flipTwoThree.LocBasis}
\leanref{Formalism/SharpRestriction/GlobalSchurBounds/FlipTwoThree.lean}{Formalism.MainNotions.flipTwoThree.actAll}
\leanref{Formalism/SharpRestriction/GlobalSchurBounds/FlipTwoThree.lean}{Formalism.MainNotions.flipTwoThree.nfAll}

\begin{lemma}\label{asm-global-count}
Let $\mathcal F$ be a family with exponent $\varepsilon$ (Definition~\ref{pd-family}), and let
$C$ and $x_0$ be real. Suppose that for every $x\ge x_0$ with $x>x_1$, every member $D$ at $x$,
with $L=x^{\varepsilon}$, the number $\Gamma_P$ and the numbers $X_u$ of
Definition~\ref{th-definite-datum}(3), the Harish-Chandra weights $w_j(\tau)$, the half root
discriminant $v$ and the level primes $g_0,g_1$ (Definition~\ref{asm-member-bounds}), and all words
$u,u'$ of $D$:
\begin{enumerate}
\item $\log T_u\le CmL$;
\item $\displaystyle\frac1f\sum_{j=1}^m\sum_{\tau\in\mathcal T}\log\frac{w_j(\tau)v}{\pi}
  +m\log(g_0g_1)-C\Big(L+\frac{m^2L}k\Big)-\frac{2\log\Gamma_P}{fk}\le\log T_u$;
\item $|\log X_u-\log X_{u'}|\le C\log k$.
\end{enumerate}
Then every member at every scale $x\ge\max(x_0,x_1+1)$ satisfies the count with
$\max(C,0)+2$ (Definition~\ref{asm-member-bounds}(1)). In particular $\mathcal F$ satisfies the
count at all large scales.
\end{lemma}
\leanref{Formalism/SharpRestriction/GlobalSchurBounds/JointsL.lean}{Formalism.MainNotions.GlobalTheoryTheory.count_at_of}
\leanref{Formalism/SharpRestriction/GlobalSchurBounds/JointsL.lean}{Formalism.MainNotions.GlobalTheoryTheory.countF_of}
\leanref{Formalism/SharpRestriction/GlobalSchurBounds/JointsL.lean}{Formalism.MainNotions.GlobalTheoryTheory.logGam_le}
\leanref{Formalism/SharpRestriction/GlobalSchurBounds/JointsL.lean}{Formalism.MainNotions.GlobalTheoryTheory.JDB.db_pi_le_v}
\begin{proof}
Let $D$ be a member at a scale $x\ge\max(x_0,x_1+1)$, let $u,u'$ be words, and put
$C'=\max(C,0)+2$.

\emph{Preliminaries.} The rank datum has $k\ge m+3\ge5$ and $f\ge1$, so $fk\ge1$ and
$\log k\ge1$. As $x>1$, $L\ge1$. The half root discriminant satisfies $v\ge\pi$, so the root
discriminant is $r_E=2v\ge1$, and $L_a=\log(ar_E)=\log(xe^{L}r_E)\ge L$. By
Definition~\ref{pd-size-data}(9), $\log\Gamma_P=C_\Gamma m^5fL$, and the sizes of the family
give $8C_\Gamma m^5fL\le k$; hence $\log\Gamma_P\le k/8\le fk$ and
$2\log\Gamma_P/(fk)\le2$.

\emph{The window of $T_u$.} Put $T^*_u=X_u\Gamma_P^{-2/(fk)}$. By
Definition~\ref{th-definite-datum}(4), $T^*_u\le T_u\le(1+1/(fk))T^*_u$, and
$\log(1+1/(fk))\le1/(fk)\le1$. So
\[
 0\le\log T_u-\Big(\log X_u-\frac{2\log\Gamma_P}{fk}\Big)\le1,
\]
and the same holds for $u'$, with the same $\Gamma_P$, which does not depend on the word.

\emph{Lower bound.} Since $L\le L_a$, $C(L+m^2L/k)\le\max(C,0)(L_a+m^2L_a/k)$. Since
$1\le L\le L_a$, $2\log\Gamma_P/(fk)\le2\le2(L_a+m^2L_a/k)$. With (2) this gives
\[
 \frac1f\sum_{j=1}^m\sum_\tau\log\frac{w_j(\tau)v}{\pi}+m\log(g_0g_1)
 -C'\Big(L_a+\frac{m^2L_a}k\Big)\le\log T_u.
\]

\emph{Upper bound.} By (1), $\log T_u\le CmL\le C'mL$.

\emph{Variation.} By the window and (3),
$|\log T_u-\log T_{u'}|\le|\log X_u-\log X_{u'}|+1\le\max(C,0)\log k+2\log k=C'\log k$, using
$\log k\ge1$.

These are the three conditions of the count with $C'$.
\end{proof}

\begin{definition}\label{asm-global-numerical}
Let $D$ be a definite-side datum (Definition~\ref{th-definite-datum}) with ranks $k,m,t$, $f=2n$
and $N(m)=k+m(t+1)$, host field with root discriminant $r_E$ and $v=r_E/2$, auxiliary prime $d$
with quadratic character $\chi_d$, level primes $g_0,g_1$ and top classes $T$, and let $u$ be a
word, with auxiliary prime $T_u$ and parameters $\lambda_u(\tau)$ with entries $u_i(\tau)$ and
$e_j(\tau)$. For $\tau\in\mathcal T$ put $\lambda_i(\tau)=\lambda_{u,0}(\tau)_i=u_i(\tau)-\frac{k+1}2+i$
for $1\le i\le k$ (Definition~\ref{th-definite-datum}(7)) and
$\omega_j(\tau)=e_j(\tau)-\frac k2-\frac{(2j-1)(t+1)}2$ for $1\le j\le m$. The place $\tau$ is
\emph{top} if it is the place $\tau_1(t)$ of a top class $t$. Let $C_n$, $x$, $b$, $L_0$, $b'$,
$S_c$ and $\eta$ be real numbers.
\begin{enumerate}
\item The \emph{numerical conditions} with $C_n$ hold at $u$ for $(x,b,L_0,b',S_c,\eta)$ if
  $x\ge2$, $b>0$, $v\ge1$, $b'\le b$, $S_c\ge1$, $\eta\ge1$,
  \[
   \log(xbv)\le L_0,\quad \log T_u\le C_nmL_0,\quad \log g_0\le C_nL_0,\quad
   \log g_1\le C_nL_0,\quad \Big|\operatorname{Re}\frac{L'}{L}(1,\chi_d)\Big|\le C_nL_0,
  \]
  \[
   S_c+\eta+mt\le\frac{b'}x,\qquad N(m)^2(L_0+\log T_u)\le N(m)b',\qquad
   N(m)L_0+\log T_u\le b',
  \]
  \[
   m^2N(m)^2\big(L_0+\log T_u+\log(g_0g_1)\big)\le N(m)b'L_0,\qquad(9+4C_n)m^3N(m)\le b'.
  \]
\item The \emph{weight conditions} with $C_n$ hold at $u$ for $(x,b,b',S_c,\eta)$ if for every
  $\tau\in\mathcal T$: every $\lambda_i(\tau)$ is an integer; if $\tau$ is top, then
  $xb\le\lambda_1(\tau)\le2xb$, $b'\le-\lambda_k(\tau)<2b'$, and $0\le\lambda_i(\tau)\le S_c$ for
  $1<i<k$; if $\tau$ is not top, then $0\le\lambda_i(\tau)\le C_nS_c$ for $1\le i\le k$;
  $\omega_1(\tau)\le\omega_2(\tau)\le\cdots\le\omega_m(\tau)$; if $\tau$ is not a thick place,
  then $\omega_j(\tau)\le2b'$ for all $j$; if $\tau$ is a thick place, then
  $b-N(m)\le\omega_j(\tau)\le b+\eta+1$ for all $j$; $-\lambda_k(\tau)\le\omega_1(\tau)$; and
  $\omega_m(\tau)-\omega_1(\tau)\le\eta$.
\end{enumerate}
A family $\mathcal F$ with exponent $\varepsilon$ (Definition~\ref{pd-family}) satisfies the
\emph{column slope bounds} if for every real $C_n$ there are $C_5$, $C_6$ and $x_0$ such that the
following holds for every $x\ge x_0$ with $x>x_1$, every member $D$ at $x$, every word $u$ of $D$
whose columns have values (Definition~\ref{asm-global-column-values}), and all real
$b,L_0,b',S_c,\eta$ such that the numerical conditions and the weight conditions with $C_n$ hold at
$u$ for $(x,b,L_0,b',S_c,\eta)$: the slopes $\mu(M_u)$ and $\mu_{\max}(M_u)$ of the first
lattice $\Lambda_{u,0}$ of the chain of $u$ (Definition~\ref{asm-tied}(3)) satisfy
\[
 -\sum_{\tau\ \mathrm{top}}\frac{\lambda_1(\tau)}k\log T_u-C_5|T|xb\log k
 -C_5fk(b'+S_c)\log k-C_5fS_c\log T_u-C_5fk^2L_0\le\mu(M_u)
\]
and
\[
 \mu_{\max}(M_u)\le C_6|T|xb\log(kv)+C_6f(b'+S_c)\big(k\log(kv)+\log T_u\big)+C_6fk^2L_0.
\]
\end{definition}
\leanref{Formalism/SharpRestriction/GlobalSchurBounds/FlipTwoThree.lean}{Formalism.MainNotions.flipTwoThree.Nums}
\leanref{Formalism/SharpRestriction/GlobalSchurBounds/FlipTwoThree.lean}{Formalism.MainNotions.flipTwoThree.Wts}
\leanref{Formalism/SharpRestriction/GlobalSchurBounds/GlobalRankLatticeProofs/SevenLogLeLOfPowerBound.lean}{Formalism.MainNotions.flipTwoThree.lamCol_eq}
\leanref{Formalism/SharpRestriction/GlobalSchurBounds/GlobalRankLatticeProofs/SevenLogLeLOfPowerBound.lean}{Formalism.MainNotions.flipTwoThree.wJ_eq}
\leanref{Formalism/SharpRestriction/GlobalSchurBounds/FlipTwoThree.lean}{Formalism.MainNotions.flipTwoThree.Conclusion5}
\leanref{Formalism/SharpRestriction/GlobalSchurBounds/FlipTwoThree.lean}{Formalism.MainNotions.flipTwoThree.Conclusion6}
\leanref{Formalism/SharpRestriction/GlobalSchurBounds/FlipTwoThree.lean}{Formalism.MainNotions.flipTwoThree.ConclusionF5}
\leanref{Formalism/SharpRestriction/GlobalSchurBounds/FlipTwoThree.lean}{Formalism.MainNotions.flipTwoThree.ConclusionF6}

\begin{lemma}\label{asm-global-numerical-hold}
Let $\mathcal F$ be a family with exponent $\varepsilon$ (Definition~\ref{pd-family}), and let
$C_c$ be real. Suppose that for all large $x$ every member at $x$ satisfies
$\log T_u\le C_cmx^{\varepsilon}$ for every word $u$ and
$|\operatorname{Re}\frac{L'}{L}(1,\chi_d)|\le x^{\varepsilon}$. Put
$C_{n,0}=\max(3,C_c,C_{YL}+1)$. Then for every $C_n\ge C_{n,0}$ and every $c_L\ge3$ there is
$x_0$ such that, at every member $D$ at every scale $x\ge x_0$, with the quantities $b$ and $b'$,
the cloud size $S_c$ and the width $\eta$ of its size data (Definition~\ref{pd-size-data}), the
numerical conditions and the weight conditions with $C_n$ (Definition~\ref{asm-global-numerical})
hold at every word for $(x,b,c_Lx^{\varepsilon},b',S_c,\eta)$.
\end{lemma}
\leanref{Formalism/SharpRestriction/GlobalSchurBounds/GlobalRankLatticeProofs/SevenLogLeLOfPowerBound.lean}{Formalism.MainNotions.flipTwoThree.nums_wts_of_fam}

This is a computation from the parameter choices and the sizes of the family, at the scales $x$
with $x\ge4$, $x>x_1$, $x^{14}\le X'$ and $m\ge c_L+\max(C_c,0)+6C_n+9$; all large $x$ have these
properties.

\begin{lemma}\label{asm-global-slopes-m}
Let $\mathcal F$ be a family with exponent $\varepsilon$ (Definition~\ref{pd-family}). Suppose
that $\mathcal F$ satisfies the column slope bounds (Definition~\ref{asm-global-numerical}) and the
count at all large scales (Definition~\ref{asm-member-bounds}(1)), and that for all large $x$, at
every member at $x$, $|\operatorname{Re}\frac{L'}{L}(1,\chi_d)|\le x^{\varepsilon}$ and the columns
of every word have values (Definition~\ref{asm-global-column-values}). Then $\mathcal F$
satisfies the slope bounds for $M$ at all large scales (Definition~\ref{asm-member-bounds}(2)).
\end{lemma}
\leanref{Formalism/SharpRestriction/GlobalSchurBounds/GlobalRankLatticeProofs/SlopesMFOfCountFam.lean}{Formalism.MainNotions.flipTwoThree.slopesMF_of_count_fam}
\leanref{Formalism/SharpRestriction/GlobalSchurBounds/GlobalRankLatticeProofs/SlopesMFOfCountFam.lean}{Formalism.MainNotions.flipTwoThree.numsL_of_fam}
\leanref{Formalism/SharpRestriction/GlobalSchurBounds/JointsL.lean}{Formalism.MainNotions.GlobalTheoryTheory.slopesMF_of_count}
\leanref{Formalism/SharpRestriction/GlobalSchurBounds/JointsL.lean}{Formalism.MainNotions.GlobalTheoryTheory.JSL.slopesMF_of}
\leanref{Formalism/SharpRestriction/GlobalSchurBounds/JointsL.lean}{Formalism.MainNotions.GlobalTheoryTheory.JSL.slopesM_at_of}
\leanref{Formalism/SharpRestriction/GlobalSchurBounds/JointsL.lean}{Formalism.MainNotions.GlobalTheoryTheory.JSL.jnt_sizes_slopes}
\leanref{Formalism/SharpRestriction/GlobalSchurBounds/JointsL.lean}{Formalism.MainNotions.GlobalTheoryTheory.JSL.jnt_slope_low}
\leanref{Formalism/SharpRestriction/GlobalSchurBounds/JointsL.lean}{Formalism.MainNotions.GlobalTheoryTheory.JSL.jnt_slope_up}
\begin{proof}
\emph{Constants.} The count gives $C\ge0$ with $\log T_u\le CmL$, $L=x^{\varepsilon}$, at every
word of every member at all large scales. With the bound on the logarithmic derivative,
Lemma~\ref{asm-global-numerical-hold} with $C_c=C$, $C_n=C_{n,0}$ and $c_L=3$ gives, at all
large scales, the numerical conditions and the weight conditions with $C_n$ at every word of every
member, for $(x,b,3L,b',S_c,\eta)$. The column slope bounds for this $C_n$ give constants $C_5$ and
$C_6$. Put $C_p=\max(C_5,C_6,0)$ and $C'=2C+C_p(15+C)\ge0$. Let $x$ be a scale beyond all these
thresholds, with $x\ge2$, $x>x_1$ and $x^{14}\le X'=\exp(x^{\varepsilon/3})$; all large $x$
qualify. Let $D$ be a member at $x$ and $u$ a word, whose columns have values.

\emph{Sizes.} The sizes of the family give $n\le x^3$, $k\le x^5$, $|T|\ge1$ and
$kS_c\le b'/(xn)$. Hence $\log k\le5\log x$ and $S_c\le kS_c\le b'$. Since $x^6\ge64$,
$2fk=4nk\le4x^8\le x^{14}\le X'$, and as $X'b'=a$,
\[
 fk(b'+S_c)\le2fkb'\le X'b'=a\le|T|a.
\]
Similarly $fk^2=2nk^2\le2x^{13}\le x^{14}\le X'\le b\le a\le|T|a$. As $m\le k$, also
$fm(b'+S_c)\le|T|a$.

\emph{Logarithms.} Since $a=xe^{L}$, $L_a=\log(ar_E)=\log x+L+\log r_E$. The numerical
conditions give $v\ge1$, so $r_E=2v\ge2$, and $\log x$, $L$, $\log r_E$ are nonnegative. Hence
$L\le L_a$, $\log k\le5\log x\le5L_a$, and $\log(kv)\le5\log x+\log r_E\le6L_a$. The count gives
$\log T_u\le CmL\le CmL_a$.

\emph{Lower bound.} By the weight conditions, $\lambda_1(\tau)\le2xb=2a$ at every top place, and
there are $|T|$ top places, one for each top class. So
\[
 \sum_{\tau\ \mathrm{top}}\frac{\lambda_1(\tau)}k\log T_u\le|T|\frac{2a}k\log T_u
 \le2C|T|aL_a\frac mk.
\]
Since $C_5\le C_p$, $C_p\ge0$ and each factor below is nonnegative, the size bounds give
\[
 C_5|T|a\log k\le5C_p|T|aL_a,\qquad C_5fk(b'+S_c)\log k\le5C_p|T|aL_a,
\]
\[
 C_5fS_c\log T_u\le C_pCfm(b'+S_c)L_a\le CC_p|T|aL_a,\qquad
 C_5fk^2\cdot3L\le3C_p|T|aL_a.
\]
The lower column slope bound, with $L_0=3L$, therefore gives
\[
 \mu(M_u)\ge-|T|aL_a\Big(2C\frac mk+C_p(13+C)\Big)\ge-C'|T|aL_a\Big(1+\frac mk\Big).
\]

\emph{Upper bound.} In the same way, $C_6|T|a\log(kv)\le6C_p|T|aL_a$,
\[
 C_6f(b'+S_c)\big(k\log(kv)+\log T_u\big)\le C_p\big(fk(b'+S_c)\cdot6L_a+fm(b'+S_c)CL_a\big)
 \le C_p(6+C)|T|aL_a,
\]
and $C_6fk^2\cdot3L\le3C_p|T|aL_a$. The upper column slope bound gives
$\mu_{\max}(M_u)\le C_p(15+C)|T|aL_a\le C'|T|aL_a$.

These are the slope bounds for $M$ with the constant $C'$, at every member at every scale beyond
the thresholds above.
\end{proof}

\begin{lemma}\label{asm-global-losses}
Let $\mathcal F$ be a family with exponent $\varepsilon$ (Definition~\ref{pd-family}). Suppose
that there are $C\ge0$ and $x_0$ such that every member at every scale $x\ge x_0$ whose weight
vectors have small shares (Definition~\ref{asm-duty}(5)) satisfies the loss bounds of
Definition~\ref{asm-member-bounds}(6) with $(C,0)$, and that for all large $x$ the weight vectors
of every member at $x$ have small shares. Then $\mathcal F$ satisfies the loss bounds with $e_L=0$
(Definition~\ref{asm-member-bounds}).
\end{lemma}
\leanref{Formalism/SharpRestriction/GlobalSchurBounds/LossesSmallShares.lean}{Formalism.MainNotions.Lnk.lossF_of_7}
\leanref{Formalism/SharpRestriction/GlobalSchurBounds/LossesSmallShares.lean}{Formalism.MainNotions.Lnk.LossesSmallShares}
\leanref{Formalism/SharpRestriction/GlobalSchurBounds/WeightVectorSmallSharesEight.lean}{Formalism.MainNotions.chosenWS_of_wall}
\leanref{Formalism/SharpRestriction/GlobalSchurBounds/WeightVectorSmallSharesSeven.lean}{Formalism.MainNotions.ChosenWS}
\leanref{Formalism/SharpRestriction/FamCore.lean}{Formalism.MainNotions.FamCore.coverU_hyp_self}
\begin{proof}
Let $x'_0$ be such that the weight vectors of every member at every $x\ge x'_0$ have small shares,
and let $x\ge\max(x_0,x'_0,x_1+1)$ be a scale. There is a member $i_0$ at $x$. Its exponents
$a_t,b_t$, indexed by the images of its top classes in $\mathbf Z$, lie in the window of
Definition~\ref{pd-size-data}(2) for the size data of $i_0$, and satisfy
$b'\le b_t\le b_{\mathrm{top}}(x)$ for the upper bound $b_{\mathrm{top}}$ of the family. So some
member at $x$, namely $i_0$, has the exponents $(a_t,b_t)$. Every member at $x$ has small shares,
so every member at $x$ with these exponents satisfies (6) with $(C,0)$. This is the loss bounds
with $e_L=0$, with the constant $C$ and the threshold $\max(x_0,x'_0,x_1+1)$.
\end{proof}

\begin{lemma}\label{asm-global-slope-index}
Let $K$ be a totally complex number field, $V$ a Hermitian space over $K$ and $M\subset V$ a
nonzero Hermitian lattice (Definition~\ref{hs-hermitian-lattice}), with the slopes of
Definition~\ref{hs-slopes}. For a finite place $v$ of $K$ with prime ideal $\mathfrak p_v$ put
$\ell_v=\log\mathrm N\mathfrak p_v/[K:\mathbf Q]$.
\begin{enumerate}
\item Let $N\subset M$ be a nonzero sublattice with $\operatorname{rk}N=\operatorname{rk}M$, and let
  $\operatorname{len}_v(M/N)$ be the length of the localization of $M/N$ at $\mathfrak p_v$. Then
  every $\operatorname{len}_v(M/N)$ is finite, all but finitely many are $0$, and
  \[
   \mu(N)=\mu(M)-\frac2{\operatorname{rk}M}\sum_v\ell_v\operatorname{len}_v(M/N).
  \]
\item Let $I$ be a nonzero fractional ideal of $K$ and $IM$ the $\mathcal O_K$-span of the products
  $ay$, $a\in I$, $y\in M$. Then
  \[
   \mu(IM)=\mu(M)-\frac{2\log\mathrm N(I)}{[K:\mathbf Q]},\qquad
   \mu_{\max}(IM)\le\mu_{\max}(M)-\frac{2\log\mathrm N(I)}{[K:\mathbf Q]}.
  \]
\item If $N\subset M$ is a nonzero sublattice, then $\mu_{\max}(N)\le\mu_{\max}(M)$.
\end{enumerate}
\end{lemma}
\leanref{Formalism/SizeWindowBrackets/ChainSlopeWindows/ExistsNsmulMem.lean}{Formalism.LFunctionChainParameters.slopeLatticeData.muDef_index}
\leanref{Formalism/SizeWindowBrackets/ChainSlopeWindows/ExistsNsmulMem.lean}{Formalism.LFunctionChainParameters.slopeLatticeData.index_finsupp}
\leanref{Formalism/LocalToGlobal/GlobalTheory/PlacesP.lean}{Formalism.L2.GlobalTheoryTheory.Places.mu_index_sum}
\leanref{Formalism/SizeWindowBrackets/ChainSlopeWindows/ExistsNsmulMem.lean}{Formalism.LFunctionChainParameters.slopeLatticeData.muDef_fractional}
\leanref{Formalism/SizeWindowBrackets/ChainSlopeWindows/ExistsNsmulMem.lean}{Formalism.LFunctionChainParameters.slopeLatticeData.muMaxDef_fractional_le}
\leanref{Formalism/SizeWindowBrackets/ChainSlopeWindows/ExistsNsmulMem.lean}{Formalism.LFunctionChainParameters.slopeLatticeData.muMaxDef_mono}
\begin{proof}
Put $n=\operatorname{rk}M$ and $r=[K:\mathbf Q]/2$.

(1) The ranks agree, so $M/N$ is a finitely generated torsion $\mathcal O_K$-module, hence finite.
It is the direct sum of its localizations, and a module of length $l$ over the localization at
$\mathfrak p_v$ has $\mathrm N\mathfrak p_v^{\,l}$ elements. So the lengths are finite, almost all
$0$, and $[M:N]=\prod_v\mathrm N\mathfrak p_v^{\operatorname{len}_v(M/N)}$. The covolume of $N$ in
$(KM)\otimes_{\mathbf Q}\mathbf R$ is $[M:N]$ times that of $M$, so
Definition~\ref{hs-slopes}(1) gives $\mu(N)=\mu(M)-\log[M:N]/(rn)$, which is the formula.

(2) The lattice $IM$ has rank $n$, and its covolume is $\mathrm N(I)^n$ times that of $M$: for
integral $I$, $IM\subset M$ has index $\mathrm N(I)^n$, since $M$ is locally free of rank $n$, and
multiplying a lattice of rank $n$ by a nonzero $c\in K$ multiplies its covolume by
$|\mathrm N_{K/\mathbf Q}(c)|^n$; write $I=c^{-1}J$ with $J$ integral. So
Definition~\ref{hs-slopes}(1) gives $\mu(IM)=\mu(M)-\log\mathrm N(I)/r$. A nonzero sublattice
$Q\subset IM$ is
$IQ'$ for the nonzero sublattice $Q'=I^{-1}Q$ of $M$, so
$\mu(Q)=\mu(Q')-2\log\mathrm N(I)/[K:\mathbf Q]\le\mu_{\max}(M)-2\log\mathrm N(I)/[K:\mathbf Q]$.

(3) A nonzero sublattice of $N$ is a nonzero sublattice of $M$, with the same norms.
\end{proof}

\begin{lemma}\label{asm-global-unit-weight}
Let $W$ be a finite set, $W_T\subset W$ and $J$ a finite set, and for $w\in W$ and $j\in J$ let
$l_w\ge0$, $\mathrm{gain}_j(w)$, $F_j(w)$, $E_j(w)$, $\tau_w$, $\iota_w$, $\operatorname{len}_w$ and
$s'_w$ be real numbers. Put
\[
 s_w=\sum_{j\in J}\big(\mathrm{gain}_j(w)+F_j(w)-E_j(w)\big),\qquad
 \Sigma'=\sum_{w\in W}2l_w(s'_w-\tau_w),\qquad I=\sum_{w\in W}l_w\iota_w,
\]
\[
 G=2\sum_{w\in W}l_w\Big(\sum_{j\in J}\mathrm{gain}_j(w)-\tau_w\Big),\qquad
 F=2\sum_{j\in J}\sum_{w\in W}l_wF_j(w),\qquad
 E=2\sum_{j\in J}\sum_{w\in W}l_wE_j(w).
\]
\begin{enumerate}
\item Let $\mu^+_M,\mu^+_R,\mu^+_N$ be real numbers with $s'_w\le s_w$ for $w\in W$,
  $\mu^+_N\le\mu^+_R+\Sigma'$ and $\mu^+_R\le\mu^+_M$. Then $\mu^+_N\le\mu^+_M+G+F-E$.
\item Let $d_M>0$, $0\le\delta<1$, $d_N=(1-\delta)d_M$, $f\ge0$, $T>2$ and
  $b,P,R,\mu_M,\mu_R,\mu_N$ be real numbers. Suppose that $\tau_w=0$ for $w\in W_T$; that
  $\tau_w\ge0$ and $\iota_w\ge0$ for $w\in W$; that $\operatorname{len}_w\le d_M(\tau_w+\iota_w)$
  for $w\in W\smallsetminus W_T$; that
  \[
   \sum_{w\in W_T}\frac{2l_w\operatorname{len}_w}{d_N}\le P,\qquad
   \sum_{w\in W}2l_w(s_w-s'_w)\le R;
  \]
  that
  \[
   \mu_N=\mu_R+\Sigma'-\sum_{w\in W}2l_w\frac{\operatorname{len}_w-d_N\tau_w}{d_N},\qquad
   \mu_R=\mu_M-\frac{f\log(T-1)}{T-2};
  \]
  and that $\delta\le2f/T\le1/2$, $2\sum_{w\in W}l_w\tau_w\le b^2$ and
  $4fb^2/T+f\log(T-1)/(T-2)\le1$. Then
  \[
   \mu_M+G+F-E-(4I+P+R+1)\le\mu_N.
  \]
\end{enumerate}
\end{lemma}
\leanref{Formalism/SizeWindowBrackets/ChainSlopeWindows/WeightedSlopeExcessDecomposition.lean}{Formalism.LFunctionChainParameters.UnitWeightPartition.weightedSlopeExcess_decomposition}
\leanref{Formalism/SizeWindowBrackets/ChainSlopeWindows/WeightedSlopeExcessDecomposition.lean}{Formalism.LFunctionChainParameters.UnitWeightPartition.claimHigh_valid}
\leanref{Formalism/SizeWindowBrackets/ChainSlopeWindows/WeightedSlopeExcessDecomposition.lean}{Formalism.LFunctionChainParameters.UnitWeightPartition.lower_bound_simplified}
\leanref{Formalism/SizeWindowBrackets/ClaimGaFree.lean}{Formalism.LFunctionChainParameters.UnitWeightPartition.HypLow}
\leanref{Formalism/SizeWindowBrackets/ClaimGaFree.lean}{Formalism.LFunctionChainParameters.UnitWeightPartition.HypLow2}
\leanref{Formalism/SizeWindowBrackets/ClaimGaFree.lean}{Formalism.LFunctionChainParameters.UnitWeightPartition.HypHigh}
\begin{proof}
\emph{The decomposition.} Expanding $s_w$ and exchanging the sums over $j$ and $w$,
$\sum_{w\in W}2l_w(s_w-\tau_w)=G+F-E$.

(1) Since $l_w\ge0$ and $s'_w\le s_w$, $\Sigma'\le\sum_w2l_w(s_w-\tau_w)=G+F-E$. Hence
$\mu^+_N\le\mu^+_R+\Sigma'\le\mu^+_M+G+F-E$.

(2) By the decomposition and the bound with $R$, $\Sigma'\ge G+F-E-R$. Write
$\beta_w=2l_w(\operatorname{len}_w-d_N\tau_w)/d_N$; here $d_N>0$. For $w\in W_T$ we have
$\tau_w=0$, so $\sum_{w\in W_T}\beta_w\le P$. For $w\in W\smallsetminus W_T$, the bound on
$\operatorname{len}_w$ and $d_N=(1-\delta)d_M$ give
\[
 \beta_w\le\frac{2l_w\big(d_M(\tau_w+\iota_w)-(1-\delta)d_M\tau_w\big)}{(1-\delta)d_M}
 =\frac{2l_w\iota_w}{1-\delta}+\frac{\delta}{1-\delta}\,2l_w\tau_w.
\]
Summing, with $l_w\iota_w\ge0$ for all $w$ and $\tau_w=0$ on $W_T$,
\[
 \sum_{w\in W}\beta_w\le P+\frac{2I}{1-\delta}+\frac{\delta}{1-\delta}\,2\sum_{w\in W}l_w\tau_w.
\]
Hence $\mu_N\ge\mu_R+G+F-E-R-P-2I/(1-\delta)-\frac{\delta}{1-\delta}2\sum_wl_w\tau_w$. Since
$\delta\le1/2$, we have $1/(1-\delta)\le2$ and $\delta/(1-\delta)\le2\delta\le4f/T$. As $I\ge0$ and
$0\le2\sum_wl_w\tau_w\le b^2$, the last two terms are at most $4I$ and $4fb^2/T$. With
$\mu_R=\mu_M-f\log(T-1)/(T-2)$ and $4fb^2/T+f\log(T-1)/(T-2)\le1$ this gives the claim.
\end{proof}

\begin{lemma}\label{asm-global-level-t}
Let $\mathcal F$ be a family with exponent $\varepsilon$ (Definition~\ref{pd-family}) and $D$ a
member at a scale $x>x_1$, with ranks $k,m,t,n$, $f=2n$, coefficient field $L$, top classes $T$,
and size data (Definition~\ref{pd-size-data}) with the quantities $a,b,b'$, cloud size $S_c$, width
$\eta$ and numbers $\operatorname{len}_{u,w}$. Let $u$ be a word, with the notation of
Definition~\ref{asm-level-setting}, and let $N^\circ_u\subset V(N_u)$, $\nu_u$ and
$M^{\mathrm{ram}}_u$ be as in Lemma~\ref{asm-level-commensurable}, with $M^{\mathrm{ram}}_u$
carrying the norms of $N_u$. For a finite place $w$ of $L$ put
$\ell_w=\log\mathrm N\mathfrak p_w/[L:\mathbf Q]$ and $\mathrm{sz}_w(y)=2\ell_we_wy$, and for
$\mathcal O_L$-submodules $A,B$ of $V(N_u)$ let $\operatorname{len}_w(B/A)$ be the length of the
localization of $B/(A\cap B)$ at $\mathfrak p_w$. Suppose that:
\begin{enumerate}
\item the setting of the level primes, the division side conditions and the congruence length
  bound hold at $u$ (Definition~\ref{asm-level-setting});
\item the numerical conditions and the weight conditions with some $C_n$ hold at $u$ for
  $(x,b,3x^{\varepsilon},b',S_c,\eta)$ (Definition~\ref{asm-global-numerical});
\item for every place $w$ over $T_u$, a vector $y\in V(N_u)$ has $\sigma y\in N^\circ_u$ for some
  $\sigma\in\mathcal O_L\smallsetminus\mathfrak p_w$ exactly when $\nu_u(y)\in N_\lambda(u,w)$; and
  if $\operatorname{len}_w(M^{\mathrm{ram}}_u/N^\circ_u)$ is finite, it equals
  $e_w\operatorname{len}_{u,w}$.
\end{enumerate}
Then:
\begin{enumerate}
\item $\dim M_{u,\mathrm{ram}}=\big(1-\frac1{T_u-1}\big)^f\dim M_u$;
\item $\mu(M^{\mathrm{ram}}_u)=\mu(\Lambda_{u,0})-f\log(T_u-1)/(T_u-2)$ and
  $\mu_{\max}(M^{\mathrm{ram}}_u)\le\mu_{\max}(\Lambda_{u,0})$;
\item at every place $w$ over $T_u$,
  $\operatorname{len}_w(M^{\mathrm{ram}}_u/N^\circ_u)$ is the length of the
  $\mathcal O_{L,w}$-module $S^{\mathrm{ram}}_w/N_\lambda(u,w)$;
\item if $\dim M_{u,\mathrm{ram}}>0$, then for every set $S$ of places of $L$ over $T_u$
  \[
   \frac{\sum_{w\in S}\mathrm{sz}_w(\operatorname{len}_{u,w})}{\dim M_{u,\mathrm{ram}}}\le
   36\,|T|\,m\Big(\frac ak+b'\Big)\log T_u .
  \]
\end{enumerate}
\end{lemma}
\leanref{Formalism/SharpRestriction/GlobalSchurBounds/LinkMin.lean}{Formalism.MainNotions.J.division_rank_slopes_of_lines}
\leanref{Formalism/SharpRestriction/GlobalSchurBounds/LinkMin.lean}{Formalism.MainNotions.J.RamifiedRankLines}
\leanref{Formalism/LocalRingCardinality/Lemmas.lean}{Formalism.BasicNotions.HLat.slopes_of_isom}
\leanref{Formalism/LocalRingCardinality/Lemmas.lean}{Formalism.BasicNotions.D7.lenLoc_atT}
\leanref{Formalism/SharpRestriction/AdmissibleParameterJoins/DIVfromGLo.lean}{Formalism.MainNotions.J.fifthLevelSizes_of_weights}
\leanref{Formalism/SharpRestriction/AdmissibleParameterJoins/DIVfromGLd.lean}{Formalism.MainNotions.J.LevelPrimeSizeBounds}
\leanref{Formalism/SharpRestriction/AdmissibleParameterJoins/DIVfromGLd.lean}{Formalism.MainNotions.J.sum_lwP_e_le}
\leanref{Formalism/SharpRestriction/Divisibility/LocalQuadraticParts.lean}{Formalism.MainNotions.LocalQuadraticParts.adm26}
\begin{proof}
(1) This is the first division side condition.

(2) The map $\nu_u$ is injective, preserves the norms at the infinite places, and carries
$V(N_u)$ onto $V^{\mathrm{ram}}_{u,0}$; since $M^{\mathrm{ram}}_u=\nu_u^{-1}(\Lambda_{u,0})$, it
carries $M^{\mathrm{ram}}_u$ onto $\Lambda_{u,0}\cap V^{\mathrm{ram}}_{u,0}=\Lambda^{\mathrm{ram}}_u$.
So $M^{\mathrm{ram}}_u$ and $\Lambda^{\mathrm{ram}}_u$ are isometric and have the same slopes, and
the other two division side conditions give (2).

(3) The localization of $M^{\mathrm{ram}}_u$ at $\mathfrak p_w$ is its $\mathcal O_{L,w}$-span,
which $\nu_u$ carries onto the $\mathcal O_{L,w}$-span of $\Lambda^{\mathrm{ram}}_u$, that is onto
$S_w\cap V^{\mathrm{ram}}_{u,0}=S^{\mathrm{ram}}_w$. The localization of $N^\circ_u$ is the set of
$y$ with $\sigma y\in N^\circ_u$ for some $\sigma\notin\mathfrak p_w$, which $\nu_u$ carries onto
$N_\lambda(u,w)\cap V^{\mathrm{ram}}_{u,0}$ by hypothesis (3). As $\nu_u$ is injective, it
identifies the localization of $M^{\mathrm{ram}}_u/(N^\circ_u\cap M^{\mathrm{ram}}_u)$ with
$S^{\mathrm{ram}}_w/(N_\lambda(u,w)\cap S^{\mathrm{ram}}_w)$, and $N_\lambda(u,w)\subset
S^{\mathrm{ram}}_w$ by the congruence length bound.

(4) \emph{One place.} Let $w\in S$. By the congruence length bound and (3),
$\operatorname{len}_w(M^{\mathrm{ram}}_u/N^\circ_u)$ is a finite number $n_w$ with
\[
 n_w\le e_w\,m\dim M_u\sum_{v\in V_{\mathrm{top}}(w)}\Big(\frac{a_{\max}(u)}k
 +2W_{\tau_w(v)}(u)+k+b_{\mathrm{top}}(\tau_w(v))\Big),
\]
and $n_w=e_w\operatorname{len}_{u,w}$ by hypothesis (3). So
$\mathrm{sz}_w(\operatorname{len}_{u,w})=2\ell_wn_w$.

\emph{The sum over $V_{\mathrm{top}}(w)$.} The numerical and weight conditions give
$N(m)\le b'$, $a_{\max}(u)\le2a$, and, at every top embedding $\tau$, $b'\le
b_{\mathrm{top}}(\tau)=-\lambda_{u,0}(\tau)_k<2b'$ and $|\lambda_{u,i}(\tau)_c|\le2b'$ for
$0\le i\le m$ and $c\ge2$; as $i(t+1)/2\le N(m)\le b'$, this gives $W_\tau(u)\le3b'$. With
$k\le N(m)\le b'$, each term of the sum is at most $2a/k+6b'+b'+2b'\le9(a/k+b')$. The map
$v\mapsto\tau_w(v)$ is injective on $V_{\mathrm{top}}(w)$ with values among the top embeddings,
which are at most $|T|$ in number. So the sum is at most $9|T|(a/k+b')$, and
\[
 \mathrm{sz}_w(\operatorname{len}_{u,w})\le2\ell_we_w\cdot m\dim M_u\cdot9|T|\Big(\frac ak+b'\Big).
\]

\emph{The sum over $S$.} For a prime $p$, the numbers $\ell_we_w$ of the places $w$ over $p$ add
up to $\log p$. So
$\sum_{w\in S}\mathrm{sz}_w(\operatorname{len}_{u,w})\le18\,m\dim M_u\,|T|(a/k+b')\log T_u$.

\emph{The rank.} The setting gives $2fN(m)+1<T_u$, so $f/(T_u-1)\le1/2$, and Bernoulli's
inequality gives $(1-\frac1{T_u-1})^f\ge1-f/(T_u-1)\ge1/2$. By (1),
$\dim M_{u,\mathrm{ram}}\ge\dim M_u/2$, and dividing the last bound by $\dim M_{u,\mathrm{ram}}>0$
gives (4).
\end{proof}

\begin{lemma}\label{asm-global-slopes-n-member}
Let $\mathcal F$ be a family with exponent $\varepsilon$ (Definition~\ref{pd-family}) and $D$ a
member at a scale $x>x_1$, with the notation of Definition~\ref{asm-member-bounds}. Let its size
data have support $\operatorname{supp}$, depths $\tau_w$, index excesses $\iota_w$, numbers
$\operatorname{len}_{u,w}$ and the numbers $s(u,w)$ of Definition~\ref{pd-size-data}(4). For a
word $u$ let $N^\circ_u$, $\nu_u$ and $M^{\mathrm{ram}}_u$ be as in
Lemma~\ref{asm-level-commensurable}, with $N^\circ_u$ and $M^{\mathrm{ram}}_u$ carrying the norms
of $N_u$, put $I_u=\prod_{w\in\operatorname{supp}}\mathfrak p_w^{-\lfloor e_ws(u,w)\rfloor}$, and let
$\operatorname{len}_w$ be as in Lemma~\ref{asm-global-level-t}. Let $C\ge36$. Suppose that every
$\tau_w$ is a nonnegative integer, that $\iota_w\ge0$ for every $w$, that
$\sum_{w\in\operatorname{supp}}\mathrm{sz}_w(\tau_w)\le b^2$, and that for every word $u$:
\begin{enumerate}
\item $N_u$ is the $\mathcal O_L$-span of $I_uN^\circ_u$; it is finitely generated, spans $V(N_u)$
  and has positive rank; $N^\circ_u\subset M^{\mathrm{ram}}_u$, and
  $cM^{\mathrm{ram}}_u\subset N^\circ_u$ for some nonzero $c\in\mathcal O_L$;
\item at every finite place $w$, if $\operatorname{len}_w(M^{\mathrm{ram}}_u/N^\circ_u)$ is finite,
  it equals $e_w\operatorname{len}_{u,w}$; and if $w$ is not in the support, then
  $\operatorname{len}_w(M^{\mathrm{ram}}_u/N^\circ_u)=0$ or $\operatorname{len}_{u,w}=0$;
\item conclusions (1), (2) and (4) of Lemma~\ref{asm-global-level-t} hold at $u$;
\item $\big|\mu(N_u)-\mu(N^\circ_u)-\sum_{w\in\operatorname{supp}}\mathrm{sz}_w(s(u,w))\big|\le CX''$;
\item at every place $w$ with $p_w\ne T_u$: locally at $w$,
  $N^\circ_u\subset\mathfrak p_w^{r_w}M^{\mathrm{ram}}_u$ with $r_w=\lfloor e_w\tau_w\rfloor$, that
  is, every $y\in N^\circ_u$ has $\sigma y\in\mathfrak p_w^{r_w}M^{\mathrm{ram}}_u$ for some
  $\sigma\in\mathcal O_L\smallsetminus\mathfrak p_w$; and
  $\operatorname{len}_{u,w}\le\dim M_u(\tau_w+\iota_w)$;
\item $\tau_w=0$ at every place $w$ over $T_u$; $T_u>2$, $2f/T_u\le1/2$ and
  $4fb^2/T_u+f\log(T_u-1)/(T_u-2)\le1$.
\end{enumerate}
Then $D$ satisfies the slope bounds for $N$ with $C$ (Definition~\ref{asm-member-bounds}(3)).
\end{lemma}
\leanref{Formalism/LocalToGlobal/GlobalTheory/PlacesP.lean}{Formalism.L2.GlobalTheoryTheory.Places.places_at_of}
\leanref{Formalism/LocalToGlobal/GlobalTheory/PlacesP.lean}{Formalism.L2.GlobalTheoryTheory.Places.mu_eq}
\leanref{Formalism/LocalToGlobal/GlobalTheory/PlacesP.lean}{Formalism.L2.GlobalTheoryTheory.Places.muMax_le}
\leanref{Formalism/LocalToGlobal/GlobalTheory/PlacesP.lean}{Formalism.L2.GlobalTheoryTheory.Places.le_prod_smul}
\leanref{Formalism/SizeWindowBrackets/ChainSlopeWindows/WeightedSlopeExcessDecomposition.lean}{Formalism.LFunctionChainParameters.UnitWeightPartition.deficit_bounds}
\leanref{Formalism/SharpRestriction/GlobalSchurBounds/TopDensityNGeSixteen.lean}{Formalism.MainNotions.Lnk.Places_at}
\leanref{Formalism/SharpRestriction/GlobalSchurBounds/TopDensityNGeSixteen.lean}{Formalism.MainNotions.Lnk.slopesN_at_of}
\leanref{Formalism/SharpRestriction/GlobalSchurBounds/QuadraticLogDerivCongr.lean}{Formalism.MainNotions.Lnk.SlopesN_at}
\begin{proof}
Fix a word $u$, and write $d_M=\dim M_u$, $d_N=\operatorname{rk}N_u$ and $T=T_u$.

\emph{Step 1: ranks.} The map $j_N$ carries a basis of $V(N_u)$ to a basis of
$\mathcal M_u\cong M_{u,\mathrm{ram}}$ (Definition~\ref{pd-period-data}(7)), and $N_u$ spans
$V(N_u)$, so $d_N=\dim M_{u,\mathrm{ram}}=(1-\frac1{T-1})^fd_M$ by (3). As $d_N>0$, also $d_M>0$.
Put $\delta=1-d_N/d_M$. Then $d_N=(1-\delta)d_M$, and since $T>2$ and
$0<1-\frac1{T-1}<1$, Bernoulli's inequality gives $0\le\delta<1$ and $\delta\le f/(T-1)\le2f/T$.
The lattices $N^\circ_u\subset M^{\mathrm{ram}}_u$ both have rank $d_N$, because
$cM^{\mathrm{ram}}_u\subset N^\circ_u$ and $N_u$ spans $V(N_u)$.

\emph{Step 2: the slope of $N_u$.} Off the support,
$\operatorname{len}_w(M^{\mathrm{ram}}_u/N^\circ_u)=0$: by (2) either this is so, or
$\operatorname{len}_{u,w}=0$, and then the finite length equals $e_w\cdot0$. In the support it is
$e_w\operatorname{len}_{u,w}$ by (2). Lemma~\ref{asm-global-slope-index}(1) gives
$\mu(N^\circ_u)=\mu(M^{\mathrm{ram}}_u)-\frac2{d_N}\sum_{w\in\operatorname{supp}}\ell_we_w
\operatorname{len}_{u,w}$, and Lemma~\ref{asm-global-slope-index}(2), with
$\log\mathrm N(I_u)=-\sum_{w\in\operatorname{supp}}\lfloor e_ws(u,w)\rfloor\log\mathrm N\mathfrak p_w$,
gives $\mu(N_u)=\mu(N^\circ_u)+\sum_{w\in\operatorname{supp}}2\ell_w\lfloor e_ws(u,w)\rfloor$. So
\[
 \mu(N_u)=\mu(M^{\mathrm{ram}}_u)+\sum_{w\in\operatorname{supp}}2\ell_w\lfloor e_ws(u,w)\rfloor
 -\frac2{d_N}\sum_{w\in\operatorname{supp}}\ell_we_w\operatorname{len}_{u,w}.
\]

\emph{Step 3: the maximal slope of $N_u$.} Put $k_w=\lfloor e_w\tau_w\rfloor=e_w\tau_w$ for $w$ in
the support. At a place $w$ of the support with $p_w\ne T$, (5) gives $N^\circ_u\subset\mathfrak
p_w^{k_w}M^{\mathrm{ram}}_u$ locally at $w$. At the other places, $N^\circ_u\subset
M^{\mathrm{ram}}_u$ gives the same with exponent $0$, and $k_w=0$ at the places of the support over
$T$, by (6). Let $\mathfrak a=\prod_{w\in\operatorname{supp}}\mathfrak p_w^{k_w}$ and $y\in
N^\circ_u$. For each place $v$ choose $t\in\prod_{w\in\operatorname{supp},\,w\ne v}\mathfrak
p_w^{k_w}$ with $t\notin\mathfrak p_v$, which exists by prime avoidance, and $\sigma\notin\mathfrak
p_v$ with $\sigma y\in\mathfrak p_v^{k_v}M^{\mathrm{ram}}_u$ (with $k_v=0$ off the support). Then
$t\sigma y\in\mathfrak aM^{\mathrm{ram}}_u$ and $t\sigma\notin\mathfrak p_v$. The ideal of the
$\rho\in\mathcal O_L$ with $\rho y\in\mathfrak aM^{\mathrm{ram}}_u$ is therefore contained in no
maximal ideal, so $y\in\mathfrak aM^{\mathrm{ram}}_u$. Hence $N_u\subset I_u\mathfrak
aM^{\mathrm{ram}}_u$, and Lemma~\ref{asm-global-slope-index}(3) and (2) give
\[
 \mu_{\max}(N_u)\le\mu_{\max}(M^{\mathrm{ram}}_u)+\sum_{w\in\operatorname{supp}}2\ell_w
 \big(\lfloor e_ws(u,w)\rfloor-e_w\tau_w\big).
\]

\emph{Step 4: Lemma~\ref{asm-global-unit-weight}.} Apply it with $W=\operatorname{supp}$, $W_T$
the places of the support over $T$, $J$ the lifts $1\le j\le m$, $l_w=\ell_we_w$,
$\mathrm{gain}_j(w)=\mathrm{gain}^\flat_j(w)$, the terms $F_j(w)$ and $E_{u,j}(w)$, $\tau_w$,
$\iota_w$, $\operatorname{len}_w=\operatorname{len}_{u,w}$ and
$s'_w=\lfloor e_ws(u,w)\rfloor/e_w$. Then $s_w=s(u,w)$, and $G,F,E,I$ are $G_a,F_o,E_u,I$ of
Definition~\ref{asm-member-bounds}. Take $d_M,d_N,\delta$ of Step 1, $f$, $T$, $b$,
\[
 P=\frac1{d_N}\sum_{w\in W_T}\mathrm{sz}_w(\operatorname{len}_{u,w}),\qquad
 R=\sum_{w\in\operatorname{supp}}\mathrm{sz}_w(s(u,w))-\sum_{w\in\operatorname{supp}}2\ell_w\lfloor
 e_ws(u,w)\rfloor,
\]
$\mu_M=\mu(\Lambda_{u,0})$, $\mu_R=\mu(M^{\mathrm{ram}}_u)$, $\mu_N=\mu(N_u)$, and the
corresponding maximal slopes. The hypotheses of part (1) are $s'_w\le s(u,w)$, Step 3 and
conclusion (2) of Lemma~\ref{asm-global-level-t}. Those of part (2) are Step 1, (6), the bounds on
$\tau_w$ and $\iota_w$, (5) at the places of the support not over $T$, Step 2 for $\mu_N$,
conclusion (2) of Lemma~\ref{asm-global-level-t} for $\mu_R$, and
$2\sum_wl_w\tau_w=\sum_w\mathrm{sz}_w(\tau_w)\le b^2$. So
\[
 \mu_{\max}(N_u)\le\mu_{\max}(\Lambda_{u,0})+G_a+F_o-E_u,\qquad
 \mu(\Lambda_{u,0})+G_a+F_o-E_u-(4I+P+R+1)\le\mu(N_u).
\]

\emph{Step 5: the error terms.} Conclusion (4) of Lemma~\ref{asm-global-level-t}, with
$d_N=\dim M_{u,\mathrm{ram}}>0$ and $C\ge36$, gives $P\le C|T|m(a/k+b')\log T$. By Step 2,
$R=-\big(\mu(N_u)-\mu(N^\circ_u)-\sum_w\mathrm{sz}_w(s(u,w))\big)$, so $R\le CX''$ by (4). As
$\iota_w\ge0$, $I\ge0$, and $4I\le CI$. Hence $4I+P+R+1\le\mathrm{Lo}_C(u)$, which is the lower
bound. Since $C\ge0$ and $B_0=fmbx^{\varepsilon}\ge0$, the upper bound gives
$\mu_{\max}(N_u)\le\mu_{\max}(M_u)+G_a+F_o-E_u+CmB_0$.
\end{proof}

\begin{lemma}\label{asm-global-slopes-n}
Let $\mathcal F$ be a family with exponent $\varepsilon$ (Definition~\ref{pd-family}) that
satisfies the size condition on $T_u$ (Definition~\ref{asm-designed}(2)), has the level data
(Definition~\ref{asm-level-setting}), and satisfies the gain link (Definition~\ref{pd-gain-link})
and the slope link (Definition~\ref{asm-global-inputs}(1)). Suppose that for all large $x$, at
every member at $x$ and every word $u$, with $N^\circ_u$ and $M^{\mathrm{ram}}_u$ as in
Lemma~\ref{asm-level-commensurable}: $N^\circ_u\subset M^{\mathrm{ram}}_u$, $N^\circ_u$ is
finitely generated and $cM^{\mathrm{ram}}_u\subset N^\circ_u$ for some nonzero $c\in\mathcal O_L$;
and $N_u$ is finitely generated, spans $V(N_u)$ and has positive rank. Suppose also the \emph{depth
bound}: there are $C_T$, $x_T$ and a function $\delta$ with $\delta(x)\to0$ as $x\to\infty$ such
that, with the notation of Definition~\ref{asm-member-bounds},
\[
 \sum_{w\in\operatorname{supp}}\mathrm{sz}_w(\tau_w)\le C_T\,m\,|T|\,a\log X'+b^{\delta(x)}
\]
at every member at every scale $x\ge x_T$. Then $\mathcal F$ satisfies the slope bounds for $N$ at
all large scales (Definition~\ref{asm-member-bounds}(3)).
\end{lemma}
\leanref{Formalism/SharpRestriction/GlobalSchurBounds/TopDensityNGeSixteen.lean}{Formalism.MainNotions.Lnk.slopesNF_of}
\leanref{Formalism/SharpRestriction/GlobalSchurBounds/TopDensityNGeSixteen.lean}{Formalism.MainNotions.Lnk.PlacesF}
\leanref{Formalism/LocalToGlobal/GlobalTheory/PlacesMin.lean}{Formalism.L2.GlobalTheoryTheory.Places.placesF_min'}
\leanref{Formalism/SharpRestriction/GlobalSchurBounds/LinkMin.lean}{Formalism.MainNotions.J.lemma_five_one_of_theorem_two}
\leanref{Formalism/SharpRestriction/Divisibility/LocalQuadraticParts.lean}{Formalism.MainNotions.LocalQuadraticParts.LevelTFamily}
\leanref{Formalism/LocalToGlobal/GlobalTheory/PlacesP.lean}{Formalism.L2.GlobalTheoryTheory.Places.sizes_at}
\leanref{Formalism/LocalToGlobal/GlobalTheory/PlacesP.lean}{Formalism.L2.GlobalTheoryTheory.Places.size_alt}
\leanref{Formalism/LocalToGlobal/GlobalTheory/PlacesP.lean}{Formalism.L2.GlobalTheoryTheory.Places.ta_le}
\leanref{Formalism/LocalToGlobal/GlobalTheory/PlacesP.lean}{Formalism.L2.GlobalTheoryTheory.Places.lenLoc_atT_zero}
\leanref{Formalism/SharpRestriction/GrowthBounds/CornerGLContinuity.lean}{Formalism.MainNotions.CornerGLContinuity.FamilyWeightSumBound}
\begin{proof}
\emph{Thresholds.} The size condition gives $C_{T_u}$ with $\log T_u\le C_{T_u}mx^{\varepsilon}$,
and the size data give $|\operatorname{Re}\frac{L'}{L}(1,\chi_d)|\le x^{\varepsilon}$; so
Lemma~\ref{asm-global-numerical-hold} with $C_c=C_{T_u}$ gives $C_n$ such that, for all large $x$,
the numerical and weight conditions with $C_n$ hold at every word of every member for
$(x,b,3x^{\varepsilon},b',S_c,\eta)$. By Lemma~\ref{asm-level-setting-holds} and the level data,
for all large $x$ the setting of the level primes, the division side conditions and the congruence
length bound hold at every word of every member. Let $C_G$ and $C_\mu$ be constants of the gain link
and the slope link, and put $C=\max(36,C_\mu)$. Let $x$ be a scale beyond all these thresholds and
those of the hypotheses, with $x>x_1$, $x\ge5$, $x\ge4|C_T|+1$, $\delta(x)\le1$, $x^9\le b$ and
$m\ge6$; since $b=\exp(x^{\varepsilon})$, and since the sizes of the family give $m\ge x^{2/5}$,
all large $x$ qualify. Let $D$ be a member at $x$ and $u$ a word.

\emph{The choices of the family.} Among the choices of Definition~\ref{pd-family}(4): $N_u$ is the
$\mathcal O_L$-span of $I_uN^\circ_u$; whenever
$\operatorname{len}_w(M^{\mathrm{ram}}_u/N^\circ_u)$ is finite it equals
$e_w\operatorname{len}_{u,w}$; at a place $w$ over $T_u$, $\sigma y\in N^\circ_u$ for some
$\sigma\notin\mathfrak p_w$ exactly when $\nu_u(y)\in N_\lambda(u,w)$; every place $w$ with $p_w\le
X''$ is in the support; the size data are filled by the gain data at every place
(Definition~\ref{pd-gain-quantities}(6)), so $\tau_w=\tau(p_w)$ and $\iota_w=\iota(p_w)$; if
$V(p_w)=\emptyset$ and $p_w\ne T_u$, then $\operatorname{len}_{u,w}=0$; and the local errors are
explicit (Definition~\ref{asm-division-kernel-errors}(1)).

\emph{Sizes.} The sizes of the family give $f,m,|T|\le2x^3$, $0\le\log X'\le x$, $X'\le b$,
$X'\le X''$, $16f\le b$ and $4\le b$. The size condition gives $b^3\le b^{m/2}\le T_u$, so
$X'\le b\le T_u$. From $16f\le b$, $4\le b$ and $b^3\le T_u$: $b^2\le T_u$, so $2f/T_u\le1/2$ and
$T_u>2$; $4fb^2/T_u\le4f/b\le1/4$; and $\log(T_u-1)\le2\sqrt{T_u}$ with $16f\le b\le\sqrt{T_u}$
gives $f\log(T_u-1)/(T_u-2)\le1/4$. So (6) of Lemma~\ref{asm-global-slopes-n-member} holds, apart
from $\tau_w=0$.

\emph{Depths and index excesses.} By Definition~\ref{pd-gain-quantities}(3), $\tau(p)$ is a
nonnegative integer, $\tau(p)=0$ when $V(p)=\emptyset$, and $\iota(p)\ge0$. The set $V(p)$ is empty
for $p\ge X'$, since candidate primes are below $X'$. As $T_u\ge X'$, $\tau_w=0$ at every place
over $T_u$. The depth bound gives, as $m|T|a\log X'\le(2x^3)^2\cdot xb\cdot x$ and
$b^{\delta(x)}\le b$,
\[
 \sum_{w\in\operatorname{supp}}\mathrm{sz}_w(\tau_w)\le4|C_T|x^8b+b\le b^2,
\]
because $(4|C_T|+1)x^8\le x^9\le b$.

\emph{The level at $T_u$.} Lemma~\ref{asm-global-level-t} applies at $u$, with the choices above
for its hypothesis (3).

\emph{Places off the support.} Let $w$ be outside the support. Then $p_w>X''\ge X'$, so
$V(p_w)=\emptyset$. If $p_w\ne T_u$, then $\operatorname{len}_{u,w}=0$. If $p_w=T_u$, then
$E_{u,1}(w)=0$, as $w$ is outside the support, while the explicit errors give
$E_{u,1}(w)=E_T(u,w)$, which is at least $fN(m)>0$ if $w$ sees $\pi_u$
(Definition~\ref{pd-hecke-at-T}(4)). So $w$ does not see $\pi_u$, hence
$N_\lambda(u,w)=S^{\mathrm{ram}}_w$ as in the proof of Lemma~\ref{asm-level-congruence}, and
$\operatorname{len}_w(M^{\mathrm{ram}}_u/N^\circ_u)=0$ by Lemma~\ref{asm-global-level-t}(3).

\emph{Conclusion.} The remaining hypotheses of Lemma~\ref{asm-global-slopes-n-member} are those of
the present lemma, the gain link (1), and the slope link with $C_\mu\le C$ and $X''\ge0$. That
lemma gives the
slope bounds for $N$ with $C$ at every member at $x$.
\end{proof}

\subsection{Column values and local facts}\label{subsection-asm-column-values}

\begin{definition}\label{asm-global-link-column-values-local-basis}
Let $E$ be a number field, $k\ge1$, $\ell$ a prime and $M\subset E^k$ a subset.
A matrix $Y\in GL_k(E)$ is a \emph{local basis of $M$ at $\ell$} if for every
$x\in E^k$: $sx\in M$ for some integer $s$ prime to $\ell$ if and only if
$s\,Y^{-1}x\in\mathcal O_E^k$ for some integer $s$ prime to $\ell$.
\end{definition}
\leanref{Formalism/LocalToGlobal/HermitianForms/IsLocBasis.lean}{Formalism.L2.Forms.IsLocBasis}

\begin{lemma}\label{asm-global-link-column-values-local-basis-facts}
Let $E$ be a number field, $k\ge1$ and $\ell$ a prime.
\begin{enumerate}
\item Every finitely generated $\mathcal O_E$-submodule $M\subset E^k$ that
  spans $E^k$ has a local basis at $\ell$.
\item Let $M$ be as in (1), let $\sigma$ be an automorphism of $E$ and $G\in
  M_k(E)$, and put $M^\vee=\{y\in E^k:\sigma(x)^{\mathsf T}Gy\in\mathcal O_E
  \text{ for all }x\in M\}$. If $Y$ is a local basis of $M$ at $\ell$ and
  $\sigma(Y)^{\mathsf T}G$ is invertible, then $(\sigma(Y)^{\mathsf T}G)^{-1}$
  is a local basis of $M^\vee$ at $\ell$.
\item If $Y$ and $Y'$ are local bases at $\ell$ of the same subset of $E^k$,
  then every entry of $Y'^{-1}Y$ is $a/s$ with $a\in\mathcal O_E$ and $s$ an
  integer prime to $\ell$. In particular, for a ring homomorphism $\phi$ from
  $E$ to a number field $L$ and a finite place $w$ of $L$ over $\ell$, the
  entries of $\phi(Y'^{-1}Y)$ lie in the valuation ring $\mathcal O_{L,w}$.
\end{enumerate}
\end{lemma}
\leanref{Formalism/LocalToGlobal/GlobalTheory/GL6/PBwLocGen.lean}{Formalism.L2.GlobalTheoryTheory.LocalLatticeBases.LocB.exists_isLocBasis}
\leanref{Formalism/LocalToGlobal/GlobalTheory/ColMemOfIsLocBasis.lean}{Formalism.L2.GlobalTheoryTheory.ColValS.isLocBasis_dual}
\leanref{Formalism/LocalToGlobal/GlobalTheory/ColMemOfIsLocBasis.lean}{Formalism.L2.GlobalTheoryTheory.ColValS.change_mem}
\begin{proof}
Let $S$ be the set of integers prime to $\ell$. Two facts are used throughout.
If $Y$ is a local basis of a set $M$ at $\ell$, then every column $Ye_a$ has a
multiple $t\,Ye_a\in M$ with $t\in S$, since $Y^{-1}Ye_a=e_a$ is integral. If
moreover $M$ is an $\mathcal O_E$-module generated by $x_1,\dots,x_r$, there is
$d\in S$ with $d\,Y^{-1}x\in\mathcal O_E^k$ for all $x\in M$: take for $d$ the
product of integers $s_i\in S$ with $s_iY^{-1}x_i$ integral.

(1) The ring $R=S^{-1}\mathcal O_E$ is a Dedekind domain with finitely many
maximal ideals, hence a principal ideal domain. The $R$-module $S^{-1}M$ is
finitely generated and torsion-free and spans $E^k$, so it is free of rank $k$.
Let $Y$ have a basis of $S^{-1}M$ as its columns. Then $sx\in M$ for some
$s\in S$ if and only if $x\in S^{-1}M$, if and only if $Y^{-1}x\in R^k$, if and
only if $sY^{-1}x\in\mathcal O_E^k$ for some $s\in S$.

(2) Put $Z=\sigma(Y)^{\mathsf T}G$, and note that
$\sigma(Yv)^{\mathsf T}Gy=\sum_a\sigma(v_a)(Zy)_a$ for $v\in E^k$. Suppose
$sy\in M^\vee$ with $s\in S$. For each $a$ choose $t_a\in S$ with $t_aYe_a\in
M$; then $st_a(Zy)_a=\sigma(t_aYe_a)^{\mathsf T}G(sy)\in\mathcal O_E$, so
$s\prod_bt_b\cdot Zy$ is integral. Conversely, suppose $s'Zy$ is integral with
$s'\in S$, and let $d$ be as above for $M$. For $x\in M$ write $x=Yv$ with
$v=Y^{-1}x$; then
$\sigma(x)^{\mathsf T}G(ds'y)=\sum_a\sigma(dv_a)\cdot s'(Zy)_a\in\mathcal O_E$,
as $dv_a\in\mathcal O_E$ and $\sigma(\mathcal O_E)=\mathcal O_E$. So
$ds'y\in M^\vee$.

(3) For a column $Ye_b$ choose $t\in S$ with $tYe_b$ in the set. As $Y'$ is a
local basis of it, $s'Y'^{-1}(tYe_b)$ is integral for some $s'\in S$; this is
the column $b$ of $Y'^{-1}Y$ times $s't$. A ring homomorphism $\phi$ carries
$\mathcal O_E$ into $\mathcal O_L\subset\mathcal O_{L,w}$, and an integer prime
to $\ell$ is a unit at $w$.
\end{proof}

\begin{lemma}\label{asm-global-link-column-values-weyl-stable}
Let $K$ be a field of characteristic $0$, $R\subset K$ a local subring,
$\mathcal T$ a finite set, $n\ge1$, and $\lambda=(\lambda_\tau)_{\tau\in\mathcal
T}$ dominant weights with $\nu_{\tau,i}=\lambda_{\tau,i}-\lambda_{\tau,n}$. Let
$\Delta_{\lambda,R}\subset K[X_\tau:\tau\in\mathcal T]$ be the $R$-span of the
products $\prod_\tau f_\tau(X_\tau)$ with $f_\tau\in\Delta_{\mathbf Z}(\nu_\tau)$
(Definition~\ref{pd-weight-lattices}(1)), and let $GL_n(K)^{\mathcal T}$ act on
$K[X_\tau:\tau\in\mathcal T]$ by the formula for $\rho_\lambda$ of
Definition~\ref{pd-weight-lattices}(4). Then
$\rho_\lambda(g)\Delta_{\lambda,R}\subset\Delta_{\lambda,R}$ for every
$g=(g_\tau)$ with $g_\tau\in GL_n(R)$.
\end{lemma}
\leanref{Formalism/LocalToGlobal/GlobalTheory/GL6/PBwLocInd.lean}{Formalism.L2.GlobalTheoryTheory.LocalLatticeBases.local_GL_stabilises_lattice}
\leanref{Formalism/LocalToGlobal/GlobalTheory/GL6/PBwLocInd.lean}{Formalism.L2.GlobalTheoryTheory.LocalLatticeBases.bwG_intSpan_allSubP}
\leanref{Formalism/LocalToGlobal/GlobalTheory/GL6/PBwLocElim.lean}{Formalism.L2.GlobalTheoryTheory.LocalLatticeBases.bwE_elim}
\leanref{Formalism/LocalToGlobal/GlobalTheory/GL6/BwAValIntegralNonneg.lean}{Formalism.L2.GlobalTheoryTheory.LocalLatticeBases.bwA_span_closure}
\leanref{Formalism/LocalToGlobal/GlobalTheory/ActAllMemIntSpanAtG.lean}{Formalism.L2.GlobalTheoryTheory.ColValS.actAll_mem_intSpanAtG}
\begin{proof}
The operator $\rho_\lambda(g)$ is the composite over $\tau$ of the operators
$\rho_\lambda(g^{(\tau)})$, where $g^{(\tau)}$ is $g_\tau$ at $\tau$ and $1$
elsewhere, and $\rho_\lambda(g^{(\tau)})$ multiplies a product
$\prod_{\tau'}f_{\tau'}(X_{\tau'})$ by $\det(g_\tau)^{\lambda_{\tau,n}}$ and
replaces its factor $f_\tau(X_\tau)$ by $f_\tau({}^tg_\tau X_\tau)$. As
$\det(g_\tau)$ is a unit of $R$, it suffices to show that the $R$-span of
$\Delta_{\mathbf Z}(\nu_\tau)$ is stable under $f\mapsto f({}^tg_\tau X)$.

Since $R$ is local, Gaussian elimination writes the invertible matrix
${}^tg_\tau$ as a product of matrices $1+cE_{ba}$ with $c\in R$ and $1\le
a,b\le n$, equal or not. For such a matrix, $f((1+cE_{ba})X)$ is $f$ with the
$b$-th row of $X$ replaced by its sum with $c$ times the $a$-th row, which is
$\sum_{j\ge0}c^jE^{(j)}_{ab}f$. By Definition~\ref{pd-weight-lattices}(1),
$\Delta_{\mathbf Z}(\nu_\tau)$ is stable under the operators $E^{(j)}_{ab}$, so
its $R$-span is stable under $f\mapsto f((1+cE_{ba})X)$, and hence under
$f\mapsto f({}^tg_\tau X)$.
\end{proof}

\begin{lemma}\label{asm-global-link-column-values-criterion}
Let $D$ be a definite-side datum (Definition~\ref{th-definite-datum}) with tied
columns, as in the proof of Lemma~\ref{asm-proviso-bottom-dimension}, and let
$u$ be a word, with the notation of Definition~\ref{asm-global-column-values}.
Let $G$ be the Gram matrix of $h$ in the basis of $X_0$, $\iota\colon L\to\mathbf
C$ the embedding of $L$, $\iota_E\colon E\to L$ the embedding of the datum,
$\hat\tau=\iota\circ\iota_E\circ\sigma_\tau\colon E\to\mathbf C$ for
$\tau\in\mathcal T$, where $\sigma_\tau$ is the automorphism of $E$ given by the
element of $G_\Phi$ of $\tau$, and $s_\tau=\lambda_{u,0}(\tau)_k$; identify
$Y_0$ with the tuples $z=(z_\tau)$ of invertible complex $k\times k$ matrices,
as in the proof of Lemma~\ref{asm-proviso-bottom-dimension}. Write
$\iota(\Delta_{\lambda,w})$ for the image of $\Delta_{\lambda,w}$ under $\iota$
applied to the coefficients. For $c\in\mathrm{cls}_0$ let $\Lambda^\vee_c\subset
E^k$ be the set of coordinate vectors of the dual lattice $L_u(0,c)^\vee$. Let
$w$ be a finite place of $L$ over the prime $\ell$. Call $y\in V_{u,0}$
\emph{integral at} $(w,c)$ if for every local basis $B$ of $\Lambda^\vee_c$ at
$\ell$ (Definition~\ref{asm-global-link-column-values-local-basis}) there is
$P\in\iota(\Delta_{\lambda,w})$ with
\[
 \operatorname{emb}_0(y)(c,z)=\prod_{\tau}\det\big(\hat\tau(B)^{-1}z_\tau\big)^{s_\tau}
 P\big((\hat\tau(B)^{-1}z_\tau)_\tau\big)\qquad\text{for all }z\in Y_0.
\]
Then $y\in(\Lambda_{u,0})_w$ if and only if $y$ is integral at $(w,c)$ for every
$c\in\mathrm{cls}_0$.
\end{lemma}
\leanref{Formalism/LocalToGlobal/GroupLattices/ComparisonMaps/ExistsPrimeToLe.lean}{Formalism.L2.G.Cmp.locWCrit_proved}
\leanref{Formalism/LocalToGlobal/GroupLattices/ComparisonMaps/ExistsPrimeToLe.lean}{Formalism.L2.G.Cmp.locWBack_proved}
\leanref{Formalism/LocalToGlobal/GroupLattices/ComparisonMaps/CmpCMemMIff.lean}{Formalism.L2.G.Cmp.locW_intAtW}
\leanref{Formalism/LocalToGlobal/GroupLattices/ComparisonMaps/CmpCMemMIff.lean}{Formalism.L2.G.Cmp.IntAtW}
\begin{proof}
Among the tied columns, the forms in degree $0$ are \emph{integral}: an element
$y\in V_{u,0}$ lies in $\Lambda_{u,0}$ if and only if it is integral at $(w',c)$
for every finite place $w'$ of $L$ and every $c\in\mathrm{cls}_0$. Note that
$\operatorname{emb}_0(ay)=\iota(a)\operatorname{emb}_0(y)$ for $a\in L$, and that
$\iota(\Delta_{\lambda,w'})$ is an $\iota(\mathcal O_{L,w'})$-module.

Suppose $y\in(\Lambda_{u,0})_w$, and let $c\in\mathrm{cls}_0$. By the
integrality of the tied forms, every element of $\Lambda_{u,0}$ is integral at
$(w,c)$. The elements integral at $(w,c)$ form a group closed under
multiplication by $\mathcal O_{L,w}$: the polynomials $P$ add and scale with
them. So $y$ is integral at $(w,c)$.

Conversely, let $y$ be integral at $(w,c)$ for every $c$. As $\Lambda_{u,0}$
spans $V_{u,0}$, there is $d\in\mathcal O_L\smallsetminus\{0\}$ with $dy\in
\Lambda_{u,0}$. By the Chinese remainder theorem in $\mathcal O_L$ there is
$s\in\mathcal O_L$ outside the prime of $w$ with $v_{w'}(s)\ge v_{w'}(d)$ at every
finite place $w'\neq w$. We show that $sy$ is integral at $(w',c)$ for every
$w'$ and $c$, so that $sy\in\Lambda_{u,0}$. At $w'=w$ this follows from the
hypothesis, as $s\in\mathcal O_{L,w}$. At $w'\neq w$, $dy$ is integral at
$(w',c)$, as $dy\in\Lambda_{u,0}$, and $sy=(s/d)\,dy$ with $s/d\in\mathcal
O_{L,w'}$. Finally $s$ is a unit at $w$, so $y=s^{-1}(sy)\in(\Lambda_{u,0})_w$.
\end{proof}
% Lean: LocWCrit = locWCrit_of T (locWBack_proved T). The forward direction is locW_intAtW (span
% induction over locW, with ((T u).val0).lat and intSpan_smul); the backward direction is
% locWBack_proved with exists_den (lat_good: the lattice spans), exists_prime_to_le (the Chinese
% remainder theorem) and val0.lat in both directions. The integrality of the tied forms is the field
% lat of ValTied (Forms0.val0).

\begin{lemma}\label{asm-global-link-column-values}
Let $D$ be a definite-side datum (Definition~\ref{th-definite-datum}) with tied
columns, as in the proof of Lemma~\ref{asm-proviso-bottom-dimension}, such that
for every word $u$ every class $c\in\mathrm{cls}_0$ of the chain of $u$ has
trivial stabilizer $\Gamma_c$. Then the columns of every word of $D$ have
values (Definition~\ref{asm-global-column-values}).
\end{lemma}
\leanref{Formalism/LocalToGlobal/GlobalTheory/TransposeBaseMatrixInfinity.lean}{Formalism.L2.GlobalTheoryTheory.ColValS.colValOfTiedS_holds}
\leanref{Formalism/SharpRestriction/GlobalSchurBounds/LinkGLOBAL7Tied4.lean}{Formalism.MainNotions.RT.ColumnValuationFromTwistCompatibility}
\begin{proof}
Let $u$ be a word, $X_0$ the bottom space of its chain, with Hermitian form $h$
and the basis of its coordinates, and $k=N(0)=\dim X_0$. For
$\tau\in\mathcal T$ let $\rho_\tau\colon E\to L$ be the composite of the
automorphism $\sigma_\tau$ of $E$ given by the element of $G_\Phi$ of $\tau$,
complex conjugation of $E$, and the datum's embedding $E\to L$. We give the
data of Definition~\ref{asm-global-column-values} in turn.

(1) The tied columns give, at level $0$, representatives $y_c$
($c\in\mathrm{cls}_0$) of the double cosets
$U(X_0)(F)\backslash U(X_0)(\mathbf A_{F,f})/U$, each double coset once (as in
the proof of Lemma~\ref{asm-proviso-bottom-dimension}). The double coset of $1$
has a representative, so $\mathrm{cls}_0\neq\emptyset$. The embeddings
$\rho_\tau$ are those above.

(2) Among the tied columns, the bottom space is standard: $X_0$ has an
$E$-basis $e_1,\dots,e_k$ with $h(e_a,e_{a'})=\delta_{aa'}$. For a field
embedding $\sigma\colon E\to\mathbf C$ let $Y_\sigma$ be the image under
$\sigma$ of the matrix that carries the coordinates of a vector of $X_0$ to its
coordinates in the basis $(e_a)$; it is invertible. As $E$ is a CM field,
$\sigma(\bar a)=\overline{\sigma(a)}$, so for $y=\sum_az_ae_a$,
\[
 \sum_i\big|\big(Y_\sigma\,\sigma(y)\big)_i\big|^2=\sum_a|\sigma(z_a)|^2
 =\operatorname{Re}\sigma\big(h(y,y)\big),
\]
where $\sigma(y)$ is the image under $\sigma$ of the coordinate vector of $y$.
Now let $c\in\mathrm{cls}_0$ and $\tau\in\mathcal T$. Let
$\beta_{c,\tau}\colon X_0\to L^k$ apply $\rho_\tau$ to the coordinates, so that
$\beta_{c,\tau}(ay)=\rho_\tau(a)\beta_{c,\tau}(y)$, and let $Q(c,\tau)$ be the
$\mathcal O_L$-span of $\beta_{c,\tau}(L_u(0,c))$ in $L^k$, with the standard
basis $b_1,\dots,b_k$ and, at an infinite place $v$ of $L$, the metric
$\|z\|_v^2=\sum_i|(Y_{v\circ\rho_\tau}v(z))_i|^2$, whose matrix is
$G_{c,\tau,v}=Y_{v\circ\rho_\tau}^{-1}$. As $L_u(0,c)$ is finitely generated
and spans $X_0$, $Q(c,\tau)$ is finitely generated, and it spans $L^k$, since
$\beta_{c,\tau}$ carries the $i$-th coordinate vector to $b_i$. The display with
$\sigma=v\circ\rho_\tau$ gives
$\|\beta_{c,\tau}(y)\|_v^2=\operatorname{Re}v(\rho_\tau(h(y,y)))$.

(3) Among the tied columns are level data at level $0$: a lattice
$\Lambda_b\subset X_0$ with $h(\Lambda_b,\Lambda_b)\subset\mathcal O_E$, and,
for each $c$, $L_u(0,c)$ is the lattice whose completion at every finite place
is $y_c\Lambda_b$. For a lattice $\Lambda$ let $\Lambda^\vee=\{y\in
X_0:h(\Lambda,y)\subset\mathcal O_E\}$ be its dual, and for $g$ in the unitary
group of $h$ over the finite adeles let $g\Lambda$ be the lattice with
completions $g\Lambda$. As $g$ preserves $h$, $g(\Lambda^\vee)\subset
(g\Lambda)^\vee$. So $L_u(0,c)=y_c\Lambda_b\subset
y_c(\Lambda_b^\vee)\subset(y_c\Lambda_b)^\vee=L_u(0,c)^\vee$. Among the tied
columns is also the discriminant at $T_u$ in degree $0$: an isomorphism of
$L_u(0,c)^\vee/L_u(0,c)$ with $\mathcal O_E/T_u\mathcal O_E$, which has
$T_u^{[E:\mathbf Q]}=T_u^{2f}$ elements. This is (3).

(4) Let $\lambda=\lambda_{u,0}$, $s_\tau=\lambda_{u,0}(\tau)_k$ its
determinant shifts, $\iota\colon L\to\mathbf C$ the embedding of $L$, and
identify $Y_0$ with the tuples $z=(z_\tau)$ of invertible complex $k\times k$
matrices, as in the proof of Lemma~\ref{asm-proviso-bottom-dimension}. Among
the tied columns, the forms in degree $0$ are \emph{rational}: a function $F$
on $\mathrm{cls}_0\times Y_0$ is $\operatorname{emb}_0(y)$ for some $y\in
V_{u,0}$ if and only if $F$ is a form of weight $\lambda$ and level $(U,\chi)$
and for every $c$ there is $P\in V_\lambda$ with
$F(c,z)=\prod_\tau\det(z_\tau)^{s_\tau}\,\iota(P)(z)$ for all $z$. Such a $P$
is determined by these values, as $\iota$ is injective and a complex polynomial
is determined by its values at the invertible tuples. So for $y\in V_{u,0}$ and $c\in\mathrm{cls}_0$
there is a unique $P_c(y)\in V_\lambda$ with
\[
 \operatorname{emb}_0(y)(c,z)=\prod_\tau\det(z_\tau)^{s_\tau}\,
 \iota\big(P_c(y)\big)(z)\qquad\text{for all }z,
\]
and by uniqueness $y\mapsto P_c(y)$ is $L$-linear. Let $G$ be the Gram matrix
of $h$ in the basis of coordinates, $M_{c,\tau}$ the matrix of
$\beta_{c,\tau}$ in that basis and the basis $(b_i)$, and put
$g_{c,\tau}=M_{c,\tau}\,\rho_\tau(G)^{-1}$, which is invertible. Define
\[
 \operatorname{co}(y)(c)=\rho_\lambda\big((g_{c,\tau})_\tau\big)\,P_c(y),
\]
with $\rho_\lambda$ as in Definition~\ref{pd-weight-lattices}(4); this is
$L$-linear in $y$. It is injective: if $\operatorname{co}(y)=\operatorname{co}(y')$,
applying $\rho_\lambda((g_{c,\tau}^{-1})_\tau)$ gives $P_c(y)=P_c(y')$ for all
$c$, so $\operatorname{emb}_0(y)=\operatorname{emb}_0(y')$, and
$\operatorname{emb}_0$ is injective. Its values lie in $V_\lambda$, since
$P_c(y)\in V_\lambda$ and $V_\lambda$ is stable under $\rho_\lambda(g)$ for
$g\in GL_k(L)^{\mathcal T}$. Conversely, let $F_c\in V_\lambda$ for
$c\in\mathrm{cls}_0$, and put
$F'_c=\rho_\lambda((g_{c,\tau}^{-1})_\tau)F_c\in V_\lambda$. As the
stabilizers $\Gamma_c$ are trivial, every element of $U(X_0)(F)$ that lies in
$y_cUy_c^{-1}$ is $1$, and, as in the proof of
Lemma~\ref{asm-proviso-bottom-dimension}, the forms of weight $\lambda$ and
level $(U,\chi)$ are the functions on $\mathrm{cls}_0\times Y_0$ whose value at
each class is the function attached to an element of the complex span of
$V_\lambda$. So
$(c,z)\mapsto\prod_\tau\det(z_\tau)^{s_\tau}\,\iota(F'_c)(z)$ is such a form;
by rationality it is $\operatorname{emb}_0(y)$ for some $y\in V_{u,0}$. By
uniqueness $P_c(y)=F'_c$, and $\operatorname{co}(y)(c)=F_c$. This is (4).

(5) Let $w$ be a finite place of $L$ over the prime $\ell$, and use the
notation of Lemma~\ref{asm-global-link-column-values-criterion}: $\hat\tau=\iota\circ\iota_E\circ\sigma_\tau$, and $\Lambda^\vee_c$
the coordinate vectors of $L_u(0,c)^\vee$; let $\Lambda_c\subset E^k$ be those
of $L_u(0,c)$, a finitely generated $\mathcal O_E$-module that spans $E^k$, and
$\bar Y$ the matrix of complex conjugates of the entries of a matrix $Y$ over $E$.
For each $c\in\mathrm{cls}_0$ choose a local basis $Y_c$ of $\Lambda_c$ at
$\ell$ (Lemma~\ref{asm-global-link-column-values-local-basis-facts}(1)), and put
$\gamma_{c,\tau}=M_{c,\tau}\,\rho_\tau(Y_c)$, which is invertible.

\emph{The $\gamma_{c,\tau}$ are local bases of $Q(c,\tau)$ at $w$.} Fix $c$ and
$\tau$, and for $1\le a\le k$ let $y_a\in X_0$ be the vector whose coordinate
vector is the column $Y_ce_a$. As $\beta_{c,\tau}$ is $\rho_\tau$-semilinear with
matrix $M_{c,\tau}$, $\beta_{c,\tau}(y_a)=\sum_i(\gamma_{c,\tau})_{ia}b_i$. Let
$y\in L_u(0,c)$, with coordinate vector $x$. There is an integer $s$ prime to
$\ell$ with $sY_c^{-1}x$ integral, and $y=\sum_a(Y_c^{-1}x)_ay_a$, so
$\beta_{c,\tau}(y)=\sum_a\rho_\tau\big((Y_c^{-1}x)_a\big)\beta_{c,\tau}(y_a)$
with coefficients in $\mathcal O_{L,w}$, as $\rho_\tau(\mathcal
O_E)\subset\mathcal O_L$ and $s$ is a unit at $w$. Hence $Q(c,\tau)_w$ lies in
the $\mathcal O_{L,w}$-span of the $\beta_{c,\tau}(y_a)$. Conversely, each
column of $Y_c$ has a multiple $tY_ce_a\in\Lambda_c$ with $t$ prime to $\ell$, so
$t\beta_{c,\tau}(y_a)\in Q(c,\tau)$ and $\beta_{c,\tau}(y_a)\in Q(c,\tau)_w$.
So $Q(c,\tau)_w$ is the $\mathcal O_{L,w}$-span of the vectors
$\sum_i(\gamma_{c,\tau})_{ia}b_i$.

\emph{The criterion for these $\gamma_{c,\tau}$.} Fix $c$, write $Y=Y_c$, and
put $B_0=({}^t\bar YG)^{-1}$. As ${}^t\bar YG$ is invertible, $B_0$ is a local
basis of $\Lambda^\vee_c$ at $\ell$, by
Lemma~\ref{asm-global-link-column-values-local-basis-facts}(2) with $\sigma$
complex conjugation. For $\tau\in\mathcal T$ put
$h_\tau=\iota(\gamma_{c,\tau}^{-1}g_{c,\tau})$ and
$N_\tau=\hat\tau(B_0^{-1})=\hat\tau({}^t\bar YG)$. Then ${}^th_\tau N_\tau=1$.
Indeed, $\gamma_{c,\tau}^{-1}g_{c,\tau}=\rho_\tau(Y)^{-1}\rho_\tau(G)^{-1}=\rho_\tau(GY)^{-1}$;
as $E$ is a CM field and the embeddings $\iota_E$ and $\iota$ carry complex
conjugation to the involution of $L$ and to complex conjugation,
$\iota\circ\rho_\tau=\overline{\hat\tau}$ and
$\hat\tau(\bar a)=\overline{\hat\tau(a)}$; and $G$ is Hermitian,
${}^t\bar G=G$. So ${}^t\iota(\rho_\tau(GY))={}^t\overline{\hat\tau(GY)}
=\hat\tau({}^t\bar Y\,{}^t\bar G)=N_\tau$, and ${}^th_\tau=N_\tau^{-1}$.

Let $y\in V_{u,0}$ and put $Q=\rho_\lambda(h)\,\iota(P_c(y))$, where
$h=(h_\tau)_\tau$. As $\rho_\lambda$ is multiplicative,
$\rho_\lambda((\gamma_{c,\tau}^{-1})_\tau)\operatorname{co}(y)(c)
=\rho_\lambda((\gamma_{c,\tau}^{-1}g_{c,\tau})_\tau)P_c(y)$, whose image under
$\iota$ is $Q$; as $\iota$ is injective, it lies in $\Delta_{\lambda,w}$ if and
only if $Q\in\iota(\Delta_{\lambda,w})$. By the formula for $\rho_\lambda$ and
${}^th_\tau N_\tau=1$, for every $z\in Y_0$
\[
 \prod_\tau\det(N_\tau z_\tau)^{s_\tau}\,Q\big((N_\tau z_\tau)_\tau\big)
 =\prod_\tau\det({}^th_\tau N_\tau z_\tau)^{s_\tau}\,
 \iota(P_c(y))\big(({}^th_\tau N_\tau z_\tau)_\tau\big)
 =\operatorname{emb}_0(y)(c,z).\tag{$*$}
\]
Suppose $y$ is integral at $(w,c)$
(Lemma~\ref{asm-global-link-column-values-criterion}). With $B=B_0$ this gives
$P\in\iota(\Delta_{\lambda,w})$ with
$\operatorname{emb}_0(y)(c,z)=\prod_\tau\det(N_\tau z_\tau)^{s_\tau}P((N_\tau
z_\tau)_\tau)$ for all $z$. By $(*)$, $P$ and $Q$ take the same values at all
tuples $(N_\tau z_\tau)_\tau$ of invertible matrices, which are all the
invertible tuples, so $P=Q$ and $Q\in\iota(\Delta_{\lambda,w})$. Conversely,
let $Q\in\iota(\Delta_{\lambda,w})$ and let $B$ be a local basis of
$\Lambda^\vee_c$ at $\ell$. Put $K_\tau=\hat\tau(B_0^{-1}B)$, so that
$N_\tau=K_\tau\hat\tau(B)^{-1}$. By
Lemma~\ref{asm-global-link-column-values-local-basis-facts}(3), applied to $B$
and $B_0$ in both orders with $\phi=\iota_E\circ\sigma_\tau$, the entries of $K_\tau$ and of
$K_\tau^{-1}$ lie in $\iota(\mathcal O_{L,w})$. Put
$P'(X)=\prod_\tau\det(K_\tau)^{s_\tau}Q((K_\tau X_\tau)_\tau)$, which is
$\rho_\lambda(({}^tK_\tau)_\tau)Q$; it lies in $\iota(\Delta_{\lambda,w})$ by
Lemma~\ref{asm-global-link-column-values-weyl-stable} with $K=\mathbf C$ and
$R=\iota(\mathcal O_{L,w})$. By $(*)$,
$\operatorname{emb}_0(y)(c,z)=\prod_\tau\det(\hat\tau(B)^{-1}z_\tau)^{s_\tau}
P'((\hat\tau(B)^{-1}z_\tau)_\tau)$ for all $z$. So $y$ is integral at $(w,c)$.
Thus $y$ is integral at $(w,c)$ if and only if
$\rho_\lambda((\gamma_{c,\tau}^{-1})_\tau)\operatorname{co}(y)(c)\in\Delta_{\lambda,w}$,
and Lemma~\ref{asm-global-link-column-values-criterion} gives the equivalence of
(5) for these $\gamma_{c,\tau}$.

\emph{Other local bases.} Let $\gamma'_{c,\tau}\in GL_k(L)$ be such that
$Q(c,\tau)_w$ is the $\mathcal O_{L,w}$-span of the vectors
$\sum_i(\gamma'_{c,\tau})_{ia}b_i$. Each vector
$\sum_i(\gamma_{c,\tau})_{ia}b_i$ is an $\mathcal O_{L,w}$-combination of these,
and conversely; as the $b_i$ are linearly independent, $\gamma_{c,\tau}=\gamma'_{c,\tau}\kappa$
and $\gamma'_{c,\tau}=\gamma_{c,\tau}\kappa'$ with $\kappa,\kappa'$ over $\mathcal O_{L,w}$, and
$\kappa'\kappa=1$ as $\gamma_{c,\tau}$ is invertible. So
$\gamma'^{-1}_{c,\tau}=\kappa\,\gamma_{c,\tau}^{-1}$ with $\kappa\in GL_k(\mathcal
O_{L,w})$, and
$\rho_\lambda((\gamma'^{-1}_{c,\tau})_\tau)=\rho_\lambda((\kappa_{c,\tau})_\tau)
\rho_\lambda((\gamma_{c,\tau}^{-1})_\tau)$. By
Lemma~\ref{asm-global-link-column-values-weyl-stable} with $K=L$ and
$R=\mathcal O_{L,w}$, $\rho_\lambda((\kappa_{c,\tau})_\tau)$ and its inverse map
$\Delta_{\lambda,w}$ into itself. So the condition of (5) is the same for
$\gamma'$ as for $\gamma$, and (5) holds.
% Lean for (5): finCo_of (TransposeBaseMatrixInfinity.lean): Y c from LocB.exists_isLocBasis for
% clsCrd (bwB_clsCrd_lat), gamma = Mb B c tau * (Y c)^rho_tau, locBasis_gamma (NatInvMemIntW.lean:
% sum_gamma, rhoC_mem_intW, natInv_mem_intW, col_mem_of_isLocBasis), Formalism.L2.G.Cmp.locWCrit_proved
% (Lemma asm-global-link-column-values-criterion) and clause_iff: B0 = (BaseBlockInverse Y)^(-1) with
% isLocBasis_dual, h_T_mul (transpose_base_matrix_infinity, embLL_rhoC, embT_sigD, gramN_sig), hfn
% (fnOf_mul, actAll_eq_subP), fnOf_injective with isUnit_base_inverse_infinity for the forward
% direction, and basis_change (change_mem, bwG_intSpan_allSubP, local_GL_stabilises_lattice) for the
% backward one; the identification of Delta_{lambda,w} with IntSpan is mem_intSpanAt_iff_map with
% map_actAll and coT_apply, actAll_comp. Then colTie_of_colTieW with mem_iff_of_locBasis
% (exists_change, exists_change_inv, actAll_mem_intSpanAtG) passes from these gamma to all local bases.

(6) Let $v$ be an infinite place of $L$, with embedding also written $v$, and
for $\tau\in\mathcal T$ let $G^v_\tau$ be the image of $G$ under the embedding
$v\circ\iota_E\circ\sigma_\tau$ of $E$, where $\iota_E\colon E\to L$ is the
datum's embedding. Call a tuple $o=(o_\tau)$ of invertible complex $k\times k$
matrices \emph{orthonormal at} $v$ if $o_\tau^*G^v_\tau o_\tau=\pm1$ for every
$\tau$, and for a complex polynomial $P$ let $P\circ o$ be the polynomial
$P\big((o_\tau X_\tau)_\tau\big)$. Among the tied columns is the \emph{norm tie}
in degree $0$: for every $y\in V_{u,0}$ and every $o$ orthonormal at $v$,
\[
 \|y\|_v^2=\sum_{c\in\mathrm{cls}_0}\operatorname{term}_c(y,o),\qquad
 \operatorname{term}_c(y,o)=\frac1{|\Gamma_c|}\prod_\tau|\det o_\tau|^{2s_\tau}
 \operatorname{Re}\big\{v(P_c(y))\circ o,v(P_c(y))\circ o\big\}_\lambda .
\]
Here $|\Gamma_c|=1$. For $c\in\mathrm{cls}_0$ and $\tau\in\mathcal T$ put
$Z_{c,\tau}=G_{c,\tau,v}^{-1}\,v(M_{c,\tau})\,v(\rho_\tau(G))^{-1}$; the
tuple $o^{(c)}=({}^tZ_{c,\tau})_\tau$ is orthonormal at $v$, by a computation
from the identity $\|\beta_{c,\tau}(y)\|_v^2=\operatorname{Re}v(\rho_\tau(h(y,y)))$
of (2) and $\overline{v\circ\iota_E\circ\sigma_\tau}=v\circ\rho_\tau$.

The term of a class does not depend on the orthonormal tuple. Indeed, given $c$
and $y$, by (4) there is $y'\in V_{u,0}$ with $P_c(y')=P_c(y)$ and
$P_{c'}(y')=0$ for $c'\neq c$. The norm tie for $y'$ has only the term of
$c$, since $\{0,0\}_\lambda=0$, so $\operatorname{term}_c(y,o)=\|y'\|_v^2
=\operatorname{term}_c(y,o')$ for any two tuples $o,o'$ orthonormal at $v$.

Fix $c_0\in\mathrm{cls}_0$. The norm tie with $o^{(c_0)}$ and the independence
just shown give $\|y\|_v^2=\sum_c\operatorname{term}_c(y,o^{(c)})$. By
Definition~\ref{pd-weight-lattices}(4),
$\rho_\lambda(Z_c)P=\prod_\tau\det(Z_{c,\tau})^{s_\tau}\,P\circ o^{(c)}$, so,
as $\{aP,aP\}_\lambda=|a|^2\{P,P\}_\lambda$,
\[
 \operatorname{term}_c(y,o^{(c)})=\operatorname{Re}\big\{\rho_\lambda(Z_c)\,
 v(P_c(y)),\rho_\lambda(Z_c)\,v(P_c(y))\big\}_\lambda .
\]
Finally $G_{c,\tau,v}^{-1}\,v(g_{c,\tau})=Z_{c,\tau}$, so
$\rho_\lambda\big((G_{c,\tau,v}^{-1})_\tau\big)\,v(\operatorname{co}(y)(c))
=\rho_\lambda(Z_c)\,v(P_c(y))$, and the sum is the formula of (6).
% Lean, for (1) and (2): classes_level_zero_nonempty ((F.herm.base.lev0 u).y_bij.2 at the class of
% 1); colBase_of_sym = ColBase_exists D u (posFrame_of_std D (F.herm.std.1 0 _)) (rhoC D), with
% base_exists (mkQ with X v = (Y v)^(-1), bcFun, mkQ_fg, mkQ_span, mkQ_norm), hT_self_of_on and
% host_emb_star; rhoC D tau = D.embE . sigD D . sigmaOf (eTau tau), with sigD D = star.
% (3) latT_of_sym (ClassesLevelZeroNonempty.lean): class_lattice_integral from F.herm.base.lev0
% (class_lattice_eq_global with eq_glob_of_hat: L_u(0,c) = globLat (y c) Lb; Lb_le_dl_lv from
% Lb_int; glob_dl_le for unitary y c), and the index from (F.herm.atT u).disc 0 c with
% card_ringOfIntegers_quot_natCast_eq_pow; about 840 lines.
% (4) co = coT (F.forms u).to0 (gC B) (ActAllCancelInv.lean), gC B c tau = Mb B c tau *
% (gramN^rho_tau)^(-1) (GramNDetNe.lean, Mb in SigDEqStar.lean); P_c(y) = pol T y c, co0 T y c
% (CoefficientMapSurjective.lean), from T.val0.rat and rd_unique (fnOf_injective); coT_injective
% (coefficient_map_injective from emb_inj, actAll_inv_cancel); coT_range: coefficient_map_mem_span
% with actAll_mem_TSpan, and coefficient_map_surjective (stabI_range and the trivial stabilizers,
% fiC_eq_pi with y_bij, valC_eq, map_mem_CSpan, then T.val0.rat), actAll_cancel_inv. About 320
% theorems and 4,600 lines, overlapping (5) and (6).
% (6) inf_coT (WeightLevelZeroEqOne.lean): normTie_at (the field norm0 : NormTie of Forms0, with
% termN), oc B v c = (Zc)^T with oc_on (Zc_on from gram_Y, hB, rhoC_conj), termN_indep (exists_single,
% termN_single, normTie_at), termN_oc (weight_level_zero_eq_one, actAll_eq_subP, nfAll_smul_self_re,
% norm_sq_prod_zpow) and coT_term (R_gC: X_v^(-1) gC = Zc). About 360 theorems and 5,100 lines,
% mostly in the norm tie's supporting lemmas.
% Lean: colValOfTiedS_holds = colValOfTiedS_of_fin finCo_of (ColValOfTiedSOfFin.lean), through
% colValOfTiedS_of_rest' (ColValOfTiedSOfRest.lean: colBase_of_sym, which is ColBase_exists with
% posFrame_of_std from an orthonormal basis of the bottom space, F.herm.std) and colRest_of_fin:
% the coordinates coT with coT_injective and coT_range (this is where the trivial stabilizers
% enter), FinCo from finCo_of (local bases), the norms inf_coT, and classes_level_zero_nonempty.
% Size: proof closure of colValOfTiedS_holds, about 340 theorems and 4,800 lines in 61 files.
\end{proof}

\begin{lemma}\label{asm-global-link-local-facts-free}
Let $D$ be a definite-side datum (Definition~\ref{th-definite-datum}) and $u$ a
word that carries tied forms (the data that tied columns, as in the proof of
Lemma~\ref{asm-proviso-bottom-dimension}, give at $u$), and suppose that every
class $c\in\mathrm{cls}_0$ of the chain of $u$ has trivial stabilizer
$\Gamma_c$. Then $u$ has a free basis at $T_u$
(Definition~\ref{asm-level-local-facts}(1)).
\end{lemma}
\leanref{Formalism/LocalToGlobal/GlobalTheory/GL6/DiagonalSlopes/SlopeLineRamifiedFormula.lean}{Formalism.L2.GlobalTheoryTheory.LocalLatticeBases.S0.freeT_of_tied}
\leanref{Formalism/LocalToGlobal/GlobalTheory/GL6/DiagonalSlopes/SlopeLineRamifiedFormula.lean}{Formalism.L2.GlobalTheoryTheory.LocalLatticeBases.S0.freeT_of_model}
\leanref{Formalism/LocalToGlobal/GlobalTheory/GL6/RamifiedRankEqFinrank.lean}{Formalism.L2.GlobalTheoryTheory.LocalLatticeBases.FreeRankModel}
\begin{proof}
The proof has three steps. First, $E$ splits at the places over $T_u$: for every
$v\in V_u$, with $F_v$ the completion of $F$ at the place of $F$ attached to $v$,
there is an idempotent $e$ among the integral elements of $E\otimes_FF_v$ with
$\bar e=1-e$; it is built from the element $\pi_u$ with $T_u=\pi_u\bar\pi_u$ of
Definition~\ref{th-definite-datum}(4). Second, this splitting gives
\emph{diamond elements} at the places over $T_u$. Third, with the trivial
stabilizers, the diamond elements give a \emph{free rank model}: a finite group
$H$ with isomorphisms $e_v\colon\mathbf F_v^\times\to H$, a finite set $X$ with
a free action of $H^{V_u}$, and a linear isomorphism $\Theta\colon M_u\to\mathbf
C^X$ with $\Theta(\rho_v(a)\varphi)(x)=\Theta(\varphi)(e_v(a)_vx)$. Letting
$\prod_v\mathbf F_v^\times$ act on $X$ through the $e_v$ gives a free basis at
$T_u$.
\openstep{The first three steps are not yet written out: the splitting of $E$ at
the places over $T_u$, the diamond elements, and the free rank model built from
them.}
% Lean: S0.freeT_of_tied Tf hS = freeT_of_model (free_rank_model_of_split Tf hS (splitAtT_of_tied
% Tf.atT)), with free_rank_model_of_split = free_rank_model_of_diamonds Tf hS (diamond_elements_of_split
% Tf hs). Sizes: splitAtT_of_tied about 12 theorems and 330 lines; diamond_elements_of_split about 37
% theorems and 390 lines; free_rank_model_of_diamonds about 94 theorems and 1,390 lines; total about
% 160 theorems and 2,430 lines.
\end{proof}

\begin{lemma}\label{asm-global-link-local-facts-slope}
Let $D$ be a definite-side datum (Definition~\ref{th-definite-datum}) with
coefficient field $L$ and $u$ a word that carries tied forms (the data that tied
columns, as in the proof of Lemma~\ref{asm-proviso-bottom-dimension}, give at
$u$). Suppose that every class $c\in\mathrm{cls}_0$ of the chain of $u$ has
trivial stabilizer $\Gamma_c$, that the setting of the level primes holds at
$u$ (Definition~\ref{asm-level-setting}(1)), and that the character $\chi$ of
the level of the forms in degree $0$ has finite order with values in $L$. Then
the slope line holds at $u$ (Definition~\ref{asm-level-local-facts}(2)).
\end{lemma}
\leanref{Formalism/LocalToGlobal/GlobalTheory/GL6/DiagonalSlopes/SlopeLineRamifiedFormula.lean}{Formalism.L2.GlobalTheoryTheory.LocalLatticeBases.S0.slopeLine_of_tied}
\leanref{Formalism/LocalToGlobal/GlobalTheory/GL6/DiagonalSlopes/SlopeFormulaOfCharacterFinite.lean}{Formalism.L2.GlobalTheoryTheory.LocalLatticeBases.S0.slope_formula_of_character_finite}
\begin{proof}
The proof has two steps. First, as $\chi$ has finite order and the stabilizers
are trivial, the diamond operators at the places $v\in V_u$, transported to
$V_{u,0}$, are isometries of the metric of $\Lambda_{u,0}$ at every infinite
place of $L$.
Second, with the setting of the level primes, this isometry gives the slope
line.
\openstep{The two steps are not yet written out: the isometry of the diamond
operators at the infinite places, and the slope line from it.}
% Lean: S0.slopeLine_of_tied hS Tf hC hA = slopeLine_iff_ramifiedSlopeFormula.2
% (slope_formula_of_character_finite hS Tf hC hA), which is slope_formula_of_isometry hS Tf hA
% (diamond_isometry_at_infinity Tf hC hA : DiamondIsometry D u). Sizes: diamond_isometry_at_infinity
% about 121 theorems and 1,680 lines; slope_formula_of_isometry about 309 theorems and 4,930 lines;
% total about 354 theorems and 5,740 lines.
\end{proof}

\begin{lemma}\label{asm-global-link-local-facts-integral}
Let $D$, $u$ and $\chi$ be as in Lemma~\ref{asm-global-link-local-facts-slope},
with its hypotheses. Then Hecke integrality at level zero holds at $u$
(Definition~\ref{asm-level-local-facts}(3)).
\end{lemma}
\leanref{Formalism/PeriodIntegrals/OrbitalBoundsConclusion/IntegralAtT0Lattice.lean}{PI.IntegralAtBasePointLattice.integral_at_t0_lattice}
\leanref{Formalism/PeriodIntegrals/UnitaryGroupMassFormula/IntegralityAtPrimeLattices/IntegralityAtFormsOfLattices.lean}{PI.UnitaryGroupMassFormula.IAT.integrality_at_forms_of_lattices}
\begin{proof}
The proof has three steps. First, the lattices $L_u(0,c)$ are local at the places
over $T_u$: for places $v\neq v'$ of $V_u$, a change of representative
$\gamma y_{c'}=y_c\eta\kappa$ with $\eta$ trivial away from the place of $v$ gives
equal localizations of $\gamma L_u(0,c')$ and $L_u(0,c)$ at $v'$. Second, as
$\chi$ has finite order, the Hecke operators at level zero are rational, and the
locality gives the local lattice statements for the operators $U_v$ and $D_v$.
Third, these give the integrality of the increments of $U_v$ and $D_v$, which is
Hecke integrality at level zero.
\openstep{The three steps are not yet written out: the locality of the lattices
$L_u(0,c)$, the rationality of the Hecke operators with the local lattice
statements, and the integrality of the increments.}
% Lean: integral_at_t0_lattice = IAT.integrality_at_forms_of_away (lattices_away_from_place: AwayLat0)
% = integrality_at_forms_of_away_up with dnLatLoc0_of_away, then integrality_at_forms_of_lattices
% (rationality_of_finite_character: HeckeRationalLevelZero; upLatLoc0_of: UpLatLoc0) and
% integrality_at_forms_of_increments (upIncInt0_of, dnIncInt0_of). Sizes: lattices_away_from_place
% about 163 theorems and 2,680 lines; dnLatLoc0_of_away 144 and 2,400; rationality_of_finite_character
% 67 and 1,200; upLatLoc0_of 13 and 360; integrality_at_forms_of_increments 21 and 310; total about 346
% theorems and 6,130 lines (overlapping).
\end{proof}

\begin{lemma}\label{asm-global-link-local-facts-ramified}
Let $D$, $u$ and $\chi$ be as in Lemma~\ref{asm-global-link-local-facts-slope},
with its hypotheses. Then the ramified level zero holds at $u$
(Definition~\ref{asm-level-local-facts}(4)).
\end{lemma}
\leanref{Formalism/LocalToGlobal/GlobalTheory/UnramifiedTrace/HeckeRamifiedZeroAll.lean}{Formalism.L2.GlobalTheoryTheory.UnramifiedTrace.hecke_ramified_zero_all}
\leanref{Formalism/LocalToGlobal/GlobalTheory/UnramifiedDiagonal/HeckeRamifiedZeroOfMatchingsAndDu.lean}{Formalism.L2.GlobalTheoryTheory.UnramifiedDiagonal.hecke_ramified_zero_of_matchings_and_du}
\begin{proof}
The proof has two inputs. First, the Hecke operator $U_v$ and the down operator
$D_v$ match the diamond operators at the same place $v$. Second, on the forms in
degree
$0$,
\[
 D_vU_v\varphi=T_u^{N(0)-1}\varphi-T_u^{N(0)-2}\operatorname{Nm}_v\varphi
 +T_u^{N(0)-2}\operatorname{tr}_v\varphi,
\]
where $\operatorname{Nm}_v$ is the sum over $a\in\mathbf F_v^\times$ of the
\emph{diamond operators} at $v$, given by the dilations of the forms, and
$\operatorname{tr}_v=T_u^{1-N(0)}E_vU_v$, with $E_v$ defined as $D_v$ but with
the sum over all nonzero linear forms $g$. With the matchings, this formula
gives the ramified level zero.
\openstep{The two inputs and their combination are not yet written out: the
matchings of the Hecke, down and diamond operators at one place, and the formula
for $D_vU_v$ on the forms in degree $0$.}
% Lean: hecke_ramified_zero_all = UnramifiedDiagonal.hecke_ramified_zero_of_matchings_and_du
% SameDiaZero.diamond_same_place_matchings dUAll_holds, through hecke_ramified_zero_of_same_place
% (CrossDiaZero.upMatch, dnMatchAll_holds, diamond_same_place_of_matchings, trpDia0_of). The formula is
% DUFormulaLevelZero. Sizes: diamond_same_place_matchings about 254 theorems and 4,910 lines; dUAll_holds
% about 326 theorems and 5,630 lines; total about 441 theorems and 9,040 lines (overlapping).
\end{proof}

\begin{lemma}\label{asm-global-link-local-facts-branch}
Let $D$, $u$ and $\chi$ be as in Lemma~\ref{asm-global-link-local-facts-slope},
with its hypotheses. Then the second branch inclusion holds at $u$
(Definition~\ref{asm-level-local-facts}(5)).
\end{lemma}
\leanref{Formalism/LocalToGlobal/GlobalTheory/CoordinateLengths/SecondBranchAllMembers.lean}{Formalism.L2.GlobalTheoryTheory.LemB.second_branch_all_members}
\leanref{Formalism/LocalToGlobal/GlobalTheory/CoordinateLengths/SecondBranchOfCoordinatesAndCount.lean}{Formalism.L2.GlobalTheoryTheory.LemB.second_branch_of_coordinates_and_count}
\begin{proof}
The proof has two inputs. First, the coordinates of the second branch lattices
at $T_u$, taken in local bases of the lattices $L_u(0,c)$, are surjective; this
rests on the rationality of the forms in degree $0$. Second, the count
$\dim M_u=|\mathrm{cls}_0|\prod_\tau\operatorname{rk}\Delta(\lambda_{u,0}(\tau))$.
The surjectivity gives the inclusion $T_u^bS_w\subset B_w(v)$ and the length of
the quotient as a sum over the classes, and the count turns that sum into the
length of Definition~\ref{asm-level-local-facts}(5).
\openstep{The two inputs and the length computation are not yet written out: the
surjectivity of the second branch coordinates, the count of $\dim M_u$, and the
length of $B_w(v)/T_u^bS_w$.}
% Lean: second_branch_all_members = second_branch_of_rationality ratInC_holds ratGL_holds =
% second_branch_of_surjectivity (ontoAll_of) duCount_holds, through
% second_branch_of_coordinates_and_count (second_branch_coordinates_of_surjectivity, lenLen2_of_coordOnto
% and lemmaB_second_branch_of_length_formula). Sizes: ontoAll_of about 78 theorems and 1,140 lines with
% ratInC_holds (7, 70) and ratGL_holds (34, 620); duCount_holds 91 and 1,220;
% second_branch_of_coordinates_and_count 198 and 3,660; total about 407 theorems and 7,150 lines.
\end{proof}

\begin{lemma}\label{asm-global-link-local-facts}
Let $D$ be a definite-side datum (Definition~\ref{th-definite-datum}) with
coefficient field $L$, and let $u$ be a word of $D$ that carries tied forms:
the data that tied columns, as in the proof of
Lemma~\ref{asm-proviso-bottom-dimension}, give at the word $u$. Let $\chi$ be
the character of the level of its forms in degree $0$. Then $u$ has a free
basis at $T_u$ (Definition~\ref{asm-level-local-facts}). Suppose moreover that
the setting of the level primes holds at $u$
(Definition~\ref{asm-level-setting}(1)) and that $\chi$ has finite order with
values in $L$, as in Lemma~\ref{asm-global-second-layer}. Then the slope line
holds at $u$. If also every class $c\in\mathrm{cls}_0$ of the chain of $u$ has
trivial stabilizer $\Gamma_c$, then the other three local facts of
Definition~\ref{asm-level-local-facts} hold at $u$: Hecke integrality at level
zero, the ramified level zero and the second branch inclusion.
\end{lemma}
\leanref{Formalism/SharpRestriction/GlobalSchurBounds/LinkTied.lean}{Formalism.MainNotions.freeTT_holds}
\leanref{Formalism/SharpRestriction/GlobalSchurBounds/LinkTied.lean}{Formalism.MainNotions.slopeLineT_holds}
\leanref{Formalism/PeriodIntegrals/OrbitalBoundsConclusion/IntegralAtT0Lattice.lean}{PI.IntegralAtBasePointLattice.integral_at_t0_lattice}
\leanref{Formalism/LocalToGlobal/GlobalTheory/UnramifiedTrace/HeckeRamifiedZeroAll.lean}{Formalism.L2.GlobalTheoryTheory.UnramifiedTrace.hecke_ramified_zero_all}
\leanref{Formalism/LocalToGlobal/GlobalTheory/CoordinateLengths/SecondBranchAllMembers.lean}{Formalism.L2.GlobalTheoryTheory.LemB.second_branch_all_members}
\leanref{Formalism/SharpRestriction/GlobalSchurBounds/IntegralityOfRationalityLocal.lean}{Formalism.MainNotions.RT.FormsTwistCompatibilityFamily}
\leanref{Formalism/SharpRestriction/GlobalSchurBounds/LemmaBSecondBranchOfLength.lean}{Formalism.MainNotions.RT.ConcreteFormsLevelZero}
\begin{proof}
Let the tied forms at $u$ be given, and let $X_0$, $U$, $\chi$ and the
representatives $y_c$ ($c\in\mathrm{cls}_0$) be their data in degree $0$, as in
the proof of Lemma~\ref{asm-proviso-bottom-dimension}.

\emph{Trivial stabilizers.} The tied forms give, for each $c$, an injective
homomorphism from $\Gamma_c$ to $U(X_0)(F)\cap y_cUy_c^{-1}$, and, at the level
primes $g_0\neq g_1$, local data under which the components of $U$ at the places
over $g_i$ lie in groups that are pro-$g_i$. As $X_0$ is definite and $U$
compact, an element of $U(X_0)(F)\cap y_cUy_c^{-1}$ has finite order, which is a
power of $g_0$ and a power of $g_1$, as in the proof of
Lemma~\ref{pd-gain-coherence-consequences}(1); so it is $1$. Hence every
$\Gamma_c$ is trivial, whether or not this is assumed.

\emph{The five facts.} With the trivial stabilizers, the free basis at $T_u$ is
Lemma~\ref{asm-global-link-local-facts-free}. When moreover the setting of the
level primes holds at $u$ and $\chi$ has finite order with values in $L$, the
slope line, Hecke integrality at level zero, the ramified level zero and the
second branch inclusion are
Lemmas~\ref{asm-global-link-local-facts-slope},
\ref{asm-global-link-local-facts-integral},
\ref{asm-global-link-local-facts-ramified}
and~\ref{asm-global-link-local-facts-branch}.
% Lean, trivial stabilizers: S0.admissible_stabilisers_of_compatible_forms (cls0.stabI_inj,
% cls0.stabI_range and adms_triv through adms_core and Formalism.L2.G.Triv.gen_triv, with the local
% data Tf.gl and Tf.glev0 at the level primes); it gives admS from the tied forms alone.
% (1) S0.freeT_of_tied = freeT_of_model (free_rank_model_of_split Tf hS (splitAtT_of_tied Tf.atT)).
% (2) S0.slopeLine_of_tied = slopeLine_iff_ramifiedSlopeFormula.2 (slope_formula_of_character_finite).
% (3) integrality_at_forms_of_away = integrality_at_forms_of_away_up with lattices_away_from_place
% (AwayLat0: for v' != v, gamma y_c' = y_c eta kappa with eta trivial away from the place of v
% gives equal localizations of gamma L(c') and L(c) at v') and dnLatLoc0_of_away, then
% integrality_at_forms_of_lattices (rationality_of_finite_character for HeckeRationalLevelZero,
% upLatLoc0_of, and integrality_at_forms_of_increments with upIncInt0_of and dnIncInt0_of).
% (4) hecke_ramified_zero_of_matchings_and_du: diamond_same_place_matchings (UpDiaSame0F,
% DnDiaSame0F, EDiaSame0F) and dUAll_holds (DUFormulaLevelZero: D_v U_v phi = T^(N0-1) phi -
% T^(N0-2) Nm_v phi + T^(N0-2) trp_v phi on the forms of degree 0), through
% hecke_ramified_zero_of_same_place (CrossDiaZero.upMatch, dnMatchAll_holds,
% diamond_same_place_of_matchings, trpDia0_of).
% (5) second_branch_of_rationality ratInC_holds ratGL_holds = second_branch_of_surjectivity
% (ontoAll_of: OntoAll, the coordinates at T_u onto for local bases of L(0,c) and of the lattices
% Lk at v) with duCount_holds (DuCount: duS u = |cls 0| prod_tau rank delSq(lam tau)), through
% second_branch_of_coordinates_and_count and lenLen2_of_coordOnto.
% Lean: the free basis is LocalLatticeBases.freeT_0 T.to0 (= S0.freeT_tied: free_rank_model_of_split,
% splitAtT_of_tied), about 160 theorems and 2,500 lines in 32 files; the slope line is
% LocalLatticeBases.slopeLine_0 hS T.to0 (chiFin0_to0S T hC) (= S0.slopeLine_tied:
% slope_formula_of_character_finite), about 350 theorems and 5,900 lines in 53 files; Hecke
% integrality is integral_at_t0_lattice (IAT.integrality_at_forms_of_away with
% lattices_away_from_place), about 350 theorems and 6,200 lines in 78 files; the ramified level
% zero is hecke_ramified_zero_all (UnramifiedDiagonal.hecke_ramified_zero_of_matchings_and_du with
% SameDiaZero.diamond_same_place_matchings and dUAll_holds), about 500 theorems and 10,000 lines in
% 101 files; the second branch inclusion is second_branch_all_members (second_branch_of_rationality
% ratInC_holds ratGL_holds), about 470 theorems and 7,900 lines in 95 files. The closures overlap.
% Tied forms with a character of finite order are FormsTwistCompatibilityFamily D u, and with
% trivial stabilizers (admS) they are ConcreteFormsLevelZero D u.
\end{proof}

\subsection{Weight data and the bounds for the count}\label{subsection-asm-weight-data}

\begin{definition}\label{asm-global-link-count-bounds-weights}
Let $D$ be a definite-side datum (Definition~\ref{th-definite-datum}) with
ranks $k,m$, $f=2n$ and $r=k+m-1$, Harish-Chandra parameters $\lambda(\tau)$
with entries $v_i(\tau)$ and $b_j(\tau)$, and parameters $\lambda_u(\tau)$ with
entries $u_i(\tau)$ for the words $u$ (Definition~\ref{th-hc-parameter}); let
$v$ be the half root discriminant of its host field and $g_0,g_1$ its level
primes, and call $\tau\in\mathcal T$ \emph{top} if it is a top embedding
(Definition~\ref{pd-hecke-at-T}). Let $x>1$ and $\varepsilon>0$, and put
$L=x^\varepsilon$, $b=\exp(x^\varepsilon)$ and $a=xb$. For real numbers
$C_0$, $\Delta$ and $w_j(\tau)$ ($1\le j\le m$, $\tau\in\mathcal T$), $D$ has
\emph{weight data} $(C_0,\Delta,w)$ at $(x,\varepsilon)$ if:
\begin{enumerate}
\item $\pi\le v$, $\log v\le L$, $f\le x^{C_0}$, $m^2L\le k\le x^{C_0}$, and
  $(kb)^8\le g_i$ and $\log g_i\le C_0L$ for $i=0,1$;
\item $\Delta\ge1$; $0\le v_i(\tau)\le4k\Delta$ for $1\le i\le k-1$, and
  $v_i(\tau)-v_{i+1}(\tau)=4\Delta$ for $1\le i\le k-2$;
\item for every word $u$ and every $\tau$: $u_i(\tau)$ is
  $v_{i-1}(\tau)-\frac12$ or $v_{i-1}(\tau)-\frac12-\Delta$ for
  $2\le i\le k-1$; if $\tau$ is top, $u_1(\tau)\le2a$ and $u_k(\tau)\ge-a$;
  otherwise $4k\Delta\le u_1(\tau)-v_1(\tau)\le8k\Delta$ and
  $u_k(\tau)=v_{k-1}(\tau)-\frac12$;
\item for every $\tau$: $0<w_j(\tau)$, $b_j(\tau)<w_j(\tau)$ and
  $|w_j(\tau)-b_j(\tau)|\le b_j(\tau)/x$ for $1\le j\le m$;
  $w_j(\tau)<b_{j+1}(\tau)$ and $r^{19}f^5L/C_0\le b_{j+1}(\tau)-b_j(\tau)$ for
  $1\le j<m$; $4k\Delta\le b_1(\tau)/x$; and $w_m(\tau)\le2b$.
\end{enumerate}
\end{definition}
\leanref{Formalism/SharpRestriction/GlobalSchurBounds/StandardLattice.lean}{Formalism.MainNotions.StandardLattice.DataB}
\leanref{Formalism/SharpRestriction/GlobalSchurBounds/StandardLattice.lean}{Formalism.MainNotions.StandardLattice.GapA}

\begin{definition}\label{asm-global-link-count-bounds-weights-root-ratios}
Let $N\ge1$ be an integer and $U$ a subgroup of $(\mathbf Z/N\mathbf Z)^\times$. For a divisor $f$
of $N$ let $K_N(f)\subset(\mathbf Z/N\mathbf Z)^\times$ be the subgroup of the classes
$u\equiv1\pmod f$, the kernel of reduction to $(\mathbf Z/f\mathbf Z)^\times$. The \emph{outside
ratio} of $U$ at $f$ is
\[
 \rho_U(f)=\frac{|K_N(f)\setminus U|\,\varphi(f)}{\varphi(N)},
\]
and the \emph{ramification sum} of $U$ is
\[
 S(U)=\sum_{m\mid N}\Lambda(m)\,\rho_U(N/m),
\]
where $\varphi$ is Euler's function and $\Lambda$ is von Mangoldt's function. For
$u\in\mathbf Z/N\mathbf Z$ we write $n_u$ for the additive order of $u-1$ in $\mathbf Z/N\mathbf Z$,
and $\bar u$ for the representative of $u$ in $\{0,\dots,N-1\}$.
\end{definition}
\leanref{Formalism/SharpRestriction/RamificationFrobenius/TotientGaloisClasses/InertiaGaloisSymmetry/RdArith.lean}{Formalism.HC.Der.Rd.KK}
\leanref{Formalism/SharpRestriction/RamificationFrobenius/TotientGaloisClasses/InertiaGaloisSymmetry/RdArith.lean}{Formalism.HC.Der.Rd.rho}
\leanref{Formalism/SharpRestriction/RamificationFrobenius/TotientGaloisClasses/InertiaGaloisSymmetry/RdBound.lean}{Formalism.HC.Der.Rd.SS}

\begin{lemma}\label{asm-global-link-count-bounds-weights-root-ratio-facts}
Let $N\ge1$ be an integer, $U$ a subgroup of $(\mathbf Z/N\mathbf Z)^\times$ and $f$ a divisor of
$N$, with $K_N(f)$ and $\rho_U(f)$ as in
Definition~\ref{asm-global-link-count-bounds-weights-root-ratios}. Then:
\begin{enumerate}
\item $|K_N(f)|\,\varphi(f)=\varphi(N)$;
\item $\rho_U(f)=1-|K_N(f)\cap U|\,\varphi(f)/\varphi(N)$;
\item if some element of $K_N(f)$ is not in $U$, then $\rho_U(f)\ge\frac12$;
\item if $g$ is a divisor of $N$ with $K_N(f)\subset U$ and $K_N(g)\subset U$, then
  $K_N(\gcd(f,g))\subset U$.
\end{enumerate}
\end{lemma}
\leanref{Formalism/SharpRestriction/RamificationFrobenius/TotientGaloisClasses/InertiaGaloisSymmetry/RdArith.lean}{Formalism.HC.Der.Rd.card_KK}
\leanref{Formalism/SharpRestriction/RamificationFrobenius/TotientGaloisClasses/InertiaGaloisSymmetry/RdArith.lean}{Formalism.HC.Der.Rd.rho_eq}
\leanref{Formalism/SharpRestriction/RamificationFrobenius/TotientGaloisClasses/InertiaGaloisSymmetry/RdArith.lean}{Formalism.HC.Der.Rd.half_le_rho}
\leanref{Formalism/SharpRestriction/RamificationFrobenius/TotientGaloisClasses/InertiaGaloisSymmetry/RdArith.lean}{Formalism.HC.Der.Rd.KK_gcd}
\begin{proof}
(1) Reduction $(\mathbf Z/N\mathbf Z)^\times\to(\mathbf Z/f\mathbf Z)^\times$ is surjective with
kernel $K_N(f)$.

(2) $K_N(f)$ is the disjoint union of $K_N(f)\setminus U$ and $K_N(f)\cap U$, so by (1) the two
quotients $|K_N(f)\setminus U|\varphi(f)/\varphi(N)$ and $|K_N(f)\cap U|\varphi(f)/\varphi(N)$ add
up to $1$.

(3) Let $u_0\in K_N(f)\setminus U$. Multiplication by $u_0$ maps $K_N(f)\cap U$ injectively into
$K_N(f)\setminus U$. So $|K_N(f)\cap U|\le|K_N(f)\setminus U|$, and by (1) and (2)
$\rho_U(f)=|K_N(f)\setminus U|/|K_N(f)|\ge\frac12$.

(4) Let $u\in K_N(\gcd(f,g))$. As $\bar u\equiv1\pmod{\gcd(f,g)}$, the Chinese remainder theorem
gives an integer $k$ with $k\equiv1\pmod f$ and $k\equiv\bar u\pmod g$. Then $k$ is prime to $f$
and to $g$, as $\bar u$ is prime to $N$, so $k$ is prime to $\operatorname{lcm}(f,g)$, which
divides $N$. Let $v\in(\mathbf Z/N\mathbf Z)^\times$ reduce to the class of $k$ modulo
$\operatorname{lcm}(f,g)$. Then $v\equiv1\pmod f$, so $v\in U$, and $v^{-1}u\equiv1\pmod g$, so
$v^{-1}u\in U$. Hence $u\in U$.
\end{proof}

\begin{lemma}\label{asm-global-link-count-bounds-weights-root-weights}
Let $N\ge1$ be an integer and $U$ a subgroup of $(\mathbf Z/N\mathbf Z)^\times$, with $n_u$ and
$S(U)$ as in Definition~\ref{asm-global-link-count-bounds-weights-root-ratios}. Then:
\begin{enumerate}
\item for every divisor $m$ of $N$ and every $u\in\mathbf Z/N\mathbf Z$, $n_u$ divides $m$ if and
  only if $\bar u\equiv1\pmod{N/m}$;
\item for every divisor $m\ne1$ of $N$,
  \[
   \sum_{m\mid m'\mid N}\Lambda(m')\,\varphi(N/m')=\varphi(N)\,\frac{\Lambda(m)}{\varphi(m)};
  \]
\item $\displaystyle\sum_{u\in(\mathbf Z/N\mathbf Z)^\times\setminus U}\frac{\Lambda(n_u)}{\varphi(n_u)}=S(U)$.
\end{enumerate}
\end{lemma}
\leanref{Formalism/SharpRestriction/RamificationFrobenius/TotientGaloisClasses/InertiaGaloisSymmetry/RdArith.lean}{Formalism.HC.Der.Rd.order_dvd_iff}
\leanref{Formalism/SharpRestriction/RamificationFrobenius/TotientGaloisClasses/InertiaGaloisSymmetry/RdArith.lean}{Formalism.HC.Der.Rd.weight_sum}
\leanref{Formalism/SharpRestriction/RamificationFrobenius/TotientGaloisClasses/InertiaGaloisSymmetry/RdArith.lean}{Formalism.HC.Der.Rd.abel}
\begin{proof}
(1) Write $N=mc$. Then $n_u\mid m$ if and only if $m(\bar u-1)\equiv0\pmod{mc}$, that is, $c$
divides $\bar u-1$.

(2) Suppose first that $m$ is not a prime power. If a multiple $m'$ of $m$ were a power of a prime
$p$, then so would be $m$, as $m\ne1$. So every term on the left vanishes, and so does $\Lambda(m)$.
Now let $m=p^j$ with $p$ prime and $j\ge1$, and write $N=p^vM$ with $p\nmid M$, so $j\le v$. The
$m'$ with $p^j\mid m'\mid N$ and $\Lambda(m')\ne0$ are the $p^{j+i}$ with $0\le i\le v-j$; for
these, $\Lambda(m')=\log p$ and $\varphi(N/m')=\varphi(p^{v-j-i})\varphi(M)$. As
$\sum_{i=0}^{v-j}\varphi(p^{v-j-i})=p^{v-j}$, the left side is $p^{v-j}\varphi(M)\log p$. And
$\varphi(N)=\varphi(p^v)\varphi(M)=p^{v-j}\varphi(p^j)\varphi(M)$, so the right side is the same.

(3) Let $u\in(\mathbf Z/N\mathbf Z)^\times\setminus U$. Then $u\ne1$, so $n_u\ne1$, and $n_u$
divides $N$. By (2) with $m=n_u$, and by (1),
\[
 \frac{\Lambda(n_u)}{\varphi(n_u)}=\frac1{\varphi(N)}\sum_{n_u\mid m\mid N}\Lambda(m)\varphi(N/m)
 =\frac1{\varphi(N)}\sum_{m\mid N,\ \bar u\equiv1\ (N/m)}\Lambda(m)\varphi(N/m).
\]
Sum over $u\notin U$ and exchange the two sums. The coefficient of
$\Lambda(m)\varphi(N/m)/\varphi(N)$ is the number of $u\in K_N(N/m)\setminus U$, so the sum is
$\sum_{m\mid N}\Lambda(m)\rho_U(N/m)=S(U)$.
\end{proof}

\begin{lemma}\label{asm-global-link-count-bounds-weights-root-cyclotomic}
\begin{enumerate}
\item Let $L$ be a number field, $m>1$ an integer and $\eta\in L$ a primitive $m$-th root of
  unity. Then $\varphi(m)\log|N_{L/\mathbf Q}(\eta-1)|=[L:\mathbf Q]\,\Lambda(m)$.
\item Let $N\ge1$ be an integer, $L$ a number field generated over $\mathbf Q$ by a primitive
  $N$-th root of unity $\zeta$, and $u\in\mathbf Z/N\mathbf Z$ with $u\ne1$, with $n_u$ and $\bar u$
  as in Definition~\ref{asm-global-link-count-bounds-weights-root-ratios}. Then
  \[
   \log|N_{L/\mathbf Q}(\zeta-\zeta^{\bar u})|=\varphi(N)\,\frac{\Lambda(n_u)}{\varphi(n_u)}.
  \]
\end{enumerate}
\end{lemma}
\leanref{Formalism/SharpRestriction/RamificationFrobenius/TotientGaloisClasses/InertiaGaloisSymmetry/RdDiff.lean}{Formalism.HC.Der.Rd.log_norm_root}
\leanref{Formalism/SharpRestriction/RamificationFrobenius/TotientGaloisClasses/InertiaGaloisSymmetry/RdDiff.lean}{Formalism.HC.Der.Rd.log_norm_weight}
\begin{proof}
(1) The minimal polynomial of $\eta$ over $\mathbf Q$ is the cyclotomic polynomial $\Phi_m$, of
degree $\varphi(m)$ \cite[Chapter~I, \S10]{Neu99}. So $\eta-1$ generates $\mathbf Q(\eta)$ and has
minimal polynomial $\Phi_m(X+1)$, whose constant term is $\Phi_m(1)$. Hence
$N_{\mathbf Q(\eta)/\mathbf Q}(\eta-1)=(-1)^{\varphi(m)}\Phi_m(1)$, and
$N_{L/\mathbf Q}(\eta-1)$ is its $[L:\mathbf Q(\eta)]$-th power.

Next, $\log|\Phi_m(1)|=\Lambda(m)$. If $m=p^k$ with $p$ prime and $k\ge1$, then
$\Phi_m(X)=\sum_{i<p}X^{ip^{k-1}}$ and $\Phi_m(1)=p$. In general, $X^m-1=\prod_{c\mid m}\Phi_c(X)$;
dividing by $X-1$ and putting $X=1$ gives $m=\prod_{c\mid m,\,c>1}\Phi_c(1)$. The prime powers
$c>1$ dividing $m$ contribute $\prod_pp^{v_p(m)}=m$, so the integers $\Phi_c(1)$ for the divisors
$c>1$ of $m$ that are not prime powers have product $\pm1$, and each is $\pm1$.

So $\log|N_{L/\mathbf Q}(\eta-1)|=[L:\mathbf Q(\eta)]\Lambda(m)$; multiply by
$\varphi(m)=[\mathbf Q(\eta):\mathbf Q]$.

(2) Let $w$ be the representative of $1-u$ in $\{0,\dots,N-1\}$. As $\bar u+w\equiv1\pmod N$,
$\zeta-\zeta^{\bar u}=\zeta^{\bar u}(\zeta^w-1)$. The norm of $\zeta$ is a rational number whose
$N$-th power is $1$, so it is $\pm1$, and
$\log|N_{L/\mathbf Q}(\zeta-\zeta^{\bar u})|=\log|N_{L/\mathbf Q}(\zeta^w-1)|$. The order of
$\zeta^w$ is the additive order of $w$ in $\mathbf Z/N\mathbf Z$, which is that of $-w=u-1$; so
$\zeta^w$ is a primitive $n_u$-th root of unity, and $n_u>1$ as $u\ne1$. Now (1) with
$[L:\mathbf Q]=\varphi(N)$ \cite[Chapter~I, \S10]{Neu99} gives the claim.
\end{proof}

\begin{lemma}\label{asm-global-link-count-bounds-weights-root-different}
Let $N\ge1$ be an integer, $L$ a number field generated over $\mathbf Q$ by a primitive $N$-th root
of unity $\zeta$, and $E$ a number field with a field embedding $E\to L$, through which we regard
$E$ as a subfield of $L$. Then
\[
 [L:E]\log|d_E|=\sum_\sigma\log|N_{L/\mathbf Q}(\zeta-\sigma\zeta)|,
\]
the sum over the $\sigma\in\operatorname{Gal}(L/\mathbf Q)$ that are not the identity on $E$.
\end{lemma}
\leanref{Formalism/SharpRestriction/RamificationFrobenius/TotientGaloisClasses/InertiaGaloisSymmetry/RdDiff.lean}{Formalism.HC.Der.Rd.log_discr_sub}
\leanref{Formalism/SharpRestriction/RamificationFrobenius/TotientGaloisClasses/InertiaGaloisSymmetry/RdDiff.lean}{Formalism.HC.Der.Rd.log_discr_tower}
\leanref{Formalism/SharpRestriction/RamificationFrobenius/TotientGaloisClasses/InertiaGaloisSymmetry/RdDiff.lean}{Formalism.HC.Der.Rd.absNorm_diff}
\begin{proof}
Let $E'$ be $E$ or $\mathbf Q$. The ring of integers of $L$ is $\mathbf Z[\zeta]$
\cite[Chapter~I, \S10]{Neu99}, so it is $\mathcal O_{E'}[\zeta]$, and $L=E'(\zeta)$. An element
$\tau$ of $\operatorname{Gal}(L/E')$ is determined by $\tau\zeta$, so the $\tau\zeta$ are
$[L:E']$ distinct roots of the minimal polynomial $f$ of $\zeta$ over $E'$, which has degree
$[L:E']$. Hence $f=\prod_{\tau\in\operatorname{Gal}(L/E')}(X-\tau\zeta)$, with coefficients in
$\mathcal O_{E'}$, and
\[
 f'(\zeta)=\prod_{\tau\ne1}(\zeta-\tau\zeta)\ne0.
\]
As $\mathcal O_L=\mathcal O_{E'}[\zeta]$, the different $\mathfrak D_{L/E'}$ is generated by
$f'(\zeta)$. Differents are multiplicative in towers,
$\mathfrak D_{L/\mathbf Q}=\mathfrak D_{L/E'}\cdot\mathfrak D_{E'/\mathbf Q}\mathcal O_L$, and the
absolute norm of $\mathfrak D_{K/\mathbf Q}$ is $|d_K|$ for every number field $K$
\cite[Chapter~III, \S2]{Neu99}. The absolute norm of $\mathfrak D_{E'/\mathbf Q}\mathcal O_L$ is that
of $\mathfrak D_{E'/\mathbf Q}$ raised to the power $[L:E']$, since for a prime $\mathfrak p$ of
$E'$ the primes $\mathfrak P$ of $L$ over it satisfy $\sum_{\mathfrak P}e_{\mathfrak P}f_{\mathfrak
P}=[L:E']$ \cite[Chapter~I, \S8]{Neu99}. Taking absolute norms and logarithms,
\[
 \log|d_L|=\sum_{\tau\in\operatorname{Gal}(L/E'),\ \tau\ne1}\log|N_{L/\mathbf Q}(\zeta-\tau\zeta)|
 +[L:E']\log|d_{E'}|.
\]
For $E'=\mathbf Q$, $d_{\mathbf Q}=1$, and the sum runs over all $\sigma\ne1$ in
$\operatorname{Gal}(L/\mathbf Q)$. Subtract the identity for $E'=E$: the elements of
$\operatorname{Gal}(L/E)$ are the $\sigma\in\operatorname{Gal}(L/\mathbf Q)$ that are the identity
on $E$, so what remains is the claim.
\end{proof}

\begin{lemma}\label{asm-global-link-count-bounds-weights-root-discriminant}
Let $N\ge1$ be an integer, $L$ a number field generated over $\mathbf Q$ by a primitive $N$-th root
of unity $\zeta$, and $E$ a number field with a field embedding $E\to L$, through which we regard
$E$ as a subfield of $L$. Let $U$ be a subgroup of $(\mathbf Z/N\mathbf Z)^\times$ such that for
every $\sigma\in\operatorname{Gal}(L/\mathbf Q)$ and every $u\in(\mathbf Z/N\mathbf Z)^\times$ with
$\sigma\zeta=\zeta^{\bar u}$, $\sigma$ is the identity on $E$ if and only if $u\in U$. Then, with
$n_u$ as in Definition~\ref{asm-global-link-count-bounds-weights-root-ratios},
\[
 \frac{\log|d_E|}{[E:\mathbf Q]}=\sum_{u\in(\mathbf Z/N\mathbf Z)^\times\setminus U}
 \frac{\Lambda(n_u)}{\varphi(n_u)}.
\]
\end{lemma}
\leanref{Formalism/SharpRestriction/RamificationFrobenius/TotientGaloisClasses/InertiaGaloisSymmetry/RdDiff.lean}{Formalism.HC.Der.Rd.log_rd}
\begin{proof}
Every $\sigma\in\operatorname{Gal}(L/\mathbf Q)$ maps $\zeta$ to $\zeta^{\bar u}$ for a unique
$u=u_\sigma\in(\mathbf Z/N\mathbf Z)^\times$, and $\sigma$ is determined by $\sigma\zeta$. So
$\sigma\mapsto u_\sigma$ is injective, and it is bijective, as both sets have
$[L:\mathbf Q]=\varphi(N)$ elements \cite[Chapter~I, \S10]{Neu99}. By hypothesis, $\sigma$ is not the
identity on $E$ if and only if $u_\sigma\notin U$; then $u_\sigma\ne1$, as $1\in U$, and
Lemma~\ref{asm-global-link-count-bounds-weights-root-cyclotomic}(2) gives
$\log|N_{L/\mathbf Q}(\zeta-\sigma\zeta)|=\varphi(N)\Lambda(n_{u_\sigma})/\varphi(n_{u_\sigma})$.
Lemma~\ref{asm-global-link-count-bounds-weights-root-different} then gives
\[
 [L:E]\log|d_E|=\varphi(N)\sum_{u\notin U}\frac{\Lambda(n_u)}{\varphi(n_u)},
\]
and $\varphi(N)=[L:E][E:\mathbf Q]$.
\end{proof}

\begin{lemma}\label{asm-global-link-count-bounds-weights-root-two}
Let $N\ge1$ and $e\ge0$ be integers, let $M$ be an odd positive integer with $2^{e+2}M\mid N$, and
let $U$ be a subgroup of $(\mathbf Z/N\mathbf Z)^\times$ all of whose elements are $\equiv\pm1\pmod
{2^{e+2}}$. Then for every integer $0\le i\le e+2$, with $\rho_U$ as in
Definition~\ref{asm-global-link-count-bounds-weights-root-ratios},
\[
 \rho_U(2^iM)\ge1-\frac{\varphi(2^i)+\delta_i}{2^{e+1}},
\]
where $\delta_i=1$ if $i<2$ and $\delta_i=0$ otherwise.
\end{lemma}
\leanref{Formalism/SharpRestriction/RamificationFrobenius/TotientGaloisClasses/InertiaGaloisSymmetry/RdBound.lean}{Formalism.HC.Der.Rd.rho_two}
\begin{proof}
Put $f=2^iM$ and $F=2^{e+2}M$, both divisors of $N$, and $K'=K_N(F)$. Let $c_i=2$ if $i<2$ and
$c_i=1$ otherwise. We show $|K_N(f)\cap U|\le c_i|K'|$.

Let $u\in K_N(f)\cap U$. Then $u\equiv1\pmod M$ and $u\equiv\pm1\pmod{2^{e+2}}$. If
$u\equiv1\pmod{2^{e+2}}$, then $u\in K'$, as $2^{e+2}$ and $M$ are coprime. If $i\ge2$, then
$u\equiv1\pmod4$, so $u\not\equiv-1\pmod{2^{e+2}}$; hence $K_N(f)\cap U\subset K'$. If $i<2$ and
some $w\in K_N(f)\cap U$ has $w\equiv-1\pmod{2^{e+2}}$, then $w^{-1}u\in K'$ for every
$u\in K_N(f)\cap U$ with $u\equiv-1\pmod{2^{e+2}}$; hence $K_N(f)\cap U\subset K'\cup wK'$. If
there is no such $w$, then $K_N(f)\cap U\subset K'$.

By Lemma~\ref{asm-global-link-count-bounds-weights-root-ratio-facts}(1),
$\varphi(N)=|K'|\varphi(F)=|K'|\,2^{e+1}\varphi(M)$, and $\varphi(f)=\varphi(2^i)\varphi(M)$. By
part (2) of that lemma,
\[
 \rho_U(f)=1-\frac{|K_N(f)\cap U|\,\varphi(2^i)}{|K'|\,2^{e+1}}\ge1-\frac{c_i\varphi(2^i)}{2^{e+1}},
\]
and $c_i\varphi(2^i)=\varphi(2^i)+\delta_i$, as $\varphi(2^i)=1$ for $i<2$.
\end{proof}

\begin{lemma}\label{asm-global-link-count-bounds-weights-root-sum}
Let $N\ge1$ and $e\ge0$ be integers with $2^{e+2}\mid N$, let $d$ be an odd prime dividing $N$, and
let $U$ be a subgroup of $(\mathbf Z/N\mathbf Z)^\times$ all of whose elements are
$\equiv\pm1\pmod{2^{e+2}}$, with $\rho_U(N/d)\ge\frac12$ (Definition~\ref{asm-global-link-count-bounds-weights-root-ratios}).
Then
\[
 S(U)\ge(e+1-2^{-e})\log2+\tfrac12\log d.
\]
\end{lemma}
\leanref{Formalism/SharpRestriction/RamificationFrobenius/TotientGaloisClasses/InertiaGaloisSymmetry/RdBound.lean}{Formalism.HC.Der.Rd.le_SS}
\begin{proof}
Write $N=2^vM$ with $M$ odd. As $2^{e+2}\mid N$, $v\ge e+2$, so $2^{e+2}M\mid N$. For $0\le i\le e+1$
put $m_i=2^{v-i}$: it divides $N$, $N/m_i=2^iM$, and $\Lambda(m_i)=\log2$ as $v-i\ge1$. The $m_i$
are distinct and even, so they differ from $d$. Every term of $S(U)$ is nonnegative, so
\[
 S(U)\ge\log2\sum_{i=0}^{e+1}\rho_U(2^iM)+\log d\cdot\rho_U(N/d).
\]
By Lemma~\ref{asm-global-link-count-bounds-weights-root-two}, with
$\sum_{i=0}^{e+1}\varphi(2^i)=2^{e+1}$ and $\sum_{i=0}^{e+1}\delta_i=2$,
\[
 \sum_{i=0}^{e+1}\rho_U(2^iM)\ge(e+2)-\frac{2^{e+1}+2}{2^{e+1}}=e+1-2^{-e}.
\]
And $\log d\cdot\rho_U(N/d)\ge\frac12\log d$, as $\log d\ge0$.
\end{proof}

\begin{lemma}\label{asm-global-link-count-bounds-weights-root}
Let $\Phi$ be a Galois datum with $2$-power exponent $e\ge2$
(Definition~\ref{hf-frame}), and let $(q,d,Q_S,E,\gamma,j,\operatorname{art})$
be a host field on $\Phi$ (Definition~\ref{hf-framed-host-field}), with root
discriminant $\mathrm{rd}_E=|d_E|^{1/[E:\mathbf Q]}$ and half root
discriminant $v=\mathrm{rd}_E/2$. Then:
\begin{enumerate}
\item $\mathrm{rd}_E\ge2^{e+1-2^{-e}}\sqrt d$;
\item $v\ge4>\pi$.
\end{enumerate}
\end{lemma}
\leanref{Formalism/SharpRestriction/RamificationFrobenius/TotientGaloisClasses/InertiaGaloisSymmetry/DvdDiscrOfInertia.lean}{Formalism.HC.Der.rd_ge}
\leanref{Formalism/SharpRestriction/GlobalSchurBounds/JointsL.lean}{Formalism.MainNotions.GlobalTheoryTheory.JDB.db_pi_le_v}
\begin{proof}
(1) Put $N=Q_Sqd$ and let $U\subset(\mathbf Z/N\mathbf Z)^\times$ be the kernel of
$\operatorname{art}$; we use the notation $K_N$, $\rho_U$, $S(U)$, $n_u$ and
$\bar u$ of Definition~\ref{asm-global-link-count-bounds-weights-root-ratios}.
Let $\zeta$ be a primitive $N$-th root of unity in
$\mathbf Q(\zeta_N)$, and regard $E$ as a subfield of $\mathbf Q(\zeta_N)$ through
$j$. If $\sigma\in\operatorname{Gal}(\mathbf Q(\zeta_N)/\mathbf Q)$ and
$u\in(\mathbf Z/N\mathbf Z)^\times$ satisfy $\sigma\zeta=\zeta^{\bar u}$, then
$\sigma=\sigma_u$, and $\sigma_u\circ j=j\circ\gamma^{-1}(\operatorname{art}(u))$
by Definition~\ref{hf-framed-host-field}(3). As $j$ and $\gamma$ are injective,
$\sigma$ is the identity on $E$ if and only if $\operatorname{art}(u)=1$, that
is, $u\in U$. So
Lemmas~\ref{asm-global-link-count-bounds-weights-root-discriminant}
and~\ref{asm-global-link-count-bounds-weights-root-weights}(3) give
\[
 \log\mathrm{rd}_E=\frac{\log|d_E|}{[E:\mathbf Q]}
 =\sum_{u\notin U}\frac{\Lambda(n_u)}{\varphi(n_u)}=S(U).
\]

We check the hypotheses of
Lemma~\ref{asm-global-link-count-bounds-weights-root-sum} for $U$, $e$ and the
prime $d$. First, $2^{e+2}\mid Q_S\mid N$, and $d\equiv7\pmod{16}$ is an odd
prime dividing $N$. Next, let $u\in U$. Then $\operatorname{art}(u)=1$, whose
$A$-component is $0$, and by Definition~\ref{hf-framed-host-field}(4) $h$ maps
it to the class of $u$ in $H_e=(\mathbf Z/2^{e+2}\mathbf Z)^\times/\{\pm1\}$. So
that class is trivial, and $u\equiv\pm1\pmod{2^{e+2}}$.

Finally, $\rho_U(N/d)\ge\frac12$, where $N/d=Q_Sq$. By
Lemma~\ref{asm-global-link-count-bounds-weights-root-ratio-facts}(3) it suffices
to find $u\equiv1\pmod{Q_Sq}$ with $\operatorname{art}(u)\ne1$. Suppose there is
none, and let $U'$ be the kernel of $u\mapsto\epsilon_2(\operatorname{art}(u))$.
Then $K_N(Q_Sq)\subset U\subset U'$. Also $K_N(d)\subset U'$: by
Proposition~\ref{hf-host-field-properties}(2),
$\epsilon_2(\operatorname{art}(u))=\left(\frac ud\right)=1$ for
$u\equiv1\pmod d$. As $\gcd(Q_S,qd)=1$ and $q\ne d$ (for $q\equiv3$ and
$d\equiv7\pmod8$), $\gcd(Q_Sq,d)=1$, and
Lemma~\ref{asm-global-link-count-bounds-weights-root-ratio-facts}(4) gives
$(\mathbf Z/N\mathbf Z)^\times=K_N(1)\subset U'$. But
$\operatorname{art}(-1)=\iota$ has second sign $-1$, so $-1\notin U'$. This
contradiction gives the element $u$.

Lemma~\ref{asm-global-link-count-bounds-weights-root-sum} now gives
\[
 \log\mathrm{rd}_E=S(U)\ge(e+1-2^{-e})\log2+\tfrac12\log d
 =\log\big(2^{e+1-2^{-e}}\sqrt d\big).
\]

(2) As $e\ge2$,
$e+1-2^{-e}\ge3-\frac14\ge2$, so $2^{e+1-2^{-e}}\ge4$; and $d\equiv7\bmod16$
gives $d\ge7$, so $\sqrt d\ge2$. By (1), $\mathrm{rd}_E\ge8$, and
$v\ge4>\pi$.
\end{proof}

\begin{lemma}\label{asm-global-link-count-bounds-weights-hold}
Let $\mathcal F$ be a family with exponent $\varepsilon$
(Definition~\ref{pd-family}), and let $C_{YL}$ be its constant of that name.
Then every member $D$ at every scale $x\ge3$ with $x>x_1$ has weight data
$(\max(5,C_{YL}+1),\Delta_S,w)$ at $(x,\varepsilon)$
(Definition~\ref{asm-global-link-count-bounds-weights}), where $\Delta_S$ is
the number of Definition~\ref{th-definite-datum}(8) and $w_j(\tau)$ is the
Harish-Chandra weight of the lift $j$ at $\tau$.
\end{lemma}
\leanref{Formalism/SharpRestriction/GlobalSchurBounds/JointsL.lean}{Formalism.MainNotions.GlobalTheoryTheory.JDB.dataB_fam}
\leanref{Formalism/SharpRestriction/GlobalSchurBounds/JointsL.lean}{Formalism.MainNotions.GlobalTheoryTheory.JDB.dataB_of}
\begin{proof}
Let $D$ be a member at a scale $x\ge3$ with $x>x_1$, and put
$C_0=\max(5,C_{YL}+1)$ and $\vartheta=\varepsilon/100$. Use the notation of
Definition~\ref{pd-size-data} for its size data, with the width $\eta>0$, the
cloud size $S_c=4k\Delta_S$, $\delta_2=|T_2|/n$ and $h_0=\eta/m$; write
$u_i(\tau)$, $v_i(\tau)$, $b_j(\tau)$ and $e_j(\tau)=w_j(\tau)$ for the entries
of $\lambda(\tau)$. The family's choices at $x$ include parameter choices and
the size bounds (Definition~\ref{asm-dependence-functional-witness-sizes}(1)
and~(2)); we number the size bounds as in
Lemma~\ref{asm-dependence-functional-witness-numbers-bounds-computations}. The
proof also uses the following further conditions of the size data, for every
$\tau\in\mathcal T$ and $1\le i,j\le m$:
\begin{enumerate}
\item[(i)] $0\le v_l(\tau)\le S_c$ for $1\le l\le k-1$, and
  $|b_i(\tau)-e_j(\tau)|\le\eta$;
\item[(ii)] if $\tau$ is not top, then $S_c\le u_1(\tau)-v_1(\tau)$;
\item[(iii)] if the element of $G_\Phi$ of $\tau$ is $\tau_1(t)$ for a top
  class $t$, then $u_1(\tau)=a_t+(k-1)/2$ and $u_k(\tau)=-b_t-(k-1)/2$;
\item[(iv)] the windows of the weights: if the element of $G_\Phi$ of $\tau$
  is $\tau_1(t)$ for a top class $t$, then $b_t+b'/x\le e_j(\tau)\le
  b_t+b'/x+\eta$; if it is a thick place, then $b\le e_j(\tau)\le b+\eta+1$;
  otherwise $w''\le e_j(\tau)\le w''+\eta$;
\item[(v)] $0<\delta_2\le1$;
\end{enumerate}
and the following parameter choices:
\begin{enumerate}
\item[(vi)] if $\tau$ is not top, then $u_1(\tau)-v_1(\tau)\le2S_c$ and
  $u_k(\tau)=v_{k-1}(\tau)-\frac12$;
\item[(vii)] there is $Y>0$ with $Y\le g_i\le2Y$ for $i=0,1$ and $\log Y\le
  C_{YL}L$;
\item[(viii)] $h_0\le b_{j+1}(\tau)-b_j(\tau)$ for $1\le j<m$.
\end{enumerate}
By Definition~\ref{th-definite-datum}(2) the parameter $\lambda_u(\tau)$ of a
word $u$ has the same $v$, $b$, $e$, $u_1$ and $u_k$ as $\lambda(\tau)$, so
(ii), (iii) and (vi) also hold for $\lambda_u(\tau)$.

\emph{Basic sizes.} By Definition~\ref{th-ranks}, $m\ge2$, $k\ge m+3\ge5$,
$t\ge k+m+4\ge k$, $f=2n\ge2$ and $r=k+m-1\ge k$. As $x>1$, $L=x^\varepsilon\ge1$,
so $b=e^L\ge1+L\ge2$ and $\sqrt b\le b$; as $x\ge3>e$, $\log x\ge1$. Put
$Q=Q_Sqd\ge1$. Size bound (3) gives $Q^{36}t^{150}\le\sqrt b$, so $k\le t\le
Q^{36}t^{150}\le b$; and $\mathrm{rd}_E\le\frac{Q_S}2\sqrt{qd}$, so
$v=\mathrm{rd}_E/2\le Q\le b$. Size bound (5) gives $\eta\le b/(t^3fk^3)\le
b/27$, and $t\eta\le w''/x$, so $(x+1)\eta\le tx\eta\le w''$, as $x+1\le3x\le
tx$. Also $w''=b'/(xn)$ with $b'>0$ and $xn\ge1$, so $0<w''\le b'$.

\emph{The exponents $b_t$.} Size bound (4) gives $x^\varepsilon\le\Lambda_G$, so
$\Lambda_G>0$, and
$4N^*x^\varepsilon\Lambda_G/k+x^\vartheta\Lambda_G\log x/k\le x^{-2}\le1$; the
second term is nonnegative, so $N^*x^\varepsilon\Lambda_G/k\le\frac14$, and the
bound $b_t\le b'+N^*x^\varepsilon\Lambda_Gb'/k$ of (4) gives $b_t\le\frac54b'$.
Size bound (5) gives $r(b'+S_c)\ell_A\log x\le\delta_2mb$ with
$\ell_A=kx^\varepsilon/(m\delta_2)$, that is $r(b'+S_c)k\,x^\varepsilon\log x\le
\delta_2^2m^2b$. As $r,x^\varepsilon,\log x\ge1$, $S_c>0$, $\delta_2\le1$ and
$m^2\le m^2L\le k$ (size bound (2)), this gives $b'k\le m^2b\le kb$. So $b'\le
b$ and $b_t\le\frac54b$ for every top class $t$.

\emph{The windows.} For every $\tau$ and $j$, $w''\le e_j(\tau)\le2b$. Indeed,
by (iv): at the top embedding of a top class $t$, $e_j(\tau)\ge b_t\ge b'\ge
w''$ by Definition~\ref{pd-size-data}(2), and $e_j(\tau)\le\frac54b+\frac
b3+\frac b{27}\le2b$ as $b'/x\le b/3$; at a thick place, $e_j(\tau)\ge b\ge
w''$ and $e_j(\tau)\le b+\frac b{27}+\frac b2\le2b$, as $1\le b/2$; and
otherwise $w''\le e_j(\tau)\le w''+\eta\le b+\frac b{27}\le2b$.

\emph{Clause (1).} By Lemma~\ref{asm-global-link-count-bounds-weights-root},
$\pi\le v$, and $\log v\le\log b=L$. Size bound (1) gives $n\le
x^{2+8\varepsilon+\vartheta}$, so $f=2n\le x\cdot x^{2+8\varepsilon+\vartheta}\le
x^5\le x^{C_0}$, as $2\le x$, $3+8\varepsilon+\vartheta\le5$ and $C_0\ge5$; in the
same way $k\le x^{9/2+10\varepsilon+\vartheta}\le x^{C_0}$. Size bound (2) gives
$m^2L\le k$, and size bound (10) gives $(kb)^8\le g_i$. By (vii), $\log
g_i\le\log2+\log Y\le\log2+C_{YL}L\le(C_{YL}+1)L\le C_0L$, as $\log2<1\le L$.

\emph{Clause (2).} Size bound (8) gives $N(m)m^{10}\le\Delta_S$, and $N(m),m\ge1$,
so $\Delta_S\ge1$. By (i), $0\le v_l(\tau)\le S_c=4k\Delta_S$, and
$v_l(\tau)-v_{l+1}(\tau)=4\Delta_S$ by Definition~\ref{th-definite-datum}(8).

\emph{Clause (3).} By Definition~\ref{th-definite-datum}(8), the entry
$u_{l+1}(\tau)$ of $\lambda_u(\tau)$ is $v_l(\tau)-\frac12-\varepsilon_l(u)\Delta_S$
with $\varepsilon_l(u)\in\{0,1\}$, for $1\le l\le k-2$. If $\tau$ is not top,
then (ii) and (vi) give $4k\Delta_S=S_c\le u_1(\tau)-v_1(\tau)\le2S_c=8k\Delta_S$
and $u_k(\tau)=v_{k-1}(\tau)-\frac12$. If $\tau$ is top, the top embedding of
the top class $t$, then by (iii), Definition~\ref{pd-size-data}(2) and
$kx^{13}\ge3$,
\[
 u_1(\tau)=a_t+\frac{k-1}2\le a+\frac a3+\frac b2\le2a,\qquad
 -u_k(\tau)=b_t+\frac{k-1}2\le\frac54b+\frac b2\le a,
\]
since $k\le b$ and $3b\le xb=a$.

\emph{Clause (4).} By the windows, $0<w''\le w_j(\tau)$ and $w_m(\tau)\le2b$. By
Definition~\ref{th-hc-parameter}(4), $e_j(\tau)-b_j(\tau)\ge(t+3)/2>0$ and
$b_{j+1}(\tau)-e_j(\tau)\ge(t+3)/2>0$. By (i), $0<e_j(\tau)-b_j(\tau)\le\eta$, so
\[
 x\big(e_j(\tau)-b_j(\tau)\big)\le x\eta\le w''-\eta\le e_j(\tau)-\eta\le
 b_j(\tau),
\]
which is $|w_j(\tau)-b_j(\tau)|\le b_j(\tau)/x$. By (viii) and size bound (6),
as $C_0\ge1$ and $x^\vartheta\ge1$,
\[
 \frac{r^{19}f^5L}{C_0}\le x^\vartheta r^{19}f^5L\le h_0\le
 b_{j+1}(\tau)-b_j(\tau).
\]
Finally, size bound (1) gives $k\ge x^{9/2+10\varepsilon-\vartheta}\ge x^2\ge3x$,
and size bound (8) gives $kS_c\le w''$, so $xS_c\le w''/3$. As $\eta\le
w''/(x+1)\le w''/4$, $b_1(\tau)\ge e_1(\tau)-\eta\ge\frac34w''\ge xS_c$, which is
$4k\Delta_S\le b_1(\tau)/x$.
\end{proof}
% Lean: dataB_fam = dataB_of (Fm.L x i) (Fm.params x hx i) (Fm.sizes x hx i) with the threshold 3,
% and Fm.params_CYL for C_YL. dataB_of collects: sizes_clause_seventeen (clause (1) and Delta >= 1:
% db_pi_le_v, SizeN, SizeK, SizeKRel, SizeG, P.g_Y, SizeDel); db_cloud (L.cloud_v with L.Sc_eq,
% D.s.cloud_gap, D.s.u_letter, L.cloud_gap, P.cloud_place_bound, P.uk_extreme, D.hcW_v/first/last);
% db_top (L.top_u, L.aT_win with A/2 = 13, db_bT from SizeBot, SizeRho, SizeLa, and db_b'le from the
% last clause of SizeEta); sizes_clause_thirteen to sixteen (db_win from L.win_top, L.win_thick,
% L.win_thin and L.bT_ge; db_be from gap_be and L.be_close; db_xeta and eta_size_le_b_div_27 from
% SizeEta; db_xS from SizeDel and SizeK; gap_eb; P.nodes_gap with h0 = eta/m and SizeNodes). The basic
% sizes are db_rk and db_basic (SizeCond, SizeQD). About 23 theorems and 480 lines in JointsL.lean.

\begin{lemma}\label{asm-global-link-count-bounds-explicit-rows}
Let $\mathcal F$ be a family with exponent $\varepsilon$
(Definition~\ref{pd-family}) and $C_0$ real. Let $D$ be a member at a scale
$x\ge44^4$ with $x>x_1$, with the notation of
Definition~\ref{asm-global-link-count-bounds-weights}, and suppose that $D$ has
weight data $(C_0,\Delta,w)$ at $(x,\varepsilon)$. Then $D$ satisfies the
\emph{row conditions} with $(C_0,\Delta)$: $\log(2v)\le L$; $\log g_i\le C_0L$
for $i=0,1$; $44\le m\le k$ and $m\le x$; $\log(2k)\le L$; and for every word $u$
and every $\tau\in\mathcal T$, with $u_i(\tau)$ the entries of $\lambda_u(\tau)$,
$u_1(\tau)-u_k(\tau)\le2a$, and, if $\tau$ is top, $4k\Delta\le
u_1(\tau)-v_1(\tau)$ and $u_k(\tau)\le-4k\Delta$.
\end{lemma}
\leanref{Formalism/SharpRestriction/GlobalSchurBounds/GlobalRankLatticeProofs/FamilyPartFiveOfExplicitConstants.lean}{Formalism.MainNotions.GlobalTheoryTheory.JDB.rowsA_of_dataB_fam}
\leanref{Formalism/SharpRestriction/GlobalSchurBounds/GlobalRankLatticeProofs/FamilyPartFiveOfExplicitConstants.lean}{Formalism.MainNotions.GlobalTheoryTheory.rowsAF_fam}
\leanref{Formalism/SharpRestriction/GlobalSchurBounds/TheoremThreeOfExplicitAndCompact.lean}{Formalism.MainNotions.StandardLattice.RowsA}
\begin{proof}
Use the notation of the proof of
Lemma~\ref{asm-global-link-count-bounds-weights-hold}: the size data with
$\vartheta=\varepsilon/100$, the cloud size $S_c=4k\Delta_S$, the size bounds of
the family numbered as in
Lemma~\ref{asm-dependence-functional-witness-numbers-bounds-computations}, the
conditions (i) to (iii) of the size data and the parameter choice (vi) listed
there. As $x\ge44^4\ge3$, the basic sizes established there hold: $L\ge1$,
$b\ge2$, $k\le t$, $k\le b$, and $b_t\le\frac54b'$ for every top class $t$.

\emph{The number $\Delta$.} By Definition~\ref{asm-global-link-count-bounds-weights}(2)
and Definition~\ref{th-definite-datum}(8), $4\Delta=v_1(\tau)-v_2(\tau)=4\Delta_S$, so
$\Delta=\Delta_S$ and $4k\Delta=S_c$.

\emph{Sizes.} Put $Q=Q_Sqd\ge1$. By size bound (3), $Q^{36}t^{150}\le\sqrt b\le b$
and $2v=\mathrm{rd}_E\le\frac{Q_S}2\sqrt{qd}\le Q$; also $t\ge k+6\ge2$
(Definition~\ref{th-ranks}), so $2k\le2t\le t^2\le t^{150}$. So $2v\le b$ and $2k\le b$, that is
$\log(2v)\le L$ and $\log(2k)\le L$. By size bound (1), with $\delta_1=
x^{-\varepsilon/8}\le1$, $m\ge x^{-\vartheta}\sqrt{x\delta_1}\ge x^{1/4}\ge44$ and
$m\le x^\vartheta\sqrt x\le x$; and $m\le k$ by Definition~\ref{th-ranks}. The
bound $\log g_i\le C_0L$ is Definition~\ref{asm-global-link-count-bounds-weights}(1).
Next, $S_c\le kS_c\le w''\le b'$ by size bound (8). Size bounds (1) and (8) give
$m^{10}\le N(m)m^{10}\le\Delta_S$, $\log\Delta_S\le x^{\varepsilon/4}/2$ and $\log
m\ge\frac14\log x$, so $\frac{10}4\log x\le x^{\varepsilon/4}/2$ and $\log x\le
x^{\varepsilon/3}$; hence $x\le X'$ and $b'=xb/X'\le b$. Finally $100b\le xb=a$.

\emph{The entries.} Let $u$ be a word and $\tau\in\mathcal T$. If $\tau$ is the
top embedding of a top class $t$, then by (iii) $u_1(\tau)=a_t+(k-1)/2$ and
$u_k(\tau)=-b_t-(k-1)/2$, with $a\le a_t\le\frac43a$ as $kx^{13}\ge3$. So
$u_1(\tau)-u_k(\tau)=a_t+b_t+k-1\le\frac43a+\frac54b+b\le2a$; by (i),
$v_1(\tau)\le S_c\le b$, so $u_1(\tau)-v_1(\tau)\ge a-b\ge S_c=4k\Delta$; and
$u_k(\tau)\le-b_t\le-b'\le-S_c=-4k\Delta$. If $\tau$ is not top, then by (i),
(ii) and (vi), $0\le v_i(\tau)\le S_c$, $u_1(\tau)\le v_1(\tau)+2S_c\le3S_c$ and
$u_k(\tau)=v_{k-1}(\tau)-\frac12\ge-\frac12$, so $u_1(\tau)-u_k(\tau)\le3S_c+
\frac12\le4b\le2a$. The parameter $\lambda_u(\tau)$ has the same $v$, $u_1$ and
$u_k$ as $\lambda(\tau)$ (Definition~\ref{th-definite-datum}(2)).
\end{proof}
% Lean: rowsA_of_dataB_fam = rowsA_of_dataB (Fm.L x i) (Fm.params x hx i) (Fm.sizes x hx i) with
% x >= 44^4: Delta = Del from the clause v_1 - v_2 = 4 Delta of DataB and db_cloud, then rowsA_of' with
% the bound on log g_i of DataB: ra_cond (2v <= b, 2k <= b), ra_44, rank_m_le_x, ra_Sc (Sc <= b',
% b' <= b, 100 b <= a), db_basic, db_bT, db_cloud, L.top_u, L.aT_win. About 13 theorems and 410 lines.

\begin{lemma}\label{asm-global-link-count-bounds-explicit-log-x}
Let $\mathcal F$ be a family with exponent $\varepsilon$
(Definition~\ref{pd-family}) such that for all large $x$ every member at $x$
has the properties (a), (b) and (c) of
Lemma~\ref{asm-global-link-count-bounds}, and let $C_0,C_Y$ be real. Then there
is $x_0$ with the following property. Let $D$ be a member at a scale $x\ge x_0$
with $x>x_1$, with the notation of
Definition~\ref{asm-global-link-count-bounds-weights} and with $X_u$ as in
Definition~\ref{th-definite-datum}(3), and suppose that $D$ has weight data
$(C_0,\Delta,w)$ at $(x,\varepsilon)$ and satisfies the row conditions with
$(C_Y,\Delta)$ (Lemma~\ref{asm-global-link-count-bounds-explicit-rows}). Then
for every word $u$
\[
 S(w)\le\log X_u\le(15+3C_Y)\,mL,\qquad
 S(w)=\frac1f\sum_{j=1}^m\sum_{\tau\in\mathcal T}\log\frac{w_j(\tau)v}\pi
 +m\log(g_0g_1)-11L.
\]
\end{lemma}
\leanref{Formalism/SharpRestriction/GlobalSchurBounds/TheoremThreeOfExplicitAndCompact.lean}{Formalism.MainNotions.StandardLattice.Conclusion5Sharp}
\leanref{Formalism/SharpRestriction/GlobalSchurBounds/LinkMin.lean}{Formalism.MainNotions.StandardLattice.gapB_ev}
\leanref{Formalism/SharpRestriction/GlobalSchurBounds/LinkMin.lean}{Formalism.MainNotions.StandardLattice.ramified_rank_lines_of_free}
\leanref{Formalism/SharpRestriction/GlobalSchurBounds/LinkMin.lean}{Formalism.MainNotions.StandardLattice.part_five_family_of_mass_volume}
\begin{proof}
Let $x_0$ be beyond the scales from which (a), (b) and (c) hold and the
thresholds of the three steps below, which depend only on $\varepsilon$, $C_0$
and $C_Y$. Let $D$, $\Delta$ and $w$ be as in the statement.

\emph{The spectral gap.} The weight data give, at every place $\tau$, a spectral
gap of the parameter of $\tau$ large enough for $\epsilon_0=1/3$, with the
constant $8192$. By $v_i(\tau)\ge0$, $b_1(\tau)>0$ and the gaps
$b_{j+1}(\tau)-b_j(\tau)\ge r^{19}f^5L/C_0$ of
Definition~\ref{asm-global-link-count-bounds-weights}, entries of opposite
signs of the parameter differ by at least $r^{19}f^5L/C_0$; the required gap is
at most $9\cdot8192\,r^9$, which is smaller once $L\ge9\cdot8192\,C_0$.

\emph{The lower bound.} With the spectral gap, the lower degree bound of (c),
the mass ratio (a) and the base volume (b) give $S(w)\le\log X_u$.

\emph{The upper bound.} With the spectral gap, the upper degree bound of (c),
the mass formula behind (a) and the base volume (b) give $\log X_u\le
(15+3C_Y)mL$.
\openstep{The three steps are not yet written out: the spectral gap at every
place from the weight data, and the lower and upper bounds for $\log X_u$ from
the masses of the standard lattices, the base volume and the formal degrees.}
% Lean: Conclusion5Sharp D sc C0 CY, assembled in division_level_prime_family from gapB_ev eps 8192 C0
% (GapAtLeast (rowSgn D) (rowLam D rho) (gap3 D 8192 (1/3)), through parameter_gap_under_section_B,
% gap_bound_mono and gap_bound_at_one_third; about 9 theorems and 160 lines), ramified_rank_lines_of_free
% (through upper_weight_bound_eventually; about 235 theorems and 5,270 lines) and
% part_five_family_of_mass_volume (through family_of_global_statements, massIdentity_holds and
% log_zetaR_le_one; about 260 theorems and 5,640 lines). The last two closures overlap in most of their
% lines; together the three are about 287 theorems and 6,450 lines in 46 files. Here log X_u is logXuD D u, identified with the datum's X_u by jnt_logXu with factDuStar_fam
% and factDstar_fam.
\end{proof}

\begin{lemma}\label{asm-global-link-count-bounds-explicit-variation}
Let $\mathcal F$ be a family with exponent $\varepsilon$
(Definition~\ref{pd-family}) and $C_0,C_Y$ real. Let $D$ be a member at a scale
$x>x_1$, with $X_u$ as in Definition~\ref{th-definite-datum}(3), that has weight
data $(C_0,\Delta,w)$ at $(x,\varepsilon)$
(Definition~\ref{asm-global-link-count-bounds-weights}) and satisfies the row
conditions with $(C_Y,\Delta)$
(Lemma~\ref{asm-global-link-count-bounds-explicit-rows}). Then
$|\log X_u-\log X_{u'}|\le1+\log k+2/k$ for all words $u,u'$.
\end{lemma}
\leanref{Formalism/SharpRestriction/GlobalSchurBounds/GlobalRankLatticeProofs/LogVandermondeVariationGeneric.lean}{Formalism.MainNotions.StandardLattice.vandermonde_variation_holds}
\leanref{Formalism/SharpRestriction/GlobalSchurBounds/TheoremThreeOfExplicitAndCompact.lean}{Formalism.MainNotions.StandardLattice.VandermondeVariationExplicitBound}
\begin{proof}
\openstep{The proof is not yet written out: it writes $d^*_u$ as a constant
times the product over the places of the Vandermonde determinants of the
entries $u_i(\tau)$, and bounds the variation of each Vandermonde determinant
between words, as the words change only the interior entries by at most
$\Delta$.}
% Lean: vandermonde_variation_holds (hdu: duStarU D u = massI D k * idxPi D g0 g1 / prod facProd * prod_rho
% vand (uColR D u rho); vandermonde_variation_at_place). Size: about 5 theorems and 380 lines in
% LogVandermondeVariationGeneric.lean.
\end{proof}

\begin{lemma}\label{asm-global-link-count-bounds-explicit}
Let $\mathcal F$ be a family with exponent $\varepsilon$
(Definition~\ref{pd-family}). Suppose that for all large $x$ every member $D$ at
$x$ has the following properties: (a) for every $M\ge3$ the ratio
$\operatorname{mass}(I_M)/\operatorname{mass}(I_{M-1})$ of the masses of
consecutive standard lattices is the one given by Siegel's mass formula
(Theorem~\ref{cv-siegel-mass-ratio}); (b)
$\operatorname{vol}_D(Y_0)=\operatorname{mass}(I_r)$; (c) for every place
$\tau$ and every $\epsilon_0$ with $e^{-r}\le\epsilon_0<1/2$, if the spectral
gap of the parameter of $\tau$ is large enough for $\epsilon_0$, the formal
degree of the discrete series of $\tau$ satisfies
$(1-\epsilon_0)\dim V(\lambda_\tau)\le d_\sigma(\tau)\le
4d(\lambda_\tau)(3pq|\lambda_\tau|)^{pq}$, with $V(\lambda_\tau)$ and $(p,q)$
as in Lemma~\ref{asm-global-link-count-bounds}(c). Then for every
$C_0$ there are $C$ and $x_0$ such that every member $D$ at every scale $x\ge
x_0$ with $x>x_1$, with the notation of Definition~\ref{th-definite-datum}(3)
and~(4), has the following property: for all $\Delta$ and $w$ such that $D$
has weight data $(C_0,\Delta,w)$ at $(x,\varepsilon)$
(Definition~\ref{asm-global-link-count-bounds-weights}),
\begin{enumerate}
\item for every $\Gamma\ge1$ and every word $u$ with
  $X_u\Gamma^{-2/(fk)}\le T_u\le X_u\Gamma^{-2/(fk)}(1+1/(fk))$,
  \[
   \log T_u\le CmL,\qquad
   \frac1f\sum_{j=1}^m\sum_{\tau\in\mathcal T}\log\frac{w_j(\tau)v}\pi
   +m\log(g_0g_1)-C\Big(L+\frac{m^2L}k\Big)-\frac{2\log\Gamma}{fk}
   \le\log T_u;
  \]
\item $|\log X_u-\log X_{u'}|\le C\log k$ for all words $u,u'$.
\end{enumerate}
\end{lemma}
\leanref{Formalism/SharpRestriction/GlobalSchurBounds/LinkMin.lean}{Formalism.MainNotions.StandardLattice.division_level_prime_family}
\leanref{Formalism/SharpRestriction/GlobalSchurBounds/GlobalRankLatticeProofs/FamilyPartFiveOfExplicitConstants.lean}{Formalism.MainNotions.GlobalTheoryTheory.family_part_five_of_explicit_constants}
\leanref{Formalism/SharpRestriction/GlobalSchurBounds/StandardLattice.lean}{Formalism.MainNotions.StandardLattice.TorusPairsWithVariant}
\leanref{Formalism/SharpRestriction/GlobalSchurBounds/StandardLattice.lean}{Formalism.MainNotions.StandardLattice.VandermondeVariationBound}
\begin{proof}
Let $C_0$ be given, and put $C=\max(11,16+3C_0)$. Let $x_0$ be at least $44^4$
and the $x_0$ of Lemma~\ref{asm-global-link-count-bounds-explicit-log-x} for
$(C_0,C_Y)=(C_0,C_0)$. Let $D$ be a member at a scale $x\ge x_0$ with $x>x_1$,
and let $D$ have weight data $(C_0,\Delta,w)$. By
Lemma~\ref{asm-global-link-count-bounds-explicit-rows}, $D$ satisfies the row
conditions with $(C_0,\Delta)$; in particular $m\ge44$, and $L\ge1$ as $x>1$, so
$mL\ge1$. By Lemma~\ref{asm-global-link-count-bounds-explicit-log-x}, $S(w)\le\log
X_u\le(15+3C_0)mL$ for every word $u$.

(1) Let $\Gamma\ge1$ and $u$ be as in (1). Taking logarithms in the window,
\[
 \log X_u-\frac{2\log\Gamma}{fk}\le\log T_u\le\log X_u-\frac{2\log\Gamma}{fk}
 +\log\Big(1+\frac1{fk}\Big)\le\log X_u+1,
\]
as $\log\Gamma\ge0$ and $\log(1+1/(fk))\le1$. So $\log T_u\le(15+3C_0)mL+1\le
(16+3C_0)mL\le CmL$. As $C\ge11$, $C(L+m^2L/k)\ge11L$, so the left side of the
second inequality of (1) is at most $S(w)-2\log\Gamma/(fk)\le\log
X_u-2\log\Gamma/(fk)\le\log T_u$.

(2) By Lemma~\ref{asm-global-link-count-bounds-explicit-variation},
$|\log X_u-\log X_{u'}|\le1+\log k+2/k\le3\log k\le C\log k$, as $k\ge5$ gives
$\log k\ge1$ and $2/k\le1$.
% Lean: division_level_prime_family Fm h2b h2v h3 : FamilyVandermondeCountExplicit Fm (for all C0,
% CY), from ramified_rank_lines_of_free (the lower bound lowS <= logXuD, through
% upper_weight_bound_eventually), part_five_family_of_mass_volume (logXuD <= (15 + 3 CY) m L,
% through family_of_global_statements and massIdentity_holds), gapB_ev (the spectral gap
% GapAtLeast with gap3 D 8192 (1/3)), which give Conclusion5Sharp at every member; then
% part_five_explicit_of_propositions (corollary_five_six_of_explicit, corollary_five_four_of_explicit)
% with vandermonde_variation_holds gives VandermondeCountBound. family_part_five_of_explicit_constants
% (part_five_family_of_explicit with rowsAF_fam, which supplies RowsA) turns it into
% PartFiveFamilyBounds: Conclusion5₂ = Conclusion5, TorusPairsWithVariant and
% VandermondeVariationBound. The lemma states the clauses that countF_of uses: the second and third
% clauses of TorusPairsWithVariant (lowB D sc C w is the left side of the second display) and the
% second clause of VandermondeVariationBound, with log X_u for logXuD (jnt_logXu with factDuStar_fam
% and factDstar_fam). Sizes: division_level_prime_family about 310 theorems and 7,200 lines in 47
% files; family_part_five_of_explicit_constants about 20 theorems and 520 lines in 3 files.
\end{proof}

\begin{lemma}\label{asm-global-link-count-bounds}
Let $\mathcal F$ be a family with exponent $\varepsilon$
(Definition~\ref{pd-family}). Suppose that for all large $x$ every member $D$ at
$x$ has the following properties.
\begin{enumerate}
\item[(a)] For every $M\ge3$ the ratio
  $\operatorname{mass}(I_M)/\operatorname{mass}(I_{M-1})$ of the masses of
  consecutive standard lattices is the one given by Siegel's mass formula
  (Theorem~\ref{cv-siegel-mass-ratio}).
\item[(b)] The volume of the base space equals the mass of the standard lattice
  of rank $r$: $\operatorname{vol}_D(Y_0)=\operatorname{mass}(I_r)$.
\item[(c)] For every place $\tau$ and every $\epsilon_0$ with
  $e^{-r}\le\epsilon_0<1/2$, if the spectral gap of the parameter of $\tau$ is
  large enough for $\epsilon_0$, the formal degree $d_\sigma(\tau)$ of the
  discrete series of $\tau$ satisfies $(1-\epsilon_0)\dim V(\lambda_\tau)\le
  d_\sigma(\tau)\le4d(\lambda_\tau)(3pq|\lambda_\tau|)^{pq}$, where
  $V(\lambda_\tau)$ is the representation of highest weight $\lambda_\tau$ and
  $(p,q)$ the signature.
\end{enumerate}
Then there are $C$ and $x_0$ such that the bounds (1) to (3) of
Lemma~\ref{asm-global-count} hold for every $x\ge x_0$ with $x>x_1$, every
member at $x$ and all its words.
\end{lemma}
\leanref{Formalism/SharpRestriction/GlobalSchurBounds/LinkMin.lean}{Formalism.MainNotions.StandardLattice.division_level_prime_family}
\leanref{Formalism/SharpRestriction/GlobalSchurBounds/GlobalRankLatticeProofs/FamilyPartFiveOfExplicitConstants.lean}{Formalism.MainNotions.GlobalTheoryTheory.family_part_five_of_explicit_constants}
\leanref{Formalism/SharpRestriction/GlobalSchurBounds/StandardLattice.lean}{Formalism.MainNotions.StandardLattice.PartFiveFamilyBounds}
\leanref{Formalism/SharpRestriction/GlobalSchurBounds/JointsL.lean}{Formalism.MainNotions.GlobalTheoryTheory.countF_fam}
\begin{proof}
Put $C_0=\max(5,C_{YL}+1)$, with $C_{YL}$ the constant of the family of
Lemma~\ref{asm-global-link-count-bounds-weights-hold}.
Lemma~\ref{asm-global-link-count-bounds-explicit} gives $C$ and $x_0'$ for this
$C_0$; put $x_0=\max(x_0',3)$. Let $D$ be a member at a scale $x\ge x_0$ with
$x>x_1$. By Lemma~\ref{asm-global-link-count-bounds-weights-hold}, $D$ has
weight data $(C_0,\Delta_S,w)$ at $(x,\varepsilon)$, with $w_j(\tau)$ the
Harish-Chandra weights. Let $u$ be a word. By
Definition~\ref{th-definite-datum}(4), $T^*_u\le T_u\le(1+1/(fk))T^*_u$ with
$T^*_u=X_u\Gamma_P^{-2/(fk)}$ and $\Gamma_P>1$. So
Lemma~\ref{asm-global-link-count-bounds-explicit}(1) applies with
$\Gamma=\Gamma_P$: its first inequality is (1) of Lemma~\ref{asm-global-count},
and its second, with $\log g_0+\log g_1=\log(g_0g_1)$, is (2). Part (2) of
Lemma~\ref{asm-global-link-count-bounds-explicit} is (3).
\end{proof}
% Lean: countF_fam Fm (family_part_five_of_explicit_constants Fm (division_level_prime_family Fm
% h2b h2v h3)), with countF_of at C0 = max 5 (CYL + 1): (1) to (3) are clauses of
% TorusPairsWithVariant and VandermondeVariationBound in Conclusion5₂ (Lemma
% asm-global-link-count-bounds-explicit), read with DataB from JDB.dataB_fam (x >= 3, Lemma
% asm-global-link-count-bounds-weights-hold), the window D.T_win, and log X_u = logXuD by jnt_logXu
% (factDuStar_fam, factDstar_fam). In the Lean the reading happens inside count_at_of (Lemma
% asm-global-count), which takes the clauses of Conclusion5₂ directly; the book states (1) to (3)
% as their consequences at Gamma = Gamma_P.

\subsection{The depth bound}\label{subsection-asm-depth-bound}

\begin{lemma}\label{asm-global-link-depth-sum}
Let $D$ be a definite-side datum with ranks $k,m,n$, $f=2n$, top classes $T$,
size data with the quantities $a$ and $X'$ (Definition~\ref{pd-size-data}),
and gain data (Definition~\ref{pd-gain-data}) with the sets $V(p)$, the numbers
$N(p,t_0)$, the Weyl exponents $e_p$ and the gain quantities $t(p,t_0)$, $D(p)$
and $\tau(p)$ of Definition~\ref{pd-gain-quantities}. Suppose that the size
data are filled by the gain data at every place
(Definition~\ref{pd-gain-quantities}(6)), and let $\mathrm{sz}_w$ be as in
Definition~\ref{asm-member-bounds}. Let $C\ge0$, and suppose that at every
$p<X'$: for every $t_0\in V(p)$,
\[
 \lceil t(p,t_0)\rceil\le C\,N(p,t_0)\frac ap,\qquad
 D(p)\le C\,N(p,t_0)\frac ap,\qquad 1\le N(p,t_0)\frac ap;
\]
and, if $V(p)\neq\emptyset$, $e_p\le Cfk^2\log_pa$. Then
\[
 \sum_{w\in\operatorname{supp}}\mathrm{sz}_w(\tau_w)
 \le4(2C+1)\,m\,|T|\,a\log X'+2Cfk^2\lceil X'\rceil\log a.
\]
\end{lemma}
\leanref{Formalism/SharpRestriction/GrowthBounds/TopChoiceBounds.lean}{Formalism.MainNotions.G.TopChoiceBounds.bt_Ta_at}
\leanref{Formalism/SharpRestriction/GrowthBounds/TopChoiceBounds.lean}{Formalism.MainNotions.G.TopChoiceBounds.bt_tau_pt}
\leanref{Formalism/SharpRestriction/GrowthBounds/TopChoiceBounds.lean}{Formalism.MainNotions.G.TopChoiceBounds.bt_sumNv}
\leanref{Formalism/SharpRestriction/GrowthBounds/TopChoiceBounds.lean}{Formalism.MainNotions.G.TopChoiceBounds.bt_sum_by_prime}
\leanref{Formalism/SharpRestriction/GrowthBounds/CornerGLContinuity.lean}{Formalism.MainNotions.CornerGLContinuity.CutSizes}
\begin{proof}
As $a=xb>1$, $\log a\ge0$; put $A=Cfk^2\log a\ge0$ and
$B=(2C+1)m|T|a\ge0$.

\emph{One prime.} Let $p$ be a prime with $V(p)\neq\emptyset$. Then $p$ is a
candidate prime, so $p<X'$. For $t_0\in V(p)$ the hypotheses give
$\lceil t(p,t_0)\rceil+D(p)+1\le(2C+1)N(p,t_0)a/p$. For distinct $t_0$ the sets
of $t\in T$ with $t^{s_p}=t_0$ are disjoint, so
$\sum_{t_0\in V(p)}N(p,t_0)\le|T|$. Hence
\[
 \tau(p)=e_p+m\sum_{t_0\in V(p)}\big(\lceil t(p,t_0)\rceil+D(p)+1\big)
 \le Cfk^2\log_pa+(2C+1)m|T|\frac ap,
\]
and, as $\log p\cdot\log_pa=\log a$, $\log p\cdot\tau(p)\le A+B\log p/p$. If
$V(p)=\emptyset$, then $\tau(p)=0$. So for every prime $p$, $\log p\cdot\tau(p)$
is at most $A+B\log p/p$ if $p<X'$, and $0$ otherwise.

\emph{The sum.} Since $\tau_w=\tau(p_w)$ and $\mathrm{sz}_w(y)=2\ell_we_wy$,
\[
 \sum_{w\in\operatorname{supp}}\mathrm{sz}_w(\tau_w)
 =2\sum_{w\in\operatorname{supp}}\ell_we_w\tau(p_w).
\]
For a prime $p$, the numbers $\ell_we_w$ of the places $w$ over $p$ are
nonnegative and add up to $\log p$, and $\tau(p)\ge0$. Grouping the places of
the support by $p_w$ gives
$\sum_{w\in\operatorname{supp}}\ell_we_w\tau(p_w)\le\sum_p\log p\cdot\tau(p)$,
over the primes $p_w$ with $w$ in the support. By the first step and since
$A+B\log p/p\ge0$, this is at most $\sum_{p<X'}(A+B\log p/p)$, over the primes
$p<X'$. There are at most $\lceil X'\rceil$ of them, and
Lemma~\ref{asm-division-kernel-sizes-mertens} with $y=X'\ge1$ gives
$\sum_{p<X'}\log p/p\le2\log X'$. So
\[
 \sum_{w\in\operatorname{supp}}\mathrm{sz}_w(\tau_w)\le2\big(\lceil X'\rceil A+2B\log X'\big)
 =4(2C+1)\,m\,|T|\,a\log X'+2Cfk^2\lceil X'\rceil\log a.
\]
\end{proof}
% Lean: bt_Ta_at X K C hC hM hF hT hNB: s1 (sz w y = 2 lwP w ramIdx w y and hF), s2
% (Places.bt_sum_by_prime, with factEtaMean), hpt (bt_tau_pt for VV nonempty, with bt_cand,
% bt_sumNv; tauP = 0 otherwise), s3, s4 (card of the primes below ceil X' and Places.mertens).
% The hypotheses are CutSizes D L X K C (stated there for p < X') and hNB. Here ceil is the
% ceiling in the natural numbers.

\begin{definition}\label{asm-global-link-depth-sizes}
Let $D$ be a definite-side datum with ranks $k,m,t,n$, top classes $T$, size
data with scale $(x,\varepsilon,J)$ and exponents $b_t$
(Definition~\ref{pd-size-data}), and gain data with constants $K_G$
(Definition~\ref{pd-gain-data}). Let $r_\varepsilon$, $\delta_1$, $G'$ and
$N_1$ be as in Definition~\ref{hf-good-orbit}, $L_2$ as in
Definition~\ref{pd-gain-data}(b), $p_0=k^2x^{26}$,
$\Lambda=K_\Lambda N_1^2\log a$ and $c_G=4S_c$, and let $K_{D,\mathrm I}$ and
$K_2$ be the constants of $K_G$ so named in
Lemma~\ref{pd-gain-first-hypotheses}. Let $C_e>0$ and $S$ be real numbers. The
datum has \emph{sizes bounded by $S$ with constant $C_e$} if:
\begin{enumerate}
\item $S\ge40$ and $C_e\log x>0$; each of $C_e\log x$, $1/\delta_1$, $k$, $n$,
  $r_\varepsilon$, $2\log a$, $c_G$, $x^{26}$, $(X')^{G'}$, $K_{D,\mathrm I}$
  and $K_2$ is at most $S$; and $S^{20}\le b'$;
\item $N_1\le C_e\log x$, $N_1\Lambda\le k$, $N_1\le p_0$ and
  $120\,C_e\log x\cdot p_0\le X'$;
\item $b_t\le2b'$ for every $t\in T$;
\item for every integer $n'\ge x^{\varepsilon/4}$ and every tiled set of
  positions $P\subset\{0,\dots,n'-1\}$ for $(x,\varepsilon,L_2)$
  (Definition~\ref{hf-good-orbit}),
  \[
   \frac{n'\delta_1}{C_e\log x}\le|P|\le40\,n'\delta_1.
  \]
\end{enumerate}
\end{definition}
\leanref{Formalism/SharpRestriction/GrowthBounds/TopChoiceBounds.lean}{Formalism.MainNotions.G.TopChoiceBounds.bt_H}

\begin{lemma}\label{asm-global-link-depth-sizes-hold}
Let $\mathcal F$ be a family with exponent $\varepsilon$
(Definition~\ref{pd-family}). There are $C_e>0$ and $x_0>x_1$ such that for
every $x\ge x_0$ every member at $x$ has sizes bounded by
$\exp(x^{\varepsilon/2})$ with constant $C_e$
(Definition~\ref{asm-global-link-depth-sizes}).
\end{lemma}
\leanref{Formalism/SharpRestriction/GrowthBounds/TopChoiceBounds.lean}{Formalism.MainNotions.G.TopChoiceBounds.bt_H_ev}
\leanref{Formalism/SharpRestriction/GrowthBounds/TopChoiceBounds.lean}{Formalism.MainNotions.G.TopChoiceBounds.bt_H_of}
\begin{proof}
\emph{Thresholds.} Let $B_\varepsilon$ be as in
Lemma~\ref{pd-gain-holds-orbit-sizes} and put $C_e=20B_\varepsilon>0$. The
gain data of all members have the same nonnegative constants $K_G$, among the
constants of the family; let $J$ be the number of the family. Choose
$x_0>x_1$ with $x_0\ge3$, at least the threshold of
Lemma~\ref{pd-kappa-orbit-large-scale} for $(\varepsilon,J)$, such that for
$x\ge x_0$, with $y=x^{\varepsilon/2}$:
\[
 48\le y,\quad 26\log x\le y,\quad C_e\log x\le x^{\varepsilon/24},\quad
 r_\varepsilon\le y,\quad K_{D,\mathrm I}\le y,\quad K_2\le y,\quad
 \tfrac32\le x^{\varepsilon/4},
\]
\[
 K_\Lambda\le x^{1/2},\quad C_e\log x\le x^{1/2},\quad 120\,C_e\log x\le x,
 \quad 37\log x\le x^{\varepsilon/3}.
\]
All large $x$ qualify, since a positive power of $x$ eventually exceeds every
constant multiple of $\log x$.

\emph{The member.} Let $x\ge x_0$, let $D$ be a member at $x$, and put
$y=x^{\varepsilon/2}$ and $S=e^y\ge1+y$. As $x>x_1$, the size data of $D$
have scale $(x,\varepsilon,J)$, and the size bounds of the family
(Definition~\ref{pd-family}(4)) give, with $N^*=2n\delta_1$:
$n\le x^{2+8\varepsilon+\varepsilon/100}$;
$x^{9/2+10\varepsilon-\varepsilon/100}\le k\le
x^{9/2+10\varepsilon+\varepsilon/100}$; $\Lambda_G\ge x^{\varepsilon}$ and
$4N^*x^{\varepsilon}\Lambda_G/k+x^{\varepsilon/100}\log x\,\Lambda_G/k\le
x^{-2}$; $b_t\le b'+N^*x^{\varepsilon}\Lambda_Gb'/k$ for $t\in T$; and
$\Delta_S\ge N(m)m^{10}>0$ with $\log\Delta_S\le x^{\varepsilon/4}/2$. As
$\varepsilon<1/52$, $x^4\le k\le x^5$ and $n\le x^3$. The size data give
$S_c=4k\Delta_S$.

\emph{Part (1).} First $S\ge1+y>40$. As $x>1$, $C_e\log x>0$, and
$C_e\log x\le x^{\varepsilon/24}\le y\le S$. Next
$1/\delta_1=x^{\varepsilon/8}\le y\le S$, and $r_\varepsilon$,
$K_{D,\mathrm I}$ and $K_2$ are at most $y\le S$. As $26\log x\le y$,
$x^{26}\le S$, $k\le x^5\le S$ and $n\le x^3\le S$. Since $a=xb$,
$\log a=\log x+x^{\varepsilon}=\log x+y^2$; as $y\ge48$, $48y^2\le y^3\le6S$,
and $\log x\le y/26\le y^2$, so $2\log a\le4y^2\le S$. Next
$c_G=4S_c=16k\Delta_S$, with $16\le1+y/3\le e^{y/3}$, $k\le x^5\le e^{y/3}$
since $15\log x\le y$, and, with $q=x^{\varepsilon/4}\ge3/2$ and $q^2=y$,
$\log\Delta_S\le q/2\le q^2/3=y/3$; so $c_G\le S$.
Lemma~\ref{pd-gain-holds-orbit-sizes}(1) gives $G'\delta_1\le B_\varepsilon\log
x$, so $G'\le C_e\log x\cdot x^{\varepsilon/8}\le
x^{\varepsilon/24+\varepsilon/8}$, hence $G'x^{\varepsilon/3}\le y$ and
$(X')^{G'}=\exp(G'x^{\varepsilon/3})\le S$. Finally
$b'=a/X'=x\exp(x^{\varepsilon}-x^{\varepsilon/3})$; as $x^{\varepsilon/3}\le
y$ and $21y\le y^2=x^{\varepsilon}$, $20y\le x^{\varepsilon}-x^{\varepsilon/3}$,
so $S^{20}\le b'$ since $x\ge1$.

\emph{Part (2).} Lemma~\ref{pd-gain-holds-orbit-sizes}(1) gives $N_1\le16.8\,
G'\delta_1\le C_e\log x$. Put $z=x^{1/2}$. Then $N_1\le C_e\log x\le z$,
$K_\Lambda\le z$ and $\log a=\log x+x^{\varepsilon}\le2x=2z^2$, so
$N_1\Lambda=K_\Lambda N_1^3\log a\le2z^6\le z^8=x^4\le k$, as $z^2=x\ge2$.
Next $N_1\le z\le x^{26}\le k^2x^{26}=p_0$. As $k^2\le x^{10}$,
$120\,C_e\log x\cdot p_0\le x\cdot x^{10}\cdot x^{26}=\exp(37\log x)\le
\exp(x^{\varepsilon/3})=X'$.

\emph{Part (3).} The term $x^{\varepsilon/100}\log x\,\Lambda_G/k$ is
nonnegative, so $N^*x^{\varepsilon}\Lambda_G/k\le x^{-2}/4\le1$, and
$b_t\le b'+b'=2b'$, as $b'\ge0$.

\emph{Part (4).} Let $n'\ge x^{\varepsilon/4}$. Then $3G'\le n'$ by
Lemma~\ref{pd-kappa-orbit-large-scale}(3), whose hypothesis on $k$ holds. For a
tiled set of positions $P$ for $(x,\varepsilon,L_2)$,
Lemma~\ref{pd-gain-holds-orbit-sizes}(2) gives $n'\delta_1/(20B_\varepsilon\log
x)\le|P|\le40n'\delta_1$, which is (4) as $C_e=20B_\varepsilon$.
\end{proof}
% Lean: bt_H_ev Fm: Ce = 20 logBoundConstant r h from master_orbit_sizes_hold (hA0 =
% segment_count_log_bounds, hB0 = period_ge_three_gap_eventually), and the thresholds n1 to n9 from
% const_le_rpow_eventually and mul_log_le_rpow_eventually_alt; then bt_H_of with Fm.sizes (SizeN,
% SizeK, SizeLa, SizeRho, SizeBot, SizeDel), L.Sc_eq, bt_rG and bt_hG. The book's (1) to (4) are the
% fields of bt_H: forty_le_bound, cl0, cl, de, k, n, r, la, c, xA, Q, KDI,
% integral_constant_le_bound, b'; N1, La, firstLength_le_prime_bound, pX; bot; cnt.

\begin{lemma}\label{asm-global-link-depth-representative}
Let $D$ be a definite-side datum with size data and gain data that has sizes
bounded by $S$ with constant $C_e$
(Definition~\ref{asm-global-link-depth-sizes}), with the notation there and in
Definition~\ref{pd-gain-data}. Let $p<X'$ be a prime and $t_0\in V(p)$, and put
$N=N(p,t_0)$ and $N^+=40n_p\delta_1$. Then $p$ is a candidate prime, and:
\begin{enumerate}
\item $S^{20}\le a/p$ and $b'\le a/p$;
\item $n_p\le S$, $1\le N\le S^2$ and $n_p\le S^2N$;
\item $N^+\le40SN$, $N^+\le40S$ and $N^+\le40\,C_e\log x\cdot N$;
\item $0\le\log_pa\le S$ and $p^{G'}\le S$;
\item $n_pc_G\le Na/p$ and $Na/p\ge1$.
\end{enumerate}
\end{lemma}
\leanref{Formalism/SharpRestriction/GrowthBounds/TopChoiceBounds.lean}{Formalism.MainNotions.G.TopChoiceBounds.bt_O_of}
\leanref{Formalism/SharpRestriction/GrowthBounds/TopChoiceBounds.lean}{Formalism.MainNotions.G.TopChoiceBounds.vertex_density_lower_bound}
\leanref{Formalism/SharpRestriction/GrowthBounds/TopChoiceBounds.lean}{Formalism.MainNotions.G.TopChoiceBounds.bt_NB}
\begin{proof}
By Definition~\ref{pd-gain-data}(e), $V(p)$ is empty unless $p$ is a candidate
prime, and then it consists of good representatives. So $p$ is a candidate
prime: $p<X'$, $p$ is of inert class and $n_p\ge x^{\varepsilon/4}$; and
$P(p,t_0)$ is a tiled set of positions for $(x,\varepsilon,L_2)$. As $p$ is of
inert class, $N=|P(p,t_0)|$ (Step 1 of the proof of
Lemma~\ref{pd-gain-holds-orbit-condition}). By (4) of
Definition~\ref{asm-global-link-depth-sizes},
$n_p\delta_1/(C_e\log x)\le N\le40n_p\delta_1$.

(1) As $p<X'$, $b'=a/X'\le a/p$, and $S^{20}\le b'$.

(2) The order $n_p$ of $h_p$ divides the order $n$ of $A$, so $n_p\le n\le S$.
As $N\ge n_p\delta_1/(C_e\log x)>0$ and $N$ is an integer, $N\ge1$. As
$\delta_1\le1$ and $S\ge40$, $N\le40n_p\delta_1\le40S\le S^2$. The lower bound
on $N$ gives $n_p\delta_1\le C_e\log x\cdot N\le SN$, and with $1/\delta_1\le
S$, $n_p\le S^2N$.

(3) The same bound gives $N^+\le40\,C_e\log x\cdot N\le40SN$, and $n_p\le S$
and $\delta_1\le1$ give $N^+\le40S$.

(4) As $a>1$ and $p\ge2$, $\log_pa\ge0$; as $\log p\ge\log2\ge1/2$,
$\log_pa\le2\log a\le S$. As $p<X'$, $p^{G'}\le(X')^{G'}\le S$.

(5) As $c_G=4S_c\ge0$ and $c_G\le S$, (2) and (1) give $n_pc_G\le
S^2N\cdot S\le NS^{20}\le Na/p$. Also $Na/p\ge S^{20}\ge1$.
\end{proof}
% Lean: bt_O_of X K Ce S H p hpX t0 ht0 gives bt_O (fields hp, hcl, hrep, SB, bB, nS, N1, NS, nN, Np,
% NpS, lb0, lb, Q) and NplusF <= 40 Ce log x Nv; it uses bt_cand, bt_good, H.cnt and
% vertex_count_eq_card_tops, addOrderOf_dvd_card with card_A, bt_nN and bt_logb_le. Part (5) is
% vertex_density_lower_bound and bt_NB. N^+ is NplusF L.sc (addOrderOf (hP D p)).

\begin{lemma}\label{asm-global-link-depth-cuts}
Let $D$ be a definite-side datum with ranks $k,m,t,n$ and $f=2n$, size data,
and gain data whose constants $K_G$ are nonnegative, with the gain quantities
$t(p,t_0)$ and $D(p)$ and the Weyl exponents $e_p$ of
Definition~\ref{pd-gain-quantities}. Let $K_t,K_D$ be the constants of $K_G$
so named in Lemma~\ref{pd-gain-second-hypotheses}, and
$K_{D,\mathrm I},K_1,K_2$ those so named in
Lemma~\ref{pd-gain-first-hypotheses}. Suppose that $D$ has sizes bounded by $S$
with constant $C_e$ (Definition~\ref{asm-global-link-depth-sizes}), and that
$C_9\ge0$ satisfies $e_p\le C_9fk^2\log_pa$ for every prime $p\le a$. Put
\[
 C'=6K_tr_\varepsilon+5r_\varepsilon+K_D+2K_1+K_2+K_{D,\mathrm I}+6+C_9.
\]
Then for every prime $p<X'$ and every $t_0\in V(p)$,
\[
 \lceil t(p,t_0)\rceil\le C'N(p,t_0)\frac ap,\qquad
 D(p)\le C'N(p,t_0)\frac ap;
\]
and for every prime $p<X'$ with $V(p)\neq\emptyset$, $e_p\le C'fk^2\log_pa$.
\end{lemma}
\leanref{Formalism/SharpRestriction/GrowthBounds/TopChoiceBounds.lean}{Formalism.MainNotions.G.TopChoiceBounds.bt_at}
\leanref{Formalism/SharpRestriction/GrowthBounds/TopChoiceBounds.lean}{Formalism.MainNotions.G.TopChoiceBounds.bt_II}
\leanref{Formalism/SharpRestriction/GrowthBounds/TopChoiceBounds.lean}{Formalism.MainNotions.G.TopChoiceBounds.bt_I}
\leanref{Formalism/SharpRestriction/GrowthBounds/TopChoiceBounds.lean}{Formalism.MainNotions.G.TopChoiceBounds.bt_C}
\begin{proof}
Write $r=r_\varepsilon$. Let $p<X'$ be a prime and $t_0\in V(p)$, and use the
notation of Lemma~\ref{asm-global-link-depth-representative}, with $Z=Na/p$.
That lemma shows that $p$ is a candidate prime, so of inert class, and its
orbit length (Definition~\ref{pd-orbit-places}(3)) is $n_p$. By
Definition~\ref{pd-gain-quantities}, $D(p)=D_{\mathrm{II}}(p,n_p)$ and
$t(p,t_0)=t_{\mathrm{II}}(p,n_p,N)$ if $p<p_0$; otherwise
$D(p)=D_{\mathrm I}(p,n_p)$ and
$t(p,t_0)=t_{\mathrm I}(p,n_p,P,a_w,b_w)$ with $P=P(p,t_0)$,
$a_w(i)=\mathbf a(t_0+ih_p)$ and $b_w(i)=\mathbf b(t_0+ih_p)$. By part (5) of
that lemma, $Z\ge1$ and $n_pc_G\le Z$. So if $s\le cZ$ with $c\ge0$ and
$c+1\le C'$, then $\lceil s\rceil\le cZ+1\le(c+1)Z\le C'Z$.

\emph{Common bounds.} As $K_\Lambda\ge0$, $a>1$ and $N_1\ge1$,
$0\le\Lambda\le N_1\Lambda\le k$, so $\Lambda a/(pk)\le a/p$. As $40S\le S^2$,
$2c_G\le2S\le S^{20}\le b'$, so $c_0=2b'+2c_G\le3b'\le3a/p$. With $Q=p^{G'}$,
$kQ\le S^2\le a/p$. If $0\le E\le S^8$, then, by parts (3) and (4),
$N^+E\log_pa\le40SN\cdot S^8\cdot S\le S^{20}N\le Z$.

\emph{Primes $p<p_0$.} As $kx^{26}>1$, $L_1(p)\ge1$ and $ap^{-L_1(p)}\le a/p$.
So
\[
 t_{\mathrm{II}}(p,n_p,N)=K_trN\Big(\frac{\Lambda a}{pk}+ap^{-L_1(p)}+c_0+kQ\Big)
 +2rn_pc_G\le(6K_tr+2r)Z,
\]
and $\lceil t(p,t_0)\rceil\le C'Z$, as $6K_tr+2r+1\le C'$. Next
$N_1^3+r^5+k^2\le S^3+S^5+S^2\le S^8$, so the first term of
$D_{\mathrm{II}}(p,n_p)$ is at most $K_DZ$, and its second term $2rn_pc_G$ is
at most $2rZ$. For the third, $N^+\le40\,C_e\log x\cdot N$, $c_0\le3b'=3a/X'$
and $120\,C_e\log x\cdot p\le120\,C_e\log x\cdot p_0\le X'$ give
$rN^+c_0\le rN\cdot120\,C_e\log x\cdot a/X'\le rZ$. So
$D_{\mathrm{II}}(p,n_p)\le\lceil(K_D+3r)Z\rceil\le C'Z$.

\emph{Primes $p\ge p_0$.} Put $E=x^{26}(N_1^5+k^3)$; then $0\le E\le
S(S^5+S^3)\le S^8$. So $D_{\mathrm I}(p,n_p)=\lceil K_{D,\mathrm
I}N^+E\log_pa\rceil\le\lceil K_{D,\mathrm I}Z\rceil\le C'Z$, and also
$K_{D,\mathrm I}N^+E\log_pa\le S\cdot40S\cdot S^8\cdot S\le S^{12}$, so
$D_{\mathrm I}(p,n_p)\le S^{13}$. In $t_{\mathrm I}(p,n_p,P,a_w,b_w)$
(Lemma~\ref{pd-gain-first-hypotheses}), each term of the sum over the odd
$y\in P$ is $0$ or at most $b_w(y)=b_{t_0+yh_p}\le2b'\le2a/p$, as
$t_0+yh_p\in T$; there are at most $|P|=N$ of them, so the sum is at most
$2Z$. As $N_1\le p_0\le p$ and $N_1\Lambda\le k$,
$N_1(a/p^2+\Lambda a/(pk))\le2a/p$. In
$c_6=8(D_{\mathrm I}+1)+K_2N^+k\log_pa+4N^+c_G+2n_pc_G$ the bounds
$D_{\mathrm I}\le S^{13}$, $N^+\le S^2$ and $K_2,k,\log_pa,c_G,n_p,Q\le S$ give
\[
 Qc_6+kQ+k\log_pa+c_G\le8S^{14}+S^6+4S^4+2S^3+2S^2+9S\le S^{15}\le\frac ap,
\]
using $40S^i\le S^{i+1}$. With $2n_pc_G\le2Z$,
\[
 t_{\mathrm I}(p,n_p,P,a_w,b_w)\le2Z+2K_1Z+K_2Z+2Z=(4+2K_1+K_2)Z,
\]
and $\lceil t(p,t_0)\rceil\le C'Z$, as $5+2K_1+K_2\le C'$.

\emph{Weyl exponents.} Let $p<X'$ with $V(p)\neq\emptyset$. Then $p$ is a
candidate prime, and $X'\le a$ since $b'=a/X'\ge S^{20}\ge1$; so $p\le a$ and
$e_p\le C_9fk^2\log_pa\le C'fk^2\log_pa$, as $\log_pa\ge0$ and $C_9\le C'$.
\end{proof}
% Lean: bt_at X K hK Ce S Ce9 H h90 h9, with constant bt_C K (rG L.sc) Ce9; the cut bounds are bt_II
% (p < p0: bt_fin, bt_real_tII, bt_real_DII, bt_real_NE, pow_sum_le_pow_eight, mul_div_le_div_of_le,
% mul_rpow_neg_le_div with one_le_first_length_at_scale, initial_constant_bounds and
% mul_le_of_density_bounds) and bt_I (p >= p0: bt_DIS, mul_pow_add_pow_le_pow_eight, bt_real_tI,
% density_sum_le_two_div, bt_real_br, H.bot); tv and Dp use nP D p, which is addOrderOf (hP D p) at
% primes of inert class (period_eq_addOrderOf). The constants K_t, K_D, K_{D,I}, K_1, K_2, K_Lambda
% are K.tII, K.DII, K.DI, K.kappaThresholdOne, K.kappaThresholdTwo, K.La; their nonnegativity is
% K.Nonneg. The sum over odd y in t_I is over TT.filter (not IsTypeV), with the term 0 when dMinus = 0.

\begin{lemma}\label{asm-global-link-depth}
Every family $\mathcal F$ with exponent $\varepsilon$
(Definition~\ref{pd-family}) satisfies the depth bound of
Lemma~\ref{asm-global-slopes-n}: there are $C_T$, $x_T$ and a function $\delta$
with $\delta(x)\to0$ as $x\to\infty$ such that, with the notation of
Definition~\ref{asm-member-bounds},
\[
 \sum_{w\in\operatorname{supp}}\mathrm{sz}_w(\tau_w)\le C_T\,m\,|T|\,a\log X'+b^{\delta(x)}
\]
at every member at every scale $x\ge x_T$.
\end{lemma}
\leanref{Formalism/SharpRestriction/GrowthBounds/TopChoiceBounds.lean}{Formalism.MainNotions.G.TopChoiceBounds.tau_sum_family_of_mertens}
\leanref{Formalism/SharpRestriction/GrowthBounds/TopChoiceBounds.lean}{Formalism.MainNotions.G.TopChoiceBounds.stmtFTa_proved}
\leanref{Formalism/SharpRestriction/GrowthBounds/CornerGLContinuity.lean}{Formalism.MainNotions.CornerGLContinuity.WeightSumBound}
\begin{proof}
At a scale $x$ put $S=\exp(x^{\varepsilon/2})$; then, with
$\delta(x)=8x^{-\varepsilon/2}$, $b^{\delta(x)}=\exp(8x^{\varepsilon/2})=S^8$,
and $\delta(x)\to0$.

\emph{Constants.} There are $C\ge0$ and $x_0$ such that at every member at
every scale $x\ge x_0$ with $x>x_1$: the hypotheses of
Lemma~\ref{asm-global-link-depth-sum} on $t(p,t_0)$, $D(p)$, $N(p,t_0)$ and
$e_p$ hold with $C$; and $S\ge40$, $n\le S$, $k\le S$ and $2\log a\le S$.
Indeed, by Lemma~\ref{asm-global-link-depth-sizes-hold} there are $C_e>0$ and
$x_0>x_1$ such that every member at every scale $x\ge x_0$ has sizes bounded
by $S$ with constant $C_e$; this gives $S\ge40$, $n\le S$, $k\le S$ and
$2\log a\le S$. Among the choices of the family
(Definition~\ref{pd-family}(4)), the gain data of all members have the same
nonnegative constants $K_G$, and their Weyl exponents are
$e_p=\lceil C_wfk^2\log_pa\rceil$ for a constant $C_w$ of the family. Put
$C_9=\max(C_w,0)+1$. For a prime $p\le a$, $\log_pa\ge1$ and $f,k\ge1$, so
$Z=fk^2\log_pa\ge1$; if $C_wZ\ge0$, then $e_p<C_wZ+1\le\max(C_w,0)Z+Z$, and
otherwise $e_p=0$; so $e_p\le C_9fk^2\log_pa$. Let $C$ be the constant $C'$ of
Lemma~\ref{asm-global-link-depth-cuts} for these $K_G$ and $C_9$; it depends
only on the family, as $r_\varepsilon$ does, and $C\ge6$. That lemma gives the
bounds on $t(p,t_0)$, $D(p)$ and $e_p$ with $C$, and
Lemma~\ref{asm-global-link-depth-representative}(5) gives $N(p,t_0)a/p\ge1$
for $t_0\in V(p)$ with $p<X'$.
% Lean: stmtFTa_proved takes hT from family_exponent_sizes_hold Fm (bt_H_ev, then bt_at with
% Ce9 = max Fm.cX.Ce 0 + 1 from (Fm.ex x hx1 i).ep_def, and bt_rG), passes to max C 0 in hT',
% takes hNB from bt_NB with bt_O_of, and the sizes H.forty_le_bound, H.n, H.k, H.la from bt_H_ev.
Enlarging $x_0$, also $C\le x^{\varepsilon/2}\le S$ for $x\ge x_0$, since
$x^{\varepsilon/2}\to\infty$ and $x^{\varepsilon/2}+1\le\exp(x^{\varepsilon/2})$.

\emph{The bound.} Let $D$ be a member at a scale $x\ge x_0$ with $x>x_1$. Among
the choices of the family (Definition~\ref{pd-family}(4)), its size data are
filled by its gain data at every place. So
Lemma~\ref{asm-global-link-depth-sum} gives
\[
 \sum_{w\in\operatorname{supp}}\mathrm{sz}_w(\tau_w)
 \le4(2C+1)\,m\,|T|\,a\log X'+2Cfk^2\lceil X'\rceil\log a.
\]
As $x\ge1$, $1\le X'=\exp(x^{\varepsilon/3})\le S$, so $\lceil X'\rceil\le
X'+1\le2S$. With $f=2n\le2S$, $k\le S$, $0\le\log a\le S$ and $0\le C\le S$,
\[
 2Cfk^2\lceil X'\rceil\log a\le2\cdot S\cdot2S\cdot S^2\cdot2S\cdot S=8S^6\le S^7\le S^8,
\]
since $S\ge40$. This is the depth bound with $C_T=4(2C+1)$, $x_T=\max(x_0,x_1+1)$
and $\delta(x)=8x^{-\varepsilon/2}$.
\end{proof}
% Lean: tau_sum_family_of_mertens Fm (fun _ _ => mertens) = stmtFTa_proved Fm hM
% (family_exponent_sizes_hold Fm): C_T = 4(2 max(C,0) + 1), delta = 8 x^(-eps/2) (tendsto via
% tendsto_rpow_neg_atTop), hCy from const_le_rpow_eventually, hmain = bt_Ta_at (Lemma
% asm-global-link-depth-sum) with hF from Fm.fills, hb (b^delta = exp(x^(eps/2))^8), and bt_bo with
% bt_X'S and bt_chain. Size of the whole: proof closure of tau_sum_family_of_mertens, about 60
% theorems and 1,100 lines in 3 files, most of it in family_exponent_sizes_hold and bt_H_ev.

\subsection{Losses and the global link}\label{subsection-asm-losses-global-link}

\begin{lemma}\label{asm-global-link-bad-primes}
Let $\mathcal F$ be a family with exponent $\varepsilon$
(Definition~\ref{pd-family}), and let $C_b\ge0$. For a member at a scale
$x>x_1$, with $X'=\exp(x^{\varepsilon/3})$, let $P$ be its set of large-order
residue primes, which are primes $p<X'$, and for $p\in P$ let $\beta(p)\ge0$ be
the number of its bad tops at $p$, as in the bound on the bad tops of
Definition~\ref{asm-designed}(1). Suppose that for all large $x$ every member
at $x$ satisfies $\delta_1n/2\le|T|$ and the bound on the bad tops with $C_b$:
\[
 \sum_{p\in P}\frac{\log p}p\,\beta(p)
 \le4C_bx^{-2}\delta_1n\sum_{p\in P}\frac{\log p}p.
\]
Then for all large $x$ every member at $x$ satisfies
\[
 \sum_{p\in P}\frac{\log p}p\,\beta(p)\le|T|\,x^{-2+\varepsilon}.
\]
\end{lemma}
\leanref{Formalism/SharpRestriction/GlobalSchurBounds/LossesSmallShares.lean}{Formalism.MainNotions.Lnk.chosenBadp_of_old}
\leanref{Formalism/SharpRestriction/GlobalSchurBounds/Chosen.lean}{Formalism.MainNotions.ChBad_at}
\leanref{Formalism/SharpRestriction/GlobalSchurBounds/WeightVectorSmallSharesSeven.lean}{Formalism.MainNotions.ChBadp_at}
\begin{proof}
Let $x$ be large enough that the hypotheses hold at $x$, that $x>x_1$ and
$x\ge1$, and that $16C_b\le x^{2\varepsilon/3}$; all large $x$ qualify. Fix a
member at $x$. As $P$ consists of primes below $X'$ and every $\log p/p$ is
nonnegative, Lemma~\ref{asm-division-kernel-sizes-mertens} gives
\[
 0\le\sum_{p\in P}\frac{\log p}p\le\sum_{p<X'}\frac{\log p}p\le2\log X'=2x^{\varepsilon/3}.
\]
As $0\le\delta_1n\le2|T|$ and $C_b\ge0$, the bound on the bad tops gives
\[
 \sum_{p\in P}\frac{\log p}p\,\beta(p)\le4C_bx^{-2}\cdot2|T|\cdot2x^{\varepsilon/3}
 =|T|\,x^{-2}\cdot16C_bx^{\varepsilon/3}\le|T|\,x^{-2}x^{2\varepsilon/3}x^{\varepsilon/3}
 =|T|\,x^{-2+\varepsilon}.
\]
\end{proof}
% Lean: chosenBadp_of_old Fm Cb hCb cTT cBad: threshold x1 + 1 + |a0| + |a1| + |a2| with a2 from
% ev_le_rpow (2 eps / 3) (16 Cb); hX' (X' = exp(x^(eps/3))), Places.mertens with hsub, hS; hdn from
% the lower bound of ChTT_at; step and hfin. P is largeOrderResiduePrimes (primes p < ceil X' with
% IsResidueClassZero and order of hP at least x^(eps/4)) and beta(p) is the card of badTops X p.

\begin{definition}\label{asm-global-link-losses-numbers}
Let $\mathcal F$ be a family with exponent $\varepsilon$
(Definition~\ref{pd-family}), let $C_G$ be real and $\delta$ a real function,
and let $D$ be a member at a scale $x>x_1$, with the notation of
Definitions~\ref{asm-member-bounds}, \ref{pd-size-data}, \ref{pd-gain-data}
and~\ref{pd-gain-link}: ranks $k,m,n,t$, $f=2n$, $N(m)=k+m(t+1)$, top classes
$T$, the quantities $b,a,b',w'',X',X'',\delta_1,p_0,\eta,S_c,\Lambda_G$, the
number $\rho$ of Definition~\ref{pd-gain-data}(c) and the deviations
$\Sigma(p)$. Let $e$ be the $2$-power exponent of the Galois datum of $D$
(Definition~\ref{th-ranks}), $d$ the auxiliary prime of its host, $C_r\ge1$
the constant of the family that bounds the dropped primes, and
$r(\varepsilon)=\lceil8(4+8\varepsilon)/\varepsilon\rceil+1$. The \emph{loss
numbers} of $D$ for $(C_G,\delta)$ are
\[
 C_P=\max(C_G,0),\quad b_o=b^{\delta(x)},\quad p_0'=\lceil p_0\rceil,\quad
 \theta_h=\frac1{kx^{13}},\quad K_0=k^2x^{78},\quad S_k=S_c+k+\eta,
\]
\[
 Q_1=C_Pr(\varepsilon)b'\log x,\qquad
 Q_2=C_P\Big(\frac{b'}x+\eta+S_c+\frac{(w''+t)\log x}{\delta_1}\Big)+b_o,
\]
\[
 Y=X'(1+C_r\rho),\qquad Y'=\max(Y,X''),\qquad
 J_o=\lfloor x^{\varepsilon/4}\rfloor,\qquad I_b=\lfloor x^{\varepsilon/3}\rfloor.
\]
Let $\mathrm{Pr}$ be the set of inert primes
$p<\max\big(\lfloor a+a/(kx^{13})\rfloor,\lfloor X''\rfloor\big)+1$. A prime
$p\in\mathrm{Pr}$ is \emph{kept} if it is not dropped
(Definition~\ref{pd-size-data}(4)). A dropped prime $p\in\mathrm{Pr}$ is
\emph{high} if $p>X''$; of the \emph{second kind} if $p\le X''$, $p$ is a
candidate prime and $\Sigma(p)>x^{-\varepsilon}$; and of the \emph{first kind}
otherwise.
\end{definition}
\leanref{Formalism/LocalToGlobal/GlobalTheory/SuppSat.lean}{Formalism.L2.GlobalTheoryTheory.Loss.par}
\leanref{Formalism/LocalToGlobal/GlobalTheory/SuppSat.lean}{Formalism.L2.GlobalTheoryTheory.Loss.ParTie}
\leanref{Formalism/LocalToGlobal/GlobalTheory/SuppSat.lean}{Formalism.L2.GlobalTheoryTheory.Loss.PsAt}
\leanref{Formalism/LocalToGlobal/GlobalTheory/SuppSat.lean}{Formalism.L2.GlobalTheoryTheory.Loss.clAt}
\leanref{Formalism/LocalToGlobal/GlobalTheory/SuppSat.lean}{Formalism.L2.GlobalTheoryTheory.Loss.ldAt}
\leanref{Formalism/SizeWindowBrackets/WindowSettingSatisfiable.lean}{Formalism.LFunctionChainParameters.PrimeWindowSizeBounds.Par}
\leanref{Formalism/SizeWindowBrackets/SecondKindBucket.lean}{Formalism.LFunctionChainParameters.DeficitSplitLemmas.PClass}

\begin{lemma}\label{asm-global-link-losses-terms}
Let $a\ge1$, $0\le a'\le\frac32a$ and $w''>0$ be real numbers. For a real
number $q\ge2$, a finite set $I$ of positive integers and real numbers
$w_i\ge0$ ($i\in I$) put $A(q,I,w)=\sum_{i\in I}\max(a'/q^i-w_i,0)$.
\begin{enumerate}
\item If $p\ge2$ is an integer and every $i\in I$ is odd, then $2\log
  p\cdot A(p,I,w)\le4a\log p/p$.
\item If $d\ge4$ is an integer, then $\log d\cdot A(d,I,w)\le2a\log d/d$.
\item Let $0\le\theta\le1$ and $X''\ge2$ with $X''w''=a$ and $a'\le a(1+\theta)$,
  and let $S$ be a finite set of integers $p>X''$, with finite sets $I_p$ of
  positive integers and weights $w_{p,i}\ge w''$. Then
  \[
   \sum_{p\in S}2\log p\cdot A(p,I_p,w_p)\le2\log(2X'')\big(a\theta^2+w''\theta\big).
  \]
\item Let $C,f,S_k,\mu\ge0$ and $N$ be real, and let $S$ be a finite set of
  integers $p\ge2$ with $|S|\le X''$. Then
  \[
   \sum_{p\in S}2\log p\,\big(CfN^2\log_pa+Cf(\log_pa)^2\big)\le2CfX''\big(N^2\log
   a+2(\log a)^2\big),\qquad
   \sum_{p\in S}2\log p\cdot\mu CS_k\log_p(2a)\le2CS_k\log(2a)\,\mu X''.
  \]
\item Let $X'\ge1$, $Q_1,Q_2\ge0$ and $p_0\ge3$ an integer, and let $S$ be a
  finite set of primes $p<X'$. Then
  \[
   \sum_{p\in S,\ p<p_0}2Q_1\log p+\sum_{p\in S,\ p\ge p_0}2Q_2\log p
   \le2\log4\cdot(p_0Q_1+X'Q_2).
  \]
  If moreover $C,l_x,\Lambda\ge0$, $k,K_0>0$, and $L_1(p)$ are real numbers with
  $K_0\le p^{L_1(p)-1}$ for the $p\in S$ with $p<p_0$, then
  \[
   \sum_{p\in S}2\log p\cdot Cl_x\Big(\frac{\Lambda a}{pk}+o(p)\Big)\le
   2Cl_xa\Big(\Big(\frac\Lambda k+\frac1{K_0}\Big)(\log X'+\log4)
   +\frac{2(\log(p_0-1)+1)}{p_0-1}\Big),
  \]
  where $o(p)=ap^{-L_1(p)}$ if $p<p_0$ and $o(p)=a/p^2$ otherwise.
\end{enumerate}
\end{lemma}
\leanref{Formalism/SizeWindowBrackets/ChainSlopeWindows/AdmissibleSchurLemma.lean}{Formalism.LFunctionChainParameters.DeficitSplitLemmas.top_loss_bound}
\leanref{Formalism/SizeWindowBrackets/ChainSlopeWindows/AdmissibleSchurLemma.lean}{Formalism.LFunctionChainParameters.DeficitSplitLemmas.d_term_bound}
\leanref{Formalism/SizeWindowBrackets/ChainSlopeWindows/AdmissibleSchurLemma.lean}{Formalism.LFunctionChainParameters.DeficitSplitLemmas.high_bucket_bound}
\leanref{Formalism/SizeWindowBrackets/ChainSlopeWindows/AdmissibleSchurLemma.lean}{Formalism.LFunctionChainParameters.DeficitSplitLemmas.floor_bucket_bound}
\leanref{Formalism/SizeWindowBrackets/ChainSlopeWindows/AdmissibleSchurLemma.lean}{Formalism.LFunctionChainParameters.DeficitSplitLemmas.cp_bucket_bound}
\leanref{Formalism/SizeWindowBrackets/ChainSlopeWindows/AdmissibleSchurLemma.lean}{Formalism.LFunctionChainParameters.DeficitSplitLemmas.pr_bucket_bound}
\leanref{Formalism/SizeWindowBrackets/ChainSlopeWindows/AdmissibleSchurLemma.lean}{Formalism.LFunctionChainParameters.DeficitSplitLemmas.or_bucket_bound}
\begin{proof}
These are computations, given here with their results. (1) and (2) bound
$A(q,I,w)$ by the geometric sum $a'\sum_iq^{-i}$ over the odd, respectively all
positive, exponents: $a'p/(p^2-1)\le\frac43a'/p$ and $a'/(d-1)$. (3) uses that
a term is nonzero only if $p<X''(1+\theta)$, and then only for $i=1$, with
$a'/p-w_{p,1}\le a(1+\theta)/X''-w''=w''\theta$ and the number of such
integers $p$, at most $X''\theta+1$. (4) bounds each term by its value with
$\log p\ge\log2\ge\frac12$, and multiplies by $|S|\le X''$. (5) uses Chebyshev's
bound $\sum_{p\le y}\log p\le y\log4$, Mertens' bound $\sum_{p\le
y}\log p/p\le\log y+\log4$ and $\sum_{n\ge p_0}\log n/n^2\le
2(\log(p_0-1)+1)/(p_0-1)$.
\end{proof}

\begin{lemma}\label{asm-global-link-losses-explicit}
Let $\mathcal F$ be a family with exponent $\varepsilon$
(Definition~\ref{pd-family}) that satisfies the gain link
(Definition~\ref{pd-gain-link}) with constant $C_G$ and function $\delta$. Then
for all large $x$, every member $D$ at $x$, with its loss numbers for
$(C_G,\delta)$ and its classes of primes
(Definition~\ref{asm-global-link-losses-numbers}), with $P$ and the bad tops as
in Lemma~\ref{asm-global-link-bad-primes}, $T(p)$ the counted classes at $p$
(Definition~\ref{pd-gain-quantities}(2)), and the notation of
Definition~\ref{asm-member-bounds}, has the following properties.
\begin{enumerate}
\item For every integer $N$ with $a_t/w_j(\rho)\le N$ for all $t\in T$, all
  lifts $j$ and all $\rho\in\Psi_A$,
  \begin{align*}
   \sum_{t\in T}\sum_{j=1}^m\mathrm{ASK}_N(a_t,t,w_j)-G_a
   &\le2mC_PfX''\big(N(m)^2\log a+2(\log a)^2\big)
   +2C_PS_k\log(2a)\,m|T|X''\\
   &\quad+2\log4\cdot m|T|\big(p_0'Q_1+X'Q_2\big)\\
   &\quad+2C_Pm|T|a\log x\Big(\Big(\frac{\Lambda_G}k+\frac1{K_0}\Big)
   (\log X'+\log4)+\frac{2(\log(p_0'-1)+1)}{p_0'-1}\Big)\\
   &\quad+\sum_{p\text{ kept}}4ma\frac{\log p}p\,|T\smallsetminus T(p)|
   +4m|T|a\sum_{p\text{ first kind}}\frac{\log p}p
   +4m|T|a\sum_{p\text{ second kind}}\frac{\log p}p\\
   &\quad+2m|T|\log(2X'')\big(a\theta_h^2+w''\theta_h\big)
   +2m|T|\frac{a\log d}d.
  \end{align*}
\item If $p$ is kept and $t\in T\smallsetminus T(p)$, then $p\in P$ and $t$ is
  a bad top at $p$.
\item Every prime $p$ of the first kind satisfies one of: $p\le Y'$ and the
  class of $p$ in $(\mathbf Z/2^{e+2}\mathbf Z)^\times/\{\pm1\}$ has order less
  than $x^{\varepsilon/4}$; or $p\ge p_0'$ and $X'(1-\rho)\le p^i\le
  X'(1+\rho)$ for some integer $i\ge1$.
\end{enumerate}
\end{lemma}
\leanref{Formalism/SizeWindowBrackets/ChainSlopeWindows/AdmissibleSchurLemma.lean}{Formalism.LFunctionChainParameters.DeficitSplitLemmas.loss_bound}
\leanref{Formalism/SizeWindowBrackets/SecondKindBucket.lean}{Formalism.LFunctionChainParameters.DeficitSplitLemmas.SecondKindLoss}
\leanref{Formalism/LocalToGlobal/GlobalTheory/LossRequirementsOfSmallShares.lean}{Formalism.L2.GlobalTheoryTheory.Loss7.loss_setting_family}
\leanref{Formalism/LocalToGlobal/GlobalTheory/PowFourLtCeilThresholdPrime.lean}{Formalism.L2.GlobalTheoryTheory.Loss.ASK_at}
\leanref{Formalism/LocalToGlobal/GlobalTheory/PowFourLtCeilThresholdPrime.lean}{Formalism.L2.GlobalTheoryTheory.Loss.Ga_at}
\leanref{Formalism/SharpRestriction/GlobalSchurBounds/FrameCodeOfAdditive.lean}{Formalism.MainNotions.Lnk.askIsF}
\leanref{Formalism/LocalToGlobal/GlobalTheory/PowFourLtCeilThresholdPrime.lean}{Formalism.L2.GlobalTheoryTheory.Loss.unserved_bad}
\leanref{Formalism/LocalToGlobal/GlobalTheory/PowFourLtCeilThresholdPrime.lean}{Formalism.L2.GlobalTheoryTheory.Loss.firstHyp_at}
\begin{proof}
Let $x$ be large enough for the gain link and the thresholds below, and let $D$
be a member at $x$, with the gain data of the family. For a prime $p$ write
$h_p$, $n_p$, $V(p)$ and $\Sigma(p)$ as in Definition~\ref{pd-gain-data}, and
$g_j(p)$ and $\tau(p)$ for the gain and the depth of
Definition~\ref{pd-gain-quantities}. Among the choices of the family is the
rule for the dropped primes: a prime $p$ is dropped if and only if it is no
auxiliary prime $T_u$ and one of the following holds:
\begin{enumerate}
\item[(i)] $p<X'(1+C_r\rho)$, $p$ is of inert class, and either
  $n_p<x^{\varepsilon/4}$, or $p\ge p_0$ and $|p^i-X'|\le\rho X'$ for some
  integer $i\ge1$, or $p$ is a candidate prime with $\Sigma(p)>x^{-\varepsilon}$;
\item[(ii)] $p>X''$;
\item[(iii)] $p$ is neither of inert class nor a signal prime;
\item[(iv)] $p$ is of inert class and $h_p\neq0$ has odd order.
\end{enumerate}
The auxiliary primes are not inert: $T_u=\pi\bar\pi$ with $\pi$ in
$K=\mathbf Q(\sqrt{-d})$ (Definition~\ref{th-definite-datum}(4)), so $-d$ is a
square modulo $T_u$. So an inert prime is dropped exactly when one of (i) to
(iv) holds. Recall that an inert prime is a signal prime or of inert class
(Definition~\ref{pd-size-data}), and that $\rho\ge0$ and $C_r\ge1$.

(2) Let $p$ be kept and $t\in T\smallsetminus T(p)$. As $T(p)=T$ when $p$ is a
signal prime or $p\ge X'$ (Definition~\ref{pd-gain-quantities}(2)), $p$ is of
inert class and $p<X'\le X'(1+C_r\rho)$. As $p$ is not dropped, (i) fails:
$n_p\ge x^{\varepsilon/4}$; $p$ is not in band, since a prime in band has
$p\ge p_0$ and $|p^i-X'|\le\rho X'$; so $p$ is a candidate prime
(Definition~\ref{pd-gain-data}(c)), and then $\Sigma(p)\le x^{-\varepsilon}$. So
$V(p)$ is the set of good representatives at $p$, and $t\notin T(p)$ says that
the representative $t^{s_p}$ of $t$ is not good: $t$ is a bad top at $p$. As
$p<X'$ is of inert class with $n_p\ge x^{\varepsilon/4}$, $p\in P$.

(3) Let $p$ be of the first kind: $p\in\mathrm{Pr}$ is dropped, $p\le X''$, and
$p$ is not a candidate prime with $\Sigma(p)>x^{-\varepsilon}$. Then (ii) and
(iii) fail, so (i) or (iv) holds. The homomorphism $h$ of the Galois datum maps
$h_p$ to the class of $p$ in $H_e=(\mathbf Z/2^{e+2}\mathbf Z)^\times/\{\pm1\}$
(Definitions~\ref{hf-frame} and~\ref{hf-framed-host-field}(4)), so the order of
that class divides $n_p$. In case (i) with $n_p<x^{\varepsilon/4}$, $p<Y\le Y'$
and the class of $p$ has order less than $x^{\varepsilon/4}$. In case (i) with
$p\ge p_0$ and $|p^i-X'|\le\rho X'$, $p\ge\lceil p_0\rceil=p'_0$, as $p$ is an
integer, and $X'(1-\rho)\le p^i\le X'(1+\rho)$. The third alternative of (i) is
excluded. In case (iv), $h_p$ has odd order and $H_e$ has order $2^e$, so the
class of $p$ is trivial, of order $1<x^{\varepsilon/4}$, and $p\le X''\le Y'$.

(1) Let $N$ be as in (1), and let $m$ be the number of lifts. For $p\in\mathrm{Pr}$,
$t\in T$ and a lift $j$ let $A_{j,t}(p)=\sum_i\max(a_t/p^i-w_{j,t,p,i},0)$, the
sum over the odd $1\le i\le\log_pa_t+1$, with the weights of the size data; and
for $t\in T$ and $j$ let $A^d_{j,t}=\log d\sum_i\max(a_t/d^i-w^d_{j,t,i},0)$,
the sum over $1\le i\le\log_da_t+1$, with the weights $w^d_{j,t,i}\ge0$ of the
lift $j$ at the places of the term at $d$. Put $g^\flat_j(p)=0$ if $p$ is
dropped and $g^\flat_j(p)=g_j(p)$ otherwise.
\openstep{The identification of the left side of (1) and the side conditions
are not yet written out: that
$\sum_{t}\sum_j\mathrm{ASK}_N(a_t,t,w_j)-G_a$ equals
$\sum_{p\in\mathrm{Pr}}2\log p\big(\sum_j\sum_{t\in T}A_{j,t}(p)-\sum_jg^\flat_j(p)+\tau(p)\big)
+\sum_j\sum_tA^d_{j,t}$; that $A_{j,t}(p)$ is the class sum $\mathrm{VM}_{j,t}(p)$
of the gain link; that $0\le a_t\le a(1+\theta_h)$, that the weights
$w_{j,t,p,i}$ are at least $w''$, that $d\ge4$, $p'_0\ge3$ and $X'\ge1$; and that the lower and floor
terms of the gain link at kept primes are at most
$C_PS_k\log_p(2a)+[p<X'](\mathrm{pr}(p)+\mathrm{or}(p))$ and $C_PfN(m)^2\log_pa+C_Pf(\log_pa)^2$,
with $\mathrm{pr}$ equal to $Q_1$ for $p<p'_0$ and $Q_2$ otherwise, and
$\mathrm{or}(p)=C_P\log x(\Lambda_Ga/(pk)+o(p))$ as in
Lemma~\ref{asm-global-link-losses-terms}(5) with $L_1(p)$ such that $K_0\le
p^{L_1(p)-1}$.}
% Lean for this open step: Lnk.askIsF (askRec_at, inert_sum), Loss.ASK_at (pb_primes, pb_atD, pb_sum_TTv),
% Loss.Ga_at (sum_places), unitRule_at and ask_at (A_{j,t}(p) = vonMangoldtClassSum), and side_at
% (pd_at, pd_wz, pd_wd, pd_fl, pd_lo, loss_constant_le_prime_pow) with side_condition_of_windows. Size:
% about 144 theorems and 1,480 lines for the identification, 46 theorems and 450 lines for the side
% conditions.

\emph{The deficit at a prime.} Put
$\mathrm{Def}(p)=\sum_j\sum_{t\in T}A_{j,t}(p)-\sum_jg^\flat_j(p)+\tau(p)$. Let
$p$ be kept. Then $p\le X''$, by (ii), and the gain link (3)(b), at a place $w$
over $p$, at which the size data are filled by the gain data
(Definition~\ref{pd-gain-quantities}(6): $T_w=T(p)\subset T$,
$\mathrm{gain}_j(w)=g_j(p)$, $\tau_w=\tau(p)$), gives
$\sum_j\sum_{t\in T(p)}(A_{j,t}(p)-\mathrm{lo}(p))-m\,\mathrm{fl}(p)\le\sum_jg_j(p)-\tau(p)$.
Hence
\[
 \mathrm{Def}(p)\le\sum_j\sum_{t\in T\smallsetminus T(p)}A_{j,t}(p)+m|T(p)|\,\mathrm{lo}(p)
 +m\,\mathrm{fl}(p).
\]
Let $p$ be dropped. Then $g^\flat_j(p)=0$, and $V(p)=\emptyset$. Indeed, a
prime with $V(p)\neq\emptyset$ is a candidate prime with $\Sigma(p)\le
x^{-\varepsilon}$: it lies below $X'\le X''$, is of inert class with
$n_p\ge x^{\varepsilon/4}$, and is not in band, so (i), (ii) and (iii) fail;
and (iv) fails as $n_p$ is even. For an element of $A$ of odd order has trivial
image in $H_e$, which has order $2^e$; as $h$ is injective on $A_0$, its order
divides the index $o$, and $o\le x^{\varepsilon/100}<x^{\varepsilon/4}$ by the
size data. So $\tau(p)=0$ and $\mathrm{Def}(p)=\sum_j\sum_{t\in T}A_{j,t}(p)$.

\emph{The terms.} As $\theta_h\le\frac12$, $0\le a_t\le\frac32a$, and $a\ge1$.
By Lemma~\ref{asm-global-link-losses-terms}(1), $2\log p\cdot A_{j,t}(p)\le4a\log
p/p$ for all $p,j,t$. So the kept primes contribute at most
$\sum_{p\text{ kept}}4ma\frac{\log p}p|T\smallsetminus T(p)|$ through their first
sums, and the primes of the first and second kinds contribute at most
$4m|T|a\sum_p\frac{\log p}p$ over each kind. By
Lemma~\ref{asm-global-link-losses-terms}(3), with $\theta=\theta_h$ and
$X''w''=a$, the high primes contribute at most
$m|T|\cdot2\log(2X'')(a\theta_h^2+w''\theta_h)$, and by (2) of that lemma the
terms $A^d_{j,t}$ add up to at most $m|T|\cdot2a\log d/d$. The kept primes are
at most $X''$ in number and lie in $[2,X'']$; with $|T(p)|\le|T|$ and the bounds
on the lower and floor terms, Lemma~\ref{asm-global-link-losses-terms}(4) and~(5),
the latter for the kept primes below $X'$, with $\mu=m|T|$, $l_x=\log x$,
$\Lambda=\Lambda_G$ and $p_0=p'_0$, bound their remaining
contributions by the first four terms of (1). Adding up gives (1).
% Lean. (1): loss_bound (SecondKindLoss), which is split_classes_count (splitEq_valid, deficit_le,
% deficit_dropped: the deficit at a prime) with the bounds of Lemma asm-global-link-losses-terms,
% applied to the record ldAt X K C bo (W = PsAt L, the inert
% primes below bndP, with l = log, cl = clAt, TTp = image of G.TTP, Iex the odd exponents below
% log_p a_t + 2, wz the fills at the partner positions, Id and wd the d-term, gain = gainQ, tau =
% tauP, lo = loG L C bo, fl = flG L C, L1 = leastPowerExponent) at P = par Fm x i C bo 0 Cord; its
% hypotheses are ParametersHyp = lossHyp_at (UnitRule, a_t <= a(1 + theta_h), the weights at least
% w'', served p <= X'', high p > X'', SizeData2 at served primes = loss_bound_at from the gain link
% (3b) LinkG_at, zero gains off the served primes, fl <= flF and lo <= loF, K0 <= p^(L1 p - 1)) and
% the 18 side conditions side_condition_of_windows. The left side is rewritten by Lnk.askIsF (the
% ASK sums of the host's arithmetic data are LocalQuadraticParts.ASK L, through askRec_at and
% inert_sum), ASK_at (= LD.ASK, through pb_primes and pb_atD) and Ga_at (sum_places, grouping the
% places of the support over each prime). loss_bound itself is split_classes_count with
% top_loss_bound, high_bucket_bound, d_term_bound and the floor, cp, pr and or buckets, about 1,400
% lines in AdmissibleSchurLemma.lean and SecondKindBucket.lean; loss_setting_family with lossHyp_at
% about 1,400 lines; Lnk.askIsF about 1,200 lines; ASK_at and Ga_at about 650 lines (overlapping).
% (2) is unserved_bad (pc_drop_iff from PtH.drop_iff, pc_cand, the NoGain cases); (3) is
% firstHyp_at (NoGain, pc_ord through J.two_readings_of_Frobenius, and
% even_addOrderOf_hP_of_pow_le for the odd-order case).
\end{proof}

\begin{lemma}\label{asm-global-link-losses-orders}
Let $0\le j\le e$ be integers and $l$ an odd integer, and let $h$ be the class
of $l$ in $(\mathbf Z/2^{e+2}\mathbf Z)^\times/\{\pm1\}$. Then $h^{2^j}=1$ if
and only if $l\equiv\pm1\bmod2^{e+2-j}$.
\end{lemma}
\leanref{Formalism/RankGapBounds/RankGapProofs/FourNotDvdOddSqAddOne.lean}{Formalism.RankGapBounds.LocalQuadraticParts.powTwoPowEqOneIff_holds}
\leanref{Formalism/RankGapBounds/RandomSubsetCounterexample.lean}{Formalism.RankGapBounds.LocalQuadraticParts.OrdersIff}

\begin{lemma}\label{asm-global-link-losses-residues}
Let $M\ge2$ be an integer, $Y\ge1$ real, and $S$ a finite set of integers $l$
with $3\le l\le Y$ and $l\equiv\pm1\bmod M$. Then
\[
 \sum_{l\in S}\frac{\log l}l\le\frac{4\log Y(1+\log Y)}M.
\]
\end{lemma}
\leanref{Formalism/RankGapBounds/RankGapProofs/PmModulusFacts.lean}{Formalism.RankGapBounds.LocalQuadraticParts.pm_y_sum}

\begin{lemma}\label{asm-global-link-losses-band}
Let $X'\ge2$ and $0\le\rho\le1$ be real.
\begin{enumerate}
\item If $p\ge8$ and $i\ge0$ are integers with $p^i\le X'(1+\rho)$, then
  $i\le\log X'$.
\item If $p\ge0$ is real and $i\ge1$ an integer with $X'(1-\rho)\le p^i\le
  X'(1+\rho)$, then $X'^{1/i}(1-\rho)\le p\le X'^{1/i}(1+\rho)$.
\item If $P\ge1$, $\rho\le1/2$ and $p_0\ge1$, and $S$ is a finite set of
  integers $n\ge p_0$ with $P(1-\rho)\le n\le P(1+\rho)$, then
  \[
   \sum_{n\in S}\frac{\log n}n\le\Big(4\rho+\frac1{p_0}\Big)\log(2P).
  \]
\end{enumerate}
\end{lemma}
\leanref{Formalism/SizeWindowBrackets/ChainSlopeWindows/AdmissibleSchurLemma.lean}{Formalism.LFunctionChainParameters.DeficitSplitLemmas.band_exponent_bound}
\leanref{Formalism/SizeWindowBrackets/ChainSlopeWindows/AdmissibleSchurLemma.lean}{Formalism.LFunctionChainParameters.DeficitSplitLemmas.band_root_bound}
\leanref{Formalism/SizeWindowBrackets/ChainSlopeWindows/AdmissibleSchurLemma.lean}{Formalism.LFunctionChainParameters.DeficitSplitLemmas.band_sum}

\begin{lemma}\label{asm-global-link-losses-first}
Let $e,J\ge1$ be integers and $Y\ge1$, $0\le\rho\le1/2$, $p_0\ge8$ and
$X'\ge2$ real numbers. For an odd prime $p$ let $h(p)$ be the class of $p$ in
$(\mathbf Z/2^{e+2}\mathbf Z)^\times/\{\pm1\}$. Let $S$ be a finite set of odd
primes, each of which satisfies one of:
\begin{enumerate}
\item[(a)] $p\le Y$ and $h(p)$ has order at most $J$;
\item[(b)] $p\ge p_0$, and $X'(1-\rho)\le p^i\le X'(1+\rho)$ for some integer
  $i\ge1$.
\end{enumerate}
Then
\[
 \sum_{p\in S}\frac{\log p}p\le\frac{12J\log(2^eY)^2}{2^e}
 +\lfloor\log X'\rfloor\Big(4\rho+\frac1{p_0}\Big)\log(2X').
\]
\end{lemma}
\leanref{Formalism/SizeWindowBrackets/ChainSlopeWindows/OneLeThresholdMulPowers.lean}{Formalism.LFunctionChainParameters.PrimeWindowSizeBounds.first_kind_bound}
\leanref{Formalism/SizeWindowBrackets/ChainSlopeWindows/AdmissibleSchurLemma.lean}{Formalism.LFunctionChainParameters.DeficitSplitLemmas.first_kind_bucket_bound}
\leanref{Formalism/RankGapBounds/RankGapProofs/PmModulusFacts.lean}{Formalism.RankGapBounds.LocalQuadraticParts.ordersSumY_of}
\begin{proof}
Let $S_a$ be the set of $p\in S$ that satisfy (a), and $S_b=S\smallsetminus S_a$;
the primes of $S_b$ satisfy (b).

\emph{The primes of (a).} Put $j=\min(e,\lfloor\log_2J\rfloor)$ and
$M=2^{e+2-j}$. Then $M\ge2$, and $2^e=2^j2^{e-j}\le JM$ as $2^j\le J$. Let
$p\in S_a$. As $p$ is odd, $p\equiv\pm1\bmod4$, so $h(p)^{2^e}=1$ by
Lemma~\ref{asm-global-link-losses-orders} with $j=e$. So the order of $h(p)$ is
$2^i$ with $i\le e$, and $2^i\le J$ gives $i\le\lfloor\log_2J\rfloor$. Hence
$i\le j$, $h(p)^{2^j}=1$, and $p\equiv\pm1\bmod M$ by
Lemma~\ref{asm-global-link-losses-orders}. Every $p\in S_a$ is thus an integer
with $3\le p\le Y$ and $p\equiv\pm1\bmod M$, and
Lemma~\ref{asm-global-link-losses-residues} gives
\[
 \sum_{p\in S_a}\frac{\log p}p\le\frac{4\log Y(1+\log Y)}M.
\]
Put $A=e\log2\ge\log2\ge1/2$ and $B=\log Y\ge0$. Then $\log(2^eY)=A+B$ and
$B(1+B)=B+B^2\le6AB+B^2\le3(A+B)^2$, since $6A\ge1$. So, with
$2^e\le JM$,
\[
 \sum_{p\in S_a}\frac{\log p}p\le\frac{12\log(2^eY)^2}M
 \le\frac{12J\log(2^eY)^2}{2^e}.
\]

\emph{The primes of (b).} Let $p\in S_b$ and $i\ge1$ with
$X'(1-\rho)\le p^i\le X'(1+\rho)$. As $X'\ge2$, $\rho\le1$ and $p\ge8$,
$i\le\log X'$, so $i\le\lfloor\log X'\rfloor$; and, as $0\le\rho\le1$,
$X'^{1/i}(1-\rho)\le p\le X'^{1/i}(1+\rho)$
(Lemma~\ref{asm-global-link-losses-band}(1) and~(2)). So $p$ lies in the set
$B_i$ of primes $p'\ge p_0$ with $P_i(1-\rho)\le p'\le P_i(1+\rho)$, where
$P_i=X'^{1/i}$ and $1\le P_i\le X'$. Every term of the sum over $S_b$ is
nonnegative and occurs in the sum over some $B_i$, $1\le i\le\lfloor\log
X'\rfloor$, so by Lemma~\ref{asm-global-link-losses-band}(3)
\[
 \sum_{p\in S_b}\frac{\log p}p\le\sum_{i=1}^{\lfloor\log X'\rfloor}
 \sum_{p\in B_i}\frac{\log p}p\le\lfloor\log X'\rfloor
 \Big(4\rho+\frac1{p_0}\Big)\log(2X').
\]
Adding the two bounds gives the lemma.
\end{proof}
% Lean: first_kind_bound = first_kind_bucket_bound with I = floor(log X') and P i = X'^(1/(i+1)),
% the band of p found by band_exponent_bound (i <= log X') and band_root_bound; the constant 12 is
% that of ordersSumY_of (OrdersSumY), whose proof is pm_resid (with OrdersIff =
% powTwoPowEqOneIff_holds), pm_modulus_facts (M >= 2, 2^e <= J M), pm_y_sum and the inequality
% log Y (1 + log Y) <= 3 log(2^e Y)^2 from e log 2 >= 1/2. The band part is band_total with band_sum.

\begin{lemma}\label{asm-global-link-losses-sizes}
Let $\mathcal F$ be a family with exponent $\varepsilon$
(Definition~\ref{pd-family}), let $C_G$ be real and $\delta$ a function with
$\delta(x)\to0$ as $x\to\infty$, and put $K=13+|C_G|$. Then for all large $x$
every member $D$ at $x$, with its loss numbers for $(C_G,\delta)$ and the
notation of Definition~\ref{asm-global-link-losses-numbers}, has the following
properties, where $\mathrm{sg}=100Kx^{4\varepsilon}/x$.
\begin{enumerate}
\item $x^2\le2^e\le n$, $x^4+1\le p_0'$, $0\le\rho\le x^{-2}$,
  $0\le C_P\le K$, $1\le Y\le2X'$, $x\ge100$ and $\mathrm{sg}\le1$.
\item If $T$ is nonempty, then
  \begin{align*}
   &2mC_PfX''\big(N(m)^2\log a+2(\log a)^2\big)+2C_PS_k\log(2a)\,m|T|X''
   +2\log4\cdot m|T|\big(p_0'Q_1+X'Q_2\big)\\
   &\quad+2C_Pm|T|a\log x\Big(\Big(\frac{\Lambda_G}k+\frac1{K_0}\Big)
   (\log X'+\log4)+\frac{2(\log(p_0'-1)+1)}{p_0'-1}\Big)
   +4am|T|x^{-2+\varepsilon}\\
   &\quad+4m|T|a\Big(\frac{12J_o\log(2^eY')^2}{2^e}
   +I_b\Big(4\rho+\frac1{p_0'}\Big)\log(2X')\Big)\\
   &\quad+2m|T|\log(2X'')\big(a\theta_h^2+w''\theta_h\big)
   +2m|T|\frac{a\log d}d+2am|T|x^{-9}
   \le\big(2\log4\cdot C_P+13\,\mathrm{sg}\big)|T|mb.
  \end{align*}
\end{enumerate}
\end{lemma}
\leanref{Formalism/LocalToGlobal/GlobalTheory/PowFourLtCeilThresholdPrime.lean}{Formalism.L2.GlobalTheoryTheory.Loss.windows_ev}
\leanref{Formalism/LocalToGlobal/GlobalTheory/PowFourLtCeilThresholdPrime.lean}{Formalism.L2.GlobalTheoryTheory.Loss.pe_windows_at}
\leanref{Formalism/SizeWindowBrackets/WindowSettingSatisfiable.lean}{Formalism.LFunctionChainParameters.PrimeWindowSizeBounds.Windows}
\leanref{Formalism/SizeWindowBrackets/WindowSettingSatisfiable.lean}{Formalism.LFunctionChainParameters.PrimeWindowSizeBounds.LargeX}
\leanref{Formalism/SizeWindowBrackets/ChainSlopeWindows/OneLeThresholdMulPowers.lean}{Formalism.LFunctionChainParameters.PrimeWindowSizeBounds.large_x_eventually}
\leanref{Formalism/SizeWindowBrackets/ChainSlopeWindows/DiscriminantSettingSatisfied.lean}{Formalism.LFunctionChainParameters.SmallOrderPrimeLosses.size_core}
\leanref{Formalism/SizeWindowBrackets/ChainSlopeWindows/OneLeThresholdMulPowers.lean}{Formalism.LFunctionChainParameters.PrimeWindowSizeBounds.size_floors_bound}
\leanref{Formalism/SizeWindowBrackets/ChainSlopeWindows/OneLeThresholdMulPowers.lean}{Formalism.LFunctionChainParameters.PrimeWindowSizeBounds.size_cp_bucket_bound}
\leanref{Formalism/SizeWindowBrackets/ChainSlopeWindows/OneLeThresholdMulPowers.lean}{Formalism.LFunctionChainParameters.PrimeWindowSizeBounds.size_prime_low_bound}
\leanref{Formalism/SizeWindowBrackets/ChainSlopeWindows/OneLeThresholdMulPowers.lean}{Formalism.LFunctionChainParameters.PrimeWindowSizeBounds.size_prime_high_bound}
\leanref{Formalism/SizeWindowBrackets/ChainSlopeWindows/OneLeThresholdMulPowers.lean}{Formalism.LFunctionChainParameters.PrimeWindowSizeBounds.size_or_bucket_bound}
\leanref{Formalism/SizeWindowBrackets/ChainSlopeWindows/DiscriminantSettingSatisfied.lean}{Formalism.LFunctionChainParameters.SmallOrderPrimeLosses.sizeOut_core}
\leanref{Formalism/SizeWindowBrackets/ChainSlopeWindows/OneLeThresholdMulPowers.lean}{Formalism.LFunctionChainParameters.PrimeWindowSizeBounds.size_orders_bucket_bound}
\leanref{Formalism/SizeWindowBrackets/ChainSlopeWindows/OneLeThresholdMulPowers.lean}{Formalism.LFunctionChainParameters.PrimeWindowSizeBounds.size_band_bound}
\leanref{Formalism/SizeWindowBrackets/ChainSlopeWindows/OneLeThresholdMulPowers.lean}{Formalism.LFunctionChainParameters.PrimeWindowSizeBounds.size_high_bucket_bound}
\leanref{Formalism/SizeWindowBrackets/ChainSlopeWindows/OneLeThresholdMulPowers.lean}{Formalism.LFunctionChainParameters.PrimeWindowSizeBounds.size_discriminant_bound}
\leanref{Formalism/SizeWindowBrackets/ChainSlopeWindows/DiscriminantSettingSatisfied.lean}{Formalism.LFunctionChainParameters.SmallOrderPrimeLosses.sizeSecond_core}
\begin{proof}
\emph{The windows.} Let $C_t$ be the constant of the size data with $t\le
x^{C_t}$ and put $e_W=\max(40,C_t+3)$. Take $x$ beyond $x_1+1$ and $3$, with
$C_r\le x$, $r(\varepsilon)\le x$, $x^7\le\exp(x^{\varepsilon/4}/2)$,
$x^2\le\exp(x^{\varepsilon/(5J)-\varepsilon/(100J)})$ and $\delta(x)<1/2$; then
$b_o=b^{\delta(x)}\le\sqrt b$. With the size bounds of the family and the bound
$\rho\le x^{-2}$ among its choices, every member at $x$ satisfies the
\emph{windows}:
$x^2\le2^e\le n\le x^{e_W}$, $x^2\le k\le x^{e_W}$, $N(m)\le x^{e_W}$,
$t\le x^{e_W}$, $x^4+1\le p_0'\le x^{e_W}$, $0\le\Lambda_G\le x^{2\varepsilon}$,
$r(\varepsilon)\le x^{e_W}$, $0\le S_c,\eta\le\exp(x^{\varepsilon/4})$,
$0\le\rho\le x^{-2}$, $0\le\theta_h\le1/x$, $x^2\le K_0$, $x^2\le d$,
$0\le C_P\le K$, $12\le K$, $1\le Y\le2X'$ and $0\le b_o\le\sqrt b$. Each is a
size bound of the family or a computation from them; for instance
$Y=X'(1+C_r\rho)\le X'(1+x\cdot x^{-2})\le2X'$.

\emph{The large scale.} As $0<\varepsilon<1/4$, for large $x$ also $x\ge1$,
$100Kx^{4\varepsilon}\le x$, $100Kx^{4e_W+4}\le X'$ and
$100Kx^{4e_W+4}X'\exp(x^{\varepsilon/4})\le\sqrt b$. As $K\ge1$ and
$x^{4\varepsilon}\ge1$, the second gives $x\ge100$ and $\mathrm{sg}\le1$, which
completes (1).

\emph{The terms.} Put $U=|T|mb$, and suppose $|T|\ge1$. The following are
computations from the windows, the large scale, $a=xb$ and $b=(\sqrt b)^2$,
term by term: the term in
$N(m)$ is at most $\mathrm{sg}U$; the term in $S_k$ is at most $\mathrm{sg}U$;
$2\log4\cdot m|T|p_0'Q_1\le\mathrm{sg}U$;
$2\log4\cdot m|T|X'Q_2\le(2\log4\cdot C_P+3\,\mathrm{sg})U$; the term in
$\log x$ is at most $\mathrm{sg}U$; $4am|T|x^{-2+\varepsilon}\le\mathrm{sg}U$;
the term in $J_o$ and the term in $I_b$ are each at most $\mathrm{sg}U$; the term
in $\theta_h$, the term in $d$ and $2am|T|x^{-9}$ are each at most
$\mathrm{sg}U$. Adding these eleven bounds gives (2).
\end{proof}
% Lean: (1) and the windows are windows_ev Fm C 0 Cord hCord delta hdelta with Cord = 12 (the constant
% of joint_core, max C1 0 with C1 = 12 from first_kind_bound and ordersSumY_of), so K = 1 + |C| +
% |0| + 12, and e = max 40 (Fm.Ct + 3); its windows are Windows (par Fm x i C (b^delta x) 0 Cord) e K
% from pe_windows_at (the size bounds SizesAt = Fm.sizes, Fm.rho_le, pe_Y), here with beta = s0 = 0
% (clr P, discriminant_windows_iff). The large scale is LargeX x eps e K (large_x_eventually); x >= 100
% and sg <= 1 are Facts.hx and Facts.sg1 of facts. (2) is size_core: tDetD = tFloors + tCp + tPrLow +
% tPrHigh + tOr + tOutD + tOrders + tBand + tHigh + tD + tSecondD, bounded by size_floors_bound,
% size_cp_bucket_bound, size_prime_low_bound, size_prime_high_bound (with 2 log 4 C + 3 sg),
% size_or_bucket_bound, sizeOut_core, size_orders_bucket_bound, size_band_bound,
% size_high_bucket_bound, size_discriminant_bound and sizeSecond_core, each a computation from Facts.
% Sizes: windows_ev about 430 lines; the eleven size bounds about 500 lines in
% OneLeThresholdMulPowers.lean.

\begin{lemma}\label{asm-global-link-losses}
Let $\mathcal F$ be a family with exponent $\varepsilon$
(Definition~\ref{pd-family}) that satisfies the gain link
(Definition~\ref{pd-gain-link}). Suppose that for all large $x$ every member at
$x$ satisfies $\delta_1n/2\le|T|\le2\delta_1n$ and, with $P$ and $\beta(p)$ as
in Lemma~\ref{asm-global-link-bad-primes},
$\sum_{p\in P}\frac{\log p}p\,\beta(p)\le|T|\,x^{-2+\varepsilon}$. Then there
are $C\ge0$ and $x_0$ such that every member at every scale $x\ge x_0$ whose
weight vectors have small shares (Definition~\ref{asm-duty}(5)) satisfies the
loss bounds of Definition~\ref{asm-member-bounds}(6) with $(C,0)$.
\end{lemma}
\leanref{Formalism/LocalToGlobal/GlobalTheory/LossRequirementsOfSmallShares.lean}{Formalism.L2.GlobalTheoryTheory.Loss7.losses_of_small_share_weight}
\leanref{Formalism/LocalToGlobal/GlobalTheory/LossRequirementsOfSmallShares.lean}{Formalism.L2.GlobalTheoryTheory.Loss7.stLoss7}
\begin{proof}
\emph{Constants.} Let $C_G$, $x_G$ and $\delta$ be a constant, a threshold and
a function of the gain link, with $\delta(x)\to0$. Put $C_P=\max(C_G,0)$ and
$K=13+|C_G|$, so that $0\le C_P\le K$ and $K\ge1$. We show the lemma with
$C=3K+13$.

Let $x$ be large enough that $x>x_1$, $x\ge3$, $x\ge x_G$, the hypotheses hold
at $x$, and Lemmas~\ref{asm-global-link-losses-explicit}
and~\ref{asm-global-link-losses-sizes} (for $C_G$, $\delta$ and $K$) apply at
$x$. Put $\mathrm{sg}=100Kx^{4\varepsilon}/x$; by
Lemma~\ref{asm-global-link-losses-sizes}(1), $x\ge100$ and $\mathrm{sg}\le1$.
Let $D$ be a member at $x$ whose weight vectors have small shares. As
$|T|\ge\delta_1n/2>0$, $T$ is nonempty. Let $\mathrm{Cand}$ be the set of candidate primes of $D$
(Definition~\ref{pd-gain-data}(c)), all of them below $X'$, and let $\Sigma(p)$,
$V(p)$, $S(p,t_0)$ and $\iota(p)$ be as in Definitions~\ref{pd-gain-data}
and~\ref{pd-gain-quantities}. The small shares say
\[
 \sum_{p\in\mathrm{Cand}}\frac{\log p}p\Big(2\cdot\mathbf 1\big[\Sigma(p)>x^{-\varepsilon}\big]
 +4\Sigma(p)\Big)\le x^{-9}.
\]

\emph{The bound on $I$.} Among the choices of the family, the size data are
filled by the gain data, so $\iota_w=\iota(p_w)$, and every place $w$ with
$p_w\le X''$ is in the support. If $p$ is not a candidate prime, then
$V(p)=\emptyset$ and $\iota(p)=0$; every candidate prime is below $X'\le X''$,
so all places over it are in the support. For a prime $p$ the numbers $\ell_we_w$
of the places $w$ over $p$ add up to $\log p$. Grouping the places of the
support by $p_w$ therefore gives
\[
 I=\sum_{w\in\operatorname{supp}}\ell_we_w\iota(p_w)
 =\sum_{p\in\mathrm{Cand}}\log p\cdot\iota(p).
\]
Each $\kappa_{p,t_0}$ is nonnegative: for $p<p_0$ it is the minimum of $1$ and
a product of nonnegative numbers, as $K_G\ge0$; for $p\ge p_0$ the local
constants of the family are least, so $\kappa_I$ is the infimum of a set of
nonnegative numbers, as in the first paragraph of the proof of
Lemma~\ref{pd-gain-first-hypotheses}. So $\Sigma(p)\ge0$. For $t_0\in V(p)$ the gain link (2) gives $S(p,t_0)\le2m\,
N(p,t_0)a/p\le2m|T|a/p$, as $N(p,t_0)\le|T|$; so
$\iota(p)\le2\Sigma(p)\cdot2m|T|a/p$, also when $V(p)=\emptyset$. Hence, as the
indicator term above is nonnegative,
\[
 I\le m|T|a\sum_{p\in\mathrm{Cand}}\frac{\log p}p\,4\Sigma(p)\le m|T|a\,x^{-9}.
\]
Now $a=xb$, so
\[
 2a\,x^{-9}=b\cdot\frac{2x^{-7}}{100}\cdot\frac{100}x\le\mathrm{sg}\cdot b,
\]
since $2x^{-7}/100\le1$ and $100/x\le100Kx^{4\varepsilon}/x=\mathrm{sg}$, as
$K\ge1$ and $x^{4\varepsilon}\ge1$. So $a\,x^{-9}\le\mathrm{sg}\cdot b$, and, as
$\mathrm{sg}\le1$,
\[
 I\le m|T|a\,x^{-9}\le\mathrm{sg}\cdot|T|mb\le|T|mb.
\]
% Lean for the bound on I: the second component hc2 of joint_discriminant_eventually, through
% joint_core: io_discriminant_bound (with lhsW cand theta Sig <= x^(-9) from ChW_at, theta =
% x^(-eps)), sizeIo_core (a x^(-9) <= sg P K * b from size_ninth_power_bound, 2 a x^(-9) <= sg b by
% F.ha, F.xr and F.close_sg with F.zr_sg at i = 0) and F.sg1 (sg <= 1). The hypothesis io p <= 2 Sig
% p (2 m |T| a / p) is Lnk.ioP_le with the gain link hG.2.1 and Lnk.Sig_nonneg, in loss_setting_family;
% I = sum over sAt (= cands) of log p io(p) is IO_at (fills, supp_sat, sum_places,
% ioP_eq_zero_of_not_cand₂).

\emph{The inert-prime sums.} Let $N$ be an integer as in
Definition~\ref{asm-member-bounds}(6), and use the loss numbers of $D$ for
$(C_G,\delta)$ and its classes of primes
(Definition~\ref{asm-global-link-losses-numbers}). By
Lemma~\ref{asm-global-link-losses-explicit}(1), the sum
$\sum_{t\in T}\sum_{j=1}^m\mathrm{ASK}_N(a_t,t,w_j)-G_a$ is at most the sum of
the nine terms there. Three of them are sums over primes, which we bound in
turn.

The kept primes. By Lemma~\ref{asm-global-link-losses-explicit}(2), for a kept
prime $p$ every $t\in T\smallsetminus T(p)$ is a bad top at $p$, and $p\in P$
if $T\smallsetminus T(p)$ is nonempty. So $|T\smallsetminus T(p)|\le\beta(p)$
for $p\in P$, and $T\smallsetminus T(p)$ is empty for $p\notin P$. As the terms
$4ma\log p/p\cdot\beta(p)$ are nonnegative and the kept primes are distinct,
the hypothesis on the bad tops gives
\[
 \sum_{p\text{ kept}}4ma\frac{\log p}p|T\smallsetminus T(p)|
 \le4ma\sum_{p\in P}\frac{\log p}p\,\beta(p)\le4am|T|x^{-2+\varepsilon}.
\]

The primes of the first kind. An inert prime is odd, since $-d$ is a square
modulo $2$. By Lemma~\ref{asm-global-link-losses-explicit}(3) every prime of
the first kind satisfies (a) of Lemma~\ref{asm-global-link-losses-first} with
$J=J_o$ and $Y'$ in place of $Y$, as an order less than $x^{\varepsilon/4}$ is
at most $J_o$, or (b) with $p_0'$, $\rho$ and $X'$. The hypotheses of that
lemma hold: $e\ge1$ as $x^2\le2^e$, $J_o\ge1$ as $x\ge1$, $Y'\ge Y\ge1$,
$0\le\rho\le x^{-2}\le1/2$ and $p_0'\ge x^4+1\ge8$
(Lemma~\ref{asm-global-link-losses-sizes}(1)), and $X'=\exp(x^{\varepsilon/3})\ge
e\ge2$. As $\lfloor\log
X'\rfloor=\lfloor x^{\varepsilon/3}\rfloor=I_b$, it gives
\[
 4m|T|a\sum_{p\text{ first kind}}\frac{\log p}p
 \le4m|T|a\Big(\frac{12J_o\log(2^eY')^2}{2^e}
 +I_b\Big(4\rho+\frac1{p_0'}\Big)\log(2X')\Big).
\]

The primes of the second kind. These are candidate primes with
$\Sigma(p)>x^{-\varepsilon}$. For such $p$,
$4m|T|a\log p/p\le2am|T|\frac{\log p}p\big(2+4\Sigma(p)\big)$, as
$\Sigma(p)\ge0$; for the other candidate primes
$0\le2am|T|\frac{\log p}p\cdot4\Sigma(p)$. Summing over the candidate primes
and using the small shares,
\[
 4m|T|a\sum_{p\text{ second kind}}\frac{\log p}p\le2am|T|\sum_{p\in\mathrm{Cand}}
 \frac{\log p}p\Big(2\cdot\mathbf 1\big[\Sigma(p)>x^{-\varepsilon}\big]
 +4\Sigma(p)\Big)\le2am|T|x^{-9}.
\]

With these three bounds the nine terms add up to at most the left side of
Lemma~\ref{asm-global-link-losses-sizes}(2), so
\[
 \sum_{t\in T}\sum_{j=1}^m\mathrm{ASK}_N(a_t,t,w_j)-G_a
 \le\big(2\log4\cdot C_P+13\,\mathrm{sg}\big)|T|mb\le(3C_P+13)|T|mb,
\]
since $2\log4\le3$, $C_P\ge0$ and $\mathrm{sg}\le1$.
% Lean: the first component hc1 of joint_discriminant_eventually, through joint_core: hA from
% loss_bound (Lemma asm-global-link-losses-explicit(1)) with the 18 side conditions
% side_condition_of_windows; out_discriminant_bound with unserved_bad and ChBadp_at (bad_card_at
% identifies beta(p)); first_le with first_kind_bound (C1 = 12 <= Cord) and firstHyp_at, where
% floor(log X') = Ib is hIb; second_discriminant_bound (second_term_le with second_spec, and lhsW <=
% x^(-9) from ChW_at); size_core (Lemma asm-global-link-losses-sizes(2)); and 2 log 4 <= 3
% (log_four_le_three_halves) with F.sg1. Two is not inert: two_not_inert.

\emph{Conclusion.} Let $x_0$ be beyond the thresholds above. As
$C=3K+13\ge3C_P+13\ge1$ and $|T|mb\ge0$, every member at $x\ge x_0$ whose
weight vectors have small shares satisfies $I\le|T|mb\le C|T|mb\,x^{0}$ and
$\sum_{t\in T}\sum_{j=1}^m\mathrm{ASK}_N(a_t,t,w_j)-G_a\le
C|T|mb\,x^{0}$ for every $N$ as in Definition~\ref{asm-member-bounds}(6); these
are the loss bounds with $(C,0)$.
% Lean: losses_of_small_share_weight hTT hBad hG, with constant 3|K| + 13 for the K = 1 + |C| + |0| +
% Cord of windows_ev, Cord = 12, so K = 13 + |C_G| >= 0 and 3|K| + 13 = 3K + 13; P.C = max C 0 <= K.
% joint_discriminant_eventually (SmallOrderPrimeLosses, DiscriminantSettingSatisfied.lean), windows_ev
% and parTie (PowFourLtCeilThresholdPrime.lean), loss_setting_family, with IO_at, Ga_at, ASK_at and
% Lnk.askIsF. T nonempty is hne from ChTT_at. The open step left is the proof of Lemma
% asm-global-link-losses-explicit. Size: proof closure of losses_of_small_share_weight, about 380
% theorems and 5,600 lines in 28 files.
\end{proof}

\begin{theorem}\label{asm-global-link}
Let $\mathcal A$ be a class of families in which every family carries its duty
data (Definition~\ref{asm-duty}), and let
$\mathcal D_3=\operatorname{Des}(\mathcal A_N,3)$
(Definition~\ref{asm-designed}). Suppose that the second-layer conditions at
the real place hold eventually along $\mathcal D_3$, that the mass ratio under
volume holds for every definite-side datum, and that the mass formula for
genera of rank $2$ holds for every definite-side datum. Then, eventually along
$\mathcal D_3$:
\begin{enumerate}
\item if $\mathcal F$ has chosen constants $(3,C_H)$, then the rank
  conditions hold;
\item if both size conditions hold, then the global link holds with
  $C_{\mathrm{abs}}=2$, $C_0=3$, $e_D=2$ and $e_L=0$.
\end{enumerate}
\end{theorem}
\begin{proof}
\emph{Step 1: facts along $\mathcal D_3$.} Eventually along $\mathcal D_3$,
every family has the following properties at every member at every large
scale.
\begin{enumerate}
\item[(a)] The ratio of the masses of consecutive standard lattices is given by
  Siegel's mass formula (Theorem~\ref{cv-siegel-mass-ratio}). This holds for
  every family.
\item[(b)] The volume of the base space equals the mass of the standard
  lattice of rank $r$: $\operatorname{vol}_D(Y_0)=\operatorname{mass}(I_r)$.
  The second-layer conditions at the real place say in particular that the
  theta realization of every member at a large scale is tied at the real
  place, and this tie contains the volume tie
  (Definition~\ref{asm-facts-volume-tie}(1)). With the mass ratio under volume
  and the mass formula for genera of rank $2$, which hold for every
  definite-side datum by hypothesis, Lemma~\ref{asm-facts-base-volume} gives
  the equality.
\item[(c)] For every place $\tau$ and every $\epsilon_0$ with
  $e^{-r}\le\epsilon_0<1/2$, if the spectral gap of the parameter of $\tau$ is
  large enough for $\epsilon_0$, the formal degree $d_\sigma(\tau)$ of the
  discrete series of $\tau$ satisfies $(1-\epsilon_0)\dim V(\lambda_\tau)\le
  d_\sigma(\tau)\le4d(\lambda_\tau)(3pq|\lambda_\tau|)^{pq}$, where
  $V(\lambda_\tau)$ is the representation of highest weight $\lambda_\tau$ and
  $(p,q)$ the signature. This holds for every family.
\item[(d)] At every word $u$ whose columns have values, the slope $\mu(M_u)$
  is bounded below by the sizes of the columns and $\log T_u$, and
  $\mu_{\max}(M_u)$ is bounded above in the same way with the constant $4$.
  These are the slope inequality of Theorem~\ref{hs-bost-slope-inequality}
  applied to the coordinates of the columns, and they hold for every family.
  In the form used in Step 4, every family satisfies the column slope bounds
  of Definition~\ref{asm-global-numerical}, by
  Lemma~\ref{asm-division-kernel-first-lattice-slopes}, whose proof goes
  through these two bounds.
% Lean: fact (d) is the fields hZa and hZb of Wire.FamGLmin, that is the per-member bounds
% SlopeLowerOpenPart and SlopeMaxOpenPart 4 (SlopeLow and SlopeMax at every word with ColVal), given by
% lower_slope_inequality_family and slope_inequality_family_four. The column slope bounds ConclusionF5
% and ConclusionF6 are derived from them inside linkGLOBAL7_min by familySlopeLower_of_open and
% familySlopeMax_of_open. Lemma asm-division-kernel-first-lattice-slopes states the composite, and its
% open step covers both the per-member bounds and the derivation.
\item[(e)] The columns of every word have values
  (Definition~\ref{asm-global-column-values}). By
  Lemma~\ref{asm-global-second-layer}, the second-layer conditions at the real
  place give, at every member at every large scale, tied columns, and trivial
  stabilizers $\Gamma_c$ for the classes $c\in\mathrm{cls}_0$ of every word.
  Tied columns with these stabilizers give values to the columns of every
  word, by Lemma~\ref{asm-global-link-column-values}.
\item[(f)] At every word at which the setting of the level primes holds
  (Definition~\ref{asm-level-setting}(1)), the five local facts of
  Definition~\ref{asm-level-local-facts} hold: a free basis at $T_u$, the slope
  line, Hecke integrality at level zero, the ramified level zero and the second
  branch inclusion. By Lemma~\ref{asm-global-second-layer}, at every member at
  every large scale, every word $u$ has the forms in degree $0$ of the tied
  columns, whose character has finite order with values in $L$, and the classes
  in degree $0$ of $u$ have trivial stabilizers. At a word where the setting
  holds, these give the five local facts, by
  Lemma~\ref{asm-global-link-local-facts}.
\end{enumerate}

\emph{Step 2: part (1).} Let $\mathcal F\in\mathcal D_3$ be large, with
chosen constants $(3,C_H)$. As a designed member it satisfies the size
condition on $T_u$, so by Lemma~\ref{asm-level-setting-holds}, for all large
$x$, the setting of the level primes holds at every word of every member at
$x$. At these words the local facts of (f) hold. With the free basis and the
slope line, Lemma~\ref{asm-level-division-side} gives the division side
conditions, and with Hecke integrality, the ramified level zero and the second
branch inclusion, Lemma~\ref{asm-level-congruence} gives the congruence length
bound. So $\mathcal F$ has the level data
(Definition~\ref{asm-level-setting}). They include the equality
\[
 \dim M_{u,\mathrm{ram}}=\dim M_u\Big(1-\frac1{T_u-1}\Big)^f
\]
at every word $u$, and the setting includes $T_u>2fN(m)+1\ge3$.

Among the choices of the family (Definition~\ref{pd-family}(4)), at every
member at a scale $x>x_1$ and every word $u$: $N^\circ_u\subset
M^{\mathrm{ram}}_u$ locally at every finite place not over $T_u$, and
$M^{\mathrm{ram}}_u\subset N^\circ_u$ locally at those places whose residue
characteristic is dropped or not inert; at a place $w$ over $T_u$, a vector of
$V(N_u)$ lies in $N^\circ_u$ locally at $w$ exactly when its image under
$\nu_u$ lies in $N_\lambda(u,w)$; and $N_u$ is the $\mathcal O_L$-span of
$I_uN^\circ_u$ for the nonzero fractional ideal $I_u$ of the proof of
Lemma~\ref{asm-slope-link}. With the inclusions $N_\lambda(u,w)\subset
S^{\mathrm{ram}}_w$ of the congruence length bound,
Lemma~\ref{asm-level-commensurable} shows that, for all large $x$, at every
member and word, $N^\circ_u$ is commensurable with $M^{\mathrm{ram}}_u$:
$N^\circ_u\subset M^{\mathrm{ram}}_u$, $N^\circ_u$ is finitely generated, and
$cM^{\mathrm{ram}}_u\subset N^\circ_u$ for a nonzero $c\in\mathcal O_L$. The
same lemma shows that every $N_u$ is finitely generated and spans its space
$V(N_u)$. The map $j_N$ carries a basis
of $V(N_u)$ to a basis of $\mathcal M_u\cong M_{u,\mathrm{ram}}$
(Definition~\ref{pd-period-data}(7)), so
$d_u=\operatorname{rk}N_u=\dim M_{u,\mathrm{ram}}$. The proviso
$3^{-1}\le\dim M_u\cdot\Gamma^2/D^{*}\le3$ of the chosen constants gives $\dim
M_u\ge D^{*}/(3\Gamma^2)>0$, and $1-1/(T_u-1)>0$. So $d_u>0$. These are the
rank conditions of Definition~\ref{asm-twisted-links}(3).

\emph{Step 3: part (2), the inputs.} Let $\mathcal F\in\mathcal D_3$ be large
and satisfy both size conditions; assume that it has chosen constants
$(3,C_H)$, with constants $C_5,c_6,C_b>0$, and the global inputs with
$(2,C_H,3)$ (Definition~\ref{asm-global-inputs}). Largeness gives $C_H>0$. By
Step 2, $\mathcal F$ has the level data, the commensurability of $N^\circ_u$
with $M^{\mathrm{ram}}_u$, and the rank conditions.

\emph{Step 4: the bounds of Definition~\ref{asm-member-bounds}.}
\begin{enumerate}
\item The rank of $RR$ is positive at all large scales, by
  Lemma~\ref{asm-bounds-glued-rank} with the orthogonality link with
  $C_0=3$, the proviso, the irreducibility link and the division link.
\item The count holds at all large scales. The mass ratio (a), the base volume
  (b) and the degree bounds (c) give the explicit bounds of the Vandermonde
  count: a constant $C$ with the bounds (1) to (3) of
  Lemma~\ref{asm-global-count} at every member at every large scale, by
  Lemma~\ref{asm-global-link-count-bounds}. That lemma gives the count.
\item The bounds for $\chi_{-d}$ in the chosen constants give
  $|\operatorname{Re}\frac{L'}{L}(1,\chi_d)|\le C_5\log x$. Since
  $\log x\le(2/\varepsilon)x^{\varepsilon/2}$, this is at most
  $\max(C_5,0)(2/\varepsilon)x^{\varepsilon/2}\le x^{\varepsilon}$ once
  $\max(C_5,0)(2/\varepsilon)\le x^{\varepsilon/2}$, that is, for large $x$.
  With the column slope bounds (d), the values of the columns (e) and the count
  of item 2, Lemma~\ref{asm-global-slopes-m} gives the slope bounds for $M$ at
  all large scales.
\item The size condition on $T_u$, the level data, the commensurability of
  $N^\circ_u$ with $M^{\mathrm{ram}}_u$ and the rank conditions of Step 2, the
  gain link, the slope link and the depth bound give the slope bounds for $N$
  at all large scales, by Lemma~\ref{asm-global-slopes-n}. The depth bound
  of that lemma holds for every family, by
  Lemma~\ref{asm-global-link-depth}.
\item The division link, the rank conditions and the positivity of the rank of
  $RR$ give the lattice conditions, by Lemma~\ref{asm-bounds-lattice}.
\item Lemma~\ref{asm-bounds-frame} gives the frame bound with $C_H$. Its
  hypotheses are the orthogonality link with $C_0=3$, the irreducibility and
  division links, the proviso and the bound $3\Gamma^2x^{C_H}\le|W|$ of the
  chosen constants, the rank conditions, the equality
  $d_u=\dim M_{u,\mathrm{ram}}=\dim M_u(1-1/(T_u-1))^f$ of Step 2, the size
  condition on $T_u$, the lattice conditions of item 5, and $C_H>0$.
\item The node bounds hold at all large scales and the character conditions
  hold for every family, by Lemma~\ref{asm-bounds-nodes-characters}.
\item Let $c_6>0$. The five inequalities between powers of $x$ and
  exponentials in Lemma~\ref{asm-bounds-windows} hold for all large $x$, and
  the lower bound for $|T_2|$ in the chosen constants gives
  $|T_2|\ge x^{-\varepsilon}n$ for all large $x$. So that lemma gives the
  window conditions with $e_K=5$.
\item The bound $\delta_1n/2\le|T|\le2\delta_1n$ and the bound on the bad tops
  with $C_b>0$ in the chosen constants give, by
  Lemma~\ref{asm-global-link-bad-primes}, the bound
  $\sum_{p\in P}\frac{\log p}p\beta(p)\le|T|x^{-2+\varepsilon}$ of that lemma,
  over the large-order residue primes, at every member for all large $x$.
  With the bound on $|T|$ and the gain link,
  Lemma~\ref{asm-global-link-losses} gives $C\ge0$ and $x_0$ such that every
  member at every scale $x\ge x_0$ whose weight vectors have small shares
  satisfies the loss bounds of Definition~\ref{asm-member-bounds}(6) with
  $(C,0)$. The duty data give these small shares at every member for all large
  $x$ (Definition~\ref{asm-duty}(5)). Lemma~\ref{asm-global-losses} gives the
  loss bounds with $e_L=0$.
\end{enumerate}
% Lean. family_conditions_of_layer_two and stabF (Lemma asm-global-second-layer), count_at_of and
% countF_of (Lemma asm-global-count), slopesMF_of_count_fam with nums_wts_of_fam (Lemmas
% asm-global-numerical-hold and asm-global-slopes-m), lossF_of_7 with chosenWS_of_wall (Lemma
% asm-global-losses), and item 4: slopesNF_of with Places.placesF_min' and
% J.lemma_five_one_of_theorem_two (Lemmas asm-global-slope-index, asm-global-unit-weight,
% asm-global-level-t, asm-global-slopes-n-member and asm-global-slopes-n).
% (e) is RT.tCV_of_tied with ColValS.colValOfTiedS_holds (Lemma asm-global-link-column-values).
% (f) is slope_free_branch_family with freeTT_holds, slopeLineT_holds and
%   LemB.second_branch_all_members; hI_tied with PI.IntegralAtBasePointLattice.integral_at_t0_lattice;
%   and ramified_level_zero_holds with UnramifiedTrace.hecke_ramified_zero_all (Lemma
%   asm-global-link-local-facts).
% Item 2: bounds (1)-(3) are StandardLattice.PartFiveFamilyBounds at C0 = max 5 (CYL+1)
%   (Conclusion5₂: TorusPairsWithVariant at Gamma = GamP with the window D.T_win, and
%   VandermondeVariationBound), with DataB from JDB.dataB_fam and log X_u = logXuD by
%   jnt_logXu (factDuStar_fam, factDstar_fam); PartFiveFamilyBounds is
%   family_part_five_of_explicit_constants (division_level_prime_family h2b h2v h3)
%   (Lemma asm-global-link-count-bounds).
% Item 4: the depth bound is CornerGLContinuity.FamilyWeightSumBound, given for every
%   family by G.TopChoiceBounds.tau_sum_family_of_mertens Fm (fun _ _ => mertens) (Lemma
%   asm-global-link-depth).
% Item 9: Loss7.losses_of_small_share_weight (Lemma asm-global-link-losses) with
%   Lnk.chosenBadp_of_old for the bad primes (Lemma asm-global-link-bad-primes).
% Lemmas asm-global-link-column-values and asm-global-link-count-bounds-weights-hold are written
% out. The open steps that remain are in the proofs of Lemmas asm-global-link-local-facts-free,
% -slope, -integral, -ramified and -branch (the five local facts), asm-global-link-count-bounds-
% explicit-rows, -log-x and -variation, and asm-global-link-losses-explicit (the identification of
% the left side of (1) and its side conditions); Lemma asm-global-link-depth is written out through Lemmas
% asm-global-link-depth-sizes-hold, asm-global-link-depth-representative and
% asm-global-link-depth-cuts.

\emph{Step 5: the tie.} Lemma~\ref{asm-tie-from-bounds} applies with
$C_{\mathrm{abs}}=2$, $e_L=0$ and the constants $C_5$ and $c_6$ of the chosen
constants: its remaining hypotheses are the bounds on $|T|$, on
$\chi_{-q},\chi_{-d},\chi_{-q}\chi_{-d}$ and the thin-scale bound, which are
chosen constants, and the irreducibility and division links, which are global
inputs. So $\mathcal F$ is tied at every large scale with constants
$(2,0,K,C_H,e_K,C_5)$. This is the global link with $C_{\mathrm{abs}}=2$,
$C_0=3$, $e_D=2$ and $e_L=0$.
\end{proof}
\leanref{Development/AbelianVarietyTheory/Proofs/GlobalInterpolation/LinkGl.lean}{GlobalTheory.Top.link_gl}
\leanref{Formalism/SharpRestriction/GlobalSchurBounds/TuSizeOfRealPlaceSecondLayer.lean}{Formalism.MainNotions.Wire.famGLmin_of_realPlaceSecondLayer}
\leanref{Formalism/SharpRestriction/GlobalSchurBounds/LogXLowerFamilyOfTheorems.lean}{Formalism.MainNotions.Wire.FamGLmin}
\leanref{Formalism/SharpRestriction/GlobalSchurBounds/LinkMin.lean}{Formalism.MainNotions.linkGLOBAL7_items}
\leanref{Formalism/SharpRestriction/GlobalSchurBounds/LogXLowerFamilyOfTheorems.lean}{Formalism.MainNotions.Wire.FamGLmin.muRank}
\leanref{Formalism/SharpRestriction/GlobalSchurBounds/LogXLowerFamilyOfTheorems.lean}{Formalism.MainNotions.Wire.family_theorems_two_and_five}
\leanref{Formalism/SharpRestriction/GlobalSchurBounds/LinkMin.lean}{Formalism.MainNotions.lvlF_of}
\leanref{Formalism/SharpRestriction/GlobalSchurBounds/LinkMin.lean}{Formalism.MainNotions.nuCommensF_min}

\subsection{Existence of designed families}\label{subsection-asm-designed-exist}

\begin{definition}\label{asm-linked-aligned}
Let $\mathcal F$ be a family with exponent $\varepsilon$ (Definition~\ref{pd-family}). A member
of $\mathcal F$ \emph{carries tied data} if it has tied data in the sense of the twist records of
Definition~\ref{asm-definite-side-hypotheses}.
\begin{enumerate}
\item $\mathcal F$ is \emph{linked and aligned} if, at all large scales, every member carries tied
  data and has distinct tame data (Definition~\ref{asm-box-row-tame}(3)), and $\mathcal F$
  satisfies items (1), (2) and (4) of Definition~\ref{asm-duty}.
\item $\mathcal F$ is a \emph{linked and aligned family of the second layer} if it is linked and
  aligned and, at all large scales, every member carries a twist record.
\end{enumerate}
\end{definition}
\leanref{Development/AutomorphicArithmetic/ThetaCorrespondence/ThetaPairLifts/ThirdPlaceSeries/LinkedFamiliesFixedValues.lean}{Formalism.L2.ExtraDataSelfDual}
\leanref{Development/AutomorphicArithmetic/ThetaCorrespondence/ThetaPairLifts/ThirdPlaceSeries/LinkedFamiliesFixedValues.lean}{Formalism.L2.SecondChernLinkedClass}
\leanref{Development/AutomorphicArithmetic/ThetaCorrespondence/ThetaPairLifts/ThirdPlaceSeries/LinkedFamiliesFixedValues.lean}{Formalism.L2.LinkedFamiliesFixedValues'}
\leanref{Development/AutomorphicArithmetic/ThetaCorrespondence/ThetaPairLifts/ThirdPlaceSeries/LinkedFamiliesFixedValues.lean}{Formalism.L2.LinkedFamiliesAligned'}
\leanref{Development/AutomorphicArithmetic/ThetaCorrespondence/ThetaPairLifts/ThirdPlaceSeries/LayerTwoFamilyClass.lean}{Formalism.L2.SecondLayerFamilyClass}

\begin{lemma}\label{asm-aligned-second-layer}
Let $\mathcal A$ be a class of families with a twist selection: for every family
$\mathcal F\in\mathcal A$ a threshold $x_0$ and, at every member of $\mathcal F$ at every scale
$x>x_1$ with $x\ge x_0$, a twist record (Definition~\ref{asm-definite-side-hypotheses}). Suppose
that every family of $\mathcal A$ carries its duty data (Definition~\ref{asm-duty}). Then every
family of $\mathcal A$ is a linked and aligned family of the second layer
(Definition~\ref{asm-linked-aligned}).
\end{lemma}
\leanref{Development/AutomorphicArithmetic/ThetaCorrespondence/ThetaPairLifts/ThirdPlaceSeries/ClassOfRecordOfChoice.lean}{Formalism.L2.Top3.layerTwoFamilyClass_of_linked_aligned}
\leanref{Development/AutomorphicArithmetic/ThetaCorrespondence/ThetaPairLifts/ThirdPlaceSeries/ClassOfRecordOfChoice.lean}{Formalism.L2.Top3.linkedFamiliesAligned_of_duty}
\leanref{Development/AutomorphicArithmetic/ThetaCorrespondence/ThetaPairLifts/ThirdPlaceSeries/ClassOfRecordOfChoice.lean}{Formalism.L2.Top3.tied_families_of_choice_alt}
\leanref{Development/AutomorphicArithmetic/ThetaCorrespondence/ThetaPairLifts/ThirdPlaceSeries/LayerTwoFamilyClass.lean}{Formalism.L2.SecondLayerFamilyClass.c2s}
\begin{proof}
Let $\mathcal F\in\mathcal A$, and let $x_0$ be the threshold of its twist selection. At every
member at every scale $x\ge x_0$ with $x>x_1$ the selection chooses a twist record, and the tied
data of that record are tied data of the member. So, at all large scales, every member carries a
twist record and tied data. The duty data of $\mathcal F$ contain items (1), (2) and (4) of
Definition~\ref{asm-duty}, and by item (5) every member has distinct tame data for all large $x$.
So $\mathcal F$ is linked and aligned, and it is a linked and aligned family of the second layer.
\end{proof}

\begin{lemma}\label{asm-proviso-splitting}
Let $D$ be a definite-side datum (Definition~\ref{th-definite-datum}) with host fields
$F\subset E$, $f=[F:\mathbf Q]$, and level primes $g_0,g_1$. Then $g_0$, $g_1$ and the auxiliary
prime $T_u$ of every word $u$ split completely in $F$: each such prime $p$ has $f$ places of $F$
above it, and the residue field of each of them has $p$ elements.
\end{lemma}
\leanref{Formalism/SharpRestriction/AdmissibleParameterJoins/SfromGLu.lean}{Formalism.MainNotions.J.splitsF_g}
\leanref{Formalism/SharpRestriction/AdmissibleParameterJoins/DIVfromGLz.lean}{Formalism.MainNotions.J.level_prime_splits_totally_real}
\begin{proof}
Among the host data of $D$, the primes $g_0$ and $g_1$ split completely in $E$. By
Definition~\ref{th-definite-datum}(4), $T_u=\pi\bar\pi$ with $\pi\equiv1$ modulo $N\mathcal O_E$,
where $N=Q_Sqd$; then $\bar\pi\equiv1$ as well, and $T_u\equiv1\pmod N$ since
$N\mathcal O_E\cap\mathbf Z=N\mathbf Z$. So $T_u$ is prime to $N$, and by
Proposition~\ref{hf-host-field-properties}(c) the Frobenius of every prime of $E$ over $T_u$ is the
image of $\operatorname{art}(T_u)=\operatorname{art}(1)=1$; hence $T_u$ splits completely in $E$.
A prime $p$ that splits completely in $E$ splits completely in $F\subset E$: the places of $F$ over
$p$ have degree $1$, since the places of $E$ over them do, and each has at most $[E:F]$ places of
$E$ above it, so there are at least $[E:\mathbf Q]/[E:F]=f$ of them, and at most $f$.
\end{proof}

\begin{lemma}\label{asm-proviso-trivial-stabilizers}
Let $D$ be a definite-side datum (Definition~\ref{th-definite-datum}) with ranks $k$ and $f$, host
fields $F\subset E$ and level primes $g_0,g_1$, and suppose that $g_i>2fk+1$ for some $i$. Let $u$
be a word, let $X_0$ be the bottom space of its chain, of rank $k$, and let $\Gamma_c$ be the
stabilizer of a class $c\in\mathrm{cls}_0$ of the chain. Let level data for $u$ be given: a compact
open subgroup $U_u$ of $U(X_0)(\mathbf A_{F,f})$ such that, for every place $v$ of $F$ over $g_0$
or $g_1$, under an isomorphism $E\otimes_FF_v\cong F_v\times F_v$ the first components of the
elements of the component of $U_u$ at $v$ lie in a conjugate of the group $I_v$ of matrices in
$GL_k(\mathcal O_{F_v})$ that are upper unitriangular modulo $v$; elements
$y_c\in U(X_0)(\mathbf A_{F,f})$ for $c\in\mathrm{cls}_0$; and injective homomorphisms
$\Gamma_c\to U(X_0)(F)\cap y_cU_uy_c^{-1}$. Then $U(X_0)(F)\cap y_cU_uy_c^{-1}=1$, and so
$\Gamma_c=1$, for every $c\in\mathrm{cls}_0$.
\end{lemma}
\leanref{Formalism/SharpRestriction/GlobalSchurBounds/TheoremThreeFamilyOfFormalDegreeInput.lean}{Formalism.MainNotions.StandardLattice.HarishChandraReduction.column_stabilisers_trivial}
\leanref{Formalism/SharpRestriction/GlobalSchurBounds/ThickPlacesCorollaryOfAnyTranslate.lean}{Formalism.MainNotions.StandardLattice.levelDimension_of_admissible}
\leanref{Formalism/SharpRestriction/GlobalSchurBounds/Joints.lean}{Formalism.MainNotions.StandardLattice.column_lemma_of_definite_form}
\leanref{Formalism/LocalToGlobal/GlobalTheory/GL5/LocalEndomorphismContinuous.lean}{Formalism.L2.GlobalTheoryTheory.GenusMassVolumes.levels_pro_g}
\begin{proof}
Let $v$ be a place of $F$ over $g_i$. The group $I_v$ is a pro-$g_i$ group: the kernel of
reduction modulo $v$ is a pro-$g_i$ group, as $v$ has residue characteristic $g_i$, and the upper
unitriangular matrices over the residue field form a $g_i$-group. An element of
$U(X_0)(F_v)$ is determined by its first component, since the second is its inverse adjoint for
the Hermitian form. So the component of $U_u$ at $v$ is a pro-$g_i$ group.

Let $\gamma\in U(X_0)(F)\cap y_cU_uy_c^{-1}$. Since $X_0$ is definite and $U_u$ is compact,
$\gamma$ has finite order. Its image at $v$ lies in a conjugate of the component of $U_u$ at $v$,
so its order is a power of $g_i$. If $\gamma\ne1$, a power of $\gamma$ has order $g_i$, and a
primitive $g_i$-th root of unity is an eigenvalue of $\gamma$ acting on $E^k$, a vector space of
dimension $[E:\mathbf Q]k=2fk$ over $\mathbf Q$. That root of unity has degree $g_i-1$ over
$\mathbf Q$, so $g_i-1\le2fk$, against the hypothesis. Hence the intersection is trivial, and so is
$\Gamma_c$, which embeds in it. The proof of Proposition~\ref{cv-class-number-formula} uses the
same argument.
\end{proof}

\begin{lemma}\label{asm-proviso-bottom-dimension}
Let $D$ be a definite-side datum (Definition~\ref{th-definite-datum}) with ranks $k$ and $f$, set
$\mathcal T$ of $f$ embeddings and parameters $\lambda_u(\tau)$, and suppose that $D$ carries tied
data (Definition~\ref{asm-linked-aligned}). For $y\in\mathbf R^k$ put
\[
 \operatorname{wdim}(y)=\prod_{1\le i<j\le k}\frac{y_i-y_j}{j-i},
\]
and for a word $u$ and $\tau\in\mathcal T$ let $u(\tau)=(u_1(\tau),\dots,u_k(\tau))$ be the first
$k$ entries of $\lambda_u(\tau)$. Let $u$ be a word such that every class $c\in\mathrm{cls}_0$ of
its chain $C_u$ has trivial stabilizer $\Gamma_c$. Then
\[
 \dim M_u=|\mathrm{cls}_0|\prod_{\tau\in\mathcal T}\operatorname{wdim}(u(\tau)).
\]
\end{lemma}
\leanref{Formalism/SharpRestriction/GlobalSchurBounds/GlobalRankLatticeProofs/UnitTraceSumBounds/RankConditionsOfLayerTwo.lean}{Formalism.MainNotions.Ext.Lw.mu0Tie_of_sym}
\leanref{Formalism/SharpRestriction/GlobalSchurBounds/StandardLattice.lean}{Formalism.MainNotions.StandardLattice.Mu0Tie}
\leanref{Formalism/SharpRestriction/GlobalSchurBounds/GlobalRankLatticeProofs/UnitTraceSumBounds/RankConditionsOfLayerTwo.lean}{Formalism.MainNotions.Ext.Lw.finrank_lowest_forms_space}
\leanref{Formalism/SharpRestriction/GlobalSchurBounds/GlobalRankLatticeProofs/UnitTraceSumBounds/RankConditionsOfLayerTwo.lean}{Formalism.MainNotions.Ext.Lw.fiC_eq_pi}
\leanref{Formalism/SharpRestriction/GlobalSchurBounds/GlobalRankLatticeProofs/UnitTraceSumBounds/RankConditionsOfLayerTwo.lean}{Formalism.MainNotions.Ext.Lw.card_ssIdx_eq}
\leanref{Formalism/SharpRestriction/GlobalSchurBounds/GeneralLinearActions/GapAtLeast.lean}{Formalism.MainNotions.StandardLattice.dimV}
\begin{proof}
Among the tied data are tied columns. At level $0$ and the word $u$ they give: a compact open
subgroup $U$ of $U(X_0)(\mathbf A_{F,f})$, where $X_0$ is the bottom space of $C_u$, with a
character $\chi$ of $U$; elements $y_c$, $c\in\mathrm{cls}_0$, that represent each double coset in
$U(X_0)(F)\backslash U(X_0)(\mathbf A_{F,f})/U$ once; for each $c$ a homomorphism from $\Gamma_c$
onto $U(X_0)(F)\cap y_cUy_c^{-1}$; and an identification of $Y_0$ with the set of tuples
$(z_\tau)_{\tau\in\mathcal T}$ of invertible complex $k\times k$ matrices. Under the last
identification, $M_u$ is the space of tuples $(\varphi(y_c))_{c\in\mathrm{cls}_0}$, where $\varphi$
runs through the forms of weight $\lambda_{u,0}$ and level $(U,\chi)$: the functions $\varphi$ on
$U(X_0)(\mathbf A_{F,f})$ with values in $V=\bigotimes_{\tau}V_{\lambda_{u,0}(\tau)}$ such that
$\varphi(\gamma g\kappa)=\chi(\kappa)\,\gamma\cdot\varphi(g)$ for $\gamma\in U(X_0)(F)$ and
$\kappa\in U$. Here $V_\lambda$ is the irreducible representation of $GL_k(\mathbf C)$ of highest
weight $\lambda$, realized on functions of $z_\tau$, and $U(X_0)(F)$ acts on $V$ through the
embeddings $\tau$.

Since the $\Gamma_c$ are trivial, so are the groups $U(X_0)(F)\cap y_cUy_c^{-1}$. So every element
of $U(X_0)(\mathbf A_{F,f})$ is $\gamma y_c\kappa$ for a unique triple $(c,\gamma,\kappa)$ with
$\gamma\in U(X_0)(F)$ and $\kappa\in U$, and every tuple $(v_c)\in V^{\mathrm{cls}_0}$ is
$(\varphi(y_c))_c$ for a unique form $\varphi$, namely
$\varphi(\gamma y_c\kappa)=\chi(\kappa)\,\gamma\cdot v_c$. Hence
$\dim M_u=|\mathrm{cls}_0|\prod_\tau\dim V_{\lambda_{u,0}(\tau)}$.

The space $V_\lambda$ has a basis indexed by the semistandard tableaux of shape $\lambda$, shifted
to be nonnegative, with entries at most $k$, and by Weyl's dimension formula
\cite[Theorem~6.3]{FH91} their number is $\prod_{i<j}(\lambda_i-\lambda_j+j-i)/(j-i)$. By
Definition~\ref{th-definite-datum}(7), $\lambda_{u,0}(\tau)_i=u_i(\tau)-\frac{k+1}2+i$ for
$1\le i\le k=N(0)$, so $\lambda_i-\lambda_j+j-i=u_i(\tau)-u_j(\tau)$, and
$\dim V_{\lambda_{u,0}(\tau)}=\operatorname{wdim}(u(\tau))$.
\end{proof}

\begin{lemma}\label{asm-proviso-window}
Let $f,k\ge1$ be integers with $2f\le2^k$, and let $T\ge2$, $\Gamma\ge1$, $D^*>0$ and $d^*>0$ be
real numbers. Put $X=(D^*/d^*)^{1/(fk)}$ and $T^*=X\Gamma^{-2/(fk)}$, and suppose that
$T^*\le T\le(1+1/(fk))T^*$. Then
\[
 \frac{D^*}{2\Gamma^2}\le d^*(T^k-1)^f\le e\,\frac{D^*}{\Gamma^2}
 \qquad\text{and}\qquad
 \log X-\frac{2\log\Gamma}{fk}\le\log T\le\log X+\log2 .
\]
\end{lemma}
\leanref{Formalism/GL3MeanSquares/HookFormulaTorsion/OneAddInvPowLeExpOne.lean}{Formalism.HermitianLattices.GenusMassVolumes.mass_corollary}
\leanref{Formalism/GL3MeanSquares/RankMulPowThreeHalvesLeFourthPower.lean}{Formalism.HermitianLattices.GenusMassVolumes.PowerProductTwoSidedBounds}
\begin{proof}
By definition $d^*(T^*)^{fk}=D^*/\Gamma^2$. For the upper bound,
$(T^k-1)^f\le T^{fk}\le(1+1/(fk))^{fk}(T^*)^{fk}\le e\,(T^*)^{fk}$, since $(1+1/N)^N\le e$ for
$N\ge1$. For the lower bound, $T^k\ge2^k\ge2f$, so $fT^{-k}\le1/2$, and Bernoulli's inequality
gives $(1-T^{-k})^f\ge1-fT^{-k}\ge1/2$. As $T\ge T^*$,
\[
 (T^k-1)^f=T^{fk}(1-T^{-k})^f\ge\tfrac12(T^*)^{fk}.
\]
Multiplying both bounds by $d^*$ gives the first assertion. For the second,
$\log T\ge\log T^*=\log X-2\log\Gamma/(fk)$; and since $\Gamma\ge1$ and $fk\ge1$,
$T\le2T^*\le2X$.
\end{proof}

\begin{proposition}\label{asm-proviso-reduction}
Let $\mathcal F$ be a family with exponent $\varepsilon$ (Definition~\ref{pd-family}) such that, at
all large scales, every member carries tied data (Definition~\ref{asm-linked-aligned}). Suppose
that the class number formula of Proposition~\ref{cv-class-number-formula} holds at every member
of $\mathcal F$ at every scale $x>x_1$, and let $C_0\ge e$. Then there is $x_2>x_1$ such that for
every $x\ge x_2$, every member $D$ at $x$, with $\Gamma=\exp(C_\Gamma m^5fx^{\varepsilon})$, and
every word $u$ of $D$,
\[
 C_0^{-1}\le\dim M_u\cdot\Gamma^2/D^{*}\le C_0 .
\]
\end{proposition}
\leanref{Development/AutomorphicArithmetic/ThetaCorrespondence/ThetaPairLifts/ThirdPlaceSeries/ClosedRankTiesFC2S.lean}{Formalism.MainNotions.Ext.Lw.second_chern_linked_closed_rank_ties}
\leanref{Development/AutomorphicArithmetic/ThetaCorrespondence/ThetaPairLifts/ThirdPlaceSeries/ClosedRankTiesFC2S.lean}{Formalism.MainNotions.Ext.Lw.rankTiesF_of_C2S_alt}
\leanref{Formalism/SharpRestriction/GlobalSchurBounds/MainSup.lean}{Formalism.MainNotions.RankTiesF}
\leanref{Formalism/SharpRestriction/GlobalSchurBounds/TheoremThreeFamilyOfFormalDegreeInput.lean}{Formalism.MainNotions.StandardLattice.HarishChandraReduction.chosenProviso_min}
\leanref{Formalism/SharpRestriction/GlobalSchurBounds/TheoremThreeFamilyOfFormalDegreeInput.lean}{Formalism.MainNotions.StandardLattice.HarishChandraReduction.chosenProviso_4f}
\leanref{Formalism/SharpRestriction/GlobalSchurBounds/TheoremThreeFamilyOfFormalDegreeInput.lean}{Formalism.MainNotions.StandardLattice.HarishChandraReduction.rankFormula_of_classNumber}
\leanref{Formalism/SharpRestriction/GlobalSchurBounds/StandardLattice.lean}{Formalism.MainNotions.StandardLattice.level_rank_of_trivial_stabilisers}
\leanref{Formalism/SharpRestriction/GlobalSchurBounds/MainSup.lean}{Formalism.MainNotions.proviso_e}
\leanref{Formalism/SharpRestriction/GlobalSchurBounds/TopDensityNGeSixteen.lean}{Formalism.MainNotions.LinkB.proviso_mono}
\begin{proof}
\emph{Step 1: the rank ties.} Let $x_0$ be such that every member at every scale $x\ge x_0$ with
$x>x_1$ carries tied data. At such a member, the primes $g_0$, $g_1$ and every $T_u$ split
completely in $F$ by Lemma~\ref{asm-proviso-splitting}, and
$\dim M_u=|\mathrm{cls}_0|\prod_\tau\operatorname{wdim}(u(\tau))$ for every word $u$ whose classes
$c\in\mathrm{cls}_0$ have trivial stabilizers, by Lemma~\ref{asm-proviso-bottom-dimension}.

\emph{Step 2: a large level prime.} Let $D$ be a member at a scale $x>x_1$. The sizes of the family
(Definition~\ref{pd-family}(4)) give $(kb)^8\le g_0$. Since $b=\exp(x^\varepsilon)\ge2$ and
$f\le k$ (Definition~\ref{th-ranks}), $g_0\ge(2k)^8>2k^2+1\ge2fk+1$.

\emph{Step 3: the rank formula.} Put $x_2=\max(x_0,x_1+1)$, let $x\ge x_2$, let $D$ be a member
at $x$ and let $u$ be a word. By Step 2 and Lemma~\ref{asm-proviso-trivial-stabilizers}, applied
with the level data among the further choices of the family (Definition~\ref{pd-family}(4)), the
classes $c\in\mathrm{cls}_0$ of $u$ have trivial stabilizers. So Step 1 gives
$\dim M_u=|\mathrm{cls}_0|\prod_\tau\operatorname{wdim}(u(\tau))$. By Steps 1 and 2 the hypotheses
of the class number formula hold at $D$ and $u$, and it gives
\[
 |\mathrm{cls}_0|=\operatorname{mass}(I_k)\prod_{i=0,1}
 \Big(\frac{|GL_k(\mathbf F_{g_i})|}{|N_k(\mathbf F_{g_i})|}\Big)^{f}(T_u^k-1)^{f}.
\]
The further choices of the family fix
\[
 d^*_u=\operatorname{mass}(I_k)\prod_{i=0,1}
 \Big(\frac{|GL_k(\mathbf F_{g_i})|}{|N_k(\mathbf F_{g_i})|}\Big)^{f}
 \prod_{\tau\in\mathcal T}\operatorname{wdim}(u(\tau)).
\]
Hence $\dim M_u=d^*_u(T_u^k-1)^f$.

\emph{Step 4: the proviso.} Keep $x$, $D$ and $u$ as in Step 3. By
Definition~\ref{pd-size-data}(9), $\Gamma=\Gamma_P$. By Definition~\ref{th-definite-datum}(3)
and~(4), $T^*_u=(D^*/d^*_u)^{1/(fk)}\Gamma_P^{-2/(fk)}\le T_u\le(1+1/(fk))T^*_u$, with $D^*>0$,
$d^*_u\ge1$ and $\Gamma_P>1$. The rank datum has $1\le f\le k$, so $2f\le2k\le2^k$, and the prime
$T_u$ is at least $2$. Lemma~\ref{asm-proviso-window} with $T=T_u$, $d^*=d^*_u$ and
$\Gamma=\Gamma_P$, together with Step 3, gives
\[
 \frac{D^*}{2\Gamma^2}\le\dim M_u\le e\,\frac{D^*}{\Gamma^2},
\]
that is, $1/2\le\dim M_u\cdot\Gamma^2/D^{*}\le e$. Since $e\ge2$ and $C_0\ge e$,
$C_0^{-1}\le e^{-1}\le1/2$, and the proviso with $C_0$ follows.
\end{proof}

\begin{theorem}\label{asm-design-exists}
Let $\mathcal A$ be a class of families with a \emph{twist selection}, which
chooses a twist record (Definition~\ref{asm-definite-side-hypotheses}) at every
member of every family of $\mathcal A$ at every scale $x>x_1$ beyond a threshold
of the family, and in which every
family carries its duty data (Definition~\ref{asm-duty}). Let $c_T$ be a
function of two real variables. Write $\mathcal A_{N,\mathrm{tw}}$ for
$(\mathcal A_N)_{\mathrm{tw}}$ (Definition~\ref{asm-designed}). Suppose that:
\begin{enumerate}
\item if every finite subgroup of $GL_n(\mathbf Q)$ has order dividing
  Minkowski's number $M(n)$, for every $n$, then design choices exist in the
  class of families of $\mathcal A$ that satisfy the twist condition, with
  constant $3$ and function $c_T$;
\item eventually along $\mathcal A_{N,\mathrm{tw}}$, the size condition on
  $T_u$ holds whenever $c_T(\varepsilon,C_{YL})\le C_{T_u}$;
\item eventually along $\mathcal A_{N,\mathrm{tw}}$, the class number formula
  of Proposition~\ref{cv-class-number-formula} holds at every member at every
  scale $x>x_1$;
\item the zero-density and zero-free-region estimates of
  Chapter~\ref{chapter-host-fields} hold.
\end{enumerate}
Then designed families exist in $\mathcal A_{N,\mathrm{tw}}$ with constant
$3$. In particular, for every small $\varepsilon$ and every largeness
threshold there is such a family with chosen constants $(3,C_H)$,
$b^{m/2}\le T_u$, $\log T_u\le C_{T_u}mL$, and
$\log D_\Gamma(u)\le12fN_m\log(fN_m)$.
\end{theorem}
\begin{proof}
\emph{Step 1.} By Minkowski's theorem \cite{Min87}, every finite subgroup of
$GL_n(\mathbf Q)$ has order dividing $M(n)$, for every $n$. By (1), design
choices exist in the class $\mathcal A_{\mathrm{tw}}$ of families of
$\mathcal A$ that satisfy the twist condition
(Definition~\ref{asm-twist-condition}), with constant $3$ and function $c_T$.

\emph{Step 2.} By definition, a design choice satisfies the conditions that
define $\mathcal C_N$ (Definition~\ref{asm-designed}(5)), so it lies in
$(\mathcal A_{\mathrm{tw}})_N$. In the product defining $M(2fN_m)$ the
exponent of $p$ is $0$ unless $p-1\le2fN_m$, so every prime factor of
$M(2fN_m)$ is at most $2fN_m+1$, and so is every prime factor of a stabilizer
denominator $D_\Gamma(u)$ that divides $M(2fN_m)$. Hence the design choices of
Step 1 lie in $\mathcal A_{N,\mathrm{tw}}$, with the same thresholds, and
design choices exist in $\mathcal A_{N,\mathrm{tw}}$.

\emph{Step 3: the proviso.} By Lemma~\ref{asm-aligned-second-layer}, every
family of $\mathcal A$ is a linked and aligned family of the second layer, so
at all large scales its members carry tied data; this holds in particular for
the families of $\mathcal A_{N,\mathrm{tw}}$, which lie in $\mathcal A$. By
(3), there is $\varepsilon_3>0$ and, for every
$\varepsilon\in(0,\varepsilon_3]$, a largeness threshold such that every large
family of $\mathcal A_{N,\mathrm{tw}}$ with exponent $\varepsilon$ satisfies
the class number formula at every member at every scale $x>x_1$. Since
$e\le3$, Proposition~\ref{asm-proviso-reduction} with $C_0=3$ gives such a
family the proviso $3^{-1}\le\dim M_u\cdot\Gamma^2/D^{*}\le3$ for all large
$x$. So the proviso holds eventually along $\mathcal A_{N,\mathrm{tw}}$, with
the thresholds of (3).
% Lean: hRec is layerTwoFamilyClass_of_linked_aligned with linkedFamiliesAligned_of_duty;
% hProv is hn4.and (Pk.Closed.sub second_chern_linked_closed_rank_ties, through
% SecondLayerFamilyClass.c2s), then HarishChandraReduction.chosenProviso_min with exp 1 <= 3.

\emph{Step 4: the thin-scale bound.} By the estimates (4) and
Lemma~\ref{hf-host-chebyshev}, every family that satisfies, for all large $x$,
the box condition and the row condition of Definition~\ref{asm-box-row-tame}
at some $Y\in[x^{1/J},x]$ has a constant $c_6>0$ with
$|g_\chi(y)|\le y\exp(-c_6x^{\varepsilon/8})$ for every character $\chi$ of
the host field, every $y\ge X'/3$ and all large $x$.

\emph{Step 5: chosen constants.} Property (2) and Step 3 hold eventually along
$\mathcal A_{N,\mathrm{tw}}$; let $\varepsilon_1$ be a threshold for both, and
$\varepsilon_0$ that of the design choices. Let
$0<\varepsilon\le\min(\varepsilon_0,\varepsilon_1)$ and let $\kappa$ be a
largeness threshold. Let $\kappa'$ be the largeness threshold of (2) and
Step 3 at $\varepsilon$, and take a design choice $\mathcal F$, with constants
$C_H\ge0$ and $C_5,C_b>0$, that is large for the larger of the lower bounds of
$\kappa$ and $\kappa'$. Then $\mathcal F$ is large for $\kappa$ and for
$\kappa'$. Since $c_T(\varepsilon,C_{YL})\le C_{T_u}$, property (2) gives the
size condition on $T_u$. Step 3 gives the proviso. The design choice satisfies
the box and row conditions, so Step 4 gives $c_6>0$ and the thin-scale bound.
With the other conditions of chosen constants, which hold for a design
choice with $C_5$ and $C_b$, the family has chosen constants $(3,C_H)$.

\emph{Step 6: stabilizer denominators.} $M(n)$ divides $(2n)!$, so
$\log M(n)\le2n\log(2n)$. With $n=2y$ and $y=fN_m\ge2$ this gives
$\log M(2fN_m)\le4y\log(4y)\le12y\log y$, since $\log4\le2\log y$. As
$D_\Gamma(u)$ divides $M(2fN_m)$, we get $\log D_\Gamma(u)\le12fN_m\log(fN_m)$, the
size condition on stabilizer denominators with constant $12$. So
$\mathcal F$ has the properties required of a designed family.
\end{proof}
\leanref{Development/AutomorphicArithmetic/ThetaCorrespondence/ThetaPairLifts/ThirdPlaceSeries/LinkEx.lean}{GlobalTheory.Top.designExists_of_tuSize_aux4n_litThin}
\leanref{Formalism/SharpRestriction/DesignExistsOfChoices.lean}{Formalism.MainNotions.designExists_of_choices}
\leanref{Formalism/PeriodIntegrals/OrbitalBoundsConclusion/MinkowskiQ.lean}{minkowski_bound_rat}

\section{Definite-side hypotheses and tied families}\label{section-asm-definite-side}

\subsection{Hypotheses and consequences of the ties}\label{subsection-asm-hypotheses-ties}

\begin{definition}\label{asm-host-requirement}
A \emph{host requirement} is a property $G(s,\Phi,H,T,t)$ of a scale $s$
(Definition~\ref{hf-windows}), a Galois datum $\Phi$
(Definition~\ref{hf-frame}), a host field $H$ on $\Phi$
(Definition~\ref{hf-framed-host-field}), a finite subset $T$ of the group of
$\Phi$ and an integer $t$. For a member $D$ of a family, $G$ \emph{holds at
$D$} if it holds for the scale of $D$, its Galois datum, its host field, its set of top classes and its twist parameter $t$.
\begin{enumerate}
\item For $\lambda>1$, the class $\mathcal C_G(\lambda)$ consists of the
  families with windows of exponent $\lambda$ such that $G$ holds at every
  member at every scale $x>x_1$.
\item \emph{Hosts satisfying $G$ exist on windows of exponent $\lambda$} if
  hosts satisfying the condition
  $P(C_5,s,R,H,q_0)=G(s,\Phi_R,H,T_{\mathrm{top}},t_R)$ exist on windows of
  exponent $\lambda$ in the sense of Definition~\ref{hf-hosts-on-windows-def},
  where $\Phi_R$ is the Galois datum of the rank datum $R$, $t_R$ their last integer, and
  $T_{\mathrm{top}}$ the top classes of the host $H$.
\item For a primitive Dirichlet character $\psi$, the \emph{host requirement
  $G_\psi$ of $\psi$} holds at $(s,\Phi,H,T,t)$, with $s=(x,\varepsilon,J)$, if
  the host field $H$ satisfies the host conditions for the cyclic level
  $K_\psi$ (Definition~\ref{wl-cyclic-level}) at $(x,\varepsilon,4)$ with top
  classes $T$ (Definition~\ref{hf-host-conditions}), and its field $E$
  satisfies the Tate-Shafarevich bound $(t,e^{x^{\varepsilon}/2})$ of
  Definition~\ref{sha-bounded-field}.
\end{enumerate}
\end{definition}
\leanref{Development/AutomorphicArithmetic/CharacterBoxes/Supply.lean}{UpperTheory.Top.FReq}
\leanref{Development/AutomorphicArithmetic/CharacterBoxes/Supply.lean}{UpperTheory.Top.FReq.cls}
\leanref{Development/AbelianVarietyTheory/Definitions/GlobalTwisting/ThinScalesIffHostBound.lean}{GlobalTheory.EX.HostExW}
\leanref{Development/AutomorphicArithmetic/CharacterBoxes/Supply.lean}{UpperTheory.Top.GChi}

\begin{definition}\label{asm-tied-everywhere}
\emph{Tied families exist for all host requirements} if the following holds
for every $\lambda>1$ and every host requirement $G$
(Definition~\ref{asm-host-requirement}). Suppose that hosts satisfying $G$
exist on windows of exponent $\lambda$, and that every family in
$\mathcal C_G(\lambda)$ satisfies the Tate-Shafarevich bound
(Definition~\ref{asm-sha-bound}). Then $\mathcal C_G(\lambda)$ has tied
members (Definition~\ref{asm-tied}): there are $e_D,e_L$ and
$\varepsilon_1>0$ such that for every $\varepsilon\in(0,\varepsilon_1)$ some
family in $\mathcal C_G(\lambda)$ with exponent $\varepsilon$ is tied at all
large scales of its set of scales.
\end{definition}
\leanref{Development/AutomorphicArithmetic/CharacterBoxes/ZeroFreeAllOfTopG.lean}{UpperTheory.Top.TopG}

\begin{definition}\label{asm-definite-side-hypotheses}
A \emph{collection of definite-side hypotheses} consists of a class
$\mathcal A$ of families, a raw selection for $\mathcal A$, a function $c_T$ of
two real variables, counting constants attached to finite lists of reals, and
a real number $c_I>0$, with the following properties. The \emph{raw selection}
chooses, at every member $D$ of every family in $\mathcal A$ at every scale
$x>x_1$, with host field $E$ and $K=\mathbf Q(\sqrt{-d})$
(Definition~\ref{hf-host}):
\begin{itemize}
\item a \emph{twist record} for $D$, which consists of: \emph{tied data}, which
  tie to one another the host and level structure of $D$, its columns, its theta
  data, its twisting data (Definition~\ref{th-twisting-data}), its characters,
  its rows and its group coherence, give kernel entries of $D$
  (Definition~\ref{pd-kernel-entries}), and give for every word $u$ and every
  level $i\le m$ a matrix $A_{u,i}\in GL_{N(i)}(E)$ such that
  $A_{u,i}^*A_{u,i}$ is the Gram matrix of level $i$ and
  $A_{u,i}^{-1}\mathcal O_E^{N(i)}$ is the self-dual lattice of level $i$; a
  \emph{row pin}, which ties the lower maps of the columns and of the twisted
  columns to the tied data; \emph{joints}, which tie the determinant
  characters, the uniformizers and characters at the level primes, and the
  divisor characters to the tied data; \emph{class data}: representatives
  of the ideal classes of $K$ prime to the auxiliary primes $T_u$ and to the
  places in the sets $Q$ of the twist choices, and a Hilbert class field of
  $E$ inside the field of the rows; and the equivariance of the Weil
  representation of the theta data under the Heisenberg and unitary groups,
  with the integral formula for its theta lifts;
\item a \emph{spectral condition} for the twist record: an automorphic setting,
  with a standard additive character, tied to $D$ in its automorphy, its
  measures and its local data, in which the columns of every word of $D$ are
  realized as Poincaré series, with dictionaries that identify the columns and
  the rows of $D$ with those of the setting;
\item a row-side family $P$ for $D$ (Definition~\ref{is-row-side-family}), the
  \emph{selected base level}, whose row side is determined by the twist record,
  with rows $S_0$, columns $S_1$, definite columns $D_c$, identifications
  $e_{N,u}$, $j_W(i)$ and $u\mapsto w_u$, and tuples $t$ with characters
  $\mu_t$;
\item a \emph{definite height}: integers $h_1,h_2$ with $\Delta_S=h_1$
  (Definition~\ref{th-definite-datum}(8)) and $\eta=mh_2$ for the width $\eta$
  of the size data;
\item \emph{square-integrable columns} for $P$, as in the bundles of
  Definition~\ref{is-definite-side-condition}: a column record over a field of
  columns embedded in $L$, with periods on every definite subspace of $P$, from
  which the group, the places, the inert indices, the tuples and the family of
  definite forms of $P$ are built, and which is tied to the columns of $P$, to
  the words of $D$ and to the kernel entries of the twist record; the family
  of definite forms of $P$ satisfies the family conditions, with the conditions
  on its inert indices (Lemma~\ref{is-eigenvalue-base-inputs}(1));
\item a \emph{selection bundle}: an indefinite realization $Q$ of $P$ for the
  automorphic setting of the spectral condition
  (Definition~\ref{is-realisation-join}), with realization $R_Q$ and face
  $F_Q$ (Definition~\ref{is-realisation}), together with a core of $D$, weight
  families $\kappa_i$, vectors $v_0(i,\sigma)$ and one-map data for $R_Q$
  (Definition~\ref{is-one-map-data}).
\end{itemize}
The twist records chosen by the raw selection form a twist selection for
$\mathcal A$ (Theorem~\ref{asm-design-exists}), with threshold $x_1$. Write
$\mathcal A_{N,\mathrm{tw}}=(\mathcal A_N)_{\mathrm{tw}}$ and
$\mathcal D_3=\operatorname{Des}(\mathcal A_N,3)$
(Definition~\ref{asm-designed}). The rows data of $D$ have a finite set of
\emph{types} $\sigma$; each of the $D_0$ base rows of $D$ has a type, and
$n_0(\sigma)$ is the number of base rows of type $\sigma$; and for every level
$j$ and type $\sigma$ there are an \emph{isotypic part} $I_j(\sigma)$ of the
lifted space of $D$ at level $j$ and, for every tuple $\mu$, a real
\emph{parameter} $\lambda_D(\mu,j,\sigma)$.
\begin{enumerate}
\item (Design.) If every finite subgroup of $GL_n(\mathbf Q)$ has order
  dividing $M(n)$ for every $n$, then for every $\lambda>1$ and every host
  requirement $G$ such that hosts satisfying $G$ exist on windows of
  exponent $\lambda$, design choices exist in the class of families of
  $\mathcal A\cap\mathcal C_G(\lambda)$ that satisfy the twist condition,
  with constant $3$ and function $c_T$ (Definition~\ref{asm-designed}).
\item (Duty.) Every family in $\mathcal A$ carries its duty data
  (Definition~\ref{asm-duty}).
\item (Selected rows.) At every member $D$ at a scale $x>x_1$, with the
  selected base level $P$ and the selection bundle:
  \begin{enumerate}
  \item $R_Q$ has bounded rows (Definition~\ref{is-bounded-rows}) for its
    one-map data, with the phases of $R_Q$;
  \item the twists are recorded in the size data: every twisted place of the
    rows $S_0$ is a place $w$ of the support of the size data with
    $\mu_w\ne\mu_0$;
  \item $R_Q$ has primitive rows at $F_Q$ (Definition~\ref{is-primitive-rows}),
    and the family of definite forms of $P$ has lift-distinct types
    (Definition~\ref{is-realisation}): for every lift $j$, the transfer at $j$
    sends distinct types of the base of $P$ to types that differ at some place
    or at some inert index of the lift forms at $(j,\mu_0)$;
  \item the rows hold with $\Lambda_3$: there is a map $\sigma\mapsto c(\sigma)$
    from the types of $D$ to the types $\Pi$ of the base of $P$ such that, for
    every lift $0\le i<m$ of $S_0$ and every $c\in\Pi$, the isotypic parts
    $I_{i+1}(\sigma)$ with $c(\sigma)=c$ are independent and their sum is the
    complex span of $j_W(i)(W_c(i))$, for the part $W_c(i)$ of $S_0$
    (Definition~\ref{is-indefinite-datum}); for every lift $i$, every tuple $t$
    of $P$ with $\mu_t$ used (Definition~\ref{pd-period-data}(6)) and every type
    $\sigma$ with $I_{i+1}(\sigma)\ne0$, the parameter
    $\lambda_D(\mu_t,i+1,\sigma)$ is the scalar $\lambda_{i,t,c(\sigma),v_0}$ of
    $S_0$; and for every $c\in\Pi$ the numbers $n_0(\sigma)$ with $c(\sigma)=c$
    add up to $\dim_LE_c$, for the part $E_c$ of $S_0$;
  \item the division rows hold for $P$ (Definition~\ref{is-division-rows});
  \item the division data are compatible with the periods: for every word $u$,
    every tuple $t$ of $P$ with $\mu_t$ used and all coordinate vectors $x$
    and $F$, the image under $L\to\mathbf C$ of the period of the member $w_u$
    of $D_c$ at $(t,x,F)$ is the period
    $\operatorname{per}_{u,\mu_t}(j_Ne_{N,u}x,F')$ of
    Definition~\ref{th-definite-datum}(5), where $F'$ is the top form for
    $\mu_t$ with coordinates $F$.
  \end{enumerate}
\item (Sizes.) Eventually along $\mathcal A_{N,\mathrm{tw}}$, the size
  condition on $T_u$ holds whenever $c_T(\varepsilon,C_{YL})\le C_{T_u}$, and
  the class number formula of Proposition~\ref{cv-class-number-formula} holds at
  every member.
\item (Analytic input.) The zero-density and zero-free-region estimates of
  Theorem~\ref{hf-zero-density-regions} hold, and so do the explicit formula
  and the bound for zeros in windows of Theorem~\ref{asm-dependence-zero-sums}.
\item (Orthogonality inputs.) The decay of Wu type for gaps
  (Definition~\ref{is-wu-decay}), the finite
  dimensionality of rows, the frame condition eventually along
  $\mathcal D_3$, the counting constants with their point counts, the
  positivity of row types, the row non-degeneracy, the row main-term
  identity, the loop rows and the selected form of the orthogonality theorem,
  as in Theorem~\ref{is-orthogonality-link-holds}.
\item (Volumes.) At every member $D$ at a scale $x>x_1$, with host
  field $E$ and totally real subfield $F$:
  \begin{enumerate}
  \item the volume is independent of the signature: if $W$ and $W_0$ are
    Hermitian spaces over $E$ of the same rank, with signature $(P,Q)$ at every
    complex place for $W$ and $(P_0,Q_0)$ for $W_0$, then fundamental domains for
    $U(W)(F)$ in $U(W)(\mathbf A_F)$ and for $U(W_0)(F)$ in
    $U(W_0)(\mathbf A_F)$ have the same volume for the volume measures of $D$;
  \item the class items hold in the rank $n_0$ of the space $Y_0$ of $D$: for
    every Hermitian space $W$ over $E$ of rank $n_0$, positive definite at every
    place, the mass of its standard lattice equals the mass of $I_{n_0}$
    (Definition~\ref{cv-mass}); $U(W)(\mathbf A_F)$ is the disjoint union of
    finitely many double cosets $U(W)(F)g_iK$, where $K$ is the stabilizer of
    the standard lattice in the finite adelic group; and for any such
    representatives $\sum_i|U(W)(F)\cap g_iKg_i^{-1}|^{-1}$ is the mass of the
    standard lattice.
  \end{enumerate}
\item (Joins.) At every member at a scale $x>x_1$, the indefinite realization
  $Q$ of the selection bundle is joined (Definition~\ref{is-realisation-join})
  for the automorphic setting of the spectral condition. Eventually along
  $\mathcal D_3$, for all large $x$, the base
  vector image is full: the map $\iota_{\mu_0}$ from the ambient space of the
  base tuple to the top forms $\mathcal E$ carries the base space
  $V_0(\mu_0)$ onto $\mathcal E$. For every family in $\mathcal A$ and all
  large $x$, every twisted lattice $\tilde M_\mu$ of every member is stable
  under multiplication by $\mathcal O_L$, for the coefficient field $L$.
\item (Definite rows.) The definite-side condition
  (Definition~\ref{is-definite-side-condition}) holds for the selected
  bundles, eventually along $\mathcal D_3$, in the form of
  Theorem~\ref{is-definite-rows}.
\item (Integrality inputs.) Eventually along $\mathcal D_3$, hypotheses (1)
  to (6) of Theorem~\ref{pd-multiplicity-from-integrality} hold: point
  integrality, twisted point integrality, period integrality
  (Definition~\ref{pd-period-integrality}), twisted period integrality
  (Definition~\ref{pd-twisted-period-integrality}), the sharp column bounds at
  the auxiliary prime and the bounds at the level primes; the second
  complement inputs hold (Definition~\ref{asm-division-kernel-errors}(2));
  and the twist selection is served
  (Definition~\ref{pd-served-selection}), which is hypothesis (7).
\item (Irreducibility inputs.) Eventually along $\mathcal D_3$, every member
  at every large scale has a labelling of all rows that is injective on good
  labels, joins all rows, and has the row count with constant $c_I$
  (Definition~\ref{is-row-labelling}), as in
  Theorem~\ref{is-irreducibility-link-holds}.
\item (Projections.) Eventually along $\mathcal D_3$, at every member at every
  large scale, the projection matrices of the selected base level $P$ at the
  base tuple to the lifts are surjective: for every lift $i$, the map $\pi$ of
  Lemma~\ref{asm-dependence-projections-onto} for the lift forms at
  $(i,\mu_0)$ of the family of $P$, the characters of that family, the subgroup
  $P_0$ of $P$ with the character $\psi_0^{-1}$, and the definite subspace of $P$
  at $(i,\mu_0)$.
\item (Real place and masses.) The second-layer conditions at the real place
  hold eventually along $\mathcal D_3$; the mass ratio under volume holds for
  every definite-side datum; and the mass formula for genera of rank $2$ holds
  for every definite-side datum.
\item (Tate-Shafarevich link.) The Tate-Shafarevich link
  (Definition~\ref{asm-twisted-links}) holds eventually along $\mathcal D_3$.
\end{enumerate}
\end{definition}
\leanref{Development/AutomorphicArithmetic/ThetaCorrespondence/ThetaPairLifts/ThirdPlaceSeries/Hyps.lean}{GlobalTheory.Fc.HypsAlt}
\leanref{Development/AutomorphicArithmetic/ThetaCorrespondence/ThetaPairLifts/ThirdPlaceSeries/DefiniteSideCondition.lean}{GlobalTheory.Fc.DefiniteSideCondition}
\leanref{Development/AbelianVarietyTheory/Definitions/GlobalTwisting/GeometricWeakTiesOfStrong.lean}{GlobalTheory.Fc.S0.NL}
\leanref{Development/AutomorphicArithmetic/ThetaCorrespondence/ThetaPairLifts/ThirdPlaceSeries/LayerTwoFamilyClass.lean}{Formalism.L2.RecTAlt}
\leanref{Development/AutomorphicArithmetic/ThetaCorrespondence/ThetaPairLifts/ThirdPlaceSeries/STiesR.lean}{Formalism.L2.SG.SpectralTiesReal.STiesRAlt}
\leanref{Development/AutomorphicArithmetic/ThetaCorrespondence/ThetaPairLifts/ThirdPlaceSeries/SquareIntegrableColumnsTied.lean}{GlobalTheory.Cp.SquareIntegrableColumns}
\leanref{Formalism/LocalToGlobal/Divisibility/InertRationalAdjoint.lean}{Formalism.L2.DIV.DefFam3.TypesSep3}

\begin{definition}\label{asm-support-local-modules}
Let $D$ be a definite-side datum (Definition~\ref{th-definite-datum}) with host field $E$, totally
real subfield $F$, coefficient field $L$ with its embedding $\iota\colon L\to\mathbf C$, rank $m$,
base tuple $\mu_0$ and tuples $\mu_w$ of the finite places $w$ of $L$
(Definition~\ref{pd-period-data}). For a finite place $w$ of $L$ let $\mu^{\mathrm{tw}}_w$ be its
twisting tuple, so that $\mu_w=\mu^{\mathrm{tw}}_w$ when $\mu_w\ne\mu_0$. For a tuple $\mu$ and a
level $1\le j\le m$, the rows data of $D$ give the \emph{space of forms} $\mathcal M_{\mu,j}$ of $D$
for $\mu$ at level $j$, a finite-dimensional complex Hilbert space, together with an $L$-vector space
$\mathcal M^L_{\mu,j}$, an $\iota$-semilinear map $\mathcal M^L_{\mu,j}\to\mathcal M_{\mu,j}$ that
carries a basis to a basis, and a lattice in $\mathcal M_{\mu,j}$. The \emph{local modules} of $D$
consist of a finite set of \emph{local indices} $i$, each with the following data.
\begin{enumerate}
\item A prime $\ell_i$, a level $1\le j(i)\le m$ and a finite place $w(i)$ of $L$ over $\ell_i$ with
  $\mu_{w(i)}\ne\mu_0$.
\item A complete discrete valuation ring $\mathcal O_i$, finite over $\mathbf Z_{\ell_i}$, with
  fraction field $K_i$, and an embedding $e_i\colon L\to K_i$ such that $e_i(x)\in\mathcal O_i$ if
  and only if $v_{w(i)}(x)\ge0$.
\item A finite nonempty set $X_i$ and a local $\mathcal O_i$-algebra $\mathbf T_i$ of
  $\mathcal O_i$-valued functions on $X_i$ that separates the points of $X_i$, with \emph{Hecke
  elements} $h_i(v,k)\in\mathbf T_i$ for the finite places $v$ of $E$ and the integers $k$.
\item A free $\mathcal O_i$-module $\tilde M_i$ of finite rank with a faithful
  $\mathcal O_i$-linear action of $\mathbf T_i$, and a finite set $Q_i$ of finite places of $E$.
\item An $e_i$-semilinear map
  $\mathrm{pr}_i\colon\mathcal M^L_{\mu_{w(i)},j(i)}\to K_i\otimes_{\mathcal O_i}\tilde M_i$ whose
  image spans over $K_i$, and which carries the vectors whose image in
  $\mathcal M_{\mu_{w(i)},j(i)}$ lies in the lattice to a set that spans $1\otimes\tilde M_i$ over
  $\mathcal O_i$.
\end{enumerate}
The \emph{local indices at} $(j,w)$ are those with $j(i)=j$ and $w(i)=w$, and the modules
$\tilde M_i$ for these $i$ are the local spaces at $w$ for the lift $j$. The support condition of
Definition~\ref{is-division-link} holds at $D$ if, for every lift $1\le j\le m$ and every finite
place $w$ of $L$ with $\mu_w\ne\mu_0$,
\[
 \sum_i\operatorname{rank}_{\mathcal O_i}\tilde M_i=\dim_L\mathcal M^L_{\mu_w,j},
\]
the sum over the local indices $i$ at $(j,w)$.
\end{definition}
\leanref{Formalism/SharpRestriction/LevelLetterCocycles/DataN.lean}{Formalism.BasicNotions.NData}
\leanref{Formalism/SharpRestriction/LevelLetterCocycles/LocsCoverTw.lean}{Formalism.MainNotions.N.LocsCoverTw}

\begin{definition}\label{asm-support-hecke-realization}
Let $E$ be a CM field with totally real subfield $F$, and let $n\ge1$. Let $U_n$ be the unitary group
over $F$ of the Hermitian form $\sum_ix_i\bar y_i$ on $E^n$, with rational points $U_n(F)$, finite
adelic points $U_n(\mathbf A_{F,f})$, and the subgroup $U_n(\widehat{\mathcal O})$ of its matrices
with integral entries.
\begin{enumerate}
\item An \emph{algebraic weight} for $U_n$ is an irreducible representation $\rho$ of $U_n(F)$ on a
  finite-dimensional complex vector space $W$ whose matrix coefficients, in some basis, are
  polynomials in the images of the matrix entries under the complex embeddings of $E$.
\item A \emph{level} is a subgroup $J\subset U_n(\mathbf A_{F,f})$ with a character
  $\chi\colon J\to\mathbf C^\times$ such that, for some $N\ge1$, every element of
  $U_n(\widehat{\mathcal O})$ congruent to $1$ modulo $N$ lies in $J$ and in the kernel of $\chi$.
  Its \emph{type space} for $\rho$ is the space of functions $f\colon U_n(\mathbf A_{F,f})\to W$
  with $f(\gamma g)=\rho(\gamma)f(g)$ for $\gamma\in U_n(F)$, invariant on the right under some
  principal congruence subgroup, and with $f(g\kappa)=\chi(\kappa)f(g)$ for $\kappa\in J$. For a set
  $B$ of finite places of $F$, the level is \emph{unramified outside $B$} if every
  $\kappa\in U_n(\mathbf A_{F,f})$ whose components at the places of $B$ are $1$ lies in $J$ if and
  only if it lies in $U_n(\widehat{\mathcal O})$, and then $\chi(\kappa)=1$.
\item A \emph{split place} $s$ of $E/F$ is a finite place of $F$ with an $F$-embedding of $E$ into
  the completion $F_s$. It gives a place $s_E$ of $E$ over $s$, and an identification of the
  component of $U_n(\mathbf A_{F,f})$ at $s$ with $GL_n(F_s)$. For a uniformizer $\varpi$ of $F_s$
  and $1\le k\le n$, let $T_{s,\varpi,k}$ be the set of elements of $U_n(\mathbf A_{F,f})$ with
  component $1$ away from $s$ and component at $s$ in
  $GL_n(\mathcal O_s)\operatorname{diag}(\varpi,\dots,\varpi,1,\dots,1)GL_n(\mathcal O_s)$, with
  $k$ entries $\varpi$. The \emph{normalized Hecke operator} $t_{s,\varpi,k}$ of the level $J$
  sends a function $f$ to
  \[
   (t_{s,\varpi,k}f)(g)=q_s^{-k(n-1)}\sum_xf(gx),
  \]
  the sum over one element $x$ of each coset $xJ$ that meets $T_{s,\varpi,k}$, where $q_s$ is the
  cardinality of the residue field of $F_s$.
\item Let $D$ be a definite-side datum with host field $E$ and ranks $m$, $r$ and $t$, with the
  spaces of forms $\mathcal M_{\mu,j}$ of Definition~\ref{asm-support-local-modules}, and put
  $n_V=r+t$. A \emph{Hecke realization} of $D$ assigns to every tuple $\mu$ and level
  $1\le j\le m$ an algebraic weight $\rho$ for $U_{n_V}$ on a space $W$, a level $(J,\chi)$, and an
  injective linear map from $\mathcal M_{\mu,j}$ to the functions
  $U_{n_V}(\mathbf A_{F,f})\to W$ whose image is the type space of $(J,\chi)$ for $\rho$.
\end{enumerate}
\end{definition}
\leanref{Formalism/AlgebraicWeights/IsAlgWeight.lean}{Formalism.AUTREP.IsAlgWeight}
\leanref{Formalism/AlgebraicWeights/IsAlgWeight.lean}{Formalism.AUTREP.tSet}
\leanref{Formalism/LocalToGlobal/GaloisLevelStructure/DefsAut.lean}{Formalism.L2.NG.TypeSp}
\leanref{Formalism/LocalToGlobal/GaloisLevelStructure/DefsAut.lean}{Formalism.L2.NG.HypOutside}
\leanref{Formalism/LocalToGlobal/GaloisLevelStructure/DefsAut.lean}{Formalism.L2.NG.tN}
\leanref{Formalism/LocalToGlobal/GaloisLevelStructure/TiedHecke.lean}{Formalism.L2.MtReal}
\leanref{Formalism/LocalToGlobal/GaloisLevelStructure/TiedHecke.lean}{Formalism.L2.MtFam}

\begin{definition}\label{asm-support-local-compatibility}
Let $D$ be a definite-side datum with local modules (Definition~\ref{asm-support-local-modules}) and
a Hecke realization (Definition~\ref{asm-support-hecke-realization}), with host field $E$, totally
real subfield $F$, coefficient field $L$, rank $m$, level primes $g_0,g_1$, further level primes
$p_{de},p'_{de}$, and $n_V=r+t$. For a prime $\ell$ put $p^*(\ell)=p'_{de}$ if $\ell=p_{de}$ and
$p^*(\ell)=p_{de}$ otherwise. For a prime $\ell$ and a finite set $Q$ of finite places of $E$, a pair
$(v,k)$ of a finite place $v$ of $E$ and an integer $k$ is \emph{generic} for $(\ell,Q)$ if $v$ is
split over $F$, $v\notin Q$, $v$ lies over none of $\ell$, $g_0$, $g_1$ and $p^*(\ell)$, and
$1\le k\le n_V$. For a local index $i$, let $B_i$ be the set of finite places of $F$ below a place of
$E$ that lies in $Q_i$ or over one of $g_0$, $g_1$ and $p^*(\ell_i)$.

\emph{Normalized local compatibility} of $D$ with its Hecke realization consists of the following
data, for every level $1\le j\le m$ and every finite place $w$ of $L$, with $\mu=\mu^{\mathrm{tw}}_w$
and $\ell=p_w$.
\begin{enumerate}
\item A discrete valuation ring $\mathcal O_{j,w}$ with fraction field $K_{j,w}$, and an embedding
  $e_{j,w}\colon L\to K_{j,w}$ with $e_{j,w}(x)\in\mathcal O_{j,w}$ if and only if $v_w(x)\ge0$.
\item A $K_{j,w}$-vector space $V_{j,w}$ with an $e_{j,w}$-semilinear map
  $\iota_{j,w}\colon\mathcal M^L_{\mu,j}\to V_{j,w}$ that carries every basis to a basis.
\item A finite set $Q_{j,w}$ of finite places of $E$, stable under complex conjugation, that
  contains with a place every place of $E$ over the same place of $F$.
\item For every finite place $v$ of $E$ and integer $k$, an $L$-linear operator $h(v,k)$ of
  $\mathcal M^L_{\mu,j}$ and a $K_{j,w}$-linear operator $h_K(v,k)$ of $V_{j,w}$ with
  $h_K(v,k)\circ\iota_{j,w}=\iota_{j,w}\circ h(v,k)$, such that, whenever $(v,k)$ is generic for
  $(\ell,Q_{j,w})$, the Hecke realization of $\mathcal M_{\mu,j}$ carries $h(v,k)$ to the normalized
  Hecke operator $t_{s,\varpi,k}$ of its level for every split place $s$ with $s_E=\bar v$ and every
  uniformizer $\varpi$ of $F_s$, and $V_{j,w}$ is the sum of the generalized eigenspaces of
  $h_K(v,k)$.
\end{enumerate}
Let $\mathbf A_{j,w}$ be the $\mathcal O_{j,w}$-subalgebra of the endomorphisms of $V_{j,w}$
generated by the $h_K(v,k)$ with $(v,k)$ generic for $(\ell,Q_{j,w})$. For an
$\mathcal O_{j,w}$-algebra homomorphism $\psi\colon\mathbf A_{j,w}\to\mathcal O_{j,w}$, let
$V(\psi)$ be the intersection over $a\in\mathbf A_{j,w}$ of the generalized eigenspaces of $a$ for
the eigenvalues $\psi(a)$, and $\mathfrak m_\psi$ the kernel of $\psi$ followed by the reduction to
the residue field. For an ideal $\mathfrak m$ of $\mathbf A_{j,w}$, the \emph{part}
$V_{j,w}[\mathfrak m]$ is the sum of the $V(\psi)$ with $\mathfrak m_\psi=\mathfrak m$.
\begin{enumerate}\setcounter{enumi}{4}
\item For every local index $i$ at $(j,w)$: $Q_i=Q_{j,w}$; a maximal ideal $\mathfrak m(i)$ of
  $\mathbf A_{j,w}$; a ring isomorphism $\mathcal O_i\cong\mathcal O_{j,w}$ whose extension
  $K_i\cong K_{j,w}$ carries $e_i$ to $e_{j,w}$; and an injective map
  $\Psi_i\colon K_i\otimes_{\mathcal O_i}\tilde M_i\to V_{j,w}$, semilinear for this isomorphism,
  with image $V_{j,w}[\mathfrak m(i)]$.
\item The map $i\mapsto\mathfrak m(i)$ is injective on the local indices at $(j,w)$, and if
  $\mu_w\ne\mu_0$, every maximal ideal of $\mathbf A_{j,w}$ is $\mathfrak m(i)$ for some local index
  $i$ at $(j,w)$.
\end{enumerate}
The compatibility also identifies $\Psi_i\circ\mathrm{pr}_i$ with $\iota_{j,w}$ modulo the
$V(\psi)$ with $\mathfrak m_\psi\ne\mathfrak m(i)$, the points of $X_i$ with the homomorphisms $\psi$
with $\mathfrak m_\psi=\mathfrak m(i)$ and $V(\psi)\ne0$, the values of the Hecke elements
$h_i(v,k)$ at these points with the values of these $\psi$ at $h_K(v,k)$, and the complex embeddings
of the rings $\mathcal O_i$ with one of $\mathcal O_{j,w}$.

The \emph{Hecke ties} of $D$ with its Hecke realization contain the condition that, for every local
index $i$, the level of the Hecke realization of $\mathcal M_{\mu_{w(i)},j(i)}$ is unramified
outside $B_i$. Their other conditions tie the Hecke elements, the rings $\mathcal O_i$ and the sets
$Q_i$ of different local indices to one another, and are not used below. The tied data of a twist
record (Definition~\ref{asm-definite-side-hypotheses}) contain a Hecke realization of the member
with normalized local compatibility and the Hecke ties.
\end{definition}
\leanref{Formalism/LocalToGlobal/NormalizedHeckeFactors/HeckeStdTwOfStd.lean}{Formalism.L2.NLocBase}
\leanref{Formalism/LocalToGlobal/NormalizedHeckeFactors/HeckeStdTwOfStd.lean}{Formalism.L2.NormalizedLocalCompatibility}
\leanref{Formalism/LocalToGlobal/NormalizedHeckeFactors/HeckeStdTwOfStd.lean}{Formalism.L2.GenAt}
\leanref{Formalism/LocalToGlobal/NormalizedHeckeFactors/HeckeStdTwOfStd.lean}{Formalism.L2.heckeAlg}
\leanref{Formalism/LocalToGlobal/NormalizedHeckeFactors/CoverAlg.lean}{Formalism.L2.NCov.part}
\leanref{Formalism/LocalToGlobal/GaloisLevelStructure/TiedHecke.lean}{Formalism.L2.HeckeTied}
\leanref{Formalism/LocalToGlobal/GaloisLevelStructure/TiedHecke.lean}{Formalism.L2.badPl}
\leanref{Development/AutomorphicArithmetic/ThetaCorrespondence/ThetaPairLifts/ThirdPlaceSeries/RingSpecNormalisedCharacterEqAlt.lean}{Formalism.L2.NTiedAllAlt}

\begin{proposition}[Hecke operators at split places \cite{Gro99,CHT08}]\label{asm-support-hecke-standard}
Let $E$ be a CM field with totally real subfield $F$, let $n\ge1$, and let $\rho$ be an algebraic
weight and $(J,\chi)$ a level for $U_n$ (Definition~\ref{asm-support-hecke-realization}), with type
space $S$ for $\rho$. Let $B$ be a set of finite places of $F$ such that $(J,\chi)$ is unramified
outside $B$. Let $s$ and $s'$ be split places of $E/F$ not in $B$, let $\varpi$ and $\varpi'$ be
uniformizers of $F_s$ and $F_{s'}$, and let $1\le k,k'\le n$. Then:
\begin{enumerate}
\item $t_{s,\varpi,k}$ maps $S$ to itself and is linear on $S$;
\item $t_{s,\varpi,k}$ and $t_{s',\varpi',k'}$ commute on $S$;
\item if $f\in S$ is nonzero and $t_{s,\varpi,k}f=\lambda f$, then for every prime $\ell$ such that
  $s$ does not lie over $\ell$ there is an integer $d$ prime to $\ell$ such that $d\lambda$ is an
  algebraic integer.
\end{enumerate}
\end{proposition}
\leanref{CompanionTheory/AutomorphicTheory/HeckeStandardModules/HeckeStdStatement.lean}{ReciprocityTheory.ChartRestriction.L3H.HeckeStdStatement}
\leanref{CompanionTheory/AutomorphicTheory/HeckeStandardModules/InvMemRl.lean}{ReciprocityTheory.ChartRestriction.L3H.heckeStd}

\begin{lemma}\label{asm-support-eigenspaces}
Let $K$ be a field, $\mathcal O\subset K$ a local subring, $V$ a finite-dimensional $K$-vector space
and $S$ a set of pairwise commuting $K$-linear endomorphisms of $V$. Suppose that $V$ is the sum of
the generalized eigenspaces of each $a\in S$, and that every eigenvalue of every $a\in S$ lies in
$\mathcal O$. Let $\mathbf A$ be the $\mathcal O$-subalgebra of the endomorphisms of $V$ generated by
$S$. For an $\mathcal O$-algebra homomorphism $\psi\colon\mathbf A\to\mathcal O$, let $V(\psi)$ be the
intersection over $a\in\mathbf A$ of the generalized eigenspaces of $a$ for the eigenvalues
$\psi(a)$, and $\mathfrak m_\psi$ the kernel of $\psi$ followed by the reduction to the residue field
of $\mathcal O$; for an ideal $\mathfrak m$ of $\mathbf A$, let $V[\mathfrak m]$ be the sum of the
$V(\psi)$ with $\mathfrak m_\psi=\mathfrak m$. Then for every finite set $\mathcal I$ of ideals of
$\mathbf A$ that contains every maximal ideal,
\[
 \sum_{\mathfrak m\in\mathcal I}\dim_KV[\mathfrak m]=\dim_KV.
\]
\end{lemma}
\leanref{Formalism/LocalToGlobal/NormalizedHeckeFactors/CoverAlg.lean}{Formalism.L2.NCov.sum_finrank_part}
\begin{proof}
The elements of $S$ commute and each is split on $V$, so $V$ is the sum of the joint generalized
eigenspaces $\bigcap_{a\in S}V_a(\lambda_a)$, for families of eigenvalues $\lambda_a\in\mathcal O$;
as $V$ is finite-dimensional, finitely many elements of $S$ cut them out. On a nonzero joint
generalized eigenspace, a sum or a product of two commuting operators with single generalized
eigenvalues $\alpha$ and $\beta$ has the single generalized eigenvalue $\alpha+\beta$ or
$\alpha\beta$. So every element of $\mathbf A$ has a single generalized eigenvalue on it, in
$\mathcal O$, and these eigenvalues form a homomorphism $\psi$ with the space inside $V(\psi)$. Hence
the $V(\psi)$ span $V$, and finitely many of them are nonzero. Two distinct homomorphisms differ at
some element of $S$, which generates $\mathbf A$, and generalized eigenspaces of one operator for
distinct eigenvalues are independent; so the $V(\psi)$ are independent, and
$\dim V=\sum_\psi\dim V(\psi)$. Each $\mathfrak m_\psi$ is the kernel of a surjection onto the
residue field, so it is maximal and lies in $\mathcal I$. Grouping the $\psi$ by $\mathfrak m_\psi$
gives the formula.
\end{proof}

\begin{lemma}\label{asm-support-from-ties}
Let $D$ be a definite-side datum with local modules (Definition~\ref{asm-support-local-modules}) and
a Hecke realization (Definition~\ref{asm-support-hecke-realization}), with normalized local
compatibility and the Hecke ties (Definition~\ref{asm-support-local-compatibility}). Let $L$ be its
coefficient field, $m$ its rank and $\mu_0$ its base tuple. Then the support condition holds at $D$:
for every lift $1\le j\le m$ and every finite place $w$ of $L$ with $\mu_w\ne\mu_0$,
\[
 \sum_i\operatorname{rank}_{\mathcal O_i}\tilde M_i=\dim_L\mathcal M^L_{\mu_w,j},
\]
the sum over the local indices $i$ at $(j,w)$.
\end{lemma}
\leanref{Formalism/LocalToGlobal/NormalizedHeckeFactors/CoverThm.lean}{Formalism.L2.locsCoverTw_of_tie}
\leanref{Formalism/LocalToGlobal/NormalizedHeckeFactors/CoverThm.lean}{Formalism.L2.sum_finrank_eq_of}
\leanref{Formalism/LocalToGlobal/NormalizedHeckeFactors/CoverStdB.lean}{Formalism.L2.heckeStd_of_plain}
\leanref{Formalism/LocalToGlobal/NormalizedHeckeFactors/CoverStdB.lean}{Formalism.L2.genOutBad_of}
\leanref{Formalism/LocalToGlobal/NormalizedHeckeFactors/CoverSplit.lean}{Formalism.L2.splitEx_of_star}
\leanref{Formalism/LocalToGlobal/NormalizedHeckeFactors/CoverF.lean}{Formalism.L2.locsCoverTwF_of}
\begin{proof}
The numbered items are those of Definition~\ref{asm-support-local-compatibility}. Fix $j$ and $w$
with $\mu_w\ne\mu_0$, and write $\mu=\mu_w$, $\ell=p_w$, $\mathcal O=\mathcal O_{j,w}$, $K=K_{j,w}$,
$V=V_{j,w}$, $Q=Q_{j,w}$ and $\mathbf A=\mathbf A_{j,w}$. Call $(v,k)$ generic if it is generic for
$(\ell,Q)$, and let $R$ be the Hecke realization of $\mathcal M_{\mu,j}$, composed with the map
$\mathcal M^L_{\mu,j}\to\mathcal M_{\mu,j}$; it is injective.

If there is no local index at $(j,w)$, then $\mathbf A$ has no maximal ideal by (6), so $\mathbf A=0$
and $V=0$. By (2), $\mathcal M^L_{\mu,j}=0$, and both sides vanish. Otherwise let $i_0$ be a local
index at $(j,w)$. Then $\ell_{i_0}=\ell$ and $\mu_{w(i_0)}=\mu$, so by the Hecke ties the level
$(J,\chi)$ of $R$ is unramified outside $B_{i_0}$.

\emph{Step 1: the Hecke operators apply.} Let $(v,k)$ be generic. As $v$ is split over $F$, there are
a split place $s$ of $E/F$ with $s_E=\bar v$, namely the place of $F$ below $v$ with the embedding
of $E$ given by $\bar v$, and a uniformizer $\varpi$ of $F_s$. If a place of $E$ over $s$ lay in
$Q_{i_0}=Q$, then by (3) so would $s_E=\bar v$, and so would $v$; but $v\notin Q$. The places of $E$
over $s$ lie over the same prime as $v$, so none of them lies over $g_0$, $g_1$ or $p^*(\ell)$. So
$s\notin B_{i_0}$, and $s$ does not lie over $\ell$. By (4), $R\circ h(v,k)=t_{s,\varpi,k}\circ R$,
and the image of $R$ is the type space of $(J,\chi)$, to which
Proposition~\ref{asm-support-hecke-standard} applies.

\emph{Step 2: the generators commute.} Let $(v,k)$ and $(v',k')$ be generic, with split places and
uniformizers as in Step 1. By Proposition~\ref{asm-support-hecke-standard}(2),
\[
 R\,h(v,k)h(v',k')=t_{s,\varpi,k}t_{s',\varpi',k'}R=t_{s',\varpi',k'}t_{s,\varpi,k}R
 =R\,h(v',k')h(v,k),
\]
so $h(v,k)$ and $h(v',k')$ commute, as $R$ is injective. Since $\iota_{j,w}$ carries a basis to a
basis and intertwines $h$ with $h_K$, the operators $h_K(v,k)$ and $h_K(v',k')$ commute.

\emph{Step 3: the eigenvalues of the generators lie in $\mathcal O$.} Let $(v,k)$ be generic, and let
$c$ be a coefficient of the characteristic polynomial of $h(v,k)$. Its image under $\iota$ is the
corresponding coefficient for the complex-linear extension $h_{\mathbf C}$ of $h(v,k)$ to
$\mathcal M_{\mu,j}$. As $t_{s,\varpi,k}$ is linear on the type space by
Proposition~\ref{asm-support-hecke-standard}(1), the realization carries $h_{\mathbf C}$ to
$t_{s,\varpi,k}$. If $h_{\mathbf C}y=\lambda y$ with $y\ne0$, then the realization of $y$ is a
nonzero eigenvector of $t_{s,\varpi,k}$ for $\lambda$, so by
Proposition~\ref{asm-support-hecke-standard}(3) $d\lambda$ is an algebraic integer for some integer
$d$ prime to $\ell$. Such complex numbers form a ring, which contains the coefficients of every monic
polynomial whose roots it contains. So $d'c$ is an algebraic integer for some integer $d'$ prime to
$\ell$, and $v_w(c)=v_w(d'c)\ge0$, since $w$ lies over $\ell$. By (1), $e_{j,w}(c)\in\mathcal O$.
Since $\iota_{j,w}$ carries a basis to a basis, the characteristic polynomial of $h_K(v,k)$ is the
image under $e_{j,w}$ of that of $h(v,k)$, a monic polynomial with coefficients in $\mathcal O$. An
eigenvalue of $h_K(v,k)$ is a root of it, so it is integral over the discrete valuation ring
$\mathcal O$ and lies in $\mathcal O$.

\emph{Step 4: counting.} By Steps 2 and 3 and (4), Lemma~\ref{asm-support-eigenspaces} applies to the
generators of $\mathbf A$, with $\mathcal I$ the set of the $\mathfrak m(i)$ for the local indices $i$
at $(j,w)$, which contains every maximal ideal by (6). As $i\mapsto\mathfrak m(i)$ is injective,
\[
 \sum_i\dim_KV[\mathfrak m(i)]=\dim_KV.
\]
By (5), $\Psi_i$ is injective with image $V[\mathfrak m(i)]$ and semilinear for an isomorphism
$K_i\cong K$, so
$\dim_KV[\mathfrak m(i)]=\dim_{K_i}(K_i\otimes_{\mathcal O_i}\tilde M_i)=\operatorname{rank}_{\mathcal O_i}\tilde M_i$.
By (2), $\dim_KV=\dim_L\mathcal M^L_{\mu,j}$. Together these give the support condition.
\end{proof}

\begin{lemma}\label{asm-rank-from-ties}
Let $D$ be a definite-side datum (Definition~\ref{th-definite-datum}) with coefficient field $L$,
rank $m$ and base tuple $\mu_0$, and let its rows data have $D_0$ base rows, $n_0(\sigma)$ of them of
type $\sigma$, and for each level $1\le j\le m$ a lifted space $\mathcal V_j$ with theta lifts
$\vartheta_{\mu,j}\colon\mathcal M_{\mu,j}\to\mathcal V_j$ from the spaces of forms
(Definition~\ref{asm-support-local-modules}), linear maps of finite-dimensional complex Hilbert
spaces. Let $d_j$ be the dimension of the image of the adjoint $\vartheta_{\mu_0,j}^*$, the space of
$j$-th lifts at $\mu_0$ (Definition~\ref{asm-twisted-links}). Let $P$ be a row-side family for $D$
(Definition~\ref{is-row-side-family}), with rows $S_0$ over $L$ of dimension $D(S_0)$ (the integer
$D$ of Definition~\ref{is-indefinite-datum}), integers $d_W(i)$, matrices $\iota_t$ and
$\theta_{i,t}$, parts $E_c$ and $W_c(i)$ and subspaces $Q_c(t)$ and $Z_c(i,t)$ of its definite side,
for the lifts $0\le i<m$ of $P$, the tuples $t$ and the parts $c$ of its base. The lift $i$ of $P$
corresponds to the level $i+1$ of $D$; let $\mu_t$, $e_M(i,t)$ and $j_W(i)$ be the identifications
of $P$, and $\Pi_{i,t}$ its projection matrix at $(i,t)$. Let $T$ be tied data for $D$, as in the
twist records of Definition~\ref{asm-definite-side-hypotheses}, and suppose that the row side of $P$
is determined by $T$, as for the selected base level there. Suppose that:
\begin{enumerate}
\item for every lift $i$, the images $j_W(i)(e_k)$ of the standard basis vectors,
  $1\le k\le d_W(i)$, form a basis of $\mathcal V_{i+1}$;
\item the matrices $\iota_t$ and $\theta_{i,t}$ are bijective; $\iota_{\mu_0}Q_c(\mu_0)=E_c$,
  $\theta_{i,\mu_0}Z_c(i,\mu_0)=W_c(i)$ and $\dim W_c(i)=\dim E_c$ for every lift $i$ and part $c$;
  the $Q_c(\mu_0)$ are independent and span; and for every lift $i$ the $Z_c(i,\mu_0)$ are
  independent and span;
\item the rows hold with $\Lambda_3$ (item (3)(d) of Definition~\ref{asm-definite-side-hypotheses}):
  there is a map $\sigma\mapsto c(\sigma)$ from the types of $D$ to the parts of $P$ such that, for
  every lift $i$ and part $c$, the isotypic parts $I_{i+1}(\sigma)$ with $c(\sigma)=c$ are
  independent with sum the complex span of $j_W(i)(W_c(i))$; the parameters
  $\lambda_D(\mu_t,i+1,\sigma)$ are the scalars of $S_0$, as there; and
  $\sum_{c(\sigma)=c}n_0(\sigma)=\dim_LE_c$ for every part $c$;
\item for every lift $i$, tuple $t$ and coordinate vector $x$, the theta lift $\vartheta_{\mu_t,i+1}$
  maps the image of $e_M(i,t)x$ in $\mathcal M_{\mu_t,i+1}$ to $j_W(i)(\theta_{i,t}\Pi_{i,t}x)$;
\item there is a tuple $t$ with $\mu_t=\mu_0$ such that $\Pi_{i,t}$ is surjective for every lift
  $i$.
\end{enumerate}
Then $\vartheta_{\mu_0,j}^*$ is injective and $d_j=D_0$ for every level $1\le j\le m$.
\end{lemma}
\leanref{Development/AutomorphicArithmetic/ThetaCorrespondence/ThetaPairLifts/ThirdPlaceSeries/BaseRowsEqRecordRows.lean}{Formalism.L2.N.ExtNRows3.thetaT_injective_and_finrank_Zj_eq_D0_of_surjective}
\leanref{Formalism/LocalToGlobal/NormalizedHeckeFactors/BaseRowsEqRecordRows.lean}{Formalism.L2.N.ExtNRows3.thetaP_surj}
\leanref{Formalism/LocalToGlobal/IndefiniteForms/IndefUse.lean}{Formalism.L2.Indef.IndefRows.dW_eq}
\leanref{Development/AutomorphicArithmetic/ThetaCorrespondence/ThetaPairLifts/ThirdPlaceSeries/BaseRowsEqRecordRows.lean}{Formalism.L2.N.ExtNRows3.base_rows_eq_record_rows_alt}
\leanref{Development/AutomorphicArithmetic/ThetaCorrespondence/ThetaPairLifts/ThirdPlaceSeries/RankZeroOfRowConditions.lean}{Formalism.L2.N.RankZeroOfLevelThree.rank_zero_of_row_conditions_alt}
\begin{proof}
The proof does not use $T$, the determination of the row side by $T$, or the first two conditions
of (3). Fix a lift $0\le i<m$, put $j=i+1$, and let $t$ be a tuple as in (5).

\emph{Step 1: $\vartheta_{\mu_0,j}$ is surjective.} By (1) it suffices that every $j_W(i)(e_k)$ is in
its image. By (2) there is $z$ with $\theta_{i,t}z=e_k$, and by (5) there is $x$ with
$\Pi_{i,t}x=z$. By (4), $\vartheta_{\mu_0,j}$ maps the image of $e_M(i,t)x$ to
$j_W(i)(\theta_{i,t}z)=j_W(i)(e_k)$.

\emph{Step 2: the adjoint is injective.} If $\vartheta_{\mu_0,j}^*z=0$, choose $y$ with
$\vartheta_{\mu_0,j}y=z$; then
$\langle z,z\rangle=\langle\vartheta_{\mu_0,j}^*z,y\rangle=0$, so $z=0$. Hence
$d_j=\dim\mathcal V_j$, which is $d_W(i)$ by (1).

\emph{Step 3: $d_W(i)=D(S_0)$.} By (2), $\theta_{i,\mu_0}$ is bijective and the $Z_c(i,\mu_0)$ are
independent and span, so their images $W_c(i)$ are independent and span $L^{d_W(i)}$. In the same
way the $E_c=\iota_{\mu_0}Q_c(\mu_0)$ are independent and span $L^{D(S_0)}$. As
$\dim W_c(i)=\dim E_c$,
\[
 d_W(i)=\sum_c\dim W_c(i)=\sum_c\dim E_c=D(S_0).
\]

\emph{Step 4: $D(S_0)=D_0$.} Every base row has one type, so by Step 3 and (3)
\[
 D(S_0)=\sum_c\dim E_c=\sum_c\sum_{c(\sigma)=c}n_0(\sigma)=\sum_\sigma n_0(\sigma)=D_0.
\]
Thus $d_j=D_0$.
\end{proof}

\begin{lemma}\label{asm-primitive-rows-parts}
Let a family of definite forms over a number field $L$ be given (Definition~\ref{is-realisation}),
with finite group $G$, tuples $t$, base tuple $\mu_0$, lifts $j$, transfers $\operatorname{tr}_j$ and
set $\Pi$ of types of the base, and suppose that it satisfies the family conditions with the
conditions on its inert indices (Lemma~\ref{is-eigenvalue-base-inputs}(1)). Let $e_X\in L[G]$ be the
sum of the central idempotents of the characters of $G$ in the family. On the row forms at a tuple
$t$, and on the lift forms at $(j,t)$, a system $x$ of eigenvalues of the Hecke operators has a
Hecke polynomial at each place of these forms, and the forms have at each place a block product, a
polynomial with constant term $1$. A system $x$ of eigenvalues of the Hecke operators and the inert
Hecke operators has \emph{type} $p$ if at every place of the forms the Hecke polynomial of $x$ is the
Hecke polynomial of $p$ times the block product, and at every inert index of the forms the value of
$x$ is that of $p$. The Hecke block of the forms consists of the vectors fixed by $e_X$ in the sum of
the joint generalized eigenspaces of the Hecke operators for the systems whose Hecke polynomial at
every place is divisible by the block product. Write $R_t(p)$ and $L_{j,t}(p)$ for the sums of the
joint generalized eigenspaces of the Hecke and inert Hecke operators, on the row forms at $t$ and on
the lift forms at $(j,t)$, for the systems of type $p$.

Let $P\subset G$ be a subgroup with a character $\psi\colon P\to L^\times$, put
$e_{P,\psi}=|P|^{-1}\sum_{h\in P}\psi(h)^{-1}h$, and let $\Gamma(t)$ and $\Gamma(j,t)$ be the images
of $e_Xe_{P,\psi}$ on the row forms at $t$ and of $e_Xe_{P,\psi^{-1}}$ on the lift forms at $(j,t)$.
Let subspaces $Z_t$ of the row forms at $t$ and $Z_{j,t}$ of the lift forms at $(j,t)$ be given,
stable under the Hecke operators, the inert Hecke operators and $G$. Consider the definite side that
the family builds from $(P,\psi)$ and these subspaces (Definition~\ref{is-complex-data}), with
coordinate maps $c_t$ and $c_{j,t}$, lattices $M_t$, matrices $P_t$ and parts
\[
 Q_p(t)=c_t^{-1}\big(R_t(p)\big),\qquad Z_p(j,t)=c_{j,t}^{-1}\big(L_{j,t}(\operatorname{tr}_jp)\big)
 \qquad(p\in\Pi).
\]
\begin{enumerate}
\item For every tuple $t$, $P_t$ is surjective and $P_tM_t$ is a lattice in $L^{d_Q(t)}$.
\item If the $Z_t$ lie in the Hecke blocks, then the $Q_p(\mu_0)$, $p\in\Pi$, are independent and
  span $L^{d_Q(\mu_0)}$.
\item If the family has lift-distinct types, then for every lift $j$ the $Z_p(j,\mu_0)$, $p\in\Pi$,
  are independent.
\item Suppose that the family has lift-distinct types, and let $j$ be a lift such that every form in
  $Z_{j,\mu_0}\cap\Gamma(j,\mu_0)$ lies in the sum of the $L_{j,\mu_0}(\operatorname{tr}_jp)$,
  $p\in\Pi$. Then the $Z_p(j,\mu_0)$ span $L^{d_Z(j,\mu_0)}$.
\end{enumerate}
\end{lemma}
\leanref{Formalism/LocalToGlobal/Divisibility/InertPartitionLePartF.lean}{Formalism.L2.DIV.DefFam3.side_well_formed}
\leanref{Formalism/LocalToGlobal/Divisibility/InertPartitionLePartF.lean}{Formalism.L2.DIV.InertDat3.inert_parts_independent}
\leanref{Formalism/LocalToGlobal/Divisibility/InertPartitionLePartF.lean}{Formalism.L2.DIV.InertDat3.inert_parts_span_top}
\leanref{Formalism/LocalToGlobal/Divisibility/InertPartitionLePartF.lean}{Formalism.L2.DIV.InertDat3.sameTypes_self3}
\leanref{Formalism/LocalToGlobal/Divisibility/InertPartitionLePartF.lean}{Formalism.L2.DIV.InertDat3.excl3}
\leanref{Formalism/LocalToGlobal/Divisibility/RowLatLemmas.lean}{Formalism.L2.DIV.BlockDat.Lam_lat}
\begin{proof}
By the family conditions with the conditions on the inert indices, on each space of row or lift forms
the Hecke operators and the inert Hecke operators commute in pairs and with $G$, and each of them
splits over $L$. So the space is the direct sum of the joint generalized eigenspaces of these
operators, for the finitely many systems that occur. For a set $C$ of systems let $\pi_C$ be the
projection onto the sum of the joint generalized eigenspaces of the systems in $C$, along those of
the other systems. It is a polynomial in the operators, so it commutes with every element of $L[G]$
and preserves every subspace that the operators preserve.

Two facts about types are used. First, a system has at most one type among types that differ at some
place or at some inert index of the forms: the block product at a place is not zero, so there the
Hecke polynomial of the system determines that of its type, and at the inert indices the values of
the type are those of the system. Second, two distinct types of the base differ at some place or
inert index of the row forms at $\mu_0$, since both have Hecke polynomial $1$ and value $0$ at the
places and inert indices of the family that are not those of these forms. If the family has
lift-distinct types, then the types $\operatorname{tr}_jp$, $p\in\Pi$, differ pairwise at some place
or inert index of the lift forms at $(j,\mu_0)$.

As in the proof of Lemma~\ref{is-baselinks-definite-coordinates}, $c_t$ sends a coordinate vector $y$
to the form in $\Gamma(t)$ with coordinates $E_ty$ in a basis of $\Gamma(t)$, where $E_t$ is a basis
matrix of the coordinates of $Z_t\cap\Gamma(t)$; so $c_t$ is injective with image $Z_t\cap\Gamma(t)$.
The matrix $P_t$ sends a vector to the coordinates, in the columns of $E_t$, of its orthogonal
projection onto the column space of $E_t$ for the Hermitian form $H_V(t)$, which is positive definite
at every infinite place by the family conditions; and $M_t$ is a lattice. The same holds for the lift
forms.

(1) Since $P_tE_t=1$, $P_t$ is surjective, and the image of the lattice $M_t$ under a surjective
linear map is a lattice.

(2) and (3), independence. By the two facts, the $R_{\mu_0}(p)$, $p\in\Pi$, are sums of joint
generalized eigenspaces over disjoint sets of systems, so they are independent. As $c_{\mu_0}$ is
injective, the $Q_p(\mu_0)$ are independent. If the family has lift-distinct types, the same argument
applies to the $L_{j,\mu_0}(\operatorname{tr}_jp)$ and the $Z_p(j,\mu_0)$.

Spanning. Let $V$ be the row forms at $\mu_0$, with $Z=Z_{\mu_0}$, $\Gamma=\Gamma(\mu_0)$,
$c=c_{\mu_0}$ and $\tau_p=p$, or the lift forms at $(j,\mu_0)$, with $Z=Z_{j,\mu_0}$,
$\Gamma=\Gamma(j,\mu_0)$, $c=c_{j,\mu_0}$ and $\tau_p=\operatorname{tr}_jp$, where in the second case
the family has lift-distinct types. Write $V(\tau_p)$ for $R_{\mu_0}(p)$ or
$L_{j,\mu_0}(\operatorname{tr}_jp)$. Suppose that every form in $Z\cap\Gamma$ lies in the sum of the
$V(\tau_p)$. Let $y$ be a coordinate vector and $f=c(y)$, and let $C_p$ be the set of systems of type
$\tau_p$. These sets are disjoint by the two facts, and $f$ lies in the sum of the $V(\tau_p)$, so
$f=\sum_p\pi_{C_p}f$. The element $e_Xe_{P,\psi}$, or $e_Xe_{P,\psi^{-1}}$, of $L[G]$ whose image is
$\Gamma$ is idempotent by the family conditions; it fixes $f$ and commutes with $\pi_{C_p}$, so it fixes
$\pi_{C_p}f$; and $\pi_{C_p}f\in Z$ because $Z$ is stable. So
$\pi_{C_p}f=c(y_p)$ for some $y_p$, which lies in $c^{-1}(V(\tau_p))$, and $y=\sum_py_p$ because $c$
is injective. So the $c^{-1}(V(\tau_p))$ span. This is (4).

For (2) it remains to show that every form $f\in Z_{\mu_0}\cap\Gamma(\mu_0)$ lies in the sum of the
$R_{\mu_0}(p)$, $p\in\Pi$. As $Z_{\mu_0}$ lies in the Hecke block, $f=\sum_xf_x$, where $x$ runs over
systems whose Hecke polynomials are divisible by the block products and $f_x$ lies in the joint
generalized eigenspace of $x$. The idempotent $e_X$ fixes $f$, because
$e_Xe_Xe_{P,\psi}=e_Xe_{P,\psi}$, and it commutes with the projections, so it fixes each $f_x$. Let
$f_x\ne0$, and let $p_x$ be the type whose Hecke polynomial at a place of the row forms at $\mu_0$ is
that of $x$ divided by the block product, and $1$ at the other places of the family, and whose value
at an inert index of these forms is that of $x$, and $0$ at the other inert indices. Then $x$ has type
$p_x$, and the eigenspace of $p_x$ contains $f_x$, a nonzero vector fixed by $e_X$, that is, a nonzero
vector of the base space at the full level. So $p_x\in\Pi$ and $f_x\in R_{\mu_0}(p_x)$.
\end{proof}

\begin{lemma}\label{asm-primitive-rows-base}
Let $D$ be a definite-side datum (Definition~\ref{th-definite-datum}) with coefficient field $L$ and
the embedding $\sigma_0$ of its base place. Let a family of definite forms of $D$ for a finite group
$G$ be given that satisfies the family conditions with the conditions on its inert indices
(Lemma~\ref{is-eigenvalue-base-inputs}(1)) and has lift-distinct types, and let $R$ be a realization
of $D$ for it with a face $F$, with subgroups $P_0,P_1$ and characters $\psi_0,\psi_1$
(Definition~\ref{is-realisation}). Suppose that $R$ satisfies the weak characterization at $F$ and has
primitive rows at $F$ (Definition~\ref{is-primitive-rows}). Let $X$ be the complex data of $R$ and $F$
(Definition~\ref{is-complex-data}), with definite subspaces $Z_t$ and $Z_{j,t}$, maps $i_b(t)$, levels
$E^b$ and parts $E_p\subset E$ and $W_p(j)\subset W_j$. For each level $b$ let $S_b$ be the
indefinite-side datum of $X$ at level $b$, with matrices $\iota_t$, on the definite side that the
family builds from $(P_b,\psi_b)$ and the definite subspaces, with lattices $M_t$, matrices $P_t$ and
the parts $Q_p(t)$ and $Z_p(j,t)$ of Lemma~\ref{asm-primitive-rows-parts}. Write $E^L_p\subset L^D$
and $W^L_p(j)\subset L^{d_W(j)}$ for the parts of the rows $S_0$, the preimages of $E_p$ and $W_p(j)$
(Definition~\ref{is-complex-data}(b) and (c)). Then:
\begin{enumerate}
\item the definite subspaces are stable under the Hecke operators, the inert Hecke operators and $G$,
  and lie in the Hecke blocks;
\item at both levels $b$ and for every tuple $t$, the map $i_b(t)$ is injective with image $E^b$, the
  matrix that expresses $i_b(t)$ through $i_b(\mu_0)$ is the image under $\sigma_0$ of the matrix
  $\iota_t$ of $S_b$, the matrices $\iota_t$ are bijective, $\iota_{\mu_0}=1$, and $P_tM_t$ is a
  lattice;
\item the rows $S_0$, with matrices $\theta_{j,t}$, satisfy: the matrices $\theta_{j,t}$ are
  bijective; $\iota_{\mu_0}Q_p(\mu_0)=E^L_p$, $\theta_{j,\mu_0}Z_p(j,\mu_0)=W^L_p(j)$ and
  $\dim_LW^L_p(j)=\dim_LE^L_p$ for every lift $j$ and every type $p$ of the base; the $Q_p(\mu_0)$
  are independent and span; and for every lift $j$ the $Z_p(j,\mu_0)$ are independent and span.
\end{enumerate}
\end{lemma}
\leanref{Development/AutomorphicArithmetic/ThetaCorrespondence/ThetaPairLifts/ThirdPlaceSeries/IndefRRealTP.lean}{Formalism.IndefR.IndefiniteRealisation.rowsHold0}
\leanref{Formalism/LocalToGlobal/IndefiniteReal/BaseLiftsPackageOfPrimitiveLinks.lean}{Formalism.IndefR.Links.complex_data_conditions_base_level_of_primitive_links}
\leanref{Formalism/LocalToGlobal/IndefiniteReal/PrimitiveStatementsBaseLevel.lean}{Formalism.IndefR.Cx.indefiniteDataAtBaseLevel}
\leanref{Formalism/LocalToGlobal/IndefiniteReal/CxCoreRG.lean}{Formalism.IndefR.Cx.CxData.rowsOfR}
\leanref{Formalism/LocalToGlobal/IndefiniteReal/MemRangeAdjointOfKerStable.lean}{Formalism.IndefR.Cx.CxData.cutStable}
\leanref{Formalism/LocalToGlobal/IndefiniteReal/EqZeroOfMemRangeAdjoint.lean}{Formalism.IndefR.Cx.cutRows}
\leanref{Formalism/LocalToGlobal/IndefiniteReal/RationalAdjoints/RatMuOfRowTest.lean}{Formalism.IndefR.Rat.ratMuOfRowTest}
\leanref{Formalism/LocalToGlobal/IndefiniteReal/LiftCoordinatesInGoodPartFamily.lean}{Formalism.IndefR.Cx.part_iff_rows3}
\leanref{Formalism/LocalToGlobal/IndefiniteReal/FinrankISupLeSum.lean}{Formalism.IndefR.Cx.CxData.cutDims}
\leanref{Formalism/LocalToGlobal/IndefiniteReal/MajAtEqZeroOfMemQpOfNe.lean}{Formalism.IndefR.Cx.types_of_lift_eigenvalues}
\leanref{Formalism/LocalToGlobal/IndefiniteReal/NormalDecomposition/PairingCommutesLiftOperators.lean}{Formalism.IndefR.U2.typ0Of}
\begin{proof}
Write $\Theta_t\colon F_t\to E$, $j_t$, $\theta'_{j,t}\colon F_{j,t}\to W_j$, $j_{j,t}$ and
$W_p(j)\subset W_j$ for the complex data, $c_{b,t}$ and $c_{b,j,t}$ for the coordinate maps at level
$b$, $t_b(j,t)$ and $w_b(j)$ for the maps of Definition~\ref{is-complex-data}(c), and
$\Gamma_b(t)$, $\Gamma_b(j,t)$ for the spaces $\Gamma(t)$, $\Gamma(j,t)$ of
Lemma~\ref{asm-primitive-rows-parts} for $(P_b,\psi_b)$. The numbered conditions are those of
Definition~\ref{is-primitive-rows}.

\emph{Step 0: inputs.} The weak characterization contains the conditions on the models used in Step 1
of the proof of Lemma~\ref{is-eigenvalue-base-inputs}: $j_t$ carries a basis of the row forms at $t$
to a basis of $F_t$, and $\langle j_t(x),j_t(y)\rangle$ is the image under $\sigma_0$ of the
class-sum pairing of $x$ and $y$; likewise for the lift forms. So every operator $T$ of the row forms
at $t$ induces an operator $T_F$ of $F_t$ with $j_tT=T_Fj_t$, and if $T'$ is adjoint to $T$ for the
class-sum pairing, then $T'_F=T_F^*$. Steps 1 to 5 of the proof of
Lemma~\ref{is-eigenvalue-base-inputs} use, of the characterization of $R$, only these conditions, and
of the phases only that $\theta'_{j,t}\theta'^*_{j,t}=\mathrm{th}_G(j,t)\mathrm{th}_G(j,t)^*$, which
holds because the weak characterization gives the phases absolute value $1$. They give, for every lift
$j$ and tuple $t$: $w_0(j)$ is injective; $t_0(j,t)$ is injective with the same image $W^0_j$; the
matrix that expresses $t_0(j,t)$ through $w_0(j)$ is $\sigma_0(\theta_{j,t})$ for a matrix
$\theta_{j,t}$ over $L$; and $Z_{j,\mu_0}\cap\Gamma_0(j,\mu_0)$ is the intersection of
$\Gamma_0(j,\mu_0)$ with the sum of the joint eigenspaces of the systems $y$ in the set $C_j$ of
Step 5 there, the systems in $L$ of the operators on the lift forms at $(j,\mu_0)$ whose images under
$\sigma_0$ are lifted.

\emph{Step 1: stability, part (1) of the lemma.} Let $T$ be a Hecke operator or an inert Hecke
operator on the row forms at $t$, or
an element of $G$. By (1), or for $T\in G$ because $G$ preserves the class-sum pairing, $T$ has an
adjoint $T'$ among these operators for the class-sum pairing, so $T'_F=T_F^*$ by Step 0. By (2) and
(3) there is an operator $h$ of $E$ with $\Theta_tT'_F=h\Theta_t$, since both sides agree on the image
of $j_t$, which spans $F_t$. So $\ker\Theta_t$ is stable under $T'_F$, and its orthogonal complement
$\Theta_t^*(E)$ is stable under $T_F$. If $x\in Z_t$, then $j_t(Tx)=T_Fj_t(x)\in\Theta_t^*(E)$, so
$Tx\in Z_t$. The lift forms are treated in the same way, and the definite subspaces lie in the Hecke
blocks by (4).

\emph{Step 2: the maps $i_b(t)$.} Fix a level $b$ and a tuple $t$. Apply
Lemma~\ref{is-baselinks-lift-cut} to $\theta=\Theta_t$, the vectors $u_k=j_t(c_{b,t}(e_k))$, the
span $\Gamma$ of $j_t(\Gamma_b(t))$, $W_0=E^b$, and the parts $E_p$ with
$\lambda_p=\kappa\prod_i\lambda^{(3)}_{i,t,p,v_0}$, which is positive by (7). The map
$c\mapsto\Theta_t(\sum_kc_ku_k)$ of that lemma is $i_b(t)$. The Gram matrix of the $u_k$ is $|P_b|(X^{-1})^*X^{-1}$ for the invertible matrix
$X=X_Q(t)_{v_0}$, by Step 0 and the construction of the definite side, as in
Lemma~\ref{is-baselinks-definite-coordinates}(2); so the $u_k$ are linearly independent.

Hypothesis (1) of Lemma~\ref{is-baselinks-lift-cut}. By (5) and Lemma~\ref{is-baselinks-rational-tests}(2), $\Theta_t^*(E)$ is spanned by
$j_t(Z_t)$. The family conditions make $e=e_Xe_{P_b,\psi_b}$ act as an idempotent on the row forms,
so Lemma~\ref{is-baselinks-group-algebra}(3) and (1), with the action of $G$ on $F_t$ induced through
$j_t$ (Step 0) and on $E$ by (3), give that $\Theta_t^*(E)\cap\Gamma$ is spanned by
$j_t(Z_t\cap\Gamma_b(t))$. As $Z_t\cap\Gamma_b(t)$ is the image of $c_{b,t}$ and $j_t\circ c_{b,t}$ is
$\sigma_0$-semilinear, this span is that of the $u_k$.

Hypothesis (2) of that lemma. Let $u\in\Gamma$. Lemma~\ref{is-baselinks-group-algebra}(2) gives
$u'\in\Theta_t^*(E)\cap\Gamma$ with $\Theta_t(u')=\Theta_t(u)$, and $u'$ is a combination of the
$u_k$ by its hypothesis (1). Each $c_{b,t}(e_k)$ lies in $Z_t$ and in $\Gamma_b(t)$, which lies in the
image of $e_X$, the space $\Gamma(t)$ for the trivial subgroup and character; so it lies in the image
of the coordinate map for the face with the trivial subgroup and character, and $\Theta_t(u)$ lies in
the image of $i^{\mathrm{pr}}(t)$, which is that of $i^{\mathrm{pr}}(\mu_0)$ by (8). Every $h\in P_b$
acts on $\Gamma_b(t)$ by $\psi_b(h)$, so on $\Theta_t(u)$ by $\sigma_0(\psi_b(h))$, by
Lemma~\ref{is-baselinks-group-algebra}(5). By (8), $\Theta_t(u)\in E^b$. Conversely, let $v\in E^b$.
By (8), $v$ lies in the image of $i^{\mathrm{pr}}(\mu_0)$, which is spanned by vectors
$\Theta_{\mu_0}j_{\mu_0}(x)$ with $x$ in the image of $e_X$; these are fixed by $e_X$ by
Lemma~\ref{is-baselinks-group-algebra}(1), so $e_Xv=v$, and $P_b$ acts on $v$ through
$\sigma_0\circ\psi_b$. Lemma~\ref{is-baselinks-group-algebra}(4) and (1) give
$\Theta_t^*v\in\Gamma$.

Hypothesis (3) of that lemma. By (9), $\Theta_t\Theta_t^*=\lambda_p$ on $E_p$, and by (10) the spaces $E_p\cap E^b$
span $E^b$.

So $i_b(t)$ is injective with image $E^b$, for every $t$.

\emph{Step 3: the matrices.} By (5) there are linear functionals $\ell_1,\dots,\ell_N$ separating the
points of $E$ whose values at the vectors $\Theta_tj_t(x)$, for row forms $x$, lie in $\sigma_0(L)$. Fix
$t$ and $k$. By Step 2, $i_b(t)(e_k)=i_b(\mu_0)(c)$ for a unique $c\in\mathbf C^{D}$, the $k$-th column
of the matrix that expresses $i_b(t)$ through $i_b(\mu_0)$. The values of the $\ell_i$ at the vectors
$i_b(\mu_0)(e_l)$ and at $i_b(\mu_0)(c)=\Theta_tj_t(c_{b,t}(e_k))$ lie in $\sigma_0(L)$, so
$c\in\sigma_0(L)^{D}$ by Lemma~\ref{is-baselinks-rational-tests}(1). So this matrix is
$\sigma_0(\iota_t)$. It is bijective, since $i_b(t)$ and $i_b(\mu_0)$ are injective with the same
image, and it is $1$ for $t=\mu_0$. A matrix over $L$ whose image under $\sigma_0$ is bijective is
bijective, so $\iota_t$ is bijective and $\iota_{\mu_0}=1$. In the same way, by Step 0, $\theta_{j,t}$
is bijective. Lemma~\ref{asm-primitive-rows-parts}(1) gives the lattices $P_tM_t$, which completes
part (2) of the lemma.

\emph{Step 4: the parts.} Let $y\in L^{D}$ and $f=c_{0,\mu_0}(y)$. By semilinearity
$i_0(\mu_0)(\sigma_0y)=\varphi(f)$ for the $\sigma_0$-semilinear map
$\varphi=\Theta_{\mu_0}\circ j_{\mu_0}$. By Step 0 and (2), $\varphi$ carries each Hecke and inert
Hecke operator $T$ to the operator $h_T$ of $E$ of (2), which commute in pairs, and by (2) $E_p$ is the
sum of the joint generalized eigenspaces of the $h_T$ for the systems $\sigma_0(x)$, $x$ of type $p$.
The map $\varphi$ is injective on $Z_{\mu_0}$: if $z\in Z_{\mu_0}$ and $\varphi(z)=0$, then $j_{\mu_0}(z)$
lies in $\Theta_{\mu_0}^*(E)$ and in $\ker\Theta_{\mu_0}$, so it is $0$. We show that $y\in Q_p(\mu_0)$,
that is $f\in R_{\mu_0}(p)$, if and only if $\varphi(f)\in E_p$. As $\varphi$ is semilinear and
intertwines $T$ with $h_T$, it carries the joint generalized eigenspace of a system $x$ into that of
$\sigma_0(x)$; this gives one direction. Conversely, let $\varphi(f)\in E_p$, and let $\pi$ be the
projection of the proof of Lemma~\ref{asm-primitive-rows-parts} onto $R_{\mu_0}(p)$. Then
$f-\pi f\in Z_{\mu_0}$ by Step 1, and $\varphi(f-\pi f)$ lies in $E_p$ and in the sum of the joint
generalized eigenspaces of the $h_T$ for the systems $\sigma_0(x)$ with $x$ not of type $p$. These two
sums meet in $0$, since $\sigma_0$ is injective and the $h_T$ commute. So $\varphi(f-\pi f)=0$,
$f=\pi f$ and $f\in R_{\mu_0}(p)$. As $E^L_p$ is the set of $y$ with
$i_0(\mu_0)(\sigma_0y)\in E_p$ and $\iota_{\mu_0}=1$, this gives $\iota_{\mu_0}Q_p(\mu_0)=E^L_p$.

In the same way, $y\in Z_p(j,\mu_0)$ if and only if $t_0(j,\mu_0)(\sigma_0y)\in W_p(j)$, using the
operators on $W_j$ of (2) and the map $\theta'_{j,\mu_0}\circ j_{j,\mu_0}$. By the definition of
$\theta_{j,\mu_0}$, $w_0(j)(\sigma_0(\theta_{j,\mu_0}y))=t_0(j,\mu_0)(\sigma_0y)$. So
$\theta_{j,\mu_0}y$ lies in $W^L_p(j)$, the set of $z$ with $w_0(j)(\sigma_0z)\in W_p(j)$, if and
only if $y\in Z_p(j,\mu_0)$. As $\theta_{j,\mu_0}$ is bijective,
$\theta_{j,\mu_0}Z_p(j,\mu_0)=W^L_p(j)$.

\emph{Step 5: independence and spanning.} By Step 1 and
Lemma~\ref{asm-primitive-rows-parts}(2) and (3), the $Q_p(\mu_0)$ are independent and span, and for
every lift $j$ the $Z_p(j,\mu_0)$ are independent. For the spanning of the $Z_p(j,\mu_0)$ we check the
hypothesis of Lemma~\ref{asm-primitive-rows-parts}(4), using Step 0: it suffices that every $y\in C_j$
has type $\operatorname{tr}_jp$ for some $p\in\Pi$. As $\sigma_0(y)$ is lifted, there is a nonzero
$w\in W^0_j$ that is an eigenvector of the operators on $W_j$ of (2) for the eigenvalues
$\sigma_0(y)$. By (10), $W^0_j$ is the sum of the $W_p(j)\cap W^0_j$, and by (2) each $W_p(j)$ is the
sum of the joint generalized eigenspaces of these operators for the systems $\sigma_0(z)$ with $z$ of
type $\operatorname{tr}_jp$. As joint generalized eigenspaces of distinct systems are independent,
$\sigma_0(y)=\sigma_0(z)$ for such a $z$, so $y=z$.

\emph{Step 6: the dimensions.} The map $y\mapsto i_0(\mu_0)(\sigma_0y)$ carries $E^L_p$ into
$E_p\cap E^0$ and $L$-linearly independent vectors to $\mathbf C$-linearly independent ones, so
$\dim_LE^L_p\le\dim_{\mathbf C}(E_p\cap E^0)$. Summing over $p$: the $E^L_p=Q_p(\mu_0)$ span $L^D$,
so $\sum_p\dim_LE^L_p\ge D$; and the $E_p\cap E^0$ are orthogonal by (10) and lie in
$E^0=i_0(\mu_0)(\mathbf C^D)$, so $\sum_p\dim(E_p\cap E^0)\le D$. Hence
$\dim_LE^L_p=\dim(E_p\cap E^0)$ for every $p$. In the same way, since $w_0(j)$ is injective,
$\dim_LW^L_p(j)\le\dim(W_p(j)\cap W^0_j)$; $\theta_{j,\mu_0}$ carries $Z_p(j,\mu_0)$ injectively into
$W^L_p(j)$ and the $Z_p(j,\mu_0)$ span, so $\sum_p\dim_LW^L_p(j)\ge d_Z(j,\mu_0)$; and the
$W_p(j)\cap W^0_j$ are orthogonal by (10) and lie in $W^0_j$, the image of the injective map
$t_0(j,\mu_0)$ on $\mathbf C^{d_Z(j,\mu_0)}$. Hence $\dim_LW^L_p(j)=\dim(W_p(j)\cap W^0_j)$, which is
$\dim(E_p\cap E^0)=\dim_LE^L_p$ by (10).
\end{proof}

\subsection{Error terms and the Bessel tie}\label{subsection-asm-errors-bessel-tie}

\begin{lemma}\label{asm-division-kernel-glued}
Let $S$ be an indefinite-side datum over a number field $K$ (Definition~\ref{is-indefinite-datum}),
with base tuple $\mu_0$, rank $D$, lattices $M_t$, matrices $P_t$ and $\iota_t$, tuple map
$w\mapsto\mu_w$ and glued lattice $RR$. Suppose that every $\iota_t$ is bijective, that every
$P_tM_t$ is a lattice in $K^{d_Q(t)}$, and that $\mu_w=\mu_0$ for all but finitely many finite places
$w$. Then $RR$ is a lattice in $K^D$: it is a finitely generated $\mathcal O_K$-module that spans
$K^D$.
\end{lemma}
\leanref{Formalism/LocalToGlobal/Divisibility/DivisibilityRemainder/RationalColumns/ColFacts.lean}{Formalism.L2.DIV.RatFullPt.ColFacts.rr_lattice}
\leanref{Formalism/QuotientMetricDivisors/HermitianTwistDivisors/SecondStepGluing.lean}{Formalism.DivisionValues.lattice_gluing}
\begin{proof}
For every tuple $t$, $\Lambda_t=\iota_tP_tM_t$ is a lattice in $K^D$, the image of a lattice under a
bijective linear map. By definition $RR$ is the set of vectors that lie, for every finite place $w$,
in the localization at $w$ of $\Lambda_{\mu_w}$, and $\Lambda_{\mu_w}=\Lambda_{\mu_0}$ for all but
finitely many $w$. By \cite[81:14]{OMe63} there is a lattice $\Lambda$ in $K^D$ whose localization at
every $w$ is that of $\Lambda_{\mu_w}$. A finitely generated torsion-free $\mathcal O_K$-module is the
intersection of its localizations at the finite places, so $\Lambda=RR$.
\end{proof}

\begin{definition}\label{asm-division-kernel-errors}
Let $\mathcal F$ be a family of exponent $\varepsilon$ (Definition~\ref{pd-family}), and let $D$ be
a member at a scale $x>x_1$, with ranks $k,m,t,n$, $f=2n$ and $N(j)=k+j(t+1)$, coefficient field
$L$, host field $E$ containing $K=\mathbf Q(\sqrt{-d})$, auxiliary prime $d$, level primes
$g_0,g_1$, top classes $T$, and embeddings $\mathcal T$ with top embeddings
$\mathcal T_{\mathrm{top}}$ (Definition~\ref{pd-hecke-at-T}). Let its size data
(Definition~\ref{pd-size-data}) have support $\operatorname{supp}$, exponents $a_t$, width $\eta$,
cloud size $S_c$, quantities $b'$ and $X''$, and local errors $E_{u,j}(w)$. For a word $u$ let
$T_u$ be its auxiliary prime and $D_\Gamma(u)$ its stabilizer denominator. For a finite place $w$
of $L$ let $p_w$ be its residue characteristic, $e_w$ its ramification index,
$\ell_w=\log\mathrm N\mathfrak p_w/[L:\mathbf Q]$ and $\mathrm{sz}_w(y)=2\ell_we_wy$, and put
$B_0=fmbx^{\varepsilon}$.
\begin{enumerate}
\item (\emph{Explicit errors}.) The error data of the member (Definition~\ref{pd-family}(4))
  contain elements $\pi_0,\pi_1\in\mathcal O_E\cap K$ with $\pi_i\overline{\pi_i}=g_i$ and constants
  among which $c_g\geq10$. With the constants of $\mathcal F$ they determine, for every lift
  $1\le j\le m$, a real number $E_d(j)$, the \emph{loss at $d$}, and for every lift $j$ and every
  prime $p$ a real number $E_{\mathrm{in}}(j,p)$, the \emph{inert loss at $p$}. For a word $u$, a
  lift $j$ and $\tau\in\mathcal T$, let $n_{u,j}(\tau)=\sum_{r<N(j-1)}n_{u,j,r}(\tau)$ be the sum of
  the row exponents of Definition~\ref{pd-orbit-places}(2), and put
  \[
   \nu^{(i)}_{u,j}(w)=e_w^{-1}\sum_{\tau\in\mathcal T}n_{u,j}(\tau)\,
   v_w\big(\tau(\overline{\pi_i})\big),\qquad c_{\mathrm{er}}=c_gf\big(N(m)^2+m^2\eta\big).
  \]
  Let $F(j)=\sum_{\tau\in\mathcal T}\big(\sum_a|\kappa^+_{j,\tau,a}|+\sum_c|\kappa^-_{j,\tau,c}|\big)$
  be the \emph{Fock degree} of the lift $j$, where $\kappa^{+}_{j,\tau,a}$ and $\kappa^-_{j,\tau,c}$
  are the integer weights that the rows data of $D$ attach to the lift $j$ at $\tau$, for
  $1\le a\le k+j-1$ and $1\le c\le m-j$, the two parts of $\operatorname{sig}W(j)$
  (Definition~\ref{th-ranks}). The local errors are \emph{explicit}: for every word $u$,
  lift $j$ and finite place $w$, with $p=p_w$, $E_{u,j}(w)$ is the exponent $E_T(u,w)$ at $T_u$
  (Definition~\ref{pd-hecke-at-T}(4)) if $p=T_u$, and otherwise
  $E_{u,j}(w)=v_p(D_\Gamma(u))+E'_{u,j}(w)$ with
  \[
   E'_{u,j}(w)=\begin{cases}
   -\frac12F(j)+\nu^{(i)}_{u,j}(w)+c_{\mathrm{er}}&\text{if }p=g_i,\ i\in\{0,1\},\\
   E_d(j)&\text{if }p=d,\\
   E_{\mathrm{in}}(j,p)&\text{if }p\notin\{g_0,g_1\}\text{ is inert for }D\text{ and }p\le X'',\\
   0&\text{otherwise.}
   \end{cases}
  \]
\item (\emph{Second complement inputs}.) Let $C_1,\dots,C_6$ be real numbers, and for a lift $j$ put
  \[
   \mathrm{err}(j)=2\sum_{\tau\in\mathcal T}e_j(\tau)
   +C_1\Big(f\big(kS_c+mt\eta+kN(m)\big)+\frac{|T|\,b'}{x}\Big),
  \]
  with the weights $e_j(\tau)$ of the lifts (Definition~\ref{pd-orbit-places}(1)). The \emph{second
  complement inputs} hold at $D$ with constants $C_1,\dots,C_6$ if for every word $u$, every lift
  $j$ and $i\in\{0,1\}$:
  \begin{enumerate}
  \item $g_i\neq T_u$;
  \item $\sum_w\mathrm{sz}_w\big(\nu^{(i)}_{u,j}(w)\big)=\big(\sum_\tau n_{u,j}(\tau)\big)\log g_i$,
    the sum over the places $w$ over $g_i$;
  \item $\big|\sum_\tau n_{u,j}(\tau)-F(j)-\sum_{t\in T}a_t\big|\le\mathrm{err}(j)$;
  \item $\big(\mathrm{err}(j)+2v_{g_i}(D_\Gamma(u))+2c_{\mathrm{er}}\big)\log g_i\le C_2B_0$;
  \item for every finite set $S$ of places over $T_u$,
    $0\le\sum_{w\in S}\mathrm{sz}_w(E_T(u,w))\le E^{\pi}_T(u)\log T_u$, and
    $E^{\pi}_T(u)\log T_u\le C_3B_0$, where
    $E^{\pi}_T(u)=2\sum_\tau W_\tau(u)+fN(m)+\frac14|\mathcal T_{\mathrm{top}}|(t+1)(m+1)$ is the
    value of $E_T(u,w)$ at the places $w$ that see $\pi_u$;
  \item $2\log D_\Gamma(u)\le C_4B_0$;
  \item $2\log d\,|E_d(j)|\le C_5B_0$;
  \item $\sum_{p\in P}2\log p\,|E_{\mathrm{in}}(j,p)|\le C_6B_0$ for every finite set $P$ of primes
    $p\le X''$ that are inert for $D$.
  \end{enumerate}
  The family $\mathcal F$ satisfies the second complement inputs if there are $C_1,\dots,C_6$ with
  $C_2\geq0$ and $x_0$ such that they hold with these constants at every member at every scale
  $x\geq x_0$ with $x>x_1$.
\end{enumerate}
\end{definition}
\leanref{Formalism/LocalRingCardinality/DataErr.lean}{Formalism.BasicNotions.SizeData₂.EDef}
\leanref{Formalism/LocalRingCardinality/DataErr.lean}{Formalism.BasicNotions.Eval}
\leanref{Formalism/LocalRingCardinality/DataErr.lean}{Formalism.BasicNotions.EatG}
\leanref{Formalism/LocalRingCardinality/DataErr.lean}{Formalism.BasicNotions.vNN}
\leanref{Formalism/LocalRingCardinality/DataErr.lean}{Formalism.BasicNotions.nuAbs}
\leanref{Formalism/LocalRingCardinality/DataForced.lean}{Formalism.BasicNotions.fockDeg}
\leanref{Formalism/SharpRestriction/Divisibility/WeightDegreeSizes.lean}{Formalism.MainNotions.WeightDegreeSizes.errG}
\leanref{Formalism/SharpRestriction/GlobalSchurBounds/ProductFirstZerothParts.lean}{Formalism.MainNotions.product_first_zeroth_parts.WeightBudgetBound}
\leanref{Formalism/SharpRestriction/GlobalSchurBounds/ProductFirstZerothParts.lean}{Formalism.MainNotions.product_first_zeroth_parts.ETpi}
\leanref{Formalism/SharpRestriction/Divisibility/WeightDegreeSizes.lean}{Formalism.MainNotions.WeightDegreeSizes.SecondComplementInputs}
\leanref{Formalism/SharpRestriction/Divisibility/WeightDegreeSizes.lean}{Formalism.MainNotions.WeightDegreeSizes.FamilySecondComplementInputs}

\begin{lemma}\label{asm-division-kernel-complement}
Let $\mathcal F$ be a family of exponent $\varepsilon$ (Definition~\ref{pd-family}) that satisfies
the second complement inputs (Definition~\ref{asm-division-kernel-errors}(2)) with constants
$C_1,\dots,C_6$, $C_2\geq0$, and threshold $x_0$, and put $C=2C_2+C_3+C_4+C_5+C_6$. For a member
at a scale $x>x_1$ let $\operatorname{supp}$, $a_t$, $T$, $g_0$, $g_1$, $E_{u,j}(w)$, $p_w$,
$\mathrm{sz}_w$ and $B_0$ be as in Definition~\ref{asm-division-kernel-errors}. Then at every
member at every scale $x\geq x_0$ with $x>x_1$, for every word $u$ and every lift $1\le j\le m$,
\[
 \Big|\sum_{w\in\operatorname{supp}}\mathrm{sz}_w(E_{u,j}(w))-\Big(\sum_{t\in T}a_t\Big)\log(g_0g_1)\Big|
 \le CB_0
 \quad\text{and}\quad
 \sum_{w\in\operatorname{supp},\ p_w\notin\{g_0,g_1\}}|\mathrm{sz}_w(E_{u,j}(w))|\le CB_0.
\]
In particular $\mathcal F$ satisfies part (6) of the division kernel bounds
(Definition~\ref{is-division-kernel-bounds}).
\end{lemma}
\leanref{Formalism/LocalToGlobal/ClassOfRecordOfChoice.lean}{Formalism.L2.Top3.complement_two_of_items}
\leanref{Formalism/SharpRestriction/Divisibility/WeightDegreeSizes.lean}{Formalism.MainNotions.WeightDegreeSizes.secondComplement_family_of_inputs}
\leanref{Formalism/SharpRestriction/Divisibility/WeightDegreeSizes.lean}{Formalism.MainNotions.WeightDegreeSizes.secondComplement_of_inputs}
\leanref{Formalism/SharpRestriction/Divisibility/WeightDegreeSizes.lean}{Formalism.MainNotions.WeightDegreeSizes.secondComplement_of_split}
\leanref{Formalism/SharpRestriction/Divisibility/WeightDegreeSizes.lean}{Formalism.MainNotions.WeightDegreeSizes.C2atG_of}
\leanref{Formalism/SharpRestriction/Divisibility/WeightDegreeSizes.lean}{Formalism.MainNotions.WeightDegreeSizes.C2off_of}
\leanref{Formalism/LocalRingCardinality/RankBoundMonotoneInDimension.lean}{Formalism.BasicNotions.factEtaMean}
\leanref{Formalism/SharpRestriction/Divisibility/WeightDegreeSizes.lean}{Formalism.MainNotions.WeightDegreeSizes.sum_by_prime}
\leanref{Formalism/SharpRestriction/Divisibility/WeightDegreeSizes.lean}{Formalism.MainNotions.WeightDegreeSizes.sum_val_log_le}
\begin{proof}
Fix a member at a scale $x\geq x_0$ with $x>x_1$, a word $u$ and a lift $j$, and write
$E(w)=E_{u,j}(w)$. The local errors are explicit (Definition~\ref{asm-division-kernel-errors}(1)),
and they vanish off the support. For every prime $p$, the places $w$ over $p$ satisfy
$\sum_{w\mid p}\ell_we_w=\log p$, since $\mathrm N\mathfrak p_w=p^{f_w}$ and
$\sum_{w\mid p}e_wf_w=[L:\mathbf Q]$. Also $B_0\geq0$.

\emph{Step 1: the places over a level prime.} Let $i\in\{0,1\}$, let $S_i$ be the set of places
over $g_i$, and put $v=v_{g_i}(D_\Gamma(u))$ and $n=\sum_\tau n_{u,j}(\tau)$. We show
\[
 \Big|\sum_{w\in\operatorname{supp},\ p_w=g_i}\mathrm{sz}_w(E(w))-\Big(\sum_ta_t\Big)\log g_i\Big|
 \le C_2B_0.
\]
As $E$ vanishes off the support, the sum may be taken over $S_i$. By input (a), $g_i\neq T_u$, so
for $w\in S_i$
\[
 \mathrm{sz}_w(E(w))=2\big(v-\tfrac12F(j)+c_{\mathrm{er}}\big)\ell_we_w
 +\mathrm{sz}_w\big(\nu^{(i)}_{u,j}(w)\big).
\]
Sum over $S_i$, using $\sum_{w\in S_i}\ell_we_w=\log g_i$ and input (b):
\[
 \sum_{w\in S_i}\mathrm{sz}_w(E(w))-\Big(\sum_ta_t\Big)\log g_i
 =\Big(\big(n-F(j)-\textstyle\sum_ta_t\big)+2v+2c_{\mathrm{er}}\Big)\log g_i.
\]
Here $v\geq0$, $c_{\mathrm{er}}\geq0$ because $c_g\geq10$ and $\eta>0$, and $\log g_i\geq0$. By
input (c) the absolute value is at most $\big(\mathrm{err}(j)+2v+2c_{\mathrm{er}}\big)\log g_i$,
which is at most $C_2B_0$ by input (d).

\emph{Step 2: the other places.} We show
\[
 \sum_{w\in\operatorname{supp},\ p_w\notin\{g_0,g_1\}}|\mathrm{sz}_w(E(w))|
 \le(C_3+C_4+C_5+C_6)B_0.
\]
Split the places of this sum into the set $S_T$ of those over $T_u$ and the set $S'$ of the others.
For $w\in S_T$, $E(w)=E_T(u,w)\geq0$, so $|\mathrm{sz}_w(E(w))|=\mathrm{sz}_w(E_T(u,w))$, and
input (e) bounds the sum over $S_T$ by $E^{\pi}_T(u)\log T_u\le C_3B_0$.

For a prime $p$ put
\[
 h(p)=v_p(D_\Gamma(u))+\delta_p\,|E_d(j)|+\delta'_p\,|E_{\mathrm{in}}(j,p)|,
\]
where $\delta_p$ is $1$ if $p=d$ and $0$ otherwise, and $\delta'_p$ is $1$ if $p$ is inert for $D$
and $p\le X''$, and $0$ otherwise. For $w\in S'$ the prime $p_w$ is not $T_u$, $g_0$ or $g_1$, so
the explicit errors give $|E(w)|\le h(p_w)$. Let $P$ be the set of primes below the places of $S'$.
Since $|\mathrm{sz}_w(y)|=2\ell_we_w|y|$ and $h\geq0$, grouping the places of $S'$ by their prime
and using $\sum_{w\mid p}\ell_we_w=\log p$ gives
\[
 \sum_{w\in S'}|\mathrm{sz}_w(E(w))|\le2\sum_{p\in P}\log p\;h(p).
\]
The product of the $p^{v_p(D_\Gamma(u))}$, $p\in P$, divides the positive integer $D_\Gamma(u)$, so
$\sum_{p\in P}\log p\,v_p(D_\Gamma(u))\le\log D_\Gamma(u)$. Only $p=d$ contributes to the second
term, so $\sum_{p\in P}\log p\,\delta_p|E_d(j)|\le\log d\,|E_d(j)|$. Input (h), applied to the set of
$p\in P$ with $\delta'_p=1$, gives $\sum_{p\in P}2\log p\,\delta'_p|E_{\mathrm{in}}(j,p)|\le C_6B_0$.
With inputs (f) and (g), the sum over $S'$ is at most $(C_4+C_5+C_6)B_0$.

\emph{Step 3.} Since $g_0\neq g_1$, $\log(g_0g_1)=\log g_0+\log g_1$. Split the sum over the
support into the places over $g_0$, the places over $g_1$ and the rest. By Steps 1 and 2 and the
triangle inequality, the left side of the first inequality is at most
$(2C_2+C_3+C_4+C_5+C_6)B_0=CB_0$. The second inequality is Step 2, since $2C_2B_0\geq0$.
\end{proof}

\begin{definition}\label{asm-division-kernel-section}
Let $D$ be a definite-side datum (Definition~\ref{th-definite-datum}) with top forms $\mathcal E$,
words $u$ with spaces $\mathcal M_u$, periods $\operatorname{per}_{u,\mu}$, Bessel pairings $B_u$ and
base tuple $\mu_0$. For a tuple $\mu$ let $e_\mu$ be the idempotent of the rows data of $D$ on the
ambient space $\mathcal A_\mu$, whose image is the base space $V_0(\mu)$, and let
$\iota_\mu\colon\mathcal A_\mu\to\mathcal E$ be the map of the rows data.
\begin{enumerate}
\item Put $V_0^{0}(\mu)=V_0(\mu)\cap\ker\iota_\mu$ and $V_0^{\perp}(\mu)=V_0(\mu)\cap
  V_0^{0}(\mu)^{\perp}$. The map $\iota_\mu$ is injective on $V_0^{\perp}(\mu)$, since
  $V_0^{\perp}(\mu)\cap\ker\iota_\mu\subset V_0^{0}(\mu)\cap V_0^{0}(\mu)^{\perp}=0$. The
  \emph{section} $\operatorname{sec}_\mu\colon\mathcal E\to\mathcal A_\mu$ is the orthogonal
  projection of $\mathcal E$ onto $\iota_\mu(V_0^{\perp}(\mu))$ followed by the inverse of the
  bijection $V_0^{\perp}(\mu)\to\iota_\mu(V_0^{\perp}(\mu))$ given by $\iota_\mu$. So
  $\operatorname{sec}_\mu(\iota_\mu q)=q$ for every $q\in V_0^{\perp}(\mu)$.
\item $D$ satisfies the \emph{Bessel tie} if
  $B_u(\varphi,x)=\operatorname{per}_{u,\mu_0}\big(\varphi,\operatorname{sec}_{\mu_0}(x)\big)$ for
  every word $u$, every $\varphi\in\mathcal M_u$ and every $x\in\mathcal E$.
\end{enumerate}
\end{definition}
\leanref{Formalism/LocalRingCardinality/DataRows.lean}{Formalism.BasicNotions.Rows.Data.Usp}
\leanref{Formalism/LocalRingCardinality/DataRows.lean}{Formalism.BasicNotions.Rows.Data.Qsp}
\leanref{Formalism/LocalRingCardinality/DataBu.lean}{Formalism.BasicNotions.Rows.Data.sec}
\leanref{Formalism/LocalRingCardinality/DataBu.lean}{Formalism.BasicNotions.Rows.Data.sec_iota}
\leanref{Formalism/LocalRingCardinality/DataBu.lean}{Formalism.BasicNotions.Rows.BuDefQ}
\leanref{Formalism/LocalRingCardinality/DataTies.lean}{Formalism.BasicNotions.Ties}

\begin{lemma}\label{asm-division-kernel-descent}
Let $D$ be a definite-side datum (Definition~\ref{th-definite-datum}) with coefficient field $L$, its
embedding $\sigma_0\colon L\to\mathbf C$, top forms $\mathcal E$, base tuple $\mu_0$, base space
$V_0(\mu_0)\subset\mathcal A_{\mu_0}$, map $\iota_{\mu_0}$ and section
$\operatorname{sec}=\operatorname{sec}_{\mu_0}$ (Definition~\ref{asm-division-kernel-section}). Let
$V_R$ be the $L$-span of its division lattice $R$, and $j_R\colon V_R\to\mathcal E$ the
$\sigma_0$-semilinear map of its division data, which carries a basis of $V_R$ to a basis of
$\mathcal E$ (Definition~\ref{th-bessel-identity}). Let $u$ be a word and $\varphi\in\mathcal M_u$.
Let $n,d\ge0$, let $a\colon L^n\to\mathcal A_{\mu_0}$ be a map, $e\colon L^d\to V_R$ an $L$-linear
isomorphism, $p\colon L^n\to L^d$ a surjective map and $\beta\colon L^d\to L$ an $L$-linear form,
such that:
\begin{enumerate}
\item $V_0(\mu_0)$ is the complex span of the vectors $a(y)$, $y\in L^n$;
\item $\iota_{\mu_0}(a(y))=j_R(e(p(y)))$ for every $y\in L^n$;
\item $\operatorname{per}_{u,\mu_0}(\varphi,a(y))=\sigma_0(\beta(p(y)))$ for every $y\in L^n$.
\end{enumerate}
Then $\operatorname{per}_{u,\mu_0}\big(\varphi,\operatorname{sec}(j_R(e(z)))\big)=\sigma_0(\beta(z))$
for every $z\in L^d$.
\end{lemma}
\leanref{Formalism/LocalToGlobal/Divisibility/DivisibilityRemainder/RationalColumns/Descent.lean}{Formalism.L2.DIV.PerDescent.per_sec}
\leanref{Formalism/LocalToGlobal/Divisibility/DivisibilityRemainder/RationalColumns/Descent.lean}{Formalism.L2.DIV.PerDescent.exists_ext}
\begin{proof}
The period $\operatorname{per}_{u,\mu_0}(\varphi,\cdot)$ is a linear form on $\mathcal A_{\mu_0}$,
being the period of an iterated lift (Definition~\ref{th-definite-datum}(5)). Let $b_1,\dots,b_r$
be a basis of $V_R$ whose image under $j_R$ is a basis of $\mathcal E$, and let $\Lambda$ be the linear
form on $\mathcal E$ with $\Lambda(j_R(b_i))=\sigma_0(\beta(e^{-1}(b_i)))$. Since $j_R$ is
$\sigma_0$-semilinear, $\Lambda(j_R(v))=\sigma_0(\beta(e^{-1}(v)))$ for every $v\in V_R$.

By (2) and (3), the linear forms $\operatorname{per}_{u,\mu_0}(\varphi,\cdot)$ and
$\Lambda\circ\iota_{\mu_0}$ agree at every vector $a(y)$, so by (1) they agree on $V_0(\mu_0)$.

Let $z\in L^d$, and choose $y$ with $p(y)=z$. Let $h$ be the orthogonal projection of $a(y)$ onto
$V_0^{0}(\mu_0)$, and put $q=a(y)-h$. Then $q\in V_0(\mu_0)$, since $a(y)\in V_0(\mu_0)$ by (1),
and $q$ is orthogonal to $V_0^{0}(\mu_0)$, so $q\in V_0^{\perp}(\mu_0)$. As $\iota_{\mu_0}(h)=0$,
(2) gives $\iota_{\mu_0}(q)=j_R(e(z))$. Hence $\operatorname{sec}(j_R(e(z)))=q$, and
\[
 \operatorname{per}_{u,\mu_0}(\varphi,q)=\Lambda(\iota_{\mu_0}q)=\Lambda\big(j_R(e(z))\big)
 =\sigma_0(\beta(z)).
\]
\end{proof}

\begin{definition}\label{asm-division-kernel-multiplicity}
Let $R$ be a realization of a definite-side datum for a family of definite forms over $L$, with base
place $v_0$ and its embedding $\sigma_0$, and let $F$ be a face of it with the subgroup $P_1$ and the
character $\psi_1$ at level $1$ (Definition~\ref{is-realisation}). Let $X$ be the complex data of $R$
and $F$ (Definition~\ref{is-complex-data}), with theta lifts $\Theta_t\colon F_t\to E$, maps $j_t$,
level $E^1$, number $n_1$, maps $i_1(t)$ and the map $w\mapsto\mu_w$, and let $S_1$ be its columns,
with matrices $P_t$ and $\iota_t$, rank $d_Q(\mu_0)$ and metric $X_E$. Let $D_c$ be the definite
columns of the face, whose members $u$ are called words as in
Definition~\ref{is-division-package}(3), with integers $a(u)$, metrics $X_c(u)$ and period pairings
$\mathrm{Pt}_u(t)$. The \emph{periods} of $D_c$ consist of a Hilbert space $H$, linear maps
$P_u\colon\mathbf C^{a(u)}\to H$ and a subspace $H_r\subset H$. Let $\Gamma_1(t)$ be the space
$\Gamma(t)$ of Lemma~\ref{asm-primitive-rows-parts} for $(P_1,\psi_1)$, and for $y\in L^{n_V(t)}$
let $e_t(y)$ be the row form in $\Gamma_1(t)$ with coordinates $y$ in the basis of $\Gamma_1(t)$ of
the definite side at level $1$. Let $s_E\colon E\to H$ be a conjugate-linear map.
\begin{enumerate}
\item The \emph{isometry of the row space} holds for $s_E$ if $n_1\|s_E(x)\|^2=\|x\|^2$ and
  $s_E(x)\in H_r$ for every $x\in E^1$, and $\dim H_r=\dim E^1$.
\item The \emph{multiplicity condition} holds for $s_E$ if there are a linear map $V$ of $E$, and for
  the words $u$ and the tuples $t$ numbers $a_u\in\mathbf C$, complex matrices $U_u$ of size $a(u)$
  and numbers $\zeta^C_u(t)\in\mathbf C$, such that:
  \begin{enumerate}
  \item $V$ maps $E^1$ into $E^1$ and preserves norms on $E^1$, each $U_u$ preserves the norm
    $|\cdot|_{X_c(u),v_0}$, and $a_u\ne0$;
  \item $\zeta^C_u(\mu_0)=1$, $\zeta^C_u(t)\ne0$, and $v_w(k)=0$ for every finite place $w$ of $L$
    and every $k\in L$ with $\sigma_0(k)=\zeta^C_u(\mu_w)$;
  \item for every word $u$, tuple $t$, $\varphi\in L^{a(u)}$ and $y\in L^{n_V(t)}$,
    \[
     \zeta^C_u(t)\,\sigma_0\big(\mathrm{Pt}_u(t)(\varphi,y)\big)
     =a_u\big\langle s_E\big(V\Theta_tj_te_t(y)\big),P_u\big(U_u\sigma_0(\varphi)\big)\big\rangle.
    \]
  \end{enumerate}
\item Let $P^{+}$ be a right inverse of $P_{\mu_0}$, chosen once and for all ($0$ if there is
  none). The \emph{column data} of
  $R$ at $F$ for $s_E$ consist of the pairings $B_u(\varphi)(y)=\mathrm{Pt}_u(\mu_0)(\varphi,P^{+}y)$,
  the map $s_E\circ i_1(\mu_0)$, and scalars $\zeta_u(t)\in L$ chosen, when such a choice exists,
  so that they satisfy the factorization of Definition~\ref{is-division-package}(3) and the column
  data are \emph{isometric}:
  \begin{enumerate}
  \item[(i)] $\|s_E(i_1(\mu_0)y)\|=|y|_{X_E,v_0}$ and $s_E(i_1(\mu_0)y)\in H_r$ for
    $y\in\mathbf C^{d_Q(\mu_0)}$, and $\dim H_r=d_Q(\mu_0)$;
  \item[(ii)] there is a complex matrix $W$ preserving $|\cdot|_{X_E,v_0}$ such that for every word
    $u$ there are $a_u\ne0$ and a complex matrix $U_u$ preserving $|\cdot|_{X_c(u),v_0}$ with
    \[
     \sum_{k,l}x_ky_l\,\sigma_0\big(B_u(e_k)(e_l)\big)
     =a_u\big\langle s_E\big(i_1(\mu_0)(Wy)\big),P_u(U_ux)\big\rangle
     \qquad(x\in\mathbf C^{a(u)},\ y\in\mathbf C^{d_Q(\mu_0)}).
    \]
  \end{enumerate}
\end{enumerate}
In an indefinite realization of a row-side family $P$ (Definition~\ref{is-realisation-join}), the
column data of $P$ are the column data of its realization at its face for its map $s_E$.
\end{definition}
\leanref{Formalism/LocalToGlobal/IndefiniteForms/IndefSide.lean}{Formalism.L2.Indef.SPer}
\leanref{Formalism/LocalToGlobal/IndefiniteReal/ZeroPartEigenvalueConvention.lean}{Formalism.IndefR.Links.Face3.eV}
\leanref{Formalism/LocalToGlobal/IndefiniteReal/IsometricPeriodsOfHypotheses.lean}{Formalism.IndefR.Cx.CxData.SEEIso}
\leanref{Formalism/LocalToGlobal/IndefiniteReal/IsometricPeriodsOfHypotheses.lean}{Formalism.IndefR.Cx.CxData.MIDat}
\leanref{Formalism/LocalToGlobal/IndefiniteReal/IsometricPeriodsOfHypotheses.lean}{Formalism.IndefR.Cx.CxData.MIid}
\leanref{Formalism/LocalToGlobal/IndefiniteReal/IsometricPeriodsOfUnitHypotheses.lean}{Formalism.IndefR.Cx.CxData.MIUnit2}
\leanref{Formalism/LocalToGlobal/IndefiniteReal/IsometricPeriodsOfUnitHypotheses.lean}{Formalism.IndefR.Cx.CxData.MIstar2}
\leanref{Formalism/LocalToGlobal/IndefiniteReal/CutTuples2.lean}{Formalism.IndefR.Cx.CxData.Bdef}
\leanref{Formalism/LocalToGlobal/IndefiniteReal/CutTuples2.lean}{Formalism.IndefR.Cx.CxData.colDef}
\leanref{Formalism/LocalToGlobal/IndefiniteReal/PartsOrth.lean}{Formalism.IndefR.Cx.CxData.colOf}
\leanref{Formalism/LocalToGlobal/IndefiniteReal/PackageFromPrimitiveStatements.lean}{Formalism.IndefR.Cx.zetaEps}
\leanref{Formalism/LocalToGlobal/IndefiniteForms/IndefSide.lean}{Formalism.L2.Indef.ColData.L1e}

\begin{lemma}\label{asm-division-kernel-evaluation}
Let $D$, the family of definite forms, $R$, $F$ and the complex data $X$ be as in
Lemma~\ref{asm-primitive-rows-base}, with its hypotheses, and let $\sigma_0$, $\Theta_t$, $j_t$,
$E^1$, $i_1(t)$, $\Gamma_1(t)$ and $e_t$ be as in Definition~\ref{asm-division-kernel-multiplicity}.
Let $c_t$, $Z_t$ and $P_t$ be the coordinate maps, the definite subspaces and the matrices of the
definite side at level $1$, and $\iota_t$ the matrices of the columns $S_1$. Then for every tuple
$t$ and every $y\in L^{n_V(t)}$,
\[
 \Theta_tj_te_t(y)=\Theta_tj_tc_t(P_ty)=i_1(\mu_0)\big(\sigma_0(\iota_tP_ty)\big).
\]
\end{lemma}
\leanref{Formalism/LocalToGlobal/IndefiniteReal/IsometricPeriodsOfHypotheses.lean}{Formalism.IndefR.Cx.CxData.theta_eV}
\leanref{Formalism/LocalToGlobal/IndefiniteReal/IsometricPeriodsOfHypotheses.lean}{Formalism.IndefR.Cx.CxData.iC_cpt'}
\leanref{Formalism/LocalToGlobal/IndefiniteReal/CxLemmas.lean}{Formalism.IndefR.Cx.CxData.key_row}
\leanref{Formalism/LocalToGlobal/IndefiniteReal/EvTieP.lean}{Formalism.IndefR.Links.EvTie.evTie_of_ok}
\begin{proof}
By the construction of $P_t$ (proof of Lemma~\ref{asm-primitive-rows-parts}), $c_t(P_ty)$ is the
orthogonal projection of $e_t(y)$ onto $Z_t\cap\Gamma_1(t)$, the image of $c_t$, for the class-sum
pairing, which is positive definite at $v_0$ by the family conditions. So
$d=e_t(y)-c_t(P_ty)$ lies in $\Gamma_1(t)$ and is orthogonal to $Z_t\cap\Gamma_1(t)$.

Let $\Gamma$ be the complex span of $j_t(\Gamma_1(t))$. Step 2 of the proof of
Lemma~\ref{asm-primitive-rows-base}, at level $1$, shows that $\Theta_t(\Gamma)\subset E^1$, that
$\Theta_t^*(E^1)\subset\Gamma$, and that $\Theta_t^*(E)\cap\Gamma$ is spanned by
$j_t(Z_t\cap\Gamma_1(t))$. So $\Theta_tj_t(d)\in E^1$. Let $x\in E^1$. Then $\Theta_t^*x$ is a
combination of vectors $j_t(z)$ with $z\in Z_t\cap\Gamma_1(t)$, and $\langle j_t(z),j_t(d)\rangle$ is
the image under $\sigma_0$ of the class-sum pairing of $z$ and $d$ (Step 0 there), which is $0$. So
$\langle x,\Theta_tj_t(d)\rangle=\langle\Theta_t^*x,j_t(d)\rangle=0$. Taking $x=\Theta_tj_t(d)$
gives $\Theta_tj_t(d)=0$, which is the first equality.

For $x\in L^{d_Q(t)}$, $\Theta_tj_tc_t(x)=i_1(t)(\sigma_0(x))$, since $j_t\circ c_t$ is
$\sigma_0$-semilinear and $i_1(t)$ sends the $k$-th basis vector to $\Theta_tj_tc_t(e_k)$. By
Lemma~\ref{asm-primitive-rows-base}(2), $i_1(t)(\sigma_0(x))=i_1(\mu_0)(\sigma_0(\iota_tx))$. With
$x=P_ty$ this is the second equality.
\end{proof}

\begin{lemma}\label{asm-division-kernel-factorization}
Let $D$, the family of definite forms, $R$, $F$ and $X$ be as in
Lemma~\ref{asm-division-kernel-evaluation}, with its notation, and let the definite columns of the
face have period pairings $\mathrm{Pt}_u(t)$ and periods $(H,P_u,H_r)$. Let $s_E\colon E\to H$ be a
conjugate-linear map for which the multiplicity condition holds
(Definition~\ref{asm-division-kernel-multiplicity}), and let $B_u$ be the pairings of the column
data of $R$ at $F$. Then:
\begin{enumerate}
\item there are $\zeta_u(t)\in L$, for the words $u$ and the tuples $t$, such that for all $u$,
  $t$, $\varphi\in L^{a(u)}$, $y\in L^{n_V(t)}$ and finite places $w$ of $L$,
  \[
   \mathrm{Pt}_u(t)(\varphi,y)=\zeta_u(t)\,B_u(\varphi)(\iota_tP_ty),\qquad\zeta_u(\mu_0)=1,\qquad
   \zeta_u(t)\ne0,\qquad v_w(\zeta_u(\mu_w))=0;
  \]
\item if the isometry of the row space also holds for $s_E$, then the column data are isometric,
  whatever their scalars, and the column data of $R$ at $F$ for $s_E$ satisfy the factorization of
  Definition~\ref{is-division-package}(3).
\end{enumerate}
\end{lemma}
\leanref{Formalism/LocalToGlobal/IndefiniteReal/IsometricPeriodsOfUnitHypotheses.lean}{Formalism.IndefR.Cx.CxData.exists_zeta_pairing_factors_through_periods}
\leanref{Formalism/LocalToGlobal/IndefiniteReal/IsometricPeriodsOfUnitHypotheses.lean}{Formalism.IndefR.Cx.CxData.btVanish_of_miStar2}
\leanref{Formalism/LocalToGlobal/IndefiniteReal/IsometricPeriodsOfUnitHypotheses.lean}{Formalism.IndefR.Cx.CxData.embedding_pairing_eq_mul_inner}
\leanref{Formalism/LocalToGlobal/IndefiniteReal/IsometricPeriodsOfUnitHypotheses.lean}{Formalism.IndefR.Cx.CxData.base_tuple_vanishing_of_hypotheses}
\leanref{Formalism/LocalToGlobal/IndefiniteReal/IsometricPeriodsOfUnitHypotheses.lean}{Formalism.IndefR.Cx.CxData.zetaC_mul_embedding_Pt}
\leanref{Formalism/LocalToGlobal/IndefiniteReal/IsometricPeriodsOfUnitHypotheses.lean}{Formalism.IndefR.Cx.CxData.matrix_identity_of_unit_hypotheses}
\leanref{Formalism/LocalToGlobal/IndefiniteReal/IsometricPeriodsOfUnitHypotheses.lean}{Formalism.IndefR.Cx.CxData.zetaOK_of_miUnit2}
\leanref{Formalism/LocalToGlobal/IndefiniteReal/IsometricPeriodsOfUnitHypotheses.lean}{Formalism.IndefR.Cx.CxData.isometric_periods_of_unit_hypotheses}
\leanref{Formalism/LocalToGlobal/IndefiniteReal/CutPassage.lean}{Formalism.IndefR.Cx.CxData.cutPairing}
\begin{proof}
Fix multiplicity data $V$, $a_u$, $U_u$ and $\zeta^C_u(t)$, and write $i_1=i_1(\mu_0)$. By
Lemma~\ref{asm-primitive-rows-base}(2), $i_1$ is injective with image $E^1$ and $\iota_{\mu_0}=1$.
By Lemma~\ref{asm-primitive-rows-parts}(1), $P_{\mu_0}$ is surjective, so $P_{\mu_0}P^{+}=1$.

\emph{Step 1: the pairings $B_u$.} Let $y\in L^{d_Q(\mu_0)}$. Condition (c) at $t=\mu_0$ and
$P^{+}y$, with $\zeta^C_u(\mu_0)=1$, and Lemma~\ref{asm-division-kernel-evaluation}, with
$\iota_{\mu_0}=1$ and $P_{\mu_0}P^{+}y=y$, give
\[
 \sigma_0\big(B_u(\varphi)(y)\big)
 =a_u\big\langle s_E\big(V i_1(\sigma_0(y))\big),P_u\big(U_u\sigma_0(\varphi)\big)\big\rangle.
 \tag{$*$}
\]

\emph{Step 2: the base tuple.} If $P_{\mu_0}y=0$, then $\Theta_{\mu_0}j_{\mu_0}e_{\mu_0}(y)=0$ by
Lemma~\ref{asm-division-kernel-evaluation}, and (c) at $t=\mu_0$ gives
$\mathrm{Pt}_u(\mu_0)(\varphi,y)=0$.

\emph{Step 3: every tuple.} For every tuple $t$, (c), Lemma~\ref{asm-division-kernel-evaluation} and
($*$) for the vector $\iota_tP_ty$ give
\[
 \zeta^C_u(t)\,\sigma_0\big(\mathrm{Pt}_u(t)(\varphi,y)\big)
 =\sigma_0\big(B_u(\varphi)(\iota_tP_ty)\big).
\]
Put $\xi_u(t)=k$ if $\zeta^C_u(t)=\sigma_0(k)$ with $k\in L$, and $\xi_u(t)=1$ if
$\zeta^C_u(t)\notin\sigma_0(L)$. Then $\xi_u(t)\,\mathrm{Pt}_u(t)(\varphi,y)=B_u(\varphi)(\iota_tP_ty)$.
In the first case this holds because $\sigma_0$ is injective. In the second,
$\mathrm{Pt}_u(t)(\varphi,y)=0$, since otherwise
$\zeta^C_u(t)=\sigma_0\big(B_u(\varphi)(\iota_tP_ty)/\mathrm{Pt}_u(t)(\varphi,y)\big)$; then both
sides vanish. By (b), $\xi_u(\mu_0)=1$, $\xi_u(t)\ne0$ and $v_w(\xi_u(\mu_w))=0$ for every finite
place $w$.

\emph{Step 4: part (1).} Put $\zeta_u(t)=\xi_u(t)^{-1}$. Then $\zeta_u(\mu_0)=1$, $\zeta_u(t)\ne0$
and $v_w(\zeta_u(\mu_w))=-v_w(\xi_u(\mu_w))=0$. For $t\ne\mu_0$ the identity of part (1) is Step 3.
For $t=\mu_0$, $\iota_{\mu_0}=1$, and $y-P^{+}P_{\mu_0}y$ lies in the kernel of $P_{\mu_0}$, so by
Step 2,
\[
 B_u(\varphi)(P_{\mu_0}y)=\mathrm{Pt}_u(\mu_0)(\varphi,P^{+}P_{\mu_0}y)=\mathrm{Pt}_u(\mu_0)(\varphi,y).
\]

\emph{Step 5: part (2).} The metric $X_E$ of $S_1$ is at $v_0$ a metric for the Gram matrix
$n_1^{-1}\big(\langle i_1e_k,i_1e_l\rangle\big)_{k,l}$ (Definition~\ref{is-complex-data}(b)), so
$|y|_{X_E,v_0}=n_1^{-1/2}\|i_1(y)\|$ for $y\in\mathbf C^{d_Q(\mu_0)}$. By the isometry of the row
space, $\|s_E(i_1y)\|=n_1^{-1/2}\|i_1y\|=|y|_{X_E,v_0}$ and $s_E(i_1y)\in H_r$, and
$\dim H_r=\dim E^1=d_Q(\mu_0)$. This is (i). As $V$ maps $E^1$ into itself, there is a linear map
$W$ of $\mathbf C^{d_Q(\mu_0)}$ with $i_1(Wy)=V(i_1y)$, and $|Wy|_{X_E,v_0}=|y|_{X_E,v_0}$ because
$V$ preserves norms on $E^1$. The two forms
\[
 (x,y)\mapsto\sum_{k,l}x_ky_l\,\sigma_0\big(B_u(e_k)(e_l)\big),\qquad
 (x,y)\mapsto a_u\big\langle s_E\big(i_1(Wy)\big),P_u(U_ux)\big\rangle
\]
on $\mathbf C^{a(u)}\times\mathbf C^{d_Q(\mu_0)}$ are complex bilinear, the second because $s_E$ is
conjugate-linear and the inner product is conjugate-linear in its first variable, and they agree at
pairs of basis vectors by ($*$). So they agree, and with $a_u\ne0$ and the matrices $U_u$ this is
(ii). The scalars of the column data of $R$ at $F$ are chosen, when such a choice exists, so that
the factorization and the isometry hold; by part (1) and (i), (ii) the choice exists.
\end{proof}

\begin{lemma}\label{asm-division-kernel-rational}
Let $D$ be a definite-side datum (Definition~\ref{th-definite-datum}) with coefficient field $L$, its
embedding $\sigma_0$, top forms $\mathcal E$, base tuple $\mu_0$, Bessel pairings $B_u$, division
lattices $N_u$ and $R$ with the maps $j_N$ and $j_R$, and division pairings $B^{\mathrm{div}}_u$
(Definition~\ref{th-bessel-identity}), and suppose that $D$ satisfies the Bessel tie
(Definition~\ref{asm-division-kernel-section}). Let $P$ be a row-side family for $D$
(Definition~\ref{is-row-side-family}) with base tuple $t_0$, columns $S_1$ on the definite side of
$(P_s,\psi_s)$ with matrices $P_t$ and $\iota_t$, definite columns with period pairings
$\mathrm{Pt}_w(t)$, column data with pairings $B_w$, and identifications $e_{N,u}$, $e_{RR}$ and
$u\mapsto w_u$. Its identifications are tied to $D$; among the ties are (1) to (3) of
Lemma~\ref{asm-dependence-base-vectors}, with the maps $j_A$, $E$ and $V$ there, and $j_A$ carries the
standard basis to a basis of $\mathcal A_{\mu_0}$. Let $Q$ be an indefinite realization of $P$
(Definition~\ref{is-realisation-join}), with realization $R_Q$, face $F_Q$ and map $s_E$. Suppose
that:
\begin{enumerate}
\item the family of definite forms of $P$ satisfies the family conditions with the conditions on
  its inert indices (Lemma~\ref{is-eigenvalue-base-inputs}(1)) and has lift-distinct types, and
  $R_Q$ has primitive rows at $F_Q$ (Definition~\ref{is-primitive-rows});
\item the isometry of the row space and the multiplicity condition hold for $R_Q$ at $F_Q$ and $s_E$
  (Definition~\ref{asm-division-kernel-multiplicity});
\item $B^{\mathrm{div}}_u(e_{N,u}x,e_{RR}y)=B_{w_u}(x)(y)$ for every word $u$ and all $x,y$
  (Definition~\ref{is-division-rows}(5));
\item $\sigma_0\big(\mathrm{Pt}_{w_u}(t_0)(x,y)\big)=\operatorname{per}_{u,\mu_0}\big(j_Ne_{N,u}x,\,
  j_A(Vy)\big)$ for every word $u$ and all coordinate vectors $x$ and $y$.
\end{enumerate}
Then the Bessel pairings of $D$ are rational: $B_u(j_N\varphi,j_Rx)=\sigma_0(B^{\mathrm{div}}_u(\varphi,x))$
for every word $u$, every $\varphi$ in the $L$-span of $N_u$ and every $x$ in the $L$-span of $R$.
\end{lemma}
\leanref{Formalism/LocalToGlobal/Divisibility/DivisibilityRemainder/RationalColumns/BuRationalOf.lean}{Formalism.L2.DIV.RatFullPt.buRational_of}
\leanref{Development/AutomorphicArithmetic/ThetaCorrespondence/ThetaPairLifts/ThirdPlaceSeries/BurrbColFacts.lean}{Formalism.L2.DIV.RatFullPt.indefinite_realisation_column_facts}
\leanref{Formalism/LocalToGlobal/Divisibility/DivisibilityRemainder/RationalColumns/ColFacts.lean}{Formalism.L2.DIV.RatFullPt.range_op_eq_span}
\leanref{Formalism/LocalToGlobal/Divisibility/DivisibilityRemainder/RationalColumns/ColFacts.lean}{Formalism.L2.DIV.RatFullPt.lvlOp_range}
\begin{proof}
\emph{Step 1: the columns.} The realization $R_Q$ satisfies the weak characterization at $F_Q$, and
its complex data at $F_Q$ have the definite subspaces, the columns $S_1$, the definite columns,
their periods and the column data of $P$; the level $1$ of $F_Q$ has the subgroup $P_s$ and the
character $\psi_s$ (Definitions~\ref{is-realisation-join}
and~\ref{asm-division-kernel-multiplicity}). By (1) and Lemma~\ref{asm-primitive-rows-base}(2) at
level $1$, the matrices $\iota_t$ are bijective, and by Lemma~\ref{asm-primitive-rows-parts}(1) the
matrices $P_t$ are surjective. By (2) and Lemma~\ref{asm-division-kernel-factorization}(2), the column
data of $P$ satisfy the factorization: $\mathrm{Pt}_w(t)(x,y)=\zeta_w(t)\,B_w(x)(\iota_tP_ty)$ with
$\zeta_w(t_0)=1$.

\emph{Step 2: the base space.} We show that $V_0(\mu_0)$ is the complex span of the vectors
$j_A(Vy)$. Let $\Gamma$ be the image of $e_{\mathcal X}e_{P_s,\psi_s}$ on the row forms at $t_0$, as
in Lemma~\ref{asm-dependence-base-vectors}. As in the proof of that lemma, $EVy=Vy$ for every $y$.
Let $z$ be the coordinate vector of a level form $f$. Then $f$ is fixed by the idempotent
$e_{P_s,\psi_s}$, whose image is the space of level forms
(Lemma~\ref{asm-dependence-projections-onto}), and $e_{\mathcal X}$ is central, so
$e_{\mathcal X}f=e_{\mathcal X}e_{P_s,\psi_s}f$ lies in $\Gamma$; hence $Ez=Vy$ for some $y$. By
tie (2), $e_{\mu_0}(j_A(z))=j_A(Ez)$; the tuple of $t_0$ is $\mu_0$ by tie (1). As $j_A$ carries
the standard basis to a basis of $\mathcal A_{\mu_0}$, the image $V_0(\mu_0)$ of $e_{\mu_0}$ is
spanned by the vectors $j_A(Ez)$, each of which is some $j_A(Vy)$. Conversely
$j_A(Vy)=j_A(EVy)=e_{\mu_0}(j_A(Vy))$ lies in $V_0(\mu_0)$.

\emph{Step 3.} Let $u$ be a word, $\varphi$ in the $L$-span of $N_u$ and $x$ in the $L$-span of
$R$, and write $\varphi=e_{N,u}(x')$ and $x=e_{RR}(z)$. By the Bessel tie and (3),
\[
 B_u(j_N\varphi,j_Rx)=\operatorname{per}_{u,\mu_0}\big(j_N\varphi,\operatorname{sec}_{\mu_0}(j_R(e_{RR}(z)))\big),
 \qquad B^{\mathrm{div}}_u(\varphi,x)=B_{w_u}(x')(z).
\]
Apply Lemma~\ref{asm-division-kernel-descent} to $j_N\varphi$, with $a(y)=j_A(Vy)$, $e=e_{RR}$,
$p=\iota_{t_0}P_{t_0}$ and $\beta=B_{w_u}(x')$. The map $p$ is surjective by Step 1. Hypothesis (1)
of that lemma is Step 2, and (2) is tie (3) of Lemma~\ref{asm-dependence-base-vectors}. Hypothesis
(3) holds because, by (4) and Step 1,
\[
 \operatorname{per}_{u,\mu_0}\big(j_N\varphi,j_A(Vy)\big)=\sigma_0\big(\mathrm{Pt}_{w_u}(t_0)(x',y)\big)
 =\sigma_0\big(B_{w_u}(x')(\iota_{t_0}P_{t_0}y)\big).
\]
So $B_u(j_N\varphi,j_Rx)=\sigma_0(B_{w_u}(x')(z))=\sigma_0(B^{\mathrm{div}}_u(\varphi,x))$.
\end{proof}

\begin{lemma}\label{asm-division-kernel-join}
Let $D$ be a definite-side datum (Definition~\ref{th-definite-datum}) with coefficient field $L$, its
embedding $\sigma$, top forms $\mathcal E$, spaces $\mathcal M_u$ with their isomorphisms onto
$M_{u,\mathrm{ram}}$, periods $\operatorname{per}_{u,\mu}$, base tuple $\mu_0$, base space
$V_0(\mu_0)$ and map $\iota_{\mu_0}$. Let $R$ be a realization of $D$ with a face, with the
notation of Definition~\ref{asm-division-kernel-multiplicity}, and let $s_E\colon E\to H$ be a
conjugate-linear map for which the multiplicity condition holds. Suppose that $\sigma_0=c\circ\sigma$,
where $c$ is the identity of $\mathbf C$ or complex conjugation, also applied to vectors entry by
entry. Let an automorphic setting attached to $D$ be given (Definition~\ref{is-wu-decay}), with, for
every word $u$ of $D$, Poincaré data $T$, their column dictionaries, which are functions on
$\mathrm{cls}_0\times Y_0$, and their Poincaré series $\mathcal P(T)$. Let $j$ be a linear isometry
from $\mathcal E$ to the automorphic forms of the setting, and write $\overline f$ for the complex
conjugate of an automorphic form $f$. Let $u\mapsto m(u)$ be a bijection from the words of $D$ to the
words of the definite columns, let $c_\varphi(u)\colon\mathcal M_u\to\mathbf C^{a(m(u))}$ be linear
isomorphisms, $c_F\colon\mathbf C^{n_V(\mu_0)}\to\mathcal A_{\mu_0}$ a linear map and
$e_R\colon\mathcal E\to E$ a $c$-semilinear map, such that:
\begin{enumerate}
\item (Chain.) $\operatorname{per}_{u,\mu_0}\big(c_\varphi(u)^{-1}\sigma(x),c_F(\sigma(y))\big)
  =\sigma\big(\mathrm{Pt}_{m(u)}(\mu_0)(x,y)\big)$ for every word $u$ and all $x,y$ with entries in
  $L$;
\item (Base vectors.) $V_0(\mu_0)\subset c_F\big(\mathbf C^{n_V(\mu_0)}\big)$;
\item (Rows.) $e_R$ maps $\mathcal E$ onto $E^1$, there is $c_0>0$ with $\|e_R(x)\|=c_0\|x\|$ for
  $x\in\mathcal E$, and there is $\lambda\ne0$ with
  $\Theta_{\mu_0}j_{\mu_0}e_{\mu_0}(y)=\lambda\,e_R\big(\iota_{\mu_0}c_F(\sigma(y))\big)$ for
  $y\in L^{n_V(\mu_0)}$;
\item (Columns.) for every word $u$ there is $c_u>0$ with
  $\big|c(c_\varphi(u)\varphi)\big|_{X_c(m(u)),v_0}=c_u\|\varphi\|$ for $\varphi\in\mathcal M_u$;
\item (Pairings.) for every word $u$ there is $\kappa_u\ne0$ such that
  \[
   \big\langle s_E(e_R(x)),P_{m(u)}\big(c(c_\varphi(u)\varphi)\big)\big\rangle
   =c\Big(\kappa_u\big\langle\overline{j(x)},\mathcal P(T)\big\rangle\Big)
  \]
  for all $x\in\mathcal E$, $\varphi\in\mathcal M_u$ and Poincaré data $T$ of $u$ whose column
  dictionary is the image of $\varphi$ in $M_{u,\mathrm{ram}}$.
\end{enumerate}
Then the \emph{kernel join} holds for $j$: there is a linear isometry $V$ of $\mathcal E$ such that
for every word $u$ there are $a_u\ne0$ and a linear isometry $U_u$ of $\mathcal M_u$ with
\[
 \operatorname{per}_{u,\mu_0}(\varphi,f)=a_u\big\langle\overline{j(V\iota_{\mu_0}f)},\mathcal P(T)\big\rangle
\]
for all $\varphi\in\mathcal M_u$, $f\in V_0(\mu_0)$ and Poincaré data $T$ of $u$ whose column
dictionary is the image of $U_u\varphi$ in $M_{u,\mathrm{ram}}$.
\end{lemma}
\leanref{Formalism/LocalToGlobal/SiegelLatticeTheory/FlagStabilizerTensors/AdmissibleOfTransported.lean}{Formalism.L2.SG.JoinSI.kernelJoin_of_multiplicity}
\leanref{Formalism/LocalToGlobal/SiegelLatticeTheory/FlagStabilizerTensors/JoinSI.lean}{Formalism.L2.SG.JoinSI.exists_conj_op}
\leanref{Formalism/LocalToGlobal/SiegelLatticeTheory/FlagStabilizerTensors/JoinSI.lean}{Formalism.L2.SG.JoinSI.JoinAt}
\leanref{Formalism/LocalToGlobal/SiegelLatticeTheory/FlagStabilizerTensors/KerJoinS.lean}{Formalism.L2.SG.KerJoinS.KerJoinS}
\begin{proof}
Let $V'$, $a'_u$, $U'_u$ and $\zeta^C_u(t)$ be multiplicity data for $s_E$.

\emph{Step 1: the isometry $V$.} The map $e_R$ is injective, since it multiplies norms by $c_0>0$,
and it maps $\mathcal E$ onto $E^1$, which $V'$ maps into itself. So for every $x\in\mathcal E$ there
is a unique $V(x)\in\mathcal E$ with $e_R(V(x))=V'(e_R(x))$. The map $V$ is linear, because $e_R$ is
injective and $c$-semilinear and $V'$ is linear. As $V'$ preserves norms on $E^1$,
$c_0\|V(x)\|=\|V'(e_R(x))\|=\|e_R(x)\|=c_0\|x\|$, so $V$ is a linear isometry of $\mathcal E$.

\emph{Step 2: the isometries $U_u$.} Fix a word $u$, write $m=m(u)$, and put
$U_u(\varphi)=c_\varphi(u)^{-1}\big(c\big(U'_m\,c(c_\varphi(u)\varphi)\big)\big)$, so that
$c(c_\varphi(u)U_u\varphi)=U'_m\,c(c_\varphi(u)\varphi)$. The map $U_u$ is linear, as $c$ is applied
twice. By (4) and since $U'_m$ preserves $|\cdot|_{X_c(m),v_0}$,
$c_u\|U_u\varphi\|=|U'_m\,c(c_\varphi(u)\varphi)|_{X_c(m),v_0}=c_u\|\varphi\|$.

\emph{Step 3: a bilinear identity.} For $x\in\mathbf C^{a(m)}$ and $y\in\mathbf C^{n_V(\mu_0)}$ put
\[
 B_1(x,y)=\operatorname{per}_{u,\mu_0}\big(c_\varphi(u)^{-1}x,c_F(y)\big),\qquad
 B_2(x,y)=c\Big(a'_m\lambda\big\langle s_E\big(V'e_R(\iota_{\mu_0}c_F(y))\big),P_m\big(U'_mc(x)\big)\big\rangle\Big).
\]
Both are complex bilinear: $\operatorname{per}_{u,\mu_0}$ is bilinear, and in $B_2$ the maps $c$,
$e_R$, $s_E$ and the first variable of the inner product combine to linear dependence on $x$ and
on $y$. Let $e_k$ and $e_l$ be standard basis vectors, which have entries in $L$ and are fixed by
$\sigma$ and $c$. By (1), $B_1(e_k,e_l)=\sigma(\mathrm{Pt}_m(\mu_0)(e_k,e_l))$. By (3),
$e_R(\iota_{\mu_0}c_F(e_l))=\lambda^{-1}\Theta_{\mu_0}j_{\mu_0}e_{\mu_0}(e_l)$, and as $s_E$ is
conjugate-linear,
\[
 B_2(e_k,e_l)=c\Big(a'_m\big\langle s_E\big(V'\Theta_{\mu_0}j_{\mu_0}e_{\mu_0}(e_l)\big),P_m(U'_me_k)\big\rangle\Big)
 =c\Big(\sigma_0\big(\mathrm{Pt}_m(\mu_0)(e_k,e_l)\big)\Big)=\sigma\big(\mathrm{Pt}_m(\mu_0)(e_k,e_l)\big),
\]
by condition (c) of the multiplicity condition at $\mu_0$, where $\zeta^C_m(\mu_0)=1$, and since
$c\circ\sigma_0=c\circ c\circ\sigma=\sigma$. So $B_1=B_2$.

\emph{Step 4.} Let $\varphi\in\mathcal M_u$, $f\in V_0(\mu_0)$, and $T$ a Poincaré datum of $u$ whose
column dictionary is the image of $U_u\varphi$ in $M_{u,\mathrm{ram}}$. By (2), $f=c_F(y)$ for some
$y$. By Step 3 with $x=c_\varphi(u)\varphi$, then Steps 1 and 2, and then (5) for $V\iota_{\mu_0}f$,
$U_u\varphi$ and $T$,
\begin{align*}
 \operatorname{per}_{u,\mu_0}(\varphi,f)
 &=c\Big(a'_m\lambda\big\langle s_E\big(e_R(V\iota_{\mu_0}f)\big),P_m\big(c(c_\varphi(u)U_u\varphi)\big)\big\rangle\Big)\\
 &=c(a'_m\lambda)\,\kappa_u\big\langle\overline{j(V\iota_{\mu_0}f)},\mathcal P(T)\big\rangle.
\end{align*}
So the kernel join holds with $a_u=c(a'_{m(u)}\lambda)\kappa_u$, which is not $0$.
\end{proof}

\begin{lemma}\label{asm-division-kernel-bessel}
Let $D$ be a definite-side datum (Definition~\ref{th-definite-datum}) with top forms $\mathcal E$,
spaces $\mathcal M_u$ with their isomorphisms onto $M_{u,\mathrm{ram}}$, Bessel pairings $B_u$,
standard Bessel periods $\mathcal B_u$, base tuple $\mu_0$, and the Bessel-period data of
Definition~\ref{th-definite-datum}(8): the conjugate-linear isometry $s\colon\mathcal E\to\mathcal H$,
the spaces $\mathcal C_u$ with their isometries onto $M_u$, and the maps
$P_u\colon\mathcal C_u\to\mathcal H$. Suppose that $D$ satisfies the Bessel tie
(Definition~\ref{asm-division-kernel-section}) and that $\iota_{\mu_0}$ maps $V_0(\mu_0)$ onto
$\mathcal E$. Let an automorphic setting attached to $D$ be given, with Poincaré data of the words,
their column dictionaries and their Poincaré series $\mathcal P(T)$ in a Hilbert space
$\mathcal H_X$ of automorphic forms, and suppose given:
\begin{enumerate}
\item an isomorphism of the Bessel-period data of $D$ with those of the setting, consisting of a
  linear isometry $e_{\mathcal H}\colon\mathcal H\to\mathcal H_X$ and linear isometries $e_u$ from
  $\mathcal C_u$ onto the Poincaré data of $u$, with $e_{\mathcal H}(P_uy)=\mathcal P(e_u(y))$;
\item a linear isometry $j$ from $\mathcal E$ to the automorphic forms of the setting with
  $e_{\mathcal H}(s(x))=\overline{j(x)}$ for every $x\in\mathcal E$;
\item for every word $u$ and every $y\in\mathcal C_u$, the column dictionary of $e_u(y)$ is the image
  of $y$ in $M_u$;
\item the kernel join for $j$ (Lemma~\ref{asm-division-kernel-join}).
\end{enumerate}
Then the Bessel pairings of $D$ agree with the standard Bessel periods: there is a linear isometry
$V$ of $\mathcal E$ such that for every word $u$ there are $a_u\ne0$ and a linear isometry $U_u$ of
$\mathcal M_u$ with $B_u(\varphi,x)=a_u\mathcal B_u(U_u\varphi,Vx)$ for all $\varphi\in\mathcal M_u$
and $x\in\mathcal E$.
\end{lemma}
\leanref{Formalism/LocalToGlobal/SiegelLatticeTheory/FlagStabilizerTensors/KerJoinS.lean}{Formalism.L2.SG.KerJoinS.besselId_of_joins}
\leanref{Formalism/LocalToGlobal/SiegelLatticeTheory/FlagStabilizerTensors/KerJoinS.lean}{Formalism.L2.SG.KerJoinS.exists_Q_iota}
\leanref{Formalism/LocalToGlobal/SiegelLatticeTheory/FlagStabilizerTensors/KerJoinS.lean}{Formalism.L2.SG.KerJoinS.besS_conc}
\leanref{Formalism/LocalToGlobal/SiegelLatticeTheory/FlagStabilizerTensors/KerJoinS.lean}{Formalism.L2.SG.KerJoinS.colJoinS_of_pin}
\leanref{Formalism/LocalToGlobal/SiegelLatticeTheory/FlagStabilizerTensors/AdmissibleOfTransported.lean}{Formalism.L2.SG.JoinSI.adm13_of_colJoin}
\leanref{Development/AutomorphicArithmetic/ThetaCorrespondence/ThetaPairLifts/ThirdPlaceSeries/AdmissibleOfTransported.lean}{Formalism.L2.SG.JoinSI.adm13_of_kernel_join}
\begin{proof}
Let $V$, $a_u$ and $U_u$ be as in the kernel join, and let $u$ be a word, $\varphi\in\mathcal M_u$
and $x\in\mathcal E$. As $\iota_{\mu_0}$ maps $V_0(\mu_0)$ onto $\mathcal E$, $x=\iota_{\mu_0}(f)$
with $f\in V_0(\mu_0)$. Subtracting from $f$ its orthogonal projection onto $V_0^{0}(\mu_0)$, which
lies in the kernel of $\iota_{\mu_0}$, gives $q\in V_0^{\perp}(\mu_0)$ with $\iota_{\mu_0}(q)=x$. By
the Bessel tie, $B_u(\varphi,x)=\operatorname{per}_{u,\mu_0}(\varphi,\operatorname{sec}_{\mu_0}(\iota_{\mu_0}q))
=\operatorname{per}_{u,\mu_0}(\varphi,q)$.

Let $y\in\mathcal C_u$ correspond to the image of $U_u\varphi$ in $M_{u,\mathrm{ram}}\subset M_u$,
and put $T=e_u(y)$. By (3), the column dictionary of $T$ is the image of $U_u\varphi$ in
$M_{u,\mathrm{ram}}$, so the kernel join gives
$\operatorname{per}_{u,\mu_0}(\varphi,q)=a_u\langle\overline{j(Vx)},\mathcal P(T)\rangle$. By (2) and (1),
and as $e_{\mathcal H}$ is an isometry,
\[
 \big\langle\overline{j(Vx)},\mathcal P(T)\big\rangle
 =\big\langle e_{\mathcal H}(s(Vx)),e_{\mathcal H}(P_uy)\big\rangle=\langle s(Vx),P_uy\rangle
 =\mathcal B_u(U_u\varphi,Vx).
\]
So $B_u(\varphi,x)=a_u\mathcal B_u(U_u\varphi,Vx)$.
\end{proof}

\begin{lemma}\label{asm-division-kernel-used}
Let $P$ be a row-side family for a definite-side datum $D$ (Definition~\ref{is-row-side-family}),
with base point $p_0$, rows $S_0$ and tuples $t$ with characters $\mu_t$, and let $Q$ be an
indefinite realization of $P$ (Definition~\ref{is-realisation-join}). Write $D(S_0)$ for the integer
$D$ of the indefinite-side datum $S_0$ (Definition~\ref{is-indefinite-datum}) and $D_0$ for the
dimension of the base rows of $D$. Then:
\begin{enumerate}
\item every tuple $t$ of $P$ is used: $\mu_t$ is the base tuple $\mu_0$ of $D$ or the tuple $\mu_w$
  of a finite place $w$;
\item $D_0=D(S_0)$, and $D(S_0)$ is the dimension $d_Q(\mu_0)$ of the definite side of $P$ at the
  base tuple, for the subgroup $P_0$ with the character $\psi_0$.
\end{enumerate}
\end{lemma}
\leanref{Development/AutomorphicArithmetic/ThetaCorrespondence/ThetaPairLifts/ThirdPlaceSeries/BaseRankEqOfRealisation.lean}{Formalism.IndefR.LevelRealisationRankEquality.tup_all_used_of_realisation}
\leanref{Development/AutomorphicArithmetic/ThetaCorrespondence/ThetaPairLifts/ThirdPlaceSeries/BaseRankEqOfRealisation.lean}{Formalism.IndefR.LevelRealisationRankEquality.indefinite_realisation_base_rank_eq}
\leanref{Formalism/LocalToGlobal/BaseRankEqOfRealisation.lean}{Formalism.IndefR.LevelRealisationRankEquality.base_rank_eq_of_tied_level}
\leanref{Formalism/LocalToGlobal/Divisibility/TupUsed.lean}{Formalism.L2.DIV.tupAllUsed_of_cover}
\begin{proof}
(1) Let $R$ be the realization of $Q$ and $F$ the face of $P$. The weak characterization of $R$ at
$F$, which is part of $Q$, says that every tuple $t$ of $R$ is its base tuple or the tuple that $R$
attaches to some finite place $w$. The rows of the complex data of $R$ at $F$ are $S_0$, so the
tuple that $R$ attaches to $w$ is the tuple of $S_0$ at $w$. The identifications of $P$ carry the
base tuple to $\mu_0$ and the tuple of $S_0$ at $w$ to $\mu_w$ (Definition~\ref{is-row-side-family}).
So $\mu_t$ is $\mu_0$ or $\mu_w$.

(2) The identification $e_{RR,0}$ of $P$ is an isomorphism from $L^{D(S_0)}$ onto the space of the
glued lattice of $p_0$. The tie of $p_0$ to $D$ gives a basis of that space and a basis of the
complex space of base rows of $D$ indexed by the same set, and the latter space has dimension $D_0$.
So $D(S_0)=D_0$. The rows of the complex data of $R$ at $F$ are $S_0$, on the definite subspaces of
$P$, and the integer $D$ of these rows is $d_Q(\mu_0)$ for $(P_0,\psi_0)$
(Definition~\ref{is-complex-data}(a)).
\end{proof}

\subsection{The division package and its inputs}\label{subsection-asm-division-package}

\begin{lemma}\label{asm-division-kernel-package-passage}
Let $P$, $Q$, $R$, $F$, $C$ and $C_H$ satisfy the hypotheses of
Lemma~\ref{asm-division-kernel-package}, with the numerical bounds $N=N(R,F;C,C_H)$, and let $X$ be
the complex data of $R$ at $F$. Then the indefinite-side data of $X$ at $(P_0,\psi_0)$ and at
$(P_s,\psi_s)$, on the definite subspaces of $P$, satisfy the passage of
Definition~\ref{is-division-package}(4) for $N$: if $0\le\delta_e\le\frac13$ and the glued lattice
$RR_0$ of the first has $\mu_{\max}(RR_0)\le H$, then the glued lattice $RR_1$ of the second has
$\mu_{\max}(RR_1)\le H+14\ell_a$ and
$\mu(CC_1)\le\mu(CC_0)+3\delta_e\big(\mu(CC_0)+H+28\ell_a\big)+28\ell_a$, where $CC_i=RR_i^\vee$.
\end{lemma}
\leanref{Formalism/LocalToGlobal/IndefiniteReal/NearConstantMultiplicities.lean}{Formalism.IndefR.Cx.passage3_of}
\leanref{Formalism/LocalToGlobal/IndefiniteReal/CutPassage3.lean}{Formalism.IndefR.CxPass.cutPassage3S}
\leanref{Formalism/LocalToGlobal/IndefiniteReal/CutPassage3S.lean}{Formalism.IndefR.CxPass.CutPassage3S}
\leanref{Formalism/LocalToGlobal/IndefiniteReal/CutPassage3SOfSteps.lean}{Formalism.IndefR.CxPass.In3.passage_of_steps}
\begin{proof}
The passage data of hypothesis (4) of Lemma~\ref{asm-division-kernel-package} are: the rationality
at $v_0$ of the row matrices of $X$ for the trivial subgroup; a finite family of representations
$\tau_s$ of $G$ over $L$, good and with the characters of the family; the condition $M_1$ for
$(P_0,\psi_0)$; a number $\bar n>0$ such that the multiplicity of $(P_s,\psi_s)$ in every $\tau_s$
lies between $(1-\delta_e)\bar n$ and $(1+\delta_e)\bar n$, for the number $\delta_e$ of $N$; and
$\ell_a=\log|G|$.
\openstep{The proof, which compares the glued lattices and metrics of the two levels through the
idempotents of $(P_0,\psi_0)$ and $(P_s,\psi_s)$ and the representations $\tau_s$, is not yet written
out.}
% Lean: passage3_of = CxPass.cutPassage3S = cutPassage3S_of_steps steps3 (CutPassage3SOfSteps.lean),
% that is In3.passage_of_steps with the ten steps of steps3: pRho, pJ, pRange, pMet_of, pFrames,
% pUni_of, pLat_of, pGLat_of, pSand_of and pGood_of (RhoSpec.lean, MajorantAtBaseEqInnerOfPackage.lean,
% RowGSurj.lean, PGoodOfRho.lean). passage_of_steps builds a representation of G on L^D from pRho,
% identifies the level idempotents with epP for characters chi0, chi1 lifting psi0, psi1 (epM_eq,
% liftChar), applies LatticeTransfer.levelAmb_transfer to the frames, lattices, G-stability and
% unitarity, and concludes with passage_from_levels. Size: proof closure of passage_of_steps and the
% ten steps, about 15,500 lines in 227 files (text count).
\end{proof}

\begin{situation}\label{asm-division-kernel-package-situation}
Let $K$ be a number field with an involution, and let complex data on a definite ambient over $K$
with a face be given (Definition~\ref{is-complex-data}), with infinite place $v_0$ and its embedding
$\sigma_0$, tuples $t$ with base tuple $\mu_0$, parts $p\in\Pi$, lifts $j$, parts $E_p\subset E$
and $W_p(j)\subset W_j$, positive numbers $n_b$ and scalars $\kappa_v$ and $\lambda_{j,t,p,v}$. Fix
a level $b$. Let $S$ be the indefinite-side datum at level $b$, with $D=d_Q(\mu_0)$, the maps
$i_b(t)$, $t_b(j,t)$ and $w_b(j)$, the matrices $\iota_t$ and $\theta_{j,t}$, the metrics $X_E$,
$X_W(j)$, $X_Q(t)$ and $X_T(j,t)$, the parts $Q_p(t)$ of the definite side, and the parts
$E^K_p=\{y\in K^D:i_b(\mu_0)(\sigma_0y)\in E_p\}$ and $W^K_p(j)$ of $S$. For a metric $X$ on $K^d$
and an infinite place $v$ write $h_{X,v}(y,z)=(X_v^{-1}\sigma_v(y))^*X_v^{-1}\sigma_v(z)$ for its
Hermitian form at $v$.
\end{situation}
\leanref{Formalism/LocalToGlobal/IndefiniteReal/PartsOrth.lean}{Formalism.IndefR.Cx.CxData.data}

\begin{lemma}\label{asm-division-kernel-package-rows}
In Situation~\ref{asm-division-kernel-package-situation}, suppose that:
\begin{enumerate}
\item $\kappa_v>0$ and $\lambda_{j,t,p,v}>0$ for all lifts $j$, tuples $t$, parts $p$ and infinite
  places $v$;
\item $X_Q(\mu_0)$ is invertible at every infinite place, and the $Q_p(\mu_0)$ are independent and
  span $K^D$;
\item $i_b(\mu_0)$ is injective, and for every tuple $t$ the map $i_b(t)$ is injective with the same
  image as $i_b(\mu_0)$, and the matrix that expresses $i_b(t)$ through $i_b(\mu_0)$ has entries in
  $\sigma_0(K)$;
\item $Q_p(\mu_0)=E^K_p$ for every part $p$;
\item for every tuple $t$, part $p$ and $x,y\in K^{d_Q(t)}$ with $i_b(t)(\sigma_0x)\in E_p$,
  \[
   \langle i_b(t)(\sigma_0x),i_b(t)(\sigma_0y)\rangle
   =n_b\kappa_{v_0}\prod_j\lambda_{j,t,p,v_0}\,h_{X_Q(t),v_0}(x,y);
  \]
\item at every infinite place $v\ne v_0$, $h_{X_Q(\mu_0),v}(x,y)=0$ whenever $x\in Q_p(\mu_0)$ and
  $y\in Q_q(\mu_0)$ with $p\ne q$.
\end{enumerate}
Then $X_E$ is invertible at every infinite place, and whenever $t=\mu_0$ or $v=v_0$,
\[
 h_{X_E,v}(\iota_tx,\iota_ty)=\kappa_v\prod_j\lambda_{j,t,p,v}\,h_{X_Q(t),v}(x,y)
\]
for every part $p$, every $x\in\iota_t^{-1}(E^K_p)$ and every $y\in K^{d_Q(t)}$.
\end{lemma}
\leanref{Formalism/LocalToGlobal/IndefiniteReal/CxCoreRG.lean}{Formalism.IndefR.Cx.CxData.XE_specR}
\leanref{Formalism/LocalToGlobal/IndefiniteReal/CxCoreRG.lean}{Formalism.IndefR.Cx.CxData.scal_rowsR}
\leanref{Formalism/LocalToGlobal/Divisibility/RowLatLemmas.lean}{Formalism.L2.DIV.rowXE_spec}
\leanref{Formalism/LocalToGlobal/IndefiniteReal/CxLemmas.lean}{Formalism.IndefR.Cx.rowXE_part}
\leanref{Formalism/LocalToGlobal/IndefiniteReal/CxLemmas.lean}{Formalism.IndefR.Cx.CxData.key_row}
\begin{proof}
By Definition~\ref{is-complex-data}(a), $\sigma_0(\iota_t)$ expresses $i_b(t)$ through
$i_b(\mu_0)$:
\[
 i_b(\mu_0)\big(\sigma_0(\iota_tx)\big)=i_b(t)(\sigma_0x)\qquad(x\in K^{d_Q(t)}).\tag{$*$}
\]
For $t=\mu_0$ this gives $\iota_{\mu_0}=1$, as $i_b(\mu_0)$ is injective.

\emph{The metric at $v_0$.} By (3) the vectors $i_b(\mu_0)(e_k)$ are linearly independent, so their
Gram matrix $A$ is positive definite, and so is $n_b^{-1}A$. By Definition~\ref{is-complex-data}(b),
$X_{E,v_0}$ is a metric for $n_b^{-1}A$, that is $X_{E,v_0}=(n_b^{-1}A)^{-1/2}$. So $X_{E,v_0}$ is
invertible, and
\[
 h_{X_E,v_0}(y,z)=n_b^{-1}\big\langle i_b(\mu_0)(\sigma_0y),i_b(\mu_0)(\sigma_0z)\big\rangle.
\]

\emph{The metric at $v\ne v_0$.} By (2), let $\pi_p$ be the projection of $K^D$ onto $Q_p(\mu_0)$
along the sum of the other parts, so that $\sum_p\pi_p=1$, and put
$c_p(v)=\kappa_v\prod_j\lambda_{j,\mu_0,p,v}$, which is positive by (1). By
Definition~\ref{is-complex-data}(b), $X_{E,v}$ is the metric $X_Q(\mu_0)$ scaled by $c_p(v)$ on the
coordinates of the $p$-th part: it is a metric for the matrix
$A_v=\sum_pc_p(v)\,\sigma_v(\pi_p)^*M_v\sigma_v(\pi_p)$, where
$M_v=(X_Q(\mu_0)_v^{-1})^*X_Q(\mu_0)_v^{-1}$ is positive definite by (2). The matrix $A_v$ is
positive semidefinite, and if $z^*A_vz=0$, then $\sigma_v(\pi_p)z=0$ for every $p$, so
$z=\sum_p\sigma_v(\pi_p)z=0$. So $A_v$ is positive definite, $X_{E,v}$ is invertible, and
\[
 h_{X_E,v}(y,z)=\sum_pc_p(v)\,h_{X_Q(\mu_0),v}(\pi_py,\pi_pz).
\]

\emph{The scalar conditions at $v_0$.} Let $x\in\iota_t^{-1}(E^K_p)$ and $y\in K^{d_Q(t)}$. By
($*$), $i_b(t)(\sigma_0x)=i_b(\mu_0)(\sigma_0(\iota_tx))\in E_p$. The formula for $h_{X_E,v_0}$,
($*$) and (5) give
\[
 h_{X_E,v_0}(\iota_tx,\iota_ty)=n_b^{-1}\big\langle i_b(t)(\sigma_0x),i_b(t)(\sigma_0y)\big\rangle
 =\kappa_{v_0}\prod_j\lambda_{j,t,p,v_0}\,h_{X_Q(t),v_0}(x,y).
\]

\emph{The scalar conditions at $v\ne v_0$, for $t=\mu_0$.} As $\iota_{\mu_0}=1$, $x\in E^K_p$, so
$x\in Q_p(\mu_0)$ by (4). Then $\pi_px=x$ and $\pi_qx=0$ for $q\ne p$, so
$h_{X_E,v}(x,y)=c_p(v)\,h_{X_Q(\mu_0),v}(x,\pi_py)$. By (6), $h_{X_Q(\mu_0),v}(x,\pi_qy)=0$ for
$q\ne p$, so $h_{X_Q(\mu_0),v}(x,\pi_py)=\sum_qh_{X_Q(\mu_0),v}(x,\pi_qy)=h_{X_Q(\mu_0),v}(x,y)$.
% Lean: XE_specR is rowXE_spec (spliceFrame at v0 with frameOf (gramE), gramE_unit; weightFrame
% elsewhere, weightFrameSpec with weightMet_posDef and sum_partPr) for the frame rowXE. scal_rowsR: at
% v0, majAt_gram_smul (the Hermitian form of frameOf (c Gram)), lin_basis, key_row (relMat_spec with
% kMat_map, from H2 and RATmu) and H4; at v != v0, rowMat_base, H3 and rowXE_part (weightFramePart with
% PartsOrth). Hypothesis (1) is H4pos (with nP > 0), (2) is DefSide.WF (XQ_quot, Qp_indep, Qp_top),
% (3) is H2 and RATmu, (4) is H3, (5) is H4 and (6) is PartsOrth.
\end{proof}

\begin{lemma}\label{asm-division-kernel-package-lifts}
In Situation~\ref{asm-division-kernel-package-situation}, let $j$ be a lift, with the Hilbert spaces
$W_{b,j,v}$, their vectors $w_{b,j,v}(k)$ and the map $e_{b,j}\colon W_{b,j,v_0}\to W_j$ of
Definition~\ref{is-complex-data}(4). For $c\in\mathbf C^{d_W(j)}$ write
$w_{b,j,v}(c)=\sum_kc_kw_{b,j,v}(k)$. Suppose that:
\begin{enumerate}
\item for every infinite place $v$ the vectors $w_{b,j,v}(k)$ are linearly independent;
\item $\langle e_{b,j}a,e_{b,j}a'\rangle=n_b\langle a,a'\rangle$ for all $a,a'\in W_{b,j,v_0}$;
\item for every tuple $t$ the map $t_b(j,t)$ is injective with the same image as $w_b(j)$, and the
  matrix that expresses $t_b(j,t)$ through $w_b(j)$ has entries in $\sigma_0(K)$;
\item for every tuple $t$, part $p$ and $x,y\in K^{d_Z(j,t)}$ with $t_b(j,t)(\sigma_0x)\in W_p(j)$,
  \[
   \langle t_b(j,t)(\sigma_0x),t_b(j,t)(\sigma_0y)\rangle=n_b\lambda_{j,t,p,v_0}\,h_{X_T(j,t),v_0}(x,y);
  \]
\item at every infinite place $v\ne v_0$, for every part $p$ and all $x,y\in K^{d_Z(j,\mu_0)}$ with
  $t_b(j,\mu_0)(\sigma_0x)\in W_p(j)$,
  $h_{X_W(j),v}(\theta_{j,\mu_0}x,\theta_{j,\mu_0}y)=\lambda_{j,\mu_0,p,v}\,h_{X_T(j,\mu_0),v}(x,y)$.
\end{enumerate}
Then $X_W(j)$ is invertible at every infinite place, and whenever $t=\mu_0$ or $v=v_0$,
\[
 h_{X_W(j),v}(\theta_{j,t}x,\theta_{j,t}y)=\lambda_{j,t,p,v}\,h_{X_T(j,t),v}(x,y)
\]
for every part $p$, every $x\in\theta_{j,t}^{-1}(W^K_p(j))$ and every $y\in K^{d_Z(j,t)}$.
\end{lemma}
\leanref{Formalism/LocalToGlobal/IndefiniteReal/CxCoreRG.lean}{Formalism.IndefR.Cx.CxData.scal_liftsR}
\leanref{Formalism/LocalToGlobal/IndefiniteReal/CxLemmas.lean}{Formalism.IndefR.Cx.frame_gram_unit}
\leanref{Formalism/LocalToGlobal/IndefiniteReal/CxLemmas.lean}{Formalism.IndefR.Cx.majAt_gram}
\leanref{Formalism/LocalToGlobal/IndefiniteReal/CxLemmas.lean}{Formalism.IndefR.Cx.CxData.key_lift}
\begin{proof}
By Definition~\ref{is-complex-data}(c), $X_W(j)$ is at each $v$ a metric for the Gram matrix of the
$w_{b,j,v}(k)$, which is positive definite by (1). So $X_W(j)$ is invertible at $v$, and
$h_{X_W(j),v}(y,z)=\langle w_{b,j,v}(\sigma_vy),w_{b,j,v}(\sigma_vz)\rangle$. By the same definition,
$w_b(j)(c)=e_{b,j}(w_{b,j,v_0}(c))$, and $\sigma_0(\theta_{j,t})$ expresses $t_b(j,t)$ through
$w_b(j)$, which is possible by (3):
\[
 w_b(j)\big(\sigma_0(\theta_{j,t}x)\big)=t_b(j,t)(\sigma_0x)\qquad(x\in K^{d_Z(j,t)}).\tag{$*$}
\]
Let $x\in\theta_{j,t}^{-1}(W^K_p(j))$, so that $w_b(j)(\sigma_0(\theta_{j,t}x))\in W_p(j)$; by ($*$),
$t_b(j,t)(\sigma_0x)\in W_p(j)$.

At $v_0$, by (2), ($*$) and (4),
\[
 n_b\,h_{X_W(j),v_0}(\theta_{j,t}x,\theta_{j,t}y)
 =\big\langle w_b(j)(\sigma_0\theta_{j,t}x),w_b(j)(\sigma_0\theta_{j,t}y)\big\rangle
 =\big\langle t_b(j,t)(\sigma_0x),t_b(j,t)(\sigma_0y)\big\rangle
 =n_b\lambda_{j,t,p,v_0}\,h_{X_T(j,t),v_0}(x,y),
\]
and $n_b>0$. At $v\ne v_0$ with $t=\mu_0$, the claim is (5).
% Lean: scal_liftsR: majAt_gram with kForm (the frame of XW is frameOf (gram (bW b j v))), at v0 EWnorm,
% key_lift (relMat_spec with kMat_map, from H5 and RATth) and ThmA, then cancellation of nP > 0; at
% v != v0, GalD. The invertibility of X_W(j) is frame_gram_unit with kForm. Hypothesis (1) is KForm, (2)
% EWnorm, (3) H5 and RATth, (4) ThmA and (5) GalD.
\end{proof}

\begin{lemma}\label{asm-division-kernel-package-hecke}
Let $\sigma\colon K\to\mathbf C$ be an embedding of a number field, let $V$ be a finite-dimensional
complex Hilbert space, $f\colon\mathbf C^d\to V$ an injective linear map with image $V'$, and
$(V_p)_{p\in\Pi}$ a finite family of subspaces of $V$ such that the spaces $V_p\cap V'$ are
orthogonal to each other and span $V'$. Put $V^K_p=\{y\in K^d:f(\sigma y)\in V_p\}$, and suppose
that every $y\in K^d$ is a sum of vectors $y_p\in V^K_p$. Let $I$ be a set and
$c\colon\Pi\times I\to K$, $(p,T)\mapsto c_p(T)$, a map with $c_p\ne c_q$ for $p\ne q$. Then there
are matrices $A_T\in M_d(K)$, for $T\in I$, with
\[
 V^K_p=\bigcap_{T\in I}\ker\big(A_T-c_p(T)\big)
\]
for every $p$.
\end{lemma}
\leanref{Formalism/LocalToGlobal/IndefiniteReal/ExistsSumOfMemISup.lean}{Formalism.IndefR.Cx.CxData.heckeRows_spec}
\leanref{Formalism/LocalToGlobal/IndefiniteReal/ExistsSumOfMemISup.lean}{Formalism.IndefR.Cx.CxData.heckeW_spec}
\leanref{Formalism/LocalToGlobal/IndefiniteReal/ExistsSumOfMemISup.lean}{Formalism.IndefR.Cx.spec_eig}
\leanref{Formalism/LocalToGlobal/IndefiniteReal/ExistsSumOfMemISup.lean}{Formalism.IndefR.Cx.spec_rat}
\leanref{Formalism/LocalToGlobal/IndefiniteReal/CxHecke.lean}{Formalism.IndefR.Cx.eig_transport}
\leanref{Formalism/LocalToGlobal/IndefiniteReal/CxHecke.lean}{Formalism.IndefR.Cx.relMat_op_key}
\leanref{Formalism/LocalToGlobal/IndefiniteReal/CxHecke.lean}{Formalism.IndefR.Cx.CxData.heckeB_of}
\leanref{Formalism/LocalToGlobal/IndefiniteReal/CxHecke.lean}{Formalism.IndefR.Cx.CxData.heckeC_of}
\begin{proof}
Let $\pi_p$ be the orthogonal projection of $V$ onto $V_p\cap V'$, and for $T\in I$ put
$h_T=\sum_p\sigma(c_p(T))\pi_p$, an operator of $V$ with image in $V'$. As the $V_p\cap V'$ are
orthogonal, $h_T$ acts on $V_q\cap V'$ as $\sigma(c_q(T))$, and $\pi_qh_T=\sigma(c_q(T))\pi_q$.

\emph{Eigenspaces.} We show that
$V_p\cap V'=\bigcap_T\ker\big(h_T-\sigma(c_p(T))\big)\cap V'$. The inclusion $\subset$ is the
remark above. Conversely, let $z\in V'$ with $h_Tz=\sigma(c_p(T))z$ for every $T$. For $q\ne p$
choose $T$ with $c_q(T)\ne c_p(T)$; then
$\sigma(c_q(T))\pi_qz=\pi_qh_Tz=\sigma(c_p(T))\pi_qz$, and $\sigma$ is injective, so $\pi_qz=0$.
As the $V_q\cap V'$ are orthogonal and span $V'$, $z=\sum_q\pi_qz=\pi_pz\in V_p$.

\emph{Rationality.} Write the $k$-th basis vector of $K^d$ as $e_k=\sum_py_{k,p}$ with
$y_{k,p}\in V^K_p$. Then $f(\sigma y_{k,p})\in V_p\cap V'$, so
$h_Tf(e_k)=\sum_p\sigma(c_p(T))f(\sigma y_{k,p})=f(\sigma a_k)$ with
$a_k=\sum_pc_p(T)y_{k,p}\in K^d$. Let $A_T$ be the matrix with columns $a_1,\dots,a_d$. By
linearity $h_Tf(c)=f(\sigma(A_T)c)$ for $c\in\mathbf C^d$, so $h_Tf(\sigma y)=f(\sigma(A_Ty))$ for
$y\in K^d$.

\emph{Transport.} Let $y\in K^d$. Then $y\in V^K_p$ if and only if $f(\sigma y)\in V_p\cap V'$, that
is, by the first step, if and only if $h_Tf(\sigma y)=\sigma(c_p(T))f(\sigma y)$ for every $T$. By
the second step this says $f(\sigma(A_Ty))=f(\sigma(c_p(T)y))$ for every $T$, and as $f$ and
$\sigma$ are injective, it says $A_Ty=c_p(T)y$ for every $T$.
% Lean: heckeRows_spec and heckeW_spec take hC T = specOp (the sum of sigma(ev p T) times the
% orthogonal projections onto the parts meet the level), with spec_eig (the first step, through
% proj_specOp, specOp_of_mem and sum_proj) and spec_rat (the second, through exists_sum_of_mem_iSup);
% heckeB_of and heckeC_of transport them to K^d by HeckeRowsC.eigE and HeckeWC.eigW, which are
% eig_transport with relMat_op_key (the third step). The orthogonality and spanning are PartsE and
% PartsW, the decompositions of K^d are rowsK and liftsK, and c is ConvFace.ev with ev_inj.
\end{proof}

\begin{lemma}\label{asm-division-kernel-package-normal}
Let $A\in M_n(\mathbf C)$ be Hermitian and positive definite, and write $\langle u,u'\rangle_A=u^*Au'$.
Let $(N_i)_{i\in I}$ be matrices in $M_n(\mathbf C)$ that commute in pairs, such that for every $i$
there is $i'\in I$ with $AN_i=N_{i'}^*A$. For a map $a\colon I\to\mathbf C$ let $V(a)$ be the set of
$u\in\mathbf C^n$ such that, for every $i$, some power of $N_i-a(i)$ kills $u$, and for a set $C$ of
such maps let $V(C)=\sum_{a\in C}V(a)$. Let $C$ and $C'$ be disjoint sets of maps $I\to\mathbf C$,
and let $u\in V(C)$ and $u'\in V(C')$.
\begin{enumerate}
\item $\langle u,u'\rangle_A=0$.
\item Let $P\in M_{d\times n}(\mathbf C)$ and $M_i\in M_d(\mathbf C)$ satisfy $PN_i=M_iP$ for every
  $i$, and suppose that $R=PA^{-1}P^*$ is invertible. Then $(Pu)^*R^{-1}(Pu')=0$.
\end{enumerate}
\end{lemma}
\leanref{Formalism/LocalToGlobal/IndefiniteReal/CxCutEP.lean}{Formalism.IndefR.Cx.EigenprojectionMatrices.cx_proj_adj}
\leanref{Formalism/LocalToGlobal/IndefiniteReal/CxCutEP.lean}{Formalism.IndefR.Cx.EigenprojectionMatrices.cx_eig_orth}
\leanref{Formalism/LocalToGlobal/IndefiniteReal/CxCutEP.lean}{Formalism.IndefR.Cx.EigenprojectionMatrices.cx_quot_orth}
\begin{proof}
(1) Let first $a\ne b$ be maps $I\to\mathbf C$, $x\in V(a)$ and $y\in V(b)$. Choose $i$ with
$a(i)\ne b(i)$, and $i'$ with $AN_i=N_{i'}^*A$. Taking adjoints gives $AN_{i'}=N_i^*A$, as $A$ is
Hermitian. So for every $\lambda\in\mathbf C$ and all $z,z'$,
\[
 \langle z,(N_i-\lambda)z'\rangle_A=\langle(N_{i'}-\bar\lambda)z,z'\rangle_A,\qquad
 \langle z,(N_{i'}-\bar\lambda)z'\rangle_A=\langle(N_i-\lambda)z,z'\rangle_A.
\]
Suppose that $(N_i-\lambda)z=0$, and put $z'=(N_{i'}-\bar\lambda)z$. As $N_i$ and $N_{i'}$ commute,
$(N_i-\lambda)z'=0$, so $\langle z',z'\rangle_A=\langle(N_i-\lambda)z',z\rangle_A=0$, and $z'=0$ as
$A$ is positive definite. So every $z$ killed by $N_i-\lambda$ is killed by $N_{i'}-\bar\lambda$. Next
suppose that
$(N_i-\lambda)^2z=0$, and put $w=(N_i-\lambda)z$. Then $w$ is killed by $N_i-\lambda$, hence by
$N_{i'}-\bar\lambda$, and $\langle w,w\rangle_A=\langle(N_{i'}-\bar\lambda)w,z\rangle_A=0$; so $w=0$.
By induction on $k$, $(N_i-\lambda)^kz=0$ implies $(N_i-\lambda)z=0$.

Hence $(N_i-a(i))x=0$, so $(N_{i'}-\overline{a(i)})x=0$, and $(N_i-b(i))y=0$. Therefore
\[
 0=\langle x,(N_i-b(i))y\rangle_A=\langle(N_{i'}-\overline{b(i)})x,y\rangle_A
 =\big(a(i)-b(i)\big)\langle x,y\rangle_A,
\]
and $\langle x,y\rangle_A=0$. As $\langle\,,\rangle_A$ is sesquilinear, it follows that
$\langle z,z'\rangle_A=0$ whenever $z\in V(D)$ and $z'\in V(D')$ for disjoint sets $D$ and $D'$ of
maps $I\to\mathbf C$.

Let $\neg C$ be the set of maps $I\to\mathbf C$ outside $C$. As $\mathbf C$ is algebraically closed,
the generalized eigenspaces of each $N_i$ span $\mathbf C^n$. As the $N_i$ commute, the spaces $V(a)$
are independent and span $\mathbf C^n$, so $\mathbf C^n=V(C)\oplus V(\neg C)$. Let $\Pi$ be the
projection of $\mathbf C^n$ onto $V(C)$ along $V(\neg C)$. It is self-adjoint for
$\langle\,,\rangle_A$: for $z,z'\in\mathbf C^n$, the vectors $\Pi z$ and $\Pi z'$ lie in $V(C)$ and
$z-\Pi z$ and $z'-\Pi z'$ lie in $V(\neg C)$, so by the last paragraph
\[
 \langle\Pi z,z'\rangle_A=\langle\Pi z,\Pi z'\rangle_A=\langle z,\Pi z'\rangle_A,
\]
that is, $\Pi^*A=A\Pi$. Now $\Pi u=u$, and $\Pi u'=0$ as $V(C')\subset V(\neg C)$. So
$\langle u,u'\rangle_A=\langle\Pi u,u'\rangle_A=\langle u,\Pi u'\rangle_A=0$.
% Lean: the first paragraph is dotProduct_mulVec_eigAt_eq_zero_of_ne (K1, K2, K3) and the extension to
% sums dotProduct_mulVec_eigSum_eq_zero_of_disjoint. The projection is Col.eigProj, which exists by
% Col.isCompl_eigSum (iSupIndep_eigAt and iSup_eigAt_eq_top; the splitting hs is supplied in
% part_orthg by Module.End.iSup_maxGenEigenspace_eq_top); its self-adjointness is cx_proj_adj, and the
% last step is cx_eig_orth (eigProj_apply_of_mem and eigProj_apply_of_mem_not with eigSum_mono).

(2) Put $Q=A^{-1}P^*R^{-1}P$; we show that $Q$ commutes with every $N_i$. Let $i'$ be as above. From
$AN_i=N_{i'}^*A$ we get $N_iA^{-1}=A^{-1}N_{i'}^*$, and taking adjoints in $PN_{i'}=M_{i'}P$ gives
$N_{i'}^*P^*=P^*M_{i'}^*$. So
\[
 M_iR=PN_iA^{-1}P^*=PA^{-1}N_{i'}^*P^*=PA^{-1}P^*M_{i'}^*=RM_{i'}^*,
\]
hence $R^{-1}M_i=M_{i'}^*R^{-1}$, and
\[
 QN_i=A^{-1}P^*R^{-1}M_iP=A^{-1}P^*M_{i'}^*R^{-1}P=A^{-1}N_{i'}^*P^*R^{-1}P=N_iQ.
\]
A matrix that commutes with every $N_i$ maps each $V(a)$ into itself, so $Qu'\in V(C')$. Finally
$(Pu)^*R^{-1}(Pu')=u^*AA^{-1}P^*R^{-1}Pu'=\langle u,Qu'\rangle_A$, which is $0$ by (1).
\end{proof}

\begin{lemma}\label{asm-division-kernel-package-orth}
Let a family of definite forms over a number field $L$ be given (Definition~\ref{is-realisation}),
with finite group $G$, base tuple $\mu_0$ and set $\Pi$ of types of the base, that satisfies the
family conditions with the conditions on its inert indices (Lemma~\ref{is-eigenvalue-base-inputs}(1)).
Let $L_0\subset L$ be the subfield on which the involution is complex conjugation at every
embedding, and suppose that the Hecke operators and the inert Hecke operators on the row forms at
$\mu_0$ have matrices with entries in $L_0$ in the rational basis, and that each of them has an
adjoint among them for the class-sum pairing. Let $P\subset G$ be a subgroup with a character
$\psi\colon P\to L^\times$, and let subspaces $Z_t$ of the row forms and $Z_{j,t}$ of the lift forms
be given, stable under the Hecke operators, the inert Hecke operators and $G$. Consider the definite
side that the family builds from $(P,\psi)$ and these subspaces, with the metric $X_Q(\mu_0)$ on
$L^{d_Q(\mu_0)}$ and the parts $Q_p(\mu_0)$ of Lemma~\ref{asm-primitive-rows-parts}, and for an
infinite place $v$ of $L$ write
$h_v(x,y)=(X_Q(\mu_0)_v^{-1}\sigma_v(x))^*X_Q(\mu_0)_v^{-1}\sigma_v(y)$. Then $h_v(x,y)=0$ for every
infinite place $v$, all types $p\ne q$ of the base, $x\in Q_p(\mu_0)$ and $y\in Q_q(\mu_0)$.
\end{lemma}
\leanref{Formalism/LocalToGlobal/IndefiniteReal/MajAtEqZeroOfMemQpOfNe.lean}{Formalism.IndefR.Cx.majAt_eq_zero_of_mem_Qp_of_ne}
\leanref{Formalism/LocalToGlobal/IndefiniteReal/EquivFunOM.lean}{Formalism.IndefR.Cx.GroupActionMatrices.part_orthg}
\leanref{Formalism/LocalToGlobal/IndefiniteReal/CxCutEP.lean}{Formalism.IndefR.Cx.EigenprojectionMatrices.Pq_tV}
\leanref{Formalism/LocalToGlobal/IndefiniteReal/CxCutEP.lean}{Formalism.IndefR.Cx.EigenprojectionMatrices.mem_eigSum_of_emb}
\begin{proof}
Fix $v$ and write $\sigma=\sigma_v$. Let $B$ be the space of row forms at $\mu_0$, with the matrix
$H_B$ of its class-sum pairing in the rational basis. As in the proof of
Lemma~\ref{asm-primitive-rows-parts} and in Lemma~\ref{is-baselinks-definite-coordinates}, the
definite side at $\mu_0$ is built from a basis of the image $\Gamma$ of $e_Xe_{P,\psi}$ on $B$, with
matrix $E_\Gamma$ in the rational basis and entries in $L_0$; the matrix
$H_\Gamma=|P|^{-1}E_\Gamma^*H_BE_\Gamma$; a basis matrix $E$ of the coordinate space $Z_0$ of
$Z_{\mu_0}\cap\Gamma$; the matrix $P_{\mu_0}$ that sends a vector to the coordinates, in the columns of
$E$, of its orthogonal projection onto $Z_0$ for $H_\Gamma$, so that $P_{\mu_0}E=1$; and the quotient
metric $X_Q(\mu_0)$ of $H_\Gamma$ through $P_{\mu_0}$. The coordinate map sends $x$ to the form with
coordinates $E_\Gamma Ex$, and $Q_p(\mu_0)$ is the set of $x$ for which this form lies in $R_{\mu_0}(p)$,
the sum of the joint generalized eigenspaces of the Hecke and inert Hecke operators for the systems of
type $p$. By the family conditions, $H_B$ is Hermitian with entries in $L_0$, $\sigma(H_\Gamma)$ is
positive definite, the Hecke and inert Hecke operators commute in pairs and with $G$, and $E_\Gamma$ is
injective.

\emph{Step 1: the operators on $\Gamma$.} Let $(T_i)_{i\in I}$ be the Hecke and inert Hecke operators of
$B$, with matrices $O_i$ in the rational basis. Each $T_i$ commutes with $G$, so with $e_Xe_{P,\psi}$,
and preserves its image $\Gamma$: there is a matrix $N_i$ over $L$ with $E_\Gamma N_i=O_iE_\Gamma$. As
the $T_i$ commute in pairs and $E_\Gamma$ is injective, so do the $N_i$. Let $T_{i'}$ be the adjoint of
$T_i$ for the class-sum pairing; as $H_B$ is Hermitian, $H_BO_i=O_{i'}^*H_B$. Then
\[
 E_\Gamma^*H_BE_\Gamma N_i=E_\Gamma^*H_BO_iE_\Gamma=(O_{i'}E_\Gamma)^*H_BE_\Gamma
 =(E_\Gamma N_{i'})^*H_BE_\Gamma=N_{i'}^*E_\Gamma^*H_BE_\Gamma,
\]
so $H_\Gamma N_i=N_{i'}^*H_\Gamma$. The same computation after applying $\sigma$, in which
$\sigma(E_\Gamma^*)=\sigma(E_\Gamma)^*$, $\sigma(O_{i'}^*)=\sigma(O_{i'})^*$ and
$\sigma(H_B)$ is Hermitian because these matrices have entries in $L_0$, gives
$\sigma(H_\Gamma)\sigma(N_i)=\sigma(N_{i'})^*\sigma(H_\Gamma)$.

\emph{Step 2: the projection.} The $T_i$ preserve $Z_{\mu_0}$, so the $N_i$ preserve $Z_0$, and
$N_iE=EM_i$ with $M_i=P_{\mu_0}N_iE$, as $P_{\mu_0}E=1$. Since $1-EP_{\mu_0}$ projects onto the
orthogonal complement of $Z_0$ for $H_\Gamma$, $E^*H_\Gamma(1-EP_{\mu_0})=0$ and
$P_{\mu_0}=(E^*H_\Gamma E)^{-1}E^*H_\Gamma$. By Step 1,
\[
 E^*H_\Gamma N_i(1-EP_{\mu_0})=(N_{i'}E)^*H_\Gamma(1-EP_{\mu_0})
 =M_{i'}^*E^*H_\Gamma(1-EP_{\mu_0})=0,
\]
so $P_{\mu_0}N_i(1-EP_{\mu_0})=0$, that is, $P_{\mu_0}N_i=M_iP_{\mu_0}$.

\emph{Step 3: the parts.} For a type $p$ let $C_p$ be the set of systems $a\colon I\to L$ of type $p$,
and $\sigma C_p$ the set of the $\sigma\circ a$, $a\in C_p$. For types $p\ne q$ of the base, $C_p$ and
$C_q$ are disjoint, by the two facts about types in the proof of
Lemma~\ref{asm-primitive-rows-parts}, and so are $\sigma C_p$ and $\sigma C_q$, as $\sigma$ is
injective. Write $V(\,\cdot\,)$ for the sums of joint generalized eigenspaces of the $\sigma(N_i)$, as in
Lemma~\ref{asm-division-kernel-package-normal}. Let $x\in Q_p(\mu_0)$ and $u=\sigma(Ex)$; we show that
$u\in V(\sigma C_p)$. The coordinates $E_\Gamma Ex$ lie in the sum of the joint generalized eigenspaces
of the $O_i$ for the systems in $C_p$, so $\sigma(E_\Gamma)u$ lies in the sum $U_1$ of those of the
$\sigma(O_i)$ for the systems in $\sigma C_p$. The $\sigma(N_i)$ commute, and their joint generalized
eigenspaces span $\mathbf C^{n}$, $n$ the dimension of $\Gamma$; write $u=u_1+u_2$ with
$u_1\in V(\sigma C_p)$ and $u_2$ in the sum of the joint generalized eigenspaces for the other systems.
As $\sigma(E_\Gamma)\sigma(N_i)=\sigma(O_i)\sigma(E_\Gamma)$, $\sigma(E_\Gamma)$ maps the joint
generalized eigenspace of the $\sigma(N_i)$ for a system into that of the $\sigma(O_i)$ for the same
system. So $\sigma(E_\Gamma)u_2$ lies in the sum $U_2$ of the joint generalized eigenspaces of the
$\sigma(O_i)$ for the systems outside $\sigma C_p$, and also in $U_1$, being
$\sigma(E_\Gamma)u-\sigma(E_\Gamma)u_1$. As the $\sigma(O_i)$ commute, $U_1\cap U_2=0$. So
$\sigma(E_\Gamma)u_2=0$, hence $u_2=0$ and $u\in V(\sigma C_p)$.

\emph{Step 4.} Let $p\ne q$, $x\in Q_p(\mu_0)$ and $y\in Q_q(\mu_0)$, and put $u=\sigma(Ex)$ and
$u'=\sigma(Ey)$, which lie in $V(\sigma C_p)$ and $V(\sigma C_q)$ by Step 3. As $P_{\mu_0}E=1$,
$\sigma(x)=\sigma(P_{\mu_0})u$ and $\sigma(y)=\sigma(P_{\mu_0})u'$. By the identity that defines the
quotient metric, its Hermitian form is $h_v(x,y)=\sigma(x)^*R^{-1}\sigma(y)$, where
$R=\sigma(P_{\mu_0})\sigma(H_\Gamma)^{-1}\sigma(P_{\mu_0})^*$ and the inverse of a singular matrix is
taken to be $0$. If $R$ is singular, then $R^{-1}=0$ and $h_v(x,y)=0$. If $R$ is invertible,
Lemma~\ref{asm-division-kernel-package-normal}(2), with $A=\sigma(H_\Gamma)$, which is positive
definite, the matrices $\sigma(N_i)$, which satisfy its hypotheses by Step 1, the matrices
$\sigma(P_{\mu_0})$ and $\sigma(M_i)$ of Step 2, and the sets $\sigma C_p$ and $\sigma C_q$, gives
$h_v(x,y)=(\sigma(P_{\mu_0})u)^*R^{-1}(\sigma(P_{\mu_0})u')=0$.
% Lean: majAt_eq_zero_of_mem_Qp_of_ne = part_orthg hB hO (the stability hst.1 at mu0) v with
% C = HasRes3 p, C' = HasRes3 q and hSS from excl3 and diffAt_of_mem_ixFin3; hO = okOp_opAll (rat and
% adj from OkH and InertRationalAdjoint, comm and split from Ok3, opG). In part_orthg: hA = lf_HV_pos,
% hc = tVg_comm, adjCg (adjC: opMat_adj, HV_map, adj_transfer), hPN = Pq_tVg (Pq_tV: tV_adj, tV_zE,
% Pq_orth), key (Step 3: eigSum_le_comap_of, emb_mem_eigSum and mem_eigSum_of_emb with vE_tVg and
% lf_v_proj), hPu (pq_orth), hq = XQ_quot (IsQuotMetric, whose identity uses Lean's nonsingular
% inverse), rewritten by majBasic_majAt. The case split is by_cases on IsUnit (det R): cx_quot_orth in
% the invertible case, Matrix.nonsing_inv_apply_not_isUnit in the singular one.
% Size: MajAtEqZeroOfMemQpOfNe.lean, EquivFunOM.lean and CxCutEP.lean, 970 lines (text count).
\end{proof}

\begin{lemma}\label{asm-division-kernel-package}
Let $P$ be a row-side family for a definite-side datum $D$ with coefficient field $L$
(Definition~\ref{is-row-side-family}), and let $Q$ be an indefinite realization of $P$
(Definition~\ref{is-realisation-join}), with realization $R$ at the face $F$ of $P$. Let $C,C_H$ be
real, write $S_0(R,F)$ and $S_1(R,F)$ for the rows and columns of $R$ at $F$, and
$N=N(R,F;C,C_H)$ for their numerical bounds (Definition~\ref{is-realisation}). Suppose that:
\begin{enumerate}
\item the family of definite forms of $P$ satisfies the family conditions at some level and has
  lift-distinct types, and $R$ has primitive rows at $F$ (Definition~\ref{is-primitive-rows});
\item every finite place $w$ at which an index $I_w(j,\mu_0)$ or $I_{\bar w}(j,\mu_0)$ of
  $S_0(R,F)$, a relative index of $T_{j,\mu_0}$ and its dual, or a value $F_0(w,j)$ or
  $F_0(\bar w,j)$ of $N$ is nonzero, is twisted;
\item $S_0(R,F)$ satisfies the covolume, surjectivity and functional bounds for $N$, and
  $S_1(R,F)$ the lower norm bound for $N$ (Definitions~\ref{is-numerical-bounds},
  \ref{is-surjectivity-bound}, \ref{is-functional-bound} and~\ref{is-lower-norm-bound});
\item the isometry of the row space, the multiplicity condition and the passage from rows to
  columns hold for $R$ at $F$, as in conclusion (5) of Theorem~\ref{is-definite-rows}.
\end{enumerate}
Then there is a division package over $L$ for $N$ (Definition~\ref{is-division-package}) on the
definite subspaces of $R$ at $F$, whose rows are $S_0(R,F)$, whose columns are $S_1(R,F)$ and
whose column data are those of $R$ at $F$, together with Hecke data for its rows and columns. The
definite subspaces, rows, columns and column data are those of $P$.
\end{lemma}
\leanref{Development/AutomorphicArithmetic/ThetaCorrespondence/ThetaPairLifts/ThirdPlaceSeries/PtKOfReal.lean}{Formalism.L2.RowDataTransport.packageOfPrimitiveLinks}
\leanref{Formalism/LocalToGlobal/IndefiniteReal/PrimitivePackageFullHecke.lean}{Formalism.IndefR.Links.primitivePackageFullHecke}
\leanref{Formalism/LocalToGlobal/IndefiniteReal/PrimitivePackageFullHecke.lean}{Formalism.IndefR.Links.PassIn3}
\leanref{Formalism/LocalToGlobal/IndefiniteReal/BaseLiftsPackageOfPrimitiveLinks.lean}{Formalism.IndefR.Links.CovolumeOntoIsometryConditions}
\leanref{Formalism/LocalToGlobal/IndefiniteReal/BaseLiftsPackageOfPrimitiveLinks.lean}{Formalism.IndefR.Links.base_lifts_package_of_primitive_links}
\leanref{Formalism/LocalToGlobal/IndefiniteForms/PairingConditionOfInverseMatrixIdentity.lean}{Formalism.L2.Indef.IndefPkgR}
\leanref{Formalism/LocalToGlobal/IndefiniteReal/NearConstantMultiplicities.lean}{Formalism.IndefR.Cx.packageFullHecke}
\leanref{Formalism/LocalToGlobal/IndefiniteReal/NearConstantMultiplicities.lean}{Formalism.IndefR.Cx.passage3_of}
\leanref{Formalism/LocalToGlobal/IndefiniteReal/PackageFromPrimitiveStatements.lean}{Formalism.IndefR.Cx.basePackageFullHecke}
\leanref{Formalism/LocalToGlobal/IndefiniteReal/CxCoreRG.lean}{Formalism.IndefR.Cx.CxData.cxPkgR}
\begin{proof}
Let $X$ be the complex data of $R$ at $F$, with definite subspaces $Z_s$. The package is assembled
from $X$ in three steps.

\emph{Step 1: the base-level statements.} The weak characterization of $R$ at $F$, which is part of
$Q$, gives the models of the forms on the rows and the lifts and their identifications. Hypothesis
(1) gives the family conditions and the lift-distinct types, and the primitive rows give the Hecke
operators and theta maps of the rows and the lifts, the unit and equivariance conditions of the
actions of $G$, the levels, the positivity of the scalars, the row and lift tests, and the
decomposition into parts with the dimensions of the parts. Hypotheses (2), (3) and (4) give the
condition on twisted places, the covolume, surjectivity, lower norm and functional bounds, and the
isometry and multiplicity conditions. Together these are the base-level statements of $X$. Among
them are the eigenvalue and type conditions of the lifts at the base level: for every lift $j$,
$Z_{j,\mu_0}\cap\Gamma_0(j,\mu_0)$ is the intersection of $\Gamma_0(j,\mu_0)$ with the sum of the
joint generalized eigenspaces of the Hecke and inert Hecke operators for the systems $y$ in the set
$C_j$ of systems in $L$ whose images under $\sigma_0$ are lifted, and every $y\in C_j$ has type
$\operatorname{tr}_jp$ for some type $p$ of the base. These are shown in Steps 0 and 5 of the proof
of Lemma~\ref{asm-primitive-rows-base}, whose hypotheses are hypothesis (1) and the weak
characterization of $R$ at $F$.

\emph{Step 2: the passage.} By Lemma~\ref{asm-division-kernel-package-passage}, the
indefinite-side data of $X$ at $(P_0,\psi_0)$ and at $(P_s,\psi_s)$ on $Z_s$ satisfy the passage of
Definition~\ref{is-division-package}(4) for $N$.

\emph{Step 3: the package.} The rows $S_0$ are the indefinite-side data of $X$ at $(P_0,\psi_0)$ on
$Z_s$, and the columns $S_1$ those at $(P_s,\psi_s)$. We check conditions (1) to (3) of
Definition~\ref{is-division-package}.

The rows and the columns are built from the same complex data, so they have the same map
$w\mapsto\mu_w$, the same $\kappa_v$ and the same $\lambda_{j,t,p,v}$; their definite sides are well
formed by the family conditions, as in the proof of Lemma~\ref{asm-primitive-rows-parts}. The weak
characterization of $R$ at $F$ gives $\mu_{\bar w}=\mu_w$ for every finite place $w$, and makes $w$
twisted exactly when $\mu_w\ne\mu_0$. By Lemma~\ref{asm-primitive-rows-base}(2), at both
levels the matrices $\iota_t$ are bijective, $\iota_{\mu_0}=1$, and the $P_tM_t$ are lattices; by
Lemma~\ref{asm-primitive-rows-base}(3) the matrices $\theta_{j,t}$ of $S_0$ are bijective,
$\iota_{\mu_0}Q_p(\mu_0)=E^L_p$ and $\theta_{j,\mu_0}Z_p(j,\mu_0)=W^L_p(j)$, the dimension count
$\dim_LW^L_p(j)=\dim_LE^L_p$ holds, and the $Q_p(\mu_0)$ and the $Z_p(j,\mu_0)$ are independent and
span. Condition (7) of Definition~\ref{is-primitive-rows} gives $d_Q(\mu_0)>0$ at both levels,
$\kappa_v>0$ and $\lambda_{j,t,p,v}>0$, and condition (12) makes the $M_j$ lattices. Step 5 of the
proof of Lemma~\ref{is-eigenvalue-base-inputs}, which applies here as in Step 0 of the proof of
Lemma~\ref{asm-primitive-rows-base}, gives the third base link: $H_T(j,\mu_0)$ has entries in $K_0$
and $X_T(j,\mu_0)$ is a metric for it at every infinite place. Hypothesis (3) gives the covolume,
surjectivity and functional bounds of $S_0$ and the lower norm bound of $S_1$.

The rows are untwisted. Let $w$ be a finite place that is not twisted and $j$ a lift. Then
$\mu_w=\mu_0$, so the indices $I_w(j,\mu_w)$, $I_{\bar w}(j,\mu_w)$ and
$j_w(H_T(j,\mu_w),T_{j,\mu_w})$ are those at $\mu_0$. By hypothesis (2) they vanish, and so does
$F_0(w,j)$; so their sum vanishes.

The column data are those of $R$ at $F$ for the isometry $s_E$ of $Q$
(Definition~\ref{asm-division-kernel-multiplicity}(3)). By hypothesis (4) the isometry of the row space
and the multiplicity condition hold for $s_E$, so Lemma~\ref{asm-division-kernel-factorization} gives
scalars $\zeta_u(t)\in L^\times$ with the factorization of Definition~\ref{is-division-package}(3),
and the column data for these scalars are isometric.

\emph{The metrics and the scalar conditions.} The rows and the columns are the indefinite-side data
$S_b$ of $X$ at the levels $b=0,1$, and the number $n_b$ of $X$ is $|P_b|$. The scalars of $X$ are
$\kappa_v=\kappa$ and $\lambda_{j,t,p,v}=\lambda^{(3)}_{j,t,p,v}$ (Definition~\ref{is-realisation}),
and $\Theta_t=\mathrm{Th}_t$. At each level $b$, Lemma~\ref{asm-division-kernel-package-rows}
applies to $X$ and $S_b$, with $K=L$. Its hypothesis (1) is condition (7) of
Definition~\ref{is-primitive-rows}. Hypothesis (2) holds because the definite side is well formed:
$X_Q(\mu_0)$ is the quotient metric of $H_V(\mu_0)$, which is positive definite, and the
$Q_p(\mu_0)$ are independent and span by Lemma~\ref{asm-primitive-rows-parts}(2). Hypothesis (3) is
Lemma~\ref{asm-primitive-rows-base}(2), and (4) is Step 4 of its proof, which applies at either
level. For (5), let $x,y\in L^{d_Q(t)}$ with $i_b(t)(\sigma_0x)\in E_p$, and put
$u(x)=j_t(c_{b,t}(x))$. Since $c_{b,t}(x)$ lies in the definite subspace
$Z_t=j_t^{-1}(\Theta_t^*E)$ (Lemma~\ref{is-baselinks-definite-coordinates}(1)),
$u(x)\in\Theta_t^*(E)$; and $i_b(t)(\sigma_0x)=\Theta_tu(x)$ by semilinearity. By condition (9),
$\Theta_t\Theta_t^*$ acts on $E_p$ as $\lambda_p=\kappa\prod_j\lambda^{(3)}_{j,t,p,v_0}$. So the
vector $z=\Theta_t^*\Theta_tu(x)-\lambda_pu(x)$ lies in $\Theta_t^*(E)$, and
$\Theta_tz=\Theta_t\Theta_t^*\big(\Theta_tu(x)\big)-\lambda_p\Theta_tu(x)=0$, as
$\Theta_tu(x)\in E_p$. As $\Theta_t^*(E)$ is the orthogonal complement of $\ker\Theta_t$, $z=0$.
Hence
\[
 \langle\Theta_tu(x),\Theta_tu(y)\rangle=\langle\Theta_t^*\Theta_tu(x),u(y)\rangle
 =\lambda_p\langle u(x),u(y)\rangle=\lambda_p|P_b|\,h_{X_Q(t),v_0}(x,y),
\]
the last by Step 0 of the proof of Lemma~\ref{asm-primitive-rows-base} and
Lemma~\ref{is-baselinks-definite-coordinates}(2). This is (5), as $n_b=|P_b|$. Hypothesis (6) is
Lemma~\ref{asm-division-kernel-package-orth} for $(P_b,\psi_b)$ and the definite subspaces of $X$,
whose hypotheses are hypothesis (1), condition (1) of Definition~\ref{is-primitive-rows} at $\mu_0$, and
the stability of the definite subspaces, Lemma~\ref{asm-primitive-rows-base}(1).
% Lean: PartsOrth is ComplexDataConditionsBaseLevel.partsOrth = majAt_eq_zero_of_mem_Qp_of_ne (Ok3,
% heckeR at mu0 and the stability ComplexDataConditionsBaseLevel.stable of the definite subspaces).
So $X_E$ is invertible at every infinite place, at both levels, and the rows and the columns satisfy
the scalar conditions for the $\iota_t$ after Definition~\ref{is-indefinite-datum}.

At level $0$, Lemma~\ref{asm-division-kernel-package-lifts} applies to every lift $j$. Its hypothesis
(1) holds because the $w_{0,j,v}(k)$ are a common nonzero multiple of the images, under the rational
structure $FW_j$ at $\sigma_v$, of $L$-linearly independent vectors, and these images are linearly
independent by condition (11).
Hypothesis (2) holds because $e_{0,j}$ is $\sqrt{n_0}$ times an inclusion, by the construction of
the complex data of $R$. Hypothesis (3) is Step 0 of the proof of
Lemma~\ref{asm-primitive-rows-base}. Hypothesis (4) follows as (5) above, with $\theta'_{j,t}$,
$j_{j,t}$, $c_{0,j,t}$ and $Z_{j,t}=j_{j,t}^{-1}(\theta'^*_{j,t}W_j)$ in place of $\Theta_t$, $j_t$,
$c_{b,t}$ and $Z_t$: condition (9) gives
$\mathrm{th}_G(j,t)\mathrm{th}_G(j,t)^*=\lambda^{(3)}_{j,t,p,v_0}$ on $W_p(j)$, and
$\theta'_{j,t}\theta'^*_{j,t}=\mathrm{th}_G(j,t)\mathrm{th}_G(j,t)^*$ by Step 0; and
Lemma~\ref{is-baselinks-definite-coordinates}(2) applies to the lift forms, with $X_Z=X_T(j,t)$.
Hypothesis (5) is condition (12). So the $X_W(j)$ are invertible and the rows satisfy the scalar
conditions for the $\theta_{j,t}$.
% Lean: cxPkgR builds rowsOfR and colsOfR from RowHyp (wf, h4pos, H2, RATmu, H3, H4, PartsOrth) and,
% at level false, LiftHyp (H5, RATth, KForm, EWnorm, ThmA, GalD, H3W, MLat), given by
% ComplexDataConditionsBaseLevel.rowHyp and liftHyp0. There H2 and H4 are cutRows (H4 through
% eq_zero_of_mem_range_adjoint, ThmBop = thmBop_of hP.thmB, and definite_isometry_of_norms with npTie:
% nP b = card P_b), RATmu is ratMuOfRowTest, H3 is definite_equivariance_of_primitive (part_iff_rows3 at
% either level), H5 and ThmA are cutLifts (ThmAop = thmAop_of with the unit phase), KForm is
% RealDat3.kForm_of hP.unif (condition (11); bW = (sqrt nP)^-1 times the images of a basis), EWnorm is
% RealDat.ewNorm (eW = sqrt nP times the inclusion) and GalD is hP.scalConj (condition (12)).

\emph{The Hecke data.} Let $I$ be a set and $c\colon\Pi\times I\to L$, $(p,T)\mapsto c_p(T)$, a map
with $c_p\ne c_q$ for $p\ne q$. Hecke data for $(I,c)$ on the rows are matrices $A_T$ of size $D$
and $B_{j,T}$ of size $d_W(j)$ over $L$, for $T\in I$ and lifts $j$, with
$E^L_p=\bigcap_T\ker(A_T-c_p(T))$ and $W^L_p(j)=\bigcap_T\ker(B_{j,T}-c_p(T))$ for every part $p$;
on the columns they are matrices $A_T$ with the first property. Lemma~\ref{asm-division-kernel-package-hecke}
gives them. For the $A_T$ at level $b$, apply it with $f=i_b(\mu_0)$, $V=E$, $V'=E^b$ and
$V_p=E_p$: $i_b(\mu_0)$ is injective with image $E^b$ by Lemma~\ref{asm-primitive-rows-base}(2), the
$E_p\cap E^b$ are orthogonal and span $E^b$ by condition (10), and every $y\in L^D$ is a sum of
vectors of the $Q_p(\mu_0)=E^L_p$, by hypotheses (2) and (4) of
Lemma~\ref{asm-division-kernel-package-rows} above. For the $B_{j,T}$, apply it with $f=w_0(j)$,
$V=W_j$, $V'$ the image of $w_0(j)$ and $V_p=W_p(j)$: $w_0(j)$ is injective by Step 0 of the proof of
Lemma~\ref{asm-primitive-rows-base}, the $W_p(j)$ meet $V'$ in orthogonal spaces that span it by
condition (10), and every $z\in L^{d_W(j)}$ is $\theta_{j,\mu_0}x$ with $x=\sum_px_p$,
$x_p\in Z_p(j,\mu_0)$ and $\theta_{j,\mu_0}x_p\in W^L_p(j)$, by Lemma~\ref{asm-primitive-rows-base}(3).
% Lean: heckeB_of (h2, kForm, ewNorm, h4pos with HeckeRowsC and HeckeWC at level false) and heckeC_of
% (HeckeRowsC at level true) of CxHecke.lean, with heckeRB_of, heckeWB_of and heckeRC_of
% (ExistsSumOfMemISup.lean): heckeRows_spec with PartsE and rowsK (from WF and H3), heckeW_spec with
% PartsW and liftsK (liftMat_bij, Types and H3W). The pair (I, c) is a ConvFace (Hk, ev, ev_inj).

With the passage of Step 2 these are a division package over $L$ for $N$, on $Z_s$. Its rows,
columns and column data are those of $R$ at $F$ by construction, and the last sentence of the
statement is part of $Q$ (Definition~\ref{is-realisation-join}).
% Lean of the three constructions: U2.unitary_two_package (eig0 = eig0Of with eigCOf, typ0 = typ0Of,
% Lemma asm-primitive-rows-base) and
% complex_data_conditions_base_level_of_primitive_links (heckeRB, heckeWB, heckeRC) in
% base_lifts_package_of_primitive_links; passage3_of = CxPass.cutPassage3S (cutPassage3S_of_steps
% steps3, CutPassage3.lean); in cxPkgR: rowsOfR and colsOfR (CxCoreRG.lean), CxData.cutDims
% (FinrankISupLeSum.lean), CxData.cutTuples2 (CutPassage.lean, with hecke_tuple_condition), colDef with
% zetaEps chosen by Classical.epsilon from the field cols, heckeB_of and heckeC_of.
% Lean: packageOfPrimitiveLinks Q hOk hsep C CH hP hI = primitivePackageFullHecke hOk hsep Q.Rl
% (faceT ...) Q.sEE (hkCan A3) C CH Q.charW hP hI = packageFullHecke (base_lifts_package_of_primitive_links
% hOk hsep hC hP) with the passage data of PassIn3 (CxPassIn3: ratMuP, isGood, chars, m1, nbarPos,
% near_constant_multiplicities, laTie). base_lifts_package_of_primitive_links reads the fields of
% PrimRows3 and of CovolumeOntoIsometryConditions (w1, covol, onto, lamLow, seeIso, miStar, fun0) into
% PrimitiveStatementsBaseLiftsOnly, with U2.unitary_two_package for eig0 and typ0 and
% complex_data_conditions_base_level_of_primitive_links for the Hecke data. The ties of the last
% sentence are Q.zs, Q.s0, Q.s1 and Q.col. Size: proof closure of primitivePackageFullHecke and
% packageFullHecke, 343 theorems, about 4,900 lines in 68 files.
\end{proof}

\begin{lemma}\label{asm-division-kernel-record-scale}
Let $D$ be a member at a scale $x>x_1$ of a family $\mathcal F$ with exponent $\varepsilon$, with
the quantities $b=\exp(x^{\varepsilon})$ and $X'=\exp(x^{\varepsilon/3})$ of its size data, and let
$\Pi$ and $c_t$ be as in Lemma~\ref{asm-division-kernel-sizes}.
\begin{enumerate}
\item If $x\ge16$ and $48c_t/\varepsilon+6/\varepsilon+2\le x^{\varepsilon/6}$, then
  $\Pi^e\le X'\le b$ for every integer $0\le e\le8$, and $\log x\le x^{\varepsilon}$.
\item If $c$ and $e\ge0$ are real with $c\cdot\lceil(e+1)/\varepsilon\rceil!\le x$, then $cx^e\le b$.
\end{enumerate}
\end{lemma}
\leanref{Formalism/LocalToGlobal/Divisibility/EigenvalueSumLeClosedOf.lean}{Formalism.L2.DIV.Arith4S.P_le_X'}
\leanref{Formalism/LocalToGlobal/Divisibility/EigenvalueSumLeClosedOf.lean}{Formalism.L2.DIV.Arith4S.logx_le_L}
\leanref{Formalism/LocalToGlobal/Divisibility/EigenvalueSumLeClosedOf.lean}{Formalism.L2.DIV.Arith4S.mul_rpow_le_exp_rpow}
\begin{proof}
For $y\ge1$ and $e'>0$, $\log y\le(2/e')y^{e'/2}$, since $\log z\le z$ for $z=y^{e'/2}$.

(1) As $x\ge9$, Lemma~\ref{asm-division-kernel-sizes}(1) gives $\Pi^e\le x^{8c_t}$. With $e'=\varepsilon/3$,
$8c_t\log x\le(48c_t/\varepsilon)x^{\varepsilon/6}\cdot x^{\varepsilon/6}\le x^{\varepsilon/3}$, so
$x^{8c_t}\le X'$; and $X'\le b$. Similarly, with $e'=\varepsilon$ and $T=x^{\varepsilon/2}\ge
2/\varepsilon$, $\log x\le(2/\varepsilon)T\le T^2=x^{\varepsilon}$.

(2) We may take $c>0$. Put $N=\lceil(e+1)/\varepsilon\rceil$, so $x^{e+1}\le x^{\varepsilon N}$. Then
$cx^e\cdot N!\le x\cdot x^e\le(x^{\varepsilon})^N$, and $(x^{\varepsilon})^N/N!\le\exp(x^{\varepsilon})=b$.
\end{proof}

\begin{lemma}\label{asm-division-kernel-record-numbers}
Let $D$ be a member at a scale $x>x_1$ of a family $\mathcal F$ with exponent $\varepsilon$, size
constant $C_{sz}\ge0$, the constant $C_c$ of its size data, the constant $C_{YL}$ of its parameter
choices and $c_t$ as in Lemma~\ref{asm-division-kernel-sizes}, with definite height $(h_1,h_2)$
(Definition~\ref{asm-definite-side-hypotheses}), and suppose that $x\ge16$,
$x\ge3+|C_c|$ and $48c_t/\varepsilon+6/\varepsilon+2\le x^{\varepsilon/6}$. Let $G$ be a group of order
$\prod_{k\in\{0,1\}}|GL_r(\mathbf F_{g_k})|^f$, and let $c_0\le c_e$ be real numbers with
$c_0\ge\max\big(1,\,C_{sz}(3+2C_{YL}+c_t),\,2(1+|C_c|)+1\big)$. For every word $u$, form the
record of $u$ with $\lambda_0=\log|G|$, $(C_0,c_e,c)=(c_0,c_e,C_{sz})$ and $S$, $\eta$ given by the
definite height, and define $\epsilon_{j,\rho}$ and $\delta_{j,\rho}$ by the two relations of
Definition~\ref{is-record-numbers}(1). Then conditions (1) to (6) of
Definition~\ref{is-record-numbers} hold. Moreover, with $S_w=(1+|C_c|)S_c$ and $C_1=2(1+|C_c|)$:
$2\le C_1$, $2S_w\le C_1S_c$, $C_1+1\le c_0$; the row exponents $n_{u,j,i}(\tau)$, $i<N(j-1)$, are
nonnegative, and those with $i\ge k$ add up to at most $m(t+1)\eta$; and
$\frac12(k+(2j-1)(t+1))\le e_j(\tau)$, $e_j(\tau)+S_w\le ar_E$ and $a_t+e_j(\tau)\le ar_E$ at the top
embedding of $t\in T$, for every lift $j$ and $\tau\in\mathcal T$.
\end{lemma}
\leanref{Formalism/LocalToGlobal/Divisibility/EigenvalueSumLeClosedOf.lean}{Formalism.L2.DIV.Arith4S.numFacts_rec}
\leanref{Formalism/LocalToGlobal/Divisibility/EigenvalueSumLeClosedOf.lean}{Formalism.L2.DIV.Arith4S.numFacts_rows}
\leanref{Formalism/LocalToGlobal/Divisibility/NumbersLink.lean}{Formalism.L2.DIV.RecNums.ofNum}
\leanref{Formalism/LocalToGlobal/Divisibility/EigenvalueSumLeClosedOf.lean}{Formalism.L2.DIV.Arith4S.nf_logrd}
\leanref{Formalism/LocalToGlobal/Divisibility/EigenvalueSumLeClosedOf.lean}{Formalism.L2.DIV.Arith4S.nf_Om}
\leanref{Formalism/LocalToGlobal/Divisibility/EigenvalueSumLeClosedOf.lean}{Formalism.L2.DIV.Arith4S.nf_lam}
\leanref{Formalism/LocalToGlobal/Divisibility/EigenvalueSumLeClosedOf.lean}{Formalism.L2.DIV.Arith4S.nf_eta}
\leanref{Formalism/LocalToGlobal/Divisibility/EigenvalueSumLeClosedOf.lean}{Formalism.L2.DIV.Arith4S.size_le_two_b}
\leanref{Formalism/LocalToGlobal/Divisibility/EigenvalueSumLeClosedOf.lean}{Formalism.L2.DIV.Arith4S.Lp_le_La}
\leanref{Formalism/LocalToGlobal/Divisibility/EigenvalueSumLeClosedOf.lean}{Formalism.L2.DIV.Arith4S.Nm_le_w''}
\begin{proof}
We use Lemma~\ref{asm-division-kernel-explicit-entries} and Lemma~\ref{asm-division-kernel-record-scale}(1):
$\Pi^e\le X'\le b$ for $e\le8$ and $\log x\le x^{\varepsilon}$. The record of $u$ has $N=N(m)$, the
places $z\in\{0,1\}\times A$, of which there are $2n=f$, weights $w_{j,z}=e_j(\tau_z)$, the row
exponents $\nu_{j,z,i}=n_{u,j,i}(\tau_z)$ for $i<N(j-1)$, $\log H_{j,z}=\log\mathrm{hook}_j(\tau_z)$,
$q_j=(k+j-1)-(m-j)$, $S=S_c=4k\Delta_S$ and the width $\eta$; the definite height gives $\Delta_S=h_1$
and $\eta=mh_2$.

\emph{Conditions (1) to (3).} Condition (1) holds by the definition of $\epsilon_{j,\rho}$ and
$\delta_{j,\rho}$. In (2), $r_E\ge1$ for the host, and the sizes give
$r_E\le\frac{Q_S}2\sqrt{qd}$, $Q_S/2\le2x^3$ and $\log q,\log d\le x^{\varepsilon/2}$; so, with
$T=x^{\varepsilon/2}\ge6/\varepsilon+2$, $\log r_E\le\log2+3\log x+T\le1+(6/\varepsilon)T+T\le T^2=x^{\varepsilon}$.
Also $N_c\ge0$. In (3), $k+m\le N(m)$ and $S,t,\eta,b',\lambda_0,L_p\ge0$.

\emph{Condition (4).} $|q_j-k|=|2j-1-m|\le m$. For the weights: the parameter choices give
$kS_c\le w''$ with $\Delta_S\ge1$, $t\le\eta$ and $m(t+1)\eta\le w''$, so $N(m)+k+1\le w''$; and
$w''\le b'/x\le b'\le b$. The windows of the size data put $e_j(\tau)$ above $b$, $b'/x$ or $w''$, so
$e_j(\tau)\ge N(m)\ge\frac12(k+(2j-1)(t+1))\ge0$. For the upper bound $e_j(\tau)\le2b$: the sizes give
$\eta\le b/(t^3fk^3)\le b/2$, $b'\le b/4$ (as $4x\le\Pi^2\le X'$), $w''\le b'$, and
$b_t\le b'+(2n\delta_1x^{\varepsilon}\Lambda_G/k)b'\le2b'$ for $t\in T$; the windows then give
$e_j(\tau)\le2b$. Finally $\log\Omega=\gamma+\log\pi+\log2-\log r_E-\operatorname{Re}\frac{L'}{L}(1,\chi_d)$
and $L_a=\log x+x^{\varepsilon}+\log r_E$, with $|\frac{L'}{L}(1,\chi_d)|\le x^{\varepsilon}$, $0\le\gamma<\frac23$,
$0\le\log\pi\le2\log2$ and $\log x\ge4\log2$; so $|\log\Omega|\le L_a$.

\emph{Condition (5).} The order of $G$ gives $\log|G|\le2fr^2\log(2g_0)$, and the sizes give
$\log(2g_0)\le2x^{2\varepsilon}$, $r\le t$; so $m\lambda_0\le4\Pi^6\le\Pi^7\le b$, as $\Pi\ge k\ge x\ge16$. Next $N^2\le\Pi^2\le b$, $m^3N^2x^{\varepsilon}\le\Pi^6\le b$ and $kN\le\Pi^2\le b$. The sizes give
$\eta t^3fk^3\le b$ and $m^2x^{\varepsilon}\le k$, and $N(m)\le3kt$, which give $m\eta Nx^{\varepsilon}\le b$.
The bounds $kS_c\le w''$ and $m(t+1)\eta\le w''$ are sizes. As $w''=b/(X'n)$ and $N\le\Pi\le X'$,
$Nw''\le b$; as $b'=a/X'$, $b'N\le a$ and, with $mk\le\Pi^2\le X'$, $mkb'\le a$. Finally $n_V\le N$.

\emph{Condition (6).} $1\le c_0\le c_e$ and $0\le C_{sz}\le c_0$. The parameter choices give
$\log g_i\le\log2+C_{YL}x^{\varepsilon}$, and $\log n_V\le c_t\log x$ by
Lemma~\ref{asm-division-kernel-sizes}(1); so
$L_p=x^{\varepsilon}+\log r_E+\log g_0+\log g_1+\log n_V\le(3+2C_{YL}+c_t)L_a$, and $C_{sz}L_p\le c_0L_a$.

\emph{The further facts.} $C_1=2(1+|C_c|)\ge2$, $2S_w=C_1S_c$, and $C_1+1\le c_0$. The lower bound
for $e_j(\tau)$ was shown under (4).

For $i\ge k$, $n_{u,j,i}(\tau)=e_j(\tau)-e_q(\tau)-(j-1-q)(t+1)$ for some $q<j-1$. The gaps of
Definition~\ref{th-hc-parameter}(4) give $e_j(\tau)-e_q(\tau)\ge(j-1-q)(t+3)$, so this is
nonnegative; and $e_j(\tau)-b_q(\tau)\le\eta$ with $e_q(\tau)-b_q(\tau)\ge(t+3)/2$ gives
$n_{u,j,i}(\tau)\le\eta$. There are $(j-1)(t+1)\le m(t+1)$ such indices, so their sum is at most
$m(t+1)\eta$. For $i<k$, $n_{u,j,i}(\tau)=u_{i+1}+e_j(\tau)-\frac12(N(j-1)+N(j)-1-2i)\ge
u_{i+1}+e_j(\tau)-N(m)+\frac12$, for the entries $u$ of $\lambda_u(\tau)$, and $u_{i+1}+e_j(\tau)\ge
N(m)$: for $i<k-1$, $u_{i+1}\ge0$ (as $u_1>v_1\ge0$ and the other entries lie in $[0,S_c]$) and
$e_j(\tau)\ge N(m)$; for $i=k-1$ at the top embedding of $t$, $u_k=-b_t-\frac{k-1}2$ and
$e_j(\tau)\ge b_t+b'/x\ge b_t+N(m)+k$; and at the other embeddings $u_k=v_{k-1}-\frac12\ge-\frac12$ and
$e_j(\tau)\ge N(m)+1$. So the row exponents are nonnegative.

Finally $r_E\ge2$, $S_c\le b$ and $a=xb\ge(3+|C_c|)b$. With $e_j(\tau)\le2b$ and $S_w\le(1+|C_c|)b$
this gives $e_j(\tau)+S_w\le ar_E$; and the window $a_t\le a+a/(kx^{13})\le a+b$ with $x\ge3$ gives
$a_t+e_j(\tau)\le a+3b\le ar_E$.
\end{proof}

\begin{lemma}\label{asm-division-kernel-record-hook}
Let $\lambda$ be a Young diagram with at most $n$ rows, of lengths
$\lambda_0\ge\lambda_1\ge\dots\ge\lambda_{n-1}\ge0$, and let $H(\lambda)$ be the product of its hook
lengths. Then
\[
 \sum_{i<n}\log(\lambda_i!)\le\log H(\lambda)
 =\sum_{i<n}\log\Gamma(\lambda_i+n-i)-\sum_{0\le i<l<n}\log(\lambda_i-\lambda_l+l-i).
\]
\end{lemma}
\leanref{Formalism/QuotientMetricDivisors/HermitianTwistDivisors/YoungDiagramHook.lean}{Formalism.DivisionValues.young_diagram_hook}
\leanref{Formalism/LocalToGlobal/Divisibility/NumbersLemmas.lean}{Formalism.L2.DIV.NumDatum.hookLower_of}
\leanref{Formalism/LocalToGlobal/Divisibility/RecBridgeNum.lean}{Formalism.L2.DIV.NumTieH.logHook_eq}
\begin{proof}
The cell in row $i$ and column $c<\lambda_i$ has hook length $(\lambda_i-c)+(\lambda'_c-i)-1$, where
$\lambda'_c>i$ is the length of column $c$; so it is at least $\lambda_i-c$. The product of the
$\lambda_i-c$ over the cells of row $i$ is $\lambda_i!$, so $\prod_i\lambda_i!\le H(\lambda)$. The
equality is the hook length formula \cite[Chapter~I, \S1, Example~1]{Mac95}, with
$\Gamma(\lambda_i+n-i)=(\lambda_i+n-1-i)!$.
\end{proof}

\begin{lemma}\label{asm-division-kernel-record-phi}
Let $\phi(y)=y\log y-y$ for $y>0$ and $\phi(0)=0$. Let $A$ be real with $L=\log A\ge1$, and let
$y,y'\ge0$ be integers. Then:
\begin{enumerate}
\item $\phi(y)\le\log(y!)$ and $-\phi(y)\le1$;
\item if $y\le A$, then $\phi(y)\le yL$;
\item if $y\le A$ and $|y-y'|\le d$, then $\phi(y)-\phi(y')\le dL$.
\end{enumerate}
\end{lemma}
\leanref{Formalism/QuotientMetricDivisors/HermitianTwistDivisors/YoungDiagramHook.lean}{Formalism.DivisionValues.W8.ph_nat_le}
\leanref{Formalism/QuotientMetricDivisors/HermitianTwistDivisors/YoungDiagramHook.lean}{Formalism.DivisionValues.W8.ph_nat_lt}
\leanref{Formalism/QuotientMetricDivisors/HermitianTwistDivisors/YoungDiagramHook.lean}{Formalism.DivisionValues.W8.pair_abs}
\leanref{Formalism/QuotientMetricDivisors/HermitianTwistDivisors/YoungDiagramHook.lean}{Formalism.DivisionValues.W8.log_nat_le}
\begin{proof}
(1) For $y\ge1$, $y^y/y!\le e^y$, a term of the exponential series, so $y\log y-\log(y!)\le y$; and
$\log y\ge1-1/y$, so $y\log y\ge y-1$.

(2) If $y\ge1$, then $\log y\le L$, so $\phi(y)\le y\log y\le yL$.

(3) If $y=0$, then $\phi(0)-\phi(y')=y'-y'\log y'\le y'\le d\le dL$. Let $y\ge1$. For $u>0$ and
$v\ge0$, $\phi(u)-\phi(v)\le(u-v)\log u$, because $v(\log u-\log v)\le u-v$ by $\log z\le z-1$ at
$z=u/v$ when $v>0$. With $0\le\log y\le L$ and $y-y'\le d$ this gives
$\phi(y)-\phi(y')\le d\log y\le dL$.
\end{proof}

\begin{definition}\label{asm-division-kernel-record-shapes}
Let $n,w\ge0$ be integers and $S,B,h$ real numbers, and let $\nu$ and $\alpha$ be finite multisets
of nonnegative integers. Write $|\nu|$ for the sum of $\nu$.
\begin{enumerate}
\item $(\nu,\alpha)$ has \emph{lower shape} $(n,w;S,B,h)$ if there are nonnegative integers
  $c_0,\dots,c_n$, $c'_0,\dots,c'_{n-1}$ and $r_0,\dots,r_{n-1}$, and multisets $\beta,\sigma$ of
  nonnegative integers, with
  \[
   c_i\le S,\qquad c_{i+1}\le c'_i\le c_i,\qquad |r_i-(w+c'_i)|\le h,\qquad |\beta|\le B,
  \]
  $\nu=\{w+c_0,\dots,w+c_n\}+\beta$ and $\alpha=\{r_0,\dots,r_{n-1}\}+\sigma$.
\item For an integer $a\ge0$, $(\nu,\alpha)$ has \emph{top shape} $(n,a,w;S,B,h)$ if there are
  nonnegative integers $c_0,\dots,c_{n-1}$, $c'_0,\dots,c'_n$, $r_0,\dots,r_n$ and $r'\le w$, and
  multisets $\beta,\sigma$ of nonnegative integers, with
  \[
   c'_i\le S,\qquad c'_{i+1}\le c_i\le c'_i,\qquad |r_i-(w+c'_i)|\le h,\qquad |\beta|\le B,
  \]
  $\nu=\{a+w,r'\}+\{w+c_0,\dots,w+c_{n-1}\}+\beta$ and $\alpha=\{r_0,\dots,r_n\}+\sigma$.
\end{enumerate}
\end{definition}
\leanref{Formalism/QuotientMetricDivisors/TheoremDivIneqRows.lean}{Formalism.DivisionValues.IsRowsOff}
\leanref{Formalism/QuotientMetricDivisors/TheoremDivIneqRows.lean}{Formalism.DivisionValues.IsRowsTop}
\leanref{Formalism/LogFactorialIndices/StatementSize.lean}{Formalism.LatticeIndexTer.Interlaced}

\begin{lemma}\label{asm-division-kernel-record-rows}
Let $n,w\ge0$ be integers, $S,B,h$ real numbers and $\nu,\alpha$ finite multisets of nonnegative
integers. Let $A$ and $F$ be real numbers with $L=\log A\ge1$, $w+S\le A$, $B\le A$ and
$\sum_{y\in\alpha}\log(y!)\le F$, and put $\Phi(\nu)=\sum_{y\in\nu}\phi(y)$, with $\phi$ as in
Lemma~\ref{asm-division-kernel-record-phi}.
\begin{enumerate}
\item If $(\nu,\alpha)$ has lower shape $(n,w;S,B,h)$
  (Definition~\ref{asm-division-kernel-record-shapes}), then
  \[
   \Phi(\nu)-F\le(w+2S+nh+B)L\qquad\text{and}\qquad0\le|\nu|-(n+1)w\le(n+1)S+B.
  \]
\item If $(\nu,\alpha)$ has top shape $(n,a,w;S,B,h)$ and $a+w\le A$, then
  \[
   \Phi(\nu)-F\le\phi(a)+(2w+S+nh+B)L+1\qquad\text{and}\qquad-w\le|\nu|-a-(n+2)w\le nS+B.
  \]
\end{enumerate}
\end{lemma}
\leanref{Formalism/QuotientMetricDivisors/HermitianTwistDivisors/YoungDiagramHook.lean}{Formalism.DivisionValues.young_diagram_factorial}
\leanref{Formalism/QuotientMetricDivisors/HermitianTwistDivisors/YoungDiagramHook.lean}{Formalism.DivisionValues.young_diagram_period}
\leanref{Formalism/QuotientMetricDivisors/TheoremDivIneqRows.lean}{Formalism.DivisionValues.RowsFactorial}
\leanref{Formalism/QuotientMetricDivisors/TheoremDivIneqRows.lean}{Formalism.DivisionValues.RowsPeriod}
\begin{proof}
We use Lemma~\ref{asm-division-kernel-record-phi}. By its part (1), $\phi(y)\le\log(y!)$ and
$\log(y!)\ge0$, so $\sum_i\phi(r_i)\le\sum_{y\in\alpha}\log(y!)\le F$ for the integers $r_i$ of the
shape. Every $y\in\beta$ has $y\le|\beta|\le B\le A$, so $\Phi(\beta)\le|\beta|L\le BL$ by part (2).

(1) For $i<n$, $|(w+c_i)-r_i|\le(c_i-c'_i)+h$ and $w+c_i\le w+S\le A$, so part (3) gives
$\phi(w+c_i)-\phi(r_i)\le(c_i-c'_i+h)L$. The sum telescopes:
$\sum_{i<n}(c_i-c'_i)\le\sum_{i<n}(c_i-c_{i+1})=c_0-c_n\le S$. By part (2),
$\phi(w+c_n)\le(w+S)L$. Adding,
$\Phi(\nu)-F\le(S+nh)L+(w+S)L+BL$. Also $|\nu|-(n+1)w=\sum_ic_i+|\beta|$, which lies between $0$ and
$(n+1)S+B$.

(2) For $i<n$, $|(w+c_i)-r_{i+1}|\le(c_i-c'_{i+1})+h$ and $w+c_i\le w+c'_i\le A$, so
$\phi(w+c_i)-\phi(r_{i+1})\le(c_i-c'_{i+1}+h)L$, and $\sum_{i<n}(c_i-c'_{i+1})\le\sum_{i<n}(c'_i-c'_{i+1})
=c'_0-c'_n\le S$. Part (3) with $|(a+w)-a|=w$ gives $\phi(a+w)-\phi(a)\le wL$; part (2) gives
$\phi(r')\le r'L\le wL$; and part (1) gives $-\phi(r_0)\le1$. Adding,
$\Phi(\nu)-F\le\phi(a)+2wL+(S+nh)L+BL+1$. Also $|\nu|-a-(n+2)w=(r'-w)+\sum_ic_i+|\beta|$, with
$-w\le r'-w\le0$ and $0\le\sum_ic_i\le nS$, as $c_i\le c'_i\le S$.
\end{proof}

\begin{lemma}\label{asm-division-kernel-record-origin}
Let $k\ge2$, $t\ge0$ and $1\le j\le m$ be integers, put $N(i)=k+i(t+1)$, $n_V=k+m-1+t$,
$s=\frac12\big(k+(2j-1)(t+1)\big)$ and $T=(j-1)t$, and suppose that $s$ is an integer. Let $w\ge s$
be an integer and put $w_r=w-s$. Let $u_1,\dots,u_k$ and $v_1,\dots,v_{k-1}$ be real numbers with
$v_i<u_i$ and $u_{i+1}<v_i$ for $1\le i\le k-1$, and let $\nu_0,\dots,\nu_{k-1}$ and
$\alpha_0,\dots,\alpha_{k-2}$ be nonnegative integers with
\[
 \nu_i=u_{i+1}+w-\tfrac12\big(N(j-1)+N(j)-1-2i\big),\qquad
 \alpha_i=v_{i+1}+w-\tfrac12(n_V-1-2i)-\tfrac12\big((k+j-1)-(m-j)\big).
\]
Let $\beta$ and $\sigma$ be finite multisets of nonnegative integers with $|\beta|\le B$, let $S$ and
$h\ge T$ be real, and put $\nu=\{\nu_0,\dots,\nu_{k-1}\}+\beta$ and
$\alpha=\{\alpha_0,\dots,\alpha_{k-2}\}+\sigma$.
\begin{enumerate}
\item If $u_k+\frac{k-1}2\ge0$ and $u_1-\frac{k-1}2\le S$, then $(\nu,\alpha)$ has lower shape
  $(k-1,w_r;S,B,h)$ (Definition~\ref{asm-division-kernel-record-shapes}).
\item If $u_1=a+\frac{k-1}2$ and $u_k=-b-\frac{k-1}2$ for integers $a,b\ge0$, and
  $v_{k-1}+\frac{k-2}2\ge0$ and $v_1-\frac{k-2}2\le S$, then $(\nu,\alpha)$ has top shape
  $(k-2,a,w_r;S,B,h)$.
\end{enumerate}
\end{lemma}
\leanref{Formalism/QuotientMetricDivisors/HermitianTwistDivisors/YoungDiagramTopExact.lean}{Formalism.DivisionValues.young_diagram_origin_off}
\leanref{Formalism/QuotientMetricDivisors/HermitianTwistDivisors/YoungDiagramTopExact.lean}{Formalism.DivisionValues.young_diagram_origin_top}
\leanref{Formalism/QuotientMetricDivisors/HermitianTwistDivisors/YoungDiagramTopExact.lean}{Formalism.DivisionValues.RowsOrigin.origin_rewrite}
\leanref{Formalism/QuotientMetricDivisors/HermitianTwistDivisors/YoungDiagramTopExact.lean}{Formalism.DivisionValues.RowsOrigin.origin_interlace}
\leanref{Formalism/QuotientMetricDivisors/HermitianTwistDivisors/YoungDiagramTopExact.lean}{Formalism.DivisionValues.RowsOrigin.off_core}
\leanref{Formalism/QuotientMetricDivisors/HermitianTwistDivisors/YoungDiagramTopExact.lean}{Formalism.DivisionValues.RowsOrigin.top_core}
\begin{proof}
Since $w=w_r+s$, $N(j-1)+N(j)=2s+k$ and $s-\frac12(n_V-1)-\frac12\big((k+j-1)-(m-j)\big)=T+1-\frac k2$,
\[
 \nu_i=w_r+u_{i+1}-\tfrac12(k-1-2i),\qquad\alpha_i=w_r+T+v_{i+1}-\tfrac12(k-2-2i).
\]
So $\alpha_i-\nu_i-T=v_{i+1}-u_{i+1}+\frac12<\frac12$ and $\nu_{i+1}+T-\alpha_i=u_{i+2}-v_{i+1}+\frac12
<\frac12$ for $0\le i\le k-2$. These are integers, so $\alpha_i\le\nu_i+T$ and
$\nu_{i+1}+T\le\alpha_i$. In particular the $\nu_i$ and the $\alpha_i$ are non-increasing in $i$.

(1) Here $\nu_{k-1}=w_r+u_k+\frac{k-1}2\ge w_r$ and $\nu_0=w_r+u_1-\frac{k-1}2\le w_r+S$. Put
$c_i=\nu_i-w_r$, $c'_i=\alpha_i-w_r-T$ and $r_i=\alpha_i$. Then $0\le c_{k-1}\le c_i\le c_0\le S$,
$c_{i+1}\le c'_i\le c_i$, so the $c'_i$ are nonnegative integers, and $|r_i-(w_r+c'_i)|=T\le h$.

(2) Here $\nu_0=w_r+a$, $\nu_{k-1}=w_r-b\le w_r$, $\alpha_{k-2}=w_r+T+v_{k-1}+\frac{k-2}2\ge w_r+T$
and $\alpha_0=w_r+T+v_1-\frac{k-2}2\le w_r+T+S$. Put $c'_i=\alpha_i-w_r-T$ and $r_i=\alpha_i$ for
$i\le k-2$, $c_i=\nu_{i+1}-w_r$ for $i\le k-3$, and $r'=\nu_{k-1}$. Then
$0\le c'_{k-2}\le c'_i\le c'_0\le S$, $c'_{i+1}\le c_i\le c'_i$ by the two inequalities above, and
$|r_i-(w_r+c'_i)|=T\le h$; and $\nu=\{w_r+a,r'\}+\{w_r+c_i\}+\beta$.
\end{proof}

\begin{lemma}\label{asm-division-kernel-record}
Let $\mathcal F$ be a family with exponent $\varepsilon$ (Definition~\ref{pd-family}) and size
constant $C_{sz}\ge0$, with the constants $C_c$ and $C_{YL}$ of its size data and parameter choices
and the exponent $C_t$ of the bound $t\le x^{C_t}$ of its size data, and put $c_t=\max(C_t,5)+3$.
Let $C\ge0$, $C_H$, $C_e$ and $C_d\ge0$ be real, and put
\[
 c_0=\max\big(28,\,2(1+|C_c|)+1,\,C_{sz}(3+2C_{YL}+c_t)\big),\qquad
 c_e=\max\big(c_0,\,2C+C_e+9C_d+5\big).
\]
There is $x_0$, depending on $\varepsilon$, $C_t$, $C_c$, $C_H$ and $c_e$, with the following
property. Let $D$ be a member at a scale $x\ge x_0$ with $x>x_1$, with coefficient field $L$, ranks
$r$ and $f$, level primes $g_0,g_1$, definite height $(h_1,h_2)$
(Definition~\ref{asm-definite-side-hypotheses}), and $\log D_\Gamma(u)\le C_dfN_m\log(fN_m)$ for
every word $u$. Let $P$ be a row-side family for $D$ with finitely many tuples, whose row side is
determined by a twist record of $D$, with group $G$ of order
$\prod_{k\in\{0,1\}}|GL_r(\mathbf F_{g_k})|^f$. Let a division package over $L$
(Definition~\ref{is-division-package}) be given on definite subspaces of $P$, for the numerical
bounds $N(C,C_H)$ of the scale data of the member with constant $C_{sz}$. Then there are real numbers
$S_w$ and $C_1$ and an integer $d_g>0$, divisible by the stabilizer denominators of the lift
lattices of $P$, such that, with $\lambda_0=\log|G|$:
\begin{enumerate}
\item the record numbers of every word $u$ (Definition~\ref{is-record-numbers}), built with the
  definite height, hold with constants $(\lambda_0,c_0,c_e,C_{sz},S_w,C_1)$;
\item at every word $u$, the record numbers of $u$ and the package satisfy inputs (5), (7) and (8)
  of Definition~\ref{is-division-inputs}, with $L_G=\lambda_0$,
  $E_2(j)=C_ex^{\varepsilon}(B_0+N_c)+\log d_g$ and $E_3=1$ in (7) and (8).
\end{enumerate}
\end{lemma}
\leanref{Formalism/LocalToGlobal/Divisibility/ArithClauseRec.lean}{Formalism.L2.DIV.Arith4S.arithClause_rec}
\leanref{Formalism/LocalToGlobal/Divisibility/ConstantGeTwentyEightAlt.lean}{Formalism.L2.DIV.ArithAt}
\leanref{Formalism/LocalToGlobal/Divisibility/NumbersLink.lean}{Formalism.L2.DIV.NumDatum.NumFacts}
\leanref{Development/AutomorphicArithmetic/ThetaCorrespondence/ThetaPairLifts/ThirdPlaceSeries/GammaSizeOfParameters.lean}{Formalism.L2.DIV.Cons.gamma_size_of_sizeData_alt}
\leanref{Formalism/LocalToGlobal/Divisibility/EigenvalueSumLeClosedOf.lean}{Formalism.L2.DIV.Arith4S.arithAt_rec}
\leanref{Formalism/LocalToGlobal/Divisibility/EigenvalueSumLeClosedOf.lean}{Formalism.L2.DIV.arithAt_scale3}
\leanref{Formalism/LocalToGlobal/Divisibility/EigenvalueSumLeClosedOf.lean}{Formalism.L2.DIV.Arith4S.q_rec}
\leanref{Formalism/LocalToGlobal/Divisibility/DivisionHeightBoundStatement.lean}{Formalism.L2.DIV.errB_of_units}
\leanref{Formalism/LocalToGlobal/Divisibility/DivisibilityConstruction/RBLam.lean}{Formalism.L2.DIV.Cons.log_card_le_unitE}
\leanref{Formalism/LocalToGlobal/Divisibility/NumbersLink.lean}{Formalism.L2.DIV.NumDatum.stepE_num}
\leanref{Formalism/QuotientMetricDivisors/HermitianTwistDivisors/YoungDiagramTopExact.lean}{Formalism.DivisionValues.young_diagram_step_origin}
\leanref{Formalism/QuotientMetricDivisors/HermitianTwistDivisors/YoungDiagramTopExact.lean}{Formalism.DivisionValues.young_diagram_step_shift}
\leanref{Formalism/QuotientMetricDivisors/HermitianTwistDivisors/YoungDiagramHook.lean}{Formalism.DivisionValues.young_diagram_step}
\begin{proof}
Put $\lambda_0=\log|G|$, $S_w=(1+|C_c|)S_c$ and $C_1=2(1+|C_c|)$, and let $d_g$ be the least common
multiple of the stabilizer denominators of the lift lattices of $P$ at $(P_0,\psi_0)$, over the lifts
and the tuples; it is positive and divisible by each of them. Let $x_0$ be the largest of $16$,
$(48c_t/\varepsilon+6/\varepsilon+2)^{6/\varepsilon}$, $3+|C_c|$, $272+2c_e$,
$\lceil(9c_t+2)/\varepsilon\rceil!$ and $\lceil(\max(C_H,0)+1)/\varepsilon\rceil!$. Fix a member $D$ at a
scale $x\ge x_0$ with $x>x_1$ and $P$ and the package as in the statement, and put
$U=x^{\varepsilon}(B_0+N_c)$. Then $x\ge16$, $x\ge3+|C_c|$ and
$48c_t/\varepsilon+6/\varepsilon+2\le x^{\varepsilon/6}$, so
Lemmas~\ref{asm-division-kernel-record-scale} and~\ref{asm-division-kernel-record-numbers} apply, and
so does Lemma~\ref{asm-division-kernel-explicit-entries}.

\emph{Sizes of the units.} The sizes give $B_0\le\Pi^3b$ and $N_c\le\Pi^3b$, so $U\le2\Pi^4b$; and
$U\ge1$, as $B_0\ge1$ and $x^{\varepsilon}\ge1$. The order of $G$ and the size bound $g_1\le2g_0$
give $\log|G|\le2fr^2\log(2g_0)$, and the sizes give $\log(2g_0)\le2x^{2\varepsilon}$ and $2r^2\le b\le
m^2b$; so $\lambda_0\le2f\,m^2b\,x^{2\varepsilon}=2mx^{\varepsilon}B_0\le2mU\le4\Pi^5b$. The size
condition on stabilizer denominators
of the member, with the orders of the stabilizers of the lift lattices dividing $D_\Gamma(u_0)$, gives
$\log d_g\le(2+9C_d)U$.

\emph{Conclusion (1).} The constants satisfy $c_0\ge\max(1,C_{sz}(3+2C_{YL}+c_t),2(1+|C_c|)+1)$ and
$c_0\le c_e$, so Lemma~\ref{asm-division-kernel-record-numbers} gives conditions (1) to (6) of
Definition~\ref{is-record-numbers} for the record of every word $u$, with the further facts listed
there.

\emph{Condition (7).} Fix a word $u$, a lift $j$ and a place $z$, with embedding $\tau=\tau_z$, and
write $w=e_j(\tau)=w_{j,z}$, $N=N(m)$ and $A=ar_E$. By condition (2),
$L_a=\log A=\log x+x^{\varepsilon}+\log r_E\ge1$. Let $\nu$ be the multiset of the row exponents
$n_{u,j,i}(\tau)$, $i<N(j-1)$, which are nonnegative integers (Definition~\ref{pd-orbit-places}(2)).
The other $\nu_{j,z,i}$ of the record are $0$, and $\phi(0)=0$; so, with $\Phi$ as in
Lemma~\ref{asm-division-kernel-record-rows},
$\epsilon_{j,z}=\Phi(\nu)-\log H_{j,z}-[z\in T]\phi(a_z)$ and $\delta_{j,z}=|\nu|-[z\in T]a_z-kw$.

Let $\alpha$ be the multiset of the numbers $a_i$, $i<k+j-1$, of
Definition~\ref{asm-division-kernel-explicit-heights}(1) at $j$ and $\tau$. They are integers, since
$t$ is odd; they are non-increasing, since $\sigma_i(\tau)\ge\sigma_{i+1}(\tau)+1$ (proof of
Lemma~\ref{asm-division-kernel-explicit-terms}(3)); and the last is
$e_j(\tau)-b_j(\tau)-\frac{t+1}2\ge1$, by Definition~\ref{th-hc-parameter}(4). So
Lemma~\ref{asm-division-kernel-record-hook}, for the Young diagram with rows $a_i$, gives
$\sum_{y\in\alpha}\log(y!)\le\log\mathrm{hook}_j(\tau)=\log H_{j,z}$.

Put $s=\frac12\big(k+(2j-1)(t+1)\big)$, an integer as $k$ is even and $t$ odd, $w_r=w-s$ and
$B=m(t+1)\eta$. Lemma~\ref{asm-division-kernel-record-numbers} gives $s\le w$, and $s\le N$ and
$(j-1)t\le N$. For $i<k$, $n_{u,j,i}(\tau)=u_{i+1}(\tau)+w-\frac12\big(N(j-1)+N(j)-1-2i\big)$, with the
entries $u_i(\tau)$ of $\lambda_u(\tau)$; for $i\le k-2$, $\sigma_i(\tau)=v_{i+1}(\tau)$, so
$a_i=v_{i+1}(\tau)+w-\frac12(n_V-1-2i)-\frac12\big((k+j-1)-(m-j)\big)$. The parameter $\lambda_u(\tau)$
has the $v$-entries, $u_1$ and $u_k$ of $\lambda(\tau)$ (Definition~\ref{th-definite-datum}(2)), and
its $u$- and $v$-entries interlace (Definition~\ref{th-hc-parameter}(2)). The row exponents with
$i\ge k$ form a multiset of sum at most $B$, by Lemma~\ref{asm-division-kernel-record-numbers}. If
$z$ is not a top place, the parameter choices give $u_k(\tau)=v_{k-1}(\tau)-\frac12$, and the size
data give $0\le v_i(\tau)\le S_c$ and $u_1(\tau)-v_1(\tau)\le C_cS_c$; so $u_k(\tau)+\frac{k-1}2\ge0$
and $u_1(\tau)-\frac{k-1}2\le(1+|C_c|)S_c=S_w$. If $z$ is the top place of $t'\in T$, the size data
give $u_1(\tau)=a_{t'}+\frac{k-1}2$ and $u_k(\tau)=-b_{t'}-\frac{k-1}2$ with integers $a_{t'},b_{t'}\ge0$,
and $v_{k-1}(\tau)+\frac{k-2}2\ge0$ and $v_1(\tau)-\frac{k-2}2\le S_c\le S_w$. So
Lemma~\ref{asm-division-kernel-record-origin}, with $S=S_w$ and $h=N$, shows that $(\nu,\alpha)$ has
lower shape $(k-1,w_r;S_w,B,N)$ if $z$ is not a top place, and top shape $(k-2,a_z,w_r;S_w,B,N)$ if
it is.

Lemma~\ref{asm-division-kernel-record-numbers} gives $w_r+S_w\le w+S_w\le A$, $a_z+w_r\le a_z+w\le
A$ at the top places, and $B\le w''\le b\le A$. So Lemma~\ref{asm-division-kernel-record-rows}, with
$F=\log H_{j,z}$, gives
\[
 \epsilon_{j,z}\le(w_r+2S_w+(k-1)N+B)L_a,\qquad0\le|\nu|-kw_r\le kS_w+B
\]
if $z$ is not a top place, and
\[
 \epsilon_{j,z}\le(2w_r+S_w+(k-2)N+B)L_a+1,\qquad-w_r\le|\nu|-a_z-kw_r\le(k-2)S_w+B
\]
if it is. By Lemma~\ref{asm-division-kernel-record-numbers}, $C_1\ge2$, $2S_w=C_1S_c$ and
$C_1+1\le c_0$; and $2\le k\le N$, so $(k-1)N\le N^2$ and $1\le N^2L_a$. In both cases
\[
 \epsilon_{j,z}\le C_1\big(w_r+S_c+N^2+B\big)L_a\le c_0\big(w+S_c+N^2+m(t+1)\eta\big)L_a.
\]
In both cases also $\big||\nu|-[z\in T]a_z-kw_r\big|\le w_r+kS_w+B\le C_1(w_r+kS_c+B)$. Since
$\delta_{j,z}=\big(|\nu|-[z\in T]a_z-kw_r\big)-ks$ and $0\le ks\le kN$,
\[
 |\delta_{j,z}|\le C_1\big(w+k(N+S_c)+B+[z\in T]kb'\big)+kN
 \le c_0\big(w+k(N+S_c)+m(t+1)\eta+[z\in T]kb'\big).
\]
This is condition (7), with $C_0=c_0$, $S=S_c$ and $t_0=t$.
% Lean: the bounds sE and sdelta of RecNums.ofNum are NumDatum.stepE_num (NumbersLink.lean), which
% applies Formalism.DivisionValues.rowsStepEOrigin = young_diagram_step_origin
% (YoungDiagramTopExact.lean; statement RowsStepEOrigin in TheoremDivIneqRowsShift.lean) to the
% facts sw, rowsA, topsA, nu0, blk, inter, off and top of NumFacts, with C1two,
% constant_weight_sum_bound and constant_succ_le, and to the hook bounds hookLower_of rowsHook
% (Lemma asm-division-kernel-record-hook). young_diagram_step_origin turns the facts into the shapes
% (young_diagram_origin_off and young_diagram_origin_top, Lemma asm-division-kernel-record-origin) with
% cs = 1 and the shift wr = w - sN j, then applies young_diagram_step_shift, which applies
% young_diagram_step (RowsStepE: young_diagram_factorial, young_diagram_period, arith_fact_top,
% arith_fact_off and arith_per) with the constant C1 and adds k (w - wr) <= k N to delta. The record's
% rows are nuPad (zero beyond N(j-1)); phRows_nuRows and sum_nuRows identify the sums. Size: proof
% closure of stepE_num, about 700 lines in 5 files.

Fix a word $u$, let $R$ be its record numbers, and let $N=N(C,C_H)$ be the numerical bounds of the
scale data of the member, formed with $C_{sz}$. By the construction of these scale data, their
functions are $\ell_6^{m}(j)=\sum_\tau\log\mathrm{hook}_j(\tau)-C_{sz}fn_V^2L_p$, $\ell_6^{u}=0$,
$\ell_3^{m}(j)=\sum_\tau q_je_j(\tau)\log\Omega$ and $\ell_3^{u}=B_0\log\log x$, and $E_1^{u}=U$,
$\ell_a=\log|G|$, $\delta_e=8fr/\min(g_0,g_1)$, and $x$, $b$ are those of the member; $F_0$ is the local
term of the size data (Step 5 of the proof of Lemma~\ref{asm-division-kernel-assembly}).

\emph{Input (5).} $h_X(j)=\sum_\tau\log\mathrm{hook}_j(\tau)-C_{sz}fn_V^2L_p=\ell_6^{m}(j)-C\ell_6^{u}
=\ell_6(j)$ for every lift $j$.

\emph{Input (7).} Put $E_{\ell_3}=mC\ell_3^{u}$. Then
$\sum_j\ell_3(j)=\sum_j\big(\ell_3^{m}(j)+C\ell_3^{u}\big)=\sum_j\sum_\tau q_jw_{j,\tau}\log\Omega+E_{\ell_3}$.
For the budget: $\log\log x\le x^{\varepsilon}$, as $\log x\le x^{\varepsilon}$, so $C\ell_3^{u}\le CU$;
$L_G=\lambda_0\le2mU$; $E_1(j)=CE_1^{u}=CU$; $E_2(j)\le(C_e+2+9C_d)U$; and $E_3=1\le mU$. Hence
\[
 E_{\ell_3}+L_G+\sum_j\big(E_1(j)+E_2(j)\big)+E_3\le(2C+C_e+9C_d+5)mU\le c_emU=\mathrm{err}_B.
\]

\emph{Input (8).} Put $B=\sum_j\sum_\tau q_jw_{j,\tau}\log\Omega+2F^{\mathrm{pk}}_o+\mathrm{err}_B$, which by
(7) is at least $B_{\mathrm{rows}}$, and $Q=\max(1,B+h_N+28\ell_a)$. Then $Q\ge1$, $\delta_e\ge0$ and
$B_{\mathrm{rows}}+h_N+28\ell_a\le Q$. By conditions (3) and (4) of the record numbers,
$|q_j|\le k+m$, $0\le w_{j,\tau}\le2b$ and $|\log\Omega|\le L_a\le4\Pi$, so the main sum is at most
$m\cdot f\cdot2\Pi\cdot2b\cdot4\Pi\le16\Pi^4b$. As in Step 5 of the proof of
Lemma~\ref{asm-division-kernel-assembly}, $2F^{\mathrm{pk}}_o=F_o$, which is at most $256\Pi^9b$ by the
proof of Lemma~\ref{asm-division-kernel-sizes}(3). And $\mathrm{err}_B=c_emU\le2c_e\Pi^5b$. So
$B\le(272+2c_e)\Pi^9b\le b^2$, by Lemma~\ref{asm-division-kernel-sizes}(1) with $e=9$ and
Lemma~\ref{asm-division-kernel-record-scale}(2) with $c=1$ and $e=9c_t+1$. Also
$h_N=x^{C_H}b\le x^{\max(C_H,0)}b\le b^2$ by Lemma~\ref{asm-division-kernel-record-scale}(2), and
$28\ell_a\le28\cdot4\Pi^5b\le b^2$. So $Q\le3b^2$. Since $fr\le\Pi^2\le\Pi^9b$,
$72fr\le(272+2c_e)\Pi^9b\le b^2\le b^6$, and
\[
 3\delta_eQ=\frac{24frQ}{\min(g_0,g_1)}\le\frac{72frb^2}{\min(g_0,g_1)}\le\frac{b^8}{(kb)^8}\le1,
\]
by the size bound $(kb)^8\le g_i$. Finally $\ell_a=\log|G|=\lambda_0$, and $28\le c_0$.
% Lean: arithClause_rec with lam0 = log |G|, Sw = swRec L, C1 = constantOneChoice L, NumFacts from
% numFacts_rec, dg = dgOf, and ArithAt from arithAt_rec, at x >= x0F eps Ct Cc CH ce, where c0 = c0F
% Csz Cc CYL Ct and ce = ceF C Ce cD c0; the hypothesis DGamSize at level 1 with cD = 2 + 9 C_d comes from
% the size condition on stabilizer denominators by gamma_size_of_sizeData_alt (dGamSize_record_one).
% ArithAt is eig6 (input (5)), hG0, l3 (input (7) with L_G = lgDef R 1 = lam0 + log 1,
% E_2 = etaSumBound (divSize L Ce) dg 1 and E_3 = 1), q (input (8)), la_le and
% constant_ge_twenty_eight. Size: proof closure of arithClause_rec and gamma_size_of_sizeData_alt,
% 259 theorems, about 3,900 lines in 31 files.
\end{proof}

\begin{lemma}\label{asm-division-kernel-inputs-reduction}
Let $K$ be a totally complex number field with an involution $a\mapsto\bar a$, let $K_0\subset K$ be a
subfield with $\sigma(\bar a)=\overline{\sigma(a)}$ for every complex embedding $\sigma$ of $K$ and
every $a\in K_0$, and let $v_0$ be an infinite place of $K$ with
$\sigma_{v_0}(\bar a)=\overline{\sigma_{v_0}(a)}$ for all $a\in K$. Let $H$ be an $n\times n$ matrix
with entries in $K_0$ that is positive definite at every infinite place, let $P$ be a $D\times n$
matrix over $K$ with $PK^n=K^D$, and put $Q=(PH^{-1}\bar P^{\mathsf T})^{-1}$. Let $M\subset K^n$ be a
lattice (Definition~\ref{is-local-indices}), and let $l>0$ and $g\ge0$ be integers such that $l^gx\in M$ for every $x$
in the dual $M^\vee=\{x\in K^n:\bar z^{\mathsf T}Hx\in\mathcal O_K\text{ for all }z\in M\}$. Let $y$
be a finite place of $K$ and $e_1,\dots,e_D$ a local basis of $PM$ at $y$
(Definition~\ref{is-local-indices}). Then its Gram determinant $G$ for $Q$ satisfies
$v_y(G)\le gD\,v_y(l)$.
\end{lemma}
\leanref{Formalism/QuotientMetricDivisors/HermitianTwistDivisors/ReductionLemmaFirst.lean}{Formalism.DivisionValues.reduction_lemma_first}
\leanref{Formalism/QuotientMetricDivisors/Statement.lean}{Formalism.DivisionValues.LemmaRa}
\leanref{Development/AutomorphicArithmetic/IntegralCasts/CapelliIdentities/SelectionLemmas.lean}{UpperTheory.Reformulated.B.SelectionLemmas.main}
\begin{proof}
First, $Q$ is defined and invertible. The image $\sigma_{v_0}(H)$ is positive definite, so $H$ is
invertible and $\sigma_{v_0}(H^{-1})=\sigma_{v_0}(H)^{-1}$ is positive definite. Put
$A=\sigma_{v_0}(P)$. Since the involution is complex conjugation at $v_0$,
$\sigma_{v_0}(PH^{-1}\bar P^{\mathsf T})=A\,\sigma_{v_0}(H)^{-1}A^*$. As $P$ has rank $D$, $A^*z\ne0$
for every $z\ne0$, so $z^*A\,\sigma_{v_0}(H)^{-1}A^*z>0$. Hence $PH^{-1}\bar P^{\mathsf T}$ is
positive definite at $v_0$, and invertible.
% Lean: reduction_lemma_first (the statement LemmaRa) calls ReductionFirst.main = SelectionLemmas.main,
% which takes any H and P, with M only finitely generated, and uses neither K_0, v_0, the positivity of
% H nor the surjectivity of P. With Mathlib's conventions A^-1 = 0 for a singular A and v_y(0) = 0 it
% settles D = 0 (G = 1) and G = 0 at once, and otherwise derives from G != 0 that E, P H^-1 P^* and H
% are invertible. The paragraph above obtains these invertibilities from the hypotheses of LemmaRa
% instead; the rest of the proof is that of SelectionLemmas.main (exists_uniform, hcompl_spec, core,
% det_int).

Let $E$ be the matrix with columns $e_1,\dots,e_D$; it is invertible, and
$G=\det(\bar E^{\mathsf T}QE)\ne0$. The $\mathcal O_K$-span of the $e_a$ has the same localization
at $y$ as $PM$, and $M$ is finitely generated, so there is $s\in\mathcal O_K$ with $s\notin y$ and
$\bar s\notin y$ such that $sPz$ lies in that span for every $z\in M$. Put $B=E^{-1}P$; then the
vector $sBz$ is integral for every $z\in M$.

Let $B_j$ be the $j$-th row of $B$ and put $x_j=\bar sH^{-1}\bar B_j^{\mathsf T}$. For $z\in M$,
$\bar z^{\mathsf T}Hx_j=\bar s\,\bar z^{\mathsf T}\bar B_j^{\mathsf T}=\overline{s(Bz)_j}$, which is
integral; so $x_j\in M^\vee$, and $l^gx_j\in M$. Hence $sB(l^gx_j)$ is integral, and its $i$-th entry is
$s\bar sl^g\,(BH^{-1}\bar B^{\mathsf T})_{ij}$. So the matrix $s\bar sl^g\,BH^{-1}\bar B^{\mathsf T}$
has integral entries, and its determinant $a$ is integral. Now
$BH^{-1}\bar B^{\mathsf T}\cdot\bar E^{\mathsf T}QE=E^{-1}(PH^{-1}\bar P^{\mathsf T})QE=1$, so
$aG=(s\bar sl^g)^D$. Taking valuations, with $v_y(s)=v_y(\bar s)=0$ and $v_y(a)\ge0$, gives
$v_y(G)\le gD\,v_y(l)$.
\end{proof}

\begin{definition}\label{asm-division-kernel-inputs-gram-spaces}
Let $L$ be a number field with an involution $a\mapsto\bar a$, and let $L_0$ be the set of $a\in L$
with $\sigma(\bar a)=\overline{\sigma(a)}$ for every complex embedding $\sigma$ of $L$, a subfield
stable under the involution. For a matrix $E$ over $L$ write $E^*=\bar E^{\mathsf T}$.
\begin{enumerate}
\item A \emph{class-sum space} over $L$ consists of: a finite set $Y$; a finite-dimensional
  $L$-vector space $\mathcal V$ with a basis and a pairing $\langle\ ,\ \rangle$ that is
  conjugate-linear in the first variable and linear in the second; for each $y\in Y$ a finite group
  $\Gamma_y$ acting linearly on $\mathcal V$ by $\rho_y$; and $\mathcal O_L$-submodules
  $\mathcal N_y$ and $\mathcal D_y$ of $\mathcal V$, the \emph{integral} and the \emph{Weyl vectors}
  at $y$. Its \emph{forms} are the families $f=(f(y))_{y\in Y}$ of vectors $f(y)\in\mathcal V$
  fixed by $\Gamma_y$, with the \emph{class-sum pairing}
  $\langle f,g\rangle=\sum_y|\Gamma_y|^{-1}\langle f(y),g(y)\rangle$. We fix a basis of the forms,
  taken with all coordinates of its values in $L_0$ if the forms have such a basis, and write $n$ for
  its size, $\phi\colon L^n\to\{\text{forms}\}$ for its coordinate map, $H$ for the matrix of the
  class-sum pairing in it, so that $x^*Hx'=\langle\phi(x),\phi(x')\rangle$, and $d_\Gamma$ for the
  least common multiple of the $|\Gamma_y|$. The \emph{integral lattice} and the \emph{Weyl lattice}
  are
  \[
   \mathcal L=\{x\in L^n:\phi(x)(y)\in\mathcal N_y\text{ for all }y\},\qquad
   \mathcal L^{\mathrm W}=\{x\in L^n:\phi(x)(y)\in\mathcal D_y\text{ for all }y\}.
  \]
\item A \emph{representation} of a finite group $G$ on the space is a representation $r$ of $G$ on
  its forms. We write $r(g)$ also for its matrix in the basis, and $r(a)=\sum_ga_gr(g)$ for
  $a=\sum_ga_gg\in L[G]$.
\item Let $s\ge1$ be an integer and $r$ a representation of $G$. The space satisfies the
  \emph{conditions at level $s$} if:
  \begin{enumerate}
  \item[(a)] $\langle x',x\rangle=\overline{\langle x,x'\rangle}$ for all $x,x'\in\mathcal V$, the
    pairing is nondegenerate, and every $\rho_y(\gamma)$ preserves it;
  \item[(b)] every $\mathcal N_y$ is a lattice in $\mathcal V$; $\mathcal D_y\subset\mathcal N_y$,
    $\mathcal D_y$ is stable under $\Gamma_y$, and $s\mathcal N_y^\vee\subset\mathcal D_y$ and
    $s\mathcal D_y\subset\mathcal N_y^\vee$, where
    $\mathcal N_y^\vee=\{x\in\mathcal V:\langle z,x\rangle\in\mathcal O_L\text{ for all }z\in\mathcal
    N_y\}$;
  \item[(c)] $H$ has entries in $L_0$ and is positive definite at every infinite place of $L$;
  \item[(d)] every $r(g)$ preserves the class-sum pairing and maps the forms with values in the
    $\mathcal N_y$, and those with values in the $\mathcal D_y$, to such forms; and the matrices
    $r(g)$ have entries in $L_0$;
  \item[(e)] the Hecke operators of the space commute in pairs and with $G$ and split over $L$, its
    blocks are well formed, its norms are positive and its places map injectively to the places of
    the family. These conditions are not used below.
  \end{enumerate}
\item A row-side family (Definition~\ref{is-row-side-family}) presents its row forms at every tuple
  and its lift forms at every lift and tuple as class-sum spaces over its coefficient field $L$, with
  representations of its group $G$. Its family of definite forms satisfies the \emph{family
  conditions at level $s$} if the involution of $L$ is complex conjugation at the base place $v_0$,
  there are finitely many tuples, the functions in the set $\mathcal X$ of characters of the family
  are the characters of absolutely irreducible representations of $G$, all these class-sum spaces
  satisfy the conditions at level $s$, and the conditions on the inert indices of the family hold.
\end{enumerate}
\end{definition}
\leanref{Formalism/LocalToGlobal/Divisibility/ClassSum.lean}{Formalism.L2.DIV.CSD.Good}
\leanref{Formalism/LocalToGlobal/Divisibility/GoodPartHermitianSpec.lean}{Formalism.L2.DIV.CSD.ActGood}
\leanref{Formalism/LocalToGlobal/Divisibility/GoodPartHermitianSpec.lean}{Formalism.L2.DIV.CSD.ActRat}
\leanref{Formalism/LocalToGlobal/Divisibility/ClassSum.lean}{Formalism.L2.DIV.CSD.DGam}
\leanref{Formalism/LocalToGlobal/Divisibility/RowLat.lean}{Formalism.L2.DIV.BlockDat.Ok}
\leanref{Formalism/LocalToGlobal/Divisibility/RowLat.lean}{Formalism.L2.DIV.DefFam.Ok}
\leanref{Formalism/LocalToGlobal/Divisibility/InertRationalAdjoint.lean}{Formalism.L2.DIV.DefFam3.Ok3}

\begin{definition}\label{asm-division-kernel-inputs-gram-side}
Let a class-sum space over $L$ with a representation $r$ of a finite group $G$ be given, with the
notation of Definition~\ref{asm-division-kernel-inputs-gram-spaces}. Let $\mathcal X$ be a finite set
of functions $G\to L$, $P\subset G$ a subgroup with a character $\psi\colon P\to L^\times$, and
$e=e_{\mathcal X}e_{P,\psi}\in L[G]$, with $e_{\mathcal X}$ and $e_{P,\psi}$ as in
Lemma~\ref{is-baselinks-group-algebra}(4).
\begin{enumerate}
\item The \emph{good part} at $(\mathcal X,P,\psi)$ is the image of $r(e)$ in $L^n$. Let $V_0$ be the
  matrix of a basis of it, taken with entries in $L_0$ if the good part has such a basis, and $n_V$
  its number of columns. Put
  \[
   H_V=|P|^{-1}V_0^*HV_0,\qquad M=\{x\in L^{n_V}:V_0x\in\mathcal L\}.
  \]
\item For a subspace $Z$ of the forms let $Z_0=\{x\in L^{n_V}:\phi(V_0x)\in Z\}$, let $E_Z$ be the
  matrix of a basis of $Z_0$, taken with entries in $L_0$ if possible, $d_Q$ its number of columns,
  and
  \[
   P_Z=(E_Z^*H_VE_Z)^{-1}E_Z^*H_V,
  \]
  with the inverse of a singular matrix taken to be $0$.
\item The \emph{definite side} of a row-side family at a level $(P_0,\psi_0)$ on definite subspaces
  $(Z_t)$ of its row forms has at every tuple $t$ the $n_V(t)$, $H_V(t)$, $M_t$, $d_Q(t)$ and
  $P_t=P_{Z_t}$ of the row forms at $t$ for $(\mathcal X,P_0,\psi_0)$.
\end{enumerate}
\end{definition}
\leanref{Formalism/LocalToGlobal/Divisibility/GoodPartHermitianSpec.lean}{Formalism.L2.DIV.CSD.V0E}
\leanref{Formalism/LocalToGlobal/Divisibility/GoodPartHermitianSpec.lean}{Formalism.L2.DIV.CSD.goodPartHermitianForm}
\leanref{Formalism/LocalToGlobal/Divisibility/GoodPartHermitianSpec.lean}{Formalism.L2.DIV.CSD.Mt}
\leanref{Formalism/LocalToGlobal/Divisibility/GoodPartHermitianSpec.lean}{Formalism.L2.DIV.orthPr}
\leanref{Formalism/LocalToGlobal/Divisibility/RowLat.lean}{Formalism.L2.DIV.BlockDat.Pq}
\leanref{Formalism/LocalToGlobal/Divisibility/RowLat.lean}{Formalism.L2.DIV.DefFam.side}

\begin{lemma}\label{asm-division-kernel-inputs-gram-idempotent}
Let $L$ be a number field with an involution $a\mapsto\bar a$ that is complex conjugation at an
infinite place $v_0$, with $L_0$ as in Definition~\ref{asm-division-kernel-inputs-gram-spaces}. Let
$G$ be a finite group, $\mathcal X$ a finite set of absolutely irreducible characters of $G$ with
values in $L$ (Lemma~\ref{is-onto-step-theta-j-idempotent}), $P\subset G$ a subgroup with a character
$\psi\colon P\to L^\times$, and $e_{\mathcal X}$ and $e_{P,\psi}$ as in
Lemma~\ref{is-baselinks-group-algebra}(4). For $a=\sum_ga_gg\in L[G]$ put $a^*=\sum_g\bar a_gg^{-1}$.
Then $e_{\mathcal X}$ is central, $he_{P,\psi}=\psi(h)e_{P,\psi}$ for $h\in P$, and each of
$a=e_{\mathcal X}$, $e_{P,\psi}$ and $e_{\mathcal X}e_{P,\psi}$, with $N$ equal to $|G|$, $|P|$ and
$|G||P|$ respectively, satisfies $a^2=a$ and $a^*=a$, has its coefficients $a_g$ in $L_0$, and has
$Na_g$ integral for every $g$.
\end{lemma}
\leanref{Formalism/LocalToGlobal/Divisibility/ProjLemmas.lean}{Formalism.L2.DIV.good_at_spec}
\leanref{Formalism/LocalToGlobal/Divisibility/ProjLemmas.lean}{Formalism.L2.DIV.idempotent_spec}
\leanref{Formalism/LocalToGlobal/Divisibility/ProjLemmas.lean}{Formalism.L2.DIV.level_idempotent_spec}
\leanref{Formalism/LocalToGlobal/Divisibility/GoodPartHermitianSpec.lean}{Formalism.L2.DIV.IsProjElt}
\begin{proof}
\emph{Values.} Let $\chi\in\mathcal X$ be the character of $\rho$, let $g\in G$ and let $\sigma$ be an
embedding of $L$ into $\mathbf C$. Extend $\sigma$ to the algebraic closure, and let $B$ be the image
of the matrix of $\rho(g)$, of some finite order $m$. The matrix $C=\sum_{k<m}(B^k)^*B^k$ is
positive definite and $B^*CB=C$, so $B^{-1}=C^{-1}B^*C$ and
$\sigma(\chi(g^{-1}))=\operatorname{tr}B^{-1}=\operatorname{tr}B^*=\overline{\sigma(\chi(g))}$. The
eigenvalues of $B$ are roots of unity, so $\chi(g)$ is integral. At $v_0$ this gives
$\overline{\chi(g)}=\chi(g^{-1})$, and then $\sigma(\overline{\chi(g)})=\overline{\sigma(\chi(g))}$
for every $\sigma$, so $\chi(g)\in L_0$. In the same way, for $h\in P$ the value $\psi(h)$ is a root
of unity, as $\psi(h)^{|P|}=1$, so it is integral, $\overline{\psi(h)}=\psi(h)^{-1}$, and
$\psi(h)\in L_0$.

\emph{$e_{\mathcal X}$.} It is a central idempotent by
Lemma~\ref{is-onto-step-theta-j-idempotent}(1) and~(3). Its coefficient at $g$ is
$\sum_\chi\chi(1)\chi(g^{-1})/|G|$, which lies in $L_0$, and $|G|$ times it is integral. As
$\overline{\chi(1)\chi(g^{-1})}=\chi(1)\chi(g)$, each $e_\chi$ satisfies $e_\chi^*=e_\chi$.

\emph{$e_{P,\psi}$.} For $h\in P$, $he_{P,\psi}=|P|^{-1}\sum_{h'}\psi(h')^{-1}hh'=\psi(h)e_{P,\psi}$,
so $e_{P,\psi}^2=|P|^{-1}\sum_h\psi(h)^{-1}\psi(h)e_{P,\psi}=e_{P,\psi}$. Its coefficients
$|P|^{-1}\psi(h)^{-1}$ lie in $L_0$, and $|P|$ times them are roots of unity. As
$\overline{\psi(h)^{-1}}=\psi(h)=\psi(h^{-1})^{-1}$, $e_{P,\psi}^*=e_{P,\psi}$.

\emph{The product.} As $e_{\mathcal X}$ is central,
$(e_{\mathcal X}e_{P,\psi})^2=e_{\mathcal X}^2e_{P,\psi}^2=e_{\mathcal X}e_{P,\psi}$, and as
$(ab)^*=b^*a^*$, $(e_{\mathcal X}e_{P,\psi})^*=e_{P,\psi}e_{\mathcal X}=e_{\mathcal X}e_{P,\psi}$. Its
coefficient at $t$ is $\sum_s(e_{\mathcal X})_s(e_{P,\psi})_{s^{-1}t}$, which lies in $L_0$, and
$|G||P|$ times it is $\sum_s\big(|G|(e_{\mathcal X})_s\big)\big(|P|(e_{P,\psi})_{s^{-1}t}\big)$, which
is integral.
\end{proof}

\begin{lemma}\label{asm-division-kernel-inputs-gram-projections}
Let $L$ be a number field with an involution.
\begin{enumerate}
\item Suppose that the involution is complex conjugation at an infinite place $v_0$. Let $H$ be an
  $n\times n$ matrix over $L$ that is positive definite at $v_0$, and $E$ an $n\times d$ matrix with
  independent columns. Then $H^*=H$, $H$ and $E^*HE$ are invertible, and $\Pi=(E^*HE)^{-1}E^*H$
  satisfies $\Pi E=1$ and $E^*H(1-E\Pi)=0$; it is the only $d\times n$ matrix with these two
  properties.
\item Let a class-sum space over $L$ with a representation $r$ of a finite group $G$ be given
  (Definition~\ref{asm-division-kernel-inputs-gram-spaces}), such that every $r(g)$ preserves the
  class-sum pairing and maps the forms with values in the $\mathcal N_y$, and those with values in
  the $\mathcal D_y$, to such forms (these two conditions are not used). Then $r(a)^*H=Hr(a^*)$ for
  every $a\in L[G]$, with $a^*$ as in Lemma~\ref{asm-division-kernel-inputs-gram-idempotent}.
\item If moreover the involution is complex conjugation at an infinite place $v_0$, $H$ is positive
  definite at $v_0$, $a\in L[G]$ satisfies $a^2=a$ and $a^*=a$, and $E$ is the matrix of a basis of
  the image of $r(a)$, then $E\Pi=r(a)$ for the $\Pi$ of (1).
\end{enumerate}
\end{lemma}
\leanref{Formalism/QuotientMetricDivisors/HermitianTwistDivisors/AuxStatement14.lean}{Formalism.DivisionValues.orthogonal_projection_exists}
\leanref{Formalism/LocalToGlobal/Divisibility/ProjLemmas.lean}{Formalism.L2.DIV.group_action_matrix_adjoint}
\leanref{Formalism/LocalToGlobal/Divisibility/ProjLemmas.lean}{Formalism.L2.DIV.orthogonal_projection_spec}
\begin{proof}
Let $\sigma_0$ be the embedding of $v_0$, so that $\sigma_0(\bar a)=\overline{\sigma_0(a)}$.

(1) $\sigma_0(H^*)=\sigma_0(H)^*=\sigma_0(H)$, as positive definite matrices are Hermitian, so
$H^*=H$; and $\det\sigma_0(H)>0$, so $H$ is invertible. If $E^*HEz=0$, then $u=Ez$ has
$\sigma_0(u)^*\sigma_0(H)\sigma_0(u)=\sigma_0(u^*Hu)=0$, so $u=0$ and $z=0$; so $E^*HE$ is
invertible. Then $\Pi E=1$ and $E^*HE\Pi=E^*H$. Any $\Pi'$ with the two properties has
$E^*HE\Pi'=E^*H$, so $\Pi'=\Pi$.

(2) As $r(g)$ preserves the class-sum pairing, $r(g)^*Hr(g)=H$, and $r(g)^{-1}=r(g^{-1})$, so
$r(g)^*H=Hr(g^{-1})$. Hence $r(a)^*H=\sum_g\bar a_gr(g)^*H=H\sum_g\bar a_gr(g^{-1})=Hr(a^*)$.

(3) Put $A=r(a)$. Then $A^2=A$, and $A^*H=HA$ by (2). As $A$ and $E$ have the same image, $A=EC$
and $E=AC'$ for some matrices $C,C'$, so $AE=A^2C'=E$. By (1), $E^*H=E^*HE\Pi$. Hence
\[
 HA=A^*H=C^*E^*H=C^*E^*HE\Pi=A^*HE\Pi=HAE\Pi=HE\Pi,
\]
and $H$ is invertible.
\end{proof}

\begin{lemma}\label{asm-division-kernel-inputs-gram-duality}
Let $L$ be a number field with an involution. For a finite-dimensional $L$-vector space $W$ with a
pairing $h$, conjugate-linear in the first variable, and an $\mathcal O_L$-submodule $A\subset W$,
write $A^\vee=\{x\in W:h(z,x)\in\mathcal O_L\text{ for all }z\in A\}$.
\begin{enumerate}
\item Let $V$ and $W$ carry such pairings $h_V$ and $h_W$ with $h(x',x)=\overline{h(x,x')}$, $h_V$
  nondegenerate, and let $E\colon W\to V$ be injective and $\Pi\colon V\to W$ surjective linear maps
  with $h_W(\Pi x,z)=h_V(x,Ez)$. Then $(E^{-1}A)^\vee=\Pi(A^\vee)$ for every lattice $A\subset V$.
\item Let a finite group $\Gamma$ act on $V$ preserving $h_V$, and let
  $\operatorname{av}=|\Gamma|^{-1}\sum_\gamma\gamma$. Then for every lattice $A\subset V$, the dual of
  $A\cap V^\Gamma$ in $V^\Gamma$, for $h_V$ restricted, is $\operatorname{av}(A^\vee)$.
\item Let $H$ be an invertible $n\times n$ matrix with $H^*=H$, and $E$, $\Pi$ matrices with $\Pi E=1$
  and $E^*H(1-E\Pi)=0$. For the pairings $x^*Hx'$ on $L^n$ and $z^*E^*HEz'$, and every lattice
  $A\subset L^n$, $(E^{-1}A)^\vee=\Pi(A^\vee)$.
\item For a matrix $E$ with independent columns and a lattice $A$, $E^{-1}A$ is a lattice. For a
  class-sum space whose $\mathcal N_y$ are lattices, $\mathcal L$ is a lattice.
\end{enumerate}
\end{lemma}
\leanref{Formalism/QuotientMetricDivisors/HermitianTwistDivisors/AuxStatement14.lean}{Formalism.DivisionValues.exponent_index_lattice.dual_comap}
\leanref{Formalism/QuotientMetricDivisors/HermitianTwistDivisors/AuxStatement14.lean}{Formalism.DivisionValues.exponent_index_lattice.inv_dual}
\leanref{Formalism/QuotientMetricDivisors/HermitianTwistDivisors/AuxStatement14.lean}{Formalism.DivisionValues.sectioned_dual_lattice}
\leanref{Formalism/QuotientMetricDivisors/HermitianTwistDivisors/AuxStatement14.lean}{Formalism.DivisionValues.exponent_index_lattice.cut_isLattice}
\leanref{Formalism/QuotientMetricDivisors/HermitianTwistDivisors/ClassSumDualGeOfFamily.lean}{Formalism.DivisionValues.ValuationLatticeField.valLatF_isLattice}
\leanref{Formalism/HermitianLattices/HermitianLatticesProofs/DualLattice.lean}{Formalism.HermitianLatticesTheory.dualLattice}
\begin{proof}
We use that for a nondegenerate pairing $h$ with $h(x',x)=\overline{h(x,x')}$ the dual of a lattice is
a lattice and $(A^\vee)^\vee=A$: localized at a finite place, $A$ is free, its dual is spanned by the
dual basis, and a lattice is determined by its localizations.

(1) For $x\in A^\vee$ and $z\in E^{-1}A$,
$h_W(z,\Pi x)=\overline{h_W(\Pi x,z)}=\overline{h_V(x,Ez)}=h_V(Ez,x)\in\mathcal O_L$, so
$\Pi(A^\vee)\subset(E^{-1}A)^\vee$. The pairing $h_W$ is nondegenerate: if $h_W(w,z)=0$ for all $w$,
then $h_V(x,Ez)=0$ for all $x$, as $\Pi$ is onto, so $Ez=0$ and $z=0$. The lattice $\Pi(A^\vee)$ has
dual $\{z:h_V(x,Ez)\in\mathcal O_L\text{ for }x\in A^\vee\}=E^{-1}((A^\vee)^\vee)=E^{-1}A$, so
$(E^{-1}A)^\vee=\Pi(A^\vee)^{\vee\vee}=\Pi(A^\vee)$.

(2) Apply (1) with $W=V^\Gamma$, $E$ the inclusion and $\Pi=\operatorname{av}$, which is onto as it
fixes $V^\Gamma$: for $z\in V^\Gamma$,
$h_V(\operatorname{av}x,z)=|\Gamma|^{-1}\sum_\gamma h_V(\gamma x,\gamma z)=h_V(x,z)$.

(3) Apply (1) with $E$ and $\Pi$: $E$ is injective and $\Pi$ onto as $\Pi E=1$, and taking $^*$ of
$E^*H=E^*HE\Pi$ gives $HE=\Pi^*E^*HE$, so $(\Pi x)^*E^*HEz=x^*HEz$.

(4) $E^{-1}A$ is finitely generated, as $E$ embeds it in $A$, and spans, since for every $z$ some
nonzero integer $m$ has $mEz\in A$. For $\mathcal L$ the same argument applies to the injective map
$x\mapsto(\phi(x)(y))_y$ and the lattice $\prod_y\mathcal N_y$.
\end{proof}

\begin{lemma}\label{asm-division-kernel-inputs-gram-hermitian}
Let $L$ be a number field with an involution that is complex conjugation at an infinite place
$v_0$, and let a class-sum space over $L$ with a representation $r$ of a finite group $G$ be given
(Definition~\ref{asm-division-kernel-inputs-gram-spaces}), such that $H$ has entries in $L_0$ and is
positive definite at every infinite place and the matrices $r(g)$ have entries in $L_0$. Let
$\mathcal X$ be a finite set of absolutely irreducible characters of $G$ with values in $L$, and
$P\subset G$ a subgroup with a character $\psi$. Then the matrices $V_0$ and $H_V$ of
Definition~\ref{asm-division-kernel-inputs-gram-side}(1) have entries in $L_0$, and $H_V$ is positive
definite at every infinite place.
\end{lemma}
\leanref{Formalism/LocalToGlobal/Divisibility/ProjLemmas.lean}{Formalism.L2.DIV.good_part_hermitian_spec}
\leanref{Formalism/LocalToGlobal/Divisibility/ProjLemmas.lean}{Formalism.L2.DIV.image_idempotent_rational}
\leanref{Formalism/LocalToGlobal/Divisibility/ProjLemmas.lean}{Formalism.L2.DIV.cut_posDef}
\begin{proof}
By Lemma~\ref{asm-division-kernel-inputs-gram-idempotent}, the coefficients of
$e=e_{\mathcal X}e_{P,\psi}$ lie in $L_0$, so $r(e)$ has entries in $L_0$. Its columns span the good
part, so some of them form a basis of it with entries in $L_0$, and $V_0$ has entries in $L_0$. As
$L_0$ is a subfield stable under the involution, $H_V=|P|^{-1}V_0^*HV_0$ has entries in $L_0$. Let
$\sigma$ be the embedding of an infinite place. As $V_0$ has entries in $L_0$,
$\sigma(V_0^*)=\sigma(V_0)^*$, so $\sigma(H_V)=|P|^{-1}\sigma(V_0)^*\sigma(H)\sigma(V_0)$; as
$\sigma(H)$ is positive definite and $\sigma(V_0)$ has independent columns, $\sigma(H_V)$ is positive
definite.
\end{proof}

\begin{lemma}\label{asm-division-kernel-inputs-gram-lattice}
Let $L$ be a number field with an involution that is complex conjugation at an infinite place $v_0$,
let $s\ge1$, and let a class-sum space over $L$ with a representation $r$ of a finite group $G$ be
given (Definition~\ref{asm-division-kernel-inputs-gram-spaces}) that satisfies conditions (a) and
(b) at level $s$, such that $H$ is positive definite at $v_0$ and every $r(g)$ preserves the
class-sum pairing and maps the forms with values in the $\mathcal N_y$, and those with values in the
$\mathcal D_y$, to such forms. Let $\mathcal X$ be a finite set of absolutely irreducible characters
of $G$ with values in $L$, $P\subset G$ a subgroup with a character $\psi$, and $V_0$, $H_V$ and $M$
as in Definition~\ref{asm-division-kernel-inputs-gram-side}(1). Then $\mathcal L$ and $M$ are
lattices, and $|G|s\,x\in M$ for every $x\in L^{n_V}$ with $z^*H_Vx\in\mathcal O_L$ for all $z\in M$.
(The inclusion $s\mathcal D_y\subset\mathcal N_y^\vee$ of (b), and the condition that $r(g)$ maps the
forms with values in the $\mathcal N_y$ to such forms, are not used.)
\end{lemma}
\leanref{Formalism/LocalToGlobal/Divisibility/ProjLemmas.lean}{Formalism.L2.DIV.local_bounds_first}
\leanref{Formalism/QuotientMetricDivisors/HermitianTwistDivisors/ClassSumDualGeOfFamily.lean}{Formalism.DivisionValues.class_sum_dual_le_fixed}
\leanref{Formalism/QuotientMetricDivisors/HermitianTwistDivisors/AuxStatement14.lean}{Formalism.DivisionValues.projection_integral}
\leanref{Formalism/QuotientMetricDivisors/HermitianTwistDivisors/AuxStatement14.lean}{Formalism.DivisionValues.divisor_bounds_first}
\begin{proof}
Write $\mathcal L^\vee$ for the dual of $\mathcal L$ for the pairing $x^*Hx'$.

\emph{Step 1: $s\mathcal L^\vee\subset\mathcal L^{\mathrm W}$.} Let $x\in\mathcal L^\vee$, put
$f=\phi(x)$, and fix $y_0\in Y$, with $\Gamma=\Gamma_{y_0}$. The vector $u=|\Gamma|^{-1}f(y_0)$ is
fixed by $\Gamma$. For every $\Gamma$-fixed $z\in\mathcal N_{y_0}$, the family $g$ with
$g(y_0)=z$ and $g(y)=0$ for $y\ne y_0$ is a form with coordinates in $\mathcal L$, so
$\langle g,f\rangle=\langle z,u\rangle\in\mathcal O_L$. Thus $u$ lies in the dual of
$\mathcal N_{y_0}\cap\mathcal V^\Gamma$ in $\mathcal V^\Gamma$, which by
Lemma~\ref{asm-division-kernel-inputs-gram-duality}(2) is $\operatorname{av}(\mathcal N_{y_0}^\vee)$:
$u=|\Gamma|^{-1}\sum_\gamma\rho_{y_0}(\gamma)\delta$ with $\delta\in\mathcal N_{y_0}^\vee$. So
$sf(y_0)=\sum_\gamma\rho_{y_0}(\gamma)(s\delta)$, which lies in $\mathcal D_{y_0}$, as
$s\delta\in\mathcal D_{y_0}$ and $\mathcal D_{y_0}$ is stable under $\Gamma$. Hence
$sx\in\mathcal L^{\mathrm W}$.

\emph{Step 2: the projection.} Let $e=e_{\mathcal X}e_{P,\psi}$, and let
$\Pi=(V_0^*HV_0)^{-1}V_0^*H$. By Lemma~\ref{asm-division-kernel-inputs-gram-idempotent}, $e^2=e$ and
$e^*=e$, so Lemma~\ref{asm-division-kernel-inputs-gram-projections}(3) gives $V_0\Pi=r(e)$. By the same
lemma the numbers $c_g=|G||P|e_g$ are integral. As every $r(g)$ maps $\mathcal L^{\mathrm W}$ into
itself, $|G||P|V_0\Pi w=\sum_gc_gr(g)w\in\mathcal L^{\mathrm W}\subset\mathcal L$ for every
$w\in\mathcal L^{\mathrm W}$, using $\mathcal D_y\subset\mathcal N_y$.

\emph{Step 3.} By Lemma~\ref{asm-division-kernel-inputs-gram-duality}(4), $\mathcal L$ and
$M=V_0^{-1}\mathcal L$ are lattices. Let $x$ be as in the statement. Then $x'=|P|^{-1}x$ satisfies
$z^*(V_0^*HV_0)x'\in\mathcal O_L$ for all $z\in M$. By
Lemma~\ref{asm-division-kernel-inputs-gram-projections}(1), $\Pi V_0=1$ and $V_0^*H(1-V_0\Pi)=0$, so
by Lemma~\ref{asm-division-kernel-inputs-gram-duality}(3) $x'=\Pi w$ with $w\in\mathcal L^\vee$. By
Step 1, $sw\in\mathcal L^{\mathrm W}$, and by Step 2,
$V_0(|G||P|s\,x')=|G||P|V_0\Pi(sw)\in\mathcal L$. As $|G||P|s\,x'=|G|s\,x$, this says
$|G|s\,x\in M$.
\end{proof}

\begin{lemma}\label{asm-division-kernel-inputs-gram}
Let $P$ be a row-side family for a definite-side datum (Definition~\ref{is-row-side-family}) whose
family of definite forms satisfies the family conditions at level $s\ge1$
(Definition~\ref{asm-division-kernel-inputs-gram-spaces}(4)), with group $G$, and let $S$ be an
indefinite-side datum whose definite side is that of the family at a level $(P_0,\psi_0)$ on definite
subspaces (Definition~\ref{asm-division-kernel-inputs-gram-side}(3)), with integer $D$, lattices $M_t$,
matrices $P_t$ and Hermitian matrices $H_V(t)$ (Definition~\ref{is-indefinite-datum}), and with
$d_Q(t)=D$ for every tuple $t$. For a tuple $t$ and a
finite place $y$ let $G_t(y)$ be the Gram determinant of a local basis of $\Lambda_t=P_tM_t$ at $y$
for $(P_tH_V(t)^{-1}\bar P_t^{\mathsf T})^{-1}$. Then $v_y(G_t(y))\le D\,v_y(s|G|)$.
\end{lemma}
\leanref{Formalism/LocalToGlobal/Divisibility/TwistedPlaceEtaSumLeExtended.lean}{Formalism.L2.DIV.gram_valuation_le_extended}
\leanref{Formalism/LocalToGlobal/Divisibility/RowLatLemmas.lean}{Formalism.L2.DIV.DefFam.GmOf_le}
\begin{proof}
Let $L$ be the coefficient field, $v_0$ the base place and $\mathcal X$ the characters of the family,
and fix a tuple $t$. By the family conditions at level $s$, the involution of $L$ is complex
conjugation at $v_0$, the elements of $\mathcal X$ are absolutely irreducible characters, and the row
forms at $t$ are a class-sum space with a representation of $G$ that satisfies the conditions at level
$s$. So:
\begin{enumerate}
\item[(i)] by Lemma~\ref{asm-division-kernel-inputs-gram-hermitian}, $H_V(t)$ has its entries in the
  field $L_0$ of Definition~\ref{asm-division-kernel-inputs-gram-spaces} and is positive definite at
  every infinite place;
\item[(ii)] by Lemma~\ref{asm-division-kernel-inputs-gram-projections}(1), applied to $H_V(t)$, which
  is positive definite at $v_0$, and to the matrix $E_{Z_t}$, whose columns are independent,
  $P_tE_{Z_t}=1$, so $P_t$ is surjective;
\item[(iii)] by Lemma~\ref{asm-division-kernel-inputs-gram-lattice}, $M_t$ is a lattice and
  $s|G|\,x\in M_t$ for every $x$ in the dual of $M_t$ for $H_V(t)$.
\end{enumerate}
By (ii) and (iii), $\Lambda_t=P_tM_t$ is a lattice, so it has a local basis at $y$. The coefficient
field $L$ is totally complex. So Lemma~\ref{asm-division-kernel-inputs-reduction}, with $K=L$,
$K_0=L_0$, this $v_0$, $H=H_V(t)$, $P=P_t$, $M=M_t$, $l=s|G|$ and $g=1$, gives
$v_y(G_t(y))\le d_Q(t)\,v_y(s|G|)=D\,v_y(s|G|)$.
% Lean: GmOf_le applies reduction_lemma_first with BlockDat.HV_rat and HV_pos (good_part_hermitian_spec,
% Lemma asm-division-kernel-inputs-gram-hermitian), Pq_surj (Pq_orth, from orthogonal_projection_exists),
% Mt_lat and Mt_dual (local_bounds_first, Lemma asm-division-kernel-inputs-gram-lattice, with the factor
% Fintype.card G * s) and Lam_lat (GluingInclusion.matLat_lattice') for the local basis locBasisOf_spec.
\end{proof}

\begin{definition}\label{asm-division-kernel-inputs-index-lifts}
Let a class-sum space over $L$ with a representation $r$ of a finite group $G$ be given, and let
$\mathcal X$, $P$, $\psi$ and a subspace $Z$ of the forms be given, with $V_0$, $H_V$, $E_Z$, $d_Q$ and
$P_Z$ as in Definition~\ref{asm-division-kernel-inputs-gram-side}. The \emph{level part} is the image of
$r(e_{P,\psi})$ in $L^n$; let $E_\ell$ be the matrix of a basis of it, taken with entries in $L_0$ if
possible, and put $\Pi_V=(V_0^*HV_0)^{-1}V_0^*H$. The \emph{lift data} of $(\mathcal X,P,\psi,Z)$ are
$d_Z=d_Q$, the Hermitian matrix $H_T=E_Z^*H_VE_Z$, the \emph{lift lattice}
$T=\{P_Z\Pi_VE_\ell x:E_\ell x\in\mathcal L\}$, the \emph{cut lattice}
$C=\{z\in L^{d_Z}:V_0E_Zz\in\mathcal L\}$, the \emph{Weyl cut lattice}
$D=\{z\in L^{d_Z}:V_0E_Zz\in\mathcal L^{\mathrm W}\}$, and the \emph{stabilizer denominator}
$d_G=|P|\,d_\Gamma$. The definite side of a row-side family at a level $(P_0,\psi_0)$ on definite
subspaces $(Z_{j,t})$ of its lift forms has at every lift $j$ and tuple $t$ the lift data
$d_Z(j,t)$, $H_T(j,t)$, $T_{j,t}$, $C_{j,t}$, $D_{j,t}$ and $d_G(j,t)$ of the lift forms at $(j,t)$ for
$(\mathcal X,P_0,\psi_0^{-1},Z_{j,t})$.
\end{definition}
\leanref{Formalism/LocalToGlobal/Divisibility/GoodPartHermitianSpec.lean}{Formalism.L2.DIV.CSD.lvlE}
\leanref{Formalism/LocalToGlobal/Divisibility/GoodPartHermitianSpec.lean}{Formalism.L2.DIV.CSD.lvlH}
\leanref{Formalism/LocalToGlobal/Divisibility/RowLat.lean}{Formalism.L2.DIV.BlockDat.HZ}
\leanref{Formalism/LocalToGlobal/Divisibility/RowLat.lean}{Formalism.L2.DIV.BlockDat.lPr}
\leanref{Formalism/LocalToGlobal/Divisibility/RowLat.lean}{Formalism.L2.DIV.BlockDat.Tl}
\leanref{Formalism/LocalToGlobal/Divisibility/RowLat.lean}{Formalism.L2.DIV.BlockDat.cutZ}
\leanref{Formalism/LocalToGlobal/Divisibility/LiftLat.lean}{Formalism.L2.DIV.BlockDat.DG}

\begin{lemma}\label{asm-division-kernel-inputs-index-general}
Let $L$ be a number field with an involution, let $H$ be an invertible $n\times n$ matrix with
$H^*=H$, and let $E$ and $\Pi$ be $n\times d$ and $d\times n$ matrices with $\Pi E=1$ and
$E^*H(1-E\Pi)=0$. Duals in $L^n$ and $L^d$ are taken for the pairings $x^*Hx'$ and $z^*E^*HEz'$. Let
$A,B\subset L^n$ be $\mathcal O_L$-submodules, $A$ a lattice, and $N\ge1$ an integer with
$NB\subset A^\vee$. Let $w$ be a finite place, $\varpi\in\mathcal O_L$ with $v_w(\varpi)=1$, and
$J\ge0$ an integer with $\varpi^JE^{-1}A\subset(E^{-1}B)_{(w)}$. Then $T=\Pi A$ and $C=E^{-1}A$ are
lattices, $C\subset T$, and
\[
 j_w(E^*HE,T)\le\operatorname{len}_w(T/C)+d\big(J+v_w(N)\big).
\]
\end{lemma}
\leanref{Formalism/QuotientMetricDivisors/HermitianTwistDivisors/AuxStatement14.lean}{Formalism.DivisionValues.bounds_fourth}
\leanref{Formalism/QuotientMetricDivisors/HermitianTwistDivisors/AuxStatement14.lean}{Formalism.DivisionValues.ValuationLatticeDual}
\leanref{Formalism/QuotientMetricDivisors/HermitianTwistDivisors/AuxStatement14.lean}{Formalism.DivisionValues.exponent_index_local}
\leanref{Formalism/QuotientMetricDivisors/HermitianTwistDivisors/AuxStatement14.lean}{Formalism.DivisionValues.projection_dual_lattice}
\begin{proof}
For lattices $X,Y$ write $\operatorname{idx}_w(X,Y)$ as in Definition~\ref{is-local-indices}. Writing
both lengths relative to a lattice contained in $X\cap Y\cap Y'$ shows
$\operatorname{idx}_w(X,Y')=\operatorname{idx}_w(X,Y)+\operatorname{idx}_w(Y,Y')$, and
$\operatorname{idx}_w(X,Y)=\operatorname{len}_w(X/Y)$ if $Y\subset X$.

\emph{The lattices.} $T$ is a lattice, as $\Pi$ is onto ($\Pi E=1$), and $C$ is a lattice by
Lemma~\ref{asm-division-kernel-inputs-gram-duality}(4), as $E$ is injective. For $z\in C$,
$z=\Pi(Ez)$ with $Ez\in A$, so $C\subset T$.

\emph{The dual of $T$.} Taking $^*$ of $E^*H=E^*HE\Pi$ gives $HE=\Pi^*E^*HE$, so
$(\Pi a)^*E^*HEz=a^*HEz$. Hence $z\in T^\vee$ if and only if $Ez\in A^\vee$: $T^\vee=E^{-1}(A^\vee)$.
It is a lattice, as the dual of a lattice for the pairing $z^*E^*HEz'$, which is nondegenerate as in
the proof of Lemma~\ref{asm-division-kernel-inputs-gram-duality}(1).

\emph{The exponent.} Put $a=\varpi^JN$. For $z\in C$ there is $u\in\mathcal O_L\smallsetminus w$ with
$u\varpi^Jz\in E^{-1}B$, so $uaEz=N\,E(u\varpi^Jz)\in NB\subset A^\vee$, and $uaz\in T^\vee$. So
$aC\subset(T^\vee)_{(w)}$. Then $(aC\cap T^\vee)_{(w)}=(aC)_{(w)}$, so
$\operatorname{idx}_w(aC,T^\vee)=-\operatorname{len}_w\big(T^\vee/(aC\cap T^\vee)\big)\le0$. As $a$ is
integral, $aC\subset C$, and $C_{(w)}$ is free of rank $d$ over the discrete valuation ring at $w$, so
$\operatorname{idx}_w(C,aC)=\operatorname{len}_w(C/aC)=d\,v_w(a)=d(J+v_w(N))$. Hence
$\operatorname{idx}_w(C,T^\vee)\le d(J+v_w(N))$.

\emph{The bound.} $j_w(E^*HE,T)=\operatorname{idx}_w(T,T^\vee)=\operatorname{idx}_w(T,C)+
\operatorname{idx}_w(C,T^\vee)\le\operatorname{len}_w(T/C)+d(J+v_w(N))$.
\end{proof}

\begin{lemma}\label{asm-division-kernel-inputs-index-classsum}
Let $L$ be a number field with an involution, $\mathcal V$ an $L$-vector space with a pairing
$\langle\ ,\ \rangle$, conjugate-linear in the first variable, $Y$ a finite set, and
$\phi\colon L^n\to\mathcal V^Y$ a linear map. Let $g_y>0$ be integers and $H$ an $n\times n$ matrix with
$x^*Hx'=\sum_yg_y^{-1}\langle\phi(x)(y),\phi(x')(y)\rangle$ for all $x,x'$. Let $\mathcal N_y$ and
$\mathcal D_y$ be $\mathcal O_L$-submodules of $\mathcal V$ and $s\ge1$ an integer with
$s\mathcal D_y\subset\mathcal N_y^\vee$ for every $y$, and let $d\ge0$ be an integer divisible by every
$g_y$. Then $ds\,x'$ lies in the dual, for $H$, of $\{x:\phi(x)(y)\in\mathcal N_y\text{ for all }y\}$
for every $x'$ with $\phi(x')(y)\in\mathcal D_y$ for all $y$.
\end{lemma}
\leanref{Formalism/QuotientMetricDivisors/HermitianTwistDivisors/ClassSumDualGeOfFamily.lean}{Formalism.DivisionValues.class_sum_dual_ge_fixed}
\begin{proof}
For such $x$ and $x'$,
$ds\,x^*Hx'=\sum_y(d/g_y)\langle\phi(x)(y),s\phi(x')(y)\rangle$, and each term is integral, as
$s\phi(x')(y)\in\mathcal N_y^\vee$ and $\phi(x)(y)\in\mathcal N_y$.
\end{proof}

\begin{lemma}\label{asm-division-kernel-inputs-index-level}
Let $L$ be a number field with an involution that is complex conjugation at an infinite place $v_0$,
let $s\ge1$, and let a class-sum space over $L$ with a representation $r$ of a finite group $G$ be
given that satisfies the conditions at level $s$
(Definition~\ref{asm-division-kernel-inputs-gram-spaces}), of which (a), (b) and (e) are not used.
Let $\mathcal X$ be a finite set of
absolutely irreducible characters of $G$ with values in $L$, $P\subset G$ a subgroup with a character
$\psi$, and $Z$ a subspace of the forms, with the notation of
Definitions~\ref{asm-division-kernel-inputs-gram-side} and~\ref{asm-division-kernel-inputs-index-lifts}.
Put $H_\ell=|P|^{-1}E_\ell^*HE_\ell$, $\Pi_\ell=(E_\ell^*HE_\ell)^{-1}E_\ell^*H$, $E=\Pi_\ell V_0E_Z$ and
$\Pi=P_Z\Pi_VE_\ell$. Then:
\begin{enumerate}
\item $x^*H_\ell x'=\sum_y(|P||\Gamma_y|)^{-1}\langle\phi(E_\ell x)(y),\phi(E_\ell x')(y)\rangle$;
\item $H_\ell$ is positive definite at every infinite place;
\item $E_\ell E=V_0E_Z$, $\Pi E=1$, $E^*H_\ell(1-E\Pi)=0$ and $E^*H_\ell E=H_T$;
\item $T=\Pi A$, $C=E^{-1}A$ and $D=E^{-1}B$, where $A=E_\ell^{-1}\mathcal L$ and
  $B=E_\ell^{-1}\mathcal L^{\mathrm W}$.
\end{enumerate}
\end{lemma}
\leanref{Formalism/LocalToGlobal/Divisibility/ProjLemmas.lean}{Formalism.L2.DIV.level_class_sum}
\leanref{Formalism/LocalToGlobal/Divisibility/LiftLatLemmas.lean}{Formalism.L2.DIV.BlockDat.lf_lvlH_pos}
\leanref{Formalism/LocalToGlobal/Divisibility/LiftLatLemmas.lean}{Formalism.L2.DIV.BlockDat.lf_lvl_vE}
\leanref{Formalism/LocalToGlobal/Divisibility/LiftLatLemmas.lean}{Formalism.L2.DIV.BlockDat.lf_lvlE_lE}
\leanref{Formalism/LocalToGlobal/Divisibility/LiftLatLemmas.lean}{Formalism.L2.DIV.BlockDat.lf_lPr_lE}
\leanref{Formalism/LocalToGlobal/Divisibility/LiftLatLemmas.lean}{Formalism.L2.DIV.BlockDat.lf_l_proj}
\leanref{Formalism/LocalToGlobal/Divisibility/LiftLatLemmas.lean}{Formalism.L2.DIV.BlockDat.lf_lE_form}
\leanref{Formalism/LocalToGlobal/Divisibility/LiftLatLemmas.lean}{Formalism.L2.DIV.BlockDat.lf_cut_nab}
\leanref{Formalism/LocalToGlobal/Divisibility/LiftLatLemmas.lean}{Formalism.L2.DIV.BlockDat.lf_cut_del}
\begin{proof}
(1) $x^*H_\ell x'=|P|^{-1}\langle\phi(E_\ell x),\phi(E_\ell x')\rangle$, and the class-sum pairing has
weights $|\Gamma_y|^{-1}$.

(2) By Lemma~\ref{asm-division-kernel-inputs-gram-idempotent} the coefficients of $e_{P,\psi}$ lie in
$L_0$, so $r(e_{P,\psi})$ and $E_\ell$ have entries in $L_0$, and the argument of
Lemma~\ref{asm-division-kernel-inputs-gram-hermitian} applies to $H_\ell$.

(3) Let $e=e_{\mathcal X}e_{P,\psi}$. By Lemmas~\ref{asm-division-kernel-inputs-gram-idempotent}
and~\ref{asm-division-kernel-inputs-gram-projections}(3), $E_\ell\Pi_\ell=r(e_{P,\psi})$ and
$V_0\Pi_V=r(e)$, and by Lemma~\ref{asm-division-kernel-inputs-gram-projections}(1), $\Pi_VV_0=1$,
$V_0^*H(1-V_0\Pi_V)=0$, and, $H_V$ being positive definite at $v_0$ by
Lemma~\ref{asm-division-kernel-inputs-gram-hermitian}, $P_ZE_Z=1$ and $E_Z^*H_V(1-E_ZP_Z)=0$. As
$e_{\mathcal X}$ is central and $e_{P,\psi}^2=e_{P,\psi}$, $e_{P,\psi}e=e$, so
\[
 E_\ell\Pi_\ell V_0=r(e_{P,\psi})V_0\Pi_VV_0=r(e_{P,\psi}e)V_0=r(e)V_0=V_0\Pi_VV_0=V_0.
\]
Hence $E_\ell E=V_0E_Z$, and $\Pi E=P_Z\Pi_VV_0E_Z=P_ZE_Z=1$. Put $W=\Pi_\ell V_0$, so that
$E_\ell W=V_0$ and $W^*E_\ell^*HE_\ell=V_0^*HE_\ell$; then $E^*H_\ell=|P|^{-1}E_Z^*V_0^*HE_\ell$, and
\[
 E^*H_\ell E\Pi=|P|^{-1}E_Z^*V_0^*HV_0E_ZP_Z\Pi_VE_\ell=E_Z^*H_VE_ZP_Z\Pi_VE_\ell
 =E_Z^*H_V\Pi_VE_\ell=|P|^{-1}E_Z^*V_0^*HE_\ell=E^*H_\ell.
\]
Finally $E^*H_\ell E=|P|^{-1}(E_\ell E)^*H(E_\ell E)=|P|^{-1}E_Z^*V_0^*HV_0E_Z=E_Z^*H_VE_Z=H_T$.

(4) The first equality is the definition of $T$, and the others hold because $E_\ell Ez=V_0E_Zz$.
\end{proof}

\begin{lemma}\label{asm-division-kernel-inputs-index}
Let $P$ be a row-side family for a definite-side datum whose family of definite forms satisfies the
family conditions at level $s\ge1$ (Definition~\ref{asm-division-kernel-inputs-gram-spaces}(4)). For a
lift $j$ and a tuple $t$ let $d_Z(j,t)$, $H_T(j,t)$, $T_{j,t}$, $C_{j,t}$, $D_{j,t}$ and $d_G(j,t)$ be
the lift data of the definite side of the family at a level $(P_0,\psi_0)$ on definite subspaces
(Definition~\ref{asm-division-kernel-inputs-index-lifts}): the dimension, the Hermitian matrix, the
lift lattice, the cut lattice, the Weyl cut lattice and the stabilizer denominator. For a finite place
$y$ let
$\eta^{\mathrm{fam}}_j(t,y)=\operatorname{len}_y\big(T_{j,t}/(T_{j,t}\cap C_{j,t})\big)$, and let
$J_y(j,t)$ be the least integer $J\ge0$ with $az\in(D_{j,t})_{(y)}$ for all $a\in y^J$ and
$z\in C_{j,t}$, as in Definition~\ref{is-onto-step-exp-bound-lift-exponent}. Then
\[
 j_y\big(H_T(j,t),T_{j,t}\big)\le\eta^{\mathrm{fam}}_j(t,y)
 +d_Z(j,t)\big(J_y(j,t)+v_y(d_G(j,t))+v_y(s)\big).
\]
\end{lemma}
\leanref{Formalism/LocalToGlobal/Divisibility/TwistedPlaceEtaSumLeExtended.lean}{Formalism.L2.DIV.lift_index_bound_extended}
\leanref{Formalism/LocalToGlobal/Divisibility/LiftLatLemmas.lean}{Formalism.L2.DIV.DefFam.jw_le}
\leanref{Formalism/LocalToGlobal/Divisibility/LiftLatLemmas.lean}{Formalism.L2.DIV.local_bounds_from_exponent}
\leanref{Formalism/LocalToGlobal/Divisibility/LiftLat.lean}{Formalism.L2.DIV.BoundsIVAt}
\begin{proof}
Fix $j$, $t$ and $y$. By the family conditions at level $s$, the involution of the coefficient field
$L$ is complex conjugation at the base place $v_0$, the characters in $\mathcal X$ are absolutely
irreducible, and the lift forms at $(j,t)$ are a class-sum space with a representation of $G$ that
satisfies the conditions at level $s$. Use the notation of
Lemma~\ref{asm-division-kernel-inputs-index-level} for these forms, $(\mathcal X,P_0,\psi_0^{-1})$ and
$Z=Z_{j,t}$, with $A=E_\ell^{-1}\mathcal L$ and $B=E_\ell^{-1}\mathcal L^{\mathrm W}$.

\emph{The exponent.} $C_{j,t}$ and $D_{j,t}$ are lattices, by
Lemma~\ref{asm-division-kernel-inputs-gram-duality}(4) applied to $\mathcal L$ and to
$\mathcal L^{\mathrm W}$; the $\mathcal D_y$ are lattices, as they lie in $\mathcal N_y$ and contain
$s\mathcal N_y^\vee$. So some nonzero $\delta\in\mathcal O_L$ has $\delta C_{j,t}\subset D_{j,t}$.
Write $\delta\mathcal O_L=y^c\mathfrak a$ with $y\nmid\mathfrak a$ and choose $u\in\mathfrak a\smallsetminus
y$; for $a\in y^c$, $au\in\delta\mathcal O_L$, so $u(az)\in D_{j,t}$ for $z\in C_{j,t}$. Hence an integer
$J$ as in the definition of $J_y(j,t)$ exists, and $J_y(j,t)$ is the least one. Choose
$\varpi\in\mathcal O_L$ with $v_y(\varpi)=1$. Then $\varpi^{J_y(j,t)}\in y^{J_y(j,t)}$, so
$\varpi^{J_y(j,t)}C_{j,t}\subset(D_{j,t})_{(y)}$.

\emph{The general bound.} By Lemma~\ref{asm-division-kernel-inputs-index-level}(2) and
Lemma~\ref{asm-division-kernel-inputs-gram-projections}(1), $H_\ell^*=H_\ell$ and $H_\ell$ is
invertible, and by part (3) of the former, $\Pi E=1$ and $E^*H_\ell(1-E\Pi)=0$. The module $A$ is a
lattice by Lemma~\ref{asm-division-kernel-inputs-gram-duality}(4). By
Lemma~\ref{asm-division-kernel-inputs-index-level}(1), $H_\ell$ is the class sum of the
$\phi(E_\ell x)$ with weights $|P_0||\Gamma_y|$, which divide $d_G(j,t)=|P_0|d_\Gamma$; with
$s\mathcal D_y\subset\mathcal N_y^\vee$, Lemma~\ref{asm-division-kernel-inputs-index-classsum} gives
$d_G(j,t)s\,B\subset A^\vee$ for $H_\ell$. By Lemma~\ref{asm-division-kernel-inputs-index-level}(4),
$E^{-1}A=C_{j,t}$ and $E^{-1}B=D_{j,t}$. So Lemma~\ref{asm-division-kernel-inputs-index-general}, with
$H=H_\ell$, this $E$ and $\Pi$, $N=d_G(j,t)s$, $w=y$, $\varpi$ and $J=J_y(j,t)$, gives, as
$\Pi A=T_{j,t}$ and $E^*H_\ell E=H_T(j,t)$,
\[
 j_y\big(H_T(j,t),T_{j,t}\big)\le\operatorname{len}_y(T_{j,t}/C_{j,t})
 +d_Z(j,t)\big(J_y(j,t)+v_y(d_G(j,t))+v_y(s)\big).
\]
As $C_{j,t}\subset T_{j,t}$, the first term is $\eta^{\mathrm{fam}}_j(t,y)$.
% Lean: DefFam.jw_le = local_bounds_from_exponent = boundsIVAt_of exponent_weight_spectral, for the
% block data F.L j t at the level lvlA P psi = (X, P, psi^-1) on Zl j t. boundsIVAt_of takes the
% uniformizer lf_exists_unif, H, det from orthogonal_projection_exists with lf_lvlH_pos, level_class_sum
% (weights card P * wt c), lf_exp_incl (from exponent_weight_spectral, whose existence part is
% expW_nonempty, with lf_cutZ_isLattice and lf_delZ_isLattice) rewritten by lf_cut_nab and lf_cut_del,
% and bounds_fourth_class_fixed (class_sum_dual_ge_fixed and bounds_fourth) with lf_l_proj, then
% rewrites by lf_cut_nab and lf_lE_form. bounds_fourth_class_fixed splits vw (DG * s) as vw DG + vw s.
\end{proof}

\begin{lemma}\label{asm-division-kernel-inputs}
Let $D$ be a member at a scale $x>x_1$ of a family $\mathcal F$ with exponent $\varepsilon$, with
size data, and let $P$ and $Q$ be as in Lemma~\ref{asm-division-kernel-package}, whose hypotheses
hold with the family conditions at level $1$, with $P$ as in Lemma~\ref{asm-division-kernel-record}.
Let $\mathcal P$ be the division package
of Lemma~\ref{asm-division-kernel-package}, for the scale data of the member, and let $\lambda_0$,
$S_w$, $C_1$, $d_g$ and the record numbers of the words satisfy the conclusions of
Lemma~\ref{asm-division-kernel-record} for $\mathcal P$, with a constant $C_e$. Suppose that:
\begin{enumerate}
\item the division bound at twisted places holds at $D$ with $C_e$: for every lift $j$,
  $\sum_{w\in\operatorname{supp},\ \mu_w\ne\mu_0}\ell_we_w\eta_j(w)\le C_eD_0x^{\varepsilon}(B_0+N_c)$
  (Definition~\ref{is-twisted-division-bound});
\item every twisted place of the rows of $\mathcal P$ is a place $w$ of the support of the size data
  with $\mu_w\ne\mu_0$;
\item $D_0$ is the dimension of the definite subspace of $\mathcal P$ at the base tuple, for $P_0$
  and $\psi_0$;
\item for all division data tied to $\mathcal P$ whose identifications are those of $P$, the rows
  condition of Definition~\ref{is-definite-side-condition} holds at every word $u$ for the record
  numbers of $u$.
\end{enumerate}
Then for all division data tied to $\mathcal P$ whose identifications are those of $P$, every word
$u$ has division inputs other than the valuations (9) (Definition~\ref{is-division-inputs}) for
$\mathcal P$, the record numbers of $u$, $s(w)=\lfloor e_ws(u,w)\rfloor$ and the module
$N^\flat_u=e_{N,u}^{-1}(N^\circ_u)$; their chain of lifts is tied to the theta kernels of the
realization and has its first columns.
\end{lemma}
\leanref{Development/AutomorphicArithmetic/ThetaCorrespondence/ThetaPairLifts/ThirdPlaceSeries/InpAtOf.lean}{Formalism.L2.RowDataTransport.inpAt_of_alt}
\leanref{Development/AutomorphicArithmetic/ThetaCorrespondence/ThetaPairLifts/ThirdPlaceSeries/PtKOfReal.lean}{Formalism.L2.RowDataTransport.InpAtAlt}
\leanref{Formalism/LocalToGlobal/Divisibility/LinkDef.lean}{Formalism.L2.DIV.lift_exponent_of_least}
\leanref{Formalism/LocalToGlobal/Divisibility/ConstantGeTwentyEightAlt.lean}{Formalism.L2.DIV.DivisionInputsWithoutNorm.ofRows}
\leanref{Formalism/LocalToGlobal/Divisibility/TwistedPlaceEtaSumLeExtended.lean}{Formalism.L2.DIV.DivInputsDef3.toDivInputs0R}
\leanref{Formalism/LocalToGlobal/Divisibility/SumLeEtaSumNum.lean}{Formalism.L2.DIV.inEta_num}
\leanref{Formalism/LocalToGlobal/Divisibility/TwistedPlaceEtaSumLeExtended.lean}{Formalism.L2.DIV.twisted_place_eta_sum_le_extended}
\leanref{Formalism/LocalToGlobal/Divisibility/LinkDef.lean}{Formalism.L2.DIV.bLG_def}
\leanref{Formalism/DivisorInequality/DivisorInequalityProofs/BaseLevelBoundsWithCovolumeHolds.lean}{Formalism.DivisorInequality.places_ell_weighted_sum_le_one}
\begin{proof}
Fix division data tied to $\mathcal P$ whose identifications are those of $P$, and a word $u$. Let
$R$ be the record numbers of $u$, and $S_0$ the rows of $\mathcal P$, with twisted places
$\mathrm{Tw}$ and integer $D$. Put $U=x^{\varepsilon}(B_0+N_c)$.

\emph{Inputs (1) to (4).} By hypothesis (4), the rows condition holds at $u$ for $R$. Its parts (6),
(3), (1) and (2), and (4) and (5), are inputs (1), (2), (3) and (4) for $N^\flat_u$, except the
finiteness in (1): $\lfloor s(w)\rfloor=0$ for $w$ off the support, which is finite. Part (2) of the
rows condition says that the chain is tied to the theta kernels and has its first columns.

\emph{Inputs (5), (7) and (8).} These are the conclusion (2) of
Lemma~\ref{asm-division-kernel-record}, with $L_G=\lambda_0$, $E_2(j)=C_eU+\log d_g$ and $E_3=1$.

\emph{Input (6).} For every tuple $t$ and finite place $y$ take any local basis $e_{t,y}$ of
$\Lambda_t$ at $y$. By Lemma~\ref{is-gram-covolume}, which applies to the well-formed definite side
of the rows, only finitely many places $y$ have $v_y(G_{\mu_0}(y))\ne0$; let $S_y$ be the union of
these places with $\mathrm{Tw}$. Off $S_y$, $\mu_y=\mu_0$, because the twisted places are those with
$\mu_y\ne\mu_0$, so $v_y(G_{\mu_y}(y))=0$. Put $g_y=v_y(|G|)$, $L_G=\lambda_0$, $E_3=1$, $E_2(j)=C_eU+\log d_g$ and
\[
 \eta_y(j)=\eta^{\mathrm{fam}}_j(y)+D\,v_y(d_G(j,\mu_y)),
\]
where $\eta^{\mathrm{fam}}_j(y)$ is $e_y$ times the local defect of the family of $P$ at $y$ for the
lift $j$ and the tuple $\mu_y$, and $d_G(j,t)$ is the
stabilizer denominator of the lift lattice of $P$ at $(j,t)$, which divides $d_g$. We check the three
conditions of Lemma~\ref{is-glued-covolume}.
\begin{enumerate}
\item For $y\in S_y$, $\sum_{y\in S_y}\ell_yv_y(|G|)\le\log|G|=\lambda_0$, by
  Lemma~\ref{is-lattice-facts}(7) applied to the positive integer $|G|$, whose valuations are
  nonnegative.
\item For $y\in\mathrm{Tw}$, $\sum_{y\in\mathrm{Tw}}\ell_y\eta_y(j)/D$ splits into
  $\frac1D\sum_y\ell_y\eta^{\mathrm{fam}}_j(y)$ and $\sum_y\ell_yv_y(d_G(j,\mu_y))\le\log d_g$, the
  latter by Lemma~\ref{is-lattice-facts}(7) again. The division rows identify the local defects of the
  family at the twisted places with those of $D$ (Definition~\ref{is-division-rows}(4)), so
  $\eta^{\mathrm{fam}}_j(y)=e_y\eta_j(y)$, and by
  hypothesis (2) every twisted place lies in the support with $\mu_y\ne\mu_0$; the terms are
  nonnegative, so hypothesis (1) gives
  $\sum_{y\in\mathrm{Tw}}\ell_y\eta^{\mathrm{fam}}_j(y)\le C_eD_0U$. By hypothesis (3), $D_0=D$. So
  $\sum_{y\in\mathrm{Tw}}\ell_y\eta_y(j)/D\le C_eU+\log d_g=E_2(j)$.
\item $\sum_y\ell_y\,2e_y/p_y^2\le1=E_3$ over any finite set of places: grouping the places over a
  prime $p$, $\sum_{y\mid p}\ell_ye_y\le\log p$, and $\sum_p2\log p/p^2\le1$.
\end{enumerate}
The two local bounds of conditions (1) and (2) hold as follows. The rows $S_0$ are the
indefinite-side datum of the family of $P$ at $(P_0,\psi_0)$ on the definite subspaces of $\mathcal P$,
and $d_Q(t)=d_Z(j,t)=D$ for all $j$ and $t$, as in the proof of Lemma~\ref{is-glued-covolume}. By
Lemma~\ref{asm-division-kernel-inputs-gram} at level $1$,
$v_y(G_{\mu_y}(y))\le D\,v_y(|G|)=D\,g_y$ for every $y$. The function $J_0$ of the scale data of the
member is the lift exponent of the family of $P$
(Definition~\ref{is-onto-step-exp-bound-lift-exponent}), so $e_yJ_0(\mu_y,y,j)=J_y(j,\mu_y)$ in the
notation of Lemma~\ref{asm-division-kernel-inputs-index}, and $\eta^{\mathrm{fam}}_j(y)$ is the number
$\eta^{\mathrm{fam}}_j(\mu_y,y)$ of that lemma. That lemma at level $1$ gives, for $y\in\mathrm{Tw}$
and every lift $j$,
\[
 j_y\big(H_T(j,\mu_y),T_{j,\mu_y}\big)\le\eta^{\mathrm{fam}}_j(y)+D\big(e_yJ_0(\mu_y,y,j)
 +v_y(d_G(j,\mu_y))\big)=e_yJ_0(\mu_y,y,j)D+\eta_y(j).
\]
% Lean: gram_valuation_le_extended (from FamIn.GmOf_le) and lift_index_bound_extended (from
% FamIn.jw_le, lift_lattices and exponent_weight_spectral, with the lift exponents at the twisted
% places LiftExponentAtTwisted, which lift_exponent_of_least gives from the family conditions with J0 the
% least exponent JJ of the scale data; TwistedPlaceEtaSumLeExtended.lean and LinkDef.lean).
% The sum over Tw is sum_le_etaSum_num (SumLeEtaSumNum.lean) with NumTie and TieTwRec, then DIVinN_at.
% The bound sum_y ell_y 2 e_y / p_y^2 <= 1 is places_ell_weighted_sum_le_one
% (ramified_places_weighted_sum_le_one).
With these, Lemma~\ref{is-glued-covolume}'s hypotheses are the data of input (6), and the division
inputs other than (9) hold.
% Lean: inpAt_of_alt. Inputs (1) to (4) come from RowsU (the rows condition, hypothesis what2), (5), (7)
% and (8) from ArithAt, and (6) from inEta3_num (the eta sums from DIVinN_at, hypothesis (1), with
% TieTwRec, hypothesis (2), and the rank equality, hypothesis (3)) through
% DivisionInputsWithoutNorm.ofRows, toDef and DivInputsDef3.toDivInputs0R; the lift exponents at the
% twisted places, LiftExponentAtTwisted, come from the family conditions by lift_exponent_of_least
% (in wire4F_alt). Size: proof closure of inpAt_of_alt, toDivInputs0R, toDef and
% lift_exponent_of_least, 211 theorems, about 3,500 lines in 31 files.
\end{proof}

\subsection{Height bounds}\label{subsection-asm-height-bounds}

\begin{lemma}\label{asm-division-kernel-assembly}
Let $\mathcal F$ be a family with exponent $\varepsilon$ and size constant $C_{sz}$, and let $D$ be a
member at a scale $x>x_1$, with coefficient field $L$, a twist record, a spectral condition for it,
and division kernel data (Definition~\ref{is-division-kernel-bounds}) built from square-integrable
columns for a row-side family $P$ whose row side is determined by the twist record, with the scale
data of the member. Let $Q$ be an indefinite realization of $P$ for the spectral condition. Let $C'$
and $C_H$ be real numbers for which the hypotheses of Lemma~\ref{asm-division-kernel-package} hold,
and let $\mathcal P$ be the package it gives. Suppose that:
\begin{enumerate}
\item the division rows hold for $P$ (Definition~\ref{is-division-rows}) at every tuple, the
  division data are compatible with the periods at every tuple, as in item (3)(f) of
  Definition~\ref{asm-definite-side-hypotheses}, and the archimedean scalars of the columns of $P$
  are those of the division data of $D$;
\item $\lambda_0$, $S_w$, $C_1$ and the record numbers of the words satisfy conclusion (1) of
  Lemma~\ref{asm-division-kernel-record} with constants $(\lambda_0,c_0,c_e,C_{sz},S_w,C_1)$, and
  the conclusion of Lemma~\ref{asm-division-kernel-inputs} holds for $\mathcal P$ and these record
  numbers;
\item $N_u$ is the $\mathcal O_L$-span of $I_uN^\circ_u$ for every word $u$, with $I_u$ as in the
  proof of Lemma~\ref{asm-slope-link}, and the heights $h_u$ of $D$ are the explicit heights with
  constant $C_{sz}$;
\item $C\ge\max(0,50c_0+10,60c_e)$, $C_{\mathrm{abs}}\ge1$ and $L_a\ge0$.
\end{enumerate}
Then the division kernel data of $D$ carry a division package with constants $(C,C_{\mathrm{abs}})$
as in part (1) of Definition~\ref{is-division-kernel-bounds}: its tuples, evaluations and words are
those of the data, and it comes with the record numbers, the modules $N^\flat_u$, the division inputs
other than the valuations and the three ties of the proof of Theorem~\ref{is-division-link-holds}.
\end{lemma}
\leanref{Development/AutomorphicArithmetic/ThetaCorrespondence/ThetaPairLifts/ThirdPlaceSeries/PtKOfReal.lean}{Formalism.L2.RowDataTransport.ptK_of_real_alt}
\leanref{Development/AutomorphicArithmetic/ThetaCorrespondence/ThetaPairLifts/ThirdPlaceSeries/DIVRecBridgeSel.lean}{Formalism.L2.DIV.pkg_of_rowsK_alt}
\leanref{Development/AutomorphicArithmetic/ThetaCorrespondence/ThetaPairLifts/ThirdPlaceSeries/DIVRecBridgeSel.lean}{Formalism.L2.DIV.DivisionPackageKernel'}
\leanref{Formalism/LocalToGlobal/Divisibility/RecBridgeMkF.lean}{Formalism.L2.DIV.DivisionPackageFace.ofOpen}
\leanref{Formalism/LocalToGlobal/Divisibility/EigenvalueSumLeClosedOf.lean}{Formalism.L2.DIV.fixedScale_tied}
\leanref{Formalism/LocalToGlobal/Divisibility/RecBridgeSel.lean}{Formalism.L2.DIV.ColTie3.transport}
\leanref{Formalism/LocalToGlobal/GroupLattices/BrTied3Hold.lean}{Formalism.L2.G.ServJointF.brTied_of}
\leanref{Formalism/LocalToGlobal/Divisibility/DivisibilityConstruction/ShiftOfColTie.lean}{Formalism.L2.DIV.Cons.shift_nuOf}
\leanref{Formalism/LocalToGlobal/Divisibility/RecBridge.lean}{Formalism.L2.DIV.hTie_of}
\leanref{Formalism/LocalToGlobal/Divisibility/RecBridgeNum.lean}{Formalism.L2.DIV.hnum_of_numRecTieH}
\leanref{Formalism/LocalToGlobal/Divisibility/SmallTies.lean}{Formalism.L2.DIV.two_sumTw_eq_Fo}
\leanref{Formalism/LocalToGlobal/Divisibility/RecBridge.lean}{Formalism.L2.DIV.rhs_le_record}
\begin{proof}
Let $\mathcal P$ have rows $S_0$, columns $S_1$ and numerical bounds $N=N(C',C_H)$ of the scale data
of the member, with function $F_0$. By Lemma~\ref{asm-division-kernel-package}, its definite
subspaces, rows, columns and column data are those of $P$.

\emph{Step 1: division data on $\mathcal P$.} The division kernel data identify their column record
with the columns of $P$. The division rows of $P$, with the ties of $P$ to $D$
(Definition~\ref{is-row-side-family}), give division data tied to $P$: the tuple map $t\mapsto\mu_t$,
the identifications $e_{RR}$, $e_{RR,0}$, $e_{N,u}$ and $e_M(i,t)$, the maps $j_W(i)$ and the word
bijection $u\mapsto w_u$ of $P$, with conditions (1) to (5) of Definition~\ref{is-division-rows}.
Since the rows, columns and column data of $\mathcal P$ are those of $P$, these are division data
tied to $\mathcal P$ whose identifications are those of $P$. In particular the tuples, the
evaluations $e_M$ and the words of these data are those of the division kernel data.

\emph{Step 2: the columns are tied to the periods.} By the compatibility with the periods in
hypothesis (1), for every word $u$, tuple $t$, $x\in L^{a(w_u)}$ and $F$, the period
$\mathrm{Pt}_{w_u}(t)(x,F)$ of the column data of $\mathcal P$ is, under the embedding of $L$, the
period of $D$ of $e_{N,u}(x)$ at the tuple $\mu_t$ against the image of $F$ under $e_M$. With the
tuple map of Step 1 and the fact that the level operators fix the images of the forms of the columns,
this is the tie of the columns to the periods of $D$ used in the proof of
Theorem~\ref{is-division-link-holds}.

\emph{Step 3: record numbers, modules and inputs.} By hypothesis (2), the record numbers of every
word $u$ hold with constants $(\lambda_0,c_0,c_e,C_{sz},S_w,C_1)$, and, for the division data of
Step 1, the conclusion of Lemma~\ref{asm-division-kernel-inputs} gives the division inputs of $u$
other than the valuations, for $\mathcal P$, these record numbers,
$s(w)=\lfloor e_ws(u,w)\rfloor$ and $N^\flat_u=e_{N,u}^{-1}(N^\circ_u)$, with the chain of lifts tied
to the theta kernels and its first columns. The division rows identify $N_u$ with a submodule
$N^C_u$ of the column of $w_u$ under $e_{N,u}$. By hypothesis (3), $N_u$ is the $\mathcal O_L$-span
of $I_uN^\circ_u$, where $I_u=\prod_w\mathfrak p_w^{-\lfloor e_ws(u,w)\rfloor}$; as $e_{N,u}$ is an
$L$-linear isomorphism, $N^C_u$ is the $\mathcal O_L$-span of $I_uN^\flat_u$. This is the relation
between $N_u$ and $N^\flat_u$ in the proof of Theorem~\ref{is-division-link-holds}.

\emph{Step 4: the height tie.} The record numbers of $u$ are built from the member, with $c=C_{sz}$
and the places $z\in\{0,1\}\times A$ identified with $\mathcal T$ by $z\mapsto\tau_z$
(Definition~\ref{pd-orbit-places}(1)). Comparing the definitions term by term, using
$S_c=4k\Delta_S$ and the width $\eta$ of the size data, gives, in the notation of
Definitions~\ref{is-record-numbers} and~\ref{asm-division-kernel-explicit-heights}, $h_A=S(u)$,
$h_B(j)=K_j(u)$ and $h_X(j)=\ell_j$; so $h^{\mathrm{rec}}=S(u)+\sum_jK_j(u)-\sum_j\ell_j$. The
scalars $\kappa_v$ of the columns of $\mathcal P$ are those of $P$, which by hypothesis (1) are the
scalars $\kappa_R(\sigma)$ at the embeddings $\sigma$ of the place $v$. Since $L$ is totally complex,
the mean of $\log\kappa_v$ over the infinite places is the mean $\overline{\log\kappa_R}$ over the
embeddings. By hypothesis (3), $h^{\mathrm{rec}}-\overline{\log\kappa}=h_u$.

\emph{Step 5: the tie of the local terms.} By the construction of the scale data of the member,
$F_0(w,j)$ is the explicit local term $F^{\mathrm{ex}}_j(p_w)$ of
Definition~\ref{asm-division-kernel-sizes-local-terms} at the lift $j$; the scale data compute the
Fock degree from the entries $a_i$ of Definition~\ref{asm-division-kernel-explicit-heights}(1) and
their counterparts for the negative part, and the error data of the member identify these with the
integer weights $\kappa^{\pm}_{j,\tau,a}$ of its rows data. The local terms of the member are
explicit, so $F_0(w,j)=F_j(w)$.
The rows of $\mathcal P$ are untwisted (Definition~\ref{is-division-package}(1)), so $F_0(w,j)=0$
off the twisted places, and $F_j(w)=0$ off the support. Extending both sums to the union of the
twisted places and the support,
\[
 2F^{\mathrm{pk}}_o=2\sum_j\sum_{w\in\mathrm{Tw}}\ell_we_wF_0(w,j)
 =\sum_j\sum_{w\in\operatorname{supp}}2\ell_we_wF_j(w)=F_o.
\]

\emph{Step 6: the right side.} The record numbers have
$\sum_{t\in T}\big(m(\phi(a_t)-a_t\log\Omega)-a_t\log T_u\big)=\mathrm{MAIN}(u)$, and their $|T|$,
$a$, $L_a$, $B_0$, $N_c$, $m$, $x$ and $\varepsilon$ are those of the member. So
\[
 \mathrm{rhs}=\mathrm{MAIN}(u)+(50c_0+10)|T|aL_a+60c_emx^{\varepsilon}(B_0+N_c).
\]
By hypothesis (4), $C\ge50c_0+10$, $C\ge60c_e$ and $C\ge0$; also $|T|aL_a\ge0$ and $B_0+N_c\ge0$,
and $x^{\varepsilon}\le x^{C_{\mathrm{abs}}\varepsilon}$ because $x>1$ and $C_{\mathrm{abs}}\ge1$. So
$\mathrm{rhs}\le\mathrm{MAIN}(u)+C|T|aL_a+Cm(B_0+N_c)x^{C_{\mathrm{abs}}\varepsilon}$.

Steps 1 to 6 are the data of part (1) of Definition~\ref{is-division-kernel-bounds}.
% Lean: ptK_of_real_alt = pkg_of_rowsK_alt (ColTie3.transport, divTied3_congr, divTieDat3_heq,
% Formalism.L2.G.ServJointF.brTied_of with dcPt_move, which uses the compatibility with periods)
% followed by DivisionPackageColumns.ofTied = DivisionPackageFace.ofOpen, whose ties are hTie
% (hTie_of with hnum_of_numRecTieH and the explicit heights hh), foTie (two_FoOfR_eq_Fo from FTie, which
% fixedScale_tied gives for the scale data of the member) and rhs (rhs_le_record with
% Cons.numRecTie_record and the constants). Size: proof closure of ptK_of_real_alt, ofTied, ofOpen and
% fixedScale_tied, 155 theorems, about 2,000 lines in 30 files.
\end{proof}

\begin{lemma}\label{asm-division-kernel-sharp}
Let $P$, $Q$, $R$, $F$, $C$ and $C_H$ satisfy the hypotheses of
Lemma~\ref{asm-division-kernel-package}, let the division package it gives be built on the scale
data of $F$, with the numbers $x$, $b$, $\ell_a$ and $\delta_e$ of these scale data, and let $RR$ be
the glued lattice of $D$. Suppose that $e_{RR}$ identifies the glued lattice of the columns $S_1$ of
$P$ with $RR$, as lattices and as metrics (Definition~\ref{is-division-rows}(1)), that the
archimedean scalars $\kappa_v$ of $S_1$ are the scalars $\kappa_R$ of the division data of $D$, and
that $0\le\delta_e\le\frac13$. Then
\[
 \mu_{\max}(RR)+\overline{\log\kappa_R}\le x^{C_H}b+14\ell_a,
\]
where $\overline{\log\kappa_R}$ is the mean of $\log\kappa_R(\sigma)$ over the embeddings $\sigma$
of $L$ into $\mathbf C$.
\end{lemma}
\leanref{Development/AutomorphicArithmetic/ThetaCorrespondence/ThetaPairLifts/ThirdPlaceSeries/HefunPointB.lean}{Formalism.L2.DIV.FunSharpTop.funSharp_pt_alt}
\begin{proof}
Let $RR_0$ and $RR_1$ be the glued lattices of the rows and of the columns of the package, and put
$\overline{\log\kappa}=\frac1{r_\infty}\sum_v\log\kappa_v$, the mean over the infinite places of $L$.
The rows satisfy the functional bound for $N(R,F;C,C_H)$, whose number $h_N$ is $x^{C_H}b$
(Definition~\ref{is-numerical-bounds}), so $\mu_{\max}(RR_0)\le x^{C_H}b-\overline{\log\kappa}$. The
passage (Definition~\ref{is-division-package}(4)), with $H=x^{C_H}b-\overline{\log\kappa}$ and
$0\le\delta_e\le\frac13$, gives
\[
 \mu_{\max}(RR_1)\le x^{C_H}b-\overline{\log\kappa}+14\ell_a.
\]
The columns of the package are $S_1$, so $e_{RR}$ carries $RR_1$ with its metric isometrically onto
$RR$, and $\mu_{\max}(RR)=\mu_{\max}(RR_1)$. The rows and the columns of the package have the same
scalars $\kappa_v$, and by hypothesis these are the $\kappa_R(\sigma)$ at the embeddings $\sigma$ of
the place $v$. Since $L$ is totally complex, each infinite place has two embeddings, so
$\overline{\log\kappa_R}=\overline{\log\kappa}$. Adding $\overline{\log\kappa_R}$ to both sides
gives the bound.
\end{proof}

\begin{lemma}\label{asm-division-kernel-mid}
Let $\mathcal C$ be a class of families, and let $P(\mathcal F,C_H)$ be a property of families and
real numbers that holds eventually along $\mathcal C$ (Definition~\ref{pd-family}). Then there is a
real function $f$ of a real number and a finite list of real numbers such that the property
$P\big(\mathcal F,f(\varepsilon,c_{\mathrm{out}}(\mathcal F))\big)$ of $\mathcal F$, where
$\varepsilon$ is the exponent of $\mathcal F$, holds eventually along $\mathcal C$.
\end{lemma}
\leanref{Formalism/LocalToGlobal/Divisibility/DivisibilityRemainder/SharpFunctions/MidExp.lean}{Formalism.L2.DIV.FunSharpTop.closed_atMid}
\leanref{Formalism/LocalToGlobal/Divisibility/DivisibilityRemainder/SharpFunctions/MidExp.lean}{Formalism.L2.DIV.FunSharpTop.large_of_large_thresholds}
\begin{proof}
Let $\varepsilon_1$ be as in the definition of \emph{eventually along}, and for
$\varepsilon\in(0,\varepsilon_1]$ let $\kappa(\varepsilon)=(\kappa_{\mathrm{pre}},\kappa_{\mathrm{mid}},
\kappa_{\mathrm{last}})$ be a largeness threshold for $\varepsilon$. Put
$f(\varepsilon,l)=\kappa_{\mathrm{mid}}(l)$ for $\varepsilon\in(0,\varepsilon_1]$, and $f=0$
otherwise. For a list $l$ that begins with a list $l_{\mathrm{out}}$ of the length of
$c_{\mathrm{out}}$, let $l'$ be $l$ with its next entry replaced by
$\kappa_{\mathrm{mid}}(l_{\mathrm{out}})$, and let $\kappa'(\varepsilon)$ be the threshold
$(\kappa_{\mathrm{pre}},\kappa_{\mathrm{mid}},l\mapsto\kappa_{\mathrm{last}}(l'))$.

Let $\mathcal F\in\mathcal C$ have exponent $\varepsilon\in(0,\varepsilon_1]$ and be large for
$(\kappa'(\varepsilon),C_H)$, and put $C'_H=\kappa_{\mathrm{mid}}(c_{\mathrm{out}})$. The list
$c_{\mathrm{mid}}(\mathcal F,C'_H)$ begins with $C'_H$, which is at least
$\kappa_{\mathrm{mid}}(c_{\mathrm{out}})$, and its other entries are those of
$c_{\mathrm{mid}}(\mathcal F,C_H)$. The list $c_{\mathrm{out}}$ followed by
$c_{\mathrm{mid}}(\mathcal F,C_H)$, with its entry $C_H$ replaced by $C'_H$, is $c_{\mathrm{out}}$
followed by $c_{\mathrm{mid}}(\mathcal F,C'_H)$, so every constant of $c_{\mathrm{last}}$ is at least
$\kappa_{\mathrm{last}}$ of the latter list. Hence $\mathcal F$ is large for $(\kappa(\varepsilon),C'_H)$,
and $P(\mathcal F,C'_H)$ holds. So the thresholds $\kappa'(\varepsilon)$ show that
$P(\mathcal F,f(\varepsilon,c_{\mathrm{out}}))$ holds eventually along $\mathcal C$.
\end{proof}

\begin{lemma}\label{asm-division-kernel-functional}
Let a collection of definite-side hypotheses be given (Definition~\ref{asm-definite-side-hypotheses}),
with class $\mathcal A$ and $\mathcal D_3=\operatorname{Des}(\mathcal A_N,3)$, and let $e_b$ and
$e_c$ be real functions of finite lists of real numbers. For a list $l$ with first entry $l_1$ put
$e_d(l)=9(\max(l_1,5)+3)+1$. Eventually along $\mathcal D_3$, every family $\mathcal F$ that satisfies
both size conditions (Definition~\ref{asm-designed}), with height constant $C_H$, has a real function
$e_F$ of finite lists of real numbers such that:
\begin{enumerate}
\item $\mathcal F$ is large for $C_H$ and the threshold $(0,\kappa_{HH},0)$ with
  $\kappa_{HH}=\max(e_b,e_c,e_F,e_d)+1$ (Definition~\ref{pd-family}); in particular
  $C_H\ge\kappa_{HH}(c_{\mathrm{out}}(\mathcal F))$;
\item if $\mathcal F$ satisfies the division bound at twisted places, then at all large scales, at
  every member,
  \[
   \mu_{\max}(RR)+\overline{\log\kappa_R}\le x^{e_F(c_{\mathrm{out}}(\mathcal F))}b,
  \]
  with $\overline{\log\kappa_R}$ as in Lemma~\ref{asm-division-kernel-sharp}.
\end{enumerate}
\end{lemma}
\leanref{Development/AutomorphicArithmetic/ThetaCorrespondence/ThetaPairLifts/ThirdPlaceSeries/HefunFunTop.lean}{Formalism.L2.DIV.FunSharpTop.closed_funSharp_thr_alt}
\leanref{Development/AutomorphicArithmetic/ThetaCorrespondence/ThetaPairLifts/ThirdPlaceSeries/HefunFunTop.lean}{Formalism.L2.DIV.FunSharpTop.closed_funSharp_alt}
\leanref{Development/AutomorphicArithmetic/ThetaCorrespondence/ThetaPairLifts/ThirdPlaceSeries/HefunFunTop.lean}{Formalism.L2.DIV.FunSharpTop.funSharp_fam_alt}
\leanref{Formalism/SharpRestriction/Divisibility/HookProducts/LambdaZeroLeFourB.lean}{Formalism.MainNotions.DIV.HookProducts.FunSharp.lambda_zero_le_four_b}
\leanref{Formalism/SharpRestriction/Divisibility/SlopeExcess/FamilyStatementOfScaledLemmas.lean}{Formalism.MainNotions.DIV.Proofs.ScaledLatticeNorms.scaled_lattice_nonzero_holds}
\leanref{Formalism/LocalToGlobal/Divisibility/DivisibilityRemainder/SharpFunctions/FunTop.lean}{Formalism.L2.DIV.FunSharpTop.pow_absorb}
\leanref{Formalism/LocalToGlobal/Divisibility/DivisibilityRemainder/SharpFunctions/MidExp.lean}{Formalism.L2.DIV.FunSharpTop.closed_largeE}
\begin{proof}
\emph{Step 1: a height constant for the rows.} The definite rows of the collection hold eventually
along $\mathcal D_3$, as a property of $(\mathcal F,C_H)$: there is $C\ge0$ such that at every member
at all large scales the realization $R_Q$ of the selection bundle satisfies conclusions (1) to (6) of
Theorem~\ref{is-definite-rows} with $(C,C_H)$. Eventually along every class $C_{sz}\ge0$, since
$C_{sz}$ is a constant of the middle list. Lemma~\ref{asm-division-kernel-mid} gives a function $f$
such that, eventually along $\mathcal D_3$, both hold with $C_H$ replaced by
$C'_H=f(\varepsilon,c_{\mathrm{out}})$. Put $e_F(l)=\max(f(\varepsilon,l),1)+1$.

\emph{Step 2: the bound.} Let $\mathcal F\in\mathcal D_3$ satisfy both size conditions, the definite
rows with $C'_H$, $C_{sz}\ge0$ and the division bound at twisted places. Let $C\ge0$ be the constant
of the definite rows and $C_d$ that of the size condition on stabilizer denominators. Let $x$ be
beyond the scale of the definite rows, the threshold of the twist selection, the scale of the size
condition on stabilizer denominators, the $x_0$ of Lemma~\ref{asm-division-kernel-record} for $C$,
$C'_H$, $C_e=0$ and $\max(C_d,0)$, and $56$. Fix a member $D$ at $x$, with selected base level $P$,
selection bundle $Q$ and realization $R_Q$ at the face $F_Q$ of $P$.

The bounded rows (3)(a) give every lift of $Q$ a single row map (Definition~\ref{is-bounded-rows}(2)),
so conclusions (1) to (5) of Theorem~\ref{is-definite-rows} hold for $R_Q$ with $(C,C'_H)$. With the
family conditions of the square-integrable columns and with item (3)(c), the primitive rows and the
lift-distinct types, Lemma~\ref{asm-division-kernel-package} gives a division package on the scale
data of the member. Lemma~\ref{asm-division-kernel-record}, at one word, gives input (8) of
Definition~\ref{is-division-inputs}: a number $Q_0\ge1$ with $0\le\delta_e$ and $3\delta_eQ_0\le1$, so
$\delta_e\le\frac1{3Q_0}\le\frac13$. The division rows (3)(e) give
Definition~\ref{is-division-rows}(1), and conclusion (5) of Theorem~\ref{is-definite-rows} gives the
equality of the archimedean scalars. Lemma~\ref{asm-division-kernel-sharp} gives
\[
 \mu_{\max}(RR)+\overline{\log\kappa_R}\le x^{C'_H}b+14\ell_a,
\]
where $\ell_a=\log|G|=\lambda_0$ for the group $G$ of $P$, of order
$\prod_{k\in\{0,1\}}|GL_r(\mathbf F_{g_k})|^f$.

Since $|GL_r(\mathbf F_g)|\le g^{r^2}$, $\lambda_0\le fr^2(\log g_0+\log g_1)$. The sizes of the
family give $f,r\ge1$, $g_1\le2g_0$, $\log g_0\le A$ and $x^{\varepsilon/100}r^{19}f^5x^{\varepsilon}\le
b$, where $A=x^{\varepsilon/100}x^{\varepsilon}\ge1$. So $fr^2\le fr^2A\le b$, and
\[
 \lambda_0\le2fr^2\log(2g_0)\le2fr^2(1+A)\le4b.
\]
As $x\ge56$, $b\ge1$ and $x^{C'_H}\le x^{\max(C'_H,1)}$, and as $x^{\max(C'_H,1)}\ge x\ge56$ gives
$x^{\max(C'_H,1)}+56\le x^{\max(C'_H,1)}x$,
\[
 x^{C'_H}b+14\lambda_0\le\big(x^{\max(C'_H,1)}+56\big)b\le x^{\max(C'_H,1)+1}b
 =x^{e_F(c_{\mathrm{out}})}b.
\]

\emph{Step 3: largeness.} Thresholds for $\mathcal D_3$ may be enlarged. Taking for each $\varepsilon$
the larger of the thresholds of Steps 1 and 2 and of $(0,\kappa_{HH},0)$, with the function $e_F$ of
Step 1, every family that is large for the result has (1) and (2).
\end{proof}

\begin{definition}\label{asm-division-kernel-explicit-heights}
Let $\mathcal F$ be a family with exponent $\varepsilon$ (Definition~\ref{pd-family}), and let $D$ be
a member at a scale $x>x_1$, with ranks $k,m,t,n$, $f=2n$, $N(j)=k+j(t+1)$ and $n_V=k+m-1+t$
(Definition~\ref{th-ranks}), coefficient field $L$, top classes $T$, auxiliary primes $T_u$, the set
$\mathcal T$ of $f$ embeddings with Harish-Chandra parameters $\lambda(\tau)=(u,v,b,e)$ and the
parameters $\lambda_u(\tau)$ of the words (Definition~\ref{th-definite-datum}(2)). Let its size data
have exponents $a_t$, width $\eta$, cloud size $S_c$ and quantities $b,b'$
(Definition~\ref{pd-size-data}). Let $L_a$, $\log\Omega$ and $\phi(y)=y\log y-y$ be as in
Definition~\ref{is-division-link}, $L_p=\log(b\,r_E\,g_0g_1\,n_V)$ the logarithmic size of
Lemma~\ref{is-hook-lower-bound}, $n_{u,j,r}(\tau)$ the row exponents of
Definition~\ref{pd-orbit-places}(2) and $n_{u,j}(\tau)=\sum_{r<N(j-1)}n_{u,j,r}(\tau)$. Let $C$ be
real.
\begin{enumerate}
\item For $\tau\in\mathcal T$ put $\sigma_i(\tau)=v_{i+1}(\tau)$ for $0\le i<k-1$ and
  $\sigma_i(\tau)=-b_{i-k+2}(\tau)$ for $k-1\le i<k-1+m$. For a lift $1\le j\le m$ put $n=k+j-1$,
  \[
   a_i=\sigma_i(\tau)+e_j(\tau)-\tfrac12(n_V-1-2i)-\tfrac12\big((k+j-1)-(m-j)\big)\qquad(0\le i<n),
  \]
  and
  \[
   \log\mathrm{hook}_j(\tau)=\sum_{i<n}\log\Gamma(a_i+n-i)-\sum_{0\le i<l<n}\log(a_i-a_l+l-i).
  \]
\item For a word $u$ and a lift $1\le j\le m$ put
  \[
   S(u)=-\sum_{t\in T}a_t\log T_u+C\big(|T|\,b'N(m)+f(m\eta+S_c)N(m)+fN(m)^2\big)L_a,
  \]
  \[
   K_j(u)=\sum_{\tau\in\mathcal T}\Big(\sum_{r<N(j-1)}\phi\big(n_{u,j,r}(\tau)\big)
   -n_{u,j}(\tau)\log\Omega\Big)+CfN(m)^2L_a+Cfm^3N(m)^2x^{\varepsilon},
  \]
  \[
   \ell_j=\sum_{\tau\in\mathcal T}\log\mathrm{hook}_j(\tau)-Cfn_V^2L_p.
  \]
\item The \emph{explicit height} of a word $u$ with constant $C$ is
  \[
   h_u=S(u)+\sum_{j=1}^mK_j(u)-\sum_{j=1}^m\ell_j-\overline{\log\kappa_R},
  \]
  where $\overline{\log\kappa_R}$ is the mean of $\log\kappa_R(\sigma)$ over the embeddings $\sigma$
  of $L$ into $\mathbf C$, for the archimedean scalars $\kappa_R$ of the division data of $D$. The
  \emph{explicit heights} of the family are those with the size constant $C_{sz}$ of $\mathcal F$.
\end{enumerate}
\end{definition}
\leanref{Formalism/LocalRingCardinality/KernelPlaceMeanBound.lean}{Formalism.BasicNotions.DV.hExpl6}
\leanref{Formalism/LocalRingCardinality/DataDIVDefs.lean}{Formalism.BasicNotions.DV.SZbP}
\leanref{Formalism/LocalRingCardinality/KernelPlaceMeanBound.lean}{Formalism.BasicNotions.DV.kernelPlaceMeanBound}
\leanref{Formalism/LocalRingCardinality/DataDIVDefs.lean}{Formalism.BasicNotions.DV.kerPlace}
\leanref{Formalism/LocalRingCardinality/KernelPlaceMeanBound.lean}{Formalism.BasicNotions.DV.L6lo}
\leanref{Formalism/LocalRingCardinality/KernelPlaceMeanBound.lean}{Formalism.BasicNotions.DV.logHook}
\leanref{Formalism/LocalRingCardinality/DataKap.lean}{Formalism.BasicNotions.DV.aPlus}
\leanref{Formalism/LocalRingCardinality/DataKap.lean}{Formalism.BasicNotions.DV.sigPar}

\begin{lemma}\label{asm-division-kernel-explicit-entries}
Let $D$ be a member at a scale $x>x_1$ of a family $\mathcal F$, with the notation of
Definition~\ref{asm-division-kernel-explicit-heights}, and put $\Pi=1+f+k+m+t+N(m)+x^{\varepsilon}$.
Then $\Pi\ge1$ and $b\ge1$, and for every $\tau\in\mathcal T$, word $u$ and lift $1\le j\le m$:
\begin{enumerate}
\item $|T|\le\Pi$, $n_V\le\Pi$, $N(j)\le\Pi$, $0\le a_t\le2\Pi b$ for $t\in T$, $0\le b'\le\Pi b$,
  $0<S_c\le\Pi b$ and $0<\eta\le b$;
\item $0\le e_j(\tau)\le4\Pi^5b$, $|b_j(\tau)|\le5\Pi^5b$, $|\sigma_i(\tau)|\le5\Pi^5b$ for
  $0\le i<k-1+m$, and the entries $u_i$ of $\lambda_u(\tau)$ satisfy $|u_i|\le3\Pi^5b$;
\item $|n_{u,j,r}(\tau)|\le11\Pi^5b$ for $0\le r<N(j-1)$;
\item $0\le L_a\le4\Pi$, $0\le L_p\le6\Pi^2$ and $|\log\Omega|\le7\Pi$.
\end{enumerate}
\end{lemma}
\leanref{Formalism/SharpRestriction/Divisibility/SlopeExcess/ConstRankBaseLePolyFourth.lean}{Formalism.MainNotions.DIV.Proofs.HHParts.aT_bounds}
\leanref{Formalism/SharpRestriction/Divisibility/SlopeExcess/ConstRankBaseLePolyFourth.lean}{Formalism.MainNotions.DIV.Proofs.HHParts.bT_le}
\leanref{Formalism/SharpRestriction/Divisibility/SlopeExcess/ConstRankBaseLePolyFourth.lean}{Formalism.MainNotions.DIV.Proofs.HHParts.e_bounds}
\leanref{Formalism/SharpRestriction/Divisibility/SlopeExcess/ConstRankBaseLePolyFourth.lean}{Formalism.MainNotions.DIV.Proofs.HHParts.b_abs}
\leanref{Formalism/SharpRestriction/Divisibility/SlopeExcess/ConstRankBaseLePolyFourth.lean}{Formalism.MainNotions.DIV.Proofs.HHParts.u_abs}
\leanref{Formalism/SharpRestriction/Divisibility/SlopeExcess/ConstRankBaseLePolyFourth.lean}{Formalism.MainNotions.DIV.Proofs.HHParts.sigPar_abs}
\leanref{Formalism/SharpRestriction/Divisibility/SlopeExcess/ConstRankBaseLePolyFourth.lean}{Formalism.MainNotions.DIV.Proofs.HHParts.nuRow_abs}
\leanref{Formalism/SharpRestriction/Divisibility/SlopeExcess/ConstRankBaseLePolyFourth.lean}{Formalism.MainNotions.DIV.Proofs.HHParts.La_bounds}
\leanref{Formalism/SharpRestriction/Divisibility/SlopeExcess/ConstRankBaseLePolyFourth.lean}{Formalism.MainNotions.DIV.Proofs.HHParts.Lprime_bounds}
\leanref{Formalism/SharpRestriction/Divisibility/SlopeExcess/ConstRankBaseLePolyFourth.lean}{Formalism.MainNotions.DIV.Proofs.HHParts.logOm_abs}
\begin{proof}
Every summand of $\Pi$ is nonnegative and the first is $1$, so each of $f,k,m,t,N(m),x^{\varepsilon}$ is
at most $\Pi$, and $b=\exp(x^{\varepsilon})\ge1$. The sizes of the family give the following.

(1) $|T|\le2n\delta_1\le f$ and $n_V\le1+k+m+t$, so $|T|,n_V\le\Pi$; and $N(j)\le N(m)$. The sizes
give $x\le k$, so $a=xb\le\Pi b$ and $b'=a/X'\le a$; the window $a\le a_t\le a+a/(kx^{13})$ gives
$0\le a_t\le2a$. The sizes give $kS_c\le w''$, $w''=b'/(xn)\le b'$ and
$\eta\le b/(t^3fk^3)\le b$, with $S_c,\eta>0$.

(2) The size data place the weights in windows: $e_j(\tau)\in[b_t+b'/x,\,b_t+b'/x+\eta]$ if $\tau$
is the top embedding of $t\in T$, $e_j(\tau)\in[b,\,b+\eta+1]$ if $\tau$ is thick, and
$e_j(\tau)\in[w'',\,w''+\eta]$ otherwise. The sizes give $b'\le b_t\le
b'+2n\delta_1x^{\varepsilon}\Lambda_Gb'/k$ and $\Lambda_G\le x^{2\varepsilon}$, so
$0\le b_t\le2\Pi^5b$. In each case $0\le e_j(\tau)\le4\Pi^5b$. The size data give
$|b_j(\tau)-e_j(\tau)|\le\eta$, and $e_j(\tau)-b_j(\tau)\ge(t+3)/2$ by
Definition~\ref{th-hc-parameter}(4); so $-b\le b_j(\tau)\le e_j(\tau)$ and $|b_j(\tau)|\le5\Pi^5b$.
They give $0\le v_i(\tau)\le S_c$, which with the bound for $b_j$ gives the bound for $\sigma_i$.
The entries $u_i$ of $\lambda_u(\tau)$ with $1<i<k$ lie in $[0,S_c]$ by the size data, and $u_1,u_k$
are those of $\lambda(\tau)$. If $\tau$ is the top embedding of $t$, the size data give
$u_1=a_t+(k-1)/2$ and $u_k=-b_t-(k-1)/2$, so $|u_1|,|u_k|\le3\Pi^5b$. Otherwise the parameter choices
give $u_1-v_1\le2S_c$ and $u_k=v_{k-1}-\frac12$; as $u_1>v_1\ge0$, $0\le u_1\le3S_c$ and
$|u_k|\le S_c+\frac12$.
% Lean: u_abs takes u_1 - v_1 <= 2 S_c from the parameter choices (ParameterChoice'.cloud_place_bound,
% the field P5 of hExpl6_lower_at, which hExpl6_lower takes from Fm.params). The size data's own bound
% S_c <= u_1 - v_1 <= C_c S_c (SizeData₂.cloud_gap) is not used here; it is the one used in the proof of
% Lemma asm-division-kernel-record (nf_off_rec). Only |u_1| <= 3 Pi^5 b is used downstream (nuRow_abs).

(3) By Definition~\ref{pd-orbit-places}(2),
$n_{u,j,r}(\tau)=\pi_{u,r}(\tau)+e_j(\tau)-\frac12(N(j-1)+N(j)-1-2r)$, where $\pi_{u,r}(\tau)$ is an
entry $u_{r+1}$ of $\lambda_u(\tau)$, or $-e_q(\tau)+\frac12(t+2-2s)$ with $1\le q\le m$ and
$1\le s\le t+1$; in both cases $|\pi_{u,r}(\tau)|\le5\Pi^5b$ by (2). Since $0\le r<N(j-1)\le\Pi$ and
$N(j)\le\Pi$, the last term has absolute value at most $\Pi$, so
$|n_{u,j,r}(\tau)|\le10\Pi^5b\le11\Pi^5b$.

(4) The host gives $1\le r_E$ and $\log r_E\le2\log n+\frac12(\log q+\log d)$, and the sizes give
$\log q,\log d\le x^{\varepsilon}$; so $0\le\log r_E\le2\Pi$. Then $L_a=\log x+x^{\varepsilon}+\log
r_E$ with $0\le\log x\le x\le\Pi$, so $0\le L_a\le4\Pi$. The sizes give $1\le g_i$ and $\log g_i\le
x^{\varepsilon/100}x^{\varepsilon}\le\Pi^2$, and $1\le n_V$, $\log n_V\le\Pi$; so $0\le
L_p=x^{\varepsilon}+\log r_E+\log g_0+\log g_1+\log n_V\le6\Pi^2$. Finally
$\log\Omega=\gamma+\log\pi-\log v-\operatorname{Re}\frac{L'}{L}(1,\chi_d)$ with $v=r_E/2$. Here
$\frac12<\gamma<\frac23$, $0\le\log\pi\le3$, $\log v\ge-\log2\ge-1$, the sizes give
$x^{\varepsilon/100}m\log v\le k$, hence $\log v\le\Pi$, and
$|\frac{L'}{L}(1,\chi_d)|\le x^{\varepsilon}\le\Pi$. So $|\log\Omega|\le7\Pi$.
\end{proof}

\begin{definition}\label{asm-division-kernel-sizes-local-terms}
Let $D$ be a member at a scale $x>x_1$ of a family $\mathcal F$, with ranks $k,m,n$, $f=2n$, Galois
datum with group $A$, level primes $g_0,g_1$, auxiliary prime $d$, integer $Q_S$ of the host and
quantity $X''$ of its size data. Let $z\in\{0,1\}\times A$ be the places of
Definition~\ref{pd-orbit-places}, with weights $w_j(z)$, positions $\operatorname{pos}_p(z,s)$ and
orbit lengths $\ell_p$, let $h_p$ be the elements of Definition~\ref{pd-gain-data}, and let $F(j)$ be
the Fock degree of the lift $j$ (Definition~\ref{asm-division-kernel-errors}(1)). For a lift $j$ and
a prime $p$, the \emph{explicit local term} is
\[
 F^{\mathrm{ex}}_j(p)=\begin{cases}
 -\frac12F(j)&\text{if }p\in\{g_0,g_1\},\\
 -\frac{k-1}2\sum_{t\in A}\min\big(w_j((0,t)),w_j((1,t))\big)&\text{if }p=d\nmid Q_S
 \text{ and }p\notin\{g_0,g_1\},\\
 F_{\mathrm{in}}(j,p)&\text{if }p\notin\{g_0,g_1,d\}\text{ is inert for }D\text{ and }p\le X'',\\
 0&\text{otherwise,}
 \end{cases}
\]
where $F_{\mathrm{in}}(j,p)=\sum_{\tau\in\mathcal T}(k-1)e_j(\tau)/(p^2-1)$ if $p$ is of inert class
and $h_p=0$, and otherwise
\[
 F_{\mathrm{in}}(j,p)=\frac1{\ell_p}\sum_{z\in\{0,1\}\times A}F(D'_{j,p,z}),\qquad
 D'_{j,p,z}(s)=(k-1)\,w_j\big(\operatorname{pos}_p(z,s)\big)\quad(s\in\mathbf Z/\ell_p\mathbf Z),
\]
with $F(D')=\gamma_p\sum_sD'(s)-M(D')$ formed as in Definition~\ref{pd-odd-level-programme} with
$l=p$ and $d=\ell_p$. The local terms of the size data are \emph{explicit} if
$F_j(w)=F^{\mathrm{ex}}_j(p_w)$ for every lift $j$ and every finite place $w$ of the coefficient
field.
\end{definition}
\leanref{Formalism/LocalRingCardinality/DataForced.lean}{Formalism.BasicNotions.SizeData₂.FDef}
\leanref{Formalism/LocalRingCardinality/DataForced.lean}{Formalism.BasicNotions.FvalP}
\leanref{Formalism/LocalRingCardinality/DataForced.lean}{Formalism.BasicNotions.FInert}
\leanref{Formalism/LocalRingCardinality/DataForced.lean}{Formalism.BasicNotions.FRamD}
\leanref{Formalism/LocalRingCardinality/DataForced.lean}{Formalism.BasicNotions.FTypeG}
\leanref{Formalism/LocalRingCardinality/DataForced.lean}{Formalism.BasicNotions.FOrb}

\begin{lemma}\label{asm-division-kernel-sizes-mertens}
For every real $y\ge1$, $\sum_{p<y}\frac{\log p}p\le2\log y$, the sum over the primes $p<y$.
\end{lemma}
\leanref{Formalism/SharpRestriction/GrowthBounds/JointsLit.lean}{Formalism.MainNotions.G.Joints.Lit.mertens_bound}
\leanref{Formalism/SharpRestriction/GrowthBounds/AdmissibleIntegers/FirstTermSizeBound.lean}{Formalism.MainNotions.WeilJacquetShalikaLaw.Proof.sz_mertens}
\begin{proof}
For an integer $n\ge1$, $\sum_{p\le n}\log p/p\le\log n+\log4$, a form of Mertens' first theorem
\cite{Mer74}. Indeed $n/p\le\lfloor n/p\rfloor+1$, so $n\sum_{p\le n}\log p/p$ is at most
$\sum_{p\le n}\lfloor n/p\rfloor\log p+\sum_{p\le n}\log p$. The first sum is at most $\log n!\le
n\log n$, since $v_p(n!)\ge\lfloor n/p\rfloor$, and the second is at most $n\log4$, since
$\prod_{p\le n}p\le4^n$.

If $y\ge4$, put $n=\lceil y\rceil-1$. Then $3\le n\le y$ and the primes $p<y$ are the primes $p\le
n$, so the sum is at most $\log n+\log4\le2\log y$. If $1\le y<4$, every prime $p<y$ has $\log p/p\le
\frac12\log y$, and there are at most four integers below $\lceil y\rceil\le4$; so the sum is at most
$2\log y$.
\end{proof}

\begin{lemma}\label{asm-division-kernel-sizes}
Let $\mathcal F$ be a family with exponent $\varepsilon$ (Definition~\ref{pd-family}), with the
exponent $C_t$ of the bound $t\le x^{C_t}$ of its size data, and put $c_t=\max(C_t,5)+3$. For a member
at a scale $x>x_1$, with ranks $k,m,t,n$, $f=2n$, $N(m)=k+m(t+1)$, $b=\exp(x^\varepsilon)$ and
$B_0=fmbx^{\varepsilon}$, put $\Pi=1+f+k+m+t+N(m)+x^{\varepsilon}$.
\begin{enumerate}
\item If $x\ge9$, then $\Pi\le x^{c_t}$; hence, for every real $C\le x$ and integer $e\ge0$,
  $C\,\Pi^eb\le x^{ec_t+1}b$.
\item For every real $C$, $CmB_0\le|C|\,\Pi^4b$.
\item There is $x_0$ such that, at every member at every scale $x\ge x_0$ with $x>x_1$,
  $G_a\le x^{4c_t+1}b$ and $F_o\le x^{9c_t+1}b$, with $G_a$ and $F_o$ as in
  Definition~\ref{asm-member-bounds}.
\end{enumerate}
\end{lemma}
\leanref{Formalism/SharpRestriction/Divisibility/SlopeExcess/ConstRankBaseLePolyFourth.lean}{Formalism.MainNotions.DIV.Proofs.szPol}
\leanref{Formalism/SharpRestriction/Divisibility/SlopeExcess/ConstRankBaseLePolyFourth.lean}{Formalism.MainNotions.DIV.Proofs.szPol_le_rpow}
\leanref{Formalism/SharpRestriction/Divisibility/SlopeExcess/ConstRankBaseLePolyFourth.lean}{Formalism.MainNotions.DIV.Proofs.bPol_le}
\leanref{Formalism/SharpRestriction/Divisibility/SlopeExcess/ConstRankBaseLePolyFourth.lean}{Formalism.MainNotions.DIV.Proofs.const_rank_base_le_poly_fourth}
\leanref{Formalism/SharpRestriction/Divisibility/HookProducts/CardLiftsEq.lean}{Formalism.MainNotions.DIV.HookProducts.Row4.ga_pow}
\leanref{Formalism/SharpRestriction/Divisibility/HookProducts/LwPNonneg.lean}{Formalism.MainNotions.DIV.HookProducts.HookProductWeightBounds.fo_pow}
\leanref{Formalism/SharpRestriction/Divisibility/HookProducts/CardLiftsEq.lean}{Formalism.MainNotions.DIV.HookProducts.Row4.ga_le_pol}
\leanref{Formalism/SharpRestriction/Divisibility/HookProducts/CardLiftsEq.lean}{Formalism.MainNotions.DIV.HookProducts.Row4.ga_le_at}
\leanref{Formalism/SharpRestriction/Divisibility/HookProducts/CardLiftsEq.lean}{Formalism.MainNotions.DIV.HookProducts.Row4.gainD_sum_le}
\leanref{Formalism/SharpRestriction/Divisibility/HookProducts/CardLiftsEq.lean}{Formalism.MainNotions.DIV.HookProducts.Row4.logX_le}
\leanref{Formalism/SharpRestriction/GrowthBounds/TopChoiceBounds.lean}{Formalism.MainNotions.G.TopChoiceBounds.bt_sum_by_prime}
\leanref{Formalism/SharpRestriction/Divisibility/HookProducts/LwPNonneg.lean}{Formalism.MainNotions.DIV.HookProducts.HookProductWeightBounds.fo_le_pol}
\leanref{Formalism/SharpRestriction/Divisibility/HookProducts/LwPNonneg.lean}{Formalism.MainNotions.DIV.HookProducts.HookProductWeightBounds.F_le}
\leanref{Formalism/SharpRestriction/Divisibility/HookProducts/LwPNonneg.lean}{Formalism.MainNotions.DIV.HookProducts.HookProductWeightBounds.FInert_le}
\leanref{Formalism/SharpRestriction/Divisibility/HookProducts/LwPNonneg.lean}{Formalism.MainNotions.DIV.HookProducts.HookProductWeightBounds.sumSz_le}
\begin{proof}
(1) Put $X=x^{c_t-1}$. The sizes of the family give $f,k,m,t,x^{\varepsilon}\le X$ and $N(m)\le2X$,
and $X\ge1$. So $\Pi\le8X\le xX=x^{c_t}$. For the consequence, $\Pi^e\le x^{ec_t}$ and $C\Pi^e\le
x\Pi^e$, as $\Pi^e\ge0$, and $b\ge0$.

(2) Since $B_0=fmbx^{\varepsilon}$, $CmB_0\le|C|\,m\cdot fmbx^{\varepsilon}$, and $m,f,x^{\varepsilon}
\le\Pi$ while $b\ge0$. So $CmB_0\le|C|\,\Pi^4b$.

(3) Take $x_0=256$, fix a member at a scale $x\ge x_0$ with $x>x_1$, and use
Lemma~\ref{asm-division-kernel-explicit-entries}. For a function $h\ge0$ on the primes,
\[
 \sum_{w\in\operatorname{supp}}\ell_we_wh(p_w)\le\sum_p h(p)\log p,\tag{$*$}
\]
the right side over the residue characteristics of the support, because the places $w$ over a
prime $p$, with residue degrees $f_w$, have
$\sum_{w\mid p}\ell_we_w=\sum_{w\mid p}e_wf_w\log p/[L:\mathbf Q]=\log p$. Also
$\log(X''+1)\le4\Pi$: since $X''=xnX'$, $X''+1\le(x+1)(n+1)X'$, so $\log(X''+1)\le
x+n+x^{\varepsilon/3}$, and the sizes give $x,n\le k$.

\emph{The bound for $G_a$.} Let $w$ be a place and $p=p_w$. If $p>X''$, then
$\mathrm{gain}^\flat_j(w)=0$ for every lift $j$: $\mathrm{gain}_j(w)=0$ off the support and when $p$
is not inert (Definition~\ref{pd-size-data}(3)), and otherwise $\mathrm{gain}_j(w)=g_j(p)$, since the
gain data of the family fill its size data (Definition~\ref{pd-gain-quantities}), and $g_j(p)=0$ for
$p>X''$. If $p\le X''$, the same reasons and the upper bounds in the proof of
Lemma~\ref{pd-gain-holds-gain-bound}, which use only the parameter choices and the windows of the
exponents, give $\mathrm{gain}^\flat_j(w)\le4|T|a/p$. So $\sum_j\mathrm{gain}^\flat_j(w)\le
h_G(p)$, where $h_G(p)=4m|T|a/p$ for $p\le X''$ and $h_G(p)=0$ otherwise. The fill also gives
$\tau_w=\tau(p)\ge0$ (Definition~\ref{pd-gain-quantities}(3)). As $\ell_we_w\ge0$, ($*$) and
Lemma~\ref{asm-division-kernel-sizes-mertens} with $y=X''+1$ give
\[
 G_a\le2\sum_{w\in\operatorname{supp}}\ell_we_wh_G(p_w)\le8m|T|a\sum_{p\le X''}\frac{\log p}p
 \le16m|T|a\log(X''+1).
\]
The sizes give $|T|\le2n$ and $nx\le k$, so $m|T|x\le2mk\le(m+k)^2\le\Pi^2$; and $a=xb$. So
$G_a\le16\Pi^2b\cdot4\Pi\le64\Pi^4b$, and (1) with $C=64\le x$ and $e=4$ gives
$G_a\le x^{4c_t+1}b$.

\emph{The bound for $F_o$.} The members of the families of Definition~\ref{pd-family} have explicit
local terms (Definition~\ref{asm-division-kernel-sizes-local-terms}); this is part of their error
data. Put $W=4\Pi^5b$, so that $0\le w_j(z)\le W$ for every place $z$. Let $j$ be a lift, $w$ a place
and $p=p_w$. Then $F_j(w)\le h_F(p)$, where $h_F(p)=16\Pi^7b/p$ for $p\le X''$ and $h_F(p)=0$
otherwise:
\begin{itemize}
\item if $p\in\{g_0,g_1\}$, $F_j(w)=-\frac12F(j)\le0$, as the Fock degree is a sum of absolute
  values;
\item if $p=d\nmid Q_S$, $F_j(w)\le0$, as $k\ge1$ and the weights are nonnegative;
\item if $p$ is inert and $p\le X''$, $F_j(w)=F_{\mathrm{in}}(j,p)\le4\Pi(k-1)W/p\le16\Pi^7b/p$. In
  the first case of $F_{\mathrm{in}}$ this holds because $p\le p^2-1$ and $|\mathcal T|=f\le\Pi$. In
  the second, the zero family is a flow for $D'_{j,p,z}\ge0$, so $M(D'_{j,p,z})\ge0$ and
  $F(D'_{j,p,z})\le\gamma_p\sum_sD'_{j,p,z}(s)\le\frac2p\ell_p(k-1)W$, as $\gamma_p=p/(p^2-1)\le2/p$;
  the $2n=f\le\Pi$ places $z$ and the factor $1/\ell_p$ give $F_{\mathrm{in}}(j,p)\le2\Pi(k-1)W/p$;
\item otherwise $F_j(w)=0$.
\end{itemize}
Since $\mathrm{sz}_w(y)=2\ell_we_wy$, ($*$) and Lemma~\ref{asm-division-kernel-sizes-mertens} give
$\sum_{w\in\operatorname{supp}}\mathrm{sz}_w(F_j(w))\le2\cdot16\Pi^7b\cdot2\log(X''+1)\le256\Pi^8b$.
Summing over the $m\le\Pi$ lifts gives $F_o\le256\Pi^9b$, and (1) with $C=256\le x$ and $e=9$ gives
$F_o\le x^{9c_t+1}b$.
% Lean: (1) is szPol_le_rpow (each of f, k, m, t and L is at most X = x^(c+2), N(m) <= 2X, then
% szPol <= 8X <= x X) and bPol_le; (2) is const_rank_base_le_poly_fourth. (3) is
% HookProducts.Row4.ga_pow (ga_le_pol with the fills of the family and the Mertens bound, then bPol_le
% with e = 4, x >= 64) and HookProductWeightBounds.fo_pow (fo_le_pol with FDef and the Mertens bound,
% then bPol_le with e = 9, x >= 256). Size of (3): proof closure of ga_pow and fo_pow, 118 theorems,
% about 1,600 lines in 14 files.
\end{proof}

\begin{lemma}\label{asm-division-kernel-explicit-terms}
Let $D$ be a member at a scale $x>x_1$ of a family $\mathcal F$, with the notation of
Definition~\ref{asm-division-kernel-explicit-heights} for a real constant $C$, and let $\Pi$ be as in
Lemma~\ref{asm-division-kernel-sizes}. Let $C_T$ be real with $\log T_u\le C_Tmx^{\varepsilon}$ for
every word $u$. Then for every word $u$, lift $1\le j\le m$ and $\tau\in\mathcal T$:
\begin{enumerate}
\item $-S(u)\le(2|C_T|+16|C|)\Pi^9b$;
\item $\sum_{r<N(j-1)}\phi(n_{u,j,r}(\tau))-n_{u,j}(\tau)\log\Omega\ge-300\Pi^7b$, and
  $-K_j(u)\le(300+5|C|)\Pi^8b$;
\item $\log\mathrm{hook}_j(\tau)\le160\Pi^7b$, and $\ell_j\le(160+6|C|)\Pi^8b$.
\end{enumerate}
\end{lemma}
\leanref{Formalism/SharpRestriction/Divisibility/SlopeExcess/ConstRankBaseLePolyFourth.lean}{Formalism.MainNotions.DIV.Proofs.HHParts.negSZbP_le}
\leanref{Formalism/SharpRestriction/Divisibility/SlopeExcess/ConstRankBaseLePolyFourth.lean}{Formalism.MainNotions.DIV.Proofs.HHParts.sumT_le}
\leanref{Formalism/SharpRestriction/Divisibility/SlopeExcess/ConstRankBaseLePolyFourth.lean}{Formalism.MainNotions.DIV.Proofs.HHParts.kerPlace_lower}
\leanref{Formalism/SharpRestriction/Divisibility/SlopeExcess/ConstRankBaseLePolyFourth.lean}{Formalism.MainNotions.DIV.Proofs.HHParts.neg_size_term_le_poly_eighth}
\leanref{Formalism/SharpRestriction/Divisibility/SlopeExcess/ConstRankBaseLePolyFourth.lean}{Formalism.MainNotions.DIV.Proofs.HHParts.logGamma_le}
\leanref{Formalism/SharpRestriction/Divisibility/SlopeExcess/ConstRankBaseLePolyFourth.lean}{Formalism.MainNotions.DIV.Proofs.HHParts.sigPar_anti}
\leanref{Formalism/SharpRestriction/Divisibility/SlopeExcess/ConstRankBaseLePolyFourth.lean}{Formalism.MainNotions.DIV.Proofs.HHParts.logHook_le}
\leanref{Formalism/SharpRestriction/Divisibility/SlopeExcess/ConstRankBaseLePolyFourth.lean}{Formalism.MainNotions.DIV.Proofs.HHParts.L6lo_le}
\begin{proof}
We use Lemma~\ref{asm-division-kernel-explicit-entries} throughout, and $\Pi^a\le\Pi^ab\le\Pi^{a'}b$
for $a\le a'$.

(1) Since $m,x^{\varepsilon}\le\Pi$, $\log T_u\le|C_T|\Pi^2$. For $t\in T$, $a_t\log T_u\le0$ if
$\log T_u\le0$, and $a_t\log T_u\le2\Pi b\cdot|C_T|\Pi^2$ otherwise; summing over the at most $\Pi$
top classes gives $\sum_ta_t\log T_u\le2|C_T|\Pi^4b$. Put
$Q=|T|b'N(m)+f(m\eta+S_c)N(m)+fN(m)^2$. Then $0\le Q\le\Pi\cdot\Pi b\cdot\Pi+\Pi\cdot2\Pi
b\cdot\Pi+\Pi^3\le4\Pi^3b$, and with $0\le L_a\le4\Pi$, $-CQL_a\le|C|QL_a\le16|C|\Pi^4b$. So
$-S(u)=\sum_ta_t\log T_u-CQL_a\le(2|C_T|+16|C|)\Pi^9b$.

(2) Put $M=11\Pi^5b\ge1$. For real $y$ with $|y|\le M$, $|y\log|y||\le M\log M+1$: if $|y|\le1$ it is
at most $1$, and otherwise $0\le|y|\log|y|\le M\log M$. Hence $\phi(y)\ge-(M\log M+M+1)$. By
Lemma~\ref{asm-division-kernel-explicit-entries}(3), summing over the $N(j-1)\le\Pi$ rows,
\[
 \sum_{r<N(j-1)}\phi(n_{u,j,r}(\tau))\ge-N(j-1)(M\log M+M+1),\qquad
 |n_{u,j}(\tau)\log\Omega|\le N(j-1)\,M\cdot7\Pi.
\]
Now $\log M=\log11+5\log\Pi+x^{\varepsilon}\le10+5\Pi+\Pi\le16\Pi$, $M\le\Pi M$ and $1\le\Pi^6b$, so
$M\log M+M+1+7\Pi M\le24\Pi M+\Pi^6b=265\Pi^6b\le300\Pi^6b$. Multiplying by $N(j-1)\le\Pi$ gives the
first bound. Summing it over the $f\le\Pi$ embeddings, and using $-CfN(m)^2L_a\le|C|\Pi^3\cdot4\Pi$ and
$-Cfm^3N(m)^2x^{\varepsilon}\le|C|\Pi^7$, gives $-K_j(u)\le300\Pi^8b+5|C|\Pi^8b$.

(3) The sequence $\sigma_i=\sigma_i(\tau)$ drops by at least $1$ at each step, for $i+1<k-1+m$. For
$i+1<k-1$ the $v_i$ are integers with $v_i>u_{i+1}>v_{i+1}$; at $i+1=k-1$ the integers satisfy
$v_{k-1}>u_k>-b_1$; and for $i\ge k-1$ the gaps of Definition~\ref{th-hc-parameter}(4) give
$b_{q+1}-b_q\ge t+3$. So $\sigma_i\ge\sigma_l+(l-i)$ for $i\le l<k-1+m$. Put $n=k+j-1$. Since
$n_V=k+m-1+t$, for $i<n$
\[
 y_i:=a_i+n-i=\sigma_i+e_j(\tau)-\tfrac12(t-1).
\]
As $\sigma_i\ge\sigma_{k+j-2}+(k+j-2-i)\ge-b_j(\tau)$ and $e_j(\tau)-b_j(\tau)\ge(t+3)/2$, $y_i\ge2$;
and $y_i\le5\Pi^5b+4\Pi^5b+\frac12\le Y:=10\Pi^5b$. For $y\ge2$, $\log\Gamma(y)\le y\log y$: $\Gamma$
increases on $[2,\infty)$, so $\Gamma(y)\le\Gamma(\lceil y\rceil)=(\lceil y\rceil-1)!\le
(\lceil y\rceil-1)^{\lceil y\rceil-1}$ with $1\le\lceil y\rceil-1\le y$. So each $\log\Gamma(y_i)\le
Y\log Y$. For $i<l<n$, $a_i-a_l+l-i=\sigma_i-\sigma_l\ge l-i\ge1$, so the subtracted logarithms are
nonnegative. As $n\le N(m)\le\Pi$ and $\log Y\le\log10+5\log\Pi+x^{\varepsilon}\le16\Pi$,
\[
 \log\mathrm{hook}_j(\tau)\le nY\log Y\le\Pi\cdot10\Pi^5b\cdot16\Pi=160\Pi^7b.
\]
Summing over the $f\le\Pi$ embeddings, and using $-Cfn_V^2L_p\le|C|\Pi^3\cdot6\Pi^2$, gives
$\ell_j\le160\Pi^8b+6|C|\Pi^8b$.
\end{proof}

\begin{lemma}\label{asm-division-kernel-explicit}
Let $\mathcal F$ be a family with exponent $\varepsilon$ that satisfies the size condition on $T_u$
(Definition~\ref{asm-designed}), with explicit heights $h_u$
(Definition~\ref{asm-division-kernel-explicit-heights}) and $c_t$ as in
Lemma~\ref{asm-division-kernel-sizes}. Then there is $x_0$ such that at every member at every scale
$x\ge x_0$ with $x>x_1$ and every word $u$,
\[
 -\big(h_u+\overline{\log\kappa_R}\big)\le x^{9c_t+1}b,
\]
with $\overline{\log\kappa_R}$ as in Lemma~\ref{asm-division-kernel-sharp}.
\end{lemma}
\leanref{Formalism/SharpRestriction/Divisibility/SlopeExcess/ConstRankBaseLePolyFourth.lean}{Formalism.MainNotions.DIV.Proofs.hExpl6_lower_pow}
\leanref{Formalism/SharpRestriction/Divisibility/SlopeExcess/ConstRankBaseLePolyFourth.lean}{Formalism.MainNotions.DIV.Proofs.hExpl6_lower}
\leanref{Formalism/SharpRestriction/Divisibility/SlopeExcess/ConstRankBaseLePolyFourth.lean}{Formalism.MainNotions.DIV.Proofs.hExpl6_lower_at}
\begin{proof}
Let $C_{T_u}$ and $x_T$ be the constant and the scale of the size condition on $T_u$, so that
$\log T_u\le C_{T_u}mx^{\varepsilon}$ at every member at every scale $x\ge x_T$. Put
$C'=500+2|C_{T_u}|+30|C_{sz}|$ and $x_0=\max(x_T,9,C')$. Fix a member at a scale $x\ge x_0$ with
$x>x_1$, and a word $u$.

By Definition~\ref{asm-division-kernel-explicit-heights}, with $C=C_{sz}$,
\[
 -\big(h_u+\overline{\log\kappa_R}\big)=-S(u)-\sum_{j=1}^mK_j(u)+\sum_{j=1}^m\ell_j.
\]
Lemma~\ref{asm-division-kernel-explicit-terms}, with $C_T=C_{T_u}$, bounds the three terms, and the
lifts are $m\le\Pi$ in number:
\[
 -\big(h_u+\overline{\log\kappa_R}\big)\le(2|C_{T_u}|+16|C_{sz}|)\Pi^9b
 +\Pi(300+5|C_{sz}|)\Pi^8b+\Pi(160+6|C_{sz}|)\Pi^8b\le C'\Pi^9b.
\]
Since $x\ge9$ and $C'\le x$, Lemma~\ref{asm-division-kernel-sizes}(1) with $e=9$ gives
$C'\Pi^9b\le x^{9c_t+1}b$.
% Lean: hExpl6_lower_pow = hExpl6_lower (hExpl6_lower_at: -(hExpl6 + meanE log kapR) <= (500 + 2|CTu| +
% 30|Csz|) szPol^9 b, from negSZbP_le, neg_size_term_le_poly_eighth, L6lo_le and sum_card_le), then
% bPol_le with e = 9. hExpl6 = SZbP + sum_j kernelPlaceMeanBound - sum_j L6lo - meanE log kapR, so the
% quantity bounded is minus the record height. Size: proof closure of hExpl6_lower_pow, 93 theorems,
% about 1,100 lines in 4 files.
\end{proof}

\begin{definition}\label{asm-division-kernel-first-lattice-slopes-member}
Let $D$ be a definite-side datum with ranks $k,m,t$ and $v$ as in
Definition~\ref{asm-global-numerical}, and for a word $u$ and $\tau\in\mathcal T$ let $\lambda_i(\tau)$
($1\le i\le k$) be as there. Put
\[
 \sigma(\tau)=\sum_{i=1}^k\lambda_i(\tau),\qquad
 \sigma^+(\tau)=\sum_{i=1}^k\big(\lambda_i(\tau)-\lambda_k(\tau)\big),\qquad
 w(\tau)=\lambda_1(\tau)-\lambda_k(\tau).
\]
The datum $D$ has the \emph{member lower slope bound} if every word $u$ whose columns have values
(Definition~\ref{asm-global-column-values}) satisfies
\[
 -\sum_\tau\frac{\sigma(\tau)}k\log T_u-\sum_\tau\Big(\sigma^+(\tau)\log k+k^2\log\big(w(\tau)+k\big)\Big)
 \le\mu(M_u),
\]
and, for a real number $C'$, the \emph{member upper slope bound} with constant $C'$ if every such $u$
satisfies
\[
 \mu_{\max}(M_u)\le\sum_\tau\Big(C'\sigma^+(\tau)\log(kv)+k^2\log\big(w(\tau)+k\big)
 +|\lambda_k(\tau)|\log T_u\Big),
\]
where $M_u$ is the first lattice $\Lambda_{u,0}$ of the chain of $u$.
\end{definition}
\leanref{Formalism/SharpRestriction/GlobalSchurBounds/IntegrableOfContinuousCircleMeasure.lean}{Formalism.MainNotions.flipTwoThree.SlopeLow}
\leanref{Formalism/SharpRestriction/GlobalSchurBounds/IntegrableOfContinuousCircleMeasure.lean}{Formalism.MainNotions.flipTwoThree.SlopeMax}
\leanref{Formalism/SharpRestriction/GlobalSchurBounds/IntegrableOfContinuousCircleMeasure.lean}{Formalism.MainNotions.flipTwoThree.SlopeLowerOpenPart}
\leanref{Formalism/SharpRestriction/GlobalSchurBounds/IntegrableOfContinuousCircleMeasure.lean}{Formalism.MainNotions.flipTwoThree.SlopeMaxOpenPart}
\leanref{Formalism/SharpRestriction/GlobalSchurBounds/IntegrableOfContinuousCircleMeasure.lean}{Formalism.MainNotions.flipTwoThree.colAbs}
\leanref{Formalism/SharpRestriction/GlobalSchurBounds/IntegrableOfContinuousCircleMeasure.lean}{Formalism.MainNotions.flipTwoThree.colAbsP}
\leanref{Formalism/SharpRestriction/GlobalSchurBounds/IntegrableOfContinuousCircleMeasure.lean}{Formalism.MainNotions.flipTwoThree.firstColumnSpread}

\begin{lemma}\label{asm-division-kernel-first-lattice-slopes-members}
Let $\mathcal F$ be a family with exponent $\varepsilon$ (Definition~\ref{pd-family}). There is $x_0$
such that every member at every scale $x\ge x_0$ with $x>x_1$ has the member lower slope bound and the
member upper slope bound with constant $4$
(Definition~\ref{asm-division-kernel-first-lattice-slopes-member}).
\end{lemma}
\leanref{Formalism/LocalToGlobal/GlobalTheory/GL4/LowerSlopeInequalityFamily.lean}{Formalism.L2.GlobalTheoryTheory.TraceIntegrality.lower_slope_inequality_family}
\leanref{Formalism/LocalToGlobal/GlobalTheory/GL4/PStmtFree.lean}{Formalism.L2.GlobalTheoryTheory.TraceIntegrality.slope_inequality_family_four}
\leanref{Formalism/LocalToGlobal/GlobalTheory/GL4/PStmtFree.lean}{Formalism.L2.GlobalTheoryTheory.TraceIntegrality.slope_inequality_constant_four}
\leanref{Formalism/LocalToGlobal/GlobalTheory/GL4/ColumnLatticeExactSlope.lean}{Formalism.L2.GlobalTheoryTheory.TraceIntegrality.column_lattice_exact_slope}
\leanref{Formalism/SharpRestriction/GlobalSchurBounds/SlopeMaxOpenOfHookTwoLeV.lean}{Formalism.MainNotions.flipTwoThree.familySlopeLowerOpen_of_std}
\begin{proof}
\openstep{The two bounds, which apply the slope inequality of Theorem~\ref{hs-bost-slope-inequality} to
the coordinates of the columns, the lower one through the exact slopes of the columns and the hook
lengths and the upper one with the constant $4$, are not yet written out.}
% Lean: the lower bound is lower_slope_inequality_family = familySlopeLowerOpen_of_std
% standard_slope_lower hook_length_bound (column_exact_slope_family Fm), where column_exact_slope_family
% is column_lattice_exact_slope at every definite-side datum (x0 = 0); the upper bound is
% slope_inequality_family_four = slope_inequality_constant_four at every definite-side datum (x0 = 0),
% that is slope_bound_of_max_slope hookBound_holds D (slopeMaxGiven_free D). Size (tools/lean_size.py):
% lower about 11,600 lines (column_lattice_exact_slope about 9,900, standard_slope_lower about 2,800,
% hook_length_bound about 1,600, overlapping); upper about 19,000 lines, of which about 8,500 are not
% below the lower one; together about 20,100 lines.
\end{proof}

\begin{lemma}\label{asm-division-kernel-first-lattice-slopes-chain}
Let $C_n$ be real.
\begin{enumerate}
\item There is $C$ with the following property. Let $D$ be a definite-side datum with the member
  lower slope bound (Definition~\ref{asm-division-kernel-first-lattice-slopes-member}), let
  $x\ge\max(C_n,1)$, let $u$ be a word of $D$ whose columns have values, and let $b,L_0,b',S_c,\eta$ be
  real numbers such that the numerical conditions and the weight conditions with $C_n$ hold at $u$
  for $(x,b,L_0,b',S_c,\eta)$ (Definition~\ref{asm-global-numerical}). Then the first inequality of the
  column slope bounds holds at $u$ with $C_5=C$.
\item For every real $C'$ there is $C$ with the same property for a datum $D$ with the member upper
  slope bound with constant $C'$ and the second inequality of the column slope bounds with $C_6=C$.
\end{enumerate}
\end{lemma}
\leanref{Formalism/SharpRestriction/GlobalSchurBounds/IntegrableOfContinuousCircleMeasure.lean}{Formalism.MainNotions.flipTwoThree.ramifiedLattice_of_chain_estimates}
\leanref{Formalism/SharpRestriction/GlobalSchurBounds/IntegrableOfContinuousCircleMeasure.lean}{Formalism.MainNotions.flipTwoThree.classMap_estimate_of_partial}
\leanref{Formalism/GL3MeanSquares/HookFormulaTorsion/SquareMomentArithmeticBound.lean}{Formalism.HermitianLattices.TraceIntegrality.zeta_proposition_second}
\leanref{Formalism/GL3MeanSquares/HookFormulaTorsion/SquareMomentArithmeticBound.lean}{Formalism.HermitianLattices.TraceIntegrality.zeta_proposition_third}

This is a computation with the numerical and weight conditions, given here with its result. Put
$K=\max(C_n,1)$ and $C_0=2\max(C_n,0)+4$. The conditions give, with $a=xb$ and $S=KS_c$: the
$\lambda_i(\tau)$ are integers decreasing in $i$; at a top place $a\le\lambda_1(\tau)\le2a$,
$b'\le-\lambda_k(\tau)\le2b'$ and $\sigma(\tau)\le\lambda_1(\tau)+kS$; at the other places
$\sigma^+(\tau)\le(k+3)S$ and $\sigma(\tau)\le C_0kS$; and $\log(w(\tau)+k)\le10L_0$. For (1) these give
\[
 \sigma^+(\tau)\log k+k^2\log\big(w(\tau)+k\big)\le\epsilon_\tau2a\log k+4k(b'+S)\log k+10k^2L_0,
\]
with $\epsilon_\tau=1$ at the top places and $0$ elsewhere, and
$\sum_\tau\frac{\sigma(\tau)}k\log T_u\le\sum_{\tau\ \mathrm{top}}\frac{\lambda_1(\tau)}k\log T_u+
C_0fS\log T_u$, so the member lower slope bound gives the first inequality with a constant
$C_1=11+C_0$ and $S$ in place of $S_c$; as $b'+KS_c\le K(b'+S_c)$, $C=\max(C_1,0)K$ works. For (2) the
terms $C'\sigma^+(\tau)\log(kv)$, $k^2\log(w(\tau)+k)$ and $|\lambda_k(\tau)|\log T_u$ are bounded in
the same way, after replacing $C'$ by $K'C_0$ with $K'=\max(1,C'/C_0)$; this gives
$C=\max(C_2,0)\,K\,K'$ for the constant $C_2$ of that computation.
% Lean: ramifiedLattice_of_chain_estimates Cn applies zeta_proposition_second (2 max Cn 0 + 4) (its
% constant is 11 + C0) after col_setup (the entries lam, antitone, with the top and other bounds) and
% numbers_of_nums (Numbers with a = x b, L = L0, S = K Sc), with hZ4 rewriting the member lower bound
% SlopeLow, and finishes with arith5 (b' + K Sc <= K (b' + Sc)). classMap_estimate_of_partial Cn C5
% applies zeta_proposition_third in the same way with K' = max 1 (C5 / (2 Cn + 4)), the bounds hterm,
% and arith6. Size: about 920 lines for the two, with familySlopeLower_of_open and
% familySlopeMax_of_open about 950.

\begin{lemma}\label{asm-division-kernel-first-lattice-slopes}
Every family with exponent $\varepsilon$ (Definition~\ref{pd-family}) satisfies the column slope
bounds (Definition~\ref{asm-global-numerical}).
\end{lemma}
\leanref{Formalism/SharpRestriction/GlobalSchurBounds/IntegrableOfContinuousCircleMeasure.lean}{Formalism.MainNotions.flipTwoThree.familySlopeLower_of_open}
\leanref{Formalism/SharpRestriction/GlobalSchurBounds/IntegrableOfContinuousCircleMeasure.lean}{Formalism.MainNotions.flipTwoThree.familySlopeMax_of_open}
\leanref{Formalism/LocalToGlobal/GlobalTheory/GL4/LowerSlopeInequalityFamily.lean}{Formalism.L2.GlobalTheoryTheory.TraceIntegrality.lower_slope_inequality_family}
\leanref{Formalism/LocalToGlobal/GlobalTheory/GL4/PStmtFree.lean}{Formalism.L2.GlobalTheoryTheory.TraceIntegrality.slope_inequality_family_four}
\begin{proof}
Let $\mathcal F$ be a family with exponent $\varepsilon$ and $C_n$ real. By
Lemma~\ref{asm-division-kernel-first-lattice-slopes-members} there is $x_0$ such that every member at
every scale $x\ge x_0$ with $x>x_1$ has the member lower slope bound and the member upper slope bound
with constant $4$. Let $C$ be the constant of
Lemma~\ref{asm-division-kernel-first-lattice-slopes-chain}(1) for $C_n$, and $C'$ that of part (2)
for $C_n$ and the constant $4$. Put $C_5=C$, $C_6=C'$ and $x_0'=\max(x_0,C_n,1)$. Let $x\ge x_0'$ with
$x>x_1$, let $D$ be a member at $x$, $u$ a word of $D$ whose columns have values, and
$b,L_0,b',S_c,\eta$ real numbers for which the numerical and the weight conditions with $C_n$ hold at
$u$. Then $x\ge\max(C_n,1)$, and Lemma~\ref{asm-division-kernel-first-lattice-slopes-chain} gives the
two inequalities of the column slope bounds at $u$.
% Lean: the lower bound is familySlopeLower_of_open Fm (lower_slope_inequality_family Fm) and the upper
% bound is familySlopeMax_of_open Fm 4 (slope_inequality_family_four Fm); both take x0' = max x0 (max Cn
% 1). Fact (d) of the proof of Theorem asm-global-link states the same bounds (hZa, hZb). The open part
% is Lemma asm-division-kernel-first-lattice-slopes-members, about 20,100 lines; the derivation here and
% in Lemma asm-division-kernel-first-lattice-slopes-chain is about 950 lines.
\end{proof}

\begin{lemma}\label{asm-division-kernel-first-lattice-arith}
Let $D$ be a member at a scale $x>x_1$ of a family $\mathcal F$ with exponent $\varepsilon$, with
ranks $k,m,n$, $f=2n$, top classes $T$, the quantities $a=xb$, $b'$ and $S_c$ of its size data, root
discriminant $r_E$ of its host and $v=r_E/2$, and let $\Pi$ be as in
Lemma~\ref{asm-division-kernel-sizes}. Let $C$ and $C_T$ be real with
$\log T_u\le C_Tmx^{\varepsilon}$ for every word $u$. If a word $u$ satisfies
\[
 \mu_{\max}(M_u)\le C|T|a\log(kv)+Cf(b'+S_c)\big(k\log(kv)+\log T_u\big)+Cfk^2x^{\varepsilon},
\]
then $\mu_{\max}(M_u)\le|C|(12+2|C_T|)\Pi^4b$.
\end{lemma}
\leanref{Formalism/LocalToGlobal/Divisibility/DivisibilityRemainder/ValueBounds/SlopeMaxB.lean}{Formalism.L2.DIV.SlopeMaxB.muMax_le_pol}
\begin{proof}
By Lemma~\ref{asm-division-kernel-explicit-entries} and its proof, $|T|\le\Pi$,
$f,k,m,x^{\varepsilon}\le\Pi$, $0\le b'\le a\le\Pi b$, $0<S_c\le\Pi b$, $1\le\Pi$, $1\le b$ and
$0\le\log r_E\le2\Pi$. As $k\ge m+3\ge5$ and
$r_E\ge1$, $kv\ge1$, and $\log(kv)\le\log k+\log r_E\le k+2\Pi\le3\Pi$; so $0\le\log(kv)\le3\Pi$.
Also $0\le\log T_u\le|C_T|mx^{\varepsilon}\le|C_T|\Pi^2$. Hence the three products on the right are
nonnegative, and
\[
 |T|a\log(kv)\le3\Pi^4b,\qquad f(b'+S_c)\big(k\log(kv)+\log T_u\big)\le\Pi\cdot2\Pi b\cdot(3+|C_T|)\Pi^2,
 \qquad fk^2x^{\varepsilon}\le\Pi^4\le3\Pi^4b.
\]
Their sum $S$ is at most $(12+2|C_T|)\Pi^4b$, and $\mu_{\max}(M_u)\le CS\le|C|S$.
\end{proof}

\begin{lemma}\label{asm-division-kernel-first-lattice}
Let $\mathcal F$ be a family with exponent $\varepsilon$ and a twist selection, that satisfies the
size condition on $T_u$ (Definition~\ref{asm-designed}), and let $c_t$ be as in
Lemma~\ref{asm-division-kernel-sizes}. Then at all large scales, at every member and every word $u$,
$\mu_{\max}(M_u)\le x^{4c_t+1}b$, where $M_u$ is the first lattice of the chain of $u$
(Definition~\ref{asm-member-bounds}).
\end{lemma}
\leanref{Development/AutomorphicArithmetic/ThetaCorrespondence/ThetaPairLifts/ThirdPlaceSeries/HbmuSlopeMaxB.lean}{Formalism.L2.DIV.SlopeMaxB.eventually_muMax_le_pow_ebMu_mul_of_tuSize}
\leanref{Formalism/LocalToGlobal/Divisibility/DivisibilityRemainder/ValueBounds/SlopeMaxB.lean}{Formalism.L2.DIV.SlopeMaxB.boundB_of_values}
\leanref{Formalism/LocalToGlobal/Divisibility/DivisibilityRemainder/ValueBounds/SlopeMaxB.lean}{Formalism.L2.DIV.SlopeMaxB.boundB_of_slopes}
\begin{proof}
Let $C_{T_u}$ be the constant of the size condition on $T_u$, so that
$\log T_u\le C_{T_u}mx^{\varepsilon}$ at every word of every member at all large scales.

\emph{Step 1: the columns have values.} At every member $D$ at every scale beyond the threshold of the
twist selection, the tied data of the twist record that it chooses give $D$ tied columns, as in the
proof of Lemma~\ref{asm-proviso-bottom-dimension}. At every scale $x\ge3$ with $x>x_1$, every class
$c\in\mathrm{cls}_0$ of the chain of every word has trivial stabilizer $\Gamma_c$, by
Lemma~\ref{asm-proviso-trivial-stabilizers}, as in the proof of
Lemma~\ref{asm-global-second-layer}(3), which uses only the sizes of the family. So, by
Lemma~\ref{asm-global-link-column-values}, at all large scales the columns of every word of every
member have values (Definition~\ref{asm-global-column-values}).

\emph{Step 2: the numerical and weight conditions.} The sizes of the family give
$|\operatorname{Re}\frac{L'}{L}(1,\chi_d)|\le x^{\varepsilon}$ at every member. So
Lemma~\ref{asm-global-numerical-hold}, with $C_c=C_{T_u}$, $C_n=C_{n,0}$ and $c_L=3$, gives at all large
scales the numerical and the weight conditions with $C_n$ at every word of every member, for
$(x,b,3x^{\varepsilon},b',S_c,\eta)$. In particular $v\ge1$.

\emph{Step 3: the slope bound.} By Lemma~\ref{asm-division-kernel-first-lattice-slopes}, the column
slope bounds for this $C_n$ give constants $C_5$ and $C_6$; put $C=3\max(C_5,C_6,0)$. At all large
scales, Steps 1 and 2 and the upper column slope bound give, at every word $u$ of every member, with
$a=xb$,
\[
 \mu_{\max}(M_u)\le C_6|T|a\log(kv)+C_6f(b'+S_c)\big(k\log(kv)+\log T_u\big)+3C_6fk^2x^{\varepsilon}.
\]
As $k\ge1$, $v\ge1$, $T_u\ge2$ and $b'+S_c\ge0$, the three products are nonnegative, and $C_6\le C$,
$3C_6\le C$; so the hypothesis of Lemma~\ref{asm-division-kernel-first-lattice-arith} holds with $C$
and $C_T=C_{T_u}$.
% The Lean forms both halves (FrobeniusExpansionCongruence, by J.orderBound_of_fifth_sixth from
% ConclusionF5 and ConclusionF6 with C = max (max C5 C6) 0 * 3) and uses only the upper one here.

\emph{Step 4.} Lemma~\ref{asm-division-kernel-first-lattice-arith} gives
$\mu_{\max}(M_u)\le C'\Pi^4b$ with $C'=|C|(12+2|C_{T_u}|)$. At every scale $x\ge\max(9,C')$,
Lemma~\ref{asm-division-kernel-sizes}(1) with $e=4$ gives $C'\Pi^4b\le x^{4c_t+1}b$.
% Lean: eventually_muMax_le_pow_ebMu_mul_of_tuSize = boundB_of_values Fm hT (tCV_of_ties Fm
% (colTies_of_sel_alt s)), where colTies_of_sel_alt takes the tied columns (s.2 ...).T.cols of the twist
% selection and stabF; boundB_of_values = boundB_of_slopes (slopes_of_values: J.slopes_family_of_column_values
% with top2c2_dlSize, familySlopeLower_of_open and familySlopeMax_of_open with
% slope_inequality_family_four), and boundB_of_slopes is muMax_le_pol followed by bPol_le with e = 4.
% The column values are fact (e) of the proof of Theorem asm-global-link (RT.tCV_of_tied). Size: proof
% closure of eventually_muMax_le_pow_ebMu_mul_of_tuSize, 1,701 theorems, about 29,000 lines in 267 files.
% Step 1 is tCV_of_ties = tCV_of_tied with ColValS.colValOfTiedS_holds (Lemma
% asm-global-link-column-values) and colTies_of_sel_alt (stabF); Step 2 is J.numsL_wtsL_fam =
% nums_wts_of_tuSize = nums_wts_of_fam (Lemma asm-global-numerical-hold) with top2c2_dlSize; Step 3 is
% J.slopes_family_of_column_values = family_orderBound_of_fifth_sixth (orderBound_of_fifth_sixth); Step 4
% is boundB_of_slopes (muMax_le_pol, bPol_le, Lemma asm-division-kernel-first-lattice-arith). The open
% part is Lemma asm-division-kernel-first-lattice-slopes.
\end{proof}

\begin{lemma}\label{asm-division-kernel-gauge}
Let $\mathcal F$ be a family with exponent $\varepsilon$ that satisfies part (6) of the division
kernel bounds (Definition~\ref{is-division-kernel-bounds}), and let $c_t$ be as in
Lemma~\ref{asm-division-kernel-sizes}. Then there is $x_0$ such that at every member at every scale
$x\ge x_0$ with $x>x_1$ and every word $u$, with $G_a$, $F_o$ and $E_u$ as in
Definition~\ref{asm-member-bounds},
\[
 G_a+F_o-E_u\le x^{9c_t+2}b.
\]
\end{lemma}
\leanref{Formalism/SharpRestriction/Divisibility/HookProducts/GaugeFormBoundOfUniformSize.lean}{Formalism.MainNotions.DIV.HookProducts.Row4C.gauge_form_bound_of_uniform_size}
\leanref{Formalism/SharpRestriction/Divisibility/HookProducts/GaugeFormBoundOfUniformSize.lean}{Formalism.MainNotions.DIV.HookProducts.Row4C.negEe_le}
\begin{proof}
Let $C$ and $x_a$ be the constant and the scale of part (6), and take $x_0$ at least $x_a$, the $x_0$
of Lemma~\ref{asm-division-kernel-sizes}(3), $9$ and $2+|C|$. Fix a member at a scale $x\ge x_0$ and
a word $u$.

For every lift $j$, the first inequality of part (6) gives
$\sum_w\mathrm{sz}_w(E_{u,j}(w))\ge(\sum_{t\in T}a_t)\log(g_0g_1)-CB_0\ge-CB_0$, because every $a_t\ge
a>0$ and $\log(g_0g_1)\ge0$. Summing over the $m$ lifts gives $-E_u\le CmB_0$. By
Lemma~\ref{asm-division-kernel-sizes}(2) and (1), $CmB_0\le|C|\,\Pi^4b\le|C|x^{4c_t}b\le
|C|x^{9c_t+1}b$, since $x\ge9$ and $c_t\ge0$. By Lemma~\ref{asm-division-kernel-sizes}(3),
$G_a\le x^{4c_t+1}b\le x^{9c_t+1}b$ and $F_o\le x^{9c_t+1}b$. Hence
\[
 G_a+F_o-E_u\le(2+|C|)x^{9c_t+1}b\le x\cdot x^{9c_t+1}b=x^{9c_t+2}b.
\]
\end{proof}

\begin{lemma}\label{asm-division-kernel-height}
Let $\mathcal F$ be a family with exponent $\varepsilon$ that satisfies the size condition on $T_u$
(Definition~\ref{asm-designed}), with explicit heights $h_u$ and $c_t$ as in
Lemma~\ref{asm-division-kernel-sizes}. Let $C_H$, $e_b$, $e_c$ and $e_F$ be real numbers with
$\max(e_b,e_c,e_F,9c_t+1)+1\le C_H$. Suppose that at all large scales, at every member:
\begin{enumerate}
\item the slope bounds for $N$ hold (Definition~\ref{asm-member-bounds}(3));
\item every $N_u$ is finitely generated, spans its space and has positive rank $d_u$;
\item $\mu_{\max}(M_u)\le x^{e_b}b$ for every word $u$;
\item $G_a+F_o-E_u\le x^{e_c}b$ for every word $u$;
\item $\mu_{\max}(RR)+\overline{\log\kappa_R}\le x^{e_F}b$, with $\overline{\log\kappa_R}$ as in
  Lemma~\ref{asm-division-kernel-sharp}.
\end{enumerate}
Then at all large scales, at every member,
\[
 s+\mu_{\max}(RR)\le x^{C_H}b,\qquad s=\frac{\sum_ud_u\big(\mu(N_u)-h_u\big)}{\sum_ud_u}.
\]
\end{lemma}
\leanref{Formalism/SharpRestriction/Divisibility/SlopeExcess/HeightBoundOfComponentBounds.lean}{Formalism.MainNotions.DIV.Proofs.height_bound_of_component_bounds}
\leanref{Formalism/SharpRestriction/Divisibility/SlopeExcess/HeightBoundOfComponentBounds.lean}{Formalism.MainNotions.DIV.Proofs.hhRest_of_parts}
\leanref{Formalism/SharpRestriction/Divisibility/SlopeExcess/HeightBoundOfComponentBounds.lean}{Formalism.MainNotions.DIV.Proofs.height_bound_of_slope_fullness}
\leanref{Formalism/SharpRestriction/Divisibility/SlopeExcess/HeightBoundOfComponentBounds.lean}{Formalism.MainNotions.DIV.Proofs.height_bound_of_word_slopes}
\leanref{Formalism/SharpRestriction/Divisibility/LocalQuadraticParts.lean}{Formalism.MainNotions.LocalQuadraticParts.FrameNumberBound}
\begin{proof}
Put $E=\max(e_b,e_c,e_F,9c_t+1)$, so $E+1\le C_H$ and $E\ge4c_t+1$.

\emph{Step 1: for every real $C$, at all large scales and every word $u$,}
\[
 \mu_{\max}(M_u)+G_a+F_o-E_u+CmB_0-h_u+\mu_{\max}(RR)\le x^{C_H}b.
\]
Take $x$ at least $\max(9,|C|)$ and beyond the scales of (3), (4), (5) and of
Lemma~\ref{asm-division-kernel-explicit}. Since $x\ge1$, $x^eb\le x^Eb$ for every $e\le E$. Then (3)
bounds $\mu_{\max}(M_u)$ by $x^Eb$; (4) bounds $G_a+F_o-E_u$ by $x^Eb$; by
Lemma~\ref{asm-division-kernel-sizes}(2) and (1), $CmB_0\le|C|\,\Pi^4b\le x^{4c_t+1}b\le x^Eb$;
Lemma~\ref{asm-division-kernel-explicit} bounds $-(h_u+\overline{\log\kappa_R})$ by $x^Eb$; and (5)
bounds $\mu_{\max}(RR)+\overline{\log\kappa_R}$ by $x^Eb$. The sum of these five quantities is the
left side, so it is at most $5x^Eb\le x\cdot x^Eb=x^{E+1}b\le x^{C_H}b$.

\emph{Step 2: $\mu(N_u)-h_u+\mu_{\max}(RR)\le x^{C_H}b$ at all large scales.} Let $C$ be the constant
of the slope bounds for $N$. At all large scales $N_u$ is finitely generated and nonzero, so
$\mu(N_u)\le\mu_{\max}(N_u)$ (Definition~\ref{hs-slopes}), and the slope bounds give
$\mu_{\max}(N_u)\le\mu_{\max}(M_u)+G_a+F_o-E_u+CmB_0$. Step 1 gives the claim.

\emph{Step 3.} The ranks $d_u$ are positive, so $\sum_ud_u>0$. Multiplying the inequality of Step 2
by $d_u\ge0$ and summing over the words gives
$\sum_ud_u(\mu(N_u)-h_u)\le\big(\sum_ud_u\big)\big(x^{C_H}b-\mu_{\max}(RR)\big)$, so
$s\le x^{C_H}b-\mu_{\max}(RR)$.
\end{proof}

\subsection{Folding and Poincaré series}\label{subsection-asm-poincare-series}

\begin{situation}\label{asm-division-kernel-finite-situation}
Let $D$ be a definite-side datum (Definition~\ref{th-definite-datum}), and let an automorphic setting
attached to $D$ be given with the spectral data of a spectral condition for $D$: the rows of every
word, column representatives with the column spaces $\mathcal C^\sharp_u$, fundamental domains, a
level subgroup with a character, and reductions, as in Definition~\ref{is-unfolding-evaluation}(1).
Write $\mathcal H$ for the space of automorphic forms of the setting.
\end{situation}
\leanref{Formalism/LocalToGlobal/SiegelLatticeTheory/FlagStabilizerTensors/RootLatticeThresholds/NeatTies.lean}{Formalism.L2.SG.Pf.finX_of_neat}

\begin{lemma}\label{asm-division-kernel-finite-folding-domain}
Let $G$ be a second countable topological group with its Borel $\sigma$-algebra, and $S\subset G$ a
discrete subgroup. There is a measurable set $F\subset G$ that meets every orbit $Sy$ in exactly one point.
It is a fundamental domain of $S$ for every measure on $G$
(Lemma~\ref{asm-division-kernel-finite-series-fold-unfold}).
\end{lemma}
\leanref{Formalism/LocalToGlobal/SiegelLatticeTheory/HaarTransversals/UnitaryHaarProducts/ConcreteLevelAdelicSubset.lean}{Formalism.L2.SG.Pf.fund_nhd}
\leanref{Formalism/LocalToGlobal/SiegelLatticeTheory/HaarTransversals/UnitaryHaarProducts/ConcreteLevelAdelicSubset.lean}{Formalism.L2.SG.Pf.exists_strictFund}
\leanref{Formalism/LocalToGlobal/SiegelLatticeTheory/HaarTransversals/UnitaryHaarProducts/ConcreteLevelAdelicSubset.lean}{Formalism.L2.SG.Pf.isFund_of_strict}
\leanref{Formalism/LocalToGlobal/SiegelLatticeTheory/HaarTransversals/UnitaryHaarProducts/ConcreteLevelAdelicSubset.lean}{Formalism.L2.SG.Pf.exists_isFund}
\begin{proof}
As $S$ is discrete and $(a,b)\mapsto ab^{-1}$ is continuous, there is an open $V\ni1$ such that $ab^{-1}\in S$
with $a,b\in V$ forces $ab^{-1}=1$. As $G$ is second countable, countably many translates $Vx_0,Vx_1,\dots$
cover $G$, and $S$, being separable and discrete, is countable. Put
$F=\bigcup_n\big(Vx_n\setminus\bigcup_{m<n}SVx_m\big)$, a measurable set. For $y\in G$ let $n$ be the least
index with $Sy\cap Vx_n$ nonempty; then $Sy$ meets $F$ only in $Vx_n$, and if $sy$ and $y'=s'y$ both lie in
$Vx_n$, then $(sy)(s'y)^{-1}=ss'^{-1}$ is of the form $ab^{-1}$ with $a,b\in V$, so $s=s'$. Hence $F$ meets $Sy$
in exactly one point. Then every $y$ has $s$ with $sy\in F$, and $sF\cap s'F=\emptyset$ for $s\ne s'$.
\end{proof}

\begin{lemma}\label{asm-division-kernel-finite-folding-multiplicity}
In Situation~\ref{asm-division-kernel-finite-situation}, with the notation of
Definition~\ref{asm-division-kernel-finite-series-poincare} and of
Lemma~\ref{asm-division-kernel-finite-series-fold}, suppose that the decay of $W_u$
(Definition~\ref{is-wu-decay}) holds for every word $u$. Let $u$ be a word, $T$ a Poincaré datum of $u$,
$\Omega\subset G_{\mathbf A}$ compact and $F_u$ a fundamental domain of $\Gamma_u$ for $\mu$. Then the
multiplicity integral $\int_{F_u}m(y)|\Phi_T(y)|\,d\mu(y)$, with
$m(y)=\sup_{x}\#\{\gamma\in\Gamma:y^{-1}\gamma x\in\Omega\}$, is finite.
\end{lemma}
\leanref{Formalism/LocalToGlobal/SiegelLatticeTheory/HaarTransversals/UnitaryHaarProducts/InHHCut.lean}{Formalism.L2.SG.PoincareInHH.Cut.IUnfold}
\leanref{Formalism/LocalToGlobal/SiegelLatticeTheory/HaarTransversals/UnitaryHaarProducts/InHHCut.lean}{Formalism.L2.SG.PoincareInHH.Cut.MultPoly}
\leanref{Formalism/LocalToGlobal/SiegelLatticeTheory/ClosedVolumeStatements/WuDefiniteHolds.lean}{Formalism.L2.SG.Closed.MultPoly_holds}
\leanref{Formalism/LocalToGlobal/SiegelLatticeTheory/HaarTransversals/UnitaryHaarProducts/RightInvOfLattice.lean}{Formalism.L2.SG.Pf.iUnfold_of_NO}
\leanref{Formalism/LocalToGlobal/SiegelLatticeTheory/ClosedLocalOrbits/GaugeLatticeShifts/SlabNDensePf.lean}{Formalism.L2.SG.Pf.SlabNDensePf.nilradical_strong_approximation}
\begin{proof}
\openstep{The polynomial bound for the multiplicity in the archimedean variable, the finiteness of the
Levi quotients, and the unfolding of the multiplicity integral by strong approximation in the nilradical,
are not yet written out.}
% Lean: IUnfold X u hK (MultPoly X -> WuDecay X -> PhWellDef X u hK -> finiteness of multiplicityLIntegral),
% from iUnfoldAll_of hFL noCocompactC_holds nilradical_strong_approximation = iUnfold_of_NO X u hK (hFL D X u)
% (hNO D X u) (nDenseAt_of X u hND) = iUnfold_holds_of' X u hK (sCompact_of hFL hNO hND) psiUnit_holds
% ColBound_holds (sinfSection_of hNO); hFL = FinLeviPf.finLeviC_holds; hNO = noCocompactC_holds (Lemma
% asm-division-kernel-finite-volumes-cocompact). MultPoly X is Closed.MultPoly_holds (multPoly_holds_of''
% with OmCover_holds). Sizes (tools/lean_size.py): MultPoly_holds 6,731 lines; iUnfold_of_NO 9,670;
% finLeviC_holds 546; nilradical_strong_approximation 2,706; the union of these with noCocompactC_holds
% 18,120.
\end{proof}

\begin{lemma}\label{asm-division-kernel-finite-folding}
In Situation~\ref{asm-division-kernel-finite-situation}, the level subgroup of the setting is
normalized (Definition~\ref{is-unfolding-evaluation}(3)). Suppose that the decay of $W_u$
(Definition~\ref{is-wu-decay}) holds for every word $u$. Then the Poincaré series of every word $u$
fold: the rational stabilizer of the Bessel data of $u$ has a fundamental domain, and there are a
compact set
$\Omega_0$ containing $1$ and a constant $a\ge0$ such that, for every Poincaré datum $T$ of $u$, the
Poincaré function of $T$ has folding data with constant $a$ on every compact set
$\Omega\supset\Omega_0$, and its multiplicity integral over every compact set and every fundamental
domain is finite.
\end{lemma}
\leanref{Formalism/LocalToGlobal/SiegelLatticeTheory/ClosedVolumeStatements/LevelSubgroupProperties.lean}{Formalism.L2.SG.Closed.level_subgroup_properties}
\leanref{Formalism/LocalToGlobal/SiegelLatticeTheory/ClosedVolumeStatements/UnipotentInLevelSubgroup.lean}{Formalism.L2.SG.Closed.FoldOKC_holds}
\leanref{Formalism/LocalToGlobal/SiegelLatticeTheory/HaarTransversals/UnitaryHaarProducts/FinLevi.lean}{Formalism.L2.SG.Pf.FinLeviPf.finLeviC_holds}
\leanref{Formalism/LocalToGlobal/SiegelLatticeTheory/HaarTransversals/WuDecay.lean}{Formalism.L2.SG.WuDecay.FoldOK}
\begin{proof}
Use the notation of Definitions~\ref{asm-division-kernel-finite-series-poincare},
\ref{asm-division-kernel-finite-series-function-bessel}
and~\ref{asm-division-kernel-finite-series-function-inputs-level}, and let $K_g$ ($g\ge1$) be as in
Lemma~\ref{asm-division-kernel-finite-series-function-group-level}, with $g=g_0g_1$ below.

\emph{The level subgroup is normalized:} $K_g\subset K_2$ and $\chi_2=1$ on $K_g$, $K_2\subset K_{\mathrm{set}}$,
and $0<\mu_{\mathrm f}(K_2)<\infty$. (This part does not use
Lemma~\ref{asm-division-kernel-finite-series-function}.) Let $x\in K_g$, so $x\in K_{\mathrm{set}}$ and
$(x-1)\hat L\subset g\hat L$. Then $(x-1)e^0_a$, $(x-1)z^0$ and $(x-1)b_i$ lie in $g\hat L$, so $x$ satisfies the
conditions defining $K_2$ (Definition~\ref{asm-division-kernel-finite-neat-level}). For each vector $w$ of the
level basis, $xw=w+(x-1)w$ with $(x-1)w\in g\hat L$, so, writing $c\equiv0$ for $c\in g\hat{\mathcal O}$,
$h(f^0_a,xe^{0+}_a)\equiv h(f^0_a,e^{0+}_a)\equiv0$ and $h(b'_i,xb_{i+1})\equiv h(b'_i,b_{i+1})\equiv0$, by the
congruences of the level basis ($a\ne a+1$ and $i\ne i+1$). So the argument of $\psi_{E,\mathrm f}$ in
$\chi_2(x)$ is $g^{-1}$ times an element of $g\hat{\mathcal O}$, that is, lies in $\hat{\mathcal O}$, and
$\chi_2(x)=1$. By definition $K_2\subset K_{\mathrm{set}}$. Finally $K_g$ is open and contains $1$, so
$\mu_{\mathrm f}(K_2)\ge\mu_{\mathrm f}(K_g)>0$, and $K_{\mathrm{set}}$ is compact
(Lemma~\ref{asm-division-kernel-finite-series-function-group-level}(1)), so
$\mu_{\mathrm f}(K_2)\le\mu_{\mathrm f}(K_{\mathrm{set}})<\infty$.

\emph{The fundamental domain.} By Lemma~\ref{asm-division-kernel-finite-series-function-group}(1),
$G_{\mathbf A}$ is second countable with its Borel $\sigma$-algebra and $\Gamma$ is discrete, so the subgroup
$\Gamma_u$ is discrete, and Lemma~\ref{asm-division-kernel-finite-folding-domain} gives a fundamental domain
of $\Gamma_u$.

\emph{The folding data.} By the decay of $W_u$ there are a compact set $\Omega_i\ni1$ in $G_\infty$, a
constant $a_0\ge0$ and a measurable $\alpha$ on $G_\infty$, vanishing off $\Omega_i$, with $|\alpha|\le a_0$ and
$\int_{G_\infty}\alpha(z)W_u(xz)\,dz=W_u(x)$ for every $x$, the integrand being integrable. As in Step 6 of the
proof of Lemma~\ref{asm-division-kernel-finite-series-function}, with $\alpha$ in place of $\varphi$,
$\Phi_T(y)=\int_{G_\infty}\alpha(z)\Phi_T(y(z,1))\,dz$, so
$|\Phi_T(y)|\le a_0\int_{\Omega_i}|\Phi_T(y(z,1))|\,dz$. The group $C=K_g$ is open, hence closed, in the compact
group $K_{\mathrm{set}}$, so it is a compact neighbourhood of $1$, with $C\subset K_2$ and
$0<\mu_{\mathrm f}(C)<\infty$. For $\kappa\in C$, $|\Phi_T(y(z,\kappa))|=|\chi_2(\kappa)||\Phi_T(y(z,1))|
=|\Phi_T(y(z,1))|$ by Lemma~\ref{asm-division-kernel-finite-series-function}(2). With
$\Omega_0=\Omega_i\times C$, compact and containing $1$, Fubini's theorem gives
$\int_{\Omega_0}|\Phi_T(y\omega)|\,d\mu(\omega)=\mu_{\mathrm f}(C)\int_{\Omega_i}|\Phi_T(y(z,1))|\,dz$. So, with
$a=a_0/\mu_{\mathrm f}(C)$, which does not depend on $T$,
\[
 |\Phi_T(y)|\le a\int_{\Omega_0}|\Phi_T(y\omega)|\,d\mu(\omega)\le a\int_{\Omega}|\Phi_T(y\omega)|\,d\mu(\omega)
\]
for every compact $\Omega\supset\Omega_0$. As $\Phi_T$ is measurable and invariant under $\Gamma_u$ on the left
(Lemma~\ref{asm-division-kernel-finite-series-function}), $\Phi_T$ has folding data with constant $a$ on every
such $\Omega$ (Lemma~\ref{asm-division-kernel-finite-series-fold}).

\emph{The multiplicity integrals} are finite by Lemma~\ref{asm-division-kernel-finite-folding-multiplicity}.
% Lean: level_subgroup_properties = levelSubgroupNormalized_of_compact X AutKCompact_holds (cong_Kg,
% Afun_Kg with reference_frame_congruences, psi_Z, topf_open, maximal_compact_isCompact). FoldOK is
% foldOK_of_steps = foldOK_of (phProps_of_steps: Lemma -series-function) (iFinite_of: sratFund_holds =
% exists_fund_srat (Lemma -folding-domain), with IUnfold and MultPoly: Lemma -folding-multiplicity)
% (phReprod_of: WuReprod, ph_conv, exists_compact_nhds_level, ph_right_K, lintegral_prod). The folding
% is FoldOKC_holds hFL = foldOKC_holds_of with KsetSplit_holds, unipotent_in_level_subgroup,
% RatLevi_holds, psiProd_all and iUnfoldAll_of hFL noCocompactC_holds
% SlabNDensePf.nilradical_strong_approximation, with hFL = FinLeviPf.finLeviC_holds. The book takes
% C = K_g for the compact neighbourhood that exists_compact_nhds_level supplies. What remains open is in
% Lemmas -folding-multiplicity and -series-function-group-level.
\end{proof}

\begin{definition}\label{asm-division-kernel-finite-neat-level}
Let $E$ be a CM field with totally real subfield $F$, let $\mathbf A_{E,\mathrm f}=E\otimes_F\mathbf
A_{F,\mathrm f}$, with the involution $\sigma$ induced by that of $E$ and the subring
$\hat{\mathcal O}$ of integral elements, and let $\hat V=\mathbf A_{E,\mathrm f}^n$ with a Hermitian
form $h$, linear in the second variable, with $h(y,x)=\sigma(h(x,y))$. Let $\hat L\subset\hat V$ be an
$\hat{\mathcal O}$-submodule on which $h$ takes values in $\hat{\mathcal O}$, and let $g\ge1$ and
$k,r_S\ge0$ be integers. For $t\in\mathbf A_{E,\mathrm f}$ write $t\equiv0$ if $t\in g\hat{\mathcal
O}$, and for $S\subset\hat V$ write $\langle S\rangle$ for the sum of the $\hat{\mathcal O}$-span of
$S$ and $g\hat L$. A \emph{level basis} for $(\hat L,g)$ consists of vectors $e^0_a,f^0_a$
($1\le a\le r_S$), $z^0$ and $b_i,b'_i$ ($1\le i\le k$) of $\hat L$ and an element
$\epsilon\in\hat{\mathcal O}$ with $\epsilon\sigma(\epsilon)=0$, $\epsilon b_i=b_i$ and
$\sigma(\epsilon)b'_i=b'_i$, such that $\hat L=\langle e^0_a,f^0_a,z^0,b_i,b'_i\rangle$ and, for all
$a,c$, all $i,j$ and $x\in\{e^0_a,f^0_a\}$,
\[
 h(e^0_a,e^0_c)\equiv h(f^0_a,f^0_c)\equiv h(e^0_a,f^0_c)-\delta_{ac}\equiv h(z^0,z^0)+1\equiv0,
\]
\[
 h(x,z^0)\equiv h(x,b_i)\equiv h(x,b'_i)\equiv h(z^0,b_i)\equiv h(z^0,b'_i)
 \equiv h(b'_i,b_j)-\delta_{ij}\epsilon\equiv0.
\]
Its \emph{level subgroup} is the set $K_2$ of the $\mathbf A_{E,\mathrm f}$-linear isometries
$\kappa$ of $\hat V$ with $\kappa\hat L=\hat L$ such that, for all $a$ and $i$,
\[
 (\kappa-1)e^0_a\in\langle e^0_c:c<a\rangle,\qquad (\kappa-1)z^0\in\langle e^0_c:c\le r_S\rangle,
 \qquad (\kappa-1)b_i\in\langle e^0_c,b_j:c\le r_S,\ j<i\rangle.
\]
\end{definition}
\leanref{Formalism/LocalToGlobal/SiegelLatticeTheory/HaarTransversals/HidaSetting.lean}{Formalism.L2.SG.HidaSetting.HidaSetting}
\leanref{Formalism/LocalToGlobal/SiegelLatticeTheory/HaarTransversals/HidaSetting.lean}{Formalism.L2.SG.HidaSetting.K2set}
\leanref{Formalism/LocalToGlobal/SiegelLatticeTheory/HaarTransversals/HidaSetting.lean}{Formalism.L2.SG.HidaSetting.K2cong}
\leanref{Formalism/LocalToGlobal/SiegelLatticeTheory/HaarTransversals/HidaSetting.lean}{Formalism.L2.SG.HidaSetting.ZMod'}
\leanref{Formalism/LocalToGlobal/SiegelLatticeTheory/HaarTransversals/HidaSetting.lean}{Formalism.L2.SG.HidaSetting.mL}

\begin{lemma}\label{asm-division-kernel-finite-neat-projection}
Let a level basis for $(\hat L,g)$ be given (Definition~\ref{asm-division-kernel-finite-neat-level}),
and for $v\in\hat V$ put
\[
 \pi(v)=v-\sum_a\big(h(f^0_a,v)e^0_a+h(e^0_a,v)f^0_a\big)+h(z^0,v)z^0
 -\sum_i\big(h(b'_i,v)b_i+h(b_i,v)b'_i\big).
\]
Then $\pi(v)\in g\hat L$ for every $v\in\hat L$.
\end{lemma}
\leanref{Formalism/LocalToGlobal/SiegelLatticeTheory/HaarTransversals/UnitaryHaarProducts/InHHGamUnipKuProj.lean}{Formalism.L2.SG.PoincareInHH.Cut.GU.projModG_holds}
\leanref{Formalism/LocalToGlobal/SiegelLatticeTheory/HaarTransversals/UnitaryHaarProducts/InHHGamUnipKuProj.lean}{Formalism.L2.SG.PoincareInHH.Cut.GU.piG_reduce}
\leanref{Formalism/LocalToGlobal/SiegelLatticeTheory/HaarTransversals/UnitaryHaarProducts/InHHGamUnipKuDefs.lean}{Formalism.L2.SG.PoincareInHH.Cut.GU.piG}
\begin{proof}
The map $\pi$ is $\hat{\mathcal O}$-linear, so the $v$ with $\pi(v)\in g\hat L$ form an
$\hat{\mathcal O}$-submodule. As $\hat L$ is spanned by the vectors of the level basis and $g\hat L$,
it suffices to treat these vectors and the elements of $g\hat L$. If $t\equiv0$ and $x\in\hat L$, then
$tx\in g\hat L$. Let $b\in\hat L$, and suppose that
\[
 d=b-\Big(\sum_a(\alpha_ae^0_a+\beta_af^0_a)-\gamma z^0+\sum_i(\mu_ib_i+\nu_ib'_i)\Big)\in g\hat L
\]
with $\alpha_a-h(f^0_a,b)$, $\beta_a-h(e^0_a,b)$, $\gamma-h(z^0,b)$, $\mu_i-h(b'_i,b)$ and
$\nu_i-h(b_i,b)$ all $\equiv0$. Then $\pi(b)$ is $d$ plus these five differences times $e^0_a$,
$f^0_a$, $-z^0$, $b_i$ and $b'_i$, summed, so $\pi(b)\in g\hat L$. It remains to choose the
coefficients. For $b\in g\hat L$ take them all $0$: $h(x,b)\equiv0$ for $x\in\hat L$, as $h$ is
integral on $\hat L$. For $b=e^0_c$ take $\alpha_a=\delta_{ac}$ and the others $0$; for $b=f^0_c$,
$\beta_a=\delta_{ac}$; for $b=z^0$, $\gamma=-1$; for $b=b_c$, $\mu_i=\delta_{ic}\epsilon$, as
$\epsilon b_c=b_c$; and for $b=b'_c$, $\nu_i=\delta_{ic}\sigma(\epsilon)$. In each case $d=0$, and the
five congruences are those of the level basis, read with $h(y,x)=\sigma(h(x,y))$, together with
$h(b_i,b_j)=h(\epsilon b_i,\epsilon b_j)=\sigma(\epsilon)\epsilon\,h(b_i,b_j)=0$ and, in the same way,
$h(b'_i,b'_j)=0$.
\end{proof}

\begin{lemma}\label{asm-division-kernel-finite-neat-unipotent}
Let a level basis for $(\hat L,g)$ be given (Definition~\ref{asm-division-kernel-finite-neat-level}),
with level subgroup $K_2$, and let $\kappa$ be an element of $K_2$ with $\kappa^{-1}\in K_2$. Let
$Y\subset\hat L$ be an $\hat{\mathcal O}$-submodule with $\kappa Y\subset Y$ such that
$h(e^0_a,v)\equiv h(f^0_a,v)\equiv h(z^0,v)\equiv0$ for every $v\in Y$ and every $a$. Then
$(\kappa-1)^kv\in g\hat L$ for every $v\in Y$.
\end{lemma}
\leanref{Formalism/LocalToGlobal/SiegelLatticeTheory/HaarTransversals/UnitaryHaarProducts/InHHGamUnipKuCore.lean}{Formalism.L2.SG.PoincareInHH.Cut.GU.coreOfProj}
\leanref{Formalism/LocalToGlobal/SiegelLatticeTheory/HaarTransversals/UnitaryHaarProducts/InHHGamUnipKuDefs.lean}{Formalism.L2.SG.PoincareInHH.Cut.GU.UnipCore}
\leanref{Formalism/LocalToGlobal/SiegelLatticeTheory/HaarTransversals/UnitaryHaarProducts/InHHGamUnipKuCore.lean}{Formalism.L2.SG.PoincareInHH.Cut.GU.flag}
\leanref{Formalism/LocalToGlobal/SiegelLatticeTheory/HaarTransversals/UnitaryHaarProducts/InHHGamUnipKuCore.lean}{Formalism.L2.SG.PoincareInHH.Cut.GU.decomp}
\leanref{Formalism/LocalToGlobal/SiegelLatticeTheory/HaarTransversals/UnitaryHaarProducts/InHHGamUnipKuCore.lean}{Formalism.L2.SG.PoincareInHH.Cut.GU.dual_zero}
\leanref{Formalism/LocalToGlobal/SiegelLatticeTheory/HaarTransversals/UnitaryHaarProducts/LevelSubgroupCharacterUnitary.lean}{Formalism.L2.SG.Pf.unitaryCongruence.b_step}
\begin{proof}
\emph{Step 1: the flag.} For $0\le m\le k$ put $G_m=\langle e^0_c,b_j:c\le r_S,\ j\le m\rangle$. Let
$\lambda\in K_2$. Then $(\lambda-1)G_m\subset G_{m-1}$ for $1\le m\le k$: the map $\lambda-1$ is
$\hat{\mathcal O}$-linear and maps $g\hat L$ into $g\hat L$, as $\lambda\hat L=\hat L$, and by the
definition of $K_2$ it maps each $e^0_c$ into $G_0$ and each $b_j$ with $j\le m$ into $G_{j-1}\subset
G_{m-1}$. Hence $(\lambda-1)^mG_m\subset G_0$. If $x\in\hat L$ and $h(x,e^0_c)\equiv0$ for every $c$,
then $h(x,s)\equiv0$ for every $s\in G_0$, as $h$ is integral on $\hat L$.

\emph{Step 2: two identities.} As $\kappa$ is an isometry, $h(x,(\kappa-1)y)=h((\kappa^{-1}-1)x,y)$,
and so $h(x,(\kappa-1)^ky)=h((\kappa^{-1}-1)^kx,y)$, for all $x,y\in\hat V$. If $w\in\hat V$ satisfies
$\sigma(\epsilon)w=w$, then $h(b'_i,w)=0$: since $\sigma(\epsilon)b'_i=b'_i$,
$h(b'_i,w)=h(\sigma(\epsilon)b'_i,w)=\epsilon\,h(b'_i,w)$, and
$h(b'_i,w)=h(b'_i,\sigma(\epsilon)w)=\sigma(\epsilon)h(b'_i,w)$, so
$h(b'_i,w)=\epsilon\sigma(\epsilon)h(b'_i,w)=0$.

\emph{Step 3: a decomposition.} If $v\in\hat L$ and $h(e^0_a,v)\equiv h(f^0_a,v)\equiv h(z^0,v)\equiv0$
for every $a$, then
\[
 v-\sum_ih(b'_i,v)b_i-\sum_ih(b_i,v)b'_i\in g\hat L.
\]
Indeed this vector is $\pi(v)+\sum_a\big(h(f^0_a,v)e^0_a+h(e^0_a,v)f^0_a\big)-h(z^0,v)z^0$, which lies
in $g\hat L$ by Lemma~\ref{asm-division-kernel-finite-neat-projection} and the congruences.

\emph{Step 4.} Let $v\in Y$ and $z=(\kappa-1)^kv$. As $\kappa Y\subset Y$, $z\in Y$; so $z\in\hat L$
and $z$ satisfies the congruences of $Y$. First, $h(b_i,z)\equiv0$ for every $i$. By Step 2,
$h(b_i,z)=h((\kappa^{-1}-1)^kb_i,v)=\sigma\big(h(v,(\kappa^{-1}-1)^kb_i)\big)$. As $b_i\in G_k$ and
$\kappa^{-1}\in K_2$, Step 1 gives $(\kappa^{-1}-1)^kb_i\in G_0$; and $h(v,e^0_c)=\sigma(h(e^0_c,v))\equiv0$,
so $h(v,(\kappa^{-1}-1)^kb_i)\equiv0$ by Step 1.

Second, $h(b'_i,z)\equiv0$ for every $i$. Write $v=v_1+v_2+v_3$ with
$v_1=\sum_jh(b'_j,v)b_j$, $v_2=\sum_jh(b_j,v)b'_j$ and $v_3\in g\hat L$, by Step 3. The coefficients of
$v_1$ lie in $\hat{\mathcal O}$, so $v_1\in G_k$ and $(\kappa-1)^kv_1\in G_0$ by Step 1; as
$h(b'_i,e^0_c)=\sigma(h(e^0_c,b'_i))\equiv0$, $h(b'_i,(\kappa-1)^kv_1)\equiv0$. Next
$\sigma(\epsilon)v_2=v_2$, and $\kappa$ is $\mathbf A_{E,\mathrm f}$-linear, so
$w=(\kappa-1)^kv_2$ satisfies $\sigma(\epsilon)w=w$ and $h(b'_i,w)=0$ by Step 2. Finally
$(\kappa-1)^kv_3\in g\hat L$, as $\kappa\hat L=\hat L$, so $h(b'_i,(\kappa-1)^kv_3)\equiv0$.

By Step 3 applied to $z$, $z-\sum_ih(b'_i,z)b_i-\sum_ih(b_i,z)b'_i\in g\hat L$, and the two sums lie in
$g\hat L$ by the two congruences just shown. So $z\in g\hat L$.
\end{proof}

\begin{lemma}\label{asm-division-kernel-finite-neat-orders}
Let $E$ be a number field, $W$ a finite-dimensional $E$-vector space and $M\subset W$ a finitely
generated $\mathcal O_E$-submodule that spans $W$. Let $\beta$ be an $E$-linear automorphism of $W$ of
finite order with $\beta M=M$. Let $g_1\ne g_2$ be prime numbers, and for $i=1,2$ let $\mathfrak p_i$ be
a proper ideal of $\mathcal O_E$ with $g_i\in\mathfrak p_i$ such that $(\beta-1)^{r_i}M\subset\mathfrak
p_iM$ for some $r_i\ge0$. Then $\beta=1$.
\end{lemma}
\leanref{Formalism/LocalToGlobal/SiegelLatticeTheory/SiegelMassRatios/LemmaN.lean}{Formalism.L2.SG.LemmaN.lemmaNcGlob_of_exists}
\leanref{Formalism/LocalToGlobal/Columns/DecorationStabilizerTrivialUnipotent.lean}{Formalism.L2.Col.UnipMod.exists_pow_eq_one}
\leanref{Formalism/LocalToGlobal/Columns/DecorationStabilizerTrivialUnipotent.lean}{Formalism.L2.Col.UnipMod.exists_pow_mem_princLev}
\leanref{Formalism/LocalToGlobal/Columns/DecorationStabilizerTrivialUnipotent.lean}{Formalism.L2.Col.princLev_torsion}
\leanref{Formalism/LocalToGlobal/SiegelLatticeTheory/SiegelMassRatios/LemmaN.lean}{Formalism.L2.SG.LemmaN.eq_one_of_dvd_two_primes}
\begin{proof}
Fix $i$, and write $p=g_i$, $\mathfrak p=\mathfrak p_i$ and $r=r_i$. We show that the order of $\beta$
divides a power of $p$. Then the order divides a power of $g_1$ and a power of $g_2$, so it is $1$.

\emph{Step 1: $(\beta^{p^r}-1)M\subset\mathfrak pM$.} Put $X=\beta-1$, which maps $M$ into $M$. As $p$
divides the binomial coefficients $\binom{p^r}{j}$ for $0<j<p^r$,
$\beta^{p^r}=(1+X)^{p^r}=1+X^{p^r}+pS$ with $S$ a polynomial in $X$ with integer coefficients. As
$r\le p^r$, $X^{p^r}M=X^r(X^{p^r-r}M)\subset X^rM\subset\mathfrak pM$, and $pSM\subset pM\subset\mathfrak
pM$.

\emph{Step 2.} Let $\gamma$ be an automorphism of $W$ with $\gamma M=M$, $(\gamma-1)M\subset\mathfrak
pM$ and $\gamma^m=1$ for some $m$ prime to $p$. Then $\gamma=1$. Put $Y=\gamma-1$, and choose an integer
$a$ with $am\equiv1\bmod p$, so that $am-1\in\mathfrak p$. We show by induction that
$YM\subset\mathfrak p^{s}M$ for every $s\ge1$. Suppose that $YM\subset\mathfrak p^sM$. The operator
$(1+Y)^m-1-mY$ is $Y^2$ times a polynomial in $Y$ with integer coefficients, so it maps $M$ into
$\mathfrak p^{2s}M\subset\mathfrak p^{s+1}M$. As $(1+Y)^m=1$, this says that $mYM\subset\mathfrak
p^{s+1}M$. Then $Y=a(mY)-(am-1)Y$ maps $M$ into $\mathfrak p^{s+1}M$. By Krull's intersection theorem,
$\bigcap_s\mathfrak p^sM=0$, as $M$ is finitely generated and torsion-free over the Noetherian domain
$\mathcal O_E$ and $\mathfrak p$ is proper. So $YM=0$, and $Y=0$ as $M$ spans $W$.

\emph{Step 3.} Write the order of $\beta$ as $p^em$ with $p\nmid m$, and put
$\gamma=\beta^{p^{r+e}}=(\beta^{p^r})^{p^e}$. The automorphisms $\lambda$ with $\lambda M=M$ and
$(\lambda-1)M\subset\mathfrak pM$ are closed under products, so by Step 1 $\gamma$ is one of them;
and $\gamma^m=1$. By Step 2, $\gamma=1$, so the order of $\beta$ divides $p^{r+e}$.
\end{proof}

\begin{lemma}\label{asm-division-kernel-finite-neat}
In Situation~\ref{asm-division-kernel-finite-situation}, let $u$ be a word, with level subgroup
$K_u$. Then $K_u$ is neat: the classes of the column representatives of $u$ have trivial stabilizers
in $K_u$. Moreover, if $K_u$ is neat, then the map $\sharp$ from the space of Poincaré data of $u$ to
$\mathcal C^\sharp_u$, evaluation at the column representatives, is bijective.
\end{lemma}
\leanref{Formalism/LocalToGlobal/SiegelLatticeTheory/ClosedLocalOrbits/GaugeLatticeShifts/SlabFinal.lean}{Formalism.L2.SG.Pf.neatKu_full}
\leanref{Formalism/LocalToGlobal/SiegelLatticeTheory/FlagStabilizerTensors/RootLatticeThresholds/NeatTies.lean}{Formalism.L2.SG.Pf.sharp_bij_of_neat}
\leanref{Formalism/LocalToGlobal/SiegelLatticeTheory/HaarTransversals/UnitaryHaarProducts/InHHNeat.lean}{Formalism.L2.SG.PoincareInHH.Cut.NeatKu}
\leanref{Formalism/LocalToGlobal/SiegelLatticeTheory/SiegelMassRatios/ColsDim.lean}{Formalism.L2.SG.ColsDim.columns_decomposition}
\leanref{Formalism/LocalToGlobal/SiegelLatticeTheory/HaarTransversals/UnitaryHaarProducts/InHHFinOrd.lean}{Formalism.L2.SG.PoincareInHH.Cut.gamFinKu_of}
\leanref{Formalism/LocalToGlobal/SiegelLatticeTheory/HaarTransversals/UnitaryHaarProducts/InHHFinOrd.lean}{Formalism.L2.SG.PoincareInHH.Cut.delta_finOrd}
\leanref{Formalism/LocalToGlobal/SiegelLatticeTheory/HaarTransversals/UnitaryHaarProducts/InHHGamUnipKu.lean}{Formalism.L2.SG.PoincareInHH.Cut.gamUnipKu_holds}
\leanref{Formalism/LocalToGlobal/SiegelLatticeTheory/HaarTransversals/UnitaryHaarProducts/InHHGamUnipKuGlob.lean}{Formalism.L2.SG.PoincareInHH.Cut.GU.globStep_holds}
\begin{proof}
\emph{Notation.} The space of the setting is $V=Y_0$ of $D$ over the CM field $E$ of its host, with its
form $h$ (Definition~\ref{th-spaces}(5)), which has signature $(k+r_S,r_S+1)$ at every embedding. The
setting gives hyperbolic pairs $e_a,f_a$ ($1\le a\le r_S$) in $V$, with $h(e_a,e_b)=h(f_a,f_b)=0$ and
$h(e_a,f_b)=\delta_{ab}$, and for the word $u$ a vector $z_u$ orthogonal to them with
$h(z_u,z_u)=-T_u$; let $U$ be the span of the $e_a$, $f_a$ and $z_u$, and $W=U^\perp$. The group
$\Gamma^H_u$ consists of the rational points of $U(V)$ whose archimedean and finite components fix every
vector of $U$ (Definition~\ref{asm-division-kernel-finite-series-function-bessel}(4)), $H_{\mathrm f}$ is
the group of finite adelic points fixing every vector of $U$, and $\iota\colon\Gamma^H_u\to H_{\mathrm f}$
is the finite component. The level subgroup $K_2$ of the setting
stabilizes the completion $\hat L$ of its lattice in $\hat V=V\otimes_F\mathbf A_{F,\mathrm f}$, on which
$h$ is integral. In the coordinates of a basis of $V$, the setting gives a level basis
$e^0_a,f^0_a,z^0,b_i,b'_i,\epsilon$ for $(\hat L,g)$, where $g=g_0g_1$ is the product of the two level
primes $g_0\ne g_1$ of the host, and $K_2$ is its level subgroup
(Definition~\ref{asm-division-kernel-finite-neat-level}), a group. The group $K_u$ is the group of
$y\in H_{\mathrm f}$ with $x_u^{-1}yx_u\in K_2$, for the base point $x_u$ of $u$, which fixes every
$e_a$ and $f_a$; and for each of the vectors $e^0_a-e_a$, $f^0_a-f_a$ and $x_u^{-1}z_u-z^0$, its
product with some integer prime to $g$ lies in $g\hat L$. Neatness says:
if $y\in H_{\mathrm f}$ and $\gamma\in\Gamma^H_u$ satisfy $y^{-1}\iota(\gamma)y\in K_u$, then $\gamma=1$.
Fix such $y$ and $\gamma$.

\emph{Step 1: $\gamma$ has finite order.} Put $x=yx_u$. The vectors $v\in V$ with $x^{-1}v\in\hat L$
form a lattice $\Lambda$ of $V$, as for every finite adelic point. Let $\gamma_0$ be the rational
isometry of $V$ that $\gamma$ is; its finite component is $\iota(\gamma)$, and
$x^{-1}\iota(\gamma)x=x_u^{-1}\big(y^{-1}\iota(\gamma)y\big)x_u$ lies in $K_2$, so it preserves
$\hat L$, and $\gamma_0$ preserves $\Lambda$. As $\gamma_0$ fixes $U$ pointwise and is an isometry, it
preserves $W$. Since $h(z_u,z_u)\ne0$, every $v\in V$ is the sum of
$\sum_a\big(h(f_a,v)e_a+h(e_a,v)f_a\big)+h(z_u,z_u)^{-1}h(z_u,v)z_u\in U$ and a vector of $W$, so
$V=U+W$. The trace form $\operatorname{Tr}_{E/\mathbf Q}h$ is positive definite on $W$: for
$0\ne w\in W$, $w$ is orthogonal to the $e_a$ and to $z_u$ and is not in the span $X$ of the $e_a$ (a
vector of $X$ orthogonal to every $f_a$ is $0$), so $\tau(h(w,w))>0$ for every embedding $\tau$, by the
argument of Step 4 of the proof of Lemma~\ref{asm-division-kernel-finite-fibres} with $e_a$, $z_u$ in
place of $f_a$, $z'$. The intersection $\Lambda\cap W$ is a lattice of $W$ preserved by $\gamma_0$, and the
isometries of a lattice for a positive definite form on a finite-dimensional $\mathbf Q$-vector space form
a finite group. So $\gamma_0^m=1$ on $W$ for some $m\ge1$, and $\gamma_0^m=1$ on $U$; hence
$\gamma_0^m=1$. A rational point of $U(V)$ is determined by its finite component, so $\gamma^m=1$.

\emph{Step 2: $\gamma=1$.} Put $\kappa=x^{-1}\iota(\gamma)x$, which lies in $K_2$ by Step 1, as does
$\kappa^{-1}$. For $w\in V$ put $\hat w=x^{-1}w\in\hat V$; so $\Lambda$ is the set of $w$ with
$\hat w\in\hat L$, and $\widehat{\gamma_0w}=\kappa\hat w$. Let $Y$ be the set of $v\in\hat L$ with
$h(e_a,v)=h(f_a,v)=h(x_u^{-1}z_u,v)=0$ for every $a$, an $\hat{\mathcal O}$-submodule of $\hat L$.

(a) $\kappa Y\subset Y$. Put $\eta=y^{-1}\iota(\gamma)y\in H_{\mathrm f}$, so that
$\kappa^{-1}=x_u^{-1}\eta^{-1}x_u$. As $\eta$ fixes $U$ and $x_u$ fixes the $e_a$ and $f_a$,
$\kappa^{-1}$ fixes $e_a$, $f_a$ and $x_u^{-1}z_u$. As $\kappa$ is an isometry,
$h(e_a,\kappa v)=h(\kappa^{-1}e_a,v)=h(e_a,v)$, and likewise for $f_a$ and $x_u^{-1}z_u$; and
$\kappa\hat L=\hat L$.

(b) Every $v\in Y$ has $h(e^0_a,v)\equiv h(f^0_a,v)\equiv h(z^0,v)\equiv0$. Let $c$ be an integer prime
to $g$ with $c(e^0_a-e_a)\in g\hat L$. As $h$ is integral on $\hat L$,
$c\,h(e^0_a,v)=h(c(e^0_a-e_a),v)\equiv0$; as $h(e^0_a,v)\in\hat{\mathcal O}$ and $c$ is invertible modulo
$g$, $h(e^0_a,v)\equiv0$. The same argument applies to $f^0_a$, and to $z^0$ with $x_u^{-1}z_u$ in
place of $e_a$.

(c) By (a), (b) and Lemma~\ref{asm-division-kernel-finite-neat-unipotent},
$(\kappa-1)^kv\in g\hat L$ for every $v\in Y$.

(d) Let $w\in\Lambda\cap W$. Then $\hat w\in Y$: $\hat w\in\hat L$, and $x^{-1}=x_u^{-1}y^{-1}$ maps
$e_a$, $f_a$ and $z_u$ to $e_a$, $f_a$ and $x_u^{-1}z_u$, as $y$ fixes $U$; so
$h(e_a,\hat w)=h(x^{-1}e_a,x^{-1}w)=h(e_a,w)=0$, and likewise for $f_a$ and $x_u^{-1}z_u$. By (c),
$\widehat{(\gamma_0-1)^kw}=(\kappa-1)^k\hat w\in g\hat L$, so $g^{-1}(\gamma_0-1)^kw\in\Lambda$. It lies
in $W$, as $\gamma_0$ preserves $W$. So $(\gamma_0-1)^k(\Lambda\cap W)\subset g(\Lambda\cap W)$.

(e) Let $M=\Lambda\cap W$, a lattice of $W$, and $\beta$ the restriction of $\gamma_0$ to $W$. By Step 1,
$\beta M=M$ and $\beta^m=1$. Let $\mathfrak p_0$ and $\mathfrak p_1$ be prime ideals of $\mathcal O_E$
containing $g_0$ and $g_1$. They contain $g$, so $(\beta-1)^kM\subset\mathfrak p_iM$ by (d). As
$g_0\ne g_1$, Lemma~\ref{asm-division-kernel-finite-neat-orders} gives $\beta=1$.

(f) So $\gamma_0$ is the identity on $W$ and on $U$, hence on $V=U+W$. Its finite component
$\iota(\gamma)$ is $1$, and $\gamma=1$.
% Lean: neatKu_full = neatKu_of_gamUnipKu X u Y.hK (latOfLattice_holds X u) (gamUnipKu_holds X u Y.hK)
% (SlabFinal.lean); neatKu_of_gamUnipKu = GamUnipKu applied to gamFinKu_of (InHHFinOrd.lean). Step 1 is
% gamFinKu_of: latOfLattice_holds (PIW.AdelClassFin.mLat_of with lattice_intersection' and
% lattice_denominator) gives Lambda = latOfC; matEC gamma0 preserves it through Kset; delta_finOrd
% (decomp for V = U + W, posW_holds for the positivity on W, which is the argument of pos_perp with the
% trace, and Formalism.L2.Col.isOfFinOrder_of_isom, from HermitianLatticesTheory.autFinite) gives
% delta^m = 1; gam_pow_eq_one (fmat_eq_one_of_fix, gf_eq_one_of_snd) gives gamma^m = 1. Step 2 is
% gamUnipKu_holds = GU.gamUnipKu_of_projModG with GU.projModG_holds (Lemma
% asm-division-kernel-finite-neat-projection), GU.coreOfProj (Lemma -neat-unipotent) and
% GU.globStep_holds (InHHGamUnipKuGlob.lean): kappa^-1 in K2 is hK.2.2.1; Y is Yu (YP); (a) is
% yu_stable (conj_fix, xu_inv_fix); (b) is hermitian_e0_congruent, hermitian_f0_congruent and
% hermitian_z0_congruent (basisApprox, xu_z, Z_left, zmod_cancel); (d) is hatv_Y (P_fix, hatv = x^-1 in
% coordinates) with hApow, mL_exists and inv_g_cancel (hkey); (e) is qLat_isLattice, q_aut, q_pow_one
% and the unipotence hunip on qLat Lambda W bot, exists_prime_over and LemmaN.lemmaNcGlob_of_exists
% (Lemma -neat-orders) with D.host.g_prime and D.host.g_ne; (f) is q_fix, Cut.decomp,
% fmat_eq_one_of_fix and gf_eq_one_of_snd. The level basis is the fields e0, f0, z0, bw, bwd, ew,
% vectorsInLattice, gramIntegral, formsAtZero, weightVectorsIntegral, bw_w, bwd_w, ew_orth, bw_dual
% and latticeSpan of HidaSetting, with g_eq; xu_fix, basisApprox and xu_z relate it to e_a, f_a, z_u.
% Size: the closure of gamUnipKu_holds without latOfLattice_holds, about 2,100 lines in 18 files of
% Formalism/LocalToGlobal (text count); Step 1, with latOfLattice_holds, about 1,450 lines in 11 files.

Suppose that $K_u$ is neat, and let $\Gamma$ be the rational group of the setting acting on the left.
The column representatives $h_\varrho$ represent the double cosets $\Gamma\backslash G_f/K_u$, each
once, so $\varrho\mapsto\Gamma h_\varrho K_u$ is a bijection onto the double cosets. Since the
stabilizers are trivial, a Poincaré datum $T$ is determined by the functionals $\xi\mapsto
T(\xi)(1,h_\varrho)$ on $V_{\pi_u}$, and every family of functionals, one for each double coset,
comes from a Poincaré datum. The map $\sharp$ sends $T$ to the vectors of $V_{\pi_u}$ that represent
these functionals for the inner product, so it is bijective.
\end{proof}

\begin{definition}\label{asm-division-kernel-finite-series-poincare}
In Situation~\ref{asm-division-kernel-finite-situation}, let $G_{\mathbf A}=G_\infty\times G_{\mathrm f}$ be
the adelic group of the setting, with its Haar measure $\mu$, its group $\Gamma$ of rational points
and its fundamental domain $\mathcal F$, and let $\varphi$ be the archimedean matrix coefficient and
$e$ the function on $G_{\mathrm f}$ equal to $\overline{\chi_2(g)}/\mathrm{vol}(K_2)$ on $K_2$ and to $0$
elsewhere. Let $u$ be a word, with base point $x_u$, archimedean Bessel function $W_u$, the group
$N_{\mathrm f}$ of Lemma~\ref{asm-division-kernel-finite-volumes-cocompact} with its character
$\psi_{\mathrm f}$, and the group $H_{\mathrm f}$ of finite adelic points that fix the vectors $e_a$,
$f_a$ and $z_u$ of the setting, and let
$\Gamma_u\subset\Gamma$ be the rational stabilizer of the Bessel data of $u$. Let $T$ be a Poincaré
datum of $u$.
\begin{enumerate}
\item The \emph{Poincaré function} of $T$ is the function $\Phi_T$ on $G_{\mathbf A}$ with
  $\Phi_T(y)=\psi_{\mathrm f}(n)\chi_2(k)\,T(W_u(y_\infty))(1,h)$ if $y_{\mathrm f}=nhx_uk$ with
  $n\in N_{\mathrm f}$, $h\in H_{\mathrm f}$ and $k\in K_2$, for a choice of such a decomposition, and
  $\Phi_T(y)=0$ if there is none.
\item A \emph{set of representatives} is a set $D\subset\Gamma$ such that every $\gamma\in\Gamma$ is
  $s\delta$ for exactly one pair $(s,\delta)\in\Gamma_u\times D$. The \emph{Poincaré series} of $T$ is
  $P_{T,D}(x)=\sum_{\delta\in D}\Phi_T(\delta x)$ where it converges.
\end{enumerate}
\end{definition}
\leanref{Formalism/LocalToGlobal/SiegelLatticeTheory/HaarTransversals/PoincareInHH.lean}{Formalism.L2.SG.PoincareInHH.PoincareSetting.Ph}
\leanref{Formalism/LocalToGlobal/SiegelLatticeTheory/HaarTransversals/PoincareInHH.lean}{Formalism.L2.SG.PoincareInHH.PoincareSetting.Srat}
\leanref{Formalism/LocalToGlobal/SiegelLatticeTheory/HaarTransversals/PoincareInHH.lean}{Formalism.L2.SG.PoincareInHH.PoincareSetting.IsReps}
\leanref{Formalism/LocalToGlobal/SiegelLatticeTheory/HaarTransversals/PoincareInHH.lean}{Formalism.L2.SG.PoincareInHH.PoincareSetting.Pser}
\leanref{Formalism/LocalToGlobal/SiegelLatticeTheory/HaarTransversals/PoincareInHH.lean}{Formalism.L2.SG.PoincareInHH.PoincareSetting.InHH}

\begin{lemma}\label{asm-division-kernel-finite-series-function-group-level}
In Situation~\ref{asm-division-kernel-finite-situation}, let $G_{\mathrm f}$ be the group of finite adelic
points of the unitary group of the setting, $\hat L$ the adelic lattice of the setting and
$K_{\mathrm{set}}$ its stabilizer in $G_{\mathrm f}$, and for $g\ge1$ let $K_g$ be the group of
$x\in K_{\mathrm{set}}$ with $(x-1)\hat L\subset g\hat L$. The topology of $G_{\mathrm f}$ is the one in which
the $K_g$ are open and form a basis of neighbourhoods of $1$.
\begin{enumerate}
\item $K_{\mathrm{set}}$ is open and compact.
\item There is a countable set $C\subset G_{\mathrm f}$ with $G_{\mathrm f}=\bigcup_{c\in C}cK_{\mathrm{set}}$.
\item $\bigcap_{g\ge1}K_g=\{1\}$.
\item For every $x\in G_{\mathrm f}$, $[K_{\mathrm{set}}:K_{\mathrm{set}}\cap xK_{\mathrm{set}}x^{-1}]
  =[K_{\mathrm{set}}:K_{\mathrm{set}}\cap x^{-1}K_{\mathrm{set}}x]$.
\item If $G_{\mathrm f}$ carries a $\sigma$-finite Haar measure, then $G_{\mathrm f}$ is second countable.
\item The Haar measure of every indefinite unitary group $U(J_v)$ of the setting is right invariant.
\end{enumerate}
\end{lemma}
\leanref{Formalism/LocalToGlobal/SiegelLatticeTheory/HaarTransversals/UnitaryHaarProducts/ConcreteLevelAdelicSubset.lean}{Formalism.L2.SG.Pf.kset_open}
\leanref{Formalism/PeriodIntegrals/UnitaryGroupMassFormula/AdelicUnimodularity/ConstructionA.lean}{PI.WAU.maximal_compact_isCompact}
\leanref{Formalism/PeriodIntegrals/UnitaryGroupMassFormula/AdelicUnimodularity/ConstructionB.lean}{PI.WAU.compact_countable}
\leanref{Formalism/LocalToGlobal/SiegelLatticeTheory/HaarTransversals/UnitaryHaarProducts/ConcreteLevelAdelicSubset.lean}{Formalism.L2.SG.Pf.sepK}
\leanref{Formalism/PeriodIntegrals/AdIndexEq.lean}{PI.AdIndexEq.IndexEq}
\leanref{Formalism/PeriodIntegrals/OrbitalBoundsConclusion/AdIndexEq.lean}{PI.adelic_index_eq}
\leanref{Formalism/SharpRestriction/CuspSpanNormBounds/UnimodularSchurMeasures/RrsSecondCountableAbs.lean}{Formalism.MainNotions.LeviCoordinates.rr_secondCountable}
\leanref{Formalism/RealPlaceAnalysis/RealPlaceTaylor/ModularCharacterEqOneOnExpHermitian.lean}{Formalism.REAL.indefinite_unitary_unimodular}
\begin{proof}
\openstep{The compactness of the stabilizer of the adelic lattice and the countability of its cosets, the
separation of $1$ by the principal congruence subgroups, the equality of the two indices, the second
countability, and the unimodularity of the indefinite unitary groups are not yet written out.}
% Lean: (1) kset_open (proved in ConcreteLevelAdelicSubset.lean) and PI.WAU.maximal_compact_isCompact, 789
% lines; (2) PI.WAU.compact_countable, 560 lines; (3) sepK (trC_nondeg, af_iInter, nt_Lf_int), 305 lines;
% (4) PI.adelic_index_eq = indexEq_of_traceInt trace_integral', 4,017 lines; (5) rr_secondCountable
% (rrs_secondCountable_abs with the K_g), 90 lines; (6) indefinite_unitary_unimodular
% (modularCharacter_eq_one_alt), 735 lines (sizes from tools/lean_size.py, each on its own). The
% topology of Gf is the AutInst fields topf_open and topf_basis.
\end{proof}

\begin{lemma}\label{asm-division-kernel-finite-series-function-group}
In Situation~\ref{asm-division-kernel-finite-situation}, let $G_{\mathbf A}=G_\infty\times G_{\mathrm f}$,
$\mu=\mu_\infty\times\mu_{\mathrm f}$ and $\Gamma$ be as in
Definition~\ref{asm-division-kernel-finite-series-poincare}: $G_\infty=\prod_vU(J_v)$ is the product over
the infinite places $v$ of the indefinite unitary groups of the setting, with the product $\mu_\infty$ of
their Haar measures, and $G_{\mathrm f}$ is the group of finite adelic points of the unitary group of the
setting, with a $\sigma$-finite Haar measure $\mu_{\mathrm f}$. The $\sigma$-algebra of $G_{\mathbf A}$ is
the product of the Borel $\sigma$-algebras of $G_\infty$ and $G_{\mathrm f}$.
\begin{enumerate}
\item The groups $G_{\mathrm f}$ and $G_{\mathbf A}$ are locally compact, second countable and Hausdorff,
  the $\sigma$-algebra of $G_{\mathbf A}$ is its Borel $\sigma$-algebra, and $\Gamma$ is a discrete
  subgroup of $G_{\mathbf A}$.
\item Every $\sigma$-finite Haar measure on $G_{\mathrm f}$ is right invariant, and the Haar measure of
  every indefinite unitary group $U(J_v)$ is right invariant.
\item The group $\Gamma$ is countable.
\end{enumerate}
\end{lemma}
\leanref{Formalism/LocalToGlobal/SiegelLatticeTheory/HaarTransversals/ExplicitFormulaBoundHolds.lean}{Formalism.L2.SG.Count.FoldReady}
\leanref{Formalism/LocalToGlobal/SiegelLatticeTheory/HaarTransversals/UnitaryHaarProducts/ConcreteLevelAdelicSubset.lean}{Formalism.L2.SG.Pf.foldReady_holds}
\leanref{Formalism/LocalToGlobal/SiegelLatticeTheory/HaarTransversals/UnitaryHaarProducts/ConcreteLevelAdelicSubset.lean}{Formalism.L2.SG.Pf.gf_secondCountable}
\leanref{Formalism/LocalToGlobal/SiegelLatticeTheory/HaarTransversals/UnitaryHaarProducts/RowsProj.lean}{Formalism.L2.SG.Pf.GfUnimod}
\leanref{Formalism/LocalToGlobal/SiegelLatticeTheory/ClosedVolumeStatements/LevelSubgroupProperties.lean}{Formalism.L2.SG.Closed.GfUnimod_holds}
\leanref{Formalism/RealPlaceAnalysis/RealPlaceTaylor/ModularCharacterEqOneOnExpHermitian.lean}{Formalism.REAL.indefinite_unitary_unimodular}
\leanref{Formalism/SharpRestriction/CuspSpanNormBounds/UnimodularSchurMeasures/RrgCountableField.lean}{Formalism.MainNotions.LeviCoordinates.rr_GF_countable}
\begin{proof}
(3) An element of $\Gamma$ is determined by the matrix $\gamma\in GL_n(E)$ that defines it, as its
components at all places are images of $\gamma$; and $GL_n(E)$ is countable, as $E$ is.

We use the notation and the facts of Lemma~\ref{asm-division-kernel-finite-series-function-group-level}.

(1) \emph{The group $G_{\mathrm f}$.} By part (1) of that lemma, $K_{\mathrm{set}}$ is a compact neighbourhood
of $1$, so the topological group $G_{\mathrm f}$ is locally compact. Each $K_g$ is an open subgroup, hence
closed, and by part (3) their intersection is $\{1\}$; so for $x\ne1$ some open $K_g$ contains $1$ and not
$x$, and $G_{\mathrm f}$ is Hausdorff. By part (2) it is the union of countably many compact sets
$cK_{\mathrm{set}}$, so it is $\sigma$-compact; a Haar measure on it is then $\sigma$-finite, and $G_{\mathrm
f}$ is second countable by part (5).

\emph{The group $G_{\mathbf A}$.} Each $U(J_v)$ is a closed subgroup of $GL_n(\mathbf C)$, so it is locally
compact, second countable and Hausdorff, and so is the finite product $G_\infty$, and then $G_{\mathbf A}$.
For second countable spaces the Borel $\sigma$-algebra of a product is the product of the Borel
$\sigma$-algebras.

\emph{Discreteness.} Let $s$ be a finite set of generators of the lattice of the setting, which spans $V$,
and let $c\ne0$ be an integer such that $c$ times every coordinate (in the basis of the setting) of every
$v\in V\cap\hat L$ is an algebraic integer. Let $U$ be the set of $x\in G_{\mathbf A}$ with
$x_{\mathrm f}\in K_{\mathrm{set}}$ such that, at every infinite place $v$ and for every $l\in s$, every
coordinate of $c(x_v-1)l$ in the basis of the setting has absolute value less than $1$. It is open and
contains $1$. Let $\gamma\in\Gamma\cap U$, with matrix $M\in GL_n(E)$. For $l\in s$, $Ml\in\hat L$ as
$\gamma_{\mathrm f}\in K_{\mathrm{set}}$, so every coordinate $a$ of $c(Ml-l)$ is an algebraic integer of $E$;
its images under the embeddings $\tau_v$ are the coordinates of $c(\gamma_v-1)l$, of absolute value less than
$1$, and so are their complex conjugates. So $|N_{E/\mathbf Q}(a)|<1$, and $a=0$. Hence $M$ fixes $s$, which
spans $V$, so $M=1$ and $\gamma=1$. So $\Gamma\cap U=\{1\}$, and $\Gamma$ is discrete.

(2) Let $\mu_{\mathrm f}$ be a $\sigma$-finite Haar measure on $G_{\mathrm f}$, and $x\in G_{\mathrm f}$. Put
$K=K_{\mathrm{set}}$, which is compact and open, so $0<\mu_{\mathrm f}(K)<\infty$. The measure
$A\mapsto\mu_{\mathrm f}(Ax)$ is a left Haar measure, so it is $\Delta\mu_{\mathrm f}$ for a number
$\Delta>0$, and it suffices to show $\mu_{\mathrm f}(Kx)=\mu_{\mathrm f}(K)$, that is
$\mu_{\mathrm f}(x^{-1}Kx)=\mu_{\mathrm f}(K)$ by left invariance. With $K'=K\cap x^{-1}Kx$, $K$ is the union of
$[K:K']$ disjoint left cosets of $K'$, and $x^{-1}Kx$ of $[x^{-1}Kx:K']=[K:xKx^{-1}\cap K]$ of them, conjugating
by $x$. By part (4) these indices are equal, so $\mu_{\mathrm f}(x^{-1}Kx)=\mu_{\mathrm f}(K)$. The Haar
measures of the $U(J_v)$ are right invariant by part (6).
% Lean: (1) foldReady_holds I = foldReady_of I (PI.WAU.maximal_compact_isCompact I)
% (PI.WAU.compact_countable I), with gf_lc, finite_adeles_group_t2, gf_sc and gf_discrete, about 1,490
% lines (proof closure, text count); gf_secondCountable = gf_sc, about 940 of them. (2) at the finite
% places: GfUnimod_holds = PI.WAU.adUnimod_of kg_subgroup maximal_compact_isCompact compact_countable
% PI.adelic_index_eq, through PI.UnimodIdxWork.unimodIdx_main (the indices [K : K cap gKg^-1] and
% [gKg^-1 : K cap gKg^-1] agree for the principal congruence subgroups K, AdelicIndexEquality), about
% 3,080 lines; at the infinite places: indefinite_unitary_unimodular (modularCharacter_eq_one_alt), about
% 650 lines, used through the instance rr_muG_right. All of (1) and (2) together: about 4,400 lines.
% FoldReady is also used by Lemma asm-division-kernel-finite-folding (foldReadyX_holds).
% The proof follows foldReady_of (gf_lc, finite_adeles_group_t2 via t2Space_of_one_sep, gf_sigmaCompact,
% gf_sc, gf_discrete with nt_Lf_denom, psi_eq, nt_int_eq_zero and mat_eq_one_of_fix) and, for (2),
% adUnimod_of (unimodIdx_main with Ksub, kset_open). What remains open is in Lemma
% asm-division-kernel-finite-series-function-group-level.
\end{proof}

\begin{definition}\label{asm-division-kernel-finite-series-function-bessel}
In Situation~\ref{asm-division-kernel-finite-situation}, with the notation of
Definition~\ref{asm-division-kernel-finite-series-poincare}, let $V$ be the Hermitian space of the
setting over the CM field $E$, with totally real subfield $F$ and form $h$, linear in the second
variable, and let $r=r_S$. The setting gives vectors $e_1,\dots,e_r$ and $f_1,\dots,f_r$ of $V$ with
$h(e_a,e_b)=h(f_a,f_b)=0$ and $h(e_a,f_b)=\delta_{ab}$, and for each word $u$ a vector $z_u$ orthogonal
to them. Its lattice is $L_0=\sum_a\mathfrak a_ae_a+L_b+\sum_a\mathfrak a_a^{-1}f_a$, for fractional
ideals $\mathfrak a_a$ of $E$ and an $\mathcal O_E$-submodule $L_b$ orthogonal to every $e_a$ and $f_a$;
$\hat L$ is the $\hat{\mathcal O}$-span of $L_0$ in $\hat V=V\otimes_F\mathbf A_{F,\mathrm f}$, and
$K_{\mathrm{set}}$ is the group of $x\in G_{\mathrm f}$ with $x\hat L=\hat L$. The setting also gives: the
level subgroup $K_2\subset K_{\mathrm{set}}$ of Definition~\ref{asm-division-kernel-finite-neat-level}, a
group on which $\chi_2$ is a character with $|\chi_2|=1$; the locally constant function $e$; the base
points $x_u\in K_{\mathrm{set}}$; and an additive character $\psi_{E,\mathrm f}$ of
$\mathbf A_{E,\mathrm f}=E\otimes_F\mathbf A_{F,\mathrm f}$, trivial on the integral elements, with
$\psi_{E,\mathrm f}(x)\psi_{E,\infty}(x)=1$ for $x\in E$, where
$\psi_{E,\infty}(x)=\prod_v\exp\big(2\pi i(\tau_v(x)+\overline{\tau_v(x)})\big)$, the product over the
infinite places $v$ of $F$ with the embeddings $\tau_v$ of $E$ given by the setting. Let $g=g_0g_1$ be
the product of the level primes. Let $u$ be a word. Put $e^+_a=e_{a+1}$ for $a<r$ and $e^+_r=z_u$, and
for $0\le m\le r$ let $U_m$ be the span of $e_1,\dots,e_m$, in $V_v=V\otimes_{E,\tau_v}\mathbf C$ or in
$\hat V$.
\begin{enumerate}
\item $S_\infty$ is the set of $s\in G_\infty$ such that, at every infinite place $v$, $(s_v-1)e_a\in
  U_{a-1}$ for every $a$ and $(s_v-1)z_u\in U_r$; and $N_\infty$ is the set of $s\in S_\infty$ with
  $(s_v-1)y\in U_r$ for every $v$ and every $y\in V_v$ orthogonal to all $e_a$. The subsets $S_{\mathrm f}$
  and $N_{\mathrm f}\subset S_{\mathrm f}$ of $G_{\mathrm f}$ are defined by the same conditions in
  $\hat V$. The subgroups $H_\infty\subset G_\infty$ and $H_{\mathrm f}\subset G_{\mathrm f}$ are those
  that fix every $e_a$, $f_a$ and $z_u$.
\item For $s\in G_\infty$ and $x\in G_{\mathrm f}$,
  \[
   \psi_\infty(s)=\prod_v\exp\Big(\frac{2\pi i}{g}\sum_{a=1}^r\big(c_{a,v}+\overline{c_{a,v}}\big)\Big),
   \quad c_{a,v}=h(f_a,s_ve^+_a),\qquad
   \psi_{\mathrm f}(x)=\psi_{E,\mathrm f}\Big(\frac1g\sum_{a=1}^rh(f_a,xe^+_a)\Big).
  \]
\item The Bessel data of $u$ are a finite-dimensional Hilbert space $V_{\pi_u}$ with operators
  $\pi_u(s)$, $s\in G_\infty$, such that $\pi_u(1)=1$, $\pi_u(st)=\pi_u(s)\pi_u(t)$ and $\pi_u(s)$ is
  unitary for $s,t\in S_\infty$, and $\pi_u(s)=1$ for $s\in N_\infty$; and the archimedean Bessel function
  $W_u\colon G_\infty\to V_{\pi_u}$, which is continuous, satisfies
  $W_u(sx)=\psi_\infty(s)\pi_u(s)W_u(x)$ for $s\in S_\infty$, and is reproduced by $\varphi$: for every
  $x$, $z\mapsto\varphi(z)W_u(xz)$ is integrable on $G_\infty$ with integral $W_u(x)$.
\item $\Gamma_u$ is the set of elements of $\Gamma$ in the subgroup of $G_{\mathbf A}$ generated by
  $S_\infty\times S_{\mathrm f}$, and $\Gamma^H_u=\Gamma\cap(H_\infty\times H_{\mathrm f})$.
\item $K_u$ is the group of $h\in H_{\mathrm f}$ with $x_u^{-1}hx_u\in K_2$. A \emph{Poincaré datum} of
  $u$ is a linear map $T$ from $V_{\pi_u}$ to the functions on $H_\infty\times H_{\mathrm f}$ such that,
  for $\xi\in V_{\pi_u}$ and $p\in H_\infty\times H_{\mathrm f}$: $T\xi((\gamma_\infty,\gamma_{\mathrm
  f})p)=T\xi(p)$ for $\gamma\in\Gamma^H_u$; $T\xi(p(1,\kappa))=\chi_2(x_u^{-1}\kappa x_u)T\xi(p)$ for
  $\kappa\in K_u$; and $T(\pi_u(\eta)\xi)(p)=T\xi(p(\eta,1))$ for $\eta\in H_\infty$.
\end{enumerate}
\end{definition}
\leanref{Formalism/LocalToGlobal/SiegelLatticeTheory/HaarTransversals/HidaSetting.lean}{Formalism.L2.SG.HidaSetting.HidaSetting}
\leanref{Formalism/LocalToGlobal/SiegelLatticeTheory/HaarTransversals/HidaSetting.lean}{Formalism.L2.SG.HidaSetting.Sinf}
\leanref{Formalism/LocalToGlobal/SiegelLatticeTheory/HaarTransversals/HidaSetting.lean}{Formalism.L2.SG.HidaSetting.Ninf}
\leanref{Formalism/LocalToGlobal/SiegelLatticeTheory/HaarTransversals/HidaSetting.lean}{Formalism.L2.SG.HidaSetting.Nfin}
\leanref{Formalism/LocalToGlobal/SiegelLatticeTheory/HaarTransversals/HidaSetting.lean}{Formalism.L2.SG.HidaSetting.psiInf}
\leanref{Formalism/LocalToGlobal/SiegelLatticeTheory/HaarTransversals/HidaSetting.lean}{Formalism.L2.SG.HidaSetting.psiFin}
\leanref{Formalism/LocalToGlobal/SiegelLatticeTheory/HaarTransversals/HidaSetting.lean}{Formalism.L2.SG.HidaSetting.nxt}
\leanref{Formalism/LocalToGlobal/SiegelLatticeTheory/HaarTransversals/HidaSetting.lean}{Formalism.L2.SG.HidaSetting.IsRepS}
\leanref{Formalism/LocalToGlobal/SiegelLatticeTheory/HaarTransversals/HidaSetting.lean}{Formalism.L2.SG.HidaSetting.IsWu}
\leanref{Formalism/LocalToGlobal/SiegelLatticeTheory/HaarTransversals/PoincareInHH.lean}{Formalism.L2.SG.PoincareInHH.Sfin}
\leanref{Formalism/LocalToGlobal/SiegelLatticeTheory/HaarTransversals/PoincareInHH.lean}{Formalism.L2.SG.PoincareInHH.PoincareSetting.Gam}
\leanref{Formalism/LocalToGlobal/SiegelLatticeTheory/HaarTransversals/PoincareInHH.lean}{Formalism.L2.SG.PoincareInHH.PoincareSetting.Ku}
\leanref{Formalism/LocalToGlobal/SiegelLatticeTheory/HaarTransversals/PoincareInHH.lean}{Formalism.L2.SG.PoincareInHH.PoincareSetting.Cu}
\leanref{Formalism/LocalToGlobal/SiegelLatticeTheory/SiegelMassRatios/ColsDim.lean}{Formalism.L2.SG.ColsDim.Cols}

\begin{lemma}\label{asm-division-kernel-finite-series-function-unipotent}
In Situation~\ref{asm-division-kernel-finite-situation}, with the notation of
Definition~\ref{asm-division-kernel-finite-series-function-bessel}, let $u$ be a word.
\begin{enumerate}
\item $S_\infty\times S_{\mathrm f}$ is a subgroup of $G_{\mathbf A}$, so
  $\Gamma_u=\Gamma\cap(S_\infty\times S_{\mathrm f})$.
\item $N_{\mathrm f}$ is a subgroup of $G_{\mathrm f}$, and $\psi_{\mathrm f}(nn')=\psi_{\mathrm
  f}(n)\psi_{\mathrm f}(n')$ for $n,n'\in N_{\mathrm f}$.
\item For $h\in H_{\mathrm f}$ and $n\in N_{\mathrm f}$, $hnh^{-1}\in N_{\mathrm f}$ and
  $\psi_{\mathrm f}(hnh^{-1})=\psi_{\mathrm f}(n)$.
\item $H_\infty\subset S_\infty$, and $\psi_\infty(s\eta)=\psi_\infty(s)$ for $s\in G_\infty$ and
  $\eta\in H_\infty$.
\end{enumerate}
\end{lemma}
\leanref{Formalism/LocalToGlobal/SiegelLatticeTheory/HaarTransversals/UnitaryHaarProducts/InHHAlg.lean}{Formalism.L2.SG.PoincareInHH.Cut.SSub}
\leanref{Formalism/LocalToGlobal/SiegelLatticeTheory/HaarTransversals/UnitaryHaarProducts/InHHAlg.lean}{Formalism.L2.SG.PoincareInHH.Cut.sGroup_holds}
\leanref{Formalism/LocalToGlobal/SiegelLatticeTheory/HaarTransversals/UnitaryHaarProducts/InHHAlg.lean}{Formalism.L2.SG.PoincareInHH.Cut.sratComp_holds}
\leanref{Formalism/LocalToGlobal/SiegelLatticeTheory/HaarTransversals/UnitaryHaarProducts/InHHAlg.lean}{Formalism.L2.SG.PoincareInHH.Cut.C1_mul}
\leanref{Formalism/LocalToGlobal/SiegelLatticeTheory/HaarTransversals/UnitaryHaarProducts/InHHAlg.lean}{Formalism.L2.SG.PoincareInHH.Cut.C1_of_left_inverse}
\leanref{Formalism/LocalToGlobal/SiegelLatticeTheory/HaarTransversals/UnitaryHaarProducts/InHHAlg.lean}{Formalism.L2.SG.PoincareInHH.Cut.nfinGroup_holds}
\leanref{Formalism/LocalToGlobal/SiegelLatticeTheory/HaarTransversals/UnitaryHaarProducts/InHHAlg.lean}{Formalism.L2.SG.PoincareInHH.Cut.nfin_ca_mul}
\leanref{Formalism/LocalToGlobal/SiegelLatticeTheory/HaarTransversals/UnitaryHaarProducts/InHHAlg.lean}{Formalism.L2.SG.PoincareInHH.Cut.nfinConj_holds}
\leanref{Formalism/LocalToGlobal/SiegelLatticeTheory/HaarTransversals/UnitaryHaarProducts/InHHCut.lean}{Formalism.L2.SG.PoincareInHH.Cut.NfinGroup}
\leanref{Formalism/LocalToGlobal/SiegelLatticeTheory/HaarTransversals/UnitaryHaarProducts/InHHCut.lean}{Formalism.L2.SG.PoincareInHH.Cut.NfinConj}
\leanref{Formalism/LocalToGlobal/SiegelLatticeTheory/HaarTransversals/PoincareInHH.lean}{Formalism.L2.SG.PoincareInHH.PoincareSetting.hinf_le_sinf}
\leanref{Formalism/LocalToGlobal/SiegelLatticeTheory/HaarTransversals/UnitaryHaarProducts/InHHAlg.lean}{Formalism.L2.SG.PoincareInHH.Cut.psiInf_mul_H}
\begin{proof}
\emph{Step 1.} Work at an infinite place in $V_v$, or in $\hat V$, over the ring $A=\mathbf C$ or
$\mathbf A_{E,\mathrm f}$. Let $M$ be an invertible $A$-linear map with $(M-1)e_a\in U_{a-1}$ for every
$a$. Then $(M-1)U_m\subset U_{m-1}$ for $1\le m\le r$, so $MU_m\subset U_m$; and $M^{-1}U_m\subset U_m$, by
induction on $m$, as $M^{-1}y=y-M^{-1}(M-1)y$. For two such maps, $MN-1=M(N-1)+(M-1)$ and
$M^{-1}-1=-M^{-1}(M-1)$.

(1) Let $s,t$ satisfy the conditions defining $S_\infty$ at a place $v$, or those defining
$S_{\mathrm f}$. By Step 1, $(st-1)e_a=s(t-1)e_a+(s-1)e_a\in U_{a-1}$ and
$(s^{-1}-1)e_a=-s^{-1}(s-1)e_a\in U_{a-1}$, and in the same way $(st-1)z_u,(s^{-1}-1)z_u\in U_r$. So
$S_\infty$ and $S_{\mathrm f}$ are subgroups, and so is $S_\infty\times S_{\mathrm f}$, which is then the
subgroup it generates.

(2) The computation of (1) applies to $N_{\mathrm f}$; for $y$ orthogonal to every $e_a$ it gives
$(nn'-1)y=n(n'-1)y+(n-1)y\in U_r$ and $(n^{-1}-1)y=-n^{-1}(n-1)y\in U_r$, as $nU_r\subset U_r$. Fix $a$.
Then $(n'-1)e^+_a\in U_a$, so $(n-1)(n'-1)e^+_a\in U_{a-1}$, which is orthogonal to $f_a$; and
$h(f_a,e^+_a)=0$. Writing $nn'e^+_a=ne^+_a+(n-1)(n'-1)e^+_a+n'e^+_a-e^+_a$ gives
$h(f_a,nn'e^+_a)=h(f_a,ne^+_a)+h(f_a,n'e^+_a)$. As $\psi_{E,\mathrm f}$ is additive, $\psi_{\mathrm
f}(nn')=\psi_{\mathrm f}(n)\psi_{\mathrm f}(n')$.

(3) The element $h^{-1}$ fixes $e_a$ and $z_u$, and $h$ preserves each $U_m$, as it fixes the $e_b$. So
$(hnh^{-1}-1)e_a=h(n-1)e_a\in U_{a-1}$ and $(hnh^{-1}-1)z_u=h(n-1)z_u\in U_r$. If $y$ is orthogonal to
every $e_a$, so is $h^{-1}y$, as $h$ is an isometry fixing the $e_a$, and
$(hnh^{-1}-1)y=h(n-1)h^{-1}y\in U_r$. Finally $h(f_a,hnh^{-1}e^+_a)=h(h^{-1}f_a,nh^{-1}e^+_a)=h(f_a,ne^+_a)$,
as $h^{-1}$ fixes $f_a$ and $e^+_a$.

(4) An element of $H_\infty$ fixes the $e_a$ and $z_u$, so it lies in $S_\infty$. For $\eta\in H_\infty$,
$\eta_ve^+_a=e^+_a$, so $h(f_a,s_v\eta_ve^+_a)=h(f_a,s_ve^+_a)$ at every $v$.
\end{proof}

\begin{definition}\label{asm-division-kernel-finite-series-function-inputs-level}
In Situation~\ref{asm-division-kernel-finite-situation}, with the notation of
Definition~\ref{asm-division-kernel-finite-series-function-bessel}, the lattice, the level subgroup and the
base points of the setting are of the following form.
\begin{enumerate}
\item $\mathfrak a_a=\mathfrak d^{a-r-1}$ for $1\le a\le r$, where $\mathfrak d=\mathfrak d_F\mathcal O_E$ is the
  ideal of $\mathcal O_E$ generated by the different $\mathfrak d_F$ of $F$; and $z_u\in L_b$ for every word
  $u$.
\item The setting gives a level basis $e^0_a,f^0_a$ ($1\le a\le r$), $z^0$, $b_i,b'_i$ ($1\le i\le k$),
  $\epsilon$ for $(\hat L,g)$ (Definition~\ref{asm-division-kernel-finite-neat-level}), where $k$ is the
  rank $k$ of $D$. The group $K_2$ is its level subgroup, and for $\kappa\in G_{\mathrm f}$
  \[
   \chi_2(\kappa)=\psi_{E,\mathrm f}\Big(\frac1g\Big(\sum_{a=1}^rh(f^0_a,\kappa e^{0+}_a)
   +\sum_{i=1}^{k-1}h(b'_i,\kappa b_{i+1})\Big)\Big),
  \]
  with $e^{0+}_a=e^0_{a+1}$ for $a<r$ and $e^{0+}_r=z^0$.
\item For every word $u$, $x_u$ fixes every $e_a$ and $f_a$, and each of the vectors $e^0_a-e_a$,
  $f^0_a-f_a$ ($1\le a\le r$) and $x_u^{-1}z_u-z^0$ has a multiple by an integer prime to $g$ in $g\hat L$.
\end{enumerate}
\end{definition}
\leanref{Formalism/LocalToGlobal/SiegelLatticeTheory/HaarTransversals/HidaSetting.lean}{Formalism.L2.SG.HidaSetting.HidaSetting}
\leanref{Formalism/LocalToGlobal/SiegelLatticeTheory/HaarTransversals/HidaSetting.lean}{Formalism.L2.SG.HidaSetting.aaD}
\leanref{Formalism/SharpRestriction/CuspSpanNormBounds/BlockEnds.lean}{Formalism.MainNotions.S.ddS}
\leanref{Formalism/LocalToGlobal/SiegelLatticeTheory/HaarTransversals/HidaSetting.lean}{Formalism.L2.SG.HidaSetting.chi2f}
\leanref{Formalism/LocalToGlobal/SiegelLatticeTheory/HaarTransversals/HidaSetting.lean}{Formalism.L2.SG.HidaSetting.HidaSetting.chi2}

\begin{lemma}\label{asm-division-kernel-finite-series-function-inputs-trace}
Let $E$ be a CM field with totally real subfield $F$, let $\tau_v$ be an embedding of $E$ above each
infinite place $v$ of $F$, and put
$\psi_{E,\infty}(x)=\prod_v\exp\big(2\pi i(\tau_v(x)+\overline{\tau_v(x)})\big)$ for $x\in E$. Let
$\hat{\mathcal O}$ be the ring of integral elements of $\mathbf A_{E,\mathrm f}=E\otimes_F\mathbf
A_{F,\mathrm f}$, and let $\psi$ be a function on $\mathbf A_{E,\mathrm f}$ with
$\psi(x+x')=\psi(x)\psi(x')$, $\psi(t)=1$ for $t\in\hat{\mathcal O}$, and $\psi(x)\psi_{E,\infty}(x)=1$
for $x\in E$.
\begin{enumerate}
\item If $y\in E$ and $\operatorname{Tr}_{E/\mathbf Q}(oy)\in\mathbf Z$ for every $o\in\mathcal O_E$, then
  $\psi(ty)=1$ for every $t\in\hat{\mathcal O}$.
\item Let $E'/F'$ be any finite extension of number fields, let $\mathfrak d_{F'}$ be the different of
  $F'$, and let $y\in(\mathfrak d_{F'}\mathcal O_{E'})^{-1}$. Then $\operatorname{Tr}_{E'/\mathbf Q}(oy)\in\mathbf
  Z$ for every $o\in\mathcal O_{E'}$.
\end{enumerate}
\end{lemma}
\leanref{Formalism/LocalToGlobal/SiegelLatticeTheory/HaarTransversals/UnitaryHaarProducts/LevelCore.lean}{Formalism.L2.SG.Pf.PsiTrCore}
\leanref{Formalism/LocalToGlobal/SiegelLatticeTheory/HaarTransversals/UnitaryHaarProducts/LevelCore.lean}{Formalism.L2.SG.Pf.DdTraceCore}
\leanref{Formalism/LocalToGlobal/SiegelLatticeTheory/TraceIntegralEstimates/PsiTrCore.lean}{Formalism.L2.S.Pv.PsiTrCore.psiTrCore_holds}
\leanref{Formalism/LocalToGlobal/SiegelLatticeTheory/TraceIntegralEstimates/PsiTrCore.lean}{Formalism.L2.S.Pv.PsiTrCore.intAf_le_congSub}
\leanref{Formalism/LocalToGlobal/SiegelLatticeTheory/TraceIntegralEstimates/DdTraceCore.lean}{Formalism.L2.S.Pv.DdTraceCore.ddTraceCore_holds}
\leanref{Formalism/LocalToGlobal/SiegelLatticeTheory/TraceIntegralEstimates/DdTraceCore.lean}{Formalism.L2.S.Pv.DdTraceCore.relTrace_mem_inv_different}
\begin{proof}
(1) Choose an integer $N\ne0$ with $Ny\in\mathcal O_E$. Every $t\in\hat{\mathcal O}$ is $o+Ns$ with
$o\in\mathcal O_E$ and $s\in\hat{\mathcal O}$. Indeed, the elements of this form make up a subring of
$\mathbf A_{E,\mathrm f}$. It contains $\mathcal O_E$, and it contains every integral finite adele $c$ of
$F$: there is $b\in\mathcal O_F$ with $c-b\in N\hat{\mathcal O}_F$, by approximation at the places dividing
$N$. These generate $\hat{\mathcal O}$. Then $ty=oy+s(Ny)$, and $\psi(s(Ny))=1$ as $s(Ny)\in\hat{\mathcal
O}$. The numbers $\tau_v(x)$ and $\overline{\tau_v(x)}$ are the images of $x$ under the embeddings of $E$
into $\mathbf C$, so $\psi_{E,\infty}(oy)=\exp\big(2\pi i\operatorname{Tr}_{E/\mathbf Q}(oy)\big)=1$, and
$\psi(oy)=1$. So $\psi(ty)=\psi(oy)\psi(s(Ny))=1$.

(2) For $d\in\mathfrak d_{F'}$, $yd\in\mathcal O_{E'}$, so
$d\operatorname{Tr}_{E'/F'}(oy)=\operatorname{Tr}_{E'/F'}(oyd)\in\mathcal O_{F'}$. Hence
$\operatorname{Tr}_{E'/F'}(oy)\in\mathfrak d_{F'}^{-1}$, the dual of $\mathcal O_{F'}$ for the trace form,
and $\operatorname{Tr}_{E'/\mathbf Q}(oy)=\operatorname{Tr}_{F'/\mathbf Q}\big(\operatorname{Tr}_{E'/F'}(oy)\big)\in\mathbf Z$.
\end{proof}

\begin{lemma}\label{asm-division-kernel-finite-series-function-inputs-congruence}
In Situation~\ref{asm-division-kernel-finite-situation}, with the notation of
Definitions~\ref{asm-division-kernel-finite-series-function-bessel}
and~\ref{asm-division-kernel-finite-series-function-inputs-level}, write $\langle S\rangle$ for the sum of
the $\hat{\mathcal O}$-span of $S\subset\hat V$ and $g\hat L$.
\begin{enumerate}
\item Let $M\subset\hat V$ be an $\hat{\mathcal O}$-submodule containing $g\hat L$, let $y\in\hat L$, and let
  $c$ be an integer prime to $g$ with $cy\in M$. Then $y\in M$.
\end{enumerate}
Let $u$ be a word and $C$ an integer prime to $g$ such that $C(e^0_a-e_a)$, $C(f^0_a-f_a)$ and
$C(x_u^{-1}z_u-z^0)$ lie in $g\hat L$ for every $a$.
\begin{enumerate}\setcounter{enumi}{1}
\item $Ce_a\in\hat L$ and $Cf_a\in\hat L$ for every $a$.
\item Let $T\subset\{1,\dots,r\}$, and let $s\in\hat L$ lie in the $\mathbf A_{E,\mathrm f}$-span of
  $e_1,\dots,e_r$, with $h(f_b,s)=0$ for $b\notin T$. Then $s\in\langle e^0_b:b\in T\rangle$.
\end{enumerate}
\end{lemma}
\leanref{Formalism/LocalToGlobal/SiegelLatticeTheory/HaarTransversals/UnitaryHaarProducts/HermitianETransportInvariant.lean}{Formalism.L2.SG.Pf.Level.bez_vec}
\leanref{Formalism/LocalToGlobal/SiegelLatticeTheory/HaarTransversals/UnitaryHaarProducts/HermitianETransportInvariant.lean}{Formalism.L2.SG.Pf.Level.C_e_Lf}
\leanref{Formalism/LocalToGlobal/SiegelLatticeTheory/HaarTransversals/UnitaryHaarProducts/HermitianETransportInvariant.lean}{Formalism.L2.SG.Pf.Level.coef}
\leanref{Formalism/LocalToGlobal/SiegelLatticeTheory/HaarTransversals/UnitaryHaarProducts/HermitianETransportInvariant.lean}{Formalism.L2.SG.Pf.Level.span_repr}
\begin{proof}
(1) Choose integers $\alpha,\beta$ with $\alpha c+\beta g=1$. Then $y=\alpha(cy)+\beta(gy)$, with $cy\in M$
and $gy\in g\hat L\subset M$.

(2) $Ce_a=Ce^0_a-C(e^0_a-e_a)$, with $e^0_a\in\hat L$ and $g\hat L\subset\hat L$; the same for $f_a$.

(3) As $h(f_b,e_c)=\delta_{bc}$, a vector $s$ of the span of the $e_c$ is $\sum_bh(f_b,s)e_b$, and here the
terms with $b\notin T$ vanish. Put $s'=\sum_{b\in T}h(f^0_b,s)e^0_b$; its coefficients lie in
$\hat{\mathcal O}$, as $h$ is integral on $\hat L$, so $s'$ lies in the $\hat{\mathcal O}$-span of the
$e^0_b$, $b\in T$. For $b\in T$,
\[
 C^2h(f_b,s)e_b-C^2h(f^0_b,s)e^0_b=h\big(C(f_b-f^0_b),s\big)\,Ce_b+h(f^0_b,s)\,C\,C(e_b-e^0_b).
\]
The first term lies in $g\hat L$, as $h(m,s)\in g\hat{\mathcal O}$ for $m\in g\hat L$ and $s\in\hat L$, and
$Ce_b\in\hat L$ by (2); the second lies in $g\hat L$, as $C(e_b-e^0_b)$ does. So $C^2s-C^2s'\in g\hat L$ and
$C^2s\in\langle e^0_b:b\in T\rangle$. As $C^2$ is prime to $g$, (1) gives the claim.
\end{proof}

\begin{lemma}\label{asm-division-kernel-finite-series-function-inputs-membership}
In Situation~\ref{asm-division-kernel-finite-situation}, with the notation of
Definitions~\ref{asm-division-kernel-finite-series-function-bessel}
and~\ref{asm-division-kernel-finite-series-function-inputs-level}, and with $\langle S\rangle$ as in
Lemma~\ref{asm-division-kernel-finite-series-function-inputs-congruence}, let $u$ be a word, let $n\in
N_{\mathrm f}\cap K_{\mathrm{set}}$, and put $\kappa=x_u^{-1}nx_u$. Then $\kappa\in K_{\mathrm{set}}$, and for
every $a$ and every $i$
\[
 (\kappa-1)e^0_a\in\langle e^0_c:c<a\rangle,\qquad(\kappa-1)z^0\in\langle e^0_c:c\le r\rangle,\qquad
 (\kappa-1)b_i\in\langle e^0_c:c\le r\rangle.
\]
In particular $\kappa\in K_2$.
\end{lemma}
\leanref{Formalism/LocalToGlobal/SiegelLatticeTheory/HaarTransversals/UnitaryHaarProducts/HermitianETransportInvariant.lean}{Formalism.L2.SG.Pf.Level.unipotent_conj_mem_level}
\leanref{Formalism/LocalToGlobal/SiegelLatticeTheory/HaarTransversals/UnitaryHaarProducts/HermitianETransportInvariant.lean}{Formalism.L2.SG.Pf.Level.cong_e}
\leanref{Formalism/LocalToGlobal/SiegelLatticeTheory/HaarTransversals/UnitaryHaarProducts/HermitianETransportInvariant.lean}{Formalism.L2.SG.Pf.Level.cong_alt}
\leanref{Formalism/LocalToGlobal/SiegelLatticeTheory/HaarTransversals/UnitaryHaarProducts/HermitianETransportInvariant.lean}{Formalism.L2.SG.Pf.Level.cong_b}
\leanref{Formalism/LocalToGlobal/SiegelLatticeTheory/HaarTransversals/UnitaryHaarProducts/HermitianETransportInvariant.lean}{Formalism.L2.SG.Pf.Level.exists_goodC}
\begin{proof}
Let $C$ be the product of the integers prime to $g$ of
Definition~\ref{asm-division-kernel-finite-series-function-inputs-level}(3). It is prime to $g$, and
$C(e^0_a-e_a)$, $C(f^0_a-f_a)$ and $C(x_u^{-1}z_u-z^0)$ lie in $g\hat L$, so
Lemma~\ref{asm-division-kernel-finite-series-function-inputs-congruence} applies to $C$. As $x_u$ and $n$
lie in $K_{\mathrm{set}}$, so does $\kappa$, and $\kappa$ maps $\hat L$ and $g\hat L$ into themselves. As
$x_u$ fixes every $e_a$, $x_u^{-1}$ fixes the span $U_r$ of the $e_a$ in $\hat V$ pointwise.

\emph{The vectors $e^0_a$.} Put $s=(n-1)e_a$, which lies in $U_{a-1}$ as $n\in N_{\mathrm f}$. Since $x_u$
fixes $e_a$, $(\kappa-1)e_a=x_u^{-1}s=s$. So $Cs=(\kappa-1)(Ce_a)\in\hat L$, as $Ce_a\in\hat L$. As
$Cs\in U_{a-1}$, $h(f_b,Cs)=0$ for $b\ge a$, and Lemma~\ref{asm-division-kernel-finite-series-function-inputs-congruence}(3)
gives $Cs\in\langle e^0_c:c<a\rangle$. Then
\[
 C(\kappa-1)e^0_a=Cs+(\kappa-1)\big(C(e^0_a-e_a)\big)\in\langle e^0_c:c<a\rangle,
\]
and, as $(\kappa-1)e^0_a\in\hat L$, part (1) of that lemma gives $(\kappa-1)e^0_a\in\langle e^0_c:c<a\rangle$.

\emph{The vector $z^0$.} Put $s=(n-1)z_u$, which lies in $U_r$, and in $\hat L$ as $z_u\in L_0$. Then
$(\kappa-1)x_u^{-1}z_u=x_u^{-1}s=s$, and $Cs\in\langle e^0_c:c\le r\rangle$ by part (3) of that lemma.
So $C(\kappa-1)z^0=Cs-(\kappa-1)\big(C(x_u^{-1}z_u-z^0)\big)$ lies in $\langle e^0_c:c\le r\rangle$, and part
(1) gives the claim.

\emph{The vectors $b_i$.} Put $\beta=x_ub_i\in\hat L$, so that $(\kappa-1)b_i=x_u^{-1}(n-1)\beta$. Write
$\beta=\sum_ah(f_a,\beta)e_a+p+\sum_ah(e_a,\beta)f_a$, where $p$ is orthogonal to every $e_a$ and $f_a$. Put
\[
 s=(n-1)\Big(\sum_ah(f_a,\beta)e_a+p\Big),\qquad m=C^2\sum_ah(e_a,\beta)(n-1)f_a,
\]
so that $C^2(n-1)\beta=C^2s+m$. By the conditions defining $N_{\mathrm f}$, $s\in U_r$. Next, $m\in g\hat L$.
As $x_u$ is an isometry fixing $e_a$, $h(e_a,\beta)=h(e_a,b_i)$, so
$Ch(e_a,\beta)=h(Ce^0_a,b_i)-h\big(C(e^0_a-e_a),b_i\big)$ lies in $g\hat{\mathcal O}$, by the congruence
$h(e^0_a,b_i)\equiv0$ of the level basis and as $C(e^0_a-e_a)\in g\hat L$. Writing $Ch(e_a,\beta)=gt_a$
with $t_a\in\hat{\mathcal O}$, the term $C^2h(e_a,\beta)(n-1)f_a=t_ag(n-1)(Cf_a)$ lies in $g\hat L$, as
$Cf_a\in\hat L$. So $x_u^{-1}m\in g\hat L$. As $x_u^{-1}$ fixes $s$,
\[
 C^2(\kappa-1)b_i=C^2s+x_u^{-1}m.
\]
Hence $C^2s\in\hat L\cap U_r$, so $C^2s\in\langle e^0_c:c\le r\rangle$ by part (3) of that lemma, and so is
$C^2(\kappa-1)b_i$. Part (1), with the integer $C^2$, gives the claim.

Finally $\langle e^0_c:c\le r\rangle\subset\langle e^0_c,b_j:c\le r,\ j<i\rangle$, so $\kappa$ satisfies the
conditions defining $K_2$ (Definition~\ref{asm-division-kernel-finite-neat-level}).
\end{proof}

\begin{lemma}\label{asm-division-kernel-finite-series-function-inputs-levi}
In Situation~\ref{asm-division-kernel-finite-situation}, with the notation of
Definition~\ref{asm-division-kernel-finite-series-function-bessel}, let $u$ be a word, let $X_m$ be the
$E$-span of $e_1,\dots,e_m$ in $V$ and $X=X_r$, and let $G$ be an $E$-linear isometry of $V$ with
$(G-1)e_a\in X_{a-1}$ for every $a$ and $(G-1)z_u\in X$. Put $t=h(z_u,z_u)$, which is not $0$, and for
$y\in V$
\[
 P(y)=y-\sum_ah(f_a,y)e_a-\sum_ah(e_a,y)f_a-t^{-1}h(z_u,y)z_u,\qquad H(y)=P\big(GP(y)\big)+y-P(y).
\]
Then $H$ is an isometry of $V$ that fixes every $e_a$, $f_a$ and $z_u$, and $N=GH^{-1}$ satisfies
$(N-1)e_a\in X_{a-1}$ for every $a$, $(N-1)z_u\in X$, and $(N-1)y\in X$ for every $y\in V$ orthogonal to
every $e_a$.
\end{lemma}
\leanref{Formalism/LocalToGlobal/SiegelLatticeTheory/HaarTransversals/UnitaryHaarProducts/HermitianETransportInvariant.lean}{Formalism.L2.SG.Pf.Level.SRat}
\leanref{Formalism/LocalToGlobal/SiegelLatticeTheory/HaarTransversals/UnitaryHaarProducts/HermitianETransportInvariant.lean}{Formalism.L2.SG.Pf.Level.PW}
\leanref{Formalism/LocalToGlobal/SiegelLatticeTheory/HaarTransversals/UnitaryHaarProducts/HermitianETransportInvariant.lean}{Formalism.L2.SG.Pf.Level.Hm}
\leanref{Formalism/LocalToGlobal/SiegelLatticeTheory/HaarTransversals/UnitaryHaarProducts/HermitianETransportInvariant.lean}{Formalism.L2.SG.Pf.Level.HGL}
\leanref{Formalism/LocalToGlobal/SiegelLatticeTheory/HaarTransversals/UnitaryHaarProducts/HermitianETransportInvariant.lean}{Formalism.L2.SG.Pf.Level.SRat.hE_H}
\leanref{Formalism/LocalToGlobal/SiegelLatticeTheory/HaarTransversals/UnitaryHaarProducts/HermitianETransportInvariant.lean}{Formalism.L2.SG.Pf.Level.SRat.nrat}
\begin{proof}
The vector $z_u$ is orthogonal to the $e_a$ and $f_a$, and $t=-T_u\ne0$. So $P$ is $E$-linear, vanishes on
the $e_a$, $f_a$ and $z_u$, fixes the vectors orthogonal to all of them, and its values are orthogonal to
all of them; hence $P(P(y))=P(y)$, and $y-P(y)$, a combination of the $e_a$, $f_a$ and $z_u$, is orthogonal to
every value of $P$. So
\[
 h(y,y')=h\big(P(y),P(y')\big)+h\big(y-P(y),y'-P(y')\big)\qquad(y,y'\in V).
\]

\emph{Step 1: $G^{-1}X\subset X$ and $(G^{-1}-1)z_u\in X$.} By induction on $a$:
$G^{-1}e_a=e_a-G^{-1}(G-1)e_a$, and $(G-1)e_a\in X_{a-1}$, which $G^{-1}$ maps into $X$ by induction. Then
$(G^{-1}-1)z_u=-G^{-1}(G-1)z_u\in X$.

\emph{Step 2: if $w\in V$ is orthogonal to every $e_a$ and to $z_u$, so is $Gw$, and $Gw-P(Gw)\in X$.}
As $G$ is an isometry, $h(e_a,Gw)=h(G^{-1}e_a,w)=0$, as $G^{-1}e_a\in X$; and
$h(z_u,Gw)=h(z_u,w)+h\big((G^{-1}-1)z_u,w\big)=0$. So in $Gw-P(Gw)=\sum_ah(f_a,Gw)e_a+\sum_ah(e_a,Gw)f_a
+t^{-1}h(z_u,Gw)z_u$ only the first sum remains.

\emph{Step 3: $H$ is an isometry.} Let $w,w'$ be values of $P$. By Step 2, $Gw-P(Gw)\in X$, while
$Gw'-P(Gw')$ is orthogonal to every $e_a$, and so to $X$. By the identity above for $Gw$ and $Gw'$,
$h\big(P(Gw),P(Gw')\big)=h(Gw,Gw')=h(w,w')$. Now $P(H(y))=P(GP(y))$ and $H(y)-P(H(y))=y-P(y)$, so the identity
above gives
\[
 h\big(H(y),H(y')\big)=h\big(PGP(y),PGP(y')\big)+h\big(y-P(y),y'-P(y')\big)=h(y,y').
\]
So $H$ is an $E$-linear map preserving the nondegenerate form $h$; it is injective, hence bijective. It fixes
the $e_a$, $f_a$ and $z_u$, on which $P$ vanishes.

\emph{Step 4: $N$.} As $H^{-1}$ fixes $e_a$ and $z_u$, $(N-1)e_a=(G-1)e_a\in X_{a-1}$ and
$(N-1)z_u=(G-1)z_u\in X$. Let $y$ be orthogonal to every $e_a$, and put $y_1=H^{-1}y$, which is orthogonal to
every $e_a$, as $H$ is an isometry fixing them. Put $w=P(y_1)$ and $b=y_1-P(y_1)$, so that
$y=H(y_1)=P(Gw)+b$ and
\[
 (N-1)y=Gy_1-y=\big(Gw-P(Gw)\big)+(G-1)b.
\]
The vector $w$ is orthogonal to every $e_a$ and to $z_u$, so $Gw-P(Gw)\in X$ by Step 2. As $h(e_a,y_1)=0$,
$b=\sum_ah(f_a,y_1)e_a+t^{-1}h(z_u,y_1)z_u$, and $(G-1)b\in X$ by the hypotheses on $G$.
% Lean: P is PW, H is Hm G = toMatrix' (Hl G) with Hf G y = PW (G PW y) + (y - PW y), and HGL makes it an
% element of GL. The identity for h is hE_sub_PW_self (with hE_perp_sub, PW_fix, hE_e_PW, hE_f_PW, hE_z_PW).
% Step 1 is SRat.inv_e and SRat.coordinate_inverse_zeta; Step 2 is SRat.perp and SRat.sub_PW_mem; Step 3 is
% SRat.hE_H (hG, hPH, hdH) with unit_of_form and det_ne_of_unit; Step 4 is SRat.nrat (Hm_e, Hm_f,
% levi_fixes_z, inv_fix). The hypotheses on G are the structure SRat (unit, inv1, inv2, c1, c2).
\end{proof}

\begin{lemma}\label{asm-division-kernel-finite-series-function-inputs-transfer}
In Situation~\ref{asm-division-kernel-finite-situation}, with the notation of
Definition~\ref{asm-division-kernel-finite-series-function-bessel}, let $u$ be a word and let $X_m$ and $X$
be as in Lemma~\ref{asm-division-kernel-finite-series-function-inputs-levi}.
\begin{enumerate}
\item Let $T\subset\{1,\dots,r\}$ and $s\in V$ be such that the image of $s$ in $\hat V$ lies in the
  $\mathbf A_{E,\mathrm f}$-span of the $e_b$, $b\in T$. Then $s$ lies in the $E$-span of the $e_b$,
  $b\in T$.
\item Let $\phi\colon E\to R$ be a homomorphism to a commutative ring $R$, let $V_R$ be an $R$-module and
  $\iota\colon V\to V_R$ an additive map with $\iota(cx)=\phi(c)\iota(x)$ whose image spans $V_R$, and let
  $h_R\colon V_R\times V_R\to R$ be a map with $h_R(\iota x,\iota y)=\phi(h(x,y))$ that is additive and
  $R$-linear in the second variable when the first lies in $\iota(V)$. Let $N$ be an $E$-linear map of $V$
  with $(N-1)e_a\in X_{a-1}$ for every $a$, $(N-1)z_u\in X$, and $(N-1)y\in X$ for every $y\in V$
  orthogonal to every $e_a$, and let $M$ be an $R$-linear map of $V_R$ with $M\iota(x)=\iota(Nx)$. Then
  $(M-1)y$ lies in the $R$-span of the $\iota(e_a)$ for every $y\in V_R$ with $h_R(\iota(e_a),y)=0$ for
  every $a$. (The first two conditions on $N$ are not used.)
\end{enumerate}
\end{lemma}
\leanref{Formalism/LocalToGlobal/SiegelLatticeTheory/HaarTransversals/UnitaryHaarProducts/HermitianETransportInvariant.lean}{Formalism.L2.SG.Pf.Level.descend}
\leanref{Formalism/LocalToGlobal/SiegelLatticeTheory/HaarTransversals/UnitaryHaarProducts/HermitianETransportInvariant.lean}{Formalism.L2.SG.Pf.Level.transfer3}
\begin{proof}
(1) As $h(f_b,e_c)=\delta_{bc}$, in $\hat V$ the image of $s$ is $\sum_bh(f_b,s)e_b$, and $h(f_b,s)=0$ for
$b\notin T$. The map $E\to\mathbf A_{E,\mathrm f}$ is injective, so $s=\sum_{b\in T}h(f_b,s)e_b$ in $V$.

(2) For $y\in V_R$ put $Q(y)=y-\sum_bh_R(\iota e_b,y)\iota f_b$ and
$S(y)=(M-1)Q(y)-\sum_bh_R\big(\iota f_b,(M-1)Q(y)\big)\iota e_b$. Both are additive and $R$-linear. For
$x\in V$, $Q(\iota x)=\iota(q)$ with $q=x-\sum_bh(e_b,x)f_b$, which is orthogonal to every $e_a$; so
$(M-1)\iota(q)=\iota\big((N-1)q\big)$ with $(N-1)q\in X$, and a vector $x'\in X$ is
$\sum_bh(f_b,x')e_b$, so $S(\iota x)=0$. As the image of $\iota$ spans $V_R$, $S=0$. If
$h_R(\iota e_a,y)=0$ for every $a$, then $Q(y)=y$, and $S(y)=0$ says that
$(M-1)y=\sum_bh_R\big(\iota f_b,(M-1)y\big)\iota e_b$.
\end{proof}

\begin{lemma}\label{asm-division-kernel-finite-series-function-inputs}
In Situation~\ref{asm-division-kernel-finite-situation}, with the notation of
Definitions~\ref{asm-division-kernel-finite-series-function-bessel}
and~\ref{asm-division-kernel-finite-series-function-inputs-level}, let $u$ be a word.
\begin{enumerate}
\item If $n\in N_{\mathrm f}$ and $h\in H_{\mathrm f}$ satisfy $nh\in K_{\mathrm{set}}$, then
  $h\in K_{\mathrm{set}}$.
\item If $n\in N_{\mathrm f}\cap K_{\mathrm{set}}$, then $x_u^{-1}nx_u\in K_2$ and
  $\chi_2(x_u^{-1}nx_u)=\psi_{\mathrm f}(n)$.
\item Every $\gamma\in\Gamma$ with $\gamma_\infty\in S_\infty$ and $\gamma_{\mathrm f}\in S_{\mathrm f}$
  is $\nu\eta$ with $\nu\in\Gamma$, $\nu_\infty\in N_\infty$, $\nu_{\mathrm f}\in N_{\mathrm f}$ and
  $\eta\in\Gamma^H_u$. (The condition $\gamma_\infty\in S_\infty$ is not used.)
\item If $\nu\in\Gamma$ with $\nu_\infty\in N_\infty$ and $\nu_{\mathrm f}\in N_{\mathrm f}$, then
  $\psi_\infty(\nu_\infty)\psi_{\mathrm f}(\nu_{\mathrm f})=1$. (The conditions $\nu_\infty\in N_\infty$
  and $\nu_{\mathrm f}\in N_{\mathrm f}$ are not used.)
\end{enumerate}
\end{lemma}
\leanref{Formalism/LocalToGlobal/SiegelLatticeTheory/HaarTransversals/UnitaryHaarProducts/InHHAlg.lean}{Formalism.L2.SG.PoincareInHH.Cut.KsetSplit}
\leanref{Formalism/LocalToGlobal/SiegelLatticeTheory/HaarTransversals/UnitaryHaarProducts/InHHAlg.lean}{Formalism.L2.SG.PoincareInHH.Cut.UnipotentConjugateInLevel}
\leanref{Formalism/LocalToGlobal/SiegelLatticeTheory/HaarTransversals/UnitaryHaarProducts/InHHAlg.lean}{Formalism.L2.SG.PoincareInHH.Cut.RatLevi}
\leanref{Formalism/LocalToGlobal/SiegelLatticeTheory/HaarTransversals/UnitaryHaarProducts/InHHAlg.lean}{Formalism.L2.SG.PoincareInHH.Cut.PsiProd}
\leanref{Formalism/LocalToGlobal/SiegelLatticeTheory/HaarTransversals/UnitaryHaarProducts/HermitianETransportInvariant.lean}{Formalism.L2.SG.Pf.ksetSplit_holds}
\leanref{Formalism/LocalToGlobal/SiegelLatticeTheory/HaarTransversals/UnitaryHaarProducts/HermitianETransportInvariant.lean}{Formalism.L2.SG.Pf.Level.kset_split_fwd}
\leanref{Formalism/LocalToGlobal/SiegelLatticeTheory/HaarTransversals/UnitaryHaarProducts/HermitianETransportInvariant.lean}{Formalism.L2.SG.Pf.Level.prA_Lf}
\leanref{Formalism/LocalToGlobal/SiegelLatticeTheory/HaarTransversals/UnitaryHaarProducts/HermitianETransportInvariant.lean}{Formalism.L2.SG.Pf.unipotent_in_level}
\leanref{Formalism/LocalToGlobal/SiegelLatticeTheory/HaarTransversals/UnitaryHaarProducts/HermitianETransportInvariant.lean}{Formalism.L2.SG.Pf.ratLevi_holds}
\leanref{Formalism/LocalToGlobal/SiegelLatticeTheory/HaarTransversals/UnitaryHaarProducts/HermitianETransportInvariant.lean}{Formalism.L2.SG.Pf.psiProd_holds}
\leanref{Formalism/LocalToGlobal/SiegelLatticeTheory/HaarTransversals/UnitaryHaarProducts/HermitianETransportInvariant.lean}{Formalism.L2.SG.Pf.Level.psiProd_holds}
\leanref{Formalism/LocalToGlobal/SiegelLatticeTheory/HaarTransversals/UnitaryHaarProducts/HermitianETransportInvariant.lean}{Formalism.L2.SG.Pf.Level.unipotent_in_level_of_psi}
\leanref{Formalism/LocalToGlobal/SiegelLatticeTheory/HaarTransversals/UnitaryHaarProducts/HermitianETransportInvariant.lean}{Formalism.L2.SG.Pf.Level.Z_term}
\leanref{Formalism/LocalToGlobal/SiegelLatticeTheory/HaarTransversals/UnitaryHaarProducts/HermitianETransportInvariant.lean}{Formalism.L2.SG.Pf.Level.Z_bw}
\leanref{Formalism/LocalToGlobal/SiegelLatticeTheory/HaarTransversals/UnitaryHaarProducts/HermitianETransportInvariant.lean}{Formalism.L2.SG.Pf.Level.psi_B}
\leanref{Formalism/LocalToGlobal/SiegelLatticeTheory/HaarTransversals/UnitaryHaarProducts/HermitianETransportInvariant.lean}{Formalism.L2.SG.Pf.Level.coord_dd}
\leanref{Formalism/LocalToGlobal/SiegelLatticeTheory/HaarTransversals/UnitaryHaarProducts/HermitianETransportInvariant.lean}{Formalism.L2.SG.Pf.psiDd_of_cores}
\leanref{Formalism/LocalToGlobal/SiegelLatticeTheory/HaarTransversals/UnitaryHaarProducts/HermitianETransportInvariant.lean}{Formalism.L2.SG.Pf.Level.ratLevi_holds}
\begin{proof}
(1) Let $\mathrm{pr}(y)=y-\sum_ah(f_a,y)e_a-\sum_ah(e_a,y)f_a$, on $V$ and on $\hat V$. It vanishes on the
$e_a$ and $f_a$ and fixes the vectors orthogonal to all of them. So it maps $L_0$ into $L_0$: it kills
$\mathfrak a_ae_a$ and $\mathfrak a_a^{-1}f_a$ and fixes $L_b$. As it is $\mathbf A_{E,\mathrm f}$-linear on
$\hat V$, it maps $\hat L$ into $\hat L$.

Let $y\in\hat L$, $p=\mathrm{pr}(y)\in\hat L$ and $q=hp$. As $h$ is an isometry fixing the $e_a$ and $f_a$,
$h(e_a,q)=h(e_a,p)=0$ and $h(f_a,q)=0$. Also $nq=(nh)p\in\hat L$. As $q$ is orthogonal to every $e_a$ and
$n\in N_{\mathrm f}$, $(n-1)q\in U_r$, which $\mathrm{pr}$ kills; and $\mathrm{pr}(q)=q$. So
$q=\mathrm{pr}(nq)\in\hat L$. The vector $y-p$ is a combination of the $e_a$ and $f_a$, so $h$ fixes it,
and $hy=(y-p)+q\in\hat L$. So $h\hat L\subset\hat L$. By Lemma~\ref{asm-division-kernel-finite-series-function-unipotent}(2)
and (3), $n'=h^{-1}n^{-1}h$ lies in $N_{\mathrm f}$, and $n'h^{-1}=(nh)^{-1}\in K_{\mathrm{set}}$; the same
argument gives $h^{-1}\hat L\subset\hat L$. So $h\hat L=\hat L$.

(4) The element $\nu$ comes from an isometry $\gamma$ of $V$. Put
$c=\sum_ah(f_a,\gamma e^+_a)\in E$ and $q=c/g$. Then $\psi_{\mathrm f}(\nu_{\mathrm f})=\psi_{E,\mathrm
f}(q)$. At every infinite place $c_{a,v}=\tau_v(h(f_a,\gamma e^+_a))$, so $\sum_a(c_{a,v}+\overline{c_{a,v}})/g
=\tau_v(q)+\overline{\tau_v(q)}$, and $\psi_\infty(\nu_\infty)=\psi_{E,\infty}(q)$. So the product is
$\psi_{E,\mathrm f}(q)\psi_{E,\infty}(q)=1$.

(2) Let $n\in N_{\mathrm f}\cap K_{\mathrm{set}}$ and $\kappa=x_u^{-1}nx_u$. By
Lemma~\ref{asm-division-kernel-finite-series-function-inputs-membership}, $\kappa\in K_2$. Write
$\psi=\psi_{E,\mathrm f}$, and $c\equiv0$ for $c\in g\hat{\mathcal O}$; then $h(m,s)\equiv0$ and
$h(s,m)\equiv0$ for $m\in g\hat L$ and $s\in\hat L$, as $h$ is integral on $\hat L$. Let $C$ be the product of
the integers of Definition~\ref{asm-division-kernel-finite-series-function-inputs-level}(3), so that
Lemma~\ref{asm-division-kernel-finite-series-function-inputs-congruence} applies to $C$. Put
\[
 A=\sum_ah(f^0_a,\kappa e^{0+}_a)+\sum_{i=1}^{k-1}h(b'_i,\kappa b_{i+1}),\qquad
 B=\sum_ah(f_a,ne^+_a),
\]
so that $\chi_2(\kappa)=\psi(A/g)$ and $\psi_{\mathrm f}(n)=\psi(B/g)$.

\emph{Step 1: $C^2\big(h(f^0_a,\kappa e^{0+}_a)-h(f_a,ne^+_a)\big)\equiv0$ for every $a$.} Put $w^0=e^{0+}_a$
and $w=e^+_a$. As $x_u$ is an isometry, $h(f^0_a,\kappa w^0)=h(x_uf^0_a,nx_uw^0)$. As $x_u$ fixes $f_a$ and
maps $g\hat L$ into itself, $C(x_uf^0_a-f_a)=x_uC(f^0_a-f_a)\in g\hat L$. Likewise $C(x_uw^0-w)\in g\hat L$:
for $a<r$ it is $x_uC(e^0_{a+1}-e_{a+1})$, and for $a=r$ it is $-x_uC(x_u^{-1}z_u-z^0)$. Now
\[
 C^2\big(h(x_uf^0_a,nx_uw^0)-h(f_a,nw)\big)=C\,h\big(C(x_uf^0_a-f_a),nx_uw^0\big)
 +h\big(Cf_a,nC(x_uw^0-w)\big),
\]
and both terms are $\equiv0$, as $nx_uw^0\in\hat L$, $Cf_a\in\hat L$
(Lemma~\ref{asm-division-kernel-finite-series-function-inputs-congruence}(2)) and $n$ maps $g\hat L$ into
itself.

\emph{Step 2: $h(b'_i,\kappa b_{i+1})\equiv0$ for $1\le i<k$.} Write $\kappa b_{i+1}=b_{i+1}+(\kappa-1)b_{i+1}$.
Then $h(b'_i,b_{i+1})\equiv0$ by the level basis, as $i\ne i+1$. By
Lemma~\ref{asm-division-kernel-finite-series-function-inputs-membership}, $(\kappa-1)b_{i+1}$ lies in the sum
of the $\hat{\mathcal O}$-span of the $e^0_c$ and $g\hat L$; and
$h(b'_i,e^0_c)=\sigma\big(h(e^0_c,b'_i)\big)\equiv0$ by the level basis, as $\sigma$ preserves $g\hat{\mathcal
O}$.

By Steps 1 and 2, $C^2(A-B)\equiv0$, so $\psi\big(C^2(A-B)/g\big)=1$, as $\psi$ is $1$ on $\hat{\mathcal O}$.

\emph{Step 3: $\psi(A-B)=1$.} First, $A\in\hat{\mathcal O}$, as $\kappa\hat L=\hat L$, the vectors of the level
basis lie in $\hat L$ and $h$ is integral on $\hat L$; so $\psi(A)=1$. Next, $h(f_a,e^+_a)=0$ for every $a$, so
$B=\sum_ac_a$ with $c_a=h\big(f_a,(n-1)e^+_a\big)$. For $y\in\hat L$, $h(f_a,y)$ lies in the $\hat{\mathcal
O}$-span $\hat{\mathcal O}\mathfrak a_a$ of $\mathfrak a_a$: this holds on
$L_0=\sum_b\mathfrak a_be_b+L_b+\sum_b\mathfrak a_b^{-1}f_b$, as $h(f_a,e_b)=\delta_{ab}$ and
$h(f_a,L_b)=h(f_a,f_b)=0$, and $\hat L$ is the $\hat{\mathcal O}$-span of $L_0$. If $a=r$, then $e^+_r=z_u\in
L_b\subset L_0$, so $(n-1)z_u\in\hat L$ and $c_r\in\hat{\mathcal O}\mathfrak a_r=\hat{\mathcal
O}\mathfrak d^{-1}$ (Definition~\ref{asm-division-kernel-finite-series-function-inputs-level}(1)). If $a<r$,
then for $t\in\mathfrak a_{a+1}$, $te_{a+1}\in L_0$, so $tc_a=h\big(f_a,(n-1)(te_{a+1})\big)\in\hat{\mathcal
O}\mathfrak a_a$. Writing $1=\sum_jt_js_j$ with $t_j\in\mathfrak a_{a+1}$ and $s_j\in\mathfrak a_{a+1}^{-1}$,
$c_a=\sum_js_j(t_jc_a)\in\hat{\mathcal O}\mathfrak a_{a+1}^{-1}\mathfrak a_a=\hat{\mathcal O}\mathfrak d^{-1}$.
By Lemma~\ref{asm-division-kernel-finite-series-function-inputs-trace}, (2) for $E/F$ and then (1) for
$\psi=\psi_{E,\mathrm f}$, $\psi$ is $1$ on $\hat{\mathcal O}\mathfrak d^{-1}$. So $\psi(B)=1$, and
$\psi(A-B)=\psi(A)\psi(-B)=1$.

\emph{Step 4.} Put $\Delta=A-B$, and choose integers $\alpha,\beta$ with $\alpha C^2+\beta g=1$. Then
$\Delta/g=\alpha\big(C^2\Delta/g\big)+\beta\Delta$, so $\psi(\Delta/g)=1$, as $\psi(C^2\Delta/g)=1$ and
$\psi(\Delta)=1$. Hence
$\chi_2(\kappa)=\psi(A/g)=\psi(B/g)\psi(\Delta/g)=\psi_{\mathrm f}(n)$.

(3) Let $\gamma\in\Gamma$ with $\gamma_{\mathrm f}\in S_{\mathrm f}$, and let $G$ be the isometry of $V$ that
defines $\gamma$, so that the components of $\gamma$ are the extensions of $G$ to the $V_v$ and to $\hat V$.
The vectors $(\gamma_{\mathrm f}-1)e_a$ and $(\gamma_{\mathrm f}-1)z_u$ are the images in $\hat V$ of
$(G-1)e_a$ and $(G-1)z_u$, and lie in $U_{a-1}$ and $U_r$. With $X_m$ and $X$ as in
Lemma~\ref{asm-division-kernel-finite-series-function-inputs-levi},
Lemma~\ref{asm-division-kernel-finite-series-function-inputs-transfer}(1) gives $(G-1)e_a\in X_{a-1}$ and
$(G-1)z_u\in X$. Let $H$ and $N=GH^{-1}$ be as in
Lemma~\ref{asm-division-kernel-finite-series-function-inputs-levi}, and let $\eta$ and $\nu=\gamma\eta^{-1}$ be
the rational points defined by $H$ and $N$. The components of $\eta$ fix every $e_a$, $f_a$ and $z_u$, as $H$
does, so $\eta\in\Gamma^H_u$. At every infinite place $v$, and at the finite places, $(\nu-1)e_a$ and
$(\nu-1)z_u$ are the images of $(N-1)e_a\in X_{a-1}$ and $(N-1)z_u\in X$, so they lie in $U_{a-1}$ and $U_r$;
and $(\nu-1)y\in U_r$ for every $y$ orthogonal to every $e_a$, by
Lemma~\ref{asm-division-kernel-finite-series-function-inputs-transfer}(2), applied to $V\to V_v$ and to
$V\to\hat V$, whose images span. So $\nu_\infty\in N_\infty$, $\nu_{\mathrm f}\in N_{\mathrm f}$ and
$\gamma=\nu\eta$.
% Lean: (2) unipotent_in_level = unipotent_in_level_of_psi (psiDd_holds X). The membership is
% unipotent_conj_mem_level (Lemma -inputs-membership); Step 1 is Z_term (with xw_diff, hA_fmat, Z_left,
% Z_right, C_f_Lf), Step 2 is Z_bw (cong_b, bw_dual, Z_sig, mL_Qs), then unitaryCongruence.psi_Z; Step 3 is
% unitaryCongruence.Afun_int, psi_B with hA_f_nxt, coord_dd (hA_f_Lf, hermitian_f_lattice_mem_ideal,
% Iad_mul, the ideal identity aaD (a+1)^-1 * aaD a = ddS^-1) and psi_Iad, with psiDd_holds =
% psiDd_of_cores (Lemma -inputs-trace); Step 4 is psi_bez. chi2 is chi2f (flagLevelCharacter_eq_psi_divN,
% Afun) and psi_f is psiFin with divN. (3) Level.ratLevi_holds: descend and descend_range (Lemma
% -inputs-transfer(1)) give SRat for the matrix gamma0 of gamma; H = HGL hS with HGL_mem (Lemma
% -inputs-levi); eta = ratGA X H hHU, which lies in Gam by Hm_e, Hm_f and levi_fixes_z; nu = gamma *
% eta^-1, whose components lie in NinfU and NfinU by transfer_span and transfer3 (Lemma
% -inputs-transfer(2)) with rvHom, rv_span, hJ_rv at the infinite places and inclHom, incl_span, hA_incl
% at the finite ones. ratLevi_holds never uses its hypothesis hgammaI (gamma_inf in SinfU).
\end{proof}

\begin{lemma}\label{asm-division-kernel-finite-series-function}
In Situation~\ref{asm-division-kernel-finite-situation}, let $u$ be a word and $T$ a Poincaré datum of
$u$, with the notation of Definitions~\ref{asm-division-kernel-finite-series-poincare}
and~\ref{asm-division-kernel-finite-series-function-bessel}. Then the value of
$\Phi_T$ does not depend on the choice of the decomposition, $\Phi_T$ is measurable, and for every
$y\in G_{\mathbf A}$:
\begin{enumerate}
\item $\Phi_T(\gamma y)=\Phi_T(y)$ for $\gamma\in\Gamma_u$;
\item $\Phi_T(y(1,\kappa))=\chi_2(\kappa)\Phi_T(y)$ for $\kappa\in K_2$;
\item $z\mapsto\varphi(z)\Phi_T(y(z,1))$ is integrable on $G_\infty$, with integral $\Phi_T(y)$.
\end{enumerate}
\end{lemma}
\leanref{Formalism/LocalToGlobal/SiegelLatticeTheory/HaarTransversals/UnitaryHaarProducts/InHHFinal.lean}{Formalism.L2.SG.PoincareInHH.Cut.phProps_closed}
\leanref{Formalism/LocalToGlobal/SiegelLatticeTheory/HaarTransversals/UnitaryHaarProducts/InHHFinal.lean}{Formalism.L2.SG.PoincareInHH.Cut.poincare_well_defined}
\leanref{Formalism/LocalToGlobal/SiegelLatticeTheory/HaarTransversals/PoincareInHH.lean}{Formalism.L2.SG.PoincareInHH.PoincareSetting.PhProps}
\leanref{Formalism/LocalToGlobal/SiegelLatticeTheory/HaarTransversals/PoincareInHH.lean}{Formalism.L2.SG.PoincareInHH.PoincareSetting.PhWellDef}
\begin{proof}
Write $\Phi=\Phi_T$. Call $(n,h,k)$ a \emph{decomposition} of $x\in G_{\mathrm f}$ if $n\in N_{\mathrm f}$,
$h\in H_{\mathrm f}$, $k\in K_2$ and $x=nhx_uk$. For $h\in H_{\mathrm f}$ let $\ell_h(\xi)=T\xi(1,h)$, a
linear form on $V_{\pi_u}$.

\emph{Step 1: if $n\in N_{\mathrm f}$, $h\in H_{\mathrm f}$ and $x_u^{-1}nhx_u\in K_2$, then
$x_u^{-1}hx_u\in K_2$ and $\chi_2(x_u^{-1}nhx_u)=\psi_{\mathrm f}(n)\chi_2(x_u^{-1}hx_u)$.} As $x_u$ and
$K_2$ lie in $K_{\mathrm{set}}$, $nh\in K_{\mathrm{set}}$. By
Lemma~\ref{asm-division-kernel-finite-series-function-inputs}(1), $h\in K_{\mathrm{set}}$, so
$n=(nh)h^{-1}\in K_{\mathrm{set}}$. By Lemma~\ref{asm-division-kernel-finite-series-function-inputs}(2),
$x_u^{-1}nx_u\in K_2$ and $\chi_2(x_u^{-1}nx_u)=\psi_{\mathrm f}(n)$. So
$x_u^{-1}hx_u=(x_u^{-1}nx_u)^{-1}(x_u^{-1}nhx_u)\in K_2$, and
$\chi_2(x_u^{-1}nhx_u)=\chi_2(x_u^{-1}nx_u)\chi_2(x_u^{-1}hx_u)$.

\emph{Step 2: independence of the decomposition.} Let $(n,h,k)$ and $(n',h',k')$ be decompositions of
$y_{\mathrm f}$. By Lemma~\ref{asm-division-kernel-finite-series-function-unipotent}(2) and (3),
$m=h'^{-1}n'^{-1}nh'$ lies in $N_{\mathrm f}$, with $\psi_{\mathrm f}(m)=\psi_{\mathrm
f}(n'^{-1})\psi_{\mathrm f}(n)$. From $nhx_uk=n'h'x_uk'$,
\[
 x_u^{-1}\,m\,(h'^{-1}h)\,x_u=k'k^{-1}\in K_2.
\]
By Step 1, applied to $m$ and $h'^{-1}h$, the element $h'^{-1}h$ lies in $K_u$ and
$\chi_2(k'k^{-1})=\psi_{\mathrm f}(m)\chi_2(x_u^{-1}h'^{-1}hx_u)$. The second property of $T$, at
$p=(1,h')$ with $\kappa=h'^{-1}h$, gives $\ell_h=\chi_2(x_u^{-1}h'^{-1}hx_u)\ell_{h'}$. By
Lemma~\ref{asm-division-kernel-finite-series-function-unipotent}(2),
$\psi_{\mathrm f}(n')\psi_{\mathrm f}(n'^{-1})\psi_{\mathrm f}(n)=\psi_{\mathrm f}(1)\psi_{\mathrm
f}(n)=\psi_{\mathrm f}(n)$; and $\chi_2(k')=\chi_2(k'k^{-1})\chi_2(k)$. Hence
$\psi_{\mathrm f}(n')\chi_2(k')=\psi_{\mathrm f}(n)\chi_2(k)\chi_2(x_u^{-1}h'^{-1}hx_u)$, and
\[
 \psi_{\mathrm f}(n)\chi_2(k)\ell_h=\psi_{\mathrm f}(n')\chi_2(k')\ell_{h'}.
\]
Evaluated at $W_u(y_\infty)$, this is the independence. So
$\Phi(y)=\psi_{\mathrm f}(n)\chi_2(k)\ell_h(W_u(y_\infty))$ for every decomposition $(n,h,k)$ of
$y_{\mathrm f}$, and $\Phi(y)=0$ if there is none.

\emph{Step 3: property (2).} Let $\kappa\in K_2$. If $(n,h,k)$ is a decomposition of $y_{\mathrm f}$, then
$(n,h,k\kappa)$ is one of $y_{\mathrm f}\kappa$, and $\chi_2(k\kappa)=\chi_2(k)\chi_2(\kappa)$. If
$y_{\mathrm f}$ has none, neither has $y_{\mathrm f}\kappa$, since a decomposition $(n,h,k')$ of
$y_{\mathrm f}\kappa$ gives the decomposition $(n,h,k'\kappa^{-1})$ of $y_{\mathrm f}$; then both sides
are $0$.

\emph{Step 4: $\Phi$ is continuous, hence measurable.} By
Lemma~\ref{asm-division-kernel-finite-folding} the level subgroup is normalized, so
$0<\mathrm{vol}(K_2)<\infty$. As $e=\overline{\chi_2}/\mathrm{vol}(K_2)$ on $K_2$ and $e=0$ elsewhere, and
$|\chi_2|=1$, $K_2$ is the set where $e\ne0$, and $\chi_2(\kappa)=\mathrm{vol}(K_2)\overline{e(\kappa)}$ for
$\kappa\in K_2$. As $e$ is locally constant, $K_2$ is open and $e$ is continuous. Fix $y$, and let $N$ be
the open set of $y'$ with $\kappa(y')=y_{\mathrm f}^{-1}y'_{\mathrm f}\in K_2$, which contains $y$. For
$y'\in N$, $y'=(y'_\infty,y_{\mathrm f})(1,\kappa(y'))$. If $(n,h,k)$ is a decomposition of $y_{\mathrm f}$,
Steps 2 and 3 give, for $y'\in N$,
\[
 \Phi(y')=\mathrm{vol}(K_2)\,\overline{e(\kappa(y'))}\,\psi_{\mathrm f}(n)\chi_2(k)\,\ell_h(W_u(y'_\infty)),
\]
a continuous function of $y'$, as $W_u$ is continuous and $\ell_h$ is a linear form on a
finite-dimensional space. If $y_{\mathrm f}$ has no decomposition, then by Step 3 no $y'_{\mathrm f}$ with
$y'\in N$ has one, and $\Phi=0$ on $N$. So $\Phi$ is continuous at $y$. By
Lemma~\ref{asm-division-kernel-finite-series-function-group}(1) the $\sigma$-algebra of $G_{\mathbf A}$ is
its Borel $\sigma$-algebra, so $\Phi$ is measurable.

\emph{Step 5: property (1).} Let $\gamma\in\Gamma_u$; by
Lemma~\ref{asm-division-kernel-finite-series-function-unipotent}(1), $\gamma_\infty\in S_\infty$ and
$\gamma_{\mathrm f}\in S_{\mathrm f}$. By Lemma~\ref{asm-division-kernel-finite-series-function-inputs}(3),
$\gamma=\nu\eta$ with $\nu\in\Gamma$, $\nu_\infty\in N_\infty$, $\nu_{\mathrm f}\in N_{\mathrm f}$ and
$\eta\in\Gamma^H_u$. If $(n,h,k)$ is a decomposition of $y_{\mathrm f}$, then
\[
 (\gamma y)_{\mathrm f}=\nu_{\mathrm f}\eta_{\mathrm f}nhx_uk
 =\big(\nu_{\mathrm f}\cdot\eta_{\mathrm f}n\eta_{\mathrm f}^{-1}\big)(\eta_{\mathrm f}h)\,x_uk,
\]
and by Lemma~\ref{asm-division-kernel-finite-series-function-unipotent}(2) and (3) this is a decomposition,
with $\psi_{\mathrm f}(\nu_{\mathrm f}\cdot\eta_{\mathrm f}n\eta_{\mathrm f}^{-1})=\psi_{\mathrm
f}(\nu_{\mathrm f})\psi_{\mathrm f}(n)$. As $\Gamma_u$ is a group, $y_{\mathrm f}$ has a decomposition if
and only if $(\gamma y)_{\mathrm f}$ has one; if neither has, $\Phi(\gamma y)=0=\Phi(y)$. Suppose $(n,h,k)$
is a decomposition of $y_{\mathrm f}$. As $\gamma_\infty=\nu_\infty\eta_\infty\in S_\infty$, with
$\nu_\infty\in N_\infty\subset S_\infty$ and $\eta_\infty\in H_\infty\subset S_\infty$
(Lemma~\ref{asm-division-kernel-finite-series-function-unipotent}(4)),
\[
 W_u(\gamma_\infty y_\infty)=\psi_\infty(\nu_\infty\eta_\infty)\pi_u(\nu_\infty)\pi_u(\eta_\infty)W_u(y_\infty)
 =\psi_\infty(\nu_\infty)\pi_u(\eta_\infty)W_u(y_\infty),
\]
by Lemma~\ref{asm-division-kernel-finite-series-function-unipotent}(4) and $\pi_u(\nu_\infty)=1$. By the
third and the first properties of $T$,
$\ell_{\eta_{\mathrm f}h}(\pi_u(\eta_\infty)\xi)=T\xi\big((1,\eta_{\mathrm f}h)(\eta_\infty,1)\big)
=T\xi\big((\eta_\infty,\eta_{\mathrm f})(1,h)\big)=\ell_h(\xi)$. Hence
\[
 \Phi(\gamma y)=\psi_{\mathrm f}(\nu_{\mathrm f})\psi_{\mathrm f}(n)\chi_2(k)\psi_\infty(\nu_\infty)
 \ell_h(W_u(y_\infty))=\psi_\infty(\nu_\infty)\psi_{\mathrm f}(\nu_{\mathrm f})\,\Phi(y)=\Phi(y)
\]
by Lemma~\ref{asm-division-kernel-finite-series-function-inputs}(4).

\emph{Step 6: property (3).} If $(n,h,k)$ is a decomposition of $y_{\mathrm f}$, it is one of
$(y(z,1))_{\mathrm f}=y_{\mathrm f}$ for every $z$, so $\Phi(y(z,1))=c\,\ell_h(W_u(y_\infty z))$ with
$c=\psi_{\mathrm f}(n)\chi_2(k)$. As $z\mapsto\varphi(z)W_u(y_\infty z)$ is integrable with integral
$W_u(y_\infty)$ and $\ell_h$ is a continuous linear form, $z\mapsto\varphi(z)\Phi(y(z,1))
=c\,\ell_h(\varphi(z)W_u(y_\infty z))$ is integrable with integral $c\,\ell_h(W_u(y_\infty))=\Phi(y)$. If
$y_{\mathrm f}$ has no decomposition, $\Phi(y(z,1))=0=\Phi(y)$ for every $z$.
% Lean: poincare_well_defined = phWellDef_of_steps = phWellDef_of X u hK nfinGroup_holds nfinConj_holds
% (levelDec_of X u hK hS hN): Step 1 is levelDec_of (kset_of_conj, hS = ksetSplit_holds, hN =
% unipotent_in_level), Step 2 is phWellDef_of (col_equiv, the linear_combination with hpn and hp1).
% phProps_closed = poincare_properties = phProps_of_steps = phProps_of with ph_right_K (Step 3),
% phMeas_of = ph_continuous (level_subgroup_isClopen, level_character_eq_conj_idempotent, with
% level_subgroup_properties and FoldReadyX) (Step 4), phLeftInv_of (sratComp_holds, hL = ratLevi_holds,
% hP = psiProd_holds, psiInf_mul_H, X.Wu_spec, X.piU_rep, T.2) (Step 5) and ph_Rph with wuRph_of_spec
% (evL.integral_comp_comm) (Step 6). hK = Y.hK is the level subgroup character of the setting. Size:
% about 680 lines (proof closure, text count) without the inputs of Lemmas -series-function-group and
% -series-function-inputs and without level_subgroup_properties.
\end{proof}

\begin{lemma}\label{asm-division-kernel-finite-series-fold-multiplicity}
Let $G$ be a second countable Hausdorff topological group with its Borel $\sigma$-algebra, let
$\Gamma\subset G$ be a discrete subgroup, and let $\Omega\subset G$ be a compact set containing $1$. For
$x,y\in G$ put $C(x,y)=\{\gamma\in\Gamma:x^{-1}\gamma y\in\Omega\}$, and let
$m(x)=\sup_{y\in G}\#C(x,y)$ in $[0,\infty]$, the \emph{multiplicity} of $\Omega$ at $x$. Then
$1\le m(x)<\infty$ and $m(\gamma x)=m(x)$ for all $x\in G$ and $\gamma\in\Gamma$, and $m$ is
measurable.
\end{lemma}
\leanref{Formalism/LocalToGlobal/SiegelLatticeTheory/HaarTransversals/FoldBound.lean}{Formalism.L2.SG.FoldBound.mult}
\leanref{Formalism/PeriodIntegrals/FoldMult.lean}{PI.FoldMult.FoldMult}
\leanref{Formalism/PeriodIntegrals/OrbitalBoundsConclusion/FoldMult.lean}{overlapMultiplicityProperties}
\leanref{Formalism/PeriodIntegrals/UnitaryGroupMassFormula/MultiplicityMeasurability/MeasurableMult.lean}{PI.FoldMult.UnitaryGroupMassFormula.cnt_encard_le}
\leanref{Formalism/PeriodIntegrals/UnitaryGroupMassFormula/MultiplicityMeasurability/MeasurableMult.lean}{PI.FoldMult.UnitaryGroupMassFormula.mult_mul_left}
\leanref{Formalism/PeriodIntegrals/UnitaryGroupMassFormula/MultiplicityMeasurability/MeasurableMult.lean}{PI.FoldMult.UnitaryGroupMassFormula.isClosed_hitSet}
\leanref{Formalism/PeriodIntegrals/UnitaryGroupMassFormula/MultiplicityMeasurability/MeasurableMult.lean}{PI.FoldMult.UnitaryGroupMassFormula.mult_eq_iSup}
\leanref{Formalism/PeriodIntegrals/UnitaryGroupMassFormula/MultiplicityMeasurability/MeasurableMult.lean}{PI.FoldMult.UnitaryGroupMassFormula.measurable_mult}
\begin{proof}
\emph{Lower bound.} $x^{-1}\cdot1\cdot x=1\in\Omega$, so $1\in C(x,x)$ and $m(x)\ge1$.

\emph{Finiteness.} The set $K=\{x\omega\omega'^{-1}x^{-1}:\omega,\omega'\in\Omega\}$ is compact, as the
image of $\Omega\times\Omega$ under a continuous map. A discrete subgroup of a Hausdorff topological
group is closed, so $\Gamma\cap K$ is compact and discrete, hence finite, with $N$ elements say. Let
$y\in G$ with $C(x,y)$ nonempty, and fix $\gamma_0\in C(x,y)$. For $\gamma\in C(x,y)$,
\[
 \gamma\gamma_0^{-1}=x\,(x^{-1}\gamma y)(x^{-1}\gamma_0y)^{-1}x^{-1}\in\Gamma\cap K,
\]
so $C(x,y)\subset(\Gamma\cap K)\gamma_0$ has at most $N$ elements. Hence $m(x)\le N$.

\emph{Invariance.} Let $\gamma\in\Gamma$. As $(\gamma x)^{-1}\delta y=x^{-1}(\gamma^{-1}\delta)y$, the
map $\delta\mapsto\gamma^{-1}\delta$ is a bijection of $C(\gamma x,y)$ onto $C(x,y)$ for every $y$. So
$m(\gamma x)=m(x)$.

\emph{Measurability.} For a finite set $A\subset\Gamma$ let $H_A$ be the set of $x\in G$ for which some
$y\in G$ has $x^{-1}ay\in\Omega$ for every $a\in A$. Then
\[
 m(x)=\sup_A\,\#A\cdot\mathbf 1_{H_A}(x),
\]
the supremum over the finite subsets $A$ of $\Gamma$. Indeed, if $x\in H_A$ through $y$, then
$A\subset C(x,y)$, so $\#A\le m(x)$; and each $C(x,y)$ is finite by the finiteness, with $x\in
H_{C(x,y)}$. Each $H_A$ is closed. If $A$ is empty, $H_A=G$. Otherwise fix $a_0\in A$. Writing
$y=a_0^{-1}x\omega$, with $\omega=x^{-1}a_0y\in\Omega$, shows that $H_A$ is the image under the first
projection of the set of $(x,\omega)\in G\times\Omega$ with $x^{-1}aa_0^{-1}x\omega\in\Omega$ for every
$a\in A$. That set is closed, as $\Omega$ is closed, and the projection $G\times\Omega\to G$ is a closed
map, as $\Omega$ is compact. Finally, $\Gamma$ is separable, as $G$ is second countable, and discrete, so
it is countable. So $\Gamma$ has countably many finite subsets, and $m$ is a countable supremum of
measurable functions.
\end{proof}

\begin{lemma}\label{asm-division-kernel-finite-series-fold-unfold}
Let $G$ be a group with a $\sigma$-algebra for which left and right multiplication by each element are
measurable, and with a left-invariant measure $\mu$. Let $\Gamma\subset G$ be a countable subgroup and
$S\subset\Gamma$ a subgroup. For a subgroup $H$ of $G$, a \emph{fundamental domain} of $H$ for $\mu$ is a
set $F\subset G$, measurable up to a $\mu$-null set, such that for $\mu$-almost every $y\in G$ there is
$h\in H$ with $hy\in F$, and $\mu(hF\cap h'F)=0$ for $h\ne h'$ in $H$. Let $F$ be a fundamental domain of
$S$ and $F_\Gamma$ one of $\Gamma$, and let $D\subset\Gamma$ be a set such that every $\gamma\in\Gamma$ is
$s\delta$ for exactly one pair $(s,\delta)\in S\times D$. Let $g\colon G\to[0,\infty]$ be measurable with
$g(sy)=g(y)$ for $s\in S$ and $y\in G$. Then
\[
 \int_{F_\Gamma}\sum_{\delta\in D}g(\delta x)\,d\mu(x)=\int_Fg\,d\mu.
\]
\end{lemma}
\leanref{Formalism/PeriodIntegrals/FoldUnfold.lean}{PI.FoldUnfold.FoldUnfold}
\leanref{Formalism/PeriodIntegrals/OrbitalBoundsConclusion/FoldUnfold.lean}{lintegralTransversalSum}
\leanref{Formalism/PeriodIntegrals/UnitaryGroupMassFormula/UnfoldingIdentity/FoldUnfoldMain.lean}{PI.FoldUnfold.UnitaryGroupMassFormula.foldUnfold_main}
\leanref{Formalism/PeriodIntegrals/UnitaryGroupMassFormula/UnfoldingIdentity/FoldUnfoldMain.lean}{PI.FoldUnfold.UnitaryGroupMassFormula.count_ae_one}
\leanref{Formalism/LocalToGlobal/SiegelLatticeTheory/HaarTransversals/FoldBound.lean}{Formalism.L2.SG.FoldBound.IsRightTransversal}
\begin{proof}
The groups $S$ and $D$ are countable, as subsets of $\Gamma$.

\emph{Step 1: for $\mu$-almost every $y$, $\sum_{s\in S}\mathbf 1_F(sy)=1$.} Almost every $y$ has some
$s\in S$ with $sy\in F$. For $s\ne t$ in $S$, the $y$ with $sy\in F$ and $ty\in F$ form
$s^{-1}F\cap t^{-1}F$, a null set. As $S$ is countable, almost every $y$ lies in none of these sets and
has some $s$ with $sy\in F$; for such $y$ exactly one term of the sum is $1$ and the others are $0$.

\emph{Step 2.} Left translations preserve $\mu$, and $D$ is countable. So by Step 1, almost every
$x\in G$ has $\sum_{s\in S}\mathbf 1_F(s\delta x)=1$ for every $\delta\in D$. For such $x$, as
$(s,\delta)\mapsto s\delta$ is a bijection of $S\times D$ onto $\Gamma$ and $g(s\delta x)=g(\delta x)$,
\[
 \sum_{\gamma\in\Gamma}\mathbf 1_F(\gamma x)g(\gamma x)
 =\sum_{\delta\in D}g(\delta x)\sum_{s\in S}\mathbf 1_F(s\delta x)=\sum_{\delta\in D}g(\delta x).
\]

\emph{Step 3.} As $F_\Gamma$ is a fundamental domain of $\Gamma$ and $\mathbf 1_Fg$ is non-negative and
measurable up to a null set,
\[
 \int_Fg\,d\mu=\int_G\mathbf 1_Fg\,d\mu=\sum_{\gamma\in\Gamma}\int_{F_\Gamma}(\mathbf 1_Fg)(\gamma x)\,d\mu(x)
 =\int_{F_\Gamma}\sum_{\gamma\in\Gamma}(\mathbf 1_Fg)(\gamma x)\,d\mu(x),
\]
the last by monotone convergence. By Step 2 the integrand is $\sum_{\delta\in D}g(\delta x)$ for almost
every $x$.
\end{proof}

\begin{lemma}\label{asm-division-kernel-finite-series-fold-bound}
Let $G$ be a locally compact, second countable topological group with its Borel $\sigma$-algebra and a
Haar measure $\mu$, let $\Gamma\subset G$ be a discrete subgroup and $S\subset\Gamma$ a subgroup, let
$\Omega\subset G$ be measurable and $a\ge0$, and let $f\colon G\to\mathbf C$ be measurable with
$f(sy)=f(y)$ and $|f(y)|\le a\int_\Omega|f(y\omega)|\,d\mu(\omega)$ for $s\in S$ and $y\in G$. Let $F$ be
a fundamental domain of $S$ for $\mu$ (Lemma~\ref{asm-division-kernel-finite-series-fold-unfold}), and
$D\subset\Gamma$ a set such that every $\gamma\in\Gamma$ is $s\delta$ for exactly one pair
$(s,\delta)\in S\times D$. For $x\in G$ let $m(x)=\sup_{y\in G}\#\{\gamma\in\Gamma:x^{-1}\gamma
y\in\Omega\}$ in $[0,\infty]$. Then, for every $x\in G$,
\[
 \sum_{\delta\in D}|f(\delta x)|\le a\,m(x)\int_F|f|\,d\mu
\]
in $[0,\infty]$. (The local compactness of $G$ is not used.)
\end{lemma}
\leanref{Formalism/LocalToGlobal/SiegelLatticeTheory/HaarTransversals/FoldBound.lean}{Formalism.L2.SG.FoldBound.FoldBound}
\leanref{Formalism/LocalToGlobal/SiegelLatticeTheory/HaarTransversals/FoldBound.lean}{Formalism.L2.SG.FoldBound.foldBound_holds}
\leanref{Formalism/LocalToGlobal/SiegelLatticeTheory/HaarTransversals/FoldBound.lean}{Formalism.L2.SG.FoldBound.tsum_transversal}
\leanref{Formalism/LocalToGlobal/SiegelLatticeTheory/HaarTransversals/FoldBound.lean}{Formalism.L2.SG.FoldBound.quotientLIntegral}
\begin{proof}
As $G$ is second countable, $\Gamma$ is separable; being discrete, it is countable, and so are $S$ and
$D$. Fix $x\in G$.

\emph{Step 1.} For $g\in G$, substituting $y=g\omega$ in the bound for $f$, which is allowed as $\mu$ is
left-invariant, gives
\[
 |f(g)|\le a\int_G\mathbf 1_\Omega(g^{-1}y)|f(y)|\,d\mu(y).
\]

\emph{Step 2.} Apply Step 1 to $g=\delta x$ and sum over $\delta\in D$. Exchanging the sum and the
integral by monotone convergence,
\[
 \sum_{\delta\in D}|f(\delta x)|\le a\int_Gh\,d\mu,\qquad
 h(y)=\sum_{\delta\in D}\mathbf 1_\Omega(x^{-1}\delta^{-1}y)|f(y)|.
\]

\emph{Step 3.} As $F$ is a fundamental domain of $S$ and $h$ is non-negative and measurable,
$\int_Gh\,d\mu=\int_F\sum_{s\in S}h(sy)\,d\mu(y)$.

\emph{Step 4.} The map $D\times S\to\Gamma$, $(\delta,s)\mapsto\delta^{-1}s$, is a bijection. If
$\delta^{-1}s=\delta'^{-1}s'$, then $s^{-1}\delta=s'^{-1}\delta'$, and the uniqueness of the
decomposition of this element of $\Gamma$ gives $s=s'$ and $\delta=\delta'$. If $\gamma\in\Gamma$, write
$\gamma^{-1}=s\delta$; then $\gamma=\delta^{-1}s^{-1}$.

\emph{Step 5.} Fix $y$. As $|f(sy)|=|f(y)|$, Step 4 gives
\[
 \sum_{s\in S}h(sy)=|f(y)|\sum_{s\in S}\sum_{\delta\in D}\mathbf 1_\Omega(x^{-1}\delta^{-1}sy)
 =|f(y)|\,\#\{\gamma\in\Gamma:x^{-1}\gamma y\in\Omega\}\le m(x)|f(y)|.
\]
With Steps 2 and 3, $\sum_{\delta\in D}|f(\delta x)|\le a\int_Fm(x)|f(y)|\,d\mu(y)=a\,m(x)\int_F|f|\,d\mu$.
\end{proof}

\begin{lemma}\label{asm-division-kernel-finite-series-fold}
Let $G$ be a locally compact, second countable Hausdorff group with a Haar measure $\mu$, let
$\Gamma\subset G$ be a discrete subgroup and $S\subset\Gamma$ a subgroup, let $\Omega\subset G$ be a
compact set containing $1$, let $a\ge0$, and let $f\colon G\to\mathbf C$ be measurable with $f(sy)=f(y)$
for $s\in S$, $y\in G$, and with \emph{folding data with constant $a$ on $\Omega$}:
\[
 |f(y)|\le a\int_\Omega|f(y\omega)|\,d\mu(\omega)\qquad(y\in G).
\]
Let $F_S$ be a fundamental domain of $S$ for $\mu$, and suppose that the \emph{multiplicity integral}
\[
 M=\int_{F_S}m(y)|f(y)|\,d\mu(y),\qquad m(y)=\sup_{x\in G}\#\{\gamma\in\Gamma:y^{-1}\gamma x\in\Omega\},
\]
is finite. Let $D\subset\Gamma$ be such that every $\gamma\in\Gamma$ is $s\delta$ for exactly one pair
$(s,\delta)\in S\times D$. Then $\sum_{\delta\in D}|f(\delta x)|<\infty$ for every $x\in G$, and for every
fundamental domain $F$ of $\Gamma$ for $\mu$ the function $P(x)=\sum_{\delta\in D}f(\delta x)$ is square
integrable on $F$, with
\[
 \int_F|P|^2\,d\mu\le aQM\le aM^2,\qquad Q=\int_{F_S}|f|\,d\mu.
\]
(The local compactness of $G$ is not used.)
\end{lemma}
\leanref{Formalism/LocalToGlobal/SiegelLatticeTheory/HaarTransversals/PoincareSeriesSquareIntegrable.lean}{Formalism.L2.SG.PoincareSeriesSquareIntegrable.PoincareSeriesFinite}
\leanref{Formalism/LocalToGlobal/SiegelLatticeTheory/HaarTransversals/PoincareSeriesSquareIntegrable.lean}{Formalism.L2.SG.PoincareSeriesSquareIntegrable.PoincareSeriesMemLp}
\leanref{Formalism/LocalToGlobal/SiegelLatticeTheory/ClosedVolumeStatements/LevelSubgroupProperties.lean}{Formalism.L2.SG.Closed.fold_L2_first}
\leanref{Formalism/LocalToGlobal/SiegelLatticeTheory/ClosedVolumeStatements/LevelSubgroupProperties.lean}{Formalism.L2.SG.Closed.fold_L2_second}
\begin{proof}
Fundamental domains are meant as in Lemma~\ref{asm-division-kernel-finite-series-fold-unfold}. As
$G$ is second countable, $\Gamma$ is separable; being discrete, it is countable, and so is $D$. By
Lemma~\ref{asm-division-kernel-finite-series-fold-multiplicity}, $1\le m(y)<\infty$ and
$m(\gamma y)=m(y)$ for $y\in G$ and $\gamma\in\Gamma$, and $m$ is measurable. As $m\ge1$, $Q\le M<\infty$.

\emph{Convergence.} The set $\Omega$ is compact, hence closed and measurable. By
Lemma~\ref{asm-division-kernel-finite-series-fold-bound}, for every $x\in G$,
\[
 \sum_{\delta\in D}|f(\delta x)|\le a\,m(x)\,Q<\infty.
\]

\emph{The bound.} Put $g(y)=m(y)|f(y)|$, a measurable function with $g(sy)=g(y)$ for $s\in S$, as
$S\subset\Gamma$. For $x\in G$ and $\delta\in D\subset\Gamma$, $m(\delta x)=m(x)$. So
\[
 |P(x)|^2\le\Big(\sum_{\delta\in D}|f(\delta x)|\Big)^2
 \le\Big(\sum_{\delta\in D}|f(\delta x)|\Big)\,a\,m(x)\,Q=aQ\sum_{\delta\in D}g(\delta x).
\]
By Lemma~\ref{asm-division-kernel-finite-series-fold-unfold}, for $\Gamma$, $S$, $F_S$, $F$, $D$ and $g$,
the integral of the last sum over $F$ is $\int_{F_S}g\,d\mu=M$. So $\int_F|P|^2\,d\mu\le aQM$, and
$aQM\le aM^2$ as $Q\le M$.

\emph{Square integrability.} The series defining $P$ converges absolutely at every point, so $P$ is the
pointwise limit of its partial sums over the finite subsets of $D$. These are measurable, as $f$ and
left translations are, so $P$ is measurable. With $\int_F|P|^2\,d\mu\le aQM<\infty$, $P$ is square
integrable on $F$.
% Lean: fold_L2_first = poincareSeries_finite = PI.PoincareSeriesFinite.UnitaryGroupMassFormula.
% poincareSeriesFinite: foldBound_holds (Lemma -series-fold-bound), mult_lt_top and
% quotientLIntegral_le_multiplicityLIntegral (one_le_mult). fold_L2_second = fold_l2_bound_b =
% PI.PoincareSeriesMemLp.UnitaryGroupMassFormula.fold_L2_double_coset: overlapMultiplicityProperties
% (Lemma -series-fold-multiplicity) gives hm1, hmtop, hminv, hmmeas; hT is foldBound_holds; hpt is the
% pointwise bound with g = mult * |Ph|; hunf is lintegralTransversalSum (Lemma -series-fold-unfold);
% hbound and hlast are the two inequalities; measurable_poinc (measurable_of_tendsto_metrizable' along
% the finite subsets of D) gives the measurability. Size: about 430 lines (proof closure, text count).
\end{proof}

\begin{lemma}\label{asm-division-kernel-finite-series-hecke-representative}
Let $G$ be a group with a $\sigma$-algebra for which left and right multiplication by each element are
measurable, and with a left- and right-invariant measure $\mu$. Let $\Gamma\subset G$ be a countable
subgroup and $F\subset G$ a measurable set such that $\Gamma y$ meets $F$ for every $y\in G$ and
$\mu(\gamma F\cap F)=0$ for $\gamma\in\Gamma$, $\gamma\ne1$. Let $N$ be the set of $y\in G$ for which
$\Gamma y\cap F$ has at least two points.
\begin{enumerate}
\item $N$ is measurable and $\mu(N)=0$.
\item There is a measurable map $r\colon G\to G$ with $r(y)\in\Gamma y\cap F$ for every $y\in G$.
\item Let $\rho\colon G\to G$ send each $y$ to $\gamma_yy$, for a chosen $\gamma_y\in\Gamma$ with
  $\gamma_yy\in F$. Let $r\colon G\to G$ be a map with $r(y)\in\Gamma y\cap F$ for every $y$, and let
  $g\in G$. Then $\rho(xg)=r(xg)$ for almost every $x\in G$, and
  \[
   \mu(\{x\in F:r(xg)\in A\})=\mu(A\cap F)
  \]
  for every measurable $A\subset G$. The map $x\mapsto\rho(xg)$ agrees almost everywhere on $F$ with a
  measurable map, and it carries $\mu$ restricted to $F$ to $\mu$ restricted to $F$. If moreover $\mu$ is
  $\sigma$-finite and $(x,y)\mapsto xy$ and $x\mapsto x^{-1}$ are measurable, then
  \[
   \int_Fh(r(xg))\,d\mu(x)=\int_Fh\,d\mu
  \]
  for every measurable $h\colon G\to[0,\infty]$. (These three further hypotheses are not used.)
\end{enumerate}
\end{lemma}
\leanref{Formalism/SharpRestriction/CuspSpanNormBounds/UnimodularSchurMeasures/RrExistsMeasRep.lean}{Formalism.MainNotions.LeviCoordinates.rr_FD}
\leanref{Formalism/SharpRestriction/CuspSpanNormBounds/UnimodularSchurMeasures/RrExistsMeasRep.lean}{Formalism.MainNotions.LeviCoordinates.rr_rep}
\leanref{Formalism/SharpRestriction/CuspSpanNormBounds/UnimodularSchurMeasures/RrExistsMeasRep.lean}{Formalism.MainNotions.LeviCoordinates.rr_N_null}
\leanref{Formalism/SharpRestriction/CuspSpanNormBounds/UnimodularSchurMeasures/RrExistsMeasRep.lean}{Formalism.MainNotions.LeviCoordinates.rr_exists_meas_rep}
\leanref{Formalism/SharpRestriction/CuspSpanNormBounds/UnimodularSchurMeasures/RrExistsMeasRep.lean}{Formalism.MainNotions.LeviCoordinates.rr_rep_ae_eq}
\leanref{Formalism/SharpRestriction/CuspSpanNormBounds/UnimodularSchurMeasures/RrExistsMeasRep.lean}{Formalism.MainNotions.LeviCoordinates.rr_meas_pre}
\leanref{Formalism/SharpRestriction/CuspSpanNormBounds/UnimodularSchurMeasures/RrExistsMeasRep.lean}{Formalism.MainNotions.LeviCoordinates.rr_map}
\leanref{Formalism/SharpRestriction/CuspSpanNormBounds/UnimodularSchurMeasures/RrkbLeOneAddSq.lean}{Formalism.MainNotions.LeviCoordinates.rrkb_lint}
\begin{proof}
(1) The set $N$ is the union over $a\ne b$ in $\Gamma$ of the sets $a^{-1}F\cap b^{-1}F$, which are
measurable; there are countably many. If $y$ lies in $a^{-1}F\cap b^{-1}F$, then
$ay=(ab^{-1})(by)\in ab^{-1}F\cap F$. So $a^{-1}F\cap b^{-1}F\subset a^{-1}(ab^{-1}F\cap F)$, which is
null by left invariance, as $ab^{-1}\ne1$.

(2) Enumerate $\Gamma=\{\gamma_0,\gamma_1,\dots\}$. For $y\in G$ let $n(y)$ be the least $n$ with
$\gamma_ny\in F$, and put $r(y)=\gamma_{n(y)}y$. The set where $n(y)=n$ is the set of $y$ with
$\gamma_ny\in F$ and $\gamma_my\notin F$ for $m<n$, which is measurable; on it $r$ is left multiplication
by $\gamma_n$. So $r$ is measurable.

(3) By (1) and right invariance, $Ng^{-1}$ is null, so $xg\notin N$ for almost every $x$; then
$\Gamma xg\cap F$ is one point, which is both $\rho(xg)$ and $r(xg)$.

Let $r_0$ be measurable as in (2), $A$ measurable, and $S$ the set of $x$ with $xg\notin N$ and
$r_0(xg)\in A$, a measurable set. As $\Gamma cy=\Gamma y$ for $c\in\Gamma$, whether $y\in N$, and for
$y\notin N$ the point $r_0(y)$ of $\Gamma y\cap F$, depend only on $\Gamma y$. So $\gamma S=S$ for
$\gamma\in\Gamma$. The set
$F'=\{x:xg\in F\}$ has the same two properties as $F$, by right invariance. For two such sets $F_1,F_2$
and a measurable set $S$ with $\Gamma S=S$, $\mu(S\cap F_1)=\mu(S\cap F_2)$: as $\Gamma$ is countable
and the $\gamma F_2$ cover $G$ and meet in null sets,
\[
 \mu(S\cap F_1)=\sum_{\gamma\in\Gamma}\mu(S\cap F_1\cap\gamma F_2)
 =\sum_{\gamma\in\Gamma}\mu(S\cap\gamma^{-1}F_1\cap F_2)=\mu(S\cap F_2),
\]
using left invariance in the middle. So $\mu(S\cap F)=\mu(S\cap F')$. For almost every $x$, $xg\notin
N$; then $r(xg)=r_0(xg)$, and if moreover $xg\in F$, then $r_0(xg)=xg$. So, up to null sets,
$\{x\in F:r(xg)\in A\}$ is $S\cap F$, and $S\cap F'$ is $\{x:xg\in A\cap F\}$, whose measure is
$\mu(A\cap F)$ by right invariance. This gives the identity of measures.

By the first claim, applied to $r_0$, $\rho(xg)=r_0(xg)$ for almost every $x$, and $x\mapsto r_0(xg)$ is
measurable. The identity of measures, applied to $\rho$ in place of $r$, says that the image of $\mu$
restricted to $F$ under $x\mapsto\rho(xg)$ is $\mu$ restricted to $F$. Finally, by the first claim,
$h(r(xg))=h(\rho(xg))$ for almost every $x$, so
$\int_Fh(r(xg))\,d\mu(x)=\int_Fh(\rho(xg))\,d\mu(x)=\int_Fh\,d\mu$, the last by the image of the measure.
\end{proof}

\begin{lemma}\label{asm-division-kernel-finite-series-hecke-fix}
Let $G$ be a group with a $\sigma$-algebra for which $(x,y)\mapsto xy$ and $x\mapsto x^{-1}$ are
measurable, with a $\sigma$-finite, left- and right-invariant measure $\mu$. Let $\Gamma$ and $F$ be as in
Lemma~\ref{asm-division-kernel-finite-series-hecke-representative}, and let $\rho\colon G\to G$ be the
map of Lemma~\ref{asm-division-kernel-finite-series-hecke-representative}(3), with $\rho(y)$ a chosen
point of $\Gamma y\cap F$. Let $R$ assign to every $g\in G$ a bounded operator $R(g)$ on
$L^2(F)=L^2(F,\mu)$ such that $(R(g)f)(x)=f(\rho(xg))$ for almost every $x\in F$, for every $f\in L^2(F)$.
Let $H$ be a set with a $\sigma$-algebra and a $\sigma$-finite measure $\nu$, let $\iota\colon H\to G$ be
measurable and $\phi\colon H\to\mathbf C$ integrable, and let $T$ be a bounded operator on $L^2(F)$ with
\[
 \langle u,Tv\rangle=\int_H\phi(k)\langle u,R(\iota(k))v\rangle\,d\nu(k)\qquad(u,v\in L^2(F)).
\]
Let $f\colon G\to\mathbf C$ be measurable, square integrable on $F$, with $f(\gamma y)=f(y)$ for
$\gamma\in\Gamma$ and $y\in G$, and with $\int_H\phi(k)f(y\iota(k))\,d\nu(k)=f(y)$ for almost every
$y\in G$. Then $Tf=f$ in $L^2(F)$.
\end{lemma}
\leanref{Formalism/LocalToGlobal/SiegelLatticeTheory/HaarTransversals/UnitaryHaarProducts/InHHLeaf.lean}{Formalism.L2.SG.Pf.InHHLeaf.fix_of}
\leanref{Formalism/LocalToGlobal/SiegelLatticeTheory/HaarTransversals/UnitaryHaarProducts/InHHLeaf.lean}{Formalism.L2.SG.Pf.InHHLeaf.step_sub}
\leanref{Formalism/SharpRestriction/CuspSpanNormBounds/UnimodularSchurMeasures/RrExistsMeasRep.lean}{Formalism.MainNotions.LeviCoordinates.rr_IsReg}
\begin{proof}
The function $\phi$ agrees almost everywhere with a measurable function, and replacing it changes
neither hypothesis; so assume $\phi$ measurable. Let $r$ be measurable as in
Lemma~\ref{asm-division-kernel-finite-series-hecke-representative}(2). As $f$ is invariant under
$\Gamma$, $f(r(z))=f(z)$ for every $z$.

\emph{Step 1: for almost every $x\in F$, $(Tf)(x)=\int_H\phi(k)f(r(x\iota(k)))\,d\nu(k)$.} The function
$\Phi(k,x)=\phi(k)f(r(x\iota(k)))$ on $H\times G$ is measurable. Let $s\subset F$ be measurable with
$\mu(s)<\infty$. As $|t|\le1+|t|^2$, Lemma~\ref{asm-division-kernel-finite-series-hecke-representative}(3)
gives, for every $k$,
\[
 \int_s|f(r(x\iota(k)))|\,d\mu(x)\le\mu(s)+\int_F|f(r(x\iota(k)))|^2\,d\mu(x)=\mu(s)+\int_F|f|^2\,d\mu=C,
\]
so $\int_H\int_s|\Phi|\,d\mu\,d\nu\le C\int_H|\phi|\,d\nu<\infty$ and $\Phi$ is integrable on $H\times s$.
For $g\in G$, $\rho(xg)=r(xg)$ for almost every $x$, and the class of $f$ in $L^2(F)$ is $f$ almost
everywhere on $F$, a property kept under $x\mapsto\rho(xg)$, which preserves $\mu$ on $F$
(Lemma~\ref{asm-division-kernel-finite-series-hecke-representative}(3)). So
$\langle\mathbf 1_s,R(g)f\rangle=\int_sf(r(xg))\,d\mu(x)$. With the hypothesis on $T$ and Fubini's
theorem,
\[
 \int_sTf\,d\mu=\langle\mathbf 1_s,Tf\rangle=\int_H\phi(k)\int_sf(r(x\iota(k)))\,d\mu(x)\,d\nu(k)
 =\int_s\int_H\Phi(k,x)\,d\nu(k)\,d\mu(x).
\]
Both functions of $x$ are integrable on every such $s$ and have the same integrals there; as $\mu$ is
$\sigma$-finite, they agree almost everywhere on $F$.

\emph{Step 2.} For almost every $x$, $\int_H\phi(k)f(r(x\iota(k)))\,d\nu(k)=\int_H\phi(k)f(x\iota(k))\,d\nu(k)
=f(x)$. With Step 1, $Tf=f$ almost everywhere on $F$.
\end{proof}

\begin{lemma}\label{asm-division-kernel-finite-series-hecke-union}
Let $G$, $\mu$, $\Gamma$ and $F$ be as in Lemma~\ref{asm-division-kernel-finite-series-hecke-fix}. Let
$S\subset\Gamma$ be a subgroup, $F_S$ a fundamental domain of $S$ for $\mu$
(Lemma~\ref{asm-division-kernel-finite-series-fold-unfold}), and $D\subset\Gamma$ a set such that every
$\gamma\in\Gamma$ is $s\delta$ for exactly one pair $(s,\delta)\in S\times D$. Let $f\colon G\to[0,\infty]$
be measurable with $f(sy)=f(y)$ for $s\in S$ and $y\in G$.
\begin{enumerate}
\item $\sum_{\delta\in D}f(\delta\gamma x)=\sum_{\delta\in D}f(\delta x)$ for $\gamma\in\Gamma$ and $x\in G$.
\item $\int_F\sum_{\delta\in D}f(\delta y)\,d\mu(y)=\int_{F_S}f\,d\mu$.
\end{enumerate}
\end{lemma}
\leanref{Formalism/LocalToGlobal/SiegelLatticeTheory/HaarTransversals/UnitaryHaarProducts/InHHLeaf.lean}{Formalism.L2.SG.Pf.InHHLeaf.poincE_left_inv}
\leanref{Formalism/LocalToGlobal/SiegelLatticeTheory/HaarTransversals/UnitaryHaarProducts/InHHLeaf.lean}{Formalism.L2.SG.Pf.InHHLeaf.unfold_lint}
\leanref{Formalism/LocalToGlobal/SiegelLatticeTheory/HaarTransversals/UnitaryHaarProducts/InHHLeaf.lean}{Formalism.L2.SG.Pf.InHHLeaf.rr_FD_union}
\begin{proof}
(1) For $\delta\in D$ write $\delta\gamma=s_\delta\pi(\delta)$ with $s_\delta\in S$ and $\pi(\delta)\in D$. The
map $\pi$ is injective: if $\pi(\delta_1)=\pi(\delta_2)$, then
$\delta_1=(s_{\delta_1}s_{\delta_2}^{-1})\delta_2$, and the uniqueness of the decomposition of
$\delta_1=1\cdot\delta_1$ gives $\delta_1=\delta_2$. It is surjective: for $d\in D$ write $d\gamma^{-1}=s\delta$;
then $\delta\gamma=s^{-1}d$ and $\pi(\delta)=d$. So
$\sum_\delta f(\delta\gamma x)=\sum_\delta f(s_\delta\pi(\delta)x)=\sum_\delta f(\pi(\delta)x)=\sum_df(dx)$.

(2) The set $D$ is countable. Let $U=\bigcup_{\delta\in D}\delta F$, a measurable set. Every $y\in G$ has
$s\in S$ with $sy\in U$: take $\gamma\in\Gamma$ with $\gamma y\in F$ and write $\gamma^{-1}=s'\delta$; then
$\delta\gamma y=s'^{-1}y$ lies in $\delta F$. For $\delta,\delta'\in D$ and $s\in S$, $s\delta F\cap\delta'F$ lies in
$\delta'\big((\delta'^{-1}s\delta)F\cap F\big)$, which is null unless $\delta'^{-1}s\delta=1$; and
$s\delta=\delta'=1\cdot\delta'$ forces $s=1$ and $\delta=\delta'$. So $sU\cap U$ is null for $s\ne1$, the sets
$\delta F$ are pairwise almost disjoint, and $U$ is a fundamental domain of $S$, the translates $sU$ being
almost disjoint by left invariance. An $S$-invariant function has the same integral over two fundamental
domains of $S$, by the computation in the proof of
Lemma~\ref{asm-division-kernel-finite-series-hecke-representative}(3). Hence
\[
 \int_{F_S}f\,d\mu=\int_Uf\,d\mu=\sum_{\delta\in D}\int_{\delta F}f\,d\mu=\sum_{\delta\in D}\int_Ff(\delta y)\,d\mu(y)
 =\int_F\sum_{\delta\in D}f(\delta y)\,d\mu(y).
\]
\end{proof}

\begin{lemma}\label{asm-division-kernel-finite-series-hecke}
In Situation~\ref{asm-division-kernel-finite-situation}, with the notation of
Definition~\ref{asm-division-kernel-finite-series-poincare}, recall that $\mathcal F$ is measurable, that
$\Gamma y$ meets $\mathcal F$ for every $y\in G_{\mathbf A}$, and that $\mu(\gamma\mathcal F\cap\mathcal F)=0$
for $\gamma\in\Gamma$, $\gamma\ne1$. Let $R$ be the right regular representation of the setting: for
$g\in G_{\mathbf A}$ and $f\in L^2(\mathcal F)$, $(R(g)f)(x)=f(\rho(xg))$ for almost every $x\in\mathcal F$,
where $\rho(y)$ is a chosen point of $\Gamma y\cap\mathcal F$, as in
Lemma~\ref{asm-division-kernel-finite-series-hecke-representative}(3). Let $R_\varphi$ and $R_e$ be the operators
of the setting on $L^2(\mathcal F)$, so that $\mathcal H$ is the image of $R_eR_\varphi$: the function $e$
is integrable, and for $u,v\in L^2(\mathcal F)$,
\[
 \langle u,R_\varphi v\rangle=\int_{G_\infty}\varphi(z)\langle u,R(z,1)v\rangle\,dz,\qquad
 \langle u,R_ev\rangle=\int_{G_{\mathrm f}}e(g)\langle u,R(1,g)v\rangle\,dg.
\]
\begin{enumerate}
\item $\mathcal F$ is a fundamental domain of $\Gamma$ for $\mu$.
\item Let $f$ be a measurable function on $G_{\mathbf A}$ with $f(\gamma y)=f(y)$ for $\gamma\in\Gamma$,
  square integrable on $\mathcal F$. If $\int_{G_\infty}\varphi(z)f(y(z,1))\,dz=f(y)$ for almost every
  $y$, and $\varphi$ is integrable, then $R_\varphi f=f$. If $\int_{G_{\mathrm f}}e(g)f(y(1,g))\,dg=f(y)$ for
  almost every $y$, then $R_ef=f$.
\item Suppose that $\varphi$ is integrable, that the Poincaré series of $u$ fold
  (Lemma~\ref{asm-division-kernel-finite-folding}), and that for every Poincaré datum $T$ of $u$ the
  function $\Phi_T$ is measurable and satisfies (1), (2) and (3) of
  Lemma~\ref{asm-division-kernel-finite-series-function}. For every Poincaré datum $T$ of $u$ and every
  set of representatives $D$, $\int_{G_\infty}\varphi(z)P_{T,D}(y(z,1))\,dz=P_{T,D}(y)$ for almost every
  $y$.
\end{enumerate}
\end{lemma}
\leanref{Formalism/LocalToGlobal/SiegelLatticeTheory/HaarTransversals/UnitaryHaarProducts/InHHLeaf.lean}{Formalism.L2.SG.Pf.fundBridge_holds}
\leanref{Formalism/LocalToGlobal/SiegelLatticeTheory/HaarTransversals/UnitaryHaarProducts/InHHLeaf.lean}{Formalism.L2.SG.Pf.regular_fixes_invariant_L2}
\leanref{Formalism/LocalToGlobal/SiegelLatticeTheory/HaarTransversals/UnitaryHaarProducts/InHHLeaf.lean}{Formalism.L2.SG.Pf.idempotent_fixes_invariant_L2}
\leanref{Formalism/LocalToGlobal/SiegelLatticeTheory/HaarTransversals/UnitaryHaarProducts/InHHLeaf.lean}{Formalism.L2.SG.Pf.pserRph_holds}
\begin{proof}
The measure $\mu=\mu_\infty\times\mu_{\mathrm f}$ is a product of Haar measures, so it is left invariant.

(1) Every $y$ has some $\gamma\in\Gamma$ with $\gamma y\in\mathcal F$. For $\gamma\ne\gamma'$ in $\Gamma$,
$\mu(\gamma\mathcal F\cap\gamma'\mathcal F)=\mu(\gamma'^{-1}\gamma\mathcal F\cap\mathcal F)=0$ by left
invariance. So $\mathcal F$ is a fundamental domain of $\Gamma$ in the sense of
Lemma~\ref{asm-division-kernel-finite-series-fold-unfold}.

(2) By Lemma~\ref{asm-division-kernel-finite-series-function-group}, $G_{\mathbf A}$ is second
countable with its Borel $\sigma$-algebra, so multiplication and inversion are measurable; $\mu$ is
$\sigma$-finite and right invariant, as its factors are; and $\Gamma$ is countable. So
Lemma~\ref{asm-division-kernel-finite-series-hecke-fix} applies to $G_{\mathbf A}$, $\mu$, $\Gamma$,
$\mathcal F$, $\rho$ and $R$. For $R_\varphi$ take $H=G_\infty$, $\nu=\mu_\infty$, $\iota(z)=(z,1)$,
$\phi=\varphi$ and $T=R_\varphi$. For $R_e$ take $H=G_{\mathrm f}$, $\nu=\mu_{\mathrm f}$,
$\iota(g)=(1,g)$, $\phi=e$, which is integrable, and $T=R_e$.

(3) Use the facts listed in (2). Fix $T$ and $D$, and let $\Gamma_u$, its fundamental domain $F_u$, the
compact set $\Omega_0\ni1$ and the constant $a$ be given by the folding. Assume, as in the proof of
Lemma~\ref{asm-division-kernel-finite-series-hecke-fix}, that $\varphi$ is measurable. Let $r$ be
measurable as in Lemma~\ref{asm-division-kernel-finite-series-hecke-representative}(2) for $\mathcal F$.

\emph{Step 1.} Put $f=|\Phi_T|$ and $\Sigma(w)=\sum_{\delta\in D}f(\delta w)$, with values in $[0,\infty]$.
The function $f$ is measurable and invariant under $\Gamma_u\subset\Gamma$, by (1) of
Lemma~\ref{asm-division-kernel-finite-series-function}, so $\Sigma$ is measurable, and
$\Sigma(\gamma w)=\Sigma(w)$ for $\gamma\in\Gamma$ by
Lemma~\ref{asm-division-kernel-finite-series-hecke-union}(1). The multiplicity
$m(y)=\sup_x\#\{\gamma\in\Gamma:y^{-1}\gamma x\in\Omega_0\}$ is at least $1$, as $\gamma=1$ and $x=y$
qualify. So, by Lemma~\ref{asm-division-kernel-finite-series-hecke-union}(2),
\[
 \int_{\mathcal F}\Sigma\,d\mu=\int_{F_u}f\,d\mu\le\int_{F_u}mf\,d\mu<\infty,
\]
the last integral being the multiplicity integral over $\Omega_0$ and $F_u$, finite by the folding.

\emph{Step 2.} Put $J(y)=\int_{G_\infty}|\varphi(z)|\,\Sigma(y(z,1))\,dz$, a measurable function of $y$ with
$J(\gamma y)=J(y)$ for $\gamma\in\Gamma$. As $r(w)\in\Gamma w$, $\Sigma(r(w))=\Sigma(w)$. By Tonelli's theorem
and Lemma~\ref{asm-division-kernel-finite-series-hecke-representative}(3),
\[
 \int_{\mathcal F}J\,d\mu=\int_{G_\infty}|\varphi(z)|\int_{\mathcal F}\Sigma(r(y(z,1)))\,d\mu(y)\,dz
 =\int_{G_\infty}|\varphi|\,dz\cdot\int_{\mathcal F}\Sigma\,d\mu<\infty.
\]
So $J<\infty$ almost everywhere on $\mathcal F$. The set $B$ where $J=\infty$ is contained in
$\bigcup_{\gamma\in\Gamma}\gamma^{-1}(B\cap\mathcal F)$, as $J$ is invariant under $\Gamma$ and $\Gamma y$
meets $\mathcal F$; this is a countable union of null sets. So $J<\infty$ almost everywhere.

\emph{Step 3.} Let $y$ have $J(y)<\infty$. By (3) of Lemma~\ref{asm-division-kernel-finite-series-function}
at $\delta y$, the function $z\mapsto\varphi(z)\Phi_T(\delta y(z,1))$ is integrable with integral
$\Phi_T(\delta y)$, for every $\delta\in D$. The sum over $\delta$ of the integrals of their absolute values
is $J(y)<\infty$. So the sum over $\delta$ and the integral can be exchanged, and
\[
 \int_{G_\infty}\varphi(z)P_{T,D}(y(z,1))\,dz=\sum_{\delta\in D}\int_{G_\infty}\varphi(z)\Phi_T(\delta y(z,1))\,dz
 =\sum_{\delta\in D}\Phi_T(\delta y)=P_{T,D}(y).
\]
% Lean: (1) fundBridge_holds (IsFundamentalDomain.mk'' from X.B.fund_spec, with muA_isHaar). (2)
% idempotent_fixes_invariant_L2 and regular_fixes_invariant_L2 are InHHLeaf.fix_of (Lemma
% -series-hecke-fix) with rr_FD_of X.B.fund_spec, rr_IsReg_of X.B.R_spec, iota = (1, g) or (y, 1), and
% X.B.Re_spec or X.Rph_spec (with PhInt); the instances come from gf_secondCountable, muf_right
% (GfUnimod_holds), rr_muA_sigmaFinite, rr_muA_left and rr_muA_right (Lemma -series-function-group).
% (3) pserRph_holds: f, F (= Sigma), hFinv (poincE_left_inv), hC (unfold_lint and one_le_mult), J, hJint
% (lintegral_lintegral_swap, rrkb_lint), hJae, hJinv, hJall, and integral_tsum. Size: about 710 lines
% (proof closure, text count) without the facts of Lemma -series-function-group.
\end{proof}

\begin{lemma}\label{asm-division-kernel-finite-series}
In Situation~\ref{asm-division-kernel-finite-situation}, let $u$ be a word, and suppose that the
archimedean matrix coefficient is integrable and that the Poincaré series of $u$ fold (as in
Lemma~\ref{asm-division-kernel-finite-folding}). Then for every Poincaré datum $T$ of $u$ the
Poincaré series of $T$ converges absolutely, does not depend on the choice of representatives, is
square integrable on the fundamental domain and lies in $\mathcal H$. The Poincaré series depends
linearly on $T$.
\end{lemma}
\leanref{Formalism/LocalToGlobal/SiegelLatticeTheory/ClosedVolumeStatements/UnipotentInLevelSubgroup.lean}{Formalism.L2.SG.Closed.InHHC_holds}
\leanref{Formalism/LocalToGlobal/SiegelLatticeTheory/HaarTransversals/UnitaryHaarProducts/InHHCut.lean}{Formalism.L2.SG.PoincareInHH.Cut.puLinear_holds}
\leanref{Formalism/LocalToGlobal/SiegelLatticeTheory/HaarTransversals/UnitaryHaarProducts/InHHCut.lean}{Formalism.L2.SG.PoincareInHH.Cut.inHH_of}
\leanref{Formalism/LocalToGlobal/SiegelLatticeTheory/HaarTransversals/UnitaryHaarProducts/InHHCut.lean}{Formalism.L2.SG.PoincareInHH.Cut.poinc_indep}
\leanref{Formalism/LocalToGlobal/SiegelLatticeTheory/HaarTransversals/UnitaryHaarProducts/InHHCut.lean}{Formalism.L2.SG.PoincareInHH.Cut.pser_Re}
\begin{proof}
Use the notation of Definition~\ref{asm-division-kernel-finite-series-poincare}. By the folding, $\Gamma_u$
has a fundamental domain $F_u$, and there are a compact set $\Omega_0\ni1$ and $a\ge0$ such that for every
Poincaré datum $T$, $\Phi_T$ has folding data with constant $a$ on $\Omega_0$ and its multiplicity
integral over $\Omega_0$ and $F_u$ is finite. The group $\Gamma$ is discrete in $G_{\mathbf A}$, and
$\Gamma_u\subset\Gamma$. Fix $T$ and a set of representatives $D$.

\emph{Convergence.} By Lemma~\ref{asm-division-kernel-finite-series-function}, $\Phi_T$ is measurable
and invariant under $\Gamma_u$ on the left. Lemma~\ref{asm-division-kernel-finite-series-fold}, for
$G_{\mathbf A}$, $\Gamma$, $S=\Gamma_u$, $\Omega_0$, $a$, $f=\Phi_T$ and $F_S=F_u$, gives
$\sum_{\delta\in D}|\Phi_T(\delta x)|<\infty$ for every $x$.

\emph{Independence.} Let $D'$ be another set of representatives. For $\delta\in D$ write
$\delta=s\delta'$ with $s\in\Gamma_u$ and $\delta'\in D'$; this is a bijection $D\to D'$, and
$\Phi_T(\delta x)=\Phi_T(s\delta'x)=\Phi_T(\delta'x)$ by
Lemma~\ref{asm-division-kernel-finite-series-function}(1). So $P_{T,D}=P_{T,D'}$.

\emph{Square integrability.} By Lemma~\ref{asm-division-kernel-finite-series-hecke}(1), $\mathcal F$ is a
fundamental domain of $\Gamma$, so Lemma~\ref{asm-division-kernel-finite-series-fold} shows that $P_{T,D}$
is square integrable on $\mathcal F$.

\emph{Membership in $\mathcal H$.} The group $\Gamma$ is countable, so $D$ is, and $P=P_{T,D}$ is the
pointwise limit of the finite partial sums, which are measurable; so $P$ is measurable. For
$\gamma\in\Gamma$, $D\gamma$ is a set of representatives, so $P(\gamma y)=\sum_{\delta\in
D}\Phi_T(\delta\gamma y)=P_{T,D\gamma}(y)=P(y)$ by the independence; $P$ is invariant under $\Gamma$ on
the left. By Lemma~\ref{asm-division-kernel-finite-series-hecke}(3),
$\int\varphi(z)P(y(z,1))\,dz=P(y)$ for almost every $y$, so $R_\varphi P=P$ by
Lemma~\ref{asm-division-kernel-finite-series-hecke}(2). For $\kappa\in K_2$,
$P(y(1,\kappa))=\chi_2(\kappa)P(y)$ by Lemma~\ref{asm-division-kernel-finite-series-function}(2), term
by term. As $|\chi_2(\kappa)|=1$ and $0<\mathrm{vol}(K_2)<\infty$, the level subgroup being normalized
(Lemma~\ref{asm-division-kernel-finite-folding}),
\[
 \int_{G_{\mathrm f}}e(g)P(y(1,g))\,dg=\int_{K_2}\frac{\overline{\chi_2(g)}\chi_2(g)}{\mathrm{vol}(K_2)}P(y)\,dg
 =P(y)
\]
for every $y$, so $R_eP=P$ by Lemma~\ref{asm-division-kernel-finite-series-hecke}(2). Hence
$P=R_eR_\varphi P$ lies in $\mathcal H$.

\emph{Linearity.} The function $\Phi_T$ is linear in $T$. As the series converge absolutely,
$P_{T+T',D}=P_{T,D}+P_{T',D}$ and $P_{cT,D}=cP_{T,D}$ pointwise, so also in $L^2(\mathcal F)$.
% Lean: Closed.InHHC_holds = Pf.inHHC_holds = inHHC_holds_of'' with ksetSplit_holds, unipotent_in_level,
% ratLevi_holds and psiProd_holds; it is inHH_of X u Y.hK (foldReadyX_holds X) fold_L2_first
% fold_L2_second (fundBridge_holds X) hOK (phProps_closed ...) (poincare_well_defined ...)
% (pserMeas_holds) (pserRph_holds) hI (regular_fixes_invariant_L2 X) (idempotent_fixes_invariant_L2 X)
% (level_subgroup_properties X) (level_subgroup_measurable ...). In inHH_of: hsum from fold_L2_first
% (Convergence), poinc_indep with transversal_bij (Independence), hmem from fold_L2_second
% (Square integrability), hmeas from pserMeas_holds, hinv from poinc_left_inv, h1' from RF and pserRph,
% h2' from EF and pser_Re (pser_right_K, X.e_eq) (Membership). Linearity is puLinear_holds (ph_add,
% ph_smul, Summable.tsum_add, tsum_mul_left, MemLp.toLp_add). The Lean's InHH states the membership as
% X.B.Re (X.Rph (PuFun ...)) = PuFun .... The measurability is pserMeas_holds (countable rr_GF,
% measurable_of_tendsto_metrizable' over finite partial sums); the left invariance is poinc_left_inv.
\end{proof}

\subsection{Finite spectral data and tied families}\label{subsection-asm-finite-spectral-data}

\begin{situation}\label{asm-division-kernel-finite-dimension-situation}
Let $n\ge1$, let $\epsilon\in\{\pm1\}^n$ have $p$ entries $+1$ and $q$ entries $-1$, and put
$q_0=\min(p,q)$. Let $J$ be the diagonal matrix with entries $\epsilon$,
$G=U(J)=\{g\in GL_n(\mathbf C):g^*Jg=J\}$ with a Haar measure $\nu$, $K=G\cap U(n)$, and
$\mathfrak p$ the space of Hermitian $X\in M_n(\mathbf C)$ with $XJ=-JX$. For a complex matrix $g$
write $|g|$ for its Hilbert--Schmidt norm, $|g|^2=\sum_{i,j}|g_{ij}|^2$, and $\|g\|$ for its operator
norm.
\end{situation}
\leanref{Formalism/RealPlaceAnalysis/Complexify.lean}{Formalism.REAL.UJ}
\leanref{ReciprocityTheory/ReductiveGroups/DiscreteSeries/DefsRep.lean}{ReciprocityTheory.ChartRestriction.DS.inP}

\begin{lemma}\label{asm-division-kernel-finite-dimension-decay-ktype-highest}
Let $n\ge1$ and $\epsilon\in\{\pm1\}^n$, and let $H$ be a complex inner product space with operators
$E_{ab}$, for the pairs $(a,b)$ with $\epsilon_a=\epsilon_b$, such that
\[
 E_{ab}E_{cd}-E_{cd}E_{ab}=\delta_{bc}E_{ad}-\delta_{da}E_{cb},\qquad
 \langle E_{ab}x,y\rangle=\langle x,E_{ba}y\rangle
\]
for all such pairs $(a,b)$ and $(c,d)$ and all $x,y\in H$. A vector $x$ has \emph{weight}
$\mu\in\mathbf R^n$ if $E_{aa}x=\mu_ax$ for every $a$; it is a \emph{highest weight vector} of weight
$\mu$ if moreover $E_{ab}x=0$ whenever $a<b$ and $\epsilon_a=\epsilon_b$. For $x\in H$ let $H_x$ be the
smallest subspace of $H$ that contains $x$ and is stable under every $E_{ab}$.
\begin{enumerate}
\item Every nonzero finite-dimensional subspace stable under every $E_{ab}$ contains a nonzero highest
  weight vector.
\item Let $x$ be a highest weight vector of weight $\mu$, and $y\in H_x$ a highest weight vector of
  weight $\nu$. Then $y=0$, or $\nu=\mu$ and $y\in\mathbf Cx$.
\item Let $x$ be a highest weight vector with $H_x$ finite-dimensional. Then the only subspaces of
  $H_x$ stable under every $E_{ab}$ are $0$ and $H_x$.
\end{enumerate}
\end{lemma}
\leanref{Formalism/RealPlaceAnalysis/QuotientSettingHolds.lean}{Formalism.REAL.kIrr}
\leanref{Formalism/RealPlaceAnalysis/StatementKT.lean}{Formalism.REAL.KIrr}
\leanref{Formalism/RealPlaceAnalysis/RealPlaceTaylor/IsCompactInvariantCompactSpan.lean}{Formalism.REAL.kIrr_of}
\leanref{Formalism/RealPlaceAnalysis/RealPlaceTaylor/IsCompactInvariantCompactSpan.lean}{Formalism.REAL.highestWeightExists_holds}
\leanref{Formalism/RealPlaceAnalysis/RealPlaceTaylor/IsCompactInvariantCompactSpan.lean}{Formalism.REAL.kHiUniq_of}
\leanref{Formalism/RealPlaceAnalysis/RealPlaceTaylor/CompactCasimirHighestWeight.lean}{Formalism.REAL.compact_lowering_holds}
\leanref{Formalism/RealPlaceAnalysis/RealPlaceTaylor/CompactCasimirHighestWeight.lean}{Formalism.REAL.kOrth_of}
\leanref{Formalism/RealPlaceAnalysis/RealPlaceTaylor/CompactCasimirHighestWeight.lean}{Formalism.REAL.compact_weight_orthogonal}
\begin{proof}
\emph{Weights.} If $y$ has weight $\nu$, then $E_{cd}y$ has weight $\nu+e_c-e_d$, as the relation gives
$E_{aa}E_{cd}-E_{cd}E_{aa}=(\delta_{ac}-\delta_{ad})E_{cd}$. Vectors $x$, $y$ of weights $\mu\ne\nu$ are
orthogonal: if $\mu_a\ne\nu_a$, then $\mu_a\langle x,y\rangle=\langle E_{aa}x,y\rangle=\langle
x,E_{aa}y\rangle=\nu_a\langle x,y\rangle$.

(1) Let $S$ be such a subspace. The $E_{aa}$ commute with each other and preserve $S$, so they have a
common eigenvector $v\ne0$ in $S$; its eigenvalues are real, as $\langle E_{aa}v,v\rangle=\langle
v,E_{aa}v\rangle$. So $S$ contains nonzero weight vectors, and their weights form a finite set, since
each coordinate of a weight is an eigenvalue of $E_{aa}$ on $S$. Choose one, $\nu$, with the largest
value of $-\sum_a(a-1)\nu_a$, and a nonzero $v\in S$ of weight $\nu$. If $E_{ab}v\ne0$ for some $a<b$
with $\epsilon_a=\epsilon_b$, then $E_{ab}v\in S$ is a nonzero vector of weight $\nu+e_a-e_b$, whose
value is larger by $b-a>0$. So $v$ is a highest weight vector.

(2) Let $x$ be a highest weight vector of weight $\mu$. A \emph{lowering word} is a product
$w=E_{c_1d_1}\cdots E_{c_kd_k}$ with $k\ge0$, $c_i>d_i$ and $\epsilon_{c_i}=\epsilon_{d_i}$; let $L_x$ be
the span of the $wx$. Then $L_x=H_x$. Clearly $L_x\subset H_x$, and it suffices to show that
$E_{ab}wx\in L_x$ for every lowering word $w$, by induction on its length. If $a>b$, $E_{ab}w$ is a
lowering word. If $a=b$, $wx$ is a weight vector, so $E_{aa}wx\in\mathbf Cwx$. Let $a<b$. If $w$ is
empty, $E_{ab}x=0$. Otherwise $w=E_{cd}w'$ with $c>d$, and
\[
 E_{ab}E_{cd}w'x=E_{cd}E_{ab}w'x+\delta_{bc}E_{ad}w'x-\delta_{da}E_{cb}w'x,
\]
whose terms lie in $L_x$ by induction, since $L_x$ is stable under $E_{cd}$.

Next, the vectors of weight $\mu$ in $L_x$ are the multiples of $x$. Let $T=\sum_a(a-1)E_{aa}$ and
$\phi_0=\sum_a(a-1)\mu_a$. For a lowering word $w$ of length $k\ge1$, $wx$ has weight
$\mu+\sum_i(e_{c_i}-e_{d_i})$, so it is an eigenvector of $T$ for the eigenvalue
$\phi_0+\sum_i(c_i-d_i)\ne\phi_0$. So every $y\in L_x$ is $cx+y'$ with $y'$ in the sum of the
eigenspaces of $T$ for eigenvalues other than $\phi_0$. If $y$ has weight $\mu$, then $y'=y-cx$ lies in
the eigenspace for $\phi_0$ as well, so $y'=0$.

Next, if $x$ and $y$ are highest weight vectors of weights $\mu\ne\nu$, then $H_x\perp H_y$. First
$x\perp H_y=L_y$: $x\perp y$, and $\langle x,E_{cd}wy\rangle=\langle E_{dc}x,wy\rangle=0$ for $c>d$.
Then, for a lowering word $E_{cd}w$ and $z\in H_y$, $\langle E_{cd}wx,z\rangle=\langle wx,E_{dc}z\rangle$
with $E_{dc}z\in H_y$, so $wx\perp H_y$ for every lowering word $w$ by induction, and $H_x=L_x\perp H_y$.

Now let $y\in H_x$ be a highest weight vector of weight $\nu$. If $\nu\ne\mu$, then $y\in H_x\cap H_y$
is orthogonal to itself, so $y=0$. If $\nu=\mu$, then $y\in L_x$ has weight $\mu$, so $y\in\mathbf Cx$.

(3) Let $S\subset H_x$ be nonzero and stable. By (1) it contains a nonzero highest weight vector $y$,
and by (2) $y=cx$ with $c\ne0$. So $x\in S$, and $H_x\subset S$.
% Lean: kIrr = kIrr_of kHiExist kHiUniq. (1) is highestWeightExists_holds (Module.End.exists_eigenvalue
% for the common eigenvector, eigenvalue_real, the finite set hfinW of weights and
% Set.exists_max_image for sum_i -(i) nu_i, with i counted from 0). (2) is kHiUniq_of kLow kWtOrth
% kOrth: kLow = compact_lowering_holds (compact_span_eq_lowering with lowering_span_stable and
% lowering_word_raises_in_span; lowering_span_weight_scalar with heightOperator and
% Module.End.eigenspaces_iSupIndep), kWtOrth = compact_weight_orthogonal, kOrth = kOrth_of. The weight
% shift is compact_weight_shift. The Lean states (3) with S a submodule of kSpan x and kSpan x
% finite-dimensional, and the relations as KRel and the adjointness as KUnit.
\end{proof}

\begin{lemma}\label{asm-division-kernel-finite-dimension-decay-ktype-derived}
In Situation~\ref{asm-division-kernel-finite-dimension-situation}, let $\mathfrak k$ be the space of
$X\in M_n(\mathbf C)$ with $X^*=-X$ and $XJ=JX$. Let $\tau$ be a continuous unitary representation of
$K$ on a finite-dimensional Hilbert space $T$, and $d\tau$ a complex-linear map from $M_n(\mathbf C)$
to the operators on $T$ such that for every $X\in\mathfrak k$ and every map $\gamma\colon\mathbf R\to
K$ with $\gamma(t)=\exp(tX)$, the derivative of $t\mapsto\tau(\gamma(t))$ at $0$ is $d\tau(X)$. Then:
\begin{enumerate}
\item every $k\in K$ is $\exp Y$ for some $Y\in\mathfrak k$, and $\tau(\exp X)=\exp(d\tau(X))$ for
  every $X\in\mathfrak k$;
\item the operators $d\tau(E_{ab})$, for $\epsilon_a=\epsilon_b$, satisfy the relations of
  Lemma~\ref{asm-division-kernel-finite-dimension-decay-ktype-highest};
\item a subspace of $T$ is stable under $\tau(K)$ if and only if it is stable under these
  $d\tau(E_{ab})$.
\end{enumerate}
\end{lemma}
\leanref{ReciprocityTheory/ReductiveGroups/DiscreteSeries/HighestWeights/OpEqExp.lean}{ReciprocityTheory.ChartRestriction.DS.HW.exists_log}
\leanref{Formalism/RealPlaceAnalysis/RealPlaceTaylor/CompactExponentialSurjective.lean}{Formalism.REAL.compact_exponential_surjective}
\leanref{ReciprocityTheory/ReductiveGroups/DiscreteSeries/HighestWeights/ModelGen.lean}{ReciprocityTheory.ChartRestriction.DS.HW.pre_kRel}
\leanref{ReciprocityTheory/ReductiveGroups/DiscreteSeries/HighestWeights/ModelGen.lean}{ReciprocityTheory.ChartRestriction.DS.HW.pre_kUnit}
\leanref{ReciprocityTheory/ReductiveGroups/DiscreteSeries/HighestWeights/ModelGen.lean}{ReciprocityTheory.ChartRestriction.DS.HW.pre_kStable_of_stable}
\leanref{ReciprocityTheory/ReductiveGroups/DiscreteSeries/HighestWeights/ModelGen.lean}{ReciprocityTheory.ChartRestriction.DS.HW.pre_kSpan_stable}
\begin{proof}
(1) Let $k\in K$. It is unitary, so normal with eigenvalues of absolute value $1$. Let $Y=\log k$,
formed by the functional calculus of $k$ with the principal logarithm. Then $\exp Y=k$; $Y^*=-Y$, as
$\log z$ is purely imaginary for $|z|=1$; and $Y$ commutes with $J$, as $k$ and $k^*$ do. So
$Y\in\mathfrak k$. For $X\in\mathfrak k$, $F(t)=\tau(\exp tX)$ is a continuous one-parameter group with
$F'(0)=d\tau(X)$, so $F'(t)=F(t)\,d\tau(X)$, and $F(t)\exp(-t\,d\tau(X))$ has derivative $0$; it is $1$.

(2) For $X\in\mathfrak k$, differentiating $\tau(\exp tX)^*\tau(\exp tX)=1$ gives
$d\tau(X)^*=-d\tau(X)$. For $k\in K$ and $Y\in\mathfrak k$, $k\exp(tY)k^{-1}=\exp(t\,kYk^{-1})$, so (1)
and differentiation give $\tau(k)\,d\tau(Y)\,\tau(k)^{-1}=d\tau(kYk^{-1})$; taking $k=\exp(sX)$ and
differentiating in $s$ gives $d\tau(X)d\tau(Y)-d\tau(Y)d\tau(X)=d\tau(XY-YX)$ for $X,Y\in\mathfrak k$.
For $\epsilon_a=\epsilon_b$, $E_{ab}=\frac12X_1-\frac i2X_2$ with $X_1=E_{ab}-E_{ba}$ and
$X_2=i(E_{ab}+E_{ba})$ in $\mathfrak k$. By complex linearity the commutator of $d\tau(E_{ab})$ and
$d\tau(E_{cd})$ is $d\tau(E_{ab}E_{cd}-E_{cd}E_{ab})=\delta_{bc}d\tau(E_{ad})-\delta_{da}d\tau(E_{cb})$,
and $d\tau(E_{ab})^*=-\frac12d\tau(X_1)-\frac i2d\tau(X_2)=d\tau(E_{ba})$.

(3) Let $W$ be stable under $\tau(K)$. It is closed, and for $X\in\mathfrak k$ and $v\in W$,
$d\tau(X)v=\lim_{t\to0}t^{-1}(\tau(\exp tX)v-v)\in W$; so $W$ is stable under
$d\tau(E_{ab})=\frac12d\tau(X_1)-\frac i2d\tau(X_2)$. Conversely, let $W$ be stable under the
$d\tau(E_{ab})$. Every $Y\in\mathfrak k$ commutes with $J$, so its entries $Y_{ab}$ with
$\epsilon_a\ne\epsilon_b$ vanish, and $W$ is stable under $d\tau(Y)$, hence under
$\exp(d\tau(Y))=\tau(\exp Y)$; by (1) it is stable under $\tau(K)$.
% Lean: (1) is exists_log (Formalism.REAL.kExp = compact_exponential_surjective, with cfc Complex.log k)
% and pre_op_curve (the product F t * E t is constant); (2) is pre_kRel (pre_dtau_bracket_single from
% pre_dtau_bracket, pre_conj_dtau and pre_star_dtau, with single_eq_kp) and pre_kUnit
% (pre_star_dtau_single); (3) is pre_kStable_of_stable (W.closed_of_finiteDimensional,
% hasDerivAt_iff_tendsto_slope) and pre_dtau_stable_of_kStable with pre_kSpan_stable
% (exp_mem_of_stable). Here dtau stands for the Lean's d-tau.
\end{proof}

\begin{lemma}\label{asm-division-kernel-finite-dimension-decay-ktype-model}
In Situation~\ref{asm-division-kernel-finite-dimension-situation}, let $\mathfrak k$ be as in
Lemma~\ref{asm-division-kernel-finite-dimension-decay-ktype-derived}, and let $m\in\mathbf Z^n$ satisfy
$m_a\ge0$ for every $a$ and $m_b\le m_a$ whenever $a<b$ and $\epsilon_a=\epsilon_b$. For each $a$ let
$\Delta_a(x)$ be the determinant of the square submatrix of $x\in M_n(\mathbf C)$ on the rows and
columns $b\le a$ with $\epsilon_b=\epsilon_a$; let $e_a=m_a-m_{a'}$ for the least $a'>a$ with
$\epsilon_{a'}=\epsilon_a$, and $e_a=m_a$ if there is none; and put $f_0(x)=\prod_a\Delta_a(x)^{e_a}$.
Let $V$ be the space of functions on $K$ spanned by the functions $k\mapsto f_0(kk')$ for $k'\in K$,
with the inner product $\langle f,g\rangle=\int_K\overline fg\,dk$ for a Haar measure $dk$ of $K$, and
let $\tau(k')f(k)=f(kk')$. Then:
\begin{enumerate}
\item $V$ is finite-dimensional, each $f\in V$ is the restriction to $K$ of a polynomial $F$ in the
  entries, and $\tau$ is a continuous unitary representation of $K$ on $V$;
\item the complex-linear map $d\tau$ with $d\tau(Z)f(k)=DF_k(kZ_K)$, where $Z_K$ keeps the entries
  $Z_{ab}$ with $\epsilon_a=\epsilon_b$ of $Z$ and replaces the others by $0$, satisfies the hypothesis
  of Lemma~\ref{asm-division-kernel-finite-dimension-decay-ktype-derived};
\item $w_0=f_0/\|f_0\|$ is a unit vector with $d\tau(E_{aa})w_0=m_aw_0$ for every $a$ and
  $d\tau(E_{ab})w_0=0$ whenever $a<b$ and $\epsilon_a=\epsilon_b$;
\item the only subspace of $V$ stable under $\tau(K)$ that contains $w_0$ is $V$.
\end{enumerate}
\end{lemma}
\leanref{ReciprocityTheory/ReductiveGroups/DiscreteSeries/HighestWeights/ModelTau.lean}{ReciprocityTheory.ChartRestriction.DS.HW.modelKType}
\leanref{ReciprocityTheory/ReductiveGroups/DiscreteSeries/HighestWeights/ModelFun.lean}{ReciprocityTheory.ChartRestriction.DS.HW.f0}
\leanref{ReciprocityTheory/ReductiveGroups/DiscreteSeries/HighestWeights/ModelFun.lean}{ReciprocityTheory.ChartRestriction.DS.HW.blkMinor}
\leanref{ReciprocityTheory/ReductiveGroups/DiscreteSeries/HighestWeights/ModelFun.lean}{ReciprocityTheory.ChartRestriction.DS.HW.mexp}
\leanref{ReciprocityTheory/ReductiveGroups/DiscreteSeries/HighestWeights/ModelSpace.lean}{ReciprocityTheory.ChartRestriction.DS.HW.Vsp}
\leanref{ReciprocityTheory/ReductiveGroups/DiscreteSeries/HighestWeights/ModelSpace.lean}{ReciprocityTheory.ChartRestriction.DS.HW.finiteDimensional_Vsp}
\leanref{ReciprocityTheory/ReductiveGroups/DiscreteSeries/HighestWeights/ModelTau.lean}{ReciprocityTheory.ChartRestriction.DS.HW.tauC_inner}
\leanref{ReciprocityTheory/ReductiveGroups/DiscreteSeries/HighestWeights/ModelFun.lean}{ReciprocityTheory.ChartRestriction.DS.HW.rdc_f0K_wt}
\leanref{ReciprocityTheory/ReductiveGroups/DiscreteSeries/HighestWeights/ModelFun.lean}{ReciprocityTheory.ChartRestriction.DS.HW.rdc_f0K_hw}
\leanref{ReciprocityTheory/ReductiveGroups/DiscreteSeries/HighestWeights/ModelTau.lean}{ReciprocityTheory.ChartRestriction.DS.HW.gen_wN}
\begin{proof}
(1) The entries of $kk'$ are linear in those of $k$, so each right translate of $f_0$ is the restriction
of a polynomial of degree at most that of $f_0$, and $V$ lies in the finite-dimensional space of such
restrictions. The functions in $V$ are continuous, and a Haar measure charges every nonempty open set,
so $\langle f,f\rangle=0$ only for $f=0$. As $K$ is compact, its Haar measure is right invariant, so
$\tau$ is unitary; it is continuous, as the translates depend continuously on $k'$.

(2) For $X\in\mathfrak k$ and $f=F|_K$, the derivative of $t\mapsto f(k\exp tX)=F(k\exp tX)$ at $0$ is
$DF_k(kX)$ by the chain rule, and $X_K=X$. So $d\tau$ is the complex-linear extension of the derivative,
and for $\epsilon_a=\epsilon_b$, $d\tau(E_{ab})f(k)$ is the derivative at $s=0$ of
$F(k(1+sE_{ab}))$.

(3) Let $x\in M_n(\mathbf C)$. For $a<b$ with $\epsilon_a=\epsilon_b$, $x(1+sE_{ab})$ is $x$ with $s$
times its column $a$ added to its column $b$. If $\epsilon_c=\epsilon_b$ and $b\le c$, the submatrix of
$\Delta_c$ contains the columns $a$ and $b$, and $\Delta_c$ does not change; otherwise column $b$ does
not occur in it. So $f_0(x(1+sE_{ab}))=f_0(x)$, and $d\tau(E_{ab})f_0=0$. Next, $x(1+sE_{aa})$ is $x$
with its column $a$ multiplied by $1+s$, so $\Delta_c$ is multiplied by $1+s$ if $\epsilon_c=\epsilon_a$
and $a\le c$ and is unchanged otherwise. Hence $f_0(x(1+sE_{aa}))=f_0(x)(1+s)^{N_a}$, where
$N_a=\sum_ce_c$ over the $c\ge a$ with $\epsilon_c=\epsilon_a$, a telescoping sum equal to $m_a$; so
$d\tau(E_{aa})f_0=m_af_0$. The exponents $e_a$ are $\ge0$ by the hypotheses on $m$. As $f_0(1)=1$,
$f_0\ne0$, and $w_0$ is a unit vector with the same properties.

(4) A subspace stable under $\tau(K)$ that contains $w_0$ contains $f_0$ and its right translates, which
span $V$.
% Lean: modelKType (ModelTau.lean) = kTypeOfGen (MV eps m) (tauM eps m) ... (wN eps m). (1) is
% Vsp_le_polyK and finiteDimensional_Vsp (rtr_polyK, totalDegree_bind1_le), mvCore (mvInner_definite
% through the continuity of the functions and the Haar measure muK), tauC_inner (muK_rightInvariant),
% continuous_tauM. (2) is dTau = cextL (derOp ...), hasDeriv_dTau, Vsp_holo and derOp_fn_eq, with
% hwmf_rdc and hwmf_line for the derivative along 1 + s E_ab and hwmf_bd_single (bdPart). (3) is
% rdc_f0K_hw (hwmf_blk_tv, det_updateCol_add_smul_self), rdc_f0K_wt (hwmf_blk_dg, f0_dg_lemma, the
% telescoping hwmf_tele, mexp_nonneg), highest_weight_fn_K_one, norm_wN, dTau_wt and dTau_hw. (4) is
% gen_wN.
\end{proof}

\begin{lemma}\label{asm-division-kernel-finite-dimension-decay-ktype}
In Situation~\ref{asm-division-kernel-finite-dimension-situation}, let $\lambda\in\mathbf R^n$ be
strictly decreasing with $\lambda_a-\frac{n-1}2\in\mathbf Z$ for every $a$, and let
$\Lambda=\lambda+\rho_n-\rho_c$ with $\rho_n,\rho_c$ as in
Definition~\ref{is-eigenvalue-conjugate-k-types}(3) for the signs $\epsilon$. Then there are an
irreducible continuous unitary representation $\tau$ of $K$ on a finite-dimensional Hilbert space $T$
and a unit vector $w_0\in T$ such that, with $d\tau$ the complex-linear extension of the derivative of
$\tau$ and $E_{ab}$ the elementary matrices, $d\tau(E_{aa})w_0=\Lambda_aw_0$ for every $a$ and
$d\tau(E_{ab})w_0=0$ whenever $a<b$ and $\epsilon_a=\epsilon_b$.
\end{lemma}
\leanref{ReciprocityTheory/ReductiveGroups/DiscreteSeries/HighestWeights/Exists.lean}{ReciprocityTheory.ChartRestriction.DS.HW.kType_exists}
\leanref{ReciprocityTheory/ReductiveGroups/DiscreteSeries/CasimirCoefficients/PsiExistsStatement.lean}{ReciprocityTheory.ChartRestriction.DS.S7.KType}
\leanref{ReciprocityTheory/ReductiveGroups/DiscreteSeries/HighestWeights/Exists.lean}{ReciprocityTheory.ChartRestriction.DS.HW.kType_exists_of_dom}
\leanref{ReciprocityTheory/ReductiveGroups/DiscreteSeries/HighestWeights/HcParamGapGeIndex.lean}{ReciprocityTheory.ChartRestriction.DS.HW.lowK_int_dom}
\leanref{ReciprocityTheory/ReductiveGroups/DiscreteSeries/HighestWeights/HcParamGapGeIndex.lean}{ReciprocityTheory.ChartRestriction.DS.HW.lowest_k_type_formula}
\leanref{ReciprocityTheory/ReductiveGroups/DiscreteSeries/HighestWeights/ModelGen.lean}{ReciprocityTheory.ChartRestriction.DS.HW.irred_of_gen}
\leanref{ReciprocityTheory/ReductiveGroups/DiscreteSeries/HighestWeights/ModelTwist.lean}{ReciprocityTheory.ChartRestriction.DS.HW.kType_twist}
\begin{proof}
\emph{Step 1: $\Lambda$ is integral and dominant for $K$.} For each $a$ let $G_a$ and $L_a$ be the
numbers of $b>a$ and of $b<a$ with $\epsilon_b=\epsilon_a$. The other indices are the $n-a-G_a$ indices
$b>a$ and the $a-1-L_a$ indices $b<a$ with $\epsilon_b\ne\epsilon_a$, so
Definition~\ref{is-eigenvalue-conjugate-k-types}(3) gives
\[
 \Lambda_a=\lambda_a+\frac{n+1}2-a-G_a+L_a,
\]
an integer, as $\lambda_a-\frac{n-1}2$ is. Let $a<b$ with $\epsilon_a=\epsilon_b$. The $\lambda_c$ are
strictly decreasing and differ by integers, so $\lambda_a-\lambda_b\ge b-a$. Every $c>a$ with
$\epsilon_c=\epsilon_a$ is either $>b$ or in $\{a+1,\dots,b\}$, so $G_a\le G_b+(b-a)$; likewise
$L_b\le L_a+(b-a)$. Hence
\[
 \Lambda_a-\Lambda_b=(\lambda_a-\lambda_b)+(b-a)-(G_a-G_b)+(L_a-L_b)\ge0.
\]

\emph{Step 2: a model.} Put $C=\sum_a|\Lambda_a|$ and $m_a=\Lambda_a+C$. Then $m\in\mathbf Z^n$,
$m_a\ge0$, and $m_b\le m_a$ for $a<b$ with $\epsilon_a=\epsilon_b$, by Step 1. Let $V$, $\tau$, $d\tau$
and $w_0$ be as in Lemma~\ref{asm-division-kernel-finite-dimension-decay-ktype-model} for $m$. Then
$\tau$ is irreducible. Indeed, let $W\subset V$ be stable under $\tau(K)$. By
Lemma~\ref{asm-division-kernel-finite-dimension-decay-ktype-derived}, the operators $d\tau(E_{ab})$
with $\epsilon_a=\epsilon_b$ satisfy the relations of
Lemma~\ref{asm-division-kernel-finite-dimension-decay-ktype-highest}, and $W$ is stable under them. The
smallest subspace $V_{w_0}$ that contains $w_0$ and is stable under them is stable under $\tau(K)$, by
Lemma~\ref{asm-division-kernel-finite-dimension-decay-ktype-derived}(3), so $V_{w_0}=V$ by
Lemma~\ref{asm-division-kernel-finite-dimension-decay-ktype-model}(4). As $w_0$ is a highest weight
vector of weight $m$, Lemma~\ref{asm-division-kernel-finite-dimension-decay-ktype-highest}(3) gives
$W=0$ or $W=V$.

\emph{Step 3: the twist.} Put $\tau'(k)=\det(k)^{-C}\tau(k)$ and
$d\tau'(Z)=d\tau(Z)-C\operatorname{tr}(Z)$. As $|\det k|=1$, $\tau'$ is a continuous unitary
representation of $K$ on $T=V$ with the same stable subspaces as $\tau$, so it is irreducible. The
derivative of $\det(\exp tX)^{-C}=e^{-Ct\operatorname{tr}X}$ at $t=0$ is $-C\operatorname{tr}X$, so
$d\tau'$ is the complex-linear extension of the derivative of $\tau'$. As $\operatorname{tr}E_{aa}=1$
and $\operatorname{tr}E_{ab}=0$ for $a\ne b$, $d\tau'(E_{aa})w_0=(m_a-C)w_0=\Lambda_aw_0$, and
$d\tau'(E_{ab})w_0=0$ for $a<b$ with $\epsilon_a=\epsilon_b$.
% Lean: kType_exists = kType_exists_of_dom eps (lowK eps lam) (lowK_int_dom ...). Step 1 is lowK_int_dom
% with lowest_k_type_formula (the formula, with a counted from 0), hc_param_gap_ge_index and the two
% counting bounds c1, c2. Step 2 is model_exists with NonnegDom (C is the sum of the |m a|) and
% modelKType = kTypeOfGen, whose irreducibility is irred_of_gen (pre_kSpan_stable, gen_wN,
% Formalism.REAL.kIrr). Step 3 is kType_twist T (-C): tau' = det^c tau, dtau' = dtau + c tr,
% twist_hasDerivAt_det_exp. This is the existence part of the theorem of the highest weight for
% K = U(p) x U(q); the book follows the Lean's construction, which realizes the representation in
% functions on K spanned by the translates of a product of leading minors.
\end{proof}

\begin{lemma}\label{asm-division-kernel-finite-dimension-decay-spherical}
In Situation~\ref{asm-division-kernel-finite-dimension-situation}, suppose that $q_0\ge1$. Let $\lambda$,
$\Lambda$, $\tau$ and $T$ be as in Lemma~\ref{asm-division-kernel-finite-dimension-decay-ktype}, and let
$h$ be a real number with $|\lambda_a-\lambda_c|\ge h$ whenever $\epsilon_a\ne\epsilon_c$ and $8nq_0<h$.
Then there is a function $\Psi$ from $G$ to the operators on $T$, a \emph{$\tau$-spherical function},
such that:
\begin{enumerate}
\item $\Psi(k_1gk_2)=\tau(k_1)\Psi(g)\tau(k_2)$ for $k_1,k_2\in K$ and $g\in G$, and $\Psi(1)=1$;
\item $\Psi$ is smooth, and $\Omega\Psi=(|\lambda|^2-|\rho|^2)\Psi$, where
  $\Omega=\sum_{a,b}E_{ab}E_{ba}$ is the Casimir operator, the $E_{ab}$ acting as complexified right
  derivatives, and $\rho_a=\frac{n-1}2-(a-1)$;
\item $\Psi$ is square integrable for $\nu$;
\item for every $k\in K$ and $X\in\mathfrak p$,
  \[
   \|\Psi(k\exp X)\|\le\exp\Big(-\frac h8\min\big(|X|^2,|X|/\sqrt{2q_0}\big)\Big).
  \]
\end{enumerate}
\end{lemma}
\leanref{ReciprocityTheory/ReductiveGroups/DiscreteSeries/DiscreteSeriesExistence/EqZeroOfDatumEmbeddingEqZero.lean}{ReciprocityTheory.ChartRestriction.DS.exists_isPsiFor}
\leanref{ReciprocityTheory/ReductiveGroups/DiscreteSeries/CasimirCoefficients/PsiExistsStatement.lean}{ReciprocityTheory.ChartRestriction.DS.S7.PsiExistsStatement}
\leanref{ReciprocityTheory/ReductiveGroups/DiscreteSeries/DiscreteSeriesExistence/EqZeroOfDatumEmbeddingEqZero.lean}{ReciprocityTheory.ChartRestriction.DS.IsPsiFor}
\begin{proof}
\openstep{The construction of the $\tau$-spherical function as a solution of the Casimir equation, and
the proof of its exponential decay along $\exp\mathfrak p$, are not yet written out.}
% Lean: exists_isPsiFor S7.psiExists, with S7.psiExists : PsiExistsStatement n (PsiGelfand.lean). The
% Lean's statement also gives the integrability of Psi for 16 n q0 < h and of |Psi|^2 |g|^s under a
% stronger gap, which are not used here. Size: import-restricted proof closure of psiExists, 847
% declarations, about 11,300 lines in 97 files (text count).
\end{proof}

\begin{lemma}\label{asm-division-kernel-finite-dimension-decay-embedding}
In Situation~\ref{asm-division-kernel-finite-dimension-situation}, let $\lambda$, $\Lambda$, $\tau$ and
$T$ be as in Lemma~\ref{asm-division-kernel-finite-dimension-decay-ktype}, let $h$ be a real number with
$|\lambda_a-\lambda_c|\ge h$ whenever $\epsilon_a\ne\epsilon_c$ and $npq+1<h$, and let $\Psi$ be a
function from $G$ to the operators on $T$ with properties (1) to (3) of
Lemma~\ref{asm-division-kernel-finite-dimension-decay-spherical}. Let $\sigma$ be an irreducible unitary
square-integrable representation of $G$ on a nonzero Hilbert space $H$ that has a $K$-type of weight
$\Lambda$ (Definition~\ref{is-eigenvalue-conjugate-k-types}(2)), with
$|\Lambda+2\rho_c|\le|\mu+2\rho_c|$ for the weight $\mu$ of every $K$-type of $\sigma$, and let its
lowest $K$-type be the sum of its $K$-types of weight $\Lambda$. Then there is a linear isometry
$j\colon T\to H$ with $\sigma(k)j=j\tau(k)$ for $k\in K$, whose image is the lowest $K$-type of
$\sigma$, and with $j^*\sigma(g)j=\Psi(g)$ for every $g\in G$.
\end{lemma}
\leanref{ReciprocityTheory/ReductiveGroups/DiscreteSeries/DiscreteSeriesExistence/RangeDatumEmbedding.lean}{ReciprocityTheory.ChartRestriction.DS.anyF}
\leanref{ReciprocityTheory/ReductiveGroups/DiscreteSeries/IsInfUnitary.lean}{ReciprocityTheory.ChartRestriction.DS.StdGap}
\begin{proof}
\openstep{The proof, which realizes $\sigma$ in the space spanned by the translates of the matrix
coefficients of $\Psi$ and identifies the lowest $K$-types, is not yet written out.}
% Lean: anyF hDS hsg hP: equiv_sigLam_of_isDS gives a unitary equivalence U of sigma with the
% representation sigLam on Hlam (the closed span of the uniformisationL2 vectors of Psi), and
% j = c^-1 U datumEmbedding, with c the norm of uniformisationL2 at T.w (datumEmbedding_highest_ne_zero);
% formulaF_sigLam gives <j v', sigma(g) j v> = <v', Psi(g) v>, rightRegular_datumEmbedding the K-
% equivariance, sigLam_lowK and range_datumEmbedding the image. Size: import-restricted proof closure
% of anyF, 1,777 declarations, about 22,600 lines in 175 files (text count), part of it shared with
% Lemmas -ktype and -spherical.
\end{proof}

\begin{lemma}\label{asm-division-kernel-finite-dimension-decay}
In Situation~\ref{asm-division-kernel-finite-dimension-situation}, suppose that $q_0\ge1$. Let
$\lambda\in\mathbf R^n$ be strictly decreasing with $\lambda_a-\frac{n-1}2\in\mathbf Z$ for every
$a$, and let $h$ be a real number with $|\lambda_a-\lambda_c|\ge h$ whenever $\epsilon_a\ne\epsilon_c$,
$8nq_0<h$ and $npq+1<h$. Let $\sigma$ be an irreducible unitary square-integrable representation of
$G$ on a Hilbert space whose lowest $K$-type has the highest weight attached to $(\lambda,\epsilon)$,
and let $w$ be a unit vector in that $K$-type. Then for every $k\in K$ and $X\in\mathfrak p$,
\[
 |\langle w,\sigma(k\exp X)w\rangle|\le\exp\Big(-\frac h8\min\big(|X|^2,|X|/\sqrt{2q_0}\big)\Big).
\]
\end{lemma}
\leanref{ReciprocityTheory/ReductiveGroups/DiscreteSeries/FormalDegree/CoefDec.lean}{ReciprocityTheory.ChartRestriction.DS.coef_dec_of_isDS}
\begin{proof}
Let $\Lambda=\lambda+\rho_n-\rho_c$, and let $\tau$ and $T$ be given by
Lemma~\ref{asm-division-kernel-finite-dimension-decay-ktype}. As $q_0\ge1$ and $8nq_0<h$,
Lemma~\ref{asm-division-kernel-finite-dimension-decay-spherical} gives a $\tau$-spherical function
$\Psi$. Its properties (1) to (3), the gap $h$ with $npq+1<h$ and the strictly decreasing $\lambda$ are
the hypotheses of Lemma~\ref{asm-division-kernel-finite-dimension-decay-embedding}, which gives an
isometry $j\colon T\to H$ onto the lowest $K$-type of $\sigma$ with $j^*\sigma(g)j=\Psi(g)$ for $g\in G$.
As $w$ lies in the lowest $K$-type, $w=jv$ for some $v\in T$, and $|v|=|w|=1$. Let $k\in K$ and
$X\in\mathfrak p$. Then $\exp X\in G$, so $g=k\exp X\in G$, and
\[
 \langle w,\sigma(g)w\rangle=\langle jv,\sigma(g)jv\rangle=\langle v,j^*\sigma(g)jv\rangle
 =\langle v,\Psi(g)v\rangle.
\]
So $|\langle w,\sigma(g)w\rangle|\le|v|\,\|\Psi(g)\|\,|v|=\|\Psi(g)\|$, and property (4) of
Lemma~\ref{asm-division-kernel-finite-dimension-decay-spherical} gives the bound.
% Lean: coef_dec_of_isDS: HW.kType_exists gives T; exists_isPsiFor S7.psiExists gives Psi with
% IsPsiFor and the decay; StdGap from hlam.1, hgap and hstd; anyF gives j (hrange for w = j v, hj for
% |v| = 1, hadj for j^* sigma(g) j = Psi(g)); rq_exp_mem and rq_UJ_mul (ffa_toMK_mem) put k exp X in
% UJ; norm_inner_le_norm and ContinuousLinearMap.le_opNorm give the last bound.
% Size: import-restricted proof closure of coef_dec_of_isDS, 2,581 declarations, about 33,800 lines in
% 263 files (text count); the open parts are Lemmas -spherical (about 11,300) and -embedding (about
% 22,600), which overlap; Lemma -ktype (about 2,800) is written out.
\end{proof}

\begin{lemma}\label{asm-division-kernel-finite-dimension-growth}
In Situation~\ref{asm-division-kernel-finite-dimension-situation}:
\begin{enumerate}
\item there is $C\ge0$ with $\nu\{g\in G:|g|\le R\}\le CR^{2nq_0}$ for every $R\ge1$;
\item let $f$ be a measurable function on $G$ with values in a normed space, and let $M$ and
  $a>2nq_0$ be real numbers such that $|f(g)|\le M$ whenever $|g|^2\le ne^2+1$, and such that for
  every $g$ with $|g|^2>ne^2+1$ there is a real $t$ with $|g|^2\le ne^{2t}$ and $|f(g)|\le Me^{-at}$.
  Then $f$ is integrable for $\nu$.
\end{enumerate}
\end{lemma}
\leanref{ReciprocityTheory/ReductiveGroups/DiscreteSeries/CasimirCoefficients/IntegrableOfDecayVG.lean}{ReciprocityTheory.ChartRestriction.DS.S7.volGrowthVG}
\leanref{ReciprocityTheory/ReductiveGroups/DiscreteSeries/CasimirCoefficients/IntegrableOfDecayVG.lean}{ReciprocityTheory.ChartRestriction.DS.S7.integrable_of_decayVG}
\leanref{ReciprocityTheory/ReductiveGroups/DiscreteSeries/CasimirCoefficients/SignedIdentityColumns.lean}{ReciprocityTheory.ChartRestriction.DS.S7.hs_le_colsV}
\leanref{ReciprocityTheory/ReductiveGroups/DiscreteSeries/CasimirCoefficients/SignedIdentityColumns.lean}{ReciprocityTheory.ChartRestriction.DS.S7.smallRV}
\begin{proof}
For $\rho>0$ the set $\{g\in G:|g|\le\rho\}$ is closed and bounded in $M_n(\mathbf C)$, and $G$ is
closed, so it is compact and has finite measure.

(1) Choose the sign $s$ that occurs $q_0$ times in $\epsilon$, let $Q$ be the set of its $q_0$
positions, $V=M_{n\times Q}(\mathbf C)$, a real vector space of dimension $2nq_0$ with a Lebesgue
measure $\lambda$, and $E\in V$ the columns of the identity matrix at the positions $Q$, so that
$|E|=\sqrt{q_0}\le\sqrt n$. Put $O=\{u\in V:|u-E|<\frac14\}$, open and bounded, so
$0<\lambda(O)<\infty$.

\emph{A bound on $|k|$.} For $k\in G$, $|k|^2\le n+2|kE|^2$. Indeed $kJk^*=J$, as $k^{-1}=Jk^*J$; its
$i$-th diagonal entry gives $\sum_j\epsilon_j|k_{ij}|^2=\epsilon_i$, so that
$\sum_j|k_{ij}|^2=\mp\epsilon_i+2\sum_{j\in Q}|k_{ij}|^2\le1+2\sum_{j\in Q}|k_{ij}|^2$, with the sign
$-$ if $s=+1$ and $+$ if $s=-1$; sum over $i$. Now let $k\in G$ and $u,v\in O$ with $ku=v$. Then
$kE=v-k(u-E)$, so $|kE|\le|v|+\frac14|k|\le\sqrt n+\frac14+\frac14|k|$. With $|k|^2\le n+2|kE|^2$, an
elementary computation gives $|kE|\le\frac52\sqrt n+\frac12$, and so $|k|^2\le r_0$, where
$r_0=n+2(\frac52\sqrt n+\frac12)^2$.

\emph{Fubini.} Left multiplication by $g\in G$ is a real linear map of $V$ of determinant
$|\det g|^{2q_0}=1$, as $|\det g|=1$. Let $R\ge1$, $S=\{g\in G:|g|\le R\}$ and
$T=\{(g,x)\in S\times V:g^{-1}x\in O\}$. For $g\in S$ the slice of $T$ over $g$ is $gO$, of measure
$\lambda(O)$; so $(\nu\times\lambda)(T)=\nu(S)\lambda(O)$. For $x\in V$, suppose the slice
$T^x=\{g\in S:g^{-1}x\in O\}$ contains some $h$. Then $x=hu$ with $u\in O$, so every entry of $x$ has
absolute value at most $|h||u|\le R(\sqrt n+1)$. And for $g\in T^x$, $h^{-1}g$ carries $g^{-1}x\in O$
to $h^{-1}x\in O$, so $|h^{-1}g|^2\le r_0$ by the bound on $|k|$: $T^x\subset hW$, where
$W=\{k\in G:|k|^2\le r_0\}$. Hence, with $B_\rho$ the set of $x\in V$ whose entries have absolute
value at most $\rho$,
\[
 \nu(S)\lambda(O)=(\nu\times\lambda)(T)\le\nu(W)\,\lambda\big(B_{R(\sqrt n+1)}\big)
 =\nu(W)\lambda(B_1)(\sqrt n+1)^{2nq_0}R^{2nq_0},
\]
and (1) holds with $C=\nu(W)\lambda(B_1)(\sqrt n+1)^{2nq_0}/\lambda(O)$.

(2) Put $N=2nq_0$. As $|1|^2=n$, $M\ge|f(1)|\ge0$. Let $B_0=\{g:|g|^2\le ne^2+1\}$ and
$S_j=\{g:|g|^2\le4^j\}$. Let $g\notin B_0$, with $t$ as in the hypothesis, and choose $j\ge0$ with
$4^j\le|g|^2<4^{j+1}$. Then $g\in S_{j+1}$, and $4^j\le ne^{2t}$ gives $t\ge j\log2-\frac12\log n$, so
$|f(g)|\le Me^{-at}\le Mn^{a/2}2^{-aj}$. Hence
\[
 |f|\le M\,\mathbf 1_{B_0}+\sum_{j\ge0}Mn^{a/2}2^{-aj}\,\mathbf 1_{S_{j+1}}.
\]
By (1) with $R=2^{j+1}$, $\nu(S_{j+1})\le C2^N2^{Nj}$, so the $j$-th term has integral at most
$MCn^{a/2}2^N(2^{N-a})^j$, and $2^{N-a}<1$ as $a>N$. So the right side has finite integral, and $f$
is integrable.
% Lean: (1) is volGrowthVG: exists_signVG chooses s, O = {x : hsRV (x - signedIdentityColumns) <
% 1/16} (so |u - E| < 1/4), RLV_det (the determinant), ballFiniteVG (W), Measure.prod_apply and
% prod_apply_symm (the two slicings), entryRV (the entries of x), smallRV with hs_le_colsV (the bound on
% |k|, through normRV_u and normRV_lt), and addHaar_closedBall for the sup-norm ball. (2) is
% integrable_of_decayVG: the dominating function F with exists_nat_pow_near, the terms bounded through
% volGrowthVG, and summable_geometric_of_lt_one. Size: IntegrableOfDecayVG.lean and
% SignedIdentityColumns.lean, about 620 lines in the closure (text count).
\end{proof}

\begin{lemma}\label{asm-division-kernel-finite-dimension-weighted}
In Situation~\ref{asm-division-kernel-finite-dimension-situation}, let $k\ge1$ and $r_S\ge0$ be
integers with $n=2r_S+k+1$, and let $\epsilon$ consist of $k-1$ entries $+1$ followed by $2r_S+2$
entries $+1,-1,+1,-1,\dots,+1,-1$. Let $v\in\mathbf R^{k-1}$ and $b\in\mathbf R^{2r_S+2}$, and put
$\lambda=(v_1,\dots,v_{k-1},-b_1,\dots,-b_{2r_S+2})$. Suppose that $\lambda$ is strictly decreasing
with $\lambda_a-\frac{n-1}2\in\mathbf Z$ for every $a$, and that there is an odd integer $t$ with
$n+5\le t$, $1801\le n+t$ and
\[
 b_{j+1}-b_j>2M^3+t+3\qquad(1\le j\le2r_S+1),\qquad M=n+(2r_S+2)t.
\]
Let $\sigma$ be an irreducible unitary square-integrable representation of $G$ on a Hilbert space
whose lowest $K$-type has the highest weight attached to $(\lambda,\epsilon)$, and let $w$ be a unit
vector in that $K$-type. Then, with $N=4n^2\sqrt{2(r_S+1)}$,
\[
 \int_G|\langle w,\sigma(g)w\rangle|\,\|g\|^{N}\,d\nu(g)<\infty.
\]
\end{lemma}
\leanref{ReciprocityTheory/ReductiveGroups/DiscreteSeries/FormalDegree/CoefLoc.lean}{ReciprocityTheory.ChartRestriction.DS.coefWeightedLoc}
\leanref{ReciprocityTheory/ReductiveGroups/DiscreteSeries/FormalDegree/CoefLocGap.lean}{ReciprocityTheory.ChartRestriction.DS.cayley_weight_gap}
\leanref{ReciprocityTheory/ReductiveGroups/DiscreteSeries/FormalDegree/CoefWeightedG.lean}{ReciprocityTheory.ChartRestriction.DS.coefWeightedN}
\leanref{ReciprocityTheory/ReductiveGroups/DiscreteSeries/FormalDegree/CoefIntVG.lean}{ReciprocityTheory.ChartRestriction.DS.integrable_opPow_of_dec}
\begin{proof}
Here $p=k+r_S$ and $q=r_S+1$, so $q_0=r_S+1\ge1$. Put $h=2M^3+t+3$ and $c=\sqrt{2q_0}\ge1$.

\emph{Step 1: the gaps.} Every increment of $b$ exceeds $h>0$, so $b_j-b_i>h$ for $i<j$. Let
$\epsilon_a\ne\epsilon_c$. If $a,c>k-1$, then $|\lambda_a-\lambda_c|=|b_{a-k+1}-b_{c-k+1}|>h$.
Otherwise one of them, say $a$, is at most $k-1$, with sign $+1$, and $c=k-1+j$ with $j$ even, so
$j\ge2$; as $\lambda$ is decreasing, $\lambda_a>\lambda_k=-b_1$, and
$\lambda_a-\lambda_c>b_j-b_1>h$. So $|\lambda_a-\lambda_c|\ge h$ whenever $\epsilon_a\ne\epsilon_c$.
Next, $p,q\le n\le M$, so $npq+1\le M^3+1<h$. As $t\ge n+5$, $M\ge3n$ and $M\ge10(r_S+1)$, so
$96n^2(r_S+1)\le2M^3<h$; in particular $8nq_0<h$. Using $c^2=2(r_S+1)$ and $q_0\le n$,
\[
 8c\,(N+2nq_0)=32n^2c^2+16nq_0c\le64n^2(r_S+1)+32n(r_S+1)^2\le96n^2(r_S+1)<h.
\]

\emph{Step 2: the coefficient.} Let $\varphi(g)=\langle w,\sigma(g)w\rangle$. It is continuous, as
$\sigma$ is strongly continuous, and $|\varphi(g)|\le1$, as $\sigma$ is unitary and $|w|=1$. By Step 1,
Lemma~\ref{asm-division-kernel-finite-dimension-decay} applies with $h$:
$|\varphi(k\exp X)|\le\exp(-\frac h8\min(|X|^2,|X|/c))$ for $k\in K$ and $X\in\mathfrak p$.

\emph{Step 3: the polar decomposition.} Every $g\in G$ is $k\exp X$ with $k\in K$ and $X\in\mathfrak p$.
Let $L$ be the Hermitian matrix with $\exp(2L)=g^*g$, which is positive definite. As
$g^{-1}=Jg^*J$, $\exp(-2L)=(g^*g)^{-1}=J\,g^*g\,J=\exp(2JLJ)$, and the exponential is injective on
Hermitian matrices, so $JLJ=-L$ and $L\in\mathfrak p$. Then $k=g\exp(-L)$ satisfies
$k^*k=\exp(-L)g^*g\exp(-L)=1$ and $k^*Jk=\exp(-L)J\exp(-L)=J\exp(L)\exp(-L)=J$, so $k\in K$, and
$g=k\exp L$.

\emph{Step 4: integration.} Let $F(g)=|\varphi(g)|\,\|g\|^N$, continuous on $G$, and
$a=h/(8c)-N$, so that $a>2nq_0$ by Step 1. Let $g=k\exp X$ as in Step 3 and $r=|X|$. As $k$ is
unitary, $|g|=|\exp X|\le\sqrt n\,e^r$, by the exponential series, as $|1|=\sqrt n$ and
$|X^m|\le|X|^m\le\sqrt n|X|^m$ for $m\ge1$; and $\|g\|\le\|\exp X\|\le e^{\|X\|}\le e^r$; also
$\|g\|\le|g|$.
If $|g|^2\le ne^2+1$, then $F(g)\le(ne^2+1)^{N/2}$. If $|g|^2>ne^2+1$, then $ne^2<ne^{2r}$, so $r>1$
and $\min(r^2,r/c)=r/c$; Step 2 gives $|\varphi(g)|\le e^{-hr/(8c)}$, so $F(g)\le e^{-ar}$, and
$|g|^2\le ne^{2r}$. Lemma~\ref{asm-division-kernel-finite-dimension-growth}(2), with
$M=(ne^2+1)^{N/2}\ge1$ and $t=r$, shows that $F$ is integrable.
% Lean: coefWeightedLoc = cayley_weight_gap (Step 1: gapAtLeast_of_increments with
% lt_sub_of_increments and cayley_weight_high_index, q_le_of_cayley_signature, cubic_bound_of_gap_nat)
% and coefWeightedN. coefWeightedN takes coef_continuousOn_le_one (Step 2) and splits on q0: for
% q0 = 0, integrable_opPow_of_compact; for q0 >= 1, integrable_opPow_of_dec with coef_dec_of_isDS
% (Lemma asm-division-kernel-finite-dimension-decay). integrable_opPow_of_dec uses S7.polar_exists
% (polar_UJ: polar_L, polar_exp_LL, polar_JLJ with polar_exp_inj, polar_L_inP, polar_k_mem; Step 3),
% hs_kmulVG and hs_exp_leVG (|g|^2 <= n e^{2r}), coefOp_kexp and coefOp_le_sqrt_hs, and
% integrable_of_decayVG (Lemma asm-division-kernel-finite-dimension-growth(2)). The Lean states the
% lemma for every k, with k - 1 the truncated difference; here q0 >= 1 needs k >= 1.
% Size: closures of cayley_weight_gap, integrable_opPow_of_dec and integrable_opPow_of_compact,
% about 1,900 lines (text count), besides the decay bound.
\end{proof}

\begin{lemma}\label{asm-division-kernel-finite-dimension}
In Situation~\ref{asm-division-kernel-finite-situation}, suppose that $D$ satisfies the gap
hypothesis (Definition~\ref{is-wu-decay}), and assume the finite dimensionality of rows. Then
$\mathcal H$ is finite-dimensional.
\end{lemma}
\leanref{Formalism/LocalToGlobal/SiegelLatticeTheory/FlagStabilizerTensors/RootLatticeThresholds/RowDictPf.lean}{Formalism.L2.SG.Pf.RowDictPf.finiteDimensional_HHX_of_hypotheses}
\leanref{Formalism/LocalToGlobal/SiegelLatticeTheory/FlagStabilizerTensors/RootLatticeThresholds/FinXHolds.lean}{Formalism.L2.SG.Pf.FinXP.phLoc_int_of_gap}
\begin{proof}
For each infinite place $v$ let $\sigma_v$ be the discrete series of the setting at $v$, $w_v$ the
unit vector of its lowest $K$-type given by the setting, $G_v$ the unitary group at $v$ with its Haar
measure, and
$\phi_v(g)=d_v\overline{\langle w_v,\sigma_v(g)w_v\rangle}$, where $d_v$ is the formal degree of
$\sigma_v$ computed from $w_v$. The finite dimensionality of rows, applied to the setting with the
$\sigma_v$, the $w_v$ and the $d_v$, gives that $\mathcal H$ is finite-dimensional, provided that each
$d_v$ is a formal degree of $\sigma_v$ and each $\phi_v$ is integrable on $G_v$.

\emph{Formal degrees.} By Schur orthogonality for square-integrable irreducible unitary
representations \cite[Chapter~14]{Dix77}, $\sigma_v$ has a formal degree $d$. As $w_v\ne0$, the
number computed from $w_v$ is $d$, so $d_v=d$.

\emph{Integrability.} Let $\|g\|$ be the operator norm of the matrix of $g\in G_v$, and $N_1$ the
exponent of Definition~\ref{is-wu-decay}. The group $G_v$ preserves a Hermitian form with matrix $J$,
$g^*Jg=J$, where $J^2=1$ and $\|J\|\le1$; so $1=\|Jg^*Jg\|\le\|g\|^2$ and $\|g\|\ge1$. Hence
\[
 |\langle w_v,\sigma_v(g)w_v\rangle|\le|\langle w_v,\sigma_v(g)w_v\rangle|\,\|g\|^{N_1}
 \le|\langle w_v,\sigma_v(g)w_v\rangle|\,\|g\|^{N_2},\qquad N_2=4n^2\sqrt{2(r_S+1)},
\]
where $n=2r_S+k+1$ is the size of the matrices; the second inequality holds because $N_1\le N_2$, as
every local invariant $q_v$ of the setting is at most $r_S+1$. At $v$ the signs of $J$ and the
Harish-Chandra parameter of $\sigma_v$ are the $\epsilon$ and $\lambda$ of
Lemma~\ref{asm-division-kernel-finite-dimension-weighted} for the ranks $k$ and $r_S$ of $D$ and the
parts $(v_i)$ and $(b_j)$ of the Harish-Chandra parameter of $D$ at that place
(Definition~\ref{th-hc-parameter}), with $n=r$; this $\lambda$ is strictly decreasing with integral
entries, and $\frac{n-1}2$ is an integer as $r$ is odd. The gap hypothesis gives the condition on the $b_j$
there with $t=r+6$, which is odd as $r$ is: $n+5\le t$ and $n+t\ge1801$, as $r\ge1801$, and
$b_{j+1}-b_j\ge16r^{19}>2\big(r+(2r_S+2)(r+6)\big)^3+r+9$, as $2r_S+2=m\le r$. So
Lemma~\ref{asm-division-kernel-finite-dimension-weighted} shows that the right side is integrable on
$G_v$. So the matrix coefficient is integrable, and so is $\phi_v$.
% Lean: finiteDimensional_HHX_of_hypotheses = finiteDimensional_HHX_of X hRF (fdeg_of_schur X hS) hint,
% which is the hypothesis hRF : RowsFinDim applied to the fields of the setting; fdeg_of_schur takes d
% from schur_orthogonality_discrete_series (PIW.HCFdeg.formal_degree_abstract_form) and fdegOf_eq with
% w != 0 (X.low v). hint is phLoc_int_of_gap = coef_integrable_of_weighted (one_le_weight: 1 <=
% opn(uMat g)^N1) applied to coefWeighted_of_gap, which bounds the N1-weighted integral by the one with
% exponent 4 n^2 sqrt(2(r+1)) (one_le_opn, first_count_le_log: N1 <= that exponent, through qAt_le)
% and takes its finiteness from coefW_cast with gapBig_of_gapHyp and the hypothesis CoefWeightedLoc.
% CoefWeightedLoc is a Fact instance in this file, proved as
% ReciprocityTheory.ChartRestriction.DS.coefWeightedLoc (CoefLoc.lean, Lemma
% asm-division-kernel-finite-dimension-weighted) and supplied by WeightedGaugeFactors/FactsDS.lean.
% gapBig_of_gapHyp takes t = r + 6 (r_odd) and the inequality numNV; coefW_cast identifies the signs and
% parameter of the setting at v with eps48 k rS and lam48 k rS (v48 D) (b48 D) (eps_cast, lam_cast,
% hdim).
\end{proof}

\begin{lemma}\label{asm-division-kernel-finite-fibres-induced}
Let $E$ be a number field with an involution $c\mapsto\bar c$ such that
$\tau(\bar c)=\overline{\tau(c)}$ for every embedding $\tau\colon E\to\mathbf C$, let $V$ be a
finite-dimensional $E$-vector space with a Hermitian form $h$, conjugate-linear in the first
variable, and let $\Lambda\subset V$ be an $\mathcal O_E$-submodule that spans $V$ over $E$ and on
which $h$ takes values in $\mathcal O_E$. Let $X\subset P\subset V$ be subspaces with $h(x,y)=0$ for
$x\in X$ and $y\in P$, and let $z\in P\cap\Lambda$ be such that for every embedding $\tau$:
$\tau(h(z,z))<0$, and $\tau(h(y,y))>0$ for every $y\in P$ with $h(z,y)=0$ and $y\notin X$. Let $A$
be the set of maps $\varphi\colon P/X\to P/X$ for which there is an $E$-linear bijection $\gamma$ of
$V$ that preserves $h$, with $\gamma z=z$, $\gamma\Lambda\subset\Lambda$, $\gamma P\subset P$ and
$\varphi(y+X)=\gamma y+X$ for every $y\in P$. Then $A$ is finite.
\end{lemma}
\leanref{Formalism/PeriodIntegrals/UnitaryGroupMassFormula/LatticeAutomorphismFiniteness/IndSetFinite.lean}{PI.AutLzFinite.SlopeTestRemainingChoices.indSet_finite}
\leanref{Formalism/PeriodIntegrals/UnitaryGroupMassFormula/LatticeAutomorphismFiniteness/MjSmulRight.lean}{PI.AutLzFinite.SlopeTestRemainingChoices.mj_pos}
\leanref{Formalism/PeriodIntegrals/UnitaryGroupMassFormula/LatticeAutomorphismFiniteness/MjSmulRight.lean}{PI.AutLzFinite.SlopeTestRemainingChoices.mj_cs}
\leanref{Formalism/PeriodIntegrals/UnitaryGroupMassFormula/LatticeAutomorphismFiniteness/FiniteOfIntegralBound.lean}{PI.AutLzFinite.SlopeTestRemainingChoices.finite_of_integral_bound}
\begin{proof}
The values $h(y,y)$ are fixed by the involution, so the $\tau(h(y,y))$ are real. Put $t=h(z,z)$,
which is nonzero, and
\[
 m(x,y)=h(x,y)-2t^{-1}h(x,z)h(z,y)\qquad(x,y\in V).
\]
It is conjugate-linear in $x$, linear in $y$, and $m(y,x)=\overline{m(x,y)}$. Fix an embedding $\tau$.

\emph{Step 1: positivity.} For $y\in P$, $\tau(m(y,y))\ge0$, with equality only if $y\in X$. Put
$c=t^{-1}h(z,y)$ and $y_0=y-cz\in P$. Then $h(z,y_0)=0$, and expanding,
$m(y,y)=h(y_0,y_0)-\bar cc\,t$. If $y_0\in X$, then $h(y_0,y_0)=0$; otherwise
$\tau(h(y_0,y_0))>0$. And $-|\tau(c)|^2\tau(t)\ge0$. So $\tau(m(y,y))\ge0$, and if it is $0$, then
$\tau(c)=0$ and $\tau(h(y_0,y_0))=0$, so $c=0$, $y_0\in X$ and $y=y_0\in X$.

\emph{Step 2: Cauchy--Schwarz.} For $x,y\in P$,
$|\tau(m(x,y))|^2\le\tau(m(x,x))\,\tau(m(y,y))$. This is the Cauchy--Schwarz inequality for the
form $\tau\circ m$, which is positive semidefinite on $P$ by Step 1: if $m(x,x)\ne0$, apply Step 1 to
$sx+y$ with $s=-m(x,y)/m(x,x)$; if $m(x,x)=0$, apply it to $sx+y$ with $s=-Nm(x,y)$ for large
integers $N$.

\emph{Step 3: a basis in $\Lambda$.} As $\Lambda$ spans $V$, every vector of $V$ has a nonzero
multiple by an element of $\mathcal O_E$ in $\Lambda$. So $P$ has a basis $b_1,\dots,b_r$ with
$b_i\in\Lambda$. Put $C=\sum_i\sum_\tau|\tau(m(b_i,b_i))|$, so that $\tau(m(b_i,b_i))\le C$ for all
$i$ and $\tau$.

\emph{Step 4: a finite set of classes.} Let $Y$ be the set of classes $y+X$ of the $y\in P\cap\Lambda$
with $\tau(m(y,y))\le C$ for every $\tau$. The map $g\colon P/X\to E^r$,
$y+X\mapsto(m(b_j,y))_j$, is well defined, since $h(b_j,y)=h(z,y)=0$ for $y\in X$, and it is
injective: if $m(b_j,w)=0$ for every $j$, then $m(v,w)=0$ for every $v\in P$, so $m(w,w)=0$ and
$w\in X$ by Step 1. Let $y+X\in Y$ and $a=m(b_j,y)$. Then
$ta=t\,h(b_j,y)-2h(b_j,z)h(z,y)$ lies in $\mathcal O_E$, as $b_j,y,z\in\Lambda$, and
$|\tau(a)|^2\le\tau(m(b_j,b_j))\,\tau(m(y,y))\le C^2$ by Step 2. The set of $a\in E$ with
$ta\in\mathcal O_E$ and $|\tau(a)|\le C$ for every $\tau$ is finite, since the $ta$ are then
algebraic integers of $E$ whose conjugates are at most $C\sum_\tau|\tau(t)|$ in absolute value. So
$g(Y)$ is finite, and so is $Y$.

\emph{Step 5.} Let $\varphi\in A$, induced by $\gamma$. Then $\varphi(b_i+X)=\gamma b_i+X$ with
$\gamma b_i\in P\cap\Lambda$, and $m(\gamma b_i,\gamma b_i)=m(b_i,b_i)$, as $\gamma$ preserves $h$ and
fixes $z$; so $\varphi(b_i+X)\in Y$. If $\varphi,\varphi'\in A$, induced by $\gamma,\gamma'$, take the
same values at the $b_i+X$, then the linear maps $y\mapsto\gamma y+X$ and $y\mapsto\gamma'y+X$ from
$P$ to $P/X$ agree on a basis, so they agree, and $\varphi=\varphi'$. So $A$ injects into $Y^r$, which
is finite.
% Lean: indSet_finite, with mj (MjSmulRight.lean), mj_pos (Step 1), mj_cs (Step 2), exists_smul_mem
% and Basis.unitsSMul (Step 3), the map g = Xc.liftQ Pm with finite_of_integral_bound
% (Embeddings.finite_of_norm_le, Step 4), and Set.Finite.of_finite_image with the map
% phi -> (phi (b_i mod X))_i (Step 5). Size: the directory LatticeAutomorphismFiniteness, 657 lines in
% six files (text count), including autLzFin_main, pos_perp and sylvester_le used in
% Lemma asm-division-kernel-finite-fibres.
\end{proof}

\begin{lemma}\label{asm-division-kernel-finite-fibres}
In Situation~\ref{asm-division-kernel-finite-situation}, over each term of the unfolding
(Definition~\ref{is-gram-terms}) lie finitely many terms of the unfolded term data of
Definition~\ref{is-unfolding-evaluation}.
\end{lemma}
\leanref{Formalism/LocalToGlobal/SiegelLatticeTheory/FlagStabilizerTensors/RootLatticeThresholds/XLeSpanEXLam.lean}{Formalism.L2.SG.Pf.FinXP.fibFin_holds}
\leanref{Formalism/PeriodIntegrals/OrbitalBoundsConclusion/AutLzFinite.lean}{automorphism_lattice_finite}
\leanref{Formalism/LocalToGlobal/SiegelLatticeTheory/FlagStabilizerTensors/RootLatticeThresholds/XLeSpanEXLam.lean}{Formalism.L2.SG.Pf.FinXP.fib_decomp}
\leanref{Formalism/LocalToGlobal/SiegelLatticeTheory/FlagStabilizerTensors/RootLatticeThresholds/XLeSpanEXLam.lean}{Formalism.L2.SG.Pf.FinXP.conj_fix_mem}
\begin{proof}
\emph{Notation.} Let $u,u'$ be words. The flag $F$ of the rows of $u$ consists of linearly
independent isotropic vectors $e_1,\dots,e_r$ and a vector $z$ orthogonal to them; put
$X_j=\langle e_1,\dots,e_j\rangle$ and $X=X_r$. For a lattice $\Lambda$ and such a flag $G$, with
spaces $X_j(G)$ and $X(G)$, let $N_\Lambda(G)$ be the set of isometries $n$ with $n\Lambda=\Lambda$,
$(n-1)X_{j+1}(G)\subset X_j(G)$ for every $j$, and $(n-1)X(G)^\perp\subset X(G)$; these are the
automorphisms of $\Lambda$ unipotent along $G$ of Definition~\ref{is-unfolding-evaluation}(1), and
they form a group. The stabilizer $\operatorname{St}(F)$ is the group of isometries that preserve
every $X_j$ and fix $z$, and the rows of $u$ contain exactly one lattice in each orbit of
$\operatorname{St}(F)$ on the row lattices for $F$. Write $F'$, with vectors $f_1,\dots,f_r$ and
$z'$ and spaces $X'_j$ and $X'$, for the flag of the rows of $u'$. A term $s$ of $T_{uu'}$ has rows
$\Lambda_s$ of $u$ and $\Lambda'_s$ of $u'$ and the double coset
$N_{\Lambda_s}(F)\gamma_sN_{\Lambda'_s}(F')$ of an isometry $\gamma_s$ with
$\gamma_s\Lambda'_s=\Lambda_s$. The term of the unfolding below $s$ is $\Lambda_s$ with the flag
$\gamma_sF'$; two flags $G_1,G_2$, each with a further vector $\tilde z_i$, give the same term over
$\Lambda$ when there are $n\in N_\Lambda(F)$ and $m\in N_\Lambda(G_1)$ such that $nm$ carries the
vectors $e_1,\dots,e_r,z$ of $G_1$ and $\tilde z_1$ to those of $G_2$ and $\tilde z_2$.

\emph{Step 1: two terms over the same term.} Let $s,s'$ lie over the same term of the unfolding.
Then $\Lambda_s=\Lambda_{s'}=:\Lambda$, and there are $n\in N_\Lambda(F)$ and
$m\in N_\Lambda(\gamma_sF')$ with $\gamma_{s'}f=nm\gamma_sf$ for $f=f_1,\dots,f_r,z'$. Put
$\mu=\gamma_s^{-1}m\gamma_s$; as $\gamma_s$ is an isometry carrying $\Lambda'_s$ to $\Lambda$ and $F'$
to $\gamma_sF'$, $\mu\in N_{\Lambda'_s}(F')$. Put $c=\mu^{-1}\gamma_s^{-1}n^{-1}\gamma_{s'}$. For
$f=f_1,\dots,f_r,z'$,
\[
 cf=\mu^{-1}\gamma_s^{-1}n^{-1}nm\gamma_sf=\mu^{-1}\mu f=f.
\]
So $c$ is an isometry that fixes the $f_a$ and $z'$; in particular $c\in\operatorname{St}(F')$, and
$c\Lambda'_{s'}=\mu^{-1}\gamma_s^{-1}n^{-1}\Lambda=\mu^{-1}\Lambda'_s=\Lambda'_s$. As the rows of $u'$
meet each orbit of $\operatorname{St}(F')$ once, $\Lambda'_{s'}=\Lambda'_s$. Finally
$\gamma_{s'}=n\gamma_s\mu c$.

\emph{Step 2: the map.} Fix a term of the unfolding. If no term of $T_{uu'}$ lies over it, there is
nothing to show. Otherwise fix one, $s_0$, and put $\Lambda'=\Lambda'_{s_0}$ and
$\gamma_0=\gamma_{s_0}$. For each $s$ over the same term, Step 1 for $(s_0,s)$ gives
$n_s\in N_\Lambda(F)$, $\mu_s\in N_{\Lambda'}(F')$ and an isometry $c_s$ fixing $f_1,\dots,f_r,z'$
with $c_s\Lambda'=\Lambda'$ and $\gamma_s=n_s\gamma_0\mu_sc_s$. As $c_s$ fixes $X'$ pointwise, it
preserves $X'^\perp$, and it induces a map $\varphi(s)$ of $X'^\perp/X'$. This map lies in the set
$A(\Lambda')$ of the maps of $X'^\perp/X'$ induced by elements of $\operatorname{St}(F')$ that
preserve $\Lambda'$.

\emph{Step 3: $\varphi$ is injective.} Let $\varphi(s)=\varphi(s')$, so that
$c_sx-c_{s'}x\in X'$ for every $x\in X'^\perp$. Put $c=c_s^{-1}c_{s'}$. It is an isometry with
$c\Lambda'=\Lambda'$, it fixes $X'$ pointwise, and for $x\in X'^\perp$,
$cx-x=c_s^{-1}(c_{s'}x-c_sx)\in X'$, since $c_s^{-1}$ fixes $X'$. So $c\in N_{\Lambda'}(F')$.
Conjugation by $c_s$ preserves $N_{\Lambda'}(F')$: for $\nu$ in it, $c_s^{-1}\nu c_s$ preserves
$\Lambda'$, for $x\in X'_{j+1}$ one has $c_s^{-1}\nu c_sx-x=c_s^{-1}(\nu x-x)\in X'_j$, and for
$x\in X'^\perp$, $c_sx\in X'^\perp$ and $c_s^{-1}\nu c_sx-x=c_s^{-1}(\nu c_sx-c_sx)\in X'$. Hence
\[
 \gamma_{s'}=n_{s'}\gamma_0\mu_{s'}c_{s'}
 =\big(n_{s'}n_s^{-1}\big)\,\gamma_s\,\big(c_s^{-1}\mu_s^{-1}\mu_{s'}c_s\big)\big(c_s^{-1}c_{s'}\big)
 \in N_\Lambda(F)\,\gamma_s\,N_{\Lambda'}(F').
\]
With $\Lambda_s=\Lambda_{s'}$ and $\Lambda'_s=\Lambda'_{s'}$ by Step 1, this gives $s=s'$.

\emph{Step 4: $A(\Lambda')$ is finite.} The isometries are those of the space $V=Y_0$ of $D$ over the
CM field $E$ of its host, with its form $h$ (Definition~\ref{th-spaces}(5)); $V$ has signature
$(k+r_S,r_S+1)$ at every embedding $\tau$ of $E$, and $r=r_S$. The vector $z'$ has
$h(z',z')=-T_{u'}$, where $T_{u'}$ is the prime of the word $u'$, and the row $\Lambda'$ is a lattice
of $V$, unimodular for $h$, that contains $z'$. We apply
Lemma~\ref{asm-division-kernel-finite-fibres-induced} to $V$, $h$, $\Lambda'$, $X=X'$, $P=X'^\perp$
and $z=z'$. Complex conjugation on $E$ commutes with every embedding. As $\Lambda'$ is unimodular,
$h$ takes values in $\mathcal O_E$ on $\Lambda'$. Clearly $h(X',X'^\perp)=0$, and $z'\in X'^\perp$, as
$z'$ is orthogonal to the $f_a$. For every $\tau$, $\tau(h(z',z'))=-T_{u'}<0$.

For the last hypothesis, let $y\in X'^\perp$ with $h(z',y)=0$ and $y\notin X'$, and suppose that
$\tau(h(y,y))\le0$ for some $\tau$. The vectors $f_1,\dots,f_r,z',y$ are linearly independent. Indeed
$y\ne0$, as $y\notin X'$; and if $w=sz'+s'y$ lies in $X'$, then $0=h(z',w)=s\,h(z',z')$, so $s=0$,
and then $s'=0$, as $y\notin X'$; so $w=0$, and the span of $z',y$ meets $X'$ in $0$. On the span of
these $r+2$ vectors the form $\tau\circ h$ is negative semidefinite: the $f_a$ are isotropic and
orthogonal to $f_b$, $z'$ and $y$, and $h(z',y)=0$, so the vector with coordinates $c_a$, $c$, $c'$
has $\tau h$-norm $|c|^2\tau(h(z',z'))+|c'|^2\tau(h(y,y))\le0$. By Sylvester's law of inertia, a
subspace on which $\tau\circ h$ is negative semidefinite has dimension at most the number $r_S+1$ of
negative eigenvalues; but $r+2=r_S+2$. This contradiction gives $\tau(h(y,y))>0$.

Every element of $\operatorname{St}(F')$ that preserves $\Lambda'$ fixes $z'$, preserves $X'$ and so
$X'^\perp$, and induces its map of $X'^\perp/X'$ as in the lemma. So $A(\Lambda')$ lies in the finite
set of Lemma~\ref{asm-division-kernel-finite-fibres-induced}, and by Step 3 finitely many terms lie
over the fixed term of the unfolding.
% Lean: Steps 1 to 3 are fib_decomp (dc_of_toTerm, mem_NLam_map, stabF_of_fix and the uniqueness in
% IsRowReps = Y.rowsOK), fib_data and cOf_spec, phi (indW, indW_mem: the induced map lies in autLz
% D (X.Fu u') lam'), and phi_inj (indW_eq, conj_fix_mem, nlam_mul, nlam_inv, gterm_eq); fibFin_holds
% concludes with Finite.of_injective. Step 4 is automorphism_lattice_finite D (D.sp.Y_sig 0) _
% (D.mem u').A.T_prime.pos (X.Fu u') _ hrow = autLzFin_main (AutLzFinMain.lean): IsRowLat gives the
% lattice, unimodularity (dual = Lambda) and z in Lambda (FixedHyp: hlat, hunimod, hlift.1); hstar is
% IsCMField.complexEmbedding_complexConj; hpos is pos_perp (MulVecMapInjective.lean) with the signature
% sigY0 = (k + rS, rS + 1) (D.sp.baseSpace_signature_eq) and sylvester_le (SylvesterLe.lean); then
% indSet_finite (Lemma asm-division-kernel-finite-fibres-induced), and autLz lies in indSet.
\end{proof}

\begin{definition}\label{asm-division-kernel-finite-pairs-invariants}
Let $V$ be a finite-dimensional Hermitian space over a CM field $E$, with form $h$, let
$\Lambda\subset V$ be an $\mathcal O_E$-lattice, and let $k\ge0$ be an integer. A \emph{flag} $F$ of
length $r$ consists of linearly independent, mutually orthogonal isotropic vectors $e_1,\dots,e_r$ and a
vector $z_F$ orthogonal to them. Put $X_j(F)=\langle e_1,\dots,e_j\rangle$, $X(F)=X_r(F)$,
$X_+(F)=X(F)+Ez_F$ and $X_\Lambda(F)=\Lambda\cap X(F)$, and let $N_\Lambda(F)$ be the group of
isometries $n$ of $V$ with $n\Lambda=\Lambda$, $(n-1)X_{j+1}(F)\subset X_j(F)$ for every $j$, and
$(n-1)X(F)^\perp\subset X(F)$. For flags $F$ and $F'$ let
\[
 U(F,F')=X_\Lambda(F)+\mathcal O_Ez_F+X_\Lambda(F')+\mathcal O_Ez_{F'},
\]
$\bar U(F,F')$ its $E$-span and $H(F,F')=\bar U(F,F')^\perp$. Let $F,F',F''$ be flags and $\beta$ an
isometry of $V$.
\begin{enumerate}
\item $\Theta(F,F',F'',\beta)$ is the space of the $y\in H(F,F')$ for which there is $y''\in H(F,F'')$
  with $y-y''\in X(F)$ and $\beta y-y''\in X(F'')$; and $s_0(F,F',F'',\beta)=\max\big(k-\dim
  \Theta(F,F',F'',\beta),0\big)$.
\item $(F',F'',\beta)$ is \emph{scalar} for $F$ if for all $y_1\in\bar U(F,F')\cap X_+(F)^\perp$,
  $y_2\in\bar U(F,F'')\cap X_+(F)^\perp$, $y'_1\in\bar U(F,F')\cap X_+(F')^\perp$ and $y''_2\in\bar
  U(F,F'')\cap X_+(F'')^\perp$ with $y_1-y_2\in X(F)$ and $y''_2-\beta y'_1\in X(F'')$, and all
  $\alpha,\alpha'\in E$,
  \[
   y_1-y'_1\in\alpha z_F+\alpha'z_{F'}+X(F)+X(F')\iff y_2-y''_2\in\alpha z_F+\alpha'z_{F''}+X(F)+X(F'').
  \]
\item In Situation~\ref{asm-division-kernel-finite-situation}, let $u,u'$ be words with flags $F$ and
  $F'$, and let $t,t'$ be terms of $T_{uu'}$ with the same rows $\Lambda$ of $u$ and $\Lambda'$ of $u'$
  and with isometries $\gamma$ and $\gamma'$ (Definition~\ref{is-unfolding-evaluation}). Their
  \emph{corank} is $s_0(t,t')=s_0(F,\gamma F',\gamma'F',\gamma'\gamma^{-1})$, with $k$ the rank $k$ of $D$,
  and $(t,t')$ is \emph{scalar} if $s_0(t,t')=1$ and $(\gamma F',\gamma'F',\gamma'\gamma^{-1})$ is scalar
  for $F$.
\end{enumerate}
\end{definition}
\leanref{Formalism/SharpRestriction/CuspSpanNormBounds/IsPrimitiveSub.lean}{Formalism.MainNotions.S.Lat.NLam}
\leanref{Formalism/SharpRestriction/CuspSpanNormBounds/IsPrimitiveSub.lean}{Formalism.MainNotions.S.Lat.UU}
\leanref{Formalism/SharpRestriction/CuspSpanNormBounds/IsPrimitiveSub.lean}{Formalism.MainNotions.S.Lat.Hp}
\leanref{Formalism/SharpRestriction/CuspSpanNormBounds/IsPrimitiveSub.lean}{Formalism.MainNotions.S.Lat.Th}
\leanref{Formalism/SharpRestriction/CuspSpanNormBounds/IsPrimitiveSub.lean}{Formalism.MainNotions.S.Lat.s0}
\leanref{Formalism/SharpRestriction/CuspSpanNormBounds/IsPrimitiveSub.lean}{Formalism.MainNotions.S.Lat.Ylift}
\leanref{Formalism/SharpRestriction/CuspSpanNormBounds/IsPrimitiveSub.lean}{Formalism.MainNotions.S.Lat.SgRel}
\leanref{Formalism/SharpRestriction/CuspSpanNormBounds/IsPrimitiveSub.lean}{Formalism.MainNotions.S.Lat.IsScalar}
\leanref{Formalism/LocalToGlobal/SiegelLatticeTheory/ClosedLocalOrbits/TermOp.lean}{Formalism.L2.SG.TermOp.s0T}
\leanref{Formalism/LocalToGlobal/SiegelLatticeTheory/ClosedLocalOrbits/TermOp.lean}{Formalism.L2.SG.TermOp.scalT}

\begin{lemma}\label{asm-division-kernel-finite-pairs-move}
Let $V$, $\Lambda$ and $k$ be as in Definition~\ref{asm-division-kernel-finite-pairs-invariants}. Let
$F$, $G_1$ and $G_2$ be flags with $z_F,z_{G_1},z_{G_2}\in\Lambda$, let $n_1,n_2\in N_\Lambda(F)$ and
$m_i\in N_\Lambda(G_i)$, and let $K_i$ be the flag with vectors $n_im_ie_a$ and $n_im_iz_{G_i}$, where
$e_a$ are the vectors of $G_i$. Let $\beta$ be an isometry with $\beta X(G_1)=X(G_2)$, and $\tau$ an
isometry with $\tau X(G_1)\subset X(G_1)$ and $(\tau-1)X(G_1)^\perp\subset X(G_1)$. Put
$\beta'=n_2m_2\beta\tau m_1^{-1}n_1^{-1}$.
\begin{enumerate}
\item If $X(F)$ and $X(G_1)$ lie in the $E$-spans of $X_\Lambda(F)$ and $X_\Lambda(G_1)$, then
  $\Theta(F,K_1,K_2,\beta')=n_1\Theta(F,G_1,G_2,\beta)$.
\item If $(G_1,G_2,\beta)$ is scalar for $F$, so is $(K_1,K_2,\beta')$.
\end{enumerate}
\end{lemma}
\leanref{Formalism/LocalToGlobal/SiegelLatticeTheory/FlagStabilizerTensors/RootLatticeThresholds/XLeSpanEXLam.lean}{Formalism.L2.SG.Pf.FinXP.Th_move}
\leanref{Formalism/LocalToGlobal/SiegelLatticeTheory/FlagStabilizerTensors/RootLatticeThresholds/XLeSpanEXLam.lean}{Formalism.L2.SG.Pf.FinXP.scal_move}
\leanref{Formalism/LocalToGlobal/SiegelLatticeTheory/FlagStabilizerTensors/RootLatticeThresholds/XLeSpanEXLam.lean}{Formalism.L2.SG.Pf.FinXP.key_mod}
\leanref{Formalism/LocalToGlobal/SiegelLatticeTheory/FlagStabilizerTensors/RootLatticeThresholds/XLeSpanEXLam.lean}{Formalism.L2.SG.Pf.FinXP.sg_move}
\leanref{Formalism/LocalToGlobal/SiegelLatticeTheory/ClosedLocalOrbits/GaugeLatticeShifts/SaturatedPf.lean}{Formalism.L2.SG.Pf.UU_DC}
\begin{proof}
\emph{Step 1: facts on $N_\Lambda$.} Let $G$ be a flag with $z_G\in\Lambda$ and $n\in N_\Lambda(G)$.
Then $nX(G)\subset X(G)$, as $(n-1)e_a\in X(G)$, so $nX(G)=X(G)$; and $n^{-1}\in N_\Lambda(G)$. As
$z_G\in X(G)^\perp$, $nz_G-z_G\in X(G)$, and $nz_G\in\Lambda$, so $nz_G-z_G\in X_\Lambda(G)$. Hence $n$
maps $X_\Lambda(G)$ onto itself and $X_\Lambda(G)+\mathcal O_Ez_G$ onto itself, and $nX_+(G)=X_+(G)$.

\emph{Step 2: transport.} Let $K=K_i$, $G=G_i$, $n=n_i$, $m=m_i$. Then $X(K)=nmX(G)=nX(G)$ by Step 1, so
$X_\Lambda(K)=nX_\Lambda(G)$, as $n\Lambda=\Lambda$; and $z_K=nmz_G$, so by Step 1 for $m$ and $G$, and
for $n$ and $F$,
\[
 U(F,K)=n\big(X_\Lambda(F)+\mathcal O_Ez_F\big)+n\big(X_\Lambda(G)+\mathcal O_Emz_G\big)=nU(F,G).
\]
So $\bar U(F,K)=n\bar U(F,G)$ and, $n$ being an isometry, $H(F,K)=nH(F,G)$. Also
$z_K-nz_G=n(mz_G-z_G)\in nX(G)=X(K)$; and $nm$ maps $X_+(G)$ into $X_+(K)$, so $n$ does too, as
$mX_+(G)=X_+(G)$.

\emph{Step 3: the key congruence.} For $x\in X(G_1)^\perp$,
$m_2\beta\tau m_1^{-1}x-\beta x\in X(G_2)$. Indeed $m_1^{-1}x-x\in X(G_1)$, so
$\tau m_1^{-1}x-x=(\tau x-x)+\tau(m_1^{-1}x-x)\in X(G_1)$ and $w=\beta\tau m_1^{-1}x$ satisfies
$w-\beta x\in X(G_2)$. As $\beta$ is an isometry with $\beta X(G_1)=X(G_2)$, $\beta x\in X(G_2)^\perp$,
and $X(G_2)\subset X(G_2)^\perp$, so $w\in X(G_2)^\perp$ and $m_2w-w\in X(G_2)$. Adding,
$m_2w-\beta x\in X(G_2)$.

(1) By hypothesis $X(F)\subset\bar U(F,G)$ for every flag $G$, so $H(F,G)\subset X(F)^\perp$; likewise
$H(F,G_1)\subset X(G_1)^\perp$. For $x_1,x_2\in X(F)^\perp$, $n_1x_1-n_2x_2\in X(F)$ if and only if
$x_1-x_2\in X(F)$, as $n_ix_i-x_i\in X(F)$. Let $y\in V$, $x_1=n_1^{-1}y$, and for $y_2\in V$ put
$x_2=n_2^{-1}y_2$. By Step 2, $y\in H(F,K_1)$ and $y_2\in H(F,K_2)$ if and only if $x_1\in H(F,G_1)$ and
$x_2\in H(F,G_2)$; then $y-y_2\in X(F)$ if and only if $x_1-x_2\in X(F)$; and
\[
 \beta'y-y_2=n_2\big(m_2\beta\tau m_1^{-1}x_1-x_2\big)
\]
lies in $X(K_2)=n_2X(G_2)$ if and only if $m_2\beta\tau m_1^{-1}x_1-x_2\in X(G_2)$, that is, by Step 3,
if and only if $\beta x_1-x_2\in X(G_2)$. So $y\in\Theta(F,K_1,K_2,\beta')$ if and only if
$x_1\in\Theta(F,G_1,G_2,\beta)$.

(2) For $y,y'\in V$, $\alpha,\alpha'\in E$ and flags $G$, $K$ with $X(K)=nX(G)$ and $z_K-nz_G\in X(K)$
for some $n\in N_\Lambda(F)$: $ny-ny'\in\alpha z_F+\alpha'z_K+X(F)+X(K)$ if and only if
$y-y'\in\alpha z_F+\alpha'z_G+X(F)+X(G)$, by applying $n^{-1}$ or $n$ and using $n^{\pm1}z_F-z_F\in
X(F)$, $n^{\pm1}X(F)=X(F)$ and $z_K-nz_G\in X(K)$. Now let $y_1,y_2,y'_1,y''_2$ and $\alpha,\alpha'$ be
as in Definition~\ref{asm-division-kernel-finite-pairs-invariants}(2) for $(K_1,K_2,\beta')$. By Steps 1
and 2, $n_1^{-1}y_1\in\bar U(F,G_1)\cap X_+(F)^\perp$, $n_2^{-1}y_2\in\bar U(F,G_2)\cap
X_+(F)^\perp$, $n_1^{-1}y'_1\in\bar U(F,G_1)\cap X_+(G_1)^\perp$ and $n_2^{-1}y''_2\in\bar U(F,G_2)\cap
X_+(G_2)^\perp$, since $n_i$ is an isometry mapping $X_+(F)$ to $X_+(F)$ and $X_+(G_i)$ into $X_+(K_i)$.
As in (1), $n_1^{-1}y_1-n_2^{-1}y_2\in X(F)$, and
$n_2^{-1}y''_2-\beta n_1^{-1}y'_1=\big(n_2^{-1}y''_2-m_2\beta\tau m_1^{-1}n_1^{-1}y'_1\big)
+\big(m_2\beta\tau m_1^{-1}n_1^{-1}y'_1-\beta n_1^{-1}y'_1\big)$ lies in $X(G_2)$ by the hypothesis and
Step 3. The scalarity of $(G_1,G_2,\beta)$ applied to these four vectors, with the equivalence above for
$(G_1,K_1,n_1)$ and $(G_2,K_2,n_2)$, gives the required equivalence.
% Lean: (1) is Th_move, with Hp_move (UUE_move from UU_DC = UU_eq, Apart_DC, NLam_Apart, NLam_XLam,
% XLam_DC, NLamAlt, sup_span_shift), Xsp_DC (NLam_map_X), Hp_le_orthF, Hp_le_orthG, key_mod (Step 3)
% and hXX; mem_Th and mem_ThPairs unfold Theta. (2) is scal_move, with UUE_move, Xsp_DC, hzK, hpF and
% hpG (xplus_nlam, xplus_frame, orth_pull), key_mod and sg_move (the equivalence). The Lean's flags
% carry a vector zt with IsLift (z in Lambda and zt = z); the book's z_F is that vector. The Lean states
% (1) with the spans as hypotheses hsF and hs1, as the book does.
\end{proof}

\begin{lemma}\label{asm-division-kernel-finite-pairs-kernel-inputs-lattice}
In Situation~\ref{asm-division-kernel-finite-situation}, with the notation of
Definition~\ref{asm-division-kernel-finite-series-function-bessel}, regard $V$ as a subset of $\hat V$, and
for $y\in G_{\mathrm f}$ put $\Lambda(y)=\{v\in V:y^{-1}v\in\hat L\}$.
\begin{enumerate}
\item $\hat L\cap V=L_0$.
\item If $a,b\in G_{\mathrm f}$ and $\Lambda(a)=\Lambda(b)$, then $a^{-1}b\in K_{\mathrm{set}}$.
\item Let $u$ be a word, $h\in H_{\mathrm f}$ and $y=hx_u$, and put
  $\mathrm{pr}(v)=v-\sum_ah(f_a,v)e_a-\sum_ah(e_a,v)f_a$. Then $v\in\Lambda(y)$ if and only if
  $v-\mathrm{pr}(v)\in L_0$ and $\mathrm{pr}(v)\in\Lambda(y)$.
\end{enumerate}
\end{lemma}
\leanref{Formalism/PeriodIntegrals/UnitaryGroupMassFormula/AdelicClassFiniteness/DictAll.lean}{PIW.AdelClassFin.lattice_intersection'}
\leanref{Formalism/LocalToGlobal/SiegelLatticeTheory/FlagStabilizerTensors/RootLatticeThresholds/LatDictK.lean}{Formalism.L2.SG.Pf.LatDictK.dict_K}
\leanref{Formalism/LocalToGlobal/SiegelLatticeTheory/FlagStabilizerTensors/RootLatticeThresholds/LatDict.lean}{Formalism.L2.SG.Pf.LatDict.mem_latOf_split}
\begin{proof}
(3) The map $\mathrm{pr}$ is $\mathbf A_{E,\mathrm f}$-linear on $\hat V$ and maps $\hat L$ into $\hat L$, by
the proof of Lemma~\ref{asm-division-kernel-finite-series-function-inputs}(1). As $h$ and $x_u$ fix every
$e_a$ and $f_a$ (Definition~\ref{asm-division-kernel-finite-series-function-inputs-level}(3)), $y^{-1}$ is an
isometry fixing them; so $y^{-1}$ commutes with $\mathrm{pr}$ and fixes $v-\mathrm{pr}(v)$, a combination of
the $e_a$ and $f_a$. If $y^{-1}v\in\hat L$, then $y^{-1}\mathrm{pr}(v)=\mathrm{pr}(y^{-1}v)\in\hat L$, and
$v-\mathrm{pr}(v)=y^{-1}v-y^{-1}\mathrm{pr}(v)$ lies in $\hat L\cap V=L_0$ by (1). Conversely
$y^{-1}v=\big(v-\mathrm{pr}(v)\big)+y^{-1}\mathrm{pr}(v)$ lies in $\hat L$ if both terms do.

(1) and (2): \openstep{That the rational points of the adelic lattice form the lattice $L_0$, and that an
adelic lattice $a\hat L$ is determined by its rational points, are not yet written out.}
% Lean: lattice_intersection' (DictAll.lean) and dict_K (LatDictK.lean) through
% PIW.AdelClassFin.lattice_map_injective (abelian2 a = abelian2 b gives k in Kset with b = a k). Size
% (tools/lean_size.py, the union of lattice_intersection' and dict_K): 978 lines. (3) is mem_latOf_split
% with prA_comm, fix_sub_prA, prA_Lf, mem_base_lattice_of_adelic and adelic_mem_of_base_lattice.
\end{proof}

\begin{lemma}\label{asm-division-kernel-finite-pairs-kernel-inputs-reduction}
In Situation~\ref{asm-division-kernel-finite-situation}, with the notation of
Definitions~\ref{asm-division-kernel-finite-series-function-bessel}
and~\ref{asm-division-kernel-finite-series-function-inputs-level}, and with $\langle S\rangle$ as in
Lemma~\ref{asm-division-kernel-finite-series-function-inputs-congruence}, let $\mathrm{res}_v\colon\hat{\mathcal
O}\to\mathbf F_{q_v}$ be the reduction maps of the setting (Definition~\ref{is-unfolding-evaluation}(1)). For
$\kappa\in G_{\mathrm f}$ let $M_v(\kappa)$ be the $k\times k$ matrix
$\big(\mathrm{res}_v\,h(b'_i,\kappa b_j)\big)_{ij}$ over $\mathbf F_{q_v}$, and let $\lambda_v(\kappa)$ be
$M_v(\kappa)$ if it is invertible and $1$ otherwise; put $\lambda(\kappa)=(\lambda_v(\kappa))_v\in I$. Call
$\kappa$ \emph{parabolic} if $\kappa\in K_{\mathrm{set}}$, $\kappa e^0_a\in\langle e^0_c:c\le r\rangle$ for
every $a$ and $\kappa b_j\in\langle e^0_c,b_l:c\le r,\ l\le k\rangle$ for every $j$, and \emph{good} if
$\kappa$ and $\kappa^{-1}$ are parabolic.
\begin{enumerate}
\item If $\kappa$ and $\kappa'$ are good, so are $\kappa\kappa'$ and $\kappa^{-1}$, and
  $\lambda(\kappa\kappa')=\lambda(\kappa)\lambda(\kappa')$ and $\lambda(\kappa^{-1})=\lambda(\kappa)^{-1}$.
\item Let $u$ be a word and $h,h'\in H_{\mathrm f}$ with $\kappa=(hx_u)^{-1}(h'x_u)\in K_{\mathrm{set}}$. Then
  $\kappa$ is good, $(\kappa-1)e^0_a\in\langle e^0_c:c<a\rangle$ for every $a$, and
  $(\kappa-1)z^0\in\langle e^0_c:c\le r\rangle$.
\item Let $\kappa$ be good, with $(\kappa-1)e^0_a\in\langle e^0_c:c<a\rangle$ for every $a$ and
  $(\kappa-1)z^0\in\langle e^0_c:c\le r\rangle$, and suppose that every $\lambda_v(\kappa)$ is upper
  unitriangular. Then $\kappa\in K_2$.
\end{enumerate}
\end{lemma}
\leanref{Formalism/LocalToGlobal/SiegelLatticeTheory/ClosedLocalOrbits/TermOp.lean}{Formalism.L2.SG.TermOp.leviM}
\leanref{Formalism/LocalToGlobal/SiegelLatticeTheory/ClosedLocalOrbits/TermOp.lean}{Formalism.L2.SG.TermOp.levi}
\leanref{Formalism/LocalToGlobal/SiegelLatticeTheory/ClosedLocalOrbits/TermOp.lean}{Formalism.L2.SG.TermOp.TypedRed}
\leanref{Formalism/LocalToGlobal/SiegelLatticeTheory/ClosedLocalOrbits/GaugeLatticeShifts/LeviColsPf.lean}{Formalism.L2.SG.NormCut.levi_mul}
\leanref{Formalism/LocalToGlobal/SiegelLatticeTheory/ClosedLocalOrbits/GaugeLatticeShifts/LeviColsPf.lean}{Formalism.L2.SG.NormCut.levi_inv}
\leanref{Formalism/LocalToGlobal/SiegelLatticeTheory/ClosedLocalOrbits/GaugeLatticeShifts/LeviColsPf.lean}{Formalism.L2.SG.NormCut.colGood}
\leanref{Formalism/LocalToGlobal/SiegelLatticeTheory/ClosedLocalOrbits/GaugeLatticeShifts/PhaseAtPf.lean}{Formalism.L2.SG.NormCut.ncong_colDiff}
\leanref{Formalism/LocalToGlobal/SiegelLatticeTheory/ClosedLocalOrbits/GaugeLatticeShifts/PhaseAtPf.lean}{Formalism.L2.SG.NormCut.mem_level_of_unitriangular}
\begin{proof}
\openstep{The multiplicativity of the reduction of the Levi block on good elements, the goodness and
congruences of the differences of column representatives, and the criterion for membership in the level
subgroup, are not yet written out.}
% Lean: (1) levi_mul and levi_inv (LeviColsPf.lean) with leviM_mul_par, levi_coe, good_mul, good_inv,
% good_one (ParCols.lean: Par and Good); (2) colGood (colStab, par_of_stab) and ncong_colDiff (NCong);
% (3) mem_level_of_unitriangular (PhaseAtPf.lean: dual_coords, Z_pair_Mo, the kernel condition res_ker
% of TypedRed and ew_orth). Size (tools/lean_size.py, union of levi_mul, levi_inv, colGood, ncong_colDiff
% and mem_level_of_unitriangular): 1,117 lines.
\end{proof}

\begin{lemma}\label{asm-division-kernel-finite-pairs-kernel-inputs-terms}
In Situation~\ref{asm-division-kernel-finite-situation}, with the notation of
Definition~\ref{asm-division-kernel-finite-series-function-bessel} and of
Lemma~\ref{asm-division-kernel-finite-pairs-kernel-inputs-reduction}, let $u,u'$ be words. For a term $t$
of $T_{uu'}$ with rows $\Lambda$ and $\Lambda'$, base column representatives $h_\Lambda$ and
$h_{\Lambda'}$, and an adelic point $x$ linked to a representative of its double coset
(Definition~\ref{is-unfolding-evaluation}), put $k(x)=(h_\Lambda x_u)^{-1}x_{\mathrm f}(h_{\Lambda'}x_{u'})$.
\begin{enumerate}
\item $k(x)$ is good, and $\lambda(k(x))$ does not depend on the choice of the representative and of $x$;
  it is $\iota(t)$.
\item Let $\Lambda'$ be a row of $u'$, put $y=h_{\Lambda'}x_{u'}$, let $\delta\in\Gamma^H_{u'}$ with
  $y^{-1}\delta_{\mathrm f}y\in K_{\mathrm{set}}$ and let $\bar\delta$ be the isometry of $V$ that defines
  $\delta$. Let $s$ be a term of $T_{uu'}$ whose row of $u'$ is $\Lambda'$, with representative $\gamma_s$
  linked to $x_s$. Then $\gamma_s\bar\delta^{-1}$ is linked to $x_s\delta^{-1}$, and there is a term
  $\delta\cdot s$ of $T_{uu'}$ with the rows of $s$ and the double coset of $\gamma_s\bar\delta^{-1}$.
\end{enumerate}
\end{lemma}
\leanref{Formalism/LocalToGlobal/SiegelLatticeTheory/ClosedLocalOrbits/GaugeLatticeShifts/SatGramActRRep.lean}{Formalism.L2.SG.Pf.SatGramRep.IotaRepC}
\leanref{Formalism/LocalToGlobal/SiegelLatticeTheory/ClosedLocalOrbits/GaugeLatticeShifts/AmpIota.lean}{Formalism.L2.SG.Pf.iotaRepC_holds}
\leanref{Formalism/LocalToGlobal/SiegelLatticeTheory/FlagStabilizerTensors/RootLatticeThresholds/ToTermAct.lean}{Formalism.L2.SG.Pf.KernelOfEmbedding.act}
\leanref{Formalism/LocalToGlobal/SiegelLatticeTheory/FlagStabilizerTensors/RootLatticeThresholds/ToTermAct.lean}{Formalism.L2.SG.Pf.KernelOfEmbedding.isoG_lift}
\begin{proof}
\openstep{That the Levi reduction of a term does not depend on the representative of its double coset,
that the adelic points of terms are good, and that the rational points fixing the flag and the row act on
the terms, are not yet written out.}
% Lean: (1) iotaRepC_holds (AmpIota.lean: rep_split, nilradical_conj_level', good_of_ncong, termK_car,
% adm_norm_of_carrier, cong_b and the Levi of a unipotent element) and Closed.AdmPar_holds; (2) act =
% SatGramActR.tmap with autFConjN_holds, isoG_lift (SatGramActL.betaLift_holds) and tmap_cls. Size
% (tools/lean_size.py, union of iotaRepC_holds, AdmPar_holds, satGramActLIota_kB_good, betaLift_holds,
% tmap, autFConjN_holds and tmap_cls): 3,262 lines; iotaRepC_holds alone 1,877, AdmPar_holds alone 1,420.
\end{proof}

\begin{lemma}\label{asm-division-kernel-finite-pairs-kernel-inputs}
In Situation~\ref{asm-division-kernel-finite-situation}, let $u'$ be a word, with base point $x_{u'}$, the
groups $H_{\mathrm f}$ and $\Gamma^H_{u'}$ of
Definition~\ref{asm-division-kernel-finite-series-function-bessel}, and the flag $F'$ whose vectors are the
vectors $e_1,\dots,e_r$ and $z_{u'}$ of that definition; put $X'=X(F')$, the span of the $e_a$
(Definition~\ref{asm-division-kernel-finite-pairs-invariants}). Let $\Lambda'$ be
a row of $u'$, with base column representative $h_{\Lambda'}\in H_{\mathrm f}$, and put
$y_{\Lambda'}=h_{\Lambda'}x_{u'}$. Let $G(\Lambda')$ be the set of $\delta\in\Gamma^H_{u'}$ with
$\kappa(\delta)=y_{\Lambda'}^{-1}\delta_{\mathrm f}y_{\Lambda'}\in K_{\mathrm{set}}$. For $\delta\in
G(\Lambda')$ write $\bar\delta$ for the isometry of $V$ that defines $\delta$, and
$\lambda(\delta)=\lambda(\kappa(\delta))\in I=\prod_vGL_k(\mathbf F_{q_v})$ for the reduction at the places
$v$ of the Levi block of $\kappa(\delta)$
(Lemma~\ref{asm-division-kernel-finite-pairs-kernel-inputs-reduction}). Let $\operatorname{St}(F')$ be the
group of isometries of $V$ that preserve every $X_j(F')$ and fix $z_{u'}$.
\begin{enumerate}
\item $G(\Lambda')$ is a group, $\overline{\delta\delta'}=\bar\delta\bar\delta'$, and $\bar\delta$ fixes the
  vectors of $F'$ and preserves $\Lambda'$. The map sending $\delta$ to the map induced by $\bar\delta$ on
  $X'^\perp/X'$ is a bijection of $G(\Lambda')$ onto the set $A(\Lambda')$ of the maps of $X'^\perp/X'$
  induced by the elements of $\operatorname{St}(F')$ that preserve $\Lambda'$.
\item Let $u$ be a word, $s$ a term of $T_{uu'}$ whose row of $u'$ is $\Lambda'$, and $\delta\in
  G(\Lambda')$. There is a term $\delta\cdot s$ of $T_{uu'}$ with the rows of $s$ and the double coset of
  $\gamma_s\bar\delta^{-1}$, and $\iota(\delta\cdot s)=\iota(s)\lambda(\delta)^{-1}$.
\item The map $\delta\mapsto\lambda(\delta)$ is injective on $G(\Lambda')$.
\end{enumerate}
\end{lemma}
\leanref{Formalism/LocalToGlobal/SiegelLatticeTheory/ClosedLocalOrbits/GaugeLatticeShifts/ActMass.lean}{Formalism.L2.SG.Pf.Act.GamL}
\leanref{Formalism/LocalToGlobal/SiegelLatticeTheory/ClosedLocalOrbits/GaugeLatticeShifts/ActMass.lean}{Formalism.L2.SG.Pf.Act.isoG}
\leanref{Formalism/LocalToGlobal/SiegelLatticeTheory/ClosedLocalOrbits/GaugeLatticeShifts/ActMass.lean}{Formalism.L2.SG.Pf.Act.gamL_autLz_bij}
\leanref{Formalism/LocalToGlobal/SiegelLatticeTheory/ClosedLocalOrbits/GaugeLatticeShifts/ActMass.lean}{Formalism.L2.SG.Pf.Act.gamL_levi_bij}
\leanref{Formalism/LocalToGlobal/SiegelLatticeTheory/ClosedLocalOrbits/GaugeLatticeShifts/SatGramActRRep.lean}{Formalism.L2.SG.Pf.SatGramRep.IotaRepC}
\leanref{Formalism/LocalToGlobal/SiegelLatticeTheory/ClosedLocalOrbits/GaugeLatticeShifts/AmpIota.lean}{Formalism.L2.SG.Pf.iotaRepC_holds}
\leanref{Formalism/LocalToGlobal/SiegelLatticeTheory/FlagStabilizerTensors/RootLatticeThresholds/ToTermAct.lean}{Formalism.L2.SG.Pf.KernelOfEmbedding.act}
\leanref{Formalism/LocalToGlobal/SiegelLatticeTheory/FlagStabilizerTensors/RootLatticeThresholds/ToTermAct.lean}{Formalism.L2.SG.Pf.KernelOfEmbedding.iota_act}
\leanref{Formalism/LocalToGlobal/SiegelLatticeTheory/FlagStabilizerTensors/RootLatticeThresholds/ToTermAct.lean}{Formalism.L2.SG.Pf.KernelOfEmbedding.isoG_lift}
\leanref{Formalism/LocalToGlobal/SiegelLatticeTheory/ClosedLocalOrbits/TermOp.lean}{Formalism.L2.SG.TermOp.iotaOf}
\leanref{Formalism/LocalToGlobal/SiegelLatticeTheory/ClosedLocalOrbits/TermOp.lean}{Formalism.L2.SG.TermOp.levi}
\begin{proof}
Write $y=y_{\Lambda'}$, $h=h_{\Lambda'}$ and $\mathrm{pr}(v)=v-\sum_ah(f_a,v)e_a-\sum_ah(e_a,v)f_a$, with the
vectors $f_a$ of Definition~\ref{asm-division-kernel-finite-series-function-bessel}. By the column system of
the setting, $\Lambda'=\Lambda(y)$ in the notation of
Lemma~\ref{asm-division-kernel-finite-pairs-kernel-inputs-lattice}. For $\delta\in\Gamma^H_{u'}$,
$\kappa(\delta)=(hx_{u'})^{-1}(\delta_{\mathrm f}hx_{u'})$ with $\delta_{\mathrm f}h\in H_{\mathrm f}$.

(1) The map $\kappa$ is multiplicative, and $K_{\mathrm{set}}$ is a group, so $G(\Lambda')$ is a group. The
isometry $\bar\delta$ is the rational point $\delta$ itself, so $\delta\mapsto\bar\delta$ is multiplicative,
and $\bar\delta$ fixes the $e_a$, $f_a$ and $z_{u'}$; in particular $\bar\delta\in\operatorname{St}(F')$. For
$v\in V$, $y^{-1}\bar\delta v=\kappa(\delta)y^{-1}v$, and $\kappa(\delta)^{\pm1}$ preserve $\hat L$; so
$\bar\delta v\in\Lambda'$ if and only if $v\in\Lambda'$, and $\bar\delta\Lambda'=\Lambda'$. So the induced map
lies in $A(\Lambda')$.

\emph{Injectivity.} Let $\delta_1,\delta_2$ induce the same map, and let $v\in V$. Both $\bar\delta_i$ fix
$v-\mathrm{pr}(v)$, a combination of the $e_a$ and $f_a$. The vector $\mathrm{pr}(v)$ is orthogonal to the
$e_a$, so $w=\bar\delta_1\mathrm{pr}(v)-\bar\delta_2\mathrm{pr}(v)\in X'$; and
$h(f_a,\bar\delta_i\mathrm{pr}(v))=h(\bar\delta_i^{-1}f_a,\mathrm{pr}(v))=h(f_a,\mathrm{pr}(v))=0$, so $w$ is
orthogonal to every $f_a$. A vector of $X'$ orthogonal to every $f_a$ is $0$, as $h(f_a,e_b)=\delta_{ab}$.
So $\bar\delta_1=\bar\delta_2$, the finite component of $\delta_1^{-1}\delta_2$ is $1\in K_{u'}$, and
$\delta_1=\delta_2$ by the neatness of $K_{u'}$ (Lemma~\ref{asm-division-kernel-finite-neat}).

\emph{Surjectivity.} Let $\varphi\in A(\Lambda')$ be induced by $\gamma_0\in\operatorname{St}(F')$ with
$\gamma_0\Lambda'=\Lambda'$, and put $c(v)=\big(v-\mathrm{pr}(v)\big)+\mathrm{pr}\big(\gamma_0\mathrm{pr}(v)\big)$.

(a) For $w\in V$, $\gamma_0\mathrm{pr}(w)-\mathrm{pr}(\gamma_0\mathrm{pr}(w))\in X'$: it is
$\sum_ah(f_a,\gamma_0\mathrm{pr}(w))e_a+\sum_ah(e_a,\gamma_0\mathrm{pr}(w))f_a$, and
$h(e_a,\gamma_0\mathrm{pr}(w))=h(\gamma_0^{-1}e_a,\mathrm{pr}(w))=0$, as $\gamma_0^{-1}$ preserves $X'$.

(b) $c$ is an isometry fixing the $e_a$, $f_a$ and $z_{u'}$. As $h(v,v')=h(v-\mathrm{pr}(v),v'-\mathrm{pr}(v'))
+h(\mathrm{pr}(v),\mathrm{pr}(v'))$, $\mathrm{pr}(c(v))=\mathrm{pr}(\gamma_0\mathrm{pr}(v))$ and
$c(v)-\mathrm{pr}(c(v))=v-\mathrm{pr}(v)$, it suffices that
$h\big(\mathrm{pr}(\gamma_0\mathrm{pr}(v)),\mathrm{pr}(\gamma_0\mathrm{pr}(v'))\big)=h(\mathrm{pr}(v),\mathrm{pr}(v'))$.
This follows from the same decomposition of $h(\gamma_0\mathrm{pr}(v),\gamma_0\mathrm{pr}(v'))$, by (a), as
$X'$ is isotropic and $\gamma_0$ is an isometry. The map $\mathrm{pr}$ vanishes on the $e_a$ and $f_a$, and
fixes $z_{u'}$, which $\gamma_0$ fixes. The same construction for $\gamma_0^{-1}$ gives the inverse of $c$:
by (a) and $\gamma_0X'=X'$, $\mathrm{pr}(\gamma_0^{-1}\mathrm{pr}(\gamma_0\mathrm{pr}(v)))=\mathrm{pr}(v)$.

(c) $c\Lambda'\subset\Lambda'$, and likewise for the inverse. Let $v\in\Lambda'$. By
Lemma~\ref{asm-division-kernel-finite-pairs-kernel-inputs-lattice}(3), $v-\mathrm{pr}(v)\in L_0$ and
$\mathrm{pr}(v)\in\Lambda'$; then $\gamma_0\mathrm{pr}(v)\in\Lambda'$, so
$\mathrm{pr}(\gamma_0\mathrm{pr}(v))\in\Lambda'$ by the same lemma; and the same lemma gives $c(v)\in\Lambda'$.

(d) By (b), $c$ is the isometry of a rational point $\delta\in\Gamma^H_{u'}$. By (c),
$\Lambda(\delta_{\mathrm f}y)=c\Lambda'=\Lambda'=\Lambda(y)$, so
$\kappa(\delta)=y^{-1}\delta_{\mathrm f}y\in K_{\mathrm{set}}$ by
Lemma~\ref{asm-division-kernel-finite-pairs-kernel-inputs-lattice}(2), and $\delta\in G(\Lambda')$.

(e) For $x\in X'^\perp$, $x-\mathrm{pr}(x)=\sum_ah(f_a,x)e_a\in X'$, so
$c(x)-\gamma_0x=\big(x-\mathrm{pr}(x)\big)-\gamma_0\big(x-\mathrm{pr}(x)\big)-\big(\gamma_0\mathrm{pr}(x)-\mathrm{pr}(\gamma_0\mathrm{pr}(x))\big)$
lies in $X'$ by (a). So $\bar\delta=c$ induces $\varphi$.

(2) Let $\Lambda$ be the row of $s$ of $u$, and $x_s$ the adelic point linked to $\gamma_s$. By
Lemma~\ref{asm-division-kernel-finite-pairs-kernel-inputs-terms}(2), the term $\delta\cdot s$ exists, with
$\gamma_s\bar\delta^{-1}$ linked to $x_s\delta^{-1}$. In the notation of that lemma,
$k(x_s\delta^{-1})=k(x_s)\,y^{-1}\delta_{\mathrm f}^{-1}y=k(x_s)\kappa(\delta)^{-1}$. Here $k(x_s)$ is good by
part (1) of that lemma, and $\kappa(\delta)$ is good by
Lemma~\ref{asm-division-kernel-finite-pairs-kernel-inputs-reduction}(2). By part (1) of the two lemmas,
$\iota(\delta\cdot s)=\lambda(k(x_s)\kappa(\delta)^{-1})=\lambda(k(x_s))\lambda(\kappa(\delta))^{-1}
=\iota(s)\lambda(\delta)^{-1}$.

(3) Let $\delta_1,\delta_2\in G(\Lambda')$ with $\lambda(\delta_1)=\lambda(\delta_2)$, and put
$h_i=\delta_{i,\mathrm f}h\in H_{\mathrm f}$. Then $\kappa(\delta_i)=(hx_{u'})^{-1}(h_ix_{u'})$, and
$m=\kappa(\delta_1)^{-1}\kappa(\delta_2)=(h_1x_{u'})^{-1}(h_2x_{u'})$ lies in $K_{\mathrm{set}}$. By
Lemma~\ref{asm-division-kernel-finite-pairs-kernel-inputs-reduction}(2), $\kappa(\delta_1)$, $\kappa(\delta_2)$
and $m$ are good, and $m$ satisfies the congruences of part (3) of that lemma; by its part (1),
$\lambda(m)=\lambda(\delta_1)^{-1}\lambda(\delta_2)=1$, which is unitriangular. By its part (3), $m\in K_2$.
As $m=x_{u'}^{-1}\big(h^{-1}\iota(\delta_1^{-1}\delta_2)h\big)x_{u'}$, $h^{-1}\iota(\delta_1^{-1}\delta_2)h\in
K_{u'}$, and the neatness of $K_{u'}$ (Lemma~\ref{asm-division-kernel-finite-neat}) gives
$\delta_1^{-1}\delta_2=1$.
% Lean: G(Lambda') is GamL (Y.cs u') q, kappa(delta) is SatGramActL.kB = colD h (iotaF g * h) (aK_eq),
% bar delta is isoG (gamIsom_of, isoG_spec); Lambda' = latOf is lam_iff (cs.lat_eq). (1) gamL_autLz_bij:
% isoG_e, isoG_f, isoG_fix, isoG_lat (kset_iff); injectivity is eq_of_indW (sub_prV_eq, eq_zero_of_Xs) with
% gf_eq and NeatKu at 1; surjectivity is gW (LatDict.lean): (a) stab_perp, (b) gW_isom (hT_decomp,
% hT_prV_stab, hT_Xs_Xs), gW_e, gW_f, stabilizer_fix_anisotropic, gW_gW, (c) gW_mem (mem_latOf_split),
% (d) ratGL_UJ, ratGam, rfinU, dict_K, (e) sub_prV_mem_Xs, stab_Xs, stab_perp. (2) iota_act: act, isoG_lift,
% iotaOf_of_mk (IotaRepC), kOf_mul_right, AdmPar_holds, satGramActLIota_kB_good, levi_mul, levi_inv.
% (3) gamL_levi_bij: colD_mul, colD_inv, good_colD, levi_mul, levi_inv, isUni_one,
% mem_level_of_unitriangular, ncong_colDiff, colD_ku and NeatKu (Closed.NeatKu_holds, Lemma
% asm-division-kernel-finite-neat). The remaining inputs are the open steps of Lemmas -pairs-kernel-inputs-
% lattice, -reduction and -terms.
\end{proof}

\begin{lemma}\label{asm-division-kernel-finite-pairs-kernel}
In Situation~\ref{asm-division-kernel-finite-situation}, let $u,u'$ be words, with flags $F$ of $u$ and
$F'$ of $u'$, the latter with vectors $f_1,\dots,f_r$ and $z'$, and put $X'=X(F')$. Let
$s_1,s_2,t_1,t_2$ be terms of $T_{uu'}$, all with rows $\Lambda$ of $u$ and $\Lambda'$ of $u'$, with
isometries $\gamma_{s_i}$ and $\gamma_{t_i}$. Suppose that
$\gamma_{t_i}=n_im_i\gamma_{s_i}c_i$ for $i=1,2$, with $n_i\in N_\Lambda(F)$,
$m_i\in N_\Lambda(\gamma_{s_i}F')$ and isometries $c_i$ with $c_i\Lambda'=\Lambda'$ that fix the $f_a$
and $z'$ (Definition~\ref{asm-division-kernel-finite-pairs-invariants}). If $\iota(s_1)=\iota(s_2)$ and
$\iota(t_1)=\iota(t_2)$, then $(c_2c_1^{-1}-1)X'^\perp\subset X'$.
\end{lemma}
\leanref{Formalism/LocalToGlobal/SiegelLatticeTheory/FlagStabilizerTensors/RootLatticeThresholds/IotaKerHolds.lean}{Formalism.L2.SG.Pf.KernelOfEmbedding.IotaKer_holds}
\leanref{Formalism/LocalToGlobal/SiegelLatticeTheory/FlagStabilizerTensors/RootLatticeThresholds/XLeSpanEXLam.lean}{Formalism.L2.SG.Pf.FinXP.IotaKer}
\leanref{Formalism/LocalToGlobal/SiegelLatticeTheory/FlagStabilizerTensors/RootLatticeThresholds/IotaKerCore.lean}{Formalism.L2.SG.Pf.KernelOfEmbedding.iotaKer_core}
\begin{proof}
Use the notation $G(\Lambda')$, $\bar\delta$ and $\lambda$ of
Lemma~\ref{asm-division-kernel-finite-pairs-kernel-inputs}.

\emph{Step 1.} For $i=1,2$, $c_i$ fixes the vectors of $F'$, so it preserves $X'$ and $X'^\perp$. As
$\gamma_{s_i}$ is an isometry carrying $\Lambda'$ to $\Lambda$ and $F'$ to $\gamma_{s_i}F'$,
$\mu_i=\gamma_{s_i}^{-1}m_i\gamma_{s_i}$ lies in $N_{\Lambda'}(F')$, as in Step 1 of the proof of
Lemma~\ref{asm-division-kernel-finite-fibres}. So $\gamma_{t_i}=n_i\gamma_{s_i}\mu_ic_i$.

\emph{Step 2: there is $\delta_i\in G(\Lambda')$ with $t_i=\delta_i\cdot s_i$ and $c_iy-\bar\delta_i^{-1}y\in X'$
for every $y\in X'^\perp$.} The map induced by $c_i$ on $X'^\perp/X'$ lies in $A(\Lambda')$, so by
Lemma~\ref{asm-division-kernel-finite-pairs-kernel-inputs}(1) there is $g\in G(\Lambda')$ with
$c_iy-\bar gy\in X'$ for $y\in X'^\perp$. Put $\delta_i=g^{-1}$, so that $\bar\delta_i^{-1}=\bar g$. The isometry
$c_i^{-1}\bar g$ preserves $\Lambda'$, fixes $X'$ pointwise, and $c_i^{-1}\bar gx-x=c_i^{-1}(\bar gx-c_ix)\in X'$
for $x\in X'^\perp$, as $c_i^{-1}$ fixes $X'$; so $c_i^{-1}\bar g\in N_{\Lambda'}(F')$. Conjugation by $\bar g$,
which fixes $F'$ and preserves $\Lambda'$, preserves $N_{\Lambda'}(F')$, as in Step 3 of the proof of
Lemma~\ref{asm-division-kernel-finite-fibres}; so $\bar g^{-1}\mu_i\bar g\in N_{\Lambda'}(F')$. Hence
\[
 \gamma_{t_i}=n_i\,\big(\gamma_{s_i}\bar g\big)\,\big(\bar g^{-1}\mu_i\bar g\big)\big(c_i^{-1}\bar g\big)^{-1}
 \in N_\Lambda(F)\,\gamma_{s_i}\bar\delta_i^{-1}\,N_{\Lambda'}(F').
\]
The term $t_i$ has the rows of $s_i$, so it is $\delta_i\cdot s_i$
(Lemma~\ref{asm-division-kernel-finite-pairs-kernel-inputs}(2)).

\emph{Step 3: $\delta_1=\delta_2$.} By Lemma~\ref{asm-division-kernel-finite-pairs-kernel-inputs}(2),
$\iota(t_i)=\iota(s_i)\lambda(\delta_i)^{-1}$. As $\iota(s_1)=\iota(s_2)$ and $\iota(t_1)=\iota(t_2)$,
$\lambda(\delta_1)=\lambda(\delta_2)$, and Lemma~\ref{asm-division-kernel-finite-pairs-kernel-inputs}(3)
gives $\delta_1=\delta_2$. Write $\delta$ for it.

\emph{Step 4.} Let $x\in X'^\perp$, and put $y=c_1^{-1}x$, which lies in $X'^\perp$ by Step 1. By Step 2,
$c_2y-\bar\delta^{-1}y$ and $c_1y-\bar\delta^{-1}y$ lie in $X'$. Their difference is $c_2c_1^{-1}x-x$.
% Lean: IotaKer_holds = iotaKer_core (Closed.NeatKu_holds X u' Y.hK) iotaRepC_holds. iotaKer_core:
% act_of_dec (gamL_autLz_bij, the bijection of GamL with the automorphisms of the row inducing on
% X'^perp / X', and nlam_of_sub) gives gamma_i in GamL (Y.cs u') q with t_i = act Y q gamma_i s_i and
% c_i = isoG gamma_i^-1 modulo X' on X'^perp; iota_act (IotaRepC = iotaRepC_holds, AmpIota.lean:
% rep_split, nilradical_conj_level', good_of_ncong, termK_car, kOf) gives iota(t_i) = iota(s_i) *
% (levi of gamma_i)^-1; gamL_levi_bij (ActMass.lean, from NeatKu, mem_level_of_unitriangular, colD_ku)
% gives gamma_1 = gamma_2; the last lines are the computation e3. Step 1 is fix_of_ez and nlam_frame;
% Step 2 is act_of_dec (stabF_of_fix, indW_mem, nlam_of_sub, nlam_conj, nlamPf_mul, nlamPf_inv,
% tmap_cls, Quot.sound); Step 4 is perp_maps and e3. Size: without NeatKu_holds (Lemma
% asm-division-kernel-finite-neat) about 3,980 lines (proof closure, text count). The inputs of Lemma
% -pairs-kernel-inputs, measured on their own, are about 4,060 lines and contain nearly all of it; this
% proof with act_of_dec is about 1,000 lines with the construction of the action.
\end{proof}

\begin{lemma}\label{asm-division-kernel-finite-pairs}
In Situation~\ref{asm-division-kernel-finite-situation}, for terms with the same rows, the corank
$s_0(t,t')$ and the scalarity of $(t,t')$ depend only on the terms of the unfolding below $t$ and
$t'$.
\end{lemma}
\leanref{Formalism/LocalToGlobal/SiegelLatticeTheory/FlagStabilizerTensors/RootLatticeThresholds/XLeSpanEXLam.lean}{Formalism.L2.SG.Pf.FinXP.pairIndep_of}
\leanref{Formalism/LocalToGlobal/SiegelLatticeTheory/FlagStabilizerTensors/RootLatticeThresholds/IotaKerHolds.lean}{Formalism.L2.SG.Pf.KernelOfEmbedding.IotaKer_holds}
\leanref{Formalism/LocalToGlobal/SiegelLatticeTheory/FlagStabilizerTensors/RootLatticeThresholds/XLeSpanEXLam.lean}{Formalism.L2.SG.Pf.FinXP.pair_step}
\leanref{Formalism/LocalToGlobal/SiegelLatticeTheory/FlagStabilizerTensors/RootLatticeThresholds/XLeSpanEXLam.lean}{Formalism.L2.SG.Pf.FinXP.conj_triv}
\begin{proof}
The corank and the scalarity are those of
Definition~\ref{asm-division-kernel-finite-pairs-invariants}(3). Let $F$, with vectors $e_a$ and $z$,
and $F'$, with vectors $f_a$ and $z'$, be the flags of the words $u$ and $u'$, and $X'=X(F')$. Let
$s_1,s_2,t_1,t_2$ be terms of $T_{uu'}$ such that $s_i$ and $t_i$ lie over the same term of the
unfolding, $\rho(s_1)=\rho(s_2)$ and $\rho(t_1)=\rho(t_2)$. We show that $s_0(t_1,t_2)=s_0(s_1,s_2)$ and
that $(t_1,t_2)$ is scalar if $(s_1,s_2)$ is; exchanging the roles of the $s_i$ and the $t_i$ gives the
converse.

\emph{Step 1: decomposition.} By Step 1 of the proof of
Lemma~\ref{asm-division-kernel-finite-fibres}, $t_i$ has the rows of $s_i$, and
$\gamma_{t_i}=n_im_i\gamma_{s_i}c_i$ with $n_i\in N_\Lambda(F)$, $m_i\in N_\Lambda(\gamma_{s_i}F')$ and
an isometry $c_i$ with $c_i\Lambda'=\Lambda'$ that fixes the $f_a$ and $z'$; the flag $\gamma_{t_i}F'$
is the flag $K_i$ with vectors $n_im_i\gamma_{s_i}f_a$ and $n_im_i\gamma_{s_i}z'$. As the rows of $s_1$
and $s_2$, and of $t_1$ and $t_2$, agree, all four terms have rows $\Lambda$ and $\Lambda'$, and
$\iota(s_1)=\iota(s_2)$, $\iota(t_1)=\iota(t_2)$.

\emph{Step 2: the move.} Put $G_i=\gamma_{s_i}F'$, $\beta=\gamma_{s_2}\gamma_{s_1}^{-1}$, and
$\tau=\gamma_{s_1}c_2c_1^{-1}\gamma_{s_1}^{-1}$. Then
\[
 \gamma_{t_2}\gamma_{t_1}^{-1}=n_2m_2\gamma_{s_2}c_2c_1^{-1}\gamma_{s_1}^{-1}m_1^{-1}n_1^{-1}
 =n_2m_2\beta\tau m_1^{-1}n_1^{-1}.
\]
The $c_i$ fix the $f_a$, so $c_2c_1^{-1}$ fixes $X'$ pointwise, and
$(c_2c_1^{-1}-1)X'^\perp\subset X'$ by Lemma~\ref{asm-division-kernel-finite-pairs-kernel}. As
$\gamma_{s_1}$ is an isometry with $\gamma_{s_1}X'=X(G_1)$, $\tau X(G_1)\subset X(G_1)$ and
$(\tau-1)X(G_1)^\perp\subset X(G_1)$. Also $\beta X(G_1)=\gamma_{s_2}X'=X(G_2)$.

\emph{Step 3: the hypotheses of Lemma~\ref{asm-division-kernel-finite-pairs-move}.} As $\Lambda$ is a row
of $u$, $z\in\Lambda$ and the $e_a$ lie in $X_\Lambda(F)$; as $\Lambda'$ is a row of $u'$, $z'\in\Lambda'$
and the $f_a$ lie in $X_{\Lambda'}(F')$. Applying $\gamma_{s_i}$, which carries $\Lambda'$ to $\Lambda$,
$z_{G_i}=\gamma_{s_i}z'\in\Lambda$ and the vectors of $G_i$ lie in $X_\Lambda(G_i)$. So
Lemma~\ref{asm-division-kernel-finite-pairs-move} applies, and gives
$\Theta(F,K_1,K_2,\gamma_{t_2}\gamma_{t_1}^{-1})=n_1\Theta(F,G_1,G_2,\beta)$. The two spaces have the same
dimension, so $s_0(t_1,t_2)=s_0(s_1,s_2)$; and if $(G_1,G_2,\beta)$ is scalar for $F$, so is
$(K_1,K_2,\gamma_{t_2}\gamma_{t_1}^{-1})$. With the equality of the coranks, $(t_1,t_2)$ is scalar if
$(s_1,s_2)$ is.
% Lean: pairIndep_of hK = pair_step twice (A and B); pair_step: fib_decomp for Step 1 (with dc_of_toTerm),
% the rows hl, hl', hi from rowOf; hker = hK (IotaKer_holds, Lemma -pairs-kernel), hcc with fix_X, and
% conj_triv for tau (Step 2), hbt and hbX; Step 3 is lift_F, lift_frame, X_le_spanE_XLam (IsRowLat,
% basis_vector_mem_span) and hs_frame, then Th_move with LinearEquiv.finrank_map_eq for s0T and
% scal_move for liftScalar. The Lean's PairIndepAll is the statement for s1, s2, s1', s2' with
% s_i.toTerm = s_i'.toTerm and equal rows, giving equal s0T and liftScalar iff; the scalarity of
% Definition -pairs-invariants(3) (scalT) follows.
\end{proof}

\begin{definition}\label{asm-division-kernel-finite-volumes-cocompact-pattern}
Let $I$ be a finite set, $\nu\ge1$ an integer and $\ell\colon I\to\{0,1,\dots,\nu-1\}$ a function, the
\emph{levels}. Let $R$ be a commutative ring with an involution $a\mapsto\bar a$ in which $2$ is
invertible, and for $M\in M_I(R)$ let $M^*$ be its conjugate transpose.
\begin{enumerate}
\item For an integer $k\ge0$, a matrix $Z\in M_I(R)$ is \emph{in pattern $k$} if $Z_{ij}=0$ whenever
  $\ell(j)<\ell(i)+k$.
\item A matrix $J\in M_I(R)$ is \emph{graded of total $T$} if $J_{ij}\ne0$ implies $\ell(i)+\ell(j)=T$.
\item Let $J,J'\in M_I(R)$ be Hermitian matrices, graded of the same total, with $JJ'=J'J=1$. For
  $M\in M_I(R)$ put $M^\sharp=J'M^*J$ and $\pi(M)=\frac12(M-M^\sharp)$. A matrix $Z$ is \emph{skew}
  for $J$ if $Z^*J=-JZ$.
\item The \emph{Cayley transform} of $Z\in M_I(R)$ is
  $\operatorname{cay}(Z)=1+Z\sum_{j=0}^{\nu-1}(Z/2)^j$.
\end{enumerate}
\end{definition}
\leanref{Formalism/LocalToGlobal/SiegelLatticeTheory/HaarTransversals/UnitaryHaarProducts/UnipTool.lean}{Formalism.L2.SG.Pf.Unip.InPat}
\leanref{Formalism/LocalToGlobal/SiegelLatticeTheory/HaarTransversals/UnitaryHaarProducts/UnipTool.lean}{Formalism.L2.SG.Pf.Unip.Graded}
\leanref{Formalism/LocalToGlobal/SiegelLatticeTheory/HaarTransversals/UnitaryHaarProducts/UnipTool.lean}{Formalism.L2.SG.Pf.Unip.sharp}
\leanref{Formalism/LocalToGlobal/SiegelLatticeTheory/HaarTransversals/UnitaryHaarProducts/UnipTool.lean}{Formalism.L2.SG.Pf.Unip.proj}
\leanref{Formalism/LocalToGlobal/SiegelLatticeTheory/HaarTransversals/UnitaryHaarProducts/UnipTool.lean}{Formalism.L2.SG.Pf.Unip.cay}

\begin{lemma}\label{asm-division-kernel-finite-volumes-cocompact-cayley}
Let $I$, $\nu$, $\ell$, $R$, $J$ and $J'$ be as in
Definition~\ref{asm-division-kernel-finite-volumes-cocompact-pattern}(1)--(3), and let $k,l\ge0$ be
integers.
\begin{enumerate}
\item Sums, multiples and images under ring homomorphisms of matrices in pattern $k$ are in pattern
  $k$, and a matrix in pattern $l\ge k$ is in pattern $k$. The identity is in pattern $0$. The product
  of a matrix in pattern $k$ and one in pattern $l$ is in pattern $k+l$. A matrix in pattern $\nu$ is
  $0$; so $Z^\nu=0$ if $Z$ is in pattern $k\ge1$.
\item If $M$ is in pattern $k$, so are $M^\sharp$ and $\pi(M)$. For every $M$, $\pi(M)$ is skew for
  $J$. If $Z$ is skew for $J$, then $Z^\sharp=-Z$ and $\pi(Z)=Z$. Moreover $1^\sharp=1$, and
  $x^\sharp x=1$ if $x^*Jx=J$.
\item If $x\in M_I(R)$ satisfies $x^*Jx=J$ and $x-1$ is in pattern $k$, then $(x-1)-\pi(x-1)$ is in
  pattern $2k$.
\item Let $Z$ be in pattern $k\ge1$. Then $\operatorname{cay}(Z)-1$ is in pattern $k$,
  $\operatorname{cay}(Z)-1-Z$ is in pattern $2k$, and
  $\operatorname{cay}(-Z)\operatorname{cay}(Z)=1$. If $Z$ is skew for $J$, then
  $\operatorname{cay}(Z)^*J\operatorname{cay}(Z)=J$.
\item Let $k\ge1$, let $x$ satisfy $x^*Jx=J$ with $x-1$ in pattern $k$, and let $Y$ be in pattern $k$
  and skew for $J$. Then
  $\operatorname{cay}(Y)^\sharp\,x\,\operatorname{cay}\big(\pi(x-1)-Y\big)^\sharp-1$ is in pattern
  $2k$.
\end{enumerate}
\end{lemma}
\leanref{Formalism/LocalToGlobal/SiegelLatticeTheory/HaarTransversals/UnitaryHaarProducts/UnipTool.lean}{Formalism.L2.SG.Pf.Unip.inPat_mul}
\leanref{Formalism/LocalToGlobal/SiegelLatticeTheory/HaarTransversals/UnitaryHaarProducts/UnipTool.lean}{Formalism.L2.SG.Pf.Unip.inPat_sharp}
\leanref{Formalism/LocalToGlobal/SiegelLatticeTheory/HaarTransversals/UnitaryHaarProducts/UnipTool.lean}{Formalism.L2.SG.Pf.Unip.proj_skew}
\leanref{Formalism/LocalToGlobal/SiegelLatticeTheory/HaarTransversals/UnitaryHaarProducts/UnipTool.lean}{Formalism.L2.SG.Pf.Unip.sub_proj_of_unitary}
\leanref{Formalism/LocalToGlobal/SiegelLatticeTheory/HaarTransversals/UnitaryHaarProducts/UnipTool.lean}{Formalism.L2.SG.Pf.Unip.cay_sub}
\leanref{Formalism/LocalToGlobal/SiegelLatticeTheory/HaarTransversals/UnitaryHaarProducts/UnipTool.lean}{Formalism.L2.SG.Pf.Unip.cay_unitary}
\leanref{Formalism/LocalToGlobal/SiegelLatticeTheory/HaarTransversals/UnitaryHaarProducts/UnipRat.lean}{Formalism.L2.SG.Pf.UnipB.cay_neg_mul}
\leanref{Formalism/LocalToGlobal/SiegelLatticeTheory/HaarTransversals/UnitaryHaarProducts/UnipTool.lean}{Formalism.L2.SG.Pf.Unip.step_mat}
\begin{proof}
(1) If $\ell(j)<\ell(i)+k+l$, then for every $m$ either $\ell(m)<\ell(i)+k$ or $\ell(j)<\ell(m)+l$, so
every term of $\sum_mZ_{im}W_{mj}$ vanishes for $Z$ in pattern $k$ and $W$ in pattern $l$. The
off-diagonal entries of $1$ vanish, and $\ell(i)<\ell(i)$ is false. As $\ell(j)\le\nu-1<\ell(i)+\nu$,
every entry of a matrix in pattern $\nu$ vanishes; and $Z^\nu$ is in pattern $k\nu\ge\nu$. The rest is
clear.

(2) Let $T$ be the total of $J$ and $J'$. In $(M^\sharp)_{ij}=\sum_{p,q}J'_{ip}\overline{M_{qp}}J_{qj}$ a
nonzero term needs $\ell(p)=T-\ell(i)$, $\ell(q)=T-\ell(j)$ and $\ell(p)\ge\ell(q)+k$, that is
$\ell(j)\ge\ell(i)+k$. So $M^\sharp$ and $\pi(M)$ are in pattern $k$. As $J$ and $J'$ are Hermitian,
$(M-M^\sharp)^*J=M^*J-JMJ'J=M^*J-JM=-J(M-M^\sharp)$, so $\pi(M)$ is skew. If $Z^*J=-JZ$, then
$Z^\sharp=J'(-JZ)=-Z$ and $\pi(Z)=Z$. Finally $1^\sharp=J'J=1$, and $x^\sharp x=J'x^*Jx=J'J=1$.

(3) Write $x=1+Z$. Then $x^\sharp=1+Z^\sharp$, and $x^\sharp x=1$ by (2) gives
$Z^\sharp Z+Z^\sharp+Z=0$. Hence $Z-\pi(Z)=\frac12(Z+Z^\sharp)=-\frac12Z^\sharp Z$, which is in pattern
$2k$ by (1) and (2).

(4) Put $Y=Z/2$ and $S=\sum_{j=0}^{\nu-1}Y^j$, which is in pattern $0$. Then
$\operatorname{cay}(Z)-1=ZS$ is in pattern $k$, and as $S=1+Y\sum_{j=0}^{\nu-2}Y^j$,
$\operatorname{cay}(Z)-1-Z=ZY\sum_{j=0}^{\nu-2}Y^j$ is in pattern $2k$. By (1) $Y^\nu=0$, so
$(1-Y)S=1$ and $\operatorname{cay}(Z)=1+2YS=(1+Y)(1-Y)^{-1}$; likewise
$\operatorname{cay}(-Z)=(1-Y)(1+Y)^{-1}$. These are polynomials in $Y$, so they commute, and their
product is $1$. If $Z$ is skew, then $Y^*J=-JY$, so $(1+Y)^*J=J(1-Y)$ and
$\big((1-Y)^{-1}\big)^*J=J(1+Y)^{-1}$, and
$\operatorname{cay}(Z)^*J\operatorname{cay}(Z)=J(1+Y)^{-1}(1-Y)(1+Y)(1-Y)^{-1}=J$.

(5) Put $S=\pi(x-1)$, which is skew and in pattern $k$ by (2), and $a=\operatorname{cay}(Y)^\sharp$,
$c=\operatorname{cay}(S-Y)^\sharp$. By (2) and (4), $a-1=(\operatorname{cay}(Y)-1)^\sharp$ and $c-1$ are
in pattern $k$, as is $x-1$. Expanding $(1+(a-1))(1+(x-1))(1+(c-1))$,
\[
 A_1=axc-1-\big((a-1)+(x-1)+(c-1)\big)
\]
is a sum of products of at least two of $a-1$, $x-1$, $c-1$, so it is in pattern $2k$. By (4) and (2),
$A_2=(a-1)-Y^\sharp=(\operatorname{cay}(Y)-1-Y)^\sharp$ and
$A_3=(c-1)-(S-Y)^\sharp=(c-1)+S+Y^\sharp$ are in pattern $2k$, using $S^\sharp=-S$; and by (3)
$A_4=(x-1)-S$ is in pattern $2k$. Then $axc-1=A_1+A_2+A_3+A_4$ is in pattern $2k$.
% Lean: (1) inPat_add, inPat_smul, inPat_map, inPat_mono, inPat_one, inPat_mul, inPat_pow,
% eq_zero_of_inPat and nil_of_pat; (2) inPat_sharp, inPat_proj, proj_skew, sharp_of_skew,
% proj_of_skew, sharp_one and sharp_mul_self; (3) sub_proj_of_unitary; (4) cay_sub_one, cay_sub (with
% geo_succ), cay_eq and cay_neg_mul (UnipRat.lean), cay_unitary; (5) step_mat, with
% filtered_mul_mul_sub_sub (A_1), sharp_cay_sub (A_2, A_3), sharp_cay_sub_one and
% sub_proj_of_unitary (A_4). The Lean's step_mat takes the skewness of Y as a hypothesis and does not
% use it.
\end{proof}

\begin{lemma}\label{asm-division-kernel-finite-volumes-cocompact-integers}
Let $F$ be a totally real number field and $E$ a totally complex quadratic extension of $F$, and for
each infinite place $v$ of $F$ let $\tau_v\colon E\to\mathbf C$ be an embedding that extends the
embedding of $F$ at $v$. Then there is $C\ge0$ such that for every family $(t_v)_v$ of complex numbers
there is $y\in\mathcal O_E$ with $|t_v-\tau_v(y)|\le C$ for every $v$.
\end{lemma}
\leanref{Formalism/LocalToGlobal/SiegelLatticeTheory/HaarTransversals/UnitaryHaarProducts/UnipLatPf.lean}{Formalism.L2.SG.Pf.oeApprox_holds}
\leanref{Formalism/LocalToGlobal/SiegelLatticeTheory/HaarTransversals/UnitaryHaarProducts/UnipLat.lean}{Formalism.L2.SG.Pf.UnipLat.OEApprox}
\begin{proof}
For each complex place $w$ of $E$ fix an embedding $\sigma_w$ defining it. The Minkowski embedding
$j\colon E\to\prod_w\mathbf C$, $y\mapsto(\sigma_w(y))_w$, over the places of $E$, all complex, maps
$\mathcal O_E$ onto a lattice, which has a bounded fundamental domain $\Phi$: every $x$ is $j(y)+f$ with
$y\in\mathcal O_E$ and $f\in\Phi$; let $C$ bound $|f_w|$ for $f\in\Phi$ and all $w$. Let $w_v$ be the
place of $\tau_v$; it lies over $v$, so distinct $v$ give distinct $w_v$. Given $(t_v)$, let $x$ have
coordinate $t_v$ at $w_v$ if $\sigma_{w_v}=\tau_v$ and $\overline{t_v}$ if
$\sigma_{w_v}=\overline{\tau_v}$, and $0$ at the other places. Write $x=j(y)+f$. At $w_v$,
$|x_{w_v}-\sigma_{w_v}(y)|\le C$, which is $|t_v-\tau_v(y)|\le C$ in the first case and its complex
conjugate in the second.
% Lean: oeApprox_holds, with ZSpan.fundamentalDomain_isBounded (mixedEmbedding.latticeBasis E), the
% places w v = InfinitePlace.mk (phi v) (hw, hinj), the point x of the mixed space, ZSpan.fract and
% InfinitePlace.embedding_mk_eq for the two cases.
\end{proof}

\begin{lemma}\label{asm-division-kernel-finite-volumes-cocompact-approx}
Let $F$, $E$ and $(\tau_v)_v$ be as in Lemma~\ref{asm-division-kernel-finite-volumes-cocompact-integers},
and let $a\mapsto\bar a$ be the nontrivial automorphism of $E$ over $F$, so that
$\tau_v(\bar a)=\overline{\tau_v(a)}$. Let $I$, $\nu$ and $\ell$ be as in
Definition~\ref{asm-division-kernel-finite-volumes-cocompact-pattern}, let $J,J'\in M_I(E)$ be
Hermitian, graded of the same total, with $JJ'=J'J=1$, and write $J_v$ and $J'_v$ for their images under
$\tau_v$. Let $k\ge1$ and $d\ge1$ be integers. Then there is $C\ge0$ with the following property: for
every family $(P_v)_v$ of matrices $P_v\in M_I(\mathbf C)$ in pattern $k$ and skew for $J_v$, there is
$Y\in d\,M_I(\mathcal O_E)$ in pattern $k$ and skew for $J$ with $|\tau_v(Y_{ij})-(P_v)_{ij}|\le C$
for every $v$, $i$ and $j$.
\end{lemma}
\leanref{Formalism/LocalToGlobal/SiegelLatticeTheory/ClosedLocalOrbits/GaugeLatticeShifts/RCocompactPfApproxFinal.lean}{Formalism.L2.SG.TermOp.Vanish.ratApprox_holds}
\leanref{Formalism/LocalToGlobal/SiegelLatticeTheory/ClosedLocalOrbits/GaugeLatticeShifts/RCocompactPf.lean}{Formalism.L2.SG.TermOp.Vanish.RatApprox}
\leanref{Formalism/LocalToGlobal/SiegelLatticeTheory/ClosedLocalOrbits/GaugeLatticeShifts/RCocompactPfApprox.lean}{Formalism.L2.SG.TermOp.Vanish.ratApprox_of}
\leanref{Formalism/LocalToGlobal/SiegelLatticeTheory/ClosedLocalOrbits/GaugeLatticeShifts/RCocompactPfSolApprox.lean}{Formalism.L2.SG.TermOp.Vanish.ratSolApprox_holds}
\leanref{Formalism/LocalToGlobal/SiegelLatticeTheory/ClosedLocalOrbits/GaugeLatticeShifts/RCocompactPfKer.lean}{Formalism.L2.SG.Pf.RCoc.kerBaseChange_holds}
\leanref{Formalism/LocalToGlobal/SiegelLatticeTheory/HaarTransversals/UnitaryHaarProducts/UnipLatStepPf.lean}{Formalism.L2.SG.Pf.UnipB.latStep_lvl}
\begin{proof}
\emph{Step 1: an integral approximation in pattern $k$.} The matrices in pattern $k$ are the
solutions $z$ of a system $Az=0$ of linear equations over $E$ in the entries. Choose a matrix $Q$ over
$E$ with $AQA=A$, and let $\pi_k=1-QA$, an $E$-linear map of $M_I(E)$. Then $A\pi_k=0$, so $\pi_k$ takes
values in pattern $k$; and $\pi_kz=z$ whenever $Az=0$, also after extending scalars along any
$\tau_v$, so $\tau_v(\pi_k)$ fixes every complex matrix in pattern $k$. Choose integers $c_0,c\ge1$
such that $c_0\pi_k$ has coefficients in $\mathcal O_E$ and $cJ,cJ'\in M_I(\mathcal O_E)$, and put
$N_1=2dc^2$ and $N_2=N_1c_0$. Let $C_0$ be the constant of
Lemma~\ref{asm-division-kernel-finite-volumes-cocompact-integers}, and $B$ the sum of the absolute values
of the coefficients of all the $\tau_v(\pi_k)$. For each $(i,j)$ choose $y_{ij}\in\mathcal O_E$ with
$|(P_v)_{ij}/N_2-\tau_v(y_{ij})|\le C_0$ for every $v$, and put $Y_0=N_2\pi_k(y)=N_1(c_0\pi_k)(y)$. Then
$Y_0$ is in pattern $k$ and $Y_0=N_1A_0$ with $A_0\in M_I(\mathcal O_E)$. As $P_v$ is in pattern $k$,
\[
 \tau_v(Y_0)-P_v=\tau_v(\pi_k)\big(N_2\tau_v(y)-P_v\big),
\]
and the entries of $N_2\tau_v(y)-P_v$ are at most $N_2C_0$. So every entry of $\tau_v(Y_0)-P_v$ is at
most $C_1=N_2C_0B$ in absolute value.

\emph{Step 2: the skew part.} Put $Y=\pi(Y_0)=\frac12(Y_0-J'Y_0^*J)$, which is in pattern $k$ and skew
for $J$ by Lemma~\ref{asm-division-kernel-finite-volumes-cocompact-cayley}(2). As $Y_0=2dc^2A_0$,
\[
 Y=d\big(c^2A_0-(cJ')A_0^*(cJ)\big),
\]
and $A_0^*$ is integral, as $\mathcal O_E$ is stable under the automorphism; so $Y\in
d\,M_I(\mathcal O_E)$. Let $\pi_v$ be the map $\pi$ formed with $J_v$ and $J'_v$. As $P_v$ is skew,
$\pi_v(P_v)=P_v$; as $\tau_v$ commutes with the involutions, $\tau_v(Y)=\pi_v(\tau_v(Y_0))$. So
$\tau_v(Y)-P_v=\pi_v(\tau_v(Y_0)-P_v)$. If $b$ and $b'$ bound the entries of all $J_v$ and $J'_v$, the
entries of $\pi_v(X)$ are at most $C_1+|I|^2bb'C_1$ for every $X$ with entries at most $C_1$, and
$C=C_1+|I|^2bb'C_1$ works.
% Lean: latStep_lvl applies ratApprox_holds D X 1 (the case M = 1 of RatApprox, where the extra
% condition InPat 1 (Gr M^-1 * P * Gr M) is implied by InPat k P). ratApprox_holds = ratApprox_of_sol
% ratSolApprox_holds. Step 1 is ratSolApprox_holds: condMat (lvk X) k G G' (cond_iff) is A,
% kerBaseChange_holds gives P0 = 1 - Q A with A P0 = 0 and the base change property (hfixE, hfixV),
% exists_common_den gives c_0, oeApprox_holds gives y for P v / (N c), and the bound is the last calc.
% Step 2 is ratApprox_of with jbDenom_holds (c), N = 2 d c^2, star_integral, map_proj and
% norm_proj_entry_le. The Lean states the lemma for the Gram matrix JB X of the basis bv X and its
% inverse JB' X; the proof uses only the properties listed in the statement.
\end{proof}

\begin{lemma}\label{asm-division-kernel-finite-volumes-cocompact}
In Situation~\ref{asm-division-kernel-finite-situation}, let $V$ be the Hermitian space of the
setting, with the group $G_\infty$ of its archimedean isometries, let $u$ be a word, with base point
$x_u$, and let $e_1,\dots,e_{r_S}$ and $z_u$ be the isotropic vectors and the vector of $u$ given by the
setting. Let $N_\infty\subset G_\infty$ be the group of $s$ such that, at every infinite place,
$(s-1)e_a$ lies in the span of the $e_b$ with $b<a$, and $(s-1)z_u$ and $(s-1)y$, for every $y$
orthogonal to all $e_a$, lie in the span of all $e_b$; let $N_{\mathrm f}$ be the group of finite
adelic points defined by the same conditions. Let $\Delta_u\subset N_\infty$ be the set of archimedean
components $\gamma_\infty$ of the rational points $\gamma$ of $U(V)$ with $\gamma_\infty\in N_\infty$,
$\gamma_{\mathrm f}\in N_{\mathrm f}$ and $x_u^{-1}\gamma_{\mathrm f}x_u\in K_2$. Then there is a
compact set $C\subset G_\infty$ with $N_\infty\subset\Delta_uC$.
\end{lemma}
\leanref{Formalism/LocalToGlobal/SiegelLatticeTheory/HaarTransversals/UnitaryHaarProducts/UnipFinal.lean}{Formalism.L2.SG.Pf.noCocompactC_holds}
\leanref{Formalism/LocalToGlobal/SiegelLatticeTheory/HaarTransversals/UnitaryHaarProducts/UnipStmt.lean}{Formalism.L2.SG.Pf.NOCocompact}
\leanref{Formalism/LocalToGlobal/SiegelLatticeTheory/HaarTransversals/UnitaryHaarProducts/PrincipalCongruenceMemLevel.lean}{Formalism.L2.SG.Pf.UnipB.noCocompact_of}
\leanref{Formalism/LocalToGlobal/SiegelLatticeTheory/HaarTransversals/UnitaryHaarProducts/PrincipalCongruenceMemLevel.lean}{Formalism.L2.SG.Pf.UnipB.step_exists}
\leanref{Formalism/LocalToGlobal/SiegelLatticeTheory/HaarTransversals/UnitaryHaarProducts/UnipReal.lean}{Formalism.L2.SG.Pf.UnipB.pat_of_mem}
\leanref{Formalism/LocalToGlobal/SiegelLatticeTheory/HaarTransversals/UnitaryHaarProducts/UnipReal.lean}{Formalism.L2.SG.Pf.UnipB.inN_of_pat}
\leanref{Formalism/LocalToGlobal/SiegelLatticeTheory/HaarTransversals/UnitaryHaarProducts/PrincipalCongruenceMemLevel.lean}{Formalism.L2.SG.Pf.UnipB.nfin_of_pat}
\leanref{Formalism/LocalToGlobal/SiegelLatticeTheory/HaarTransversals/UnitaryHaarProducts/PrincipalCongruenceMemLevel.lean}{Formalism.L2.SG.Pf.UnipB.gam_sub}
\leanref{Formalism/LocalToGlobal/SiegelLatticeTheory/HaarTransversals/UnitaryHaarProducts/PrincipalCongruenceMemLevel.lean}{Formalism.L2.SG.Pf.UnipB.kg_of}
\leanref{Formalism/LocalToGlobal/SiegelLatticeTheory/HaarTransversals/UnitaryHaarProducts/PrincipalCongruenceMemLevel.lean}{Formalism.L2.SG.Pf.UnipB.gE_intDiv}
\leanref{Formalism/LocalToGlobal/SiegelLatticeTheory/HaarTransversals/UnitaryHaarProducts/UnipLatStepPf.lean}{Formalism.L2.SG.Pf.UnipB.latStep_holds}
\leanref{Formalism/LocalToGlobal/SiegelLatticeTheory/HaarTransversals/UnitaryHaarProducts/UnipTool.lean}{Formalism.L2.SG.Pf.Unip.cover_of_steps}
\begin{proof}
\emph{Notation.} Write $r=r_S$. The space $V$ is the space $Y_0$ of $D$ over the CM field $E$ of its
host, with totally real subfield $F$ and form $h$ (Definition~\ref{th-spaces}(5)). For each infinite
place $v$ of $F$ let $\tau_v\colon E\to\mathbf C$ be the embedding of the setting above $v$, so that
the component $s_v$ of $s\in G_\infty$ is an isometry of $V_v=V\otimes_{E,\tau_v}\mathbf C$; as $E$ is
a CM field, $\tau_v(\bar a)=\overline{\tau_v(a)}$. The setting gives vectors $f_1,\dots,f_r$ with
$h(e_a,e_b)=h(f_a,f_b)=0$ and $h(e_a,f_b)=\delta_{ab}$, and $z_u$ is orthogonal to the $e_a$ and $f_a$.
Let $W$ be the orthogonal complement of the $e_a$ and $f_a$, with a basis $w_1,\dots,w_m$. Every
$x\in V$ is $\sum_ah(f_a,x)e_a+\sum_ah(e_a,x)f_a+x_W$ with $x_W\in W$, so
$e_1,\dots,e_r,w_1,\dots,w_m,f_1,\dots,f_r$ is a basis $(b_i)_{i\in I}$ of $V$, in which the coordinate
of $x$ along $f_a$ is $h(e_a,x)$. Give it the levels $\ell(e_a)=a$, $\ell(w_i)=r+1$ and
$\ell(f_a)=2r+2-a$, which lie in $\{1,\dots,2r+1\}$, and take $\nu=2r+2$ in
Definition~\ref{asm-division-kernel-finite-volumes-cocompact-pattern}. The Gram matrix
$J=(h(b_i,b_j))_{ij}$ is Hermitian and graded of total $2r+2$, as $h(b_i,b_j)\ne0$ only if
$\{b_i,b_j\}=\{e_a,f_a\}$ for some $a$ or $b_i,b_j\in W$. Its inverse $J'$ is graded of total $2r+2$
too: the matrix $J''$ obtained from $J'$ by replacing its entries $J'_{jk}$ with
$\ell(j)+\ell(k)\ne2r+2$ by $0$ satisfies $JJ''=1$, since $J_{ij}\ne0$ forces $\ell(j)=2r+2-\ell(i)$; so
$J''=J'$. For $s\in G_\infty$ let $M_v(s)\in M_I(\mathbf C)$ be the matrix of $s_v$ in the basis
$\tau_v(b)$ of $V_v$, and let $J_v,J'_v$ be the images of $J,J'$ under $\tau_v$, with the involution
$\sharp$ and the map $\pi$ that they define. Then $M_v(st)=M_v(s)M_v(t)$ and
$M_v(s)^*J_vM_v(s)=J_v$, so $M_v(s^{-1})=M_v(s)^\sharp$ by
Lemma~\ref{asm-division-kernel-finite-volumes-cocompact-cayley}(2). For $k\ge1$ let $\mathcal P_k$ be
the set of $x\in G_\infty$ such that $M_v(x)-1$ is in pattern $k$ for every $v$.

\emph{Step 1: $N_\infty\subset\mathcal P_1$.} Let $s\in N_\infty$, let $v$ be an infinite place and
$M=M_v(s)$. The column of $M-1$ at $e_a$ is the coordinate vector of $(s_v-1)e_a$, which lies in the
span of the $e_b$ with $b<a$; so $(M-1)_{ie_a}=0$ unless $\ell(i)<a$. The $w_j$ are orthogonal to the
$e_a$, so $(s_v-1)w_j$ lies in the span of the $e_b$, whose levels are at most $r$; so
$(M-1)_{iw_j}=0$ unless $\ell(i)<r+1$. Thus $(M-1)_{ij}=0$ whenever $\ell(j)\le r+1$ and
$\ell(j)<\ell(i)+1$, and the same holds for $M_v(s^{-1})$, as $N_\infty$ is a group. Now let
$\ell(j)\ge r+2$ and $\ell(j)<\ell(i)+1$. As $M=M_v(s^{-1})^\sharp$ and $1^\sharp=1$,
\[
 (M-1)_{ij}=\sum_{p,q}(J'_v)_{ip}\,\overline{\big(M_v(s^{-1})-1\big)_{qp}}\,(J_v)_{qj},
\]
and a nonzero term needs $\ell(p)=2r+2-\ell(i)\le r$ and $\ell(q)=2r+2-\ell(j)\ge\ell(p)$; then
$(M_v(s^{-1})-1)_{qp}=0$. So $M-1$ is in pattern $1$.

\emph{Step 2: matrices in pattern $1$.} Let $\gamma$ be an isometry of $V_v$, or of
$V\otimes_F\mathbf A_{F,\mathrm f}$, whose matrix $M$ in the basis $b$ has $M-1$ in pattern $1$. Then
$\gamma$ satisfies the conditions that define $N_\infty$ at $v$, or $N_{\mathrm f}$. Indeed
$(M-1)_{ie_a}=0$ for $\ell(i)\ge a$, so $(\gamma-1)e_a$ lies in the span of the $e_b$ with $b<a$. If
$y$ is orthogonal to every $e_a$, its coordinates $h(e_a,y)$ along the $f_a$ vanish, so $y$ is a
combination of the $b_j$ with $\ell(j)\le r+1$, and $(\gamma-1)y$ is a combination of the $b_i$ with
$\ell(i)\le r$, which are the $e_b$. This applies to $y=z_u$.

\emph{Step 3: a semigroup in $\Delta_u$.} Let $\hat L$ be the completion of the lattice $L$ of the
setting, which $x_u$ stabilizes, and $g=g_0g_1$ the product of the level primes. Let $\Gamma$ be the set
of the $\gamma_\infty$ for the rational points $\gamma$ of $U(V)$ with $\gamma_\infty\in N_\infty$,
$\gamma_{\mathrm f}\in N_{\mathrm f}$, $\gamma_{\mathrm f}\hat L=\hat L$ and $(\gamma_{\mathrm
f}-1)\hat L\subset g\hat L$. It contains $1$ and is closed under products, as
$(\alpha\beta-1)y=\alpha(\beta-1)y+(\alpha-1)y$. And $\Gamma\subset\Delta_u$: for such $\gamma$,
$\kappa=x_u^{-1}\gamma_{\mathrm f}x_u$ stabilizes $\hat L$ and
$(\kappa-1)\hat L=x_u^{-1}(\gamma_{\mathrm f}-1)x_u\hat L\subset g\hat L$; so $\kappa$ satisfies the
congruences that define the level subgroup $K_2$
(Definition~\ref{asm-division-kernel-finite-neat-level}), whose right sides contain $g\hat L$, and
$\kappa\in K_2$.

\emph{Step 4: integral Cayley transforms.} Let $B$ be the matrix whose columns are the coordinates of
the $b_i$ in the coordinates of $V$ of the setting. Choose integers $c_L,d_L,d_B,d'_B\ge1$ such that $c_L$
times every coordinate of every vector of $\hat L$ is integral, $d_L$ times every standard basis
vector $\varepsilon_i$ lies in $L$, and $d_BB$ and $d'_BB^{-1}$ have integral entries, and put
$t=2^{2r+2}d_Bd'_Bgc_Ld_L$. Let $k\ge1$, and let $Y\in t\,M_I(\mathcal O_E)$ be in pattern $k$ and skew
for $J$. Then $B\operatorname{cay}(Y)B^{-1}$ is the matrix of a rational point $\gamma$ of $U(V)$ with
$\gamma_\infty\in\Gamma$ and $M_v(\gamma_\infty)=\operatorname{cay}(\tau_v(Y))$ for every $v$. Indeed,
by Lemma~\ref{asm-division-kernel-finite-volumes-cocompact-cayley}(4), $\operatorname{cay}(Y)$
preserves $J$, with inverse $\operatorname{cay}(-Y)$, so $\gamma$ is an isometry, with matrix
$\operatorname{cay}(Y)$ in the basis $b$; and $\operatorname{cay}(Y)-1$ is in pattern $k$, hence in
pattern $1$, so $\gamma_\infty\in N_\infty$ and $\gamma_{\mathrm f}\in N_{\mathrm f}$ by Step 2. As
$\tau_v$ is a ring homomorphism, $M_v(\gamma_\infty)=\operatorname{cay}(\tau_v(Y))$. As $2^{2r+2}$
divides $t$, the powers of $Y/2$ are integral, so $\operatorname{cay}(Y)-1=Y\sum_j(Y/2)^j\in
t\,M_I(\mathcal O_E)$, and
$B\operatorname{cay}(Y)B^{-1}-1=(d_BB)\big(\operatorname{cay}(Y)-1\big)(d'_BB^{-1})/(d_Bd'_B)=gc_Ld_LA$
with $A$ integral. For $y\in\hat L$ with coordinates $y_j$,
\[
 (\gamma_{\mathrm f}-1)y=g\sum_j(c_Ly_j)\Big(d_L\sum_iA_{ij}\varepsilon_i\Big),
\]
with $c_Ly_j$ integral and $d_L\sum_iA_{ij}\varepsilon_i\in L$; so $(\gamma_{\mathrm f}-1)\hat L\subset
g\hat L$, and $\gamma_{\mathrm f}\hat L\subset\hat L$. The same holds for $\gamma^{-1}$, with $-Y$ in place
of $Y$, so $\gamma_{\mathrm f}\hat L=\hat L$.

\emph{Step 5: one step down.} Let $\rho$ be the largest of the constants of
Lemma~\ref{asm-division-kernel-finite-volumes-cocompact-approx} for $J$, $J'$, $d=t$ and
$k=1,\dots,2r+1$, and let $S_\rho$ be the set of $s\in G_\infty$ such that for every $v$,
$M_v(s)=\operatorname{cay}(Z_v)$ for some $Z_v\in M_I(\mathbf C)$ whose entries have absolute value at
most $\rho$. It is compact: it is the intersection of the closed subgroup $G_\infty$ with the image of a
product of closed discs under a continuous map. Let $1\le k\le2r+1$ and $x\in\mathcal P_k$. We show that
$x=\gamma ys$ with $\gamma\in\Gamma$, $y\in\mathcal P_{k+1}$ and $s\in S_\rho$. Let
$P_v=\pi(M_v(x)-1)$; it is in pattern $k$ and skew for $J_v$ by
Lemma~\ref{asm-division-kernel-finite-volumes-cocompact-cayley}(2). By
Lemma~\ref{asm-division-kernel-finite-volumes-cocompact-approx} there is $Y\in t\,M_I(\mathcal O_E)$ in
pattern $k$ and skew for $J$ with $|\tau_v(Y_{ij})-(P_v)_{ij}|\le\rho$, and Step 4 gives
$\gamma\in\Gamma$ with $M_v(\gamma)=\operatorname{cay}(\tau_v(Y))$. The matrix $Z_v=P_v-\tau_v(Y)$ is in
pattern $k$ and skew for $J_v$, with entries at most $\rho$, so $\operatorname{cay}(Z_v)$ preserves
$J_v$ by Lemma~\ref{asm-division-kernel-finite-volumes-cocompact-cayley}(4), and there is $s\in S_\rho$
with $M_v(s)=\operatorname{cay}(Z_v)$ for every $v$. Put $y=\gamma^{-1}xs^{-1}$. Then
\[
 M_v(y)=\operatorname{cay}(\tau_v(Y))^\sharp\,M_v(x)\,\operatorname{cay}\big(\pi(M_v(x)-1)-\tau_v(Y)\big)^\sharp,
\]
and Lemma~\ref{asm-division-kernel-finite-volumes-cocompact-cayley}(5) puts $M_v(y)-1$ in pattern $2k$,
hence in pattern $k+1$.

\emph{Step 6: induction.} A matrix in pattern $2r+2$ is $0$, so $\mathcal P_{2r+2}=\{1\}$. Descending
from $k=2r+2$ to $k=1$, there are compact sets $C_k$ with $\mathcal P_k\subset\Gamma C_k$: take
$C_{2r+2}=\{1\}$ and $C_k=C_{k+1}S_\rho$, which is compact; for $x\in\mathcal P_k$, Step 5 gives
$x=\gamma ys$, and $y=\gamma'c$ with $\gamma'\in\Gamma$ and $c\in C_{k+1}$, so $x=(\gamma\gamma')(cs)$
with $\gamma\gamma'\in\Gamma$ by Step 3. By Steps 1 and 3, $N_\infty\subset\mathcal P_1\subset\Gamma
C_1\subset\Delta_uC_1$.
% Lean: noCocompactC_holds = UnipB.noCocompactC_of UnipB.latStepC_holds, and noCocompact_of X u hL.
% Notation: bv X (basis), lvk X (levels), JB and JB' (JB_graded, JB'_graded, JB_herm, JB'_herm), Mv X v
% (Mv_mul, Mv_unitary, Mv_inv), phi_star; the sets P_k are R (k - 1) in noCocompact_of. Step 1 is
% pat_of_mem (half_of_inN, half_to_full with Mv_inv). Step 2 is inN_of_pat and nfin_of_pat (lvk_notinl,
% hJ_e_Bv and hA_e_BA for the coordinates along f_a). Step 3 is Gam X u with gam_mul, gam_one, kg_mul
% and gam_sub (kg_conj with X.xu_K u, principal_congruence_mem_level). Step 4 is gA with gA_GF, gA_ninf,
% gA_nfin, Mv_gA and kg_of (gE_intDiv, kg_rat, exists_cL, exists_dL, exists_den), with t the number
% 2^N0 (dB dB' (g cL dL)) of step_exists. Step 5 is step_exists with latStep_holds (rho is the sum of
% the |C j| over j < N0 there; for k >= N0 both P and Y are 0), Sset and sset_compact (uJEmb), cay_unitary
% and step_mat. Step 6 is cover_of_steps with m = 2 rS + 1 and eq_zero_of_inPat.
% The cocompactness of a congruence subgroup of the rational points in the archimedean points of a
% unipotent group is classical; the book follows the Lean's descent.
\end{proof}

\begin{lemma}\label{asm-division-kernel-finite-volumes}
In Situation~\ref{asm-division-kernel-finite-situation}, the Haar measures of the setting have
positive scaling constants, and the volumes $\mathrm{vol}_N(u)$ and $\mathrm{vol}_H(u)$ are positive
for every word $u$.
\end{lemma}
\leanref{Formalism/PeriodIntegrals/OrbitalBoundsConclusion/CdHaarScale.lean}{cd_haar_scaling}
\leanref{Formalism/LocalToGlobal/SiegelLatticeTheory/ClosedVolumeStatements/UnipotentInLevelSubgroup.lean}{Formalism.L2.SG.Closed.VolPosC_holds}
\leanref{Formalism/LocalToGlobal/SiegelLatticeTheory/FlagStabilizerTensors/RootLatticeThresholds/FullExists.lean}{Formalism.L2.SG.Pf.volPos_of}
\leanref{Formalism/LocalToGlobal/SiegelLatticeTheory/FlagStabilizerTensors/RootLatticeThresholds/FullExists.lean}{Formalism.L2.SG.Pf.muN_pos_of}
\leanref{Formalism/LocalToGlobal/SiegelLatticeTheory/FlagStabilizerTensors/RootLatticeThresholds/FullExists.lean}{Formalism.L2.SG.Pf.muN_fin_of}
\leanref{Formalism/LocalToGlobal/SiegelLatticeTheory/HaarTransversals/UnitaryHaarProducts/ConcreteLevelAdelicSubset.lean}{Formalism.L2.SG.Pf.gf_key}
\leanref{Formalism/LocalToGlobal/SiegelLatticeTheory/HaarTransversals/UnitaryHaarProducts/SCompact.lean}{Formalism.L2.SG.Pf.hinf_compact}
\begin{proof}
\emph{Scaling constants.} Let $\epsilon\in\{\pm1\}^n$, with the Haar measure $\nu_\epsilon$ of
$U(J_\epsilon)$, $J_\epsilon$ the diagonal matrix with entries $\epsilon$, and let $\mu_{P,Q}$ be the
Haar measures of the definite-side datum on the groups $U(P,Q)$. A bijection $e$ from
$\{1,\dots,P\}\sqcup\{1,\dots,Q\}$ to $\{1,\dots,n\}$ that sends the first part to the positions with
$\epsilon=+1$ and the second to those with $\epsilon=-1$ gives, by reindexing matrices, an isomorphism
of topological groups $\phi_e\colon U(P,Q)\to U(J_\epsilon)$. The claim is that there is $c>0$ with
$c\,\nu_\epsilon=\phi_{e*}\mu_{P,Q}$ for every such $e$. If there is no such $e$, take $c=1$. Otherwise
fix one, $e_0$, with $(P_0,Q_0)$. Then $\phi_{e_0*}\mu_{P_0,Q_0}$ is a Haar measure of $U(J_\epsilon)$,
so it is $c\,\nu_\epsilon$ for some $c>0$, by the uniqueness of Haar measure. Another such $e$ has
$(P,Q)=(P_0,Q_0)$, as $P$ and $Q$ count the signs, and $\phi_e=\rho\circ\phi_{e_0}$, where $\rho$ is
the reindexing by the permutation $e\circ e_0^{-1}$ of $\{1,\dots,n\}$, which preserves $\epsilon$. So
$\rho$ is an automorphism of $U(J_\epsilon)$ with $\rho^N=1$ for some $N\ge1$. The image of a Haar
measure $\mu$ under $\rho$ is $c'\mu$ with $c'>0$; applying $\rho$ $N$ times gives $c'^N=1$, so
$c'=1$. Hence $\phi_{e*}\mu_{P,Q}=\rho_*(c\,\nu_\epsilon)=c\,\nu_\epsilon$.

\emph{Notation for the volumes.} Fix a word $u$, and let $N_\infty$, $N_{\mathrm f}$ and $\Delta_u$ be
as in Lemma~\ref{asm-division-kernel-finite-volumes-cocompact}, and $H_\infty\subset G_\infty$ the group
of $s$ that fix every $e_a$, $f_a$ and $z_u$. The setting gives measures $\mu_N$ and $\mu_H$ on
$G_\infty$, Haar measures on these closed subgroups: each is $\sigma$-finite and nonzero, vanishes off
the subgroup, and is invariant under left translation by its elements. It also gives a fundamental
domain $F_N$ for $\Delta_u$: a measurable set such that for every $x\in G_\infty$ there is
$s\in\Delta_u$ with $sx\in F_N$, and such that $x,sx\in F_N$ with $s\in\Delta_u$ imply $s=1$. The
volumes are $\mathrm{vol}_N(u)=\mathrm{vol}(K_2)\,\mu_N(F_N)$ and $\mathrm{vol}_H(u)=\mu_H(G_\infty)$. A
Haar measure on a closed subgroup is finite on compact sets.

\emph{Step 1: $\Delta_u$ is a group.} Let $s=\gamma_\infty$ and $t=\delta_\infty$ lie in $\Delta_u$.
Then $\delta\gamma^{-1}$ is rational. The conditions that define $N_\infty$ and $N_{\mathrm f}$ are
preserved by products and inverses, so $(\delta\gamma^{-1})_\infty=ts^{-1}\in N_\infty$ and
$(\delta\gamma^{-1})_{\mathrm f}\in N_{\mathrm f}$; and
$x_u^{-1}(\delta\gamma^{-1})_{\mathrm f}x_u=(x_u^{-1}\delta_{\mathrm f}x_u)(x_u^{-1}\gamma_{\mathrm
f}x_u)^{-1}\in K_2$, as $K_2$ is a group. So $ts^{-1}\in\Delta_u$.

\emph{Step 2: $\Delta_u$ is countable.} The finite components of the points of $\Delta_u$ lie in
$x_uK_2x_u^{-1}$, which stabilizes $\hat L$ since $x_u$ and $K_2$ do. Let $l_1,\dots,l_s$ span the
lattice $L$ of the setting over $\mathcal O_E$, and let $c\ne0$ be an integer such that $c$ times
every coordinate of a vector of $\hat L$ is integral. Let $W$ be the set of $y\in G_\infty$ such that,
at every infinite place $v$, every coordinate of $c\,(y_v-1)l_j$ has absolute value less than $1$; it
is an open neighbourhood of $1$. If $\gamma$ is rational, $\gamma_\infty\in W$ and $\gamma_{\mathrm f}$
stabilizes $\hat L$, then each coordinate of $c(\gamma l_j-l_j)$ is an algebraic integer of $E$ whose
conjugates all have absolute value less than $1$, so it is $0$; thus $\gamma$ fixes the $l_j$, which
span $V$, and $\gamma=1$. Hence $\Delta_u\cap W=\{1\}$, and $\Delta_u$ is a discrete subgroup of the
second countable group $G_\infty$, so it is countable.

\emph{Step 3: $H_\infty$ is compact.} At each infinite place $v$ the setting gives vectors
$w_1,\dots,w_k$ orthogonal to the $e_a$, $f_a$ and $z_u$ with $h(w_b,w_c)=\delta_{bc}$; with the
$e_a$, $f_a$ and $z_u$ they form a basis, as $2r_S+1+k=n$, and every vector $x$ is
$\sum_a\big(h(f_a,x)e_a+h(e_a,x)f_a\big)+h(z_u,z_u)^{-1}h(z_u,x)z_u+\sum_bh(w_b,x)w_b$. Let $g$ be an
isometry fixing the $e_a$, $f_a$ and $z_u$. Then $gw_b$ is orthogonal to them, so
$gw_b=\sum_ch(w_c,gw_b)w_c$ with $\sum_c|h(w_c,gw_b)|^2=h(gw_b,gw_b)=1$. Expanding the standard
basis vectors in the basis therefore bounds the entries of $g$ independently of $g$. So the component
of $H_\infty$ at $v$ is closed and bounded, hence compact, and $H_\infty$, their product, is compact.

\emph{Step 4: $\mathrm{vol}_H(u)>0$.} As $\mu_H$ vanishes off $H_\infty$, which is compact,
$\mu_H(G_\infty)=\mu_H(H_\infty)<\infty$; and $\mu_H(G_\infty)>0$, as $\mu_H\ne0$.

\emph{Step 5: $\mu_N(F_N)>0$.} The translates $s^{-1}F_N$, $s\in\Delta_u$, cover $G_\infty$, and
$\mu_N(s^{-1}F_N)=\mu_N(F_N)$ as $s\in N_\infty$. If $\mu_N(F_N)=0$, then $N_\infty$, covered by these
countably many sets (Step 2), has measure $0$, and with $\mu_N(G_\infty\setminus N_\infty)=0$ this
gives $\mu_N=0$, which is false.

\emph{Step 6: $\mu_N(F_N)<\infty$.} Let $C$ be the compact set of
Lemma~\ref{asm-division-kernel-finite-volumes-cocompact}; $\mu_N(C)<\infty$, as $N_\infty$ is closed.
For $s\in\Delta_u$ put $P_s=\{c\in C:sc\in F_N\}$ and $Q_s=\{x\in F_N:s^{-1}x\in C\}=sP_s$, so that
$\mu_N(Q_s)=\mu_N(P_s)$. Every $x\in F_N\cap N_\infty$ is $sc$ with $s\in\Delta_u$ and $c\in C$, so lies
in $Q_s$. The $P_s$ are disjoint: if $c\in P_s\cap P_t$, then $sc$ and $tc=(ts^{-1})(sc)$ lie in $F_N$
with $ts^{-1}\in\Delta_u$ by Step 1, so $ts^{-1}=1$. As $\mu_N$ vanishes off $N_\infty$ and $\Delta_u$
is countable,
\[
 \mu_N(F_N)=\mu_N(F_N\cap N_\infty)\le\sum_{s\in\Delta_u}\mu_N(Q_s)=\sum_{s\in\Delta_u}\mu_N(P_s)
 =\mu_N\Big(\bigcup_sP_s\Big)\le\mu_N(C)<\infty.
\]

\emph{Step 7.} The level subgroup is normalized (Lemma~\ref{asm-division-kernel-finite-folding}), so
$0<\mu_{\mathrm f}(K_2)<\infty$ and $\mathrm{vol}(K_2)>0$. With Steps 5 and 6,
$\mathrm{vol}_N(u)=\mathrm{vol}(K_2)\,\mu_N(F_N)>0$, and $\mathrm{vol}_H(u)>0$ by Step 4.
% Lean: the scaling constants are cd_haar_scaling (CdHaarScale.lean) for D.cd : CDHaar, with phi
% (transp, memUpq_iff), card_true and card_false, haarScalarFactor_pos_of_isHaarMeasure, and
% map_phi_indep through rho, rho_finite_order and map_eq_self_of_iterate (MapEqSelfOfIterate.lean).
% The volumes are VolPosC_holds = volPosC_of'' noCofiniteC_holds HFiniteC_holds = volPos_of with
% Full.volN = volK2 * muN (FN u) and Full.volH = muH univ: Step 1 is noSub (nfinGroup_holds, inN_mul,
% inN_inv, LevelSubgroupCharacter), Step 2 is noCount (NOsub_discrete with gf_key: nt_Lf_denom,
% nt_int_eq_zero, psi_eq, mat_eq_one_of_fix; xu_K; countable_of_discrete), Step 3 is hinf_compact
% (hv_compact, frame_compact, frame_bounded, frame_expand_lemma, card_frame, wfr_on, wfr_perp), Step 4
% is hFinite_holds and muH_haar, Step 5 is muN_pos_of, Step 6 is muN_fin_of with noCofinite_holds
% (noCocompactC_holds, ninf_closed and haarOnSet_isHaar_inst), and Step 7 uses LevelSubgroupNormalized
% from level_subgroup_properties.
% Size: import-restricted proof closure of cd_haar_scaling and VolPosC_holds, about 5,200 lines in 61
% files; about 3,000 of them are Lemma asm-division-kernel-finite-volumes-cocompact (text count).
\end{proof}

\begin{lemma}\label{asm-division-kernel-finite}
Assume the decay hypothesis for $W_u$ (Definition~\ref{is-wu-decay}) and the finite dimensionality of
rows. Let $D$ be a definite-side datum that satisfies the gap hypothesis, and let an automorphic
setting attached to $D$ be given with the spectral data of a spectral condition for $D$: the rows of
every word, column representatives with the column spaces $\mathcal C^\sharp_u$, fundamental domains,
a level subgroup with a character, and reductions, as in Definition~\ref{is-unfolding-evaluation}(1).
Then the spectral data are \emph{finite}:
\begin{enumerate}
\item for every word $u$, the map $\sharp$ from the space of Poincaré data of $u$ to
  $\mathcal C^\sharp_u$, evaluation at the column representatives, is bijective;
\item for every word $u$ and every Poincaré datum $T$ of $u$, the Poincaré series of $T$ converges
  absolutely, does not depend on the choice of representatives, is square integrable on the
  fundamental domain and lies in the space $\mathcal H$ of the setting; and it depends linearly on $T$;
\item $\mathcal H$ is finite-dimensional;
\item over each term of the unfolding (Definition~\ref{is-gram-terms}) lie finitely many terms of the
  unfolded term data of Definition~\ref{is-unfolding-evaluation};
\item for terms with the same rows, the corank $s_0(t,t')$ and the scalarity of $(t,t')$ depend only
  on the terms of the unfolding below $t$ and $t'$;
\item the invariant $G$ of every term lies in the totally real subfield $F$ of the host;
\item the Haar measures of the setting have positive scaling constants, and the volumes
  $\mathrm{vol}_N(u)$ and $\mathrm{vol}_H(u)$ are positive for every word $u$.
\end{enumerate}
\end{lemma}
\leanref{Formalism/LocalToGlobal/SiegelLatticeTheory/FlagStabilizerTensors/RootLatticeThresholds/FinXHolds.lean}{Formalism.L2.SG.Pf.FinXP.finX_holds}
\leanref{Formalism/LocalToGlobal/SiegelLatticeTheory/FlagStabilizerTensors/RootLatticeThresholds/FinXClosed.lean}{Formalism.L2.SG.Pf.finXC_of}
\leanref{Formalism/LocalToGlobal/SiegelLatticeTheory/FlagStabilizerTensors/TheoremSRecursion.lean}{Formalism.L2.SG.PoincareAmplitudeLowerBounds.FinXC}
\leanref{Formalism/LocalToGlobal/SiegelLatticeTheory/FlagStabilizerTensors/RootLatticeThresholds/NeatTies.lean}{Formalism.L2.SG.Pf.finX_of_neat}
\leanref{Formalism/LocalToGlobal/SiegelLatticeTheory/FlagStabilizerTensors/RootLatticeThresholds/FinXHolds.lean}{Formalism.L2.SG.Pf.FinXP.greal_holds}
\leanref{Formalism/LocalToGlobal/SiegelLatticeTheory/FlagStabilizerTensors/RootLatticeThresholds/FinXHolds.lean}{Formalism.L2.SG.Pf.FinXP.phLoc_int_of_gap}
\begin{proof}
The data of the statement are those of
Situation~\ref{asm-division-kernel-finite-situation}. By the decay hypothesis, the gap hypothesis
gives the decay of $W_u$ for every word $u$ of the setting and the integrability of the archimedean
matrix coefficient. By Lemma~\ref{asm-division-kernel-finite-folding}, the level subgroup of the
setting is normalized and the Poincaré series of every word fold. The seven properties then hold as
follows.
\begin{enumerate}
\item By Lemma~\ref{asm-division-kernel-finite-neat}, the level subgroup $K_u$ of every word is
  neat, and so $\sharp$ is bijective.
\item This is Lemma~\ref{asm-division-kernel-finite-series}, with the integrability and the folding.
\item This is Lemma~\ref{asm-division-kernel-finite-dimension}.
\item This is Lemma~\ref{asm-division-kernel-finite-fibres}.
\item This is Lemma~\ref{asm-division-kernel-finite-pairs}.
\item Let $t$ be a term over the flags $\mathcal F$ and $\mathcal F'$, with natural numbers $T$ and
  $T'$. Its invariant is $G=\Delta_{r+1}\Delta_{r+1}^*-TT'\Delta_r\Delta_r^*$, where $\Delta_r$ and
  $\Delta_{r+1}$ are leading minors attached to the pair of flags. So $G^*=G$, and $G$ lies in $F$,
  the subfield on which the involution is trivial.
\item This is Lemma~\ref{asm-division-kernel-finite-volumes}.
\end{enumerate}
% The open parts are in Lemma asm-division-kernel-finite-folding; in Lemmas -series-function,
% -series-fold and -series-hecke (the inputs of -series); in Lemma -pairs-kernel (the input of
% -pairs); and in -dimension-decay-spherical and -dimension-decay-embedding (the inputs of the decay
% bound behind -dimension). Lemmas -neat, -volumes-cocompact and -dimension-decay-ktype are written
% out.
% Lean: finX_holds hWu hRF hFL hK, with hFL = FinLeviPf.finLeviC_holds and hK =
% KernelOfEmbedding.IotaKer_holds (finXC_of). (1) neatKu_full and sharp_bij_of_neat; (2)
% Closed.InHHC_holds with hph from hWu and FoldOKC_holds hFL (from level_subgroup_properties), and
% PoincareInHH.Cut.puLinear_holds; (3) RowDictPf.finiteDimensional_HHX_of_hypotheses with hRF,
% schur_orthogonality_discrete_series (PIW.HCFdeg.formal_degree_abstract_form) and phLoc_int_of_gap
% (WuDecayPf.coef_integrable_of_weighted with WuDec.coefWeighted_of_gap and CoefWeightedLoc_holds);
% (4) fibFin_holds (phi_inj into autLz, automorphism_lattice_finite); (5) pairIndep_of IotaKer_holds;
% (6) greal_holds (star_Gflag, conj_eq_self_iff); (7) cd_haar_scaling and Closed.VolPosC_holds.
% Size: proof closure of finX_holds with finLeviC_holds and IotaKer_holds, 2,210 theorems, about
% 29,400 lines in 240 files. The closures of the parts overlap: folding with FinLevi about 13,900
% lines, InHHC_holds about 8,800, (5) with IotaKer about 12,200, VolPosC_holds about 3,900, (1) about
% 2,800, (3) about 1,100, the level subgroup about 1,000 and (4) about 500.
\end{proof}

\begin{theorem}\label{asm-tied-from-green}
Assume Green's parametrization of the cuspidal characters of
$GL_d(\mathbf F_g)$ in the labelled form of
Corollary~\ref{lit-green-labelled}. If a collection of definite-side hypotheses
exists (Definition~\ref{asm-definite-side-hypotheses}), then tied families
exist for all host requirements (Definition~\ref{asm-tied-everywhere}), with
$e_D=2$ and $e_L=0$.
\end{theorem}
\begin{proof}
Let $\mathcal A$ be the class of the collection, with
$\mathcal A_N$ and $\mathcal D_3=\operatorname{Des}(\mathcal A_N,3)$ as
there. Let $\lambda>1$ and let $G$ be a host requirement such that hosts
satisfying $G$ exist on windows of exponent $\lambda$ and every family in
$\mathcal C_G(\lambda)$ satisfies the Tate-Shafarevich bound. Apply
Theorem~\ref{asm-tied-from-links} with $\mathcal C=\mathcal A_N$, with
$\mathcal C'$ the class $(\mathcal A\cap\mathcal C_G(\lambda))_N$, with
$\mathcal E=\mathcal C_G(\lambda)$, and with $C_{\mathrm{abs}}=2$, $C_0=3$,
$e_D=2$, $e_L=0$. Then $\mathcal C'\subset\mathcal C$, $\mathcal C'\subset
\mathcal E$, and every family in $\mathcal E$ satisfies the Tate-Shafarevich
bound. Its hypotheses hold as follows.
\begin{itemize}
\item Designed families exist in $\mathcal C'_{\mathrm{tw}}$ with constant
  $3$, by Theorem~\ref{asm-design-exists} for the class
  $\mathcal A\cap\mathcal C_G(\lambda)$. Its hypothesis (1) is the design
  property of the collection for $\lambda$ and $G$; (2) and (3) restrict the
  size properties of the collection to the smaller class; (4) is its analytic
  input.
\item The Tate-Shafarevich link is the last property of the collection.
\item The support condition holds along $\mathcal D_3$ for every family. At
  every member beyond the threshold of the twist selection, the tied data of
  the twist record give a Hecke realization of the member with normalized
  local compatibility and the Hecke ties
  (Definition~\ref{asm-support-local-compatibility}), and
  Lemma~\ref{asm-support-from-ties} gives the support condition there.
\item The irreducibility link holds eventually along $\mathcal D_3$ by
  Theorem~\ref{is-irreducibility-link-holds}, from Green's parametrization and
  the irreducibility inputs.
\item The rank equality at small level holds eventually along $\mathcal D_3$,
  by Lemma~\ref{asm-rank-from-ties} at every member at a large scale, applied
  to the selected base level $P$ and the tied data of the twist record, which
  determine its row side. Its hypothesis (1) is among the ties of $P$
  to the member (Definition~\ref{is-row-side-family}); (2) is
  Lemma~\ref{asm-primitive-rows-base}(3) for the realization $R_Q$ at $F_Q$:
  the family of definite forms of $P$, built from the square-integrable
  columns, satisfies the family conditions with the conditions on its inert
  indices, it has lift-distinct types and $R_Q$ has primitive rows at $F_Q$ by
  item (3)(c) of the selected rows, and the indefinite realization $Q$
  includes the weak characterization of $R_Q$ at $F_Q$ and identifies the
  definite subspaces and the rows of the complex data of $R_Q$ at $F_Q$ with
  those of $P$, so that $S_0$ and the definite side of $P$ are those of that
  lemma at level $0$; (3) is item (3)(d); (4) is the
  division rows (3)(e) at the tuples of $P$, which are all used, being the
  tuples of the indefinite realization $Q$; and (5) holds at the base tuple of
  $P$ by (Projections).
\item The conditional division link holds, under the size conditions, by
  Theorem~\ref{is-division-link-holds} with $C_{\mathrm{abs}}=2$, applied to
  $\mathcal D_3$ with the division kernel data that the square-integrable
  columns and the primitive rows give at every member. Its hypothesis, the
  division kernel bounds (Definition~\ref{is-division-kernel-bounds}) with
  constants $(2,C_H)$, holds eventually along $\mathcal D_3$ for every family
  that satisfies both size conditions, part by part. Part (1), the division
  packages, holds eventually along $\mathcal D_3$, as shown after this list.
  Part (2) is the multiplicity
  bound (Definition~\ref{pd-multiplicity-bound}), which holds by
  Theorem~\ref{pd-multiplicity-from-integrality} and the integrality inputs.
  Part (3) holds because every twisted lattice $\tilde M_\mu$ is stable under
  $\mathcal O_L$ (Joins), so every used tuple is cut. Part (4) holds for every
  family. Part (6) holds by Lemma~\ref{asm-division-kernel-complement}, since
  the integrality inputs give the second complement inputs
  (Definition~\ref{asm-division-kernel-errors}(2)) eventually along
  $\mathcal D_3$. Part (7) holds eventually along $\mathcal D_3$, at every
  member at every large scale. There the multiplicity condition holds for
  $R_Q$ at $F_Q$ and the map $s_E$ of $Q$, by the definite rows
  (Theorem~\ref{is-definite-rows}(5)); the base vector image is full, by
  (Joins); and the spectral data of the spectral condition are finite, by
  Lemma~\ref{asm-division-kernel-finite}: the decay of Wu type and the finite
  dimensionality of rows are orthogonality inputs, and the gap hypothesis
  holds there, as in the proof of Lemma~\ref{is-gram-clauses}. With finite
  spectral data, the
  spectral condition gives an isomorphism of the Bessel-period data of the
  member with those of its automorphic setting, a row dictionary $j$ and the
  column pin, hypotheses (1) to (3) of
  Lemma~\ref{asm-division-kernel-bessel}. The join of $Q$ (Joins) gives
  hypotheses (1) to (5) of Lemma~\ref{asm-division-kernel-join} for this $j$
  and the embedding $\sigma$ of the coefficient field of the member; the
  realization has its base place at the place of $\sigma$, so $\sigma_0$ is
  $\sigma$ or its conjugate. That lemma gives the kernel join, and
  Lemma~\ref{asm-division-kernel-bessel}, with the Bessel tie
  (Definition~\ref{asm-division-kernel-section}), which the row pin of the
  twist record contains, gives part (7). Part (9) holds at every member at
  every large scale. The
  columns $S_1$ of $P$ are the indefinite-side datum at level $1$ of the
  complex data of $R_Q$ at $F_Q$, by the indefinite realization; by
  Lemma~\ref{asm-primitive-rows-base}(2), with the inputs listed for the rank
  equality, its matrices $\iota_t$ are bijective and the modules $P_tM_t$ are
  lattices; and the weak characterization of $R_Q$ makes every place $w$ with
  $\mu_w\ne\mu_0$ one of the finitely many twisted places. So the glued
  lattice of $S_1$ is a lattice by Lemma~\ref{asm-division-kernel-glued}, and
  the division rows (3)(e) (Definition~\ref{is-division-rows}(1)) carry it, by
  the linear isomorphism $e_{RR}$, onto the glued lattice $RR$ of $D$, which
  is therefore finitely generated and spans. Part (10) holds for every family,
  by the size conditions of its members. Part (8) holds eventually along
  $\mathcal D_3$, at every member at every scale beyond the scale of the
  definite rows and the threshold of the twist selection, by
  Lemma~\ref{asm-division-kernel-rational} for the selected base level $P$ and
  the indefinite realization $Q$ of the selection bundle. The member satisfies
  the Bessel tie, one of its ties (Definition~\ref{pd-family}(4)). Hypothesis
  (1) of the lemma is in the square-integrable columns and item (3)(c); (2) is
  given by the definite rows (Theorem~\ref{is-definite-rows}(5)); (3) is part
  (5) of the division rows (3)(e); and (4) is the compatibility with periods
  (3)(f) at the base tuple of $P$, which is used, as all the tuples of $P$
  are, being the tuples of $Q$. Part (5) holds eventually along
  $\mathcal D_3$, as shown after this list.
\item The rank conditions and the global link with $(2,C_H,3,2,0)$ hold by
  Theorem~\ref{asm-global-link}, from the duty data, the second-layer
  conditions at the real place and the mass formulas.
\item The gain link holds eventually along $\mathcal D_3$ by
  Theorem~\ref{pd-gain-link-holds}, from the duty data and the twist
  selection.
\item The orthogonality link holds with $C_0=3$, under the size condition on
  $T_u$, by Theorem~\ref{is-orthogonality-link-holds}, from the twist
  selection, the spectral conditions, the orthogonality inputs and the volumes.
\end{itemize}

\emph{Part (1) of the division kernel bounds.} Eventually along $\mathcal D_3$,
$C_{sz}\ge0$, since $C_{sz}$ is a constant of the middle list, and the
definite rows give $C\ge0$ such that, at every member at all large scales, the
realization $R_Q$ of the selection bundle satisfies conclusions (1) to (6) of
Theorem~\ref{is-definite-rows} with $(C,C_H)$; conclusion (4) applies because
the bounded rows (3)(a) give every lift a single row map
(Definition~\ref{is-bounded-rows}(2)). Let $\mathcal F\in\mathcal D_3$ be such a
family that satisfies both size conditions, and the division bound at twisted
places with a constant $C_e$. Let $C_d$ be the constant of the size condition
on stabilizer denominators, let $c_0$ and $c_e$ be as in
Lemma~\ref{asm-division-kernel-record} for $C$, $C_H$, $C_e$ and
$\max(C_d,0)$, and put $C_{\mathrm{pk}}=\max(0,50c_0+10,60c_e)$. Let $x$ be
beyond the scales of the definite rows, of the division bound and of the size
condition on stabilizer denominators, the threshold of the twist selection
and the $x_0$ of Lemma~\ref{asm-division-kernel-record}. Fix a member $D$ at
$x$, with selected base level $P$ and selection bundle $Q$, whose realization
$R_Q$ is at the face $F_Q$ of $P$.

By Lemma~\ref{asm-division-kernel-used}(1) every tuple of $P$ is used, so the
division rows (3)(e) and the compatibility with periods (3)(f) hold at every
tuple of $P$. The family of definite forms of $P$ satisfies the family
conditions, by the square-integrable columns; with item (3)(c) and
conclusions (1) to (5) of Theorem~\ref{is-definite-rows}, the hypotheses of
Lemma~\ref{asm-division-kernel-package} hold, and it gives a division package
$\mathcal P$ on the scale data of the member, with rows $S_0$ and columns
$S_1$. Lemma~\ref{asm-division-kernel-record} applies to $\mathcal P$, since
$P$ has finitely many tuples, its row side is determined by the twist record
and $C_{sz}\ge0$; it gives $\lambda_0$, $S_w$, $C_1$, $d_g$ and the record
numbers of the words. The hypotheses of Lemma~\ref{asm-division-kernel-inputs}
hold: (1) is the division bound at twisted places at $D$; (2) is the tie of
the twists to the size data (3)(b), since the rows of $\mathcal P$ are $S_0$;
(3) is Lemma~\ref{asm-division-kernel-used}(2); and (4) is the definite-side
condition, conclusion (6) of Theorem~\ref{is-definite-rows}, whose premises
are the primitive rows, conclusions (1) to (5) and the record numbers.
Lemma~\ref{asm-division-kernel-assembly} then applies with
$C=C_{\mathrm{pk}}$ and $C_{\mathrm{abs}}=2$. Its hypothesis (1) holds by the
first sentence of this paragraph and conclusion (5) of
Theorem~\ref{is-definite-rows}; (2) by
Lemmas~\ref{asm-division-kernel-record} and~\ref{asm-division-kernel-inputs};
(3) because the members of the families of Definition~\ref{pd-family} have
these division lattices, and their heights are the explicit heights with
constant $C_{sz}$; and (4) because $L_a\ge0$ by the sizes of the member. It gives the division package of
part (1) at $D$, with constants $(C_{\mathrm{pk}},2)$.

\emph{Part (5) of the division kernel bounds.} Eventually along
$\mathcal D_3$: the second-layer conditions at the real place hold; the gain
link holds, by the proof of Theorem~\ref{pd-gain-link-holds} with the group
coherence of its Step 1 read off the tied data of the twist records; and the
conclusion of Lemma~\ref{asm-division-kernel-functional} holds for
$e_b(l)=4(\max(l_1,5)+3)+1$ and $e_c(l)=9(\max(l_1,5)+3)+2$. Let
$\mathcal F\in\mathcal D_3$ be such a family that satisfies both size
conditions, part (6) and the division bound at twisted places, with height
constant $C_H$, let $e_F$ be its function from
Lemma~\ref{asm-division-kernel-functional}, and put $c_t=\max(C_t,5)+3$, so
that $C_H\ge\max(4c_t+1,9c_t+2,e_F(c_{\mathrm{out}}),9c_t+1)+1$.

As in Step 2 of the proof of Theorem~\ref{asm-global-link}, the second-layer
conditions and the size condition on $T_u$ give the level data, the
commensurability of $N^\circ_u$ with $M^{\mathrm{ram}}_u$ and the finite
generation and spanning of every $N_u$; with the proviso of the chosen
constants of $\mathcal F\in\mathcal D_3$, every $N_u$ has positive rank.
Lemma~\ref{asm-slope-link} gives the slope link: the local data vanish above
$X''(n)$ at every member at every scale $x>x_1$, the ranks are those just
found, and the two size conditions bound $\log T_u$ and $\log D_\Gamma(u)$.
Lemma~\ref{asm-global-slopes-n}, whose depth bound holds for every family,
gives the slope bounds for $N$ at all large scales. At all large scales,
Lemma~\ref{asm-division-kernel-first-lattice} gives
$\mu_{\max}(M_u)\le x^{4c_t+1}b$, Lemma~\ref{asm-division-kernel-gauge} gives
$G_a+F_o-E_u\le x^{9c_t+2}b$ from part (6), and
Lemma~\ref{asm-division-kernel-functional}(2) gives
$\mu_{\max}(RR)+\overline{\log\kappa_R}\le x^{e_F(c_{\mathrm{out}})}b$.
Lemma~\ref{asm-division-kernel-height}, with $e_b=4c_t+1$, $e_c=9c_t+2$ and
$e_F(c_{\mathrm{out}})$, gives $s+\mu_{\max}(RR)\le x^{C_H}b$ at all large
scales, with $s$ formed from the explicit heights, which are the heights of
the family. This is part (5).
% Lean: the division kernel bounds are hBR: RemainingDivisionBounds.toRows_alt.
% Part (1), pkg, is wire_sel_alt = Top4Wire2.wire4F_alt: hDivR = DivRows3G.toRows and hDcPt =
% dcPt_of_G with tup_all_used_of_realisation, and hD0 = indefinite_realisation_base_rank_eq (Lemma
% asm-division-kernel-used); hP from PrimRows3 and IndefiniteRealisationConstantBounds (rRows, with
% lamLow from the one-map condition oneMap = (hRB ...).oneMap); packageOfPrimitiveLinks (Lemma
% asm-division-kernel-package); gamma_size_of_sizeData_alt and arithClause_rec (Lemma
% asm-division-kernel-record); inpAt_of_alt with DIVinN_at, hTw (TieTwRec), hD0, lift_exponent_of_least
% and the definite-side condition hd (Lemma asm-division-kernel-inputs); ptK_of_real_alt with
% fixedScale_tied, Fm.nprimeDef, Fm.h_eq and La_nonneg (Lemma asm-division-kernel-assembly). The
% constants are C_pk = max 0 (max (50 c0F + 10) (60 ceF C Ce (2 + 9 max Cd 0) c0F)) and x0 = max of the
% scales of rRows, oneMap, DIVinN, DGamSize and x0F; the threshold of the twist selection is that of
% the division kernel data RankOfDivisionKernel. The hypothesis bu (BuRational) of Wire4Alt is not
% used by wire4F_alt.
% Part (5), hh, is RestAll.bRest_all_alt (through GlobalTheory.Top.link_brest_alt):
% closed_funSharp_thr_alt with Eb = SlopeMaxB.EbMu and Ec = Row4C.EcHH (Lemma
% asm-division-kernel-functional), the second-layer conditions tie, LinkGCl.closed_linkG_of_tuSize_dGamSize
% (closed_linkG: linkG_concrete5 with gTiedF_of_sel_alt and admissible_tops_of_class), then
% LvlWire.lvlF_of_ties, muRank_of_lvl (proviso_of_working_class), SlopeId.linkMu_of_full
% (linkMu_of_sizes = J.linkMu with muZero_free, muFG_of_full, muTu_of_size, muDGam_of_size) and
% LvlWire.hh_of_lvl = height_bound_of_component_bounds (Lemma asm-division-kernel-height) with
% slope_lemma_division_side (slopesNF_of with placesF_min', Lemma asm-global-slopes-n),
% nuFullF_of_lvl, eventually_muMax_le_pow_ebMu_mul_of_tuSize (Lemma asm-division-kernel-first-lattice),
% gauge_form_bound_of_uniform_size (Lemma asm-division-kernel-gauge) and the bound of
% closed_funSharp_thr_alt.
% Done above: part (6), c2, is complement_two_of_items (Lemma
% asm-division-kernel-complement); part (7), c5, is selection_closed_of_wu_decay_alt
% (gapEv_closed_of, closed_finX_of_gapHyp with finXC_of hWu hRF = FinXP.finX_holds, Lemma
% asm-division-kernel-finite, complementFive_closed_of_image_v0 = adm13_of_kernel_join; Lemmas
% asm-division-kernel-join and asm-division-kernel-bessel); part (8), bu, is
% RatFullFam.buRat_rrFull_closed_alt with RatFullPt.buRational_of (Lemmas
% asm-division-kernel-descent to asm-division-kernel-rational); part (9), rr, is
% RatFullPt.rrFull_of (ColFacts.rr_lattice, DivRows3 sharpRR.lat); part (10), la,
% is Wire.la_of.
Theorem~\ref{asm-tied-from-links} shows that $\mathcal C_G(\lambda)$ has tied
members with exponents $2$ and $0$.
\end{proof}
\leanref{Development/AutomorphicArithmetic/ThetaCorrespondence/ThetaPairLifts/ThirdPlaceSeries/TopShape.lean}{UpperTheory.Top.shape_closed_link_n_of_green}
\leanref{Development/AutomorphicArithmetic/ThetaCorrespondence/ThetaPairLifts/ThirdPlaceSeries/ShapeOf.lean}{GlobalTheory.Fc.forall_goal_of_shape_of_hyps}

\subsection{Stages and parameters of the construction}\label{subsection-asm-stages-parameters}

\begin{definition}\label{asm-dependence-functional-witness-sizes}
Let $D$ be a definite-side datum (Definition~\ref{th-definite-datum}) with ranks $k,m,t,n$, $f=2n$,
$r=k+m-1$ and $N(j)=k+j(t+1)$ (Definition~\ref{th-ranks}); with host field $E$, root discriminant
$\mathrm{rd}_E=|d_E|^{1/[E:\mathbf Q]}$, auxiliary primes $q,d$, integer $Q_S$, level primes $g_0,g_1$
and ideal $\mathfrak c$ of $\mathcal O_K$, $K=\mathbf Q(\sqrt{-d})$ (Definition~\ref{hf-host}); with
the prime $q_0$ of Definition~\ref{asm-twist-condition}; and with Harish-Chandra
parameters whose $b$- and $e$-parts at an embedding $\tau$ are $b_l(\tau)$ and $e_l(\tau)$,
$1\le l\le m$ (Definition~\ref{th-hc-parameter}). Let size data for $D$ be given
(Definition~\ref{pd-size-data}), at the scale $(x,\varepsilon)$, with $b=\exp(x^{\varepsilon})$ and
cloud size $S_c$, and let $\gamma(1)$ be the integer of
Definition~\ref{asm-dependence-rows-parameters}(2) for $c=1$. The \emph{rows of the size data have
the sizes of the size data} if:
\begin{enumerate}
\item there are \emph{parameter choices} for the size data, the positive constants and the
  inequalities between the ranks, the host, the Harish-Chandra parameters and the quantities of the
  size data that the members of a family carry (Definition~\ref{pd-family}(4)), whose constant
  $C_{nW}$ is $10$, with \emph{$T$-parameters}: integers $T_O,n_W,n_P,D_p$ and constants with
  $n_P=8fr+4$, $D_p=f\binom r2$, $r+t=n_P+T_On_W$, $8\mid n_W$ and the bounds between them (those
  listed in Lemma~\ref{asm-dependence-functional-witness-numbers-blocks-computations}, with the
  constants of the $T$-parameters in place of those of $C$, and lower bounds on the constants);
\item the size data satisfy the \emph{size bounds} (listed in
  Lemma~\ref{asm-dependence-functional-witness-numbers-bounds-computations}), the further conditions
  of Definition~\ref{pd-size-data}: windows of exponent $\varepsilon/100$ about
  $x^{2+8\varepsilon}$, $\sqrt{x\delta_1}$ and $x^{9/2+10\varepsilon}$ for $n$, $m$ and $k$; inequalities
  that bound $t$, $Q_Sqd$, $\log q$, $\log d$, $\mathrm{rd}_E$, the number of top classes,
  $\Lambda_G$, the exponents $b_t$, the width $\eta$, $\Delta_S$, $w''$, $g_0$, $g_1$, $p_{de}$,
  $p'_{de}$, the gaps of the $b$-parts and the number of words in terms of the ranks, $x$ and $b$;
  among them $\log q,\log d\le x^{\varepsilon/5+\varepsilon/100}$ and $\mathrm{rd}_E\ge1$; and the
  splitting bound is $B\ge11$;
\item $m^4\le f$ and $m\ge44$;
\item $x^{2+\varepsilon}m^3\log\mathrm{rd}_E\le k$ and $mfN(m)^2d\le b$;
\item $fm\le(r+1)^3$;
\item $\log g_i\le9x^{\varepsilon}+1$ for $i=0,1$, and $\log q_0\le x^{\varepsilon}$;
\item $\log\mathrm{rd}_E\le\log(Q_Sqd)\le x^{\varepsilon}/100$;
\item $S_c+(m+1)\gamma(1)+\tfrac12\le b$;
\item $b_l(\tau)>0$ and $e_l(\tau)\le4b$ for every embedding $\tau$ and every $1\le l\le m$;
\item $\mathfrak c=\mathcal O_K$.
\end{enumerate}
\end{definition}
\leanref{Development/AbelianVarietyTheory/BlockTheory/DivisorBase/LevelPlusSubgroup.lean}{GlobalTheory.IndefiniteRowSideFullHecke.LevelPlusSubgroup.SizeRows}
\leanref{Formalism/LocalRingCardinality/DataParams.lean}{Formalism.BasicNotions.ParameterChoice'}
\leanref{Formalism/LocalRingCardinality/DataParamsT.lean}{Formalism.BasicNotions.ParamsT}
\leanref{Formalism/SharpRestriction/ElevenLe.lean}{Formalism.MainNotions.SizesAt}

\begin{definition}\label{asm-dependence-functional-witness-setting}
Let $D$ be a definite-side datum with host field $E/F$, ranks $k,m,t$, $r=k+m-1$ and
$N(j)=k+j(t+1)$, level primes $p_{de},p'_{de}$, and Harish-Chandra parameters with parts $u_i(\tau)$,
$v_i(\tau)$, $b_j(\tau)$, $e_j(\tau)$; for a word $u$ write $b_j(u,\tau)$ for the $b$-parts of
$\lambda_u(\tau)$ (Definition~\ref{th-definite-datum}(2)). Let size data for $D$ be given. Let a
family of definite forms of $D$ for a finite group $G$ be given, with tuples $\mathcal T$ and base
tuple $\mu_0$, and let $R$ be a realization of $D$ for it (Definition~\ref{is-realisation}), with
level system $(\xi_t,Q_t)$, characters $\mu_{t,j}=\mu^0_j\xi_t$ and $\eta$ (Definition~\ref{th-grid}),
standard character $\psi$, map $w\mapsto\mu_w$, twisted places $\mathrm{Tw}$, places $P_g$ with the
isomorphism $G\cong\prod_{P\in P_g}GL_r(k_P)$, characters $\delta,\delta'$ of its columns, and Weil
data $\omega_{t,k}$ for $(Y_k,Y_{k+1})$, $0\le k\le m-2$, and $\omega_{t,m-1}$ for $(Y_{m-1},X')$ at
each tuple $t$. For a prime $\ell$ put $\delta_\ell=\delta'$ if $\ell=p_{de}$ and $\delta_\ell=\delta$
otherwise. Let $F_R$ be a face of the family. Then $R$ is a \emph{realization in an automorphic
setting}, with these size data, at $F_R$ if:
\begin{enumerate}
\item the rows of the size data have the sizes of the size data
  (Definition~\ref{asm-dependence-functional-witness-sizes});
\item the series idempotents of $D$ are tied to its labels;
\item at every embedding $\tau$: the parameters interlace strictly,
  $u_1>v_1>u_2>\dots>v_{k-1}>u_k>-b_1$ and $b_j<e_j<b_{j+1}$; $e_j-b_j\ge(t+3)/2$ and
  $b_{j+1}-e_j\ge(t+3)/2$; any two consecutive terms of $v_1,\dots,v_{k-1},-b_1,\dots,-b_m$ differ by at
  least $100r^4$; and $b_{j+1}(u,\tau)-b_j(u,\tau)\ge4N(m)^3+2t+6$ and $\ge256r^4$ for every word $u$
  and $1\le j<m$;
\item $F$ is the maximal totally real subfield of $E$;
\item $R$ satisfies the core conditions at $F_R$ (Definition~\ref{is-eigenvalue-lift-data});
\item the tuples are those of the twisted places: $\mu_w=\mu_0$ for $w\notin\mathrm{Tw}$ and
  $\mu_w\ne\mu_0$ for $w\in\mathrm{Tw}$; every tuple other than $\mu_0$ is some $\mu_w$;
  $\mu_{\bar w}=\mu_w$; and for $w\in\mathrm{Tw}$, with $\ell$ the residue characteristic of $w$,
  $\xi_{\mu_w}=\delta_\ell\psi_w$ for a Hecke character $\psi_w$ of $E$ of $\ell$-power order, trivial on
  $\mathbf A_F^\times$, tamely ramified, with trivial components at the places above $\ell$, and whose
  set of ramified places is a finite set of places of $E$ split over $F$;
\item $R$ \emph{records the group}: there are an isomorphism $e$ of $G$ with the group
  $G_0=\prod_{k=0,1}\prod_{v\in P_G(k)}GL_r(\mathbf F_{g_k})$ of $D$, a bijection from the places
  $v\in P_G(k)$, $k=0,1$, to $P_g$ whose image of $v$ lies over the prime of $K$ above $g_k$ chosen by
  the columns, and isomorphisms $\mathbf F_{g_k}\cong k_P$ at corresponding places, such that the
  idempotent of the family has the coefficients of the type idempotent $e_0$ of $D$ under $e$ and the
  isomorphism $G\cong\prod_PGL_r(k_P)$ of $R$ is $e$ followed by these isomorphisms;
\item the \emph{guard} holds: two places of one set $Q_t$ with the same residue characteristic lie
  over the same prime of $K$, and no place of a set $Q_t$ contains $Q_Sq\,\mathrm N(\mathfrak c)$;
\item $\mathcal T$ is finite;
\item for every tuple $t$ the Weil data $\omega_{t,k}$, with the characters
  $(\mu_{t,k+1}\eta^{-1},\mu_{t,k+1})$ if $k+1$ is odd and $(\mu_{t,k+1},\mu_{t,k+1}\eta^{-1})$ if it is
  even, and $\omega_{t,m-1}$ with these characters for $k+1=m$, are \emph{full}: besides being Weil
  data (Definition~\ref{th-weil-datum}(3)), they fix the lattice vacuum at every finite place $v$ of
  $F$ prime to $2$ and not split in $E$, at which both characters are unramified, under the
  elements of the stabilizers of self-dual lattices supported at $v$, when the two ranks differ and
  one of them is even; and at every archimedean place they act on Fock vectors by the standard
  formulas for the Lie algebras, shifted by the infinity types of the two characters;
\item $\psi$ is standard (Definition~\ref{th-weil-datum}(1));
\item the basis of $X'$ is orthonormal and spans its lattice over $\mathcal O_E$, and the adelic
  matrix of $X'$ in the level data of $R$ is the identity.
\end{enumerate}
\end{definition}
\leanref{Development/AbelianVarietyTheory/Statements/HeckeThetaIntegrals/FormulaAt.lean}{GlobalTheory.RowClaim.HidaSetting}
\leanref{Formalism/LocalToGlobal/IndefiniteReal/LinksLemmas.lean}{Formalism.IndefR.Links.HostHyp}
\leanref{Formalism/LocalToGlobal/IndefiniteReal/OneMapT.lean}{Formalism.IndefR.FWT.QGuard}
\leanref{Formalism/LocalToGlobal/IndefiniteReal/WeilFull.lean}{Formalism.IndefR.WeilFull.IsWeilFull}
\leanref{Formalism/LocalToGlobal/IndefiniteReal/LinksChar.lean}{Formalism.IndefR.Links.RealDat.grp_record}

\begin{definition}\label{asm-grid-evaluation}
Let $D$ be a definite-side datum (Definition~\ref{th-definite-datum}) with $m$ lifts, host field
$E/F$ and the spaces $Y_0,\dots,Y_{m-1}$ and $X'$ of Definition~\ref{th-spaces}. Let a family of
definite forms of $D$ for a finite group be given, and let $R$ be a realization of $D$ for it
(Definition~\ref{is-realisation}), with standard character $\psi_R$, archimedean family $\theta_R$,
the level system given by its level and character data (Definition~\ref{th-grid}), and Weil data
$\omega_{R,t,k}$ for the pairs $(Y_k,Y_{k+1})$ and $(Y_{m-1},X')$ at each tuple $t$. Let $C_F$ and
$b$ be real numbers.
\begin{enumerate}
\item An \emph{evaluation grid} for $R$ with constant $C_F$ at $b$ is a theta grid for the level
  system of $R$ with constants $(1,C_F,b)$ (Definition~\ref{th-grid}), with spaces $V_{i,j}$, $P_i$
  and $Q_j$. Its places consist of a \emph{base place} $v_0$ of $F$ above the auxiliary prime
  $q_0$ of the host of $D$ and, for every place $v$ of $F$, a place $w(v)$ of $E$; they define an
  ideal $\mathfrak m$ of $\mathcal O_E$, the \emph{modulus}, a product of powers of the places
  $w(v)$ for $v=v_0$ and for $v$ in a finite set $S_1$. Write $\Lambda_X$ for the lattice of a
  grid space $X$, $K_X$ for the principal congruence subgroup of modulus $\mathfrak m$ of the
  stabilizer of $\Lambda_X$ in $U(X)(\mathbf A_{F,f})$, and $K^+=K_{Q_0}$.
\item A \emph{Weil datum over $R$} for the grid is a grid Weil datum (Definition~\ref{th-grid})
  whose standard character is $\psi_R$, whose archimedean family is $\theta_R$, whose Cartan
  involutions and orthonormal bases on $V_{0,j}=Y_j$ ($j<m$) and on $V_{0,m}=X'$ are those of $R$,
  and whose Weil datum for the pair $(V_{0,k},V_{0,k+1})$ at the tuple $t$ is $\omega_{R,t,k}$.
\item A Haar measure on $U(X)(\mathbf A)$ is \emph{normalized at} an open compact subgroup
  $K_f\subset U(X)(\mathbf A_{F,f})$ if it is the product of the Haar measure of the finite adelic
  group that gives $K_f$ volume $1$ and of the fixed Haar measures on the groups $U(p,q)$ at the
  archimedean places. A \emph{row volume} for the grid consists of Haar measures on the groups
  $U(X)(\mathbf A)$ of the grid spaces $X$ with fundamental domains for $U(X)(F)$, and of a Haar
  measure with a fundamental domain on the unipotent radical $N$ of the stabilizer in $U(V_{m,0})$ of
  the flag attached to the embedding $Q_0\subset V_{m,0}$, such that the measure on
  $U(V_{i,0})(\mathbf A)$ is normalized at $K_{V_{i,0}}$ for $0\le i\le m$, the measure on
  $U(Q_0)(\mathbf A)$ is normalized at $K^+$, and the fundamental domain of $N$ has measure $1$.
\item Let a Weil datum over $R$ and a row volume be given. For $0\le i<m$, the grid Weil datum
  gives a coset vector $f_{i,\kappa}$ for every coset $\kappa$ of the pair lattice of
  $(V_{i+1,0},V_{i,0})$ modulo $\mathfrak m$. A \emph{real vector} is a family
  $\rho=(\rho_i)_{0\le i<m}$ of polynomials in the creation operators of the pairs
  $(V_{i+1,0},V_{i,0})$, and a \emph{coset datum} is a family $\Phi=(\Phi_i)_{0\le i<m}$ of such
  cosets. For a function $u$ on $U(Y_0)(\mathbf A)$, the \emph{iterated theta lift}
  $\Theta_{\rho,\Phi}(u)$ is the function on $U(V_{m,0})(\mathbf A)$ obtained from $u$ by applying,
  for $i=0,1,\dots,m-1$ in turn, the theta lift from $U(V_{i,0})$ to $U(V_{i+1,0})$ of
  Definition~\ref{th-weil-datum}(4) with the vector $f_{i,\Phi_i}\otimes\rho_i$ and the measures of
  the row volume. It \emph{converges} if each of these integrals converges absolutely at every point.
\item The grid gives a finite-dimensional representation $(\varrho,W)$ of $U(Q_0)$, the
  \emph{coefficients}, built from $\Lambda_{Q_0}$ and $\theta_R$. A \emph{column} is a function $\Psi$
  on $U(Q_0)(\mathbf A_{F,f})$ with values in $W$ such that $\Psi(\gamma gk)=\varrho(\gamma)\Psi(g)$
  for $\gamma\in U(Q_0)(F)$ and $k\in K^+$. For a linear functional $\xi$ on $W$ put
  $\varphi_{\xi,\Psi}(h)=\xi\bigl(\varrho(h)^{-1}\Psi(h_f)\bigr)$ for $h\in U(Q_0)(\mathbf A)$. For
  functions $\varphi$ on $U(Q_0)(\mathbf A)$ and $G$ on $U(V_{m,0})(\mathbf A)$, the \emph{Bessel
  period} is
  \[
   B^+(\varphi,G)=\int_{\mathcal F_{Q_0}}\int_{\mathcal F_N}\varphi(h)\,G\bigl(n\,\iota(h)\bigr)\,
   \overline{\psi_N(n)}\,dn\,dh,
  \]
  where $\iota\colon U(Q_0)\to U(V_{m,0})$ is the embedding as the pointwise stabilizer of
  $Q_0^{\perp}$, $\psi_N$ is the character of $N$ defined by $\psi_R$, and $\mathcal F_{Q_0}$,
  $\mathcal F_N$ are the fundamental domains of the row volume. It \emph{converges} if the inner
  integral converges absolutely for every $h$ and the outer one converges absolutely.
\item Let $w$ be a nonzero smooth vector of the model of $U(Y_0)$ in $R$. For an element $t$ of the
  space of theta multiplicities of $R$ at the trivial character, let $u_{w,t}$ be the function on
  $U(Y_0)(\mathbf A)$ obtained by applying $t$ to $w$. \emph{Functionals separate $w$} for the Weil
  datum and the row volume if there are a real vector $\rho$ and a linear functional $\xi$ on $W$
  such that for every nonzero such $t$ there are a column $\Psi$ and a coset datum $\Phi$ for which
  $\Theta_{\rho,\Phi}(u_{w,t})$ and $B^+\bigl(\varphi_{\xi,\Psi},\Theta_{\rho,\Phi}(u_{w,t})\bigr)$
  converge and $B^+\bigl(\varphi_{\xi,\Psi},\Theta_{\rho,\Phi}(u_{w,t})\bigr)\ne0$.
\item \emph{Functionals separate} if for every $\varepsilon$ and every $C_F\ge6510$ there is $x_0$
  with the following property. Let $D$ be a definite-side datum with size data of exponent
  $\varepsilon$ at a scale $x\ge x_0$, with $b=e^{x^{\varepsilon}}$, whose rows have the sizes of the
  size data (Definition~\ref{asm-dependence-functional-witness-sizes}) and whose series idempotents
  are tied to its labels. Let a family of definite forms of
  $D$ for a finite group be given, with base tuple $\mu_0$, and let $R$ be a realization of $D$ for it
  whose level characters are those of $\mu_0$ and which records the group. Then for every
  evaluation grid for $R$ with constant $C_F$ at $b$, every Weil datum over $R$ for it, every row
  volume for it and every nonzero smooth vector $w$ of the model of $U(Y_0)$ in $R$, functionals
  separate $w$.
\end{enumerate}
\end{definition}
\leanref{Development/AbelianVarietyTheory/Definitions/Divisors/BasePlaceNotMemChiGridBadPlaces.lean}{GlobalTheory.ChiGrid.Places}
\leanref{Development/AbelianVarietyTheory/Definitions/Divisors/BasePlaceNotMemChiGridBadPlaces.lean}{GlobalTheory.ChiGrid.GridData.Weil.Over}
\leanref{Development/AbelianVarietyTheory/Definitions/Divisors/BplusOfEquivariant.lean}{GlobalTheory.FunObj.Vol}
\leanref{Development/AbelianVarietyTheory/Definitions/Divisors/BplusOfEquivariant.lean}{GlobalTheory.FunObj.Vol.IsRow}
\leanref{Development/AbelianVarietyTheory/Definitions/Divisors/BplusOfEquivariant.lean}{GlobalTheory.FunObj.thetaChi}
\leanref{Development/AbelianVarietyTheory/Definitions/Divisors/BplusOfEquivariant.lean}{GlobalTheory.FunObj.Bplus}
\leanref{Development/AbelianVarietyTheory/Definitions/Divisors/BplusOfEquivariant.lean}{GlobalTheory.FunObj.fPsiPhi}
\leanref{Development/AbelianVarietyTheory/BlockTheory/DivisorBase/LevelPlusSubgroup.lean}{GlobalTheory.IndefiniteRowSideFullHecke.LevelPlusSubgroup.Cols}
\leanref{Development/AbelianVarietyTheory/BlockTheory/DivisorBase/LevelPlusSubgroup.lean}{GlobalTheory.IndefiniteRowSideFullHecke.LevelPlusSubgroup.Separates}
\leanref{Development/AbelianVarietyTheory/BlockTheory/DivisorBase/LevelPlusSubgroup.lean}{GlobalTheory.IndefiniteRowSideFullHecke.LevelPlusSubgroup.FunctionalsSeparate}

\begin{definition}\label{asm-stage-construction}
The definite-side data of the construction are assembled in \emph{stages}. A \emph{context} $K$
consists of a scale $(x,\varepsilon)$, the constants of a family (Definition~\ref{pd-family}) and a
rank datum with its level data. \emph{Base data} $B$ consist of ranks, a host field and coefficient
fields over the rank datum, with a \emph{framing} $f$; \emph{twist data} $T$ consist of a finite set
of \emph{twisting primes} and twisting characters; a \emph{column choice} $w$ consists of columns
over $(B,f,T)$ with their laws and pins; a \emph{later stage} $Z$ adds the row, division, size,
norm and jet data, and $(B,w,Z)$ determine a definite-side datum $D(w,Z)$; and a \emph{core} at a
stage consists of the twisting, level and character objects of this datum. A \emph{slot} is a
condition on the cores of the stages (Definition~\ref{is-slot-rows}).
\begin{enumerate}
\item A property of contexts holds \emph{eventually in the context} if for every small exponent
  $\varepsilon>0$ and for every choice of the constants of a context, made in the order in which
  they are chosen, it holds at every context with these constants at every large enough scale.
\item A property of base data, framings and twist data holds \emph{at large base data} in a context
  $K$ if there are finite sets of algebraic numbers and of primes, depending on the base data and
  the framing, and for the algebraic numbers also on the set of characters that the framing and the
  twist data give in (5), such
  that the property holds whenever the base data lie over the rank datum of $K$, their coefficient
  fields contain the given algebraic numbers, and the twisting primes contain the primes of the
  first set and avoid those of the second. A property $\mathcal Q$ of column choices holds \emph{at all
  large stages after} a property $\mathcal Q_0$ if, eventually in the context and at large base data, for all base
  data whose coefficient field carries an exact field datum for the framing and whose inclusion of
  $X'$ into $X_m$ carries the standard basis of $X'$ to the first vectors of that of $X_m$, and all
  twist data whose twisting primes contain every prime dividing $g_0$, $g_1$ or $d$ and every prime
  $p\le X''(n)$, every column choice that satisfies $\mathcal Q_0$ satisfies $\mathcal Q$.
\item Complex data at a face (Definition~\ref{is-complex-data}) pass the \emph{row test} if there
  are finitely many linear functionals on the top space $E$ that separate its points and whose
  values at $\Theta_tj_t(y)$ lie in $\sigma_0(L)$ for every tuple $t$ and every row form $y$.
\item The \emph{rows of the construction are rational} if the following holds at all large stages
  after the column choices that are pinned, supported and coherent with the twist data. For every
  later stage $Z$ and every realization $z$ of $D(w,Z)$ built at the stage
  (Definition~\ref{is-realisation}) such that the group of $z$ is trivial at the stage, the scalars
  $\kappa$ of $D(w,Z)$ are defined, $z$ exists at every place and lies in the geometry of the
  stage, $z$ has a core, the group and jet stages satisfy the stage conditions, and the set of
  twisted places of $z$ is the set $W_1$ of the stage and consists of places with residue
  characteristic among the twisting primes, the complex data of $z$ at its face pass the row test.
\item Let $C_F$ be real. A stage $(K,B,f,T)$ has the \emph{grid property for $C_F$} if the
  following holds for every later stage. Let $D$ be its definite-side datum, with size data at the
  scale of $K$ that satisfy the recursion conditions, and let $R$ be a realization of $D$ in an
  automorphic setting, with these size data at a face
  (Definition~\ref{asm-dependence-functional-witness-setting}), whose columns are given by the
  framing $f$, whose twisted places lie off $q_0$, and every character $\xi_t$ of whose level system
  lies in the set consisting of the trivial character, the characters $\delta$ and $\delta'$ of $f$,
  and the characters $\delta_\ell\psi_w$, where $w$ runs over the finite places of the coefficient
  field, $\psi_w$ is the character that $T$ chooses at $w$, and $\delta_\ell=\delta'$ if the residue
  characteristic $\ell$ of $w$ is $p_{de}$ and $\delta_\ell=\delta$ otherwise. Then for every
  evaluation grid $G_0$ for $R$
  with constant $C_F$ at $b$ (Definition~\ref{asm-grid-evaluation}) there is an evaluation grid $G$
  for $R$ with constant $C_F$ at $b$, with the base place, the place above it and the characters
  $\chi_1,\dots,\chi_m$ of $G_0$, which carries a \emph{full evaluation choice}: a Weil datum over
  $R$, a row volume, a real vector, a linear functional on the coefficients and the further choices
  of the evaluation, satisfying the evaluation conditions and the field conditions, and such that
  the norms at the bottom of $G$ and the ratios of its toric periods lie in the coefficient field.
  The \emph{grid with constant $C_F$ evaluates} if the grid property for $C_F$ holds eventually in
  the context and at large base data.
\end{enumerate}
\end{definition}
\leanref{Development/AbelianVarietyTheory/Definitions/AutomorphicMembership/CmQuadraticFormEqAlternatingForm.lean}{GlobalTheory.Mem.Top.EvG}
\leanref{Development/AbelianVarietyTheory/Definitions/AutomorphicMembership/CmQuadraticFormEqAlternatingForm.lean}{GlobalTheory.Mem.Top.large₂}
\leanref{Development/AbelianVarietyTheory/Definitions/AutomorphicMembership/NotBaseCondFalse.lean}{GlobalTheory.Mem.Top.GroupRetwisting.Step}
\leanref{Development/AbelianVarietyTheory/Definitions/AutomorphicMembership/PositionOfIntro.lean}{GlobalTheory.Mem.Top.Bl.StepW}
\leanref{Formalism/LocalToGlobal/IndefiniteReal/RationalAdjoints/RatZOfTests.lean}{Formalism.IndefR.Rat.RowTest}
\leanref{Development/AbelianVarietyTheory/Definitions/GlobalTwisting/GeometricWeakTiesOfStrong.lean}{GlobalTheory.Fc.S0.RowRat}
\leanref{Development/AbelianVarietyTheory/Definitions/AutomorphicMembership/GKeyFin.lean}{GlobalTheory.Mem.LiteralAutomorphy.GridAt}
\leanref{Development/AbelianVarietyTheory/Definitions/AutomorphicMembership/GKeyFin.lean}{GlobalTheory.Mem.LiteralAutomorphy.GridA}
\leanref{Development/AbelianVarietyTheory/Definitions/AutomorphicMembership/CmQuadraticFormEqAlternatingForm.lean}{GlobalTheory.Mem.Top.TwDat.set}

\begin{definition}\label{asm-dependence-parameters}
Let local Langlands maps be given: for every number field and every finite place, a map
satisfying the local Langlands laws (Definition~\ref{gal-llc-laws}).
\begin{enumerate}
\item An \emph{extension} $N$ is a condition on the full data of a stage: its context, base data,
  framing, twist data, column choice, later stage, set of twisted places, base point, theta datum,
  twist choices and core. An extension $N'$ is \emph{stronger} than $N$ if every stage that
  satisfies $N'$ satisfies $N$. An extension $N$ determines \emph{parameters} $p=p(N)$ for the
  construction of bundles (Theorem~\ref{is-definite-side-condition-holds}): a core is \emph{built
  from} $p$ if it is the core of a stage that satisfies $N$ and the standard conditions of the
  construction; a \emph{good core} of $p$ is a core built from $p$ that satisfies further conditions
  on its twists and its geometry, which do not involve $N$; and the link data $\Lambda_p$ link a
  bundle to a base point when the bundle is built from $p$ at that point, with a witness of its
  construction. Write $Q^{\mathrm b}_p$ and $Q_p$ for the slots of the built and of the good cores
  of $p$.
\item Let $Q^{\mathrm b}$ and $Q$ be slots. A property $r$ of slots is \emph{linked for
  $(Q^{\mathrm b},Q)$} if there is a slot $R$ such that, eventually in the context and at large
  base data, every core in $Q^{\mathrm b}$ satisfies $R$, and $r$ holds for the slot $Q\cap R$. For
  link data $\Lambda$, a property $H$ of classes and selections is \emph{linked for
  $(Q^{\mathrm b},Q)$ with $\Lambda$} if ``$H$ holds along $(S,\Lambda)$''
  (Definition~\ref{is-slot-rows}) is linked for $(Q^{\mathrm b},Q)$. A property $r$ of a slot and
  link data is \emph{linked at $p$} if $r(\cdot,\Lambda_p)$ is linked for
  $(Q^{\mathrm b}_p,Q_p)$: there is a slot $R$ such that, eventually in the context and at large
  base data, every core built from $p$ satisfies $R$, and $r$ holds for the slot of the good cores
  of $p$ that satisfy $R$ and the link data $\Lambda_p$. A property $H$ of classes and selections
  is \emph{linked at $p$} if it is linked for $(Q^{\mathrm b}_p,Q_p)$ with $\Lambda_p$.
\item The \emph{law extension} requires the real expansion data of the core and the rationality
  law of the datum of the stage: if the values of its map of weights and degrees are rational, so is
  the datum. The real expansion data include a condition on generators: for $k=0,1$ and every
  $d\ge1$, the generator of $\mathbf F_{g_k^d}^\times$ that the group data of the core choose at the
  level prime $g_k$ of the host is the generator $\gamma_d$ of the label datum of $g_k$, the
  labelling of $\bar{\mathbf F}_{g_k}$ (Definition~\ref{lit-labelling}) fixed for each prime. For a
  real number $c$, the \emph{list extension}
  $N_{\mathrm{list}}(c)$ requires the law extension and, in addition, that the level characters of
  the realization of the theta datum of the stage have an evaluation grid with constant $c$ at
  $b=e^{x^{\varepsilon}}$ whose characters take values in the coefficient field and whose scalars
  $\lambda$ are tied to the twisted places. The \emph{parameters of dependence $c$} are
  $p(N_{\mathrm{list}}(c))$.
\item The \emph{origin extension} $N_{\mathrm{orig}}(c)$ requires $N_{\mathrm{list}}(c)$ and, in
  addition, origin data: the stage comes from a stage of the second pass, listed for $c$, with a
  strictly Henselian acyclic level, the same set of twisted places and the same twist choices. For
  an extension $N$ and a property $A$ of stages $(K,B,f,T)$, the \emph{extension of $N$ by $A$}
  requires $N$ and $A$. Write $p_c$ for the parameters of the extension of $N_{\mathrm{orig}}(c)$ by
  the grid property for $c$ (Definition~\ref{asm-stage-construction}(5)).
\item Let $Q^{\mathrm b}$ and $Q$ be slots. The \emph{member property} for $(Q^{\mathrm b},Q)$
  consists of four conditions. \emph{Base condition}: eventually in the context, there are base data
  over the rank datum whose coefficient fields carry an exact field datum and contain any two given
  finite sets of algebraic numbers, and these conditions are stable under enlarging the base data.
  \emph{First pass}: at large base data with untwisted twist data there is a core in $Q^{\mathrm b}$.
  \emph{Export}: at large base data with untwisted twist data every core in $Q^{\mathrm b}$ exports
  its twisting objects. \emph{Second pass}: at large base data, from every exported core in
  $Q^{\mathrm b}$ with untwisted twist data one passes, for any finite sets of primes to avoid and of
  algebraic numbers to contain, to enlarged base data with the same framing and to twist data with
  the same twisting primes, avoiding the first set and with coefficient field containing the second,
  at which there is a core in $Q^{\mathrm b}\cap Q$ with all its twists. The member property
  \emph{of $p$} is the member property for $(Q^{\mathrm b}_p,Q_p)$.
\item \emph{Bundles exist at large bases} for $p$ if there is a slot $R$ such that, eventually in
  the context and at large base data, every core built from $p$ satisfies $R$, and, eventually,
  every base point of the good cores of $p$ that satisfy $R$ has, for every $c$, a bundle that
  $\Lambda_p$ links to it.
\end{enumerate}
\end{definition}
\leanref{Development/AutomorphicArithmetic/ThetaCorrespondence/ThetaPairLifts/ThirdPlaceSeries/CoreAltAltBis.lean}{GlobalTheory.Mem.Top.Bl.Parameters}
\leanref{Development/AutomorphicArithmetic/ThetaCorrespondence/ThetaPairLifts/ThirdPlaceSeries/CoreAltAltBis.lean}{GlobalTheory.Mem.Top.Bl.parametersOfExtension}
\leanref{Development/AutomorphicArithmetic/ThetaCorrespondence/ThetaPairLifts/ThirdPlaceSeries/TwistedBundle.lean}{GlobalTheory.Mem.Top.Bl.BundleLinkage}
\leanref{Development/AbelianVarietyTheory/Definitions/GlobalTwisting/AllRows.lean}{GlobalTheory.Fc.Lg1}
\leanref{Development/AutomorphicArithmetic/ThetaCorrespondence/ThetaPairLifts/ThirdPlaceSeries/RowGood.lean}{GlobalTheory.Fc.RL.LinkedAt}
\leanref{Development/AutomorphicArithmetic/ThetaCorrespondence/ThetaPairLifts/ThirdPlaceSeries/SameGaloisLevelExists.lean}{GlobalTheory.HeightCounting.Up.extensionWithLaw}
\leanref{Development/AutomorphicArithmetic/ThetaCorrespondence/ThetaPairLifts/ThirdPlaceSeries/SettingsFormsExistList.lean}{GlobalTheory.HeightCounting.Up.extensionWithList}
\leanref{Development/AbelianVarietyTheory/Definitions/AutomorphicMembership/SettingsFormsExistList.lean}{GlobalTheory.HeightCounting.Up.ListPin}
\leanref{Development/AutomorphicArithmetic/ThetaCorrespondence/ThetaPairLifts/ThirdPlaceSeries/Prm1.lean}{GlobalTheory.Mem.Top.Bl.parametersOfDependence}
\leanref{Development/AutomorphicArithmetic/ThetaCorrespondence/ThetaPairLifts/ThirdPlaceSeries/ExtensionWithOrigin.lean}{GlobalTheory.HeightCounting.Up.extensionWithOrigin}
\leanref{Development/AutomorphicArithmetic/ThetaCorrespondence/ThetaPairLifts/ThirdPlaceSeries/ExtensionWithAtProperty.lean}{GlobalTheory.Mem.LiteralAutomorphy.extensionWithAtProperty}
\leanref{Development/AbelianVarietyTheory/Definitions/AutomorphicMembership/ThetaChr.lean}{GlobalTheory.Mem.Top.Bl.MemberE}
\leanref{Development/AbelianVarietyTheory/Definitions/AutomorphicMembership/NotBaseCondFalse.lean}{GlobalTheory.Mem.Top.GroupRetwisting.SecondG}
\leanref{Development/AbelianVarietyTheory/Definitions/GlobalTwisting/EventuallyAllScaleContextsAndOfBoth.lean}{GlobalTheory.Fc.ExBL}

\begin{definition}\label{asm-open-parameter-statements}
Let $p$ be parameters (Definition~\ref{asm-dependence-parameters}). For a class $\mathcal A$ of
families and a selection $s$ for $\mathcal A$, write $P$ for the row-side family and $Q$ for the
indefinite realization of the selection bundle of the selected bundle at a member
(Definition~\ref{asm-definite-side-hypotheses}), and $\mathcal W(\mathcal A)$ for the working
subclass of $\mathcal A$ (Theorem~\ref{pd-period-integrality-along-families}). The \emph{open
parameter statements} hold at $p$ if each of the following properties of $\mathcal A$ and $s$ is
linked at $p$. In (a) to (e) and (g) the property is required at every member of every family of
$\mathcal A$ at every scale $x>x_1$; in (f) and (h) it is required eventually along
$\mathcal W(\mathcal A)$ and eventually along each family.
\begin{enumerate}
\item[(a)] The realization of $Q$ has primitive rows and its family of definite forms has
  lift-distinct types (item (3)(c) of Definition~\ref{asm-definite-side-hypotheses}).
\item[(b)] The rows hold with $\Lambda_3$, the division rows hold, and the division data are
  compatible with the periods (items (3)(d), (3)(e) and (3)(f)).
\item[(c)] $Q$ is joined (Definition~\ref{is-realisation-join}) for the automorphic setting of the
  spectral condition.
\item[(d)] The selected bundle has a comparison part.
\item[(e)] The labels of the rows are tied to an automorphic setting: there are an automorphic
  setting for the datum $D$ of the member, a good row operator on it and an isometry $\iota$ of the
  Hilbert space of rows of $D$ onto that of the setting, such that the automorphy of $D$ is tied to
  the setting, the row dictionary is tied to the operator and to $\iota$, every label of a row of
  $D$ is a label, every label tuple is good, and $\iota$ carries the row of each type $\sigma$ into
  the space of the label tuple of $\sigma$.
\item[(f)] Every row of the definite-side datum is good.
\item[(g)] The geometric weak compatibility: the selected bundle has a comparison part whose point
  has the normalized embedding of its setting and satisfies the class-sum pin, the compatibility
  conditions and the further conditions of Definition~\ref{is-point} (hypothesis (4) of
  Theorem~\ref{is-definite-rows}).
\item[(h)] The multiplicity condition of the realization of $Q$ at its face holds, and the
  archimedean scalars $\kappa$ of the columns of $P$ are those of the division data (hypotheses (7)
  and (8) of Theorem~\ref{is-definite-rows}).
\item[(i)] The definite-side condition holds for the selected bundles, in the form of hypothesis (9)
  of Theorem~\ref{is-definite-rows}.
\end{enumerate}
\end{definition}
\leanref{Development/AutomorphicArithmetic/ThetaCorrespondence/ThetaPairLifts/ThirdPlaceSeries/GeometricWeakTies.lean}{GlobalTheory.Fc.S0.OpenParameterStatements}
\leanref{Development/AutomorphicArithmetic/ThetaCorrespondence/ThetaPairLifts/ThirdPlaceSeries/RowGood.lean}{GlobalTheory.Fc.RL.RowLabelCondition}
\leanref{Development/AutomorphicArithmetic/ThetaCorrespondence/ThetaPairLifts/ThirdPlaceSeries/RowGood.lean}{GlobalTheory.Fc.RL.RowGoodnessCondition}
\leanref{Development/AutomorphicArithmetic/ThetaCorrespondence/ThetaPairLifts/ThirdPlaceSeries/GeometricWeakTies.lean}{GlobalTheory.Fc.S0.GeometricWeakCompatibility'}
\leanref{Development/AutomorphicArithmetic/ThetaCorrespondence/ThetaPairLifts/ThirdPlaceSeries/RowGood.lean}{GlobalTheory.Fc.RL.RowMultiplicityCondition}
\leanref{Development/AutomorphicArithmetic/ThetaCorrespondence/ThetaPairLifts/ThirdPlaceSeries/RowGood.lean}{GlobalTheory.Fc.RL.RowKappaCondition}

\begin{lemma}\label{asm-dependence-linked}
Let $Q^{\mathrm b}$ and $Q$ be slots, let $\Lambda$ be link data, and let $r$ and $r'$ be
properties of slots (Definition~\ref{asm-dependence-parameters}(2)).
\begin{enumerate}
\item If $r$ is linked for $(Q^{\mathrm b},Q)$, and $r(S)$ implies $r'(S)$ for every slot $S$,
  then $r'$ is linked for $(Q^{\mathrm b},Q)$. If $r(S)$ holds for every slot $S$, then $r$ is
  linked for $(Q^{\mathrm b},Q)$. More generally, let $Q^{\mathrm b}_1\subset Q^{\mathrm b}$ and
  $Q_1\subset Q$ be slots, and suppose that $r(S)$ implies $r'(S_1)$ for all slots $S_1\subset S$;
  if $r$ is linked for $(Q^{\mathrm b},Q)$, then $r'$ is linked for $(Q^{\mathrm b}_1,Q_1)$.
\item If properties $H$ and $H'$ of classes and selections are linked for $(Q^{\mathrm b},Q)$ with
  $\Lambda$, then so is ``$H$ and $H'$''.
\item Let local Langlands maps be given, and let $N$ and $N'$ be extensions with $N'$ stronger
  than $N$. Then every core built from
  $p(N')$ is built from $p(N)$, every good core of $p(N')$ is a good core of $p(N)$, and every
  bundle that $\Lambda_{p(N')}$ links to a base point is linked to it by $\Lambda_{p(N)}$. Every
  property $H$ of classes and selections that is linked at $p(N)$ is linked at $p(N')$, and if the
  open parameter statements (Definition~\ref{asm-open-parameter-statements}) hold at $p(N)$, they
  hold at $p(N')$.
\end{enumerate}
\end{lemma}
\leanref{Development/AutomorphicArithmetic/ThetaCorrespondence/ThetaPairLifts/ThirdPlaceSeries/OpenAndLinkedOfStrongerExtension.lean}{GlobalTheory.Glob.open_and_linked_of_stronger_extension_alt}
\leanref{Development/AutomorphicArithmetic/ThetaCorrespondence/ThetaPairLifts/ThirdPlaceSeries/RowsOfStrongerNExt.lean}{GlobalTheory.HeightCounting.Up.rows_of_stronger_nExt_alt}
\leanref{Development/AbelianVarietyTheory/Proofs/IwasawaDecomposition/CoreSlotConstructionMonoFewerCoresBothPasses.lean}{GlobalTheory.HeightCounting.Up.eventually_coreProperty_mono_fewerCores}
\leanref{Development/AbelianVarietyTheory/Definitions/GlobalTwisting/EventuallyAllScaleContextsAndOfBoth.lean}{GlobalTheory.Fc.RowL.and}
\begin{proof}
First, a remark on Definition~\ref{is-slot-rows}. Let $Q_1\subset Q$ be slots and
$\Lambda_1\subset\Lambda$ link data. If $H$ holds along $(Q,\Lambda)$ with a condition $E$ on
bundles, then $H$ holds along $(Q_1,\Lambda_1)$ with the same $E$: every bundle that $\Lambda_1$
links to a base point of $Q_1$ is linked by $\Lambda$ to a base point of $Q$, so it eventually
satisfies $E$; and a selection whose bundles are built from $(Q_1,\Lambda_1)$ and satisfy $E$ is
built from $(Q,\Lambda)$.

(1) Keep the slot $R$ of $r$. If $r$ holds for every slot, take for $R$ the slot of all cores. In
the general case, every core in $Q^{\mathrm b}_1$ lies in $Q^{\mathrm b}$, so it satisfies $R$ at
large base data, eventually in the context; and $Q_1\cap R\subset Q\cap R$, so $r(Q\cap R)$ gives
$r'(Q_1\cap R)$.

(2) Let $R$ and $R'$ be the slots for $H$ and $H'$, and put $R''=R\cap R'$. Every core in
$Q^{\mathrm b}$ satisfies $R''$ at large base data, eventually in the context, since it satisfies
$R$ and $R'$ there. The slot $Q\cap R''$ is contained in $Q\cap R$ and in $Q\cap R'$, so by the
remark $H$ holds along $(Q\cap R'',\Lambda)$ with a condition $E$ and $H'$ with a condition $E'$.
Then ``$H$ and $H'$'' holds along it with the condition ``$E$ and $E'$'': a bundle linked to a
base point satisfies both eventually, and a selection whose bundles satisfy both satisfies the
hypotheses for $H$ and for $H'$.

(3) A core built from $p(N')$ is the core of a stage that satisfies $N'$ and the standard
conditions. That stage satisfies $N$, so the core is built from $p(N)$. A good core of $p(N')$ is
a core built from $p(N')$ with further conditions that do not involve the extension, so it is a
good core of $p(N)$. A bundle that $\Lambda_{p(N')}$ links to a base point is built from $p(N')$ at
that point with a witness; the same stage, read as a stage of $p(N)$, carries the bundle, and the
witness gives one for $p(N)$. Now let $H$ be linked at $p(N)$, with slot $R$. At large base data,
eventually in the context, every core built from $p(N')$ is built from $p(N)$ and so satisfies $R$.
The slot $Q_{p(N')}\cap R$ is contained in $Q_{p(N)}\cap R$ and $\Lambda_{p(N')}$ in
$\Lambda_{p(N)}$, so by the remark $H$ holds along $(Q_{p(N')}\cap R,\Lambda_{p(N')})$; this is
the general form of (1), with $Q^{\mathrm b}_{p(N')}\subset Q^{\mathrm b}_{p(N)}$ and
$Q_{p(N')}\subset Q_{p(N)}$. Each of the open parameter statements is linked at $p(N)$ in this
form, so it is linked at $p(N')$.
\end{proof}

\begin{lemma}\label{asm-dependence-howe-unramified}
Let $L$ be a field of characteristic $0$, complete for a discrete valuation $v$, with valuation ring
$\mathcal O$ and uniformizer $\varpi$. Let $\sigma$ be an involution of $L$ with $v\circ\sigma=v$
and $\sigma(\varpi)=\varpi$ such that $v(\sigma(a)-a)=0$ for some $a\in\mathcal O$, let $L_0$ be its
fixed field and $\mathcal O_0=\mathcal O\cap L_0$, and suppose that the residue field of $L$ has
$q^2$ elements for an odd integer $q$. Let $n,n'\ge0$ be integers, and let $J_S\in GL_n(\mathcal O)$
and $J_T\in GL_{n'}(\mathcal O)$ be Hermitian matrices: $\sigma(J_{ij})=J_{ji}$. For such a matrix $J$
of size $k$ let $U(J)$ be the group of $g\in GL_k(L)$ with ${}^t\sigma(g)Jg=J$, and
$U(J)(\mathcal O)=U(J)\cap GL_k(\mathcal O)$. Let $\delta\in\mathcal O^\times$ with
$\sigma(\delta)=-\delta$, and let $\psi$ be a character of the additive group of $L$ that is trivial
on $\mathcal O_0$ and not on $\varpi^{-1}\mathcal O_0$. On $X=M_{n'\times n}(L)$ put
\[
 H(x,y)=\delta\operatorname{tr}\big(J_S^{-1}\,{}^t\sigma(y)\,J_T\,x\big),\qquad
 \langle x,y\rangle=\tfrac12\big(H(x,y)+\sigma(H(x,y))\big).
\]
Let $M$ be a complex vector space, $\rho$ a map from $X$ to the linear endomorphisms of $M$, and
$\omega$ a homomorphism from $U(J_S)\times U(J_T)$ to the monoid of linear endomorphisms of $M$. Call
$s\in M$ a \emph{vacuum vector} if $\rho(u)s=s$ for every $u\in M_{n'\times n}(\mathcal O)$. Suppose
that:
\begin{enumerate}
\item $\rho(0)=1$, and $\rho(x)\rho(y)=\psi(\langle x,y\rangle)\rho(x+y)$ for all $x,y\in X$;
\item $M$ is spanned by the vectors $\rho(x)s$ with $x\in X$ and $s$ a vacuum vector;
\item $\omega(g,h)\rho(x)=\rho(hxg^{-1})\omega(g,h)$ for all $g\in U(J_S)$, $h\in U(J_T)$ and
  $x\in X$;
\item $\omega(g,h)s=s$ for every vacuum vector $s$, every $g\in U(J_S)(\mathcal O)$ and every
  $h\in U(J_T)(\mathcal O)$.
\end{enumerate}
Let $s$ be a nonzero vacuum vector such that the vectors $\rho(x)s$, $x\in X$, span $M$. Then every
vector of $M$ fixed by $\omega(g,1)$ for all $g\in U(J_S)(\mathcal O)$ lies in the span of the
vectors $\omega(1,h)s$ with $h\in U(J_T)$, and every vector fixed by $\omega(1,h)$ for all
$h\in U(J_T)(\mathcal O)$ lies in the span of the vectors $\omega(g,1)s$ with $g\in U(J_S)$.
\end{lemma}
\leanref{Development/AbelianVarietyTheory/PublishedTheorems/GaloisGeometry₂/HoweUnram.lean}{GlobalTheory.Hk3.HoweUnram}
\leanref{Development/AbelianVarietyTheory/Definitions/GaloisGeometry/NormalFam.lean}{GlobalTheory.Hk3.LocPair}
\leanref{Development/AbelianVarietyTheory/Definitions/GaloisGeometry/NormalFam.lean}{GlobalTheory.Hk3.LocVac}
\leanref{Development/AbelianVarietyTheory/Proofs/SatakeParameters/RhoMemOfMkEq.lean}{GlobalTheory.Hk3.howeUnram_holds}

Here $L$ is an unramified quadratic extension of $L_0$ of odd residual characteristic. The lattice
model of the Weil representation of the unramified dual pair $(U(J_S),U(J_T))$, with $\rho$ the
action of the Heisenberg group of $X$ and $s$ its standard vector, satisfies conditions (1) to (4).
For this model the lemma is Howe's result that the vectors fixed by the integral points of one
member of the pair are generated by the standard vector under the other member
\cite[Chapitre~5]{MVW87}.

\begin{theorem}\label{asm-dependence-zero-sums}
There is a constant $C>0$ with the following properties. Let $\chi$ be a primitive Dirichlet
character of conductor $N\ge1$, with $L(s,\chi)=\zeta(s)$ if $N=1$. Put
\[
 \psi_0(u,\chi)=\sum_{1\le n\le u}\Lambda(n)\chi(n)-\delta_N u,
\]
where $\Lambda$ is von Mangoldt's function and $\delta_N$ is $1$ if $N=1$ and $0$ otherwise. For a
zero $\rho$ of $L(s,\chi)$ let $m(\rho)$ be its multiplicity. For real $T$ let $Z(T)$ be the set of
zeros $\rho$ of $L(s,\chi)$ with $0<\operatorname{Re}\rho<1$ and $|\operatorname{Im}\rho|<T$, and
for real $t$ let $Z'(t)$ be the set of those with $0<\operatorname{Re}\rho<1$ and
$|\operatorname{Im}\rho-t|\le1$.
\begin{enumerate}
\item (Explicit formula.) For real $u$ and $T$ with $2\le T\le u$, the set $Z(T)$ is finite and
  \[
   \Big|\psi_0(u,\chi)+\sum_{\rho\in Z(T)}m(\rho)\frac{u^\rho}{\rho}\Big|
   \le C\Big(\frac{u\log^2(NuT)}{T}+\log(Nu)+\sum_{\rho\in Z(1)}\frac{m(\rho)}{|\rho|}\Big).
  \]
\item (Zeros in windows.) For every real $t$, the set $Z'(t)$ is finite and
  $\sum_{\rho\in Z'(t)}m(\rho)\le C\log\big(N(|t|+2)\big)$.
\end{enumerate}
\end{theorem}
\leanref{Development/AbelianVarietyTheory/Proofs/GlobalInterpolation/TruncatedExplicitFormula.lean}{GlobalTheory.ExplicitFormula.truncated_explicit_formula}
\leanref{Development/AbelianVarietyTheory/Proofs/GlobalInterpolation/FinsumMemLeOfSubset.lean}{GlobalTheory.ZeroWindow.count_zeros_window}
\leanref{Development/AbelianVarietyTheory/Proofs/GlobalInterpolation/ExplicitFormulaAndWindowCount.lean}{GlobalTheory.ZeroBounds.explicit_formula_and_window_count}

Part (1) is the explicit formula for $\psi(u,\chi)$ truncated at height $T$, with the terms of the
zeros of height less than $1$ kept in the error, and part (2) is the bound for the number of zeros
of $L(s,\chi)$ in a window of height $2$ \cite[Chapters 15 to 17 and 19]{Dav00}.

\begin{lemma}\label{asm-dependence-genus-mass}
Suppose that for every definite-side datum (Definition~\ref{th-definite-datum}), with host field
$E$, every Hermitian space over $E$ with correct signature (Definition~\ref{th-spaces}) that is
positive definite at every complex embedding of $E$ has an orthonormal basis. (The proof does not use
this hypothesis.) Let $D$ be a definite-side datum with host fields $F\subset E$, and let $W$ be a
Hermitian space over $E$ with correct signature, of rank $M\ge0$, positive definite at every complex
embedding, with standard lattice $L_W$. Then $\operatorname{mass}(L_W)=\operatorname{mass}(I_M)$
(Definition~\ref{cv-mass}). In particular the mass formula for genera of rank $2$ holds for $D$
(Definition~\ref{asm-facts-volume-tie}(3)).
\end{lemma}
\leanref{Development/AbelianVarietyTheory/Proofs/FoundationProofs/GenusMassStdOfLandherrHasse.lean}{GlobalTheory.genusMassStd_of_landherrHasse}
\leanref{Development/AbelianVarietyTheory/Proofs/FoundationProofs/GenusMassStdOfLandherrHasse.lean}{GlobalTheory.genusMassStd}
\leanref{Development/AbelianVarietyTheory/Proofs/FoundationProofs/GenusMassStdAll.lean}{GlobalTheory.genusMassStd_all}
\leanref{Formalism/LocalToGlobal/GlobalTheory/GL5/DiscreteSeriesOfFormalDegree.lean}{Formalism.L2.GlobalTheoryTheory.GenusMassVolumes.genusMassAt_of}
\begin{proof}
By Proposition~\ref{hf-host-field-properties}, $E$ is a CM field with totally real subfield $F$, and
no finite prime of $F$ ramifies in $E$. If $M\ge1$, Lemma~\ref{cv-volume-mass-definite}(1) gives
$\operatorname{mass}(L_W)=\operatorname{mass}(I_M)$. If $M=0$, both lattices are zero and their masses
agree. The case $M=2$ is the mass formula for genera of rank $2$.
\end{proof}

\begin{lemma}\label{asm-dependence-projections-onto}
Let $B$ be a space of forms of a family of definite forms over a number field $L$ for a finite group
$G$ (Definition~\ref{is-realisation}), the row forms at a tuple or the lift forms at a lift and a
tuple, with its class-sum pairing and its action of $G$, and let $v_0$ be the base place of the
family. Let $X$ be a finite set of functions $G\to L$, let $P\subset G$ be a subgroup with a
character $\chi\colon P\to L^\times$, let $e_X$ and $e_{P,\chi}$ be as in
Lemma~\ref{is-baselinks-group-algebra}(4), and let $Z\subset B$ be a subspace. The image of $e_{P,\chi}$
on $B$ is the space of \emph{level forms}, the $f$ with $hf=\chi(h)f$ for $h\in P$; let $\Gamma$ be the
image of $e_Xe_{P,\chi}$. Let $f\colon L^s\to B$ and $c\colon L^d\to B$ be the coordinate maps of
bases of the level forms and of $Z\cap\Gamma$. Let $\pi\colon L^s\to L^d$ be the linear map such that
$c(\pi(x))$ is the orthogonal projection onto $Z\cap\Gamma$, for the class-sum pairing restricted to
$\Gamma$, of the orthogonal projection of $f(x)$ onto $\Gamma$. Suppose that the involution of $L$ is
complex conjugation at $v_0$, that the functions in $X$ are the characters of absolutely irreducible
representations of $G$, and that $B$ satisfies the conditions that the family conditions impose on its
forms (Lemma~\ref{is-eigenvalue-base-inputs}(1)). Then $\Gamma$ consists of level forms, and if
$\lambda(y)\in L^s$ denotes the coordinates of $c(y)$ in the basis of the level forms, then
$\pi(\lambda(y))=y$ for every $y\in L^d$. In particular $\pi$ is surjective.
\end{lemma}
\leanref{Formalism/LocalToGlobal/Divisibility/LiftLatLemmas.lean}{Formalism.L2.DIV.BlockDat.lf_lPr_surj}
\leanref{Formalism/LocalToGlobal/Divisibility/LiftLatLemmas.lean}{Formalism.L2.DIV.BlockDat.lf_lPr_lE}
\leanref{Formalism/LocalToGlobal/Divisibility/LiftLatLemmas.lean}{Formalism.L2.DIV.BlockDat.lf_lvlE_lE}
\leanref{Formalism/LocalToGlobal/Divisibility/LiftLatLemmas.lean}{Formalism.L2.DIV.BlockDat.lf_lvl_vE}
\begin{proof}
Work in the coordinates of a basis of $B$, and let $H$ be the matrix of the class-sum pairing. Let
$E_\ell$ and $V$ be the matrices of the bases of the level forms and of $\Gamma$, so that $f$ is
$x\mapsto E_\ell x$, and let $E_Z$ be the matrix that expresses the basis of $Z\cap\Gamma$ through
that of $\Gamma$, so that $c$ is $y\mapsto VE_Zy$. For a matrix $A$ with independent columns write
$\Pi(A,H)=(A^*HA)^{-1}A^*H$; then $A\,\Pi(A,H)$ is the orthogonal projection onto the span of the
columns of $A$ for $H$, and $\Pi(A,H)A=1$. Put $H_\Gamma=|P|^{-1}V^*HV$, the restricted pairing in
the coordinates of $\Gamma$. Then
\[
 \pi=\Pi(E_Z,H_\Gamma)\,\Pi(V,H)\,E_\ell,\qquad \lambda=\Pi(E_\ell,H)\,VE_Z.
\]

By the family conditions, $H$ is positive definite at $v_0$ and $G$ acts by isometries of the
pairing. The values of $\chi$ are roots of unity, and the involution is complex conjugation at $v_0$,
so $e_{P,\chi}$ is an idempotent that is self-adjoint for the pairing. As the functions in $X$ are
the characters of absolutely irreducible representations, $e_X$ is a central idempotent, so
$e_Xe_{P,\chi}$ is also a self-adjoint idempotent. Hence $E_\ell\Pi(E_\ell,H)$ and $V\Pi(V,H)$ are the
matrices of $e_{P,\chi}$ and of $e_Xe_{P,\chi}$.

Since $e_X$ is central and $e_{P,\chi}$ is idempotent, $e_{P,\chi}e_Xe_{P,\chi}=e_Xe_{P,\chi}$. So
$e_{P,\chi}$ fixes $\Gamma$, that is, $E_\ell\Pi(E_\ell,H)V=V$, and $\Gamma$ consists of level forms.
Multiplying by $E_Z$ gives $E_\ell\lambda=VE_Z$, that is, $f(\lambda(y))=c(y)$.

The pairing $H$ is positive definite at $v_0$ and $V$ has independent columns, so $H_\Gamma$ is
positive definite at $v_0$ and $\Pi(E_Z,H_\Gamma)$ is defined. Therefore
\[
 \pi\lambda=\Pi(E_Z,H_\Gamma)\,\Pi(V,H)\,E_\ell\lambda=\Pi(E_Z,H_\Gamma)\,\Pi(V,H)\,VE_Z
 =\Pi(E_Z,H_\Gamma)\,E_Z=1.
\]
\end{proof}

\begin{lemma}\label{asm-dependence-base-vectors}
Let $D$ be a definite-side datum (Definition~\ref{th-definite-datum}) with coefficient field $L$,
embedding $\sigma_0$, top forms $\mathcal E$ and base tuple $\mu_0$. For a tuple $\mu$ let $e_\mu$ be the
idempotent of the rows data of $D$ on the ambient space $\mathcal A_\mu$, whose image is the base space
$V_0(\mu)$, and $\iota_\mu\colon\mathcal A_\mu\to\mathcal E$ the map of the rows data. Let
$P$ be a row-side family for $D$ (Definition~\ref{is-row-side-family}), with base tuple $t_0$ of its
family of definite forms, characters $\mathcal X$ of that family and $e_{\mathcal X}$ as in
Lemma~\ref{is-baselinks-group-algebra}(4), subgroups $P_0,P_s$ with characters $\psi_0,\psi_s$, and
columns $S_1$ on the definite side of $(P_s,\psi_s)$, with matrices $P_t$ of that side and $\iota_t$
of $S_1$. Write $\Gamma$ for the image of $e_{\mathcal X}e_{P_s,\psi_s}$ on the row forms at $t_0$,
$E$ for the matrix of $e_{\mathcal X}$ on the coordinates of the level forms of the row forms at $t_0$
for $(P_s,\psi_s)$ (Lemma~\ref{asm-dependence-projections-onto}), and $V$ for the matrix that
expresses a basis of $\Gamma$ through the basis of those level forms. The identifications of $P$ are
tied to the data of $D$; among the ties are:
\begin{enumerate}
\item the tuple $\mu$ of Hecke characters of $D$ attached to $t_0$ is $\mu_0$;
\item a $\sigma_0$-semilinear map $j_A$ from the coordinates of the level forms at $t_0$ to
  $\mathcal A_\mu$ with $e_\mu(j_A(x))=j_A(Ex)$ for every $x$;
\item for every $F$,
  $\iota_\mu(j_A(VF))=j_E\big(e_{RR}(\iota_{t_0}P_{t_0}F)\big)$, where $j_E$ is the
  $\sigma_0$-semilinear map from the space of the glued lattice of the sharp row side of $D$ to
  $\mathcal E$, which carries the basis of that lattice to a basis of $\mathcal E$.
\end{enumerate}
Let $Q$ be an indefinite realization of $P$ (Definition~\ref{is-realisation-join}), with realization $R$
and face $F_Q$. Suppose that the family of definite forms of $P$ satisfies the family conditions with
the conditions on its inert indices (Lemma~\ref{is-eigenvalue-base-inputs}(1)) and has lift-distinct
types, and that $R$ has primitive rows at $F_Q$ (Definition~\ref{is-primitive-rows}). Then
$\iota_{\mu_0}$ maps $V_0(\mu_0)$ onto $\mathcal E$.
\end{lemma}
\leanref{Development/AutomorphicArithmetic/ThetaCorrespondence/ThetaPairLifts/ThirdPlaceSeries/V0MapIotaEqTopOfReal.lean}{GlobalTheory.V0_map_iota_eq_top_of_real_alt}
\leanref{Development/AbelianVarietyTheory/Proofs/FlagMeasures/RowMapSurjectiveOfPrimitiveRows.lean}{GlobalTheory.rowMap_surjective_of_primitiveRows_alt}
\leanref{Development/AbelianVarietyTheory/Proofs/FlagMeasures/PrimitiveRowsMapBijectiveColumnLevel.lean}{GlobalTheory.prim_rows_surjective}
\leanref{Development/AbelianVarietyTheory/Proofs/FlagMeasures/V0MapIotaEqTopDefinite.lean}{GlobalTheory.V0_map_iota_eq_top_of_definite}
\leanref{Development/AbelianVarietyTheory/Proofs/FlagMeasures/RowSideIotaMapsFixedColumnsOntoOfMatrices.lean}{GlobalTheory.rowSide_iota_maps_fixedColumns_onto_of_matrices}
\leanref{Formalism/LocalToGlobal/GroupLattices/JointServices/LevelOperatorFixesImage.lean}{Formalism.L2.G.ServJoint.level_operator_fixes_image}
\begin{proof}
The realization $R$ satisfies the weak characterization at $F_Q$, and its complex data at $F_Q$ have
the definite subspaces and the columns $S_1$ of $P$; the level $1$ of $F_Q$ has the subgroup $P_s$ and
the character $\psi_s$. By Lemma~\ref{asm-primitive-rows-base}(2), at the level $1$, the matrix
$\iota_{t_0}$ of $S_1$ is bijective. By Lemma~\ref{asm-primitive-rows-parts}(1), for $(P_s,\psi_s)$,
the matrix $P_{t_0}$ is surjective.

Let $\mu$ be the tuple of $t_0$. The vectors $j_E(b)$, for $b$ in the basis of the glued lattice,
form a basis of $\mathcal E$, so it suffices that each of them lies in $\iota_\mu(V_0(\mu))$. Choose $y$ with $\iota_{t_0}y=e_{RR}^{-1}(b)$, and then $F$ with $P_{t_0}F=y$.
By (3), $\iota_\mu(j_A(VF))=j_E(b)$.

It remains that $j_A(VF)\in V_0(\mu)$. The vector $VF$ is the coordinate vector of a form in
$\Gamma$. By the family conditions, Lemma~\ref{asm-dependence-projections-onto} applies to the row
forms at $t_0$, and $\Gamma$ consists of level forms; and $e_{\mathcal X}$ is an idempotent whose image
contains $\Gamma$. So $EVF=VF$, and by (2), $e_\mu(j_A(VF))=j_A(EVF)=j_A(VF)$, which lies in
$V_0(\mu)$. Hence $\iota_\mu$ maps $V_0(\mu)$ onto $\mathcal E$, and $\mu=\mu_0$ by (1).
\end{proof}

\begin{lemma}\label{asm-dependence-complement-norms}
Let $\mathcal F$ be a family (Definition~\ref{pd-family}), and let $D$ be a member at a scale
$x>x_1$, with level primes $g_0,g_1$, the elements $\pi_0,\pi_1$ of its error data and the notation of
Definition~\ref{asm-division-kernel-errors}. Then conditions (a) and (b) of
Definition~\ref{asm-division-kernel-errors}(2) hold at $D$ for every word $u$, every lift $j$ and
$i\in\{0,1\}$: $g_i\ne T_u$, and
$\sum_w\mathrm{sz}_w\big(\nu^{(i)}_{u,j}(w)\big)=\big(\sum_\tau n_{u,j}(\tau)\big)\log g_i$, the sum over
the places $w$ of $L$ over $g_i$.
\end{lemma}
\leanref{Formalism/LocalRingCardinality/IntCastDvdOfDvdIntegers.lean}{Formalism.BasicNotions.D7.T_ne_g}
\leanref{Development/AbelianVarietyTheory/Proofs/DivisorProofs/TypePrimeValuationSumEqLogGapFamily.lean}{GlobalTheory.TypeValNN.typePrime_valuationSum_eq_logGap_all}
\leanref{Development/AbelianVarietyTheory/Proofs/DivisorProofs/MemHeightOneIdealIffAdditiveValuationPos.lean}{GlobalTheory.TypeLoss.sum_half}
\begin{proof}
(a) By the conditions at the auxiliary primes (Definition~\ref{th-definite-datum}), $T_u=\pi\bar\pi$
with $\pi\in\mathcal O_E$ and $\pi\equiv1$ modulo $M=Q_Sqd(g_0g_1)^2$. Then
$T_u-1=(\pi-1)(\bar\pi-1)+(\pi-1)+(\bar\pi-1)$ lies in $M\mathcal O_E\cap\mathbf Z=M\mathbf Z$, so
$T_u\equiv1$ modulo $g_i$. As $g_i>1$, $T_u\ne g_i$.

(b) For $\tau\in\mathcal T$ put $y_\tau=\tau(\overline{\pi_i})\in L$. Since $\tau$ carries complex
conjugation to the involution of $L$, $y_\tau\overline{y_\tau}=\tau(\pi_i\overline{\pi_i})=g_i$. Let
$w$ run over the places of $L$ over $g_i$. Then $v_w(y_\tau)+v_w(\overline{y_\tau})=v_w(g_i)=e_w$, and
$\sum_w\ell_we_w=\log g_i$, since $\sum_we_wf_w=[L:\mathbf Q]$ for the residue degrees $f_w$. The
involution of $L$ permutes the places over $g_i$, preserves $\ell_w$, and carries $v_w(y_\tau)$ to
$v_{\bar w}(\overline{y_\tau})$; so $\sum_w\ell_wv_w(\overline{y_\tau})=\sum_w\ell_wv_w(y_\tau)$, and
\[
 \sum_w2\ell_wv_w(y_\tau)=\sum_w\ell_w\big(v_w(y_\tau)+v_w(\overline{y_\tau})\big)=\log g_i.
\]
By the definitions, $\mathrm{sz}_w(\nu^{(i)}_{u,j}(w))=2\ell_we_w\nu^{(i)}_{u,j}(w)
=\sum_\tau n_{u,j}(\tau)\,2\ell_wv_w(y_\tau)$. Summing over $w$ and exchanging the sums gives
$\sum_\tau n_{u,j}(\tau)\log g_i$.
\end{proof}

\begin{lemma}\label{asm-dependence-findim}
Let $E$ be a CM field with totally real subfield $F$, let $J\in GL_n(E)$ be a Hermitian matrix, and
let $L\subset E^n$ be an $\mathcal O_E$-lattice that is self-dual for the Hermitian form of $J$. Let
$G=U(J)$ over $F$, let $K\subset G(\mathbf A_{F,f})$ be the stabilizer of the closure of $L$, and give
$G_f=G(\mathbf A_{F,f})$ the topology for which the subgroups $\{k\in K:(k-1)L\subset gL\}$, $g\ge1$,
form a basis of neighbourhoods of $1$. For every real place $v$ of $F$ let an embedding of $E$ over $v$
and a matrix $P_v$ identify $G(F_v)$ with a group $U(p_v,q_v)$, let $\mu_v$ be a Haar measure on
$U(p_v,q_v)$, and put $G_\infty=\prod_vU(p_v,q_v)$, with $\mu_\infty=\prod_v\mu_v$, and
$G_{\mathbf A}=G_\infty\times G_f$, in which $G(F)$ is embedded diagonally. Let $\mu_f$ be a Haar
measure on $G_f$, let $\mathcal F\subset G_{\mathbf A}$ be a measurable fundamental domain for $G(F)$:
every $G(F)$-orbit meets $\mathcal F$, and $\gamma\mathcal F\cap\mathcal F$ is a null set for
$\mu_\infty\times\mu_f$ for every $\gamma\ne1$ in $G(F)$; and let $R$ be the right regular
representation of $G_{\mathbf A}$ on $L^2(\mathcal F)$, a function on $\mathcal F$ being extended to a
$G(F)$-invariant function on $G_{\mathbf A}$. Let $e\colon G_f\to\mathbf C$ be locally constant and
supported in $K$, and let $R(e)$ be the operator with
$\langle x,R(e)y\rangle=\int e(g)\langle x,R(1,g)y\rangle\,d\mu_f(g)$. For every real place $v$, let
$\sigma_v$ be an irreducible unitary representation of $U(p_v,q_v)$ on a Hilbert space $H_v$ that is
square-integrable with formal degree $d_v$ for $\mu_v$, and $w_v\in H_v$ a unit vector whose matrix
coefficient $g\mapsto\langle w_v,\sigma_v(g)w_v\rangle$ is integrable. Let $T$ be the operator with
\[
 \langle x,Ty\rangle=\int_{G_\infty}\Big(\prod_vd_v\Big)\,
 \overline{\prod_v\langle w_v,\sigma_v(g_v)w_v\rangle}\;\langle x,R(g,1)y\rangle\,d\mu_\infty(g).
\]
Then the image of $R(e)T$ is finite-dimensional.
\end{lemma}
\leanref{Formalism/LocalToGlobal/SiegelLatticeTheory/HaarTransversals/RowsDef.lean}{Formalism.L2.SG.RowsDef.RowsFinDim}
\leanref{Formalism/SharpRestriction/CuspSpanNormBounds/GlobAut.lean}{Formalism.MainNotions.S.AutInst}
\leanref{CompanionTheory/HarishChandraModules/KFiniteAnalyticVectors/FundamentalDomainCusps/AutomorphicFormsFiniteDimensional.lean}{CompanionTheory.D.gk.RF.automorphic_forms_finite_dimensional}

This is the finiteness of the multiplicity of the representation $\bigotimes_v\sigma_v$ of $G_\infty$
in the square-integrable functions on $G(F)\backslash G(\mathbf A_F)$ that are invariant under an open
compact subgroup of $G_f$, which follows from Harish-Chandra's finiteness theorem \cite{HC68,BJ79}. By
Schur orthogonality the image of $T$ consists of the images of $\bigotimes_vw_v$ under the
$G_\infty$-equivariant maps from $\bigotimes_v\sigma_v$ to $L^2(\mathcal F)$; the operator $R(e)$
commutes with $G_\infty$, and its image is invariant under an open compact subgroup of $G_f$, as $e$ is
locally constant with compact support.

\subsection{Row models and Bessel periods}\label{subsection-asm-row-models}

\begin{definition}\label{asm-dependence-rows-parameters}
Let $D$ be a definite-side datum (Definition~\ref{th-definite-datum}) with rank datum $(k,m,t)$,
$r=k+m-1$ (Definition~\ref{th-ranks}) and set $\mathcal T$ of embeddings, and for
$\tau\in\mathcal T$ let $v(\tau)$, $b(\tau)$ and $e(\tau)$ be the $v$-, $b$- and $e$-parts of its
Harish-Chandra parameter at $\tau$ (Definition~\ref{th-hc-parameter}).
\begin{enumerate}
\item A \emph{signed set} is a finite set $S$ of pairs $(x,\sigma)$ of a rational number $x$ and a
  sign $\sigma\in\{+,-\}$. It is \emph{regular} if no number occurs in it with both signs, and
  $n_-(S)$ is the number of its pairs with sign $-$. For a rational $h$, $S+h$ is the set of the
  pairs $(x+h,\sigma)$ and $-S$ that of the pairs $(-x,\sigma)$. An \emph{enumeration} of $S$
  consists of real numbers $\lambda_1>\dots>\lambda_n$ and signs $\epsilon_1,\dots,\epsilon_n$ such
  that $S$ is the set of the pairs $(\lambda_a,\epsilon_a)$; it is \emph{integral} if
  $\lambda_a-(n-1)/2\in\mathbf Z$ for every $a$.
\item For a rational $y$ and an integer $L\ge0$ put $B(y,L)=\{y+(L-1)/2-a:0\le a<L\}$, and put
  $s=r+6$. For a real number $c$, let the integers $\ell$ and $\gamma(c)$, the numbers $a_l(\tau)$ for
  $0\le l\le m$ and the centres $c_l(\tau)$ for $1\le l\le m$ be as in Definition~\ref{th-grid-centres}.
\item For $0\le i\le m$, the \emph{core parameter} $S^{\mathrm{core}}_i(\tau)$ of $D$ and $c$ is the
  signed set of the following pairs: $(a_l(\tau),-)$ for $0\le l<m-i$ with $m-i-l$ odd, and
  $(a_l(\tau),+)$ for the other $l$ with $0\le l\le m$; $(v_j(\tau)-\tfrac12,+)$ for $1\le j\le k-1$;
  and $(x,+)$ for $x\in B(c_l(\tau),s-1)$ and $m-i+1\le l\le m$.
\item For $0\le i,j\le m$, let $P_{i,j}(\tau)$ be the signed set of the following pairs:
  $(v_l(\tau),+)$ for $1\le l\le k-1$; $(-b_l(\tau),-)$ for $j+1<l\le m$ with $l-j-1$ odd, and
  $(-b_l(\tau),+)$ for the other $l$ with $1\le l\le m$; $(x,+)$ for $x\in B(-e_l(\tau),t)$ and
  $1\le l\le j$; and $(x,+)$ for $x\in B(c_l(\tau),s)$ and $m-i+1\le l\le m$. The \emph{row
  parameter} $S^{\mathrm{row}}_i(\tau)$ of $D$ and $c$ is $-P_{i,0}(\tau)$ if $i$ is even and
  $-(P_{i,0}(\tau)-\tfrac12)$ if $i$ is odd.
\end{enumerate}
\end{definition}
\leanref{Development/AbelianVarietyTheory/Definitions/Foundations/StrictInterlaceIffDescEnum.lean}{GlobalTheory.SignedHC.IsRegular}
\leanref{Development/AbelianVarietyTheory/Definitions/Foundations/StrictInterlaceIffDescEnum.lean}{GlobalTheory.SignedHC.nMinus}
\leanref{Development/AbelianVarietyTheory/Definitions/Foundations/StrictInterlaceIffDescEnum.lean}{GlobalTheory.SignedHC.shift}
\leanref{Formalism/LocalToGlobal/IndefiniteReal/ScalarOnPartsFull.lean}{Formalism.IndefR.Mult.dualPar}
\leanref{Formalism/LocalToGlobal/BesselHermitianTheta/DefsReal.lean}{Formalism.L2.Theta.IsEnum}
\leanref{Formalism/FormalDegreeBessel/StatementDS.lean}{Formalism.QuantumDeterminants.IsHCParam}
\leanref{Formalism/LocalRingCardinality/DataTheta.lean}{Formalism.BasicNotions.TH.TPar.block}
\leanref{Development/AbelianVarietyTheory/Definitions/Foundations/StrictInterlaceIffDescEnum.lean}{GlobalTheory.SignedHC.GridData}
\leanref{Development/AbelianVarietyTheory/Definitions/Foundations/StrictInterlaceIffDescEnum.lean}{GlobalTheory.SignedHC.GridData.ofHCParam}
\leanref{Development/AbelianVarietyTheory/Definitions/Foundations/StrictInterlaceIffDescEnum.lean}{GlobalTheory.SignedHC.GridData.parPi}
\leanref{Development/AbelianVarietyTheory/Definitions/Foundations/StrictInterlaceIffDescEnum.lean}{GlobalTheory.SignedHC.GridData.parD}
\leanref{Development/AbelianVarietyTheory/Definitions/Divisors/LeastSplitAuxiliaryPrimeThetaChainSpec.lean}{GlobalTheory.ChiGrid.gap}
\leanref{Development/AbelianVarietyTheory/Definitions/Divisors/LeastSplitAuxiliaryPrimeThetaChainSpec.lean}{GlobalTheory.ChiGrid.cen}
\leanref{Development/AbelianVarietyTheory/Definitions/Divisors/LeastSplitAuxiliaryPrimeThetaChainSpec.lean}{GlobalTheory.ChiGrid.num}
\leanref{Development/AbelianVarietyTheory/Definitions/Divisors/LeastSplitAuxiliaryPrimeThetaChainSpec.lean}{GlobalTheory.ChiGrid.parV}

\begin{definition}\label{asm-dependence-rows-models}
Let $D$ be a definite-side datum (Definition~\ref{th-definite-datum}) with host field $E/F$ and set
$\mathcal T$ of embeddings, and let $V$ be a Hermitian space over $E$ of rank $n$
(Definition~\ref{th-spaces}). For $\tau\in\mathcal T$ let $G_V(\tau)$ be the image under $\tau$ of
the Gram matrix of $V$, and $U(G_V(\tau))$ the group of $g\in GL_n(\mathbf C)$ with
$g^*G_V(\tau)g=G_V(\tau)$. An archimedean family $\theta$, such as that of a grid Weil datum, embeds
$U(G_V(\tau))$ in $U(V)(\mathbf A)$ at the place $\tau$.
\begin{enumerate}
\item Let $S=(S_\tau)_{\tau\in\mathcal T}$ be finite sets of pairs of a rational number and a sign.
  A \emph{model of discrete series} for $U(V)$ with parameter $S$ consists, for each
  $\tau\in\mathcal T$, of a Hilbert space $H_\tau$, a strictly decreasing
  $\lambda_\tau\in\mathbf Q^n$ and signs $\epsilon_\tau\in\{\pm1\}^n$ with
  $S_\tau=\{(\lambda_{\tau,a},\epsilon_{\tau,a}):1\le a\le n\}$, an irreducible unitary
  square-integrable representation $\sigma_\tau$ on $H_\tau$ of the unitary group of the diagonal
  matrix $J_\tau$ with entries $\epsilon_\tau$, whose lowest $K$-type has the highest weight
  attached to $(\lambda_\tau,\epsilon_\tau)$, and an invertible matrix $f_\tau$ with $f_\tau^*J_\tau
  f_\tau=G_V(\tau)$. The group $U(G_V(\tau))$ acts on $H_\tau$ by $g\mapsto\sigma_\tau(f_\tau
  gf_\tau^{-1})$, and $f_\tau^{-1}J_\tau f_\tau$ is the \emph{Cartan involution} of the model at
  $\tau$. Write $M^\infty=\bigotimes_\tau H_\tau^\infty$ for the tensor product of the spaces of
  smooth vectors, and $M^K\subset M^\infty$ for the span of the tensors of $K$-finite vectors.
\item Let $M$ be such a model, $\theta$ an archimedean family, $J\subset U(V)(\mathbf A)$ and $\chi$
  a function on $U(V)(\mathbf A)$. A \emph{multiplicity map} of level $(J,\chi)$ is a map $T$ from
  $\prod_\tau H_\tau^\infty$ to functions on $U(V)(\mathbf A)$, linear in each variable, such that
  for every $v$: $T(v)$ is a cuspidal automorphic form, smooth at the archimedean places, finite
  under the centres of the enveloping algebras and of moderate growth; $T(v)(gk)=\chi(k)T(v)(g)$ for
  $k\in J$; replacing $v_\tau$ by its translate under $g\in U(G_V(\tau))$ replaces $T(v)$ by its
  right translate by the image of $g$ under $\theta$; and for every point $x$ and every $\tau$ the
  functional $v_\tau\mapsto T(v)(x)$ is continuous for a Sobolev norm of the model. The
  \emph{multiplicity space} $\operatorname{Mult}_\theta(M;J,\chi)$ is the span of the multiplicity
  maps, each read as a linear map on $M^\infty$, and
  $\operatorname{Mult}_\theta(M)=\operatorname{Mult}_\theta(M;\emptyset,1)$.
\item A \emph{row setting} consists of a definite-side datum $D$ in which $p_{de}$ and $p'_{de}$ are
  prime to $Q_Sq$, a level system for $D$ with base tuple, a theta grid for it with constants
  $(c,C_F,b)$ and spaces $V_{i,j}$, $P_i$, $Q_j$ (Definition~\ref{th-grid}), a grid Weil datum with
  archimedean family $\theta$, and a row volume (Definition~\ref{asm-grid-evaluation}(3)). For $0\le
  i\le m$, $D$ and the first constant $c$ of the grid determine sets $S^{\mathrm{row}}_i$ and
  $S^{\mathrm{core}}_i$ as in (1), the \emph{row parameter} of $V_{i,0}$ and the \emph{core
  parameter} of $P_i$ (Definition~\ref{asm-dependence-rows-parameters});
  $S^{\mathrm{row}}_0$ does not depend on $c$, and it is the parameter of the
  models of $U(Y_0)$ of the realizations of $D$ (Definition~\ref{is-realisation}(3)). A \emph{row
  model} at $i$ is a model of discrete series for $U(V_{i,0})$ with parameter $S^{\mathrm{row}}_i$
  whose Cartan involutions are those of the grid Weil datum, and a \emph{core model} at $i$ is one
  for $U(P_i)$ with parameter $S^{\mathrm{core}}_i$ and the same property.
\item In a row setting, let $K_{\mathrm{row}}$ be the principal congruence subgroup modulo $g_0g_1$
  of the stabilizer of the lattice of $Y_0$ of the level system, and $K_{P_0}$ the group of
  Definition~\ref{asm-grid-evaluation}(1). For a model $M_0$ of $U(Y_0)$ with parameter
  $S^{\mathrm{row}}_0$, the \emph{row multiplicities} are the elements of
  $\operatorname{Mult}_\theta(M_0;K_{\mathrm{row}},1)$; for a model $M^P$ of $U(P_0)$ with parameter
  $S^{\mathrm{core}}_0$, the \emph{core multiplicities} are those of
  $\operatorname{Mult}_\theta(M^P;K_{P_0},1)$. The \emph{core period} of functions $F$ on
  $U(P_0)(\mathbf A)$ and $u$ on $U(Y_0)(\mathbf A)$ is
  \[
   \operatorname{per}(F,u)=\int_{\mathcal F_{Y_0}}F(\iota_0h)\,u(h)\,dh,
  \]
  with $\iota_0\colon U(Y_0)\to U(P_0)$ the embedding of the grid and the measure and fundamental
  domain of the row volume. For Haar measures and fundamental domains on the groups
  $U(P_i)(\mathbf A)$, a function $F$ on $U(P_0)(\mathbf A)$ is \emph{reachable} if there are a
  column chain of core forms $(\Psi_0,\ldots,\Psi_m)$ (Definition~\ref{th-core-chain}(2)) with
  $\Psi_0=F$, a column $\Psi$ and a linear functional $\xi$ on the coefficients
  (Definition~\ref{asm-grid-evaluation}(5)) with $\Psi_m=\varphi_{\xi,\Psi}$.
\item The datum $D$ has the \emph{gap sizes} if $r>1800$ and $b_{j+1}(\tau)-b_j(\tau)>2N^3$ for
  every $\tau\in\mathcal T$ and $1\le j<m$, where $N=k+m(t+r+7)$ and the $b_j(\tau)$ are the
  $b$-parts of the Harish-Chandra parameters of $D$ (Definition~\ref{th-hc-parameter}). A \emph{gap
  row setting} is a row setting in which $D$ has the gap sizes and the first constant $c$ of the
  grid is $1$.
\end{enumerate}
\end{definition}
\leanref{Formalism/LocalToGlobal/IndefiniteReal/ScalarOnPartsFull.lean}{Formalism.IndefR.Mult.DSModel}
\leanref{Formalism/LocalToGlobal/IndefiniteReal/ScalarOnPartsFull.lean}{Formalism.IndefR.Mult.IsMultMap}
\leanref{Formalism/LocalToGlobal/IndefiniteReal/ScalarOnPartsFull.lean}{Formalism.IndefR.Mult.Mult}
\leanref{Formalism/LocalToGlobal/IndefiniteReal/ScalarOnPartsFull.lean}{Formalism.IndefR.Mult.MultT}
\leanref{Development/AbelianVarietyTheory/Definitions/TwistedInnerProducts/OffQOfSizeRows.lean}{GlobalTheory.SepRows.OffQ}
\leanref{Development/AbelianVarietyTheory/Definitions/TwistedInnerProducts/OffQOfSizeRows.lean}{GlobalTheory.SepRows.All}
\leanref{Development/AbelianVarietyTheory/Definitions/TwistedInnerProducts/OffQOfSizeRows.lean}{GlobalTheory.SepRows.Gap}
\leanref{Development/AbelianVarietyTheory/Definitions/TwistedInnerProducts/ParRowZero.lean}{GlobalTheory.SepRows.parRow}
\leanref{Development/AbelianVarietyTheory/Definitions/Divisors/BasePlaceNotMemChiGridBadPlaces.lean}{GlobalTheory.ChiGrid.Node.par}
\leanref{Development/AbelianVarietyTheory/Definitions/TwistedInnerProducts/OffQOfSizeRows.lean}{GlobalTheory.SepRows.levRow}
\leanref{Development/AbelianVarietyTheory/Definitions/TwistedInnerProducts/OffQOfSizeRows.lean}{GlobalTheory.SepRows.Rows}
\leanref{Development/AbelianVarietyTheory/Definitions/TwistedInnerProducts/OffQOfSizeRows.lean}{GlobalTheory.SepRows.CoreForms}
\leanref{Development/AbelianVarietyTheory/Definitions/TwistedInnerProducts/OffQOfSizeRows.lean}{GlobalTheory.SepRows.KF}
\leanref{Development/AbelianVarietyTheory/Definitions/TwistedInnerProducts/OffQOfSizeRows.lean}{GlobalTheory.SepRows.per0}
\leanref{Development/AbelianVarietyTheory/Definitions/TwistedInnerProducts/OffQOfSizeRows.lean}{GlobalTheory.SepRows.Reach}
\leanref{Development/AbelianVarietyTheory/BlockTheory/TwistedDiscreteSeries/CuspAlongSmoothRows.lean}{GlobalTheory.IndefiniteRowSideFullHecke.CuspAlongSmoothLifts.Sizes}

\begin{lemma}\label{asm-dependence-rows-sizes}
Let $D$ be a definite-side datum with size data at a scale $x$ whose rows have the sizes of the size
data (Definition~\ref{asm-dependence-functional-witness-sizes}). Then $p_{de}$ and $p'_{de}$ are prime to $Q_Sq$, and
$D$ has the gap sizes (Definition~\ref{asm-dependence-rows-models}(5)).
\end{lemma}
\leanref{Development/AbelianVarietyTheory/Definitions/TwistedInnerProducts/OffQOfSizeRows.lean}{GlobalTheory.SepRows.offQ_of_sizeRows}
\leanref{Development/AutomorphicArithmetic/CyclotomicBounds/RamifiedDiscriminants/FunKit.lean}{UpperTheory.Ram.Fun.coprime_QSq_of_splitK}
\leanref{Development/AbelianVarietyTheory/Proofs/TwistedOperators/SizesOfSizeRows.lean}{GlobalTheory.RowSizes.sizes_of_sizeRows}
\leanref{Development/AbelianVarietyTheory/Proofs/TwistedOperators/LargeScaleSizesImplyGapConditions.lean}{GlobalTheory.SepRows.sizeGap_of_sizes}
\leanref{Development/AbelianVarietyTheory/Definitions/TwistedInnerProducts/OffQOfSizeRows.lean}{GlobalTheory.SepRows.SizeGap}
\begin{proof}
The sizes of the rows include parameter choices for the size data. These make every level prime
split in the subfield $K$ of the host, and a prime that splits in $K$ does not divide $Q_Sqd$;
$p_{de}$ and $p'_{de}$ are level primes. The parameter choices also give a number $h_0$ with
$t+2N^3+1\le h_0\le b_{j+1}(\tau)-b_j(\tau)$ for every $\tau$ and $1\le j<m$, so
$b_{j+1}(\tau)-b_j(\tau)>2N^3$. Finally, the sizes give $m\geq44$ and $m^4\le f$, and $f=2n\le k$
for a rank datum (Definition~\ref{th-ranks}), so $r=k+m-1\geq44^4>1800$.
\end{proof}

\begin{lemma}\label{asm-dependence-rows-counts}
Let $D$ be a definite-side datum with rank datum $(k,m,t)$ and $r=k+m-1$, let $c$ be real, let $\tau$
be one of its embeddings and let $0\le i\le m$. Take the signed sets of
Definition~\ref{asm-dependence-rows-parameters} for $D$ and $c$.
\begin{enumerate}
\item The core parameter $S^{\mathrm{core}}_i(\tau)$ is regular, has $r+1+i(r+5)$ elements and
  $n_-(S^{\mathrm{core}}_i(\tau))=\lfloor(m-i+1)/2\rfloor$, and every enumeration of it is integral.
\item The row parameter $S^{\mathrm{row}}_i(\tau)$ is regular, has $r+i(r+6)$ elements and
  $n_-(S^{\mathrm{row}}_i(\tau))=\lfloor m/2\rfloor$. If $c=1$, every enumeration of it is integral.
\item The row parameter $S^{\mathrm{row}}_0(\tau)$ does not depend on $c$, and it has an integral
  enumeration with $q_\tau$ signs $-$, where $(p_\tau,q_\tau)$ is the signature of $Y_0$ at $\tau$
  (Definition~\ref{th-spaces}).
\end{enumerate}
\end{lemma}
\leanref{Development/AbelianVarietyTheory/Proofs/TwistedOperators/GridParamRegularCardMinusCount.lean}{GlobalTheory.PiCounts.isRegular_parPi}
\leanref{Development/AbelianVarietyTheory/Proofs/TwistedOperators/GridParamRegularCardMinusCount.lean}{GlobalTheory.PiCounts.card_parPi}
\leanref{Development/AbelianVarietyTheory/Proofs/TwistedOperators/GridParamRegularCardMinusCount.lean}{GlobalTheory.PiCounts.nMinus_parPi}
\leanref{Development/AbelianVarietyTheory/Proofs/TwistedOperators/BasicModelsOfFunctoriality.lean}{GlobalTheory.SepRows.isHCParam_parPi}
\leanref{Development/AbelianVarietyTheory/Definitions/Divisors/LeastSplitAuxiliaryPrimeThetaChainSpec.lean}{GlobalTheory.ChiGrid.num_nearMid}
\leanref{Development/AbelianVarietyTheory/Definitions/Divisors/LeastSplitAuxiliaryPrimeThetaChainSpec.lean}{GlobalTheory.ChiGrid.sLen_lt_gap}
\leanref{Development/AbelianVarietyTheory/Proofs/TwistedOperators/GridParameterRegularCardEqRankFormula.lean}{GlobalTheory.GridParDCount.isRegular_parD}
\leanref{Development/AbelianVarietyTheory/Proofs/TwistedOperators/GridParameterRegularCardEqRankFormula.lean}{GlobalTheory.GridParDCount.card_parD}
\leanref{Development/AbelianVarietyTheory/Definitions/Foundations/UnitarySixTwoHarishChandraParameter.lean}{GlobalTheory.SignedHC.Model.nMinus_parD}
\leanref{Development/AbelianVarietyTheory/Proofs/TwistedOperators/ExistsPermOfCardFalse.lean}{GlobalTheory.GridJoin.dualPar_shift_counts}
\leanref{Development/AbelianVarietyTheory/Proofs/TwistedOperators/LevelModelsOfFunctoriality.lean}{GlobalTheory.SepRows.isHCParam_dual_parV}
\leanref{Development/AbelianVarietyTheory/Definitions/Divisors/LeastSplitAuxiliaryPrimeThetaChainSpec.lean}{GlobalTheory.ChiGrid.parV_zero}
\leanref{Development/AbelianVarietyTheory/Definitions/TwistedInnerProducts/ParRowZero.lean}{GlobalTheory.SepRows.parRow_zero}
\leanref{Development/AbelianVarietyTheory/Proofs/IwasawaDecomposition/GoodParParBLast.lean}{GlobalTheory.ChainModels.goodPar_parB}

\begin{lemma}\label{asm-dependence-rows-model-exists}
Assume the existence of the discrete series of $U(p,q)$ with their lowest $K$-types
(Theorem~\ref{lit-bundle-theta}(2)). Let $D$ be a definite-side datum with set $\mathcal T$ of
embeddings, let $V$ be a Hermitian space over its host field with correct signature
(Definition~\ref{th-spaces}), of rank $n$ and with signature $(p_\tau,q_\tau)$ at $\tau$, and let
$S=(S_\tau)_{\tau\in\mathcal T}$ be signed sets (Definition~\ref{asm-dependence-rows-parameters}).
For $\tau\in\mathcal T$ let $O_\tau$ and $C_\tau$ be complex $n\times n$ matrices and $\epsilon'_\tau$
signs with $O_\tau^*G_V(\tau)O_\tau=J_{\epsilon'_\tau}$ and $C_\tau O_\tau=O_\tau J_{\epsilon'_\tau}$,
where $J_\epsilon$ is the diagonal matrix with entries $\epsilon$. Suppose that every $S_\tau$ has an
integral enumeration with $q_\tau$ signs $-$. Then there is a model of discrete series for $U(V)$ with
parameter $S$ (Definition~\ref{asm-dependence-rows-models}(1)) whose Cartan involution at $\tau$ is
$C_\tau$ for every $\tau$.
\end{lemma}
\leanref{Development/AbelianVarietyTheory/Proofs/TwistedOperators/LevelModelsOfFunctoriality.lean}{GlobalTheory.SepRows.exists_model_of_enum}
\leanref{Development/AbelianVarietyTheory/Proofs/TwistedOperators/ExistsDiscreteSeriesModelOfOrthonormalBasisCarta.lean}{GlobalTheory.dsModel_of_onb_cart}
\leanref{Development/AbelianVarietyTheory/Proofs/TwistedOperators/ExistsPermOfCardFalse.lean}{GlobalTheory.GridJoin.card_false_of_gram}
\leanref{Development/AbelianVarietyTheory/Proofs/TwistedOperators/ExistsPermOfCardFalse.lean}{GlobalTheory.GridJoin.exists_perm_of_card_false}
\leanref{Development/AbelianVarietyTheory/PublishedTheorems/AutomorphicForms₂/DiscreteSeriesUpqExistsIffDsExists.lean}{GlobalTheory.Published.DiscreteSeriesUpqExists}
\begin{proof}
Fix $\tau$, and let $(\lambda,\epsilon)$ be an integral enumeration of $S_\tau$ with $q_\tau$ signs
$-$. The matrix $J_{\epsilon'_\tau}$ is congruent to $G_V(\tau)$, which has $q_\tau$ negative
eigenvalues since the signature of $V$ is correct. By the uniqueness of the signature (Sylvester's
law of inertia), $\epsilon'_\tau$ has $q_\tau$ signs $-$, as $\epsilon$ does. So there is a
permutation $\pi$ of $\{1,\dots,n\}$ with $\epsilon'_{\tau,\pi(a)}=\epsilon_a$ for every $a$. Let
$O'$ be the matrix whose $a$-th column is the $\pi(a)$-th column of $O_\tau$. Then
$O'^*G_V(\tau)O'=J_\epsilon$ and $C_\tau O'=O'J_\epsilon$.

Since $\lambda$ is strictly decreasing and integral, the existence of the discrete series gives an
irreducible unitary square-integrable representation $\sigma_\tau$ of the unitary group of
$J_\epsilon$ on a Hilbert space $H_\tau$ whose lowest $K$-type has the highest weight attached to
$(\lambda,\epsilon)$. The matrix $O'$ is invertible, since $(J_\epsilon O'^*G_V(\tau))O'=J_\epsilon^2=1$.
Put $f_\tau=O'^{-1}$. Then $f_\tau^*J_\epsilon f_\tau=(O'^{-1})^*O'^*G_V(\tau)O'O'^{-1}=G_V(\tau)$,
and the Cartan involution $f_\tau^{-1}J_\epsilon f_\tau=O'J_\epsilon O'^{-1}$ is $C_\tau$. The data
$(H_\tau,\lambda,\epsilon,\sigma_\tau,f_\tau)$, $\tau\in\mathcal T$, form the model.
\end{proof}

\begin{definition}\label{asm-dependence-rows-step-pairing}
Use the signed sets of Definition~\ref{asm-dependence-rows-parameters}(1). For finite lists $l_1,l_2$
of rational numbers and a rational $c$, let $[l_1;l_2]_c$ be the signed set of the pairs $(x+c,+)$ for
$x$ in $l_1$ and $(x+c,-)$ for $x$ in $l_2$. Write $|l|$ for the length of a list $l$, $l\,l'$ for the
concatenation of lists, and $B_t$ for the list $\tfrac{t-1}2,\tfrac{t-3}2,\dots,-\tfrac{t-1}2$ of
length $t$. A pair of lists $(l_1;l_2)$ has \emph{gaps above} $h$ if $|x-y|>h$ for all $x$ in $l_1$
and $y$ in $l_2$, and \emph{gaps at least} $h$ if $|x-y|\ge h$ for all of them.
\begin{enumerate}
\item Let $t$ be an integer and $\alpha,\beta,\gamma,\delta$ lists of rational numbers, and put
  $n=|\alpha|+|\beta|+|\gamma|+|\delta|$. They satisfy the \emph{Rallis conditions of the first kind}
  if: $t\ge1$ and $n\ge1$; each of the four lists is strictly decreasing, and no number occurs twice
  in $\alpha\beta\gamma\delta$; the entries of $\alpha$ and $\gamma$ are positive and those of $\beta$
  and $\delta$ negative; every entry $e$ has $e-\tfrac{t+1}2\in\mathbf Z$ and
  $|e|\ge\tfrac{t+1}2$; the pair $(\alpha\beta;\gamma\delta)$ has gaps above $(n+t)^3$ and gaps at
  least $120\,n\min(|\alpha|+|\beta|,|\gamma|+|\delta|)^{3/2}$; and the pair
  $(\alpha B_t\delta;\gamma\beta)$ has gaps above $(n+t)^3$. They satisfy the \emph{Rallis
  conditions of the second kind} if the same holds with $(\alpha\delta;\gamma B_t\beta)$ in place of
  $(\alpha B_t\delta;\gamma\beta)$.
\item Let $P$ and $P'$ be signed sets, $t$ an integer and $c,c'$ rational. Then $P'$ is a \emph{theta
  partner} of $P$ of length $t$ with centres $(c,c')$ if there are lists $\alpha,\beta,\gamma,\delta$
  with $P=[\alpha\beta;\gamma\delta]_c$ such that either they satisfy the Rallis conditions of the
  first kind and $P'=[\alpha B_t\delta;\gamma\beta]_{c'}$, or they satisfy those of the second kind
  and $P'=[\alpha\delta;\gamma B_t\beta]_{c'}$.
\item Let $S_1$ and $S_2$ be signed sets and $a,a'$ rational. Then $S_1$ and $S_2$ are \emph{theta
  paired} for $(a,a')$ if there are an integer $t\ge0$ and signed sets $P_1,P_2$ with $S_1=-P_1$ and
  $S_2=-P_2$ such that either $P_2$ is a theta partner of $P_1$ of length $t$ with centres $(a',a)$
  and $|x-a'|\ge\tfrac{t+3}2$ for every number $x$ occurring in $P_1$, or $P_1$ is a theta partner of
  $P_2$ of length $t$ with centres $(a,a')$.
\end{enumerate}
\end{definition}
\leanref{Development/AbelianVarietyTheory/Proofs/TwistedOperators/CartFrSgn.lean}{GlobalTheory.PairLift.PairPar}
\leanref{Development/AbelianVarietyTheory/Proofs/TwistedOperators/CartFrSgn.lean}{GlobalTheory.PairLift.ThetaPar}
\leanref{Formalism/SharpRestriction/Divisibility/SchurIrreducible.lean}{Formalism.MainNotions.SchurIrreducible.DataAR}
\leanref{Formalism/SharpRestriction/Divisibility/SchurIrreducible.lean}{Formalism.MainNotions.SchurIrreducible.DataARm}
\leanref{Formalism/LocalRingCardinality/DataTheta.lean}{Formalism.BasicNotions.TH.HypA}
\leanref{Formalism/LocalRingCardinality/DataTheta.lean}{Formalism.BasicNotions.TH.HypG}
\leanref{Formalism/LocalRingCardinality/DataTheta.lean}{Formalism.BasicNotions.TH.HypGm}
\leanref{Formalism/LocalRingCardinality/DataTheta.lean}{Formalism.BasicNotions.TH.GapGt}
\leanref{Formalism/LocalRingCardinality/DataTheta.lean}{Formalism.BasicNotions.TH.GapGe}
\leanref{Formalism/LocalRingCardinality/DataTheta.lean}{Formalism.BasicNotions.TH.block}
\leanref{Development/AbelianVarietyTheory/Definitions/TwistedInnerProducts/IsStepWeil.lean}{GlobalTheory.Chain.signed}
\leanref{Development/AbelianVarietyTheory/Definitions/Foundations/StrictInterlaceIffDescEnum.lean}{GlobalTheory.SignedHC.FarFrom}

\begin{lemma}\label{asm-dependence-rows-step-negated}
Let $D$ be a definite-side datum with set $\mathcal T$ of embeddings, $\theta$ an archimedean family,
$V$ a Hermitian space over its host field and $V^-$ the same space with the negated form, so that
$G_{V^-}(\tau)=-G_V(\tau)$, $U(G_{V^-}(\tau))=U(G_V(\tau))$ and $U(V^-)(\mathbf A)=U(V)(\mathbf A)$. Let
$M$ be a model of discrete series for $U(V)$ (Definition~\ref{asm-dependence-rows-models}(1)) with
parameter $S$ and data $(H_\tau,\lambda_\tau,\epsilon_\tau,\sigma_\tau,f_\tau)$. Then
$(H_\tau,\lambda_\tau,-\epsilon_\tau,\sigma_\tau,f_\tau)$ is a model $M^-$ of discrete series for
$U(V^-)$, whose parameter at $\tau$ is $S_\tau$ with all its signs changed and whose Cartan involution
at $\tau$ is the negative of that of $M$. Its space $(M^-)^\infty$ is $M^\infty$, and every element of
$\operatorname{Mult}_\theta(M^-)$ is an element of $\operatorname{Mult}_\theta(M)$.
\end{lemma}
\leanref{Development/AbelianVarietyTheory/Proofs/TwistedOperators/ExistsNegatedSpaceDiscreteSeriesModelTransport.lean}{GlobalTheory.negModel_tr}
\leanref{Development/AbelianVarietyTheory/Proofs/TwistedOperators/ExistsNegatedSpaceDiscreteSeriesModelTransport.lean}{GlobalTheory.NegModel.negM}
\leanref{Development/AbelianVarietyTheory/Proofs/TwistedOperators/ExistsNegatedSpaceDiscreteSeriesModelTransport.lean}{GlobalTheory.NegModel.cart_negM}
\leanref{Development/AbelianVarietyTheory/Proofs/TwistedOperators/ExistsNegatedSpaceDiscreteSeriesModelTransport.lean}{GlobalTheory.NegModel.multT_ofNeg}
\leanref{Development/AbelianVarietyTheory/Proofs/TwistedOperators/IsDiscreteSeriesNegSignatureOfIsDiscreteSeries.lean}{GlobalTheory.PairLift.isDS_flip}
\leanref{Development/AbelianVarietyTheory/Proofs/TwistedOperators/AutomorphicCuspFormsHermitianSpaceSignTwistStatement.lean}{GlobalTheory.T.automorphic_cusp_forms_hermitianSpaceSignTwist}
\begin{proof}
Fix $\tau$. The pairs $(\lambda_{\tau,a},-\epsilon_{\tau,a})$ form $S_\tau$ with its signs changed, and
$\lambda_\tau$ is still strictly decreasing. The diagonal matrix with entries $-\epsilon_\tau$ is $-J_\tau$,
whose unitary group is that of $J_\tau$, with the same maximal compact subgroup. So $\sigma_\tau$ is an
irreducible unitary square-integrable representation of it, with the same $K$-types; the highest
weight attached to $(\lambda_\tau,-\epsilon_\tau)$ is the one attached to $(\lambda_\tau,\epsilon_\tau)$,
and the norms that decide which $K$-type is lowest do not change. Further
$f_\tau^*(-J_\tau)f_\tau=-G_V(\tau)=G_{V^-}(\tau)$, and the Cartan involution
$f_\tau^{-1}(-J_\tau)f_\tau$ is the negative of that of $M$. So $M^-$ is a model with the stated
parameter and Cartan involutions.

Both models let $U(G_V(\tau))=U(G_{V^-}(\tau))$ act on $H_\tau$ by $g\mapsto\sigma_\tau(f_\tau gf_\tau^{-1})$,
and the two groups have the same Lie algebra, so the smooth vectors and the actions agree, and
$(M^-)^\infty=M^\infty$. The archimedean family embeds $U(G_{V^-}(\tau))$ in $U(V^-)(\mathbf A)$ as it
embeds $U(G_V(\tau))$ in $U(V)(\mathbf A)$. A function on $U(V)(\mathbf A)=U(V^-)(\mathbf A)$ is a
cuspidal automorphic form, smooth at the archimedean places, finite under the centres and of moderate
growth for $V^-$ if and only if it is one for $V$, since the groups, their rational points and their
Lie algebras are the same. So a multiplicity map for $M^-$ (Definition~\ref{asm-dependence-rows-models}(2))
is one for $M$, with the same values, and the span of the first is contained in that of the second.
\end{proof}

\begin{lemma}\label{asm-dependence-rows-step-frames}
Let $D$ be a definite-side datum, $\psi$ a standard character, $\theta$ an archimedean family, $V$ and
$W$ Hermitian spaces over its host field, $\Phi,\Phi'$ frames of $V$ and $\Xi,\Xi'$ frames of $W$ at
the real places such that $\Phi'$ has the Cartan involutions of $\Phi$ and $\Xi'$ those of $\Xi$. There
is a linear automorphism $j$ of the algebra of polynomials in the creation operators of the pair
$(V,W)$ such that for every theta function $f$ for $(V,W)$ and every polynomial $P$, the vector
$f\otimes j(P)$ formed with the frames $\Phi'$ and $\Xi'$ equals the vector $f\otimes P$ formed with
$\Phi$ and $\Xi$.
\end{lemma}
\leanref{Development/AbelianVarietyTheory/Proofs/TwistedOperators/FockVectorFrameChangeSameCartan.lean}{GlobalTheory.FrameChange.ptF_frame_change}
\leanref{Development/AbelianVarietyTheory/Proofs/TwistedOperators/FockVectorFrameChangeSameCartan.lean}{GlobalTheory.FrameChange.sub}
\leanref{Development/AbelianVarietyTheory/Proofs/TwistedOperators/FockVectorFrameChangeSameCartan.lean}{GlobalTheory.FrameChange.aPlusAll_change}
\leanref{Development/AbelianVarietyTheory/Proofs/TwistedOperators/FockVectorFrameChangeSameCartan.lean}{GlobalTheory.FrameChange.ptF_sub}

This is a computation. At each $\tau$ the matrices that change one orthonormal basis into the other
commute with the common Cartan involution, so their entries between basis vectors of opposite signs
vanish. Each creation operator for $(\Phi',\Xi')$ is then, on theta functions, the linear combination
of the creation operators for $(\Phi,\Xi)$ with coefficients the products of entries of the two
matrices, conjugated where the signs differ. The automorphism $j$ substitutes these combinations for
the variables; its inverse is the substitution for the opposite change of frames.

\begin{lemma}\label{asm-dependence-rows-step-models}
Assume the results of Theorem~\ref{lit-bundle-theta}. Let $D$ be a definite-side datum with set
$\mathcal T$ of embeddings, let $V$ and $W$ be Hermitian spaces over its host field, and let $M$ and
$M'$ be models of discrete series for $U(V)$ and $U(W)$ (Definition~\ref{asm-dependence-rows-models}(1))
with parameters $S$ and $S'$. Let $a,a'\colon\mathcal T\to\mathbf Q$, and suppose that $S_\tau$ and
$S'_\tau$ are theta paired for $(a_\tau,a'_\tau)$ (Definition~\ref{asm-dependence-rows-step-pairing}(3))
for every $\tau$. Let $\theta$ be an archimedean family and $\psi$ a standard character, and form the
vectors $f\otimes P$ with the frames of the models: at $\tau$ the orthonormal bases given by $f_\tau^{-1}$
and the Cartan involutions of $M$ and $M'$. Then there is a map $I_0$ from pairs $(P,v)$ of a
polynomial $P$ in the creation operators of the pair $(V,W)$ and a vector $v\in M'^\infty$ to
$M^\infty$, linear in each variable, with the following property. Let $\omega$ be a Weil datum for
$(V,W)$ whose characters have infinity exponents $a$ and $a'$ and which has a nonzero Fock vacuum
vector for these frames, and let $\nu$ be a Haar measure on $U(W)(\mathbf A)$ with a fundamental
domain for $U(W)(F)$. For every Fock vacuum vector $f$ of $\omega$ for these frames and every
$T\in\operatorname{Mult}_\theta(M')$ there is $T_1\in\operatorname{Mult}_\theta(M)$ such that for all
$P$ and $v$ the theta lift of $T(v)$ by $\omega$ with the vector $f\otimes P$ and the measure $\nu$
converges at every point and equals $T_1(I_0(P,v))$.
\end{lemma}
\leanref{Development/AbelianVarietyTheory/Definitions/TwistedInnerProducts/ParRowZero.lean}{GlobalTheory.SepRows.PairEx}
\leanref{Development/AutomorphicArithmetic/ThetaCorrespondence/ThetaPairLifts/LocalTypeDictionary/ExistsSeesawPairOfThetaData.lean}{GlobalTheory.T.exists_seesaw_pair_of_theta_data}
\leanref{Development/AbelianVarietyTheory/Proofs/TwistedOperators/CartFrSgn.lean}{GlobalTheory.PairLift.PairLifts}
\leanref{Development/AbelianVarietyTheory/Proofs/TwistedOperators/PairThetaLiftFactorsThroughIntertwiner.lean}{GlobalTheory.PairLift.RuleLift.pairLifts}
\leanref{Development/AbelianVarietyTheory/Proofs/DecoratedLatticeProofs/FirstOccurrenceRuleThetaChainStepBothDirections.lean}{GlobalTheory.StepRule.rulePar_of_pairPar}
\leanref{Development/AbelianVarietyTheory/Proofs/DecoratedLatticeProofs/ExistsRealDiscreteSeriesDataThetaProperties.lean}{GlobalTheory.RealPlaceTheta.exists_record_theta_all}
\leanref{Formalism/LocalToGlobal/IndefiniteReal/LinksChar.lean}{Formalism.IndefR.Links.IsItW}
\begin{proof}
Alpöge's proof has two parts. First, from a seesaw pair of theta data at the real places, it
constructs $I_0$ so that, for every Weil datum $\omega$ as in the lemma and every nonzero Fock vacuum
vector $f_0$, the map $I_0$ is an intertwiner for the vectors $f_0\otimes P$: it is nonzero; on pure
tensors $\otimes_\tau v_\tau$ its pairing with a fixed pure tensor is bounded by a constant times
$\prod_\tau\|v_\tau\|$; if $f_0\otimes P'$ is the translate of $f_0\otimes P$ under $U(G_V(\tau))$, or
its derivative along the Lie algebra, then $I_0(P',v)$ is the translate, or the derivative, of
$I_0(P,v)$; and $I_0(P,v)$ is unchanged when $U(G_W(\tau))$ moves $f_0\otimes P$ and $v$ together,
with the corresponding identity for the Lie algebra. Second, the parameters satisfy the Rallis
conditions (Definition~\ref{asm-dependence-rows-step-pairing}(1)), under which the results on the
archimedean theta correspondence among those of Theorem~\ref{lit-bundle-theta} apply to the local
models; with these, the theta lift of a multiplicity map with the vectors $f\otimes P$ is shown to be
a multiplicity map evaluated at $I_0(P,v)$.
\openstep{The construction of the intertwiner $I_0$ from a seesaw pair of theta data at the real
places, and the factorization of the theta lifts of multiplicity maps through it under the Rallis
conditions, are not yet written out.}
% Open: PairEx = exists_seesaw_pair_of_theta_data = PairIt.exists_seesaw_pair (ItPairObj.exists_record,
% thetaModels_eq_coverDiscreteSeries, exists_locIt_pair, exists_globalFockIntertwiner_tensor_realPlaces,
% thetaVacuumCharacters_atRealPlaces, fockVectorsOverVacuumIsometry_of_skeletonDatum, isItW_of_globIt),
% a depth-6 count of about 13,600 lines of Lean proofs; PairLifts = RuleLift.pairLifts (cuspAll_of_lowAll,
% lowAll_of_lowFock, LowerMembersRule.lowFock, exists_local, Inputs from inputs_of_chosenSeries), about
% 11,600 lines; StepRule.rulePar_of_pairPar from the Rallis inputs ArchimedeanRallisUpDownLine, Fourth and
% Fifth of exists_record_theta_all (lowest_kTypes_jointShape_block_unitary,
% realThetaRule_exchange_pair_members), about 14,200 lines. The intertwiner condition is Links.IsItW.
\end{proof}

\begin{lemma}\label{asm-dependence-rows-step-lift}
Assume the results of Theorem~\ref{lit-bundle-theta}. Let $D$ be a definite-side datum with set
$\mathcal T$ of embeddings, let $A$ and $B$ be Hermitian spaces over its host field, and let $M_A$
and $M_B$ be models of discrete series for $U(A)$ and $U(B)$ (Definition~\ref{asm-dependence-rows-models}(1))
with parameters $S$ and $S'$. Let $a,a'\colon\mathcal T\to\mathbf Q$, and suppose that for every
$\tau$ the signed set $S'_\tau$ with all its signs changed and $S_\tau$ are theta paired for
$(a_\tau,a'_\tau)$ (Definition~\ref{asm-dependence-rows-step-pairing}(3)). Let $\theta$ be an
archimedean family, $\psi$ a standard character, and frames of $A$ and $B$ at the real places, that
is, orthonormal bases with Cartan involutions, whose Cartan involutions are those of $M_A$ and $M_B$.
Let $\omega$ be a Weil datum (Definition~\ref{th-weil-datum}) for the pair formed by $B$ with the
negated form and $A$, whose characters have infinity exponents $a$ and $a'$, let $f_0$ be a nonzero
Fock vacuum vector of $\omega$ for these frames, and let $\nu$ be a Haar measure on $U(A)(\mathbf A)$
with a fundamental domain for $U(A)(F)$. Then there is a map $I$ from pairs $(P,v)$ of a polynomial
$P$ in the creation operators of the pair and a vector $v\in M_A^\infty$ to $M_B^\infty$ with the
following property. For every Fock vacuum vector $f$ of $\omega$ for these frames and every
$T\in\operatorname{Mult}_\theta(M_A)$ there is $T'\in\operatorname{Mult}_\theta(M_B)$ such that for all
$P$ and $v$ the theta lift of $T(v)$ by $\omega$ with the vector $f\otimes P$ and the measure $\nu$
converges at every point and equals $T'(I(P,v))$.
\end{lemma}
\leanref{Development/AbelianVarietyTheory/Proofs/TwistedOperators/LevelInductionFromPairs.lean}{GlobalTheory.RowArrow.lifts_neg}
\leanref{Development/AbelianVarietyTheory/Proofs/TwistedOperators/CartFrSgn.lean}{GlobalTheory.PairLift.PairPar}
\leanref{Development/AbelianVarietyTheory/Proofs/TwistedOperators/CartFrSgn.lean}{GlobalTheory.PairLift.IsWeilExp}
\leanref{Development/AbelianVarietyTheory/Definitions/TwistedInnerProducts/IsVacFOfIsVacVec.lean}{GlobalTheory.ChainFr.IsVacF}
\leanref{Development/AbelianVarietyTheory/Proofs/TwistedOperators/CartFrSgn.lean}{GlobalTheory.PairLift.frSgn}
\begin{proof}
Write $B^-$ for $B$ with the negated form. Lemma~\ref{asm-dependence-rows-step-negated} gives a model
$N=M_B^-$ of discrete series for $U(B^-)$ with parameter $S'$ with all its signs changed, with
$N^\infty=M_B^\infty$, and whose Cartan involution at each $\tau$ is the negative of that of $M_B$. The
frame of $B$ at $\tau$ with its signs changed is a frame of $B^-$, with the same orthonormal basis; its
Cartan involution is the negative of that of the frame of $B$, which is that of $M_B$, so it is that
of $N$. The Cartan involutions of the frame of $A$ are those of $M_A$. So
Lemma~\ref{asm-dependence-rows-step-frames}, for the pair $(B^-,A)$, the frames of $B^-$ and $A$ in the
roles of $\Phi$ and $\Xi$ and the frames of the models $N$ and $M_A$ in the roles of $\Phi'$ and $\Xi'$,
gives a linear automorphism $j$ of the polynomials in the creation operators such that for every theta
function $f$ the vector $f\otimes j(P)$ formed with the frames of the models equals the vector
$f\otimes P$ formed with the given frames.

A Fock vacuum vector for a pair of frames is a theta function annihilated by the annihilation
operators at the real places, which depend only on the Cartan involutions. So the Fock vacuum vectors
of $\omega$ for the given frames are those for the frames of $N$ and $M_A$; in particular $f_0$ is a
nonzero one. By hypothesis the parameter of $N$ at $\tau$, which is $S'_\tau$ with its signs changed,
and $S_\tau$ are theta paired for $(a_\tau,a'_\tau)$. Lemma~\ref{asm-dependence-rows-step-models},
applied to $V=B^-$, $W=A$, $M=N$ and $M'=M_A$, gives a map $I_0$, and $\omega$ satisfies its
hypotheses. Put $I(P,v)=I_0(j(P),v)$, an element of $N^\infty=M_B^\infty$.

Let $f$ be a Fock vacuum vector of $\omega$ for the given frames and
$T\in\operatorname{Mult}_\theta(M_A)$. Then $f$ is a Fock vacuum vector for the frames of the models,
so Lemma~\ref{asm-dependence-rows-step-models} gives $T_1\in\operatorname{Mult}_\theta(N)$ such that for
all $P'$ and $v$ the theta lift of $T(v)$ with the vector $f\otimes P'$, formed with the frames of the
models, converges at every point and equals $T_1(I_0(P',v))$. By
Lemma~\ref{asm-dependence-rows-step-negated}, $T_1$ is an element $T'$ of
$\operatorname{Mult}_\theta(M_B)$, with the same values on $U(B^-)(\mathbf A)=U(B)(\mathbf A)$. For
$P'=j(P)$ the vector $f\otimes j(P)$ is the vector $f\otimes P$ formed with the given frames. So the
theta lift of $T(v)$ with $f\otimes P$ converges at every point and equals
$T_1(I_0(j(P),v))=T'(I(P,v))$.
\end{proof}

\begin{lemma}\label{asm-dependence-rows-local}
Assume the published results on the local theta correspondence (Theorem~\ref{lit-bundle-theta}).
\begin{enumerate}
\item In every row setting (Definition~\ref{asm-dependence-rows-models}(3)), for every $i\le m$
  there is a core model at $i$.
\item In every gap row setting (Definition~\ref{asm-dependence-rows-models}(5)), for every $i\le m$
  there is a row model at $i$.
\item Let a gap row setting be given, let $0\le i<m$, and let $M_i$ and $M_{i+1}$ be row models at
  $i$ and $i+1$. There is a map $I$ from pairs $(P,v)$ of a polynomial $P$ in the creation operators
  of the pair $(V_{i+1,0},V_{i,0})$ and a vector $v\in M_i^\infty$ to $M_{i+1}^\infty$ with the
  following property. For every $T\in\operatorname{Mult}_\theta(M_i)$ and every coset $\kappa$ of
  the pair lattice of $(V_{i+1,0},V_{i,0})$ modulo $\mathfrak m$ there is
  $T'\in\operatorname{Mult}_\theta(M_{i+1})$ such that for all $P$ and $v$ the theta lift of $T(v)$
  from $U(V_{i,0})$ to $U(V_{i+1,0})$ with the vector $f_{i,\kappa}\otimes P$ and the measures of
  the row volume (Definition~\ref{asm-grid-evaluation}(4)) converges at every point and equals
  $T'(I(P,v))$.
\end{enumerate}
\end{lemma}
\leanref{Development/AutomorphicArithmetic/ThetaCorrespondence/ThetaPairLifts/LocalTypeDictionary/SeparatingFamilyHasModels.lean}{GlobalTheory.T.exists_separating_models}
\leanref{Development/AbelianVarietyTheory/Definitions/TwistedInnerProducts/OffQOfSizeRows.lean}{GlobalTheory.SepRows.CoreModels}
\leanref{Development/AutomorphicArithmetic/ThetaCorrespondence/ThetaPairLifts/LocalTypeDictionary/SeparatingModelsOfSeparatingData.lean}{GlobalTheory.T.exists_models_of_separating_data}
\leanref{Development/AbelianVarietyTheory/Definitions/TwistedInnerProducts/ParRowZero.lean}{GlobalTheory.SepRows.RowModels}
\leanref{Development/AutomorphicArithmetic/ThetaCorrespondence/ThetaPairLifts/LocalTypeDictionary/SeparatingFamilyExtendsByTwistedRow.lean}{GlobalTheory.T.separating_family_extend_by_twist}
\leanref{Development/AutomorphicArithmetic/ThetaCorrespondence/ThetaPairLifts/LocalTypeDictionary/ExistsSeesawPairOfThetaData.lean}{GlobalTheory.T.exists_seesaw_pair_of_theta_data}
\leanref{Development/AbelianVarietyTheory/Definitions/TwistedInnerProducts/ParRowZero.lean}{GlobalTheory.SepRows.RowStep}
\leanref{Development/AbelianVarietyTheory/Definitions/TwistedInnerProducts/ParRowZero.lean}{GlobalTheory.SepRows.PairEx}
\leanref{Development/AutomorphicArithmetic/ThetaCorrespondence/ThetaPairLifts/LocalTypeDictionary/SeparatingFamilyExtendsByTwistedRow.lean}{GlobalTheory.RowArrow.separatingFamily_extends_by_twisted_row}
\leanref{Development/AbelianVarietyTheory/Proofs/TwistedOperators/ContragredientThetaParamsAddBlockArrow.lean}{GlobalTheory.GridParD.pairPar_X}
\leanref{Development/AbelianVarietyTheory/Proofs/TwistedOperators/LevelInductionFromPairs.lean}{GlobalTheory.RowArrow.exp_att}

\begin{proof}
For every grid space $X$ and every $\tau$, the grid Weil datum gives an orthonormal basis $O_\tau$ of
$X$ at $\tau$ and a Cartan involution $C_\tau$ with $O_\tau^*G_X(\tau)O_\tau=J_{\epsilon'}$ and
$C_\tau O_\tau=O_\tau J_{\epsilon'}$ for some signs $\epsilon'$. A core or row model at $i$ is a model
with these Cartan involutions, so it suffices in (1) and (2) to check the hypothesis on the
parameter in Lemma~\ref{asm-dependence-rows-model-exists}.

(1) Let $i\le m$ and $\tau$ be given. By Definition~\ref{th-grid}(3), $P_i$ has correct signature,
rank $r+1+i(r+5)$ and $q_\tau(P_i)=\lfloor(m-i+1)/2\rfloor$. By Lemma~\ref{asm-dependence-rows-counts}(1),
$S^{\mathrm{core}}_i(\tau)$ is regular with that many elements and that many signs $-$. Listing its
numbers in decreasing order, each with its sign, gives an enumeration with $q_\tau(P_i)$ signs $-$,
and it is integral by the same lemma. Lemma~\ref{asm-dependence-rows-model-exists} gives a core
model at $i$.

(2) Let $i\le m$. If $i=0$, then $V_{0,0}=Y_0$, which has correct signature, and
Lemma~\ref{asm-dependence-rows-counts}(3) gives for every $\tau$ an integral enumeration of
$S^{\mathrm{row}}_0(\tau)$ with $q_\tau(Y_0)$ signs $-$. If $i\ge1$, the first constant of the grid is
$c=1$, since the row setting is a gap row setting. By Definition~\ref{th-grid}(2), $V_{i,0}$ has
correct signature, rank $r+i(r+6)$ and $q_\tau(V_{i,0})=\lfloor m/2\rfloor$. By Lemma~\ref{asm-dependence-rows-counts}(2),
$S^{\mathrm{row}}_i(\tau)$ is regular with that many elements and that many signs $-$, so it has an
enumeration with $\lfloor m/2\rfloor$ signs $-$, which is integral since $c=1$. In both cases
Lemma~\ref{asm-dependence-rows-model-exists} gives a row model at $i$.

(3) Let $0\le i<m$ and the row models $M_i$ and $M_{i+1}$ be given, with parameters
$S^{\mathrm{row}}_i$ and $S^{\mathrm{row}}_{i+1}$, and write $c(\tau)$ for the centre of the grid at
the lift $m-i$ and the embedding $\tau$. For a rational $c$ and an integer $n$ put
$\operatorname{att}(c,n)=c-\tfrac12$ if $n$ is odd and $\operatorname{att}(c,n)=c$ if $n$ is even, and
let $a(\tau)=\operatorname{att}(c(\tau),\operatorname{rk}V_{i+1,0})$ and
$a'(\tau)=\operatorname{att}(c(\tau),\operatorname{rk}V_{i,0})$. We check the hypotheses of
Lemma~\ref{asm-dependence-rows-step-lift} for $A=V_{i,0}$, $B=V_{i+1,0}$, $M_A=M_i$, $M_B=M_{i+1}$, the
archimedean family $\theta$ and the standard character $\psi$ of the grid Weil datum.

The grid Weil datum gives orthonormal bases and Cartan involutions of $V_{i,0}$ and $V_{i+1,0}$ at the
real places, which form frames; the Cartan involutions are those of the row models. Since the first
constant of the grid is $1$ and $D$ has the gap sizes, at every $\tau$ the signed set
$S^{\mathrm{row}}_{i+1}(\tau)$ with its signs changed and $S^{\mathrm{row}}_i(\tau)$ are theta paired
for $(a(\tau),a'(\tau))$ (Definition~\ref{asm-dependence-rows-step-pairing}(3)); this is a computation with the signed sets of
Definition~\ref{asm-dependence-rows-parameters}. The character $\eta_L$ of the level system has
infinity exponent $\tfrac12$, and $\chi_{m-i}$ has infinity exponent $c(\tau)$ at every $\tau$, by the
conditions on the characters of the grid. The characters of the pair $(V_{i+1,0},V_{i,0})$ are
$\chi_{m-i}$ times $\eta_L^{-1}$ for a space of odd rank and $\chi_{m-i}$ for one of even rank
(Definition~\ref{th-grid}), so the Weil datum $\omega$ of the grid Weil datum at this pair, for
$V_{i+1,0}$ with the negated form and $V_{i,0}$, has characters with infinity exponents $a$ and $a'$.
The vacuum vector of $\omega$ in the grid Weil datum is a nonzero theta function that is a Fock vacuum
vector at every real place for the Cartan involutions of the grid, hence for the two frames. Every
coset vector $f_{i,\kappa}$ is a theta function for $\omega$ with the same components at the real
places as the vacuum vector, so it is a Fock vacuum vector for the two frames as well. Finally, the
row volume gives a Haar measure on $U(V_{i,0})(\mathbf A)$ with a fundamental domain.

Lemma~\ref{asm-dependence-rows-step-lift} now gives a map $I$. Let $T\in\operatorname{Mult}_\theta(M_i)$
and a coset $\kappa$ be given. Applying the lemma to the Fock vacuum vector $f_{i,\kappa}$ gives
$T'\in\operatorname{Mult}_\theta(M_{i+1})$ such that for all $P$ and $v$ the theta lift of $T(v)$ with
the vector $f_{i,\kappa}\otimes P$, formed with the frames of the grid Weil datum, converges at every
point and equals $T'(I(P,v))$. This vector is the vector $f_{i,\kappa}\otimes P$ of the grid Weil
datum, whose orthonormal bases and Cartan involutions are those of the frames.
\end{proof}

\begin{lemma}\label{asm-dependence-rows-tower}
Assume the results of Theorem~\ref{lit-bundle-theta}. Let a gap row setting be given
(Definition~\ref{asm-dependence-rows-models}(5)), and let $M_0$ be a row model at $0$. For every
$n\le m$ there are a row model $M_n$ at $n$ and a map $v_n$ from pairs $(\rho,w)$ of a real vector
$\rho$ (Definition~\ref{asm-grid-evaluation}(4)) and a vector $w\in M_0^\infty$ to $M_n^\infty$ with
the following property. For every $T\in\operatorname{Mult}_\theta(M_0)$ and every coset datum $\Phi$
there is $T_n\in\operatorname{Mult}_\theta(M_n)$ such that, for every real vector $\rho$ and every
$w\in M_0^\infty$, the partial row lifts of $T(w)$ with the row data $(\rho,\Phi)$
(Definition~\ref{th-core-chain}(1)) converge up to $n$ and
\[
 \theta^{(n)}_{\rho,\Phi}\bigl(T(w)\bigr)=T_n\bigl(v_n(\rho,w)\bigr).
\]
\end{lemma}
\leanref{Development/AbelianVarietyTheory/Proofs/TwistedOperators/FirstRowTowerMultiplicitySpacesOfArrowStatement.lean}{GlobalTheory.T.firstRowTower_multiplicitySpaces_of_arrow}
\leanref{Development/AbelianVarietyTheory/Proofs/TwistedOperators/FirstRowTowerMultiplicitySpacesOfArrow.lean}{GlobalTheory.SepRows.rowTower_of_step}
\leanref{Development/AbelianVarietyTheory/Definitions/TwistedInnerProducts/ParRowZero.lean}{GlobalTheory.SepRows.RowTower}
\begin{proof}
Induction on $n$. For $n=0$ take $M_0$ itself, $v_0(\rho,w)=w$ and $T_0=T$; the lift $\theta^{(0)}$
is the identity, and there is nothing to converge.

Let $n<m$, and let $M_n$ and $v_n$ be as in the lemma. Lemma~\ref{asm-dependence-rows-local}(2)
gives a row model $M_{n+1}$ at $n+1$, and Lemma~\ref{asm-dependence-rows-local}(3) with $i=n$ gives
a map $I$. Put $v_{n+1}(\rho,w)=I(\rho_n,v_n(\rho,w))$. Let $T$ and $\Phi$ be given, and let $T_n$
be as in the lemma. Lemma~\ref{asm-dependence-rows-local}(3), applied to $T_n$ and the coset
$\kappa=\Phi_n$, gives $T_{n+1}\in\operatorname{Mult}_\theta(M_{n+1})$. Let $\rho$ and $w$ be given.
The partial lifts up to $n$ converge, and $\theta^{(n)}_{\rho,\Phi}(T(w))=T_n(v_n(\rho,w))$. The
lift $\theta^{(n+1)}_{\rho,\Phi}$ is the theta lift of $\theta^{(n)}_{\rho,\Phi}$ with the vector
$f_{n,\Phi_n}\otimes\rho_n$. By the choice of $T_{n+1}$ it converges, and it equals
$T_{n+1}(I(\rho_n,v_n(\rho,w)))=T_{n+1}(v_{n+1}(\rho,w))$.
\end{proof}

\begin{lemma}\label{asm-dependence-rows-cusp}
Assume rapid decay of cusp forms (Theorem~\ref{lit-bundle-theta}(3)). Let $D$ be a definite-side
datum, $V$ a Hermitian space over its host field, $M$ a model of discrete series for $U(V)$
(Definition~\ref{asm-dependence-rows-models}(1)), $\theta$ an archimedean family,
$J\subset U(V)(\mathbf A)$ and $\chi$ a function on $U(V)(\mathbf A)$. For every
$T\in\operatorname{Mult}_\theta(M;J,\chi)$ and every $v\in M^\infty$, the function $T(v)$ is a
cuspidal automorphic form on $U(V)(\mathbf A)$, smooth at the archimedean places, finite under the
centres of the enveloping algebras and of moderate growth, and it is rapidly decreasing.
\end{lemma}
\leanref{Development/AbelianVarietyTheory/Proofs/DecoratedLatticeProofs/MultTApplyRapidDec.lean}{GlobalTheory.ChainOnto.MultRapid.multT_apply_rapidDec}
\leanref{Development/AbelianVarietyTheory/Proofs/TwistedOperators/LowerTowerLiftsVanishOfVanishOnFock.lean}{GlobalTheory.PairLift.LowAllPf.multT_apply_smooth_cusp}
\leanref{Development/AbelianVarietyTheory/PublishedTheorems/TwistedTraceFormula/UnifModGrowthOfIsUnifModerate.lean}{GlobalTheory.Grow.rapidDec_of_cuspForm}
\begin{proof}
Write $T$ as a finite linear combination of multiplicity maps of level $(J,\chi)$, and $v$ as a
finite sum of tensors $\bigotimes_\tau v_\tau$. Then $T(v)$ is a finite linear combination of values
of multiplicity maps at families $(v_\tau)$, each of which has the properties listed in
Definition~\ref{asm-dependence-rows-models}(2). These properties are preserved by sums and scalar
multiples, so $T(v)$ has them. A smooth cuspidal automorphic form of moderate growth is rapidly
decreasing, by rapid decay of cusp forms.
\end{proof}

\begin{lemma}\label{asm-dependence-rows-bessel-conv}
Let a row setting be given (Definition~\ref{asm-dependence-rows-models}(3)), with grid spaces
$V_{i,j}$, $P_i$ and $Q_j$ and archimedean family $\theta$. Let $F$ be a continuous function on
$U(V_{m,0})(\mathbf A)$ such that $F(\gamma g)=F(g)$ for $\gamma\in U(V_{m,0})(F)$ and $F$ is
invariant on the right under an open compact subgroup of the finite adelic group. Then for every
column $\Psi$ and every linear functional $\xi$ on the coefficients
(Definition~\ref{asm-grid-evaluation}(5)), the Bessel period $B^+(\varphi_{\xi,\Psi},F)$ converges.
\end{lemma}
\leanref{Development/AbelianVarietyTheory/Proofs/TwistedOperators/BesselPeriodCorankConvergesColumn.lean}{GlobalTheory.RowBplus.bplusConv_col}
\leanref{Development/AbelianVarietyTheory/Proofs/TwistedOperators/BesselPeriodCorankConvergesColumn.lean}{GlobalTheory.RowBplus.integrable_of_haar}
\leanref{Development/AbelianVarietyTheory/Proofs/TwistedOperators/BesselPeriodCorankConvergesColumn.lean}{GlobalTheory.RowBplus.pairVec_isAutFun}
\leanref{Development/AbelianVarietyTheory/Proofs/FlagMeasures/ExistsAllTiedPerTied.lean}{GlobalTheory.Per.Inhab.bes_inv}
\begin{proof}
Write $\varphi=\varphi_{\xi,\Psi}$, $\iota$ for the embedding of $U(Q_0)$ in $U(V_{m,0})$, $N$ for
the unipotent radical and $\psi_N$ for the character of Definition~\ref{asm-grid-evaluation}(5). The
function $\varphi$ is continuous, invariant on the right under an open compact subgroup, and
invariant on the left under $U(Q_0)(F)$: for rational $\gamma$,
$\varphi(\gamma h)=\xi(\varrho(h)^{-1}\varrho(\gamma)^{-1}\Psi(\gamma h_f))=\varphi(h)$, since
$\Psi(\gamma g)=\varrho(\gamma)\Psi(g)$. Put
$\Phi(h,n)=\varphi(h)F(n\iota(h))\overline{\psi_N(n)}$ for $h\in U(Q_0)(\mathbf A)$ and
$n\in N(\mathbf A)$; it is continuous. For $\gamma\in U(Q_0)(F)$, the element $\iota(\gamma)$
normalizes $N$ and preserves $\psi_N$, and
$\Phi(\gamma h,n)=\Phi(h,\iota(\gamma)^{-1}n\iota(\gamma))$, by the left invariance of $\varphi$ and
of $F$. For $\delta\in N(F)$, $\Phi(h,\delta n)=\Phi(h,n)$, since $F$ is left invariant and $\psi_N$
is trivial on $N(F)$.

The space $Q_0$ is totally definite (Definition~\ref{th-grid}(4)), hence anisotropic, so
$U(Q_0)(\mathbf A)=U(Q_0)(F)K$ for a compact set $K$ \cite{BHC62}; and $N(\mathbf A)=N(F)K'$ for a
compact set $K'$. Given $(h,n)$, write $h=\gamma k$ with $\gamma$ rational and $k\in K$, and
$\iota(\gamma)^{-1}n\iota(\gamma)=\delta k'$ with $\delta\in N(F)$ and $k'\in K'$. Then
$\Phi(h,n)=\Phi(k,k')$, so $|\Phi|$ is bounded by its maximum on $K\times K'$. The fundamental domain
$\mathcal F_{Q_0}$ of the row volume is covered by the sets $\gamma K$, $\gamma\in U(Q_0)(F)$, so its
measure is at most $\sum_\gamma\mu(\gamma K\cap\mathcal F_{Q_0})=\mu(K)<\infty$; and $\mathcal F_N$
has measure $1$. So the inner integral of $B^+(\varphi,F)$ converges absolutely for every $h$, and
the outer one converges absolutely.
\end{proof}

\begin{definition}\label{asm-dependence-rows-bessel-forms}
Let $J$ be an invertible Hermitian $n\times n$ matrix and $J_0$ an invertible $m\times m$ matrix,
and write $U(J)$ for the group of $g\in GL_n(\mathbf C)$ with $g^*Jg=J$.
\begin{enumerate}
\item A \emph{Bessel frame} for $(J,J_0)$ consists of an $n\times m$ matrix $F$ with $F^*JF=J_0$, an
  $n\times r$ matrix $e$ with linearly independent columns $e_1,\dots,e_r$ such that $e^*Je=0$ and
  $F^*Je=0$, and a vector $z$ with $F^*Jz=0$, $e^*Jz=0$ and $z^*Jz\ne0$, where $m+2r+1=n$. Its
  \emph{unipotent group} $N$ is the set of $u\in U(J)$ such that $(u-1)e_i$ lies in the span of
  $e_1,\dots,e_{i-1}$ for every $i$, and $(u-1)x$ lies in the span of the $e_i$ for every $x$ with
  $e^*Jx=0$. Put $f_{j}=e_{j}$ for $j\le r$ and $f_{r+1}=z$. A function $\chi$ on $N$ is a
  \emph{generic character} if there are nonzero $c_1,\dots,c_r\in\mathbf C$ such that
  $\chi(u)=\exp\bigl(i\operatorname{Re}\sum_jc_ja_{jj}\bigr)$ whenever $u\in N$ and
  $(u-1)f_{j+1}=\sum_ka_{jk}e_k$ for $1\le j\le r$.
\item Let $\sigma$ be a representation of $U(J)$ on a Hilbert space $H$, $V\subset H$ a subspace,
  $\pi_0$ a representation of $U(J_0)$ on a finite-dimensional space $W_0$, $(F,e,z)$ a Bessel frame
  and $\chi$ a function on $N$. A \emph{Bessel form} is a bilinear form $\beta$ on $W_0\times V$ such
  that: for every $w$, $\beta(w,\cdot)$ is the restriction of a linear form on $H$ that is continuous
  for a Sobolev norm of $\sigma$; $\beta(\pi_0(h)w,\sigma(g)x)=\beta(w,x)$ whenever $h\in U(J_0)$,
  $g\in U(J)$, $gF=Fh$ and $g$ fixes every $x$ with $F^*Jx=0$; and
  $\beta(w,\sigma(u)x)=\chi(u)\beta(w,x)$ for $u\in N$.
\end{enumerate}
\end{definition}
\leanref{Development/AbelianVarietyTheory/PublishedTheorems/TwistedTraceFormula/BesselFormsTrivialNeBot.lean}{GlobalTheory.RealBessel.IsFrame}
\leanref{Development/AbelianVarietyTheory/PublishedTheorems/TwistedTraceFormula/BesselFormsTrivialNeBot.lean}{GlobalTheory.RealBessel.unipRad}
\leanref{Development/AbelianVarietyTheory/PublishedTheorems/TwistedTraceFormula/BesselFormsTrivialNeBot.lean}{GlobalTheory.RealBessel.IsCharWith}
\leanref{Development/AbelianVarietyTheory/PublishedTheorems/TwistedTraceFormula/BesselFormsTrivialNeBot.lean}{GlobalTheory.RealBessel.IsGeneric}
\leanref{Development/AbelianVarietyTheory/PublishedTheorems/TwistedTraceFormula/BesselFormsTrivialNeBot.lean}{GlobalTheory.RealBessel.IsEmb}
\leanref{Development/AbelianVarietyTheory/PublishedTheorems/TwistedTraceFormula/BesselFormsTrivialNeBot.lean}{GlobalTheory.RealBessel.IsBesselForm}
\leanref{Development/AbelianVarietyTheory/PublishedTheorems/TwistedTraceFormula/BesselFormsTrivialNeBot.lean}{GlobalTheory.RealBessel.besselForms}

\begin{lemma}\label{asm-dependence-rows-bessel-unique}
Let $(F,e,z)$ be a Bessel frame for $(J,J_0)$ (Definition~\ref{asm-dependence-rows-bessel-forms})
with $m\ge1$, $z^*Jz<0$ and $J_0$ positive definite, and $\chi$ a generic character of its unipotent
group. Let $\sigma$ be an irreducible unitary representation of $U(J)$ on a Hilbert space $H$ with
smooth vectors $H^\infty$. Suppose that there are signs $\epsilon\in\{\pm1\}^n$, an invertible matrix
$f$ with $f^*J_\epsilon f=J$, where $J_\epsilon$ is the diagonal matrix with entries $\epsilon$, and
a strictly decreasing $\lambda\in\mathbf R^n$ with $\lambda_a-(n-1)/2\in\mathbf Z$ for every $a$ and
$|\lambda_a-\lambda_b|\ge h$ whenever $\epsilon_a\ne\epsilon_b$, for some $h\ge120n^3$, such that
$g\mapsto\sigma(f^{-1}gf)$ is the discrete series of $U(J_\epsilon)$ with Harish-Chandra parameter
$(\lambda,\epsilon)$. Let $\pi_0$ be an irreducible representation of $U(J_0)$ on a
finite-dimensional space $W_0$ with continuous matrix coefficients. Then the Bessel forms on
$W_0\times H^\infty$ form a space of dimension at most $1$.
\end{lemma}
\leanref{Development/AutomorphicArithmetic/ThetaCorrespondence/BesselModels/BurGapOfHyp.lean}{UpperTheory.Lit.Bur.BurGapPos}
\leanref{Development/AutomorphicArithmetic/ThetaCorrespondence/BesselModels/BurGapOfHyp.lean}{UpperTheory.Lit.Bur.IsGapDS}
\leanref{Development/AutomorphicArithmetic/ThetaCorrespondence/BesselModels/Separate.lean}{UpperTheory.Lit.burGapPos}
\leanref{CompanionTheory/AutomorphicTheory/BesselOrbitBounds/BesselMapUniqueness/BesselMapUniqGapAlt.lean}{CompanionTheory.Bmu.besselMapUniqGap}

This is a case of the uniqueness of Bessel models for real unitary groups \cite{JSZ10}: a Bessel form
is a continuous functional on $W_0\otimes H^\infty$ that transforms under the Bessel subgroup
$U(J_0)\ltimes N$ by the character $1\otimes\chi$.

\begin{lemma}\label{asm-dependence-rows-bessel-line}
Let $D$ be a definite-side datum with set $\mathcal T$ of embeddings, let $V$ and $A$ be Hermitian
spaces over its host field, with $n=\operatorname{rk}V$ and $\operatorname{rk}A\ge1$, and let $M$
be a model of discrete series for $U(V)$ (Definition~\ref{asm-dependence-rows-models}(1)) whose
parameter $(\lambda_\tau,\epsilon_\tau)$ at every $\tau$ has $\lambda_{\tau,a}-(n-1)/2\in\mathbf Z$
for all $a$ and $|\lambda_{\tau,a}-\lambda_{\tau,b}|\ge h_\tau$ whenever
$\epsilon_{\tau,a}\ne\epsilon_{\tau,b}$, for some $h_\tau\ge120n^3$. For every $\tau$ let
$(F_\tau,e_\tau,z_\tau)$ be a Bessel frame for $(G_V(\tau),G_A(\tau))$ with
$z_\tau^*G_V(\tau)z_\tau<0$ and $G_A(\tau)$ positive definite, $\chi_\tau$ a generic character of
its unipotent group, and $\pi_\tau$ an irreducible representation of $U(G_A(\tau))$ on a
finite-dimensional space $W_\tau$ with continuous matrix coefficients. Then there is a function
$\ell$ of families $x=(x_\tau)\in\prod_\tau W_\tau$ and $v=(v_\tau)\in\prod_\tau H_\tau^\infty$ such that
every function $\beta$ on these families that is \emph{separately Bessel} is a multiple of $\ell$:
$\beta$ is separately Bessel if for every $\tau$ and all $x$ and $v$, the form
$(y,u)\mapsto\beta(x',v')$, where $x'$ and $v'$ are $x$ and $v$ with $x_\tau$ and $v_\tau$ replaced by
$y$ and $u$, is a Bessel form for the action of $U(G_V(\tau))$ on $H_\tau$ given by $M$, $\pi_\tau$,
$(F_\tau,e_\tau,z_\tau)$ and $\chi_\tau$.
\end{lemma}
\leanref{Development/AutomorphicArithmetic/ThetaCorrespondence/BesselModels/SepGap.lean}{UpperTheory.Lit.Bur.sepLine_gap}
\leanref{Development/AutomorphicArithmetic/ThetaCorrespondence/BesselModels/SepGap.lean}{UpperTheory.Lit.Bur.burGap_model}
\leanref{Development/AutomorphicArithmetic/ThetaCorrespondence/BesselModels/SepGap.lean}{UpperTheory.Lit.Bur.HasGap}
\leanref{Development/AbelianVarietyTheory/Definitions/TwistedInnerProducts/ParRowZero.lean}{GlobalTheory.SepRows.SepBes}
\begin{proof}
Fix $\tau$, and let $(H_\tau,\lambda_\tau,\epsilon_\tau,\sigma_\tau,f_\tau)$ be the data of $M$ at
$\tau$, so that $U(G_V(\tau))$ acts on $H_\tau$ by $g\mapsto\sigma_\tau(f_\tau gf_\tau^{-1})$. This action
is irreducible and unitary, since $\sigma_\tau$ is, and composing it with $g\mapsto f_\tau^{-1}gf_\tau$
gives $\sigma_\tau$, the discrete series with parameter $(\lambda_\tau,\epsilon_\tau)$, which is
strictly decreasing. So Lemma~\ref{asm-dependence-rows-bessel-unique} applies with $f=f_\tau$: the
Bessel forms at $\tau$ form a space of dimension at most $1$. Let $b_\tau$ span it (or be $0$), and put
$\ell(x,v)=\prod_\tau b_\tau(x_\tau,v_\tau)$.

Let $\beta$ be separately Bessel. For every $\tau$ and every pair $p=(x,v)$, the form of the
definition is a multiple $c\,b_\tau$; so as a function of the pair $(x_\tau,v_\tau)$, with the other
coordinates fixed, $\beta$ is a multiple of $b_\tau$. Changing the coordinates one $\tau$ at a time
shows that $\beta(x,v)=l\prod_\tau b_\tau(x_\tau,v_\tau)$ for a constant $l$: if some $b_\tau$ vanishes
identically then so does $\beta$; otherwise fix $(x^0_\tau,v^0_\tau)$ with $b_\tau(x^0_\tau,v^0_\tau)\ne0$
for every $\tau$, and $l=\beta(x^0,v^0)/\ell(x^0,v^0)$.
\end{proof}

\begin{lemma}\label{asm-dependence-rows-gap}
Let $D$ be a definite-side datum with size data at a scale $x$ whose rows have the sizes of the size
data, with a level system with base tuple and a theta grid with constants $(1,C_F,b)$, and let $M$ be
a model of discrete series for $U(V_{m,0})$ with parameter $S^{\mathrm{row}}_m$
(Definition~\ref{asm-dependence-rows-parameters}), $n=\operatorname{rk}V_{m,0}$. Then at every
$\tau$ the parameter $(\lambda_\tau,\epsilon_\tau)$ of $M$ has $\lambda_{\tau,a}-(n-1)/2\in\mathbf Z$
for every $a$, and $|\lambda_{\tau,a}-\lambda_{\tau,b}|\ge120n^3$ whenever
$\epsilon_{\tau,a}\ne\epsilon_{\tau,b}$.
\end{lemma}
\leanref{Development/AutomorphicArithmetic/ThetaCorrespondence/BesselModels/GapArith.lean}{UpperTheory.Lit.Bur.Gap.fit_gap}

This is a computation with the signed set $S^{\mathrm{row}}_m$ for $c=1$: its enumerations are
integral by Lemma~\ref{asm-dependence-rows-counts}(2), and the sizes of the rows give the lower bound
for the gaps between the blocks of opposite signs.

\begin{lemma}\label{asm-dependence-rows-coefficients}
Let $D$ be a definite-side datum with set $\mathcal T$ of embeddings, let $\theta$ be an
archimedean family, and let $(\varrho,W)$ be the coefficients of a theta grid built with $\theta$
(Definition~\ref{asm-grid-evaluation}(5)). There are irreducible representations $\pi_\tau$ of $U(G_{Q_0}(\tau))$ on
finite-dimensional spaces $W_\tau$, $\tau\in\mathcal T$, with continuous matrix coefficients, and an
isomorphism $e\colon\bigotimes_\tau W_\tau\to W^*$ with
$e(\otimes_\tau y_\tau)=e(\otimes_\tau x_\tau)\circ\varrho(\theta_\tau(g))^{-1}$ whenever
$g\in U(G_{Q_0}(\tau))$, $y_\tau=\pi_\tau(g)x_\tau$ and $y_{\tau'}=x_{\tau'}$ for $\tau'\ne\tau$.
\end{lemma}
\leanref{Development/AbelianVarietyTheory/Definitions/TwistedInnerProducts/ParRowZero.lean}{GlobalTheory.SepRows.CoefTensor}
\leanref{Development/AbelianVarietyTheory/Proofs/TwistedOperators/DualCoefSpaceIsTensorOverRealPlaces.lean}{GlobalTheory.CoefTensorProof.coefTensor}
\leanref{Development/AbelianVarietyTheory/Proofs/TwistedOperators/DualCoefSpaceIsTensorOverRealPlacesStatement.lean}{GlobalTheory.T.dualCoefSpace_isTensorOverRealPlaces}
\leanref{Development/AbelianVarietyTheory/Proofs/TwistedOperators/DualIrreducibleCoefficientSpaceRealPlacesStatement.lean}{GlobalTheory.T.dual_irreducible_coefficientSpace_realPlaces}
\leanref{Development/AbelianVarietyTheory/Proofs/TwistedOperators/ExistsTensorDecompositionCommutingFamiliesStatement.lean}{GlobalTheory.T.exists_tensorDecomposition_commutingFamilies}
\begin{proof}
For $g\in U(G_{Q_0}(\tau))$ let $a_\tau(g)$ be the operator $f\mapsto f\circ\varrho(\theta_\tau(g))^{-1}$
on $W^*$. The images under $\theta$ of the groups at distinct embeddings commute, so $a_\tau(g)$ and
$a_{\tau'}(g')$ commute for $\tau\ne\tau'$. The space $W^*$ is irreducible under all the $a_\tau(g)$
together. A finite-dimensional complex space that is irreducible under commuting families of
operators is the tensor product of spaces $W_\tau$, irreducible under the $\tau$-th family, embedded
in it compatibly with the actions. This gives $W_\tau$, $\pi_\tau$ and $e$; the matrix coefficients of
$\pi_\tau$ are matrix coefficients of the contragredient of $\varrho\circ\theta_\tau$, which are
continuous. The irreducibility of $W^*$ under the real places and the continuity of these
coefficients are part of the open step in the proof of Theorem~\ref{asm-dependence-inputs}.
\end{proof}

\begin{lemma}\label{asm-dependence-rows-bessel-period}
Let a row setting be given (Definition~\ref{asm-dependence-rows-models}(3)), with archimedean family
$\theta$, and let $\iota\colon U(Q_0)\to U(V_{m,0})$, $N$ and $\psi_N$ be as in
Definition~\ref{asm-grid-evaluation}(5). Call a function on $U(X)(\mathbf A)$ \emph{admissible} if it
is continuous and invariant on the left under $U(X)(F)$.
\begin{enumerate}
\item For every $\tau$, the matrix of $\iota$ at $\tau$, the flag vectors of the Bessel datum of the
  grid and its last vector $z_\tau$ form a Bessel frame for $(G_{V_{m,0}}(\tau),G_{Q_0}(\tau))$
  (Definition~\ref{asm-dependence-rows-bessel-forms}) with $z_\tau^*G_{V_{m,0}}(\tau)z_\tau<0$;
  $G_{Q_0}(\tau)$ is positive definite; and $u\mapsto\psi_N(\theta_\tau(u))$ is a generic character of
  its unipotent group, whose image under $\theta_\tau$ lies in $N(\mathbf A)$.
\item There is a compact set $\Omega\subset U(V_{m,0})(\mathbf A)$ that meets $U(V_{m,0})(F)\,n\,\iota(h)$
  for every $n\in N(\mathbf A)$ and $h\in U(Q_0)(\mathbf A)$. For every admissible $\varphi$ on
  $U(Q_0)(\mathbf A)$ there is $C$ such that $B^+(\varphi,\cdot)$ is additive on admissible functions,
  $B^+$ is additive in its first argument on admissible functions, and
  $|B^+(\varphi,G)|\le C\sup_\Omega|G|$ for every admissible $G$.
\item For admissible $\varphi$ and $G$ and $a\in U(Q_0)(\mathbf A)$, the right translates satisfy
  $B^+(\varphi(\cdot\,a),G(\cdot\,\iota(a)))=B^+(\varphi,G)$; and for $u\in N(\mathbf A)$,
  $B^+(\varphi,G(\cdot\,u))=\psi_N(u)B^+(\varphi,G)$.
\item Let $M$ be a model of discrete series for $U(V_{m,0})$, $T\in\operatorname{Mult}_\theta(M)$,
  $\Omega$ compact, $\tau$ an embedding and $v\in M^\infty$. There are a constant $C'$ and finitely
  many products $Y_1\cdots Y_k$ of elements of the Lie algebra at $\tau$ such that
  $|T(v')(g)|\le C'\sum\|Y_1\cdots Y_k\,u\|$ for $g\in\Omega$, $u\in H_\tau^\infty$ and $v'$ the vector
  $v$ with $v_\tau$ replaced by $u$.
\end{enumerate}
\end{lemma}
\leanref{Development/AutomorphicArithmetic/ThetaCorrespondence/BesselModels/SepGap.lean}{UpperTheory.Lit.Bur.lastVertex_besselDatum_isFrame_isGeneric}
\leanref{Development/AutomorphicArithmetic/ThetaCorrespondence/BesselModels/SepGap.lean}{UpperTheory.Lit.Bur.topCfg}
\leanref{Development/AbelianVarietyTheory/Proofs/TwistedOperators/ExistsCompactMeetsBesselTopRationalOrbitsStatement.lean}{GlobalTheory.T.exists_compact_meets_besselTopRationalOrbits}
\leanref{Development/AbelianVarietyTheory/Proofs/TwistedOperators/BesselPeriodCorankOneAdditiveBoundedCompactStatement.lean}{GlobalTheory.T.besselPeriodCorankOne_additive_bounded_compact}
\leanref{Development/AbelianVarietyTheory/Proofs/TwistedOperators/BesselPeriodCorankOneSimultaneousRightTranslateStatement.lean}{GlobalTheory.T.besselPeriodCorankOne_simultaneous_rightTranslate}
\leanref{Development/AbelianVarietyTheory/Proofs/TwistedOperators/BesselPeriodCorankOneRightTranslateUnipotentStatement.lean}{GlobalTheory.T.besselPeriodCorankOne_rightTranslate_unipotent}
\leanref{Development/AbelianVarietyTheory/Proofs/DecoratedLatticeProofs/BesselPeriodRightTranslateSobolevSepContinuous.lean}{GlobalTheory.Core.Mi.SepPer.mult_bound}
\leanref{Development/AbelianVarietyTheory/Proofs/TwistedOperators/IsBesselFrameOfBesselDatumRealPlaceStatement.lean}{GlobalTheory.T.isBesselFrame_of_besselDatum_realPlace}
\leanref{Development/AbelianVarietyTheory/Proofs/TwistedOperators/BesselCharacterRealPlaceUnipotentEqTraceStatement.lean}{GlobalTheory.T.besselCharacter_realPlace_unipotent_eq_trace}
\leanref{Development/AbelianVarietyTheory/Proofs/TwistedOperators/StandardCharTraceSingleRealPlaceStatement.lean}{GlobalTheory.T.standardChar_trace_single_realPlace}
\leanref{Development/AbelianVarietyTheory/Proofs/TwistedOperators/AtRealPlaceMemUnipotentFlagAdelicStatement.lean}{GlobalTheory.T.atRealPlace_mem_unipotentFlag_adelic}
\leanref{Development/AutomorphicArithmetic/ThetaCorrespondence/BesselModels/SepGap.lean}{UpperTheory.Lit.Bur.besselDatum_middleVector_norm_eq}
\leanref{Development/AbelianVarietyTheory/Proofs/IwasawaDecomposition/IsotropicUnipotentRadicalsRationalCocompact.lean}{GlobalTheory.Mem.UnipCocompact.radicals_cocompact}
\leanref{Development/AbelianVarietyTheory/Proofs/TwistedOperators/ExistsCompactMeetsBesselTopRationalOrbits.lean}{GlobalTheory.BesCover.exists_compact_cover}
\leanref{Development/AbelianVarietyTheory/Proofs/TwistedOperators/BesselPeriodCorankOneSimultaneousRightTranslate.lean}{GlobalTheory.BplusShift.inner_left}
\leanref{Development/AbelianVarietyTheory/Proofs/TwistedOperators/BesselPeriodCorankOneRightTranslateUnipotent.lean}{GlobalTheory.BplusUnip.setIntegral_mul_right}
\leanref{Development/AbelianVarietyTheory/Proofs/DecoratedLatticeProofs/CuspFormMapValuesCompactSobolevBound.lean}{GlobalTheory.Core.Mi.sobolev_unif_bound}
\begin{proof}
(1) Fix $\tau$. The matrix of $\iota$, the flag vectors and the last vector of the Bessel datum form a
Bessel frame at $\tau$: this is a computation from the definition of the Bessel datum, whose
embedding is isometric, whose flag vectors span an isotropic subspace orthogonal to the image of
$Q_0$, and whose last vector is orthogonal to both. The last vector has norm $-1$, so
$z_\tau^*G_{V_{m,0}}(\tau)z_\tau=-1<0$. The space $Q_0$ has rank at least $1$ and signature
$(\operatorname{rk}Q_0,0)$ at every embedding (Definition~\ref{th-grid}), so $G_{Q_0}(\tau)$ is
positive definite. Let $u$ lie in the unipotent group of the frame, with
$(u-1)f_{j+1}=\sum_ka_{jk}e_k$ for $1\le j\le r$. Then $\theta_\tau(u)$ lies in $N(\mathbf A)$, and
$\psi_N(\theta_\tau(u))=\psi(\operatorname{tr}t)$ for the adele $t$ of $E$ with finite part $0$,
archimedean components $0$ away from $\tau$ and component $\sum_j\tau(c_j)a_{jj}$ at $\tau$, where the
$c_j\neq0$ are the constants of the Bessel datum; both are computations with the Bessel datum. As $\psi$ is
standard, $\psi(\operatorname{tr}t)=\exp\bigl(i\operatorname{Re}\sum_j(-4\pi\tau(c_j))a_{jj}\bigr)$. So the
character is generic, with the nonzero constants $-4\pi\tau(c_j)$.

(2) The space $Q_0$ is totally definite, hence anisotropic, so $U(Q_0)(\mathbf A)=U(Q_0)(F)K$ for a
compact set $K$ \cite{BHC62}; and $N(\mathbf A)=N(F)K'$ for a compact set $K'$, since $N$ is unipotent.
For $\delta\in U(Q_0)(F)$ the element $\iota(\delta)$ normalizes $N(\mathbf A)$ and $N(F)$ and preserves
$\psi_N$. Put $\Omega=K'\iota(K)$, which is compact. Given $n$ and $h$, write $h=\delta k$ with
$\delta\in U(Q_0)(F)$ and $k\in K$, and $\iota(\delta)^{-1}n\iota(\delta)=\delta'k'$ with $\delta'\in N(F)$
and $k'\in K'$. Then $(\iota(\delta)\delta')^{-1}\,n\,\iota(h)=k'\iota(k)\in\Omega$.

Let $\varphi$ be admissible. It is bounded, by $C_\varphi$ say, since it is continuous and left
invariant and $U(Q_0)(\mathbf A)=U(Q_0)(F)K$. The fundamental domain $\mathcal F_{Q_0}$ has finite
measure, as in the proof of Lemma~\ref{asm-dependence-rows-bessel-conv}, and $\mathcal F_N$ has
measure $1$. Let $G$ be admissible; it is bounded on the compact set $\Omega$, and if $|G|\le K_G$ on
$\Omega$ then $|G(n\iota(h))|\le K_G$ for all $n$ and $h$, by the left invariance of $G$ and the choice of
$\Omega$. Since $|\psi_N|=1$, the integrand $\varphi(h)G(n\iota(h))\overline{\psi_N(n)}$ is continuous
and bounded by $C_\varphi K_G$. So both integrals of $B^+(\varphi,G)$ converge absolutely, and
$|B^+(\varphi,G)|\le C_\varphi\operatorname{vol}(\mathcal F_{Q_0})\sup_\Omega|G|$. Absolutely convergent
iterated integrals are additive in the integrand, so $B^+$ is additive in each argument on admissible
functions. Take $C=C_\varphi\operatorname{vol}(\mathcal F_{Q_0})$.

(3) Let $a\in U(Q_0)(\mathbf A)$. Since $\iota$ is a homomorphism, the integrand of
$B^+(\varphi(\cdot\,a),G(\cdot\,\iota(a)))$ at $(h,n)$ is that of $B^+(\varphi,G)$ at $(ha,n)$. Let
$I(h)$ be the inner integral of $B^+(\varphi,G)$. The Haar measure of $U(Q_0)(\mathbf A)$ is right
invariant, so the outer integral of the translate is the integral of $I$ over $\mathcal F_{Q_0}a$,
which is again a fundamental domain for $U(Q_0)(F)$. The function $I$ is left invariant under
$U(Q_0)(F)$: for rational $\gamma$, the substitution $n\mapsto\iota(\gamma)n\iota(\gamma)^{-1}$ maps $N(F)$
to itself and preserves the measure and $\psi_N$, and
$\varphi(\gamma h)G(\iota(\gamma)n\iota(\gamma)^{-1}\iota(\gamma h))=\varphi(h)G(n\iota(h))$ by the left
invariance of $\varphi$ and $G$; the integral over $\mathcal F_N$ of a function left invariant under
$N(F)$ does not depend on the fundamental domain. The integrals of the left invariant integrable
function $I$ over two fundamental domains agree, which gives the first identity.

For $u\in N(\mathbf A)$ and fixed $h$ put $y=\iota(h)u\iota(h)^{-1}$, which lies in $N(\mathbf A)$ with
$\psi_N(y)=\psi_N(u)$. Then
$G(n\iota(h)u)\overline{\psi_N(n)}=\psi_N(u)\,G(ny\,\iota(h))\overline{\psi_N(ny)}$. The function
$n\mapsto\varphi(h)G(n\iota(h))\overline{\psi_N(n)}$ is left invariant under $N(F)$, since $\psi_N$ is
trivial there, and the Haar measure of $N(\mathbf A)$ is right invariant, since $N(F)$ is a discrete
cocompact subgroup; so its integral over $\mathcal F_N$ does not change under $n\mapsto ny$. The inner
integral is therefore multiplied by $\psi_N(u)$, and so is $B^+$.

(4) Write $T$ as a finite linear combination of multiplicity maps; bounds of the stated form add, so
it suffices to treat one multiplicity map $T$. For $g\in\Omega$ the linear form
$\ell_g\colon u\mapsto T(v')(g)$ on $H_\tau^\infty$ is continuous for a Sobolev norm, by
Definition~\ref{asm-dependence-rows-models}(2). The smooth vectors, with the seminorms
$u\mapsto\|Y_1\cdots Y_ku\|$, form a Fr\'echet space: the families of derivatives of smooth vectors
form a closed subspace of a product of Hilbert spaces. For each $u$ the numbers $\ell_g(u)$,
$g\in\Omega$, are bounded, because $T(v')$ is continuous and $\Omega$ is compact. By the
Banach-Steinhaus theorem \cite[Chapter~2]{Rud91} the seminorm $u\mapsto\sup_{g\in\Omega}|\ell_g(u)|$ is
continuous, hence bounded by a constant times a finite sum of the defining seminorms. These are
derivatives along elements of the Lie algebra of the unitary group of $J_\tau$; conjugating by $f_\tau$
turns them into derivatives along elements of the Lie algebra at $\tau$ for the action of the model.
\end{proof}

\begin{lemma}\label{asm-dependence-rows-bessel-separate}
Assume the results of Theorem~\ref{lit-bundle-theta}. Let a row setting be given, with Weil datum
archimedean family $\theta$ and row volume, let $M$ be a model of discrete series for $U(V_{m,0})$,
$T'\in\operatorname{Mult}_\theta(M)$ and $\Psi$ a column. Then
$(\xi,v)\mapsto B^+(\varphi_{\xi,\Psi},T'(v))$ is bilinear on $W^*\times M^\infty$. For every family
$(W_\tau,\pi_\tau)$ and isomorphism $e$ as in Lemma~\ref{asm-dependence-rows-coefficients}, the
function $(x,v)\mapsto B^+(\varphi_{e(\otimes x),\Psi},T'(\otimes v))$ is separately Bessel
(Lemma~\ref{asm-dependence-rows-bessel-line}) for the Bessel frames and generic characters of
Lemma~\ref{asm-dependence-rows-bessel-period}(1) and the representations $\pi_\tau$.
\end{lemma}
\leanref{Development/AbelianVarietyTheory/Definitions/TwistedInnerProducts/ParRowZero.lean}{GlobalTheory.SepRows.TopSep}
\leanref{Development/AutomorphicArithmetic/ThetaCorrespondence/ThetaPairLifts/LocalTypeDictionary/ThetaLiftsSeparateOfSeparatingData.lean}{GlobalTheory.SepRows.SepProof.theta_lifts_separate_of_separating_data}
\leanref{Development/AutomorphicArithmetic/ThetaCorrespondence/ThetaPairLifts/LocalTypeDictionary/ThetaLiftsSeparateOfSeparatingData.lean}{GlobalTheory.T.theta_lifts_separate_points}
\begin{proof}
Let $\Omega$ be the compact set of Lemma~\ref{asm-dependence-rows-bessel-period}(2). Each
$\varphi_{\xi,\Psi}$ is admissible, since $\Psi(\gamma g)=\varrho(\gamma)\Psi(g)$ for rational $\gamma$,
and it is linear in $\xi$. Each $T'(v)$ is admissible, being a finite sum of cuspidal automorphic
forms, and it is linear in $v$. So by part (2) of that lemma the function
$(\xi,v)\mapsto B^+(\varphi_{\xi,\Psi},T'(v))$ is additive in both variables; it is homogeneous in
both, since $B^+$ is an integral. So it is bilinear.

Fix $\tau$, $x$ and $v$, and let $\beta(y,u)$ be the value at $x$ and $v$ with the $\tau$-th
coordinates replaced by $y$ and $u$; it is bilinear. \emph{Continuity.} For each $y$, part (4) of
that lemma bounds $T'(v')$ on $\Omega$ by a Sobolev norm of $u$, and part (2) bounds $\beta(y,u)$ by a
constant times $\sup_\Omega|T'(v')|$; the linear form $u\mapsto\beta(y,u)$ extends to a linear form
on $H_\tau$, bounded by the same Sobolev norm. \emph{Invariance.} Let $h\in U(G_{Q_0}(\tau))$ and
$g\in U(G_{V_{m,0}}(\tau))$ with $gF_\tau=F_\tau h$ and $g$ fixing the orthogonal complement of the
image of $F_\tau$. Then $\theta_\tau(g)=\iota(\theta_\tau(h))$. By the property of $e$,
$\varphi_{e(\otimes x'),\Psi}$ for $x'_\tau=\pi_\tau(h)y$ is the right translate of
$\varphi_{e(\otimes x),\Psi}$ by $\theta_\tau(h)$, since
$\varphi_{\xi\circ\varrho(a)^{-1},\Psi}(k)=\varphi_{\xi,\Psi}(ka)$. A multiplicity map is equivariant,
so replacing $u$ by its translate under $g$ replaces $T'(v')$ by its right translate by
$\theta_\tau(g)$. Part (3) of that lemma gives
$\beta(\pi_\tau(h)y,g\cdot u)=\beta(y,u)$. \emph{Equivariance.} For $u'$ in the unipotent group at
$\tau$, $\theta_\tau(u')\in N(\mathbf A)$, and part (3) gives
$\beta(y,u'\cdot u)=\psi_N(\theta_\tau(u'))\beta(y,u)$. So $\beta$ is a Bessel form.
\end{proof}

\begin{lemma}\label{asm-dependence-rows-global}
Let $\theta$ denote the archimedean family of the grid Weil datum.
\begin{enumerate}
\item (Convergence.) Assume the results of Theorem~\ref{lit-bundle-theta}. Let a gap row setting be
  given (Definition~\ref{asm-dependence-rows-models}(5)) and let $M_0$ be a row model at $0$. Let
  $T$ be a row multiplicity of $M_0$, $w\in M_0^\infty$, $\rho$ a real vector and $\Phi$ a coset
  datum. Then the partial row lifts $\theta^{(n)}_{\rho,\Phi}(T(w))$, $n\le m$, are rapidly
  decreasing; and for every column $\Psi$ and every linear functional $\xi$ on the coefficients, the
  iterated theta lift $\Theta_{\rho,\Phi}(T(w))$ and the Bessel period
  $B^+\bigl(\varphi_{\xi,\Psi},\Theta_{\rho,\Phi}(T(w))\bigr)$ converge
  (Definition~\ref{asm-grid-evaluation}(4) and~(5)).
\item (Core periods.) Assume the description of irreducible unitary representations of $U(p,q)$ by
  their Harish-Chandra modules, the Casselman-Wallach theorem (Theorem~\ref{lit-casselman-wallach})
  and rapid decay of cusp forms, which are among the results of Theorem~\ref{lit-bundle-theta}. For
  every $\varepsilon$ and every $C_F\ge6510$ there is $x_0$ with the following property. Let $D$ be
  a definite-side datum with size data of exponent $\varepsilon$ at a scale $x\ge x_0$ whose rows
  have the sizes of the size data, with a level system with base tuple, a theta grid with constants
  $(1,C_F,e^{x^\varepsilon})$, a grid Weil datum and a row volume. Let $M_0$ be a model of discrete
  series for $U(Y_0)$ with parameter $S^{\mathrm{row}}_0$ and $M^P$ one for $U(P_0)$ with parameter
  $S^{\mathrm{core}}_0$. For every nonzero row multiplicity $T$ of $M_0$ and every nonzero
  $w\in M_0^\infty$ there are a core multiplicity $y$ of $M^P$ and $v\in(M^P)^K$ with
  $\operatorname{per}(y(v),T(w))\ne0$.
\item (Core surjectivity.) Assume the results of Theorem~\ref{lit-bundle-theta}, Howe duality for
  unramified data (Lemma~\ref{asm-dependence-howe-unramified}) and the Siegel-Weil formula
  (Theorem~\ref{lit-siegel-weil}). Let a gap row setting be given, with Haar measures and
  fundamental domains on the groups $U(P_i)(\mathbf A)$, and let $M^P$ be a core model at $0$. For
  every core multiplicity $y$ of $M^P$ and every $v\in(M^P)^K$, the function $y(v)$ is a finite
  linear combination of reachable functions.
\item (Bessel functionals at the real places.) Assume the results of Theorem~\ref{lit-bundle-theta}.
  Let $D$ be a definite-side datum with size data at a scale $x$ whose rows have the sizes of the
  size data, with a level system with base tuple, a theta grid with constants $(1,C_F,b)$, a grid
  Weil datum and a row volume. Let $M$ be a model of discrete series for $U(V_{m,0})$ with parameter
  $S^{\mathrm{row}}_m$. There is a function $\ell$ of a linear functional $\xi$ on the coefficients
  and a vector $v\in M^\infty$ such that for every $T'\in\operatorname{Mult}_\theta(M)$ and every
  column $\Psi$ there is $l\in\mathbf C$ with
  \[
   B^+\bigl(\varphi_{\xi,\Psi},T'(v)\bigr)=l\,\ell(\xi,v)\qquad\text{for all }\xi\text{ and }v.
  \]
\end{enumerate}
\end{lemma}
\leanref{Development/AutomorphicArithmetic/ThetaCorrespondence/ThetaPairLifts/LocalTypeDictionary/RapidDecayOfSeparatingFamily.lean}{GlobalTheory.T.rapid_decay_of_separating_family}
\leanref{Development/AbelianVarietyTheory/Definitions/TwistedInnerProducts/OffQOfSizeRows.lean}{GlobalTheory.SepRows.RowRapid}
\leanref{Development/AutomorphicArithmetic/ThetaCorrespondence/ThetaPairLifts/LocalTypeDictionary/ConvergenceOfSeparatingFamily.lean}{GlobalTheory.T.convergence_of_separating_family}
\leanref{Development/AbelianVarietyTheory/Definitions/TwistedInnerProducts/ParRowZero.lean}{GlobalTheory.SepRows.RowConv}
\leanref{Development/AutomorphicArithmetic/ThetaCorrespondence/RealBands/BranchPeriod.lean}{UpperTheory.Lit.Branch.periodNe_of_eigen}
\leanref{Development/AutomorphicArithmetic/ThetaCorrespondence/FockModels/RealBranch.lean}{UpperTheory.Lit.realInputE_holds}
\leanref{Development/AbelianVarietyTheory/Definitions/TwistedInnerProducts/OffQOfSizeRows.lean}{GlobalTheory.SepRows.PeriodNe}
\leanref{Development/AutomorphicArithmetic/ThetaCorrespondence/ThetaPairLifts/LocalTypeDictionary/TwistCoreSurjectiveOfLocalComponents.lean}{GlobalTheory.T.twist_core_surjective}
\leanref{Development/AbelianVarietyTheory/Definitions/TwistedInnerProducts/OffQOfSizeRows.lean}{GlobalTheory.SepRows.CoreOnto}
\leanref{Development/AutomorphicArithmetic/ThetaCorrespondence/BesselModels/SepGap.lean}{UpperTheory.Lit.Bur.topLine_gap}
\leanref{Development/AutomorphicArithmetic/ThetaCorrespondence/BesselModels/GapArith.lean}{UpperTheory.Lit.Bur.Gap.fit_gap}
\leanref{Development/AutomorphicArithmetic/ThetaCorrespondence/BesselModels/Separate.lean}{UpperTheory.Lit.burGapPos}
\leanref{Development/AbelianVarietyTheory/Proofs/TwistedOperators/DualCoefSpaceIsTensorOverRealPlacesStatement.lean}{GlobalTheory.T.dualCoefSpace_isTensorOverRealPlaces}
\leanref{Development/AutomorphicArithmetic/ThetaCorrespondence/ThetaPairLifts/LocalTypeDictionary/ThetaLiftsSeparateOfSeparatingData.lean}{GlobalTheory.T.theta_lifts_separate_points}

\begin{proof}
(1) A multiplicity map of level $(K_{\mathrm{row}},1)$ is one with no level condition, so
$T\in\operatorname{Mult}_\theta(M_0)$. Let $n\le m$. Lemma~\ref{asm-dependence-rows-tower} gives a row
model $M_n$ at $n$, a map $v_n$ and $T_n\in\operatorname{Mult}_\theta(M_n)$ such that the partial row
lifts of $T(w)$ with the row data $(\rho,\Phi)$ converge up to $n$ and
$\theta^{(n)}_{\rho,\Phi}(T(w))=T_n(v_n(\rho,w))$. By Lemma~\ref{asm-dependence-rows-cusp} this
function is rapidly decreasing. For $n=m$, the iterated theta lift
$\Theta_{\rho,\Phi}(T(w))=\theta^{(m)}_{\rho,\Phi}(T(w))$ converges, and by
Lemma~\ref{asm-dependence-rows-cusp} it is an automorphic form on $U(V_{m,0})(\mathbf A)$: it is
continuous, invariant on the left under $U(V_{m,0})(F)$ and on the right under an open compact
subgroup of the finite adelic group. Lemma~\ref{asm-dependence-rows-bessel-conv} shows that
$B^+\bigl(\varphi_{\xi,\Psi},\Theta_{\rho,\Phi}(T(w))\bigr)$ converges for every column $\Psi$ and
every $\xi$.

(4) By Lemma~\ref{asm-dependence-rows-sizes}, $p_{de}$ and $p'_{de}$ are prime to $Q_Sq$, so the
data form a row setting. Let $n=\operatorname{rk}V_{m,0}$. By Lemma~\ref{asm-dependence-rows-gap},
the parameter of $M$ at every $\tau$ is integral, with gaps at least $120n^3$ between entries of
opposite signs. Let $(W_\tau,\pi_\tau)$ and $e$ be as in Lemma~\ref{asm-dependence-rows-coefficients},
and let $\ell_0$ be the function that Lemma~\ref{asm-dependence-rows-bessel-line} gives for $M$, the
space $A=Q_0$ of rank at least $1$, the Bessel frames and generic characters of
Lemma~\ref{asm-dependence-rows-bessel-period}(1) and the representations $\pi_\tau$. The vectors
$e(\otimes_\tau x_\tau)$ span $W^*$, since $e$ is an isomorphism, and the tensors $\otimes_\tau v_\tau$
span $M^\infty$.

If $B^+(\varphi_{\xi,\Psi},T'(v))=0$ for all $T'$, $\Psi$, $\xi$ and $v$, take $\ell=0$. Otherwise fix
$T'_0$, $\Psi_0$, $\xi_0$ and $v_0$ with $B^+(\varphi_{\xi_0,\Psi_0},T'_0(v_0))\ne0$, and put
$\ell(\xi,v)=B^+(\varphi_{\xi,\Psi_0},T'_0(v))$. Let $T'$ and $\Psi$ be given, and put
$b(\xi,v)=B^+(\varphi_{\xi,\Psi},T'(v))$ and $b_0=\ell$. By
Lemma~\ref{asm-dependence-rows-bessel-separate}, $b$ and $b_0$ are bilinear, and
$(x,v)\mapsto b(e(\otimes x),\otimes v)$ and $(x,v)\mapsto b_0(e(\otimes x),\otimes v)$ are
separately Bessel. By Lemma~\ref{asm-dependence-rows-bessel-line} there are $l_1$ and $l_0$ with
$b(e(\otimes x),\otimes v)=l_1\ell_0(x,v)$ and $b_0(e(\otimes x),\otimes v)=l_0\ell_0(x,v)$ for all
$x$ and $v$. Since $b_0\ne0$ and both sides are bilinear, $l_0\ne0$. The bilinear forms $b$ and
$(l_1/l_0)b_0$ agree on spanning sets, so they are equal:
$B^+(\varphi_{\xi,\Psi},T'(v))=(l_1/l_0)\,\ell(\xi,v)$ for all $\xi$ and $v$.

The proofs of (2) and (3) are not yet written out; they are part of the open step in the proof of
Theorem~\ref{asm-dependence-inputs}. Part (2) rests on an analysis at the real place of the period of
a cusp form against a theta function; part (3) on the local components of the theta lifts between
the core spaces.
\end{proof}

\begin{lemma}\label{asm-dependence-rows-good}
Let $D$ be a definite-side datum with a family of definite forms for a finite group, with base tuple
$\mu_0$, and let $R$ be a realization of $D$ for it whose level characters are those of $\mu_0$
(Definition~\ref{is-realisation}), with archimedean family $\theta_R$ and model $M_0$ of $U(Y_0)$.
Then every theta multiplicity of $R$ at the trivial character, read as a linear map on $M_0^\infty$,
is a row multiplicity of $M_0$ for $\theta_R$ (Definition~\ref{asm-dependence-rows-models}(4)).
\end{lemma}
\leanref{Development/AbelianVarietyTheory/Proofs/TwistedOperators/GoodRowsPrincipalLevelAreRowsStatement.lean}{GlobalTheory.T.goodRows_principalLevel_are_rows}
\leanref{Development/AbelianVarietyTheory/Proofs/TwistedOperators/GoodRowsPrincipalLevelAreRows.lean}{GlobalTheory.SepRows.goodRows}
\leanref{Development/AbelianVarietyTheory/Proofs/TwistedOperators/GoodRowsPrincipalLevelAreRows.lean}{GlobalTheory.SepRows.Good.levRow_sub}
\leanref{Development/AbelianVarietyTheory/Proofs/TwistedOperators/GoodRowsPrincipalLevelAreRows.lean}{GlobalTheory.SepRows.Good.principalRowLevelCharacter_trivial}
\leanref{Development/AbelianVarietyTheory/Definitions/TwistedInnerProducts/ParRowZero.lean}{GlobalTheory.SepRows.GoodRows}
\begin{proof}
Let $J_1$ and $\chi_1$ be the level of $R$ at $Y_0$ for the trivial subgroup and its character. The
theta multiplicities of $R$ at the trivial character lie in
$\operatorname{Mult}_{\theta_R}(M_0;J_1,\chi_1)$. The character $\chi_1$ is trivial. Since the level
characters are those of $\mu_0$, every place of $P_g$ divides $g_0$ or $g_1$, so the principal
congruence subgroup modulo $g_0g_1$ of the stabilizer of the lattice of $Y_0$ lies in $J_1$:
$K_{\mathrm{row}}\subset J_1$. Hence every multiplicity map of level $(J_1,\chi_1)$ is one of level
$(K_{\mathrm{row}},1)$, and the span of the first lies in that of the second.
\end{proof}

\begin{lemma}\label{asm-dependence-rows-factor}
Assume the results of Theorem~\ref{lit-bundle-theta}. Let $D$ be a definite-side datum with size
data at a scale $x$ whose rows have the sizes of the size data, with a level system with base tuple,
a theta grid with constants $(1,C_F,b)$, a grid Weil datum and a row volume. Let $M_0$ be a row
model at $0$ (Definition~\ref{asm-dependence-rows-models}(3)) and $w\in M_0^\infty$. There is a
function $a$ of a real vector $\rho$ and a linear functional $\xi$ on the coefficients such that for
every row multiplicity $T$ of $M_0$, every column $\Psi$ and every coset datum $\Phi$ there is
$l\in\mathbf C$ with
\[
 B^+\bigl(\varphi_{\xi,\Psi},\Theta_{\rho,\Phi}(T(w))\bigr)=a(\rho,\xi)\,l
 \qquad\text{for all }\rho\text{ and }\xi.
\]
\end{lemma}
\leanref{Development/AutomorphicArithmetic/ThetaCorrespondence/BesselModels/SepGap.lean}{UpperTheory.Lit.Bur.factor_gap}
\leanref{Development/AbelianVarietyTheory/Proofs/TwistedOperators/FactorisationFromTopLine.lean}{GlobalTheory.SepRows.factor_multT_le_empty}
\begin{proof}
By Lemma~\ref{asm-dependence-rows-sizes} the data form a gap row setting.
Lemma~\ref{asm-dependence-rows-tower} with $n=m$ gives a row model $M_m$ at $m$ and a map $v_m$, and
Lemma~\ref{asm-dependence-rows-global}(4) applied to $M_m$ gives a function $\ell$. Put
$a(\rho,\xi)=\ell(\xi,v_m(\rho,w))$.

Let $T$, $\Psi$ and $\Phi$ be given. A multiplicity map of level $(K_{\mathrm{row}},1)$ is one with
no level condition, so $T\in\operatorname{Mult}_\theta(M_0)$. Lemma~\ref{asm-dependence-rows-tower}
gives $T_m\in\operatorname{Mult}_\theta(M_m)$ with
$\Theta_{\rho,\Phi}(T(w))=\theta^{(m)}_{\rho,\Phi}(T(w))=T_m(v_m(\rho,w))$ for all $\rho$, and
Lemma~\ref{asm-dependence-rows-global}(4) gives $l$ with
$B^+(\varphi_{\xi,\Psi},T_m(v))=l\,\ell(\xi,v)$ for all $\xi$ and $v$. Taking $v=v_m(\rho,w)$ gives
$B^+(\varphi_{\xi,\Psi},\Theta_{\rho,\Phi}(T(w)))=l\,\ell(\xi,v_m(\rho,w))=a(\rho,\xi)\,l$.
\end{proof}

\begin{lemma}\label{asm-dependence-rows-weak}
Assume the results of Theorem~\ref{lit-bundle-theta}, Howe duality for unramified data
(Lemma~\ref{asm-dependence-howe-unramified}) and the Siegel-Weil formula
(Theorem~\ref{lit-siegel-weil}). For every $\varepsilon$ and every $C_F\ge6510$ there is $x_0$ with
the following property. Let $D$ be a definite-side datum with size data of exponent $\varepsilon$ at
a scale $x\ge x_0$ whose rows have the sizes of the size data, with a level system with base tuple,
a theta grid with constants $(1,C_F,e^{x^\varepsilon})$, a grid Weil datum and a row volume, and let
$M_0$ be a row model at $0$. For every nonzero row multiplicity $T$ of $M_0$ and every nonzero
$w\in M_0^\infty$ there are a real vector $\rho$, a linear functional $\xi$ on the coefficients, a
column $\Psi$ and a coset datum $\Phi$ with
$B^+\bigl(\varphi_{\xi,\Psi},\Theta_{\rho,\Phi}(T(w))\bigr)\ne0$.
\end{lemma}
\leanref{Development/AbelianVarietyTheory/Proofs/TwistedOperators/RowSeparationByRealFunctionalsOfChainStepsStatement.lean}{GlobalTheory.T.rowSeparation_by_realFunctionals_of_chain_steps}
\leanref{Development/AbelianVarietyTheory/Proofs/TwistedOperators/RowSeparationByRealFunctionalsOfChainSteps.lean}{GlobalTheory.SepRows.weakSep_of_steps}
\leanref{Development/AbelianVarietyTheory/Proofs/TwistedOperators/RowSeparationByRealFunctionalsOfChainSteps.lean}{GlobalTheory.SepRows.WeakSepPf.exists_reach_of_span}
\leanref{Development/AbelianVarietyTheory/Proofs/TwistedOperators/RowSeparationByRealFunctionalsOfChainSteps.lean}{GlobalTheory.SepRows.WeakSepPf.integrableOn_corePeriod_rapidlyDecreasing}
\leanref{Development/AbelianVarietyTheory/Definitions/TwistedInnerProducts/OffQOfSizeRows.lean}{GlobalTheory.SepRows.WeakSep}
\begin{proof}
Let $x_0$ be the scale of Lemma~\ref{asm-dependence-rows-global}(2) for $\varepsilon$ and $C_F$, and
let $D$ and the rest be as in the lemma, with $T\ne0$ and $w\ne0$. Put $u=T(w)$. By
Lemma~\ref{asm-dependence-rows-sizes} the data form a gap row setting.

Lemma~\ref{asm-dependence-rows-local}(1) gives a core model $M^P$ at $0$. By
Lemma~\ref{asm-dependence-rows-global}(2) there are a core multiplicity $y$ of $M^P$ and
$v\in(M^P)^K$ with $\operatorname{per}(y(v),u)\ne0$. Choose Haar measures and fundamental domains on
the groups $U(P_i)(\mathbf A)$. By Lemma~\ref{asm-dependence-rows-global}(3), $y(v)=\sum_kc_kF_k$
with reachable functions $F_k$.

By Lemma~\ref{asm-dependence-rows-global}(1) all partial row lifts of $u$ are rapidly decreasing; in
particular $u=\theta^{(0)}(u)$ is. A reachable function is the first member of a column chain of
core forms, so it is rapidly decreasing; its pullback by $\iota_0$ is then bounded and continuous,
and its product with $u$ is integrable on $\mathcal F_{Y_0}$. So $\operatorname{per}(\cdot,u)$ is
defined and linear on the span of the reachable functions, and
$\operatorname{per}(y(v),u)=\sum_kc_k\operatorname{per}(F_k,u)$. Hence $\operatorname{per}(F,u)\ne0$
for a reachable $F$. Let $(\Psi_0,\ldots,\Psi_m)$ be a column chain of core forms with $\Psi_0=F$
and $\Psi_m=\varphi_{\xi,\Psi}$.

By Lemma~\ref{asm-dependence-rows-sizes}, $p_{de}$ and $p'_{de}$ are prime to $Q_Sq$. The ascent of
nonvanishing (Theorem~\ref{th-bessel-ascent}) applies to the chain and to $u$, all of whose partial
row lifts are rapidly decreasing, with
$\int_{\mathcal F_{Y_0}}\Psi_0(\iota_0h)u(h)\,dh=\operatorname{per}(F,u)\ne0$. It gives row data
$(\rho,\Phi)$ with $B^+(\Psi_m,\theta^{(m)}_{\rho,\Phi}(u))\ne0$, that is,
$B^+(\varphi_{\xi,\Psi},\Theta_{\rho,\Phi}(T(w)))\ne0$.
\end{proof}

\subsection{Linking the open parameter statements}\label{subsection-asm-open-statements}

\begin{lemma}\label{asm-dependence-orbit-witness}
Let local Langlands maps be given and let $p$ be parameters
(Definition~\ref{asm-dependence-parameters}). Then open parameter statement (a)
(Definition~\ref{asm-open-parameter-statements}) is linked at $p$, and so is the compatibility of the
division data with the periods at every member, which is item (3)(f) of
Definition~\ref{asm-definite-side-hypotheses} and the third part of statement (b).
\end{lemma}
\leanref{Development/AutomorphicArithmetic/ThetaCorrespondence/ThetaPairLifts/ThirdPlaceSeries/HermWCellLastMulMulVecEqOfInWb.lean}{GlobalTheory.Fc.RL.PrimAt.linkedAt_hPA}
\leanref{Development/AutomorphicArithmetic/ThetaCorrespondence/ThetaPairLifts/ThirdPlaceSeries/HermWCellLastMulMulVecEqOfInWb.lean}{GlobalTheory.Fc.RL.PrimAt.herm_wCell_last_mul_mulVec_eq_of_inWb}
\leanref{Development/AutomorphicArithmetic/ThetaCorrespondence/ThetaPairLifts/ThirdPlaceSeries/RowEDcPtLinked.lean}{GlobalTheory.Fc.RL.linkedAt_hDcPtG}
\leanref{Development/AutomorphicArithmetic/ThetaCorrespondence/ThetaPairLifts/ThirdPlaceSeries/RowEDcPtLinked.lean}{GlobalTheory.Fc.RL.rowE_dcPt_linked_alt}
\leanref{Development/AbelianVarietyTheory/Proofs/IwasawaDecomposition/RowConditionsAtCompanionFromReadings.lean}{GlobalTheory.Mem.R3.rowConditions_at_companion_from_readings}
\begin{proof}
Let $Q$ be a slot, and let a bundle be linked by $\Lambda_p$ to a base point of $Q$. Then the bundle
is built from $p$ at that point, with a witness of its construction. The witness contains a twisted
input, which consists of a realization with a face, a family of definite forms, and readings at the
face, among them a reading of the definite columns by the words of the datum. It also contains a
twisted bundle over this input. The twisted bundle records that the realization of the input has primitive rows at its face
and that its family of definite forms has lift-distinct types, and that the division pair, the
indefinite realization and the family of definite forms of the bundle are those built from the
input. So the realization of the indefinite realization of the bundle has primitive rows, and its
family of definite forms has lift-distinct types. The reading of the definite columns gives, for
every word $u$, every tuple $t$ and all coordinate vectors, the equality of periods of item (3)(f); in
particular they give it for the tuples whose character is used. So both properties hold at every
bundle that $\Lambda_p$ links to a base point of $Q$.

Take their conjunction as the condition on bundles in Definition~\ref{is-slot-rows}. A selection
whose bundles satisfy it has both properties at every member. So (a) and item (3)(f) at every member
hold along $(Q,\Lambda_p)$, for every slot $Q$. Taking for the slot of
Definition~\ref{asm-dependence-parameters}(2) the slot of all cores, which every core satisfies,
both properties are linked at $p$.
\end{proof}

\begin{lemma}\label{asm-dependence-orbit-inputs-transfer}
Let local Langlands maps be given, let $p$ be parameters
(Definition~\ref{asm-dependence-parameters}), and let $X$ and $A$ be properties of a bundle at a
definite-side datum. For a class $\mathcal A$ of families and a selection $s$ for $\mathcal A$, say
that \emph{$X$ holds at every member} if the bundle that $s$ selects at every member of every family
of $\mathcal A$, at every scale $x>x_1$, satisfies $X$.
\begin{enumerate}
\item If, eventually in the context, every bundle that $\Lambda_p$ links to a base point of the good
  cores of $p$ satisfies $X$, then ``$X$ holds at every member'' is linked at $p$.
\item Suppose that ``$A$ holds at every member'' is linked at $p$, and that every bundle at the
  datum of a stage built from $p$ that has a witness of its construction from that stage and
  satisfies $A$ also satisfies $X$. Then ``$X$ holds at every member'' is linked at $p$.
\item Let $Q$ be a slot and $\Lambda$ link data. If ``$X$ holds at every member'' holds along
  $(Q,\Lambda)$ (Definition~\ref{is-slot-rows}), then so does the property ``eventually along
  $\mathcal W(\mathcal A)$, eventually along each family, the selected bundle satisfies $X$''.
\end{enumerate}
\end{lemma}
\leanref{Development/AutomorphicArithmetic/ThetaCorrespondence/ThetaPairLifts/ThirdPlaceSeries/RowGood.lean}{GlobalTheory.Fc.RL.LinkedAt.of_row}
\leanref{Development/AbelianVarietyTheory/Definitions/GlobalTwisting/ValidForConstructedPointClassesEventuallyWorking.lean}{GlobalTheory.Fc.RowE.row}
\leanref{Development/AutomorphicArithmetic/ThetaCorrespondence/ThetaPairLifts/ThirdPlaceSeries/JoinLinkedRealization.lean}{GlobalTheory.Fc.JoinLinked.linked_atAll_of_imp}
\leanref{Development/AutomorphicArithmetic/ThetaCorrespondence/ThetaPairLifts/ThirdPlaceSeries/JoinLinkedRealization.lean}{GlobalTheory.Fc.JoinLinked.of_built_of_bWit}
\leanref{Development/AbelianVarietyTheory/Proofs/DivisorProofs/RowOfCovolumeOntoGrid.lean}{GlobalTheory.Fc.FourBoundsAtBundles.evW_of_atAll}
\begin{proof}
(1) Take for the condition $E$ on bundles of Definition~\ref{is-slot-rows} the property $X$. By
hypothesis, eventually in the context every bundle that $\Lambda_p$ links to a base point of the good
cores of $p$ satisfies $E$. A selection whose bundles are built from these good cores and $\Lambda_p$
and satisfy $E$ has $X$ at every member. So ``$X$ holds at every member'' holds along
$(Q_p,\Lambda_p)$, and so along $(Q_p\cap R,\Lambda_p)$ for the slot $R$ of all cores, by the remark
at the start of the proof of Lemma~\ref{asm-dependence-linked}. Every core satisfies $R$, so the
property is linked at $p$.

(2) Let $R$ be a slot and $E$ a condition on bundles for ``$A$ holds at every member''. Let $s$ be a
selection for a class $\mathcal A$ whose bundles are built from $(Q_p\cap R,\Lambda_p)$ and satisfy
$E$. Then $A$ holds at every member. The bundle selected at a member is linked by $\Lambda_p$ to a
base point of a good core of $p$, so it is built from $p$ at a stage, with a witness of its
construction from that stage; hence it satisfies $X$. So ``$X$ holds at every member'' holds along
$(Q_p\cap R,\Lambda_p)$, with the same $E$, and it is linked at $p$ with the slot $R$.

(3) Let $s$ be a selection for a class $\mathcal A$ whose bundles are built from $(Q,\Lambda)$ and
satisfy the condition $E$ of the hypothesis. For every family of $\mathcal W(\mathcal A)$, of any
exponent, the selected bundle at every member at every scale $x>x_1$ satisfies $X$. This is the
eventual statement with $\varepsilon_1=1$, every largeness threshold $0$, and the scale $x_0=0$ for
each family: every scale $x\ge0$ with $x>x_1$ is a scale $x>x_1$.
\end{proof}
% Lean: (1) is LinkedAt.of_row (the slot fun _ => True, largeBaseProperty_true, Row.strengthen)
% after RowE.row; (2) is linked_atAll_of_imp through of_built_of_bWit and Built.ind; (3) is
% evW_of_atAll (Closed.of_all, Ind.ev_of_atAll), used as Row.imp, and for (6) and (7) of Lemma
% asm-dependence-orbit-inputs in the explicit form <1, one_pos, ... Ind.ev_of_atAll key ...>.

\begin{lemma}\label{asm-dependence-orbit-inputs-comparison}
Let local Langlands maps be given, let $p$ be parameters
(Definition~\ref{asm-dependence-parameters}) and let $Q$ be a slot. Eventually in the context, every
bundle that $\Lambda_p$ links to a base point of $Q$ has a comparison part.
\end{lemma}
\leanref{Development/AutomorphicArithmetic/ThetaCorrespondence/ThetaPairLifts/ThirdPlaceSeries/RowEHasTie.lean}{GlobalTheory.Fc.TieOfLink.hasTie_of_linked}
\begin{proof}
Eventually in the context, two things hold. First, by
Lemma~\ref{asm-dependence-test-assembly-host}, the frames of the spaces $V_i$ and $W_i$ of the datum
of every stage whose base data lie over the context are diagonal. Second, let a stage be built from
$p$, with its witness, and let a bundle at its datum have a witness of its construction from that
stage, at a core whose size data have the scale of the context. Then five identities relate the
witness to the bundle: (i) the scale data of the twist arguments of the witness have, as their
$F_0$, the local terms of the datum at the scale of the context (as in the proof of
Lemma~\ref{asm-dependence-functional-witness-twisted}); (ii) the character data of the realization
of the twisted input (its choices at the level primes and its Hecke characters $\mu_0$, $\eta$,
$\chi_\psi$, $\delta$ and $\delta'$) are those of the realization of the bundle; (iii) the places $P_g$
of the realization of the twisted input are the images of the places of the kernel of the
realization of the bundle; (iv) the levels and level characters of the realization of the twisted
input at level $0$ match, after conjugation, those of the spectral matching of the bundle; and (v)
they are pinned to the local Galois data of the realization of the bundle. The five identities hold
in every context; they are proved below.

Fix such a context. A bundle that $\Lambda_p$ links to a base point of $Q$ is a bundle at the datum
of a stage built from $p$, with a witness of its construction from that stage, and its core has the
scale of the context, since the base point has it. The realization of the twisted input is
\emph{normalized}: its character data are those of the realization of the bundle, by (ii); the frames
of the $V_i$ and $W_i$ are diagonal; its places $P_g$ are the images of the places of the kernel, by
(iii); the point $t_S$ of the real input lies in the finite adelic group of $U(Y_0)$, and conjugation
by it carries the level and the level character at level $0$ onto those of the spectral matching, at
which the rows match, by (iv); and the twisted input records the remaining conditions in its pins
and its reading: the guard, $F$ is the maximal real subfield of $E$, the phases are $1$, the map
$\mathrm{pr}_X$ is the standard one, the group of the realization is pinned at the level primes, and
the scalars $\beta$ are $1$. Two further conditions, on the group isomorphisms and on the readings of
the vectors $b_R$, hold by construction. By (i) and
Lemma~\ref{asm-dependence-functional-witness-twisted}, the realization of the twisted input is
normalized in degree three: its twisted places are the places of that lemma.

The real input carries a recursion of its realization. The compatibility of a normalization with
this recursion consists of the standard character $\psi$ of the reading of the twisted input, which
the reading records, and of the pin (v); so the normalization is compatible with the recursion. The
pins of the twisted input tie the level groups and characters $(P_0,\psi_0)$ and $(P_s,\psi_s)$ of its
division pair to the group of the datum. A comparison part of a bundle consists of a recursion of
its indefinite realization, the tie of the characters of its division pair, the normalization in
degree three and a normalization compatible with the recursion. The witness identifies the real
package of the bundle with that of its real input, so these data give a comparison part of the
bundle.

It remains to prove the five identities. Let the stage, its witness, the bundle and its witness be
as above, with the scale of the context. The witness records that the realization of its twisted
input is the realization built at the stage, with the scalars $\lambda$, the error terms and the scale
data of its twist arguments, and that its real input has the real package of the bundle. It also pins
the realization of the bundle: its columns are those of the stage, its local Galois data are those of
the core at the weighted point of the stage, its character data are those of the column choice, and
its kernel is the standard kernel of the column choice.

(i) is the first paragraph of the proof of Lemma~\ref{asm-dependence-functional-witness-twisted}.

(ii) The level data of the realization built at the stage record that its character data are those
that the column choice gives at $Z$. The realization of the twisted input has the character data of
the realization built at the stage, and the witness records that the character data of the
realization of the bundle are those of the column choice.

(iii) The witness records that the kernel of the realization of the bundle is the standard kernel of
the column choice, with its places, its seam and its map $\mathrm{Pr}$. The level data of the
realization built at the stage record that its places $P_g$ are the images of the places of this
kernel under the map of places of the seam. The realization of the twisted input has these places.

(iv) The stage carries a saturation of its row package at its weighted point, with a point $t$ of
$U(Y_0)(\mathbf A)$. Its seeing condition says that $t$ lies in the finite adelic group, that
conjugation by $t$ carries the level and the level character at level $0$ of the realization built at
the stage onto those of the Hida setting of the saturation, and that the rows match at $t$. The
witness records that $t_S=t$, and the columns of the bundle are those of the stage, so the Hida
setting of the spectral matching of the bundle is that of the saturation, with the same
identification of groups. Replacing the scalars, the error terms and the scale data does not change
the levels and level characters at level $0$, or the matching of the rows. So (iv) holds for the
realization of the twisted input.

(v) The core of the stage, at its weighted point, carries row data with the local row pin of the
realization built at the stage for the local Galois data of the core. The witness records that these
are the local Galois data of the realization of the bundle. The realization of the twisted input has
the same archimedean family, row characters, models, levels, level characters and row frame as the
realization built at the stage, so the pin passes to it.
\end{proof}
% Lean: hasTie_of_linked (RowEHasTie.lean): EvP of bWit_zsc_F0_heq_F0K_local_of_sc_eq,
% bWit_C_eq_cf_of_sc_eq, bWit_Pg_eq_image_plMap_of_sc_eq, theta_level_conjugate_character_alt,
% local_level_character_pin_of_witness and Mem.HostFrames.hostHyp_frames_ev; then
% R3.realPins_realTOf_alt (the normalization) and TieOfLink.nonempty_tie_of_wit. In the text, (i) to
% (v) are the five EvP statements in this order (F0 of w.a.Zsc = F0K, w.x.R.C = cf (a CharFrame: P, W,
% mu0, eta, chiPsi, del, del'), w.x.R.Pg = image
% of vGs under plMap, tS in GfC with ConjLevel and RowLevelMatchAt, RowLocPinAlt), all at
% q.O.L.sc = K.s; the diagonal frames are the third component of hostHyp_frames_ev. realPins_realTOf_alt
% takes (ii), the frames, (iii) and the three parts of (iv), with qGuard, hF, phs, prX, grp and beta
% from x.pins and x.Y; nonempty_tie_of_wit takes (i) (realPins3_realTOf_alt, degree three), the
% normalization and (v), with recursion_local_row_condition_alt (RecPinsTie.main_alt_alt: compatibility,
% PartsRow, PsiPPTie), and tie_of_rpk_alt along w.y.eq. The five identities: (i)
% bWit_zsc_F0_heq_F0K_of_sc_eq (forcedContent_eq_closedContent_numericalDatum = ofRecord_Fj_succ, with
% b.R.T.what2.ka, along y.eq and the pack of w.a); (ii) bwit_cf_eq_alt (w.recPin.s.chr and
% q.x.what2.r.z.lev.C); (iii) PgTie.pgEq_of_linked_alt (w.recPin.k.std and lev.Pg); (iv)
% LinkLev.tS_fin_conj_match_alt (q.saturationAtWeightedPoint.see: fin, conj, mat; w.pt.1 for tS;
% w.recPin.s.cols for the Hida setting; levels_of_twArgs = levels_withErr_updS); (v)
% locPin_level_character_of_bWit (w.recPin.s.gal, q.core_data_of_point.row.rowLoc,
% rowLocPin_twArgs_alt = rowLocPin_congr_alt). The normalization is the structure
% IndefiniteRealisationNormalisation (cfEq, qGuard, hF, phsOne, prXunit, grpRecPin, fdiag, pgEq, tSfin,
% tSpin, eGtie = rfl, betaStd, bR_read = rfl); the compatibility is main_alt_alt's <x.Y.stdPsi, h>;
% PartsRow is not used. Size of the step written out here: proof closure of hasTie_of_linked, about
% 49 theorems and 710 lines in 28 files.

\begin{lemma}\label{asm-dependence-orbit-inputs-join}
Assume the covering of discrete series, the finiteness of Harish-Chandra modules, the
Casselman-Wallach theorem, Harish-Chandra's finiteness theorem for automorphic forms on unitary
groups and rapid decay of cusp forms (Theorem~\ref{lit-bundle-unitary}). Let local Langlands maps be
given and let $p$ be parameters (Definition~\ref{asm-dependence-parameters}). Then every bundle at
the datum of a stage built from $p$, with a witness of its construction from that stage and a
comparison part, has its indefinite realization joined (Definition~\ref{is-realisation-join}).
\end{lemma}
\leanref{Development/AutomorphicArithmetic/ThetaCorrespondence/ThetaPairLifts/ThirdPlaceSeries/JoinLinkedRealization.lean}{GlobalTheory.Fc.JoinLinked.selJ_of_bWit}
\leanref{Development/AutomorphicArithmetic/ThetaCorrespondence/ThetaPairLifts/ThirdPlaceSeries/JoinLinkedRealization.lean}{GlobalTheory.Fc.JoinLinked.join_linked_realization_alt}
\begin{proof}
The witness identifies the bundle with the real input of its twisted input, so the indefinite
realization of the bundle is the realization of the real input. The comparison part makes the
former normalized, so the latter is normalized. The point data of the stage transfer to the
spectral stage of the witness, as follows. The saturation of the stage at its weighted point gives
spectral data (a spectral package, a restriction, an element $t$ and a number $c$) at which three
conditions hold for the realization built at the stage and its face: the condition on what the
point sees (its finiteness, conjugation and matrix parts), the tie of the evaluations and the
condition on the rows. The twist arguments of the witness give the realization of its twisted input
as this realization with the scalars $\kappa$, $\lambda$ and the error terms replaced, and its face
as this face with the definite columns, periods and scales replaced; the three conditions involve
none of these data, so they hold for the realization of the twisted input at its face. The witness
records that the column forms of the bundle are those of the column data of the stage, and that its
spectral stage, with the element and the number of the real input, is the spectral data of the
saturation read with these forms; so the three conditions hold at the spectral stage of the witness
and the element of the real input. From the five results, the twisted input, the normalization, the
conditions that the twisted bundle records for the realization of the twisted input at its face
(primitive rows, the existence of its theta maps and further conditions, a family ok at level $1$ and
lift-distinct types) and these three conditions, the realization of the real input is joined; so
the indefinite realization of the bundle is joined.
\openstep{The join of the realization of the real input, from the five results, the normalization,
the recorded conditions and the three conditions at the spectral stage, is not yet written out.}
\end{proof}
% Lean: selJ_of_bWit: hPins from t.pins and w.y.eq; StageMove.spectralStage_point_rowPlacement_transfer
% with w.a and q.saturationAtWeightedPoint (s.see, s.ev, s.rowJ; stageAt_upd copies SeeAt fin, conj,
% mat, EvalTieT and RowJ from t.r.z.R, F.lv2 to ((R.updS k lam).withErr e), (reFace F Dc Sp Zsc).lv2,
% along a.pack and a.eE); EntriesCombinedReal.move along hF (w.recPin.s.cols) with w.s (sP, sS, hP,
% hS) and w.pt; then join_linked_realization_alt with w.y.intersectionNumber_baseChange, which is
% R3.Holds of the twisted input (prim, aux = UnusedChoiceStatements with ex, ok, sep; nothing to do
% with base change). Size of the open step, join_linked_realization_alt following imports and
% stopping at statements the book cites: about 28,600 lines of theorems in 408 files.

\begin{lemma}\label{asm-dependence-orbit-inputs-multiplicity}
Assume the following results of Theorem~\ref{lit-bundle-unitary}: the existence of the Weil
representation of a unitary dual pair, Harish-Chandra's finiteness theorem for automorphic forms on
unitary groups, the covering of discrete series, the finiteness of Harish-Chandra modules, rapid
decay of cusp forms, the Casselman-Wallach theorem, the classification of the discrete series,
Theorems~\ref{lit-weil-siegel}, \ref{lit-siegel-weil} and~\ref{lit-doubling}, the unramified
computation of the doubling integral for $GL_n$, strong approximation for special unitary groups,
and Siegel's formula for the mean of the representation numbers of a Hermitian form with the
finiteness of the volume of the automorphic quotient. Assume also the conclusion of
Lemma~\ref{asm-dependence-howe-unramified}. Let local Langlands maps be given and let $p$ be
parameters (Definition~\ref{asm-dependence-parameters}). Eventually in the context, every bundle
that $\Lambda_p$ links to a base point of the good cores of $p$ satisfies the multiplicity condition
of the realization of its indefinite realization at its face.
\end{lemma}
\leanref{Development/AutomorphicArithmetic/ThetaCorrespondence/ThetaPairLifts/ThirdPlaceSeries/MainTheoremPublishedInputsVariant.lean}{GlobalTheory.T.main_theorem_published_inputs_variant_alt_alt}
\leanref{Development/AbelianVarietyTheory/Proofs/DecoratedLatticeProofs/DiagonalRowOfSiegelWeilUnitaryStatement.lean}{GlobalTheory.T.diagonalRow_of_siegelWeil_unitary}
\begin{proof}
Siegel's formula for the mean of the representation numbers of a Hermitian form, with the finiteness
of the volume of the automorphic quotient, gives the diagonal rows of the definite-side data. The
covering of discrete series, the finiteness of Harish-Chandra modules, the Casselman-Wallach theorem,
Harish-Chandra's finiteness theorem and rapid decay of cusp forms give, eventually in the context,
the point conditions of every bundle that $\Lambda_p$ links to a base point of the good cores of $p$.
The witness of every stage built from $p$ records that the rows of the stage are rows of the
construction. From these and the remaining results, eventually in the context, every bundle that
$\Lambda_p$ links to a base point of the good cores of $p$ satisfies the multiplicity condition.
\openstep{The diagonal rows from Siegel's formula, the point conditions of the linked bundles, and the
multiplicity condition at the linked bundles are not yet written out.}
\end{proof}
% Lean: main_theorem_published_inputs_variant_alt_alt = miStarRowClosing_alt_alt, with
% RowDiag (IsScaleWithBoundedGamma and RankMu) from diagonalRow_of_siegelWeil_unitary =
% Core.Rk.rowDiag_of_siegelWeilUnitary, and RXs at every weighted point from q.x.rx ('the rows are
% rows of the construction'); miStarRowClosing_alt_alt applies main_theorem_published_results_final_alt_alt
% with T.point_lines_of_cusp_rapid_decay (RowE of Fam.BundlePointConditions). Sizes (text
% dependencies, overcounts): miStarRowClosing_alt_alt about 37,600 theorems and 607,000 lines;
% rowDiag_of_siegelWeilUnitary about 42,600 theorems and 695,000 lines; they overlap.

\begin{lemma}\label{asm-dependence-orbit-inputs-division}
Assume the covering of discrete series, the finiteness of Harish-Chandra modules, the
Casselman-Wallach theorem, Harish-Chandra's finiteness theorem for automorphic forms on unitary
groups and rapid decay of cusp forms (Theorem~\ref{lit-bundle-unitary}). Let local Langlands maps be
given, let $p$ be parameters (Definition~\ref{asm-dependence-parameters}), and let a bundle be given
at the datum $D(w,Z)$ of a stage built from $p$, with a witness of its construction from that stage.
Suppose that $D(w,Z)$ satisfies the rationality law (Definition~\ref{asm-dependence-parameters}(3))
and that the bundle satisfies the multiplicity condition of the realization of its indefinite
realization at its face. Then the division rows of the bundle hold.
\end{lemma}
\leanref{Development/AutomorphicArithmetic/ThetaCorrespondence/ThetaPairLifts/ThirdPlaceSeries/AtHDivGOfMi.lean}{GlobalTheory.Fc.DivLinked.hDivG_of_mIstar2}
\begin{proof}
The witness identifies the bundle with the real input of its twisted input: the division pair, the
indefinite realization and the realization of the bundle are those of the real input. Through this
identification, the multiplicity condition holds for the realization of the twisted input at its
face. With the conditions that the twisted bundle records for this realization (a family ok at level
$1$, lift-distinct types and primitive rows) and its weak characterization, the multiplicity
condition makes the pairing on the definite base vanish. Indeed, these are the hypotheses of
Lemma~\ref{asm-primitive-rows-base} for this realization at its face, and Step 2 of the proof of
Lemma~\ref{asm-division-kernel-factorization}, which uses that lemma and
Lemma~\ref{asm-division-kernel-evaluation}, gives $\mathrm{Pt}_u(\mu_0)(\varphi,y)=0$ for every word
$u$, every $\varphi$ and every $y$ with $P_{\mu_0}y=0$. With the five results, these recorded
conditions and the rationality law of $D(w,Z)$, the vanishing gives the division rows of the division
pair of the real input, which is that of the bundle.
\openstep{The division rows that the rationality law gives from the vanishing of the pairing on the
definite base are not yet written out.}
\end{proof}
% Lean: hDivG_of_mIstar2: Div.pairing_definite_base_vanish_of_miStar2 (at level 1, with ok, sep, prim
% of intersectionNumber_baseChange = R3.Holds, pins.char = the weak characterization, and w.x.sEE),
% then divisionRows_pairing_vanishes_of_rational_pairing_law hC w.x (Holds) BkLaw, transported along
% w.y.eq; hC is the five results. The vanishing: pairing_definite_base_vanish_of_miStar2 =
% complex_data_conditions_base_level_of_primitive_links (Lemma asm-primitive-rows-base) and
% CxData.base_tuple_vanishing_of_hypotheses (Step 2 of asm-division-kernel-factorization) with
% EvTie.evTie_of_ok and definite_coordinates_in_good_part; its conclusion is CxData.BtVanish. Size of
% the open step, divisionRows_pairing_vanishes_of_rational_pairing_law, following imports and
% stopping at statements the book cites: about 30,300 lines of theorems in 436 files (through
% T.divisionPairing_rational_on_rationalPoints and the theta steps of the row side).

\begin{lemma}\label{asm-dependence-orbit-inputs-multiplicity-one}
Assume Harish-Chandra analyticity, the classification of the discrete series, rapid decay of cusp
forms, Harish-Chandra's finiteness theorem for automorphic forms on unitary groups, the trace class
property of the cuspidal spectrum, the factorisations of Dixmier-Malliavin and of Flath, Zelevinsky's
polynomial ring, the geometric lemma and Green's description of the cuspidal characters
(Theorem~\ref{lit-bundle-unitary}). Let $D$ be a definite-side datum whose host satisfies the host
frames hypothesis and has $\mathfrak c=\mathcal O_K$, and which has size data with a parameter choice
(Definition~\ref{asm-dependence-functional-witness-sizes}). Let $R$ be a realization of $D$ for a
family of definite forms, with a face $F$ (Definition~\ref{is-realisation}), such that the family is
ok at some level, $R$ satisfies the weak characterization at $F$, the theta maps of $R$ exist, the
guard holds (Definition~\ref{is-bounded-rows}) and the types of $R$ are good. Then \emph{cuspidal
multiplicities are equal} for $R$: for every level $i$, every discrete series model at infinity for
the base space $Y_0$, every isomorphism of the finite adelic groups of $U(Y_0)$ and $U(W_i)$ that is
compatible with the lattices and the group of $R$, and every irreducible admissible representation
$\sigma$ of the finite adelic group of $U(Y_0)$ that is good for $R$, the cuspidal multiplicity of
$\sigma$ with that model equals the cuspidal multiplicity of the transport of $\sigma$ to $U(W_i)$ with
the model of $R$ at level $i$.
\end{lemma}
\leanref{Development/AutomorphicArithmetic/ThetaCorrespondence/ThetaPairLifts/LocalTypeDictionary/TwistedMultiplicityOne.lean}{GlobalTheory.T.twisted_multiplicity_one}
\leanref{Development/AbelianVarietyTheory/Proofs/DivisorProofs/FinRepApply.lean}{GlobalTheory.CuspMult.MultEq}
\begin{proof}
\openstep{The proof of the equality of cuspidal multiplicities is not yet written out.}
\end{proof}
% Lean: twisted_multiplicity_one = MultEqFinal.twisted_multiplicity_eq, for RealDat3 R and Face3 F
% with Ok3 s, R.WeakCharacterisation F, R.toRealDat.ExAt (the theta maps exist), HostHyp D,
% a parameter choice for some SizeData2, D.host.ccK = top, FWT.QGuard and GoodTypes. Size of the
% open step: proof closure of twisted_multiplicity_eq, about 6,900 theorems and 114,000 lines in
% 1,790 files (text dependencies, an overcount).

\begin{lemma}\label{asm-dependence-orbit-inputs-lambda}
Assume the existence of the Weil representation of a unitary dual pair and of the discrete series of
$U(p,q)$, together with the results named in
Lemma~\ref{asm-dependence-orbit-inputs-multiplicity-one}. Let local Langlands maps be given and let
$p$ be parameters (Definition~\ref{asm-dependence-parameters}). Eventually in the context, every
bundle that $\Lambda_p$ links to a base point of the good cores of $p$ has rows that hold with
$\Lambda_3$ (item (3)(d) of Definition~\ref{asm-definite-side-hypotheses}).
\end{lemma}
\leanref{Development/AbelianVarietyTheory/Proofs/GlobalInterpolation/RealGroupInputsCombined.lean}{GlobalTheory.Fc.RL.LamLinked.evP_params}
\leanref{Development/AutomorphicArithmetic/ThetaCorrespondence/ThetaPairLifts/LocalTypeDictionary/RealGroupInputsCombined.lean}{GlobalTheory.Fc.RL.LamLinked.isotRank_of_wit_alt}
\leanref{Development/AutomorphicArithmetic/ThetaCorrespondence/ThetaPairLifts/ThirdPlaceSeries/RealGroupInputsCombined.lean}{GlobalTheory.Fc.RL.at_row_hLamP_alt_alt}
\begin{proof}
By Lemma~\ref{asm-dependence-functional-witness-numbers}, eventually in the context the size data of
the datum of every base point of a core carry a parameter choice. Fix such a context, and let a bundle
be linked by $\Lambda_p$ to a base point; it is built from $p$ at a stage, with a witness of its
construction. Lemma~\ref{asm-dependence-orbit-inputs-multiplicity-one} applies to the realization of
the theta datum of the stage, with the host frames hypothesis that the real input of the witness
records and the parameter choice: the family is ok at level $1$, and the stage records the weak
characterization, the theta maps, the guard, $\mathfrak c=\mathcal O_K$ and, through its true group,
the good types. So cuspidal multiplicities are equal for it. The good idempotent of the stage is
stable, by Green's description of the cuspidal characters applied to the true group. The cuspidal
spectrum decomposes, by Harish-Chandra's finiteness theorem and rapid decay of cusp forms. From the
existence of the Weil representation and of the discrete series, the pins and the recorded conditions
of the realization of the stage, the host frames hypothesis, the stability, the decomposition and the
equality of cuspidal multiplicities, the datum of the stage has isotypic rank. From the isotypic rank,
the pins of the realization of the twisted input, the conditions that the twisted bundle records for
it and the host frames hypothesis, the readings of the twisted input satisfy the condition
$\Lambda_3$. By the first part of the readings at the companion, the rows of the division pair of the
twisted input, which is that of the bundle, then hold with $\Lambda_3$.
\openstep{The isotypic rank that the equality of cuspidal multiplicities gives, the stability of the
good idempotent, and the condition $\Lambda_3$ of the readings are not yet written out.}
\end{proof}
% Lean: at_row_hLamP_alt_alt: LamLinked.evP_params (szRow_numAt, through Lemma
% asm-dependence-functional-witness-numbers), isotRank_of_wit_alt (goodStable_of_coeff with
% TameTie.idemTie, twisted_multiplicity_one, R3.isotRank_of_multEq with CuspDec.cuspDecomp),
% finiteTypes_via_full_heckeAlgebra_eigensystems (whose conclusion is Lam3 of the readings) and
% rowConditions_at_companion_from_readings (its first part). Sizes of the open step (text
% dependencies, overcounts): isotRank_of_multEq about 2,180 theorems and 34,500 lines;
% finiteTypes_via_full_heckeAlgebra_eigensystems about 2,980 theorems and 49,400 lines.

\begin{lemma}\label{asm-dependence-orbit-inputs-kappa}
Assume the covering of discrete series, the finiteness of Harish-Chandra modules and the
Casselman-Wallach theorem, together with the results named in
Lemma~\ref{asm-dependence-orbit-inputs-multiplicity-one}. Let local Langlands maps be given and let
$p$ be parameters (Definition~\ref{asm-dependence-parameters}). Let a bundle be given at the datum
$D$ of a stage built from $p$, with a witness of its construction from that stage, such that the
size data of $D$ carry a parameter choice and the rows of the bundle hold with $\Lambda_3$. Then the
archimedean scalars $\kappa$ of the columns of the row-side family of the bundle are those of the
division data.
\end{lemma}
\leanref{Development/AutomorphicArithmetic/ThetaCorrespondence/ThetaPairLifts/ThirdPlaceSeries/HarishChandraInputsCombined.lean}{GlobalTheory.Fc.KapLinked.row_kap_of_row_lam_alt_alt}
\begin{proof}
The witness identifies the row-side family of the bundle with that of the division pair of the
twisted input, and the rows with $\Lambda_3$ with the condition $\Lambda_3$ of the readings of the
twisted input. Lemma~\ref{asm-dependence-orbit-inputs-multiplicity-one} applies to the realization
of the theta datum of the stage, as in the proof of Lemma~\ref{asm-dependence-orbit-inputs-lambda},
with the parameter choice of $D$ and with the good types that the rows of the stage give; so
cuspidal multiplicities are equal for it. By Harish-Chandra's finiteness theorem and Green's
description of the cuspidal characters, the isotypic parts of $D$ at the types $\sigma$ and the
levels $i+1$, $0\le i<m$, are then nonzero. The stage records that the archimedean scalars of its
division data at all embeddings are the scalar $\kappa^\sharp$ of $D$. From the five results of
Lemma~\ref{asm-dependence-orbit-inputs-join}, the pins of the realization of the twisted input and the
conditions that the twisted bundle records for it, the condition $\Lambda_3$ of its readings, the
nonvanishing of the isotypic parts and this record, the archimedean scalar of the division data at
every embedding $e$ of the coefficient field is the scalar $\kappa$ at the infinite place of $e$ of
the rows $S_1$ of the division pair of the twisted input. Through the identification of the first
sentence, this is the conclusion.
\openstep{The nonvanishing of the isotypic parts and the equality of the archimedean scalars with
the scalars $\kappa$ of the rows are not yet written out.}
\end{proof}
% Lean: inside row_kap_of_row_lam_alt_alt: twisted_multiplicity_one (with goodTypes_of_rows),
% KapLinked.isot_ne_bot_of_rows (HarishChandraFinite, GreenCuspidal; isot (i+1) (ft sigma) /= bot),
% T.divisorEmbedding_kappa_eq_rowsKappa_at_infinitePlace (the five results, Pins, Holds, Lam3 Z, the
% nonvanishing and Xd.kapR sigma = kapSharp D from q.x.obj.kapR; conclusion Xd.kapR e = (s1Of R F).kappa
% at InfinitePlace.mk e, which is RowsHoldKap3 of the division pair). Sizes of the open step (text
% dependencies, overcounts): isot_ne_bot_of_rows about
% 1,080 theorems and 17,100 lines; divisorEmbedding_kappa_eq_rowsKappa_at_infinitePlace about
% 2,830 theorems and 46,200 lines.

\begin{lemma}\label{asm-dependence-orbit-inputs-good}
Assume Green's parametrization (Corollary~\ref{lit-green-labelled}) and Zelevinsky's polynomial
ring. Let local Langlands maps be given and let $N$ be an extension stronger than the law extension
(Definition~\ref{asm-dependence-parameters}).
\begin{enumerate}
\item Eventually in the context and at large base data, every core built from $p(N)$ is the core of
  a stage at whose datum every row is good.
\item Suppose that at every stage that satisfies $N$ the generators that the group data of the core
  choose at the level primes of the host are those of the label data of these primes, as in the
  condition on generators of Definition~\ref{asm-dependence-parameters}(3). Then the datum $D(w,Z)$
  of every stage built from $p(N)$
  at which every row is good satisfies open parameter statement (e): its labels are tied to an
  automorphic setting (Definition~\ref{asm-open-parameter-statements}).
\end{enumerate}
\end{lemma}
\leanref{Development/AutomorphicArithmetic/ThetaCorrespondence/ThetaPairLifts/ThirdPlaceSeries/LinkedBundlesGoodOfLocalLanglands.lean}{GlobalTheory.Fc.GoodLinked.largeBasePropertyG_goodRec_alt_alt}
\leanref{Development/AutomorphicArithmetic/ThetaCorrespondence/ThetaPairLifts/ThirdPlaceSeries/RowLabOfLayer.lean}{GlobalTheory.Fc.LabLinked.ties_of_built_genPin}
\begin{proof}
(1) By Green's parametrization and Zelevinsky's polynomial ring, eventually in the context and at
large base data, every row of the datum is good at every stage whose core has real expansion data,
whose rows are rows of the construction and whose columns are pinned. A core built from $p(N)$ is the
core of a stage that satisfies $N$ and the standard conditions of the construction. So the stage
satisfies the law extension, which requires the real expansion data of its core; its rows are rows
of the construction and its columns are pinned. Hence every row of its datum is good.

(2) Let the stage be built from $p(N)$, with its witness, and let every row of $D(w,Z)$ be good. The
stage satisfies $N$, so the generators that the group data of its core choose at the level primes
$g_0,g_1$ of the host are those of the label data of $g_0,g_1$. The true group of the rows of the
stage gives that the Teichmüller characters and the series idempotents of $D(w,Z)$ at $g_k$ are those
of the label datum of $g_k$ (Definition~\ref{asm-dependence-test-assembly-true}(1)). From these, the
norm-choice data of the core, the saturation of the stage and the goodness of the rows, Green's
parametrization and Zelevinsky's polynomial ring give the tie of the labels of $D(w,Z)$ to an
automorphic setting.
\openstep{The goodness of the rows at the stages of (1), and the tie of the labels of (2) at a datum
whose rows are good, are not yet written out.}
\end{proof}
% Lean: (1) largeBasePropertyG_goodRec_alt_alt, from GoodLinked.goodFTAll_at_members_of_depthZero
% (MemberAtLayerPoint of GoodFTAll) with the real expansion data (first component of the law
% extension of b.nx.2), the rows of the construction (RXs, b.rx) and the pin (b.pin). (2)
% LabLinked.ties_of_built_genPin = ties_of_genPin_goodFTAll' = ties_of_nAtData_real hG hZ with the core
% row data q.core_data_of_point.row, q.saturationAtWeightedPoint, lab_of_rows q.x.what2 (teich and ser
% of the true group, SerSys.LabTie with labDat), the generators GenPinAlt (c.a.nt.A.grp.zeta k =
% LabelDatum.gen of labDat (g k), the field zeta of the real expansion data) and GoodFTAll.
% Sizes of the open step (text dependencies, overcounts): largeBasePropertyG_goodRec_alt_alt about
% 1,930 theorems and 32,500 lines; ties_of_built_genPin about 1,400 theorems and 23,000 lines.

\begin{lemma}\label{asm-dependence-orbit-inputs}
Let local Langlands maps be given, let $N$ be an extension
(Definition~\ref{asm-dependence-parameters}) and let $p=p(N)$. In the open parameter statements
(Definition~\ref{asm-open-parameter-statements}), write (b1), (b2) and (b3) for the three parts of
(b): the rows hold with $\Lambda_3$, the division rows hold, and the division data are compatible
with the periods; and write (h1) and (h2) for the two parts of (h): the multiplicity condition and
the equality of the archimedean scalars. The results named below are those of
Theorem~\ref{lit-bundle-unitary}. Write $\mathcal R$ for the existence of the Weil representation of
a unitary dual pair, the existence of the discrete series of $U(p,q)$, Harish-Chandra's finiteness
theorem for automorphic forms on unitary groups, and the four results on automorphic forms: the
covering of discrete series, the finiteness of Harish-Chandra modules, rapid decay of cusp forms
and the Casselman-Wallach theorem. Write $\mathcal G$ for the existence of the Weil representation,
the existence and the classification of the discrete series, Harish-Chandra's finiteness theorem,
rapid decay of cusp forms, Harish-Chandra analyticity, the trace class property of the cuspidal
spectrum, the factorisations of Dixmier-Malliavin and of Flath, Zelevinsky's polynomial ring, the
geometric lemma and Green's description of the cuspidal characters
(Theorem~\ref{lit-green-cuspidal}).
\begin{enumerate}
\item Assume $\mathcal R$. If $N$ is stronger than the law extension, statement (d) is linked at
  $p$. The results of $\mathcal R$ and the hypothesis on $N$ are not used in the proof.
\item Assume $\mathcal R$. If statements (a), (b3) and (d) are linked at $p$, so is statement (c).
  The hypotheses on (a) and (b3), and the existence of the Weil representation and of the discrete
  series, are not used in the proof.
\item Assume the results of $\mathcal R$ other than the existence of the discrete series, together
  with the classification of the discrete series, Theorems~\ref{lit-weil-siegel},
  \ref{lit-siegel-weil} and~\ref{lit-doubling}, the unramified computation of the doubling integral
  for $GL_n$, strong approximation for special unitary groups, Siegel's formula for the mean of the
  representation numbers of a Hermitian form with the finiteness of the volume of the automorphic
  quotient, and the conclusion of Lemma~\ref{asm-dependence-howe-unramified}. Then the multiplicity
  condition of the realization of $Q$ at its face, at every member, is linked at $p$, and so is
  (h1).
\item Assume $\mathcal R$, and suppose that the datum $D(w,Z)$ of every stage built from $p$ satisfies
  the rationality law (Definition~\ref{asm-dependence-parameters}(3)). If the first property of (3) is
  linked at $p$, so is (b2). The existence of the Weil representation and of the discrete series is
  not used in the proof.
\item Assume $\mathcal G$. Then (b1) is linked at $p$.
\item Assume the results of $\mathcal G$ other than the existence of the Weil representation and of
  the discrete series, together with the covering of discrete series, the finiteness of
  Harish-Chandra modules and the Casselman-Wallach theorem. If (b1) is linked at $p$, so is (h2).
\item Assume Green's parametrization (Corollary~\ref{lit-green-labelled}) and Zelevinsky's
  polynomial ring. If $N$ is stronger than the law extension, then the property ``at every member,
  every row of the definite-side datum is good'' and statement (f) are linked at $p$.
\item Assume Green's parametrization and Zelevinsky's polynomial ring, and let $N$ be stronger than
  the law extension. If the first property of (7) is linked at $p$, so is statement (e).
\end{enumerate}
\end{lemma}
\leanref{Development/AutomorphicArithmetic/ThetaCorrespondence/ThetaPairLifts/ThirdPlaceSeries/AtTieStatement.lean}{GlobalTheory.T.linkedAt_tie_of_nonempty_extensionWithLaw}
\leanref{Development/AutomorphicArithmetic/ThetaCorrespondence/ThetaPairLifts/ThirdPlaceSeries/AtTie.lean}{GlobalTheory.Fc.RL.at_tie_alt}
\leanref{Development/AutomorphicArithmetic/ThetaCorrespondence/ThetaPairLifts/ThirdPlaceSeries/JoinLinkedRealizationStatement.lean}{GlobalTheory.T.entries_combined_real_alt}
\leanref{Development/AutomorphicArithmetic/ThetaCorrespondence/ThetaPairLifts/ThirdPlaceSeries/JoinLinkedRealization.lean}{GlobalTheory.Fc.JoinLinked.at_selJ_alt}
\leanref{Development/AutomorphicArithmetic/ThetaCorrespondence/ThetaPairLifts/ThirdPlaceSeries/MainTheoremPublishedInputsThird.lean}{GlobalTheory.T.main_theorem_published_inputs_third_alt_alt}
\leanref{Development/AutomorphicArithmetic/ThetaCorrespondence/ThetaPairLifts/ThirdPlaceSeries/MainTheoremPublishedInputsThird.lean}{GlobalTheory.miStarRow_alt_alt}
\leanref{Development/AutomorphicArithmetic/ThetaCorrespondence/ThetaPairLifts/ThirdPlaceSeries/AtHDivGOfMiStatement.lean}{GlobalTheory.T.bk_law_entries_alt}
\leanref{Development/AutomorphicArithmetic/ThetaCorrespondence/ThetaPairLifts/ThirdPlaceSeries/AtHDivGOfMi.lean}{GlobalTheory.Fc.DivLinked.at_hDivG_main_alt}
\leanref{Development/AutomorphicArithmetic/ThetaCorrespondence/ThetaPairLifts/ThirdPlaceSeries/RealGroupInputsCombined.lean}{GlobalTheory.T.real_group_inputs_combined_alt_alt}
\leanref{Development/AutomorphicArithmetic/ThetaCorrespondence/ThetaPairLifts/ThirdPlaceSeries/RealGroupInputsCombined.lean}{GlobalTheory.Fc.RL.at_row_hLamP_alt_alt}
\leanref{Development/AutomorphicArithmetic/ThetaCorrespondence/ThetaPairLifts/ThirdPlaceSeries/HarishChandraInputsCombined.lean}{GlobalTheory.T.harish_chandra_inputs_combined_alt_alt}
\leanref{Development/AutomorphicArithmetic/ThetaCorrespondence/ThetaPairLifts/ThirdPlaceSeries/HarishChandraInputsCombined.lean}{GlobalTheory.Fc.KapLinked.row_kap_of_row_lam_alt_alt}
\leanref{Development/AutomorphicArithmetic/ThetaCorrespondence/ThetaPairLifts/ThirdPlaceSeries/LinkedBundlesGoodOfLocalLanglands.lean}{GlobalTheory.T.linked_bundles_good_of_local_langlands_alt_alt}
\leanref{Development/AutomorphicArithmetic/ThetaCorrespondence/ThetaPairLifts/ThirdPlaceSeries/LinkedBundlesGoodOfLocalLanglands.lean}{GlobalTheory.Fc.GoodLinked.row_good_alt_alt}
\leanref{Development/AutomorphicArithmetic/ThetaCorrespondence/ThetaPairLifts/ThirdPlaceSeries/RowLabOfLayer.lean}{GlobalTheory.T.good_ft_entries_alt_alt}
\leanref{Development/AutomorphicArithmetic/ThetaCorrespondence/ThetaPairLifts/ThirdPlaceSeries/RowLabOfLayer.lean}{GlobalTheory.Fc.LabLinked.row_lab_of_layer_alt_alt}
\begin{proof}
Write $Q_p$ for the slot of the good cores of $p$, and say that a property of bundles holds at every
member as in Lemma~\ref{asm-dependence-orbit-inputs-transfer}.

(1) By Lemma~\ref{asm-dependence-orbit-inputs-comparison} for the slot $Q_p$, eventually in the
context every bundle that $\Lambda_p$ links to a base point of a good core of $p$ has a comparison
part. By Lemma~\ref{asm-dependence-orbit-inputs-transfer}(1), statement (d) is linked at $p$.

(2) Statement (d) is linked at $p$. By Lemma~\ref{asm-dependence-orbit-inputs-join}, every bundle at
the datum of a stage built from $p$, with a witness of its construction from that stage and a
comparison part, has its indefinite realization joined. So
Lemma~\ref{asm-dependence-orbit-inputs-transfer}(2), with $A$ the existence of a comparison part,
shows that statement (c) is linked at $p$.

(3) By Lemma~\ref{asm-dependence-orbit-inputs-multiplicity}, eventually in the context every bundle
that $\Lambda_p$ links to a base point of a good core of $p$ satisfies the multiplicity condition.
Taking this as the condition on bundles, the multiplicity condition at every member holds along
$(Q_p,\Lambda_p)$, as in the proof of Lemma~\ref{asm-dependence-orbit-inputs-transfer}(1). By
Lemma~\ref{asm-dependence-orbit-inputs-transfer}(3), its eventual form (h1) holds along
$(Q_p,\Lambda_p)$ with the same condition on bundles. With the slot of all cores, both are linked at
$p$.

(4) The first property of (3) is linked at $p$, and by hypothesis the datum of every stage built
from $p$ satisfies the rationality law. By Lemma~\ref{asm-dependence-orbit-inputs-division}, every
bundle at the datum of a stage built from $p$, with a witness of its construction from that stage,
that satisfies the multiplicity condition at its face satisfies the division rows. So (b2) is linked
at $p$, by Lemma~\ref{asm-dependence-orbit-inputs-transfer}(2).

(5) By Lemmas~\ref{asm-dependence-orbit-inputs-lambda}
and~\ref{asm-dependence-orbit-inputs-transfer}(1), (b1) is linked at $p$.

(6) Let $R$ and $E$ be a slot and a condition on bundles for (b1), and let $s$ be a selection for a
class $\mathcal A$ whose bundles are built from $(Q_p\cap R,\Lambda_p)$ and satisfy $E$. Let $D$ be a
member at a scale $x>x_1$. The selected bundle at $D$ is linked by $\Lambda_p$ to a base point of a
good core of $p$, so it lies at the datum $D$ of a stage built from $p$, with a witness of its
construction from that stage. The size data of $D$ carry a parameter choice
(Definition~\ref{pd-family}(4)), and the rows of the bundle hold with $\Lambda_3$, by (b1). By
Lemma~\ref{asm-dependence-orbit-inputs-kappa}, the scalars $\kappa$ of the columns of its row-side
family are those of the division data. This holds at every member, so by
Lemma~\ref{asm-dependence-orbit-inputs-transfer}(3) (h2) holds along $(Q_p\cap R,\Lambda_p)$ with
$E$, and (h2) is linked at $p$ with the slot $R$.

(7) Let $R$ be the slot of the cores of stages at whose datum every row is good. By
Lemma~\ref{asm-dependence-orbit-inputs-good}(1), eventually in the context and at large base data,
every core built from $p$ satisfies $R$. A bundle that $\Lambda_p$ links to a base point of a core in
$Q_p\cap R$ lies at the datum of that base point, at which every row is good. Taking this as the
condition on bundles, ``at every member, every row of the definite-side datum is good'' holds along
$(Q_p\cap R,\Lambda_p)$, so it is linked at $p$. By
Lemma~\ref{asm-dependence-orbit-inputs-transfer}(3), so is statement (f).

(8) The real expansion data that the law extension requires include the condition on generators of
Definition~\ref{asm-dependence-parameters}(3); so the hypothesis of
Lemma~\ref{asm-dependence-orbit-inputs-good}(2) holds at every stage that satisfies $N$. Let $R$ and
$E$ be a slot and a condition on bundles for the first property of (7),
and let $s$ be a selection for a class $\mathcal A$ whose bundles are built from
$(Q_p\cap R,\Lambda_p)$ and satisfy $E$. At a member $D$, the selected bundle lies at the datum $D$ of
a stage built from $p$, and every row of $D$ is good. By
Lemma~\ref{asm-dependence-orbit-inputs-good}(2), statement (e) holds at $D$. So (e) is linked at
$p$, with the slot $R$.
\end{proof}
% Lean, part by part: (1) at_tie_alt = LinkedAt.of_row (RowE.row (TieOfLink.hasTie_of_linked));
% (2) at_selJ_alt = linked_atAll_of_imp (A := HasCompatibility) with selJ_of_bWit; (3)
% miStarRow_alt_alt, LinkedAt.of_row with hr.imp evW_of_atAll; (4) at_hDivG_main_alt =
% hDivG_linked_of_mIstar2 = linked_atAll_of_imp with hDivG_of_mIstar2; (5) at_row_hLamP_alt_alt
% (EvP.mono of LamLinked.evP_params); (6) row_kap_of_row_lam_alt_alt (Row.imp, Built.ind with
% Fm.params, and Closed <1, ...>); (7) row_good_alt_alt = linkedAt_goodFTAll = LgModel.rowL_of_good
% (largeBasePropertyG_goodRec_alt_alt) (row_goodRec_alt), where rowL_of_good passes from the
% LargeBaseConditionG of Lemma -good(1) to the slot of Definition asm-dependence-parameters(2) through
% largeBaseProperty_conditional (the cores that satisfy R in the contexts where that condition holds)
% and the condition on bundles 'built from (Q_p and R, Lambda_p), and E'; (8) row_lab_of_layer_alt_alt =
% row_lab_of_good_alt_alt with the field zeta of RealExpansionExtraData (e.1.zeta), through
% ties_of_built_genPin.

\subsection{The covolume bound and bundles}\label{subsection-asm-covolume-bundles}

\begin{lemma}\label{asm-dependence-covolume-ties-points-host}
Eventually in the context (Definition~\ref{asm-stage-construction}(1)), the following holds at every
base point of a core. Let $D$ be its definite-side datum, with ranks $k$, $m$, $t$, $f$ and
$r=k+m-1$, host field $E$ with discriminant $d_E$, level primes $g_0,g_1$ and primes
$p_{de},p'_{de}$, and with its size data at the scale $(x,\varepsilon)$ of the context; put
$L=x^{\varepsilon}=\log b$. Then:
\begin{enumerate}
\item the host satisfies the \emph{size bound} for $L$: $r\ge62$, $L\ge4$, $\log f\le L$, $\log
  g_i\le9L+1$ for $i=0,1$, $\log p_{de}\le L$, $\log p'_{de}\le L$, and $\log|d_E|\le4fL$;
\item let $F_O\subset\mathbf C$ be the field generated over $\mathbf Q$ by the values of the map
  $\varphi_O$ of $D$, and put $P(L')=19fr^2L'+198f^2L'$. If $\log[F_O:\mathbf Q]\le P(L')$ and
  $\log|d_{F_O}|\le[F_O:\mathbf Q]\,P(L')$ for every $L'$ for which the host satisfies the size bound,
  then these hold for $L'=L$.
\end{enumerate}
\end{lemma}
\leanref{Development/AbelianVarietyTheory/Proofs/HeightProofs/EventuallyConstructedPointHostSizeParametersNumericaStatement.lean}{GlobalTheory.T.eventually_constructedPoint_hostSize_numerical}
\leanref{Development/AbelianVarietyTheory/Proofs/HeightProofs/EventuallyConstructedPointHostSizeNumerical.lean}{GlobalTheory.Plug.eventually_hostSize_and_numAt}
\leanref{Development/AbelianVarietyTheory/Proofs/HeightProofs/HostSzOfSizes.lean}{GlobalTheory.FieldSize.hostSz_of_sizes}
\leanref{Development/AbelianVarietyTheory/Definitions/HeightFunctions/JointQ.lean}{GlobalTheory.FieldSize.HostSz}
\begin{proof}
By Lemma~\ref{asm-dependence-functional-witness-numbers}, eventually in the context the size data of
the datum of every base point of a core carry parameter choices and satisfy the size bounds, and
$g_0,g_1\le\exp(9L)$. Take a context beyond its threshold with $x\ge10^{60}$, and a base point.

(1) Since $g_i\le e^{9L}$, $\log g_i\le9L\le9L+1$. The size bounds give the window
$x^{9/2+10\varepsilon-\varepsilon/100}\le k$, so $x\le k\le r$, and $r\ge x\ge10^{60}\ge62$. Put
$Q=Q_Sqd$. The parameter choices give $Q^{36}t^{150}\le b^{1/2}$; as $Q,t\ge1$, this gives $\log
Q\le L/72$ and $\log t\le L/300$. The ranks satisfy $6\le r\le t$, so $\log2\le\log r\le L/300$, and
$L\ge300\log2>4$. The field $E$ has degree $2f$ and lies in the cyclotomic field of conductor $Q$, so
$2f\le\varphi(Q)\le Q$, and $\log f\le\log Q\le L/72\le L$. The parameter choices give $\log
p_{de}\le L$ and $\log p'_{de}\le L$. Finally $\log|d_E|=2f\log\mathrm{rd}_E$, and the host has
$\mathrm{rd}_E\le\frac{Q_S}2\sqrt{qd}\le Q\le Q^{36}t^{150}\le b^{1/2}\le b$
(Lemma~\ref{asm-dependence-functional-witness-numbers-bounds-computations}(3)), so $\log
\mathrm{rd}_E\le L$ and $\log|d_E|\le2fL\le4fL$.

(2) By (1) the host satisfies the size bound for $L$, so the hypothesis applies with $L'=L$.
\end{proof}
% Lean: T.eventually_constructedPoint_hostSize_numerical = Plug.eventually_hostSize_and_numAt: EvP.of_ev
% with EX.szRow_numAt (Lemma asm-dependence-functional-witness-numbers) at the base point (Val from
% beta.p.pt_val, hostArith_bpt), threshold max(x0, CP^20, 10^60); hg from N.g_le (g <= exp(9L)); then
% FieldSize.hostSz_of_sizes: rank from scale_le_rank (SizesAt, Pm.k_le_r) and 10^60 <= x; len from
% Pm.log_r_le (log_cond_le from P.condE_le, r_le_alt_alt) and six_le_r; deg from Pm.log_f_le
% (f_le_condE: E embeds in the cyclotomic field of QS q d); p from P.pde_log; rd from
% logdE_le_two_f_L (logdE_eq_rd, log_rdE_le_scL with rdE_pos_le = factRd). The HostSz field rd is
% log|d_E|/(2f) <= 2L, stated here as log|d_E| <= 4fL. Part (2) is hS (Sz -> FieldSmall, with
% PO D L = poR D.R L) applied to FieldTies.small. The second component NumAtL also gives DV.KapDef
% and PageNum (1000 <= CP from Std, CP^20 <= zBase L through x_le_zBase_of_pow_le_n); the book states
% the parts that rowE_part_ties_alt uses (FieldSmall is the third component of NumAtL).

\begin{lemma}\label{asm-dependence-covolume-ties-points-core}
Let $D$ be a definite-side datum with host field $E$, imaginary quadratic subfield $K$, level primes
$g_0\ne g_1$, primes $p_{de}\ne p'_{de}$ and coefficient field $L$, let $F_O$ be the field of
Lemma~\ref{asm-dependence-covolume-ties-points-host}(2), let size data of $D$ be given, with their
number $b$, and let $P$ be the polynomial of that lemma. Let error data of $D$ be given, with elements
$\pi_0,\pi_1\in\mathcal O_K$ such that $\pi_k\bar\pi_k=g_k$. Let $R$ be a realization of $D$
(Definition~\ref{is-realisation}), with set $P_g$ of places of $E$, ramified sets $Q_t$ of its tuples
$t$, map $v\mapsto\mu_v$ and twisted places, and with a frame consisting of primes
$\mathfrak p_0,\mathfrak p_1$ of $K$ with $g_k\in\mathfrak p_k$, characters $\delta,\delta'$ of $E$ and
further characters. Suppose that the level system of $R$ satisfies, among its conditions:
\begin{enumerate}
\item[(i)] a place $P$ of $E$ lies in $P_g$ exactly when it lies over $\mathfrak p_0$ or
  $\mathfrak p_1$;
\item[(ii)] for $p\in\{p_{de},p'_{de}\}$ there is a prime $\mathfrak p$ of $K$ such that every place
  $P\in Q_t$ that contains $p$, for every tuple $t$, lies over $\mathfrak p$.
\end{enumerate}
Let a core of $D$ be given, with primes $\mathfrak q_0,\mathfrak q_1$ of $K$, level $M$ and level ideal
$M_b$ (Definition~\ref{is-onto-step-model-level-condition}). Suppose that the core is \emph{tied} to
the frame; besides the laws of the core data, this means:
\begin{enumerate}
\item[(a)] the core satisfies its laws, its quotient data are canonical, the signature $(p_j,q_j)$ of
  each of its levels has $p_j+q_j\ge1$, and $\mathfrak q_k=\mathfrak p_k$ for $k=0,1$;
\item[(b)] the \emph{field ties} hold: $F_O$ is \emph{good} at $p_{de}$ and at $p'_{de}$, the values
  of $\delta$ and $\delta'$ lie in $F_O$, and $F_O$ satisfies the hypothesis of
  Lemma~\ref{asm-dependence-covolume-ties-points-host}(2). Here $F_O$ is good at $p$ if, among other
  conditions, it contains the roots of unity of orders $g_0g_1$ and $p$, every automorphism of
  $\mathbf C$ over $F_O$ fixes the curve data of the host at the modulus $g_0g_1p$, and $p$ splits
  in $K$;
\item[(c)] $\bar\pi_k\in\mathfrak p_k\mathcal O_E$ for $k=0,1$.
\end{enumerate}
Suppose also that $F_O$ is \emph{small} for the size data: $\log[F_O:\mathbf Q]\le P(\log b)$ and
$\log|d_{F_O}|\le[F_O:\mathbf Q]\,P(\log b)$. Then the \emph{core ties} hold: the conditions of (a)
other than the last, condition (c) with $\mathfrak q_k$ in place of $\mathfrak p_k$, the goodness of
$F_O$ at $p_{de}$ and $p'_{de}$, its smallness, the values of $\delta$ and $\delta'$ in $F_O$, and:
\begin{enumerate}
\item $M_b=\prod_{P\in P_g}P$;
\item for $p\in\{p_{de},p'_{de}\}$ there is a prime $\mathfrak P_p$ of $K$ containing $p$ such that
  $\mathfrak P_p\mathcal O_E\subset P$ for every tuple $t$ and every $P\in Q_t$ that contains $p$. In
  particular, for every twisted place $v$ of $L$ with own prime $p$
  (Definition~\ref{is-onto-step-runs}(1)), $\mathfrak P_p\mathcal O_E$ lies in every place of
  $Q_{\mu_v}$ that contains $p$.
\end{enumerate}
\end{lemma}
\leanref{Development/AbelianVarietyTheory/Proofs/HeightProofs/CoreTiesOfTied.lean}{GlobalTheory.Plug.coreTies_of_tied}
\leanref{Development/AbelianVarietyTheory/Proofs/DivisorProofs/ExistsCoreTiesOfHyp.lean}{GlobalTheory.BigOn.coreTies_of_is}
\leanref{Development/AbelianVarietyTheory/Proofs/DivisorProofs/ExistsCoreTiesOfHyp.lean}{GlobalTheory.BigOn.level_of_is}
\leanref{Development/AbelianVarietyTheory/Proofs/DivisorProofs/ExistsCoreTiesOfHyp.lean}{GlobalTheory.BigOn.oneHalf_of_is}
\leanref{Development/AbelianVarietyTheory/Definitions/Divisors/IsDeltaDelAtOwnAt.lean}{GlobalTheory.BigOn.CoreTies}
\leanref{Development/AbelianVarietyTheory/Definitions/HeightFunctions/RunsAlong.lean}{GlobalTheory.Plug.Tied}
\leanref{Development/AbelianVarietyTheory/Definitions/AutomorphicMembership/FieldTies.lean}{GlobalTheory.AtMember.FieldTies}
\leanref{Development/AbelianVarietyTheory/Definitions/HeightFunctions/NormComparisonConstants.lean}{GlobalTheory.FinalJoin2.FieldOK0}
\begin{proof}
The conditions of (a) other than the last, the goodness of $F_O$, its smallness and the values of
$\delta$ and $\delta'$ are hypotheses. Since $\mathfrak q_k=\mathfrak p_k$, condition (c) holds with
$\mathfrak q_k$.

(1) The laws of the core contain its pins and its half condition. The pins give $M=g_0g_1$, and the
half condition gives $g_k\in\mathfrak q_k$ and $M_b=\mathfrak q_0\mathcal O_E\cdot\mathfrak q_1\mathcal
O_E$. The level of the core has $M_b\bar M_b=M\mathcal O_E=g_0g_1\mathcal O_E$. The datum records that
$g_0$ and $g_1$ split completely in $E$: there are $[E:\mathbf Q]$ places of $E$ over $g_k$. Since the
sum of the products $e_Pf_P$ of ramification index and residue degree over these places is
$[E:\mathbf Q]$, every $e_P$ is $1$, and $g_k\mathcal O_E$ is a product of distinct primes. The ideals
$g_0\mathcal O_E$ and $g_1\mathcal O_E$ are coprime, as $g_0\ne g_1$ are primes. So
$g_0g_1\mathcal O_E$ is squarefree, and so is its divisor $M_b$. A nonzero squarefree ideal of the
Dedekind domain $\mathcal O_E$ is the product of the primes that contain it. A place $P$ of $E$
contains $M_b=\mathfrak q_0\mathcal O_E\cdot\mathfrak q_1\mathcal O_E$ exactly when it contains some
$\mathfrak q_k\mathcal O_E$, that is, when $P\cap\mathcal O_K$ contains the maximal ideal
$\mathfrak q_k$, that is, when $P$ lies over $\mathfrak q_k=\mathfrak p_k$. By (i) this says $P\in
P_g$. Hence $M_b=\prod_{P\in P_g}P$.

(2) Let $p\in\{p_{de},p'_{de}\}$, and let $\mathfrak p$ be the prime of $K$ given by (ii). Suppose
first that some place $P\in Q_t$ contains $p$. It lies over $\mathfrak p$, so $p\in P\cap\mathcal
O_K=\mathfrak p$. Every place $P'\in Q_{t'}$ that contains $p$ also lies over $\mathfrak p$, so
$\mathfrak p\mathcal O_E\subset P'$. Take $\mathfrak P_p=\mathfrak p$. Otherwise no place of any $Q_t$
contains $p$, and $\mathfrak P_p$ may be any prime of $K$ over the prime number $p$. The own prime of
a finite place $v$ of $L$ is $p_{de}$ or $p'_{de}$, so the last statement is the case $t=\mu_v$ with
$p$ the own prime of $v$.
\end{proof}
% Lean: Plug.coreTies_of_tied (D, EX, J, R, mu0; Tied EX J R.C, R.toLevCh.Is mu0, FieldSmall D L):
% HalfEX from hT.core.pb and hT.half, then BigOn.coreTies_of_is with hT.core (laws, canon, rank,
% pb) and hT.field (one, two, del, del'); the field small of FieldTies is not used. The tie is
% Plug.Tied = (lawsJ, CorePin with pb, FieldTies, AtMember.HalfAt EX C); (i) is LevCh.Is.Pg (Div k
% PP: PP.under K = pb k), (ii) is LevCh.Is.Q_de. (1) level_of_is: T.Q.l.prod with pins.2.1.1 (PinLevel
% M = g0 g1), CoverHalf.squarefree_span_g (GSplit.g_splits = split_of_ncard with D.nd.g_split, and
% coprimality of g0 g1), Squarefree.of_mul_left, CoverHalf.eq_prod_of_squarefree, and for the
% membership hL.Pg, hT.half.2 and CoverHalf.prod_map_le_iff. (2) oneHalf_of_is with exists_half_of_is
% (hL.Q_de; the fallback is exists_heightOneSpectrum_natCast_mem) and ownAt_special; CoreTies.oneHalf
% quantifies over the twisted places R.twPl and placesOver R (R.mu v) (ownAt D v). Size: about 210
% lines in five files (ExistsCoreTiesOfHyp, ExistsHalfOfBasePins, GSplit,
% ExistsInjectiveCharacterOfIsCyclic, CoreTiesOfTied).

\begin{lemma}\label{asm-dependence-covolume-ties-points-fields-weber}
Let base data $B$ be given, with host field $E$, embedding $\varphi_0$, prime $d$, level primes $g_0,g_1$ and primes
$p_{de},p'_{de}$; put $N=g_0g_1p_{de}p'_{de}$ and $\theta=\varphi_0(\sqrt{-d})$. Let $\delta$ be a field datum for $B$,
that is the data of an exact field datum of Lemma~\ref{asm-dependence-covolume-ties-points-fields} without the
condition of exactness, with its ray class field $R_\delta\subset\mathbf C$ of $E$ of modulus $N\mathcal O_E$. Let $n\ge3$
be an odd divisor of $N$. Then the Weber values of level $n$ of the lattice $\mathbf Z+\mathbf Z\frac{1+\theta}2$
(Lemma~\ref{cv-deep-ray-field}) lie in $R_\delta$.
\end{lemma}
\leanref{Development/AbelianVarietyTheory/Proofs/IwasawaDecomposition/ContainsRayFieldSub.lean}{GlobalTheory.FieldSize.WeberRay.weberSet_subset_ray}
\leanref{Development/AbelianVarietyTheory/Definitions/HeightFunctions/FiniteDimensionalOfSz.lean}{GlobalTheory.FieldSize.weberSet}
\begin{proof}
\openstep{The proof, which compares the field generated over $R_\delta$ by the Weber values of level $n$ with $R_\delta$
through the primes of $\mathbf Q(\theta)$ that split completely in it, is not yet written out.}
% Lean: FieldSize.WeberRay.weberSet_subset_ray, 5,095 lines (tools/lean_size.py). Its first steps are
% FieldSize.exists_hostK (the subfield K of C with O_K the lattice), containsRayField_sub (the ray class field of K in
% R_delta) and T.principalRayPrimes_splitCompletely_weberValues; the rest compares splitting sets in the field
% generated by R_delta and the Weber values (WeberIn.numberField_adjoin).
\end{proof}

\begin{lemma}\label{asm-dependence-covolume-ties-points-fields-size}
Assume the existence of ray class fields (Theorem~\ref{lit-bundle-unitary}(5)). Let base data $B$ with ranks $f$ and
$r=k+m-1$ and a field datum $\delta$ for $B$ (Lemma~\ref{asm-dependence-covolume-ties-points-fields-weber}) be given, and
let $L'$ be a number for which the host of $B$ satisfies the size bound of
Lemma~\ref{asm-dependence-covolume-ties-points-host}(1). Then the field $\Omega_\delta$ of $\delta$ is a number field, and
with $P(L')=19fr^2L'+198f^2L'$,
\[
 \log[\Omega_\delta:\mathbf Q]\le P(L'),\qquad\log|d_{\Omega_\delta}|\le[\Omega_\delta:\mathbf Q]\,P(L').
\]
\end{lemma}
\leanref{Development/AbelianVarietyTheory/Proofs/IwasawaDecomposition/OmegaFieldDegreeRootDiscrBoundedByRingPolynomialStatement.lean}{GlobalTheory.T.omegaField_degree_rootDiscr_bounded_by_ringPolynomial}
\begin{proof}
\openstep{The proof is not yet written out.}
% Lean: T.omegaField_degree_rootDiscr_bounded_by_ringPolynomial = FieldSize.sz_Om_poR, 3,592 lines (tools/lean_size.py),
% with HostSz for the size bound and Sz Om (poR L) (poR L) for the conclusion.
\end{proof}

\begin{lemma}\label{asm-dependence-covolume-ties-points-fields}
Assume the existence of ray class fields (Theorem~\ref{lit-bundle-unitary}(5)). Let base data $B$
with a framing $f$ and twist data be given, with host field $E$, CM embedding $\varphi_0$, imaginary
quadratic subfield $K$, level primes $g_0,g_1$ and primes $p_{de},p'_{de}$; put
$N=g_0g_1p_{de}p'_{de}$, and let $L_s\subset\mathbf C$ be the image of the coefficient field of $B$.
Let a column choice over these data carry an \emph{exact field datum} $\delta$ for $f$: subfields of
$\mathbf C$, among them a ray class field $R_\delta$ of $E$ of modulus $N\mathcal O_E$, into which
$E$ embeds through $\varphi_0$, and a field $\Omega_\delta$ that contains all of them, such that
$L_s$ is the compositum of $\varphi_0(E)$, $R_\delta$, the field of the roots of unity of $B$
(which contains those of order $g_k(g_k-1)$ for $k=0,1$), the field of the values of the
characters of $f$, and further fields of $\delta$. Let $D$ be the datum of a later stage of the
column choice, with coefficient field $L$, the subfield $L_0\subset L$ through which its
coefficients are read, and the frame $C_f$ that $f$ gives at the stage, with characters
$\delta_f,\delta'_f$. Then:
\begin{enumerate}
\item the field ties of Lemma~\ref{asm-dependence-covolume-ties-points-core}(b) hold for $D$ and
  $C_f$;
\item $D$ and $C_f$ satisfy the \emph{condition on small places}: the characters of $C_f$ take
  their values in $L_0$, in the normalization of the frame (weight $1/2$ for its character $\eta$,
  weight $0$ for the others); there is a ray class field $R_K$ of $K$ of modulus $N\mathcal O_K$ with an
  embedding into $L_0$ that restricts to $\varphi_0$ on $K$; and $L_0$ contains the roots of unity
  of order $g_k(g_k-1)$ for $k=0,1$.
\end{enumerate}
\end{lemma}
\leanref{Development/AbelianVarietyTheory/Proofs/IwasawaDecomposition/IndexBoundCoefficientFieldOfExactSmallField.lean}{GlobalTheory.AtMember.fieldTies_of_ksExact}
\leanref{Development/AbelianVarietyTheory/Proofs/HeightProofs/SmallFrameAtColumn.lean}{GlobalTheory.SmallExact.smallHas_of_rayField}
\leanref{Development/AbelianVarietyTheory/Statements/PlaceShiftStatements/SingleIndexPlacesAreTwisted.lean}{GlobalTheory.BigOn.SmallHas}
\leanref{Development/AbelianVarietyTheory/Definitions/AutomorphicMembership/FixedCoefficientFieldInComplex.lean}{GlobalTheory.Mem.Base.KsExact}
\leanref{Development/AbelianVarietyTheory/Definitions/AutomorphicMembership/FixedCoefficientFieldInComplex.lean}{GlobalTheory.Mem.Fld.FieldDat}
\begin{proof}
The field $L_s$ lies in $L_0$. The field $F_O$ generated by the values of the map $\varphi_O$ of $D$
is $L_s$: the values of $\varphi_O$ are the images of the integers of the coefficient field of $B$.

(2) The column choice records that the characters of $C_f$ take their values in $L_0$. The
existence of ray class fields gives a ray class field $R_K$ of $K$ of modulus $N\mathcal O_K$. The ray
class field of $K$ of modulus $N\mathcal O_K$ embeds over $K$ into every ray class field of $E$ of
modulus $N\mathcal O_E$ \cite{Neu99}; so there is an embedding $j$ of $R_K$ into $R_\delta$ that
restricts to $\varphi_0$ on $K$. By exactness $R_\delta\subset L_s\subset L_0$. A root of unity of
order dividing $g_k(g_k-1)$ lies in the cyclotomic field of that order, which lies in the field of
the roots of unity of $B$, hence in $L_s\subset L_0$.

(1) Let $p\in\{p_{de},p'_{de}\}$. The field $F_O$ is good at $p$: $p$ is a prime among the level
primes, and splits in $K$, as the base data require of their level primes; the datum built at the
stage carries its ring of coefficients onto the ring of integers of its coefficient field, has the
roots of unity of order $g_0g_1$ among the values of $\varphi_O$, and has principal type norms; $0$
is not a level prime; and since $p$ divides $N$, the ray
class field $R_\delta$ contains the roots of unity of order $p$, so they lie in $F_O=L_s$. Finally,
every automorphism $\tau$ of $\mathbf C$ that fixes $F_O$ pointwise fixes a CM curve with its points of
order $n=g_0g_1p$ (Definition~\ref{cv-components-conjugation}(3), for $\theta=\varphi_0(\sqrt{-d})$).
Indeed, $n$ is odd, as the level primes are, it is at least $3$, and it divides $N$. The field
$F_O=L_s$ contains $\theta$, as it contains $\varphi_0(E)$, and the Weber values of level $n$ of the
lattice $\mathbf Z+\mathbf Z\frac{1+\theta}2$, as they lie in $R_\delta\subset L_s$ by
Lemma~\ref{asm-dependence-covolume-ties-points-fields-weber}. So $\tau$ fixes $\theta$ and these values,
and Lemma~\ref{cv-components-weber-curve} applies.

The characters $\delta_f$ and $\delta'_f$ of the frame have finite order, so they are trivial on the
archimedean part of the ideles, and by exactness the field generated by their values lies in
$L_s=F_O$; so their values lie in $F_O$.

For the hypothesis of Lemma~\ref{asm-dependence-covolume-ties-points-host}(2), let $L'$ be a number
for which the host satisfies the size bound. The field $F_O=L_s$ lies in $\Omega_\delta$, and
$\log[\Omega_\delta:\mathbf Q]\le P(L')$ and $\log|d_{\Omega_\delta}|\le[\Omega_\delta:\mathbf
Q]\,P(L')$ by Lemma~\ref{asm-dependence-covolume-ties-points-fields-size}. The same bounds hold for its
subfield $F_O$: its degree divides that of $\Omega_\delta$, and its root discriminant
$|d_{F_O}|^{1/[F_O:\mathbf Q]}$ is at most that of $\Omega_\delta$, since
$d_{F_O}^{[\Omega_\delta:F_O]}$ divides $d_{\Omega_\delta}$.
\end{proof}
% Lean: (1) AtMember.fieldTies_of_ksExact (X, f, delta, KsExact): hFO from mem_FO_assembled_iff and
% mem_KsC (FO = KsC); KsExact.lower gives ray <= KsC, charField <= KsC. FieldOK0 at p: own; full =
% hermitianCodeWordStages_ringOfIntegers_full; zetaM = hermitianCodeWordStages_has_rootOfUnity;
% zetaP from RayRoots.rootOne_mem_of_isRayFieldIn (p | modulus = g0 g1 pde pde'); curve from
% FieldSize.curveOver_of_weberSet (through FinalJoin.CurveField.curveOver_of_weberIn and
% CompRat.Weber.fixesCurve_of_weber, Lemma cv-components-weber-curve) with imField_le_KsC and
% WeberRay.weberSet_subset_ray (n = g0 g1 p odd, >= 3, dividing the modulus; Lemma -fields-weber, open);
% kpsi = typeNorms_principal; bad from level_cond; split from level_split. del, del' from ValFin.mem_of_orderDvd
% (isDel.fin, charField <= KsC). small from listField_le_Om, FieldSize.sz_mono_and_of_le.2 and
% T.omegaField_degree_rootDiscr_bounded_by_ringPolynomial = FieldSize.sz_Om_poR (Lemma -fields-size, open). (2) SmallExact.smallHas_of_rayField: val = w.val; ray from RayField
% at (modulus) of K, rayLeC (through rayLe, proved with Bauer's theorem,
% LeOfFiniteSplitNotSplitBauer.lean; the book cites the inclusion of ray class fields), w.exact
% (ray <= KsC) and KsC_le_LL0C; roots from mem_cycField, rootsField_le_KsC. The open parts are now in
% Lemmas -fields-weber (weberSet_subset_ray, 5,095 lines) and -fields-size (sz_Om_poR, 3,592 lines);
% curveOver_of_weberSet itself (203 lines) is the reduction to Lemma cv-components-weber-curve written above.

\begin{lemma}\label{asm-dependence-covolume-ties-points-weighted}
Assume the existence of ray class fields. Let local Langlands maps be given, let $N$ be an extension
stronger than the law extension (Definition~\ref{asm-dependence-parameters}), and let a stage
$(K,B,f,T,w,Z)$ satisfy $N$, with rows that are rows of the construction, norm choice with
characters $\psi_v$, core object $O$ with its error data, and datum $D=D(w,Z)$ with coefficient field
$L$. Let the side of the stage at its weighted point be formed with the reading of $N$ of the proof of
Lemma~\ref{asm-dependence-covolume-ties-points}(2), and let $C_f$ be the frame that $f$ gives at the
stage. Then the side satisfies its \emph{ties}:
\begin{enumerate}
\item the frame of its realization and the frame of its core are $C_f$, and the level system of the
  realization satisfies its conditions for some base tuple;
\item the core is tied to $C_f$: conditions (a)--(c) of
  Lemma~\ref{asm-dependence-covolume-ties-points-core} hold, with the error data of $O$, and the core
  data satisfy their own laws;
\item the real expansion has a series tie, the group isomorphism carries the level group and the
  character $\psi_0^{-1}$ of the side to the level data of the datum, and the theta data are complete;
\item the choice of the characters $\psi_v$ is compatible with the norm choice, and $\psi_v$ takes
  values in $L$ when $\mu_v\ne\mu_0$;
\item $D$ and $C_f$ satisfy the condition on small places
  (Lemma~\ref{asm-dependence-covolume-ties-points-fields}(2));
\item the types have the tame tie: for $k=0,1$, the Teichmüller character and the series idempotents
  of $D$ at $g_k$ are those of the label datum of $g_k$, and for every place $v$ of the group at $g_k$
  and every $j$ the label datum carries the character of the types at $(k,v,j)$ to
  $\omega_k^{2(j+1)}$, where $\omega_k$ is the character of $(\mathbf Z/g_k)^\times$ of $C_f$.
\end{enumerate}
\end{lemma}
\leanref{Development/AutomorphicArithmetic/ThetaCorrespondence/ThetaPairLifts/ThirdPlaceSeries/SideWTies.lean}{GlobalTheory.PlugPt.side_weight_ties_of_ray_field}
\leanref{Development/AbelianVarietyTheory/Definitions/HeightFunctions/OntoAlong.lean}{GlobalTheory.PlugPt.Side.Ties}
\leanref{Development/AbelianVarietyTheory/Proofs/HeightProofs/WPtFrame.lean}{GlobalTheory.Plug.wPt_frame}
\leanref{Development/AbelianVarietyTheory/Proofs/HeightProofs/CoreFrameCompatibleConstructedPoint.lean}{GlobalTheory.Plug.tied_cfOf}
\leanref{Development/AbelianVarietyTheory/Proofs/IwasawaDecomposition/ConjGeneratorsMemFramePrimesTrueColumns.lean}{GlobalTheory.AtMember.conj_gen_mem_frame}
\leanref{Development/AutomorphicArithmetic/ThetaCorrespondence/ThetaPairLifts/ThirdPlaceSeries/SideWChr.lean}{GlobalTheory.PlugPt.side_weight_psi_values_tame_tie}
\leanref{Development/AbelianVarietyTheory/Proofs/HeightProofs/TwDatPsiValuesInLL.lean}{GlobalTheory.twDat_psi_valuesIn_LL}
\leanref{Development/AutomorphicArithmetic/ThetaCorrespondence/ThetaPairLifts/ThirdPlaceSeries/SideWExportSmall.lean}{GlobalTheory.sideW_export_small_alt}
\begin{proof}
The column choice $w$ carries an exact field datum for $f$, so
Lemma~\ref{asm-dependence-covolume-ties-points-fields} applies to $D$ and $C_f$.

(1) The realization built at the stage records that its frame is the frame of its level data at the
column choice, and the column choice records that this frame is the one $f$ gives; the frame of the
core of the side is by definition the frame $f$ gives at the stage. The level data of the
realization record that its level system satisfies its conditions for a base tuple.

(2) The geometry of the realization records the laws of the core data. The rows are rows of the
construction, which require that the core satisfies its laws, that its quotient data are canonical,
that its levels have $p_j+q_j\ge1$, and that its primes are the primes $\mathfrak p_0,\mathfrak p_1$ of
the frame $f$. The field ties are Lemma~\ref{asm-dependence-covolume-ties-points-fields}(1).

For (c), the error data of $O$ are those of the column data of the stage, whose elements
$\pi_0,\pi_1\in\mathcal O_K$ are those of the column choice, with $\pi_k\bar\pi_k=g_k$; and the column
choice records that every place of $E$ that divides $\bar\pi_k$ lies over $\mathfrak p_k$. Let
$k\in\{0,1\}$, put $\pi=\pi_k$ and $\beta=\bar\pi$, an element of $\mathcal O_K$, and choose a prime
$Q$ of $\mathcal O_E$ over $\mathfrak p_k$. Then $\pi\beta=g_k\in\mathfrak p_k\subset Q$, so $\pi\in Q$
or $\beta\in Q$. If $\beta\in Q$, then $\beta\in Q\cap\mathcal O_K=\mathfrak p_k$. If $\pi\in Q$, then
$\beta=\bar\pi$ lies in the conjugate prime $\bar Q$, which therefore divides $\bar\pi_k$ and lies over
$\mathfrak p_k$; so $\beta\in\bar Q\cap\mathcal O_K=\mathfrak p_k$. In both cases
$\bar\pi_k\in\mathfrak p_k\mathcal O_E$.

(3) The reading of $N$ gives the norm choice of the stage an exponential form, and that form carries
the series tie, the level tie and the completeness of the theta data at the weighted point.

(4) The norm choice records that the choice of the characters $\psi_v$ is compatible with it, and
its characters $\psi_v$ are those that the twist data $T$ choose. Let $v$ be a place, with residue
characteristic $\ell$ and exponent $k$ of $T$ at $v$. The character $\psi_v$ has order dividing
$\ell^k$, and the coherence of $T$ with the base data gives a primitive root of unity $z$ of order
$\ell^k$ in $L$. Every value of $\psi_v$ is a root of unity of order dividing
$\ell^k$, hence a power of the image of $z$ in $\mathbf C$, and so lies in the image of $L$.

(5) This is Lemma~\ref{asm-dependence-covolume-ties-points-fields}(2).

(6) The group data of the stage are true for $f$ (Definition~\ref{asm-dependence-test-assembly-true}(1)):
the Teichmüller characters and the series idempotents at $g_k$ are those of the label datum of $g_k$,
and the character at $(k,v,j)$ is the $2(j+1)$-th power of the label of the character $\omega_k$ that
$f$ chooses, that is, the label of $\omega_k^{2(j+1)}$. The label datum carries the label of a
character of $(\mathbf Z/g_k)^\times$ back to that character. By (1) the frame of the side is $C_f$,
whose character at $g_k$ is $\omega_k$; this is the tame tie.
\end{proof}
% Lean: side_weight_ties_of_ray_field. (1) Plug.wPt_frame: p.x.what2.r.z.lev.C with w.C_eq (cf is
% TC.cf = Fr.at (st Z) f by rfl), and lev.is. (2) Plug.tied_cfOf: lawsJ = z.geo.fw.laws; CorePin from
% the CoreSlot hypothesis, which coreCondition_RX gives for the slot of p (RXs: laws, canon, rank,
% pb); field = T.fieldTies_of_rayField_of_ksExact with w.exact (Lemma -fields (1)); half =
% T.conjGenerators_mem_framePrimes_trueColumns = AtMember.conj_gen_mem_frame after p.x.errX (errX of
% O is toErrXS of the column data; its piG is w.px.piG through levelZeroCompatibility, t.C.piG and
% t.std), with w.ok.pibar (Q' in vG k iff Q'.under K = pb k) and Ideal.exists_ideal_over_prime_of_isIntegral.
% (3) the Export (ExportTie, LevelIs, AllData) is (nu.rd ... p.x.nx).2.1 of the exponential norm
% choice of the reading. (4) PsiChoiceCompatible = p.x.n.chrPsi; values from twDat_psi_valuesIn_LL
% (Coh.roots: a primitive root of order resChar^k; T.psi_ord; eq_pow_of_pow_eq_one) after p.x.n.psi.
% (5) sideW_export_small_alt with small_frame_at_column = SmallExact.smallHas_of_rayField. (6)
% types_of_true (tameTie_of_typesTrue: teich, ser, isCharLab through thetaOf_of_isCharLab) with the
% true group p.x.what2.grp and the frame equality of (1).

\begin{lemma}\label{asm-dependence-covolume-ties-points}
Let local Langlands maps be given, let $N$ be an extension stronger than the law extension
(Definition~\ref{asm-dependence-parameters}), and let $p=p(N)$. Say that the core of a stage
satisfies the \emph{ray condition} if, for every finite place $w$ and every finite set $S$ of primes
$P$ of the host field $E$ of its datum, each of which divides $p_{de}p'_{de}$ and has a residue field
of prime order or lies in the set $Q_w$ of the norm choice of the core at a place $w$ with
$\mu_w\ne\mu_0$, the coefficient field $L$ contains the ray class field of $K=\mathbf Q(\sqrt{-d})$ of
modulus $M=g_0g_1\prod_{P\in S}|k_P|$, in this sense: every prime $\mathfrak P$ of $L$ prime to $M$,
over the prime $\mathfrak p$ of $K$, has $N\mathfrak P=N\mathfrak p^f$ for some $f\ge1$ such that
$\mathfrak p^f$ has a generator congruent to $1$ modulo $M$.
\begin{enumerate}
\item Assume the existence of ray class fields (Theorem~\ref{lit-bundle-unitary}(5)). Eventually in
  the context and at large base data, every core built from $p$ satisfies the ray condition.
\item Assume the existence of ray class fields. Eventually in the context, for every bundle that
  $\Lambda_p$ links to a base point of a good core of $p$ that satisfies the ray condition, and every
  comparison part of the bundle, the point of the part satisfies the intrinsic conditions of its side
  and its rays, it satisfies the further conditions of Definition~\ref{is-point}, and the datum of
  the base point with its size data satisfies the coefficient field bound (Situation~\ref{cv-point}).
\item For every bundle at the datum of a stage built from $p$, with a witness of its construction
  from that stage, and every comparison part of the bundle, the point of the part has the normalized
  embedding of its setting, and if the divisor conditions of the part hold, the point satisfies the
  class-sum pin.
\item Let a bundle be linked to a base point of a core, and let a comparison part of it be given whose
  point satisfies the intrinsic conditions of its side and its rays, at a datum that satisfies the
  coefficient field bound. Then the point satisfies the compatibility conditions of
  Definition~\ref{is-point}.
\end{enumerate}
\end{lemma}
\leanref{Development/AbelianVarietyTheory/Proofs/HeightProofs/LargeBaseRayCondition.lean}{GlobalTheory.Plug.largeBaseProperty_rayProperty}
\leanref{Development/AutomorphicArithmetic/ThetaCorrespondence/ThetaPairLifts/ThirdPlaceSeries/HasExpOfLayer.lean}{GlobalTheory.ExpTiesOfLayer.hasExp_of_layer_alt}
\leanref{Development/AutomorphicArithmetic/ThetaCorrespondence/ThetaPairLifts/ThirdPlaceSeries/SidePtPartOfPackStatement.lean}{GlobalTheory.T.exponent_reading_of_hasExp_bundle}
\leanref{Development/AutomorphicArithmetic/ThetaCorrespondence/ThetaPairLifts/ThirdPlaceSeries/RowEPartTies.lean}{GlobalTheory.PlugPt.rowE_part_ties_alt}
\leanref{Development/AutomorphicArithmetic/ThetaCorrespondence/ThetaPairLifts/ThirdPlaceSeries/ClassSumPinOfReadAtStatement.lean}{GlobalTheory.T.class_sum_divisibility_of_extensionWithLaw}
\leanref{Development/AutomorphicArithmetic/ThetaCorrespondence/ThetaPairLifts/ThirdPlaceSeries/TiesOfTiesB.lean}{GlobalTheory.PlugPt.ties_of_tiesB_alt}
\begin{proof}
(1) Let base data with a framing be given, and let $\Xi$ be the set of characters that the framing
and the twist data give (Definition~\ref{asm-stage-construction}(2)). The primes of $E$ that divide
$p_{de}p'_{de}$, together with the primes at which some character of $\Xi$ is ramified, form a finite
set $\mathcal R$, since $p_{de}p'_{de}\ne0$ and $\Xi$ is a finite set of characters, each ramified at
finitely many primes. Let $\mathcal M$ be the finite set of the moduli $M_S=g_0g_1\prod_{P\in
S}|k_P|$ for the subsets $S$ of $\mathcal R$; they are nonzero.

Let the core of a stage over these data be given, with twist data whose characters are $\Xi$, and
let $w$ and $S$ be as in the ray condition. Every $P\in S$ lies in $\mathcal R$. Either $P$ divides
$p_{de}p'_{de}$, or $P$ lies in the set $Q_w$ of the norm choice at a place $w$ with $\mu_w\ne\mu_0$.
Such a place is twisted, and there the norm choice of the core takes the set $Q_w$ of the twist data.
A prime $P\in Q_w$ prime to $p_{de}p'_{de}$ is ramified for the twisting character $\xi_w$ of $w$,
which lies in $\Xi$: $\xi_w$ is $\delta\psi_w$ or $\delta'\psi_w$, the character $\psi_w$ is ramified at
the primes of $Q_w$, and $\delta$ and $\delta'$ are unramified at $P$. So $M_S\in\mathcal M$.

For $M\in\mathcal M$, the existence of ray class fields gives the ray class field $R_M$ of $K$ of
modulus $M$: an abelian extension of $K$, unramified outside $M$, such that an ideal $I$ prime to $M$
is killed by every character of its Galois group exactly when $I\beta=\alpha\mathcal O_K$ for some
$\alpha,\beta$ prime to $M$ with $\alpha\equiv\beta\pmod M$. Such an ideal has a generator congruent to
$1$ modulo $M$. Embed $R_M$ in $\mathbf C$ over the embedding of $K$ through
$E$, and let $A_M$ be the image of a basis of $R_M$ over $\mathbf Q$, a finite set of algebraic
numbers. If the coefficient field $L$ contains $A_M$, then its image in $\mathbf C$ contains $R_M$,
and $L$ contains a copy of $R_M$ over $K$. Then $L$ satisfies the condition for $M$: for a prime
$\mathfrak P$ of $L$ prime to $M$, over $\mathfrak p$ of $K$, with residue degree $f$, the prime $\mathfrak
p$ is unramified in $R_M$ and its Frobenius has order dividing $f$, so every character kills
$\mathfrak p^f$; and $N\mathfrak P=N\mathfrak p^f$. Take for the large-base condition the algebraic
numbers $\bigcup_{M\in\mathcal M}A_M$, which depend only on the base data, the framing and $\Xi$, and
no primes. Every core of a stage whose coefficient field contains them satisfies the ray condition,
in every context.

(2) A stage that satisfies $N$ satisfies the law extension. Its real expansion data give its norm
choice an exponential form, built from their series, the levels and level characters of the group of
the core and the theta data of the later stage; this gives a reading of the extension. The point of
every comparison part of a bundle that $\Lambda_p$ links to a base point is the pinned point of this
reading at the stage of the bundle. Indeed, let the bundle be built at a stage with its witness, and
let a comparison part be given. The side of its point consists of the data of the real package of the
bundle: the group, the family, the core, the realization, the level group and character at level $0$,
the group isomorphism, the base place, the rational structure and the real expansion, the character
data and the tied tuple of the part. The witness identifies the real package of the bundle with that
of its real input. The realization of its twisted input is the realization built at the stage, with
the scalars $\lambda$ and the error terms of the twist arguments; its base place is that of the stage;
the witness records that the rational structure, the forms and the maps of the real expansion of the
real input are those of the row data of the core; and the tied tuple of the part is that of the stage,
since the character data of the realization of the bundle are those of the column choice. So the
side of the point is the side of the stage at its weighted point, with the archimedean scalar
$\kappa$ of the realization built at the stage, the scalars $\lambda$ and the error terms of the twist
arguments in place of its own: the point is the pinned point.

The intrinsic conditions of a side are its ties: the frame of the realization is the frame of the
side; the core is tied to this frame, which includes the field ties
of the datum; the real expansion has a series tie; the group isomorphism carries the level group and
the character $\psi_0^{-1}$ of the side to the subgroup $U$ and the characters $\psi_U$ of the
twisting data; the theta data are complete; the choice of the characters $\psi_w$ is compatible with
the norm choice, and $\psi_w$ takes values in $L$ when $\mu_w\ne\mu_0$; the small places are as the
frame requires; and the types have the tame tie. The rays of a side say that for every tuple $t$ of
the realization, $L$ contains the ray class field of $K$ of modulus $g_0g_1\prod q$, the product over the
distinct orders $q$ of the residue fields of the places of $Q_t$, in the sense of the ray condition.

Let the good core satisfy the ray condition, and let the point be the pinned point. Its side is the
side of the stage at its weighted point with $\kappa$, $\lambda$ and the error terms replaced. The
rays and the ties do not involve these data, except the series tie, whose conditions on the scalars
are copied unchanged; so it suffices to prove the ties and the rays at the weighted point. The ray
condition of the core is the ray condition for the norm choice of the base point.

Rays. Let $t$ be a tuple of the realization. If $t$ is the base tuple, $\xi_t=1$ is unramified, so
$Q_t$ is empty, and the ray condition, for any place and the empty set, gives the ray class field of
modulus $g_0g_1$. Otherwise $t=\mu_w$ for a place $w$ of the set of twisted places of the stage. Let
$P\in Q_t$; then $\xi_t$ is ramified at $P$, and $k_P$ has prime order, as the level data of the
realization require. If $P$ divides $p_{de}p'_{de}$, it is of the first kind in the ray condition.
Otherwise, $\xi_t$ is the twisting character $\delta_\ell\psi_w$ of $w$ for the framing, with
$\delta_\ell\in\{\delta,\delta'\}$ unramified at $P$; so $\psi_w$ is ramified at $P$, and $P$ lies in
the set $Q_w$ of the twist data, which is the set $Q_w$ of the norm choice at $w$. Moreover
$\mu_w\ne\mu_0$, since otherwise the twisting character of $w$ would be $\xi_{\mu_0}=1$. So $P$ is of
the second kind, and the ray condition for $w$ and the set $Q_t$ gives the ray class field of modulus
$g_0g_1\prod q$.

Ties. The stage satisfies $N$, and its rows are rows of the construction, as the cores built from $p$
require. So Lemma~\ref{asm-dependence-covolume-ties-points-weighted} gives the ties at the weighted
point.

By Lemma~\ref{asm-dependence-covolume-ties-points-host}(1), eventually in the context the host of the
base point satisfies the size bound for $L=x^{\varepsilon}$. The field ties, part of the tie of the
core to the frame, give the hypothesis of part (2) of that lemma, which gives the first part of the
coefficient field bound: $\log[F_O:\mathbf Q]\le P(L)$ and $\log|d_{F_O}|\le[F_O:\mathbf Q]P(L)$. The
second part, the tie of the series idempotents of the datum to its labels, is part of the tame tie of
the types. The further conditions are that the places of the single indices are twisted, which the
part records; the core ties, which Lemma~\ref{asm-dependence-covolume-ties-points-core} gives from
the tie of the core to the frame, the conditions of the level system of the realization of the
point that its weak characterization records, and the first part of the coefficient field bound;
and the remaining ties and the rays.

(4) The point of a comparison part is realized, and the base point records the twist conditions at
all places for its norm choice. The intrinsic conditions, with the first part of the coefficient
field bound, give the core condition of the point, which is the core ties of
Lemma~\ref{asm-dependence-covolume-ties-points-core}, its level and data conditions and one further
agreement of the point with the columns of its realization, the equality of the primes of the core
with those of the frame; and they contain its remaining conditions on types,
series, characters, rays and roots. With the realization and the twist conditions, these are the
compatibility conditions.

(3) Let the bundle be built at a stage, with its witness, and let a comparison part be given. The
witness carries a spectral package for the bundle, whose pin says that the embedding $\iota_E$ of the
Hida setting of the spectral matching of the bundle is the identity; this is the normalized embedding
of the setting of the point.

Now suppose that the divisor conditions of the part hold. The stage satisfies $N$, hence the law
extension, and its real expansion data contain a theta condition for the reading of the rows of the
stage and the real expansion of the row data of the core. It says: at every place $w$ with
$\mu_w\ne\mu_0$, every level $i+1$ and every $x$ in the index space of level $i+1$ at the tuple
$\mu_w$, if an element $\sum_ka_kf_k$ of the rational forms of level $i+1$ is carried by the real
expansion and the map $F_W$ of level $i$ to the image of $\theta_P(j_{M_t}(x))$ under the reading,
then the theta value of $e_M(x)$ for the norm choice at $w$ is $\sum_k\mathrm{emb}_w(a_k)\otimes
e_F^{-1}(1\otimes f_k)$. The twisted input has the reading of the stage, and the witness records that
the real expansion of its real input has the rational structure, the forms and the maps of the row
data of the core; so the theta condition holds for the reading of the twisted input and the real
expansion of the real input. The lift of the twisted input reads the maps $j_W$ of its division pair
through the maps $j_{D,L}$ of the realization. The witness identifies the real package of the bundle
with that of its real input, so the realization of the part has a reading with these two properties.
The class-sum pin at $(w,i+1)$, a tuple $t$ with $t=\mu_w$ and coefficients $c$, for
$x=e_{M_t}(c)$ in the index space and an element $\sum_ka_kf_k$ whose image under the real expansion
is $j_{D,L}$ of $y=\theta_i(t)\,\mathrm{pr}(c)$, asks for the same equality of theta values. The
divisor conditions give that $\theta_P(j_{M_t}(x))$ is $j_W$ of $\theta_i(t)\,\mathrm{pr}(c)$, and the
reading identity carries $j_W$ to $F_W\circ j_{D,L}$; so the hypothesis of the theta condition holds,
and it gives the pin.
\end{proof}
% Lean: (1) Plug.largeBaseProperty_rayProperty (raySets, raySets_alg, rayRec_of_has); (2)
% ExpTiesOfLayer.hasExp_of_layer_alt (the exponential norm choice from y.ser.ex, hg.lvl0, hg.psi0
% and y.xd), ExtensionReading.of RX hE, exponent_reading_of_hasExp_bundle =
% pin_pt_of_bundle_witness, and rowE_part_ties_alt (PlugPt.rowE_tiesB_alt,
% T.eventually_constructedPoint_hostSize_numerical for FieldSmall, pathCoreConditions_of_intrinsicConditions
% for ESerTie, more_of_tiesB with the field muDef of the part); (3)
% class_sum_divisibility_of_extensionWithLaw = ClassSumPinLink.pIota_classSumPin_alt; (4)
% ties_of_tiesB_alt with Ind.pt_realised_and_pb_eq and GeoTies.twAll_pt. The intrinsic conditions
% are TiesB (Ties and Rays of the side point). In rowE_part_ties_alt, FieldSmall comes from
% (H.2 beta).2 applied to ht.1.tied.field, RecTied is FieldSmall with the third component of
% pathCoreConditions_of_intrinsicConditions, and More from more_of_tiesB with muDef 0 0; in
% ties_of_tiesB_alt, pathCoreConditions_of_intrinsicConditions gives core, pb, level and data, and
% TiesB gives types, series, psiVal, ray, roots and rayQ. The pinned point is pin_pt_of_bundle_witness
% (sidePt_part_of_pack_alt along y.eq, side_eq_upd with a.pack, a.v0, sx, fi, eF, lev.C and
% recPin.s.chr; kap = q.x.what2.r.z.R.kap). (3) is pIota_classSumPin_alt: the normalized embedding is
% sP.pin of the spectral package hb.s; the pin is classSumPin_of_readAt_of_pDiv = classSumPin_of_read
% (hth, theta_cast, hD.thTheta, hjW) after readAt_rpkX_of_eMtTheta (x.Z.jW_read) and
% eMtTheta_of_args_alt (eMtTheta_of_pack, hb.sx, hb.eF, hb.fi) from the field th of the real
% expansion data (m.1.th). (1) is RayLarge: TwistLevels.levels_bounded (ramSet finite: pde pde' and
% ram_finite of the characters of T.set f; not_unrAt_xiTw), raySet from
% RayOnto.exists_finset_containsRayField (a Q-basis of the ray field Rf; containsRayField_of_ray with
% exists_frob_pow_eq_one, hart and exists_gen_of_ray; containsRayField_comp), rayRec_of_has with
% c.mu.off and c.nch_Q; the ray condition is RayRec, i.e. ContainsRayField (embK) M, stated in the
% lemma. Written out in (2): rowE_tiesB_alt top level (RowE.of_all; rayClassFields_in_coefficients
% from the slot RayPosition; pinned point = (sideW nu p).upd kap lam err; side_upd_ties_rays =
% exportTie_upd for the series field, the other fields copied, Rays unchanged); the rays are
% Plug.RaysW.rays_of_rayRec_alt (mem_range_ixOf, trueColumnChoices_heckeChar_reference_eq_one, Q_ram,
% Q_pl = IsTwistPlaces, xiTw with isDel.unr/isDel'.unr, psi_is.ram, n.Q, mu.tw/on/mu0); the size
% bound is Lemma -ties-points-host; FieldSmall is (H.2 beta).2 at ht.1.tied.field, ESerTie is
% hT.types.eSerTie, More is more_of_tiesB. The ties at the weighted point,
% PlugPt.side_weight_ties_of_ray_field (from rowE_tiesB_alt with the CoreSlot hypothesis that
% coreCondition_RX gives for the slot of p), are Lemma -ties-points-weighted; its open part is in
% Lemma -ties-points-fields (1). The core ties, Plug.coreTies_of_tied in more_of_tiesB and
% ties_core_of_tiesB (b.charW.base.lev for the level system), are Lemma -ties-points-core.

\begin{lemma}\label{asm-dependence-covolume-ties}
Assume the existence of ray class fields (Theorem~\ref{lit-bundle-unitary}(5)). Let local Langlands
maps be given, let $N$ be an extension stronger than the law extension
(Definition~\ref{asm-dependence-parameters}), and let $p=p(N)$. If the division rows of open
parameter statement (b) are linked at $p$, then so is the following property of a class and a
selection: at every member, for every comparison part of the selected bundle, the point of the part
has the normalized embedding of its setting and satisfies the class-sum pin, the compatibility
conditions and the further conditions of Definition~\ref{is-point}, and the datum of the member
with its size data satisfies the coefficient field bound (Situation~\ref{cv-point}).
\end{lemma}
\leanref{Development/AutomorphicArithmetic/ThetaCorrespondence/ThetaPairLifts/ThirdPlaceSeries/LinkedTiesStatement.lean}{GlobalTheory.T.class_sum_divisor_all_alt}
\leanref{Development/AutomorphicArithmetic/ThetaCorrespondence/ThetaPairLifts/ThirdPlaceSeries/LinkedTies.lean}{GlobalTheory.GeoTies.linked_ties_alt}
\begin{proof}
Use the ray condition of Lemma~\ref{asm-dependence-covolume-ties-points}, and let $T_0$ be the
property: at every member, for every comparison part of the selected bundle, the point of the part
satisfies the intrinsic conditions of its side and its rays and the further conditions, and the
datum of the member with its size data satisfies the coefficient field bound. Let $Q_p$ be the slot of
the good cores of $p$ and $R_1$ the slot of the cores that satisfy the ray condition. By part (2) of
that lemma, taking for the condition on bundles that every comparison part has these properties,
$T_0$ holds along $(Q_p\cap R_1,\Lambda_p)$; by part (1), every core built from $p$ satisfies $R_1$
eventually in the context and at large base data. So $T_0$ is linked at $p$. By hypothesis the
division rows are linked at $p$, so by Lemma~\ref{asm-dependence-linked}(2) the conjunction of the
division rows at every member and $T_0$ is linked at $p$, with a slot $R$ and a condition $E$ on
bundles.

Let $s$ be a selection for a class $\mathcal A$ whose bundles are built from $(Q_p\cap R,\Lambda_p)$
and satisfy $E$, let $D$ be a member at a scale $x>x_1$, and let a comparison part of the selected
bundle be given. The bundle is linked by $\Lambda_p$ to a base point of a good core of $p$, so it lies
at the datum $D$ of a stage built from $p$, with a witness of its construction. Every tuple of the
part is the base tuple or the tuple $\mu_w$ of a place, so the division rows of the bundle give the
divisor conditions of the part. By Lemma~\ref{asm-dependence-covolume-ties-points}(3), the point of
the part has the normalized embedding of its setting and satisfies the class-sum pin. By $T_0$ it
satisfies the intrinsic conditions and the further conditions, and $D$ satisfies the coefficient
field bound; so by part (4) of that lemma the point satisfies the compatibility conditions. Hence the
property of the lemma holds along $(Q_p\cap R,\Lambda_p)$ with $E$, and it is linked at $p$ with the
slot $R$.
\end{proof}
% Lean: linked_ties_alt: hE from hasExp_of_layer_alt, nu = ExtensionReading.of RX hE, hRow =
% rowE_part_ties_alt for IsGoodE, hlink through exponent_reading_of_hasExp_bundle, hR2 with the slot
% RayPosition and largeBaseProperty_rayProperty, RowL.and with hDiv, then Built.ind with
% FourBoundsAtBundles.p_div_of_div3g (through tupAllUsed_part_alt), class_sum_divisibility_of_extensionWithLaw,
% ties_of_tiesB_alt, Ind.pt_realised_and_pb_eq and twAll_pt.

\begin{lemma}\label{asm-dependence-orbit-geometric}
Assume the existence of ray class fields. Let local Langlands maps be given, let $N$ be an extension
stronger than the law extension, and let $p=p(N)$. If the division rows of open parameter statement
(b) and statement (d) are linked at $p$, then so is statement (g).
\end{lemma}
\leanref{Development/AutomorphicArithmetic/ThetaCorrespondence/ThetaPairLifts/ThirdPlaceSeries/RowGeo0AtStatement.lean}{GlobalTheory.T.linkedAt_geometricWeakTies_of_tie}
\leanref{Development/AutomorphicArithmetic/ThetaCorrespondence/ThetaPairLifts/ThirdPlaceSeries/RowGeo0At.lean}{GlobalTheory.Fc.GeoZeroAt.geometricWeakTies_of_linked_ties}
\begin{proof}
By Lemma~\ref{asm-dependence-covolume-ties}, the property $T$ of that lemma is linked at $p$. By
Lemma~\ref{asm-dependence-linked}(2), so is the property ``$T$ and (d)''. Let $S$ be a slot along
which, with $\Lambda_p$, this property holds with a condition $E$ on bundles. Let $s$ be a selection
for a class $\mathcal A$ whose bundles are built from $(S,\Lambda_p)$ and satisfy $E$, and let $D$
be a member at a scale $x>x_1$. By (d) the selected bundle at $D$ has a comparison part, and by $T$
the point of this part has the normalized embedding of its setting and satisfies the class-sum pin,
the compatibility conditions and the further conditions. This is (g) at $D$. So ``$T$ and (d)''
along $(S,\Lambda_p)$ implies (g) along $(S,\Lambda_p)$, for every slot $S$, and (g) is linked at
$p$ by Lemma~\ref{asm-dependence-linked}(1).
\end{proof}

\begin{lemma}\label{asm-dependence-covolume-points-rows}
Assume the covering of discrete series, the finiteness of Harish-Chandra modules, rapid decay of cusp
forms and the Casselman-Wallach theorem (Theorem~\ref{lit-bundle-unitary}). Let a comparison part of a
bundle at a definite-side datum $D$ be given, with indefinite realization $Q$, realization $R$ and
places $P_g$ (Definition~\ref{is-realisation}). Then:
\begin{enumerate}
\item the theta maps of $R$ exist; $R$ satisfies the evaluation conditions at the embeddings of the
  coefficient field that are not good and has the geometric action; and the evaluations of $R$ at
  the good embeddings are injective;
\item if $8fr+4\le n_V$, where $f=2n$, $r=k+m-1$ and $n_V$ are the numbers of the ranks of $D$ (as in
  Definition~\ref{th-ranks} and Lemma~\ref{is-onto-step-numerics}), then $Q$ has bounded rows (Definition~\ref{is-bounded-rows});
\item the spaces $Y_j$, $X'$, $V_i$ and $W_i$ of $D$ have lattices with adelic bases that satisfy
  the normal conditions at the places of $P_g$.
\end{enumerate}
\end{lemma}
\leanref{Development/AutomorphicArithmetic/ThetaCorrespondence/ThetaPairLifts/ThirdPlaceSeries/DeliveryFourRows.lean}{GlobalTheory.Fc.Ind.atAll_fourRows'}
\leanref{Development/AutomorphicArithmetic/ThetaCorrespondence/ThetaPairLifts/ThirdPlaceSeries/FireStatement.lean}{GlobalTheory.T.indefinite_scale_theta_equality_alt}
\leanref{Development/AutomorphicArithmetic/ThetaCorrespondence/ThetaPairLifts/ThirdPlaceSeries/FramesExOfRealisationStatement.lean}{GlobalTheory.T.framesEx_of_realPins_indefinite_scale}
\begin{proof}
(1) The four results give these four conditions for the indefinite realization of every comparison
part.

(2) By rapid decay of cusp forms, an indefinite realization of a datum with $8fr+4\le n_V$ whose
realization has its theta maps has bounded rows. The theta maps of $R$ exist by (1).

(3) The comparison part records a normalization of $Q$, and the normalization gives these lattices
and bases.
\openstep{The derivation of the four conditions of (1) from the four results, of the bounded rows in
(2) from rapid decay of cusp forms, and of the lattices and bases in (3) from the normalization is not
yet written out.}
\end{proof}
% Lean: atAll_fourRows' (DeliveryFourRows.lean): (1) indefinite_scale_theta_equality_alt (ExAt,
% JEvBar, GeoAct, EvInj); (2) indefinite_scale_theta_bound_right_alt (IndefiniteRealisationBoundedRows,
% from CuspRapidDecay, (Fm.parT x hx i).nV_ge and (1)); (3) framesEx_of_realPins_indefinite_scale
% (Q.FramesEx = Seed.Frm.FramesExistDiagonal, from the field pins of the part); the fourth
% conjunct, indefinite_scale_theta_identity_alt (ThetaFireRow), is not used. Sizes of the open step
% (text dependencies, overcounts): indefinite_scale_theta_equality_alt about 18,400 theorems and
% 320,000 lines; indefinite_scale_theta_bound_right_alt about 18,500 theorems and 329,000 lines;
% framesEx_of_realPins_indefinite_scale about 500 theorems and 10,000 lines; the first two overlap.

\begin{lemma}\label{asm-dependence-covolume-points-export-span}
Let $\mathcal O$ be a domain with fraction field $K$, let $K''$ be a field containing $K$, let $L$ be a
field of characteristic $0$ with an embedding $\sigma\colon L\to K''$, let $F$ be a $K$-vector space
and $W$ an $L$-vector space, and let $\phi\colon W\to K''\otimes_{\mathcal O}F$ be an injective
$\sigma$-semilinear map whose image spans $K''\otimes_{\mathcal O}F$ over $K''$. Let $\Gamma$ be a
finite group acting $K$-linearly on $F$, and let $a(h)$, $h\in\Gamma$, be $L$-linear endomorphisms
of $W$ with $\phi(a(h)x)=(1\otimes h)\phi(x)$ for all $h$ and $x$. Let $P\subset\Gamma$ be a subgroup
with characters $t\colon P\to L^\times$ and $t_F\colon P\to K^\times$, and let $c\colon\Gamma\to L$
and $c_F\colon\Gamma\to K$ be functions, such that $\sigma(t(u))=t_F(u)$ and $\sigma(c(h))=c_F(h)$ in
$K''$. Suppose that $e_F=\sum_hc_F(h)h$ acts on $F$ as an idempotent that commutes with the elements
of $P$. Let $F^{P,t_F}$ be the space of $y\in F$ with $uy=t_F(u)y$ for $u\in P$, and $V=e_F(F^{P,t_F})$;
let $W_{\mathrm g}$ be the space of the $x\in W$ with $a(u)x=t(u)x$ for $u\in P$ that lie in the image
of $\sum_hc(h)a(h)$. Then the $K''$-span of $\phi(W_{\mathrm g})$ is the $K''$-span of $1\otimes V$ in
$K''\otimes_{\mathcal O}F$.
\end{lemma}
\leanref{Development/AbelianVarietyTheory/Proofs/TwistedOperators/LatticeFamilyAdoptionSecond.lean}{GroupAverageFixedPoints.GlobalTheory.CompanionTheory.LatSpan.span_image_good}
\leanref{Development/AbelianVarietyTheory/Proofs/TwistedOperators/LatticeFamilyAdoptionSecond.lean}{GroupAverageFixedPoints.GlobalTheory.CompanionTheory.LatSpan.mem_range_iff_type}
\leanref{Development/AbelianVarietyTheory/Proofs/TwistedOperators/LatticeFamilyAdoptionSecond.lean}{GroupAverageFixedPoints.GlobalTheory.CompanionTheory.LatSpan.mem_range_iff_image}
\leanref{Development/AbelianVarietyTheory/Proofs/TwistedOperators/LatticeFamilyAdoptionSecond.lean}{GroupAverageFixedPoints.GlobalTheory.CompanionTheory.LatSpan.range_ext}
\begin{proof}
Since $\phi$ is injective, $a$ is a homomorphism: $\phi(a(hh')x)=(1\otimes hh')\phi(x)=\phi(a(h)a(h')x)$,
and $\phi(a(1)x)=\phi(x)$. The order of $P$ is nonzero in $L$, which has characteristic $0$, and so in
$K$, since $\sigma$ is injective. On $W$ put $A=|P|^{-1}\sum_{u\in P}t(u)^{-1}a(u)$ and
$E=\sum_hc(h)a(h)$; on $F$ put $A_F=|P|^{-1}\sum_{u\in P}t_F(u)^{-1}u$. Since $\phi$ is
$\sigma$-semilinear and $\sigma$ carries $t$ and $c$ to $t_F$ and $c_F$, $\phi\circ A=(1\otimes
A_F)\circ\phi$ and $\phi\circ E=(1\otimes e_F)\circ\phi$. Since $\phi$ is injective, $E$ is an
idempotent that commutes with the $a(u)$, $u\in P$, as $e_F$ is and does.

Call a vector of $W$ of type $t$ if $a(u)x=t(u)x$ for $u\in P$. The average $A$ fixes the vectors of
type $t$, and $Ax$ is of type $t$ for every $x$: reindexing the sum gives $a(v)A=t(v)A$ for $v\in P$.
So the image of $EA$ is $E$ of the vectors of type $t$. It is also the set of vectors of type $t$ in
the image of $E$: $EAx$ is of type $t$, since $E$ commutes with $P$; and if $x=Ey$ is of type $t$,
then $EAx=Ex=x$. Hence $W_{\mathrm g}$ is the image of $EA$. In the same way $V$ is the image of
$e_FA_F$.

For a $K$-linear endomorphism $T$ of $F$, the image of $1\otimes T$ on $K''\otimes_{\mathcal O}F$ is
the $K''$-span of $1\otimes T(F)$. So the $K''$-span of $1\otimes V$ is the image of $1\otimes
e_FA_F$. As $\phi(W)$ spans, this image is the $K''$-span of $(1\otimes e_FA_F)\phi(W)=\phi(EA(W))=
\phi(W_{\mathrm g})$.
\end{proof}

\begin{lemma}\label{asm-dependence-covolume-points-export-lattice}
Let $\mathcal O'$ be an integrally closed domain with fraction field $K''$, let $L$ be a number field
with an embedding $\sigma\colon L\to K''$, let $M$ be an $L$-vector space with a finitely generated
$\mathcal O_L$-submodule $\Lambda$ that spans $M$ over $L$, let $T$ be a $K''$-vector space, and let
$\phi\colon M\to T$ be $\sigma$-semilinear. Let $V''$ be the $K''$-span of $\phi(M)$. Then the
$\mathcal O'$-span $N$ of $\phi(\Lambda)$ is a lattice in $V''$: it is a finitely generated
$\mathcal O'$-module contained in $V''$, and for every $x\in V''$ there is a nonzero $c\in\mathcal O'$
with $cx\in N$.
\end{lemma}
\leanref{Development/AbelianVarietyTheory/Proofs/TwistedOperators/LatticeFamilyAdoptionSecond.lean}{GroupAverageFixedPoints.GlobalTheory.CompanionTheory.LatSpan.isLatIn_span}
\leanref{Development/AbelianVarietyTheory/Proofs/TwistedOperators/LatticeFamilyAdoptionSecond.lean}{GroupAverageFixedPoints.GlobalTheory.CompanionTheory.LatSpan.integers_embed}
\leanref{Development/AbelianVarietyTheory/Proofs/TwistedOperators/ImportsFurther.lean}{GlobalTheory.SmulMemOfMemSpanExists.exists_smul_mem_of_mem_span}
\begin{proof}
For $a\in\mathcal O_L$, $\sigma(a)$ is integral over $\mathbf Z$, hence over $\mathcal O'$, so it lies
in $\mathcal O'$. So $\phi(a\lambda)=\sigma(a)\phi(\lambda)$ lies in the $\mathcal O'$-span of
$\phi(\lambda)$, and $N$ is spanned by the images of finitely many generators of $\Lambda$. It is
contained in $V''$. Since $\Lambda$ spans $M$ over $L$ and $\phi$ is semilinear, $\phi(M)$ lies in the
$K''$-span of $N$, so $V''$ is the $K''$-span of $N$. An element $\sum_k(a_k/b_k)n_k$ of it, with
$a_k,b_k\in\mathcal O'$, $b_k\ne0$ and $n_k\in N$, is carried into $N$ by $c=\prod_kb_k$.
\end{proof}

\begin{lemma}\label{asm-dependence-covolume-points-export}
Let a comparison part of a bundle at a definite-side datum be given, and suppose that its divisor
conditions hold and that its point satisfies the class-sum pin (Definition~\ref{is-point}). Then the
point satisfies the export condition: the real expansion of its indefinite realization has a majorant
for the norm choice of the core.
\end{lemma}
\leanref{Development/AutomorphicArithmetic/ThetaCorrespondence/ThetaPairLifts/ThirdPlaceSeries/PNmjOfPDivOfClassSumPin.lean}{GlobalTheory.Fc.Ind.pNmj_of_pDiv_of_classSumPin_alt}
\leanref{Development/AutomorphicArithmetic/ThetaCorrespondence/ThetaPairLifts/ThirdPlaceSeries/NmjTAtTiedStatement.lean}{GlobalTheory.T.lattice_family_adoption_fourth_alt}
\begin{proof}
The norm choice of the core is tied to the datum, and the comparison part records the normalization
and the recursion of its indefinite realization $Q$. The real expansion of $Q$ consists of a rational
structure, with spaces $W_L(j)$ over the coefficient field $L$ and $K_s$-subspaces $A_j$, of the forms
$f_i$, which identify $W_L(i+1)$ with the rows of $Q$ at level $i$, and of maps $e_F$ at the places
of the norm choice. A majorant for the norm choice consists of linear maps $j_D$ from the spaces $M_j$
of the division data to $W_L(j)$ that agree with the maps $j_{D,L}$ of $Q$ through the forms $f_i$
and the identifications $e_{W,i}$ of the lift, and of maps $e_M(w,j)$, for every place $w$ and level
$1\le j\le m$, from $M_j$ to the space $\operatorname{Frac}(\mathcal O_w)\otimes F_m(w,j)$ of the norm
choice at $w$, semilinear over its embedding $\mathrm{emb}_w$ of $L$, such that:
\begin{enumerate}
\item[(a)] each $A_j$ is a $K_s$-form of $W_L(j)$;
\item[(b)] if $j_D(x)=\sum_ka_kf_k$ with $a_k\in L$ and $f_k\in A_j$, then
  $e_M(w,j)(x)=\sum_k\mathrm{emb}_w(a_k)\otimes e_F^{-1}(1\otimes f_k)$;
\item[(c)] on $K_s$, $\mathrm{emb}_w$ is the embedding of $K_s$ into the fraction field of the
  twisting ring at $w$;
\end{enumerate}
and, at the places $w$ with $\mu_w\ne\mu_0$:
\begin{enumerate}
\item[(d)] $e_M(w,j)$ carries a basis of $M_j$ to linearly independent vectors;
\item[(e)] the span of the image of the lattice of $M_j$ is a lattice in the base change of the
  space $V_0$ of the norm choice at $w$;
\item[(f)] $e_M(w,j)(y)$ is the theta value of $e_{M}(x)$ for the norm choice at $w$ whenever $x$
  lies in the index space of level $j$ at the tuple of $w$ and $j_M(y)=\theta_P(j_{M_t}(x))$.
\end{enumerate}

\emph{The maps.} On $M_{i+1}$ put $j_D=f_i^{-1}\circ j_{D,L}\circ e_{W,i}^{-1}$, read through the
equality of the dimensions of the rows of $Q$ and of the lift. It agrees with $j_{D,L}$ by
construction, and it is injective, since $j_{D,L}$ is the coordinate map of a basis and the other
factors are isomorphisms. The real expansion carries a $K_s$-form of the rows of $Q$ at level $i$, and
$f_i$ carries it to $A_{i+1}$; this is (a). The tie of the norm choice to the datum gives (c). Fix $w$
and $j$. Since $A_j$ is a $K_s$-form and (c) holds, the map $f\mapsto e_F^{-1}(1\otimes f)$ on $A_j$
extends to a map $e$ on $W_L(j)=L\otimes_{K_s}A_j$, semilinear over $\mathrm{emb}_w$, with
$e(\sum_ka_kf_k)=\sum_k\mathrm{emb}_w(a_k)\otimes e_F^{-1}(1\otimes f_k)$. It carries a $K_s$-basis
$(b_k)$ of $A_j$, which is an $L$-basis of $W_L(j)$, to the vectors $1\otimes e_F^{-1}(1\otimes b_k)$,
which form a basis of $\operatorname{Frac}(\mathcal O_w)\otimes F_m(w,j)$; so $e$ carries bases to
bases, and linearly independent families to linearly independent families. Put
$e_M(w,j)=e\circ j_D$. Then (b) holds, and (d) holds since $j_D$ is injective.

\emph{(f).} Let $j=i+1$, and let $x$ and $y$ be as in (f). The recursion of $Q$ gives a tuple $t$ of
the family whose tuple is that of $w$, and $x=e_{M_t}(c)$ for coefficients $c$. Since $A_j$ spans,
$j_{D,L}(\theta_i(t)\,\mathrm{pr}(c))=f_i(\sum_ka_kf_k)$ for some $a_k$ and $f_k$. The class-sum pin
says that the theta value of $e_M(x)$ is $\sum_k\mathrm{emb}_w(a_k)\otimes e_F^{-1}(1\otimes f_k)$.
The divisor conditions and the tie of the division data to the lift give
$j_M(e_{W,i}(\theta_i(t)\,\mathrm{pr}(c)))=\theta_P(j_{M_t}(x))=j_M(y)$, and $j_M$ is injective, so
$y=e_{W,i}(\theta_i(t)\,\mathrm{pr}(c))$ and $j_D(y)=\sum_ka_kf_k$. By (b), $e_M(w,j)(y)$ is the same
sum.

\emph{(e).} Fix $w$ with $\mu_w\ne\mu_0$ and $j=i+1$. Let $\iota$ and $b$ be the twisting index and
the sign that the norm choice gives at $w$, let $\mathcal O=\mathcal O_\iota$, with fraction field
$K$, and let $K''$ be the fraction field of $\mathcal O_w$. The space $F=F_m(w,j)$ is the space
$F_{\iota,j,b}$ of the twisting data (Definition~\ref{th-twisting-data}(5)), with its action of $G_0$,
the subgroup $U$, the character $\psi_U=\psi_{U,\iota,b}$ and the central idempotent $e_{\iota,b}$;
and $V_0=e_{\iota,b}F^0$ is the space of item (7) there. By (b), written in the basis $(b_k)$,
$e_M(w,j)=e\circ j_D$ for the extension $e$ above, which is injective and whose image spans
$K''\otimes_{\mathcal O}F$.

The extension $e$ is equivariant. The real expansion has a group $G_{\mathrm r}$ with an isomorphism
$s\mapsto g_s$ onto $G_0$, which acts on $W_L(j)$ by $\rho$, preserving $A_j$, with a character
$\chi$ of $G_{\mathrm r}$ with values in $L^0$; and $e_F$ carries the action of $g_s$ on $F$ to
$\chi(s)^{-1}\otimes\rho(s)$ on $K\otimes_{K_s}A_j$, where $\chi(s)^{-1}$ is read in $K$ through the
embedding of $L^0$ of the sign $b$ (Definition~\ref{th-twisting-data}(1)). The tie of the norm
choice to the datum gives that on $L^0$, $\mathrm{emb}_w$ is this embedding followed by the
inclusion $K\subset K''$. So $g_se_F^{-1}(1\otimes b_k)=\chi(s)^{-1}e_F^{-1}(1\otimes\rho(s)b_k)$, and
for $x=\sum_ka_kb_k$,
\[
e\bigl(\chi(s)^{-1}\rho(s)x\bigr)=\sum_k\mathrm{emb}_w\bigl(\chi(s)^{-1}a_k\bigr)\otimes
e_F^{-1}\bigl(1\otimes\rho(s)b_k\bigr)=(1\otimes g_s)\,e(x).
\]

The forms $f_i$ consist of a linear isomorphism of $W_L(j)$ onto the space of the forms $F_W$ of
level $i$ and an isomorphism of $G_{\mathrm r}$ onto the group $G$ of the realization of $Q$, which
carry $\rho$ and $\chi$ to the action $\rho'$ and the character $\chi'$ of $F_W$. By the
normalization, the isomorphism $I\colon G\to G_0$ of the realization, composed with the isomorphism of
groups of $f_i$, is $s\mapsto g_s$. So $\phi=e\circ f_i^{-1}$ is injective, its image spans
$K''\otimes_{\mathcal O}F$, and $\phi(\chi'(h)^{-1}\rho'(h)y)=(1\otimes I(h))\phi(y)$ for $h\in G$.

Apply Lemma~\ref{asm-dependence-covolume-points-export-span} with $\sigma=\mathrm{emb}_w$, $\Gamma=G$
acting on $F$ through $I$, $a(h)=\chi'(h)^{-1}\rho'(h)$, $P=P_0$, $t=\psi_0^{-1}$, $t_F=\psi_U\circ
I$, $c(h)=e_X(h)$, the coefficients of the idempotent of the characters of the family, and $c_F(h)$
the coefficient of $e_{\iota,b}$ at $I(h)$. The division pair is tied to the group of the datum: $h\in
P_0$ exactly when $I(h)\in U$; for $u\in P_0$, $\psi_0(u)^{-1}$ lies in $L^0$ and its image in $K$
under the embedding of the sign $b$ is $\psi_U(I(u))$; and the image of $e_X(h)$ under the embedding
of $L$ in $\mathbf C$ is the coefficient of the type idempotent $e_0$ of the datum at $I(h)$. The
conditions on the datum give elements $c_0(g)\in L^0$, $g\in G_0$, whose images in $\mathbf C$ are the
coefficients of $e_0$ and whose images in $K$ under the embedding of the sign $b$ are those of
$e_{\iota,b}$; so $e_X(h)=c_0(I(h))$. With the tie of $\mathrm{emb}_w$ on $L^0$, this gives
$\mathrm{emb}_w(t(u))=t_F(u)$ and $\mathrm{emb}_w(c(h))=c_F(h)$. The operator $\sum_hc_F(h)I(h)$ is
the action of $e_{\iota,b}$, a central idempotent of $K[G_0]$. Since $h\in P_0$ exactly when $I(h)\in
U$, the space $F^{P_0,t_F}$ is $F^0$, and $V=V_0$. The space $W_{\mathrm g}$ consists of the forms $y$
with $\rho'(u)y=\chi'(u)\psi_0(u)^{-1}y$ for $u\in P_0$ that lie in the image of
$\sum_he_X(h)\chi'(h)^{-1}\rho'(h)$: it is the good part of $F_W$ for $(P_0,\psi_0^{-1})$ and the
idempotent of the characters of the family, which is the space of rows of $Q$ at level $i$. So the
$K''$-span of $\phi$ of the rows of $Q$ at level $i$ is the $K''$-span of $1\otimes V_0$.

The image of $j_D$ is $f_i^{-1}$ of the rows of $Q$ at level $i$, since $f_i\circ
j_D=j_{D,L}\circ e_{W,i}^{-1}$, $j_{D,L}$ is the coordinate map of a basis of the rows and $e_{W,i}$
is an isomorphism. So the image of $e_M(w,j)=\phi\circ f_i\circ j_D$ is $\phi$ of the rows, and its
$K''$-span is the base change of $V_0$. The lattice of $M_j$ is a finitely generated module over the
integers of $L$ that spans $M_j$, and $\mathcal O_w$, a discrete valuation ring, is integrally closed.
By Lemma~\ref{asm-dependence-covolume-points-export-lattice}, the $\mathcal O_w$-span of the image of
the lattice of $M_j$ is a lattice in the base change of $V_0$. This is (e).
\end{proof}
% Lean: pNmj_of_pDiv_of_classSumPin_alt = lattice_family_adoption_fourth_alt with b.pins, b.Rr and
% O.ex.nch_tie (NChTie); the conclusion NormJumpPointSummands is NX.RealExpansionMajorant, the
% field export_ of PointConditions' (PtAt.ExportRow). Its proof is PtAt.nmjT_at_tied_alt: jD, eM and
% (a) to (d) are lattice_family_adoption_first_alt = ReadingExt.reading_ext_alt (embeddingAtLevel,
% jDa_injective_alt, isKsForm_of_formsIso with B.fw.ks, exists_ext_alt = SemilinearExt.
% semilinearExtension with li_of_carries, embW_Ks_alt from FractionFieldEmbeddingCompat.embCompat and
% NChTie); (f) is lattice_family_adoption_third_alt = ClassSumLift.eM_eq_thPe_alt = at_index_alt
% (Qr.tupTie, exists_sum_smul_eq_of_isKsForm, the class-sum pin hP, Idx.jM_inj with P.tied.mjH.j_eq
% and hD.thTheta); (e) is lattice_family_adoption_second_alt = LatSpan.lattice_of_lifts_alt
% (LatticeFamilyAdoptionSecondAltAlt.lean; the same proof as LatSpan.lattice_of_lifts): hC from
% P.tied.div.grp.e0 and D.nStage_trivialCharacter_coeff_eq; LatSpan.stage2 with P.I.eG, B.Q.Sx,
% B.Q.nexp.fi i, hgeq from hPin.eGtie and B.Q.nexp.geq, rho, chi, eF, act, P.P0, P.psi0^-1, hgrp.lvl0
% and hgrp.psi0; stage2 is stage1 (semilinearExtension, inj_span_of_carries, ext_intertwine with
% embW_eq_of_tie) followed by span_image_good (Lemma -export-span) with hEq (NForms.eG), eAlg_idem,
% eAlg_central, NForms.Fm0 and FormsL.mem_typeP; heM from hSum (b); hrange from hjD and hWg (rfl:
% the rows at level i are (FW i).Wgood P0 psi0^-1 (idemOf X).coeff); then isLatIn_span with
% integers_embed (Lemma -export-lattice) and Y.MjH_good. Size of (e): LatticeFamilyAdoptionSecond.lean,
% 773 lines, and the 190-line AltAlt file.

\begin{lemma}\label{asm-dependence-covolume-points}
Assume the covering of discrete series, the finiteness of Harish-Chandra modules, rapid decay of cusp
forms and the Casselman-Wallach theorem (Theorem~\ref{lit-bundle-unitary}). Let $\mathcal A$ be a
class of families, with working subclass $\mathcal W(\mathcal A)$
(Theorem~\ref{pd-period-integrality-along-families}), and choose at every member of every family of
$\mathcal A$, at every scale beyond its threshold, a bundle and a comparison part of it, the
\emph{selected part}. Suppose that at every member the realization of the selected part has
primitive rows and its family of definite forms has lift-distinct types, the divisor conditions of
the selected part hold, and the point of the selected part satisfies the class-sum pin and the
compatibility conditions of Definition~\ref{is-point}. Suppose also that eventually along
$\mathcal W(\mathcal A)$, eventually along each family, every row of the definite-side datum is good.
Then eventually along $\mathcal W(\mathcal A)$, eventually along each family, the point of the
selected part satisfies the point conditions (Situation~\ref{cv-point}).
\end{lemma}
\leanref{Development/AutomorphicArithmetic/ThetaCorrespondence/ThetaPairLifts/ThirdPlaceSeries/EvWPRowsOfRows.lean}{GlobalTheory.Fc.Ind.evW_pRows_of_rows_alt}
\begin{proof}
The point conditions of the point of a comparison part with indefinite realization $Q$ and
realization $R$, at a datum $D$, are: (i) the four conditions of
Lemma~\ref{asm-dependence-covolume-points-rows}(1); (ii) the two conditions of bounded rows of $Q$;
(iii) primitive rows of $R$ and lift-distinct types of its family; (iv) the family is ok at level
$1$; (v) the basis conditions of $R$ at the places $P_g$; (vi) the divisor conditions of the part;
(vii) the tie of the characters $\psi_0$ and $\psi_s$ of the part to $D$; (viii) the dimension $D_0$
of the base rows of $D$ is that of the definite subspace of the part at the base tuple, at the level
of $(P_0,\psi_0)$; (ix) the export condition of
Lemma~\ref{asm-dependence-covolume-points-export}; (x) every row of $D$ is good; (xi) a finite place
$w$ is twisted whenever, for some lift $j$, one of $I_w(j,\mu_0)$, $I_{\bar w}(j,\mu_0)$, the relative
indices of $T_{j,\mu_0}$ and its dual at $w$ and at $\bar w$, $F_0(w,j)$ or $F_0(\bar w,j)$ is nonzero,
with $F_0$ that of $N(0,0)$, as in the further conditions of Definition~\ref{is-point}; and (xii)
the compatibility conditions.

Let a family of $\mathcal W(\mathcal A)$ and a member at a scale $x>x_1$ be given, with its selected
part. The parameter choices of the member give $8fr+4\le n_V$ (as at the end of the proof of
Lemma~\ref{is-onto-step-numerics}). By Lemma~\ref{asm-dependence-covolume-points-rows}, (i) and (ii)
hold, and the lattices with normal adelic bases of part (3) of that lemma give (v). Conditions (iii), (vi) and (xii) are
hypotheses, and by Lemma~\ref{asm-dependence-covolume-points-export}, (vi) and the class-sum pin give
(ix). Conditions (iv), (vii) and (xi) are recorded by the comparison part. For (viii): the base rows
of $D$ have an orthonormal basis whose coordinates identify them with the row side $S_0$ of the
base level of the part, so $D_0$ is the dimension of $S_0$; and $S_0$ is the row side of $Q$, defined
on the definite subspace of the part at level $(P_0,\psi_0)$, so its dimension is that of this
definite subspace at the base tuple. So all conditions other than (x) hold at every member of every
family of $\mathcal A$ at every scale $x>x_1$.

By hypothesis, eventually along $\mathcal W(\mathcal A)$ every family has a threshold $x_0$ such that
(x) holds at every member at every scale $x\ge x_0$, $x>x_1$. Beyond it, all the point conditions
hold, which is the conclusion, with the same $\varepsilon_1$ and largeness thresholds.
\end{proof}
% Lean: evW_pRows_of_rows_alt: atAll_fourRows' (Lemma asm-dependence-covolume-points-rows), then
% Closed.imp of the hypothesis on GoodFTAll, with bundle_lines_of_primitive_frames_div_nmj =
% pRows_of_parts_alt, the structure PointConditions' with the fields exAt, jEvBar, geoAct, evInj
% (i), ioExt, oneMap (ii), prim, sep (iii), ok = b.ok1 (iv), framesOK = b.B.Q.frames hFr (v),
% divRows (vi), psi = b.psi (vii), d0 = ind_piece_degree_zero = base_level_rows_eq_dimension_alt
% (rowSide_card_eq_of_coord with onb0 and eRR0, then b.B.Q.zs and b.B.Q.s0) (viii), export_ from
% pNmj_of_pDiv_classSumPin (ix), goodFT (x), w1 from b.muDef 0 0 (xi), ties (xii).

\begin{lemma}\label{asm-dependence-covolume-bound}
Assume the conclusion of Theorem~\ref{cv-covolume-bound}. Let $\mathcal A$ be a class of families
each of which carries its duty data (Definition~\ref{asm-duty}), and choose selected parts as in
Lemma~\ref{asm-dependence-covolume-points}, with the working subclass $\mathcal W(\mathcal A)$.
Suppose that eventually along $\mathcal W(\mathcal A)$, eventually along each family, the point of
the selected part satisfies the point conditions, and the datum of the member with its size data
satisfies the coefficient field bound. Then eventually along $\mathcal W(\mathcal A)$, every family
has a constant $C_0\ge0$ such that for every $C\ge C_0$, eventually along the family, the rows of the
realization of the selected part satisfy the covolume bound for $N(C,C_H)$
(Definition~\ref{is-numerical-bounds}).
\end{lemma}
\leanref{Development/AutomorphicArithmetic/ThetaCorrespondence/ThetaPairLifts/ThirdPlaceSeries/EventualL3OfCovolume.lean}{GlobalTheory.Fc.Ind.eventual_L3_of_covolume_alt}
\leanref{Development/AutomorphicArithmetic/ThetaCorrespondence/ThetaPairLifts/ThirdPlaceSeries/EvWCOnPtOfGuardedStatement.lean}{GlobalTheory.T.eventually_onPt_of_pDiv_ties_and_onPt}
\leanref{Development/AutomorphicArithmetic/ThetaCorrespondence/ThetaPairLifts/ThirdPlaceSeries/EvWCOnPtOfGuarded.lean}{GlobalTheory.Fc.Ind.evWC_onPt_of_guarded_alt}
\leanref{Development/AbelianVarietyTheory/Proofs/DivisorProofs/FamilyHypothesesOfWorkingClassSevenDuties.lean}{GlobalTheory.BigTop.famHyp_of_working_duty}
\begin{proof}
A family of $\mathcal W(\mathcal A)$ that carries its duty data satisfies the family hypotheses of
Theorem~\ref{cv-covolume-bound}: the duty data give the first five of them, and the conditions
defining the working subclass give the others. So that theorem gives statement (i) below eventually
along $\mathcal W(\mathcal A)$ (Definition~\ref{pd-family}), and the hypotheses are statements (ii)
and (iii) eventually along $\mathcal W(\mathcal A)$.
\begin{enumerate}
\item[(i)] The family has $C_0\ge0$ and $x_0$ such that, at every scale $x\ge x_0$ and every point of
  the family at $x$ that satisfies the point conditions and the coefficient field bound, the covolume
  bound holds for every $C\ge C_0$.
\item[(ii)] There is $x_0'$ beyond which the point of the selected part satisfies the point
  conditions.
\item[(iii)] There is $x_0''$ beyond which the datum of the member satisfies the coefficient field
  bound.
\end{enumerate}
Each of the three comes with its own $\varepsilon_1$ and largeness thresholds. Take the least of the
three numbers $\varepsilon_1$ and, for each exponent, a threshold that is at least each of the three;
then (i), (ii) and (iii) hold together for every family of $\mathcal W(\mathcal A)$ that is large
for it. For such a family and $x\ge\max(x_0,x_0',x_0'')$, the covolume bound holds at the point of
the selected part at every member at $x$, for every $C\ge C_0$. The rows of the realization of the
selected part are those of this point, so they satisfy the covolume bound for $N(C,C_H)$.
\end{proof}

\begin{lemma}\label{asm-dependence-bundles-member}
Let local Langlands maps be given, let $N$ be an extension, and let $A$ be a property of stages
$(K,B,f,T)$ that holds eventually in the context and at large base data
(Definition~\ref{asm-stage-construction}(1) and (2)). If $p(N)$ has the member property
(Definition~\ref{asm-dependence-parameters}(5)), then so do the parameters of the extension of $N$ by
$A$.
\end{lemma}
\leanref{Development/AutomorphicArithmetic/ThetaCorrespondence/ThetaPairLifts/ThirdPlaceSeries/IsBuiltEPrmAIffStatement.lean}{GlobalTheory.T.memberE_extensionWithAtProperty_of_evLgG}
\leanref{Development/AutomorphicArithmetic/ThetaCorrespondence/ThetaPairLifts/ThirdPlaceSeries/IsBuiltEPrmAIff.lean}{GlobalTheory.Mem.LiteralAutomorphy.memberE_extensionWithAtProperty_of_memberE}
\leanref{Development/AbelianVarietyTheory/Proofs/IwasawaDecomposition/AddLargeBaseCoreSlot.lean}{GlobalTheory.Mem.Top.GroupRetwisting.SecondG.sAnd_of_largeBase}
\begin{proof}
Write $N_A$ for the extension of $N$ by $A$. A core is built from $p(N_A)$ exactly when it is built
from $p(N)$ at a stage $(K,B,f,T)$ that satisfies $A$. The same holds for good cores, since the
further conditions on a good core do not involve the extension.

The base condition does not involve the extension. First pass: eventually in the context, at large
base data with untwisted twist data, there is a core built from $p(N)$. The property $A$ holds at
large base data as well, and two properties that hold at large base data hold together at large
base data, for the unions of their finite sets of algebraic numbers and of primes. So, eventually in
the context, at large base data with untwisted twist data, there is a core built from $p(N)$ at a
stage that satisfies $A$, which is a core built from $p(N_A)$. Export: a core built from $p(N_A)$ is
built from $p(N)$, so it exports its twisting objects.

Second pass: let a core built from $p(N_A)$, with untwisted twist data, export its twisting objects,
and let finite sets of primes to avoid and of algebraic numbers to contain be given. The core is built
from $p(N)$. Apply the second pass of $p(N)$ to it, with the given sets enlarged by the finite sets
that the large-base condition for $A$ attaches to the base data, the framing and the coefficient
fields. It gives enlarged base data with the same framing and the same coefficient fields, twist data
with the same twisting primes that avoid the enlarged set of primes and whose coefficient field
contains the enlarged set of algebraic numbers, and at them a core built from $p(N)$ and good for
$p(N)$, with all its twists. The base condition is stable under enlarging the base data. So the new
stage lies over the rank datum and satisfies the conditions of the large-base condition for $A$, and
$A$ holds at it. Hence the new core is built from $p(N_A)$ and good for $p(N_A)$.
\end{proof}

\begin{lemma}\label{asm-dependence-bundles-inputs-move}
Let local Langlands maps be given, and let a stage $(K,B,f,T,w,Z)$ be given with a set of twisted
places, a core object $O$, rows $t$ at the base level of $O$, and a norm choice and theta data of the
core. Let new division lattices $N_u$ and heights $h$ be chosen at the later stage $Z$, let $Z'$ be
the later stage with these division data, and $O'$ the transport of $O$ to $Z'$. Then:
\begin{enumerate}
\item there are rows $t'$ at the base level of $O'$ with the same row data as $t$, whose geometry and
  row package are the transports of those of $t$, such that the core object conditions, and the
  condition that the rows are rows of the construction, pass from $t$ to $t'$;
\item the norm choice and the theta data transport to $Z'$, the transport keeps the norm choice, and
  for rows $t'$ as in (1) and every norm choice of the moved stage, a norm choice of the law extension
  at $(Z,t)$ gives one at $(Z',t')$;
\item for rows $t'$ as in (1), the list pin of the level characters of the realization passes from
  $t$ to $t'$, for every constant;
\item for rows $t'$ as in (1), the spectral stage data pass from $Z$ to $Z'$, for every column datum at
  the split primes of the moved stage;
\item for rows $t'$ as in (1), the scalars of the division geometry pass from $(Z,O)$ to $(Z',O')$.
\end{enumerate}
\end{lemma}
\leanref{Development/AbelianVarietyTheory/Proofs/IwasawaDecomposition/RowsPassToMovedDivisionLatticesHeightsStatement.lean}{GlobalTheory.T.rows_pass_to_moved_divisionLatticesHeights}
\leanref{Development/AutomorphicArithmetic/ThetaCorrespondence/ThetaPairLifts/ThirdPlaceSeries/LayerUpStatement.lean}{GlobalTheory.T.exists_kept_choice_upper_lattice_alt}
\leanref{Development/AbelianVarietyTheory/Proofs/IwasawaDecomposition/GridStaircaseLatticeFirstColumnTransportStages.lean}{GlobalTheory.Mem.GridMove.listPin_rowsSame}
\leanref{Development/AbelianVarietyTheory/Proofs/IwasawaDecomposition/SpectralStageDataTransferLatticesHeightsStatement.lean}{GlobalTheory.T.spectralStageData_transfer_lattices_heights}
\leanref{Development/AbelianVarietyTheory/Proofs/IwasawaDecomposition/DivisionGeometryScalarsTransferSetDivisionStatement.lean}{GlobalTheory.T.divisionGeometryScalars_transfer_setDivision}
\begin{proof}
In each part the moved data are the given data, read at $D(w,Z')$, and each condition is checked
clause by clause. The move replaces the division lattices, the heights and the bound on the heights.
The group of the later stage and the realizations built at the stage, with their theta functions,
multiplicity spaces and Weil data, are the same at $Z$ and at $Z'$, and so are the value $\eta$ of the
division datum and the scalar $\kappa^\sharp$.

(1) The rows $t$ consist of a realization $z$ of $D(w,Z)$ built at the stage, with its core,
coefficient embedding, group isomorphism, geometry, pins and level data; of a reading of the base
level of $O$ by $z$, with its lift; and of conditions. These are the conditions that $z$ records at
its face, the existence of its theta maps, its lying in the geometry, the group conditions, the twist
conditions, the equality of the set of twisted places with the set of forced twisted places at the
scale of the context, the vanishing of the indices off it, the inner products of the maps $e_{M_t}$,
the value $\eta=6x^{-C_F}$, the integrality of the rows, and the index conditions at every division
datum tied to the lift. Let $z'$ be $z$ read at $Z'$: its core, realization and geometry are the
transports of those of $z$, and its other data are those of $z$. Let the reading and the lift of
$t'$ be those of $t$, transported along the change of division data. They have the same tuples, the
same characters at the levels, the same row representations, the same embeddings of the spaces $E$
and $M_t$, and the same coordinate maps. The conditions pass. Those on the realization, its theta maps
and the geometry are transported with $z$. The set of forced twisted places of $z'$ is that of $z$.
A division datum of $D(w,Z')$ tied to the lift of $t'$ moves back to one of $D(w,Z)$ tied to the lift
of $t$, so the index conditions pass. The other conditions are copied. This defines $t'$. It has the
same row data as $t$, and its geometry and row package are the transports of those of $t$.

The core object conditions of $O$ at $t$ are the ties of the ambient forms of $O$ at the levels $1$
and $0$ to the reading, the tie of the division data of $O$ to the lift, and the equality of the
archimedean scalars $\kappa_R(\sigma)$ of the division data of $O$ with $\kappa^\sharp$. The ambient
forms of $O'$ are those of $O$ and the readings agree, so the first ties pass. The division data of
$O'$ are the transport of those of $O$, so the second tie passes. The scalars $\kappa_R(\sigma)$ and
$\kappa^\sharp$ are unchanged. The conditions that the rows are rows of the construction are the data
at all levels with their maps, forms and lattices, their agreement with the reading, and conditions
on the core, its rank, its pullback and its standard form. They are copied, and the laws of the core
pass to the transported core.

(2) The moved norm choice and theta data are the given ones, read at $D(w,Z')$, and the move keeps
every component of the norm choice. A tie of the norm choice to the stage gives one of the moved norm
choice, with the same tuple, the same characters $\xi$ and $\psi$, the same set $Q$ and the same
character data. The good twists pass, since the twisting data of the moved theta data are those of
the given ones.

Let rows $t'$ be as in (1), let a tie of the moved norm choice to $(w,Z')$ be given, and let a norm
choice of the law extension be given at $(Z,t)$: data at the norm choice, realization data for the
rows, real expansion data and the rationality law. The tie of the moved norm choice has the same
characters $\xi$ and $\psi$ as the given tie, so its tuple is the transport of the given tuple. The
data at the norm choice move: the constructed norm choice is rebuilt from the moved datum with the
same choices, and the other components are kept. The rows $t'$ have the same realization, coefficient
embedding and reading as $t$, and in particular the same embeddings of the spaces $M_t$. So the
realization data of $(Z,t)$, namely the expansion of the realization and its Hecke, group, level and
local data, are copied to $(Z',t')$. The series condition of the real expansion data passes, with the
same series and characters and with the scalars $\lambda$ of the core moved, and the other real
expansion data are copied. The rationality law does not involve the division data, so it holds at
$D(w,Z')$.

(3) Rows with the same row data have the same realization data, field embedding and core, so the
realization of $t'$ is the transport of that of $t$. The list pin of $t$ is an evaluation grid for the
level system of its realization with constant $c$ at $b$ (Definition~\ref{asm-grid-evaluation}(1)),
whose characters take values in the coefficient field, with the roots of unity required there, and
whose scalars $\lambda$ are tied to the twisted places. Its places, its characters and its spaces
$V_{i,j}$, $P_i$ and $Q_j$, with their lattices and the orthogonal decompositions that relate them,
are read unchanged as an evaluation grid for the transported level system at $Z'$. The two field
conditions and the tie of the scalars $\lambda$ are the same conditions after the move.

(4) The spectral stage data at $(Z,t)$ for a column datum at the split primes consist of a saturation
of the row package of $t$ for the forms of that column datum: a spectral package with its
restrictions, a point $t_S$ of $U(Y_0)(\mathbf A)$, a scalar $c>0$, and the seeing, evaluation and row
conditions at them. Let a column datum at the split primes of the moved stage be given. Copy the forms
of the given column datum to forms $F_1$ at $D(w,Z')$ with the same levels, characters,
representatives and identifications. The saturation moves along the change of division data to a
saturation of the transported row package for $F_1$, with the same point and scalar: the spectral
package, its restrictions and the three conditions move with the data. Both column data are standard.
At level $0$ the representatives and the identification of a standard form are the prescribed ones,
and its level lattice, level group and level character are those that its reading at the ends
determines: the lattice and the group are read off from the reading, and the character on the group
as well. So the forms of the two column data, and hence $F_1$ and the forms of the new column datum,
have the same level lattice, level group, level character, representatives and identification. The saturation passes from $F_1$ to the
forms of the new column datum, since the transports of the spectral package, the column dictionary
and the level match use only these data. By (1), the row package of $t'$ is the transported row
package, so this is spectral stage data at $(Z',t')$.

(5) The scalars of the division geometry at $(Z,O)$ consist of the geometry condition on the core and
the division data of $O$, and of a labelling $\mathrm{lab}$ of the types by indices of the family. It
picks for each type $\sigma$ an index over $\sigma$ whenever there is one, and at every level $i+1$,
type $\sigma$ and embedding $e$ of the coefficient field other than the given one, the logarithm of
the scalar $\lambda$ of the division data at $(i+1,\sigma,e)$ equals $\log\lambda_i(\mu_0,
\mathrm{lab}(\sigma))$ at the infinite place of $e$, for the rows $S_0$ of the realization, when the
isotypic part of the datum at $(i+1,\sigma)$ is nonzero, and $0$ otherwise. The geometry condition is the closure of a condition under the moves of the division data,
so it passes to the transported core and division data. The labelling condition involves only the
lift of the rows, the realization and the logarithms of the scalars $\lambda$ of the division data;
the rows $t'$ have the same lift and realization, and the move keeps these logarithms. So it passes.
\end{proof}
% Lean: (1) rows_pass_to_moved_divisionLatticesHeights (RowsSame, geoHalf, the pack (SRe.pkOf t).up,
% ObjT and RXs); (2) exists_kept_choice_upper_lattice_alt with S' = Z.S (nchTr, xdTr, Kept, NchT,
% GoodO and normChoiceOfExtension of extensionWithLaw); (3) listPin_rowsSame, which identifies Xg, tr,
% core, alg and R and applies listPin_upL; (4) spectralStageData_transfer_lattices_heights
% (SRe.SxMv); (5) divisionGeometryScalars_transfer_setDivision (OAt (geoCl Gamma)). In the proof:
% (1) is DivMv.rows_move_geo = readSame (readB along Refill.ofDiv, fields tup ... eRR0), with
% ReadSame.what2 (Mv.cpDat with sameG; cp_holds, cp_exAt, cp_inGeo, forcedContentPlaces_eq_transport,
% idx through mj_back), rowsSame, pk, objT (amb1/amb0 by AmbW.congr, mjH by mjW_fwd, kapR with
% EX.OB.kapSharp_tr) and rxs (allMv, lawsCore_tr); (2) is LayerMove.exists_kept_nChoice_thData_upL
% (nchUp, xdUp, kept_up, nchT_up, goodO_up with twEOf; tupTied_ext, Mv.exists_at_alt,
% exists_nonempty_nMore_upL with JmtUp.jMt_rowsSame and geo_heq_geoHalf_of_same for SerAt, Mv.bkLaw,
% which is the identity); (4) is SxPass.sxMv (exists_sameF_ofDiv, sAtL_exists_ofDiv,
% StdUnique.sameStd = std_pair and reads_unique, satCar); (5) is OxMove.oAt_setDiv (GeoCl.step and
% lamAre_copy). All are structure copies apart from reads_unique (about 15 lines). Sizes (text
% dependencies): about 33 theorems and 270 lines, 34 and 820, 9 and 180, 36 and 530, 15 and 300.

\begin{lemma}\label{asm-dependence-bundles-inputs}
Let local Langlands maps be given. The results named below are those of
Theorem~\ref{lit-bundle-unitary}.
\begin{enumerate}
\item Assume Harish-Chandra analyticity, the classification of the discrete series, rapid decay of
  cusp forms, Harish-Chandra's finiteness theorem for automorphic forms on unitary groups, the trace
  class property of the cuspidal spectrum, the factorisations of Dixmier-Malliavin and of Flath,
  Zelevinsky's polynomial ring, the geometric lemma and Green's description of the cuspidal
  characters. Let $N$ be an extension. Then \emph{cuspidal multiplicities are equal} at $p(N)$:
  eventually in the context, for every stage that satisfies $N$ and the standard conditions of the
  construction, with realization $R$ of its theta datum, the following holds. For every level $i$,
  every discrete series model at infinity for the base space $Y_0$, every isomorphism of the finite
  adelic groups of $U(Y_0)$ and $U(W_i)$ that is compatible with the lattices and the group of $R$,
  and every irreducible admissible representation $\sigma$ of the finite adelic group of $U(Y_0)$ that
  is good for $R$, the cuspidal multiplicity of $\sigma$ with that model equals the cuspidal
  multiplicity of the transport of $\sigma$ to $U(W_i)$ with the model of $R$ at level $i$.
\item For every real $c$, the positions of the parameters $p_{\mathrm{orig}}=p(N_{\mathrm{orig}}(c))$
  \emph{move with the division data}: for every stage built from $p_{\mathrm{orig}}$, with its
  witness, every choice of new division lattices $N_u$ and heights at its later stage, every column
  datum at the split primes and every norm choice of the moved stage, there are rows at the moved
  stage that carry the transported core, with which the moved stage satisfies the conditions of
  $p_{\mathrm{orig}}$ on its rows, the extension $N_{\mathrm{orig}}(c)$, the conditions on its twist
  choices and those on its core and geometry.
\item Assume the published results on unitary groups and that the rows of the construction are
  rational (Definition~\ref{asm-stage-construction}(4)). Let $c$ be real. If cuspidal multiplicities
  are equal at $p_{\mathrm{orig}}=p(N_{\mathrm{orig}}(c))$, then $p_{\mathrm{orig}}$ has the
  member property and bundles exist at large bases for $p_{\mathrm{orig}}$
  (Definition~\ref{asm-dependence-parameters}(5) and (6)).
\item Assume the published results on unitary groups and that the rows of the construction are
  rational. Let $N$ be an extension stronger than the law extension such that the positions of
  $p(N)$ move with the division data, as in (2), and let $A$ be a property of stages $(K,B,f,T)$.
  Then bundles exist at large bases for the parameters of the extension of $N$ by $A$.
\end{enumerate}
\end{lemma}
\leanref{Development/AutomorphicArithmetic/ThetaCorrespondence/ThetaPairLifts/ThirdPlaceSeries/MultiplicativeClosurePrimeMem.lean}{GlobalTheory.T.member_multiplicatively_closed_alt}
\leanref{Development/AutomorphicArithmetic/ThetaCorrespondence/ThetaPairLifts/ThirdPlaceSeries/MultiplicativeClosurePrimeMem.lean}{GlobalTheory.Mem.Top.Bl.MultQClose.multiplicative_closure_prime_mem_alt}
\leanref{Development/AutomorphicArithmetic/ThetaCorrespondence/ThetaPairLifts/ThirdPlaceSeries/MultQ.lean}{GlobalTheory.Mem.Top.Bl.MultiplicityCondition}
\leanref{Development/AbelianVarietyTheory/Definitions/AutomorphicMembership/PositionMv.lean}{GlobalTheory.Mem.Top.Bl.PositionMv}
\leanref{Development/AutomorphicArithmetic/ThetaCorrespondence/ThetaPairLifts/ThirdPlaceSeries/SlotMvPrmRNelStatement.lean}{GlobalTheory.T.slot_mv_extension_with_origin_geo_cl}
\leanref{Development/AutomorphicArithmetic/ThetaCorrespondence/ThetaPairLifts/ThirdPlaceSeries/SlotMvPrmRNel.lean}{GlobalTheory.HeightCounting.Up.movable_prmR_nel_alt}
\leanref{Development/AutomorphicArithmetic/ThetaCorrespondence/ThetaPairLifts/ThirdPlaceSeries/ExistsBaseLatticeOfPrime.lean}{GlobalTheory.T.member_exists_baseLattice_alt}
\leanref{Development/AutomorphicArithmetic/ThetaCorrespondence/ThetaPairLifts/ThirdPlaceSeries/ExistsBaseLatticeOfPrime.lean}{GlobalTheory.Mem.Top.Bl.exists_baseLattice_of_prime_alt}
\leanref{Development/AutomorphicArithmetic/ThetaCorrespondence/ThetaPairLifts/ThirdPlaceSeries/ExistsWeilBlockNeAtMember.lean}{GlobalTheory.T.compatibleSystem_weilBlock_ne_alt}
\leanref{Development/AutomorphicArithmetic/ThetaCorrespondence/ThetaPairLifts/ThirdPlaceSeries/ExistsWeilBlockNeAtMember.lean}{GlobalTheory.Mem.LiteralAutomorphy.exists_weilBlock_ne_at_member_alt}
\begin{proof}
(1) Eventually in the context, two things hold. First, for every later stage over base data of the
context whose division datum has $\eta=6x^{-C_F}$, with $C_F$ the constant of the context, and all of
whose twisting exponents $A_j(w)$ vanish, the datum carries size data at the scale of the context
with a parameter choice; this follows from
Lemma~\ref{asm-dependence-functional-witness-numbers-values} and the explicit size data. Second, the
datum of every stage over the context satisfies the host frames hypothesis. Fix such a context and a
stage that satisfies $N$ and the standard conditions, with later stage $Z$ and realization $R$. By
Green's description of the cuspidal characters and Lemma~\ref{asm-dependence-test-assembly-bare},
there is a later stage $Z'$ of the same column choice, whose division datum has this value of $\eta$,
whose twisting exponents $A_j(w)=I_j(\mu_w,w)+I_j(\mu_w,\bar w)+j_w(\mu_w)$ vanish, and whose group
data are true for the framing (Definition~\ref{asm-dependence-test-assembly-true}(1)). It has the
bare rows of the rows of $Z$, the division datum of $Z$ with $\eta$ replaced, the norm data over the
bare rows that Lemma~\ref{asm-dependence-test-assembly-bare-norm} builds from Green's description,
and the size and jet data of $Z$. The group data of $Z$ are also true for the framing, which is one
of the group conditions on the rows of the stage, so $Z$ and $Z'$ have the same group
(Definition~\ref{asm-dependence-test-assembly-true}(2) and the remark after it): their places
correspond, through the map of places of their automorphism data, and so do the coefficients of their
type idempotents. The realization $R$ has a copy $R'$ at $Z'$
(Lemma~\ref{asm-dependence-test-assembly-copy}), with the same realization data, the transported
core and the group isomorphism composed with this identification; it exists at every place, since
$R$ does. If cuspidal multiplicities are equal for $R'$, they are equal for $R$, by the same lemma.
Indeed, transporting $R'$ back to $Z$ gives $R$. The statement for a realization involves its archimedean
family, its models, its lattices and frames, its places $P_g$, its group isomorphism and its level
character at level $0$, which the transport keeps, and the cuspidal multiplicities of its models,
which do not involve the group stage. The realization $R'$ satisfies the
hypotheses of Lemma~\ref{asm-dependence-orbit-inputs-multiplicity-one} for the datum of $Z'$: its
theta maps exist and its family is ok at level $1$, since this holds for $R$; it records the weak
characterization and the guard; the host has $\mathfrak c=\mathcal O_K$; the host frames hypothesis
and the parameter choice are the two facts above; and its types are good, since the group of $Z'$ is
true. So cuspidal multiplicities are equal for $R'$, and hence for $R$.

(2) Let a stage built from $p_{\mathrm{orig}}$ be given, with its witness, new division lattices and
heights, a column datum at the split primes and a norm choice of the moved stage. Let $t'$ be the rows
at the moved stage that Lemma~\ref{asm-dependence-bundles-inputs-move}(1) gives for the rows of the
witness. By that lemma they carry the transported core object and are rows of the construction,
which are the conditions of $p_{\mathrm{orig}}$ on rows; the transported norm choice is kept, so the norm choice of
the law extension (part (2)), the list pin (part (3)) and the origin data, by the transitivity of
keeping, pass to the moved stage, which is the extension $N_{\mathrm{orig}}(c)$; the spectral stage
data pass (part (4)), which are the conditions on the twist choices; and the scalars of the division
geometry pass (part (5)), which are the conditions on the core and its geometry.

(3) From the results of $P$ and the rationality of the rows follow the hypotheses on the member and
on the rational weights of the multiplicity cover, the saturated state of weight one, and, from the
first and second collections of item (6) of Theorem~\ref{lit-bundle-unitary}, the elliptic and local
Langlands hypotheses. With these and the equality of cuspidal multiplicities at $p_{\mathrm{orig}}$,
$p_{\mathrm{orig}}$ has the member property and bundles exist at large bases for it.

(4) Write $N_A$ for the extension of $N$ by $A$. A stage that satisfies $N_A$ and moves to a stage with
new division lattices and heights keeps its context, base data, framing and twist data, so $A$ still
holds after the move; hence the positions of $p(N_A)$ move with the division data, as those of
$p(N)$ do. By (1), cuspidal multiplicities are equal at $p(N_A)$. With the hypotheses that $P$ and the
rationality of the rows give, as in (3), and the law extension, bundles exist at large bases for
$p(N_A)$.
\openstep{In (3), the derivation of the member property and of the bundles from the equality of
cuspidal multiplicities, and in (4), the existence of the bundles for an extension whose positions
move and whose cuspidal multiplicities are equal, are not yet written out.}
\end{proof}
% Lean: (1) multiplicative_closure_prime_mem_alt: EvP of LedRec.ev_exists_sizeData (szRow_numAt and
% EX.exists_sizeData_explicit) and HostFrames.hostHyp_frames_ev; LedRec.exists_bare (Green; bare
% a.W with etaF = 6 x^(-CF) and Aj = 0, through exists_nStage_grpTrue) and LedRec.sameGrp;
% equalMultiplicities_rowTest_of_sameGroupTransfer = Twist.multEq_rowTest_of_copy (DatCopy.twice_eq and
% DatCopy.multEq_tr, with mult_eq); exists_bare = <bare a.W, a.V with etaF, a.S, N, a.G> with N from
% Top.exists_nStage_grpTrue (trueLayer, GaloisStage.nAut, TypesInhab.types0 with isSer_eSerC from
% GreenCuspidal, exists_char, retype; written out as Lemmas asm-dependence-test-assembly-bare-norm and
% -bare); LedRec.sameGrp (plEquiv, plGw_plEquiv, typeIdempotent_coeff_across_layers; the remark after
% Definition asm-dependence-test-assembly-true); DatCopy.copy_exAt (Lemma -test-assembly-copy);
% twisted_multiplicity_one with s = 1. (2) movable_prmR_nel_alt, with the fields Rx = hrx, Nx =
% (c.1, c.2, ListPin, origin data with kept_trans hk hkept), Sx = s and Ox = o. (3)
% member_exists_baseLattice_alt = exists_baseLattice_of_prime_alt (w_of_weight_residue_char,
% fineAdelicDatum_extends_to_larger_field_of_hypotheses, all_stages_at_member_of_specializations_alt,
% sFino_and_secondNo_of_sFinl', has3_neo_of_localLanglands_alt, lv3Mv_extensionWithOrigin_and_NEog,
% exists_baseLattice_of_pair_alt), with the hypotheses as in
% StageOneU.member_bundles_large_unramified_exponents_alt: compatibleSystem_member_hypotheses (HypAAlt),
% compatibleSystem_member_rationalWeights (MultiplicityCoverHypotheses, with S0.RowRat),
% compatibleSystem_weightOne_saturated, and hyps_e_third_layer_ser_n_of_geometric_galois (P.litN,
% P.toHypothesesX). (4) compatibleSystem_weilBlock_ne_alt = exists_weilBlock_ne_at_member_alt:
% slot_mv_extension_with_at_property (the A of the witness, x.nx.2.2), member_multiplicatively_closed_alt
% and geometric_parameter_uniform_alt. Sizes of the open step (text dependencies, overcounts):
% exists_baseLattice_of_pair_alt about 56,000 theorems and 978,000 lines;
% geometric_parameter_uniform_alt about 5,200 theorems and 85,000 lines.

\subsection{Evaluation grids and the functional bound}\label{subsection-asm-functional-bound}

\begin{definition}\label{asm-dependence-functional-twist-bounds}
Let $D$ be a definite-side datum with host field $E/F$, and take its twisting data
(Definition~\ref{th-twisting-data}): twisting indices $i$ with primes $\ell_i$, sets $S_0(i)$ of places
of $E$, rings $\mathcal O_{i,k}$, ray class groups $\mathrm{RCG}_i(Q)$, and the sets $\Lambda_i$ whose
elements $\lambda$ carry signs $b_\lambda$, with the theta values $\vartheta_i(Q,\lambda,\psi)$ twisted
by characters $\psi$ of $\mathrm{RCG}_i(Q)$ and the bounds $c_0(\lambda)$. Write $v$ for the valuation
normalized by $v(\ell_i)=1$, and for a character $\psi$ and a sign $b$ put $\psi^{(b)}=\psi$ if $b=+$
and $\psi^{(b)}=\psi^{-1}$ if $b=-$. Let $i$ be a twisting index and $Q$ a finite set of places of $E$.
\begin{enumerate}
\item $Q$ is a \emph{stable split set} for $i$ if every place of $Q$ is split in $E/F$, $Q$ is
  disjoint from $S_0(i)$, and the conjugate of every place of $Q$ lies in $Q$.
\item Let $N_0$ be a function from the twisting indices to the integers $\ge0$, let $k'\ge0$, and let
  $\psi$ be a character of $\mathrm{RCG}_i(Q)$ with values in $\mathcal O_{i,k'}^\times$. The
  \emph{theta value bounds} hold for $N_0$, $i$, $Q$, $k'$ and $\psi$ if for every
  $\lambda\in\Lambda_i$ the theta value $\vartheta=\vartheta_i(Q,\lambda,\psi^{(b_\lambda)})$ of
  $\lambda$ twisted by $\psi^{(b_\lambda)}$ is nonzero and
  $-N_0(i)\le v(\vartheta)\le c_0(\lambda)+\ell_i^{-2}$.
\end{enumerate}
\end{definition}
\leanref{Formalism/SharpRestriction/LevelLetterCocycles/BlochKatoPhases.lean}{Formalism.MainNotions.BlochKatoReduction.Gen.QOk}
\leanref{Formalism/SharpRestriction/LevelLetterCocycles/BlochKatoPhases.lean}{Formalism.MainNotions.BlochKatoReduction.Gen.partC}
\leanref{Formalism/SharpRestriction/LevelLetterCocycles/BlochKatoPhases.lean}{Formalism.MainNotions.BlochKatoReduction.Gen.psiAt}
\leanref{Formalism/SharpRestriction/LevelLetterCocycles/TwGood.lean}{Formalism.MainNotions.twEps}

For the character $\psi^{(b)}$ of a twist choice $(i,b,Q,k',\psi)$, the theta value bounds are
condition (2)(e) of the twist condition of a family (Definition~\ref{asm-twist-condition}), and a set
$Q$ as in condition (2)(b) there is a stable split set for $i$ if it contains the conjugates of its
places.

\begin{definition}\label{asm-dependence-functional-fitted}
Let $p$ be a point of a family (Definition~\ref{is-point}), with definite-side datum $D$, coefficient
field $L\subset\mathbf C$, the subfield $L_0$ of $L$ that $D$ carries, realization $R$, and twist
choices of $D$: for every twisted place $w$ a twisting index $\iota_w$, a finite set $Q_w$ of places
of the host field, an integer $k'_w$ and a character $\psi'_w$. Let $C_F$ be real and let $G$ be an
evaluation grid for $R$ with constant $C_F$ at $b$ (Definition~\ref{asm-grid-evaluation}(1)), with
base place $v_0$, place $w_0=w(v_0)$ above it and characters $\chi_1,\dots,\chi_m$.
\begin{enumerate}
\item $G$ is \emph{valued in $L_0$} if every $\chi_i$ takes values in $L_0$ and $L_0$ contains a
  primitive root of unity of order $N(v)^{d(v)}$ for every place $v$ of the set $S_0$ of $G$, where
  $d(v)$ is the depth of $G$ at $v$. The \emph{scalars} of $G$ are the columns of $R$, the auxiliary
  prime $q_0$, the place $w_0$, the centres of $G$ at the lifts and the characters $\chi_i$; they
  are \emph{tied to the twisted places} as in the list extension
  (Definition~\ref{asm-dependence-parameters}(3)).
\item The twist choices \emph{lie over their pairs} if for every twisted place $w$ the twisting
  index $\iota_w$ has the residue characteristic of $w$ and lies under $w$ or under $\bar w$, and
  twisted places $w,w'$ of the same pair with $Q_w\ne\emptyset$ have $\iota_w=\iota_{w'}$. They
  \emph{satisfy the twist conditions} if at every twisted place $w$ the set $Q_w$ is a stable split
  set for $\iota_w$ (Definition~\ref{asm-dependence-functional-twist-bounds}(1)), and there is a
  function $N_0$ on the twisting indices such that at every twisted place $w$ and for every integer
  $a$ prime to the residue characteristic of $\iota_w$, the theta value bounds
  (Definition~\ref{asm-dependence-functional-twist-bounds}(2)) hold for $N_0$, $\iota_w$, $Q_w$,
  $k'_w$ and $(\psi'_w)^a$.
\item $G$ is \emph{fitted to} $p$ if the characters of the columns of $R$ take values in $L_0$, $G$
  is valued in $L_0$, the twist choices lie over their pairs and satisfy the twist conditions, the
  character choice of $p$ is compatible with the twist choices, and the scalars of $G$ are tied to
  the twisted places.
\item Let $W$ be a Weil datum over $R$ for $G$. \emph{Remaining choices of the slope test} for $G$
  and $W$ consist of a row volume, a real vector, a linear functional on the coefficients and the
  further choices of an evaluation. They are \emph{admissible} if together with $G$ and $W$ they
  form a full evaluation choice (Definition~\ref{asm-stage-construction}(5)) and the series
  idempotents of $D$ are tied to its labels.
\item The point $p$ \emph{has a fitted evaluation grid with constant $C_F$} if there are an
  evaluation grid $G$ for $R$ with constant $C_F$ at $b$ fitted to $p$, a Weil datum $W$ over $R$
  for $G$ and admissible remaining choices of the slope test for $G$ and $W$.
\end{enumerate}
\end{definition}
\leanref{Development/AbelianVarietyTheory/Definitions/Divisors/TwistedPlaceChoiceLiesOverPair.lean}{GlobalTheory.BigTop.Fun.FieldHas}
\leanref{Development/AbelianVarietyTheory/Definitions/Divisors/TwistedPlaceChoiceLiesOverPair.lean}{GlobalTheory.BigTop.Fun.lamOf}
\leanref{Development/AbelianVarietyTheory/Definitions/NonGaloisDescent/NotTeichTrivOfIsDelta.lean}{GlobalTheory.LamTie.IsLamTieFull}
\leanref{Development/AbelianVarietyTheory/Definitions/Divisors/TwistedPlaceChoiceLiesOverPair.lean}{GlobalTheory.BigTop.Fun.TwistedPlaceChoiceLiesOverPair}
\leanref{Formalism/LocalToGlobal/UnitaryNorms/CharTied.lean}{Formalism.L2.Host.ValLL0}
\leanref{Formalism/LocalToGlobal/UnitaryNorms/CharGal.lean}{Formalism.L2.Host.PsiChoiceCompatible}
\leanref{Development/AbelianVarietyTheory/Definitions/Divisors/TwistedPlaceChoiceLiesOverPair.lean}{GlobalTheory.BigTop.Fun.GridFit}
\leanref{Development/AbelianVarietyTheory/Definitions/Divisors/SlopeTestRemainingChoicesAlt.lean}{GlobalTheory.BigTop.Fun.SlopeTestRemainingChoices}
\leanref{Development/AbelianVarietyTheory/Definitions/Divisors/SlopeTestRemainingChoicesAlt.lean}{GlobalTheory.BigTop.Fun.SlopeTestRemainingChoices.OK}
\leanref{Development/AbelianVarietyTheory/Definitions/Divisors/SlopeTestRemainingChoicesAlt.lean}{GlobalTheory.BigTop.Fun.Fld}
\leanref{Development/AbelianVarietyTheory/Definitions/AutomorphicMembership/GKeyFin.lean}{GlobalTheory.Mem.LiteralAutomorphy.Full}

\begin{lemma}\label{asm-dependence-functional-twists}
Let local Langlands maps be given and let $p$ be parameters
(Definition~\ref{asm-dependence-parameters}). Let a bundle be linked by $\Lambda_p$ to a base point
of a core, and let $D$ be its definite-side datum. Then the twist choices of $D$ lie over their pairs
and satisfy the twist conditions (Definition~\ref{asm-dependence-functional-fitted}(2)).
\end{lemma}
\leanref{Development/AutomorphicArithmetic/ThetaCorrespondence/ThetaPairLifts/ThirdPlaceSeries/TwistsAtLinked.lean}{GlobalTheory.Mem.LiteralAutomorphy.TwAtLink.twists_at_linked_alt}
\leanref{Development/AutomorphicArithmetic/ThetaCorrespondence/ThetaPairLifts/ThirdPlaceSeries/TwistsAtLinked.lean}{GlobalTheory.Mem.LiteralAutomorphy.TwAtLink.Q_eq_alt}
\leanref{Development/AbelianVarietyTheory/Proofs/IwasawaDecomposition/TwistsAtLinked.lean}{GlobalTheory.Mem.LiteralAutomorphy.TwAtLink.twI_eq}
\leanref{Development/AbelianVarietyTheory/Proofs/IwasawaDecomposition/TwistsAtLinked.lean}{GlobalTheory.Mem.LiteralAutomorphy.TwAtLink.cond_powers}
\begin{proof}
The link data give a stage built from $p$, with a witness, whose datum and twist choices are those
of the bundle; its core is good. The conditions of a good core on its twists say that at every
twisted place $w$ the twist choices are good for a twist input $I$ at every exponent prime to the
residue characteristic of $\iota_w$, and that twisted places with distinct twisting indices have
disjoint sets $Q$.

Let $w$ be twisted. The twisting index $\iota_w$ has the residue characteristic of $w$ by the
construction of the twist choices. It lies under $w$ or under $\bar w$ by the tie of the twist
choices to the places, since the subfield $L_0$ of the coefficient field is Galois over $\mathbf Q$.

Let $w$ and $w'$ be twisted places of the same pair. Every twisted place lies in the set $W_t$ of
twisted places of the stage, since the tuple at a place outside $W_t$ is trivial. If $w$ and $w'$
lie over the same place of $F$, the twist data of the stage on $W_t$ and the coherence of the column
choice give $Q_w=Q_{w'}$ and $k'_w=k'_{w'}$. Otherwise $\bar w$ and $w'$ lie over the same place,
$\bar w$ is twisted, and the twist choices of the core are symmetric under conjugation, so
$Q_w=Q_{\bar w}=Q_{w'}$ and $k'_w=k'_{w'}$. If moreover $Q_w\ne\emptyset$ and
$\iota_w\ne\iota_{w'}$, then $Q_w$ and $Q_{w'}=Q_w$ would be disjoint and nonempty, which is
impossible. So $\iota_w=\iota_{w'}$.

The goodness at $w$ for the exponent $1$ shows that $Q_w$ is a stable split set for $\iota_w$: its
places are split, it is disjoint from $S_0(\iota_w)$, and it is stable under conjugation because the
character is ramified at the places of $Q_w$ and nowhere else and is inverted by the conjugation.
Let $N_0$ be the function
of the twist input $I$. For an integer $a$ prime to the residue characteristic of $\iota_w$, let
$b=\pm a$, with the sign given by the side of $w$, so that the character attached to $(\psi'_w)^b$ at
that side is $(\psi'_w)^a$. Then $|b|=|a|$ is prime to that characteristic, and the goodness at $w$
for the exponent $b$ gives the theta value bounds for $(\psi'_w)^a$.
\end{proof}

\begin{lemma}\label{asm-dependence-functional-pinned}
Let $C_F$ be real, let local Langlands maps be given, let $N$ be an extension stronger than
$N_{\mathrm{list}}(C_F)$, and let $p$ be the parameters of the extension of $N$ by the grid property
for $C_F$ (Definition~\ref{asm-dependence-parameters}). Let a bundle be linked by $\Lambda_p$ to a
base point of a core in a context with scale $(x,\varepsilon)$, and put $b=e^{x^{\varepsilon}}$. For
every comparison part of the bundle whose point satisfies the compatibility conditions of
Definition~\ref{is-point}, the realization $R$ of the point has an evaluation grid with constant
$C_F$ at $b$ that is valued in $L_0$ and whose scalars are tied to the twisted places
(Definition~\ref{asm-dependence-functional-fitted}(1)).
\end{lemma}
\leanref{Development/AutomorphicArithmetic/ThetaCorrespondence/ThetaPairLifts/ThirdPlaceSeries/EvPSizeRowsDlog.lean}{GlobalTheory.Mem.LiteralAutomorphy.GridRow.rowE_pinned_grid_alt}
\leanref{Development/AbelianVarietyTheory/Proofs/IwasawaDecomposition/RowEListPinPt.lean}{GlobalTheory.Mem.LiteralAutomorphy.PinPt.listPin_twArgs}
\leanref{Development/AutomorphicArithmetic/ThetaCorrespondence/ThetaPairLifts/ThirdPlaceSeries/ListPinRpk.lean}{GlobalTheory.Mem.LiteralAutomorphy.PinPt.listPin_rpk_alt}
\begin{proof}
The link data give a stage built from $p$, with a witness. The stage satisfies the extension of $N$
by the grid property, so it satisfies $N$ and hence $N_{\mathrm{list}}(C_F)$. So the level characters
of the realization of the theta datum of the stage have an evaluation grid with constant $C_F$ at $b$
that is valued in $L_0$ and whose scalars are tied to the twisted places, where $b$ is that of the
scale of the context. The witness identifies these level characters with those of the realization
of its twist input, and the realization of its real input with the indefinite realization of the
bundle. The realization of the point of a comparison part is the indefinite realization of the
bundle, and its size data are at the scale of the context. So the grid is one for $R$ with constant
$C_F$ at $b$. The compatibility conditions are not used.
\end{proof}

\begin{definition}\label{asm-dependence-functional-witness-numbers-explicit}
Let $C$ be standard constants of a context (Definition~\ref{asm-stage-construction}), among them
$C_8$, the constant $C_F$ of the parameter choices, $C_\Gamma$, $C_c$, $C_t$, a constant $C_T$ for
the host and a constant $K_\Lambda$; and let $c_D,c_h$ be real numbers. Let $D$ be a definite-side
datum with ranks $k,m,t,n$, $f=2n$ and $N(j)=k+j(t+1)$, with the number $\Delta_S$ of
Definition~\ref{th-definite-datum}(8), with a host $H$ in the notation of Definition~\ref{hf-host},
$N=Q_Sqd$, and with Harish-Chandra parameters whose parts at an embedding $\tau$ are $u_i(\tau)$,
$v_i(\tau)$, $b_j(\tau)$ and $e_j(\tau)$. Let size data for $D$ be given (Definition~\ref{pd-size-data})
at a scale $s=(x,\varepsilon,J)$, with $b$, $b'$, $\delta_1$, $w''(n)=b/(X'n)$, the number $N_1$ of
Definition~\ref{hf-good-orbit}, exponents $b_t$, width $\eta$, $\Lambda_G$, a depth $j_0$, constants
$C_t$, $C_c$, $C_\Gamma$ and the logarithmic derivative $L'/L$ at $1$. Put $\Delta=\lceil x^{c_D}\rceil$ and
$h=\lceil x^{c_h}\rceil$. The size data take the \emph{explicit values} for $(C,c_D,c_h)$ if:
\begin{enumerate}
\item the rank datum of $D$ is adapted to $s$ (Definition~\ref{hf-adapted-ranks});
\item the host has a pair scale (Definition~\ref{hf-good-host}), $\delta_1n/2\le|T_{\mathrm{top}}|\le
  2\delta_1n$, $|T_2|\ge x^{-\varepsilon}n$, $p_{de},p'_{de}\le N^{C_8}$, $(kb)^8\le g_i\le2(kb)^8$ for
  $i=0,1$, $L_w=10\lceil405\log x\rceil+2\lceil30/\varepsilon\rceil+22$, $\mathfrak c=\mathcal O_K$,
  $B\ge\lfloor x^{\varepsilon/(5J)}\rfloor$, $J\le C_8$, $o\le C_8$ and $Q_S\le C_8n$;
\item at every embedding $\tau$: $v_i(\tau)=(k-1-i)\cdot4\Delta$ for $1\le i\le k-1$; if $\tau$ is not
  the place of a top class, $u_1(\tau)=(k-2)\cdot4\Delta+4k\Delta+\frac12$ and $u_k(\tau)=-\frac12$;
  $b_j(\tau)=b_1(\tau)+2(j-1)h$ and $e_j(\tau)=b_j(\tau)+h$, where $b_1(\tau)=b_t+\lfloor b'/x\rfloor$ at
  the place of a top class $t$, $\lceil b\rceil\le b_1(\tau)<\lceil b\rceil+w_E$ at a thick place, $w_E$
  the number of roots of unity in $E$, and $b_1(\tau)=\lceil w''(n)\rceil$ at a thin place; and for
  every $j$, $\prod_\tau\tau(z)^{e_j(\tau)}=1$ for every root of unity $z\in E$;
\item $|L'/L(1)|\le x^{\varepsilon}$;
\item $\Delta_S=\Delta$, $\eta=2mh$, $\Lambda_G=K_\Lambda N_1^2\log(xb)$,
  $b_t\le b'+2n\delta_1x^{\varepsilon}\Lambda_Gb'/k$ for every top class $t$, and
  $j_0=\lceil\log_5(nN(m)/m)\rceil$;
\item the constants $C_t$, $C_c$ and $C_\Gamma$ of the size data and the constant $C_T$ of the host are
  those of $C$;
\item $D$ has $\lceil3\Gamma^2x^{C_F}\rceil$ words, where $\Gamma=\exp(C_\Gamma m^5fx^{\varepsilon})$;
\item the division data of $D$ have $\eta_F=6x^{-C_F}$.
\end{enumerate}
A property of definite-side data with size data holds \emph{for explicit values} if for every
$0<\varepsilon<1/52$ and all standard constants $C$ there is $c_1$ such that for every $c_D\ge c_1$
there is $c_2$ such that for every $c_h\ge c_2$ there is $x_0$ such that the property holds for every
definite-side datum with size data at a scale $(x,\varepsilon)$, $x\ge x_0$, that take the explicit
values for $(C,c_D,c_h)$. Two properties that hold for explicit values hold together for explicit
values: take the larger $c_1$, then the larger $c_2$, then the larger $x_0$. The host of $D$
satisfies the \emph{arithmetic conditions} if $g_i=\pi\bar\pi$ with $\pi\in\mathcal O_E\cap K$ and
$g_i\equiv1\pmod N$ for $i=0,1$, the primes $p_{de}$ and $p'_{de}$ split completely in $E$, and every
level prime splits in $K$. Standard constants have some fixed values, among them $C_{nW}=10$, and
satisfy lower bounds, among them $C_8\ge1$ and $K_\Lambda\ge1$.
\end{definition}
\leanref{Development/AbelianVarietyTheory/Definitions/GlobalTwisting/PairRowIffQdRows.lean}{GlobalTheory.EX.Val}
\leanref{Development/AbelianVarietyTheory/Definitions/GlobalTwisting/PairRowIffQdRows.lean}{GlobalTheory.EX.HostVal}
\leanref{Development/AbelianVarietyTheory/Definitions/GlobalTwisting/PairRowIffQdRows.lean}{GlobalTheory.EX.HCVal}
\leanref{Development/AbelianVarietyTheory/Definitions/GlobalTwisting/PairRowIffQdRows.lean}{GlobalTheory.EX.HostArith}
\leanref{Development/AbelianVarietyTheory/Definitions/GlobalTwisting/PairRowIffQdRows.lean}{GlobalTheory.EX.SzRow}
\leanref{Development/AbelianVarietyTheory/Proofs/GlobalInterpolation/SzRowAndImp.lean}{GlobalTheory.EX.szRow_and_imp}
\leanref{Development/AbelianVarietyTheory/Definitions/GlobalTwisting/PairRowIffQdRows.lean}{GlobalTheory.EX.Std}

\begin{lemma}\label{asm-dependence-functional-witness-numbers-parameters-computations}
In the notation of Definition~\ref{asm-dependence-functional-witness-numbers-explicit}, with the scale
$(x,\varepsilon,J)$, put $L=x^{\varepsilon}$, $\delta_2=|T_2|/n$, $N^*=2n\delta_1$, $R_0=nx^{-\varepsilon}$,
$h_0=\eta/m$, $w''=w''(n)$, $X''=X''(n)$, $Q=Q_Sqd$, $\ell_A=kx^{\varepsilon}/(m\delta_2)$, $\ell^*$ as in
Lemma~\ref{asm-dependence-functional-witness-numbers-bounds-computations},
$\nu(c)=cr(f+f^2(k-1)^3L^2/m)$ and $\gamma=C_{\mathrm{gap}}km^4+N(m)m^{10}$, where $C_t$ is the constant of
the size data and $C_{k\mathrm{Low}}$, $C_{k6}$, $C_{fL}$, $c_{LL}$, $C_{nW}$, $C_\rho$, $C_{\Delta G}$,
$C_{Lw}$, $C_{j0}$, $C_w$, $c_w$, $C_\Delta$, $C_{\mathrm{count}}$, $C_{h0r}$, $C_{h0k}$, $C_{h0N}$,
$C_{\mathrm{gap}}$, $C_{\eta F}$, $C_F$, $C_Y$, $C_{YL}$, $C_8$ and $C_{\mathfrak c}$ are constants of $C$.
Indices $i,j$ of the parts $u_i,v_i,b_j,e_j$ run over their ranges, $\tau$ over the embeddings, and
``at a top place'' means at the place of a top class $t$. Each of the following twenty-one statements
holds for explicit values.
\begin{enumerate}
\item $x^{2C_t+5}<q,d$; $x^{\varepsilon/(6J)}\le\log q$; $|L'/L(1)|\le x^{\varepsilon}$; $Q^{36}t^{150}\le
  b^{1/2}$; and $\log p_{de},\log p'_{de}\le L$.
\item $C_{k\mathrm{Low}}m^5fLx^{\varepsilon}\le k\le x^{1/2}nR_0$ and $\sqrt{k/m}\le R_0$.
\item $m^2L\le k$, $mx^{1/2}\log x\le k$, $C_{k6}m^6L\le k\log k$ and $C_{fL}L\le f\log k$.
\item $\log f\le\sqrt{\lfloor m/2\rfloor-1}$, $\log^2f\le m^2L$, $\log\log b\le c_{LL}rf\log r$ and
  $(4f+10)r\le\nu(C_{nW})$.
\item $\eta\le b/(t^3fk^3)$, $\log\eta\le x^{\varepsilon/4}$, $t\le\eta$, $t\eta\le w''/x$ and
  $m(t+1)\eta\le w''$.
\item $m(S_c+\eta)\le w''L$, $N(m)\eta\le\delta_2mw''$, $\eta^2rL\le\delta_2mw''^2$ and $kS_c\le w''$.
\item $|T_{\mathrm{top}}|\le N^*$; $4N^*x^{\varepsilon}\Lambda_G/k+C_\rho\Lambda_G\log x/k\le x^{-2}$; and
  $b_t\le b'+N^*x^{\varepsilon}\Lambda_Gb'/k$ for every top class $t$.
\item $4S_c\le\Lambda_Gb'/k$, $4C_{\Delta G}k^2\Delta_S\le\Lambda_Gb'$ and $\log\Delta_S\le
  x^{\varepsilon/4}/2$.
\item For every integer $l\ge11$ there is $j$ with $\ell_A\le(l-1)^j$ and $2j+2\le L_w$; $L_w\le
  C_{Lw}\log x$; and $j_0\le C_{j0}\log x$.
\item At a top place $b_1(\tau)=b_t+\lfloor b'/x\rfloor$; away from the top places $u_1-v_1\le2S_c$ and
  $u_k=v_{k-1}-\frac12$; $|e_j-b_j-h_0/2|\le2$ for every $j$ (in particular for $j=m$); $h_0\le
  b_{j+1}-b_j\le h_0+1$; $\prod_\tau\tau(z)^{e_j(\tau)}=1$ for every root of unity $z\in E$;
  $\eta_F=C_{\eta F}x^{-C_F}$; $p_{de}\equiv p'_{de}\equiv1\pmod{15}$; $p_{de},p'_{de}\notin\{g_0,g_1\}$;
  $\log\mathrm N(\mathfrak c)\le L$; and $\mathrm N(\mathfrak c)\le C_{\mathfrak c}d$.
\item $2b/\ell^*+t/2+3\le|b_i(\tau)-e_j(\tau)|$ for all $i,j$.
\item $e_j(\tau)\le X''e_{j'}(\tau')$ for all $j,j',\tau,\tau'$, and $b_m(\tau)\le2b_1(\tau)$.
\item $k^{C_w}\le w''$ and $b^{c_w}\le w''$.
\item $k^{C_\Delta}(C_\Delta m)^3/\eta_F^{C_\Delta}$, $C_\Delta mN(m)$, $k^{C_\Delta}(C_\Delta m^4L)^4$,
  $N(m)m^{10}$ and $(N(m)fL)^{C_{\mathrm{count}}}$ are at most $\Delta_S$.
\item $12b/\ell^*+2t+8$, $ktL$, $C_{h0r}r^{19}f^5L$, $(N(m)fL)^{C_{\mathrm{count}}}$, $4N(m)^3+2t+6$,
  $C_{h0k}km^{9/2}$, $2(r+mt)^3+t+3$, $120rm^{3/2}$, $N(m)m^{10}$ and $16mN(m)$ are at most $h_0$.
\item $t+2\bigl(k+m(t+r+7)\bigr)^3+1\le h_0$ and $fN(m)^{C_{h0N}}\le h_0$.
\item $\gamma$ is less than $v_i-v_{i+1}$, than $v_{k-1}+b_1$ and than $b_{j+1}-b_j$, at every $\tau$.
\item $r(b'+S_c)\ell_A\log x\le\delta_2mb$.
\item $b'+S_c+\eta+b/\ell^*<\ell^*$ and $2\eta+t+8b/\ell^*+7<\ell^*$.
\item $D$ has at most $2^{k-3}$ words.
\item There is $Y$ with $(kb)^8\le Y$, $N(m)^4Q^{C_Y}\le Y^{1/4}$, $\log Y\le C_{YL}L$ and $Y\le
  g_i\le2Y$ for $i=0,1$.
\end{enumerate}
\end{lemma}
\leanref{Development/AbelianVarietyTheory/Proofs/GlobalInterpolation/SzRowPrimesConductorStatement.lean}{GlobalTheory.T.sizes_host_degree_and_log_q_and_conductor_bounds}
\leanref{Development/AbelianVarietyTheory/Proofs/GlobalInterpolation/ExplicitChoiceRankKLowerUpperBounds.lean}{GlobalTheory.EX.explicitChoice_rankK_lower_upper_bounds}
\leanref{Development/AbelianVarietyTheory/Proofs/GlobalInterpolation/RankRowsOfCrudeBounds.lean}{GlobalTheory.EX.szRow_rank_lower_bounds}
\leanref{Development/AbelianVarietyTheory/Proofs/GlobalInterpolation/LogfSqLeRS.lean}{GlobalTheory.EX.szRow_logf}
\leanref{Development/AbelianVarietyTheory/Proofs/GlobalInterpolation/SzRowEtaBoundsStatement.lean}{GlobalTheory.T.sizes_etaSz_le_b_div_and_log_etaSz_le_and_t_le}
\leanref{Development/AbelianVarietyTheory/Proofs/GlobalInterpolation/SzRowEtaLeWStatement.lean}{GlobalTheory.T.sizes_m_mul_Sc_add_etaSz_le_w_mul_L_and_bounds}
\leanref{Development/AbelianVarietyTheory/Proofs/GlobalInterpolation/SzRowTopsBandBottomsStatement.lean}{GlobalTheory.T.sizes_card_TT_le_Nstar_and_rhoP_le_and_bT_le}
\leanref{Development/AbelianVarietyTheory/Proofs/GlobalInterpolation/SzRowCloudLe.lean}{GlobalTheory.EX.szRow_cloud_le}
\leanref{Development/AbelianVarietyTheory/Proofs/GlobalInterpolation/SzRowWindowLength.lean}{GlobalTheory.EX.szRow_windowLength}
\leanref{Development/AbelianVarietyTheory/Proofs/GlobalInterpolation/SzRowExact.lean}{GlobalTheory.EX.szRow_exact}
\leanref{Development/AbelianVarietyTheory/Proofs/GlobalInterpolation/HalfLeAbsNodeSubMid.lean}{GlobalTheory.EX.szRow_lowsFar}
\leanref{Development/AbelianVarietyTheory/Proofs/GlobalInterpolation/ExplicitChoiceRowWeightRatioAndLastCloudLeTwiceFStatement.lean}{GlobalTheory.T.explicitChoice_rowWeight_ratio_and_lastCloud_le_twice_first}
\leanref{Development/AbelianVarietyTheory/Proofs/GlobalInterpolation/SzRowThinSizeGe.lean}{GlobalTheory.EX.szRow_thinSize_ge}
\leanref{Development/AbelianVarietyTheory/Proofs/GlobalInterpolation/SzRowCloudUnitLowerStatement.lean}{GlobalTheory.T.sizes_del_ge_cloudUnit_and_count_bounds}
\leanref{Development/AbelianVarietyTheory/Proofs/GlobalInterpolation/SzRowNodeGapStatement.lean}{GlobalTheory.T.sizes_node_bounds_le_h0}
\leanref{Development/AbelianVarietyTheory/Proofs/GlobalInterpolation/SzRowGapFunFNStatement.lean}{GlobalTheory.T.sizes_t_add_cube_le_h0_and_f_mul_N_pow_le_h0}
\leanref{Development/AbelianVarietyTheory/Proofs/GlobalInterpolation/SzRowSigGaps.lean}{GlobalTheory.EX.szRow_sig_gaps}
\leanref{Development/AbelianVarietyTheory/Proofs/GlobalInterpolation/RbSNumeric.lean}{GlobalTheory.EX.szRow_rbS}
\leanref{Development/AbelianVarietyTheory/Proofs/GlobalInterpolation/SzRowLStarStatement.lean}{GlobalTheory.T.sizes_b_add_Sc_add_etaSz_add_div_lt_lStar}
\leanref{Development/AbelianVarietyTheory/Proofs/GlobalInterpolation/SzRowCardWordsLe.lean}{GlobalTheory.EX.szRow_card_words_le}
\leanref{Development/AbelianVarietyTheory/Proofs/GlobalInterpolation/SzRowTypePrimesStatement.lean}{GlobalTheory.T.sizes_exists_Y_bounds_host_g_mem_Icc_Y_two_mul_Y}

\begin{lemma}\label{asm-dependence-functional-witness-numbers-parameters}
The following holds for explicit values
(Definition~\ref{asm-dependence-functional-witness-numbers-explicit}). If the host satisfies the
arithmetic conditions, the size data carry parameter choices whose constants are those of $C$, and
whose constant $C_{nW}$ is that of $C$.
\end{lemma}
\leanref{Development/AbelianVarietyTheory/Proofs/GlobalInterpolation/ExplicitConstantsParameterConditionsLarge.lean}{GlobalTheory.EX.explicitConstants_parameterConditions_large}
\begin{proof}
The twenty-one statements of
Lemma~\ref{asm-dependence-functional-witness-numbers-parameters-computations} hold for explicit values,
so they hold together; take the data beyond the combined thresholds, and let the host satisfy the
arithmetic conditions. Take the constants of the parameter choices to be those of $C$; standard
constants are positive, as the parameter choices require. The conditions of parameter choices other
than the positivity of the constants are of two kinds. The arithmetic ones are the arithmetic
conditions of the host: $g_i=\pi\bar\pi$ with $\pi\in\mathcal O_E\cap K$, $g_i\equiv1\pmod Q$,
$p_{de}$ and $p'_{de}$ split completely in $E$, and every level prime splits in $K$. The others are
inequalities and identities among the ranks, the host, the Harish-Chandra parameters and the
quantities of the size data, and they are the conclusions of the twenty-one statements, with
$C_{nW}$ that of $C$.
% Lean: EX.explicitConstants_parameterConditions_large: SzRow.imp of the SzRow.and of the twenty-one
% statements; the structure ParameterChoice' takes the constants C.P.* (positive by hC.parConsts_pos),
% the fields qd_pow, ..., g_Y from the destructured conclusions (e_mu from MUQ by .num), and g_norm,
% g_mod, pde_split and level_split from HostArith hA; parConsts P = C.PK and CnW = C.P.CnW by rfl. The
% twenty-one statements are computations from Val, stated with their results and not proved in the
% book. About 6,900 lines of Lean proofs at depth 5.
\end{proof}

\begin{lemma}\label{asm-dependence-functional-witness-numbers-blocks-computations}
In the notation of Definition~\ref{asm-dependence-functional-witness-numbers-explicit}, let $C_t$ be
the constant of the size data, $\delta_2=|T_2|/n$, $h_0=\eta/m$, and let $C_P$, $e_P$, $C_p$, $C_b$,
$c_{TO}$, $C_{hT}$, $C_{bN}$ and $C_{nW}$ be constants of $C$. Put
\[
 z=frx^{\varepsilon},\quad P_P=C_Pz^{e_P},\quad D_p=f\tbinom r2,\quad T_O=\lceil x^{377}\rceil,\quad
 \nu(c)=cr\Bigl(f+\frac{f^2(k-1)^3x^{2\varepsilon}}m\Bigr),\quad n_W=8\lceil\nu(10)/8\rceil.
\]
Each of the following nine statements holds for explicit values.
\begin{enumerate}
\item $P_P,D_p,T_O\le x^{C_t}$.
\item $\nu(C_{nW})\le n_W<\nu(C_{nW})+8$; $8\mid n_W$; $r+t=8fr+4+T_On_W$; $T_O=\lceil x^{c_{TO}}\rceil$;
  $r+6$ is odd; and $fr(r+1)$, $(k-1)f+2r$ and $8fr+4$ are at most $n_W$.
\item $(2fm)^4n_W\le P_P$ and $n_W\le P_P$.
\item $P_P^4D_p^3kfx^{C_p}/(m\delta_2)\le T_O$ and $D_pk/m\le T_O$.
\item $(D_p+2)n_W+2r$, $40(D_p+2)$, $11D_p^2$ and $r+2500r(D_p+1)$ are at most $T_O$.
\item $T_OP_Px^{C_b}\le b$, $T_OP_P\le b^{1/2}$, $P_P^3T_O^2x^{C_b}\le b$ and $P_P^3T_O^2\le b^{1/2}$.
\item $T_O^2\le\Delta_S$ and $(T_O+D_p+2)(T_O+D_p+2+2N(m))\le\Delta_S$.
\item $4T_O^2+2t$, $P_PD_pfm^2N(m)T_O$, $P_P(fmk)^3k/m$ and $C_{hT}T_ON(m)f^2m^2kD_pP_P$ are at most
  $h_0$.
\item $T_OD_p^2x^{C_b}\log b\le b$, $T_OD_p^2\log b\le b^{1/2}$ and $C_{bN}rmkP_P^2D_p^3\le b$.
\end{enumerate}
\end{lemma}
\leanref{Development/AbelianVarietyTheory/Proofs/GlobalInterpolation/SzRowPpDimToLeRpowStatement.lean}{GlobalTheory.T.sizes_PP_and_dim_and_TOv_le_x_rpow_Ct}
\leanref{Development/AbelianVarietyTheory/Proofs/GlobalInterpolation/ExplicitChoiceWindowRankConditionsStatement.lean}{GlobalTheory.T.explicitChoice_windowRank_conditions}
\leanref{Development/AbelianVarietyTheory/Proofs/GlobalInterpolation/ExplicitChoiceWindowRankLePolyBoundStatement.lean}{GlobalTheory.T.explicitChoice_windowRank_le_polyBound}
\leanref{Development/AbelianVarietyTheory/Proofs/GlobalInterpolation/SzRowTOLowerStatement.lean}{GlobalTheory.T.sizes_PP_pow_mul_dim_pow_div_le_TOv}
\leanref{Development/AbelianVarietyTheory/Proofs/GlobalInterpolation/SzRowRankPolysLeTOStatement.lean}{GlobalTheory.T.sizes_rank_and_poly_counts_le_TOv}
\leanref{Development/AbelianVarietyTheory/Proofs/GlobalInterpolation/SzRowTOPPLeBStatement.lean}{GlobalTheory.T.sizes_TOv_mul_PP_le_b_and_sqrt_b}
\leanref{Development/AbelianVarietyTheory/Proofs/GlobalInterpolation/SzRowDelTOGapStatement.lean}{GlobalTheory.T.sizes_TOv_sq_le_del_and_gap_mul_le_del}
\leanref{Development/AbelianVarietyTheory/Proofs/GlobalInterpolation/SzRowNodeGapBlockLength.lean}{GlobalTheory.EX.szRow_nodeGap_blockLength}
\leanref{Development/AbelianVarietyTheory/Proofs/GlobalInterpolation/ExplicitChoicePolynomialLeBaseSizeStatement.lean}{GlobalTheory.T.explicitChoice_polynomial_le_baseSize}

\begin{lemma}\label{asm-dependence-functional-witness-numbers-blocks}
The following holds for explicit values
(Definition~\ref{asm-dependence-functional-witness-numbers-explicit}). For all parameter choices for
the size data whose constant $C_{nW}$ is that of $C$, there are $T$-parameters
(Definition~\ref{asm-dependence-functional-witness-sizes}(1)) whose constants are those of $C$.
\end{lemma}
\leanref{Development/AbelianVarietyTheory/Proofs/GlobalInterpolation/SzRowNodeGapBlockLength.lean}{GlobalTheory.EX.szRow_blockLengthParams}
\begin{proof}
The nine statements of
Lemma~\ref{asm-dependence-functional-witness-numbers-blocks-computations} hold for explicit values,
so they hold together; take the data beyond the combined thresholds, and let parameter choices for
the size data be given whose constant $C_{nW}$ is that of $C$. Take $T_O$, $n_W$, $P_P$ and $D_p$ as
there, $n_P=8fr+4$, $s=r+6$, and the constants $C_P$, $e_P$, $C_p$, $C_b$, $c_{TO}$, $C_{hT}$ and $C_{bN}$
of $C$. Standard constants have $C_P\ge1000$, $e_P=12$, $C_p=1$, $c_{TO}=377$ and $C_b,C_{hT},C_{bN}\ge1$,
which gives the lower bounds required of these constants. The equalities $n_P=8fr+4$, $D_p=f\binom r2$
and $P_P=C_Pz^{e_P}$ hold by the choices, and $r+5<s\le r+7$. The remaining conditions of
$T$-parameters are the conclusions of the nine statements: the window for $n_W$ about
$\nu(C_{nW})$, with $C_{nW}$ that of the parameter choices, $8\mid n_W$, the splitting
$r+t=n_P+T_On_W$, $T_O=\lceil x^{c_{TO}}\rceil$, the parity of $s$, and the bounds (1) and (3) to (9).
% Lean: EX.szRow_blockLengthParams, from T.sizes_PP_and_dim_and_TOv_le_x_rpow_Ct,
% T.explicitChoice_windowRank_conditions, T.explicitChoice_windowRank_le_polyBound,
% T.sizes_PP_pow_mul_dim_pow_div_le_TOv, T.sizes_rank_and_poly_counts_le_TOv,
% T.sizes_TOv_mul_PP_le_b_and_sqrt_b, T.sizes_TOv_sq_le_del_and_gap_mul_le_del,
% szRow_nodeGap_blockLength and T.explicitChoice_polynomial_le_baseSize; about 2,000 lines of Lean
% proofs at depth 5. The structure ParamsT: TO = TOv, nWn = nWv, nP = 8 f r + 4, PP = CP zBase^eP
% (zBase = f r L), Dp = f (r choose 2), s = r + 6, the constants of C.T; one_le_CP, eP_ge, Cp_ge,
% Cb_nonneg, ChT_ge, CbN_ge from Std C (1000 <= CP, eP = 12, Cp = 1, cLast); the remaining fields are
% the conjuncts of the nine statements (nW_win with hP : P.CnW = C.P.CnW). The nine statements are
% computations from Val, stated with their results and not proved in the book.
\end{proof}

\begin{lemma}\label{asm-dependence-functional-witness-numbers-bounds-computations}
In the notation of Definition~\ref{asm-dependence-functional-witness-numbers-explicit}, with the scale
$(x,\varepsilon,J)$, put $\vartheta=\varepsilon/100$, $L=x^{\varepsilon}=\log b$, $\delta_2=|T_2|/n$,
$N^*=2n\delta_1$, $h_0=\eta/m$, $w''=w''(n)$, $X''=X''(n)$, $Q=Q_Sqd$, $\ell_A=kx^{\varepsilon}/(m\delta_2)$ and
\[
 \ell^*=\frac{mb}{nt^2(N(m)+t)L^2+kn^2(S_c+t+L)},
\]
and let $C_t$ and $C_\Gamma$ be the constants of the size data. Each of the following eleven statements
holds for explicit values.
\begin{enumerate}
\item $x^{2+8\varepsilon-\vartheta}\le n\le x^{2+8\varepsilon+\vartheta}$,
  $x^{-\vartheta}\sqrt{x\delta_1}\le m\le x^{\vartheta}\sqrt{x\delta_1}$,
  $x^{9/2+10\varepsilon-\vartheta}\le k\le x^{9/2+10\varepsilon+\vartheta}$, $1\le t\le x^{C_t+\vartheta}$,
  $x^{-\vartheta}n\delta_1\le|T_{\mathrm{top}}|\le N^*$, $x^{-\varepsilon}\le\delta_2$ and $L_w\le\log^2x$.
\item $x^{\vartheta}m^5fLx^{\varepsilon}\le k$, $8C_\Gamma m^5fL\le k$, $k\le\sqrt x\,n\cdot nx^{-\varepsilon}$,
  $\sqrt{k/m}\le nx^{-\varepsilon}$, $m^2L\le k$, $m\sqrt x\log x\le k$, $x^{\vartheta}m^6L\le k\log k$,
  $x^{\vartheta}m^3L\log\mathrm{rd}_E\le k\log^2k$, $\log f\le\sqrt{\lfloor m/2\rfloor-1}$, $\log^2f\le
  m^2L$, $x^{\vartheta}\log\log b\le rf\log r$, $x^{\vartheta}L\le f\log k$, $f\le m^2n$ and $f\le r^2$.
\item $Q^{36}t^{150}\le\sqrt b$; $x^{\varepsilon/(6J)}\le\log q$; $x^{\varepsilon/(5J)-\varepsilon/(100J)}$ is
  at most $\log q$ and $\log d$, which are at most $x^{\varepsilon/5+\vartheta}$; $|L'/L(1)|\le
  x^{\varepsilon}$; $1\le\mathrm{rd}_E\le\frac{Q_S}2\sqrt{qd}$; and $\log p_{de},\log p'_{de}\le L$.
\item $x^{\varepsilon}\le\Lambda_G\le x^{\varepsilon+\vartheta}$,
  $4N^*x^{\varepsilon}\Lambda_G/k+x^{\vartheta}\Lambda_G\log x/k\le x^{-2}$, and $b_t\le
  b'+N^*x^{\varepsilon}\Lambda_Gb'/k$ for every top class $t$.
\item $t\le\eta\le b/(t^3fk^3)$, $\log\eta\le x^{\varepsilon/4}$, $t\eta\le w''/x$, $m(t+1)\eta\le w''$,
  $m(S_c+\eta)\le w''L$, $N(m)\eta\le\delta_2mw''$, $\eta^2rL\le\delta_2mw''^2$ and
  $r(b'+S_c)\ell_A\log x\le\delta_2mb$.
\item $12b/\ell^*+2t+8$, $ktL$, $x^{\vartheta}r^{19}f^5L$, $4N(m)^3+2t+6$, $x^{\vartheta}km^{9/2}$,
  $2(r+mt)^3+t+3$ and $N(m)m^{10}$ are at most $h_0$.
\item $b'+S_c+\eta+b/\ell^*<\ell^*$ and $2\eta+t+8b/\ell^*+7<\ell^*$.
\item $N(m)m^{10}\le\Delta_S$, $x^{\vartheta}mN(m)\le\Delta_S$, $\log\Delta_S\le x^{\varepsilon/4}/2$,
  $kS_c\le w''$, $x^{\vartheta}mk\Delta_S\le w''$, $4S_c\le\Lambda_Gb'/k$, $4x^{\vartheta}k^2\Delta_S\le
  \Lambda_Gb'$ and $x^{\vartheta}m\log(\mathrm{rd}_E/2)\le k$.
\item $\sqrt b\le w''$; $b_m(\tau)\le2b_1(\tau)$ for every $\tau$; and $e_j(\tau)\le X''e_{j'}(\tau')$ for all
  $1\le j,j'\le m$ and all $\tau,\tau'$.
\item $(kb)^8\le g_i\le2g_{i'}$ and $\log g_i\le x^{\vartheta}L$ for $i,i'\in\{0,1\}$, and $D$ has at most
  $2^{k-3}$ words.
\item The splitting bound is $B\ge11$.
\end{enumerate}
\end{lemma}
\leanref{Development/AbelianVarietyTheory/Proofs/GlobalInterpolation/SzRowRankSizesStatement.lean}{GlobalTheory.T.sizes_sizeN_sizeM_sizeK_sizeT_sizeTT_sizeDe2_sizeLw}
\leanref{Development/AbelianVarietyTheory/Proofs/GlobalInterpolation/SzRowSizeKRelStatement.lean}{GlobalTheory.T.sizes_sizeKRel}
\leanref{Development/AbelianVarietyTheory/Proofs/GlobalInterpolation/SzRowCondQdPdeStatement.lean}{GlobalTheory.T.sizes_sizeCond_sizeQD_sizePde}
\leanref{Development/AbelianVarietyTheory/Proofs/GlobalInterpolation/SzRowLaRhoBotStatement.lean}{GlobalTheory.T.sizes_sizeLa_sizeRho_sizeBot}
\leanref{Development/AbelianVarietyTheory/Proofs/GlobalInterpolation/ExplicitChoiceSizeEtaStatement.lean}{GlobalTheory.T.explicitChoice_sizeEta}
\leanref{Development/AbelianVarietyTheory/Proofs/GlobalInterpolation/ExplicitChoiceSizeNodesStatement.lean}{GlobalTheory.T.explicitChoice_sizeNodes}
\leanref{Development/AbelianVarietyTheory/Proofs/GlobalInterpolation/ExplicitChoiceSizeLstarStatement.lean}{GlobalTheory.T.explicitChoice_sizeLstar}
\leanref{Development/AbelianVarietyTheory/Proofs/GlobalInterpolation/SzRowSizeDelStatement.lean}{GlobalTheory.T.sizes_sizeDel}
\leanref{Development/AbelianVarietyTheory/Proofs/GlobalInterpolation/SzRowSizeWSizeBmSizeXppStatement.lean}{GlobalTheory.T.sizes_sizeW_sizeBm_sizeXpp}
\leanref{Development/AbelianVarietyTheory/Proofs/GlobalInterpolation/SzRowSizeGSizeWordsStatement.lean}{GlobalTheory.T.sizes_sizeG_sizeWords}
\leanref{Development/AbelianVarietyTheory/Proofs/GlobalInterpolation/SzRowElevenLeSplitBound.lean}{GlobalTheory.EX.szRow_eleven_le_splitBound}
\leanref{Formalism/SharpRestriction/ElevenLe.lean}{Formalism.MainNotions.SizesAt}

\begin{lemma}\label{asm-dependence-functional-witness-numbers-bounds}
The size bounds (Definition~\ref{asm-dependence-functional-witness-sizes}(2)) hold for explicit
values (Definition~\ref{asm-dependence-functional-witness-numbers-explicit}).
\end{lemma}
\leanref{Development/AbelianVarietyTheory/Proofs/GlobalInterpolation/SzRowElevenLeSplitBound.lean}{GlobalTheory.EX.szRow_sizesAt}
\begin{proof}
The size bounds are the inequalities of the eleven statements of
Lemma~\ref{asm-dependence-functional-witness-numbers-bounds-computations}. Each of them holds for
explicit values, so they hold together.
% Lean: EX.szRow_sizesAt, by SzRow.and of the eleven statements and the reordering of their conjuncts
% into SizesAt (SizeN, SizeM, SizeK, SizeKRel, SizeT, SizeCond, SizeQD, SizeTT, SizeLa, SizeRho, SizeBot,
% SizeDe2, SizeLw, SizeEta, SizeNodes, SizeLstar, SizeDel, SizeW, SizeBm, SizeXpp, SizeG, SizePde,
% SizeWords, 11 <= splitBound). The eleven statements are computations from Val, stated with their
% results and not proved in the book.
% Lean: EX.szRow_sizesAt, from the eleven SzRow statements
% T.sizes_sizeN_sizeM_sizeK_sizeT_sizeTT_sizeDe2_sizeLw, T.sizes_sizeKRel,
% T.sizes_sizeCond_sizeQD_sizePde, T.sizes_sizeLa_sizeRho_sizeBot, T.explicitChoice_sizeEta,
% T.explicitChoice_sizeNodes, T.explicitChoice_sizeLstar, T.sizes_sizeDel, T.sizes_sizeW_sizeBm_sizeXpp,
% T.sizes_sizeG_sizeWords and szRow_eleven_le_splitBound; about 4,500 lines of Lean proofs at depth 5.
\end{proof}

\begin{lemma}\label{asm-dependence-functional-witness-numbers-band}
The following holds for explicit values
(Definition~\ref{asm-dependence-functional-witness-numbers-explicit}):
$g_0,g_1\le\exp(9x^{\varepsilon})$, $p_{de},p'_{de}<X'$ and $x^{\varepsilon}\le\Lambda_G$.
\end{lemma}
\leanref{Development/AbelianVarietyTheory/Proofs/GlobalInterpolation/BandWidthLeInvSqOfRankBounds.lean}{GlobalTheory.EX.szRow_band_lambda_primes}
\leanref{Development/AbelianVarietyTheory/Proofs/GlobalInterpolation/GoodTopCountLeConstMulLog.lean}{GlobalTheory.topCount_and_lambda_bounds}
\leanref{Development/AbelianVarietyTheory/Proofs/GlobalInterpolation/HostBoundsTwoMulRpowLogSqLe.lean}{GlobalTheory.EX.hostBounds}
\leanref{Development/AbelianVarietyTheory/Proofs/GlobalInterpolation/ScalePolynomialQuantitiesComparison.lean}{GlobalTheory.scale_poly_comparison}
\begin{proof}
The constants $c_D$ and $c_h$ play no role: take $c_1=c_2=0$. Let $x_0\ge512$ be so large that for
$x\ge x_0$
\[
 x^{41}X'\le b\qquad\text{and}\qquad
 \exp\bigl(2C_8x^{\varepsilon/5}\log^2x\bigr)\,x^{2C_8}\le X';
\]
both hold for large $x$, since $41\log x+x^{\varepsilon/3}\le x^{\varepsilon}$ and
$2C_8(x^{\varepsilon/5}\log^2x+\log x)\le x^{\varepsilon/3}$ for large $x$. Let the size data take the
explicit values at a scale $x\ge x_0$.

\emph{The level primes $g_i$.} By condition (1), $k\le2x^{9/2+10\varepsilon}\le2x^5$, since
$\varepsilon<1/20$. By condition (2), $g_i\le2(kb)^8$. Hence
\[
 g_i\le2(2x^5b)^8=512\,x^{40}b^8\le x^{41}b^8\le b^9=\exp(9x^{\varepsilon}),
\]
using $x\ge512$ and $x^{41}\le x^{41}X'\le b$.

\emph{The primes $p_{de}$ and $p'_{de}$.} By condition (2) the host has a pair scale $Y$, with
$x^{1/J}\le Y\le x$ and $\log N\le2Y^{\varepsilon/5}\log^2Y$; since $1\le Y\le x$, this gives $\log
N\le2x^{\varepsilon/5}\log^2x$. Also $p_{de}\le N^{C_8}$, so $\log p_{de}\le2C_8x^{\varepsilon/5}\log^2x$,
and as $C_8\ge1$ and $x>1$,
\[
 p_{de}\le\exp\bigl(2C_8x^{\varepsilon/5}\log^2x\bigr)<\exp\bigl(2C_8x^{\varepsilon/5}\log^2x\bigr)\,x^{2C_8}\le X'.
\]
The same argument applies to $p'_{de}$.

\emph{The width $\Lambda_G$.} By condition (5), $\Lambda_G=K_\Lambda N_1^2\log(xb)$, and
$\log(xb)=\log x+x^{\varepsilon}\ge x^{\varepsilon}$. The number $N_1$ of
Definition~\ref{hf-good-orbit} is an integer part plus $1$, so $N_1\ge1$, and $K_\Lambda\ge1$. So
$x^{\varepsilon}\le\Lambda_G$.
% Lean: EX.szRow_band_lambda_primes with c1 = c2 = 0. The threshold is the maximum of 512 and of the
% numbers of topCount_and_lambda_bounds (N1 >= 1 and x^eps <= G.La when K.La >= 1), EX.hostBounds
% (log pde, log pde' <= 2 Cpde x^(eps/5) log^2 x, from PairRow's log N bound, PairWin and HostVal's
% pde <= N^Cpde), scale_poly_comparison eps 41 1 (x^41 X' <= b) and scale_poly_comparison eps (2 Cpde) 1
% (exp(2 Cpde x^(eps/5) log^2 x) x^(2 Cpde) <= X'), and two numbers for rhoB. The g bound is the
% calc from (hg j).2 and hV.ranks; La_eq gives Lambda_G. Std C gives 1 <= Cpde and 1 <= GK.La (in
% cMid). The conclusion also contains rhoB <= x^(-2), ChTT_at L and ChCode_at L 3 C.P.CF, which the
% book does not use. About 1,850 lines of Lean proofs at depth 5.
\end{proof}

\begin{lemma}\label{asm-dependence-functional-witness-numbers-values}
The following holds for explicit values
(Definition~\ref{asm-dependence-functional-witness-numbers-explicit}), with the notation there. If
the host satisfies the arithmetic conditions, then the size data carry parameter choices with
$T$-parameters whose constants are those of $C$, among them $C_{nW}$, and satisfy the size bounds
(Definition~\ref{asm-dependence-functional-witness-sizes}(1) and (2)); and
$x^{2+8\varepsilon}\le n\le2x^{2+8\varepsilon}$, $\sqrt{x\delta_1}\le m\le2\sqrt{x\delta_1}$,
$x^{9/2+10\varepsilon}\le k\le2x^{9/2+10\varepsilon}$, $g_0,g_1\le\exp(9x^{\varepsilon})$ and
$\mathfrak c=\mathcal O_K$.
\end{lemma}
\leanref{Development/AbelianVarietyTheory/Proofs/GlobalInterpolation/SzRowNumAt.lean}{GlobalTheory.EX.szRow_numAt}
\leanref{Formalism/SharpRestriction/InertJacobiLevels/AsmIface.lean}{Formalism.MainNotions.Inhab.NumAt}
\begin{proof}
Lemmas~\ref{asm-dependence-functional-witness-numbers-parameters},
\ref{asm-dependence-functional-witness-numbers-blocks},
\ref{asm-dependence-functional-witness-numbers-bounds}
and~\ref{asm-dependence-functional-witness-numbers-band} hold for explicit values, so they hold
together. Let the size data take the explicit values, beyond the combined thresholds, and let the
host satisfy the arithmetic conditions. The first lemma gives parameter choices whose constants are
those of $C$, among them $C_{nW}$; the second gives $T$-parameters for them whose constants are those
of $C$; the third gives the size bounds; and the fourth gives $g_0,g_1\le\exp(9x^{\varepsilon})$. The
windows for $n$, $m$ and $k$ are those of an adapted rank datum, condition (1) of the explicit
values, and $\mathfrak c=\mathcal O_K$ is part of condition (2).
% Lean: szRow_numAt = szRow_and_imp of the four SzRow statements; NumAt C L with params, parT,
% sizes, g_le from the four, win from Val.ranks and ccK_top from Val.host. The Lean's NumAt also
% records rho_le, LaG_ge, pde_lt, bT_le, Lw_eq (from Val.host and C.Ct = 401), cloud1 (from Val.hc),
% logFiveDepth_eq and the constants, and szRow_numAt also gives ChTT_at L and ChCode_at L 3 C.P.CF;
% the book states the parts it uses.
\end{proof}

\begin{lemma}\label{asm-dependence-functional-witness-numbers}
Eventually in the context (Definition~\ref{asm-stage-construction}(1)), the following holds at every
base point of a core. Let $D$ be its definite-side datum, with ranks $k,m,n$ and level primes
$g_0,g_1$ and ideal $\mathfrak c\subset\mathcal O_K$ of its host, and with its size data, at the scale
$(x,\varepsilon)$ of the context; put $\delta_1=x^{-\varepsilon/8}$. Then the size data carry parameter
choices with $T$-parameters whose constant $C_{nW}$ is $10$, and they satisfy the size bounds
(Definition~\ref{asm-dependence-functional-witness-sizes}(1) and (2)); moreover
$x^{2+8\varepsilon}\le n$, $\sqrt{x\delta_1}\le m\le2\sqrt{x\delta_1}$, $x^{9/2+10\varepsilon}\le k$,
$g_0,g_1\le\exp(9x^{\varepsilon})$ and $\mathfrak c=\mathcal O_K$.
\end{lemma}
\leanref{Development/AbelianVarietyTheory/Proofs/GlobalInterpolation/SzRowNumAt.lean}{GlobalTheory.EX.szRow_numAt}
\leanref{Development/AbelianVarietyTheory/Definitions/GlobalTwisting/PairRowIffQdRows.lean}{GlobalTheory.EX.Val}
\leanref{Development/AbelianVarietyTheory/Definitions/GlobalTwisting/PairRowIffQdRows.lean}{GlobalTheory.EX.SzRow}
\begin{proof}
A base point records that its size data are at the scale of the context and take the explicit
values for the constants $(C,c_D,c_h)$ of the context
(Definition~\ref{asm-dependence-functional-witness-numbers-explicit}), and that its host satisfies
the arithmetic conditions. The order in which the constants of a context are chosen is that of
Definition~\ref{asm-dependence-functional-witness-numbers-explicit}: the exponent and the standard
constants $C$, then $c_D$, then $c_h$, then the scale. So
Lemma~\ref{asm-dependence-functional-witness-numbers-values}, which holds for explicit values, holds
at every base point of a core eventually in the context. It gives the conclusion, with $C_{nW}$ the
constant of $C$; the constants of a context are standard, and their fixed values include
$C_{nW}=10$.
% Lean: AtSize.eventually_sizeRows_alt applies EX.szRow_numAt (Lemma
% asm-dependence-functional-witness-numbers-values) to the base point, with Val from
% beta.p.pt_val beta.val and HostArith from Plug.hostArith_bpt; CnW = 10 is the field of FC.Fixed in
% Std C.
\end{proof}

\begin{lemma}\label{asm-dependence-functional-witness-computation}
Let $D$ be a definite-side datum with size data at the scale $(x,\varepsilon)$, in the notation of
Definition~\ref{asm-dependence-functional-witness-sizes}, put $\delta_1=x^{-\varepsilon/8}$ and
$X'=\exp(x^{\varepsilon/3})$, and suppose that the size data carry parameter choices with
$T$-parameters whose constant $C_{nW}$ is $10$ and satisfy the size bounds. Suppose that
$x^{2+8\varepsilon}\le n$, $\sqrt{x\delta_1}\le m\le2\sqrt{x\delta_1}$, $x^{9/2+10\varepsilon}\le k$,
$g_0,g_1\le\exp(9x^{\varepsilon})$, $q_0<X'$, $\mathfrak c=\mathcal O_K$, $x^{\varepsilon/4}\ge8$
and $x\ge44^4$. Then the rows of the size data have the sizes of the size data.
\end{lemma}
\leanref{Development/AbelianVarietyTheory/Statements/PlaceShiftStatements/SetAtOfSizeAt.lean}{GlobalTheory.BigTop.Fun.SetAtModel.sizeRows_of}
\leanref{Development/AbelianVarietyTheory/Statements/PlaceShiftStatements/SetAtOfSizeAt.lean}{GlobalTheory.BigTop.Fun.SetAtModel.hc_bounds}
\leanref{Development/AbelianVarietyTheory/Statements/PlaceShiftStatements/SetAtOfSizeAt.lean}{GlobalTheory.BigTop.Fun.SetAtModel.gap_le}
\leanref{Development/AbelianVarietyTheory/Statements/PlaceShiftStatements/SetAtOfSizeAt.lean}{GlobalTheory.BigTop.Fun.SetAtModel.cloudLen_le}
\leanref{Development/AutomorphicArithmetic/CyclotomicBounds/RamifiedDiscriminants/FourMulQLeCondN.lean}{UpperTheory.Ram.Rd.rdE_le_cond}
\leanref{Development/AutomorphicArithmetic/CyclotomicBounds/RamifiedDiscriminants/FourMulQLeCondN.lean}{UpperTheory.Ram.Rd.log_cond_le}
\begin{proof}
Put $z=x^{\varepsilon/4}$ and $s=x^{1/4}$. Then $x^{\varepsilon}=z^4$, $x=s^4$, $z\ge8$ and $s\ge44$;
so $z^3\ge512$ and $512z\le x^{\varepsilon}$.

\emph{Ranks.} Since $\delta_1\le1$, $m\le2\sqrt x=2s^2$. Since $s^2=x^{1/2}\le x^{1-\varepsilon/8}=x\delta_1$,
$m\ge\sqrt{x\delta_1}\ge s\ge44$. Also $n\ge x^{2+8\varepsilon}\ge s^8x^{\varepsilon}$ and
$k\ge x^{9/2+10\varepsilon}\ge s^{18}x^{2\varepsilon}$. The inequalities of a rank datum
(Definition~\ref{th-ranks}) give $f\le k\le t$, $n\le t$, $m+1\le t$ and $t>3m\ge132$.

\emph{Host.} The size bounds give $\log q,\log d\le x^{\varepsilon/5+\varepsilon/100}\le z$ and
$\mathrm{rd}_E\ge1$, and the parameter choices give $t\le\eta$ and $\log\eta\le x^{\varepsilon/4}$, so
$\log t\le z$. The root discriminant satisfies $\mathrm{rd}_E\le Q_Sqd$, and
$\log(Q_Sqd)\le2\log n+\log q+\log d\le4z\le x^{\varepsilon}/100$. Put $B=\sqrt b$. The parameter
choices give $(Q_Sqd)^{36}t^{150}\le B$, so $dt^9\le B$, $t^9\le B$ and $1\le B\le b$. Since $x\le k\le t$
and $\log t\le z\le x^{\varepsilon/3}$, $x\le X'$; so $w''=xb/(X'xn)\le b$, and the bound $kS_c\le w''$ of
the parameter choices, with $k\ge4$, gives $4S_c\le b$.

\emph{The conditions.} Conditions (1), (2) and (10) of
Definition~\ref{asm-dependence-functional-witness-sizes} are hypotheses, and (7) was shown above. For
(3), $m^4\le16s^8\le2s^8x^{\varepsilon}\le2n=f$, and $m\ge44$. For (4),
$x^{2+\varepsilon}m^3\log\mathrm{rd}_E\le s^8(2s^2)^3x^{2\varepsilon}=8s^{14}x^{2\varepsilon}\le
s^{18}x^{2\varepsilon}\le k$, since $\log\mathrm{rd}_E\le x^{\varepsilon}$ and $8\le s^4$; and
$N(m)\le t+t(t+1)\le t^3$, so $mfN(m)^2d\le dt^8\le dt^9\le B\le b$. For (5), $f\le k\le r+1$ and $m\le
r+1$. For (6), $\log g_i\le9x^{\varepsilon}$, and $\log q_0<\log X'=x^{\varepsilon/3}\le
x^{\varepsilon}$. For (8), the number $\ell$ of Definition~\ref{asm-dependence-rows-parameters}(2) is at
most $S_c\le b/4$, so $\sqrt\ell\le B$, and as $m\ge44$ this gives $\gamma(1)\le4t^7B$. With
$16t^8\le t^9\le B$, $(m+1)\gamma(1)\le4t^8B\le B^2/4=b/4$, and $S_c+(m+1)\gamma(1)+\frac12\le
b/2+\frac12\le b$. For (9), the windows of the size data place $e_l(\tau)$ between $w''$ and $4b$, by
$w''\le b'/x\le b'\le b$ and the bound $b_t\le2b'$ at the top classes that the parameter choices give;
and $|b_l(\tau)-e_l(\tau)|\le\eta$ by the size data, while $t\eta\le w''/x$ by the parameter choices,
so $\eta<w''$ and $b_l(\tau)>0$.
\end{proof}

\begin{lemma}\label{asm-dependence-functional-witness-twisted}
Let $p$ be parameters (Definition~\ref{asm-dependence-parameters}), let a bundle be linked by
$\Lambda_p$ to a base point of a core in a context with scale $(x,\varepsilon)$, built at a stage
$(K,B,f,T)$ with a witness, and let $R$ be the indefinite realization of the bundle, with the face
$F_R$ of the bundle. Let $C,C_H$ be real. A finite place $w$ is a twisted place of $R$ if and only
if, for some lift $j$, one of the following is nonzero: the index $I_w(j,\mu_0)$ or $I_{\bar
w}(j,\mu_0)$ of the rows $S_0(R,F_R)$; the relative index of $T_{j,\mu_0}$ and its dual for
$H_T(j,\mu_0)$ at $w$ or at $\bar w$; or $F_0(w,j)$ or $F_0(\bar w,j)$, for the $F_0$ of the numerical
bounds $N(R,F_R;C,C_H)$ (Definition~\ref{is-realisation}).
\end{lemma}
\leanref{Development/AutomorphicArithmetic/ThetaCorrespondence/ThetaPairLifts/ThirdPlaceSeries/RealPins3RealTOf.lean}{GlobalTheory.Mem.R3.realPins3_realTOf_alt}
\leanref{Development/AutomorphicArithmetic/ThetaCorrespondence/ThetaPairLifts/ThirdPlaceSeries/RealPins3Iff.lean}{Formalism.IndefR.IndefiniteRealisationNormalisationDegreeThree}
\leanref{Development/AutomorphicArithmetic/ThetaCorrespondence/ThetaPairLifts/ThirdPlaceSeries/WeightOnePlacesFiniteOfWitness.lean}{GlobalTheory.Mem.Top.Bl.zscale_f0_of_witness_alt}
\leanref{Development/AbelianVarietyTheory/Definitions/AutomorphicMembership/NonzeroPlacesOverPrimes.lean}{GlobalTheory.Mem.Top.Bl.F0K}
\leanref{Development/AbelianVarietyTheory/Proofs/IwasawaDecomposition/ForcedContentEqClosedContentNumericalDatum.lean}{GlobalTheory.Mem.Top.Bl.forcedContent_eq_closedContent_numericalDatum}
\leanref{Development/AbelianVarietyTheory/Definitions/AutomorphicMembership/ComplexSidePlacesIffMemNonzeroIndexPlaces.lean}{GlobalTheory.Mem.R3.Dat.twisted_places_eq_complex_side_of_eq_nonzero_index}
\leanref{Development/AbelianVarietyTheory/Definitions/AutomorphicMembership/DefiniteSideNonzeroIndexPlacesAnyColumnsScale.lean}{GlobalTheory.Mem.R3.complexSide_nonzeroIndexPlaces_iff}
\leanref{Development/AbelianVarietyTheory/Definitions/AutomorphicMembership/DefiniteSideNonzeroIndexPlacesAnyColumnsScale.lean}{GlobalTheory.Mem.R3.complexSide_numberContents_eq_scaleAtSubspaces}
\begin{proof}
Write $D$ for the datum of the stage, and for a finite place $y$ of its coefficient field and a lift
$j$ let $F^K_0(y,j)$ be the local term of $D$ at the scale of the context for the lift $j+1$ and
the residue characteristic of $y$.

\emph{The local terms.} The bundle records that the archimedean scalars of $D$ are defined, so
$F^K_0(y,j)$ is the local term
$F_{j+1}(y)$ of the numerical datum of $D$ and its size data at that scale. The real package of the
bundle has scale data with these local terms as its $F_0$, at every choice of definite subspaces.
The witness identifies this package with the package of its real input, so the scale data of its
twist arguments, at the definite subspaces of their realization, have $F_0=F^K_0$.

\emph{The stage.} The rows of the stage are exact: its set $W_t$ of twisted places is the set of
places $w$ at which, for some lift $j$, the index $I(j,\mu_0)$, the relative index of
$T_{j,\mu_0}$ for $H_T(j,\mu_0)$ or $F^K_0(\cdot,j)$ is nonzero at $w$ or at $\bar w$, these indices
being those of the realization built at the stage at its face. The twisted places of that
realization are the places of $W_t$.

\emph{The face.} The realization of the twist arguments is the realization built at the stage, with
the scalars $\lambda$ and the error terms of the twist arguments, at the face obtained by replacing
the definite columns, periods and scale data by those of the twist arguments. The numerical bounds
$N(\cdot\,;C,C_H)$ of this face have the $F_0$ of its scale data, which is $F^K_0$. Replacing the
scalars and the error terms does not change the twisted places; with the change of the columns it
does not change the indices of the rows either. So the
places at which an index of the realization of the twist arguments, or its $F_0$, is nonzero are
those of the realization built at the stage. With the two previous paragraphs, the twisted places of
the realization of the twist arguments are the places of the statement for it.

The witness identifies the realization and the face of its real input, which are those of its twist
arguments, with $R$ and $F_R$. So the statement holds for $R$.
\end{proof}

\begin{lemma}\label{asm-dependence-functional-witness-core}
Let $R$ be a realization of a definite-side datum for a family of definite forms
(Definition~\ref{is-realisation}), with a face $F_R$. If $R$ satisfies the weak characterization at
$F_R$ and exists at every place, then $R$ satisfies the core conditions at $F_R$
(Definition~\ref{is-eigenvalue-lift-data}).
\end{lemma}
\leanref{Development/AbelianVarietyTheory/Proofs/DivisorProofs/CoreClausesOfWeakCharacterisation.lean}{GlobalTheory.CoreOfWeak.core_of_weak}
\leanref{Formalism/LocalToGlobal/IndefiniteReal/ZeroPartEigenvalueConvention.lean}{Formalism.IndefR.Links.RealDat3.characterisation_of_weak}
\leanref{Formalism/LocalToGlobal/IndefiniteReal/LinksW.lean}{Formalism.IndefR.Links.RealDat.chain_of_weak}
\leanref{Formalism/LocalToGlobal/IndefiniteReal/LinksW.lean}{Formalism.IndefR.Links.RealDat.thP_of_weak}
\leanref{Formalism/LocalToGlobal/IndefiniteReal/LinksW.lean}{Formalism.IndefR.Links.RealDat.ExAt}
\begin{proof}
The core conditions are the conditions of the characterization of $R$ on its level characters, its
group, its family, its models, its chain of theta maps and its theta lifts. The weak
characterization contains the first four. For the chain of theta maps and the theta lifts it
contains weak forms: for each iterated theta map between the models of $Y_{k+1}$ and $Y_k$ and from
$X'$ to $Y_{m-1}$, each map $\mathrm{io}$, each theta lift $\mathrm{Th}_t$ and each theta map
$\mathrm{th}_F(i,t)$ and $\mathrm{th}_G(i,t)$, the condition of the characterization holds if the map
exists. That $R$ exists at every place says that each of these maps exists. So the conditions on
the chain and on the theta lifts hold.
\end{proof}

\begin{lemma}\label{asm-dependence-functional-witness}
Let local Langlands maps be given and let $p$ be parameters
(Definition~\ref{asm-dependence-parameters}).
\begin{enumerate}
\item Eventually in the context (Definition~\ref{asm-stage-construction}(1)), at every base point of
  a good core of $p$, the rows of the size data of the datum have the sizes of the size data
  (Definition~\ref{asm-dependence-functional-witness-sizes}), and the logarithmic derivative at $1$
  of the size data has absolute value at most $x^{\varepsilon}$.
\item Let a bundle be linked by $\Lambda_p$ to a base point of a core, built at a stage
  $(K,B,f,T)$ with a witness, and let $R$ be the indefinite realization of the bundle, with the face
  $F_R$ of the bundle. Then the columns of $R$ are given by the framing $f$; the twisted places of
  $R$ lie off $q_0$; every character $\xi_t$ of the level system of $R$ lies in the set consisting of
  the trivial character, the characters $\delta$ and $\delta'$ of the framing $f$, and the characters
  $\delta_\ell\psi_w$, where $w$ runs over the finite places of the coefficient field, $\psi_w$ is the
  character that $T$ chooses at $w$, and $\delta_\ell=\delta'$ if the residue characteristic $\ell$ of $w$ is $p_{de}$ and
  $\delta_\ell=\delta$ otherwise; every finite place at which an index of $R$ at $F_R$ is nonzero, in
  the sense of Lemma~\ref{asm-dependence-functional-witness-twisted} with $C=C_H=0$, is a twisted
  place of $R$; $R$ exists at every place; and the set of tuples of the row-side family of the bundle
  is finite.
\item In (2), suppose that the rows of the size data of the datum have the sizes of the size data,
  that the series idempotents of the datum are tied to its labels, and that a comparison part of the
  bundle is given, with its normalization of the realization. Then $R$, with these size data, is a
  realization in an automorphic setting at $F_R$
  (Definition~\ref{asm-dependence-functional-witness-setting}).
\end{enumerate}
\end{lemma}
\leanref{Development/AutomorphicArithmetic/ThetaCorrespondence/ThetaPairLifts/ThirdPlaceSeries/EvPSizeRowsDlog.lean}{GlobalTheory.Mem.LiteralAutomorphy.GridRow.eventually_sizeRows_and_logDeriv_le_alt}
\leanref{Development/AutomorphicArithmetic/ThetaCorrespondence/ThetaPairLifts/ThirdPlaceSeries/EvPSizeRows.lean}{GlobalTheory.Mem.AtSize.eventually_sizeRows_alt}
\leanref{Development/AutomorphicArithmetic/ThetaCorrespondence/ThetaPairLifts/ThirdPlaceSeries/BwitFrameTwists.lean}{GlobalTheory.RowEx.frame_eq_and_twistPlacesOffPrimeQZero_of_bWit}
\leanref{Development/AutomorphicArithmetic/ThetaCorrespondence/ThetaPairLifts/ThirdPlaceSeries/WeightOnePlacesFiniteOfWitness.lean}{GlobalTheory.Mem.Top.Bl.weight_one_places_finite_of_witness}
\leanref{Development/AutomorphicArithmetic/ThetaCorrespondence/ThetaPairLifts/ThirdPlaceSeries/RowTestSetting.lean}{GlobalTheory.PrimIn.rowTest_setting_alt}
\leanref{Development/AbelianVarietyTheory/Statements/HeckeThetaIntegrals/FormulaAt.lean}{GlobalTheory.RowClaim.HidaSetting}
\leanref{Development/AbelianVarietyTheory/Proofs/IwasawaDecomposition/BwitFrameTwists.lean}{GlobalTheory.RowEx.FrmW.frame_of_args}
\leanref{Development/AbelianVarietyTheory/Proofs/TwistedOperators/TwistedDataOffBaseLevel.lean}{GlobalTheory.RowEx.twisted_data_off_base_level}
\leanref{Development/AbelianVarietyTheory/Proofs/TwistedOperators/TwistedDataOffBaseLevel.lean}{GlobalTheory.RowEx.DatTw.unrAt_of_mem_set}
\leanref{Development/AbelianVarietyTheory/Proofs/IwasawaDecomposition/WeightOnePlacesAreTwistPlacesOfWitness.lean}{GlobalTheory.Mem.Top.Bl.finite_T_of_twArgs}
\leanref{Development/AbelianVarietyTheory/Definitions/AutomorphicMembership/CmQuadraticFormEqAlternatingForm.lean}{GlobalTheory.Mem.Top.TwDat.set}
\leanref{Development/AbelianVarietyTheory/Proofs/IwasawaDecomposition/LeastSplitAuxiliaryPrimeLtScaleBound.lean}{GlobalTheory.Mem.AtSize.leastSplitAuxiliaryPrime_lt_scaleBound}
\begin{proof}
(1) The size data of the datum at a base point are at the scale $(x,\varepsilon)$ of the context, and
a base point records the bound $|L'/L|\le x^{\varepsilon}$ for the logarithmic derivative at $1$ of its
size data. So the second conclusion holds at every base point.

For the first, let $\varepsilon$ and the constants of a context be given, and let $x_0$ be a
threshold beyond which Lemma~\ref{asm-dependence-functional-witness-numbers} holds for them. Let the
scale $x$ of the context be at least $\max(x_0,44^4,8^{4/\varepsilon})$, and let a base point of a good
core of $p$ be given, with datum $D$. Its size data are at the scale of the context, so
Lemma~\ref{asm-dependence-functional-witness-numbers} gives the parameter choices with
$C_{nW}=10$, the size bounds, the bounds on $n$, $m$ and $k$, $g_0,g_1\le\exp(9x^{\varepsilon})$ and
$\mathfrak c=\mathcal O_K$. The base point records that its context is good, and the conditions of a
good context make $q_0$ a level prime of the host different from $g_0$ and $g_1$, and
every level prime other than $g_0$ and $g_1$ less than $X'$; so $q_0<X'$. Since
$x\ge8^{4/\varepsilon}$, $x^{\varepsilon/4}\ge8$, and $x\ge44^4$. So
Lemma~\ref{asm-dependence-functional-witness-computation} applies, and the rows of the size data
have the sizes of the size data.

(2) The witness consists of twist arguments, built at the later stage from the twist data $T$, whose
package contains a realization of the datum together with its row-side family, and of a real input
whose package the witness identifies with the bundle's. So $R$ is the realization of the twist
arguments, and the row-side family of the bundle is theirs. The Hecke characters $\xi_t$ of the level
system of this realization are characters of the column choice, whose values form the set of the
statement: the trivial character, the characters $\delta,\delta'$ of the framing and the characters
$\delta_\ell\psi_w$ of $T$; and its columns are given by $f$. In a level system with base tuple,
$\xi_t$ is ramified at the places of $Q_t$. The characters of that set are unramified at the places
above $q_0$. So no place of $Q_t$ lies above $q_0$: the twisted places of $R$ lie off $q_0$. The
tuples of the twist arguments are the base tuple and the tuples of the places they use, so the set
of tuples of the row-side family is finite. By
Lemma~\ref{asm-dependence-functional-witness-twisted} with $C=C_H=0$, every place at which an index
of $R$ at $F_R$ is nonzero is a twisted place of $R$. The witness records that the realization of its
twist arguments, at their face, satisfies the conditions of a realization built at a stage, among
them that it exists at every place; the identification of the packages carries this to $R$.

(3) The conditions of Definition~\ref{asm-dependence-functional-witness-setting} hold as follows.
Conditions (1) and (2) are hypotheses. The indefinite realization of the bundle records the gap
conditions (3), the standard character (11), and the orthonormal basis of $X'$ spanning its lattice
in (12). The normalization of the realization in the comparison part gives (4), the guard (8), and
the identity adelic matrix of $X'$ in (12). The indefinite realization satisfies the weak
characterization at $F_R$, whose conditions on the tuples are (6), whose conditions on the group
record the group, (7), and which contains the fullness of the Weil data, (10). By (2), $R$ exists at
every place and the set of tuples is finite, (9); so
Lemma~\ref{asm-dependence-functional-witness-core} gives the core conditions, (5).
% Lean: (1) AtSize.eventually_sizeRows_alt: szRow_numAt (Lemma asm-dependence-functional-witness-numbers),
% the threshold max x0 (max (44^4) (8^(4/eps))), leastSplitAuxiliaryPrime_lt_scaleBound (beta.good: ChOK.small
% puts q0 in the level primes with q0 != g i, ChOK.lev bounds the other level primes by X'), and
% SetAtModel.sizeRows_of (Lemma asm-dependence-functional-witness-computation). (2) weight_one_places_finite_of_witness:
% zscale_f0_of_witness_alt and R3.realPins3_realTOf_alt (Lemma asm-dependence-functional-witness-twisted)
% with Pins3.w1 at (0, 0); ExAt is the field aux.ex (UnusedChoiceStatements) of R3.Holds, the field
% intersectionNumber_baseChange of the TwistedBundle w.y, carried along w.y.eq; Finite from finite_T_of_twArgs.
% (3) PrimIn.rowTest_setting_alt: size, eser from the hypotheses; host, std, stdX from B.Q; hF, guard
% (qGuard) and betaX (betaStd.2) from the normalization hp; mu_off, mu_used, mu_ne, mu_conj, mu_on,
% grp (base.grp.record), fullY, fullX from B.Q.charW; core from CoreOfWeak.core_of_weak; fin from (2).
% The hypothesis on the places of nonzero index is not used in (3).
\end{proof}

\begin{lemma}\label{asm-dependence-functional-full}
Let $C_F$ be real, let local Langlands maps be given, let $N$ be an extension, and let $p$ be the
parameters of the extension of $N$ by the grid property for $C_F$
(Definition~\ref{asm-dependence-parameters}). Eventually in the context, every bundle that
$\Lambda_p$ links to a base point of a good core of $p$ has the following property. For every
comparison part of the bundle whose point satisfies the compatibility conditions of
Definition~\ref{is-point}, with realization $R$, and every evaluation grid $G_0$ for $R$ with
constant $C_F$ at $b$, there is an evaluation grid $G$ for $R$ with constant $C_F$ at $b$, with the
base place, the place above it and the characters $\chi_1,\dots,\chi_m$ of $G_0$, which carries a
full evaluation choice (Definition~\ref{asm-stage-construction}(5)).
\end{lemma}
\leanref{Development/AutomorphicArithmetic/ThetaCorrespondence/ThetaPairLifts/ThirdPlaceSeries/EvPSizeRowsDlog.lean}{GlobalTheory.Mem.LiteralAutomorphy.GridRow.rowE_full_grid_alt}
\leanref{Development/AbelianVarietyTheory/Statements/PlaceShiftStatements/NormSmall.lean}{GlobalTheory.SqTop.Rec}
\begin{proof}
By Lemma~\ref{asm-dependence-functional-witness}(1), eventually in the context, at every base point
of a good core of $p$ the rows of the size data have the sizes of the size data and the logarithmic
derivative at $1$ has absolute value at most $x^{\varepsilon}$. Let $K$ be such a context, with scale
$(x,\varepsilon)$, let a bundle be linked to a base point of a good core of $p$ in $K$, and let the
link data give the stage $(K,B,f,T)$, its later stage and the witness. The size data of the bundle
are at the scale of $K$. The stage satisfies the extension of $N$ by the grid property for $C_F$, so
it has the grid property for $C_F$ (Definition~\ref{asm-stage-construction}(5)).

Let a comparison part be given whose point satisfies the compatibility conditions, and let $R$ be
the realization of the point, which is the indefinite realization of the bundle. By
Lemma~\ref{asm-dependence-functional-witness}(2), the columns of $R$ are given by $f$, its twisted
places lie off $q_0$ and the characters $\xi_t$ of its level system lie in the set described there,
of the trivial character, $\delta$, $\delta'$ and the characters $\delta_\ell\psi_w$, which is the set
of the grid property. The compatibility
conditions include that the series idempotents of the datum are tied to its labels, so
Lemma~\ref{asm-dependence-functional-witness}(3) shows that $R$ is a realization in an automorphic
setting. The automorphic setting gives the host conditions and that $F$ is the maximal real subfield
of $E$; with the sizes of the rows and the bound on the logarithmic derivative, the size data satisfy
the recursion conditions. The grid property for $C_F$, applied to the datum of the later stage, to
$R$ and to $G_0$, gives $G$.
\end{proof}

\begin{lemma}\label{asm-dependence-functional-grid}
Let $C_F$ be real, let local Langlands maps be given, let $N$ be an extension stronger than
$N_{\mathrm{list}}(C_F)$, and let $p$ be the parameters of the extension of $N$ by the grid property
for $C_F$ (Definition~\ref{asm-dependence-parameters}). Then every bundle that $\Lambda_p$ links to
a base point of a good core of $p$ eventually has the following property: for every comparison part
of the bundle whose point satisfies the compatibility conditions of Definition~\ref{is-point}, the
point has a fitted evaluation grid with constant $C_F$
(Definition~\ref{asm-dependence-functional-fitted}(5)), as in part (3) of
Theorem~\ref{asm-dependence-inputs}.
\end{lemma}
\leanref{Development/AutomorphicArithmetic/ThetaCorrespondence/ThetaPairLifts/ThirdPlaceSeries/RowEHasGridFldStatement.lean}{GlobalTheory.T.grid_field_has_lattices_alt}
\leanref{Development/AutomorphicArithmetic/ThetaCorrespondence/ThetaPairLifts/ThirdPlaceSeries/RowEHasGridFld.lean}{GlobalTheory.Mem.LiteralAutomorphy.GridRow.rowE_hasGrid_fld_alt}
\leanref{Development/AbelianVarietyTheory/Definitions/GlobalTwisting/ValidForConstructedPointClassesEventuallyWorking.lean}{GlobalTheory.Fc.RowE}
\leanref{Development/AbelianVarietyTheory/Definitions/Divisors/TwistedPlaceChoiceLiesOverPair.lean}{GlobalTheory.BigTop.Fun.HasGrid}
\leanref{Development/AutomorphicArithmetic/ThetaCorrespondence/ThetaPairLifts/ThirdPlaceSeries/RowEHasGridFld.lean}{GlobalTheory.Mem.LiteralAutomorphy.GridRow.hasGrid_of_full_field_lamTie}
\leanref{Development/AbelianVarietyTheory/Proofs/IwasawaDecomposition/GridFieldHasLattices.lean}{GlobalTheory.Mem.LiteralAutomorphy.GridRow.gridFit_of_parts}
\leanref{Development/AbelianVarietyTheory/Proofs/IwasawaDecomposition/GridFieldHasLattices.lean}{GlobalTheory.Mem.LiteralAutomorphy.GridRow.fieldHas_of_same}
\leanref{Development/AbelianVarietyTheory/Proofs/IwasawaDecomposition/GridFieldHasLattices.lean}{GlobalTheory.Mem.LiteralAutomorphy.GridRow.lamOf_of_same}
\begin{proof}
By Lemma~\ref{asm-dependence-functional-full}, eventually in the context every bundle that
$\Lambda_p$ links to a base point of a good core of $p$ has the property of that lemma, and by
Lemmas~\ref{asm-dependence-functional-twists} and~\ref{asm-dependence-functional-pinned}, which
hold in every context, it also has the properties of those lemmas. Let such a bundle be given, with
datum $D$, and a comparison part whose point $q$ satisfies the compatibility conditions, with
realization $R$.

Lemma~\ref{asm-dependence-functional-pinned} gives an evaluation grid $G_0$ for $R$ with constant
$C_F$ at $b$, valued in $L_0$ and with scalars tied to the twisted places.
Lemma~\ref{asm-dependence-functional-full} gives an evaluation grid $G$ for $R$ with constant $C_F$ at
$b$, with the base place $v_0$, the place $w(v_0)$ and the characters $\chi_1,\dots,\chi_m$ of $G_0$,
and a full evaluation choice for it, with a Weil datum $W$ over $R$. The set $S_0$ and the depths of
a grid are determined by its base place, so $G$ and $G_0$ have the same set $S_0$ and the same
depths; with the same characters, $G$ is valued in $L_0$. The scalars of a grid are determined by
$R$, $q_0$, the place $w(v_0)$, the first constant and the characters, so $G$ and $G_0$ have the same
scalars, and those of $G$ are tied to the twisted places.

By Lemma~\ref{asm-dependence-functional-twists}, the twist choices of $D$ lie over their pairs and
satisfy the twist conditions. The compatibility conditions of $q$ include that the characters of the
columns of $R$ take values in $L_0$, that the character choice of $q$ is compatible with the twist
choices, and that the series idempotents of $D$ are tied to its labels. So $G$ is fitted to $q$, and
the choices of the full evaluation choice other than $G$ and $W$ are admissible remaining choices of
the slope test for $G$ and $W$. Hence $q$ has a fitted evaluation grid with constant $C_F$.
\end{proof}

\begin{lemma}\label{asm-dependence-functional-geometric}
Assume the existence of ray class fields (Theorem~\ref{lit-bundle-unitary}(5)). Let $C_F$ be real,
let local Langlands maps be given, let $N$ be an extension stronger than $N_{\mathrm{list}}(C_F)$,
and let $p$ be the parameters of the extension of $N$ by the grid property for $C_F$
(Definition~\ref{asm-dependence-parameters}). If the division rows of open parameter statement (b)
and statement (d) (Definition~\ref{asm-open-parameter-statements}) are linked at $p$, then so is the
following property $G$ of a class and a selection: at every member, the selected bundle has a
comparison part whose point has the normalized embedding of its setting and satisfies the class-sum
pin, the compatibility conditions and the further conditions of Definition~\ref{is-point}, and whose
point has an evaluation grid with constant $C_F$ at $b$ fitted to it, with a Weil datum over its
realization and the remaining choices of the slope test; and the datum of the member with its size
data satisfies the coefficient field bound (Situation~\ref{cv-point}).
\end{lemma}
\leanref{Development/AutomorphicArithmetic/ThetaCorrespondence/ThetaPairLifts/ThirdPlaceSeries/RowGeoAtGridAStatement.lean}{GlobalTheory.T.geometric_condition_large_grid_alt}
\leanref{Development/AutomorphicArithmetic/ThetaCorrespondence/ThetaPairLifts/ThirdPlaceSeries/RowGeoAtGridA.lean}{GlobalTheory.Mem.LiteralAutomorphy.rowGeo_at_gridA_alt}
\leanref{Development/AutomorphicArithmetic/ThetaCorrespondence/ThetaPairLifts/ThirdPlaceSeries/RowGood.lean}{GlobalTheory.Fc.RL.RowGeometricCondition}
\leanref{Development/AutomorphicArithmetic/ThetaCorrespondence/ThetaPairLifts/ThirdPlaceSeries/RowGood.lean}{GlobalTheory.Fc.RL.LinkedAt.of_row}
\begin{proof}
The extension of $N$ by the grid property is stronger than $N$, hence than $N_{\mathrm{list}}(C_F)$,
which requires the law extension. So Lemma~\ref{asm-dependence-covolume-ties} applies at $p$: with
the division rows of (b), the property $T$ of that lemma is linked at $p$.

By Lemma~\ref{asm-dependence-functional-grid}, the condition $E$ of that lemma is satisfied
eventually by every bundle that $\Lambda_p$ links to a base point of the slot $Q_p$ of the good
cores of $p$. Take it as the condition on bundles: for every class and every selection whose bundles
are built from $(Q_p,\Lambda_p)$ and satisfy $E$, at every member, every comparison part of the
selected bundle whose point satisfies the compatibility conditions has a fitted evaluation grid as
in $G$. This property $H$ holds along $(Q_p,\Lambda_p)$, so it is linked at $p$, with the slot of
all cores for $R$ (Definition~\ref{asm-dependence-parameters}(2)).

By Lemma~\ref{asm-dependence-linked}(2), the conjunction of $T$, $H$ and statement (d) is linked at
$p$. It implies $G$: at a member, (d) gives a comparison part of the selected bundle, $T$ gives the
properties of its point and of the member, among them the compatibility conditions, and then $H$
gives the evaluation grid. So $G$ is linked at $p$ by Lemma~\ref{asm-dependence-linked}(1).
\end{proof}

\begin{lemma}\label{asm-dependence-functional-bound}
Assume the conclusion of part (3) of Theorem~\ref{asm-dependence-inputs}, with its constant $C'$.
Let $C_F\ge C'$. Let $\mathcal A$ be a class of families each of which carries its duty data
(Definition~\ref{asm-duty}), and choose selected parts as in
Lemma~\ref{asm-dependence-covolume-points}, with the working subclass $\mathcal W(\mathcal A)$.
Suppose that eventually along $\mathcal W(\mathcal A)$, eventually along each family, the point of
the selected part satisfies the point conditions (Situation~\ref{cv-point}) and has an evaluation
grid with constant $C_F$ at $b$ fitted to it, with a Weil datum over its realization and the
remaining choices of the slope test. Then eventually along $\mathcal W(\mathcal A)$, every family
that satisfies the division
bound at twisted places (Definition~\ref{is-twisted-division-bound}) has a constant $C_0\ge0$ such
that for every $C\ge C_0$, eventually along the family, the rows of the realization of the selected
part satisfy the functional bound (Definition~\ref{is-functional-bound}) for $N(C,C_H)$.
\end{lemma}
\leanref{Development/AutomorphicArithmetic/ThetaCorrespondence/ThetaPairLifts/ThirdPlaceSeries/EvWCFunOfFun0Fit.lean}{GlobalTheory.Fc.Ind.evWC_fun_of_fun0Fit_alt}
\leanref{Development/AutomorphicArithmetic/ThetaCorrespondence/ThetaPairLifts/ThirdPlaceSeries/SummandFunctionalBound.lean}{GlobalTheory.Fc.Ind.summandFunctionalBound}
\leanref{Development/AutomorphicArithmetic/ThetaCorrespondence/ThetaPairLifts/ThirdPlaceSeries/LeafInfinitePlaceBound.lean}{GlobalTheory.BigTop.Fun.ConstantFunctionFits}
\begin{proof}
A family of $\mathcal W(\mathcal A)$ that carries its duty data is a standard family
(Section~\ref{section-is-points}): the duty data give the first five family hypotheses, and the
conditions defining the working subclass give the others. So part (3) of
Theorem~\ref{asm-dependence-inputs} gives statement (i) below eventually along
$\mathcal W(\mathcal A)$, with $C_0=0$, and the hypotheses are statements (ii) and (iii) eventually
along $\mathcal W(\mathcal A)$.
\begin{enumerate}
\item[(i)] If the family satisfies the division bound at twisted places, it has $x_0$ such that, at
  every scale $x\ge x_0$ and every point of the family at $x$ that satisfies the point conditions
  and has an evaluation grid as above, the rows satisfy the functional bound for $N(C,C_H)$ for
  every $C\ge0$.
\item[(ii)] There is $x_0'$ beyond which the point of the selected part satisfies the point
  conditions.
\item[(iii)] There is $x_0''$ beyond which the point of the selected part has an evaluation grid as
  above.
\end{enumerate}
Each of the three comes with its own $\varepsilon_1$ and largeness thresholds. Take the least of the
three numbers $\varepsilon_1$ and, for each exponent, a threshold that is at least each of the
three; then (i), (ii) and (iii) hold together for every family of $\mathcal W(\mathcal A)$ that
is large for it. For such a family with the division bound and $x\ge\max(x_0,x_0',x_0'')$, the point
of the selected part at every member at $x$ satisfies the point conditions and has an evaluation
grid as above, so its rows satisfy the functional bound for $N(C,C_H)$ for every $C\ge0$. The rows
of the realization of the selected part are those of this point.
\end{proof}

\subsection{Test choices}\label{subsection-asm-test-choices}

\begin{definition}\label{asm-dependence-test-assembly-statements}
Let $D$ be a definite-side datum with $m$ lifts and coefficient field $L$, with size data at the
scale $(x,\varepsilon)$, and put $b=e^{x^{\varepsilon}}$. Let $R$ be a realization of $D$ for a
family of definite forms, with base tuple $\mu_0$, which with these size data is a realization in
an automorphic setting at a face (Definition~\ref{asm-dependence-functional-witness-setting}), and
let $E$, $\mathrm{Th}_t$ and $j_t$ be its top space, its theta lifts and its maps from the row forms
(Definition~\ref{is-realisation}). Let $C_F$ be real.
\begin{enumerate}
\item \emph{Test choices} for $C_F$ consist of an evaluation grid $G$ for $R$ with constant $C_F$ at
  $b$, a Weil datum $W$ over $R$ for $G$, a row volume, a real vector $\rho$ and a linear functional
  $\xi$ on the coefficients (Definition~\ref{asm-grid-evaluation}), a smooth vector $w$ of the model
  of $U(Y_0)$ in $R$, and the further choices of Definition~\ref{asm-dependence-grid-agree}(2),
  among them a sum choice extending a grid choice, an enumeration $a=1,\dots,n$ of the test indices
  of $G$, which gives for each $a$ a coset datum $\Phi_a$ and, with class representatives, a column
  $\Psi_a$, and the \emph{square constants} $c(i,j,t)\in\mathbf C$ for $1\le i,j\le m$ and tuples $t$.
  The \emph{base conditions} are conditions on these choices other than $\rho$ and $\xi$; they include
  that $w\ne0$, that $W$ is a Weil datum over $R$ and the measures form a row volume
  (Definition~\ref{asm-grid-evaluation}(2) and (3)), and that $c(i,j,t)$ is the constant of the square
  $(i,j)$ at $t$.
\item For test choices and $u\in E$, let $u_w$ be the function on $U(Y_0)(\mathbf A)$ obtained by
  applying the theta multiplicity of $u$ to $w$, as in Definition~\ref{asm-grid-evaluation}(6). The
  \emph{test functionals} are
  $\psi_a(u)=B^+\bigl(\varphi_{\xi,\Psi_a},\Theta_{\rho,\Phi_a}(u_w)\bigr)$, and
  $c_{\mathrm{tot}}(t)=\prod_{i,j}c(i,j,t)$. The sum choice gives \emph{sum values}
  $v(t,y,a)\in\mathbf C$ for tuples $t$, row forms $y$ at $t$ and test indices $a$. The
  \emph{evaluation conditions} are the base conditions and the \emph{real condition}: for every
  $u\in E$ with $\psi_a(u)=0$ for every $a$, and for every real vector $\rho'$ and linear functional
  $\xi'$ on the coefficients, the test functionals of the test choices with $\rho'$ and $\xi'$ in place
  of $\rho$ and $\xi$ vanish at $u$.
\item Each of the following seven \emph{statements on test choices} asserts that there is
  $C_{F,0}$ such that for every $\varepsilon$ and every $C_F\ge C_{F,0}$ there is $x_0$ with the
  following property: for every $D$, its size data and $R$ as above with $x\ge x_0$, and all test
  choices for $C_F$ that satisfy the evaluation conditions,
  \begin{enumerate}
  \item[(a)] for every $a$ and $u$, the iterated theta lift $\Theta_{\rho,\Phi_a}(u_w)$ and the Bessel
    period $\psi_a(u)$ converge;
  \item[(b)] if $\psi_a(u)=0$ for every $a$, then $u=0$;
  \item[(c)] for all $i$, $j$ and $t$, the square $(i,j)$ at $t$ is nonzero and $c(i,j,t)\ne0$;
  \item[(d)] for every $t$ and $a$, the sum choice satisfies the sum conditions at $t$ and $\Phi_a$;
  \item[(e)] if the field conditions hold, there are $\kappa\in\mathbf C$ and roots of unity
    $\zeta_t\in L$ with $\zeta_t\,c_{\mathrm{tot}}(t)\,\psi_a(\mathrm{Th}_tj_ty)=\kappa\,v(t,y,a)$ for
    every tuple $t$, every row form $y$ at $t$ in the good part and every $a$;
  \item[(f)] if the field conditions hold, every sum value $v(t,y,a)$ lies in $L$;
  \item[(g)] if the field conditions hold, $c(i,j,t)/c(i,j,\mu_0)$ lies in $L$ for all $i$, $j$ and $t$.
  \end{enumerate}
\item The \emph{number field statement} asserts that there is $C_{F,0}$ such that for every
  $\varepsilon$ and every $C_F\ge C_{F,0}$ there is $x_0$ with the following property. For every
  scale $(x,\varepsilon)$ with $x\ge x_0$, all base data $B$ with a framing $f$ and every finite set
  $\Xi$ of Hecke characters of the host field of $B$, there is a finite set $S$ of algebraic numbers
  such that every realization $R$ of the datum of a later stage over $(B,f)$ whose coefficient field
  contains $S$, which with size data at the scale $(x,\varepsilon)$ is in an automorphic setting at a
  face, whose columns are given by $f$, whose twisted places lie off $q_0$ and whose characters
  $\xi_t$ lie in $\Xi$, has test choices for $C_F$ that satisfy the evaluation conditions and the field
  conditions.
\end{enumerate}
\end{definition}
\leanref{Development/AbelianVarietyTheory/Definitions/TwistedInnerProducts/ExistsMatchOver.lean}{GlobalTheory.RowTest.Choice}
\leanref{Development/AbelianVarietyTheory/Definitions/TwistedInnerProducts/ExistsMatchOver.lean}{GlobalTheory.RowTest.Choice.Is}
\leanref{Development/AbelianVarietyTheory/Definitions/TwistedInnerProducts/MemHalfCosetSelf.lean}{GlobalTheory.RowTest.Choice.FieldOK}
\leanref{Development/AbelianVarietyTheory/Definitions/TwistedInnerProducts/ExistsMatchOver.lean}{GlobalTheory.RowTest.psiK}
\leanref{Development/AbelianVarietyTheory/Definitions/TwistedInnerProducts/ExistsMatchOver.lean}{GlobalTheory.RowTest.ConvK}
\leanref{Development/AbelianVarietyTheory/Statements/HeckeThetaIntegrals/FormulaAt.lean}{GlobalTheory.RowClaim.Big}
\leanref{Development/AbelianVarietyTheory/Statements/HeckeThetaIntegrals/FormulaAt.lean}{GlobalTheory.RowClaim.LinK}
\leanref{Development/AbelianVarietyTheory/Statements/HeckeThetaIntegrals/FormulaAt.lean}{GlobalTheory.RowClaim.InjK}
\leanref{Development/AbelianVarietyTheory/Statements/HeckeThetaIntegrals/FormulaAt.lean}{GlobalTheory.RowClaim.ConstK}
\leanref{Development/AbelianVarietyTheory/Statements/HeckeThetaIntegrals/FormulaAt.lean}{GlobalTheory.RowClaim.SumOKK}
\leanref{Development/AbelianVarietyTheory/Statements/HeckeThetaIntegrals/FormulaAt.lean}{GlobalTheory.RowClaim.Formula2K}
\leanref{Development/AbelianVarietyTheory/Statements/HeckeThetaIntegrals/FormulaAt.lean}{GlobalTheory.RowClaim.SumRatK}
\leanref{Development/AbelianVarietyTheory/Statements/HeckeThetaIntegrals/FormulaAt.lean}{GlobalTheory.RowClaim.RatioK}
\leanref{Development/AbelianVarietyTheory/Definitions/TwistedInnerProducts/TwistPlacesOffPrimeQZero.lean}{GlobalTheory.RowEx.ExField}
\leanref{Development/AbelianVarietyTheory/Definitions/TwistedInnerProducts/TwistPlacesOffPrimeQZero.lean}{GlobalTheory.RowEx.FieldFor}

\begin{lemma}\label{asm-dependence-test-assembly-first-layer-rows-pieces}
Assume the results of Theorem~\ref{lit-bundle-theta} and the theta hypotheses
(Definition~\ref{th-theta-hypotheses}).
\begin{enumerate}
\item The theta integrals of the twisted Weil representations converge.
\item Statement (c) on test choices holds (Definition~\ref{asm-dependence-test-assembly-statements}(3)).
\item Statement (d) on test choices holds.
\item Statement (e) on test choices holds.
\item If statement (d) on test choices holds, so does statement (f).
\end{enumerate}
\end{lemma}
\leanref{Development/AutomorphicArithmetic/ThetaCorrespondence/ThetaPairLifts/LocalTypeDictionary/TwistedWeilRepresentationConvergence.lean}{GlobalTheory.WeilConv.twisted_weil_representation_convergence}
\leanref{Development/AutomorphicArithmetic/ThetaCorrespondence/ThetaPairLifts/LocalTypeDictionary/ThetaLiftRowPropertiesConj.lean}{GlobalTheory.RowClaim.Join.thetaLift_row_properties_conj}
\leanref{Development/AbelianVarietyTheory/Proofs/TwistedOperators/RowTestLiftFormulaPrincipalLevelStatement.lean}{GlobalTheory.T.rowTestLiftFormula_principalLevel}
\leanref{Development/AbelianVarietyTheory/Proofs/TwistedOperators/RowTestRationalSumsOfFiniteSumsMeaningful.lean}{GlobalTheory.RowTest.sumRatK_of_sumOKK}
\begin{proof}
Part (2) rests on (1), the matching of the twisted local factors
(Lemma~\ref{th-first-lift-factor}(2)) and the agreement of the twisted square theta integrals
(Lemma~\ref{th-first-lift-inputs}(3)). Part (4) is derived from (2), (3) and four further statements:
the period identity of the projections, the invariance of the twisted lifts under the flip (which
uses (1), (2), the theta hypotheses and the cuspidality of the theta lifts of pairs), the unfolding of
the twisted theta integral (which uses (3) and that cuspidality) and the assembly of the jets (which
uses (3)).
\openstep{The convergence of the theta integrals of the twisted Weil representations from Siegel
sets, the nonvanishing of the squares and of their constants, the sum conditions, the period
identity, the flip invariance, the unfolding and the jets with the derivation of the formula (e) from
them, and the rationality of the sum values from the sum conditions, are not yet written out.}
% Lean: (1) WeilConv.twisted_weil_representation_convergence (weilConvergence_of_siegelSet; about 5,300
% lines within five dependency levels by the count used here). (2) ConstNZ.constK_alt P K LM s1 s2 s3
% (about 5,300). (3) T.sumOKK_of_theta_local_inputs P (about 3,000). (4)
% T.rowTestLiftFormula_principalLevel (about 3,800) from T.twisted_period_identity P (StepProj, about
% 3,000), T.twisted_lift_flip_invariant P cusp s1 K c (StepFlip, about 5,700),
% T.twisted_theta_integral_unfolds P cusp sum (StepUnfold, about 4,700), T.twisted_jets_assemble P sum
% (StepJets, about 3,000), c and sum, where cusp = T.twisted_thetaLift_cuspidal P is
% thetaPairLift_cuspidal_of_hypotheses, open in Lemma -convergence-lift. (5)
% T.rowTestRationalSums_of_finiteSumsMeaningful = RowTest.sumRatK_of_sumOKK (SumRat.vK_mem: sAA_mem,
% coset expansions, class sums, unfoldedFlipTensor_values_mem_coeffField,
% topRow_normalisedKernelChain_preserves_coeffs, lastColumn_normalizedKernelChain_coeffField,
% branchingPairing_mem_coefficientField; about 5,300). The counts overlap.
\end{proof}

\begin{lemma}\label{asm-dependence-test-assembly-first-layer-rows}
Assume the results of Theorem~\ref{lit-bundle-theta} and the theta hypotheses
(Definition~\ref{th-theta-hypotheses}). Then statements (a), (c), (d), (e) and (f) on test choices
hold (Definition~\ref{asm-dependence-test-assembly-statements}(3)).
\end{lemma}
\leanref{Development/AutomorphicArithmetic/ThetaCorrespondence/ThetaPairLifts/LocalTypeDictionary/ThetaLiftRowPropertiesConj.lean}{GlobalTheory.RowClaim.Join.thetaLift_row_properties_conj}
\leanref{Development/AutomorphicArithmetic/ThetaCorrespondence/ThetaPairLifts/LocalTypeDictionary/TwistedWeilRepresentationConvergence.lean}{GlobalTheory.WeilConv.twisted_weil_representation_convergence}
\leanref{Development/AutomorphicArithmetic/ThetaCorrespondence/ThetaPairLifts/LocalTypeDictionary/TwistedSquareLocalMatching.lean}{GlobalTheory.SquareLM.twisted_square_local_matching}
\leanref{Development/AutomorphicArithmetic/ThetaCorrespondence/ThetaPairLifts/LocalTypeDictionary/TwistedSquareThetaIntegralsEq.lean}{GlobalTheory.SquareAgree.twisted_theta_integrals_agree}
\begin{proof}
Three statements on the twisted Weil representations are used. The matching of the twisted local
factors is Lemma~\ref{th-first-lift-factor}(2), and the agreement of the twisted square theta
integrals is Lemma~\ref{th-first-lift-inputs}(3); both rest on Howe duality at the real place and the
theta hypotheses. The third is the convergence of the theta integrals of the twisted Weil
representations, Lemma~\ref{asm-dependence-test-assembly-first-layer-rows-pieces}(1).

Statement (a): test choices that satisfy the evaluation conditions satisfy the base conditions, so
(a) is Lemma~\ref{asm-dependence-test-assembly-first-layer-convergence}. Statements (c), (d) and (e) are
parts (2), (3) and (4) of Lemma~\ref{asm-dependence-test-assembly-first-layer-rows-pieces}, and (f)
follows from (d) by part (5).
% Lean: LinK (a) is T.twisted_characters_linearly_independent P cusp =
% RowTest.twist_linear_independence_test, with C_F0 = 0 and x0 = 0, from
% T.twisted_theta_lift_automorphic (= twisted_base_change_lift_exists at ch.toIsBase, Lemma
% -convergence-lift) and besselPeriodCorank_converges_column; the same argument as twisted_linear_bound.
% s1 = WeilConv.twisted_weil_representation_convergence (weilConvergence_of_siegelSet; about
% 1,600 lines at depth 5), s2 = SquareLM.twisted_square_local_matching, s3 =
% SquareAgree.twisted_theta_integrals_agree; RowClaim.Join.thetaLift_row_properties_conj P K LM s1 s2
% s3 (ConstNZ.constK_alt and the rest) gives SumOKK (d), ConstK (c), LinK (a), Formula2K (e) and
% SumRatK (f); about 23,700 lines of Lean proofs at depth 5 in all, counted from
% thetaLift_row_properties_conj.
\end{proof}

\begin{lemma}\label{asm-dependence-test-assembly-first-layer-convergence-lift-grid}
Assume the results of Theorem~\ref{lit-bundle-theta}. Let $D$, its size data and $R$ be as in
Definition~\ref{asm-dependence-test-assembly-statements}, let $C_F$ be real, and let test choices
for $C_F$ satisfy the base conditions, with the evaluation grid $V_{i,j}$, the Weil datum and the row
volume. Then there is a grid of models of discrete series over $R$, satisfying the conditions of
such a grid, that \emph{matches} the test choices: for $0\le i\le m$ its space at $(i,0)$ is
$V_{i,0}$ and its measure and fundamental domain there are those of the row volume; and for every
real vector $\rho$, every coset datum $\Phi$ and every $i<m$ there are a finite vector $f$ of the pair
$(V_{i+1,0},V_{i,0})$ and a polynomial $P$ in the creation operators such that the theta kernel of
the Weil datum of the grid of models at the Fock vector $f\otimes P$ is the theta kernel of the Weil
datum of the test choices at the vector $f_{i,\Phi_i}\otimes\rho_i$.
\end{lemma}
\leanref{Development/AutomorphicArithmetic/ThetaCorrespondence/ThetaPairLifts/LocalTypeDictionary/ExistsRowLiftMatching.lean}{GlobalTheory.T.twist_row_lift_exists}
\leanref{Development/AutomorphicArithmetic/ThetaCorrespondence/ThetaPairLifts/LocalTypeDictionary/ExistsRowLiftMatching.lean}{GlobalTheory.RowLift.exists_row_lift_matching}
\leanref{Development/AbelianVarietyTheory/Proofs/TwistedOperators/IsAutFunTransportS.lean}{GlobalTheory.RowLift.RowMatchS}
\begin{proof}
The grid of models is built from the evaluation grid with the results of
Theorem~\ref{lit-bundle-theta}, its measures at the spaces $V_{i,0}$ are replaced by those of the row
volume, and the kernels are matched.
\openstep{The construction of a grid of models of discrete series on the spaces of the evaluation
grid from the existence and covering of discrete series and the finiteness of Harish-Chandra
modules, the replacement of its measures by those of the row volume, and the matching of its theta
kernels with those of the Weil datum of the test choices are not yet written out.}
% Lean: T.twist_row_lift_exists = RowLift.exists_row_lift_matching hP: GridJoin.exists_gridPre
% (discreteSeriesUpqExists, discreteSeriesCover, harishChandraModules, hS.host, hS.size, hS.core, cGap = 1,
% ch.W, hb.weil), the frames from hS.core.lev.betaY and the vertices, ColMeas.measCol_of_choice for the
% measures, GridJoin.withMeas, and RowKer.kernel_row for the kernels. About 9,600 lines of Lean proofs
% at depth 5 by the count used here. The Lean's grid of models is a GridPre R with G'.Is; matching is
% RowMatchS ch G' (fields V, nu, fd, ker).
\end{proof}

\begin{lemma}\label{asm-dependence-test-assembly-first-layer-convergence-lift-cuspidal}
Assume the results of Theorem~\ref{lit-bundle-theta}. Let $D$ be a definite-side datum with the gap
sizes (Definition~\ref{asm-dependence-rows-models}(5)), let $R$ be a realization of $D$ that satisfies
the core conditions on its level characters, whose archimedean family is one and whose measures on
the groups $U(Y_j)$ and $U(X')$ are Haar measures with fundamental domains, and let a grid of models
of discrete series over $R$ satisfy the conditions of such a grid. Then along the grid the theta lifts
of pairs carry cuspidal spaces to cuspidal spaces: for every pair $(A,B)$ of adjacent spaces of the
grid in the direction of its theta lifts, every finite vector $f$ of the pair and every $F$ in the
cuspidal space of the model of $A$, there is $F'$ in the cuspidal space of the model of $B$ such that
for every polynomial $P$ in the creation operators, every vector $v$ of the model of $A$ and every
$h\in U(B)(\mathbf A)$ the function $g\mapsto\theta(h,g;f\otimes P)F(v)(g)$ is integrable on the
fundamental domain of $U(A)$ and its integral is $F'(I(P,v))(h)$. Here $\theta(h,g;\cdot)$ is the theta
kernel of the Weil datum of the pair at $(h,g)$, its complex conjugate for the pairs of core spaces
lifted upward, and $I$ is the iterated theta map of the models of the pair.
\end{lemma}
\leanref{Development/AutomorphicArithmetic/ThetaCorrespondence/ThetaPairLifts/LocalTypeDictionary/ThetaPairLiftCuspidalOfHypotheses.lean}{GlobalTheory.PairLift.thetaPairLift_cuspidal_of_hypotheses}
\leanref{Development/AbelianVarietyTheory/Proofs/TwistedOperators/LiftsCuspidalOfFunctoriality.lean}{GlobalTheory.PairLift.liftsCuspidal_of_functoriality}
\leanref{Development/AbelianVarietyTheory/BlockTheory/TwistedDiscreteSeries/CuspAlongSmoothRows.lean}{GlobalTheory.IndefiniteRowSideFullHecke.CuspAlongSmoothLifts.LiftsCuspidal}
\leanref{Development/AbelianVarietyTheory/BlockTheory/TwistedDiscreteSeries/CuspAlongSmoothRows.lean}{GlobalTheory.IndefiniteRowSideFullHecke.CuspAlongSmoothLifts.CuspAlong}
\begin{proof}
The results used are the covering of discrete series, the finiteness of Harish-Chandra modules,
rapid decay of cusp forms and the Casselman-Wallach theorem, all among those of
Theorem~\ref{lit-bundle-theta}.
\openstep{The proof is not yet written out: real discrete series data for the theta
correspondence at the archimedean places, the analytic properties of the theta kernels of the pairs
built from them, and the deduction that the theta lift of a cuspidal multiplicity map converges and is
again a cuspidal multiplicity map.}
% Lean: PairLift.thetaPairLift_cuspidal_of_hypotheses = liftsCuspidal_of_functoriality (discreteSeriesCover,
% harishChandraModules, cuspRapidDecay, casselmanWallach): T.exists_real_discreteSeries_data_theta_properties,
% T.thetaPairAnalyticHypotheses_of_discreteSeries (bridge, vacChar, lie, cpt, dense, laws, const, ext) and
% liftsCuspidal. The pairs are the five fields of CuspAlong: X for (V_{i,j}, V_{i+1,j}), th for
% (V_{i,j+1}, V_{i,j}), up and dn for the core spaces (P_i, P_{i+1}) in both directions, the first with the
% complex conjugate kernel, and col for (X_j, X_{j+1}); the statement is LiftsMult for each. Its
% hypotheses are Sizes D (the gap sizes) and RowBase R (R.toLevCh.Is, IsInfFam R.theta, models_nuY,
% models_nuY', models_nuX). About 7,900 lines of Lean proofs at depth 5 by the count used here.
\end{proof}

\begin{lemma}\label{asm-dependence-test-assembly-first-layer-convergence-lift}
Assume the results of Theorem~\ref{lit-bundle-theta}. Let $D$, its size data and $R$ be as in
Definition~\ref{asm-dependence-test-assembly-statements}, let $C_F$ be real, and let test choices
for $C_F$ satisfy the base conditions. Then for every test index $a$ and every $u\in E$ the iterated
theta lift $\Theta_{\rho,\Phi_a}(u_w)$ converges, and it is a continuous function on
$U(V_{m,0})(\mathbf A)$, invariant on the left under $U(V_{m,0})(F)$ and on the right under an open
compact subgroup of the finite adelic group.
\end{lemma}
\leanref{Development/AutomorphicArithmetic/ThetaCorrespondence/ThetaPairLifts/LocalTypeDictionary/TwistedBaseChangeLiftExists.lean}{GlobalTheory.T.twisted_base_change_lift_exists}
\leanref{Development/AutomorphicArithmetic/ThetaCorrespondence/ThetaPairLifts/LocalTypeDictionary/TwistedBaseChangeLiftExists.lean}{GlobalTheory.RowLift.twisted_lift_from_base_field}
\leanref{Development/AutomorphicArithmetic/ThetaCorrespondence/ThetaPairLifts/LocalTypeDictionary/ExistsRowLiftMatching.lean}{GlobalTheory.T.twist_row_lift_exists}
\leanref{Development/AbelianVarietyTheory/Proofs/TwistedOperators/IsAutFunOfMemCusp.lean}{GlobalTheory.RowLift.row_induction}
\leanref{Development/AbelianVarietyTheory/Proofs/TwistedOperators/IsAutFunOfMemCusp.lean}{GlobalTheory.RowLift.isAutFun_of_mem_cusp}
\begin{proof}
By Lemma~\ref{asm-dependence-test-assembly-first-layer-convergence-lift-grid} there is a grid of
models of discrete series over $R$ that matches the test choices: its spaces and measures at the
spaces $V_{i,0}$ are those of the grid and the row volume of the test choices, and the theta kernel of
the vector $f_{i,\Phi_{a,i}}\otimes\rho_i$ is a kernel of its models. The rows of the size data have their
sizes, so $D$ has the gap sizes (Lemma~\ref{asm-dependence-rows-sizes}), and $R$ satisfies the core
conditions. So Lemma~\ref{asm-dependence-test-assembly-first-layer-convergence-lift-cuspidal} applies:
along this grid the theta lifts of pairs carry the cuspidal spaces to cuspidal spaces, and for each
pair of adjacent spaces of the grid the theta lift of a vector of the cuspidal space of the one
converges and is a vector of the cuspidal space of the other.

For $0\le i\le m$ let $C_i$ be the set of functions on $U(V_{i,0})(\mathbf A)$ that are values of
vectors of the cuspidal space of the grid at $V_{i,0}$. The function $u_w$ lies in $C_0$, by the core
conditions of the automorphic setting. If $\varphi\in C_i$ with $i<m$, the theta lift of $\varphi$ from
$U(V_{i,0})$ to $U(V_{i+1,0})$ with the vector $f_{i,\Phi_{a,i}}\otimes\rho_i$ converges and lies in
$C_{i+1}$, by the matching of the kernels. By induction on $i$, all $m$ integrals of
$\Theta_{\rho,\Phi_a}(u_w)$ converge, and $\Theta_{\rho,\Phi_a}(u_w)\in C_m$. The cuspidal space is
spanned by maps whose values are automorphic forms, and these properties of a function, continuity,
left invariance under the rational points and right invariance under an open compact subgroup, are
kept by sums and scalar multiples; so the values of its vectors have them.
% Lean: T.twisted_base_change_lift_exists = RowLift.twisted_lift_from_base_field hP hL: G' from
% T.twist_row_lift_exists hP (Lemma -convergence-lift-grid, open), CuspAlong G' from hL =
% PairLift.thetaPairLift_cuspidal_of_hypotheses (Lemma -convergence-lift-cuspidal, open) with
% RowSizes.sizes_of_sizeRows and rowBase_of_core, the sets C i of HEq-values of (G'.V i 0).Cusp, the step from hC.X with hM.ker and
% transport_stepS, start_mem from hS.core, row_induction, and isAutFun_transportS with
% isAutFun_of_mem_cusp (a span induction of the IsAut fields). In all about 3,100 lines at depth 5
% for twisted_linear_bound without the cuspidality.
\end{proof}

\begin{lemma}\label{asm-dependence-test-assembly-first-layer-convergence}
Assume the results of Theorem~\ref{lit-bundle-theta}. There is $C_{F,0}$ such that for every
$\varepsilon$ and every $C_F\ge C_{F,0}$ there is $x_0$ with the following property. For every $D$, its
size data and $R$ as in Definition~\ref{asm-dependence-test-assembly-statements} with $x\ge x_0$, and
all test choices for $C_F$ that satisfy the base conditions, the iterated theta lift
$\Theta_{\rho,\Phi_a}(u_w)$ and the test functional $\psi_a(u)$ converge for every $a$ and every $u\in E$.
\end{lemma}
\leanref{Development/AbelianVarietyTheory/Definitions/TwistedInnerProducts/TwistPlacesOffPrimeQZero.lean}{GlobalTheory.RowEx.LinB}
\leanref{Development/AutomorphicArithmetic/ThetaCorrespondence/ThetaPairLifts/LocalTypeDictionary/TwistedLinearBound.lean}{GlobalTheory.RowEx.twisted_linear_bound}
\leanref{Development/AutomorphicArithmetic/ThetaCorrespondence/ThetaPairLifts/LocalTypeDictionary/ThetaPairLiftCuspidalOfHypotheses.lean}{GlobalTheory.PairLift.thetaPairLift_cuspidal_of_hypotheses}
\begin{proof}
Take $C_{F,0}=0$ and $x_0=0$. Let $D$, its size data, $R$ and test choices satisfying the base
conditions be given, and let $a$ be a test index and $u\in E$. By
Lemma~\ref{asm-dependence-test-assembly-first-layer-convergence-lift}, $\Theta_{\rho,\Phi_a}(u_w)$
converges and is a continuous function on $U(V_{m,0})(\mathbf A)$, invariant on the left under the
rational points and on the right under an open compact subgroup. The rows of the size data have
their sizes, so $p_{de}$ and $p'_{de}$ are prime to $Q_Sq$ (Lemma~\ref{asm-dependence-rows-sizes}),
and the level system of $R$, the grid, its Weil datum and the row volume of the test choices form a
row setting (Definition~\ref{asm-dependence-rows-models}(3)). The function $\Psi_a$ is a column. So
Lemma~\ref{asm-dependence-rows-bessel-conv} shows that $\psi_a(u)=B^+\bigl(\varphi_{\xi,\Psi_a},
\Theta_{\rho,\Phi_a}(u_w)\bigr)$ converges.
% Lean: RowEx.twisted_linear_bound P (PairLift.thetaPairLift_cuspidal_of_hypotheses P), with C_F0 = 0
% and x0 = 0: T.twisted_base_change_lift_exists (Lemma -convergence-lift) gives ThetaConv and IsAutFun,
% and T.besselPeriodCorank_converges_column = RowBplus.bplusConv_col (Lemma
% asm-dependence-rows-bessel-conv) with hch.vol and col_mem ch a gives BplusConv. bplusConv_col uses
% only the grid, its Weil datum and the row volume. The open part is in Lemma -convergence-lift.
\end{proof}

\begin{lemma}\label{asm-dependence-test-assembly-first-layer-columns}
Let test choices be given as in Definition~\ref{asm-dependence-test-assembly-statements}(1), and let
$\Phi$ be a coset datum. Then the columns $\Psi_a$ with $\Phi_a=\Phi$ span the space of columns of
the grid (Definition~\ref{asm-grid-evaluation}(5)).
\end{lemma}
\leanref{Development/AbelianVarietyTheory/Proofs/TwistedOperators/EqAdFormsOfValues.lean}{GlobalTheory.RowTest.span_col_eq_cols}
\leanref{Development/AbelianVarietyTheory/Proofs/TwistedOperators/EqAdFormsOfValues.lean}{GlobalTheory.RowTest.span_colOf_psi}
\leanref{Development/AbelianVarietyTheory/Proofs/TwistedOperators/EqAdFormsOfValues.lean}{GlobalTheory.RowTest.eq_adForms_of_values}
\leanref{Development/AbelianVarietyTheory/Definitions/Divisors/GridTestLatticeLargestSlopeBound.lean}{GlobalTheory.FunObj.psi}
\begin{proof}
The test indices $a$ are the triples $(\Phi_a,c,i)$ of a coset datum, a class $c$ of
$U(Q_0)(F)\backslash U(Q_0)(\mathbf A_{F,f})/K^+$ and an index $i$ of a basis $(e_i)$ of the rational
coefficients, whose images $w_i$ span $W$. The class representatives $y_c$ are neat: an element of
$U(Q_0)(F)$ that lies in $y_cK^+y_c^{-1}$ is $1$. So for each $c$ and $i$ there is a column that
vanishes on the classes other than $c$ and takes the value $w_i$ at $y_c$, and this column is
$\Psi_{(\Phi,c,i)}$. Every $g$ is $\gamma y_ck$ with $\gamma\in U(Q_0)(F)$ and $k\in K^+$, so a column is
determined by its values at the $y_c$. Hence a column $\Psi$ is the sum of
$\Psi_{(\Phi,c,i)}$ with the coefficients of $\Psi(y_c)$ in the $w_i$, summed over $c$ and $i$.
% Lean: span_col_eq_cols = span_colOf_psi (exists_gf_reps, eq_adForms_of_values, colOf_mem,
% gridClassReps_delta_apply_rep with ClassReps.neat, span_mapVV). About 1,550 lines of Lean proofs at
% depth 5, mostly the adelic class decomposition.
\end{proof}

\begin{lemma}\label{asm-dependence-test-assembly-first-layer-injective}
Assume that functionals separate (Definition~\ref{asm-grid-evaluation}(7)) and the conclusion of
Lemma~\ref{asm-dependence-test-assembly-first-layer-convergence}. Then statement (b) on test choices
holds (Definition~\ref{asm-dependence-test-assembly-statements}(3)).
\end{lemma}
\leanref{Development/AbelianVarietyTheory/Proofs/TwistedOperators/InjKOfSeparatesLinB.lean}{GlobalTheory.RowEx.injK_of_separates_linB}
\leanref{Development/AbelianVarietyTheory/Proofs/TwistedOperators/InjKOfSeparates.lean}{GlobalTheory.RowTest.injK_of_separates}
\leanref{Development/AbelianVarietyTheory/Proofs/TwistedOperators/InjKOfSeparates.lean}{GlobalTheory.RowTest.inj_of_separates}
\leanref{Development/AbelianVarietyTheory/Definitions/TwistedInnerProducts/ExistsMatchOver.lean}{GlobalTheory.RowTest.Choice.RealOK}
\leanref{Development/AbelianVarietyTheory/Definitions/TwistedInnerProducts/ExistsMatchOver.lean}{GlobalTheory.RowTest.Choice.IsBase}
\begin{proof}
Let $C_{F,0}=\max(6510,C_1)$, where $C_1$ is the constant of
Lemma~\ref{asm-dependence-test-assembly-first-layer-convergence}, and for $\varepsilon$ and $C_F\ge
C_{F,0}$ let $x_0$ be the larger of the numbers that the separation of functionals and that lemma
give. Let $D$, its size data and $R$ be as in Definition~\ref{asm-dependence-test-assembly-statements}
with $x\ge x_0$, fix test choices for $C_F$ that satisfy the evaluation conditions, and let $u\in E$
with $\psi_a(u)=0$ for every $a$. Suppose $u\ne0$.

\emph{Separation.} The automorphic setting gives that the rows of the size data have their sizes,
that the series idempotents are tied to the labels, that the level characters of $R$ are those of
its base tuple, and that $R$ records the group. The test choices give an evaluation grid $G$ with
constant $C_F\ge6510$ at $b$, a Weil datum $W$ over $R$, a row volume and a smooth vector $w\ne0$. So
functionals separate $w$: there are a real vector $\rho'$ and a linear functional $\xi'$ on the
coefficients such that for every nonzero element $t$ of the space of theta multiplicities at the
trivial character there are a column $\Psi$ and a coset datum $\Phi$ with
$B^+\bigl(\varphi_{\xi',\Psi},\Theta_{\rho',\Phi}(u_{w,t})\bigr)\ne0$, both integrals converging.

\emph{The element $u$.} By the core conditions on the models, $E$ is identified, by an injective map,
with the good part at the trivial level of the multiplicity maps of the model of $U(Y_0)$; let $t$ be
the image of $u$, which is nonzero, so that $u_w=u_{w,t}$. Let $\Psi$ and $\Phi$ be as above. Write
$\psi'_a$ for the test functionals of the test choices with $\rho'$ and $\xi'$ in place of $\rho$ and
$\xi$. These test choices still satisfy the base conditions, so by
Lemma~\ref{asm-dependence-test-assembly-first-layer-convergence} every $\psi'_a(u)$ converges; and by
the real condition, $\psi'_a(u)=0$ for every $a$.

\emph{Contradiction.} By Lemma~\ref{asm-dependence-test-assembly-first-layer-columns}, $\Psi$ is a
linear combination of the columns $\Psi_a$ with $\Phi_a=\Phi$. The map $\Psi'\mapsto\varphi_{\xi',\Psi'}$
is linear, and where the integrals converge the Bessel period $B^+(\varphi,G)$ is linear in $\varphi$.
For each such $a$, $B^+\bigl(\varphi_{\xi',\Psi_a},\Theta_{\rho',\Phi}(u_w)\bigr)=\psi'_a(u)$ converges and
vanishes. So $B^+\bigl(\varphi_{\xi',\Psi},\Theta_{\rho',\Phi}(u_w)\bigr)=0$, a contradiction. Hence $u=0$.
% Lean: injK_of_separates_linB: LinB at (ch.withReal rho xi) with IsBase.withReal;
% injK_of_separates (the constants max 6510 C1 and max x1 x2; the hypotheses of FunctionalsSeparate
% from hset.size, hset.eser, hset.core.lev, hset.grp, hch.weil, hch.vol, hch.w_ne) and
% inj_of_separates (hIs.real, hM.HE, span_col_eq_cols, and span induction with bplusConv_add,
% bplusConv_smul, bplusConv_zero). About 1,000 lines of Lean proofs at depth 5.
\end{proof}

\begin{lemma}\label{asm-dependence-test-assembly-first-layer-ratio-inputs-nonvanishing}
Assume the functional equation of Hecke $L$-functions of CM type (Theorem~\ref{asm-dependence-inputs}).
Let $E$ be a CM field with a CM type $\Phi$, and let $\nu$ be a unitary Hecke character of $E$ whose
archimedean component satisfies $\nu_\infty(k)=\prod_{\tau\in\Phi}\bigl(\tau(k)/|\tau(k)|\bigr)^{2m(\tau)}$
for $k\in E^\times$, with integers $m(\tau)>0$. Then $\nu$ has an $L$-function
(Definition~\ref{th-large-data}), and $\Lambda(s)\ne0$ for every $L$-function $\Lambda$ of $\nu$ and
every integer $s$ with $s+m(\tau)>0$ for all $\tau\in\Phi$.
\end{lemma}
\leanref{Development/AbelianVarietyTheory/Proofs/DivisorProofs/SquareBottomHalfRankDiffLeHalfRank.lean}{GlobalTheory.LValNe.heckeL_ne_zero_int}
\leanref{Development/AbelianVarietyTheory/Proofs/DivisorProofs/SquareBottomHalfRankDiffLeHalfRank.lean}{GlobalTheory.LValNe.exists_heckeL}
\leanref{Development/AbelianVarietyTheory/Proofs/DivisorProofs/SquareBottomHalfRankDiffLeHalfRank.lean}{GlobalTheory.LValNe.completed_eq}
\leanref{Development/AbelianVarietyTheory/PublishedTheorems/Divisors₂/HeckeFunctionalEquationAbs.lean}{GlobalTheory.Published.HeckeFunctionalEquationAbs}
\leanref{Development/AbelianVarietyTheory/BlockTheory/DivisorBase/CompletedLStripBounds.lean}{GlobalTheory.IndefiniteRowSideFullHecke.CompletedLStripBounds.CompletedLBoundedInStrips}
\begin{proof}
The infinity type of $\nu$ is nonzero, so $\nu$ is not a power of the norm, and by Hecke's theorem
(Theorem~\ref{lit-bundle-theta}(3)) its $L$-function $L(s,\nu)$, the Euler product over the places
where $\nu$ is unramified, continues to an entire function and has no zeros on $\operatorname{Re}s\ge1$.
This continuation is an $L$-function of $\nu$, and any two $L$-functions of $\nu$ agree on
$\operatorname{Re}s>1$, hence everywhere. The character $\nu^{-1}$ is unitary of infinity type $-2m$, so
the same holds for it, with an $L$-function $\Lambda'$. For $\nu$ and $\nu^{-1}$ put
\[
 \gamma(z)=\prod_{\tau\in\Phi}2(2\pi)^{-(z+m(\tau))}\,\Gamma\bigl(z+m(\tau)\bigr).
\]
By Hecke's theorem \cite{Hec20,Tat67} the completed $L$-functions $\Lambda_c$ of $\nu$ and $\Lambda'_c$
of $\nu^{-1}$, entire functions equal to $\gamma(z)L(z,\nu)$ and $\gamma(z)L(z,\nu^{-1})$ for
$\operatorname{Re}z>1$, exist, and by the functional equation
$\Lambda_c(z)=\varepsilon A^{1/2-z}\Lambda'_c(1-z)$ for all $z$, with $|\varepsilon|=1$ and $A$ a positive
integer, the absolute discriminant of $E$ times the norm of the conductor of $\nu$. Where all the
$\Gamma(z+m(\tau))$ are finite, $\Lambda_c(z)=\gamma(z)\Lambda(z)$, since $\Lambda_c/\gamma$ is entire,
as $1/\Gamma$ is, and agrees with $\Lambda$ for $\operatorname{Re}z>1$; the same holds for $\nu^{-1}$.

Let $s$ be an integer with $s+m(\tau)>0$ for all $\tau$. If $s\ge1$, $\Lambda(s)=L(s,\nu)\ne0$. Let
$s\le0$. Then $1-s\ge1$, so $\Lambda'(1-s)\ne0$, and $1-s+m(\tau)>0$, so
$\Lambda'_c(1-s)=\gamma(1-s)\Lambda'(1-s)\ne0$. By the functional equation $\Lambda_c(s)\ne0$. As
$s+m(\tau)>0$, the factors $\Gamma(s+m(\tau))$ are finite, so $\Lambda_c(s)=\gamma(s)\Lambda(s)$ and
$\Lambda(s)\ne0$.
% Lean: LValNe.heckeL_ne_zero_int: the case s >= 1 is IsHeckeL.ne_zero_nat (eulerProduct_ne_zero_one_of_entire
% at s = 1, the Euler product for s >= 2); for s <= 0, CompletedLBoundedInStrips (proved in the archive
% as HeckeStrip.completedLBoundedInStrips) gives the completed functions for chi with n = 2m and for
% chi^{-1} with -2m, Published.HeckeFEModel.zero_iff the functional equation in the form Lambda_c(s) = 0
% iff Lambda'_c(1 - s) = 0, and completed_eq (entire_eq_of_halfPlane) the relation Lambda_c = gamma Lambda
% where the gamma factors are finite. The existence of Lambda is LValNe.exists_heckeL
% (ToricDistribution.exists_isHeckeL_of_hasInf). The Lean's archFactor is the book's gamma with
% |n(tau)|/2 = m(tau). The book cites Hecke's continuation and the functional equation as in the proof
% of Lemma asm-theta-agree-nonvanishing-hecke.
\end{proof}

\begin{lemma}\label{asm-dependence-test-assembly-first-layer-ratio-inputs-characters}
Let $D$ be a definite-side datum with host field $E/F$, embedding $\varphi_0$, set $\mathcal T$ and
Harish-Chandra parameters $\lambda(\tau)$ with $e$-parts $e_1(\tau),\dots,e_m(\tau)$
(Definition~\ref{th-definite-datum}), and let a level system with base tuple $\mu_0$ and a theta grid
for it be given (Definition~\ref{th-grid}), with the two characters $\delta$ and $\delta'$ of finite
order of the level system (condition (L3) there). For $k\in E^\times$ let $k_\infty$ be
the idele with archimedean components $k$ and finite components $1$. Let $0\le i,j<m$, let
$c_{m-i}(\tau)$ be the integers, fixed by the grid, with
$\chi_{m-i}(k_\infty)=\prod_{\tau\in\mathcal T}\bigl(\tau(k)/|\tau(k)|\bigr)^{2c_{m-i}(\tau)}$ for
$k\in E^\times$, and put $\kappa(\tau)=e_{j+1}(\tau)+c_{m-i}(\tau)$.
\begin{enumerate}
\item $E$ is a CM field, and every embedding $\sigma\colon E\to\mathbf C$ is $\tau$ or $\bar\tau$ for
  exactly one $\tau\in\mathcal T$. So $\Phi_D=\{\bar\tau:\tau\in\mathcal T\}$ is a CM type of $E$.
\item For every tuple $t$, $\nu_{t,i,j}=\nu_{\mu_0,i,j}\xi_t$. Each $\nu$ among the characters
  $\nu_{t,i,j}$, $\nu_{\mu_0,i,j}\delta$ and $\nu_{\mu_0,i,j}\delta'$ is unitary, and
  \[
   \nu(k_\infty)=\prod_{\tau\in\mathcal T}\Bigl(\frac{\bar\tau(k)}{|\tau(k)|}\Bigr)^{2\kappa(\tau)}
   \qquad(k\in E^\times).
  \]
\item $\kappa(\tau)>h_1(i,j)$ for every $\tau\in\mathcal T$.
\end{enumerate}
\end{lemma}
\leanref{Development/AbelianVarietyTheory/Proofs/DivisorProofs/NatCastMemBelow.lean}{GlobalTheory.SqTop.FieldDatPf.isCMType_cmPhi}
\leanref{Development/AbelianVarietyTheory/Definitions/Divisors/IsValCharTwist.lean}{GlobalTheory.FunObj.SqNum.Model.isValChar_twist}
\leanref{Development/AbelianVarietyTheory/Proofs/IwasawaDecomposition/HeckeCharacterIdeleCharacterMatchUnique.lean}{GlobalTheory.ChrEqUnique.isUnitary_of_chrEq}
\leanref{Development/AbelianVarietyTheory/Proofs/FoundationProofs/IntValuationEqOfMem.lean}{GlobalTheory.TypeValues.inf_eq_one}
\leanref{Development/AbelianVarietyTheory/Definitions/Divisors/IsValCharTwist.lean}{GlobalTheory.FunObj.SqNum.Model.kap_add}
\leanref{Development/AbelianVarietyTheory/Definitions/Divisors/IsValCharTwist.lean}{GlobalTheory.FunObj.SqNum.Model.gap_le}
\begin{proof}
(1) $E$ is a CM field by Proposition~\ref{hf-host-field-properties}(a). The elements of $\mathcal T$ are
the elements $g$ of $\Psi$, read as the embeddings $\varphi_0\circ\gamma^{-1}(g)$
(Definitions~\ref{th-definite-datum}(2) and~\ref{hf-host}). For $g\in\Psi$,
$\gamma^{-1}(g)(\sqrt{-d})=\epsilon_2(g)\sqrt{-d}=\sqrt{-d}$ by
Definition~\ref{hf-framed-host-field}(ii), so $\tau(\sqrt{-d})=\varphi_0(\sqrt{-d})=i\sqrt d$ for every
$\tau\in\mathcal T$; and distinct $g$ give distinct embeddings. Conversely, an embedding $\sigma$ with
$\sigma(\sqrt{-d})=i\sqrt d$ is $\varphi_0\circ\gamma^{-1}(g)$ for some $g\in G_\Phi$, as $E/\mathbf Q$ is
Galois, and then $\epsilon_2(g)=1$, so $g\in\Psi$. So the elements of $\mathcal T$ are the embeddings
$\sigma$ with $\sigma(\sqrt{-d})=i\sqrt d$. For every $\sigma$, exactly one of $\sigma$ and $\bar\sigma$
has this property, since complex conjugation changes the sign of $\sigma(\sqrt{-d})$. Hence $\sigma$ is
$\tau$ or $\bar\tau$ for some $\tau\in\mathcal T$; and if $\sigma\in\{\tau,\bar\tau\}\cap\{\tau',\bar\tau'\}$
with $\tau,\tau'\in\mathcal T$, then $\tau'\in\{\tau,\bar\tau\}$, and $\bar\tau\notin\mathcal T$, so
$\tau'=\tau$.

(2) By the definitions, $\nu_{t,i,j}=\mu^0_{j+1}\xi_t/\chi_{m-i}$, and $\xi_{\mu_0}=1$; so
$\nu_{t,i,j}=\nu_{\mu_0,i,j}\xi_t$.

\emph{Unitarity.} The level system and the grid give $\mu^0_{j+1}$, $\xi_t$ and $\chi_{m-i}$ together
with characters of the idele class group of $E$ with values in the unit circle that have the same
components at all finite places. A Hecke character is determined by its components at the finite
places, since $E^\times\mathbf A_{E,f}^\times$ is dense in $\mathbf A_E^\times$ by weak approximation at
the archimedean places. So $\mu^0_{j+1}$, $\xi_t$ and $\chi_{m-i}$ take values in the unit circle, and
so does $\nu_{t,i,j}$. The characters $\delta$ and $\delta'$ have finite order by (L3) of
Definition~\ref{th-grid}, so they also take values in the unit circle, and $\nu_{\mu_0,i,j}\delta$ and $\nu_{\mu_0,i,j}\delta'$ are unitary.

\emph{The archimedean component.} Let $\xi$ be a Hecke character of finite order $n$, and $k\in
E^\times$. The field $E$ is totally complex, so at each archimedean place $w$ the component $k_w$ is
$z_w^n$ for some $z_w\in\mathbf C^\times$, and the component of $\xi$ at $w$ takes the value
$\xi_w(z_w)^n=\xi_w^n(z_w)=1$ at $k_w$. So $\xi(k_\infty)=1$; this applies to $\xi_t$, $\delta$ and
$\delta'$. The level system gives, by (L2),
$\mu^0_{j+1}(k_\infty)=\prod_{\tau}\bigl(\tau(k)/|\tau(k)|\bigr)^{-2e_{j+1}(\tau)}$, and $\chi_{m-i}$ has
the exponents $2c_{m-i}(\tau)$. Hence each $\nu$ of (2) satisfies
\[
 \nu(k_\infty)=\prod_{\tau\in\mathcal T}\Bigl(\frac{\tau(k)}{|\tau(k)|}\Bigr)^{-2e_{j+1}(\tau)-2c_{m-i}(\tau)}
 =\prod_{\tau\in\mathcal T}\Bigl(\frac{\bar\tau(k)}{|\tau(k)|}\Bigr)^{2\kappa(\tau)},
\]
since $\bar\tau(k)/|\tau(k)|=\bigl(\tau(k)/|\tau(k)|\bigr)^{-1}$.

(3) Put $N=k+m(t+r+7)$, where $k$, $m$, $t$ and $r=k+m-1$ are the ranks of $D$. The grid builds its
centres $c_l(\tau)$ from $\lambda(\tau)$ with a gap $g>2N^3$, and a computation from the definition of
the centres and the interlacing of $\lambda(\tau)$ gives $g\le e_{j+1}(\tau)+c_{m-i}(\tau)=\kappa(\tau)$.
As $\operatorname{rk}V_{i,j}<\operatorname{rk}V_{i+1,j+1}=r+(j+1)t+(i+1)(r+6)\le r+m(t+r+6)<N$, we have
$h_1(i,j)<N\le N^3<g\le\kappa(\tau)$.
% Lean: (1) SqTop.FieldDatPf.isCMType_cmPhi (with D.host.isCMField_hostE): LamTie.CMType.phi_cmType from
% Mem.HostCM.host_cm_embedding_properties, whose bijectivity of g -> InfinitePlace.mk (embPsi H g) rests on
% FactEmb (Host.emb g = phi0 o sigmaOf g is bijective, and emb g lies in cmType = {psi : psi sqd = i sqrt d}
% iff g in Psi), then isCMType_conj; the Lean's cmPhi D tau is the conjugate of embTau D tau.
% (2) SqTop.lNe_square and FunObj.SqNum.Model.isValChar_twist (with xi = 1, delta = L.C.del or
% L.C.del', IsDelta): nuH G t i j = nuH G mu0 i j * Lc.xiH t from hL.xi_base; unitarity from
% ChrEqUnique.isUnitary_of_chrEq with FunObj.SqNum.chrEq_nu (nuH matches the Circle-valued
% L.mu0 (j+1) * L.xi t / G.ch.chi (m - i); hchar_unique) and HChar.isUnitary_of_order for delta;
% the archimedean component from hm.inf of IsMu0 (exponent -2 e_{j+1} at embTau, through ofHCParam_spec),
% G.ch.inf (exponent 2 cen at embTau), TypeValues.inf_eq_one for the characters of finite order, and
% cmPhi = conjugate of embTau. The Lean's HasInf in the ratio form has exponent kap D cGap i j 0 =
% e_{j+1} + cen, the book's kappa. (3) kap_add D cGap i j tau at s = halfRank, from gap_le (gap <= e +
% cen, from e_add_v and the definition of cen as the floor of the midpoint), isLeast_gap
% (2 bigN^3 + 2 sLen + 2 < gap) and halfRank <= bigN; bigN D = k + m (t + sLen + 1) with sLen = r + 6.
\end{proof}

\begin{lemma}\label{asm-dependence-test-assembly-first-layer-ratio-inputs-definite}
Let $D$ be a definite-side datum with a standard character $\psi$
(Definition~\ref{th-weil-datum}), let $S$ and $W$ be Hermitian spaces of $D$ with $S$ totally positive
definite, and let $\omega$ be a Weil datum for $(S,W)$ or for $(-S,W)$. Let $\nu$ be a Haar measure on
$U(S)(\mathbf A)$, $\mathcal F$ a fundamental domain for $U(S)(F)$ with respect to $\nu$, and $f$ a
theta function. Then:
\begin{enumerate}
\item for every $h\in U(W)(\mathbf A)$ the function $g\mapsto\theta_\omega(g,h;f)$ is integrable on
  $\mathcal F$;
\item the function $I(h)=\int_{\mathcal F}\theta_\omega(g,h;f)\,d\nu(g)$ is continuous on
  $U(W)(\mathbf A)$;
\item $I(\gamma h)=I(h)$ for all $\gamma\in U(W)(F)$ and $h$.
\end{enumerate}
\end{lemma}
\leanref{Development/AbelianVarietyTheory/Proofs/TwistedOperators/ThetaIntegralDefiniteSourceConvergesContinuousInStatement.lean}{GlobalTheory.T.thetaIntegral_definiteSource_converges_continuous_invariant}
\leanref{Development/AbelianVarietyTheory/Proofs/TwistedOperators/ThetaIntegralDefiniteSourceConvergesContinuousIn.lean}{GlobalTheory.DefSourceInt.theta_integral_source_definite}
\begin{proof}
Since $S$ is totally definite, $U(S)$ is anisotropic over $F$, and there is a compact set $C$ with
$U(S)(\mathbf A)=U(S)(F)\,C$ \cite{Bor69}. The theta kernel
$\theta_\omega(g,h;f)=(\omega(g,h)f)(0)$ is continuous in $(g,h)$, and invariant on the left under
$U(S)(F)\times U(W)(F)$, by the conditions on a Weil datum (Definition~\ref{th-weil-datum}(3)).

(1) For fixed $h$, $g\mapsto\theta_\omega(g,h;f)$ is continuous and invariant on the left under
$U(S)(F)$, so it is bounded by its maximum on $C$. The fundamental domain $\mathcal F$ has measure at
most $\nu(C)<\infty$, since $C$ meets every orbit of $U(S)(F)$. So the function is integrable on
$\mathcal F$.

(2) Let $h_0\in U(W)(\mathbf A)$ and let $K$ be a compact neighbourhood of $h_0$. The kernel is bounded on
the compact set $C\times K$, hence, by its invariance in $g$, on $U(S)(\mathbf A)\times K$. By dominated
convergence on $\mathcal F$, which has finite measure, $I$ is continuous at $h_0$.

(3) $\theta_\omega(g,\gamma h;f)=\theta_\omega(g,h;f)$ for every $g$, so the integrals agree.
% Lean: DefSourceInt.theta_integral_source_definite: definite_unitary_cocompact (the compact C),
% thetaKO_continuous_and_rat_invariant, CocompactIntegrable.integrableOn_fundamentalDomain for (1),
% DefLiftCont.continuous_integral_of_cocompact for (2), and hIW for (3). The sign epsilon of the Lean
% gives the two cases (S, W) and (-S, W).
\end{proof}

\begin{lemma}\label{asm-dependence-test-assembly-first-layer-ratio-inputs-vanish}
Let $D$ be a definite-side datum and $W$ a totally positive definite Hermitian space of $D$. Let
$\phi$ be a continuous function
on $U(W)(\mathbf A)$ with $\phi(\gamma h)=\phi(h)$ for all $\gamma\in U(W)(F)$ and $h$, that vanishes at
every element whose archimedean components are $1$. Then $\phi=0$.
\end{lemma}
\leanref{Development/AbelianVarietyTheory/Proofs/TwistedOperators/ContinuousAutomorphicEqZeroOfVanishOnFiniteAdeliStatement.lean}{GlobalTheory.T.continuous_automorphic_eq_zero_of_vanish_on_finiteAdelic}
\begin{proof}
Let $h\in U(W)(\mathbf A)$, and write $h=uf$ with $u$ in the product $U(W)(F_\infty)$ of the archimedean
groups and $f$ with archimedean components $1$. The function $\psi(x)=\phi(xf)$ is continuous, and for
$\gamma\in U(W)(F)$ and $k$ with archimedean components $1$,
$\psi(\gamma k)=\phi(\gamma kf)=\phi(kf)=0$, as $kf$ has archimedean components $1$. If $\psi(u)\ne0$,
the set where $\psi$ does not vanish is open and contains $u$; by weak approximation for $U(W)$ at the
archimedean places \cite[Chapter~7]{PR94} it contains an element $\gamma k$ of this kind, which is
impossible. So $\phi(h)=\psi(u)=0$.
% Lean: Adopt.continuous_automorphic_eq_zero_of_vanish_on_finiteAdelic: UnitarySplit.real_finite_split
% and RealApprox.weak_approx_real. The hypothesis that W is positive definite is not used.
\end{proof}

\begin{lemma}\label{asm-dependence-test-assembly-first-layer-ratio-inputs-standard-waves}
Let $G$ be an abelian group in which $kt=0$ implies $t=0$ for every integer $k\ge1$, and let $X,Y\subset G$
be subgroups such that for every $k\ge1$ every element of $X$ is $k$ times an element of $X$, and every
element of $Y$ is $k$ times an element of $Y$. Let $\chi\colon G\times G\to\mathbf C$ satisfy
$\chi(a+b,y)=\chi(a,y)\chi(b,y)$, $\chi(a,y+z)=\chi(a,y)\chi(a,z)$ and $\chi(0,y)=\chi(a,0)=1$. Let
$L\subset X$ and $L'\subset Y$ be subgroups such that:
\begin{enumerate}
\item[(i)] $\chi(u,y)=1$ for $u\in L$ and $y\in L'$;
\item[(ii)] an element $x\in X$ with $\chi(x,y)=1$ for every $y\in L'$ lies in $L$, and an element $y\in Y$
  with $\chi(u,y)=1$ for every $u\in L$ lies in $L'$;
\item[(iii)] for every $k\ge1$ there is a finite set $R\subset G$ with $L=R+kL$.
\end{enumerate}
Call a function $\Phi$ on $G$ \emph{Schwartz--Bruhat} if it vanishes off $X$ and there is an integer
$k\ge1$ with $\Phi(t+ku)=\Phi(t)$ for all $t\in G$ and $u\in L$, and with $kt\in L$ whenever $\Phi(t)\ne0$.
For $x\in X$ and $y\in Y$ let $w_{x,y}(t)=\chi(t+x,y)1_L(t+x)$, where $1_L$ is the indicator function of
$L$, and let $B$ be the space spanned by these functions.
\begin{enumerate}
\item If $l\in L$ and $l'\in L'$, then $w_{x+l,y+l'}=\chi(l,y)\,w_{x,y}$. The functions $w_{x,y}$, with $x$
  and $y$ running over sets of representatives of $X/L$ and $Y/L'$, are linearly independent.
\item Every Schwartz--Bruhat function lies in $B$.
\item Let $V$ be a complex vector space and $\rho$ a map from $G$ to the linear endomorphisms of $V$ with
  $\rho(x)\rho(x')=\rho(x+x')$ for $x,x'\in X$ and for $x,x'\in Y$, and
  $\rho(x)\rho(y)=\chi(x,y)\rho(y)\rho(x)$ for $x\in X$ and $y\in Y$. Let $f_0\in V$ with $\rho(u)f_0=f_0$ for
  every $u\in L$ and every $u\in L'$. Then there is a linear map $T$ from the functions on $G$ to $V$ with
  $T(1_L)=f_0$ such that, for every Schwartz--Bruhat function $\Phi$, every $x\in X$ and every $y\in Y$,
  the translate $t\mapsto\Phi(t+x)$ goes to $\rho(x)T(\Phi)$ and the product $t\mapsto\chi(t,y)\Phi(t)$
  goes to $\rho(y)T(\Phi)$.
\item Suppose moreover that every element of $Y$ has a positive multiple in $L'$, and let $V$, $\rho$, $f_0$
  and $T$ be as in (3). Then $T(w_{x,y})=\rho(x)\rho(y)f_0$ for all $x\in X$ and $y\in Y$. So if
  $\Phi=\sum_ic_iw_{x_i,y_i}$ is Schwartz--Bruhat, then $T(\Phi)=\sum_ic_i\rho(x_i)\rho(y_i)f_0$.
\end{enumerate}
\end{lemma}
\leanref{Development/AbelianVarietyTheory/Proofs/AdelicLatticeProofs/LinearIndependentWave.lean}{GlobalTheory.PlaceModel.linearIndependent_wave}
\leanref{Development/AbelianVarietyTheory/Proofs/AdelicLatticeProofs/LinearIndependentWave.lean}{GlobalTheory.PlaceModel.exists_map_wave}
\leanref{Development/AbelianVarietyTheory/Proofs/AdelicLatticeProofs/LinearIndependentWave.lean}{GlobalTheory.PlaceModel.ind_nsmul_mem}
\leanref{Development/AbelianVarietyTheory/Proofs/AdelicLatticeProofs/LinearIndependentWave.lean}{GlobalTheory.PlaceModel.mem_of_finite_cosets}
\leanref{Development/AbelianVarietyTheory/Proofs/AdelicLatticeProofs/LinearIndependentWave.lean}{GlobalTheory.PlaceModel.exists_intertwiner}
\leanref{Development/AbelianVarietyTheory/Proofs/IwasawaDecomposition/IdempotentMulEqZeroOfOrthogonal.lean}{GlobalTheory.PlaceDatum.mem_waveSpan}
\begin{proof}
The value $\chi(a,y)$ is nonzero, since $\chi(a,y)\chi(-a,y)=\chi(0,y)=1$.

(1) Let $l\in L$, $l'\in L'$ and $t\in G$. If $t+x\notin L$, then $t+x+l\notin L$ and both functions vanish
at $t$. If $t+x\in L$, then $\chi(t+x+l,l')=1$ by (i), so
$w_{x+l,y+l'}(t)=\chi(t+x+l,y)=\chi(l,y)\chi(t+x,y)=\chi(l,y)w_{x,y}(t)$.

Fix representatives $x_a$ of the classes $a\in X/L$ and $y_c$ of the classes $c\in Y/L'$, and let
$\sum_{a,c}g(a,c)\,w_{x_a,y_c}=0$, a finite sum. Fix $a$, and let $u\in L$. For $b\ne a$, $u-x_a+x_b\notin L$,
so $w_{x_b,y}(u-x_a)=0$; and $w_{x_a,y}(u-x_a)=\chi(u,y)$. Evaluating at $u-x_a$ gives
$\sum_cg(a,c)\chi(u,y_c)=0$ for every $u\in L$. The maps $u\mapsto\chi(u,y_c)$ are homomorphisms from $L$ to
$\mathbf C^\times$, and they are distinct: if $\chi(u,y_c)=\chi(u,y_{c'})$ for every $u\in L$, then
$\chi(u,y_{c'}-y_c)=1$ for every $u\in L$, so $y_{c'}-y_c\in L'$ by (ii), and $c=c'$. By the linear
independence of characters, $g(a,c)=0$ for every $c$.

(2) The space $B$ contains $1_L=w_{0,0}$. It is stable under translation by the elements of $X$ and under
multiplication by the functions $\chi(\cdot,y)$ with $y\in Y$, since $w_{x_0,y_0}(t+x)=w_{x+x_0,y_0}(t)$ and
$\chi(t,y)w_{x_0,y_0}(t)=\chi(-x_0,y)\,w_{x_0,y+y_0}(t)$, as $\chi(t,y)=\chi(-x_0,y)\chi(t+x_0,y)$. So it is
stable under multiplication by the functions $\alpha\chi(\cdot,y)+\beta$ with $\alpha,\beta\in\mathbf C$.
\begin{enumerate}
\item[(a)] \emph{The indicator function of $kL$ lies in $B$ for $k\ge1$.} First, for $t\in L$ with $t\notin
  kL$ there is $y\in Y$ with $ky\in L'$ and $\chi(t,y)\ne1$. Indeed, write $t=kx$ with $x\in X$, and suppose
  $\chi(t,y)=1$ whenever $y\in Y$ and $ky\in L'$. For $y'\in L'$, write $y'=ky$ with $y\in Y$; then
  $\chi(x,y')=\chi(kx,y)=\chi(t,y)=1$. So $x\in L$ by (ii), and $t\in kL$, a contradiction. Let $R$ be a
  finite set with $L=R+kL$, let $F$ be the set of $r\in R$ with $r\in L$ and $r\notin kL$, and choose such a
  $y_r$ for each $r\in F$. Then
  \[
   1_{kL}(t)=1_L(t)\prod_{r\in F}\frac{\chi(t,y_r)-\chi(r,y_r)}{1-\chi(r,y_r)}.
  \]
  Both sides vanish off $L$. Let $t\in L$, and write $t=r_0+ku$ with $r_0\in R$ and $u\in L$. For $r\in F$,
  $\chi(t,y_r)=\chi(r_0,y_r)\chi(u,ky_r)=\chi(r_0,y_r)$ by (i). If $t=kv$ with $v\in L$, then
  $\chi(t,y_r)=\chi(v,ky_r)=1$, and every factor is $1$. Otherwise $r_0=t-ku$ lies in $L$ and not in $kL$, so
  $r_0\in F$, and its factor is $0$. As $1_L\in B$, the right side lies in $B$.
\item[(b)] \emph{Every Schwartz--Bruhat function lies in $B$.} Let $\Phi$ be one, with the integer $k$. It
  is invariant under translation by $kL$. It vanishes off finitely many cosets of $kL$: let $R$ be finite
  with $L=R+kL$. If $\Phi(t)\ne0$, then $kt\in L$, so $kt=r+ku$ with $r\in R$ and $u\in L$, and $s=t-u$
  satisfies $ks=r$; writing $u=r'+ku'$ with $r'\in R$ and $u'\in L$ gives $t=s+r'+ku'$. As multiplication by
  $k$ is injective, $s$ lies in the finite set of the $s'$ with $ks'\in R$. So the cosets of $kL$ on which
  $\Phi$ does not vanish are finitely many, $t_i+kL$ with $\Phi(t_i)\ne0$, $1\le i\le n$; then $t_i\in X$, and
  $\Phi=\sum_i\Phi(t_i)\,1_{kL}(\cdot-t_i)$, a sum of translates of $1_{kL}$ by elements of $X$, which lies in
  $B$ by (a).
\end{enumerate}

(3) By (1), the functions $w_{x_a,y_c}$ form a basis of $B$. Let $T$ be a linear map on the functions on $G$
with $T(w_{x_a,y_c})=\rho(x_a)\rho(y_c)f_0$ for all $a$ and $c$, extended from $B$ to all functions by a
complement of $B$. Then $T(w_{x,y})=\rho(x)\rho(y)f_0$ for all $x\in X$ and $y\in Y$. Indeed, write
$x=x_a+l$ and $y=y_c+l'$ with $l\in L$ and $l'\in L'$. By (1), $w_{x,y}=\chi(l,y_c)w_{x_a,y_c}$; and
\[
 \rho(x)\rho(y)f_0=\rho(x_a)\rho(l)\rho(y_c)\rho(l')f_0=\rho(x_a)\rho(l)\rho(y_c)f_0
 =\chi(l,y_c)\rho(x_a)\rho(y_c)\rho(l)f_0=\chi(l,y_c)\rho(x_a)\rho(y_c)f_0.
\]
In particular $T(1_L)=T(w_{0,0})=f_0$. For $x\in X$ and $y\in Y$,
$T(w_{x_0,y_0}(\cdot+x))=T(w_{x+x_0,y_0})=\rho(x)\rho(x_0)\rho(y_0)f_0=\rho(x)T(w_{x_0,y_0})$, and
\[
 T(\chi(\cdot,y)w_{x_0,y_0})=\chi(-x_0,y)\rho(x_0)\rho(y)\rho(y_0)f_0
 =\chi(-x_0,y)\chi(x_0,y)\rho(y)\rho(x_0)\rho(y_0)f_0=\rho(y)T(w_{x_0,y_0}).
\]
Both identities are linear in the function, so they hold on $B$, and so for every Schwartz--Bruhat function
by (2).

(4) Let $y\in Y$, and $k\ge1$ with $ky\in L'$. The function $\chi(\cdot,y)1_L$ is Schwartz--Bruhat with the
integer $k$: it vanishes off $L\subset X$; for $u\in L$, $\chi(t+ku,y)1_L(t+ku)=\chi(t,y)\chi(u,ky)1_L(t)=
\chi(t,y)1_L(t)$; and $kt\in L$ for $t\in L$. So is $1_L$, with the integer $1$. The function $w_{x,y}$ is
the translate of $\chi(\cdot,y)1_L$ by $x$. So $T(w_{x,y})=\rho(x)T(\chi(\cdot,y)1_L)=\rho(x)\rho(y)T(1_L)
=\rho(x)\rho(y)f_0$. The last assertion follows by linearity.
% Lean: (1) wave_add and linearIndependent_wave (the characters psi d of L, linearIndependent_monoidHom);
% (2) wave_trans and wave_modul (trans_mem, modul_mem), ind_nsmul_mem for (a) (hsep, the product hprod
% over F) and mem_of_finite_cosets for (b), combined in PlaceDatum.mem_waveSpan with the finite set for
% k * k, and in exists_intertwiner with the finite set for k used twice; (3) exists_map_wave
% (Module.Basis.span, constr, LinearMap.exists_extend; the computation hrho) and exists_intertwiner
% (htrans, hmod by span induction); (4) is the computation in PlaceDatum.placeModel_unique_of_dual and
% SchrTheta.placeModel_induction (s0 = isSBHalf_halfInd, s1 with the integer k * 1, and h.trans, h.modul,
% h.vac). LinearIndependentWave.lean has about 450 lines.
\end{proof}

\begin{lemma}\label{asm-dependence-test-assembly-first-layer-ratio-inputs-standard}
Let $D$ be a definite-side datum with host field $E/F$ and a standard character $\psi$
(Definition~\ref{th-weil-datum}), let $S$ and $W$ be Hermitian spaces of $D$ with self-dual lattices, let
$\Lambda\subset\mathbb W$ be the lattice of the pair, and let $f_0$ be a Fock vacuum vector for $(S,W)$.
Let $\varepsilon$ be a sum of the idempotents of $\mathbf A_E$ at the places of a finite set, with
$\varepsilon\bar\varepsilon=0$, and let $a$ be $\varepsilon$ or $\bar\varepsilon$. Call a function $\Phi$ on
$\mathbb W$ \emph{Schwartz--Bruhat on $a\mathbb W$} if it vanishes off $a\mathbb W$ and there is an integer
$k\ge1$ with $\Phi(x+ku)=\Phi(x)$ for all $x\in\mathbb W$ and $u\in\Lambda\cap a\mathbb W$, and with $kx\in\Lambda$
whenever $\Phi(x)\ne0$. A \emph{place model} of $f_0$ at $a$ is a linear map $T$ on functions on $\mathbb W$
that carries, for every such $\Phi$, the translate $x'\mapsto\Phi(x'+x)$ by $x\in a\mathbb W$ to $\rho(x)T(\Phi)$
and the product of $\Phi$ with $x'\mapsto\psi(\langle x',y\rangle)$ for $y\in\bar a\mathbb W$ to $\rho(y)T(\Phi)$,
and the indicator function of $\Lambda\cap a\mathbb W$ to $f_0$.
\begin{enumerate}
\item The vector $f_0$ has a place model at $a$, and for every function $\Phi$ that is Schwartz--Bruhat on
  $a\mathbb W$ there is a function $f$ with $f=T(\Phi)$ for every place model $T$ of $f_0$ at $a$; this $f$ is a
  theta function.
\item Let $(A,A',W)$ be spaces of a bottom triple for $(i,j)$ (Definition~\ref{th-bottom-triple}), with
  their standard lattices, and $t$ a tuple. The standard function at $\varepsilon_t$ on $\mathbb W(-A,W)$ is
  Schwartz--Bruhat on $\varepsilon_t\mathbb W(-A,W)$, and the standard function at $\bar\varepsilon_t$ on
  $\mathbb W(A',W)$ is Schwartz--Bruhat on $\bar\varepsilon_t\mathbb W(A',W)$.
\end{enumerate}
\end{lemma}
\leanref{Development/AbelianVarietyTheory/Proofs/TwistedOperators/FirstStandardPairMemberHasThetaFunctionStatement.lean}{GlobalTheory.T.first_standard_pair_member_has_theta_function}
\leanref{Development/AbelianVarietyTheory/Proofs/TwistedOperators/ExistsThetaFunctionStandardPairSecondStatement.lean}{GlobalTheory.T.exists_thetaFunction_standardPairSecond}
\leanref{Development/AbelianVarietyTheory/Proofs/TwistedOperators/FirstStandardPairMemberHasThetaFunction.lean}{StandardFirstTheta.GlobalTheory.CompanionTheory.StdFstTheta.stdFst_isSB}
\leanref{Development/AbelianVarietyTheory/Proofs/TwistedOperators/ExistsThetaFunctionStandardPairSecond.lean}{GlobalTheory.CompanionTheory.StdSndTheta.stdSnd_isSB}
\leanref{Development/AbelianVarietyTheory/Proofs/IwasawaDecomposition/EqZeroOrEqOfPlaceIdem.lean}{GlobalTheory.PlaceModelSum.placeModelExistsSum}
\begin{proof}
(2) Let $\Phi$ be the first standard function, with $\varepsilon=\varepsilon_t$, the finite set $S_t$, the
ideal $\mathfrak f=\mathfrak f_t$ and $\nu=\nu_{t,i,j}$. For $x\in\mathbb W$ let $N(x)$ be its matrix of
finite coordinates in the standard bases, so that $\Phi(x)\ne0$ only if $x\in\varepsilon\mathbb W$, $N(x)$ is
integral and its last $a'$ rows lie in $\mathfrak f$; then $\Phi(x)$ is $\nu(z)^{-1}$ if the finite idele
$z=1-\varepsilon+\varepsilon\det N_0(x)$, with $N_0(x)$ the top $a\times a$ block, is a unit at every finite
place, and $0$ otherwise. Choose the integers:
\begin{itemize}
\item $n\ge1$ with $\nu(z)=1$ for every finite idele $z$ that is a unit at every finite place, with
  archimedean components $1$ and $z-1\in n\hat{\mathcal O}_E$, which exists as $\nu$ is continuous;
\item $n_0\ge1$ in $\mathfrak f$, the norm of $\mathfrak f$;
\item $c\ge1$ lying in every place of $S_t$, the product of their norms;
\item $M\ge1$ with $MN(u)$ integral for every $u\in\Lambda$, and $K\ge1$ with $Kx\in\Lambda$ for every
  $x\in\varepsilon\mathbb W$ with $N(x)$ integral, which exist as $\Lambda$ is a lattice.
\end{itemize}
Put $k=Kn_0ncM$. By definition $\Phi$ vanishes off $\varepsilon\mathbb W$, and if $\Phi(x)\ne0$, then $N(x)$ is
integral, so $Kx\in\Lambda$ and $kx\in\Lambda$. Let $u\in\Lambda\cap\varepsilon\mathbb W$ and $x\in\mathbb W$. Then
$x+ku\in\varepsilon\mathbb W$ exactly when $x\in\varepsilon\mathbb W$, and $N(x+ku)=N(x)+n_0ncZ$ with
$Z=KMN(u)$ integral. So $N(x+ku)$ is integral exactly when $N(x)$ is, and as $n_0\in\mathfrak f$ its last
rows lie in $\mathfrak f$ exactly when those of $N(x)$ do. If they are integral, the determinants of the top
blocks differ by $ncw$ with $w$ integral, so the ideles $z$ and $z_1$ formed from the top blocks of $N(x)$ and
$N(x+ku)$ differ by $\varepsilon ncw$; as $c$ lies in the places of $S_t$, where $\varepsilon$ is supported,
one is a unit exactly when the other is, and then $z_1/z-1\in n\hat{\mathcal O}_E$, so $\nu(z_1)=\nu(z)$.
Hence $\Phi(x+ku)=\Phi(x)$.

The second standard function is treated in the same way, with $\bar\varepsilon_t$, the dual coordinates,
the first $a$ columns in $\mathfrak f$ and the last $a'\times a'$ block.

(1) Put $X=a\mathbb W$, $Y=\bar a\mathbb W$, $L=\Lambda\cap X$, $L'=\Lambda\cap Y$ and
$\chi(x,y)=\psi(\langle x,y\rangle)$ on $\mathbb W$, and let $\rho$ be the operators of
Definition~\ref{th-weil-datum} on the functions on $\mathbb W$. These data satisfy the hypotheses of
Lemma~\ref{asm-dependence-test-assembly-first-layer-ratio-inputs-standard-waves}.
\begin{enumerate}
\item[(a)] $\mathbb W$ is a module over $\mathbf A_E$, which contains $\mathbf Q$. So $kt=0$ implies $t=0$ for
  $k\ge1$, and every $x\in X$ is $k(k^{-1}x)$ with $k^{-1}x\in X$, as $a(k^{-1}x)=k^{-1}ax$; the same holds
  for $Y$. As $\psi$ is a character and the pairing is biadditive, $\chi$ is multiplicative in each variable.
\item[(b)] For $c\in\mathbf A_E$, $\langle cx,y\rangle=\langle x,\bar cy\rangle$, and $\langle
  y,x\rangle=-\langle x,y\rangle$, by the Hermitian symmetry of the Gram matrices and $\bar\xi_0=-\xi_0$. As
  $a\bar a=0$, $\langle x,x'\rangle=\langle ax,ax'\rangle=\langle x,\bar aax'\rangle=0$ for $x,x'\in X$, so
  $\rho(x)\rho(x')=\psi(\frac12\langle x,x'\rangle)\rho(x+x')=\rho(x+x')$; the same holds on $Y$. For $x\in X$
  and $y\in Y$, $\rho(x)\rho(y)=\psi(\frac12\langle x,y\rangle)\rho(x+y)$ and
  $\rho(y)\rho(x)=\psi(-\frac12\langle x,y\rangle)\rho(x+y)$, so $\rho(x)\rho(y)=\chi(x,y)\rho(y)\rho(x)$.
\item[(c)] The idempotent $a$ has archimedean component $0$ and finite components $0$ and $1$. So the
  elements of $X$ and $Y$ have archimedean component $0$, and $a\Lambda\subset\Lambda$ and $\bar
  a\Lambda\subset\Lambda$, as $\Lambda$ is stable under multiplication by the elements of $\mathbf A_E$ with
  integral finite components.
\item[(d)] By Lemma~\ref{th-lattice-signs}(1), $\chi(u,y)=1$ for $u\in L$ and $y\in L'$. Let $x\in X$ with
  $\chi(x,y)=1$ for every $y\in L'$. For $u\in\Lambda$, $\bar au\in L'$ by (c), so
  $\psi(\langle x,u\rangle)=\psi(\langle ax,u\rangle)=\psi(\langle x,\bar au\rangle)=1$; by
  Lemma~\ref{th-lattice-signs}(1), $x\in\Lambda$, so $x\in L$. In the same way, with $\langle
  y,u\rangle=-\langle u,y\rangle$, an element $y\in Y$ with $\chi(u,y)=1$ for every $u\in L$ lies in $L'$.
\item[(e)] $\Lambda$ is generated, as a module over the integral finite adeles of $F$, by finitely many
  elements, so for $k\ge1$ there is a finite set $R$ with $\Lambda=R+k\Lambda$. For $t\in L$, write $t=r+ku$
  with $r\in R$ and $u\in\Lambda$; then $t=at=ar+k(au)$ with $au\in L$. So $L=aR+kL$.
\item[(f)] $f_0$ is a theta function and a lattice vacuum for $\Lambda$. The idempotent $a$ is the sum of
  the idempotents $e_w$ of $\mathbf A_E$ at the places $w$ of a finite set (for $a=\bar\varepsilon$, the
  conjugates of those of $\varepsilon$). These satisfy $e_we_{w'}=0$ for $w\ne w'$, so $e_w=e_wa$ and
  $e_w\bar e_w=e_wa\bar a\bar e_w=0$; hence $e_w$ is a split idempotent
  (Lemma~\ref{is-eigenvalue-conjugate-split-form}). Let $u\in L$. Then $u=au=\sum_we_wu$, and $e_wu\in\Lambda$ lies in the half attached to $e_w$, so
  $\rho(e_wu)f_0=f_0$. The $e_wu$ lie in $X$, so $\rho(u)=\prod_w\rho(e_wu)$ by (b), and $\rho(u)f_0=f_0$.
  The same holds for $u\in L'$, with the $\bar e_w$ and $\bar a$.
\item[(g)] Every $y\in Y$ has a positive multiple in $L'$: it has archimedean component $0$; some positive
  integer multiple of it has integral finite components; and there is $K\ge1$ with $Kt\in\Lambda$ for
  every $t$ of archimedean component $0$ with integral finite components.
\end{enumerate}
A function is Schwartz--Bruhat on $a\mathbb W$ exactly when it is Schwartz--Bruhat in the sense of
Lemma~\ref{asm-dependence-test-assembly-first-layer-ratio-inputs-standard-waves}, since an $x\in X$ with
$kx\in\Lambda$ has $kx\in L$. So a place model of $f_0$ at $a$ is a map $T$ as in part (3) of that lemma for
$V$ the space of functions on $\mathbb W$, and part (3) gives one. Let $\Phi$ be Schwartz--Bruhat on
$a\mathbb W$. By part (2) of the lemma, $\Phi=\sum_ic_iw_{x_i,y_i}$ with $x_i\in X$ and $y_i\in Y$, and by
part (4) and (g), every place model $T$ of $f_0$ at $a$ satisfies
$T(\Phi)=\sum_ic_i\rho(x_i)\rho(y_i)f_0$. This function $f$ does not depend on $T$. It is a theta function:
$f_0$ is one, $\rho(x_i)$ and $\rho(y_i)$ carry theta functions to theta functions
(Lemma~\ref{is-eigenvalue-conjugate-theta-closed}(1)), and the conditions of
Definition~\ref{th-weil-datum}(2) are kept by sums, with the intersection of the two open subgroups, and
by scalar multiples.
% Lean (1): PlaceModelSum.placeModelExistsSum (isOrthSum_of_isIdemSum, with eq_or_mul_eq_zero for
% e_w e_w' = 0; IsOrthSum.exists_model for e and, by IsOrthSum.conj, for cA e) =
% PlaceModel.exists_placeModel (exists_intertwiner, Lemma -standard-waves(3)) with
% LatVac.latO_selfDual_approx_sign for (d), IsOrthSum.mem_of_dual (symplO_smul, symplO_skew),
% LatFinMod.latO_finite_mod and idem_smul_half for (e), IsOrthSum.heisO_eq for (f) (the lattice vacuum
% field hf.2 at each split idempotent, heisL_half), half_div, heisL_half and heisL_comm for (a) and (b),
% and LatticeStable.smul_mem_latO with AdIdem.idempotent_mem_intAE for (c). The uniqueness is
% PlaceDatum.exists_isThetaAt (placeModel_unique_of_dual, with exists_nsmul_mem_latO' for (g)), and the
% theta property SchrTheta.isTheta_placeModel (placeModel_induction with isTheta_heis_fin, the Lean's
% form of Lemma is-eigenvalue-conjugate-theta-closed(1) for elements of archimedean component 0, and
% WeilVacMove.isTheta_zero, isTheta_add, isTheta_smul). The Lean proves the uniqueness and the theta
% property by the same induction over the span of the waves; the book states that computation once, as
% Lemma -standard-waves(4). About 1,100 lines with Lemma -standard-waves.
% Lean: (2) StdFstTheta.stdFst_isSB with stdFst_shift (unitChar_congr: exists_isUnitDet and nu_eq through
% det_cong_int; inId_shift), CharFin.char_finUnit_finite_order for n, Ideal.absNorm_mem for n0, the
% product of the absNorm of the places for c, PlaceDatum.exists_den_frMapO for M and exists_nsmul_of_frMapO
% for K; StdSndTheta.stdSnd_isSB in the same way with dualCoord, exists_den_dualCoord and exists_mul_dualCoord.
% (1) PlaceModelSum.placeModelExistsSum (isOrthSum_of_isIdemSum, exists_model, for e and for cA e),
% PlaceDatum.exists_isThetaAt (and PlaceVec.exists_isThetaAt), and SchrTheta.isTheta_placeModel. The Lean's
% IsThetaAt at e with c = 1 is the book's f = T(Phi) for every place model. In
% StdFstTheta.exists_theta the space is -A with its lattice, self-dual by neg_gram.
\end{proof}

\begin{lemma}\label{asm-dependence-test-assembly-first-layer-ratio-inputs-weil-root}
Let $D$ be a definite-side datum with host field $E/F$, and let $W$ be a neighbourhood of $1$ in
$\mathbf A_E^\times$. Then there is a neighbourhood $N$ of $1$ in $\mathbf A_E^\times$ such that every
$u\in N$ with $u\bar u=1$ is $x/\bar x$ for some $x\in W$.
\end{lemma}
\leanref{Development/AbelianVarietyTheory/Proofs/UnitaryExistence/ExistsRootNearOne.lean}{GlobalTheory.Member.exists_root_near_one}

\begin{lemma}\label{asm-dependence-test-assembly-first-layer-ratio-inputs-weil-twist}
Let $D$ be a definite-side datum with host field $E/F$ and a standard character $\psi$
(Definition~\ref{th-weil-datum}), let $S$ and $T$ be Hermitian spaces of $D$, and let $\omega$ be a Weil
datum for $(S,T)$ with characters $(\chi_S,\chi_T)$. Let $\nu_S$ and $\nu_T$ be continuous unitary characters
of $\mathbf A_E^\times$ that are trivial on $E^\times$ and on $\mathbf A_F^\times$. Then there is a Weil
datum $\omega'$ for $(S,T)$ with characters $(\chi_S\nu_S,\chi_T\nu_T)$ such that
$\omega'(g,h)=\nu_T(x)^{-1}\nu_S(y)\,\omega(g,h)$ whenever $x$ and $y$ are determinant roots of $g$ and $h$
(Lemma~\ref{is-eigenvalue-conjugate-frame-det-roots}).
\end{lemma}
\leanref{Development/AbelianVarietyTheory/Proofs/DivisorProofs/ExistsWeilDatumTwistHeckeCharactersStatement.lean}{GlobalTheory.T.exists_weilDatum_twist_heckeCharacters}
\leanref{Development/AbelianVarietyTheory/Proofs/UnitaryExistence/ExistsUnitaryAdelicCharOfHeckeChar.lean}{GlobalTheory.Member.exists_muVal}
\begin{proof}
\emph{A character of a unitary group.} Let $V$ be a Hermitian space of $D$ and $\nu$ a continuous unitary
character of $\mathbf A_E^\times$, trivial on $E^\times$ and on $\mathbf A_F^\times$. Every
$g\in U(V)(\mathbf A)$ has a determinant root, and two of them differ by a factor in $\mathbf A_F^\times$
(Lemma~\ref{is-eigenvalue-conjugate-frame-root-values}(1) and (2)), on which $\nu$ is trivial. So
$a_\nu(g)=\nu(y)$, for any determinant root $y$ of $g$, is well defined, and $|a_\nu(g)|=1$. By
Lemma~\ref{is-eigenvalue-conjugate-frame-det-roots}(1), $a_\nu$ is multiplicative. A point
$\gamma\in U(V)(F)$ has a determinant root in $E^\times$
(Lemma~\ref{is-eigenvalue-conjugate-frame-root-values}(4)), so $a_\nu(\gamma)=1$. The homomorphism
$a_\nu$ is continuous at $1$, hence continuous: given a neighbourhood $U$ of $1$ in the unit circle, let
$N$ be the neighbourhood that
Lemma~\ref{asm-dependence-test-assembly-first-layer-ratio-inputs-weil-root} gives for $\nu^{-1}(U)$. The
$g$ with $\det g\in N$ form a neighbourhood of $1$, and for such $g$, as $\det(g)\overline{\det(g)}=1$,
there is $x\in\nu^{-1}(U)$ with $x=\det(g)\bar x$, a determinant root of $g$; so $a_\nu(g)=\nu(x)\in U$.

\emph{The twist.} Put $a=a_{\nu_T^{-1}}$ on $U(S)(\mathbf A)$, $b=a_{\nu_S}$ on $U(T)(\mathbf A)$ and
$\omega'(g,h)=a(g)b(h)\,\omega(g,h)$. As $a$ and $b$ are multiplicative, $\omega'$ is a homomorphism, and
$\omega'(g,h)=\nu_T(x)^{-1}\nu_S(y)\,\omega(g,h)$ for determinant roots $x$ of $g$ and $y$ of $h$. The
conditions of Definition~\ref{th-weil-datum}(3) hold for $\omega'$:
\begin{itemize}
\item $\chi_S\nu_S$ and $\chi_T\nu_T$ are trivial on $E^\times$ and restrict to $\eta^{\operatorname{rk}S}$
  and $\eta^{\operatorname{rk}T}$ on $\mathbf A_F^\times$, as $\nu_S$ and $\nu_T$ are trivial on
  $E^\times$ and on $\mathbf A_F^\times$;
\item $\omega'(g,h)$ is a multiple of $\omega(g,h)$, so it preserves theta functions and commutes with
  the operators $\rho(w_0)$ as $\omega(g,h)$ does;
\item $a(\gamma)=b(\gamma')=1$ for rational points, so $(\omega'(\gamma,\gamma')f)(0)=f(0)$;
\item $a$ and $b$ are continuous, so $(g,h)\mapsto(\omega'(g,h)f)(w)$ is continuous;
\item $|a(g)b(h)|=1$, so $\omega'(g,h)$ preserves the $L^2$-norms;
\item at a place of $F$ split in $E$, let $g$ and $h$ be supported there, with determinants $x$ and $y$
  there, the ideles with these determinants at the place selected by the split idempotent and $1$
  elsewhere. These are determinant roots of $g$ and $h$, by a computation with the idempotent. So
  $\omega'(g,h)$ acts on the invariant functionals through
  $a(g)b(h)\,r\chi_T(x)^{-1}\chi_S(y)=r\,(\chi_T\nu_T)(x)^{-1}(\chi_S\nu_S)(y)$.\qedhere
\end{itemize}
% Lean: WeilChange.exists_weilDatum_twist: a and b from Member.exists_muVal for (S, nu_T^-1) and (T, nu_S)
% (isDetRoot_laws: existence, uniqueness up to A_F^x, rationality and multiplicativity of roots; continuity
% from exists_root_near_one and continuous_det_matHom), then the seven clauses of IsWeilDatumO: hecke_mul and
% resEta_mul, thetaO_smul, heisO_smul with om_heis, arat and brat, acont and bcont, anorm and bnorm, and
% om_split with isDetAt_root and splitFacO.
\end{proof}

\begin{lemma}\label{asm-dependence-test-assembly-first-layer-ratio-inputs-weil-eta}
Let $D$ be a definite-side datum with host field $E/F$ and set $\mathcal T$ of embeddings, and let $\eta$ be the
quadratic character of $\mathbf A_F^\times$ attached to $E/F$. There is a continuous character $\tilde\eta$ of
$\mathbf A_E^\times$ with values in the unit circle, trivial on $E^\times$ and on the units of $E_w$ for every
finite place $w$, whose restriction to $\mathbf A_F^\times$ is $\eta$, and which is $z\mapsto\tau(z)/|\tau(z)|$ at the
place of each $\tau\in\mathcal T$.
\end{lemma}
\leanref{Development/AbelianVarietyTheory/Proofs/DivisorProofs/ExistsIsEtaTildeUnramified.lean}{GlobalTheory.WeilParts.exists_isEtaTilde_unramified}
\begin{proof}
Let $\tilde\eta$ be the Hecke character of Lemma~\ref{is-eigenvalue-conjugate-frame-phase-eta}. It is continuous and
trivial on $E^\times$; it is unramified at every finite place $w$, that is, trivial on the units of $E_w$; it
restricts to $\eta$ on $\mathbf A_F^\times$; and it has archimedean exponent $\frac12$ at every $\tau\in\mathcal T$,
that is, it is $z\mapsto\tau(z)/|\tau(z)|$ at the place of $\tau$ (Definition~\ref{th-first-lift-matching}(2)). Its
values lie in the unit circle: $|\tilde\eta|$ is a continuous homomorphism to $\mathbf R_{>0}$ that is trivial on
$E^\times$, on the units of every $E_w$ and on the archimedean components, so it factors through the ideal class group
of $E$, which is finite, and it is $1$.
% Lean: WeilParts.exists_isEtaTilde_unramified reorders the clauses of Mem.EtaTilde.exists_etaTilde (Lemma
% is-eigenvalue-conjugate-frame-phase-eta, which follows it) as IsEtaTilde: IsHeckeChar, UnramEverywhere, ResEta 1
% and HasInfExp 1/2 at every tau. The Lean's character takes values in the circle by its type.
\end{proof}

\begin{lemma}\label{asm-dependence-test-assembly-first-layer-ratio-inputs-weil-powers}
Let $D$ be a definite-side datum with host field $E/F$ and a standard character $\psi$
(Definition~\ref{th-weil-datum}), let $S$ and $T$ be Hermitian spaces of $D$, and let $\tilde\eta$ be as in
Lemma~\ref{asm-dependence-test-assembly-first-layer-ratio-inputs-weil-eta}. Then there is a Weil datum for
$(S,T)$ with characters $(\tilde\eta^{\operatorname{rk}S},\tilde\eta^{\operatorname{rk}T})$.
\end{lemma}
\leanref{Development/AbelianVarietyTheory/Proofs/DivisorProofs/ExistsWeilDatumPowersUnramifiedEtaTildeStatement.lean}{GlobalTheory.T.exists_weilDatum_powers_unramifiedEtaTilde}
\leanref{Development/AbelianVarietyTheory/Proofs/DivisorProofs/ExistsWeilDatumPowersUnramifiedEtaTilde.lean}{GlobalTheory.WeilProduct.exists_weilDatum_etaTilde}
\begin{proof}
Give $S$ and $T$ the self-dual lattices of their frames. For a Hermitian space $V$ let $G_f(V)$ and $G_\infty(V)$
be the subgroups of $U(V)(\mathbf A)$ of the elements with archimedean, respectively finite, component $1$; every
element of $U(V)(\mathbf A)$ is a product of commuting elements of the two. Write $p$ for a pair in
$U(S)(\mathbf A)\times U(T)(\mathbf A)$. The characters $\tilde\eta^{\operatorname{rk}S}$ and
$\tilde\eta^{\operatorname{rk}T}$ are continuous, unitary and trivial on $E^\times$, and restrict to
$\eta^{\operatorname{rk}S}$ and $\eta^{\operatorname{rk}T}$ on $\mathbf A_F^\times$. The construction has five
inputs:
\begin{enumerate}
\item[(a)] a homomorphism $\omega_f$ from $U(S)(\mathbf A)\times U(T)(\mathbf A)$ to the endomorphisms of the
  functions on $\mathbb W$ that preserves theta functions, is continuous and preserves their $L^2$-norms as in
  Definition~\ref{th-weil-datum}(3), is the identity on theta functions at the pairs in $G_\infty(S)\times
  G_\infty(T)$, satisfies $\omega_f(p)\rho(w_0)=\rho(hw_0g^{-1})\omega_f(p)$ on theta functions for
  $p=(g,h)\in G_f(S)\times G_f(T)$, and acts at the split places through the scalars of
  Definition~\ref{th-weil-datum}(3) for the characters $(\tilde\eta^{\operatorname{rk}S},
  \tilde\eta^{\operatorname{rk}T})$: the lattice model at the finite places;
\item[(b)] a homomorphism $\omega_\infty$ with the same properties for the archimedean places, with the roles
  of $G_f$ and $G_\infty$ exchanged and without the condition at the split places: the Fock model;
\item[(c)] $\omega_f(p)\omega_\infty(q)f=\omega_\infty(q)\omega_f(p)f$ for all $p$, $q$ and theta functions $f$;
\item[(d)] the continuity of $p\mapsto(\omega_f(p)\omega_\infty(p)f)(w)$ for theta functions $f$;
\item[(e)] the invariance of the theta distribution: $(\omega_f(p)\omega_\infty(p)f)(0)=f(0)$ for $p$ in
  $U(S)(F)\times1$ and in $1\times U(T)(F)$ and every theta function $f$.
\end{enumerate}
Given these, let $\omega(p)$ be $\omega_f(p)\omega_\infty(p)$ on the theta functions, extended by the identity on a
complement of the space of theta functions. By (c), $\omega$ is a homomorphism. It preserves theta functions
and their norms, and is continuous by (d). For $p=(g,h)$, write $g=ak$ and $h=bl$ with $a,b$ archimedean and
$k,l$ finite. Then $\omega(p)=\omega_\infty(a,b)\omega_f(k,l)$ on theta functions, and the Heisenberg condition
for $p$ follows from those of (a) at $(k,l)$ and of (b) at $(a,b)$, with $hw_0g^{-1}=b(lw_0k^{-1})a^{-1}$. The
invariance at the rational points is (e), for $(\gamma,\gamma')=(\gamma,1)(1,\gamma')$. A pair supported at a
split place lies in $G_f(S)\times G_f(T)$, where $\omega=\omega_f$, so the condition at the split places is that
of (a). So $\omega$ is a Weil datum with characters $(\tilde\eta^{\operatorname{rk}S},
\tilde\eta^{\operatorname{rk}T})$.
\openstep{The lattice model $\omega_f$ and the Fock model $\omega_\infty$ with the properties (a) and (b), their
commutation (c), the continuity (d) of their product and the invariance (e) of the theta distribution under the
rational points of the two unitary groups are not yet written out.}
% Lean: WeilProduct.exists_weilDatum_etaTilde: (a) FinWeilO.exists_isFinWeil (IsFinWeil: pres, cont, unit,
% triv, Heisenberg2, split, inert), (b) RealWeil.exists_isInfWeil (IsInfWeil: pres, cont, unit, triv,
% Heisenberg2, vac), (c) fin_inf_comm (through the Schur lemma for theta functions,
% T.schurLemma_thetaFunctions_orientedPair, and the lattice representation exists_isLatRep), (d)
% continuous_product, (e) T.thetaDistribution_automorphy_source_group and
% T.thetaDistribution_invariant_rationalTarget; exists_parts (UnitarySplit.real_finite_split), the
% extension exists_extend from thetaSub, mem_GfC_of_suppAt for the split clause. The Lean's statement also
% gives the inert clause FixesVacAt and VacCharPlus, not used here. About 25,000 lines at depth 5, almost all
% in (a), (b) and (e).
\end{proof}

\begin{lemma}\label{asm-dependence-test-assembly-first-layer-ratio-inputs-weil}
Let $D$ be a definite-side datum with host field $E/F$ and a standard character $\psi$
(Definition~\ref{th-weil-datum}), and let $S$ and $T$ be Hermitian spaces of $D$. Let $\chi_S$ and $\chi_T$
be continuous unitary characters of $\mathbf A_E^\times$, trivial on $E^\times$, whose restrictions to
$\mathbf A_F^\times$ are $\eta^{\operatorname{rk}S}$ and $\eta^{\operatorname{rk}T}$, where $\eta$ is the
quadratic character of $E/F$. Then there is a Weil datum for $(S,T)$ with characters $(\chi_S,\chi_T)$.
\end{lemma}
\leanref{Development/AbelianVarietyTheory/Proofs/TwistedOperators/RatInOfDefinite.lean}{GlobalTheory.ConstRatio.exists_isWeil}
\leanref{Development/AbelianVarietyTheory/Proofs/DivisorProofs/ExistsWeilDatumOfSplittingCharactersStatement.lean}{GlobalTheory.T.exists_weilDatum_of_splittingCharacters}
\leanref{Development/AbelianVarietyTheory/Proofs/DivisorProofs/ExistsWeilDatumOfSplittingCharacters.lean}{GlobalTheory.WeilExist.exists_isWeilDatumO}
\leanref{Development/AbelianVarietyTheory/Proofs/DivisorProofs/ExistsIsEtaTildeUnramified.lean}{GlobalTheory.WeilParts.exists_isEtaTilde_unramified}
\begin{proof}
Let $\tilde\eta$ be the character of
Lemma~\ref{asm-dependence-test-assembly-first-layer-ratio-inputs-weil-eta}: a Hecke character of $E$,
unramified at every finite place, whose restriction to $\mathbf A_F^\times$ is $\eta$. By
Lemma~\ref{asm-dependence-test-assembly-first-layer-ratio-inputs-weil-powers} there is a Weil datum
$\tilde\omega$ for $(S,T)$ with characters $(\tilde\eta^{\operatorname{rk}S},\tilde\eta^{\operatorname{rk}T})$.
The characters
$\nu_S=\tilde\eta^{-\operatorname{rk}S}\chi_S$ and $\nu_T=\tilde\eta^{-\operatorname{rk}T}\chi_T$ are
continuous, unitary and trivial on $E^\times$, and on $\mathbf A_F^\times$ they are
$\eta^{-\operatorname{rk}S}\eta^{\operatorname{rk}S}=1$ and $1$. By
Lemma~\ref{asm-dependence-test-assembly-first-layer-ratio-inputs-weil-twist}, twisting $\tilde\omega$ by
$\nu_S$ and $\nu_T$ gives a Weil datum with characters
$(\tilde\eta^{\operatorname{rk}S}\nu_S,\tilde\eta^{\operatorname{rk}T}\nu_T)=(\chi_S,\chi_T)$.
% Lean: ConstRatio.exists_isWeil (the sign of S through ofSgn) = T.exists_weilDatum_of_splittingCharacters =
% WeilExist.exists_isWeilDatumO: the eta-tilde from WeilParts.exists_isEtaTilde_unramified
% (Mem.EtaTilde.exists_etaTilde: EtaT.etaT_exists of the host and CharTransport.exists_chrEq), the Weil datum
% om from T.exists_weilDatum_powers_unramifiedEtaTilde (WeilProduct.exists_weilDatum_etaTilde:
% FinWeilO.exists_isFinWeil, RealWeil.exists_isInfWeil, fin_inf_comm, exists_extend,
% thetaDistribution_automorphy_source_group, thetaDistribution_invariant_rationalTarget; about 25,000 lines
% at depth 5), with the self-dual lattices PlaceModel.isSelfDualC_latFrame, then
% T.exists_weilDatum_twist_heckeCharacters (Lemma -weil-twist) with isHeckeChar_div_pow and resEta_div_pow.
% The Lean's existence proof does not use the Weil representation of Theorem lit-bundle-theta.
\end{proof}

\begin{lemma}\label{asm-dependence-test-assembly-first-layer-ratio-inputs-factorization}
Assume the results of Theorem~\ref{lit-bundle-theta} and the theta hypotheses
(Definition~\ref{th-theta-hypotheses}). For all sufficiently large data
(Definition~\ref{th-large-data}) the following holds. Let a grid Weil datum with standard character,
$0\le i,j<m$ and a bottom triple $(A,A',W)$ for $(i,j)$ (Definition~\ref{th-bottom-triple}) be given.
Then there are a complex vector space $\Pi$, linear maps $p$ and $p'$ to $\Pi$ from the polynomials in
the creation operators of $(-A,W)$ and of $(A',W)$, with $p$ surjective, and for each tuple $t$ a
linear map $J_t$ from $\Pi$ to functions on $U(W)(\mathbf A)$ that is $0$ or injective and a number
$\lambda(t)$, such that $I_t(P,\cdot)=\lambda(t)J_t(p(P))$ and $I'_t(P',\cdot)=J_t(p'(P'))$ for all $P$
and $P'$.
\end{lemma}
\leanref{Development/AutomorphicArithmetic/ThetaCorrespondence/ThetaPairLifts/LocalTypeDictionary/BottomTripleIntegralFactorizationOfReal.lean}{GlobalTheory.T.bottom_triple_integral_factorization_of_real}
\leanref{Development/AutomorphicArithmetic/ThetaCorrespondence/ThetaPairLifts/LocalTypeDictionary/BottomTripleIntegralFactorizationOfReal.lean}{GlobalTheory.BotFac.bottom_triple_integral_factorization_real}
\begin{proof}
Take the data large enough for Lemma~\ref{th-first-lift-inputs}(5) and
Lemma~\ref{th-first-lift-finite}, and fix a grid Weil datum, $(i,j)$ and a bottom triple
$(A,A',W)$, with the characters $\alpha_t$ and $\beta$ of the square $(i,j)$.

\emph{The archimedean data.} By Lemma~\ref{th-first-lift-inputs}(5) there are archimedean data on $A$,
$A'$ and $W$ and $a_\alpha,a_\beta\colon\mathcal T\to\mathbf Q$ such that $A$, $A'$ and $W$ are totally
positive definite and, for every tuple $t$, $(A,A',W;\alpha_t,\beta)$ is a square triple
(Definition~\ref{th-first-lift-matching}) in which $\alpha_t$ and $\beta$ have archimedean exponents
$a_\alpha$ and $a_\beta$, $f_t$ and $f'_t$ are archimedean vacua, and
$I_t(P,\cdot)=I_{\omega_t}(f_t\otimes P)$ and $I'_t(P',\cdot)=I_{\omega'_t}(f'_t\otimes P')$. Let $\Pi$, $q$
and $q'$ be the space and the surjective maps that Lemma~\ref{th-first-lift-factor}(2) gives for
$a_\alpha$ and $a_\beta$; they do not depend on $t$.

\emph{Three facts at a tuple $t$.} The square triple $(A,A',W;\alpha_t,\beta)$ has a matching by
Lemma~\ref{th-first-lift-factor}(2), and so, by Lemma~\ref{th-first-lift-factor}(1), a matching
$(\omega_t,\omega'_t,\Pi_t,\varrho,p_t,p'_t)$ with the Weil data of the bottom triple. For it,
Lemma~\ref{th-first-lift-factor}(2) gives a space $\Pi_f$, linear maps $r$ and $r'$ and an
isomorphism $e\colon\Pi\otimes\Pi_f\to\Pi_t$ with $p_t(f\otimes P)=e(q(P)\otimes r(f))$ and
$p'_t(f'\otimes P')=e(q'(P')\otimes r'(f'))$ for all archimedean vacua $f$ and $f'$. As $A'$ is totally
positive definite, Lemma~\ref{th-first-lift-comparison} gives $c\ne0$ with
$I_{\omega_t}(\Phi)=c\,I_{\omega'_t}(\Psi)$ whenever $p_t(\Phi)=p'_t(\Psi)$, and $I_{\omega'_t}(\Psi)=0$
whenever $p'_t(\Psi)=0$.
\begin{enumerate}
\item[(a)] If $q'(P')=0$, then $p'_t(f'_t\otimes P')=e(0)=0$, so $I'_t(P',\cdot)=0$.
\item[(b)] Either $I'_t(P',\cdot)=0$ for every $P'$, or $I'_t(P',\cdot)=0$ implies $q'(P')=0$. Indeed, let
  $I'_t(P',\cdot)=0$ with $q'(P')\ne0$, and let $P''$ be a polynomial. If $r'(f'_t)=0$, then
  $p'_t(f'_t\otimes P'')=0$ and $I'_t(P'',\cdot)=0$. Otherwise
  $p'_t(f'_t\otimes P')=e(q'(P')\otimes r'(f'_t))\ne0$, and as $A$, $A'$ and $W$ are totally positive
  definite, Lemma~\ref{th-first-lift-irreducible} gives $I_{\omega'_t}=0$ on $\mathcal A(A',W)$, so
  $I'_t(P'',\cdot)=0$.
\item[(c)] There is $\lambda(t)$ with $I_t(P,\cdot)=\lambda(t)I'_t(P',\cdot)$ whenever $q(P)=q'(P')$. Indeed,
  Lemma~\ref{th-first-lift-finite} gives $\kappa$ with $r(f_t)=\kappa\,r'(f'_t)$. If $q(P)=q'(P')$, then
  $p_t(f_t\otimes P)=e(q'(P')\otimes\kappa\,r'(f'_t))=p'_t(\kappa f'_t\otimes P')$, so
  $I_t(P,\cdot)=c\kappa\,I'_t(P',\cdot)$; put $\lambda(t)=c\kappa$.
\end{enumerate}

\emph{The maps $J_t$.} The map $P'\mapsto I'_t(P',\cdot)$ is linear, and by (a) it vanishes on the
kernel of $q'$. As $q'$ is surjective, there is a linear map $J_t$ on $\Pi$ with
$J_t(q'(P'))=I'_t(P',\cdot)$ for every $P'$. By (b), $J_t$ is $0$ or injective. For $P$, choose $P'$ with
$q'(P')=q(P)$; by (c), $I_t(P,\cdot)=\lambda(t)I'_t(P',\cdot)=\lambda(t)J_t(q(P))$. So $\Pi$, $p=q$,
$p'=q'$, the $J_t$ and the $\lambda(t)$ have the stated properties.
% Lean: BotFac.bottom_triple_integral_factorization_real: definiteBottom_dictionary_frames_exponents_kernel
% (Lemma th-first-lift-inputs(5): FA, FA', FW, a, b, positivity hdA, hdef, hdW and hD), twisted_local_factors_match
% (Lemma th-first-lift-factor(2): Pr, q, q', hq, hq', hLMt), triple_weil_transfer_matching (Lemma
% th-first-lift-factor(1)), twisted_square_theta_integrals_eq with theta_integral_vanishes_on_finite_relations
% (the steps same and kill; Lemma th-first-lift-comparison), theta_lift_vanishes_of_irreducible_defect (Lemma
% th-first-lift-irreducible) and theta_bottom_finite_place_matching (Lemma th-first-lift-finite, hF); then J t
% through liftQ on the kernel of q' and quotKerEquivOfSurjective. It is the assembly of
% BottomFactor.bottom_theta_norm_factorization.
\end{proof}

\begin{definition}\label{asm-dependence-test-assembly-first-layer-ratio-inputs-towers-stages}
Let $D$ be a definite-side datum with a level system, a theta grid and a grid Weil datum with character
$\psi$ (Definition~\ref{th-grid}), and let $0\le i,j<m$. Put $p=\lfloor(m-j-1)/2\rfloor$,
$q=\lfloor(m-j)/2\rfloor$ and $e=(\operatorname{rk}V_{i+1,j+1}-\operatorname{rk}V_{i,j})/2$, an integer.
\begin{enumerate}
\item A \emph{triple} for $(i,j)$ consists of Hermitian spaces $A$, $A'$ and $W$ over $E$ with
  $\operatorname{rk}A+\operatorname{rk}A'=\operatorname{rk}W$, with self-dual lattices and archimedean data,
  and the data (1)--(4) of Definition~\ref{th-bottom-triple} for them, with Fock vacuum vectors for these
  lattices and archimedean data; it has the theta integrals $I_t$ and $I'_t$ defined there. A bottom triple
  for $(i,j)$ is a triple.
\item A triple is a \emph{stage} $(p',q')$, for integers $0\le p'\le p$ and $0\le q'\le q$, if
  $\operatorname{rk}A=\operatorname{rk}V_{i+1,j+1}-2(p-p')$, $\operatorname{rk}A'=\operatorname{rk}V_{i,j}-2(q-q')$,
  and $A$, $A'$ and $W$ have signatures $(\operatorname{rk}A-p',p')$, $(\operatorname{rk}A'-q',q')$ and
  $(\operatorname{rk}W-p'-q',p'+q')$ at every embedding.
\item \emph{Constants} of a triple are numbers $c(t)$, one for each tuple $t$ and not all $0$, for which there
  are polynomials $P$ and $P'$ in the creation operators such that for every $t$ the integrals $I_t(P,h)$ and
  $I'_t(P',h)$ converge for every $h$, $I'_t(P',\cdot)\ne0$, and $I_t(P,\cdot)=c(t)I'_t(P',\cdot)$.
\item A \emph{factorization} of a triple consists of nonzero numbers $\lambda(t)$, a complex vector space $\Pi$,
  linear maps $p_0$ and $p'_0$ to $\Pi$ from the polynomials in the creation operators of $(-A,W)$ and of
  $(A',W)$, with $p_0$ surjective, and injective linear maps $J_t$ from $\Pi$ to functions on $U(W)(\mathbf A)$,
  such that the integrals $I_t(P,h)$ converge and $I_t(P,\cdot)=\lambda(t)J_t(p_0(P))$ for all $P$,
  $I'_t(P',\cdot)=J_t(p'_0(P'))$ for all $P'$ whose integrals converge, and $p'_0(P')\ne0$ for some such $P'$.
\end{enumerate}
\end{definition}
\leanref{Development/AbelianVarietyTheory/Proofs/DivisorProofs/NegIdxOfIsStd.lean}{GlobalTheory.SqTower.IsStageC}
\leanref{Development/AbelianVarietyTheory/BlockTheory/TwistedDiscreteSeries/UnitCharTriangle.lean}{GlobalTheory.IndefiniteRowSideFullHecke.UnitCharTriangle.Tri.IsConst}
\leanref{Development/AbelianVarietyTheory/Proofs/TwistedOperators/RatInOfUnit.lean}{GlobalTheory.ConstRatio.Factor}
\leanref{Development/AbelianVarietyTheory/Statements/PlaceShiftStatements/NormSmall.lean}{GlobalTheory.SqTop.pp}
\leanref{Development/AbelianVarietyTheory/Statements/PlaceShiftStatements/NormSmall.lean}{GlobalTheory.SqTop.qq}
\leanref{Development/AbelianVarietyTheory/Statements/PlaceShiftStatements/NormSmall.lean}{GlobalTheory.SqTop.eTop}
\leanref{Development/AbelianVarietyTheory/Statements/PlaceShiftStatements/AuxiliaryPrimeNeTypePrime.lean}{GlobalTheory.SqTop.Model.tower_numbers}

\begin{lemma}\label{asm-dependence-test-assembly-first-layer-ratio-inputs-towers-inputs}
Assume the results of Theorem~\ref{lit-bundle-theta} and the theta hypotheses
(Definition~\ref{th-theta-hypotheses}). For all sufficiently large data (Definition~\ref{th-large-data}), with
a theta grid with base tuple $\mu_0$, the following hold for every grid Weil datum with standard character,
every grid choice on all tuples and every real vector, whose numbers at $\mu_0$ lie in the coefficient field
$L$ of $D$, with its embedding $\sigma$, and for all $0\le i,j<m$, with $p$, $q$ and $e$ as in
Definition~\ref{asm-dependence-test-assembly-first-layer-ratio-inputs-towers-stages}. Let $\Lambda_t$ be an
$L$-function of $\nu_{t,i,j}$ for every tuple $t$.
\begin{enumerate}
\item (Top.) There are a stage $(p,q)$ with constants $c_T(t)$ and nonzero numbers $\kappa(t)$ with
  $\kappa(t)/\kappa(\mu_0)\in\sigma(L)$ for every $t$, such that every constant $c$ of a nonzero square $(i,j)$ at
  a tuple $t$ is $\kappa(t)c_T(t)$.
\item (Factorization.) Every stage has a factorization.
\item (Small step.) Let $X$ be a stage $(p',q'+1)$ with a factorization with numbers $\lambda(t)$, and with
  $\operatorname{rk}A=\operatorname{rk}A'+2e'$. Then there are a stage $(p',q')$ with constants $c(t)$ and
  nonzero numbers $y(t)$ with $y(t)/y(\mu_0)\in\sigma(L)$ and $\lambda(t)\Lambda_t(e'+1)=y(t)c(t)$ for every $t$.
\item (Large step.) Let $X$ be a stage $(p'+1,0)$ with a factorization with numbers $\lambda(t)$, and with
  $\operatorname{rk}A=\operatorname{rk}A'+2e'$. Then there are a stage $(p',0)$ with constants $c(t)$ and nonzero
  numbers $y(t)$ with $y(t)/y(\mu_0)\in\sigma(L)$ and $\lambda(t)=y(t)\Lambda_t(e')c(t)$ for every $t$.
\item (Bottom.) Let $X$ be a bottom triple for $(i,j)$ and $Y$ a stage $(0,0)$ with constants $c(t)$. Then there
  are nonzero numbers $v(t)$ with $v(t)/v(\mu_0)\in\sigma(L)$ such that the $v(t)c(t)$ are constants of $X$.
\end{enumerate}
\end{lemma}
\leanref{Development/AutomorphicArithmetic/ThetaCorrespondence/ThetaPairLifts/LocalTypeDictionary/WittTowerTopStageThetaNormFormula.lean}{GlobalTheory.T.witt_tower_top_stage_theta_norm_formula}
\leanref{Development/AutomorphicArithmetic/ThetaCorrespondence/ThetaPairLifts/LocalTypeDictionary/ThetaTowerStageNormFactorization.lean}{GlobalTheory.T.exists_factor_of_isStageC}
\leanref{Development/AutomorphicArithmetic/ThetaCorrespondence/ThetaPairLifts/LocalTypeDictionary/WittTowerNormStepOfSmallRank.lean}{GlobalTheory.T.witt_tower_norm_step_of_small_rank}
\leanref{Development/AutomorphicArithmetic/ThetaCorrespondence/ThetaPairLifts/LocalTypeDictionary/WittTowerNormStepOfLargeRank.lean}{GlobalTheory.T.witt_tower_norm_step_of_large_rank}
\leanref{Development/AutomorphicArithmetic/ThetaCorrespondence/ThetaPairLifts/LocalTypeDictionary/WittTowerConstantRatioEqBottom.lean}{GlobalTheory.T.witt_tower_constant_ratio_eq_bottom}
\begin{proof}
\openstep{The five inputs of the tower are not yet written out: the theta norm formula at the top stage, which
compares the constant of a square with that of the top stage; the factorization of the theta integrals of every
stage; the norm steps of small and of large rank, which pass from one stage to the next by adding a hyperbolic
plane and produce one $L$-value; and the comparison of the constants of the stage $(0,0)$ with those of a bottom
triple.}
% Lean: (1) SqTower.theta_tower_top_stage_norm_eq, (2) StageFactor.theta_tower_stage_norm_factorization, (3)
% SqStep.theta_tower_norm_step_small_rank, (4) SqStepBig.theta_tower_norm_step_large_rank, (5)
% SqTower.theta_tower_bottom_constant_ratio, each under ThetaLocalInputs, ThetaIntegralsAgree and LocalMatching
% and inside GBig. The Lean gives more than the book states: in (1) kappa t = embLL zeta * kappa mu0 with zeta an
% IsAwayUnit (Qn G t), in (3) and (4) y t and y mu0 related by Square.StepRatio, in (5) v t = embLL zeta * v mu0
% with zeta an IsAwayUnit; the book keeps only that the ratios lie in sigma(L), which is what the assembly uses.
% A stage is IsStageC (NegIdx for the signatures; the triple data are Tri.Is with the characters of the square
% and the places locPl G), constants are Tri.IsConst and a factorization ConstRatio.Factor. About 8,100 lines of
% Lean proofs at depth 5, counted from SqTower.tower_alt_bis; almost all of it is in these five inputs.
\end{proof}

\begin{lemma}\label{asm-dependence-test-assembly-first-layer-ratio-inputs-towers-descent}
In the situation of Lemma~\ref{asm-dependence-test-assembly-first-layer-ratio-inputs-towers-inputs}, assume
its conclusions (2), (3) and (4), and let $\Lambda_t$ be an $L$-function of $\nu_{t,i,j}$ for every tuple $t$.
\begin{enumerate}
\item If a triple has a factorization with numbers $\lambda(t)$ and constants $c(t)$, then there is $\kappa\ne0$
  with $c(t)=\kappa\lambda(t)$ for every $t$.
\item Let $X$ be a stage $(p',q')$ with constants $c(t)$, with $\operatorname{rk}A=\operatorname{rk}A'+2e'$ and
  $p'\le e'$. Then there are a stage $(0,0)$ with constants $c_Y(t)$, nonzero numbers $u_k(t)$ for $0\le
  k<p'+q'$ with $u_k(t)/u_k(\mu_0)\in\sigma(L)$, and $\kappa\ne0$ such that for every tuple $t$
  \[
   c(t)\prod_{k=0}^{q'-1}\Lambda_t(e'+k+1)=\kappa\Bigl(\prod_{k=0}^{p'+q'-1}u_k(t)\Bigr)c_Y(t)
   \prod_{k=0}^{p'-1}\Lambda_t(e'+q'-k).
  \]
\end{enumerate}
\end{lemma}
\leanref{Development/AbelianVarietyTheory/Proofs/TwistedOperators/RatInOfUnit.lean}{GlobalTheory.ConstRatio.Factor.eq_mul}
\leanref{Development/AbelianVarietyTheory/Proofs/DivisorProofs/TowerOfThetaMatching.lean}{GlobalTheory.SqTower.descent}
\begin{proof}
(1) Let $P$ and $P'$ be the polynomials of the constants. For every $t$, the integrals of $P'$ converge, so
$\lambda(t)J_t(p_0(P))=I_t(P,\cdot)=c(t)I'_t(P',\cdot)=c(t)J_t(p'_0(P'))$, and as $J_t$ is injective,
$\lambda(t)p_0(P)=c(t)p'_0(P')$. Here $p'_0(P')\ne0$, as $I'_t(P',\cdot)=J_t(p'_0(P'))\ne0$. So
$\kappa=c(t)/\lambda(t)$ does not depend on $t$, and $\kappa\ne0$ as some $c(t)\ne0$.

(2) By induction on $q'$, and for $q'=0$ on $p'$. If $p'=q'=0$, take $X$ itself, $c_Y=c$ and $\kappa=1$.

\emph{The case $(p'+1,0)$.} By (2) of the lemma, $X$ has a factorization with numbers $\lambda(t)$, and by part
(1), $c(t)=\kappa_1\lambda(t)$ with $\kappa_1\ne0$. The large step gives a stage $Y$ $(p',0)$ with constants
$c'(t)$ and numbers $y(t)$ with $\lambda(t)=y(t)\Lambda_t(e')c'(t)$. Its ranks satisfy $\operatorname{rk}B=
\operatorname{rk}A-2$ and $\operatorname{rk}B'=\operatorname{rk}A'$, so $\operatorname{rk}B=\operatorname{rk}B'+2(e'-1)$,
and $p'\le e'-1$. By induction, $c'(t)=\kappa(\prod_{k<p'}u_k(t))c_Y(t)\prod_{k<p'}\Lambda_t(e'-1-k)$. Then
\[
 c(t)=\kappa_1\kappa\,y(t)\Bigl(\prod_{k<p'}u_k(t)\Bigr)c_Y(t)\,\Lambda_t(e')\prod_{k<p'}\Lambda_t(e'-1-k),
\]
which is the assertion with the numbers $y,u_0,\dots,u_{p'-1}$ and the constant $\kappa_1\kappa$.

\emph{The case $(p',q'+1)$.} As before $c(t)=\kappa_1\lambda(t)$. The small step gives a stage $Y$ $(p',q')$ with
constants $c'(t)$ and numbers $y(t)$ with $\lambda(t)\Lambda_t(e'+1)=y(t)c'(t)$. Now $\operatorname{rk}B=
\operatorname{rk}A$ and $\operatorname{rk}B'=\operatorname{rk}A'-2$, so $\operatorname{rk}B=\operatorname{rk}B'+2(e'+1)$,
and $p'\le e'+1$. By induction,
$c'(t)\prod_{k<q'}\Lambda_t(e'+k+2)=\kappa(\prod_{k<p'+q'}u_k(t))c_Y(t)\prod_{k<p'}\Lambda_t(e'+1+q'-k)$. So
\[
 c(t)\prod_{k<q'+1}\Lambda_t(e'+k+1)=\kappa_1\lambda(t)\Lambda_t(e'+1)\prod_{k<q'}\Lambda_t(e'+k+2)
 =\kappa_1\kappa\,y(t)\Bigl(\prod_{k<p'+q'}u_k(t)\Bigr)c_Y(t)\prod_{k<p'}\Lambda_t(e'+(q'+1)-k),
\]
which is the assertion with the numbers $y,u_0,\dots,u_{p'+q'-1}$.
% Lean: Factor.eq_mul (key: lam t p P = c t p' P' through hJ, hne, then the ratio); SqTower.descent, by
% induction on q and then on p, with cons y u for the new numbers (the Lean's StepRatio, of which the book keeps
% that u k t / u k mu0 lies in sigma(L)) and the rank relations from IsStageC by omega. About 150 lines.
\end{proof}

\begin{lemma}\label{asm-dependence-test-assembly-first-layer-ratio-inputs}
Assume the results of Theorem~\ref{lit-bundle-theta}, the theta hypotheses
(Definition~\ref{th-theta-hypotheses}) and the functional equation of
Theorem~\ref{asm-dependence-inputs}. For all sufficiently large data
(Definition~\ref{th-large-data}), with a theta grid with base tuple $\mu_0$, the following hold for all
$0\le i,j<m$. Write $e_0=e_0(i,j)$ and $h_1=h_1(i,j)$ (Definition~\ref{th-bottom-triple}), and let
$p$, $q$ and $e$ be the integers of
Definition~\ref{asm-dependence-test-assembly-first-layer-ratio-inputs-towers-stages}; they satisfy
$e_0+p=e+q$ and $e+q\le h_1$.
\begin{enumerate}
\item For every tuple $t$ there is an $L$-function $\Lambda$ of $\nu_{t,i,j}$ with $\Lambda(s)\ne0$ for
  the integers $1-e_0\le s\le h_1$; and each of $\nu_{\mu_0,i,j}$, $\nu_{\mu_0,i,j}\delta$ and
  $\nu_{\mu_0,i,j}\delta'$, with the characters $\delta$ and $\delta'$ of the level system
  (Theorem~\ref{th-first-lift-nonzero}), has an $L$-function that does not vanish at the integers
  $s\ge1$.
\item For every grid Weil datum with standard character there is a bottom triple for $(i,j)$.
\item For every grid Weil datum with standard character, every bottom triple $(A,A',W)$ for $(i,j)$
  has a \emph{factorization}:
  nonzero numbers $\lambda(t)$, a complex vector space $\Pi$, linear maps $p_0$ and $p'_0$ to $\Pi$ from
  the polynomials in the creation operators of $(-A,W)$ and of $(A',W)$, with $p_0$ surjective, and
  injective linear maps $J_t$ from $\Pi$ to functions on $U(W)(\mathbf A)$, such that
  $I_t(P,\cdot)=\lambda(t)J_t(p_0(P))$ for all $P$, $I'_t(P',\cdot)=J_t(p'_0(P'))$ for all $P'$ whose
  integrals converge, and $p'_0(P')\ne0$ for some such $P'$.
\item (Towers.) Let a grid Weil datum with standard character, a grid choice on all tuples and a real
  vector be given whose numbers at $\mu_0$ lie in the coefficient field $L$ of $D$. Let $t$ be a tuple,
  let the squares $(i,j)$ at $t$ and at $\mu_0$ be nonzero with constants $c$ and $c_0$, let $(A,A',W)$ be
  a bottom triple for $(i,j)$ and $\lambda(t')$ numbers, not all zero, such that for some $P,P'$ the
  integrals $I_{t'}(P,\cdot)$ and $I'_{t'}(P',\cdot)$ converge, $I'_{t'}(P',\cdot)\ne0$ and
  $I_{t'}(P,\cdot)=\lambda(t')I'_{t'}(P',\cdot)$ at every tuple $t'$, and let $\Lambda$
  and $\Lambda_0$ be $L$-functions of $\nu_{t,i,j}$ and $\nu_{\mu_0,i,j}$. Then there are numbers
  $u_1,\dots,u_{p+q}$, each the image of an element of $L$, and $\zeta\in L$ with
  \[
   \frac c{c_0}\prod_{k=0}^{q-1}\frac{\Lambda(e+k+1)}{\Lambda_0(e+k+1)}
   =\zeta\Bigl(\prod_ku_k\Bigr)\frac{\lambda(t)}{\lambda(\mu_0)}
   \prod_{k=0}^{p-1}\frac{\Lambda(e+q-k)}{\Lambda_0(e+q-k)}.
  \]
\item Let a bottom triple $(A,A',W)$ for $(i,j)$ have a factorization with numbers $\lambda(t)$, and
  let $L$ be a field with an embedding $\sigma\colon L\to\mathbf C$. Suppose that the \emph{bottom sums}
  of the triple \emph{lie in} $L$: there are nonzero $\Omega,\Omega'\in\mathbf C$ with
  $I_t(P,h)\in\Omega\,\sigma(L)$ and $I'_t(P',h)\in\Omega'\sigma(L)$ for every tuple $t$, all monomials
  $P$ and $P'$ in the creation operators and every $h\in U(W)(\mathbf A)$ with archimedean components
  $1$. Then for every tuple $t$ there is $y\in L^\times$ with $\sigma(y)=\lambda(t)/\lambda(\mu_0)$.
\end{enumerate}
\end{lemma}
\leanref{Development/AbelianVarietyTheory/Statements/PlaceShiftStatements/NormSmall.lean}{GlobalTheory.SqTop.LNe}
\leanref{Development/AbelianVarietyTheory/Statements/PlaceShiftStatements/NormSmall.lean}{GlobalTheory.SqTop.TriEx}
\leanref{Development/AbelianVarietyTheory/Statements/PlaceShiftStatements/NormSmall.lean}{GlobalTheory.SqTop.Tower}
\leanref{Development/AbelianVarietyTheory/Proofs/TwistedOperators/RatInOfUnit.lean}{GlobalTheory.ConstRatio.Factor}
\leanref{Development/AutomorphicArithmetic/ThetaCorrespondence/ThetaPairLifts/LocalTypeDictionary/BottomTripleIntegralEqConstMulLocal.lean}{GlobalTheory.T.bottom_triple_integral_eq_const_mul_local}
\leanref{Development/AbelianVarietyTheory/Proofs/TwistedOperators/RatInOfDefinite.lean}{GlobalTheory.ConstRatio.ratIn_of_definite}
\leanref{Development/AbelianVarietyTheory/Statements/PlaceShiftStatements/AuxiliaryPrimeNeTypePrime.lean}{GlobalTheory.SqTop.Model.tower_numbers}
\begin{proof}
(1) Let $\kappa$ be as in Lemma~\ref{asm-dependence-test-assembly-first-layer-ratio-inputs-characters}.
By its part (1), $E$ is a CM field with the CM type $\Phi_D=\{\bar\tau:\tau\in\mathcal T\}$. By its part
(2), each of $\nu_{t,i,j}$, $\nu_{\mu_0,i,j}\delta$ and $\nu_{\mu_0,i,j}\delta'$ is unitary, and its
archimedean component is that of
Lemma~\ref{asm-dependence-test-assembly-first-layer-ratio-inputs-nonvanishing} for $\Phi_D$ with
$m(\bar\tau)=\kappa(\tau)$, since $|\bar\tau(k)|=|\tau(k)|$; and $\kappa(\tau)>h_1\ge0$ by its part (3).
As $0\le a'(i,j)$ and $a(i,j)\le\operatorname{rk}V_{i+1,j+1}$, also $e_0=(a-a')/2\le h_1$. So every
integer $s\ge1-e_0$ has $s+\kappa(\tau)>s+e_0\ge1$, and
Lemma~\ref{asm-dependence-test-assembly-first-layer-ratio-inputs-nonvanishing} gives an $L$-function
$\Lambda$ of $\nu_{t,i,j}$ with $\Lambda(s)\ne0$ for these $s$, in particular for $1-e_0\le s\le h_1$.
For the characters $\nu_{\mu_0,i,j}$ (the case $t=\mu_0$), $\nu_{\mu_0,i,j}\delta$ and
$\nu_{\mu_0,i,j}\delta'$ and an integer $s\ge1$, $s+\kappa(\tau)>0$, and the same lemma applies.

(2) Let a grid Weil datum with standard character $\psi$ be given, and put $a=a(i,j)$ and $a'=a'(i,j)$.
Take for $A$, $A'$ and $W$ the spaces $E^a$, $E^{a'}$ and $E^{a+a'}$ with identity Gram matrices,
standard lattices $\mathcal O_E^n$, which are self-dual, and trivial archimedean data; they are totally
positive definite. Let $t$ be a tuple.
\begin{enumerate}
\item[(i)] A computation from the restrictions to $\mathbf A_F^\times$ of the characters of the level
  system and of the grid shows that $\alpha_t$ and $\beta$ restrict to $\eta^{a'}$ and $\eta^a$, where
  $\eta$ is the quadratic character of $E/F$. So $\beta/\alpha_t$ and $\alpha_t/\beta$, which are
  continuous, unitary and trivial on $E^\times$, restrict to $\eta^{a+a'}$, as $\eta^2=1$. The spaces
  $-A$, $A'$ and $W$ have ranks $a$, $a'$ and $a+a'$, so
  Lemma~\ref{asm-dependence-test-assembly-first-layer-ratio-inputs-weil} gives Weil data $\omega_t$ for
  $(-A,W)$ with characters $(\beta,\beta/\alpha_t)$ and $\omega'_t$ for $(A',W)$ with characters
  $(\alpha_t,\alpha_t/\beta)$.
\item[(ii)] Fock vacuum vectors for $(-A,W)$ and $(A',W)$ exist by
  Lemma~\ref{is-eigenvalue-conjugate-vacuum-exists}(2).
\item[(iii)] Every place $w\in S_t$ is a place of the finite set of the grid or a place of $Q_t$, so it
  is not conjugate to itself, and distinct places of $S_t$ lie over distinct places of $F$. So
  $\varepsilon_t=\sum_{w\in S_t}1_w$ is an idempotent of $\mathbf A_E$ supported at $S_t$. The character
  $\nu_{t,i,j}$ has a conductor exponent $c_t(w)$ at each $w$, the least $n$ such that it is trivial on
  the units congruent to $1$ modulo $w^n$.
\item[(iv)] The groups $U(A)(F)$ and $U(A')(F)$ are discrete subgroups of the second countable locally
  compact groups $U(A)(\mathbf A)$ and $U(A')(\mathbf A)$, so for Haar measures on these groups they
  have measurable fundamental domains.
\item[(v)] By Lemma~\ref{asm-dependence-test-assembly-first-layer-ratio-inputs-standard}(2), the standard
  functions at $\varepsilon_t$ on $\mathbb W(-A,W)$ and at $\bar\varepsilon_t$ on $\mathbb W(A',W)$ are
  Schwartz--Bruhat on the halves. By part (1) of that lemma, applied to the Fock vacuum vectors of (ii),
  there are theta functions $f_t$ and $f'_t$ equal to $T(\Phi)$ for every place model $T$ of the vacuum at
  $\varepsilon_t$, respectively $\bar\varepsilon_t$, and the standard function $\Phi$ there.
\end{enumerate}
With these functions $f_t$ and $f'_t$, those of Definition~\ref{th-bottom-triple}(3), these data form a bottom
triple for $(i,j)$.

(3) Take the data large enough for Theorem~\ref{th-first-lift-nonzero}, whose hypothesis on $L$-values
is part (1), and fix a grid Weil datum with standard character, a bottom triple $(A,A',W)$ for $(i,j)$
and its characters. Take the data also large enough for
Lemma~\ref{asm-dependence-test-assembly-first-layer-ratio-inputs-factorization}, which gives a complex
vector space $\Pi$, linear maps $q$ and $q'$ from the polynomials in the creation operators of $(-A,W)$
and of $(A',W)$ to $\Pi$, with $q$ surjective, and for each tuple $t$ a linear map $J_t$ from $\Pi$ to
functions on $U(W)(\mathbf A)$ that is $0$ or injective, and a number $\lambda(t)$, with
$I_t(P,\cdot)=\lambda(t)J_t(q(P))$ and $I'_t(P',\cdot)=J_t(q'(P'))$ for all $P$ and $P'$. By
Theorem~\ref{th-first-lift-nonzero}, for each $t$ there are $P$ and $h$ with $I_t(P,h)\ne0$; so $J_t\ne0$,
hence $J_t$ is injective, and $\lambda(t)\ne0$. For $t=\mu_0$ there are $P'$ and $h'$ with
$I'_{\mu_0}(P',h')\ne0$, so $q'(P')\ne0$. Since $A$ and $A'$ are totally positive definite and the
functions $f_t\otimes P$ and $f'_t\otimes P'$ are theta functions, all the theta integrals converge
(Lemma~\ref{asm-dependence-test-assembly-first-layer-ratio-inputs-definite}(1)). So $\lambda$, $\Pi$,
$p_0=q$, $p'_0=q'$ and the $J_t$ form a factorization.

(5) Let $\Pi$, $p_0$, $p'_0$ and $J_t$ be those of the factorization, and choose $P'$ with
$p'_0(P')\ne0$. Since $p_0$ is surjective, choose $P$ with $p_0(P)=p'_0(P')$.
\begin{enumerate}
\item[(a)] \emph{Analytic facts.} The spaces $A$ and $A'$ are totally positive definite. By
  Lemma~\ref{asm-dependence-test-assembly-first-layer-ratio-inputs-definite}, all the theta integrals
  $I_t(Q,\cdot)$ and $I'_t(Q',\cdot)$ converge, so $I'_t(Q',\cdot)=J_t(p'_0(Q'))$ for every $Q'$; and
  they are continuous and invariant on the left under $U(W)(F)$. Let $U_f$ be the set of
  $h\in U(W)(\mathbf A)$ with archimedean components $1$, and write $\ell_{t,h}(Q)=I_t(Q,h)$ and
  $\ell'_{t,h}(Q')=I'_t(Q',h)$.
\item[(b)] \emph{Separation.} If $a,b\in\mathbf C$ and $a\ell_{t,h}(Q)=b\ell'_{t,h}(Q')$ for every
  $h\in U_f$, then $a\lambda(t)p_0(Q)=bp'_0(Q')$. Indeed, $aI_t(Q,\cdot)-bI'_t(Q',\cdot)$ is continuous,
  invariant on the left under $U(W)(F)$ and zero on $U_f$, so it is zero by
  Lemma~\ref{asm-dependence-test-assembly-first-layer-ratio-inputs-vanish}; it equals
  $J_t(a\lambda(t)p_0(Q)-bp'_0(Q'))$, and $J_t$ is injective. In particular, if $p'_0(Q')\ne0$, then
  $\ell'_{t,h}(Q')\ne0$ for some $h\in U_f$, by the case $a=0$, $b=1$.
\item[(c)] \emph{Rational parts.} View $\mathbf C$ as a vector space over $L$ through $\sigma$. For an
  $L$-linear map $\varphi\colon\mathbf C\to L$, a number $w$ and a polynomial $Q=\sum_kq_kx^k$ in the
  creation operators, put $Q^{\varphi,w}=\sum_k\sigma(\varphi(wq_k))x^k$. For $h\in U_f$ the values of
  $\ell_{t,h}$ at monomials lie in $\Omega\sigma(L)$, so by linearity
  $\ell_{t,h}(Q^{\varphi,w\Omega})=\Omega\,\sigma(\varphi(w\,\ell_{t,h}(Q)))$, and
  $\ell_{t,h}(Q^{\varphi,w})\in\Omega\sigma(L)$; the same holds for $\ell'_{t,h}$ with $\Omega'$.
\end{enumerate}
Fix a tuple $t_0$. By (b) there is $h_0\in U_f$ with $\ell'_{t_0,h_0}(P')\ne0$, and there is an $L$-linear
$\varphi\colon\mathbf C\to L$ with $\varphi(\ell'_{t_0,h_0}(P'))\ne0$. Since $p_0(P)=p'_0(P')$,
$\lambda(t_0)^{-1}\ell_{t_0,h}(P)=\ell'_{t_0,h}(P')$ for every $h$. Put
$Q=P^{\varphi,\lambda(t_0)^{-1}\Omega}$ and $Q'=P'^{\varphi,\Omega'}$. By (c), for $h\in U_f$,
\[
 \Omega'\ell_{t_0,h}(Q)=\Omega'\Omega\,\sigma\bigl(\varphi(\lambda(t_0)^{-1}\ell_{t_0,h}(P))\bigr)
 =\Omega\Omega'\sigma\bigl(\varphi(\ell'_{t_0,h}(P'))\bigr)=\Omega\,\ell'_{t_0,h}(Q'),
\]
so by (b), $\Omega'\lambda(t_0)p_0(Q)=\Omega p'_0(Q')$. Also $p'_0(Q')\ne0$, since
$\ell'_{t_0,h_0}(Q')=\Omega'\sigma(\varphi(\ell'_{t_0,h_0}(P')))\ne0$. Now let $t$ be a tuple. By (b) there
is $h\in U_f$ with $\ell'_{t,h}(Q')\ne0$, and by (c), $\ell_{t,h}(Q)=\Omega\sigma(z)$ and
$\ell'_{t,h}(Q')=\Omega'\sigma(z')$ with $z,z'\in L$ and $z'\ne0$. Applying $\lambda(t)J_t$ and evaluating
at $h$,
\[
 \Omega'\lambda(t_0)\,\ell_{t,h}(Q)=\lambda(t)\,\Omega\,\ell'_{t,h}(Q'),\quad\text{that is,}\quad
 \Omega\Omega'\lambda(t_0)\sigma(z)=\Omega\Omega'\lambda(t)\sigma(z').
\]
So $\sigma(z/z')=\lambda(t)/\lambda(t_0)$, which is nonzero; with $t_0=\mu_0$, $y=z/z'$ works.

(4) Take the data large enough for
Lemma~\ref{asm-dependence-test-assembly-first-layer-ratio-inputs-towers-inputs}, and let the data of (4) be
given. For each tuple $t'$ let $\Lambda_{t'}$ be an $L$-function of $\nu_{t',i,j}$, which exists by (1), with
$\Lambda_t=\Lambda$ and $\Lambda_{\mu_0}=\Lambda_0$; this is possible, as two $L$-functions of one character agree
on $\operatorname{Re}s>1$ and so everywhere. By (1), $\Lambda_0(s)\ne0$ for the integers $1\le s\le h_1$. The
numbers $\lambda(t')$ are constants of the bottom triple $(A,A',W)$
(Definition~\ref{asm-dependence-test-assembly-first-layer-ratio-inputs-towers-stages}(3)).
\begin{enumerate}
\item[(a)] \emph{Top.} By part (1) of
  Lemma~\ref{asm-dependence-test-assembly-first-layer-ratio-inputs-towers-inputs} there are a stage $(p,q)$
  with constants $c_T(t')$ and nonzero numbers $\kappa(t')$ with $\kappa(t')/\kappa(\mu_0)\in\sigma(L)$, such that
  $c=\kappa(t)c_T(t)$ and $c_0=\kappa(\mu_0)c_T(\mu_0)$. By part (2) the stage has a factorization, so by
  Lemma~\ref{asm-dependence-test-assembly-first-layer-ratio-inputs-towers-descent}(1) its constants are a
  nonzero multiple of the nonzero numbers of the factorization; so $c_T(t')\ne0$ for every $t'$. The stage has
  $\operatorname{rk}A_T=\operatorname{rk}V_{i+1,j+1}$ and $\operatorname{rk}A'_T=\operatorname{rk}V_{i,j}$, so
  $\operatorname{rk}A_T=\operatorname{rk}A'_T+2e$. Also $p\le e$, as $p\le m$ and
  $2e=\operatorname{rk}V_{i+1,j+1}-\operatorname{rk}V_{i,j}$ is the rank $t$ of $D$ plus $r+6$
  (Definition~\ref{th-grid}), and $t>3m$ (Definition~\ref{th-ranks}).
\item[(b)] \emph{Descent.} By
  Lemma~\ref{asm-dependence-test-assembly-first-layer-ratio-inputs-towers-descent}(2) there are a stage $(0,0)$
  with constants $c_Y(t')$, nonzero numbers $u_k(t')$, $0\le k<p+q$, with $u_k(t')/u_k(\mu_0)\in\sigma(L)$, and
  $\kappa'\ne0$ with
  \[
   c_T(t')\prod_{k=0}^{q-1}\Lambda_{t'}(e+k+1)=\kappa'\Bigl(\prod_ku_k(t')\Bigr)c_Y(t')
   \prod_{k=0}^{p-1}\Lambda_{t'}(e+q-k)
  \]
  for every $t'$.
\item[(c)] \emph{Bottom.} By part (5) of
  Lemma~\ref{asm-dependence-test-assembly-first-layer-ratio-inputs-towers-inputs} there are nonzero numbers
  $v(t')$ with $v(t')/v(\mu_0)\in\sigma(L)$ such that the $v(t')c_Y(t')$ are constants of $(A,A',W)$. The bottom
  triple is a stage $(0,0)$, so it has a factorization, and by
  Lemma~\ref{asm-dependence-test-assembly-first-layer-ratio-inputs-towers-descent}(1) both $\lambda$ and $vc_Y$
  are nonzero multiples of its numbers. Hence
  $\lambda(t)/\lambda(\mu_0)=v(t)c_Y(t)/(v(\mu_0)c_Y(\mu_0))$.
\item[(d)] \emph{Nonvanishing at $\mu_0$.} The integers $e+k+1$ with $k<q$ and $e+q-k$ with $k<p$ lie between
  $1$ and $e+q\le h_1$, as $p\le e$; so the products of the values of $\Lambda_0$ in (b) are nonzero. With
  $c_T(\mu_0)\ne0$ and $u_k(\mu_0)\ne0$, the identity (b) at $\mu_0$ gives $c_Y(\mu_0)\ne0$.
\end{enumerate}
Dividing the identity (b) at $t$ by the one at $\mu_0$ and using (a) and (c),
\[
 \frac c{c_0}\prod_{k=0}^{q-1}\frac{\Lambda(e+k+1)}{\Lambda_0(e+k+1)}
 =\frac{\kappa(t)}{\kappa(\mu_0)}\Bigl(\frac{v(t)}{v(\mu_0)}\Bigr)^{-1}
 \Bigl(\prod_k\frac{u_k(t)}{u_k(\mu_0)}\Bigr)\frac{\lambda(t)}{\lambda(\mu_0)}
 \prod_{k=0}^{p-1}\frac{\Lambda(e+q-k)}{\Lambda_0(e+q-k)}.
\]
The numbers $u_k(t)/u_k(\mu_0)$ are images of elements of $L$, and so is
$\kappa(t)v(\mu_0)/(\kappa(\mu_0)v(t))=\sigma(\zeta)$; this is the assertion.
% Lean (4): SqTower.tower_alt_bis, whose proof is followed: the inputs htop, hfac, hsm, hbig, hbot of Lemma
% -towers-inputs and hlne = LNe (Lf with IsHeckeL.unique for Lambda and Lambda_0), ranks_exact, t_odd,
% odd_sLen and three_m_lt_t for p <= e and e + q <= h_1, descent (Lemma -towers-descent(2)), hlamT.eq_mul and
% hlam0.eq_mul (Lemma -towers-descent(1)), stepRatio_div for u k t / u k mu0, and zeta = zetaT * zetav^-1
% (isAwayUnit_inv). The open part is in Lemma -towers-inputs.
% Lean: (1) SqTop.lNe_square (via lNe_of_functionalEquation; FieldDat from SqTop.fieldDat is used only for
% IsCMField and IsCMType (cmPhi D)), with Lemma -ratio-inputs-characters (isCMType_cmPhi, isValChar_twist,
% chrEq_nu, isUnitary_of_chrEq, inf_eq_one, kap_add), squareBottomHalfRankDiff_le_halfRank for
% e0 <= halfRank, LValNe.exists_heckeL and LValNe.heckeL_ne_zero_int (Lemma -ratio-inputs-nonvanishing);
% the second clause, for fam G mu0 i j c with c in {none, del, del'} and s >= 1, is exists_heckeL with
% IsHeckeL.ne_zero_nat. (2) The book's (2) is the existence of the triple in TriEx, which
% SqTop.triEx_holds_alt takes from RowTest.exists_bottomTri: exists_stdRig (standard spaces of ranks
% botBig, botSmall and their sum), bottom_chars and ConstRatio.exists_isWeil (Lemma -ratio-inputs-weil,
% through T.exists_weilDatum_of_splittingCharacters, with WeilParts.exists_isEtaTilde_unramified),
% VacuumO.vacExists, FunObj.LocTensor.exists_idemSum_locPl (Good.locPl_cases, Good.under_injOn,
% GridIdem.exists_isIdemSum), ConstRatio.exists_expAt, T.first_standard_pair_member_has_theta_function and
% T.exists_thetaFunction_standardPairSecond (Lemma -ratio-inputs-standard), and ConstRatio.exists_meas
% (FundDomDiscrete.exists_fundDom). The Weil data are built in the archive its own way, followed in Lemma
% -ratio-inputs-weil; the Weil representation of Theorem lit-bundle-theta is not used here. (3)
% T.bottom_triple_integral_eq_const_mul_local = ConstRatio.bottom_triple_integral_factorization:
% T.bottom_triple_integral_factorization_of_real (Lemma -ratio-inputs-factorization) and
% T.first_occurrence_theta_lift_ne_zero (Theorem th-first-lift-nonzero), with hX.integrable for the kernels. (5) ConstRatio.ratIn_of_definite with
% posDef_of_std: T.thetaIntegral_definiteSource_converges_continuous_invariant (Lemma -ratio-inputs-definite),
% T.continuous_automorphic_eq_zero_of_vanish_on_finiteAdelic (Lemma -ratio-inputs-vanish), and
% Factor.ratIn (sep, nz, exists_dual for phi, ratPart for Q^{phi,w}, ratPart_eq and ratPart_rat for (c)).
% (4) Tower = SqTower.tower_alt_bis (about 8,100 lines at depth 5): witt_tower_top_stage_theta_norm_formula,
% exists_factor_of_isStageC, witt_tower_norm_step_of_small_rank, witt_tower_norm_step_of_large_rank,
% witt_tower_constant_ratio_eq_bottom and descent, with LNe for the L-functions.
% Lean: (1) is LNe = SqTop.lNe_of_functionalEquation a12.abs
% HeckeStrip.completedLBoundedInStrips SqTop.fieldDat; (2) TriEx = SqTop.triEx_holds_alt, which also
% gives constants for the triple (IsConst), not used by sqRatio_mem_grid; (3)
% T.bottom_triple_integral_eq_const_mul_local hP hTA hLM hLN (ConstRatio.Factor); (4) Tower =
% SqTower.tower_alt_bis (StepRatio for u k, IsAwayUnit for zeta); (5) ConstRatio.ratIn_of_definite with
% posDef_of_std. The numbers e = eTop, p = pp, q = qq satisfy tower_numbers (e0 + pp = eTop + qq,
% eTop + qq <= halfRank, 1 <= e0). The sizes in the proof of Lemma -ratio: about 2,300, 10,200, 8,100
% and 2,100 lines at depth 5, overlapping.
\end{proof}

\begin{lemma}\label{asm-dependence-test-assembly-first-layer-ratio}
Assume the results of Theorem~\ref{lit-bundle-theta}, the theta hypotheses
(Definition~\ref{th-theta-hypotheses}) and the functional equation of
Theorem~\ref{asm-dependence-inputs}. Then statement (g) on test choices holds
(Definition~\ref{asm-dependence-test-assembly-statements}(3)).
\end{lemma}
\leanref{Development/AutomorphicArithmetic/ThetaCorrespondence/ThetaPairLifts/LocalTypeDictionary/TwistedLocalFactorRatioHolds.lean}{GlobalTheory.RowClaim.twisted_local_factor_ratio_holds}
\leanref{Development/AbelianVarietyTheory/Proofs/DivisorProofs/SquareBottomHalfRankDiffLeHalfRank.lean}{GlobalTheory.SqTop.lNe_of_functionalEquation}
\leanref{Development/AutomorphicArithmetic/ThetaCorrespondence/ThetaPairLifts/LocalTypeDictionary/TripleExistenceOfTheta.lean}{GlobalTheory.SqTop.triEx_holds_alt}
\leanref{Development/AutomorphicArithmetic/ThetaCorrespondence/ThetaPairLifts/LocalTypeDictionary/TowerOfThetaMatching.lean}{GlobalTheory.SqTower.tower_alt_bis}
\leanref{Development/AbelianVarietyTheory/Proofs/TwistedOperators/RatioOfLocalMatchingTowers.lean}{GlobalTheory.RowClaim.sqRatio_mem_grid}
\begin{proof}
By Lemma~\ref{asm-dependence-test-assembly-first-layer-rows}, statement (c) holds. Let $C_{F,0}$ be
at least its constant and that of Lemma~\ref{asm-dependence-test-assembly-first-layer-ratio-inputs},
for $\varepsilon$ and $C_F\ge C_{F,0}$ let $x_0$ be at least the numbers that they give, and let $D$,
its size data, $R$ and test choices for $C_F$ that satisfy the evaluation conditions be given, with
$x\ge x_0$. The automorphic setting gives the size conditions on the rows, the host conditions and
that $F$ is the maximal real subfield of $E$, so these data are sufficiently large, and the level
system of $R$ has base tuple $\mu_0$. Suppose that the
field conditions hold, and fix $i$, $j$ and a tuple $t$; put $c=c(i,j,t)$ and $c_0=c(i,j,\mu_0)$.

By (c), the squares $(i,j)$ at $t$ and at $\mu_0$ are nonzero, with constants $c$ and $c_0$, both
nonzero. By Lemma~\ref{asm-dependence-test-assembly-first-layer-ratio-inputs}(2) there is a bottom
triple for $(i,j)$; the field conditions then give a bottom triple $(A,A',W)$ for $(i,j)$ whose bottom
sums lie in $L$. By part (3) it has a factorization, with numbers $\lambda(t)\ne0$. If
$p'_0(P')\ne0$ with convergent integrals, choose $P$ with $p_0(P)=p'_0(P')$; then
$I_t(P,\cdot)=\lambda(t)I'_t(P',\cdot)$ at every tuple, and $I'_t(P',\cdot)=J_t(p'_0(P'))\ne0$ since $J_t$
is injective. By part (5), $\lambda(t)/\lambda(\mu_0)\in L$. By part (1) there are $L$-functions
$\Lambda$ of $\nu_{t,i,j}$ and $\Lambda_0$ of $\nu_{\mu_0,i,j}$ that do not vanish at the integers
$1-e_0\le s\le h_1$. The Weil datum of the test choices has the standard character of $R$, and the
field conditions say that the numbers of the grid choice at $\mu_0$ lie in $L$. So part (4) gives
$u_1,\dots,u_{p+q}$ and $\zeta$ as there.

The integers $e+k+1$ for $0\le k<q$ and $e+q-k$ for $0\le k<p$ lie between $e_0+1$ or $e+1$ and $e+q\le
h_1$, so between $1$ and $h_1$. At these integers $\Lambda$ and $\Lambda_0$ do not vanish, and the field
conditions give $\Lambda(n)/\Lambda_0(n)\in L$. So in the identity of part (4) every factor other than
$c/c_0$ lies in $L$, and the product on the left multiplying $c/c_0$ is nonzero. Hence $c/c_0\in L$.
% Lean: RowClaim.twisted_local_factor_ratio_holds P K LM j.2.1 l (SqTop.triEx_holds_alt hw P K LM l)
% (SqTower.tower_alt_bis P K LM l), with l = SqTop.lNe_of_functionalEquation a12.abs
% HeckeStrip.completedLBoundedInStrips SqTop.fieldDat. About 2,300, 10,200, 8,100 and 2,100 lines of
% Lean proofs at depth 5; they overlap. The constants max c1 c5 and max x1 x5 come from ConstK (c1)
% and gBig_and of LNe, TriEx, Tower and bottom_triple_integral_eq_const_mul_local (c5); Rec from
% hs.size and hs.host (rec_iff); the derivation is sqRatio_mem_grid: hN.bottom (NumbersIn) for the
% triple with RationalSums, hlam.isConst, ratIn_of_definite, hN.ratios for the L-value ratios,
% tower_numbers for the range of the arguments, StepRatio for the u k, and eq_div_of_mul_eq. The
% open parts are in Lemma -ratio-inputs.
\end{proof}

\subsection{Theta grids}\label{subsection-asm-theta-grids}

\begin{definition}\label{asm-dependence-test-assembly-first-layer-grid-pieces-places}
Let $D$ be a definite-side datum with host field $E/F$ and $K=\mathbf Q(\sqrt{-d})\subset E$, let $q_0$ be
the auxiliary prime of Definition~\ref{th-grid-places}, and let $S_1$ be the set of places of $F$ given
there. Let a level system for $D$ be given (Definition~\ref{th-grid}), with its set
$P_g$ of places and its sets $Q_t$. \emph{Places of a theta grid} for it are places of a theta grid for $P_g$
and the sets $Q_t$ as in Definition~\ref{th-grid-places}: a place $v_0$ of $F$ that contains $q_0$, the
\emph{base place}, and a place $w(v)$ of $E$ for every place $v$ of $F$, with conditions (1) to (5) there.
\end{definition}
\leanref{Development/AbelianVarietyTheory/Definitions/Divisors/BasePlaceNotMemChiGridBadPlaces.lean}{GlobalTheory.ChiGrid.Places}
\leanref{Development/AbelianVarietyTheory/Definitions/Divisors/LeastSplitAuxiliaryPrimeThetaChainSpec.lean}{GlobalTheory.ChiGrid.S1}
\leanref{Development/AbelianVarietyTheory/Definitions/Divisors/LeastSplitAuxiliaryPrimeThetaChainSpec.lean}{GlobalTheory.ChiGrid.q0}
\leanref{Development/AbelianVarietyTheory/Definitions/Divisors/LeastSplitAuxiliaryPrimeThetaChainSpec.lean}{GlobalTheory.ChiGrid.leastSplitAuxiliaryPrimeThetaChain_spec}

\begin{lemma}\label{asm-dependence-test-assembly-first-layer-grid-pieces-spaces}
Let $D$ be a definite-side datum with host field $E$.
\begin{enumerate}
\item For integers $0\le q\le\nu$ there is a Hermitian space of $D$ of rank $\nu$ with a self-dual lattice and
  signature $(\nu-q,q)$ at every embedding.
\item Let $X$ be a Hermitian space of $D$ of rank $\nu$ with a self-dual lattice $\Lambda_X$ and signature
  $(\nu-q,q)$ at every embedding. Let $c\ge0$, $\epsilon\in\{1,-1\}$, and let $\mathfrak b_1,\dots,\mathfrak b_c$
  be nonzero fractional ideals of $E$ with $\mathfrak b_{c+1-k}=\bar{\mathfrak b}_k^{-1}$ for every $k$. Then there
  are a Hermitian space $Y$ of $D$ and an isometric embedding $X\subset Y$ whose orthogonal complement has a
  basis $a_1,\dots,a_c$ with $\langle a_k,a_l\rangle=1$ if $k+l=c+1$ and $k\ne l$, $\langle a_k,a_k\rangle=\epsilon$
  if $2k=c+1$, and $\langle a_k,a_l\rangle=0$ otherwise, such that $\Lambda_X+\sum_k\mathfrak b_ka_k$ is a
  self-dual lattice of $Y$. The space $Y$ has rank $\nu+c$ and signature $(\nu-q+\lceil c/2\rceil,q+\lfloor
  c/2\rfloor)$ if $\epsilon=1$ and $(\nu-q+\lfloor c/2\rfloor,q+\lceil c/2\rceil)$ if $\epsilon=-1$.
\end{enumerate}
\end{lemma}
\leanref{Development/AbelianVarietyTheory/Proofs/DivisorProofs/ExistsSpaceOfSignature.lean}{GlobalTheory.ChiGrid.exists_isV}
\leanref{Development/AbelianVarietyTheory/Definitions/Divisors/NonemptyGridDataIff.lean}{GlobalTheory.ChiGridModel.exists_sum}
\leanref{Development/AbelianVarietyTheory/Proofs/DivisorProofs/HermitianSpaceAddSignatureExists.lean}{GlobalTheory.SignatureAddition.exists_hermitian_space_add_signature}
\leanref{Development/AbelianVarietyTheory/Proofs/DivisorProofs/AntiDiagonalEpsilonGramMulVec.lean}{GlobalTheory.ChiGrid.SumDual.isSelfDualC_lat}
\begin{proof}
(1) Take the diagonal form with $\nu-q$ entries $1$ and $q$ entries $-1$, with the self-dual lattice of its frame.

(2) Let $Y$ be the orthogonal sum of $X$ and the space $E^c$ with the Gram matrix of the $a_k$ in the statement.
That matrix is the sum of $\lfloor c/2\rfloor$ hyperbolic planes and, if $c$ is odd, of the line $\langle
a_k,a_k\rangle=\epsilon$ with $2k=c+1$; so its signature is $(\lceil c/2\rceil,\lfloor c/2\rfloor)$ if $\epsilon=1$
and $(\lfloor c/2\rfloor,\lceil c/2\rceil)$ if $\epsilon=-1$ at every embedding, and the signature of $Y$ is the sum.
The lattice $\Lambda_X+\sum_k\mathfrak b_ka_k$ is self-dual: $\Lambda_X$ is self-dual in $X$, and the dual of
$\sum_k\mathfrak b_ka_k$ in the complement is $\sum_k\bar{\mathfrak b}_{c+1-k}^{-1}a_k$, by the shape of the Gram
matrix with $\epsilon\bar\epsilon=1$, which is $\sum_k\mathfrak b_ka_k$ by the hypothesis on the ideals.
% Lean: (1) ChiGrid.exists_isV (diagSp with signs, latFrame, isSelfDualC_latFrame, hasSignature_diag_signs).
% (2) ChiGridModel.exists_sum: SignatureAddition.exists_hermitian_space_add_signature (Doubling.hsum of A with
% the standard space of the antidiagonal Gram matrix antiDiagonalSignGramMatrix c eps, sigOK_hsum, and the
% signature count of UnitPerp.ofKind_stdSpace) and SumDual.isSelfDualC_lat (selfDual_sum, the dual computed
% coordinatewise from the Gram matrix; the Lean's hypothesis is x in b(rev p) iff star x in (b p)^-1).
\end{proof}

\begin{lemma}\label{asm-dependence-test-assembly-first-layer-grid-pieces-characters-wild}
Let $L$ be a number field, $v$ a finite place of $L$ and $p$ a prime in $v$. Then there are a character $\theta$ of
$L_v^\times$ whose values are $p$-th roots of unity and which is trivial on the units of $L_v$ congruent to $1$ modulo
$v^2$, and an element $k\in L^\times$, a unit at $v$ congruent to $1$ modulo $v$, with $\theta(k)\ne1$.
\end{lemma}
\leanref{Development/AbelianVarietyTheory/Proofs/DivisorProofs/ExistsPrimeOfMemChiGridBadPlaces.lean}{GlobalTheory.GridChars.exists_wild}
\leanref{Development/AbelianVarietyTheory/Proofs/DivisorProofs/ExistsPrimeOfMemChiGridBadPlaces.lean}{GlobalTheory.GridChars.logOne}
\begin{proof}
Let $\kappa=\mathcal O_L/v$, a finite field of characteristic $p$ with $\mathrm Nv$ elements, and put $n=\mathrm Nv-1$,
which is prime to $p$. Fix $\varpi\in L^\times$ with $\operatorname{ord}_v(\varpi)=1$, and for $x\in L_v^\times$ let
$u(x)=x\varpi^{-\operatorname{ord}_v(x)}$, a unit. Then $u$ is multiplicative, and $u(x)=x$ for a unit $x$. The residue
of $u(x)^n$ is $1$, as $\kappa^\times$ has order $n$; so $u(x)^n$ lies in the group $U^1$ of the units congruent to $1$
modulo $v$. For $y\in U^1$ let $\ell(y)\in\kappa$ be the residue of $(y-1)/\varpi$. Then $\ell$ is a homomorphism to the
additive group of $\kappa$, as
\[
 \frac{yz-1}\varpi=\frac{y-1}\varpi+\frac{z-1}\varpi+(y-1)\frac{z-1}\varpi
\]
and the last term lies in $v$. Let $\psi$ be a character of the additive group of $\kappa$ with $\psi(1)\ne1$; its
values are $p$-th roots of unity, as $p\kappa=0$. Put $\theta(x)=\psi(\ell(u(x)^n))$, a character of $L_v^\times$ whose
values are $p$-th roots of unity. If $x$ is a unit congruent to $1$ modulo $v^2$, then so is $u(x)^n=x^n$, so
$\ell(x^n)=0$ and $\theta(x)=1$. Let $k=1+\varpi$, a unit at $v$ congruent to $1$ modulo $v$. Then
$\ell(k^n)=n\,\ell(k)=n$, and $\theta(k)=\psi(1)^n\ne1$, as $\psi(1)$ has order $p$ and $p\nmid n$.
\end{proof}
% Lean: exists_wild (60 lines), followed: unitPart and pow_unitPart_mem for u(x)^n, oneUnits for U^1, logOne (resInt of
% (x - 1) / unifC v) for ell, CommGroup.exists_apply_ne_one_of_hasEnoughRootsOfUnity for psi, coprime_card and dvd_card
% for n prime to p, and k = 1 + unif v. About 175 lines with logOne and its helpers.

\begin{lemma}\label{asm-dependence-test-assembly-first-layer-grid-pieces-characters-tame}
Let $L$ be a number field with group of roots of unity $\mu_L$, let $v$ be a finite place of $L$ and $p$ a prime in
$v$ with $|\mu_L|+1<p$, and let $\varepsilon$ be a character of $L^\times$. Then there are a character $\theta$ of
$L_v^\times$ with values in the unit circle, trivial on the units congruent to $1$ modulo $v$, with
$\theta(\zeta)=\varepsilon(\zeta)$ for every $\zeta\in\mu_L$, and an element $k\in L^\times$, a unit at $v$, with
$\theta(k)\ne1$.
\end{lemma}
\leanref{Development/AbelianVarietyTheory/Proofs/DivisorProofs/ExistsPrimeOfMemChiGridBadPlaces.lean}{GlobalTheory.GridChars.exists_tame}
\leanref{Development/AutomorphicArithmetic/CyclotomicBounds/RootsOfUnity/ResidueEisenstein.lean}{UpperTheory.Mu.tame_ext_prime}
\leanref{Development/AutomorphicArithmetic/CyclotomicBounds/RootsOfUnity/TorsionReduceInjective.lean}{UpperTheory.Mu.tame_ext}
\leanref{Development/AbelianVarietyTheory/Proofs/IwasawaDecomposition/ExistsTameLocOfResidue.lean}{GlobalTheory.Mem.Chr.exists_tameLoc_of_residue}
\begin{proof}
Let $\kappa=\mathcal O_L/v$. Its order $\mathrm Nv$ is a power of $p$, so $\mathrm Nv\ge p>|\mu_L|+1$. The roots of unity
of $L$ are units, and reduction modulo $v$ is injective on $\mu_L$: $|\mu_L|<p$, so $|\mu_L|\notin v$, and distinct
$|\mu_L|$-th roots of unity are distinct modulo a prime that does not contain $|\mu_L|$ (as in the proof of
Lemma~\ref{is-eigenvalue-conjugate-frame-phase-eta-roots}). So the image $M$ of $\mu_L$ in $\kappa^\times$ has order
$|\mu_L|$ and index $(\mathrm Nv-1)/|\mu_L|\ge2$. The characters of the finite abelian group $\kappa^\times$ that
extend the character of $M$ given by $\varepsilon$ form a coset of the group of the characters trivial on $M$, whose
order is this index; so one of them, $\theta_0$, is not trivial. Fix $\varpi\in L^\times$ with
$\operatorname{ord}_v(\varpi)=1$, and let $\theta(x)$ be $\theta_0$ of the residue of the unit
$x\varpi^{-\operatorname{ord}_v(x)}$. Then $\theta$ is a character of $L_v^\times$ whose values are roots of unity, it
is trivial on the units congruent to $1$ modulo $v$, and $\theta(\zeta)=\varepsilon(\zeta)$ for $\zeta\in\mu_L$. As
$\theta_0\ne1$, some $b\in\kappa^\times$ has $\theta_0(b)\ne1$; an element $k\in\mathcal O_L$ with residue $b$ is a
unit at $v$ with $\theta(k)=\theta_0(b)\ne1$.
\end{proof}
% Lean: exists_tame (15 lines), followed: tame_ext_prime (|mu_L| not in v from |mu_L| < p, and Nv >= p through
% absNorm dividing p^[L:Q]) and tame_ext (torsion_reduce_injective for M, MonoidHom.domRestrict_surjective for an
% extension theta0, and a nontrivial character trivial on M, from CommGroup.card_domRestrictHom_ker, when theta0 = 1);
% then exists_tameLoc_of_residue (theta0 composed with resPart, the residue of the unit part), exists_lift for k and
% exists_unit_of_root for the roots of unity. About 150 lines with the helpers.

\begin{lemma}\label{asm-dependence-test-assembly-first-layer-grid-pieces-characters-extension}
Let $G$ be an abelian topological group with an open subgroup $U$, let $A$ be an abelian group with homomorphisms
$\alpha\colon A\to G$ and $a\colon A\to\mathbf C^\times$, with values of $a$ in the unit circle, and let $\psi$ be a
continuous character of $G$ with values in the unit circle such that $a(y)=\psi(\alpha(y))$ whenever $\alpha(y)\in U$.
Then there is a continuous character $\chi$ of $G$ with values in the unit circle such that $\chi=\psi$ on $U$ and
$\chi(\alpha(y))=a(y)$ for every $y\in A$.
\end{lemma}
\leanref{Development/AbelianVarietyTheory/Proofs/IwasawaDecomposition/ExistsUnitaryCharOfOpenSubgroupAndHom.lean}{GlobalTheory.exists_unitary_char_of_open_subgroup_and_hom}
\begin{proof}
Give $A$ the discrete topology, and let $\Gamma\subset G\times A$ be the subgroup of the elements $(\alpha(y),y)$. The
character $(x,y)\mapsto\psi(x)a(y)^{-1}$ of the open subgroup $U\times A$ is continuous, with values in the unit
circle, and it is trivial on $(U\times A)\cap\Gamma$ by the hypothesis. So it induces a character of the subgroup
$(U\times A)\Gamma/\Gamma$ of $(G\times A)/\Gamma$. The unit circle is a divisible abelian group, hence an injective one
\cite[Chapter~XX, \S4]{Lan02}, so this character extends to $(G\times A)/\Gamma$. This gives a character $\tilde\chi$ of
$G\times A$ with values in the unit circle, trivial on $\Gamma$ and equal to $\psi(x)a(y)^{-1}$ on $U\times A$; it is
continuous, as it is continuous on an open subgroup. Put $\chi(x)=\tilde\chi(x,1)$. Then $\chi=\psi$ on $U$, and
$\chi(\alpha(y))=\tilde\chi(\alpha(y),y)\,\tilde\chi(1,y)^{-1}=a(y)$.
\end{proof}
% Lean: exists_unitary_char_of_open_subgroup_and_hom (about 80 lines), followed, through
% exists_continuous_circle_char_extension (circle_hom_extend_of_injective). Chapter 8 gives the same argument inside the
% proof of Lemma is-eigenvalue-conjugate-frame-phase-eta-extension.

\begin{lemma}\label{asm-dependence-test-assembly-first-layer-grid-pieces-characters-anti}
Let $L$ be a CM field with complex conjugation $c$, and let $C$ be the automorphism of $\mathbf A_L^\times$ induced by
$c$, which carries the component at a place $w$ to the component at $c(w)$. Let $\psi$ be a continuous character of
$\mathbf A_L^\times$ with values in the unit circle such that $\psi(\zeta)=1$ for every root of unity $\zeta\in L$. Then
there is a continuous character $\chi$ of $\mathbf A_L^\times$ with values in the unit circle, trivial on $L^\times$,
such that $\chi(z)=\psi(z/C(z))$ for every idele $z$ for which $z/C(z)$ is a unit at every finite place.
\end{lemma}
\leanref{Development/AbelianVarietyTheory/Proofs/IwasawaDecomposition/ExistsHeckeCharTrivOnRealSubfieldInfinityType.lean}{GlobalTheory.Mem.Chr.AntiExist.exists_anti}
\leanref{Development/AbelianVarietyTheory/Proofs/IwasawaDecomposition/ExistsUnitOfRoot.lean}{GlobalTheory.Mem.Chr.normOne_unit_isRoot}
\begin{proof}
Let $U$ be the open subgroup of the ideles that are units at every finite place, and $A$ the group of the
$k\in L^\times$ with $k\,c(k)=1$. Apply
Lemma~\ref{asm-dependence-test-assembly-first-layer-grid-pieces-characters-extension} to $G=\mathbf A_L^\times$, $U$, the
inclusion $\alpha$ of $A$ in $L^\times\subset\mathbf A_L^\times$, $a=1$ and $\psi$. Its hypothesis holds: if $k\in A$ lies
in $U$, then $k$ is an algebraic integer, and $|\varphi(k)|^2=\varphi(k\,c(k))=1$ for every embedding $\varphi$ of $L$, as
$\varphi\circ c$ is the complex conjugate of $\varphi$ on the CM field $L$; so $k$ is a root of unity by Kronecker's
theorem \cite{Kro57}, and $\psi(k)=1$. Let $\nu$ be the character that the lemma gives, and put $\chi(z)=\nu(z/C(z))$.
It is continuous, with values in the unit circle. For $k\in L^\times$, $k/c(k)$ lies in $A$, so
$\chi(k)=\nu(k/c(k))=1$. If $z/C(z)\in U$, then $\chi(z)=\psi(z/C(z))$.
\end{proof}
% Lean: exists_anti (25 lines), followed: exists_unitary_char_of_open_subgroup_and_hom with A = normOne c, a = 1 and
% U = unitIdl; the compatibility is normOne_unit_isRoot (Kronecker through
% NumberField.Embeddings.pow_eq_one_of_norm_eq_one) and the hypothesis h1; chi = nu composed with id / C, and hC.diag
% for the principal ideles.

\begin{lemma}\label{asm-dependence-test-assembly-first-layer-grid-pieces-characters-hecke}
Let $L$ be a CM field with maximal totally real subfield $F$ and complex conjugation $x\mapsto\bar x$, and let
$(\varphi_i)_{i\in I}$ be a family of embeddings $L\to\mathbf C$ whose places are the archimedean places of $L$, each
once. For a finite place $v$ of $L$ and $x\in L_v^\times$ write $x_{(v)}$ for the idele with component $x$ at $v$ and
$1$ elsewhere. Let $e\colon I\to\mathbf Z$, let $S$ be a finite set of finite places of $L$ such that $\bar w\notin S$
for every $w\in S$, and for $w\in S$ let $\theta_w$ be a continuous character of $L_w^\times$ with values in the unit
circle. Suppose that
\[
 \prod_{i\in I}\varphi_i(\zeta)^{e(i)}\prod_{w\in S}\theta_w(\zeta)=1
\]
for every root of unity $\zeta\in L$. Then there is a continuous character $\chi$ of $\mathbf A_L^\times$ with values in
the unit circle, trivial on $L^\times$, such that:
\begin{enumerate}
\item $\chi$ is trivial on $\mathbf A_F^\times$;
\item at the place of each $\varphi_i$, $\chi$ is $z\mapsto(\varphi_i(z)/|\varphi_i(z)|)^{2e(i)}$;
\item $\chi(x_{(v)})=1$ for every finite place $v$ of $L$ over a place of $F$ that lies under no place of $S$, and
  every unit $x$ of $L_v$;
\item let $w\in S$, let $v$ be $w$ or $\bar w$, and let $\sigma\colon L_v\to L_w$ be the identity or the isomorphism
  induced by complex conjugation. Then $\chi(x_{(v)})=1$ if and only if $\theta_w(\sigma(x))=1$, for every unit $x$ of
  $L_v$.
\end{enumerate}
\end{lemma}
\leanref{Development/AbelianVarietyTheory/Proofs/DivisorProofs/ExistsPrimeOfMemChiGridBadPlaces.lean}{GlobalTheory.GridChars.exists_hecke_level}
\leanref{Development/AbelianVarietyTheory/Proofs/IwasawaDecomposition/ContinuousOfTameLoc.lean}{GlobalTheory.ModelChar.exists_modelChar}
\begin{proof}
Let $\psi_0$ be the continuous character of $\mathbf A_L^\times$ that is trivial on the finite ideles and is
$z\mapsto(\varphi_i(z)/|\varphi_i(z)|)^{e(i)}$ at the place of $\varphi_i$, and put
$\psi(z)=\psi_0(z)\prod_{w\in S}\theta_w(z_w)$. It is a continuous character with values in the unit circle, and for a
finite place $v$ and $x\in L_v^\times$, $\psi(x_{(v)})=\theta_v(x)$ if $v\in S$ and $\psi(x_{(v)})=1$ otherwise. For a
root of unity $\zeta\in L$, $|\varphi_i(\zeta)|=1$, so $\psi(\zeta)=\prod_i\varphi_i(\zeta)^{e(i)}\prod_{w\in S}
\theta_w(\zeta)=1$. Let $\chi$ be the character of
Lemma~\ref{asm-dependence-test-assembly-first-layer-grid-pieces-characters-anti} for $\psi$, with the automorphism $C$
there, so that $\chi(z)=\psi(z/C(z))$ whenever $z/C(z)$ is a unit at every finite place. For a finite place $v$ and a
unit $x$ of $L_v$, $C(x_{(v)})=\bar x_{(\bar v)}$, where $\bar x\in L_{\bar v}$ is the image of $x$, and
$x_{(v)}/\bar x_{(\bar v)}$ is a unit at every finite place; so
\[
 \chi(x_{(v)})=\psi(x_{(v)})\,\psi(\bar x_{(\bar v)})^{-1}.
\]

(1) For $z\in\mathbf A_F^\times$, $C(z)=z$, so $\chi(z)=\psi(1)=1$.

(2) Let $z$ be an idele whose components away from the place of $\varphi_i$ are $1$. As $c$ is complex conjugation
through every embedding of $L$, $C(z)$ is the idele with component $\bar z$ there; so $z/C(z)$ is a unit at every finite
place, and, reading components through $\varphi_i$, $\chi(z)=\psi_0(z/C(z))=(\varphi_i(z)/\overline{\varphi_i(z)})^{e(i)}
=(\varphi_i(z)/|\varphi_i(z)|)^{2e(i)}$.

(3) If $v$ lies over a place of $F$ under no place of $S$, then $v,\bar v\notin S$, and $\chi(x_{(v)})=1$ by the
displayed formula.

(4) Let $w\in S$; then $\bar w\notin S$. For a unit $x$ of $L_w$ the formula gives $\chi(x_{(w)})=\theta_w(x)$. For a unit
$y$ of $L_{\bar w}$ it gives $\chi(y_{(\bar w)})=\theta_w(\bar y)^{-1}$, with $\bar y\in L_w$ the image of $y$. In both
cases $\chi(x_{(v)})=1$ if and only if $\theta_w(\sigma(x))=1$.
\end{proof}
% Lean: exists_hecke_level (about 85 lines), followed: psi0 is ModelChar.exists_modelChar with S empty (the model
% character of infinity type e), psi = psi0 times the theta v at the components in S (hpsiS, hpsiN), h1 for the roots of
% unity from hpar, and AntiExist.exists_anti; hfix for (1) (conj_locF, conj_upF, conj_inf), phase_div with
% InfTypeLocal.infType_iff_local for (2), conj_loc and hloc for (3) and (4). The Lean states (4) as the existence of a
% map sigma that is surjective onto the image of L^x and preserves units, the valuation of x - 1 and roots of unity; the
% book names it, as the proof does.

\begin{lemma}\label{asm-dependence-test-assembly-first-layer-grid-pieces-characters}
Let $D$ be a definite-side datum, let a level system for $D$ be given with places of a theta grid
(Definition~\ref{asm-dependence-test-assembly-first-layer-grid-pieces-places}) whose base place $v_0$ lies
outside $S_1$, and let $c$ be a real number. Then there are unitary Hecke characters $\chi_1,\dots,\chi_m$ of $E$
with the conditions (a)--(f) of Definition~\ref{th-grid} for these places and the constant $c$.
\end{lemma}
\leanref{Development/AbelianVarietyTheory/Proofs/DivisorProofs/ExistsPrimeOfMemChiGridBadPlaces.lean}{GlobalTheory.GridChars.nonempty_chars}
\leanref{Development/AbelianVarietyTheory/Proofs/DivisorProofs/ExistsPrimeOfMemChiGridBadPlaces.lean}{GlobalTheory.GridChars.exists_one}
\begin{proof}
Let $E/F$ be the host field of $D$, with complex conjugation $x\mapsto\bar x$ and set $\mathcal T$ of embeddings. It is a
CM field with maximal totally real subfield $F$, and its archimedean places are the places of the $\tau\in\mathcal T$,
each once (Lemma~\ref{asm-dependence-test-assembly-first-layer-ratio-inputs-characters}(1)). Put $S_0=\{v_0\}\cup S_1$,
$w_0=w(v_0)$ and $W_0=\{w(v):v\in S_0\}$; each $w(v)$ lies over $v$. Fix $1\le i\le m$, and let $e(\tau)=c_i(\tau)$ be the
centres of Definition~\ref{th-grid-centres} for the constant $c$.

\emph{The local characters at $S_1$.} Let $v\in S_1$. It contains one of $g_0,g_1,p_{de},p'_{de}$
(Definition~\ref{th-grid-places}), a prime $p_v$, and so does $w(v)$. By
Lemma~\ref{asm-dependence-test-assembly-first-layer-grid-pieces-characters-wild} there are a character $\theta_{w(v)}$ of
$E_{w(v)}^\times$ whose values are $p_v$-th roots of unity, trivial on the units congruent to $1$ modulo $w(v)^2$, and
$k_v\in E^\times$, a unit congruent to $1$ modulo $w(v)$, with $\theta_{w(v)}(k_v)\ne1$.

\emph{The local character at $v_0$.} The place $w_0$ contains $q_0$, as $v_0$ does. The group $\mu_E$ of roots of unity
of $E$ has $|\mu_E|+1<q_0$: the orders of its elements divide $N=Q_Sqd$ (proof of
Lemma~\ref{is-eigenvalue-conjugate-frame-phase-eta-roots}), so $|\mu_E|\le N$, while $q_0\equiv1$ modulo $r_\ast N$
(Definition~\ref{th-grid-places}), so $q_0\ge2N+1$. Let $\varepsilon$ be the character
\[
 \varepsilon(k)=\Bigl(\prod_{\tau\in\mathcal T}\tau(k)^{e(\tau)}\prod_{v\in S_1}\theta_{w(v)}(k)\Bigr)^{-1}
\]
of $E^\times$. By Lemma~\ref{asm-dependence-test-assembly-first-layer-grid-pieces-characters-tame}, for $w_0$ and the
prime $q_0$, there are a character $\theta_{w_0}$ of $E_{w_0}^\times$ with values in the unit circle, trivial on the
units congruent to $1$ modulo $w_0$ and equal to $\varepsilon$ on $\mu_E$, and $k_0\in E^\times$, a unit at $w_0$, with
$\theta_{w_0}(k_0)\ne1$.

\emph{The character $\chi_i$.} Apply
Lemma~\ref{asm-dependence-test-assembly-first-layer-grid-pieces-characters-hecke} to $L=E$, the family
$(\tau)_{\tau\in\mathcal T}$, $e$, $S=W_0$ and the characters $\theta_w$. Its hypotheses hold.
\begin{itemize}
\item The places $w(v)$, $v\in S_0$, lie over the distinct places $v$ of $F$. If $\overline{w(v)}=w(v')$, then $v=v'$, as
  conjugate places lie over the same place of $F$, and $w(v)=\overline{w(v)}$, against condition (1) of
  Definition~\ref{th-grid-places}. So $\bar w\notin W_0$ for $w\in W_0$.
\item Each $\theta_w$ has values in the unit circle and is continuous, as it is trivial on an open subgroup of the
  units.
\item As $v_0\notin S_1$, $W_0$ is the disjoint union of $\{w_0\}$ and the $w(v)$, $v\in S_1$; so for $\zeta\in\mu_E$,
  $\prod_\tau\tau(\zeta)^{e(\tau)}\prod_{w\in W_0}\theta_w(\zeta)=\varepsilon(\zeta)^{-1}\theta_{w_0}(\zeta)=1$.
\end{itemize}
Let $\chi_i$ be the character that the lemma gives. It is a unitary Hecke character of $E$. As the places of $F$ under
the places of $W_0$ are those of $S_0$, (1), (2) and (3) of the lemma give the conditions (a), (b) and (c) of
Definition~\ref{th-grid}.

\emph{The conditions (d), (e) and (f).} Let $P$ be a place of $E$ over a place $v\in S_0$. As $E/F$ is Galois of degree
$2$, $P$ is $w(v)$ or $\overline{w(v)}$; let $\sigma\colon E_P\to E_{w(v)}$ be the identity or the isomorphism induced by
complex conjugation. It maps the units congruent to $1$ modulo $P^n$ onto those congruent to $1$ modulo $w(v)^n$, for
$n\ge0$, and $E^\times$ onto $E^\times$. Write $x_{(P)}$ for the idele with component $x\in E_P^\times$ at $P$ and $1$
elsewhere. By (4) of the lemma, $\chi_i(x_{(P)})=1$ if and only if
$\theta_{w(v)}(\sigma(x))=1$, for every unit $x$ of $E_P$.
\begin{itemize}
\item Let $v=v_0$. Then $\chi_i$ is trivial on the units of $E_P$ congruent to $1$ modulo $P$, as $\theta_{w_0}$ is on
  those modulo $w_0$; and $\chi_i(x_{(P)})\ne1$ for the unit $x=\sigma^{-1}(k_0)$, so $\chi_i$ is ramified at $P$. This is
  (d).
\item Let $v\in S_1$. Then $\chi_i$ is trivial on the units of $E_P$ congruent to $1$ modulo $P^2$, and
  $\chi_i(x_{(P)})\ne1$ for $x=\sigma^{-1}(k_v)$, a unit congruent to $1$ modulo $P$. This is (e). Let $x\in E_P^\times$ with
  $x^{\mathrm NP-1}=1$. Then $x$ is a unit, and $t=\theta_{w(v)}(\sigma(x))$ satisfies $t^{\mathrm NP-1}=1$ and
  $t^{p_v}=1$. As $\mathrm NP$ is a power of $p_v$, $\mathrm NP-1$ is prime to $p_v$; so $t=1$, and $\chi_i(x_{(P)})=1$.
  Every place of $E$ that contains $p_{de}$ or $p'_{de}$ lies over a place of $S_1$, so this gives (f).
\end{itemize}
Doing this for $i=1,\dots,m$ gives the characters $\chi_1,\dots,\chi_m$.
% Lean: GridChars.nonempty_chars (exists_one for each i with e = cen D C tau i, then CharTransport.exists_chrEq, which
% passes to characters of the units of A_F tensor E; the book identifies the two). exists_one (about 95 lines),
% followed: Mem.HostCM.host_cm for the CM data with D.eTau; exists_prime_of_mem_chiGridBadPlaces and exists_wild for the
% theta1 w at the places over S1; exists_tame_par at w0 (exists_tame for the character eps of the book, with
% basePlace_over_qZero and torsion_order_lt_auxiliary_prime); exists_hecke_level with S = W0, continuity from
% continuous_of_tameLoc and continuous_of_level, and not_conjPl_self (the Lean's form of condition (1) through the
% split idempotent) for the conjugate places; then lev1 and wild for the places over S1 (coprime_card for N P - 1 prime
% to p) and lev0 for the places over v0. The bound |mu_E| + 1 < q0 is UpperTheory.Mu.w_add_one_lt, which uses that
% |mu_E| divides Q_S (w_dvd_QS); the book uses the weaker |mu_E| <= N. About 480 lines in
% ExistsPrimeOfMemChiGridBadPlaces.lean, and about 400 more in the helpers of the lemmas above.
\end{proof}

\begin{lemma}\label{asm-dependence-test-assembly-first-layer-grid-pieces-weil-vectors}
Let $D$ be a definite-side datum with host field $E/F$ and set $\mathcal T$ of embeddings, let $\psi$ be a standard
character (Definition~\ref{th-weil-datum}) and $\theta$ an archimedean family whose member at each $\tau\in\mathcal T$
extends $\tau$. Let $S$ and $T$ be Hermitian spaces of $D$ with correct signatures (Definition~\ref{th-spaces}),
Cartan involutions $s_S$ and $s_T$ at the $\tau\in\mathcal T$, and self-dual lattices $\Lambda_S$ and $\Lambda_T$; let
$\Lambda\subset\mathbb W$ be the lattice of the pair and $M=\operatorname{Hom}(\Lambda_S,\Lambda_T)$. Let $e$ be a sum
of the idempotents of $\mathbf A_E$ at the places of a finite set, with $e\bar e=0$, and let $\mathfrak m$ be a nonzero
ideal of $\mathcal O_E$. Then there is a Fock vacuum vector $f$ for $s_S$, $s_T$ and $\Lambda$, and for every class
$\kappa$ of $M$ modulo $\mathfrak mM$ there is a theta function $f_\kappa$ with $f_\kappa=T(\Phi_\kappa)$ for every place
model $T$ of $f$ at $e$ (Lemma~\ref{asm-dependence-test-assembly-first-layer-ratio-inputs-standard}), where
$\Phi_\kappa$ is the indicator function of the set of the $x\in e\mathbb W$ with $x-ex_\kappa$ in the closure of
$\mathfrak mM$, for an element $x_\kappa\in M$ of $\kappa$.
\end{lemma}
\leanref{Development/AbelianVarietyTheory/Proofs/DivisorProofs/ExistsWeilOver.lean}{GlobalTheory.ChiGrid.WeilOver.exists_vac_cos}
\leanref{Development/AbelianVarietyTheory/Proofs/DivisorProofs/ExistsWeilOver.lean}{GlobalTheory.ChiGrid.WeilOver.isSBHalf_cosInd}
\begin{proof}
Lemma~\ref{is-eigenvalue-conjugate-vacuum-exists}(2) gives the Fock vacuum vector $f$. The function $\Phi_\kappa$ is
Schwartz--Bruhat on $e\mathbb W$ in the sense of
Lemma~\ref{asm-dependence-test-assembly-first-layer-ratio-inputs-standard}. Indeed, as $\Lambda_S$ and $\Lambda_T$ are
lattices, $\Lambda$ and the closure of $M$ are commensurable: there are integers $a,b\ge1$ with $ay\in\Lambda$ for
every $y$ in the closure of $M$ and $bu$ in the closure of $M$ for every $u\in\Lambda$. Put $k=ab\,\mathrm N\mathfrak m$.
For $u\in\Lambda\cap e\mathbb W$, $ku$ lies in the closure of $\mathfrak mM$, so $\Phi_\kappa(x+ku)=\Phi_\kappa(x)$. If
$\Phi_\kappa(x)\ne0$, then $x=ex_\kappa+y$ with $y$ in the closure of $\mathfrak mM$; $ay\in\Lambda$, and
$aex_\kappa\in\Lambda$, as $e$ is integral and $ax_\kappa\in\Lambda$; so $kx\in\Lambda$. And $\Phi_\kappa$ vanishes off
$e\mathbb W$. So Lemma~\ref{asm-dependence-test-assembly-first-layer-ratio-inputs-standard}(1), at $a=e$, gives
$f_\kappa$.
\end{proof}
% Lean: exists_vac_cos (15 lines), followed: VacuumO.vacExists for f, and for each coset PlaceVec.exists_isThetaAt with
% PlaceModelSum.placeModelExistsSum (the place model of f at e) and isSBHalf_cosInd (about 55 lines, with exists_hom_le
% and exists_le_hom, about 60 lines, for the integers a and b). The Lean's IsThetaAt ... e ... 1 f (cosInd e kappa) is
% the book's f_kappa = T(Phi_kappa) for every place model.

\begin{lemma}\label{asm-dependence-test-assembly-first-layer-grid-pieces-weil-assemble}
Let $D$ be a definite-side datum with set $\mathcal T$ of embeddings, and let a level system for $D$ and a theta grid
for it be given (Definition~\ref{th-grid}), with places $w(v)$, $v\in S_0$, and modulus $\mathfrak m$. Let $\psi$ be a
standard character, $\theta$ an archimedean family whose member at each $\tau\in\mathcal T$ extends $\tau$, and $e$ a
sum of the idempotents of $\mathbf A_E$ at the places $w(v)$, $v\in S_0$, with $e\bar e=0$. Suppose given, for every
space $X$ of the grid, Cartan involutions and orthonormal bases at the archimedean places, and for every grid pair $p$
of (G1) to (G4) in Definition~\ref{th-grid} a Weil datum with standard character $\psi$ and characters
$(\chi_{S,p},\chi_{T,p})$, the first space of the pairs (G1) and (G3) taken with the negated form. Then the grid has a
grid Weil datum with standard
character $\psi$, archimedean family $\theta$, these Cartan involutions and bases and these Weil data, whose Fock
vacuum vectors and coset vectors for the pairs $(V_{i+1,0},V_{i,0})$ are taken at $e$ and $\mathfrak m$.
\end{lemma}
\leanref{Development/AbelianVarietyTheory/Proofs/DivisorProofs/ExistsWeilOver.lean}{GlobalTheory.ChiGrid.WeilOver.exists_weil}
\begin{proof}
Let $0\le i<m$. The spaces $V_{i+1,0}$ and $V_{i,0}$ have correct signatures and self-dual lattices, as spaces of a
theta grid; so do $-V_{i+1,0}$, with the same lattice, and the negated Cartan involutions of $V_{i+1,0}$ are Cartan
involutions of it. Lemma~\ref{asm-dependence-test-assembly-first-layer-grid-pieces-weil-vectors}, for $S=-V_{i+1,0}$,
$T=V_{i,0}$, $e$ and $\mathfrak m$, gives a Fock vacuum vector of the pair and a coset vector for every class modulo
$\mathfrak m$. With $\psi$, $\theta$, the Cartan involutions, the bases and the Weil data, these vectors form a grid
Weil datum (Definition~\ref{th-grid}).
\end{proof}
% Lean: WeilOver.exists_weil (about 30 lines), followed: exists_vac_cos for each i < m with sigOK_negSp,
% isCartanInv_neg and isSelfDualC_neg; the hypothesis hV is the signature and self-duality of the spaces V_{i,0},
% and modul_ne_bot gives m != 0. The structure G.Weil records psi, theta, the Cartan involutions and bases (IsMaj), the
% Weil data, the vacuum vectors, the coset vectors and the idempotent e.

\begin{lemma}\label{asm-dependence-test-assembly-first-layer-grid-pieces-weil-exists}
Assume the existence of the Weil representation of a unitary dual pair, which is among the results of
Theorem~\ref{lit-bundle-theta}. Let $D$ be a definite-side datum, let a level system for $D$ with base tuple $\mu_0$
and a theta grid for it be given (Definition~\ref{th-grid}), and let $\psi$ be a standard character. Then the grid has
a grid Weil datum with standard character $\psi$.
\end{lemma}
\leanref{Development/AbelianVarietyTheory/Proofs/DivisorProofs/ExistsWeilRepresentationVacuumCosetVectorsGridStatement.lean}{GlobalTheory.T.exists_weilRepresentation_vacuum_cosetVectors_grid}
\leanref{Development/AbelianVarietyTheory/Proofs/DivisorProofs/ExistsWeilRepresentationVacuumCosetVectorsGrid.lean}{GlobalTheory.GridWeil.exists_weil}
\leanref{Development/AbelianVarietyTheory/Proofs/DivisorProofs/ExistsWeilRepresentationVacuumCosetVectorsGrid.lean}{GlobalTheory.GridWeil.attChr_ok}
\leanref{Development/AbelianVarietyTheory/Proofs/DivisorProofs/ExistsIdempotentSumAuxiliaryPlaces.lean}{GlobalTheory.GridIdem.exists_idempotentSum_auxiliaryPlaces}
\begin{proof}
Let $E/F$ be the host field and $\eta$ the quadratic character of $E/F$, and let $\theta$ be an archimedean family whose
member at each $\tau\in\mathcal T$ extends $\tau$.

\emph{Cartan involutions.} Every space $X$ of the grid has a correct signature and a self-dual lattice: the spaces
$V_{0,j}$ are the spaces $Y_j$ and $X'$ of Definition~\ref{th-spaces} with the lattices of (L7) in
Definition~\ref{th-grid}, and the others are given with these properties.
At each $\tau$, the spectral theorem diagonalizes the Hermitian matrix $\tau(G_X)$ by a matrix $B_\tau$ with
$\bar B_\tau^{\mathsf T}\tau(G_X)B_\tau$ a diagonal matrix $J_\tau$ of signs; this gives a Cartan involution with an
orthonormal basis.

\emph{Weil data.} The characters $(\chi_{S,p},\chi_{T,p})$ of a grid pair $p$ are $c_p$ or $c_p\,\eta_L^{-1}$, according
as the rank of the space is even or odd, where $c_p$ is the character $\chi_{m-i}$ of the grid for the pairs (G1) and
(G3), and the character $\mu_{t,j+1}=\mu^0_{j+1}\xi_t$ of the level system for the pairs (G2) and (G4)
(Definition~\ref{th-grid}). Each $c_p$ is a unitary Hecke character trivial on $\mathbf A_F^\times$: the $\chi_i$ by
condition (a) of Definition~\ref{th-grid}, the $\mu^0_j$ by (L2) and the $\xi_t$, of finite order, by (L4). The character
$\eta_L$ is a unitary Hecke character whose restriction to $\mathbf A_F^\times$ is $\eta$, by (L3). As $\eta^2=1$, $c\,\eta_L^{-1}$ restricts to
$\eta^{-1}=\eta$ and $c$ to $1$; in both cases to $\eta^n$, with $n$ the rank. So the Weil representation of the unitary
dual pair gives a Weil datum for $p$ with standard character $\psi$ and these characters.

\emph{The idempotent.} Each place $w(v)$, $v\in S_0$, is not conjugate to itself
(Definition~\ref{th-grid-places}(1)), so $E\otimes_FF_v$ is the product of $E_{w(v)}$ and $E_{\overline{w(v)}}$, and
the unit of the first factor is an idempotent $e_{w(v)}$ of $\mathbf A_E$. No two of these places are conjugate, as
they lie over distinct places of $F$. So $e=\sum_{v\in S_0}e_{w(v)}$ satisfies $e\bar e=0$.

Lemma~\ref{asm-dependence-test-assembly-first-layer-grid-pieces-weil-assemble}, for $\psi$, $\theta$, $e$ and these
Cartan involutions, bases and Weil data, gives the grid Weil datum.
\end{proof}
% Lean: GridWeil.exists_weil (about 95 lines), followed: Member.exists_isInfFam_family for theta; G.isV and
% ChiGrid.colY_isV for the signatures and lattices (hV); GridWeil.exists_cartan for the Cartan involutions and bases;
% the characters G.chr p are attChr c eta_L n (Lift.attChr: c * eta_L^-1 for odd n, c for even n), with c = chi (m - i)
% (hchi, from chi_eq and triv) or L.mu t (j + 1) (hmu, from LevCh.Is.mu0, LevCh.C.isMu0 and xi_triv); the restriction
% of eta_L is Mem.EtaTilde.isEta_res for LevCh.C.isEta; attChr_ok (with IdeleNormChar.existsUnique_isEta) gives
% IsHeckeChar and ResEta n; the Weil data are the hypothesis hW, Published.WeilRepresentationUnitaryPair; the idempotent
% is GridIdem.exists_idempotentSum_auxiliaryPlaces (MovedIdem through P.w_split). The Lean builds the vacuum and coset
% vectors inline (VacuumO.vacExists, PlaceVec.exists_isThetaAt, placeModelExistsSum, isSBHalf_cosInd), as
% Lemma -weil-assemble does. The grid pairs and their characters are (G1)-(G4) of Definition th-grid, and the
% conditions on mu^0_j, eta_L and xi_t are (L2), (L3) and (L4) there.

\begin{definition}\label{asm-dependence-test-assembly-first-layer-grid-pieces-choice-standard}
Let $D$ be a definite-side datum with host field $E/F$ and with integers $Q_S$, $q$ and $d$, ideal $\mathfrak c$ of
$\mathcal O_K$ and level primes $g_0,g_1$ of its host, and put $N=Q_Sqd$. Let a level system for $D$, with the
characters $\mu_{t,j}$ of (L4) and the sets $Q_t$ of (L5), and a theta grid for it, with distinguished places $S_0$ and
lattices $\Lambda_X$ of its spaces $X$, be given (Definition~\ref{th-grid}), and let $U$ be a set of tuples. For an
$\mathcal O_E$-module $L\subset E^n$ write $\hat L\subset\mathbf A_{E,f}^n$ for its closure, the $\hat{\mathcal
O}_E$-module that it spans.
\begin{enumerate}
\item The \emph{definite} grid spaces are $V_{i,m}$ for $0\le i\le m$ and $Q_j$ for $0\le j\le m$. The \emph{definite} grid
  pairs are the pairs $(V_{i+1,m},V_{i,m})$, $0\le i<m$, of (G1) with $j=m$, and the pairs $(Q_j,Q_{j+1})$ at a tuple,
  $0\le j<m$, of (G4).
\item Let $\Sigma$ be a set of primes. A place of $E$ is \emph{off} $\Sigma$ if it contains no prime of $\Sigma$. For a
  Hermitian space $X$ of $D$, an element $u$ of $U(X)(\mathbf A_{F,f})$ is \emph{supported off} $\Sigma$ if there is a
  finite set $P$ of places of $E$ off $\Sigma$, none of them conjugate to itself or to another place of $P$, such that $u$
  has component $1$ at every place of $F$ below no place of $P$.
\item \emph{Standard data} on $U$ consist of a finite set $\Sigma$ of primes and, for every grid space $X$, a matrix
  $\phi_X\in GL_{\operatorname{rk}X}(E)$ and an element $u_X\in U(X)(\mathbf A_{F,f})$, such that:
  \begin{enumerate}
  \item[(a)] every place of $E$ off $\Sigma$ lies over no place of $S_0$ and over no place of $F$ below a place of a set
    $Q_t$ with $t\in U$, and does not contain $Ng_0g_1\mathrm N\mathfrak c$;
  \item[(b)] for every place $w$ of $E$ off $\Sigma$ that is not conjugate to itself, every $t\in U$ and every $1\le j\le
    m$, the character $\mu_{t,j}$ is trivial on the units of $E_w$;
  \item[(c)] for every definite grid space $X$: the columns of $\phi_X$ form an orthonormal basis of $X$, that is
    $\bar\phi_X^{\mathsf T}G_X\phi_X=1$; $u_X$ is supported off $\Sigma$; and $u_X\hat L_X=\hat\Lambda_X$, where $L_X$ is the
    $\mathcal O_E$-span of the columns of $\phi_X$;
  \item[(d)] for the embedding $\iota\colon V_{m,m}\to Q_m$ of the grid and the vector $a$ of the grid that spans the
    orthogonal complement of its image: the first $\operatorname{rk}V_{m,m}$ columns of $\phi_{Q_m}$ are the images under
    $\iota$ of the columns of $\phi_{V_{m,m}}$, and its last column is $a$; $u_{Q_m}\iota=\iota u_{V_{m,m}}$; and
    $u_{Q_m}a=a$.
  \end{enumerate}
\end{enumerate}
\end{definition}
\leanref{Development/AbelianVarietyTheory/Definitions/Divisors/KernelNormConstantSpec.lean}{GlobalTheory.FunObj.Std}
\leanref{Development/AbelianVarietyTheory/Definitions/Divisors/KernelNormConstantSpec.lean}{GlobalTheory.FunObj.DefNode}
\leanref{Development/AbelianVarietyTheory/Definitions/Divisors/KernelNormConstantSpec.lean}{GlobalTheory.FunObj.DefPair}
\leanref{Development/AbelianVarietyTheory/Definitions/Divisors/KernelNormConstantSpec.lean}{GlobalTheory.FunObj.Off}
\leanref{Development/AbelianVarietyTheory/Definitions/Divisors/KernelNormConstantSpec.lean}{GlobalTheory.FunObj.UnrIdem}

\begin{definition}\label{asm-dependence-test-assembly-first-layer-grid-pieces-choice-normal}
Let $D$ be a definite-side datum with host field $E$ of degree $2f$, host element $\xi_0$, embedding $\varphi_0$, set
$\mathcal T$ of types, and ideals $\mathfrak c$ of $\mathcal O_K$ and $\mathfrak b$ of $\mathcal O_E$ of its host
(Definition~\ref{hf-host}). Use the notation of Definition~\ref{is-vacuum-theta-vectors}, with $\mathfrak b^{-1}$ for
its ideal $\mathfrak A$: the spaces $V_{a,b}$ with coordinates $z_{\tau,\beta,\alpha}$, the subgroups $\Lambda_{L,L'}$ and
$\Lambda_v$, theta functions, norms, averaging operators $I_{\Gamma\to\Gamma'}$, tensor powers $g_v$, the functions
$\vartheta_{\Lambda,t}$, the numbers $n(\Lambda,t)$, orthonormal bases, unkilled two-torsion points, theta vectors up to
units and CM lattices of the host.
\begin{enumerate}
\item Let $\Lambda=\delta\varphi_0(\mathfrak c')$ be a CM lattice of the host, $t_0$ an unkilled two-torsion point of
  $\Lambda$, and $a,b\ge0$ integers. An \emph{odd vacuum function} of size $(a,b)$ for $(\mathfrak c',\delta,t_0)$ is a
  function $e$ on $V_{a,b}$ such that:
  \begin{enumerate}
  \item[(a)] $e$ is a theta function for $\Lambda_{\mathcal O_E^a,\mathcal O_E^b}$, and its Taylor coefficients at $0$ are
    algebraic numbers;
  \item[(b)] $\|e\|^2_{\Lambda_{\mathcal O_E^a,\mathcal O_E^b}}=n(\Lambda,t_0)^{fab}$, and $e$ is a theta vector up to units
    for $(t_0,\mathcal O_E^a,\mathcal O_E^b)$;
  \item[(c)] for every prime $p$ prime to $\mathrm N\mathfrak c'$ and every embedding $\lambda$ of $\overline{\mathbf Q}$
    into $\mathbf C_p$ there are an orthonormal basis $v=(v_k)$ and an algebraic number $\varepsilon$ with
    $|\lambda(\varepsilon)|_p=1$ such that: the $v_k$ lie in $Y\otimes\mathbf Z_{(p)}$ and span it over
    $\mathcal O_K\otimes\mathbf Z_{(p)}$, where $Y$ is the $\mathcal O_E$-span of the matrices $sB$ with
    $s\in\mathfrak b^{-1}$ and $B\in M_{b\times a}(\mathcal O_E)$; the index of
    $\Lambda_v\cap\Lambda_{\mathcal O_E^a,\mathcal O_E^b}$ in $\Lambda_{\mathcal O_E^a,\mathcal O_E^b}$ is prime to $p$; and
    $e=\varepsilon\,I_{\Lambda_v\to\Lambda_{\mathcal O_E^a,\mathcal O_E^b}}\big((\vartheta_{\Lambda,t_0})_v\big)$.
  \end{enumerate}
\item Let a theta grid for a level system for $D$, a grid Weil datum $W$ on it with Cartan involutions at the archimedean
  places, and standard data $(\Sigma,\phi_X,u_X)$ on a set $U$ of tuples be given (Definitions~\ref{th-grid}
  and~\ref{asm-dependence-test-assembly-first-layer-grid-pieces-choice-standard}). Let $(S,T)$ be a definite grid pair,
  the first space of $(V_{i+1,m},V_{i,m})$ taken with the negated form and Cartan involution, as for the pairs (G1) in
  Definition~\ref{th-grid}. For a function $g$ on $\mathbb W(S,T)$ and a polynomial $P$ in the coordinates
  $z_{\tau,\beta,\alpha}$ of $V_{\operatorname{rk}S,\operatorname{rk}T}$, write $g\otimes_\phi P$ for the Fock vector
  $g\otimes P$ formed with the Cartan involutions of $W$ on $S$ and $T$ and, at each $\tau\in\mathcal T$, the bases
  $\tau(\phi_S)$ of $S$ and $\tau(\phi_T)$ of $T$, the coordinate $z_{\tau,\beta,\alpha}$ standing for the creation
  operator of the entry $(\beta,\alpha)$ at $\tau$. Put $\pi_\tau(\delta)=\delta/\sqrt{2\pi|\tau(\xi_0)|}$, and let
  $\epsilon_{S,T}$ be complex conjugation if $(S,T)=(V_{i+1,m},V_{i,m})$ and the identity if $(S,T)=(Q_j,Q_{j+1})$.
\item A \emph{vacuum normalization} for these standard data consists of $\delta\in\mathbf C$ with
  $\delta^{12}=\Delta(\varphi_0(\mathfrak c))$, an unkilled two-torsion point $t_0$ of $\delta\varphi_0(\mathfrak c)$, and,
  for every grid pair $(S,T)$, at a tuple for the pairs that have one, a function $e_{S,T}$ on
  $V_{\operatorname{rk}S,\operatorname{rk}T}$, a function $g_{S,T}$ on $\mathbb W(S,T)$ and $c_{S,T}\in\mathbf C^\times$,
  such that these three at the pairs $(Q_j,Q_{j+1})$ do not depend on the tuple, and for every definite grid pair $(S,T)$:
  \begin{enumerate}
  \item[(a)] $e_{S,T}$ is an odd vacuum function of size $(\operatorname{rk}S,\operatorname{rk}T)$ for
    $(\mathfrak c,\delta,t_0)$;
  \item[(b)] $g_{S,T}$ is a Fock vacuum vector for the Cartan involutions of $W$ on $S$ and $T$ and the lattice of the pair
    given by $L_S$ and $L_T$, the $\mathcal O_E$-spans of the columns of $\phi_S$ and $\phi_T$;
  \item[(c)] for every monomial $z^k=\prod_{\tau,\beta,\alpha}z_{\tau,\beta,\alpha}^{k_{\tau,\beta,\alpha}}$,
    \[
     (g_{S,T}\otimes_\phi z^k)(0)=c_{S,T}\,\epsilon_{S,T}\Big(\prod_{\tau,\beta,\alpha}
     \pi_\tau(\delta)^{k_{\tau,\beta,\alpha}}\cdot(\partial^ke_{S,T})(0)\Big),\qquad
     \partial^k=\prod_{\tau,\beta,\alpha}\partial_{z_{\tau,\beta,\alpha}}^{k_{\tau,\beta,\alpha}}.
    \]
  \end{enumerate}
\end{enumerate}
\end{definition}
\leanref{Formalism/SharpRestriction/GlobalSchurBounds/StandardLatticesGL3.lean}{Formalism.MainNotions.StandardLatticesGL3.oddE0}
\leanref{Formalism/SharpRestriction/GlobalSchurBounds/StandardLatticesGL3.lean}{Formalism.MainNotions.StandardLatticesGL3.LocForm}
\leanref{Formalism/SharpRestriction/GlobalSchurBounds/StandardLatticesGL3.lean}{Formalism.MainNotions.StandardLatticesGL3.Crv}
\leanref{Development/AbelianVarietyTheory/Definitions/Divisors/KernelNormConstantSpec.lean}{GlobalTheory.FunObj.Normal}
\leanref{Development/AbelianVarietyTheory/Definitions/Divisors/KernelNormConstantSpec.lean}{GlobalTheory.FunObj.fockIn}
\leanref{Development/AbelianVarietyTheory/Definitions/Divisors/KernelNormConstantSpec.lean}{GlobalTheory.FunObj.period}
\leanref{Development/AbelianVarietyTheory/Definitions/Divisors/KernelNormConstantSpec.lean}{GlobalTheory.FunObj.opp}

\begin{definition}\label{asm-dependence-test-assembly-first-layer-grid-pieces-choice-grid}
Let $D$ be a definite-side datum with host field $E/F$ and with ideals $\mathfrak c$ of $\mathcal O_K$ and $\mathfrak b$ of
$\mathcal O_E$ of its host, and let $\mathfrak A=\mathfrak c\mathcal O_E\,\mathfrak b^{-1}$ be the ideal of the host. Let a
level system for $D$ with sets $Q_t$, a theta grid for it with places $w(v)$ and depths $d(v)$, $v\in S_0$, and lattices
$\Lambda_X$, and a grid Weil datum $W$ on the grid be given (Definition~\ref{th-grid}), and let $U$ be a set of tuples.
Write $v_w(\mathfrak a)$ for the exponent of a finite place $w$ of $E$ in a fractional ideal $\mathfrak a$, and $\hat L$ as
in Definition~\ref{asm-dependence-test-assembly-first-layer-grid-pieces-choice-standard}.
\begin{enumerate}
\item The \emph{local places} at a tuple $t$ are the places $w(v)$, $v\in S_0$, and the places of $Q_t$. For a finite
  place $w$ of $E$ put $d_w=d(v)$ if $w=w(v)$ with $v\in S_0$ and $d_w=0$ otherwise, and let $\varpi_w$ be a fixed
  uniformizer of $E_w$.
\item An \emph{adelic basis} of an $\mathcal O_E$-lattice $\Lambda\subset E^n$ is a matrix $\beta\in
  GL_n(\mathbf A_{E,f})$ with $\hat\Lambda=\beta\hat{\mathcal O}_E^n$. Let $\iota\colon A\to B$ be an isometric embedding
  of Hermitian spaces of $D$ whose orthogonal complement has the basis $a_1,\dots,a_c$, let $\beta_A\in
  GL_{\operatorname{rk}A}(\mathbf A_{E,f})$ and $\beta_B\in GL_{\operatorname{rk}B}(\mathbf A_{E,f})$, let $w$ be a finite
  place of $E$ and $u_1,\dots,u_c\in E_w$. Then $\beta_B$ is the \emph{sum} of $\beta_A$ and $(u_ka_k)_k$ at $w$ if the
  first $\operatorname{rk}A$ columns of the component $\beta_{B,w}$ are the images under $\iota$ of the columns of
  $\beta_{A,w}$, and its column $\operatorname{rk}A+k$ is $u_ka_k$ for $1\le k\le c$.
\item A \emph{grid choice} for $W$ on $U$ consists of adelic bases $\beta_{i,j}$ of $\Lambda_{V_{i,j}}$ for $1\le i\le m$
  and $0\le j\le m$, adelic bases $\beta^+_j$ of $\Lambda_{Q_j}$ for $0\le j\le m$, standard data on $U$
  (Definition~\ref{asm-dependence-test-assembly-first-layer-grid-pieces-choice-standard}), a vacuum normalization for them
  (Definition~\ref{asm-dependence-test-assembly-first-layer-grid-pieces-choice-normal}(3)), and a finite set
  $P_{\mathfrak A}$ of places of $E$, such that:
  \begin{enumerate}
  \item[(a)] for every $t\in U$, every local place $w$ at $t$ and $0\le j<m$, $\beta_{m,j}$ is the sum of $\beta^+_j$ and
    $(\epsilon_k\varpi_w^{e_k}a_k)_{1\le k\le m-1-j}$ at $w$, for the embedding $Q_j\subset V_{m,j}$ of the grid ($P_m\subset
    V_{m,0}$ if $j=0$) and the basis $(a_k)$ of its complement, where $\epsilon_k=-1$ if $2k>m-j$ and $\epsilon_k=1$
    otherwise, and $e_k=kd_w$ if $j=0$ and $e_k=0$ if $j\ge1$;
  \item[(b)] for every $t\in U$ and every local place $w$ at $t$, $\beta^+_m$ is the sum of $\beta_{m,m}$ and $a$ at $w$, for
    the embedding $V_{m,m}\subset Q_m$ of the grid and the vector $a$ that spans its complement;
  \item[(c)] every $w\in P_{\mathfrak A}$ has $v_w(\mathfrak A)\neq0$ or $v_w(\bar{\mathfrak A})\neq0$, and for every
    finite place $w$ with $v_w(\mathfrak A)\neq0$ exactly one place of $P_{\mathfrak A}$ lies over the same place of $F$ as
    $w$.
  \end{enumerate}
\end{enumerate}
\end{definition}
\leanref{Development/AbelianVarietyTheory/Definitions/Divisors/InvVolConjLevelOne.lean}{GlobalTheory.FunObj.GridChoice}
\leanref{Development/AbelianVarietyTheory/Definitions/Divisors/InvVolConjLevelOne.lean}{GlobalTheory.FunObj.IsSumFrame}
\leanref{Development/AbelianVarietyTheory/Definitions/Divisors/InvVolConjLevelOne.lean}{GlobalTheory.FunObj.frameSign}
\leanref{Development/AbelianVarietyTheory/Definitions/Divisors/ExistsGivenStageEqOne.lean}{GlobalTheory.FunObj.Frames}
\leanref{Development/AbelianVarietyTheory/Definitions/Divisors/ExistsGivenStageEqOne.lean}{GlobalTheory.FunObj.locPl}
\leanref{Development/AbelianVarietyTheory/Definitions/Divisors/ExistsGivenStageEqOne.lean}{GlobalTheory.FunObj.cosExp}
\leanref{Formalism/LocalToGlobal/Divisibility/DefGroup.lean}{Formalism.L2.DIV.IsFrame}

\begin{lemma}\label{asm-dependence-test-assembly-first-layer-grid-pieces-choice-primes}
Let $D$ be a definite-side datum with host field $E/F$ and with integers $Q_S$, $q$ and $d$, ideal $\mathfrak c$ of
$\mathcal O_K$ and level primes $g_0,g_1$ of its host, let a level system for $D$ with sets $Q_t$ and a theta grid for it
with distinguished places $S_0$ be given (Definition~\ref{th-grid}), and let $U$ be a finite set of tuples. Then there is
a finite set $\Sigma$ of primes such that every place of $E$ off $\Sigma$ lies over no place of $S_0$ and over no place of
$F$ below a place of a set $Q_t$ with $t\in U$, and does not contain $Q_Sqdg_0g_1\mathrm N\mathfrak c$.
\end{lemma}
\leanref{Development/AbelianVarietyTheory/Proofs/DivisorProofs/ExistsFinitePrimeSetAvoidingGridPlacesStatement.lean}{GlobalTheory.T.exists_finite_primeSet_avoiding_gridPlaces}
\leanref{Development/AbelianVarietyTheory/Proofs/DivisorProofs/ExistsFinitePrimeSetAvoidingGridPlaces.lean}{GlobalTheory.FunObj.exists_std_primes}
\begin{proof}
Let $M$ be the product of $Q_Sqdg_0g_1\mathrm N\mathfrak c$, of the norms $\mathrm Nv$ for $v\in S_0$, and of the norms
$\mathrm Nv$ of the places $v$ of $F$ below the places of the sets $Q_t$, $t\in U$. It is a nonzero integer, as these sets
are finite. Let $\Sigma$ be the set of prime factors of $M$, and $w$ a place of $E$ off $\Sigma$. No divisor $n$ of $M$
lies in $w$: otherwise $M\in w$, and the prime below $w$ divides $M$ and lies in $\Sigma$. If $w$ lay over a place $v$ of
$F$ in $S_0$ or below a place of $Q_t$, $t\in U$, then $\mathrm Nv\in v\subset w$, and $\mathrm Nv$ divides $M$. Finally
$Q_Sqdg_0g_1\mathrm N\mathfrak c$ divides $M$.
\end{proof}

\begin{lemma}\label{asm-dependence-test-assembly-first-layer-grid-pieces-choice-unramified}
Let $E$ be the host field of a definite-side datum and $\chi$ a continuous homomorphism from $\mathbf A_E^\times$ to the
unit circle. Then there is a finite set $\Sigma$ of primes such that $\chi$ is trivial on the units of $E_w$ for every
place $w$ of $E$ off $\Sigma$ that is not conjugate to itself.
\end{lemma}
\leanref{Development/AbelianVarietyTheory/Proofs/DivisorProofs/CharacterUnramifiedOffFinitePrimesStatement.lean}{GlobalTheory.T.character_unramified_off_finite_primes}
\leanref{Development/AbelianVarietyTheory/Proofs/DivisorProofs/CharacterUnramifiedOffFinitePrimes.lean}{GlobalTheory.UnrOff.unr_off}
\leanref{Development/AbelianVarietyTheory/Proofs/AdelicLatticeProofs/ContinuousCharacterFiniteOrderIntegralUnits.lean}{GlobalTheory.CharFin.char_finUnit_finite_order}
\begin{proof}
The unit circle has a neighbourhood of $1$ that contains no subgroup other than $\{1\}$. By continuity its inverse image
contains the open subgroup of the ideles with archimedean component $1$ and finite part in $1+N\hat{\mathcal O}_E$, for
some integer $N\ge1$; the image of this subgroup is a subgroup in that neighbourhood, so $\chi$ is trivial on it. Let
$\Sigma$ be the set of prime factors of $N$, let $w$ be a place off $\Sigma$ with $w\neq\bar w$, and let $x$ be a unit of
$E_w$, viewed as an idele with component $1$ at the other places. No prime factor of $N$ lies in $w$, so $N$ is a unit of
$\mathcal O_{E_w}$ and $x-1=N\cdot N^{-1}(x-1)$ with $N^{-1}(x-1)$ integral. So $x$ lies in that subgroup, and
$\chi(x)=1$.
\end{proof}

\begin{lemma}\label{asm-dependence-test-assembly-first-layer-grid-pieces-choice-orthonormal-pieces}
Let $D$ be a definite-side datum with host field $E/F$, and $X$ a Hermitian space of $D$.
\begin{enumerate}
\item If $X$ has correct signature and signature $(\operatorname{rk}X,0)$ at every $\tau\in\mathcal T$, then $X$ has an
  orthonormal basis: there is $\phi\in GL_{\operatorname{rk}X}(E)$ with $\bar\phi^{\mathsf T}G_X\phi=1$.
\item Let $\phi\in GL_{\operatorname{rk}X}(E)$ with $\bar\phi^{\mathsf T}G_X\phi=1$, let $L$ be the $\mathcal O_E$-span of
  its columns, and let $\Lambda$ be a self-dual lattice of $X$. Then there is $x\in U(X)(\mathbf A_{F,f})$ with
  $x\hat L=\hat\Lambda$.
\item Let $\phi\in GL_{\operatorname{rk}X}(E)$, let $L$ be the $\mathcal O_E$-span of its columns, let $\Sigma$ be a finite
  set of primes and $x\in U(X)(\mathbf A_{F,f})$. Then $x=\gamma hk$ with $\gamma\in U(X)(F)$, $h\in
  U(X)(\mathbf A_{F,f})$ supported off $\Sigma$
  (Definition~\ref{asm-dependence-test-assembly-first-layer-grid-pieces-choice-standard}(2)), and $k\in
  U(X)(\mathbf A_{F,f})$ with $k\hat L=\hat L$.
\end{enumerate}
\end{lemma}
\leanref{Development/AbelianVarietyTheory/Proofs/DivisorProofs/ExistsOrthonormalBasisDefiniteHermitianSpaceStatement.lean}{GlobalTheory.T.exists_orthonormalBasis_definite_hermitianSpace}
\leanref{Development/AbelianVarietyTheory/Proofs/DivisorProofs/ExistsFiniteAdelicUnitaryMovesSelfDualLatticeStatement.lean}{GlobalTheory.T.exists_finiteAdelicUnitary_moves_selfDualLattice}
\leanref{Development/AbelianVarietyTheory/Proofs/DivisorProofs/IdempotentSupportHermitianDecompositionStatement.lean}{GlobalTheory.T.idempotent_support_hermitian_decomposition}
\begin{proof}
\openstep{The proofs of the three statements are not yet written out: the orthonormal basis of a totally definite
Hermitian space, the finite adelic unitary element carrying the lattice of an orthonormal basis to a self-dual lattice,
and the decomposition of a finite adelic unitary element into a rational element, an element supported off $\Sigma$ and
an element fixing $\hat L$.}
% Lean, sizes from tools/lean_size.py: (1) T.exists_orthonormalBasis_definite_hermitianSpace, 3,191 lines; (2)
% T.exists_finiteAdelicUnitary_moves_selfDualLattice, 4,130 lines; (3) T.idempotent_support_hermitian_decomposition,
% 1,549 lines; 6,845 lines together (--union).
\end{proof}

\begin{lemma}\label{asm-dependence-test-assembly-first-layer-grid-pieces-choice-orthonormal}
Let $D$ be a definite-side datum, $X$ a Hermitian space of $D$ of correct signature with signature
$(\operatorname{rk}X,0)$ at every $\tau\in\mathcal T$ and with a self-dual lattice $\Lambda$, and $\Sigma$ a finite set of
primes. Then there are $\phi\in GL_{\operatorname{rk}X}(E)$ with $\bar\phi^{\mathsf T}G_X\phi=1$ and $u\in
U(X)(\mathbf A_{F,f})$ supported off $\Sigma$ with $u\hat L=\hat\Lambda$, where $L$ is the $\mathcal O_E$-span of the
columns of $\phi$.
\end{lemma}
\leanref{Development/AbelianVarietyTheory/Proofs/DivisorProofs/ExistsOrthonormalBasisSplitSupportedAdelicPointStatement.lean}{GlobalTheory.T.exists_orthonormalBasis_splitSupported_adelicPoint}
\leanref{Development/AbelianVarietyTheory/Proofs/DivisorProofs/ExistsOrthonormalBasisSplitSupportedAdelicPoint.lean}{GlobalTheory.FunObj.StdVertex.exists_std_vertex}
\leanref{Development/AbelianVarietyTheory/Proofs/DivisorProofs/RationalUnitaryConjSplitSupportedFiniteAdelicPoiStatement.lean}{GlobalTheory.T.rationalUnitary_conj_splitSupported_finiteAdelicPoint}
\begin{proof}
By Lemma~\ref{asm-dependence-test-assembly-first-layer-grid-pieces-choice-orthonormal-pieces}(1) there is
$\phi_0\in GL_{\operatorname{rk}X}(E)$ with $\bar\phi_0^{\mathsf T}G_X\phi_0=1$; let $L_0$ be the span of its columns. By
part (2) there is $x\in U(X)(\mathbf A_{F,f})$ with $x\hat L_0=\hat\Lambda$, and by part (3), for $\phi_0$ and $\Sigma$,
$x=\gamma hk$ with $\gamma$, $h$ and $k$ as there. Put $\phi=\gamma\phi_0$ and $u=\gamma h\gamma^{-1}$. As $\gamma$ is
unitary, $\bar\phi^{\mathsf T}G_X\phi=\bar\phi_0^{\mathsf T}G_X\phi_0=1$. As $\gamma$ is rational, $u$ has component $1$
wherever $h$ has, so $u$ is supported off $\Sigma$, with the same set of places. The span of the columns of $\phi$ is
$\gamma L_0$, so $\hat L=\gamma\hat L_0$, and $u\hat L=\gamma h\hat L_0=\gamma hk\hat L_0=x\hat L_0=\hat\Lambda$, as
$k\hat L_0=\hat L_0$.
% Lean: StdVertex.exists_std_vertex, followed: exists_orthonormalBasis_definite_hermitianSpace (fr0),
% exists_finiteAdelicUnitary_moves_selfDualLattice (x), idempotent_support_hermitian_decomposition (gamma, h, k),
% rationalUnitary_conj_splitSupported_finiteAdelicPoint (u, 223 lines), mem_completedLattice_columnLattice_mul_iff
% (hat of the span of gamma fr0, 109 lines).
\end{proof}

\begin{lemma}\label{asm-dependence-test-assembly-first-layer-grid-pieces-choice-top}
Let $D$ be a definite-side datum, $\iota\colon A\to B$ an isometric embedding of Hermitian spaces of $D$ whose
orthogonal complement is spanned by a vector $a$ with $\langle a,a\rangle=1$, $\Lambda_A\subset A$ an
$\mathcal O_E$-module, $\Lambda_B=\iota(\Lambda_A)+\mathcal O_Ea$, and $\Sigma$ a finite set of primes. Let
$\phi_A\in GL_{\operatorname{rk}A}(E)$ with $\bar\phi_A^{\mathsf T}G_A\phi_A=1$, and let $u_A\in U(A)(\mathbf A_{F,f})$
be supported off $\Sigma$ with $u_A\hat L_A=\hat\Lambda_A$, where $L_A$ is the $\mathcal O_E$-span of the columns of
$\phi_A$. Then there are $\phi_B\in GL_{\operatorname{rk}B}(E)$ with $\bar\phi_B^{\mathsf T}G_B\phi_B=1$ and $u_B\in
U(B)(\mathbf A_{F,f})$ supported off $\Sigma$ with $u_B\hat L_B=\hat\Lambda_B$, for the span $L_B$ of the columns of
$\phi_B$, such that the first $\operatorname{rk}A$ columns of $\phi_B$ are the images under $\iota$ of the columns of
$\phi_A$, its last column is $a$, $u_B\iota=\iota u_A$ and $u_Ba=a$.
\end{lemma}
\leanref{Development/AbelianVarietyTheory/Proofs/DivisorProofs/GridTopVertexFramePointTransportLatticeStatement.lean}{GlobalTheory.T.gridTopVertex_frame_point_transport_lattice}
\leanref{Development/AbelianVarietyTheory/Proofs/DivisorProofs/GridTopVertexFramePointTransportLattice.lean}{GlobalTheory.FunObj.top_transport}
\leanref{Development/AbelianVarietyTheory/Proofs/DivisorProofs/GridTopVertexFramePointTransportLattice.lean}{GlobalTheory.TopExt.top_abs}
\begin{proof}
The space $B$ is the orthogonal sum of $\iota(A)$ and $Ea$. Let $\phi_B$ have the columns $\iota(\phi_Ae_k)$ and then $a$.
It is invertible, and its Gram matrix is that of the columns of $\phi_A$, which is $1$, with the entry $\langle a,a\rangle=1$
added. Let $u_B$ act as $\iota u_A\iota^{-1}$ on $\iota(A)$ and as the identity on $Ea$. It lies in
$U(B)(\mathbf A_{F,f})$, has component $1$ wherever $u_A$ has, so it is supported off $\Sigma$ with the set of places of
$u_A$, and $u_B\iota=\iota u_A$, $u_Ba=a$. Finally $L_B=\iota(L_A)+\mathcal O_Ea$, so
$u_B\hat L_B=\iota(u_A\hat L_A)+\hat{\mathcal O}_Ea=\iota(\hat\Lambda_A)+\hat{\mathcal O}_Ea=\hat\Lambda_B$.
% Lean: TopExt.top_abs for d : HermitianOrthogonalExtensionData A B 1 1 and Lambda_B = d.lat Lambda_A (fun _ => 1):
% frG d fr0 (frM_on, frM_iC, frM_iL), the point g of exists_pt (extM of u_0 with the identity on a; extM_supp,
% extM_mul, extM_mulVec) and hat_transport for the lattices.
\end{proof}

\begin{lemma}\label{asm-dependence-test-assembly-first-layer-grid-pieces-choice-standard-exists}
Let $D$ be a definite-side datum, let a level system for $D$ with base tuple $\mu_0$ and a theta grid for it be given
(Definition~\ref{th-grid}), and let $U$ be a finite set of tuples. Then there are standard data on $U$
(Definition~\ref{asm-dependence-test-assembly-first-layer-grid-pieces-choice-standard}).
\end{lemma}
\leanref{Development/AbelianVarietyTheory/Proofs/DivisorProofs/DefNodeIsV.lean}{GlobalTheory.FunObj.exists_std}
\leanref{Development/AbelianVarietyTheory/Proofs/DivisorProofs/DefNodeIsV.lean}{GlobalTheory.FunObj.defNode_isV}
\begin{proof}
\emph{The primes.} Lemma~\ref{asm-dependence-test-assembly-first-layer-grid-pieces-choice-primes} gives a finite set
$\Sigma_1$ of primes with condition (a). The characters $\mu^0_j$, $1\le j\le m$, and $\xi_t$, $t\in U$, of the level
system are unitary Hecke characters, by (L2) and (L4) of Definition~\ref{th-grid}, so continuous with values in the unit
circle, and
Lemma~\ref{asm-dependence-test-assembly-first-layer-grid-pieces-choice-unramified} gives finite sets $\Sigma(\mu^0_j)$ and
$\Sigma(\xi_t)$. Let $\Sigma$ be the union of $\Sigma_1$ and these sets; it is finite, as $U$ is. A place off $\Sigma$ is
off $\Sigma_1$, so (a) holds. For a place $w$ off $\Sigma$ with $w\neq\bar w$, the characters $\mu^0_j$ and $\xi_t$ are
trivial on the units of $E_w$, hence so is $\mu_{t,j}=\mu^0_j\xi_t$; this is (b).

\emph{The definite spaces.} Every definite grid space $X$ has correct signature and signature $(\operatorname{rk}X,0)$ at
every $\tau$, and its lattice $\Lambda_X$ is self-dual: $V_{0,m}=X'$ has the signature of $Y_m$ by the conditions of the
datum and the self-dual lattice $\Lambda_{X'}$ of (L7); $V_{i,m}$, $i\ge1$, is a space of Definition~\ref{th-grid}(2)
with $j=m$; and $Q_j$ is a space of Definition~\ref{th-grid}(4). For each of them
Lemma~\ref{asm-dependence-test-assembly-first-layer-grid-pieces-choice-orthonormal}, for $\Sigma$, gives $\phi^0_X$ and
$u^0_X$ with the properties of (c).

\emph{The top space.} The complement of $\iota(V_{m,m})$ in $Q_m$ is spanned by a vector $a$ with $\langle a,a\rangle=1$,
the last embedding of Definition~\ref{th-grid}, and $\Lambda_{Q_m}=\iota(\Lambda_{V_{m,m}})+\mathcal O_Ea$, by the
lattices there ($I_1=\mathcal O_E$).
Lemma~\ref{asm-dependence-test-assembly-first-layer-grid-pieces-choice-top}, for $\phi^0_{V_{m,m}}$ and
$u^0_{V_{m,m}}$, gives $\phi_{Q_m}$ and $u_{Q_m}$ with the properties of (c) and (d).

\emph{The data.} Put $\phi_X=\phi^0_X$ and $u_X=u^0_X$ for the definite spaces $X\neq Q_m$, and $\phi_X=1$ and $u_X=1$
for the spaces that are not definite. Then (a) to (d) hold.
% Lean: FunObj.exists_std (82 lines), followed: T.exists_finite_primeSet_avoiding_gridPlaces (S1),
% T.character_unramified_off_finite_primes applied to L.mu0 j (continuity hL.mu0) and L.xi t (hL.xi) for Sc, off_mono and
% unrIdem_mul for (a) and (b); defNode_isV (colY_isV with hL.betaX, G.isV) and
% T.exists_orthonormalBasis_splitSupported_adelicPoint for fr0, u0; T.gridTopVertex_frame_point_transport_lattice (G.sp.topDec
% with c = 1, eps = 1, and G.sp.topLat) for the top; Function.update for fr and u. The Lean takes fr X = 1 and u X = 1 at the
% nodes that are not definite.
\end{proof}

\begin{lemma}\label{asm-dependence-test-assembly-first-layer-grid-pieces-choice-odd}
Let $D$ be a definite-side datum, with the notation of
Definition~\ref{asm-dependence-test-assembly-first-layer-grid-pieces-choice-normal}. Let $\Lambda=\delta\varphi_0(\mathfrak
c')$ be a CM lattice of the host, $t_0$ an unkilled two-torsion point of $\Lambda$, and $a,b\ge0$ integers with $ab$ even.
Then there is a nonzero odd vacuum function of size $(a,b)$ for $(\mathfrak c',\delta,t_0)$.
\end{lemma}
\leanref{Development/AbelianVarietyTheory/Proofs/FoundationProofs/NormSqMulGaussianLatticeTranslateEqShifted.lean}{GlobalTheory.Vacuum.exists_vacuum}
\begin{proof}
\openstep{The proof is not yet written out. It takes
$e=\kappa\,I_{\Lambda_v\to\Lambda_{\mathcal O_E^a,\mathcal O_E^b}}\big((\vartheta_{\Lambda,t_0})_v\big)$ for an orthonormal
basis $v$ that spans $Y\otimes\mathbf Z_{(2)}$ over $\mathcal O_K\otimes\mathbf Z_{(2)}$ and an algebraic number $\kappa$
all of whose conjugates have absolute value $1$, and checks conditions (a) to (c) of
Definition~\ref{asm-dependence-test-assembly-first-layer-grid-pieces-choice-normal}(1) and $e\ne0$.}
% Lean: Vacuum.exists_vacuum, 9,134 lines (tools/lean_size.py). Its steps: InducedType.induced_type,
% evenTheta_properties_curveLattice, T.standardPairLattice_locally_orthonormal_basis (b0 at p = 2), exists_kappa,
% T.intertwinerUnitary_all (theta space, norm and e != 0), T.thetaTaylorCoeffs_algebraic_at_complexMultiplication,
% normSq_pullB, and for (c) at a prime p: standardPairLattice_locally_orthonormal_basis (b), T.thetaPullbackConstant_sq_eq_det,
% T.admissible_characters_exist, eps = kappa / (c mu), T.localBasisLattice_index_coprime and
% T.theta_space_selfDual_pair_at_most_line.
\end{proof}

\begin{lemma}\label{asm-dependence-test-assembly-first-layer-grid-pieces-choice-jets}
Let $D$ be a definite-side datum, and let a theta grid for a level system for $D$, a grid Weil datum $W$ on it and standard
data on a set $U$ of tuples be given, with the notation of
Definition~\ref{asm-dependence-test-assembly-first-layer-grid-pieces-choice-normal}(2). Let $\delta\neq0$ with
$\delta^{12}=\Delta(\varphi_0(\mathfrak c))$, $t_0$ an unkilled two-torsion point of $\delta\varphi_0(\mathfrak c)$, $(S,T)$
a definite grid pair, $e$ an odd vacuum function of size $(\operatorname{rk}S,\operatorname{rk}T)$ for $(\mathfrak
c,\delta,t_0)$, and $g$ a Fock vacuum vector for the Cartan involutions of $W$ on $S$ and $T$ and the lattice of the pair
given by $L_S$ and $L_T$. Then there is $c\in\mathbf C^\times$ such that for every monomial $z^k$
\[
 (g\otimes_\phi z^k)(0)=c\,\epsilon_{S,T}\Big(\prod_{\tau,\beta,\alpha}\pi_\tau(\delta)^{k_{\tau,\beta,\alpha}}\cdot
 (\partial^ke)(0)\Big).
\]
\end{lemma}
\leanref{Development/AbelianVarietyTheory/Proofs/DivisorProofs/FockVacEqJets.lean}{GlobalTheory.FunObj.fock_vac_eq_jets}
\begin{proof}
\openstep{The proof, made separately for the pairs $(V_{i+1,m},V_{i,m})$ and $(Q_j,Q_{j+1})$, is not yet written out.}
% Lean: FunObj.fock_vac_eq_jets, 8,827 lines (tools/lean_size.py): the case chi i m is
% T.fockJets_origin_oppositeSignVacuum_conjPeriod (6,673 lines) and the case top t j is
% T.fockJet_vacuum_eq_const_mul_periodScaledJet (6,195 lines); they overlap.
\end{proof}

\begin{lemma}\label{asm-dependence-test-assembly-first-layer-grid-pieces-choice-normal-exists}
Let $D$ be a definite-side datum, let a level system for $D$ with base tuple $\mu_0$, a theta grid for it, a grid Weil
datum $W$ on the grid and standard data on a set $U$ of tuples be given
(Definitions~\ref{th-grid} and~\ref{asm-dependence-test-assembly-first-layer-grid-pieces-choice-standard}); the conditions
of the level system at $\mu_0$ are not used. Then there is a vacuum normalization for these standard data
(Definition~\ref{asm-dependence-test-assembly-first-layer-grid-pieces-choice-normal}(3)).
\end{lemma}
\leanref{Development/AbelianVarietyTheory/Proofs/DivisorProofs/FockVacEqJets.lean}{GlobalTheory.FunObj.normal_nonempty}
\leanref{Development/AbelianVarietyTheory/Proofs/DivisorProofs/FockVacEqJets.lean}{GlobalTheory.FunObj.exists_delta}
\leanref{Development/AbelianVarietyTheory/Proofs/DivisorProofs/FockVacEqJets.lean}{GlobalTheory.FunObj.even_rank_mul}
\leanref{Development/AbelianVarietyTheory/Proofs/DivisorProofs/ExistsUnitVacuumStandardLatticesDefinitePairStatement.lean}{GlobalTheory.T.exists_unitVacuum_standardLattices_definitePair}
\begin{proof}
\emph{The lattice.} Let $\delta$ be a twelfth root of $\Delta(\varphi_0(\mathfrak c))$. It is nonzero, as the
discriminant $g_2^3-27g_3^2$ of a lattice in $\mathbf C$ is nonzero. The ideal $\mathfrak c$ is nonzero, with norm prime to
$2$ and $d$ (Definition~\ref{hf-host}(5)), so $\Lambda=\delta\varphi_0(\mathfrak c)$ is a CM lattice of the host. It has
complex multiplication by $\mathcal O_K$, as $\mathfrak c$ is an ideal of $\mathcal O_K$, so it has an unkilled two-torsion
point $t_0$ by Lemma~\ref{is-vacuum-genus-one}(1).

\emph{The odd functions.} For a definite grid pair $(S,T)$ the product $\operatorname{rk}S\cdot\operatorname{rk}T$ is
even: the rank $r+1+m(r+5)+j(t+1)$ of $Q_j$ is even, as $r$ and $t$ are odd; and the ranks of $V_{i+1,m}$ and $V_{i,m}$
differ by the odd number $r+6$, so one of them is even. So
Lemma~\ref{asm-dependence-test-assembly-first-layer-grid-pieces-choice-odd} gives an odd vacuum function $e_{S,T}$ of size
$(\operatorname{rk}S,\operatorname{rk}T)$ for $(\mathfrak c,\delta,t_0)$.

\emph{The vacuum vectors.} For a definite grid pair $(S,T)$ the spaces $S$ and $T$ have correct signatures, and the
lattices $L_S$ and $L_T$ are self-dual, as they are spanned by orthonormal bases. So
Lemma~\ref{is-eigenvalue-conjugate-vacuum-exists}(2), for the standard character and the archimedean family of $W$ and the
Cartan involutions of $W$ on $S$ and $T$, gives a Fock vacuum vector $g_{S,T}$ for these Cartan involutions and the lattice
of the pair given by $L_S$ and $L_T$.

\emph{The constants.} Lemma~\ref{asm-dependence-test-assembly-first-layer-grid-pieces-choice-jets}, for $e_{S,T}$ and
$g_{S,T}$, gives $c_{S,T}\in\mathbf C^\times$ with condition (c).

At the pairs that are not definite take $e_{S,T}=0$, $g_{S,T}=0$ and $c_{S,T}=1$. The choices at a pair $(Q_j,Q_{j+1})$
are made for the two spaces, their bases $\phi_{Q_j}$ and $\phi_{Q_{j+1}}$ and the Cartan involutions of $W$ on them, none
of which depends on the tuple; so they do not depend on the tuple.
% Lean: FunObj.normal_nonempty (34 lines), followed: exists_delta (IsAlgClosed.exists_pow_nat_eq, jnt_ccK_curve,
% CurveLattice.discr_ne_zero), TwoTorsionPt.exists_point (its first part exists_lamOf, through
% isTwoTorsionCMPoint_of_odd_multiple, one of the declarations of Lemma is-vacuum-genus-one), Vacuum.exists_vacuum with
% even_rank_mul (rankX even and sLen odd), T.exists_unitVacuum_standardLattices_definitePair (vac_exists_of_defPair:
% VacuumO.vacExists with DefNode.sigOK and N.selfDual), fock_vac_eq_jets for c0, and rfl for the field free. The hypotheses
% mu0 and hL are not used.
\end{proof}

\begin{lemma}\label{asm-dependence-test-assembly-first-layer-grid-pieces-choice-adelic}
Let $E$ be the host field of a definite-side datum $D$, with the notation of
Definition~\ref{asm-dependence-test-assembly-first-layer-grid-pieces-choice-grid}.
\begin{enumerate}
\item Every finitely generated $\mathcal O_E$-module $\Lambda\subset E^n$ that spans $E^n$ has an adelic basis.
\item Let $I$ be a nonzero fractional ideal of $E$, $P$ a finite set of finite places of $E$, and $u_w\in E_w$ with
  $v_w(u_w)=v_w(I)$ for $w\in P$. Then there is $y\in\mathbf A_{E,f}^\times$ with $v_w(y_w)=v_w(I)$ at every finite place
  $w$ and $y_w=u_w$ for $w\in P$.
\item Let $\iota\colon A\to B$ be an isometric embedding of Hermitian spaces of $D$, and $a_1,\dots,a_c\in B$ orthogonal
  to $\iota(A)$, with $\operatorname{rk}A+c=\operatorname{rk}B$ and $\langle a_k,a_l\rangle=1$ if $k+l=c+1$ and $k\neq
  l$, $\langle a_k,a_k\rangle=\epsilon$ if $2k=c+1$, and $\langle a_k,a_l\rangle=0$ otherwise, for some
  $\epsilon\in E^\times$. Let $\Lambda$ be an $\mathcal O_E$-lattice of $A$ with an adelic basis $\beta$, let
  $\mathfrak b_1,\dots,\mathfrak b_c$ be nonzero fractional ideals and $y_1,\dots,y_c\in\mathbf A_{E,f}^\times$ with
  $v_w(y_{k,w})=v_w(\mathfrak b_k)$ at every finite place $w$. Then the lattice
  $\iota(\Lambda)+\sum_k\mathfrak b_ka_k$ has an adelic basis $\beta'$ that is the sum of $\beta$ and $(y_{k,w}a_k)_k$ at
  every finite place $w$.
\end{enumerate}
\end{lemma}
\leanref{Development/AbelianVarietyTheory/Proofs/IwasawaDecomposition/ExistsSelfDualAdelicFrameLevelOne.lean}{GlobalTheory.AdelicFrame.exists_isFrame}
\leanref{Development/AbelianVarietyTheory/Proofs/DivisorProofs/AuxiliaryVectorDecompositionEqZero.lean}{GlobalTheory.FunObj.exists_gen}
\leanref{Development/AbelianVarietyTheory/Proofs/DivisorProofs/AuxiliaryVectorDecompositionEqZero.lean}{GlobalTheory.FunObj.exists_sumFrame}
\begin{proof}
(1) At each finite place $w$ the closure $\Lambda_w$ of $\Lambda$ is a finitely generated torsion-free module over the
discrete valuation ring $\mathcal O_{E_w}$ that spans $E_w^n$, so $\Lambda_w=\beta_w\mathcal O_{E_w}^n$ with $\beta_w\in
GL_n(E_w)$. As $\Lambda$ and $\mathcal O_E^n$ each contain a nonzero integer multiple of the other, $\Lambda_w=\mathcal
O_{E_w}^n$ at almost every $w$, where $\beta_w=1$ can be taken. Then $\beta=(\beta_w)_w$ is an adelic basis.

(2) The ideal $I$ is generated at every finite place $w$ by an element of $E_w$ of valuation $v_w(I)$, which is $0$ at
almost every $w$; so there is $t\in\mathbf A_{E,f}^\times$ with $v_w(t_w)=v_w(I)$ for every $w$. Replace its components at
the places of $P$ by the $u_w$, which are nonzero.

(3) Let $C$ be the matrix whose columns are the images under $\iota$ of the basis vectors of $A$, followed by
$a_1,\dots,a_c$. It is invertible: if $\iota(x)+\sum_kc_ka_k=0$, pairing with each $a_l$ gives $\sum_kc_k\langle
a_l,a_k\rangle=0$, and the Gram matrix of the $a_k$ is invertible as $\epsilon\neq0$, so every $c_k=0$; then $\iota(x)=0$
and $x=0$. Put $\beta'=C\cdot\operatorname{diag}(\beta,y_1,\dots,y_c)$. Then
$\beta'\hat{\mathcal O}_E^{\operatorname{rk}B}=\iota(\beta\hat{\mathcal O}_E^{\operatorname{rk}A})+\sum_ky_k\hat{\mathcal
O}_Ea_k=\iota(\hat\Lambda)+\sum_k\hat{\mathfrak b}_ka_k$, the closure of $\iota(\Lambda)+\sum_k\mathfrak b_ka_k$, and the
component of $\beta'$ at $w$ has the columns $\iota(\beta_w)$ followed by the $y_{k,w}a_k$.
% Lean: (1) AdelicFrame.exists_isFrame (exists_basis_hat), (2) FunObj.exists_gen (Idele.exists_idealOf_eq,
% exists_unit_cpt), (3) FunObj.exists_sumFrame (C = fromCols (frameMat B A d.iota) d.a, injective by
% auxiliaryVectorDecomposition_eq_zero, reordered by frame d.rank). The Lean proofs take 979, 716 and 1,638 lines, mostly
% on the adele ring.
\end{proof}

\begin{lemma}\label{asm-dependence-test-assembly-first-layer-grid-pieces-choice-rows}
Let $D$ be a definite-side datum with host field $E$, let a level system for $D$ and a theta grid for it with modulus
$\mathfrak m$ be given (Definition~\ref{th-grid}), and let $\mathfrak A$ be the ideal of the host, with the notation of
Definition~\ref{asm-dependence-test-assembly-first-layer-grid-pieces-choice-grid}. Put $\mathfrak a=\mathfrak
A\bar{\mathfrak m}^{-1}$, $\mathfrak a'=\bar{\mathfrak A}\mathfrak m^{-1}$, and for $0\le j<m$ and $1\le k\le m-1-j$ let
$\mathfrak b_{0,k}=\mathfrak a^{m-k}\mathfrak a'^{-k}$ and $\mathfrak b_{j,k}=\mathfrak A^{m-j-k}\bar{\mathfrak A}^{-k}$
for $j\ge1$, the ideals $I_k$ of the lattices of Definition~\ref{th-grid} for the embedding $Q_j\subset V_{m,j}$. Let $w$
be a local place at a tuple $t$ with $v_w(\mathfrak A)=v_w(\bar{\mathfrak A})=0$. Then the
$\mathfrak b_{j,k}$ are nonzero, $v_w(\mathfrak b_{0,k})=kd_w$, and $v_w(\mathfrak b_{j,k})=0$ for $j\ge1$.
\end{lemma}
\leanref{Development/AbelianVarietyTheory/Proofs/DivisorProofs/AuxiliaryVectorDecompositionEqZero.lean}{GlobalTheory.FunObj.rowId}
\leanref{Development/AbelianVarietyTheory/Proofs/DivisorProofs/AuxiliaryVectorDecompositionEqZero.lean}{GlobalTheory.FunObj.count_rowId}
\leanref{Development/AbelianVarietyTheory/Proofs/DivisorProofs/AuxiliaryVectorDecompositionEqZero.lean}{GlobalTheory.FunObj.rowId_ne_zero}
\leanref{Development/AbelianVarietyTheory/Proofs/DivisorProofs/AuxiliaryVectorDecompositionEqZero.lean}{GlobalTheory.FunObj.count_modul}
\leanref{Development/AbelianVarietyTheory/Proofs/DivisorProofs/AuxiliaryVectorDecompositionEqZero.lean}{GlobalTheory.FunObj.count_modulBar}
\begin{proof}
The ideals $\mathfrak A$ and $\mathfrak m$ are nonzero, so the $\mathfrak b_{j,k}$ are. As
$\mathfrak m=\prod_{v\in S_0}w(v)^{d(v)}$ and the $w(v)$ are distinct, lying over distinct places $v$, $v_w(\mathfrak m)=d_w$
for every finite place $w$. Next $v_w(\bar{\mathfrak m})=v_{\bar w}(\mathfrak m)=0$ for the local place $w$: otherwise
$\bar w=w(v)$ with $v\in S_0$, and $v$ is the place of $F$ below $w$. A local place over a place $v\in S_0$ is $w(v)$:
this holds for the places $w(v)$, and a place of $Q_t$ over $v$ is $w(v)$ by conditions (4) and (5) of
Definition~\ref{th-grid-places}. So $w=w(v)=\bar w$, which condition (1) there excludes. Hence
$v_w(\mathfrak b_{0,k})=(m-k)\big(v_w(\mathfrak A)-v_w(\bar{\mathfrak m})\big)-k\big(v_w(\bar{\mathfrak A})-v_w(\mathfrak
m)\big)=kd_w$, and $v_w(\mathfrak b_{j,k})=(m-j-k)v_w(\mathfrak A)-kv_w(\bar{\mathfrak A})=0$ for $j\ge1$.
% Lean: rowId G j p = stairId (stA, stAbar, m - 1) p for j = 0 and stairId (AA, AAbar, m - 1 - j) p otherwise, with p = k - 1;
% count_rowId (conj_AA_and_stA for the nonvanishing, count_stairId, count_modul, count_modulBar with Good.locPl_cases and
% conjPrime_ne_of_isIdempotentAt), rowId_ne_zero.
\end{proof}

\begin{lemma}\label{asm-dependence-test-assembly-first-layer-grid-pieces-choice-frames}
Let $D$ be a definite-side datum with host field $E$, let a level system for $D$, a theta grid for it and a grid Weil datum
$W$ on the grid be given (Definition~\ref{th-grid}), let $U$ be a set of tuples, and let standard data on $U$ and a vacuum
normalization for them be given
(Definitions~\ref{asm-dependence-test-assembly-first-layer-grid-pieces-choice-standard}
and~\ref{asm-dependence-test-assembly-first-layer-grid-pieces-choice-normal}). Suppose that the local places at the tuples
of $U$ form a finite set and that $v_w(\mathfrak A)=v_w(\bar{\mathfrak A})=0$ at each of them, for the ideal $\mathfrak A$
of the host. Then there is a grid choice for $W$ on $U$
(Definition~\ref{asm-dependence-test-assembly-first-layer-grid-pieces-choice-grid}).
\end{lemma}
\leanref{Development/AbelianVarietyTheory/Proofs/DivisorProofs/AuxiliaryVectorDecompositionEqZero.lean}{GlobalTheory.FunObj.gridChoice_nonempty}
\leanref{Development/AbelianVarietyTheory/Proofs/DivisorProofs/AuxiliaryVectorDecompositionEqZero.lean}{GlobalTheory.FunObj.rowLat_eq}
\leanref{Development/AbelianVarietyTheory/Proofs/DivisorProofs/AuxiliaryVectorDecompositionEqZero.lean}{GlobalTheory.FunObj.valued_frameSign}
\begin{proof}
Let $P_U$ be the finite set of the local places at the tuples of $U$, and write
Lemma~\ref{asm-dependence-test-assembly-first-layer-grid-pieces-choice-adelic} as ``the adelic lemma''. The lattices of the
grid are finitely generated and span their spaces, as they are self-dual.

\emph{The top bases.} By part (1) of the adelic lemma, $\Lambda_{Q_j}$ has an adelic basis $b_j$ for $0\le j\le m$.

\emph{The last row.} Let $0\le j<m$, and let $a_1,\dots,a_{m-1-j}$ be the basis of the complement of $Q_j$ in $V_{m,j}$
given with the grid, with the $\epsilon_k$ and $e_k$ of
Definition~\ref{asm-dependence-test-assembly-first-layer-grid-pieces-choice-grid}(3)(a). The lattice of $V_{m,j}$ is
$\iota(\Lambda_{Q_j})+\sum_k\mathfrak b_{j,k}a_k$ for the ideals $\mathfrak b_{j,k}$ of
Lemma~\ref{asm-dependence-test-assembly-first-layer-grid-pieces-choice-rows}, which are the ideals $I_k$ of the lattices
of Definition~\ref{th-grid} for the embedding $Q_j\subset V_{m,j}$, and the middle complement vector, if there is one,
has $\langle a_k,a_k\rangle=-1$. For
each $k$ and each $w\in P_U$, $\epsilon_k\varpi_w^{e_k}$ has valuation $e_k=v_w(\mathfrak b_{j,k})$, by that lemma. So
part (2) of the adelic lemma gives $y_k\in\mathbf A_{E,f}^\times$ with $v_w(y_{k,w})=v_w(\mathfrak b_{j,k})$ at every $w$
and $y_{k,w}=\epsilon_k\varpi_w^{e_k}$ for $w\in P_U$. Part (3), with $\epsilon=-1$, gives an adelic basis
$\beta_{m,j}$ of $\Lambda_{V_{m,j}}$ that is the sum of $b_j$ and $(\epsilon_k\varpi_w^{e_k}a_k)_k$ at every $w\in P_U$.

\emph{The other spaces.} For $1\le i\le m$ and $0\le j\le m$ with $i<m$ or $j=m$, part (1) gives an adelic basis
$\beta_{i,j}$ of $\Lambda_{V_{i,j}}$.

\emph{The last top space.} The lattice of $Q_m$ is $\iota(\Lambda_{V_{m,m}})+\mathcal O_Ea$ for the vector $a$, with
$\langle a,a\rangle=1$, spanning the complement of $V_{m,m}$, by the embeddings and lattices of
Definition~\ref{th-grid}. Part (3), with
$c=1$, $\epsilon=1$, $\mathfrak b_1=\mathcal O_E$ and $y_1=1$, gives an adelic basis $\beta^+_m$ of $\Lambda_{Q_m}$ that is
the sum of $\beta_{m,m}$ and $a$ at every place. Put $\beta^+_j=b_j$ for $j<m$.

\emph{The places.} The ideal $\mathfrak A$ is nonzero, so the set of finite places $w$ with $v_w(\mathfrak A)\neq0$ is
finite. For each place of $F$ below one of them choose one of them above it, and let $P_{\mathfrak A}$ be the set of the
chosen places. Every place of $P_{\mathfrak A}$ has $v_w(\mathfrak A)\neq0$, and for a place $w$ with $v_w(\mathfrak
A)\neq0$ the place chosen over the place of $F$ below $w$ is the only place of $P_{\mathfrak A}$ over it.

These bases, the given standard data and vacuum normalization, and $P_{\mathfrak A}$ form a grid choice for $W$ on $U$.
% Lean: FunObj.gridChoice_nonempty (117 lines), followed: AdelicFrame.exists_isFrame for b0 j and for F i j off the last row;
% for the last row exists_gen with rowId_ne_zero, valued_frameSign and count_rowId, then exists_sumFrame with eps = -1 and
% rowLat_eq (G.sp.colLat for j = 0, G.sp.rowLat otherwise); for the top exists_sumFrame with c = 1, y = 1 and G.sp.topLat;
% for PA Idele.exists_idealOf_eq with StairConj.conj_AA_and_stA and Idele.ord_vw_eventually_zero, and the choice rep.
\end{proof}

\begin{lemma}\label{asm-dependence-test-assembly-first-layer-grid-pieces-choice-different}
Let $D$ be a definite-side datum with integers $Q_S$ and $q$ and ideal $\mathfrak b$ of its host
(Definition~\ref{hf-host}), and let $p$ be a prime. Then $p$ divides $\mathrm N\mathfrak b$ if and only if $p$ divides
$Q_Sq$.
\end{lemma}
\leanref{Development/AutomorphicArithmetic/CyclotomicBounds/RamifiedDiscriminants/HalfNorm.lean}{Formalism.BasicNotions.Host.prime_dvd_absNorm_half_iff}
\begin{proof}
Two facts on the host field are used: (i) a prime $\mathfrak P$ of $E$ divides $\mathfrak b$ or $\bar{\mathfrak b}$ if and
only if it lies over a prime $\ell\neq d$ that divides the discriminant of $E$; (ii) a prime $\ell\neq d$ divides the
discriminant of $E$ if and only if it divides $Q_Sq$. If $p\mid\mathrm N\mathfrak b$, a prime $\mathfrak P$ of $E$ over $p$
divides $\mathfrak b$, so $p\mid Q_Sq$ by (i) and (ii). Conversely, if $p\mid Q_Sq$, then $p\ne d$, and a prime $\mathfrak
P$ over $p$ divides $\mathfrak b$ or $\bar{\mathfrak b}$ by (i) and (ii); in both cases $p$ divides $\mathrm N\mathfrak
b=\mathrm N\bar{\mathfrak b}$.
\openstep{Facts (i) and (ii) are not yet written out; (i) rests on
$\mathfrak D_{E/K}=\mathfrak b\bar{\mathfrak b}$ and $\mathfrak b+\bar{\mathfrak b}=\mathcal O_E$
(Proposition~\ref{hf-host-field-properties}(f)), and (ii) on the primes dividing the discriminant of $E$.}
% Lean: Host.prime_dvd_absNorm_half_iff, 1,616 lines (tools/lean_size.py): (i) is UpperTheory.HF.half_supp (1,581 lines,
% through diff_EK_supp), (ii) is Host.mem_ramI_iff with Host.dvd_natAbs_discr_iff (p divides disc E iff p divides
% Q_S q d) and d_not_dvd_QSq.
\end{proof}

\begin{lemma}\label{asm-dependence-test-assembly-first-layer-grid-pieces-choice-host-ideal}
Let $D$ be a definite-side datum with host field $E$ and integers $Q_S$ and $q$ and ideal $\mathfrak c$ of its host, and let
$\mathfrak A$ be the ideal of the host. Let $w$ be a finite place of $E$ that does not contain $Q_Sq\,\mathrm N\mathfrak
c$. Then $v_w(\mathfrak A)=v_w(\bar{\mathfrak A})=0$.
\end{lemma}
\leanref{Development/AbelianVarietyTheory/Proofs/DivisorProofs/AuxiliaryVectorDecompositionEqZero.lean}{GlobalTheory.FunObj.count_AA_AAbar}
\leanref{Development/AbelianVarietyTheory/Proofs/FoundationProofs/PrimeDvdAbsNormHalfDifferent.lean}{GlobalTheory.HalfCompare.primeToAA_of_not_mem}
\begin{proof}
Let $\mathfrak P$ be a finite place of $E$ not containing $M=Q_Sq\,\mathrm N\mathfrak c$. It does not divide $\mathfrak
c\mathcal O_E$, which contains $\mathrm N\mathfrak c$. It does not divide $\mathfrak b$: otherwise the prime $p$ below
$\mathfrak P$ divides $\mathrm N\mathfrak b$, hence $Q_Sq$ by
Lemma~\ref{asm-dependence-test-assembly-first-layer-grid-pieces-choice-different}, and $M\in\mathfrak P$. So
$v_{\mathfrak P}(\mathfrak A)=v_{\mathfrak P}(\mathfrak c\mathcal O_E)-v_{\mathfrak P}(\mathfrak b)=0$. This applies to
$w$ and to $\bar w$, which does not contain the rational integer $M$ either. Hence $v_w(\mathfrak A)=0$ and
$v_w(\bar{\mathfrak A})=v_{\bar w}(\mathfrak A)=0$.
% Lean: FunObj.count_AA_AAbar (25 lines): HalfCompare.primeToAA_of_not_mem at w and at cP w (natCast_mem_cP), with
% primeToAA_iff (absNorm_ccK_mem, two_q_mem through prime_dvd_absNorm_halfDifferent), ccE_ne_bot, bb_ne_bot and
% count_conj_smul for AAbar.
\end{proof}

\begin{lemma}\label{asm-dependence-test-assembly-first-layer-grid-pieces-choice}
Let $D$ be a definite-side datum with host field $E$, integers $Q_S$ and $q$ and the ideal $\mathfrak c$ of its host,
let a level system $(\xi_t,Q_t)$ for $D$ with base tuple $\mu_0$, a theta grid for it with places $w(v)$, $v\in S_0$,
and a grid Weil datum $W$ on the grid be given (Definition~\ref{th-grid}), and let $U$ be a finite set of tuples. The
local places of the grid at a tuple $t$ are the places $w(v)$, $v\in S_0$, and the places of $Q_t$
(Definition~\ref{asm-dependence-test-assembly-first-layer-grid-pieces-choice-grid}(1)). If no local place at a tuple of
$U$ contains $Q_Sq\,\mathrm N(\mathfrak c)$, then there is a grid choice for $W$ on $U$
(Definition~\ref{asm-dependence-test-assembly-first-layer-grid-pieces-choice-grid}(3)).
\end{lemma}
\leanref{Development/AbelianVarietyTheory/Proofs/DivisorProofs/AuxiliaryVectorDecompositionEqZero.lean}{GlobalTheory.FunObj.gridChoice_nonempty}
\leanref{Development/AbelianVarietyTheory/Proofs/DivisorProofs/DefNodeIsV.lean}{GlobalTheory.FunObj.exists_std}
\leanref{Development/AbelianVarietyTheory/Proofs/DivisorProofs/FockVacEqJets.lean}{GlobalTheory.FunObj.normal_nonempty}
\leanref{Development/AbelianVarietyTheory/Proofs/DivisorProofs/AuxiliaryVectorDecompositionEqZero.lean}{GlobalTheory.FunObj.count_AA_AAbar}
\begin{proof}
By Lemma~\ref{asm-dependence-test-assembly-first-layer-grid-pieces-choice-standard-exists} there are standard data on $U$,
and by Lemma~\ref{asm-dependence-test-assembly-first-layer-grid-pieces-choice-normal-exists} a vacuum normalization for
them. The local places at the tuples of $U$ form a finite set, as $U$, $S_0$ and the sets $Q_t$ are finite. None of them
contains $Q_Sq\,\mathrm N(\mathfrak c)$, so $v_w(\mathfrak A)=v_w(\bar{\mathfrak A})=0$ at each of them, for the ideal
$\mathfrak A$ of the host, by Lemma~\ref{asm-dependence-test-assembly-first-layer-grid-pieces-choice-host-ideal}. So
Lemma~\ref{asm-dependence-test-assembly-first-layer-grid-pieces-choice-frames}, for these standard data and this vacuum
normalization, gives a grid choice for $W$ on $U$.
% Lean: the steps of FunObj.gridChoice_nonempty_of after Good.badInt_notMem: exists_std G mu0 hL U hU (Lemma
% -choice-standard-exists), normal_nonempty W mu0 hL U N (Lemma -choice-normal-exists), and gridChoice_nonempty W N Z
% (Lemma -choice-frames) with hU.biUnion for the finiteness and count_AA_AAbar (Lemma -choice-host-ideal) at each local
% place. What remains open is in Lemmas -choice-orthonormal-pieces (6,845 lines in three parts), -choice-odd (9,134),
% -choice-jets (8,827) and -choice-different (1,616), sizes from tools/lean_size.py; they overlap.
\end{proof}

\begin{lemma}\label{asm-dependence-test-assembly-first-layer-grid-pieces}
Let $D$ be a definite-side datum with host field $E$, auxiliary prime $q_0$ and level primes
$g_0,g_1,p_{de},p'_{de}$, and let a level system $(\xi_t,Q_t)$ for $D$ with base tuple $\mu_0$ be given
(Definition~\ref{th-grid}), with the set $S_1$ and the places of a theta grid of
Definition~\ref{asm-dependence-test-assembly-first-layer-grid-pieces-places}.
\begin{enumerate}
\item Say that a prime $p$ \emph{splits in} $K$ if $p$ is prime to $N=Q_Sqd$ and $-d$ is a square
  modulo $p$. If $g_0$, $g_1$, $p_{de}$ and $p'_{de}$ split in $K$, $q_0\ne g_0,g_1$, and no place of any
  $Q_t$ lies above $q_0$, then the level system has places of a theta grid whose base place lies
  outside $S_1$.
\item For all constants $c$, the level system has a theta grid with constants $c$ if and only if it
  has places of a theta grid whose base place lies outside $S_1$.
\item Assume the results of Theorem~\ref{lit-bundle-theta}, among them the Weil representation of
  unitary dual pairs. Let $R$ be a realization of $D$ with this level system that satisfies the chain and model conditions of its
  core. Then every theta grid for the level system of $R$ carries a Weil datum over $R$.
\item Let $W$ be a grid Weil datum on a theta grid for the level system, and $U$ a finite set of
  tuples. If $p_{de}$ and $p'_{de}$ are prime to $Q_Sq$ and no place of $Q_t$, $t\in U$, contains
  $\mathrm N(\mathfrak c)$, then there is a grid choice for $W$ on $U$.
\end{enumerate}
\end{lemma}
\leanref{Development/AbelianVarietyTheory/Proofs/TwistedOperators/GridExistenceOfFunctoriality.lean}{GlobalTheory.RowEx.GridExPf.exists_places_of_cls}
\leanref{Development/AbelianVarietyTheory/Definitions/Divisors/NonemptyGridDataIff.lean}{GlobalTheory.ChiGridModel.nonempty_gridData_iff}
\leanref{Development/AbelianVarietyTheory/Proofs/DivisorProofs/ExistsWeilOverOfClauses.lean}{GlobalTheory.ChiGrid.WeilOver.exists_weil_over_of_clauses}
\leanref{Development/AbelianVarietyTheory/Proofs/DivisorProofs/AuxiliaryVectorDecompositionEqZero.lean}{GlobalTheory.FunObj.gridChoice_nonempty_of}
\begin{proof}
(1) A prime that splits in $K$ is different from $2$ and $d$, as $2$ divides $Q_S$. Four facts are used.
\begin{enumerate}
\item[(a)] If $p$ splits in $K$ and a place $P$ of $E$ contains $p$, then $P$ and $\bar P$ lie over
  different primes of $K$. Indeed, choose $s\in\mathbf Z$ with $s^2\equiv-d\pmod p$. Then
  $(s+\sqrt{-d})(s-\sqrt{-d})=s^2+d\in P$, so $P$ contains $s+\sqrt{-d}$, say. If $P\cap K=\bar P\cap K$,
  then $P$ also contains the conjugate $s-\sqrt{-d}$, hence $2\sqrt{-d}$ and $4d$, so $p\in\{2,d\}$.
\item[(b)] If $v$ is a place of $F$ and $\mathfrak p$ a prime of $K$, both containing the prime $p$,
  there is a place of $E$ over both. Indeed, let $w$ lie over $v$ and $u$ over $\mathfrak p$; as $E$ is
  Galois over $\mathbf Q$, $w=\sigma u$ for an automorphism $\sigma$. If $\sigma(\sqrt{-d})=\sqrt{-d}$,
  then $w\cap K=\sigma(u\cap K)=\mathfrak p$. Otherwise $\bar w$ lies over $v$, as complex conjugation
  fixes $F$, and $\bar w=c\sigma u$ with $c\sigma$ fixing $\sqrt{-d}$, so $\bar w$ lies over $\mathfrak p$.
\item[(c)] Two places of $E$ over the same place of $F$ are equal or conjugate, as $E/F$ is Galois of
  degree $2$; and a place contains at most one rational prime.
\item[(d)] The prime $q_0$ splits in $K$: $q_0\equiv1$ modulo $N$, so $q_0$ is prime to $N$, and as
  $q_0\equiv1$ modulo $4$ and modulo $d$, quadratic reciprocity gives
  $\bigl(\frac{-d}{q_0}\bigr)=\bigl(\frac{q_0}d\bigr)=1$. Also $q_0\ne p_{de},p'_{de}$, since the
  prime $r_\ast$ of Definition~\ref{th-grid-places} divides $q_0-1$ but not $(p_{de}-1)(p'_{de}-1)$.
\end{enumerate}
For each prime $p$ choose a prime $\mathfrak p_p$ of $K$ over $p$: for $p=g_k$ the prime chosen above
$g_k$ by the level system (condition (L1) of Definition~\ref{th-grid}), so that $P_g$ consists of the places of
$E$ over $\mathfrak p_{g_0}$ and $\mathfrak p_{g_1}$; for $p\in\{p_{de},p'_{de}\}$, if some place of a set $Q_t$
contains $p$, the prime of $K$ under all such places, which the level system provides by (L5); and any prime over
$p$ otherwise. Let $v_0$
be a place of $F$ containing $q_0$. For a place $v$ of $F$ containing the prime $p$, let $w(v)$ be a place of
$E$ over $v$ and over $\mathfrak p_p$, by (b).

Condition (i): every $v\in\{v_0\}\cup S_1$ contains one of $q_0$, $g_0$, $g_1$, $p_{de}$, $p'_{de}$, which
split in $K$ by (d) and the hypothesis, so $w(v)\ne\overline{w(v)}$ by (a). Condition (ii): if $v$
contains $g_k$, then $w(v)$ lies over $\mathfrak p_{g_k}$, so it lies in $P_g$. Condition (iii) holds with
$\mathfrak p=\mathfrak p_p$. Condition (iv): let $P\in Q_t$ lie over $v\in S_1$. The level system gives a
prime $q\ne g_0,g_1$ in $P$, the order of its residue field (L5), and $q$ lies in $v$, so $q\in\{p_{de},p'_{de}\}$ by (c), and $P$ lies over
$\mathfrak p_q$. By (c), $P=w(v)$ or $P=\overline{w(v)}$; in the second case $P$ and $w(v)$ would both lie
over $\mathfrak p_q$, against (a). Condition (v): a place of $Q_t$ over $v_0$ would contain $q_0$, which
the hypothesis excludes. Finally $v_0\notin S_1$: by (c), $v_0$ contains no prime other than $q_0$, and
$q_0$ is none of $g_0$, $g_1$ (hypothesis), $p_{de}$, $p'_{de}$ (by (d)).

(2) A theta grid with constants $c$ comes with places of a theta grid, and their base place lies outside
$S_1$. Indeed, the place $w(v_0)$ of $E$ lies over the base place $v_0$. By condition (d) of
Definition~\ref{th-grid}, the character $\chi_1$ is trivial on the units of $E_{w(v_0)}$ congruent to $1$
modulo $w(v_0)$. If $v_0$ were in $S_1$, condition (e) would make $\chi_1$ nontrivial on these units. Conversely, let places of a theta grid with
base place outside $S_1$ be given. A theta grid on them consists of these places, characters
$\chi_1,\dots,\chi_m$ with the conditions of a theta grid at its places, and the spaces of the grid with their
embeddings. The characters exist by
Lemma~\ref{asm-dependence-test-assembly-first-layer-grid-pieces-characters}, for the constant $c$ of the triple of
constants. The spaces are built with
Lemma~\ref{asm-dependence-test-assembly-first-layer-grid-pieces-spaces}. Let $\mathfrak m$ be the modulus of the
grid, $\mathfrak A$ the ideal of the host, and for nonzero fractional ideals $\mathfrak a,\mathfrak a'$ with
$\mathfrak a'=\bar{\mathfrak a}$ and $n\ge0$ call the ideals $\mathfrak a^{n-k}\mathfrak a'^{-(k+1)}$,
$0\le k<n$, the \emph{staircase ideals} of $(\mathfrak a,n)$; they satisfy the condition of part (2) of that
lemma, and they are the ideals $I_{k+1}$ of the lattices of a theta grid (Definition~\ref{th-grid}). Write $\operatorname{ind}$ for the number of negative entries of a signature.
\begin{enumerate}
\item[(a)] $P_0$ is the sum of $Y_0$, of rank $r$ and index $m/2$, with a line ($c=1$, $\epsilon=1$,
  $\mathfrak b_1=\mathcal O_E$): rank $r+1$ and index $m/2=\lfloor(m+1)/2\rfloor$.
\item[(b)] For $1\le i\le m$, $P_i$ is a space of rank $\operatorname{rk}P_i$ and index
  $\lfloor(m-i+1)/2\rfloor$ (part (1)), and $V_{i,0}$ is the sum of $P_i$ with $c=i-1$, $\epsilon=-1$ and the
  staircase ideals of $(\mathfrak A\bar{\mathfrak m}^{-1},i-1)$: rank $\operatorname{rk}P_i+i-1=r+i(r+6)$ and index
  $\lfloor(m-i+1)/2\rfloor+\lfloor i/2\rfloor=m/2$, as $m$ is even.
\item[(c)] For $1\le j<m$, $Q_j$ is a space of rank $\operatorname{rk}Q_j$ and signature $(\operatorname{rk}Q_j,0)$,
  and $V_{m,j}$ is the sum of $Q_j$ with $c=m-1-j$, $\epsilon=-1$ and the staircase ideals of $(\mathfrak A,m-1-j)$:
  rank $\operatorname{rk}Q_j+m-1-j=\operatorname{rk}V_{m,j}$ and index $\lceil(m-1-j)/2\rceil=\lfloor(m-j)/2\rfloor$.
\item[(d)] $V_{m,m}$ is a space of rank $\operatorname{rk}V_{m,m}$ and signature $(\operatorname{rk}V_{m,m},0)$, and
  $Q_m$ its sum with a line: rank $\operatorname{rk}V_{m,m}+1=\operatorname{rk}Q_m$ and signature
  $(\operatorname{rk}Q_m,0)$.
\item[(e)] The other spaces are spaces of the ranks and signatures of Definition~\ref{th-grid}, by part (1).
\end{enumerate}
These are the spaces and embeddings of a theta grid, with complements of the ranks and Gram matrices of
Definition~\ref{th-grid} and the lattices $\Lambda+\sum_kI_ka_k$ that it asks for.

(3) Let $G$ be a theta grid for the level system of $R$, with places $w(v)$, $v\in S_0$. By the chain condition of
the core, the standard character $\psi_R$ of $R$ is standard, and for every tuple $t$ and $0\le k<m$ the Weil datum
$\omega_{R,t,k}$ of $R$ is a Weil datum for $(Y_k,Y_{k+1})$, or for $(Y_{m-1},X')$ if $k=m-1$, with standard character
$\psi_R$ and the characters that the level system gives to this step at $t$; these are the characters of the grid
pair $(V_{0,k},V_{0,k+1})$ at $t$, the pair (G2) of Definition~\ref{th-grid} with $i=0$. By the model condition, the models of discrete series of $R$ on the $Y_j$ and on
$X'$ carry Cartan involutions with orthonormal bases at the archimedean places, and the archimedean family $\theta_R$
of $R$ has at each $\tau\in\mathcal T$ a member that extends $\tau$. By
Lemma~\ref{asm-dependence-test-assembly-first-layer-grid-pieces-weil-exists}, $G$ has a grid Weil datum $W$ with
standard character $\psi_R$; let $e$ be the idempotent of its coset vectors, a sum of the idempotents of
$\mathbf A_E$ at the places $w(v)$ with $e\bar e=0$. Take on $V_{0,j}=Y_j$ ($j<m$) and on $V_{0,m}=X'$ the Cartan
involutions and bases of the models of $R$, and on the other grid spaces those of $W$; take at the pairs
$(V_{0,k},V_{0,k+1})$ at $t$ the Weil data $\omega_{R,t,k}$, and at the other pairs those of $W$.
Lemma~\ref{asm-dependence-test-assembly-first-layer-grid-pieces-weil-assemble}, for $\psi_R$, $\theta_R$ and $e$,
gives a grid Weil datum $W'$ with these data. Its standard character, archimedean family, Cartan involutions and bases
on $V_{0,j}$ and Weil data at the pairs $(V_{0,k},V_{0,k+1})$ are those of $R$, so $W'$ is a Weil datum over $R$
(Definition~\ref{asm-grid-evaluation}(2)).

(4) The local places of the grid at the tuples of $U$
(Definition~\ref{asm-dependence-test-assembly-first-layer-grid-pieces-choice-grid}(1)) form a finite set, as $U$ and the sets $Q_t$ are finite. None of them contains $Q_Sq\,\mathrm N(\mathfrak c)$.
\begin{itemize}
\item A place over $v_0$ contains $q_0$, which is prime to $Q_Sq$, as $q_0\equiv1$ modulo $Q_Sq$, and does not divide
  $\mathrm N(\mathfrak c)$ (Definition~\ref{th-grid-places}).
\item A place over a place of $S_1$ contains one of $g_0$, $g_1$, $p_{de}$, $p'_{de}$. The primes $g_0$ and $g_1$ are
  congruent to $1$ modulo $N=Q_Sqd$, a condition of the datum, and, as level primes, prime to $\mathrm N(\mathfrak c)$
  (Definition~\ref{hf-host}(5)); $p_{de}$ and $p'_{de}$ are prime to $Q_Sq$ by the hypothesis and to
  $\mathrm N(\mathfrak c)$ as level primes.
\item A place $P$ of $Q_t$, $t\in U$, does not contain $\mathrm N(\mathfrak c)$, by the hypothesis, nor $Q_Sq$, as it
  contains none of $Q_S$, $q$, $d$ and the prime $T$ of the system of spaces, by condition (L5) of
  Definition~\ref{th-grid}.
  As $P$ is prime, it does not contain the product.
\end{itemize}
So Lemma~\ref{asm-dependence-test-assembly-first-layer-grid-pieces-choice} gives a grid choice for $W$ on $U$.
% Lean: (3) ChiGrid.WeilOver.exists_weil_over_of_clauses (about 55 lines), followed: hch.omegaX mu0 .std for psi_R,
% T.exists_weilRepresentation_vacuum_cosetVectors_grid (Lemma -weil-exists) for W, W.cosVec_is for e, exists_maj_heq with
% model_spec for the Cartan involutions and bases of R.MY j and R.MX (s, B), exists_om_heq with hch.omegaY,
% hch.omegaX and GridData.chr_mu_zero for the Weil data (om), and WeilOver.exists_weil (Lemma -weil-assemble) with
% hmod.inf for theta_R and, for hV, T.gridFirstColumnSpace_rank_signature_selfDual at i = 0 and G.sp.V_is otherwise.
% (4) FunObj.gridChoice_nonempty_of (10 lines), followed: Good.badInt_notMem (badInt D = Q_S q N(c); the three cases
% locPl_cases: basePlace_over_qZero with leastSplitAuxiliaryPrimeThetaChain_spec, CoverHalf.g_coprime with
% ntypes.g_cond, the hypothesis hde with ccK_coprime, and LevCh.Is.Q_off with the hypothesis hc) and, for the rest,
% Lemma -grid-pieces-choice (exists_std, normal_nonempty, gridChoice_nonempty with count_AA_AAbar).
% Lean: (2) ChiGridModel.nonempty_gridData_iff. The direction from a grid to places is
% ChiGrid.Chars.basePlace_not_mem_chiGridBadPlaces, from the fields lev0 (LevelLe 1 and not UnrAt at the places
% over v0) and lev1 (LevelLe 2 and not LevelLe 1 at the places over S1) of ChiGrid.Chars at i = 1; these
% are conditions (d) and (e) on the characters in Definition th-grid. The
% converse is exists_gridData: GridChars.nonempty_chars (Lemma -grid-pieces-characters) and nonempty_vertices, followed in (a)-(e): exists_line for P_0 from Y_0 (hY0) and for Q_m from V_{m,m},
% exists_stair with the ideals P.stA sc = AA * modulBar^-1 and P.stAbar sc for V_{i,0} (colDec, colLat) and with
% AA, AAbar for V_{m,j}, 1 <= j < m - 1 (rowDecS, rowLat), StairConj.conj_AA_and_stA and
% stairId_ne_zero_and_mem_rev_iff for the condition on the ideals, and ChiGrid.exists_isV (choose S hS) for the
% other spaces (Lemma -grid-pieces-spaces). The self-dual lattice of Y_0 is colY_isV ... .sd.
% Lean: (1) GridExPf.exists_places_of_cls, written out: GridExPf.exists_primes (level system fields Pg,
% Q_de), leastSplitAuxiliaryPrimeThetaChain_spec and Mem.cls_eq_splitK_of_cond for (d),
% leastSplitAuxiliaryPrime_ne_typePrimes (with rPr_spec), exists_over_both for (b),
% under_conjugatePrime_ne_of_class for (a), eq_or_eq_conjP and prime_eq_of_mem for (c),
% MovedIdem.exists_isIdemAt for the idempotent at w(v) (the Lean's form of w(v) != conj w(v)),
% exists_prime_mem (Q_pl) and S1_mem_OmF. The Lean's "cls p = splitK" is the book's "p splits in K"
% (primeClass_eq_splitK_iff).
% Lean: (1) GridExPf.exists_places_of_cls, (2) ChiGridModel.nonempty_gridData_iff, (3)
% ChiGrid.WeilOver.exists_weil_over_of_clauses (from Published.WeilRepresentationUnitaryPair, a field of
% ThetaLocalInputs), (4) FunObj.gridChoice_nonempty_of. In all about 9,300 lines of Lean proofs at
% depth 5, counted from twisted_grid_existence_of_identity.
\end{proof}

\begin{lemma}\label{asm-dependence-test-assembly-first-layer-grid}
Assume the results of Theorem~\ref{lit-bundle-theta}. There is $C_{F,0}$ such that for every
$\varepsilon$ and every $C_F\ge C_{F,0}$ there is $x_0$ with the following property. For every
definite-side datum $D$ with size data of exponent $\varepsilon$ at a scale $x\ge x_0$, with
$b=e^{x^{\varepsilon}}$, and every realization $R$ of $D$ in an automorphic setting, with these size
data at a face (Definition~\ref{asm-dependence-functional-witness-setting}), whose twisted places lie
off $q_0$, there are an evaluation grid $G$ for $R$ with constant $C_F$ at $b$, a Weil datum $W$ over
$R$ for $G$ and a grid choice for $W$ (Definition~\ref{asm-dependence-grid-agree}(2)).
\end{lemma}
\leanref{Development/AbelianVarietyTheory/Definitions/TwistedInnerProducts/TwistPlacesOffPrimeQZero.lean}{GlobalTheory.RowEx.GridEx}
\leanref{Development/AutomorphicArithmetic/ThetaCorrespondence/ThetaPairLifts/LocalTypeDictionary/TwistedGridExistenceOfIdentity.lean}{GlobalTheory.RowEx.twisted_grid_existence_of_identity}
\begin{proof}
Take $C_{F,0}=0$ and $x_0=0$, and let $D$, its size data and $R$ be as in the statement. The rows of
the size data have their sizes, so the size data carry parameter choices; these make every level
prime split in $K$, in particular $g_0$, $g_1$, $p_{de}$ and $p'_{de}$, and the sizes give $q_0\ne
g_0,g_1$. The core conditions give a level system with base tuple $\mu_0$. The twisted places of $R$ lie
off $q_0$. So Lemma~\ref{asm-dependence-test-assembly-first-layer-grid-pieces}(1) gives places of a
theta grid whose base place lies outside $S_1$, and part (2) a theta grid $G$ for the level system of
$R$ with constants $(1,C_F,b)$, that is, an evaluation grid for $R$ with constant $C_F$ at $b$. The
core of $R$ satisfies the chain and model conditions, so part (3) gives a Weil datum $W$ over $R$ for
$G$. Finally $p_{de}$ and $p'_{de}$ are prime to $Q_Sq$ (Lemma~\ref{asm-dependence-rows-sizes}), the
set of tuples is finite, and by the guard no place of a set $Q_t$ contains $\mathrm N(\mathfrak c)$; so
part (4), applied to all tuples, gives a grid choice for $W$.
% Lean: RowEx.twisted_grid_existence_of_identity (with C_F0 = 0 and x_0 = 0): Pm.level_split,
% GridExPf.exists_places_of_cls with leastSplitAuxiliaryPrime_ne_typePrime_of_sizes,
% ChiGridModel.nonempty_gridData_iff, ChiGrid.WeilOver.exists_weil_over_of_clauses and
% FunObj.gridChoice_nonempty_of (with RowClaim.Test.pde_ne_q and RowClaim.Test.guard). About 9,300
% lines of Lean proofs at depth 5.
\end{proof}

\begin{lemma}\label{asm-dependence-test-assembly-first-layer-field}
Assume the results of Theorem~\ref{lit-bundle-theta}, the join statements and the functional
equation of Theorem~\ref{asm-dependence-inputs}, the theta hypotheses
(Definition~\ref{th-theta-hypotheses}), and that functionals separate
(Definition~\ref{asm-grid-evaluation}(7)). Then the number field statement holds
(Definition~\ref{asm-dependence-test-assembly-statements}(4)).
\end{lemma}
\leanref{Development/AutomorphicArithmetic/ThetaCorrespondence/ThetaPairLifts/LocalTypeDictionary/ThetaLiftExtensionFieldExists.lean}{GlobalTheory.T.theta_lift_extension_field_exists}
\leanref{Development/AutomorphicArithmetic/ThetaCorrespondence/ThetaPairLifts/LocalTypeDictionary/ThetaLiftExtensionFieldExists.lean}{GlobalTheory.RowEx.exists_extension_field_of_join}
\leanref{Development/AbelianVarietyTheory/Proofs/TwistedOperators/ExFieldOfKeyFinExKey.lean}{GlobalTheory.RowEx.exField_of_keyFin_exKey}
\leanref{Development/AbelianVarietyTheory/Proofs/TwistedOperators/ExKeyOf.lean}{GlobalTheory.RowEx.exKey_of}
\begin{proof}
The hypotheses of parts (2) and (3) of Lemma~\ref{asm-dependence-grid-pieces} hold, as in the proof of
Lemma~\ref{asm-dependence-grid-exists}. Let $C_{F,0}$ be the largest of the constants of
Lemma~\ref{asm-dependence-test-assembly-first-layer-grid} and of these two parts, let $\varepsilon$
and $C_F\ge C_{F,0}$ be given, and let $x_0$ be the largest of the numbers that they give.

Let $(x,\varepsilon)$ be a scale with $x\ge x_0$, put $b=e^{x^{\varepsilon}}$, and let $B$, $f$ and a
finite set $\Xi$ of Hecke characters be given. By Lemma~\ref{asm-dependence-grid-keys}(1) there is a
finite set $\mathcal K$ of keys such that every realization $R$ as in the number field statement has
its key in $\mathcal K$. Say that $(X,R)$ \emph{represent} $\kappa\in\mathcal K$ if $X$ is a later
stage over $(B,f)$ whose datum has size data at $(x,\varepsilon)$, and $R$ is a realization of this
datum in an automorphic setting with these size data at a face, whose columns are given by $f$,
whose twisted places lie off $q_0$ and whose key is $\kappa$.

For each $\kappa$ choose a finite set $S_\kappa$ of algebraic numbers. If no data represent $\kappa$,
put $S_\kappa=\emptyset$. Otherwise fix $(X_0,R_0)$ that represent it, with base tuple $\mu_0$.
Lemma~\ref{asm-dependence-test-assembly-first-layer-grid} gives an evaluation grid $G_0$ for $R_0$ with
constant $C_F$ at $b$, a Weil datum $W_0$ over $R_0$ and a grid choice $S_0$ for $W_0$.
Lemma~\ref{asm-dependence-grid-pieces}(2) gives further choices extending $S_0$ that satisfy the
evaluation conditions, and part (3) a number field $N\subset\mathbf C$ in which the numbers of
$(G_0,W_0,S_0)$ at $\mu_0$ lie. Let $S_\kappa$ be a basis of $N$ over $\mathbf Q$, and $S$ the union of
the $S_\kappa$, which is finite.

Now let $R$ be as in the number field statement, with base tuple $\mu'_0$, for a later stage $X$
whose coefficient field $L$ contains $S$, and let $\kappa$ be its key. Then $(X,R)$ represent $\kappa$,
so $(X_0,R_0)$ were fixed. As in the proof of Lemma~\ref{asm-dependence-grid-exists}, the level
systems of $R_0$ and $R$ agree through a map $\iota$ of the tuples with $\iota(\mu'_0)=\mu_0$. By
Lemma~\ref{asm-dependence-grid-transport}(1) to (3) there are an evaluation grid $G$ for $R$ with
constant $C_F$ at $b$ agreeing with $G_0$, a Weil datum $W$ over $R$ for $G$ and a grid choice $S'$ for
$W$, agreeing with $W_0$ and $S_0$ through $\iota$; the automorphic setting of $R$ gives the chain and
model conditions that part (2) needs. Lemma~\ref{asm-dependence-grid-pieces}(2), applied to $R$, $G$,
$W$ and $S'$, gives further choices extending $S'$ that satisfy the evaluation conditions; with $G$
and $W$ they are test choices for $C_F$. By Lemma~\ref{asm-dependence-grid-transport}(4), the numbers
of $(G,W,S')$ at $\mu'_0$ lie in $N$, and $L$ contains a basis of $N$, hence $N$. So the test choices
satisfy the field conditions.
% Lean: T.theta_lift_extension_field_exists = RowEx.exists_extension_field_of_join:
% exField_of_keyFin_exKey keyFin_holds (exKey_of (twisted_grid_existence_of_identity P)
% (isEx_of_separates ...) (nfAll_of ...) carryG carryW carryS trEasy trPhase trJet trBot). exKey_of:
% the constants max a1 (max a2 a3), the case split on representing data, the basis of N, LAgree from
% the keys (xi_inj, xi_base, betaX), and hAll'.numbersIn for FieldOK. The open parts are in the cited
% lemmas.
\end{proof}

\begin{lemma}\label{asm-dependence-test-assembly-first-layer}
Assume the results of Theorem~\ref{lit-bundle-theta}, the join statements and the functional
equation of Theorem~\ref{asm-dependence-inputs}, the theta hypotheses
(Definition~\ref{th-theta-hypotheses}), and that functionals separate
(Definition~\ref{asm-grid-evaluation}(7)). Then the seven statements on test choices and the number
field statement hold (Definition~\ref{asm-dependence-test-assembly-statements}).
\end{lemma}
\leanref{Development/AutomorphicArithmetic/ThetaCorrespondence/ThetaPairLifts/LocalTypeDictionary/TestVectorAtDatumOfJoin.lean}{GlobalTheory.RowTestAssembly.test_vector_at_datum_of_join}
\leanref{Development/AutomorphicArithmetic/ThetaCorrespondence/ThetaPairLifts/LocalTypeDictionary/ThetaLiftRowPropertiesConj.lean}{GlobalTheory.RowClaim.Join.thetaLift_row_properties_conj}
\leanref{Development/AutomorphicArithmetic/ThetaCorrespondence/ThetaPairLifts/LocalTypeDictionary/ThetaLiftExtensionFieldExists.lean}{GlobalTheory.T.theta_lift_extension_field_exists}
\begin{proof}
Statements (a), (c), (d), (e) and (f) are
Lemma~\ref{asm-dependence-test-assembly-first-layer-rows}. Statement (b) is
Lemma~\ref{asm-dependence-test-assembly-first-layer-injective}, whose hypotheses are the separation of
functionals and Lemma~\ref{asm-dependence-test-assembly-first-layer-convergence}. Statement (g) is
Lemma~\ref{asm-dependence-test-assembly-first-layer-ratio}, and the number field statement is
Lemma~\ref{asm-dependence-test-assembly-first-layer-field}. Each of these lemmas assumes only results
among those assumed here.
% Lean: test_vector_at_datum_of_join. s1, s2, s3 and j = RowClaim.Join.thetaLift_row_properties_conj
% (Lemma asm-dependence-test-assembly-first-layer-rows) give SumOKK (j.1), ConstK (j.2.1), LinK
% (j.2.2.1), Formula2K and SumRatK (j.2.2.2); InjK = RowEx.injK_of_separates_linB T3
% (twisted_linear_bound P (thetaPairLift_cuspidal_of_hypotheses P)) (Lemmas
% asm-dependence-test-assembly-first-layer-injective and -convergence); RatioK =
% RowClaim.twisted_local_factor_ratio_holds P K LM j.2.1 l (triEx_holds_alt) (tower_alt_bis) (Lemma
% asm-dependence-test-assembly-first-layer-ratio); ExField = T.theta_lift_extension_field_exists
% (Lemma asm-dependence-test-assembly-first-layer-field).
\end{proof}

\subsection{The row test}\label{subsection-asm-row-test}

\begin{definition}\label{asm-dependence-test-assembly-true}
Let $B$ be base data with a framing $f$, whose host has level primes $g_0,g_1$ and
$K=\mathbf Q(\sqrt{-d})$, and let $Z$ and $Z'$ be later stages of the same column choice
(Definition~\ref{asm-stage-construction}), with the groups
$G_0(Z)=\prod_{k=0,1}\prod_{v\in P_G(k)}GL_r(\mathbf F_{g_k})$ of their definite-side data.
\begin{enumerate}
\item The group data of $Z$ are \emph{true} for $f$ if the place of $E$ of every $v\in P_G(k)$ lies
  over the prime of $K$ above $g_k$ that $f$ chooses; the Teichmüller characters and the series
  idempotents of $Z$ at $g_k$ are those of the label datum of $g_k$; and its characters at the
  places of the group are the powers $\omega_k^{2(j+1)}$ of the character $\omega_k$ of
  $(\mathbf Z/g_k)^\times$ that $f$ chooses.
\item $Z$ and $Z'$ have the \emph{same group} if there are bijections from $P_G(k)$ for $Z$ to
  $P_G(k)$ for $Z'$, $k=0,1$, that preserve the places of $E$ of the group, and under the induced
  isomorphism $G_0(Z)\cong G_0(Z')$ the type idempotents of the two data have the same coefficients.
\end{enumerate}
\end{definition}
\leanref{Development/AbelianVarietyTheory/Proofs/IwasawaDecomposition/TypeIdempotentCoeffAcrossLayers.lean}{GlobalTheory.Mem.Top.GrpTrue}
\leanref{Development/AbelianVarietyTheory/Proofs/IwasawaDecomposition/TypeIdempotentCoeffAcrossLayers.lean}{GlobalTheory.Mem.Top.TypesTrue}
\leanref{Development/AbelianVarietyTheory/Definitions/AutomorphicMembership/WeakCharInertPartRealAdelicDataTransport.lean}{GlobalTheory.Mem.R3.SameGrp}
\leanref{Development/AbelianVarietyTheory/Proofs/IwasawaDecomposition/ExistsBare.lean}{GlobalTheory.Mem.R3.LedRec.sameGrp}

Two later stages whose group data are both true for $f$ have the same group: the sets $P_G(k)$ of
both are identified with one set of places over $g_k$, and the true data give the same places of $E$
and the same coefficients of the type idempotents.

\begin{definition}\label{asm-dependence-test-assembly-realization-equivariant-levels}
Let $D$ be a definite-side datum with host field $E/F$, and let $V$ be a Hermitian space of $D$ of
rank $n$ (Definition~\ref{th-spaces}) with an $\mathcal O_E$-lattice $\Lambda\subset E^n$ and an
adelic matrix $\beta\in GL_n(E\otimes_F\mathbf A_{F,f})$ that is a frame of $\Lambda$: the completion
$\hat\Lambda$ consists of the vectors $x$ with $\beta^{-1}x$ integral. Let $\operatorname{St}(\beta)$
be the group of the matrices $\beta y\beta^{-1}$ with $y\in GL_n(\hat{\mathcal O}_E)$, and for
$x\in\operatorname{St}(\beta)$ and a finite place $P$ of $E$ with residue field $k_P$ let
$\operatorname{red}_P(x)\in GL_n(k_P)$ be the reduction of $\beta^{-1}x\beta$ at $P$. A \emph{level
datum} $L$ on $V$ consists of integers $0\le r\le n$, $t\ge1$ and $s\ge0$, finite sets $P_g$ and $Q$
of finite places of $E$, a subgroup of $G_0=\prod_{P\in P_g}GL_r(k_P)$ with a character, and
characters $\chi_{c,P}$ ($1\le c\le s$, $P\in P_g$) and $\xi_P$ ($P\in Q$) of the groups $k_P^\times$.
Call the sets $\{1,\dots,r\}$ and $\{r+(c-1)t+1,\dots,r+ct\}$, $c\ge1$, of indices the \emph{flag
blocks}, and $\{1,\dots,r\}$ and $\{r+1,\dots,n\}$ the \emph{two blocks}.
\begin{enumerate}
\item The \emph{parahoric level} $K$ of $(\beta,L)$ is the group of the $\kappa\in U(V)(\mathbf A)$
  with archimedean components $1$ whose finite part lies in $\operatorname{St}(\beta)$ and whose
  reduction $\operatorname{red}_P(\kappa)$ is block upper triangular for the flag blocks at every
  $P\in P_g$ and for the two blocks at every $P\in Q$.
\item The map $\operatorname{red}\colon K\to G_0$ that takes $\kappa$ to the upper left $r\times r$
  blocks of the matrices $\operatorname{red}_P(\kappa)$, $P\in P_g$, is a homomorphism, since these
  matrices are block upper triangular. A \emph{level lift} of $g_0\in G_0$ is an element
  $\kappa\in K$ with $\operatorname{red}(\kappa)=g_0$. The kernel $J$ of $\operatorname{red}$ is the
  \emph{level} of $(\beta,L)$.
\item The \emph{level character} of $(\beta,L)$ is the character $\chi$ of $K$ given by
  \[
   \chi(\kappa)=\prod_{P\in P_g}\prod_{c=1}^{s}\chi_{c,P}\bigl(\det d_c(\operatorname{red}_P\kappa)\bigr)
   \prod_{P\in Q}\xi_P\bigl(\det d'(\operatorname{red}_P\kappa)\bigr),
  \]
  where $d_c(x)$ is the diagonal block of $x$ on the indices $r+(c-1)t+1,\dots,r+ct$ and $d'(x)$
  the one on the indices $r+1,\dots,n$.
\item A function $\Phi$ on $U(V)(\mathbf A)$ has the \emph{type} of $(\beta,L)$ if
  $\Phi(h\kappa)=\chi(\kappa)\Phi(h)$ for all $\kappa\in J$ and $h\in U(V)(\mathbf A)$. Fix a level
  lift $\kappa_{g_0}$ of each $g_0\in G_0$ that has one. The \emph{level action} of $g_0$ on such
  functions is $a_{g_0}\Phi(h)=\chi(\kappa_{g_0})^{-1}\Phi(h\kappa_{g_0})$, and $a_{g_0}\Phi=0$ if
  $g_0$ has no level lift.
\item Let $R$ be a realization of $D$ (Definition~\ref{is-realisation}) with places $P_g$, ramified
  sets $Q_t$ of its tuples and spaces $Y_0,\dots,Y_{m-1}$ and $Y_m=X'$ of ranks $r+jt$. The level
  data of $R$ give, for every tuple $t$ and $0\le j\le m$, a level datum on $Y_j$ with the lattice
  and adelic matrix of $R$ at $Y_j$, the integers $r$, $t$ and $s=j$, the sets $P_g$ and $Q_t$, the
  trivial subgroup of $G_0$, and residue characters of the characters of $R$. Its parahoric level,
  level, level character and level lifts are those \emph{of $Y_j$ at $t$}. The \emph{trivial level
  datum} of $R$ on $Y_0$ and on $W_i$ has the integers $r$, $t$ and $s=0$, the sets $P_g$ and
  $Q=\emptyset$ and the trivial subgroup; its level character is trivial.
\end{enumerate}
\end{definition}
\leanref{Formalism/LocalToGlobal/Divisibility/DefGroup.lean}{Formalism.L2.DIV.parLev}
\leanref{Formalism/LocalToGlobal/Divisibility/DefGroup.lean}{Formalism.L2.DIV.rTup}
\leanref{Formalism/LocalToGlobal/Divisibility/DefGroup.lean}{Formalism.L2.DIV.lev}
\leanref{Formalism/LocalToGlobal/Divisibility/DefGroup.lean}{Formalism.L2.DIV.levFn}
\leanref{Formalism/LocalToGlobal/IndefiniteReal/SubgroupCharExtensionOneApply.lean}{Formalism.IndefR.Lev.IsLevelLift}
\leanref{Formalism/LocalToGlobal/IndefiniteReal/SubgroupCharExtensionOneApply.lean}{Formalism.IndefR.Lev.levG}
\leanref{Formalism/LocalToGlobal/IndefiniteReal/SubgroupCharExtensionOneApply.lean}{Formalism.IndefR.Lev.levChi}
\leanref{Formalism/LocalToGlobal/IndefiniteReal/SubgroupCharExtensionOneApply.lean}{Formalism.IndefR.Lev.HasType}
\leanref{Formalism/LocalToGlobal/IndefiniteReal/SubgroupCharExtensionOneApply.lean}{Formalism.IndefR.Lev.actG0}
\leanref{Formalism/LocalToGlobal/IndefiniteReal/LinksDefs.lean}{Formalism.IndefR.Links.LevCh.datYt}
\leanref{Formalism/LocalToGlobal/IndefiniteReal/LinksDefs.lean}{Formalism.IndefR.Links.LevCh.datXt}
\leanref{Formalism/LocalToGlobal/IndefiniteReal/LinksDefs.lean}{Formalism.IndefR.Links.LevCh.datWp}

\begin{lemma}\label{asm-dependence-test-assembly-realization-equivariant-lifts}
\begin{enumerate}
\item Let $V$, $\Lambda$, $\beta$ and a level datum $L$ with places $P_g$ and $Q$ be as in
  Definition~\ref{asm-dependence-test-assembly-realization-equivariant-levels}. Suppose that
  $\beta$ is a \emph{base frame}: $\Lambda$ is self-dual for the Hermitian form of $V$, and
  $N\Lambda_0\subset\Lambda$ and $N\Lambda\subset\Lambda_0$ for some integer $N>0$, where $\Lambda_0$
  is the lattice of the vectors with coordinates in $\mathcal O_E$ in the basis of $V$. Suppose also
  that no place of $P_g$ is complex conjugate to a place of $P_g\cup Q$, itself included, and that
  $P_g\cap Q=\emptyset$. Then every element of $G_0$ has a level lift.
\item Let $R$ be a realization of a definite-side datum $D$ (Definition~\ref{is-realisation}) that
  satisfies the core conditions on its level characters (Definition~\ref{is-eigenvalue-lift-data}).
  Then for every tuple $t$ and every $0\le j<m$, every element of $G_0$ has a level lift for $Y_j$ at
  $t$.
\end{enumerate}
\end{lemma}
\leanref{Development/AbelianVarietyTheory/Proofs/TwistedOperators/ParahoricLevelLiftsExistStatement.lean}{GlobalTheory.T.parahoricLevel_liftsExist}
\leanref{Development/AbelianVarietyTheory/Proofs/TwistedOperators/ParahoricLevelLiftsExist.lean}{GlobalTheory.ParahoricLift.exists_local_lift}
\leanref{Development/AbelianVarietyTheory/Proofs/TwistedOperators/ParahoricLevelLiftsExist.lean}{GlobalTheory.ParahoricLift.exists_lift_prod}
\leanref{Development/AbelianVarietyTheory/Proofs/TwistedOperators/ChainCompositeLiftIntertwinesFiniteGroupActions.lean}{GlobalTheory.ChainLiftEquiv.liftsExist_chain}
\begin{proof}
(1) Let $g_0\in G_0$, and for $P\in P_g$ let $k_P\in GL_n(k_P)$ be the block diagonal matrix with
blocks $g_{0,P}$ and $1$.

\emph{One place.} Let $P\in P_g$. As $P$ is not conjugate to itself, it lies over a place $v$ of $F$
that splits in $E$, and the idempotent $1_P$ of $\mathbf A_E$ at $P$ satisfies $1_P\bar1_P=0$, with
$1_P+\bar1_P$ the idempotent of $v$. Lift $k_P$ to a matrix $Y$ over $\mathcal O_E$. Its determinant
does not lie in $P$, as it reduces to $\det k_P\ne0$, so there is $u\in\hat{\mathcal O}_E$ with
$u\det Y$ equal to the finite part of $1_P$, and $u\operatorname{adj}Y$ is an integral inverse of $Y$
at $P$. So there is $\kappa_P\in U(V)(\mathbf A)$ that is $1$ away from $v$ and whose component at $P$
is $\beta Y\beta^{-1}$; its archimedean components are $1$. The components of $\kappa_P$ and
$\kappa_P^{-1}$ at $P$ preserve $\hat\Lambda$, as $Y$ and its inverse at $P$ are integral. The
component at $\bar P$ of a unitary element is the inverse of the adjoint of its component at $P$,
and $\hat\Lambda$ is self-dual; so $\kappa_P$ and $\kappa_P^{-1}$ preserve $\hat\Lambda$, and
$\kappa_P\in\operatorname{St}(\beta)$. Its reduction at $P$ is $k_P$, since $1_P$ reduces to $1$ at
$P$. At every place $P'\ne P$ not conjugate to $P$ the element $\kappa_P$ is $1$, so its reduction
there is $1$.

\emph{The product.} Put $\kappa=\prod_{P\in P_g}\kappa_P$. Reduction at a place is multiplicative on
$\operatorname{St}(\beta)$. For $P\in P_g$ every factor $\kappa_{P'}$ with $P'\ne P$ reduces to $1$ at
$P$, since $P$ is not conjugate to $P'$; so $\operatorname{red}_P(\kappa)=k_P$, which is block
diagonal, hence block upper triangular for the flag blocks. For $P\in Q$, $P$ is neither a place of
$P_g$ nor conjugate to one, so every factor reduces to $1$ at $P$, and $\operatorname{red}_P(\kappa)=1$.
Hence $\kappa\in K$, and the upper left blocks of its reductions at the places of $P_g$ are the
$g_{0,P}$: $\kappa$ is a level lift of $g_0$.

(2) By the core conditions on the level characters, the adelic matrix of $R$ at $Y_j$ is a base
frame of its lattice; every place of $P_g$ lies over one of the primes $\mathfrak p_0,\mathfrak p_1$
of the imaginary quadratic subfield $\mathbf Q(\sqrt{-d})$ of $E$ above $g_0,g_1$ chosen by the data
of $R$; and every place of $Q_t$ contains a prime number different from $g_0$ and $g_1$. Complex
conjugation fixes the rational integers, so a place conjugate to a place that contains a prime number
contains it too.

Let $P,P'\in P_g$ lie over $\mathfrak p_k$ and $\mathfrak p_{k'}$, and suppose $P'=\bar P$. Then $P'$
contains $g_k$ and $g_{k'}$, which are distinct primes if $k\ne k'$; so $k=k'$. As $g_k$ splits in
$\mathbf Q(\sqrt{-d})$ into $\mathfrak p_k\bar{\mathfrak p}_k$ with $\mathfrak p_k\ne\bar{\mathfrak
p}_k$, a place of $E$ over $g_k$ lies over $\mathfrak p_k$ if and only if its conjugate does not, while
$P$ and $\bar P=P'$ both lie over $\mathfrak p_k$; this is impossible. Let $P\in P_g$ lie over
$\mathfrak p_k$ and let $P'\in Q_t$ contain the prime $q\ne g_0,g_1$. If $P'$ were $P$ or $\bar P$, it
would contain the two primes $g_k$ and $q$. So no place of $P_g$ is conjugate to a place of $P_g\cup
Q_t$, and $P_g\cap Q_t=\emptyset$, and (1) applies to the level datum of $Y_j$ at $t$.
% Lean: (1) ParahoricLift.liftsExist = exists_lift_prod (the product over Finset.univ.toList of the
% kappa P, redFn_list_prod, prod_map_single, diagOne_mem, topLeftBlock_blockDiagWithIdentityGL) with
% exists_local_lift: exists_isIdemAt (IsIdemAt e P from not IsConjPl P P), exists_idem_int,
% exists_matrix_lift, finSw_div_int (u * det Y = p), conj_inv_part, split_lift_fin
% (SplitLift.split_lift, SuppAt e), mem_St_of_part with hatL_selfDual (from the base frame) and
% SplitStab.stab_of_half, redFn_of_smul_eq at P, and redFn_eq_one with exists_sep and idem_div_int
% away from P and its conjugate. (2) ChainLiftEquiv.liftsExist_chain: hL.betaY j (IsSDFrame: a base
% frame with N^c), ProjectedLift.not_isConjPl (Chr.splitsC_main for the two places over pb k) with
% hL.Pg, LevelRamif.exists_prime_mem (hL.Q_pl), not_two_primes, natCast_mem_conj, typePrime_mem, and
% ramif_at_levels for the disjointness. About 6,800 lines of Lean proofs at depth 5 by the count used
% here.
\end{proof}

\begin{lemma}\label{asm-dependence-test-assembly-realization-equivariant-shift}
Let $V$, $\beta$ and a level datum $L$ be as in
Definition~\ref{asm-dependence-test-assembly-realization-equivariant-levels}, with parahoric level
$K$, level $J$, level character $\chi$ and level action $a$. Let $\theta$ be an archimedean family
and $M$ a model of discrete series for $U(V)$ (Definition~\ref{asm-dependence-rows-models}), and
let $\kappa$ be a level lift of some $g_0\in G_0$.
\begin{enumerate}
\item For every function $\Phi$ of the type of $(\beta,L)$ and every $h$,
  $a_{g_0}\Phi(h)=\chi(\kappa)^{-1}\Phi(h\kappa)$.
\item For every $u\in\operatorname{Mult}_\theta(M;J,\chi)$, the map $u'$ with
  $u'(v)(h)=\chi(\kappa)^{-1}u(v)(h\kappa)$ lies in $\operatorname{Mult}_\theta(M;J,\chi)$.
\end{enumerate}
\end{lemma}
\leanref{Development/AbelianVarietyTheory/Proofs/TwistedOperators/ChainCompositeLiftIntertwinesFiniteGroupActions.lean}{GlobalTheory.ChainLiftEquiv.exists_shift}
\leanref{Formalism/LocalToGlobal/IndefiniteReal/SubgroupCharExtensionOneApply.lean}{Formalism.IndefR.Lev.level_action_spec}
\leanref{Development/AbelianVarietyTheory/Proofs/TwistedOperators/FiniteGroupActionPreservesPrincipalFlagMultiplic.lean}{GlobalTheory.FlagAction.finiteGroupAction_preserves_principalFlagMultiplicity}
\begin{proof}
(1) Let $\kappa_0=\kappa_{g_0}$ be the level lift chosen in the level action. Then
$\kappa_0^{-1}\kappa\in K$ and $\operatorname{red}(\kappa_0^{-1}\kappa)=g_0^{-1}g_0=1$, so
$\kappa_0^{-1}\kappa\in J$. As $\chi$ is a character of $K$,
$\chi(\kappa)=\chi(\kappa_0)\chi(\kappa_0^{-1}\kappa)$. Hence
\[
 \chi(\kappa)^{-1}\Phi(h\kappa)=\chi(\kappa)^{-1}\Phi(h\kappa_0\cdot\kappa_0^{-1}\kappa)
 =\chi(\kappa)^{-1}\chi(\kappa_0^{-1}\kappa)\Phi(h\kappa_0)=\chi(\kappa_0)^{-1}\Phi(h\kappa_0)
 =a_{g_0}\Phi(h).
\]

(2) The multiplicity maps of level $(J,\chi)$ form a subspace, as their conditions are linear, so
$u$ is one. Its values have the type of $(\beta,L)$, so by (1), $u'(v)=\chi(\kappa_0)^{-1}T(v)$ with
$T(v)(h)=u(v)(h\kappa_0)$. The map $T$ is a multiplicity map of level $(J,\chi)$. Right translation
by $\kappa_0$, whose archimedean components are $1$, carries cuspidal automorphic forms that are
smooth at the archimedean places, finite under the centres of the enveloping algebras and of
moderate growth to such forms; it commutes with right translation by the images of the archimedean
groups under $\theta$, so $T$ has the equivariance at the archimedean places; and the functional
$v_\tau\mapsto T(v)(x)$ is $v_\tau\mapsto u(v)(x\kappa_0)$, continuous for a Sobolev norm. For the
level, let $k\in J$. Then $\kappa_0^{-1}k\kappa_0\in K$ and
$\operatorname{red}(\kappa_0^{-1}k\kappa_0)=\operatorname{red}(\kappa_0)^{-1}\operatorname{red}(k)\operatorname{red}(\kappa_0)=1$,
so $\kappa_0^{-1}k\kappa_0\in J$; and $\chi(\kappa_0^{-1}k\kappa_0)=\chi(k)$, as $\chi$ is a character
of $K$. So $T(v)(hk)=u(v)(h\kappa_0\cdot\kappa_0^{-1}k\kappa_0)=\chi(k)T(v)(h)$. Hence $T$, and with it
$u'$, is a multiplicity map of level $(J,\chi)$.
% Lean: ChainLiftEquiv.exists_shift = exists_translate: (1) is level_action_spec (lift_div_mem,
% levFn_flag_mul, levFn_flag1); (2) is FlagAction.finiteGroupAction_preserves_principalFlagMultiplicity
% (U = levChi(kappa_0)^{-1} heckeOn M {kappa_0} T, SphericalOps.isMultMap_heckeOn with isMultMap_single
% for the conditions other than the level, and conj_mem_flag1 for the level), with MultMem.mem_mult_iff
% for the span. The Lean's level J is levG of flag1L L and its character levChi of flag1L L, which
% equals levChi of flagL L by levFn_flag1. About 1,900 lines of Lean proofs at depth 5.
\end{proof}

\begin{lemma}\label{asm-dependence-test-assembly-realization-equivariant-integral}
Let $V$, $\beta$ and a level datum $L$ be as in
Definition~\ref{asm-dependence-test-assembly-realization-equivariant-levels}, with the group $G_0$
of $L$, and let $\kappa$ be a level lift of some $g_0\in G_0$. Let $\nu$ be a Haar measure on
$U(V)(\mathbf A)$ and $\mathcal F$ a fundamental domain for $U(V)(F)$ with respect to $\nu$. Let $f$ be
a function on $U(V)(\mathbf A)$ that is invariant on the left under $U(V)(F)$ and integrable on
$\mathcal F$. Then
\[
 \int_{\mathcal F}f(h\kappa)\,d\nu(h)=\int_{\mathcal F}f(h)\,d\nu(h).
\]
\end{lemma}
\leanref{Development/AbelianVarietyTheory/Proofs/TwistedOperators/LiftLowestKTypeCompatibility.lean}{GlobalTheory.ThetaProjEquiv.integral_shift}
\leanref{Development/AbelianVarietyTheory/Proofs/TwistedOperators/IsCompactParahoricLevelAndContainsLifts.lean}{GlobalTheory.ParahoricCompact.isCompact_levG}
\leanref{Development/AbelianVarietyTheory/Proofs/TwistedOperators/HaarMapMulRightOfMemCompactSubgroup.lean}{GlobalTheory.haar_map_mul_right_of_mem_compact_subgroup}
\leanref{Development/AbelianVarietyTheory/Proofs/TwistedOperators/FundamentalDomainSetIntegralMulRight.lean}{GlobalTheory.fundamentalDomain_setIntegral_mul_right}
\begin{proof}
Let $K$ be the parahoric level of $(\beta,L)$; by definition the level lift $\kappa$ lies in $K$. The
group $K$ is compact: the elements of $U(V)(\mathbf A)$ with archimedean components $1$ and
finite part in the compact group $\operatorname{St}(\beta)$ form a compact set, and $K$ is a closed
subset of it, consisting of the elements of this set whose finite part lies in a subgroup of
$\operatorname{St}(\beta)$ defined by congruence conditions, which is open, hence closed. So right
translation by $\kappa$ preserves $\nu$. Indeed, the image of $\nu$ under $h\mapsto h\kappa$ is
invariant under left translations, so it is $c\nu$ for some $c\ge0$. For a compact neighbourhood $C$
of $1$ the set $CK$ is compact and contains a neighbourhood of $1$, so $0<\nu(CK)<\infty$; and its
preimage under $h\mapsto h\kappa$ is $CK\kappa^{-1}=CK$. Hence $c=1$.

The function $h\mapsto f(h\kappa)$ is invariant on the left under $U(V)(F)$, and the preimage
$\mathcal F\kappa^{-1}$ of $\mathcal F$ under $h\mapsto h\kappa$ is again a fundamental domain, since
this map preserves $\nu$ and commutes with left translations. For a function invariant on the left
under $U(V)(F)$, integrability and the integral are the same over any two fundamental domains. So
$h\mapsto f(h\kappa)$ is integrable on $\mathcal F$ if and only if it is on $\mathcal F\kappa^{-1}$,
that is, if and only if $f$ is integrable on $\mathcal F$, and
\[
 \int_{\mathcal F}f(h\kappa)\,d\nu(h)=\int_{\mathcal F\kappa^{-1}}f(h\kappa)\,d\nu(h)
 =\int_{\mathcal F}f(h)\,d\nu(h).
\]
% Lean: ThetaProjEquiv.integral_shift (hκ : IsLevelLift V β L g0 κ; the conclusion is the equality of
% the integrals only, the .2 of fundamentalDomain_setIntegral_mul_right, whose .1 is the integrability
% used in the second paragraph): isCompact_flag_and_lift (isCompact_levG for the parahoric level levG of
% flagL L, and that a level lift lies in it), haar_map_mul_right_of_mem_compact_subgroup and
% fundamentalDomain_setIntegral_mul_right. The proof uses only that kappa lies in the parahoric
% level. About 4,600 lines of Lean proofs at depth 5 by the count used here, most of them the
% topology of the adelic unitary group.
\end{proof}

\begin{lemma}\label{asm-dependence-test-assembly-realization-equivariant-scalars}
Let $R$ be a realization of a definite-side datum $D$ for a family of definite forms for a finite
group $G$ (Definition~\ref{is-realisation}) that satisfies the core conditions on its level
characters, its models and its chain (Definition~\ref{is-eigenvalue-lift-data}), and identify $G$
with $G_0=\prod_{P\in P_g}GL_r(k_P)$ by the isomorphism of $R$. Let $t$ be a tuple and $0\le j<m$,
write $Y_m=X'$, and let $\omega$ be the Weil datum of the pair $(Y_j,Y_{j+1})$ at $t$ and $v$ its test
vector, with $\chi_j$ and $\chi_{j+1}$ the level characters of $Y_j$ and $Y_{j+1}$ at $t$
(Definition~\ref{asm-dependence-test-assembly-realization-equivariant-levels}). Let $g\in G$, and
let $\kappa$ and $\kappa'$ be level lifts of $g$ for $Y_j$ and for $Y_{j+1}$ at $t$. Then
\[
 \omega(\kappa,\kappa')(v\otimes P)=\chi_j(\kappa)\,\chi_{j+1}(\kappa')^{-1}\,(v\otimes P)
\]
for every polynomial $P$ in the creation operators, where $v\otimes P$ is formed with the Cartan
involutions and bases of the models of $R$ (Lemma~\ref{is-eigenvalue-conjugate-theta-facts}).
\end{lemma}
\leanref{Development/AbelianVarietyTheory/Proofs/TwistedOperators/StepVectorLiftPairFlagLevelRealizedChainStatement.lean}{GlobalTheory.T.stepVector_lift_pair_flagLevel_realized_chain}
\leanref{Development/AbelianVarietyTheory/Proofs/TwistedOperators/LastStepVectorLiftPairFlagLevelCharactersStatement.lean}{GlobalTheory.T.lastVector_lift_pair_flagLevel_characters}
\leanref{Development/AbelianVarietyTheory/Proofs/TwistedOperators/OmPureTensorOOfFinite.lean}{GlobalTheory.FiniteFock.om_pureTensorO_of_finite}
\begin{proof}
Write $c=\chi_j(\kappa)\chi_{j+1}(\kappa')^{-1}$. The case $P=1$ is the identity
$\omega(\kappa,\kappa')v=cv$. For general $P$: the elements $\kappa$ and $\kappa'$ have archimedean
components $1$, which commute with the Cartan involutions and fix the bases. So
Lemma~\ref{is-eigenvalue-conjugate-theta-facts}(4), with the same Cartan involutions and bases on
both sides, gives $\omega(\kappa,\kappa')(v\otimes P)=(\omega(\kappa,\kappa')v)\otimes P=(cv)\otimes
P=c\,(v\otimes P)$.
\openstep{The action of the Weil datum of a pair of the chain at a pair of level lifts on the test
vector of the pair, the case $P=1$, is not yet written out: on the Fock vacuum it is a product of
residue characters of the determinants of the reductions, computed place by place, and the test
vector differs from the vacuum at the places of $P_g$ and $Q_t$ by a function whose translates under
the block triangular reductions give the level characters.}
% Lean: the case P = 1 is StepVecLaw.stepVec_lifts for j + 1 < m and StepLaw.last_step_law for
% j = m - 1, both through step_law: om_vac (eigen_of_local_components_finiteAdelic_pair; om_vac_split
% through om_vacuum_eq_splitFac_smul_sum at the split places of P_g u Q_t,
% weilOperator_fixes_latticeVacuum_nonsplit_unramified elsewhere),
% schrodingerIntertwiner_schwartzBruhat_isTheta and weilRep_theta_datum_half_eq_smul for the test vector
% (IsStepVecOf), typePlaceDatum_blockTriangular_move at P_g, twistPlaceDatum_parabolic_move at Q_t,
% levFn_flag_datY and unit_product_rescale_disjoint_finsets for the product of characters, with the
% restrictions of the characters of the pair from levRes_chrStep. About 13,700 lines (stepVec_lifts)
% and 11,300 (last_step_law) of Lean proofs at depth 5 by the count used here, overlapping. The
% passage to v (x) P is om_pureTensorO_of_finite (OrientedHeis.om_pureTensorO with compAt_mt_of_finite).
\end{proof}

\begin{lemma}\label{asm-dependence-test-assembly-realization-equivariant-projection}
Let $R$ be a realization of a definite-side datum $D$ for a family of definite forms for a finite
group $G$, whose idempotent is $e=\sum_{g\in G}e(g)\,g$ with $e(g)$ in the coefficient field $L$, and
suppose that $R$ records the group
(Definition~\ref{asm-dependence-functional-witness-setting}(7)). Identify $G$ with
$G_0=\prod_{P\in P_g}GL_r(k_P)$ by the isomorphism of $R$. Let $V$ be $Y_0$ or one of the spaces
$W_i$, with the trivial level datum of $R$ on it and its level action $a$
(Definition~\ref{asm-dependence-test-assembly-realization-equivariant-levels}). For an embedding
$\sigma$ of $L$ put $e_\sigma=\sum_{y\in G}\sigma(e(y))\,a_y$. Then
\[
 e_\sigma(a_g\Phi)=a_g(e_\sigma\Phi)
\]
for every $g\in G$ and every function $\Phi$ on $U(V)(\mathbf A)$ of the type of the trivial level
datum.
\end{lemma}
\leanref{Development/AbelianVarietyTheory/Proofs/TwistedOperators/LevelGroupActionEqZeroOfNoLift.lean}{GlobalTheory.ProjComm.proj_comm}
\leanref{Development/AbelianVarietyTheory/Proofs/TwistedOperators/LevelGroupActionEqZeroOfNoLift.lean}{GlobalTheory.ProjComm.gaL_comm}
\leanref{Development/AbelianVarietyTheory/Proofs/TwistedOperators/LevelGroupActionEqZeroOfNoLift.lean}{GlobalTheory.ProjComm.actL_conj}
\begin{proof}
\emph{The coefficients.} The function $y\mapsto e(y)$ is constant on the conjugacy classes of $G$.
Indeed, by the record of the group, $\sigma_D(e(y))$ is the coefficient of the type idempotent $e_0$
of $D$ at the image of $y$ in $\prod_{k=0,1}\prod_{v\in P_G(k)}GL_r(\mathbf F_{g_k})$. The coefficient
of $e_0$ at an element is a sum of products, over the factors, of coefficients of series idempotents,
which are central in the group algebras of the factors; and the coefficients of a central element
of a group algebra are constant on conjugacy classes. As $\sigma_D$ is injective,
$e(gyg^{-1})=e(y)$, and hence $\sigma(e(gyg^{-1}))=\sigma(e(y))$.

\emph{Conjugation.} Let $g,h\in G_0$ and let $\Phi$ have the type. Then
$a_{ghg^{-1}}a_g\Phi=a_ga_h\Phi$. Since $\operatorname{red}$ is a homomorphism, products and inverses
of level lifts are level lifts of the products and inverses. If $g$ and $g'$ have level lifts, then
$a_ga_{g'}\Phi=a_{gg'}\Phi$: with $\kappa'$ a level lift of $g'$, $\kappa_g\kappa'$ is one of $gg'$, and
by Lemma~\ref{asm-dependence-test-assembly-realization-equivariant-shift}(1) and the
multiplicativity of the level character on the parahoric level,
\[
 a_ga_{g'}\Phi(x)=\chi(\kappa_g)^{-1}\chi(\kappa')^{-1}\Phi(x\kappa_g\kappa')
 =\chi(\kappa_g\kappa')^{-1}\Phi(x\kappa_g\kappa')=a_{gg'}\Phi(x).
\]
So if $g$ and $h$ have level lifts, so does $ghg^{-1}$, and $a_{ghg^{-1}}a_g\Phi=a_{gh}\Phi=a_ga_h\Phi$.
If $g$ has a level lift $\kappa$ but $h$ has none, then $ghg^{-1}$ has none, since for a level lift
$\mu$ of it $\kappa^{-1}\mu\kappa$ would be one of $h$; both sides vanish. If $g$ has none, both sides
vanish.

\emph{The sum.} Substituting $y=gzg^{-1}$,
\[
 e_\sigma(a_g\Phi)=\sum_{z\in G}\sigma(e(gzg^{-1}))\,a_{gzg^{-1}}a_g\Phi
 =\sum_{z\in G}\sigma(e(z))\,a_ga_z\Phi=a_g(e_\sigma\Phi).
\]
% Lean: ProjComm.proj_comm: hcls from the record (eR, he) and goodProjector_coeff_classFunction
% (coeff_conj_of_central with eSer_central), injectivity of embLL; gaL_comm (Fintype.sum_equiv along
% MulAut.conj g) with actL_conj (levelGroupAction_mul_of_lifts through level_action_spec,
% levChi_flagL_mul_of_mem_parLev, isLift_mul, isLift_inv, levelGroupAction_eq_zero_of_no_lift). The
% Lean's e0 is gaL at Y0 with sigma = embLL, and eWf i sigma at W_i. About 1,300 lines of Lean proofs
% at depth 5 for gaL_comm.
\end{proof}

\begin{lemma}\label{asm-dependence-test-assembly-realization-equivariant-schur-slot-closable-schur}
Let $H$ be a complex Hilbert space and $\mathcal S$ a set of bounded operators on $H$ that contains the
adjoint of each of its members and leaves no closed subspace of $H$ other than $0$ and $H$ invariant.
Let $T$ be a bounded operator on $H$ with $T\sigma=\sigma T$ for every $\sigma\in\mathcal S$. Then $T$ is
a scalar multiple of the identity.
\end{lemma}
\leanref{Development/AbelianVarietyTheory/Proofs/IwasawaDecomposition/HilbertSchurOfAdjointClosed.lean}{GlobalTheory.hilbert_schur_of_adjoint_closed}
\begin{proof}
\emph{Self-adjoint operators.} Let $A$ be a bounded self-adjoint operator with $A\sigma=\sigma A$ for
every $\sigma\in\mathcal S$, and let $\Sigma\subset\mathbf R$ be its spectrum. Suppose that $x<y$ lie in
$\Sigma$, and choose $m$ with $x<m<y$. Put $f(t)=\max(t-m,0)$ and $k(t)=\max(m-t,0)$, and let
$F=f(A)$ and $K=k(A)$ be given by the continuous functional calculus. Then $FK=(fk)(A)=0$, as $fk=0$.
A continuous function of $A$ is $0$ only if the function vanishes on $\Sigma$; as $k(x)>0$, $K\ne0$,
and as $f(y)>0$, $F\ne0$. The operator $F$ commutes with every $\sigma\in\mathcal S$, since $A$ does, so
its kernel is a closed subspace invariant under $\mathcal S$. Choose $w$ with $Kw\ne0$. Then $Kw$ lies
in the kernel of $F$, so this kernel is $H$ and $F=0$, a contradiction. Hence $\Sigma\subset\{c\}$ for
some $c\in\mathbf R$. The functions $t$ and $c$ agree on $\Sigma$, so $A=c$.

\emph{The general case.} Let $\sigma\in\mathcal S$, and let $\sigma'\in\mathcal S$ be its adjoint.
Taking adjoints in $T\sigma'=\sigma'T$ gives $\sigma T^*=T^*\sigma$. So the self-adjoint operators
$\frac12(T+T^*)$ and $\frac1{2i}(T-T^*)$ commute with every member of $\mathcal S$, and by the first part
they are real scalars $a$ and $b$. Hence $T=a+ib$.
% Lean: hilbert_schur_of_adjoint_closed: exists_real_scalar_of_selfAdjoint_commute_irreducible (the key
% step with cfc of max (t - m) 0 and max (m - t) 0, eqOn_of_cfc_eq_cfc, Commute.cfc_real for the
% invariance of the kernel, then cfc_id, cfc_congr and cfc_const) applied to realPart T and imaginaryPart T,
% and realPart_add_I_smul_imaginaryPart.
\end{proof}

\begin{lemma}\label{asm-dependence-test-assembly-realization-equivariant-schur-slot-closable}
Let $H$ and $E$ be complex Hilbert spaces, and let $\mathcal R$ be a set of unitary operators on $H$
that contains the inverse of each of its members and leaves no closed subspace of $H$ other than $0$
and $H$ invariant. Let $D\subset H$ be a subspace with $\rho(D)\subset D$ for every $\rho\in\mathcal R$,
and let $A,A'\colon D\to E$ be linear maps such that for all $\rho\in\mathcal R$ and $w\in D$
\[
 \|A\rho w\|=\|Aw\|,\qquad\|A'\rho w\|=\|A'w\|,\qquad\langle A\rho w,A'\rho w\rangle=\langle Aw,A'w\rangle,
\]
and that are \emph{closable}: if $w_k\in D$, $w_k\to0$ and $Aw_k\to f$, then $f=0$, and the same for
$A'$. Then there is $c\in\mathbf C$ with $\langle Aw,A'w'\rangle=c\langle w,w'\rangle$ for all
$w,w'\in D$.
\end{lemma}
\leanref{Development/AbelianVarietyTheory/Proofs/IwasawaDecomposition/SchurInnerEqPairClosableNormInvariantStatement.lean}{GlobalTheory.T.schur_inner_eq_pair_closable_normInvariant}
\leanref{Development/AbelianVarietyTheory/Proofs/IwasawaDecomposition/SchurInnerEqPairClosableNormInvariant.lean}{GlobalTheory.ClosableSchur.inner_eq_pair}
\leanref{Development/AbelianVarietyTheory/Proofs/IwasawaDecomposition/InnerOfNormEq.lean}{GlobalTheory.ClosableSchur.norm_sq_eq}
\begin{proof}
\emph{One map.} Let $B\colon D\to E$ be linear, closable, with $\|B\rho w\|=\|Bw\|$ for all $\rho$ and
$w$. We show that $\|Bw\|^2=b\|w\|^2$ on $D$ for some $b\in\mathbf R$. Let $K$ be the closure in $H\oplus E$
of the graph $\{(w,Bw):w\in D\}$, and $J\colon K\to H$ the first projection, which is bounded. It is
injective: if $(0,f)\in K$, it is the limit of points $(w_k,Bw_k)$ with $w_k\to0$ and $Bw_k\to f$, so
$f=0$. For $\rho\in\mathcal R$ the map $(w,Bw)\mapsto(\rho w,B\rho w)$ preserves norms, so it extends to
an isometry $U_\rho$ of $K$, with $JU_\rho=\rho J$, and $U_\rho U_{\rho^{-1}}=1$. For $x\in H$ and
$\gamma\in K$,
\[
 \langle J^*\rho x,\gamma\rangle=\langle\rho x,\rho JU_{\rho^{-1}}\gamma\rangle
 =\langle J^*x,U_{\rho^{-1}}\gamma\rangle=\langle U_\rho J^*x,\gamma\rangle,
\]
so $J^*\rho=U_\rho J^*$, and $T=JJ^*$ is a bounded operator on $H$ with $T\rho=\rho T$ for every
$\rho\in\mathcal R$. The adjoint of each $\rho\in\mathcal R$ is $\rho^{-1}\in\mathcal R$, so by
Lemma~\ref{asm-dependence-test-assembly-realization-equivariant-schur-slot-closable-schur}, $T=c$ for a
scalar $c$. Then
$J(J^*J\gamma)=cJ\gamma$, so $J^*J=c$ by injectivity, and $\|J\gamma\|^2=c\|\gamma\|^2$. At
$\gamma=(w,Bw)$ this reads $\|w\|^2=c(\|w\|^2+\|Bw\|^2)$. If $c=0$, then $D=0$ and any $b$ works;
otherwise $b=(1-c)/c$ works.

\emph{Two maps.} Let $a$ and $a'$ be the numbers given for $A$ and $A'$. For $z\in\mathbf C$ the map
$B_z=zA+A'$ satisfies
$\|B_zw\|^2=|z|^2\|Aw\|^2+2\operatorname{Re}(\bar z\langle Aw,A'w\rangle)+\|A'w\|^2$, so it is
norm-invariant by the hypotheses. It is closable: if $w_k\to0$, then $\|Aw_k\|^2=a\|w_k\|^2\to0$ and
$\|A'w_k\|\to0$, so $B_zw_k\to0$. So $\|B_zw\|^2=b_z\|w\|^2$. With $z=1$ and $z=i$ this gives
$\operatorname{Re}\langle Aw,A'w\rangle=\tfrac12(b_1-a-a')\|w\|^2$ and
$\operatorname{Im}\langle Aw,A'w\rangle=\tfrac12(b_i-a-a')\|w\|^2$. So the sesquilinear form
$\langle Aw,A'w'\rangle-c\langle w,w'\rangle$, with $c=\tfrac12(b_1-a-a')+\tfrac i2(b_i-a-a')$, vanishes on
the diagonal of $D$, hence on $D\times D$ by polarization.
% Lean: ClosableSchur.inner_eq_pair: norm_sq_eq (the graph closure K in WithLp 2 (H x E), J = fst,
% injectivity from closability, and ClosableSchur.abstract: the isometries U i from
% exists_isometry_extension_of_dense_of_norm_eq, hJU, hUinv, hadj, and hilbert_schur_of_adjoint_closed
% for T = J J^*), then key for z A + A' at z = 1 and z = I (tendsto_zero_of_norm_sq for closability),
% and inner_eq_of_diag. Schur's lemma for a family closed under adjoints is
% hilbert_schur_of_adjoint_closed (Lemma -schur-slot-closable-schur), applied with hadjrho (rho i^* = rho i^-1).
\end{proof}

\begin{lemma}\label{asm-dependence-test-assembly-realization-equivariant-schur-slot-smoothing}
Let $J_0$, $G=U(J_0)$, $\sigma$, $H$, $H^\infty$ and the Sobolev norms be as in
Theorem~\ref{lit-casselman-wallach}, and fix a Haar measure $dg$ on $G$. For a smooth function $\varphi$
with compact support on $G$ and $w\in H$ put $\sigma(\varphi)w=\int_G\varphi(g)\sigma(g)w\,dg$. Then
$\sigma(\varphi)w\in H^\infty$. Let $\ell$ be a linear functional on $H^\infty$ that is continuous for a
Sobolev norm. Then:
\begin{enumerate}
\item there is $C$ with $|\ell(\sigma(\varphi)w)|\le C\|w\|$ for all $w\in H$;
\item for $w\in H^\infty$ the function $g\mapsto\ell(\sigma(g)w)$ is continuous, and
  $\ell(\sigma(\varphi)w)=\int_G\varphi(g)\,\ell(\sigma(g)w)\,dg$.
\end{enumerate}
\end{lemma}
\leanref{Development/AbelianVarietyTheory/Proofs/IwasawaDecomposition/SobolevContinuousFormGardingVectorStatement.lean}{GlobalTheory.T.sobolevContinuous_form_gardingVector}
\leanref{Development/AbelianVarietyTheory/Proofs/IwasawaDecomposition/SobolevContinuousFormGardingVector.lean}{GlobalTheory.SobGarding.sobCont_gard}
\begin{proof}
By Gårding \cite{Gar47}, $\sigma(\varphi)w$ is smooth, and for every word $X_1\cdots X_k$ in elements of
$\mathfrak g$ there is a smooth function $\psi$ with compact support, obtained from $\varphi$ by
differentiation, with $d\sigma(X_1)\cdots d\sigma(X_k)\sigma(\varphi)w=\sigma(\psi)w$ for all $w\in H$;
and $\|\sigma(\psi)w\|\le\|\psi\|_{L^1}\|w\|$. If $\ell$ is continuous for the Sobolev norm of a finite
set $S$ of words, then $\ell(w)=\sum_{l\in S}\langle a_l,d\sigma(l)w\rangle$ for some vectors $a_l\in H$,
by the Hahn-Banach and Riesz theorems applied to the map $w\mapsto(d\sigma(l)w)_{l\in S}$ into $H^S$.

(1) With $\psi_l$ the function for the word $l$,
$|\ell(\sigma(\varphi)w)|\le\sum_l\|a_l\|\,\|\psi_l\|_{L^1}\|w\|$.

(2) For smooth $w$ and a word $l$, $g\mapsto d\sigma(l)\sigma(g)w$ is continuous, since
$d\sigma(X)\sigma(g)=\sigma(g)\,d\sigma(\operatorname{Ad}(g^{-1})X)$ on $H^\infty$. So
$g\mapsto\ell(\sigma(g)w)$ is continuous, and $d\sigma(l)$ may be taken under the integral defining
$\sigma(\varphi)w$, which gives
$\langle a_l,d\sigma(l)\sigma(\varphi)w\rangle=\int_G\varphi(g)\langle a_l,d\sigma(l)\sigma(g)w\rangle\,dg$.
Summing over $l$ gives the formula.
% Lean: SobGarding.sobCont_gard: T.sobolevContinuousForm_sum_lieWord_inner (the sum of inner products),
% word_gard and norm_gard_le for (1), T.smooth_vector_word_action_continuous and inner_word_gard for (2);
% gard_mem for the smoothness of the Garding vector. The Lean's test functions (IsTest) are smooth with
% compact support on the matrices, evaluated on G.
\end{proof}

\begin{lemma}\label{asm-dependence-test-assembly-realization-equivariant-schur-slot}
Let $J_0$ be a diagonal $n\times n$ matrix with entries $\pm1$,
$U(J_0)=\{g\in GL_n(\mathbf C):g^*J_0g=J_0\}$, and $\sigma$ an irreducible unitary representation of
$U(J_0)$ on a Hilbert space $H$, with space of smooth vectors $H^\infty$ and the Sobolev norms of
Theorem~\ref{lit-casselman-wallach}. Let $X$ be a second countable topological space with a finite
Borel measure $\nu$ and a continuous action $(x,g)\mapsto x\cdot g$ of $U(J_0)$, and let $a$ and $a'$
be linear maps from $H$ to functions on $X$ with the following properties.
\begin{enumerate}
\item For every $w\in H^\infty$, $a(w)$ and $a'(w)$ are continuous and bounded, and
  $a(\sigma(g)w)(x)=a(w)(x\cdot g)$ and $a'(\sigma(g)w)(x)=a'(w)(x\cdot g)$ for all $g$ and $x$.
\item For every $x\in X$ the functionals $w\mapsto a(w)(x)$ and $w\mapsto a'(w)(x)$ on $H^\infty$ are
  continuous for a Sobolev norm.
\item For $b,b'\in\{a,a'\}$ the pairing $\int_X\overline{b(u)}\,b'(w)\,d\nu$ of $u,w\in H^\infty$ does
  not change when $u$ and $w$ are replaced by $\sigma(g)u$ and $\sigma(g)w$, for every $g\in U(J_0)$.
\end{enumerate}
Then there is $c\in\mathbf C$ with $\int_X\overline{a(w)}\,a'(w')\,d\nu=c\langle w,w'\rangle$ for all
$w,w'\in H^\infty$.
\end{lemma}
\leanref{Development/AbelianVarietyTheory/Proofs/IwasawaDecomposition/SchurLemmaEquivariantPairBoundedFunctionsStatement.lean}{GlobalTheory.T.schurLemma_equivariant_pair_boundedFunctions_alt}
\leanref{Development/AbelianVarietyTheory/Proofs/IwasawaDecomposition/SchurLemmaEquivariantPairBoundedFunctions.lean}{GlobalTheory.SmVec.schurLemma_equivariant_pair_boundedFunctions}
\begin{proof}
The maps $a$ and $a'$ carry smooth vectors to bounded continuous functions, so they define linear
maps $A$ and $A'$ from $H^\infty$ to $L^2(X,\nu)$. By (3),
$\langle A(\sigma(g)u),A(\sigma(g)w)\rangle=\langle A(u),A(w)\rangle$ for all $g$ and all smooth $u,w$,
and likewise for $A'$ and for the mixed pairing $\langle A(u),A'(w)\rangle$. In particular
$\|A(\sigma(g)w)\|=\|A(w)\|$ and $\|A'(\sigma(g)w)\|=\|A'(w)\|$. The space $H^\infty$ is invariant under
$\sigma$, and the operators $\sigma(g)$ are unitary, contain their inverses and leave no closed
subspace other than $0$ and $H$ invariant. So by
Lemma~\ref{asm-dependence-test-assembly-realization-equivariant-schur-slot-closable}, with $D=H^\infty$,
it suffices to show that $A$ and $A'$ are closable; we do this for $A$. Write $\sigma(\varphi)$ for the
operators of Lemma~\ref{asm-dependence-test-assembly-realization-equivariant-schur-slot-smoothing}, for
smooth functions $\varphi$ with compact support on $U(J_0)$.
\begin{enumerate}
\item[(a)] For smooth $u$ and $f\in L^2(X,\nu)$, the function $g\mapsto\langle f,A(\sigma(g)u)\rangle
  =\int_X\overline{f(x)}\,a(u)(x\cdot g)\,d\nu(x)$ is continuous, by dominated convergence: $a(u)$ is
  bounded and continuous, the action is continuous, and $f$ is integrable, as $\nu$ is finite.
\item[(b)] For each $\varphi$, the map $w\mapsto A(\sigma(\varphi)w)$ from $H$ to $L^2(X,\nu)$ is
  continuous. By the closed graph theorem it suffices that its graph is closed. Let $u_k\to w$ in $H$
  with $A(\sigma(\varphi)u_k)\to f$ in $L^2$. A subsequence converges to $f$ almost everywhere. For
  every $x$, the functional $w\mapsto a(w)(x)$ is continuous for a Sobolev norm, by (2), so by
  Lemma~\ref{asm-dependence-test-assembly-realization-equivariant-schur-slot-smoothing}(1),
  $a(\sigma(\varphi)u_k)(x)\to a(\sigma(\varphi)w)(x)$. So $f=A(\sigma(\varphi)w)$.
\item[(c)] For smooth $w$ and $f\in L^2(X,\nu)$,
  $\langle f,A(\sigma(\varphi)w)\rangle=\int\varphi(g)\langle f,A(\sigma(g)w)\rangle\,dg$. Indeed, by
  Lemma~\ref{asm-dependence-test-assembly-realization-equivariant-schur-slot-smoothing}(2) at each $x$,
  $a(\sigma(\varphi)w)(x)=\int\varphi(g)\,a(w)(x\cdot g)\,dg$, and Fubini's theorem applies, as
  $|\varphi(g)\overline{f(x)}a(w)(x\cdot g)|\le C|\varphi(g)|\,|f(x)|$ with $C$ a bound of $a(w)$.
\end{enumerate}
Now let $u_k\in H^\infty$ with $u_k\to0$ and $A(u_k)\to f$. Let $v$ be smooth, and put
$F(g)=\langle A(\sigma(g^{-1})v),f\rangle$, which is continuous by (a). For each $\varphi$, by (b),
$\langle A(v),A(\sigma(\varphi)u_k)\rangle\to0$. By (c) and the invariance of the pairing,
\[
 \langle A(v),A(\sigma(\varphi)u_k)\rangle=\int\varphi(g)\langle A(v),A(\sigma(g)u_k)\rangle\,dg
 =\int\varphi(g)\langle A(\sigma(g^{-1})v),A(u_k)\rangle\,dg,
\]
which differs from $\int\varphi(g)F(g)\,dg$ by at most $\|\varphi\|_{L^1}\|A(v)\|\,\|A(u_k)-f\|$, since
$\|A(\sigma(g^{-1})v)\|=\|A(v)\|$. So $\int\varphi F\,dg=0$ for every $\varphi$. As $F$ is continuous, $F=0$,
and $\langle A(v),f\rangle=F(1)=0$. Hence $\|f\|^2=\lim_k\langle A(u_k),f\rangle=0$, and $A$ is closable.
% Lean: SmVec.schurLemma_equivariant_pair_boundedFunctions: exists_l2_lift for A and A' (l2_inner_of_ae),
% the invariance hAinv, hA'inv, hxinv, closable_of_smoothing with continuous_inner_l2_translate,
% T.smooth_vectors_continuous_l2 and T.smooth_vectors_inner_product, and
% T.schur_inner_eq_pair_closable_normInvariant. The Lean states it for the representation transported
% to U(G) through P with P^* J P = G, whose smooth vectors and Sobolev norms are those of sigma moved
% by the same conjugation. About 5,200 lines of Lean proofs at depth 5 by the count used here
% (closable_of_smoothing about 1,800, schur_inner_eq_pair_closable_normInvariant about 2,500).
% The proof follows them: (a) is continuous_inner_l2_translate, (b) T.smooth_vectors_continuous_l2 =
% SmVec.continuous_toL2_gard (continuous_of_seq_closed_graph, exists_seq_tendsto_ae), (c)
% T.smooth_vectors_inner_product = SmVec.inner_l2_gard (integral_integral_swap), and the closability is
% closable_of_smoothing (eq_zero_of_forall_test for F); Lemma -schur-slot-closable is
% T.schur_inner_eq_pair_closable_normInvariant, and the smoothing facts are Lemma -schur-slot-smoothing.
\end{proof}

\begin{lemma}\label{asm-dependence-test-assembly-realization-equivariant-schur}
Assume rapid decay of cusp forms (Theorem~\ref{lit-bundle-theta}(3)). Let $D$ be a definite-side
datum, $V$ a Hermitian space of $D$, $\theta$ an archimedean family and $M$ a model of discrete
series for $U(V)$ with spaces $H_\tau$ (Definition~\ref{asm-dependence-rows-models}). Let
$J\subset U(V)(\mathbf A)$, let $\chi$ be a function on $U(V)(\mathbf A)$, let $\nu$ be a Haar measure on
$U(V)(\mathbf A)$ and $\mathcal F$ a fundamental domain for $U(V)(F)$ with respect to $\nu$.
\begin{enumerate}
\item For $T,T'\in\operatorname{Mult}_\theta(M;J,\chi)$ there is $c\in\mathbf C$ such that for all
  $v,v'\in\prod_\tau H_\tau^\infty$ the function $\overline{T(v)}\,T'(v')$ is integrable on $\mathcal
  F$ and
  \[
   \int_{\mathcal F}\overline{T(v)}\,T'(v')\,d\nu=c\prod_\tau\langle v_\tau,v'_\tau\rangle.
  \]
\item Let $u\in\operatorname{Mult}_\theta(M;J,\chi)$, read as a linear map on $M^\infty$. If $u(x)=0$
  for some nonzero $x\in M^\infty$, then $u=0$.
\end{enumerate}
\end{lemma}
\leanref{Development/AbelianVarietyTheory/Proofs/IwasawaDecomposition/HcFiniteAndSchurMult.lean}{GlobalTheory.MultFinite.schurMult}
\leanref{Development/AbelianVarietyTheory/Proofs/IwasawaDecomposition/HcFiniteAndSchurMult.lean}{GlobalTheory.MultFinite.schur_multMap}
\leanref{Formalism/LocalToGlobal/IndefiniteReal/ScalarOnPartsFull.lean}{Formalism.IndefR.Mult.SchurMult}
\leanref{Development/AbelianVarietyTheory/Proofs/IwasawaDecomposition/ExistsEqConstMulProdInnerOfSlotwise.lean}{GlobalTheory.exists_eq_const_mul_prod_inner_of_positionwise}
\leanref{Development/AbelianVarietyTheory/Proofs/TwistedOperators/MultiplicitySpaceEqZeroOfApplyNonzeroTensorZero.lean}{GlobalTheory.multT_eq_zero_of_apply_eq_zero}
\leanref{Development/AbelianVarietyTheory/Proofs/TwistedOperators/MultiplicitySpaceEqZeroOfApplyNonzeroTensorZero.lean}{GlobalTheory.eq_zero_of_integral_fundDomain}
\begin{proof}
(1) The multiplicity maps of level $(J,\chi)$ form a subspace, so $T$ and $T'$ are multiplicity
maps. By Theorem~\ref{asm-theta-agree-definite-volume}, $\nu(\mathcal F)$ is finite and $\nu$ is
invariant under right translations. Every value $T(v)$ is a cuspidal automorphic form, so it is
continuous and, by rapid decay, bounded; the same holds for $T'$. Hence every function
$\overline{T(v)}\,T'(v')$ is integrable on $\mathcal F$. For a function $f$ invariant on the left under
$U(V)(F)$ and integrable on $\mathcal F$, and any $y\in U(V)(\mathbf A)$, the integral of $h\mapsto
f(hy)$ over $\mathcal F$ is that of $f$, by the argument of the second paragraph of the proof of
Lemma~\ref{asm-dependence-test-assembly-realization-equivariant-integral}, right translation by $y$
preserving $\nu$.

Fix $\tau$ and $v,v'$, and for $w\in H_\tau^\infty$ write $v[\tau\leftarrow w]$ for $v$ with its
component at $\tau$ replaced by $w$. Extend $w\mapsto T(v[\tau\leftarrow w])$ and $w\mapsto
T'(v'[\tau\leftarrow w])$ to linear maps $a$ and $a'$ from $H_\tau$ to functions on
$X=U(V)(\mathbf A)$, with the finite measure $\nu$ restricted to $\mathcal F$. With the model's
representation $\sigma_\tau$ of $U(J_\tau)$ and frame $f_\tau$, the group $U(J_\tau)$ acts on
$U(V)(\mathbf A)$ by right translation by the image of $f_\tau^{-1}gf_\tau$ under $\theta$, and the
equivariance of multiplicity maps gives $a(\sigma_\tau(g)w)(x)=a(w)(x\cdot g)$ for smooth $w$. The
values are bounded and continuous, the functionals at points are continuous for a Sobolev norm, and
by the invariance of the integrals under right translation the three pairings of
Lemma~\ref{asm-dependence-test-assembly-realization-equivariant-schur-slot}(3) are invariant. That
lemma gives $c_\tau(v,v')$ with
$\int_{\mathcal F}\overline{T(v[\tau\leftarrow w])}\,T'(v'[\tau\leftarrow w'])\,d\nu=c_\tau(v,v')\langle
w,w'\rangle$ for all smooth $w,w'$.

Put $B(v,v')=\int_{\mathcal F}\overline{T(v)}\,T'(v')\,d\nu$. If some $H_\tau^\infty$ is $0$, the last
identity with $w=v_\tau$ and $w'=v'_\tau$ gives $B=0$, and $c=0$ works. Otherwise choose unit vectors
$e_\tau\in H_\tau^\infty$. For every $\tau$, the identity with $w=v_\tau$, $w'=v'_\tau$ and with
$w=w'=e_\tau$ gives $B(v,v')=B(v[\tau\leftarrow e_\tau],v'[\tau\leftarrow e_\tau])\langle
v_\tau,v'_\tau\rangle$. Applying this at every $\tau$ in turn gives
$B(v,v')=B(e,e)\prod_\tau\langle v_\tau,v'_\tau\rangle$, and $c=B(e,e)$ works.

(2) Let $T$ be the multiplicity map with $T(v)=u(\otimes_\tau v_\tau)$, and let $c$ be the constant
of (1) for $T=T'$. By linearity in the second variable, $\int_{\mathcal F}\overline{T(w)}\,u(y)\,d\nu=
c\,\langle\otimes_\tau w_\tau,y\rangle$ for all $w$ and all $y\in M^\infty$, with the inner product of
the tensor product. At $y=x$ the left side vanishes. If $c\ne0$, then $x$ is orthogonal to every
pure tensor, so $x=0$, which is not the case. So $c=0$, and $\int_{\mathcal F}|T(v)|^2\,d\nu=0$ for
every $v$. The function $T(v)$ is continuous and invariant on the left under $U(V)(F)$; the open set
where it does not vanish is invariant under $U(V)(F)$ and meets $\mathcal F$ in a null set, so it is
null, and hence empty, as $\nu$ charges every nonempty open set. So $T=0$, and $u=0$.
% Lean: (1) MultFinite.schurMult: schur_multMap at each tau (exists_continuousLinearMap_tens_update for a
% and a', CuspDecay.bounded_of_cuspForm under CuspRapidDecay, Mem.VolFin.finVol for the finite measure,
% AdelicPointsUnimod.isMulRightInvariant_of_isHaarMeasure and fundamentalDomain_integral_mul_right for
% the invariance, and T.schurLemma_equivariant_pair_boundedFunctions_alt, Lemma -schur-slot), then
% T.exists_eq_const_mul_prod_inner_of_slotwise' = exists_eq_const_mul_prod_inner_of_positionwise, with
% CuspDecay.integrableOn_star_mul for the integrability. (2) multT_eq_zero_of_apply_eq_zero (hall by
% induction on tensors, eq_zero_of_tprodInnerFun, eq_zero_of_integral_fundDomain). The Lean's finVol and
% the right invariance are proved in the archive; the book cites Theorem asm-theta-agree-definite-volume.
% About 8,800 lines of Lean proofs at depth 5 for schurMult by the count used here, about 5,200 of them
% in Lemma -schur-slot.
\end{proof}

\begin{lemma}\label{asm-dependence-test-assembly-realization-equivariant}
Assume rapid decay of cusp forms (Theorem~\ref{lit-bundle-theta}(3)). Let $R$ be a realization of a
definite-side datum $D$ for a family of definite forms for a finite group $G$
(Definition~\ref{is-realisation}), with coefficient field $L$, the embedding $\sigma_0$ of its base
place, the top space $H_E$ with its representation of $G$, the theta lifts $\mathrm{Th}_t$ and the
$\sigma_0$-semilinear maps $j_t$ from the row forms. Suppose that $R$ satisfies the core conditions
at a face (Definition~\ref{is-eigenvalue-lift-data}). Let $e=\sum_{g\in G}e(g)\,g$, $e(g)\in L$, be the
idempotent of the family, and let $G$ act on the row forms through the family.
\begin{enumerate}
\item For every tuple $t$, every $g\in G$ and every row form $y$ at $t$,
  $\mathrm{Th}_tj_t(gy)=g\,\mathrm{Th}_tj_t(y)$.
\item For every $u\in H_E$, $\sum_{g\in G}\sigma_0(e(g))\,g\,u=u$.
\item For every tuple $t$ and every row form $y$ at $t$, the form $ey$ lies in the good part of the
  row forms at $t$, and $\mathrm{Th}_tj_t(ey)=\mathrm{Th}_tj_t(y)$.
\end{enumerate}
\end{lemma}
\leanref{Development/AbelianVarietyTheory/Proofs/TwistedOperators/ThThGAct.lean}{GlobalTheory.LiftGroupEquiv.th_thG_act}
\leanref{Development/AbelianVarietyTheory/Proofs/TwistedOperators/ActEqTranslate.lean}{GlobalTheory.FormAction.act_eq_translate}
\leanref{Development/AbelianVarietyTheory/Proofs/TwistedOperators/ChainCompositeLiftIntertwinesFiniteGroupActionsStatement.lean}{GlobalTheory.T.chainCompositeLift_intertwines_finiteGroupActions}
\leanref{Development/AbelianVarietyTheory/Proofs/TwistedOperators/ChainCompositeLiftIntertwinesFiniteGroupActions.lean}{GlobalTheory.ChainLiftEquiv.th_intertwines}
\leanref{Development/AbelianVarietyTheory/Proofs/TwistedOperators/ChainCompositeLiftIntertwinesFiniteGroupActions.lean}{GlobalTheory.ChainLiftEquiv.step}
\leanref{Development/AbelianVarietyTheory/Proofs/TwistedOperators/ThetaLiftTranslateOfKernelFixed.lean}{GlobalTheory.thetaLift_translate_of_kernel_fixed}
\leanref{Development/AbelianVarietyTheory/Proofs/TwistedOperators/MultiplicitySpaceEqZeroOfApplyNonzeroTensorZero.lean}{GlobalTheory.multT_eq_zero_of_apply_eq_zero}
\leanref{Development/AbelianVarietyTheory/Proofs/IwasawaDecomposition/HcFiniteAndSchurMult.lean}{GlobalTheory.MultFinite.schurMult}
\leanref{Development/AbelianVarietyTheory/Proofs/TwistedOperators/RowTestHilbertModelsOfChoiceClaims.lean}{GlobalTheory.RowTestFromClaims.sum_gE}
\leanref{Development/AbelianVarietyTheory/Proofs/TwistedOperators/RowTestHilbertModelsOfChoiceClaims.lean}{GlobalTheory.RowTestFromClaims.exists_good}
\begin{proof}
(1) The family presents the row forms at $t$ as functions on the finite adelic points of $U(X')$,
and $G$ acts on them through a surjection $\mathrm{red}\colon Q\to G$ from an open compact subgroup
$Q$: $g$ acts by right translation by any $\kappa\in Q$ with $\mathrm{red}(\kappa)=g$, multiplied by a
character of $Q$. The core conditions on the group identify $\mathrm{red}$ with the isomorphism of
$G$ that $R$ records, and those on the models carry the presentation to $j_t$. So there is a level
lift $\kappa$ of $g$ for $X'$ at $t$
(Definition~\ref{asm-dependence-test-assembly-realization-equivariant-levels}) such that
$j_t(gy)(h)=\chi(\kappa)^{-1}j_t(y)(h\kappa)$ for every $h$, where $\chi$ is the level character.

The theta lift $\mathrm{Th}_t$ is a chain: writing $Y_m=X'$, the map $j_t(y)$ gives a multiplicity map
$F_m$ at $X'$, the theta lift $\iota_j$ of the pair $(Y_j,Y_{j+1})$ at $t$ carries multiplicity maps
at $Y_{j+1}$ to multiplicity maps at $Y_j$ for the level of $Y_j$ at $t$, and $\mathrm{Th}_tj_t(y)$
corresponds to $F_0=\iota_0\cdots\iota_{m-1}(F_m)$. Put $\kappa_m=\kappa$; for $j<m$ the element $g$ has
a level lift $\kappa_j$ for $Y_j$ at $t$, by
Lemma~\ref{asm-dependence-test-assembly-realization-equivariant-lifts}(2). Write $\chi_j$ for the
level character of $Y_j$ at $t$ and $a_j=\chi_j(\kappa_j)^{-1}$, and call $F'$ the \emph{$\kappa_j$-shift} of a multiplicity map $F$ at
$Y_j$ if $F'(v)(h)=a_jF(v)(h\kappa_j)$ for all vectors $v$ and all $h$. So $F_m(gy)$ is the
$\kappa_m$-shift of $F_m(y)$.

\emph{Step.} Let $j<m$, let $F$ be a multiplicity map at $Y_{j+1}$ and $F'$ its $\kappa_{j+1}$-shift;
then $\iota_j(F')$ is the $\kappa_j$-shift of $\iota_j(F)$. Let $u$ be the $\kappa_j$-shift of
$\iota_j(F)$, which is again a multiplicity map, by
Lemma~\ref{asm-dependence-test-assembly-realization-equivariant-shift}(2). For a vector $\varphi$ of
the chain at the pair and
a vector $v$, with $I(\varphi,v)$ the vector of the model of $U(Y_j)$ they give,
$\iota_j(F)(I(\varphi,v))(g')=\int_{\mathcal F}\theta(g',h;\varphi)F(v)(h)\,dh$ over a fundamental domain
$\mathcal F$ for the rational points of $U(Y_{j+1})$. The Weil datum of the pair maps $\varphi$ at
$(\kappa_j,\kappa_{j+1})$ to $a_j^{-1}a_{j+1}\varphi$
(Lemma~\ref{asm-dependence-test-assembly-realization-equivariant-scalars}), so
$a_j\theta(g'\kappa_j,h\kappa_{j+1};\varphi)=a_{j+1}\theta(g',h;\varphi)$; and the integral over
$\mathcal F$ of an integrable function that is left invariant under the rational points is unchanged
by the right translation $h\mapsto h\kappa_{j+1}$
(Lemma~\ref{asm-dependence-test-assembly-realization-equivariant-integral}). Hence
\[
 u(I(\varphi,v))(g')=a_j\int_{\mathcal F}\theta(g'\kappa_j,h\kappa_{j+1};\varphi)F(v)(h\kappa_{j+1})\,dh
 =\int_{\mathcal F}\theta(g',h;\varphi)\,a_{j+1}F(v)(h\kappa_{j+1})\,dh=\iota_j(F')(I(\varphi,v))(g').
\]
So $\iota_j(F')-u$ vanishes at every $I(\varphi,v)$. Some $I(\varphi_0,v_0)$ is nonzero, and by Schur's lemma
for the multiplicity maps at $Y_j$, Lemma~\ref{asm-dependence-test-assembly-realization-equivariant-schur}(2),
a multiplicity map that vanishes at a nonzero vector is zero. So $\iota_j(F')=u$.

By the step and descending induction from $j=m$, the multiplicity map $F_0(gy)$ is the
$\kappa_0$-shift of $F_0(y)$. At $Y_0$ the lift $\kappa_0$ of $g$ also lies in the level of the trivial
tuple with the same value of the level character, and there the $\kappa_0$-shift of a function of the
type of that level is the action of $g$
(Lemma~\ref{asm-dependence-test-assembly-realization-equivariant-shift}(1)). The identification of
$H_E$ with the multiplicity maps at $Y_0$ carries the action of $g$ on $H_E$ to this action, and the
projection to the good part commutes with it
(Lemma~\ref{asm-dependence-test-assembly-realization-equivariant-projection}). So
$\mathrm{Th}_tj_t(gy)=g\,\mathrm{Th}_tj_t(y)$.
% Lean: ChainLiftEquiv.th_intertwines: liftsExist_chain (Lemma -lifts(2)), IsShift, thDn_shift by
% dnFrom_rel from stepX (at X', toMap_shift) and stepY, both through the generic step: exists_translate
% (Lemma -shift(2)), thetaKO_rat_left_and_eigen_right, om_pureTensorO_of_finite with
% T.stepVector_lift_pair_flagLevel_realized_chain or T.lastVector_lift_pair_flagLevel_characters (Lemma
% -scalars, whose case P = 1 is open), ThetaProjEquiv.integral_shift (Lemma -integral),
% thetaLift_translate_of_kernel_fixed (the computation above, written out) and
% multT_eq_zero_of_apply_eq_zero with SchurMult (MultFinite.schurMult, Lemma -schur, whose slotwise
% Schur lemma -schur-slot is written out). Then lift_zero, actL_zero (level_action_spec, Lemma -shift(1)) and
% th_of_thDn (models_HE, models_gE, chain_Th, ProjComm.proj_comm, Lemma -projection).
% Lean: FormAction.act_eq_translate (act_tie.red_surj, act_eq, exists_lift; about 40 lines) gives
% kappa; T.chainCompositeLift_intertwines_finiteGroupActions = ChainLiftEquiv.th_intertwines (about
% 2,400 lines at depth 5) with SchurMult from MultFinite.schurMult (about 5,600). In all about 7,200
% lines of Lean proofs at depth 5. The Lean's th_thG_act also proves the same for the lift forms and
% the theta maps th_G (act_eq_translate_lift_V, thetaProjection_intertwines_levelLifts).

(2) By the core conditions on the models, an injective map identifies $H_E$ with the good part at
the trivial level of the multiplicity maps of the model of $U(Y_0)$, and carries the action of $g$
on $H_E$ to the action of $g$ on these maps through the level data. The good part consists of maps
fixed by the operator of $\sum_g\sigma_0(e(g))\,g$. So the image of $\sum_g\sigma_0(e(g))\,g\,u$ is
the image of $u$, and injectivity gives (2).

(3) The good part at the trivial subgroup and the trivial character is the image of $e$, so $ey$
lies in it. Since $j_t$ is $\sigma_0$-semilinear and $\mathrm{Th}_t$ is linear, (1) and (2) give
\[
 \mathrm{Th}_tj_t(ey)=\sum_g\sigma_0(e(g))\,\mathrm{Th}_tj_t(gy)
 =\sum_g\sigma_0(e(g))\,g\,\mathrm{Th}_tj_t(y)=\mathrm{Th}_tj_t(y).
\]
% Lean: sum_gE (hE : models_HE, hg : models_gE, mem_goodPart) and exists_good (goodAt_bot,
% th_thG_act).
\end{proof}

\begin{lemma}\label{asm-dependence-test-assembly-realization}
Assume rapid decay of cusp forms (Theorem~\ref{lit-bundle-theta}(3)) and the seven statements on test
choices (Definition~\ref{asm-dependence-test-assembly-statements}). There is $C_{F_1}$ such that for
every $\varepsilon$ and every $C_F\ge C_{F_1}$ there is $x_0$ with the following property. Let $D$ be
a definite-side datum with size data of exponent $\varepsilon$ at a scale $x\ge x_0$, and let $R$ be a
realization of $D$ for a family of definite forms which with these size data is a realization in an
automorphic setting at a face $F_R$, and which has test choices for $C_F$ that satisfy the
evaluation conditions and the field conditions. Then the complex data of $R$ at $F_R$ pass the row
test (Definition~\ref{asm-stage-construction}(3)).
\end{lemma}
\leanref{Development/AbelianVarietyTheory/Proofs/TwistedOperators/RowTestAtRealisationOfHilbertModelsStatement.lean}{GlobalTheory.T.rowTest_at_realisation_of_hilbertModels}
\leanref{Development/AbelianVarietyTheory/Proofs/TwistedOperators/RowTestAtRealisationOfHilbertModels.lean}{GlobalTheory.RowTest.rowTest_realisation}
\leanref{Development/AbelianVarietyTheory/Proofs/TwistedOperators/RowTestHilbertModelsOfChoiceClaims.lean}{GlobalTheory.RowTestFromClaims.rowTestModel_of_claims}
\leanref{Development/AbelianVarietyTheory/Proofs/TwistedOperators/RowTestHilbertModelsOfChoiceClaims.lean}{GlobalTheory.RowTestFromClaims.rowTest_of_choice}
\leanref{Development/AbelianVarietyTheory/BlockTheory/TwistedDiscreteSeries/RationalRowTestKernel.lean}{GlobalTheory.IndefiniteRowSideFullHecke.RationalRowTestKernel.test_tw_iff}
\leanref{Development/AbelianVarietyTheory/Proofs/TwistedOperators/RowTestComplexSideOfHilbertModels.lean}{GlobalTheory.Twist.rowTest_complexSide_of_hilbert_models.rowTest_toCx3}
\begin{proof}
Let $C_{F_1}$ be the largest of the constants $C_{F,0}$ of statements (a), (b), (c), (e), (f) and
(g), and for $\varepsilon$ and $C_F\ge C_{F_1}$ let $x_0$ be the largest of the numbers $x_0$ that
they give. Let $D$, $R$ and $F_R$ be as in the statement, with $x\ge x_0$, and fix test choices for
$C_F$ that satisfy the evaluation conditions and the field conditions. The six statements apply to
them. By condition (5) of an automorphic setting, $R$ satisfies the core conditions at $F_R$. Write
$\sigma$ for the embedding of the coefficient field $L$ of $D$, which is $\sigma_0$ for $R$.

\emph{Step 1: the functionals.} By (a), the test functional $\psi_a(u)$ converges for every $a$ and
every $u\in E$. Where they converge, the iterated theta lift and the Bessel period are linear in $u$,
so each $\psi_a$ is a linear functional on $E$. By (b), the $\psi_a$ separate the points of $E$.

\emph{Step 2: the values.} By (c), $c(i,j,t)\ne0$ for all $i$, $j$ and $t$, so $c_{\mathrm{tot}}(t)\ne0$.
By (e) there are $\kappa\in\mathbf C$ and roots of unity $\zeta_t\in L$ with
$\sigma(\zeta_t)\,c_{\mathrm{tot}}(t)\,\psi_a(\mathrm{Th}_tj_ty')=\kappa\,v(t,y',a)$ for every tuple
$t$, every row form $y'$ at $t$ in the good part and every $a$; and $\sigma(\zeta_t)\ne0$. Let $t$,
a row form $y$ at $t$ and $a$ be given. By
Lemma~\ref{asm-dependence-test-assembly-realization-equivariant}(3) there is $y'$ in the good part with
$\mathrm{Th}_tj_ty'=\mathrm{Th}_tj_ty$, so
\[
 \sigma(\zeta_t)\,c_{\mathrm{tot}}(t)\,\psi_a(\mathrm{Th}_tj_ty)=\kappa\,v(t,y',a).
\]
If $\kappa=0$, then $\psi_a(\mathrm{Th}_tj_ty)=0$ for all $t$, $y$ and $a$, and the functionals
$\psi_a$ have values $0\in\sigma(L)$ at these vectors. If $\kappa\ne0$, put
$k=c_{\mathrm{tot}}(\mu_0)\kappa^{-1}\ne0$; the functionals $k\psi_a$ still separate the points of $E$,
and
\[
 k\,\psi_a(\mathrm{Th}_tj_ty)=v(t,y',a)\cdot\sigma(\zeta_t)^{-1}\cdot
 \Bigl(\frac{c_{\mathrm{tot}}(t)}{c_{\mathrm{tot}}(\mu_0)}\Bigr)^{-1}.
\]
By (f), $v(t,y',a)\in\sigma(L)$. By (g), $c_{\mathrm{tot}}(t)/c_{\mathrm{tot}}(\mu_0)=
\prod_{i,j}c(i,j,t)/c(i,j,\mu_0)$ lies in $\sigma(L)$, and it is nonzero. So
$k\psi_a(\mathrm{Th}_tj_ty)\in\sigma(L)$. In both cases there are finitely many linear functionals on
$E$ that separate its points and take values in $\sigma(L)$ at every $\mathrm{Th}_tj_ty$.

\emph{Step 3: the complex data.} The complex data of $R$ at $F_R$ have the top space $E$, the theta
lifts $\Theta_t=\mathrm{Th}_t$ and the maps $j_t$ (Definition~\ref{is-realisation}). Their embedding
at the base place is $\sigma$ or its complex conjugate $\bar\sigma$. In the second case $E$ carries the
conjugate complex structure; a linear functional $\psi$ for the old structure gives the linear
functional $\overline{\psi}$ for the new one, these still separate points, and
$\overline{\psi(v)}\in\bar\sigma(L)$ whenever $\psi(v)\in\sigma(L)$. So the functionals of Step 2 give
the row test. Statement (d) is not used.
% Lean: RowTest.rowTest_realisation = RowTestFromClaims.rowTestModel_of_claims (Big.and of LinK,
% InjK, ConstK, Formula2K, SumRatK, RatioK; SumOKK is the unused hypothesis) and
% rowTest_complexSide_of_hilbert_models (test_tw_iff with cj s). rowTest_of_choice:
% exists_linearMap_eq_psiK (fPsiPhi_add, fPsiPhi_smul), the case kappa = 0, and
% rowTestModel_of_scaled_values_rational with k = cTot mu0 / kappa and rat_of_formula;
% exists_embLL_mul_psiK_eq_vK_of_formulaAt uses exists_good (Lemma
% asm-dependence-test-assembly-realization-equivariant(3)). About 820 lines of Lean proofs at depth 5
% outside th_thG_act.
\end{proof}

\begin{lemma}\label{asm-dependence-test-assembly-existence}
Assume the number field statement (Definition~\ref{asm-dependence-test-assembly-statements}(4)).
There is $C_{F_2}$ such that for every $C_F\ge C_{F_2}$ the following holds eventually in the
context and at large base data (Definition~\ref{asm-stage-construction}(1) and (2)). Let $w$ be a
column choice over the base data, framing and twist data, $Z$ a later stage, $z$ a realization of
$D(w,Z)$ built at the stage, and size data for $D(w,Z)$ at the scale of the context. If $z$ with
these size data is a realization in an automorphic setting at its face, then it has test choices for
$C_F$ that satisfy the evaluation conditions and the field conditions.
\end{lemma}
\leanref{Development/AbelianVarietyTheory/Proofs/TwistedOperators/TestChoicesExistLargeCoeffFieldAdelicDataStatement.lean}{GlobalTheory.T.testChoices_exist_largeCoeffField_adelicData}
\leanref{Development/AbelianVarietyTheory/Proofs/TwistedOperators/TestChoicesExistLargeCoeffFieldAdelicData.lean}{GlobalTheory.RowEx.dat_choice_exists}
\leanref{Development/AbelianVarietyTheory/Proofs/TwistedOperators/TwistedDataOffBaseLevel.lean}{GlobalTheory.RowEx.twisted_data_off_base_level}
\leanref{Development/AbelianVarietyTheory/Proofs/TwistedOperators/TwistedDataOffBaseLevel.lean}{GlobalTheory.RowEx.DatTw.unrAt_of_mem_set}
\begin{proof}
Let $C_{F_2}$ be the constant $C_{F,0}$ of the number field statement, and let $C_F\ge C_{F_2}$. For
an exponent $\varepsilon$ let $x_0$ be the number that the number field statement gives for
$\varepsilon$ and $C_F$. It suffices to show that the property holds at large base data in every
context with exponent $\varepsilon$ and scale $x\ge x_0$, whatever its constants.

Let $K$ be such a context. For base data $B$ with a framing $f$ and twist data $T$, let $\Xi$ be the
set of characters that $f$ and $T$ give (Definition~\ref{asm-stage-construction}(5)): the trivial
character, the characters $\delta$ and $\delta'$ of $f$, and the characters $\delta_\ell\psi_w$. For
the condition at large base data take no primes, and as the algebraic numbers the finite set $S$
that the number field statement gives for the scale of $K$, $B$, $f$ and $\Xi$ if $\Xi$ is finite,
and the empty set otherwise.

Let $B$, $f$ and $T$ lie over the context, with coefficient field containing these numbers, and let
$w$, $Z$, $z$ and the size data be as in the statement. The characters $\xi_t$ of the level system
of $z$ are characters of the column choice $w$, and the characters of $w$ form the set $\Xi$, the
image of a finite set. So $\Xi$ is finite, the coefficient field contains $S$, and every $\xi_t$ lies
in $\Xi$. The columns of $z$ are given by $f$. The twisted places of $z$ lie off $q_0$: in a level
system with base tuple, $\xi_t$ is ramified at every place of $Q_t$, while every character in $\Xi$
is unramified at the places above $q_0$. Indeed $\delta$ is unramified away from the places above
$p_{de}$, and $\delta'$ away from those above $p'_{de}$, two primes different from $q_0$; and each
$\psi_w$ is a twisting character,
ramified only at places whose norm is a prime outside a set that contains $q_0$. So $z$, with its size
data at the scale $(x,\varepsilon)$, $x\ge x_0$, satisfies the hypotheses of the number field
statement for $B$, $f$ and $\Xi$, which gives test choices for $C_F$ that satisfy the evaluation
conditions and the field conditions.
% Lean: RowEx.dat_choice_exists, with epsilon_0 = 1 and every constant of the context free; the set
% S_2 of large2 is (P, f, Xi) |-> choose (hF P f Xi h) if Xi is finite and empty otherwise, applied
% at Xi = T.set f. twisted_data_off_base_level: z.lev.xiH with w.xi_range (the characters of w form
% T.set f), finite as a range; z.lev.C with w.C_eq; Q_ram of the level system with
% DatTw.unrAt_of_mem_set (f.isDel.unr and f.isDel'.unr off pde and pde', not_overP_pde;
% unrAt_of_isTwist, q0H lies in offSet). About 160 lines of Lean proofs at depth 5.
\end{proof}

\begin{lemma}\label{asm-dependence-test-assembly-sizes-explicit}
Let $0<\varepsilon<1/52$, let $C$ be standard constants of a context, and let $c_D\ge0$ and
$c_h\ge402$. There is $x_0$ such that the following holds for every scale $s=(x,\varepsilon,J)$ with
$x\ge x_0$, in the notation of Definition~\ref{asm-dependence-functional-witness-numbers-explicit}.
Let $D$ be a definite-side datum and let $a_t,b_t$ be integers for the top classes $t$ of its host.
Suppose that the rank datum of $D$ is adapted to $s$; that its host satisfies condition (2) of the
explicit values and $|\operatorname{Re}\frac{L'}{L}(1,\chi_{-d})|\le x^{\varepsilon}$; that its
Harish-Chandra parameters satisfy condition (3) with these $b_t$, and at the place of a top class $t$
have $u_1=a_t+\frac{k-1}2$ and $u_k=-b_t-\frac{k-1}2$; that $\Delta_S=\lceil x^{c_D}\rceil$; that the
constant $C_T$ of the host is that of $C$; that $\log\Gamma_P=C_\Gamma m^5fx^{\varepsilon}$ at every word
and $D$ has $\lceil3\Gamma_P^2x^{C_F}\rceil$ words; that $\eta_F=6x^{-C_F}$; that $a\le a_t\le
a+a/(kx^{13})$ and $b'\le b_t\le b'+2n\delta_1x^{\varepsilon}\Lambda b'/k$ for every top class, with
$\Lambda=K_\Lambda N_1^2\log(xb)$; and that $I_j(\mu_w,w)+I_j(\mu_w,\bar w)+j_w(\mu_w)=0$ for every lift
$1\le j\le m$ and every place $w$ with $\mu_w=\mu_0$. Then there are size data for $D$ at the scale $s$,
with exponents $a_t$ and $b_t$, that take the explicit values for $(C,c_D,c_h)$.
\end{lemma}
\leanref{Development/AbelianVarietyTheory/Proofs/GlobalInterpolation/ExistsExplicitParameters.lean}{GlobalTheory.EX.exists_sizeData_explicit}
\leanref{Development/AbelianVarietyTheory/Proofs/GlobalInterpolation/ExistsExplicitParameters.lean}{GlobalTheory.EX.bottom_window_le}
\leanref{Development/AbelianVarietyTheory/Proofs/GlobalInterpolation/ParametersWindowsOfExplicitParametersStatement.lean}{GlobalTheory.T.windows_of_explicitParameters}
\leanref{Development/AbelianVarietyTheory/Proofs/GlobalInterpolation/ExistsExplicitParameters.lean}{GlobalTheory.EX.middle_between}
\begin{proof}
Standard constants have $C_t=401$, $C_c=2$, $K_\Lambda\ge1$ and $C_\Gamma\ge1$. Put $\Delta=\lceil
x^{c_D}\rceil$, $h=\lceil x^{c_h}\rceil$, $\eta=2mh$, $S=4k\Delta$ and $\Lambda=K_\Lambda N_1^2\log(xb)$.
Two computations are used, each valid for $x$ beyond a threshold depending on $\varepsilon$ and the
constants.
\begin{enumerate}
\item[(A)] If the rank datum is adapted to the scale, then $2n\delta_1x^{\varepsilon}\Lambda b'/k\le b'$.
\item[(B)] Suppose that the rank datum is adapted to the scale, that the host satisfies condition (2)
  of the explicit values, that $b'\le b_t\le2b'$ for every top class $t$, and that the Harish-Chandra
  parameters satisfy condition (3) with these $b_t$. Then: $b_1(\tau)-b_t>\frac{k-1}2+3$ at the place
  of a top class $t$; $b_j(\tau)\le b^2$; $b\le e_j(\tau)\le b+\eta+1$ at a thick place,
  $b_t+b'/x\le e_j(\tau)\le b_t+b'/x+\eta$ at the place of a top class $t$, and $w''\le e_j(\tau)\le
  w''+\eta$ at a thin place; $S\le u_1(\tau)-v_1(\tau)\le2S$ at an embedding that is not the place of a
  top class; $0\le v_i(\tau)\le S$; $|b_i(\tau)-e_j(\tau)|\le\eta$; $1\le t\le x^{401}$;
  $10\lceil\log_5(nN(m)/m)\rceil+20\le L_w$; and $0<|T_2|\le n$.
\end{enumerate}
Let $x_0$ be the larger of the two thresholds and $\max(C_8,1)^{100/\varepsilon}$, so that $C_8\le
x^{\varepsilon/100}$ for $x\ge x_0$. Let $s$, $D$, $a_t$ and $b_t$ be as in the statement. By (A),
$b_t\le b'+2n\delta_1x^{\varepsilon}\Lambda b'/k\le2b'$, so (B) applies.

\emph{The size data.} Take the scale $s$, the exponents $a_t$ and $b_t$, the empty support, so that
all gains, local terms, errors, $\tau_w$ and $\iota_w$ vanish, no dropped primes, the zero submodules
$N^\circ_u$, $T_w=T_{\mathrm{top}}$ at every place, all weights $1$ and all lengths $0$, the depth
$j_0=\lceil\log_5(nN(m)/m)\rceil$, the logarithmic derivative at $1$ of the $L$-function of
$\chi_{-d}$, the thick density $|T_2|/n$, $\Lambda_G=\Lambda$, the cloud size $S_c=4k\Delta_S=S$, the
width $\eta$, and the constants $C_c$, $C_t$ and $C_\Gamma$ of $C$. The conditions of size data hold.
The windows for $a_t$, the bound $b_t\ge b'$ and the values of $u_1$ and $u_k$ at the places of top
classes are hypotheses. The conditions on the gains, the local terms and the errors hold since all of
them vanish, and at a place $w$ with $\mu_w=\mu_0$ the indices $I_j$ vanish by hypothesis. There is
no dropped prime. The number $\log\Gamma_P$ and the splitting bound are given by the hypotheses and
condition (2). Computation (B) gives the windows of $b_1$ at the places of top classes, $b_j\le b^2$,
the windows of the $e_j$ at the thick, top and thin places, the cloud conditions $S_c\le
u_1-v_1\le C_cS_c$ and $0\le v_i\le S_c$, the closeness $|b_i-e_j|\le\eta$, $1\le t\le x^{C_t}$,
$10j_0+20\le L_w$ and $0<|T_2|/n\le1$. For every word $u$ the parameter $\lambda_u(\tau)$ has the
parts $v_i$ of $\lambda(\tau)$ (Definition~\ref{th-definite-datum}(2)) and interlaces, so
$v_i<u_i<v_{i-1}$ for $2\le i\le k-1$; hence $0\le u_i\le S_c$ for these $i$. Finally $o\le C_8\le x^{\varepsilon/100}$ and $Q_S\le C_8n\le
4nx^{\varepsilon/100}$ by condition (2).

\emph{The explicit values.} Conditions (1), (2), (3), (7) and (8) are hypotheses, and so is (4), since
the logarithmic derivative of the size data is that of $\chi_{-d}$. In (5), $\Delta_S=\Delta$ is a
hypothesis, and $\eta=2mh$, $\Lambda_G=\Lambda$ and $j_0$ hold by the choices; the bound on $b_t$ is a
hypothesis. The constants in (6) are those of $C$ by the choices, and $C_T$ by hypothesis.
% Lean: EX.exists_sizeData_explicit. Std C gives Ct = 401 (C.Fixed), Cc = 2, and 1 <= GK.La, 1 <=
% CGam (cMid). (A) is bottom_window_le eps C.GK, (B) is T.windows_of_explicitParameters eps cD ch Cpde
% (its conclusions w1, ..., w13 are top_b, b_le, win_thick, win_top, win_thin, cloud_gap, cloud_v,
% be_close, the two halves of t_win, weightLength_ge_depth and the two bounds on T2.card); the
% threshold is max (max x1 x2) ((max Cpde 1)^(100/eps)). The structure L: supp = empty, gain = F = E =
% tau = io = 0, drop = empty, NuLat = bot, TTw = univ, wPart = 1, lenMN = 0, j0 the ceiling,
% logDerivativeAtOne = dlogL D 1, de2 = T2.card / n, LaG = G.La C.GK s, Sc = k (4 Delv), etaSz = m h0v
% with h0v = 2 ceil(x^ch); cloud_u from middle_between with D.hcW_v; o_tie and QS_tie from hH.o_le,
% hH.QS_le and Cpde <= x^(eps/100); split_tie from hH. Val: ranks, host, hc, dlog, Del_eq, words,
% etaF_eq and secondConstant_eq from the hypotheses, the rest rfl. About 2,300 lines of Lean proofs at
% depth 5.
\end{proof}

\begin{lemma}\label{asm-dependence-test-assembly-sizes}
The following holds eventually in the context, for contexts that are good. Let $C_P$ be the
constant $C_F$ of the parameter choices among the constants of the context, and $(x,\varepsilon)$
its scale. Let $X$ be a stage whose base data lie over the context, whose division data have
$\eta_F=6x^{-C_P}$, and whose rows have vanishing indices: $I_j(\mu_w,w)+I_j(\mu_w,\bar w)+j_w(\mu_w)=0$
for every lift $j$ and every place $w$. Then there are size data for the definite-side datum of $X$,
at the scale of the context, whose rows have the sizes of the size data
(Definition~\ref{asm-dependence-functional-witness-sizes}).
\end{lemma}
\leanref{Development/AbelianVarietyTheory/Proofs/TwistedOperators/EventuallyStagesParametersParameterSizeConditionsStatement.lean}{GlobalTheory.T.eventually_stages_parameterSizeConditions}
\leanref{Development/AbelianVarietyTheory/Proofs/TwistedOperators/EventuallyStagesParamSizeConditions.lean}{GlobalTheory.Mem.R3.LedSize.evG_exists_sizeRows}
\begin{proof}
Let $\varepsilon$ and standard constants $C$ be given. Let $c_1$ be the number of
Lemma~\ref{asm-dependence-functional-witness-numbers-values}, and take $c_D\ge\max(c_1,0)$; let $c_2$
be its number for $c_D$, and take $c_h\ge\max(c_2,402)$; let $x_1$ be its number for $c_h$, and $x_0$
the number of Lemma~\ref{asm-dependence-test-assembly-sizes-explicit} for $\varepsilon$, $C$, $c_D$ and
$c_h$. Consider good contexts with these constants at a scale $x\ge\max(x_0,x_1,44^4,8^{4/\varepsilon})$,
and let $X$ be a stage as in the statement, with datum $D$.

\emph{Explicit size data.} Base data over the context carry the rank datum, the host and the
Harish-Chandra parameters of the context. The context records that the rank datum is adapted to
the scale, that the host satisfies condition (2) of the explicit values, the arithmetic conditions
and $|\operatorname{Re}\frac{L'}{L}(1,\chi_{-d})|\le x^{\varepsilon}$, that the Harish-Chandra
parameters at the base word satisfy condition (3) with its integers $b_t$ and have the values of $u_1$
and $u_k$ at the top classes, with its integers $a_t$, that $a_t$ and $b_t$ lie in their windows, that
$C_T$ is that of $C$, the number of words, and $\log\Gamma_P=C_\Gamma m^5fx^{\varepsilon}$. The number
$\Delta_S$ of $D$ is $\lceil x^{c_D}\rceil$: both $v_1-v_2=4\Delta_S$ (Definition~\ref{th-definite-datum}(8))
and $v_1-v_2=4\lceil x^{c_D}\rceil$ (condition (3)). The division data have $\eta_F=6x^{-C_P}$, and the
indices vanish at every place. So Lemma~\ref{asm-dependence-test-assembly-sizes-explicit} gives size
data $L$ for $D$ at the scale of the context that take the explicit values for $(C,c_D,c_h)$.

\emph{The sizes.} By Lemma~\ref{asm-dependence-functional-witness-numbers-values}, $L$ carries
parameter choices with $T$-parameters whose constant $C_{nW}$ is that of $C$, which is $10$, and
satisfies the size bounds, the windows for $n$, $m$ and $k$, $g_0,g_1\le\exp(9x^{\varepsilon})$ and
$\mathfrak c=\mathcal O_K$. The context is good, so, as in the proof of
Lemma~\ref{asm-dependence-functional-witness}(1), $q_0$ is a level prime different from
$g_0$ and $g_1$, and every such level prime is less than $X'$; so $q_0<X'$. Since
$x\ge8^{4/\varepsilon}$, $x^{\varepsilon/4}\ge8$, and $x\ge44^4$. So
Lemma~\ref{asm-dependence-functional-witness-computation} applies, and the rows of $L$ have the sizes
of the size data.
% Lean: Mem.R3.LedSize.evG_exists_sizeRows: szRow_numAt (c1, c2, x1), cD >= max c1 0, ch >=
% max c2 402, exists_sizeData_explicit (x0), threshold max (max x0 x1) (max (44^4) (8^(4/eps)));
% hdel from the cloud gap D.s.cloud_gap and HCVal; the hypotheses of exists_sizeData_explicit from
% K.ok (rank, val, hc u0, dlog, ct2, card, win, top), Cc.A_GamP with K.Gam_eq (GamS), heta and hA;
% hq from hG.ok.small and hG.ok.lev; SetAtModel.sizeRows_of. The indices are Rows.Data.Aj; eta_F is the
% field etaF of the division data (DivX). The open steps are in the lemmas cited.
\end{proof}

\begin{lemma}\label{asm-dependence-test-assembly-host-gaps}
Let $D$ be a definite-side datum with ranks $k,m,t$, $r=k+m-1$ and $N(m)=k+m(t+1)$
(Definition~\ref{th-ranks}), and for a word $u$ and an embedding $\tau$ write $u_i$, $v_i$, $b_j$, $e_j$
for the parts of the Harish-Chandra parameter $\lambda_u(\tau)$ (Definition~\ref{th-hc-parameter});
at the base word these are the parameters of $D$. Let $h$ and $\Delta$ be integers with
$256N(m)^4\le h$ and $256N(m)^4\le\Delta$. Suppose that for every word $u$ and every $\tau$:
$e_j-b_j=h$ for $1\le j\le m$; $b_{j+1}-e_j=h$ for $1\le j<m$; $v_i=(k-1-i)\cdot4\Delta$ for
$1\le i\le k-1$; and $b_1\ge100r^4$. Then $D$ satisfies the gap conditions (3) of
Definition~\ref{asm-dependence-functional-witness-setting}.
\end{lemma}
\leanref{Development/AbelianVarietyTheory/Proofs/IwasawaDecomposition/StandardLatticesDiagonal.lean}{GlobalTheory.Mem.HostFrames.hostHyp_of_gaps}
\begin{proof}
The interlacing and the bounds $e_j-b_j\ge(t+3)/2$, $b_{j+1}-e_j\ge(t+3)/2$ are conditions (2) to (4) of
a Harish-Chandra parameter. They give $b_j<e_j<b_{j+1}$. For every word, $b_{j+1}-b_j=2h$. Since
$N(m)\ge t$ and $N(m)\ge r\ge1$, $2h\ge512N(m)^4\ge4N(m)^3+2t+6$ and $2h\ge256r^4$; these are the gaps
between the $b$-parts. In the sequence $v_1,\dots,v_{k-1},-b_1,\dots,-b_m$ the consecutive $v_i$
differ by $4\Delta\ge100r^4$, $v_{k-1}-(-b_1)=b_1\ge100r^4$ since $v_{k-1}=0$, and the consecutive
$-b_j$ differ by $2h\ge100r^4$.
\end{proof}

\begin{lemma}\label{asm-dependence-test-assembly-host-sizes}
Let $0<\varepsilon<1/52$. There is $x_0$ such that for every scale $(x,\varepsilon)$ with $x\ge x_0$
and every rank datum adapted to it (Definition~\ref{hf-adapted-ranks}), with $n$, $r$ and $N(m)$ as
there, and with $w''(n)=b/(X'n)$:
\begin{enumerate}
\item $256\,x^{\varepsilon}N(m)^4\le\lceil x^y\rceil$ for every $y\ge1610$;
\item $100r^4+1\le w''(n)$.
\end{enumerate}
\end{lemma}
\leanref{Development/AbelianVarietyTheory/Proofs/IwasawaDecomposition/DiagSpacesOfOver.lean}{GlobalTheory.Mem.HostFrames.size_le}
\leanref{Development/AbelianVarietyTheory/Proofs/IwasawaDecomposition/StandardLatticesDiagonal.lean}{GlobalTheory.Mem.HostFrames.thin_size}
\leanref{Development/AbelianVarietyTheory/Proofs/GlobalInterpolation/ExplicitRanksBasicBoundsLargeScale.lean}{GlobalTheory.EX.rank_basic_bounds}
\begin{proof}
For $x$ large, an adapted rank datum has $N(m)\le x^{402}$, $r\le x^5$ and $n\le2x^3$
(Proposition~\ref{hf-adapted-ranks-exist}(2) and the windows of Definition~\ref{hf-adapted-ranks}).
For (1), with $x\ge256$, $256\,x^{\varepsilon}N(m)^4\le x\cdot x\cdot x^{1608}\le x^y\le\lceil x^y\rceil$.
For (2), $w''(n)=\exp(x^{\varepsilon}-x^{\varepsilon/3})/n$, and
$24\log x+x^{\varepsilon/3}\le x^{\varepsilon}$ for $x$ large, so $w''(n)\ge x^{24}/n$. With $x\ge202$,
$(100r^4+1)\,n\le101x^{20}\cdot2x^3\le x^{24}$.
% Lean: size_le eps 256 4 (with rank_basic_bounds: N(m) <= x^402, L = x^eps <= x) and thin_size
% (rank_basic_bounds: r <= x^5; log_add_rpow_le; the threshold max x1 x2 202).
\end{proof}

\begin{lemma}\label{asm-dependence-test-assembly-host}
Eventually in the context, every stage whose base data lie over the context satisfies, with its
definite-side datum $D$: the gap conditions (3) of
Definition~\ref{asm-dependence-functional-witness-setting}; $F$ is the maximal totally real subfield
of $E$; the basis of $X'$ and those of the spaces $V_i$ are orthonormal and span their lattices; and
the spaces $V_i$ are positive definite, the frames of the $V_i$ and $W_i$ are diagonal, and the
lattices of the $Y_j$, $V_j$, $W_j$ and $X'$ are standard.
\end{lemma}
\leanref{Development/AbelianVarietyTheory/Proofs/IwasawaDecomposition/StandardLatticesDiagonal.lean}{GlobalTheory.Mem.HostFrames.hostHyp_frames_ev}
\leanref{Development/AbelianVarietyTheory/Proofs/IwasawaDecomposition/StandardLatticesDiagonal.lean}{GlobalTheory.Mem.HostFrames.first_node}
\leanref{Development/AbelianVarietyTheory/Proofs/IwasawaDecomposition/DiagSpacesOfOver.lean}{GlobalTheory.Mem.HostFrames.frames_of_over}
\leanref{Development/AbelianVarietyTheory/Proofs/IwasawaDecomposition/StandardLatticesDiagonal.lean}{GlobalTheory.Mem.HostFrames.stdLat_diag}
\begin{proof}
Let $0<\varepsilon<1/52$, let $x_0$ be the number of
Lemma~\ref{asm-dependence-test-assembly-host-sizes} for $\varepsilon$, and consider the contexts of
exponent $\varepsilon$ whose constants $c_D$ and $c_h$ are at least $1610$ and whose scale is
$x\ge x_0$; the other constants are arbitrary. Let $X$ be a stage whose base data lie over such a
context $K$, and $D$ its datum. Base data over $K$ carry the rank datum of $K$, which is adapted to
the scale, the host of $K$ and the Harish-Chandra parameters of $K$, and their spaces $Y_j$, $V_j$,
$W_j$ and $X'$ are \emph{diagonal standard spaces}: $E^N$ with the standard basis, the lattice
$\mathcal O_E^N$ and a diagonal form with entries $\pm1$, positive for the $V_i$, $1\le i\le m$, and
for $X'$. The context records that its Harish-Chandra parameters take explicit values: with
$h=\lceil x^{c_h}\rceil$ and $\Delta=\lceil x^{c_D}\rceil$, for every word $u$ and embedding $\tau$,
$v_i=(k-1-i)\cdot4\Delta$, $b_j=b_1+2(j-1)h$ and $e_j=b_j+h$, where $b_1$ is $\lceil w''(n)\rceil$ at the
thin places, at least $\lceil b\rceil$ at the thick places, and $b_T(t)+\lfloor b'/x\rfloor$ at the place
of a top class $t$, whose window gives $b_T(t)\ge b'$.

\emph{The gaps.} By Lemma~\ref{asm-dependence-test-assembly-host-sizes}(1) with $y=c_h$ and $y=c_D$,
and $x^{\varepsilon}\ge1$, $256N(m)^4\le h$ and $256N(m)^4\le\Delta$. By part (2), $100r^4+1\le w''(n)$,
and $w''(n)\le b'/x\le b$. So $b_1\ge100r^4$ at every $\tau$: at a thin place
$b_1\ge w''(n)$; at a thick place $b_1\ge b$; at a top class $b_1\ge b'+b'/x-1\ge w''(n)-1$. Now
Lemma~\ref{asm-dependence-test-assembly-host-gaps} gives the gap conditions.

\emph{The spaces.} $F$ is the maximal totally real subfield of $E$ by the construction of the host.
In a diagonal standard space the frame is diagonal, the lattice is the standard one, and the basis
spans the lattice; where the form is positive, the basis is orthonormal. This gives the remaining
conditions for the $V_i$, the $W_i$, the $Y_j$ and $X'$.
% Lean: hostHyp_frames_ev: ev_of_large with y = 402 * 4 + 2 = 1610 and x0 = max x0 x1 of size_le
% eps 256 4 and thin_size; val_of_over (RankVal and HCVal, with h0v = 2 ceil(x^ch) and Delv =
% ceil(x^cD)); first_node_over = first_node (thin_le, WinOK for b_T >= b'); hostHyp_of_gaps;
% diagSpaces_of_over (the DiagSpaces clause of Ctx.Over), frames_of_over (hF by rfl; signFrame_W,
% signFrame_V, unitFrame_V with V_pos, unitFrame_X') and stdLat_diag (stdLat_std). IsFDiag from
% SignFrame.gram_eq.
\end{proof}

\begin{lemma}\label{asm-dependence-test-assembly-setting}
The following holds eventually in the context. Let $B$ be base data with a framing $f$ and twist
data $T$, lying over the context, let $w$ be a column choice, $Z$ a later stage whose group data are
true for $f$ (Definition~\ref{asm-dependence-test-assembly-true}), and $z$ a realization of $D(w,Z)$
built at the stage that exists at every place. Then for all size data for $D(w,Z)$ whose rows have
the sizes of the size data, $z$ with these size data is a realization in an automorphic setting at
its face (Definition~\ref{asm-dependence-functional-witness-setting}).
\end{lemma}
\leanref{Development/AbelianVarietyTheory/Proofs/TwistedOperators/RowTestHypothesesAtMemberAdelicDataStatement.lean}{GlobalTheory.T.hidaSetting_of_sizeRows}
\leanref{Development/AbelianVarietyTheory/Proofs/TwistedOperators/RowTestHypothesesAtMemberAdelicData.lean}{GlobalTheory.PrimIn.rowTest_setting_dat}
\leanref{Development/AbelianVarietyTheory/Proofs/IwasawaDecomposition/TypeIdempotentCoeffAcrossLayers.lean}{GlobalTheory.Mem.Top.GrpTrue.eSerTie}
\begin{proof}
Take the contexts in which Lemma~\ref{asm-dependence-test-assembly-host} holds. The conditions of
Definition~\ref{asm-dependence-functional-witness-setting} hold as follows. Condition (1) is the
hypothesis. The series idempotents of true group data are those of the label data of $g_0$ and
$g_1$, so they are tied to the labels, (2). Lemma~\ref{asm-dependence-test-assembly-host} gives (3),
(4) and the orthonormal basis of $X'$ in (12). A realization built at a stage records the weak
characterization at its face, whose conditions on the tuples are (6), whose conditions on the group
record the group, (7), and which contains the fullness of the Weil data, (10); it also records the
guard (8), the standard character (11) and the identity adelic matrix of $X'$ in (12). Its tuples
are the base tuple and the tuples of the used places, finitely many, (9). Since $z$ exists at every
place, Lemma~\ref{asm-dependence-functional-witness-core} gives (5).
% Lean: PrimIn.rowTest_setting_dat: hostHyp_frames_ev (host, hF, stdX), hG.eSerTie (eser),
% z.pins.char (core via core_of_weak, mu_*, grp, fullY, fullX), z.pins.keep (guard, std, beta.2),
% and hfin from tuOf of the used places.
\end{proof}

\begin{lemma}\label{asm-dependence-test-assembly-bare-norm-labels-series-induction}
Let $g$ be a prime and write $G_n=GL_n(\mathbf F_g)$. For a finite group $H$ and functions $\alpha,\beta$ on
$H$ write $\langle\alpha,\beta\rangle=|H|^{-1}\sum_{x\in H}\alpha(x)\beta(x^{-1})$, and for functions $\alpha$
on $G_a$ and $\beta$ on $G_b$ write $\alpha\times\beta$ for $(h_1,h_2)\mapsto\alpha(h_1)\beta(h_2)$. For
$a,b\ge0$ and $n=a+b$ let $P_{a,b}\subset G_n$ be the subgroup of block upper triangular matrices with
diagonal blocks of sizes $a$ and $b$, and $\ell\colon P_{a,b}\to G_a\times G_b$ the map to the diagonal
blocks. For a function $\Phi$ on $G_a\times G_b$ and a function $\chi$ on $G_n$ put
\[
 \operatorname{Ind}_{a,b}\Phi(x)=\frac1{|P_{a,b}|}\sum_{\substack{y\in G_n\\y^{-1}xy\in P_{a,b}}}\Phi\bigl(\ell(y^{-1}xy)\bigr),
 \qquad
 \operatorname{Res}_{a,b}\chi(h_1,h_2)=\frac1{g^{ab}}\sum_{\substack{x\in P_{a,b}\\\ell(x)=(h_1,h_2)}}\chi(x).
\]
Call an irreducible character of $G_d$, $d\ge1$, cuspidal as in Theorem~\ref{lit-green-cuspidal}.
\begin{enumerate}
\item (Frobenius reciprocity.) For every class function $\chi$ on $G_n$ and every function $\Phi$ on
  $G_a\times G_b$, $\langle\chi,\operatorname{Ind}_{a,b}\Phi\rangle=\langle\operatorname{Res}_{a,b}\chi,\Phi\rangle$.
\item If $\alpha$ and $\beta$ are characters of $G_a$ and $G_b$, then $\operatorname{Ind}_{a,b}(\alpha\times\beta)$
  is a character of $G_n$. If $\chi$ is a character of $G_n$, then $\operatorname{Res}_{a,b}\chi$ is a finite
  sum of functions $\alpha\times\beta$ with $\alpha$ and $\beta$ irreducible characters of $G_a$ and $G_b$.
\item If $\chi$ is a cuspidal irreducible character of $G_d$ and $0<a<d$, then
  $\operatorname{Res}_{a,d-a}\chi=0$.
\item Let $\alpha$, $\beta$ and $\gamma$ be class functions on $G_a$, $G_b$ and $G_c$. Then
  \[
   \operatorname{Ind}_{a+b,c}\bigl(\operatorname{Ind}_{a,b}(\alpha\times\beta)\times\gamma\bigr)
   =\operatorname{Ind}_{a,b+c}\bigl(\alpha\times\operatorname{Ind}_{b,c}(\beta\times\gamma)\bigr),
  \]
  $\operatorname{Ind}_{a,b}(\alpha\times\beta)=\operatorname{Ind}_{b,a}(\beta\times\alpha)$, and
  $\operatorname{Ind}_{0,a}(1\times\alpha)=\operatorname{Ind}_{a,0}(\alpha\times1)=\alpha$.
\item (Mackey formula.) Let $n=a+b=c+d$, and let $\Phi$ be a function on $G_a\times G_b$ that is a class
  function in each variable. For $0\le i\le a$ with $i\le c$ and $c-i\le b$ put
  $\Phi_i(y_1,y_2,y_3,y_4)=\operatorname{Res}_{i,a-i}\bigl(x\mapsto\operatorname{Res}_{c-i,b-c+i}(\Phi(x,\cdot))(y_3,y_4)\bigr)(y_1,y_2)$
  for $y_1\in G_i$, $y_2\in G_{a-i}$, $y_3\in G_{c-i}$ and $y_4\in G_{b-c+i}$. Then for $h_1\in G_c$ and
  $h_2\in G_d$
  \[
   \operatorname{Res}_{c,d}\operatorname{Ind}_{a,b}\Phi(h_1,h_2)=\sum_i\operatorname{Ind}_{i,c-i}\Bigl((y_1,y_3)\mapsto
   \operatorname{Ind}_{a-i,b-c+i}\bigl((y_2,y_4)\mapsto\Phi_i(y_1,y_2,y_3,y_4)\bigr)(h_2)\Bigr)(h_1),
  \]
  the sum over the integers $i$ with $0\le i\le a$, $i\le c$ and $c-i\le b$.
\end{enumerate}
\end{lemma}
\leanref{Development/AbelianVarietyTheory/Proofs/InertiaEigenvalues/HarishChandraFrobeniusReciprocityStatement.lean}{GlobalTheory.T.harishChandra_frobenius_reciprocity}
\leanref{Development/AbelianVarietyTheory/Proofs/InertiaEigenvalues/HarishChandraInductionCharacterRestrictionDecompStatement.lean}{GlobalTheory.T.harishChandra_induction_character_restriction_decomp}
\leanref{Development/AbelianVarietyTheory/Proofs/InertiaEigenvalues/HarishChandraRestrictionTwoBlockCuspidalEqZeroStatement.lean}{GlobalTheory.T.harishChandraRestrictionTwoBlock_cuspidal_eq_zero}
\leanref{Development/AbelianVarietyTheory/Proofs/InertiaEigenvalues/HarishChandraInductionExteriorProductAssocCommStatement.lean}{GlobalTheory.T.harishChandraInductionExteriorProduct_assoc_comm}
\leanref{Development/AbelianVarietyTheory/Proofs/InertiaEigenvalues/MackeyFormulaHarishChandraRestrictionInductionStatement.lean}{GlobalTheory.T.mackey_formula_harishChandra_restriction_induction}

\begin{lemma}\label{asm-dependence-test-assembly-bare-norm-labels-series-mackey}
In the notation of Lemma~\ref{asm-dependence-test-assembly-bare-norm-labels-series-induction}, for class
functions $\tau_1,\dots,\tau_s$ on $G_{d_1},\dots,G_{d_s}$ put $R()=1$ and
$R(\tau_1,\dots,\tau_s)=\operatorname{Ind}_{d_1,d_2+\dots+d_s}\bigl(\tau_1\times R(\tau_2,\dots,\tau_s)\bigr)$, and
for $I\subset\{1,\dots,s\}$ let $R_I$ be $R$ of the $\tau_i$ with $i\in I$, in increasing order of $i$. Let
$\tau_1,\dots,\tau_s$ be cuspidal irreducible characters of $G_{d_1},\dots,G_{d_s}$, and let $a,b\ge0$ with
$a+b=d_1+\dots+d_s$. Then
\[
 \operatorname{Res}_{a,b}R(\tau_1,\dots,\tau_s)=\sum_IR_I\times R_{I^c},
\]
the sum over the subsets $I$ of $\{1,\dots,s\}$ with $\sum_{i\in I}d_i=a$, where $I^c$ is the complement
of $I$.
\end{lemma}
\leanref{Development/AbelianVarietyTheory/Proofs/InertiaEigenvalues/HarishChandraRestrictionCuspidalProductSumSubsetStatement.lean}{GlobalTheory.T.harishChandraRestriction_cuspidal_product_sum_subsets}
\leanref{Development/AbelianVarietyTheory/Proofs/InertiaEigenvalues/HarishChandraRestrictionCuspidalProductSumSubset.lean}{GlobalTheory.GLq.ResCuspProd.main}
\begin{proof}
Two facts follow from the definitions. For a function $\chi$ on $G_n$,
$\operatorname{Res}_{0,n}\chi(h_1,h_2)=\chi(h_2)$ and $\operatorname{Res}_{n,0}\chi(h_1,h_2)=\chi(h_1)$, since
$P_{0,n}=P_{n,0}=G_n$. And $\operatorname{Ind}_{a,b}$ and $\operatorname{Res}_{a,b}$ are linear. Each
$R(\tau_1,\dots,\tau_s)$ is a class function, being induced or $1$.

We argue by induction on $s$. If $s=0$, then $a=b=0$, $\operatorname{Res}_{0,0}1=1$, and the only subset is
$I=\emptyset$, with $R_\emptyset\times R_\emptyset=1$.

Let $s\ge1$, put $e=d_1$, $m=d_2+\dots+d_s$ and $R'=R(\tau_2,\dots,\tau_s)$, so that
$R(\tau_1,\dots,\tau_s)=\operatorname{Ind}_{e,m}(\tau_1\times R')$, and $e\ge1$ as $\tau_1$ is cuspidal. The
function $\tau_1\times R'$ is a class function in each variable. By the Mackey formula
(Lemma~\ref{asm-dependence-test-assembly-bare-norm-labels-series-induction}(5)), with $(a,b,c,d)$ there equal
to $(e,m,a,b)$, and by linearity, $\operatorname{Res}_{a,b}R(\tau_1,\dots,\tau_s)(h_1,h_2)$ is the sum over the
$i$ with $0\le i\le e$, $i\le a$ and $a-i\le m$ of
\[
 T_i=\operatorname{Ind}_{i,a-i}\Bigl((y_1,y_3)\mapsto\operatorname{Ind}_{e-i,m-a+i}\bigl((y_2,y_4)\mapsto
 \operatorname{Res}_{i,e-i}\tau_1(y_1,y_2)\,\operatorname{Res}_{a-i,m-a+i}R'(y_3,y_4)\bigr)(h_2)\Bigr)(h_1).
\]
If $0<i<e$, then $\operatorname{Res}_{i,e-i}\tau_1=0$ by
Lemma~\ref{asm-dependence-test-assembly-bare-norm-labels-series-induction}(3), so $T_i=0$. There remain
$i=0$, present when $a\le m$, and $i=e$, present when $e\le a$; they are distinct as $e\ge1$.

\emph{The term $i=0$.} Here $\operatorname{Res}_{0,e}\tau_1(y_1,y_2)=\tau_1(y_2)$, and by induction
$\operatorname{Res}_{a,m-a}R'=\sum_JR_J\times R_{J'}$, over the subsets $J$ of $\{2,\dots,s\}$ with
$\sum_{j\in J}d_j=a$, where $J'=\{2,\dots,s\}\setminus J$. By linearity the inner induction is
$\sum_JR_J(y_3)\operatorname{Ind}_{e,m-a}(\tau_1\times R_{J'})(h_2)$, a class function of $y_3$ that does not
depend on $y_1\in G_0$, and $\operatorname{Ind}_{0,a}(1\times\alpha)=\alpha$ for class functions $\alpha$
(Lemma~\ref{asm-dependence-test-assembly-bare-norm-labels-series-induction}(4)). So
$T_0=\sum_JR_J(h_1)\operatorname{Ind}_{e,m-a}(\tau_1\times R_{J'})(h_2)=\sum_JR_J(h_1)R_{\{1\}\cup J'}(h_2)$.

\emph{The term $i=e$.} Here $\operatorname{Res}_{e,0}\tau_1(y_1,y_2)=\tau_1(y_1)$, and by induction
$\operatorname{Res}_{a-e,m-a+e}R'=\sum_JR_J\times R_{J'}$, over the subsets $J$ of $\{2,\dots,s\}$ with
$\sum_{j\in J}d_j=a-e$. The inner induction $\operatorname{Ind}_{0,m-a+e}$ is applied to a class function of
$y_4$ alone, so it gives that function at $h_2$, and
$T_e=\sum_J\operatorname{Ind}_{e,a-e}(\tau_1\times R_J)(h_1)R_{J'}(h_2)=\sum_JR_{\{1\}\cup J}(h_1)R_{J'}(h_2)$.

If $a>m$ there is no $J\subset\{2,\dots,s\}$ with $\sum_{j\in J}d_j=a$, and if $e>a$ there is none with sum
$a-e$; so in all cases $T_0$ is the sum of $R_I(h_1)R_{I^c}(h_2)$ over the subsets $I$ with
$\sum_{i\in I}d_i=a$ that do not contain $1$, and $T_e$ the sum over those that contain $1$.
% Lean: ResCuspProd.main, by induction on L: T.mackey_formula_harishChandra_restriction_induction with
% res_const_mul and res_mul_const; collapse_mid (harishChandraRestrictionTwoBlock_cuspidal_eq_zero) for
% 0 < i < e; collapse_left (res_left_zero, hcT_sum, hcT_mul_const, hcT_left_zero) for i = 0 and
% collapse_right (res_right_zero, hcT_left_zero) for i = e; sum_subsets_fin_succ splits the subsets of
% {0, ..., s} by whether they contain 0, with sort_map_succ, sort_insert_map_succ, compl_map_succ and
% compl_insert_map_succ; isClass_hcL for the class functions.
\end{proof}

\begin{theorem}\label{asm-dependence-test-assembly-bare-norm-labels-series}
Let $g$ be a prime, write $G_n=GL_n(\mathbf F_g)$ and
$\langle\alpha,\beta\rangle=|G_n|^{-1}\sum_{x\in G_n}\alpha(x)\beta(x^{-1})$ for class functions on
$G_n$, and call an irreducible character of $G_d$, $d\ge1$, cuspidal as in
Theorem~\ref{lit-green-cuspidal}. For class functions $\tau_1,\dots,\tau_s$ on $G_{d_1},\dots,G_{d_s}$
let $R(\tau_1,\dots,\tau_s)$ be their iterated Harish-Chandra induction to $G_{d_1+\dots+d_s}$: $R()=1$,
and $R(\tau_1,\dots,\tau_s)$ is induced from the parabolic subgroup of block upper triangular matrices
with diagonal blocks of sizes $d_1$ and $d_2+\dots+d_s$ of the function $\tau_1\times
R(\tau_2,\dots,\tau_s)$ of the diagonal blocks. Let $n\ge0$.
\begin{enumerate}
\item If $\tau_1,\dots,\tau_s$ are characters, then $R(\tau_1,\dots,\tau_s)$ is a character.
\item For every irreducible character $\chi$ of $G_n$ there are cuspidal irreducible characters $\tau_i$
  of $G_{d_i}$, $1\le i\le s$, with $d_1+\dots+d_s=n$ and $\langle\chi,R(\tau_1,\dots,\tau_s)\rangle\ne0$.
\item Let $(\tau_1,\dots,\tau_s)$ and $(\tau'_1,\dots,\tau'_r)$ be lists of cuspidal irreducible
  characters of groups $G_{d_i}$ and $G_{d'_j}$ with $\sum_id_i=\sum_jd'_j=n$. If the second list is not a
  rearrangement of the first, then $\langle R(\tau_1,\dots,\tau_s),R(\tau'_1,\dots,\tau'_r)\rangle=0$.
\end{enumerate}
\end{theorem}
\leanref{Development/AbelianVarietyTheory/Proofs/InertiaEigenvalues/HarishChandraProductListIsCharacterAndPermInvariStatement.lean}{GlobalTheory.T.harishChandraProductList_isCharacter_and_perm_invariant}
\leanref{Development/AbelianVarietyTheory/Proofs/InertiaEigenvalues/ExistsCuspidalListHarishChandraPairingNeZeroStatement.lean}{GlobalTheory.T.exists_cuspidal_list_harishChandra_pairing_ne_zero}
\leanref{Development/AbelianVarietyTheory/Proofs/InertiaEigenvalues/CharacterPairingHarishChandraCuspidalMackeyStatement.lean}{GlobalTheory.T.characterPairing_harishChandra_cuspidal_mackey}
\leanref{Development/AbelianVarietyTheory/Proofs/InertiaEigenvalues/ExistsCuspidalListHarishChandraPairingNeZero.lean}{GlobalTheory.GLq.SuppCusp.exists_cusp_list}
\leanref{Development/AbelianVarietyTheory/Proofs/InertiaEigenvalues/CharacterPairingHarishChandraCuspidalMackey.lean}{GlobalTheory.GLq.IpCuspProd.finiteRing_sq_sub_mul_sq_eq_of_card}
\begin{proof}
Write $\operatorname{Ind}$ and $\operatorname{Res}$ as in
Lemma~\ref{asm-dependence-test-assembly-bare-norm-labels-series-induction}; the function $R$ of the
statement is that of Lemma~\ref{asm-dependence-test-assembly-bare-norm-labels-series-mackey}. The pairing
satisfies $\langle\alpha,\beta\rangle=\langle\beta,\alpha\rangle$, by the substitution $x\mapsto x^{-1}$, and
$\langle\alpha\times\beta,\gamma\times\delta\rangle=\langle\alpha,\gamma\rangle\langle\beta,\delta\rangle$ on
$G_a\times G_b$. For irreducible characters $\chi$ and $\psi$ of $G_n$, $\langle\chi,\psi\rangle$ is $1$ if
$\chi=\psi$ and $0$ otherwise, so $\langle\chi,A\rangle$ is a nonnegative integer for every character $A$.

(1) By induction on $s$: $R()=1$ is a character, and if $R(\tau_2,\dots,\tau_s)$ is one, so is
$R(\tau_1,\dots,\tau_s)=\operatorname{Ind}_{d_1,d_2+\dots+d_s}(\tau_1\times R(\tau_2,\dots,\tau_s))$, by
Lemma~\ref{asm-dependence-test-assembly-bare-norm-labels-series-induction}(2).

\emph{Two consequences of
Lemma~\ref{asm-dependence-test-assembly-bare-norm-labels-series-induction}(4).} For a class function $\tau$ on
$G_d$, $R(\tau)=\operatorname{Ind}_{d,0}(\tau\times1)=\tau$. Let $L_1=(\tau_1,\dots,\tau_k)$ and
$L_2=(\tau_{k+1},\dots,\tau_s)$ be lists of characters with $\sum_{i\le k}d_i=m_1$ and $\sum_{i>k}d_i=m_2$,
and write $(L_1,L_2)$ for the joined list. Then
$R(L_1,L_2)=\operatorname{Ind}_{m_1,m_2}(R(L_1)\times R(L_2))$. This is shown by induction on $k$. For $k=0$
it is $R(L_2)=\operatorname{Ind}_{0,m_2}(1\times R(L_2))$, as $R(L_2)$ is a character by (1), hence a class
function. For $k\ge1$, by the definition, the induction hypothesis and associativity,
\[
 R(L_1,L_2)=\operatorname{Ind}_{d_1,m_1-d_1+m_2}\bigl(\tau_1\times\operatorname{Ind}_{m_1-d_1,m_2}(R(\tau_2,\dots,\tau_k)\times
 R(L_2))\bigr)=\operatorname{Ind}_{m_1,m_2}\bigl(R(L_1)\times R(L_2)\bigr).
\]

(2) By induction on $n$. Let $\chi$ be an irreducible character of $G_n$; then $\langle\chi,\chi\rangle=1$.
If $n=0$, $G_0$ is trivial and the empty list works: $\langle\chi,R()\rangle=\chi(1)$, and
$\chi(1)^2=\langle\chi,\chi\rangle=1$. If $\chi$ is cuspidal, the list $(\chi)$ works, as $R(\chi)=\chi$. Let
$n\ge1$ and $\chi$ not cuspidal. Then for some $0<a<n$ the unipotent radical $U$ of $P_{a,b}$, $b=n-a$, has a
nonzero fixed vector in the representation with character $\chi$. The dimension of the space of fixed
vectors is $|U|^{-1}\sum_{u\in U}\chi(u)=\operatorname{Res}_{a,b}\chi(1,1)$, as $|U|=g^{ab}$; so
$\operatorname{Res}_{a,b}\chi(1,1)\ne0$. By
Lemma~\ref{asm-dependence-test-assembly-bare-norm-labels-series-induction}(2),
$\operatorname{Res}_{a,b}\chi=\sum_k\alpha_k\times\beta_k$ with irreducible characters $\alpha_k$ of $G_a$ and
$\beta_k$ of $G_b$; the sum is not empty. Fix one $k$. By induction there are lists $L_1$ and $L_2$ of
cuspidal irreducible characters of total sizes $a$ and $b$ with $\langle\alpha_k,R(L_1)\rangle\ne0$ and
$\langle\beta_k,R(L_2)\rangle\ne0$. The joined list has total size $n$, and by Frobenius reciprocity
(Lemma~\ref{asm-dependence-test-assembly-bare-norm-labels-series-induction}(1)), as $\chi$ is a class
function,
\[
 \langle\chi,R(L_1,L_2)\rangle=\langle\chi,\operatorname{Ind}_{a,b}(R(L_1)\times R(L_2))\rangle
 =\sum_{k'}\langle\alpha_{k'},R(L_1)\rangle\langle\beta_{k'},R(L_2)\rangle.
\]
By (1), $R(L_1)$ and $R(L_2)$ are characters, so every term is a nonnegative integer, and the term at $k$ is
not $0$. So the sum is not $0$.

(3) By induction on $r$, for all lists $(\tau_1,\dots,\tau_s)$ at once. If $r=0$, then $n=0$, so $s=0$ as
every $d_i\ge1$, and the second list is a rearrangement of the first. Let $r\ge1$, put $\psi=\tau'_1$,
$a=d'_1$, $M=(\tau'_2,\dots,\tau'_r)$ and $m=n-a$. With the notation $R_I$ of
Lemma~\ref{asm-dependence-test-assembly-bare-norm-labels-series-mackey}, Frobenius reciprocity, applied to
the class function $R(\tau_1,\dots,\tau_s)$, and that lemma give
\[
 \langle R(\tau_1,\dots,\tau_s),R(\tau'_1,\dots,\tau'_r)\rangle
 =\langle\operatorname{Res}_{a,m}R(\tau_1,\dots,\tau_s),\psi\times R(M)\rangle
 =\sum_I\langle R_I,\psi\rangle\langle R_{I^c},R(M)\rangle,
\]
over the $I$ with $\sum_{i\in I}d_i=a$. Such an $I$ is not empty, as $a\ge1$. If $I$ has at least two
elements and $i_1$ is the least, then $0<d_{i_1}<a$ and
$R_I=\operatorname{Ind}_{d_{i_1},a-d_{i_1}}(\tau_{i_1}\times R_{I\setminus\{i_1\}})$, so
$\langle R_I,\psi\rangle=\langle\psi,R_I\rangle
=\langle\operatorname{Res}_{d_{i_1},a-d_{i_1}}\psi,\tau_{i_1}\times R_{I\setminus\{i_1\}}\rangle=0$
by Frobenius reciprocity and Lemma~\ref{asm-dependence-test-assembly-bare-norm-labels-series-induction}(3),
$\psi$ being cuspidal. If $I=\{i\}$, then $R_I=\tau_i$ is an irreducible character of $G_a$, and
$\langle\tau_i,\psi\rangle$ is $1$ if $\tau_i=\psi$ and $0$ otherwise. Hence the pairing is the sum, over the
$i$ with $\tau_i=\psi$, of $\langle R(L_i),R(M)\rangle$, where $L_i$ is the list $(\tau_1,\dots,\tau_s)$ with
$\tau_i$ removed, of total size $m$. If the second list is not a rearrangement of the first, then for each
such $i$ the list $M$ is not a rearrangement of $L_i$, and by induction every term is $0$.
% Lean: (1) is GLqPerm.isChar_hcL (harishChandra_induction_character_restriction_decomp, first half). The two
% consequences are hcL_single and SuppCusp.hcL_append (harishChandraInductionExteriorProduct_assoc_comm,
% class_hcL). (2) is SuppCusp.exists_cusp_list (strong induction on n; the case n = 0 through hc1 and hself;
% IsCusp chi with hcL_cons and assoc_comm; otherwise hne, the decomposition hRes of
% harishChandra_induction_character_restriction_decomp, ih twice, hcL_append, harishChandra_frobenius_reciprocity,
% ip_prod_list and sum_nat_ne_zero with ip_irr_char_nat). (3) is the first half of
% IpCuspProd.finiteRing_sq_sub_mul_sq_eq_of_card (induction on L'): key (Frobenius reciprocity,
% harishChandraRestriction_cuspidal_product_sum_subsets, ip_sum_left, ip_prod), ip_sub_cusp (ip_long_cusp
% through ip_comm and harishChandraRestrictionTwoBlock_cuspidal_eq_zero for |I| >= 2, ip_single_cusp through
% hcL_single and the orthogonality centralIdempotentOfIrrCharacter_orthogonal_central for |I| = 1), hrest and
% ih. The Lean's IsCusp asks that Res_{a,d-a} chi (1,1) = 0, the dimension of the fixed vectors.
\end{proof}

\begin{lemma}\label{asm-dependence-test-assembly-bare-norm-labels-uncertainty}
Let $C$ be a finite abelian group, $S$ a finite set of homomorphisms $C\to\mathbf C^\times$ and
$c\colon S\to\mathbf C$ a function. If $f=\sum_{\theta\in S}c(\theta)\theta$ is not identically zero,
then $|C|\le|S|\cdot|\{z\in C:f(z)\ne0\}|$.
\end{lemma}
\leanref{Development/AbelianVarietyTheory/Proofs/FlagMeasures/CardLeCardMulCardSupport.lean}{GlobalTheory.GLm.card_le_card_mul_card_support}
\begin{proof}
Let $T=\{z:f(z)\ne0\}$ and $M=\max_z|f(z)|>0$. For $\eta\in S$, the orthogonality of characters gives
$|C|\,c(\eta)=\sum_{z\in T}f(z)\eta(z^{-1})$, so $|C|\,|c(\eta)|\le|T|M$, as the values of $\eta$ have
absolute value $1$. Hence $M\le\sum_{\eta\in S}|c(\eta)|\le|S|\,|T|\,M/|C|$, and $|C|\le|S|\,|T|$.
\end{proof}

\begin{lemma}\label{asm-dependence-test-assembly-bare-norm-labels-count}
Let $g\ge5$ be a prime and $d\ge1$, and let $B_d$ be the set of the $z\in\bar{\mathbf F}_g^\times$ with
$z^{g^d}=z$ whose degree over $\mathbf F_g$ is not $d$. Then $2d\,|B_d|<g^d-1$.
\end{lemma}
\leanref{Development/AbelianVarietyTheory/Proofs/FlagMeasures/CardProperSubfieldElementsLt.lean}{GlobalTheory.GLm.card_proper_subfield_elements_lt}
\begin{proof}
An element $z\in B_d$ lies in $\mathbf F_{g^d}$, so its degree $e$ divides $d$, and $e\ne d$, so $2e\le
d$; and $z^{g^e-1}=1$. So $|B_d|\le\sum_{e=1}^{\lfloor d/2\rfloor}(g^e-1)$. Put
$T_m=\sum_{e=1}^m(g^e-1)$ and $x=g^m$ for $m\ge1$. Induction on $m$, with $g\ge5$, gives
$4T_m+4m+5\le5x$ and $5m\le x$. Multiplying the first by $m$ and using the second,
$4mT_m+1<x^2=g^{2m}$. For $d=2m$ this is $2d\,T_m<g^d-1$. For $d=2m+1$,
$2d\,T_m=(4m+2)T_m\le5\cdot4mT_m<5(g^{2m}-1)<g^d-1$. For $d=1$, $B_d$ is empty.
% Lean: card_proper_subfield_elements_lt with SubfieldCount.num_bound (sum_bound, five_mul_le) and
% natDegree_dvd_iff.
\end{proof}

\begin{lemma}\label{asm-dependence-test-assembly-bare-norm-labels-cuspidal}
Let $g\ge5$ be a prime with a labelling of $\bar{\mathbf F}_g$ (Definition~\ref{lit-labelling}), let
$d\ge1$ and $P,P'\in\mathcal P_d$ (Corollary~\ref{lit-green-labelled}), and let $\chi$ be a cuspidal
irreducible character of $GL_d(\mathbf F_g)$ that satisfies the character formula of
Corollary~\ref{lit-green-labelled}(1) both for $P$ and for $P'$. Then $P=P'$.
\end{lemma}
\leanref{Development/AbelianVarietyTheory/Proofs/InertiaEigenvalues/CuspidalCharacterLabelUniqueGLnStatement.lean}{GlobalTheory.T.cuspidalCharacter_label_unique_GLn}
\leanref{Development/AbelianVarietyTheory/Proofs/InertiaEigenvalues/CuspidalCharacterLabelUniqueGLn.lean}{GlobalTheory.GLq.LabelUnique.label_eq}
\begin{proof}
Let $t$ and $\gamma_d$ be those of the labelling, put $m=g^d-1$, and let $x$ and $x'$ be roots of $P$ and
$P'$. They lie in $\mathbf F_{g^d}^\times$, so $x^m=x'^m=1$. The generator $\gamma_d$ of
$\mathbf F_{g^d}^\times$ has order $m$. Let $B\subset\mathbf Z/m$ be the set of $k$ such that $\gamma_d^k$
does not have degree $d$; as $k\mapsto\gamma_d^k$ is injective on $\mathbf Z/m$ and carries $B$ into
the set $B_d$ of Lemma~\ref{asm-dependence-test-assembly-bare-norm-labels-count}, $2d|B|<m$.

Let $k\notin B$. Then $\gamma_d^k$ has degree $d$, and it is a root of the characteristic polynomial of
the companion matrix $y\in GL_d(\mathbf F_g)$ of its minimal polynomial, which is irreducible. The
character formula for $P$ and for $P'$ at $y$, with $z=\gamma_d^k$, gives
$(-1)^{d-1}\sum_{i<d}\theta_{x,d}(\gamma_d^{kg^i})=\chi(y)=(-1)^{d-1}\sum_{i<d}\theta_{x',d}(\gamma_d^{kg^i})$.
Since $\theta_{x,d}(\gamma_d^{kg^i})=t(x)^{kg^i}=t(x^{g^i})^k$, with $\zeta_i=t(x^{g^i})$ and
$\zeta'_i=t(x'^{g^i})$, which are $m$-th roots of unity,
\[
 \sum_{i<d}\zeta_i^k=\sum_{i<d}\zeta_i'^k\qquad(k\in\mathbf Z/m,\ k\notin B).
\]
The maps $k\mapsto\zeta^k$, for $m$-th roots of unity $\zeta$, are characters of $\mathbf Z/m$, and a
character determines $\zeta$ as its value at $1$. Let $S$ be the set of the characters $k\mapsto\zeta_i^k$
and $k\mapsto\zeta_i'^k$, so $|S|\le2d$, and for $\theta\in S$ let $c(\theta)$ be the number of $i<d$ with
$\zeta_i=\theta(1)$ minus the number of $i<d$ with $\zeta'_i=\theta(1)$. Then
$f=\sum_\theta c(\theta)\theta$ is $k\mapsto\sum_i\zeta_i^k-\sum_i\zeta_i'^k$, which vanishes off $B$. If
$f$ were not zero, Lemma~\ref{asm-dependence-test-assembly-bare-norm-labels-uncertainty} would give
$m\le|S|\,|B|\le2d|B|<m$. So $f=0$, and by the linear independence of characters every $c(\theta)$
vanishes. For $\theta=(k\mapsto\zeta_0'^k)$ the second count is at least $1$, so $\zeta'_0=\zeta_i$ for
some $i$, that is, $t(x')=t(x^{g^i})$. As $t$ is injective, $x'=x^{g^i}$, which is a root of $P$. So $P$
and $P'$ are the minimal polynomial of $x'$, and $P=P'$.
% Lean: T.cuspidalCharacter_label_unique_GLn (first clause) = GLq.isCuspOf_label_unique through
% LabelUnique.label_eq: card_proper_subfield_elements_lt for hB, exists_irreducible_charpoly_mem_rootsBar
% (the left multiplication matrix of a power basis, whose characteristic polynomial is the minimal
% polynomial), thetaOf_gen_pow, exists_eq_of_sum_eq (card_le_card_mul_card_support and
% linearIndependent_monoidHom), teich_inj and aeval_pow_eq_zero. The second clause of the Lean
% statement, that a polynomial has only one cuspidal character, is not used here.
\end{proof}

\begin{lemma}\label{asm-dependence-test-assembly-bare-norm-labels}
Assume Green's parametrization (Corollary~\ref{lit-green-labelled}). Let $g\ge5$ be a prime with a
labelling of $\bar{\mathbf F}_g$ (Definition~\ref{lit-labelling}), let $n\ge0$, and write
$G_n=GL_n(\mathbf F_g)$ and $\langle\alpha,\beta\rangle=|G_n|^{-1}\sum_{x\in G_n}\alpha(x)\beta(x^{-1})$
for class functions. For $P\in\mathcal P_d$, a \emph{cuspidal character of $P$} is a cuspidal
irreducible character of $G_d$ that satisfies the character formula of
Corollary~\ref{lit-green-labelled}(1) for $P$. For class functions $\tau_1,\dots,\tau_s$ on
$G_{d_1},\dots,G_{d_s}$ let $R(\tau_1,\dots,\tau_s)$ be their iterated Harish-Chandra induction to
$G_{d_1+\dots+d_s}$: $R()=1$, and $R(\tau_1,\dots,\tau_s)$ is induced from the parabolic subgroup of
block upper triangular matrices with diagonal blocks of sizes $d_1$ and $d_2+\dots+d_s$ of the
function $\tau_1\times R(\tau_2,\dots,\tau_s)$ of the diagonal blocks. A class function $\chi$ on $G_n$
has \emph{label} $f\in\mathbf F_g[X]$ if there are $P_i\in\mathcal P_{d_i}$ and cuspidal characters
$\tau_i$ of $P_i$, $1\le i\le s$, with $d_1+\dots+d_s=n$, $f=P_1\cdots P_s$ and
$\langle\chi,R(\tau_1,\dots,\tau_s)\rangle\ne0$.
\begin{enumerate}
\item Every irreducible character of $G_n$ has exactly one label.
\item Every label of a class function on $G_n$ is monic of degree $n$ with nonzero constant term.
\item Every monic polynomial of degree $n$ over $\mathbf F_g$ with nonzero constant term is the
  label of an irreducible character of $G_n$.
\end{enumerate}
\end{lemma}
\leanref{Development/AbelianVarietyTheory/Definitions/NonGaloisDescent/LabelDatum.lean}{GlobalTheory.GLq.HasLabel}
\leanref{Development/AbelianVarietyTheory/Proofs/InertiaEigenvalues/HarishChandraSeriesLabelExistsUniqueGLnStatement.lean}{GlobalTheory.T.harishChandraSeriesLabel_exists_unique_GLn}
\leanref{Development/AbelianVarietyTheory/Proofs/InertiaEigenvalues/HarishChandraSeriesLabelExistsUniqueGLn.lean}{GlobalTheory.GLq.Label.label_main}
\leanref{Development/AbelianVarietyTheory/Proofs/FlagMeasures/HarishChandraSeriesCentralIdempotentNeZero.lean}{GlobalTheory.GLm.exists_irr_hasLabel}
\leanref{Development/AbelianVarietyTheory/Proofs/FlagMeasures/HarishChandraSeriesCentralIdempotentNeZero.lean}{GlobalTheory.GLm.exists_cusp_list}
\begin{proof}
(2) If $f=P_1\cdots P_s$ with $P_i\in\mathcal P_{d_i}$ and $\sum d_i=n$, then $f$ is monic of degree
$n$, and its constant term is the product of the nonzero constant terms of the $P_i$.

(3) Let $f$ be monic of degree $n$ with $f(0)\ne0$, and factor it as $P_1\cdots P_s$ into monic
irreducible factors, of degrees $d_i$ with $\sum d_i=n$; their constant terms are nonzero, so
$P_i\in\mathcal P_{d_i}$. By Corollary~\ref{lit-green-labelled}(1) each $P_i$ has a cuspidal character
$\tau_i$. The function $R=R(\tau_1,\dots,\tau_s)$ is a character, induced from a parabolic subgroup
of a character of its Levi quotient, and its degree satisfies
\[
 R(1)\prod_{i=1}^s\prod_{j=1}^{d_i}(g^j-1)=\prod_{j=1}^n(g^j-1)\prod_{i=1}^s\tau_i(1),
\]
whose right side is nonzero. So $R\ne0$. As $R$ is the sum of the irreducible characters $\chi$ with
multiplicities $\langle\chi,R\rangle$, some irreducible $\chi$ has $\langle\chi,R\rangle\ne0$, and this
$\chi$ has label $f$.

(1) \emph{Existence.} Let $\chi$ be an irreducible character of $G_n$. By
Theorem~\ref{asm-dependence-test-assembly-bare-norm-labels-series}(2) there are cuspidal irreducible
characters $\tau_i$ of $G_{d_i}$, $1\le i\le s$, with $\sum d_i=n$ and
$\langle\chi,R(\tau_1,\dots,\tau_s)\rangle\ne0$. By Corollary~\ref{lit-green-labelled}(3) each $\tau_i$
is $\tau_{P_i}$ for some $P_i\in\mathcal P_{d_i}$, and by part (1) of that corollary it satisfies the
character formula for $P_i$; so it is a cuspidal character of $P_i$. Hence $\chi$ has the label
$P_1\cdots P_s$.

\emph{Uniqueness.} Let $\chi$ have the labels $f=P_1\cdots P_s$ and $f'=P'_1\cdots P'_r$, given by
cuspidal characters $\tau_i$ of $P_i$ and $\tau'_j$ of $P'_j$, and put $R=R(\tau_1,\dots,\tau_s)$ and
$R'=R(\tau'_1,\dots,\tau'_r)$. These are characters, by
Theorem~\ref{asm-dependence-test-assembly-bare-norm-labels-series}(1); write $R=\sum_\rho m_\rho\rho$ and
$R'=\sum_\rho m'_\rho\rho$ over the irreducible characters $\rho$, with integers $m_\rho,m'_\rho\ge0$. By
the orthogonality of irreducible characters, $m_\chi=\langle\chi,R\rangle\ne0$, $m'_\chi\ne0$, and
$\langle R,R'\rangle=\sum_\rho m_\rho m'_\rho\ge m_\chi m'_\chi>0$. So, by
Theorem~\ref{asm-dependence-test-assembly-bare-norm-labels-series}(3), $r=s$ and there is a permutation
$\pi$ with $\tau'_{\pi(i)}=\tau_i$ for every $i$. Then $\tau_i$ is a cuspidal character of $P_i$ and of
$P'_{\pi(i)}$, so $P_i=P'_{\pi(i)}$ by Lemma~\ref{asm-dependence-test-assembly-bare-norm-labels-cuspidal},
as $g\ge5$. Hence $f=f'$.
% Lean: (1) is GLq.Label.label_main: label_exists (T.exists_cuspidal_list_harishChandra_pairing_ne_zero,
% Green for each member, lab) and label_unique (T.harishChandraProductList_isCharacter_and_perm_invariant
% for IsChar, ip_ne_zero_of_common through character_decomposition_and_classFunction_expansion and
% centralIdempotentOfIrrCharacter_orthogonal_central, T.characterPairing_harishChandra_cuspidal_mackey
% for the permutation, and map_fst_eq with lab_eq, T.cuspidalCharacter_label_unique_GLn). The three
% Harish-Chandra statements are proved in the archive, and Theorem -labels-series follows those proofs
% down to Lemmas -labels-series-induction and -labels-series-mackey.
% Lean: (1) and (2) are the first two clauses of T.harishChandraSeriesLabel_exists_unique_GLn =
% GLq.Label.label_main (under 5 <= g and Green; its third clause, that for squarefree f the irreducible
% character with label f is unique and equals R, is not used here). About 4,200 lines of Lean proofs at
% depth 5 by the count used here. (3) is GLm.exists_irr_hasLabel: exists_cusp_list (strong induction on
% the degree, Polynomial.exists_monic_irreducible_factor, Green for each factor),
% T.harishChandraProductList_isCharacter_and_perm_invariant, T.harishChandraProductList_degree, and
% T.character_decomposition_and_classFunction_expansion; the Lean writes the pairing ip chi R and uses
% its symmetry for characters (ip_symm_of_inv).
\end{proof}

\begin{lemma}\label{asm-dependence-test-assembly-bare-norm-series}
Assume Green's parametrization (Corollary~\ref{lit-green-labelled}). Let $g\ge5$ be a prime with a
labelling of $\bar{\mathbf F}_g$, let $n\ge0$, and let $N>0$ be an integer divisible by
$|GL_n(\mathbf F_g)|$. For an irreducible character $\chi$ of $G_n=GL_n(\mathbf F_g)$ put
$z_\chi=\frac{\chi(1)}{|G_n|}\sum_{x\in G_n}\chi(x^{-1})\,x\in\mathbf C[G_n]$, and for
$f\in\mathbf F_g[X]$ let the \emph{series idempotent} $e(f)$ be the sum of the $z_\chi$ over the
irreducible characters $\chi$ with label $f$
(Lemma~\ref{asm-dependence-test-assembly-bare-norm-labels}). Then:
\begin{enumerate}
\item each $e(f)$ is a central idempotent of $\mathbf C[G_n]$, and $e(f)e(f')=0$ for $f\ne f'$;
\item $e(f)=0$ if and only if $f$ is not monic of degree $n$ with nonzero constant term;
\item the sum of the $e(f)$ over the monic $f$ of degree $n$ with nonzero constant term is $1$;
\item the coefficients of every $e(f)$ lie in $\mathbf Q(e^{2\pi i/N})$.
\end{enumerate}
\end{lemma}
\leanref{Development/AbelianVarietyTheory/Proofs/FlagMeasures/ZOfCoeffMem.lean}{GlobalTheory.GLm.isSer_eSerC}
\leanref{Development/AbelianVarietyTheory/Definitions/AutomorphicMembership/GoodChi.lean}{GlobalTheory.MemN.IsSer}
\leanref{Development/AbelianVarietyTheory/Definitions/NonGaloisDescent/LabelDatum.lean}{GlobalTheory.GLq.eSerC}
\leanref{Development/AbelianVarietyTheory/Proofs/InertiaEigenvalues/HarishChandraSeriesCentralIdempotentLaws.lean}{GlobalTheory.GLq.eSerC_laws}
\leanref{Development/AbelianVarietyTheory/Proofs/FlagMeasures/HarishChandraSeriesCentralIdempotentNeZero.lean}{GlobalTheory.GLm.eSerC_ne_zero}
\leanref{Development/AbelianVarietyTheory/Proofs/FlagMeasures/IsCharacterOfFDRepMemCyclotomicAdjoin.lean}{GlobalTheory.GLm.isChar_mem_adjoin_exp}
\begin{proof}
The $z_\chi$, for the finitely many irreducible characters $\chi$ of $G_n$, are central idempotents of
$\mathbf C[G_n]$ with $z_\chi z_{\chi'}=0$ for $\chi\ne\chi'$ and $\sum_\chi z_\chi=1$, and $z_\chi\ne0$
\cite[\S2.6, Theorem~8]{Ser77}. By Lemma~\ref{asm-dependence-test-assembly-bare-norm-labels}(1), the
sets of irreducible characters with label $f$, for distinct $f$, are disjoint and cover all of them.

(1) $e(f)$ is central as a sum of central elements; $e(f)e(f')$ is the sum of the $z_\chi z_{\chi'}$
over $\chi$ with label $f$ and $\chi'$ with label $f'$, which is $e(f)$ if $f=f'$ and $0$ otherwise.

(2) If $f$ is not monic of degree $n$ with nonzero constant term, no irreducible character has label
$f$, by Lemma~\ref{asm-dependence-test-assembly-bare-norm-labels}(2), and $e(f)=0$. Otherwise some
irreducible $\rho$ has label $f$, by part (3) of that lemma, and $z_\rho e(f)=z_\rho\ne0$.

(3) Every irreducible character has its label among these $f$, so the sum is $\sum_\chi z_\chi=1$.

(4) For $x\in G_n$, $x^N=1$, so the eigenvalues of $x$ in a representation are $N$-th roots of
unity, and every character value $\chi(x)$, a sum of them, lies in $\mathbf Q(e^{2\pi i/N})$. The
coefficient of $z_\chi$ at $x$ is $\chi(1)\chi(x^{-1})/|G_n|$.
% Lean: GLm.isSer_eSerC (IsSer r g N (eSerC L r)): T.harishChandraSeriesCentralIdempotent_laws =
% GLq.eSerC_laws (finite_irrCharacters_and_finsum_centralIdempotents_eq_one,
% centralIdempotentOfIrrCharacter_orthogonal_central, label_main, eSerC_eq_sum) for idem, central, orth,
% the vanishing and the sum; T.harishChandraSeriesCentralIdempotent_ne_zero = GLm.eSerC_ne_zero
% (exists_irr_hasLabel, zOf_ne_zero_of_isIrr) for the nonvanishing; and
% T.isCharacterOfFDRep_mem_cyclotomic_adjoin = GLm.isChar_mem_adjoin_exp with zOf_coeff_mem for the
% coefficients. About 2,700 lines of Lean proofs at depth 5 for eSerC_laws, most of them in label_main.
\end{proof}

\begin{lemma}\label{asm-dependence-test-assembly-bare-norm}
Assume Green's parametrization (Corollary~\ref{lit-green-labelled}). Let $B$ be base data with a
framing $f$, whose host has level primes $g_0,g_1$ and rank $r$, and let rows over columns over $B$ be
given whose Sha lengths, for $0<e<t$, are the Sha lengths of local Galois cohomology data of the host
field $E$, whose local defects $\eta_j(w)$ vanish for
$1\le j\le m$, and whose operators $e_{\mathrm{Op}}(\mu)$ and $e_{\mathrm{Mt}}(\mu,j)$ vanish. Then there
are norm data over these rows with no local places and no twisted places, whose group data are true
for $f$ (Definition~\ref{asm-dependence-test-assembly-true}(1)), whose sets $P_G(k)$ are the sets of
places of $F$ above $g_k$, $k=0,1$, and with a character $\omega$ of the centre of
$\mathbf C[G_0]$ such that the central operators of the types on the ambient and on the lifts are the
scalars $\omega(z)$.
\end{lemma}
\leanref{Development/AbelianVarietyTheory/Proofs/IwasawaDecomposition/TypeIdempotentCoeffAcrossLayers.lean}{GlobalTheory.Mem.Top.exists_nStage_grpTrue}
\leanref{Development/AbelianVarietyTheory/Proofs/IwasawaDecomposition/TypeIdempotentCoeffAcrossLayers.lean}{GlobalTheory.Mem.Top.trueLayer}
\leanref{Development/AbelianVarietyTheory/Proofs/IwasawaDecomposition/ThetaOfOfIsCharLab.lean}{GlobalTheory.Mem.TypePl.above}
\leanref{Development/AbelianVarietyTheory/Proofs/IwasawaDecomposition/ThetaOfOfIsCharLab.lean}{GlobalTheory.Mem.TypePl.existsUnique_div}
\leanref{Development/AbelianVarietyTheory/Proofs/IwasawaDecomposition/NcardPrimesOverOfStabilizer.lean}{GlobalTheory.HostSplit.split_of_one_mod}
\leanref{Development/AbelianVarietyTheory/Proofs/IwasawaDecomposition/ExistsHalfPlacesTypePrime.lean}{GlobalTheory.MemN.HalfPlaces.exists_places_cover}
\leanref{Development/AbelianVarietyTheory/Proofs/IwasawaDecomposition/GoodTupleConditionOfIdempotentZero.lean}{GlobalTheory.Mem.GaloisStage.nAut}
\leanref{Development/AbelianVarietyTheory/Proofs/FoundationProofs/ExistsTameGenAtOfValuationSubring.lean}{GlobalTheory.TameInertiaGen.exists_tameGenAt}
\leanref{Development/AbelianVarietyTheory/Proofs/IwasawaDecomposition/TypeIdempotentCoeffAcrossLayers.lean}{GlobalTheory.Mem.Top.five_le_type}
\leanref{Development/AbelianVarietyTheory/Proofs/FlagMeasures/ZOfCoeffMem.lean}{GlobalTheory.GLm.isSer_eSerC}
\leanref{Development/AbelianVarietyTheory/Proofs/IwasawaDecomposition/ExistsTypesIdempotentCoefficients.lean}{GlobalTheory.MemN.TypesInhab.types0}
\leanref{Development/AbelianVarietyTheory/Proofs/IwasawaDecomposition/ExistsTypesIdempotentCoefficients.lean}{GlobalTheory.MemN.TypesInhab.goodTupleIdempotent_central_and_action}
\leanref{Development/AbelianVarietyTheory/Proofs/IwasawaDecomposition/ExistsTypesIdempotentCoefficients.lean}{GlobalTheory.MemN.TypesInhab.exists_char}
\leanref{Development/AbelianVarietyTheory/Proofs/IwasawaDecomposition/ExistsTypesIdempotentCoefficients.lean}{GlobalTheory.MemN.TypesInhab.not_good_unip}
\leanref{Development/AbelianVarietyTheory/Proofs/IwasawaDecomposition/ExistsTypesIdempotentCoefficients.lean}{GlobalTheory.MemN.TypesInhab.goodTupleIdempotent_coeff_mem_subring}
\leanref{Development/AbelianVarietyTheory/Proofs/IwasawaDecomposition/ExistsExponentThetaData.lean}{GlobalTheory.MemN.exists_nexp_nth_of_isEmpty}
\begin{proof}
The construction has seven steps.

\emph{Step 1: the layer.} The \emph{true layer} consists of norm data on the rows with no local
places, whose set $P_G(k)$ is the set of the $f$ places of $F$ above $g_k$, whose labels are
$(X-1)^r$, and whose characters at the places of the group are the powers $\omega_k^{2(j+1)}$ of the
character $\omega_k$ of $(\mathbf Z/g_k)^\times$ that $f$ chooses; it is built from the local Galois
cohomology data whose Sha lengths are those of the rows, an algebraic closure $\Omega$ of $E$ with a
place, and the fact that $g_k$ splits completely in $E$. The last holds since $g_k\equiv1\pmod N$
with $N=Q_Sqd$, so that the Frobenius of every prime of $E$ above $g_k$ is
$\gamma^{-1}(\operatorname{art}(1))=1$ (Proposition~\ref{hf-host-field-properties}(c)). So $g_k$ has
$[E:\mathbf Q]$ places in $E$, permuted simply transitively by $\operatorname{Gal}(E/\mathbf Q)$.
Complex conjugation $c\ne1$ (Proposition~\ref{hf-host-field-properties}(a)) therefore moves every one
of them, and the places of $E$ above $g_k$ fall into conjugate pairs, one pair over each place of
$F$ above $g_k$; so there are $[E:\mathbf Q]/2=f$ of the latter. The layer also carries local data at
every prime $\ell$: Fontaine's rings of periods for $\ell$, with their Galois actions, Frobenius,
filtrations and period $t$, and the period data of the finite extensions of $\mathbf Q_\ell$. They
exist by Fontaine's construction
(Lemma~\ref{sha-automorphy-unramified-tangents-unit-crystalline-periods}).

\emph{Step 2: the places.} Let $\mathfrak p_k$ be the prime of $K$ above $g_k$ that $f$ chooses. The
prime $g_k$ splits in $K$ as $\mathfrak p_k\bar{\mathfrak p}_k$ with $\mathfrak p_k\ne\bar{\mathfrak
p}_k$, so a place $w$ of $E$ above $g_k$ lies over $\mathfrak p_k$ if and only if its conjugate $\bar
w$ does not. For $v\in P_G(k)$ the places of $E$ above $v$ are a place $w$ and $\bar w$, and exactly
one of them lies over $\mathfrak p_k$; call it $A(v)$. If $A(v)=A(v')$ then $v=v'$, since $A(v)$ lies
above $v$. If $A(v)$ and $A(v')$ were conjugate, they could not both lie over $\mathfrak p_k$. So the
places $A(v)$ are distinct and no two are complex conjugate. The automorphic data of the layer are
built on these places with unramified parameters only: a local parameter of dimension $d$ is a
multiset of $d$ nonzero complex numbers, its Frobenius eigenvalues, with inertia polynomial
$(X-1)^d$; since there are no local places, their remaining laws are direct checks. At each $A(v)$
choose a generator of tame inertia: an element $\gamma$ of the inertia group of the place of $\Omega$ chosen over $A(v)$ whose
closed subgroup together with wild inertia contains the inertia group. It exists since the tame
inertia group is procyclic \cite[Chapter~IV, \S2]{Ser79}.

\emph{Step 3: the series idempotents.} By
Lemma~\ref{asm-dependence-test-assembly-bare-norm-series}, the series idempotents of the label datum
of $g_k$ have the properties (1) to (4) of that lemma for $n=r$ and $N=|G_0|$, which make them series
idempotents of rank $r$ for $g_k$ and the order $|G_0|$; here $|GL_r(\mathbf F_{g_k})|$ divides
$|G_0|$, and $g_k\ge5$ since $g_k$ is a prime with $g_k\equiv1\pmod 4$, as $4\mid Q_S$. They give
types $T_0$ on the places of Step 2, with the following data. The abelian extension of $K$ of the
types is $K$ itself. Its discriminant is $-d$, so it is unramified outside the primes dividing
$Q_Sqd$; and $g_k$ splits in it, since $g_k\equiv1\pmod{Q_Sqd}$, $4\mid Q_S$ and $d\equiv7\pmod{16}$ give
$g_k\equiv1\pmod4$ and $g_k\equiv1\pmod d$, hence $\bigl(\frac{-d}{g_k}\bigr)=1$. The congruences
$g_k\equiv1\pmod{Q_Sqd}$ and the disjointness $\mathcal B_r(g_0)\cap\mathcal B_r(g_1)=\emptyset$ are
those of $B$. Every type is given the multiset $\{(r,1)\}$, which makes it cuspidal, and trivial
residual representations. The inertia elements at the places $A(v)$ are the tame generators of
Step 2. The Teichmüller character is a homomorphism $t\colon\bar{\mathbf F}_{g_k}\to\mathbf C$,
injective on nonzero elements, from an embedding in $\mathbf C$ of the fraction field of the Witt
vectors of $\bar{\mathbf F}_{g_k}$. At every place of the rows, with residue characteristic $\ell$,
the types carry a homomorphism from $\bar{\mathbf F}_{g_k}$ to $\bar{\mathbf F}_\ell$ that is injective
on the nonzero elements of order prime to $\ell$; a Teichmüller character followed by a reduction of
the algebraic integers modulo a prime above $\ell$ is one. The inertia polynomials $(X-1)^r$ of the layer are the products of the $X-t(x)$ over
the roots $x$ of the labels $(X-1)^r$, as these roots are $1$ and $t(1)=1$. The series idempotents
are those of the label data, the centre of $\mathbf C[G_0]$ acts on the ambient and on the lifts
through the augmentation, and there are no parameters.

\emph{Step 4: a character of the centre.} The good tuple idempotent $e_0$ of $T_0$ is central, and for
every tuple of monic polynomials of degree $r$ with nonzero constant term there is $t\ne0$ with
$e_0t=t$ if the tuple is good and $e_0t=0$ otherwise. The unipotent tuple $(X-1)^r$ is not good, since
$(X-1)^r$ is not squarefree for $r\ge2$. So $e_0t=0$ for some $t\ne0$; then $e_0$ is not a unit of the
centre, and a maximal ideal of the centre containing it gives a character $\omega$ with
$\omega(e_0)=0$.

\emph{Step 5: the types.} Retype $T_0$ with the label data and $\omega$; this keeps $e_0$. The
operators $e_{\mathrm{Op}}(\mu)$ and $e_{\mathrm{Mt}}(\mu,j)$ of the rows vanish, and so equal
$\omega(e_0)=0$: the types match the action of the rows.

\emph{Step 6: coefficients.} The coefficients of the series idempotents lie in the field generated
by a primitive $|G_0|$-th root of unity, which lies in the subfield $L_0$ of the coefficient field; so do those of
$e_0$, a sum of products of them.

\emph{Step 7: the rest.} There are no twisted places, and exponent data and theta data exist for
the empty set of twisted places, with the coefficients of Step 6. The remaining conditions on norm
data concern local or twisted places, of which there are none, except the condition on the local
defects, which holds since they vanish. The group data are true for $f$: the places $A(v)$ lie over
the primes that $f$ chooses, the Teichmüller characters and the characters at the places are those
of Step 1, and the series idempotents are those of Step 3.
% Step 1: locStage_inhabited is MemLoc.exists_locPrime_tied (a LocPrime, LocAinf and LocTied at every
% prime l, built from Fontaine's period rings; about 4,100 lines at depth 5 by the count used here),
% cited through Lemma sha-automorphy-unramified-tangents-unit-crystalline-periods, whose leanref it is.
% Step 3: GLm.isSer_eSerC is Lemma -bare-norm-series, which rests on Lemma -bare-norm-labels(1)
% (GLq.Label.label_main), written out down to Lemma -bare-norm-labels-series-induction. TypesInhab.types0 is a structure built field by field: Mray = K with
% GaloisStage.discr_K for Mray_unr and GaloisStage.split_K for g_Mray, g_cond and BB_disj from B,
% bc = {(r, 1)}, rbar = 1, tg from the tame generators, iota from GaloisStage.exists_iota
% (exists_teichC with nonempty_reduction), teichC from GaloisStage.exists_teichC, inertiaCharpoly_type
% from GaloisStage.rootsBar_unip, eSer with the fields of IsSer, cenAmb and cenMt = GaloisStage.cenTriv
% (the augmentation), Par empty, and the fields indexed by the empty Loc vacuous.
% Written out: emptyLayer's g_split is HostSplit.split_of_one_mod (Frobenius art(1) = 1, through
% H.der.frob); card_placesOver through HalfPlaces.exists_places_cover (Gal(E/Q) acts simply
% transitively, conj acts freely); Step 2 is TypePl.existsUnique_div (Chr.splitsC_main for w | pb k
% iff not conj w | pb k, under_eq_iff), above_injective, not_isConjPl, div_over, and GaloisStage.nAut
% (141 lines of definitional fields: UnrPar parameters, inePol = (X-1)^card, the Loc-indexed fields
% empty); TameInertiaGen.exists_tameGenAt (about 550 lines) is cited to Serre; five_le_type from
% g mod QS q d = 1 and 4 | QS.
% Lean: Mem.Top.exists_nStage_grpTrue. Step 1: trueLayer (emptyLayer with omLab Om k ^ (2(j+1)), PlG
% = placesOver F (g k), lab = (X-1)^r; about 710 lines at depth 5) and locStage_inhabited (about
% 2,700). Step 2: GaloisStage.nAut with TypePl.above f.P, div_over, above_injective, not_isConjPl
% (about 5,300); TameInertiaGen.exists_tameGenAt (about 550). Step 3: GLm.isSer_eSerC with
% five_le_type (about 1,300) and TypesInhab.types0 with B.g_cond and B.BB_disj (about 2,100). Step 4:
% goodTupleIdempotent_central_and_action, not_good_unip, exists_char. Step 5: retype, htie. Step 6:
% goodTupleIdempotent_coeff_mem_subring with (he k).rat and adjoin_le. Step 7: GaloisStage.nTw,
% exists_nexp_nth_of_isEmpty, the fields indexed by the empty Loc and Tw, and etaJ_tie from heta. In
% all about 7,700 lines of Lean proofs at depth 5.
\end{proof}

\begin{lemma}\label{asm-dependence-test-assembly-bare}
Assume Green's parametrization (Corollary~\ref{lit-green-labelled}). Let $B$ be base data with a
framing $f$, let $Z$ be a later stage of a column choice with columns over $B$, and let $\eta$ be
real. There is a later stage $Z'$ of the same column choice whose group data are true for $f$
(Definition~\ref{asm-dependence-test-assembly-true}), whose division data have $\eta_F=\eta$, and whose
rows have vanishing indices: $I_j(\mu_w,w)+I_j(\mu_w,\bar w)+j_w(\mu_w)=0$ for every lift $j$ and
every place $w$.
\end{lemma}
\leanref{Development/AbelianVarietyTheory/Proofs/IwasawaDecomposition/ExistsBare.lean}{GlobalTheory.Mem.R3.LedRec.exists_bare}
\leanref{Development/AbelianVarietyTheory/Proofs/IwasawaDecomposition/TypeIdempotentCoeffAcrossLayers.lean}{GlobalTheory.Mem.Top.exists_nStage_grpTrue}
\leanref{Development/AbelianVarietyTheory/Proofs/IwasawaDecomposition/ExistsBare.lean}{GlobalTheory.Mem.R3.LedRec.bare}
\leanref{Development/AbelianVarietyTheory/Proofs/IwasawaDecomposition/GoodTupleConditionOfIdempotentZero.lean}{GlobalTheory.Mem.GaloisStage.adapt}
\leanref{Development/AbelianVarietyTheory/Proofs/IwasawaDecomposition/ExistsNonnegShaOK.lean}{GlobalTheory.MemN.shaOK_shaTrue}
\begin{proof}
Take the rows of $Z$ with their indices $I_j$ and $j_w$, their local defects $\eta_j(w)$ and their
operators $e_{\mathrm{Op}}$ and $e_{\mathrm{Mt}}$ replaced by $0$, and their Sha lengths replaced by
those of the local Galois cohomology data of the host field; their indices vanish. Take the division
data of $Z$ with $\eta_F$ replaced by $\eta$, and the size and jet data of $Z$. These rows satisfy the
hypotheses of Lemma~\ref{asm-dependence-test-assembly-bare-norm}, which gives norm data over them
whose group data are true for $f$. With these norm data the rows, division, size and jet data form the
later stage $Z'$.
% Lean: LedRec.exists_bare: bare a.W = GaloisStage.adapt a.W with Ij = jw = 0 (adapt sets shaLen to
% shaTrue and etaJ, eOp, eMt to 0), shaOK_shaTrue, and Mem.Top.exists_nStage_grpTrue (Lemma
% asm-dependence-test-assembly-bare-norm).
\end{proof}

\begin{lemma}\label{asm-dependence-test-assembly-copy}
Let $Z$ and $Z'$ be later stages of the same column choice that have the same group
(Definition~\ref{asm-dependence-test-assembly-true}), and let $z$ be a realization of $D(w,Z)$ built
at the stage. Its data, with the isomorphism of groups, form a realization $z'$ of $D(w,Z')$ built at
the stage, its \emph{copy}. If $z$ exists at every place, so does $z'$. If the cuspidal
multiplicities of $z'$ are equal, so are those of $z$; and if the complex data of $z'$ at its face
pass the row test, so do those of $z$.
\end{lemma}
\leanref{Development/AbelianVarietyTheory/Proofs/IwasawaDecomposition/CopyLIntegral.lean}{GlobalTheory.Mem.R3.DatCopy.copy}
\leanref{Development/AbelianVarietyTheory/Proofs/IwasawaDecomposition/CopyLIntegral.lean}{GlobalTheory.Mem.R3.DatCopy.copy_exAt}
\leanref{Development/AbelianVarietyTheory/Proofs/TwistedOperators/EqualMultiplicitiesRowTestOfSameGroupTransferStatement.lean}{GlobalTheory.T.equalMultiplicities_rowTest_of_sameGroupTransfer}
\leanref{Development/AbelianVarietyTheory/Proofs/TwistedOperators/EqualMultiplicitiesRowTestOfSameGroupTransfer.lean}{GlobalTheory.Twist.multEq_rowTest_of_copy}
\leanref{Development/AbelianVarietyTheory/Proofs/IwasawaDecomposition/CopyLIntegral.lean}{GlobalTheory.Mem.R3.DatCopy.twice_eq}
\leanref{Development/AbelianVarietyTheory/Proofs/IwasawaDecomposition/CopyLIntegral.lean}{GlobalTheory.Mem.R3.DatCopy.multEq_tr}
\leanref{Development/AbelianVarietyTheory/Proofs/IwasawaDecomposition/CopyLIntegral.lean}{GlobalTheory.Mem.R3.DatCopy.rowTest_tr}
\leanref{Development/AbelianVarietyTheory/Definitions/AutomorphicMembership/PlaceModelImp.lean}{GlobalTheory.Mem.R3.Tr.mult_eq}
\leanref{Development/AbelianVarietyTheory/Definitions/AutomorphicMembership/Wp3ApplyEqSqMulSubSqMul.lean}{GlobalTheory.Mem.R3.Tr.toCx3_eq}
\begin{proof}
The definite-side data $D(w,Z)$ and $D(w,Z')$ have the same base data and columns, so the same
host, spaces, coefficient field and models of discrete series; the multiplicity spaces of a model
for the two data are the same. The copy $z'$ keeps every datum of $z$: its level and character
data, models, Hecke modules and their maps, measures, Weil data, theta maps, scalars and face. Only
two things are moved: the isomorphism of its group with that of the datum is composed with the
isomorphism $G_0(Z)\cong G_0(Z')$, and the maps $\mathrm{io}$ out of the spaces of theta
multiplicities are carried along the identification of the multiplicity spaces.

Copying from $Z'$ back to $Z$, along the same group in the other direction, sends $z'$ to $z$: the
maps $\mathrm{io}$, carried along the identification and back, are those of $z$, and nothing else
was moved.

Copying preserves each of the three properties. A realization exists at every place if its iterated
theta maps, its maps $\mathrm{io}$, its theta lifts and its theta maps exist; these are the maps of
$z$, so $z'$ exists at every place if $z$ does. The condition that cuspidal multiplicities are equal
compares the cuspidal multiplicities of the models of $U(Y_0)$ and of the $U(W_i)$ at the
representations that are good for the level data; it reads the models, the level data and the
group of the realization, which the copy keeps, and multiplicity spaces, which are identified. The
complex data of a realization at its face are built from its theta lifts, its Hecke modules and the
parts and levels of the face, which the copy keeps, so the complex data of a realization at its
face and those of its copy at the copied face are the same, and one passes the row test exactly when
the other does.

Apply this to $z'$ and its copy at $Z$, which is $z$. If the cuspidal multiplicities of $z'$ are
equal, so are those of $z$, and if the complex data of $z'$ pass the row test, so do those of $z$.
% Lean: DatCopy.copy (real3 b z.R, cG.trans s.iso, ioYB and ioXB cast along multT_eq) and
% DatCopy.copy_exAt = Tr.exAt_tr. Twist.multEq_rowTest_of_copy: twice_eq (R = twice b R, through
% twiceY_eq and twiceX_eq), multEq_tr (through Tr.mult_eq: the multiplicity maps of b.D are those of
% a.D at the transported model ds) and rowTest_tr (through Tr.toCx3_eq: the complex data agree), both
% applied with the roles of a and b exchanged. About 250 lines of Lean proofs at depth 5.
\end{proof}

\begin{lemma}\label{asm-dependence-test-assembly-second-layer}
Assume Green's parametrization (Corollary~\ref{lit-green-labelled}), rapid decay of cusp forms
(Theorem~\ref{lit-bundle-theta}(3)), the seven statements on test choices and the number field
statement (Definition~\ref{asm-dependence-test-assembly-statements}). Then the following holds at
all large stages after the column choices that are pinned
(Definition~\ref{asm-stage-construction}(2)). Let $Z$ be a later stage whose group data are true for
the framing (Definition~\ref{asm-dependence-test-assembly-true}), and $z$ a realization of $D(w,Z)$
built at the stage. Suppose that the group data of the stage are recorded, and that $z$ exists at
every place and lies in the geometry of the stage. Then the complex data of $z$ at its face pass the
row test.
\end{lemma}
\leanref{Development/AbelianVarietyTheory/Proofs/IwasawaDecomposition/AtDatRowTest.lean}{GlobalTheory.Mem.R3.Rw.RowAt.atDat_rowTest}
\leanref{Development/AbelianVarietyTheory/Proofs/TwistedOperators/AtDatRowTestOfExField.lean}{GlobalTheory.Twist.atDat_rowTest_of_exField}
\begin{proof}
Let $C_{F_1}$ and $C_{F_2}$ be the constants of
Lemmas~\ref{asm-dependence-test-assembly-realization} and~\ref{asm-dependence-test-assembly-existence},
and put $C_F=\max(C_{F_1},C_{F_2})$.

\emph{Step 1: the eventual properties.} By Lemma~\ref{asm-dependence-test-assembly-existence} for
$C_F$, at large base data in every context beyond a threshold, every realization built at a stage
that is in an automorphic setting with size data at the scale of the context has test choices for
$C_F$ that satisfy the evaluation conditions and the field conditions. The following three hold
eventually in the context, hence together: Lemma~\ref{asm-dependence-test-assembly-realization} for
$C_F$ at the exponent $\varepsilon$ of the context, with its $x_0$ at most the scale $x$ of the
context; Lemma~\ref{asm-dependence-test-assembly-sizes}; and
Lemma~\ref{asm-dependence-test-assembly-setting}.

\emph{Step 2: a bare stage.} Let a context with these properties be given, with the constant $C_P$
of Lemma~\ref{asm-dependence-test-assembly-sizes}, and base data $B$ lying over the context, with
framing $f$ and twist data $T$, that are large for it, let $w$ be a pinned column choice, and let $Z$ and $z$ be as in the
statement. By Lemma~\ref{asm-dependence-test-assembly-bare} there is a later stage $Z'$ of the column
choice $w$ whose group data are true for $f$, whose division data have $\eta_F=6x^{-C_P}$, and whose
rows have vanishing indices. By Lemma~\ref{asm-dependence-test-assembly-sizes}, the datum $D(w,Z')$
has size data at the scale of the context whose rows have the sizes of the size data.

\emph{Step 3: the copy.} The group data of $Z$ and $Z'$ are both true for $f$, so $Z$ and $Z'$ have
the same group. By Lemma~\ref{asm-dependence-test-assembly-copy}, it suffices that the complex data of
the copy $z'$ of $z$ at $Z'$ pass the row test, and $z'$ exists at every place.

\emph{Step 4: the row test for the copy.} By Lemma~\ref{asm-dependence-test-assembly-setting}, $z'$
with the size data of Step 2 is a realization in an automorphic setting at its face. By Step 1 it has
test choices for $C_F$ that satisfy the evaluation conditions and the field conditions, and
Lemma~\ref{asm-dependence-test-assembly-realization} applies to it, since its size data are at the
scale of the context. So the complex data of $z'$, and hence those of $z$, pass the row test. The
hypotheses that the group data of the stage are recorded and that $z$ lies in the geometry of the
stage are not used.
% Lean: Rw.RowAt.atDat_rowTest: rowTest_at_realisation_of_hilbertModels (CF1),
% testChoices_exist_largeCoeffField_adelicData (CF2) at max CF1 CF2 (Step.of_evLgG), the eventual
% conjunction of H1, eventually_stages_parameterSizeConditions and hidaSetting_of_sizeRows
% (Step.of_evG), LedRec.exists_bare with eta = 6 x^(-K.C.P.CF), LedRec.sameGrp,
% equalMultiplicities_rowTest_of_sameGroupTransfer and DatCopy.copy_exAt. The hypotheses GrpAt and
% InGeo of the property are not used.
\end{proof}

\begin{lemma}\label{asm-dependence-test-assembly}
Assume the results of Theorem~\ref{lit-bundle-theta}, Green's parametrization
(Corollary~\ref{lit-green-labelled}), the join statements and the functional equation of
Theorem~\ref{asm-dependence-inputs}, the theta hypotheses (Definition~\ref{th-theta-hypotheses}),
and that functionals separate (Definition~\ref{asm-grid-evaluation}(7)). Then the following holds at
all large stages after the column choices that are pinned
(Definition~\ref{asm-stage-construction}(2)). Let $Z$ be a later stage and $z$ a realization of
$D(w,Z)$ built at the stage (Definition~\ref{is-realisation}). Suppose that the group data of $Z$ are
true for the framing (Definition~\ref{asm-dependence-test-assembly-true}), that the group data of the
stage are recorded (its places are those of the
group of $D(w,Z)$, and the type idempotent of $D(w,Z)$ is a central idempotent of the group
algebra), and that $z$ exists at every place and lies in the geometry of the stage. Then the complex
data of $z$ at its face pass the row test (Definition~\ref{asm-stage-construction}(3)).
\end{lemma}
\leanref{Development/AutomorphicArithmetic/ThetaCorrespondence/ThetaPairLifts/LocalTypeDictionary/TestVectorAtDatumOfJoin.lean}{GlobalTheory.RowTestAssembly.test_vector_at_datum_of_join}
\leanref{Development/AbelianVarietyTheory/Proofs/IwasawaDecomposition/HoldsUIff.lean}{GlobalTheory.Mem.Top.Bl.AtDat}
\leanref{Development/AbelianVarietyTheory/Definitions/AutomorphicMembership/IsAbsIrrCharCompIff.lean}{GlobalTheory.Mem.R3.GrpAt}
\leanref{Development/AbelianVarietyTheory/Proofs/TwistedOperators/AtDatRowTestOfExField.lean}{GlobalTheory.Twist.atDat_rowTest_of_exField}
\leanref{Development/AbelianVarietyTheory/Proofs/IwasawaDecomposition/AtDatRowTest.lean}{GlobalTheory.Mem.R3.Rw.RowAt.atDat_rowTest}
\leanref{Development/AbelianVarietyTheory/Proofs/TwistedOperators/RowTestAtRealisationOfHilbertModelsStatement.lean}{GlobalTheory.T.rowTest_at_realisation_of_hilbertModels}
\leanref{Development/AbelianVarietyTheory/Proofs/TwistedOperators/TestChoicesExistLargeCoeffFieldAdelicDataStatement.lean}{GlobalTheory.T.testChoices_exist_largeCoeffField_adelicData}
\leanref{Development/AbelianVarietyTheory/Proofs/TwistedOperators/EventuallyStagesParametersParameterSizeConditionsStatement.lean}{GlobalTheory.T.eventually_stages_parameterSizeConditions}
\leanref{Development/AbelianVarietyTheory/Proofs/TwistedOperators/RowTestHypothesesAtMemberAdelicDataStatement.lean}{GlobalTheory.T.hidaSetting_of_sizeRows}
\leanref{Development/AbelianVarietyTheory/Proofs/IwasawaDecomposition/ExistsBare.lean}{GlobalTheory.Mem.R3.LedRec.exists_bare}
\leanref{Development/AbelianVarietyTheory/Proofs/TwistedOperators/EqualMultiplicitiesRowTestOfSameGroupTransferStatement.lean}{GlobalTheory.T.equalMultiplicities_rowTest_of_sameGroupTransfer}
\begin{proof}
Alpöge's proof has two layers. By Lemma~\ref{asm-dependence-test-assembly-first-layer}, the inputs
give the seven statements on test choices and the number field statement
(Definition~\ref{asm-dependence-test-assembly-statements}). Rapid decay of cusp forms is among the
results of Theorem~\ref{lit-bundle-theta}. So
Lemma~\ref{asm-dependence-test-assembly-second-layer} applies, with Green's parametrization, and its
conclusion is that of the lemma.
% Lean: test_vector_at_datum_of_join supplies LinK, InjK, ConstK, SumOKK, Formula2K, SumRatK,
% RatioK and ExField (Lemma asm-dependence-test-assembly-first-layer) and P.cuspRapidDecay to
% Twist.atDat_rowTest_of_exField = Rw.RowAt.atDat_rowTest (Lemma asm-dependence-test-assembly-second-layer).
\end{proof}

\subsection{Existence of a collection of hypotheses}\label{subsection-asm-hypotheses-exist}

\begin{definition}\label{asm-dependence-grid-agree}
Let $D$ be the definite-side datum of a later stage over base data $B$, with host field $E/F$, and
let $R$ be a realization of $D$ (Definition~\ref{is-realisation}) whose level system
(Definition~\ref{th-grid}) has tuples $t$ with Hecke characters $\xi_t$, finite sets $Q_t$ of places
of $E$ and base tuple $\mu_0$. Let $G$ be an evaluation grid for $R$ with constant $C_F$ at $b$
(Definition~\ref{asm-grid-evaluation}(1)) and $W$ a Weil datum over $R$ for $G$.
\begin{enumerate}
\item The \emph{key} of $R$ is the set of the pairs $(\xi_t,Q_t)$. The \emph{grid key} of $G$ is
  $(v_0,w(v_0),\chi_1,\dots,\chi_m)$, its base place, the place above it and its characters.
\item A \emph{grid choice} for $W$ is a grid choice for $W$ on the set of all tuples
  (Definition~\ref{asm-dependence-test-assembly-first-layer-grid-pieces-choice-grid}(3)): adelic bases
  of the lattices of the grid spaces, adapted at the local places of $G$ to the row and top
  decompositions, standard data with a vacuum normalization
  (Definitions~\ref{asm-dependence-test-assembly-first-layer-grid-pieces-choice-standard}
  and~\ref{asm-dependence-test-assembly-first-layer-grid-pieces-choice-normal}), and a finite set of
  places of $E$ at which the ideal $\mathfrak A$ of the host or $\bar{\mathfrak A}$ has nonzero
  exponent, with exactly one place above each place of $F$ below a place where $\mathfrak A$ has. The
  further choices of a full evaluation choice with Weil datum $W$
  (Definition~\ref{asm-stage-construction}(5)) contain a grid choice for $W$, which they
  \emph{extend}, and the \emph{evaluation conditions} are conditions on all the choices.
\item Let $\sigma\colon M\to\mathbf C$ be an embedding of a field and $S$ a grid choice for $W$. The
  \emph{numbers} of $(G,W,S)$ lie in $M$ if $\sigma(M)$ contains the values of the characters of
  the level system, of its columns and of $G$, the roots of unity of order $8$ and of order $N(w)^{d(w)}$ for the
  local places $w$ of $G$ and their depths $d(w)$, the square roots of the norms of these places and
  of the places of $S$, the phases and the jets of $S$, the ratios of the $L$-values of $G$ at
  $\mu_0$, and the sums at the bottom of $W$. The \emph{field conditions} of a full evaluation
  choice say that the numbers of its grid, Weil datum and grid choice lie in the coefficient field
  of $D$. The \emph{square numbers} of $G$ are the norms at its bottom and the ratios of its toric
  periods at $\mu_0$.
\item Let $R'$ be a realization of the datum of another later stage over $B$, with level system
  $(\xi'_t,Q'_t)$ and base tuple $\mu'_0$, and let $\iota$ be a map from the tuples of $R'$ to those of
  $R$. The level systems of $R$ and $R'$ \emph{agree through} $\iota$ if the columns of both are given
  by the same framing of $B$, $\iota(\mu'_0)=\mu_0$, $\xi'_t=\xi_{\iota(t)}$ and $Q'_t=Q_{\iota(t)}$ for
  every tuple $t$, and the adelic frames of $X'$ of both are the identity. Evaluation grids $G$ for
  $R$ and $G'$ for $R'$ \emph{agree} if they have the same base place, the same places $w(v)$, the same
  characters, the same spaces $V_{i,j}$ for $i\ge1$ and top spaces with the same lattices, and the
  same row and top decompositions. Weil data $W$ for $G$ and $W'$ for $G'$ \emph{agree through}
  $\iota$ if they have the same standard character, archimedean family and Cartan involutions, and the
  Weil datum of $W'$ at every pair is that of $W$ at its image under $\iota$. Grid choices $S$ for $W$
  and $S'$ for $W'$ \emph{agree through} $\iota$ if they have the same adelic bases $\beta_{i,j}$
  and $\beta^+_j$, the same matrices $\phi_X$ and elements $u_X$ of their standard data at the
  definite spaces $X$, the same number $\delta$, the same vacuum vectors $g$ and constants $c$ of
  their vacuum normalizations at corresponding definite pairs, and the same set of places.
\end{enumerate}
\end{definition}
\leanref{Development/AbelianVarietyTheory/Definitions/TwistedInnerProducts/TwistPlacesOffPrimeQZero.lean}{GlobalTheory.RowEx.keyOf}
\leanref{Development/AbelianVarietyTheory/Definitions/AutomorphicMembership/GKeyFin.lean}{GlobalTheory.Mem.LiteralAutomorphy.gkeyOf}
\leanref{Development/AbelianVarietyTheory/Definitions/Divisors/InvVolConjLevelOne.lean}{GlobalTheory.FunObj.GridChoice}
\leanref{Development/AbelianVarietyTheory/Definitions/TwistedInnerProducts/ExistsMatchOver.lean}{GlobalTheory.RowTest.Choice}
\leanref{Development/AbelianVarietyTheory/Definitions/TwistedInnerProducts/ExistsMatchOver.lean}{GlobalTheory.RowTest.Choice.Is}
\leanref{Development/AbelianVarietyTheory/Definitions/TwistedInnerProducts/MemHalfCosetSelf.lean}{GlobalTheory.RowTest.NumbersIn}
\leanref{Development/AbelianVarietyTheory/Definitions/TwistedInnerProducts/MemHalfCosetSelf.lean}{GlobalTheory.RowTest.Choice.FieldOK}
\leanref{Development/AbelianVarietyTheory/Definitions/TwistedInnerProducts/TwistPlacesOffPrimeQZero.lean}{GlobalTheory.RowEx.AllIn}
\leanref{Development/AbelianVarietyTheory/Definitions/Divisors/SlopeTestRemainingChoicesAlt.lean}{GlobalTheory.BigTop.Fun.SqIn}
\leanref{Development/AbelianVarietyTheory/Definitions/TwistedInnerProducts/TwistPlacesOffPrimeQZero.lean}{GlobalTheory.RowEx.LAgree}
\leanref{Development/AbelianVarietyTheory/Definitions/TwistedInnerProducts/TwistPlacesOffPrimeQZero.lean}{GlobalTheory.RowEx.GAgree}
\leanref{Development/AbelianVarietyTheory/Definitions/TwistedInnerProducts/TwistPlacesOffPrimeQZero.lean}{GlobalTheory.RowEx.WAgree}
\leanref{Development/AbelianVarietyTheory/Definitions/TwistedInnerProducts/TwistPlacesOffPrimeQZero.lean}{GlobalTheory.RowEx.SAgree}
\leanref{Development/AbelianVarietyTheory/Definitions/TwistedInnerProducts/TwistPlacesOffPrimeQZero.lean}{GlobalTheory.RowEx.Agree}

\begin{lemma}\label{asm-dependence-grid-keys}
\begin{enumerate}
\item Let $B$ be base data with host field $E$ and $\Xi$ a finite set of Hecke characters of $E$.
  There are finitely many sets that are the key of the level system, with base tuple, of a
  realization of the datum of a later stage over $B$ whose characters $\xi_t$ lie in $\Xi$.
\item Let $B$ be base data and $c$ real. There are finitely many grid keys of theta grids, with
  first constant $c$, for level systems of the data of later stages over $B$.
\end{enumerate}
\end{lemma}
\leanref{Development/AbelianVarietyTheory/Definitions/TwistedInnerProducts/TwistPlacesOffPrimeQZero.lean}{GlobalTheory.RowEx.KeyFin}
\leanref{Development/AbelianVarietyTheory/Proofs/TwistedOperators/KeyFinHolds.lean}{GlobalTheory.RowEx.keyFin_holds}
\leanref{Development/AbelianVarietyTheory/Definitions/AutomorphicMembership/GKeyFin.lean}{GlobalTheory.Mem.LiteralAutomorphy.GKeyFin}
\leanref{Development/AbelianVarietyTheory/Proofs/IwasawaDecomposition/KeyFinitenessGlobal.lean}{GlobalTheory.Mem.LiteralAutomorphy.gKeyFin}
\leanref{Development/AbelianVarietyTheory/Proofs/IwasawaDecomposition/KeyFinitenessGlobal.lean}{GlobalTheory.Mem.LiteralAutomorphy.KeyFiniteness.KS_finite}
\leanref{Development/AbelianVarietyTheory/Proofs/IwasawaDecomposition/KeyFinitenessGlobal.lean}{GlobalTheory.Mem.LiteralAutomorphy.heckeFin}
\begin{proof}
(1) Let $U$ be the set of pairs $(\xi,Q)$ with $\xi\in\Xi$ ramified at only finitely many places and
$Q$ a set of places at which $\xi$ is ramified. For each $\xi$ there are finitely many such $Q$, so
$U$ is finite, and it suffices to show that every such key is a subset of $U$. Let $t$ be a tuple.
Then $\xi_t\in\Xi$, and in a level system with base tuple the set of ramified places of $\xi_t$ is
$Q_t\cup\bar Q_t$. So $\xi_t$ is ramified at the places of $Q_t$ and at only finitely many places,
and $(\xi_t,Q_t)\in U$.

(2) The base place $v_0$ of a grid lies above the auxiliary prime $q_0$, and $w(v_0)$ lies above
$v_0$; there are finitely many such places. For $1\le a\le m$, the character $\chi_a$ has the
infinity type $2c(\tau)$ given by the centres of the grid at $a$, is unramified outside the places of
$E$ above $v_0$ or above the places of $S_1$, and has level at most $2$ at those places. For a given
$v_0$ there are finitely many such characters, since there are finitely many Hecke characters of a
number field with a given infinity type, unramified outside a given finite set of places and of
bounded level there: they differ by characters of a ray class group, which is finite. These
conditions depend only on $B$, $c$ and $v_0$.
\end{proof}

\begin{lemma}\label{asm-dependence-grid-transport}
Assume the existence of the Weil representations of unitary dual pairs, which is among the results
of Theorem~\ref{lit-bundle-theta}. Let $R$ and $R'$ be realizations whose level systems agree
through $\iota$ (Definition~\ref{asm-dependence-grid-agree}(4)), with base tuples $\mu_0$ and $\mu'_0$.
\begin{enumerate}
\item Every theta grid $G$ for the level system of $R$ has an agreeing theta grid $G'$ for that of
  $R'$, with the same constants.
\item If $R'$ satisfies the chain and model conditions of a core, then for agreeing grids $G$ and
  $G'$, every grid Weil datum $W$ for $G$ has a Weil datum $W'$ over $R'$ for $G'$ that agrees with
  it through $\iota$.
\item For agreeing grids and Weil data, every grid choice $S$ for $W$ has a grid choice $S'$ for $W'$
  that agrees with it through $\iota$.
\item Let $(G,W,S)$ and $(G',W',S')$ agree through $\iota$, and let $\sigma\colon M\to\mathbf C$ be
  the inclusion of a subfield. If the numbers of $(G,W,S)$ at $\mu_0$ lie in $M$, so do those of
  $(G',W',S')$ at $\mu'_0$; and if the square numbers of $G$ at $\mu_0$ lie in $M$, so do those of
  $G'$ at $\mu'_0$.
\end{enumerate}
\end{lemma}
\leanref{Development/AbelianVarietyTheory/Definitions/TwistedInnerProducts/TwistPlacesOffPrimeQZero.lean}{GlobalTheory.RowEx.CarryG}
\leanref{Development/AbelianVarietyTheory/Definitions/TwistedInnerProducts/TwistPlacesOffPrimeQZero.lean}{GlobalTheory.RowEx.CarryW}
\leanref{Development/AbelianVarietyTheory/Definitions/TwistedInnerProducts/TwistPlacesOffPrimeQZero.lean}{GlobalTheory.RowEx.CarryS}
\leanref{Development/AbelianVarietyTheory/Definitions/TwistedInnerProducts/TwistPlacesOffPrimeQZero.lean}{GlobalTheory.RowEx.TrEasy}
\leanref{Development/AbelianVarietyTheory/Definitions/TwistedInnerProducts/TwistPlacesOffPrimeQZero.lean}{GlobalTheory.RowEx.TrPhase}
\leanref{Development/AbelianVarietyTheory/Definitions/TwistedInnerProducts/TwistPlacesOffPrimeQZero.lean}{GlobalTheory.RowEx.TrJet}
\leanref{Development/AbelianVarietyTheory/Definitions/TwistedInnerProducts/TwistPlacesOffPrimeQZero.lean}{GlobalTheory.RowEx.TrBot}
\leanref{Development/AbelianVarietyTheory/Definitions/AutomorphicMembership/GKeyFin.lean}{GlobalTheory.Mem.LiteralAutomorphy.TrSq}
\leanref{Development/AbelianVarietyTheory/Proofs/TwistedOperators/TwoLeMOf.lean}{GlobalTheory.RowEx.CarryGrid.carryG}
\leanref{Development/AbelianVarietyTheory/Proofs/TwistedOperators/CarryWeilUnitaryPair.lean}{GlobalTheory.RowEx.carryW_of_weil}
\leanref{Development/AbelianVarietyTheory/Proofs/TwistedOperators/CarrySecond.lean}{GlobalTheory.RowEx.carryS}
\leanref{Development/AbelianVarietyTheory/Proofs/TwistedOperators/TrEasy.lean}{GlobalTheory.RowEx.trEasy}
\leanref{Development/AbelianVarietyTheory/Proofs/TwistedOperators/TrPhaseOfTrFlip.lean}{GlobalTheory.RowEx.PhaseCarry.trPhase_of_trFlip}
\leanref{Development/AbelianVarietyTheory/Proofs/TwistedOperators/TrJetOfTrFlip.lean}{GlobalTheory.RowEx.Jet.trJet_of_trFlip}
\leanref{Development/AbelianVarietyTheory/Proofs/TwistedOperators/TransferFlip.lean}{GlobalTheory.RowEx.Carry.flippedTensorsTransferToGrid_proof}
\leanref{Development/AbelianVarietyTheory/Proofs/TwistedOperators/TransferBottom.lean}{GlobalTheory.RowEx.trBot_holds}
\leanref{Development/AbelianVarietyTheory/Proofs/IwasawaDecomposition/TraceSquare.lean}{GlobalTheory.Mem.LiteralAutomorphy.trSq_holds}

The proof is not yet written out; it is part of the open step in the proof of
Theorem~\ref{asm-dependence-inputs}. In (2), $W'$ takes its Weil data at the pairs of $R'$ from $R'$
and is matched with $W$ by the uniqueness of the standard character and of the archimedean family;
in (4), the phases and the jets carry over because the local coset sets at a local place do not
depend on the base lattice.

\begin{lemma}\label{asm-dependence-grid-pieces}
Assume the results of Theorem~\ref{lit-bundle-theta}. In (2) assume also the theta hypotheses
(Definition~\ref{th-theta-hypotheses}), that functionals separate
(Definition~\ref{asm-grid-evaluation}(7)), the convergence of the theta integrals of the twisted Weil
representations, the matching of the twisted local factors and the agreement of the twisted square
theta integrals; in (3) and (4) assume also the join statements and the functional equation
(Theorem~\ref{asm-dependence-inputs}). Each of the following holds with some constant $C_{F_0}$ in
place of $C$: for every $\varepsilon$ and every $C_F\ge C$ there is $x_0$ such that the statement holds
for every definite-side datum $D$ with size data of exponent $\varepsilon$ at a scale $x\ge x_0$, with
$b=e^{x^{\varepsilon}}$.
\begin{enumerate}
\item For every realization $R$ of $D$ in an automorphic setting, with these size data at a face
  (Definition~\ref{asm-dependence-functional-witness-setting}), and every evaluation grid $G$ for
  $R$ with constant $C_F$ at $b$, there are a Weil datum $W$ over $R$ for $G$ and a grid choice for
  $W$.
\item For every realization $R$ of $D$ in an automorphic setting, every evaluation grid $G$ for $R$
  with constant $C_F$ at $b$, every Weil datum $W$ over $R$ for $G$ and every grid choice $S_0$ for
  $W$, there are further choices with Weil datum $W$ that extend $S_0$ and satisfy the evaluation
  conditions.
\item For every realization $R$ of $D$ in an automorphic setting and every choice of an evaluation
  grid $G$ for $R$ with constant $C_F$ at $b$, a Weil datum $W$ over $R$ and further choices that
  satisfy the evaluation conditions, with grid choice $S$, there is a number field $M\subset\mathbf C$
  in which the numbers of $(G,W,S)$ at the base tuple lie.
\item If the size data of $D$ satisfy the recursion conditions, then for every level system for $D$
  with base tuple $\mu_0$ and every theta grid $G$ for it with constants $(1,C_F,b)$, there is a
  number field $M\subset\mathbf C$ in which the square numbers of $G$ at $\mu_0$ lie.
\end{enumerate}
\end{lemma}
\leanref{Development/AbelianVarietyTheory/Definitions/AutomorphicMembership/GKeyFin.lean}{GlobalTheory.Mem.LiteralAutomorphy.WeilEx}
\leanref{Development/AutomorphicArithmetic/ThetaCorrespondence/ThetaPairLifts/LocalTypeDictionary/ExistsWeilDeligneOfLocalInputs.lean}{GlobalTheory.Mem.LiteralAutomorphy.exists_weilDeligne_of_localInputs}
\leanref{Development/AbelianVarietyTheory/Definitions/TwistedInnerProducts/TwistPlacesOffPrimeQZero.lean}{GlobalTheory.RowEx.IsEx}
\leanref{Development/AbelianVarietyTheory/Proofs/TwistedOperators/IsExOfSeparates.lean}{GlobalTheory.RowEx.isEx_of_separates}
\leanref{Development/AbelianVarietyTheory/Definitions/TwistedInnerProducts/TwistPlacesOffPrimeQZero.lean}{GlobalTheory.RowEx.NFAll}
\leanref{Development/AbelianVarietyTheory/Proofs/TwistedOperators/RangeAlgebraMapMono.lean}{GlobalTheory.RowEx.nfAll_of}
\leanref{Development/AbelianVarietyTheory/Definitions/AutomorphicMembership/GKeyFin.lean}{GlobalTheory.Mem.LiteralAutomorphy.NFSq}
\leanref{Development/AbelianVarietyTheory/Proofs/IwasawaDecomposition/NormFieldSquareClass.lean}{GlobalTheory.Mem.LiteralAutomorphy.nfSq_of_fieldDat}
\leanref{Development/AbelianVarietyTheory/Definitions/TwistedInnerProducts/TwistPlacesOffPrimeQZero.lean}{GlobalTheory.RowEx.BigR}
\leanref{Development/AbelianVarietyTheory/Statements/HeckeThetaIntegrals/FormulaAt.lean}{GlobalTheory.RowClaim.Big}

The proof is not yet written out; it is part of the open step in the proof of
Theorem~\ref{asm-dependence-inputs}. Part (1) rests on the existence of Weil-Deligne data from the
local inputs; part (2) on the linear bound and the base and constant existence for the twisted
Weil representations, through the cuspidality of the theta lifts of the pairs; part (3) on six
number fields, for the characters, the roots of unity, the phases, the jets, the ratios of
$L$-values (through the conjugates of the toric ratios and the nonvanishing of the $L$-values given
by the functional equation) and the bottom sums; and part (4) on the field data of the squares, the
existence of the toric periods and the smallness of the bottom.

\begin{lemma}\label{asm-dependence-grid-exists}
Assume the published results on the local theta correspondence (Theorem~\ref{lit-bundle-theta}),
the join statements and the functional equation (Theorem~\ref{asm-dependence-inputs}), the theta
hypotheses (Definition~\ref{th-theta-hypotheses}), and that functionals separate
(Definition~\ref{asm-grid-evaluation}(7)). There is $C_{F_1}$ with the following property. For every
$\varepsilon$ and every $C_F\ge C_{F_1}$ there is $x_0$ such that for every scale $(x,\varepsilon)$
with $x\ge x_0$, all base data $B$ with a framing $f$ and every finite set $\Xi$ of Hecke characters
of the host field of $B$, there is a finite set $S$ of algebraic numbers with the following property.
Let $D$ be the definite-side datum of a later stage over $(B,f)$ whose coefficient field contains $S$,
with size data at the scale $(x,\varepsilon)$ that satisfy the recursion conditions, and let $R$ be a
realization of $D$ in an automorphic setting, with these size data at a face
(Definition~\ref{asm-dependence-functional-witness-setting}), whose columns are given by the
framing, whose twisted places lie off $q_0$, and every character $\xi_t$ of whose level system lies
in $\Xi$. Then for every evaluation grid $G_0$
for $R$ with constant $C_F$ at $b$ there is an evaluation grid $G$ for $R$ with constant $C_F$ at $b$,
with the base place, the place above it and the characters $\chi_1,\dots,\chi_m$ of $G_0$, which
carries a full evaluation choice (Definition~\ref{asm-stage-construction}(5)).
\end{lemma}
\leanref{Development/AbelianVarietyTheory/Definitions/AutomorphicMembership/GKeyFin.lean}{GlobalTheory.Mem.LiteralAutomorphy.ExGrid}
\leanref{Development/AbelianVarietyTheory/Definitions/AutomorphicMembership/GKeyFin.lean}{GlobalTheory.Mem.LiteralAutomorphy.GridAt}
\leanref{Development/AbelianVarietyTheory/Proofs/IwasawaDecomposition/GridExistenceOfPieces.lean}{GlobalTheory.Mem.LiteralAutomorphy.exGrid_of}
\leanref{Development/AutomorphicArithmetic/ThetaCorrespondence/ThetaPairLifts/LocalTypeDictionary/EvaluationGridOfJoin.lean}{GlobalTheory.Mem.LiteralAutomorphy.evaluation_grid_of_join}
\begin{proof}
The hypotheses of Lemma~\ref{asm-dependence-grid-pieces} hold. The convergence of the theta integrals
of the twisted Weil representations follows from Siegel sets for unitary groups, which are among the
results of Theorem~\ref{lit-bundle-theta}. The matching of the twisted local factors is
Lemma~\ref{th-first-lift-factor}(2), and the agreement of the twisted square theta integrals is
Lemma~\ref{th-first-lift-inputs}(3); both rest on Howe duality at the real place and the theta
hypotheses. Let $C_{F_1}$ be the largest of the four constants of
Lemma~\ref{asm-dependence-grid-pieces}. Let $\varepsilon$ and $C_F\ge C_{F_1}$ be given, and let
$x_0$ be the largest of the four numbers that the lemma gives for them.

Let $(x,\varepsilon)$ be a scale with $x\ge x_0$, put $b=e^{x^{\varepsilon}}$, and let $B$, $f$ and
$\Xi$ be given. By Lemma~\ref{asm-dependence-grid-keys}, there are finite sets $\mathcal K$ of keys
and $\mathcal G$ of grid keys such that every realization $R$ as in the lemma has its key in
$\mathcal K$, and every evaluation grid for it with constant $C_F$ at $b$ has its grid key in
$\mathcal G$. Say that data $(X,R,G)$ \emph{represent} $(\kappa,\gamma)\in\mathcal K\times\mathcal G$
if $X$ is a later stage over $(B,f)$ whose datum has size data at $(x,\varepsilon)$ satisfying the
recursion conditions, $R$ is a realization of this datum in an automorphic setting whose columns are
given by $f$, whose twisted places lie off $q_0$ and whose key is $\kappa$, and $G$ is an evaluation
grid for $R$ with constant $C_F$ at $b$ with grid key $\gamma$.

For each pair $(\kappa,\gamma)$ choose a finite set $S_{\kappa,\gamma}$ of algebraic numbers as
follows. If no data represent the pair, put $S_{\kappa,\gamma}=\emptyset$. Otherwise fix data
$(X_0,R_0,G_{00})$ that represent it, and let $\mu_0$ be the base tuple of $R_0$.
Lemma~\ref{asm-dependence-grid-pieces}(1) gives a Weil datum $W_0$ over $R_0$ for $G_{00}$ and a grid
choice $S_0$ for $W_0$, and part (2) of that lemma gives further choices extending $S_0$ that satisfy
the evaluation conditions. Part (3) then gives a number field $N_1\subset\mathbf C$ in which the
numbers of $(G_{00},W_0,S_0)$ at $\mu_0$ lie, and part (4) a number field $N_2\subset\mathbf C$ in which
the square numbers of $G_{00}$ at $\mu_0$ lie. Let $S_{\kappa,\gamma}$ be the union of a basis of $N_1$
over $\mathbf Q$ and a basis of $N_2$ over $\mathbf Q$; it consists of algebraic numbers. Let $S$ be the
union of the sets $S_{\kappa,\gamma}$, which is finite.

Now let $D$, $R$ and $G_0$ be as in the lemma, with base tuple $\mu'_0$, and let $L$ be the coefficient
field of $D$, so that $S\subset L$. Let $\kappa$ be the key of $R$ and $\gamma$ the grid key of $G_0$.
Then $(X,R,G_0)$ represent $(\kappa,\gamma)$, where $X$ is the later stage of $D$, so data
$(X_0,R_0,G_{00})$ were fixed for this pair. Since $R$ and $R_0$ have the same key, every tuple $t$ of
$R$ has a tuple $\iota(t)$ of $R_0$ with $\xi_{\iota(t)}=\xi_t$ and $Q_{\iota(t)}=Q_t$. The base tuples
are the tuples with trivial character, and $t\mapsto\xi_t$ is injective on the tuples of $R_0$, so
$\iota(\mu'_0)=\mu_0$. Both stages lie over $B$, the columns of both realizations are given by $f$,
and in an automorphic setting the adelic frame of $X'$ is the identity. So the level systems of $R_0$
and $R$ agree through $\iota$.

By Lemma~\ref{asm-dependence-grid-transport}(1), $G_{00}$ has an agreeing grid $G$ for $R$, with
constant $C_F$ at $b$. Its base place, the place above it and its characters are those of $G_{00}$,
so $G$ has the grid key $\gamma$ of $G_0$. The automorphic setting of $R$ gives the chain and model
conditions of its core, so by parts (2) and (3) of that lemma there are a Weil datum $W$ over $R$ for
$G$ and a grid choice $S'$ for $W$ that agree with $W_0$ and $S_0$ through $\iota$.
Lemma~\ref{asm-dependence-grid-pieces}(2), applied to $D$, $R$, $G$, $W$ and $S'$, gives further
choices that extend $S'$ and satisfy the evaluation conditions. By
Lemma~\ref{asm-dependence-grid-transport}(4), the numbers of $(G,W,S')$ at $\mu'_0$ lie in $N_1$, and
the square numbers of $G$ at $\mu'_0$ lie in $N_2$. The field $L$, embedded in $\mathbf C$, contains a
basis of $N_1$ and a basis of $N_2$, hence $N_1$ and $N_2$. So the full evaluation choice of $G$ with
$W$ and these further choices satisfies the field conditions, and the square numbers of $G$ lie in
$L$: $G$ carries a full evaluation choice.
\end{proof}

\begin{lemma}\label{asm-dependence-inputs-sha-published}
Assume the first and second collections of item (6) of Theorem~\ref{lit-bundle-unitary}. Then
Alpöge's published inputs of the third layer, his series statements and his elliptic hypothesis
hold, and there are local Langlands maps that satisfy his local Langlands hypothesis.
\end{lemma}
\leanref{Development/AutomorphicArithmetic/ThetaCorrespondence/ThetaPairLifts/ThirdPlaceSeries/ThetaLiftHypothesesOfNonvanishing.lean}{GlobalTheory.T.hyps_e_third_layer_ser_n_of_geometric_galois}
\begin{proof}
\openstep{The derivation of the published inputs of the third layer, the series statements, the
elliptic hypothesis and the local Langlands hypothesis from the first and second collections is not
yet written out.}
\end{proof}
% Lean: hyps_e_third_layer_ser_n_of_geometric_galois = MemberClose.exists_hyp_n_of_geometric_galois_compat,
% from NonEndoscopicHypothesesAlt and GeometricGaloisRepCompatibility; conclusion EllipticHypothesis,
% PublishedInputsThirdLayerAlt, SerN and (exists hL, LocalLanglandsHypothesis hL). Size of the open
% step: proof closure of exists_hyp_n_of_geometric_galois_compat, about 59,000 theorems and
% 1,025,000 lines (text dependencies, an overcount).

\begin{lemma}\label{asm-dependence-inputs-sha-preconditions}
Let local Langlands maps be given. Say that a property of a definite-side datum and a norm choice
for it holds \emph{at the layer stages} if, eventually in the context and at large base data
(Definition~\ref{asm-stage-construction}), it holds for the datum $D(w,Z)$ and the norm choice of
every stage whose core has real expansion data, whose rows are rows of the construction and whose
columns are pinned.
\begin{enumerate}
\item Assume Alpöge's published inputs of the third layer, the second collection of item (6) of
  Theorem~\ref{lit-bundle-unitary}, his series statements, the passage of depth zero to the cuspidal
  support (item (2) there), and local finiteness, local Tate duality and the local Euler
  characteristic formula (item (5) there). Then his Galois preconditions on a datum and its norm
  choice hold at the layer stages.
\item Assume the published inputs of the third layer, the count of Hodge-Tate weights in the second
  collection and the passage of depth zero to the cuspidal support. Then his fourth tie of the datum
  holds at the layer stages.
\item If two properties hold at the layer stages, so does their conjunction.
\end{enumerate}
\end{lemma}
\leanref{Development/AutomorphicArithmetic/ThetaCorrespondence/ThetaPairLifts/ThirdPlaceSeries/ThetaDataExists.lean}{GlobalTheory.HeightCounting.Up.MemberAtLayerPoint}
\leanref{Development/AutomorphicArithmetic/ThetaCorrespondence/ThetaPairLifts/ThirdPlaceSeries/PreAtMem.lean}{UpperTheory.GalS.galois_preconditions_at_mem}
\leanref{Development/AutomorphicArithmetic/ThetaCorrespondence/ThetaPairLifts/ThirdPlaceSeries/LocalTypeDeterminedAtAllPrimes.lean}{GlobalTheory.T.tieChi_at_members_of_depthZero}
\leanref{Development/AutomorphicArithmetic/ThetaCorrespondence/ThetaPairLifts/ThirdPlaceSeries/AtMemberImp.lean}{GlobalTheory.HeightCounting.Up.Model.AtMem.and_z}
\begin{proof}
(3) Each property holds, eventually in the context, at large base data. Two properties that hold
eventually in the context hold together eventually in the context, and two that hold at large base
data hold together at large base data, for the unions of their finite sets of algebraic numbers and
of primes.

(1) Three properties hold at the layer stages: the existence of inputs with their pins and the count
of their jumps, from the published inputs and the second collection; the condition on finite
triples, from the published inputs and the series statements; and the goodness of every row of the
datum, from the published inputs. By (3) they hold together. At a layer stage where they hold, the
inputs and their jumps, with three of the published inputs, the passage of depth zero to the cuspidal
support and the three local results, give the core of the Galois preconditions. The other
preconditions are a statement on the local data of the inputs and their Hecke operators, the
largeness of the datum, which the real expansion data of the core record, the tie of the norm choice
to the stage, the condition on finite triples, the agreement of the twists, and the goodness of the
rows.
\openstep{In (1), the three properties at the layer stages, the core of the Galois preconditions, the
statement on the local data of the inputs and the agreement of the twists; and part (2), are not yet
written out.}
\end{proof}
% Lean: MemberAtLayerPoint hL p = EvLgG over defaultExtension, RXs and ColPin. (3) AtMem.and_z (EvG.inter,
% large2.inter). (1) galois_preconditions_at_mem: NRows.AtLayer.exists_inputs_pinAt_jumpsCard_at_mem,
% TripFin.finiteTriples_satisfy_condition_alt (with SerN), T.goodFTAll_at_mem_of_inputs, then
% GalS.core0_of_local_duality_inputs (prn.what2, prn.hc, prn.decay, DepthZeroSupport, LocalFiniteness,
% LocalTateDuality, LocalEulerChar) and the fields c2 (LocCoincide.main_alt_alt), large, nct, trip,
% same (SameTw.sameTwist_at_alt) and good of GaloisPreconditions (eighteen fields in all). (2)
% tieChi_at_members_of_depthZero with HtCard = hX.ht; its conclusion is Tie4 (tieChi, admDZ with
% EigSpan, TieInePoly, SerLoc). Sizes of the open step (text dependencies, overcounts):
% galois_preconditions_at_mem about 37,400 theorems and 622,000 lines; tieChi_at_members_of_depthZero
% about 11,200 theorems and 181,000 lines.

\begin{lemma}\label{asm-dependence-inputs-sha-transport-parts}
Let a stage $(K,B,f,T,w,Z)$ be given with a core object $O$, let new division lattices $N_u$ and
heights $h$ be chosen at the later stage $Z$, and let $Z'$ be the later stage with these division
data and $O'$ the transport of $O$ to $Z'$ (Lemma~\ref{asm-dependence-bundles-inputs-move}). The datum
$D(w,Z')$ is $D(w,Z)$ with its division lattices replaced by the $N_u$, its heights by $h$, and the
bound on the heights recomputed from them.
\begin{enumerate}
\item Column data at the split primes of the stage, whose Hermitian form and forms satisfy the
  recursion conditions, transport to column data at the split primes of $(K,B,f,T,w,Z')$ whose
  transported Hermitian form and forms satisfy the recursion conditions.
\item If a kernel Taylor series of a kernel family of $D(w,Z)$ satisfies the entry condition, its
  transport is a kernel Taylor series of the transported kernel family that satisfies the entry
  condition.
\item A tie of a norm choice to the stage transports to a tie of the transported norm choice to
  $(w,Z')$. If division local data and theta data of $D(w,Z)$ hold with the first tie, their
  transports hold with the second.
\item If the twists of a norm choice and theta data of $D(w,Z)$ are good, so are the twists of their
  transports.
\end{enumerate}
\end{lemma}
\leanref{Development/AbelianVarietyTheory/Definitions/GlobalTwisting/SetUnipotentTranslate.lean}{GlobalTheory.EX.OB.setDiv}
\leanref{Development/AbelianVarietyTheory/Proofs/IwasawaDecomposition/ConcreteSplitPrimeAutomorphicDataDivisionLatticeStatement.lean}{GlobalTheory.T.concreteSplitPrimeAutomorphicData_divisionLatticeChange}
\leanref{Development/AbelianVarietyTheory/Proofs/IwasawaDecomposition/ColumnLocalDataTransportDivision.lean}{GlobalTheory.Mem.Top.Bl.WitTr.columnLocalDataTransportDivision}
\leanref{Development/AbelianVarietyTheory/Proofs/IwasawaDecomposition/KernelTaylorSeriesPreservedDivisionLatticeChangeStatement.lean}{GlobalTheory.T.kernelTaylorSeries_preserved_divisionLatticeChange}
\leanref{Development/AbelianVarietyTheory/Proofs/IwasawaDecomposition/ColumnLocalDataTransportDivision.lean}{GlobalTheory.Mem.Top.Bl.WitTr.nchTTr}
\leanref{Development/AbelianVarietyTheory/Proofs/IwasawaDecomposition/DivisionLocalDataPreservedDivisionLatticeChangeStatement.lean}{GlobalTheory.T.divisionLocalData_preserved_divisionLatticeChange}
\leanref{Development/AbelianVarietyTheory/Proofs/IwasawaDecomposition/ThetaAutomorphicWeilDataPreservedDivisionLatticeStatement.lean}{GlobalTheory.T.thetaAutomorphicWeilData_preserved_divisionLatticeChange}
\leanref{Development/AbelianVarietyTheory/Proofs/IwasawaDecomposition/GoodTwistsThetaLatticesPreservedDivisionChangeStatement.lean}{GlobalTheory.T.goodTwistsThetaLattices_preserved_divisionChange}
\begin{proof}
In each part the transported data are the given data, read at the moved datum, and each condition
is checked clause by clause.

(1) The recursion conditions say that the ties, the comparison of lattices, the Hermitian form and
the forms of the column data are those built recursively from a datum $D_0$ with column levels,
local identifications and characters at the places above the twisting primes, a Hermitian form and
forms, a basis in standard position and embeddings at these places, subject to conditions on these
choices that do not involve the division data. For the transported column data take $D_0$ moved in
the same way, with the same other choices. The recursive construction from the moved $D_0$ gives the
transported ties, Hermitian form and forms, and its comparison of lattices is the transport of the
given one. The other column data at the split primes are kept.

(2) The transported series has the same entries, the same place data and the same matrices at the
places. The entry condition is an equation between the entries of the series and those built from
the place data and matrices; it is the same equation after the move.

(3) The tie of the norm choice, the division local data and the theta data have the same fields after
the move. The geometric tie of the division local data at a moved datum is the geometric tie at the
datum it was moved from. The automorphic functions, theta functions and polynomial Fock vectors that
the theta data use are defined in the same way for both data, so their conditions hold after the move.

(4) The good twists are a condition on the twisting data of the norm choice and of the theta data,
which the move keeps.
\end{proof}
% Lean: setDiv D N h = D with div.Nu u having lattice N u, div.h = h and div.HH recomputed;
% setDivZ moves the later stage. (1) concreteSplitPrimeAutomorphicData_divisionLatticeChange =
% WitTr.isRecG_setDiv (IsRecG.intro with the same psi, hd, levels, iK, lam, hK, hPos, stdB_setDiv,
% cmpHyp0P_setDiv and sigma_setDiv/latCmpC_setDiv, which is rfl), and columnLocalDataTransportDivision
% { colTTr with isRec, om }; (2) kerT_entC_setDiv = (DivTr.ktTr, DivTr.entC_tr), structure copies of
% KerT and of the entry data (EntPl, place matrices, KerEntHyp); (3) nchTTr { n with tp, tie, chrPsi };
% divT_setDiv (Qp, chr, loc, inK copied, geo = GeoDivTr.geoTiedP_divTr through datBk and famT);
% ThTTr.thT_setDiv (allTiedTr, weilTiedEq_tr, with isAutFun_tr, isThetaFun_tr, isAutFunSm_tr and
% isPolFock_tr, iffs between the two data); (4) WitTr.goodO_setDiv (the twists twEOf xd copied). Size:
% about 460 lines in seven files (WitTr, DivTr), all structure copies.

\begin{lemma}\label{asm-dependence-inputs-sha-transport}
Let local Langlands maps be given, let $c$ be real, let $N=N_{\mathrm{list}}(c)$ and $p=p(N)$, the
parameters of dependence $c$ (Definition~\ref{asm-dependence-parameters}).
\begin{enumerate}
\item The positions of $p$ move with the division data, in the sense of
  Lemma~\ref{asm-dependence-bundles-inputs}(2).
\item The witnesses of $p$ move with the division data: for every stage built from $p$ with a witness
  whose kernel satisfies the entry condition, and every choice of new division lattices and heights at
  its later stage, the moved stage has a witness whose kernel satisfies the entry condition.
\item Let $P$ be a property of a definite-side datum and a norm choice that holds at the layer stages
  (Lemma~\ref{asm-dependence-inputs-sha-preconditions}). Then the property ``eventually in the
  context, every bundle that $\Lambda_p$ links to a base point of a slot satisfies $P$ for its datum
  and the norm choice of its core'' is linked at $p$.
\end{enumerate}
\end{lemma}
\leanref{Development/AutomorphicArithmetic/ThetaCorrespondence/ThetaPairLifts/ThirdPlaceSeries/SlotMvPin.lean}{GlobalTheory.NLink.Final.slot_mv_extension_with_list_geo_cl}
\leanref{Development/AbelianVarietyTheory/Proofs/IwasawaDecomposition/WitnessesPassAlongSetDivisionStatement.lean}{GlobalTheory.T.witnesses_pass_along_set_division}
\leanref{Development/AutomorphicArithmetic/ThetaCorrespondence/ThetaPairLifts/ThirdPlaceSeries/OverNEOfNEb.lean}{GlobalTheory.Mem.AtPtTransport.atMem_at_wit_alt}
\leanref{Development/AbelianVarietyTheory/Proofs/IwasawaDecomposition/WitnessesPassAlongSetDivision.lean}{GlobalTheory.Mem.Top.Bl.witMv_of_positionMv}
\leanref{Development/AutomorphicArithmetic/ThetaCorrespondence/ThetaPairLifts/ThirdPlaceSeries/Lg1EntryOfMembership.lean}{GlobalTheory.NRows.AtPt.lg1_entry_of_atWit}
\leanref{Development/AutomorphicArithmetic/ThetaCorrespondence/ThetaPairLifts/ThirdPlaceSeries/Lg1EntryOfMembership.lean}{GlobalTheory.NRows.AtPt.exists_wPt_entC_atWit}
\begin{proof}
(1) Let a stage $(K,B,f,T,w,Z)$ built from $p$ be given, with core object $O$ and a witness, new
division lattices $N_u$ and heights $h$ at $Z$, a column datum at the split primes and a norm choice
of the moved stage $(K,B,f,T,w,Z')$. Let $t'$ be the rows at the base level of the transport $O'$ of
$O$ that Lemma~\ref{asm-dependence-bundles-inputs-move}(1) gives for the rows $t$ of the witness. By
that lemma they satisfy the core object conditions and are rows of the construction, which are the
conditions of $p$ on rows. The norm choice and theta data of $O'$ are the transports of those of $O$,
by part (2) of that lemma, so the data of the law extension that the witness carries give data of the
law extension at $(Z',t')$; and the list pin passes to $t'$, by part (3). These are the conditions of
$N_{\mathrm{list}}(c)$. The spectral stage data pass, by part (4), which are the conditions on the
twist choices; and the scalars of the division geometry pass, by part (5), which are the conditions
on the core and its geometry.

(2) Let $x$ be a witness that the stage $(K,B,f,T,w,Z)$, with core object $O$, is built from $p$,
whose kernel satisfies the entry condition. Let new division lattices $N_u$ and heights $h$ be chosen
at $Z$, and let $Z'$ and $O'$ be the moved later stage and the transported core object. A witness for
$(K,B,f,T,w,Z')$ with core object $O'$ is assembled from $x$ as follows.
\begin{itemize}
\item The parts of $x$ that do not involve the division data are kept: that the base data lie over
  the context, the forced conditions, the set $W_t$ of twisted places, the pinned columns, the
  coherence of the twist data, the volume and real ties of the datum, the values $d_\sigma$ of the
  size stage, the gauge, and the equality of the levels of $O$ with those of the column data.
\item The column data at the split primes of $x$ transport to $Z'$, with their recursion conditions,
  and the tie of the norm choice of $O$ transports to a tie of that of $O'$
  (Lemma~\ref{asm-dependence-inputs-sha-transport-parts}(1) and (3)).
\item Part (1), applied to $x$, $N_u$, $h$, the transported column data and the transported tie,
  gives rows $t'$ at the base level of $O'$ that satisfy the core object conditions, the conditions
  of $p$ on rows, the extension $N$ with the transported tie, the conditions on the twist choices
  with the transported column data, and the conditions on the core and its geometry.
\item The kernel family of $x$ transports; the kernel and the error term of $O'$ are the transported
  kernel family and the error term of the transported column data, since those of $O$ are the kernel
  family of $x$ and the error term of its column data. The kernel Taylor series of $x$ transports to
  one of the transported kernel family that satisfies the entry condition
  (Lemma~\ref{asm-dependence-inputs-sha-transport-parts}(2)).
\item The division local data and the theta data of $x$ transport, with the transported tie, and the
  good twists of $O$ give good twists of $O'$
  (Lemma~\ref{asm-dependence-inputs-sha-transport-parts}(3) and (4)).
\end{itemize}
These are the parts of a witness that $(K,B,f,T,w,Z')$ is built from $p$, and its kernel satisfies
the entry condition.

(3) Let $R_P$ be the slot of the cores of the stages $(K,B,f,T)$ such that $P$ holds for the datum
$D(w,Z)$ and the norm choice of every stage $(K,B,f,T,w,Z)$ whose core has real expansion data,
whose rows are rows of the construction and whose columns are pinned. Since $P$ holds at the layer
stages, eventually in the context and at large base data every core satisfies $R_P$.

Let $\beta$ be a base point, in a context $K$, of a good core of $p$ that satisfies $R_P$. The good
core is the core of a stage $(K,B,f,T,w,Z)$ with core object $O$ that is built from $p$, with a
record of its construction whose kernel satisfies the entry condition, and the twists of $O$ are
good; the record with the good twists is a witness $x$ that the stage is built from $p$. The base
point $\beta$ is given by a layer of $D(w,Z)$ whose object is $O$ and by size data at the scale of
$K$ with further choices; its datum is $D(w,Z')$ and its core object is the transport $O'$ of $O$,
where $Z'$ is $Z$ with the division lattices and the heights that the size data determine. By (2),
the stage $(K,B,f,T,w,Z')$ has a witness $x'$ whose kernel satisfies the entry condition. The witness
$x'$ carries data of $N$ for the norm choice of $O'$, and $N$ is stronger than the law extension, so
the core of $x'$ has real expansion data; $x'$ also records that its rows are rows of the
construction and that its columns are pinned. The good core satisfies $R_P$, and $(K,B,f,T,w,Z')$ is
one of the stages in the definition of $R_P$ for $(K,B,f,T)$; so $P$ holds for $D(w,Z')$ and the norm
choice of $O'$, which are the datum of $\beta$ and the norm choice of its core.

So, in every context, $P$ holds for the datum and the norm choice of the core of every base point of
a good core of $p$ that satisfies $R_P$, and every bundle that $\Lambda_p$ links to such a base point
satisfies the property of (3). With the slot $R_P$ of Definition~\ref{asm-dependence-parameters}(2),
the property of (3) is linked at $p$.
\end{proof}
% Lean: (1) NLink.Final.slot_mv_extension_with_list_geo_cl: rows_pass_to_moved_divisionLatticesHeights
% (t', RowsSame, geoHalf, the pack, ObjT, RXs), exists_kept_choice_upper_lattice_alt with S' = Z.S
% (nch' = nchTr, xd' = xdTr, and hmv applied to <x.nx.1, x.nx.2.1>), listPin_rowsSame on x.nx.2.2,
% spectralStageData_transfer_lattices_heights with the pack and x.sx.1, divisionGeometryScalars_transfer
% with x.ox; Nx = (c.1, c.2, ListPin), Sx = (s, unit), Ox = o. (2) witnesses_pass_along_set_division =
% witMv_of_positionMv: hR = isRecG_setDiv, col = columnLocalDataTransportDivision, (k, hk) =
% kerT_entC_setDiv, n = nchTTr, dv = divT_setDiv, th = thT_setDiv, then PositionMv (part (1)) at
% O' = O.tr N h with col and n, hker from x.ker_eq, herr from x.errX, and the fields isOver, forced, Wt,
% pin, coh, vol, lev, real, dsig_eq, gauge copied; good = goodO_setDiv. (3) the second part of
% AtPtTransport.atMem_at_wit_alt = lg1_entry_of_atWit with hN' = overNE_of_NEb_alt (N gives the
% defaultExtension, the real expansion data): lgSlotG_memSlot' (R_P is membershipLocation P, a condition
% on K, B, f, T only), memSlot_entry_of_atWit (RowE.of_all, every context), exists_wPt_entC_atWit
% (positionOf_ind on GoodE = BuiltE and GoodO; Built.wit), exists_wit_pt_alt (WitMv at nPrime p.L and
% p.ht of the LedPart p of the base point), exists_wPt_entC_atWit_of_heq (D_setDivZ and pointObj_pt,
% both rfl; hP with x.nx.1, hN' x.nx.2, x.rx and x.pin), and linkedBundleProperty_after_largeBaseCoreProperty,
% whose slot is largeBaseProperty_conditional of R_P (R_P in the contexts where the large-base condition
% holds). The property holds for any link data, not only Lambda_p. Size of the step written out here:
% witMv_of_positionMv (75 lines), the transfers of Lemma -transport-parts (about 460 lines of structure
% copies), and Lg1EntryOfMembership.lean with OverNEOfNEb.lean (about 110 lines).

\begin{lemma}\label{asm-dependence-inputs-sha-cut}
Assume the Chebotarev density theorem of Alpöge's published inputs of the third layer, local
finiteness, local Tate duality and the local Euler characteristic formula
(Theorem~\ref{lit-bundle-unitary}(5)). Let $D$ be a definite-side datum with a norm choice that
satisfies the Galois preconditions, and suppose that the fourth tie holds for $D$
(Lemma~\ref{asm-dependence-inputs-sha-preconditions}). Then $D$ has a good cut, at which Alpöge's
norm-choice facts hold for the cut of the norm choice, and at which the fourth tie holds.
\end{lemma}
\leanref{Development/AutomorphicArithmetic/ThetaCorrespondence/ThetaPairLifts/ThirdPlaceSeries/CFacts.lean}{UpperTheory.GalC.exists_good_cut_facts_tie4}
\leanref{Development/AutomorphicArithmetic/ThetaCorrespondence/ThetaPairLifts/ThirdPlaceSeries/LinkNFam.lean}{UpperTheory.GalC.CutGivesAlt}
\begin{proof}
The central elements given by the coefficients of the central idempotents, with the Galois
preconditions and the fourth tie, give a good cut of $D$. The Galois preconditions pass to the cut,
and the cut is recursive; with the Chebotarev density theorem and the three local results, the
norm-choice facts hold at the cut. The fourth tie of $D$ passes to the cut.
\openstep{The good cut, the passage of the Galois preconditions to it, the norm-choice facts at it and
the passage of the fourth tie to it are not yet written out.}
\end{proof}
% Lean: exists_good_cut_facts_tie4: exists_good_alt (central_elements_from_idempotent_coeffs
% centralIdemCoeff), facts_of_galois_preconditions_of_cutRec (prN.cheb, lf, ltd, lec) with
% pre_cut_alt and cutRec_cut, and compatibilityTies_eigenRestriction. Size of the open step: proof
% closure of exists_good_cut_facts_tie4, about 2,400 theorems and 38,400 lines (text dependencies, an
% overcount).

\begin{lemma}\label{asm-dependence-inputs-sha-link}
Assume the third collection of item (6) of Theorem~\ref{lit-bundle-unitary}, and suppose that every
definite-side datum with a norm choice that satisfies the Galois preconditions and the fourth tie has
a good cut as in Lemma~\ref{asm-dependence-inputs-sha-cut}. Let local Langlands maps be given and let
$p$ be parameters (Definition~\ref{asm-dependence-parameters}). Suppose that the property
``eventually in the context, every bundle that $\Lambda_p$ links to a base point of a slot satisfies
the Galois preconditions for its datum and the norm choice of its core, and the fourth tie for its
datum'' is linked at $p$, and that the property ``the local bound at the twisted places above $2$ and
$3$ (Definition~\ref{asm-local-bound-two-three}(4)) holds eventually along
$\mathcal A^{\mathrm{des}}$'' is linked at $p$. Then the property ``the Tate-Shafarevich link
(Definition~\ref{asm-twisted-links}) holds eventually along $\mathcal A^{\mathrm{des}}$'' is linked
at $p$. The proof does not use the local bound at $2$ and $3$, only the number $\varepsilon_1$ and
the largeness thresholds with which it holds eventually along $\mathcal A^{\mathrm{des}}$.
\end{lemma}
\leanref{Development/AutomorphicArithmetic/ThetaCorrespondence/ThetaPairLifts/ThirdPlaceSeries/LinkNFam.lean}{UpperTheory.GalC.linkedAt_linkN_of_boundAtTwoThree}
\leanref{Development/AutomorphicArithmetic/ThetaCorrespondence/ThetaPairLifts/ThirdPlaceSeries/LinkNFam.lean}{UpperTheory.GalC.linkN_of_boundAtTwoThree_slot}
\begin{proof}
Let $R_1$ and $R_3$ be slots for the two hypotheses. Every core built from $p$ satisfies $R_1$ and
$R_3$ eventually in the context and at large base data, so it satisfies $R_1\cap R_3$ there. Let $S$
be the slot of the good cores of $p$ that satisfy $R_1\cap R_3$. Restricting from $R_1$, eventually in
the context every bundle that $\Lambda_p$ links to a base point of $S$ satisfies the Galois
preconditions and the fourth tie; restricting from $R_3$, the local bound at $2$ and $3$ holds
eventually along $\mathcal A^{\mathrm{des}}$, along $(S,\Lambda_p)$, with a condition $E$ on bundles.
Take for the condition on bundles $E$ together with the first property. Let $s$ be a selection for a
class $\mathcal A$ whose bundles are built from $(S,\Lambda_p)$ and satisfy it. At every member, the
datum and the norm choice of the selected bundle satisfy the Galois preconditions and the fourth tie.
Take $\varepsilon_1$ and the largeness thresholds of the local bound at $2$ and $3$, and let
$\mathcal F$ be a family of $\mathcal A^{\mathrm{des}}$ that is large for them. It satisfies the twist
condition (Definition~\ref{asm-twist-condition}), as every family of $\mathcal A^{\mathrm{des}}$ does
(Definition~\ref{asm-designed}(4)). From the twist condition, the Galois preconditions and
the fourth tie at its members, the third collection and the good cuts, the Tate-Shafarevich link
holds for $\mathcal F$.
\openstep{The Tate-Shafarevich link of a family that satisfies the twist condition and whose members
satisfy the Galois preconditions and the fourth tie, from the third collection and the good cuts, is
not yet written out.}
\end{proof}
% Lean: linkedAt_linkN_of_boundAtTwoThree (r1.and r3, EvP.mono, Row.strengthen), then
% linkN_of_boundAtTwoThree_slot ((RowE.row hE).and h23, Closed with the eps_1 and thresholds K of
% h23, whose bound hK is not used), whose content is linkN_of_twGood_galoisPreconditions hP hG Fm with
% hc.1.2 = TwGood Fm, a condition of the designed class Des (WithChN A) 3. Size of the open step:
% proof closure of linkN_of_twGood_galoisPreconditions, about 17,700 theorems and 311,000 lines (text
% dependencies, an overcount).

\begin{theorem}\label{asm-dependence-inputs}
Write $P$ for the published results on unitary groups (Theorem~\ref{lit-bundle-unitary}) and $F$ for
those on the local theta correspondence (Theorem~\ref{lit-bundle-theta}), and
$\mathcal A^{\mathrm{des}}=\operatorname{Des}(\mathcal A_N,3)$ for a class $\mathcal A$ of families.
Each part assumes the results named in it, and only those. The \emph{join statements} and the
\emph{functional equation} are Alpöge's five statements on lattices with complex multiplication and
on denominators and reflection identities of zeta and $L$-values, and the functional equation of
Hecke $L$-functions of CM type, in the first collection of item (6) of
Theorem~\ref{lit-bundle-unitary}; the \emph{theta hypotheses} are the agreement of global and local
theta integrals and the local matching of models of the Weil representation
(Definition~\ref{th-theta-hypotheses}); and the \emph{four results on automorphic forms} are the
covering of discrete series, the finiteness of Harish-Chandra modules, rapid decay of cusp forms and
the Casselman-Wallach theorem (Theorem~\ref{lit-bundle-unitary}).
\begin{enumerate}
\item (Separation.) Assume $F$ and the Siegel-Weil formula (Theorem~\ref{lit-siegel-weil}). Then
  functionals separate (Definition~\ref{asm-grid-evaluation}(7)).
\item (Rows.) Assume $F$, Green's parametrization (Corollary~\ref{lit-green-labelled}), the join
  statements, the functional equation and the theta hypotheses, and that functionals separate. Then
  the rows of the construction are rational (Definition~\ref{asm-stage-construction}(4)).
\item (Fourier fit.) Assume $F$, the join statements, the functional equation and the theta
  hypotheses, that functionals separate, and the conclusion of part (2) of
  Theorem~\ref{is-surjectivity-and-lower-norm}. Then there is $C'$ such that for every $C_F\ge C'$,
  eventually along the class of
  standard families (Section~\ref{section-is-points}), every family that satisfies the division
  bound at twisted places (Definition~\ref{is-twisted-division-bound}) has a scale $x_0$ with the
  following property: for every point $p$ at a scale $x\ge x_0$ that satisfies the point conditions
  (Definition~\ref{is-point}) and has a fitted evaluation grid with constant $C_F$
  (Definition~\ref{asm-dependence-functional-fitted}(5)), the rows $S_0(p)$ satisfy the functional bound (Definition~\ref{is-functional-bound}) for
  $N_p(C,C_H)$ for every real $C$.
\item (Evaluation.) Assume $F$, the join statements, the functional equation and the theta
  hypotheses, and that functionals separate. Then there is $C_{F_1}$ such that for every
  $C_F\ge C_{F_1}$ the grid with constant $C_F$ evaluates (Definition~\ref{asm-stage-construction}(5)).
\item (Functional condition.) Assume the four results on automorphic forms, the existence of ray
  class fields (Theorem~\ref{lit-bundle-unitary}(5)) and the conclusion of part (3). Then there is
  $C_{F_0}$ with the following property. Let $C_F\ge C_{F_0}$, let local Langlands maps be given
  (Definition~\ref{asm-dependence-parameters}), let $N$ be an extension stronger than
  $N_{\mathrm{list}}(C_F)$, and let $p$ be the parameters of the extension of $N$ by the grid
  property for $C_F$. If the open parameter statements
  (Definition~\ref{asm-open-parameter-statements}) hold at $p$, then hypothesis (6) of
  Theorem~\ref{is-definite-rows}, the functional bound for the rows of the selected bundles, is
  linked at $p$.
\item (Open statements.) Assume $P$. For all local Langlands maps and every real $c$, the open
  parameter statements hold at the parameters of dependence $c$.
\item (Covolume.) Assume the four results on automorphic forms, the existence of ray class fields
  and the covolume bound of Theorem~\ref{cv-covolume-bound}. Let local Langlands maps be given and
  let $N$ be an extension stronger than the law extension. If the open parameter statements hold at
  $p(N)$, then hypothesis (5) of Theorem~\ref{is-definite-rows}, the covolume bound for the rows of
  the selected bundles (Remark~\ref{cv-bound-as-used}), is linked at $p(N)$.
\item (Bundles.) Assume $P$ and that the rows of the construction are rational. Let local Langlands
  maps be given, let $c$ be real, and let $A$ be a property of stages that holds eventually in the
  context and at large base data. Then Siegel's mass ratio for CM fields
  (Theorem~\ref{cv-siegel-mass-ratio}) holds, the parameters $p$ of the extension of
  $N_{\mathrm{orig}}(c)$ by $A$ have the member property, and bundles exist at large bases for $p$.
\item (Tate-Shafarevich link.) Assume the three collections of item (6) of
  Theorem~\ref{lit-bundle-unitary}, the passage of depth zero to the cuspidal support (item (2)
  there), and local Tate duality, the finiteness of local Galois cohomology and the local Euler
  characteristic formula (item (5) there). Let local Langlands maps be given and let $c$ be real. If
  the property ``eventually along
  $\mathcal A^{\mathrm{des}}$, the local bound at the twisted places above $2$ and $3$
  (Definition~\ref{asm-local-bound-two-three}(4)) holds'' is linked at the parameters of dependence
  $c$, then so is the property ``the Tate-Shafarevich link (Definition~\ref{asm-twisted-links}) holds
  eventually along $\mathcal A^{\mathrm{des}}$''.
\item (Constructed points.) Let $Q$ be a slot and $\Lambda$ link data, and let
  $\mathcal W(\mathcal A)$ be the working subclass of a class $\mathcal A$
  (Theorem~\ref{pd-period-integrality-along-families}). It is the class $\mathcal A^{\mathrm{des}}$:
  both consist of the families of $\mathcal A$ that satisfy the conditions defining $\mathcal C_N$
  and $\mathcal C_{\mathrm{tw}}$, have chosen constants $(3,C_H)$ for some $C_H$ and satisfy both
  size conditions of Definition~\ref{asm-designed}. The following hold, each under the hypothesis
  named with it.
  \begin{enumerate}
  \item The class number formula of Proposition~\ref{cv-class-number-formula} at every member,
    eventually along $\mathcal A_{N,\mathrm{tw}}$, and the second complement inputs
    (Definition~\ref{asm-division-kernel-errors}(2)), eventually along $\mathcal W(\mathcal A)$, hold
    along $(Q,\Lambda)$.
  \item Let local Langlands maps be given and let $p$ be parameters. Point integrality, twisted point
    integrality, period integrality, twisted period integrality, the second-layer conditions at the
    real place, the bounds at the level primes and the sharp column bounds at the auxiliary prime,
    eventually along $\mathcal W(\mathcal A)$, are linked for $(Q^{\mathrm b}_p,Q_p)$ with $\Lambda$.
  \item Assume the four results on automorphic forms. Then item (3)(a) of
    Definition~\ref{asm-definite-side-hypotheses}, the bounded rows, at every member, holds along
    $(Q,\Lambda)$.
  \item Assume that the mass ratio under volume holds for every definite-side datum
    (Definition~\ref{asm-facts-volume-tie}(2)). Then the volume items (7) of
    Definition~\ref{asm-definite-side-hypotheses} at every member hold along $(Q,\Lambda)$.
  \item Assume the Siegel-Weil formula. Then the frame condition eventually along
    $\mathcal A^{\mathrm{des}}$, the loop rows and the selected form of the orthogonality theorem hold
    along $(Q,\Lambda)$.
  \item Assume global Artin reciprocity (Theorem~\ref{lit-bundle-unitary}(5)). Let local Langlands
    maps be given and let $p$ be parameters. Then the selections are served
    (Definition~\ref{pd-served-selection}) along $(Q_p,\Lambda_p)$.
  \item If the property of open parameter statement (d) holds along $(Q,\Lambda)$, then so does item
    (3)(b), the tie of the twists to the size data, at every member.
  \item If the property of open parameter statement (a) holds along $(Q,\Lambda)$, then so do the
    fullness of the image of the base vectors, eventually along $\mathcal W(\mathcal A)$, and the
    surjectivity of the projections (Projections).
  \item Assume the Siegel-Weil formula, and let $Q^{\mathrm b}$ be a slot. If the property of open
    parameter statement (e) is linked for $(Q^{\mathrm b},Q)$ with $\Lambda$, then so is the
    existence of counting constants with their point counts and of a number $c_I>0$ with the
    irreducibility inputs (items (6) and (11) of Definition~\ref{asm-definite-side-hypotheses}).
  \end{enumerate}
\item (Design.) Assume Siegel's mass ratio for CM fields (Theorem~\ref{cv-siegel-mass-ratio}) and
  that the mass ratio under volume holds for every definite-side datum. Let $Q^{\mathrm b}$ and $Q$
  be slots, and for the constants $\kappa$ of every context let $O_\kappa$ be a property of base
  points such that, eventually in the context, every base point of $Q$ in a context with constants
  $\kappa$ satisfies $O_\kappa$. Let $\mathcal A_O$ be the class of families that are linked to the
  constants $\kappa$ of a context, satisfy the class conditions of the construction, and whose
  members at every scale $x>x_1$ satisfy $O_\kappa$. If every finite subgroup of
  $GL_n(\mathbf Q)$ has order dividing Minkowski's number $M(n)$ for every $n$, and
  $(Q^{\mathrm b},Q)$ has the member property (Definition~\ref{asm-dependence-parameters}(5)), then
  for every function $c_T$ of two real variables, every $\lambda>1$ and every host requirement $G$
  such that hosts satisfying $G$ exist on windows of exponent $\lambda$
  (Definition~\ref{asm-host-requirement}), design choices exist in the class of families of
  $\mathcal A_O\cap\mathcal C_G(\lambda)$ that satisfy the twist condition, with constant $3$ and
  function $c_T$ (Definition~\ref{asm-designed}).
\item (Sizes.) Assume that the mass ratio under volume and the mass formula for genera of rank $2$
  hold for every definite-side datum (Definition~\ref{asm-facts-volume-tie}). For every slot $Q$
  and all link data $\Lambda$, the following holds along $(Q,\Lambda)$, with $c_T$ the function
  given by the Harish-Chandra reduction: eventually along $\mathcal A_{N,\mathrm{tw}}$, the size
  condition on $T_u$ holds whenever $c_T(\varepsilon,C_{YL})\le C_{T_u}$.
\item (Proved facts.) The analytic input (item (5) of Definition~\ref{asm-definite-side-hypotheses}:
  the estimates of Theorem~\ref{hf-zero-density-regions}, the explicit formula and the bound for
  zeros in windows of Theorem~\ref{asm-dependence-zero-sums}), the decay of Wu type for gaps
  (Definition~\ref{is-wu-decay}), the finite dimensionality of rows (the statement of
  Lemma~\ref{asm-dependence-findim} for all its data), the positivity of row types
  and the mass formula for genera of rank $2$ for every definite-side datum hold. Under the
  Siegel-Weil formula, the row non-degeneracy, the row main-term identity and the mass ratio under
  volume for every definite-side datum hold. Under $P$, the conclusions of parts (1) and (2) of
  Theorem~\ref{is-surjectivity-and-lower-norm} hold.
\end{enumerate}
\end{theorem}
\leanref{Development/AutomorphicArithmetic/ThetaCorrespondence/ThetaPairLifts/LocalTypeDictionary/FunctionalsSeparateOfHoweUnramOfSiegelWeil.lean}{UpperTheory.Lit.functionalsSeparate_of_howeUnram_of_siegelWeil}
\leanref{Development/AbelianVarietyTheory/Proofs/SatakeParameters/RhoMemOfMkEq.lean}{GlobalTheory.Hk3.howeUnram_holds}
\leanref{Development/AutomorphicArithmetic/ThetaCorrespondence/ThetaPairLifts/LocalTypeDictionary/TestVectorAtDatumOfJoin.lean}{GlobalTheory.T.test_vector_assembly_holds}
\leanref{Development/AbelianVarietyTheory/Proofs/IwasawaDecomposition/RowRatWOfRowTest.lean}{GlobalTheory.Mem.LiteralAutomorphy.rowRatW_of_rowTest}
\leanref{Development/AutomorphicArithmetic/ThetaCorrespondence/ThetaPairLifts/ThirdPlaceSeries/FourierCoefficientFitOfJoin.lean}{GlobalTheory.BigTop.Fun.Close.fourier_coefficient_fit_of_join_alt}
\leanref{Development/AutomorphicArithmetic/ThetaCorrespondence/ThetaPairLifts/LocalTypeDictionary/EvaluationGridOfJoin.lean}{GlobalTheory.Mem.LiteralAutomorphy.evaluation_grid_of_join}
\leanref{Development/AutomorphicArithmetic/ThetaCorrespondence/ThetaPairLifts/LocalTypeDictionary/TwistedWeilRepresentationConvergence.lean}{GlobalTheory.WeilConv.twisted_weil_representation_convergence}
\leanref{Development/AutomorphicArithmetic/ThetaCorrespondence/ThetaPairLifts/ThirdPlaceSeries/RowFunAtGridA.lean}{GlobalTheory.Fc.StageOne.rowFun_at_gridA_alt}
\leanref{Development/AutomorphicArithmetic/ThetaCorrespondence/ThetaPairLifts/ThirdPlaceSeries/OpenOrbitOfUnramifiedExponents.lean}{GlobalTheory.T.open0_of_unitary_group_hypotheses}
\leanref{Development/AutomorphicArithmetic/ThetaCorrespondence/ThetaPairLifts/ThirdPlaceSeries/LinkedBundlesL3OfOpen.lean}{GlobalTheory.Glob.CovolRow.linked_bundles_l3_of_open_alt}
\leanref{Development/AutomorphicArithmetic/ThetaCorrespondence/ThetaPairLifts/ThirdPlaceSeries/MemberBundlesLargeUnramifiedExponents.lean}{GlobalTheory.T.large_unramified_exponent_bundles_exist_alt}
\leanref{Development/AutomorphicArithmetic/ThetaCorrespondence/ThetaPairLifts/ThirdPlaceSeries/ThetaLiftHypothesesOfNonvanishing.lean}{GlobalTheory.T.hyps_e_third_layer_ser_n_of_geometric_galois}
\leanref{Development/AutomorphicArithmetic/ThetaCorrespondence/ThetaPairLifts/ThirdPlaceSeries/CFinal.lean}{UpperTheory.GalC.final_alt_alt}
\leanref{Development/AbelianVarietyTheory/Proofs/GlobalInterpolation/ClassCountBoundParamsLambdaConstructedPoints.lean}{GlobalTheory.Fc.R0.classCountBound_paramsLambda_constructedPoints}
\leanref{Development/AbelianVarietyTheory/Proofs/GlobalInterpolation/DivisionSecondComplementConstructedPoints.lean}{GlobalTheory.Fc.ClassTwoEntries.divisionSecondComplement_constructedPoints}
\leanref{Development/AutomorphicArithmetic/ThetaCorrespondence/ThetaPairLifts/ThirdPlaceSeries/SixRowsAtMemberStatement.lean}{GlobalTheory.T.period_entries_tie_and_level_valuations_at_built}
\leanref{Development/AutomorphicArithmetic/ThetaCorrespondence/ThetaPairLifts/ThirdPlaceSeries/RowAAtGAtBuilt.lean}{GlobalTheory.Fc.row_aAtG_at_built_alt}
\leanref{Development/AutomorphicArithmetic/ThetaCorrespondence/ThetaPairLifts/ThirdPlaceSeries/RowHRBUniv.lean}{GlobalTheory.Fc.R1.row_hRB_univ_alt}
\leanref{Development/AbelianVarietyTheory/Proofs/GlobalInterpolation/RowHInstROfHU.lean}{GlobalTheory.row_hInstR_of_hU}
\leanref{Development/AutomorphicArithmetic/ThetaCorrespondence/ThetaPairLifts/ThirdPlaceSeries/AtChoice.lean}{GlobalTheory.LocalRows.lg1_sLp_sO_hTF}
\leanref{Development/AutomorphicArithmetic/ThetaCorrespondence/ThetaPairLifts/ThirdPlaceSeries/RowHServOfLinkStatement.lean}{GlobalTheory.T.hServ_of_linkR}
\leanref{Development/AutomorphicArithmetic/ThetaCorrespondence/ThetaPairLifts/ThirdPlaceSeries/RowHTwOfTie.lean}{GlobalTheory.Fc.row_hTw_of_tie_alt}
\leanref{Development/AutomorphicArithmetic/ThetaCorrespondence/ThetaPairLifts/ThirdPlaceSeries/RowJOnOfRowHPA.lean}{GlobalTheory.Fc.JOnBundle.row_jOn_of_row_hPA_alt}
\leanref{Development/AutomorphicArithmetic/ThetaCorrespondence/ThetaPairLifts/ThirdPlaceSeries/LPrjSurjAt.lean}{GlobalTheory.Fc.R1.ZLAt.zL_of_hPA}
\leanref{Development/AutomorphicArithmetic/ThetaCorrespondence/ThetaPairLifts/ThirdPlaceSeries/RowCountsOfLabels.lean}{GlobalTheory.CountRowByName.row_counts_of_labels_alt}
\leanref{Development/AbelianVarietyTheory/Proofs/IwasawaDecomposition/MemberTopFinalContextParameters.lean}{GlobalTheory.Fc.memberTop_final_contextParameters}
\leanref{Development/AutomorphicArithmetic/CharacterBoxes/ObjCTrue.lean}{UpperTheory.Top.recS}
\leanref{Development/AutomorphicArithmetic/ThetaCorrespondence/ThetaPairLifts/ThirdPlaceSeries/RowHTuCTminStatement.lean}{GlobalTheory.T.hTu_cTmin_of_genusMass}
\leanref{Development/AbelianVarietyTheory/Proofs/GlobalInterpolation/LitThin.lean}{GlobalTheory.ThinScaleLiterature.litThin}
\leanref{Development/AbelianVarietyTheory/Proofs/FlagMeasures/WuDecayOfGap.lean}{GlobalTheory.wuDecayOfGap}
\leanref{Development/AbelianVarietyTheory/Proofs/FlagMeasures/RowsFinDimLink.lean}{GlobalTheory.rowsFinDim_link}
\leanref{Development/AbelianVarietyTheory/Proofs/FlagMeasures/RowsPosTypeP.lean}{GlobalTheory.rowsPos_typeP}
\leanref{Development/AbelianVarietyTheory/Proofs/FlagMeasures/RowPositivityOfSiegelWeil.lean}{GlobalTheory.rowNGP_of_siegelWeil}
\leanref{Development/AbelianVarietyTheory/Proofs/FlagMeasures/RowMainEqPOfSiegelWeil.lean}{GlobalTheory.rowMainEqP_of_siegelWeil}
\leanref{Development/AbelianVarietyTheory/Proofs/FlagMeasures/UnitaryHyperplaneVolumeRatioSiegelWeil.lean}{GlobalTheory.VolRatioFromSiegelWeil.volumeRatio}
\leanref{Development/AbelianVarietyTheory/Proofs/GlobalInterpolation/GenusMassAtTwo.lean}{GlobalTheory.genusMassAt_two}
\begin{proof}
Among the results in $F$ are Siegel sets for unitary groups and Howe duality at the real place
(Theorem~\ref{lit-howe-duality-real}). Each part combines the results it assumes as follows.

(1) Let $\varepsilon$ and $C_F\ge6510$ be given, and let $x_0$ be the scale of
Lemma~\ref{asm-dependence-rows-weak}. Let $D$, $R$, an evaluation grid for $R$ with constant $C_F$ at
$b=e^{x^\varepsilon}$, a Weil datum over $R$, a row volume and a nonzero smooth vector $w$ of the
model $M_0$ of $U(Y_0)$ in $R$ be as in Definition~\ref{asm-grid-evaluation}(7), with $x\ge x_0$. The
grid is a theta grid with constants $(1,C_F,b)$, and a Weil datum over $R$ has the archimedean family
$\theta_R$ of $R$ and its Cartan involutions on $Y_0$; so $M_0$, whose parameter is
$S^{\mathrm{row}}_0$, is a row model at $0$ (Definition~\ref{asm-dependence-rows-models}(3)). Let $a$
be the function of Lemma~\ref{asm-dependence-rows-factor} for $w$. By
Lemma~\ref{asm-dependence-rows-good}, every theta multiplicity $t$ of $R$ at the trivial character is
a row multiplicity of $M_0$, nonzero as a linear map if $t\ne0$, and $u_{w,t}=t(w)$.

If every such $t$ is $0$, any $\rho$ and $\xi$ will do. Otherwise fix $t_0\ne0$.
Lemma~\ref{asm-dependence-rows-weak} gives $\rho_0$, $\xi_0$, $\Psi_0$ and $\Phi_0$ with
$B^+(\varphi_{\xi_0,\Psi_0},\Theta_{\rho_0,\Phi_0}(u_{w,t_0}))\ne0$. By
Lemma~\ref{asm-dependence-rows-factor} this number is $a(\rho_0,\xi_0)l_0$ for some $l_0$, so
$a(\rho_0,\xi_0)\ne0$. Now let $t\ne0$. Lemma~\ref{asm-dependence-rows-weak} gives $\rho$, $\xi$,
$\Psi$ and $\Phi$ with $B^+(\varphi_{\xi,\Psi},\Theta_{\rho,\Phi}(u_{w,t}))\ne0$, and
Lemma~\ref{asm-dependence-rows-factor} gives $l$ with
$B^+(\varphi_{\xi',\Psi},\Theta_{\rho',\Phi}(u_{w,t}))=a(\rho',\xi')l$ for all $\rho'$ and $\xi'$.
Hence $l\ne0$ and $B^+(\varphi_{\xi_0,\Psi},\Theta_{\rho_0,\Phi}(u_{w,t}))=a(\rho_0,\xi_0)l\ne0$. The
lift and the period converge by Lemma~\ref{asm-dependence-rows-global}(1). So functionals separate $w$,
with $\rho_0$ and $\xi_0$. The tie of the series idempotents to the labels and the record of the group
are not used. The lemmas assume the results of Theorem~\ref{lit-bundle-theta}, the Siegel-Weil formula
and Howe duality for unramified data, which holds by Lemma~\ref{asm-dependence-howe-unramified}.

(2) By Lemma~\ref{asm-dependence-test-assembly}, the row test holds at all large stages after the
pinned column choices, for the realizations described there. A column choice that is pinned,
supported and coherent with the twist data is pinned, so the property of that lemma also holds at
all large stages after these column choices. Let $Z$ and $z$ be as in
Definition~\ref{asm-stage-construction}(4). Among the conditions imposed there are that the group of
$z$ is trivial at the stage, that the group stage satisfies its conditions, which include the record
of the group data, and that $z$ exists at every place and lies in the geometry of the stage. So the
complex data of $z$ at its face pass the row test; the further conditions on $\kappa$, on the core,
on the jet stage and on the twisted places are not used. This is (2).

(3) $F$, the join statements, the functional equation, the agreement of theta integrals, the local
matching, part (1) and the lower norm bound of Theorem~\ref{is-surjectivity-and-lower-norm}(2),
through the fitted function at zero.

(4) Lemma~\ref{asm-dependence-grid-exists} assumes the results named in (4); let $C_{F_1}$ be its
constant, and let $C_F\ge C_{F_1}$. Let $0<\varepsilon\le1/100$, let the constants of a context be
given, let $x_0$ be the number of that lemma for $\varepsilon$ and $C_F$, and let $K$ be a context
with these constants at a scale $x\ge x_0$. For base data $B$ with a framing $f$ and twist data $T$,
let $\Xi$ be the set consisting of the trivial character, the characters $\delta$ and $\delta'$ of
$f$, and the characters $\delta_\ell\psi_w$, where $w$ runs over the finite places of the coefficient
field, $\psi_w$ is the character that $T$ chooses at $w$, and $\delta_\ell=\delta'$ if the residue
characteristic $\ell$ of $w$ is $p_{de}$ and $\delta_\ell=\delta$ otherwise. It is the set of
Definition~\ref{asm-stage-construction}(5). It is finite: $\psi_w$ is trivial unless the residue
characteristic of $w$ is a twisting prime, so that then $\delta_\ell\psi_w$ is $\delta$ or $\delta'$,
and only finitely many places lie above the finitely many twisting primes. Let $S$ be the finite set of algebraic numbers
that the lemma gives for $B$, $f$ and $\Xi$, and take $S$, with no primes, as the finite sets of the
condition at large base data (Definition~\ref{asm-stage-construction}(2)). If the coefficient field
of $B$ contains $S$, the lemma gives the conclusion of the grid property for $C_F$ at every later
stage of $(K,B,f,T)$. So the grid property for $C_F$ holds eventually in the context and at large
base data: the grid with constant $C_F$ evaluates.

(5) Let $C_{F_0}$ be the constant $C'$ of the conclusion of part (3). Let $C_F\ge C_{F_0}$, local
Langlands maps, $N$ and $p$ be as in (5), and suppose that the open parameter statements hold at $p$.
By Lemma~\ref{asm-dependence-functional-geometric}, from the ray class field statement, the division
rows of (b) and statement (d), the property $G$ of that lemma is linked at $p$. By
Lemma~\ref{asm-dependence-linked}(2), the conjunction of statement (a), the division rows, statement
(f) and $G$ is linked at $p$: there is a slot $R$ that, eventually in the context and at large base
data, every core built from $p$ satisfies, and the conjunction holds along $(S,\Lambda_p)$, with a
condition $E$ on bundles, for the slot $S$ of the good cores of $p$ that satisfy $R$. It suffices to
show that hypothesis (6) of Theorem~\ref{is-definite-rows} holds along $(S,\Lambda_p)$ with the
same $E$.

Let $s$ be a selection for a class $\mathcal A$ whose bundles are built from $(S,\Lambda_p)$ and
satisfy $E$. By $G$, choose at every member a comparison part of the selected bundle whose point has
the normalized embedding of its setting, satisfies the class-sum pin, the compatibility conditions
and the further conditions, and has an evaluation grid with constant $C_F$ at $b$ fitted to it, with
a Weil datum over its realization and the remaining choices of the slope test; this is the selected
part. Since the selected bundles are built from $(S,\Lambda_p)$, the families of $\mathcal A$ satisfy
the class conditions of the construction, so each of them carries its duty data
(Definition~\ref{asm-duty}). By (a), the realization of the selected part has primitive rows and its
family of definite forms has lift-distinct types; as in the proof of Theorem~\ref{is-definite-rows},
the division rows give the divisor conditions of the selected part. With (f), Lemma~\ref{asm-dependence-covolume-points}
shows that eventually along $\mathcal W(\mathcal A)$, eventually along each family, the point of
the selected part satisfies the point conditions. The evaluation grid of the selected part exists at
every member, hence eventually. Lemma~\ref{asm-dependence-functional-bound} then gives, eventually
along $\mathcal W(\mathcal A)$, for every family with the division bound at twisted places, a
constant $C_0\ge0$ such that for every $C\ge C_0$, eventually along the family, the rows of the
realization of the selected part satisfy the functional bound for $N(C,C_H)$. That realization is
the realization of the selected bundle, and $\mathcal W(\mathcal A)$ is the class
$\mathcal A^{\mathrm{des}}$ (part (10)), so this is hypothesis (6) of
Theorem~\ref{is-definite-rows}.

(6) Let $c$ be real, let $N=N_{\mathrm{list}}(c)$, and let $p=p(N)$ be the parameters of dependence
$c$. The list extension requires the law extension, so $N$ is stronger than the law extension, and
the datum of every stage built from $p$ satisfies the rationality law. The results named in
Lemmas~\ref{asm-dependence-orbit-inputs} and~\ref{asm-dependence-orbit-geometric} are in $P$; Green's
parametrization follows from Green's description of the cuspidal characters
(Corollary~\ref{lit-green-labelled}), and Lemma~\ref{asm-dependence-howe-unramified} gives the
remaining hypothesis of part (3) of Lemma~\ref{asm-dependence-orbit-inputs}. By
Lemma~\ref{asm-dependence-orbit-witness}, statement (a) and the compatibility of the division data
with the periods are linked at $p$. In Lemma~\ref{asm-dependence-orbit-inputs}, part (1) gives (d),
and part (2) then gives (c). Part (3) gives the multiplicity condition at every member and the
multiplicity condition of (h); with the first, part (4) gives the division rows. Part (5) gives the
rows with $\Lambda_3$, and part (6) then gives the equality of the archimedean scalars of (h). Part
(7) gives the goodness of the rows at every member and statement (f), and part (8) then gives (e). By
Lemma~\ref{asm-dependence-orbit-geometric}, the division rows and (d) give (g). Statement (i) holds
at any parameters, by Theorem~\ref{is-definite-side-condition-holds}. So the open parameter
statements hold at $p$.

(7) Let local Langlands maps be given, let $N$ be an extension stronger than the law extension, let
$p=p(N)$, and suppose that the open parameter statements hold at $p$. By
Lemma~\ref{asm-dependence-covolume-ties}, from the ray class field statement and the division rows of
(b), the property $T$ of that lemma is linked at $p$. By Lemma~\ref{asm-dependence-linked}(2), the
conjunction of statement (a), the division rows, statements (f) and (g) and the property $T$ is
linked at $p$: there is a slot $R$ that, eventually in the context and at large base data, every
core built from $p$ satisfies, and the conjunction holds along $(S,\Lambda_p)$, with a condition $E$
on bundles, for the slot $S$ of the good cores of $p$ that satisfy $R$. It suffices to show that
hypothesis (5) of Theorem~\ref{is-definite-rows} holds along $(S,\Lambda_p)$ with the same $E$.

Let $s$ be a selection for a class $\mathcal A$ whose bundles are built from $(S,\Lambda_p)$ and
satisfy $E$. By (g), choose at every member a comparison part of the selected bundle whose point has
the normalized embedding of its setting and satisfies the class-sum pin, the compatibility conditions
and the further conditions; this is the selected part. Since the selected bundles are built from
$(S,\Lambda_p)$, the families of $\mathcal A$ satisfy the class conditions of the construction, so
each of them carries its duty data (Definition~\ref{asm-duty}). The realization of the selected part
is that of the selected bundle, so by (a) it has primitive rows and its family of definite forms has
lift-distinct types. As in the proof of Theorem~\ref{is-definite-rows}, every tuple of the selected
part is the base tuple or the tuple $\mu_w$ of a place, so its character is used, and the division
rows give the divisor conditions of the selected part. With (f) and the class-sum pin and the
compatibility conditions of the selected part, Lemma~\ref{asm-dependence-covolume-points} shows
that eventually along $\mathcal W(\mathcal A)$, eventually along each family, the point of
the selected part satisfies the point conditions. By $T$, at every member, and so eventually along
$\mathcal W(\mathcal A)$, the datum with its size data satisfies the coefficient field bound.
Lemma~\ref{asm-dependence-covolume-bound}, with the covolume bound of Theorem~\ref{cv-covolume-bound},
then gives the covolume bound for the rows of the realization of the selected part, which is that of
the selected bundle. Since $\mathcal W(\mathcal A)$ is the class $\mathcal A^{\mathrm{des}}$ (part
(10)), this is hypothesis (5) of Theorem~\ref{is-definite-rows} (Remark~\ref{cv-bound-as-used}).

(8) Let local Langlands maps be given, let $c$ be real, and let $A$ be a property of stages that
holds eventually in the context and at large base data. Write $p_{\mathrm{orig}}=p(N_{\mathrm{orig}}(c))$,
and $p$ for the parameters of the extension of $N_{\mathrm{orig}}(c)$ by $A$. Siegel's mass ratio for
CM fields holds by Theorem~\ref{cv-siegel-mass-ratio}. The results named in
Lemma~\ref{asm-dependence-bundles-inputs} are in $P$, and the rows of the construction are rational
by hypothesis. By part (1) of that lemma, cuspidal
multiplicities are equal at $p_{\mathrm{orig}}$, so by part (3) $p_{\mathrm{orig}}$ has the member
property, and by Lemma~\ref{asm-dependence-bundles-member} so has $p$. The origin extension requires
the list extension, and so the law extension, and by part (2) the positions of $p_{\mathrm{orig}}$
move with the division data. So part (4), applied to $N_{\mathrm{orig}}(c)$ and $A$, shows that
bundles exist at large bases for $p$.

(9) Let local Langlands maps be given, let $c$ be real, let $N=N_{\mathrm{list}}(c)$, and let $p=p(N)$
be the parameters of dependence $c$; $N$ is stronger than the law extension. Suppose that the local
bound at the twisted places above $2$ and $3$, eventually along $\mathcal A^{\mathrm{des}}$, is
linked at $p$. Alpöge's argument has five steps.

\emph{Step 1.} By Lemma~\ref{asm-dependence-inputs-sha-published}, the first and second collections
of item (6) of Theorem~\ref{lit-bundle-unitary} give Alpöge's published inputs of the third layer and
his series statements.

\emph{Step 2.} By Lemma~\ref{asm-dependence-inputs-sha-preconditions}, from these, the second
collection, the passage of depth zero to the cuspidal support and the three local results, the Galois
preconditions on a datum and its norm choice and the fourth tie of the datum hold at the layer stages,
and so does their conjunction.

\emph{Step 3.} By Lemma~\ref{asm-dependence-inputs-sha-transport}(3), applied to this conjunction,
the property ``eventually in the context, every bundle that $\Lambda_p$ links to a base point of a slot
satisfies the Galois preconditions and the fourth tie'' is linked at $p$.

\emph{Step 4.} By Lemma~\ref{asm-dependence-inputs-sha-cut}, with the Chebotarev density theorem
among the published inputs of the third layer and the three local results, every datum with a norm
choice that satisfies the Galois preconditions and the fourth tie has a good cut, at which the
norm-choice facts and the fourth tie hold.

\emph{Step 5.} By Lemma~\ref{asm-dependence-inputs-sha-link}, with the third collection, Steps 3 and
4 and the hypothesis on the local bound at $2$ and $3$, of which only the number $\varepsilon_1$ and
the largeness thresholds are used, the Tate-Shafarevich link eventually along
$\mathcal A^{\mathrm{des}}$ is linked at $p$.

(10) Parts (a) and (d) hold along every slot, (d) by the mass ratio under volume. In (a), the class
number formula of Proposition~\ref{cv-class-number-formula} holds for every definite-side datum and
all its level data. So, for every class $\mathcal A$, it holds at every member of every family, and
hence eventually along $\mathcal A_{N,\mathrm{tw}}$, with any $\varepsilon_1>0$ and any largeness
threshold. A property of classes that holds for every class holds along every $(Q,\Lambda)$, with
the condition on bundles that every bundle satisfies (Definition~\ref{is-slot-rows}).
For the second complement inputs in (a), let $\mathcal A$ be a class and $\mathcal F$ a family of
$\mathcal W(\mathcal A)$ with exponent $0<\varepsilon\le1/100$; it satisfies both size conditions of
Definition~\ref{asm-designed}. Conditions (a) and (b) of
Definition~\ref{asm-division-kernel-errors}(2) hold at every member at every scale $x>x_1$, by
Lemma~\ref{asm-dependence-complement-norms}. Since $\varepsilon<1/52$, there are constants
$C_1,\dots,C_6$, with $C_2\ge0$, and thresholds beyond which conditions (c) to (h) hold at every
member: (c), (g) and (h) for every family of exponent $\varepsilon$, (d) under the size condition on
stabilizer denominators, with $C_2$ chosen after $C_1$, (e) under the size condition on $T_u$, and (f)
under both. With $x_0$ the largest threshold, $\mathcal F$ satisfies the second complement inputs.
So they hold eventually along $\mathcal W(\mathcal A)$, with $\varepsilon_1=1/100$ and any largeness
threshold, for every class, and hence along every $(Q,\Lambda)$ as above.
In part (b), each property is linked with the slot $R$ of
Definition~\ref{asm-dependence-parameters}(2) that every core satisfies, so it suffices that it holds
along $(Q_p,\Lambda)$. The good cores of $p$ lie over true columns
(Remark~\ref{pd-integrality-inputs}). So period integrality and twisted period integrality hold,
eventually along $\mathcal W(\mathcal A)$, along $(Q_p,\Lambda)$ by
Theorem~\ref{pd-period-integrality-along-families}. The other properties of (b) are proved along the
slots over true columns as well. Part (c) is proved from the covering of discrete series,
the finiteness of Harish-Chandra modules, the rapid decay of cusp forms and the Casselman-Wallach
theorem; part (e) from the Siegel-Weil formula; part (f) from global Artin reciprocity and the
reciprocity records of the linked bundles; part (g) from the comparison parts; part (h) from the
primitive rows; and part (i) from the Siegel-Weil formula and the labels.

In (d), the volume items hold at every definite-side datum, so for every class they hold at every
member, and they hold along every $(Q,\Lambda)$ as in (a). Let $D$ be a definite-side datum with host
fields $F\subset E$, and let $D'$ be the datum $D$ with each volume $\operatorname{vol}_D(W)$
replaced by the volume $\operatorname{vol}(W)$ of Definition~\ref{cv-volume-mass-normalized}. The
volume measures of $D'$ are those of $D$, and by Lemma~\ref{cv-volume-mass-invariance}(1) the volume
tie holds for $D'$: every volume measure gives every fundamental domain of $U(W)(F)$ the measure
$\operatorname{vol}(W)$. For item (7)(a), let $W$ and $W_0$ be Hermitian spaces over $E$ with
correct signature and the same rank $M$, of signatures $(P,Q)$ and $(P_0,Q_0)$ at every complex
embedding, with volume measures $\mu$ and $\mu_0$ and fundamental domains $\mathcal F$ and
$\mathcal F_0$. Then $\mu(\mathcal F)=\operatorname{vol}(W)$ and
$\mu_0(\mathcal F_0)=\operatorname{vol}(W_0)$. If $M=0$, then $W$ and $W_0$ are isometric and
$\operatorname{vol}(W)=\operatorname{vol}(W_0)$ by Lemma~\ref{cv-volume-mass-invariance}(3). If
$M\ge1$, Steps 1 to 3 of the proof of Theorem~\ref{cv-volume-mass}, with the ratio of volumes used
there supplied by the mass ratio under volume for $D'$, give
$\operatorname{vol}(W)=\operatorname{mass}(I_{P+Q})$ and
$\operatorname{vol}(W_0)=\operatorname{mass}(I_{P_0+Q_0})$, and $P+Q=M=P_0+Q_0$. In both cases
$\mu(\mathcal F)=\mu_0(\mathcal F_0)$. For item (7)(b), let $W$ be positive definite of rank $n_0$,
with standard lattice $L$ and $\Gamma=U(W)(F)$. Then $\operatorname{mass}(L)=\operatorname{mass}(I_{n_0})$
by Lemma~\ref{cv-volume-mass-definite}(1). As in the proof of that lemma, $U(W)(\mathbf A_F)$ is the
disjoint union of finitely many double cosets $\Gamma g_iK$, and for these representatives
$\sum_i|\Gamma\cap g_iKg_i^{-1}|^{-1}=\operatorname{mass}(L)$. The sum does not depend on the
representatives: any other system can be numbered so that its $i$-th element is $\gamma_ig_ik_i$
with $\gamma_i\in\Gamma$ and $k_i\in K$, and then its $i$-th group is
$\gamma_i(\Gamma\cap g_iKg_i^{-1})\gamma_i^{-1}$.

In (g), let the property of open parameter statement (d) hold along $(Q,\Lambda)$ with a condition
$E$ on bundles; the same $E$ serves for item (3)(b). Let $s$ be a selection for a class $\mathcal A$
whose bundles are built from $(Q,\Lambda)$ and satisfy $E$, and let $D$ be a member of a family of
$\mathcal A$ at a scale $x>x_1$, with size data $L$, selected base level $P$ with rows $S_0$, and
indefinite realization $Q'$ of the selection bundle, with realization $R$. By (d) the selected
bundle has a comparison part. It carries a recursion for $Q'$, which gives a map $t$ from the
tuples of $D$ to the tuples of the family of definite forms of $P$ with $t(\mu_w)$ the tuple that
$R$ attaches to $w$, for every place $w$, and with $t(\mu_0)$ the base tuple of that family: both
have the trivial Hecke character, and the Hecke character determines the tuple. The character data
of $R$ attach to every twisted place of $R$ a tuple different from the base tuple. So if $w$ is
twisted for $R$, then $t(\mu_w)\ne t(\mu_0)$ and $\mu_w\ne\mu_0$. The twisted places of $R$ are
those of $S_0$. Finally, a family records that at every member at a scale $x>x_1$ every place $w$
with $\mu_w\ne\mu_0$ lies in the support of the size data. Hence every twisted place of $S_0$ is a
place $w$ of the support of $L$ with $\mu_w\ne\mu_0$, which is item (3)(b).

In (h), the surjectivity of the projections holds along every $(Q,\Lambda)$ without the hypothesis,
with the condition on bundles that every bundle satisfies. Let $s$ be a selection for a class
$\mathcal A$, and let $D$ be a member of a family of $\mathcal A$ at a scale $x>x_1$ beyond the
threshold of $s$, with selected base level $P$. The square-integrable columns of the selected bundle
give the family conditions for the family of definite forms of $P$. So for every lift $i$,
Lemma~\ref{asm-dependence-projections-onto}, applied to the lift forms at $(i,\mu_0)$, the characters
of the family, the subgroup $P_0$ with the character $\psi_0^{-1}$ and the definite subspace of $P$ at
$(i,\mu_0)$, shows that the projection matrix of $P$ at $(i,\mu_0)$ is surjective. This holds at every
member at every such scale, so eventually along $\mathcal A^{\mathrm{des}}$. For the fullness of the
image of the base vectors, let the property of open parameter statement (a) hold along $(Q,\Lambda)$
with a condition $E$ on bundles; the same $E$ serves. Let $s$ be a selection for a class $\mathcal A$
whose bundles are built from $(Q,\Lambda)$ and satisfy $E$, and let $D$ be a member of a family of
$\mathcal W(\mathcal A)$ at a scale $x\ge x_1+1$. The selected bundle has a row-side family $P$,
whose family of definite forms satisfies the family conditions with the conditions on its inert
indices, and an indefinite realization of $P$, whose realization has primitive rows at its face; and
the family has lift-distinct types, by (a). Lemma~\ref{asm-dependence-base-vectors} shows that
$\iota_{\mu_0}$ maps $V_0(\mu_0)$ onto $\mathcal E$.

(11) Siegel's mass ratio for CM fields and the mass ratio under volume, through the existence of
the objects of the construction at every large scale.

(12) The mass ratio under volume and the mass formula for genera of rank $2$.

(13) The analytic input consists of the three estimates of Theorem~\ref{hf-zero-density-regions},
with its constants $C_1,c$, and the two parts of Theorem~\ref{asm-dependence-zero-sums}, with its
constant $C$; none of them depends on a definite-side datum. The finite dimensionality of rows is
Lemma~\ref{asm-dependence-findim}.

The mass formula for genera of rank $2$ holds for every definite-side datum $D$. The hypothesis of
Lemma~\ref{asm-dependence-genus-mass} holds: the host field of every definite-side datum is a CM
field over which no finite prime of its totally real subfield ramifies
(Proposition~\ref{hf-host-field-properties}), so Theorem~\ref{cv-volume-mass-landherr} gives every
positive definite Hermitian space over it with correct signature an orthonormal basis. That lemma
gives the mass formula for genera of rank $2$ for $D$.

Under the Siegel-Weil formula, the mass ratio under volume holds for every definite-side datum. The
Siegel-Weil formula gives Siegel's formula (Theorem~\ref{cv-volume-mass-siegel-weil}) for every CM
field and every nondegenerate Hermitian form, and Proposition~\ref{cv-volume-mass-ratio}, whose
proof rests on Siegel's formula, gives, for a datum
whose volume tie holds, a vector $x$ of norm $\pm1$ as in Definition~\ref{asm-facts-volume-tie}(2)
and the ratio of the volumes of $W$ and $x^\perp$; the ratio required there is its reciprocal, since
$\big((M-1)!\,(v/\pi)^M\big)^{-f}(2v)^{f/2}=\big((2\pi)^M/(M-1)!\big)^f(2v)^{f(1/2-M)}$.

Under $P$, the conclusions of parts (1) and (2) of Theorem~\ref{is-surjectivity-and-lower-norm} are
that theorem.
\openstep{In parts (1) to (5), these proofs are not yet written out: the factorization of the theta
lifts of multiplicity maps through an intertwiner of the local models,
Lemma~\ref{asm-dependence-rows-step-models}; the core periods and the core surjectivity, parts (2)
and (3) of Lemma~\ref{asm-dependence-rows-global}; the irreducibility of the dual of the coefficients
under the real places with the continuity of its matrix coefficients in
Lemma~\ref{asm-dependence-rows-coefficients}; the convergence of the twisted theta integrals, the
nonvanishing of the squares, the sum conditions, the formula for the test functionals and the
rationality of the sum values, Lemma~\ref{asm-dependence-test-assembly-first-layer-rows-pieces}; the
construction of a grid of models matching the test choices,
Lemma~\ref{asm-dependence-test-assembly-first-layer-convergence-lift-grid}, and the cuspidality of the
theta lifts of pairs along it,
Lemma~\ref{asm-dependence-test-assembly-first-layer-convergence-lift-cuspidal}; that the standard
functions of the bottom triples are theta functions, and the towers at the squares,
Lemma~\ref{asm-dependence-test-assembly-first-layer-ratio-inputs}; the theta grids on given places, Weil data
and grid choices of Lemma~\ref{asm-dependence-test-assembly-first-layer-grid-pieces}; the action of
the Weil data of the chain at pairs of level lifts on the test vectors,
Lemma~\ref{asm-dependence-test-assembly-realization-equivariant-scalars}; the
existence of Weil data, of evaluation
choices and of the number fields of their numbers, Lemma~\ref{asm-dependence-grid-pieces}, the
transport of grids, Weil data, grid choices and numbers, Lemma~\ref{asm-dependence-grid-transport},
and the convergence of the theta integrals of the twisted Weil representations used in the proof of
Lemma~\ref{asm-dependence-grid-exists}. The proof of part (3), the fit of the constant function at
zero, is not yet broken into stated lemmas.}
% Parts (6) to (9) are written out down to Lemmas asm-dependence-orbit-inputs-*,
% asm-dependence-covolume-ties-points, asm-dependence-covolume-points-*, asm-dependence-bundles-inputs-move,
% the proof of Lemma asm-dependence-bundles-inputs and Lemmas asm-dependence-inputs-sha-*, which carry
% the remaining open steps. In (9), final_alt_alt = final_of_alt: Step 1
% hyps_e_third_layer_ser_n_of_geometric_galois (PublishedInputsThirdLayerAlt, SerN); Step 2
% galois_preconditions_at_mem and tieChi_at_members_of_depthZero, combined by Model.AtMem.and_z; Step 3
% witnesses_pass_along_set_division of slot_mv_extension_with_list_geo_cl and the second part of
% atMem_at_wit_alt; Step 4 exists_good_cut_facts_tie4 (CutGivesAlt); Step 5
% linkedAt_linkN_of_boundAtTwoThree.
\openstep{In parts (10) to (13) these intermediate results are not yet stated as lemmas: conditions
(c) to (h) of the second complement inputs in (10)(a); point integrality, twisted point integrality,
the second-layer conditions at the real place, the sharp column bounds at the auxiliary prime and the
bounds at the level primes in (10)(b); the bounded rows of (10)(c); the frame condition, the loop
rows and the selected orthogonality theorem of (10)(e); the served selections of (10)(f); the
counting constants and the irreducibility inputs of (10)(i); the two passes that give the families
with their design choices in (11); the Harish-Chandra reduction for the size condition in (12); and,
in (13), the decay of Wu type, the positivity of row types, the row non-degeneracy and the row
main-term identity.}
% Lean, part by part: (1) real_functionals_separate_of_hypothesis, written out from Lemmas
% asm-dependence-rows-* (rowTower_of_step, weakSep_of_steps, factor_gap, goodRows, separate_of_burGapPos);
% written out: Lemma asm-dependence-rows-local (1) and (2) (separating_family_has_models =
% coreModels_of_upq and separating_models_of_separating_data, through Lemmas asm-dependence-rows-counts
% and asm-dependence-rows-model-exists) and Lemma asm-dependence-rows-global (1)
% (rapid_decay_of_separating_family and convergence_of_separating_family, through Lemmas
% asm-dependence-rows-cusp and asm-dependence-rows-bessel-conv); Lemma asm-dependence-rows-local (3)
% written out from separatingFamily_extends_by_twisted_row with pairPar_X, exp_att and the vacuum and
% coset vectors of the grid Weil datum, down to Lemma asm-dependence-rows-step-lift (lifts_neg), written
% out from Lemmas asm-dependence-rows-step-negated (negModel_tr, proved), asm-dependence-rows-step-frames
% (ptF_frame_change, a computation) and asm-dependence-rows-step-models, which is open: PairEx =
% exists_seesaw_pair (ItPairObj.exists_record,
% thetaModels_eq_coverDiscreteSeries, exists_locIt_pair, exists_globalFockIntertwiner_tensor_realPlaces,
% thetaVacuumCharacters_atRealPlaces, fockVectorsOverVacuumIsometry_of_skeletonDatum, isItW_of_globIt),
% PairLifts = RuleLift.pairLifts (cuspAll_of_lowAll, lowAll_of_lowFock, lowFock) and rulePar_of_pairPar
% with the Rallis inputs of
% exists_record_theta_all and inputs_of_chosenSeries; Lemma asm-dependence-rows-global (4) written out
% (topLine_gap, Gap.fit_gap = Lemma asm-dependence-rows-gap, sepLine_gap and burGap_model = Lemma
% asm-dependence-rows-bessel-line, BurGapPos cited to JSZ10, bilin_prop), with Lemma
% asm-dependence-rows-bessel-separate (theta_lifts_separate_of_separating_data) written out from Lemma
% asm-dependence-rows-bessel-period, written out (bplus_cover, bplus_add_bound, bplus_rTr, bplus_rTr_unip,
% mult_bound via sobolev_unif_bound, lastVertex_besselDatum_isFrame_isGeneric with its real-place
% computations, topCfg), and Lemma asm-dependence-rows-coefficients (coefTensor) written out but for
% dual_irreducible_coefficientSpace_realPlaces and contragredient_matrix_coefficient_continuousOn_unitary;
% open: Lemma asm-dependence-rows-global (2) and (3) (Branch.periodNe_of_eigen with realInputE_holds,
% about 82,700 lines; twist_core_surjective); (2)
% rowRatW_of_rowTest, written out; Lemma asm-dependence-test-assembly written out from
% test_vector_at_datum_of_join (Lemma asm-dependence-test-assembly-first-layer, written out from its
% -rows, -convergence, -columns, -injective, -ratio, -grid and -field lemmas; written out: -columns,
% -injective (injK_of_separates_linB) and -field (exists_extension_field_of_join); -rows,
% -convergence, -ratio and -grid written out down to -rows-pieces (weilConvergence_of_siegelSet,
% constK_alt, sumOKK_of_theta_local_inputs, rowTestLiftFormula_principalLevel and its four steps,
% sumRatK_of_sumOKK), -convergence-lift (twist_row_lift_exists, thetaPairLift_cuspidal_of_hypotheses),
% -ratio-inputs (LNe, written out through Lemma -ratio-inputs-nonvanishing (heckeL_ne_zero_int) but
% for fieldDat and isValChar_nu; TriEx, bottom_triple_integral_eq_const_mul_local, Tower,
% ratIn_of_definite) and
% -grid-pieces (exists_places_of_cls, nonempty_gridData_iff, exists_weil_over_of_clauses,
% gridChoice_nonempty_of), which are open) and Rw.RowAt.atDat_rowTest (Lemma
% asm-dependence-test-assembly-second-layer, written out, with rowTest_realisation, dat_choice_exists,
% evG_exists_sizeRows, hostHyp_frames_ev, exists_bare and multEq_rowTest_of_copy written out;
% EX.exists_sizeData_explicit (Lemma -sizes-explicit) written out down to the computations
% bottom_window_le and windows_of_explicitParameters; LiftGroupEquiv.th_thG_act (Lemma
% -realization-equivariant) written out with Lemmas -realization-equivariant-levels, -lifts, -shift,
% -integral, -projection and -schur; open: the case P = 1 of Lemma -realization-equivariant-scalars
% (StepVecLaw.stepVec_lifts, StepLaw.last_step_law); Lemma -realization-equivariant-schur-slot
% (SmVec.schurLemma_equivariant_pair_boundedFunctions) is written out down to Schur's lemma (-closable-schur) and
% Garding [Gar47]; and Mem.Top.exists_nStage_grpTrue (Lemma
% -bare-norm, written out with locStage_inhabited cited to Fontaine, GLm.isSer_eSerC as Lemma
% -bare-norm-series and TypesInhab.types0 in its Step 3; Lemma -bare-norm-labels(1),
% GLq.Label.label_main, written out down to Lemmas -bare-norm-labels-series-induction and -mackey)); (3)
% fittedFunction_at_zero_close_alt, open: square_leading_term_bound,
% sum_size_bound_all_places, local_sum_estimates_joined, injK_of_separates with
% twisted_convergence_with_real_places, and fun0Fit_of_leaves_alt with its twenty leaves
% (existsFunctionIn_of_sqRatio_range, leafKernel_tie_of_claim_inputs, leafInfinitePlaceBound_of_lamLow and
% the rest); (4) evaluation_grid_of_join = evLgG_gridA_of_exGrid, written out from Lemma
% asm-dependence-grid-exists (ExGrid; epsilon_1 = 1/100, the algebraic numbers chosen from B.toPre, f and
% T.set f, finite by tw_fin); its proof exGrid_of is written out, with KeyFin and GKeyFin (Lemma
% asm-dependence-grid-keys) proved; open: Lemma asm-dependence-grid-pieces, that is WeilEx
% (exists_weilDeligne_of_localInputs), IsEx (isEx_of_separates, twisted_linear_bound,
% twisted_base_existence, twisted_constant_existence), NFAll (nfAll_of: nfChars, nfRoot,
% nfPhase_of_sumOKK, nfJet_holds, twistRatio_mem_numberField, bottomField_isNumberField) and NFSq
% (nfSq_of_fieldDat); Lemma asm-dependence-grid-transport (carryG, carryW_of_weil, carryS, trEasy,
% trPhase_of_trFlip and trJet_of_trFlip through flippedTensorsTransferToGrid_proof, trBot_holds,
% trSq_holds); and twisted_weil_representation_convergence (weilConvergence_of_siegelSet); (5)
% rowFun_at_gridA_alt, written out from Lemmas asm-dependence-functional-* and covolume-points
% (evW_pRows_of_rows_alt); Lemma asm-dependence-functional-grid written out (rowE_hasGrid_fld_alt with
% twists_at_linked_alt, rowE_pinned_grid_alt, rowE_full_grid_alt, hasGrid_of_full_field_lamTie) down to
% Lemma asm-dependence-functional-witness, written out (eventually_sizeRows_alt with sizeRows_of,
% weight_one_places_finite_of_witness with realPins3_realTOf_alt and zscale_f0_of_witness_alt,
% rowTest_setting_alt with core_of_weak) and szRow_numAt (Lemma
% asm-dependence-functional-witness-numbers-values, written from its four SzRow pieces, which are
% written out: explicitConstants_parameterConditions_large, szRow_blockLengthParams and szRow_sizesAt
% down to their SzRow computations, stated in Lemmas -parameters-computations, -blocks-computations
% and -bounds-computations, and szRow_band_lambda_primes); (6)
% Open0ProofU.open_orbit_of_unramified_exponents_alt, written out from Lemmas asm-dependence-orbit-*;
% (7) linked_bundles_l3_of_open_alt, written out (class_sum_divisor_all_alt, evW_pRows_of_rows_alt,
% eventual_L3_of_covolume_alt); (8) StageOneU.member_bundles_large_unramified_exponents_alt, written
% out from Lemmas asm-dependence-bundles-*; (9) final_alt_alt, its five steps outlined; the open
% parts of (6) to (9) are listed after their open step; (10) the theorems
% named in the leanref lines; in (a) the class count is ColClassNumber.classNumber_closed
% (classNumber_all, Proposition cv-class-number-formula); the second complement inputs are
% admissibleConstants_closed_under_sizeBounds with familySecondComplementInputs_of_tuSize_of_dGamSize,
% written out down to its eight conjuncts, of which (a) T_ne_g and (b) normDecompositionCond_famP are
% Lemma asm-dependence-complement-norms; open: the six bounds (c)
% AdmissibleDegree.admissibleDegreeCount_along_family, (d) TypeSize.typePrimeSizeBound_of_eps_lt_one_div_fiftyTwo,
% (e) LevelBudget.admissible_off_type_level_budget, (f)
% exists_nonneg_const_discriminant_term_le_budget_family, (g)
% RamifiedOffD.offDiagonalRamified_bound_along_family, (h)
% InertLossAdm.inertPrimes_loss_log_sum_admissible; (b) six_rows_at_member_alt over
% goodTwistCores_over_trueColumns, with period integrality and twisted period integrality written out
% (Theorem pd-period-integrality-along-families); open: referenceChainEvaluationAtBundlePoints_of_overCols,
% twistedChainEvaluation_rational_integral_bundlePoints, realSideColumnsCompatible_bundlePoints_overColumns,
% levelPrimeColumnValuations_bundlePoints_overColumns and row_aAtG_at_built_alt
% (holds_at_bundlePoints_of_trueColumns, rowFormsMatchPeriodMap_of_furtherProperties); (c) open:
% FireAtRealisation.fire_alt and RowsAtRealTBClose.rowsB_of_decay_alt; (e) open:
% LocalRows.at_choice_alt; (f) open: Svd.servH_row_of_artin (cmGoodModel_of_globalArtinReciprocity,
% bundle_bothParts_of_standardKernel_and_six_hypotheses) with builtC_hasPart and
% LinkStd.linkR_stdTie_entry, then R1.row_hServ_of_kerServ_alt; (h) is written out (Lemmas
% asm-dependence-projections-onto and asm-dependence-base-vectors); (i) open:
% CountConstants.count_constants_of_siegelWeil and irreducibilityRow_of_concreteAutomorphicData;
% (d) is written out from row_hInstR_of_hU,
% Glob.volSigIndep_and_classItemsAt_of_massRatioUnderVolume, Glob.volMeas_fundDomain_eq_of_tiedIndep
% (retie D, retie_volTieG_and_cdScaleTie), Glob.vol_eq_massI_all_sig with volumeMass_of_unfolding,
% classItemsAt_all (genusMassStd_all, ClassNumberMass.classFin_and_classMass,
% lattice_stabilizer_doubleCosets_classRest_mass, WeightedDoubleCoset.sum_inv_card_stab_eq);
% (g) is written out from row_hTw_of_tie_alt,
% Top4.tie_tw_rec_of_real_rec, Top4.real_t_tw_pl_eq, IndefiniteRealisationRecursion.tw_iff (only its
% forward half), tupleRead_baseTuple_eq_family, recR.tupAt, charW.tup.mu_ne and FamP.tw_supp;
% (11) memberTop_final_contextParameters with recS (hostWW, designS,
% kitS); (12) row_hTu_cTmin_alt, through HarishChandraReduction.tuSize_min and
% GappedMultiplicityBounds.explicit_bound_family (theorem_three_family_of_gapped_multiplicity with
% BesselDiscreteSeries.lowest_K_type_mult_one_of_gaps), with explicitConstants_family_of_volume_eq_mass,
% rowsAF_fam, tuSize_fam and part_five_all_of_explicit; (13) done: litThin, genusMassAt_two (Lemma
% asm-dependence-genus-mass), rowsFinDim_link (Lemma asm-dependence-findim, cited to Harish-Chandra),
% VolRatioFromSiegelWeil.volumeRatio, ontoX_and_lamLow_of_unitaryGroupHypotheses'; open:
% wuDecayOfGap (Pf.wuDecayOfGap_of_uniqGap with Bmu.besselMapUniqGap), rowsPos_typeP
% (typeBound_of_largeGap_discreteSeries_bessel with bessel_multiplicity_one_large_gap_and_whittaker_bessel),
% rowNGP_of_siegelWeil and rowMainEqP_of_siegelWeil (stepChoiceDataC_of_siegelWeil_inInfinity with
% realObj_inInfinityX_jsz, analytic_pinnedTheoremB_uniform and bDat_theoremB_global).
\end{proof}

\begin{lemma}\label{asm-dependence-all-rows}
Assume global Artin reciprocity, the Siegel-Weil formula (Theorem~\ref{lit-siegel-weil}), the
covering and the classification of discrete series, the finiteness of Harish-Chandra modules, rapid
decay of cusp forms and the Casselman-Wallach theorem (Theorem~\ref{lit-bundle-unitary}), and the
conclusions of parts (1) and (2) of Theorem~\ref{is-surjectivity-and-lower-norm}. Let local
Langlands maps be given, let $N$ be an extension and $p=p(N)$
(Definition~\ref{asm-dependence-parameters}), and let $M$ be a property of classes of families. Suppose that the open parameter statements hold at $p$
(Definition~\ref{asm-open-parameter-statements}), that hypotheses (5) and (6) of
Theorem~\ref{is-definite-rows} are linked at $p$, and that the property ``$M$ holds for
$\mathcal A^{\mathrm{des}}=\operatorname{Des}(\mathcal A_N,3)$'' is linked at $p$. Then each of the
following properties of a class $\mathcal A$ and a selection is linked at $p$, with the items of
Definition~\ref{asm-definite-side-hypotheses}:
\begin{enumerate}
\item the selected rows, items (3)(a) to (3)(f), at every member;
\item (Joins): the join of the indefinite realization at every member, and the fullness of the
  image of the base vectors eventually along the working subclass $\mathcal W(\mathcal A)$
  (Theorem~\ref{pd-period-integrality-along-families}), which is $\mathcal A^{\mathrm{des}}$; and
  (Projections);
\item the volume items (7) at every member, and the class number formula of (Sizes) at every
  member;
\item the integrality inputs (10);
\item the frame condition, the loop rows, the selected form of the orthogonality theorem, and
  counting constants with their point counts and a number $c_I>0$ with the irreducibility inputs
  (items (6) and (11));
\item (Definite rows): the conclusion of Theorem~\ref{is-definite-rows} for every class and
  selection, in the form of item (9);
\item $M$ holds for $\mathcal A^{\mathrm{des}}$.
\end{enumerate}
\end{lemma}
\leanref{Development/AutomorphicArithmetic/ThetaCorrespondence/ThetaPairLifts/ThirdPlaceSeries/LinkROfIsLinked.lean}{GlobalTheory.Fc.AllAtMemberStage0.allRows_of_open0_alt}
\leanref{Development/AutomorphicArithmetic/ThetaCorrespondence/ThetaPairLifts/ThirdPlaceSeries/AllRows.lean}{GlobalTheory.Fc.AllRowsAlt}
\begin{proof}
Each property is linked at $p$ by one of the following, with
Lemma~\ref{asm-dependence-linked}(1) to pass from a property along every slot, or from a linked
property that implies it, to a linked one.
\begin{enumerate}
\item Item (3)(a) is Theorem~\ref{asm-dependence-inputs}(10)(c). Item (3)(b) is
  Theorem~\ref{asm-dependence-inputs}(10)(g), applied with open parameter statement (d). Items
  (3)(c) to (3)(f) are open parameter statements (a) and (b).
\item The join is open parameter statement (c). The fullness of the image of the base vectors and
  the surjectivity of the projections are Theorem~\ref{asm-dependence-inputs}(10)(h), applied with
  open parameter statement (a).
\item The volume items are Theorem~\ref{asm-dependence-inputs}(10)(d), with the mass ratio under
  volume, which holds by the Siegel-Weil formula and Theorem~\ref{asm-dependence-inputs}(13); the
  class number formula is (10)(a).
\item Point integrality, twisted point integrality, period integrality, twisted period integrality,
  the second-layer conditions at the real place, the bounds at the level primes and the sharp
  column bounds at the auxiliary prime are Theorem~\ref{asm-dependence-inputs}(10)(b)
  (Remark~\ref{pd-integrality-inputs}); the second complement inputs are (10)(a); and the selections
  are served by (10)(f), from global Artin reciprocity.
\item The frame condition, the loop rows and the selected orthogonality theorem are
  Theorem~\ref{asm-dependence-inputs}(10)(e). The counting constants and $c_I$ are (10)(i), applied
  with open parameter statement (e).
\item By Lemma~\ref{asm-dependence-linked}(2), the conjunction of open parameter statements (a),
  (b), (f), (g), (h) and (i) and of hypotheses (5) and (6) of Theorem~\ref{is-definite-rows} is
  linked at $p$, with a slot $R$. For every slot $S$, if these properties hold along
  $(S,\Lambda_p)$, then the hypotheses of Theorem~\ref{is-definite-rows} hold for $(S,\Lambda_p)$:
  (1) is statement (a), (2) is statement (b), (3) is statement (f), (4) is statement (g), (5) and
  (6) are the two hypotheses, (7) and (8) are statement (h), and (9) is statement (i). Its
  published inputs and the conclusions of parts (1) and (2) of
  Theorem~\ref{is-surjectivity-and-lower-norm} are among the assumptions. So
  its conclusion holds along $(S,\Lambda_p)$, and by Lemma~\ref{asm-dependence-linked}(1) it is
  linked at $p$.
\item This is a hypothesis.
\end{enumerate}
\end{proof}

\begin{lemma}\label{asm-class-selection}
Assume Siegel's mass ratio for CM fields (Theorem~\ref{cv-siegel-mass-ratio}) and that the mass
ratio under volume holds for every definite-side datum (Definition~\ref{asm-facts-volume-tie}(2)).
Let local Langlands maps be given, and let $p$ be parameters for which bundles exist at large bases
(Definition~\ref{asm-dependence-parameters}) and which have the member property if every finite
subgroup of $GL_n(\mathbf Q)$ has order dividing $M(n)$ for every $n$. Let $H$ be a property of
classes and selections that is linked at $p$, and let $c_T$ be a function of two real variables.
Then there are a class $\mathcal A$ of families and a raw selection $s$ for $\mathcal A$ such that:
\begin{enumerate}
\item if every finite subgroup of $GL_n(\mathbf Q)$ has order dividing $M(n)$ for every $n$, then
  for every $\lambda>1$ and every host requirement $G$ such that hosts satisfying $G$ exist on
  windows of exponent $\lambda$, design choices exist in the class of families of
  $\mathcal A\cap\mathcal C_G(\lambda)$ that satisfy the twist condition, with constant $3$ and
  function $c_T$;
\item every family in $\mathcal A$ carries its duty data (Definition~\ref{asm-duty}), and for all
  large $x$ every twisted lattice $\tilde M_\mu$ of every member is stable under $\mathcal O_L$;
\item the bundles of $s$ are built from $(Q_p,\Lambda_p)$;
\item $H(\mathcal A,s)$ holds.
\end{enumerate}
\end{lemma}
\leanref{Development/AutomorphicArithmetic/CharacterBoxes/ClassSel.lean}{UpperTheory.Top.classSel}
\leanref{Development/AbelianVarietyTheory/Definitions/GlobalTwisting/ValidForConstructedPointClassesEventuallyWorking.lean}{GlobalTheory.Fc.cls}
\leanref{Development/AbelianVarietyTheory/Definitions/GlobalTwisting/ValidForConstructedPointClassesEventuallyWorking.lean}{GlobalTheory.Fc.pick}
\leanref{Development/AbelianVarietyTheory/Definitions/GlobalTwisting/ValidForConstructedPointClassesEventuallyWorking.lean}{GlobalTheory.Fc.sevenDuties_of_class}
\begin{proof}
Let $R_0$ be the slot for which bundles exist at large bases, and $R_1$ the slot with which $H$ is
linked at $p$. Put $R=R_0\cap R_1$ and $Q=Q_p\cap R$. Every core built from $p$ satisfies $R$ at
large base data, eventually in the context. By the remark at the start of the proof of
Lemma~\ref{asm-dependence-linked}, $H$ holds along $(Q,\Lambda_p)$, with a condition $E$ on
bundles, and eventually every base point of $Q$ has, for every $c$, a bundle linked to it. Every
linked bundle eventually satisfies $E$, so eventually every base point of $Q$ has, for every $c$, a
bundle linked to it that satisfies $E$.

Let $\mathcal A$ be the class of families $\mathcal F$ that are linked to the constants $\kappa$ of
a context, satisfy the class conditions of the construction, and whose members at every scale
$x>x_1$ have, for every $c$, a bundle built from $(Q,\Lambda_p)$ in a context with constants
$\kappa$ that satisfies $E$. At every member of such a family let $s$ choose such a bundle for the
size constant of the family. Then $s$ is a selection whose bundles are built from $(Q,\Lambda_p)$
and satisfy $E$, so $H(\mathcal A,s)$ holds, which is (4). Since $Q\subset Q_p$, (3) holds. The class
conditions of the construction include the duty data and the stability of the twisted lattices,
which is (2).

For (1), assume the divisibility, so that $p$ has the member property. It passes from
$(Q^{\mathrm b}_p,Q_p)$ to $(Q^{\mathrm b}_p,Q)$: the base
condition, the first pass and the export do not involve $Q_p$, and in the second pass the core in
$Q^{\mathrm b}_p\cap Q_p$ with all its twists is built from $p$ at large base data, so it satisfies
$R$. Take for $O_\kappa$ the property that a base point has, for every $c$, a bundle built from
$(Q,\Lambda_p)$ in a context with constants $\kappa$ that satisfies $E$; by the first paragraph
every base point of $Q$ satisfies it eventually in the context. The class $\mathcal A_O$ of
Theorem~\ref{asm-dependence-inputs}(11) is then $\mathcal A$, and that part gives (1).
\end{proof}

\begin{theorem}\label{asm-definite-side-exists}
Assume the published results on unitary groups
(Theorem~\ref{lit-bundle-unitary}) and on the local theta correspondence
(Theorem~\ref{lit-bundle-theta}), the agreement of global and local theta
integrals and the local matching of models of the Weil representation
(Definition~\ref{th-theta-hypotheses}), the covolume bound of
Theorem~\ref{cv-covolume-bound}, and the existence, for every
number field and every finite place, of a map satisfying the local Langlands
laws (Definition~\ref{gal-llc-laws}). Suppose that for every real $c$ the
property ``eventually along $\operatorname{Des}(\mathcal A_N,3)$, the local
bound at the twisted places above $2$ and $3$
(Definition~\ref{asm-local-bound-two-three}(4)) holds'' is linked at the
parameters of dependence $c$ (Definition~\ref{asm-dependence-parameters}). Then
a collection of definite-side hypotheses exists
(Definition~\ref{asm-definite-side-hypotheses}).
\end{theorem}
\begin{proof}
Write $P$ for the published results on unitary groups and $F$ for those on the
local theta correspondence, and fix the local Langlands maps. For a real number
$c$ we use the list extension $N_{\mathrm{list}}(c)$, the origin extension
$N_{\mathrm{orig}}(c)$, the parameters of dependence $c$ and the parameters
$p_c$ of the extension of $N_{\mathrm{orig}}(c)$ by the grid property for $c$
(Definition~\ref{asm-dependence-parameters}). The published results that the
parts of Theorem~\ref{asm-dependence-inputs} and
Lemma~\ref{asm-dependence-all-rows} assume are in $P$ and $F$, and the theta
hypotheses are assumed. By Theorem~\ref{is-surjectivity-and-lower-norm}, the
conclusions of its parts (1) and (2) hold, and by
Theorem~\ref{asm-dependence-inputs}(13), with the Siegel-Weil formula, the mass
ratio under volume and the mass formula for genera of rank $2$ hold for every
definite-side datum.

\emph{Step 1: separation of functionals.} Functionals separate, by
Theorem~\ref{asm-dependence-inputs}(1).

\emph{Step 2: rationality of rows.} The rows of the construction are rational,
by Theorem~\ref{asm-dependence-inputs}(2) and Step 1.

\emph{Step 3: the Fourier fit.} Theorem~\ref{asm-dependence-inputs}(3), with
Step 1, gives a constant $C'$ such that for every $C_F\ge C'$ the rows of every point with an
evaluation grid of constant $C_F$, fitted to the point, satisfy the functional
bound, eventually along the standard families with the division bound at
twisted places.

\emph{Step 4: the evaluation grid.} Theorem~\ref{asm-dependence-inputs}(4),
with Step 1, gives a constant $C_{F_1}$ such that for every $C_F\ge C_{F_1}$ the grid with
constant $C_F$ evaluates.

\emph{Step 5: the functional condition.} Theorem~\ref{asm-dependence-inputs}(5),
with Step 3, gives a constant $C_{F_0}$. Put $c=\max(C_{F_0},C_{F_1})$.

\emph{Step 6: open parameter statements.} By
Theorem~\ref{asm-dependence-inputs}(6), the open parameter statements
(Definition~\ref{asm-open-parameter-statements}) hold at the parameters of
dependence $c$.

\emph{Step 7: passing to the extension.} The origin extension
$N_{\mathrm{orig}}(c)$ requires $N_{\mathrm{list}}(c)$, so its extension by the
grid property for $c$ is stronger than $N_{\mathrm{list}}(c)$. By
Lemma~\ref{asm-dependence-linked}(3), every property of classes and selections
that is linked at the parameters of dependence $c$ is linked at $p_c$, and by
Step 6 the open parameter statements hold at $p_c$.

\emph{Step 8: functional and covolume conditions.} Apply
Theorem~\ref{asm-dependence-inputs}(5) with $C_F=c\ge C_{F_0}$ and
$N=N_{\mathrm{orig}}(c)$, which is stronger than $N_{\mathrm{list}}(c)$; the
extension of $N$ by the grid property for $c$ gives the parameters $p_c$, at
which the open parameter statements hold by Step 7. So hypothesis (6) of
Theorem~\ref{is-definite-rows} is linked at $p_c$. The extension that defines
$p_c$ requires $N_{\mathrm{list}}(c)$, hence the law extension, so
Theorem~\ref{asm-dependence-inputs}(7), with the covolume bound of
Theorem~\ref{cv-covolume-bound} and Step 7, shows that hypothesis (5) of
Theorem~\ref{is-definite-rows} is linked at $p_c$.

\emph{Step 9: bundles.} Since $c\ge C_{F_1}$, the grid property for $c$ holds
eventually in the context and at large base data, by Step 4. So
Theorem~\ref{asm-dependence-inputs}(8), with Step 2, applied with $A$ the grid property for
$c$, gives Siegel's mass ratio for CM fields
(Theorem~\ref{cv-siegel-mass-ratio}), the member property of $p_c$, and the
existence of bundles at large bases for $p_c$.

\emph{Step 10: the Tate-Shafarevich link.} By hypothesis, the property
``eventually along $\operatorname{Des}(\mathcal A_N,3)$, the local bound at the
twisted places above $2$ and $3$ holds'' is linked at the parameters of
dependence $c$. By Theorem~\ref{asm-dependence-inputs}(9), so is the property
``the Tate-Shafarevich link (Definition~\ref{asm-twisted-links}) holds
eventually along $\operatorname{Des}(\mathcal A_N,3)$''. By Step 7 it is linked
at $p_c$.

\emph{Step 11: all rows at $p_c$.} Apply Lemma~\ref{asm-dependence-all-rows} to
$p_c$, with $M$ the property ``the Tate-Shafarevich link holds eventually along
the class''. Its hypotheses are Steps 7, 8 and 10. So the properties (1) to (7)
of that lemma are linked at $p_c$. By Theorem~\ref{asm-dependence-inputs}(12)
and Lemma~\ref{asm-dependence-linked}(1), the size condition on $T_u$ of
(Sizes), with $c_T$ the function given by the Harish-Chandra reduction, is
linked at $p_c$ as well. Let $H$ be the conjunction of all these properties;
by Lemma~\ref{asm-dependence-linked}(2), $H$ is linked at $p_c$.

\emph{Step 12: the collection.} By Step 9, $p_c$ has the member property and
bundles exist at large bases for $p_c$. Lemma~\ref{asm-class-selection},
applied to $p_c$, $H$ and $c_T$, gives a class $\mathcal A$ and a raw selection
$s$ for $\mathcal A$; let the raw selection of the collection be $s$, with
$c_T$ and with the number $c_I>0$ and the counting constants given by $H$. The
items of Definition~\ref{asm-definite-side-hypotheses} hold as follows. (Design)
is Lemma~\ref{asm-class-selection}(1). (Duty) and the stability of the twisted
lattices $\tilde M_\mu$ in (Joins) are Lemma~\ref{asm-class-selection}(2). By
Lemma~\ref{asm-class-selection}(4), $H(\mathcal A,s)$ holds; it gives (Selected
rows), the rest of (Joins), (Volumes), (Definite rows), (Integrality inputs),
(Irreducibility inputs), (Projections), the second-layer conditions at the
real place in (Real place and masses), the frame condition, the loop rows, the
selected orthogonality theorem and the counting constants of (Orthogonality
inputs), (Sizes) and (Tate-Shafarevich link); here the designed families
$\operatorname{Des}(\mathcal A_N,3)$ lie in the working subclass
$\mathcal W(\mathcal A)$ (Theorem~\ref{pd-period-integrality-along-families}),
along which the properties of Lemma~\ref{asm-dependence-all-rows}(4) hold. The
remaining items are
Theorem~\ref{asm-dependence-inputs}(13): the analytic input
(Analytic input); the decay of Wu type, the
finite dimensionality of rows, the positivity of row types, the row
non-degeneracy and the row main-term identity (Orthogonality inputs); and the
mass ratio under volume and the mass formula for genera of rank $2$ (Real place
and masses). These give a collection of definite-side hypotheses.
\end{proof}
\leanref{Development/AutomorphicArithmetic/ThetaCorrespondence/ThetaPairLifts/ThirdPlaceSeries/SiegelEisensteinSecondTermHypotheses.lean}{UpperTheory.Lit.nonempty_hyps_definite_side_of_closed_bound}
\leanref{Development/AutomorphicArithmetic/ThetaCorrespondence/ThetaPairLifts/ThirdPlaceSeries/HypothesesNonemptyOfUnramifiedExponents.lean}{GlobalTheory.Fc.S0.StageU.hypotheses_nonempty_of_unramified_exponents_alt}

\section{The half-plane}\label{section-asm-half-plane}

\subsection{Zero-free half-planes and local data at 2 and 3}\label{subsection-asm-half-plane-local-data}

\begin{theorem}\label{asm-zero-free-primitive}
Let $\lambda>1$. Suppose that tied families exist for all host requirements
(Definition~\ref{asm-tied-everywhere}), and that for every primitive Dirichlet
character $\psi$ hosts satisfying the host requirement $G_\psi$ of $\psi$
(Definition~\ref{asm-host-requirement}) exist on windows of exponent
$\lambda$. Then $L(s,\psi)\ne0$ for every primitive Dirichlet character
$\psi$ and every $s$ with $\operatorname{Re}s>1/2$. For the trivial character
modulo $1$ and $s=1$ this is the assertion $\zeta(1)\ne0$ about the value of
Remark~\ref{mt-value-at-one}.
\end{theorem}
\begin{proof}
Let $\psi$ be primitive and let $K_\psi$ be its cyclic level
(Definition~\ref{wl-cyclic-level}). Apply Corollary~\ref{wl-zero-free-cyclic}
with $Q=\{\psi\}$, in Situation~\ref{wl-families} with the families of
exponent $\varepsilon$ (Definition~\ref{pd-family}) as the class
$\mathcal F_\varepsilon$ and the tie of Definition~\ref{asm-tied} as the
condition $T_F$. Its condition (1), the Chebyshev bound along tied families,
holds for every level by Theorem~\ref{ib-window-bound-on-scales}, and its
condition (3) is Theorem~\ref{wl-zero-free-level}.

For condition (2) we need $e_D,e_L$ and $\varepsilon_1>0$ such that for every
$\varepsilon\in(0,\varepsilon_1)$ some family with exponent $\varepsilon$
satisfies the host conditions for $K_\psi$, is tied on its scales with
exponents $(e_D,e_L)$, and has windows of exponent $\lambda$. Every family in
$\mathcal C_{G_\psi}(\lambda)$ satisfies the Tate-Shafarevich bound
(Definition~\ref{asm-sha-bound}), since $G_\psi$ includes it at every member
at every scale $x>x_1$. Hosts satisfying $G_\psi$ exist on windows of
exponent $\lambda$ by hypothesis. Since tied families exist for all host
requirements, $\mathcal C_{G_\psi}(\lambda)$ has tied members: there are $e_D,e_L,\varepsilon_1$, and for every
$\varepsilon\in(0,\varepsilon_1)$ a family $\mathcal F\in\mathcal
C_{G_\psi}(\lambda)$ with exponent $\varepsilon$ that is tied at all its large
scales. This family has windows of exponent $\lambda$. At a scale $x>x_1$
every member has scale $(x,\varepsilon,J)$, so $G_\psi$ gives the host
conditions for $K_\psi$ at $(x,\varepsilon,4)$ at every member at every scale
$x\ge x_1+1$.
\end{proof}
\leanref{Development/AutomorphicArithmetic/CharacterBoxes/ZeroFreeAllOfTopG.lean}{UpperTheory.Top.zeroFreeAll_of_topG}

The case $\psi=\mathbf 1_1$, $s=1$ of Theorem~\ref{asm-zero-free-primitive}
is the assertion $\zeta(1)\ne0$. It holds only because Mathlib assigns to
$\zeta$ at its pole the value $\zeta(1)=(\gamma-\log4\pi)/2\ne0$
(Remark~\ref{mt-value-at-one}); for the meromorphic function $\zeta$ the point
$s=1$ is a pole and the assertion has no content. In the proof it enters
through Corollary~\ref{wl-zero-free-cyclic} for $K_{\mathbf 1}=\{\mathbf 1\}$:
the proof of Proposition~\ref{wl-pole-data-family} uses $\zeta(1)\ne0$ to see
that $1$ is not in the zero set. The later theorems of this chapter exclude
the trivial character at $s=1$, and Proposition~\ref{mt-imprimitive}, which
they use, assumes this case as part of its hypothesis.

\begin{theorem}\label{asm-every-character-green}
Assume Green's parametrization of the cuspidal characters of
$GL_d(\mathbf F_g)$ in the labelled form of
Corollary~\ref{lit-green-labelled}, and the weakening of Vandiver's conjecture
(Definition~\ref{sha-vandiver-weakening}). Suppose that a collection of
definite-side hypotheses exists (Definition~\ref{asm-definite-side-hypotheses}).
Then for every $N\ge1$, every Dirichlet character $\chi$ modulo $N$ and every
$s\in\mathbf C$ with $\operatorname{Re}s>1/2$, $L(s,\chi)\ne0$ unless $\chi$
is the trivial character modulo $N$ and $s=1$.
\end{theorem}
\begin{proof}
By Theorem~\ref{asm-tied-from-green}, tied families exist for all host
requirements. For every primitive $\psi$, Theorem~\ref{hf-hosts-on-windows}
and the weakening of Vandiver's conjecture for $\psi$ show that hosts
satisfying $G_\psi$ exist on windows of exponent $3/2$.
Theorem~\ref{asm-zero-free-primitive} with $\lambda=3/2$ gives
$L(s,\psi)\ne0$ for every primitive $\psi$ and $\operatorname{Re}s>1/2$.
Proposition~\ref{mt-imprimitive} passes to all characters.
\end{proof}
\leanref{Development/AutomorphicArithmetic/ThetaCorrespondence/ThetaPairLifts/ThirdPlaceSeries/TopPlug.lean}{UpperTheory.Top.every_character_of_green_v_win}
\leanref{Development/AutomorphicArithmetic/ThetaCorrespondence/ThetaPairLifts/ThirdPlaceSeries/TopPlug.lean}{UpperTheory.Top.zero_free_all_of_green_v_win}

\begin{definition}\label{asm-two-three-local-data}
Let $D$ be a definite-side datum (Definition~\ref{th-definite-datum}) with
ranks $k,m,t$ and $r=k+m-1$, host field $E$ with totally real subfield $F$,
integer $Q_S$, level primes $g_0,g_1$ and $p_{de},p'_{de}$,
and coefficient field with tuples $\mu_w$ at its places $w$ and base tuple
$\mu_0$. Let $\Omega$ be an algebraic closure of $E$, put
$G_E=\operatorname{Gal}(\Omega/E)$, and for a prime $\ell$ let
$\varepsilon_\ell$ be the $\ell$-adic cyclotomic character of $G_E$. The
\emph{local data at the twisted places} of $D$, part of the local data of
Definition~\ref{th-definite-datum}(9), consist of a finite set of \emph{local
indices} $i$, each with:
\begin{enumerate}
\item a prime $\ell_i$, a lift $1\le j_i\le m$ and a place $w_i$ of the
  coefficient field over $\ell_i$ with $\mu_{w_i}\ne\mu_0$;
\item a complete discrete valuation ring $\mathcal O_i$, finite over
  $\mathbf Z_{\ell_i}$, with fraction field $K_i$;
\item a finite nonempty set $X_i$ of \emph{Hecke points}, a local
  $\mathcal O_i$-subalgebra $\mathbf T_i$ of the $\mathcal O_i$-valued
  functions on $X_i$ that separates the points of $X_i$, and two subsets
  $Z_i$ and $Z^{\mathrm{el}}_i$ of $X_i$; the points of $Z^{\mathrm{el}}_i$
  are the \emph{eligible} points;
\item disjoint finite sets $S_i$ and $Q_i$ of places of $E$ that are split
  over $F$ and do not lie over $\ell_i$, where $S_i$ is the set of places over
  $g_0$, $g_1$ and $p_{de}$ (over $p'_{de}$ if
  $\ell_i=p_{de}$) that do not lie over $\ell_i$;
\item a continuous character $\chi_i\colon G_E\to\mathcal O_i^\times$ that is
  ramified at every place of $Q_i$, and for every place $v$ of $E$ an element
  $\gamma_{i,v}$ of an inertia group at $v$ which, when $v\in S_i\cup Q_i$,
  topologically generates this inertia group together with its wild inertia;
  put $\zeta_{i,v}=\chi_i(\gamma_{i,v})$;
\item an integer $s_0(i)$, equal to $v_2(Q_S)$ if $\ell_i=2$ and to
  $\max(1,v_{\ell_i}(Q_S))$ if $\ell_i$ is odd.
\end{enumerate}
For a local index $i$ and an $\mathcal O_i$-module $M$ put
$\operatorname{len}_iM=\operatorname{len}_{\mathcal O_i}M/
\operatorname{len}_{\mathcal O_i}(\mathcal O_i/\ell_i\mathcal O_i)$ if $M$ has
finite length, and $\operatorname{len}_iM=0$ otherwise.
\begin{enumerate}\setcounter{enumi}{6}
\item The \emph{congruence length} of $\pi\in X_i$ is
  $c_i(\pi)=\operatorname{len}_i(\mathcal O_i/I_i(\pi))$, where $I_i(\pi)$ is
  the ideal of the values at $\pi$ of the functions in $\mathbf T_i$ that
  vanish on $X_i\setminus Z_i$.
\item Let $\rho\colon G_E\to GL_r(\mathcal O_i)$ be a representation. Its
  \emph{Selmer group} $\operatorname{Sel}_i(\rho)$ is the
  group of classes in $H^1(G_E,(K_i/\mathcal O_i)^r)$, for the action through
  $\rho\otimes\chi_i^{-1}$, whose restriction to a decomposition group at a
  place $v$ of $E$ lies in the Bloch-Kato subgroup defined by the comparison
  data of $D$ at $v$ if $v$ lies over $\ell_i$, is arbitrary if
  $v\in S_i\cup Q_i$, and is unramified otherwise. Its \emph{Selmer length}
  is $\operatorname{sel}_i(\rho)=\operatorname{len}_i\operatorname{Sel}_i(\rho)$.
\item For a prime $\ell$ let $\nu(\ell)$ be the least $\nu\ge0$ such that
  $\varepsilon_\ell(g)^{tr}\not\equiv1\pmod{\ell^{\nu+1}}$ for some $g\in G_E$,
  as in Definition~\ref{th-twisting-data}(11). A family
  $\rho=(\rho_{i,\pi})$ of representations
  $\rho_{i,\pi}\colon G_E\to GL_r(\mathcal O_i)$, one for each local index $i$
  and each $\pi\in X_i$, satisfies the \emph{killing condition} if for every
  local index $i$: $\operatorname{Sel}_i(\rho_{i,\pi})$ has finite length and
  $\operatorname{sel}_i(\rho_{i,\pi})\le\nu(\ell_i)+\ell_i^{-2}$ for every
  eligible $\pi$; and
  $t\sum_{v\in Q_i}\operatorname{len}_i\big(\mathcal O_i/(\zeta_{i,v}-1)\big)
  \le\ell_i^{-2}$.
\end{enumerate}
\end{definition}
\leanref{Formalism/SharpRestriction/LevelLetterCocycles/DataN.lean}{Formalism.BasicNotions.NData}
\leanref{Formalism/SharpRestriction/LevelLetterCocycles/DataN.lean}{Formalism.BasicNotions.NData.vcc}
\leanref{Formalism/SharpRestriction/LevelLetterCocycles/QuotientHomNonvanishingBound.lean}{Formalism.MainNotions.N.Br.selPi}
\leanref{Formalism/SharpRestriction/LevelLetterCocycles/QuotientHomNonvanishingBound.lean}{Formalism.MainNotions.N.Br.sPi}
\leanref{Development/AutomorphicArithmetic/ResidueCharacters/TwistedPlaces.lean}{Formalism.MainNotions.TwistedPlaces.nuL}
\leanref{Formalism/SharpRestriction/LevelLetterCocycles/JointConditionsZpowers/CommutingEndomorphismCondition.lean}{Formalism.MainNotions.N.KillingConditions.KillAt}

\begin{definition}\label{asm-local-bound-two-three}
Let $D$ be a definite-side datum with local data at the twisted places
(Definition~\ref{asm-two-three-local-data}), ranks $m,t,r,n$ and $f=2n$, host
field $E$ with class number $h_E$, integer $Q_S$ and splitting bound $B$
(Definition~\ref{hf-host-field-rows}). For every prime $\ell$ and $0<k<t$ the
datum records a length for $\text{Ш}^1_\ell(E,k)$, equal to its length when
this is finite; let $V(\ell)$ be the sum of these $t-1$ recorded lengths. Let
size data for $D$ be given (Definition~\ref{pd-size-data}), with scale
$(x,\varepsilon,J)$, $b=\exp(x^{\varepsilon})$, $b'$, $w''$, exponents $b_t$
at the top classes, and numbers $\eta$ and $S_c$.
\begin{enumerate}
\item Let $\rho=(\rho_{i,\pi})$ be a family of representations
  $\rho_{i,\pi}\colon G_E\to GL_r(\mathcal O_i)$, one for each local index $i$
  and each $\pi\in X_i$, and let $C$ be real. The \emph{bound at $2$ and $3$
  with constant $C$} holds for $\rho$ if for every local index $i$ with
  $\ell_i=2$, or with $\ell_i=3$ and $3\nmid Q_S$, and every eligible
  $\pi\in Z^{\mathrm{el}}_i$,
  \[
   c_i(\pi)\le\operatorname{sel}_i(\rho_{i,\pi})+C\,f\,m\,b\log x
   +C\,x^{C}x^{2\varepsilon}+V(\ell_i)+v_{\ell_i}(h_E)+\ell_i^{-2}.
  \]
\item \emph{Hub data} for $D$ consist of two families of representations
  $r_{i,\pi}\colon G_E\to GL_{r+t}(\mathcal O_i)$ and
  $\rho_{i,\pi}\colon G_E\to GL_r(\mathcal O_i)$, for the local indices $i$
  and the points $\pi\in X_i$, together with further local data: lengths at
  the places of $E$, error bounds, a partition of the places of $E$ over each
  $\ell_i$ into row places and column places, constants, band conditions, a
  finite set of exceptional primes and an integer.
\item Fix twist choices of $D$ at the places of its coefficient field
  (Definition~\ref{asm-twist-condition}). A \emph{setup} for $D$, its size
  data and these choices consists of local Hodge-Tate data, Fontaine data and
  comparison data of $D$ at every prime, hub data and a real number $C_2$,
  such that:
  \begin{enumerate}
  \item the datum and its twist choices satisfy the \emph{norm-choice
    conditions}, a list of conditions on the datum, its local data and its
    twist choices; among them are residual irreducibility in rank $r$, de Rham
    conditions on the families of representations attached to the Hecke
    points, the uniqueness of the local index with a given lift and place, and
    ties between the twist choices and the local data;
  \item the hub data are \emph{canonical} for the size data, the Hodge-Tate
    data and the row constants $(C_2,320,25,23,32,32)$: each $r_{i,\pi}$ is
    continuous, semisimple over $K_i$ and attached to $\pi$, in that its
    characteristic polynomials of Frobenius at the split places of $E$ outside
    a finite set given by the datum are the Hecke polynomials of $\pi$; each
    $\rho_{i,\pi}$ is continuous; for $\pi\in Z_i$, $r_{i,\pi}$ is the
    semisimple block sum of $\rho_{i,\pi}$ and the characters
    $\chi_i\varepsilon_{\ell_i}^{1-c}$, $1\le c\le t$; and the numerical data
    of the hub are those that the size data, the Hodge-Tate data and the row
    constants determine;
  \item the size data are admissible: $\eta+t+S_c<w''$, $3w''\le b$,
    $b_t+b'/x\le b$ for every exponent $b_t$, $t\le\eta$, and
    $S_c+\eta+t+2\le b$;
  \item $x^{\varepsilon}\ge8$ and $B\ge60$;
  \item the representations $\rho=(\rho_{i,\pi})$ of the hub data satisfy
    the killing condition (Definition~\ref{asm-two-three-local-data});
  \item $\text{Ш}^1_p(E,k)$ has finite length for every prime $p$ and every
    $0<k<t$.
  \end{enumerate}
\item A family $\mathcal F$ with exponent $\varepsilon$
  (Definition~\ref{pd-family}) satisfies the \emph{local bound at the twisted
  places above $2$ and $3$} if there are $C$ and $x_0$ such that for every
  $x\ge x_0$ with $x>x_1$, every member at $x$, with the datum, size data and
  twist choices that $\mathcal F$ gives it, and every setup for them, the
  bound at $2$ and $3$ with constant $C$ holds for the representations $\rho$
  of the hub data of the setup.
\end{enumerate}
\end{definition}
\leanref{Development/AbelianVarietyTheory/Definitions/NonGaloisDescent/BoundAtTwoThreeFamily.lean}{GlobalTheory.NLink.BoundAtTwoThree}
\leanref{Development/AbelianVarietyTheory/Definitions/NonGaloisDescent/BoundAtTwoThreeFamily.lean}{GlobalTheory.NLink.Bd}
\leanref{Development/AbelianVarietyTheory/Definitions/NonGaloisDescent/BoundAtTwoThreeFamily.lean}{GlobalTheory.NLink.midF}
\leanref{Formalism/SharpRestriction/LevelLetterCocycles/JointConditionsZpowers/CommutingEndomorphismCondition.lean}{Formalism.MainNotions.N.KillingConditions.HubData}
\leanref{Development/AutomorphicArithmetic/ThetaCorrespondence/ThetaPairLifts/ThirdPlaceSeries/BoundAtTwoThreeFamilyFon.lean}{GlobalTheory.NLink.SetupAlt}
\leanref{Development/AutomorphicArithmetic/ThetaCorrespondence/ThetaPairLifts/ThirdPlaceSeries/BoundAtTwoThreeFamilyFon.lean}{GlobalTheory.NLink.BoundAtTwoThreeFamilyFonAlt}

\begin{lemma}\label{asm-two-three-cost}
Let $D$ be a definite-side datum with ranks $m,t,r,n$, $f=2n$, host field with
integer $Q_S$ and auxiliary primes $q,d$, and $N=Q_Sqd$. Let size data for $D$
be given, with scale $(x,\varepsilon,J)$, $b=\exp(x^{\varepsilon})$ and
numbers $\eta$ and $S_c$, and parameter choices for them
(Definition~\ref{pd-family}(4)). Put $h_B=b+S_c+\eta+t+2$ and, for a prime
$\ell$, $s_\ell=\max(1,v_\ell(N))$ and $\lambda_\ell=s_\ell+\log_\ell(2h_B)+1$;
for real $A,A'$ put
\[
 \kappa_{A,A'}(\ell)=A\,f\,m\,h_B\,s_\ell+A'f^2t^2(r+t)\,\lambda_\ell.
\]
Let $A,A'\ge0$ and $C_0\ge1$, and suppose that $\log x\ge1$, $n\le x^{C_0}$ and
$t\le x^{C_0}$. Let $\ell$ be a prime with $v_\ell(N)=v_\ell(Q_S)$, and let
$C\ge18AC_0+2A'(3C_0+9)+5C_0+1$. Then
\[
 \kappa_{A,A'}(\ell)\le C\,f\,m\,b\log x+C\,x^{C}x^{2\varepsilon}.
\]
\end{lemma}
\leanref{Development/AutomorphicArithmetic/ResidueCharacters/NumCost.lean}{UpperTheory.N.Num.cost_le'}
\leanref{Development/AutomorphicArithmetic/ResidueCharacters/RowDef.lean}{UpperTheory.N.costC}
\begin{proof}
The parameter choices give $S_c,\eta,t\le b$, and $b\ge1$, so
$b\le h_B\le6b$. The host field has $Q_S\le n^2$
(Definition~\ref{hf-host-field-rows}(3)), so
$2^{v_\ell(Q_S)}\le Q_S\le x^{2C_0}$ and $v_\ell(N)=v_\ell(Q_S)\le
2C_0\log x/\log2\le3C_0\log x$. Since $C_0\log x\ge1$, also
$s_\ell\le3C_0\log x$. Hence
\[
 A\,f\,m\,h_B\,s_\ell\le A\,f\,m\,(6b)(3C_0\log x)=18AC_0\,f\,m\,b\log x
 \le C\,f\,m\,b\log x.
\]
Next, $\log(2h_B)\le\log(12b)\le x^{\varepsilon}+3$ and $\log\ell\ge\log2\ge1/2$,
so $\log_\ell(2h_B)\le2x^{\varepsilon}+6$. With $\log x\le x$ and
$x,x^{\varepsilon}\ge1$ this gives $\lambda_\ell\le(3C_0+9)\,x\,x^{\varepsilon}$.
The ranks satisfy $m\le f\le k\le r\le t$ (Definition~\ref{th-ranks}), so
$f\le x^{C_0}$ and $r+t\le2x^{C_0}$. Therefore
\[
 A'f^2t^2(r+t)\lambda_\ell\le2A'(3C_0+9)\,x^{5C_0+1}x^{\varepsilon}
 \le C\,x^{C}x^{2\varepsilon},
\]
since $C\ge2A'(3C_0+9)$, $C\ge5C_0+1$ and $x>1$. Adding the two bounds gives
the lemma.
\end{proof}

\begin{lemma}\label{asm-two-three-split-numerics}
Let $D$ be a definite-side datum with local data at the twisted places
(Definition~\ref{asm-two-three-local-data}), ranks $m,t,r,n$, $f=2n$, host
field $E$ with class number $h_E$, integer $Q_S$, auxiliary primes $q,d$ and
$N=Q_Sqd$, with size data and parameter choices for them. Let $b$, $h_B$,
$s_\ell$, $\lambda_\ell$ and $\kappa_{A,A'}$ be as in
Lemma~\ref{asm-two-three-cost}, and $V(\ell)$ as in
Definition~\ref{asm-local-bound-two-three}. Let $\rho=(\rho_{i,\pi})$ be a
family of representations $G_E\to GL_r(\mathcal O_i)$ that satisfies the
killing condition, and suppose that $\text{Ш}^1_p(E,k)$ has finite length
for every prime $p$ and $0<k<t$. Let $i$ be a local index, put $\ell=\ell_i$,
suppose that $s_0(i)=s_\ell$, and let $\pi\in Z^{\mathrm{el}}_i$.

For a place $u$ of $E$ over $\ell$ write $d_u=[E_u:\mathbf Q_\ell]$, and
$h^0_u(\pi)$ for the normalized length $\operatorname{len}_i$ of the
invariants of a decomposition group at $u$ in $(K_i/\mathcal O_i)^r$, with the
action through $\rho_{i,\pi}\otimes\chi_i^{-1}$. Put
$\delta_i=\sum_{0<k<t}\big(s_0(i)+v_\ell(k)\big)$. For a set $P$ of places over
$\ell$ let $h^0_i(P)$ be the largest normalized length of the invariants of a
decomposition group at $u$ in $K_i/\mathcal O_i$, with the action through
$\varepsilon_\ell^{-k}$, over $u\in P$ and $0<k<t$, and $0$ if $P$ is empty.
For a partition of the places of $E$ over $\ell$ into sets $P_{\mathrm{row}}$
and $P_{\mathrm{col}}$ and integers $e',c_0,c_B\ge0$ put
\[
 \begin{aligned}
 B_{\mathrm{split}}(e',c_0,c_B)={}&\sum_{u\in P_{\mathrm{col}}}\big(e'
 +d_um(e'+\delta_i)+2\delta_i+h^0_u(\pi)+d_ume'+rc_0\big)
 +\operatorname{sel}_i(\rho_{i,\pi})+V(\ell)+v_\ell(h_E)\\
 &+t\sum_{w\in P_{\mathrm{row}}}(d_w+2)c_B+4(r+t)(|P_{\mathrm{col}}|+1)\delta_i
 +t|P_{\mathrm{col}}|\,h^0_i(P_{\mathrm{col}})
 +t\sum_{v\in Q_i}\operatorname{len}_i\big(\mathcal O_i/(\zeta_{i,v}-1)\big).
 \end{aligned}
\]
For an integer $y\ge0$ let $\lfloor\log_\ell y\rfloor$ be the largest $a\ge0$
with $\ell^a\le y$, and $0$ if $y=0$. For an integer $h\ge0$ put
$\lambda'_\ell(h)=s_\ell+\lfloor\log_\ell2h\rfloor+\lfloor\log_\ell N\rfloor+1$,
and put $n'=4n(r+t)$. Let such a partition be
given, let $h_1,h_2,h_3\ge0$ be integers with $h_1,h_2,h_3\le4h_B$, put
\[
 e'=32(2h_1s_\ell+1)+\lambda'_\ell(h_1),\qquad
 c_0=20(n'+1)\lambda'_\ell(h_2),\qquad c_B=40(t+1)(n'+1)\lambda'_\ell(h_3),
\]
and suppose that:
\begin{enumerate}
\item $h^0_u(\pi)\le rc_0$ for every $u\in P_{\mathrm{col}}$;
\item $h^0_i(P_{\mathrm{col}})\le s_0(i)+\lfloor\log_\ell t\rfloor$;
\item $\zeta_{i,v}\ne1$ for every $v\in Q_i$;
\item if $\operatorname{Sel}_i(\rho_{i,\pi})$ has finite length, the groups
  $\text{Ш}^1_\ell(E,k)$ with $0<k<t$ have finite length, and
  $\zeta_{i,v}\ne1$ for every $v\in Q_i$, then
  $c_i(\pi)\le B_{\mathrm{split}}(e',c_0,c_B)$.
\end{enumerate}
Then
\[
 c_i(\pi)\le\operatorname{sel}_i(\rho_{i,\pi})+\kappa_{6000,30000}(\ell)
 +V(\ell)+v_\ell(h_E)+\ell^{-2}.
\]
\end{lemma}
\leanref{Development/AutomorphicArithmetic/ResidueCharacters/NumRow.lean}{UpperTheory.N.bound_smallPrime_of_splitBound_cyclotomic_of_level_eq}
\leanref{Formalism/SharpRestriction/LevelLetterCocycles/EmbeddingTypes.lean}{Formalism.MainNotions.EmbeddingTypes.boundSplit}
\leanref{Formalism/SharpRestriction/LevelLetterCocycles/EmbeddingTypes.lean}{Formalism.MainNotions.EmbeddingTypes.ShaBoundTheoremB}
\begin{proof}
The parameter choices give $1\le b\le h_B$, $N\le b$ and $t\le b$. With
$h_1,h_2,h_3\le4h_B$ and $N\le8h_B$ this gives $e'\le1000\,h_Bs_\ell$,
$c_0\le1000\,f(r+t)\lambda_\ell$ and $c_B\le1000\,tf(r+t)\lambda_\ell$, and also
$\lfloor\log_\ell t\rfloor\le\lambda_\ell$, $s_0(i)+1\le\lambda_\ell$ and
$\delta_i\le t\lambda_\ell$. The killing condition makes
$\operatorname{Sel}_i(\rho_{i,\pi})$ of finite length and gives
$t\sum_{v\in Q_i}\operatorname{len}_i(\mathcal O_i/(\zeta_{i,v}-1))\le\ell^{-2}$.
With (3) and the finiteness of the groups $\text{Ш}^1_\ell(E,k)$, hypothesis
(4) gives $c_i(\pi)\le B_{\mathrm{split}}(e',c_0,c_B)$. The local degrees $d_u$
sum to at most $2f$ over $P_{\mathrm{row}}$ and over $P_{\mathrm{col}}$, and
these sums bound $|P_{\mathrm{row}}|$ and $|P_{\mathrm{col}}|$. With these
bounds, (1), (2), $1\le m\le f$ and $r\le t$, the terms of
$B_{\mathrm{split}}$ other than $\operatorname{sel}_i(\rho_{i,\pi})$, $V(\ell)$,
$v_\ell(h_E)$ and the sum over $Q_i$ add up to at most
$6c\,fmh_Bs_\ell+30c\,f^2t^2(r+t)\lambda_\ell$ with $c=1000$, which is
$\kappa_{6000,30000}(\ell)$.
\end{proof}

\begin{definition}\label{asm-two-three-local-killing}
Let $\ell$ be a prime, fix an algebraic closure $\bar{\mathbf Q}_\ell$ of
$\mathbf Q_\ell$, and for a finite extension $K\subset\bar{\mathbf Q}_\ell$ of
$\mathbf Q_\ell$ put $G_K=\operatorname{Gal}(\bar{\mathbf Q}_\ell/K)$.
\begin{enumerate}
\item Let $D$ be a definite-side datum with local data at the twisted places
  (Definition~\ref{asm-two-three-local-data}), host field $E$ and
  $N=Q_Sqd$ for the integer $Q_S$ and the auxiliary primes $q,d$ of $E$. Let
  $i$ be a local index with $\ell_i=\ell$ and $w$ a place of $E$ over $\ell$.
  The \emph{comparison data} of $D$ at $w$ consist of: a finite extension
  $E_w\subset\bar{\mathbf Q}_\ell$ of $\mathbf Q_\ell$ and a field embedding
  $\Omega\to\bar{\mathbf Q}_\ell$ that maps $E$ into $E_w$ with dense image; an
  isomorphism of topological groups from the decomposition group at $w$ in
  $G_E$ onto $G_{E_w}$, compatible with this embedding; and a finite extension
  $L_w\subset\bar{\mathbf Q}_\ell$ of $\mathbf Q_\ell$ that contains
  $\tau(E_w)$ for every $\mathbf Q_\ell$-embedding
  $\tau\colon E_w\to\bar{\mathbf Q}_\ell$, with a ring isomorphism
  $\mathcal O_i\to\mathcal O_{L_w}$ over $\mathbf Z_\ell$ that is a
  homeomorphism. The datum requires $L_w$ to contain every
  $\operatorname{lcm}(\ell,N)$-th root of unity in $\bar{\mathbf Q}_\ell$.
\item Let $N\geq1$ be an integer and $\mu_N^\ell$ the group of
  $\operatorname{lcm}(\ell,N)$-th roots of unity in $\bar{\mathbf Q}_\ell$. A
  pair $(K,L)$ of finite extensions of $\mathbf Q_\ell$ in
  $\bar{\mathbf Q}_\ell$ is \emph{$N$-cyclotomic} if $4\mid N$ when $\ell=2$,
  $K\subset\mathbf Q_\ell(\mu_N^\ell)$ and $\mu_N^\ell\subset L$. For integers
  $h,n\geq0$ put
  \[
   c_\ell(N,h,n)=20(n+1)\big(\max(1,v_\ell(N))+\lfloor\log_\ell2h\rfloor
   +\lfloor\log_\ell N\rfloor+1\big),
  \]
  with $\lfloor\log_\ell y\rfloor$ as in
  Lemma~\ref{asm-two-three-split-numerics}.
\item The \emph{invariant-killing property} holds at $\ell$ if the following
  is true for all integers $N\geq1$ and $h,n\geq0$, every $N$-cyclotomic pair
  $(K,L)$, every finite-dimensional $L$-vector space $V$ with an $L$-linear
  action of $G_K$ such that $g\mapsto gx$ is continuous for every $x\in V$,
  every $G_K$-stable $\mathcal O_L$-lattice $T\subset V$ and every set $U$ of
  at most $n$ integers in $[-h,h]$: if $V$ is crystalline and, for every
  $\mathbf Q_\ell$-embedding $\tau\colon K\to L$, the Hodge-Tate weights of $V$
  along $\tau$ lie in $U$ and do not include $0$, then $\ell^{c_\ell(N,h,n)}$
  kills every $G_K$-invariant element of $V/T$. The Hodge-Tate weights along
  $\tau$ are the integers $a$ with $\operatorname{Fil}^{a+1}\neq
  \operatorname{Fil}^a$ on the $\tau$-part of $D_{\mathrm{dR}}(V)$
  (Definition~\ref{gal-p-adic-hodge}).
\end{enumerate}
\end{definition}
\leanref{Formalism/SharpRestriction/LevelLetterCocycles/DataNLoc.lean}{Formalism.BasicNotions.LocBridge}
\leanref{Development/AutomorphicArithmetic/ResidueCharacters/InvariantsKilledOverCyclotomicAlt.lean}{UpperTheory.N.Cyc}
\leanref{Development/AutomorphicArithmetic/ResidueCharacters/InvariantsKilledOverCyclotomicAlt.lean}{UpperTheory.N.c0L}
\leanref{Development/AutomorphicArithmetic/ResidueCharacters/InvariantsKilledOverCyclotomicAlt.lean}{UpperTheory.N.InvariantsKilledOverCyclotomic}

\begin{lemma}\label{asm-two-three-local-invariants}
Let $D$ be a definite-side datum with local data at the twisted places
(Definition~\ref{asm-two-three-local-data}), ranks $r,t,n$ and host field $E$.
Let $i$ be a local index, put $\ell=\ell_i$, and let $N\geq1$ be an integer
such that for every place $w$ of $E$ over $\ell$, with comparison data $E_w$
and $L_w$, and every finite extension $L'\subset\bar{\mathbf Q}_\ell$ of $L_w$,
the pair $(E_w,L')$ is $N$-cyclotomic. Suppose that the invariant-killing
property holds at $\ell$ (Definition~\ref{asm-two-three-local-killing}). Let
$\rho\colon G_E\to GL_r(\mathcal O_i)$ be a homomorphism such that
$W'=\rho\otimes\chi_i^{-1}$ is continuous, and let $h\geq0$ be an integer. For
a place $w$ of $E$ over $\ell$, with comparison data $E_w$ and $L_w$, and a
finite extension $L'\subset\bar{\mathbf Q}_\ell$ of $L_w$, let $W'_{w,L'}$ be
the representation of $G_{E_w}$ on $L'^r$ obtained from the restriction of
$W'$ to the decomposition group at $w$ through the comparison data and the
inclusion $\mathcal O_{L_w}\subset L'$. Suppose that for every place $w$ of $E$
over $\ell$ and every such $L'$, the representation $W'_{w,L'}$ is crystalline
and, for every $\mathbf Q_\ell$-embedding $\tau\colon E_w\to L'$, the
Hodge-Tate weights of $W'_{w,L'}$ along $\tau$ lie in a set of at most $r$
integers in $[-h,h]$ and do not include $0$. Then for every place $u$ of $E$
over $\ell$, the normalized length $\operatorname{len}_i$ of the invariants of
the decomposition group at $u$ in $(K_i/\mathcal O_i)^r$, with the action
through $W'$, is at most $r\,c_\ell(N,h,4n(r+t))$.
\end{lemma}
\leanref{Development/AutomorphicArithmetic/ResidueCharacters/InvariantsLengthAdmissibleOfNonzeroJumps.lean}{UpperTheory.N.length_invariants_Wprime_divisible_le_of_nonzero_jumps}
\leanref{Development/AutomorphicArithmetic/ResidueCharacters/InvariantsLengthAdmissibleOfNonzeroJumps.lean}{UpperTheory.N.kill_of}
\begin{proof}
Put $n'=4n(r+t)$ and $c=c_\ell(N,h,n')$. Let $u$ be a place of $E$ over
$\ell$, with comparison data $E_u$ and $L_u$, and put $W'_u=W'_{u,L_u}$; the
proof uses the hypotheses only for $w=u$ and $L'=L_u$.

First, $(E_u,L_u)$ is $N$-cyclotomic, by hypothesis.

Second, since $L_u$ contains the images of all embeddings of $E_u$ in
$\bar{\mathbf Q}_\ell$, the number of $\mathbf Q_\ell$-embeddings
$\tau\colon E_u\to L_u$ is $[E_u:\mathbf Q_\ell]$. Since $E$ is dense in $E_u$,
$[E_u:\mathbf Q_\ell]\leq[E:\mathbf Q]=4n$
(Proposition~\ref{hf-host-field-properties}(1)).

Third, for each $\tau$ choose a set $U_\tau$ of at most $r$ integers in
$[-h,h]$ that contains the Hodge-Tate weights of $W'_u$ along $\tau$, and let
$U$ be the union of the $U_\tau$. Then $U$ has at most $4nr\leq n'$ elements,
all in $[-h,h]$, it contains the Hodge-Tate weights of $W'_u$ along every
$\tau$, and none of these is $0$.

Fourth, apply the invariant-killing property at $\ell$ to $N$, $h$, $n'$, the
pair $(E_u,L_u)$, the representation $W'_u$ on $V=L_u^r$, the lattice
$T=\mathcal O_{L_u}^r$ and the set $U$. It shows that $\ell^c$ kills every
$G_{E_u}$-invariant element of $L_u^r/\mathcal O_{L_u}^r$.

Fifth, the isomorphism $\mathcal O_i\to\mathcal O_{L_u}$ of the comparison
data extends to the fraction fields and induces an injective map
$(K_i/\mathcal O_i)^r\to L_u^r/\mathcal O_{L_u}^r$. It is compatible with the
action of the decomposition group at $u$ through $W'$ and the action of
$G_{E_u}$ through $W'_u$, under the isomorphism of these groups. So it maps an
invariant $x$ of the decomposition group to a $G_{E_u}$-invariant element,
which $\ell^c$ kills; by injectivity, $\ell^cx=0$.

Finally, the invariants form an $\mathcal O_i$-submodule of
$(K_i/\mathcal O_i)^r$ killed by $\ell^c$. It lies in
$(\ell^{-c}\mathcal O_i/\mathcal O_i)^r$, whose length is
$rc\operatorname{len}_{\mathcal O_i}(\mathcal O_i/\ell\mathcal O_i)$, so its
normalized length is at most $rc$.
\end{proof}

\begin{lemma}\label{asm-two-three-local-zeta}
Let $D$ be a definite-side datum with local data at the twisted places
(Definition~\ref{asm-two-three-local-data}), let $i$ be a local index and
$v\in Q_i$, and let $I_v$ be the inertia group at $v$ that contains
$\gamma_{i,v}$. Suppose that $\chi_i$ is trivial on the wild inertia subgroup
of $I_v$. Then $\zeta_{i,v}\neq1$.
\end{lemma}
\leanref{Formalism/SharpRestriction/LevelLetterCocycles/JointConditionsZpowers/GaloisConditionFiveExtraOf.lean}{Formalism.MainNotions.N.galois_condition_five_extra_of}
\leanref{Formalism/SharpRestriction/LevelLetterCocycles/JointConditionsZpowers/GaloisConditionFiveExtraOf.lean}{Formalism.MainNotions.N.inertia_character_mem_zpowers}
\begin{proof}
Let $P_v\subset I_v$ be the wild inertia subgroup, and suppose that
$\zeta_{i,v}=1$. The subgroup $H$ of $G_E$ of the $g$ with $\chi_i(g)$ in the
finite group generated by $\zeta_{i,v}$ is the preimage of a finite set under
the continuous map $\chi_i$, so it is closed, since $\mathcal O_i$ is
Hausdorff. It contains $\gamma_{i,v}$, and it contains $P_v$ because $\chi_i$
is trivial there. So $H$ contains the closure of the subgroup generated by
$\gamma_{i,v}$ and $P_v$, which contains $I_v$
(Definition~\ref{asm-two-three-local-data}(5)). Since $\zeta_{i,v}=1$, this
says that $\chi_i$ is trivial on $I_v$. But $\chi_i$ is ramified at every
place of $Q_i$ (Definition~\ref{asm-two-three-local-data}(5)), so
$\chi_i(g)\neq1$ for some $g\in I_v$, a contradiction.
\end{proof}

\subsection{Cocycles of the twisted action}\label{subsection-asm-twisted-cocycles}

\begin{definition}\label{asm-two-three-killing-wach-frame}
Let $\ell$ be a prime and fix an algebraic closure $\bar{\mathbf Q}_\ell$ of
$\mathbf Q_\ell$. Let $\mathbf C_\ell$ be its completion, with the absolute
value normalized by $|\ell|=\ell^{-1}$, valuation ring $\mathcal O_{\mathbf C}$
and the continuous action of $G_{\mathbf Q_\ell}=\operatorname{Gal}(\bar{\mathbf
Q}_\ell/\mathbf Q_\ell)$, which acts on $\mathbf C_\ell^r$ and on matrices over
$\mathbf C_\ell$ entry by entry, and let $\chi\colon G_{\mathbf Q_\ell}\to\mathbf
Z_\ell^\times$ be the cyclotomic character. For a subfield
$F\subset\bar{\mathbf Q}_\ell$ put $G_F=\operatorname{Gal}(\bar{\mathbf Q}_\ell/F)$.
For a finite extension $K\subset\bar{\mathbf Q}_\ell$ of $\mathbf Q_\ell$ let
$K_0$ be its maximal unramified subextension, $K_j$ the field generated over
$K_0$ by the $\ell^j$-th roots of unity for $j\geq0$, and $K_\infty$ the union
of the $K_j$. Let $m\geq1$ be an integer, with $m\geq2$ if $\ell=2$, and let
$r\geq0$.
\begin{enumerate}
\item A \emph{frame of rank $r$ at level $m$} for $K$ is a pair $(U,g)$ of a map
  $U\colon G_{K_m}\to M_r(\mathbf C_\ell)$ and an element $g\in G_{K_m}$ such
  that: the entries of $U$ are continuous and take values in the valuation ring
  of $K_m$; $U(\sigma\sigma')=U(\sigma)U(\sigma')$ for
  $\sigma,\sigma'\in G_{K_m}$; $U(\sigma)=1$ for $\sigma\in G_{K_\infty}$;
  $|\chi(g)-1|=\ell^{-m}$; and every entry of $U(g)-1$ has absolute value less
  than $\ell^{-1/(\ell-1)}$.
\item Let $(U,g)$ be such a frame and $i\in\mathbf Z$. The \emph{twisted action}
  of $\sigma\in G_{K_m}$ on $\mathbf C_\ell^r$ is
  $\sigma*_iy=\chi(\sigma)^iU(\sigma)\sigma(y)$. A \emph{cocycle} for it is a
  continuous map $c\colon G_{K_m}\to\mathbf C_\ell^r$ with
  $c(\sigma\sigma')=c(\sigma)+\sigma*_ic(\sigma')$, and $c$ is the
  \emph{coboundary} of $x\in\mathbf C_\ell^r$ if $c(\sigma)=\sigma*_ix-x$ for
  every $\sigma\in G_{K_m}$.
\item A \emph{Hodge-Tate basis} of the frame $(U,g)$ with weights in a set
  $J\subset\mathbf Z$ is a basis $e_1,\ldots,e_r$ of $\mathbf C_\ell^r$ together
  with integers $u_1,\ldots,u_r\in J$ such that $\sigma*_{u_k}e_k=e_k$ for every
  $\sigma\in G_{K_m}$ and every $k$.
\item For a finite set $J\subset\mathbf Z$ and $i\in\mathbf Z$ put
  \[
   \nu_m(J,i)=\max\Big(1,\sum_{u\in J,\,u\neq i}\big(m+v_\ell(u-i)\big)\Big).
  \]
\end{enumerate}
\end{definition}
\leanref{Development/AbelianVarietyTheory/Definitions/Foundations/HTBasis.lean}{GlobalTheory.Cyc.Frame}
\leanref{Development/AbelianVarietyTheory/Definitions/Foundations/HTBasis.lean}{GlobalTheory.Cyc.twAct}
\leanref{Development/AbelianVarietyTheory/Definitions/Foundations/HTBasis.lean}{GlobalTheory.Cyc.HTBasis}
\leanref{Development/AbelianVarietyTheory/Definitions/Foundations/HTBasis.lean}{GlobalTheory.Cyc.CocOn}
\leanref{Development/AbelianVarietyTheory/Definitions/Foundations/HTBasis.lean}{GlobalTheory.Cyc.CobOn}

\begin{theorem}[Tate, Sen {\cite[Theorem~2.2.7]{BC09}}]\label{asm-two-three-killing-wach-tate-sen}
Let $\ell$ be a prime and use the notation of
Definition~\ref{asm-two-three-killing-wach-frame}. Let
$K\subset\bar{\mathbf Q}_\ell$ be a finite extension of $\mathbf Q_\ell$, let
$m\geq1$ be an integer, with $m\geq2$ if $\ell=2$, let $i\in\mathbf Z$, and let
$y\in\mathbf C_\ell$ satisfy $\chi(\sigma)^i\sigma(y)=y$ for every
$\sigma\in G_{K_m}$. If $i\neq0$, then $y=0$. If $i=0$, then $y\in K_m$.
\end{theorem}
\leanref{Development/AbelianVarietyTheory/Proofs/InertiaEigenvalues/TateInvariantsCyclotomicTwist.lean}{GlobalTheory.CyclotomicTowerSubgroupInclusions.invariants_twist}

\begin{theorem}[Tate {\cite[Theorem~13.1.2]{BC09}}]\label{asm-two-three-killing-wach-torsion-trace}
Let $\ell$ be a prime and use the notation of
Definition~\ref{asm-two-three-killing-wach-frame}. Let
$K\subset\bar{\mathbf Q}_\ell$ be a finite extension of $\mathbf Q_\ell$ and put
$H=G_{K_\infty}$. Let $L\subset\bar{\mathbf Q}_\ell$ be a field with
$K_\infty\subset L$ and $h(L)\subset L$ for every $h\in H$, and let $S\subset H$
be a finite set such that for every $h\in H$ there is exactly one $s\in S$ with
$s^{-1}h\in G_L$. Then for every $a\in K_\infty$ with $|a|<1$ there is $b\in L$
with $|b|\leq1$ and $\sum_{s\in S}s(b)=a$.
\end{theorem}
\leanref{Development/AbelianVarietyTheory/Proofs/InertiaEigenvalues/CyclotomicTowerTraceAlmostSurjectiveStatement.lean}{GlobalTheory.T.cyclotomicTower_trace_almost_surjective}

\begin{lemma}\label{asm-two-three-killing-wach-torsion-average}
Let $\ell$ be a prime and use the notation of
Definition~\ref{asm-two-three-killing-wach-frame}. Let
$K\subset\bar{\mathbf Q}_\ell$ be a finite extension of $\mathbf Q_\ell$, put
$H=G_{K_\infty}$, and let $a\in K_\infty$ with $|a|<1$. Let
$c\colon H\to\mathbf C_\ell$ be a continuous map with $c(hh')=c(h)+h(c(h'))$ for
$h,h'\in H$, and let $t\geq0$ and $\varepsilon>0$ be real numbers with
$|c(h)|\leq t$ for every $h\in H$. Then there is $x\in\mathbf C_\ell$ with
$|x|\leq t$ and $|a\,c(h)-(h(x)-x)|\leq\varepsilon$ for every $h\in H$.
\end{lemma}
\leanref{Development/AbelianVarietyTheory/Proofs/InertiaEigenvalues/CyclotomicTateMaximalIdealKillsFirstCohomology.lean}{GlobalTheory.Cyc.AveragingRankOne.step}
\leanref{Development/AbelianVarietyTheory/Proofs/InertiaEigenvalues/ImportsRest.lean}{GlobalTheory.PermutationExistence.main}
\begin{proof}
First, a field and a trace. The cocycle relation at $h=h'=1$ gives $c(1)=0$, so
by continuity there is an open neighbourhood $W$ of $1$ in
$G_{\mathbf Q_\ell}$ with $|c(h)|<\varepsilon$ for $h\in H\cap W$. Choose a
finite Galois extension $E\subset\bar{\mathbf Q}_\ell$ of $\mathbf Q_\ell$ with
$G_E\subset W$, and put $L=K_\infty E$. Every $h\in H$ maps $K_\infty$ and $E$
onto themselves, so $h(L)=L$. The group $G_L=H\cap G_E$ has finite index in
$H$; let $S\subset H$ be a set of representatives of the cosets $hG_L$ with
$h\in H$, so that every $h\in H$ has exactly one $s\in S$ with $s^{-1}h\in G_L$.
Theorem~\ref{asm-two-three-killing-wach-torsion-trace} gives $b\in L$ with
$|b|\leq1$ and $\sum_{s\in S}s(b)=a$.

Second, the average. Put $x=-\sum_{s\in S}s(b)\,c(s)$. Then $|x|\leq t$, since
$|s(b)|=|b|\leq1$ and $|c(s)|\leq t$. Let $h\in H$. For $s\in S$ let
$\sigma(s)\in S$ be the representative of the coset $hsG_L$, and put
$v_s=\sigma(s)^{-1}hs\in G_L\cap H$. Left multiplication by $h$ permutes the
cosets of $G_L$ in $H$, so $\sigma$ is a permutation of $S$. Since $v_s$ fixes
$b$, $h(s(b))=\sigma(s)(b)$. The cocycle relation for $hs=\sigma(s)v_s$ gives
$c(h)+h(c(s))=c(\sigma(s))+\sigma(s)(c(v_s))$. Hence
\[
 h\big(s(b)\,c(s)\big)=\sigma(s)(b)\,c(\sigma(s))
 +\sigma(s)\big(b\,c(v_s)\big)-\sigma(s)(b)\,c(h).
\]
Summing over $s\in S$, and using that $\sigma$ permutes $S$ and that
$\sum_{s\in S}\sigma(s)(b)=a$, we get
$h(x)=x-\sum_{s\in S}\sigma(s)(b\,c(v_s))+a\,c(h)$. So
\[
 a\,c(h)-\big(h(x)-x\big)=\sum_{s\in S}\sigma(s)\big(b\,c(v_s)\big).
\]
Each term has absolute value $|b|\,|c(v_s)|\leq\varepsilon$, as
$v_s\in H\cap G_E\subset H\cap W$, and the absolute value is ultrametric.
\end{proof}

\begin{lemma}\label{asm-two-three-killing-wach-torsion-tate}
Let $\ell$ be a prime and use the notation of
Definition~\ref{asm-two-three-killing-wach-frame}. Let
$K\subset\bar{\mathbf Q}_\ell$ be a finite extension of $\mathbf Q_\ell$, put
$H=G_{K_\infty}$, let $a\in K_\infty$ with $|a|<1$, and let $r\geq0$. Let
$c\colon H\to\mathcal O_{\mathbf C}^r$ be a continuous map with
$c(hh')=c(h)+h(c(h'))$ for $h,h'\in H$. Then there is
$x\in\mathcal O_{\mathbf C}^r$ with $a\,c(h)=h(x)-x$ for every $h\in H$.
\end{lemma}
\leanref{Development/AbelianVarietyTheory/Proofs/InertiaEigenvalues/CyclotomicTateMaximalIdealKillsFirstCohomologyStatement.lean}{GlobalTheory.T.cyclotomic_tate_maximalIdeal_kills_first_cohomology}
\leanref{Development/AbelianVarietyTheory/Proofs/InertiaEigenvalues/CyclotomicTateMaximalIdealKillsFirstCohomology.lean}{GlobalTheory.Cyc.AveragingRankOne.rank_one}
\begin{proof}
Each coordinate of $c$ satisfies the hypotheses with $r=1$, so we may assume
$r=1$. If $a=0$, take $x=0$. Let $a\neq0$. Since $a\in K_\infty$, $h(a)=a$ for
$h\in H$.

We construct $y_0,y_1,\ldots\in\mathbf C_\ell$ with $|y_n|\leq1$,
$|a\,c(h)-(h(y_n)-y_n)|\leq|a|2^{-n}$ for every $h\in H$, and
$|y_{n+1}-y_n|\leq2^{-n}$. Take $y_0=0$, which works since $|c(h)|\leq1$. Given
$y_n$, put $c_n(h)=c(h)-a^{-1}(h(y_n)-y_n)$. It is continuous, since the orbit
maps $h\mapsto h(y_n)$ are, it satisfies the cocycle relation, since $h$ fixes
$a$, and $|c_n(h)|=|a|^{-1}|a\,c(h)-(h(y_n)-y_n)|\leq2^{-n}$.
Lemma~\ref{asm-two-three-killing-wach-torsion-average}, with $t=2^{-n}$ and
$\varepsilon=|a|2^{-n-1}$, gives $x_n$ with $|x_n|\leq2^{-n}$ and
$|a\,c_n(h)-(h(x_n)-x_n)|\leq|a|2^{-n-1}$ for $h\in H$. Put
$y_{n+1}=y_n+x_n$. Then $|y_{n+1}|\leq\max(|y_n|,|x_n|)\leq1$, and
$a\,c(h)-(h(y_{n+1})-y_{n+1})=a\,c_n(h)-(h(x_n)-x_n)$.

The sequence $(y_n)$ is Cauchy; let $x$ be its limit, so $|x|\leq1$. For
$h\in H$ the map $h$ is continuous on $\mathbf C_\ell$, so
$|a\,c(h)-(h(x)-x)|$ is the limit of $|a\,c(h)-(h(y_n)-y_n)|\leq|a|2^{-n}$,
which is $0$.
\end{proof}

\begin{lemma}\label{asm-two-three-killing-wach-torsion-fixed}
Let $\ell$ be a prime and use the notation of
Definition~\ref{asm-two-three-killing-wach-frame}. Let
$K\subset\bar{\mathbf Q}_\ell$ be a finite extension of $\mathbf Q_\ell$, put
$H=G_{K_\infty}$, and let $\hat{\mathcal O}$ be the closure in $\mathbf C_\ell$ of
the set of $x\in K_\infty$ with $|x|\leq1$. Let $y\in\mathbf C_\ell$ with
$|y|\leq1$. Then $h(y)=y$ for every $h\in H$ if and only if
$y\in\hat{\mathcal O}$.
\end{lemma}
\leanref{Development/AbelianVarietyTheory/Proofs/InertiaEigenvalues/CyclotomicTowerFixedIffMemCompletedIntegersStatement.lean}{GlobalTheory.T.cyclotomicTower_fixed_iff_mem_completedIntegers}
\leanref{Development/AbelianVarietyTheory/Proofs/InertiaEigenvalues/CyclotomicTowerFixedIffMemCompletedIntegers.lean}{GlobalTheory.CycInvariants.fixed_iff_mem_Ohat}
\begin{proof}
If $y\in\hat{\mathcal O}$, then $y$ is a limit of elements of $K_\infty$, which
$h$ fixes, and $h$ is continuous, so $h(y)=y$.

Conversely, let $y$ be fixed by $H$ and let $0<\delta\leq1$. Choose
$x\in\bar{\mathbf Q}_\ell$ with $|y-x|<\delta/\ell$. The stabilizer of $x$ in
$G_{\mathbf Q_\ell}$ has finite index, so it contains a normal subgroup $N$ of
finite index. Let $L$ be the fixed field of $H\cap N$. Then $x\in L$,
$K_\infty\subset L$, $h(L)\subset L$ for $h\in H$ since $N$ is normal, and
$G_L$, which
contains $H\cap N$, has finite index in $H$. Let $S\subset H$ be a set of
representatives of the cosets $hG_L$ with $h\in H$.
Theorem~\ref{asm-two-three-killing-wach-torsion-trace} with $a=\ell$ gives
$b\in L$ with $|b|\leq1$ and $\sum_{s\in S}s(b)=\ell$. Put
$w=\sum_{s\in S}s(bx)$. Since $bx\in L$, $s(bx)$ depends only on the coset
$sG_L$, and every $h\in H$ permutes these cosets, so $h(w)=w$; as $K_\infty$ is
the fixed field of $H$ in $\bar{\mathbf Q}_\ell$, $w\in K_\infty$. Since every
$s\in S$ fixes $y$,
\[
 \ell y-w=\sum_{s\in S}s\big(b(y-x)\big),
\]
so $|\ell y-w|\leq\delta/\ell$, that is, $|y-w/\ell|\leq\delta$. Then
$w/\ell\in K_\infty$ and $|w/\ell|\leq\max(|y|,\delta)\leq1$. As $\delta$ was
arbitrary, $y\in\hat{\mathcal O}$.
\end{proof}

\begin{lemma}\label{asm-two-three-killing-wach-torsion-cyclotomic}
Let $\ell$ be a prime and use the notation of
Definition~\ref{asm-two-three-killing-wach-frame}, and let
$K\subset\bar{\mathbf Q}_\ell$ be a finite extension of $\mathbf Q_\ell$.
\begin{enumerate}
\item The character $\chi$ is continuous and $|\chi(\sigma)|=1$ for every
  $\sigma$. Let $\sigma\in G_{K_0}$. For $j\geq0$, $\sigma\in G_{K_j}$ if and
  only if $|\chi(\sigma)-1|\leq\ell^{-j}$; and $\sigma\in G_{K_\infty}$ if and
  only if $\chi(\sigma)=1$.
\item Let $j\geq1$, with $j\geq2$ if $\ell=2$, and let $u\in\mathbf Q_\ell$
  with $|u-1|=\ell^{-j}$. Then $|u^n-1|=\ell^{-j-v_\ell(n)}$ for every nonzero
  integer $n$.
\item If $w\in\mathbf Q_\ell$ and $|w|\leq1$, then
  $|w^k-1-k(w-1)|\leq|w-1|^2$ for every integer $k\geq0$.
\item If $u\in\mathbf C_\ell$ and $|u-1|<1$, then $|u^n-1|\leq|u-1|$ for every
  integer $n$.
\item If $\zeta\in\mathbf C_\ell$ is a primitive $\ell$-th root of unity, then
  $|\zeta-1|=\ell^{-1/(\ell-1)}$.
\end{enumerate}
\end{lemma}
\leanref{Development/AbelianVarietyTheory/Proofs/InertiaEigenvalues/CyclotomicCharacterPadicBasicPropertiesStatement.lean}{GlobalTheory.T.cyclotomicCharacterPadic_basic_properties}
\leanref{Development/AbelianVarietyTheory/Proofs/InertiaEigenvalues/PadicNormZpowSubOnePadicValIntStatement.lean}{GlobalTheory.T.padic_norm_zpow_sub_one_padicValInt}
\leanref{Development/AbelianVarietyTheory/Proofs/InertiaEigenvalues/PowSubOneLe.lean}{GlobalTheory.UltrametricPowerBound.main}
\leanref{Development/AbelianVarietyTheory/Proofs/InertiaEigenvalues/NormZpowSubOneLe.lean}{GlobalTheory.UnitNormPowerSubOne.norm_zpow_sub_one_le}
\leanref{Development/AbelianVarietyTheory/Proofs/InertiaEigenvalues/NormPrimitiveRootPrimePowSubOnePadicComplexStatement.lean}{GlobalTheory.T.norm_primitiveRoot_primePow_sub_one_padicComplex}
\begin{proof}
(1) For every $j\geq0$ and every $\ell^j$-th root of unity $\zeta$,
$\sigma(\zeta)=\zeta^a$ for an integer $a\geq0$ with
$|\chi(\sigma)-a|\leq\ell^{-j}$, and $\zeta^a$ depends only on $a$ modulo
$\ell^j$. So $\sigma$ fixes the $\ell^j$-th roots of unity if and only if
$|\chi(\sigma)-1|\leq\ell^{-j}$; as $K_j$ is generated over $K_0$ by them and
$K_\infty$ is the union of the $K_j$, this gives both claims. Parts (2) to (5)
are self-contained computations with $\ell$-adic absolute values.
\end{proof}

\begin{lemma}\label{asm-two-three-killing-wach-torsion-dense}
Let $\ell$ be a prime and use the notation of
Definition~\ref{asm-two-three-killing-wach-frame}. Let
$K\subset\bar{\mathbf Q}_\ell$ be a finite extension of $\mathbf Q_\ell$, put
$H=G_{K_\infty}$, let $m\geq1$ be an integer, with $m\geq2$ if $\ell=2$, and let
$g\in G_{K_m}$ with $|\chi(g)-1|=\ell^{-m}$. Then:
\begin{enumerate}
\item for every $\sigma\in G_{K_m}$ and every $n\geq0$ there is an integer
  $a\geq0$ with $\sigma g^{-a}\in G_{K_n}$;
\item the set of products $hg^a$ with $h\in H$ and $a\geq0$ is dense in
  $G_{K_m}$.
\end{enumerate}
\end{lemma}
\leanref{Development/AbelianVarietyTheory/Proofs/InertiaEigenvalues/ExistsCyclotomicTowerLevelGeneratorProperties.lean}{GlobalTheory.Cyc.gen_dense}
\leanref{Development/AbelianVarietyTheory/Proofs/InertiaEigenvalues/ExistsCyclotomicTowerLevelGeneratorProperties.lean}{GlobalTheory.Cyc.exists_norm_sub_pow_le}
\begin{proof}
Put $u=\chi(g)$, so $|u|=1$ and $|u-1|=\ell^{-m}$.

First, let $s\in\mathbf Q_\ell$ with $|s-1|\leq\ell^{-m}$; we show that for
every $n\geq0$ there is $a\geq0$ with $|s-u^a|\leq\ell^{-m-n}$, by induction on
$n$. For $n=0$ take $a=0$. Let $a$ work for $n$, and put $w=u^{\ell^n}$ and
$t=u^a$. Then $|t|=1$ and $|w-1|=\ell^{-m-n}$ by
Lemma~\ref{asm-two-three-killing-wach-torsion-cyclotomic}(2), so
$(s-t)/(t(w-1))$ has absolute value at most $1$, and there is an integer
$k\geq0$ with $|(s-t)/(t(w-1))-k|\leq\ell^{-1}$. Now
\[
 s-u^{a+\ell^nk}=t(w-1)\Big(\frac{s-t}{t(w-1)}-k\Big)
 -t\big(w^k-1-k(w-1)\big).
\]
The first term has absolute value at most $\ell^{-m-n-1}$. By
Lemma~\ref{asm-two-three-killing-wach-torsion-cyclotomic}(3) the second has
absolute value at most $|w-1|^2=\ell^{-2(m+n)}\leq\ell^{-m-n-1}$, as $m\geq1$.
So $a+\ell^nk$ works for $n+1$.

(1) Let $\sigma\in G_{K_m}$ and $n\geq0$. Since $\sigma\in G_{K_0}$,
Lemma~\ref{asm-two-three-killing-wach-torsion-cyclotomic}(1) gives
$|\chi(\sigma)-1|\leq\ell^{-m}$. By the first step there is $a\geq0$ with
$|\chi(\sigma)-u^a|\leq\ell^{-m-n}$. Then
$|\chi(\sigma g^{-a})-1|=|\chi(\sigma)-u^a|\,|u|^{-a}\leq\ell^{-n}$, and
$\sigma g^{-a}\in G_{K_0}$, so $\sigma g^{-a}\in G_{K_n}$ by
Lemma~\ref{asm-two-three-killing-wach-torsion-cyclotomic}(1).

(2) Let $\sigma\in G_{K_m}$ and let $V$ be an open set containing $\sigma$. The
set $W$ of $x\in G_{\mathbf Q_\ell}$ with $hx^{-1}\sigma\in V$ for some
$h\in H$ is open, and it contains $H$. The groups $G_{K_n}$ are closed and
decrease with $n$, and their intersection is $H$: an element of all of them lies
in $G_{K_0}$ and has $|\chi-1|\leq\ell^{-n}$ for every $n$, so it lies in $H$ by
Lemma~\ref{asm-two-three-killing-wach-torsion-cyclotomic}(1). Since
$G_{\mathbf Q_\ell}$ is compact, $G_{K_n}\subset W$ for some $n$. By (1) there
is $a\geq0$ with $\sigma g^{-a}\in G_{K_n}\subset W$, so there is $h\in H$ with
$h(\sigma g^{-a})^{-1}\sigma=hg^a\in V$.
\end{proof}

\begin{lemma}\label{asm-two-three-killing-wach-torsion-procyclic}
Let $\ell$ be a prime and use the notation of
Definition~\ref{asm-two-three-killing-wach-frame}. Let
$K\subset\bar{\mathbf Q}_\ell$ be a finite extension of $\mathbf Q_\ell$, let
$m\geq1$ be an integer, with $m\geq2$ if $\ell=2$, let $(U,g)$ be a frame of
rank $r$ at level $m$ for $K$, and let $i\in\mathbf Z$. Let $c$ be a cocycle for
the twisted action $*_i$. If $c(h)=0$ for every $h\in G_{K_\infty}$ and
$c(g)=0$, then $c(\sigma)=0$ for every $\sigma\in G_{K_m}$.
\end{lemma}
\leanref{Development/AbelianVarietyTheory/Proofs/InertiaEigenvalues/CocycleVanishingOnFixingSubgroupDeterminedByGeneStatement.lean}{GlobalTheory.T.cocycle_vanishing_on_fixingSubgroup_determined_by_generator}
\leanref{Development/AbelianVarietyTheory/Proofs/InertiaEigenvalues/CocycleVanishingOnFixingSubgroupDeterminedByGene.lean}{GlobalTheory.Cyc.Procyclic.procyclic}
\begin{proof}
Since $\sigma*_i0=0$, the cocycle relation gives
$c(g^{a+1})=c(g^a)+g^a*_ic(g)=c(g^a)$, so $c(g^a)=c(1)=0$ for every $a\geq0$.
For $h\in G_{K_\infty}$, $c(hg^a)=c(h)+h*_ic(g^a)=0$. By
Lemma~\ref{asm-two-three-killing-wach-torsion-dense}(2) these products are dense
in $G_{K_m}$, and $c$ is continuous on $G_{K_m}$, so $c$ vanishes on $G_{K_m}$.
\end{proof}

\begin{theorem}[{\cite[Chapter~IV, \S4]{Ser79}}]\label{asm-two-three-killing-wach-torsion-decomposition-ramification}
Let $\ell$ be a prime and use the notation of
Definition~\ref{asm-two-three-killing-wach-frame}. Let
$K\subset\bar{\mathbf Q}_\ell$ be a finite extension of $\mathbf Q_\ell$, let
$j\geq1$ be an integer, and put $\varphi_j=\ell^{j-1}(\ell-1)$. Then
$|\zeta-1|=\ell^{-1/\varphi_j}$ for every primitive $\ell^j$-th root of unity
$\zeta\in\mathbf C_\ell$, and every nonzero $x\in K_j$ has
$|x|=\ell^{a/\varphi_j}$ for some integer $a$.
\end{theorem}
\leanref{Development/AbelianVarietyTheory/Proofs/InertiaEigenvalues/NormPrimitiveRootPrimePowSubOnePadicComplex.lean}{GlobalTheory.Cyc.norm_primitiveRoot_sub_one}
\leanref{Development/AbelianVarietyTheory/Proofs/InertiaEigenvalues/NormMemCyclotomicLayerPrimeRpowStatement.lean}{GlobalTheory.T.norm_mem_cyclotomic_layer_prime_rpow}
\leanref{Development/AbelianVarietyTheory/Proofs/InertiaEigenvalues/NormMemCyclotomicLayerPrimeRpow.lean}{GlobalTheory.Cyc.norm_of_mem_Ks}

\begin{lemma}\label{asm-two-three-killing-wach-torsion-decomposition-layer}
Let $\ell$ be a prime and use the notation of
Definition~\ref{asm-two-three-killing-wach-frame}. Let
$K\subset\bar{\mathbf Q}_\ell$ be a finite extension of $\mathbf Q_\ell$, let
$n\geq2$ be an integer, and let $\zeta\in\bar{\mathbf Q}_\ell$ be a primitive
$\ell^n$-th root of unity.
\begin{enumerate}
\item Every $x\in K_n$ with $|x|\leq1$ is $\sum_{a=0}^{\ell-1}\zeta^ay_a$ for
  some $y_0,\ldots,y_{\ell-1}\in K_{n-1}$ with $|y_a|\leq1$.
\item If $y_0,\ldots,y_{\ell-1}\in K_{n-1}$ and
  $\sum_{a=0}^{\ell-1}\zeta^ay_a=0$, then $y_a=0$ for every $a$.
\end{enumerate}
\end{lemma}
\leanref{Development/AbelianVarietyTheory/Proofs/InertiaEigenvalues/IntegersEqSumZetaPowStatement.lean}{GlobalTheory.T.integers_eq_sum_zeta_pow_unique}
\leanref{Development/AbelianVarietyTheory/Proofs/InertiaEigenvalues/IntegersEqSumZetaPow.lean}{GlobalTheory.Cyc.integers_eq_sum_zeta_pow}
\begin{proof}
Put $\varphi=\ell^{n-1}(\ell-1)$, so that $\ell^{n-2}(\ell-1)=\varphi/\ell$.

First, the terms of an expansion in powers of $\zeta-1$ have distinct absolute
values. Let $z_0,\ldots,z_{\ell-1}\in K_{n-1}$. By
Theorem~\ref{asm-two-three-killing-wach-torsion-decomposition-ramification},
applied with $j=n-1\geq1$ and with $j=n$, every nonzero $z_b$ has
$|z_b|=\ell^{k_b\ell/\varphi}$ for an integer $k_b$, and
$|\zeta-1|=\ell^{-1/\varphi}$. So $|z_b(\zeta-1)^b|=\ell^{(k_b\ell-b)/\varphi}$.
For $0\leq b<c<\ell$ the integers $k_b\ell-b$ and $k_c\ell-c$ are different, as
$b\not\equiv c\pmod\ell$. So the nonzero terms $z_b(\zeta-1)^b$ have distinct
absolute values, and since the absolute value is ultrametric,
\[
 |z_b(\zeta-1)^b|\leq\Big|\sum_{c=0}^{\ell-1}z_c(\zeta-1)^c\Big|
 \qquad(0\leq b<\ell).\tag{$*$}
\]
If moreover $z_b\neq0$ and $|z_b(\zeta-1)^b|\leq1$, then $k_b\ell\leq b<\ell$,
so $k_b\leq0$ and $|z_b|\leq1$.

Second, $\zeta^\ell$ is an $\ell^{n-1}$-th root of unity, so it lies in
$K_{n-1}$, and $K_n=K_{n-1}(\zeta)$, since every $\ell^n$-th root of unity is a
power of $\zeta$. As $\zeta$ is a root of $X^\ell-\zeta^\ell$, every $x\in K_n$
is $p(\zeta)$ for a polynomial $p$ over $K_{n-1}$ of degree less than $\ell$.
Such a $p$ can be written $p(X)=\sum_{b=0}^{\ell-1}q_b(X-1)^b$ with
$q_b\in K_{n-1}$. Expanding $(X-1)^b$, and $X^a=((X-1)+1)^a$, shows that the
coefficients of $p$ are combinations of the $q_b$ with integer coefficients, and
conversely.

(1) Let $x\in K_n$ with $|x|\leq1$, and write $x=p(\zeta)$ with $p$ and the
$q_b$ as in the second step. Then $x=\sum_bq_b(\zeta-1)^b$, so $(*)$ gives
$|q_b(\zeta-1)^b|\leq|x|\leq1$, and $|q_b|\leq1$ by the first step. The
coefficients $y_a$ of $p$ are combinations of the $q_b$ with integer
coefficients, so $|y_a|\leq1$, and $x=\sum_a\zeta^ay_a$.

(2) Put $p(X)=\sum_ay_aX^a$, and let the $q_b$ be as in the second step. Then
$\sum_bq_b(\zeta-1)^b=p(\zeta)=0$, so by $(*)$ every term $q_b(\zeta-1)^b$ is
$0$. As $\zeta\neq1$, every $q_b$ is $0$, so $p=0$ and every $y_a$ is $0$.
% Lean: GlobalTheory.Cyc.IntBasis: norm_term and term_le (first step, from
% norm_mem_cyclotomic_layer_prime_rpow at j = n - 1 and norm_primitiveRoot_sub_one),
% coeff_le (its last sentence), exists_poly (second step: the power basis of
% K_{n-1}(zeta) with minpoly dividing X^l - zeta^l), taylor 1 and taylor (-1) with
% norm_coeff_taylor_le_one for (1), and taylor_injective for (2).
\end{proof}

\begin{lemma}\label{asm-two-three-killing-wach-torsion-decomposition-complement}
Let $\ell$ be a prime and use the notation of
Definition~\ref{asm-two-three-killing-wach-frame}. Let
$K\subset\bar{\mathbf Q}_\ell$ be a finite extension of $\mathbf Q_\ell$, and
for a field $F\subset\bar{\mathbf Q}_\ell$ let $\mathcal O_F$ be the set of
$y\in F$ with $|y|\leq1$. For $n\geq1$ let $X^{(n)}\subset\mathbf C_\ell$ be the
additive group generated by the products $\zeta y$ with $\zeta$ a primitive
$\ell^n$-th root of unity and $y\in\mathcal O_{K_{n-1}}$. For $m\geq1$ let
$X'_m$ be the additive group generated by the groups $X^{(n)}$ with $n>m$, and
$X_m$ its closure in $\mathbf C_\ell$. Let $m\geq1$. Then:
\begin{enumerate}
\item $X^{(n)}\subset\mathcal O_{K_n}$ for every $n\geq1$, so
  $X'_m\subset\mathcal O_{K_\infty}$, and $|x|\leq1$ for every $x\in X_m$;
\item for every $z\in\mathcal O_{K_\infty}$ there is $y\in\mathcal O_{K_m}$
  with $z-y\in X'_m$;
\item $X_m\cap K_m=0$;
\item $yx\in X_m$ and $\sigma(x)\in X_m$ for $x\in X_m$, $y\in\mathcal O_{K_m}$
  and $\sigma\in G_{K_0}$.
\end{enumerate}
\end{lemma}
\leanref{Development/AbelianVarietyTheory/Proofs/InertiaEigenvalues/TateDecompositionCompletedCyclotomicIntegers.lean}{GlobalTheory.TateDecomp.decomp}
\leanref{Development/AbelianVarietyTheory/Proofs/InertiaEigenvalues/TateDecompositionCompletedCyclotomicIntegers.lean}{GlobalTheory.TateDecomp.Xj_inter}
\leanref{Development/AbelianVarietyTheory/Proofs/InertiaEigenvalues/TateDecompositionCompletedCyclotomicIntegers.lean}{GlobalTheory.TateDecomp.mul_Xj}
\leanref{Development/AbelianVarietyTheory/Proofs/InertiaEigenvalues/TateDecompositionCompletedCyclotomicIntegers.lean}{GlobalTheory.TateDecomp.act_Xj}
\begin{proof}
First, the elements of $X^{(n)}$, $n\geq1$. A generator $\zeta y$ lies in
$K_n$, as $\zeta\in K_n$ and $K_{n-1}\subset K_n$, and $|\zeta y|=|y|\leq1$, as
$|\zeta|=1$. Since $\mathcal O_{K_n}$ is a group, $X^{(n)}\subset
\mathcal O_{K_n}$. Hence $X'_m\subset\mathcal O_{K_\infty}$, and the elements of
its closure $X_m$ have absolute value at most $1$. This is (1). Now fix a
primitive $\ell^n$-th root of unity $\zeta_n$. A generator $\xi y$ of
$X^{(n)}$ has $\xi=\zeta_n^k$ with $\ell\nmid k$; write $k=q\ell+a$ with
$0<a<\ell$. Then $\xi y=\zeta_n^a(\zeta_n^\ell)^qy$, and
$(\zeta_n^\ell)^qy\in K_{n-1}$, as $\zeta_n^\ell$ is an $\ell^{n-1}$-th root of
unity. So every element of $X^{(n)}$ is $\sum_{a=1}^{\ell-1}\zeta_n^ay_a$ with
$y_a\in K_{n-1}$.

Second, $X^{(n)}\cap K_{n-1}=0$ for $n\geq2$. Let
$x=\sum_{a=1}^{\ell-1}\zeta_n^ay_a$ as in the first step, and suppose that
$x\in K_{n-1}$. Then $(-x)+\sum_{a=1}^{\ell-1}\zeta_n^ay_a=0$, and
Lemma~\ref{asm-two-three-killing-wach-torsion-decomposition-layer}(2), with the
coefficient $-x$ at $a=0$, gives $x=0$.

Third, (2). Let $n\geq2$ and $z\in\mathcal O_{K_n}$.
Lemma~\ref{asm-two-three-killing-wach-torsion-decomposition-layer}(1) writes
$z=\sum_{a=0}^{\ell-1}\zeta_n^ay_a$ with $y_a\in\mathcal O_{K_{n-1}}$. For
$0<a<\ell$, $\zeta_n^a$ is a primitive $\ell^n$-th root of unity, so
$z-y_0\in X^{(n)}$, with $y_0\in\mathcal O_{K_{n-1}}$. Now let
$z\in\mathcal O_{K_\infty}$. It lies in $K_{m+c}$ for some $c\geq0$. By
induction on $c$ there is $y\in\mathcal O_{K_m}$ with $z-y$ in the group
generated by $X^{(m+1)},\ldots,X^{(m+c)}$: for $c=0$ take $y=z$; for $c\geq1$,
the case $n=m+c\geq2$ just shown gives $y_0\in\mathcal O_{K_{m+c-1}}$ with
$z-y_0\in X^{(m+c)}$, and the induction hypothesis applies to $y_0$.

Fourth, $X'_m\cap K_m=0$. An element $x$ of $X'_m$ lies in the group generated
by $X^{(m+1)},\ldots,X^{(m+c)}$ for some $c\geq0$. We show by induction on $c$
that $x\in K_m$ forces $x=0$. For $c=0$ the group is $0$. For $c\geq1$ write
$x=x'+x''$ with $x'$ in the group generated by
$X^{(m+1)},\ldots,X^{(m+c-1)}$, which lies in $K_{m+c-1}$ by the first step,
and $x''\in X^{(m+c)}$. Then $x''=x-x'$ lies in $K_{m+c-1}$, as
$K_m\subset K_{m+c-1}$, so $x''=0$ by the second step, as $m+c\geq2$. So
$x=x'$, and $x=0$ by induction.

Fifth, let $y\in K_m$, $x\in X'_m$ and $k\geq0$ with $|y-x|\leq\ell^{-k}$;
then $|y|\leq\ell^{-k}$. Indeed, $z=(y-x)/\ell^k$ lies in
$\mathcal O_{K_\infty}$, so by (2) there are $y'\in\mathcal O_{K_m}$ and
$x'\in X'_m$ with $z=y'+x'$. Then $y-\ell^ky'=\ell^kx'+x$ lies in
$X'_m\cap K_m$, which is $0$ by the fourth step, so $|y|=|\ell^ky'|\leq\ell^{-k}$.

(3) Let $x\in X_m\cap K_m$, and suppose that $x\neq0$. Choose $k\geq0$ with
$\ell^{-k}<|x|$. As $x$ lies in the closure of $X'_m$, there is $x'\in X'_m$
with $|x-x'|\leq\ell^{-k}$, and the fifth step gives $|x|\leq\ell^{-k}$, a
contradiction.

(4) Let $f$ be multiplication by $y\in\mathcal O_{K_m}$, or
$\sigma\in G_{K_0}$. In both cases $f$ is additive and continuous on
$\mathbf C_\ell$. It suffices to show that $f$ maps each generator $\zeta w$ of
$X^{(n)}$, $n>m$, into $X^{(n)}$: then $f(X'_m)\subset X'_m$, and
$f(X_m)\subset X_m$ by continuity. For the multiplication,
$y\zeta w=\zeta(yw)$, and $yw\in\mathcal O_{K_{n-1}}$, as
$K_m\subset K_{n-1}$. For $\sigma$, $\sigma(\zeta w)=\sigma(\zeta)\sigma(w)$.
Here $\sigma(\zeta)$ is a primitive $\ell^n$-th root of unity, and
$\sigma(w)\in\mathcal O_{K_{n-1}}$: $\sigma$ fixes $K_0$ and permutes the
$\ell^{n-1}$-th roots of unity, so it maps $K_{n-1}$ into itself, and it
preserves absolute values.
% Lean: GlobalTheory.TateDecomp, with X^(n) = Xlev K n, X'_m = Xsum K m and X_m = Xj K m:
% Xlev_le, Xsum_le and Xlev_form (first step), Xlev_inter (second step, from
% integers_eq_sum_zeta_pow_unique), step and decomp with exists_level (third step),
% Xsum_inter with Xfin (fourth step), norm_le_of_close (fifth step), Xj_inter (3), and
% map_Xj with mul_Xj, act_Xj and gal_Ks (4).
\end{proof}

\begin{lemma}\label{asm-two-three-killing-wach-torsion-decomposition-operator}
Let $\ell$ be a prime and use the notation of
Definition~\ref{asm-two-three-killing-wach-frame}. Let
$K\subset\bar{\mathbf Q}_\ell$ be a finite extension of $\mathbf Q_\ell$, let
$m\geq1$ be an integer, with $m\geq2$ if $\ell=2$, and let $X_m$ be as in
Lemma~\ref{asm-two-three-killing-wach-torsion-decomposition-complement}. Let
$g\in G_{K_m}$ with $|\chi(g)-1|=\ell^{-m}$, and let $\zeta$ be a primitive
$\ell$-th root of unity. Then:
\begin{enumerate}
\item $|\zeta-1|\,|x|\leq\ell\,|g(x)-x|$ for every $x\in X_m$;
\item for every $x\in X_m$ there is $y\in X_m$ with $g(y)-y=(\zeta-1)x$.
\end{enumerate}
\end{lemma}
\leanref{Development/AbelianVarietyTheory/Proofs/InertiaEigenvalues/ExistsTateOperatorLevelComplement.lean}{GlobalTheory.Cyc.TateS.est}
\leanref{Development/AbelianVarietyTheory/Proofs/InertiaEigenvalues/ExistsTateOperatorLevelComplement.lean}{GlobalTheory.Cyc.TateS.solve}
\begin{proof}
Let $\mathcal O_F$, $X^{(n)}$ and $X'_m$ be as in
Lemma~\ref{asm-two-three-killing-wach-torsion-decomposition-complement}; the
proof uses its parts (1) to (4). The field $K_\infty$ is the union of the
$K_n$. Since $g\in G_{K_0}$, it maps every $K_n$ into itself, as in the proof
of part (4) of that lemma, and it fixes $K_m$, which contains $\zeta$, as
$m\geq1$. Also $\zeta-1\in\mathcal O_{K_m}$.

First, if $x\in X_m\cap K_\infty$ and $g(x)=x$, then $x=0$. Let $x\in K_n$.
Every power of $g$ fixes $x$. Let $\sigma\in G_{K_m}$. By
Lemma~\ref{asm-two-three-killing-wach-torsion-dense}(1) there is $a\geq0$ with
$\sigma g^{-a}\in G_{K_n}$, which fixes $x$, so
$\sigma(x)=\sigma g^{-a}(g^a(x))=x$. So $x$ is fixed by $G_{K_m}$, and by Galois
theory $x\in K_m$. By part (3) of
Lemma~\ref{asm-two-three-killing-wach-torsion-decomposition-complement},
$x=0$.

Second, let $b\geq0$, put $g'=g^{\ell^b}$, and let $z$ be a primitive
$\ell^{m+b+1}$-th root of unity. Then $g'\in G_{K_{m+b}}$, and $g'(z)=\eta z$
for a primitive $\ell$-th root of unity $\eta$. Indeed, $g'\in G_{K_0}$, and
$|\chi(g')-1|=|\chi(g)^{\ell^b}-1|=\ell^{-m-b}$ by
Lemma~\ref{asm-two-three-killing-wach-torsion-cyclotomic}(2), so
$g'\in G_{K_{m+b}}$ by part (1) of that lemma. The root of unity $z^\ell$ lies in
$K_{m+b}$, so $g'(z)^\ell=g'(z^\ell)=z^\ell$, and $\eta=g'(z)/z$ satisfies
$\eta^\ell=1$. Suppose that $\eta=1$. By the definition of the cyclotomic
character there is an integer $a\geq0$ with $g'(z)=z^a$ and
$|\chi(g')-a|\leq\ell^{-m-b-1}$. Then $z^{a-1}=1$, so $\ell^{m+b+1}$ divides
$a-1$, and $|\chi(g')-1|\leq\ell^{-m-b-1}$, a contradiction. So $\eta\neq1$,
and $\eta$ has order $\ell$.

Third, for every $x\in X'_m$ there is $y\in X_m\cap K_\infty$ with
$g(y)-y=(\zeta-1)x$. The set of such $x$ is an additive group, so it suffices
to treat a generator $x=zw$ of $X^{(n)}$ with $n>m$, where $z$ is a primitive
$\ell^n$-th root of unity and $w\in\mathcal O_{K_{n-1}}$. Write $n=m+b+1$ with
$b\geq0$, and let $g'$ and $\eta$ be as in the second step. The roots $\zeta$
and $\eta$ lie in $K_1\subset K_{n-1}$, and $|\zeta-1|=|\eta-1|$ by
Lemma~\ref{asm-two-three-killing-wach-torsion-cyclotomic}(5). So
$w'=\frac{\zeta-1}{\eta-1}w$ lies in $\mathcal O_{K_{n-1}}$, and $u=zw'$ lies in
$X^{(n)}\subset X'_m$. Put $y=\sum_{k=0}^{\ell^b-1}g^k(u)$. Each term lies in
$X_m$ by part (4) of
Lemma~\ref{asm-two-three-killing-wach-torsion-decomposition-complement}, and in
$K_n$, so $y\in X_m\cap K_\infty$. The sum telescopes:
$g(y)-y=g'(u)-u$. As $g'$ fixes $w'\in K_{m+b}$, $g'(u)=\eta zw'$, so
$g(y)-y=(\eta-1)zw'=(\zeta-1)zw=(\zeta-1)x$.

Fourth, let $d\in X'_m$ and $k\geq0$ with $|g(d)-d|\leq\ell^{-k}$; then
$|(\zeta-1)d|\leq\ell^{-k}$. Put $v=g(d)-d$. It lies in $X_m$ by part (4), and
in $K_\infty$, as $d$ does. So $v/\ell^k\in\mathcal O_{K_\infty}$, and part (2)
gives $o\in\mathcal O_{K_m}$ and $d'\in X'_m$ with $v/\ell^k=o+d'$. Then
$\ell^ko=v-\ell^kd'$ lies in $X_m\cap K_m$, so $o=0$ by part (3), and
$v=\ell^kd'$. The third step gives $y\in X_m\cap K_\infty$ with
$g(y)-y=(\zeta-1)d'$. Put $x=(\zeta-1)d-\ell^ky$. It lies in $X_m$, by part (4),
and in $K_\infty$, and since $g$ fixes $\zeta-1$,
$g(x)-x=(\zeta-1)v-\ell^k(\zeta-1)d'=0$. By the first step $x=0$, so
$|(\zeta-1)d|=\ell^{-k}|y|\leq\ell^{-k}$, by part (1).

(1) The set of $x\in\mathbf C_\ell$ with $|\zeta-1|\,|x|\leq\ell\,|g(x)-x|$ is
closed, as $g$ is continuous, and $X_m$ is the closure of $X'_m$, so it
suffices to treat $d\in X'_m$. If $g(d)=d$, the fourth step gives
$|(\zeta-1)d|\leq\ell^{-k}$ for every $k$, so $|\zeta-1|\,|d|=0$. Otherwise
$0<|g(d)-d|\leq1$, by parts (1) and (4), and there is $k\geq0$ with
$\ell^{-k-1}<|g(d)-d|\leq\ell^{-k}$. The fourth step gives
$|\zeta-1|\,|d|\leq\ell^{-k}<\ell\,|g(d)-d|$.

(2) Put $f(y)=g(y)-y$ for $y\in X_m$. By (1), applied to $y-y'\in X_m$,
$|y-y'|\leq\ell|\zeta-1|^{-1}|f(y)-f(y')|$ for $y,y'\in X_m$. So $f(X_m)$ is
closed: if $f(y_k)$ converges, then $(y_k)$ is a Cauchy sequence, it converges
to some $y$ in the closed subset $X_m$ of the complete field $\mathbf C_\ell$,
and $f(y_k)$ converges to $f(y)$, as $f$ is continuous. Hence the set $Y$ of
$x\in\mathbf C_\ell$ with $(\zeta-1)x\in f(X_m)$ is closed. By the third step
$X'_m\subset Y$, so $X_m\subset Y$, which is (2).
% Lean: GlobalTheory.Cyc.TateS, with K_infinity in C_l as algC K: fixed_alg_eq_zero (first
% step, with the density part of Cyc.gen_dense and InfiniteGalois.fixedField_fixingSubgroup),
% gen_root with the first part of gen_dense, IsGen chi K (m + b) (g^(l^b)) from
% norm_pow_pow_sub_one (second step), solve_alg (third step), int_decomp and est_pow
% (fourth step; int_decomp is the third step of the proof of Lemma
% asm-two-three-killing-wach-torsion-decomposition-complement, for elements of K_n),
% est (1), with exists_nat_pow_near_of_lt_one, and solve (2), with
% AntilipschitzWith.isClosed_range.
\end{proof}

\begin{lemma}\label{asm-two-three-killing-wach-torsion-decomposition}
Let $\ell$ be a prime and use the notation of
Definition~\ref{asm-two-three-killing-wach-frame}. Let
$K\subset\bar{\mathbf Q}_\ell$ be a finite extension of $\mathbf Q_\ell$, let
$m\geq1$ be an integer, let $\hat{\mathcal O}$ be as
in Lemma~\ref{asm-two-three-killing-wach-torsion-fixed}, and let
$\mathcal O_{K_m}$ be the set of $y\in K_m$ with $|y|\leq1$. Let
$X_m\subset\mathbf C_\ell$ be the closure of the additive group generated by the
products $\zeta y$ with $n>m$, $\zeta$ a primitive $\ell^n$-th root of unity and
$y\in K_{n-1}$ with $|y|\leq1$. Then:
\begin{enumerate}
\item every $z\in\hat{\mathcal O}$ is $y+x$ for some $y\in\mathcal O_{K_m}$ and
  $x\in X_m$;
\item if $y\in K_m$, $x\in X_m$ and $\ell^ny=x$ for some $n\geq0$, then $y=0$;
\item $X_m\subset\hat{\mathcal O}$;
\item $yx\in X_m$ for $y\in\mathcal O_{K_m}$ and $x\in X_m$;
\item $\sigma(X_m)\subset X_m$ for $\sigma\in G_{K_0}$;
\item if $x\in\hat{\mathcal O}$, $y\in\mathcal O_{K_m}$ is nonzero and
  $yx\in X_m$, then $x\in X_m$;
\item suppose that $m\geq2$ if $\ell=2$, let $g\in G_{K_m}$ with
  $|\chi(g)-1|=\ell^{-m}$, and let $\zeta$ be a primitive $\ell$-th root of
  unity. There is a map $S\colon X_m\to X_m$ with
  $S(x+x')=S(x)+S(x')$, $S(yx)=yS(x)$, $g(S(x))-S(x)=(\zeta-1)x$ and
  $S(g(x)-x)=(\zeta-1)x$ for $x,x'\in X_m$ and $y\in\mathcal O_{K_m}$.
\end{enumerate}
\end{lemma}
\leanref{Development/AbelianVarietyTheory/Definitions/Foundations/HTBasis.lean}{GlobalTheory.Cyc.Xj}
\leanref{Development/AbelianVarietyTheory/Proofs/InertiaEigenvalues/TateDecompositionCompletedCyclotomicIntegersStatement.lean}{GlobalTheory.T.tate_decomposition_completed_cyclotomic_integers}
\leanref{Development/AbelianVarietyTheory/Proofs/InertiaEigenvalues/ExistsTateOperatorLevelComplement.lean}{GlobalTheory.Cyc.exists_tate_operator}
\begin{proof}
Let $X'_m$ be as in
Lemma~\ref{asm-two-three-killing-wach-torsion-decomposition-complement}, so that
$X_m$ is its closure; we use the parts of that lemma.

(1) The set $\mathcal O_{K_m}$ is compact, as $K_m$ is finite over
$\mathbf Q_\ell$, and $X_m$ is closed, so $\mathcal O_{K_m}+X_m$ is closed in
$\mathbf C_\ell$. It contains the integers of $K_\infty$, by part (2), and hence
their closure $\hat{\mathcal O}$.

(2) Here $\ell^ny=x$ lies in $X_m\cap K_m$, which is $0$ by part (3). So $y=0$.

(3) By part (1), $X'_m$ lies in the set of integers of $K_\infty$, and $X_m$ is
its closure.

(4) and (5) are part (4).

(6) By (1) write $x=y'+x'$ with $y'\in\mathcal O_{K_m}$ and $x'\in X_m$. Then
$yy'=yx-yx'$ lies in $X_m$, by (4), and in $K_m$, so $yy'=0$ by part (3). As
$y\neq0$, $y'=0$, and $x=x'\in X_m$.

(7) By
Lemma~\ref{asm-two-three-killing-wach-torsion-decomposition-operator}(2), for
every $x\in X_m$ there is $y\in X_m$ with $g(y)-y=(\zeta-1)x$, and it is
unique: if $y,y'\in X_m$ and $g(y)-y=g(y')-y'$, then part (1) of that lemma,
applied to $y-y'$, gives $|\zeta-1|\,|y-y'|\leq0$, so $y=y'$. Let $S(x)$ be
this $y$. For $x,x'\in X_m$, $S(x)+S(x')\in X_m$ and
$g(S(x)+S(x'))-(S(x)+S(x'))=(\zeta-1)(x+x')$, so $S(x+x')=S(x)+S(x')$. For
$y\in\mathcal O_{K_m}$, $yS(x)\in X_m$ by (4), and $g$ fixes $y$, so
$g(yS(x))-yS(x)=(\zeta-1)yx$ and $S(yx)=yS(x)$. Finally, $g(x)-x\in X_m$ by (5),
$(\zeta-1)x\in X_m$ by (4), as $\zeta-1\in\mathcal O_{K_m}$, and $g$ fixes
$\zeta-1$, so $g((\zeta-1)x)-(\zeta-1)x=(\zeta-1)(g(x)-x)$ and
$S(g(x)-x)=(\zeta-1)x$.
% Lean: GlobalTheory.T.tate_decomposition_completed_cyclotomic_integers
% (= TateDecomp.tate_decomposition, for 1 <= m) gives (1) to (6) for Xj K m: Ohat_decomp
% (isCompact_intC and IsClosed.add_left_of_isCompact) for (1), Xj_inter for (2),
% closure_mono with Xsum_le for (3), mul_Xj and act_Xj for (4) and (5), and Ohat_decomp,
% mul_Xj and Xj_inter for (6). GlobalTheory.Cyc.exists_tate_operator (from
% exists_tate_operator_bound: S chosen by TateS.solve, uniqueness huniq from TateS.est)
% gives (7); the Lean also records |S(x)| <= l|x|. Proof closure of the two: 65 theorems,
% about 1,140 lines in 12 files. Part (6) is not used in this chapter.
\end{proof}

\begin{lemma}\label{asm-two-three-killing-wach-torsion-perturbation}
Let $\ell$ be a prime and use the notation of
Definition~\ref{asm-two-three-killing-wach-frame}. Let
$K\subset\bar{\mathbf Q}_\ell$ be a finite extension of $\mathbf Q_\ell$, let
$m\geq1$ be an integer, with $m\geq2$ if $\ell=2$, and let $X_m$ and
$\mathcal O_{K_m}$ be as in
Lemma~\ref{asm-two-three-killing-wach-torsion-decomposition}. Let
$g\in G_{K_m}$ with $|\chi(g)-1|=\ell^{-m}$, let $r\geq0$, let $A$ be an
$r\times r$ matrix with entries in $\mathcal O_{K_m}$ such that every entry of
$A-1$ has absolute value less than $\ell^{-1/(\ell-1)}$, let $i\in\mathbf Z$, and
put $g*y=\chi(g)^iA\,g(y)$ for $y\in\mathbf C_\ell^r$.
\begin{enumerate}
\item If $y\in X_m^r$ and $g*y=y$, then $y=0$.
\item If $\zeta$ is a primitive $\ell$-th root of unity and $w\in X_m^r$, there
  is $y\in X_m^r$ with $g*y-y=(\zeta-1)w$.
\end{enumerate}
\end{lemma}
\leanref{Development/AbelianVarietyTheory/Proofs/InertiaEigenvalues/NormZpowSubOneLeBis.lean}{GlobalTheory.CycEigen.twAct_fixed_eq_zero}
\leanref{Development/AbelianVarietyTheory/Proofs/InertiaEigenvalues/NormZpowSubOneLeTer.lean}{GlobalTheory.CycSeries.perturbation}
\begin{proof}
Let $\zeta$ be a primitive $\ell$-th root of unity, the one given in (2). It lies
in $K_m$, as $m\geq1$, and $|\zeta-1|=\rho$ with $\rho=\ell^{-1/(\ell-1)}$ by
Lemma~\ref{asm-two-three-killing-wach-torsion-cyclotomic}(5). We have
$\ell^{-m}<\rho$: for $\ell\geq3$ because $m\geq1>1/(\ell-1)$, and for $\ell=2$
because $m\geq2$. By
Lemma~\ref{asm-two-three-killing-wach-torsion-cyclotomic}(4),
$|\chi(g)^i-1|\leq|\chi(g)-1|=\ell^{-m}<\rho$, and $|\chi(g)|=1$. Put
$B=(\zeta-1)^{-1}(\chi(g)^iA-1)$. Since
$\chi(g)^iA-1=\chi(g)^i(A-1)+(\chi(g)^i-1)\cdot1$, every entry of $B$ has
absolute value less than $1$; the entries lie in $K_m$, hence in
$\mathcal O_{K_m}$. Let $c<1$ be the largest absolute value of an entry of $B$.
Since $\chi(g)^iA=1+(\zeta-1)B$,
\[
 g*y-y=g(y)-y+(\zeta-1)B\,g(y)\qquad(y\in\mathbf C_\ell^r).\tag{$*$}
\]
Let $S$ be the map of
Lemma~\ref{asm-two-three-killing-wach-torsion-decomposition}(7) for $g$ and
$\zeta$, applied to vectors coordinate by coordinate. By parts (4) and (5) of that
lemma, $g$ maps $X_m$ to itself and $X_m$ is stable under multiplication by
$\mathcal O_{K_m}$; by part (3) its elements have absolute value at most $1$.
The map $g$ fixes $\mathcal O_{K_m}$.

(1) Let $y\in X_m^r$ with $g*y=y$. By $(*)$, $g(y)-y=-(\zeta-1)B\,g(y)$. Apply
$S$: as $S$ is additive and $\mathcal O_{K_m}$-linear on $X_m$, and $\zeta-1$
and the entries of $B$ lie in $\mathcal O_{K_m}$, this gives
$(\zeta-1)y=-(\zeta-1)B\,T(y)$ with $T(x)=S(g(x))$. So $y=-B\,T(y)$. The map $T$
sends $X_m^r$ to itself, is additive, and commutes with multiplication by
matrices over $\mathcal O_{K_m}$. By induction on $n$, $y=(-B)^nz_n$ for some
$z_n\in X_m^r$: if $y=(-B)^nz_n$, then
$y=-B\,T((-B)^nz_n)=(-B)^{n+1}T(z_n)$. Since $|z_n|\leq1$, $|y|\leq c^n$ for
every $n$, so $y=0$.

(2) Let $w\in X_m^r$. If $B=0$ put $q=\ell$ and $B'=0$; otherwise let $q$ be an
entry of $B$ of largest absolute value and $B'=q^{-1}B$. Then
$q\in\mathcal O_{K_m}$, $0<|q|<1$, $B=qB'$, and $B'$ has entries in
$\mathcal O_{K_m}$. Put $v_0=w$ and $v_{n+1}=-B'g(S(v_n))$, which lie in
$X_m^r$, and $y=\sum_{n\geq0}q^nS(v_n)$. The series converges, as
$|q^nS(v_n)|\leq|q|^n$, and $y\in X_m^r$, as $X_m$ is closed. Let
$Y_N=\sum_{n<N}q^nS(v_n)$. By induction on $N$, using $g(q)=q$ and
$g(S(v_N))-S(v_N)=(\zeta-1)v_N$,
\[
 g(Y_N)-Y_N+(\zeta-1)qB'g(Y_N)=(\zeta-1)w-(\zeta-1)q^Nv_N.
\]
As $N\to\infty$ the right side tends to $(\zeta-1)w$, since $|q^Nv_N|\leq|q|^N$,
and the left side tends to $g(y)-y+(\zeta-1)B\,g(y)$, since $g$ is continuous.
By $(*)$ this is $g*y-y=(\zeta-1)w$.
\end{proof}

\begin{lemma}\label{asm-two-three-killing-wach-torsion-eigen}
Let $\ell$ be a prime and use the notation of
Definition~\ref{asm-two-three-killing-wach-frame}. Let
$K\subset\bar{\mathbf Q}_\ell$ be a finite extension of $\mathbf Q_\ell$, let
$m\geq1$ be an integer, with $m\geq2$ if $\ell=2$, and let $(U,g)$ be a frame of
rank $r$ at level $m$ for $K$ with a Hodge-Tate basis with weights in a finite
set $J\subset\mathbf Z$. Then
\[
 \prod_{u\in J}\big(U(g)-\chi(g)^{-u}\big)=0.
\]
\end{lemma}
\leanref{Development/AbelianVarietyTheory/Proofs/InertiaEigenvalues/AevalProdEqZeroHodgeTateCyclotomicFrameStatement.lean}{GlobalTheory.T.aeval_prod_eq_zero_hodgeTate_cyclotomicFrame}
\leanref{Development/AbelianVarietyTheory/Proofs/InertiaEigenvalues/AevalProdEqZeroHodgeTateCyclotomicFrame.lean}{GlobalTheory.CycEigen.aeval_prod_eq_zero}
\begin{proof}
Let $e_1,\ldots,e_r$ with weights $u_1,\ldots,u_r\in J$ be the Hodge-Tate basis.
It suffices to show $U(g)e_k=\chi(g)^{-u_k}e_k$ for every $k$: then the factor
for $u=u_k$ kills $e_k$, the factors commute, and the $e_k$ form a basis.

Fix $k$, put $u=u_k$ and $y=\ell^Ne_k$, with $N$ so large that $|y|\leq1$. Then
$\sigma*_uy=y$ for $\sigma\in G_{K_m}$, and it suffices to show
$U(g)y=\chi(g)^{-u}y$. For $h\in G_{K_\infty}$ we have $U(h)=1$ and
$\chi(h)=1$, so $h(y)=y$, and $y\in\hat{\mathcal O}^r$ by
Lemma~\ref{asm-two-three-killing-wach-torsion-fixed}. Write $y=y_0+x$ with
$y_0\in\mathcal O_{K_m}^r$ and $x\in X_m^r$
(Lemma~\ref{asm-two-three-killing-wach-torsion-decomposition}(1)). Since $g$
fixes $y_0$, $g*_uy_0-y_0=\chi(g)^uU(g)y_0-y_0$ lies in $K_m^r$. Since $g$ maps
$X_m$ to itself and the entries of $U(g)$ and $\chi(g)^u$ lie in
$\mathcal O_{K_m}$, $g*_ux\in X_m^r$ (parts (4) and (5) of that lemma). From
$g*_uy=y$ we get $g*_uy_0-y_0=x-g*_ux\in X_m^r$, so $g*_uy_0=y_0$ by part (2)
with $n=0$, and then $g*_ux=x$. The entries of $U(g)$ lie in
$\mathcal O_{K_m}$ and those of $U(g)-1$ have absolute value less than
$\ell^{-1/(\ell-1)}$, so
Lemma~\ref{asm-two-three-killing-wach-torsion-perturbation}(1) gives $x=0$.
Hence $y=y_0$ is fixed by $g$, and $y=g*_uy=\chi(g)^uU(g)y$.
\end{proof}

\begin{lemma}\label{asm-two-three-killing-wach-torsion-coinvariants}
Let $\ell$ be a prime and use the notation of
Definition~\ref{asm-two-three-killing-wach-frame}. Let
$K\subset\bar{\mathbf Q}_\ell$ be a finite extension of $\mathbf Q_\ell$, let
$m\geq1$ be an integer, with $m\geq2$ if $\ell=2$, let $(U,g)$ be a frame of
rank $r$ at level $m$ for $K$ with a Hodge-Tate basis with weights in a finite
set $J\subset\mathbf Z$, let $i\in\mathbf Z$, and let $\hat{\mathcal O}$ be as in
Lemma~\ref{asm-two-three-killing-wach-torsion-fixed}. Let
$w,w'\in\hat{\mathcal O}^r$ and $k\geq0$ satisfy $g*_iw'-w'=\ell^kw$. Then there
is $y\in\hat{\mathcal O}^r$ with $g*_iy-y=\ell^{\nu_m(J,i)}w$.
\end{lemma}
\leanref{Development/AbelianVarietyTheory/Proofs/InertiaEigenvalues/TorsionCoinvariantsTwistedCyclotomicGeneratorStatement.lean}{GlobalTheory.T.torsion_coinvariants_twistedCyclotomicGenerator}
\leanref{Development/AbelianVarietyTheory/Proofs/InertiaEigenvalues/TorsionCoinvariantsTwistedCyclotomicGenerator.lean}{GlobalTheory.CycCoinv.torsion_coinvariants}
\begin{proof}
Put $\nu=\nu_m(J,i)$, $\gamma=\chi(g)$ and $A=U(g)$, and let $X_m$ and
$\mathcal O_{K_m}$ be as in
Lemma~\ref{asm-two-three-killing-wach-torsion-decomposition}. The entries of
$A$ lie in $\mathcal O_{K_m}$, and $|\gamma|=1$, so the powers of $\gamma$ lie in
$\mathcal O_{K_m}$. For $y\in\mathcal O_{K_m}^r$, $g$ fixes $y$, so
\[
 g*_iy-y=\gamma^i\big(Ay-\gamma^{-i}y\big).\tag{$*$}
\]

First, write $w=w_0+w_x$ and $w'=v_0+v_x$ with $w_0,v_0\in\mathcal O_{K_m}^r$
and $w_x,v_x\in X_m^r$
(Lemma~\ref{asm-two-three-killing-wach-torsion-decomposition}(1)). The
hypothesis gives
$(g*_iv_0-v_0)-\ell^kw_0=\ell^kw_x-(g*_iv_x-v_x)$. The left side lies in
$\mathcal O_{K_m}^r$ by $(*)$, and the right side in $X_m^r$ by parts (4) and (5)
of that lemma, so both vanish by its part (2). With $(*)$ this gives
$Av_0-\gamma^{-i}v_0=\gamma^{-i}\ell^kw_0$.

Second, put $P(X)=\prod_{u\in J}(X-\gamma^{-u})$ and
$Q(X)=\prod_{u\in J,\,u\neq i}(X-\gamma^{-u})$, so that $P(A)=0$ by
Lemma~\ref{asm-two-three-killing-wach-torsion-eigen}. By induction on the number
of factors of $Q$, there is $y'\in\mathcal O_{K_m}^r$ with
$Q(A)w_0=Q(\gamma^{-i})w_0+(Ay'-\gamma^{-i}y')$. Indeed, if this holds for $Q$
and $u\neq i$ is a further element of $J$, then
\[
 (A-\gamma^{-u})Q(A)w_0=Q(\gamma^{-i})(\gamma^{-i}-\gamma^{-u})w_0
 +(Ay''-\gamma^{-i}y''),\qquad y''=Q(\gamma^{-i})w_0+Ay'-\gamma^{-u}y',
\]
and $y''\in\mathcal O_{K_m}^r$. The
polynomial $P$ divides $Q(X)(X-\gamma^{-i})$, so
$Q(A)(A-\gamma^{-i})=0$, and applying it to $v_0$ gives
$\gamma^{-i}\ell^kQ(A)w_0=0$, so $Q(A)w_0=0$. With $y=-y'$ and
$\beta=Q(\gamma^{-i})=\prod_{u\in J,\,u\neq i}(\gamma^{-i}-\gamma^{-u})$, this
is $Ay-\gamma^{-i}y=\beta w_0$.

Third, $\gamma^{-i}-\gamma^{-u}=\gamma^{-u}(\gamma^{u-i}-1)$, and
Lemma~\ref{asm-two-three-killing-wach-torsion-cyclotomic}(2) gives
$|\gamma^{u-i}-1|=\ell^{-m-v_\ell(u-i)}$. So
$|\beta|=\ell^{-\sum_{u\in J,u\neq i}(m+v_\ell(u-i))}\geq\ell^{-\nu}$, and
$t=\ell^\nu/\beta$ lies in $\mathbf Q_\ell$ with $|t|\leq1$. Put
$y_0=t\gamma^{-i}y\in\mathcal O_{K_m}^r$. By $(*)$,
$g*_iy_0-y_0=t(Ay-\gamma^{-i}y)=t\beta w_0=\ell^\nu w_0$.

Fourth, let $\zeta$ be a primitive $\ell$-th root of unity. It lies in $K_m$, and
$|\zeta-1|=\ell^{-1/(\ell-1)}\geq\ell^{-1}\geq|\ell^\nu|$, as $\nu\geq1$, so
$\ell^\nu/(\zeta-1)\in\mathcal O_{K_m}$ and
$(\ell^\nu/(\zeta-1))w_x\in X_m^r$.
Lemma~\ref{asm-two-three-killing-wach-torsion-perturbation}(2), with $A=U(g)$,
gives $y_x\in X_m^r$ with $g*_iy_x-y_x=\ell^\nu w_x$.

Finally, $y=y_0+y_x$ lies in $\hat{\mathcal O}^r$, since $\mathcal O_{K_m}$ and
$X_m$ lie in $\hat{\mathcal O}$ and $\hat{\mathcal O}$ is closed under addition,
and $g*_iy-y=\ell^\nu w_0+\ell^\nu w_x=\ell^\nu w$.
\end{proof}

\begin{lemma}\label{asm-two-three-killing-wach-torsion}
Let $\ell$ be a prime and use the notation of
Definition~\ref{asm-two-three-killing-wach-frame}. Let
$K\subset\bar{\mathbf Q}_\ell$ be a finite extension of $\mathbf Q_\ell$, let
$m\geq1$ be an integer, with $m\geq2$ if $\ell=2$, let $(U,g)$ be a frame of
rank $r$ at level $m$ for $K$ with a Hodge-Tate basis with weights in a finite
set $J\subset\mathbf Z$, and let $i\in\mathbf Z$. Let $c$ be a cocycle for the
twisted action $*_i$ with values in $\mathcal O_{\mathbf C}^r$, and suppose
that $\ell^kc$ is the coboundary of an element of $\mathcal O_{\mathbf C}^r$ for
some integer $k\geq0$. Then $\ell^{1+\nu_m(J,i)}c$ is the coboundary of an
element of $\mathcal O_{\mathbf C}^r$.
\end{lemma}
\leanref{Development/AbelianVarietyTheory/Proofs/InertiaEigenvalues/CocycleTwistCoboundaryOnBall.lean}{GlobalTheory.Cyc.FrameTorsion.torsion_killed}
\begin{proof}
Put $\nu=\nu_m(J,i)$ and $H=G_{K_\infty}$, a closed normal subgroup of
$G_{K_m}$, and let $\hat{\mathcal O}$ be the closure in $\mathbf C_\ell$ of the
ring of integers of $K_\infty$. On $H$ the twisted action is the action on
coordinates, since $U=1$ there and $\chi=1$ on $H$. For $\sigma\in G_{K_m}$ the
map $y\mapsto\sigma*_iy$ preserves $\mathcal O_{\mathbf C}^r$, since the
entries of $U(\sigma)$ are integral and $|\chi(\sigma)|=1$. The proof uses four
facts.
\begin{enumerate}
\item[(T1)] If $c'\colon H\to\mathcal O_{\mathbf C}^r$ is continuous and
  $c'(hh')=c'(h)+h(c'(h'))$ for $h,h'\in H$, then there is
  $x\in\mathcal O_{\mathbf C}^r$ with $\ell c'(h)=h(x)-x$ for every $h\in H$.
  This is Lemma~\ref{asm-two-three-killing-wach-torsion-tate} with $a=\ell$.
\item[(T2)] An element of $\mathcal O_{\mathbf C}$ is fixed by $H$ if and only
  if it lies in $\hat{\mathcal O}$
  (Lemma~\ref{asm-two-three-killing-wach-torsion-fixed}).
\item[(T3)] If $w,w'\in\hat{\mathcal O}^r$ and $g*_iw'-w'=\ell^kw$ for some
  $k\geq0$, then there is $y\in\hat{\mathcal O}^r$ with $g*_iy-y=\ell^\nu w$
  (Lemma~\ref{asm-two-three-killing-wach-torsion-coinvariants}).
\item[(T4)] A cocycle for $*_i$ that vanishes on $H$ and at $g$ vanishes on
  $G_{K_m}$ (Lemma~\ref{asm-two-three-killing-wach-torsion-procyclic}).
\end{enumerate}
% Lean: (T1) is GlobalTheory.T.cyclotomic_tate_maximalIdeal_kills_first_cohomology
% (= Cyc.AveragingRankOne.cyclotomic_maximalIdeal_smul_cocycle_isCoboundary, applied with
% a = l), (T2) is GlobalTheory.T.cyclotomicTower_fixed_iff_mem_completedIntegers
% (= CycInvariants.fixed_iff_mem_Ohat, Ohat K the closure of intC (Kinf K)), (T3) is
% GlobalTheory.T.torsion_coinvariants_twistedCyclotomicGenerator (= CycCoinv.torsion_coinvariants)
% and (T4) is the first part of
% GlobalTheory.T.cocycle_vanishing_on_fixingSubgroup_determined_by_generator (= Cyc.Procyclic.procyclic).
% Tate's decomposition and operator are Lemma asm-two-three-killing-wach-torsion-decomposition
% and the lemmas before it. Proof closure of torsion_killed: 246 theorems, about 3,870 lines
% in 37 files.

First, restricted to $H$, $c$ satisfies the hypothesis of (T1), which gives
$x_1\in\mathcal O_{\mathbf C}^r$ with $\ell c(h)=h(x_1)-x_1$ for $h\in H$. The
map $c_1(\sigma)=\ell c(\sigma)-(\sigma*_ix_1-x_1)$ is a cocycle for $*_i$ with
values in $\mathcal O_{\mathbf C}^r$, and it vanishes on $H$.

Second, $c_1(g)\in\hat{\mathcal O}^r$. Let $h\in H$ and put
$h'=g^{-1}hg\in H$. Then $c_1(gh')=c_1(g)+g*_ic_1(h')=c_1(g)$, and also
$c_1(gh')=c_1(hg)=c_1(h)+h*_ic_1(g)=h(c_1(g))$. So $c_1(g)$ is fixed by $H$, and
(T2) applies to its coordinates.

Third, let $x_0\in\mathcal O_{\mathbf C}^r$ with $\ell^kc(\sigma)=\sigma*_ix_0-x_0$
for every $\sigma$, and put $x_2=\ell x_0-\ell^kx_1\in\mathcal O_{\mathbf C}^r$.
Then $\ell^kc_1(\sigma)=\sigma*_ix_2-x_2$ for every $\sigma\in G_{K_m}$. At
$\sigma\in H$ the left side vanishes, so $x_2$ is fixed by $H$ and
$x_2\in\hat{\mathcal O}^r$ by (T2).

Fourth, (T3) with $w=c_1(g)$ and $w'=x_2$ gives $y\in\hat{\mathcal O}^r$ with
$g*_iy-y=\ell^\nu c_1(g)$. By (T2), $y$ is fixed by $H$. So the cocycle
$d(\sigma)=\ell^\nu c_1(\sigma)-(\sigma*_iy-y)$ vanishes on $H$ and at $g$, and
by (T4) it vanishes on $G_{K_m}$.

Finally, for $\sigma\in G_{K_m}$,
\[
 \ell^{1+\nu}c(\sigma)=\ell^\nu c_1(\sigma)+\ell^\nu(\sigma*_ix_1-x_1)
 =\sigma*_iz-z,\qquad z=y+\ell^\nu x_1\in\mathcal O_{\mathbf C}^r.
\]
\end{proof}

\begin{lemma}\label{asm-two-three-killing-wach-twisted}
Let $\ell$ be a prime and use the notation of
Definition~\ref{asm-two-three-killing-wach-frame}. Let
$K\subset\bar{\mathbf Q}_\ell$ be a finite extension of $\mathbf Q_\ell$, let
$m\geq1$ be an integer, with $m\geq2$ if $\ell=2$, and let $(U,g)$ be a frame of
rank $r$ at level $m$ for $K$ with a Hodge-Tate basis with weights in a finite
set $J\subset\mathbf Z$. Let $i\in\mathbf Z$ with $i\notin J$, and let
$\Lambda\subset\mathbf C_\ell^r$ be an additive subgroup with
$\mathcal O_{\mathbf C}^r\subset\Lambda$ and
$\ell\Lambda\subset\mathcal O_{\mathbf C}^r$. If $y\in\mathbf C_\ell^r$ satisfies
$\sigma*_iy-y\in\Lambda$ for every $\sigma\in G_{K_m}$, then
$\ell^{2+\nu_m(J,i)}y\in\Lambda$.
\end{lemma}
\leanref{Development/AbelianVarietyTheory/Proofs/InertiaEigenvalues/PrimePowSmulMemOfTwistedInvariantModLatticeStatement.lean}{GlobalTheory.T.prime_pow_smul_mem_of_twistedInvariant_mod_lattice}
\leanref{Development/AbelianVarietyTheory/Proofs/InertiaEigenvalues/PrimePowSmulMemOfTwistedInvariantModLattice.lean}{GlobalTheory.FrameInv.pow_smul_mem_of_invariant_mod}
\begin{proof}
Put $\nu=\nu_m(J,i)$. First, $\sigma*_i(ay)=\sigma(a)\,(\sigma*_iy)$ for
$a\in\mathbf C_\ell$, and $\sigma*_i(\sigma'*_iy)=(\sigma\sigma')*_iy$ for
$\sigma,\sigma'\in G_{K_m}$, since $\sigma$ fixes the entries of $U(\sigma')$,
which lie in $K_m$, and the number $\chi(\sigma')^i\in\mathbf Q_\ell$. The map
$\sigma\mapsto\sigma*_iy$ is continuous on $G_{K_m}$, since $U$, $\chi$ and the
orbit maps of $G_{\mathbf Q_\ell}$ on $\mathbf C_\ell$ are.

Second, $c(\sigma)=\ell(\sigma*_iy-y)$ is a cocycle for $*_i$ with values in
$\ell\Lambda\subset\mathcal O_{\mathbf C}^r$. Choose $n\geq0$ with
$\ell^{n+1}y\in\mathcal O_{\mathbf C}^r$. Since $\sigma$ fixes $\ell^{n+1}$,
$\ell^nc(\sigma)=\sigma*_i(\ell^{n+1}y)-\ell^{n+1}y$, so $\ell^nc$ is the
coboundary of an element of $\mathcal O_{\mathbf C}^r$.
Lemma~\ref{asm-two-three-killing-wach-torsion} gives
$z\in\mathcal O_{\mathbf C}^r$ with $\ell^{1+\nu}c(\sigma)=\sigma*_iz-z$ for every
$\sigma\in G_{K_m}$.

Third, $w=\ell^{2+\nu}y-z$ satisfies $\sigma*_iw=w$ for every
$\sigma\in G_{K_m}$, since $\sigma*_iw-w=\ell^{1+\nu}c(\sigma)-(\sigma*_iz-z)$.
Let $e_1,\ldots,e_r$ with weights $u_1,\ldots,u_r$ be the Hodge-Tate basis, and
write $w=\sum_ka_ke_k$ with $a_k\in\mathbf C_\ell$. Since
$\sigma*_iy'=\chi(\sigma)^{i-u_k}(\sigma*_{u_k}y')$ for every $y'$,
\[
 \sigma*_iw=\sum_k\chi(\sigma)^{i-u_k}\sigma(a_k)\,e_k,
\]
so $\chi(\sigma)^{i-u_k}\sigma(a_k)=a_k$ for every $k$ and every
$\sigma\in G_{K_m}$. As $u_k\in J$ and $i\notin J$, $i-u_k\neq0$, and
Theorem~\ref{asm-two-three-killing-wach-tate-sen} gives $a_k=0$. Hence $w=0$, and
$\ell^{2+\nu}y=z\in\mathcal O_{\mathbf C}^r\subset\Lambda$.
\end{proof}

\subsection{Wach modules}\label{subsection-asm-wach-modules}

\begin{definition}\label{asm-two-three-killing-wach-frame-exists-data}
Let $\ell$ be a prime and use the notation of
Definition~\ref{asm-two-three-killing-wach-frame}. Let $s\geq1$ be an integer,
with $s\geq2$ if $\ell=2$, let $K\subset\bar{\mathbf Q}_\ell$ be a finite
extension of $\mathbf Q_\ell$ with $K\subset K_s$, and let
$L\subset\bar{\mathbf Q}_\ell$ be a finite extension of $\mathbf Q_\ell$ that
contains $\tau(K_s)$ for every $\mathbf Q_\ell$-embedding
$\tau\colon K_s\to\bar{\mathbf Q}_\ell$. Put $\mathcal O=\mathcal O_L$ and
$n=[K_0:\mathbf Q_\ell]$. Lattice representations $(V,T)$ of $G_K$ over $L$ are
as in Definition~\ref{asm-two-three-killing-holds-lift-block}; a morphism
$(V,T)\to(V',T')$ is a continuous $G_K$-equivariant $L$-linear map
$f\colon V\to V'$ with $f(T)\subset T'$. \emph{Wach data} for $(K,s,L)$
consist of the following.
\begin{enumerate}
\item The ring $A=\prod_{i\in\mathbf Z/n}\mathcal O[[\pi]]$, with its
  idempotents $e_i$ and the element $\pi$ all of whose components are $\pi$; an
  $\mathcal O$-algebra endomorphism $\varphi_0$ of $\mathcal O[[\pi]]$ and
  $q_0\in\mathcal O[[\pi]]$ with $\pi q_0=(1+\pi)^\ell-1$ and
  $\varphi_0(\pi)=\pi q_0$. The endomorphism $\varphi$ of $A$ is
  $(\varphi a)_i=\varphi_0(a_{i-1})$, the element $q\in A$ has all components
  $q_0$, and $Q=\varphi^{s-1}(q)$.
\item A commutative group $\Gamma$ with a surjective homomorphism
  $G_K\to\Gamma$, $g\mapsto\bar g$; an injective character
  $\chi_\Gamma\colon\Gamma\to\mathcal O^\times$ with $\chi_\Gamma(\bar g)=\chi(g)$
  for $g\in G_K$; and an action of $\Gamma$ on $A$ by $\mathcal O$-algebra
  endomorphisms that commutes with $\varphi$, in which $\gamma\in\Gamma$ acts on
  every factor $\mathcal O[[\pi]]$ by the same endomorphism $\gamma_0$, and
  $\gamma_0(\pi)$ is a power series with constant term $0$ and coefficient of
  $\pi$ equal to $\chi_\Gamma(\gamma)$.
\item An $A$-algebra $E$ into which $A$ embeds and in which $\pi$ is a unit,
  with extensions of $\varphi$ and of the action of $\Gamma$ to $E$.
\item For each lattice representation $(V,T)$: a free $E$-module $D(T)$ of rank
  the rank of $T$ over $\mathcal O$, with a $\varphi$-semilinear endomorphism
  $\varphi_D$ and a semilinear action of $\Gamma$ that commutes with
  $\varphi_D$; an $A$-submodule $N(T)$ of $D(T)$, the \emph{Wach module} of
  $T$, finitely generated over $A$ and projective over $\mathbf Z_\ell[[\pi]]$,
  acting through $A$, such that $E\otimes_AN(T)\to D(T)$ is an isomorphism when
  $V$ is crystalline, with $\gamma x\in N(T)$ and $\gamma x-x\in\pi N(T)$ for $\gamma\in\Gamma$
  and $x\in N(T)$, and such that for every $x\in N(T)$ there is $k\geq0$ with
  $Q^k\varphi_D(x)\in N(T)$ and $Q^kx$ in the $A$-submodule generated by
  $\varphi_D(N(T))$; and for each morphism $f\colon(V,T)\to(V',T')$ an $E$-linear
  map $D(f)\colon D(T)\to D(T')$ that commutes with $\varphi_D$ and $\Gamma$, is
  functorial, and maps $N(T)$ into $N(T')$ when $V$ and $V'$ are crystalline.
\item A ring $R$ with an endomorphism $\varphi_R$ and an action of $G_K$, a ring
  map $E\to R$, and for each $(V,T)$ an injective map
  $j_T\colon D(T)\to R\otimes_{\mathcal O}T$, semilinear for $E\to R$, which
  carries $\varphi_D$ to $\varphi_R\otimes1$ and the action of $\bar g$ to that
  of $g\otimes g$ for $g\in G_K$, and whose image generates
  $R\otimes_{\mathcal O}T$ over $R$.
\item A lattice representation $(V_1,T_1)$ on which $G_K$ acts trivially, with
  $T_1=\mathcal Oe_0$ free of rank one, and an element $e\in D(T_1)$ with
  $j_{T_1}(e)=1\otimes e_0$; a generator $\gamma_0$ of $\Gamma$ with
  $v(\chi_\Gamma(\gamma_0)-1)=s$, where $v$ is the valuation of $L$ with
  $v(\ell)=1$; and further data which this chapter does not use:
  the subgroup of torsion elements of $\Gamma$, completeness and density properties of
  $E$, comparisons of $N(T)$ with $D_{\mathrm{cris}}(V)$ and with
  $V\otimes B_{\mathrm{dR}}$ beyond the expansions of
  Definition~\ref{asm-two-three-killing-wach-frame-exists-inputs-heights-expansions},
  and cohomology classes attached to pairs in
  $D(T)$ that are cocycles for $\varphi_D$ and $\Gamma$.
\end{enumerate}
The local Hodge-Tate data and Fontaine data of a definite-side datum at $\ell$
(Definition~\ref{asm-local-bound-two-three}(3)) give Wach data for every such
triple $(K,s,L)$, with $R=\mathcal O\otimes_{\mathbf Z_\ell}W(\mathbf C_\ell^\flat)$,
$\varphi_R$ the Frobenius of the second factor and $G_K$ acting on the second
factor: the \emph{Wach data of the datum}. They come with an $\mathcal O$-algebra
map $\iota_A$ from $A$ to the subring
$R_A=\mathcal O\otimes_{\mathbf Z_\ell}A_{\mathrm{inf}}$ of $R$,
$A_{\mathrm{inf}}=W(\mathcal O_{\mathbf C}^\flat)$, such that the map $E\to R$
restricted to $A$ is $\iota_A$, and $\iota_A(\pi)=1\otimes([\varepsilon^{(s)}]-1)$.
Here $(\zeta_{\ell^j})_{j\geq0}$ is a fixed system of primitive $\ell^j$-th
roots of unity in $\mathcal O_{\mathbf C}$ with $\zeta_{\ell^{j+1}}^\ell=
\zeta_{\ell^j}$, $\varepsilon\in\mathcal O_{\mathbf C}^\flat$ is the element it
defines, $\varepsilon^{(j)}$ is the $\ell^j$-th root of $\varepsilon$, and
$[\,\cdot\,]$ is the Teichmüller lift.
\end{definition}
\leanref{Formalism/SharpRestriction/LevelLetterCocycles/DataNLoc.lean}{Formalism.BasicNotions.LocWach}
\leanref{Formalism/SharpRestriction/LevelLetterCocycles/DataNLoc.lean}{Formalism.BasicNotions.EtTie}
\leanref{Formalism/SharpRestriction/LevelLetterCocycles/DataNLoc.lean}{Formalism.BasicNotions.LocHW}

\begin{definition}\label{asm-two-three-killing-wach-frame-exists-clauses}
Let $\ell$, $s$, $K$ and $L$ be as in
Definition~\ref{asm-two-three-killing-wach-frame-exists-data}, with Wach data
for $(K,s,L)$, let $(V,T)$ be a lattice representation of $G_K$ over $L$, and
let $M_0,h_0\geq0$ be integers. An $A$-submodule $N$ of $D(T)$ satisfies the
\emph{Wach clauses} for $M_0$ and $h_0$ if:
\begin{enumerate}
\item[(a)] $N$ has an $A$-basis whose image under $j_T$ is an $R$-basis of
  $R\otimes_{\mathcal O}T$;
\item[(b)] $Q^{M_0}\varphi_D(x)\in N$ for every $x\in N$;
\item[(c)] $Q^{h_0}y$ lies in the $A$-submodule generated by $\varphi_D(N)$ for
  every $y\in N$;
\item[(d)] $\bar gx\in N$ and $\bar gx-x\in\pi N$ for every $x\in N$ and
  $g\in G_{K_s}$;
\item[(e)] $D(f)(x)\in N$ for every $x\in N$ and every morphism
  $f\colon(V,T)\to(V,T)$.
\end{enumerate}
The \emph{Wach hypothesis} holds for $(V,T)$, $M_0$ and $h_0$ if $V$ is
crystalline as a representation of $G_{K_s}$, every integer $a$ with
$\operatorname{Fil}^{a+1}\neq\operatorname{Fil}^a$ on $D_{\mathrm{dR}}(V)$ lies
in $[-M_0,h_0]$, and some $A$-submodule of $D(T)$ satisfies the Wach clauses for
$M_0$ and $h_0$.
\end{definition}
\leanref{Formalism/SharpRestriction/LevelLetterCocycles/LocalSituation.lean}{Formalism.MainNotions.LocalSituation.WClauses}
\leanref{Formalism/SharpRestriction/LevelLetterCocycles/LocalSituation.lean}{Formalism.MainNotions.LocalSituation.HypW}
\leanref{Formalism/SharpRestriction/LevelLetterCocycles/LocalSituation.lean}{Formalism.MainNotions.LocalSituation.IsCrysOn}

\begin{lemma}\label{asm-two-three-killing-wach-frame-exists-inputs-basis}
Let $\ell$, $s$, $K$ and $L$ be as in
Definition~\ref{asm-two-three-killing-wach-frame-exists-data}, let $D$ be a
definite-side datum with local Hodge-Tate data, Fontaine data and comparison
data at every prime (Definition~\ref{asm-local-bound-two-three}(3)), with its
Wach data for $(K,s,L)$, and let $(V,T)$ be a lattice representation of $G_K$
over $L$ with $V$ crystalline. Then $N(T)$ has an $A$-basis, and every
$A$-basis of $N(T)$ is an $E$-basis of $D(T)$.
\end{lemma}
\leanref{Development/AbelianVarietyTheory/Proofs/InertiaEigenvalues/NBasisLoc.lean}{GlobalTheory.NL.N_basis_loc}
\leanref{Development/AbelianVarietyTheory/Proofs/InertiaEigenvalues/MemXSmulOfProjective.lean}{GlobalTheory.WachFree.basis_of_regular}
\leanref{Development/AbelianVarietyTheory/Proofs/InertiaEigenvalues/MemXSmulOfProjective.lean}{GlobalTheory.WachFree.free_of_regular}
\leanref{Development/AbelianVarietyTheory/Proofs/InertiaEigenvalues/MemXSmulOfProjective.lean}{GlobalTheory.WachFree.basis_of_baseChange}
\begin{proof}
Put $N=N(T)$, and let $\mathcal O[[\pi]]$ act on $N$ through the diagonal map
to $A$; then $N$ is finitely generated over $\mathcal O[[\pi]]$, as $A$ is. As
$V$ is crystalline, $E\otimes_AN\to D(T)$ is an isomorphism
(Definition~\ref{asm-two-three-killing-wach-frame-exists-data}(4)), so every
$A$-basis of $N$ is an $E$-basis of $D(T)$.

First, multiplication by $\pi$ is injective on $N$, since $N\subset D(T)$ and
$\pi$ is a unit in $E$.

Second, if $x,y\in N$ and $\ell x=\pi y$, then $x\in\pi N$. Indeed $N$ is a
direct summand of a free $\mathbf Z_\ell[[\pi]]$-module $F$, with projection
$p\colon F\to N$. In $F$ every coordinate $f$ of $x$ satisfies $\ell f\in\pi
\mathbf Z_\ell[[\pi]]$, so $f\in\pi\mathbf Z_\ell[[\pi]]$; thus $x=\pi w$ with
$w\in F$, and $x=p(x)=\pi p(w)$.

Third, $N$ is free over $\mathcal O[[\pi]]$. By the second step applied to
$\ell^k$, multiplication by every nonzero element of $\mathcal O$, which divides
a power of $\ell$, is injective on $N/\pi N$, a finitely generated module over
$\mathcal O[[\pi]]/(\pi)=\mathcal O$. So $N/\pi N$ is free over $\mathcal O$, and
as $\pi$ is regular on $N$ by the first step and lies in the maximal ideal of
the local ring $\mathcal O[[\pi]]$, $N$ is free.

Finally, $A$ is the product of $n$ copies of $\mathcal O[[\pi]]$, so $N$ is a
finitely generated flat $A$-module, and $e_iN$ is free over the $i$-th factor.
Its rank is the rank $r$ of $D(T)$ over $E$: the ring $E/(1-e_i)E$ is nonzero,
as $A$ embeds in $E$ and $e_i\neq0$, and the base change of $e_iN$ to it is the
base change of $D(T)$, free of rank $r$. Choosing bases $(b_{ij})_{j\leq r}$ of
the $e_iN$, the elements $\sum_ib_{ij}$ form an $A$-basis of $N$.
% Lean: N_basis_loc: W_proj (projectivity over Z_l[[X]] through zlA), N_fg, N_base, EE_π and
% the DVR O_L with c = l; WachFree.basis_of_regular = free_of_regular (via
% mem_smul_of_projective_comp = mem_X_smul_of_projective for the second step) and
% basis_of_baseChange (finrank_fibre at the maximal ideals of A). 259 lines.
\end{proof}

\begin{definition}\label{asm-two-three-killing-wach-frame-exists-inputs-heights-expansions}
Let $\ell$, $s$, $K$, $L$ and $D$ be as in
Lemma~\ref{asm-two-three-killing-wach-frame-exists-inputs-basis}, with the Wach
data of $D$ for $(K,s,L)$, $n=[K_0:\mathbf Q_\ell]$ and $\mathcal O=\mathcal O_L$.
For an $L$-vector space $M$, a sequence $f\colon\mathbf Z\to M$ that vanishes
below some integer, and $\eta=\sum_t\eta_tX^t\in L[[X]]$, put
$(\eta*f)(u)=\sum_{t\geq0}\eta_tf(u-t)$, a finite sum. For a lattice representation
$(V,T)$ of $G_K$ over $L$, let
$\iota\colon V\otimes_{\mathbf Q_\ell}B_{\mathrm{cris}}\to V\otimes_{\mathbf
Q_\ell}B_{\mathrm{dR}}$ be the inclusion and $\varphi$ the Frobenius of
$V\otimes B_{\mathrm{cris}}$, acting on the second factor. The local Hodge-Tate
data and comparison data of $D$ give, with these Wach data:
\begin{enumerate}
\item a bijection $i\mapsto\tau_i$ from $\mathbf Z/n$ to the set of
  $\mathbf Q_\ell$-embeddings $K_0\to L$, with $\tau_{i+1}\circ\sigma=\tau_i$
  for every $i$, where $\sigma$ is the arithmetic Frobenius of $K_0$, the
  restriction to $K_0\subset B_{\mathrm{cris}}$ of the Frobenius of
  $B_{\mathrm{cris}}$. For $(V,T)$ and $i\in\mathbf Z/n$,
  $D_i(V)$ is the subspace of $D_{\mathrm{cris},K}(V)=(V\otimes
  B_{\mathrm{cris}})^{G_K}$ on which every $c\in K_0\subset B_{\mathrm{cris}}$
  acts as $\tau_i(c)\in L$ acting on $V$;
\item for $i\in\mathbf Z/n$ and $\zeta\in L$, a map of $\mathcal O$-algebras
  $\lambda_{i,\zeta}\colon A\to L[[X]]$, the \emph{expansion} at $(i,\zeta)$,
  with $\lambda_{i,\zeta}(a)=0$ when the $i$-th component of $a$ is $0$, and,
  when $\zeta$ is a root of unity of $\ell$-power order,
  $1+\lambda_{i,\zeta}(\pi)=\zeta\exp(X)$;
\item for $(V,T)$, $i$ and $\zeta$, an additive map $\varepsilon_{i,\zeta}$ from
  $N(T)$ to the sequences $\mathbf Z\to V\otimes B_{\mathrm{cris}}$, whose
  values vanish below some integer and lie in $D_i(V)$, such that
  $\iota\circ\varepsilon_{i,\zeta}(ax)=\lambda_{i,\zeta}(a)*\big(\iota\circ
  \varepsilon_{i,\zeta}(x)\big)$ for $a\in A$ and $x\in N(T)$;
\item for $(V,T)$ with $V$ crystalline, $i$, a root of unity $\zeta\in L$ of
  $\ell$-power order, $x\in N(T)$ and $m\geq0$ with $Q^m\varphi_D(x)\in N(T)$:
  if $\lambda_{i,\zeta}(Q^m)=\sum_t\mu_tX^t$, then for every $u\in\mathbf Z$
  \[
   \varepsilon_{i,\zeta}\big(Q^m\varphi_D(x)\big)(u)=\sum_{t\geq0}\mu_t\,
   \ell^{u-t}\varphi\big(\varepsilon_{i-1,\zeta^\ell}(x)(u-t)\big).
  \]
\item for $(V,T)$ with $V$ crystalline, $\gamma\in\Gamma$, $i$, a primitive
  $\ell^j$-th root of unity $\zeta\in L$ for some $j\geq0$, an integer
  $a\geq0$ with $v(\chi_\Gamma(\gamma)-a)\geq j$ ($v$ as in
  Definition~\ref{asm-two-three-killing-wach-frame-exists-data}(6)),
  $x\in N(T)$ and $u\in\mathbf Z$:
  $\varepsilon_{i,\zeta}(\gamma x)(u)=\chi_\Gamma(\gamma)^u\,
  \varepsilon_{i,\zeta^a}(x)(u)$.
\item for $(V,T)$, an additive map $\beta_T\colon R_A\otimes_{\mathcal O}T\to
  V\otimes_{\mathbf Q_\ell}B_{\mathrm{dR}}$, with $R_A\subset R$ as in
  Definition~\ref{asm-two-three-killing-wach-frame-exists-data}, such that
  $\beta_T((o\otimes w)\otimes t)=ot\otimes w$ for $o\in\mathcal O$, $w\in
  A_{\mathrm{inf}}\subset B_{\mathrm{dR}}$ and $t\in T$; and, for $(V,T)$ with
  $V$ crystalline, $i$, $\zeta\in L$, a map of $\mathbf Q_\ell$-algebras
  $\kappa\colon L\to B_{\mathrm{dR}}$ with $\kappa(\tau_i(c))=c$ for $c\in K_0$
  and $\kappa(\zeta)=\zeta_{\ell^s}$, $x\in N(T)$, $k\geq0$, an element $y$ of
  $R_A\otimes_{\mathcal O}T$ whose image in $R\otimes_{\mathcal O}T$ is
  $(1\otimes p_1)^kj_T(x)$, where $p_1=[\varepsilon^{(1)}]-1\in A_{\mathrm{inf}}$,
  and $a\geq0$ with $\varepsilon_{i,\zeta}(x)(r)=0$ for $r<-a$: there is $e\in\mathbf
  Z$ such that for every $M\geq0$
  \[
   t_s^a\beta_T(y)-p_1^k\sum_{0\leq r<M}t_s^r\,\iota\big(\varepsilon_{i,\zeta}(x)
   (r-a)\big)\in V\otimes\operatorname{Fil}^{M-e}B_{\mathrm{dR}}+K_\kappa.
  \]
  Here $t_s=\ell^{-s}t$, with $t\in\operatorname{Fil}^1B_{\mathrm{dR}}$
  Fontaine's element $\log[\varepsilon]$; elements of $B_{\mathrm{dR}}$ act on
  $V\otimes B_{\mathrm{dR}}$ through the second factor; and $K_\kappa$ is the
  $L$-span of the elements $cz-\kappa(c)z$ of $V\otimes B_{\mathrm{dR}}$ with
  $c\in L$, where $c$ acts on $V$ and $\kappa(c)$ on $B_{\mathrm{dR}}$.
\item $A_{\mathrm{inf}}\subset\operatorname{Fil}^0B_{\mathrm{dR}}$, and the
  element $p_1$ of (6) is a unit of $\operatorname{Fil}^0B_{\mathrm{dR}}$: there
  is $u\in\operatorname{Fil}^0B_{\mathrm{dR}}$ with $p_1u=1$. The first assertion
  is part of the local data of $D$ at $\ell$, the second part of its comparison
  data. For Fontaine's ring both hold. Indeed
  $\operatorname{Fil}^0B_{\mathrm{dR}}=B^+_{\mathrm{dR}}$, which contains
  $A_{\mathrm{inf}}=W(\mathcal O_{\mathbf C}^\flat)$
  \cite[(4.4.1), Corollary~4.4.5 and Definition~4.4.7]{BC09}. It is a complete
  discrete valuation ring whose maximal ideal is the kernel of Fontaine's map
  $\theta\colon B^+_{\mathrm{dR}}\to\mathbf C_\ell$, so that an element is a unit
  if and only if its image under $\theta$ is nonzero
  \cite[Proposition~4.4.6 and its proof]{BC09}. On Teichmüller lifts,
  $\theta([x])=\lim_n\hat x_n^{\ell^n}$ for $x=(x_n)_n$ and any lifts $\hat x_n\in
  \mathcal O_{\mathbf C}$ of the $x_n$ \cite[\S4.4, the definition of $\theta$
  and Lemma~4.4.1]{BC09}. For $x=\varepsilon^{(1)}$ one may take $\hat
  x_n=\zeta_{\ell^{n+1}}$, so $\theta(p_1)=\zeta_\ell-1\neq0$.
\end{enumerate}
For such $(V,T)$, $i$ and $\zeta$, $\Lambda_{i,\zeta}(T)$ is the $L$-subspace of
the sequences $\mathbf Z\to V\otimes B_{\mathrm{cris}}$ spanned by the
$\eta*\varepsilon_{i,\zeta}(x)$ with $\eta\in L[[X]]$ and $x\in N(T)$. For a
$\mathbf Q_\ell$-embedding $\tau\colon K\to L$ let $D_\tau$ be the $\tau$-part of
$D_{\mathrm{dR}}(V)$, the $z$ with $cz=\tau(c)z$ for $c\in K$, where $c$ acts on
$B_{\mathrm{dR}}$ and $\tau(c)$ on $V$; for $a\in\mathbf Z$, $F^a_\tau(V)\subset
V\otimes B_{\mathrm{dR}}$ is the sum of $D_\tau\cap(V\otimes\operatorname{Fil}^a
B_{\mathrm{dR}})$ and the $D_{\tau'}$ with $\tau'\neq\tau$.
\end{definition}
\leanref{Formalism/SharpRestriction/LevelLetterCocycles/DataNLoc.lean}{Formalism.BasicNotions.CmpTie}
\leanref{Development/AbelianVarietyTheory/Definitions/NonGaloisDescent/NotFullBot.lean}{GlobalTheory.NL.conv}
\leanref{Development/AbelianVarietyTheory/Definitions/NonGaloisDescent/NotFullBot.lean}{GlobalTheory.NL.Lat}
\leanref{Development/AbelianVarietyTheory/Definitions/NonGaloisDescent/NotFullBot.lean}{GlobalTheory.NL.FilAt}
% Lean: (1) is LocWach.emb with emb_bij, emb_σ (σ is PD.σ, with PD.φ_K₀: φ = σ on K_0) and
% DcrisL PD rho (emb i) (GlobalTheory.NL.Di); (2) is
% LocWach.lev with lev_letter and lev_π (the Lean's E with E(0) = 1 and E' = E is exp(X));
% (3) is CmpTie.levB (= GlobalTheory.NL.ex) with
% levB_mem, levB_supp and levN_eq, together with LocWach.levN_smul for levN = ιV ∘ levB;
% (4) is CmpTie.levB_φ (PD.p = l); (5) is CmpTie.levB_Γ (v is P.v LL, with LocWach.v_l).
% (6) is CmpTie.βdR with βdR_tmul and levN_dR (written with levN = ιV ∘ levB by levN_eq),
% for CmpSit.tie: ιdR = ιB ∘ ιA (A_inf in B_dR), p1 = AI.π₁ = [ε1] - 1, tS = l^(-s) ιB(tB)
% and zS = ζ_{l^s}; K_κ is compKer T κ and V ⊗ Fil^a is FilV a.
% Λ is Lat C h T i ζ = latSpan (ex C h T i ζ), and F_τ^a is
% FilAt PD T τ a = FildR τ a ⊔ the DdRE τ' for τ' ≠ τ.

\begin{lemma}\label{asm-two-three-killing-wach-frame-exists-inputs-heights-points}
Let $\ell$, $s$, $K$, $L$ and $D$ be as in
Definition~\ref{asm-two-three-killing-wach-frame-exists-inputs-heights-expansions},
with its notation, let $i\in\mathbf Z/n$, and let $\zeta\in L$ be a primitive
$\ell^s$-th root of unity. Then $\lambda_{i,\zeta}(Q)=Xw$ for a unit $w$ of
$L[[X]]$.
\end{lemma}
\leanref{Development/AbelianVarietyTheory/Proofs/InertiaEigenvalues/WachLevelExpansionAtRootOfUnity.lean}{GlobalTheory.NL.lev_at_root_of_unity}
\leanref{Development/AbelianVarietyTheory/Proofs/InertiaEigenvalues/WachLevelExpansionAtRootOfUnity.lean}{GlobalTheory.NL.LevPt.lev_values}
\begin{proof}
Write $\lambda=\lambda_{i,\zeta}$ and $U=\lambda(1+\pi)=\zeta\exp(X)$
(Definition~\ref{asm-two-three-killing-wach-frame-exists-inputs-heights-expansions}(2)),
so that $U$ has constant term $\zeta$ and $U'=U$. In $A$,
$\varphi(1+\pi)=1+\pi q=(1+\pi)^\ell$
(Definition~\ref{asm-two-three-killing-wach-frame-exists-data}(1)), so
$X_s=\varphi^{s-1}(\pi)$ is $(1+\pi)^{\ell^{s-1}}-1$, and
$X_sQ=\varphi^{s-1}(\pi q)=(1+\pi)^{\ell^s}-1$. Hence
$\lambda(X_s)=U^{\ell^{s-1}}-1$ and $\lambda(Q)\lambda(X_s)=U^{\ell^s}-1$. The
constant term of $\lambda(X_s)$ is $\zeta^{\ell^{s-1}}-1\neq0$, so $\lambda(X_s)$
is a unit of $L[[X]]$. As $(U^k)'=kU^k$, the coefficient of $X$ in $U^k$ is
$k\zeta^k$; so $U^{\ell^s}-1$ has constant term $\zeta^{\ell^s}-1=0$ and
coefficient of $X$ equal to $\ell^s\neq0$, and it is $Xw'$ for a unit $w'$. Then
$\lambda(Q)=Xw'\lambda(X_s)^{-1}$.
% Lean: lev_at_root_of_unity, fourth part, from LevPt.lev_values (U = C ζ * E from lev_π,
% X_eq, Q_eq, hXQ by geom_sum_mul, coeff_one_pow_of_deriv, simple_zero). The datum's l is
% PD.p. About 220 lines.
\end{proof}

\begin{lemma}\label{asm-two-three-killing-wach-frame-exists-inputs-heights-division}
Let $\ell$, $s$, $K$, $L$ and $D$ be as in
Definition~\ref{asm-two-three-killing-wach-frame-exists-inputs-heights-expansions},
with its notation, and let $a\in A$ and $k\geq0$. If $X^k$ divides
$\lambda_{i,\zeta}(a)$ for every $i\in\mathbf Z/n$ and every primitive
$\ell^s$-th root of unity $\zeta\in L$, then $Q^k$ divides $a$ in $A$.
\end{lemma}
\leanref{Development/AbelianVarietyTheory/Proofs/InertiaEigenvalues/WachWeierstrassTestAndDensity.lean}{GlobalTheory.NL.WDiv.weierstrass_dense}
\begin{proof}
The ring $\mathcal O=\mathcal O_L$ is a complete discrete valuation ring with
fraction field $L$, and $\ell$ lies in its maximal ideal $\mathfrak m$. As $L$
contains $\tau(K_s)$ for every embedding $\tau$ of $K_s$ in $\bar{\mathbf
Q}_\ell$, it contains the primitive $\ell^s$-th roots of unity of
$\bar{\mathbf Q}_\ell$; they are integral, so they lie in $\mathcal O$. Let $P$
be the set of them. For $\zeta\in P$, $\zeta-1\in\mathfrak m$: in the residue
field, of characteristic $\ell$, $(\zeta-1)^{\ell^s}=\zeta^{\ell^s}-1=0$.

Fix $i\in\mathbf Z/n$, let $a_i\in\mathcal O[[\pi]]$ be the $i$-th component of
$a$, and let $\Delta\colon\mathcal O[[\pi]]\to A$ be the diagonal map. For
$\zeta\in P$ put $\jmath_\zeta=\lambda_{i,\zeta}\circ\Delta$, a map of
$\mathcal O$-algebras $\mathcal O[[\pi]]\to L[[X]]$. The $i$-th component of
$a-\Delta(a_i)$ is $0$, so
$\lambda_{i,\zeta}(a)=\jmath_\zeta(a_i)$
(Definition~\ref{asm-two-three-killing-wach-frame-exists-inputs-heights-expansions}(2)),
and $X^k$ divides $\jmath_\zeta(a_i)$. By the same definition,
$\jmath_\zeta(\pi)=\zeta\exp(X)-1$, so
\[
 \jmath_\zeta\big(\pi-(\zeta-1)\big)=\zeta\big(\exp(X)-1\big)=X\,w_\zeta,
 \qquad w_\zeta=\zeta\sum_{j\geq0}\frac{X^j}{(j+1)!},
\]
and $w_\zeta$ is a unit of $L[[X]]$, with constant term $\zeta\neq0$.
Lemma~\ref{is-onto-step-exp-bound-link-depth-norm-descent}(4), for the ideal
$\mathfrak m$, $x=\zeta-1$ and $\jmath=\jmath_\zeta$, gives
$(\pi-(\zeta-1))^k\mid a_i$ for every $\zeta\in P$. The elements $\zeta-1$,
$\zeta\in P$, are distinct elements of $\mathfrak m$, so by
Lemma~\ref{is-onto-step-exp-bound-link-depth-norm-division}(3)
\[
 \Big(\prod_{\zeta\in P}\big(1+\pi-\zeta\big)\Big)^k\ \text{divides}\ a_i.
\]

It remains to see that the product is the $i$-th component $Q_i$ of $Q$. In
$\mathcal O[[\pi]]$, $\pi q_0=(1+\pi)^\ell-1=\pi\sum_{j<\ell}(1+\pi)^j$, so
$q_0=\sum_{j<\ell}(1+\pi)^j$, as $\mathcal O[[\pi]]$ is a domain. In $A$,
$\varphi(1+\pi)=1+\pi q=(1+\pi)^\ell$
(Definition~\ref{asm-two-three-killing-wach-frame-exists-data}(1)), so
$\varphi^{s-1}(1+\pi)=(1+\pi)^{\ell^{s-1}}$ and
$Q=\varphi^{s-1}(q)=\sum_{j<\ell}(1+\pi)^{j\ell^{s-1}}$. This is
$\Phi(1+\pi)$ for the cyclotomic polynomial
$\Phi(Y)=\sum_{j<\ell}Y^{j\ell^{s-1}}$ of order $\ell^s$, and
$\Phi(Y)=\prod_{\zeta\in P}(Y-\zeta)$ in $\mathcal O[Y]$, as the domain
$\mathcal O$ contains a primitive $\ell^s$-th root of unity. Hence
$Q_i=\prod_{\zeta\in P}(1+\pi-\zeta)$, and $Q_i^k\mid a_i$. As this holds for
every $i$, $Q^k$ divides $a$ in $A=\prod_i\mathcal O[[\pi]]$.
% Lean: Q_test (WachWeierstrassTestAndDensity.lean): Q_eq and prod_prim (cyclotomic_eq_prod_
% X_sub_primitiveRoots for μ₀ from exists_prim and root_integral), test = Per.weierstrass_test
% (X_pow_dvd_ev_iff, prime_X_sub_C, prod_prime_pow_dvd, with Mathlib's Weierstrass division in
% exists_div_X_sub_C), ev_X (from lev_π), lev_eq (from lev_letter), not_isUnit_sub_one;
% integers_completeDVR for O_L. The book cites the same facts from Chapter is (Lemmas
% is-onto-step-exp-bound-link-depth-norm-division and -norm-descent), proved there from other
% copies in the Lean. About 300 lines.
\end{proof}

\begin{lemma}\label{asm-two-three-killing-wach-frame-exists-inputs-heights-letters}
Let $\ell$, $s$, $K$, $L$ and $D$ be as in
Definition~\ref{asm-two-three-killing-wach-frame-exists-inputs-heights-expansions},
with its notation, and let $(V,T)$ be a lattice representation of $G_K$ over $L$
with $V$ crystalline. Then for every $i\in\mathbf Z/n$, $\varphi$ maps
$D_{i-1}(V)$ bijectively onto $D_i(V)$, and $D_i(V)$ has dimension $\dim_LV$
over $L$.
\end{lemma}
\leanref{Development/AbelianVarietyTheory/Proofs/InertiaEigenvalues/CrystallineFrobeniusBijOnLettersDirectSum.lean}{GlobalTheory.NL.PhiD.phi_letters}
\leanref{Development/AbelianVarietyTheory/Proofs/InertiaEigenvalues/FinrankDcrisLetterDeRhamComponentEqDim.lean}{GlobalTheory.Per.LetterDim.tied_letter_component_finrank}
\begin{proof}
Put $r=\dim_LV$ and $D=D_{\mathrm{cris},K}(V)$, and let $\sigma$ be as in
Definition~\ref{asm-two-three-killing-wach-frame-exists-inputs-heights-expansions}(1).
For a $\mathbf Q_\ell$-embedding $\tau\colon K_0\to L$ let $D_\tau\subset D$ be
the subspace on which every $c\in K_0$ acts as $\tau(c)$; so $D_i(V)=D_{\tau_i}$.
By Definition~\ref{asm-two-three-killing-wach-frame-exists-inputs-heights-expansions}(1)
there are $n$ such embeddings.

First, $\dim_LD=nr$. As $V$ is crystalline, $\dim_{K_0}D=\dim_{\mathbf
Q_\ell}V=[L:\mathbf Q_\ell]\,r$, where $K_0\subset B_{\mathrm{cris}}$ acts on
the second factor; and
$[L:\mathbf Q_\ell]\dim_LD=\dim_{\mathbf Q_\ell}D=n\dim_{K_0}D$.

Second, $D$ is the direct sum of the $D_\tau$. The elements of $K_0$ are
$G_K$-invariant in $B_{\mathrm{cris}}$, so $K_0$ acts on $D$, commuting with
the action of $L$ through $V$, and $D$ is a module over $K_0\otimes_{\mathbf
Q_\ell}L$. As there are $n=[K_0:\mathbf Q_\ell]$ embeddings, the map
$c\otimes a\mapsto(\tau(c)a)_\tau$ from $K_0\otimes_{\mathbf Q_\ell}L$ to
$L^n$ is an isomorphism, and its idempotents $e_\tau$ split $D$ as
$\bigoplus_\tau e_\tau D$, with $e_\tau D=D_\tau$. Hence
$\sum_\tau\dim_LD_\tau=nr$.

Third, for every $\tau$, $\varphi$ maps $D_{\tau\circ\sigma}$ injectively
into $D_\tau$. The Frobenius of $B_{\mathrm{cris}}$ is injective and commutes
with $G_K$ \cite[\S9.1]{BC09}, and $V$ is free over $\mathbf Q_\ell$, so
$\varphi$ is injective on $V\otimes B_{\mathrm{cris}}$ and maps $D$ into $D$.
For $c\in K_0$ and $x\in V\otimes B_{\mathrm{cris}}$,
$\varphi(cx)=\sigma(c)\varphi(x)$. If $x\in D_{\tau\circ\sigma}$ and $c\in
K_0$, write $c=\sigma(c')$; then
$c\,\varphi(x)=\varphi(c'x)=\tau(\sigma(c'))\varphi(x)=\tau(c)\varphi(x)$.

Fourth, every embedding $\tau'$ is $\tau\circ\sigma^k$ for some $k\geq0$. The
fixed field of $\sigma$ in $K_0$ is $\mathbf Q_\ell$: if $\sigma(c)=c$, then
$\varphi(c)=c$ and $c\in K\subset\operatorname{Fil}^0B_{\mathrm{dR}}$, and
$B_{\mathrm{cris}}^{\varphi=1}\cap\operatorname{Fil}^0B_{\mathrm{dR}}=\mathbf
Q_\ell$ \cite{Fon94a}. So the group generated by $\sigma$ has $n$ elements,
and the $n$ embeddings $\tau\circ\sigma^k$, $0\leq k<n$, are distinct; they
are all of them.

By the third and fourth steps, $\dim D_{\tau'}\leq\dim D_\tau$ for all
$\tau,\tau'$, so all the $D_\tau$ have the same dimension, which is $r$ by
the first two steps. Finally,
$\tau_i\circ\sigma=\tau_{i-1}$, so $\varphi$ maps $D_{i-1}(V)$ injectively into
$D_i(V)$; both have dimension $r$, so it maps $D_{i-1}(V)$ onto $D_i(V)$.
% Lean: phi_letters = PhiD.phi_bijOn (phi_mem, phi_injective, emb_σ) with
% Per.LetterDim.letter_eq (IsCrys; DeRhamSplit.finrank_eq_sum through evalEmb_bijective and
% FontCmp.mK; RegularRing.letter_le; RegularRing.exists_pow_sigma, which uses PD.fund and
% IntermediateField.finrank_fixedField_eq_card), with LevInj.card_emb for the count of
% embeddings. About 400 lines; the rest of the closure is the datum's period rings
% (isDomain_of_tied, which Fontaine's B_cris satisfies).
\end{proof}

\begin{lemma}\label{asm-two-three-killing-wach-frame-exists-inputs-heights-standard-values}
Let $\ell$, $s$, $K$, $L$ and $D$ be as in
Definition~\ref{asm-two-three-killing-wach-frame-exists-inputs-heights-expansions},
with its notation, let $i\in\mathbf Z/n$, and let $\zeta\in L$ be a root of
unity of $\ell$-power order. Then $\lambda_{i,\zeta}(Q)\neq0$, and if
$\zeta^{\ell^{s-1}}=1$, then $\lambda_{i,\zeta}(Q)$ is a unit of $L[[X]]$.
\end{lemma}
\leanref{Development/AbelianVarietyTheory/Proofs/InertiaEigenvalues/WachLevelExpansionAtRootOfUnity.lean}{GlobalTheory.NL.lev_at_root_of_unity}
\begin{proof}
Write $\lambda=\lambda_{i,\zeta}$ and $U=\lambda(1+\pi)=\zeta\exp(X)$
(Definition~\ref{asm-two-three-killing-wach-frame-exists-inputs-heights-expansions}(2)),
so that $U$ has constant term $\zeta$ and $U'=U$. As in the proofs of
Lemmas~\ref{asm-two-three-killing-wach-frame-exists-inputs-heights-points}
and~\ref{asm-two-three-killing-wach-frame-exists-inputs-heights-division},
$Q=\sum_{j<\ell}(1+\pi)^{j\ell^{s-1}}$ and $X_sQ=(1+\pi)^{\ell^s}-1$ with
$X_s=(1+\pi)^{\ell^{s-1}}-1$. So $\lambda(Q)=\sum_{j<\ell}U^{j\ell^{s-1}}$ and
$\lambda(Q)\lambda(X_s)=U^{\ell^s}-1$. As $(U^k)'=kU^k$, the coefficient of $X$
in $U^k$ is $k\zeta^k$; so the coefficient of $X$ in $U^{\ell^s}-1$ is
$\ell^s\zeta^{\ell^s}\neq0$, and $\lambda(Q)\neq0$. If
$\zeta^{\ell^{s-1}}=1$, the constant term of $\lambda(Q)$ is
$\sum_{j<\ell}\zeta^{j\ell^{s-1}}=\ell\neq0$, so $\lambda(Q)$ is a unit.
% Lean: the third and last parts of lev_at_root_of_unity, through LevPt.lev_values (hQ, hXQ,
% hc1, hQone). About 100 lines.
\end{proof}

\begin{lemma}\label{asm-two-three-killing-wach-frame-exists-inputs-heights-standard-generator}
Let $\ell$, $s$, $K$, $L$ and $D$ be as in
Definition~\ref{asm-two-three-killing-wach-frame-exists-inputs-heights-expansions},
with the generator $\gamma_0$ of $\Gamma$ and the valuation $v$ of
Definition~\ref{asm-two-three-killing-wach-frame-exists-data}(6), and put
$c=\chi_\Gamma(\gamma_0)$. Then $c\neq0$, and $c^k\neq1$ for every integer
$k\neq0$.
\end{lemma}
\leanref{Development/AbelianVarietyTheory/Proofs/InertiaEigenvalues/WachExpansionLatticeGradedOfLevelLe.lean}{GlobalTheory.NL.GradedLat.zpow_ne_one}
\begin{proof}
As $v(c-1)=s$ is finite, $c\neq1$, and $c\neq0$ as $v(-1)=0\neq s$. It suffices
to show that $c$ is not a root of unity. Suppose it is, of order $N>1$; let $p$
be a prime dividing $N$ and put $y=c^{N/p}-1\neq0$. As $c\in\mathcal O$,
$c^m-1=(c-1)(1+c+\cdots+c^{m-1})$ has valuation at least $s$ for $m\geq1$; so
$v(y)\geq s\geq1$. Expanding $(1+y)^p=1$ gives $y\big(p+\binom p2y+y^2z\big)=0$
with $v(z)\geq0$, so $p=-\binom p2y-y^2z$. If $p\neq\ell$, the right side has
valuation at least $1$, while $v(p)=0$. If $p=\ell$ is odd, $\ell$ divides
$\binom\ell2$ and the right side has valuation at least $2$; if $p=\ell=2$, then
$s\geq2$ and again the right side has valuation at least $2$. Both contradict
$v(\ell)=1$.
% Lean: GradedLat.zpow_ne_one = pow_ne_one (expand, Nat.Prime.dvd_choose_self) with
% ne_zero_of_val, from gam_gen and v_l; AtZero.lower uses the same fact through
% chi_gam_zpow_ne_one. About 110 lines.
\end{proof}

\begin{lemma}\label{asm-two-three-killing-wach-frame-exists-inputs-heights-standard-descent}
Let $\ell$, $s$, $K$, $L$ and $D$ be as in
Definition~\ref{asm-two-three-killing-wach-frame-exists-inputs-heights-expansions},
with its notation, let $(V,T)$ be a lattice representation of $G_K$ over $L$
with $V$ crystalline, let $i\in\mathbf Z/n$, and let $\zeta\in L$ be a root of
unity of $\ell$-power order. Suppose that there is $m_0\geq0$ such that every
sequence $f\colon\mathbf Z\to D_i(V)$ with $f(u)=0$ for $u<m_0$ lies in
$\Lambda_{i,\zeta}(T)$. Then there is $m_0'\geq0$ such that every sequence
$g\colon\mathbf Z\to D_{i-1}(V)$ with $g(u)=0$ for $u<m_0'$ lies in
$\Lambda_{i-1,\zeta^\ell}(T)$.
\end{lemma}
\leanref{Development/AbelianVarietyTheory/Proofs/InertiaEigenvalues/WachExpansionLatticeFullDescendsStatement.lean}{GlobalTheory.T.wachExpansionLattice_full_descends}
\leanref{Development/AbelianVarietyTheory/Proofs/InertiaEigenvalues/WachExpansionLatticeFullDescends.lean}{GlobalTheory.NL.FullDown.full_down}
\begin{proof}
Write $\lambda=\lambda_{i,\zeta}$, $\varepsilon=\varepsilon_{i,\zeta}$,
$\varepsilon'=\varepsilon_{i-1,\zeta^\ell}$, $\Lambda=\Lambda_{i,\zeta}(T)$ and
$\Lambda'=\Lambda_{i-1,\zeta^\ell}(T)$, and $N=N(T)$. For a sequence $g$ that
vanishes below some integer put $(\Theta g)(u)=\ell^u\varphi(g(u))$, and for
$\eta\in L[[X]]$ let $\eta_\ell(X)=\eta(\ell X)$. By
Definition~\ref{asm-two-three-killing-wach-frame-exists-inputs-heights-expansions}(3),
as $\iota$ is injective \cite[Theorem~9.1.5]{BC09},
$\varepsilon(ax)=\lambda(a)*\varepsilon(x)$ for $a\in A$ and $x\in N$.

(a) $\Theta(\eta*g)=\eta_\ell*\Theta(g)$, by a direct computation:
$\ell^u\eta_t\varphi(g(u-t))=(\ell^t\eta_t)\,\ell^{u-t}\varphi(g(u-t))$. As
$\eta\mapsto\eta_\ell$ is bijective, $\Theta(\Lambda')$ is stable under $\eta*$
for every $\eta\in L[[X]]$. The Frobenius $\varphi$ is injective on
$V\otimes B_{\mathrm{cris}}$ (proof of
Lemma~\ref{asm-two-three-killing-wach-frame-exists-inputs-heights-letters}), so
$\Theta$ is injective.

(b) There is $k\geq0$ with $\varepsilon(Q^ky)\in\Theta(\Lambda')$ for every
$y\in N$. The $y\in N$ with $\varepsilon(y)\in\Theta(\Lambda')$ form an
$A$-submodule $N'$, by (a) and the semilinearity of $\varepsilon$. For $x\in N$
and $m$ with $Q^m\varphi_D(x)\in N$,
Definition~\ref{asm-two-three-killing-wach-frame-exists-inputs-heights-expansions}(4)
says $\varepsilon(Q^m\varphi_D(x))=\lambda(Q^m)*\Theta(\varepsilon'(x))$, which
lies in $\Theta(\Lambda')$ by (a); so $Q^m\varphi_D(x)\in N'$. Let $z\in N$. By
Definition~\ref{asm-two-three-killing-wach-frame-exists-data}(4),
$Q^jz=\sum_ra_r\varphi_D(x_r)$ for some $j\geq0$ and finitely many $a_r\in A$
and $x_r\in N$; with $m$ such that $Q^m\varphi_D(x_r)\in N$ for every $r$,
$Q^{m+j}z=\sum_ra_rQ^m\varphi_D(x_r)\in N'$. As $N$ is finitely generated and
$N'$ is a submodule, one $k$ works for every $y\in N$.

(c) Put $q=\lambda(Q^k)$, nonzero by
Lemma~\ref{asm-two-three-killing-wach-frame-exists-inputs-heights-standard-values},
and write $q=X^bw$ with $w$ a unit. For $f\in\Lambda$, $q*f\in\Theta(\Lambda')$:
$\Lambda$ is spanned by the $\eta*\varepsilon(x)$, and
$q*(\eta*\varepsilon(x))=\eta*\varepsilon(Q^kx)$ lies in $\Theta(\Lambda')$ by
(b) and (a).

Put $m_0'=m_0+b$, and let $g\colon\mathbf Z\to D_{i-1}(V)$ vanish below $m_0'$.
Let $h(u)=(\Theta g)(u+b)$. By
Lemma~\ref{asm-two-three-killing-wach-frame-exists-inputs-heights-letters},
$\varphi$ maps $D_{i-1}(V)$ into $D_i(V)$, so $h$ takes values in $D_i(V)$ and
vanishes below $m_0$, and so does $w^{-1}*h$. So $w^{-1}*h\in\Lambda$, and by
(c) $q*(w^{-1}*h)=X^b*h=\Theta g$ lies in $\Theta(\Lambda')$. As $\Theta$ is
injective, $g\in\Lambda'$.
% Lean: full_down = full_down_generic: map_conv and Tw_conv for (a), phi_inj, pow_mem for
% (b) (N_φ, N_φinv, N_fg), lev_Q_ne for q ≠ 0 (the Lean proves it from CmpTie.lev_φ; the
% book uses lev_values instead), conv_latSpan, phi_mem and full_of_conv. About 400 lines.
\end{proof}

\begin{lemma}\label{asm-two-three-killing-wach-frame-exists-inputs-heights-standard-graded}
Let $\ell$, $s$, $K$, $L$ and $D$ be as in
Definition~\ref{asm-two-three-killing-wach-frame-exists-inputs-heights-expansions},
with its notation, let $(V,T)$ be a lattice representation of $G_K$ over $L$
with $V$ crystalline, let $i\in\mathbf Z/n$, and let $\zeta\in L$ with
$\zeta^{\ell^s}=1$. Suppose that there is $m_0\geq0$ such that every sequence
$f\colon\mathbf Z\to D_i(V)$ with $f(u)=0$ for $u<m_0$ lies in
$\Lambda=\Lambda_{i,\zeta}(T)$. For $u\in\mathbf Z$ and $d\in
V\otimes B_{\mathrm{cris}}$ let $\delta_ud$ be the sequence with value $d$ at
$u$ and $0$ elsewhere. Then a sequence $f\colon\mathbf Z\to V\otimes
B_{\mathrm{cris}}$ lies in $\Lambda$ if and only if $f$ vanishes below some
integer and $\delta_uf(u)\in\Lambda$ for every $u\in\mathbf Z$.
\end{lemma}
\leanref{Development/AbelianVarietyTheory/Proofs/InertiaEigenvalues/WachExpansionLatticeGradedOfLevelLeStatement.lean}{GlobalTheory.T.wachExpansionLattice_graded_of_level_le}
\leanref{Development/AbelianVarietyTheory/Proofs/InertiaEigenvalues/WachExpansionLatticeGradedOfLevelLe.lean}{GlobalTheory.NL.GradedLat.graded_of_level_le}
\begin{proof}
Write $\varepsilon=\varepsilon_{i,\zeta}$, let $\gamma_0$ be as in
Definition~\ref{asm-two-three-killing-wach-frame-exists-data}(6), and put
$c=\chi_\Gamma(\gamma_0)$. The elements of $\Lambda$ take values in $D_i(V)$
and vanish below some integer, by
Definition~\ref{asm-two-three-killing-wach-frame-exists-inputs-heights-expansions}(3).

(a) For $x\in N(T)$, $\gamma_0x\in N(T)$
(Definition~\ref{asm-two-three-killing-wach-frame-exists-data}(4)) and
$\varepsilon(\gamma_0x)(u)=c^u\varepsilon(x)(u)$, by
Definition~\ref{asm-two-three-killing-wach-frame-exists-inputs-heights-expansions}(5)
with $a=1$: $\zeta$ is a primitive $\ell^j$-th root of unity with $j\leq s$, and
$v(c-1)=s$.

(b) If $f\in\Lambda$, then $(c^uf(u))_u\in\Lambda$. For a generator,
$c^u(\eta*\varepsilon(x))(u)=\sum_t\eta_tc^t\,c^{u-t}\varepsilon(x)(u-t)
=(\eta_c*\varepsilon(\gamma_0x))(u)$ with $\eta_c(X)=\eta(cX)$, by (a).

(c) If $f\in\Lambda$ vanishes outside a finite set $S$, then $\delta_uf(u)\in
\Lambda$ for every $u$. By induction on $S$; the case $S=\emptyset$ is clear.
Let $S=S'\cup\{a\}$ with $a\notin S'$. By (b),
$g=(c^uf(u))_u-c^af\in\Lambda$, and $g(u)=(c^u-c^a)f(u)$ vanishes outside
$S'$. By induction $\delta_u((c^u-c^a)f(u))\in\Lambda$ for every $u$, and for
$u\neq a$, $c^u-c^a\neq0$ by
Lemma~\ref{asm-two-three-killing-wach-frame-exists-inputs-heights-standard-generator},
so $\delta_uf(u)\in\Lambda$. Finally $\delta_af(a)=f-\sum_{u\in S'}\delta_uf(u)\in
\Lambda$.

Now let $f\in\Lambda$, vanishing below $a$. The sequence
$f-\sum_{a\leq u<m_0}\delta_uf(u)$ takes values in $D_i(V)$ and vanishes below
$m_0$, so it lies in $\Lambda$; hence so does the finite sum, and by (c)
$\delta_uf(u)\in\Lambda$ for $a\leq u<m_0$. For $u<a$, $\delta_uf(u)=0$, and for
$u\geq m_0$, $\delta_uf(u)\in\Lambda$ by the hypothesis. Conversely, let $f$
vanish below $a$ with every $\delta_uf(u)\in\Lambda$. Then $f$ takes values in
$D_i(V)$, so $f-\sum_{a\leq u<m_0}\delta_uf(u)\in\Lambda$ by the hypothesis, and
$f\in\Lambda$.
% Lean: graded_of_level_le = graded_Lat: ex_gam (CmpTie.levB_Γ with a = 1, gam_gen),
% latSpan_stable, Per.single_mem_of_support_subset, graded_of_full (tail), zpow_ne_one,
% Model.Lat_coeff. About 250 lines.
\end{proof}

\begin{lemma}\label{asm-two-three-killing-wach-frame-exists-inputs-heights-full-embedding}
Let $\ell$, $s$, $K$, $L$ and $D$ be as in
Definition~\ref{asm-two-three-killing-wach-frame-exists-inputs-heights-expansions},
with its notation and the roots of unity $\zeta_{\ell^j}\in\bar{\mathbf Q}_\ell$
of Definition~\ref{asm-two-three-killing-wach-frame-exists-data}, and regard
$\bar{\mathbf Q}_\ell$ as a subfield of $B_{\mathrm{dR}}$. Let $i\in\mathbf Z/n$
and let $\zeta\in L$ be a primitive $\ell^s$-th root of unity.
\begin{enumerate}
\item There is a $\mathbf Q_\ell$-embedding $\tau_s\colon K_s\to L$ with
  $\tau_s=\tau_i$ on $K_0$ and $\tau_s(\zeta_{\ell^s})=\zeta$.
\item For such $\tau_s$, let $\tau$ be its restriction to $K$. Then $\tau=\tau_i$
  on $K_0$, and there is a map of $\mathbf Q_\ell$-algebras $\kappa\colon L\to
  B_{\mathrm{dR}}$ with $\kappa(\tau(c))=c$ for $c\in K$ and
  $\kappa(\zeta)=\zeta_{\ell^s}$; in particular $\kappa(\tau_i(c))=c$ for $c\in
  K_0$.
\end{enumerate}
\end{lemma}
\leanref{Development/AbelianVarietyTheory/Proofs/InertiaEigenvalues/UnramifiedSubfieldLeCyclotomicExt.lean}{GlobalTheory.NL.Points.points}
\begin{proof}
(1) Let $g_0\in G_{\mathbf Q_\ell}$ extend the embedding $\tau_i$ of $K_0$ into
$L\subset\bar{\mathbf Q}_\ell$. Then $g_0^{-1}(\zeta)$ is a primitive
$\ell^s$-th root of unity. As $K_0$ is unramified over $\mathbf Q_\ell$, the
cyclotomic polynomial of order $\ell^s$ is irreducible over $K_0$
\cite[Chapter~IV, \S4]{Ser79}, and it is the minimal polynomial over $K_0$ of
both $g_0^{-1}(\zeta)$ and $\zeta_{\ell^s}$. So there is $\sigma\in G_{K_0}$
with $\sigma(\zeta_{\ell^s})=g_0^{-1}(\zeta)$. Put $g=g_0\sigma$: it equals
$\tau_i$ on $K_0$ and maps $\zeta_{\ell^s}$ to $\zeta$. As $L$ contains $g(K_s)$
(Definition~\ref{asm-two-three-killing-wach-frame-exists-data}), the
restriction of $g$ to $K_s$ is an embedding $\tau_s$ into $L$ as required.

(2) Let $g\in G_{\mathbf Q_\ell}$ extend $\tau_s$, and let $\kappa$ be the
restriction of $g^{-1}$ to $L$, followed by
$\bar{\mathbf Q}_\ell\subset B_{\mathrm{dR}}$. For $c\in K$,
$\kappa(\tau(c))=g^{-1}(g(c))=c$, and
$\kappa(\zeta)=g^{-1}(\tau_s(\zeta_{\ell^s}))=\zeta_{\ell^s}$. As $K_0\subset
K_s$, $\tau=\tau_s=\tau_i$ on $K_0$.
% Lean: points = exists_pt (exists_aut, Cyc.minpoly_unramifiedSubfield_eq_cyclotomic for the
% irreducibility, IsConjRoot.exists_algEquiv, exists_restr), pt_props (aut_of_pt, zeta_prim)
% and card_emb. The Lean's Pt (HW l) h i ζ τ is the existence of τ_s as in (1) restricting to
% τ, and Comp the properties of ι in (2). About 140 lines; the rest of the closure is the
% minimal polynomial over K_0.
\end{proof}

\begin{lemma}\label{asm-two-three-killing-wach-frame-exists-inputs-heights-full-evaluation}
Let $\ell$ be a prime and $s\geq1$, let $B_{\mathrm{dR}}$ be Fontaine's ring of
de Rham periods with its filtration $\operatorname{Fil}^\bullet$, let
$t\in\operatorname{Fil}^1B_{\mathrm{dR}}$ be Fontaine's element
$\log[\varepsilon]$, and put $t_s=\ell^{-s}t$. Let $L$ be a finite extension of
$\mathbf Q_\ell$ and $\kappa\colon L\to B_{\mathrm{dR}}$ a map of
$\mathbf Q_\ell$-algebras. Then:
\begin{enumerate}
\item $t_s$ is a unit of $B_{\mathrm{dR}}$; for $a\in\mathbf Z$ and $b\in
  B_{\mathrm{dR}}$, $t_sb\in\operatorname{Fil}^{a+1}$ if and only if
  $b\in\operatorname{Fil}^a$; the only element of every $\operatorname{Fil}^a$
  is $0$; and $\kappa(L)\subset\operatorname{Fil}^0$;
\item there is a ring map $\mathrm{ev}\colon L[[X]]\to B_{\mathrm{dR}}$ such that
  for $f=\sum_kf_kX^k$ and $M\geq0$,
  $\mathrm{ev}(f)-\sum_{k<M}\kappa(f_k)t_s^k\in\operatorname{Fil}^M$, and
  $\mathrm{ev}(f)\in\operatorname{Fil}^M$ if and only if $X^M$ divides $f$.
\end{enumerate}
\end{lemma}
\leanref{Development/AbelianVarietyTheory/Proofs/InertiaEigenvalues/ExistsPowerSeriesEvaluation.lean}{GlobalTheory.NL.EvdR.exists_ev_dR}
\begin{proof}
The ring $B_{\mathrm{dR}}^+=\operatorname{Fil}^0B_{\mathrm{dR}}$ is a complete
discrete valuation ring with uniformizer $t$ and fraction field
$B_{\mathrm{dR}}$, and $\operatorname{Fil}^a=t^aB_{\mathrm{dR}}^+$
\cite[Proposition~4.4.6, Definition~4.4.7 and Proposition~4.4.8]{BC09}. As
$\ell$ is a unit of $B_{\mathrm{dR}}^+$, $t_s$ is also a
uniformizer, which gives the first three assertions of (1). An element $x$ of
$B_{\mathrm{dR}}$ algebraic over $\mathbf Q_\ell$ lies in $B_{\mathrm{dR}}^+$:
otherwise, in a relation $x^d+\sum_{j<d}a_jx^j=0$ with $a_j\in\mathbf Q_\ell$,
the term $x^d$ has smaller valuation than all the others. This gives
$\kappa(L)\subset\operatorname{Fil}^0$.

(2) As $B_{\mathrm{dR}}^+$ is $t$-adically complete, the series
$\mathrm{ev}(f)=\sum_k\kappa(f_k)t_s^k$ converges, and $\mathrm{ev}$ is a ring map
with the first property. If $f=X^Mg$, then $\mathrm{ev}(f)=t_s^M\mathrm{ev}(g)\in
\operatorname{Fil}^M$. If $X^M$ does not divide $f$, let $k<M$ be least with
$f_k\neq0$; then $\mathrm{ev}(f)\equiv\kappa(f_k)t_s^k$ modulo
$\operatorname{Fil}^{k+1}$, and $\kappa(f_k)$ is a unit of $B_{\mathrm{dR}}^+$,
as $\kappa(f_k)^{-1}=\kappa(f_k^{-1})$ also lies in $B_{\mathrm{dR}}^+$. So
$\mathrm{ev}(f)\notin\operatorname{Fil}^{k+1}\supset\operatorname{Fil}^M$.
% Lean: exists_ev_tied = Bridge.exists_power_series_evaluation, for the datum's B_dR read in
% Fontaine's sense (tS P s = l^(-s) ιB(tB)): exists_ev (exists_lim, dvd_part_sub), ev_sub,
% ev_dvd_iff, mem_zero_of_isIntegral, tS_mem_iff, reduction. About 300 lines; the rest of the
% closure is the datum's B_dR.
\end{proof}

\begin{lemma}[{\cite[\S4.3, Lemma~4.3.3 and Theorem~4.3.5]{BC09}}]\label{asm-two-three-killing-wach-frame-exists-decompletion-frobenius-tilt}
Let $\ell$ be a prime and use the notation of
Definition~\ref{asm-two-three-killing-wach-frame}. Let
$\mathcal O_{\mathbf C}^\flat$ be the ring of sequences $(x_n)_{n\geq0}$ in
$\mathcal O_{\mathbf C}/\ell$ with $x_{n+1}^\ell=x_n$ for every $n$, with the
operations taken term by term, and let $\mathbf C_\ell^\flat$ be its fraction
field. Then $\mathcal O_{\mathbf C}^\flat$ is a valuation ring of characteristic
$\ell$ on which $x\mapsto x^\ell$ is bijective, and $x\mapsto x^\ell$ is bijective
on $\mathbf C_\ell^\flat$.
\end{lemma}
\leanref{Development/AbelianVarietyTheory/Proofs/InertiaEigenvalues/FrobeniusIntegralityWittTiltAinf.lean}{GlobalTheory.FrobInt.valuationRing_tilt}
\leanref{Development/AbelianVarietyTheory/Proofs/InertiaEigenvalues/FrobeniusIntegralityWittTiltAinf.lean}{GlobalTheory.FrobInt.dvd_total_tilt}
\leanref{Development/AbelianVarietyTheory/Proofs/InertiaEigenvalues/FrobeniusIntegralityWittTiltAinf.lean}{GlobalTheory.FrobInt.perfectRing_Cb}
% Lean: the tilt is Mathlib's PreTilt (OC l) l, a perfect domain. valuationRing_tilt is
% ValuationRing.iff_dvd_total with dvd_total_tilt: for nonzero a, b the sharp roots a^(n)
% (sharpRoot, the untilt of the n-th Frobenius root; sharpRoot_pow_succ, mk_sharpRoot) are
% nonzero (sharpRoot_ne_zero); as O_C is a valuation ring, a^(0) | b^(0) or b^(0) | a^(0),
% and this passes to every n since O_C is integrally closed; the ratios b^(n)/a^(n) then
% define c with b = ac (dvd_of_sharpRoot_dvd). This is the argument of BC09 Lemma 4.3.3.
% perfectRing_Cb: a/b has the l-th root a^(1/l)/b^(1/l). About 110 lines.

\begin{lemma}\label{asm-two-three-killing-wach-frame-exists-decompletion-frobenius-residue}
Let $\ell$ be a prime, $R$ a valuation ring and $F$ its fraction field. Let $z$
be a $p\times q$ matrix over $F$ and $s$, $s'$ and $y$ matrices over $R$ of sizes
$p\times p$, $q\times q$ and $p\times q$, such that $z_{ab}^\ell=(szs'+y)_{ab}$
for every entry. Then $z$ has entries in $R$.
\end{lemma}
\leanref{Development/AbelianVarietyTheory/Proofs/InertiaEigenvalues/FrobeniusIntegralityWittTiltAinf.lean}{GlobalTheory.FrobInt.residue_integral}
\begin{proof}
Let $v$ be the valuation of $F$ with valuation ring $R$, written
multiplicatively, so that $x\in R$ if and only if $v(x)\leq1$. Choose an entry
$z_c$ of $z$ with $v(z_c)$ largest, and suppose that $v(z_c)>1$. Every entry of
$szs'$ is a sum of products $s_{ai}z_{ij}s'_{jb}$, each of value at most
$v(z_c)$, as the entries of $s$ and $s'$ lie in $R$; and $v(y_c)\leq1<v(z_c)$.
So $v(z_c)^\ell=v(z_c^\ell)\leq v(z_c)$, which contradicts $v(z_c)>1$ and
$\ell\geq2$. Hence $v(z_{ab})\leq v(z_c)\leq1$ for every entry.
\end{proof}

\begin{lemma}\label{asm-two-three-killing-wach-frame-exists-decompletion-frobenius-witt}
Let $\ell$ be a prime, $R$ a valuation ring of characteristic $\ell$ on which
$x\mapsto x^\ell$ is bijective, and $F$ its fraction field, on which
$x\mapsto x^\ell$ is also bijective. Let $W(R)\subset W(F)$ be their rings of
Witt vectors, and $\varphi$ the endomorphism of $W(F)$ induced by
$x\mapsto x^\ell$. Let $Z$ be a $p\times q$ matrix over $W(F)$, and $S$, $S'$ and
$Y$ matrices over $W(R)$ of sizes $p\times p$, $q\times q$ and $p\times q$, with
$\varphi(Z)=SZS'+Y$. Then $Z$ has entries in $W(R)$.
\end{lemma}
\leanref{Development/AbelianVarietyTheory/Proofs/InertiaEigenvalues/FrobeniusIntegralityWittTiltAinf.lean}{GlobalTheory.FrobInt.integral_of_frobenius}
\leanref{Development/AbelianVarietyTheory/Proofs/InertiaEigenvalues/FrobeniusIntegralityWittTiltAinf.lean}{GlobalTheory.FrobInt.step}
\leanref{Development/AbelianVarietyTheory/Proofs/InertiaEigenvalues/FrobeniusIntegralityWittTiltAinf.lean}{GlobalTheory.FrobInt.approx}
\begin{proof}
For a ring $k$ of characteristic $\ell$ on which $x\mapsto x^\ell$ is bijective,
an element $x$ of $W(k)$ has components $x_0,x_1,\ldots\in k$, and
\cite[Chapter~II, \S\S5--6]{Ser79}: $x\mapsto x_0$ is a ring map $W(k)\to k$,
which carries the endomorphism induced by $x\mapsto x^\ell$ to $x\mapsto x^\ell$;
the Teichmüller lift $[r]$ of $r\in k$ has $[r]_0=r$; $x\in\ell^nW(k)$ if and
only if $x_0,\ldots,x_{n-1}$ vanish; and $x$ and $x+\ell^nw$ have the same
components $x_0,\ldots,x_{n-1}$. Moreover $\ell$ is not a zero divisor in
$W(F)$, and $W(R)\to W(F)$ is injective, applied to the components. Matrices
are treated entry by entry.

First, if $Z$ is a matrix over $W(F)$ with $\varphi(Z)=SZS'+Y$ for a matrix $Y$
over $W(R)$, then $Z=[r]+\ell Z'$ for matrices $r$ over $R$ and $Z'$ over $W(F)$
with $\varphi(Z')=SZ'S'+Y'$ for some matrix $Y'$ over $W(R)$. Indeed, the
matrices $z$, $s$, $s'$ and $y$ of first components satisfy
$z_{ab}^\ell=(szs'+y)_{ab}$, so $z$ is a matrix $r$ over $R$ by
Lemma~\ref{asm-two-three-killing-wach-frame-exists-decompletion-frobenius-residue}.
The entries of $Z-[r]$ have first component $0$, so $Z=[r]+\ell Z'$ for a matrix
$Z'$ over $W(F)$. The matrix $B=S[r]S'+Y-\varphi([r])$ has entries in $W(R)$,
and as $\varphi(Z)=\varphi([r])+\ell\varphi(Z')$ and
$SZS'=S[r]S'+\ell SZ'S'$, the relation gives $B=\ell(\varphi(Z')-SZ'S')$ in
$W(F)$. So the entries of $B$ have first component $0$, in $F$ and hence in $R$,
and $B=\ell Y'$ for a matrix $Y'$ over $W(R)$. As $\ell$ is not a zero divisor
in $W(F)$, $\varphi(Z')=SZ'S'+Y'$.

Second, by induction on $n$, every such $Z$ is $A+\ell^nU$ with $A$ over $W(R)$
and $U$ over $W(F)$. For $n=0$ take $A=0$. For the inductive step write
$Z=[r]+\ell Z'$ as in the first step and $Z'=A'+\ell^nU'$; then
$Z=([r]+\ell A')+\ell^{n+1}U'$.

Finally, let $n\geq0$ and write $Z=A+\ell^{n+1}U$ as in the second step. The
components $0,\ldots,n$ of each entry of $Z$ are those of the same entry of $A$,
which lie in $R$. So every component of every entry of $Z$ lies in $R$, and $Z$
has entries in $W(R)$.
% Lean: FrobInt.step (residue_integral on the constant coefficients, Teichmüller Z0,
% exists_p_mul through WittVector.mem_span_p_pow_iff_le_coeff_eq_zero, B and Y', cancellation
% by WittVector.p_nonzero), approx (the induction) and integral_of_frobenius (the
% coefficients through the truncation, WittVector.mem_ker_truncate). About 140 lines.
\end{proof}

\begin{lemma}\label{asm-two-three-killing-wach-frame-exists-decompletion-frobenius}
Let $\ell$ be a prime, $L$ a finite extension of $\mathbf Q_\ell$ and
$\mathcal O=\mathcal O_L$, and let $R=\mathcal O\otimes_{\mathbf Z_\ell}
W(\mathbf C_\ell^\flat)$, with $\varphi_R$ the Frobenius of the second factor,
and $R_A=\mathcal O\otimes_{\mathbf Z_\ell}A_{\mathrm{inf}}$ as in
Definition~\ref{asm-two-three-killing-wach-frame-exists-data}. Let $Z$ be a
$p\times q$ matrix over $R$, and $S$, $S'$ and $Y$ matrices over $R_A$ of sizes
$p\times p$, $q\times q$ and $p\times q$, with $\varphi_R(Z)=SZS'+Y$. Then $Z$
has entries in $R_A$.
\end{lemma}
\leanref{Development/AbelianVarietyTheory/Proofs/InertiaEigenvalues/FrobeniusIntegralityTensorAinfCoefficientsStatement.lean}{GlobalTheory.T.frobenius_integrality_tensor_ainf_coefficients}
\leanref{Development/AbelianVarietyTheory/Proofs/InertiaEigenvalues/FrobeniusIntegralityTensorAinfCoefficients.lean}{GlobalTheory.FrobIntegralCoeff.integrality_baseChange}
\leanref{Development/AbelianVarietyTheory/Proofs/InertiaEigenvalues/FrobeniusIntegralityWittTiltAinfStatement.lean}{GlobalTheory.T.frobeniusIntegrality_wittTilt_ainf}
\begin{proof}
Choose a $\mathbf Z_\ell$-basis $o_1,\ldots,o_f$ of $\mathcal O$. For a ring
$W$ and $z\in\mathcal O\otimes_{\mathbf Z_\ell}W$ let $\rho(z)\in M_f(W)$ be the
matrix of multiplication by $z$ in the basis $(o_k\otimes1)$. Then $\rho$ is a
ring map, it commutes with maps $W\to W'$ applied to the second factor, and
$z=\sum_{i,j}c_j\,o_i\otimes\rho(z)_{ij}$, where $1=\sum_jc_jo_j$. So $z$ lies in
$\mathcal O\otimes A_{\mathrm{inf}}$ as soon as the entries of $\rho(z)$ lie in
$A_{\mathrm{inf}}$. Replacing each entry of $Z$, $S$, $S'$ and $Y$ by its
$f\times f$ block $\rho(\cdot)$ turns the hypothesis into the same relation for
matrices over $W(\mathbf C_\ell^\flat)$ and $A_{\mathrm{inf}}$, with the
Frobenius of $W(\mathbf C_\ell^\flat)$. So it suffices to treat
$\mathcal O=\mathbf Z_\ell$.

In that case,
Lemma~\ref{asm-two-three-killing-wach-frame-exists-decompletion-frobenius-tilt}
shows that $\mathcal O_{\mathbf C}^\flat$ is a valuation ring of characteristic
$\ell$ on which $x\mapsto x^\ell$ is bijective, as it is on its fraction field
$\mathbf C_\ell^\flat$, and $A_{\mathrm{inf}}=W(\mathcal O_{\mathbf C}^\flat)$.
Lemma~\ref{asm-two-three-killing-wach-frame-exists-decompletion-frobenius-witt},
for $\mathcal O_{\mathbf C}^\flat$ and the matrices $Z$, $S$, $S'$ and $Y$, shows
that $Z$ has entries in $A_{\mathrm{inf}}$.
% Lean: FrobIntegralCoeff.integrality_baseChange (regRep, blow, mem_range_of_regRep) for the
% reduction; the case O = Z_l is T.frobeniusIntegrality_wittTilt_ainf
% (= FrobInt.frobenius_integrality: valuationRing_tilt, perfectRing_Cb and
% integral_of_frobenius). The closure of frobenius_integrality_tensor_ainf_coefficients is
% 37 theorems, about 500 lines (with the definitions of the rings).
\end{proof}

\begin{lemma}\label{asm-two-three-killing-wach-frame-exists-inputs-heights-full-image-sandwich}
Let $\ell$, $s$, $K$, $L$ and $D$ be as in
Definition~\ref{asm-two-three-killing-wach-frame-exists-inputs-heights-expansions},
with its notation, $R_A\subset R$ and $\iota_A$ as in
Definition~\ref{asm-two-three-killing-wach-frame-exists-data}, and put
$\varpi=\iota_A(X)$ with $X=\varphi^{s-1}(\pi)\in A$. Let $(V,T)$ be a lattice
representation of $G_K$ over $L$ with $V$ crystalline. Then
$\varpi=1\otimes([\varepsilon^{(1)}]-1)$, and there is $m\geq0$ such that
$\varpi^mj_T(y)\in R_A\otimes_{\mathcal O}T$ for every $y\in N(T)$, and for
every $t\in T$, $\varpi^m(1\otimes t)$ lies in the $\mathcal O$-submodule of
$R\otimes_{\mathcal O}T$ generated by the $a\,j_T(y)$ with $a\in R_A$ and
$y\in N(T)$.
\end{lemma}
\leanref{Development/AbelianVarietyTheory/Proofs/InertiaEigenvalues/MapMulLeft.lean}{GlobalTheory.NL.SandArb.sandwich}
\leanref{Development/AbelianVarietyTheory/Proofs/InertiaEigenvalues/WachLatticeAinfSandwichInvertPi.lean}{GlobalTheory.HWL.SandProof.sand_of}
\leanref{Development/AbelianVarietyTheory/Proofs/InertiaEigenvalues/WachInclusionXQFrobenius.lean}{GlobalTheory.NL.IoXQ.io_X}
\begin{proof}
Write $\iota=\iota_A$ and $N=N(T)$. As $T$ is free over $\mathcal O$ and
$R_A\subset R$, $R_A\otimes_{\mathcal O}T$ is a submodule of
$R\otimes_{\mathcal O}T$. The map $j_T$ satisfies $j_T(ay)=\iota(a)j_T(y)$ for
$a\in A$, since $E\to R$ restricts to $\iota$ on $A$, and
$j_T(\varphi_Dy)=(\varphi_R\otimes1)(j_T(y))$
(Definition~\ref{asm-two-three-killing-wach-frame-exists-data}(5)).

First, $\varpi=1\otimes([\varepsilon^{(1)}]-1)$ and
$\varphi_R(\varpi)=\varpi\,\iota(Q)$. In $A$,
$\varphi(1+\pi)=1+\pi q=(1+\pi)^\ell$
(Definition~\ref{asm-two-three-killing-wach-frame-exists-data}(1)), so
$1+X=(1+\pi)^{\ell^{s-1}}$ and $XQ=\varphi^{s-1}(\pi q)=(1+\pi)^{\ell^s}-1$. In
$R$, $1+\iota(\pi)=1\otimes[\varepsilon^{(s)}]$, whose $\ell^{s-1}$-th power is
$1\otimes[\varepsilon^{(1)}]$; this gives the formula for $\varpi$. And
$\varphi_R(1+\iota(\pi))=(1+\iota(\pi))^\ell$, as $\varphi_R$ is the Frobenius
of $W(\mathbf C_\ell^\flat)$ on the second factor, so
$\varphi_R(\varpi)=(1+\iota(\pi))^{\ell^s}-1=\iota(XQ)=\varpi\,\iota(Q)$.

Second, $N$ has an $A$-basis $b_1,\ldots,b_r$ whose images $n_i=j_T(b_i)$ form
an $R$-basis of $R\otimes_{\mathcal O}T$, where $r$ is the rank of $T$ over
$\mathcal O$. By
Lemma~\ref{asm-two-three-killing-wach-frame-exists-inputs-basis}, $N$ has an
$A$-basis $b_1,\ldots,b_r$, which is an $E$-basis of $D(T)$, so $r$ is the rank
of $T$ (Definition~\ref{asm-two-three-killing-wach-frame-exists-data}(4)). The
$n_i$ generate $R\otimes T$ over $R$: the image of $j_T$ does, and
$j_T(\sum_ie_ib_i)=\sum_i\bar e_in_i$ for $e_i\in E$, with $\bar e_i$ the image
of $e_i$ in $R$. As $R\otimes T$ is free of rank $r$ over the commutative ring
$R$, these $r$ generators form a basis.

Third, there is $m\geq0$ such that, for every $i$, $Q^m\varphi_D(b_i)\in N$ and
$Q^mb_i$ lies in the $A$-submodule generated by $\varphi_D(N)$.
Definition~\ref{asm-two-three-killing-wach-frame-exists-data}(4) gives such an
exponent for each $b_i$ and each condition, and the largest of them works for
all, as both sets are $A$-submodules.

Let $t_1,\ldots,t_r$ be an $\mathcal O$-basis of $T$, and $Y$ the matrix over $R$
with $n_i=\sum_kY_{ik}(1\otimes t_k)$, which is invertible by the second step.
For a matrix $Z$ over $R$ write $\varphi_R(Z)$ for the matrix of the
$\varphi_R(Z_{ik})$; then $(\varphi_R\otimes1)(n_i)=\sum_k\varphi_R(Y_{ik})
(1\otimes t_k)$.

Fourth, the first assertion. By the third step,
$Q^m\varphi_D(b_i)=\sum_jP_{ij}b_j$ with $P_{ij}\in A$; applying $j_T$,
$\iota(Q)^m\varphi_R(Y)=\iota(P)Y$. By the first step, $Z=\varpi^mY$ satisfies
$\varphi_R(Z)=\varpi^m\iota(Q)^m\varphi_R(Y)=\iota(P)Z$, so $Z$ has entries in
$R_A$ by
Lemma~\ref{asm-two-three-killing-wach-frame-exists-decompletion-frobenius}. So
$\varpi^mn_i\in R_A\otimes T$, and for $y=\sum_ic_ib_i\in N$,
$\varpi^mj_T(y)=\sum_i\iota(c_i)\varpi^mn_i\in R_A\otimes T$.

Fifth, the second assertion. As $\varphi_D(ay)=\varphi(a)\varphi_D(y)$, the
$A$-submodule generated by $\varphi_D(N)$ is generated by the $\varphi_D(b_j)$,
so by the third step $Q^mb_i=\sum_jU_{ij}\varphi_D(b_j)$ with $U_{ij}\in A$, and
$\iota(Q)^mY=\iota(U)\varphi_R(Y)$. Let $Y'=Y^{-1}$, so that
$1\otimes t_k=\sum_iY'_{ki}n_i$. As $\varphi_R(Y)\varphi_R(Y')=1$,
$\iota(Q)^m\varphi_R(Y')=Y'\iota(U)$, and $Z'=\varpi^mY'$ satisfies
$\varphi_R(Z')=Z'\iota(U)$. By
Lemma~\ref{asm-two-three-killing-wach-frame-exists-decompletion-frobenius}, $Z'$
has entries in $R_A$, so $\varpi^m(1\otimes t_k)=\sum_iZ'_{ki}\,j_T(b_i)$ with
$Z'_{ki}\in R_A$. For $t=\sum_kc_kt_k$ with $c_k\in\mathcal O$,
$\varpi^m(1\otimes t)=\sum_kc_k\varpi^m(1\otimes t_k)$.
% Lean: SandArb.sandwich = of_layers applied to sand_layers, which is
% HWL.SandProof.sandwich for the datum's layers (apply_layers): the basis from
% wach_lattice_free_of_crystalline (HWL.lattice_free, the argument WachFree.basis_of_regular
% of Lemma -inputs-basis), frobenius_inclusion_X_eq_mul_Q with Sandwich.φRW_teichmuller and
% x_eq (first step; IoXQ.io_X for the formula for ϖ), and sand_of: hsp and hli (second step),
% uniform (third), Sandwich.upper (fourth) and Sandwich.lower (fifth), with
% frobenius_integrality_coeff (Lemma -decompletion-frobenius). The layers are built from the
% datum through AllLaws.laws (the expansion laws, the equivariance of E -> R, the
% crystallinity of the trivial lattice), which sand_of does not use. sand_of with upper and
% lower is about 250 lines; the closure of SandArb.sandwich is 310 theorems, about 3,900
% lines, of which AllLaws.laws is about 1,900 and wachInclusion_X_Q_frobenius about 1,700.
\end{proof}

\begin{lemma}\label{asm-two-three-killing-wach-frame-exists-inputs-heights-full-image}
Let $\ell$, $s$, $K$, $L$ and $D$ be as in
Definition~\ref{asm-two-three-killing-wach-frame-exists-inputs-heights-expansions},
with its notation, let $(V,T)$ be a lattice representation of $G_K$ over $L$
with $V$ crystalline, let $i\in\mathbf Z/n$, and let $\zeta\in L$ be a primitive
$\ell^s$-th root of unity. Let $\tau\colon K\to L$ be a $\mathbf
Q_\ell$-embedding and $\kappa\colon L\to B_{\mathrm{dR}}$ a map of $\mathbf
Q_\ell$-algebras with $\kappa(\tau_i(c))=c$ for $c\in K_0$,
$\kappa(\zeta)=\zeta_{\ell^s}$ and $\kappa(\tau(c))=c$ for $c\in K$, and let $t_s$
be as in
Lemma~\ref{asm-two-three-killing-wach-frame-exists-inputs-heights-full-evaluation}.
Let $V_\kappa=V\otimes_{L,\kappa}B_{\mathrm{dR}}$, the quotient of
$V\otimes_{\mathbf Q_\ell}B_{\mathrm{dR}}$ by the $L$-span $K_\kappa$ of the
elements $cy-\kappa(c)y$ with $c\in L$, where $\kappa(c)$ acts on the second
factor, and for $a\in\mathbf Z$ let $V_\kappa^a$ be the image of
$V\otimes\operatorname{Fil}^aB_{\mathrm{dR}}$. Then there is a map
$\Psi\colon N(T)\to V\otimes_{\mathbf Q_\ell}\operatorname{Fil}^0B_{\mathrm{dR}}$
such that:
\begin{enumerate}
\item the images of the $\Psi(x)$, $x\in N(T)$, generate $V_\kappa^0$ as a
  module over $\operatorname{Fil}^0B_{\mathrm{dR}}$;
\item for $x\in N(T)$ and $a\geq0$ with $\varepsilon_{i,\zeta}(x)(u)=0$ for
  $u<-a$, there is $c\in\mathbf Z$ such that for every $M\geq0$ the image in
  $V_\kappa$ of
  $t_s^a\Psi(x)-\sum_{0\leq r<M}t_s^r\,\iota\big(\varepsilon_{i,\zeta}(x)(r-a)\big)$
  lies in $V_\kappa^{M-c}$.
\end{enumerate}
\end{lemma}
\leanref{Development/AbelianVarietyTheory/Proofs/InertiaEigenvalues/WachImageDeRhamPlusComponentLevelStatement.lean}{GlobalTheory.T.wach_image_deRham_plus_component_level}
\leanref{Development/AbelianVarietyTheory/Proofs/InertiaEigenvalues/WachImageDeRhamPlusComponentLevel.lean}{GlobalTheory.NL.PerS.image}
\leanref{Development/AbelianVarietyTheory/Proofs/InertiaEigenvalues/WachImageDeRhamPlusComponentLevel.lean}{GlobalTheory.NL.PerS.core_gen}
\begin{proof}
Write $N=N(T)$, $\varepsilon=\varepsilon_{i,\zeta}$,
$\operatorname{Fil}^a=\operatorname{Fil}^aB_{\mathrm{dR}}$ and
$V^a=V\otimes\operatorname{Fil}^a$, and let $b\in B_{\mathrm{dR}}$ act on
$V\otimes B_{\mathrm{dR}}$ through the second factor; so $V^a_\kappa$ is the
image of $V^a$, and an element has image in $V^a_\kappa$ if and only if it lies
in $V^a+K_\kappa$. Let $p_1=[\varepsilon^{(1)}]-1\in A_{\mathrm{inf}}$, let $m$ be
as in
Lemma~\ref{asm-two-three-killing-wach-frame-exists-inputs-heights-full-image-sandwich},
where $\varpi=1\otimes p_1$, and let $\beta=\beta_T$ be as in
Definition~\ref{asm-two-three-killing-wach-frame-exists-inputs-heights-expansions}(6).

First, four facts.

(a) $A_{\mathrm{inf}}\subset\operatorname{Fil}^0$, and there is
$u\in\operatorname{Fil}^0$ with $p_1u=1$, by
Definition~\ref{asm-two-three-killing-wach-frame-exists-inputs-heights-expansions}(7).

(b) $\kappa(L)\subset\operatorname{Fil}^0$, by
Lemma~\ref{asm-two-three-killing-wach-frame-exists-inputs-heights-full-evaluation}(1).

(c) Multiplication by $b\in\operatorname{Fil}^0$ maps $K_\kappa$ into itself, as
$b(cz-\kappa(c)z)=c(bz)-\kappa(c)(bz)$, and maps every $V^a$ into itself.

(d) For $r\in R_A$ there is $\beta_r\in\operatorname{Fil}^0$ with
$\beta(ry)-\beta_r\beta(y)\in K_\kappa$ for every $y\in R_A\otimes_{\mathcal
O}T$. Write $r=\sum_jo_j\otimes w_j$ with $o_j\in\mathcal O$ and $w_j\in
A_{\mathrm{inf}}$, and put $\beta_r=\sum_j\kappa(o_j)w_j$, which lies in
$\operatorname{Fil}^0$ by (a) and (b). As $\beta$ is additive, it suffices to
check the claim for one term $r=o\otimes w$, with $\beta_r=\kappa(o)w$, and for
$y=(o'\otimes w')\otimes t$; then
$\beta(ry)-\beta_r\beta(y)=o(o't\otimes ww')-\kappa(o)(o't\otimes ww')$, which
lies in $K_\kappa$.

For $x\in N$ choose, by
Lemma~\ref{asm-two-three-killing-wach-frame-exists-inputs-heights-full-image-sandwich},
$Y(x)\in R_A\otimes_{\mathcal O}T$ whose image in $R\otimes_{\mathcal O}T$ is
$\varpi^mj_T(x)$, and put $\Psi(x)=u^m\beta(Y(x))$. By its formula and (a),
$\beta$ takes values in $V^0$, so $\Psi(x)\in V^0$ by (c). As $p_1u=1$,
$\beta(Y(x))=p_1^m\Psi(x)$.

(2): Let $x\in N$ and $a\geq0$ with $\varepsilon(x)(u')=0$ for $u'<-a$. By
Definition~\ref{asm-two-three-killing-wach-frame-exists-inputs-heights-expansions}(6),
with $k=m$ and $y=Y(x)$, there is $c\in\mathbf Z$ such that for every $M\geq0$
\[
 t_s^a\beta(Y(x))-p_1^m\sum_{0\leq r<M}t_s^r\,\iota\big(\varepsilon(x)(r-a)\big)
 \in V^{M-c}+K_\kappa.
\]
Multiplying by $u^m\in\operatorname{Fil}^0$, and using (c), $u^mp_1^m=1$ and
$u^m\beta(Y(x))=\Psi(x)$, gives
$t_s^a\Psi(x)-\sum_{0\leq r<M}t_s^r\,\iota(\varepsilon(x)(r-a))\in
V^{M-c}+K_\kappa$, which is (2).

(1): Let $G$ be the sum of $K_\kappa$ and the $L$-span of the $b\Psi(x)$ with
$b\in\operatorname{Fil}^0$ and $x\in N$. By (c), multiplication by
$b\in\operatorname{Fil}^0$ maps $G$ into itself. Every element of $G$ is
congruent modulo $K_\kappa$ to a sum $\sum_jb_j\Psi(x_j)$ with
$b_j\in\operatorname{Fil}^0$, since $c\,b\Psi(x)\equiv\kappa(c)b\Psi(x)$ modulo
$K_\kappa$ for $c\in L$, and $\kappa(c)b\in\operatorname{Fil}^0$ by (b). So (1)
holds once $V^0\subset G$.

For $r\in R_A$, $x\in N$ and $b\in\operatorname{Fil}^0$, $b\beta(rY(x))\in G$: by
(d) and (c) it is congruent modulo $K_\kappa$ to
$b\beta_r\beta(Y(x))=b\beta_rp_1^m\Psi(x)$, and $b\beta_rp_1^m\in
\operatorname{Fil}^0$ by (a). Let $S\subset R_A\otimes_{\mathcal O}T$ be the
additive subgroup generated by the $rY(x)$; it is an $\mathcal O$-submodule, as
$o\,rY(x)=((o\otimes1)r)Y(x)$, and as $\beta$ is additive, $b\beta(y)\in G$ for
$y\in S$ and $b\in\operatorname{Fil}^0$.

Let $t\in T$. By
Lemma~\ref{asm-two-three-killing-wach-frame-exists-inputs-heights-full-image-sandwich},
$\varpi^m(1\otimes t)=\sum_jo_jr_jj_T(x_j)$ with $o_j\in\mathcal O$,
$r_j\in R_A$ and $x_j\in N$. Multiplying by $\varpi^m$, the image of
$y=\sum_jo_jr_jY(x_j)\in S$ in $R\otimes T$ is $\varpi^{2m}(1\otimes t)$, which is
also the image of $(1\otimes p_1^{2m})\otimes t$. As $R_A\otimes T\to R\otimes T$
is injective, $y=(1\otimes p_1^{2m})\otimes t$, and $\beta(y)=t\otimes p_1^{2m}$.
Hence $t\otimes1=u^{2m}\beta(y)\in G$. As $T$ spans $V$ over $L$ and $G$ is an
$L$-subspace, $v\otimes1\in G$ for every $v\in V$, and so $v\otimes
b=b(v\otimes1)\in G$ for $b\in\operatorname{Fil}^0$. So $V^0\subset G$.
% Lean: PerS.image = PerS.core = core_gen with SandArb.sandwich (Lemma -full-image-sandwich),
% hX and hinj from IoXQ.main (io_X; inclR injective by HWL.TensorWitt.main, which the book takes
% from R_A being a subring of R, Definition -data), hι from EmbFil.emb_mem_fil_zero (fact (b)),
% and the datum's fields AI.ainf_image_mem_filtration_zero (of LocAinf, D.locA) and
% CmpSit.π₁_inv (of the comparison data): fact (a), stated as Definition -expansions(7), with
% the BC09 locators there for Fontaine's B_dR. core_gen: Ψ = u^m βdR(Y x), βdR_mem_FilV and
% mulVR_FilV (Ψ in V^0), Ct.levN_dR and mulVR_sup (2), and for (1) Gen, mulVR_gen, gen_form,
% act_βdR (fact (d)), hA', hB (the span S), ht1 (with TWincl_injective), hv. Facts (c) are
% mulVR_compKer and mulVR_FilV. core_gen with its helpers is about 420 lines; the rest of the
% closure (547 theorems, about 6,800 lines by the count used here) is the sandwich, the image
% of X and the injectivity of inclR, and emb_mem_fil_zero.
\end{proof}

\begin{lemma}\label{asm-two-three-killing-wach-frame-exists-inputs-heights-full}
Let $\ell$, $s$, $K$, $L$ and $D$ be as in
Definition~\ref{asm-two-three-killing-wach-frame-exists-inputs-heights-expansions},
with its notation, let $(V,T)$ be a lattice representation of $G_K$ over $L$
with $V$ crystalline, let $i\in\mathbf Z/n$, and let $\zeta\in L$ be a primitive
$\ell^s$-th root of unity. Then there is $m_0\geq0$ such that every sequence
$f\colon\mathbf Z\to D_i(V)$ with $f(u)=0$ for $u<m_0$ lies in
$\Lambda_{i,\zeta}(T)$.
\end{lemma}
\leanref{Development/AbelianVarietyTheory/Proofs/InertiaEigenvalues/WachExpansionsSpanFullLattice.lean}{GlobalTheory.NL.FullLevel.full_lat}
\begin{proof}
Write $\varepsilon=\varepsilon_{i,\zeta}$ and $\Lambda=\Lambda_{i,\zeta}(T)$.
Choose $\tau_s$, $\tau$ and $\kappa$ by
Lemma~\ref{asm-two-three-killing-wach-frame-exists-inputs-heights-full-embedding},
let $t_s$ and $\mathrm{ev}$ be as in
Lemma~\ref{asm-two-three-killing-wach-frame-exists-inputs-heights-full-evaluation}
for $\kappa$, and $V_\kappa$, $V_\kappa^a$ and $\Psi$ as in
Lemma~\ref{asm-two-three-killing-wach-frame-exists-inputs-heights-full-image}.
Put $r=\dim_LV$, which is $\dim_LD_i(V)$ by
Lemma~\ref{asm-two-three-killing-wach-frame-exists-inputs-heights-letters}; fix
an $L$-basis $v_1,\ldots,v_r$ of $V$ and an $L$-basis $d_1,\ldots,d_r$ of
$D_i(V)$. For $y\in V\otimes_{\mathbf Q_\ell}B_{\mathrm{dR}}$ let
$p_1(y),\ldots,p_r(y)\in B_{\mathrm{dR}}$ be the coordinates of its image in
$V_\kappa=\bigoplus_aB_{\mathrm{dR}}v_a$. Each $p_a$ is linear for
$B_{\mathrm{dR}}$ acting on the second factor, $p_a(cy)=\kappa(c)p_a(y)$ for
$c\in L$, $p_a$ vanishes on $K_\kappa$, and $p_a$ maps $V\otimes
\operatorname{Fil}^m$ into $\operatorname{Fil}^m$, as
$\kappa(L)\subset\operatorname{Fil}^0$.

First, the vectors $g(x)=(p_a(\Psi(x)))_a\in B_{\mathrm{dR}}^r$, $x\in N(T)$,
span $B_{\mathrm{dR}}^r$. Indeed, by
Lemma~\ref{asm-two-three-killing-wach-frame-exists-inputs-heights-full-image}(1),
$v_a\otimes1\equiv\sum_jb_j\Psi(x_j)$ modulo $K_\kappa$ for some $b_j\in
B_{\mathrm{dR}}$ and $x_j\in N(T)$; applying the $p_{a'}$ gives the $a$-th unit
vector as $\sum_jb_jg(x_j)$. So there are $x_1,\ldots,x_r\in N(T)$ such that
the matrix $G=(g(x_j)_a)_{j,a}$ has nonzero determinant: otherwise the
determinant, an alternating $r$-linear form, would vanish on a spanning set of
$B_{\mathrm{dR}}^r$, and so everywhere.

Second, choose $w\geq0$ with $\varepsilon(x_j)(u)=0$ for every $j$ and $u<-w$,
and let $U_{ja}\in L[[X]]$ be the series with
$\varepsilon(x_j)(n-w)=\sum_a\operatorname{coeff}_n(U_{ja})\,d_a$ for $n\geq0$.
Put $P_{aa'}=p_{a'}(\iota(d_a))$. Then
$t_s^wg(x_j)_{a'}=\sum_a\mathrm{ev}(U_{ja})P_{aa'}$. Indeed, by
Lemma~\ref{asm-two-three-killing-wach-frame-exists-inputs-heights-full-image}(2)
there is $c$ with
\[
 t_s^wg(x_j)_{a'}-\sum_{n<M}t_s^n\sum_a\kappa(\operatorname{coeff}_nU_{ja})P_{aa'}
 \in\operatorname{Fil}^{M-c}
\]
for every $M\geq0$; by
Lemma~\ref{asm-two-three-killing-wach-frame-exists-inputs-heights-full-evaluation}(2),
$\mathrm{ev}(U_{ja})-\sum_{n<M}\kappa(\operatorname{coeff}_nU_{ja})t_s^n\in
\operatorname{Fil}^M$, and every $P_{aa'}$ lies in some $\operatorname{Fil}^b$.
So the difference of the two sides lies in $\operatorname{Fil}^{M-c'}$ for a
fixed $c'$ and every $M$, and it is $0$ by
Lemma~\ref{asm-two-three-killing-wach-frame-exists-inputs-heights-full-evaluation}(1).

Third, $\det U\neq0$ in $L[[X]]$, for $U=(U_{ja})$. By the second step,
$t_s^{wr}\det G=\mathrm{ev}(\det U)\det P$, and $t_s$ is a unit while
$\det G\neq0$.

Fourth, let $m_0$ be the order of $\det U$, so that $\det U=X^{m_0}e_0$ with
$e_0$ a unit, and let $f\colon\mathbf Z\to D_i(V)$ vanish below $m_0$. Let
$F_a\in L[[X]]$ be the series whose coefficient of $X^n$ is the $a$-th
coordinate of $f(n-w)$ in the basis $d$. As $f(n-w)=0$ for $n<m_0$, $X^{m_0}$
divides $F_a$, so $F_a=\det U\cdot F'_a$ with $F'_a\in L[[X]]$. Let
$e=\operatorname{adj}(U^{\mathrm t})F'$, so that $U^{\mathrm t}e=F$, that is
$\sum_jU_{ja}e_j=F_a$ for every $a$. Then $f=\sum_je_j*\varepsilon(x_j)$: both
sides vanish at $u<-w$, and at $n-w$ with $n\geq0$ the right side is
\[
 \sum_j\sum_{p+q=n}\operatorname{coeff}_p(e_j)\,\varepsilon(x_j)(q-w)
 =\sum_a\operatorname{coeff}_n\Big(\sum_jU_{ja}e_j\Big)d_a
 =\sum_a\operatorname{coeff}_n(F_a)\,d_a=f(n-w).
\]
So $f\in\Lambda$.
% Lean: full_lat = FullLevel.full_of_span (form, form_ker, form_fil, exists_det_ne_zero,
% exists_coord, eq_sum_ev, full_of_det) with Points.points (h1, h2), the image Ψ
% (wach_image_deRham_plus_component_level), EvdR.exists_ev_dR and the dimension of
% tied_letter_component_finrank. About 250 lines besides the lemmas cited.
\end{proof}

\begin{lemma}\label{asm-two-three-killing-wach-frame-exists-inputs-heights-basis}
Let $\ell$, $s$, $K$, $L$ and $D$ be as in
Definition~\ref{asm-two-three-killing-wach-frame-exists-inputs-heights-expansions},
with its notation, let $(V,T)$ be a lattice representation of $G_K$ over $L$
with $V$ crystalline, let $i\in\mathbf Z/n$, and let $\zeta\in L$ be a primitive
$\ell^{j_0}$-th root of unity for some $j_0\geq0$. Suppose that there is
$m_0\geq0$ such that every sequence $f\colon\mathbf Z\to D_i(V)$ with $f(u)=0$
for $u<m_0$ lies in $\Lambda_{i,\zeta}(T)$. Let $b_1,\ldots,b_r$ be an $A$-basis
of $N(T)$, and for $y\in N(T)$ write $y=\sum_jc_j(y)b_j$ with $c_j(y)\in A$.
Then:
\begin{enumerate}
\item if $\eta_1,\ldots,\eta_r\in L[[X]]$ and
  $\sum_j\eta_j*\varepsilon_{i,\zeta}(b_j)=0$, then every $\eta_j$ is $0$;
\item every $f\in\Lambda_{i,\zeta}(T)$ is $\sum_j\eta_j*\varepsilon_{i,\zeta}(b_j)$
  for some $\eta_1,\ldots,\eta_r\in L[[X]]$;
\item for $y\in N(T)$ and $k\geq0$, there is $f\in\Lambda_{i,\zeta}(T)$ with
  $\varepsilon_{i,\zeta}(y)(u)=f(u-k)$ for every $u$ if and only if $X^k$
  divides $\lambda_{i,\zeta}(c_j(y))$ for every $j$.
\end{enumerate}
\end{lemma}
\leanref{Development/AbelianVarietyTheory/Proofs/InertiaEigenvalues/WachExpansionBasisDivisibility.lean}{GlobalTheory.NL.LevInj.expansion_basis}
\begin{proof}
Write $\lambda=\lambda_{i,\zeta}$, $\varepsilon=\varepsilon_{i,\zeta}$ and
$\Lambda=\Lambda_{i,\zeta}(T)$. Let $S$ be the set of sequences $\mathbf Z\to
V\otimes B_{\mathrm{cris}}$ that vanish below some integer. Identifying $f\in S$
with the Laurent series $\sum_uf(u)X^u$, $S$ is a vector space over the field
$L((X))$, and for $\eta\in L[[X]]$ the product $\eta f$ is $\eta*f$. So $*$ makes
$S$ a module over $L[[X]]$, and $S$ has no $L[[X]]$-torsion.

First, $\varepsilon(ax)=\lambda(a)*\varepsilon(x)$ for $a\in A$ and $x\in N(T)$,
by Definition~\ref{asm-two-three-killing-wach-frame-exists-inputs-heights-expansions}(3),
as $\iota$ is injective \cite[Theorem~9.1.5]{BC09}. Hence, for $y\in N(T)$,
$\varepsilon(y)=\sum_j\lambda(c_j(y))*\varepsilon(b_j)$.

Second, every $f\in\Lambda$ is $\sum_je_j*\varepsilon(b_j)$ with $e_j\in
L[[X]]$. The sums of this form make up an $L$-subspace, since
$cf=c*f$ for $c\in L$, and they contain the generators of $\Lambda$:
$\eta*\varepsilon(x)=\sum_j(\eta\lambda(c_j(x)))*\varepsilon(b_j)$ by the first
step.

Third, $r\leq\dim_LD_i(V)$. By
Lemma~\ref{asm-two-three-killing-wach-frame-exists-inputs-basis},
$b_1,\ldots,b_r$ is an $E$-basis of $D(T)$, which is free of rank the rank of
$T$ over $\mathcal O$
(Definition~\ref{asm-two-three-killing-wach-frame-exists-data}(4)); a basis of
$T$ over $\mathcal O$ is linearly independent over the fraction field $L$, so
this rank is at most $\dim_LV$, which is $\dim_LD_i(V)$ by
Lemma~\ref{asm-two-three-killing-wach-frame-exists-inputs-heights-letters}.

Fourth, the $\varepsilon(b_j)$ are linearly independent over $L[[X]]$. Choose
$d_1,\ldots,d_r\in D_i(V)$, linearly independent over $L$, by the third step.
For each $a$, the sequence $\delta_a$ with value $d_a$ at $m_0$ and $0$
elsewhere lies in $\Lambda$ by hypothesis, so by the second step
$\delta_a=\sum_jE_{aj}*\varepsilon(b_j)$ with $E_{aj}\in L[[X]]$. The
$\delta_a$ are linearly independent over $L[[X]]$: if
$\sum_ag_a*\delta_a=0$, its value at $m_0+t$ is $\sum_ag_{a,t}d_a=0$, where
$g_{a,t}$ is the coefficient of $X^t$ in $g_a$, so every $g_{a,t}$ is $0$. Thus
the $L[[X]]$-span of the $r$ elements $\varepsilon(b_j)$ contains $r$ linearly
independent elements. Over the field $L((X))$, $r$ vectors whose span contains
$r$ linearly independent vectors are linearly independent; as $S$ has no
torsion, a relation over $L[[X]]$ is one over $L((X))$, and the claim follows.

Now (1) is the fourth step and (2) is the second. For (3), suppose that
$\varepsilon(y)(u)=f(u-k)$ for every $u$, with $f\in\Lambda$; then
$\varepsilon(y)=X^k*f$. By the second step $f=\sum_je_j*\varepsilon(b_j)$, so
$\sum_j\big(\lambda(c_j(y))-X^ke_j\big)*\varepsilon(b_j)=0$ by the first step,
and the fourth step gives $\lambda(c_j(y))=X^ke_j$ for every $j$. Conversely,
if $\lambda(c_j(y))=X^ke_j$ with $e_j\in L[[X]]$ for every $j$, then
$f=\sum_je_j*\varepsilon(b_j)$ lies in $\Lambda$, and
$X^k*f=\sum_j\lambda(c_j(y))*\varepsilon(b_j)=\varepsilon(y)$ by the first
step; that is, $\varepsilon(y)(u)=f(u-k)$ for every $u$.
% Lean: expansion_basis = LevInj.core' (the Hahn-module embedding of S, coeff_smul_hahn) and
% LevInj.core (hsum, hy, h2, hli by exists_linearIndependent_of_le_finrank, conv_single and
% indep_of_span through linearIndependent_iff_card_le_finrank_span), with hr from card_eq_rank
% (DD_free, N_base), rank_le and the dimension of Lemma -heights-letters, and lift_smul for the
% first step. About 250 lines; the rest of the closure is Lemma -heights-letters.
\end{proof}

\begin{lemma}\label{asm-two-three-killing-wach-frame-exists-inputs-heights-standard-one}
Let $\ell$, $s$, $K$, $L$ and $D$ be as in
Definition~\ref{asm-two-three-killing-wach-frame-exists-inputs-heights-expansions},
with its notation, let $(V,T)$ be a lattice representation of $G_K$ over $L$
with $V$ crystalline, and let $i\in\mathbf Z/n$. Then $\Lambda_{i,1}(T)$
consists of the sequences $f\colon\mathbf Z\to D_i(V)$ with $f(u)=0$ for $u<0$.
\end{lemma}
\leanref{Development/AbelianVarietyTheory/Proofs/InertiaEigenvalues/WachLatticeEqStdPowerSeriesAtZeroStatement.lean}{GlobalTheory.T.wach_lattice_eq_std_powerSeries_at_zero}
\leanref{Development/AbelianVarietyTheory/Proofs/InertiaEigenvalues/WachLatticeEqStdPowerSeriesAtZero.lean}{GlobalTheory.NL.AtZero.lat_one_std}
\begin{proof}
Write $\varepsilon=\varepsilon_{i,1}$, $\Lambda=\Lambda_{i,1}(T)$, $N=N(T)$, and
$\delta_ud$ as in
Lemma~\ref{asm-two-three-killing-wach-frame-exists-inputs-heights-standard-graded}.
Let $\gamma_0$ be as in
Definition~\ref{asm-two-three-killing-wach-frame-exists-data}(6) and
$c=\chi_\Gamma(\gamma_0)$.

First, there is $m_0\geq0$ such that every $f\colon\mathbf Z\to D_i(V)$ vanishing
below $m_0$ lies in $\Lambda$. Let $\zeta\in L$ be a primitive $\ell^s$-th root of
unity. By
Lemma~\ref{asm-two-three-killing-wach-frame-exists-inputs-heights-full} this
holds for $\Lambda_{i+s,\zeta}(T)$, and $s$ applications of
Lemma~\ref{asm-two-three-killing-wach-frame-exists-inputs-heights-standard-descent}
carry it to $\Lambda_{i+s-j,\zeta^{\ell^j}}(T)$ for $j=1,\ldots,s$, and so to
$\Lambda_{i,1}(T)$. By
Lemma~\ref{asm-two-three-killing-wach-frame-exists-inputs-heights-standard-graded},
$f\in\Lambda$ if and only if $f$ vanishes below some integer and every
$\delta_uf(u)$ lies in $\Lambda$.

Second, there is $m_1$ with $\varepsilon(x)(u)=0$ for every $x\in N$ and $u<m_1$:
take the least of the bounds for a finite set of generators of $N$ over $A$; it
works for all of $N$, since $\varepsilon(ax)=\lambda_{i,1}(a)*\varepsilon(x)$.

Third, let $S\subset V\otimes B_{\mathrm{cris}}$ be an $L$-subspace and $k\neq0$
an integer with $\varepsilon(y)(j)\in S$ for every $y\in N$ and $j<k$. Then
$\varepsilon(x)(k)\in S$ for every $x\in N$. Indeed, $\gamma_0x\in N$ and
$\gamma_0x=x+\pi y$ with $y\in N$
(Definition~\ref{asm-two-three-killing-wach-frame-exists-data}(4)). By
Definition~\ref{asm-two-three-killing-wach-frame-exists-inputs-heights-expansions}(5)
with $\zeta=1$, $j=0$ and $a=0$,
$\varepsilon(\gamma_0x)(k)=c^k\varepsilon(x)(k)$. Also
$\varepsilon(\gamma_0x)(k)=\varepsilon(x)(k)+(\lambda_{i,1}(\pi)*\varepsilon(y))(k)$,
and $\lambda_{i,1}(\pi)=\exp(X)-1$ has constant term $0$, so the last term is a
combination of the $\varepsilon(y)(k-t)$ with $t\geq1$, which lie in $S$. Hence
$(c^k-1)\varepsilon(x)(k)\in S$, and $c^k\neq1$ by
Lemma~\ref{asm-two-three-killing-wach-frame-exists-inputs-heights-standard-generator}.

Fourth, $\varepsilon(y)(j)=0$ for $y\in N$ and $j<0$: apply the third step with
$S=0$ to $k=m_1,m_1+1,\ldots,-1$ in turn. Let $S_0$ be the $L$-span of the
$\varepsilon(x)(0)$, $x\in N$. Then every $\varepsilon(y)(j)$ lies in $S_0$: for
$j\leq0$ by what was just shown, and for $j\geq1$ by the third step applied to
$k=1,2,\ldots,j$ in turn.

Fifth, $D_i(V)\subset S_0$: for $d\in D_i(V)$, $\delta_{m_0}d\in\Lambda$ by the
first step, and the elements of $\Lambda$ take values in $S_0$, by the fourth
step. And $\delta_nd\in\Lambda$ for $d\in S_0$ and $n\geq0$: for
$d=\varepsilon(x)(0)$, the sequence $X^n*\varepsilon(x)$ lies in $\Lambda$ and
has value $d$ at $n$, so $\delta_nd\in\Lambda$ by the first step, and the
general case follows by linearity.

Finally, an element of $\Lambda$ takes values in $D_i(V)$ and vanishes below $0$,
by the fourth step. Conversely, let $f\colon\mathbf Z\to D_i(V)$ vanish below $0$.
For $u<0$, $\delta_uf(u)=0$, and for $u\geq0$, $f(u)\in D_i(V)\subset S_0$, so
$\delta_uf(u)\in\Lambda$ by the fifth step. By the first step, $f\in\Lambda$.
% Lean: lat_one_std: down (wachExpansionLattice_full_descends s times from
% wach_expansions_span_full_lattice at a primitive l^s-th root from exists_isPt_primitiveRoot),
% wachExpansionLattice_graded_of_level_le, uniform_supp, AtZero.lower (N_Γ, N_π,
% CmpTie.levB_Γ with ζ = ζ' = 1, j = a = 0, Model.pi_van_one, chi_gam_zpow_ne_one) and
% std_of_lower (up, hneg, hS, hDS, hsingle). About 300 lines besides the lemmas cited.
\end{proof}

\begin{lemma}\label{asm-two-three-killing-wach-frame-exists-inputs-heights-standard}
Let $\ell$, $s$, $K$, $L$ and $D$ be as in
Definition~\ref{asm-two-three-killing-wach-frame-exists-inputs-heights-expansions},
with its notation, let $(V,T)$ be a lattice representation of $G_K$ over $L$
with $V$ crystalline, let $i\in\mathbf Z/n$, and let $\zeta\in L$ with
$\zeta^{\ell^{s-1}}=1$. Then $\Lambda_{i,\zeta}(T)$ consists of the sequences
$f\colon\mathbf Z\to D_i(V)$ with $f(u)=0$ for $u<0$.
\end{lemma}
\leanref{Development/AbelianVarietyTheory/Proofs/InertiaEigenvalues/WachLatticeExpansionsStdBelowLevel.lean}{GlobalTheory.NL.LatLow.lat_low}
\begin{proof}
For $j\geq0$ we show, by induction on $j$, that the assertion holds for every
$i\in\mathbf Z/n$ and every $\zeta$ with $\zeta^{\ell^j}=1$ and
$\zeta^{\ell^{s-1}}=1$; the case $j=s-1$ is the lemma. For $j=0$, $\zeta=1$, and
this is
Lemma~\ref{asm-two-three-killing-wach-frame-exists-inputs-heights-standard-one}.

Let $j\geq0$, and let $\zeta^{\ell^{j+1}}=1$ and $\zeta^{\ell^{s-1}}=1$. Then
$\zeta^\ell$ satisfies the conditions for $j$, so by induction
$\Lambda'=\Lambda_{i-1,\zeta^\ell}(T)$ consists of the sequences in $D_{i-1}(V)$
that vanish below $0$. Write $\lambda=\lambda_{i,\zeta}$,
$\varepsilon=\varepsilon_{i,\zeta}$, $\varepsilon'=\varepsilon_{i-1,\zeta^\ell}$,
$\Lambda=\Lambda_{i,\zeta}(T)$, $N=N(T)$, and $\Theta$ as in the proof of
Lemma~\ref{asm-two-three-killing-wach-frame-exists-inputs-heights-standard-descent}.
By Lemma~\ref{asm-two-three-killing-wach-frame-exists-inputs-heights-standard-values},
$\lambda(Q)$ is a unit of $L[[X]]$.

(a) For $x\in N$ and $m$ with $Q^m\varphi_D(x)\in N$,
$\varepsilon(Q^m\varphi_D(x))=\lambda(Q^m)*\Theta(\varepsilon'(x))$
(Definition~\ref{asm-two-three-killing-wach-frame-exists-inputs-heights-expansions}(4)),
and it vanishes below $0$, as $\varepsilon'(x)\in\Lambda'$ does.

(b) Every $\varepsilon(y)$, $y\in N$, vanishes below $0$. By
Definition~\ref{asm-two-three-killing-wach-frame-exists-data}(4),
$Q^jy=\sum_ra_r\varphi_D(x_r)$ for some $j$ and finitely many $a_r\in A$ and
$x_r\in N$; choose $m$ with $Q^m\varphi_D(x_r)\in N$ for every $r$. Then
$\lambda(Q)^{m+j}*\varepsilon(y)=\varepsilon(Q^{m+j}y)=\sum_r\lambda(a_r)*
\varepsilon(Q^m\varphi_D(x_r))$ vanishes below $0$ by (a), and so does
$\varepsilon(y)=\lambda(Q)^{-m-j}*\big(\lambda(Q)^{m+j}*\varepsilon(y)\big)$.

(c) $\Theta(\Lambda')\subset\Lambda$. For a generator $\eta*\varepsilon'(x)$
of $\Lambda'$, choose $m$ with $Q^m\varphi_D(x)\in N$; then, with (a) of the
proof of
Lemma~\ref{asm-two-three-killing-wach-frame-exists-inputs-heights-standard-descent},
$\Theta(\eta*\varepsilon'(x))=\eta_\ell*\Theta(\varepsilon'(x))
=(\eta_\ell\lambda(Q^m)^{-1})*\varepsilon(Q^m\varphi_D(x))\in\Lambda$.

Now an element of $\Lambda$ takes values in $D_i(V)$
(Definition~\ref{asm-two-three-killing-wach-frame-exists-inputs-heights-expansions}(3))
and vanishes below $0$ by (b). Conversely, let $f\colon\mathbf Z\to D_i(V)$
vanish below $0$. By
Lemma~\ref{asm-two-three-killing-wach-frame-exists-inputs-heights-letters},
$\varphi$ maps $D_{i-1}(V)$ onto $D_i(V)$; choose $d_u\in D_{i-1}(V)$ with
$\varphi(d_u)=f(u)$, and $d_u=0$ for $u<0$. The sequence $g(u)=\ell^{-u}d_u$
lies in $\Lambda'$, and $\Theta g=f$; so $f\in\Lambda$ by (c).
% Lean: lat_low (induction on j, with wach_lattice_eq_std_powerSeries_at_zero for ζ = 1) and
% LatLow.std_step (hspan, hv, hmap, tw_conv, conv_mul) with the third part of
% wachLevel_expansion_at_root_of_unity, the first part of
% crystallineFrobenius_bijOn_letters_directSum, N_φ, N_φinv and CmpTie.levB_φ. About 250
% lines besides the lemmas cited.
\end{proof}

\begin{lemma}\label{asm-two-three-killing-wach-frame-exists-inputs-heights-filtration-component-split}
Let $F$ be a field, $E$ a finite extension of $F$, and $L'$ a field containing
$F$ such that there are $[E:F]$ maps of $F$-algebras $E\to L'$. Let $M$ be an
$F$-vector space, and let $\alpha\colon L'\to\operatorname{End}_F(M)$ and
$\beta\colon E\to\operatorname{End}_F(M)$ be maps of $F$-algebras such that every
$\alpha(a)$ commutes with every $\beta(c)$. Then there are $F$-linear maps
$e_\sigma\colon M\to M$, one for each map of $F$-algebras $\sigma\colon E\to L'$,
such that:
\begin{enumerate}
\item $\sum_\sigma e_\sigma(x)=x$ for $x\in M$;
\item $\beta(c)e_\sigma(x)=\alpha(\sigma(c))e_\sigma(x)$ for every $\sigma$,
  $c\in E$ and $x\in M$;
\item every $e_\sigma$ commutes with every $\alpha(a)$ and every $\beta(c)$;
\item every $e_\sigma$ maps into itself each additive subgroup of $M$ that every
  $\alpha(a)$ and every $\beta(c)$ map into itself.
\end{enumerate}
\end{lemma}
\leanref{Development/AbelianVarietyTheory/Proofs/InertiaEigenvalues/DeRhamComponentWachLetterBasisFiltration.lean}{GlobalTheory.NL.CompD.exists_proj}
\leanref{Development/AbelianVarietyTheory/Proofs/FoundationProofs/FinrankDdRFinsum.lean}{GlobalTheory.DeRhamSplit.evalEmb_bijective}
\begin{proof}
Let $\Sigma$ be the set of maps of $F$-algebras $E\to L'$. The map of
$L'$-algebras $\theta\colon L'\otimes_FE\to\prod_{\sigma\in\Sigma}L'$,
$a\otimes c\mapsto(a\sigma(c))_\sigma$, is surjective: the elements of $\Sigma$
are distinct characters of $E^\times$, hence linearly independent over $L'$
\cite[Chapter~VI, Theorem~4.1]{Lan02}, so the vectors $(\sigma(c))_\sigma$,
$c\in E$, span $\prod_\sigma L'$ over $L'$. Both sides have dimension $[E:F]$
over $L'$, so $\theta$ is an isomorphism. Let $\epsilon_\sigma\in L'\otimes_FE$
be the element that $\theta$ maps to the $\sigma$-th unit vector. Then
$\sum_\sigma\epsilon_\sigma=1$, and $(1\otimes c)\epsilon_\sigma=
(\sigma(c)\otimes1)\epsilon_\sigma$ for $c\in E$, since both sides have the same
image under $\theta$.

As the actions commute, $\Psi(a\otimes c)=\alpha(a)\beta(c)$ defines a map of
$F$-algebras $\Psi\colon L'\otimes_FE\to\operatorname{End}_F(M)$. Put
$e_\sigma=\Psi(\epsilon_\sigma)$. Applying $\Psi$ to the two identities gives (1)
and (2). (3) holds as $L'\otimes_FE$ is commutative, $\alpha(a)=\Psi(a\otimes1)$
and $\beta(c)=\Psi(1\otimes c)$. (4) holds as $\epsilon_\sigma$ is a sum of
elements $a\otimes c$, and $\Psi(a\otimes c)=\alpha(a)\beta(c)$.
% Lean: CompD.exists_proj (Ψ = Algebra.TensorProduct.lift α β, e σ = Ψ(Φ⁻¹(Pi.single σ 1)),
% hsum, hmul, and the induction on tensors for (4)) with DeRhamSplit.evalEmb_bijective
% (linearIndependent_monoidHom for the span; injective_iff_surjective_of_finrank_eq_finrank).
% About 90 lines.
\end{proof}

\begin{lemma}\label{asm-two-three-killing-wach-frame-exists-inputs-heights-filtration-component-descent}
Let $\ell$, $s$, $K$, $L$ and $D$ be as in
Definition~\ref{asm-two-three-killing-wach-frame-exists-inputs-heights-expansions},
with its notation, and let $(V,T)$ be a lattice representation of $G_K$ over
$L$. Let $\tau\colon K\to L$ be a $\mathbf Q_\ell$-embedding, $\kappa\colon L\to
B_{\mathrm{dR}}$ a map of $\mathbf Q_\ell$-algebras with $\kappa(\tau(c))=c$ for
$c\in K$, and $K_\kappa$ as in
Lemma~\ref{asm-two-three-killing-wach-frame-exists-inputs-heights-full-image}.
Let $a\in\mathbf Z$ and $y\in D_\tau$. If $y\in
V\otimes\operatorname{Fil}^aB_{\mathrm{dR}}+K_\kappa$, then $y\in
V\otimes\operatorname{Fil}^aB_{\mathrm{dR}}$.
\end{lemma}
\leanref{Development/AbelianVarietyTheory/Proofs/InertiaEigenvalues/DeRhamComponentWachLetterBasisFiltration.lean}{GlobalTheory.NL.CompD.mem_FilV_local}
\leanref{Development/AbelianVarietyTheory/Proofs/InertiaEigenvalues/DeRhamComponentWachLetterBasisFiltration.lean}{GlobalTheory.NL.CompD.mem_FilV_of_comp}
\leanref{Development/AbelianVarietyTheory/Proofs/InertiaEigenvalues/DeRhamComponentWachLetterBasisFiltration.lean}{GlobalTheory.NL.CompD.exists_factor}
\leanref{Development/AbelianVarietyTheory/Proofs/InertiaEigenvalues/DeRhamComponentWachLetterBasisFiltration.lean}{GlobalTheory.NL.CompD.exists_conj}
\begin{proof}
Write $V^a=V\otimes\operatorname{Fil}^aB_{\mathrm{dR}}$, and let $G_K$ act on
$V\otimes B_{\mathrm{dR}}$ by $g\otimes g$, so that $D_\tau$ lies in
$D_{\mathrm{dR}}(V)=(V\otimes B_{\mathrm{dR}})^{G_K}$. Recall that
$B_{\mathrm{dR}}$ is the fraction field of the discrete valuation ring
$B^+_{\mathrm{dR}}=\operatorname{Fil}^0B_{\mathrm{dR}}$, with
$\operatorname{Fil}^aB_{\mathrm{dR}}$ the $a$-th power of its maximal ideal, and
that $G_K$ acts on $B_{\mathrm{dR}}$ by ring automorphisms that preserve every
$\operatorname{Fil}^aB_{\mathrm{dR}}$
\cite[Proposition~4.4.6 and Definition~4.4.7]{BC09}; in particular
$\operatorname{Fil}^0B_{\mathrm{dR}}\cdot\operatorname{Fil}^aB_{\mathrm{dR}}
\subset\operatorname{Fil}^aB_{\mathrm{dR}}$. Recall also that
$B^+_{\mathrm{dR}}$ contains a unique copy of $\bar{\mathbf Q}_\ell$ over
$\mathbf Q_\ell$, on which $G_K$ acts by its Galois action
\cite[Lemma~4.4.10]{BC09}. For a map of $\mathbf Q_\ell$-algebras $\mu\colon L\to
B_{\mathrm{dR}}$ let $K_\mu$ be the $L$-span of the elements $cz-\mu(c)z$ with
$c\in L$ and $z\in V\otimes B_{\mathrm{dR}}$, where $c$ acts on $V$ and $\mu(c)$
on $B_{\mathrm{dR}}$.

First, there is a $\mathbf Q_\ell$-embedding $\sigma_0\colon L\to\bar{\mathbf
Q}_\ell$ with $\kappa(c)=\sigma_0(c)$ in $B_{\mathrm{dR}}$ for every $c\in L$.
Apply
Lemma~\ref{asm-two-three-killing-wach-frame-exists-inputs-heights-filtration-component-split}
with $F=\mathbf Q_\ell$, $E=L$, $L'=\bar{\mathbf Q}_\ell$, into which $L$ has
$[L:\mathbf Q_\ell]$ embeddings, and $M=B_{\mathrm{dR}}$, on which $b\in\bar{\mathbf
Q}_\ell$ and $c\in L$ act by multiplication by $b$ and by $\kappa(c)$. By (1) of
that lemma $\sum_\sigma e_\sigma(1)=1$, so $e_{\sigma_0}(1)\neq0$ for some
$\sigma_0$, and by (2) $\kappa(c)e_{\sigma_0}(1)=\sigma_0(c)e_{\sigma_0}(1)$. As
$B_{\mathrm{dR}}$ is a field, $\kappa(c)=\sigma_0(c)$.

Second, apply the same lemma with $F=\mathbf Q_\ell$, $E=L$,
$L'=\bar{\mathbf Q}_\ell$ and $M=V\otimes B_{\mathrm{dR}}$, with
$\bar{\mathbf Q}_\ell$ acting by multiplication on the second factor and $L$
through $V$, and let $e_\sigma$, for the $\mathbf Q_\ell$-embeddings
$\sigma\colon L\to\bar{\mathbf Q}_\ell$, be the maps it gives. Every $e_\sigma$
maps $V^a$ into itself, by (4), as $V^a$ is stable under $L$ and under
$\bar{\mathbf Q}_\ell\subset\operatorname{Fil}^0B_{\mathrm{dR}}$. Every
$e_\sigma$ vanishes on $K_\sigma$: by (3) and (2),
$e_\sigma(cz)=ce_\sigma(z)=\sigma(c)e_\sigma(z)=e_\sigma(\sigma(c)z)$.

Write $y=u+w$ with $u\in K_\kappa$ and $w\in V^a$. By (1), $y=\sum_\sigma
e_\sigma(y)$, so it suffices to show that every $e_\sigma(y)$ lies in $V^a$.

Third, let $\sigma$ be an embedding with $\sigma(\tau(c))=c$ for every $c\in K$.
Extend $\sigma_0$ and $\sigma$ to automorphisms $g_0$ and $g_1$ of
$\bar{\mathbf Q}_\ell$, and put $g=g_1g_0^{-1}$. Then $g(\sigma_0(c))=\sigma(c)$
for $c\in L$, and $g(c)=g(\sigma_0(\tau(c)))=\sigma(\tau(c))=c$ for $c\in K$, as
$\sigma_0(\tau(c))=\kappa(\tau(c))=c$; so $g\in G_K$. The action of $g$ fixes
$y\in D_{\mathrm{dR}}(V)$, maps $V^a$ into itself, and maps $K_\kappa$ into
$K_\sigma$: as $G_K$ acts $L$-linearly on $V$,
$g(cz-\kappa(c)z)=c\,g(z)-\sigma(c)g(z)$. Hence
$e_\sigma(y)=e_\sigma(g(y))=e_\sigma(g(u))+e_\sigma(g(w))=e_\sigma(g(w))$, which
lies in $V^a$.

Fourth, let $\sigma$ be an embedding with $\sigma(\tau(k))\neq k$ for some
$k\in K$. As $y\in D_\tau$, $ky=\tau(k)y$, with $k$ acting on $B_{\mathrm{dR}}$
and $\tau(k)$ on $V$. By (2) and (3), with $\sigma(\tau(k))$ and $k$ acting on
$B_{\mathrm{dR}}$,
\[
 \sigma(\tau(k))e_\sigma(y)=\tau(k)e_\sigma(y)=e_\sigma(\tau(k)y)=e_\sigma(ky)
 =ke_\sigma(y).
\]
As $\sigma(\tau(k))-k$ is a nonzero element of $\bar{\mathbf Q}_\ell$, it is a
unit of $B_{\mathrm{dR}}$, and $e_\sigma(y)=0$.
% Lean: mem_FilV_local = mem_FilV_of_comp with E = Om l (the algebraic closure; hE =
% AlgHom.card), exists_factor (exists_proj, IsDomain BdR from RegularRing.isField_of_tied),
% exists_conj (Points.exists_aut), P.PD_actdR and P.actQdR_Om (the action on Q_l-bar inside
% BdR), act_compKer, act_FilV, om_mul_mem_Fil (Q_l-bar times Fil^a in Fil^a, through P.PD_Fil)
% and mu_FilV; the steps are hk (third) and h6, h7 (fourth). About 220 lines. The closure of
% isField_of_tied (224 theorems, about 2,500 lines) proves that the datum's BdR is a field:
% bdR_isField for the Lean's own Per.BdR, carried to the datum's ring by the ring isomorphism
% eD of the tie D.locT; this is not a field of the datum, and the book cites BC09 Proposition
% 4.4.6 and Definition 4.4.7 for it. The other recalled facts are fields of the datum
% (LocPrime): P.actQdR_Om (the action on Q_l-bar), P.PD_actdR with PeriodData's Fil stable
% under G_K, and Fil^a a K-submodule with P.PD_Fil (om_mul_mem_Fil); for Fontaine's B_dR they
% are BC09 Lemma 4.4.10 and Definition 4.4.7.
\end{proof}

\begin{lemma}\label{asm-two-three-killing-wach-frame-exists-inputs-heights-filtration-component}
Let $\ell$, $s$, $K$, $L$ and $D$ be as in
Definition~\ref{asm-two-three-killing-wach-frame-exists-inputs-heights-expansions},
with its notation, let $(V,T)$ be a lattice representation of $G_K$ over $L$
with $V$ crystalline, let $i\in\mathbf Z/n$, and let $\zeta\in L$ be a primitive
$\ell^s$-th root of unity. Let $\tau\colon K\to L$ be a $\mathbf
Q_\ell$-embedding with $\tau=\tau_i$ on $K_0$, and $\kappa\colon L\to
B_{\mathrm{dR}}$ a map of $\mathbf Q_\ell$-algebras with $\kappa(\tau_i(c))=c$
for $c\in K_0$, $\kappa(\zeta)=\zeta_{\ell^s}$ and $\kappa(\tau(c))=c$ for $c\in
K$, and let $K_\kappa$ be as in
Lemma~\ref{asm-two-three-killing-wach-frame-exists-inputs-heights-full-image}.
Then for $d\in D_i(V)$ and $a\in\mathbf Z$, $\iota(d)\in F^a_\tau(V)$ if and only
if $\iota(d)\in V\otimes\operatorname{Fil}^aB_{\mathrm{dR}}+K_\kappa$.
\end{lemma}
\leanref{Development/AbelianVarietyTheory/Proofs/InertiaEigenvalues/DeRhamComponentWachLetterBasisFiltrationStatement.lean}{GlobalTheory.T.deRhamComponent_wachLetter_basis_filtration}
\leanref{Development/AbelianVarietyTheory/Proofs/InertiaEigenvalues/DeRhamComponentWachLetterBasisFiltration.lean}{GlobalTheory.NL.CompD.filAt_iff}
\begin{proof}
Put $y_0=\iota(d)$ and $V^a=V\otimes\operatorname{Fil}^aB_{\mathrm{dR}}$.

First, $y_0\in D_{\mathrm{dR}}(V)$: $d$ is fixed by $G_K$, and $\iota$ commutes
with the action of $G_K$ \cite[\S9.1]{BC09}.

Second, $D_{\tau'}\subset K_\kappa$ for every $\mathbf Q_\ell$-embedding
$\tau'\neq\tau$ of $K$ in $L$. Let $z\in D_{\tau'}$, and choose $c\in K$ with
$\tau(c)\neq\tau'(c)$. With $c$ acting on $B_{\mathrm{dR}}$, $cz=\tau'(c)z$ and
$c=\kappa(\tau(c))$, so $(\tau(c)-\tau'(c))z=\tau(c)z-\kappa(\tau(c))z$ lies in
$K_\kappa$. As $\tau(c)-\tau'(c)$ is a nonzero element of $L$ and $K_\kappa$ is
an $L$-subspace, $z\in K_\kappa$.

Third, $K$ has $[K:\mathbf Q_\ell]$ $\mathbf Q_\ell$-embeddings into $L$. It has
that many into $\bar{\mathbf Q}_\ell$, each extends to an automorphism $g$ of
$\bar{\mathbf Q}_\ell$, and $g(K)\subset g(K_s)\subset L$
(Definition~\ref{asm-two-three-killing-wach-frame-exists-data}).

Fourth, $y_0=\sum_{\tau'}y_{\tau'}$, the sum over the $\mathbf Q_\ell$-embeddings
$\tau'\colon K\to L$, with $y_{\tau'}\in D_{\tau'}$. Apply
Lemma~\ref{asm-two-three-killing-wach-frame-exists-inputs-heights-filtration-component-split}
with $F=\mathbf Q_\ell$, $E=K$, $L'=L$, which it allows by the third step, and
$M=D_{\mathrm{dR}}(V)$, with $L$ acting through $V$ and $K$ by multiplication on
$B_{\mathrm{dR}}$. The two actions commute, and they preserve $D_{\mathrm{dR}}(V)$,
as $G_K$ acts $L$-linearly on $V$ and fixes $K\subset B_{\mathrm{dR}}$. Put
$y_{\tau'}=e_{\tau'}(y_0)$. By (2) of that lemma, $cy_{\tau'}=\tau'(c)y_{\tau'}$
for $c\in K$, so $y_{\tau'}\in D_{\tau'}$; and $\sum_{\tau'}y_{\tau'}=y_0$ by (1).

Now suppose that $y_0\in F^a_\tau(V)$. This space is the sum of $D_\tau\cap
V^a\subset V^a$ and the $D_{\tau'}$ with $\tau'\neq\tau$, which lie in
$K_\kappa$ by the second step; so $y_0\in V^a+K_\kappa$. Conversely, suppose
that $y_0\in V^a+K_\kappa$, and put $y=y_\tau$. By the second step,
$y_0-y=\sum_{\tau'\neq\tau}y_{\tau'}$ lies in $K_\kappa$, so $y\in
V^a+K_\kappa$. As $y\in D_\tau$,
Lemma~\ref{asm-two-three-killing-wach-frame-exists-inputs-heights-filtration-component-descent}
gives $y\in V^a$. Hence $y\in D_\tau\cap V^a$, and
$y_0=y+\sum_{\tau'\neq\tau}y_{\tau'}$ lies in $F^a_\tau(V)$.
% Lean: the second part of T.deRhamComponent_wachLetter_basis_filtration = CompD.component,
% through component_tied: CompD.filAt_iff (iota_mem with PD.ι_comm; others_le_compKer;
% exists_part = exists_sum = exists_proj on DdR, with DeRhamSplit.mul_mem_DdR) with hkey =
% mem_FilV_local (Lemma -filtration-component-descent), and hc from Points.points (card_emb:
% exists_aut, exists_restr). The proof uses only d in D_cris and the third clause of Comp
% (iota o tau = id on K). filAt_iff is 165 lines with its helpers. The other parts of component
% (independence over B_dR, the projection of the tau-part of D_dR onto D_i, dim D_i = dim V)
% are not used by lat_eq_fil; they carry most of component's closure of 846 theorems, about
% 11,000 lines (ddR_le, isDR_of_isCrys, tied_deRham_span, tied_letter_component_finrank).
\end{proof}

\begin{lemma}\label{asm-two-three-killing-wach-frame-exists-inputs-heights-filtration-single}
Let $\ell$, $s$, $K$, $L$ and $D$ be as in
Definition~\ref{asm-two-three-killing-wach-frame-exists-inputs-heights-expansions},
with its notation, let $(V,T)$ be a lattice representation of $G_K$ over $L$
with $V$ crystalline, let $i\in\mathbf Z/n$, and let $\zeta\in L$ be a primitive
$\ell^s$-th root of unity. Let $\tau\colon K\to L$ be a $\mathbf
Q_\ell$-embedding and $\kappa\colon L\to B_{\mathrm{dR}}$ a map of $\mathbf
Q_\ell$-algebras with $\kappa(\tau_i(c))=c$ for $c\in K_0$,
$\kappa(\zeta)=\zeta_{\ell^s}$ and $\kappa(\tau(c))=c$ for $c\in K$, and let
$K_\kappa$ be as in
Lemma~\ref{asm-two-three-killing-wach-frame-exists-inputs-heights-full-image}.
Let $d\in D_i(V)$, $k\in\mathbf Z$, and let $\delta_kd$ be the sequence with
value $d$ at $k$ and $0$ elsewhere. Then $\delta_kd\in\Lambda_{i,\zeta}(T)$ if
and only if $\iota(d)\in V\otimes\operatorname{Fil}^{-k}B_{\mathrm{dR}}+K_\kappa$.
\end{lemma}
\leanref{Development/AbelianVarietyTheory/Proofs/InertiaEigenvalues/WachLatticeEqFiltrationSumAtLevel.lean}{GlobalTheory.NL.LatFil.single_mem_iff}
\begin{proof}
Write $\varepsilon=\varepsilon_{i,\zeta}$, $\Lambda=\Lambda_{i,\zeta}(T)$ and
$N=N(T)$. Let $\Psi$ be as in
Lemma~\ref{asm-two-three-killing-wach-frame-exists-inputs-heights-full-image},
$t_s$ and $\mathrm{ev}$ as in
Lemma~\ref{asm-two-three-killing-wach-frame-exists-inputs-heights-full-evaluation}
for $\kappa$, and $m_0$ as in
Lemma~\ref{asm-two-three-killing-wach-frame-exists-inputs-heights-full}. Fix an
$L$-basis $v_1,\ldots,v_{n'}$ of $V$ and let $p_1,\ldots,p_{n'}$ be the
coordinates of $V_\kappa=\bigoplus_\alpha B_{\mathrm{dR}}v_\alpha$, as in the
proof of Lemma~\ref{asm-two-three-killing-wach-frame-exists-inputs-heights-full};
then $y\in V\otimes\operatorname{Fil}^a+K_\kappa$ if and only if
$p_\alpha(y)\in\operatorname{Fil}^a$ for every $\alpha$. Let $b_1,\ldots,b_r$ be
an $A$-basis of $N$
(Lemma~\ref{asm-two-three-killing-wach-frame-exists-inputs-basis}); $r\leq n'$, as
in the third step of the proof of
Lemma~\ref{asm-two-three-killing-wach-frame-exists-inputs-heights-basis}. By
Lemma~\ref{asm-two-three-killing-wach-frame-exists-inputs-heights-basis}(2),
which applies by the choice of $m_0$, $\Lambda$ consists of the sums
$\sum_je_j*\varepsilon(b_j)$ with $e_j\in L[[X]]$. For $x\in N$ put
$\psi(x)=(p_\alpha(\Psi(x)))_\alpha\in(\operatorname{Fil}^0)^{n'}$. Choose
$a_0\geq0$ with $\varepsilon(b_j)(u)=0$ for every $j$ and $u<-a_0$.

For $a\geq0$, a sequence $f$ vanishing below $-a$, $\alpha$ and $y\in
B_{\mathrm{dR}}$, say that $y$ \emph{expands} $f$ at shift $a$ if there is $c$
with $y-\sum_{r<M}t_s^rp_\alpha(\iota(f(r-a)))\in\operatorname{Fil}^{M-c}$ for
every $M\geq0$. At most one $y$ expands a given $f$, since
$\bigcap_m\operatorname{Fil}^m=0$
(Lemma~\ref{asm-two-three-killing-wach-frame-exists-inputs-heights-full-evaluation}(1)).
If $y$ and $y'$ expand $f$ and $f'$, then $y+y'$ expands $f+f'$; and for $e\in
L[[X]]$, $\mathrm{ev}(e)y$ expands $e*f$, by multiplying the two truncated sums
and using $p_\alpha(cz)=\kappa(c)p_\alpha(z)$ for $c\in L$ and
Lemma~\ref{asm-two-three-killing-wach-frame-exists-inputs-heights-full-evaluation}(2).

(a) For $e_1,\ldots,e_r\in L[[X]]$ and $a\geq a_0$,
$t_s^a\sum_j\mathrm{ev}(e_j)\psi(b_j)_\alpha$ expands
$\sum_je_j*\varepsilon(b_j)$ at shift $a$: by
Lemma~\ref{asm-two-three-killing-wach-frame-exists-inputs-heights-full-image}(2),
$t_s^a\psi(b_j)_\alpha$ expands $\varepsilon(b_j)$.

(b) For $x\in N$, write $\varepsilon(x)=\sum_je_j*\varepsilon(b_j)$. Then
$\psi(x)=\sum_j\mathrm{ev}(e_j)\psi(b_j)$: by
Lemma~\ref{asm-two-three-killing-wach-frame-exists-inputs-heights-full-image}(2)
and (a), $t_s^{a_0}\psi(x)_\alpha$ and $t_s^{a_0}\sum_j\mathrm{ev}(e_j)
\psi(b_j)_\alpha$ both expand $\varepsilon(x)$ at shift $a_0$, and $t_s$ is a
unit.

(c) Every $z\in(\operatorname{Fil}^0)^{n'}$ is $\sum_j\beta_j\psi(b_j)$ with
$\beta_j\in\operatorname{Fil}^0$: by
Lemma~\ref{asm-two-three-killing-wach-frame-exists-inputs-heights-full-image}(1),
$z$ is a combination of the $\psi(x)$ with coefficients in $\operatorname{Fil}^0$,
and by (b) each $\psi(x)$ is a combination of the $\psi(b_j)$ with coefficients
$\mathrm{ev}(e_j)\in\operatorname{Fil}^0$.

(d) The $\psi(b_j)$ are linearly independent over $B_{\mathrm{dR}}$: every
vector of $B_{\mathrm{dR}}^{n'}$ is $t_s^{-q}$ times a vector of
$(\operatorname{Fil}^0)^{n'}$ for some $q$, so by (c) the $r\leq n'$ vectors
$\psi(b_j)$ span $B_{\mathrm{dR}}^{n'}$.

(e) If $a\geq a_0$, $q\geq0$ and $\delta_{q-a}d=\sum_je_j*\varepsilon(b_j)$,
then $t_s^qp_\alpha(\iota(d))=t_s^a\sum_j\mathrm{ev}(e_j)\psi(b_j)_\alpha$ for
every $\alpha$: the left side expands $\delta_{q-a}d$ at shift $a$, and so does
the right side, by (a).

Now write $k=q-a$ with $a\geq a_0$ and $q\geq0$. If $\delta_kd\in\Lambda$, then
$\delta_kd=\sum_je_j*\varepsilon(b_j)$, and by (e)
$t_s^qp_\alpha(\iota(d))\in t_s^a\operatorname{Fil}^0=\operatorname{Fil}^a$, so
$p_\alpha(\iota(d))\in\operatorname{Fil}^{-k}$ for every $\alpha$. Conversely,
suppose that $p_\alpha(\iota(d))\in\operatorname{Fil}^{-k}$ for every $\alpha$.
By (c), $t_s^{q-a}p(\iota(d))=\sum_j\beta_j\psi(b_j)$ with
$\beta_j\in\operatorname{Fil}^0$, where $p=(p_\alpha)_\alpha$. Put $m=m_0+a$. The
sequence $\delta_{q+m-a}d$ lies in $\Lambda$, as $q+m-a\geq m_0$, so it is
$\sum_je_j*\varepsilon(b_j)$, and by (e)
$t_s^a\sum_j\mathrm{ev}(e_j)\psi(b_j)=t_s^{q+m}p(\iota(d))=
t_s^{m+a}\sum_j\beta_j\psi(b_j)$. By (d), $\mathrm{ev}(e_j)=t_s^m\beta_j\in
\operatorname{Fil}^m$, so $e_j=X^me'_j$ with $e'_j\in L[[X]]$
(Lemma~\ref{asm-two-three-killing-wach-frame-exists-inputs-heights-full-evaluation}(2)).
Then $\delta_{q+m-a}d=X^m*\sum_je'_j*\varepsilon(b_j)$, and shifting back by
$m$, $\delta_kd=\sum_je'_j*\varepsilon(b_j)\in\Lambda$.
% Lean: single_mem_iff (Rep, Rep.unique, Rep.add, Rep.sum, Rep.single, Rep.conv,
% pow_mul_mem_iff; key, hA, hspan', hB, hli by linearIndependent_of_top_le_span_of_card_le_finrank)
% with mem_sup_iff, crd, finrank_le, as applied in lat_eq_fil. About 300 lines.
\end{proof}

\begin{lemma}\label{asm-two-three-killing-wach-frame-exists-inputs-heights-filtration}
Let $\ell$, $s$, $K$, $L$ and $D$ be as in
Definition~\ref{asm-two-three-killing-wach-frame-exists-inputs-heights-expansions},
with its notation, let $(V,T)$ be a lattice representation of $G_K$ over $L$
with $V$ crystalline, let $i\in\mathbf Z/n$, and let $\zeta\in L$ be a primitive
$\ell^s$-th root of unity. Then there is a $\mathbf Q_\ell$-embedding
$\tau\colon K\to L$ such that a sequence $f\colon\mathbf Z\to V\otimes
B_{\mathrm{cris}}$ lies in $\Lambda_{i,\zeta}(T)$ if and only if $f$ vanishes
below some integer and, for every $u\in\mathbf Z$, $f(u)\in D_i(V)$ and
$\iota(f(u))\in F^{-u}_\tau(V)$.
\end{lemma}
\leanref{Development/AbelianVarietyTheory/Proofs/InertiaEigenvalues/UnramifiedSubfieldLeCyclotomicExt.lean}{GlobalTheory.NL.Points.points}
\leanref{Development/AbelianVarietyTheory/Proofs/InertiaEigenvalues/WachLatticeEqFiltrationSumAtLevel.lean}{GlobalTheory.NL.LatFil.lat_eq_fil}
\begin{proof}
Choose $\tau_s$, $\tau$ and $\kappa$ by
Lemma~\ref{asm-two-three-killing-wach-frame-exists-inputs-heights-full-embedding};
then $\tau=\tau_i$ on $K_0$, $\kappa(\tau_i(c))=c$ for $c\in K_0$,
$\kappa(\zeta)=\zeta_{\ell^s}$ and $\kappa(\tau(c))=c$ for $c\in K$. Write
$\Lambda=\Lambda_{i,\zeta}(T)$, and $\delta_ud$ as in
Lemma~\ref{asm-two-three-killing-wach-frame-exists-inputs-heights-filtration-single}.

First, for $d\in D_i(V)$ and $u\in\mathbf Z$, $\delta_ud\in\Lambda$ if and only
if $\iota(d)\in F^{-u}_\tau(V)$. By
Lemma~\ref{asm-two-three-killing-wach-frame-exists-inputs-heights-filtration-single},
$\delta_ud\in\Lambda$ if and only if $\iota(d)\in
V\otimes\operatorname{Fil}^{-u}B_{\mathrm{dR}}+K_\kappa$, and by
Lemma~\ref{asm-two-three-killing-wach-frame-exists-inputs-heights-filtration-component}
this holds if and only if $\iota(d)\in F^{-u}_\tau(V)$.

Second, by
Lemma~\ref{asm-two-three-killing-wach-frame-exists-inputs-heights-full} there is
$m_0$ as in
Lemma~\ref{asm-two-three-killing-wach-frame-exists-inputs-heights-standard-graded},
which applies as $\zeta^{\ell^s}=1$: $f\in\Lambda$ if and only if $f$ vanishes
below some integer and $\delta_uf(u)\in\Lambda$ for every $u$. The elements of
$\Lambda$ take values in $D_i(V)$
(Definition~\ref{asm-two-three-killing-wach-frame-exists-inputs-heights-expansions}(3)),
so $\delta_uf(u)\in\Lambda$ implies $f(u)\in D_i(V)$. With the first step, $f\in
\Lambda$ if and only if $f$ vanishes below some integer and, for every $u$,
$f(u)\in D_i(V)$ and $\iota(f(u))\in F^{-u}_\tau(V)$.
% Lean: the first part of Points.points (a τ with Pt (HW l) h i ζ τ) and lat_eq_fil for that
% τ: Points.points (2) for κ (Comp), wach_expansions_span_full_lattice, graded_of_level_le,
% N_basis and the span of wachExpansion_basis_divisibility, exists_ev_dR, the image Ψ, the
% second part of deRhamComponent_wachLetter_basis_filtration (hfil), mem_sup_iff and
% single_mem_iff (main). About 100 lines besides the lemmas cited.
\end{proof}

\begin{lemma}\label{asm-two-three-killing-wach-frame-exists-inputs-heights}
Let $\ell$, $s$, $K$ and $L$ be as in
Definition~\ref{asm-two-three-killing-wach-frame-exists-data}, let $D$ be a
definite-side datum with local Hodge-Tate data, Fontaine data and comparison
data at every prime (Definition~\ref{asm-local-bound-two-three}(3)), with its
Wach data for $(K,s,L)$, let $(V,T)$ be a lattice representation of $G_K$ over
$L$ with $V$ crystalline, and let $M_0,h_0\geq0$ be integers such that the
Hodge-Tate weights of $V$ along every $\mathbf Q_\ell$-embedding $K\to L$
(Definition~\ref{asm-two-three-local-killing}(3)) lie in $[-M_0,h_0]$. Then for
every $i\in\mathbf Z/n$, $Q^{M_0}\varphi_D(x)\in N(T)$ for $x\in e_{i-1}N(T)$,
and $Q^{h_0}y$ lies in the $A$-submodule generated by $\varphi_D(e_{i-1}N(T))$
for $y\in e_iN(T)$.
\end{lemma}
\leanref{Development/AbelianVarietyTheory/Proofs/InertiaEigenvalues/WachLatticeHeightsFromHodgeJumpsStatement.lean}{GlobalTheory.T.wachLattice_heights_from_hodgeJumps}
\leanref{Development/AbelianVarietyTheory/Proofs/InertiaEigenvalues/WachLatticeHeightsFromHodgeJumps.lean}{GlobalTheory.NL.Heights.heights_from_jumps}
\begin{proof}
Use the notation of
Definition~\ref{asm-two-three-killing-wach-frame-exists-inputs-heights-expansions}.
Put $N=N(T)$ and $r=\dim_LV$, the rank of $T$ over $\mathcal O$. By
Lemma~\ref{asm-two-three-killing-wach-frame-exists-inputs-basis} and the proof
of Lemma~\ref{asm-two-three-killing-wach-frame-exists-inputs}(2), $N$ has an
$A$-basis $b_1,\ldots,b_r$; for $y\in N$ write $y=\sum_jc_j(y)b_j$ with
$c_j(y)\in A$. A \emph{point} is a pair $(i,\zeta)$ with $i\in\mathbf Z/n$ and
$\zeta\in L$ a primitive $\ell^s$-th root of unity; then $\zeta^\ell$ is a
primitive $\ell^{s-1}$-th root of unity. Recall that $\varphi(e_{i-1})=e_i$ in
$A$, so $\varphi_D(e_{i-1}x)=e_i\varphi_D(x)$ for $x\in D(T)$.

First, five facts.

(a) $\varepsilon_{i,\zeta}(ax)=\lambda_{i,\zeta}(a)*\varepsilon_{i,\zeta}(x)$ for
$a\in A$ and $x\in N$, by
Definition~\ref{asm-two-three-killing-wach-frame-exists-inputs-heights-expansions}(3),
as $\iota$ is injective \cite[Theorem~9.1.5]{BC09} and additive.

(b) Let $x\in D(T)$ and $k\geq0$ with $Q^kx\in N$, and suppose that $Q^k$
divides every $c_j(Q^kx)$ in $A$. Then $x\in N$. Indeed $Q$ is a unit in $E$:
$\pi$ is a unit in $E$ and $\varphi$ extends to $E$, so $X_s=\varphi^{s-1}(\pi)$
and $\varphi(X_s)$ are units, and $\varphi(X_s)=\varphi^{s-1}(\pi q)=X_sQ$ in $A$.
So multiplication by $Q^k$ is injective on the $E$-module $D(T)$, and from
$Q^kx=\sum_jQ^kd_jb_j$, with $c_j(Q^kx)=Q^kd_j$, $x=\sum_jd_jb_j\in N$.

(c) Let $(i,\zeta)$ be a point, $x\in N$ and $m\geq0$ with $Q^m\varphi_D(x)\in
N$. Then $\varepsilon_{i,\zeta}(Q^m\varphi_D(x))(u)=0$ for $u<m$. Indeed, by
Lemma~\ref{asm-two-three-killing-wach-frame-exists-inputs-heights-points},
$\lambda_{i,\zeta}(Q^m)=X^mw^m$, whose coefficients $\mu_t$ vanish for $t<m$; and
$\varepsilon_{i-1,\zeta^\ell}(x)=1*\varepsilon_{i-1,\zeta^\ell}(x)$ lies in
$\Lambda_{i-1,\zeta^\ell}(T)$, so it vanishes in negative degrees by
Lemma~\ref{asm-two-three-killing-wach-frame-exists-inputs-heights-standard}. In
the formula of
Definition~\ref{asm-two-three-killing-wach-frame-exists-inputs-heights-expansions}(4)
for $u<m$, every term vanishes: $\mu_t=0$ if $t<m$, and $u-t<0$ if $t\geq m$.

(d) Let $(i,\zeta)$ be a point and $\tau$ as in
Lemma~\ref{asm-two-three-killing-wach-frame-exists-inputs-heights-filtration}.
Then $F^{a+1}_\tau(V)=F^a_\tau(V)$ if $a<-M_0$ or $a>h_0$, since such an $a$ is
not a Hodge-Tate weight of $V$ along $\tau$, so that
$\operatorname{Fil}^{a+1}=\operatorname{Fil}^a$ on $D_\tau$.

(e) Let $(i,\zeta)$ be a point. By
Lemma~\ref{asm-two-three-killing-wach-frame-exists-inputs-heights-full} there is
$m_0$ as in Lemma~\ref{asm-two-three-killing-wach-frame-exists-inputs-heights-basis},
which therefore applies to $(i,\zeta)$. It also applies to $(i',\zeta^\ell)$ for every
$i'\in\mathbf Z/n$, with $m_0=0$, by
Lemma~\ref{asm-two-three-killing-wach-frame-exists-inputs-heights-standard}.

Second, the first assertion. Let $i\in\mathbf Z/n$ and $x\in e_{i-1}N$, say
$x=e_{i-1}x_0$ with $x_0\in N$. By
Definition~\ref{asm-two-three-killing-wach-frame-exists-data}(4) there is
$k\geq0$ with $Q^k\varphi_D(x_0)\in N$; put $y=Q^{k+M_0}\varphi_D(x_0)\in N$.
Then $Q^k\cdot Q^{M_0}\varphi_D(x)=e_iy\in N$. By (b) it suffices to show that
$Q^k$ divides $c_j(e_iy)=e_ic_j(y)$ for every $j$, and by
Lemma~\ref{asm-two-three-killing-wach-frame-exists-inputs-heights-division} that
$X^k$ divides $\lambda_{i',\zeta}(e_ic_j(y))$ for every point $(i',\zeta)$. For
$i'\neq i$ the $i'$-th component of $e_ic_j(y)$ is $0$, so
$\lambda_{i',\zeta}(e_ic_j(y))=0$. Let $i'=i$, let $\tau$ be as in
Lemma~\ref{asm-two-three-killing-wach-frame-exists-inputs-heights-filtration}
and $m_0$ as in (e), and put $f(u)=\varepsilon_{i,\zeta}(e_iy)(u+k)$. By (a),
$\varepsilon_{i,\zeta}(e_iy)=\lambda_{i,\zeta}(e_i)*\varepsilon_{i,\zeta}(y)$,
and $\varepsilon_{i,\zeta}(y)$ vanishes below $k+M_0$ by (c); so $f$ vanishes
below $M_0$, and its values lie in $D_i(V)$. For $u\geq M_0$, the sequence with
value $f(u)$ at $\max(m_0,u)$ and $0$ elsewhere lies in $\Lambda_{i,\zeta}(T)$
by the choice of $m_0$, so
Lemma~\ref{asm-two-three-killing-wach-frame-exists-inputs-heights-filtration}
gives $\iota(f(u))\in F^{-\max(m_0,u)}_\tau(V)$, which is $F^{-u}_\tau(V)$ by
(d), as every $a$ with $-\max(m_0,u)\leq a<-u$ satisfies $a<-M_0$. So
$\iota(f(u))\in F^{-u}_\tau(V)$ for every $u$, and $f\in\Lambda_{i,\zeta}(T)$ by
the same lemma. As $\varepsilon_{i,\zeta}(e_iy)(u)=f(u-k)$ for every $u$,
Lemma~\ref{asm-two-three-killing-wach-frame-exists-inputs-heights-basis}(3)
shows that $X^k$ divides $\lambda_{i,\zeta}(c_j(e_iy))$ for every $j$.

Third, the second assertion. Let $i\in\mathbf Z/n$ and $y\in e_iN$, say
$y=e_iy_0$ with $y_0\in N$. The $A$-submodule generated by $\varphi_D(N)$ is
generated by the $\varphi_D(b_j)$, as $\varphi_D(ax)=\varphi(a)\varphi_D(x)$. By
Definition~\ref{asm-two-three-killing-wach-frame-exists-data}(4) there are
$k\geq0$ and $a_j\in A$ with $Q^ky_0=\sum_ja_j\varphi_D(b_j)$, and $m\geq0$ with
$Q^m\varphi_D(b_j)\in N$ for every $j$. Then
\[
 Q^{h_0+m+k}y_0=\sum_j\big(Q^{h_0}a_j\big)\,Q^m\varphi_D(b_j).
\]
We show that $Q^k$ divides $e_iQ^{h_0}a_j$ for every $j$. By
Lemma~\ref{asm-two-three-killing-wach-frame-exists-inputs-heights-division} it
suffices that $X^k$ divides $\lambda_{i',\zeta}(e_iQ^{h_0}a_j)$ for every point
$(i',\zeta)$; for $i'\neq i$ this is $0$. Let $i'=i$, write
$\lambda=\lambda_{i,\zeta}$ and $\varepsilon=\varepsilon_{i,\zeta}$, and put
$\beta_j=\lambda(Q^{h_0}a_j)$ and $g_j=\varepsilon(Q^m\varphi_D(b_j))$. As
$\lambda(e_iQ^{h_0}a_j)=\lambda(e_i)\beta_j$, it suffices that $X^k$ divides every
$\beta_j$.

The $g_j$ vanish below $m$ by (c), and the $g_j(m)$ are linearly independent
over $L$. Indeed, write $\lambda(Q)=Xw$ as in
Lemma~\ref{asm-two-three-killing-wach-frame-exists-inputs-heights-points}. In the
formula of
Definition~\ref{asm-two-three-killing-wach-frame-exists-inputs-heights-expansions}(4)
at $u=m$ only the term $t=m$ survives, as in (c), so
$g_j(m)=w(0)^m\varphi(\varepsilon_{i-1,\zeta^\ell}(b_j)(0))$ with $w(0)\neq0$.
The vectors $\varepsilon_{i-1,\zeta^\ell}(b_j)(0)$ span $D_{i-1}(V)$: for $d\in
D_{i-1}(V)$, the sequence with value $d$ at $0$ and $0$ elsewhere lies in
$\Lambda_{i-1,\zeta^\ell}(T)$ by
Lemma~\ref{asm-two-three-killing-wach-frame-exists-inputs-heights-standard}, so by
(e) and Lemma~\ref{asm-two-three-killing-wach-frame-exists-inputs-heights-basis}(2)
it is $\sum_j\eta_j*\varepsilon_{i-1,\zeta^\ell}(b_j)$ with $\eta_j\in L[[X]]$, and
at $0$ this gives $d=\sum_j\eta_j(0)\,\varepsilon_{i-1,\zeta^\ell}(b_j)(0)$, as the
$\varepsilon_{i-1,\zeta^\ell}(b_j)$ vanish in negative degrees. By
Lemma~\ref{asm-two-three-killing-wach-frame-exists-inputs-heights-letters},
$\varphi$ maps $D_{i-1}(V)$ onto $D_i(V)$, of dimension $r$; so the $r$ vectors
$\varphi(\varepsilon_{i-1,\zeta^\ell}(b_j)(0))$ span $D_i(V)$ and are linearly
independent.

The sequence $\varepsilon(y_0)$ vanishes below $-h_0$. Let $\tau$ be as in
Lemma~\ref{asm-two-three-killing-wach-frame-exists-inputs-heights-filtration}.
By (e) and Lemma~\ref{asm-two-three-killing-wach-frame-exists-inputs-heights-basis}(2),
every $f\in\Lambda_{i,\zeta}(T)$ is $\sum_j\eta_j*\varepsilon(b_j)$, so all of them
vanish below one integer $m_1$. Let $f\in\Lambda_{i,\zeta}(T)$, $u<-h_0$ and
$u'=\min(m_1-1,u)$. By
Lemma~\ref{asm-two-three-killing-wach-frame-exists-inputs-heights-filtration},
$f(u)\in D_i(V)$ and $\iota(f(u))\in F^{-u}_\tau(V)$, which is
$F^{-u'}_\tau(V)$ by (d), as every $a$ with $-u\leq a<-u'$ satisfies $a>h_0$. By
the same lemma, the sequence with value $f(u)$ at $u'$ and $0$ elsewhere lies in
$\Lambda_{i,\zeta}(T)$; as $u'<m_1$, $f(u)=0$. Apply this to
$f=1*\varepsilon(y_0)$.

Now, by (a) and the identity above,
$\sum_j\beta_j*g_j=\varepsilon(Q^{h_0+m+k}y_0)=\lambda(Q)^{h_0+m+k}*\varepsilon(y_0)$,
which vanishes below $m+k$, as $X^{h_0+m+k}$ divides $\lambda(Q)^{h_0+m+k}$.
Suppose that $X^t$ divides every $\beta_j$ for some $t<k$. As $g_j$ vanishes below
$m$, the value of $\sum_j\beta_j*g_j$ at $m+t$ is $\sum_j\nu_{j,t}\,g_j(m)$, where
$\nu_{j,t}$ is the coefficient of $X^t$ in $\beta_j$; it is $0$, as $m+t<m+k$, so every
$\nu_{j,t}$ is $0$ and $X^{t+1}$ divides every $\beta_j$. By induction on $t$, $X^k$
divides every $\beta_j$.

So $e_iQ^{h_0}a_j=Q^kd_j$ with $d_j\in A$, and $e_iQ^kd_j=Q^kd_j$, as
$e_i^2=e_i$. Hence
\[
 Q^k\cdot Q^{h_0}y=e_iQ^{h_0}\sum_ja_j\varphi_D(b_j)=\sum_jQ^kd_j\varphi_D(b_j)
 =Q^k\sum_jd_j\varphi_D(e_{i-1}b_j),
\]
as $\varphi_D(e_{i-1}b_j)=e_i\varphi_D(b_j)$. By the injectivity in (b),
$Q^{h_0}y=\sum_jd_j\varphi_D(e_{i-1}b_j)$, which lies in the $A$-submodule
generated by $\varphi_D(e_{i-1}N)$.
% Lean: GlobalTheory.NL.Heights.heights_from_jumps (P5 D HW F, the first two parts of p5 W T M h)
% for M i = M0 and h i = h0 (jumpsLet_const, OverLetter): N_basis (nb), hunit (IoXQ.main,
% last part, through isUnit_Q: isUnit_φA_pow, isUnit_X and φA_X), hdiv (Lemma
% -heights-division), hfull, hstd, hinj, hlev, hpts, hlat (Lemmas -heights-full, -standard,
% -basis, -points, -filtration), hsm (LevInj.lift_smul with levN_eq, fact (a)), hφ (levB_φ),
% hneg and hφ0 (fact (c)), hconst (fact (d), from hJ and OverLetter). The second step is
% pole_bound (mem_of_dvd and cancel_Q for (b), lev_letE_mul, conv_low, mem_of_low with
% const_on); the third is inv_pole_bound (N_φinv, the span of the φ_D(b_j), dvd_of_low with
% conv_lead, low_of_mem with bdd_of_gens, indep_lead and indep_const with PhiD.phi_letters,
% tied_letter_component_finrank and LevInj.rank_le, hfin with letE_mul_self). The third
% part of p5 (saturation) is not used here. The top level, without the lemmas it cites:
% about 470 lines. The whole closure: 1,373 theorems, about 19,400 lines (theorems only) in
% 218 files; an earlier count gave 960 theorems, about 14,300 lines in 188 files.
\end{proof}

\begin{theorem}[Fontaine {\cite[\S9.1]{BC09}}]\label{asm-two-three-killing-wach-frame-exists-inputs-ratio}
Let $\ell$ be a prime, fix an algebraic closure $\bar{\mathbf Q}_\ell$ of
$\mathbf Q_\ell$, and let $B_{\mathrm{cris}}$ be Fontaine's ring of crystalline
periods, with its action of $G_{\mathbf Q_\ell}=\operatorname{Gal}(\bar{\mathbf
Q}_\ell/\mathbf Q_\ell)$. Let $F\subset\bar{\mathbf Q}_\ell$ be a finite
extension of $\mathbf Q_\ell$ with maximal unramified subextension $F_0$, which
lies in $B_{\mathrm{cris}}$, and put $G_F=\operatorname{Gal}(\bar{\mathbf
Q}_\ell/F)$. Then $B_{\mathrm{cris}}$ is a domain, and if $x,y\in
B_{\mathrm{cris}}$ with $x\neq0$ satisfy $\sigma(x)\,y=x\,\sigma(y)$ for every
$\sigma\in G_F$, then $y=cx$ for some $c\in F_0$.
\end{theorem}
\leanref{Development/AbelianVarietyTheory/Proofs/InertiaEigenvalues/RatioInvariantInK0Statement.lean}{GlobalTheory.T.ratio_invariant_in_K0}
\leanref{Development/AbelianVarietyTheory/Proofs/InertiaEigenvalues/RatioInvariantInK0.lean}{GlobalTheory.Per.FracInv.tied_frac_invariant}
\leanref{Development/AbelianVarietyTheory/Proofs/InertiaEigenvalues/FontaineFormalismOfRegular.lean}{GlobalTheory.RegularRing.isDomain_of_tied}
% Lean: for the datum's period rings (LocTied), read in Fontaine's sense. tied_frac_invariant:
% the datum's B_dR is a field (Per.FontaineBound.isField_of_tied), the ratio lies in K by
% frac_dR (invariants of B_dR), and in the maximal unramified extension by
% FracAlg.mem_Kur_of_ratio (exists_fits, norm_of_ratio, UnramValue.le_Kur_of_norm). Size of
% the closure of tied_frac_invariant: 682 theorems, about 9,300 lines in 124 files (a count
% that includes the definitions it uses); isDomain_of_tied is ι_inj into the field B_dR.

\begin{lemma}\label{asm-two-three-killing-wach-frame-exists-inputs-independent}
Let $\ell$, $B_{\mathrm{cris}}$, $F$ and $F_0$ be as in
Theorem~\ref{asm-two-three-killing-wach-frame-exists-inputs-ratio}. Let $L$ be a
finite extension of $\mathbf Q_\ell$ and $V$ a finite-dimensional $L$-vector
space with an $L$-linear action of $G_F$, let $G_F$ act on
$V\otimes_{\mathbf Q_\ell}B_{\mathrm{cris}}$ by $\sigma\otimes\sigma$, and put
$D_{\mathrm{cris},F}(V)=(V\otimes_{\mathbf Q_\ell}B_{\mathrm{cris}})^{G_F}$, a
vector space over $F_0$, which acts on the second factor. Then every family in
$D_{\mathrm{cris},F}(V)$ that is linearly independent over $F_0$ is linearly
independent over $B_{\mathrm{cris}}$, and
$\dim_{F_0}D_{\mathrm{cris},F}(V)\leq\dim_{\mathbf Q_\ell}V$.
\end{lemma}
\leanref{Development/AbelianVarietyTheory/Proofs/InertiaEigenvalues/FontaineFormalismOfRegular.lean}{GlobalTheory.RegularRing.indepB}
\leanref{Development/AbelianVarietyTheory/Proofs/InertiaEigenvalues/LinearIndepOnOfFixed.lean}{GlobalTheory.LinearIndependentFixed.linearIndepOn_of_fixed}
\begin{proof}
Put $B=B_{\mathrm{cris}}$ and $M=V\otimes_{\mathbf Q_\ell}B$, a free $B$-module
of rank $d=\dim_{\mathbf Q_\ell}V$; $B$ is a domain by
Theorem~\ref{asm-two-three-killing-wach-frame-exists-inputs-ratio}, so $M$ is
torsion-free over $B$. For $\sigma\in G_F$, $b\in B$ and $z\in M$,
$\sigma(bz)=\sigma(b)\sigma(z)$.

First, the independence, for a finite family $v_1,\ldots,v_k$ in
$D_{\mathrm{cris},F}(V)$, linearly independent over $F_0$, by induction on the
number of nonzero coefficients of a relation. Suppose that $b_1,\ldots,b_k\in B$
satisfy $\sum_jb_jv_j=0$, that every such relation with fewer nonzero
coefficients is trivial, and that $b_{j_0}\neq0$. For $\sigma\in G_F$, applying
$\sigma$ and using $\sigma(v_j)=v_j$ gives $\sum_j\sigma(b_j)v_j=0$, so
\[
 \sum_j\big(b_{j_0}\sigma(b_j)-\sigma(b_{j_0})b_j\big)v_j
 =b_{j_0}\sum_j\sigma(b_j)v_j-\sigma(b_{j_0})\sum_jb_jv_j=0.
\]
Its coefficient at $j_0$ vanishes, and so does every coefficient with $b_j=0$,
as $\sigma(0)=0$. So it has fewer nonzero coefficients, and it is trivial:
$b_{j_0}\sigma(b_j)=\sigma(b_{j_0})b_j$ for every $j$ and every $\sigma\in G_F$.
By Theorem~\ref{asm-two-three-killing-wach-frame-exists-inputs-ratio} there are
$c_j\in F_0$ with $b_j=c_jb_{j_0}$. Then $b_{j_0}\sum_jc_jv_j=\sum_jb_jv_j=0$,
so $\sum_jc_jv_j=0$, as $M$ is torsion-free, and $c_{j_0}=0$ by the
independence over $F_0$. Hence $b_{j_0}=c_{j_0}b_{j_0}=0$, a contradiction. So
every relation is trivial.

Second, the dimension. A family in the free $B$-module $M$ of rank $d$ that is
linearly independent over the domain $B$ has at most $d$ members, as it stays
independent over the fraction field of $B$. So every family in
$D_{\mathrm{cris},F}(V)$ that is linearly independent over $F_0$ has at most $d$
members.
% Lean: indepB applies LinearIndependentFixed.linearIndepOn_of_fixed (induction on the number
% k of nonzero coefficients, the relation b' = b_{j0} g(b) - g(b_{j0}) b) with FracInv from
% Per.TiedRestrict.fracInv_of_tied (= T.ratio_invariant_in_K0); the dimension bound is
% RegularRing.finite_le₀ (rank_le with LinearIndependent.fintype_card_le_finrank).
% 10 theorems, about 260 lines.
\end{proof}

\begin{lemma}\label{asm-two-three-killing-wach-frame-exists-inputs-restrict}
Let $\ell$ and $B_{\mathrm{cris}}$ be as in
Theorem~\ref{asm-two-three-killing-wach-frame-exists-inputs-ratio}, let
$F\subset F'\subset\bar{\mathbf Q}_\ell$ be finite extensions of
$\mathbf Q_\ell$ with maximal unramified subextensions $F_0$ and $F'_0$, let $L$
be a finite extension of $\mathbf Q_\ell$ and $V$ a finite-dimensional
$L$-vector space with an $L$-linear action of $G_F$, and define
$D_{\mathrm{cris},F}(V)$ as in
Lemma~\ref{asm-two-three-killing-wach-frame-exists-inputs-independent}, and
$D_{\mathrm{cris},F'}(V)$ likewise for the restriction of the action to
$G_{F'}$. If $\dim_{F_0}D_{\mathrm{cris},F}(V)=\dim_{\mathbf Q_\ell}V$, then
$\dim_{F'_0}D_{\mathrm{cris},F'}(V)=\dim_{\mathbf Q_\ell}V$.
\end{lemma}
\leanref{Development/AbelianVarietyTheory/Proofs/InertiaEigenvalues/RestrictCrystallineDeRhamFiniteExtStatement.lean}{GlobalTheory.T.restrict_crystalline_deRham_finiteExt}
\leanref{Development/AbelianVarietyTheory/Proofs/InertiaEigenvalues/RestrictCrystallineDeRhamFiniteExt.lean}{GlobalTheory.Per.TiedRestrict.isCrys_restrict}
\begin{proof}
Put $d=\dim_{\mathbf Q_\ell}V$. As $G_{F'}\subset G_F$,
$D_{\mathrm{cris},F}(V)\subset D_{\mathrm{cris},F'}(V)$. Let $v_1,\ldots,v_d$ be
an $F_0$-basis of $D_{\mathrm{cris},F}(V)$. By
Lemma~\ref{asm-two-three-killing-wach-frame-exists-inputs-independent} for $F$,
the $v_j$ are linearly independent over $B_{\mathrm{cris}}$, hence over its
subfield $F'_0$. They lie in $D_{\mathrm{cris},F'}(V)$, so
$\dim_{F'_0}D_{\mathrm{cris},F'}(V)\geq d$, and
Lemma~\ref{asm-two-three-killing-wach-frame-exists-inputs-independent} for $F'$
gives the reverse inequality.
% Lean: Per.TiedRestrict.isCrys_restrict (isCrys_iff_finrank, dcris_le, indepB with
% fracInv_of_tied, LinearIndependent.restrict_scalars', finite_le₀). The Lean's IsCrys says
% dim_L D = [F_0 : Q_l] dim_L V, which is the condition here by isCrys_iff_finrank. The
% closure of isCrys_restrict is 702 theorems, about 9,600 lines, nearly all of it in
% Theorem -inputs-ratio.
\end{proof}

\begin{lemma}\label{asm-two-three-killing-wach-frame-exists-inputs}
Let $\ell$, $s$, $K$ and $L$ be as in
Definition~\ref{asm-two-three-killing-wach-frame-exists-data}, let $D$ be a
definite-side datum with local Hodge-Tate data, Fontaine data and comparison
data at every prime (Definition~\ref{asm-local-bound-two-three}(3)), with its
Wach data for $(K,s,L)$, and let $(V,T)$ be a lattice representation of $G_K$
over $L$ with $V$ crystalline. Then:
\begin{enumerate}
\item $V$ is crystalline as a representation of $G_{K_s}$;
\item $N(T)$ has an $A$-basis whose image under $j_T$ is an $R$-basis of
  $R\otimes_{\mathcal O}T$;
\item if $M_0,h_0\geq0$ are integers and the Hodge-Tate weights of $V$ along
  every $\mathbf Q_\ell$-embedding $K\to L$
  (Definition~\ref{asm-two-three-local-killing}(3)) lie in $[-M_0,h_0]$, then
  for every $i\in\mathbf Z/n$, $Q^{M_0}\varphi_D(x)\in N(T)$ for
  $x\in e_{i-1}N(T)$, and $Q^{h_0}y$ lies in the $A$-submodule generated by
  $\varphi_D(e_{i-1}N(T))$ for $y\in e_iN(T)$;
\item $N(T)$ is the sum of its $A$-submodules $e_iN(T)$, $i\in\mathbf Z/n$.
\end{enumerate}
\end{lemma}
\leanref{Development/AbelianVarietyTheory/Proofs/InertiaEigenvalues/WachModuleLocalConclusionsUsedFormsStatement.lean}{GlobalTheory.T.wachModuleLocalConclusions_usedForms}
\leanref{Development/AbelianVarietyTheory/Proofs/InertiaEigenvalues/CrysOnCondOfEveryDatum.lean}{GlobalTheory.crysOnCond_all}
\leanref{Development/AbelianVarietyTheory/Proofs/InertiaEigenvalues/IsCrysOnOfIsCrys.lean}{GlobalTheory.CrysOnLevel.isCrysOn_of_isCrys}
\leanref{Development/AbelianVarietyTheory/Proofs/InertiaEigenvalues/EtaleComparisonOfWachBasis.lean}{GlobalTheory.NL.NygaardExtension.main}
\leanref{Development/AbelianVarietyTheory/Proofs/InertiaEigenvalues/LettersFree.lean}{GlobalTheory.NL.LetterFree.letters_free}
\begin{proof}
(1) For a finite extension $F$ of $\mathbf Q_\ell$ in $\bar{\mathbf Q}_\ell$
with maximal unramified subextension $F_0$, define $D_{\mathrm{cris},F}(V)$ as in
Lemma~\ref{asm-two-three-killing-wach-frame-exists-inputs-independent}. It is a
vector space over $L$ and over $F_0$, so
$[F_0:\mathbf Q_\ell]\dim_{F_0}D_{\mathrm{cris},F}(V)=[L:\mathbf Q_\ell]
\dim_LD_{\mathrm{cris},F}(V)$, and $\dim_{F_0}D_{\mathrm{cris},F}(V)=
\dim_{\mathbf Q_\ell}V$ if and only if $D_{\mathrm{cris},F}(V)$ has dimension
$[F_0:\mathbf Q_\ell]\dim_LV$ over $L$. As $V$ is crystalline,
$\dim_{K_0}D_{\mathrm{cris},K}(V)=\dim_{\mathbf Q_\ell}V$. Since $K\subset K_s$,
Lemma~\ref{asm-two-three-killing-wach-frame-exists-inputs-restrict}, with
$F=K$ and $F'=K_s$, gives $\dim_{F'_0}D_{\mathrm{cris},K_s}(V)=
\dim_{\mathbf Q_\ell}V$ for the maximal unramified subextension $F'_0$ of $K_s$.
By the second step of the proof of
Lemma~\ref{asm-two-three-killing-wach-frame-exists-derham}, $F'_0$ has degree
$[K_0:\mathbf Q_\ell]$. So $(V\otimes_{\mathbf Q_\ell}B_{\mathrm{cris}})^{G_{K_s}}$
has dimension $[K_0:\mathbf Q_\ell]\dim_LV$ over $L$, which is the
crystallinity of $V$ as a representation of $G_{K_s}$.

(2) By Lemma~\ref{asm-two-three-killing-wach-frame-exists-inputs-basis}, $N(T)$
has an $A$-basis $b_1,\ldots,b_r$, which is an $E$-basis of $D(T)$, so $r$ is
the rank of $T$ over $\mathcal O$. The elements $j_T(b_i)$ generate
$R\otimes_{\mathcal O}T$ over $R$: the image of $j_T$ does, and
$j_T(\sum_ie_ib_i)=\sum_i\bar e_ij_T(b_i)$ for $e_i\in E$, with $\bar e_i$ the
image of $e_i$ in $R$. As $R\otimes_{\mathcal O}T$ is free of rank $r$ over the
nonzero commutative ring $R$, these $r$ generators form a basis.

(3) This is Lemma~\ref{asm-two-three-killing-wach-frame-exists-inputs-heights}.

(4) Each $e_iN(T)$ lies in $N(T)$, as $N(T)$ is an $A$-submodule, and every
$x\in N(T)$ is $\sum_ie_ix$, since $\sum_ie_i=1$ in $A$.
% Lean: (1) is GlobalTheory.crysOnCond_all = CrysOnLevel.isCrysOn_of_isCrys: crys_iff (the
% identification of IsCrysOn over GKs with crystallinity over levelField K s = K_s: fixed_eq
% and finrank_maxUnramified_levelFieldFiniteExtension, 519 lines) and the first part of
% T.restrict_crystalline_deRham_finiteExt (= Per.TiedRestrict.isCrys_restrict, Lemma
% -inputs-restrict, for the datum's period rings read in Fontaine's sense; its size is in
% Theorem -inputs-ratio). (2) is NygaardExtension.main, first part:
% N_basis (Lemma -inputs-basis) and etale_comparison_of_wach_basis with
% WachTensor.exists_basis_jD (49 lines); its second clause (the phi-fixed part of R tensor T
% is 1 tensor T, WittTensorFixed.main) is not used. (3) is P5W through
% T.wachLattice_heights_from_hodgeJumps (Lemma -inputs-heights). (4) is the first clause of
% WachLatticeFree: NL.LetterFree.letters_free, letters_loc, lettersDecompose with sum_letE_smul
% and mem_Nl_iff.
% Old summary. (1) is GlobalTheory.crysOnCond_all (CrysOnLevel.isCrysOn_of_isCrys): 814
% theorems, about 10,200 lines in 157 files. (2) and (3) come from
% GlobalTheory.T.wachModuleLocalConclusions_usedForms (= NL.usedLocalConclusion_and_prvE):
% (2) is the first part p1et of WachEtaleComparison (ConclusionE.1), (3) is P5W
% (UsedLocalConclusion.1, the fifth part of DeltaInvariantLatticeCriterion) for the constant
% bounds M i = M0, h i = h0 (JumpsLet through OverLetter and the bijective letters emb);
% 2,635 theorems, about 40,900 lines in 548 files. The two overlap: with wClauses,
% wAll_holds is 2,655 theorems, about 41,100 lines in 557 files. The Lean's P5W also
% allows bounds that depend on the letter and has a saturation clause, and
% WachLatticeFree a freeness clause for each letter; neither is used here.
\end{proof}

\subsection{Decompletion and the killing of invariants}\label{subsection-asm-killing-invariants}

\begin{lemma}\label{asm-two-three-killing-wach-frame-exists-wach}
Let $\ell$, $s$, $K$, $L$ and $D$ be as in
Lemma~\ref{asm-two-three-killing-wach-frame-exists-inputs}, with the Wach data
of $D$ for $(K,s,L)$. Let $(V,T)$ be a lattice representation of $G_K$ over
$L$ with $V$ crystalline, and let $M_0,h_0\geq0$ be integers such that every
integer $a$ with $\operatorname{Fil}^{a+1}\neq\operatorname{Fil}^a$ on
$D_{\mathrm{dR}}(V)$ lies in $[-M_0,h_0]$. Then the Wach module $N(T)$
satisfies the Wach clauses for $M_0$ and $h_0$, and the Wach hypothesis holds
for $(V,T)$, $M_0$ and $h_0$.
\end{lemma}
\leanref{Development/AbelianVarietyTheory/Proofs/InertiaEigenvalues/CrystallineLatticeWachModuleExistsProof.lean}{GlobalTheory.NLink.wAll_holds}
\leanref{Development/AbelianVarietyTheory/Proofs/InertiaEigenvalues/WachHypothesisCrystallineLattices.lean}{GlobalTheory.WachHyp.hypW_crystalline}
\leanref{Development/AbelianVarietyTheory/Proofs/InertiaEigenvalues/WachHypothesisCrystallineLattices.lean}{GlobalTheory.WachHyp.wClauses}
\begin{proof}
Put $N=N(T)$. By
Lemma~\ref{asm-two-three-killing-wach-frame-exists-inputs}(1), $V$ is
crystalline as a representation of $G_{K_s}$. Let $\tau\colon K\to L$ be a
$\mathbf Q_\ell$-embedding and $a$ a Hodge-Tate weight of $V$ along $\tau$. If
$\operatorname{Fil}^{a+1}=\operatorname{Fil}^a$ on $D_{\mathrm{dR}}(V)$, the same
holds on its $\tau$-part, so $a\in[-M_0,h_0]$, and part (3) of that lemma
applies. By part (4) of that lemma, $N=\sum_ie_iN$.

(a) This is part (2) of
Lemma~\ref{asm-two-three-killing-wach-frame-exists-inputs}.

(b) The $x\in N$ with $Q^{M_0}\varphi_D(x)\in N$ form a subgroup of $N$, since
$\varphi_D$ is additive. For every $j$ it contains $e_jN=e_{(j+1)-1}N$, by
part (3) of that lemma. So it contains $\sum_je_jN=N$.

(c) The $y\in N$ for which $Q^{h_0}y$ lies in the $A$-submodule generated by
$\varphi_D(N)$ form a subgroup of $N$. For every $j$ it contains $e_jN$: for
$y\in e_jN$, part (3) of that lemma puts $Q^{h_0}y$ in the $A$-submodule
generated by $\varphi_D(e_{j-1}N)$, which lies in the one generated by
$\varphi_D(N)$, as $e_{j-1}N\subset N$. So it contains $\sum_je_jN=N$.

(d) and (e) are properties of the Wach module
(Definition~\ref{asm-two-three-killing-wach-frame-exists-data}(4)), for
$\gamma=\bar g$ and, as $V$ is crystalline, for $D(f)$.

With the crystallinity over $G_{K_s}$ and the hypothesis on the filtration,
the Wach hypothesis holds.
% Lean: WachHyp.hypW_crystalline = HypW_of_rec with hC (crysOnCond_all), hj and wClauses,
% whose parts are nBasis (from p1et), the iSup_induction over the letter parts Nl Y i
% = e_i N, with N = sup of the Nl Y i from WachLatticeFree (part (4) of -inputs), and
% P5W (through jumpsLet_const and jump_mem) on each letter part, N_Γ, N_π and map_N.
\end{proof}

\begin{lemma}\label{asm-two-three-killing-wach-frame-exists-decompletion-sandwich}
Let $\ell$, $s$, $K$, $L$ and $D$ be as in
Lemma~\ref{asm-two-three-killing-wach-frame-exists-inputs}, with the Wach data
of $D$ for $(K,s,L)$, $R_A\subset R$ and $\iota_A$
(Definition~\ref{asm-two-three-killing-wach-frame-exists-data}). Let $(V,T)$
be a lattice representation of $G_K$ over $L$, let $M_0,h_0\geq0$ be integers,
and let $N\subset D(T)$ be an $A$-submodule that satisfies the Wach clauses for
$M_0$ and $h_0$. Put $X=\varphi^{s-1}(\pi)\in A$ and $x=\iota_A(X)$. Then
$x^{M_0}j_T(y)\in R_A\otimes_{\mathcal O}T$ for every $y\in N$, and for every
$t\in T$, $x^{h_0}(1\otimes t)$ lies in the $\mathcal O$-submodule of
$R\otimes_{\mathcal O}T$ generated by the $a\,j_T(y)$ with $a\in R_A$ and
$y\in N$.
\end{lemma}
\leanref{Development/AbelianVarietyTheory/Proofs/InertiaEigenvalues/WachLatticeSandwichAinfStatement.lean}{GlobalTheory.T.wach_lattice_sandwich_Ainf}
\leanref{Development/AbelianVarietyTheory/Proofs/InertiaEigenvalues/WachLatticeSandwichAinf.lean}{GlobalTheory.Sandwich.sandwich_of}
\leanref{Development/AbelianVarietyTheory/Proofs/InertiaEigenvalues/WachLatticeSandwichAinf.lean}{GlobalTheory.Sandwich.frobenius_inclusion_X_eq_mul_Q}
\begin{proof}
Write $\iota=\iota_A$. The map $j_T$ satisfies $j_T(ay)=\iota(a)j_T(y)$ for
$a\in A$, since the map $E\to R$ restricts to $\iota$ on $A$, and
$j_T(\varphi_Dy)=(\varphi_R\otimes1)(j_T(y))$.

First, $\varphi_R(x)=x\,\iota(Q)$. In $A$,
$\varphi(1+\pi)=1+\pi q=(1+\pi)^\ell$, so $1+X=(1+\pi)^{\ell^{s-1}}$ and
$XQ=\varphi^{s-1}(\pi q)=(1+\pi)^{\ell^s}-1$. In $R$,
$1+\iota(\pi)=1\otimes[\varepsilon^{(s)}]$, so
$\varphi_R(1+\iota(\pi))=(1+\iota(\pi))^\ell$. Hence
$\varphi_R(x)=(1+\iota(\pi))^{\ell^s}-1=\iota(XQ)=x\,\iota(Q)$.

Second, let $b_1,\ldots,b_\rho$ be an $A$-basis of $N$ whose images
$n_i=j_T(b_i)$ form an $R$-basis of $R\otimes_{\mathcal O}T$ (clause (a)), let
$t_1,\ldots,t_{\rho'}$ be an $\mathcal O$-basis of $T$, and let $Y$ be the
matrix with $n_i=\sum_kY_{ik}(1\otimes t_k)$. Then $\varphi_R\otimes1$ maps $n_i$
to $\sum_k\varphi_R(Y_{ik})(1\otimes t_k)$.

Third, the first assertion. By clause (b),
$Q^{M_0}\varphi_D(b_i)=\sum_jP_{ij}b_j$ with $P_{ij}\in A$, so
$\iota(Q)^{M_0}\varphi_R(Y)=\iota(P)Y$. The matrix $Z=x^{M_0}Y$ satisfies
$\varphi_R(Z)=x^{M_0}\iota(Q)^{M_0}\varphi_R(Y)=\iota(P)Z$ by the first step,
so it has entries in $R_A$ by
Lemma~\ref{asm-two-three-killing-wach-frame-exists-decompletion-frobenius}.
So $x^{M_0}n_i\in R_A\otimes T$, and for $y=\sum_ic_ib_i$,
$x^{M_0}j_T(y)=\sum_i\iota(c_i)x^{M_0}n_i\in R_A\otimes T$.

Fourth, the second assertion. Since $\varphi_D(ay)=\varphi(a)\varphi_D(y)$, the
$A$-submodule generated by $\varphi_D(N)$ is generated by the $\varphi_D(b_j)$,
and clause (c) gives $Q^{h_0}b_i=\sum_jR'_{ij}\varphi_D(b_j)$ with
$R'_{ij}\in A$. So $\iota(Q)^{h_0}Y=\iota(R')\varphi_R(Y)$. Let $U=Y^{-1}$, so
that $1\otimes t_k=\sum_iU_{ki}n_i$. From $\varphi_R(Y)\varphi_R(U)=1$,
$\iota(Q)^{h_0}\varphi_R(U)=U\iota(R')$, and $Z=x^{h_0}U$ satisfies
$\varphi_R(Z)=Z\iota(R')$. By
Lemma~\ref{asm-two-three-killing-wach-frame-exists-decompletion-frobenius}, $Z$
has entries in $R_A$, so
$x^{h_0}(1\otimes t_k)=\sum_iZ_{ki}\,j_T(b_i)$ with $Z_{ki}\in R_A$. For
$t=\sum_kc_kt_k$ with $c_k\in\mathcal O$,
$x^{h_0}(1\otimes t)=\sum_kc_kx^{h_0}(1\otimes t_k)$.
% Lean: Sandwich.sandwich = sandwich_of with frobenius_integrality_tensor_ainf_coefficients and
% frobenius_inclusion_X_eq_mul_Q (φRW_teichmuller); upper for the third step and lower for the
% fourth.
\end{proof}

\begin{lemma}\label{asm-two-three-killing-wach-frame-exists-decompletion-lift}
Let $\ell$, $s$, $K$, $L$, $D$, $(V,T)$, $M_0$, $h_0$, $N$ and $x$ be as in
Lemma~\ref{asm-two-three-killing-wach-frame-exists-decompletion-sandwich}, and
let $G_K$ act on $R_A\otimes_{\mathcal O}T$ and $R\otimes_{\mathcal O}T$ by
$\sigma\otimes\sigma$. Then there is an additive map
$z\colon N\to R_A\otimes_{\mathcal O}T$ such that:
\begin{enumerate}
\item $z(ay)=\iota_A(a)z(y)$ for $a\in A$ and $y\in N$;
\item $x^{M_0}\sigma(z(y))=\sigma(x)^{M_0}z(\bar\sigma y)$ for $\sigma\in G_K$
  and $y\in N$ with $\bar\sigma y\in N$;
\item for $t\in T$, $x^{M_0+h_0}(1\otimes t)$ lies in the additive subgroup
  generated by the $r\,z(y)$ with $r\in R_A$ and $y\in N$.
\end{enumerate}
\end{lemma}
\leanref{Development/AbelianVarietyTheory/Proofs/InertiaEigenvalues/ExistsDecompletionFrameOfLattice.lean}{GlobalTheory.Cyc.Decomp.exists_lift}
\begin{proof}
The map $R_A\otimes_{\mathcal O}T\to R\otimes_{\mathcal O}T$ is injective, as
$T$ is free over $\mathcal O$ and $R_A\subset R$. By
Lemma~\ref{asm-two-three-killing-wach-frame-exists-decompletion-sandwich}, for
$y\in N$ let $z(y)$ be the element of $R_A\otimes T$ with image
$x^{M_0}j_T(y)$; $z$ is additive. Each property is checked after the injective
map to $R\otimes T$. (1): $x^{M_0}j_T(ay)=\iota_A(a)x^{M_0}j_T(y)$. (2): both
sides map to $x^{M_0}\sigma(x)^{M_0}\sigma(j_T(y))$, since
$j_T(\bar\sigma y)=\sigma(j_T(y))$. (3): by the sandwich lemma,
$x^{h_0}(1\otimes t)=\sum c\,a\,j_T(y)$ with $c\in\mathcal O$, $a\in R_A$ and
$y\in N$; multiplying by $x^{M_0}$ gives
$x^{M_0+h_0}(1\otimes t)=\sum (ca)\,z(y)$ after the injective map.
\end{proof}

\begin{lemma}\label{asm-two-three-killing-wach-frame-exists-decompletion-specialization}
Let $\ell$ be a prime and use the notation of
Definition~\ref{asm-two-three-killing-wach-frame}, with $A_{\mathrm{inf}}$,
$\zeta_{\ell^j}$, $\varepsilon^{(j)}$ and $[\,\cdot\,]$ as in
Definition~\ref{asm-two-three-killing-wach-frame-exists-data}. Let
$K,L\subset\bar{\mathbf Q}_\ell$ be finite extensions of $\mathbf Q_\ell$,
$\mathcal O=\mathcal O_L$, $o_1,\ldots,o_f$ a $\mathbf Z_\ell$-basis of
$\mathcal O$, and for $w\in\mathcal O_{\mathbf C}\otimes_{\mathbf Z_\ell}\mathcal O$
let $\kappa_k(w)\in\mathcal O_{\mathbf C}$ be its coordinates in the
$\mathcal O_{\mathbf C}$-basis $(1\otimes o_k)$. Let $(V,T)$ be a lattice
representation of $G_K$ over $L$, let $b_1,\ldots,b_d$ be a
$\mathbf Z_\ell$-basis of $T$, with the matrices $M(\sigma)$ of $\sigma\in G_K$
and $A(c)$ of multiplication by $c\in\mathcal O$, and put
$\lambda(w)=\sum_k\kappa_k(w)A(o_k)$. Let $n\geq0$. Then:
\begin{enumerate}
\item there is a $G_{\mathbf Q_\ell}$-equivariant ring map
  $\theta_n\colon A_{\mathrm{inf}}\to\mathcal O_{\mathbf C}$ with
  $\theta_n([\varepsilon^{(j)}])=\zeta_{\ell^{j+n}}$ for every $j\geq0$;
\item $\lambda$ is a ring map
  $\mathcal O_{\mathbf C}\otimes_{\mathbf Z_\ell}\mathcal O\to
  M_d(\mathcal O_{\mathbf C})$ with $\lambda(1\otimes c)=A(c)$ and
  $\lambda(c'\otimes1)=c'$;
\item with $\theta_n$ applied to the second factor of
  $R_A=\mathcal O\otimes A_{\mathrm{inf}}$, there is an additive map
  $\Phi\colon R_A\otimes_{\mathcal O}T\to\mathbf C_\ell^d$ with
  $\Phi(rv)=\lambda(\theta_n(r))\Phi(v)$, $\Phi(v)\in\mathcal O_{\mathbf C}^d$,
  $\Phi(\sigma v)=M(\sigma)\sigma(\Phi(v))$ for $\sigma\in G_K$, and
  $\Phi(1\otimes b_j)$ the $j$-th unit vector.
\end{enumerate}
\end{lemma}
\leanref{Development/AbelianVarietyTheory/Proofs/InertiaEigenvalues/ExistsFontaineThetaFrobeniusInverseShiftStatement.lean}{GlobalTheory.T.exists_fontaineTheta_frobeniusInverse_shift}
\leanref{Development/AbelianVarietyTheory/Proofs/InertiaEigenvalues/ExistsAinfTensorLatticeSpecialisationStatement.lean}{GlobalTheory.T.exists_ainf_tensor_lattice_specialisation}
\leanref{Development/AbelianVarietyTheory/Proofs/InertiaEigenvalues/OcTensorCoeffActionMatrixMultiplicativeStatement.lean}{GlobalTheory.T.ocTensorCoeffActionMatrix_multiplicative}
\begin{proof}
(1) For $n=0$ take Fontaine's map $\theta$, which is equivariant and maps
$[\varepsilon^{(j)}]$ to $\zeta_{\ell^j}$ \cite[\S4.4, the definition of
$\theta$, Lemma~4.4.1 and Proposition~4.4.2]{BC09}. Given
$\theta_n$, put $\theta_{n+1}=\theta_n\circ W(\varphi^{-1})$, where $\varphi^{-1}$
is the inverse of the Frobenius of the perfect ring $\mathcal O_{\mathbf
C}^\flat$; it commutes with $G_{\mathbf Q_\ell}$ and maps $\varepsilon^{(j)}$
to $\varepsilon^{(j+1)}$, so
$\theta_{n+1}([\varepsilon^{(j)}])=\theta_n([\varepsilon^{(j+1)}])=
\zeta_{\ell^{j+1+n}}$.

(2) Identify $\mathcal O_{\mathbf C}\otimes_{\mathbf Z_\ell}T$ with
$\mathcal O_{\mathbf C}^d$ through the basis $(b_j)$. Multiplication by
$1\otimes c$ is $A(c)$, and $w=\sum_k\kappa_k(w)\otimes o_k$ acts by
$\lambda(w)$. So $\lambda(w)$ is the matrix of multiplication by $w$, which is
multiplicative and additive in $w$.

(3) For $t\in T$ let $\operatorname{co}(t)\in\mathbf Z_\ell^d$ be its
coordinates, so $\operatorname{co}(ct)=A(c)\operatorname{co}(t)$ and
$\operatorname{co}(\sigma t)=M(\sigma)\operatorname{co}(t)$. Put
$\Phi(r\otimes t)=\lambda(\theta_n(r))\operatorname{co}(t)$; this is
well defined on $R_A\otimes_{\mathcal O}T$, as
$\lambda(\theta_n(c\otimes1))=A(c)$ for $c\in\mathcal O$, and
$\Phi(rv)=\lambda(\theta_n(r))\Phi(v)$ by (2). Its values are integral, as
$\lambda(w)\operatorname{co}(t)=\sum_k\kappa_k(w)\operatorname{co}(o_kt)$. For
$\sigma\in G_K$, $\kappa_k(\theta_n(\sigma r))=\sigma(\kappa_k(\theta_n(r)))$ by
(1), and $\sigma$ fixes $\operatorname{co}(o_kt)$, so
\[
 \Phi(\sigma r\otimes\sigma t)=\sum_k\sigma(\kappa_k(\theta_n(r)))
 M(\sigma)\operatorname{co}(o_kt)=M(\sigma)\sigma(\Phi(r\otimes t)).
\]
Finally $\Phi(1\otimes b_j)=\operatorname{co}(b_j)$ is the $j$-th unit vector.
% Lean: Cyc.Decomp.exists_isThetaShift (exists_thetaShift_ringHom by induction on n from
% exists_isTheta, the datum's theta on A_inf, read in Fontaine's sense; 216 lines) and
% exists_specPhi (specPhi_co, specPhi_generic, specPhi_abstract, lam_mul_of_isMat; 333 lines).
\end{proof}

\begin{lemma}\label{asm-two-three-killing-wach-frame-exists-decompletion-tower-surjective}
Let $\ell$ be a prime and use the notation of
Definition~\ref{asm-two-three-killing-wach-frame}, and let
$K\subset\bar{\mathbf Q}_\ell$ be a finite extension of $\mathbf Q_\ell$. For
every $u\in\mathbf Q_\ell$ with $|u|=1$ there is $\sigma\in G_{K_0}$ with
$\chi(\sigma)=u$.
\end{lemma}
\leanref{Development/AbelianVarietyTheory/Proofs/InertiaEigenvalues/FinrankCyclotomicExtAndCyclotomicCharSurjectiveStatement.lean}{GlobalTheory.T.finrank_cyclotomicExt_and_cyclotomicChar_surjective}
\leanref{Development/AbelianVarietyTheory/Proofs/InertiaEigenvalues/FinrankCyclotomicExtAndCyclotomicCharSurjective.lean}{GlobalTheory.Cyc.exists_chi_eq}
\leanref{Development/AbelianVarietyTheory/Proofs/InertiaEigenvalues/FinrankCyclotomicExtAndCyclotomicCharSurjective.lean}{GlobalTheory.Cyc.minpoly_unramifiedSubfield_eq_cyclotomic}
\leanref{Development/AbelianVarietyTheory/Proofs/InertiaEigenvalues/FinrankCyclotomicExtAndCyclotomicCharSurjective.lean}{GlobalTheory.Cyc.exists_chi_approx}
\begin{proof}
First, let $j\geq1$ and let $\zeta\in\bar{\mathbf Q}_\ell$ be a primitive
$\ell^j$-th root of unity; then the minimal polynomial of $\zeta$ over $K_0$ is
the cyclotomic polynomial $\Phi_{\ell^j}$, of degree
$N=\ell^{j-1}(\ell-1)$. It divides $\Phi_{\ell^j}$, as $\Phi_{\ell^j}(\zeta)=0$.
Suppose that its degree is less than $N$. Then $p(X)=f(X+1)$, for $f$ the
minimal polynomial, is a nonzero polynomial over $K_0$ of degree less than $N$
with $p(\zeta-1)=0$. Now $|\zeta-1|=\ell^{-1/N}$, and the nonzero elements of
the unramified field $K_0$ have absolute values in $\ell^{\mathbf Z}$, so the
nonzero terms $p_i(\zeta-1)^i$, $0\leq i<N$, have absolute values
$\ell^{a_i-i/N}$ with $a_i\in\mathbf Z$, which are pairwise distinct. So
$|p(\zeta-1)|$ is the largest of them, which is not $0$: a contradiction.

Second, for every $j\geq1$ there is $\sigma\in G_{K_0}$ with
$|\chi(\sigma)-u|\leq\ell^{-j}$. Choose an integer $a\geq0$ with
$|u-a|\leq\ell^{-j}$; then $|a|=1$, so $a$ is prime to $\ell$, and for a
primitive $\ell^j$-th root of unity $\zeta$, $\zeta^a$ is again one. By the
first step $\zeta$ and $\zeta^a$ have the same minimal polynomial over $K_0$, so
there is $\sigma\in G_{K_0}$ with $\sigma(\zeta)=\zeta^a$. Also
$\sigma(\zeta)=\zeta^b$ for an integer $b\geq0$ with
$|\chi(\sigma)-b|\leq\ell^{-j}$, so $a\equiv b$ modulo $\ell^j$, and
$|\chi(\sigma)-u|\leq\max(|\chi(\sigma)-b|,|b-a|,|a-u|)\leq\ell^{-j}$.

Third, $\chi$ is continuous
(Lemma~\ref{asm-two-three-killing-wach-torsion-cyclotomic}(1)) and $G_{K_0}$ is
compact, as a closed subgroup of the profinite group $G_{\mathbf Q_\ell}$. So
$\chi(G_{K_0})$ is compact, hence closed in $\mathbf Q_\ell$, and by the second
step it contains $u$.
% Lean: Cyc.exists_chi_eq (closure of the compact image, hcl.isCompact.image hcont) with
% exists_chi_approx (PadicInt.appr, IsConjRoot.exists_algEquiv from
% minpoly_unramifiedSubfield_eq_cyclotomic, IsCyc for the second b) and
% minpoly_unramifiedSubfield_eq_cyclotomic (aeval_ne_zero_of_norm_eq with
% T.norm_eq_zpow_of_mem_maxUnramifiedSubfield and norm_primitiveRoot_primePow_sub_one).
% 27 theorems, about 570 lines with the definitions.
\end{proof}

\begin{lemma}\label{asm-two-three-killing-wach-frame-exists-decompletion-tower-equivariant}
Let $\ell$, $s$, $K$, $L$ and $D$ be as in
Lemma~\ref{asm-two-three-killing-wach-frame-exists-inputs}, with the Wach data
of $D$ for $(K,s,L)$, $R_A\subset R$ and $\iota_A$
(Definition~\ref{asm-two-three-killing-wach-frame-exists-data}), and let $G_K$
act on $R_A$ and $R$ through the second factor. Then
$\iota_A(\bar\sigma a)=\sigma(\iota_A(a))$ for every $\sigma\in G_K$ and every
$a\in A$.
\end{lemma}
\leanref{Development/AbelianVarietyTheory/Proofs/InertiaEigenvalues/IoEquivariantStatement.lean}{GlobalTheory.T.io_equivariant_alt}
\leanref{Development/AbelianVarietyTheory/Proofs/InertiaEigenvalues/IoEquivariant.lean}{GlobalTheory.Cyc.Decomp.etTie_io_act}
\leanref{Development/AbelianVarietyTheory/Proofs/InertiaEigenvalues/IoEquivariant.lean}{GlobalTheory.Cyc.Decomp.tmul_gen_cancel}
\begin{proof}
Let $(V_1,T_1)$, $e_0$ and $e\in D(T_1)$ be as in
Definition~\ref{asm-two-three-killing-wach-frame-exists-data}(6), and put
$j=j_{T_1}$. The map $j$ is semilinear for $E\to R$, whose restriction to $A$ is
$\iota_A$, so $j(ae)=\iota_A(a)\otimes e_0$ for $a\in A$. Let $\sigma\in G_K$.
As $\bar\sigma$ acts semilinearly on $D(T_1)$, $\bar\sigma(ae)=(\bar\sigma
a)(\bar\sigma e)$, and as $j$ carries $\bar\sigma$ to $\sigma\otimes\sigma$ and
$\sigma e_0=e_0$, $j(\bar\sigma e)=1\otimes e_0$. So
\[
 \iota_A(\bar\sigma a)\otimes e_0=j(\bar\sigma(ae))
 =(\sigma\otimes\sigma)(\iota_A(a)\otimes e_0)=\sigma(\iota_A(a))\otimes e_0
\]
in $R\otimes_{\mathcal O}T_1$. As $T_1$ is free over $\mathcal O$ with basis
$e_0$, $r\mapsto r\otimes e_0$ is injective on $R$, so
$\iota_A(\bar\sigma a)=\sigma(\iota_A(a))$ in $R$, hence in $R_A$.
% Lean: Decomp.io_equivariant = etTie_io_act (E.one_e, W.one_line, W.one_triv; jD_act and
% γD_smul for a • e, tmul_gen_cancel for the cancellation of e_0). 66 lines.
\end{proof}

\begin{lemma}\label{asm-two-three-killing-wach-frame-exists-decompletion-tower-idempotent}
Let $\ell$ be a prime and use the notation of
Definition~\ref{asm-two-three-killing-wach-frame}. Let
$K,L\subset\bar{\mathbf Q}_\ell$ be finite extensions of $\mathbf Q_\ell$,
$\mathcal O=\mathcal O_L$, $o_1,\ldots,o_f$ a $\mathbf Z_\ell$-basis of
$\mathcal O$, and for $w\in\mathcal O_{\mathbf C}\otimes_{\mathbf Z_\ell}\mathcal O$
let $\kappa_k(w)\in\mathcal O_{\mathbf C}$ be its coordinates in the
$\mathcal O_{\mathbf C}$-basis $(1\otimes o_k)$. Let $e$ be an idempotent of
$\mathcal O_{\mathbf C}\otimes_{\mathbf Z_\ell}\mathcal O$ with
$\sigma(\kappa_k(e))=\kappa_k(e)$ for every $\sigma\in G_K$ and every $k$. Then
$\kappa_k(e)\in K$ for every $k$.
\end{lemma}
\leanref{Development/AbelianVarietyTheory/Proofs/InertiaEigenvalues/KapIdemMemInCStatement.lean}{GlobalTheory.T.kap_mem_inC_of_invariant}
\leanref{Development/AbelianVarietyTheory/Proofs/InertiaEigenvalues/KapIdemMemInC.lean}{GlobalTheory.Cyc.Decomp.kap_idem_algebraic}
\begin{proof}
First, the $\kappa_k(e)$ lie in $\bar{\mathbf Q}_\ell$. Let
$\tau_1,\ldots,\tau_f$ be the $\mathbf Q_\ell$-embeddings of $L$ into
$\bar{\mathbf Q}_\ell$; there are $f=[L:\mathbf Q_\ell]$ of them, and the
matrix $P=(\tau_j(o_k))_{k,j}$ is invertible, as the square of its determinant
is the discriminant of the basis $(o_k)$, which is not $0$. For each $j$, the
ring map $\mathcal O_{\mathbf C}\otimes\mathcal O\to\mathbf C_\ell$,
$c\otimes a\mapsto c\,\tau_j(a)$, maps $e$ to an idempotent of the field
$\mathbf C_\ell$, that is to $\varepsilon_j\in\{0,1\}$, and $e=\sum_k\kappa_k(e)
\otimes o_k$ to $\sum_k\kappa_k(e)\tau_j(o_k)$. So the row vector
$(\kappa_k(e))_k$ is $(\varepsilon_j)_jP^{-1}$, which has entries in
$\bar{\mathbf Q}_\ell$.

Second, an element of $\bar{\mathbf Q}_\ell$ fixed by $G_K$ lies in $K$, by
Galois theory, and $G_{\mathbf Q_\ell}$ acts on $\bar{\mathbf Q}_\ell\subset
\mathbf C_\ell$ through its action on $\bar{\mathbf Q}_\ell$.
% Lean: Decomp.kap_idem_mem_inC: kap_idem_algebraic (exists_emb_det_ne_zero through
% Algebra.discr_eq_det_embeddingsMatrixReindex_pow_two, exists_evEmb, evEmb_eq_sum,
% IsIdempotentElem.iff_eq_zero_or_one) and InfiniteGalois.fixedField_fixingSubgroup.
% About 140 lines.
\end{proof}

\begin{lemma}\label{asm-two-three-killing-wach-frame-exists-decompletion-tower-gamma}
Let $\ell$, $s$, $K$ and $L$ be as in
Definition~\ref{asm-two-three-killing-wach-frame-exists-data}, with Wach data
for $(K,s,L)$, and let $I\subset A$ be the ideal $(\ell,\pi)$. Let
$g\in G_{K_s}$. Then $\bar ga-a\in I^{k+1}$ for every $k\geq0$ and every
$a\in I^k$.
\end{lemma}
\leanref{Development/AbelianVarietyTheory/Proofs/InertiaEigenvalues/GammaActionSubMemSpanPowSuccStatement.lean}{GlobalTheory.T.gammaAction_sub_mem_span_pow_succ}
\leanref{Development/AbelianVarietyTheory/Proofs/InertiaEigenvalues/GammaActionSubMemSpanPowSucc.lean}{GlobalTheory.GamFilt.gamma_sub_mem_pow_succ}
\leanref{Development/AbelianVarietyTheory/Proofs/InertiaEigenvalues/GammaActionSubMemSpanPowSucc.lean}{GlobalTheory.GamFilt.sub_mem_pow_succ}
\begin{proof}
Put $\gamma=\bar g$, acting on every factor $\mathcal O[[\pi]]$ of $A$ by the
endomorphism $\gamma_0$ of
Definition~\ref{asm-two-three-killing-wach-frame-exists-data}(2). Since
$g\in G_{K_s}\subset G_{K_0}$ and $s\geq1$,
Lemma~\ref{asm-two-three-killing-wach-torsion-cyclotomic}(1) gives
$|\chi(g)-1|\leq\ell^{-1}$, so $\chi_\Gamma(\gamma)=\chi(g)=1+\ell c$ with
$c\in\mathbf Z_\ell$, and $\gamma_0(\pi)=\pi(1+\ell c+\pi h)$ for some
$h\in\mathcal O[[\pi]]$.

First, $\gamma a-a\in\pi A\subset I$ for every $a\in A$: in each factor,
$a=a_0+\pi a'$ with $a_0\in\mathcal O$, and $\gamma_0(a)-a=\gamma_0(\pi)
\gamma_0(a')-\pi a'$ lies in $\pi\mathcal O[[\pi]]$. Second, $\gamma\ell=\ell$
and $\gamma\pi-\pi=\pi(\ell c+\pi h)\in I^2$. Third, for $y=a\ell+b\pi\in I$,
\[
 \gamma y-y=(\gamma a-a)\ell+(\gamma b-b)\pi+\gamma b\,(\gamma\pi-\pi)\in I^2,
\]
by the first two steps. Finally, induction on $k$: the case $k=0$ is the first
step, and $I^{k+1}$ is generated as an abelian group by products $xy$ with
$x\in I^k$ and $y\in I$, for which
$\gamma(xy)-xy=(\gamma x-x)\gamma y+x(\gamma y-y)$ lies in $I^{k+1}\cdot I+
I^k\cdot I^2$, as $\gamma y=y+(\gamma y-y)\in I$.
% Lean: GamFilt.gamma_sub_mem_pow_succ: GKs_def and χΓ_res with exists_eq_one_add for
% χ(g) = 1 + l c, γ₀_X, and letγ_sub_mem (X_dvd_sub, exists_X_sub,
% sub_mem_pow_succ_pair, sub_mem_pow_succ). 124 lines.
\end{proof}

\begin{lemma}\label{asm-two-three-killing-wach-frame-exists-decompletion-tower-iterate}
Let $\ell$ be a prime, $A$ a commutative ring, $I\subset A$ an ideal, $N$ an
$A$-module, $\gamma$ a ring endomorphism of $A$ with $\gamma a-a\in I^{k+1}$ for
every $k\geq0$ and $a\in I^k$, and $Y\colon N\to N$ an additive map with
$Y(ax)=\gamma(a)Y(x)$ and $Y(x)-x\in IN$ for every $a\in A$ and $x\in N$. Then
for every $b\geq1$ and $x\in N$, $Y^{\ell^b}(x)-x\in\big((\ell^2)+
I^{\ell^{b-1}}\big)N$.
\end{lemma}
\leanref{Development/AbelianVarietyTheory/Proofs/InertiaEigenvalues/SemilinearIteratePrimePowSubMemFiltrationStatement.lean}{GlobalTheory.T.semilinear_iterate_prime_pow_sub_mem_filtration}
\leanref{Development/AbelianVarietyTheory/Proofs/InertiaEigenvalues/SemilinearIteratePrimePowSubMemFiltration.lean}{GlobalTheory.semilinear_iterate_prime_pow_sub_mem}
\leanref{Development/AbelianVarietyTheory/Proofs/InertiaEigenvalues/OneAddXPowPrimePowSubOneMemSpan.lean}{GlobalTheory.one_add_X_pow_prime_pow_sub_one_mem_span}
\begin{proof}
Put $\Delta=Y-1$, an endomorphism of the abelian group $N$.

First, $\Delta(I^kN)\subset I^{k+1}N$ for every $k\geq0$, so $\Delta^k(N)\subset
I^kN$. Indeed $I^kN$ is generated by the $ay$ with $a\in I^k$ and $y\in N$, and
$\Delta(ay)=\gamma(a)\Delta(y)+(\gamma a-a)y$, where $\gamma(a)\in I^k$,
$\Delta(y)\in IN$ and $\gamma a-a\in I^{k+1}$.

Second, $(1+X)^{\ell^b}-1=\ell^2u(X)+X^{\ell^{b-1}}v(X)$ for polynomials $u,v$
with integer coefficients. Indeed the coefficient $\binom{\ell^b}{i}$ of $X^i$,
for $0<i<\ell^{b-1}$, is divisible by $\ell^2$: from
$i\binom{\ell^b}{i}=\ell^b\binom{\ell^b-1}{i-1}$ and $v_\ell(i)\leq b-2$,
$v_\ell\binom{\ell^b}{i}\geq2$. The terms with $i\geq\ell^{b-1}$ are divisible
by $X^{\ell^{b-1}}$.

Finally, substituting $\Delta$ for $X$, as $Y=1+\Delta$,
$Y^{\ell^b}(x)-x=\ell^2\,u(\Delta)(x)+\Delta^{\ell^{b-1}}\big(v(\Delta)(x)\big)$,
whose first term lies in $\ell^2N$ and whose second lies in $I^{\ell^{b-1}}N$ by
the first step.
% Lean: semilinear_iterate_prime_pow_sub_mem (semilinear_sub_mem_pow_succ_smul,
% semilinear_sub_one_pow_mem, one_add_X_pow_prime_pow_sub_one_mem_span with
% sq_dvd_choose_prime_pow, Polynomial.aeval at D = Y - 1). 111 lines.
\end{proof}

\begin{lemma}\label{asm-two-three-killing-wach-frame-exists-decompletion-tower}
Let $\ell$, $s$, $K$, $L$ and $D$ be as in
Lemma~\ref{asm-two-three-killing-wach-frame-exists-inputs}, with the Wach data
of $D$ for $(K,s,L)$, and $\iota_A$. Let $n\geq0$, put $m=2s+n+1$, and let
$\theta_n$, $o_k$ and $\kappa_k$ be as in
Lemma~\ref{asm-two-three-killing-wach-frame-exists-decompletion-specialization}.
Put $\phi=\theta_n\circ\iota_A\colon A\to\mathcal O_{\mathbf C}\otimes
\mathcal O$ and let $J\subset A$ be the ideal $(\ell^2)+(\ell,\pi)^{\ell^{s+n}}$.
Then:
\begin{enumerate}
\item there is $g_0\in G_{K_s}$ such that $g=g_0^{\ell^{s+n+1}}$ lies in
  $G_{K_m}$ and $|\chi(g)-1|=\ell^{-m}$;
\item $\bar\sigma=1$ for $\sigma\in G_{K_\infty}$;
\item $\kappa_k(\phi(a))$ lies in $K_m$ and has absolute value at most $1$, for
  every $a\in A$ and every $k$;
\item $|\kappa_k(\phi(a))|<\ell^{-1/(\ell-1)}$ for every $a\in J$ and every $k$;
\item for $g_0$ as in (1), if $N$ is an $A$-submodule of some $D(T)$ that
  satisfies clause (d) of the Wach clauses, then
  $\bar g_0^{\ell^{s+n+1}}y-y\in JN$ for every $y\in N$;
\item if $H\geq0$ is an integer and $n=n_0(\ell,H)$, the least $n'\geq0$ with
  $H\leq\ell^{n'}(\ell-1)$, then $\ell=u(\zeta_{\ell^{n+1}}-1)^H$ for some
  $u\in\mathcal O_{\mathbf C}$.
\end{enumerate}
\end{lemma}
\leanref{Development/AbelianVarietyTheory/Proofs/InertiaEigenvalues/ExistsCyclotomicTowerLevelGeneratorPropertiesStatement.lean}{GlobalTheory.T.exists_cyclotomicTowerLevelGenerator_properties}
\leanref{Development/AbelianVarietyTheory/Proofs/InertiaEigenvalues/CyclotomicLevelSubgroupMemWachSubgroupsStatement.lean}{GlobalTheory.T.cyclotomicLevelSubgroup_mem_wachSubgroups}
\leanref{Development/AbelianVarietyTheory/Proofs/InertiaEigenvalues/OcTensorCoeffBasisCoordsMemCyclotomicLevelOfPoweStatement.lean}{GlobalTheory.T.ocTensorCoeffBasisCoords_mem_cyclotomicLevel_of_powerSeries}
\leanref{Development/AbelianVarietyTheory/Proofs/InertiaEigenvalues/NormLtRho1OfMemIdealPowerStatement.lean}{GlobalTheory.T.norm_lt_rho1_of_mem_ideal_power}
\leanref{Development/AbelianVarietyTheory/Proofs/InertiaEigenvalues/GammaActionSubMemSpanPowSuccStatement.lean}{GlobalTheory.T.gammaAction_sub_mem_span_pow_succ}
\leanref{Development/AbelianVarietyTheory/Proofs/InertiaEigenvalues/ExistsPrimeEqMulPrimitiveRootSubOnePowStatement.lean}{GlobalTheory.T.exists_prime_eq_mul_primitiveRoot_sub_one_pow'}
\begin{proof}
(1) By
Lemma~\ref{asm-two-three-killing-wach-frame-exists-decompletion-tower-surjective},
for $u=1+\ell^s$, there is $g_0\in G_{K_0}$ with $\chi(g_0)=1+\ell^s$. As
$|\chi(g_0)-1|=\ell^{-s}$, $g_0\in G_{K_s}$ by
Lemma~\ref{asm-two-three-killing-wach-torsion-cyclotomic}(1). Put
$g=g_0^{\ell^{s+n+1}}$. As $s\geq1$, and $s\geq2$ if $\ell=2$, part (2) of that
lemma gives $|\chi(g)-1|=|\chi(g_0)^{\ell^{s+n+1}}-1|=\ell^{-s-(s+n+1)}=
\ell^{-m}$; and $g\in G_{K_0}$, so $g\in G_{K_m}$ by part (1).

(2) The character $\chi_\Gamma$
(Definition~\ref{asm-two-three-killing-wach-frame-exists-data}(2)) is
injective, and $\chi_\Gamma(\bar\sigma)=\chi(\sigma)=1$ for
$\sigma\in G_{K_\infty}$, so $\bar\sigma=1$ there.

(3) Put $c=\zeta_{\ell^{s+n}}-1$. First, $\phi(a')=1\otimes a'$ for
$a'\in\mathcal O$, as $\iota_A$ is a map of $\mathcal O$-algebras and $\theta_n$
acts on the factor $A_{\mathrm{inf}}$; and $\phi(\pi)=c\otimes1$, as
$\iota_A(\pi)=1\otimes([\varepsilon^{(s)}]-1)$ and
$\theta_n([\varepsilon^{(s)}])=\zeta_{\ell^{s+n}}$
(Lemma~\ref{asm-two-three-killing-wach-frame-exists-decompletion-specialization}(1)).
Since $1\leq s+n\leq m$, $c$ lies in $K_{s+n}\subset K_m$ and
$|c|=\ell^{-1/(\ell^{s+n-1}(\ell-1))}<1$.

Second, $\kappa_k(\phi(e_i))\in K_m$ for every idempotent $e_i$ of
Definition~\ref{asm-two-three-killing-wach-frame-exists-data}(1) and every $k$.
The element $\phi(e_i)$ is an idempotent, as $e_i$ is. For $\sigma\in G_K$,
$\bar\sigma e_i=e_i$, as $\bar\sigma$ acts on every factor of $A$ by the same
endomorphism, which fixes $0$ and $1$; so
$\sigma(\iota_A(e_i))=\iota_A(e_i)$ by
Lemma~\ref{asm-two-three-killing-wach-frame-exists-decompletion-tower-equivariant},
and as $\theta_n$ is $G_{\mathbf Q_\ell}$-equivariant,
$\sigma(\kappa_k(\phi(e_i)))=\kappa_k(\theta_n(\sigma\,\iota_A(e_i)))=
\kappa_k(\phi(e_i))$. By
Lemma~\ref{asm-two-three-killing-wach-frame-exists-decompletion-tower-idempotent},
$\kappa_k(\phi(e_i))\in K\subset K_s\subset K_m$.

Third, for every $N\geq0$, every $a\in A$ and every $k$,
$\kappa_k(\phi(a))=t+c^Nr$ with $t\in K_m$ and $r\in\mathcal O_{\mathbf C}$. We
argue by induction on $N$; for $N=0$ take $t=0$. For the inductive step write
$a=\sum_ie_ia_i+\pi a'$, with $a_i\in\mathcal O$ the constant term of the $i$-th
component of $a$ and $a'\in A$. Then
$\phi(a)=\sum_i\phi(e_i)(1\otimes a_i)+(c\otimes1)\phi(a')$. If
$o_ja_i=\sum_k\lambda_{ijk}o_k$ with $\lambda_{ijk}\in\mathbf Z_\ell$, then
$\kappa_k(\phi(e_i)(1\otimes a_i))=\sum_j\kappa_j(\phi(e_i))\lambda_{ijk}$, which
lies in $K_m$ by the second step, and $\kappa_k((c\otimes1)w)=c\,\kappa_k(w)$.
With $\kappa_k(\phi(a'))=t'+c^Nr'$ from the induction hypothesis,
$\kappa_k(\phi(a))=t+c^{N+1}r'$ for
$t=\sum_{i,j}\kappa_j(\phi(e_i))\lambda_{ijk}+ct'\in K_m$.

Finally, let $a\in A$, let $\delta>0$ and choose $N$ with $|c|^N<\delta$.
With $t$ and $r$ as in the third step, $|\kappa_k(\phi(a))-t|\leq|c|^N<\delta$,
and $|t|\leq\max(|\kappa_k(\phi(a))|,|c^Nr|)\leq1$, as $\kappa_k(\phi(a))\in
\mathcal O_{\mathbf C}$. So $\kappa_k(\phi(a))$ is a limit of elements of the
valuation ring of $K_m$, which is compact, as $K_m$ is a finite extension of
$\mathbf Q_\ell$, hence closed in $\mathbf C_\ell$. So $\kappa_k(\phi(a))$ lies
in it.

(4) Let $a\in J$ and put $z=\zeta_{\ell^{s+n}}-1$, so that $\phi(\pi)=z\otimes1$
by the first step of (3). As $s+n\geq1$, $|z|=\ell^{-1/(\ell^{s+n-1}(\ell-1))}
\geq\ell^{-1}=|\ell|$, so $\ell=zu$ with $u\in\mathcal O_{\mathbf C}$. So
$\phi(\ell)=(z\otimes1)(u\otimes1)$ and $\phi(\pi)$ lie in the ideal generated
by $z\otimes1$, and $\phi(a)=\ell^2x+(z^{\ell^{s+n}}\otimes1)y$ for some
$x,y\in\mathcal O_{\mathbf C}\otimes\mathcal O$. Then
$\kappa_k(\phi(a))=\ell^2\kappa_k(x)+z^{\ell^{s+n}}\kappa_k(y)$ has absolute
value at most $\max(\ell^{-2},|z|^{\ell^{s+n}})=\max(\ell^{-2},
\ell^{-\ell/(\ell-1)})$, which is less than $\ell^{-1/(\ell-1)}$, as
$1/(\ell-1)<2$ and $1/(\ell-1)<\ell/(\ell-1)$.

(5) Let $N\subset D(T)$ satisfy clause (d) and put $I=(\ell,\pi)\subset A$. As
$g_0\in G_{K_s}$, clause (d) gives $\bar g_0x\in N$ and $\bar g_0x-x\in\pi N
\subset IN$ for $x\in N$. Let $Y$ be the restriction of $\bar g_0$ to $N$; it is
additive and $Y(ax)=(\bar g_0a)Y(x)$, as $\Gamma$ acts semilinearly on $D(T)$
(Definition~\ref{asm-two-three-killing-wach-frame-exists-data}(4)). By
Lemma~\ref{asm-two-three-killing-wach-frame-exists-decompletion-tower-gamma},
$\gamma=\bar g_0$ on $A$ satisfies the hypothesis of
Lemma~\ref{asm-two-three-killing-wach-frame-exists-decompletion-tower-iterate},
which, for $b=s+n+1$, gives $Y^{\ell^{s+n+1}}(y)-y\in\big((\ell^2)+
I^{\ell^{s+n}}\big)N=JN$. As $\Gamma$ acts on $D(T)$ by a group action,
$Y^{\ell^{s+n+1}}(y)=\bar g_0^{\ell^{s+n+1}}y$.

(6) The element $\zeta_{\ell^{n+1}}-1$ has absolute value
$\ell^{-1/(\ell^n(\ell-1))}$, so $|(\zeta_{\ell^{n+1}}-1)^H|\geq\ell^{-1}$ when
$H\leq\ell^n(\ell-1)$, and $u=\ell/(\zeta_{\ell^{n+1}}-1)^H$ is integral.
% Lean: (1) is the first two parts of exists_cyclotomicTowerLevelGenerator_properties
% (= Cyc.gen_dense: exists_chi_eq for u = 1 + l^s and cyc_basic for g0, then
% norm_pow_pow_sub_one for the powers); its density parts are not used here. (2) is
% cyclotomicLevelSubgroup_mem_wachSubgroups (76 lines, through GKs_def and χΓ_inj). (3) is
% ocTensorCoeffBasisCoords_mem_cyclotomicLevel_of_powerSeries (= Decomp.Coord.coord, with
% let_split, kap_mul_one_tmul, kap_tmul_one_mul and the closedness of intC (Ks K m),
% TateDecomp.isCompact_intC; 172 lines) with hφa, hφπ, norm_zeta_sub_one_lt, zeta_sub_one_mem
% and hidem (kap_mem_inC_of_invariant with io_equivariant_alt, kap_thR_gRA and γA_letE) in
% exists_decompletion. (4) is norm_lt_rho1_of_mem_ideal_power (=
% Decomp.norm_kap_lt_rho1_of_mem_sq_sup_pow with exists_map_eq_of_mem_sq_sup_pow; 162 lines).
% (5) is Decomp.γD_pow_sub with semilinear_iterate_prime_pow_sub_mem_filtration and
% gammaAction_sub_mem_span_pow_succ, for the generator resΓ g0 (in W.GKs by
% cyclotomicLevelSubgroup_mem_wachSubgroups) and j = s + n. (6) is
% exists_prime_eq_mul_primitiveRoot_sub_one_pow' (74 lines).
\end{proof}

\begin{lemma}\label{asm-two-three-killing-wach-frame-exists-decompletion-matrix}
Let $\ell$ be a prime and use the notation of
Definition~\ref{asm-two-three-killing-wach-frame}. Let
$K\subset\bar{\mathbf Q}_\ell$ be a finite extension of $\mathbf Q_\ell$, let
$m\geq1$ be an integer, with $m\geq2$ if $\ell=2$, and let $g\in G_{K_m}$ with
$|\chi(g)-1|=\ell^{-m}$. Let $M\colon G_{K_m}\to M_d(\mathbf C_\ell)$ have
continuous entries, with $M(\sigma\tau)=M(\sigma)M(\tau)$ and
$\sigma(M(\tau))=M(\tau)$ for $\sigma,\tau\in G_{K_m}$. Let $\Delta$ be a
$d\times d$ matrix with entries in $\mathcal O_{\mathbf C}$ such that:
$\Delta X=\ell$ for some $X$ with entries in $\mathcal O_{\mathbf C}$; for
$\sigma\in G_{K_m}$, $M(\sigma)\sigma(\Delta)=\Delta X_\sigma$ for some
$X_\sigma$ with entries in the valuation ring of $K_m$; $M(\sigma)\sigma(\Delta)
=\Delta$ for $\sigma\in G_{K_\infty}$; and $M(g)g(\Delta)-\Delta=\Delta X_g$
for some $X_g$ with entries of absolute value less than $\ell^{-1/(\ell-1)}$.
Then $\Delta$ is invertible, $\ell\Delta^{-1}$ has entries in
$\mathcal O_{\mathbf C}$, and with $U_0(\sigma)=\Delta^{-1}M(\sigma)\sigma(\Delta)$,
$(U_0,g)$ is a frame of rank $d$ at level $m$ for $K$ and
$\Delta U_0(\sigma)=M(\sigma)\sigma(\Delta)$ for $\sigma\in G_{K_m}$.
\end{lemma}
\leanref{Development/AbelianVarietyTheory/Proofs/InertiaEigenvalues/DecompletionFrameOfMatrixFormStatement.lean}{GlobalTheory.T.decompletion_frame_of_matrix_form}
\leanref{Development/AbelianVarietyTheory/Proofs/InertiaEigenvalues/DecompletionFrameOfMatrixForm.lean}{GlobalTheory.Cyc.Decomp.isDec_of_matrix}
\begin{proof}
Since $\Delta(\ell^{-1}X)=1$, $\Delta$ is invertible with
$\ell\Delta^{-1}=X$. Then $U_0(\sigma)=X_\sigma$ for $\sigma\in G_{K_m}$, with
entries in the valuation ring of $K_m$; its entries are continuous, as those of
$M$ and the maps $\sigma\mapsto\sigma(y)$ are; $U_0=1$ on $G_{K_\infty}$; and
$U_0(g)-1=X_g$. For $\sigma,\tau\in G_{K_m}$, $\sigma$ fixes $X_\tau$, so
\[
 U_0(\sigma\tau)=\Delta^{-1}M(\sigma)\sigma\big(M(\tau)\tau(\Delta)\big)
 =\Delta^{-1}M(\sigma)\sigma(\Delta X_\tau)=U_0(\sigma)U_0(\tau).
\]
\end{proof}

\begin{lemma}\label{asm-two-three-killing-wach-frame-exists-decompletion}
Let $\ell$, $s$, $K$, $L$ and $D$ be as in
Lemma~\ref{asm-two-three-killing-wach-frame-exists-inputs}, with the Wach data
of $D$ for $(K,s,L)$. For an integer $H\geq0$ let $n_0(\ell,H)$ be the least
$n\geq0$ with $H\leq\ell^n(\ell-1)$. Let $(V,T)$ be a lattice representation of
$G_K$ over $L$ and $M_0,h_0\geq0$ integers such that the Wach hypothesis holds
for $(V,T)$, $M_0$ and $h_0$, and put $m=2s+n_0(\ell,M_0+h_0)+1$. Let
$b_1,\ldots,b_d$ be a $\mathbf Z_\ell$-basis of $T$, and in this basis let
$M(\sigma)$ be the matrix of $\sigma\in G_K$ and $A(c)$ that of multiplication
by $c\in\mathcal O_L$. Then there are an invertible $d\times d$ matrix $\Delta$
over $\mathbf C_\ell$, a map $U_0$ from $G_{K_m}$ to the $d\times d$ matrices
over $\mathbf C_\ell$ and $g\in G_{K_m}$ such that: the entries of $\Delta$ and
of $\ell\Delta^{-1}$ lie in $\mathcal O_{\mathbf C}$; $(U_0,g)$ is a frame of
rank $d$ at level $m$ for $K$; $\Delta U_0(\sigma)=M(\sigma)\,\sigma(\Delta)$ for
every $\sigma\in G_{K_m}$; and the entries of $\Delta^{-1}A(c)\Delta$ lie in
$K_m$ and have absolute value at most $1$, for every $c\in\mathcal O_L$.
\end{lemma}
\leanref{Development/AbelianVarietyTheory/Proofs/InertiaEigenvalues/ExistsDecompletionFrameOfLattice.lean}{GlobalTheory.Cyc.Decomp.exists_decompletion}
\leanref{Development/AbelianVarietyTheory/Definitions/NonGaloisDescent/ExpStar.lean}{GlobalTheory.Cyc.IsDec}
\begin{proof}
Let $N\subset D(T)$ be an $A$-submodule that satisfies the Wach clauses for
$M_0$ and $h_0$, and $b'_1,\ldots,b'_\rho$ an $A$-basis of $N$ as in clause (a).
Put $n=n_0(\ell,M_0+h_0)$, so that $m=s+(s+n+1)$, $X=\varphi^{s-1}(\pi)$,
$x=\iota_A(X)$ and $\varpi=\zeta_{\ell^{n+1}}-1\neq0$. Let $o_1,\ldots,o_f$ be a
$\mathbf Z_\ell$-basis of $\mathcal O=\mathcal O_L$. Since
$j_T(b'_1),\ldots,j_T(b'_\rho)$ is an $R$-basis of $R\otimes_{\mathcal O}T$,
$\rho$ is the rank of $T$ over $\mathcal O$, and $d=f\rho$; index
$\{1,\ldots,d\}$ by the pairs $(k,i)$.

First, Lemma~\ref{asm-two-three-killing-wach-frame-exists-decompletion-lift}
gives $z\colon N\to R_A\otimes_{\mathcal O}T$, and
Lemma~\ref{asm-two-three-killing-wach-frame-exists-decompletion-specialization}
gives $\theta_n$, $\lambda$ and $\Phi$. Put $\Psi=\Phi\circ z$ and
$\phi=\theta_n\circ\iota_A$. Then $\Psi(ay)=\lambda(\phi(a))\Psi(y)$ for
$a\in A$, $\Psi$ has integral values, and $\Psi(cy)=A(c)\Psi(y)$ for
$c\in\mathcal O$, as $\phi(c)=1\otimes c$.

Second, $\theta_n(x)=\varpi\otimes1$, and $\theta_n(\sigma(x))=\varpi\otimes1$ for
$\sigma\in G_{K_m}$. Indeed, as in the first step of the proof of
Lemma~\ref{asm-two-three-killing-wach-frame-exists-decompletion-sandwich},
$x=(1+\iota_A(\pi))^{\ell^{s-1}}-1=1\otimes([\varepsilon^{(1)}]-1)$, so
$\theta_n(x)=(\zeta_{\ell^{1+n}}-1)\otimes1$; and $\sigma$ fixes
$\zeta_{\ell^{n+1}}\in K_m$. Hence $\Phi(x^kv)=\varpi^k\Phi(v)$ and
$\Phi(\sigma(x)^kv)=\varpi^k\Phi(v)$, by
Lemma~\ref{asm-two-three-killing-wach-frame-exists-decompletion-specialization}(2)
and (3).

Third, for $\sigma\in G_{K_m}$ and $y\in N$, $\bar\sigma y\in N$ by clause (d),
as $G_{K_m}\subset G_{K_s}$, and $\Psi(\bar\sigma y)=M(\sigma)\sigma(\Psi(y))$.
Indeed, apply $\Phi$ to
Lemma~\ref{asm-two-three-killing-wach-frame-exists-decompletion-lift}(2): by
the second step and
Lemma~\ref{asm-two-three-killing-wach-frame-exists-decompletion-specialization}(3),
$\varpi^{M_0}M(\sigma)\sigma(\Psi(y))=\varpi^{M_0}\Psi(\bar\sigma y)$, and $\varpi\neq0$. For
$\sigma\in G_{K_\infty}$, $\bar\sigma=1$
(Lemma~\ref{asm-two-three-killing-wach-frame-exists-decompletion-tower}(2)),
so $\Psi(y)=M(\sigma)\sigma(\Psi(y))$.

Fourth, let $\Delta$ be the matrix whose column $(k,i)$ is
$\Psi(o_kb'_i)$, with entries in $\mathcal O_{\mathbf C}$. For
$y=\sum_ia_ib'_i\in N$ let $\operatorname{co}(y)$ be the vector with entry
$\kappa_k(\phi(a_i))$ at $(k,i)$. Since
$\lambda(w)=\sum_k\kappa_k(w)A(o_k)$ and $A(o_k)\Psi(b'_i)=\Psi(o_kb'_i)$,
\[
 \Psi(y)=\sum_i\lambda(\phi(a_i))\Psi(b'_i)=\Delta\operatorname{co}(y),
\]
and $\operatorname{co}(y)$ has entries in the valuation ring of $K_m$ by
Lemma~\ref{asm-two-three-killing-wach-frame-exists-decompletion-tower}(3).
Likewise $\lambda(w)\Psi(y)=\Delta\sum_k\kappa_k(w)\operatorname{co}(o_ky)$ is
$\Delta$ times an integral vector, for $w\in\mathcal O_{\mathbf C}\otimes
\mathcal O$.

Fifth, the hypotheses of
Lemma~\ref{asm-two-three-killing-wach-frame-exists-decompletion-matrix}, with
$g=g_0^{\ell^{s+n+1}}$ from
Lemma~\ref{asm-two-three-killing-wach-frame-exists-decompletion-tower}(1). The
matrices $M(\sigma)$ have entries in $\mathbf Z_\ell$, continuous in $\sigma$,
and are multiplicative. For each $j$, the unit vector $e_j$ is $\Phi(1\otimes b_j)$,
so by the second step and
Lemma~\ref{asm-two-three-killing-wach-frame-exists-decompletion-lift}(3),
$\varpi^{M_0+h_0}e_j=\Phi(x^{M_0+h_0}(1\otimes b_j))$ is a sum of terms
$\pm\Phi(r\,z(y))=\pm\lambda(\theta_n(r))\Psi(y)$, hence $\Delta w_j$ for an
integral $w_j$ by the fourth step. By
Lemma~\ref{asm-two-three-killing-wach-frame-exists-decompletion-tower}(6),
$\ell=u\varpi^{M_0+h_0}$ with $u$ integral, so $\ell e_j=\Delta(uw_j)$. For
$\sigma\in G_{K_m}$, column $(k,i)$ of $M(\sigma)\sigma(\Delta)$ is
$\Psi(\bar\sigma(o_kb'_i))=\Delta\operatorname{co}(\bar\sigma(o_kb'_i))$ by the
third and fourth steps; for $\sigma\in G_{K_\infty}$ it is $\Psi(o_kb'_i)$. For
$g$, with $y=o_kb'_i$ and $y'=\bar g_0^{\ell^{s+n+1}}y=\bar gy$, column $(k,i)$
of $M(g)g(\Delta)-\Delta$ is $\Psi(y')-\Psi(y)=\Delta\operatorname{co}(y'-y)$.
By Lemma~\ref{asm-two-three-killing-wach-frame-exists-decompletion-tower}(5),
$y'-y\in JN$, so its coordinates in the basis $(b'_i)$ lie in $J$, and the
entries of $\operatorname{co}(y'-y)$ have absolute value less than
$\ell^{-1/(\ell-1)}$ by part (4) of that lemma. So
Lemma~\ref{asm-two-three-killing-wach-frame-exists-decompletion-matrix} gives
$U_0$, with all the properties claimed for $\Delta$, $U_0$ and $g$.

Finally, for $c'\in\mathcal O_L$, column $(k,i)$ of $A(c')\Delta$ is
$\Psi(c'o_kb'_i)=\Delta\operatorname{co}(c'o_kb'_i)$, so
$\Delta^{-1}A(c')\Delta$ has entries in the valuation ring of $K_m$.
% Lean: GlobalTheory.Cyc.Decomp.exists_decompletion, with IsDec actC chi K m M Dm U g
% (Dm invertible, |Dm| <= 1, |l Dm^{-1}| <= 1, Frame chi K m U g, Dm U(sigma) =
% M(sigma) sigma(Dm) on Gj K m) and the entries of Dm^{-1} A a Dm in intC (Ks K m). Its
% proof: T.wach_lattice_sandwich_Ainf (the image of the Wach module and the lattice over
% A_inf, with E.io S.W.X raised to M0 and h0), the A-basis of N (NBasis),
% exists_fontaineTheta_frobeniusInverse_shift, exists_ainf_tensor_lattice_specialisation,
% exists_lift and the generator of level s + n0 + 1 with m = s + (s + n0 + 1). Size: 268
% theorems, about 3,900 lines in 63 files. The fourth and fifth steps are Decomp.assemble (core,
% hball, cols) feeding decompletion_frame_of_matrix_form; hmain is the third step, hsc the
% second. Lemmas -decompletion-frobenius and -decompletion-tower are written out to the floor.
\end{proof}

\begin{lemma}\label{asm-two-three-killing-wach-frame-exists-component-idempotent}
Let $F\subset K$ be fields with $K$ finite and separable over $F$.
\begin{enumerate}
\item There are $n\geq0$ and $\alpha_1,\ldots,\alpha_n,\beta_1,\ldots,\beta_n\in K$
  with $\sum_j\alpha_j\beta_j=1$ and
  $\sum_j\kappa\alpha_j\otimes\beta_j=\sum_j\alpha_j\otimes\kappa\beta_j$ in
  $K\otimes_FK$ for every $\kappa\in K$.
\item Let $C$ be a commutative $F$-algebra and $R$ a $C$-algebra, let
  $j\colon K\to C$ and $B\colon K\to R$ be maps of $F$-algebras, and let
  $\alpha_j,\beta_j$ be as in (1). Then $E=\sum_jj(\beta_j)B(\alpha_j)$
  satisfies $B(\kappa)E=j(\kappa)E$ for every $\kappa\in K$, and $E^2=E$.
\end{enumerate}
\end{lemma}
\leanref{Development/AbelianVarietyTheory/Proofs/InertiaEigenvalues/ExistsTauComponentFrameOfLattice.lean}{GlobalTheory.TauFrame.exists_sep_elem}
\leanref{Development/AbelianVarietyTheory/Proofs/InertiaEigenvalues/ExistsTauComponentFrameOfLattice.lean}{GlobalTheory.TauFrame.sep_idem}
\begin{proof}
(1) Write $K=F[X]/(f)$ with $f$ separable, and let $x\in K$ be the class of
$X$. Then $K\otimes_FK=K[Y]/(f(Y))$, and $f(Y)=(Y-x)h(Y)$ with $h(x)\neq0$, so
$K\otimes_FK\cong K\times K[Y]/(h)$, where the first projection is the
multiplication $\mu$. Let $\varepsilon=\sum_j\alpha_j\otimes\beta_j$ be the
idempotent of the first factor. Then $\mu(\varepsilon)=\sum_j\alpha_j\beta_j=1$,
and $\kappa\otimes1-1\otimes\kappa$ lies in the kernel of $\mu$, the second
factor, so $(\kappa\otimes1)\varepsilon=(1\otimes\kappa)\varepsilon$.

(2) The $F$-bilinear map $(a,c)\mapsto j(c)B(a)$ induces an $F$-linear map
$\phi\colon K\otimes_FK\to R$. Applying $\phi$ to the identity in (1) gives
$\sum_jj(\beta_j)B(\kappa)B(\alpha_j)=\sum_jj(\kappa)j(\beta_j)B(\alpha_j)$,
that is, $B(\kappa)E=j(\kappa)E$. Hence
$E^2=\sum_jj(\beta_j)B(\alpha_j)E=\sum_jj(\beta_j)j(\alpha_j)E
=j\big(\sum_j\alpha_j\beta_j\big)E=E$.
\end{proof}

\begin{lemma}\label{asm-two-three-killing-wach-frame-exists-component-factor}
Let $C$ be a field with an ultrametric absolute value, let $F\subset C$ be a
subfield, and let $e$ be a $d\times d$ matrix with entries in $F$ and $e^2=e$.
Then there are an integer $r\geq0$, an $r\times d$ matrix $R$ and a $d\times r$
matrix $S$, with entries in $F$ of absolute value at most $1$, such that
$RS=1$, $Re=R$ and $eSR=e$.
\end{lemma}
\leanref{Development/AbelianVarietyTheory/Proofs/InertiaEigenvalues/ExistsIntegralFactorOfIdempotentSubfield.lean}{GlobalTheory.exists_integral_factor_of_idempotent_subfield}
\leanref{Development/AbelianVarietyTheory/Proofs/InertiaEigenvalues/ExistsIntegralFactorOfIdempotentSubfield.lean}{GlobalTheory.exists_integral_factor_of_idempotent}
\begin{proof}
For maps $\tau,\sigma\colon\{1,\ldots,k\}\to\{1,\ldots,d\}$ let
$e_{\tau,\sigma}$ be the $k\times k$ matrix $(e_{\tau(a)\sigma(b)})$. Let $r$ be
the largest $k\leq d$ for which some $\det e_{\tau,\sigma}$ is nonzero; $r=0$
qualifies, with the empty determinant $1$. Every $\det e_{\tau,\sigma}$ of size
$r+1$ vanishes: by the choice of $r$ if $r+1\leq d$, and because $\tau$ repeats
a row otherwise. Among the pairs $(\tau,\sigma)$ of size $r$ choose one with
$|\det e_{\tau,\sigma}|$ largest, and put $A=e_{\tau,\sigma}$, which is
invertible. Let $S$ be the $d\times r$ matrix with $S_{ik}=1$ if $i=\sigma(k)$
and $0$ otherwise, so that $XS$ is the matrix of the columns $\sigma(1),\ldots,
\sigma(r)$ of $X$, and let $R=A^{-1}e_{\tau,*}$, where $e_{\tau,*}$ is the
matrix of the rows $\tau(1),\ldots,\tau(r)$ of $e$.

By Cramer's rule, $R_{kj}=\det e_{\tau,\sigma'}/\det A$, where $\sigma'$ is
$\sigma$ with $\sigma(k)$ replaced by $j$, and this has absolute value at most
$1$ by the choice of $(\tau,\sigma)$. The entries of $S$ are $0$ and $1$.
Next, $RS=A^{-1}e_{\tau,\sigma}=1$, and $Re=A^{-1}(e^2)_{\tau,*}=R$. Finally,
$(eSR)_{ij}=e_{i,\sigma}A^{-1}e_{\tau,j}$, with the row $e_{i,\sigma}$ and the
column $e_{\tau,j}$. The $(r+1)\times(r+1)$ matrix with rows $\tau(1),\ldots,
\tau(r),i$ and columns $\sigma(1),\ldots,\sigma(r),j$ of $e$ has determinant
$\det A\cdot(e_{ij}-e_{i,\sigma}A^{-1}e_{\tau,j})$, by the Schur complement, and
this vanishes. So $eSR=e$.
\end{proof}

\begin{lemma}\label{asm-two-three-killing-wach-frame-exists-component-symbols}
Let $\ell$ be a prime and use the notation of
Definition~\ref{asm-two-three-killing-wach-frame}. Let
$K,L\subset\bar{\mathbf Q}_\ell$ be finite extensions of $\mathbf Q_\ell$, let
$V$ be a finite-dimensional $L$-vector space with an $L$-linear action of
$G_K$ such that $V$ is de Rham, let $b_1,\ldots,b_d$ be a $\mathbf Q_\ell$-basis
of $V$, and in this basis let $M(\sigma)$ be the matrix of $\sigma\in G_K$ and
$B(c)$ that of multiplication by $c\in L$. Let $\tau\colon K\to L$ be a
$\mathbf Q_\ell$-embedding, let $\alpha_j,\beta_j\in K$ be as in
Lemma~\ref{asm-two-three-killing-wach-frame-exists-component-idempotent}(1)
for $F=\mathbf Q_\ell$, and put $E=\sum_j\beta_jB(\tau(\alpha_j))$. Then there
are integers $u_1,\ldots,u_d$ and a basis $e_1,\ldots,e_d$ of
$\mathbf C_\ell^d$ such that $\chi(\sigma)^{u_k}M(\sigma)\sigma(e_k)=e_k$ for
every $\sigma\in G_K$ and every $k$, and $Ee_k=0$ whenever $u_k$ is not a
Hodge-Tate weight of $V$ along $\tau$
(Definition~\ref{asm-two-three-local-killing}(3)).
\end{lemma}
\leanref{Development/AbelianVarietyTheory/Proofs/InertiaEigenvalues/ExistsSymbolBasis.lean}{GlobalTheory.DeRhamSymbolBasis.exists_symbol_basis}
\leanref{Development/AbelianVarietyTheory/Proofs/InertiaEigenvalues/TmulMemFilV.lean}{GlobalTheory.EmbeddingJump.proj_eq_zero}
\leanref{Development/AbelianVarietyTheory/Proofs/InertiaEigenvalues/SerreTateLinearIndependentOfWeightwise.lean}{GlobalTheory.SerreTateIndep.linearIndependent_of_weightwise}
\begin{proof}
Let $B_{\mathrm{dR}}$ be Fontaine's field of de Rham periods, with its
decreasing, exhaustive and separated filtration $(\operatorname{Fil}^a)_{a\in
\mathbf Z}$, the unit $t$ with $\operatorname{Fil}^a=t^a\operatorname{Fil}^0$
and $\sigma(t)=\chi(\sigma)t$, and Fontaine's map
$\theta\colon\operatorname{Fil}^0\to\mathbf C_\ell$, a
$G_{\mathbf Q_\ell}$-equivariant ring map that is the identity on
$\bar{\mathbf Q}_\ell$ and has kernel $\operatorname{Fil}^1$
\cite[Proposition~4.4.6, Definition~4.4.7, the computation of $g(t)$ before
Proposition~4.4.8, Proposition~4.4.8 and Lemma~4.4.10]{BC09}. Apply $\theta$ and
the filtration to vectors entry by entry,
and for $x\in B_{\mathrm{dR}}^d$ and $\sigma\in G_K$ put
$\sigma\cdot x=M(\sigma)\sigma(x)$.

First, $x\mapsto\sum_ib_i\otimes x_i$ is an isomorphism
$B_{\mathrm{dR}}^d\to V\otimes_{\mathbf Q_\ell}B_{\mathrm{dR}}$ that carries
$\sigma\cdot x$ to the action of $\sigma\otimes\sigma$,
$(\operatorname{Fil}^a)^d$ onto $V\otimes\operatorname{Fil}^a$, and
multiplication by $c\in K\subset B_{\mathrm{dR}}$ to multiplication on the
second factor. So it identifies
$D=\{x\in B_{\mathrm{dR}}^d:\sigma\cdot x=x\text{ for }\sigma\in G_K\}$ with
$D_{\mathrm{dR}}(V)$. As $V$ is de Rham, $D_{\mathrm{dR}}(V)$ has dimension
$[K:\mathbf Q_\ell]\dim_LV$ over $L$, hence $[K:\mathbf Q_\ell]d$ over
$\mathbf Q_\ell$, and $D$ has dimension $d$ over $K$.

Second, $D$ has a basis $x_1,\ldots,x_d$ adapted to the filtration: there are
integers $u_k$ with $x_k\in(\operatorname{Fil}^{u_k})^d$, such that for every
$a\in\mathbf Z$ and every $c\in K^d$ with $c_k=0$ whenever $u_k\neq a$,
$\sum_kc_kx_k\in(\operatorname{Fil}^{a+1})^d$ only if $c=0$. This follows by
induction on $\dim D$: let $a_0$ be the largest integer with
$D\subset(\operatorname{Fil}^{a_0})^d$, which exists as the filtration is
exhaustive and separated; choose $x_0\in D$ outside
$(\operatorname{Fil}^{a_0+1})^d$ and a $K$-linear form $f$ on
$B_{\mathrm{dR}}^d$ that vanishes on $(\operatorname{Fil}^{a_0+1})^d$ with
$f(x_0)\neq0$; and add $x_0$, with $u=a_0$, to an adapted basis of
$D\cap\ker f$.

Third, put $y_k=t^{-u_k}x_k\in(\operatorname{Fil}^0)^d$ and
$e_k=\theta(y_k)\in\mathbf C_\ell^d$. For $\sigma\in G_K$,
$\sigma\cdot y_k=\chi(\sigma)^{-u_k}t^{-u_k}(\sigma\cdot x_k)
=\chi(\sigma)^{-u_k}y_k$, and applying $\theta$, as $M(\sigma)$ has entries in
$\mathbf Q_\ell$, gives $\chi(\sigma)^{u_k}M(\sigma)\sigma(e_k)=e_k$.

Fourth, let $a\in\mathbf Z$ and $c\in K^d$ with $c_k=0$ whenever $u_k\neq a$,
and suppose that $\sum_kc_ke_k=0$. Then $\theta(\sum_kc_ky_k)=0$, so
$\sum_kc_ky_k\in(\operatorname{Fil}^1)^d$ and
$\sum_kc_kx_k=t^a\sum_kc_ky_k\in(\operatorname{Fil}^{a+1})^d$. So $c=0$.

Fifth, the $e_k$ are linearly independent over $\mathbf C_\ell$, hence a basis.
Otherwise choose a nonzero $c\in\mathbf C_\ell^d$ with $\sum_kc_ke_k=0$ and as
few nonzero entries as possible, and $k_0$ with $c_{k_0}=1$, after scaling. For
$\sigma\in G_K$, applying $x\mapsto\chi(\sigma)^{u_{k_0}}M(\sigma)\sigma(x)$
gives $\sum_k\chi(\sigma)^{u_{k_0}-u_k}\sigma(c_k)e_k=0$; subtracting
$\sum_kc_ke_k=0$ leaves a relation with fewer nonzero entries, so
$\chi(\sigma)^{u_{k_0}-u_k}\sigma(c_k)=c_k$ for every $k$ and every
$\sigma\in G_K$. By Tate's theorem \cite[Theorem~2.2.7]{BC09}, $c_k=0$ if
$u_k\neq u_{k_0}$, and $c_k\in K$ if $u_k=u_{k_0}$. The fourth step, with
$a=u_{k_0}$, then gives $c=0$, a contradiction.

Finally, suppose that $u=u_k$ is not a Hodge-Tate weight of $V$ along $\tau$.
The $\tau$-part of $D_{\mathrm{dR}}(V)$ consists of the $z$ with $cz=\tau(c)z$
for $c\in K$, where $c$ acts on $B_{\mathrm{dR}}$ and $\tau(c)$ on $V$; the
hypothesis says that its intersections with $V\otimes\operatorname{Fil}^u$ and
with $V\otimes\operatorname{Fil}^{u+1}$ coincide. In the coordinates of the
first step, multiplication by $c'\in L$ on $V$ is $x\mapsto B(c')x$, so the
element $\sum_j\tau(\alpha_j)\beta_jx_k$ is $z=\sum_j\beta_jB(\tau(\alpha_j))x_k=Ex_k$.
It lies in $D$, as multiplication by $L$ on $V$ and by $K$ on
$B_{\mathrm{dR}}$ commute with the action of $G_K$, and in
$(\operatorname{Fil}^u)^d$. Applying the $\mathbf Q_\ell$-linear map
$a\otimes c'\mapsto c'B(\tau(a))x_k$ to the identity of
Lemma~\ref{asm-two-three-killing-wach-frame-exists-component-idempotent}(1)
gives $cz=B(\tau(c))z$ for $c\in K$, so $z$ lies in the $\tau$-part. Hence
$z\in(\operatorname{Fil}^{u+1})^d$, so $Ey_k=t^{-u}z\in(\operatorname{Fil}^1)^d$
and $Ee_k=\theta(Ey_k)=0$, as the entries of $E$ lie in $K$.
% Lean: exists_symbol_basis (exists_adapted with exists_coord_tensor_dR,
% deRham_invariants_coord_finrank and exists_family_adapted_filtration; shift_act, theta_comb
% and θ_act for the third step; linearIndependent_of_weightwise with tate_invariants_twist,
% the G_K form of Tate-Sen, for the fourth and fifth) and EmbeddingJump.proj_eq_zero
% (proj_mem, coord_proj, theta_lin) for the last. The period ring, its filtration, t and θ are
% the datum's (P.R.BdR, P.PD K, tU P and the theta datum ThetaL), read in Fontaine's sense.
\end{proof}

\begin{lemma}\label{asm-two-three-killing-wach-frame-exists-component-cut}
Let $\ell$ be a prime and use the notation of
Definition~\ref{asm-two-three-killing-wach-frame}. Let
$K\subset\bar{\mathbf Q}_\ell$ be a finite extension of $\mathbf Q_\ell$, let
$m\geq1$ be an integer, with $m\geq2$ if $\ell=2$, and let $(U_0,g)$ be a frame
of rank $d$ at level $m$ for $K$. Let $\Delta$ be an invertible $d\times d$
matrix over $\mathbf C_\ell$ such that the entries of $\Delta$ and of
$\ell\Delta^{-1}$ lie in $\mathcal O_{\mathbf C}$, and let
$M\colon G_{K_m}\to M_d(\mathbf C_\ell)$ satisfy
$\Delta U_0(\sigma)=M(\sigma)\sigma(\Delta)$ for $\sigma\in G_{K_m}$. Let $E$
be a $d\times d$ matrix with $EM(\sigma)=M(\sigma)E$ and $\sigma(E)=E$ for
$\sigma\in G_{K_m}$, such that $e'=\Delta^{-1}E\Delta$ has entries in $K_m$.
Let $R$ be an $r\times d$ and $S$ a $d\times r$ matrix with entries in $K_m$ of
absolute value at most $1$, with $RS=1$, $Re'=R$ and $e'SR=e'$. Let
$e_1,\ldots,e_d$ be a basis of $\mathbf C_\ell^d$ and $u_1,\ldots,u_d$ integers
with $\chi(\sigma)^{u_k}M(\sigma)\sigma(e_k)=e_k$ for $\sigma\in G_{K_m}$, and
let $J\subset\mathbf Z$ contain $u_k$ whenever $Ee_k\neq0$. Put
$U(\sigma)=RU_0(\sigma)S$ and $P=R\Delta^{-1}$. Then:
\begin{enumerate}
\item $(U,g)$ is a frame of rank $r$ at level $m$ for $K$, with a Hodge-Tate
  basis with weights in $J$;
\item $PM(\sigma)=U(\sigma)\,\sigma(P)$ for $\sigma\in G_{K_m}$;
\item $\mathcal O_{\mathbf C}^r\subset P(\mathcal O_{\mathbf C}^d)$ and
  $\ell P(\mathcal O_{\mathbf C}^d)\subset\mathcal O_{\mathbf C}^r$;
\item $QP=E$ for the $d\times r$ matrix $Q=\Delta e'S$.
\end{enumerate}
\end{lemma}
\leanref{Development/AbelianVarietyTheory/Proofs/InertiaEigenvalues/ExistsSubBasis.lean}{GlobalTheory.TauFrame.frame_of_idem}
\begin{proof}
Elements of $G_{K_m}$ fix matrices with entries in $K_m$, such as $e'$ and $R$.

First, $e'U_0(\sigma)=U_0(\sigma)e'$ for $\sigma\in G_{K_m}$. Indeed
$\Delta e'=E\Delta$, so
\[
 \Delta e'U_0(\sigma)=EM(\sigma)\sigma(\Delta)=M(\sigma)\sigma(E)\sigma(\Delta)
 =M(\sigma)\sigma(\Delta e')=M(\sigma)\sigma(\Delta)e'=\Delta U_0(\sigma)e',
\]
and $\Delta$ is invertible.

Second, $RU_0(\sigma)SR=RU_0(\sigma)$ for $\sigma\in G_{K_m}$. By the first
step and $Re'=R$, $RU_0(\sigma)e'=Re'U_0(\sigma)=RU_0(\sigma)$, so
$RU_0(\sigma)SR=RU_0(\sigma)e'SR=RU_0(\sigma)e'=RU_0(\sigma)$.

Third, (2). From $\Delta U_0(\sigma)=M(\sigma)\sigma(\Delta)$,
$\Delta^{-1}M(\sigma)=U_0(\sigma)\sigma(\Delta^{-1})$, so
$PM(\sigma)=RU_0(\sigma)\sigma(\Delta^{-1})=RU_0(\sigma)SR\,\sigma(\Delta^{-1})
=U(\sigma)\sigma(R\Delta^{-1})$ by the second step.

Fourth, $(U,g)$ is a frame. The entries of $U$ are linear combinations of those
of $U_0$, so continuous; they lie in $K_m$ and have absolute value at most $1$,
as those of $R$, $U_0(\sigma)$ and $S$ do and the absolute value is
ultrametric. By the second step,
$U(\sigma)U(\sigma')=RU_0(\sigma)SRU_0(\sigma')S=RU_0(\sigma\sigma')S
=U(\sigma\sigma')$. On $G_{K_\infty}$, $U=RS=1$. The element $g$ is the one of
the frame $(U_0,g)$, and $U(g)-1=R(U_0(g)-1)S$ has entries of absolute value at
most the largest absolute value of an entry of $U_0(g)-1$, which is less than
$\ell^{-1/(\ell-1)}$.

Fifth, the Hodge-Tate basis. Since $P(\Delta S)=RS=1$, $P$ is surjective, so
$Pe_1,\ldots,Pe_d$ span $\mathbf C_\ell^r$; choose $k_1,\ldots,k_r$ with
$Pe_{k_1},\ldots,Pe_{k_r}$ a basis. Since $PE=Re'\Delta^{-1}=P$, $Ee_{k_i}=0$
would give $Pe_{k_i}=0$; so $Ee_{k_i}\neq0$ and $u_{k_i}\in J$. For
$\sigma\in G_{K_m}$, by (2),
$\chi(\sigma)^{u_{k_i}}U(\sigma)\sigma(Pe_{k_i})
=P\big(\chi(\sigma)^{u_{k_i}}M(\sigma)\sigma(e_{k_i})\big)=Pe_{k_i}$.

Sixth, (3): for $y\in\mathcal O_{\mathbf C}^r$, $y=P(\Delta Sy)$ with
$\Delta Sy\in\mathcal O_{\mathbf C}^d$; and for $x\in\mathcal O_{\mathbf C}^d$,
$\ell Px=R(\ell\Delta^{-1}x)\in\mathcal O_{\mathbf C}^r$. Finally,
$QP=\Delta e'SR\Delta^{-1}=\Delta e'\Delta^{-1}=E$.
\end{proof}

\begin{lemma}\label{asm-two-three-killing-wach-frame-exists-component-coordinates}
Let $\ell$ be a prime and use the notation of
Definition~\ref{asm-two-three-killing-wach-frame}. Let
$L\subset\bar{\mathbf Q}_\ell$ be a finite extension of $\mathbf Q_\ell$, let
$V$ be a finite-dimensional $L$-vector space and $T\subset V$ an
$\mathcal O_L$-lattice, let $b_1,\ldots,b_d$ be a $\mathbf Z_\ell$-basis of $T$,
let $B(c)$ be the matrix of multiplication by $c\in L$ in this basis, and let
$\iota\in G_{\mathbf Q_\ell}$. Then there are $\rho\geq0$ and a $\rho\times d$
matrix $P'$ with entries in $\mathcal O_{\mathbf C}$ such that
$P'B(c)=\iota(c)P'$ for every $c\in L$, and every $y\in\mathbf Q_\ell^d$ with
$P'y\in\mathcal O_{\mathbf C}^\rho$ lies in $\mathbf Z_\ell^d$.
\end{lemma}
\leanref{Development/AbelianVarietyTheory/Proofs/InertiaEigenvalues/NormCoeAutLeOneIff.lean}{GlobalTheory.EigenCoord.exists_matrix}
\leanref{Development/AbelianVarietyTheory/Proofs/InertiaEigenvalues/NormCoeAutLeOneIff.lean}{GlobalTheory.EigenCoord.exists_matrix_of_basis}
\begin{proof}
The lattice $T$ is free over $\mathcal O_L$; let $t_1,\ldots,t_\rho$ be a
basis. It is an $L$-basis of $V$, and $v\in V$ lies in $T$ exactly when its
coordinates $t_i^*(v)$ in this basis lie in $\mathcal O_L$. For $x\in L$,
$|\iota(x)|\leq1$ exactly when $x\in\mathcal O_L$, since $\iota$ preserves
integrality over $\mathbf Z_\ell$. Put $P'_{ij}=\iota(t_i^*(b_j))$. Its entries
are integral, as $b_j\in T$. For $c\in L$,
$(P'B(c))_{ik}=\iota\big(t_i^*(\sum_jB(c)_{jk}b_j)\big)=\iota(t_i^*(cb_k))
=\iota(c)P'_{ik}$, as $t_i^*$ is $L$-linear and $B(c)$ has entries in
$\mathbf Q_\ell$. If $y\in\mathbf Q_\ell^d$ and $P'y$ is integral, then
$\iota(t_i^*(\sum_jy_jb_j))$ is integral for every $i$, so
$\sum_jy_jb_j\in T$ and $y\in\mathbf Z_\ell^d$.
\end{proof}

\begin{lemma}\label{asm-two-three-killing-wach-frame-exists-component}
Let $\ell$ be a prime and use the notation of
Definition~\ref{asm-two-three-killing-wach-frame}. Let
$K,L\subset\bar{\mathbf Q}_\ell$ be finite extensions of $\mathbf Q_\ell$, let
$m\geq1$ be an integer, with $m\geq2$ if $\ell=2$, such that $K\subset K_m$,
let $(V,T)$ be a lattice representation of $G_K$ over $L$ with $V$ de Rham, and
let $\tau\colon K\to L$ be a $\mathbf Q_\ell$-embedding. Let $b_1,\ldots,b_d$
be a $\mathbf Z_\ell$-basis of $T$, with the matrices $M(\sigma)$ and $A(c)$ of
Lemma~\ref{asm-two-three-killing-wach-frame-exists-decompletion}, and let
$\Delta$, $U_0$ and $g$ be as in the conclusion of that lemma for this $m$. Then
the Hodge-Tate weights of $V$ along $\tau$ form a finite set, and there are an
integer $r\geq0$, a frame $(U,g)$ of rank $r$ at level $m$ for $K$, a
Hodge-Tate basis of $(U,g)$ whose weights are Hodge-Tate weights of $V$ along
$\tau$, and an $r\times d$ matrix $P$ over $\mathbf C_\ell$ with the properties
(1) to (3) of Lemma~\ref{asm-two-three-killing-wach-frame-exists}.
\end{lemma}
\leanref{Development/AbelianVarietyTheory/Proofs/InertiaEigenvalues/ExistsTauComponentFrameOfLattice.lean}{GlobalTheory.TauFrame.frame_of_dec}
\begin{proof}
For $c\in L$ let $B(c)$ be the matrix of multiplication by $c$ in the basis
$(b_j)$, so that $B(c)=A(c)$ for $c\in\mathcal O_L$; the $b_j$ form a
$\mathbf Q_\ell$-basis of $V$. Since $K\subset K_m$, $G_{K_m}\subset G_K$.

First, the Hodge-Tate weights along $\tau$ form a finite set. By Fontaine's
bound for the field $B_{\mathrm{dR}}$, $\dim_KD_{\mathrm{dR}}(V)\leq
\dim_{\mathbf Q_\ell}V$ \cite[Theorem~5.2.1(1)]{BC09}, so
$\dim_LD_{\mathrm{dR}}(V)\leq[K:\mathbf Q_\ell]\dim_LV$, and
$D_{\mathrm{dR}}(V)$ is finite-dimensional over $L$, and so is its $\tau$-part
$D_\tau$. The subspaces $D_\tau\cap(V\otimes\operatorname{Fil}^a)$ decrease as
$a$ grows, as the filtration of $B_{\mathrm{dR}}$ does, so their dimension drops
at most $\dim_LD_\tau$ times, and the weights are the $a$ where it drops.

Second, the idempotent. Let $\alpha_j,\beta_j\in K$ be as in
Lemma~\ref{asm-two-three-killing-wach-frame-exists-component-idempotent}(1)
for $F=\mathbf Q_\ell$, and put $E=\sum_j\beta_jB(\tau(\alpha_j))$. Part (2) of
that lemma, with $C=\mathbf C_\ell$, $j$ the inclusion of $K$ and $B\circ\tau$,
gives $E^2=E$. For $\sigma\in G_{K_m}$, $EM(\sigma)=M(\sigma)E$, as
multiplication by $L$ commutes with $G_K$; and $\sigma(E)=E$, as the entries of
$E$ lie in $K\subset K_m$. For $c\in L$ choose a nonzero $c'\in\mathbf Z_\ell$
with $c'c\in\mathcal O_L$; then
$\Delta^{-1}B(c)\Delta=c'^{-1}\Delta^{-1}A(c'c)\Delta$ has entries in $K_m$. So
$e'=\Delta^{-1}E\Delta=\sum_j\beta_j\Delta^{-1}B(\tau(\alpha_j))\Delta$ has
entries in $K_m$, and $e'^2=\Delta^{-1}E^2\Delta=e'$.

Third, Lemma~\ref{asm-two-three-killing-wach-frame-exists-component-factor},
for the subfield $K_m$ of $\mathbf C_\ell$ and $e'$, gives $r$, $R$ and $S$.
Lemma~\ref{asm-two-three-killing-wach-frame-exists-component-symbols} gives
integers $u_k$ and a basis $(e_k)$ of $\mathbf C_\ell^d$ with
$\chi(\sigma)^{u_k}M(\sigma)\sigma(e_k)=e_k$ for $\sigma\in G_K$, hence for
$\sigma\in G_{K_m}$, such that $u_k$ is a Hodge-Tate weight along $\tau$
whenever $Ee_k\neq0$. Lemma~\ref{asm-two-three-killing-wach-frame-exists-component-cut},
with $J$ the set of Hodge-Tate weights along $\tau$ and the frame $(U_0,g)$ of
rank $d$, gives the frame $(U,g)$ of rank $r$, its Hodge-Tate basis, the
matrix $P=R\Delta^{-1}$ with (1) and (2), and $Q$ with $QP=E$.

Finally, (3). Extend $\tau\colon K\to L\subset\bar{\mathbf Q}_\ell$ to an
automorphism of $\bar{\mathbf Q}_\ell$, and let $\iota\in G_{\mathbf Q_\ell}$
be its inverse, so that $\iota(\tau(x))=x$ for $x\in K$.
Lemma~\ref{asm-two-three-killing-wach-frame-exists-component-coordinates}, for
this $\iota$, gives $P'$. Then
$P'E=\sum_j\beta_jP'B(\tau(\alpha_j))=\sum_j\beta_j\alpha_jP'=P'$. Let
$y\in\mathbf Q_\ell^d$ with $Py=Px$ for some $x\in\mathcal O_{\mathbf C}^d$.
Then $P'y=P'QPy=P'QPx=P'x$, which lies in $\mathcal O_{\mathbf C}^\rho$, as
$P'$ and $x$ are integral. So $y\in\mathbf Z_\ell^d$.
% Lean: GlobalTheory.TauFrame.frame_of_dec, for the datum's period data P, AI, a theta
% datum Theta (ThetaTie.nonempty_thetaL) and fontCore (P.PD K) (adm7 D, from
% Per.FontaineAll.fontaine_bounds, 641 theorems, about 8,300 lines, read in Fontaine's sense):
% exists_ratBasis_of_isMat, exists_sep_elem and sep_idem (the idempotent E of the
% tau-component), finite_jumps, exists_integral_factor_of_idempotent_subfield (R, S),
% DeRhamSymbolBasis.exists_symbol_basis and EmbeddingJump.proj_eq_zero (the Hodge-Tate
% basis), frame_of_idem, and EigenCoord.exists_matrix with exists_undo for (3); P is
% R Dm^{-1}. Size: 194 theorems, about 3,300 lines in 43 files. The first step is
% finite_jumps: the second part of fontCore (P.PD K) (Fontaine's bound, Module.Finite and
% finrank D_dR <= [K:Q_l] dim_L V; from adm7 D = Per.FontaineAll.fontaine_bounds, 641
% theorems, about 8,300 lines, for the datum's period rings), DdRE (the tau-part) inside
% D_dR, FilV antitone through exists_coord_tensor_dR and Fil_anti, and
% finite_jumps_inf_filtration; hEF uses conj_mem_subC for the passage from O_L to L.
\end{proof}

\begin{lemma}\label{asm-two-three-killing-wach-frame-exists-derham-descent}
Let $\ell$ be a prime and fix an algebraic closure $\bar{\mathbf Q}_\ell$ of
$\mathbf Q_\ell$. Let $B_{\mathrm{dR}}$ be Fontaine's field of de Rham periods,
which contains $\bar{\mathbf Q}_\ell$, with its action of
$G_{\mathbf Q_\ell}=\operatorname{Gal}(\bar{\mathbf Q}_\ell/\mathbf Q_\ell)$, and
for a subfield $F\subset\bar{\mathbf Q}_\ell$ put
$G_F=\operatorname{Gal}(\bar{\mathbf Q}_\ell/F)$. Let $L$ be a finite extension
of $\mathbf Q_\ell$, let $F\subset F'\subset\bar{\mathbf Q}_\ell$ be finite
extensions of $\mathbf Q_\ell$ with $\sigma(F')\subset F'$ for every
$\sigma\in G_F$, and let $V$ be a finite-dimensional $L$-vector space with an
$L$-linear action of $G_F$. For a subgroup $H\subset G_F$ let $D_H$ be the
$L$-subspace of $V\otimes_{\mathbf Q_\ell}B_{\mathrm{dR}}$ fixed by the elements
$\sigma\otimes\sigma$, $\sigma\in H$. For $F''$ equal to $F$ or $F'$, $V$ is
\emph{de Rham} as a representation of $G_{F''}$ if
$\dim_LD_{G_{F''}}=[F'':\mathbf Q_\ell]\dim_LV$; this is Fontaine's condition
$\dim_{F''}D_{\mathrm{dR}}(V)=\dim_{\mathbf Q_\ell}V$. Then $D_{G_{F'}}$ is isomorphic to
$D_{G_F}^{[F':F]}$ as an $L$-vector space. In particular, if $V$ is de Rham as a
representation of $G_{F'}$, then it is de Rham as a representation of $G_F$.
\end{lemma}
\leanref{Development/AbelianVarietyTheory/Proofs/InertiaEigenvalues/IsDeRhamOfRestrictGaloisDescentStatement.lean}{GlobalTheory.T.isDeRham_of_restrict_galois_descent}
\leanref{Development/AbelianVarietyTheory/Proofs/InertiaEigenvalues/IsDeRhamOfRestrictGaloisDescent.lean}{GlobalTheory.DeRhamDescent.isDR_of_restrict}
\leanref{Development/AbelianVarietyTheory/Proofs/InertiaEigenvalues/IsDeRhamOfRestrictGaloisDescent.lean}{GlobalTheory.DeRhamDescent.finrank_fixed}
\begin{proof}
Put $M=V\otimes_{\mathbf Q_\ell}B_{\mathrm{dR}}$, $D=D_{G_F}$, $D'=D_{G_{F'}}$
and $n=[F':F]$.

First, $F'$ is Galois over $F$: every $F$-embedding of $F'$ into
$\bar{\mathbf Q}_\ell$ extends to an element of $G_F$, so it maps $F'$ into
$F'$. Let $\Gamma=\operatorname{Gal}(F'/F)$, a group of order $n$, and
$\varepsilon\colon G_F\to\Gamma$ the restriction, which is surjective with
kernel $G_{F'}$. For $c\in F'$ let $m_c$ be the multiplication by $c$ on the
factor $B_{\mathrm{dR}}$ of $M$; it is $L$-linear, and
$\sigma(m_cx)=m_{\varepsilon(\sigma)(c)}\sigma(x)$ for $\sigma\in G_F$ and
$x\in M$. So $m_c$ maps $D'$ into $D'$, and $D\subset D'$. For $x\in D'$ the
element $\sigma(x)$ depends only on $\varepsilon(\sigma)$; write $\gamma(x)$
for it, $\gamma\in\Gamma$.

Second, let $(b_\gamma)_{\gamma\in\Gamma}$ be a basis of $F'$ over $F$ and
$(\mu_\gamma)_{\gamma\in\Gamma}$ the dual basis for the trace form, so that
$\operatorname{Tr}_{F'/F}(\mu_\delta b_\gamma)$ is $1$ if $\delta=\gamma$ and
$0$ otherwise. The matrices $P=(\sigma(b_\gamma))_{\sigma,\gamma}$ and
$P'=(\sigma(\mu_\delta))_{\delta,\sigma}$, indexed by $\Gamma$, satisfy
$P'P=1$, since $(P'P)_{\delta\gamma}=\sum_{\sigma\in\Gamma}\sigma(\mu_\delta
b_\gamma)=\operatorname{Tr}_{F'/F}(\mu_\delta b_\gamma)$. Hence $PP'=1$, and its
row at $\sigma=1$ says that $\sum_\gamma b_\gamma\,\tau(\mu_\gamma)$ is $1$ if
$\tau=1$ and $0$ for the other $\tau\in\Gamma$.

Third, let $\Phi\colon D^\Gamma\to M$ be the $L$-linear map
$(x_\gamma)\mapsto\sum_\gamma m_{b_\gamma}x_\gamma$. It is injective: if
$\sum_\gamma m_{b_\gamma}x_\gamma=0$, apply $\sigma\in G_F$ with
$\varepsilon(\sigma)=\tau$; as the $x_\gamma$ are fixed, this gives
$\sum_\gamma m_{\tau(b_\gamma)}x_\gamma=0$ for every $\tau\in\Gamma$. Applying
$m_{\tau(\mu_\delta)}$ and summing over $\tau$ gives
$\sum_\gamma m_{\operatorname{Tr}(\mu_\delta b_\gamma)}x_\gamma=x_\delta=0$.

Fourth, the image of $\Phi$ is $D'$. It lies in $D'$ by the first step.
Conversely, let $x\in D'$ and put $x_\gamma=\sum_{\tau\in\Gamma}\tau(m_{\mu_\gamma}x)$.
For $\sigma\in G_F$, $\sigma(x_\gamma)=\sum_\tau(\varepsilon(\sigma)\tau)(m_{\mu_\gamma}x)=x_\gamma$,
so $x_\gamma\in D$. And
\[
 \sum_\gamma m_{b_\gamma}x_\gamma=\sum_{\tau\in\Gamma}
 m_{\sum_\gamma b_\gamma\tau(\mu_\gamma)}\,\tau(x)=x
\]
by the second step. So $\Phi$ is an isomorphism $D^\Gamma\to D'$.

Finally, if $\dim_LD'=[F':\mathbf Q_\ell]\dim_LV$, then $D$ is finite-dimensional
and $n\dim_LD=[F':F]\,[F:\mathbf Q_\ell]\dim_LV$, so
$\dim_LD=[F:\mathbf Q_\ell]\dim_LV$.
\end{proof}

\begin{lemma}\label{asm-two-three-killing-wach-frame-exists-derham}
Let $\ell$, $s$, $K$, $L$ and $D$ be as in
Lemma~\ref{asm-two-three-killing-wach-frame-exists-inputs}, and let $(V,T)$ be
a lattice representation of $G_K$ over $L$ such that $V$ is crystalline as a
representation of $G_{K_s}$. Then $V$ is de Rham.
\end{lemma}
\leanref{Development/AbelianVarietyTheory/Proofs/InertiaEigenvalues/FiniteDimensionalKsOfGood.lean}{GlobalTheory.Cyc.CrysDeRham.isDR_of_isCrysOn}
\leanref{Development/AbelianVarietyTheory/Proofs/InertiaEigenvalues/IsCrystallineCyclotomicLevelOfFinrankInvariants.lean}{GlobalTheory.Cyc.isCrys_Ks_of_finrank}
\leanref{Development/AbelianVarietyTheory/Proofs/InertiaEigenvalues/IsCrystallineCyclotomicLevelOfFinrankInvariants.lean}{GlobalTheory.Cyc.finrank_maxUnramified_eq_cyclotomicLevel}
\begin{proof}
First, $K\subset K_s$, and $\sigma(K_s)\subset K_s$ for every $\sigma\in G_K$,
since $\sigma$ fixes $K_0\subset K$ and permutes the $\ell^s$-th roots of unity.

Second, the maximal unramified subextension $F_0$ of $K_s$ has degree
$[K_0:\mathbf Q_\ell]$ over $\mathbf Q_\ell$. Indeed $K_0\subset F_0\subset K_s$,
so $K_s$ is generated over $F_0$, as over $K_0$, by the $\ell^s$-th roots of
unity. Over an unramified extension of $\mathbf Q_\ell$ these generate an
extension of degree $(\ell-1)\ell^{s-1}$, since $\Phi_{\ell^s}(X+1)$, for the
$\ell^s$-th cyclotomic polynomial $\Phi_{\ell^s}$, is an Eisenstein polynomial
at $\ell$. So $[K_s:\mathbf Q_\ell]=[K_0:\mathbf Q_\ell](\ell-1)\ell^{s-1}
=[F_0:\mathbf Q_\ell](\ell-1)\ell^{s-1}$. Hence the crystallinity of $V$ as a
representation of $G_{K_s}$, which says that
$(V\otimes_{\mathbf Q_\ell}B_{\mathrm{cris}})^{G_{K_s}}$ has dimension
$[K_0:\mathbf Q_\ell]\dim_LV$ over $L$, is Fontaine's crystallinity of $V$ as a
representation of $G_{K_s}$ for the base field $K_s$.

Third, a crystalline representation is de Rham
\cite[Proposition~9.1.9]{BC09}. So $V$ is de Rham as a representation of
$G_{K_s}$.

Finally, Lemma~\ref{asm-two-three-killing-wach-frame-exists-derham-descent},
with $F=K$ and $F'=K_s$, shows that $V$ is de Rham.
% Lean: CrysDeRham.isDR_of_isCrysOn: E = K_s (finiteDimensional_Ks_of_good); isCrys_Ks_of_finrank
% (Dcris_eq_iInf, mem_GKs_iff_mem_Gj, finrank_maxUnramified_eq_cyclotomicLevel with finrank_Ks)
% for the second step; the third part of fontCore (P.PD E) from adm7 D (Fontaine's
% crystalline-implies-de-Rham, which Per.FontaineAll.fontaine_bounds proves for the datum's
% period rings: fontCore_of_tied, CrisDR.isDR_of_isCrys; 641 theorems, about 8,300 lines) for
% the third; isDR_of_restrict with apply_mem_Ks for the last. 50 theorems, about 860 lines
% without fontaine_bounds.
\end{proof}

\begin{lemma}\label{asm-two-three-killing-wach-frame-exists}
Let $\ell$ be a prime and use the notation of
Definition~\ref{asm-two-three-killing-wach-frame}. For an integer $H\geq0$ let
$n_0(\ell,H)$ be the least $n\geq0$ with $H\leq\ell^n(\ell-1)$. Let $s\geq1$ be
an integer, with $s\geq2$ if $\ell=2$, let $K\subset\bar{\mathbf Q}_\ell$ be a
finite extension of $\mathbf Q_\ell$ with $K\subset K_s$, and let
$L\subset\bar{\mathbf Q}_\ell$ be a finite extension of $\mathbf Q_\ell$ that
contains $\tau(K_s)$ for every $\mathbf Q_\ell$-embedding
$\tau\colon K_s\to\bar{\mathbf Q}_\ell$. Let $D$ be a definite-side datum
with local Hodge-Tate data, Fontaine data and comparison data at every prime
(Definition~\ref{asm-local-bound-two-three}(3)). Let $V$ be a finite-dimensional
$L$-vector space with an $L$-linear action of $G_K$ such that $g\mapsto gx$ is
continuous for every $x\in V$, and let $T\subset V$ be a $G_K$-stable
$\mathcal O_L$-lattice. Let $M_0,h_0\geq0$ be integers, suppose that $V$ is
crystalline and that every integer $a$ with $\operatorname{Fil}^{a+1}\neq
\operatorname{Fil}^a$ on $D_{\mathrm{dR}}(V)$ lies in $[-M_0,h_0]$, and put
$m=2s+n_0(\ell,M_0+h_0)+1$. Let $b_1,\ldots,b_d$ be a $\mathbf Z_\ell$-basis of
$T$, and for $\sigma\in G_K$ let $M(\sigma)\in GL_d(\mathbf Z_\ell)$ be the
matrix of $\sigma$ in this basis. Let $\tau\colon K\to L$ be a
$\mathbf Q_\ell$-embedding. Then the Hodge-Tate weights of $V$ along $\tau$
(Definition~\ref{asm-two-three-local-killing}(3)) form a finite set, and there
are an integer $r\geq0$, a frame $(U,g)$ of rank $r$ at level $m$ for $K$, a
Hodge-Tate basis of $(U,g)$ whose weights are Hodge-Tate weights of $V$ along
$\tau$, and an $r\times d$ matrix $P$ over $\mathbf C_\ell$ such that:
\begin{enumerate}
\item $PM(\sigma)=U(\sigma)\,\sigma(P)$ for every $\sigma\in G_{K_m}$;
\item $\mathcal O_{\mathbf C}^r\subset P(\mathcal O_{\mathbf C}^d)$ and
  $\ell P(\mathcal O_{\mathbf C}^d)\subset\mathcal O_{\mathbf C}^r$;
\item if $y\in\mathbf Q_\ell^d$ and $Py\in P(\mathcal O_{\mathbf C}^d)$, then
  $y\in\mathbf Z_\ell^d$.
\end{enumerate}
\end{lemma}
\leanref{Development/AbelianVarietyTheory/Proofs/InertiaEigenvalues/ExistsTauComponentFrameOfLatticeStatement.lean}{GlobalTheory.T.exists_tau_component_frame_of_lattice}
\leanref{Development/AbelianVarietyTheory/Proofs/InertiaEigenvalues/ExistsTauComponentFrameOfLattice.lean}{GlobalTheory.TauFrame.exists_tau_frame}
\begin{proof}
Use the Wach data of $D$ for $(K,s,L)$
(Definition~\ref{asm-two-three-killing-wach-frame-exists-data}), and let $A(c)$
be the matrix of multiplication by $c\in\mathcal O_L$ on $T$ in the basis
$(b_j)$.

First, $(V,T)$ is a lattice representation of $G_K$ over $L$, and
Lemma~\ref{asm-two-three-killing-wach-frame-exists-wach} shows that the Wach
hypothesis holds for $(V,T)$, $M_0$ and $h_0$. In particular, $V$ is
crystalline as a representation of $G_{K_s}$.

Second, Lemma~\ref{asm-two-three-killing-wach-frame-exists-decompletion} gives
$\Delta$, $U_0$ and $g$ for this $m$.

Third, since $V$ is crystalline as a representation of $G_{K_s}$,
Lemma~\ref{asm-two-three-killing-wach-frame-exists-derham} shows that $V$ is de
Rham.

Finally, $K\subset K_s\subset K_m$, as $s\leq m$, and $m\geq3$. So
Lemma~\ref{asm-two-three-killing-wach-frame-exists-component}, applied to
$\tau$ and to $\Delta$, $U_0$ and $g$, gives the conclusion.
% Lean: GlobalTheory.TauFrame.exists_tau_frame after NLink.wAll_holds: ThetaTie.nonempty_thetaL,
% Cyc.Decomp.exists_decompletion, hKm (h.2.2.1 and Ks_mono), and frame_of_dec with adm7
% (fontaine_bounds) and CrysDeRham.isDR_of_isCrysOn applied to hW.1 (Lemma
% asm-two-three-killing-wach-frame-exists-derham; 50 theorems, about 860 lines).
% Earlier summary: the hypotheses give LocalSituation.HypW S E T M0 h0 by
% GlobalTheory.T.crystallineLatticeWachModuleExists_proof (= NLink.wAll_holds:
% T.wachHypothesis_crystallineLattices = WachHyp.hypW_crystalline, from
% T.wachModuleLocalConclusions_usedForms = NL.usedLocalConclusion_and_prvE and
% crysOnCond_all): IsCrysOn over G_{K_s}, the Jset window, and the Wach module
% S.W.N T (a submodule of S.W.DD T over LetA) with WClauses (NBasis, Q^M0 phi(N) in N,
% Q^h0 N in the span of phi(N), Gamma-stability, gamma - 1 maps N into pi N, stability
% under endomorphisms of T). Then GlobalTheory.T.exists_tau_component_frame_of_lattice
% (= TauFrame.exists_tau_frame: ThetaTie.nonempty_thetaL, Cyc.Decomp.exists_decompletion,
% frame_of_dec with Per.FontaineAll.fontaine_bounds and CrysDeRham.isDR_of_isCrysOn)
% gives r, Pm, U, g and J with the properties of the lemma; IsMat K LL T M A is the
% matrix data of the basis (Cyc.isMat_exists_and_props). The hypotheses of the lemma are
% those of WachHypothesisLocal and exists_tau_frame together; neither has published
% results among its hypotheses. Size: proof closure of exists_tau_frame, wAll_holds and
% fontaine_bounds, 2,928 theorems, about 45,400 lines in 621 files (wAll_holds alone
% about 41,100 lines in 557 files; exists_tau_frame alone about 6,700 lines in 99 files).
\end{proof}

\begin{lemma}\label{asm-two-three-killing-wach}
Assume three results of the third collection of published results in
Theorem~\ref{lit-bundle-unitary}(6): the local results of Bloch and Kato, and
local Tate duality and the local Euler characteristic formula over finite
extensions of $\mathbf Q_p$ for every prime $p$. The proof does not use these
three assumptions. Let $\ell$ be a prime and fix an
algebraic closure $\bar{\mathbf Q}_\ell$ of $\mathbf Q_\ell$. For a finite
extension $K\subset\bar{\mathbf Q}_\ell$ of $\mathbf Q_\ell$ let $K_0$ be its
maximal unramified subextension and, for $j\geq1$, let $K_j$ be the field
generated over $K_0$ by the $\ell^j$-th roots of unity in
$\bar{\mathbf Q}_\ell$. For an integer $H\geq0$ let $n_0(\ell,H)$ be the least
$n\geq0$ with $H\leq\ell^n(\ell-1)$. Let $s\geq1$ be an integer, with $s\geq2$
if $\ell=2$, let $K\subset\bar{\mathbf Q}_\ell$ be a finite extension of
$\mathbf Q_\ell$ with $K\subset K_s$, and let $L\subset\bar{\mathbf Q}_\ell$ be
a finite extension of $\mathbf Q_\ell$ that contains $\tau(K_s)$ for every
$\mathbf Q_\ell$-embedding $\tau\colon K_s\to\bar{\mathbf Q}_\ell$. Let $D$ be a
definite-side datum with local Hodge-Tate data, Fontaine data and comparison
data at every prime (Definition~\ref{asm-local-bound-two-three}(3)). Let $V$ be
a finite-dimensional $L$-vector space with an $L$-linear action of $G_K$ such
that $g\mapsto gx$ is continuous for every $x\in V$, let $T\subset V$ be a
$G_K$-stable $\mathcal O_L$-lattice, and let $J$ be the set of integers $a$
with $\operatorname{Fil}^{a+1}\neq\operatorname{Fil}^a$ on
$D_{\mathrm{dR}}(V)$. Let $M_0,h_0\geq0$ be integers. Suppose that $V$ is
crystalline, that $J\subset[-M_0,h_0]$, and that $0$ is not a Hodge-Tate
weight of $V$ along any $\mathbf Q_\ell$-embedding $K\to L$
(Definition~\ref{asm-two-three-local-killing}(3)). Put
$m=2s+n_0(\ell,M_0+h_0)+1$ and
\[
 e=v_\ell[K_m:\mathbf Q_\ell]-v_\ell[K:\mathbf Q_\ell]+2
 +\max\Big(1,\sum_{u\in J,\,u\neq0}\big(m+v_\ell(u)\big)\Big).
\]
Then $\ell^e$ kills every $G_K$-invariant element of $V/T$.
\end{lemma}
\leanref{Development/AbelianVarietyTheory/Proofs/InertiaEigenvalues/TorsionExponentKillsInvariantsNoJumpZeroStatement.lean}{GlobalTheory.T.torsion_exponent_kills_invariants_of_good}
\leanref{Development/AbelianVarietyTheory/Proofs/InertiaEigenvalues/CrystallineLatticeWachModuleExistsProofStatement.lean}{GlobalTheory.T.crystallineLatticeWachModuleExists_proof}
\leanref{Formalism/SharpRestriction/LevelLetterCocycles/DataNLoc.lean}{Formalism.BasicNotions.Loc.Good}
\leanref{Formalism/SharpRestriction/LevelLetterCocycles/DataNLoc.lean}{Formalism.BasicNotions.Loc.Ks}
\leanref{Formalism/SharpRestriction/LevelLetterCocycles/LocalSituation.lean}{Formalism.MainNotions.LocalSituation.cLL}
\leanref{Formalism/SharpRestriction/LevelLetterCocycles/LocalSituation.lean}{Formalism.MainNotions.LocalSituation.Jset}
\begin{proof}
Use the notation of Definition~\ref{asm-two-three-killing-wach-frame}, and let
$x\in V/T$ be $G_K$-invariant.

First, fix a $\mathbf Q_\ell$-embedding $\tau\colon K\to L$. One exists: every
$\mathbf Q_\ell$-embedding of $K$ in $\bar{\mathbf Q}_\ell$ extends to $K_s$, so
it maps $K$ into $L$, and there are $[K:\mathbf Q_\ell]\geq1$ of them. Let
$J_\tau$ be the set of Hodge-Tate weights of $V$ along $\tau$. If
$\operatorname{Fil}^{a+1}=\operatorname{Fil}^a$ on $D_{\mathrm{dR}}(V)$, the
same holds on its $\tau$-part, so $J_\tau\subset J$; and $0\notin J_\tau$ by
hypothesis.

Second, $T$ is free of finite rank over $\mathcal O_L$, hence over
$\mathbf Z_\ell$; let $b_1,\ldots,b_d$ be a $\mathbf Z_\ell$-basis of $T$, which
is a $\mathbf Q_\ell$-basis of $V$, and let $M(\sigma)\in GL_d(\mathbf Z_\ell)$
be the matrix of $\sigma\in G_K$. Since $K\subset K_s\subset K_m$, we have
$G_{K_m}\subset G_K$. Lemma~\ref{asm-two-three-killing-wach-frame-exists}, for
$D$, shows that $J_\tau$ is finite and gives $r$, a frame $(U,g)$ of rank $r$ at level $m$
for $K$, a Hodge-Tate basis of it with weights in a set $J'\subset J_\tau$, and
a matrix $P$ with properties (1) to (3) of that lemma. Put
$c=2+\sum_{u\in J_\tau}(m+v_\ell(u))$, which is defined since $0\notin J_\tau$.

Third, let $y\in\mathbf Q_\ell^d$ with $M(\sigma)y-y\in\mathbf Z_\ell^d$ for every
$\sigma\in G_K$; we show $\ell^cy\in\mathbf Z_\ell^d$. If $J'$ is empty, then
$r=0$, since the weights $u_1,\ldots,u_r$ of the basis lie in $J'$. Then
$Py=0\in P(\mathcal O_{\mathbf C}^d)$, so $y\in\mathbf Z_\ell^d$ by (3). Suppose
that $J'$ is not empty, and let $\Lambda=P(\mathcal O_{\mathbf C}^d)$, an
additive subgroup of $\mathbf C_\ell^r$ with $\mathcal O_{\mathbf C}^r\subset
\Lambda$ and $\ell\Lambda\subset\mathcal O_{\mathbf C}^r$ by (2). For
$\sigma\in G_{K_m}$, $\sigma$ fixes $y$, so by (1)
\[
 \sigma*_0(Py)-Py=U(\sigma)\,\sigma(P)\,y-Py=P\big(M(\sigma)y-y\big)\in\Lambda.
\]
Since $0\notin J'$, Lemma~\ref{asm-two-three-killing-wach-twisted} with $i=0$
gives $\ell^{2+\nu}Py\in\Lambda$, where $\nu=\nu_m(J',0)$. As $J'$ is not empty
and $m\geq1$, the sum in $\nu_m(J',0)$ is at least $1$, so
$\nu=\sum_{u\in J'}(m+v_\ell(u))\leq c-2$. Now (3), applied to
$\ell^{2+\nu}y$, gives $\ell^{2+\nu}y\in\mathbf Z_\ell^d$, and so
$\ell^cy\in\mathbf Z_\ell^d$.

Fourth, let $v\in V$ lift $x$, and let $y\in\mathbf Q_\ell^d$ be its
coordinates in the basis $(b_j)$. Since $x$ is invariant, $\sigma v-v\in T$ for
every $\sigma\in G_K$, and its coordinates are $M(\sigma)y-y$, which therefore
lie in $\mathbf Z_\ell^d$. By the third step $\ell^cy\in\mathbf Z_\ell^d$, that
is, $\ell^cv\in T$, and $\ell^cx=0$.

Finally, $c\leq e$. Indeed, $J_\tau\subset J\setminus\{0\}$ and the terms
$m+v_\ell(u)$ are nonnegative, so
$\sum_{u\in J_\tau}(m+v_\ell(u))\leq\max\big(1,\sum_{u\in J,\,u\neq0}(m+v_\ell(u))\big)$;
and $v_\ell[K_m:\mathbf Q_\ell]\geq v_\ell[K:\mathbf Q_\ell]$, since $K\subset
K_m$. Hence $\ell^ex=\ell^{e-c}\ell^cx=0$.
% Lean: the lemma combines two hypotheses of
% UpperTheory.N.invariants_killed_over_cyclotomic_subfield_of_admissible, as locAll_of
% supplies them: hW = GlobalTheory.T.crystallineLatticeWachModuleExists_proof gives
% LocalSituation.HypW, and hH0 = torsion_exponent_kills_invariants_of_good =
% HubX.InvariantsNoJump.torsion_exponent_kills_invariants_no_jump_zero, which ignores its
% three published hypotheses. Its proof: JsetEmb.card_emb_good for the embedding tau (first
% step); invariants_vanish_at_prime_of_admissible D HW F (Per.FontaineAll.fontaine_bounds D)
% = Cyc.InvKilled.invariants_killed at tau, which bounds by l^(c0 S T tau M0 h0) with c0 the
% c above; in it, exists_continuous_galoisLatticeMatrix_integral (Cyc.isMat_exists_and_props,
% the basis and M), exists_tau_component_frame_of_lattice (Lemma
% asm-two-three-killing-wach-frame-exists), prime_pow_smul_mem_of_twistedInvariant_mod_lattice
% with i = 0 and M = imBall P (Lemma asm-two-three-killing-wach-twisted; J empty forces r = 0),
% killed_of_coord and killed_quot (fourth step); and two_add_sum_jumps_le_intervalExponent
% (with WachHyp.jump_mem / mem_Jset_of_isJumpV) for c <= e.
\end{proof}

\subsection{The lifting, invariant-killing and block properties}\label{subsection-asm-lifting-killing-block}

\begin{lemma}\label{asm-two-three-killing-count}
Let $\ell$ be a prime, let $N\geq1$ and $h,n\geq0$ be integers, put
$s=\max(1,v_\ell(N))$ and $m=2s+n_0(\ell,2h)+1$, with $n_0$ as in
Lemma~\ref{asm-two-three-killing-wach}, and let $c_\ell(N,h,n)$ be as in
Definition~\ref{asm-two-three-local-killing}(2). Let $U$ be a set of at most
$n$ integers in $[-h,h]$, let $J\subset U$, and let $0\leq v\leq m$ be an
integer. Then
\[
 v+2+\max\Big(1,\sum_{u\in J,\,u\neq0}\big(m+v_\ell(u)\big)\Big)
 \leq c_\ell(N,h,n).
\]
\end{lemma}
\leanref{Development/AutomorphicArithmetic/ResidueCharacters/LocNum.lean}{UpperTheory.N.LocNum.latticeExponent_one_le_invariantsKillingExponent}
\leanref{Development/AutomorphicArithmetic/ResidueCharacters/LocNum.lean}{UpperTheory.N.LocNum.cLL_le}
\leanref{Development/AutomorphicArithmetic/ResidueCharacters/LocNum.lean}{UpperTheory.N.LocNum.rI_le}
\leanref{Development/AutomorphicArithmetic/ResidueCharacters/LocNum.lean}{UpperTheory.N.LocNum.mN_le}
\begin{proof}
Put $\lambda=s+\lfloor\log_\ell2h\rfloor+\lfloor\log_\ell N\rfloor+1$, with
$\lfloor\log_\ell y\rfloor$ as in Lemma~\ref{asm-two-three-split-numerics}, so
that $c_\ell(N,h,n)=20(n+1)\lambda$ and $\lambda\geq1$. Since
$2h<\ell^{\lfloor\log_\ell2h\rfloor+1}\leq\ell^{\lfloor\log_\ell2h\rfloor+1}(\ell-1)$,
we have $n_0(\ell,2h)\leq\lfloor\log_\ell2h\rfloor+1$, so
$m\leq2s+\lfloor\log_\ell2h\rfloor+2\leq2\lambda$. Let $u\in J$ with
$u\neq0$. Then $0<|u|\leq2h$, so $\ell^{v_\ell(u)}\leq|u|$ gives
$v_\ell(u)\leq\lfloor\log_\ell2h\rfloor$, and
$m+v_\ell(u)\leq2s+2\lfloor\log_\ell2h\rfloor+2\leq2\lambda$. The sum has at
most $n$ terms, so it is at most $2n\lambda$, and the sum of $2$ and the
maximum is at most $3+2n\lambda\leq3(n+1)\lambda$. With $v\leq m\leq2\lambda$,
the left side is at most $2\lambda+3(n+1)\lambda\leq20(n+1)\lambda$.
\end{proof}

\begin{definition}\label{asm-two-three-killing-holds-lift-block}
Let $\ell$ be a prime, fix an algebraic closure $\bar{\mathbf Q}_\ell$ of
$\mathbf Q_\ell$, and use the notation $G_K$, $N$-cyclotomic pairs and
Hodge-Tate weights along an embedding of
Definition~\ref{asm-two-three-local-killing}; let $\chi$ be the cyclotomic
character and $B_{\mathrm{cris}}$ Fontaine's ring of crystalline periods. For
integers $N\geq1$ and $h,n,t\geq0$ put
\[
 \begin{aligned}
 \lambda_\ell(N,h)&=\max(1,v_\ell(N))+\lfloor\log_\ell2h\rfloor
 +\lfloor\log_\ell N\rfloor+1,\\
 e_\ell(N,h)&=32\big(2h\max(1,v_\ell(N))+1\big)+\lambda_\ell(N,h),\qquad
 c^B_\ell(N,h,n,t)=40(t+1)(n+1)\lambda_\ell(N,h),
 \end{aligned}
\]
with $\lfloor\log_\ell y\rfloor$ as in Lemma~\ref{asm-two-three-split-numerics},
so that $c_\ell(N,h,n)=20(n+1)\lambda_\ell(N,h)$. Let $K,L\subset\bar{\mathbf
Q}_\ell$ be finite extensions of $\mathbf Q_\ell$. A \emph{lattice
representation} of $G_K$ over $L$ is a finite-dimensional $L$-vector space $V$
with an $L$-linear action of $G_K$ such that $g\mapsto gx$ is continuous for
every $x\in V$, together with a $G_K$-stable $\mathcal O_L$-lattice $T\subset V$.
\begin{enumerate}
\item A continuous map $c\colon G_K\to V$ with $c(gg')=c(g)+g\,c(g')$ is
  \emph{crystalline} if there is $x\in V\otimes_{\mathbf Q_\ell}B_{\mathrm{cris}}$
  with $c(g)\otimes1=gx-x$ for every $g\in G_K$. Let $Z^1_f(T)$ be the set of
  crystalline such maps with values in $T$, and $B^1(T)$ the set of maps
  $g\mapsto gy-y$ with $y\in T$. Let $B^1(V/T)$ be the set of maps
  $g\mapsto gz-z$ with $z\in V/T$, and $H^1_f(V/T)$ the set of maps
  $G_K\to V/T$ of the form $g\mapsto(c(g)\bmod T)+b(g)$ with $c$ crystalline
  and $b\in B^1(V/T)$. A \emph{divisible cocycle} is a locally constant map
  $c\colon G_K\to V/T$ with $c(gg')=c(g)+g\,c(g')$.
\item The \emph{lifting property} holds at $\ell$ if the following is true for
  all integers $N\geq1$ and $h\geq0$, every $N$-cyclotomic pair $(K,L)$, all
  lattice representations $(V_1,T_1)$ and $(V_2,T_2)$ of $G_K$ over $L$ with
  $V_1$ and $V_2$ crystalline, and every continuous $G_K$-equivariant
  $L$-linear map $f\colon V_1\to V_2$ with $f(T_1)=T_2$: if the Hodge-Tate
  weights of $V_1$ along every $\mathbf Q_\ell$-embedding $K\to L$ lie in
  $[-h,h]$, then for every $c\in Z^1_f(T_2)$ there are $c_1\in Z^1_f(T_1)$ and
  $b\in B^1(T_2)$ with $\ell^{e_\ell(N,h)}c=f\circ c_1+b$.
\item The \emph{block property} holds at $\ell$ if the following is true for
  all integers $N\geq1$ and $h,n,t\geq0$ with $t\leq h$, every $N$-cyclotomic
  pair $(K,L)$, every family $(J_\tau)$ of sets of integers indexed by the
  $\mathbf Q_\ell$-embeddings $\tau\colon K\to L$, and every set $U$ of at most
  $n$ integers in $[-h,h]$ with $J_\tau\subset U$ for every $\tau$. Put
  $c_B=c^B_\ell(N,h,n,t)$, let $I_K\subset G_K$ be the inertia subgroup, and
  for $0\leq e<t$ call $(V_e,T_e)$ a \emph{twist of weight $e$} if it is a
  lattice representation of $G_K$ over $L$ with $V_e=Lv$ one-dimensional,
  $T_e=\mathcal O_Lv$ and $gv=\chi(g)^{-e}v$. Then:
  \begin{enumerate}
  \item let $(V,T)$ be a lattice representation of $G_K$ over $L$ with $V$
    crystalline and its Hodge-Tate weights along every $\tau$ in $J_\tau$, let
    $p\colon V\to L$ be a $G_K$-invariant $L$-linear map with
    $p(T)\subset\mathcal O_L$, let $e_0\in T$ with $p(e_0)=1$, let
    $(V_e,T_e)$ be a twist of weight $e$ for some $0\leq e<t$, and let
    $\varphi\colon V\to V_e/T_e$ be an $\mathcal O_L$-linear map with
    $\varphi(gy)=g\varphi(y)$ for $g\in G_K$ and $y\in T$ with $p(y)=0$. Then
    the map $g\mapsto\ell^{c_B}\varphi(ge_0-e_0)$ lies in $H^1_f(V_e/T_e)$;
  \item let $(V_e,T_e)$ be a twist of weight $e$, $0\leq e<t$. If $e>0$, every
    divisible cocycle in $H^1_f(V_e/T_e)$ lies in $B^1(V_e/T_e)$. If $e=0$, a
    divisible cocycle lies in $H^1_f(V_e/T_e)$ if and only if it is the sum of
    an element of $B^1(V_e/T_e)$ and a map $G_K\to V_e/T_e$ that agrees on
    $I_K$ with $g\mapsto gz-z$ for some $z\in V_e/T_e$. The divisible cocycles
    $c$ with $\ell^{c_B}c\in H^1_f(V_e/T_e)$, modulo those in
    $H^1_f(V_e/T_e)$, form an $\mathcal O_L$-module of length at most
    $([K:\mathbf Q_\ell]+1)\,c_B$ times the length of
    $\mathcal O_L/\ell\mathcal O_L$.
  \end{enumerate}
\end{enumerate}
\end{definition}
\leanref{Development/AutomorphicArithmetic/ResidueCharacters/InvariantsKilledOverCyclotomicAlt.lean}{UpperTheory.N.LiftK}
\leanref{Development/AutomorphicArithmetic/ResidueCharacters/InvariantsKilledOverCyclotomicAlt.lean}{UpperTheory.N.TBK}
\leanref{Development/AutomorphicArithmetic/ResidueCharacters/InvariantsKilledOverCyclotomicAlt.lean}{UpperTheory.N.lamN}
\leanref{Development/AutomorphicArithmetic/ResidueCharacters/InvariantsKilledOverCyclotomicAlt.lean}{UpperTheory.N.eL}
\leanref{Development/AutomorphicArithmetic/ResidueCharacters/InvariantsKilledOverCyclotomicAlt.lean}{UpperTheory.N.cBL}
\leanref{Formalism/SharpRestriction/LevelLetterCocycles/QuotientHomNonvanishingBound.lean}{Formalism.MainNotions.N.TBAt}
\leanref{Formalism/WeaklyAdmissibleFiltrations/Statement.lean}{Formalism.WeaklyAdmissibleFiltrations.LatHom.LiftF}
\leanref{Formalism/WeaklyAdmissibleFiltrations/StatementTor.lean}{Formalism.WeaklyAdmissibleFiltrations.HfD}

\begin{definition}\label{asm-two-three-killing-holds-exponents}
Let $\ell$ be a prime, let $s\geq1$ and $M,H\geq0$ be integers, and let $J$ be a
finite set of integers. Put
\[
 J^+_M=\{j+m:j\in J,\ 1\leq m\leq M,\ j+m\neq0\}\cup\{1,\ldots,M\},
\]
\[
 e^W_\ell(s,H)=2Hs+v_\ell\big((2H)!\big),\qquad
 e^U_\ell(s,M,J)=\sum_{j\in J^+_M}\big(s+v_\ell(j)\big)+Ms+v_\ell(M!),
\]
and $e^P_\ell(s,M,H,J)=2e^W_\ell(s,H)+e^U_\ell(s,M,J)$.
\end{definition}
\leanref{Formalism/SharpRestriction/LevelLetterCocycles/LiftingTheory.lean}{Formalism.MainNotions.LiftingTheory.ePB}
\leanref{Formalism/SharpRestriction/LevelLetterCocycles/LiftingTheory.lean}{Formalism.MainNotions.LiftingTheory.JSet}
\leanref{Formalism/SharpRestriction/LevelLetterCocycles/LiftingTheory.lean}{Formalism.MainNotions.LiftingTheory.eWB}
\leanref{Formalism/LagrangeIdealLifts/Statement.lean}{Formalism.LIFT.eU}

\begin{lemma}\label{asm-two-three-killing-holds-exponent-bound}
Let $\ell$ be a prime, let $s\geq1$ and $M,h,H\geq0$ be integers with $M\leq H$
and $h\leq H$, and let $J$ be a finite set of integers in $[-M,h]$. Then, with
the notation of Definition~\ref{asm-two-three-killing-holds-exponents},
$e^P_\ell(s,M,H,J)\leq32Hs$.
\end{lemma}
\leanref{Formalism/SharpRestriction/LevelLetterCocycles/FrobeniusPowUniformizerFactorisation.lean}{Formalism.MainNotions.LiftingTheory.arith_sizeB}
\leanref{Formalism/WachLatticeTwists/LagrangeExtensionCocycles/ExtensionCocycleProofs/SizeBoundSecondCase.lean}{Formalism.RepresentationBoundsCohomology.N2.size_bound_second_case}
\begin{proof}
Put $H'=2H$, so that $M+h\leq H'$, and recall that $v_\ell(m!)\leq m$ for every
$m\geq0$. As $(2H)!$ divides $(2H')!$, $e^W_\ell(s,H)\leq e^W_\ell(s,H')$, and
\[
 2e^W_\ell(s,H')=4H's+2v_\ell\big((2H')!\big)\leq4H's+4H'\leq8H's.
\]
Every element of $J^+_M$ is a nonzero integer in $[1-M,h+M]$. So $J^+_M$ has at
most $h+2M\leq2H'$ elements, its negative elements are among $-1,\ldots,1-M$ and
its positive elements among $1,\ldots,h+M$, and hence
$\sum_{j\in J^+_M}v_\ell(j)\leq v_\ell(M!)+v_\ell\big((h+M)!\big)\leq M+h+M\leq2H'$.
Therefore
\[
 e^U_\ell(s,M,J)\leq2H's+2H'+H's+H'\leq6H's.
\]
Adding, $e^P_\ell(s,M,H,J)\leq2e^W_\ell(s,H')+e^U_\ell(s,M,J)\leq14H's\leq32Hs$.
% Lean: arith_sizeB = size_bound_second_case (StatSizeB, the fifth part of
% N2.lifting_size_bounds_all) at (M, h, 2H, s) with arith_eWB_mono and arith_ePB_eq; the bound
% there is 16 H' s. size_bound_second_case uses Size.JSet_card, Size.JSet_sum (sum_val_le) and
% Size.val_factorial_le_self (Legendre). About 140 lines.
\end{proof}

\begin{lemma}\label{asm-two-three-killing-holds-wach-lift}
Let $\ell$, $s$, $K$ and $L$ be as in
Definition~\ref{asm-two-three-killing-wach-frame-exists-data}, with $K=K_s$, and
let $D$ be a definite-side datum with local Hodge-Tate data, Fontaine data and
comparison data at every prime (Definition~\ref{asm-local-bound-two-three}(3)).
Use the notation $Z^1_f$ and $B^1$ of
Definition~\ref{asm-two-three-killing-holds-lift-block}(1) and $e^P_\ell$ of
Definition~\ref{asm-two-three-killing-holds-exponents}. Let $h\geq0$, let
$(V_1,T_1)$ and $(V_2,T_2)$ be lattice representations of $G_K$ over $L$ with
$V_1$ and $V_2$ crystalline, and let $f\colon V_1\to V_2$ be a continuous
$G_K$-equivariant $L$-linear map with $f(T_1)=T_2$. Suppose that the Hodge-Tate
weights of $V_1$ along every $\mathbf Q_\ell$-embedding $K\to L$ lie in
$[-h,h]$, and let $J$ be the set of these weights. Then for every $c\in
Z^1_f(T_2)$ there are $c_1\in Z^1_f(T_1)$ and $b\in B^1(T_2)$ with
$\ell^ec=f\circ c_1+b$, where $e=e^P_\ell(s,h,2h,J)$.
\end{lemma}
\leanref{Formalism/SharpRestriction/LevelLetterCocycles/LiftingTheory.lean}{Formalism.MainNotions.LiftingTheory.coreLiftB}
\leanref{Formalism/SharpRestriction/LevelLetterCocycles/BaseChangePreservesUnitaryOf.lean}{Formalism.MainNotions.LiftingTheory.lifting_case_B_of_hypotheses}
\leanref{Development/AutomorphicArithmetic/ResidueCharacters/LocRes.lean}{UpperTheory.N.LocRes.caseB}
\begin{proof}
\openstep{The proof, Alpöge's lifting theorem for crystalline lattices over the
field $K_s$, which works with their Wach modules, is not yet written out.}
% Lean: coreLiftB ((HW l).W K s LL h) under CaseB W (LocRes.caseB: W.GKs = top and [K : K_0] =
% (l - 1) l^(s - 1), both from K = K_s), read for constant bounds Mi = hi = h, so Mmax = h,
% Hmax = 2h (as n = [K_0 : Q_l] >= 1) and jumpSet = J. LiftTheoremCaseB is
% lifting_case_B_of_hypotheses (coreLiftB_of' with gam_zl, lker_HKer_D, lqb_HLeft_holds,
% hstr_B, liftb_HExt_holds, liftb_HOne_holds) from wach_statement_case_B and the admitted
% conditions a1, a2, a5, a6, a11, a13 to a17 and a21 to a23 of LiftAdmitted.all_admitted. None
% of the published results of Theorem lit-bundle-unitary(6) enters:
% localConstants_from_admitted_local_Galois_statements discards its three hypotheses. Size,
% following imports and leaving out what the book states: the proof closure of
% lifting_case_B_of_hypotheses, wach_statement_case_B and all_admitted is about 59,500 lines in
% 709 files, an overcount, since all_admitted also proves the conditions used only in the
% Lean's case A.
\end{proof}

\begin{lemma}\label{asm-two-three-killing-holds-transfer}
Let $G$ be a group, $H\subset G$ a subgroup of finite index $d$, and $X$ an
abelian group on which $G$ acts by group automorphisms. For every coset $q\in
G/H$ choose $r_q\in q$. For a map $f\colon H\to X$ define $\operatorname{tr}f
\colon G\to X$ by
\[
 \operatorname{tr}f(g)=\sum_{q\in G/H}r_{gq}\,f\big(r_{gq}^{-1}gr_q\big).
\]
Then:
\begin{enumerate}
\item $r_{gq}^{-1}gr_q\in H$ for $g\in G$ and $q\in G/H$;
\item if $f(ab)=f(a)+af(b)$ for $a,b\in H$, then
  $\operatorname{tr}f(ab)=\operatorname{tr}f(a)+a\operatorname{tr}f(b)$ for
  $a,b\in G$;
\item if $x\in X$ and $f(a)=ax-x$ for $a\in H$, then
  $\operatorname{tr}f(g)=gy-y$ for $g\in G$, where $y=\sum_qr_qx$;
\item if $f\colon G\to X$ satisfies $f(ab)=f(a)+af(b)$ for $a,b\in G$, then
  $\operatorname{tr}(f|_H)(g)=d\,f(g)+(gz-z)$ for $g\in G$, where
  $z=\sum_qf(r_q)$;
\item if $G$ is a topological group, $H$ is open, $X$ is a topological group on
  which every $g\in G$ acts continuously, and $f$ is continuous, then
  $\operatorname{tr}f$ is continuous.
\end{enumerate}
\end{lemma}
\leanref{Development/AbelianVarietyTheory/Proofs/InertiaEigenvalues/TransferCocycleFormula.lean}{GlobalTheory.transfer_cocycle_formula}
\leanref{Development/AbelianVarietyTheory/Proofs/InertiaEigenvalues/TransferContinuousOfContinuousOn.lean}{GlobalTheory.transfer_continuous_of_continuousOn}
\begin{proof}
(1) $r_{gq}$ and $gr_q$ lie in the same coset $gq$. For $g\in G$ the map
$q\mapsto gq$ is a bijection of $G/H$, which allows the sums below to be
reindexed.

(2) For $a,b\in G$ and $q\in G/H$,
$r_{abq}^{-1}abr_q=\big(r_{a(bq)}^{-1}ar_{bq}\big)\big(r_{bq}^{-1}br_q\big)$, a
product of two elements of $H$. So
$r_{abq}f(r_{abq}^{-1}abr_q)=r_{a(bq)}f(r_{a(bq)}^{-1}ar_{bq})
+ar_{bq}f(r_{bq}^{-1}br_q)$. Summing over $q$, and reindexing the first sum by
$q\mapsto bq$, gives $\operatorname{tr}f(a)+a\operatorname{tr}f(b)$.

(3) Each term is $r_{gq}(r_{gq}^{-1}gr_qx-x)=gr_qx-r_{gq}x$; summing over $q$
and reindexing the second sum gives $gy-y$.

(4) The cocycle relation gives $f(1)=0$, $f(r^{-1})=-r^{-1}f(r)$ and
$f(r^{-1}gr')=f(r^{-1})+r^{-1}f(g)+r^{-1}gf(r')$. So each term is
$r_{gq}f(r_{gq}^{-1}gr_q)=f(g)+gf(r_q)-f(r_{gq})$; summing over $q$ and
reindexing the last sum gives $d\,f(g)+gz-z$.

(5) For each $q$, the map $g\mapsto gq$ is constant on the open sets
$g_0r_qHr_q^{-1}$, $g_0\in G$, so on such a set the $q$-term is the composite of
the continuous map $g\mapsto r_{g_0q}^{-1}gr_q$ into $H$, $f$, and the action
of $r_{g_0q}$.
% Lean: transfer_cocycle_formula (the four parts, with q.out as r_q; reidx is the
% reindexing) and transfer_continuous_of_continuousOn (for f continuous on H, with
% psi g = (g q.out)^(-1) (g0 q).out). About 100 lines.
\end{proof}

\begin{lemma}\label{asm-two-three-killing-holds-descent}
Let $\ell$ be a prime, let $K\subset K'$ and $L$ be finite extensions of
$\mathbf Q_\ell$ in $\bar{\mathbf Q}_\ell$, and use the notation $Z^1_f$ and
$B^1$ of Definition~\ref{asm-two-three-killing-holds-lift-block}(1) for
lattice representations of $G_K$ and of $G_{K'}$. Let $(V_1,T_1)$ and
$(V_2,T_2)$ be lattice representations of $G_K$ over $L$, and let $f\colon
V_1\to V_2$ be a continuous $G_K$-equivariant $L$-linear map with
$f(T_1)\subset T_2$. Let $E\geq0$, and suppose that for every $c'\in Z^1_f(T_2)$,
for the restriction of $(V_2,T_2)$ to $G_{K'}$, there are $c'_1\in Z^1_f(T_1)$
and $b'\in B^1(T_2)$, for the restrictions to $G_{K'}$, with
$\ell^Ec'=f\circ c'_1+b'$. Put $v=v_\ell([K':K])$. Then for every $c\in
Z^1_f(T_2)$ there are $c_1\in Z^1_f(T_1)$ and $b\in B^1(T_2)$ with
$\ell^{E+v}c=f\circ c_1+b$.
\end{lemma}
\leanref{Formalism/SharpRestriction/LevelLetterCocycles/FrobeniusPowUniformizerFactorisation.lean}{Formalism.MainNotions.LiftingTheory.desc_LiftF}
\leanref{Development/AbelianVarietyTheory/Proofs/InertiaEigenvalues/RestrictionCorestrictionH1f.lean}{GlobalTheory.InducedSectionsSpanStatement.coreCor_of_isOpen}
\leanref{Development/AbelianVarietyTheory/Proofs/InertiaEigenvalues/RestrictionCorestrictionH1f.lean}{GlobalTheory.res_cor_H1f}
\begin{proof}
Put $G=G_K$ and $H=G_{K'}$, an open subgroup of index $d=[K':K]$, and choose
coset representatives $r_q$ as in
Lemma~\ref{asm-two-three-killing-holds-transfer}. For a lattice representation
$(V,T)$ of $G$ let $H^1(T)$ be the $\mathcal O_L$-module of continuous maps
$c\colon G\to T$ with $c(gg')=c(g)+gc(g')$, modulo $B^1(T)$, and $H^1_f(T)$
the image of $Z^1_f(T)$ in it; define $H^1(T)$ and $H^1_f(T)$ over $H$ in the
same way. A continuous $G$-equivariant $L$-linear map $f$ with $f(T_1)\subset
T_2$ induces $f_*\colon H^1(T_1)\to H^1(T_2)$, over $G$ and over $H$, as it maps
$Z^1_f(T_1)$ into $Z^1_f(T_2)$ and $B^1(T_1)$ into $B^1(T_2)$. In this notation
the hypothesis says that $\ell^E\kappa\in f_*(H^1_f(T_1))$ for every
$\kappa\in H^1_f(T_2)$ over $H$, and the conclusion says the same over $G$
with $E+v$.

First, restriction of cocycles to $H$ gives
$\operatorname{res}\colon H^1(T)\to H^1(T)$ from $G$ to $H$, and it maps
$H^1_f(T)$ into $H^1_f(T)$: if $c(g)\otimes1=gx-x$ for every $g\in G$, with
$x\in V\otimes B_{\mathrm{cris}}$, the same holds on $H$.

Second, the map $\operatorname{tr}$ of
Lemma~\ref{asm-two-three-killing-holds-transfer}, with $X=T$, induces
$\operatorname{cor}\colon H^1(T)\to H^1(T)$ from $H$ to $G$. Indeed $T$ is
stable under $G$, so $\operatorname{tr}c$ takes values in $T$; it is a
continuous cocycle by (2) and (5) of that lemma, and it maps $B^1(T)$ over $H$
into $B^1(T)$ over $G$ by (3), as $\sum_qr_qx\in T$ for $x\in T$.

Third, $\operatorname{cor}\circ\operatorname{res}$ is multiplication by $d$, by
(4) of that lemma, as $\sum_qc(r_q)\in T$.

Fourth, $\operatorname{cor}$ maps $H^1_f(T)$ over $H$ into $H^1_f(T)$ over $G$.
Let $c'$ be a continuous cocycle on $H$ with values in $T$ and $c'(a)\otimes1=
ax-x$ for $a\in H$, with $x\in V\otimes B_{\mathrm{cris}}$. Apply (3) of that
lemma to $X=V\otimes B_{\mathrm{cris}}$, with $G$ acting on both factors, and
to $a\mapsto c'(a)\otimes1$: as $\operatorname{tr}(c'\otimes1)=
(\operatorname{tr}c')\otimes1$, this gives $\operatorname{tr}c'(g)\otimes1=gy-y$
with $y=\sum_qr_qx$.

Fifth, $\operatorname{cor}\circ f_*=f_*\circ\operatorname{cor}$, since
$\operatorname{tr}(f\circ c')=f\circ\operatorname{tr}c'$ for $G$-equivariant
$f$.

Now let $c\in Z^1_f(T_2)$ and $\kappa$ its class. By the first step and the
hypothesis, $\ell^E\operatorname{res}\kappa=f_*\kappa'_1$ for some
$\kappa'_1\in H^1_f(T_1)$ over $H$. By the third and fifth steps,
$\ell^Ed\,\kappa=\operatorname{cor}(\ell^E\operatorname{res}\kappa)
=f_*(\operatorname{cor}\kappa'_1)$, and $\operatorname{cor}\kappa'_1\in
H^1_f(T_1)$ by the fourth. Write $d=\ell^vu$ with $u$ prime to $\ell$; then $u$
is a unit of $\mathbf Z_\ell\subset\mathcal O_L$, and
$\ell^{E+v}\kappa=f_*(u^{-1}\operatorname{cor}\kappa'_1)$. Choosing $c_1\in
Z^1_f(T_1)$ in the class $u^{-1}\operatorname{cor}\kappa'_1$ gives
$\ell^{E+v}c-f\circ c_1\in B^1(T_2)$.
% Lean: desc_LiftF (desc_LiftF_iff, desc_sub: the argument of the last paragraph at the level
% of LatticeCohomology and latticeFiniteCohomology) with coreCor from res_cor_H1f =
% coreCor_of_isOpen: resFun_Z1T and resFun_B1T (first step), corH1 from trCochain_Z1T and
% trCochain_B1T (second), corH1_res (third), the finite part through repVB and actVB_res
% (fourth), corH1_natural (fifth); the index is inclusion_GK_index, [G_K : G_K'] = [K' : K].
% About 420 lines with Lemma -transfer.
\end{proof}

\begin{lemma}\label{asm-two-three-killing-holds-block-top}
Let $\ell$ be a prime, let $N\geq1$ and $h,n,t,e,v\geq0$ be integers with
$e<t\leq h$ and $v\leq\lfloor\log_\ell N\rfloor+1$, and put
$s=\max(1,v_\ell(N))$. Let $K$ and $L$ be such that $\ell$, $s$, $K$ and $L$ are
as in Definition~\ref{asm-two-three-killing-wach-frame-exists-data}, with
$K=K_s$, and let $D$ be a definite-side datum with local Hodge-Tate data,
Fontaine data and comparison data at every prime
(Definition~\ref{asm-local-bound-two-three}(3)). Use the notation of
Definition~\ref{asm-two-three-killing-holds-lift-block}. Let $U$ be a set of at
most $n$ integers in $[-h,h]$, and let $(V,T)$ be a lattice representation of
$G_K$ over $L$ with $V$ crystalline and its Hodge-Tate weights along every
$\mathbf Q_\ell$-embedding $K\to L$ in $U$. Let $p\colon V\to L$ be a
$G_K$-invariant $L$-linear map with $p(T)\subset\mathcal O_L$, and $e_0\in T$
with $p(e_0)=1$. Let $(V_e,T_e)$ be a lattice representation of $G_K$ over $L$
with $V_e$ one-dimensional and $gx=\chi(g)^{-e}x$ for $g\in G_K$ and $x\in
V_e$, and let $\varphi\colon V\to V_e/T_e$ be an $\mathcal O_L$-linear map with
$\varphi(gy)=g\varphi(y)$ for $g\in G_K$ and $y\in T$ with $p(y)=0$. Then there
is an integer $c\geq0$ with $c+v\leq c^B_\ell(N,h,n,t)$ such that the map
$g\mapsto\ell^c\varphi(ge_0-e_0)$ lies in $H^1_f(V_e/T_e)$.
\end{lemma}
\leanref{Development/AutomorphicArithmetic/ResidueCharacters/TbTop.lean}{UpperTheory.N.tb_top}
\begin{proof}
\openstep{The proof is not yet written out. It applies Alpöge's bound for the
images, under equivariant maps to $V_e/T_e$, of the class of an extension of the
trivial lattice by a crystalline lattice over $K_s$ to the extension of
$p(T)=\mathcal O_L$ by the kernel of $p$ on $T$, and bounds the exponent of that
bound by $c^B_\ell(N,h,n,t)-v$.}
% Lean: tb_top D HW F l h4 hCo hW hF, with h4 = Conclusion4 from
% T.theorem_tb_of_local_blochKato_duality_euler bk tate euler (TheoremTBRecords: from
% LocInputs.ties_and_inputs_of_duality_euler tate euler and blochKatoLocal_all_of_tateDuality_eulerChar,
% through twistedTorsionClassTheorem_goodCase), hCo = crysOnCond_all (Lemma
% asm-two-three-killing-wach-frame-exists-inputs(1)) and hW = crystallineLatticeWachModuleExists_proof
% (Lemma asm-two-three-killing-wach-frame-exists). Its proof: Y = kerY (the lattice of y in T with
% p(y) = 0), cj_sub (Y crystalline with weights in U), JsetEmb.jsetSplit_all, the Wach hypothesis
% for Y, TBLoc.tb_part1 (the exact sequence 0 -> Y -> T -> O_L -> 0 and TBat, the first part of
% Conclusion4, with exponent cTBY), and LocNum.cTB_le (cTBY + v <= cBL, with LocRes.vdeg_le and
% gK = 0 as K = K_s). The Lean writes the integer v as g. Size, following imports and leaving
% out what the book states: the proof closure of tb_top, theorem_tb_of_local_blochKato_duality_euler
% and crysOnCond_of_every_datum is about 42,000 lines in 586 files.
\end{proof}

\begin{lemma}\label{asm-two-three-killing-holds-block-transfer}
Let $\ell$ be a prime, let $K\subset K'$ and $L$ be finite extensions of
$\mathbf Q_\ell$ in $\bar{\mathbf Q}_\ell$, and use the notation $H^1_f$ of
Definition~\ref{asm-two-three-killing-holds-lift-block}(1) for lattice
representations of $G_K$ and of $G_{K'}$. Let $(V,T)$ be a lattice
representation of $G_K$ over $L$, and let $c\colon G_K\to V/T$ be a map with
$c(gg')=c(g)+g\,c(g')$ whose restriction to $G_{K'}$ lies in $H^1_f(V/T)$ for
the restriction of $(V,T)$ to $G_{K'}$. Then $\ell^vc\in H^1_f(V/T)$, where
$v=v_\ell([K':K])$.
\end{lemma}
\leanref{Development/AutomorphicArithmetic/ResidueCharacters/TransferMemBlochKatoCocycles.lean}{UpperTheory.N.CorD.index_smul_mem}
\leanref{Development/AutomorphicArithmetic/ResidueCharacters/TransferMemBlochKatoCocycles.lean}{UpperTheory.N.CorD.transfer_mem_blochKato_cocycles}
\leanref{Development/AutomorphicArithmetic/ResidueCharacters/LocRes.lean}{UpperTheory.N.LocRes.pow_val_smul_mem}
\begin{proof}
Put $G=G_K$ and $H=G_{K'}$, of index $d=[K':K]$, and use the map
$\operatorname{tr}$ of Lemma~\ref{asm-two-three-killing-holds-transfer} for
coset representatives $r_q$. Write $c|_H=(z\bmod T)+(a\mapsto aw-w)$, with
$w\in V/T$ and $z\colon H\to V$ a continuous map with $z(ab)=z(a)+az(b)$ and
$z(a)\otimes1=ax-x$ for $a\in H$, for some $x\in V\otimes B_{\mathrm{cris}}$.
By (4) of that lemma, applied to $X=V/T$,
$\operatorname{tr}(c|_H)(g)=d\,c(g)+(gS'-S')$ with $S'=\sum_qc(r_q)$. By (3),
applied to $X=V/T$, and as reduction modulo $T$ commutes with the action,
$\operatorname{tr}(c|_H)=(\operatorname{tr}z\bmod T)+(g\mapsto gS-S)$ with
$S=\sum_qr_qw$. The map $\operatorname{tr}z\colon G\to V$ is a continuous cocycle
by (2) and (5), and $\operatorname{tr}z(g)\otimes1=gy-y$ with
$y=\sum_qr_qx$ by (3) applied to $X=V\otimes B_{\mathrm{cris}}$. Hence
\[
 d\,c=(\operatorname{tr}z\bmod T)+\big(g\mapsto g(S-S')-(S-S')\big)
\]
lies in $H^1_f(V/T)$. Write $d=\ell^vu$ with $u$ prime to $\ell$, a unit of
$\mathcal O_L$. As $H^1_f(V/T)$ is an $\mathcal O_L$-submodule,
$\ell^vc=u^{-1}(d\,c)\in H^1_f(V/T)$.
% Lean: index_smul_mem (hcsplit, transfer_cocycle_formula parts 3 and 4, hkey) with
% transfer_mem_blochKato_cocycles (continuity, cocycle and the crystalline element
% sum_q repVB q.out x), then pow_val_smul_mem (d = u l^v with u = ordCompl[l] d a unit of
% Z_l). About 250 lines with Lemma -transfer.
\end{proof}

\begin{lemma}\label{asm-two-three-killing-holds-block-twists}
Let $\ell$ be a prime, let $K$ and $L$ be finite extensions of $\mathbf Q_\ell$
in $\bar{\mathbf Q}_\ell$, let $D$ be a definite-side datum with local
Hodge-Tate data, Fontaine data and comparison data at every prime
(Definition~\ref{asm-local-bound-two-three}(3)), and use the notation of
Definition~\ref{asm-two-three-killing-holds-lift-block}, with $I_K\subset G_K$
the inertia subgroup. Let $e\geq0$, and let $(V_e,T_e)$ be a lattice
representation of $G_K$ over $L$ with $V_e$ one-dimensional and
$gx=\chi(g)^{-e}x$ for $g\in G_K$ and $x\in V_e$. Let $Z_{\mathrm{div}}$ be the
$\mathcal O_L$-module of divisible cocycles $G_K\to V_e/T_e$, put
$H^1_{\mathrm{div}}=Z_{\mathrm{div}}/(Z_{\mathrm{div}}\cap B^1(V_e/T_e))$, and let
$H^1_{f,\mathrm{div}}\subset H^1_{\mathrm{div}}$ be the image of
$Z_{\mathrm{div}}\cap H^1_f(V_e/T_e)$. Then:
\begin{enumerate}
\item for every $m\geq0$, the elements of $H^1_{\mathrm{div}}/H^1_{f,\mathrm{div}}$
  killed by $\ell^m$ form an $\mathcal O_L$-module of length at most
  $([K:\mathbf Q_\ell]+1)m$ times the length of $\mathcal O_L/\ell\mathcal O_L$;
\item if $e>0$, then $H^1_f(V_e/T_e)=B^1(V_e/T_e)$;
\item if $e=0$, a map $c\colon G_K\to V_e/T_e$ lies in $H^1_f(V_e/T_e)$ if and
  only if it is a divisible cocycle and there is $x\in V_e/T_e$ with
  $c(g)=gx-x$ for every $g\in I_K$.
\end{enumerate}
\end{lemma}
\leanref{Development/AutomorphicArithmetic/ResidueCharacters/TwistDef.lean}{UpperTheory.N.TwistK.TBe}
\leanref{Development/AutomorphicArithmetic/ResidueCharacters/TwistK.lean}{UpperTheory.N.TwistK.tbe_all}
\begin{proof}
\openstep{The proof, which uses the local results of Bloch and Kato, local Tate
duality and the local Euler characteristic formula, is not yet written out.}
% Lean: TwistK.tbe_all D HW F hA (blochKatoLocal_all_of_tateDuality_eulerChar bk tate euler D)
% l K LL, with hA from LocInputs.ties_and_inputs_of_duality_euler tate euler: tbe_of with the
% admitted fields hA5 (tate_fil_quot) and hA7 (fontaine_bounds), Conclusion5 (local Bloch-Kato),
% CocycleCount.cocycles_finite_ncard_le (from duality and Euler characteristic),
% LocalIntegers.good_integers, exists_topologicalClosure_zpowers_sup_inertia, exists_unr and
% Cyc.exists_mem_GK_chi_pow_ne_one. TBe is stated for the period data (D.loc l).PD K, read in
% Fontaine's sense, and I_K is IK K. Size, following imports and leaving out what the book
% states: the proof closure of tbe_all, ties_and_inputs_of_duality_euler,
% blochKatoLocal_all_of_tateDuality_eulerChar and TBLoc.tbe_clauses is about 30,000 lines in
% 431 files; most of it is shared with that of Lemma -block-top.
\end{proof}

\begin{proposition}\label{asm-two-three-killing-holds}
Assume the third collection of published results in
Theorem~\ref{lit-bundle-unitary}(6). The proof uses three of them: the local
results of Bloch and Kato, and local Tate duality and the local Euler
characteristic formula over finite extensions of $\mathbf Q_p$ for every prime
$p$. It does not use Fontaine's theory of $B_{\mathrm{cris}}$, global Artin
reciprocity, local finiteness or the integral form of the theorem of Bhatt and
Scholze, and the proofs of the invariant-killing property and of the lifting
property use none of the collection. Let $D$ be a definite-side datum with
local Hodge-Tate data, Fontaine data and comparison data at every prime
(Definition~\ref{asm-local-bound-two-three}(3)). Then at every prime $\ell$ the
lifting property, the invariant-killing property
(Definition~\ref{asm-two-three-local-killing}(3)) and the block property
(Definition~\ref{asm-two-three-killing-holds-lift-block}) hold.
\end{proposition}
\leanref{Development/AutomorphicArithmetic/ResidueCharacters/LocAll.lean}{UpperTheory.N.locAll}
\leanref{Development/AutomorphicArithmetic/ResidueCharacters/LocAllOf.lean}{UpperTheory.N.locAll_of}
\leanref{Development/AutomorphicArithmetic/ResidueCharacters/InvariantsKilledOverCyclotomic.lean}{UpperTheory.N.invariants_killed_over_cyclotomic_subfield_of_admissible}
\leanref{Development/AutomorphicArithmetic/ResidueCharacters/LocNum.lean}{UpperTheory.N.LocNum.exists_up}
\leanref{Development/AutomorphicArithmetic/ResidueCharacters/FiniteRoots.lean}{UpperTheory.N.Cyc.Ks_lcm}
\leanref{Development/AutomorphicArithmetic/ResidueCharacters/LocRes.lean}{UpperTheory.N.LocRes.res_crys}
\leanref{Development/AutomorphicArithmetic/ResidueCharacters/LocRes.lean}{UpperTheory.N.LocRes.vdeg_le}
\leanref{Development/AutomorphicArithmetic/ResidueCharacters/LiftKOf.lean}{UpperTheory.N.liftK_of}
\leanref{Development/AutomorphicArithmetic/ResidueCharacters/TbKOf.lean}{UpperTheory.N.tbK_of}
\leanref{Development/AutomorphicArithmetic/ResidueCharacters/FiniteRoots.lean}{UpperTheory.N.Cyc.val_index_le}
\begin{proof}
\emph{The invariant-killing property.}
Let $N\geq1$ and $h,n\geq0$ be integers, let $(K,L)$ be an $N$-cyclotomic
pair, and let $V$, $T$ and $U$ be as in
Definition~\ref{asm-two-three-local-killing}(3), with $V$ crystalline and its
Hodge-Tate weights along every $\mathbf Q_\ell$-embedding $K\to L$ in $U$ and
different from $0$. Let $x\in V/T$ be $G_K$-invariant. Put
$M=\operatorname{lcm}(\ell,N)$, $s=\max(1,v_\ell(N))$ and $c=c_\ell(N,h,n)$,
let $K'=\mathbf Q_\ell(\mu_N^\ell)$ be the field generated by the $M$-th roots
of unity, and use the notation $K'_j$ and $n_0$ of
Lemma~\ref{asm-two-three-killing-wach}.

First, $K\subset K'$, since $(K,L)$ is $N$-cyclotomic, and $K'$, $s$ and $L$
satisfy the hypotheses of Lemma~\ref{asm-two-three-killing-wach}. Indeed,
$s\geq1$, and if $\ell=2$ then $4\mid N$, so $s\geq2$. Next, $K'_s=K'$. We have
$v_\ell(M)=s$; write $M=\ell^sM'$ with $\ell\nmid M'$. The field $K'_s$ is
generated by $K'_0\subset K'$ and the $\ell^s$-th roots of unity, which are
$M$-th roots of unity, so $K'_s\subset K'$. Conversely, let $\zeta$ be an
$M$-th root of unity and choose integers $a,b$ with $a\ell^s+bM'=1$. Then
$\zeta=(\zeta^{\ell^s})^a(\zeta^{M'})^b$, where $\zeta^{\ell^s}$ is an $M'$-th
root of unity, so it lies in the unramified subextension $K'_0$, and
$\zeta^{M'}$ is an $\ell^s$-th root of unity. So $\zeta\in K'_s$, and
$K'\subset K'_s$. Finally, a $\mathbf Q_\ell$-embedding of $K'$ in
$\bar{\mathbf Q}_\ell$ maps the $M$-th roots of unity, which generate $K'$, to
$M$-th roots of unity, and these lie in $L$ because $(K,L)$ is $N$-cyclotomic.

Second, restrict the action to $G_{K'}\subset G_K$. By Fontaine's theory
\cite{Fon94}, $V$ is crystalline as a representation of $G_{K'}$, and for every
$\mathbf Q_\ell$-embedding $\tau'\colon K'\to L$ its Hodge-Tate weights along
$\tau'$ are those of $V$ along $\tau'|_K$. So $0$ is not a Hodge-Tate weight of
$V|_{G_{K'}}$ along any $\tau'$. Since $L$ contains the images of all the
embeddings of $K'$, there are $[K':\mathbf Q_\ell]$ of them, and
$K'\otimes_{\mathbf Q_\ell}L$ is the product of copies of $L$ indexed by the
embeddings $\tau'$. So every jump $a$ of
the Hodge filtration of $D_{\mathrm{dR}}(V|_{G_{K'}})$ is a Hodge-Tate weight
along some $\tau'$. So the set $J'$ of these jumps lies in $U$, and in
$[-h,h]$.

Third, $T$ is a $G_{K'}$-stable lattice and $x$ is $G_{K'}$-invariant.
Lemma~\ref{asm-two-three-killing-wach}, applied to $K'$, $s$, $L$, $D$, $V$,
$T$ and $M_0=h_0=h$, gives $\ell^ex=0$, where $m=2s+n_0(\ell,2h)+1$ and
\[
 e=v+2+\max\Big(1,\sum_{u\in J',\,u\neq0}\big(m+v_\ell(u)\big)\Big),\qquad
 v=v_\ell[K'_m:\mathbf Q_\ell]-v_\ell[K':\mathbf Q_\ell].
\]

Fourth, $0\leq v\leq m$. For $j\geq1$ the $\ell^j$-th cyclotomic polynomial
stays irreducible over the unramified field $K'_0$, so
$[K'_j:\mathbf Q_\ell]=[K'_0:\mathbf Q_\ell](\ell-1)\ell^{j-1}$. Since
$K'=K'_s$ and $m\geq s$, this gives $v=m-s$.

Finally, Lemma~\ref{asm-two-three-killing-count}, with $J'\subset U$ and this
$v$, gives $e\leq c$. So $\ell^cx=\ell^{c-e}\ell^ex=0$.

\emph{The lifting property.} Let $N\geq1$ and $h\geq0$ be integers, let $(K,L)$
be an $N$-cyclotomic pair, and let $(V_1,T_1)$, $(V_2,T_2)$ and $f$ be as in
Definition~\ref{asm-two-three-killing-holds-lift-block}(2), with the Hodge-Tate
weights of $V_1$ along every $\mathbf Q_\ell$-embedding $K\to L$ in $[-h,h]$. Let
$M$, $s$ and $K'$ be as in the proof of the invariant-killing property. By its
first step, $K\subset K'=K'_s$, and $\ell$, $s$, $K'$ and $L$ are as in
Definition~\ref{asm-two-three-killing-wach-frame-exists-data}.

First, put $v=v_\ell([K':K])$. Then $v\leq\lfloor\log_\ell N\rfloor+1$. Indeed
$\ell^v\leq[K':K]\leq[K':\mathbf Q_\ell]$, and $[K':\mathbf Q_\ell]\leq
M\leq\ell N$, as $K'$ is generated by a primitive $M$-th root of unity, a root
of $X^M-1$.

Second, restricted to $G_{K'}$, the representations $V_1$ and $V_2$ are
crystalline, and the Hodge-Tate weights of $V_1$ along a $\mathbf
Q_\ell$-embedding $\tau'\colon K'\to L$ are those along $\tau'|_K$
\cite{Fon94}, so they lie in $[-h,h]$; let $J'$ be the set of these weights. By
Lemma~\ref{asm-two-three-killing-holds-wach-lift} for $K'$, $s$ and $L$, the
hypothesis of Lemma~\ref{asm-two-three-killing-holds-descent} holds with
$E=e^P_\ell(s,h,2h,J')$. That lemma gives, for every $c\in Z^1_f(T_2)$, some
$c_1\in Z^1_f(T_1)$ and $b\in B^1(T_2)$ with $\ell^{E+v}c=f\circ c_1+b$.

Third, $E\leq32\cdot2h\cdot s\leq32(2hs+1)$ by
Lemma~\ref{asm-two-three-killing-holds-exponent-bound} with $M=h$ and $H=2h$, and
$v\leq\lfloor\log_\ell N\rfloor+1\leq\lambda_\ell(N,h)$. So $E+v\leq
e_\ell(N,h)$, and multiplying by $\ell^{e_\ell(N,h)-E-v}$ gives
$\ell^{e_\ell(N,h)}c=f\circ c'_1+b''$ with $c'_1\in Z^1_f(T_1)$ and $b''\in
B^1(T_2)$.
% Lean: liftK_of D HW F l h54 hres hF. LocNum.exists_up gives K' (Cyc.Kc l (lcm l N)) with
% Cyc.val_index_le (first step: padicValNat_le_nat_log and finrank_le, [K' : Q_l] <= lcm(l, N));
% LocRes.caseB gives the Lean's case B at (K', s, LL); LocRes.res_crys (hres =
% T.crystalline_restrict_finite_extension_jumps, hF = FontaineAll.fontaineAt_all, cited as
% Fontaine's theory as in the invariant-killing part) is the second step; h1 is the case B part
% of LiftGaloisDescent (lifting_descent_to_subfields = coreDesc_of: coreLiftB for the
% restriction, Lemma -wach-lift, then desc_LiftF, Lemma -descent); h2 is the second part of
% LiftExponentShapes (coreLiftO_holds from arith_sizeB, Lemma -exponent-bound); Hmax_const_le,
% LiftF_mono and omega give the third step. h54 is
% T.localConstants_from_admitted_local_Galois_statements bk tate euler D HW F Cm, which discards
% bk, tate and euler (= frobenius_lifting with LiftAdmitted.all_admitted).

\emph{The block property.} Let $N\geq1$ and $h,n,t\geq0$ be integers with $t\leq
h$, let $(K,L)$ be an $N$-cyclotomic pair, let $(J_\tau)$ and $U$ be as in
Definition~\ref{asm-two-three-killing-holds-lift-block}(3), and put
$c_B=c^B_\ell(N,h,n,t)$. Let $M$, $s$, $K'$ and $v$ be as above.

(a) Let $(V,T)$, $p$, $e_0$, a twist $(V_e,T_e)$ of weight $e<t$ and $\varphi$ be
as in (a). Restricted to $G_{K'}$, $V$ is crystalline and its Hodge-Tate weights
along every $\mathbf Q_\ell$-embedding $\tau'\colon K'\to L$ lie in
$J_{\tau'|_K}\subset U$ \cite{Fon94}; $V_e$ is one-dimensional with
$gx=\chi(g)^{-e}x$; and $p$, $e_0$ and $\varphi$ keep their properties.
Lemma~\ref{asm-two-three-killing-holds-block-top}, for $K'$, $s$, $L$ and this
$v$, gives $c\geq0$ with $c+v\leq c_B$ such that the restriction to $G_{K'}$ of
$\gamma(g)=\ell^c\varphi(ge_0-e_0)$ lies in $H^1_f(V_e/T_e)$ for the restriction
to $G_{K'}$. The map $\gamma$ satisfies $\gamma(gg')=\gamma(g)+g\gamma(g')$ on
$G_K$, since $g'e_0-e_0\in T$ and $p(g'e_0-e_0)=0$, so that
$\varphi(gg'e_0-ge_0)=g\varphi(g'e_0-e_0)$. So
Lemma~\ref{asm-two-three-killing-holds-block-transfer} gives $\ell^v\gamma\in
H^1_f(V_e/T_e)$, and $g\mapsto\ell^{c_B}\varphi(ge_0-e_0)$, which is
$\ell^{c_B-c-v}\cdot\ell^v\gamma$, lies in $H^1_f(V_e/T_e)$.

(b) Let $(V_e,T_e)$ be a twist of weight $e$, $0\leq e<t$, and use the notation
$Z_{\mathrm{div}}$, $H^1_{\mathrm{div}}$ and $H^1_{f,\mathrm{div}}$ of
Lemma~\ref{asm-two-three-killing-holds-block-twists}, which applies to it. If
$e>0$, part (2) of that lemma gives $H^1_f(V_e/T_e)=B^1(V_e/T_e)$, so every
divisible cocycle in $H^1_f(V_e/T_e)$ lies in $B^1(V_e/T_e)$. Let $e=0$. If a
divisible cocycle $c$ lies in $H^1_f(V_e/T_e)$, part (3) gives $x$ with
$c(g)=gx-x$ for $g\in I_K$, and $c=0+c$ is a sum as required. Conversely, if a
divisible cocycle is $c=b+u$ with $b(g)=gy-y$ for all $g$ and $u(g)=gx-x$ for
$g\in I_K$, then $c(g)=g(x+y)-(x+y)$ for $g\in I_K$, and $c\in H^1_f(V_e/T_e)$
by part (3). Finally, let $A$ be the module of divisible cocycles $c$ with
$\ell^{c_B}c\in H^1_f(V_e/T_e)$, and send $c\in A$ to the image of its class in
$H^1_{\mathrm{div}}/H^1_{f,\mathrm{div}}$. The image is killed by $\ell^{c_B}$.
It is $0$ if and only if $c$ differs by an element of $B^1(V_e/T_e)$ from an
element of $Z_{\mathrm{div}}\cap H^1_f(V_e/T_e)$, that is, as
$B^1(V_e/T_e)\subset H^1_f(V_e/T_e)$, if and only if $c\in H^1_f(V_e/T_e)$. So
the divisible cocycles $c$ with $\ell^{c_B}c\in H^1_f(V_e/T_e)$, modulo those in
$H^1_f(V_e/T_e)$, embed in the elements of $H^1_{\mathrm{div}}/H^1_{f,\mathrm{div}}$
killed by $\ell^{c_B}$, and part (1) with $m=c_B$ bounds their length by
$([K:\mathbf Q_\ell]+1)c_B$ times the length of $\mathcal O_L/\ell\mathcal O_L$.
% Lean: tbK_of D HW F l h4 hCo hW hE hres hF. (a): LocNum.exists_up, TBLoc.eps_lat, LocRes.res_crys,
% hchi, tb_top at K' (Lemma -block-top, with g the index valuation and hidx), hcoc (the cocycle
% relation), CorD.index_smul_mem with inclusion_GK_index (Lemma -block-transfer, with
% LocRes.pow_val_smul_mem) and Submodule.smul_mem for l^(c_B - c - v). (b): TBLoc.eps_lat,
% hE = TwistK.tbe_all (Lemma -block-twists) and TBLoc.tbe_clauses (the three statements; the
% length through len_le_torsion, with bd_le_hf for B^1 in H^1_f). The block property uses bk,
% tate and euler, through Lemmas -block-top and -block-twists. The open parts are now in Lemmas
% -wach-lift (about 59,500 lines), -block-top (about 42,000) and -block-twists (about 30,000),
% which overlap; before, this open step measured about 96,900 lines for liftK_of, tbK_of and
% their inputs.
\end{proof}

\subsection{The local bound at 2 and 3}\label{subsection-asm-bound-two-three}

\begin{definition}\label{asm-two-three-crystalline-jump-sets}
Let $D$ be a definite-side datum with ranks $k,m,t$ (Definition~\ref{th-ranks})
and Harish-Chandra parameters $\lambda(\tau)$ for $\tau\in\mathcal T$
(Definition~\ref{th-definite-datum}(2)), with entries $v_a(\tau)$, $b_a(\tau)$ and
$e_a(\tau)$. Put $c_0=(t-1)/2$ and $\operatorname{refl}(q)=t-1-q$ for
$q\in\mathbf Q$. Let $1\leq j\leq m$ and $\tau\in\mathcal T$.
\begin{enumerate}
\item The \emph{cloud set} is $C_j(\tau)=\{e_j(\tau)+v_a(\tau)+c_0:1\leq a\leq
  k-1\}$, the \emph{low set} is $\Lambda_j(\tau)=\{e_j(\tau)-b_a(\tau)+c_0:
  1\leq a\leq m\}$, and the \emph{block} is $B=\{0,1,\ldots,t-1\}$.
\item Put $J^+_j(\tau)=C_j(\tau)\cup\Lambda_j(\tau)\cup B$,
  $J^-_j(\tau)=\operatorname{refl}(J^+_j(\tau))$,
  $J^{\mathrm W,+}_j(\tau)=C_j(\tau)\cup\Lambda_j(\tau)$ and
  $J^{\mathrm W,-}_j(\tau)=\operatorname{refl}(C_j(\tau)\cup\Lambda_j(\tau))$.
\item $P_j(\tau)$ is the multiset of the numbers $v_a(\tau)$ for
  $1\leq a\leq k-1$, $-b_a(\tau)$ for $1\leq a\leq m$, and
  $-e_j(\tau)+(t+1)/2-(c+1)$ for $0\leq c<t$.
\end{enumerate}
\end{definition}
\leanref{Formalism/SharpRestriction/LevelLetterCocycles/AdmissibleParameterJ.lean}{Formalism.MainNotions.AdmissibleParameterJ.cloudJ}
\leanref{Formalism/SharpRestriction/LevelLetterCocycles/AdmissibleParameterJ.lean}{Formalism.MainNotions.AdmissibleParameterJ.lowJ}
\leanref{Formalism/SharpRestriction/LevelLetterCocycles/AdmissibleParameterJ.lean}{Formalism.MainNotions.AdmissibleParameterJ.blockJ}
\leanref{Formalism/SharpRestriction/LevelLetterCocycles/AdmissibleParameterJ.lean}{Formalism.MainNotions.AdmissibleParameterJ.JJ}
\leanref{Formalism/SharpRestriction/LevelLetterCocycles/AdmissibleParameterJ.lean}{Formalism.MainNotions.AdmissibleParameterJ.JW}
\leanref{Formalism/SharpRestriction/LevelLetterCocycles/AdmissibleParameterJ.lean}{Formalism.MainNotions.AdmissibleParameterJ.PPm}

\begin{lemma}\label{asm-two-three-crystalline-jump-numerics}
Let $D$ be a definite-side datum with ranks $k,m,t,n$, with size data
(Definition~\ref{pd-size-data}) with $x$, $b$, $b'$, $w''$, exponents $b_{t'}$ at
the top classes $t'$, width $\eta$ and cloud size $S_c$, and suppose that the
size data are admissible: $\eta+t+S_c<w''$, $3w''\leq b$, $b_{t'}+b'/x\leq b$ for
every top class $t'$, and $t\leq\eta$. Let $1\leq j\leq m$ and
$\tau\in\mathcal T$, and use the notation of
Definition~\ref{asm-two-three-crystalline-jump-sets}. Then:
\begin{enumerate}
\item the elements of $P_j(\tau)$ are distinct, and for $q\in\mathbf Q$,
  $q-e_j(\tau)-c_0\in P_j(\tau)$ if and only if $q\in J^+_j(\tau)$;
\item $B\subset J^\pm_j(\tau)$, $B$ does not meet $J^{\mathrm W,\pm}_j(\tau)$, and
  $J^\pm_j(\tau)\setminus B=J^{\mathrm W,\pm}_j(\tau)$;
\item every $q\in J^\pm_j(\tau)$ satisfies $|q|\leq2b$;
\item $0\notin J^{\mathrm W,\pm}_j(\tau)$;
\item $J^{\mathrm W,\pm}_j(\tau)$ has at most $k-1+m$ elements.
\end{enumerate}
\end{lemma}
\leanref{Development/AbelianVarietyTheory/Definitions/NonGaloisDescent/JumpSetDistinctOfRealPlaceFacts.lean}{GlobalTheory.JumpSets.jFacts_of}
\leanref{Development/AbelianVarietyTheory/Definitions/NonGaloisDescent/JumpSetDistinctOfRealPlaceFacts.lean}{GlobalTheory.JumpSets.count_PPm}
\leanref{Development/AbelianVarietyTheory/Definitions/NonGaloisDescent/JumpSetDistinctOfRealPlaceFacts.lean}{GlobalTheory.JumpSets.JJ_iff_JW}
\leanref{Development/AbelianVarietyTheory/Definitions/NonGaloisDescent/JumpSetDistinctOfRealPlaceFacts.lean}{GlobalTheory.JumpSets.jumpSetCloud_of_realPlaceFacts}
\leanref{Development/AbelianVarietyTheory/Definitions/NonGaloisDescent/JumpSetDistinctOfRealPlaceFacts.lean}{GlobalTheory.JumpSets.jumpSetEta_of_realPlaceFacts}
\leanref{Development/AbelianVarietyTheory/Proofs/InertiaEigenvalues/DecompositionLocalAbove.lean}{GlobalTheory.JumpWin.exists_finset_JW}
\begin{proof}
Write $v_a$, $b_a$, $e_a$ for the entries of $\lambda(\tau)$ and $e=e_j$. Five
facts come first.
\begin{enumerate}
\item[(i)] The size data give $0\leq v_a\leq S_c$ and $|e-b_a|\leq\eta$, with
  $S_c>0$ and $\eta>0$, and $t\geq1$.
\item[(ii)] By the interlacing of the parameter, $v_1>v_2>\cdots>v_{k-1}$. By
  its condition (4), $b_1<e_1<b_2<\cdots<b_m<e_m$, with gaps at least
  $(t+3)/2$; so $e-b_a\geq e_j-b_j\geq(t+3)/2$ for $a\leq j$, and
  $b_a-e\geq b_{j+1}-e_j\geq(t+3)/2$ for $a>j$.
\item[(iii)] $w''\leq e\leq b+\eta+1$. If $\tau$ is a thick place, the windows of
  the size data give $b\leq e\leq b+\eta+1$, and $w''\leq b/3$. If
  $\tau=\tau_1(t')$ for a top class $t'$, they give
  $b_{t'}+b'/x\leq e\leq b_{t'}+b'/x+\eta$; here $b_{t'}\geq b'>0$ and
  $n\geq1$, so $w''=b'/(xn)\leq b'/x\leq e$, and $e\leq b+\eta$ by admissibility.
  Otherwise they give $w''\leq e\leq w''+\eta\leq b+\eta$.
\item[(iv)] By (iii) and admissibility, $e>\eta+t+S_c\geq2t\geq t+1$, and
  $\eta+t+S_c<b$.
\item[(v)] $b_a>0$, since $b_a\geq e-\eta>t+S_c$.
\end{enumerate}

(1) The numbers $v_a$ are distinct and nonnegative, and the numbers $-b_a$ are
distinct and negative, by (i), (ii) and (v). The numbers
$-e+(t+1)/2-(c+1)$, $0\leq c<t$, are distinct and at most $c_0-e<0$, as
$e>t$, so none of them is a $v_a$. None of them is a $-b_a$ either, since
that would give $|e-b_a|=|(t+1)/2-(c+1)|\leq c_0$, against (ii). So the
elements of $P_j(\tau)$ are distinct. Next,
$q-e-c_0=v_a$ means $q\in C_j(\tau)$, $q-e-c_0=-b_a$ means $q\in\Lambda_j(\tau)$,
and $q-e-c_0=-e+(t+1)/2-(c+1)$ means $q=t-1-c$, which runs over $B$ as $c$ does.

(2) $B\subset J^+_j(\tau)$, and $\operatorname{refl}(B)=B$ gives
$B\subset J^-_j(\tau)$. The elements of $C_j(\tau)$ are at least $e+c_0>t-1$ by
(i) and (iv). An element $q$ of $\Lambda_j(\tau)$ has $|q-c_0|=|e-b_a|\geq(t+3)/2$
by (ii), while $|q'-c_0|\leq c_0$ for $q'\in B$. So $B$ does not meet
$C_j(\tau)\cup\Lambda_j(\tau)$, nor, applying $\operatorname{refl}$, its image.
The last claim follows from the definitions.

(3) Let $q\in J^+_j(\tau)$. If $q\in C_j(\tau)$, then
$0\leq q-c_0=e+v_a\leq b+\eta+1+S_c$ by (i) and (iii); if $q\in\Lambda_j(\tau)$,
then $|q-c_0|=|e-b_a|\leq\eta$; and if $q\in B$, then $|q-c_0|\leq c_0\leq\eta$.
Since $|\operatorname{refl}(q)-c_0|=|q-c_0|$, every $q\in J^\pm_j(\tau)$ has
$|q-c_0|\leq b+S_c+\eta+1$. So
$|q|\leq b+S_c+\eta+1+c_0\leq b+S_c+\eta+t\leq2b$, using $1+c_0\leq t$ and (iv).

(4) Elements of $C_j(\tau)$ are greater than $t+1$, by (i) and (iv). An element
$e-b_a+c_0$ of $\Lambda_j(\tau)$ is at least $c_0+(t+3)/2=t+1$ if $a\leq j$, and
at most $c_0-(t+3)/2=-2$ if $a>j$, by (ii). So $C_j(\tau)\cup\Lambda_j(\tau)$
contains neither $0$ nor $t-1=\operatorname{refl}(0)$.

(5) $C_j(\tau)$ has at most $k-1$ elements and $\Lambda_j(\tau)$ at most $m$,
and $\operatorname{refl}$ is injective.
% Lean: (i)-(v) are the fields of JumpSets.JFacts, from jFacts_of hP (AdmPar L): cloud_v,
% be_close, Sc_pos, etaSz_pos and t_odd; hc_v_anti and hc_b_mono' with hc_low_pos and
% hc_low_neg (gap_be, gap_eb, hc_b_le); sz_win (win_thick, win_top with bT_ge and
% w''_le_top, win_thin); s_gt and par_b; jp_b_pos. (1) is jp_nodup and jp_mem_iff, hence
% count_PPm. (2) is block_sub_JJ, block_not_mem_JW (jp_block_not_cloud, jp_block_not_low,
% jp_refl_block) and JJ_iff_JW. (3) is jumpSetEta_of_realPlaceFacts (JumpAbsoluteBound,
% through jb_Jp_bound and jb_refl_abs, with hJ = b + Sc + eta + t). (4) is the cloud and low
% parts of jumpSetCloud_of_realPlaceFacts (jcCondition) with zero_not_mem_JW. (5) is
% exists_finset_JW.
\end{proof}

\begin{lemma}\label{asm-two-three-crystalline}
Let $D$ be a definite-side datum with local data at the twisted places, size
data with parameter choices, and twist choices at the places of its
coefficient field, and let a setup for them be given
(Definition~\ref{asm-local-bound-two-three}), with hub representations
$r=(r_{i,x})$ and $\rho=(\rho_{i,\pi})$. Let $b=\exp(x^{\varepsilon})$ be as in
the size data and put $h_2=\lfloor2b\rfloor+1$. Let $i$ be a local index such
that $\ell=\ell_i$ is neither of the level primes $g_0,g_1$, let $\pi\in Z_i$,
and let $W'=\rho_{i,\pi}\otimes\chi_i^{-1}$, with the representations
$W'_{w,L'}$ of Lemma~\ref{asm-two-three-local-invariants}. Then for every
place $w$ of $E$ over $\ell$ and every finite extension
$L'\subset\bar{\mathbf Q}_\ell$ of $L_w$:
\begin{enumerate}
\item $W'_{w,L'}$ is crystalline;
\item for every $\mathbf Q_\ell$-embedding $\tau\colon E_w\to L'$, the
  Hodge-Tate weights of $W'_{w,L'}$ along $\tau$ lie in a set of at most $r$
  integers in $[-h_2,h_2]$ and do not include $0$.
\end{enumerate}
\end{lemma}
\leanref{Development/AutomorphicArithmetic/ThetaCorrespondence/ThetaPairLifts/ThirdPlaceSeries/MainAltAltBisStatement.lean}{GlobalTheory.T.decomposition_fontaine_at_places_over_l}
\leanref{Development/AutomorphicArithmetic/ThetaCorrespondence/ThetaPairLifts/ThirdPlaceSeries/MainAltAltBis.lean}{GlobalTheory.CrysChiOne.main_alt_alt}
\leanref{Development/AutomorphicArithmetic/ThetaCorrespondence/ThetaPairLifts/ThirdPlaceSeries/RxCrys.lean}{GlobalTheory.CrysChiOne.rx_crys_alt_alt}
\leanref{Development/AbelianVarietyTheory/Proofs/InertiaEigenvalues/DecompositionLocalAway.lean}{GlobalTheory.CrysChiOne.chi1inv_crys}
\leanref{Development/AbelianVarietyTheory/Proofs/InertiaEigenvalues/DecompositionLocalAboveStatement.lean}{GlobalTheory.T.decomposition_local_above}
\leanref{Development/AbelianVarietyTheory/Proofs/InertiaEigenvalues/DecompositionLocalAbove.lean}{GlobalTheory.jumps_WpO_T2O_window}
\begin{proof}
Fix $w$ and $L'$. For a continuous homomorphism $M$ from $G_E$ to
$GL_k(\mathcal O_i)$ or to $\mathcal O_i^\times$, let $M_{w,L'}$ be the
representation of $G_{E_w}$ on $L'^k$ obtained from $M$ as $W'_{w,L'}$ is
obtained from $W'$. The hub representations are continuous
(Definition~\ref{asm-local-bound-two-three}(3)(b)), and so is $\chi_i$.
Fontaine's theory \cite{Fon94} gives four facts: a subrepresentation of a
crystalline representation is crystalline; the twist of a crystalline
representation by a crystalline character is crystalline; the dual of a
crystalline representation, twisted by the cyclotomic character
$\varepsilon_\ell$, is crystalline; and a character $\psi$ is crystalline if
and only if $\psi\varepsilon_\ell$ is.

(1) First, $(r_{i,x})_{w,L'}$ is crystalline for every Hecke point $x$. One of
the norm-choice conditions is that, since $\ell\neq g_0,g_1$, every $r_{i,x}$,
made $\bar{\mathbf Q}_\ell$-valued by the datum's embedding of $\mathcal O_i$,
is crystalline at every place over $\ell$ (Definition~\ref{gal-p-adic-hodge}).
The datum ties this notion to its comparison data: a continuous representation
over $\mathcal O_i$ is crystalline at $w$ if and only if its representation
$(\cdot)_{w,L_w}$ is crystalline. So $(r_{i,x})_{w,L_w}$ is crystalline, and
extending the coefficients from $L_w$ to $L'$ keeps it crystalline.

Second, since $\pi\in Z_i$, the canonical hub data give an isomorphism over
$K_i$ of $r_{i,\pi}$ with the block sum of $\rho_{i,\pi}$ and the characters
$\chi_i\varepsilon_\ell^{1-c}$, $1\leq c\leq t$
(Definition~\ref{asm-local-bound-two-three}(3)(b)). So $(\rho_{i,\pi})_{w,L'}$
and $(\chi_i)_{w,L'}$, the character with $c=1$, are subrepresentations of
$(r_{i,\pi})_{w,L'}$, and both are crystalline.

Third, the dual of $(\chi_i)_{w,L'}$ twisted by $\varepsilon_\ell$ is
$(\chi_i^{-1}\varepsilon_\ell)_{w,L'}$, which is therefore crystalline. So
$(\chi_i^{-1})_{w,L'}$ is crystalline.

Finally, $W'_{w,L'}$ is the twist of $(\rho_{i,\pi})_{w,L'}$ by the
crystalline character $(\chi_i^{-1})_{w,L'}$, so it is crystalline.

(2) Let $\tau\colon E_w\to L'$ be a $\mathbf Q_\ell$-embedding, let
$\tilde\tau\colon E\to\bar{\mathbf Q}_\ell$ be the composite of $E\to E_w$,
$\tau$ and $L'\subset\bar{\mathbf Q}_\ell$, and put $j=j_i$. For a continuous
representation $M$ of $G_E$ over $\mathcal O_i$ write
$\mathrm{HT}_{\tilde\tau}(M)$ for the Hodge-Tate multiset along $\tilde\tau$ of
$M$ made $\bar{\mathbf Q}_\ell$-valued (Definition~\ref{gal-p-adic-hodge}), and
$\mu_\tau(M,a)$ for the multiplicity of the integer $a$ as a Hodge-Tate weight
of $M_{w,L'}$ along $\tau$. Let $\iota\colon\bar{\mathbf Q}_\ell\to\mathbf C$ be
the datum's isomorphism, identify $\mathcal T$ with the subgroup $\Psi$ of
$G_\Phi$, and use the notation of
Definition~\ref{asm-two-three-crystalline-jump-sets}. Since the size data are
admissible (Definition~\ref{asm-local-bound-two-three}(3)(c)), the norm-choice
conditions include the following three.
\begin{enumerate}
\item[(R)] \emph{Reading.} The embedding $\tilde\tau$ induces the place $w$,
  and $\iota\circ\tilde\tau$ is the complex embedding of $E$ attached to an
  element $g\in G_\Phi$. Let $x$ be a Hecke point such that $r_{i,x}$ and
  $\chi_i$ are de Rham at $w$ and $\chi_i$ has the single Hodge-Tate weight $c$
  along $\tilde\tau$. Then for every integer $a$, $\mu_\tau(r_{i,x}\otimes
  \chi_i^{-1},a)$ is the multiplicity of $a+c$ in $\mathrm{HT}_{\tilde\tau}(r_{i,x})$.
  If moreover $x\in Z_i$ and $r_{i,x}$ is the semisimple block sum of a
  representation $\rho'$ and the characters $\chi_i\varepsilon_\ell^{1-c'}$,
  $1\leq c'\leq t$, then $\mu_\tau(\rho'\otimes\chi_i^{-1},a)=
  \mu_\tau(r_{i,x}\otimes\chi_i^{-1},a)-1$ for $0\leq a<t$, the two
  multiplicities are equal for the other $a$, and $a$ is a Hodge-Tate weight of
  $(\rho'\otimes\chi_i^{-1})_{w,L'}$ along $\tau$ if and only if its multiplicity
  is not $0$.
\item[(H)] \emph{Weights of the hub.} Since $\ell\neq g_0,g_1$, every $r_{i,x}$
  is crystalline at $w$ and $\mathrm{HT}_{\tilde\tau}(r_{i,x})$ is the multiset
  $\mathrm{HT}^{\mathrm{aut}}_i(\tilde\tau)$ of integers that the datum attaches
  to $i$ and $\tilde\tau$ through its infinitesimal characters; and a continuous
  representation is crystalline at $w$ if and only if it is de Rham at $w$ with
  unramified Weil-Deligne representation.
\item[(D)] \emph{Dictionary.} There is a sign $\epsilon\in\{\pm1\}$, independent
  of $\tilde\tau$, with the following property. Let $g\in G_\Phi$ be such that
  $\iota\circ\tilde\tau$ is attached to $g$, let $\rho\in\mathcal T$ with
  $g=\rho$ or $g=\rho\iota$, and let $\varsigma=1$ if $g\in\Psi$ is equivalent to
  $\epsilon=1$, and $\varsigma=-1$ otherwise. Then there is an integer
  $c_{\tilde\tau}$ with
  $\mathrm{HT}^{\mathrm{aut}}_i(\tilde\tau)=\varsigma P_j(\rho)+c_{\tilde\tau}$
  as multisets, and $\chi_i$ is de Rham at $w$ with the single Hodge-Tate weight
  $\varsigma(-e_j(\rho))-c_0+c_{\tilde\tau}$ along $\tilde\tau$.
\end{enumerate}
Let $g$ be as in (R). If $g\in\Psi$ put $\rho=g$; otherwise $g\iota\in\Psi$,
since $\iota$ has second sign $-1$, and put $\rho=g\iota$. Let $\varsigma$ and
$c_{\tilde\tau}$ be as in (D), and put $c=\varsigma(-e_j(\rho))-c_0+c_{\tilde\tau}$,
the Hodge-Tate weight of $\chi_i$ along $\tilde\tau$. By (H), $r_{i,\pi}$ is
crystalline at $w$, hence de Rham there, and
$\mathrm{HT}_{\tilde\tau}(r_{i,\pi})=\varsigma P_j(\rho)+c_{\tilde\tau}$. As in
(1), $\pi\in Z_i$ and $r_{i,\pi}$ is the semisimple block sum of
$\rho_{i,\pi}$ and the characters $\chi_i\varepsilon_\ell^{1-c'}$. So (R), with
$x=\pi$ and $\rho'=\rho_{i,\pi}$, gives for every integer $a$
\[
 \mu_\tau(\rho_{i,\pi}\otimes\chi_i^{-1},a)=
 \#\{p\in P_j(\rho):\varsigma p+c_{\tilde\tau}=a+c\}-\delta_a,
\]
counted with multiplicity, where $\delta_a=1$ if $0\leq a<t$ and $\delta_a=0$
otherwise. Since $\varsigma^2=1$ and $2c_0=t-1$, the condition
$\varsigma p+c_{\tilde\tau}=a+c$ says $p=a-e_j(\rho)-c_0$ if $\varsigma=1$, and
$p=(t-1-a)-e_j(\rho)-c_0$ if $\varsigma=-1$. By
Lemma~\ref{asm-two-three-crystalline-jump-numerics}(1), the count is therefore
$1$ if $a\in J^{\varsigma}_j(\rho)$ and $0$ otherwise, where $J^{\varsigma}$
stands for $J^+$ or $J^-$. By
Lemma~\ref{asm-two-three-crystalline-jump-numerics}(2), the integers
$0\leq a<t$ lie in $J^{\varsigma}_j(\rho)$ and not in $J^{\mathrm W,\varsigma}_j(\rho)$, and the other
integers lie in $J^{\varsigma}_j(\rho)$ exactly when they lie in
$J=J^{\mathrm W,\varsigma}_j(\rho)$. Hence $\mu_\tau(\rho_{i,\pi}\otimes
\chi_i^{-1},a)$ is $1$ if $a\in J$ and $0$ otherwise. With the last part of (R)
this gives three properties of the set $J$ of rational numbers:
\begin{enumerate}
\item[(a)] the Hodge-Tate weights of $W'_{w,L'}$ along $\tau$ are the integers
  in $J$;
\item[(b)] every $q\in J$ has $|q|\leq2b$, by
  Lemma~\ref{asm-two-three-crystalline-jump-numerics}(3), as $J\subset
  J^{\varsigma}_j(\rho)$;
\item[(c)] $0\notin J$, by Lemma~\ref{asm-two-three-crystalline-jump-numerics}(4).
\end{enumerate}
% Lean: decomposition_local_above = jumps_WpO_T2O_window takes (a) from
% hodgeTateJumpsAtModulesForSign D ep (ConclusionJa), (b) from JumpAbsoluteBound with
% JW_sub_JJ and natAbs_le_of_mem_JJ, and (c) from jcCondition with zero_not_mem_JW;
% these are fields of JumpConditions L = JumpStatement.holds_all_alt D L hAJ
% (= JumpSets.stmtJ_holds_alt), with hAJ = S.facts.admJ L S.par: the norm-choice
% field admJ gives AdmJAlt L (DeRhamFamilyAwayAlt = G2a, CrysIffWD,
% DeRhamJumpsDistinctCondition, AdmDict, AdmRead, AdmPar) for admissible L. ConclusionJa
% is jumpLemma_a_of_admissible_alt_alt: (R) is AdmRead (InducesW, the g, and the
% multiplicity identities for rTwO and WpO, the latter with SemisimpleDecomposition),
% (H) is G2a (g2a', l != g_0, g_1) with CrysIffWD, (D) is AdmDict with exists_placeOf
% (D.eTau : Tau = Psi, cc = iota) and up = decide (TypeU D ep g); the count is count_key
% (count_PPm), and block_sub_JJ, block_mem_iff, block_not_mem_JW and JJ_iff_JW give the
% passage from JJ to JW. The Lean also treats l = g_0, g_1 (DeRhamJumpsDistinctCondition),
% the twist rTwO and the duals rDuO and T2O, which this lemma does not use.
So $J$ has at most $k-1+m=r$ elements, by
Lemma~\ref{asm-two-three-crystalline-jump-numerics}(5). Let $U$ be the set of integers in $J$
that lie in $[-h_2,h_2]$, a set of at most $r$ integers. By (a) and (b), every
Hodge-Tate weight $a$ of $W'_{w,L'}$ along $\tau$ lies in $J$ and satisfies
$|a|\leq\lfloor2b\rfloor<h_2$, so it lies in $U$; by (c), $0$ is not one of
them.
\end{proof}

\begin{lemma}\label{asm-two-three-wild-powers}
Let $E$ be a number field, $\Omega$ an algebraic closure of $E$ and
$G_E=\operatorname{Gal}(\Omega/E)$. Let $v$ be a finite place of $E$ with
residue characteristic $p$, let $A$ be a place of $\Omega$ over $v$, with
absolute value $|\cdot|_A$, and let the \emph{wild inertia subgroup} at $A$ be
the set $P_A$ of $\sigma\in G_E$ with $|\sigma(y)/y-1|_A<1$ for every nonzero
$y\in\Omega$. Let $\sigma\in P_A$ and let $n\geq1$ be an integer prime to $p$.
\begin{enumerate}
\item If $H$ is an open normal subgroup of $G_E$ with $\sigma^n\in H$, then
  $\sigma\in H$.
\item Let $R$ be a noetherian local ring with maximal ideal $\mathfrak m$,
  topologized by the powers of $\mathfrak m$, and let
  $\chi\colon G_E\to R^\times$ be a homomorphism that is continuous as a map to
  $R$. If $\chi(\sigma)^n=1$, then $\chi(\sigma)=1$.
\end{enumerate}
\end{lemma}
\leanref{Development/AbelianVarietyTheory/Proofs/FoundationProofs/ValuationPowApplyDivSubOneLt.lean}{GlobalTheory.mem_of_wild_of_pow_mem_open_normal}
\leanref{Development/AbelianVarietyTheory/Proofs/InertiaEigenvalues/TameDescentConditionsEveryLocalisation.lean}{GlobalTheory.TameDesign.char_eq_one_of_wild}
\begin{proof}
(1) Since $p\nmid n$, $|n|_A=1$. For $z\in\Omega$ put
$s(z)=\sum_{k=0}^{n-1}\sigma^k(z)$. Each $\sigma^k$ lies in $P_A$, so for
$z\neq0$ the element $s(z)/z-n=\sum_{k<n}(\sigma^k(z)/z-1)$ has absolute value
less than $1$, and therefore $s(z)\neq0$. Let $y\in\Omega$ be fixed by
$\sigma^n$. Then $\sigma(s(y))=s(y)$, so for $z=ny-s(y)$ we get
$\sigma^k(z)=n\sigma^k(y)-s(y)$ and $s(z)=ns(y)-ns(y)=0$. Hence $z=0$, that is
$s(y)=ny$, and applying $\sigma$ gives $n\sigma(y)=s(y)=ny$, so $\sigma(y)=y$.
Every element of the fixed field of $H$ is fixed by $\sigma^n\in H$, hence by
$\sigma$. Since $H$ is open, it is closed, so it is the group of all elements
fixing its fixed field, and $\sigma\in H$.

(2) For $k\geq1$ the kernel $H_k$ of the composite of $\chi$ with
$R^\times\to(R/\mathfrak m^k)^\times$ is an open normal subgroup of $G_E$,
since $\chi$ is continuous and $\mathfrak m^k$ is open. It contains
$\sigma^n$, so it contains $\sigma$ by (1); that is,
$\chi(\sigma)-1\in\mathfrak m^k$. By Krull's intersection theorem the
$\mathfrak m^k$ meet in $0$, so $\chi(\sigma)=1$.
\end{proof}

\begin{lemma}\label{asm-two-three-wild}
Let $D$ be a definite-side datum with local data at the twisted places
(Definition~\ref{asm-two-three-local-data}) and twisting data
(Definition~\ref{th-twisting-data}). Let $i$ be a local index, let $v$ be a
place of $E$ that does not lie over $\ell_i$, let $I_v$ be the inertia group at
$v$ that contains $\gamma_{i,v}$, at a place $A$ of $\Omega$ over $v$, and let
$P_v\subset I_v$ be the wild inertia subgroup at $A$
(Lemma~\ref{asm-two-three-wild-powers}). Then $\chi_i(\sigma)=1$ for every
$\sigma\in P_v$.
\end{lemma}
\leanref{Development/AbelianVarietyTheory/Proofs/InertiaEigenvalues/TameDescentConditionsEveryLocalisation.lean}{GlobalTheory.TameDesign.firstLocalChar_wild_eq_one_away_from_l}
\begin{proof}
Let $\sigma\in P_v$, let $p$ be the residue characteristic of $v$ and put
$\ell=\ell_i$, so that $p\neq\ell$. The twisting data attach to $i$ a twisting
index $i'$ with $\ell_{i'}=\ell$, a sign $b$, a ring homomorphism
$\iota\colon\mathcal O^0_{i'}\to\mathcal O_i$ and a homomorphism
$\psi_i\colon G_E\to\mathcal O_i^\times$ with open kernel, each of whose values
has order a power of $\ell$, such that
$\chi_i=(\iota\circ\chi_{i',j_i,b})\cdot\psi_i$, where $j_i$ is the lift of
$i$. Put $\chi'=\chi_{i',j_i,b}$, a continuous character with values in the
complete discrete valuation ring $\mathcal O^0_{i'}$. Since $\sigma\in P_v$,
$\sigma$ lies in $I_v$.

First, $\psi_i(\sigma)=1$. Let $a\geq0$ with $\psi_i(\sigma)^{\ell^a}=1$. Then
$\sigma^{\ell^a}$ lies in the open normal subgroup $\ker\psi_i$, and $\ell^a$
is prime to $p$, so $\sigma\in\ker\psi_i$ by
Lemma~\ref{asm-two-three-wild-powers}(1).

Second, $\chi'(\sigma)=1$. There are three cases.
\begin{itemize}
\item If $v\in P_D(i')$, the twisting data give a character
  $\delta\colon G_E\to(\mathcal O^0_{i'})^\times$ that maps $I_v$ onto a group
  of order $15$ and agrees with $\chi'$ or with $\chi'^{-1}$ on $I_v$,
  according to $b$, and they require $Nv\equiv1\pmod{15}$. So
  $\delta(\sigma)^{15}=1$ and $\chi'(\sigma)^{15}=1$. Since $Nv$ is a power of
  $p$ and $Nv\equiv1\pmod{15}$, $p\nmid15$, and
  Lemma~\ref{asm-two-three-wild-powers}(2) gives $\chi'(\sigma)=1$.
\item If $v\notin P_D(i')$ lies over a level prime $g_k$, then $p=g_k$ is
  prime to $g_k-1$, and $\chi'(\sigma)^{g_k-1}=1$, since $\chi'$ is tamely
  ramified of order dividing $g_k-1$ at the places over $g_k$ not over $\ell$.
  Lemma~\ref{asm-two-three-wild-powers}(2) gives $\chi'(\sigma)=1$.
\item Otherwise $v$ is not in $P_D(i')$ and lies over none of $\ell$, $g_0$
  and $g_1$. The twisting data make $\chi'$ unramified at such places, so
  $\chi'(\sigma)=1$.
\end{itemize}
Hence $\chi_i(\sigma)=\iota(\chi'(\sigma))\psi_i(\sigma)=1$.
\end{proof}

\begin{lemma}\label{asm-two-three-rows-partition}
Let $D$ be a definite-side datum with local data at the twisted places
(Definition~\ref{asm-two-three-local-data}), host field $E$ with auxiliary
prime $d$ and element $\sqrt{-d}\in E$
(Definition~\ref{hf-framed-host-field}), and coefficient field $L\supset E$ with
its embedding $\iota_L\colon L\to\mathbf C$
(Definition~\ref{th-definite-datum}(5)). Let $i$ be a local index, put
$\ell=\ell_i$, and suppose that $\ell=2$, or that $\ell$ is odd and $-d$ is a
nonzero square modulo $\ell$. For a place $w$ of $E$ over $\ell$, with
comparison data $E_w$ and $L_w$
(Definition~\ref{asm-two-three-local-killing}(1)), let
$\lambda_w\colon L\to L_w$ be the composite of the embedding $L\to K_i$ of the
local data and the isomorphism $K_i\to L_w$ of the comparison data, and view
$E$ in $E_w$ through the comparison data. For a sign $\epsilon\in\{\pm1\}$, a
$\mathbf Q_\ell$-embedding $\tau_0\colon E_w\to L_w$ has \emph{type $\epsilon$}
if $\tau_0(E)\subset\lambda_w(L)$ and there is a field embedding
$\varphi\colon E\to\mathbf C$ with $\varphi(\sqrt{-d})=\epsilon\,i\sqrt d$ such
that $\iota_L(y)=\varphi(x)$ whenever $x\in E$, $y\in L$ and
$\lambda_w(y)=\tau_0(x)$. Fix $\epsilon$, let $P_{\mathrm{col}}$ be the set of
places of $E$ over $\ell$ at which every $\mathbf Q_\ell$-embedding
$E_w\to L_w$ has type $\epsilon$, and let $P_{\mathrm{row}}$ be the set of the
other places over $\ell$. Then:
\begin{enumerate}
\item every $\mathbf Q_\ell$-embedding $E_w\to L_w$ has exactly one type; at
  each place $w$ over $\ell$ there is at least one such embedding, and all of
  them have the same type;
\item if $w$ and $w'$ are places over $\ell$ that are complex conjugate, the
  embeddings at $w$ and those at $w'$ have opposite types;
\item a place over $\ell$ lies in $P_{\mathrm{row}}$ if and only if its complex
  conjugate lies in $P_{\mathrm{col}}$, and no embedding at the complex
  conjugate of a place of $P_{\mathrm{col}}$ has type $\epsilon$.
\end{enumerate}
\end{lemma}
\leanref{Development/AutomorphicArithmetic/ResidueCharacters/PlaceTypeConj.lean}{UpperTheory.N.decl_split}
\leanref{Development/AutomorphicArithmetic/ResidueCharacters/PlaceTypeConj.lean}{UpperTheory.N.place_type_conj}
\leanref{Development/AutomorphicArithmetic/ResidueCharacters/TwoDecl.lean}{UpperTheory.N.Two.decl_split}
\leanref{Development/AutomorphicArithmetic/ResidueCharacters/TwoRoot.lean}{UpperTheory.N.Two.place_type_conj}
\leanref{Development/AbelianVarietyTheory/Counterexamples/NonGaloisSurvey/SplitRowIffOfIsEmpty.lean}{GlobalTheory.SplitRowProbe.embIn_xor}
\leanref{Development/AbelianVarietyTheory/Proofs/InertiaEigenvalues/SplitPrimePlaceCmTypeConjugate.lean}{GlobalTheory.SplitType.root_of_embIn}
\leanref{Development/AbelianVarietyTheory/Proofs/InertiaEigenvalues/SplitPrimePlaceCmTypeConjugate.lean}{GlobalTheory.SplitType.root_ne}
\leanref{Development/AutomorphicArithmetic/ResidueCharacters/TwoRoot.lean}{UpperTheory.N.Two.root_ne}
\begin{proof}
Let $|\cdot|$ be the absolute value of $\bar{\mathbf Q}_\ell$ with
$|\ell|=\ell^{-1}$. The embedding of $E$ in $E_w$ induces $w$, so an integer
$x$ of $E$ lies in $w$ if and only if its image in $E_w$ has absolute value
less than $1$.

(1) Let $\tau_0\colon E_w\to L_w$ be a $\mathbf Q_\ell$-embedding. Since $E$ is
Galois over $\mathbf Q$, the restriction of $\tau_0$ to $E$ and the restriction
of $\lambda_w$ to $E\subset L$ differ by an automorphism of $E$: there is
$g\in G_\Phi$ with $\tau_0(x)=\lambda_w(\sigma_g(x))$ for $x\in E$, where
$\sigma_g=\gamma^{-1}(g)$. So $\tau_0(E)\subset\lambda_w(L)$, and since
$\lambda_w$ is injective, $\lambda_w(y)=\tau_0(x)$ holds exactly when
$y=\sigma_g(x)$. Hence $\tau_0$ has type $\epsilon$ if and only if
$\varphi=\iota_L\circ\sigma_g$ satisfies $\varphi(\sqrt{-d})=\epsilon\,i\sqrt d$.
As $\varphi(\sqrt{-d})^2=-d$, exactly one sign has this property. There is at
least one embedding $E_w\to L_w$, since $L_w$ contains $\tau(E_w)$ for every
embedding $\tau$ of $E_w$ in $\bar{\mathbf Q}_\ell$. The last claim of (1) is
proved below.

Next, a square root of $-d$ in $\mathbf Z_\ell$. If $\ell$ is odd, $-d$ is a
nonzero square modulo $\ell$ and $\ell\nmid2d$, so Hensel's lemma gives
$s\in\mathbf Z_\ell$ with $s^2=-d$ and $|2s|=1$. If $\ell=2$, then
$d\equiv7\pmod{16}$ (Definition~\ref{hf-framed-host-field}(1)), so $-d\equiv1
\pmod8$ and Hensel's lemma gives $s\in\mathbf Z_2$ with $s^2=-d$. The image of
$\sqrt{-d}$ under $L\to K_i$ is a square root of $-d$, so it is
$s_0\in\{s,-s\}$, and $s_0\neq-s_0$ as $d\neq0$. For a place $w$ over $\ell$
the image $t_w$ of $\sqrt{-d}$ in $E_w$ is a square root of $-d$ in
$\mathbf Q_\ell$, so $t_w\in\{s_0,-s_0\}$.

Let $\tau_0$ be an embedding at $w$ of type $\epsilon'$, with $\varphi$ as in
the definition. The datum's embedding $E\to L\to\mathbf C$ is the embedding
$\varphi_0$ of the host, which maps $\sqrt{-d}$ to $i\sqrt d$
(Definition~\ref{hf-host}). If $\epsilon'=1$, the element $y\in L$ with
$\lambda_w(y)=\tau_0(\sqrt{-d})$ satisfies $\iota_L(y)=i\sqrt d$, so $y$ is the
image of $\sqrt{-d}$ in $L$; if $\epsilon'=-1$, it is its negative. As $\tau_0$
fixes $t_w\in\mathbf Q_\ell$ and $\lambda_w$ maps the image of $\sqrt{-d}$ to
$s_0$, this gives $t_w=\epsilon's_0$. So every embedding at $w$ has type $1$ if
$t_w=s_0$ and type $-1$ if $t_w=-s_0$, which completes (1).

(2) Let $w$ and $w'$ be complex conjugate places over $\ell$; by (1) it
suffices to show $t_w\neq t_{w'}$. Suppose that $t_w=t_{w'}=t$. If $\ell$ is
odd, choose an integer $n$ with $|t-n|<1$. Then $\sqrt{-d}-n\in w$, so its
conjugate $-\sqrt{-d}-n$ lies in $w'$, that is, $|-t-n|<1$. Then
$|2t|=|(t-n)-(-t-n)|<1$, a contradiction. If $\ell=2$, write $d+1=8k$. The
element $\omega=(1+\sqrt{-d})/2$ is integral, as $\omega^2-\omega+2k=0$, and its
conjugate is $1-\omega$. Since $(1+t)(1-t)=8k$, one of $1\pm t$ has absolute
value less than $1$, and then so does the other, as they differ by $2t$. So
$u=(1+t)/2\in\mathbf Z_2$, and $u(1-u)=2k$. The image of $\omega$ in $E_w$ is
$u$, and the image of $1-\omega$ in $E_{w'}$ is $1-u$. As $w$ and $w'$ are
conjugate, $\omega\in w$ if and only if $1-\omega\in w'$, that is, $|u|<1$ if
and only if $|1-u|<1$. Since $|u|\,|1-u|=|2k|<1$, one of them, and hence both,
are less than $1$, and then $|1|=|u+(1-u)|<1$, a contradiction.

(3) Every place over $\ell$ has a complex conjugate, which lies over $\ell$. By
(1) and (2), exactly one place of each conjugate pair lies in
$P_{\mathrm{col}}$. So $w\in P_{\mathrm{row}}$ if and only if its conjugate lies
in $P_{\mathrm{col}}$. If $u\in P_{\mathrm{col}}$, the embeddings at its
conjugate have type $-\epsilon$ by (2), so none has type $\epsilon$.
% Lean: for l = 3, decl_split = (decl D i hl ep, hypDecl, ambTypeD) with columns_xor and
% place_type_conj; for l = 2, Two.decl_split with Two.place_type_conj. Type epsilon is
% EmbTypeUe ep (Br.EmbTypeU for ep = true, EmbTypeUc for ep = false), equivalent to
% EmbIn ep (embTypeUe_iff), with halfType true = cmType = {phi | phi(sqd) = I sqrt d}.
% (1): HTW.exists_letter_of_local_embedding (exists_aut_comp_eq_of_normal), embIn_xor
% (halfType_xor), nonempty_emb (LL_split). The roots: SplitType.exists_sqrt (Hensel,
% Misc.isSquare_of_hensel) or Two.exists_sqrt, exists_root, root_of_embIn (with
% D.embE_phi0 and host.phi0_mem). (2): SplitType.root_ne (mem_iff_norm, bridge_place) or
% Two.root_ne (Two.exists_half). (3): GlobalTheory.SplitPlaces.exists_conjPl,
% SplitCount.mem_placesOver_of_conj, not_both_columns, mem_col and mem_row.
\end{proof}

\begin{lemma}\label{asm-two-three-rows-split-bound-lifting}
Let $D$ be a definite-side datum with local data at the twisted places
(Definition~\ref{asm-two-three-local-data}) and host field $E$, let $i$ be a
local index, put $\ell=\ell_i$, and let $N\geq1$ be an integer such that for
every place $w$ of $E$ over $\ell$, with comparison data $E_w$ and $L_w$
(Definition~\ref{asm-two-three-local-killing}(1)), and every finite extension
$L'\subset\bar{\mathbf Q}_\ell$ of $L_w$, the pair $(E_w,L')$ is
$N$-cyclotomic. Suppose that the lifting property holds at $\ell$, and use the
notation of Definition~\ref{asm-two-three-killing-holds-lift-block}. Let $w$ be
a place of $E$ over $\ell$, $L'\subset\bar{\mathbf Q}_\ell$ a finite extension
of $L_w$, and $h\geq0$ and $k\geq1$ integers. Let $(M_x)_{x\in X}$ be a family
of crystalline representations of $G_{E_w}$ on $L'^k$, with $g\mapsto M_x(g)v$
continuous for every $v$, whose Hodge-Tate weights along every $\mathbf
Q_\ell$-embedding $E_w\to L'$ lie in $[-h,h]$. Call a lattice representation
$(V,T)$ of $G_{E_w}$ over $L'$ a \emph{subquotient of sums} of the $M_x$ if
there are $x_1,\ldots,x_m\in X$, a subspace $U$ of $(L'^k)^m$ stable under
$G_{E_w}$ acting through $M_{x_a}$ on the $a$-th factor, and a surjective
$G_{E_w}$-equivariant $L'$-linear map $U\to V$. Let $(V_1,T_1)$ and
$(V_2,T_2)$ be subquotients of sums of the $M_x$, and let $f\colon V_1\to V_2$
be a continuous $G_{E_w}$-equivariant $L'$-linear map with $f(T_1)=T_2$. Then
for every $c\in Z^1_f(T_2)$ there are $c_1\in Z^1_f(T_1)$ and $b\in B^1(T_2)$
with $\ell^{e_\ell(N,h)}c=f\circ c_1+b$.
\end{lemma}
\leanref{Development/AutomorphicArithmetic/ResidueCharacters/LiftExponentAdmissibleColsRowsOfCrystalline.lean}{UpperTheory.N.liftAt_of_liftK}
\leanref{Formalism/SharpRestriction/LevelLetterCocycles/JointConditionsZpowers/BaseFont.lean}{Formalism.MainNotions.N.subquot_crys}
\begin{proof}
By Fontaine's theory \cite{Fon94}, a finite direct sum of crystalline
representations is crystalline, and so are its subrepresentations and its
quotients; along each embedding, the Hodge-Tate weights of a direct sum are
those of its summands, and those of a subrepresentation or a quotient are among
those of the representation. So $V_1$ and $V_2$ are crystalline, and the
Hodge-Tate weights of $V_1$ along every $\mathbf Q_\ell$-embedding $E_w\to L'$
lie in $[-h,h]$. The lifting property at $\ell$, applied to $N$, $h$, the
$N$-cyclotomic pair $(E_w,L')$, $(V_1,T_1)$, $(V_2,T_2)$ and $f$, gives the
lemma.
% Lean: liftAt_of_liftK (LiftAt with exponent eL l N h, for IsSubquotOfSums), with subquot_crys
% (cj_sum, cj_sub and cj_image from FontaineAt, the datum's Fontaine properties at w, read as
% Fontaine's theory) and hK at (Kw, C.LL) with hC. About 420 lines with
% lift_exponent_admissible_cols_rows_of_crystalline, which applies it at the column places to
% the r_{i,x} tensor chi_i^(-1) and at the row places to their duals twisted by chi_i eps_l.
\end{proof}

\begin{lemma}\label{asm-two-three-rows-split-bound-block}
Let $D$, $E$, $i$, $\ell$ and $N$ be as in
Lemma~\ref{asm-two-three-rows-split-bound-lifting}, let $r,t,n$ be the ranks of
$D$, and suppose that the block property holds at $\ell$
(Definition~\ref{asm-two-three-killing-holds-lift-block}). Let
$R\colon G_E\to GL_{r+t}(\mathcal O_i)$ be a homomorphism, and for a place $w$
of $E$ over $\ell$ and a finite extension $L'$ of $L_w$ let $R'_{w,L'}$ be the
representation of $G_{E_w}$ on $L'^{r+t}$ obtained from $R\otimes\chi_i^{-1}$
as in Lemma~\ref{asm-two-three-local-invariants}. Let $h\geq t$ be an integer,
and suppose that for every such $w$ and $L'$ and every $\mathbf
Q_\ell$-embedding $\tau\colon E_w\to L'$, the Hodge-Tate weights of
$R'_{w,L'}$ along $\tau$ lie in a set of at most $r+t$ integers in $[-h,h-1]$.
Then for every place $w$ of $E$ over $\ell$ and every finite extension $L'$ of
$L_w$, the conclusions (a) and (b) of the block property hold for the pair
$(E_w,L')$, with $J_\tau$ the set of Hodge-Tate weights of $R'_{w,L'}$ along
$\tau$ and with $c_B=c^B_\ell(N,h,4n(r+t),t)$.
\end{lemma}
\leanref{Development/AutomorphicArithmetic/ResidueCharacters/BlockConditionAtPlacesOfJumpWindow.lean}{UpperTheory.N.block_condition_at_places_of_jump_window}
\leanref{Development/AutomorphicArithmetic/ResidueCharacters/ExistsUnion.lean}{UpperTheory.N.exists_union}
\begin{proof}
Fix $w$ and $L'$. As in the second step of the proof of
Lemma~\ref{asm-two-three-local-invariants}, there are at most
$[E_w:\mathbf Q_\ell]\leq4n$ embeddings $\tau\colon E_w\to L'$. For each $\tau$
choose a set $U_\tau$ as in the hypothesis, and let $U$ be the union of the
$U_\tau$. It has at most $4n(r+t)$ elements, all in $[-h,h]$, and it contains
every $J_\tau$. The block property at $\ell$, applied to $N$, $h$,
$4n(r+t)$, $t\leq h$, the $N$-cyclotomic pair $(E_w,L')$, the family
$(J_\tau)$ and $U$, gives (a) and (b) with $c_B$.
% Lean: block_condition_at_places_of_jump_window (the conclusion admTB is TBAt at every place
% over l and every good coefficient extension; TBAt is (a) and (b) for fixed J and cB), with
% HubLoc.card_emb and bridge_finrank_le, exists_union, and BridgeInertia.inertEq for the
% inertia subgroup Br.inert. The window hypothesis is BlockWindow.JumpWindow, which the Lean
% takes from hodgeTateJumpWindow_of_jumpSets at h = floor(2b). About 180 lines.
\end{proof}

\begin{lemma}\label{asm-two-three-rows-split-bound-invariants}
Let $D$, $E$, $i$, $\ell$, $N$ and $r,t,n$ be as in
Lemma~\ref{asm-two-three-rows-split-bound-block}, and suppose that the
invariant-killing property holds at $\ell$
(Definition~\ref{asm-two-three-local-killing}(3)). Let $\rho\colon G_E\to
GL_r(\mathcal O_i)$ be a homomorphism such that $W'=\rho\otimes\chi_i^{-1}$ is
continuous, let $T_2$ be the representation $g\mapsto\varepsilon_\ell(g)\,
{}^{\mathrm t}W'(g^{-1})$, and for a place $w$ of $E$ over $\ell$ and a finite
extension $L'$ of $L_w$ let $T_{2,w,L'}$ be obtained from $T_2$ as in
Lemma~\ref{asm-two-three-local-invariants}. Let $h\geq0$ be an integer, and
suppose that for every such $w$ and $L'$, the representation $T_{2,w,L'}$ is
crystalline and, for every $\mathbf Q_\ell$-embedding $\tau\colon E_w\to L'$,
its Hodge-Tate weights along $\tau$ lie in a set of at most $r$ integers in
$[-h,h]$ and do not include $0$. Then for every place $w$ of $E$ over $\ell$
and every finite extension $L'$ of $L_w$, the $G_{E_w}$-invariants of
$L'^r/\mathcal O_{L'}^r$, with the action through $T_{2,w,L'}$, form an
$\mathcal O_{L'}$-module of length at most $r\,c_\ell(N,h,4n(r+t))$ times the
length of $\mathcal O_{L'}/\ell\mathcal O_{L'}$.
\end{lemma}
\leanref{Development/AutomorphicArithmetic/ResidueCharacters/InvariantsLengthAdmissibleOfNonzeroJumps.lean}{UpperTheory.N.invariants_length_admissible_of_nonzero_jumps}
\begin{proof}
Fix $w$ and $L'$, and put $c=c_\ell(N,h,4n(r+t))$. The first four steps of the
proof of Lemma~\ref{asm-two-three-local-invariants}, with $w$, $L'$ and
$T_{2,w,L'}$ in place of $u$, $L_u$ and $W'_u$, show that $\ell^c$ kills every
$G_{E_w}$-invariant element of $L'^r/\mathcal O_{L'}^r$. So the invariants lie
in $(\ell^{-c}\mathcal O_{L'}/\mathcal O_{L'})^r$, whose length is $rc$ times
the length of $\mathcal O_{L'}/\ell\mathcal O_{L'}$.
% Lean: invariants_length_admissible_of_nonzero_jumps (H0ColumnLengthBound: the Lean states it
% for w conjugate to a column place, a good coefficient extension and the standard lattice
% (IsStd); the proof uses none of the first two) = kill_of with JumpsNoZero for T2O, then
% EmbeddingTypes.len_le_of_kill. T2O = twistM (dualM WpO) epsU. About 170 lines.
\end{proof}

\begin{lemma}\label{asm-two-three-rows-split-bound}
Assume the third collection of published results in
Theorem~\ref{lit-bundle-unitary}(6). Let $D$ be a definite-side datum with
local data at the twisted places, size data with parameter choices, and twist
choices at the places of its coefficient field, and let a setup for them be
given (Definition~\ref{asm-local-bound-two-three}), with hub representations
$\rho=(\rho_{i,\pi})$. Let $r,t,n$ be the ranks of $D$, let $E$ be its host
field, with integer $Q_S$, auxiliary primes $q,d$ and level primes $g_0,g_1$,
and let $b$ and $h_B$ be as in Lemma~\ref{asm-two-three-cost}. Let $i$ be a
local index, put $\ell=\ell_i$, and let $N\geq1$ be an integer. Suppose that:
\begin{enumerate}
\item $\ell\neq g_0$ and $\ell\neq g_1$;
\item $\ell$ divides $Q_S$, or $\ell$ does not divide $Q_S$, $\ell\neq q,d$ and
  $-d$ is a nonzero square modulo $\ell$;
\item if $\ell$ does not divide $Q_S$, every place of $E$ over $\ell$ is
  unramified over $\mathbf Q$;
\item for every place $w$ of $E$ over $\ell$, with comparison data $E_w$ and
  $L_w$ (Definition~\ref{asm-two-three-local-killing}(1)), and every finite
  extension $L'\subset\bar{\mathbf Q}_\ell$ of $L_w$, the pair $(E_w,L')$ is
  $N$-cyclotomic;
\item for each sign $\epsilon$, with the types of embeddings of
  Lemma~\ref{asm-two-three-rows-partition}, let $P_\epsilon$ be the set of
  places of $E$ over $\ell$ at which every $\mathbf Q_\ell$-embedding
  $E_w\to L_w$ has type $\epsilon$; then $P_\epsilon$ and the set of complex
  conjugates of its places are disjoint and together contain every place of $E$
  over $\ell$, and no embedding at the complex conjugate of a place of
  $P_\epsilon$ has type $\epsilon$;
\item the lifting property, the invariant-killing property and the block
  property hold at $\ell$ (Definitions~\ref{asm-two-three-local-killing}(3)
  and~\ref{asm-two-three-killing-holds-lift-block}).
\end{enumerate}
Put $n'=4n(r+t)$, $e'=e_\ell(N,\lfloor h_B\rfloor)$,
$c_0=c_\ell(N,\lfloor2b\rfloor+1,n')$ and
$c_B=c^B_\ell(N,\lfloor2b\rfloor,n',t)$
(Definition~\ref{asm-two-three-killing-holds-lift-block}). Then there is a
partition of the places of $E$ over $\ell$ into sets $P_{\mathrm{row}}$ and
$P_{\mathrm{col}}$, with $P_{\mathrm{row}}$ the set of complex conjugates of the
places of $P_{\mathrm{col}}$, such that for every eligible
$\pi\in Z^{\mathrm{el}}_i$, with $h^0_u(\pi)$ and $B_{\mathrm{split}}$ as in
Lemma~\ref{asm-two-three-split-numerics} for this partition:
\begin{enumerate}
\item[(a)] if $\operatorname{Sel}_i(\rho_{i,\pi})$ has finite length, the
  groups $\text{Ш}^1_\ell(E,k)$ with $0<k<t$ have finite length, and
  $\zeta_{i,v}\neq1$ for every $v\in Q_i$, then
  $c_i(\pi)\le B_{\mathrm{split}}(e',c_0,c_B)$;
\item[(b)] $h^0_u(\pi)\le rc_0$ for every $u\in P_{\mathrm{col}}$;
\item[(c)] $\zeta_{i,v}\neq1$ for every $v\in Q_i$.
\end{enumerate}
\end{lemma}
\leanref{Development/AutomorphicArithmetic/ThetaCorrespondence/ThetaPairLifts/ThirdPlaceSeries/ConclusionSplitOfHub.lean}{UpperTheory.N.conclusion_split_alt}
\leanref{Development/AbelianVarietyTheory/Proofs/InertiaEigenvalues/SplitCaseBoundAtPointOfAdmittedItemsStatement.lean}{GlobalTheory.T.splitCase_bound_at_point_of_admitted_items_alt}
\begin{proof}
Let $\epsilon$ be a sign as in condition (D) of the proof of
Lemma~\ref{asm-two-three-crystalline}, put $P_{\mathrm{col}}=P_\epsilon$, and
let $P_{\mathrm{row}}$ be the set of complex conjugates of its places. By
hypothesis (5) these two sets are disjoint and together contain every place of
$E$ over $\ell$. Fix an eligible $\pi\in Z^{\mathrm{el}}_i$, and put
$h_2=\lfloor2b\rfloor+1$, so that $c_0=c_\ell(N,h_2,n')$.

(a) Alpöge's split-case bound gives the inequality of (a) at $\pi$ from a list
of conditions on the datum, the hub representations, the partition and $\pi$.
Three of these conditions are where hypothesis (6) enters through the lifting
and block properties and the invariant-killing property. The first is the
lifting statement with exponent $e'$ at the column places, for the
subquotients of sums of the representations $r_{i,x}\otimes\chi_i^{-1}$,
$x\in X_i$, and at the row places, for those of their duals twisted by
$\chi_i\varepsilon_\ell$; given that these representations are crystalline with
Hodge-Tate weights in $[-\lfloor h_B\rfloor,\lfloor h_B\rfloor]$, it is
Lemma~\ref{asm-two-three-rows-split-bound-lifting} with $h=\lfloor h_B\rfloor$.
The second is the bound $r\,c_0$ on the length of the invariants of the dual of
$W'=\rho_{i,\pi}\otimes\chi_i^{-1}$ twisted by $\varepsilon_\ell$ at the complex
conjugates of the column places; given crystallinity and a window
for the Hodge-Tate weights, it is
Lemma~\ref{asm-two-three-rows-split-bound-invariants} with $h=h_2$. The third
is the pair of conclusions of the block property at the row places, for the
Hodge-Tate weights of $r_{i,\pi}\otimes\chi_i^{-1}$, with constant $c_B$; given
a window for these weights, it is
Lemma~\ref{asm-two-three-rows-split-bound-block} with $h=\lfloor2b\rfloor$.
\openstep{The rest of the proof of (a) is not yet written out. It consists of
Alpöge's split-case bound, which rests on global Artin reciprocity, local Tate
duality and the local Euler characteristic formula, with the list of conditions
at $\pi$ that it needs; and of the verification at $\pi$ of the conditions
other than the three above and of the crystallinity and windows those three
take as given: conditions on the hub representations $r_{i,x}$ and
$\rho_{i,\pi}$ away from $\ell$ and on their reductions; crystallinity and
Hodge-Tate weights at the places over $\ell$ of the twists of $r_{i,x}$, of $W'$
and of its dual twisted by $\varepsilon_\ell$; and local Bloch-Kato, duality and
finiteness statements at the places over $\ell$.}
% Measured again, following imports and leaving out what the book states (the closure count
% used for the other open steps of this section): conclusion_split_alt about 81,300 lines in
% 1,187 files; splitCase_bound_at_point_of_admitted_items_alt about 35,300; the condition
% sources set_pt_of_setup_not_twisted about 48,200, admG5_of_setup about 38,400,
% blochKatoAnnihilator_all_of_tateDuality_eulerChar about 22,500, closingFamilyX_alt_alt about
% 21,800, ambT2_of_deRham_family about 21,400, admG2_of_setup_jumps about 19,100 and
% unramifiedTrivialH1f_of_decl about 17,200 (these overlap heavily);
% admissibleAtGoodPoints_of_coefRing_taken_at_live about 4,600,
% charpoly_span_and_hom_relations_of_admissible about 1,200,
% tateFinitenessAll_of_localTateDuality about 1,100, amb_of_setPt about 430 and
% hodgeTateJumpWindow_of_jumpSets about 70. The three conditions written out above are
% lift_exponent_admissible_cols_rows_of_crystalline (about 420 lines; its windows hwinT, hwinD
% come from the jump conditions hA and he, its crystallinity hT, hD from admG2_of_setup_jumps and
% decomposition_fontaine_at_places_over_l), invariants_length_admissible_of_nonzero_jumps (about
% 170; JumpsNoZero for T2O from decomposition_fontaine_at_places_over_l and
% decomposition_local_above, whose T2O parts Lemma asm-two-three-crystalline does not state) and
% block_condition_at_places_of_jump_window (about 180; JumpWindow from
% hodgeTateJumpWindow_of_jumpSets). The conditions of EmbeddingTypes.admitted are 36 in all
% (admG1 to admG6, TameGeneratorConditions, ColumnH0Vanishing, adm51, CharpolyCoeffsLemmaFamily,
% SpanQuotientLemma, admBK, LiftExponentColumn, LiftExponentDualPlaces, H0ColumnLengthBound,
% admTB, UnramifiedTrivialH1f, FamilyContinuous, admSS, SemisimpleDecomposition,
% CyclotomicInertiaImage, LocalCoefficientCompatibility, CyclotomicCompletionAdmissible,
% FontaineAtPlacesOverL, admH2, tateH1FinitenessAtPlacesOver, UnramifiedClassFieldBound,
% ambBlock, ambDual, ambT2, ambW, ambFil, H1fFiniteRankTorsion, ambSub, ambCont, ambTypeD),
% and the split-case bound also takes TruncatedSeriesCountingIdentities.admG2 (GainBoundAt).
% Several are long (admG5 and adm51 through IdempotentChoiceData and GoodCongruenceObjects,
% SpanQuotientLemma through VanishingLemma, admBK through bkBAt); stating them is the next step.
% Lean: the third part of the conclusion aux_3 of
% GlobalTheory.T.splitCase_bound_at_point_of_admitted_items_alt Pr.artin Pr.tate Pr.euler
% (= splitCase_bound_at_point_of_admitted_items: jump_and_length_of_admitted,
% globalLatticeIndexTheorem_of_lengthBoundAtPlaces = SelmerBoundAtPoint.main =
% K.Closed.theorem_n5_from_inputs, length_bound_of_theorem), applied in
% UpperTheory.N.conclusion_split_alt to EmbeddingTypes.admitted (35 conditions), which it
% assembles from set_pt_of_setup_not_twisted, admG2_of_setup_jumps, admG5_of_setup,
% decomposition_fontaine_at_places_over_l, decomposition_local_above,
% lift_exponent_admissible_cols_rows_of_crystalline,
% invariants_length_admissible_of_nonzero_jumps, hodgeTateJumpWindow_of_jumpSets,
% block_condition_at_places_of_jump_window, charpoly_span_and_hom_relations_of_admissible,
% ambT2_of_deRham_family, closingFamilyX_alt_alt (font, fin),
% admissibleAtGoodPoints_of_coefRing_taken_at_live (Chebotarev),
% blochKatoAnnihilator_all_of_tateDuality_eulerChar, tateFinitenessAll_of_localTateDuality,
% unramifiedTrivialH1f_of_decl and amb_of_setPt, with the partition dc of hdecl ep (ep the
% sign of the jump conditions, JumpStatement.holds_all_alt) and the hypothesis hK (hypothesis
% (6) of the lemma; LiftK enters through lift_exponent_admissible_cols_rows_of_crystalline,
% TBK through block_condition_at_places_of_jump_window). Size: proof closure of conclusion_split_alt without the parts already in
% the book (the two decomposition statements, the jump conditions and
% length_invariants_Wprime_divisible_le_of_nonzero_jumps), 3,536 theorems, about 60,400
% lines in 890 files; of these, splitCase_bound_at_point_of_admitted_items_alt alone is
% 1,718 theorems, about 32,600 lines in 428 files.

(b) Put $W'=\rho_{i,\pi}\otimes\chi_i^{-1}$. The hub representation
$\rho_{i,\pi}$ is continuous (Definition~\ref{asm-local-bound-two-three}(3)(b)),
and so are $\chi_i$ and $W'$. Every eligible point lies in $Z_i$, and $\ell$ is
neither of the level primes by hypothesis (1). So
Lemma~\ref{asm-two-three-crystalline} shows that for every place $w$ of $E$
over $\ell$ and every finite extension $L'$ of $L_w$, the representation
$W'_{w,L'}$ is crystalline, with Hodge-Tate weights along every
$\mathbf Q_\ell$-embedding of $E_w$ in $L'$ in a set of at most $r$ integers in
$[-h_2,h_2]$, none of them $0$. The pairs $(E_w,L')$ are $N$-cyclotomic by
hypothesis (4), and the invariant-killing property holds at $\ell$ by
hypothesis (6). So Lemma~\ref{asm-two-three-local-invariants}, with
$\rho=\rho_{i,\pi}$ and $h=h_2$, gives $h^0_u(\pi)\le r\,c_\ell(N,h_2,n')=rc_0$
for every place $u$ of $E$ over $\ell$, and in particular for every
$u\in P_{\mathrm{col}}$.
% Lean: length_invariants_Wprime_divisible_le_of_nonzero_jumps i N hK.2.1 hC (X.ρ i) π
% (hOf L) (AmbSub.WpO_cont from hfam.cont i π) hW u, with u over l by dc.cover; hW :
% JumpsNoZero for WpO from decomposition_fontaine_at_places_over_l (crystalline, with
% hg0, hg1 from hL) and decomposition_local_above (window and no zero jump) at π ∈ Z i
% (LargeLevelLemmas.zgoodSubsetZ).

(c) Let $v\in Q_i$. Then $v$ does not lie over $\ell$
(Definition~\ref{asm-two-three-local-data}(4)), so $\chi_i$ is trivial on the
wild inertia at $v$ by Lemma~\ref{asm-two-three-wild}, and
Lemma~\ref{asm-two-three-local-zeta} gives $\zeta_{i,v}\neq1$.
% Lean: (g5x v (Or.inr hv)).2.2.1 hv, with g5x the TameGeneratorConditions field of
% set_pt_of_setup_not_twisted (= SetPtOff.main_alt_alt), which is (TameDesign.main D i).2 =
% galois_condition_five_extra_of D i hc, with hc.1 v from
% TameDesign.firstLocalChar_wild_eq_one_away_from_l and D.nd.SQ_l.
\end{proof}

\begin{lemma}\label{asm-two-three-rows-cyclotomic}
Let $D$ be a definite-side datum with local data at the twisted places
(Definition~\ref{asm-two-three-local-data}), host field $E$ with integer $Q_S$
and auxiliary primes $q,d$, and $N=Q_Sqd$. Let $i$ be a local index, put
$\ell=\ell_i$, and suppose that $4\mid N$ if $\ell=2$. Then for every place $w$
of $E$ over $\ell$, with comparison data $E_w$ and $L_w$
(Definition~\ref{asm-two-three-local-killing}(1)), and every finite extension
$L'\subset\bar{\mathbf Q}_\ell$ of $L_w$, the pair $(E_w,L')$ is
$N$-cyclotomic.
\end{lemma}
\leanref{Development/AutomorphicArithmetic/ResidueCharacters/ExistsUnion.lean}{UpperTheory.N.cycAt_of}
\begin{proof}
If $\ell=2$ then $4\mid N$ by hypothesis, and $N\geq1$. The host field embeds
in $\mathbf Q(\zeta_N)$ (Definition~\ref{hf-framed-host-field}(3)), so the
image of $E$ in $\bar{\mathbf Q}_\ell$ lies in $\mathbf Q_\ell(\mu_N^\ell)$,
since $N$ divides $\operatorname{lcm}(\ell,N)$. This field is
finite-dimensional over $\mathbf Q_\ell$, hence closed, and the image of $E$ is
dense in $E_w$, so $E_w\subset\mathbf Q_\ell(\mu_N^\ell)$. The comparison data
give $\mu_N^\ell\subset L_w\subset L'$.
% Lean: cycAt_of: lvl from the hypothesis, pos from condN_ne_zero, le from
% KStar.Cyc.kw_le_roots (le_roots_of_dense and the host's image_mem_of_roots), coef from
% CoefExt.isGood (self_good and C.le).
\end{proof}

\begin{proposition}\label{asm-two-three-rows}
Assume the third collection of published results in
Theorem~\ref{lit-bundle-unitary}(6): Fontaine's theory of $B_{\mathrm{cris}}$,
global Artin reciprocity, the local results of Bloch and Kato, local Tate
duality, finiteness and the Euler characteristic over finite extensions of
$\mathbf Q_\ell$, and the integral form of the theorem of Bhatt and Scholze.
Let $D$ be a definite-side datum with local data at the twisted places,
size data with parameter choices, and twist choices at the places of its
coefficient field, and let a setup for them be given
(Definition~\ref{asm-local-bound-two-three}), with hub representations
$\rho=(\rho_{i,\pi})$. Let $\kappa_{A,A'}$ be as in
Lemma~\ref{asm-two-three-cost}. Let $i$ be a local index with $\ell_i=2$, or
with $\ell_i=3$ and $3\nmid Q_S$. Then for every eligible
$\pi\in Z^{\mathrm{el}}_i$,
\[
 c_i(\pi)\le\operatorname{sel}_i(\rho_{i,\pi})+\kappa_{6000,30000}(\ell_i)
 +V(\ell_i)+v_{\ell_i}(h_E)+\ell_i^{-2}.
\]
\end{proposition}
\leanref{Development/AutomorphicArithmetic/ThetaCorrespondence/ThetaPairLifts/ThirdPlaceSeries/NRowTwo.lean}{UpperTheory.N.onAt_eq_two_bounds_of_setup}
\leanref{Development/AutomorphicArithmetic/ThetaCorrespondence/ThetaPairLifts/ThirdPlaceSeries/NRowThree.lean}{UpperTheory.N.rowThreeAt_of_setup_alt}
\leanref{Development/AutomorphicArithmetic/ResidueCharacters/RowThreeC.lean}{UpperTheory.N.RowOnAt}
\begin{proof}
Write $\ell=\ell_i$, let $q,d$ be the auxiliary primes of the host field and
$N=Q_Sqd$, and let $b$, $h_B$ and $s_\ell$ be as in
Lemma~\ref{asm-two-three-cost}.

\emph{Step 1: the prime $\ell$.} The level primes satisfy $g_k\equiv1\pmod N$,
so neither divides $Q_S$, and neither is $3$, since $N>3$. Let $\ell=2$. The
host field has $4\mid Q_S$ (Definition~\ref{hf-framed-host-field}), so $\ell$
is not a level prime, $4\mid N$, and $s_0(i)=v_2(Q_S)\ge2$; since $q$ and $d$
are odd, $s_0(i)=v_2(N)=s_2$. Let $\ell=3$ with $3\nmid Q_S$. Then $3\ne q$
(Definition~\ref{hf-host-field-rows}(3)), $3\ne d$ since $d\equiv7\pmod{16}$, and $3$
is not a level prime, so $3\nmid N$. The setup gives $B\ge60$, so
$\left(\frac{-d}{3}\right)=1$ by Definition~\ref{hf-host-field-rows}(4); the
places of $E$ over $3$ are unramified over $\mathbf Q$, since $E$ embeds in
$\mathbf Q(\zeta_N)$; and $s_0(i)=\max(1,v_3(Q_S))=1=s_3$. In both cases $\ell$
splits in $\mathbf Q(\sqrt{-d})\subset E$ (for $\ell=2$ because
$-d\equiv1\pmod8$), so complex conjugation pairs the places of $E$ over $\ell$
without fixed points.

\emph{Step 2: the split-case bound.} Put $n'=4n(r+t)$,
$h_1=\lfloor h_B\rfloor$, $h_2=\lfloor2b\rfloor+1$ and $h_3=\lfloor2b\rfloor$.
Lemma~\ref{asm-two-three-rows-split-bound} applies to $i$ and $N$. Its
hypotheses (1) to (3) hold by Step 1, and its hypothesis (4) by
Lemma~\ref{asm-two-three-rows-cyclotomic}, since $4\mid N$ if $\ell=2$ by
Step 1. For its hypothesis (5): by Step 1, $\ell=2$ or $-d$ is a nonzero
square modulo $\ell$, so Lemma~\ref{asm-two-three-rows-partition} applies to
both signs, and its part (3) gives (5). Its hypothesis (6) is
Proposition~\ref{asm-two-three-killing-holds}. The lemma gives a partition
$(P_{\mathrm{row}},P_{\mathrm{col}})$ of the places over $\ell$, whose row
places are the complex conjugates of its column places, with its parts (a),
(b) and (c) at every eligible $\pi$. Since $s_\ell=\max(1,v_\ell(N))$, we have
$\lambda'_\ell(h)=\lambda_\ell(N,h)$, so the numbers $e'$, $c_0$ and $c_B$ of
the two lemmas agree. So parts (b), (c) and (a) are hypotheses (1), (3) and
(4) of Lemma~\ref{asm-two-three-split-numerics} for this partition and
$h_1,h_2,h_3$ at every eligible $\pi$.
% Lean: Step 2 is UpperTheory.N.conclusion_split_alt with N = host.condN,
% cycAt_of (Lemma asm-two-three-rows-cyclotomic), Two.decl_split (l = 2) or decl_split
% (l = 3) and locAll (LiftK, InvariantsKilledOverCyclotomic, TBK from the fields bk, tate,
% euler of nLink). Its three parts hB, hW and hz at pi are hypotheses (4), (1) and (3) of
% bound_smallPrime_of_splitBound_cyclotomic (Step 4). The partition is the dc of
% conclusion_split_alt, the dc of hdecl ep (Lemma asm-two-three-rows-partition) with ep the
% sign of ConclusionJa, which is the sign of AdmDict.

\emph{Step 3: invariants of the cyclotomic twists.} Let $u$ be a place over
$\ell$ and $0<k<t$. The local data of $D$ give an element $g$ of an inertia
group at $u$ with $\varepsilon_\ell(g)=1+\ell^{s_0(i)}$. The invariants of $g$
in $K_i/\mathcal O_i$, with the action through $\varepsilon_\ell^{-k}$, are
killed by $(1+\ell^{s_0(i)})^{-k}-1$, whose valuation is $s_0(i)+v_\ell(k)$,
because $s_0(i)\ge1$, and $s_0(i)\ge2$ if $\ell=2$. So the invariants of a
decomposition group at $u$ have normalized length at most
$s_0(i)+v_\ell(k)\le s_0(i)+\lfloor\log_\ell t\rfloor$, which is hypothesis (2)
of Lemma~\ref{asm-two-three-split-numerics}.
% Lean: h0_le through lenUnit_invariants_le and EpsAll.exists_generator (epsAll).

\emph{Step 4: the bound.} Since $1\le b\le h_B$, we have $h_1\le h_B$,
$h_2\le2b+1\le4h_B$ and $h_3\le2b\le4h_B$. The setup gives the killing
condition for $\rho$ and the finiteness of the groups $\text{Ш}^1_p(E,k)$, and
$s_0(i)=s_\ell$ by Step 1. Lemma~\ref{asm-two-three-split-numerics} applies to
the partition of Step 2 and gives the bound.
\end{proof}

\begin{proposition}\label{asm-local-bound-two-three-holds}
Assume the third collection of published results in
Theorem~\ref{lit-bundle-unitary}(6).
\begin{enumerate}
\item Every family satisfies the local bound at the twisted places above $2$
  and $3$ (Definition~\ref{asm-local-bound-two-three}).
\item Suppose that maps satisfying the local Langlands laws exist for every
  number field and every finite place (Definition~\ref{gal-llc-laws}). Then
  for all parameters $p$ of the construction in the proof of
  Theorem~\ref{asm-definite-side-exists}, among them the parameters of
  dependence $c$ for every real $c$, the property that the local bound at the
  twisted places above $2$ and $3$ holds eventually along $\mathcal D_3$ is
  linked at $p$.
\end{enumerate}
\end{proposition}
\leanref{Development/AutomorphicArithmetic/ThetaCorrespondence/ThetaPairLifts/ThirdPlaceSeries/At23OfAll.lean}{UpperTheory.N.at23_family_alt}
\leanref{Development/AutomorphicArithmetic/ThetaCorrespondence/ThetaPairLifts/ThirdPlaceSeries/FamilyBoundAtTwoThreeOfLocalBounds.lean}{UpperTheory.N.boundAtTwoThree_of_local_bounds_at_two_three}
\leanref{Development/AutomorphicArithmetic/ThetaCorrespondence/ThetaPairLifts/ThirdPlaceSeries/FamilyBoundAtTwoThreeOfLocalBounds.lean}{UpperTheory.N.mid_of_rowOn_alt}
\leanref{Development/AutomorphicArithmetic/ThetaCorrespondence/ThetaPairLifts/ThirdPlaceSeries/At23OfAll.lean}{UpperTheory.N.boundAtTwoThree_n_linked}
\leanref{Development/AutomorphicArithmetic/ThetaCorrespondence/ThetaPairLifts/ThirdPlaceSeries/At23OfAll.lean}{UpperTheory.N.boundAtTwoThree_linked_of_recT}
\begin{proof}
(1) Let $\mathcal F$ be a family with exponent $\varepsilon$ and threshold
$x_1$, and let $\kappa_{A,A'}$ be as in Lemma~\ref{asm-two-three-cost}. By
Proposition~\ref{asm-two-three-rows}, at every member at every scale $x>x_1$,
with the datum, size data, parameter choices and twist choices that
$\mathcal F$ gives it, and for every setup for them, every local index $i$ with
$\ell_i=2$, or with $\ell_i=3$ and $3\nmid Q_S$, and every eligible $\pi$
satisfy
\[
 c_i(\pi)\le\operatorname{sel}_i(\rho_{i,\pi})+\kappa_{6000,30000}(\ell_i)
 +V(\ell_i)+v_{\ell_i}(h_E)+\ell_i^{-2}.
\]
The sizes of the family (Definition~\ref{pd-family}(4)) give
$n\le x^{2+8\varepsilon+\varepsilon/100}$ and $t\le x^{C_t+\varepsilon/100}$
at every member at every scale $x>x_1$, for a constant $C_t$ of the family.
So there is $C_0\ge1$ with $n\le x^{C_0}$ and $t\le x^{C_0}$ at all these
members, since $x>x_1>1$. Put
\[
 C=18\cdot6000\,C_0+2\cdot30000\,(3C_0+9)+5C_0+1
\]
and $x_0=3$. Let $x\ge x_0$ with $x>x_1$, and take a member at $x$ with a
setup. Then $\log x\ge1$. For $\ell=2$ we have $v_2(N)=v_2(Q_S)$, since $q$ and
$d$ are odd; for $\ell=3$ we have $v_3(N)=v_3(Q_S)$, since $3\ne q,d$
(Step 1 of the proof of Proposition~\ref{asm-two-three-rows}).
Lemma~\ref{asm-two-three-cost} with $A=6000$ and $A'=30000$ gives
$\kappa_{6000,30000}(\ell_i)\le C\,f\,m\,b\log x+C\,x^{C}x^{2\varepsilon}$
for both primes, and the bound at $2$ and $3$ with constant $C$ follows.

(2) By (1) the local bound holds for every family. So for every class
$\mathcal A$ it holds eventually along $\mathcal D_3$, with
$\varepsilon_1=1$ and any largeness threshold (Definition~\ref{pd-family}).
This holds for every class and every selection, hence along $(Q_p,\Lambda_p)$
with the trivial condition on bundles (Definition~\ref{is-slot-rows}); that
is, it is linked at $p$.
% Lean: in (1) the constants come from exists_nonneg_twoC_constants_of_prints and
% exists_nonneg_threeC_constants_of_prints (A = 6000, A' = 30000, threshold 0); the
% polynomial bounds are GlobalTheory.T.admissible_parameters_family_const (fields
% sizes: SizeN, SizeT); the valuations are Num.factorization_condN and
% factorization_condN'. The Lean takes the larger of the two constants, which
% coincide here, through midF_mono. (2) is LinkedAt.of_row, Row.of_sel and
% closed_of_all.
\end{proof}

\begin{theorem}\label{asm-every-character-theta}
Assume the published results on unitary groups
(Theorem~\ref{lit-bundle-unitary}) and on the local theta correspondence
(Theorem~\ref{lit-bundle-theta}), the agreement of global and local theta
integrals and the local matching of models of the Weil representation
(Definition~\ref{th-theta-hypotheses}), the covolume bound of
Theorem~\ref{cv-covolume-bound}, the existence, for every
number field and every finite place, of a map satisfying the local Langlands
laws (Definition~\ref{gal-llc-laws}), and the weakening of Vandiver's
conjecture (Definition~\ref{sha-vandiver-weakening}). Then for every Dirichlet
character $\chi$ modulo $N\ge1$ and every $s$ with $\operatorname{Re}s>1/2$,
$L(s,\chi)\ne0$ unless $\chi$ is trivial and $s=1$.
\end{theorem}
\begin{proof}
The published results on unitary groups contain Green's parametrization in
the form of Theorem~\ref{lit-green-cuspidal}, which gives the labelled form by
Corollary~\ref{lit-green-labelled}. They also contain the third collection of
Theorem~\ref{lit-bundle-unitary}(6). By
Proposition~\ref{asm-local-bound-two-three-holds}, the local bound at the
twisted places above $2$ and $3$ (Definition~\ref{asm-local-bound-two-three})
holds for every family, and for every real $c$ the property that it holds
eventually along $\mathcal D_3$ is linked at the parameters of dependence $c$.
This is the hypothesis on the local bound in
Theorem~\ref{asm-definite-side-exists}.
% Lean: boundAtTwoThree_linked_of_recT applied to the field nLink of P, through
% at23_family_alt; the bound is BoundAtTwoThreeFamilyFonAlt.
Theorem~\ref{asm-definite-side-exists} gives a collection of definite-side
hypotheses, and Theorem~\ref{asm-every-character-green} concludes.
\end{proof}
\leanref{Development/AutomorphicArithmetic/ThetaCorrespondence/ThetaPairLifts/ThirdPlaceSeries/SiegelEisensteinSecondTermHypotheses.lean}{UpperTheory.Lit.every_character_of_theta_integrals_local_matching}

\subsection{Theta bases and matched quadruples}\label{subsection-asm-theta-bases}

\begin{definition}\label{asm-theta-agree-doubling}
Let $E$ be a CM field, Galois over $\mathbf Q$, with maximal totally real
subfield $F$, complex conjugation $c$ and quadratic character $\eta$ of
$\mathbf A_F^\times$, and let $(\psi,(\tau_w)_w,P,(\tau_v)_{v\in P})$ satisfy
the conditions on a good base for $E$ (Definition~\ref{gal-good-base}) other
than the condition that $E$ contain $\sqrt{-d}$. Write
$\mathbf A_E=\mathbf A_F\otimes_FE$, and for $x\in\mathbf A_E^\times$ put
$|x|=|N_{E/F}(x)|_{\mathbf A_F}$. A \emph{Hecke character} of $E$ is a
continuous homomorphism $\mathbf A_E^\times\to\mathbf C^\times$ with values of
absolute value $1$ that is trivial on $E^\times$.
\begin{enumerate}
\item A \emph{lattice space} of rank $n\ge1$ consists of a matrix
  $J\in M_n(E)$, an $\mathcal O_E$-lattice $L\subset E^n$, a matrix
  $\theta_w\in M_n(\mathbf C)$ for each infinite place $w$ of $F$, and a
  matrix $b_v\in GL_n(F_v)$ for each finite place $v$ of $F$. Put
  $h(x,y)=x^{\mathsf T}J\,c(y)$. A \emph{Cartan involution} of
  $M\in M_n(\mathbf C)$ is a matrix $\theta$ with $\theta^2=1$,
  $\theta^{\mathsf T}M\bar\theta=M$ and
  $\operatorname{Re}\bigl((\theta x)^{\mathsf T}M\bar x\bigr)>0$ for $x\neq0$.
  The lattice space is \emph{Hermitian of signature $(r,s)$} if
  $c(J)^{\mathsf T}=J$; $L=\{y\in E^n:h(y,x)\in\mathcal O_E\text{ for all }x\in L\}$;
  for every $w$ the matrix $M_w=\tau_w(J)$ has signature $(r,s)$ and
  $\theta_w$ is a Cartan involution of $M_w$; and for $v\in P$,
  $b_v\mathcal O_v^n$ is the completion of $L$ at the place of $E$ given by
  $\tau_v$. It is \emph{skew-Hermitian of signature $(r,s)$} if the same
  holds with $c(J)^{\mathsf T}=-J$, with $M_w=i\,\tau_w(J)$, and with
  $L=\{y:\operatorname{Tr}_{E/\mathbf Q}(o\,h(y,x))\in\mathbf Z\text{ for all }x\in L,\ o\in\mathcal O_E\}$.
  For a lattice space $A$ of rank $m$ with matrix $J_A$, let $U_A$ be the
  group of $g\in GL_m(\mathbf A_E)$ with $g^{\mathsf T}J_Ac(g)=J_A$, and
  $U_A(F)$ its subgroup of matrices over $E$.
\item Let $Y$ be a lattice space of rank $n$ with matrix $J$, and put
  $J^\square=\operatorname{diag}(J,-J)$ and
  $\Delta=\left(\begin{smallmatrix}1_n\\1_n\end{smallmatrix}\right)$. The
  \emph{doubled group} $G_Y$ is the group of $x\in GL_{2n}(\mathbf A_E)$ with
  $x^{\mathsf T}J^\square c(x)=J^\square$, and $G_Y(F)$ is its subgroup of
  matrices over $E$. The \emph{Siegel parabolic subgroup} $P_Y$ consists of
  the $p\in G_Y$ with $p\Delta=\Delta t(p)$ for some
  $t(p)\in GL_n(\mathbf A_E)$. The \emph{height} of
  $x\in GL_{2n}(\mathbf A_E)$ is
  $\operatorname{ht}(x)=\prod_v\max(1,\|x_v\|,\|x_v^{-1}\|)$, the product
  over the places $v$ of $F$, where $\|\cdot\|$ is the largest operator norm
  under the complex embeddings of $E$ over $v$ if $v$ is infinite and the
  largest absolute value of an entry if $v$ is finite.
\item A \emph{frame} of $Y$ is a matrix $b\in GL_{2n}(E)$ with
  $b^{\mathsf T}J^\square c(b)=\left(\begin{smallmatrix}0&1\\\epsilon&0\end{smallmatrix}\right)$
  for some $\epsilon=\pm1$, whose first $n$ columns are $\Delta t$ for some
  $t\in GL_n(E)$. Its \emph{standard compact subgroup} $K_b$ is the set of
  $k\in G_Y$ such that $\tau_w(b^{-1}kb)$ is unitary for every infinite place
  $w$, and the finite components of $b^{-1}kb$ and $b^{-1}k^{-1}b$ have
  entries in $\mathcal O_E\otimes\widehat{\mathbf Z}$.
\item Let $K\subset G_Y$ be a subgroup. A \emph{height function} for $K$ is a
  function $a\colon G_Y\to\mathbf R$ such that every element of $G_Y$ is $pk$
  with $p\in P_Y$ and $k\in K$, and $a(pk)=|\det t(p)|$ for all $p\in P_Y$
  and $k\in K$. Let $a$ be one, let $\chi$ be a Hecke character of $E$ and
  let $s_1\in\mathbf C$. A \emph{section} of degree $s_1$ for $(K,\chi)$ is a
  continuous function $f$ on $G_Y$ with
  $f(px)=\chi(\det t(p))\,|\det t(p)|^{s_1+n/2}f(x)$ for $p\in P_Y$, whose
  right translates under $K$ span a finite-dimensional space. Its \emph{flat
  family} is $f_s(x)=a(x)^{s-s_1}f(x)$, $s\in\mathbf C$, and the
  \emph{Eisenstein sum} of a function $f'$ on $G_Y$ at $x$ is the sum of
  $f'(\gamma x)$ over $\gamma\in(P_Y\cap G_Y(F))\backslash G_Y(F)$.
\item A \emph{function of moderate growth} on $G_Y$ is a finite linear
  combination of continuous functions $f$ with $|f|\le\operatorname{ht}^N$
  for some $N$. A \emph{growth functional} is a linear form $\mu$ on the space
  of functions of moderate growth such that for every $N$ there is $C$ with
  $|\mu(f)|\le C$ whenever $|f|\le\operatorname{ht}^N$.
\item Let $A$ and $Y$ be lattice spaces of ranks $m$ and $n$, with matrices
  $J_A$ and $J$, and let $\mathbb W=M_{m\times2n}(\mathbf A_E)$ carry the
  symplectic $\mathbf A_F$-form
  $\langle x,y\rangle=-\operatorname{Tr}_{\mathbf A_E/\mathbf A_F}\operatorname{tr}\bigl(x^{\mathsf T}J_A\,c(y)\,(J^\square)^{\mathsf T}\bigr)$.
  Let $\chi$ be a Hecke character of $E$. A \emph{doubling datum} for $(A,Y)$
  with character $\chi$ consists of a finite set $I$, an
  $\mathbf A_F$-linear isomorphism $d=(d_1,d_2)$ from $\mathbb W$ to
  $\mathbf A_F^I\times\mathbf A_F^I$, a representation $\omega$ of $G_Y$ and a
  representation $S$ of $U_A$ on the functions on $\mathbf A_F^I$, and a
  linear form $l$ on these functions, such that the following hold. Let
  $\mathcal S$ be the Schwartz-Bruhat space of $\mathbf A_F^I$; call a linear
  form \emph{admissible} if for each fixed finite component it is continuous
  in the archimedean Schwartz component; and put
  $(\rho(w)\phi)(x)=\psi\bigl(x\cdot d_2(w)+\tfrac12d_1(w)\cdot d_2(w)\bigr)\phi(x+d_1(w))$.
  \begin{enumerate}
  \item[(a)] $\langle w,w'\rangle=d_1(w)\cdot d_2(w')-d_1(w')\cdot d_2(w)$,
    and $d$ maps $M_{m\times2n}(E)$ into $F^I\times F^I$.
  \item[(b)] The operators $\omega(x)$ preserve $\mathcal S$ and the
    admissible forms, $\omega(x)\rho(w)=\rho(wx^{\mathsf T})\omega(x)$ on
    $\mathcal S$, they preserve the $L^2$-norm of every element of
    $\mathcal S$ for every Haar measure, and $x\mapsto\lambda(\omega(x)\phi)$
    is continuous for admissible $\lambda$ and $\phi\in\mathcal S$. For every
    admissible $\lambda$ with $\lambda(\rho(z\Delta^{\mathsf T})\phi)=\lambda(\phi)$
    for all $z\in M_{m\times n}(\mathbf A_E)$ and $\phi\in\mathcal S$,
    \[
    \lambda(\omega(p)\phi)=\chi(\det t(p))\,|\det t(p)|^{m/2}\lambda(\phi)
    \qquad(p\in P_Y,\ \phi\in\mathcal S).
    \]
  \item[(c)] The operators $S(g)$ satisfy the conditions of (b) before the
    displayed rule, with $S(g)\rho(w)=\rho(gw)S(g)$.
  \item[(d)] $l$ is admissible, and for $\Phi\in\mathcal S$, $z\in M_{m\times n}(\mathbf A_E)$
    and $g\in U_A$: $l(\rho(z\Delta^{\mathsf T})\Phi)=l(\Phi)$,
    $l(S(g)\Phi)=l(\Phi)$, and the sum of $l(\rho((y,-y))\Phi)$ over
    $y\in M_{m\times n}(E)$ converges to the \emph{theta sum}
    $\theta(\Phi)=\sum_{\xi\in F^I}\Phi(\xi)$.
  \end{enumerate}
\end{enumerate}
\end{definition}
\leanref{Development/AbelianVarietyTheory/BlockTheory/DivisorBase/WeilTensorModels.lean}{GlobalTheory.IndefiniteRowSideFullHecke.WeilTensorModels.Base.IsGood}
\leanref{Development/AbelianVarietyTheory/BlockTheory/DivisorBase/WeilTensorModels.lean}{GlobalTheory.IndefiniteRowSideFullHecke.WeilTensorModels.Spc.IsHerm}
\leanref{Development/AbelianVarietyTheory/BlockTheory/DivisorBase/WeilTensorModels.lean}{GlobalTheory.IndefiniteRowSideFullHecke.WeilTensorModels.Spc.IsSkew}
\leanref{Development/WeilBound/ThetaLifts/LatticeBasics/Hypotheses2.lean}{GlobalTheory.ThK.Dbl.G2}
\leanref{Development/WeilBound/ThetaLifts/LatticeBasics/Hypotheses2.lean}{GlobalTheory.ThK.Dbl.par2}
\leanref{Development/WeilBound/ThetaLifts/LatticeBasics/Hypotheses2.lean}{GlobalTheory.ThK.Dbl.ht2}
\leanref{Development/WeilBound/ThetaLifts/LatticeBasics/HypothesesE.lean}{GlobalTheory.ThK.Dbl.IsFrame2}
\leanref{Development/WeilBound/ThetaLifts/LatticeBasics/HypothesesE.lean}{GlobalTheory.ThK.Dbl.IsStdK}
\leanref{Development/WeilBound/ThetaLifts/LatticeBasics/Hypotheses2.lean}{GlobalTheory.ThK.Dbl.IsHgt}
\leanref{Development/WeilBound/ThetaLifts/LatticeBasics/Hypotheses2.lean}{GlobalTheory.ThK.Dbl.IsSec}
\leanref{Development/WeilBound/ThetaLifts/LatticeBasics/Hypotheses2.lean}{GlobalTheory.ThK.Dbl.flat}
\leanref{Development/WeilBound/ThetaLifts/LatticeBasics/Hypotheses2.lean}{GlobalTheory.ThK.Dbl.eisTerm}
\leanref{Development/WeilBound/ThetaLifts/LatticeBasics/HypothesesE.lean}{GlobalTheory.ThK.Dbl.modFun}
\leanref{Development/WeilBound/ThetaLifts/LatticeBasics/HypothesesE.lean}{GlobalTheory.ThK.Dbl.IsGrowForm}
\leanref{Development/WeilBound/ThetaLifts/LatticeBasics/Hypotheses2.lean}{GlobalTheory.ThK.Dbl.IsDbl}
\leanref{Development/AbelianVarietyTheory/Definitions/Divisors/HilbertNinetyQuotientHomEqIff.lean}{GlobalTheory.WeilM.IsUWeilSplit}
\leanref{Development/AbelianVarietyTheory/Definitions/Divisors/HilbertNinetyQuotientHomEqIff.lean}{GlobalTheory.WeilM.SiegelRule}

\begin{theorem}[Meromorphic continuation of Siegel Eisenstein series \cite{Lan76,MW95}]\label{asm-theta-agree-eisenstein}
Let $E$ be a CM field, Galois over $\mathbf Q$, with maximal totally real
subfield $F$, quadratic character $\eta$ and data
$(\psi,(\tau_w),P,(\tau_v))$ as in Definition~\ref{asm-theta-agree-doubling}.
Let $Y$ be a lattice space of rank $n\ge1$ that is Hermitian or
skew-Hermitian of some signature, let $b$ be a frame of $Y$ with standard
compact subgroup $K_b$, and let $a$ be a height function for $K_b$. Let
$\chi$ be a Hecke character of $E$ with $\chi|_{\mathbf A_F^\times}=\eta^c$
for some $c\ge0$, let $s_1\in\mathbf C$, and let $f$ be a section of degree
$s_1$ for $(K_b,\chi)$, with flat family $(f_s)$. Then there is a function
$\hat E$ on $\mathbf C\times G_Y$ such that:
\begin{enumerate}
\item there is $t\in\mathbf R$ such that for $\operatorname{Re}s>t$ and every
  $x\in G_Y$ the Eisenstein sum of $f_s$ at $x$ converges unconditionally to
  $\hat E(s,x)$, and $s\mapsto\hat E(s,x)$ is holomorphic on
  $\operatorname{Re}s>t$;
\item for every $s_0\in\mathbf C$ there are $k\ge0$, a neighbourhood $U$ of
  $s_0$ and functions $H_s$ of moderate growth on $G_Y$, $s\in\mathbf C$,
  with $H_s(x)=(s-s_0)^k\hat E(s,x)$ for $s\in U\setminus\{s_0\}$ and
  $x\in G_Y$, such that $s\mapsto\mu(H_s)$ is complex differentiable on $U$
  for every growth functional $\mu$.
\end{enumerate}
\end{theorem}
\leanref{Development/WeilBound/ThetaLifts/LatticeBasics/HypothesesE.lean}{GlobalTheory.ThK.Dbl.EisBL}
\leanref{Development/WeilBound/RepresentationTheoryGL/EisensteinSeries/SiegelEisensteinHalfPlane/SchwartzBruhatMajorants/SiegelMultOf.lean}{PositivityTheory.RgEis.Hd.eisBL}
\leanref{Development/AutomorphicArithmetic/ThetaCorrespondence/ThetaPairLifts/ThirdPlaceSeries/EveryCharacterOfLiteratureOfVWin.lean}{UpperTheory.LiteratureTheorems.eis_bl_of_literature}

\begin{theorem}[Siegel-Weil formula \cite{Wei65,KR88,Ich04,Ich07,Yam14}]\label{asm-theta-agree-siegel-weil}
Let $E$ be a CM field, Galois over $\mathbf Q$, with maximal totally real
subfield $F$, quadratic character $\eta$ and data
$(\psi,(\tau_w),P,(\tau_v))$ as in Definition~\ref{asm-theta-agree-doubling}.
Let $A$ and $Y$ be lattice spaces of ranks $m,n\ge1$, one Hermitian and the
other skew-Hermitian, with $A$ of signature $(r,s)$ where $\min(r,s)=0$. Let
$\chi$ be a Hecke character of $E$ with $\chi|_{\mathbf A_F^\times}=\eta^m$,
and let $(I,d,\omega,S,l)$ be a doubling datum for $(A,Y)$ with character
$\chi$. Let $b$ be a frame of $Y$ and $a$ a height function for $K_b$. Let
$\nu$ be a Haar measure on $U_A$ and $\mathcal D$ a fundamental domain for
$U_A(F)$ in $U_A$. Let $\Phi$ be a Schwartz-Bruhat function on
$\mathbf A_F^I$ such that the functions $\omega(k)\Phi$, $k\in K_b$, span a
finite-dimensional space. Put $s_0=(m-n)/2$, $f(x)=l(\omega(x)\Phi)$ and
$f_s(x)=a(x)^{s-s_0}f(x)$, and let $\kappa=1$ if $n<m$ and $\kappa=2$ if
$n\ge m$. Then for every $x\in G_Y$ there is a meromorphic function $e_x$ on
$\mathbf C$ such that for all $s$ with $\operatorname{Re}s$ large enough the
Eisenstein sum of $f_s$ at $x$ converges unconditionally to $e_x(s)$, $e_x$
is holomorphic at $s_0$, and
\[
e_x(s_0)=\frac{\kappa}{\nu(\mathcal D)}\int_{\mathcal D}\theta\bigl(\omega(x)S(g)\Phi\bigr)\,d\nu(g).
\]
\end{theorem}
\leanref{Development/WeilBound/ThetaLifts/LatticeBasics/HypothesesS.lean}{GlobalTheory.ThK.Dbl.SiegelWeilStd}
\leanref{Development/WeilBound/ThetaLifts/LatticeBasics/Hypotheses2.lean}{GlobalTheory.ThK.Dbl.EisValueAt}
\leanref{Development/WeilBound/CoreStatements/AutomorphicKernels/SiegelWeilFormula/SiegelWeilStandard/SiegelWeilStd.lean}{PositivityTheory.SiegelWeilFormula.Fin.siegelWeilStd}
\leanref{Development/AutomorphicArithmetic/ThetaCorrespondence/ThetaPairLifts/ThirdPlaceSeries/EveryCharacterOfLiteratureOfVWin.lean}{UpperTheory.LiteratureTheorems.siegel_weil_std_of_literature}

\begin{definition}\label{asm-theta-agree-models}
Let $E$ be a CM field, Galois over $\mathbf Q$, with maximal totally real
subfield $F$, complex conjugation $c$ and quadratic character $\eta$. Write
$\mathbf A=\mathbf A_F$, $\mathbf A_f$ for the finite adeles of $F$,
$F_\infty=F\otimes\mathbf R$ and $\mathbf A_E=\mathbf A\otimes_FE$. For a
finite set $I$, a \emph{Schwartz-Bruhat function} on $\mathbf A^I$ is a finite
sum of functions $x\mapsto f(x_\infty)g(x_f)$ with $f$ in the Schwartz space
$\mathcal S(F_\infty^I)$ and $g$ a Schwartz-Bruhat function on
$\mathbf A_f^I$; write $f\otimes g$ for such a product.
\begin{enumerate}
\item A \emph{theta base} for $E$ consists of data
  $(\psi,(\tau_w)_w,P,(\tau_v)_{v\in P})$ as in
  Definition~\ref{asm-theta-agree-doubling}, integers $m_\alpha(w)$ and
  $m_\beta(w)$ for each infinite place $w$ of $F$, and a finite place $v_1$ of
  $F$ with an $F$-embedding $\tau_1\colon E\to F_{v_1}$. For a finite place $v$
  and an $F$-embedding $t\colon E\to F_v$, let $\pi_t\colon\mathbf A_E\to F_v$
  be the projection to $E\otimes_FF_v$ followed by $x\otimes y\mapsto t(x)y$;
  for an infinite place $w$, let $\pi_w\colon\mathbf A_E\to\mathbf C$ be the
  component at $w$ read through $\tau_w$. These maps are applied to vectors
  and matrices entrywise.
\item Let $v$ be a finite place of $F$ or $v=\infty$, and put
  $F_v=F_\infty$ if $v=\infty$. An element of $\mathbf A$, of $\mathbf A_E$ or
  of a space of vectors or matrices over them is \emph{at $v$} if its
  components away from $v$ are $0$; an element $g$ of a unitary group $U_X$
  is at $v$ if $g-1$ is, and $U_X^{(v)}$ is the subgroup of these elements.
  Let $T$ be an operator on functions on $\mathbf A^I$, and $\phi,\phi_1$
  functions on $F_v^I$. Then $T$ \emph{acts at $v$ by} $\phi\mapsto\phi_1$ if
  $T(\phi\otimes r)=\phi_1\otimes r$ for every function $r$ on $\mathbf A^I$
  that does not depend on the $v$-component and for which
  $\phi\otimes r\colon x\mapsto\phi(x_v)r(x)$ is a Schwartz-Bruhat function. A
  family $(T_t)_{t\in\mathbf R}$ \emph{acts at $v$ to first order by}
  $\phi\mapsto\phi_1$ if for all such $r$ and all $x$ the function
  $t\mapsto T_t(\phi\otimes r)(x)$ has derivative $\phi_1(x_v)r(x)$ at $t=0$.
\item Let $X$ and $W$ be lattice spaces (Definition~\ref{asm-theta-agree-doubling}(1))
  of ranks $n$ and $d$ with matrices $J_X$ and $J_W$, put
  $I=\{1,\ldots,n\}\times\{1,\ldots,d\}$, and give
  $\mathbb W(X,W)=M_{n\times d}(\mathbf A_E)$ the form
  $\langle x,y\rangle=-\operatorname{Tr}_{\mathbf A_E/\mathbf A}h(x,y)$, where
  $h(x,y)=\sum_{i,j,k,l}x_{ij}(J_X)_{ik}(J_W)_{jl}\,c(y_{kl})$. A \emph{model}
  $M$ for $(X,W)$ consists of an $\mathbf A$-linear isomorphism
  $e=(e_1,e_2)\colon\mathbb W(X,W)\to\mathbf A^I\times\mathbf A^I$ with
  $\langle w,w'\rangle=e_1(w)\cdot e_2(w')-e_1(w')\cdot e_2(w)$, a function
  $\mathrm{vac}_v$ on $F_v^I$ for each finite place $v$, a measure $\mu$ on
  $U_X$ and a subset $\mathcal D\subset U_X$. Its \emph{Schrödinger operators}
  are
  $(\rho(w)\phi)(x)=\psi\bigl(x\cdot e_2(w)+\tfrac12e_1(w)\cdot e_2(w)\bigr)\phi(x+e_1(w))$.
  The model is \emph{regular} if $e$ maps $M_{n\times d}(E)$ into
  $F^I\times F^I$; $\mu$ is a Haar measure and $\mathcal D$ a fundamental
  domain for $U_X(F)$ in $U_X$ with respect to $\mu$; $\mathrm{vac}_v$ is the
  indicator function of $\mathcal O_v^I$ for almost all $v$; and every
  $\mathrm{vac}_v$ is a \emph{vacuum}: a nonzero Schwartz-Bruhat function such
  that, for some $\theta\in\mathcal O_E$ with $\theta+c(\theta)=1$, every
  $\gamma$ in the group generated by the matrices $a\,xy^{\mathsf T}$ with
  $a\in\mathbf A$ at $v$, $a_v\in\mathcal O_v$, $x\in L_X$ and $y\in L_W$ has
  $\rho(\gamma)$ acting at $v$ by
  $\mathrm{vac}_v\mapsto\psi\bigl(\tfrac12\operatorname{Tr}_{\mathbf A_E/\mathbf A}h(\theta\gamma,\gamma)\bigr)\mathrm{vac}_v$.
\item Let $B$ be a theta base and $M$ a model for $(X,W)$. A \emph{Gaussian} is
  an $f_0\in\mathcal S(F_\infty^I)$ for which there are signs
  $\varepsilon_w=\pm1$ such that the following holds: if $u,u'\in\mathbb W(X,W)$
  are at $\infty$ with
  $\pi_w(u')=\varepsilon_wi\,\theta_{X,w}\pi_w(u)\theta_{W,w}^{\mathsf T}$ for
  every $w$, and $t\mapsto\rho(tu)$ and $t\mapsto\rho(tu')$ act at $\infty$ to
  first order by $f_0\mapsto f$ and $f_0\mapsto f'$, then $f'=if$. The
  \emph{polynomial Fock space} $\operatorname{Pol}(M)$ is the smallest subspace
  of $\mathcal S(F_\infty^I)$ that contains every Gaussian and contains $f_1$
  whenever it contains $f$ and $t\mapsto\rho(tu)$, for some $u$ at $\infty$,
  acts at $\infty$ to first order by $f\mapsto f_1$. For $f$ on $F_\infty^I$
  and functions $\phi_v$ on $F_v^I$, write
  $(f\otimes\bigotimes_v\phi_v)(x)=f(x_\infty)\prod_v\phi_v(x_v)$. The space
  $\operatorname{PD}(M)$ of \emph{polynomial data} is spanned by the
  $f\otimes\bigotimes_v\phi_v$ with $f\in\operatorname{Pol}(M)$, every
  $\phi_v$ Schwartz-Bruhat and $\phi_v=\mathrm{vac}_v$ for almost all $v$.
\item Let $v$ be a finite place and $t\colon E\to F_v$ an $F$-embedding. For
  ring maps $\pi^{(1)},\pi^{(0)}\colon\mathbf A_E\to F_v$ and a map
  $\operatorname{co}$ from $\mathbb W(X,W)$ to $M_{n\times d}(F_v)$, a
  \emph{linear model} for $(\pi^{(1)},\pi^{(0)},\operatorname{co})$ is a
  linear map $\iota$ from functions on $F_v^I$ to functions on
  $M_{n\times d}(F_v)$, bijective between the Schwartz-Bruhat spaces, such that
  for $u$ at $v$ and Schwartz-Bruhat $\phi,\phi_1$ with $\rho(u)$ acting at
  $v$ by $\phi\mapsto\phi_1$: if $\pi^{(0)}(u)=0$ then
  $\iota\phi_1(Y)=\iota\phi(Y+\operatorname{co}(u))$ for all $Y$; and if
  $\pi^{(1)}(u)=0$ then
  $\iota\phi_1(\operatorname{co}(x))=\psi(\langle x,u\rangle)\,\iota\phi(\operatorname{co}(x))$
  for every $x$ at $v$ with $\pi^{(0)}(x)=0$. For $b\in GL_n(F_v)$ and
  $b'\in GL_d(F_v)$, a linear model \emph{of the first kind} for $(t,b,b')$ is
  one for $(\pi_t,\pi_{t\circ c},u\mapsto b^{-1}\pi_t(u)(b'^{-1})^{\mathsf T})$,
  and one \emph{of the second kind} is one for
  $(\pi_{t\circ c},\pi_t,u\mapsto-b^{\mathsf T}t(J_X)\pi_{t\circ c}(u)t(J_W)^{\mathsf T}b')$.
  A Schwartz-Bruhat function $\phi$ on $F_v^I$ has \emph{full rank at $t$} if
  some linear model $\iota$ of the second kind for $(t,1,1)$ has
  $\iota\phi(Y)=0$ for every $Y$ of rank different from $n$. The space
  $\mathcal S^\circ(M)$ of \emph{full-rank data} is spanned by the
  $f\otimes\bigotimes_v\phi_v$ with $f\in\mathcal S(F_\infty^I)$, every $\phi_v$
  Schwartz-Bruhat, $\phi_v=\mathrm{vac}_v$ for almost all $v$, and
  $\phi_{v_1}$ of full rank at $\tau_1$.
\item Let $\chi_X$ and $\chi_W$ be Hecke characters of $E$
  (Definition~\ref{asm-theta-agree-doubling}). A representation $\omega$ of
  $U_X\times U_W$ on functions on $\mathbf A^I$ is a \emph{Weil representation}
  for $M$ with characters $(\chi_X,\chi_W)$ if:
  \begin{enumerate}
  \item[(a)] the operators $\omega(1,h)$, $h\in U_W$, satisfy condition (b) of
    Definition~\ref{asm-theta-agree-doubling}(6) before the displayed rule,
    for the operators $\rho$ and with $\omega(1,h)\rho(w)=\rho(wh^{\mathsf T})\omega(1,h)$;
    and there are an isomorphism $e^\square$ for
    $M_{n\times2d}(\mathbf A_E)$ as in (3), for the matrices $J_X$ and
    $\operatorname{diag}(J_W,-J_W)$, with
    $e^\square(w_1\ w_2)=\bigl((e_1(w_1),e_1(w_2)),(e_2(w_1),-e_2(w_2))\bigr)$,
    and operators $\omega^\square(x)$ for $x$ in the doubled group $G_W$ that
    satisfy condition (b) of Definition~\ref{asm-theta-agree-doubling}(6),
    with the character $\chi_X$ and $m=n$, and
    $\omega^\square(\operatorname{diag}(h,1))(\phi_1\boxtimes\phi_2)=(\omega(1,h)\phi_1)\boxtimes\phi_2$
    for $h\in U_W$ and Schwartz-Bruhat $\phi_1,\phi_2$, where
    $(\phi_1\boxtimes\phi_2)(x,y)=\phi_1(x)\phi_2(y)$;
  \item[(b)] the same holds for the operators $\omega(g,1)$, $g\in U_X$, with
    $X$ and $W$ exchanged, the isomorphism $w\mapsto e(w^{\mathsf T})$ for
    $M_{d\times n}(\mathbf A_E)$, the character $\chi_W$ and $m=d$;
  \item[(c)] $x\mapsto\lambda(\omega(x)\Phi)$ is continuous on $U_X\times U_W$
    for every admissible linear form $\lambda$ and Schwartz-Bruhat $\Phi$.
  \end{enumerate}
\item Let $\omega$ be a representation of $U_X\times U_W$ on functions on
  $\mathbf A^I$. The \emph{theta function} of a function $\Phi$ on $\mathbf A^I$
  is $\theta_\omega(\Phi)(g,h)=\sum_{\xi\in F^I}(\omega(g,h)\Phi)(\xi)$, and its
  \emph{theta integral} is the function
  $I_{M,\omega}(\Phi)(h)=\int_{\mathcal D}\theta_\omega(\Phi)(g,h)\,d\mu(g)$ on
  $U_W$. The function $\Phi$ \emph{converges} if
  $g\mapsto\theta_\omega(\Phi)(g,h)$ is $\mu$-integrable on $\mathcal D$ for
  every $h\in U_W$.
\end{enumerate}
\end{definition}
\leanref{Development/AbelianVarietyTheory/BlockTheory/DivisorBase/WeilTensorModels.lean}{GlobalTheory.IndefiniteRowSideFullHecke.WeilTensorModels.Model}
\leanref{Development/AbelianVarietyTheory/BlockTheory/DivisorBase/WeilTensorModels.lean}{GlobalTheory.IndefiniteRowSideFullHecke.WeilTensorModels.Model.IsGood}
\leanref{Development/AbelianVarietyTheory/BlockTheory/DivisorBase/WeilTensorModels.lean}{GlobalTheory.IndefiniteRowSideFullHecke.WeilTensorModels.Model.IsVac}
\leanref{Development/AbelianVarietyTheory/BlockTheory/DivisorBase/WeilTensorModels.lean}{GlobalTheory.IndefiniteRowSideFullHecke.WeilTensorModels.Acts}
\leanref{Development/AbelianVarietyTheory/BlockTheory/DivisorBase/WeilTensorModels.lean}{GlobalTheory.IndefiniteRowSideFullHecke.WeilTensorModels.ActsD}
\leanref{Development/AbelianVarietyTheory/BlockTheory/DivisorBase/WeilTensorModels.lean}{GlobalTheory.IndefiniteRowSideFullHecke.WeilTensorModels.Model.IsGauss}
\leanref{Development/AbelianVarietyTheory/BlockTheory/DivisorBase/WeilTensorModels.lean}{GlobalTheory.IndefiniteRowSideFullHecke.WeilTensorModels.Model.Pol}
\leanref{Development/AbelianVarietyTheory/Definitions/Divisors/CTSmallCone.lean}{GlobalTheory.ThK.PolData}
\leanref{Development/AbelianVarietyTheory/BlockTheory/DivisorBase/WeilTensorModels.lean}{GlobalTheory.IndefiniteRowSideFullHecke.WeilTensorModels.Model.IsLin}
\leanref{Development/AbelianVarietyTheory/BlockTheory/DivisorBase/WeilTensorModels.lean}{GlobalTheory.IndefiniteRowSideFullHecke.WeilTensorModels.Model.IsLinW}
\leanref{Development/AbelianVarietyTheory/BlockTheory/DivisorBase/WeilTensorModels.lean}{GlobalTheory.IndefiniteRowSideFullHecke.WeilTensorModels.Model.IsLinBar}
\leanref{Development/AbelianVarietyTheory/BlockTheory/DivisorBase/WeilTensorModels.lean}{GlobalTheory.IndefiniteRowSideFullHecke.WeilTensorModels.Model.SoLoc}
\leanref{Development/AbelianVarietyTheory/BlockTheory/DivisorBase/WeilTensorModels.lean}{GlobalTheory.IndefiniteRowSideFullHecke.WeilTensorModels.Model.So}
\leanref{Development/AbelianVarietyTheory/BlockTheory/DivisorBase/WeilTensorModels.lean}{GlobalTheory.IndefiniteRowSideFullHecke.WeilTensorModels.Model.IsWeil}
\leanref{Development/AbelianVarietyTheory/Definitions/Divisors/HilbertNinetyQuotientHomEqIff.lean}{GlobalTheory.WeilM.IsUWeilPair}
\leanref{Development/AbelianVarietyTheory/Definitions/Divisors/HilbertNinetyQuotientHomEqIff.lean}{GlobalTheory.WeilM.IsUWeilB}
\leanref{Development/AbelianVarietyTheory/BlockTheory/DivisorBase/WeilTensorModels.lean}{GlobalTheory.IndefiniteRowSideFullHecke.WeilTensorModels.theta}
\leanref{Development/AbelianVarietyTheory/BlockTheory/DivisorBase/WeilTensorModels.lean}{GlobalTheory.IndefiniteRowSideFullHecke.WeilTensorModels.Model.I}
\leanref{Development/AbelianVarietyTheory/BlockTheory/DivisorBase/WeilTensorModels.lean}{GlobalTheory.IndefiniteRowSideFullHecke.WeilTensorModels.Model.Conv}

\begin{definition}\label{asm-theta-agree-triple}
Let $E$, $F$, $c$ and $\eta$ be as in Definition~\ref{asm-theta-agree-models},
and let $B$ be a theta base for $E$.
\begin{enumerate}
\item A \emph{Weil pair} of signature $(r,s;r',s')$ for $B$ consists of a
  Hermitian lattice space $X$ of rank $n$ and signature $(r,s)$, a
  skew-Hermitian lattice space $W$ of rank $d$ and signature $(r',s')$, a
  regular model $M$ for $(X,W)$, Hecke characters $\chi_X$ and $\chi_W$ of $E$
  with $\chi_X|_{\mathbf A^\times}=\eta^n$ and $\chi_W|_{\mathbf A^\times}=\eta^d$,
  and a Weil representation $\omega$ for $M$ with characters
  $(\chi_X,\chi_W)$. We write $I=I_{M,\omega}$.
\item A \emph{triple datum} $G$ of ranks $(a,a')$ consists of integers
  $p,q\geq0$, lattice spaces $V$, $V'$ and $W$ of ranks $a$, $a'$ and $a+a'$,
  a model $M$ for $(V,W)$ and a model $M'$ for $(V',W)$. It is \emph{proper}
  for $B$ if $a'<a$; $a+a'$ is even; $a+a'<m_\alpha(w)-m_\beta(w)$ for every
  infinite place $w$; $2(p+q)+1<\lfloor(a-a')/2\rfloor$ and $2(p+q)+1<a'$; $V$
  and $V'$ are Hermitian of signatures $(a-p,p)$ and $(q,a'-q)$; $W$ is
  skew-Hermitian of signature $(a+a'-p-q,p+q)$; and $M$ and $M'$ are regular.
  For lattice spaces $V^\circ$, $V'^\circ$ and $W^\circ$ of ranks $a$, $a'$ and
  $a+a'$, the \emph{re-chosen} datum $G[V^\circ,V'^\circ,W^\circ]$ is $G$ with
  the matrices $\theta_w$ and $b_v$ of $V$, $V'$ and $W$ replaced by those of
  $V^\circ$, $V'^\circ$ and $W^\circ$; its models have the same isomorphisms,
  vacua, measures and domains, so its theta integrals are those of $G$.
\item A \emph{character datum} $\mathcal X=(\alpha,\beta,S_0,Q)$ consists of
  Hecke characters $\alpha$ and $\beta$ of $E$ and finite sets $S_0$ and $Q$ of
  finite places of $F$; put $\nu=\alpha/\beta$ and $T=S_0\cup Q$. A finite
  place $v$ is \emph{clean} for $\mathcal X$ if $v\notin T$ and the rational
  prime below $v$ splits completely in $E$. The datum is \emph{proper} for
  $(B,a,a')$ if $\alpha|_{\mathbf A^\times}=\eta^{a'}$ and
  $\beta|_{\mathbf A^\times}=\eta^a$; $\alpha(z)=\prod_w(\pi_w(z)/|\pi_w(z)|)^{m_\alpha(w)}$
  and $\beta(z)=\prod_w(\pi_w(z)/|\pi_w(z)|)^{m_\beta(w)}$ for every
  $z\in\mathbf A_E^\times$ whose finite components are $1$; $S_0\cap Q=\emptyset$;
  $T\subset P$; $\nu$ is ramified at every place of $E$ above $T$; $\alpha$ and
  $\beta$ are unramified at every place of $E$ not above $T$; and $v_1$ is clean
  for $\mathcal X$.
\item A \emph{triple} $(B,G,\mathcal X,\omega,\omega')$ consists of a triple
  datum $G$ of ranks $(a,a')$ proper for $B$, a character datum $\mathcal X$
  proper for $(B,a,a')$, a Weil representation $\omega$ for $M$ with characters
  $(\beta,\beta/\alpha)$ and a Weil representation $\omega'$ for $M'$ with
  characters $(\alpha,\alpha/\beta)$. Its theta integrals are
  $I=I_{M,\omega}$ and $I'=I_{M',\omega'}$. Then $(V,W,M,\beta,\beta/\alpha,\omega)$
  is a Weil pair of signature $(a-p,p;a+a'-p-q,p+q)$ and
  $(V',W,M',\alpha,\alpha/\beta,\omega')$ one of signature
  $(q,a'-q;a+a'-p-q,p+q)$, since $\eta^2=1$.
\end{enumerate}
These are the triples of Definition~\ref{th-theta-hypotheses}.
\end{definition}
\leanref{Development/AbelianVarietyTheory/Definitions/Divisors/CTSmallCone.lean}{GlobalTheory.ThK.IsPair}
\leanref{Development/AbelianVarietyTheory/BlockTheory/DivisorBase/WeilTensorModels.lean}{GlobalTheory.IndefiniteRowSideFullHecke.WeilTensorModels.Geo}
\leanref{Development/AbelianVarietyTheory/BlockTheory/DivisorBase/WeilTensorModels.lean}{GlobalTheory.IndefiniteRowSideFullHecke.WeilTensorModels.Geo.IsGood}
\leanref{Development/AbelianVarietyTheory/Definitions/Divisors/CTSmallCone.lean}{GlobalTheory.ThK.reG}
\leanref{Development/AbelianVarietyTheory/BlockTheory/DivisorBase/WeilTensorModels.lean}{GlobalTheory.IndefiniteRowSideFullHecke.WeilTensorModels.Chars}
\leanref{Development/AbelianVarietyTheory/BlockTheory/DivisorBase/WeilTensorModels.lean}{GlobalTheory.IndefiniteRowSideFullHecke.WeilTensorModels.Chars.Good}
\leanref{Development/AbelianVarietyTheory/BlockTheory/DivisorBase/WeilTensorModels.lean}{GlobalTheory.IndefiniteRowSideFullHecke.WeilTensorModels.Chars.IsGood}
\leanref{Development/AbelianVarietyTheory/BlockTheory/DivisorBase/WeilTensorModels.lean}{GlobalTheory.IndefiniteRowSideFullHecke.WeilTensorModels.HasExp}
\leanref{Development/AbelianVarietyTheory/BlockTheory/DivisorBase/WeilTensorModels.lean}{GlobalTheory.IndefiniteRowSideFullHecke.WeilTensorModels.IsTriple}
\leanref{Development/WeilBound/ThetaLifts/ConeHypotheses/PureEqPureA.lean}{PositivityTheory.ThKG.pairV}
\leanref{Development/WeilBound/ThetaLifts/ConeHypotheses/PureEqPureA.lean}{PositivityTheory.ThKG.pairV'}

\begin{definition}\label{asm-theta-agree-matched}
Let $(B,G,\mathcal X,\omega,\omega')$ be a triple of ranks $(a,a')$
(Definition~\ref{asm-theta-agree-triple}), with models $M$ and $M'$, and
put $I_0=\{1,\ldots,a\}\times\{1,\ldots,a+a'\}$ and
$I'_0=\{1,\ldots,a'\}\times\{1,\ldots,a+a'\}$.
\begin{enumerate}
\item For a lattice space $Y$, let $K_\infty(Y)$ be the group of $k\in U_Y$ at
  $\infty$ with $\pi_w(k)\theta_{Y,w}=\theta_{Y,w}\pi_w(k)$ for every $w$, and
  call a continuous homomorphism $\gamma\colon\mathbf R\to U_Y$ with values at
  $\infty$ a \emph{line} in $U_Y$. For a finite place $v$, the
  \emph{augmentation space} $\mathfrak a_v(\omega)$ is spanned by the
  $\phi_1-\phi$ with $\phi,\phi_1$ Schwartz-Bruhat on $F_v^{I_0}$ and
  $\omega(g,1)$ acting at $v$ by $\phi\mapsto\phi_1$ for some
  $g\in U_V^{(v)}$; $\mathfrak a_v(\omega')$ is defined in the same way.
\item \emph{Real maps} $R=(\Pi,p,p')$ consist of a complex vector space $\Pi$
  and linear maps $p\colon\mathcal S(F_\infty^{I_0})\to\Pi$ and
  $p'\colon\mathcal S(F_\infty^{I'_0})\to\Pi$. They are \emph{matched} to the
  triple if:
  \begin{enumerate}
  \item[(a)] $p(\operatorname{Pol}(M))=\Pi$; $p(f_1)=0$ whenever
    $f\in\operatorname{Pol}(M)$, $\gamma$ is a line in $U_V$ and
    $t\mapsto\omega(\gamma(t),1)$ acts at $\infty$ to first order by
    $f\mapsto f_1$; and $p(f_1)=p(f)$ whenever $f\in\operatorname{Pol}(M)$,
    $k\in K_\infty(V)$ and $\omega(k,1)$ acts at $\infty$ by $f\mapsto f_1$;
    and the same holds for $p'$, $M'$, $V'$ and $\omega'$;
  \item[(b)] if $f\in\operatorname{Pol}(M)$ and $f'\in\operatorname{Pol}(M')$
    have $p(f)=p'(f')$, and $f\mapsto f_1$ and $f'\mapsto f'_1$ are the actions
    at $\infty$ to first order of $t\mapsto\omega(1,\gamma(t))$ and
    $t\mapsto\omega'(1,\gamma(t))$ for a line $\gamma$ in $U_W$, or the actions
    at $\infty$ of $\omega(1,k)$ and $\omega'(1,k)$ for some $k\in K_\infty(W)$,
    then $p(f_1)=p'(f'_1)$;
  \item[(c)] $\Pi\neq0$, and $0$ and $\Pi$ are the only subspaces $N$ of $\Pi$
    with $p(f_1)\in N$ whenever $f\in\operatorname{Pol}(M)$, $p(f)\in N$ and
    $f\mapsto f_1$ is one of the actions of $\omega$ in (b).
  \end{enumerate}
\item \emph{Finite maps} $R_f=(\Pi_v,p_v,p'_v)_v$ consist, for every finite
  place $v$, of a complex vector space $\Pi_v$ and linear maps $p_v$ and $p'_v$
  from functions on $F_v^{I_0}$, respectively $F_v^{I'_0}$, to $\Pi_v$. They
  are \emph{matched} to the triple if for every finite place $v$:
  $p_v$ and $p'_v$ map the Schwartz-Bruhat functions onto $\Pi_v$; $p_v$
  vanishes on $\mathfrak a_v(\omega)$ and $p'_v$ on $\mathfrak a_v(\omega')$;
  if $h\in U_W^{(v)}$, $\omega(1,h)$ and $\omega'(1,h)$ act at $v$ by
  $\phi\mapsto\phi_1$ and $\phi'\mapsto\phi'_1$ for Schwartz-Bruhat functions,
  and $p_v(\phi)=p'_v(\phi')$, then $p_v(\phi_1)=p'_v(\phi'_1)$; $\Pi_v\neq0$,
  and $0$ and $\Pi_v$ are the only subspaces $N$ with $p_v(\phi_1)\in N$
  whenever $p_v(\phi)\in N$ and $\omega(1,h)$ acts at $v$ by $\phi\mapsto\phi_1$
  for some $h\in U_W^{(v)}$ and Schwartz-Bruhat $\phi,\phi_1$;
  $p_v(\mathrm{vac}_v)=p'_v(\mathrm{vac}'_v)\neq0$ if $v\notin T$; and if
  $v\in P\cap T$, the following \emph{zeta condition} holds. Let $\alpha_v$,
  $\beta_v$ and $\nu_v$ be the characters of $F_v^\times$ given by $\alpha$,
  $\beta$ and $\nu$ at the place of $E$ defined by $\tau_v$, and $Z$ and $Z'$
  the zeta functions of Lemma~\ref{th-bottom-split-zeta} for $k=F_v$,
  $\alpha_v$, $\beta_v$ and a fixed uniformizer. There are a linear model
  $\iota$ of the first kind for $(\tau_v,b_{V,v},b_{W,v})$ for $M$ and a linear
  model $\iota'$ of the second kind for $(\tau_v,b_{V',v},b_{W,v})$ for $M'$,
  carrying $\mathrm{vac}_v$ and $\mathrm{vac}'_v$ to the indicator functions
  of the matrices with entries in $\mathcal O_v$, such that for Schwartz-Bruhat
  $\phi$ and $\phi'$, $p_v(\phi)=p'_v(\phi')$ if and only if
  $Z(\iota\phi)=Z'(\iota'\phi')$.
\item A \emph{matched quadruple} $(B,G,\mathcal X,\omega,\omega',R,R_f)$ is a
  triple with real and finite maps matched to it. For a matched quadruple,
  $\Phi$ and $\Phi'$ \emph{correspond} if there are a finite set $T_1\supset T$
  of finite places and finite sums
  $\Phi=\sum_if_i\otimes\bigotimes_v\phi_{i,v}$ and
  $\Phi'=\sum_jf'_j\otimes\bigotimes_v\phi'_{j,v}$ with
  $f_i\in\operatorname{Pol}(M)$, $f'_j\in\operatorname{Pol}(M')$, all
  $\phi_{i,v}$ and $\phi'_{j,v}$ Schwartz-Bruhat, $\phi_{i,v}=\mathrm{vac}_v$
  and $\phi'_{j,v}=\mathrm{vac}'_v$ for $v\notin T_1$, and
  \[
  \sum_ip(f_i)\otimes\bigotimes_{v\in T_1}p_v(\phi_{i,v})
  =\sum_jp'(f'_j)\otimes\bigotimes_{v\in T_1}p'_v(\phi'_{j,v})
  \quad\text{in }\Pi\otimes\bigotimes_{v\in T_1}\Pi_v.
  \]
  Then $\Phi\in\operatorname{PD}(M)$ and $\Phi'\in\operatorname{PD}(M')$.
  Further, $\Phi$ and $\Phi'$ \emph{match} if $\Phi'\in\mathcal S^\circ(M')$
  and there are finite sums $\Phi=\sum_if_i\otimes g_i$ and
  $\Phi'=\sum_jf'_j\otimes g'_j$, with Schwartz-Bruhat $g_i$ and $g'_j$ on
  $\mathbf A_f^{I_0}$ and $\mathbf A_f^{I'_0}$, and sequences $s_{k,i}\to f_i$
  and $s'_{k,j}\to f'_j$ in the Schwartz spaces, such that for every $k$,
  $\sum_js'_{k,j}\otimes g'_j\in\mathcal S^\circ(M')$ and corresponds to
  $\sum_is_{k,i}\otimes g_i$.
\item A complex number $c$ is a \emph{constant} of the matched quadruple if
  $I(\Phi)=c\,I'(\Phi')$ whenever $\Phi$ and $\Phi'$ match, and a
  \emph{polynomial constant} if $I(\Phi)=c\,I'(\Phi')$ whenever
  $\Phi'\in\mathcal S^\circ(M')$ and $\Phi$ and $\Phi'$ correspond.
\item The matched quadruple is \emph{nondegenerate} if it has a nonzero
  polynomial constant, $I(\Phi)\neq0$ for some $\Phi\in\operatorname{PD}(M)$,
  and $I'(\Phi')\neq0$ for some
  $\Phi'\in\operatorname{PD}(M')\cap\mathcal S^\circ(M')$. For an integer
  $n\geq0$, \emph{nondegeneracy holds below $n$} if every matched quadruple,
  over every CM field and theta base, with $p+q<n$ is nondegenerate.
\end{enumerate}
\end{definition}
\leanref{Development/AbelianVarietyTheory/BlockTheory/DivisorBase/WeilTensorModels.lean}{GlobalTheory.IndefiniteRowSideFullHecke.WeilTensorModels.KInf}
\leanref{Development/AbelianVarietyTheory/BlockTheory/DivisorBase/WeilTensorModels.lean}{GlobalTheory.IndefiniteRowSideFullHecke.WeilTensorModels.Lines}
\leanref{Development/AbelianVarietyTheory/BlockTheory/DivisorBase/WeilTensorModels.lean}{GlobalTheory.IndefiniteRowSideFullHecke.WeilTensorModels.augAt}
\leanref{Development/AbelianVarietyTheory/BlockTheory/DivisorBase/WeilTensorModels.lean}{GlobalTheory.IndefiniteRowSideFullHecke.WeilTensorModels.RealMaps}
\leanref{Development/AbelianVarietyTheory/BlockTheory/DivisorBase/WeilTensorModels.lean}{GlobalTheory.IndefiniteRowSideFullHecke.WeilTensorModels.RealQuot}
\leanref{Development/AbelianVarietyTheory/BlockTheory/DivisorBase/WeilTensorModels.lean}{GlobalTheory.IndefiniteRowSideFullHecke.WeilTensorModels.IsRealMaps}
\leanref{Development/AbelianVarietyTheory/BlockTheory/DivisorBase/WeilTensorModels.lean}{GlobalTheory.IndefiniteRowSideFullHecke.WeilTensorModels.FinMaps}
\leanref{Development/AbelianVarietyTheory/BlockTheory/DivisorBase/WeilTensorModels.lean}{GlobalTheory.IndefiniteRowSideFullHecke.WeilTensorModels.IsFinMaps}
\leanref{Development/AbelianVarietyTheory/BlockTheory/DivisorBase/WeilTensorModels.lean}{GlobalTheory.IndefiniteRowSideFullHecke.WeilTensorModels.IsZetaAt}
\leanref{Development/AbelianVarietyTheory/BlockTheory/DivisorBase/WeilTensorModels.lean}{GlobalTheory.IndefiniteRowSideFullHecke.WeilTensorModels.IsMatched}
\leanref{Development/AbelianVarietyTheory/BlockTheory/DivisorBase/WeilTensorModels.lean}{GlobalTheory.IndefiniteRowSideFullHecke.WeilTensorModels.PolMatch}
\leanref{Development/WeilBound/ThetaLifts/ConeHypotheses/PureEqPureA.lean}{PositivityTheory.ThKG.polMatch_mem}
\leanref{Development/AbelianVarietyTheory/BlockTheory/DivisorBase/WeilTensorModels.lean}{GlobalTheory.IndefiniteRowSideFullHecke.WeilTensorModels.Match}
\leanref{Development/AbelianVarietyTheory/BlockTheory/DivisorBase/WeilTensorModels.lean}{GlobalTheory.IndefiniteRowSideFullHecke.WeilTensorModels.IsConst}
\leanref{Development/AbelianVarietyTheory/Definitions/Divisors/CTSmallCone.lean}{GlobalTheory.ThK.PolConst}
\leanref{Development/WeilBound/ThetaLifts/ConeHypotheses/PolDense.lean}{PositivityTheory.ThKG.Claim}
\leanref{Development/WeilBound/ThetaLifts/ConeHypotheses/PolDense.lean}{PositivityTheory.ThKG.Below}

\begin{definition}\label{asm-theta-agree-planes}
Let $E$ and $F$ be as in Definition~\ref{asm-theta-agree-models}, let $B$ be a
theta base for $E$, and let $Y$ and $Y_1$ be lattice spaces of ranks $m$ and
$m_1$ with matrices $J$ and $J_1$ and lattices $L$ and $L_1$. Write
$h(x,y)=x^{\mathsf T}J\,c(y)$ for $x,y\in E^m$.
\begin{enumerate}
\item An \emph{adapted plane} $H$ in $Y$ over $Y_1$ at positions $i_0<i_1$
  (counting from $0$) consists of vectors $x,y\in E^m$, a matrix
  $\iota\in M_{m\times m_1}(E)$ and a fractional ideal $\mathfrak b$ of $E$,
  such that: $m=m_1+2$; $h(x,x)=h(y,y)=0$ and $h(x,y)=1$;
  $h(\iota u,x)=h(\iota u,y)=0$ for every $u$; $\iota^{\mathsf T}J\,c(\iota)=J_1$;
  $u\in L_1$ if and only if $\iota u\in L$; every $z\in L$ is $sx+\iota u+ty$
  with $sx\in L$, $ty\in L$ and $u\in L_1$; $s\in\mathfrak b$ if and only if
  $sx\in L$; for every infinite place $w$,
  $\theta_{Y,w}\tau_w(\iota)=\tau_w(\iota)\theta_{Y_1,w}$ and
  $\tau_w(y)=\lambda_w\theta_{Y,w}\tau_w(x)$ for some $\lambda_w\neq0$; and for
  every $v\in P$ the matrix $b_{Y,v}$ has column $i_0$ equal to $\tau_v(x)$,
  column $i_1$ equal to $\tau_v(y)$, and its other columns, in order, equal to
  the columns of $\tau_v(\iota)b_{Y_1,v}$.
\item Let $H=(x,y,\iota,\mathfrak b)$ be an adapted plane in $Y$ over $Y_1$,
  $t\in\mathbf A_E^\times$ and $h_1\in U_{Y_1}$. An element $m\in U_Y$ is a
  \emph{Levi element} for $(t,h_1)$ if $mx=tx$, $my=c(t)^{-1}y$ and
  $m\iota=\iota h_1$. The \emph{unipotent group} $N_H\subset U_Y$ is generated
  by the $u\in U_Y$ with $ux=x$ and $u\iota z\in\iota z+\mathbf A_Ex$ for every
  $z\in\mathbf A_E^{m_1}$. A function $f_0$ on $U_Y$ is the \emph{constant
  term} of a function $f$ along $H$ if
  $f_0(g)=\int_{\mathcal F}f(ug)\,d\mu(u)$ for every $g$, every Haar measure
  $\mu$ on $N_H$ and every fundamental domain $\mathcal F$ for
  $N_H\cap U_Y(F)$ in $N_H$ of measure $1$. A function $f$ on $U_Y$ has the
  \emph{law} $(\chi,z)$ for a character $\chi$ of $\mathbf A_E^\times$ and
  $z\in\mathbf R$ if $f(mg)=\chi(t)|t|^zf(g)$ for every Levi element $m$ for
  $(t,1)$ and every $g$, with $|t|=|N_{E/F}(t)|_{\mathbf A}$.
\item Let $(B,G,\mathcal X,\omega,\omega')$ be a triple of ranks $(a,a')$ and
  $H'$ an adapted plane in $W$ over a lattice space $W_1$. Functions $f_\alpha$
  and $f_\beta$ on $U_W$ \emph{split} a function $f$ along $H'$ if
  $f_\alpha+f_\beta$ is the constant term of $f$ along $H'$, $f_\alpha$ has
  the law $(\alpha,a'/2)$ and $f_\beta$ the law $(\beta,a/2)$. If
  $(f_\alpha,f_\beta)$ splits $f$, then $(\lambda f_\alpha,\lambda f_\beta)$
  splits $\lambda f$, and $(0,0)$ splits $0$.
\item Let $G$ and $G_1$ be triple data, $H$ an adapted plane in a space of
  $G$ over the corresponding space of $G_1$, and $H'$ one in $W$ over $W_1$.
  Then $(G,G_1,H,H')$ is a \emph{$\beta$-configuration} if $G$ and $G_1$ have
  ranks $(a,b+2)$ and $(a,b)$, $V_1=V$, $p_1=p$, $q_1+1=q$, $H$ is in $V'$
  over $V'_1$ at positions $(0,b+1)$ and $H'$ is at positions $(a,a+b+1)$; it
  is an \emph{$\alpha$-configuration} if $G$ and $G_1$ have ranks $(b+2,a')$
  and $(b,a')$, $V'_1=V'$, $q=q_1=0$, $p_1+1=p$, $H$ is in $V$ over $V_1$ at
  positions $(0,b+1)$ and $H'$ is at positions $(0,b+1)$. Write
  $k_\beta=\lfloor(a-b-2)/2\rfloor$, respectively
  $k_\alpha=\lfloor(b+2-a')/2\rfloor$.
\item Let $M$ be a model for $(X,W)$ and $M_1$ one for $(X_1,W_1)$, of ranks
  $(n,d)$ and $(n_1,d_1)$, with Schrödinger operators $\rho$ and $\rho_1$, and
  let $\iota_X\in M_{n\times n_1}(E)$, $\iota_W\in M_{d\times d_1}(E)$ and
  $\mathcal J\subset\mathbb W(X,W)$. An \emph{evaluation map} for
  $(M,M_1,\iota_X,\iota_W,\mathcal J)$ is a linear map $e$ from functions on
  $\mathbf A^{n\times d}$ to functions on $\mathbf A^{n_1\times d_1}$ that
  maps Schwartz-Bruhat functions to Schwartz-Bruhat functions, with
  $e(\rho(\iota_Xz\iota_W^{\mathsf T})\Phi)=\rho_1(z)e(\Phi)$ for
  $z\in\mathbb W(X_1,W_1)$ and $e(\rho(u)\Phi)=e(\Phi)$ for $u\in\mathcal J$,
  for Schwartz-Bruhat $\Phi$, and such that there are a nonzero continuous
  linear map $e_\infty$ between the Schwartz spaces and linear maps $e_v$ at
  the finite places, with $e_v(\mathrm{vac}_v)=\mathrm{vac}_{1,v}$, $e_v$
  mapping the Schwartz-Bruhat functions onto the Schwartz-Bruhat functions,
  and $e(f\otimes\bigotimes_v\phi_v)=e_\infty(f)\otimes\bigotimes_ve_v(\phi_v)$
  whenever the $\phi_v$ are Schwartz-Bruhat and almost all equal
  $\mathrm{vac}_v$. For adapted planes $H=(x,y,\iota,\mathfrak b)$ in $X$ and
  $H'=(x',y',\iota',\mathfrak b')$ in $W$, put
  $\mathcal J_0(H')=\{zx'^{\mathsf T}\}$,
  $\mathcal J_\alpha(H,H')=\mathcal J_0(H')\cup\{x(\iota'z')^{\mathsf T}\}$ and
  $\mathcal J_\beta(H,H')=\{xz'^{\mathsf T}\}\cup\{(\iota z)x'^{\mathsf T}\}$,
  with $z$ and $z'$ running over vectors over $\mathbf A_E$ of the appropriate
  sizes.
\item For a Hermitian lattice space $X$ of rank $n$, let $\hat L_X$ be the
  subgroup of $\mathbf A_E^n$ generated by the $a\,x$ with $a\in\mathbf A$,
  $a_\infty=0$, $a_v\in\mathcal O_v$ for every $v$, and $x\in L_X$. A subgroup
  $K\subset U_X$ is the \emph{lattice stabilizer} of $X$ if it consists of the
  $g$ with $\pi_w(g)\theta_{X,w}=\theta_{X,w}\pi_w(g)$ for every $w$,
  $g\hat L_X\subset\hat L_X$ and $g^{-1}\hat L_X\subset\hat L_X$.
\item Let $M$ be a model for $(X,W)$, $\omega$ a representation of
  $U_X\times U_W$, $e$ a linear map as in (5), $K\subset U_X$ a subgroup with a
  measure $\mu_K$, $\mu$ a measure on $\mathbf A_E^\times$,
  $u\colon\mathbf A_E^\times\to\mathbb W(X,W)$, $\chi$ a character of
  $\mathbf A_E^\times$ and $z\geq0$ an integer. The \emph{cone} of $\Phi$ is
  \[
  \operatorname{cone}(\Phi)(y)=\int_K\int_{\mathbf A_E^\times}
  e\bigl(\rho(u(l))\,\omega(k,1)\Phi\bigr)(y)\,\chi(l)|l|^z\,d\mu(l)\,d\mu_K(k).
  \]
  For adapted planes $H=(x,y,\ldots)$ in $X$ and $H'=(x',y',\ldots)$ in $W$,
  put $u_\alpha(l)=l\,xy'^{\mathsf T}$ and $u_\beta(l)=c(l)\,yx'^{\mathsf T}$.
\end{enumerate}
\end{definition}
\leanref{Development/AbelianVarietyTheory/BlockTheory/DivisorBase/WeilTensorModels.lean}{GlobalTheory.IndefiniteRowSideFullHecke.WeilTensorModels.Plane}
\leanref{Development/AbelianVarietyTheory/BlockTheory/DivisorBase/WeilTensorModels.lean}{GlobalTheory.IndefiniteRowSideFullHecke.WeilTensorModels.Plane.IsAdapted}
\leanref{Development/AbelianVarietyTheory/Definitions/Divisors/CTSmallCone.lean}{GlobalTheory.ThK.IsLevi}
\leanref{Development/AbelianVarietyTheory/Definitions/Divisors/CTSmallCone.lean}{GlobalTheory.ThK.unip}
\leanref{Development/AbelianVarietyTheory/Definitions/Divisors/CTSmallCone.lean}{GlobalTheory.ThK.IsCT}
\leanref{Development/AbelianVarietyTheory/Definitions/Divisors/CTSmallCone.lean}{GlobalTheory.ThK.HasLaw}
\leanref{Development/AbelianVarietyTheory/Definitions/Divisors/CTSmallCone.lean}{GlobalTheory.ThK.Split2}
\leanref{Development/WeilBound/ThetaLifts/ConeHypotheses/PureEqPureA.lean}{PositivityTheory.ThKG.split2_smul}
\leanref{Development/AbelianVarietyTheory/Definitions/Divisors/CTSmallCone.lean}{GlobalTheory.ThK.CfgB}
\leanref{Development/AbelianVarietyTheory/Definitions/Divisors/CTSmallCone.lean}{GlobalTheory.ThK.CfgA}
\leanref{Development/AbelianVarietyTheory/Definitions/Divisors/CTSmallCone.lean}{GlobalTheory.ThK.IsEv}
\leanref{Development/AbelianVarietyTheory/Definitions/Divisors/CTSmallCone.lean}{GlobalTheory.ThK.iso0}
\leanref{Development/AbelianVarietyTheory/Definitions/Divisors/CTSmallCone.lean}{GlobalTheory.ThK.isoA}
\leanref{Development/AbelianVarietyTheory/Definitions/Divisors/CTSmallCone.lean}{GlobalTheory.ThK.isoB}
\leanref{Development/AbelianVarietyTheory/Definitions/Divisors/CTSmallCone.lean}{GlobalTheory.ThK.latHat}
\leanref{Development/AbelianVarietyTheory/Definitions/Divisors/CTSmallCone.lean}{GlobalTheory.ThK.IsStab}
\leanref{Development/AbelianVarietyTheory/Definitions/Divisors/CTSmallCone.lean}{GlobalTheory.ThK.cone}
\leanref{Development/AbelianVarietyTheory/Definitions/Divisors/CTSmallCone.lean}{GlobalTheory.ThK.shiftA}
\leanref{Development/AbelianVarietyTheory/Definitions/Divisors/CTSmallCone.lean}{GlobalTheory.ThK.shiftB}

\begin{definition}\label{asm-theta-agree-parts}
Let $(B,G,\mathcal X,\omega,\omega')$ be a triple of ranks $(a,a')$
(Definition~\ref{asm-theta-agree-triple}), and let $\operatorname{ht}$ be the
height of Definition~\ref{asm-theta-agree-doubling}(2), for invertible
matrices of any size.
\begin{enumerate}
\item The triple has \emph{decay} if every Schwartz-Bruhat function converges
  for $(M,\omega)$, and every $\Phi'\in\mathcal S^\circ(M')$ converges for
  $(M',\omega')$, has $|I'(\Phi')(h)|\leq C\operatorname{ht}(h)^k$ for some
  $C$ and $k$, and for every $N$ satisfies
  $|\theta_{\omega'}(\Phi')(g,h)|\,\overline{\operatorname{ht}}(g)^N\leq
  C\operatorname{ht}(h)^k$ for some $C$ and $k$, where
  $\overline{\operatorname{ht}}(g)$ is the infimum of $\operatorname{ht}(\gamma g)$
  over $\gamma\in U_{V'}(F)$.
\item A matched quadruple has \emph{one constant} if $I'(\Phi')\neq0$ for some
  $\Phi'\in\mathcal S^\circ(M')$, it has exactly one constant, and every
  constant of it is nonzero.
\item A matched quadruple \emph{describes the coinvariants} if $\omega'(g,h)$
  maps $\mathcal S^\circ(M')$ into itself for all $g$ and $h$, and for every
  finite place $v$ clean for $\mathcal X$: a Schwartz-Bruhat $\phi$ has
  $p_v(\phi)=0$ if and only if $\phi\in\mathfrak a_v(\omega)$, and likewise
  for $p'_v$ and $\mathfrak a_v(\omega')$; for every $F$-embedding
  $t\colon E\to F_v$, every element of $\Pi_v$ is $p'_v(\phi)$ for some $\phi$
  of full rank at $t$; and every linear form $\ell$ that vanishes on
  $\mathfrak a_v(\omega')$ is $\ell_0\circ p'_v$ on the Schwartz-Bruhat
  functions for some linear form $\ell_0$ on $\Pi_v$.
\item For $v\in P$ let $\nu_v$ be as in
  Definition~\ref{asm-theta-agree-matched}(3), with conductor exponent
  $c_v$, the least $c\geq0$ such that $\nu_v$ is trivial on the units in
  $1+\varpi_v^c\mathcal O_v$, for a fixed uniformizer $\varpi_v$. Its Gauss sum
  and epsilon factors are
  \[
  \mathfrak g(\nu_v)=\sum_{y\in(\mathcal O_v/\varpi_v^{c_v})^\times}\nu_v(y)^{-1}\psi_v(y\varpi_v^{-c_v}),\qquad
  \epsilon(\nu_v,z)=N(v)^{-c_vz}\,\mathfrak g(\nu_v)\,\nu_v(\varpi_v)^{c_v}
  \quad(z\in\mathbf Z),
  \]
  with fixed representatives $y\in\mathcal O_v$. For a fractional ideal
  $\mathfrak a$ of $E$ put
  $\nu(\mathfrak a)=\prod_w\nu_w(\varpi_w)^{\operatorname{ord}_w\mathfrak a}$,
  for fixed uniformizers $\varpi_w$ of the places $w$ of $E$. For signs
  $e_1,e_2,s\in\{\pm1\}$, fractional ideals $\mathfrak b$ and $\mathfrak c$ of
  $E$ and $z\in\mathbf Z$, the \emph{step factor} is
  \[
  \sigma_{e_1,e_2,s}(\mathfrak b,\mathfrak c;z)=\nu(\mathfrak b)^{e_1}\nu(\mathfrak c)^{e_2}
  \prod_{v\in P\cap T}\nu_v(s)\,\epsilon(\nu_v,z).
  \]
  The \emph{partial $L$-function} $L^{S_0}(\nu,s)$, for
  $\operatorname{Re}s>1$, is the product of
  $(1-\nu_w(\varpi_w)N(w)^{-s})^{-1}$ over the finite places $w$ of $E$ not
  above $S_0$ at which $\nu$ is unramified.
\item For signs $e_1,e_2,s$, the \emph{$\beta$-step} holds if for every CM
  field $E$ as in Definition~\ref{asm-theta-agree-models}, every theta base $B$
  for $E$, every $\beta$-configuration $(G,G_1,H,H')$
  (Definition~\ref{asm-theta-agree-planes}(4)) and all real maps $R$ and $R_1$
  for $G$ and $G_1$, there is $C\neq0$ such that for all matched quadruples
  $(B,G,\mathcal X,\omega,\omega',R,R_f)$ and
  $(B,G_1,\mathcal X,\omega_1,\omega'_1,R_1,R_{f,1})$ and all constants $c$
  and $c_1$ of them, $C\sigma\neq0$ and
  $c\,L^{S_0}(\nu,k_\beta+1)=(C\sigma)^{-1}c_1$, where
  $\sigma=\sigma_{e_1,e_2,s}(\mathfrak b_H,\mathfrak b_{H'};-k_\beta)$. The
  \emph{$\alpha$-step} holds if the same holds for $\alpha$-configurations with
  $c=C\sigma_{e_1,e_2,s}(\mathfrak b_H,\mathfrak b_{H'};1-k_\alpha)\,L^{S_0}(\nu,k_\alpha)\,c_1$.
\end{enumerate}
\end{definition}
\leanref{Development/AbelianVarietyTheory/BlockTheory/DivisorBase/WeilTensorModels.lean}{GlobalTheory.IndefiniteRowSideFullHecke.WeilTensorModels.Decay}
\leanref{Development/AbelianVarietyTheory/BlockTheory/DivisorBase/WeilTensorModels.lean}{GlobalTheory.IndefiniteRowSideFullHecke.WeilTensorModels.OneConst}
\leanref{Development/AbelianVarietyTheory/BlockTheory/DivisorBase/WeilTensorModels.lean}{GlobalTheory.IndefiniteRowSideFullHecke.WeilTensorModels.Coinv}
\leanref{Development/AbelianVarietyTheory/BlockTheory/DivisorBase/WeilTensorModels.lean}{GlobalTheory.IndefiniteRowSideFullHecke.WeilTensorModels.gauss}
\leanref{Development/AbelianVarietyTheory/BlockTheory/DivisorBase/WeilTensorModels.lean}{GlobalTheory.IndefiniteRowSideFullHecke.WeilTensorModels.eps}
\leanref{Development/AbelianVarietyTheory/BlockTheory/DivisorBase/WeilTensorModels.lean}{GlobalTheory.IndefiniteRowSideFullHecke.WeilTensorModels.stepFac}
\leanref{Development/AbelianVarietyTheory/Definitions/Foundations/NotIsSplitOfIsInertAt.lean}{GlobalTheory.LFun.charL}
\leanref{Development/AbelianVarietyTheory/BlockTheory/DivisorBase/WeilTensorModels.lean}{GlobalTheory.IndefiniteRowSideFullHecke.WeilTensorModels.StepBeta}
\leanref{Development/AbelianVarietyTheory/BlockTheory/DivisorBase/WeilTensorModels.lean}{GlobalTheory.IndefiniteRowSideFullHecke.WeilTensorModels.StepAlpha}

\subsection{Common quotients of Fock spaces}\label{subsection-asm-fock-quotients}

\begin{definition}\label{asm-theta-agree-real-fock}
Let $P,Q,R,S\geq0$. Let $\mathrm{Ix}(P,Q)$ be the disjoint union of a
\emph{first block} $\{1,\ldots,P\}$ and a \emph{second block}
$\{1,\ldots,Q\}$; put $\epsilon(a)=1$ for $a$ in the first block and
$\epsilon(a)=-1$ for $a$ in the second, and let $J_{P,Q}$ be the diagonal
matrix with entries $\epsilon(a)$, $a\in\mathrm{Ix}(P,Q)$. Let
$\mathfrak u(P,Q)$ be the real Lie algebra of the $Y\in M_{P+Q}(\mathbf C)$
with $Y^*J_{P,Q}+J_{P,Q}Y=0$. Call $Y\in M_{P+Q}(\mathbf C)$
\emph{block-diagonal} if $Y_{ab}=0$ whenever $a$ and $b$ lie in different
blocks, and let $E_{ab}$ be the matrix units.
\begin{enumerate}
\item The \emph{Fock space} $\mathcal P(P,Q;R,S)$ is the polynomial ring over
  $\mathbf C$ in variables $x_l$, $l=(a,b)\in\mathrm{Ix}(P,Q)\times\mathrm{Ix}(R,S)$.
  Put $s_l=\epsilon(a)\epsilon(b)$. The \emph{lowering} and \emph{raising
  operators} are $L_l=\partial/\partial x_l$ and $R_l=x_l$ if $s_l=1$, and
  $L_l=x_l$ and $R_l=-\partial/\partial x_l$ if $s_l=-1$, where $x_l$ acts by
  multiplication. Put $\omega_{l,l'}=R_lL_{l'}+\tfrac12\delta_{ll'}$. For
  $X\in M_{P+Q}(\mathbf C)$, $Y\in M_{R+S}(\mathbf C)$ and
  $c\in M_{(P+Q)\times(R+S)}(\mathbf C)$ put
  \[
  d_1(X)=\sum_{a,a'}\sum_bX_{aa'}\,\omega_{(a,b),(a',b)},\qquad
  d_2(Y)=\sum_a\sum_{b,b'}Y_{bb'}\,\omega_{(a,b),(a,b')},\qquad
  h(c)=\sqrt{2\pi}\sum_l\bigl(c_lR_l-s_l\bar c_lL_l\bigr),
  \]
  with $a,a'\in\mathrm{Ix}(P,Q)$, $b,b'\in\mathrm{Ix}(R,S)$ and
  $c_{(a,b)}=c_{ab}$; $h(c)$ is the \emph{Heisenberg operator} of $c$.
\item A \emph{$(\mathfrak g,K)$-module} of type $(R,S)$ is a complex vector
  space $N$ with a $\mathbf C$-linear map $Y\mapsto N(Y)$ from
  $M_{R+S}(\mathbf C)$ to $\operatorname{End}N$ such that
  $N(YY'-Y'Y)=N(Y)N(Y')-N(Y')N(Y)$; every vector lies in a finite-dimensional
  subspace stable under the $N(Y)$ with $Y$ block-diagonal; and every vector is
  a sum of vectors $v$ for which there are $\mu_b\in\mathbf Q$ with
  $N(E_{bb})v=\mu_bv$ for all $b$. It is \emph{irreducible} if $N\neq0$ and $0$
  and $N$ are its only subspaces stable under every $N(Y)$, and it has
  \emph{scalar endomorphisms} if every linear map $N\to N$ that commutes with
  every $N(Y)$ is a scalar multiple of the identity.
\item Let $N$ be a $(\mathfrak g,K)$-module of type $(R,S)$ and
  $x,t\in\mathbf Q$. A \emph{lift} of type $(x,t)$ to $N$ is a surjective linear
  map $\pi\colon\mathcal P(P,Q;R,S)\to N$ with
  $\pi(d_1(X)f)=x\operatorname{tr}(X)\,\pi(f)$ and
  $N(Y)\pi(f)=\pi(d_2(Y)f)+t\operatorname{tr}(Y)\,\pi(f)$ for all $X$, $Y$ and
  $f$.
\item Let $\iota$ be a finite set and $L_i$ an operator on a complex vector
  space $V_i$ for $i\in\iota$. Then $L_i^{(i)}$ is the operator on
  $\bigotimes_{j\in\iota}V_j$ that acts as $L_i$ on the factor $V_i$ and as the
  identity on the others. An operator $L$ on a complex vector space $V$ is
  \emph{locally finite} if every vector lies in a finite-dimensional subspace
  stable under $L$; an operator $U$ on $V$ is the \emph{exponential} of $L$ if
  $\sum_{n\geq0}\ell(L^nv)/n!$ converges to $\ell(Uv)$ for every $v\in V$ and
  every linear form $\ell$ on $V$.
\item For a finite set $\alpha$ and $H\in M_\alpha(\mathbf C)$, a \emph{frame}
  of $H$ of type $(P,Q)$ is a matrix $A\in M_{\alpha\times\mathrm{Ix}(P,Q)}(\mathbf C)$
  with $A^{\mathsf T}H\bar A=J_{P,Q}$ and $AA'=1$ for some $A'$.
\item Let $B$ be a theta base for a CM field $E$
  (Definition~\ref{asm-theta-agree-models}(1)), let $Y$ be a lattice space of
  rank $m$, and let $C_w\in M_{m\times\mathrm{Ix}(R,S)}(\mathbf C)$ for every
  infinite place $w$ of $F$. A line $\gamma$ in $U_Y$
  (Definition~\ref{asm-theta-agree-matched}(1)) is \emph{generated} by a family
  $(Y_w)_w$ of matrices $Y_w\in M_{R+S}(\mathbf C)$, relative to $(C_w)_w$, if
  $\pi_w(\gamma(t))C_w=C_w\exp(tY_w)$ for every $t\in\mathbf R$ and every $w$.
\item Let $I$ be a finite set, $f,f_1\in\mathcal S(F_\infty^I)$, and $T_t$
  ($t\in\mathbf R$) and $T$ operators on functions on $\mathbf A^I$. The family
  $(T_t)$ \emph{acts weakly at $\infty$ to first order} by $f\mapsto f_1$ if for
  every Schwartz-Bruhat function $g$ on $\mathbf A_f^I$ and every admissible
  linear form $\lambda$ (Definition~\ref{asm-theta-agree-doubling}(6)) the
  function $t\mapsto\lambda(T_t(f\otimes g))$ has derivative
  $\lambda(f_1\otimes g)$ at $t=0$. The operator $T$ \emph{acts weakly at
  $\infty$} by $f\mapsto f_1$ if $T(f\otimes g)=f_1\otimes g$ for every such $g$.
\end{enumerate}
\end{definition}
\leanref{Formalism/LocalRingCardinality/DataTheta.lean}{Formalism.BasicNotions.TH.Jm}
\leanref{Formalism/LocalRingCardinality/DataTheta.lean}{Formalism.BasicNotions.TH.upq}
\leanref{Formalism/LocalRingCardinality/DataTheta.lean}{Formalism.BasicNotions.TH.IsKC}
\leanref{Formalism/LocalRingCardinality/DataTheta.lean}{Formalism.BasicNotions.TH.Pol}
\leanref{Formalism/LocalRingCardinality/DataTheta.lean}{Formalism.BasicNotions.TH.lowB}
\leanref{Formalism/LocalRingCardinality/DataTheta.lean}{Formalism.BasicNotions.TH.raiB}
\leanref{Formalism/LocalRingCardinality/DataTheta.lean}{Formalism.BasicNotions.TH.omE}
\leanref{Formalism/LocalRingCardinality/DataTheta.lean}{Formalism.BasicNotions.TH.fockAct1}
\leanref{Formalism/LocalRingCardinality/DataTheta.lean}{Formalism.BasicNotions.TH.fockAct2}
\leanref{Development/AbelianVarietyTheory/Proofs/DivisorProofs/IsFockMap.lean}{GlobalTheory.RealLift.heisOp}
\leanref{Formalism/LocalRingCardinality/DataTheta.lean}{Formalism.BasicNotions.TH.GKMod}
\leanref{Formalism/SharpRestriction/Divisibility/SchurIrreducible.lean}{Formalism.MainNotions.SchurIrreducible.IrrGK}
\leanref{Development/AbelianVarietyTheory/Proofs/DivisorProofs/IsFockMap.lean}{GlobalTheory.RealLift.ScalarEnd}
\leanref{Development/AbelianVarietyTheory/Proofs/DivisorProofs/IsFockMap.lean}{GlobalTheory.RealLift.IsLift}
\leanref{Development/AbelianVarietyTheory/Proofs/DivisorProofs/IsFockMap.lean}{GlobalTheory.RealLift.atPl}
\leanref{Development/AbelianVarietyTheory/Proofs/DivisorProofs/IsFockMap.lean}{GlobalTheory.RealLift.sumPl}
\leanref{Development/AbelianVarietyTheory/Proofs/DivisorProofs/IsFockMap.lean}{GlobalTheory.RealLift.LocFin}
\leanref{Development/AbelianVarietyTheory/Proofs/DivisorProofs/IsFockMap.lean}{GlobalTheory.RealLift.IsExp}
\leanref{Development/AbelianVarietyTheory/Proofs/DivisorProofs/IsFockMap.lean}{GlobalTheory.RealLift.IsFrame}
\leanref{Development/AbelianVarietyTheory/Proofs/DivisorProofs/IsFockMap.lean}{GlobalTheory.RealLift.RLines}
\leanref{Development/AbelianVarietyTheory/Proofs/DivisorProofs/IsFockMap.lean}{GlobalTheory.RealLift.IsGen}
\leanref{Development/AbelianVarietyTheory/Proofs/DivisorProofs/IsFockMap.lean}{GlobalTheory.RealLift.SD}
\leanref{Development/AbelianVarietyTheory/Proofs/DivisorProofs/IsFockMap.lean}{GlobalTheory.RealLift.SA}

\begin{lemma}\label{asm-theta-agree-real-fock-relations}
Let $P,Q,R,S\geq0$, and use the notation of
Definition~\ref{asm-theta-agree-real-fock}(1).
\begin{enumerate}
\item The maps $d_1$ and $d_2$ are $\mathbf C$-linear, satisfy
  $d_i(XX'-X'X)=d_i(X)d_i(X')-d_i(X')d_i(X)$, and $d_1(X)d_2(Y)=d_2(Y)d_1(X)$
  for all $X$, $X'$ and $Y$.
\item If $X$ and $Y$ are block-diagonal, then $d_1(X)$ and $d_2(Y)$ preserve the
  space of polynomials of degree at most $n$, for every $n$.
\item For a monomial $x^\nu=\prod_lx_l^{\nu_l}$,
  $d_1(E_{aa})x^\nu=\sum_bs_{(a,b)}(\nu_{(a,b)}+\tfrac12)\,x^\nu$ and
  $d_2(E_{bb})x^\nu=\sum_as_{(a,b)}(\nu_{(a,b)}+\tfrac12)\,x^\nu$.
\end{enumerate}
\end{lemma}
\leanref{Development/AbelianVarietyTheory/Proofs/DivisorProofs/FockOperatorsLieRepresentationCommute.lean}{GlobalTheory.FockAlg.fock_lie_and_commute}
\leanref{Development/AbelianVarietyTheory/Proofs/DivisorProofs/FockPolynomialsLocallyFiniteRationalWeights.lean}{GlobalTheory.FockPolGK.kStable1}
\leanref{Development/AbelianVarietyTheory/Proofs/DivisorProofs/FockPolynomialsLocallyFiniteRationalWeights.lean}{GlobalTheory.FockPolGK.kStable2}
\leanref{Development/AbelianVarietyTheory/Proofs/DivisorProofs/FockPolynomialsLocallyFiniteRationalWeights.lean}{GlobalTheory.FockPolGK.monomial_mem_weightSpace_fockAction_first}
\leanref{Development/AbelianVarietyTheory/Proofs/DivisorProofs/FockPolynomialsLocallyFiniteRationalWeights.lean}{GlobalTheory.FockPolGK.monomial_mem_weightSpace_fockAction_second}
\leanref{Development/AbelianVarietyTheory/Proofs/DivisorProofs/FockPolynomialsLocallyFiniteRationalWeights.lean}{GlobalTheory.FockPolGK.pol_locFin_weights}
\begin{proof}
Part (1) is a computation with the relations
$L_lR_{l'}-R_{l'}L_l=\delta_{ll'}$, $L_lL_{l'}=L_{l'}L_l$ and
$R_lR_{l'}=R_{l'}R_l$, and its result is stated. For (2), if $s_l=s_{l'}$, then
$\omega_{l,l'}$ is $x_l\partial/\partial x_{l'}$ or
$-\partial/\partial x_l\circ x_{l'}$, up to a scalar, and both preserve the
degree. If $X$ is block-diagonal and $X_{aa'}\neq0$, then $a$ and $a'$ lie in
the same block, so $s_{(a,b)}=s_{(a',b)}$ for every $b$; the same holds for
$d_2(Y)$. For (3), $d_1(E_{aa})=\sum_b(R_lL_l+\tfrac12)$ with $l=(a,b)$. If
$s_l=1$, then $R_lL_l=x_l\partial/\partial x_l$ has eigenvalue $\nu_l$ on
$x^\nu$; if $s_l=-1$, then
$R_lL_l=-\partial/\partial x_l\circ x_l$ has eigenvalue $-(\nu_l+1)$. In both
cases $R_lL_l+\tfrac12$ has eigenvalue $s_l(\nu_l+\tfrac12)$. The same holds
for $d_2(E_{bb})$.
\end{proof}

\begin{definition}\label{asm-theta-agree-real-quotient-lie-module}
Let $R,S\geq0$, and use the set $\mathrm{Ix}(R,S)$, its two blocks and the
matrix units $E_{ab}$ of Definition~\ref{asm-theta-agree-real-fock}. For
$a,b\in\mathrm{Ix}(R,S)$ write $a\prec b$ if $a$ and $b$ lie in the same block
and $a$ comes before $b$ in it. For $\mu\colon\mathrm{Ix}(R,S)\to\mathbf Q$ put
\[
c(\mu)=\sum_a\mu_a^2+\sum_{a\prec b}(\mu_a-\mu_b).
\]
Inner products are linear in the second variable.
\begin{enumerate}
\item A \emph{Lie module} of type $(R,S)$ is a complex vector space $T$ with a
  $\mathbf C$-linear map $Y\mapsto T(Y)$ from $M_{R+S}(\mathbf C)$ to
  $\operatorname{End}T$ such that $T(YY'-Y'Y)=T(Y)T(Y')-T(Y')T(Y)$. A subspace
  of $T$ is \emph{stable} if it is stable under every $T(Y)$, and the stable
  subspace \emph{generated} by a subset is the smallest stable subspace
  containing it. The \emph{compact Casimir operator} of $T$ is
  $\Omega_T=\sum T(E_{ab})T(E_{ba})$, the sum over the pairs $(a,b)$ with $a$
  and $b$ in the same block.
\item Let $T$ be a Lie module of type $(R,S)$ and
  $\Lambda\colon\mathrm{Ix}(R,S)\to\mathbf Q$. A \emph{highest-weight datum}
  of $T$ of weight $\Lambda$ consists of a finite-dimensional complex inner
  product space $A$, operators $\rho(a,b)$ on $A$ for $a,b$ in the same block,
  a linear map $\iota\colon A\to T$ and a vector $f\in A$ with the following
  properties, for $a,b$ in the same block and $c,d$ in the same block:
  $\langle\rho(a,b)u,v\rangle=\langle u,\rho(b,a)v\rangle$ for all $u,v\in A$;
  \[
  \rho(a,b)\rho(c,d)-\rho(c,d)\rho(a,b)=\delta_{bc}\,\rho(a,d)-\delta_{da}\,\rho(c,b);
  \]
  $\rho(a,a)f=\Lambda_af$, and $\rho(a,b)f=0$ if $a\prec b$; and
  $\iota\rho(a,b)=T(E_{ab})\iota$.
\item A Lie module $T$ of type $(R,S)$ is \emph{generated with Casimir gap} in
  weight $\Lambda$ if $T\neq0$ and $T$ has a highest-weight datum
  $(A,\rho,\iota,f)$ of weight $\Lambda$ such that $T$ is the stable subspace
  generated by $\iota(f)$, and such that for every real
  $t\leq c(\Lambda)+R+S$ the only vector $x$ with $\Omega_Tx=tx$ in the span
  of the vectors $T(E_{ab})\iota(u)$, with $a$ and $b$ in different blocks and
  $u\in A$, is $0$.
\end{enumerate}
\end{definition}
\leanref{Development/WeilBound/CoreStatements/GaussianSeminormBounds/RealSeminormFiltration/CoreDefs.lean}{PositivityTheory.KReal.IsLie}
\leanref{Development/WeilBound/CoreStatements/GaussianSeminormBounds/RealSeminormFiltration/CoreDefs.lean}{PositivityTheory.KReal.genG}
\leanref{Development/WeilBound/CoreStatements/GaussianSeminormBounds/RealSeminormFiltration/CoreDefs.lean}{PositivityTheory.KReal.casK}
\leanref{Development/WeilBound/CoreStatements/GaussianSeminormBounds/RealSeminormFiltration/CoreDefs.lean}{PositivityTheory.KReal.casW}
\leanref{Development/WeilBound/CoreStatements/GaussianSeminormBounds/RealSeminormFiltration/CoreDefs.lean}{PositivityTheory.KReal.AdjK}
\leanref{Development/WeilBound/CoreStatements/GaussianSeminormBounds/RealSeminormFiltration/CoreDefs.lean}{PositivityTheory.KReal.GlK}
\leanref{Development/WeilBound/CoreStatements/GaussianSeminormBounds/RealSeminormFiltration/CoreDefs.lean}{PositivityTheory.KReal.HwK}
\leanref{Development/WeilBound/CoreStatements/GaussianSeminormBounds/RealSeminormFiltration/CoreDefs.lean}{PositivityTheory.KReal.StarK}
\leanref{Development/WeilBound/CoreStatements/GaussianSeminormBounds/RealSeminormFiltration/CoreDefs.lean}{PositivityTheory.KReal.Side}

\begin{lemma}\label{asm-theta-agree-real-quotient-filtration}
Let $R,S\geq0$, let $T$ be a Lie module of type $(R,S)$
(Definition~\ref{asm-theta-agree-real-quotient-lie-module}), let $A$ be a
finite-dimensional complex inner product space, and let $\rho(a,b)$, for
$a,b\in\mathrm{Ix}(R,S)$ in the same block, be operators on $A$ with
$\langle\rho(a,b)u,v\rangle=\langle u,\rho(b,a)v\rangle$ for all $u,v\in A$.
Let $c_0\in\mathbf R$ with $\sum\rho(a,b)\rho(b,a)=c_0$, the sum over the pairs
in the same block, and let $\iota\colon A\to T$ be a linear map with
$\iota\rho(a,b)=T(E_{ab})\iota$. Suppose that for every real
$t\leq c_0+R+S$ the only vector $x$ with $\Omega_Tx=tx$ in the span of the
vectors $T(E_{ab})\iota(u)$, with $a$ and $b$ in different blocks and
$u\in A$, is $0$. Let $G$ be the stable subspace generated by $\iota(A)$.
\begin{enumerate}
\item If $x\in G$ and $\Omega_Tx-c_0x\in\iota(A)$, then $x\in\iota(A)$.
\item Every $x\in G$ lies in a finite-dimensional subspace stable under
  $T(E_{ab})$ for all $a,b$ in the same block.
\end{enumerate}
\end{lemma}
\leanref{Development/WeilBound/CoreStatements/GaussianSeminormBounds/RealSeminormFiltration/CoreCasGl.lean}{PositivityTheory.KReal.cas_bound_gl}
\leanref{Development/WeilBound/CoreStatements/GaussianSeminormBounds/RealSeminormFiltration/CoreCas.lean}{PositivityTheory.KReal.cas_bound}
\leanref{Development/WeilBound/CoreStatements/GaussianSeminormBounds/RealSeminormFiltration/CoreCasGl.lean}{PositivityTheory.KReal.CasGl.adC_sw}
\leanref{Development/WeilBound/CoreStatements/GaussianSeminormBounds/RealSeminormFiltration/CoreCasGl.lean}{PositivityTheory.KReal.CasGl.adC_cas}
\leanref{Development/WeilBound/CoreStatements/GaussianSeminormBounds/RealSeminormFiltration/CoreCasB.lean}{PositivityTheory.KReal.filtX}
\leanref{Development/WeilBound/CoreStatements/GaussianSeminormBounds/RealSeminormFiltration/CoreCasB.lean}{PositivityTheory.KReal.filtX_e}
\leanref{Development/WeilBound/CoreStatements/GaussianSeminormBounds/RealSeminormFiltration/CoreCasB.lean}{PositivityTheory.KReal.filtX_perm}
\leanref{Development/WeilBound/CoreStatements/GaussianSeminormBounds/RealSeminormFiltration/CoreCasB.lean}{PositivityTheory.KReal.wordX_comm}
\leanref{Development/WeilBound/CoreStatements/GaussianSeminormBounds/RealSeminormFiltration/CoreCasA.lean}{PositivityTheory.KReal.cas_one}
\leanref{Development/WeilBound/CoreStatements/GaussianSeminormBounds/RealSeminormFiltration/CoreCasA.lean}{PositivityTheory.KReal.pos_of_eig}
\leanref{Development/WeilBound/CoreStatements/GaussianSeminormBounds/RealSeminormFiltration/CoreCasC.lean}{PositivityTheory.KReal.cas_tensor}
\begin{proof}
Let $\mathcal K$ be the set of pairs $k=(a,b)$ with $a$ and $b$ in the same
block, with $\bar k=(b,a)$, and $\mathcal I$ the set of pairs $(a,b)$ with $a$
and $b$ in different blocks. Write $E_{(a,b)}=E_{ab}$, and put $e_k=T(E_k)$ and
$X_i=T(E_i)$. So $\Omega_T=\sum_ke_ke_{\bar k}$, and
$\sum_k\rho(k)\rho(\bar k)=c_0$.

\emph{Brackets.} For $k\in\mathcal K$ and $i\in\mathcal I$, the matrix
$E_kE_i-E_iE_k$ vanishes at the entries in $\mathcal K$, so
$E_kE_i-E_iE_k=\sum_{j\in\mathcal I}c_k(j,i)E_j$ for numbers $c_k(j,i)$. For
$i,j\in\mathcal I$, the matrix $E_iE_j-E_jE_i$ vanishes at the entries in
$\mathcal I$. By the bracket relation of $T$:
\begin{enumerate}
\item[(a)] $e_kX_i-X_ie_k=\sum_jc_k(j,i)X_j$;
\item[(b)] $X_iX_j-X_jX_i$ lies in the span of the $e_k$.
\end{enumerate}
A computation with matrix units gives, for the matrices
$c_k=(c_k(j,i))_{j,i}$ on $\mathcal I$:
\begin{enumerate}
\item[(c)] $c_{\bar k}=c_k^*$, the conjugate transpose;
\item[(d)] $\sum_kc_kc_{\bar k}=(R+S)\cdot1$.
\end{enumerate}
For (d), the entry of the sum at $(j,i)$ is $0$ for $j\neq i$; for $i=(e,f)$
the entry at $(i,i)$ is the number of $(a,b)\in\mathcal K$ with $a=e$ plus the
number with $b=f$, which is $R+S$, as $e$ and $f$ lie in different blocks.

\emph{The filtration.} For a word $w=(w_1,\ldots,w_N)$ in $\mathcal I$ put
$X_w=X_{w_1}\cdots X_{w_N}$, and let $F_n$ be the span of the vectors
$X_w\iota(u)$ with $w$ of length at most $n$ and $u\in A$. Then
$F_0=\iota(A)$, $F_n\subset F_{n+1}$ and $X_iF_n\subset F_{n+1}$, and $F_n$ is
finite-dimensional, since $F_{n+1}\subset F_n+\sum_iX_iF_n$. By (a) and
induction on $N$,
\[
e_kX_w=X_we_k+\sum_{s=1}^N\sum_{j\in\mathcal I}c_k(j,w_s)\,X_{w[s\leftarrow j]},
\]
where $w[s\leftarrow j]$ is $w$ with its $s$-th letter replaced by $j$. With
$e_k\iota=\iota\rho(k)$ this gives $e_kF_n\subset F_n$. The union of the $F_n$
contains $\iota(A)$ and is stable under the $e_k$ and the $X_i$, hence under
every $T(Y)$, as $Y$ is a combination of matrix units; so it contains $G$.
This proves (2). If $w'$ is a permutation of a word $w$ of length at most
$n+1$, then $X_w\iota(u)-X_{w'}\iota(u)\in F_n$, by (b), $e_kF_m\subset F_m$,
$X_iF_m\subset F_{m+1}$ and induction on the permutation. Put
$C=\Omega_T-c_0$; it preserves every $F_n$.

\emph{The main step:} $F_{n+1}\subset F_n+C(F_{n+1})$ for every $n\geq0$.
Suppose not, and let $\phi$ be a linear form on $T$ that vanishes on
$F_n+C(F_{n+1})$ but not on $F_{n+1}$. For $w\in\mathcal I^{n+1}$ let
$\Psi(w)\in A$ be the vector with $\langle\Psi(w),u\rangle=\phi(X_w\iota(u))$
for all $u\in A$. It suffices to show that $\Psi=0$, for then $\phi$ vanishes
on $F_{n+1}$.

(i) Fix a position $s$ and $w\in\mathcal I^{n+1}$, put
$\Phi_i=\Psi(w[s\leftarrow i])$ for $i\in\mathcal I$, and let
$q(s,w)=\sum_k\sum_{i,j}c_k(i,j)\langle\Phi_i,\rho(\bar k)\Phi_j\rangle$. Then
$\operatorname{Re}q(s,w)\geq0$, with equality only if every $\Phi_i$ is $0$.
First, if $\Theta\colon\mathcal I\to A$ satisfies $\sum_iX_i\iota(\Theta_i)=0$,
then $\sum_i\langle\Phi_i,\Theta_i\rangle=0$: modulo $F_n$, on which $\phi$
vanishes, $X_{w[s\leftarrow i]}\iota(\Theta_i)$ is $X_{w'}X_i\iota(\Theta_i)$,
where $w'$ is $w$ with its $s$-th letter removed. Next, give $A^{\mathcal I}$
the orthogonal sum of the inner products, let $R_k$ be the operator with
$(R_k\Theta)_i=\sum_jc_k(i,j)\Theta_j+\rho(k)\Theta_i$, and let
$m(\Theta)=\sum_iX_i\iota(\Theta_i)$. By (c) and the adjoint relation,
$\langle R_kx,y\rangle=\langle x,R_{\bar k}y\rangle$, and by (a),
$e_km(\Theta)=m(R_k\Theta)$. So every $R_k$ preserves $\ker m$, and
$N=(\ker m)^\perp$ is stable under every $R_k$; the self-adjoint operator
$\Sigma=\sum_kR_kR_{\bar k}-(c_0+R+S)$ preserves $N$. If $x\in N$ and
$\Sigma x=tx$ with $t\leq0$, then $m(x)$ lies in the span of the
$X_i\iota(u)$ and $\Omega_Tm(x)=(t+c_0+R+S)m(x)$, so $m(x)=0$ by hypothesis,
and $x=0$. So the eigenvalues of $\Sigma$ on $N$ are positive, and
$\operatorname{Re}\langle x,\Sigma x\rangle\geq0$ for $x\in N$, with equality
only for $x=0$. By the first remark, $\Phi=(\Phi_i)$ lies in $N$. By (d),
$\sum_k\rho(k)\rho(\bar k)=c_0$ and the symmetry $k\mapsto\bar k$,
$(\Sigma\Phi)_i=2\sum_k\sum_jc_k(i,j)\rho(\bar k)\Phi_j$, so
$\langle\Phi,\Sigma\Phi\rangle=2q(s,w)$.

(ii) For a matrix $c$ on $\mathcal I$ and maps $\Psi',\Psi''$ from
$\mathcal I^{n+1}$ to $A$ put
$(c\star\Psi')(w)=\sum_s\sum_jc(w_s,j)\Psi'(w[s\leftarrow j])$ and
$\langle\Psi',\Psi''\rangle=\sum_w\langle\Psi'(w),\Psi''(w)\rangle$. Then
$\langle c^*\star\Psi',\Psi''\rangle=\langle\Psi',c\star\Psi''\rangle$, by the
substitution $(w,j)\mapsto(w[s\leftarrow j],w_s)$. Write $\rho(k)\Psi$ for
$w\mapsto\rho(k)\Psi(w)$. As $\phi$ vanishes on $C(F_{n+1})$,
$\phi(\Omega_TX_w\iota(u))=c_0\,\phi(X_w\iota(u))$ for all $w$ and $u$.
Expand $e_ke_{\bar k}X_w$ by the formula for $e_kX_w$ applied twice, use
$\sum_ke_ke_{\bar k}\iota=c_0\iota$ and (c), put $u=\Psi(w)$ and sum over
$w$. The result is
\[
\sum_k\langle c_{\bar k}\star\Psi,\rho(\bar k)\Psi\rangle
+\sum_k\langle c_k\star\Psi,\rho(k)\Psi\rangle
+\sum_k\langle c_k\star(c_{\bar k}\star\Psi),\Psi\rangle=0.
\]
The first two sums are equal, and the third is
$\sum_k\langle c_{\bar k}\star\Psi,c_{\bar k}\star\Psi\rangle\geq0$, by (c).
Hence $\operatorname{Re}\sum_k\langle c_k\star\Psi,\rho(k)\Psi\rangle\leq0$.
Now put
$H(s,w)=\sum_k\sum_jc_k(w_s,j)\langle\Psi(w),\rho(\bar k)\Psi(w[s\leftarrow j])\rangle$.
Then $q(s,w)=\sum_iH(s,w[s\leftarrow i])$, so
$\sum_wq(s,w)=|\mathcal I|\sum_wH(s,w)$; and by the adjoint relation
$\sum_{s,w}H(s,w)$ is the complex conjugate of
$\sum_k\langle c_k\star\Psi,\rho(k)\Psi\rangle$. So
$\sum_{s,w}\operatorname{Re}q(s,w)\leq0$, while each term is nonnegative by
(i). Hence every $\operatorname{Re}q(s,w)$ is $0$, and by (i)
$\Psi(w[s\leftarrow i])=0$ for all $s$, $w$ and $i$. So $\Psi=0$.

\emph{Conclusion.} By the main step, $C$ induces a surjective, hence
injective, endomorphism of the finite-dimensional space $F_{n+1}/F_n$. So if
$x\in F_{n+1}$ and $Cx\in F_n$, then $x\in F_n$. Let $x\in G$ with
$Cx\in\iota(A)=F_0$. Then $x\in F_n$ for some $n$, and descending induction on
$n$ gives $x\in F_0$. This proves (1).
\end{proof}

\begin{lemma}\label{asm-theta-agree-real-quotient-local}
Let $R,S\geq0$, let $T$ be a Lie module of type $(R,S)$, let
$\Lambda\colon\mathrm{Ix}(R,S)\to\mathbf Q$, and let $(A,\rho,\iota,f)$ be a
highest-weight datum of $T$ of weight $\Lambda$
(Definition~\ref{asm-theta-agree-real-quotient-lie-module}) with
$\iota(f)\neq0$. Suppose that for every real $t\leq c(\Lambda)+R+S$ the only
vector $x$ with $\Omega_Tx=tx$ in the span of the vectors $T(E_{ab})\iota(u)$,
with $a$ and $b$ in different blocks and $u\in A$, is $0$. Let $G$ be the
stable subspace generated by $\iota(f)$. Then there is a subspace $Z\subset T$
with $\iota(f)\notin Z$ that contains every stable subspace $N\subset G$ with
$N\neq G$.
\end{lemma}
\leanref{Development/WeilBound/CoreStatements/GaussianSeminormBounds/RealSeminormFiltration/CoreMain.lean}{PositivityTheory.KReal.local_of}
\begin{proof}
Write $\Omega_\rho=\sum\rho(a,b)\rho(b,a)$, the sum over the pairs in the same
block. Let $B\subset A$ be the smallest subspace containing $f$ and stable
under the $\rho(a,b)$ with $b\prec a$.

(i) $B$ is stable under every $\rho(c,d)$ with $c,d$ in the same block. The
$u\in B$ with $\rho(c,d)u\in B$ for all such $(c,d)$ form a subspace. It
contains $f$, since $\rho(c,c)f=\Lambda_cf$, $\rho(c,d)f=0$ for $c\prec d$ and
$\rho(c,d)f\in B$ for $d\prec c$. It is stable under the $\rho(a,b)$ with
$b\prec a$, since
$\rho(c,d)\rho(a,b)u=\rho(a,b)\rho(c,d)u+\delta_{da}\rho(c,b)u-\delta_{bc}\rho(a,d)u$.
So it is $B$.

(ii) If $g\in B$, $\rho(a,b)g=0$ for all $a\prec b$, and $\langle g,f\rangle=0$,
then $g=0$. Indeed, the orthogonal complement of $g$ contains $f$ and is stable
under the $\rho(a,b)$ with $b\prec a$, as
$\langle g,\rho(a,b)x\rangle=\langle\rho(b,a)g,x\rangle=0$; so it contains $B$,
and $g$ is orthogonal to itself.

(iii) If $W\subset B$ is stable under every $\rho(a,b)$ and $f\notin W$, then
$W=0$. By the adjoint relation, $W^\perp$ is stable too. Write $f=w_0+v_0$ with
$w_0\in W$ and $v_0\in W^\perp$. For $a\prec b$,
$\rho(a,b)w_0=-\rho(a,b)v_0$ lies in $W\cap W^\perp$, so both vanish. Put
$\alpha=\|w_0\|^2$, $\beta=\|v_0\|^2$ and $g=\beta w_0-\alpha v_0$, which lies
in $B$ as $v_0=f-w_0$. Then $\rho(a,b)g=0$ for $a\prec b$ and
$\langle g,f\rangle=\beta\alpha-\alpha\beta=0$, so $g=0$ by (ii), and
$\langle w_0,g\rangle=\beta\alpha=0$. If $\beta=0$, then $f=w_0\in W$, which is
excluded. So $\alpha=0$ and $f=v_0\in W^\perp$. Then $B\subset W^\perp$, by the
definition of $B$, and $W\subset B\cap W^\perp=0$.

(iv) $\Omega_\rho$ commutes with every $\rho(c,d)$, by the commutation
relations, and $\Omega_\rho f=c(\Lambda)f$: the terms with $a=b$ give
$\Lambda_a^2f$, those with $a\prec b$ give
$\rho(a,b)\rho(b,a)f=(\Lambda_a-\Lambda_b)f$, and those with $b\prec a$ vanish.
So the kernel of $\Omega_\rho-c(\Lambda)$ contains $B$.

(v) By (i) and (iv), Lemma~\ref{asm-theta-agree-real-quotient-filtration}
applies to $B$ with the restrictions of the $\rho(a,b)$ and of $\iota$ and with
$c_0=c(\Lambda)$; its hypothesis on $\Omega_T$ holds because the span for $B$
lies in the span for $A$. Let $L=\iota(B)$, and $G_L$ the stable subspace
generated by $L$; it contains $G$, as $\iota(f)\in L$. For $v\in B$,
$\Omega_T\iota(v)=\iota(\Omega_\rho v)=c(\Lambda)\iota(v)$. By induction on $k$
and part (1) of that lemma, every $y\in G_L$ with
$(\Omega_T-c(\Lambda))^ky=0$ lies in $L$.

(vi) Let $Z$ be the set of $x\in T$ with $h(\Omega_T)x=0$ for some polynomial
$h$ with $h(c(\Lambda))\neq0$; it is a subspace, as a product of two such
polynomials is one. As $\Omega_T\iota(f)=c(\Lambda)\iota(f)$,
$h(\Omega_T)\iota(f)=h(c(\Lambda))\iota(f)\neq0$ for such $h$; so
$\iota(f)\notin Z$.

(vii) Let $N\subset G$ be a stable subspace with $N\neq G$, and $x\in N$. By
part (2) of Lemma~\ref{asm-theta-agree-real-quotient-filtration}, $x$ lies in a
finite-dimensional subspace $F$ stable under the $T(E_{ab})$ with $a,b$ in the
same block, hence under $\Omega_T$; so is $N$. Write the minimal polynomial of
$\Omega_T$ on $F\cap N$ as $(X-c(\Lambda))^kh$ with $h(c(\Lambda))\neq0$, and
put $y=h(\Omega_T)x\in N$. Then $(\Omega_T-c(\Lambda))^ky=0$ and $y\in G_L$, so
$y=\iota(v)$ with $v\in B$, by (v). Suppose $y\neq0$. The $u\in B$ with
$\iota(u)\in N$ form a subspace stable under every $\rho(a,b)$, which contains
$v\neq0$; by (iii) it contains $f$. So $\iota(f)\in N$ and $G\subset N$, which
is excluded. Hence $y=0$, and $x\in Z$.
\end{proof}

\begin{proposition}\label{asm-theta-agree-real-quotient-common}
Let $R,S\geq0$, let $\Lambda\colon\mathrm{Ix}(R,S)\to\mathbf Q$, and let $T$
and $T'$ be Lie modules of type $(R,S)$ that are both generated with Casimir
gap in weight $\Lambda$
(Definition~\ref{asm-theta-agree-real-quotient-lie-module}). Then there are a
complex vector space $\Pi\neq0$, operators $\varrho(Y)$ on $\Pi$ for
$Y\in M_{R+S}(\mathbf C)$ such that $0$ and $\Pi$ are the only subspaces stable
under every $\varrho(Y)$, and surjective linear maps $\varpi\colon T\to\Pi$
and $\varpi'\colon T'\to\Pi$ with $\varpi T(Y)=\varrho(Y)\varpi$ and
$\varpi'T'(Y)=\varrho(Y)\varpi'$ for every $Y$.
\end{proposition}
\leanref{Development/WeilBound/CoreStatements/GaussianSeminormBounds/RealSeminormFiltration/CoreMain.lean}{PositivityTheory.KReal.common_quot}
\leanref{Development/WeilBound/CoreStatements/GaussianSeminormBounds/RealSeminormFiltration/CoreMain.lean}{PositivityTheory.KReal.quot_of_local}
\begin{proof}
Let $(A,\rho,\iota,f)$ and $(A',\rho',\iota',f')$ be the highest-weight data
of $T$ and $T'$. The space $T\times T'$ with the operators $T(Y)\times T'(Y)$
is a Lie module of type $(R,S)$, and $A\oplus A'$, with the orthogonal sum of
the inner products, the operators $\rho(a,b)\oplus\rho'(a,b)$, the map
$\iota\times\iota'$ and the vector $(f,f')$, is a highest-weight datum of it
of weight $\Lambda$. Its compact Casimir operator is
$\Omega_T\times\Omega_{T'}$, and the two projections map the span of the
vectors $(T(E_{ab})\iota(u),T'(E_{ab})\iota'(u'))$, with $a,b$ in different
blocks, into the corresponding spans for $T$ and $T'$; so the condition on
$\Omega$ in Lemma~\ref{asm-theta-agree-real-quotient-local} holds for
$T\times T'$, componentwise. As $T\neq0$ is generated by $\iota(f)$,
$\iota(f)\neq0$. Let $G$ be the stable subspace of $T\times T'$ generated by
$(\iota(f),\iota'(f'))$, and $Z$ the subspace given by that lemma.

The projections $p\colon G\to T$ and $p'\colon G\to T'$ are surjective, as
their images are stable and contain $\iota(f)$ and $\iota'(f')$. Let $M$ be the
sum of the stable subspaces $N\subset G$ with $N\neq G$. It is stable, and
$M\subset Z$, so $(\iota(f),\iota'(f'))\notin M$ and $M\neq G$. Let $\Pi=G/M$,
with the operators $\varrho(Y)$ induced by the $T(Y)\times T'(Y)$; then
$\Pi\neq0$. The kernels of $p$ and $p'$ are stable, and different from $G$ as
$T$ and $T'$ are nonzero; so they lie in $M$, and $p$ and $p'$ induce
surjective linear maps $\varpi\colon T\to\Pi$ and $\varpi'\colon T'\to\Pi$
with $\varpi p=\varpi'p'$ the quotient map. These satisfy
$\varpi T(Y)=\varrho(Y)\varpi$ and $\varpi'T'(Y)=\varrho(Y)\varpi'$. A
subspace of $\Pi$ stable under every $\varrho(Y)$ is $N/M$ for a stable
subspace $N\supset M$ of $G$; if $N\neq G$, then $N\subset M$. So the only
such subspaces are $0$ and $\Pi$.
\end{proof}

\begin{definition}\label{asm-theta-agree-real-quotient-harmonics}
Let $P,Q,R,S\geq0$, and use the notation of
Definitions~\ref{asm-theta-agree-real-fock}(1)
and~\ref{asm-theta-agree-real-quotient-lie-module}; write
$\mathcal P=\mathcal P(P,Q;R,S)$.
\begin{enumerate}
\item For $d\geq0$, $\mathcal P_d$ is the space of homogeneous polynomials of
  degree $d$, and $f_d$ is the component of degree $d$ of $f\in\mathcal P$. For
  an operator $A$ on $\mathcal P$, its \emph{raising part} and \emph{lowering
  part} are the operators $A^+f=\sum_d(Af_d)_{d+2}$ and
  $A^-f=\sum_d(Af_{d+2})_d$.
\item A matrix $X$ in $M_{P+Q}(\mathbf C)$ or $M_{R+S}(\mathbf C)$ is
  \emph{off-block} if $X_{ab}=0$ whenever $a$ and $b$ lie in the same block.
  The space $H_1$ of \emph{harmonics} is the set of $f\in\mathcal P$ with
  $d_1(X)^-f=0$ for every off-block $X$, and the space $H\subset H_1$ of
  \emph{joint harmonics} is the set of $f\in H_1$ with $d_2(Y)^-f=0$ for every
  off-block $Y$.
\item For $c\in\mathbf C$, $\mathcal P^{(c)}$ is the space of
  $f\in\mathcal P$ with $d_1(X)f=c\operatorname{tr}(X)f$ for every
  block-diagonal $X$.
\item Let $\mu\colon\mathrm{Ix}(R,S)\to\mathbf Q$. A \emph{highest-weight
  vector} of weight $\mu$ for $d_2$ is an $f\neq0$ with
  $d_2(E_{bb})f=\mu_bf$ for every $b$ and $d_2(E_{bb'})f=0$ whenever
  $b\prec b'$. Highest-weight vectors for $d_1$, of weights
  $\mathrm{Ix}(P,Q)\to\mathbf Q$, are defined in the same way.
\item The \emph{Fock inner product} on $\mathcal P$ is the inner product,
  conjugate-linear in the first variable, for which the monomials are
  orthogonal and $\langle x^\nu,x^\nu\rangle=\prod_l\nu_l!$.
\end{enumerate}
\end{definition}
\leanref{Formalism/PeriodIntegrals/JointBlockShapes/MinusPart.lean}{FockH1Action.plusPart}
\leanref{Formalism/PeriodIntegrals/JointBlockShapes/MinusPart.lean}{FockH1Action.minusPart}
\leanref{Formalism/PeriodIntegrals/JointBlockShapes/MinusPart.lean}{FockH1Action.IsPC}
\leanref{Development/WeilBound/CoreStatements/GaussianSeminormBounds/RealSeminormFiltration/CoreDefs.lean}{PositivityTheory.KReal.harm1}
\leanref{Formalism/PeriodIntegrals/JointBlockShapes/JointShape.lean}{FockH1Action.Jset}
\leanref{Formalism/LocalToGlobal/ThetaOperatorCommutation/RealTh.lean}{ReciprocityTheory.ChartRestriction.G.ThK.Real.jHarm}
\leanref{Formalism/LocalToGlobal/ThetaOperatorCommutation/RealTh.lean}{ReciprocityTheory.ChartRestriction.G.ThK.Real.polChi}
\leanref{Formalism/PeriodIntegrals/JointBlockShapes/JointShape.lean}{FockH1Action.IsHW}
\leanref{Formalism/PeriodIntegrals/JointBlockShapes/MinusPart.lean}{FockH1Action.fockInner}

\begin{definition}\label{asm-theta-agree-real-quotient-shapes}
Let $R,S\geq0$, and use the notation of
Definition~\ref{asm-theta-agree-real-quotient-lie-module}.
\begin{enumerate}
\item For $c\in\mathbf Q$ and $\mu\colon\mathrm{Ix}(R,S)\to\mathbf Q$ let
  $\mu^{[c]}$ take the value $\mu_b-c$ at $b$ in the first block and $\mu_b+c$
  at $b$ in the second, and put $\deg_c(\mu)=\sum_b|\mu^{[c]}_b|$ and
  $c_{R,S}=\tfrac12(R-S)$; the same notation is used for maps on
  $\mathrm{Ix}(P,Q)$. For $b\in\mathrm{Ix}(R,S)$, $e_b$ is the map with value
  $1$ at $b$ and $0$ elsewhere.
\item A map $\mu\colon\mathrm{Ix}(R,S)\to\mathbf Q$ is \emph{dominant} if
  $\mu_{b'}\leq\mu_b$ whenever $b\prec b'$. For $P,Q\geq0$ it is
  \emph{sign-bounded} for $(P,Q)$ if $\nu=\mu^{[c_{P,Q}]}$ has $\nu_b\leq0$
  when $b$ is the $i$-th element of the first block with $i>P$ or of the second
  block with $i>Q$, and $\nu_b\geq0$ when $b$ is the $i$-th element of the first
  block with $i\leq R-Q$ or of the second block with $i\leq S-P$.
\item Let $P,Q\geq0$. Maps $\mu\colon\mathrm{Ix}(P,Q)\to\mathbf Q$ and
  $\mu'\colon\mathrm{Ix}(R,S)\to\mathbf Q$ form a \emph{joint shape} if there
  are nonincreasing finite sequences $\alpha,\beta,\gamma,\delta$ of integers,
  the terms of $\alpha$ and $\gamma$ positive and those of $\beta$ and
  $\delta$ negative, with $|\alpha|+|\beta|\leq P$, $|\gamma|+|\delta|\leq Q$,
  $|\alpha|+|\delta|\leq R$ and $|\gamma|+|\beta|\leq S$, such that
  $\mu^{[c_{R,S}]}$ is $(\alpha,0,\ldots,0,\beta)$ on the first block and
  $(\gamma,0,\ldots,0,\delta)$ on the second, and $\mu'^{[c_{P,Q}]}$ is
  $(\alpha,0,\ldots,0,\delta)$ on the first block and
  $(\gamma,0,\ldots,0,\beta)$ on the second. Here $|\alpha|$ is the length of
  $\alpha$, and $(\alpha,0,\ldots,0,\beta)$ is the sequence of the terms of
  $\alpha$, then zeros, then the terms of $\beta$, of the length of the block.
\item Let $P,Q\geq0$ and $c_0,\tau_0\in\mathbf Q$, and write $c_0$ also for
  the constant map with value $c_0$ on $\mathrm{Ix}(P,Q)$. A map
  $\Lambda\colon\mathrm{Ix}(R,S)\to\mathbf Q$ is a \emph{gap weight} for
  $(P,Q,c_0,\tau_0)$ if $c_0$ and $\Lambda$ form a joint shape, $\Lambda$ is
  the only map that forms a joint shape with $c_0$, and
  \[
  c(\Lambda+\tau_0)+R+S<c(\nu+\tau_0)
  \]
  for all $a,b$ in different blocks and every dominant $\nu$, sign-bounded for
  $(P,Q)$, with $\nu=\Lambda-e_a+e_b$ and
  $\deg_{c_{P,Q}}(\nu)=\deg_{c_{R,S}}(c_0)+2$. Here $c$ is the function of
  Definition~\ref{asm-theta-agree-real-quotient-lie-module}, and
  $\Lambda+\tau_0$ is the map $b\mapsto\Lambda_b+\tau_0$.
\end{enumerate}
\end{definition}
\leanref{Formalism/PeriodIntegrals/JointBlockShapes/JointShape.lean}{FockH1Action.shiftNu}
\leanref{Formalism/PeriodIntegrals/JointBlockShapes/JointShape.lean}{FockH1Action.degW}
\leanref{Formalism/PeriodIntegrals/JointBlockShapes/JointShape.lean}{FockH1Action.cPair}
\leanref{Formalism/PeriodIntegrals/JointBlockShapes/JointShape.lean}{FockH1Action.JointShape}
\leanref{Development/WeilBound/CoreStatements/GaussianSeminormBounds/RealSeminormFiltration/CoreDefs.lean}{PositivityTheory.KReal.Dom}
\leanref{Development/WeilBound/CoreStatements/GaussianSeminormBounds/RealSeminormFiltration/CoreDefs.lean}{PositivityTheory.KReal.Supp}
\leanref{Development/WeilBound/CoreStatements/GaussianSeminormBounds/RealSeminormFiltration/CoreDefs.lean}{PositivityTheory.KReal.RootTo}
\leanref{Development/WeilBound/CoreStatements/GaussianSeminormBounds/RealSeminormFiltration/CoreDefs.lean}{PositivityTheory.KReal.CombOK}

\begin{theorem}[Kashiwara, Vergne; Howe \cite{KV78,How89}]\label{asm-theta-agree-real-quotient-fock-facts}
Let $P,Q,R,S\geq0$, and use the notation of
Definitions~\ref{asm-theta-agree-real-fock}(1),
\ref{asm-theta-agree-real-quotient-harmonics}
and~\ref{asm-theta-agree-real-quotient-shapes}, with
$\mathcal P=\mathcal P(P,Q;R,S)$. Let $\mu\colon\mathrm{Ix}(P,Q)\to\mathbf Q$
and $\mu'\colon\mathrm{Ix}(R,S)\to\mathbf Q$.
\begin{enumerate}
\item If $X$ and $Y$ are block-diagonal, then $d_1(X)$ and $d_2(Y)$ preserve
  every $\mathcal P_d$. If $X$ and $Y$ are off-block, then
  $d_1(X)=d_1(X)^++d_1(X)^-$ and $d_2(Y)=d_2(Y)^++d_2(Y)^-$.
\item For every $d\geq0$, $\mathcal P_{d+2}$ is spanned by
  $\mathcal P_{d+2}\cap H_1$ and the $d_1(X)^+g$, and also by
  $\mathcal P_{d+2}\cap H$, the $d_1(X)^+g$ and the $d_2(Y)^+g$, with $X$ and
  $Y$ off-block and $g\in\mathcal P_d$.
\item The space of $f\in H$ that are $0$ or highest-weight vectors of weight
  $\mu$ for $d_1$ and of weight $\mu'$ for $d_2$ has dimension $1$ if $\mu$ and
  $\mu'$ form a joint shape, and is $0$ otherwise.
\item Let $f\neq0$ lie in a finite-dimensional subspace that is stable under
  the $d_1(X)$ with $X$ block-diagonal and in which every highest-weight vector
  for $d_1$ has weight $\mu$, and suppose that $d_2(Y)^-f=0$ for every
  off-block $Y$. Then $f\in\mathcal P_d$ with $d=\deg_{c_{R,S}}(\mu)$.
\item If $f\in\mathcal P_d$ is a highest-weight vector of weight $\mu$ for
  $d_1$, then $\deg_{c_{R,S}}(\mu)\leq d$.
\item If $f\in\mathcal P_e\cap H_1$ is a highest-weight vector of weight
  $\mu'$ for $d_2$, then $e=\deg_{c_{P,Q}}(\mu')$.
\item If $f\in H$ is a highest-weight vector of weight $\mu$ for $d_1$ and of
  weight $\mu'$ for $d_2$, then $f\in\mathcal P_d$ with
  $d=\deg_{c_{R,S}}(\mu)$.
\item Every finite-dimensional subspace $U\neq0$ of $\mathcal P$ that is
  stable under the $d_2(E_{bb'})$ with $b$ and $b'$ in the same block contains
  a highest-weight vector for $d_2$.
\end{enumerate}
\end{theorem}
\leanref{Formalism/PeriodIntegrals/OrbitalBoundsConclusion/FockRaiseLow.lean}{gradedOscillatorProperty}
\leanref{Formalism/PeriodIntegrals/OrbitalBoundsConclusion/FockHarmSplit.lean}{harmonicSplitting_holds}
\leanref{Formalism/PeriodIntegrals/OrbitalBoundsConclusion/FockJointMult.lean}{fock_joint_multiplicity}
\leanref{Formalism/PeriodIntegrals/UnitaryGroupMassFormula/FockJointMultiplicity/FockJointMultOfCauchy.lean}{PI.FJM.shapeOf_of_cauchy}
\leanref{Formalism/PeriodIntegrals/OrbitalBoundsConclusion/FockOneDeg.lean}{fock_degree_one}
\leanref{Formalism/PeriodIntegrals/OrbitalBoundsConclusion/FockWtDeg.lean}{highestWeightDegree_le}
\leanref{Formalism/PeriodIntegrals/UnitaryGroupMassFormula/FockDegreeOne/KroneckerOneRightSwap.lean}{PI.FockOneDegW.Swap.hwHarmDeg2}
\leanref{Formalism/PeriodIntegrals/UnitaryGroupMassFormula/FockJointMultiplicity/MinDeg.lean}{PI.FJMF.minimal_degree_joint}
\leanref{Development/WeilBound/CoreStatements/GaussianSeminormBounds/RealSeminormFiltration/CoreLie.lean}{PositivityTheory.KReal.hwExist2}

\begin{lemma}\label{asm-theta-agree-real-quotient-fock-rules}
Let $P,Q,R,S\geq0$, use the notation of
Definitions~\ref{asm-theta-agree-real-fock}(1)
and~\ref{asm-theta-agree-real-quotient-harmonics}, with
$\mathcal P=\mathcal P(P,Q;R,S)$, and put $\epsilon_{ab}=1$ if $a$ and $b$
lie in the same block and $\epsilon_{ab}=-1$ otherwise.
\begin{enumerate}
\item For $f,g\in\mathcal P$, $a,a'\in\mathrm{Ix}(P,Q)$ and
  $b,b'\in\mathrm{Ix}(R,S)$,
  $\langle d_1(E_{aa'})f,g\rangle=\epsilon_{aa'}\langle f,d_1(E_{a'a})g\rangle$
  and
  $\langle d_2(E_{bb'})f,g\rangle=\epsilon_{bb'}\langle f,d_2(E_{b'b})g\rangle$.
\item Let $A$ be an operator on $\mathcal P$ and $X$ a block-diagonal matrix.
  Then $d_1(X)A^--A^-d_1(X)=(d_1(X)A-Ad_1(X))^-$, and the same holds for $d_2$.
\item The spaces $H_1$, $H$, $\mathcal P_d$ and $\mathcal P^{(c)}$ are stable
  under the $d_2(E_{bb'})$ with $b,b'$ in the same block, and $H_1$, $H$ and
  $\mathcal P_d$ are stable under the $d_1(E_{aa'})$ with $a,a'$ in the same
  block. The homogeneous components of an element of $H$ lie in $H$.
\item Let $b,b'\in\mathrm{Ix}(R,S)$ lie in different blocks, and put
  $\Delta_{bb'}=-d_2(E_{bb'})^-$. Then $\Delta_{bb'}$ maps $\mathcal P_{e+2}$
  into $\mathcal P_e$ and kills $\mathcal P_0$ and $\mathcal P_1$;
  $d_2(E_{bb'})f+\Delta_{bb'}f\in\mathcal P_{e+2}$ for $f\in\mathcal P_e$;
  $\Delta_{bb'}$ maps $H_1$ into $H_1$ and $\mathcal P^{(c)}$ into
  $\mathcal P^{(c)}$; and for $a,a'$ in the same block,
  $d_2(E_{aa'})\Delta_{bb'}-\Delta_{bb'}d_2(E_{aa'})
  =\delta_{a'b}\Delta_{ab'}-\delta_{b'a}\Delta_{ba'}$. An $f\in H_1$ lies in
  $H$ if and only if $\Delta_{bb'}f=0$ for all $b,b'$ in different blocks.
\item Let $\mu\colon\mathrm{Ix}(P,Q)\to\mathbf Q$ be constant on each block, and
  $f$ a highest-weight vector of weight $\mu$ for $d_1$. Then
  $d_1(E_{aa'})f=\delta_{aa'}\mu_af$ for all $a,a'$ in the same block. In
  particular, for $c_0\in\mathbf Q$, $f\in\mathcal P^{(c_0)}$ if and only if
  $f=0$ or $f$ is a highest-weight vector for $d_1$ of the constant weight
  $c_0$.
\item If $f$ is a highest-weight vector of weight $\nu$ for $d_2$, then $\nu$
  is dominant and sign-bounded for $(P,Q)$
  (Definition~\ref{asm-theta-agree-real-quotient-shapes}(2)).
\end{enumerate}
\end{lemma}
\leanref{Development/WeilBound/CoreStatements/GaussianSeminormBounds/RealSeminormFiltration/CoreLie.lean}{PositivityTheory.KReal.adj1}
\leanref{Development/WeilBound/CoreStatements/GaussianSeminormBounds/RealSeminormFiltration/CoreLie.lean}{PositivityTheory.KReal.adj2}
\leanref{Development/WeilBound/CoreStatements/GaussianSeminormBounds/RealSeminormFiltration/CoreHarm.lean}{PositivityTheory.KReal.Harm.mp_comm}
\leanref{Development/WeilBound/CoreStatements/GaussianSeminormBounds/RealSeminormFiltration/CoreHarm.lean}{PositivityTheory.KReal.harmonic_stable_second_action}
\leanref{Development/WeilBound/CoreStatements/GaussianSeminormBounds/RealSeminormFiltration/CoreHarm.lean}{PositivityTheory.KReal.jHarmonic_stable_second_action}
\leanref{Development/WeilBound/CoreStatements/GaussianSeminormBounds/RealSeminormFiltration/CoreHarm.lean}{PositivityTheory.KReal.homogeneous_stable_second_action}
\leanref{Development/WeilBound/CoreStatements/GaussianSeminormBounds/RealSeminormFiltration/CoreHarm.lean}{PositivityTheory.KReal.polChi_stable_second_action}
\leanref{Development/WeilBound/CoreStatements/GaussianSeminormBounds/RealSeminormFiltration/CoreHarm.lean}{PositivityTheory.KReal.harmonic_stable_first_action}
\leanref{Development/WeilBound/CoreStatements/GaussianSeminormBounds/RealSeminormFiltration/CoreHarm.lean}{PositivityTheory.KReal.jHarmonic_stable_first_action}
\leanref{Development/WeilBound/CoreStatements/GaussianSeminormBounds/RealSeminormFiltration/CoreHarm.lean}{PositivityTheory.KReal.homogeneous_stable_first_action}
\leanref{Formalism/PeriodIntegrals/UnitaryGroupMassFormula/FockJointMultiplicity/MinDeg.lean}{PI.FJMF.jset_hc}
\leanref{Development/WeilBound/CoreStatements/GaussianSeminormBounds/RealSeminormFiltration/CoreHarm.lean}{PositivityTheory.KReal.Harm.minusPart_second_action_Eab}
\leanref{Development/WeilBound/CoreStatements/GaussianSeminormBounds/RealSeminormFiltration/CoreHarm.lean}{PositivityTheory.KReal.low_homog}
\leanref{Development/WeilBound/CoreStatements/GaussianSeminormBounds/RealSeminormFiltration/CoreHarm.lean}{PositivityTheory.KReal.low_homog_small}
\leanref{Development/WeilBound/CoreStatements/GaussianSeminormBounds/RealSeminormFiltration/CoreHarm.lean}{PositivityTheory.KReal.second_action_add_lowering}
\leanref{Development/WeilBound/CoreStatements/GaussianSeminormBounds/RealSeminormFiltration/CoreHarm.lean}{PositivityTheory.KReal.harmonic_stable_lowering}
\leanref{Development/WeilBound/CoreStatements/GaussianSeminormBounds/RealSeminormFiltration/CoreHarm.lean}{PositivityTheory.KReal.polChi_low}
\leanref{Development/WeilBound/CoreStatements/GaussianSeminormBounds/RealSeminormFiltration/CoreHarm.lean}{PositivityTheory.KReal.low_comm}
\leanref{Development/WeilBound/CoreStatements/GaussianSeminormBounds/RealSeminormFiltration/CoreHarm.lean}{PositivityTheory.KReal.mem_jHarm_iff}
\leanref{Development/WeilBound/CoreStatements/GaussianSeminormBounds/RealSeminormFiltration/CoreLie.lean}{PositivityTheory.KReal.scalar_of_hw}
\leanref{Development/WeilBound/CoreStatements/GaussianSeminormBounds/RealSeminormFiltration/CoreLie.lean}{PositivityTheory.KReal.mem_polChi_iff_hw}
\leanref{Development/WeilBound/CoreStatements/GaussianSeminormBounds/RealSeminormFiltration/CoreLie.lean}{PositivityTheory.KReal.dom2}
\leanref{Development/WeilBound/CoreStatements/GaussianSeminormBounds/RealSeminormFiltration/CoreSupp.lean}{PositivityTheory.KReal.supp2}
\begin{proof}
Parts (1) to (4), and the sign bound in (6), are computations with the
operators $L_l$ and $R_l$ and with the monomials of $f$, and their results are
stated; (2) uses that $d_1(X)$ and $d_2(X)$ commute with taking homogeneous
components (Theorem~\ref{asm-theta-agree-real-quotient-fock-facts}(1)), and
the commutation relation in (4) is (2) for $A=d_2(E_{bb'})$, with the bracket
relation of Lemma~\ref{asm-theta-agree-real-fock-relations}(1).

(5) Let $a\prec a'$. By the bracket relation and $d_1(E_{aa'})f=0$,
$d_1(E_{aa'})d_1(E_{a'a})f=(\mu_a-\mu_{a'})f=0$, so by (1)
$\langle d_1(E_{a'a})f,d_1(E_{a'a})f\rangle=\langle f,d_1(E_{aa'})d_1(E_{a'a})f\rangle=0$,
and $d_1(E_{a'a})f=0$. With $d_1(E_{aa})f=\mu_af$ and $d_1(E_{aa'})f=0$ this
gives the first claim. For the second, an $f\in\mathcal P^{(c_0)}$ has
$d_1(E_{aa'})f=\delta_{aa'}c_0f$ for $a,a'$ in the same block, and conversely
the first claim gives these equations, which give
$d_1(X)f=c_0\operatorname{tr}(X)f$ for block-diagonal $X$ by linearity.

(6) For $b\prec b'$, as in (5),
$\langle d_2(E_{b'b})f,d_2(E_{b'b})f\rangle=(\nu_b-\nu_{b'})\langle f,f\rangle$;
so $\nu_{b'}\leq\nu_b$, and $\nu$ is dominant.
\end{proof}

\begin{lemma}\label{asm-theta-agree-real-quotient-split}
Let $P,Q,R,S\geq0$ and $c_0\in\mathbf Q$, use the notation of
Definitions~\ref{asm-theta-agree-real-quotient-harmonics}
and~\ref{asm-theta-agree-real-quotient-shapes}, with
$\mathcal P=\mathcal P(P,Q;R,S)$, and let $\pi\colon\mathcal P\to\Theta$ be the
quotient map by the span of the $d_1(X)f-c_0\operatorname{tr}(X)f$, for all
$X\in M_{P+Q}(\mathbf C)$ and $f\in\mathcal P$. Put
$D_0=\deg_{c_{R,S}}(c_0)$, for the constant map $c_0$ on $\mathrm{Ix}(P,Q)$.
\begin{enumerate}
\item If $W\subset\mathcal P_d$ is stable under the $d_1(E_{aa'})$ with
  $a,a'$ in the same block, then $\pi(W)=\pi(W\cap\mathcal P^{(c_0)})$.
\item If $d<D_0$, then $\mathcal P_d\cap\mathcal P^{(c_0)}=0$ and
  $\pi(\mathcal P_d)=0$.
\item For off-block $X$ and $g\in\mathcal P$,
  $\pi(d_1(X)^+g)=-\pi(d_1(X)^-g)$.
\item If $D_0$ is an integer, then
  $\pi(\mathcal P_{D_0+2})=\pi(H_1\cap\mathcal P^{(c_0)}\cap\mathcal P_{D_0+2})$.
\end{enumerate}
\end{lemma}
\leanref{Development/WeilBound/CoreStatements/GaussianSeminormBounds/RealSeminormFiltration/CoreSplit.lean}{PositivityTheory.KReal.charSplit}
\leanref{Development/WeilBound/CoreStatements/GaussianSeminormBounds/RealSeminormFiltration/CoreSplit.lean}{PositivityTheory.KReal.polChi_small}
\leanref{Development/WeilBound/CoreStatements/GaussianSeminormBounds/RealSeminormFiltration/CoreSplit.lean}{PositivityTheory.KReal.th_small}
\leanref{Development/WeilBound/CoreStatements/GaussianSeminormBounds/RealSeminormFiltration/CoreSplit.lean}{PositivityTheory.KReal.th_plus1}
\leanref{Development/WeilBound/CoreStatements/GaussianSeminormBounds/RealSeminormFiltration/CoreSplit.lean}{PositivityTheory.KReal.th_next}
\begin{proof}
(1) Let $x\in W$. For $a,a'$ in the same block put
$D_{aa'}=d_1(E_{aa'})-\delta_{aa'}c_0$. It maps $W$ into $W$, and $\pi$ kills
its image, as $\operatorname{tr}(E_{aa'})=\delta_{aa'}$. Let $K\subset W$ be
the span of the $D_{aa'}(W)$, and $y$ the orthogonal projection of $x$ on $K$
for the Fock inner product, in the finite-dimensional space $\mathcal P_d$.
Then $\pi(y)=0$, $x-y\in W$, and $x-y$ is orthogonal to $K$. For $a,a'$ in the
same block, $D_{a'a}D_{aa'}(x-y)\in K$, so by
Lemma~\ref{asm-theta-agree-real-quotient-fock-rules}(1)
\[
\langle D_{aa'}(x-y),D_{aa'}(x-y)\rangle=\langle D_{a'a}D_{aa'}(x-y),x-y\rangle=0,
\]
and $D_{aa'}(x-y)=0$. By linearity $x-y\in\mathcal P^{(c_0)}$, and
$\pi(x)=\pi(x-y)$.

(2) A nonzero $g\in\mathcal P_d\cap\mathcal P^{(c_0)}$ is a highest-weight
vector for $d_1$ of the constant weight $c_0$
(Lemma~\ref{asm-theta-agree-real-quotient-fock-rules}(5)), so $D_0\leq d$ by
Theorem~\ref{asm-theta-agree-real-quotient-fock-facts}(5). Then (1), for
$W=\mathcal P_d$, which is stable by
Lemma~\ref{asm-theta-agree-real-quotient-fock-rules}(3), gives
$\pi(\mathcal P_d)=0$.

(3) As $X$ is off-block, $\operatorname{tr}(X)=0$, so $\pi(d_1(X)g)=0$; and
$d_1(X)g=d_1(X)^+g+d_1(X)^-g$
(Theorem~\ref{asm-theta-agree-real-quotient-fock-facts}(1)).

(4) By Theorem~\ref{asm-theta-agree-real-quotient-fock-facts}(2),
$\mathcal P_{D_0+2}$ is spanned by $\mathcal P_{D_0+2}\cap H_1$ and the
$d_1(X)^+g$ with $X$ off-block and $g\in\mathcal P_{D_0}$. By (3),
$\pi(d_1(X)^+g)=-\pi(d_1(X)^-g)$. If $D_0<2$, then $d_1(X)^-g=0$; otherwise
$d_1(X)^-g$ is the component of $d_1(X)g$ of degree $D_0-2$, which $\pi$ kills
by (2). So $\pi(\mathcal P_{D_0+2})=\pi(\mathcal P_{D_0+2}\cap H_1)$, and (1),
for $W=\mathcal P_{D_0+2}\cap H_1$, which is stable by
Lemma~\ref{asm-theta-agree-real-quotient-fock-rules}(3), gives the claim.
\end{proof}

\begin{lemma}\label{asm-theta-agree-real-quotient-generator}
Let $P,Q,R,S\geq0$, $c_0\in\mathbf Q$ and
$\Lambda\colon\mathrm{Ix}(R,S)\to\mathbf Q$, use the notation of
Definitions~\ref{asm-theta-agree-real-quotient-harmonics}
and~\ref{asm-theta-agree-real-quotient-shapes}, with
$\mathcal P=\mathcal P(P,Q;R,S)$, and suppose that the constant map $c_0$ on
$\mathrm{Ix}(P,Q)$ and $\Lambda$ form a joint shape and that $\Lambda$ is the
only map that forms a joint shape with $c_0$. Then $D_0=\deg_{c_{R,S}}(c_0)$
is an integer, $H\cap\mathcal P^{(c_0)}\subset\mathcal P_{D_0}$, and there is
$f_0\in H\cap\mathcal P^{(c_0)}$, a highest-weight vector of weight $\Lambda$
for $d_2$, such that every subspace $U\subset\mathcal P$ that contains $f_0$
and is stable under the $d_2(E_{bb'})$ with $b,b'$ in the same block contains
$H\cap\mathcal P^{(c_0)}$.
\end{lemma}
\leanref{Development/WeilBound/CoreStatements/GaussianSeminormBounds/RealSeminormFiltration/CoreGen.lean}{PositivityTheory.KReal.exists_gen}
\leanref{Development/WeilBound/CoreStatements/GaussianSeminormBounds/RealSeminormFiltration/CoreGen.lean}{PositivityTheory.KReal.gen_oneDeg}
\leanref{Development/WeilBound/CoreStatements/GaussianSeminormBounds/RealSeminormFiltration/CoreGen.lean}{PositivityTheory.KReal.jHarm_polChi_homog}
\begin{proof}
By Theorem~\ref{asm-theta-agree-real-quotient-fock-facts}(3), the $f\in H$
that are $0$ or highest-weight vectors of weight $c_0$ for $d_1$ and $\Lambda$
for $d_2$ form a line; let $f_0$ span it. By
Lemma~\ref{asm-theta-agree-real-quotient-fock-rules}(5),
$f_0\in\mathcal P^{(c_0)}$.

\emph{Degree.} Let $g\in H\cap\mathcal P^{(c_0)}$ be nonzero. The line
$\mathbf Cg$ is stable under the $d_1(X)$ with $X$ block-diagonal, as
$d_1(X)g=c_0\operatorname{tr}(X)g$, and a highest-weight vector for $d_1$ in it
is a multiple of $g$, so has the constant weight $c_0$. As $g\in H$,
$d_2(Y)^-g=0$ for off-block $Y$. By
Theorem~\ref{asm-theta-agree-real-quotient-fock-facts}(4),
$g\in\mathcal P_d$ with $d=D_0$. For $g=f_0$ this shows that $D_0$ is an
integer, and it shows $H\cap\mathcal P^{(c_0)}\subset\mathcal P_{D_0}$.

\emph{Generation.} Let $U$ be as in the statement, and put
$V=H\cap\mathcal P^{(c_0)}$, a subspace of the finite-dimensional space
$\mathcal P_{D_0}$ that is stable under the $d_2(E_{bb'})$ with $b,b'$ in the
same block (Lemma~\ref{asm-theta-agree-real-quotient-fock-rules}(3)). Let $Z$
be the space of $z\in V$ orthogonal to $U\cap V$. For $z\in Z$ and
$x\in U\cap V$,
$\langle x,d_2(E_{bb'})z\rangle=\langle d_2(E_{b'b})x,z\rangle=0$, by
Lemma~\ref{asm-theta-agree-real-quotient-fock-rules}(1), as
$d_2(E_{b'b})x\in U\cap V$; so $Z$ is stable under these operators. Suppose
$Z\neq0$. By Theorem~\ref{asm-theta-agree-real-quotient-fock-facts}(8), $Z$
contains a highest-weight vector $g$ for $d_2$, of some weight $\nu$. As
$g\in V$, it is a highest-weight vector of weight $c_0$ for $d_1$ in $H$, so by
Theorem~\ref{asm-theta-agree-real-quotient-fock-facts}(3) $c_0$ and $\nu$ form
a joint shape, $\nu=\Lambda$, and $g$ is a multiple of $f_0$. Then
$g\in U\cap V$ and $g\in Z$, so $\langle g,g\rangle=0$, which is false. So
$Z=0$. For $v\in V$ let $y$ be the orthogonal projection of $v$ on $U\cap V$;
then $v-y\in Z$, so $v=y\in U$.
\end{proof}

\begin{lemma}\label{asm-theta-agree-real-quotient-generation}
Let $P,Q,R,S\geq0$, $c_0\in\mathbf Q$ and $\tau\in\mathbf C$. Use the notation
of Lemma~\ref{asm-theta-agree-real-quotient-split}, and let $\vartheta(Y)$ be the
operator on $\Theta$ induced by $d_2(Y)+\tau\operatorname{tr}(Y)$
(Lemma~\ref{asm-theta-agree-real-fock-relations}(1)). Let $N\subset\Theta$ be a
subspace stable under every $\vartheta(Y)$ that contains
$\pi(H\cap\mathcal P^{(c_0)})$. Then $N=\Theta$.
\end{lemma}
\leanref{Development/WeilBound/CoreStatements/GaussianSeminormBounds/RealSeminormFiltration/CoreGen.lean}{PositivityTheory.KReal.th_gen_harm}
\begin{proof}
First, $\pi(\mathcal P_d\cap H)\subset N$ for every $d$: by
Lemma~\ref{asm-theta-agree-real-quotient-split}(1), for
$W=\mathcal P_d\cap H$, which is stable by
Lemma~\ref{asm-theta-agree-real-quotient-fock-rules}(3),
$\pi(W)=\pi(W\cap\mathcal P^{(c_0)})\subset N$.

We show $\pi(\mathcal P_d)\subset N$ by induction on $d$. If $d<2$, every
lowering part vanishes on $\mathcal P_d$, so $\mathcal P_d\subset H$. Let
$d=e+2$. For an operator $A$ and $g\in\mathcal P_e$, $\pi(A^-g)\in N$: it is
$0$ if $e<2$, and otherwise $A^-g$ is the component of $Ag$ of degree $e-2$, to
which the induction hypothesis applies. By
Theorem~\ref{asm-theta-agree-real-quotient-fock-facts}(2), $\mathcal P_{e+2}$
is spanned by $\mathcal P_{e+2}\cap H$, the $d_1(X)^+g$ and the $d_2(Y)^+g$,
with $X,Y$ off-block and $g\in\mathcal P_e$. The first lies in $N$ by the
first paragraph. By Lemma~\ref{asm-theta-agree-real-quotient-split}(3),
$\pi(d_1(X)^+g)=-\pi(d_1(X)^-g)\in N$. By
Theorem~\ref{asm-theta-agree-real-quotient-fock-facts}(1),
\[
d_2(Y)^+g=\bigl(d_2(Y)g+\tau\operatorname{tr}(Y)g\bigr)-\tau\operatorname{tr}(Y)g-d_2(Y)^-g,
\]
and $\pi$ maps it to
$\vartheta(Y)\pi(g)-\tau\operatorname{tr}(Y)\pi(g)-\pi(d_2(Y)^-g)\in N$, as
$\pi(g)\in N$ by induction. So $\pi(\mathcal P_{e+2})\subset N$. As every
element of $\Theta$ is $\pi(f)=\sum_d\pi(f_d)$, $N=\Theta$.
\end{proof}

\begin{lemma}\label{asm-theta-agree-real-quotient-root}
Let $P,Q,R,S\geq0$, $c_0\in\mathbf Q$ and
$\Lambda\colon\mathrm{Ix}(R,S)\to\mathbf Q$, use the notation of
Definitions~\ref{asm-theta-agree-real-quotient-harmonics}
and~\ref{asm-theta-agree-real-quotient-shapes}, with
$\mathcal P=\mathcal P(P,Q;R,S)$, and suppose that $\Lambda$ is the only map
that forms a joint shape with the constant map $c_0$ on $\mathrm{Ix}(P,Q)$. Let
$D_0\geq0$ be an integer with $D_0=\deg_{c_{R,S}}(c_0)$, and let
$g\in H_1\cap\mathcal P^{(c_0)}\cap\mathcal P_{D_0+2}$ be a highest-weight
vector of weight $\nu$ for $d_2$. Then $\deg_{c_{P,Q}}(\nu)=D_0+2$, and
$\nu=\Lambda-e_a+e_b$ for some $a,b$ in different blocks.
\end{lemma}
\leanref{Development/WeilBound/CoreStatements/GaussianSeminormBounds/RealSeminormFiltration/CoreRoot.lean}{PositivityTheory.KReal.root_of_hw}
\begin{proof}
By Theorem~\ref{asm-theta-agree-real-quotient-fock-facts}(6),
$D_0+2=\deg_{c_{P,Q}}(\nu)$. Use the operators $\Delta_{bb'}$ of
Lemma~\ref{asm-theta-agree-real-quotient-fock-rules}(4). For
$h\in\mathcal P^{(c_0)}\cap\mathcal P_{D_0}$ and $b,b'$ in different blocks,
$\Delta_{bb'}h=0$: if $D_0<2$ by that lemma, and otherwise because
$\Delta_{bb'}h\in\mathcal P^{(c_0)}\cap\mathcal P_{D_0-2}$, which is $0$ by
Lemma~\ref{asm-theta-agree-real-quotient-split}(2).

Number each block $1,2,\ldots$ and let $t(b)=-i$ for $b$ the $i$-th element of
its block, so that $t(b')<t(b)$ when $b\prec b'$. Suppose first that
$\Delta_{bb'}g\neq0$ for some $b,b'$ in different blocks, and choose such a
pair $(c,d)$ with $t(c)-t(d)$ largest. Put $g'=\Delta_{cd}g$. By
Lemma~\ref{asm-theta-agree-real-quotient-fock-rules}(4),
$g'\in H_1\cap\mathcal P^{(c_0)}\cap\mathcal P_{D_0}$, so
$\Delta_{bb'}g'=0$ for all $b,b'$ in different blocks by the first paragraph,
and $g'\in H$. By the commutation relation of that lemma,
$d_2(E_{xx})g'=(\nu_x+\delta_{xc}-\delta_{xd})g'$ for every $x$. For
$x\prec y$,
$d_2(E_{xy})g'=\delta_{yc}\Delta_{xd}g-\delta_{dx}\Delta_{cy}g$. If $y=c$, then
$x$ and $d$ lie in different blocks and $t(x)-t(d)>t(c)-t(d)$, so
$\Delta_{xd}g=0$ by the choice of $(c,d)$; if $d=x$, then $c$ and $y$ lie in
different blocks and $t(c)-t(y)>t(c)-t(d)$, so $\Delta_{cy}g=0$. So $g'$ is a
highest-weight vector of weight $\nu+e_c-e_d$ for $d_2$, and of the constant
weight $c_0$ for $d_1$
(Lemma~\ref{asm-theta-agree-real-quotient-fock-rules}(5)). By
Theorem~\ref{asm-theta-agree-real-quotient-fock-facts}(3), $c_0$ and
$\nu+e_c-e_d$ form a joint shape, so $\nu+e_c-e_d=\Lambda$; take $a=c$ and
$b=d$.

Otherwise $\Delta_{bb'}g=0$ for all $b,b'$ in different blocks, so $g\in H$ by
Lemma~\ref{asm-theta-agree-real-quotient-fock-rules}(4), and $g$ is a
highest-weight vector of weight $c_0$ for $d_1$. By
Theorem~\ref{asm-theta-agree-real-quotient-fock-facts}(7),
$g\in\mathcal P_{D_0}$. As $g\neq0$ lies in $\mathcal P_{D_0+2}$, this is
impossible.
\end{proof}

\begin{lemma}\label{asm-theta-agree-real-quotient-gap}
Let $P,Q,R,S\geq0$, $c_0,\tau_0\in\mathbf Q$, and let
$\Lambda\colon\mathrm{Ix}(R,S)\to\mathbf Q$ be a gap weight for
$(P,Q,c_0,\tau_0)$ (Definition~\ref{asm-theta-agree-real-quotient-shapes}(4)).
Use the notation of Lemma~\ref{asm-theta-agree-real-quotient-split}, let
$\vartheta(Y)$ be the operator on $\Theta$ induced by
$d_2(Y)+\tau_0\operatorname{tr}(Y)$, and let $D_0\geq0$ be an integer with
$D_0=\deg_{c_{R,S}}(c_0)$. Let $\Omega$ be the compact Casimir operator of
$\Theta$ with the $\vartheta(Y)$
(Definition~\ref{asm-theta-agree-real-quotient-lie-module}(1)). Then for every
real $t\leq c(\Lambda+\tau_0)+R+S$, the only $x$ with $\Omega x=tx$ in the span
of the vectors $\vartheta(E_{bb'})\pi(u)$, with $b,b'$ in different blocks and
$u\in\mathcal P_{D_0}$, is $0$.
\end{lemma}
\leanref{Development/WeilBound/CoreStatements/GaussianSeminormBounds/RealSeminormFiltration/CoreStar.lean}{PositivityTheory.KReal.star_th}
\leanref{Development/WeilBound/CoreStatements/GaussianSeminormBounds/RealSeminormFiltration/CoreStar.lean}{PositivityTheory.KReal.star_inj_quot}
\leanref{Development/WeilBound/CoreStatements/GaussianSeminormBounds/RealSeminormFiltration/CoreStar.lean}{PositivityTheory.KReal.star_cas_mk}
\begin{proof}
Let $x$ be in that span, with $\Omega x=tx$.

\emph{Step 1: $x\in\pi(\mathcal P_{D_0+2})$.} Let $b,b'$ lie in different
blocks and $u\in\mathcal P_{D_0}$. Then $\pi(\Delta_{bb'}u)=0$, with
$\Delta_{bb'}$ as in Lemma~\ref{asm-theta-agree-real-quotient-fock-rules}(4):
$\Delta_{bb'}u=0$ if $D_0<2$, and otherwise
$\Delta_{bb'}u\in\mathcal P_{D_0-2}$, which $\pi$ kills by
Lemma~\ref{asm-theta-agree-real-quotient-split}(2). As
$\operatorname{tr}(E_{bb'})=0$,
$\vartheta(E_{bb'})\pi(u)=\pi(d_2(E_{bb'})u+\Delta_{bb'}u)$, which lies in
$\pi(\mathcal P_{D_0+2})$ by that lemma.

\emph{Step 2.} By Lemma~\ref{asm-theta-agree-real-quotient-split}(4),
$x=\pi(y)$ with $y\in Y=H_1\cap\mathcal P^{(c_0)}\cap\mathcal P_{D_0+2}$, a
finite-dimensional space. For $b,b'$ in the same block put
$\rho(b,b')=d_2(E_{bb'})+\delta_{bb'}\tau_0$; these operators preserve $Y$
(Lemma~\ref{asm-theta-agree-real-quotient-fock-rules}(3)), satisfy
$\pi\rho(b,b')=\vartheta(E_{bb'})\pi$, and satisfy the bracket relations of the
matrix units, by Lemma~\ref{asm-theta-agree-real-fock-relations}(1). Put
$\Omega_\rho=\sum\rho(b,b')\rho(b',b)$, the sum over the pairs in the same block,
and $D=\Omega_\rho-t$. Then $D$ preserves $Y$, $\pi D=(\Omega-t)\pi$, and
$\pi(Dy)=0$.

\emph{Step 3: $D$ is injective on $Y$.} Suppose $U=Y\cap\ker D\neq0$. As in step
(iv) of the proof of Lemma~\ref{asm-theta-agree-real-quotient-local},
$\Omega_\rho$ commutes with every $\rho(b,b')$, so $U$ is stable under the
$\rho(b,b')$, hence under the $d_2(E_{bb'})$ with $b,b'$ in the same block. By
Theorem~\ref{asm-theta-agree-real-quotient-fock-facts}(8), $U$ contains a
highest-weight vector $g$ for $d_2$, of some weight $\nu$. By
Lemma~\ref{asm-theta-agree-real-quotient-root},
$\deg_{c_{P,Q}}(\nu)=D_0+2$ and $\nu=\Lambda-e_a+e_b$ with $a,b$ in different
blocks, and $\nu$ is dominant and sign-bounded for $(P,Q)$ by
Lemma~\ref{asm-theta-agree-real-quotient-fock-rules}(6). As $\Lambda$ is a gap
weight, $c(\Lambda+\tau_0)+R+S<c(\nu+\tau_0)$. For the $\rho(b,b')$, $g$ is a
vector of highest weight $\nu+\tau_0$, so
$\Omega_\rho g=c(\nu+\tau_0)g$, by the computation in step (iv) of the proof of
Lemma~\ref{asm-theta-agree-real-quotient-local}. As $Dg=0$,
$c(\nu+\tau_0)=t\leq c(\Lambda+\tau_0)+R+S$, a contradiction.

\emph{Step 4.} Let $Z=Y\cap\ker\pi$. As $\pi D=(\Omega-t)\pi$, $D$ maps $Z$ into
$Z$; it is injective there by Step 3, so surjective, as $Z$ is
finite-dimensional. As $Dy\in Z$, $Dy=Dz$ for some $z\in Z$, and $y=z$ by
Step 3. So $x=\pi(y)=0$.
\end{proof}

\begin{proposition}\label{asm-theta-agree-real-quotient-side}
Let $P,Q,R,S\geq0$, $c_0,\tau_0\in\mathbf Q$, and let
$\Lambda\colon\mathrm{Ix}(R,S)\to\mathbf Q$ be a gap weight for
$(P,Q,c_0,\tau_0)$ (Definition~\ref{asm-theta-agree-real-quotient-shapes}(4)).
Let $\Theta$ be the quotient of $\mathcal P(P,Q;R,S)$ by the span of the
polynomials $d_1(X)f-c_0\operatorname{tr}(X)f$, for all $X$ and $f$, and let
$\vartheta(Y)$ be the operator on $\Theta$ induced by
$d_2(Y)+\tau_0\operatorname{tr}(Y)$. If $\Theta\neq0$, then $\Theta$ with the
operators $\vartheta(Y)$ is a Lie module of type $(R,S)$ generated with Casimir
gap in weight $\Lambda+\tau_0$
(Definition~\ref{asm-theta-agree-real-quotient-lie-module}).
\end{proposition}
\leanref{Development/WeilBound/CoreStatements/GaussianSeminormBounds/RealSeminormFiltration/CoreSide.lean}{PositivityTheory.KReal.side_of}
\leanref{Development/WeilBound/CoreStatements/GaussianSeminormBounds/RealSeminormFiltration/CoreSide.lean}{PositivityTheory.KReal.GoodSideLattices.sd_rhoA_adj}
\leanref{Development/WeilBound/CoreStatements/GaussianSeminormBounds/RealSeminormFiltration/CoreSide.lean}{PositivityTheory.KReal.GoodSideLattices.sd_rhoA_gl}
\leanref{Development/WeilBound/CoreStatements/GaussianSeminormBounds/RealSeminormFiltration/CoreLie.lean}{PositivityTheory.KReal.lie_th}
\begin{proof}
By Lemma~\ref{asm-theta-agree-real-fock-relations}(1), $Y\mapsto d_2(Y)$ is
linear, satisfies the bracket relation and commutes with every $d_1(X)$, and
$\operatorname{tr}(YY'-Y'Y)=0$; so $\Theta$ with the $\vartheta(Y)$ is a Lie
module of type $(R,S)$. Let $\pi$ be the quotient map, $D_0$ and $f_0$ as in
Lemma~\ref{asm-theta-agree-real-quotient-generator}, and let $A=\mathcal P_{D_0}$,
with the Fock inner product. For $b,b'$ in the same block let $\rho(b,b')$ be
the restriction of $d_2(E_{bb'})+\delta_{bb'}\tau_0$ to $A$
(Lemma~\ref{asm-theta-agree-real-quotient-fock-rules}(3)), let
$\iota=\pi|_A$, and $f=f_0$. These form a highest-weight datum of weight
$\Lambda+\tau_0$: the adjoint relation holds by
Lemma~\ref{asm-theta-agree-real-quotient-fock-rules}(1), as $\tau_0$ is real;
the bracket relation holds by Lemma~\ref{asm-theta-agree-real-fock-relations}(1),
as $\delta_{bc}\delta_{ad}-\delta_{da}\delta_{cb}=0$; $\rho(b,b)f_0=(\Lambda_b+\tau_0)f_0$
and $\rho(b,b')f_0=0$ for $b\prec b'$; and
$\iota\rho(b,b')=\vartheta(E_{bb'})\iota$, as $\operatorname{tr}(E_{bb'})=\delta_{bb'}$.

\emph{Generation.} Let $N$ be the stable subspace generated by $\iota(f_0)$,
and $U=\pi^{-1}(N)$. Then $f_0\in U$, and $U$ is stable under the
$d_2(E_{bb'})$ with $b,b'$ in the same block, as
$\pi(d_2(E_{bb'})u)=\vartheta(E_{bb'})\pi(u)-\delta_{bb'}\tau_0\pi(u)$. By
Lemma~\ref{asm-theta-agree-real-quotient-generator}, $U$ contains
$H\cap\mathcal P^{(c_0)}$, so $N$ contains $\pi(H\cap\mathcal P^{(c_0)})$, and
$N=\Theta$ by Lemma~\ref{asm-theta-agree-real-quotient-generation}.

\emph{The gap} is Lemma~\ref{asm-theta-agree-real-quotient-gap}, as
$\iota(A)=\pi(\mathcal P_{D_0})$. With $\Theta\neq0$, $\Theta$ is generated
with Casimir gap in weight $\Lambda+\tau_0$.
\end{proof}

\begin{lemma}\label{asm-theta-agree-real-quotient-weights}
Let $P,Q,R,S\geq0$ and $n,m>0$ be integers, and use the notation of
Definition~\ref{asm-theta-agree-real-quotient-shapes}.
\begin{enumerate}
\item Suppose that $P\leq R$, $Q\leq S$, $n+2c_{P,Q}>0$ and $m-2c_{P,Q}>0$, and
  let $c_0=n+c_{R,S}=m-c_{R,S}$. Let $\Lambda^+$ take the value $n+c_{P,Q}$ at
  the first $P$ elements of the first block of $\mathrm{Ix}(R,S)$, $c_{P,Q}$ at
  its other elements, $m-c_{P,Q}$ at the first $Q$ elements of the second
  block and $-c_{P,Q}$ at its other elements. Then $\Lambda^+$ is a gap weight
  for $(P,Q,c_0,\tau)$ for every $\tau\in\mathbf Q$.
\item Suppose that $Q\leq R$, $P\leq S$, $m+2c_{P,Q}>0$ and $n-2c_{P,Q}>0$,
  and let $c_0=-m+c_{R,S}=-n-c_{R,S}$. Let $\Lambda^-$ take the value
  $-n+c_{P,Q}$ at the last $Q$ elements of the first block, $c_{P,Q}$ at its
  other elements, $-m-c_{P,Q}$ at the last $P$ elements of the second block and
  $-c_{P,Q}$ at its other elements. Then $\Lambda^-$ is a gap weight for
  $(P,Q,c_0,\tau)$ for every $\tau\in\mathbf Q$.
\end{enumerate}
\end{lemma}
\leanref{Development/WeilBound/CoreStatements/GaussianSeminormBounds/RealSeminormFiltration/CoreDefs.lean}{PositivityTheory.KReal.laPos}
\leanref{Development/WeilBound/CoreStatements/GaussianSeminormBounds/RealSeminormFiltration/CoreDefs.lean}{PositivityTheory.KReal.laNeg}
\leanref{Development/WeilBound/CoreStatements/GaussianSeminormBounds/RealSeminormFiltration/CoreComb.lean}{PositivityTheory.KReal.comb_pos}
\leanref{Development/WeilBound/CoreStatements/GaussianSeminormBounds/RealSeminormFiltration/CoreComb.lean}{PositivityTheory.KReal.comb_neg}
\leanref{Development/WeilBound/CoreStatements/GaussianSeminormBounds/RealSeminormFiltration/CoreShape.lean}{PositivityTheory.KReal.shape_pos}
\leanref{Development/WeilBound/CoreStatements/GaussianSeminormBounds/RealSeminormFiltration/CoreShape.lean}{PositivityTheory.KReal.shape_neg}
\leanref{Development/WeilBound/CoreStatements/GaussianSeminormBounds/RealSeminormFiltration/CoreComb.lean}{PositivityTheory.KReal.extreme}
\leanref{Development/WeilBound/CoreStatements/GaussianSeminormBounds/RealSeminormFiltration/CoreComb.lean}{PositivityTheory.KReal.cas_lt}
\leanref{Development/WeilBound/CoreStatements/GaussianSeminormBounds/RealSeminormFiltration/CoreCasW.lean}{PositivityTheory.KReal.casW_root}
\leanref{Development/WeilBound/CoreStatements/GaussianSeminormBounds/RealSeminormFiltration/CoreCasW.lean}{PositivityTheory.KReal.degW_root}
\begin{proof}
For $b$ the $i$-th element of the first block put $r_b=R+1-2i$, and for $b$ the
$i$-th element of the second $r_b=S+1-2i$; then
$c(\mu)=\sum_b\mu_b^2+\sum_br_b\mu_b$ for every $\mu$, by counting the pairs
$b\prec b'$. Hence, if $\nu=\Lambda-e_a+e_b$ with $a\neq b$,
\[
c(\nu+\tau)-c(\Lambda+\tau)=2(\Lambda_b-\Lambda_a)+2+r_b-r_a,
\]
and, for every $c$, with $\Lambda'=\Lambda^{[c]}$,
\[
\deg_c(\nu)-\deg_c(\Lambda)=\bigl(|\Lambda'_b+1|-|\Lambda'_b|\bigr)+\bigl(|\Lambda'_a-1|-|\Lambda'_a|\bigr),
\]
where each bracket is at most $1$. Suppose also that $\nu$ is dominant. If $b$
is the $j$-th element of its block, $j>1$, and $\Lambda$ takes the same value
at the first $j$ elements, then $\nu_b=\Lambda_b+1$ exceeds the value of $\nu$
at the first element, which contradicts dominance; so $b$ is the first element
of its block. In the same way, if $\Lambda$ takes the same value at $a$ and at
the elements after it, then $a$ is the last element of its block. When $a$ is
the last element of one block and $b$ the first of the other,
$r_b-r_a=R+S-2$, and the first display gives
$c(\nu+\tau)-c(\Lambda+\tau)=2(\Lambda_b-\Lambda_a)+R+S$.

(1) \emph{Joint shape.} With $\alpha=(n,\ldots,n)$ of length $P$,
$\gamma=(m,\ldots,m)$ of length $Q$ and $\beta,\delta$ empty, $c_0$ and
$\Lambda^+$ form a joint shape: $c_0^{[c_{R,S}]}$ is $n$ on the first block of
$\mathrm{Ix}(P,Q)$ and $m$ on the second, and $\Lambda^{+[c_{P,Q}]}$ is $n$ at
the first $P$ elements of the first block, $m$ at the first $Q$ elements of the
second, and $0$ elsewhere. Conversely, if $c_0$ and $\nu$ form a joint shape
with sequences $\alpha',\beta',\gamma',\delta'$, then
$(\alpha',0,\ldots,0,\beta')$ is constant equal to $n>0$, so $\beta'$ is empty
and $\alpha'=(n,\ldots,n)$ has length $P$; in the same way $\delta'$ is empty
and $\gamma'=(m,\ldots,m)$. So $\nu^{[c_{P,Q}]}=\Lambda^{+[c_{P,Q}]}$ and
$\nu=\Lambda^+$.

\emph{The gap.} Let $\nu=\Lambda^+-e_a+e_b$ be as in
Definition~\ref{asm-theta-agree-real-quotient-shapes}(4), and put
$\Lambda'=\Lambda^{+[c_{P,Q}]}$. Here $\deg_{c_{R,S}}(c_0)=Pn+Qm$, and
$\deg_{c_{P,Q}}(\Lambda^+)=Pn+Qm$, so the two brackets above, for
$c=c_{P,Q}$, are both $1$. Suppose that $a$ is the $i$-th element of the
first block and $b$ the $j$-th of the second. If $i\leq P$, then
$|\Lambda'_a-1|-|\Lambda'_a|=|n-1|-n=-1$; so $i>P$. If $j>Q$, then
$\nu^{[c_{P,Q}]}_b=\Lambda'_b+1=1>0$, which the sign bound excludes; so
$j\leq Q$. Then $\Lambda^+$ is constant on the first $j$ elements of the second
block and on the elements of the first block from $a$ on, so $b$ is the first
element of the second block and $a$ the last of the first, and
$c(\nu+\tau)-c(\Lambda^++\tau)=2(m-2c_{P,Q})+R+S>R+S$. If instead $a$ is the
$j$-th element of the second block and $b$ the $i$-th of the first, then
$j>Q$ and $i\leq P$ in the same way, $b$ is the first element of the first
block and $a$ the last of the second, and the difference is
$2(n+2c_{P,Q})+R+S>R+S$.

(2) The argument is the same. The sequences are $\beta=(-m,\ldots,-m)$ of
length $P$ and $\delta=(-n,\ldots,-n)$ of length $Q$, with $\alpha,\gamma$
empty; $c_0^{[c_{R,S}]}$ is $-m$ on the first block and $-n$ on the second,
and the uniqueness follows as in (1). Here $\deg_{c_{R,S}}(c_0)=Pm+Qn=\deg_{c_{P,Q}}(\Lambda^-)$.
If $a$ is the $i$-th element of the first block and $b$ the $j$-th of the
second, then $j\leq S-P$, as otherwise $|\Lambda'_b+1|-|\Lambda'_b|=-1$, and
$i>R-Q$, as otherwise $\nu^{[c_{P,Q}]}_a=-1<0$, which the sign bound excludes;
the difference is $2(n-2c_{P,Q})+R+S$. If $a$ is in the second block and $b$
in the first, then $b$ is among the first $R-Q$ elements of the first block
and $a$ among the last $P$ of the second, and the difference is
$2(m+2c_{P,Q})+R+S$.
\end{proof}

\begin{lemma}\label{asm-theta-agree-real-quotient-intertwiner}
Let $p,q\geq0$, and let $\mathcal A$ and $\mathcal B$ be complex inner product
spaces with linear maps $Z\mapsto\rho_{\mathcal A}(Z)$ and
$Z\mapsto\rho_{\mathcal B}(Z)$ from $M_{p+q}(\mathbf C)$ to their operators
that satisfy the bracket relation and
$\langle\rho(E_{ab})u,v\rangle=\epsilon_{ab}\langle u,\rho(E_{ba})v\rangle$ for
all $a,b\in\mathrm{Ix}(p,q)$, with $\epsilon_{ab}$ as in
Lemma~\ref{asm-theta-agree-real-quotient-fock-rules}. Let
$\lambda\colon\mathrm{Ix}(p,q)\to\mathbf C$, $x\in\mathcal A$ nonzero and
$y\in\mathcal B$, with $\rho_{\mathcal A}(E_{ab})x=\delta_{ab}\lambda_ax$ and
$\rho_{\mathcal B}(E_{ab})y=\delta_{ab}\lambda_ay$ for $a,b$ in the same block,
and $\rho_{\mathcal A}(E_{ji})x=0$ and $\rho_{\mathcal B}(E_{ji})y=0$ for $i$
in the first block and $j$ in the second. Let $G$ be the smallest subspace of
$\mathcal A$ containing $x$ and stable under every $\rho_{\mathcal A}(Z)$, and
suppose that $\mathcal A=G+G^\perp$. Then there is a linear map
$\Phi\colon\mathcal A\to\mathcal B$ with
$\Phi\rho_{\mathcal A}(Z)=\rho_{\mathcal B}(Z)\Phi$ for every $Z$ and
$\Phi(x)=y$.
\end{lemma}
\leanref{Development/WeilBound/CoreStatements/GaussianSeminormBounds/RealSeminormFiltration/CoreGraph.lean}{PositivityTheory.KReal.graph_map}
\leanref{Development/WeilBound/CoreStatements/GaussianSeminormBounds/RealSeminormFiltration/CoreGraph.lean}{PositivityTheory.KReal.graph_sub}
\leanref{Development/WeilBound/CoreStatements/GaussianSeminormBounds/RealSeminormFiltration/CoreGraph.lean}{PositivityTheory.KReal.graph_orth_stable}
\begin{proof}
On $\mathcal A\times\mathcal B$ put $\rho(Z)=\rho_{\mathcal A}(Z)\times\rho_{\mathcal B}(Z)$;
it satisfies the bracket relation. Call $E_{ij}$, with $i$ in the first block
and $j$ in the second, the \emph{raising} matrix units, and the $E_{ji}$ the
\emph{lowering} ones. Let $N$ be the smallest subspace containing $(x,y)$ and
stable under the $\rho(E_{ij})$ with $E_{ij}$ raising.

(i) $N$ is stable under $\rho(E_{ab})$ for $a,b$ in the same block. The
$n\in N$ with $\rho(E_{ab})n\in N$ for all such $(a,b)$ form a subspace that
contains $(x,y)$, an eigenvector of each $\rho(E_{ab})$, and is stable under the
raising operators, since
$\rho(E_{ab})\rho(E_{ij})=\rho(E_{ij})\rho(E_{ab})+\delta_{bi}\rho(E_{aj})-\delta_{ja}\rho(E_{ib})$,
and $E_{aj}$ with $b=i$, and $E_{ib}$ with $j=a$, are raising.

(ii) $N$ is stable under the lowering operators: the $n\in N$ with
$\rho(E_{ji})n\in N$ for all lowering $E_{ji}$ contain $(x,y)$, which they
kill, and are stable under the raising operators, as
$[E_{ji},E_{i'j'}]=\delta_{ii'}E_{jj'}-\delta_{j'j}E_{i'i}$ is block-diagonal
and (i) applies. So $N$ is stable under every $\rho(Z)$.

(iii) Let $c_x=\langle x,x\rangle$ and $c_y=\langle y,y\rangle$. Every
$(u,v)\in N$ has $c_x\langle v,y\rangle=c_y\langle u,x\rangle$: this holds for
$(x,y)$, and both sides vanish after a raising operator, since
$\langle\rho_{\mathcal B}(E_{ij})v,y\rangle=-\langle v,\rho_{\mathcal B}(E_{ji})y\rangle=0$
and likewise in $\mathcal A$. Next, every $n\in N$ has
$c_x\langle m_2,n_2\rangle=c_y\langle m_1,n_1\rangle$ for all $m\in N$, where
$m=(m_1,m_2)$: for $n=(x,y)$ this is the previous claim, and if it holds for $n$,
it holds for $\rho(E_{ij})n$, because
$\langle m_2,\rho_{\mathcal B}(E_{ij})n_2\rangle=-\langle\rho_{\mathcal B}(E_{ji})m_2,n_2\rangle$,
the same holds in $\mathcal A$, and $\rho(E_{ji})m\in N$ by (ii). For
$n=m=(0,v)\in N$ this gives $c_x\langle v,v\rangle=0$, so $v=0$, as $x\neq0$.

(iv) So $N$ is the graph of a linear map from the first projection of $N$ to
$\mathcal B$. That projection is stable under the $\rho_{\mathcal A}(Z)$ and
contains $x$, so it contains $G$. Let $\Phi$ be that map on $G$ and $0$ on
$G^\perp$. By the adjoint relation, $G^\perp$ is stable under every
$\rho_{\mathcal A}(Z)$, and $G\cap G^\perp=0$, so $\mathcal A=G\oplus G^\perp$.
For $u=g+g'$ with $g\in G$ and $g'\in G^\perp$, $(\rho_{\mathcal A}(Z)g,\rho_{\mathcal B}(Z)\Phi(g))\in N$,
so $\Phi(\rho_{\mathcal A}(Z)u)=\rho_{\mathcal B}(Z)\Phi(g)=\rho_{\mathcal B}(Z)\Phi(u)$.
Finally $(x,y)\in N$ gives $\Phi(x)=y$.
\end{proof}

\begin{lemma}\label{asm-theta-agree-real-quotient-nonzero}
Let $P,Q,R,S\geq0$ and $n,m>0$ be integers, let $c_0\in\mathbf Q$, use the
notation of Definition~\ref{asm-theta-agree-real-quotient-shapes}(1), and let
$\Theta$ be the quotient of $\mathcal P(P,Q;R,S)$ by the span of the
polynomials $d_1(X)f-c_0\operatorname{tr}(X)f$, for all $X$ and $f$
(Definition~\ref{asm-theta-agree-real-fock}(1)). Then $\Theta\neq0$ in each
of the following cases:
\begin{enumerate}
\item $P\leq R$, $Q\leq S$ and $c_0=n+c_{R,S}=m-c_{R,S}$;
\item $Q\leq R$, $P\leq S$ and $c_0=-m+c_{R,S}=-n-c_{R,S}$.
\end{enumerate}
\end{lemma}
\leanref{Development/WeilBound/CoreStatements/GaussianSeminormBounds/RealSeminormFiltration/CoreNonv.lean}{PositivityTheory.KReal.nontrivial_of_pair}
\leanref{Development/WeilBound/CoreStatements/GaussianSeminormBounds/RealSeminormFiltration/CoreDelta.lean}{PositivityTheory.KReal.nontrivial_of_form}
\leanref{Development/WeilBound/CoreStatements/GaussianSeminormBounds/RealSeminormFiltration/CoreDefs.lean}{PositivityTheory.KReal.PairData}
\leanref{Development/WeilBound/CoreStatements/GaussianSeminormBounds/RealSeminormFiltration/CorePair.lean}{PositivityTheory.KReal.pair_pos}
\leanref{Development/WeilBound/CoreStatements/GaussianSeminormBounds/RealSeminormFiltration/CorePair.lean}{PositivityTheory.KReal.pair_neg}
\leanref{Development/WeilBound/CoreStatements/GaussianSeminormBounds/RealSeminormFiltration/CorePair.lean}{PositivityTheory.KReal.exists_blk}
\leanref{Development/WeilBound/CoreStatements/GaussianSeminormBounds/RealSeminormFiltration/CoreProj.lean}{PositivityTheory.KReal.proj_gen}
\leanref{Development/WeilBound/CoreStatements/GaussianSeminormBounds/RealSeminormFiltration/CoreNonv.lean}{PositivityTheory.KReal.dualAct}
\begin{proof}
Write $\mathcal A=\mathcal P(P,Q;R,0)$ and $\mathcal B=\mathcal P(P,Q;0,S)$,
each with its operators $d_1$ and its Fock inner product, and let $i$ and $j$
run over the first and second blocks of $\mathrm{Ix}(P,Q)$.

\emph{Step 1: the vectors.} There are $\lambda\colon\mathrm{Ix}(P,Q)\to\mathbf Q$,
a nonzero homogeneous $x\in\mathcal A$ and a nonzero $y\in\mathcal B$ with
$d_1(E_{aa'})x=\delta_{aa'}\lambda_ax$ and
$d_1(E_{aa'})y=\delta_{aa'}(c_0-\lambda_a)y$ for $a,a'$ in the same block,
$d_1(E_{ji})x=0$ and $d_1(E_{ij})y=0$. In case (1), take $\lambda$ equal to
$n+c_{R,0}$ on the first block and $-c_{R,0}$ on the second. With
$\alpha=(n,\ldots,n)$ of length $P\leq R$ and $\beta,\gamma,\delta$ empty,
$\lambda$ forms a joint shape with a map on $\mathrm{Ix}(R,0)$, so by
Theorem~\ref{asm-theta-agree-real-quotient-fock-facts}(3), for $\mathcal A$,
there is a nonzero joint harmonic that is a highest-weight vector of weight
$\lambda$ for $d_1$. Let $x$ be a nonzero homogeneous component of it; it is
again one (Lemma~\ref{asm-theta-agree-real-quotient-fock-rules}(3), and
Theorem~\ref{asm-theta-agree-real-quotient-fock-facts}(1)). As $\lambda$ is
constant on each block, $d_1(E_{aa'})x=\delta_{aa'}\lambda_ax$
(Lemma~\ref{asm-theta-agree-real-quotient-fock-rules}(5)). In $\mathcal A$,
$d_1(E_{ji})=-\sum_b\partial^2/\partial x_{(j,b)}\partial x_{(i,b)}$ lowers the
degree by $2$, so it is its own lowering part, and $d_1(E_{ji})x=0$ as $x$ is
harmonic. In the same way, with $\gamma=(m,\ldots,m)$ of length $Q\leq S$ and
$\alpha,\beta,\delta$ empty, $\mathcal B$ has a nonzero joint harmonic $y$
that is a highest-weight vector for $d_1$ of the weight that is $c_{0,S}$ on
the first block and $m-c_{0,S}$ on the second; this weight is $c_0-\lambda$,
as $c_0=n+c_{R,S}=m-c_{R,S}$. In $\mathcal B$,
$d_1(E_{ij})=-\sum_b\partial^2/\partial x_{(i,b)}\partial x_{(j,b)}$, so
$d_1(E_{ij})y=0$. In case (2), take $\lambda$ equal to $c_{R,0}$ on the first
block and $-n-c_{R,0}$ on the second, with $\delta=(-n,\ldots,-n)$ of length
$Q\leq R$ for $\mathcal A$, and $\beta=(-m,\ldots,-m)$ of length $P\leq S$ for
$\mathcal B$, which gives the weight $-m-c_{0,S}$ on the first block and
$c_{0,S}$ on the second, that is $c_0-\lambda$; the rest is the same.

\emph{Step 2: a criterion.} Identify $\mathcal P(P,Q;R,S)$ with
$\mathcal A\otimes\mathcal B$, the variables $x_{(a,b)}$ with $b$ in the
first block of $\mathrm{Ix}(R,S)$ going to $\mathcal A$ and the others to
$\mathcal B$. Then $d_1(Z)$ becomes $d_1(Z)\otimes1+1\otimes d_1(Z)$, as each
$\omega_{(a,b),(a',b)}$ involves only the variables with second index $b$. Let
$\beta$ be a bilinear form on $\mathcal A\times\mathcal B$ with
\[
\beta(d_1(Z)u,v)+\beta(u,d_1(Z)v)=c_0\operatorname{tr}(Z)\beta(u,v)
\]
for all $Z$, $u$ and $v$. The linear form $u\otimes v\mapsto\beta(u,v)$ then
vanishes on every $d_1(Z)f-c_0\operatorname{tr}(Z)f$. So if $\beta(x,y)\neq0$,
the image of the product $xy$ in $\Theta$ is nonzero.

\emph{Step 3: the form.} Let $\theta(Z)$ be the matrix with entries
$\theta(Z)_{aa'}=\epsilon_{aa'}Z_{a'a}$, with $\epsilon$ as in
Lemma~\ref{asm-theta-agree-real-quotient-fock-rules}; then
$\theta(ZW)=\theta(W)\theta(Z)$. On $\mathcal B$ put
$\rho_{\mathcal B}(Z)=c_0\operatorname{tr}(Z)-d_1(\theta(Z))$. It satisfies the
bracket relation, and, as $c_0$ is real, the adjoint relation of
Lemma~\ref{asm-theta-agree-real-quotient-intertwiner}, by
Lemma~\ref{asm-theta-agree-real-quotient-fock-rules}(1); so does
$\rho_{\mathcal A}=d_1$ on $\mathcal A$. For $w\in\mathcal B$ let $\bar w$ have
the conjugate coefficients; the operators $d_1(E_{aa'})$ have real
coefficients, so they commute with $w\mapsto\bar w$. The $\lambda_a$ are
real. So $\rho_{\mathcal B}(E_{aa'})\bar y=\delta_{aa'}\lambda_a\bar y$ for
$a,a'$ in the same block, and
$\rho_{\mathcal B}(E_{ji})\bar y=\overline{d_1(E_{ij})y}=0$. Let $G$ be the
smallest subspace of $\mathcal A$ containing $x$ and stable under every
$d_1(Z)$. As $x$ is homogeneous and each $d_1(E_{aa'})$ maps homogeneous
polynomials to homogeneous ones, $G$ contains the homogeneous components of its
elements. For $u\in\mathcal A$, let $y_d$ be the orthogonal projection of
$u_d$ on the finite-dimensional space $G\cap\mathcal A_d$; then
$\sum_dy_d\in G$ and $u-\sum_dy_d\in G^\perp$, as components of different
degrees are orthogonal. So $\mathcal A=G+G^\perp$, and
Lemma~\ref{asm-theta-agree-real-quotient-intertwiner}, with $\lambda$, $x$ and
$\bar y$, gives $\Phi\colon\mathcal A\to\mathcal B$ with
$\Phi(d_1(Z)u)=c_0\operatorname{tr}(Z)\Phi(u)-d_1(\theta(Z))\Phi(u)$ and
$\Phi(x)=\bar y$. Put $\beta(u,v)=\langle\overline{\Phi(u)},v\rangle$, a
bilinear form. By Lemma~\ref{asm-theta-agree-real-quotient-fock-rules}(1) and
$\epsilon_{aa'}^2=1$,
$\langle\overline{d_1(\theta(Z))w},v\rangle=\langle\bar w,d_1(Z)v\rangle$ for
all $w$ and $v$, so $\beta$ satisfies the relation of Step 2. And
$\beta(x,y)=\langle y,y\rangle\neq0$.
\end{proof}

\begin{lemma}\label{asm-theta-agree-local-mult-archimedean-sides}
Let $a,a',p,q,k\geq0$ be integers with $1\leq a'<a$, $a-a'$ even,
$2p+1\leq a$, $2q+1\leq a'$, $a+a'<k$, $k$ even, $2(p+q)<a'$ and
$2p+a'<a$. Put $(P,Q)=(a-p,p)$, $(P',Q')=(q,a'-q)$,
$(R,S)=(a+a'-p-q,p+q)$, $n=(k+2(p+q)-a-a')/2$ and
$m=(k+a+a'-2(p+q))/2$. Let $\Theta$ be the quotient of
$\mathcal P(P,Q;R,S)$ by the span of the polynomials
$d_1(X)f-\tfrac k2\operatorname{tr}(X)f$, with the operators $\vartheta(Y)$
induced by $d_2(Y)$, and $\Theta'$ the quotient of $\mathcal P(P',Q';R,S)$ by
the span of the polynomials $d_1(X)f+\tfrac k2\operatorname{tr}(X)f$, with the
operators $\vartheta'(Y)$ induced by $d_2(Y)+\tfrac k2\operatorname{tr}(Y)$,
for all $X$, $Y$ and $f$ (Definition~\ref{asm-theta-agree-real-fock}(1) and
Lemma~\ref{asm-theta-agree-real-fock-relations}(1)). Let
$\Lambda\colon\mathrm{Ix}(R,S)\to\mathbf Q$ take the value
$n+\tfrac12(P-Q)$ at the first $P$ elements of the first block,
$\tfrac12(P-Q)$ at its other elements, $m-\tfrac12(P-Q)$ at the first $Q$
elements of the second block, and $-\tfrac12(P-Q)$ at its other elements. Let
$\Lambda'$ take the value $\tfrac k2+\tfrac12(P'-Q')-n$ at the last $Q'$
elements of the first block, $\tfrac k2+\tfrac12(P'-Q')$ at its other
elements, $\tfrac k2-\tfrac12(P'-Q')-m$ at the last $P'$ elements of the second
block, and $\tfrac k2-\tfrac12(P'-Q')$ at its other elements. Then $n$ and $m$
are positive integers; $\Theta$ with the operators $\vartheta(Y)$ is a Lie
module of type $(R,S)$ generated with Casimir gap in weight $\Lambda$, and
$\Theta'$ with the operators $\vartheta'(Y)$ is a Lie module of type $(R,S)$
generated with Casimir gap in weight $\Lambda'$
(Definition~\ref{asm-theta-agree-real-quotient-lie-module}); and
$\Lambda'=\Lambda$.
\end{lemma}
\leanref{Development/WeilBound/CoreStatements/GaussianSeminormBounds/RealSeminormFiltration/CoreComb.lean}{PositivityTheory.KReal.nm_facts}
\leanref{Development/WeilBound/CoreStatements/GaussianSeminormBounds/RealSeminormFiltration/CoreComb.lean}{PositivityTheory.KReal.combV}
\leanref{Development/WeilBound/CoreStatements/GaussianSeminormBounds/RealSeminormFiltration/CoreComb.lean}{PositivityTheory.KReal.combV'}
\leanref{Development/WeilBound/CoreStatements/GaussianSeminormBounds/RealSeminormFiltration/CoreComb.lean}{PositivityTheory.KReal.la_eq}
\begin{proof}
The numbers $n$ and $m$ are integers, as $k$ and $a+a'$ are even, and
positive, as $a+a'<k$ and $2(p+q)<a'$. Next, $\Lambda'=\Lambda$: the first
block has $R=P+Q'$ elements and the second $S=Q+P'$, and
$\tfrac k2+\tfrac12(P'-Q')=n+\tfrac12(P-Q)$ and
$\tfrac k2-\tfrac12(P'-Q')=m-\tfrac12(P-Q)$, as both sides of the first equation
are $\tfrac12(k+2q-a')$ and both sides of the second are $\tfrac12(k+a'-2q)$.

With the notation of Definition~\ref{asm-theta-agree-real-quotient-shapes}(1),
$c_{R,S}=\tfrac12(a+a')-(p+q)$, so $\tfrac k2=n+c_{R,S}=m-c_{R,S}$, and
$c_{P,Q}=\tfrac12(P-Q)$ and $c_{P',Q'}=\tfrac12(P'-Q')$. As $q\leq a'$ and
$p\leq a$, $P\leq R$, $Q\leq S$, $Q'\leq R$ and $P'\leq S$. As $k>a+a'$,
$a>2p+a'$ and $a'>2q$, the numbers $n+2c_{P,Q}=n+a-2p$,
$m-2c_{P,Q}=n-2c_{P',Q'}=\tfrac12(k-a+a'+2p-2q)$ and
$m+2c_{P',Q'}=\tfrac12(k+a-a'-2p+2q)$ are positive. The map $\Lambda$ is the
map $\Lambda^+$ of Lemma~\ref{asm-theta-agree-real-quotient-weights}(1), for
$(P,Q,R,S)$, $n$ and $m$; by that lemma it is a gap weight for
$(P,Q,\tfrac k2,0)$. By
Lemma~\ref{asm-theta-agree-real-quotient-nonzero}(1), $\Theta\neq0$. So by
Proposition~\ref{asm-theta-agree-real-quotient-side}, with $c_0=\tfrac k2$ and
$\tau_0=0$, $\Theta$ with the operators $\vartheta(Y)$ is a Lie module of type
$(R,S)$ generated with Casimir gap in weight $\Lambda$. The map $\Lambda^-$ of
Lemma~\ref{asm-theta-agree-real-quotient-weights}(2), for $(P',Q',R,S)$, $n$
and $m$, is $\Lambda'-\tfrac k2$, and $-\tfrac k2=-m+c_{R,S}=-n-c_{R,S}$. So
$\Lambda'-\tfrac k2$ is a gap weight for $(P',Q',-\tfrac k2,\tfrac k2)$,
$\Theta'\neq0$ by Lemma~\ref{asm-theta-agree-real-quotient-nonzero}(2), and by
Proposition~\ref{asm-theta-agree-real-quotient-side}, with $c_0=-\tfrac k2$
and $\tau_0=\tfrac k2$, $\Theta'$ with the operators $\vartheta'(Y)$ is a Lie
module of type $(R,S)$ generated with Casimir gap in weight $\Lambda'$.
\end{proof}
% Lean: no single declaration states this lemma. coreLM (CoreLM.lean) and one_data
% (ExistsGreatestProperInvariantSubmoduleOfSide.lean), and two_data in LocalPoly.lean, each build
% the two sides inline: s1 = side_of (kappa/2) 0 (combV h) n1, a Side for laPos + 0 on
% ThV D = ThC P Q R S (kappa/2) with thAct tau = 0, and s2 = side_of (-kappa/2) (kappa/2)
% (combV' h) n2, a Side for laNeg + kappa/2 on ThV' D; la_eq rewrites laNeg + kappa/2 to
% laPos + 0. nm_facts gives n = nD D, m = mD D, their positivity and P <= R, Q <= S, Q' <= R,
% P' <= S; rat_facts gives the positivity of the four numbers (inside combV and combV',
% CoreComb.lean, about 10 lines each). n1 and n2 are nontrivial_of_pair with pair_pos and
% pair_neg (Lemma asm-theta-agree-real-quotient-nonzero); side_of is Proposition
% asm-theta-agree-real-quotient-side, and comb_pos and comb_neg for laPos and laNeg are Lemma
% asm-theta-agree-real-quotient-weights. The shared material occupies CoreSide, CoreComb,
% CoreShape, CoreCasW, CoreNonv, CorePair, CoreDelta, CoreGraph, CoreProj, CoreGen, CoreSplit,
% CoreHarm, CoreStar, CoreRoot, CoreSupp, CoreLie and CoreCyc (about 4,700 lines), with the
% Fock statements of Theorem asm-theta-agree-real-quotient-fock-facts.

\begin{lemma}\label{asm-theta-agree-real-quotient}
Let $a,a',p,q,k\geq0$ be integers with $1\leq a'<a$, $a-a'$ even,
$2p+1\leq a$, $2q+1\leq a'$, $a+a'<k$, $k$ even, $2(p+q)<a'$ and
$2p+a'<a$. Put $(P,Q)=(a-p,p)$, $(P',Q')=(q,a'-q)$ and
$(R,S)=(a+a'-p-q,p+q)$. Let $\Theta$ be the quotient of
$\mathcal P(P,Q;R,S)$ by the span of the polynomials
$d_1(X)f-\tfrac k2\operatorname{tr}(X)f$, and $\Theta'$ the quotient of
$\mathcal P(P',Q';R,S)$ by the span of the polynomials
$d_1(X)f+\tfrac k2\operatorname{tr}(X)f$, for all $X$ and $f$
(Definition~\ref{asm-theta-agree-real-fock}(1)). By
Lemma~\ref{asm-theta-agree-real-fock-relations}(1), for
$Y\in M_{R+S}(\mathbf C)$ the operator $d_2(Y)$ induces an operator
$\vartheta(Y)$ on $\Theta$, and $d_2(Y)+\tfrac k2\operatorname{tr}(Y)$
induces an operator $\vartheta'(Y)$ on $\Theta'$. Then there are a complex
vector space $\Pi\neq0$, operators $\varrho(Y)$ on $\Pi$ for
$Y\in M_{R+S}(\mathbf C)$ such that $0$ and $\Pi$ are the only subspaces
stable under every $\varrho(Y)$, and surjective linear maps
$\varpi\colon\Theta\to\Pi$ and $\varpi'\colon\Theta'\to\Pi$ with
$\varpi\vartheta(Y)=\varrho(Y)\varpi$ and $\varpi'\vartheta'(Y)=\varrho(Y)\varpi'$
for every $Y$.
\end{lemma}
\leanref{Development/WeilBound/CoreStatements/GaussianSeminormBounds/RealSeminormFiltration/TensorIrr.lean}{PositivityTheory.KReal.CoreLM}
\leanref{Development/WeilBound/CoreStatements/GaussianSeminormBounds/RealSeminormFiltration/CoreLM.lean}{PositivityTheory.KReal.coreLM}
\begin{proof}
By Lemma~\ref{asm-theta-agree-local-mult-archimedean-sides}, $\Theta$ with the
operators $\vartheta(Y)$ and $\Theta'$ with the operators $\vartheta'(Y)$ are
Lie modules of type $(R,S)$, both generated with Casimir gap in the same
weight $\Lambda=\Lambda'$
(Definition~\ref{asm-theta-agree-real-quotient-lie-module}(3)).
Proposition~\ref{asm-theta-agree-real-quotient-common}, applied to $\Theta$ and
$\Theta'$, gives the space $\Pi$, the operators $\varrho(Y)$ and the maps
$\varpi$ and $\varpi'$.
\end{proof}
% Lean: coreLM (CoreLM.lean): s1 and s2 are the two sides of Lemma
% asm-theta-agree-local-mult-archimedean-sides (side_of with combV h and combV' h, and
% nontrivial_of_pair), la_eq rewrites the weight of s2, and common_quot concludes.

\subsection{Real maps matched to a triple}\label{subsection-asm-real-maps}

\begin{lemma}\label{asm-theta-agree-real-one-place}
Let $a,a',p,q\geq0$ and $\kappa\in\mathbf Z$, with $a+a'$ and $\kappa$ even,
$a+a'<\kappa$, $2(p+q)+1<\lfloor(a-a')/2\rfloor$ and $2(p+q)+1<a'$. Put
$(R,S)=(a+a'-p-q,p+q)$. Then there are an irreducible
$(\mathfrak g,K)$-module $N$ of type $(R,S)$ with scalar endomorphisms, a lift
$\pi\colon\mathcal P(a-p,p;R,S)\to N$ of type $(\kappa/2,0)$ and a lift
$\pi'\colon\mathcal P(q,a'-q;R,S)\to N$ of type $(-\kappa/2,\kappa/2)$
(Definition~\ref{asm-theta-agree-real-fock}).
\end{lemma}
\leanref{Development/WeilBound/CoreStatements/GaussianSeminormBounds/RealSeminormFiltration/KRealAsmOne.lean}{PositivityTheory.KRealA.onePlace}
\leanref{Development/WeilBound/CoreStatements/GaussianSeminormBounds/RealSeminormFiltration/KRealAsmOne.lean}{PositivityTheory.KRealA.onePlaceD}
\leanref{Development/WeilBound/CoreStatements/GaussianSeminormBounds/RealSeminormFiltration/KRealAsmOne.lean}{PositivityTheory.KRealA.gkImage}
\leanref{Development/WeilBound/CoreStatements/GaussianSeminormBounds/RealSeminormFiltration/KRealAsmOne.lean}{PositivityTheory.KRealA.scalarEnd_of}
\begin{proof}
As $\kappa>a+a'\geq0$, $k=\kappa$ is a nonnegative integer. The hypotheses of
Lemma~\ref{asm-theta-agree-real-quotient} hold for $(a,a',p,q,k)$: from
$1\leq2(p+q)+1<\lfloor(a-a')/2\rfloor$ we get $a-a'\geq4$ and
$4p+2<a-a'$, so $a'<a$, $2p+1\leq a$ and $2p+a'<a$; from $2(p+q)+1<a'$ we get
$a'\geq1$, $2q+1\leq a'$ and $2(p+q)<a'$; and $a-a'$ is even since $a+a'$ is.
Let $\Pi$, $\varrho$, $\varpi$ and $\varpi'$ be as there, and let $\pi$ and
$\pi'$ be the compositions of $\varpi$ and $\varpi'$ with the quotient maps
from $\mathcal P(P,Q;R,S)$ and $\mathcal P(P',Q';R,S)$, where $(P,Q)=(a-p,p)$
and $(P',Q')=(q,a'-q)$. They are surjective, and for all $Y$ and $f$,
\[
\varrho(Y)\pi(f)=\pi(d_2(Y)f),\qquad
\varrho(Y)\pi'(f)=\pi'(d_2(Y)f)+\tfrac\kappa2\operatorname{tr}(Y)\pi'(f).
\]

Let $N=\Pi$ with $N(Y)=\varrho(Y)$. It is a $(\mathfrak g,K)$-module of type
$(R,S)$; each condition is checked on vectors $\pi(f)$, as $\pi$ is
surjective. The map $Y\mapsto\varrho(Y)$ is linear and satisfies the bracket
relation because $Y\mapsto d_2(Y)$ does
(Lemma~\ref{asm-theta-agree-real-fock-relations}(1)) and
$\varrho(Y)\pi=\pi d_2(Y)$. If $f$ has degree $n$, the polynomials of degree at
most $n$ form a finite-dimensional space stable under the $d_2(Y)$ with $Y$
block-diagonal (Lemma~\ref{asm-theta-agree-real-fock-relations}(2)), and its
image under $\pi$ contains $\pi(f)$ and is stable under the corresponding
$\varrho(Y)$. Every $f$ is a sum of monomials, which are eigenvectors of every
$d_2(E_{bb})$ with rational eigenvalues
(Lemma~\ref{asm-theta-agree-real-fock-relations}(3)), and $\pi$ maps them to
eigenvectors of the $\varrho(E_{bb})$ with the same eigenvalues.

The module $N$ is irreducible by Lemma~\ref{asm-theta-agree-real-quotient}. It
has countable dimension, as the image of a polynomial ring in finitely many
variables, so it has scalar endomorphisms by
Lemma~\ref{th-first-lift-dixmier}. In $\Theta$,
$d_1(X)f$ and $\tfrac\kappa2\operatorname{tr}(X)f$ have the same image, so
$\pi(d_1(X)f)=\tfrac\kappa2\operatorname{tr}(X)\pi(f)$; with the first
displayed relation, $\pi$ is a lift of type $(\kappa/2,0)$. In the same way
$\pi'(d_1(X)f)=-\tfrac\kappa2\operatorname{tr}(X)\pi'(f)$, and with the second
displayed relation $\pi'$ is a lift of type $(-\kappa/2,\kappa/2)$.
\end{proof}

\begin{lemma}\label{asm-theta-agree-real-parity}
Let $(B,G,\mathcal X,\omega,\omega')$ be a triple of ranks $(a,a')$ for a theta
base $B$ of a CM field $E$ (Definition~\ref{asm-theta-agree-triple}). Then
$m_\alpha(w)-m_\beta(w)$ is even for every infinite place $w$ of $F$.
\end{lemma}
\leanref{Development/WeilBound/CoreStatements/GaussianSeminormBounds/RealSeminormFiltration/RLParity.lean}{PositivityTheory.RL.even_im_hermitian_sub}
\begin{proof}
Fix $w$, and let $s\in\mathbf A^\times$ have component $-1$ at $w$ and $1$ at
every other place. Its finite components are $1$, and $\pi_{w'}(s)$ is $-1$ for
$w'=w$ and $1$ otherwise. Since $\mathcal X$ is proper,
$\alpha(s)=\prod_{w'}(\pi_{w'}(s)/|\pi_{w'}(s)|)^{m_\alpha(w')}=(-1)^{m_\alpha(w)}$,
and likewise $\beta(s)=(-1)^{m_\beta(w)}$. Also $\alpha(s)=\eta(s)^{a'}$ and
$\beta(s)=\eta(s)^a$, with $\eta(s)=\pm1$; since $G$ is proper, $a+a'$ is even,
so $\eta(s)^{a'}=\eta(s)^a$. Hence
$(-1)^{m_\alpha(w)-m_\beta(w)}=\alpha(s)\beta(s)^{-1}=1$.
\end{proof}

\begin{lemma}\label{asm-theta-agree-real-frames}
Let $n,r,s\geq0$ and $H,\theta\in M_n(\mathbf C)$. Suppose that $H$ has
signature $(r,s)$, in the sense that $r+s=n$ and $B^{\mathsf T}H\bar B=J_{r,s}$
for some $B\in M_{n\times(r+s)}(\mathbf C)$, and that $\theta$ is a Cartan
involution of $H$ (Definition~\ref{asm-theta-agree-doubling}(1)). Then $H$ has
a frame $A$ of type $(r,s)$ (Definition~\ref{asm-theta-agree-real-fock}(5))
with $\theta A=AJ_{r,s}$.
\end{lemma}
\leanref{Development/AbelianVarietyTheory/Proofs/DivisorProofs/IsFockMap.lean}{GlobalTheory.RealLift.FrameEx}
\leanref{Development/WeilBound/CoreStatements/GaussianSeminormBounds/RealSeminormFiltration/RLFrame.lean}{PositivityTheory.RL.frameEx}
\leanref{Development/WeilBound/CoreStatements/GaussianSeminormBounds/RealSeminormFiltration/RLFrame.lean}{PositivityTheory.RL.Frame.frame_conj}
\leanref{Development/WeilBound/CoreStatements/GaussianSeminormBounds/RealSeminormFiltration/RLFrame.lean}{PositivityTheory.RL.Frame.diag_ex}
\leanref{Development/WeilBound/CoreStatements/GaussianSeminormBounds/RealSeminormFiltration/RLFrame.lean}{PositivityTheory.RL.Frame.inertia_le}
\begin{proof}
Write $J=J_{r,s}$, $\phi=\bar\theta$ and $B_0=\bar B$. Then $B_0^*HB_0=J$,
$\phi^2=1$ and $\phi^*H\phi=H$, since $\phi^*=\theta^{\mathsf T}$. For
$x\neq0$, the positivity condition on $\theta$ applied to $\bar x$ gives
$\operatorname{Re}(x^*\phi^*Hx)>0$.

As $J^2=1$ and $B_0$ is square, $B_0$ is invertible with inverse
$JB_0^*H$, and $H=(B_0^{-1})^*JB_0^{-1}$ is Hermitian. From $\phi^2=1$ and
$\phi^*H\phi=H$ we get $\phi^*H=H\phi$. So $N=H\phi$ is Hermitian, and it is
positive definite: $x^*Nx=x^*\phi^*Hx$ is real, with positive real part.
Write $N=P^*P$ with $P$ invertible, and put $Q=P^{-1}$ and $K=Q^*HQ$. Then $K$
is Hermitian, and $K=P\phi Q$ since $H=N\phi=P^*P\phi$; hence $K^2=P\phi^2Q=1$.
Let $K=UDU^*$ with $U$ unitary and $D$ real diagonal. Then $D^2=1$, so the
diagonal entries $d_i$ of $D$ are $\pm1$. Put $C=QU$. Then $C^*HC=U^*KU=D$, and
$\phi C=QP\phi QU=QKU=CD$.

The number of $i$ with $d_i=1$ is at most $r$. Indeed, $G=B_0^{-1}C$ satisfies
$G^*JG=D$. Let $v$ be a vector supported on $\{i:d_i=1\}$ such that $Gv$
vanishes on the first block. Then
$\sum_i|v_i|^2=v^*Dv=(Gv)^*J(Gv)\leq0$, so $v=0$. Hence
$v\mapsto(\text{first block of }Gv)$ is injective on these vectors, and the
bound follows. Applied to $-D$ and $-J$, the same argument shows that the
number of $i$ with $d_i=-1$ is at most $s$. As $n=r+s$, these numbers are $r$
and $s$, and after a permutation of the columns of $C$ we may assume $D=J$.

Put $A=\bar C$. Then $A^{\mathsf T}H\bar A=C^*HC=J$, $A$ is invertible, and
$\theta A=\overline{\phi C}=\overline{CJ}=AJ$.
\end{proof}

\begin{lemma}\label{asm-theta-agree-real-lines}
Let $B$ be a theta base for a CM field $E$, let $Y$ be a lattice space of rank
$m$ with matrix $J$, let $b\in\mathbf C$, and for every infinite place $w$ of
$F$ let $C_w$ be a frame of $b\,\tau_w(J)$ of type $(R,S)$
(Definition~\ref{asm-theta-agree-real-fock}(5)).
\begin{enumerate}
\item Every line $\gamma$ in $U_Y$ is generated, relative to $(C_w)$, by a
  family $(Y_w)$ with $Y_w\in\mathfrak u(R,S)$ for every $w$.
\item Every family $(Y_w)$ with $Y_w\in\mathfrak u(R,S)$ for every $w$
  generates a line in $U_Y$, relative to $(C_w)$.
\item Let $\theta_w\in M_m(\mathbf C)$ with $\theta_wC_w=C_wJ_{R,S}$ for every
  $w$, and let $k\in U_Y$ be at $\infty$ with $\pi_w(k)\theta_w=\theta_w\pi_w(k)$
  for every $w$. Then $k=\gamma(1)$ for a line $\gamma$ generated, relative to
  $(C_w)$, by a family of block-diagonal matrices $Y_w$.
\end{enumerate}
\end{lemma}
\leanref{Development/AbelianVarietyTheory/Proofs/DivisorProofs/IsFockMap.lean}{GlobalTheory.RealLift.GenEx}
\leanref{Development/AbelianVarietyTheory/Proofs/DivisorProofs/IsFockMap.lean}{GlobalTheory.RealLift.LineEx}
\leanref{Development/AbelianVarietyTheory/Proofs/DivisorProofs/IsFockMap.lean}{GlobalTheory.RealLift.KLine}
\leanref{Development/WeilBound/CoreStatements/GaussianSeminormBounds/RealSeminormFiltration/RLLines.lean}{PositivityTheory.RL.genEx}
\leanref{Development/WeilBound/CoreStatements/GaussianSeminormBounds/RealSeminormFiltration/RLLines.lean}{PositivityTheory.RL.lineEx}
\leanref{Development/WeilBound/CoreStatements/GaussianSeminormBounds/RealSeminormFiltration/RLLines.lean}{PositivityTheory.RL.kLine}
\leanref{Development/WeilBound/CoreStatements/GaussianSeminormBounds/RealSeminormFiltration/RLLines.lean}{PositivityTheory.RL.exists_gen_matrix}
\leanref{Development/WeilBound/CoreStatements/GaussianSeminormBounds/RealSeminormFiltration/RLLines.lean}{PositivityTheory.RL.exists_log_unitary}
\leanref{Development/WeilBound/CoreStatements/GaussianSeminormBounds/RealSeminormFiltration/RLLines.lean}{PositivityTheory.RL.exists_section}
\begin{proof}
Three facts are used. (i) For $x\in\mathbf A_E$, $\pi_w(c(x))=\overline{\pi_w(x)}$,
since $E$ is a CM field; so for $g\in U_Y$,
$\pi_w(g)^{\mathsf T}\tau_w(J)\overline{\pi_w(g)}=\tau_w(J)$. (ii) A frame is
invertible: if $C_wC'_w=1$, then $C'_wC_w=1$ since $C_w$ is square. So for
$g\in M_m(\mathbf C)$ and $h\in M_{R+S}(\mathbf C)$ with $gC_w=C_wh$, the
equality $C_w^{\mathsf T}(g^{\mathsf T}b\tau_w(J)\bar g)\bar C_w=h^{\mathsf T}J_{R,S}\bar h$
shows that $g^{\mathsf T}\tau_w(J)\bar g=\tau_w(J)$ if and only if
$h^*J_{R,S}h=J_{R,S}$. (iii) Since $E$ is a CM field, there is a continuous
map $s$ from $\mathbf C^{\{w\}}$ to the elements of $\mathbf A_E$ at $\infty$
with $\pi_w(s(z))=z_w$ for every $w$: write elements at $\infty$ in the basis
$1,\beta$ of $E$ over $F$, for some $\beta\in E$ with $c(\beta)\neq\beta$, so
that no $\tau_w(\beta)$ is real. An element of $\mathbf A_E$, or a matrix over
it, at $\infty$ with all $\pi_w$-components $0$ is $0$.

(1) Fix $w$, and let $C'_w$ be the inverse of $C_w$. Then
$t\mapsto h(t)=C'_w\pi_w(\gamma(t))C_w$ is a continuous homomorphism from
$\mathbf R$ to $GL_{R+S}(\mathbf C)$, as $\pi_w$ is a continuous ring
homomorphism. By (i) and (ii), $h(t)^*J_{R,S}h(t)=J_{R,S}$. A continuous
homomorphism from $\mathbf R$ to $GL_{R+S}(\mathbf C)$ is
$t\mapsto\exp(tY_w)$ for some $Y_w$ \cite[Theorem~2.14]{Hal15}. Differentiating
$h(t)^*J_{R,S}h(t)=J_{R,S}$ at $t=0$ gives $Y_w\in\mathfrak u(R,S)$, and
$\pi_w(\gamma(t))C_w=C_wh(t)=C_w\exp(tY_w)$.

(2) Put $g_w(t)=C_w\exp(tY_w)C'_w$, and let $G(t)$ be the matrix over
$\mathbf A_E$ equal to $1$ plus the matrix with entries
$s\bigl((g_w(t)-1)_{ij}\bigr)_w$, with $s$ as in (iii). Then $G(t)-1$ is at
$\infty$, $\pi_w(G(t))=g_w(t)$, and $t\mapsto G(t)$ is continuous. The
matrices $G(s+t)-G(s)G(t)$ and $G(t)^{\mathsf T}Jc(G(t))-J$ are at $\infty$, and
their $\pi_w$-components vanish: the first because
$g_w(s+t)=g_w(s)g_w(t)$, the second by (i) and (ii), since
$\exp(tY_w)^*J_{R,S}\exp(tY_w)=J_{R,S}$ for $Y_w\in\mathfrak u(R,S)$. By (iii),
$G(s+t)=G(s)G(t)$ and $G(t)\in U_Y$, with inverse $G(-t)$. So
$\gamma(t)=G(t)$ is a line, and $\pi_w(\gamma(t))C_w=g_w(t)C_w=C_w\exp(tY_w)$.

(3) Fix $w$ and put $h=C'_w\pi_w(k)C_w$. By (i) and (ii),
$h^*J_{R,S}h=J_{R,S}$. From $\pi_w(k)\theta_w=\theta_w\pi_w(k)$ and
$\theta_wC_w=C_wJ_{R,S}$ we get $hJ_{R,S}=J_{R,S}h$, so $h^*h=1$: $h$ is
unitary and commutes with the Hermitian matrix $J_{R,S}$. Choose $\theta\in\mathbf R$
with $-e^{i\theta}$ not an eigenvalue of $h$. Then $e^{-i\theta}h$ is unitary
without the eigenvalue $-1$, so it is $\exp(iX)$ for the Hermitian matrix $X$
obtained by applying the principal argument to its eigenvalues, and $X$
commutes with $J_{R,S}$ as $h$ does. So $h=\exp(Y_w)$ with
$Y_w=iX+i\theta$, which satisfies $Y_w^*=-Y_w$ and commutes with $J_{R,S}$;
hence $Y_w\in\mathfrak u(R,S)$ and $Y_w$ is block-diagonal. By (2) there is a
line $\gamma$ generated by $(Y_w)$. Both $\gamma(1)$ and $k$ are at $\infty$,
and $\pi_w(\gamma(1))C_w=C_w\exp(Y_w)=\pi_w(k)C_w$, so
$\pi_w(\gamma(1))=\pi_w(k)$ for every $w$, and $\gamma(1)=k$ by (iii).
\end{proof}

\begin{lemma}\label{asm-theta-agree-real-weak}
Let $I$ be a finite set, let $f,f_1,f_2\in\mathcal S(F_\infty^I)$, and let $T_t$
($t\in\mathbf R$) and $T$ be operators on functions on $\mathbf A^I$.
\begin{enumerate}
\item If $T_t(0)=0$ for all $t$ and $(T_t)$ acts weakly at $\infty$ to first
  order by $f\mapsto f_1$ (Definition~\ref{asm-theta-agree-real-fock}(7)), then
  it acts at $\infty$ to first order by $f\mapsto f_1$
  (Definition~\ref{asm-theta-agree-models}(2)).
\item If $T(0)=0$ and $T$ acts weakly at $\infty$ by $f\mapsto f_1$, then it
  acts at $\infty$ by $f\mapsto f_1$.
\item If $(T_t)$ acts at $\infty$ to first order by $f\mapsto f_1$ and by
  $f\mapsto f_2$, then $f_1=f_2$. The same holds for $T$ acting at $\infty$, and
  for $(T_t)$ acting weakly at $\infty$ to first order.
\end{enumerate}
\end{lemma}
\leanref{Development/AbelianVarietyTheory/Proofs/DivisorProofs/IsFockMap.lean}{GlobalTheory.RealLift.ActsLaws}
\leanref{Development/WeilBound/CoreStatements/GaussianSeminormBounds/RealSeminormFiltration/RLActs.lean}{PositivityTheory.RL.actsLaws}
\begin{proof}
Four facts are used. (a) There are a Schwartz-Bruhat function $g_0$ on
$\mathbf A_f^I$ and $y_0$ with $g_0(y_0)=1$: the indicator function of a
product of compact open subgroups. (b) For $x\in\mathbf A^I$, evaluation at $x$
is an admissible linear form, since $(f\otimes g)(x)=g(x_f)f(x_\infty)$. (c) If
$f\neq0$ and $r$ is a function on $\mathbf A^I$ that does not depend on the
component at $\infty$, with $f\otimes r$ Schwartz-Bruhat, then
$r(x)=g(x_f)$ for a Schwartz-Bruhat function $g$ on $\mathbf A_f^I$: choose $x_0$
with $f(x_0)\neq0$ and put $g(y)=f(x_0)^{-1}(f\otimes r)(x_0,y)$, a
Schwartz-Bruhat function since restricting a Schwartz-Bruhat function to a
slice $\{x_0\}\times\mathbf A_f^I$ gives one. (d) If $f'\otimes g=f''\otimes g$
for every Schwartz-Bruhat $g$ on $\mathbf A_f^I$, then $f'=f''$, by (a).

(1) If $f=0$, then for every $g$ and $x$ the function
$t\mapsto T_t(0)(x)=0$ has derivative $(f_1\otimes g)(x)$ by (b), so
$f_1\otimes g=0$ and $f_1=0$ by (d); then both sides of the condition in
Definition~\ref{asm-theta-agree-models}(2) vanish. If $f\neq0$, let $r$ be as
in that condition; by (c), $f\otimes r=f\otimes g$ for a Schwartz-Bruhat $g$,
and evaluation at $x$ is admissible by (b), so $t\mapsto T_t(f\otimes r)(x)$
has derivative $(f_1\otimes r)(x)$. (2) is proved in the same way. (3) Take
$r(x)=g(x_f)$ for a Schwartz-Bruhat $g$ on $\mathbf A_f^I$; then
$f_1\otimes g=f_2\otimes g$ at every point, and $f_1=f_2$ by (d). For the weak
form, use evaluation at points, which is admissible by (b).
\end{proof}

\begin{lemma}\label{asm-theta-agree-real-derivative}
Let $B$ be a theta base for a CM field $E$, let $I$ be a finite set and
$n,m\geq1$, let
$e=(e_1,e_2)\colon M_{n\times m}(\mathbf A_E)\to\mathbf A^I\times\mathbf A^I$
be an $\mathbf A$-linear isomorphism, and let $\rho(u)$,
$u\in M_{n\times m}(\mathbf A_E)$, be the operators
$(\rho(u)\phi)(x)=\psi\bigl(x\cdot e_2(u)+\tfrac12e_1(u)\cdot e_2(u)\bigr)\phi(x+e_1(u))$.
For $u$ at $\infty$, let $a$ and $b$ be the components at $\infty$ of $e_1(u)$
and $e_2(u)$, and put $D_u=D_{a,b}$, the operator
$D_{a,b}f(x)=\partial_af(x)-2\pi i\,T(x,b)f(x)$ on $\mathcal S(F_\infty^I)$ of
Lemma~\ref{asm-theta-agree-density-derivative}. Then for every $u$ at $\infty$
and $f\in\mathcal S(F_\infty^I)$, the family $t\mapsto\rho(tu)$ acts weakly at
$\infty$ to first order by $f\mapsto D_uf$.
\end{lemma}
\leanref{Development/AbelianVarietyTheory/Proofs/DivisorProofs/IsFockMap.lean}{GlobalTheory.RealLift.DerEx}
\leanref{Development/WeilBound/CoreStatements/GaussianSeminormBounds/RealSeminormFiltration/EvIHalfTrace.lean}{PositivityTheory.RL.derEx}
\leanref{Development/WeilBound/CoreStatements/GaussianSeminormBounds/RealSeminormFiltration/EvIHalfTrace.lean}{PositivityTheory.RL.Der.schr_pure}
\leanref{Development/WeilBound/CoreStatements/GaussianSeminormBounds/RealSeminormFiltration/EvIHalfTrace.lean}{PositivityTheory.RL.Der.hl_deriv}
\begin{proof}
As $e$ is $\mathbf A$-linear and $u$ is at $\infty$, so are $e_1(u)$ and
$e_2(u)$. Let $g$ be a Schwartz-Bruhat function on $\mathbf A_f^I$. Since
$\psi(z)=\exp(-2\pi i\operatorname{Tr}_{F_\infty/\mathbf R}z_\infty)$ for $z$ at
$\infty$, $\rho(tu)(f\otimes g)=h_tf\otimes g$, where
$h_tf(y)=\exp\bigl(-2\pi i\bigl(t\,T(y,b)+\tfrac12t^2T(a,b)\bigr)\bigr)f(y+ta)$.
This is a Schwartz function, and $t\mapsto h_tf$ is differentiable at $0$ as a
map into $\mathcal S(F_\infty^I)$, with derivative $D_{a,b}f$: by Taylor's
formula with remainder, the difference quotients differ from $D_{a,b}f$ by
$t$ times functions whose Schwartz seminorms are bounded for $|t|\leq1$. An
admissible form $\lambda$ is continuous on $f'\mapsto\lambda(f'\otimes g)$, so
$t\mapsto\lambda(\rho(tu)(f\otimes g))$ has derivative
$\lambda(D_{a,b}f\otimes g)$ at $0$.
\end{proof}

\begin{lemma}\label{asm-theta-agree-real-exponential}
\begin{enumerate}
\item Let $\iota$ be a finite set and $L_i$ a locally finite operator on a
  complex vector space $V_i$ for $i\in\iota$. Then $\sum_iL_i^{(i)}$ is locally
  finite on $\bigotimes_iV_i$ (Definition~\ref{asm-theta-agree-real-fock}(4)).
\item Every locally finite operator $L$ on a complex vector space $V$ has a
  unique exponential. Let $L$ and $L'$ be locally finite operators on $V$ and
  $V'$ with exponentials $U$ and $U'$, let $q\colon V\to Z$ and
  $q'\colon V'\to Z$ be linear, and let $L_0$ be an operator on $Z$ with
  $qL=L_0q$ and $q'L'=L_0q'$. If $q(v)=q'(v')$, then $q(Uv)=q'(U'v')$.
\item Let $I$ be a finite set, $V$ a complex vector space,
  $\Phi\colon V\to\mathcal S(F_\infty^I)$ linear, and $L$ a locally finite
  operator on $V$ with exponential $U$. Let $\omega_t$, $t\in\mathbf R$, be
  operators on functions on $\mathbf A^I$ with $\omega_0=1$ and
  $\omega_{s+t}=\omega_s\omega_t$, such that $\lambda\circ\omega_t$ is admissible
  for every admissible linear form $\lambda$, and $(\omega_t)$ acts weakly at
  $\infty$ to first order by $\Phi(T)\mapsto\Phi(LT)$ for every $T\in V$. Then
  $\omega_1$ acts weakly at $\infty$ by $\Phi(T)\mapsto\Phi(UT)$ for every
  $T\in V$.
\end{enumerate}
\end{lemma}
\leanref{Development/AbelianVarietyTheory/Proofs/DivisorProofs/IsLocallyFinitePiTensorSumOfFactorOperators.lean}{GlobalTheory.RealLift.locFin_sumPl}
\leanref{Development/AbelianVarietyTheory/Proofs/DivisorProofs/WeakExponentialActionOnFockVectors.lean}{GlobalTheory.ExpAct.exists_isExp}
\leanref{Development/WeilBound/CoreStatements/GaussianSeminormBounds/RealSeminormFiltration/KRealAsm.lean}{PositivityTheory.KRealA.exp_equiv}
\leanref{Development/AbelianVarietyTheory/Proofs/DivisorProofs/WeakExponentialActionOnFockVectors.lean}{GlobalTheory.ExpAct.expAct_fock}
\begin{proof}
(1) It suffices to treat pure tensors, as a sum of two vectors lies in the sum
of the subspaces for each. For $\bigotimes_iv_i$, choose finite-dimensional
$L_i$-stable subspaces $N_i\ni v_i$; the image of $\bigotimes_iN_i$ is
finite-dimensional, contains $\bigotimes_iv_i$ and is stable under every
$L_i^{(i)}$.

(2) Let $v\in V$, and $N\ni v$ a finite-dimensional $L$-stable subspace. In a
basis of $N$, $L$ acts by a matrix $A$, and for a linear form $\ell$ the series
$\sum_n\ell(L^nv)/n!$ converges to $\ell$ applied to the vector with
coordinates $\exp(A)$ times those of $v$. This vector $Uv$ does not depend on
the choice of $N$, since linear forms separate points, and $v\mapsto Uv$ is
linear for the same reason. Uniqueness also follows. For the second claim, let
$\ell$ be a linear form on $Z$. Since $qL^n=L_0^nq$ and $q'L'^n=L_0^nq'$, the
series for $\ell(q(Uv))$ and $\ell(q'(U'v'))$ have the same terms
$\ell(L_0^nq(v))/n!$. So $\ell(q(Uv))=\ell(q'(U'v'))$ for every $\ell$.

(3) Let $T\in V$, let $N\ni T$ be a finite-dimensional $L$-stable subspace with
basis $e_1,\ldots,e_m$, let $A$ be the matrix of $L$ on $N$ and $y\in\mathbf C^m$
the coordinates of $T$, so that $UT$ has coordinates $\exp(A)y$. Fix a
Schwartz-Bruhat $g$ on $\mathbf A_f^I$ and $x\in\mathbf A^I$, and put
$a_s(z)=\omega_s\bigl(\Phi(\sum_jz_je_j)\otimes g\bigr)(x)$ for $z\in\mathbf C^m$.
Evaluation at $x$ composed with $\omega_s$ is admissible, by fact (b) in the
proof of Lemma~\ref{asm-theta-agree-real-weak} and the hypothesis, and
$\omega_{s+t}=\omega_s\omega_t$. So the weak action to first order gives
$\frac d{ds}a_s(z)=a_s(Az)$ for all $s$. Then
$s\mapsto\sum_jc_j(s)a_s(e_j)$, with $c(s)=\exp((1-s)A)y$, has derivative
$0$, and comparing $s=1$ and $s=0$ gives $a_1(y)=a_0(\exp(A)y)$, that is,
$\omega_1(\Phi(T)\otimes g)(x)=(\Phi(UT)\otimes g)(x)$.
\end{proof}

\begin{lemma}\label{asm-theta-agree-real-tensor}
Let $\iota$ be a finite set, and for $i\in\iota$ let $N_i$ be an irreducible
$(\mathfrak g,K)$-module of type $(R,S)$ with scalar endomorphisms
(Definition~\ref{asm-theta-agree-real-fock}(2)). Then
$\Pi=\bigotimes_iN_i$ is nonzero, and $0$ and $\Pi$ are its only subspaces
stable under $N_i(Z)^{(i)}$ for every $i\in\iota$ and $Z\in\mathfrak u(R,S)$.
\end{lemma}
\leanref{Development/AbelianVarietyTheory/Proofs/DivisorProofs/TensorIrreducibleModules.lean}{GlobalTheory.TensIrr.tensor_irr}
\leanref{Development/AbelianVarietyTheory/Proofs/DivisorProofs/TensorIrreducibleModules.lean}{GlobalTheory.TensIrr.irr_of_stable}
\leanref{Development/AbelianVarietyTheory/Proofs/DivisorProofs/TensorIrreducibleModules.lean}{GlobalTheory.TensIrr.stable_all}
\begin{proof}
Choose $v_i\neq0$ in $N_i$ and linear forms $\ell_i$ with $\ell_i(v_i)=1$; then
$(\bigotimes_i\ell_i)(\bigotimes_iv_i)=1$, so $\Pi\neq0$. Let $N\subset\Pi$ be
a stable subspace. Every $Z\in M_{R+S}(\mathbf C)$ is $Z_1+iZ_2$ with
$Z_1,Z_2\in\mathfrak u(R,S)$, namely
$Z_1=\tfrac12(Z-J_{R,S}Z^*J_{R,S})$ and
$Z_2=-\tfrac i2(Z+J_{R,S}Z^*J_{R,S})$; so $N$ is stable under $N_i(Z)^{(i)}$
for every $Z$. Fix $i$ and an operator $f$ on $N_i$. An element $y$ of $\Pi$
involves finitely many vectors of $N_i$ in the $i$-th factor. By the Jacobson
density theorem \cite[\S11]{Lam01}, as $N_i$ is irreducible with scalar
endomorphisms, some element of the algebra generated by the $N_i(Z)$ agrees
with $f$ on these vectors; so $f^{(i)}y\in N$ if $y\in N$. Composing, $N$ is
stable under $\bigotimes_if_i$ for all operators $f_i$. Let $0\neq y\in N$.
Choosing bases, there are linear forms $\ell_i$ on $N_i$ with
$c=(\bigotimes_i\ell_i)(y)\neq0$. For any $u_i\in N_i$, the operators
$f_i(x)=\ell_i(x)u_i$ give $(\bigotimes_if_i)(y)=c\bigotimes_iu_i\in N$. So
$N$ contains every pure tensor, and $N=\Pi$.
\end{proof}

\begin{lemma}\label{asm-theta-agree-real-fock-map}
Let $B$ be a theta base for a CM field $E$, let $X$ and $W$ be lattice spaces
of ranks $n$ and $d$ with matrices $J_X$ and $J_W$, and put
$I=\{1,\ldots,n\}\times\{1,\ldots,d\}$. Let
$e=(e_1,e_2)\colon\mathbb W(X,W)\to\mathbf A^I\times\mathbf A^I$ be an
$\mathbf A$-linear isomorphism with
$\langle w,w'\rangle=e_1(w)\cdot e_2(w')-e_1(w')\cdot e_2(w)$, for instance
the isomorphism of a model for $(X,W)$, and let $\rho$ be its Schrödinger
operators (Definition~\ref{asm-theta-agree-models}(3)). Let $P,Q,R,S\geq0$ and
$a_0,b_0\in\mathbf C$ with $a_0b_0=i$, and for every infinite place $w$ of $F$
let $A_w$ be a frame of $a_0\tau_w(J_X)$ of type $(P,Q)$ and $C_w$ a frame of
$b_0\tau_w(J_W)$ of type $(R,S)$
(Definition~\ref{asm-theta-agree-real-fock}(5)). Let
$\mathcal T=\bigotimes_w\mathcal P(P,Q;R,S)$, the tensor product over the
infinite places, and for a family $c=(c_w)_w$ of matrices
$c_w\in M_{(P+Q)\times(R+S)}(\mathbf C)$ put $h(c)=\sum_wh(c_w)^{(w)}$
(Definition~\ref{asm-theta-agree-real-fock}(1) and (4)).
\begin{enumerate}
\item There is an injective linear map $\Phi\colon\mathcal T\to\mathcal S(F_\infty^I)$
  such that for every $u\in\mathbb W(X,W)$ at $\infty$ and every family $c$
  with $\pi_w(u)=A_wc_wC_w^{\mathsf T}$ for all $w$, the family
  $t\mapsto\rho(tu)$ acts weakly at $\infty$ to first order by
  $\Phi(T)\mapsto\Phi(h(c)T)$ for every $T\in\mathcal T$
  (Definition~\ref{asm-theta-agree-real-fock}(7)). Such a $\Phi$ is a
  \emph{Fock map} for $(e,(A_w),(C_w))$; if $e$ is the isomorphism of a model
  $M$, it is a \emph{Fock map} for $(M,(A_w),(C_w))$.
\item Let $M$ be a model for $(X,W)$ with isomorphism $e$. If $a_0=1$,
  $b_0=i$, and $\theta_{X,w}A_w=A_wJ_{P,Q}$ and
  $\theta_{W,w}C_w=C_wJ_{R,S}$ for every $w$, then the image of every Fock map
  for $(M,(A_w),(C_w))$ is $\operatorname{Pol}(M)$
  (Definition~\ref{asm-theta-agree-models}(4)).
\end{enumerate}
\end{lemma}
\leanref{Development/AbelianVarietyTheory/Proofs/DivisorProofs/IsFockMap.lean}{GlobalTheory.RealLift.IsFockMap}
\leanref{Development/AbelianVarietyTheory/Proofs/DivisorProofs/IsFockMap.lean}{GlobalTheory.RealLift.FockEx}
\leanref{Development/AbelianVarietyTheory/Proofs/DivisorProofs/IsFockMap.lean}{GlobalTheory.RealLift.VacLine}
\leanref{Development/AbelianVarietyTheory/Proofs/DivisorProofs/IsFockMap.lean}{GlobalTheory.RealLift.PolRange}
\leanref{Development/WeilBound/CoreStatements/GaussianSeminormBounds/RealSeminormFiltration/RLFock.lean}{PositivityTheory.RL.fockEx}
\leanref{Development/WeilBound/CoreStatements/GaussianSeminormBounds/RealSeminormFiltration/RLFock.lean}{PositivityTheory.RL.vacLine}
\leanref{Development/WeilBound/CoreStatements/GaussianSeminormBounds/RealSeminormFiltration/RLFock.lean}{PositivityTheory.RL.Fock.exists_heis}
\leanref{Development/WeilBound/CoreStatements/GaussianSeminormBounds/RealSeminormFiltration/RLFock.lean}{PositivityTheory.RL.Fock.Heis.exists_vac}
\leanref{Development/WeilBound/CoreStatements/GaussianSeminormBounds/RealSeminormFiltration/RLFock.lean}{PositivityTheory.RL.Fock.Heis.exists_fockMap}
\leanref{Development/WeilBound/CoreStatements/GaussianSeminormBounds/RealSeminormFiltration/RLFock.lean}{PositivityTheory.RL.Fock.Heis.eq_zero_of_neg}
\leanref{Development/WeilBound/ThetaLifts/EigenvalueEndgame/EigenKilling/PolRange.lean}{PositivityTheory.ThKEnd.Comp.PR.polRange_of}
\begin{proof}
Write $V=F_\infty^I$, and let $\Lambda$ be the set of pairs $(w,l)$ of an
infinite place $w$ and $l\in\mathrm{Ix}(P,Q)\times\mathrm{Ix}(R,S)$.

(1) \emph{Coordinates.} For $z\in\mathbf C^\Lambda$ let $c(z)$ be the family
with $c(z)_{w,l}=z_{w,l}$ if $s_l=1$ and $c(z)_{w,l}=\bar z_{w,l}$ if
$s_l=-1$; every family $c$ is $c(z)$ for exactly one $z$. As $E$ is a CM field,
there is a unique $u(z)\in\mathbb W(X,W)$ at $\infty$ with
$\pi_w(u(z))=A_wc(z)_wC_w^{\mathsf T}$ for every $w$, and $z\mapsto u(z)$ is
$\mathbf R$-linear (fact (iii) in the proof of
Lemma~\ref{asm-theta-agree-real-lines}). Let $v(z)$ and $y(z)$ in $V$ be the
components at $\infty$ of $e_1(u(z))$ and $e_2(u(z))$, and put
$D_z=D_{u(z)}$, the operator of Lemma~\ref{asm-theta-agree-real-derivative};
so $D_zf(x)=\partial_{v(z)}f(x)-2\pi i\,T(x,y(z))f(x)$, and $z\mapsto D_z$ is
$\mathbf R$-linear. By that lemma, $t\mapsto\rho(tu(z))$ acts weakly at
$\infty$ to first order by $f\mapsto D_zf$.

\emph{The form.} For $z,z'\in\mathbf C^\Lambda$,
\[
T(v(z),y(z'))-T(v(z'),y(z))=-2\operatorname{Im}\sum_{\lambda\in\Lambda}z_\lambda\bar z'_\lambda.
\]
Indeed, by the condition on $e$ the left side is
$\operatorname{Tr}_{F_\infty/\mathbf R}$ of the component at $\infty$ of
$\langle u(z),u(z')\rangle=-\operatorname{Tr}_{\mathbf A_E/\mathbf A}h(u(z),u(z'))$,
and for $x$ at $\infty$ this is $-\sum_w2\operatorname{Re}\pi_w(h(u(z),u(z')))$.
Put $x=\pi_w(u(z))$ and $x'=\pi_w(u(z'))$. Since $E$ is a CM field,
$\pi_w(h(u(z),u(z')))=\operatorname{tr}(x^{\mathsf T}\tau_w(J_X)\bar x'\tau_w(J_W)^{\mathsf T})$.
Substituting $x=A_wc_wC_w^{\mathsf T}$ and $x'=A_wc'_wC_w^{\mathsf T}$, with
$c=c(z)$ and $c'=c(z')$, and using $A_w^{\mathsf T}a_0\tau_w(J_X)\bar A_w=J_{P,Q}$,
$\bar C_w^{\mathsf T}b_0\tau_w(J_W)^{\mathsf T}C_w=J_{R,S}$ (the transpose of
the frame condition for $C_w$) and $a_0b_0=i$, this trace is
$-i\sum_ls_lc_{w,l}\bar c'_{w,l}$, whose real part is
$\operatorname{Im}\sum_lz_{w,l}\bar z'_{w,l}$.

\emph{Commutators.} By the symmetry of second derivatives and the form,
\[
D_zD_{z'}-D_{z'}D_z=4\pi i\operatorname{Im}\Bigl(\sum_\lambda z_\lambda\bar z'_\lambda\Bigr).
\]
The kernel of $v$ has dimension at least $\dim V$: for $y_0\in V$, the element
$u$ with $e(u)=(0,y_0)$, $y_0$ placed at $\infty$, is at $\infty$, and is
$u(z)$ for some $z$ with $v(z)=0$ and $y(z)=y_0$. If $v(k)=0$ and $k\neq0$, the
form with $z=k$ and $z'=ik$ gives $T(v(ik),y(k))=-2\sum_\lambda|k_\lambda|^2<0$.
So $k\mapsto v(ik)$ is injective on $\ker v$, hence an isomorphism onto $V$;
let $\sigma\colon V\to\ker v$ be its inverse.

\emph{The vacuum.} Call $f\in\mathcal S(V)$ a \emph{vacuum} if
$D_{iz}f=iD_zf$ for every $z$. Applied to $z=\sigma(y)$, this gives
\[
\partial_yf(x)=2\pi\bigl(T(x,y(\sigma(y)))+i\,T(x,y(i\sigma(y)))\bigr)f(x)
\qquad(x,y\in V),
\]
and conversely this equation implies that $f$ is a vacuum, by a direct
computation using $\sigma(v(iz))=z-i\sigma(v(z))$ (the right side lies in
$\ker v$, and $k\mapsto v(ik)$ maps it to $v(iz)$). The forms
$\beta(y,x)=\pi T(x,y(\sigma(y)))$ and $\beta'(y,x)=\pi T(x,y(i\sigma(y)))$
are symmetric, by the form applied to $(i\sigma(y),\sigma(x))$,
$(i\sigma(x),\sigma(y))$, $(i\sigma(y),i\sigma(x))$ and $(\sigma(y),\sigma(x))$,
and $\beta(x,x)=\pi T(v(i\sigma(x)),y(\sigma(x)))<0$ for $x\neq0$. So
$G(x)=\exp(\beta(x,x)+i\beta'(x,x))$ is a nonzero Schwartz function satisfying
the equation, hence a vacuum; and every vacuum $f$ is a multiple of $G$, since
$f(x)\exp(-\beta(x,x)-i\beta'(x,x))$ has derivative $0$.

\emph{The Fock map.} Put $r=\sqrt{2\pi}$, let $e_\lambda$ be the unit vectors of
$\mathbf C^\Lambda$, and put $Z_\lambda=(2r)^{-1}(D_{e_\lambda}-iD_{ie_\lambda})$
and $d_\lambda=-(2r)^{-1}(D_{e_\lambda}+iD_{ie_\lambda})$. By the commutator
formula the $Z_\lambda$ commute with each other and
$d_\lambda Z_\mu=Z_\mu d_\lambda+\delta_{\lambda\mu}$, and $d_\lambda G=0$ since
$G$ is a vacuum. Let $\varphi(p)=p(Z)G$ for $p\in\mathbf C[x_\lambda:\lambda\in\Lambda]$,
the polynomial evaluated at the commuting operators $Z_\lambda$. Then
$\varphi(x_\lambda p)=Z_\lambda\varphi(p)$, and
$\varphi(\partial p/\partial x_\lambda)=d_\lambda\varphi(p)$ by induction on
$p$. So $\ker\varphi$ is stable under the derivations
$\partial/\partial x_\lambda$ and contains no nonzero constant, as $G\neq0$;
hence $\varphi$ is injective. Since
$D_{e_\lambda}=r(Z_\lambda-d_\lambda)$, $D_{ie_\lambda}=ir(Z_\lambda+d_\lambda)$
and $z\mapsto D_z$ is $\mathbf R$-linear,
$D_z=r\sum_\lambda(z_\lambda Z_\lambda-\bar z_\lambda d_\lambda)$, and therefore
$D_z\varphi(p)=\varphi(\eta(z)p)$ with
$\eta(z)=r\sum_\lambda(z_\lambda x_\lambda-\bar z_\lambda\,\partial/\partial x_\lambda)$.

Let $j\colon\mathcal T\to\mathbf C[x_\lambda:\lambda\in\Lambda]$ be the
injective algebra homomorphism sending $x_l$ in the factor of $w$ to
$x_{(w,l)}$. On the factor of $w$, $h(c(z)_w)$ is
$r\sum_l(z_{w,l}x_l-\bar z_{w,l}\,\partial/\partial x_l)$: for $s_l=1$ the term
$c_{w,l}R_l-s_l\bar c_{w,l}L_l$ is
$c_{w,l}x_l-\bar c_{w,l}\,\partial/\partial x_l$ with $c_{w,l}=z_{w,l}$, and for
$s_l=-1$ it is $\bar c_{w,l}x_l-c_{w,l}\,\partial/\partial x_l$ with
$\bar c_{w,l}=z_{w,l}$. Hence $j(h(c(z))T)=\eta(z)j(T)$.
Put $\Phi=\varphi\circ j$, which is injective. If $u$ is at $\infty$ with
$\pi_w(u)=A_wc_wC_w^{\mathsf T}$ for all $w$, write $c=c(z)$; then $u=u(z)$ by
uniqueness, and $t\mapsto\rho(tu)$ acts weakly at $\infty$ to first order by
$\Phi(T)\mapsto D_z\Phi(T)=\varphi(\eta(z)j(T))=\Phi(h(c)T)$.

(2) Let $\Phi$ be a Fock map, and write $\mathbf 1=1\otimes\cdots\otimes1$.
Since $\rho(tu)0=0$, Lemma~\ref{asm-theta-agree-real-weak} shows that for $u$
at $\infty$ with $\pi_w(u)=A_wc_wC_w^{\mathsf T}$, the family
$t\mapsto\rho(tu)$ acts at $\infty$ to first order by
$\Phi(T)\mapsto\Phi(h(c)T)$, and by no other map on $\Phi(T)$. Every $u$ at
$\infty$ has this form, with $c_w=A_w^{-1}\pi_w(u)(C_w^{\mathsf T})^{-1}$.
From $\theta_{X,w}A_w=A_wJ_{P,Q}$ and $\theta_{W,w}C_w=C_wJ_{R,S}$,
\[
i\,\theta_{X,w}A_wc_wC_w^{\mathsf T}\theta_{W,w}^{\mathsf T}=A_w\bigl(iJ_{P,Q}c_wJ_{R,S}\bigr)C_w^{\mathsf T}.
\]
Also $h(iJ_{P,Q}cJ_{R,S})1=i\,h(c)1$ in $\mathcal P(P,Q;R,S)$: both sides are
$\sqrt{2\pi}\,i$ times $\sum_{s_l=1}c_lx_l+\sum_{s_l=-1}\bar c_lx_l$.

$\operatorname{Pol}(M)\subset\Phi(\mathcal T)$: the image of $\Phi$ contains
every Gaussian, and it contains $f_1$ whenever it contains $f$ and
$t\mapsto\rho(tu)$, for some $u$ at $\infty$, acts at $\infty$ to first order
by $f\mapsto f_1$, by the first paragraph. Let $f_0$ be a Gaussian with signs
$\varepsilon_w$. By the displayed identity and
Lemmas~\ref{asm-theta-agree-real-derivative}
and~\ref{asm-theta-agree-real-weak}, $D_{z'}f_0=iD_zf_0$ for every $z$, where
$z'_{w,l}=\varepsilon_wiz_{w,l}$. If every $\varepsilon_w$ is $1$, then $f_0$
is a vacuum. So is $\Phi(\mathbf 1)$: by the first paragraph
$D_z\Phi(T)=\Phi(h(c(z))T)$ for all $z$ and $T$, and
$h(c(iz))\mathbf 1=i\,h(c(z))\mathbf 1$, because $j$ is injective and
$j(h(c(z))\mathbf 1)=\eta(z)1=r\sum_\lambda z_\lambda x_\lambda$. As $\Phi$ is
injective, $\Phi(\mathbf 1)\neq0$, so $f_0$ is a multiple of $\Phi(\mathbf 1)$.
If some $\varepsilon_w=-1$, take $z=e_\lambda$ with $\lambda=(w,l)$; then
$D_{ie_\lambda}f_0=-iD_{e_\lambda}f_0$. The operators $D_z$ are skew-adjoint
for the $L^2$ inner product on $\mathcal S(V)$, so the commutator formula,
paired with $f_0$, gives
$2i\|D_{e_\lambda}f_0\|^2=-4\pi i\|f_0\|^2$, and $f_0=0$.

$\Phi(\mathcal T)\subset\operatorname{Pol}(M)$: by the first paragraph and the
displayed identity, $\Phi(\mathbf 1)$ is a Gaussian with all signs $1$, so it
lies in $\operatorname{Pol}(M)$. Let $\mathcal N$ be the set of $T$ with
$\Phi(T)\in\operatorname{Pol}(M)$. It contains $\mathbf 1$, and it is stable
under $h(c_w)^{(w)}$ for every $w$ and $c_w$: apply the first paragraph to $u$
at $\infty$ with $\pi_w(u)=A_wc_wC_w^{\mathsf T}$ and $\pi_{w'}(u)=0$ for
$w'\neq w$, which exists by fact (iii) in the proof of
Lemma~\ref{asm-theta-agree-real-lines}. Multiplication by $x_l$ is $R_l$ or
$L_l$, and $R_l=(2r)^{-1}(h(E_l)-ih(iE_l))$ and
$L_l=-(2r)^{-1}s_l(h(E_l)+ih(iE_l))$, where $E_l$ is the matrix unit at $l$.
So $\mathcal N$ contains $\mathbf 1$ and is stable under multiplication by
every variable in every factor; hence it contains every pure tensor of
monomials, and $\mathcal N=\mathcal T$.
\end{proof}

\begin{lemma}\label{asm-theta-agree-real-metaplectic}
Let $l\geq0$, give $\mathbf R^l$ its standard inner product, and write the
elements of $\mathbf R^{2l}=\mathbf R^l\times\mathbf R^l$ as $w=(p,q)$. For
$w\in\mathbf R^{2l}$ and $f\in\mathcal S(\mathbf R^l)$ put
$\rho_S(w)f(x)=\exp\bigl(2\pi i(\langle q,x\rangle+\tfrac12\langle p,q\rangle)\bigr)f(x+p)$
and $d\rho(w)f(x)=\partial_pf(x)+2\pi i\langle q,x\rangle f(x)$, and put
$\operatorname{sympl}(w,w')=\langle p,q'\rangle-\langle q,p'\rangle$. Let
$\mathfrak{sp}$ be the Lie algebra of the $X\in M_{2l}(\mathbf R)$ with
$\operatorname{sympl}(Xw,w')+\operatorname{sympl}(w,Xw')=0$. An operator $U$ on
$\mathcal S(\mathbf R^l)$ \emph{implements} $s\in M_{2l}(\mathbf R)$ if it
preserves the $L^2$-norm and $U\rho_S(w)=\rho_S(sw)U$ for every $w$.
\begin{enumerate}
\item For every $X\in M_{2l}(\mathbf R)$ there are finitely many
  $u_j,v_j\in\mathbf R^{2l}$ with $Xw=\sum_j\operatorname{sympl}(v_j,w)u_j$ for
  all $w$, and the operator
  $\operatorname{quad}(X)=(8\pi i)^{-1}\sum_j\bigl(d\rho(u_j)d\rho(v_j)+d\rho(v_j)d\rho(u_j)\bigr)$
  does not depend on their choice.
\item Let $X\in\mathfrak{sp}$. There are $\varepsilon>0$, $c\in\mathbf C$ and
  continuous linear operators $V(t)$ and $W(t)$ on $\mathcal S(\mathbf R^l)$,
  for $t\in\mathbf R$, with $V(0)=1$, such that for $|t|<\varepsilon$:
  $V(t)W(t)=W(t)V(t)=1$; $V(t)$ implements $\exp(tX)$; for every $f$ and
  continuous linear form $\lambda$, $u\mapsto\lambda(V(u)f)$ is differentiable
  at $t$; and its derivative at $0$ is $\lambda(\operatorname{quad}(X)f+cf)$.
\item If $s$ preserves $\operatorname{sympl}$, $ss'=1$, and $U$ implements
  $s$, then $U\,d\rho(w)=d\rho(sw)\,U$ for every $w$, and
  $U\operatorname{quad}(X)=\operatorname{quad}(sXs')\,U$ for every $X$.
\end{enumerate}
\end{lemma}
\leanref{Development/AbelianVarietyTheory/PublishedTheorems/TwistedTraceFormula/RealMetaplecticRepresentation.lean}{GlobalTheory.RealWeil.quad}
\leanref{Development/WeilBound/CoreStatements/GaussianSeminormBounds/RealSeminormFiltration/RLLieB.lean}{PositivityTheory.RL.Up.quad_eq}
\leanref{Development/WeilBound/ThetaLifts/Reductions/MetaDef.lean}{PositivityTheory.ThKR.LocalLift}
\leanref{Development/WeilBound/ThetaLifts/Reductions/LocalLift.lean}{PositivityTheory.ThKR.localLift_holds}
\leanref{Development/WeilBound/ThetaLifts/Reductions/RealSmulL2Coe.lean}{PositivityTheory.ThKR.Im.covD}
\leanref{Development/WeilBound/ThetaLifts/Reductions/RealSmulL2Coe.lean}{PositivityTheory.ThKR.Im.quad}
\begin{proof}
(1) The vectors $u_b=XJe_b$ and $v_b=e_b$, with $e_b$ the unit vectors and $J$
the matrix with $\operatorname{sympl}(w,w')=-w^{\mathsf T}Jw'$, satisfy the
condition. If $(u_j,v_j)$ satisfy it, then $\sum_ju_j(a)v_j(b)=(XJ)_{ab}$ for
all $a,b$, and expanding $d\rho(u_j)$ and $d\rho(v_j)$ in the unit vectors
shows that the sum depends only on these numbers. Part (2) is the existence of
the metaplectic representation near the identity, with its infinitesimal
action by quadratic operators \cite{Wei64}, \cite[Chapter~4]{Fol89}. For (3),
differentiating $U\rho_S(tw)=\rho_S(tsw)U$ at $t=0$, evaluated at each point,
gives $U\,d\rho(w)=d\rho(sw)\,U$. If $Xw=\sum_j\operatorname{sympl}(v_j,w)u_j$,
then $sXs'w=\sum_j\operatorname{sympl}(sv_j,w)su_j$, since $s$ preserves
$\operatorname{sympl}$; so (1) and the first identity give the second.
\end{proof}

\begin{lemma}\label{asm-theta-agree-real-projective}
Let $\iota$ be a finite set, $m\geq1$, and for $w\in\iota$ let
$K_w\in M_m(\mathbf C)$ be Hermitian and invertible. Let
$\mathfrak A=\prod_{w\in\iota}M_m(\mathbf C)$, let $G$ be the group of
$g\in\mathfrak A^\times$ with $g_w^*K_wg_w=K_w$ for all $w$, and $\mathfrak g$
the real Lie algebra of the $Z\in\mathfrak A$ with $Z_w^*K_w+K_wZ_w=0$ for all
$w$. Let $l\geq0$, use the notation of
Lemma~\ref{asm-theta-agree-real-metaplectic}, and let
$\rho_{\mathbf R}\colon\mathfrak A\to M_{2l}(\mathbf R)$ be an $\mathbf R$-algebra
homomorphism such that $\rho_{\mathbf R}(g)$ preserves $\operatorname{sympl}$
for every $g\in G$. Let $H\subset G$ be the subgroup of the $g$ such that
$\rho_{\mathbf R}(g)$ is implemented by an invertible operator, and let $A$ be a
homomorphism from $H$ to the continuous linear operators on
$\mathcal S(\mathbf R^l)$ such that $A(g)$ is a scalar multiple of every
invertible operator implementing $\rho_{\mathbf R}(g)$, and $g\mapsto\lambda(A(g)f)$
is continuous for every $f$ and continuous linear form $\lambda$. Then there
are $c_w\in\mathbf C$ such that for every $Z\in\mathfrak g$, $\exp(tZ)\in H$ for
all $t$, and $t\mapsto\lambda(A(\exp(tZ))f)$ has derivative
$\lambda\bigl(\operatorname{quad}(\rho_{\mathbf R}(Z))f+(\sum_wc_w\operatorname{tr}Z_w)f\bigr)$
at $0$, for every $f$ and continuous linear form $\lambda$.
\end{lemma}
\leanref{Development/WeilBound/ThetaLifts/Reductions/SymplOfMem.lean}{PositivityTheory.ThKR.mid2}
\leanref{Development/WeilBound/ThetaLifts/Reductions/SymplOfMem.lean}{PositivityTheory.ThKR.scalar_lie}
\leanref{Development/WeilBound/ThetaLifts/Reductions/RealSmulL2Coe.lean}{PositivityTheory.ThKR.exists_im}
\leanref{Development/WeilBound/ThetaLifts/Reductions/RealSmulL2Coe.lean}{PositivityTheory.ThKR.line}
\leanref{Development/WeilBound/ThetaLifts/Reductions/RealSmulL2Coe.lean}{PositivityTheory.ThKR.delta_zero}
\leanref{Development/WeilBound/ThetaLifts/Reductions/RealSmulL2Coe.lean}{PositivityTheory.ThKR.strong}
\leanref{Development/WeilBound/CoreStatements/GaussianSeminormBounds/RealSeminormFiltration/RLLieC.lean}{PositivityTheory.RL.Up.real_form_trace}
\begin{proof}
Write $\|\cdot\|$ for the $L^2$-norm on $\mathcal S(\mathbf R^l)$. By hypothesis
each $A(g)$ is a scalar multiple of an operator preserving $\|\cdot\|$.

(i) Let $Z\in\mathfrak g$ and $X=\rho_{\mathbf R}(Z)$. Then $\exp(tZ)\in G$,
$\rho_{\mathbf R}(\exp(tZ))=\exp(tX)$, and $X\in\mathfrak{sp}$, since
$\exp(tX)$ preserves $\operatorname{sympl}$ for all $t$. By
Lemma~\ref{asm-theta-agree-real-metaplectic}(2), $\exp(sX)$ is implemented by
the invertible operator $V(s)$ for $|s|<\varepsilon$; for $n>|t|/\varepsilon$,
$\exp(tX)=\exp((t/n)X)^n$ is implemented by $V(t/n)^n$. So $\exp(tZ)\in H$.

(ii) Put $S(t)=A(\exp(tZ))$, a one-parameter group. For $|t|<\varepsilon$,
$S(t)=c(t)V(t)$ for a scalar $c(t)$. Fix $f_0\neq0$ and let
$\lambda_0(f)=\langle f_0,f\rangle_{L^2}$. The operators $S(s)$, $s\in[-1,1]$,
are uniformly bounded for $\|\cdot\|$, by the Banach-Steinhaus theorem applied
to the weakly continuous functions $s\mapsto S(s)f_0$ and the fact that each
$S(s)$ is a multiple of an isometry. Choose $\delta\in[-1,1]$ with
$\mu(f_0)\neq0$, where $\mu(f)=\int_0^\delta\lambda_0(S(s)f)\,ds$; this is
possible as $\lambda_0(f_0)\neq0$. Then $\mu$ is continuous, and
$\mu(S(t)f_0)=\int_t^{t+\delta}\lambda_0(S(s)f_0)\,ds$ is differentiable in $t$.
So $c(t)=\mu(S(t)f_0)/\mu(V(t)f_0)$ is differentiable at $0$ with $c(0)=1$, and
$\lambda(S(t)f)=c(t)\lambda(V(t)f)$ has derivative
$\lambda(\operatorname{quad}(X)f)+(c'(0)+c)\lambda(f)$ at $0$. Hence there is
$z(Z)\in\mathbf C$ such that $t\mapsto\lambda(S(t)f)$ has derivative
$\lambda(\operatorname{quad}(X)f+z(Z)f)$ for all $f$ and $\lambda$; it is unique,
by $f=f_0$ and $\lambda=\lambda_0$.

(iii) $z$ is $\mathbf R$-linear. Reparametrizing gives $z(rZ)=rz(Z)$, as
$\operatorname{quad}$ is linear. A one-parameter group $S(t)$ with these
derivatives is differentiable at $0$ also for $\|\cdot\|$: the weak derivative
exists at every $t$, and with $Q=\operatorname{quad}(X)+z(Z)$, Taylor's
formula and a uniform bound $M$ on $\|S(s)Q^2f\|$ for $s\in[-1,1]$, from the
Banach-Steinhaus theorem, give $\|S(t)f-f-tQf\|\leq Mt^2$ for $|t|\leq1$. Let $X,Y\in\mathfrak g$
and $k(t)=\exp(tX)\exp(tY)\exp(-t(X+Y))$. By the product rule for these
derivatives, $t\mapsto A(k(t))f$ has derivative $df$ at $0$ for $\|\cdot\|$,
where $d=z(X)+z(Y)+z(-X-Y)$, since $\operatorname{quad}$ is linear. So
$\|A(k(t))f-e^{dt}f\|=o(t)$, and $A(k(t))$ has operator norm at most $1+C|t|$
near $0$; a telescoping estimate gives $A(k(1/n))^nf-e^df\to0$. Since $k$ has
derivative $0$ at $0$, $k(1/n)^n\to1$, and by the continuity of $A$,
$\lambda_0(A(k(1/n)^n)f_0)\to\lambda_0(f_0)$. Hence $e^d=1$. The same holds
for $tX$ and $tY$, so $e^{td}=1$ for all $t$, and $d=0$; with $z(-W)=-z(W)$ this
gives $z(X+Y)=z(X)+z(Y)$.

(iv) $z$ vanishes on commutators. Let $X,Y\in\mathfrak g$ and $g=\exp(sX)\in H$,
and write $A(g)=aU$ with $U$ implementing $\rho_{\mathbf R}(g)$. Then
$A(g)A(\exp(tY))A(g)^{-1}=A(\exp(tgYg^{-1}))$, and
$U\operatorname{quad}(\rho_{\mathbf R}(Y))U^{-1}=\operatorname{quad}(\rho_{\mathbf R}(gYg^{-1}))$
by Lemma~\ref{asm-theta-agree-real-metaplectic}(3). Differentiating at $t=0$
and using the uniqueness in (ii), $z(gYg^{-1})=z(Y)$. Differentiating
$s\mapsto z(\exp(sX)Y\exp(-sX))$ at $s=0$ gives $z(XY-YX)=0$.

(v) Fix $w$, and let $\zeta_w(Z)=z(Z')$ for $Z\in M_m(\mathbf C)$ with
$Z^*K_w+K_wZ=0$, where $Z'$ has component $Z$ at $w$ and $0$ elsewhere. This
is a real linear form vanishing on commutators. Every $A\in M_m(\mathbf C)$ is
$P(A)-iP(iA)$ with $P(A)=\tfrac12(A-K_w^{-1}A^*K_w)$ in this real Lie algebra,
so $\xi(A)=\zeta_w(P(A))-i\zeta_w(P(iA))$ is a complex linear form on
$M_m(\mathbf C)$ extending $\zeta_w$, and it vanishes on commutators. Such a
form vanishes on $E_{ab}$ for $a\neq b$ and takes the same value $c_w$ on every
$E_{aa}$, so $\xi=c_w\operatorname{tr}$. As $Z=\sum_wZ'_w$, where $Z'_w$ keeps the
component of $Z$ at $w$, $z(Z)=\sum_wc_w\operatorname{tr}Z_w$.
\end{proof}
% Lean: mid2: exists_im (from LocalLift) for (i); scalar_lie with line (exists_avg, bd_of_conf,
% bdd_of_weak) for (ii), strong (taylor2), dif_mul, norm_le_lin, tele_lim, pow_lim and
% delta_zero for (iii), Im.quad and the conjugation hAd' for (iv); real_form_trace (eq_trace)
% for (v), with lie_of_exp and exp_preserves.

\begin{lemma}\label{asm-theta-agree-real-oscillator}
Let $B$ be a theta base for a CM field $E$, let $n,m\geq1$, $J_1\in M_n(E)$
and $J_2\in M_m(E)$, and let $U_2$ be the group of $g\in GL_m(\mathbf A_E)$ with
$g^{\mathsf T}J_2c(g)=J_2$; a \emph{line} in $U_2$ is a continuous homomorphism
$\gamma\colon\mathbf R\to U_2$ with $\gamma(t)-1$ at $\infty$ for every $t$.
Let $I$ be a finite set and
$e=(e_1,e_2)\colon M_{n\times m}(\mathbf A_E)\to\mathbf A^I\times\mathbf A^I$
an $\mathbf A$-linear isomorphism with
$\langle w,w'\rangle=e_1(w)\cdot e_2(w')-e_1(w')\cdot e_2(w)$, where
$\langle x,y\rangle=-\operatorname{Tr}_{\mathbf A_E/\mathbf A}\sum_{i,j,k,l}x_{ij}(J_1)_{ik}(J_2)_{jl}c(y_{kl})$;
let $\rho$ and $D_u$ be as in Lemma~\ref{asm-theta-agree-real-derivative}, and
for $u,v$ at $\infty$ put
$\sigma(u,v)=-\operatorname{Tr}_{F_\infty/\mathbf R}\langle u,v\rangle_\infty$.
Let $\omega_2$ be a representation of $U_2$ on functions on $\mathbf A^I$ whose
operators preserve the Schwartz-Bruhat functions and the admissible linear
forms (Definition~\ref{asm-theta-agree-doubling}(6)), satisfy
$\omega_2(g)\rho(w)=\rho(wg^{\mathsf T})\omega_2(g)$ on Schwartz-Bruhat
functions, and preserve the $L^2$-norm of every Schwartz-Bruhat function for
every Haar measure, and such that $g\mapsto\lambda(\omega_2(g)\phi)$ is
continuous for admissible $\lambda$ and Schwartz-Bruhat $\phi$. For every
infinite place $w$ let $K_w\in M_m(\mathbf C)$ be Hermitian and invertible,
with $g^*K_wg=K_w$ if and only if $g^{\mathsf T}\tau_w(J_2)\bar g=\tau_w(J_2)$,
for $g\in M_m(\mathbf C)$. Then there are $c_w\in\mathbf C$ with the following
property. Let $\gamma$ be a line in $U_2$ and $Z_w\in M_m(\mathbf C)$ with
$\pi_w(\gamma(t))=\exp(tZ_w)$ for all $t$ and $w$, and let $(u_j,v_j)_j$ be a
finite family of pairs of elements of $M_{n\times m}(\mathbf A_E)$ at $\infty$
with $\pi_w(x)Z_w^{\mathsf T}=\sum_j\sigma(v_j,x)\,\pi_w(u_j)$ for every $x$ at
$\infty$ and every $w$. Then for every $f\in\mathcal S(F_\infty^I)$ the family
$t\mapsto\omega_2(\gamma(t))$ acts weakly at $\infty$ to first order
(Definition~\ref{asm-theta-agree-real-fock}(7)) by
\[
f\mapsto(8\pi i)^{-1}\sum_j\bigl(D_{u_j}D_{v_j}f+D_{v_j}D_{u_j}f\bigr)+\Bigl(\sum_wc_w\operatorname{tr}Z_w\Bigr)f.
\]
\end{lemma}
\leanref{Development/WeilBound/ThetaLifts/Reductions/SymplOfMem.lean}{PositivityTheory.ThKR.main2}
\leanref{Development/WeilBound/CoreStatements/GaussianSeminormBounds/RealSeminormFiltration/RLLieE.lean}{PositivityTheory.RL.Up.Co.schur2}
\leanref{Development/WeilBound/CoreStatements/GaussianSeminormBounds/RealSeminormFiltration/RLLieE.lean}{PositivityTheory.RL.Up.exists_co}
\leanref{Development/WeilBound/CoreStatements/GaussianSeminormBounds/RealSeminormFiltration/RLLieG.lean}{PositivityTheory.RL.Up.exists_iota}
\begin{proof}
Write $V=F_\infty^I$ and $T(x,y)=\operatorname{Tr}_{F_\infty/\mathbf R}(\sum_ix_iy_i)$
for $x,y\in V$, and use the notation of
Lemma~\ref{asm-theta-agree-real-metaplectic} for $\mathbf R^l$, $l=\dim_{\mathbf R}V$.

\emph{Coordinates.} Let $\varepsilon_1\colon V\to\mathbf R^l$ be a linear
isomorphism, and $\varepsilon_2\colon V\to\mathbf R^l$ the linear isomorphism
with $\langle\varepsilon_2(y),\varepsilon_1(x)\rangle=-T(x,y)$, which exists
since $T$ is nondegenerate. For a pair $p=(p_1,p_2)$ of elements of
$\mathbf A^I$ at $\infty$ put
$\operatorname{vec}(p)=(\varepsilon_1(p_{1,\infty}),\varepsilon_2(p_{2,\infty}))\in\mathbf R^{2l}$;
$\operatorname{vec}$ is an $\mathbf R$-linear bijection from these pairs onto
$\mathbf R^{2l}$. Let $Pf=f\circ\varepsilon_1$, an isomorphism of
$\mathcal S(\mathbf R^l)$ onto $\mathcal S(V)$, and let $\varrho(p)$ be the
operator $\varrho(p)\phi(x)=\psi(x\cdot p_2+\tfrac12p_1\cdot p_2)\phi(x+p_1)$, so
that $\rho(u)=\varrho(e(u))$. Direct computation gives, for $p$ and $p'$ at
$\infty$, $f\in\mathcal S(\mathbf R^l)$ and Schwartz-Bruhat $h$ on
$\mathbf A_f^I$:
$\varrho(p)(Pf\otimes h)=P(\rho_S(\operatorname{vec}p)f)\otimes h$;
$P(d\rho(\operatorname{vec}p)f)=D_{p_{1,\infty},p_{2,\infty}}(Pf)$; and
$\operatorname{sympl}(\operatorname{vec}p,\operatorname{vec}p')=-\operatorname{Tr}_{F_\infty/\mathbf R}(p_1\cdot p'_2-p'_1\cdot p_2)_\infty$.
In particular
$\operatorname{sympl}(\operatorname{vec}e(u),\operatorname{vec}e(v))=\sigma(u,v)$
and $P\,d\rho(\operatorname{vec}e(u))=D_uP$ for $u,v$ at $\infty$.

\emph{The group.} Let $\mathfrak A$, $G$ and $\mathfrak g$ be as in
Lemma~\ref{asm-theta-agree-real-projective} for the matrices $K_w$. For
$z=(z_w)\in\prod_wM_{n\times m}(\mathbf C)$ let $L(z)$ be the element of
$M_{n\times m}(\mathbf A_E)$ at $\infty$ with $\pi_w(L(z))=z_w$ (fact (iii) in
the proof of Lemma~\ref{asm-theta-agree-real-lines}); then
$E_0(z)=\operatorname{vec}e(L(z))$ is an $\mathbf R$-linear isomorphism onto
$\mathbf R^{2l}$. Let $\rho_{\mathbf R}\colon\mathfrak A\to M_{2l}(\mathbf R)$
be the $\mathbf R$-algebra homomorphism with
$\rho_{\mathbf R}(a)E_0(z)=E_0((z_wa_w^{\mathsf T})_w)$. For $g\in G$ let
$\iota(g)\in GL_m(\mathbf A_E)$ be the element with $\iota(g)-1$ at $\infty$ and
$\pi_w(\iota(g))=g_w$; by the hypothesis on $K_w$ and facts (i) and (iii) in
the proof of Lemma~\ref{asm-theta-agree-real-lines}, $\iota(g)\in U_2$, and
$\iota$ is a continuous homomorphism. For $u$ at $\infty$,
$\rho_{\mathbf R}(g)\operatorname{vec}e(u)=\operatorname{vec}e(u\iota(g)^{\mathsf T})$,
and as $\langle u\iota(g)^{\mathsf T},v\iota(g)^{\mathsf T}\rangle=\langle u,v\rangle$
for $\iota(g)\in U_2$, $\rho_{\mathbf R}(g)$ preserves $\operatorname{sympl}$.

\emph{Schur.} Let $g\in G$ and let $U$ be an invertible operator implementing
$\rho_{\mathbf R}(g)$. Then there is $z\in\mathbf C$ with
$\omega_2(\iota(g))(Pf\otimes h)=P(zUf)\otimes h$ for all $f$ and $h$. Indeed,
put $\Theta=\omega_2(\iota(g))$ and $\Omega=U^{-1}$. (a) If $p$ has component
$0$ at $\infty$, then $\Theta\varrho(p)=\varrho(p)\Theta$, since $p=e(u)$ with
$u\iota(g)^{\mathsf T}=u$, as $\iota(g)-1$ is at $\infty$. (b) If $p=e(u)$ is at
$\infty$, then $\Theta(P(\Omega\rho_S(\operatorname{vec}p)f)\otimes h)=\varrho(p)\Theta(P(\Omega f)\otimes h)$,
by the intertwining properties of $U$ and $\omega_2$ and the formula for
$\rho_{\mathbf R}$. Let $\lambda_0$ be evaluation at $0$ composed with
$\Theta$, an admissible form. For fixed $f$, the form $h\mapsto\lambda_0(f\otimes h)$
is invariant under multiplication by the characters
$y\mapsto\psi_v(c\,y_{v,i})$ of the finite places, by (a); these are locally
constant and separate $0$ from every other point, so on the compact support of
an $h$ with $h(0)=0$ one of them is constant and different from $1$ near each
point, and the form vanishes on such $h$. For fixed $h$, the continuous form
$f\mapsto\lambda_0(P(\Omega f)\otimes h)$ on $\mathcal S(\mathbf R^l)$ is
invariant under multiplication by $x\mapsto e^{2\pi i\langle q,x\rangle}$, by
(b) with $p=(0,y)$, $y$ at $\infty$; differentiating in $q$, it vanishes on
$x_jf$ for every $j$, so it is supported at $0$ and is a multiple of evaluation
at $0$ \cite[Theorem~2.3.4]{Hor90}. Hence
$\lambda_0(P(\Omega f)\otimes h)=z\,f(0)h(0)$ for some $z$. For
$x\in\mathbf A^I$, translating by the finite part of $x$ with (a) and by
$x_\infty$ with (b) gives
$\Theta(P(\Omega f)\otimes h)(x)=z\,f(\varepsilon_1(x_\infty))h(x_f)$. So
$\Theta(P(\Omega f)\otimes h)=zPf\otimes h$, and the claim follows with $f$
replaced by $Uf$.

\emph{The homomorphism.} Let $H\subset G$ be the subgroup of the $g$ such that
$\rho_{\mathbf R}(g)$ is implemented by an invertible operator. By the previous
step, for $g\in H$ there is a unique operator $A(g)$ on
$\mathcal S(\mathbf R^l)$ with $\omega_2(\iota(g))(Pf\otimes h)=P(A(g)f)\otimes h$
for all $f$ and $h$ (uniqueness by fact (a) in the proof of
Lemma~\ref{asm-theta-agree-real-weak}), and $A(g)=zU$ for every invertible $U$
implementing $\rho_{\mathbf R}(g)$. By uniqueness, $A$ is a homomorphism. For a
continuous linear form $\lambda$ on $\mathcal S(\mathbf R^l)$ there are an
admissible form $\lambda'$ and $h_0$ with $\lambda'(Pf\otimes h_0)=\lambda(f)$
for all $f$, so $g\mapsto\lambda(A(g)f)=\lambda'(\omega_2(\iota(g))(Pf\otimes h_0))$
is continuous. Let $c_w$ be the constants of
Lemma~\ref{asm-theta-agree-real-projective} for $(K_w)$, $\rho_{\mathbf R}$, $H$
and $A$.

\emph{Conclusion.} Let $\gamma$, $(Z_w)$ and $(u_j,v_j)$ be as in the
statement, and put $\gamma_0(t)=(\exp(tZ_w))_w$. By fact (i) in the proof of
Lemma~\ref{asm-theta-agree-real-lines} and the hypothesis on $K_w$,
$\gamma_0(t)\in G$; differentiating, $Z=(Z_w)\in\mathfrak g$, and
$\gamma(t)=\iota(\gamma_0(t))$ since both are at $\infty$ with the same
components. For $x$ at $\infty$, the hypothesis on $(u_j,v_j)$ and the
formula for $\rho_{\mathbf R}$ give
\[
\rho_{\mathbf R}(Z)\operatorname{vec}e(x)=\operatorname{vec}e\Bigl(\sum_j\sigma(v_j,x)u_j\Bigr)
=\sum_j\operatorname{sympl}(\operatorname{vec}e(v_j),\operatorname{vec}e(x))\operatorname{vec}e(u_j),
\]
and $\operatorname{vec}e$ is onto $\mathbf R^{2l}$, so by
Lemma~\ref{asm-theta-agree-real-metaplectic}(1),
$P\operatorname{quad}(\rho_{\mathbf R}(Z))=(8\pi i)^{-1}\sum_j(D_{u_j}D_{v_j}+D_{v_j}D_{u_j})P$.
Let $f'\in\mathcal S(V)$, $f=P^{-1}f'$, $h$ Schwartz-Bruhat and $\lambda$
admissible. Then
$\lambda(\omega_2(\gamma(t))(f'\otimes h))=\lambda(P(A(\gamma_0(t))f)\otimes h)$,
and $f''\mapsto\lambda(Pf''\otimes h)$ is continuous, so by
Lemma~\ref{asm-theta-agree-real-projective} the derivative at $0$ is $\lambda$
of
$\bigl((8\pi i)^{-1}\sum_j(D_{u_j}D_{v_j}+D_{v_j}D_{u_j})f'+(\sum_wc_w\operatorname{tr}Z_w)f'\bigr)\otimes h$.
\end{proof}
% Lean: main2: exists_co (Co: e1, e2, pair), Co.vec, Co.P, schr_P, P_dRho, sympl_vec,
% sympl_e; exists_iota (with exists_lf, ext_at and map_form) for iota; exists_lf, Ee and Rm
% for rho_R (hrho, hact, hsp with hTens_actB and mem_sp_of); step and uniq with Co.schur2
% (eq_zero_of_mul_invariant for the finite places, Delta.eq_delta' for the archimedean
% place); A from A0; continuity from exists_adm and hom.cont; mid2; quad_eq and IsD.P.

\begin{lemma}\label{asm-theta-agree-real-lie}
Let $B$ be a theta base for a CM field $E$, let $X$ and $W$ be lattice spaces
of ranks $n$ and $d$ with matrices $J_X$ and $J_W$, put
$I=\{1,\ldots,n\}\times\{1,\ldots,d\}$, and let
$e\colon\mathbb W(X,W)\to\mathbf A^I\times\mathbf A^I$ be an
$\mathbf A$-linear isomorphism as in
Lemma~\ref{asm-theta-agree-real-fock-map}. Let $\chi$ be a homomorphism from
$\mathbf A_E^\times$ to the complex numbers of absolute value $1$ with
\emph{exponents} $m$: integers $m(w)$, one for each infinite place $w$ of $F$,
such that $\chi(z)=\prod_w(\pi_w(z)/|\pi_w(z)|)^{m(w)}$ for every
$z\in\mathbf A_E^\times$ whose finite components are $1$. Let $P,Q,R,S\geq0$,
$a_0,b_0\in\mathbf C$ with $a_0b_0=i$, $A_w$ frames of $a_0\tau_w(J_X)$ of type
$(P,Q)$ and $C_w$ frames of $b_0\tau_w(J_W)$ of type $(R,S)$, and let $\Phi$ be
a Fock map for $(e,(A_w),(C_w))$
(Lemma~\ref{asm-theta-agree-real-fock-map}). Let $T\in\mathcal T$.
\begin{enumerate}
\item Let $\omega_W$ be a representation of $U_W$ on functions on
  $\mathbf A^I$ whose operators $\omega_W(h)$ satisfy condition (a) of
  Definition~\ref{asm-theta-agree-models}(6) in place of the $\omega(1,h)$,
  for the isomorphism $e$ and the character $\chi_X=\chi$. If $\gamma$ is a
  line in $U_W$ generated by $(Y_w)$ relative to $(C_w)$
  (Definition~\ref{asm-theta-agree-real-fock}(6)), then
  $t\mapsto\omega_W(\gamma(t))$ acts weakly at $\infty$ to first order by
  $\Phi(T)\mapsto\Phi\bigl(\sum_w(d_2(Y_w)+\tfrac12m(w)\operatorname{tr}Y_w)^{(w)}T\bigr)$.
\item Let $\omega_X$ be a representation of $U_X$ on functions on
  $\mathbf A^I$ whose operators $\omega_X(g)$ satisfy condition (b) of
  Definition~\ref{asm-theta-agree-models}(6) in place of the $\omega(g,1)$,
  for the isomorphism $e$ and the character $\chi_W=\chi$. If $\gamma$ is a
  line in $U_X$ generated by $(Y_w)$ relative to $(A_w)$, then
  $t\mapsto\omega_X(\gamma(t))$ acts weakly at $\infty$ to first order by
  $\Phi(T)\mapsto\Phi\bigl(\sum_w(d_1(Y_w)+\tfrac12m(w)\operatorname{tr}Y_w)^{(w)}T\bigr)$.
\end{enumerate}
\end{lemma}
\leanref{Development/AbelianVarietyTheory/Proofs/DivisorProofs/IsFockMap.lean}{GlobalTheory.RealLift.LieSecond}
\leanref{Development/AbelianVarietyTheory/Proofs/DivisorProofs/IsFockMap.lean}{GlobalTheory.RealLift.LieFirst}
\leanref{Development/WeilBound/ThetaLifts/Reductions/LieScalar2.lean}{PositivityTheory.ThKR.lieSecond2}
\leanref{Development/WeilBound/CoreStatements/GaussianSeminormBounds/RealSeminormFiltration/RLLieH.lean}{PositivityTheory.RL.lieFirst}
\leanref{Development/WeilBound/CoreStatements/GaussianSeminormBounds/RealSeminormFiltration/TraceBlocks.lean}{PositivityTheory.RL.Up.fock_data}
\leanref{Development/WeilBound/CoreStatements/GaussianSeminormBounds/RealSeminormFiltration/RLLieH.lean}{PositivityTheory.RL.Up.fock_eq}
\begin{proof}
Write $\rho$, $D_u$ and $\sigma$ for the operators and the form of
Lemmas~\ref{asm-theta-agree-real-derivative}
and~\ref{asm-theta-agree-real-oscillator} attached to the isomorphism $e$,
and $\rho^\square$, $D^\square_u$ and $\sigma^\square$ for those attached
to the isomorphism $e^\square$ for $M_{n\times2d}(\mathbf A_E)$ of
condition (a) of Definition~\ref{asm-theta-agree-models}(6) for $\omega_W$,
with the matrices $J_X$ and
$J_W^\square=\operatorname{diag}(J_W,-J_W)$. For $u,v$ at $\infty$,
$D_uD_v=D_vD_u+2\pi i\,\sigma(u,v)$, by the symmetry of second derivatives and
the condition on $e$; $\sigma^\square((u,u'),(v,v'))=\sigma(u,v)-\sigma(u',v')$;
and $D^\square_{(u,0)}(f\boxtimes f')=(D_uf)\boxtimes f'$, where
$(f\boxtimes f')(x,x')=f(x)f'(x')$, since
$e^\square(u,0)=((e_1(u),0),(e_2(u),0))$.

(1) Write $\omega(1,h)$ for $\omega_W(h)$, $\chi_X$ for $\chi$ and $m_X$ for
$m$.

\emph{Data.} Let $C'_w$ be the inverse of $C_w$ and
$Z_w=C_wY_wC'_w$, so $\pi_w(\gamma(t))=\exp(tZ_w)$ and
$\operatorname{tr}Z_w=\operatorname{tr}Y_w$. As in the proof of
Lemma~\ref{asm-theta-agree-real-lines}(1), $Y_w\in\mathfrak u(R,S)$; so
$\operatorname{tr}Y_w$ is imaginary, and
$\pi_w(\det\gamma(t))=\exp(t\operatorname{tr}Y_w)$ has absolute value $1$. For
$w$, $l$ and $\epsilon\in\{0,1\}$ let $u_{w,l,\epsilon}$ and $v_{w,l,\epsilon}$
be the elements at $\infty$ with $\pi_{w'}$-components $0$ for $w'\neq w$ and
\[
\pi_w(u_{w,l,1})=A_wE_lY_w^{\mathsf T}C_w^{\mathsf T},\quad
\pi_w(u_{w,l,0})=A_w(iE_lY_w^{\mathsf T})C_w^{\mathsf T},\quad
\pi_w(v_{w,l,1})=A_w(\tfrac{i}2s_lE_l)C_w^{\mathsf T},\quad
\pi_w(v_{w,l,0})=A_w(-\tfrac12s_lE_l)C_w^{\mathsf T},
\]
with $E_l$ the matrix unit at $l$. Two computations in the coordinates of the
proof of Lemma~\ref{asm-theta-agree-real-fock-map} give:
\begin{enumerate}
\item[(a)] $\pi_w(x)Z_w^{\mathsf T}=\sum_j\sigma(v_j,x)\,\pi_w(u_j)$ for every
  $x$ at $\infty$, the sum over $j=(w',l,\epsilon)$. Indeed, if
  $\pi_w(x)=A_wcC_w^{\mathsf T}$, then $\pi_w(x)Z_w^{\mathsf T}=A_wcY_w^{\mathsf T}C_w^{\mathsf T}$,
  $\sigma(v_{w,l,1},x)=\operatorname{Re}c_l$ and
  $\sigma(v_{w,l,0},x)=\operatorname{Im}c_l$ by the form computed there, and
  $c=\sum_l(\operatorname{Re}c_l+i\operatorname{Im}c_l)E_l$.
\item[(b)] $(8\pi i)^{-1}\sum_j(h(c_{u_j})h(c_{v_j})+h(c_{v_j})h(c_{u_j}))=d_2(Y_w)$
  on $\mathcal P(P,Q;R,S)$, the sum over $j=(w,l,\epsilon)$ for fixed $w$, where
  $c_{u_j}$ and $c_{v_j}$ are the matrices between $A_w$ and $C_w^{\mathsf T}$
  displayed above. This follows from $L_{l'}R_l+R_lL_{l'}=2\omega_{l,l'}$ and
  the relations $\bar Y_{b'b}=-\epsilon(b)\epsilon(b')Y_{bb'}$ expressing
  $Y_w\in\mathfrak u(R,S)$.
\end{enumerate}
By (b), the Fock map property and Lemma~\ref{asm-theta-agree-real-weak}(3),
$(8\pi i)^{-1}\sum_j(D_{u_j}D_{v_j}+D_{v_j}D_{u_j})\Phi(T)=\Phi\bigl(\sum_wd_2(Y_w)^{(w)}T\bigr)$.

\emph{Two oscillator constants.} Put $K_w=\overline{b_0\tau_w(J_W)}$; as
$C_w$ is a frame of $b_0\tau_w(J_W)$, $K_w=C_w'^*J_{R,S}C'_w$ is Hermitian and
invertible, and $g^*K_wg=K_w$ if and only if
$g^{\mathsf T}\tau_w(J_W)\bar g=\tau_w(J_W)$. Let $c_w$ be the constants of
Lemma~\ref{asm-theta-agree-real-oscillator} for $J_X$, $J_W$, $e$,
$h\mapsto\omega(1,h)$ and $(K_w)$, which satisfy its hypotheses by condition
(a) of Definition~\ref{asm-theta-agree-models}(6). Let $c'_w$ be those for
$J_X$, $J_W^\square$,
$e^\square$, the operators $\omega^\square$ of that condition, and
$K^\square_w=\operatorname{diag}(K_w,-K_w)$. By (a),
Lemma~\ref{asm-theta-agree-real-oscillator} for $\omega$ shows that
$t\mapsto\omega(1,\gamma(t))$ acts weakly at $\infty$ to first order by
$\Phi(T)\mapsto\Phi(\sum_wd_2(Y_w)^{(w)}T)+(\sum_wc_w\operatorname{tr}Y_w)\Phi(T)$.
It remains to show $\sum_wc_w\operatorname{tr}Y_w=\sum_w\tfrac12m_X(w)\operatorname{tr}Y_w$.

\emph{First comparison:} $\sum_wc_w\operatorname{tr}Y_w=\sum_wc'_w\operatorname{tr}Y_w$.
The matrices $\delta(t)=\operatorname{diag}(\gamma(t),1)$ form a line in the
doubled group $G_W$, with $\pi_w(\delta(t))=\exp(t\operatorname{diag}(Z_w,0))$.
The pairs $((u_j,0),(v_j,0))$ satisfy the hypothesis of
Lemma~\ref{asm-theta-agree-real-oscillator} for
$\operatorname{diag}(Z_w,0)$, by (a) and the formula for $\sigma^\square$. Let
$f$ be a real smooth function with compact support on $F_\infty^I$ with
$f(0)=1$, let $g$ be a Schwartz-Bruhat function on $\mathbf A_f^I$ with
$g(y_0)=1$, and let $x_0\in\mathbf A^I$ have components $0$ at $\infty$ and
$y_0$ at the finite places. By the compatibility
$\omega^\square(\operatorname{diag}(h,1))(\phi_1\boxtimes\phi_2)=(\omega(1,h)\phi_1)\boxtimes\phi_2$,
\[
\omega^\square(\delta(t))\bigl((f\otimes g)\boxtimes(f\otimes g)\bigr)(x_0,x_0)=\omega(1,\gamma(t))(f\otimes g)(x_0)\cdot(f\otimes g)(x_0),
\]
and $(f\otimes g)(x_0)=f(0)g(y_0)=1$. Differentiate at $t=0$, using
Lemma~\ref{asm-theta-agree-real-oscillator} for $\omega^\square$ on the left
and for $\omega$ on the right, with evaluation at a point as the admissible
form. Since $D^\square_{(u,0)}$ acts on the first factor of $f\boxtimes f$, the
quadratic terms agree, and what remains is
$(\sum_wc'_w\operatorname{tr}Z_w)f(0)^2=(\sum_wc_w\operatorname{tr}Z_w)f(0)$.

\emph{Second comparison:}
$2\sum_wc'_w\operatorname{tr}Y_w=\sum_wm_X(w)\operatorname{tr}Y_w$. Now let
$\delta(t)=\operatorname{diag}(\gamma(t),\gamma(t))$, a line in $G_W$ with
$\pi_w(\delta(t))=\exp(t\operatorname{diag}(Z_w,Z_w))$ and
$\delta(t)\Delta=\Delta\gamma(t)$ for
$\Delta=\left(\begin{smallmatrix}1\\1\end{smallmatrix}\right)$. By (a) and the
formula for $\sigma^\square$, the pairs
$((\tfrac12u_j,\tfrac12u_j),(v_j,-v_j))$ and
$((\tfrac12u_j,-\tfrac12u_j),(v_j,v_j))$ satisfy the hypothesis of
Lemma~\ref{asm-theta-agree-real-oscillator} for $\operatorname{diag}(Z_w,Z_w)$.
Let $\lambda$ be a linear form on the functions on
$\mathbf A^I\times\mathbf A^I$ that agrees with
$\phi\mapsto\int_{\mathbf A^I}\phi(x,x)\,dx$, for a Haar measure, on the
functions integrable on the diagonal. It is admissible; it satisfies
$\lambda(\rho^\square(u)\phi)=\lambda(\phi)$ for Schwartz-Bruhat $\phi$
whenever $e^\square(u)=((x,x),(y,-y))$ for some $x,y\in\mathbf A^I$, as the
operator translates along the diagonal and multiplies by a character that is
$1$ on it; and $\lambda(\Psi)\neq0$ for
$\Psi=(f\otimes|g|)\boxtimes(f\otimes|g|)$, with $f$ and $g$ as before. For $y$
at $\infty$, $e^\square(y,y)$ has this form, so differentiating
$t\mapsto\lambda(\rho^\square(t(y,y))\phi)$ gives
$\lambda(D^\square_{(y,y)}\phi'\otimes G)=0$ for Schwartz $\phi'$ and
Schwartz-Bruhat $G$, by Lemma~\ref{asm-theta-agree-real-derivative}. On one
side, $e^\square(z\Delta^{\mathsf T})=e^\square(z,z)$ also has this form, so
the displayed rule of Definition~\ref{asm-theta-agree-doubling}(6)(b) for
$\omega^\square$, with character $\chi_X$ and $m=n$, gives
$\lambda(\omega^\square(\delta(t))\Psi)=\chi_X(\det\gamma(t))|\det\gamma(t)|^{n/2}\lambda(\Psi)$.
Here $\det\gamma(t)-1$ is at $\infty$ and its $\pi_w$-components have
absolute value $1$, so $|\det\gamma(t)|=1$ and
$\chi_X(\det\gamma(t))=\exp\bigl(t\sum_wm_X(w)\operatorname{tr}Y_w\bigr)$. The
derivative at $0$ is $(\sum_wm_X(w)\operatorname{tr}Y_w)\lambda(\Psi)$. On the
other side, by Lemma~\ref{asm-theta-agree-real-oscillator} for $\omega^\square$
the derivative is $\lambda$ applied to the quadratic terms plus
$(\sum_wc'_w\operatorname{tr}\operatorname{diag}(Z_w,Z_w))\lambda(\Psi)$. In
each quadratic term, one of the two operators has the form
$D^\square_{(y,y)}$; moving it to the left with the commutation relation, the
two terms of the pair $j$ contribute
$2\pi i\,\sigma^\square((v_j,-v_j),(\tfrac12u_j,\tfrac12u_j))\lambda(\Psi)$ and
$2\pi i\,\sigma^\square((\tfrac12u_j,-\tfrac12u_j),(v_j,v_j))\lambda(\Psi)$,
which are $4\pi i\,\sigma(v_j,\tfrac12u_j)\lambda(\Psi)$ and its negative. So
the derivative is $2(\sum_wc'_w\operatorname{tr}Y_w)\lambda(\Psi)$, and the
comparison follows as $\lambda(\Psi)\neq0$.

So $\sum_wc_w\operatorname{tr}Y_w=\sum_w\tfrac12m_X(w)\operatorname{tr}Y_w$,
and $t\mapsto\omega(1,\gamma(t))$ acts weakly at $\infty$ to first order by
$\Phi(T)\mapsto\Phi(\sum_w(d_2(Y_w)+\tfrac12m_X(w)\operatorname{tr}Y_w)^{(w)}T)$.

(2) Let $\varsigma$ be the isomorphism from $\mathcal P(R,S;P,Q)$ to
$\mathcal P(P,Q;R,S)$ sending $x_{(b,a)}$ to $x_{(a,b)}$. It carries $L_{(b,a)}$
and $R_{(b,a)}$ to $L_{(a,b)}$ and $R_{(a,b)}$, hence $h(c)$ to $h(c^{\mathsf T})$
and $d_2(Y)$, for $Y\in M_{P+Q}(\mathbf C)$, to $d_1(Y)$. Let $\Psi$ be the
isomorphism $\bigotimes_w\varsigma$. Let $e^{\mathsf T}$ be the isomorphism
$w\mapsto e(w^{\mathsf T})$ from $\mathbb W(W,X)$ to
$\mathbf A^I\times\mathbf A^I$, with $I$ read as
$\{1,\ldots,d\}\times\{1,\ldots,n\}$ through $(i,j)\mapsto(j,i)$; it is an
isomorphism as in Lemma~\ref{asm-theta-agree-real-fock-map} for $(W,X)$.
Condition (b) for $\omega_X$, $e$ and $\chi$ is condition (a) for $\omega_X$,
$e^{\mathsf T}$ and $\chi$, with $W$ and $X$ exchanged. If $u$ is at $\infty$
with $\pi_w(u)=C_wc_wA_w^{\mathsf T}$, then
$\pi_w(u^{\mathsf T})=A_wc_w^{\mathsf T}C_w^{\mathsf T}$, so $\Phi\circ\Psi$ is
a Fock map for $(e^{\mathsf T},(C_w),(A_w))$, with $b_0a_0=i$. Part (1) for
$(W,X)$, $e^{\mathsf T}$, $\Phi\circ\Psi$ and $\omega_X$, applied to
$\Psi^{-1}T$, gives the claim, since $\Psi$ carries
$\sum_w(d_2(Y_w)+\tfrac12m(w)\operatorname{tr}Y_w)^{(w)}$ to
$\sum_w(d_1(Y_w)+\tfrac12m(w)\operatorname{tr}Y_w)^{(w)}$.
\end{proof}
% Lean: (1) is lieSecond2: main2 for (K, om) and for the doubled (dd, omd); fock_data (Z, u, v
% from cu and cv, expand, fock_eq); step2 with exists_dbl (gamma, 1), exists_box, exists_SB,
% box_D and the restriction hres; step1 with exists_dbl (gamma, gamma), exists_form, the Siegel
% rule hS, at_det, norm_eq_one, IsD.comm and sB_dbl. (2) is lieFirst through sw, sw_heis and
% sw_fock.

\begin{lemma}\label{asm-theta-agree-real}
Let $(B,G,\mathcal X,\omega,\omega')$ be a triple for a theta base $B$ of a CM
field $E$ (Definition~\ref{asm-theta-agree-triple}). Then there are real maps
matched to it (Definition~\ref{asm-theta-agree-matched}(2)).
\end{lemma}
\leanref{Development/AbelianVarietyTheory/Definitions/Divisors/MatchingOfParts.lean}{GlobalTheory.ThK.KReal}
\leanref{Development/WeilBound/ThetaLifts/Reductions/KReal.lean}{PositivityTheory.ThKR.kReal'}
\leanref{Development/WeilBound/CoreStatements/GaussianSeminormBounds/RealSeminormFiltration/KRealAsm.lean}{PositivityTheory.KRealA.kReal_of_lie}
\begin{proof}
Let $(a,a')$ be the ranks of the triple and $p$ and $q$ the integers of $G$,
and put $(P,Q)=(a-p,p)$, $(P',Q')=(q,a'-q)$ and $(R,S)=(a+a'-p-q,p+q)$; these
are the signatures of $V$, $V'$ and $W$. For an infinite place $w$ put
$\kappa_w=m_\alpha(w)-m_\beta(w)$.

\emph{Local data.} Fix $w$. By Lemma~\ref{asm-theta-agree-real-parity}
$\kappa_w$ is even, and as $G$ is proper, $a+a'$ is even, $a+a'<\kappa_w$,
$2(p+q)+1<\lfloor(a-a')/2\rfloor$ and $2(p+q)+1<a'$. By
Lemma~\ref{asm-theta-agree-real-one-place} there are an irreducible
$(\mathfrak g,K)$-module $N_w$ of type $(R,S)$ with scalar endomorphisms, a lift
$\pi_w\colon\mathcal P(P,Q;R,S)\to N_w$ of type $(\kappa_w/2,0)$ and a lift
$\pi'_w\colon\mathcal P(P',Q';R,S)\to N_w$ of type $(-\kappa_w/2,\kappa_w/2)$.
By Lemma~\ref{asm-theta-agree-real-frames}, applied to $\tau_w(J_V)$,
$\tau_w(J_{V'})$ and $i\tau_w(J_W)$ with their Cartan involutions, there are
frames $A_w$ of $\tau_w(J_V)$, $A'_w$ of $\tau_w(J_{V'})$ and $C_w$ of
$i\tau_w(J_W)$, of types $(P,Q)$, $(P',Q')$ and $(R,S)$, with
$\theta_{V,w}A_w=A_wJ_{P,Q}$, $\theta_{V',w}A'_w=A'_wJ_{P',Q'}$ and
$\theta_{W,w}C_w=C_wJ_{R,S}$.

\emph{Fock maps.} Let $\mathcal T=\bigotimes_w\mathcal P(P,Q;R,S)$ and
$\mathcal T'=\bigotimes_w\mathcal P(P',Q';R,S)$. By
Lemma~\ref{asm-theta-agree-real-fock-map} there are Fock maps $\Phi$ for
$(M,(A_w),(C_w))$ and $\Phi'$ for $(M',(A'_w),(C_w))$, with $a_0=1$ and
$b_0=i$, whose images are $\operatorname{Pol}(M)$ and $\operatorname{Pol}(M')$.
They are injective; let $\Psi$ and $\Psi'$ be linear maps with
$\Psi\Phi=1$ and $\Psi'\Phi'=1$.

\emph{Derivatives.} The characters $\beta$, $\alpha$, $\beta/\alpha$ and
$\alpha/\beta$ have exponents $m_\beta$, $m_\alpha$, $-\kappa$ and $\kappa$, as
$\mathcal X$ is proper. Lemma~\ref{asm-theta-agree-real-lie}, applied to
$h\mapsto\omega(1,h)$ and $g\mapsto\omega(g,1)$ for $\omega$ with characters
$(\beta,\beta/\alpha)$, and likewise for $\omega'$ with characters
$(\alpha,\alpha/\beta)$, gives the following, for $T\in\mathcal T$ and
$T'\in\mathcal T'$:
\begin{enumerate}
\item[(i)] if $\gamma$ is a line in $U_V$ generated by $(Y_w)$ relative to
  $(A_w)$, then $t\mapsto\omega(\gamma(t),1)$ acts weakly at $\infty$ to first
  order by $\Phi(T)\mapsto\Phi(\sum_w(d_1(Y_w)-\tfrac12\kappa_w\operatorname{tr}Y_w)^{(w)}T)$,
  and likewise for $V'$, $(A'_w)$, $\omega'$ and $\Phi'$ with
  $+\tfrac12\kappa_w$;
\item[(ii)] if $\gamma$ is a line in $U_W$ generated by $(Y_w)$ relative to
  $(C_w)$, then $t\mapsto\omega(1,\gamma(t))$ acts weakly at $\infty$ to first
  order by $\Phi(T)\mapsto\Phi(L_YT)$ and $t\mapsto\omega'(1,\gamma(t))$ by
  $\Phi'(T')\mapsto\Phi'(L'_YT')$, where
  $L_Y=\sum_w(d_2(Y_w)+\tfrac12m_\beta(w)\operatorname{tr}Y_w)^{(w)}$ and
  $L'_Y=\sum_w(d_2(Y_w)+\tfrac12m_\alpha(w)\operatorname{tr}Y_w)^{(w)}$.
\end{enumerate}
In each case the operators send $0$ to $0$, so by
Lemma~\ref{asm-theta-agree-real-weak} the family also acts at $\infty$ to first
order by the same map, and by no other map on these vectors.

\emph{The real maps.} Let $\Pi=\bigotimes_wN_w$, let
$\pi=\bigotimes_w\pi_w\colon\mathcal T\to\Pi$ and
$\pi'=\bigotimes_w\pi'_w\colon\mathcal T'\to\Pi$, which are surjective, and put
$\varpi=\pi\Psi$ and $\varpi'=\pi'\Psi'$; so $\varpi(\Phi(T))=\pi(T)$ and
$\varpi'(\Phi'(T'))=\pi'(T')$. For a family $(Y_w)$ with
$Y_w\in M_{R+S}(\mathbf C)$ put
$S_Y=\sum_w(N_w(Y_w)+\tfrac12m_\beta(w)\operatorname{tr}Y_w)^{(w)}$ on $\Pi$.
Since $\pi_w$ and $\pi'_w$ are lifts of types $(\kappa_w/2,0)$ and
$(-\kappa_w/2,\kappa_w/2)$, and $\tfrac12m_\alpha(w)=\tfrac12\kappa_w+\tfrac12m_\beta(w)$,
\[
\pi\circ\sum_w(d_1(X_w)-\tfrac12\kappa_w\operatorname{tr}X_w)^{(w)}=0,\quad
\pi'\circ\sum_w(d_1(X_w)+\tfrac12\kappa_w\operatorname{tr}X_w)^{(w)}=0,\quad
\pi\circ L_Y=S_Y\circ\pi,\quad
\pi'\circ L'_Y=S_Y\circ\pi'.
\]
We show that $(\Pi,\varpi,\varpi')$ are matched to the triple
(Definition~\ref{asm-theta-agree-matched}(2)).

(a) First, $\varpi(\operatorname{Pol}(M))=\pi(\mathcal T)=\Pi$. Let
$f=\Phi(T)\in\operatorname{Pol}(M)$, let $\gamma$ be a line in $U_V$, and let
$t\mapsto\omega(\gamma(t),1)$ act at $\infty$ to first order by
$f\mapsto f_1$. By Lemma~\ref{asm-theta-agree-real-lines}(1) $\gamma$ is
generated by some $(Y_w)$ relative to $(A_w)$, and by (i)
$f_1=\Phi(\sum_w(d_1(Y_w)-\tfrac12\kappa_w\operatorname{tr}Y_w)^{(w)}T)$; so
$\varpi(f_1)=0$ by the first relation. Next let $k\in K_\infty(V)$, and let
$\omega(k,1)$ act at $\infty$ by $f\mapsto f_1$. By
Lemma~\ref{asm-theta-agree-real-lines}(3), $k=\gamma(1)$ for a line $\gamma$
generated by block-diagonal matrices $Y_w$. The operator
$L=\sum_w(d_1(Y_w)-\tfrac12\kappa_w\operatorname{tr}Y_w)^{(w)}$ is locally
finite, by Lemma~\ref{asm-theta-agree-real-fock-relations}(2) and
Lemma~\ref{asm-theta-agree-real-exponential}(1); let $U$ be its exponential.
The operators $\omega(\gamma(t),1)$ form a one-parameter group, they preserve
the admissible forms by condition (b) of a Weil representation, and by (i)
they act weakly at $\infty$ to first order by $\Phi(T')\mapsto\Phi(LT')$ for
every $T'$. By Lemma~\ref{asm-theta-agree-real-exponential}(3), $\omega(k,1)$
acts weakly at $\infty$ by $\Phi(T)\mapsto\Phi(UT)$, so $f_1=\Phi(UT)$ by
Lemma~\ref{asm-theta-agree-real-weak}. As $\pi L=0\circ\pi$ and the exponential
of $0$ is the identity, Lemma~\ref{asm-theta-agree-real-exponential}(2) gives
$\pi(UT)=\pi(T)$, that is, $\varpi(f_1)=\varpi(f)$. The same argument applies
to $\varpi'$, with $M'$, $V'$, $\omega'$, $(A'_w)$ and the second relation.

(b) Let $f=\Phi(T)\in\operatorname{Pol}(M)$ and
$f'=\Phi'(T')\in\operatorname{Pol}(M')$ with $\varpi(f)=\varpi'(f')$, that is,
$\pi(T)=\pi'(T')$. If $\gamma$ is a line in $U_W$ and $f\mapsto f_1$ and
$f'\mapsto f'_1$ are the actions at $\infty$ to first order of
$t\mapsto\omega(1,\gamma(t))$ and $t\mapsto\omega'(1,\gamma(t))$, let $(Y_w)$
generate $\gamma$ relative to $(C_w)$
(Lemma~\ref{asm-theta-agree-real-lines}(1)). By (ii), $f_1=\Phi(L_YT)$ and
$f'_1=\Phi'(L'_YT')$, so
$\varpi(f_1)=S_Y\pi(T)=S_Y\pi'(T')=\varpi'(f'_1)$. If $k\in K_\infty(W)$ and
$\omega(1,k)$ and $\omega'(1,k)$ act at $\infty$ by $f\mapsto f_1$ and
$f'\mapsto f'_1$, let $k=\gamma(1)$ with $\gamma$ generated by block-diagonal
matrices $Y_w$ (Lemma~\ref{asm-theta-agree-real-lines}(3)). As in
(a), $L_Y$ and $L'_Y$ are locally finite with exponentials $U$ and $U'$, and
$f_1=\Phi(UT)$ and $f'_1=\Phi'(U'T')$. By the last two relations and
Lemma~\ref{asm-theta-agree-real-exponential}(2), with $L_0=S_Y$,
$\pi(UT)=\pi'(U'T')$, that is, $\varpi(f_1)=\varpi'(f'_1)$.

(c) By Lemma~\ref{asm-theta-agree-real-tensor}, $\Pi\neq0$. Let $N\subset\Pi$ be
a subspace with $\varpi(f_1)\in N$ whenever $f\in\operatorname{Pol}(M)$,
$\varpi(f)\in N$ and $f\mapsto f_1$ is one of the actions of $\omega$ in (b).
Let $w$, $Z\in\mathfrak u(R,S)$ and $y\in N$; write $y=\pi(T)$. Let
$Y_w=Z$ and $Y_{w'}=0$ for $w'\neq w$. By
Lemma~\ref{asm-theta-agree-real-lines}(2), $(Y_{w'})$ generates a line
$\gamma$ in $U_W$, and by (ii) $t\mapsto\omega(1,\gamma(t))$ acts at $\infty$
to first order by $\Phi(T)\mapsto\Phi(L_YT)$. As $\varpi(\Phi(T))=y\in N$,
\[
N\ni\varpi(\Phi(L_YT))=S_Y\pi(T)=N_w(Z)^{(w)}y+\tfrac12m_\beta(w)\operatorname{tr}(Z)\,y,
\]
so $N_w(Z)^{(w)}y\in N$. By Lemma~\ref{asm-theta-agree-real-tensor},
$N=0$ or $N=\Pi$.
\end{proof}
% Lean: KReal is discharged by kReal' localLift_holds, which is kReal_of_lie (lieSecond2
% localLift_holds derEx) (lieFirst ...). kReal_of_lie: onePlace with i_parity; i_frameex for V,
% V', W; i_fockex and i_polrange; hsec and hfst give hW2, hW2', hW1, hW1'; Psi from
% exists_leftInverse_of_injective; Rm = (tensor of the N w, map p o Psi, map p' o Psi');
% side_quot (genEx, actsLaws, isLift_fock1; kLine, locFin_fock1, locFin_sumPl, expAct_fock,
% exp_equiv, exp_id) for (a); side_wlie and side_wcpt with exp_equiv for (b); tensor_irr with
% lineEx for (c).

\subsection{Finite maps matched to a triple}\label{subsection-asm-finite-maps}

\begin{lemma}\label{asm-theta-agree-finite-unique}
Let $F$ be a number field, $I$ a finite set, $v$ a finite place of $F$ and $T$
an operator on functions on $\mathbf A_F^I$. Let $\phi$ be a Schwartz-Bruhat
function on $F_v^I$, and suppose that $T$ acts at $v$
(Definition~\ref{asm-theta-agree-models}(2)) by $\phi\mapsto\phi_1$ and by
$\phi\mapsto\phi_2$. Then $\phi_1=\phi_2$.
\end{lemma}
\leanref{Development/WeilBound/CoreStatements/LocallyConstantCoinvariants/RightOnto.lean}{PositivityTheory.KCo.ActsUniq}
\leanref{Development/WeilBound/CoreStatements/LocallyConstantCoinvariants/PfActs.lean}{PositivityTheory.KCo.action_uniqueness}
\begin{proof}
Choose $f\in\mathcal S(F_\infty^I)$ with $f(0)=1$, and let $r(x)$ be $f(x_\infty)$
times the product over the finite places $u\neq v$ of the indicator function of
$\mathcal O_u^I$ at $x_u$. Then $r$ does not depend on the $v$-component and
$\phi\otimes r$ is a Schwartz-Bruhat function, so
$\phi_1\otimes r=T(\phi\otimes r)=\phi_2\otimes r$. For $y\in F_v^I$ let
$s\in\mathbf A_F^I$ have component $y$ at $v$ and $0$ elsewhere. Then $r(s)=1$,
and evaluating at $s$ gives $\phi_1(y)=\phi_2(y)$.
\end{proof}

\begin{definition}\label{asm-theta-agree-finite-split-models}
Let $k$ be a nonarchimedean local field with ring of integers $\mathcal O$,
uniformizer $\varpi$, residue field of $q$ elements and normalized absolute
value $|\cdot|$, and let $a,a'\geq0$ with $d=a+a'$. A character $\chi$ of
$k^\times$ is \emph{unitary} if $|\chi(x)|=1$ for all $x$, \emph{smooth} if it is
trivial on $1+\varpi^e\mathcal O$ for some $e\geq1$, and \emph{unramified} if it
is trivial on $\mathcal O^\times$; otherwise it is \emph{ramified}. Let $\nu$,
$\alpha$ and $\beta$ be characters of $k^\times$.
\begin{enumerate}
\item For $g\in GL_a(k)$, $g'\in GL_{a'}(k)$, $h\in GL_d(k)$ and functions
  $\Xi$ on $M_{a\times d}(k)$ and $\Psi$ on $M_{a'\times d}(k)$, put
  \begin{align*}
  \omega_A^{\nu,\beta}(g,h)\Xi(X)&=\nu(\det g)^{-1}\beta(\det h)\,|\det g|^{-d/2}|\det h|^{-a/2}\,\Xi\bigl(g^{-1}X(h^{-1})^{\mathsf T}\bigr),\\
  \omega_B^{\nu,\alpha}(g',h)\Psi(Y)&=\nu(\det g')\,\alpha(\det h)\,|\det g'|^{d/2}|\det h|^{a'/2}\,\Psi(g'^{\mathsf T}Yh).
  \end{align*}
  For unitary smooth $\nu$ and $\alpha$, $\omega_B^{\nu,\alpha}$ is the
  representation $\omega_{\nu,\alpha}$ of
  Definition~\ref{asm-theta-agree-coinvariants-local}(1).
\item A matrix $b\in GL_d(k)$ is \emph{block upper triangular} if
  $b=\left(\begin{smallmatrix}b_A&*\\0&b_D\end{smallmatrix}\right)$ with $b_A$
  of size $a\times a$. The space $I(\alpha,\beta)$ consists of the functions $F$
  on $GL_d(k)$ that are invariant under right translation by some open subgroup
  of $GL_d(k)$ and satisfy
  \[
  F(bh)=\alpha(\det b_A)\,\beta(\det b_D)\,|\det b_A|^{a'/2}|\det b_D|^{-a/2}\,F(h)
  \]
  for every block upper triangular $b$ and every $h\in GL_d(k)$. For
  $h_1\in GL_d(k)$, $R(h_1)$ is right translation, $R(h_1)F(h)=F(hh_1)$; it
  preserves $I(\alpha,\beta)$.
\item Let $F$ be a number field with $k=F_v$ for a finite place $v$, let $I$ be a
  finite set and $n\geq0$, let $\iota$ be a linear map from functions on $F_v^I$
  to functions on $M_{n\times d}(F_v)$, $T$ an operator on functions on
  $\mathbf A_F^I$ and $S$ an operator on functions on $M_{n\times d}(F_v)$. Then
  $T$ \emph{reads as $S$ through $\iota$} if for every Schwartz-Bruhat function
  $\phi$ on $F_v^I$ there is a Schwartz-Bruhat function $\phi_1$ such that $T$
  acts at $v$ by $\phi\mapsto\phi_1$ (Definition~\ref{asm-theta-agree-models}(2))
  and $\iota\phi_1=S(\iota\phi)$.
\end{enumerate}
\end{definition}
\leanref{Development/AbelianVarietyTheory/Definitions/Divisors/RankEqCardIffExistsFrameMatrixFirstColumnBlock.lean}{GlobalTheory.SplitFlip.omA}
\leanref{Development/AbelianVarietyTheory/Definitions/Divisors/RankEqCardIffExistsFrameMatrixFirstColumnBlock.lean}{GlobalTheory.SplitFlip.omB}
\leanref{Development/AbelianVarietyTheory/Definitions/Divisors/RankEqCardIffExistsFrameMatrixFirstColumnBlock.lean}{GlobalTheory.SplitFlip.CondLE}
\leanref{Development/AbelianVarietyTheory/Definitions/Foundations/MatrixMeasurableSpace.lean}{GlobalTheory.SplitWeil.IsUnram}
\leanref{Development/WeilBound/CoreStatements/CompactAveragingSplitting/EigenvectorLifting/SplitLoc.lean}{PositivityTheory.ThK.SplitLoc.IsBlockUpper}
\leanref{Development/WeilBound/CoreStatements/CompactAveragingSplitting/EigenvectorLifting/SplitLoc.lean}{PositivityTheory.ThK.SplitLoc.indSp}
\leanref{Development/WeilBound/CoreStatements/CompactAveragingSplitting/EigenvectorLifting/SplitLoc.lean}{PositivityTheory.ThK.SplitLoc.rtr}
\leanref{Development/WeilBound/CoreStatements/LocallyConstantCoinvariants/RightOnto.lean}{PositivityTheory.KCo.Reads}

\begin{theorem}[Bernstein \cite{Ber84}]\label{asm-theta-agree-finite-bernstein}
Let $k$ be a nonarchimedean local field, let $a,a'\geq0$, and let $\alpha$ and
$\beta$ be unitary smooth characters of $k^\times$ with $a\neq a'$ or
$\alpha\neq\beta$. Then $0$ and $I(\alpha,\beta)$ are the only subspaces of
$I(\alpha,\beta)$ (Definition~\ref{asm-theta-agree-finite-split-models}) that
are stable under every $R(h_1)$, $h_1\in GL_{a+a'}(k)$.
\end{theorem}
\leanref{Development/AutomorphicSpectrum/Proofs/BlockDecomposition/FinSplitOfNarrow.lean}{GlobalTheory.CompanionTheory.BernBridge.Split.NarrowBernstein}
\leanref{Development/WeilBound/CoreStatements/TemperedRepresentationTheory/BernsteinNarrowness/RamifiedMultiplicityOneOfGeometricFiltration.lean}{PositivityTheory.BernsteinRamification.narrowBern}
\leanref{Development/AutomorphicSpectrum/Proofs/BlockDecomposition/FinSplitOfNarrow.lean}{GlobalTheory.CompanionTheory.BernBridge.Split.indIrredLt_of}
\leanref{Development/AutomorphicSpectrum/Proofs/BlockDecomposition/IsSmoothCharOfCondLE.lean}{GlobalTheory.CompanionTheory.SplitNarrow.indIrredNarrow_of}

\begin{lemma}\label{asm-theta-agree-finite-iwasawa}
Let $k$ be a nonarchimedean local field with ring of integers $\mathcal O$, let
$a,a'\geq0$ with $d=a+a'$, and let $\alpha$ and $\beta$ be characters of
$k^\times$. If $F\in I(\alpha,\beta)$
(Definition~\ref{asm-theta-agree-finite-split-models}) vanishes on
$GL_d(\mathcal O)$, then $F=0$.
\end{lemma}
\leanref{Development/WeilBound/CoreStatements/CompactAveragingSplitting/EigenvectorLifting/SplitDict.lean}{PositivityTheory.ThK.SplitDict.IndIwasawa}
\leanref{Development/WeilBound/CoreStatements/CompactAveragingSplitting/EigenvectorLifting/InducedVanishOfVanishOnIntegral.lean}{PositivityTheory.ThK.SplitDict.induced_vanish_of_vanish_on_integral}
\begin{proof}
By the Iwasawa decomposition, every $h\in GL_d(k)$ is $bk$ with $b$ block upper
triangular and $k\in GL_d(\mathcal O)$. The law of $I(\alpha,\beta)$ gives
$F(h)=c\,F(k)$ for a number $c$, and $F(k)=0$.
\end{proof}

\begin{lemma}\label{asm-theta-agree-finite-zeta}
Let $k$ be a nonarchimedean local field with ring of integers $\mathcal O$,
uniformizer $\varpi$, residue field of $q$ elements and normalized absolute value
$|\cdot|$, let $a>a'\geq0$ with $d=a+a'$ even, let $\alpha$ and $\beta$ be
unitary smooth characters of $k^\times$, and put $\nu=\alpha/\beta$,
$\omega_A=\omega_A^{\nu,\beta}$ and $\omega_B=\omega_B^{\nu,\alpha}$
(Definition~\ref{asm-theta-agree-finite-split-models}). Let $Z$ and $Z'$ be the
maps of Lemma~\ref{th-bottom-split-zeta}, on functions $\Xi$ on
$M_{a\times d}(k)$ and $\Psi$ on $M_{a'\times d}(k)$, with $X_0$, $Y_0$, $K$ and
$\varpi^n$ as there.
\begin{enumerate}
\item Suppose that $\nu$ is ramified. Then $Z$ and $Z'$ are linear on the
  Schwartz-Bruhat functions and map them into $I(\alpha,\beta)$; for
  Schwartz-Bruhat $\Xi$ and $\Psi$ and all $g\in GL_a(k)$, $g'\in GL_{a'}(k)$
  and $h\in GL_d(k)$,
  \[
  Z(\omega_A(g,1)\Xi)=Z(\Xi),\quad Z(\omega_A(1,h)\Xi)=R(h)Z(\Xi),\quad
  Z'(\omega_B(g',1)\Psi)=Z'(\Psi),\quad Z'(\omega_B(1,h)\Psi)=R(h)Z'(\Psi);
  \]
  and $Z(\Xi)\neq0$ and $Z'(\Psi)\neq0$ for some Schwartz-Bruhat $\Xi$ and
  $\Psi$.
\item Suppose that $\alpha$ and $\beta$ are unramified. For $r\geq0$ let
  $P_r(X)=\prod_{i=0}^{r-1}\bigl(1-\nu(\varpi)q^{i-d/2}X\bigr)=\sum_jc_{r,j}X^j$.
  For a function $\phi$ on $M_a(k)$ and $m\in\mathbf Z$ let $z_m(\phi)$ be the sum,
  over the nonincreasing $n\in\mathbf Z^a$ with $n_1+\cdots+n_a=m$, of
  $\int_{K\varpi^nK}\phi(g)\,\nu(\det g)|\det g|^{d/2}\,dg$, and put
  $\zeta(\phi)=\sum_{m\in\mathbf Z}\sum_{j=0}^ac_{a,j}z_{m-j}(\phi)$, each sum being
  taken to be $0$ if infinitely many of its terms are nonzero. Put
  \[
  Z^{\mathrm u}(\Xi)(h)=\beta(\det h)|\det h|^{-a/2}\,\zeta\bigl(g\mapsto\Xi(gX_0(h^{-1})^{\mathsf T})\bigr),
  \qquad Z'^{\mathrm u}(\Psi)=P_{a'}(1)\,Z'(\Psi).
  \]
  Then $Z^{\mathrm u}$ and $Z'^{\mathrm u}$ have the properties of $Z$ and $Z'$
  in (1), and
  $Z^{\mathrm u}(\mathbf 1_{M_{a\times d}(\mathcal O)})(h)=Z'^{\mathrm u}(\mathbf 1_{M_{a'\times d}(\mathcal O)})(h)=1$
  for every $h\in GL_d(\mathcal O)$.
\end{enumerate}
\end{lemma}
\leanref{Development/WeilBound/CoreStatements/CompactAveragingSplitting/EigenvectorLifting/SplitLoc.lean}{PositivityTheory.ThK.SplitLoc.HidaSetting}
\leanref{Development/WeilBound/CoreStatements/CompactAveragingSplitting/EigenvectorLifting/SplitLoc.lean}{PositivityTheory.ThK.SplitLoc.pRam}
\leanref{Development/WeilBound/CoreStatements/CompactAveragingSplitting/EigenvectorLifting/SplitLoc.lean}{PositivityTheory.ThK.SplitLoc.pRam'}
\leanref{Development/WeilBound/CoreStatements/CompactAveragingSplitting/EigenvectorLifting/SplitLoc.lean}{PositivityTheory.ThK.SplitLoc.invL}
\leanref{Development/WeilBound/CoreStatements/CompactAveragingSplitting/EigenvectorLifting/SplitLoc.lean}{PositivityTheory.ThK.SplitLoc.zCoeff}
\leanref{Development/WeilBound/CoreStatements/CompactAveragingSplitting/EigenvectorLifting/SplitLoc.lean}{PositivityTheory.ThK.SplitLoc.zetaUnr}
\leanref{Development/WeilBound/CoreStatements/CompactAveragingSplitting/EigenvectorLifting/SplitLoc.lean}{PositivityTheory.ThK.SplitLoc.pUnr}
\leanref{Development/WeilBound/CoreStatements/CompactAveragingSplitting/EigenvectorLifting/SplitLoc.lean}{PositivityTheory.ThK.SplitLoc.pUnr'}
\leanref{Development/WeilBound/CoreStatements/CompactAveragingSplitting/EigenvectorLifting/SplitLoc.lean}{PositivityTheory.ThK.SplitLoc.RamMaps}
\leanref{Development/WeilBound/CoreStatements/CompactAveragingSplitting/EigenvectorLifting/SplitDict.lean}{PositivityTheory.ThK.SplitDict.RamNonzero}
\leanref{Development/WeilBound/CoreStatements/CompactAveragingSplitting/EigenvectorLifting/SplitLoc.lean}{PositivityTheory.ThK.SplitLoc.UnrMaps}
\leanref{Development/AutomorphicSpectrum/Proofs/BlockDecomposition/IsSmoothCharOfCondLE.lean}{GlobalTheory.CompanionTheory.SplitNarrow.ramMaps_of}
\leanref{Development/WeilBound/CoreStatements/CompactAveragingSplitting/EigenvectorLifting/ExistsUpperProjectionInverseTranspose.lean}{PositivityTheory.ThK.SplitDict.ramified_map_ne_zero}
\leanref{Development/AutomorphicSpectrum/Proofs/BlockDecomposition/FinSplitOfNarrow.lean}{GlobalTheory.CompanionTheory.BernBridge.Split.unrMaps_of_lt}
\begin{proof}
This is a computation with the two zeta integrals, and its result is stated.
In both parts the equivariance under $\omega_A(1,h)$ and $\omega_B(1,h)$ holds
by definition, since $Z(\Xi)(h)$ and $Z'(\Psi)(h)$ are the values at $1$ of
$Z(\omega_A(1,h)\Xi)$ and $Z'(\omega_B(1,h)\Psi)$. The law of $I(\alpha,\beta)$
and the invariance under $\omega_A(g,1)$ and $\omega_B(g',1)$ come from the
invariance of the Haar measures: for $b$ block upper triangular,
$X_0(b^{-1})^{\mathsf T}=(b_A^{-1})^{\mathsf T}X_0$ and $Y_0b=b_DY_0$, so the
change of variables $g\mapsto g\,(b_A^{-1})^{\mathsf T}$, respectively
$g\mapsto gb_D$, gives the factor
$\alpha(\det b_A)\beta(\det b_D)|\det b_A|^{a'/2}|\det b_D|^{-a/2}$, using
$\alpha=\nu\beta$. The function $Z(\Xi)$ is invariant under right translation by
the elements $h$ of an open subgroup with $\omega_A(1,h)\Xi=\Xi$, and $Z'(\Psi)$
by those of an open subgroup with $\omega_B(1,h)\Psi=\Psi$.

In (1), the integral defining $Z'$ converges absolutely since $2a'<d$. For $Z$,
the function $\phi(g)=\Xi(gX_0(h^{-1})^{\mathsf T})$ vanishes unless the entries
of $g$ lie in $\varpi^N\mathcal O$ for some $N$, and is invariant under adding
matrices with entries in $\varpi^M\mathcal O$ for some $M$. On a double coset
$K\varpi^nK$ with some $n_i\geq M$, multiplying the $i$-th diagonal entry of
$\varpi^n$ by a unit $u$ does not change the value of $\phi$ and multiplies the
integrand by $\nu(u)$; as $\nu$ is ramified, the average over $K$ of the left
translates of the integrand vanishes there. So only finitely many double cosets
contribute, the sum defining $Z$ is the integral over $GL_a(k)$ of this average,
and it is linear. For the nonvanishing, take $\Xi$ supported near $X_0$ so that
$g\mapsto\Xi(gX_0)$ is the indicator function of a neighbourhood of $1$ in the
subgroup of $K$ on which $\nu\circ\det$ is $1$; then $Z(\Xi)(1)$ is a positive
volume. The same works for $Z'$ with $Y_0$.

In (2), $\nu$ is unramified. By the Hermite normal form, the power series
$\sum_mz_m(\phi)X^m$ of a Schwartz-Bruhat $\phi$ becomes a Laurent polynomial when
multiplied by $P_a(X)$, so $\zeta(\phi)$ is its value at $X=1$, and $\zeta$ is
linear and is multiplied by $\nu(\det g)^{-1}|\det g|^{-d/2}$ when $\phi$ is
translated on the left or the right by $g\in GL_a(k)$. For
$\phi=\mathbf 1_{M_a(\mathcal O)}$ the same count gives
$\sum_mz_m(\phi)X^m=P_a(X)^{-1}$, so $\zeta(\phi)=1$. For $a'$, the inequality
$2a'<d+2$ makes the integral absolutely convergent and gives
$P_{a'}(1)Z'(\Psi)(h)=\alpha(\det h)|\det h|^{a'/2}\zeta'(g\mapsto\Psi(gY_0h))$,
where $\zeta'$ is $\zeta$ with $a'$ in place of $a$. For $h\in GL_d(\mathcal O)$
the functions $g\mapsto\mathbf 1_{M_{a\times d}(\mathcal O)}(gX_0(h^{-1})^{\mathsf T})$
and $g\mapsto\mathbf 1_{M_{a'\times d}(\mathcal O)}(gY_0h)$ are
$\mathbf 1_{M_a(\mathcal O)}$ and $\mathbf 1_{M_{a'}(\mathcal O)}$, and
$\alpha$, $\beta$ and $|\cdot|$ are $1$ at $\det h$, which gives the values $1$.
The nonvanishing follows from these values.
% (1): pRam_add, pRam'_add, pRam_mem, pRam'_mem, pRam_inv, pRam'_inv, pRam_equiv, pRam'_equiv,
%   pRam_ne_zero, pRam'_ne_zero (ExistsUpperProjectionInverseTranspose.lean), through zint_unit,
%   zint_low, shells_left, shells_right, shellSum_of_KO, exists_bump and integrable_zint.
% (2): Unr.p_lin, p'_lin, p_mem, p'_mem, p_inv, p'_inv, p_equiv, p'_equiv, p_unit, p'_unit
%   (UpperProjectionInverseTransposeRightInverse.lean), through Jm_hnf, conv_finite, zetaUnr_add,
%   zetaUnr_mul_left, zetaUnr_mul_right, Jm_indicator, zetaUnr_indicator, zetaUnr_eq_integral and
%   pUnr'_eq. The surjectivity in RamMaps and UnrMaps is added in the proofs of
%   Lemmas asm-theta-agree-finite-ramified and -split.
\end{proof}

\begin{lemma}\label{asm-theta-agree-finite-characters}
Let $B$ be a theta base for a CM field $E$, let $a,a'\geq0$, let
$\mathcal X=(\alpha,\beta,S_0,Q)$ be a character datum proper for $(B,a,a')$
(Definition~\ref{asm-theta-agree-triple}(3)), with $\nu=\alpha/\beta$ and
$T=S_0\cup Q$, let $v$ be a finite place of $F$ and $t\colon E\to F_v$ an
$F$-embedding, and let $\alpha_t$, $\beta_t$ and $\nu_t$ be the characters of
$F_v^\times$ given by $\alpha$, $\beta$ and $\nu$ at the place of $E$ defined by
$t$. Then $\alpha_t$ and $\beta_t$ are unitary and smooth
(Definition~\ref{asm-theta-agree-finite-split-models}) and
$\nu_t=\alpha_t/\beta_t$. If $v\in T$, then $\nu_t$ is ramified; if
$v\notin T$, then $\alpha_t$ and $\beta_t$ are unramified.
\end{lemma}
\leanref{Development/WeilBound/CoreStatements/CompactAveragingSplitting/EigenvectorLifting/SplitAsm.lean}{PositivityTheory.ThK.SplitAsm.ChrFacts}
\leanref{Development/WeilBound/CoreStatements/CompactAveragingSplitting/EigenvectorLifting/SplitCharacterFacts.lean}{PositivityTheory.ThK.SplitAsm.split_character_facts}
\begin{proof}
Hecke characters are continuous with values of absolute value $1$, so
$\alpha_t$ and $\beta_t$ are unitary and their kernels are open, hence contain
$1+\varpi^e\mathcal O_v$ for some $e\geq1$. As $\nu=\alpha/\beta$,
$\nu_t=\alpha_t/\beta_t$. The place $w$ of $E$ defined by $t$ lies over $v$. If
$v\in T$, then $\nu$ is ramified at $w$ by properness, so $\nu_t$ is ramified.
If $v\notin T$, then $w$ is not above $T$, so $\alpha$ and $\beta$ are
unramified at $w$ by properness.
\end{proof}

\begin{definition}\label{asm-theta-agree-finite-local-objects}
Let $E$, $F$ and $c$ be as in Definition~\ref{asm-theta-agree-models}, let $v$ be
a finite place of $F$, and let $Y$ be a lattice space of rank $r$ with matrix
$J$ and lattice $L$ (Definition~\ref{asm-theta-agree-doubling}(1)).
\begin{enumerate}
\item Let $L_v\subset\mathbf A_E^r$ be the subgroup generated by the $a\,y$ with
  $a\in\mathbf A_F$ at $v$, $a_v\in\mathcal O_v$ and $y\in L$. The
  \emph{stabilizer at $v$} of $Y$ is the set $K_Y(v)$ of $g\in U_Y^{(v)}$ with
  $gL_v\subset L_v$ and $g^{-1}L_v\subset L_v$.
\item Let $r=2m$. A \emph{Lagrangian} of $Y$ at $v$ is a matrix
  $C\in M_{r\times m}(\mathbf A_E)$ at $v$ with $C^{\mathsf T}J\,c(C)=0$ such that
  $Ct\neq0$ for every nonzero $t\in\mathbf A_E^m$ at $v$. It is \emph{adapted} if
  there is a second Lagrangian $C'$ of $Y$ at $v$ such that every $y\in L_v$ is
  $Cs+C't$ with $Cs\in L_v$ and $C't\in L_v$. Its \emph{parabolic} $P_C$ is the
  set of $p\in U_Y^{(v)}$ with $pC=Ct$ for some $t\in M_m(\mathbf A_E)$, and
  $P_C^\times$ is the subgroup of the $p$ with $p,p^{-1}\in P_C$.
\item Let $(X,W,M,\chi_X,\chi_W,\omega)$ be a Weil pair for a theta base
  (Definition~\ref{asm-theta-agree-triple}(1)), with $X$ of rank $n$, $W$ of rank
  $2m$, $I=\{1,\ldots,n\}\times\{1,\ldots,2m\}$ and Schrödinger operators
  $\rho$, and let $C$ be a Lagrangian of $W$ at $v$. An \emph{evaluation form}
  along $C$ is a linear form $l$ on the functions on $F_v^I$ such that
  $l(\phi_1)=l(\phi)$ whenever $\phi$ is Schwartz-Bruhat,
  $z\in M_{n\times m}(\mathbf A_E)$ and $\rho(zC^{\mathsf T})$ acts at $v$ by
  $\phi\mapsto\phi_1$.
\item Let $(X,W,M,\chi_X,\chi_W,\omega)$ and $I$ be as in (3), with $W$ of any
  rank $d$. \emph{Local operators} at $v$ for the pair are linear operators
  $\omega_v(g,h)$, for $g\in U_X^{(v)}$ and $h\in U_W^{(v)}$, and $\rho_v(u)$, for
  $u\in\mathbb W(X,W)$ at $v$, on the functions on $F_v^I$, such that for all
  such $g,g',h,h',u$ and every Schwartz-Bruhat $\phi$: $\omega_v(g,h)\phi$ and
  $\rho_v(u)\phi$ are Schwartz-Bruhat; $\omega(g,h)$ acts at $v$ by
  $\phi\mapsto\omega_v(g,h)\phi$ and $\rho(u)$ by $\phi\mapsto\rho_v(u)\phi$;
  $\omega_v(1,1)\phi=\phi$ and
  $\omega_v(gg',hh')\phi=\omega_v(g,h)\omega_v(g',h')\phi$;
  $\omega_v(g,h)\rho_v(u)\phi=\rho_v(guh^{\mathsf T})\omega_v(g,h)\phi$; and there
  is a neighbourhood of $1$ in $U_W$ whose elements $h$ at $v$ have
  $\omega_v(1,h)\phi=\phi$.
\end{enumerate}
\end{definition}
\leanref{Development/WeilBound/CoreStatements/CompactAveragingSplitting/EigenvectorLifting/InertDef.lean}{PositivityTheory.ThK.latW}
\leanref{Development/WeilBound/CoreStatements/CompactAveragingSplitting/EigenvectorLifting/InertDef.lean}{PositivityTheory.ThK.KAt}
\leanref{Development/WeilBound/CoreStatements/CompactAveragingSplitting/EigenvectorLifting/InertDef.lean}{PositivityTheory.ThK.IsLagAt}
\leanref{Development/WeilBound/CoreStatements/CompactAveragingSplitting/EigenvectorLifting/InertDef.lean}{PositivityTheory.ThK.IsAdaptedAt}
\leanref{Development/WeilBound/CoreStatements/CompactAveragingSplitting/EigenvectorLifting/InertDef.lean}{PositivityTheory.ThK.PAt}
\leanref{Development/WeilBound/ThetaLifts/SphericalLattices/FunDefs.lean}{PositivityTheory.Sph.PSub}
\leanref{Development/WeilBound/CoreStatements/CompactAveragingSplitting/EigenvectorLifting/InertDef.lean}{PositivityTheory.ThK.IsEvAt}
\leanref{Development/WeilBound/CoreStatements/CompactAveragingSplitting/EigenvectorLifting/InertDef.lean}{PositivityTheory.ThK.IsLocRep}

\begin{lemma}\label{asm-theta-agree-finite-inputs-slice}
Let $F$ be a number field with adele ring $\mathbf A$ and finite adeles
$\mathbf A_f$, let $I$ be a finite set and $v$ a finite place of $F$. For
$x\in\mathbf A^I$ and $y\in F_v^I$ let $x\langle y\rangle\in\mathbf A^I$ have
component $y$ at $v$ and the components of $x$ elsewhere. For a function $\phi$
on $F_v^I$ and $a\in F_v^I$ let $P_{\phi,a}$ be the operator
$(P_{\phi,a}\Phi)(x)=\phi(x_v)\Phi(x\langle a\rangle)$ on functions on
$\mathbf A^I$, and for $x_0\in\mathbf A^I$ let $\operatorname{res}_{x_0}$ map a
function $\Phi$ on $\mathbf A^I$ to the function $y\mapsto\Phi(x_0\langle y\rangle)$
on $F_v^I$. Call a linear form on functions on $\mathbf A^I$ admissible as in
Definition~\ref{asm-theta-agree-doubling}(6), and let $\phi$ be a
Schwartz-Bruhat function on $F_v^I$ and $a\in F_v^I$.
\begin{enumerate}
\item $P_{\phi,a}$ maps Schwartz-Bruhat functions on $\mathbf A^I$
  (Definition~\ref{asm-theta-agree-models}) to Schwartz-Bruhat functions, and
  $\lambda\circ P_{\phi,a}$ is admissible for every admissible $\lambda$.
\item The Schwartz-Bruhat functions on $\mathbf A^I$ are spanned by the
  $P_{\phi',a'}\Psi$ with $\phi'$ Schwartz-Bruhat on $F_v^I$, $a'\in F_v^I$ and
  $\Psi$ Schwartz-Bruhat on $\mathbf A^I$. A linear form $\lambda$ is
  admissible if $\lambda\circ P_{\phi',a'}$ is admissible for every
  Schwartz-Bruhat $\phi'$ on $F_v^I$ and every $a'\in F_v^I$.
\item For $x_0\in\mathbf A^I$, $\operatorname{res}_{x_0}$ maps Schwartz-Bruhat
  functions on $\mathbf A^I$ to Schwartz-Bruhat functions on $F_v^I$, and
  $\lambda\circ\operatorname{res}_{x_0}$ is admissible for every linear form
  $\lambda$ on functions on $F_v^I$.
\item There are $x_0\in\mathbf A^I$ and a function $r_0$ on $\mathbf A^I$ that
  does not depend on the $v$-component, with $r_0(x_0)=1$, such that
  $\phi'\otimes r_0$ (Definition~\ref{asm-theta-agree-models}(2)) is
  Schwartz-Bruhat for every Schwartz-Bruhat function $\phi'$ on $F_v^I$.
\end{enumerate}
\end{lemma}
\leanref{Development/WeilBound/CoreStatements/LocalThetaKernels/LocrepAd.lean}{PositivityTheory.ThKH.Locrep.PA_mem}
\leanref{Development/WeilBound/CoreStatements/LocalThetaKernels/LocrepAd.lean}{PositivityTheory.ThKH.Locrep.PA_adm}
\leanref{Development/WeilBound/CoreStatements/LocalThetaKernels/LocrepSB.lean}{PositivityTheory.ThKH.Locrep.Pm_mem_SB}
\leanref{Development/WeilBound/CoreStatements/LocalThetaKernels/LocrepSB.lean}{PositivityTheory.ThKH.Locrep.mem_span_Pm}
\leanref{Development/WeilBound/CoreStatements/LocalThetaKernels/LocrepAd.lean}{PositivityTheory.ThKH.Locrep.schA_le}
\leanref{Development/WeilBound/CoreStatements/LocalThetaKernels/LocrepAd.lean}{PositivityTheory.ThKH.Locrep.admissible_of_integral_conj_le}
\leanref{Development/WeilBound/CoreStatements/LocalThetaKernels/LocrepAd.lean}{PositivityTheory.ThKH.Locrep.res_mem}
\leanref{Development/WeilBound/CoreStatements/LocalThetaKernels/LocrepAd.lean}{PositivityTheory.ThKH.Locrep.res_adm}
\leanref{Development/WeilBound/CoreStatements/LocalThetaKernels/LocrepAd.lean}{PositivityTheory.ThKH.Locrep.exists_base}
\begin{proof}
For a function $g$ on $\mathbf A_f^I$ put
$(P^f_{\phi,a}g)(x)=\phi(x_v)g(x\langle a\rangle)$; then
$P_{\phi,a}(f\otimes g)=f\otimes P^f_{\phi,a}g$ for $f$ on $F_\infty^I$. Let $g$
be Schwartz-Bruhat. Then $P^f_{\phi,a}g$ is locally constant. If
$P^f_{\phi,a}g(x)\neq0$, then $x\langle a\rangle$ lies in the support of $g$,
$x_v$ in that of $\phi$, and $x=x\langle a\rangle+j(x_v)-j(a)$, where $j(y)$ has
component $y$ at $v$ and $0$ elsewhere. So $P^f_{\phi,a}g$ vanishes outside a
compact set and is Schwartz-Bruhat. This gives (1), as
$f\mapsto\lambda(f\otimes P^f_{\phi,a}g)$ is continuous for admissible $\lambda$.

For (2), let $g$ be Schwartz-Bruhat on $\mathbf A_f^I$, and for $y\in F_v^I$ let
$g_y(x)=g(x\langle y\rangle)$. As $g$ is locally constant with compact support,
$y\mapsto g_y$ is locally constant, and $g_y=0$ for $y$ outside the compact image
$C$ of the support of $g$ under $x\mapsto x_v$. So the $g_y$ take finitely many
nonzero values $h_1,\ldots,h_k$. Let $\phi_j$ be the indicator function of
$\{y:g_y=h_j\}$, a Schwartz-Bruhat function as this set is open and closed and
lies in $C$, and choose $y_j$ with $g_{y_j}=h_j$. As $g(x)=g_{x_v}(x)$,
$g=\sum_jP^f_{\phi_j,y_j}g$, and so $f\otimes g=\sum_jP_{\phi_j,y_j}(f\otimes g)$.
This gives the spanning, and also
$\lambda(f\otimes g)=\sum_j\lambda(P_{\phi_j,y_j}(f\otimes g))$, which is continuous
in $f$ when every $\lambda\circ P_{\phi_j,y_j}$ is admissible.

For (3), $\operatorname{res}_{x_0}(P_{\phi',a'}\Psi)(z)=\phi'(z)\Psi(x_0\langle a'\rangle)$,
so $\operatorname{res}_{x_0}(P_{\phi',a'}\Psi)=\Psi(x_0\langle a'\rangle)\phi'$ is
Schwartz-Bruhat, and (2) gives the first claim. For $f\otimes g$ as above,
$\operatorname{res}_{x_0}(f\otimes g)=f(x_{0,\infty})\cdot g'$ with
$g'(y)=g(x_{0,f}\langle y\rangle)$, so $\lambda(\operatorname{res}_{x_0}(f\otimes g))$
is $\lambda(g')$ times the value of $f$ at a point, which is continuous in $f$.

For (4), let $\Phi$ be a nonzero Schwartz-Bruhat function on $\mathbf A^I$, for
example one with $\theta(\Phi)\neq0$
(Theorem~\ref{asm-theta-agree-evaluation-heisenberg}(2)), and choose $x_0$ with
$\Phi(x_0)\neq0$. Put $r_0(x)=\Phi(x_0)^{-1}\Phi(x\langle x_{0,v}\rangle)$. Then
$r_0$ does not depend on $x_v$, $r_0(x_0)=1$, and
$\phi'\otimes r_0=\Phi(x_0)^{-1}P_{\phi',x_{0,v}}\Phi$, which is Schwartz-Bruhat
by (1).
\end{proof}

\begin{lemma}\label{asm-theta-agree-finite-inputs-schur}
Let $F$, $\mathbf A$, $I$, $v$, $x\langle y\rangle$, $P_{\phi,a}$ and
$\operatorname{res}_{x_0}$ be as in
Lemma~\ref{asm-theta-agree-finite-inputs-slice}, and let $x_0$ and $r_0$ be as
in its part (4). Let $\psi$ be a continuous character of $\mathbf A$, trivial on
$F$ and not $1$, with Schrödinger operators $\rho(w)$ for
$w\in\mathbf A^I\times\mathbf A^I$ as in
Theorem~\ref{asm-theta-agree-evaluation-heisenberg}, and let $\phi$ be a
Schwartz-Bruhat function on $F_v^I$.
\begin{enumerate}
\item If an operator $T$ on functions on $\mathbf A^I$ acts at $v$ by
  $\phi\mapsto\phi_1$ (Definition~\ref{asm-theta-agree-models}(2)), then
  $\phi_1=\operatorname{res}_{x_0}(T(\phi\otimes r_0))$.
\item Let $T$ be a linear operator on functions on $\mathbf A^I$ that maps
  Schwartz-Bruhat functions to Schwartz-Bruhat functions, such that
  $\lambda\circ T$ is admissible for every admissible $\lambda$ and
  $T\rho(w)\Phi=\rho(w)T\Phi$ for every Schwartz-Bruhat $\Phi$ and every
  $w=(w_1,w_2)$ with $w_{1,v}=w_{2,v}=0$. Then $T$ acts at $v$ by
  $\phi\mapsto\operatorname{res}_{x_0}(T(\phi\otimes r_0))$.
\end{enumerate}
\end{lemma}
\leanref{Development/WeilBound/CoreStatements/LocalThetaKernels/LocrepAd.lean}{PositivityTheory.ThKH.Locrep.eq_of_acts}
\leanref{Development/WeilBound/CoreStatements/LocalThetaKernels/LocrepAd.lean}{PositivityTheory.ThKH.Locrep.acts_res}
\leanref{Development/WeilBound/CoreStatements/LocalThetaKernels/LocrepAd.lean}{PositivityTheory.ThKH.Locrep.Sop_PA}
\leanref{Development/WeilBound/CoreStatements/LocalThetaKernels/LocrepAd.lean}{PositivityTheory.ThKH.Locrep.PA_schr}
\leanref{Development/WeilBound/CoreStatements/LocalThetaKernels/LocrepAd.lean}{PositivityTheory.ThKH.Locrep.Sop_schr}
\leanref{Development/WeilBound/CoreStatements/LocalThetaKernels/LocrepAd.lean}{PositivityTheory.ThKH.Locrep.PA_tens}
\begin{proof}
(1) As $r_0$ does not depend on the $v$-component and $\phi\otimes r_0$ is
Schwartz-Bruhat, $T(\phi\otimes r_0)=\phi_1\otimes r_0$. At $x_0\langle y\rangle$
this reads $T(\phi\otimes r_0)(x_0\langle y\rangle)=\phi_1(y)r_0(x_0)=\phi_1(y)$.

(2) Put $\phi_1=\operatorname{res}_{x_0}(T(\phi\otimes r_0))$. If $\phi=0$, then
$T(\phi\otimes r)=0=\phi_1\otimes r$ for every $r$. Otherwise fix $b$ with
$\phi(b)\neq0$, fix $y\in F_v^I$, and let $S$ be the operator
$(S\Phi)(x)=T(P_{\phi,x_v}\Phi)(x\langle y\rangle)$ on functions on $\mathbf A^I$.

As $P_{\phi,c}P_{\phi',a'}=\phi'(c)P_{\phi,a'}$, one has
$S\circ P_{\phi',a'}=P_{\phi',y}\circ T\circ P_{\phi,a'}$ for every $\phi'$ and
$a'$. By Lemma~\ref{asm-theta-agree-finite-inputs-slice}(1) and (2) and the
hypotheses on $T$, $S$ maps Schwartz-Bruhat functions to Schwartz-Bruhat
functions and $\lambda\circ S$ is admissible for every admissible $\lambda$.

Let $w=(w_1,w_2)$, let $w'=(w_1\langle0\rangle,w_2\langle0\rangle)$, and let
$j(c)$ have component $c$ at $v$ and $0$ elsewhere. A direct computation gives
\[
P_{\phi,c}\,\rho(w)\Phi=\psi\bigl(j(c)\cdot w_2+\tfrac12(w_1\cdot w_2-w'_1\cdot w'_2)\bigr)\,
\rho(w')\,P_{\phi,c+w_{1,v}}\Phi
\]
for every $c\in F_v^I$. Take $c=x_v$, apply $T$, which commutes with $\rho(w')$
on Schwartz-Bruhat functions, and evaluate at $x\langle y\rangle$. As
$x\langle y\rangle+w'_1=(x+w_1)\langle y\rangle$ and
$j(x_v)\cdot w_2+x\langle y\rangle\cdot w'_2=x\cdot w_2$, this gives
$S\rho(w)\Phi=\rho(w)S\Phi$ for every $w$ and Schwartz-Bruhat $\Phi$. By
Theorem~\ref{asm-theta-agree-evaluation-heisenberg}(5) there is $c_y\in\mathbf C$
with $S\Phi=c_y\Phi$ for every Schwartz-Bruhat $\Phi$.

Let $r$ be a function on $\mathbf A^I$ that does not depend on the
$v$-component, with $\phi\otimes r$ Schwartz-Bruhat. Then
$P_{\phi,b}(\phi\otimes r)=\phi(b)\,\phi\otimes r$, and evaluating
$S(\phi\otimes r)=c_y\,\phi\otimes r$ at $x\langle b\rangle$ gives
$\phi(b)\,T(\phi\otimes r)(x\langle y\rangle)=c_y\phi(b)r(x)$. So
$T(\phi\otimes r)(x\langle y\rangle)=c_yr(x)$ for every $x$. For $r=r_0$ and
$x=x_0$ this gives $c_y=\phi_1(y)$. Taking $y=x_v$ for each $x$ gives
$T(\phi\otimes r)(x)=\phi_1(x_v)r(x)$, that is,
$T(\phi\otimes r)=\phi_1\otimes r$.
\end{proof}

\begin{lemma}\label{asm-theta-agree-finite-inputs-powers}
Let $R$ be a normed complex algebra with unit $1$, and let $a\in R$ satisfy
$\|a^p-1\|<\tfrac12$ for every integer $p\geq0$. Then $a=1$.
\end{lemma}
\leanref{Development/WeilBound/CoreStatements/LocalThetaKernels/LinearIndependentOfNotMem.lean}{PositivityTheory.ThKH.Locrep.eq_one_of_pow_near}
\leanref{Development/AutomorphicArithmetic/IntegralCasts/CapelliIdentities/ImportsSideBis.lean}{UpperTheory.Reformulated.B.Yhd7.main}
\begin{proof}
Let $N\geq1$ and $S=\sum_{k<N}a^k$. Then $S(a-1)=a^N-1$ and
$N(a-1)=(N-S)(a-1)+(a^N-1)$, where $\|N-S\|\leq\sum_{k<N}\|1-a^k\|\leq N/2$. So
$N\|a-1\|<\tfrac N2\|a-1\|+\tfrac12$, that is, $N\|a-1\|<1$, for every $N$.
Hence $a=1$.
\end{proof}

\begin{lemma}\label{asm-theta-agree-finite-inputs-smooth}
Let $G$ be a topological group with subgroups $K_0\supset K_1\supset\cdots$ such
that every neighbourhood of $1$ contains some $K_n$. Let $V$ be a complex vector
space, $V_1\subset V$ a subspace, and $\omega$ a map from $G$ to the linear
operators on $V$ such that $\omega(g)V_1\subset V_1$ for $g\in K_0$,
$\omega(1)\phi=\phi$ and $\omega(gg')\phi=\omega(g)\omega(g')\phi$ for
$g,g'\in K_0$ and $\phi\in V_1$, and $g\mapsto\lambda(\omega(g)\phi)$ is
continuous on $G$ for every linear form $\lambda$ on $V$ and every $\phi\in V_1$.
Then for every $\phi\in V_1$ there is $n$ with $\omega(k)\phi=\phi$ for all
$k\in K_n$.
\end{lemma}
\leanref{Development/WeilBound/CoreStatements/LocalThetaKernels/LinearIndependentOfNotMem.lean}{PositivityTheory.ThKH.Locrep.exists_finite_orbit}
\leanref{Development/WeilBound/CoreStatements/LocalThetaKernels/LinearIndependentOfNotMem.lean}{PositivityTheory.ThKH.Locrep.smooth_of_weak}
\begin{proof}
Let $\phi\in V_1$ and $f(g)=\omega(g)\phi-\phi$. First, there are $n$ and a
finite set $s\subset V$ with $f(k)$ in the span of $s$ for every $k\in K_n$.
Otherwise choose inductively $k_n\in K_n$ with $u_n=f(k_n)$ outside the span of
$u_0,\ldots,u_{n-1}$. The $u_n$ are linearly independent, so some linear form
$\lambda$ has $\lambda(u_n)=1$ for every $n$. As $\lambda\circ f$ is continuous
and $f(1)=0$, the set of $g$ with $\lambda(f(g))\neq1$ is a neighbourhood of
$1$, so it contains some $K_n$, and $k_n$ in it gives a contradiction.

So the span $V_0$ of the $\omega(k)\phi$, $k\in K_n$, is finite-dimensional. It
contains $\phi$, lies in $V_1$, and is stable under $\omega(k)$ for $k\in K_n$.
Choose a linear map $\pi\colon V\to V_0$ that is the identity on $V_0$, a basis
of $V_0$, and a norm on the endomorphisms of $V_0$ written as matrices in it, and
put $A(g)=\pi\circ\omega(g)|_{V_0}$. Its matrix entries are continuous in $g$,
$A(1)=1$, and $A(k)^p=A(k^p)$ for $k\in K_n$ and $p\geq0$. Choose $m$ with
$\|A(g)-1\|<\tfrac12$ for $g\in K_m$, and let $k\in K_{\max(n,m)}$. Every power of
$k$ lies in $K_m$, so $\|A(k)^p-1\|<\tfrac12$ for every $p$, and $A(k)=1$ by
Lemma~\ref{asm-theta-agree-finite-inputs-powers}. As $\omega(k)\phi\in V_0$,
$\omega(k)\phi=A(k)\phi=\phi$.
\end{proof}

\begin{lemma}\label{asm-theta-agree-finite-inputs-subgroups}
Let $E$ and $F$ be as in Definition~\ref{asm-theta-agree-models}, let $Y$ be a
lattice space of rank $d$ (Definition~\ref{asm-theta-agree-doubling}(1)), and
let $v$ be a finite place of $F$. Then there are subgroups
$K_0\supset K_1\supset\cdots$ of $U_Y^{(v)}$
(Definition~\ref{asm-theta-agree-models}(2)) such that every neighbourhood of
$1$ in $U_Y$ contains some $K_n$, and for every $n$ some neighbourhood $U$ of
$1$ in $U_Y$ has $U\cap U_Y^{(v)}\subset K_n$.
\end{lemma}
\leanref{Development/WeilBound/CoreStatements/LocalThetaKernels/ExistsSmallSubgroups.lean}{PositivityTheory.ThKH.Locrep.exists_small_lattices}
\leanref{Development/WeilBound/CoreStatements/LocalThetaKernels/ExistsSmallSubgroups.lean}{PositivityTheory.ThKH.Locrep.exists_small_subgroups}
\begin{proof}
Fix an $F$-basis $b_1,\ldots,b_e$ of $E$, write $z=\sum_iz_ib_i$ with
$z_i\in\mathbf A$ for $z\in\mathbf A_E$, and let $C\geq1$ bound the absolute
values at $v$ of the coordinates of all $b_ib_j$. Let $\Lambda_n$ be the set of
$z\in\mathbf A_E$ at $v$ with $|z_{i,v}|\leq2^{-n}/C$ for all $i$. By the
ultrametric inequality $\Lambda_n$ is a group under addition, and it is closed
under multiplication, since each coordinate of $zz'$ is a sum of terms
$z_iz'_j(b_ib_j)_k$ of absolute value at most $(2^{-n}/C)^2C\leq2^{-n}/C$. Every
neighbourhood of $0$ in $\mathbf A_E$ contains some $\Lambda_n$. For each $n$
the set $N_n$ of $z\in\mathbf A_E$ with $|z_{i,v}|\leq2^{-n}/C$ for all $i$ is a
neighbourhood of $0$, and its elements at $v$ lie in $\Lambda_n$.

Let $K_n$ be the set of $h\in U_Y$ such that all entries of $h-1$ and
$h^{-1}-1$ lie in $\Lambda_n$. As $hk-1=(h-1)(k-1)+(h-1)+(k-1)$ and
$(hk)^{-1}=k^{-1}h^{-1}$, $K_n$ is a subgroup; it lies in $U_Y^{(v)}$ and
decreases with $n$. A neighbourhood of $1$ in $U_Y$ contains the $h$ with $h-1$
and $h^{-1}-1$ near $0$, hence some $K_n$. Conversely, by continuity of
$h\mapsto h-1$ and $h\mapsto h^{-1}-1$ there is a neighbourhood $U$ of $1$ whose
elements have all entries of $h-1$ and $h^{-1}-1$ in $N_n$. For $h\in U$ at $v$,
$h-1$ and $h^{-1}-1=-h^{-1}(h-1)$ are at $v$, so $h\in K_n$.
\end{proof}

\begin{theorem}[Bruhat, Tits \cite{BT72,Car79}]\label{asm-theta-agree-finite-inputs-bruhat-tits}
Let $k$ be a nonarchimedean local field, let $k_2$ be its unramified quadratic
extension, with ring of integers $\mathcal O_2$ and nontrivial automorphism $c$
over $k$, let $d\geq1$, let $A_d\in M_d(k_2)$ have entries $1$ on the
antidiagonal and $0$ elsewhere, let $G$ be the group of $g\in GL_d(k_2)$ with
$g^{\mathsf T}A_d\,c(g)=A_d$, and let $K=G\cap GL_d(\mathcal O_2)$.
\begin{enumerate}
\item Every $g\in G$ has $g^{-1}\in KgK$.
\item Let $d=2m$, and let $P$ be the subgroup of the $g\in G$ whose lower left
  $m\times m$ block is $0$. Every $g\in G$ is $pk$ with $p\in P$ and $k\in K$.
  Every $Z\in M_{2m\times m}(k_2)$ of rank $m$ with $Z^{\mathsf T}A_{2m}c(Z)=0$
  has $kZ=\left(\begin{smallmatrix}T\\0\end{smallmatrix}\right)$ for some
  $k\in K$ and $T\in GL_m(k_2)$.
\end{enumerate}
\end{theorem}
\leanref{Development/WeilBound/ThetaLifts/SphericalLattices/InvMemKgKCart2.lean}{PositivityTheory.Sph.inv_mem_KgK_cart2}
\leanref{Development/WeilBound/CoreStatements/LocalThetaKernels/IwasLoc.lean}{PositivityTheory.ThKH.Iwas.group_main}
\leanref{Development/WeilBound/ThetaLifts/SphericalLattices/MemIntGLIff.lean}{PositivityTheory.Sph.M2.Chart.exists_move}

\begin{lemma}\label{asm-theta-agree-finite-inputs-modulus}
Let $k$, $k_2$, $\mathcal O_2$, $c$, $A_{2m}$, $G$, $K$ and $P$ be as in
Theorem~\ref{asm-theta-agree-finite-inputs-bruhat-tits} with $d=2m$, let $q$ be
the number of elements of the residue field of $k$, let $\varpi$ be a
uniformizer of $k$ and $\operatorname{ord}$ the valuation of $k_2$ with
$\operatorname{ord}(\varpi)=1$, and put $H=P\cap K$. For subgroups $H_1$ let
$[H:H_1]$ be the index of $H\cap H_1$ in $H$. Then for every $p\in P$, with upper
left $m\times m$ block $p_{11}$,
\[
[H:p^{-1}Hp]=q^{-2m\operatorname{ord}(\det p_{11})}\,[H:pHp^{-1}].
\]
\end{lemma}
\leanref{Development/WeilBound/ThetaLifts/SphericalLattices/LocModSt.lean}{PositivityTheory.Sph.M2.LocModSt}
\leanref{Development/WeilBound/ThetaLifts/SphericalLattices/LocMod.lean}{PositivityTheory.Sph.M2.locMod}
\leanref{Development/WeilBound/ThetaLifts/SphericalLattices/OrdDetUniformizerImageEq.lean}{PositivityTheory.Sph.M2.locMod_of}
\begin{proof}
Both indices are finite and nonzero for $p\in P$. The function
$\delta(p)=[H:pHp^{-1}]/[H:p^{-1}Hp]$ is the modular character of $P$ with
respect to its compact open subgroup $H$: it is multiplicative and positive on
$P$, and $1$ on $H$. So $\chi(p)=\delta(p)\,q^{-2m\operatorname{ord}(\det p_{11})}$
is multiplicative and positive on $P$ and $1$ on $H$. Let $z$ be the diagonal
matrix with entries $\varpi$ in the first $m$ places and $\varpi^{-1}$ in the
last $m$; it lies in $P$. A count of cosets gives $zHz^{-1}\subset H$ with index
$q^{2m^2}$, so $\delta(z)=q^{2m^2}$ and $\chi(z)=1$. A multiplicative positive
function on $P$ that is $1$ on $H$ and at $z$ is $1$ on $P$. So $\chi=1$.
\end{proof}
% Lean: locMod = locMod_of locComm zCount gen. locComm: the indices are nonzero; the modular character
%   is GlobalTheory.T.haarInvVolume_conj_modular_character; zCount is ZCountSt (the count for z, written
%   zW L m x); gen is GenSt (the last step).

\begin{lemma}\label{asm-theta-agree-finite-inputs-frame}
Let $B$ be a theta base for a CM field $E$ with maximal totally real subfield
$F$ (Definition~\ref{asm-theta-agree-models}(1)), let $W$ be a skew-Hermitian
lattice space of rank $d$ and some signature
(Definition~\ref{asm-theta-agree-doubling}(1)), with matrix $J$, and let $v$ be
a finite place of $F$ inert in $E$, with $w$ the place of $E$ above it. Let
$\mathcal O_w$ be the ring of integers of $E_w$, let $c$ also denote the
nontrivial automorphism of $E_w$ over $F_v$, write $g_w$ for the component at
$w$ of a matrix $g$ over $\mathbf A_E$, and let $A_d$, $G_w$ and $K_w$ be the
$A_d$, $G$ and $K$ of Theorem~\ref{asm-theta-agree-finite-inputs-bruhat-tits}
for $k=F_v$ and $k_2=E_w$.
\begin{enumerate}
\item There are $a\in GL_d(E_w)$ and $l\in E_w^\times$ with
  $a^{\mathsf T}J\,c(a)=l\,A_d$ such that a matrix $g\in M_d(\mathbf A_E)$ has
  $gL_v\subset L_v$ (Definition~\ref{asm-theta-agree-finite-local-objects}(1))
  if and only if $a^{-1}g_wa$ has entries in $\mathcal O_w$, and an element of
  $\mathbf A_E^d$ lies in $L_v$ if and only if it is at $v$ and its component at
  $w$ lies in $a\,\mathcal O_w^d$.
\item For $a$ as in (1), $\Phi(g)=a^{-1}g_wa$ is an isomorphism of groups from
  $U_W^{(v)}$ onto $G_w$ that carries $K_W(v)$ onto $K_w$.
\item Let $d=2m$, and let $P_w$ be the subgroup of $G_w$ of
  Theorem~\ref{asm-theta-agree-finite-inputs-bruhat-tits}(2). For every
  Lagrangian $C$ of $W$ at $v$
  (Definition~\ref{asm-theta-agree-finite-local-objects}(2)) there is $a$ as in
  (1) with $a^{-1}C_w=\left(\begin{smallmatrix}T\\0\end{smallmatrix}\right)$
  for some $T\in GL_m(E_w)$. For such $a$, $\Phi$ carries $P_C^\times$ onto
  $P_w$, and if $p\in P_C^\times$ and $t\in M_m(\mathbf A_E)$ have $pC=Ct$, then
  the upper left $m\times m$ block of $\Phi(p)$ has determinant $\det t_w$.
\end{enumerate}
\end{lemma}
\leanref{Development/WeilBound/CoreStatements/LocalThetaKernels/IwasFrame.lean}{PositivityTheory.ThKH.Iwas.exists_frame}
\leanref{Development/WeilBound/CoreStatements/LocalThetaKernels/IwasMain.lean}{PositivityTheory.ThKH.Iwas.stab_iff}
\leanref{Development/WeilBound/CoreStatements/LocalThetaKernels/IwasAd.lean}{PositivityTheory.ThKH.Iwas.mem_latW_iff}
\leanref{Development/WeilBound/CoreStatements/LocalThetaKernels/IwasAd.lean}{PositivityTheory.ThKH.Iwas.mat_inj}
\leanref{Development/WeilBound/CoreStatements/LocalThetaKernels/IwasMain.lean}{PositivityTheory.ThKH.Iwas.exists_lift_UA}
\leanref{Development/WeilBound/ThetaLifts/SphericalLattices/ExistsFrameLag.lean}{PositivityTheory.Sph.M2.Chart.exists_frame_lag}
\leanref{Development/WeilBound/ThetaLifts/SphericalLattices/ChartSt.lean}{PositivityTheory.Sph.M2.ChartSt}
\leanref{Development/WeilBound/ThetaLifts/SphericalLattices/ExistsFrameLag.lean}{PositivityTheory.Sph.M2.chart}
\begin{proof}
(1) As $E/F$ is unramified at $v$, the completion $L_w$ of $L$ at $w$ is a
unimodular lattice for the Hermitian form $l^{-1}h$, $h(x,y)=x^{\mathsf T}J\,c(y)$,
for a suitable $l\in E_w^\times$ with $c(l)=-l$. Over the unramified extension
$E_w/F_v$ the unimodular Hermitian lattices of rank $d$ are all isometric
\cite[Section~7]{Jac62}, so $L_w$ has a basis in which $l^{-1}h$ has matrix
$A_d$. Let $a$ have these vectors as columns. As $v$ is inert, the components
of $\mathbf A_E$ at $v$ form the field $E_w$, and an element of $\mathbf A_E^d$
lies in $L_v$ if and only if it is at $v$ and its component at $w$ lies in
$L_w=a\,\mathcal O_w^d$. The criterion for matrices follows.

(2) If $g\in U_W$, then $g^{\mathsf T}Jc(g)=J$, and so
$\Phi(g)^{\mathsf T}(l\,A_d)c(\Phi(g))=l\,A_d$; dividing by $l$, $\Phi(g)\in G_w$.
$\Phi$ is multiplicative. At the inert place $v$ the components of $\mathbf A_E$
at $v$ form the field $E_w$, so an element of $U_W^{(v)}$ is determined by its
component at $w$, and $\Phi$ is injective on $U_W^{(v)}$. For $M\in G_w$, the
matrix with component $aMa^{-1}$ at $w$ and $1$ elsewhere lies in $U_W^{(v)}$
and maps to $M$. By (1) for $g$ and $g^{-1}$, $g\in K_W(v)$ if and only if
$\Phi(g)$ and $\Phi(g)^{-1}$ have entries in $\mathcal O_w$, that is,
$\Phi(g)\in K_w$.

(3) Let $a$ be as in (1). As $C$ is a Lagrangian, $Z=a^{-1}C_w$ has rank $m$ and
$Z^{\mathsf T}(l\,A_{2m})c(Z)=C_w^{\mathsf T}Jc(C_w)=0$. By
Theorem~\ref{asm-theta-agree-finite-inputs-bruhat-tits}(2) there are $k\in K_w$
and $T\in GL_m(E_w)$ with $kZ=\left(\begin{smallmatrix}T\\0\end{smallmatrix}\right)$.
Replace $a$ by $ak^{-1}$; as $k\in K_w$, (1) still holds. For
$g\in U_W^{(v)}$ and $t$ as in the statement, $gC=Ct$ gives
$\Phi(g)\left(\begin{smallmatrix}T\\0\end{smallmatrix}\right)=\left(\begin{smallmatrix}T\\0\end{smallmatrix}\right)t_w$,
so the lower left block of $\Phi(g)$ is $0$ and its upper left block is
$Tt_wT^{-1}$. Conversely, if $\Phi(g)\in P_w$ with upper left block $u$, then
$g_wC_w=C_w\,T^{-1}uT$, and as $C$ and $g-1$ are at $v$, $gC=Ct$ for the matrix
$t$ at $v$ with $t_w=T^{-1}uT$. As $P_w$ is a group, $\Phi$ carries $P_C^\times$
onto $P_w$.
\end{proof}
% Lean: (1) is exists_frame with stab_iff; the Lean finds a from a rational frame of L at v, gram_bounds,
%   Formalism.UG.P.cartan_exists_glint and exists_congr_antiJ. (2) is iso_smul_iff, mat_inj and
%   exists_lift_UA. (3) is exists_frame_lag (exists_frame and exists_move) and chart (phi_par).

\begin{lemma}\label{asm-theta-agree-dictionary-change}
Let $B$ be a theta base for a CM field $E$, let $M$ be a model for lattice
spaces $(X,W)$ of ranks $n$ and $d$ (Definition~\ref{asm-theta-agree-models}(3)),
put $I=\{1,\ldots,n\}\times\{1,\ldots,d\}$, and let $v$ be a finite place of
$F$. Let $\pi^{(1)},\pi^{(0)}\colon\mathbf A_E\to F_v$ be ring homomorphisms,
let $\operatorname{co}$ and $\operatorname{co}_1$ be maps from
$\mathbb W(X,W)$ to $M_{n\times d}(F_v)$, and let $P\in GL_n(F_v)$,
$Q\in GL_d(F_v)$ and $s\in\mathbf C^\times$ with
$P\operatorname{co}_1(u)Q=\operatorname{co}(u)$ for every $u$. For a linear map
$\iota$ from the functions on $F_v^I$ to the functions on $M_{n\times d}(F_v)$
put $\iota_1\phi(Y)=s\,\iota\phi(PYQ)$.
\begin{enumerate}
\item If $\iota$ is a linear model for $(\pi^{(1)},\pi^{(0)},\operatorname{co})$
  (Definition~\ref{asm-theta-agree-models}(5)), then $\iota_1$ is a linear model
  for $(\pi^{(1)},\pi^{(0)},\operatorname{co}_1)$.
\item Let $T$ be an operator on the functions on $\mathbf A^I$, and let $S$ and
  $S_1$ be operators on the functions on $M_{n\times d}(F_v)$ with
  $s\,(S\Phi)(PYQ)=\bigl(S_1(Y'\mapsto s\,\Phi(PY'Q))\bigr)(Y)$ for every
  $\Phi$ and $Y$. If $T$ reads as $S$ through $\iota$
  (Definition~\ref{asm-theta-agree-finite-split-models}(3)), then $T$ reads as
  $S_1$ through $\iota_1$.
\end{enumerate}
\end{lemma}
\leanref{Development/WeilBound/CoreStatements/CompactAveragingSplitting/EigenvectorLifting/PfDict.lean}{PositivityTheory.ThK.SplitDict.lin_chg}
\leanref{Development/WeilBound/CoreStatements/CompactAveragingSplitting/EigenvectorLifting/PfDict.lean}{PositivityTheory.ThK.SplitDict.reads_chg}
\begin{proof}
(1) The map $\iota_1$ is bijective between the Schwartz-Bruhat functions, as
$\iota$ is, $Y\mapsto PYQ$ is a homeomorphism and $s\neq0$. Let $u$ be at $v$
and $\phi$, $\phi_1$ Schwartz-Bruhat with $\rho(u)$ acting at $v$ by
$\phi\mapsto\phi_1$. If $\pi^{(0)}(u)=0$, then
$\iota_1\phi_1(Y)=s\,\iota\phi(PYQ+\operatorname{co}(u))=s\,\iota\phi(P(Y+\operatorname{co}_1(u))Q)=\iota_1\phi(Y+\operatorname{co}_1(u))$.
If $\pi^{(1)}(u)=0$ and $x$ is at $v$ with $\pi^{(0)}(x)=0$, then
$\iota_1\phi_1(\operatorname{co}_1(x))=s\,\iota\phi_1(\operatorname{co}(x))=\psi(\langle x,u\rangle)\,\iota_1\phi(\operatorname{co}_1(x))$.

(2) For Schwartz-Bruhat $\phi$, let $\phi_1$ be Schwartz-Bruhat with $T$ acting
at $v$ by $\phi\mapsto\phi_1$ and $\iota\phi_1=S(\iota\phi)$. Then
$\iota_1\phi_1(Y)=s\,(S\iota\phi)(PYQ)=(S_1\iota_1\phi)(Y)$.
\end{proof}

\begin{lemma}\label{asm-theta-agree-dictionary-conjugate}
Let $E$ be a CM field with maximal totally real subfield $F$, complex
conjugation $c$ and quadratic character $\eta$, let $v$ be a finite place of
$F$ and $t\colon E\to F_v$ an $F$-embedding, and let $\chi$ be a Hecke character
of $E$ with $\chi|_{\mathbf A_F^\times}=\eta^m$ for some $m\geq0$. Then
$\chi_{t\circ c}(x)=\chi_t(x)^{-1}$ for every $x\in F_v^\times$
(Definition~\ref{asm-theta-agree-coinvariants-local}(4)).
\end{lemma}
\leanref{Development/WeilBound/CoreStatements/CompactAveragingSplitting/EigenvectorLifting/PfDict.lean}{PositivityTheory.ThK.SplitDict.chrW_conj}
\begin{proof}
By Lemma~\ref{asm-theta-agree-dictionary-split}(2), for $m=1$, there is
$z\in\mathbf A_E^\times$ with $z-1$ at $v$, $\pi_t(z)=x$ and
$\pi_{t\circ c}(z)=1$. Then $c(z)-1$ is at $v$, $\pi_{t\circ c}(c(z))=x$ and
$\pi_t(c(z))=1$, by part (1) of that lemma. Part (3), for $t$ and for $t\circ c$,
gives $\chi(z)=\chi_t(x)$ and $\chi(c(z))=\chi_{t\circ c}(x)$. As
$z\,c(z)=N_{E/F}(z)$ is a norm, $\eta(z\,c(z))=1$, so
$\chi_t(x)\chi_{t\circ c}(x)=\chi(z\,c(z))=1$.
\end{proof}
% Lean: chrW_conj, with h.exists_gl for a 1 x 1 matrix, h.idele for t and t o c, and ResEtaPow,
%   whose subgroup contains the norms.

\begin{lemma}\label{asm-theta-agree-dictionary-vacuum}
Let $B$ be a theta base for a CM field $E$, with character $\psi$, let $X$ be a
Hermitian and $W$ a skew-Hermitian lattice space, of ranks $n$ and $d$, of some
signatures and with lattices $L_X$ and $L_W$, and let $M$ be a model for $(X,W)$.
Let $v$ be a finite place of $F$ such that $\mathrm{vac}_v$ is a vacuum
(Definition~\ref{asm-theta-agree-models}(3)), let $t\colon E\to F_v$ be an
$F$-embedding, and put $\psi_v(s)=\psi(\bar s)$, where $\bar s$ has component
$s$ at $v$ and $0$ elsewhere. Let $b\in GL_n(F_v)$ and $c\in GL_d(F_v)$ be such
that $b\,\mathcal O_v^n$ and $c\,\mathcal O_v^d$ are the completions of $L_X$ and
$L_W$ at the place of $E$ given by $t$, read through $t$, and let
$\delta\in F_v^\times$ be such that, for $x\in F_v$, $\psi_v(xy)=1$ for every
$y\in\mathcal O_v$ if and only if $\delta x\in\mathcal O_v$. Then every linear
model of the first kind for $(t,b,c)$, and every linear model of the second kind
for $(t,\delta b,c)$ (Definition~\ref{asm-theta-agree-models}(5)), carries
$\mathrm{vac}_v$ to a nonzero multiple of the indicator function of
$M_{n\times d}(\mathcal O_v)$.
\end{lemma}
\leanref{Development/WeilBound/CoreStatements/CompactAveragingSplitting/EigenvectorLifting/PfDict.lean}{PositivityTheory.ThK.SplitDict.vac_ind}
\leanref{Development/WeilBound/CoreStatements/CompactAveragingSplitting/EigenvectorLifting/PfDict.lean}{PositivityTheory.ThK.SplitDict.vacuum_image_intertwiner}
\leanref{Development/WeilBound/CoreStatements/CompactAveragingSplitting/EigenvectorLifting/PfDict.lean}{PositivityTheory.ThK.SplitDict.vac_Bar}
\begin{proof}
Let $\Lambda_v$ be the group of Lemma~\ref{asm-theta-agree-evaluation-lattice},
generated by the $a\,xy^{\mathsf T}$ with $a\in\mathbf A$ at $v$,
$a_v\in\mathcal O_v$, $x\in L_X$ and $y\in L_W$, and let
$\mathbf 1_{\mathcal O}$ be the indicator function of
$M_{n\times d}(\mathcal O_v)$. The following facts are direct.
\begin{itemize}
\item[(a)] The elements of $\Lambda_v$ are at $v$.
\item[(b)] If $u$ is at $v$ and either $\pi_t(u)=\pi_t(u')=0$ or
  $\pi_{t\circ c}(u)=\pi_{t\circ c}(u')=0$, then $h(u,u')=0$: it is at $v$, and
  by Lemma~\ref{asm-theta-agree-dictionary-split}(1) its images under $\pi_t$ and
  $\pi_{t\circ c}$ are sums of terms that contain the factors
  $\pi_t(u_{ij})\pi_{t\circ c}(u'_{kl})$, respectively
  $\pi_{t\circ c}(u_{ij})\pi_t(u'_{kl})$.
\item[(c)] There is $\epsilon\in\mathbf A_E$ at $v$ with $\pi_t(\epsilon)=1$,
  $\pi_{t\circ c}(\epsilon)=0$ and $\epsilon\Lambda_v\subset\Lambda_v$: the
  idempotent $(1,0)$ of $E\otimes_FF_v=F_v\times F_v$ is integral, so it is
  $\sum_io_i\otimes a_i$ with $o_i\in\mathcal O_E$ and $a_i\in\mathcal O_v$; put
  $\epsilon=\sum_io_i\bar a_i$, and note
  $\epsilon\,a\,xy^{\mathsf T}=\sum_i(\bar a_ia)(o_ix)y^{\mathsf T}$ with
  $o_ix\in L_X$.
\item[(d)] For $\gamma\in\Lambda_v$, $b^{-1}\pi_t(\gamma)(c^{-1})^{\mathsf T}$ has
  entries in $\mathcal O_v$, as
  $\pi_t(a\,xy^{\mathsf T})=a_v\,t(x)t(y)^{\mathsf T}$ with $t(x)\in b\,\mathcal O_v^n$
  and $t(y)\in c\,\mathcal O_v^d$.
\item[(e)] For $Y\in M_{n\times d}(\mathcal O_v)$ there is $\gamma\in\Lambda_v$
  with $\pi_t(\gamma)=bYc^{\mathsf T}$ and $\pi_{t\circ c}(\gamma)=0$: the matrix
  $bYc^{\mathsf T}$ is a sum of products $y z^{\mathsf T}$ of vectors $y$ in the
  completion of $L_X$ and $z$ in that of $L_W$, so it is $\pi_t(\gamma')$ for
  some $\gamma'\in\Lambda_v$, and $\gamma=\epsilon\gamma'$ works.
\end{itemize}

Let $\iota$ be a linear model for $(\pi^{(1)},\pi^{(0)},\operatorname{co})$, where
either $\iota$ is of the first kind for $(t,b,c)$, with $\pi^{(1)}=\pi_t$,
$\pi^{(0)}=\pi_{t\circ c}$ and $\operatorname{co}(u)=b^{-1}\pi_t(u)(c^{-1})^{\mathsf T}$,
or $\iota$ is of the second kind for $(t,\delta b,c)$, with
$\pi^{(1)}=\pi_{t\circ c}$, $\pi^{(0)}=\pi_t$ and
$\operatorname{co}(u)=\delta\,b^{\mathsf T}\operatorname{co}_0(u)c$, where
$\operatorname{co}_0(u)=-t(J_X)\pi_{t\circ c}(u)t(J_W)^{\mathsf T}$. Write
$\phi=\mathrm{vac}_v$. In both cases:
\begin{itemize}
\item[(i)] every $\gamma\in\Lambda_v$ is $\gamma_0+\gamma_1$ with
  $\gamma_0,\gamma_1\in\Lambda_v$, $\pi^{(1)}(\gamma_0)=0$ and
  $\pi^{(0)}(\gamma_1)=0$, by (c): take $\{\gamma_0,\gamma_1\}=\{\epsilon\gamma,\gamma-\epsilon\gamma\}$;
\item[(ii)] every matrix $Y$ is $\operatorname{co}(x)$ for some $x$ at $v$ with
  $\pi^{(0)}(x)=0$, by Lemma~\ref{asm-theta-agree-dictionary-split}(1), as $b$,
  $c$, $\delta$, $t(J_X)$ and $t(J_W)$ are invertible;
\item[(iii)] for $x$ at $v$ with $\pi^{(0)}(x)=0$: if $x\in\Lambda_v$ holds if and
  only if $\psi(\langle x,u\rangle)=1$ for every $u\in\Lambda_v$ with
  $\pi^{(1)}(u)=0$, then $x\in\Lambda_v$ if and only if $\operatorname{co}(x)$ has
  entries in $\mathcal O_v$.
\end{itemize}
For the first kind, (iii) holds without its hypothesis: if $x\in\Lambda_v$, then
$\operatorname{co}(x)$ is integral by (d); if $\operatorname{co}(x)$ is integral,
then (e) gives $\gamma\in\Lambda_v$ with $\pi_t(\gamma)=\pi_t(x)$ and
$\pi_{t\circ c}(\gamma)=0=\pi_{t\circ c}(x)$, so $x=\gamma$. For the second kind
it suffices to show that $\psi(\langle x,u\rangle)=1$ for every $u\in\Lambda_v$
with $\pi_{t\circ c}(u)=0$ if and only if $\operatorname{co}(x)$ is integral. Put
$C=b^{\mathsf T}\operatorname{co}_0(x)c$, so that $\operatorname{co}(x)=\delta C$.
Lemma~\ref{asm-theta-agree-dictionary-coordinates}(2), for $J_X$, $J_W$ and
$\varepsilon=1$, shows that for $u$ at $v$ with $\pi_t(u)=bZc^{\mathsf T}$,
$\langle x,u\rangle$ is at $v$ with component
$-\operatorname{tr}(\operatorname{co}_0(x)^{\mathsf T}bZc^{\mathsf T})=-\sum_{i,j}C_{ij}Z_{ij}$,
as $\pi_t(x)=0$. If the condition holds, then for all $i,j$ and $y\in\mathcal O_v$,
(e) gives $u\in\Lambda_v$ with $\pi_{t\circ c}(u)=0$ and $\pi_t(u)=bZc^{\mathsf T}$,
where $Z$ has entry $-y$ at $(i,j)$ and $0$ elsewhere; so $\psi_v(C_{ij}y)=1$, and
$\delta C_{ij}\in\mathcal O_v$. Conversely, if $\delta C$ is integral and
$u\in\Lambda_v$ has $\pi_{t\circ c}(u)=0$, then
$Z=b^{-1}\pi_t(u)(c^{-1})^{\mathsf T}$ is integral by (d), and
$\psi(\langle x,u\rangle)=\prod_{i,j}\psi_v(C_{ij}(-Z_{ij}))=1$.

Now let $\theta$ be as in the definition of a vacuum. If $\gamma\in\Lambda_v$
has $\pi^{(1)}(\gamma)=0$ or $\pi^{(0)}(\gamma)=0$, then $h(\gamma,\gamma)=0$ by
(a) and (b), so $\rho(\gamma)$ acts at $v$ by $\phi\mapsto\phi$. Hence, by the
two properties of a linear model, applied with $\phi_1=\phi$:
\begin{itemize}
\item[(k1)] if $x\in\Lambda_v$ and $\pi^{(0)}(x)=0$, then
  $\iota\phi(\operatorname{co}(x))=\iota\phi(0)$;
\item[(k2)] if $x$ is at $v$, $\pi^{(0)}(x)=0$ and
  $\iota\phi(\operatorname{co}(x))\neq0$, then $\psi(\langle x,u\rangle)=1$ for
  every $u\in\Lambda_v$ with $\pi^{(1)}(u)=0$.
\end{itemize}
Moreover:
\begin{itemize}
\item[(k3)] if $x$ is at $v$, $\pi^{(0)}(x)=0$ and $\psi(\langle x,u\rangle)=1$ for
  every $u\in\Lambda_v$ with $\pi^{(1)}(u)=0$, then $x\in\Lambda_v$. Indeed, let
  $\gamma\in\Lambda_v$, written as in (i). Then $\langle\gamma_1,x\rangle=0$ by
  (b), and $\langle\gamma_0,x\rangle=-\langle x,\gamma_0\rangle$, so
  $\psi(\langle\gamma,x\rangle)=1$, and
  Lemma~\ref{asm-theta-agree-evaluation-lattice}(2) gives $x\in\Lambda_v$.
\end{itemize}
As $\phi\neq0$ and $\iota$ is injective on the Schwartz-Bruhat functions, there
is $Y$ with $\iota\phi(Y)\neq0$; by (ii), $Y=\operatorname{co}(x)$ with $x$ at $v$
and $\pi^{(0)}(x)=0$, and (k2), (k3) and (k1) give $x\in\Lambda_v$ and
$\iota\phi(0)=\iota\phi(Y)\neq0$. For every $x$ at $v$ with $\pi^{(0)}(x)=0$,
(k1), (k2) and (k3) show that $x\in\Lambda_v$ if and only if
$\psi(\langle x,u\rangle)=1$ for every $u\in\Lambda_v$ with $\pi^{(1)}(u)=0$, so
by (iii) $x\in\Lambda_v$ if and only if $\operatorname{co}(x)$ is integral. If
$\operatorname{co}(x)$ is integral, then
$\iota\phi(\operatorname{co}(x))=\iota\phi(0)$ by (k1). Otherwise
$x\notin\Lambda_v$, and $\iota\phi(\operatorname{co}(x))=0$, since otherwise (k2)
and (k3) would give $x\in\Lambda_v$. By (ii),
$\iota\phi=\iota\phi(0)\,\mathbf 1_{\mathcal O}$.
\end{proof}
% Lean: vac_ind, applied by vacuum_image_intertwiner (first kind) and vac_Bar (second kind,
%   b' = delta b with the coordinate formula co); (a) is latAtV_at_stable, (b) hTens_zero, (c)
%   lat_idem (splitIdem_mem), (d) lat_to, (e) lat_from, and k3 uses Vac.mem_latAtV. The hypothesis
%   on b and c is IsSplitBasis t L b, with splitLat the image of the completion latLoc under t.

\begin{lemma}\label{asm-theta-agree-dictionary-models}
Let $B$ be a theta base for a CM field $E$, with character $\psi$, let $X$ be a
Hermitian and $W$ a skew-Hermitian lattice space, of ranks $n$ and $d$, of some
signatures and with lattices $L_X$ and $L_W$, let $M$ be a model for $(X,W)$,
let $\chi_X$ and $\chi_W$ be Hecke characters of $E$, and let $\omega$ be a Weil
representation for $M$ with characters $(\chi_X,\chi_W)$
(Definition~\ref{asm-theta-agree-models}(6)). Let $v$ be a finite place of $F$
such that $\mathrm{vac}_v$ is a vacuum, let $t\colon E\to F_v$ be an
$F$-embedding, let $\psi_v$ be as in
Lemma~\ref{asm-theta-agree-dictionary-vacuum}, and let
$\mathbf 1_{\mathcal O}$ be the indicator function of
$M_{n\times d}(\mathcal O_v)$. Let $b\in GL_n(F_v)$ and $c\in GL_d(F_v)$ be such
that $b\,\mathcal O_v^n$ and $c\,\mathcal O_v^d$ are the completions of $L_X$ and
$L_W$ at the place of $E$ given by $t$, read through $t$, and for
$g\in U_X^{(v)}$ and $h\in U_W^{(v)}$ put $\gamma(g)=b^{-1}\pi_t(g)b$ and
$\eta(h)=c^{-1}\pi_t(h)c$.
\begin{enumerate}
\item Put $\nu=(\chi_W)_{t\circ c}$ and $\beta=(\chi_X)_{t\circ c}^{-1}$. There is
  a linear model $\iota$ of the first kind for $(t,b,c)$ with
  $\iota(\mathrm{vac}_v)=\mathbf 1_{\mathcal O}$, through which $\omega(g,1)$
  reads as $\omega_A^{\nu,\beta}(\gamma(g),1)$ and $\omega(1,h)$ reads as
  $\omega_A^{\nu,\beta}(1,\eta(h))$ for all $g\in U_X^{(v)}$ and $h\in U_W^{(v)}$
  (Definition~\ref{asm-theta-agree-finite-split-models}).
\item Let $\delta$ be as in Lemma~\ref{asm-theta-agree-dictionary-vacuum}, and
  put $\nu'=(\chi_W)_t$ and $\alpha'=(\chi_X)_t$. There is a linear model $\iota'$
  of the second kind for $(t,\delta b,c)$ with
  $\iota'(\mathrm{vac}_v)=\mathbf 1_{\mathcal O}$, through which $\omega(g,1)$
  reads as $\omega_B^{\nu',\alpha'}(\gamma(g),1)$ and $\omega(1,h)$ reads as
  $\omega_B^{\nu',\alpha'}(1,\eta(h))$ for all $g\in U_X^{(v)}$ and
  $h\in U_W^{(v)}$.
\item Matrices $b$ and $c$ as above exist, and so does $\delta$. If $\psi_v$ has
  conductor $\mathcal O_v$, then $\delta=1$ has the property of
  Lemma~\ref{asm-theta-agree-dictionary-vacuum}.
\end{enumerate}
\end{lemma}
\leanref{Development/WeilBound/CoreStatements/CompactAveragingSplitting/EigenvectorLifting/PfDict.lean}{PositivityTheory.ThK.SplitDict.model_alt}
\leanref{Development/WeilBound/CoreStatements/CompactAveragingSplitting/EigenvectorLifting/PfDict.lean}{PositivityTheory.ThK.SplitDict.model_Bar}
\leanref{Development/WeilBound/ThetaLifts/LocalTheory/StandardValuations/VacuumImageIntertwiner.lean}{PositivityTheory.ThKG.StdDict.model_alt}
\leanref{Development/WeilBound/ThetaLifts/LocalTheory/StandardValuations/VacuumImageIntertwiner.lean}{PositivityTheory.ThKG.StdDict.model_Bar}
\leanref{Development/WeilBound/CoreStatements/CompactAveragingSplitting/EigenvectorLifting/PfDict.lean}{PositivityTheory.ThK.SplitDict.exists_cond}
\leanref{Development/AbelianVarietyTheory/Definitions/Foundations/GLintEqKOAndPrincEqKOne.lean}{GlobalTheory.AdelicU.exists_isSplitBasis}
\begin{proof}
(1) The embedding $t\circ c$ is an $F$-embedding of $E$ into $F_v$, and
$(t\circ c)\circ c=t$. By Lemma~\ref{asm-theta-agree-coinvariants-dictionary} for
$t\circ c$ there is a linear model $\iota_0$ of the second kind for
$(t\circ c,1,1)$, through which $\omega(g,1)$ and $\omega(1,h)$ read as
$\omega_{\nu,\beta^{-1}}(\pi_{t\circ c}(g),1)$ and
$\omega_{\nu,\beta^{-1}}(1,\pi_{t\circ c}(h))$. It is a linear model for
$(\pi_t,\pi_{t\circ c},\operatorname{co}_0)$ with
$\operatorname{co}_0(u)=-A\,\pi_t(u)D^{\mathsf T}$, where $A=t(c(J_X))$ and
$D=t(c(J_W))$ are invertible. Put $P=-Ab$ and $Q=(Dc)^{\mathsf T}$; then
$P\,b^{-1}\pi_t(u)(c^{-1})^{\mathsf T}Q=\operatorname{co}_0(u)$. By
Lemma~\ref{asm-theta-agree-dictionary-change}(1),
$\iota_1\phi(Y)=\iota_0\phi(PYQ)$ is a linear model of the first kind for
$(t,b,c)$, and by Lemma~\ref{asm-theta-agree-dictionary-vacuum},
$\iota_1(\mathrm{vac}_v)=a\mathbf 1_{\mathcal O}$ with $a\neq0$. So
$\iota\phi(Y)=a^{-1}\iota_0\phi(PYQ)$ is a linear model of the first kind with
$\iota(\mathrm{vac}_v)=\mathbf 1_{\mathcal O}$.

Let $g\in U_X^{(v)}$, $g'=\pi_{t\circ c}(g)$ and $\gamma=\gamma(g)$. Applying
$\pi_{t\circ c}$ to $g^{\mathsf T}J_Xc(g)=J_X$ gives
$g'^{\mathsf T}A\,\pi_t(g)=A$, so $g'^{\mathsf T}P=P\gamma^{-1}$ and
$\det g'=(\det\gamma)^{-1}$. Hence, for every $\Phi$ and $Y$,
\[
a^{-1}\bigl(\omega_{\nu,\beta^{-1}}(g',1)\Phi\bigr)(PYQ)
=a^{-1}\nu(\det g')|\det g'|^{d/2}\Phi(g'^{\mathsf T}PYQ)
=\nu(\det\gamma)^{-1}|\det\gamma|^{-d/2}\,a^{-1}\Phi(P\gamma^{-1}YQ),
\]
which is $\omega_A^{\nu,\beta}(\gamma,1)$ applied to $Y'\mapsto a^{-1}\Phi(PY'Q)$
at $Y$. By Lemma~\ref{asm-theta-agree-dictionary-change}(2), $\omega(g,1)$ reads
as $\omega_A^{\nu,\beta}(\gamma,1)$ through $\iota$. Let $h\in U_W^{(v)}$,
$h'=\pi_{t\circ c}(h)$ and $\eta=\eta(h)$. In the same way
$h'^{\mathsf T}D\,\pi_t(h)=D$, so $Qh'=(\eta^{-1})^{\mathsf T}Q$ and
$\det h'=(\det\eta)^{-1}$, and
\[
a^{-1}\beta^{-1}(\det h')|\det h'|^{n/2}\Phi(PYQh')
=\beta(\det\eta)|\det\eta|^{-n/2}\,a^{-1}\Phi(PY(\eta^{-1})^{\mathsf T}Q),
\]
so $\omega(1,h)$ reads as $\omega_A^{\nu,\beta}(1,\eta)$ through $\iota$.

(2) Put $b'=\delta b$. By Lemma~\ref{asm-theta-agree-coinvariants-dictionary} for
$t$ there is a linear model $\iota_0$ of the second kind for $(t,1,1)$, through
which $\omega(g,1)$ and $\omega(1,h)$ read as
$\omega_{\nu',\alpha'}(\pi_t(g),1)$ and $\omega_{\nu',\alpha'}(1,\pi_t(h))$. Put
$P=(b'^{\mathsf T})^{-1}$ and $Q=c^{-1}$; then
$P\bigl(-b'^{\mathsf T}t(J_X)\pi_{t\circ c}(u)t(J_W)^{\mathsf T}c\bigr)Q=-t(J_X)\pi_{t\circ c}(u)t(J_W)^{\mathsf T}$.
As in (1), Lemmas~\ref{asm-theta-agree-dictionary-change}(1)
and~\ref{asm-theta-agree-dictionary-vacuum} give $a\neq0$ such that
$\iota'\phi(Y)=a^{-1}\iota_0\phi(PYQ)$ is a linear model of the second kind for
$(t,b',c)$ with $\iota'(\mathrm{vac}_v)=\mathbf 1_{\mathcal O}$. For
$g\in U_X^{(v)}$ with $g_0=\pi_t(g)$, $\gamma(g)=b'^{-1}g_0b'$, so
$g_0^{\mathsf T}P=P\gamma(g)^{\mathsf T}$ and $\det g_0=\det\gamma(g)$; for
$h\in U_W^{(v)}$ with $h_0=\pi_t(h)$, $Qh_0=\eta(h)Q$ and $\det h_0=\det\eta(h)$.
So
$a^{-1}(\omega_{\nu',\alpha'}(g_0,1)\Phi)(PYQ)=(\omega_B^{\nu',\alpha'}(\gamma(g),1)\Phi_1)(Y)$
and
$a^{-1}(\omega_{\nu',\alpha'}(1,h_0)\Phi)(PYQ)=(\omega_B^{\nu',\alpha'}(1,\eta(h))\Phi_1)(Y)$,
with $\Phi_1(Y')=a^{-1}\Phi(PY'Q)$, and
Lemma~\ref{asm-theta-agree-dictionary-change}(2) gives the readings.

(3) The completion of the lattice $L_X$ at the place $w$ of $E$ given by $t$ is a
free module of rank $n$ over the ring of integers of $E_w$, which $t$ identifies
with $\mathcal O_v$, and it spans $F_v^n$; a basis gives $b$, and likewise $c$.
As $\psi$ is trivial on $F$ and $\psi(a)=\exp(-2\pi i\operatorname{Tr}(a))$ at
$\infty$, the $x\in F_v$ with $\psi_v(x\mathcal O_v)=1$ form the completion at
$v$ of the inverse different of $F$, a fractional ideal $\delta^{-1}\mathcal O_v$.
If $\psi_v$ has conductor $\mathcal O_v$, this ideal is $\mathcal O_v$.
\end{proof}
% Lean: model_alt and model_Bar (PfDict.lean) apply DictX (here
%   Lemma asm-theta-agree-coinvariants-dictionary) at t o sigma and at t, with grams; the change of
%   coordinates is lin_chg, reads_chg, A_left, A_right, B_left and B_right with gdet_rel and gdet_conj,
%   and the rescaling is rescale with vacuum_image_intertwiner and vac_Bar. (3) is
%   exists_isSplitBasis (through a lattice frame) and exists_cond (mem_adelicIdeal_codiff_iff); a
%   conductor O_v is PsiUnr.

\begin{lemma}\label{asm-theta-agree-finite-inputs-invariant}
Let $k$ be a nonarchimedean local field in which $2\neq0$, let $\psi$ be a
continuous character of $k$ that is not $1$, let $\kappa$ be a finite set, and
let $\rho_\psi$ and $[\ ,\ ]$ be as in
Theorem~\ref{asm-theta-agree-dictionary-stone-von-neumann}. Let
$L\subset k^\kappa\times k^\kappa$ be a $k$-subspace with $[w,w']=0$ for all
$w,w'\in L$. Then there is a linear form $l$ on the functions on $k^\kappa$ such
that $l(\rho_\psi(w)\phi)=l(\phi)$ for every $w\in L$ and every Schwartz-Bruhat
$\phi$, and $l(\phi)\neq0$ for some Schwartz-Bruhat $\phi$.
\end{lemma}
\leanref{Development/AutomorphicSpectrum/Proofs/MembershipProofs/ExistsInvariantForm.lean}{GlobalTheory.CompanionTheory.LagFun.exists_invariant_form}
\leanref{Development/AutomorphicSpectrum/Proofs/MembershipProofs/ExistsInvariantForm.lean}{GlobalTheory.CompanionTheory.LagFun.exists_param}
\begin{proof}
Let $L_1\subset k^\kappa$ be the image of $L$ under $(a,b)\mapsto a$, and $r$ its
dimension. Choose an injective linear map $j\colon k^r\to k^\kappa$ with image
$L_1$, a linear map $p\colon k^\kappa\to k^r$ with $p(j(t))=t$, and a linear map
$\sigma\colon k^r\to k^\kappa$ with $(j(t),\sigma(t))\in L$ for every $t$. Put
$g(t)=\psi(\tfrac12j(t)\cdot\sigma(t))$, a continuous function on $k^r$, and let
$dt$ be a Haar measure on $k^r$. For Schwartz-Bruhat $\phi$ the function
$t\mapsto g(t)\phi(j(t))$ is continuous with compact support, as $j$ is a closed
embedding. The functions $\phi$ for which $t\mapsto g(t)\phi(j(t))$ is
integrable form a subspace, and the integral is linear on it; let $l$ be a
linear form on all functions with
\[
l(\phi)=\int_{k^r}g(t)\phi(j(t))\,dt
\]
on this subspace.

Let $w=(a,b)\in L$, and choose $t_1$ with $j(t_1)=a$. For $t\in k^r$, the
isotropy of $L$ for the pairs $(j(t),\sigma(t))$ and $w$, $(j(t_1),\sigma(t_1))$
and $w$, and $(j(t),\sigma(t))$ and $(j(t_1),\sigma(t_1))$ gives
\[
j(t)\cdot b=a\cdot\sigma(t),\qquad a\cdot b=a\cdot\sigma(t_1),\qquad
j(t)\cdot\sigma(t_1)=a\cdot\sigma(t).
\]
Hence $\tfrac12j(t+t_1)\cdot\sigma(t+t_1)=\tfrac12j(t)\cdot\sigma(t)+j(t)\cdot b+\tfrac12a\cdot b$,
and
\[
g(t)\,(\rho_\psi(w)\phi)(j(t))=g(t)\,\psi\bigl(j(t)\cdot b+\tfrac12a\cdot b\bigr)\phi(j(t)+a)=g(t+t_1)\,\phi(j(t+t_1)).
\]
As $dt$ is invariant under translation, $l(\rho_\psi(w)\phi)=l(\phi)$ for
Schwartz-Bruhat $\phi$.

As $g(p(0))=1$ and $g\circ p$ is continuous, there is a compact open
neighbourhood $U$ of $0$ in $k^\kappa$ on which $g\circ p$ has positive real
part. Its indicator function $\phi$ is Schwartz-Bruhat. If $j(t)\in U$, then
$\operatorname{Re}g(t)=\operatorname{Re}g(p(j(t)))>0$. So
$t\mapsto\operatorname{Re}g(t)\,\phi(j(t))$ is continuous, nonnegative and
positive at $t=0$, and the real part of $l(\phi)$, which is its integral, is
positive.
\end{proof}

\begin{lemma}\label{asm-theta-agree-finite-inputs-criterion}
Let $k$, $\psi$, $\kappa$, $\rho_\psi$ and $[\ ,\ ]$ be as in
Theorem~\ref{asm-theta-agree-dictionary-stone-von-neumann}, let $\phi_0\neq0$ be
a Schwartz-Bruhat function on $k^\kappa$, let $A\subset k^\kappa\times k^\kappa$
be a subset, and let $l$ be a linear form on the functions on $k^\kappa$ with
$l(\rho_\psi(a)\phi)=l(\phi)$ for $a\in A$ and Schwartz-Bruhat $\phi$. Suppose
that every $w\in k^\kappa\times k^\kappa$ is $a+b$ with $a\in A$ such that either
$\rho_\psi(b)\phi_0$ is a multiple of $\phi_0$, or there is $a'\in A$ with
$\rho_\psi(a')\phi_0=\phi_0$ and $\psi([a',b])\neq1$. If $l(\phi_0)=0$, then
$l(\phi)=0$ for every Schwartz-Bruhat $\phi$.
\end{lemma}
\leanref{Development/WeilBound/CoreStatements/LocalThetaKernels/EvlLoc.lean}{PositivityTheory.ThKH.Evl.core}
\leanref{Development/WeilBound/CoreStatements/LocalThetaKernels/EvlLoc.lean}{PositivityTheory.ThKH.Evl.cyc}
\begin{proof}
A direct computation from the formula gives
$\rho_\psi(w)\rho_\psi(w')=\psi(\tfrac12[w,w'])\rho_\psi(w+w')$, and each
$\rho_\psi(w)$ maps Schwartz-Bruhat functions to Schwartz-Bruhat functions. By
Theorem~\ref{asm-theta-agree-dictionary-stone-von-neumann}(1) it suffices to
show $l(\rho_\psi(w)\phi_0)=0$ for every $w$. Write $w=a+b$ as in the
hypothesis. As $a\in A$,
\[
l(\rho_\psi(b)\phi_0)=l(\rho_\psi(a)\rho_\psi(b)\phi_0)=\psi(\tfrac12[a,b])\,l(\rho_\psi(w)\phi_0),
\]
so it suffices to show $l(\rho_\psi(b)\phi_0)=0$. In the first case
$\rho_\psi(b)\phi_0=c\,\phi_0$, and $l(\rho_\psi(b)\phi_0)=c\,l(\phi_0)=0$. In the
second, put $s=[a',b]$, $Y=l(\rho_\psi(b)\phi_0)$ and
$Z=l(\rho_\psi(a'+b)\phi_0)$. The invariance under $\rho_\psi(a')$ gives
$Y=l(\rho_\psi(a')\rho_\psi(b)\phi_0)=\psi(\tfrac12s)Z$, and
$\rho_\psi(a')\phi_0=\phi_0$ gives
$Y=l(\rho_\psi(b)\rho_\psi(a')\phi_0)=\psi(-\tfrac12s)Z$. So
$Y=\psi(s)\psi(-\tfrac12s)Z=\psi(s)Y$, and $Y=0$ as $\psi(s)\neq1$.
\end{proof}
% Lean: core, with schr_schr for the product rule and schr_mem; the reduction to the translates of
%   phi_0 is cyc, which takes the cyclicity of the Schrödinger representation from
%   VacLoc.ModelLine.cyclic_and_line (Theorem asm-theta-agree-dictionary-stone-von-neumann(1)).

\begin{lemma}\label{asm-theta-agree-finite-inputs-uniqueness}
Let $k$, $\psi$, $\kappa$, $\rho_\psi$ and $[\ ,\ ]$ be as in
Theorem~\ref{asm-theta-agree-dictionary-stone-von-neumann}, and let
$L\subset k^\kappa\times k^\kappa$ be a $k$-subspace of dimension $|\kappa|$ with
$[w,w']=0$ for all $w,w'\in L$.
\begin{enumerate}
\item If $z\in k^\kappa\times k^\kappa$ has $[w,z]=0$ for every $w\in L$, then
  $z\in L$.
\item There is a function $\phi_0$ on $k^\kappa$ such that every linear form $l$
  on the functions on $k^\kappa$ with $l(\rho_\psi(w)\phi)=l(\phi)$ for $w\in L$
  and Schwartz-Bruhat $\phi$, and with $l(\phi_0)=0$, vanishes on the
  Schwartz-Bruhat functions.
\item Let $l$ and $l'$ be linear forms on the functions on $k^\kappa$ with
  $l(\rho_\psi(w)\phi)=l(\phi)$ and $l'(\rho_\psi(w)\phi)=l'(\phi)$ for $w\in L$
  and Schwartz-Bruhat $\phi$, and with $l(\phi)\neq0$ for some Schwartz-Bruhat
  $\phi$. Then there is $c\in\mathbf C$ with $l'(\phi)=c\,l(\phi)$ for every
  Schwartz-Bruhat $\phi$.
\end{enumerate}
\end{lemma}
\leanref{Development/WeilBound/CoreStatements/LocalThetaKernels/EvlLoc.lean}{PositivityTheory.ThKH.Evl.mem_of_orth}
\leanref{Development/WeilBound/CoreStatements/LocalThetaKernels/EvlLoc.lean}{PositivityTheory.ThKH.Evl.exists_test}
\leanref{Development/WeilBound/CoreStatements/LocalThetaKernels/EvlLoc.lean}{PositivityTheory.ThKH.Evl.line}
\leanref{Development/WeilBound/CoreStatements/LocalThetaKernels/EvlLoc.lean}{PositivityTheory.ThKH.Evl.transl}
\begin{proof}
(1) The map sending $z$ to the linear form $w\mapsto[w,z]$ is injective, as
$[\ ,\ ]$ is nondegenerate, and by isotropy it maps $L$ into the annihilator of
$L$. Both have dimension $|\kappa|$, so $L$ maps onto the annihilator. If
$[w,z]=0$ for $w\in L$, the form of $z$ lies in the annihilator, so it is the form
of some $z'\in L$, and $z=z'$.

(2) Choose $j\colon k^r\to k^\kappa$, $p$ and $\sigma$ as in the proof of
Lemma~\ref{asm-theta-agree-finite-inputs-invariant}, put
$g(t)=\psi(\tfrac12j(t)\cdot\sigma(t))$, which is locally constant and has no
zero, and let $U\subset k^\kappa$ and $V\subset k^r$ be compact open subgroups.
For a function $f$ on $k^r$ put
\[
\operatorname{ext}(f)(x)=g(p(x))^{-1}f(p(x))\,\mathbf 1_U\bigl(x-j(p(x))\bigr),
\]
and for a function $\phi$ on $k^\kappa$ put $\operatorname{res}(\phi)(t)=g(t)\phi(j(t))$.
Both maps are linear and preserve Schwartz-Bruhat functions, the support of
$\operatorname{ext}(f)$ lying in $j(\operatorname{supp}f)+U$, and
$\operatorname{res}(\operatorname{ext}(f))=f$. Put
$\phi_0=\operatorname{ext}(\mathbf 1_V)$, and let $l$ be as in (2).

First, $l(\phi)=0$ for every Schwartz-Bruhat $\phi$ that vanishes on
$j(k^r)$. If $(0,y)\in L$, then $\rho_\psi(0,y)\phi$ is $\phi$ times the locally
constant function $x\mapsto\psi(x\cdot y)$. Let $x$ have $\phi(x)\neq0$. Then
$x\neq j(p(x))$; choose $i$ with $(x-j(p(x)))_i\neq0$ and $s$ with
$\psi(s)\neq1$, and let $y\in k^\kappa$ be such that
$z\cdot y=s\,(z-j(p(z)))_i/(x-j(p(x)))_i$ for all $z$. Every $w\in L$ is
$(j(t),b)$ for some $t$ and $b$, and $[w,(0,y)]=j(t)\cdot y=0$; so $(0,y)\in L$ by
(1), and $\psi(x\cdot y)=\psi(s)\neq1$. The support of $\phi$ is compact and the
sets on which $x\mapsto\psi(x\cdot y)$ is constant are open and closed, so the
support is covered by disjoint open and closed sets $O_1,\ldots,O_N$ such that
on $O_i$ the function $\psi(\cdot\,y_i)$ is a constant $c_i\neq1$ for some $y_i$
with $(0,y_i)\in L$. Each $\phi\mathbf 1_{O_i}$ is Schwartz-Bruhat, and
$l(\phi\mathbf 1_{O_i})=l(\rho_\psi(0,y_i)(\phi\mathbf 1_{O_i}))=c_i\,l(\phi\mathbf 1_{O_i})$,
so $l(\phi\mathbf 1_{O_i})=0$ and $l(\phi)=0$.

Second, for Schwartz-Bruhat $\phi$, $\phi-\operatorname{ext}(\operatorname{res}(\phi))$
vanishes on $j(k^r)$, as $g$ has no zero; so
$l(\phi)=l(\operatorname{ext}(\operatorname{res}(\phi)))$.

Third, the linear form $l'(f)=l(\operatorname{ext}(f))$ on the functions on
$k^r$ is invariant under translations of Schwartz-Bruhat functions. Let
$t_1\in k^r$ and $w=(j(t_1),\sigma(t_1))\in L$. The identities in the proof of
Lemma~\ref{asm-theta-agree-finite-inputs-invariant}, with $b=\sigma(t_1)$, give
$\operatorname{res}(\rho_\psi(w)\operatorname{ext}(f))(t)=\operatorname{res}(\operatorname{ext}(f))(t+t_1)=f(t+t_1)$.
By the second step and the invariance of $l$,
$l'(f(\cdot+t_1))=l(\operatorname{ext}(\operatorname{res}(\rho_\psi(w)\operatorname{ext}(f))))=l(\rho_\psi(w)\operatorname{ext}(f))=l'(f)$.

Fourth, $l'$ vanishes on the Schwartz-Bruhat functions on $k^r$. Apply
Lemma~\ref{asm-theta-agree-finite-inputs-criterion} on $k^r$ to $\mathbf 1_V$,
the set $A$ of the $(t,0)$, $t\in k^r$, and $l'$, using
$(\rho_\psi(t,0)f)(x)=f(x+t)$ and $l'(\mathbf 1_V)=l(\phi_0)=0$. Let
$w=(w_1,w_2)$, and write it as $(w_1,0)+b$ with $b=(0,w_2)$. If
$\psi(u\cdot w_2)=1$ for every $u\in V$, then $\rho_\psi(b)\mathbf 1_V=\mathbf 1_V$.
Otherwise choose $u\in V$ with $\psi(u\cdot w_2)\neq1$; then $a'=(u,0)$ has
$\rho_\psi(a')\mathbf 1_V=\mathbf 1_V$, as $V$ is a group, and
$[a',b]=u\cdot w_2$.

Finally, for Schwartz-Bruhat $\phi$, $l(\phi)=l'(\operatorname{res}(\phi))=0$.

(3) Let $\phi_0$ be as in (2). Then $l(\phi_0)\neq0$, as $l$ does not vanish on
the Schwartz-Bruhat functions. Put $c=l'(\phi_0)/l(\phi_0)$. The form $l'-c\,l$
is invariant under the $\rho_\psi(w)$, $w\in L$, and vanishes at $\phi_0$, so by
(2) it vanishes on the Schwartz-Bruhat functions.
\end{proof}
% Lean: (1) is mem_of_orth; (2) is exists_test, through exists_param, exists_subgroup, ext_j, ext_mem,
%   res_mem, the vanishing step van (GlobalTheory.WeilM.eq_zero_of_mul_invariant, with mem_of_orth),
%   key, and transl, which applies core to the translations; (3) is line.

\begin{lemma}\label{asm-theta-agree-finite-inputs-averaging}
Let $k$, $\psi$, $\kappa$, $\rho_\psi$ and $[\ ,\ ]$ be as in
Theorem~\ref{asm-theta-agree-dictionary-stone-von-neumann}, with $|\cdot|$ the
normalized absolute value of $k$. Let $L\subset k^\kappa\times k^\kappa$ be a
$k$-subspace of dimension $|\kappa|$ with $[w,w']=0$ for all $w,w'\in L$, and
let $l$ be a linear form on the functions on $k^\kappa$ with
$l(\rho_\psi(w)\phi)=l(\phi)$ for $w\in L$ and Schwartz-Bruhat $\phi$, and with
$l(\phi)\neq0$ for some Schwartz-Bruhat $\phi$. Let $\iota$ be a finite set and
$X\colon k^\iota\to k^\kappa\times k^\kappa$ a linear map with $[X(x),X(x')]=0$
for all $x,x'$, such that $X(x)\in L$ only for $x=0$. Then there is a linear form
$\Lambda$ on the functions on $k^\kappa$ such that:
\begin{enumerate}
\item $\Lambda(\phi)\neq0$ for some Schwartz-Bruhat $\phi$;
\item $\Lambda(\rho_\psi(X(x))\phi)=\Lambda(\phi)$ for $x\in k^\iota$ and
  Schwartz-Bruhat $\phi$;
\item $\Lambda(\rho_\psi(a)\phi)=\Lambda(\phi)$ for Schwartz-Bruhat $\phi$ and
  every $a\in L$ with $[X(x),a]=0$ for all $x\in k^\iota$;
\item if a linear operator $T$ on the functions on $k^\kappa$ maps
  Schwartz-Bruhat functions to Schwartz-Bruhat functions, $g\in GL_\iota(k)$,
  $a\in\mathbf C$, and $T\rho_\psi(X(x))\phi=\rho_\psi(X(gx))T\phi$ and
  $l(T\phi)=a\,l(\phi)$ for $x\in k^\iota$ and Schwartz-Bruhat $\phi$, then
  $\Lambda(T\phi)=|\det g|\,a\,\Lambda(\phi)$ for Schwartz-Bruhat $\phi$.
\end{enumerate}
\end{lemma}
\leanref{Development/WeilBound/CoreStatements/CompactAveragingSplitting/EigenvectorLifting/LinearActInt.lean}{PositivityTheory.ThK.LinAct.exists_avg}
\leanref{Development/WeilBound/CoreStatements/CompactAveragingSplitting/EigenvectorLifting/LinearActInt.lean}{PositivityTheory.ThK.LinAct.exists_section}
\leanref{Development/WeilBound/CoreStatements/CompactAveragingSplitting/EigenvectorLifting/LinearActInt.lean}{PositivityTheory.ThK.LinAct.coeff_mem}
\begin{proof}
By Lemma~\ref{asm-theta-agree-finite-inputs-uniqueness}(1), if $[X(x),a]=0$ for
every $a\in L$, then $X(x)\in L$, so $x=0$. Hence the map from $L$ to the dual of
$k^\iota$ sending $a$ to $x\mapsto[X(x),a]$ is onto, and a linear right inverse
gives a linear map $S\colon k^\iota\to L$ with $[X(x),S(y)]=x\cdot y$ for all
$x,y$. By the product rule in the proof of
Lemma~\ref{asm-theta-agree-finite-inputs-criterion},
$\rho_\psi(w)\rho_\psi(w')=\psi([w,w'])\rho_\psi(w')\rho_\psi(w)$, and
$\rho_\psi(w)\rho_\psi(w')=\rho_\psi(w+w')$ when $[w,w']=0$.

For Schwartz-Bruhat $\phi$ put $F_\phi(x)=l(\rho_\psi(X(x))\phi)$. It is a
Schwartz-Bruhat function on $k^\iota$. As $\phi$ is locally constant with
compact support and $\psi$ is continuous, there is a neighbourhood $N$ of $0$
with $\rho_\psi(w)\phi=\phi$ for $w\in N$. If $X(y-x)\in N$, then
$F_\phi(y)=l(\rho_\psi(X(x))\rho_\psi(X(y-x))\phi)=F_\phi(x)$, so $F_\phi$ is
locally constant. Let $U$ be the neighbourhood of $0$ of the $y$ with
$S(y)\in N$. For $y\in U$, as $S(y)\in L$ and $\rho_\psi(S(y))\phi=\phi$,
\[
F_\phi(x)=l(\rho_\psi(X(x))\rho_\psi(S(y))\phi)=\psi(x\cdot y)\,l(\rho_\psi(S(y))\rho_\psi(X(x))\phi)=\psi(x\cdot y)F_\phi(x).
\]
So $F_\phi(x)\neq0$ implies $\psi(x\cdot y)=1$ for every $y\in U$. Let
$\mathcal O$ be the ring of integers of $k$, and choose $s$ with $\psi(s)\neq1$
and $\varepsilon\in k^\times$ with $\varepsilon\mathcal O^\iota\subset U$. Then
$\psi(x_i\varepsilon\mathcal O)=1$ for every $i$, so $s\notin x_i\varepsilon\mathcal O$
and $|x_i|<|s|/|\varepsilon|$. So $F_\phi$ has compact support.

Let $dx$ be a Haar measure on $k^\iota$, and let $\Lambda$ be a linear form on
the functions on $k^\kappa$ with $\Lambda(\phi)=\int_{k^\iota}F_\phi(x)\,dx$ for
Schwartz-Bruhat $\phi$; this is linear in $\phi$.

(2) $F_{\rho_\psi(X(x_0))\phi}(x)=F_\phi(x+x_0)$, and $dx$ is invariant under
translation.

(3) $F_{\rho_\psi(a)\phi}(x)=\psi([X(x),a])\,l(\rho_\psi(a)\rho_\psi(X(x))\phi)=F_\phi(x)$,
as $a\in L$.

(4) By the hypotheses, $F_{T\phi}(gx)=l(T\rho_\psi(X(x))\phi)=a\,F_\phi(x)$. The
change of variables $x\mapsto gx$ multiplies $dx$ by $|\det g|$, so
$\Lambda(T\phi)=|\det g|\int F_{T\phi}(gx)\,dx=|\det g|\,a\,\Lambda(\phi)$.

(1) Suppose that $\Lambda$ vanishes on the Schwartz-Bruhat functions, and let
$\phi$ be Schwartz-Bruhat with $l(\phi)\neq0$. Then $F_\phi(0)=l(\phi)\neq0$.
For $w=(w_1,w_2)\in k^\iota\times k^\iota$ and the operators $\rho_\psi$ on the
functions on $k^\iota$,
\[
(\rho_\psi(w)F_\phi)(x)=\psi(x\cdot w_2+\tfrac12w_1\cdot w_2)\,l(\rho_\psi(X(x+w_1))\phi)
=\psi(\tfrac12w_1\cdot w_2)\,F_{\phi'}(x),
\]
with $\phi'=\rho_\psi(S(w_2))\rho_\psi(X(w_1))\phi$, since
$l(\rho_\psi(X(x))\rho_\psi(S(w_2))\phi'')=\psi(x\cdot w_2)\,l(\rho_\psi(X(x))\phi'')$
for Schwartz-Bruhat $\phi''$. So a linear form $I$ on the functions on $k^\iota$
equal to $\int\cdot\,dx$ on the Schwartz-Bruhat functions has
$I(\rho_\psi(w)F_\phi)=\psi(\tfrac12w_1\cdot w_2)\Lambda(\phi')=0$ for every $w$.
By Theorem~\ref{asm-theta-agree-dictionary-stone-von-neumann}(1), $I$ vanishes on
the Schwartz-Bruhat functions on $k^\iota$; but its value at the indicator
function of a compact open subgroup is the positive volume of that subgroup, a
contradiction.
\end{proof}
% Lean: exists_avg, with exists_section for S, coeff_mem for F_phi (eventually_fix, exists_compact_ann,
%   schr_comm, schr_add), and cyc for (1); the change of variables is integral_comp_mulVec, and
%   |det g| is distribHaarChar k (det g).

\begin{definition}\label{asm-theta-agree-finite-inputs-local-data}
Let $k$ be a nonarchimedean local field of characteristic $0$ with normalized
absolute value $|\cdot|$, let $\psi$ be a continuous character of $k$ that is not
$1$, and let $K$ be a field containing $k$, with a $k$-automorphism $c$ and an
element $\xi$ such that $c(\xi)=-\xi$ and $1,\xi$ is a $k$-basis of $K$. For a
finite set $\kappa$ let $\rho_\psi$ and $[\ ,\ ]$ be as in
Theorem~\ref{asm-theta-agree-dictionary-stone-von-neumann}. For a square matrix
$J$ over $K$ of size $\gamma$, let $U_J$ be the group of $g\in GL_\gamma(K)$ with
$g^{\mathsf T}Jc(g)=J$; call $J$ \emph{Hermitian} if $c(J)^{\mathsf T}=J$ and
$\det J\neq0$, and \emph{skew-Hermitian} if $c(J)^{\mathsf T}=-J$ and
$\det J\neq0$. Let $\alpha$, $\beta$, $\kappa$ and $\nu$ be finite sets,
$J_A\in M_\alpha(K)$ and $J_B\in M_\beta(K)$, and call a homomorphism from
$K^\times$ to the complex numbers of absolute value $1$ a \emph{unitary
character} of $K^\times$.
\begin{enumerate}
\item A \emph{local Lagrangian} of $J_B$ is a $C\in M_{\beta\times\nu}(K)$ with
  $C^{\mathsf T}J_Bc(C)=0$ for which there is $C'\in M_{\beta\times\nu}(K)$ with
  $C'^{\mathsf T}J_Bc(C')=0$ such that the matrix $(C\ C')$ has a two-sided
  inverse.
\item The space $M_{\alpha\times\beta}(K)$ carries the form
  $\langle x,y\rangle=-\operatorname{Tr}_{K/k}h(x,y)$, where
  $h(x,y)=\sum_{i,j,k',l}x_{ij}(J_A)_{ik'}(J_B)_{jl}c(y_{k'l})$. A \emph{local
  Darboux isomorphism} for $(J_A,J_B)$ is a $k$-linear isomorphism
  $d=(d_1,d_2)\colon M_{\alpha\times\beta}(K)\to k^\kappa\times k^\kappa$ with
  $\langle x,y\rangle=[d(x),d(y)]$ for all $x,y$. Write
  $\rho_d(x)=\rho_\psi(d(x))$.
\item A \emph{local splitting} for $(d,J_B)$ is a homomorphism $\omega$ from
  $U_{J_B}$ to the linear operators on the functions on $k^\kappa$ such that
  every $\omega(h)$ maps Schwartz-Bruhat functions to Schwartz-Bruhat functions
  and $\omega(h)\rho_d(x)\phi=\rho_d(xh^{\mathsf T})\omega(h)\phi$ for
  $x\in M_{\alpha\times\beta}(K)$ and Schwartz-Bruhat $\phi$.
\item Let $C\in M_{\beta\times\nu}(K)$ and let $\chi$ be a unitary character of
  $K^\times$. A local splitting $\omega$ for $(d,J_B)$ satisfies the
  \emph{local Siegel rule} for $(C,\chi)$ if every linear form $\lambda$ on the
  functions on $k^\kappa$ with $\lambda(\rho_d(zC^{\mathsf T})\phi)=\lambda(\phi)$
  for $z\in M_{\alpha\times\nu}(K)$ and Schwartz-Bruhat $\phi$ satisfies
  \[
  \lambda(\omega(p)\phi)=\chi(\det t)\,\bigl|N_{K/k}(\det t)\bigr|^{|\alpha|/2}\lambda(\phi)
  \]
  whenever $\phi$ is Schwartz-Bruhat, $p\in U_{J_B}$, $t\in GL_\nu(K)$ and
  $pC=Ct$.
\item Let $\chi$ be a unitary character of $K^\times$. A \emph{local Weil datum}
  for $(J_A,J_B,d,\chi)$ is a local splitting $\omega$ for $(d,J_B)$ for which
  there are a local Darboux isomorphism $d^\square$ for $(J_A,J_B^\square)$,
  $J_B^\square=\operatorname{diag}(J_B,-J_B)$, from
  $M_{\alpha\times(\beta\sqcup\beta)}(K)$ to
  $k^{\kappa\sqcup\kappa}\times k^{\kappa\sqcup\kappa}$, with
  $d^\square(x_1\ x_2)=\bigl((d_1(x_1),d_1(x_2)),(d_2(x_1),-d_2(x_2))\bigr)$, and a
  local splitting $\omega^\square$ for $(d^\square,J_B^\square)$ that satisfies
  the local Siegel rule for $(\Delta,\chi)$,
  $\Delta=\left(\begin{smallmatrix}1\\1\end{smallmatrix}\right)\in M_{(\beta\sqcup\beta)\times\beta}(K)$,
  and has
  $\omega^\square(\operatorname{diag}(h,1))(\phi_1\boxtimes\phi_2)=(\omega(h)\phi_1)\boxtimes\phi_2$
  for $h\in U_{J_B}$ and Schwartz-Bruhat $\phi_1,\phi_2$, where
  $(\phi_1\boxtimes\phi_2)(x,y)=\phi_1(x)\phi_2(y)$.
\item Let $\chi_A$ and $\chi_B$ be unitary characters of $K^\times$. A
  \emph{local Weil pair} for $(J_A,J_B,d,\chi_A,\chi_B)$ is a homomorphism
  $\omega$ from $U_{J_A}\times U_{J_B}$ to the linear operators on the functions
  on $k^\kappa$ such that $h\mapsto\omega(1,h)$ is a local Weil datum for
  $(J_A,J_B,d,\chi_A)$, and $g\mapsto\omega(g,1)$ is a local Weil datum for
  $(J_B,J_A,d',\chi_B)$ for some local Darboux isomorphism $d'$ for
  $(J_B,J_A)$ with $d'(x^{\mathsf T})=d(x)$ for $x\in M_{\alpha\times\beta}(K)$.
\end{enumerate}
\end{definition}
\leanref{Development/AbelianVarietyTheory/Definitions/Divisors/HilbertNinetyQuotientHomEqIff.lean}{GlobalTheory.WeilM.sbClass}
\leanref{Development/AbelianVarietyTheory/Definitions/Divisors/HilbertNinetyQuotientHomEqIff.lean}{GlobalTheory.WeilM.IsQuad}
\leanref{Development/AbelianVarietyTheory/Definitions/Divisors/HilbertNinetyQuotientHomEqIff.lean}{GlobalTheory.WeilM.UG}
\leanref{Development/AbelianVarietyTheory/Definitions/Divisors/HilbertNinetyQuotientHomEqIff.lean}{GlobalTheory.WeilM.IsHerm}
\leanref{Development/WeilBound/CoreStatements/CompactAveragingSplitting/EigenvectorLifting/LocalDef.lean}{PositivityTheory.ThK.IsSkewM}
\leanref{Development/AbelianVarietyTheory/Definitions/Divisors/HilbertNinetyQuotientHomEqIff.lean}{GlobalTheory.WeilM.IsLagr}
\leanref{Development/AbelianVarietyTheory/Definitions/Divisors/HilbertNinetyQuotientHomEqIff.lean}{GlobalTheory.WeilM.IsSplitting}
\leanref{Development/AbelianVarietyTheory/Definitions/Divisors/HilbertNinetyQuotientHomEqIff.lean}{GlobalTheory.WeilM.SiegelRule}
\leanref{Development/AbelianVarietyTheory/Definitions/Divisors/HilbertNinetyQuotientHomEqIff.lean}{GlobalTheory.WeilM.IsDblCoord}
\leanref{Development/AbelianVarietyTheory/Definitions/Divisors/HilbertNinetyQuotientHomEqIff.lean}{GlobalTheory.WeilM.IsWeilB}
\leanref{Development/AbelianVarietyTheory/Definitions/Divisors/HilbertNinetyQuotientHomEqIff.lean}{GlobalTheory.WeilM.IsWeilPair}

\begin{lemma}\label{asm-theta-agree-finite-inputs-transport}
Let $k$, $\psi$, $K$, $c$ and $|\cdot|$ be as in
Definition~\ref{asm-theta-agree-finite-inputs-local-data}, let $\alpha$, $\beta$,
$\kappa$ and $\nu$ be finite sets, let $J_A\in M_\alpha(K)$ be Hermitian and
$J_B\in M_\beta(K)$ skew-Hermitian, let $d$ be a local Darboux isomorphism for
$(J_A,J_B)$, let $\chi$ be a unitary character of $K^\times$, and let $\omega$ be
a local Weil datum for $(J_A,J_B,d,\chi)$. Then $\omega$ satisfies the local
Siegel rule for $(C,\chi)$ for every local Lagrangian
$C\in M_{\beta\times\nu}(K)$ of $J_B$.
\end{lemma}
\leanref{Development/WeilBound/CoreStatements/CompactAveragingSplitting/EigenvectorLifting/LocalDef.lean}{PositivityTheory.ThK.SiegelTransport}
\leanref{Development/WeilBound/CoreStatements/CompactAveragingSplitting/EigenvectorLifting/SiegelTransport.lean}{PositivityTheory.ThK.siegel_transport}
\leanref{Development/WeilBound/CoreStatements/CompactAveragingSplitting/EigenvectorLifting/SiegelTransportAlg.lean}{PositivityTheory.ThK.SiegelT.exists_gamma}
\leanref{Development/WeilBound/CoreStatements/CompactAveragingSplitting/EigenvectorLifting/SiegelTransportAlg.lean}{PositivityTheory.ThK.SiegelT.rule_of_conj}
\begin{proof}
Let $d^\square$, $\omega^\square$ and $\Delta$ be as in the definition of a local
Weil datum, let $C'$ be as in the definition of a local Lagrangian, and write
the inverse of $(C\ C')$ as $\left(\begin{smallmatrix}m_1\\m_2\end{smallmatrix}\right)$,
so that $Cm_1+C'm_2=1$, $m_1C=1$, $m_1C'=0$, $m_2C=0$ and $m_2C'=1$. The
continuous character $\psi$ takes values of absolute value $1$, as $k$ is a
union of compact subgroups, and $c$ is an involution, as it fixes $k$ and
$c(\xi)=-\xi$.

\emph{Step 1.} A direct computation gives the following. The matrix
$A_0=C^{\mathsf T}J_Bc(C')$ has the inverse $c(m_2)J_B^{-1}m_1^{\mathsf T}$. Put
$f=C'c(A_0)^{-1}$ and $n_2=c(A_0)m_2$. Then $m_1f=0$, $n_2C=0$, $n_2f=1$ and
$Cm_1+fn_2=1$, and, as $c(J_B)^{\mathsf T}=-J_B$,
$C^{\mathsf T}J_Bc(f)=1$, $f^{\mathsf T}J_Bc(C)=-1$, $f^{\mathsf T}J_Bc(f)=0$ and
$J_B=m_1^{\mathsf T}c(n_2)-n_2^{\mathsf T}c(m_1)$. It follows that
\[
\gamma=\begin{pmatrix}1&-fn_2\\-fm_1&fm_1-Cn_2\end{pmatrix}
\]
lies in $U_{J_B^\square}$, that $s=\left(\begin{smallmatrix}m_1\\-n_2\end{smallmatrix}\right)$
has the two-sided inverse $(C\ \ {-f})$, and that
$\gamma\Delta=\operatorname{diag}(C,C)\,s$, where
$\operatorname{diag}(C,C)\in M_{(\beta\sqcup\beta)\times(\nu\sqcup\nu)}(K)$.

\emph{Step 2.} $\omega^\square$ satisfies the local Siegel rule for
$(\operatorname{diag}(C,C),\chi)$. Let $\lambda$ be a linear form with
$\lambda(\rho_{d^\square}(z\operatorname{diag}(C,C)^{\mathsf T})\Phi)=\lambda(\Phi)$
for all $z$ and Schwartz-Bruhat $\Phi$, let $p\in U_{J_B^\square}$ and
$t\in GL_{\nu\sqcup\nu}(K)$ with $p\operatorname{diag}(C,C)=\operatorname{diag}(C,C)t$,
and let $\Phi$ be Schwartz-Bruhat. For $z\in M_{\alpha\times\beta}(K)$,
$z\Delta^{\mathsf T}\gamma^{\mathsf T}=(zs^{\mathsf T})\operatorname{diag}(C,C)^{\mathsf T}$,
so the splitting property shows that $\lambda\circ\omega^\square(\gamma)$ is
invariant under the $\rho_{d^\square}(z\Delta^{\mathsf T})$ on Schwartz-Bruhat
functions. Moreover
$\gamma^{-1}p\gamma\Delta=\gamma^{-1}p\operatorname{diag}(C,C)s=\Delta\,s^{-1}ts$,
and $\det(s^{-1}ts)=\det t$. The local Siegel rule for $(\Delta,\chi)$, applied
to $\lambda\circ\omega^\square(\gamma)$, $\gamma^{-1}p\gamma$, $s^{-1}ts$ and
$\omega^\square(\gamma^{-1})\Phi$, gives
$\lambda(\omega^\square(p)\Phi)=\chi(\det t)|N_{K/k}(\det t)|^{|\alpha|/2}\lambda(\Phi)$.

\emph{Step 3.} Let $\lambda$ be a linear form with
$\lambda(\rho_d(zC^{\mathsf T})\phi)=\lambda(\phi)$ for $z\in M_{\alpha\times\nu}(K)$
and Schwartz-Bruhat $\phi$, let $p\in U_{J_B}$ and $t\in GL_\nu(K)$ with
$pC=Ct$, and let $\phi$ be Schwartz-Bruhat. If $\lambda$ vanishes on the
Schwartz-Bruhat functions, both sides of the rule are $0$, as $\omega(p)\phi$ is
Schwartz-Bruhat. Otherwise choose a Schwartz-Bruhat $\phi_0$ with
$\lambda(\phi_0)\neq0$. For a function $\Phi$ on $k^{\kappa\sqcup\kappa}$ let
$F_\Phi(y)=\lambda(x\mapsto\Phi(x,y))$ and
$\Lambda(\Phi)=\overline{\lambda(\overline{F_\Phi})}$. For Schwartz-Bruhat
$\Phi$ the functions $x\mapsto\Phi(x,y)$ and $F_\Phi$ are Schwartz-Bruhat, and
$\Lambda(\phi_1\boxtimes\phi_2)=\lambda(\phi_1)\overline{\lambda(\overline{\phi_2})}$;
in particular
$\Lambda(\phi_1\boxtimes\overline{\phi_0})=\lambda(\phi_1)\overline{\lambda(\phi_0)}$.

The form $\Lambda$ is invariant under the
$\rho_{d^\square}(z\operatorname{diag}(C,C)^{\mathsf T})$. Write $z=(z_1\ z_2)$
and $d(z_iC^{\mathsf T})=(a_i,b_i)$; then
$d^\square(z\operatorname{diag}(C,C)^{\mathsf T})=((a_1,a_2),(b_1,-b_2))$, and for
Schwartz-Bruhat $\Phi$
\[
\bigl(\rho_{d^\square}(z\operatorname{diag}(C,C)^{\mathsf T})\Phi\bigr)(x,y)
=\psi\bigl(-y\cdot b_2-\tfrac12a_2\cdot b_2\bigr)\bigl(\rho_\psi(a_1,b_1)\Phi(\cdot,y+a_2)\bigr)(x).
\]
Applying $\lambda$ in $x$, which is invariant under $\rho_\psi(a_1,b_1)$, and
using $\overline{\psi(u)}=\psi(-u)$, the function $\overline{F}$ of the left side
is $\rho_\psi(a_2,b_2)\overline{F_\Phi}$, on which $\lambda$ takes the value
$\lambda(\overline{F_\Phi})$.

The matrix $p^\square=\operatorname{diag}(p,1)$ lies in $U_{J_B^\square}$, and
$p^\square\operatorname{diag}(C,C)=\operatorname{diag}(C,C)\operatorname{diag}(t,1)$
with $\det\operatorname{diag}(t,1)=\det t$. By Step 2 and the last property of
$\omega^\square$,
\[
\lambda(\omega(p)\phi)\overline{\lambda(\phi_0)}
=\Lambda\bigl(\omega^\square(p^\square)(\phi\boxtimes\overline{\phi_0})\bigr)
=\chi(\det t)|N_{K/k}(\det t)|^{|\alpha|/2}\lambda(\phi)\overline{\lambda(\phi_0)},
\]
and dividing by $\overline{\lambda(\phi_0)}$ gives the rule.
\end{proof}
% Lean: siegel_transport. Step 1 is exists_gamma (with fromBlocks_transpose_dbl_map_eq_dbl,
%   fromBlocks_mul_diagCols_eq and exists_UG), Step 2 is rule_of_conj, and Step 3 uses
%   Lambda = boxL l (conjL l), boxL_inv, boxL_box, box_mem, star_mem, norm_psi and exists_UG for
%   p^square; the hypothesis that J_A is Hermitian is not used.

\begin{lemma}\label{asm-theta-agree-finite-inputs-fixed}
Let $k$, $\psi$, $K$, $c$ and $|\cdot|$ be as in
Definition~\ref{asm-theta-agree-finite-inputs-local-data}, let $\alpha$, $\beta$,
$\kappa$ and $\nu$ be finite sets, let $J_A\in M_\alpha(K)$ be Hermitian and
$J_B\in M_\beta(K)$ skew-Hermitian, let $C\in M_{\beta\times\nu}(K)$ be a local
Lagrangian of $J_B$, let $d$ be a local Darboux isomorphism for $(J_A,J_B)$, let
$\chi$ be a unitary character of $K^\times$, and let $\omega$ be a local Weil
pair for $(J_A,J_B,d,\chi,1)$. Let $l$ be a linear form on the functions on
$k^\kappa$ with $l(\rho_d(zC^{\mathsf T})\phi)=l(\phi)$ for
$z\in M_{\alpha\times\nu}(K)$ and Schwartz-Bruhat $\phi$, and with
$l(\phi)\neq0$ for some Schwartz-Bruhat $\phi$. Then $l(\omega(g,1)\phi)=l(\phi)$
for every $g\in U_{J_A}$ and every Schwartz-Bruhat $\phi$.
\end{lemma}
\leanref{Development/WeilBound/CoreStatements/CompactAveragingSplitting/EigenvectorLifting/LocalDef.lean}{PositivityTheory.ThK.LinearAct}
\leanref{Development/WeilBound/CoreStatements/CompactAveragingSplitting/EigenvectorLifting/LinearAct.lean}{PositivityTheory.ThK.split_character_fixes_evaluation}
\leanref{Development/WeilBound/CoreStatements/CompactAveragingSplitting/EigenvectorLifting/LinearAct.lean}{PositivityTheory.ThK.LinAct.dbl_fix}
\leanref{Development/WeilBound/CoreStatements/CompactAveragingSplitting/EigenvectorLifting/LinearAct.lean}{PositivityTheory.ThK.LinAct.weilB_fix}
\leanref{Development/WeilBound/CoreStatements/CompactAveragingSplitting/EigenvectorLifting/LinearActAlg.lean}{PositivityTheory.ThK.LinAct.exists_coord}
\begin{proof}
Let $d'$ be as in the definition of a local Weil pair, so that
$g\mapsto\omega(g,1)$ is a local Weil datum for $(J_B,J_A,d',1)$, and let
$d^\square$ and $\omega^\square$ be its data: $d^\square$ is a local Darboux
isomorphism for $(J_B,J_A^\square)$ on $M_{\beta\times(\alpha\sqcup\alpha)}(K)$,
and $\omega^\square$ is a local splitting of $U_{J_A^\square}$ satisfying the
local Siegel rule for $(\Delta,1)$, $\Delta=\left(\begin{smallmatrix}1\\1\end{smallmatrix}\right)\in M_{(\alpha\sqcup\alpha)\times\alpha}(K)$.
Let $C'$ be as in the definition of a local Lagrangian, and write the inverse of
$(C\ C')$ as $\left(\begin{smallmatrix}m_1\\m_2\end{smallmatrix}\right)$.

\emph{Step 0.} For a local Darboux isomorphism $d_1$ for $(J_B,J)$, $J$ square of
size $\gamma$, let $L(d_1)$ be the set of the $d_1(Cz)$,
$z\in M_{\nu\times\gamma}(K)$. It is a $k$-subspace with $[w,w']=0$ on it, as
$h(Cz,Cz')=\sum_{a,b}(C^{\mathsf T}J_Bc(C))_{ab}(zJc(z')^{\mathsf T})_{ab}$ with
$C^{\mathsf T}J_Bc(C)=0$. As $(C\ C')$ is invertible, $|\beta|=2|\nu|$ and
$z\mapsto Cz$ is injective, so $L(d_1)$ has dimension $2|\nu||\gamma|=|\beta||\gamma|$,
which is the number of elements of the index set of $d_1$. By hypothesis $l$ is
invariant under the $\rho_\psi(w)$, $w\in L(d')$, since
$d'(Cz)=d(z^{\mathsf T}C^{\mathsf T})$.

\emph{Step 1.} Let $l_2$ be a linear form on the functions on
$k^{\kappa\sqcup\kappa}$ invariant under the $\rho_\psi(w)$, $w\in L(d^\square)$,
on Schwartz-Bruhat functions, with $l_2(\Phi_0)\neq0$ for some Schwartz-Bruhat
$\Phi_0$. Then $l_2(\omega^\square(q)\Phi)=l_2(\Phi)$ for every
$q\in U_{J_A^\square}$ and Schwartz-Bruhat $\Phi$.

For $q\in U_{J_A^\square}$ the form $l_2\circ\omega^\square(q)$ is again
invariant under $L(d^\square)$, as $Czq^{\mathsf T}=C(zq^{\mathsf T})$, so by
Lemma~\ref{asm-theta-agree-finite-inputs-uniqueness}(3) there is $a(q)$ with
$l_2(\omega^\square(q)\Phi)=a(q)l_2(\Phi)$ for Schwartz-Bruhat $\Phi$. Evaluating
at $\Phi_0$ shows $a(1)=1$ and $a(qq')=a(q)a(q')$, so $a$ is a homomorphism to
$\mathbf C^\times$. The matrix $J_A^\square=\operatorname{diag}(J_A,-J_A)$ is
Hermitian, and $\Delta$ is a local Lagrangian of it, with second matrix
$\left(\begin{smallmatrix}1\\-1\end{smallmatrix}\right)$. By
Lemma~\ref{asm-theta-agree-invariance-character}, for $k$, $K$, $c$, $\xi$,
$J_A^\square$ and these two matrices, it suffices to show $a(p)=1$ for
$p\in U_{J_A^\square}$ and $t\in GL_\alpha(K)$ with $p\Delta=\Delta t$.

Choose a $k$-linear isomorphism $e\colon k^r\to M_{\nu\times\alpha}(K)$. For
$t\in GL_\alpha(K)$ let $g_t\in GL_r(k)$ be the matrix of
$x\mapsto e^{-1}(e(x)t^{\mathsf T})$; a computation of determinants gives
$\det g_t=N_{K/k}(\det t)^{|\nu|}$, as $Y\mapsto Yt^{\mathsf T}$ is $K$-linear
with determinant $(\det t)^{|\nu|}$. Put
$X(x)=d^\square(C'e(x)\Delta^{\mathsf T})$ for $x\in k^r$. Then
$[X(x),X(x')]=0$, as $\Delta^{\mathsf T}J_A^\square c(\Delta)=J_A-J_A=0$ gives
$\langle y\Delta^{\mathsf T},y'\Delta^{\mathsf T}\rangle=0$ for all
$y,y'\in M_{\beta\times\alpha}(K)$. If $X(x)\in L(d^\square)$, say
$C'e(x)\Delta^{\mathsf T}=Cz$, then $(e(x)\ e(x))=e(x)\Delta^{\mathsf T}=m_2Cz=0$,
so $x=0$. Lemma~\ref{asm-theta-agree-finite-inputs-averaging}, for
$L(d^\square)$, $l_2$ and $X$, gives a form $\Lambda$ with its properties (1)
to (4).

$\Lambda$ is invariant under the $\rho_{d^\square}(y\Delta^{\mathsf T})$,
$y\in M_{\beta\times\alpha}(K)$. Indeed $y=Cm_1y+C'm_2y$; choose $x$ with
$e(x)=m_2y$, and put $a'=d^\square(Cm_1y\Delta^{\mathsf T})\in L(d^\square)$. Then
$d^\square(y\Delta^{\mathsf T})=a'+X(x)$, the vectors $a'$ and the $X(x')$ are
pairwise orthogonal for $[\ ,\ ]$, so
$\rho_\psi(a'+X(x))=\rho_\psi(a')\rho_\psi(X(x))$, and properties (3) and (2)
give the invariance. Now let $p$ and $t$ be as above and $\Phi_1$
Schwartz-Bruhat with $\Lambda(\Phi_1)\neq0$. The local Siegel rule for
$(\Delta,1)$ gives
$\Lambda(\omega^\square(p)\Phi_1)=|N_{K/k}(\det t)|^{|\beta|/2}\Lambda(\Phi_1)$.
On the other hand, by the splitting property and
$C'e(x)\Delta^{\mathsf T}p^{\mathsf T}=C'e(x)t^{\mathsf T}\Delta^{\mathsf T}=C'e(g_tx)\Delta^{\mathsf T}$,
the operator $T=\omega^\square(p)$ satisfies the hypotheses of property (4) with
$g=g_t$ and $a=a(p)$, so
$\Lambda(\omega^\square(p)\Phi_1)=|N_{K/k}(\det t)|^{|\nu|}a(p)\Lambda(\Phi_1)$.
As $|\beta|=2|\nu|$, $a(p)=1$.

\emph{Step 2.} By Lemma~\ref{asm-theta-agree-finite-inputs-invariant} there is
a form $l_2$ as in Step 1. There are Schwartz-Bruhat $\phi_1$ and $\phi_2$ on
$k^\kappa$ with $l_2(\phi_1\boxtimes\phi_2)\neq0$: otherwise, for $\phi_0\neq0$
Schwartz-Bruhat, the identity
$\rho_\psi((a_1,a_2),(b_1,b_2))(\phi_0\boxtimes\phi_0)=\rho_\psi(a_1,b_1)\phi_0\boxtimes\rho_\psi(a_2,b_2)\phi_0$
gives $l_2(\rho_\psi(w)(\phi_0\boxtimes\phi_0))=0$ for all $w$, and
Theorem~\ref{asm-theta-agree-dictionary-stone-von-neumann}(1) would make $l_2$
vanish on the Schwartz-Bruhat functions. Put $l'(\phi)=l_2(\phi\boxtimes\phi_2)$.
For $z\in M_{\nu\times\alpha}(K)$ with $d'(Cz)=(a,b)$,
$d^\square(C(z\ 0))=((a,0),(b,0))$ and
$\rho_\psi((a,0),(b,0))(\phi\boxtimes\phi_2)=(\rho_\psi(a,b)\phi)\boxtimes\phi_2$,
so $l'$ is invariant under $L(d')$. By
Lemma~\ref{asm-theta-agree-finite-inputs-uniqueness}(3) there is $a'$ with
$l'(\phi)=a'l(\phi)$ for Schwartz-Bruhat $\phi$, and $a'\neq0$ as
$l'(\phi_1)\neq0$. For $g\in U_{J_A}$ the matrix $\operatorname{diag}(g,1)$ lies
in $U_{J_A^\square}$, and by the last property of $\omega^\square$ and Step 1,
\[
a'l(\omega(g,1)\phi)=l_2\bigl(\omega^\square(\operatorname{diag}(g,1))(\phi\boxtimes\phi_2)\bigr)=l_2(\phi\boxtimes\phi_2)=a'l(\phi).
\]
\end{proof}
% Lean: split_character_fixes_evaluation = weilB_fix applied to the X-side datum (d', with hd'). Step 0
%   is lagY, lagY_iso (hTens_transpose, hTens_eq_zero), lagY_dim, card_lagr and indep; Step 1 is
%   dbl_fix, with line, exists_coord, isLagr_diagCols, isHerm_dbl, exists_avg and
%   GlobalTheory.T.lagrangian_stabilizer_character_trivial; Step 2 is weilB_fix, with
%   exists_invariant_form, cyc, schr_box, line and WeilOneMember.exists_embed.

\begin{lemma}\label{asm-theta-agree-finite-inputs-local-pair}
Let $(X,W,M,\chi_X,\chi_W,\omega)$ be a Weil pair for a theta base $B$ of a CM
field $E$ (Definition~\ref{asm-theta-agree-triple}(1)), with $X$ and $W$ of ranks
$n$ and $d$, matrices $J_X$ and $J_W$, isomorphism $e$ and Schrödinger operators
$\rho$, and put $I=\{1,\ldots,n\}\times\{1,\ldots,d\}$. Let $v$ be a finite place
of $F$ inert in $E$ and $w$ the place of $E$ above it; write $u_w$ for the
component at $w$ of a matrix $u$ over $\mathbf A_E$, and $c$ also for the
nontrivial automorphism of $E_w$ over $F_v$. Put $\psi_v(s)=\psi(\bar s)$, where
$\bar s\in\mathbf A$ has component $s$ at $v$ and $0$ elsewhere, and for a Hecke
character $\chi$ of $E$ let $\chi_w(y)=\chi(z_y)$ for $y\in E_w^\times$, where
$z_y\in\mathbf A_E^\times$ has component $y$ at $w$ and $1$ elsewhere. Then $F_v$,
$\psi_v$, $E_w$ and $c$ satisfy the conditions on $k$, $\psi$, $K$ and $c$ in
Definition~\ref{asm-theta-agree-finite-inputs-local-data}, and there are a local
Darboux isomorphism $e_v$ for $(J_X,J_W)$, from $M_{n\times d}(E_w)$ to
$F_v^I\times F_v^I$, and a local Weil pair $\omega_w$ for
$(J_X,J_W,e_v,\chi_{X,w},\chi_{W,w})$, such that:
\begin{enumerate}
\item $e_v(u_w)=e(u)_v$ for $u\in M_{n\times d}(\mathbf A_E)$, and for $u$ at $v$
  the operator $\rho(u)$ acts at $v$ (Definition~\ref{asm-theta-agree-models}(2))
  by $\phi\mapsto\rho_{\psi_v}(e_v(u_w))\phi$ for every Schwartz-Bruhat $\phi$;
\item $g\mapsto g_w$ maps $U_X^{(v)}$ onto $U_{J_X}$ and $U_W^{(v)}$ onto
  $U_{J_W}$, the unitary groups over $E_w$;
\item for $g\in U_X^{(v)}$, $h\in U_W^{(v)}$ and Schwartz-Bruhat $\phi$,
  $\omega(g,h)$ acts at $v$ by $\phi\mapsto\omega_w(g_w,h_w)\phi$.
\end{enumerate}
\end{lemma}
\leanref{Development/WeilBound/CoreStatements/CompactAveragingSplitting/EigenvectorLifting/LocalDef.lean}{PositivityTheory.ThK.IsLocChar}
\leanref{Development/WeilBound/CoreStatements/CompactAveragingSplitting/EigenvectorLifting/LocalDef.lean}{PositivityTheory.ThK.IsLocPair}
\leanref{Development/WeilBound/CoreStatements/CompactAveragingSplitting/EigenvectorLifting/LocalDef.lean}{PositivityTheory.ThK.LocPairEx}
\leanref{Development/WeilBound/CoreStatements/LocalThetaKernels/Locpair.lean}{PositivityTheory.ThKH.Locpair.locPairEx}
\leanref{Development/WeilBound/CoreStatements/LocalThetaKernels/Locpair.lean}{PositivityTheory.ThKH.Locpair.locPair_core}
\begin{proof}
As $v$ is inert, $E\otimes_FF_v$ is the field $E_w$, of degree $2$ over $F_v$,
and the elements of $\mathbf A_E$ at $v$ are the elements of $E_w$ placed at
$v$. For $y\in E_w$ let $\tilde y\in\mathbf A_E$ be the element at $v$ with
component $y$, apply this entrywise to matrices, and for a square matrix $m$
over $E_w$ put $\hat m=1+\widetilde{(m-1)}$. Then $u\mapsto u_w$ is a ring
homomorphism commuting with $c$, $\widetilde{u_w}=\bar1u$, $m\mapsto\hat m$ is
multiplicative, $\hat m_w=m$, a matrix $g$ with $g-1$ at $v$ is
$\widehat{g_w}$, $\det\hat m=z_{\det m}$, and the idele $z_y$ has norm with
component $N_{E_w/F_v}(y)$ at $v$ and $1$ elsewhere.

\emph{The setting.} $F_v$ has characteristic $0$. The character $\psi_v$ is
continuous, and it is not $1$, as $\psi_v$ has conductor $\mathcal O_v$ at the
inert place $v$ (Definition~\ref{gal-good-base}). As $c\neq1$ on $E$ and $2\neq0$,
$E=F+F\xi$ for some $\xi$ with $c(\xi)=-\xi$, and $1,\xi$ is an $F_v$-basis of
$E_w$. Each $\chi_w$ is a unitary character of $E_w^\times$, as
$z_{yy'}=z_yz_{y'}$.

\emph{Unitary groups.} Let $J$ be a square matrix over $E$. If $g\in U_J$ is at
$v$, then $g_w^{\mathsf T}Jc(g_w)=J$; if $m$ lies in the unitary group of $J$
over $E_w$, then $\hat m^{\mathsf T}Jc(\hat m)=J$, as both sides have component
$J$ at $w$ and elsewhere. Since $g=\widehat{g_w}$, $g\mapsto g_w$ is a bijection
from $U_J^{(v)}$ onto the unitary group of $J$ over $E_w$, with inverse
$m\mapsto\hat m$. This gives (2).

\emph{The Darboux isomorphism.} Put $e_v(x)=e(\tilde x)_v$ for
$x\in M_{n\times d}(E_w)$. It is $F_v$-linear, as $\widetilde{sx}=\bar s\tilde x$
for $s\in F_v$. It is injective, as $e(\tilde x)$ is at $v$. It is onto: for
$(a,b)$, the element $u=e^{-1}(\bar a,\bar b)$ is at $v$, and
$e_v(u_w)=e(\bar1u)_v=(a,b)$. The same computation gives $e_v(u_w)=e(u)_v$ for
every $u$. As $J_X$ and $J_W$ have entries in $E$, $h(\tilde x,\tilde y)$ is
the element at $v$ with component $h(x,y)$, and
$\operatorname{Tr}_{\mathbf A_E/\mathbf A}$ of an element at $v$ is the element
at $v$ with component its trace from $E_w$ to $F_v$. So the component at $v$ of
$\langle\tilde x,\tilde y\rangle$ is the local form $\langle x,y\rangle$, and by
the property of $e$ it is also $[e_v(x),e_v(y)]$. So $e_v$ is a local Darboux
isomorphism for $(J_X,J_W)$, and $e_v^{\mathsf T}(y)=e_v(y^{\mathsf T})$ is one
for $(J_W,J_X)$. For $u$ at $v$, $e(u)$ is at $v$, as $e$ is $\mathbf A$-linear,
and the formula for $\rho$ gives
$\rho(u)(\phi\otimes r)=(\rho_{\psi_v}(e(u)_v)\phi)\otimes r$ for every $r$ that
does not depend on the $v$-component; this gives (1).

\emph{The operators.} Let $x_0$ and $r_0$ be as in
Lemma~\ref{asm-theta-agree-finite-inputs-slice}(4) for $I$. Let $m$ and $m'$ lie
in the unitary groups of $J_X$ and $J_W$ over $E_w$. The operator
$T=\omega(\hat m,\hat m')$ maps Schwartz-Bruhat functions to Schwartz-Bruhat
functions, keeps admissible forms admissible, and has
$T\rho(u)\Phi=\rho(\hat mu\hat m'^{\mathsf T})T\Phi$ for Schwartz-Bruhat $\Phi$, by
conditions (a) and (b) of Definition~\ref{asm-theta-agree-models}(6). If $e(u)$
has $v$-components $0$, then so does $u$, and $\hat mu\hat m'^{\mathsf T}=u$. By
Lemma~\ref{asm-theta-agree-finite-inputs-schur}(2), $T$ acts at $v$ by
$\phi\mapsto T_v(m,m')\phi=\operatorname{res}_{x_0}(T(\phi\otimes r_0))$, which is
Schwartz-Bruhat by Lemma~\ref{asm-theta-agree-finite-inputs-slice}(3). As $\omega$
is a representation and $m\mapsto\hat m$ is multiplicative,
Lemma~\ref{asm-theta-agree-finite-unique} gives $T_v(1,1)\phi=\phi$ and
$T_v(mn,m'n')\phi=T_v(m,m')T_v(n,n')\phi$ for Schwartz-Bruhat $\phi$. Choose a
linear projection $P$ of the functions on $F_v^I$ onto the Schwartz-Bruhat
functions, and put $\omega_w(m,m')=T_v(m,m')P+(1-P)$. This is a homomorphism,
equal to $T_v(m,m')$ on Schwartz-Bruhat functions, and (3) holds as
$g=\widehat{g_w}$ and $h=\widehat{h_w}$.

\emph{The splittings.} Let $m'$ be in the unitary group of $J_W$ over $E_w$,
$x\in M_{n\times d}(E_w)$ and $\phi$ Schwartz-Bruhat. The matrix
$\tilde x\hat m'^{\mathsf T}$ is at $v$ with component $xm'^{\mathsf T}$. By (1)
and (3), the two sides of
$\omega(1,\hat m')\rho(\tilde x)(\phi\otimes r_0)=\rho(\tilde x\hat m'^{\mathsf T})\omega(1,\hat m')(\phi\otimes r_0)$
are $(\omega_w(1,m')\rho_{e_v}(x)\phi)\otimes r_0$ and
$(\rho_{e_v}(xm'^{\mathsf T})\omega_w(1,m')\phi)\otimes r_0$, where
$\rho_{e_v}(x)=\rho_{\psi_v}(e_v(x))$, and evaluating at the points
$x_0\langle y\rangle$, where $r_0=1$, gives
$\omega_w(1,m')\rho_{e_v}(x)\phi=\rho_{e_v}(xm'^{\mathsf T})\omega_w(1,m')\phi$. So
$m'\mapsto\omega_w(1,m')$ is a local splitting for $(e_v,J_W)$, and in the same
way $m\mapsto\omega_w(m,1)$ is one for $(e_v^{\mathsf T},J_X)$.

\emph{The doubled data.} By Definition~\ref{asm-theta-agree-models}(6)(a) there
are an isomorphism $e^\square$ for $(J_X,J_W^\square)$ and a representation
$\omega^\square$ of the doubled group $G_W$ satisfying condition (b) of
Definition~\ref{asm-theta-agree-doubling}(6) with the character $\chi_X$ and
$m=n$, with
$\omega^\square(\operatorname{diag}(h,1))(\phi_1\boxtimes\phi_2)=(\omega(1,h)\phi_1)\boxtimes\phi_2$.
The constructions above, for $e^\square$ and $\omega^\square$, give a local
Darboux isomorphism $e^\square_v$ for $(J_X,J_W^\square)$, which by the formula
for $e^\square$ satisfies
$e^\square_v(x_1\ x_2)=((e_v(x_1)_1,e_v(x_2)_1),(e_v(x_1)_2,-e_v(x_2)_2))$, and a
local splitting $\omega^\square_w$ for $(e^\square_v,J_W^\square)$ such that
$\omega^\square(\hat q)$ acts at $v$ by $\Phi\mapsto\omega^\square_w(q)\Phi$. For
$m'$ as above, $\widehat{\operatorname{diag}(m',1)}=\operatorname{diag}(\hat m',1)$
and
$(\phi_1\boxtimes\phi_2)\otimes(r_0\boxtimes r_0)=(\phi_1\otimes r_0)\boxtimes(\phi_2\otimes r_0)$,
so the property of $\omega^\square$ and (3) give
$\bigl(\omega^\square_w(\operatorname{diag}(m',1))(\phi_1\boxtimes\phi_2)\bigr)\otimes(r_0\boxtimes r_0)
=\bigl((\omega_w(1,m')\phi_1)\boxtimes\phi_2\bigr)\otimes(r_0\boxtimes r_0)$, and
evaluating at the points $(x_0,x_0)\langle y\rangle$ gives
$\omega^\square_w(\operatorname{diag}(m',1))(\phi_1\boxtimes\phi_2)=(\omega_w(1,m')\phi_1)\boxtimes\phi_2$.

\emph{The local Siegel rule.} Let $\lambda$ be a linear form on the functions on
$F_v^{I\sqcup I}$ with $\lambda(\rho_{e^\square_v}(z\Delta^{\mathsf T})\Phi)=\lambda(\Phi)$
for $z\in M_{n\times d}(E_w)$ and Schwartz-Bruhat $\Phi$, let $p$ be in the
unitary group of $J_W^\square$ over $E_w$ and $t\in GL_d(E_w)$ with $p\Delta=\Delta t$,
and let $\Phi$ be Schwartz-Bruhat. Let
$\ell(\Xi)=\int_{\mathbf A^I}\Xi(x,x)\,dx$ for a Haar measure on $\mathbf A^I$.
For $\Xi=f\otimes g$ it is the product of $\int f(x,x)\,dx$, which is
continuous in $f$, and $\int g(x,x)\,dx$, so $\ell$ is admissible. The function
$\rho((x,x),(y,-y))\Xi$ takes the value $\Xi(x'+x,x'+x)$ at $(x',x')$, so $\ell$
is invariant under these operators, which are the
$\rho(e^\square(u\Delta^{\mathsf T}))$, $u\in M_{n\times d}(\mathbf A_E)$; and
$\ell(\Psi_1\boxtimes\Psi_1)>0$ for a nonzero Schwartz-Bruhat $\Psi_1\geq0$
(compare Lemma~\ref{asm-theta-agree-dictionary-diagonal}). By
Lemma~\ref{asm-theta-agree-finite-inputs-slice}(2) there are a Schwartz-Bruhat
$\phi_0$ on $F_v^{I\sqcup I}$, $y_0\in F_v^{I\sqcup I}$ and a Schwartz-Bruhat
$\Psi$ with $c_0=\ell(P_{\phi_0,y_0}\Psi)\neq0$. For a function $\Psi'$ on
$\mathbf A^{I\sqcup I}$ put $s(\Psi')(y)=\ell(P_{\phi_0,y}\Psi')$, and let
$l=\lambda\circ s$. As $P_{\phi_0,y}P_{\phi_1,y_1}=\phi_1(y)P_{\phi_0,y_1}$,
$s(P_{\phi_1,y_1}\Psi')=\ell(P_{\phi_0,y_1}\Psi')\,\phi_1$. So $s$ maps
Schwartz-Bruhat functions to Schwartz-Bruhat functions,
$l\circ P_{\phi_1,y_1}=\lambda(\phi_1)\,\ell\circ P_{\phi_0,y_1}$, and $l$ is
admissible by Lemma~\ref{asm-theta-agree-finite-inputs-slice}(1) and (2). Let
$u\in M_{n\times d}(\mathbf A_E)$, $a=e^\square(u\Delta^{\mathsf T})$, and $a^\circ$
the vector $a$ with its $v$-components replaced by $0$, which is
$e^\square((1-\bar1)u\Delta^{\mathsf T})$. The computation in the proof of
Lemma~\ref{asm-theta-agree-finite-inputs-schur}(2) gives
$P_{\phi_0,y}\rho(a)\Psi'=\psi_v(y\cdot a_{2,v}+\tfrac12a_{1,v}\cdot a_{2,v})\,\rho(a^\circ)P_{\phi_0,y+a_{1,v}}\Psi'$;
applying $\ell$, which is invariant under $\rho(a^\circ)$, gives
$s(\rho(a)\Psi')=\rho_{\psi_v}(a_v)s(\Psi')$, with
$a_v=e^\square_v(u_w\Delta^{\mathsf T})$. So $l$ is invariant under the
$\rho(e^\square(u\Delta^{\mathsf T}))$ on Schwartz-Bruhat functions. Finally, for
$r(x)=\Psi(x\langle y_0\rangle)$, which does not depend on the $v$-component,
$\Phi\otimes r=P_{\Phi,y_0}\Psi$ is Schwartz-Bruhat and
$l(\Phi'\otimes r)=c_0\lambda(\Phi')$ for every Schwartz-Bruhat $\Phi'$.

The element $\hat p$ of $G_W$ satisfies $\hat p\Delta=\Delta\hat t$, as both sides
have component $p\Delta=\Delta t$ at $w$ and $\Delta$ elsewhere. By the rule of
Definition~\ref{asm-theta-agree-doubling}(6)(b) for $l$,
$l(\omega^\square(\hat p)(\Phi\otimes r))=\chi_X(\det\hat t)|\det\hat t|^{n/2}l(\Phi\otimes r)$.
Here $\omega^\square(\hat p)(\Phi\otimes r)=(\omega^\square_w(p)\Phi)\otimes r$,
$\det\hat t=z_{\det t}$, so $\chi_X(\det\hat t)=\chi_{X,w}(\det t)$ and
$|\det\hat t|=|N_{E_w/F_v}(\det t)|$. Dividing by $c_0$ gives the local Siegel
rule for $(\Delta,\chi_{X,w})$. So $m'\mapsto\omega_w(1,m')$ is a local Weil
datum for $(J_X,J_W,e_v,\chi_{X,w})$, and in the same way, from
Definition~\ref{asm-theta-agree-models}(6)(b), $m\mapsto\omega_w(m,1)$ is one for
$(J_W,J_X,e_v^{\mathsf T},\chi_{W,w})$.
\end{proof}
% Lean: locPairEx = locPair_core. The lifts are tensAt, liftM, liftU and exists_liftU (LocpairAlg.lean,
%   with extE, extM, comp_liftM, aw_liftM, det_liftM and norm_liftE); the setting is PT.psi_good,
%   PT.isField_Ev and PT.isQuad_Ev; chi_w is locChr (isLocChar). e_v is exists_locD (locW, ev_symp,
%   trace_map'); (1) is acts_rho (acts_schr, SchrLoc.schr_tens_of_at). The operators are exists_locOm and
%   exists_locOm' (Locrep.acts_res, eq_of_acts, exists_base, res_mem, exists_ext), with lift_mul_away and
%   mul_lift_away; the splittings are split_loc (tens_inj); the doubled data are dbl_loc, with comp_dbl and
%   the box property through tens_inj; the local Siegel rule is siegel_loc, through exists_form
%   (Pm_Pm, schA_le, PA_adm, admissible_of_integral_conj_le, PA_schr, arg_eq, putA_zero) and
%   KCo.adjoint_statement's field diag for ell, with modu_single for |det t^|; weilB_loc assembles the
%   two data.

\begin{lemma}\label{asm-theta-agree-finite-inputs}
Let $(B,G,\mathcal X,\omega,\omega')$ be a triple of ranks $(a,a')$ for a theta
base $B$ of a CM field $E$ (Definition~\ref{asm-theta-agree-triple}), with models
$M$ and $M'$, put $d=a+a'$, and let $v$ be a finite place of $F$.
\begin{enumerate}
\item Let $t\colon E\to F_v$ be an $F$-embedding, let $\alpha_t$, $\beta_t$ and
  $\nu_t$ be as in Lemma~\ref{asm-theta-agree-finite-characters}, and put
  $\omega_A=\omega_A^{\nu_t,\beta_t}$ and $\omega_B=\omega_B^{\nu_t,\alpha_t}$
  (Definition~\ref{asm-theta-agree-finite-split-models}). Then there are
  $b\in GL_a(F_v)$, $b'\in GL_{a'}(F_v)$, $c\in GL_d(F_v)$, a linear model
  $\iota$ of the first kind for $(t,b,c)$ for $M$, a linear model $\iota'$ of the
  second kind for $(t,b',c)$ for $M'$ (Definition~\ref{asm-theta-agree-models}(5)),
  and maps $\gamma\colon U_V^{(v)}\to GL_a(F_v)$,
  $\gamma'\colon U_{V'}^{(v)}\to GL_{a'}(F_v)$ and
  $\eta\colon U_W^{(v)}\to GL_d(F_v)$ with $\eta$ onto, such that: if $v\in P$
  and $t=\tau_v$, then $b=b_{V,v}$, $b'=b_{V',v}$ and $c=b_{W,v}$; $\iota$ and
  $\iota'$ carry $\mathrm{vac}_v$ and $\mathrm{vac}'_v$ to the indicator
  functions of the matrices with entries in $\mathcal O_v$; and for
  $g\in U_V^{(v)}$, $g'\in U_{V'}^{(v)}$ and $h\in U_W^{(v)}$ the operators
  $\omega(g,1)$ and $\omega(1,h)$ read as $\omega_A(\gamma(g),1)$ and
  $\omega_A(1,\eta(h))$ through $\iota$, and $\omega'(g',1)$ and $\omega'(1,h)$
  read as $\omega_B(\gamma'(g'),1)$ and $\omega_B(1,\eta(h))$ through $\iota'$.
\item Each Weil pair for $B$ (Definition~\ref{asm-theta-agree-triple}(1)) has
  local operators $(\omega_v,\rho_v)$ at $v$
  (Definition~\ref{asm-theta-agree-finite-local-objects}(4)), and for any such
  operators and Schwartz-Bruhat $\phi$ there is a neighbourhood of $1$ in $U_X$
  whose elements $g$ at $v$ have $\omega_v(g,1)\phi=\phi$.
\end{enumerate}
Suppose now that $v$ is inert in $E$, with $w$ the place of $E$ above it, and
write $t_w$ for the component at $w$ of a matrix $t$ over $\mathbf A_E$ and
$|\cdot|$ for the normalized absolute value of $F_v$.
\begin{enumerate}
\setcounter{enumi}{2}
\item The stabilizers $K_V(v)$, $K_{V'}(v)$ and $K_W(v)$
  (Definition~\ref{asm-theta-agree-finite-local-objects}(1)) are compact, and
  $K_W(v)\cap hK_W(v)h^{-1}$ has finite index in $K_W(v)$ for every
  $h\in U_W^{(v)}$.
\item Every $h\in U_W^{(v)}$ has $h^{-1}\in K_W(v)\,h\,K_W(v)$.
\item Write $d=2m$. There is an adapted Lagrangian $C$ of $W$ at $v$ such that
  every $h\in U_W^{(v)}$ is $pk$ with $p\in P_C$ and $k\in K_W(v)$.
\item Let $(X,W,M_0,\chi_X,\chi_W,\omega_0)$ be one of the two Weil pairs of the
  triple, with $X$ of rank $n$, let $(\omega_v,\rho_v)$ be local operators at
  $v$ for it, and let $C$ be an adapted Lagrangian of $W$ at $v$. Then there is an
  evaluation form $l_0$ along $C$ with $l_0(\mathrm{vac}_v)\neq0$, where
  $\mathrm{vac}_v$ is the vacuum of $M_0$. Let $l_0$ be any evaluation form along
  $C$. If $l_0$ is not $0$ on the Schwartz-Bruhat functions, and $\chi_W(z)=1$
  for every $z\in\mathbf A_E^\times$ with $z-1$ at $v$, then
  $l_0(\omega_v(g,1)\phi)=l_0(\phi)$ for $g\in U_X^{(v)}$ and Schwartz-Bruhat
  $\phi$. If $p\in U_W^{(v)}$ and $t\in M_m(\mathbf A_E)$ have $pC=Ct$, and
  $z\in\mathbf A_E^\times$ has $z-1$ at $v$ and $z_w=\det t_w$, then
  \[
  l_0(\omega_v(1,p)\phi)=\chi_X(z)\,\bigl|N_{E_w/F_v}(\det t_w)\bigr|^{n/2}\,l_0(\phi)
  \]
  for Schwartz-Bruhat $\phi$.
\item Let $C$ be an adapted Lagrangian of $W$ at $v$, put
  $H=P_C^\times\cap K_W(v)$, let $p\in P_C^\times$ and $t\in M_m(\mathbf A_E)$
  with $pC=Ct$, and let $[H:H_1]$ denote the index of $H\cap H_1$ in $H$, taken to
  be $0$ if it is infinite. Then
  \[
  [H:p^{-1}Hp]=\bigl|N_{E_w/F_v}(\det t_w)\bigr|^m\,[H:pHp^{-1}].
  \]
\end{enumerate}
\end{lemma}
\leanref{Development/WeilBound/CoreStatements/CompactAveragingSplitting/EigenvectorLifting/SplitDict.lean}{PositivityTheory.ThK.SplitDict.DictAt}
\leanref{Development/WeilBound/CoreStatements/CompactAveragingSplitting/EigenvectorLifting/SplitDict.lean}{PositivityTheory.ThK.SplitDict.Dict}
\leanref{Development/WeilBound/CoreStatements/CompactAveragingSplitting/EigenvectorLifting/PfDict.lean}{PositivityTheory.ThK.SplitDict.split_place_dictionary}
\leanref{Development/WeilBound/CoreStatements/CompactAveragingSplitting/EigenvectorLifting/InertDef.lean}{PositivityTheory.ThK.LocRepEx}
\leanref{Development/WeilBound/CoreStatements/LocalThetaKernels/Locrep.lean}{PositivityTheory.ThKH.Locrep.locRepEx}
\leanref{Development/WeilBound/ThetaLifts/SphericalLattices/ExistsFixedSum.lean}{PositivityTheory.Sph.IsLocRep.smoothX}
\leanref{Development/WeilBound/CoreStatements/CompactAveragingSplitting/EigenvectorLifting/Kcompact.lean}{PositivityTheory.ThK.stabiliser_compact}
\leanref{Development/WeilBound/ThetaLifts/SphericalLattices/FunDefs.lean}{PositivityTheory.Sph.FinSt}
\leanref{Development/WeilBound/ThetaLifts/SphericalLattices/FunFin.lean}{PositivityTheory.Sph.finSt}
\leanref{Development/WeilBound/ThetaLifts/SphericalLattices/FunDefs.lean}{PositivityTheory.Sph.CartanSt}
\leanref{Development/WeilBound/ThetaLifts/SphericalLattices/InvMemKgKCart2.lean}{PositivityTheory.Sph.cartanSt2}
\leanref{Development/WeilBound/CoreStatements/LocalThetaKernels/IwasMain.lean}{PositivityTheory.ThKH.Iwas.iwasawa}
\leanref{Development/WeilBound/CoreStatements/CompactAveragingSplitting/EigenvectorLifting/InertDef.lean}{PositivityTheory.ThK.EvVac}
\leanref{Development/WeilBound/CoreStatements/LocalThetaKernels/EvlMain.lean}{PositivityTheory.ThKH.Evl.evVac}
\leanref{Development/WeilBound/ThetaLifts/SphericalLattices/FunDefs.lean}{PositivityTheory.Sph.EvInvSt}
\leanref{Development/WeilBound/ThetaLifts/SphericalLattices/FunEvInv.lean}{PositivityTheory.Sph.evInvSt}
\leanref{Development/WeilBound/ThetaLifts/SphericalLattices/FunDefs.lean}{PositivityTheory.Sph.SiegSt}
\leanref{Development/WeilBound/ThetaLifts/SphericalLattices/FunSieg.lean}{PositivityTheory.Sph.siegSt}
\leanref{Development/WeilBound/ThetaLifts/SphericalLattices/FunDefs.lean}{PositivityTheory.Sph.ModSt}
\leanref{Development/WeilBound/ThetaLifts/SphericalLattices/ModSt2.lean}{PositivityTheory.Sph.modSt2}
\begin{proof}
(1) By Definition~\ref{asm-theta-agree-triple}(4),
$(V,W,M,\beta,\beta/\alpha,\omega)$ and $(V',W,M',\alpha,\alpha/\beta,\omega')$ are
Weil pairs; in particular $\mathrm{vac}_v$ and $\mathrm{vac}'_v$ are vacua, as $M$
and $M'$ are regular, and $\alpha|_{\mathbf A^\times}$ and
$\beta|_{\mathbf A^\times}$ are powers of $\eta$. By
Lemma~\ref{asm-theta-agree-dictionary-conjugate},
$(\beta/\alpha)_{t\circ c}=(\beta_t/\alpha_t)^{-1}=\nu_t$ and
$\beta_{t\circ c}^{-1}=\beta_t$; and $(\alpha/\beta)_t=\nu_t$. Let $\psi_v$ be as in
Lemma~\ref{asm-theta-agree-dictionary-vacuum}. If $v\in P$ and $t=\tau_v$, let
$b=b_{V,v}$, $b_0=b_{V',v}$, $c=b_{W,v}$ and $\delta=1$: by
Definition~\ref{asm-theta-agree-doubling}(1), $b\,\mathcal O_v^a$,
$b_0\,\mathcal O_v^{a'}$ and $c\,\mathcal O_v^d$ are the completions of the
lattices of $V$, $V'$ and $W$ at the place of $E$ given by $t$, and $\psi_v$ has
conductor $\mathcal O_v$ (Definition~\ref{gal-good-base}). Otherwise choose $b$,
$b_0$, $c$ and $\delta$ by Lemma~\ref{asm-theta-agree-dictionary-models}(3). Put
$b'=\delta b_0$, which is $b_{V',v}$ in the first case. Define
$\gamma(g)=b^{-1}\pi_t(g)b$, $\gamma'(g')=b_0^{-1}\pi_t(g')b_0$ and
$\eta(h)=c^{-1}\pi_t(h)c$. Lemma~\ref{asm-theta-agree-dictionary-models}(1), for
the first Weil pair, gives $\iota$ with the stated properties, with
$(\chi_W)_{t\circ c}=(\beta/\alpha)_{t\circ c}=\nu_t$ and
$(\chi_X)_{t\circ c}^{-1}=\beta_t$. Part (2) of that lemma, for the second Weil
pair and $b_0$, $c$ and $\delta$, gives $\iota'$, of the second kind for
$(t,b',c)$, with $(\chi_W)_t=\nu_t$ and $(\chi_X)_t=\alpha_t$. Finally $\eta$ is
onto: for $x\in GL_d(F_v)$, Lemma~\ref{asm-theta-agree-coinvariants-dictionary}(2)
gives $h\in U_W^{(v)}$ with $\pi_t(h)=cxc^{-1}$, and then $\eta(h)=x$.

(2) Let $(X,W,M,\chi_X,\chi_W,\omega)$ be a Weil pair for $B$, with $X$ and $W$
of ranks $n$ and $d$, $I=\{1,\ldots,n\}\times\{1,\ldots,d\}$, isomorphism $e$
and Schrödinger operators $\rho$. Then $\rho(u)$ is the operator $\rho(e(u))$ of
Theorem~\ref{asm-theta-agree-evaluation-heisenberg} for $\psi$, which is
continuous, trivial on $F$ and not $1$ (Definition~\ref{gal-good-base}). By
conditions (a) and (b) of Definition~\ref{asm-theta-agree-models}(6), and as
$\omega(g,h)=\omega(g,1)\omega(1,h)$, each $\omega(g,h)$ maps Schwartz-Bruhat
functions to Schwartz-Bruhat functions, $\lambda\circ\omega(g,h)$ is admissible
for admissible $\lambda$, and $\omega(g,h)\rho(u)\Phi=\rho(guh^{\mathsf T})\omega(g,h)\Phi$
for Schwartz-Bruhat $\Phi$. Let $x_0$ and $r_0$ be as in
Lemma~\ref{asm-theta-agree-finite-inputs-slice}(4), and for $g\in U_X$,
$h\in U_W$, $u\in\mathbb W(X,W)$ and $\phi$ on $F_v^I$ put
\[
\omega_v(g,h)\phi=\operatorname{res}_{x_0}\bigl(\omega(g,h)(\phi\otimes r_0)\bigr),\qquad
\rho_v(u)\phi=\operatorname{res}_{x_0}\bigl(\rho(u)(\phi\otimes r_0)\bigr).
\]
Let $\phi$ be Schwartz-Bruhat. As $\phi\otimes r_0$ is Schwartz-Bruhat, so are
$\omega_v(g,h)\phi$ and $\rho_v(u)\phi$, by
Theorem~\ref{asm-theta-agree-evaluation-heisenberg}(1) and
Lemma~\ref{asm-theta-agree-finite-inputs-slice}(3).

Let $g\in U_X^{(v)}$ and $h\in U_W^{(v)}$. If $e(u)$ has $v$-components $0$, then
so does $u$, as $e$ is $\mathbf A$-linear and injective; as $g-1$ and $h-1$ are
at $v$, $guh^{\mathsf T}=u$. So $\omega(g,h)$ commutes with $\rho(w)$ on
Schwartz-Bruhat functions for every $w$ with $v$-components $0$, and by
Lemma~\ref{asm-theta-agree-finite-inputs-schur}(2) it acts at $v$ by
$\phi\mapsto\omega_v(g,h)\phi$. For $u$ at $v$, a direct computation gives
$\rho(u)(\phi\otimes r)=\rho^{(v)}(u)\phi\otimes r$ for every $r$ that does not
depend on the $v$-component, where
$(\rho^{(v)}(u)\phi)(y)=\psi\bigl(j(y)\cdot e_2(u)+\tfrac12e_1(u)\cdot e_2(u)\bigr)\phi(y+e_1(u)_v)$
and $j(y)$ has component $y$ at $v$ and $0$ elsewhere. So $\rho(u)$ acts at $v$
by $\phi\mapsto\rho^{(v)}(u)\phi$, and $\rho^{(v)}(u)\phi=\rho_v(u)\phi$ by
Lemma~\ref{asm-theta-agree-finite-inputs-schur}(1). As $\omega(1,1)$ is the
identity, the same lemma gives $\omega_v(1,1)\phi=\phi$.

For $g,g'\in U_X^{(v)}$ and $h,h'\in U_W^{(v)}$,
$\omega(gg',hh')(\phi\otimes r_0)=\omega(g,h)(\omega_v(g',h')\phi\otimes r_0)$, and
applying $\operatorname{res}_{x_0}$ gives
$\omega_v(gg',hh')\phi=\omega_v(g,h)\omega_v(g',h')\phi$. For $u$ at $v$,
\[
\omega_v(g,h)\rho_v(u)\phi
=\operatorname{res}_{x_0}\bigl(\omega(g,h)\rho(u)(\phi\otimes r_0)\bigr)
=\operatorname{res}_{x_0}\bigl(\rho(guh^{\mathsf T})(\omega_v(g,h)\phi\otimes r_0)\bigr)
=\rho_v(guh^{\mathsf T})\omega_v(g,h)\phi .
\]
For the neighbourhood in $U_W$, let $K_n$ be the subgroups of
Lemma~\ref{asm-theta-agree-finite-inputs-subgroups} for $W$, and apply
Lemma~\ref{asm-theta-agree-finite-inputs-smooth} to $G=U_W$, the
Schwartz-Bruhat functions inside all functions on $F_v^I$, and
$h\mapsto\omega_v(1,h)$. Its hypotheses hold by the above, since $K_0\subset U_W^{(v)}$,
and since $\lambda(\omega_v(1,h)\phi)=(\lambda\circ\operatorname{res}_{x_0})(\omega(1,h)(\phi\otimes r_0))$
is continuous in $h$ by condition (c) of
Definition~\ref{asm-theta-agree-models}(6), $\lambda\circ\operatorname{res}_{x_0}$
being admissible (Lemma~\ref{asm-theta-agree-finite-inputs-slice}(3)). So
$\phi$ is fixed by $\omega_v(1,k)$ for $k\in K_n$ for some $n$, and some
neighbourhood $U$ of $1$ has $U\cap U_W^{(v)}\subset K_n$. So
$(\omega_v,\rho_v)$ are local operators at $v$.

Now let $(\omega_v,\rho_v)$ be any local operators at $v$ for the pair. For
$g\in U_X^{(v)}$, $\omega(g,1)$ acts at $v$ by $\phi\mapsto\omega_v(g,1)\phi$, so
$\omega_v(g,1)\phi=\operatorname{res}_{x_0}(\omega(g,1)(\phi\otimes r_0))$ by
Lemma~\ref{asm-theta-agree-finite-inputs-schur}(1). Apply
Lemma~\ref{asm-theta-agree-finite-inputs-smooth} to $G=U_X$, the subgroups
$K_n$ of Lemma~\ref{asm-theta-agree-finite-inputs-subgroups} for $X$, and
$g\mapsto\operatorname{res}_{x_0}\circ\,\omega(g,1)\circ(\,\cdot\otimes r_0)$. On
$K_0$ this is $g\mapsto\omega_v(g,1)$, which preserves the Schwartz-Bruhat
functions, is the identity at $1$ and is multiplicative, and continuity holds
as before. So some $K_n$ fixes $\phi$, and a neighbourhood $U$ of $1$ in $U_X$
with $U\cap U_X^{(v)}\subset K_n$ is as required.

(3) Let $Y$ be one of $V$, $V'$ and $W$, of rank $r$, with lattice $L$. The
group $L_v$ of Definition~\ref{asm-theta-agree-finite-local-objects}(1) is
compact. Each standard basis vector $e_j$ of $E^r$ is an $E$-linear
combination of elements of $L$, so for every matrix $g$ with $gL_v\subset L_v$
the column $g\,e_j$ placed at $v$ lies in a compact set $C_j$ that does not
depend on $g$, built from $L_v$ by finite sums and scalar multiples. So the
matrices $g$ with $g-1$ at $v$ and $gL_v\subset L_v$ lie in a compact set $S$.
$K_Y(v)$ is closed in $U_Y$ and lies in the set of $k$ with $k,k^{-1}\in S$; as
$U_Y$ is closed in $GL_r(\mathbf A_E)$, whose topology is that of
$k\mapsto(k,k^{-1})$, $K_Y(v)$ is compact. For the index, let $a$ and $\Phi$ be
as in Lemma~\ref{asm-theta-agree-finite-inputs-frame}(1) and (2) for $W$, and
let $h\in U_W^{(v)}$. The subgroup
$GL_d(\mathcal O_w)\cap\Phi(h)GL_d(\mathcal O_w)\Phi(h)^{-1}$ is open in the
compact group $GL_d(\mathcal O_w)$, so it has finite index there, and its
intersection $N$ with $K_w$ has finite index in $K_w$. By
Lemma~\ref{asm-theta-agree-finite-inputs-frame}(2), the elements of $K_W(v)$
that $\Phi$ maps into $N$ form a subgroup of finite index in $K_W(v)$, and they
lie in $hK_W(v)h^{-1}$.

(4) Let $h\in U_W^{(v)}$ and $\Phi$ be as in (3). By
Theorem~\ref{asm-theta-agree-finite-inputs-bruhat-tits}(1),
$\Phi(h)^{-1}=\kappa_1\Phi(h)\kappa_2$ with $\kappa_1,\kappa_2\in K_w$, and
$\kappa_i=\Phi(k_i)$ with $k_i\in K_W(v)$ by
Lemma~\ref{asm-theta-agree-finite-inputs-frame}(2). As $\Phi$ is injective on
$U_W^{(v)}$, $h^{-1}=k_1hk_2$.

(5) Let $a$ and $\Phi$ be as in Lemma~\ref{asm-theta-agree-finite-inputs-frame}(1)
and (2), and let $C$ and $C'$ be the $d\times m$ matrices at $v$ whose
components at $w$ are the first and the last $m$ columns of $a$. Then
$C^{\mathsf T}Jc(C)$ and $C'^{\mathsf T}Jc(C')$ are blocks of $l\,A_{2m}$ that
vanish, and $C$ and $C'$ have rank $m$ at $w$, so they are Lagrangians of $W$
at $v$. Every $y\in L_v$ has $a^{-1}y_w\in\mathcal O_w^d$, and splitting this
vector into its first and last $m$ entries writes $y=Cs+C't$ with $Cs$ and $C't$
in $L_v$ by Lemma~\ref{asm-theta-agree-finite-inputs-frame}(1). So $C$ is
adapted. Let $h\in U_W^{(v)}$. By
Theorem~\ref{asm-theta-agree-finite-inputs-bruhat-tits}(2),
$\Phi(h)=\pi\kappa$ with $\pi\in P_w$ and $\kappa\in K_w$. Write $\pi=\Phi(p)$
and $\kappa=\Phi(k)$ with $p,k\in U_W^{(v)}$. Then $h=pk$, $k\in K_W(v)$, and as
$\Phi(p)$ carries the first $m$ standard basis vectors into their span,
$pC=Ct$ for some $t$, so $p\in P_C$.

(7) Let $a$ and $\Phi$ be as in Lemma~\ref{asm-theta-agree-finite-inputs-frame}(3)
for $C$. Then $\Phi$ is an isomorphism from $U_W^{(v)}$ onto $G_w$ carrying
$P_C^\times$ onto $P_w$ and $K_W(v)$ onto $K_w$, so it carries $H$ onto
$P_w\cap K_w$, $p^{-1}Hp$ onto $\Phi(p)^{-1}(P_w\cap K_w)\Phi(p)$ and $pHp^{-1}$
onto $\Phi(p)(P_w\cap K_w)\Phi(p)^{-1}$, preserving indices. The upper left block
of $\Phi(p)$ has determinant $\det t_w$. Let $q$ be the number of elements of
the residue field of $F_v$; as $E_w/F_v$ is unramified,
$|N_{E_w/F_v}(z)|=q^{-2\operatorname{ord}_w(z)}$. So (7) is
Lemma~\ref{asm-theta-agree-finite-inputs-modulus} for $\Phi(p)$.

(6) Let $e$, $\rho$ and $I$ be those of the pair, $J_X$ and $J_W$ the matrices
of $X$ and $W$, $\psi_v$ as in
Lemma~\ref{asm-theta-agree-finite-inputs-local-pair}, and $\rho_{\psi_v}$ the
operators of Theorem~\ref{asm-theta-agree-dictionary-stone-von-neumann} on the
functions on $F_v^I$. The character $\psi_v$ is continuous, and it is not $1$, as
it has conductor $\mathcal O_v$ (Definition~\ref{gal-good-base}). Let $u$ be at
$v$. As $e$ is $\mathbf A$-linear, $e(u)$ is at $v$, and the formula for $\rho$
gives $\rho(u)(\phi\otimes r)=(\rho_{\psi_v}(e(u)_v)\phi)\otimes r$ for every $r$
that does not depend on the $v$-component; so $\rho(u)$ acts at $v$ by
$\phi\mapsto\rho_{\psi_v}(e(u)_v)\phi$. By the property of $e$,
$\langle u,u'\rangle$ is at $v$ with component $[e(u)_v,e(u')_v]$ for every $u'$.

\emph{The form.} Let $C'$ be the second Lagrangian in the definition of an
adapted Lagrangian, and let $\Lambda_C$ be the set of the $e(zC^{\mathsf T})_v$,
$z\in M_{n\times m}(\mathbf A_E)$; each $zC^{\mathsf T}$ is at $v$, as $C$ is. By
the above and Lemma~\ref{asm-theta-agree-finite-unique}, a linear form is an
evaluation form along $C$ if and only if it is invariant under the
$\rho_{\psi_v}(a)$, $a\in\Lambda_C$, on Schwartz-Bruhat functions. The set
$\Lambda_C$ is an $F_v$-subspace, as $s\,e(zC^{\mathsf T})_v=e((\bar sz)C^{\mathsf T})_v$,
and $[a,a']=0$ for $a,a'\in\Lambda_C$, as
$h(zC^{\mathsf T},z'C^{\mathsf T})=\sum_{i,j}(z^{\mathsf T}J_Xc(z'))_{ij}(C^{\mathsf T}J_Wc(C))_{ij}=0$.
The same holds for $C'$. By Lemma~\ref{asm-theta-agree-finite-inputs-invariant}
there is an evaluation form $l$ along $C$ with $l(\phi_1)\neq0$ for some
Schwartz-Bruhat $\phi_1$.

\emph{Its value at the vacuum.} Let $\Lambda_v$ be the group of
Lemma~\ref{asm-theta-agree-evaluation-lattice}, generated by the $a\,xy^{\mathsf T}$
with $a\in\mathbf A$ at $v$, $a_v\in\mathcal O_v$, $x\in L_X$ and $y\in L_W$.
\begin{itemize}
\item[(i)] Every $\gamma\in\Lambda_v$ is $zC^{\mathsf T}+z'C'^{\mathsf T}$ with
  $zC^{\mathsf T}\in\Lambda_v$. For a generator, $ay\in L_v$, so $ay=Cs+C't$ with
  $Cs,C't\in L_v$, and $a\,xy^{\mathsf T}=(xs^{\mathsf T})C^{\mathsf T}+(xt^{\mathsf T})C'^{\mathsf T}$,
  where $x(Cs)^{\mathsf T}\in\Lambda_v$ as $Cs\in L_v$; sums and negatives follow.
\item[(ii)] Every $u$ at $v$ is $zC^{\mathsf T}+z'C'^{\mathsf T}$. The matrices
  $xy^{\mathsf T}$, $x\in L_X$, $y\in L_W$, span $M_{n\times d}(\mathbf A_E)$ over
  $\mathbf A$, as $L_X$ and $L_W$ are lattices; so $u=\bar1u$ is an
  $\mathbf A$-linear combination of the $\bar1\,xy^{\mathsf T}\in\Lambda_v$, and
  (i) applies to these.
\end{itemize}
Apply Lemma~\ref{asm-theta-agree-finite-inputs-criterion} to $F_v$,
$\psi_v$, $\phi_0=\mathrm{vac}_v$, which is a nonzero Schwartz-Bruhat function
as $M_0$ is regular, $A=\Lambda_C$ and $l$. Let $\vartheta\in F_v^I\times F_v^I$;
then $\vartheta=e(u)_v$ for the element $u=e^{-1}(\bar\vartheta)$, which is at
$v$. Write $u=zC^{\mathsf T}+z'C'^{\mathsf T}$ by (ii), and put
$a=e(zC^{\mathsf T})_v\in\Lambda_C$, $b=e(z'C'^{\mathsf T})_v$ and
$\gamma=z'C'^{\mathsf T}$, which is at $v$, so that $\vartheta=a+b$. If
$\gamma\in\Lambda_v$, then $\rho(\gamma)$ acts at $v$ by
$\mathrm{vac}_v\mapsto\psi(\tfrac12\operatorname{Tr}h(\theta\gamma,\gamma))\mathrm{vac}_v$
by the definition of a vacuum, so $\rho_{\psi_v}(b)\mathrm{vac}_v$ is a multiple of
$\mathrm{vac}_v$ by Lemma~\ref{asm-theta-agree-finite-unique}. Otherwise, by
Lemma~\ref{asm-theta-agree-evaluation-lattice}(2), some $\gamma_0\in\Lambda_v$ has
$\psi(\langle\gamma_0,\gamma\rangle)\neq1$. Write
$\gamma_0=yC^{\mathsf T}+y'C'^{\mathsf T}$ with $yC^{\mathsf T}\in\Lambda_v$ by (i).
As $\langle y'C'^{\mathsf T},\gamma\rangle=0$,
$\psi(\langle yC^{\mathsf T},\gamma\rangle)\neq1$, that is,
$\psi_v([a',b])\neq1$ for $a'=e(yC^{\mathsf T})_v\in\Lambda_C$. And
$\rho_{\psi_v}(a')\mathrm{vac}_v=\mathrm{vac}_v$, by the definition of a vacuum
for $yC^{\mathsf T}\in\Lambda_v$, as
$h(\theta yC^{\mathsf T},yC^{\mathsf T})=\theta\,h(yC^{\mathsf T},yC^{\mathsf T})=0$.
So if $l(\mathrm{vac}_v)=0$, then $l$ vanishes on the Schwartz-Bruhat functions,
which contradicts $l(\phi_1)\neq0$. So $l_0=l$ is as required.

\emph{The local pair.} Let $w$, $u\mapsto u_w$, $\chi\mapsto\chi_w$, $e_v$ and
$\omega_w$ be as in Lemma~\ref{asm-theta-agree-finite-inputs-local-pair} for the
pair, and put $\rho_{e_v}(x)=\rho_{\psi_v}(e_v(x))$. As $X$ is Hermitian and $W$
skew-Hermitian, $J_X$ and $J_W$ are Hermitian and skew-Hermitian over $E_w$ in
the sense of Definition~\ref{asm-theta-agree-finite-inputs-local-data}: their
determinants are not $0$, as the matrices $\tau_{w'}(J_X)$ and
$i\tau_{w'}(J_W)$ of Definition~\ref{asm-theta-agree-doubling}(1) have
signatures $(r,s)$ with $r+s$ the rank. The matrix $C_w$ is a local Lagrangian of
$J_W$. Indeed $C_w^{\mathsf T}J_Wc(C_w)=0$ is the component at $w$ of
$C^{\mathsf T}J_Wc(C)=0$, and $C_w$ has rank $m$, as $Ct\neq0$ for nonzero $t$ at
$v$. Let $L_0$ be a left inverse of $C_w$, put $Z=c(J_W^{-1}L_0^{\mathsf T})$, so
that $C_w^{\mathsf T}J_Wc(Z)=1$, and
$C''=Z-\tfrac12C_w(Z^{\mathsf T}J_Wc(Z))^{\mathsf T}$. Using
$c(J_W)^{\mathsf T}=-J_W$ one finds $C_w^{\mathsf T}J_Wc(C'')=1$,
$C''^{\mathsf T}J_Wc(C_w)=-1$ and $C''^{\mathsf T}J_Wc(C'')=0$, so that
$(C_w\ C'')^{\mathsf T}J_Wc(C_w\ C'')=\left(\begin{smallmatrix}0&1\\-1&0\end{smallmatrix}\right)$
and $(C_w\ C'')$ is invertible.

For $z\in M_{n\times m}(E_w)$ let $\tilde z$ be the matrix at $v$ with component
$z$. Then $\tilde zC^{\mathsf T}$ is at $v$ with component $zC_w^{\mathsf T}$, so
by Lemma~\ref{asm-theta-agree-finite-inputs-local-pair}(1)
$\rho(\tilde zC^{\mathsf T})$ acts at $v$ by
$\phi\mapsto\rho_{e_v}(zC_w^{\mathsf T})\phi$. Hence every evaluation form $l_0$
along $C$ has $l_0(\rho_{e_v}(zC_w^{\mathsf T})\phi)=l_0(\phi)$ for
$z\in M_{n\times m}(E_w)$ and Schwartz-Bruhat $\phi$. For $g\in U_X^{(v)}$ and
$h\in U_W^{(v)}$, $\omega_0(g,h)$ acts at $v$ by $\phi\mapsto\omega_v(g,h)\phi$
and, by Lemma~\ref{asm-theta-agree-finite-inputs-local-pair}(3), by
$\phi\mapsto\omega_w(g_w,h_w)\phi$; by Lemma~\ref{asm-theta-agree-finite-unique}
these agree on Schwartz-Bruhat $\phi$.

\emph{Invariance under $U_X^{(v)}$.} Suppose that $l_0$ is not $0$ on the
Schwartz-Bruhat functions and that $\chi_W(z)=1$ whenever $z-1$ is at $v$. Then
$\chi_{W,w}=1$, as each $z_y-1$ is at $v$. So $\omega_w$ is a local Weil pair for
$(J_X,J_W,e_v,\chi_{X,w},1)$, and
Lemma~\ref{asm-theta-agree-finite-inputs-fixed} gives
$l_0(\omega_w(g',1)\phi)=l_0(\phi)$ for every $g'$ in the unitary group of $J_X$
over $E_w$. For $g\in U_X^{(v)}$ this gives
$l_0(\omega_v(g,1)\phi)=l_0(\omega_w(g_w,1)\phi)=l_0(\phi)$.

\emph{The parabolic.} Let $p$, $t$ and $z$ be as in (6), and $\phi$
Schwartz-Bruhat. If $l_0$ vanishes on the Schwartz-Bruhat functions, both sides
are $0$, as $\omega_v(1,p)\phi$ is Schwartz-Bruhat. Otherwise, $p_wC_w=C_wt_w$,
and $t_w$ is invertible, as $\det t_w=z_w\neq0$. As $z-1$ is at $v$,
$z=z_{\det t_w}$ and $\chi_X(z)=\chi_{X,w}(\det t_w)$. By
Lemma~\ref{asm-theta-agree-finite-inputs-transport}, the local Weil datum
$h\mapsto\omega_w(1,h)$ for $(J_X,J_W,e_v,\chi_{X,w})$ satisfies the local Siegel
rule for $(C_w,\chi_{X,w})$, so
\[
l_0(\omega_v(1,p)\phi)=l_0(\omega_w(1,p_w)\phi)
=\chi_X(z)\,\bigl|N_{E_w/F_v}(\det t_w)\bigr|^{n/2}\,l_0(\phi).
\]
% Lean: (6) is evVac, evInvSt and siegSt. evVac: psiAt_good; isEvAt_iff and lagr_iso for the
%   evaluation forms along C (acts_iff = rho_acts_locCo with acts_unique, symp_lag); the form is
%   LagFun.exists_invariant_form; the value at the vacuum is core, with exists_locCo, space_split and
%   lat_split for (ii) and (i) (LatOpenCpt.span_top, vecMulVec_mem), the vacuum field of M.IsGood,
%   psi_form, and mem_latAtV_of_dual, which proves Lemma asm-theta-agree-evaluation-lattice(2) again
%   for psi_v of conductor O_v. evInvSt and siegSt take the local pair of ThKH.Locpair.locPairEx, with
%   PT.isHerm_loc, PT.isSkew_loc, PT.isLagr_loc (PT.isLagr_of_free), PT.isEv_of (Iwas.sgE for z~),
%   PT.actsX and PT.actsW; evInvSt applies split_character_fixes_evaluation with EvInv.locChar_one,
%   and siegSt applies siegel_transport to L.weil.1, with L.chX for chi_X(z).
% Done above: (1) split_place_dictionary (model_alt, model_Bar, chrW_conj, exists_isSplitBasis,
%   exists_cond, and dictX for the surjectivity of inB tau c); (2) locRepEx with IsLocRep.smoothX;
%   (3) stabiliser_compact (isCompact_KAt) and finSt; (4) cartanSt2; (5) ThKH.Iwas.iwasawa;
%   (7) modSt2 = M2.modSt_of M2.chart M2.locMod.
\end{proof}

\begin{lemma}\label{asm-theta-agree-finite-split-core}
Let $(B,G,\mathcal X,\omega,\omega')$ be a triple of ranks $(a,a')$
(Definition~\ref{asm-theta-agree-triple}), with models $M$ and $M'$, put
$d=a+a'$, $I_0=\{1,\ldots,a\}\times\{1,\ldots,d\}$ and
$I'_0=\{1,\ldots,a'\}\times\{1,\ldots,d\}$, let $v$ be a finite place of $F$,
and let $\nu$, $\alpha$ and $\beta$ be characters of $F_v^\times$, with
$\omega_A=\omega_A^{\nu,\beta}$ and $\omega_B=\omega_B^{\nu,\alpha}$
(Definition~\ref{asm-theta-agree-finite-split-models}). Let $I$ be a space of
functions on $GL_d(F_v)$ such that $0$ and $I$ are the only subspaces $N\subset I$
with $R(h_1)N\subset N$ for every $h_1\in GL_d(F_v)$. Let $f$ and $f'$ be maps
from the functions on $M_{a\times d}(F_v)$ and on $M_{a'\times d}(F_v)$ to the
functions on $GL_d(F_v)$ that are linear on the Schwartz-Bruhat functions and map
them onto $I$, such that for Schwartz-Bruhat $\Xi$ and $\Psi$ and all
$g\in GL_a(F_v)$, $g'\in GL_{a'}(F_v)$ and $h\in GL_d(F_v)$
\[
f(\omega_A(g,1)\Xi)=f(\Xi),\quad f(\omega_A(1,h)\Xi)=R(h)f(\Xi),\quad
f'(\omega_B(g',1)\Psi)=f'(\Psi),\quad f'(\omega_B(1,h)\Psi)=R(h)f'(\Psi),
\]
and $f(\Xi)\neq0$ for some Schwartz-Bruhat $\Xi$. Let $\iota$ and $\iota'$ be
linear maps from the functions on $F_v^{I_0}$ and $F_v^{I'_0}$ to the functions
on $M_{a\times d}(F_v)$ and $M_{a'\times d}(F_v)$, bijective between the
Schwartz-Bruhat functions, and let $\gamma\colon U_V^{(v)}\to GL_a(F_v)$,
$\gamma'\colon U_{V'}^{(v)}\to GL_{a'}(F_v)$ and
$\eta\colon U_W^{(v)}\to GL_d(F_v)$ be maps, with $\eta$ onto, such that for
$g\in U_V^{(v)}$, $g'\in U_{V'}^{(v)}$ and $h\in U_W^{(v)}$ the operators
$\omega(g,1)$ and $\omega(1,h)$ read as $\omega_A(\gamma(g),1)$ and
$\omega_A(1,\eta(h))$ through $\iota$, and $\omega'(g',1)$ and $\omega'(1,h)$
read as $\omega_B(\gamma'(g'),1)$ and $\omega_B(1,\eta(h))$ through $\iota'$.
Then there are linear maps $p_v$ and $p'_v$ from the functions on $F_v^{I_0}$ and
on $F_v^{I'_0}$ to $\Pi_v=I$ with $p_v(\phi)=f(\iota\phi)$ and
$p'_v(\phi')=f'(\iota'\phi')$ for Schwartz-Bruhat $\phi$ and $\phi'$, and these
satisfy the conditions at $v$ of Definition~\ref{asm-theta-agree-matched}(3)
other than the last two:
\begin{enumerate}
\item $p_v$ and $p'_v$ map the Schwartz-Bruhat functions onto $\Pi_v$, $p_v$
  vanishes on $\mathfrak a_v(\omega)$ and $p'_v$ on $\mathfrak a_v(\omega')$;
\item if $h\in U_W^{(v)}$, $\omega(1,h)$ and $\omega'(1,h)$ act at $v$ by
  $\phi\mapsto\phi_1$ and $\phi'\mapsto\phi'_1$ for Schwartz-Bruhat functions,
  and $p_v(\phi)=p'_v(\phi')$, then $p_v(\phi_1)=p'_v(\phi'_1)$;
\item $\Pi_v\neq0$, and $0$ and $\Pi_v$ are the only subspaces $N$ with
  $p_v(\phi_1)\in N$ whenever $p_v(\phi)\in N$ and $\omega(1,h)$ acts at $v$ by
  $\phi\mapsto\phi_1$ for some $h\in U_W^{(v)}$ and Schwartz-Bruhat
  $\phi,\phi_1$.
\end{enumerate}
\end{lemma}
\leanref{Development/WeilBound/CoreStatements/CompactAveragingSplitting/EigenvectorLifting/SplitFin.lean}{PositivityTheory.ThK.SplitAsm.core}
\leanref{Development/WeilBound/CoreStatements/CompactAveragingSplitting/EigenvectorLifting/SplitAsm.lean}{PositivityTheory.ThK.SplitAsm.mk_p}
\begin{proof}
As $f\circ\iota$ is linear on the Schwartz-Bruhat functions on $F_v^{I_0}$ and
takes values in $I$, it extends to a linear map $p_v$ on all functions; likewise
$p'_v$. Since $\iota$ is bijective between the Schwartz-Bruhat functions and $f$
maps them onto $I$, $p_v$ maps the Schwartz-Bruhat functions onto $\Pi_v$, and so
does $p'_v$.

Let $g\in U_V^{(v)}$, and let $\omega(g,1)$ act at $v$ by $\phi\mapsto\phi_1$
with $\phi,\phi_1$ Schwartz-Bruhat. Reading through $\iota$ gives a
Schwartz-Bruhat $\phi_2$ with $\omega(g,1)$ acting at $v$ by $\phi\mapsto\phi_2$
and $\iota\phi_2=\omega_A(\gamma(g),1)\iota\phi$. By
Lemma~\ref{asm-theta-agree-finite-unique}, $\phi_1=\phi_2$, so
$p_v(\phi_1)=f(\omega_A(\gamma(g),1)\iota\phi)=f(\iota\phi)=p_v(\phi)$. Hence
$p_v$ vanishes on $\mathfrak a_v(\omega)$, and in the same way $p'_v$ vanishes on
$\mathfrak a_v(\omega')$.

Let $h\in U_W^{(v)}$, let $\omega(1,h)$ and $\omega'(1,h)$ act at $v$ by
$\phi\mapsto\phi_1$ and $\phi'\mapsto\phi'_1$, all four functions
Schwartz-Bruhat, and suppose $p_v(\phi)=p'_v(\phi')$. As above,
$\iota\phi_1=\omega_A(1,\eta(h))\iota\phi$ and
$\iota'\phi'_1=\omega_B(1,\eta(h))\iota'\phi'$, so
\[
p_v(\phi_1)=R(\eta(h))f(\iota\phi)=R(\eta(h))f'(\iota'\phi')=p'_v(\phi'_1).
\]
As $f(\Xi)\neq0$ for some $\Xi$, $\Pi_v\neq0$.

Let $N\subset\Pi_v$ be a subspace with $p_v(\phi_1)\in N$ whenever
$p_v(\phi)\in N$ and $\omega(1,h)$ acts at $v$ by $\phi\mapsto\phi_1$ for some
$h\in U_W^{(v)}$ and Schwartz-Bruhat $\phi,\phi_1$. Let $y\in N$ and
$h_1\in GL_d(F_v)$. Write $y=p_v(\phi)$ with $\phi$ Schwartz-Bruhat and
$h_1=\eta(h)$ with $h\in U_W^{(v)}$. Reading $\omega(1,h)$ through $\iota$ gives
$\phi_1$ with $\omega(1,h)$ acting at $v$ by $\phi\mapsto\phi_1$ and
$p_v(\phi_1)=f(\omega_A(1,h_1)\iota\phi)=R(h_1)y$, which lies in $N$. So $N$ is
stable under every $R(h_1)$, and $N=0$ or $N=\Pi_v$ by the hypothesis on $I$.
\end{proof}

\begin{lemma}\label{asm-theta-agree-finite-ramified}
Let $(B,G,\mathcal X,\omega,\omega')$ be a triple
(Definition~\ref{asm-theta-agree-triple}), and let $v\in T$. Then there are a
complex vector space $\Pi_v$ and linear maps $p_v$ and $p'_v$ satisfying the
conditions at $v$ of Definition~\ref{asm-theta-agree-matched}(3).
\end{lemma}
\leanref{Development/WeilBound/CoreStatements/CompactAveragingSplitting/EigenvectorLifting/FinSplit.lean}{PositivityTheory.ThK.FinRam}
\leanref{Development/WeilBound/CoreStatements/CompactAveragingSplitting/EigenvectorLifting/FinSplit.lean}{PositivityTheory.ThK.IsFinAt}
\leanref{Development/WeilBound/CoreStatements/CompactAveragingSplitting/EigenvectorLifting/SplitFin.lean}{PositivityTheory.ThK.SplitAsm.finRam_of}
\leanref{Development/AutomorphicSpectrum/Proofs/BlockDecomposition/IsSmoothCharOfCondLE.lean}{GlobalTheory.CompanionTheory.SplitNarrow.finRam_of_narrow}
\leanref{Development/WeilBound/CoreStatements/TemperedRepresentationTheory/BernsteinNarrowness/BernsteinClosure/RamifiedMultiplicityOne.lean}{PositivityTheory.BernsteinRamification.Close.finRam}
\begin{proof}
Let $(a,a')$ be the ranks of the triple and $d=a+a'$. As $T\subset P$, the base
gives an $F$-embedding $\tau_v\colon E\to F_v$. Let $\alpha_v$, $\beta_v$ and
$\nu_v$ be the characters of $F_v^\times$ given by $\alpha$, $\beta$ and $\nu$ at
the place of $E$ defined by $\tau_v$. By
Lemma~\ref{asm-theta-agree-finite-characters} they are unitary and smooth,
$\nu_v=\alpha_v/\beta_v$, and $\nu_v$ is ramified. The triple datum is proper,
so $a'<a$ and $d$ is even. Let $Z$ and $Z'$ be the maps of
Lemma~\ref{th-bottom-split-zeta} for $k=F_v$, $\alpha_v$, $\beta_v$ and a fixed
uniformizer, with the properties of Lemma~\ref{asm-theta-agree-finite-zeta}(1).
The image of the Schwartz-Bruhat functions under $Z$ is a nonzero subspace of
$I(\alpha_v,\beta_v)$, stable under every $R(h)$ since
$R(h)Z(\Xi)=Z(\omega_A(1,h)\Xi)$ and $\omega_A(1,h)\Xi$ is Schwartz-Bruhat. By
Theorem~\ref{asm-theta-agree-finite-bernstein}, applied with $a\neq a'$, it is
$I(\alpha_v,\beta_v)$. The same holds for $Z'$.

Let $b,b',c,\iota,\iota',\gamma,\gamma',\eta$ be given by
Lemma~\ref{asm-theta-agree-finite-inputs}(1) for $t=\tau_v$; as $v\in P$,
$b=b_{V,v}$, $b'=b_{V',v}$ and $c=b_{W,v}$.
Lemma~\ref{asm-theta-agree-finite-split-core}, with $\nu_v,\alpha_v,\beta_v$,
$I=I(\alpha_v,\beta_v)$, $f=Z$ and $f'=Z'$, gives $\Pi_v=I(\alpha_v,\beta_v)$,
$p_v$ and $p'_v$ with all the conditions at $v$ other than the last two. The
condition on the vacua concerns only places outside $T$. For the zeta condition,
$\iota$ and $\iota'$ are linear models of the first kind for
$(\tau_v,b_{V,v},b_{W,v})$ and of the second kind for
$(\tau_v,b_{V',v},b_{W,v})$ that carry the vacua to the indicator functions of
the integral matrices, and for Schwartz-Bruhat $\phi$ and $\phi'$,
$p_v(\phi)=p'_v(\phi')$ if and only if $Z(\iota\phi)=Z'(\iota'\phi')$, as both
sides are elements of $\Pi_v$.
\end{proof}

\begin{lemma}\label{asm-theta-agree-finite-split}
Let $(B,G,\mathcal X,\omega,\omega')$ be a triple
(Definition~\ref{asm-theta-agree-triple}), and let $v\notin T$ be a finite place
of $F$ such that $E$ has an $F$-embedding into $F_v$. Then there are a complex
vector space $\Pi_v$ and linear maps $p_v$ and $p'_v$ satisfying the conditions
at $v$ of Definition~\ref{asm-theta-agree-matched}(3).
\end{lemma}
\leanref{Development/WeilBound/CoreStatements/CompactAveragingSplitting/EigenvectorLifting/FinSplit.lean}{PositivityTheory.ThK.FinSplit}
\leanref{Development/WeilBound/CoreStatements/CompactAveragingSplitting/EigenvectorLifting/SplitFin.lean}{PositivityTheory.ThK.SplitAsm.finSplit_of}
\leanref{Development/AutomorphicSpectrum/Proofs/BlockDecomposition/FinSplitOfNarrow.lean}{GlobalTheory.CompanionTheory.BernBridge.Split.finSplit_of_lt}
\leanref{Development/AutomorphicSpectrum/Proofs/BlockDecomposition/FinSplitOfNarrow.lean}{GlobalTheory.CompanionTheory.BernBridge.Split.finSplit_of_narrow}
\leanref{Development/WeilBound/CoreStatements/TemperedRepresentationTheory/BernsteinNarrowness/BernsteinClosure/RamifiedMultiplicityOne.lean}{PositivityTheory.BernsteinRamification.Close.finSplit}
\begin{proof}
Let $(a,a')$ be the ranks of the triple, $d=a+a'$, $t\colon E\to F_v$ an
$F$-embedding, and $\alpha_t$, $\beta_t$ and $\nu_t$ as in
Lemma~\ref{asm-theta-agree-finite-characters}: they are unitary and smooth, with
$\nu_t=\alpha_t/\beta_t$, and $\alpha_t$ and $\beta_t$ are unramified as
$v\notin T$; so $\nu_t$ is unramified. As in the proof of
Lemma~\ref{asm-theta-agree-finite-ramified}, $a'<a$ and $d$ is even. Let
$Z^{\mathrm u}$ and $Z'^{\mathrm u}$ be the maps of
Lemma~\ref{asm-theta-agree-finite-zeta}(2) for $k=F_v$, $\alpha_t$ and
$\beta_t$. They are not $0$, since they take the value $1$ at $1$ on the
indicator functions of the integral matrices, so as in that proof
Theorem~\ref{asm-theta-agree-finite-bernstein} shows that they map the
Schwartz-Bruhat functions onto $I(\alpha_t,\beta_t)$.

Let $\iota,\iota',\gamma,\gamma',\eta$ be given by
Lemma~\ref{asm-theta-agree-finite-inputs}(1) for $t$, and apply
Lemma~\ref{asm-theta-agree-finite-split-core} with $\nu_t,\alpha_t,\beta_t$,
$I=I(\alpha_t,\beta_t)$, $f=Z^{\mathrm u}$ and $f'=Z'^{\mathrm u}$. This gives
$\Pi_v=I(\alpha_t,\beta_t)$, $p_v$ and $p'_v$ with all the conditions at $v$
other than the last two. The zeta condition concerns only places of $T$. For the
vacua, $\iota\mathrm{vac}_v$ and $\iota'\mathrm{vac}'_v$ are the indicator
functions of the integral matrices, which are Schwartz-Bruhat, so
$F=p_v(\mathrm{vac}_v)-p'_v(\mathrm{vac}'_v)$ is an element of
$I(\alpha_t,\beta_t)$ that vanishes on $GL_d(\mathcal O_v)$ by
Lemma~\ref{asm-theta-agree-finite-zeta}(2). By
Lemma~\ref{asm-theta-agree-finite-iwasawa}, $F=0$, so
$p_v(\mathrm{vac}_v)=p'_v(\mathrm{vac}'_v)$, and this is not $0$ since its value
at $1$ is $1$.
\end{proof}

\begin{situation}\label{asm-theta-agree-finite-inert-situation}
Let $(B,G,\mathcal X,\omega,\omega')$ be a triple of ranks $(a,a')$ for a theta
base $B$ of a CM field $E$ (Definition~\ref{asm-theta-agree-triple}), with models
$M$ and $M'$, let $I_0$ and $I'_0$ be as in
Definition~\ref{asm-theta-agree-matched}, and let $v\notin T$ be a finite place
of $F$ such that $E$ has no $F$-embedding into $F_v$; as $E/F$ is unramified at
the finite places, $v$ is inert in $E$. Let $(\omega_v,\rho_v)$
and $(\omega'_v,\rho'_v)$ be local operators at $v$
(Definition~\ref{asm-theta-agree-finite-local-objects}(4)) for the Weil pairs
$(V,W,M,\beta,\beta/\alpha,\omega)$ and $(V',W,M',\alpha,\alpha/\beta,\omega')$,
and let $K_V(v)$, $K_{V'}(v)$ and $K_W(v)$ be the stabilizers at $v$.
\end{situation}

\begin{theorem}[Howe \cite{How79}]\label{asm-theta-agree-finite-howe}
In Situation~\ref{asm-theta-agree-finite-inert-situation}, let
$(X,\omega_0,\omega_{0,v},\mathrm{vac}_v)$ be $(V,\omega,\omega_v,\mathrm{vac}_v)$
or $(V',\omega',\omega'_v,\mathrm{vac}'_v)$. Then:
\begin{enumerate}
\item $\omega_{0,v}(g,1)\mathrm{vac}_v=\mathrm{vac}_v$ for $g\in K_X(v)$, and
  $\omega_{0,v}(1,k)\mathrm{vac}_v=\mathrm{vac}_v$ for $k\in K_W(v)$;
\item every Schwartz-Bruhat function fixed by every $\omega_{0,v}(g,1)$,
  $g\in K_X(v)$, lies in the span of the $\omega_{0,v}(1,h)\mathrm{vac}_v$,
  $h\in U_W^{(v)}$;
\item every Schwartz-Bruhat function fixed by every $\omega_{0,v}(1,k)$,
  $k\in K_W(v)$, lies in the span of the $\omega_{0,v}(g,1)\mathrm{vac}_v$,
  $g\in U_X^{(v)}$.
\end{enumerate}
\end{theorem}
\leanref{Development/WeilBound/ThetaLifts/SphericalLattices/HoweInert.lean}{PositivityTheory.Sph.IsHowe}
\leanref{Development/WeilBound/ThetaLifts/SphericalLattices/HoweInert.lean}{PositivityTheory.Sph.HoweInert}
\leanref{Development/WeilBound/ThetaLifts/SphericalLattices/HoweDuality.lean}{PositivityTheory.Sph.howeInert}

\begin{lemma}\label{asm-theta-agree-finite-vacuum}
In Situation~\ref{asm-theta-agree-finite-inert-situation}, let
$(X,\omega_0,\omega_{0,v},\mathrm{vac}_v)$ be $(V,\omega,\omega_v,\mathrm{vac}_v)$
or $(V',\omega',\omega'_v,\mathrm{vac}'_v)$. Then:
\begin{enumerate}
\item $\omega_{0,v}(1,k)\mathrm{vac}_v=\mathrm{vac}_v$ for every $k\in K_W(v)$;
\item every Schwartz-Bruhat $\phi$ lies in the sum of $\mathfrak a_v(\omega_0)$ and
  the span of the $\omega_{0,v}(1,h)\mathrm{vac}_v$, $h\in U_W^{(v)}$;
\item for every Schwartz-Bruhat $\phi$ there are $m\geq0$,
  $k_0,\ldots,k_m\in K_W(v)$ and $c\in\mathbf C$ with
  $\sum_i\omega_{0,v}(1,k_i)\phi-c\,\mathrm{vac}_v\in\mathfrak a_v(\omega_0)$.
\end{enumerate}
\end{lemma}
\leanref{Development/WeilBound/ThetaLifts/SphericalLattices/HoweInert.lean}{PositivityTheory.Sph.IsVacSph}
\leanref{Development/WeilBound/ThetaLifts/SphericalLattices/HoweInert.lean}{PositivityTheory.Sph.VacInert}
\leanref{Development/WeilBound/ThetaLifts/SphericalLattices/ExistsFixedSum.lean}{PositivityTheory.Sph.isVacSph_of}
\leanref{Development/WeilBound/ThetaLifts/SphericalLattices/ExistsFixedSum.lean}{PositivityTheory.Sph.vacInert}
\leanref{Development/WeilBound/ThetaLifts/SphericalLattices/ExistsFixedSum.lean}{PositivityTheory.Sph.exists_fixed_sum}
\begin{proof}
Part (1) is part of Theorem~\ref{asm-theta-agree-finite-howe}(1). For
$g\in U_X^{(v)}$ and Schwartz-Bruhat $\psi$, $\omega_0(g,1)$ acts at $v$ by
$\psi\mapsto\omega_{0,v}(g,1)\psi$, so
$\omega_{0,v}(g,1)\psi-\psi\in\mathfrak a_v(\omega_0)$.

If a compact group $K$ acts linearly on a space and the elements of $K$ fixing a
vector $\phi$ form an open subgroup $K_\phi$, then $K_\phi$ has finite index, and
for representatives $k_0,\ldots,k_m$ of $K/K_\phi$ the vector $\sum_ik_i\phi$ is
fixed by $K$. By Lemma~\ref{asm-theta-agree-finite-inputs}(3), $K_X(v)$ and
$K_W(v)$ are compact. By the smoothness of local operators
(Definition~\ref{asm-theta-agree-finite-local-objects}(4) and
Lemma~\ref{asm-theta-agree-finite-inputs}(2)), the elements of $K_X(v)$ fixing a
Schwartz-Bruhat $\phi$ under $\omega_{0,v}(\cdot,1)$ form a subgroup containing a
neighbourhood of $1$, hence open, and the same holds for $K_W(v)$ and
$\omega_{0,v}(1,\cdot)$.

For (2), let $f=\sum_{i=0}^m\omega_{0,v}(k_i,1)\phi$ be fixed by $K_X(v)$, with
$k_i\in K_X(v)$. By Theorem~\ref{asm-theta-agree-finite-howe}(2), $f$ lies in
the span of the $\omega_{0,v}(1,h)\mathrm{vac}_v$, and
$(m+1)\phi-f=-\sum_i(\omega_{0,v}(k_i,1)\phi-\phi)\in\mathfrak a_v(\omega_0)$.
Dividing by $m+1$ gives (2).

For (3), let $f=\sum_{i=0}^m\omega_{0,v}(1,k_i)\phi$ be fixed by $K_W(v)$, with
$k_i\in K_W(v)$. By Theorem~\ref{asm-theta-agree-finite-howe}(3), $f$ is a
finite sum $\sum_jc_j\,\omega_{0,v}(g_j,1)\mathrm{vac}_v$, and
$f-(\sum_jc_j)\mathrm{vac}_v\in\mathfrak a_v(\omega_0)$ by the first paragraph.
\end{proof}

\begin{lemma}\label{asm-theta-agree-finite-average}
Let $K$ be a group acting by linear operators $\sigma(k)$ on a complex vector
space $S$, let $V\subset S$ be a subspace with $\sigma(k)V\subset V$,
$\sigma(1)\phi=\phi$ and $\sigma(kk')\phi=\sigma(k)\sigma(k')\phi$ for
$k,k'\in K$ and $\phi\in V$, such that every $\phi\in V$ is fixed by a subgroup of
finite index in $K$, and let $l_0$ be a linear form on $S$. Then there is a
linear form $\lambda$ on $S$ such that $\lambda(\sigma(k)\phi)=\lambda(\phi)$ for
$k\in K$ and $\phi\in V$, and
\[
\lambda(\phi)=[K:H]^{-1}\sum_{q\in K/H}l_0(\sigma(q)\phi)
\]
for every $\phi\in V$ and every subgroup $H\subset K$ of finite index whose
elements fix $\phi$, the sum running over representatives $q$ of the cosets.
\end{lemma}
\leanref{Development/WeilBound/ThetaLifts/SphericalLattices/FunDefs.lean}{PositivityTheory.Sph.AvgSt}
\leanref{Development/WeilBound/ThetaLifts/SphericalLattices/FunAvg.lean}{PositivityTheory.Sph.avgSt}
\begin{proof}
For $\phi\in V$ and $H$ as in the statement, write $A_H(\phi)$ for the right
side; it does not depend on the representatives, since $H$ fixes $\phi$. If
$H'\subset H$ is another such subgroup, each coset of $H$ is the union of
$[H:H']$ cosets of $H'$ on which $\sigma(q)\phi$ is constant, so
$A_{H'}(\phi)=A_H(\phi)$. Any two such subgroups contain their intersection,
which is again such a subgroup, so $A_H(\phi)$ does not depend on $H$; call it
$\lambda(\phi)$. For $\phi,\psi\in V$ a single $H$ fixes both, so $\lambda$ is
linear on $V$; extend it linearly to $S$. Let $k\in K$ and $\phi\in V$, and let
$N$ be a normal subgroup of finite index fixing $\phi$ and $\sigma(k)\phi$, for
instance the intersection of the conjugates of a subgroup of finite index fixing
both. Then $\lambda(\sigma(k)\phi)=[K:N]^{-1}\sum_ql_0(\sigma(qk)\phi)$, and
$qN\mapsto qkN$ permutes $K/N$, so this is $\lambda(\phi)$.
\end{proof}

\begin{lemma}\label{asm-theta-agree-finite-count}
Let $G$ be a group with subgroups $K,P\subset U$ such that every $g\in U$ is $pk$
with $p\in P$ and $k\in K$, every $h\in U$ has $h^{-1}\in KhK$, and
$K_h=K\cap hKh^{-1}$ has finite index in $K$ for every $h\in U$. Put $H=P\cap K$,
and for subgroups $H_1$ let $[H:H_1]$ be the index of $H\cap H_1$ in $H$, taken to
be $0$ if it is infinite. Let $f$ and $f'$ be functions on $G$ with
$f(1)=f'(1)=1$, $f(pgk)=f(p)f(g)$ and $f'(pgk)=f'(p)f'(g)$ for $p\in P$,
$g\in U$ and $k\in K$, and
\[
f(p)f'(p)\,[H:pHp^{-1}]=[H:p^{-1}Hp]\qquad(p\in P).
\]
For a function $\phi$ on $G$ and $h\in U$ put
$A_\phi(h)=[K:K_h]^{-1}\sum_{q\in K/K_h}\phi(qh)$, over fixed representatives $q$.
Then $A_f(h)=A_{f'}(h)$ for every $h\in U$.
\end{lemma}
\leanref{Development/WeilBound/ThetaLifts/SphericalLattices/FunDefs.lean}{PositivityTheory.Sph.avgK}
\leanref{Development/WeilBound/ThetaLifts/SphericalLattices/FunDefs.lean}{PositivityTheory.Sph.CountSt}
\leanref{Development/WeilBound/ThetaLifts/SphericalLattices/FunCount.lean}{PositivityTheory.Sph.countSt}
\begin{proof}
Taking $p=1$ shows that $f$ and $f'$ are right $K$-invariant on $U$, and taking
$g=p^{-1}$, $k=1$ shows $f(p)f(p^{-1})=1$ for $p\in P$. With $p=g=1$, $f'$ is $1$
on $K$, so $f'(agb)=f'(g)$ for $a\in H$, $g\in U$ and $b\in K$.

Fix $h\in U$, and let $T\subset P/H$ be the set of cosets $pH$ with
$p\in P\cap KhK$. The map $qK_h\mapsto qhK$ is injective from $K/K_h$ to $G/K$,
with image $KhK/K$. The map $pH\mapsto pK$ is injective from $T$ to $G/K$, with
the same image: each $qh$ is $pk$ with $p\in P$ and $k\in K$, and then
$p\in P\cap KhK$. So $T$ is finite, and since $f$ is right $K$-invariant,
$\sum_{q\in K/K_h}f(qh)=\sum_{pH\in T}f(p)$, and the same holds for $f'$. As
$h^{-1}\in KhK$, $Kh^{-1}K=KhK$, so the set $T$ for $h^{-1}$ is $T$.

For $y\in P$ let $O(y)=\{ayH:a\in H\}\subset P/H$, the orbit of $yH$ under left
multiplication by $H$. It has $[H:yHy^{-1}]$ elements, and $O(ayb)=O(y)$ for
$a,b\in H$. For $e\in P$ put $O'(e)=O(e^{-1})$; then $e'H\in O'(e)$ if and only
if $eH\in O'(e')$. If $eH\in T$ with $e=xhy$, $x,y\in K$, then
$ae^{-1}=ay^{-1}h^{-1}x^{-1}$ lies in $P\cap Kh^{-1}K$ for $a\in H$, so
$O'(e)\subset T$; in particular these orbits are finite. Counting the pairs
$(dH,eH)$ in $T\times T$ with $eH\in O'(d)$, with weight $f'(d)/|O'(d)|$, gives
\[
\sum_{dH\in T}f'(d)=\sum_{eH\in T}\ \sum_{dH\in O'(e)}\frac{f'(d)}{|O'(d)|}.
\]
The inner sum has $|O'(e)|=[H:e^{-1}He]$ terms. For $dH\in O'(e)$, $d=ae^{-1}b$
with $a,b\in H$, so $f'(d)=f'(e^{-1})$ and $O'(d)=O(b^{-1}ea^{-1})=O(e)$ has
$[H:eHe^{-1}]$ elements. The inner sum is therefore
$[H:e^{-1}He]\,f'(e^{-1})/[H:eHe^{-1}]$. By the hypothesis at $p=e^{-1}$ and
$f(e)f(e^{-1})=1$, this is $f(e)$. Hence $\sum_Tf'=\sum_Tf$, and
$A_f(h)=A_{f'}(h)$.
\end{proof}

\begin{lemma}\label{asm-theta-agree-finite-spherical}
In Situation~\ref{asm-theta-agree-finite-inert-situation}, suppose that
$\omega_v(1,k)\mathrm{vac}_v=\mathrm{vac}_v$ and
$\omega'_v(1,k)\mathrm{vac}'_v=\mathrm{vac}'_v$ for every $k\in K_W(v)$. Then
there are linear forms $l$ on the functions on $F_v^{I_0}$ and $l'$ on the
functions on $F_v^{I'_0}$ such that: $l$ vanishes on $\mathfrak a_v(\omega)$,
$l(\omega_v(1,k)\phi)=l(\phi)$ for $k\in K_W(v)$ and Schwartz-Bruhat $\phi$, and
$l(\mathrm{vac}_v)=1$; $l'$ has the same properties for $\omega'$, $\omega'_v$ and
$\mathrm{vac}'_v$; and
$l(\omega_v(1,h)\mathrm{vac}_v)=l'(\omega'_v(1,h)\mathrm{vac}'_v)$ for every
$h\in U_W^{(v)}$.
\end{lemma}
\leanref{Development/WeilBound/ThetaLifts/SphericalLattices/HoweInert.lean}{PositivityTheory.Sph.IsSphForm}
\leanref{Development/WeilBound/ThetaLifts/SphericalLattices/HoweInert.lean}{PositivityTheory.Sph.SphFun}
\leanref{Development/WeilBound/ThetaLifts/SphericalLattices/FunAsm.lean}{PositivityTheory.Sph.exists_form}
\leanref{Development/WeilBound/ThetaLifts/SphericalLattices/FunAsm.lean}{PositivityTheory.Sph.sphFun_of}
\leanref{Development/WeilBound/ThetaLifts/SphericalLattices/FunTop.lean}{PositivityTheory.Sph.sphFun_of_modSt}
\begin{proof}
Let $w$ be the place of $E$ above $v$ and $K=K_W(v)$. As $E/F$ is unramified at
the finite places and $E$ has no $F$-embedding into $F_v$, $v$ is inert in $E$
and $E_w$ is the unramified quadratic extension of $F_v$. The rank $d=a+a'$ is
even; write $d=2m$.

\emph{Characters.} Let $z\in\mathbf A_E^\times$ with $z-1$ at $v$. As $v\notin T$,
$\alpha$ and $\beta$ are unramified at $w$. On $\mathbf A_F^\times$ the
characters $\alpha\beta$ and $\alpha/\beta$ are $\eta^{a+a'}$ and
$\eta^{a'-a}$, which are trivial as $a+a'$ is even. Both are unramified at $w$,
and $E_w^\times=F_v^\times\mathcal O_{E_w}^\times$ as $w$ is inert, so both are
trivial at $w$. Hence $\alpha(z)\beta(z)=1$ and $\alpha(z)=\beta(z)$; in
particular the characters $\beta/\alpha$ and $\alpha/\beta$ of the second spaces
of the two pairs are $1$ at $z$.

\emph{Forms.} Let $C$ be an adapted Lagrangian of $W$ at $v$ as in
Lemma~\ref{asm-theta-agree-finite-inputs}(5), with parabolic $P_C$. By
Lemma~\ref{asm-theta-agree-finite-inputs}(6) there are evaluation forms $l_0$
and $l'_0$ along $C$ for the two pairs with $l_0(\mathrm{vac}_v)\neq0$ and
$l'_0(\mathrm{vac}'_v)\neq0$, and by the characters, $l_0$ is invariant under
the $\omega_v(g,1)$, $g\in U_V^{(v)}$, and $l'_0$ under the $\omega'_v(g,1)$,
$g\in U_{V'}^{(v)}$, on Schwartz-Bruhat functions.

Apply Lemma~\ref{asm-theta-agree-finite-average} to $K$ acting by
$\sigma(k)=\omega_v(1,k)$, with $V$ the Schwartz-Bruhat functions: each of them
is fixed by an open subgroup of $K$, of finite index as $K$ is compact
(Lemma~\ref{asm-theta-agree-finite-inputs}(3)). This gives a $K$-invariant
$\lambda$. As $\omega_v(1,q)$ commutes with $\omega_v(g,1)$ and $l_0$ is invariant
under $\omega_v(g,1)$, the formula for $\lambda$ with a subgroup fixing both
$\phi$ and $\omega_v(g,1)\phi$ gives $\lambda(\omega_v(g,1)\phi)=\lambda(\phi)$,
so $\lambda$ vanishes on $\mathfrak a_v(\omega)$. With $H=K$, which fixes
$\mathrm{vac}_v$, $\lambda(\mathrm{vac}_v)=l_0(\mathrm{vac}_v)$. Put
$l=\lambda/l_0(\mathrm{vac}_v)$. For $h\in U_W^{(v)}$ the subgroup
$K_h=K\cap hKh^{-1}$ has finite index (Lemma~\ref{asm-theta-agree-finite-inputs}(3))
and fixes $\omega_v(1,h)\mathrm{vac}_v$: for $k=hk_1h^{-1}$ with $k_1\in K$,
$\omega_v(1,k)\omega_v(1,h)\mathrm{vac}_v=\omega_v(1,hk_1)\mathrm{vac}_v=\omega_v(1,h)\mathrm{vac}_v$.
Hence $l(\omega_v(1,h)\mathrm{vac}_v)=A_f(h)$ in the notation of
Lemma~\ref{asm-theta-agree-finite-count}, where
$f(g)=l_0(\omega_v(1,g)\mathrm{vac}_v)/l_0(\mathrm{vac}_v)$ for $g\in U_W^{(v)}$.
In the same way $l'$, built from $l'_0$, has the stated properties, and
$l'(\omega'_v(1,h)\mathrm{vac}'_v)=A_{f'}(h)$ with $f'$ built from $l'_0$.

\emph{The parabolic.} Let $p\in P_C$ with $pC=Ct$. Then $t_w$ is invertible: if
$t_ws=0$, then $Cs=p^{-1}Cts=0$, so $s=0$ as $C$ is a Lagrangian. So
$p^{-1}C=Ct'$ with $t'_w=t_w^{-1}$, hence $p^{-1}\in P_C$ and $P_C=P_C^\times$.
Let $z\in\mathbf A_E^\times$ be $\det t_w$ at $w$ and $1$ elsewhere, and
$u=|N_{E_w/F_v}(\det t_w)|$. By Lemma~\ref{asm-theta-agree-finite-inputs}(6),
where the first pair has $\chi_X=\beta$ and $n=a$ and the second has
$\chi_X=\alpha$ and $n=a'$,
\[
l_0(\omega_v(1,p)\phi)=\beta(z)u^{a/2}l_0(\phi),\qquad
l'_0(\omega'_v(1,p)\phi')=\alpha(z)u^{a'/2}l'_0(\phi').
\]
For $g\in U_W^{(v)}$ and $k\in K$, as $\mathrm{vac}_v$ is fixed by $K$,
$f(pgk)=\beta(z)u^{a/2}f(g)$; with $g=1$ this is $f(p)$, so $f(pgk)=f(p)f(g)$,
and likewise $f'(pgk)=f'(p)f'(g)$ with $f'(p)=\alpha(z)u^{a'/2}$. As
$\alpha(z)\beta(z)=1$, $f(p)f'(p)=u^m$, and by
Lemma~\ref{asm-theta-agree-finite-inputs}(7),
$f(p)f'(p)[H:pHp^{-1}]=[H:p^{-1}Hp]$ for $H=P_C\cap K$. Also $f(1)=f'(1)=1$.

\emph{Conclusion.} Apply Lemma~\ref{asm-theta-agree-finite-count} with
$G=U=U_W^{(v)}$, $K$, $P=P_C$, $f$ and $f'$: the decomposition $U=PK$ is
Lemma~\ref{asm-theta-agree-finite-inputs}(5), the condition $h^{-1}\in KhK$ is
part (4), and the finiteness of $[K:K_h]$ is part (3). So $A_f(h)=A_{f'}(h)$,
that is, $l(\omega_v(1,h)\mathrm{vac}_v)=l'(\omega'_v(1,h)\mathrm{vac}'_v)$.
% Lean: sphFun_of countSt avgSt pInvSt unitLiftSt charSt evInvSt siegSt modSt2 cartanSt2 finSt,
%   with inert_of_local, ThKH.Iwas.iwasawa, ThKH.Evl.evVac and stabiliser_compact; the forms are
%   exists_form (with stabK and relIndex_ne_zero_of_nhds), the products are fn_mul and fn_val.
%   The paragraph on the parabolic is pInvSt and unitLiftSt, the characters are charSt
%   (locIdele_eq_of_inert, beta_eq_alpha_of_inert).
\end{proof}

\begin{lemma}\label{asm-theta-agree-finite-inert}
Let $(B,G,\mathcal X,\omega,\omega')$ be a triple
(Definition~\ref{asm-theta-agree-triple}), and let $v\notin T$ be a finite place
of $F$ such that $E$ has no $F$-embedding into $F_v$. Then there are a complex
vector space $\Pi_v$ and linear maps $p_v$ and $p'_v$ satisfying the conditions
at $v$ of Definition~\ref{asm-theta-agree-matched}(3).
\end{lemma}
\leanref{Development/WeilBound/CoreStatements/CompactAveragingSplitting/EigenvectorLifting/FinSplit.lean}{PositivityTheory.ThK.FinInert}
\leanref{Development/WeilBound/ThetaLifts/SphericalLattices/FinInert.lean}{PositivityTheory.Sph.finAt_sph}
\leanref{Development/WeilBound/ThetaLifts/SphericalLattices/FinInert.lean}{PositivityTheory.Sph.finInert_of}
\begin{proof}
By Lemma~\ref{asm-theta-agree-finite-inputs}(2) the two Weil pairs of the triple
have local operators $(\omega_v,\rho_v)$ and $(\omega'_v,\rho'_v)$ at $v$, so
Situation~\ref{asm-theta-agree-finite-inert-situation} applies. By
Lemma~\ref{asm-theta-agree-finite-vacuum}(1) for both pairs,
Lemma~\ref{asm-theta-agree-finite-spherical} applies and gives $l$ and $l'$. For
a function $\phi$ on $F_v^{I_0}$ let $\kappa(\phi)$ be the function
$x\mapsto l(\omega_v(1,x)\phi)$ on $U_W^{(v)}$, and define $\kappa'(\phi')$ from
$l'$ and $\omega'_v$. Put $\Omega=\kappa(\mathrm{vac}_v)$, so that
$\Omega=\kappa'(\mathrm{vac}'_v)$ and $\Omega(1)=1$, and let $\Pi_v$ be the span
of the right translates $R(h)\Omega\colon x\mapsto\Omega(xh)$,
$h\in U_W^{(v)}$.

For Schwartz-Bruhat $\phi$ and $h\in U_W^{(v)}$,
$\kappa(\omega_v(1,h)\phi)=R(h)\kappa(\phi)$, since
$\omega_v(1,x)\omega_v(1,h)\phi=\omega_v(1,xh)\phi$. The space
$\mathfrak a_v(\omega)$ is stable under $\omega_v(1,h)$, which commutes with the
$\omega_v(g,1)$, and $l$ vanishes on it, so $\kappa$ vanishes on
$\mathfrak a_v(\omega)$. By Lemma~\ref{asm-theta-agree-finite-vacuum}(2),
$\kappa(\phi)\in\Pi_v$ for every Schwartz-Bruhat $\phi$, and as
$\kappa(\omega_v(1,h)\mathrm{vac}_v)=R(h)\Omega$, $\kappa$ maps the
Schwartz-Bruhat functions onto $\Pi_v$. The same holds for $\kappa'$, with the
same $\Pi_v$ since $\kappa'(\mathrm{vac}'_v)=\Omega$. Let $\pi$ be a linear
projection of the functions on $U_W^{(v)}$ onto $\Pi_v$, and put
$p_v=\pi\circ\kappa$ and $p'_v=\pi\circ\kappa'$; on Schwartz-Bruhat functions
$p_v=\kappa$ and $p'_v=\kappa'$.

The conditions at $v$ hold. The maps are onto, and they vanish on the
augmentation spaces. If $\omega(1,h)$ and $\omega'(1,h)$ act at $v$ by
$\phi\mapsto\phi_1$ and $\phi'\mapsto\phi'_1$ and $\kappa(\phi)=\kappa'(\phi')$,
then $\phi_1=\omega_v(1,h)\phi$ and $\phi'_1=\omega'_v(1,h)\phi'$ by
Lemma~\ref{asm-theta-agree-finite-unique}, so
$\kappa(\phi_1)=R(h)\kappa(\phi)=R(h)\kappa'(\phi')=\kappa'(\phi'_1)$. The space
$\Pi_v$ is not $0$ as $\Omega(1)=1$. The vacuum condition holds, as
$p_v(\mathrm{vac}_v)=\Omega=p'_v(\mathrm{vac}'_v)\neq0$, and the zeta condition
concerns only places of $T$.

For the irreducibility let $N\subset\Pi_v$ be a nonzero subspace with
$\kappa(\phi_1)\in N$ whenever $\kappa(\phi)\in N$ and $\omega(1,h)$ acts at $v$
by $\phi\mapsto\phi_1$, $h\in U_W^{(v)}$; then $\kappa(\omega_v(1,h)\phi)\in N$
whenever $\kappa(\phi)\in N$. Let $0\neq\kappa(\phi)\in N$, and choose $h$ with
$l(\omega_v(1,h)\phi)\neq0$; put $\psi=\omega_v(1,h)\phi$, so $\kappa(\psi)\in N$.
By Lemma~\ref{asm-theta-agree-finite-vacuum}(3) there are
$k_0,\ldots,k_m\in K_W(v)$ and $c_0\in\mathbf C$ with
$\sum_i\omega_v(1,k_i)\psi-c_0\mathrm{vac}_v\in\mathfrak a_v(\omega)$. Applying
$l$, which is $K_W(v)$-invariant and vanishes on $\mathfrak a_v(\omega)$, gives
$(m+1)l(\psi)=c_0$, so $c_0\neq0$. Applying $\kappa$ gives
$\sum_i\kappa(\omega_v(1,k_i)\psi)=c_0\Omega$, and the left side lies in $N$. So
$\Omega\in N$, hence every $R(h)\Omega=\kappa(\omega_v(1,h)\mathrm{vac}_v)$ lies
in $N$, and $N=\Pi_v$.
\end{proof}

\begin{lemma}\label{asm-theta-agree-finite}
Let $(B,G,\mathcal X,\omega,\omega')$ be a triple for a theta base $B$ of a CM
field $E$ (Definition~\ref{asm-theta-agree-triple}). Then there are finite
maps matched to it (Definition~\ref{asm-theta-agree-matched}(3)).
\end{lemma}
\leanref{Development/AbelianVarietyTheory/Definitions/Divisors/MatchingOfParts.lean}{GlobalTheory.ThK.KFin}
\leanref{Development/WeilBound/CoreStatements/CompactAveragingSplitting/EigenvectorLifting/FinSplit.lean}{PositivityTheory.ThK.kFin_of}
\begin{proof}
Each condition on matched finite maps concerns a single finite place $v$. So
it suffices to give, for each $v$, a space $\Pi_v$ and maps $p_v$ and $p'_v$
satisfying the conditions at $v$, and to take these as the components of
$R_f$. There are three cases: $v\in T$; $v\notin T$ and $E$ has an
$F$-embedding into $F_v$, so that $v$ splits in $E$; and $v\notin T$ and $E$
has none, so that $v$ is inert in $E$, as $E/F$ is unramified. In these three
cases such data are given by Lemmas~\ref{asm-theta-agree-finite-ramified},
\ref{asm-theta-agree-finite-split} and~\ref{asm-theta-agree-finite-inert}.
% Lean: kFin_of FinRam FinSplit FinInert, with FinRam =
% PositivityTheory.BernsteinRamification.Close.finRam, FinSplit = Close.finSplit and
% FinInert = PositivityTheory.Sph.finInert_of (vacInert howeInert) (sphFun_of_modSt modSt2)
% PositivityTheory.ThKH.Locrep.locRepEx. Lemma asm-theta-agree-finite-inputs is now proved
% in full.
\end{proof}

\subsection{Decay and the coinvariants}\label{subsection-asm-decay-coinvariants}

\begin{lemma}\label{asm-theta-agree-decay-full}
Let $(X,W,M,\chi_X,\chi_W,\omega)$ be a Weil pair for a theta base $B$ of a CM
field $E$ (Definition~\ref{asm-theta-agree-triple}(1)), and let
$\Phi\in\mathcal S^\circ(M)$ (Definition~\ref{asm-theta-agree-models}(5)). For
$g\in U_X$ let $\overline{\operatorname{ht}}(g)$ be the infimum of
$\operatorname{ht}(\gamma g)$ over $\gamma\in U_X(F)$. Then for every $N\geq0$
there are $C$ and $k$ with
\[
|\theta_\omega(\Phi)(g,h)|\,\overline{\operatorname{ht}}(g)^N\leq C\operatorname{ht}(h)^k
\]
for all $g\in U_X$ and $h\in U_W$.
\end{lemma}
\leanref{Development/AbelianVarietyTheory/Definitions/Divisors/CTSmallCone.lean}{GlobalTheory.ThK.SoDecay}
\leanref{Development/WeilBound/ThetaLifts/FrameEstimates/WeilConv.lean}{PositivityTheory.ThKEst.soDecay}
\leanref{Development/WeilBound/ThetaLifts/FrameEstimates/ThickFrames.lean}{PositivityTheory.ThKEst.soDecay_of}
\begin{proof}
By Lemma~\ref{asm-theta-agree-continuity-pieces}(1),
$\Phi=\sum_{i=1}^mf_i\otimes\phi_i$ with full-rank finite parts $\phi_i$
(Definition~\ref{asm-theta-agree-continuity-sums}(4)), and
$\theta_\omega(\Phi)=\sum_i\theta_\omega(f_i\otimes\phi_i)$ by
Lemma~\ref{asm-theta-agree-convergence-theta}(2). Let $k$ be given by
Lemma~\ref{asm-theta-agree-continuity-decay} for $N$, and $p_i$ the seminorm
given there for $\phi_i$. Then
\[
|\theta_\omega(\Phi)(g,h)|\,\overline{\operatorname{ht}}(g)^N\leq\sum_i|\theta_\omega(f_i\otimes\phi_i)(g,h)|\,\overline{\operatorname{ht}}(g)^N\leq\Bigl(\sum_ip_i(f_i)\Bigr)\operatorname{ht}(h)^k .
\]
\end{proof}

\begin{lemma}\label{asm-theta-agree-decay}
Let $(B,G,\mathcal X,\omega,\omega')$ be a triple for a theta base $B$ of a CM
field $E$ (Definition~\ref{asm-theta-agree-triple}). Then it has decay
(Definition~\ref{asm-theta-agree-parts}(1)).
\end{lemma}
\leanref{Development/AbelianVarietyTheory/Definitions/Divisors/MatchingOfParts.lean}{GlobalTheory.ThK.KDecay}
\leanref{Development/WeilBound/ThetaLifts/FrameEstimates/WeilConv.lean}{PositivityTheory.ThKEst.kDecay}
\leanref{Development/WeilBound/ThetaLifts/FrameEstimates/ThickFrames.lean}{PositivityTheory.ThKEst.kDecay_of}
\leanref{Development/WeilBound/ThetaLifts/FrameEstimates/SquareIntegrableWeilRange.lean}{PositivityTheory.ThKEst.TriplePair}
\leanref{Development/WeilBound/ThetaLifts/FrameEstimates/InitialPairs/TriplePairHolds.lean}{PositivityTheory.ThKEst.Sq.triplePair_holds}
\begin{proof}
Let the triple have ranks $(a,a')$, integers $p$ and $q$, and models $M$ and
$M'$. By Definition~\ref{asm-theta-agree-triple}(4),
$(V,W,M,\beta,\beta/\alpha,\omega)$ is a Weil pair of signature
$(a-p,p;a+a'-p-q,p+q)$, and $(V',W,M',\alpha,\alpha/\beta,\omega')$ is one of
signature $(q,a'-q;a+a'-p-q,p+q)$. As $G$ is proper,
$2(p+q)+1<\lfloor(a-a')/2\rfloor$, so $4p<a-a'$ and
\[
(a+a')+2\min(a-p,p)\leq a+a'+2p<2a .
\]
So Lemma~\ref{asm-theta-agree-convergence}(1) applies to the first Weil pair,
with $n=a$ and $d=a+a'$, and every Schwartz-Bruhat function converges for
$(M,\omega)$.

Let $\Phi'\in\mathcal S^\circ(M')$, and let $\mu'$ and $\mathcal D'$ be the
measure and the domain of $M'$. By
Lemma~\ref{asm-theta-agree-continuity-pieces}(1),
$\Phi'=\sum_if_i\otimes\phi_i$ with full-rank finite parts $\phi_i$ for $M'$.
Lemma~\ref{asm-theta-agree-continuity-domination} for the second Weil pair
gives the bounds required in Lemma~\ref{asm-theta-agree-continuity-integrals}
for the $\phi_i$. By Lemma~\ref{asm-theta-agree-continuity-integrals}(1),
$\Phi'$ converges for $(M',\omega')$, and there are $C$ and $k$ with
\[
|I'(\Phi')(h)|\leq\int_{\mathcal D'}|\theta_{\omega'}(\Phi')(g,h)|\,d\mu'(g)\leq C\operatorname{ht}(h)^k
\]
for every $h\in U_W$. By Lemma~\ref{asm-theta-agree-decay-full} for the second
Weil pair, for every $N$ there are $C$ and $k$ with
$|\theta_{\omega'}(\Phi')(g,h)|\,\overline{\operatorname{ht}}(g)^N\leq C\operatorname{ht}(h)^k$
for all $g\in U_{V'}$ and $h\in U_W$, where $\overline{\operatorname{ht}}(g)$
is the infimum of $\operatorname{ht}(\gamma g)$ over $\gamma\in U_{V'}(F)$.
These are the conditions of Definition~\ref{asm-theta-agree-parts}(1).
\end{proof}
% Lean: PositivityTheory.ThKEst.kDecay hSiegelSet ThKF0.thetaBasic = kDecay_of thetaBasic
% Sq.triplePair_holds (weilConv hSiegelSet thetaBasic) (soDecay hSiegelSet) (domSoU
% hSiegelSet) soFin_of, with hSiegelSet = ThKR.siegel_of_core Prw.siegelCore_holds. The
% inequality is the third component of triplePair_holds, from the field eBig of
% Geo.IsGood. Convergence and moderate growth of I' come from dom_sum with
% dom_of_domSoU (Lemmas asm-theta-agree-continuity-domination and -integrals); the decay
% with the quotient height is soDecay (Lemma asm-theta-agree-decay-full).

\begin{definition}\label{asm-theta-agree-coinvariants-local}
Let $k$ be a nonarchimedean local field, let $n,d\geq0$, and let $\chi_1$ and
$\chi_2$ be characters of $k^\times$ that are \emph{unitary and smooth}:
$|\chi_i(x)|=1$ for all $x$, and $\ker\chi_i$ is open.
\begin{enumerate}
\item For $g\in GL_n(k)$, $h\in GL_d(k)$ and a function $\Psi$ on
  $M_{n\times d}(k)$, put
  \[
  \omega_{\chi_1,\chi_2}(g,h)\Psi(Y)=\chi_1(\det g)\,\chi_2(\det h)\,|\det g|^{d/2}|\det h|^{n/2}\,\Psi(g^{\mathsf T}Yh).
  \]
\item The \emph{local augmentation space} $\mathfrak a(\chi_1,\chi_2)$ is the
  span of the $\omega_{\chi_1,\chi_2}(g,1)\Psi-\Psi$ with $g\in GL_n(k)$ and
  $\Psi$ Schwartz-Bruhat on $M_{n\times d}(k)$. A Schwartz-Bruhat function
  $\Psi$ on $M_{n\times d}(k)$ has \emph{full rank} if $\Psi(Y)=0$ for every $Y$
  of rank different from $n$.
\item Let $E$ be a CM field with maximal totally real subfield $F$, let $X$ and
  $W$ be lattice spaces of ranks $n$ and $d$
  (Definition~\ref{asm-theta-agree-doubling}(1)), let
  $I=\{1,\ldots,n\}\times\{1,\ldots,d\}$, let $\omega$ be a representation of
  $U_X\times U_W$ on the functions on $\mathbf A_F^I$, and let $v$ be a finite
  place of $F$, with $k=F_v$. A \emph{dictionary} at $v$ with
  characters $(\chi_1,\chi_2)$ is a linear map $\iota$ from functions on $F_v^I$
  to functions on $M_{n\times d}(F_v)$, bijective between the Schwartz-Bruhat
  functions, such that:
  \begin{enumerate}
  \item[(a)] for every $g\in U_X^{(v)}$ there is $g_0\in GL_n(F_v)$, and for every
    $g_0\in GL_n(F_v)$ there is $g\in U_X^{(v)}$, such that for every
    Schwartz-Bruhat $\phi$ there is a Schwartz-Bruhat $\phi_1$ with
    $\omega(g,1)$ acting at $v$ by $\phi\mapsto\phi_1$
    (Definition~\ref{asm-theta-agree-models}(2)) and
    $\iota\phi_1=\omega_{\chi_1,\chi_2}(g_0,1)\iota\phi$;
  \item[(b)] the same holds for $h\in U_W^{(v)}$, $h_0\in GL_d(F_v)$,
    $\omega(1,h)$ and $\omega_{\chi_1,\chi_2}(1,h_0)$.
  \end{enumerate}
\item For a Hecke character $\chi$ of $E$ and an $F$-embedding
  $t\colon E\to F_v$, $\chi_t$ is the character of $F_v^\times$ given by $\chi$
  at the place of $E$ defined by $t$.
\end{enumerate}
\end{definition}
\leanref{Development/AbelianVarietyTheory/Definitions/Divisors/RankEqCardIffExistsFrameMatrixFirstColumnBlock.lean}{GlobalTheory.SplitFlip.omB}
\leanref{Development/WeilBound/CoreStatements/LocallyConstantCoinvariants/RightOnto.lean}{PositivityTheory.KCo.USm}
\leanref{Development/WeilBound/CoreStatements/LocallyConstantCoinvariants/RightOnto.lean}{PositivityTheory.KCo.aug}
\leanref{Development/WeilBound/CoreStatements/LocallyConstantCoinvariants/RightOnto.lean}{PositivityTheory.KCo.Full}
\leanref{Development/WeilBound/CoreStatements/LocallyConstantCoinvariants/RightOnto.lean}{PositivityTheory.KCo.IsDict}
\leanref{Development/AbelianVarietyTheory/BlockTheory/DivisorBase/WeilTensorModels.lean}{GlobalTheory.IndefiniteRowSideFullHecke.WeilTensorModels.chrW}

\begin{lemma}\label{asm-theta-agree-coinvariants-step}
Let $k$ be a nonarchimedean local field, let $n,d\geq0$, and let $\chi_1,\chi_2$
be characters of $k^\times$ with $\chi_1$ unitary and smooth and $\chi_1\neq1$
(Definition~\ref{asm-theta-agree-coinvariants-local}). For
$Y\in M_{n\times d}(k)$ let $Y_1,\ldots,Y_n\in k^d$ be its rows, and for
$S\subset\{1,\ldots,n\}$ call $Y$ \emph{free on} $S$ if the rows $Y_s$,
$s\in S$, are linearly independent. Let $S\subset\{1,\ldots,n\}$ and
$i\notin S$, and let $\Phi$ be a Schwartz-Bruhat function on
$M_{n\times d}(k)$ with $\Phi(Y)=0$ for every $Y$ that is not free on $S$.
Then there is a Schwartz-Bruhat function $\Psi$ with $\Psi(Y)=0$ for every $Y$
that is not free on $S\cup\{i\}$ and with
$\Phi-\Psi\in\mathfrak a(\chi_1,\chi_2)$.
\end{lemma}
\leanref{Development/WeilBound/CoreStatements/LocallyConstantCoinvariants/PfSurj.lean}{PositivityTheory.KCo.step}
\leanref{Development/WeilBound/CoreStatements/LocallyConstantCoinvariants/PfSurj.lean}{PositivityTheory.KCo.exists_gl}
\leanref{Development/WeilBound/CoreStatements/LocallyConstantCoinvariants/PfSurj.lean}{PositivityTheory.KCo.exists_inv}
\leanref{Development/WeilBound/CoreStatements/LocallyConstantCoinvariants/PfSurj.lean}{PositivityTheory.KCo.exists_bound}
\begin{proof}
Write $|\cdot|$ for the normalized absolute value of $k$ and
$\mathfrak a=\mathfrak a(\chi_1,\chi_2)$. For $x\in k$ and $a\in k^S$ let
$m_{x,a}(Y)$ be $Y$ with its row $Y_i$ replaced by $xY_i+\sum_{s\in S}a_sY_s$.
As $i\notin S$, $m_{x,a}\circ m_{y,b}=m_{xy,xb+a}$ and $m_{1,0}$ is the
identity.

\emph{Moves.} Let $x\in k^\times$ and $a\in k^S$, and let $g\in GL_n(k)$ be the
transpose of the identity matrix with its $i$-th row replaced by
$xe_i+\sum_sa_se_s$. Then $\det g=x$ and $g^{\mathsf T}Y=m_{x,a}(Y)$, so
$\omega_{\chi_1,\chi_2}(g,1)\Phi'=z_x\,\Phi'\circ m_{x,a}$ with
$z_x=\chi_1(x)|x|^{d/2}$, for every function $\Phi'$. Hence, for
Schwartz-Bruhat $\Phi'$, $z_x\,\Phi'\circ m_{x,a}$ is Schwartz-Bruhat and
$z_x\,\Phi'\circ m_{x,a}-\Phi'\in\mathfrak a$. For a finite family
$(a_q)_{q\in Q}$, summing these with $x=1$ gives
$\sum_q\Phi'\circ m_{1,a_q}-|Q|\,\Phi'\in\mathfrak a$.

\emph{Boxes.} For $c\in k$ let $B(c)$ be the group of the $a\in k^S$ with
$|a_s|\leq|c|$ for every $s$. It is compact, and open if $c\neq0$. There is
$c_0\neq0$ with $\Phi(m_{1,a}(Y))=\Phi(Y)$ for all $a\in B(c_0)$ and all $Y$:
as $\Phi$ is locally constant with compact support, this holds for small $a$
and $Y$ in the support of $\Phi$; and if $\Phi(Y)=0\neq\Phi(m_{1,a}(Y))$,
applying it to $m_{1,a}(Y)$ and $-a$ gives $\Phi(Y)=\Phi(m_{1,a}(Y))$. The
support of $\Phi$ is compact, and its elements are free on $S$. Near a point
$Y^0$ of it, choose a $d\times S$ matrix $C$ with $R_{Y^0}C=1$, where $R_Y$ is
the $S\times d$ matrix of the rows $Y_s$, $s\in S$; then $\det(R_YC)\neq0$ near
$Y^0$, and the unique $b$ with $Y_i=\sum_sb_sY_s$, if it exists, is
$Y_iC(R_YC)^{-1}$, a continuous function of $Y$. By compactness these $b$ lie in
a compact set $L$, for $Y$ in the support of $\Phi$. Choose $c\neq0$ with
$L\subset B(c)$.

\emph{The function $\Psi$.} As $\chi_1\neq1$, there is $x\in k^\times$ with
$\chi_1(x)\neq1$, and replacing $x$ by $x^{-1}$ we may take $|x|\leq1$; then
$B(c)\subset B(x^{-1}c)$. Let $(a_q)$ be representatives of the finitely many
classes of $B(c)$ modulo $B(c)\cap B(c_0)$, and $(b_r)$ of those of
$B(x^{-1}c)$ modulo $B(c)$, with $N_1$ and $N_2$ elements. Put
\[
\Phi_1=\sum_q\Phi\circ m_{1,a_q},\qquad\Phi_2=\sum_r\Phi_1\circ m_{1,b_r},\qquad
\Phi_3=z_x\,\Phi_1\circ m_{x,0}.
\]
By the moves these are Schwartz-Bruhat, and $\Phi_1-N_1\Phi$, $\Phi_2-N_2\Phi_1$
and $\Phi_3-\Phi_1$ lie in $\mathfrak a$. If $z_xN_2=1$, then taking absolute
values, $|x|^{d/2}N_2=1$ as $\chi_1$ is unitary, and so $\chi_1(x)=1$; hence
$z_xN_2\neq1$. Put $\Psi=((z_xN_2-1)N_1)^{-1}(z_x\Phi_2-\Phi_3)$. It is
Schwartz-Bruhat, and
\[
\Phi-\Psi=-((z_xN_2-1)N_1)^{-1}\bigl(z_x(\Phi_2-N_2\Phi_1)-(\Phi_3-\Phi_1)+(z_xN_2-1)(\Phi_1-N_1\Phi)\bigr)\in\mathfrak a .
\]

\emph{Vanishing.} Let $Y$ be not free on $S\cup\{i\}$. It suffices that
$\Phi_2(Y)=\Phi_1(m_{x,0}(Y))$. If $Y$ is not free on $S$, then neither is any
$m_{y,a}(Y)$, which has the same rows $Y_s$, $s\in S$; so both sides are $0$.
Otherwise $Y_i=\sum_sb_sY_s$ for some $b\in k^S$, and $Y=m_{0,b}(Y)$. Put
$f(a)=\Phi(m_{0,a}(Y))$ and $f_1(a)=\Phi_1(m_{0,a}(Y))=\sum_qf(a+a_q)$ for
$a\in k^S$. As $m_{1,a'}\circ m_{0,a}=m_{0,a+a'}$, $f$ is invariant under
translation by $B(c_0)$. If $f(a)\neq0$, then $m_{0,a}(Y)$ lies in the support
of $\Phi$ and its row $i$ is $\sum_sa_sY_s$, so $a\in L\subset B(c)$. Hence $f_1$
is invariant under translation by $B(c)$, being a sum over the classes modulo
$B(c)\cap B(c_0)$, and vanishes outside $B(c)$. Now
$\Phi_2(Y)=\sum_rf_1(b+b_r)$ and $\Phi_1(m_{x,0}(Y))=f_1(xb)$. If
$b\in B(x^{-1}c)$, translation by $b$ permutes the classes modulo $B(c)$, so
$\sum_rf_1(b+b_r)=\sum_rf_1(b_r)=f_1(0)$, as only the class $B(c)$ contributes;
and $xb\in B(c)$, so $f_1(xb)=f_1(0)$. If $b\notin B(x^{-1}c)$, then no
$b+b_r$ lies in $B(x^{-1}c)\supset B(c)$, and $xb\notin B(c)$, so both sides are
$0$.
\end{proof}

\begin{lemma}\label{asm-theta-agree-coinvariants-kill}
Let $Z$ be a topological space, call a function on $Z$ Schwartz-Bruhat if it is
locally constant with compact support, let $N$ be a subspace of the functions
on $Z$, and let $Y_1\in Z$. Suppose that for every $Y\neq Y_1$ there is a locally
constant function $m$ on $Z$ with $m(Y)\neq m(Y_1)$ and
$(m-m(Y_1))\Phi\in N$ for every Schwartz-Bruhat $\Phi$. Then every
Schwartz-Bruhat $\Phi$ with $\Phi(Y_1)=0$ lies in $N$.
\end{lemma}
\leanref{Development/WeilBound/CoreStatements/LocallyConstantCoinvariants/IndicatorMem.lean}{PositivityTheory.KCo.kill}
\begin{proof}
For $Y\neq Y_1$ let $m_Y$ be such a function and $U_Y=\{y:m_Y(y)\neq m_Y(Y_1)\}$,
an open and closed set containing $Y$. The support of $\Phi$ is closed in a
compact set, so compact, and does not contain $Y_1$; so it is covered by finitely
many $U_{Y^1},\ldots,U_{Y^r}$. We argue by induction on $r$; for $r=0$,
$\Phi=0$. Let $U=U_{Y^1}$ and $m=m_{Y^1}$. Then
$\Phi=\mathbf 1_U\Phi+\mathbf 1_{Z\setminus U}\Phi$. The function $\Phi'$ equal to
$\Phi/(m-m(Y_1))$ on $U$ and to $0$ elsewhere is Schwartz-Bruhat, and
$\mathbf 1_U\Phi=(m-m(Y_1))\Phi'\in N$. The function $\mathbf 1_{Z\setminus U}\Phi$
is Schwartz-Bruhat, vanishes at $Y_1$, and its support is covered by
$U_{Y^2},\ldots,U_{Y^r}$, so it lies in $N$ by induction.
\end{proof}

\begin{lemma}\label{asm-theta-agree-coinvariants-eigen}
Let $B$ be a theta base for a CM field $E$ with character $\psi$
(Definition~\ref{asm-theta-agree-models}(1)), let $X$ and $W$ be lattice spaces
of ranks $n$ and $d$ whose matrices $J_X$ and $J_W$ are invertible, let
$I=\{1,\ldots,n\}\times\{1,\ldots,d\}$, and let $M$ be a model for $(X,W)$, with
isomorphism $e=(e_1,e_2)$, Schrödinger operators $\rho$ and form
$\langle\,,\rangle$ on $\mathbb W(X,W)$
(Definition~\ref{asm-theta-agree-models}(3)). Let $v$ be a finite place of $F$,
$t\colon E\to F_v$ an $F$-embedding, and $\iota$ a linear model of the second
kind for $(t,1,1)$, with
$\operatorname{co}(u)=-t(J_X)\,\pi_{t\circ c}(u)\,t(J_W)^{\mathsf T}$
(Definition~\ref{asm-theta-agree-models}(5)). For $y\in F_v^I$ let
$\bar y\in\mathbf A^I$ have component $y$ at $v$ and $0$ elsewhere.
\begin{enumerate}
\item For every $Y\in M_{n\times d}(F_v)$ there is exactly one
  $x\in\mathbb W(X,W)$ at $v$ with $\pi_t(x)=0$ and $\operatorname{co}(x)=Y$.
\item Let $u\in\mathbb W(X,W)$ be at $v$, and for a function $\phi$ on $F_v^I$
  put
  \[
  (\rho_v(u)\phi)(y)=\psi\bigl(\bar y\cdot e_2(u)+\tfrac12e_1(u)\cdot e_2(u)\bigr)\,\phi\bigl(y+e_1(u)_v\bigr).
  \]
  Then $\rho_v(u)$ maps Schwartz-Bruhat functions to Schwartz-Bruhat functions;
  $\rho(u)(\phi\otimes r)=(\rho_v(u)\phi)\otimes r$ for every function $r$ on
  $\mathbf A^I$ that does not depend on the $v$-component, so that $\rho(u)$ acts
  at $v$ by $\phi\mapsto\rho_v(u)\phi$; and
  $\operatorname{res}_z(\rho(u)\Psi)=\rho_v(u)\operatorname{res}_z\Psi$ for every
  function $\Psi$ on $\mathbf A^I$ and every $z\in\mathbf A^I$
  (Lemma~\ref{asm-theta-agree-finite-inputs-slice}).
\item Let $l$ be a linear form on the functions on $F_v^I$, and let
  $x_1\in\mathbb W(X,W)$ be at $v$ with $\pi_t(x_1)=0$, such that
  $l(\rho_v(u)\phi)=\psi(\langle x_1,u\rangle)\,l(\phi)$ for every
  Schwartz-Bruhat $\phi$ and every $u$ at $v$ with $\pi_{t\circ c}(u)=0$. Then
  $l(\phi)=0$ for every Schwartz-Bruhat $\phi$ with
  $\iota\phi(\operatorname{co}(x_1))=0$.
\end{enumerate}
\end{lemma}
\leanref{Development/WeilBound/CoreStatements/LocallyConstantCoinvariants/SoStatement.lean}{PositivityTheory.KCo.eigen_kill}
\leanref{Development/WeilBound/CoreStatements/LocallyConstantCoinvariants/SoStatement.lean}{PositivityTheory.KCo.exists_coB}
\leanref{Development/WeilBound/CoreStatements/LocallyConstantCoinvariants/CompFinLiftE.lean}{PositivityTheory.KCo.locRho}
\leanref{Development/WeilBound/CoreStatements/LocallyConstantCoinvariants/CompFinLiftE.lean}{PositivityTheory.KCo.rho_tens}
\leanref{Development/WeilBound/CoreStatements/LocallyConstantCoinvariants/CompFinLiftE.lean}{PositivityTheory.KCo.slice_rho}
\leanref{Development/WeilBound/CoreStatements/LocallyConstantCoinvariants/CompFinLiftE.lean}{PositivityTheory.KCo.exists_sep}
\begin{proof}
The embeddings $t$ and $t\circ c$ are distinct, so
$E\otimes_FF_v\to F_v\times F_v$, $x\otimes a\mapsto(t(x)a,t(c(x))a)$, is an
isomorphism, and $z\mapsto(\pi_t(z),\pi_{t\circ c}(z))$ is a bijection from the
elements of $\mathbf A_E$ at $v$ onto $F_v\times F_v$. Note
$\pi_{t\circ c}(c(z))=\pi_t(z)$.

(1) Take $x$ at $v$ with $\pi_t(x)=0$ and
$\pi_{t\circ c}(x)=-t(J_X)^{-1}Y(t(J_W)^{\mathsf T})^{-1}$, entrywise. If $x'$ is
another, then $x-x'$ is at $v$ with $\pi_t(x-x')=0$ and, as $t(J_X)$ and $t(J_W)$
are invertible, $\pi_{t\circ c}(x-x')=0$; so $x=x'$.

(2) As $e$ is $\mathbf A$-linear and $u$ is at $v$, $e_1(u)$ and $e_2(u)$ are at
$v$. The function $s\mapsto\psi(\bar s)$ on $F_v$ is a continuous character with
values of absolute value $1$, hence locally constant; so
$y\mapsto\psi(\bar y\cdot e_2(u))$ is locally constant, and $\rho_v(u)\phi$ is
Schwartz-Bruhat when $\phi$ is. For $x\in\mathbf A^I$,
$x\cdot e_2(u)=\overline{x_v}\cdot e_2(u)$, and $r(x+e_1(u))=r(x)$ when $r$ does
not depend on the $v$-component; this gives the second claim. The third follows
in the same way from $z\langle y\rangle+e_1(u)=z\langle y+e_1(u)_v\rangle$.

(3) First, $\psi(\bar s)\neq1$ for some $s\in F_v$. The character $\psi$ is
trivial on the finite adeles that are integral at every place: it is trivial at
the finite part of every $o\in\mathcal O_F$, as $\psi(o)=1$ and
$\exp(-2\pi i\operatorname{Tr}_{F/\mathbf Q}(o))=1$, and $\mathcal O_F$ is dense
in them. The elements of $F$ integral at every finite place other than $v$ form
an $\mathcal O_F$-module that is not contained in the inverse different, so one
of them, $z_0$, has $\operatorname{Tr}_{F/\mathbf Q}(z_0)\notin\mathbf Z$. Its
finite part is an integral finite adele plus $\overline{z_{0,v}}$. If $\psi$ were
trivial on every $\bar s$, it would be trivial on this finite part, and
$1=\psi(z_0)=\exp(-2\pi i\operatorname{Tr}_{F/\mathbf Q}(z_0))$, which is false.

Next, let $z\neq0$ be at $v$ with $\pi_t(z)=0$; we find $u$ at $v$ with
$\pi_{t\circ c}(u)=0$ and $\psi(\langle z,u\rangle)\neq1$. As $e(z)\neq0$ is at
$v$, some coordinate $e_1(z)_j$ or $e_2(z)_j$ has a nonzero component $\zeta$ at
$v$. In the first case let $w=e^{-1}(0,\overline{(s/\zeta)\delta_j})$, in the
second $w=e^{-1}(\overline{(-s/\zeta)\delta_j},0)$, where $\delta_j$ is the $j$-th
standard basis vector of $F_v^I$. Then $w$ is at $v$ and
$\langle z,w\rangle=e_1(z)\cdot e_2(w)-e_1(w)\cdot e_2(z)=\bar s$. Let
$\varepsilon\in\mathbf A_E$ be at $v$ with $\pi_t(\varepsilon)=1$ and
$\pi_{t\circ c}(\varepsilon)=0$, and $u=\varepsilon w$. Then $u$ is at $v$,
$\pi_{t\circ c}(u)=0$ and $\pi_t(w-u)=0$. Each term of $h(z,w-u)$ is
$z_{ij}\,c((w-u)_{kl})$ times an element of $E$; it is at $v$, its $\pi_t$
vanishes with $\pi_t(z_{ij})$ and its $\pi_{t\circ c}$ with $\pi_t((w-u)_{kl})$,
so it is $0$. Hence $\langle z,u\rangle=\langle z,w\rangle=\bar s$.

Now let $N$ be the space of the $\iota\phi'$ with $\phi'$ Schwartz-Bruhat and
$l(\phi')=0$, and $Y_1=\operatorname{co}(x_1)$; we apply
Lemma~\ref{asm-theta-agree-coinvariants-kill} on $M_{n\times d}(F_v)$. For $y$
there let $x_y$ be given by (1). Let $Y\neq Y_1$, let $u$ be as above for
$z=x_Y-x_1$, and put $m(y)=\psi(\langle x_y,u\rangle)$. For Schwartz-Bruhat
$\phi'$, $\rho(u)$ acts at $v$ by $\phi'\mapsto\rho_v(u)\phi'$ by (2), so the
property of $\iota$ for $u$ with $\pi_{t\circ c}(u)=0$, applied at each
$x_y$, gives $\iota(\rho_v(u)\phi')=m\cdot\iota\phi'$. As $\iota$ is bijective
between the Schwartz-Bruhat spaces, $m\Phi$ is Schwartz-Bruhat for every
Schwartz-Bruhat $\Phi$; multiplying by the indicator function of an open compact
neighbourhood of a point shows that $m$ is locally constant. By uniqueness in
(1), $x_{Y_1}=x_1$ and $m(Y_1)=\psi(\langle x_1,u\rangle)\neq0$, and
$m(Y)=\psi(\langle z,u\rangle)\,m(Y_1)\neq m(Y_1)$. For $\Phi=\iota\phi'$,
$(m-m(Y_1))\Phi=\iota(\rho_v(u)\phi'-m(Y_1)\phi')$ and
$l(\rho_v(u)\phi'-m(Y_1)\phi')=\psi(\langle x_1,u\rangle)l(\phi')-m(Y_1)l(\phi')=0$,
so $(m-m(Y_1))\Phi\in N$. By Lemma~\ref{asm-theta-agree-coinvariants-kill},
$\iota\phi\in N$: $\iota\phi=\iota\phi'$ with $l(\phi')=0$. As $\iota$ is injective
on Schwartz-Bruhat functions, $\phi=\phi'$ and $l(\phi)=0$.
\end{proof}
% Lean: (3) is eigen_kill, through kill, lc_of_mul, mult' (the field mult of IsLinBar) and
% exists_sep, whose first step is GlobalTheory.PsiNontriv.psiV_ne_one (exists_trace_not_int,
% psi_int, psi_fin_rat); the vanishing of <z, w - u> is symp_iso, and the splitting of
% E (x) F_v is splitPair_bijective with exists_pair and eq_zero_of_at.

\begin{lemma}\label{asm-theta-agree-coinvariants-slices}
Let $B$ be a theta base for a CM field $E$, with place $v_1$ and embedding
$\tau_1$ (Definition~\ref{asm-theta-agree-models}(1)), let $X$ and $W$ be lattice
spaces of ranks $n$ and $d$, let $I=\{1,\ldots,n\}\times\{1,\ldots,d\}$, and let
$M$ be a regular model for $(X,W)$. Let $\Psi$ be a Schwartz-Bruhat function on
$\mathbf A^I$ such that for every $z\in\mathbf A^I$ the function
$\operatorname{res}_z\Psi\colon y\mapsto\Psi(z\langle y\rangle)$ on
$F_{v_1}^I$ (Lemma~\ref{asm-theta-agree-finite-inputs-slice}) has full rank at
$\tau_1$ (Definition~\ref{asm-theta-agree-models}(5)). Then
$\Psi\in\mathcal S^\circ(M)$.
\end{lemma}
\leanref{Development/WeilBound/CoreStatements/LocallyConstantCoinvariants/SoStatement.lean}{PositivityTheory.KCo.mem_so}
\leanref{Development/WeilBound/CoreStatements/LocallyConstantCoinvariants/IndicatorMem.lean}{PositivityTheory.KCo.schA_le}
\begin{proof}
Let $\Sigma$ be the span of the functions on $F_{v_1}^I$ of full rank at
$\tau_1$, and $A$ a linear map from the functions on $F_{v_1}^I$ onto $\Sigma$
that is the identity on $\Sigma$. For a function $\Phi$ on $\mathbf A^I$ put
$(A_*\Phi)(x)=(A\operatorname{res}_x\Phi)(x_{v_1})$. Then $A_*$ is linear, and
$A_*\Psi=\Psi$, as $\operatorname{res}_x\Psi\in\Sigma$ and
$(\operatorname{res}_x\Psi)(x_{v_1})=\Psi(x)$. A Schwartz-Bruhat function on
$\mathbf A_f^I$ is a finite sum of functions $x\mapsto\prod_w\phi_w(x_w)$ with
$\phi_w$ Schwartz-Bruhat on $F_w^I$ and equal to the indicator function of
$\mathcal O_w^I$ for almost all $w$: it is locally constant with compact
support, so a finite sum of multiples of indicator functions of products of
balls. Hence $\Psi$ is a finite sum of functions $\Phi=f\otimes\bigotimes_w\phi_w$
of this kind with $f\in\mathcal S(F_\infty^I)$, and $\phi_w=\mathrm{vac}_w$ for
almost all $w$, as $M$ is regular. For such $\Phi$, $\operatorname{res}_x\Phi$ is
$\phi_{v_1}$ times the product of the other factors at $x$, so
$A_*\Phi=f\otimes A\phi_{v_1}\otimes\bigotimes_{w\neq v_1}\phi_w$. As
$A\phi_{v_1}$ is a linear combination of functions of full rank at $\tau_1$, this
lies in $\mathcal S^\circ(M)$. Summing, $\Psi=A_*\Psi\in\mathcal S^\circ(M)$.
\end{proof}
% Lean: mem_so, with the projection linearProjOfIsCompl, the operator opL v A, opL_pure and
% pure_update_mem (ExistsNhdsInvariant.lean) and schA_le with the boxes of IndicatorMem.lean.

\begin{lemma}\label{asm-theta-agree-coinvariants-fibre}
Let $k$ be a nonarchimedean local field with ring of integers $\mathcal O$ and
normalized absolute value $|\cdot|$, let $0\leq n\leq d$ and $t=d-n$, and let
$\mu$ be a character of $k^\times$ with open kernel. Call a function on a
subspace of a space of matrices over $k$ Schwartz-Bruhat if it is locally
constant with compact support.
\begin{enumerate}
\item Let $R$ be the set of $d\times n$ matrices of rank $n$, $\mathcal S(R)$ the
  space of Schwartz-Bruhat functions on $R$, and for $g\in GL_n(k)$,
  $h\in GL_d(k)$ and $\Phi\in\mathcal S(R)$ put
  \[
  \omega^R(g,h)\Phi(x)=\mu(\det g)^{-1}|\det g|^{d/2}\,\mu(\det h)|\det h|^{-n/2}\,\Phi(h^{-1}xg).
  \]
\item Let $\mathcal S(GL_n)$ be the space of Schwartz-Bruhat functions on
  $GL_n(k)$, and for $g,p_1\in GL_n(k)$, $p_2\in GL_t(k)$ and
  $\phi\in\mathcal S(GL_n)$ put
  \[
  \tau(g,p_1,p_2)\phi(y)=\mu(\det g)^{-1}|\det g|^{d/2}\,\mu(\det p_1)|\det p_1|^{-d/2}\,\mu(\det p_2)\,\phi(p_1^{-1}yg).
  \]
\item Let $\operatorname{Ind}(\tau)$ be the space of functions $f$ on $GL_d(k)$
  with values in $\mathcal S(GL_n)$ that are invariant under right translation
  by some open subgroup and satisfy
  $f(bh)=|\det b_A|^{t/2}|\det b_D|^{-n/2}\,\tau(1,b_A,b_D)f(h)$ for every
  $h\in GL_d(k)$ and every block upper triangular $b$, with blocks $b_A$ and
  $b_D$ of sizes $n$ and $t$
  (Definition~\ref{asm-theta-agree-finite-split-models}(2) with $(a,a')=(n,t)$).
  Write $R(h_1)f(h)=f(hh_1)$.
\end{enumerate}
For $y\in GL_n(k)$ let $(y;0)\in R$ have upper $n\times n$ block $y$ and lower
block $0$. Then $E(\Phi)(h)(y)=\mu(\det h)|\det h|^{-n/2}\,\Phi(h^{-1}(y;0))$
defines a linear bijection $E$ from $\mathcal S(R)$ onto
$\operatorname{Ind}(\tau)$, and
$E(\omega^R(g,h)\Phi)(x)=\tau(g,1,1)\bigl(E(\Phi)(xh)\bigr)$ for all $g$, $h$,
$x$ and $\Phi$.
\end{lemma}
\leanref{Development/AutomorphicSpectrum/Proofs/BlockDecomposition/OmegaEquivInducedFibreStatement.lean}{GlobalTheory.T.omega_equiv_induced_fibre}
\leanref{Development/AutomorphicSpectrum/Proofs/BlockDecomposition/OmegaEquivInducedFibre.lean}{GlobalTheory.CompanionTheory.ThetaGLLen.fibStr}
\leanref{Development/AbelianVarietyTheory/Proofs/TwistedOperators/SchwartzBruhatSweepSmoothlyInduced.lean}{GlobalTheory.SBInduced.sb_induced}
\leanref{Development/AbelianVarietyTheory/BlockTheory/TwistedDiscreteSeries/OpenClosedSequences.lean}{GlobalTheory.IndefiniteRowSideFullHecke.OpenClosedSequences.omegaOn}
\leanref{Development/AutomorphicSpectrum/Proofs/BlockDecomposition/FibRepApply.lean}{GlobalTheory.CompanionTheory.ThetaGLLen.fibRep}
\leanref{Development/AbelianVarietyTheory/BlockTheory/TwistedDiscreteSeries/OpenClosedSequences.lean}{GlobalTheory.IndefiniteRowSideFullHecke.OpenClosedSequences.indSnd}
\begin{proof}
Let $K=GL_d(\mathcal O)$ and $j(y)=(y;0)$. Three facts about $j$ come first.
(a) Every $x\in R$ is $w^{-1}j(y)$ with $w\in K$ and $y\in GL_n(k)$: as $x$ has
rank $n$, $x=u^{-1}j(y')$ for some $u\in GL_d(k)$ and $y'\in GL_n(k)$; by the
Iwasawa decomposition $u=pw$ with $p$ block upper triangular and $w\in K$; and
$x=w^{-1}p^{-1}j(y')=w^{-1}j(p_A^{-1}y')$. (b) For block upper triangular $p$,
$p^{-1}j(y)=j(p_A^{-1}y)$. (c) If $h^{-1}j(y)=j(y')$, then the lower left block of
$h^{-1}$ times $y$ is $0$, so, as $y$ is invertible, $h^{-1}$ and $h$ are block
upper triangular. Also $j$ is a closed embedding, with image the matrices in $R$
whose lower block is $0$.

$E(\Phi)(h)$ is Schwartz-Bruhat, being $\Phi$ composed with the closed embedding
$y\mapsto h^{-1}j(y)$. As $\Phi$ is locally constant with compact support and
$GL_d(k)$ acts continuously, $\Phi(h^{-1}x)=\Phi(x)$ for all $x$ when $h$ is near
$1$; and $\mu(\det h)=|\det h|=1$ for $h$ near $1$; so $E(\Phi)$ is invariant under
right translation by an open subgroup. For block upper triangular $b$, by (b),
\[
E(\Phi)(bh)(y)=\mu(\det b)|\det b|^{-n/2}E(\Phi)(h)(b_A^{-1}y)
=|\det b_A|^{t/2}|\det b_D|^{-n/2}\,\bigl(\tau(1,b_A,b_D)E(\Phi)(h)\bigr)(y),
\]
as $\det b=\det b_A\det b_D$. So $E(\Phi)\in\operatorname{Ind}(\tau)$. As
$(xh)^{-1}j(y)g=h^{-1}x^{-1}j(yg)$,
\[
E(\omega^R(g,h)\Phi)(x)(y)=\mu(\det g)^{-1}|\det g|^{d/2}E(\Phi)(xh)(yg)=\bigl(\tau(g,1,1)E(\Phi)(xh)\bigr)(y).
\]

If $E(\Phi)=0$, then $\Phi$ vanishes on every $w^{-1}j(y)$, so $\Phi=0$ by (a).
Let $f\in\operatorname{Ind}(\tau)$, invariant under right translation by an
open subgroup $U$. Put
$\Phi(h^{-1}j(y))=\mu(\det h)^{-1}|\det h|^{n/2}f(h)(y)$. This is well defined:
if $h^{-1}j(y)=h'^{-1}j(y')$, then $b=hh'^{-1}$ is block upper triangular by
(c), $y'=b_A^{-1}y$ by (b), and the law of $f$ at $h=bh'$ gives the same value.
By (a) it defines $\Phi$ on $R$, and $E(\Phi)=f$. The map
$(y,w)\mapsto w^{-1}j(y)$ from $GL_n(k)\times K$ onto $R$ is continuous and
closed, as $K$ is compact and $j$ a closed embedding, so it is a quotient map;
and $(y,w)\mapsto\mu(\det w)^{-1}f(w)(y)$ is locally constant, as $f$ is
invariant under $U$ and each $f(w)$ is locally constant. So $\Phi$ is locally
constant. Finitely many classes $w_iU$ cover $K$, and $\Phi$ vanishes outside the
image of the compact set $\bigcup_i\operatorname{supp}f(w_i)\times K$. So
$\Phi\in\mathcal S(R)$.
\end{proof}
% Lean: fibStr applies sb_induced with r = lAct, j = topEmbR, P = parab (two e), K = KO,
% exists_topEmbR (a), lAct_topEmbR (b) and eq_zero_of_vecMul (c); sb_induced uses
% isQuotientMap_sweep, descend and schwartzBruhat_eventually_invariant_of_continuous_act.
% The set R is RankIn k (Fin n) (Fin d) {n}, matrices of size d x n.

\begin{lemma}\label{asm-theta-agree-coinvariants-haar}
Let $k$, $n\leq d$, $\mu$, $\mathcal S(GL_n)$ and $\tau$ be as in
Lemma~\ref{asm-theta-agree-coinvariants-fibre}, let $dg$ be a Haar measure on
$GL_n(k)$, and put $\xi(g)=\mu(\det g)^{-1}|\det g|^{d/2}$ and
$l(\phi)=\int\phi(g)\xi(g)\,dg$ for $\phi\in\mathcal S(GL_n)$. Then $l\neq0$;
$l(\tau(g,p_1,p_2)\phi)=\mu(\det p_2)\,l(\phi)$ for all $g$, $p_1$, $p_2$ and
$\phi$; and every $\phi$ with $l(\phi)=0$ lies in the span of the
$\tau(g,1,1)\phi'-\phi'$.
\end{lemma}
\leanref{Development/WeilBound/CoreStatements/LocallyConstantCoinvariants/LocConstChar.lean}{PositivityTheory.KCo.exists_fun}
\leanref{Development/WeilBound/CoreStatements/LocallyConstantCoinvariants/LocConstChar.lean}{PositivityTheory.KCo.common_gen}
\begin{proof}
As $GL_n(k)$ is unimodular, the substitution $y\mapsto p_1yg^{-1}$ gives
$\int\phi(p_1^{-1}yg)\xi(y)\,dy=\xi(p_1)\xi(g)^{-1}l(\phi)$. The factor of
$\tau(g,p_1,p_2)$ is $\xi(g)\xi(p_1)^{-1}\mu(\det p_2)$, which gives the second
claim. For $\phi_0=\mathbf 1_{GL_n(\mathcal O)}\,\xi^{-1}$, $l(\phi_0)$ is the
volume of $GL_n(\mathcal O)$, which is positive.

Let $\phi$ with $l(\phi)=0$. As $\phi$ and $\phi_0$ are locally constant with
compact support, there is an open compact subgroup $K'$ with
$\phi(k'y)=\phi(y)$ and $\phi_0(k'y)=\phi_0(y)$ for $k'\in K'$; so both are finite
linear combinations of indicator functions $\mathbf 1_{K'y_j}$, which are the right
translates $y\mapsto\theta(yy_j^{-1})$ of $\theta=\mathbf 1_{K'}$. Let $\mathrm{Rel}$
be the span of the $\tau(g,1,1)\phi'-\phi'$. As
$\tau(g,1,1)\phi'(y)=\xi(g)\phi'(yg)$, the right translate
$\phi'(\cdot\,g)$ is $\xi(g)^{-1}\phi'$ modulo $\mathrm{Rel}$, and
$l(\phi'(\cdot\,g))=\xi(g)^{-1}l(\phi')$. So every $w$ in the span of the right
translates of $\theta$ is $s\theta$ modulo $\mathrm{Rel}$ for some $s\in\mathbf C$,
with $l(w)=s\,l(\theta)$. For $w=\phi_0$ this gives $l(\theta)\neq0$; for
$w=\phi$ it gives $s=0$, so $\phi\in\mathrm{Rel}$.
\end{proof}
% Lean: exists_fun with xi = detChar (mu^{-1} |.|^{(n+t)/2}) and rankHomeo identifying GL_n
% with RankIn k (Fin n) (Fin n) {n}; the second claim and the kernel in terms of tau(g,1,1)
% are hτl and the last step of irr_ind_narrow.

\begin{lemma}\label{asm-theta-agree-coinvariants-exact}
Let $k$ be a nonarchimedean local field with ring of integers $\mathcal O$, let
$G=GL_d(k)$, $K=GL_d(\mathcal O)$, and let $P$ be a closed subgroup of $G$ with
$G=PK$. For a representation $\sigma$ of $P$ on a complex vector space $V$, let
$\operatorname{Ind}\sigma$ be the space of functions $f\colon G\to V$ with
$f(pg)=\sigma(p)f(g)$ for $p\in P$ and $g\in G$ that are invariant under right
translation by some open subgroup of $G$. Let $\sigma$ and $\rho$ be
representations of $P$ on $V$ and $U$, with every vector of $V$ fixed by an open
subgroup of $P$; let $\alpha\colon V\to U$ be a surjective linear map with
$\alpha\sigma(p)=\rho(p)\alpha$; and let $(T_\gamma)_{\gamma\in\Gamma}$ be linear
maps $V\to V$ commuting with every $\sigma(p)$, such that the kernel of $\alpha$
is spanned by the $T_\gamma z-z$. Then $f\mapsto\alpha\circ f$ maps
$\operatorname{Ind}\sigma$ onto $\operatorname{Ind}\rho$, and its kernel is
spanned by the $T_\gamma\circ f-f$, $f\in\operatorname{Ind}\sigma$.
\end{lemma}
\leanref{Development/AbelianVarietyTheory/Proofs/TwistedOperators/SmIndCoinvariants.lean}{GlobalTheory.SmGL.IndCoinv.smInd_coinvariants}
\leanref{Development/AbelianVarietyTheory/Proofs/TwistedOperators/SmIndCoinvariants.lean}{GlobalTheory.SmGL.IndCoinv.mem_of_fixed}
\begin{proof}
Let $J\subset K$ be an open subgroup normalized by $K$, such as a principal
congruence subgroup, and $H=J\cap P$, a compact group. For $w\in V$ fixed by
$H$, the function $b_J(w)$ supported on $PJ$ with $b_J(w)(pj)=\sigma(p)w$ is well
defined and lies in $\operatorname{Ind}\sigma$; and $\alpha\circ b_J(w)=b_J(\alpha w)$.

(i) $\operatorname{Ind}\sigma$ is spanned by the right translates of the
$b_J(w)$. Let $f$ be invariant under right translation by such a $J$. As $K$ is
compact and $G=PK$, $G$ is the disjoint union of finitely many $Px_iJ$ with
$x_i\in K$. For $j\in H$, $\sigma(j)f(x_i)=f(x_i\cdot x_i^{-1}jx_i)=f(x_i)$, as
$K$ normalizes $J$; and $f$ is the sum over $i$ of the translates
$g\mapsto b_J(f(x_i))(gx_i^{-1})$, which is supported on $PJx_i=Px_iJ$.

(ii) Every $w'\in U$ fixed by $H$ is $\alpha(w)$ for some $w\in V$ fixed by $H$:
take a preimage $w_1$, fixed by an open subgroup $H_1$ of $H$, and let $w$ be the
average of the $\sigma(h)w_1$ over $H/H_1$.

By (i) for $\rho$, (ii), and as the image is stable under right translation, the
map is onto. Let $f$ with $\alpha\circ f=0$. Then $f$ takes values in
$\ker\alpha$, on which $\sigma$ restricts, so by (i) $f$ is a linear combination
of right translates of $b_J(w)$ with $w\in\ker\alpha$ fixed by $H$. The span $Z$
of the $T_\gamma\circ f'-f'$ is stable under right translation, so it suffices
that $b_J(w)\in Z$. We show that $w$ lies in the span $M$ of the $T_\gamma z-z$
with $z$ fixed by $H$; then $b_J(w)$ is a combination of the
$b_J(T_\gamma z)-b_J(z)=T_\gamma\circ b_J(z)-b_J(z)\in Z$. If $w\notin M$, take
a linear form on the $H$-fixed vectors vanishing on $M$ and not at $w$, extend
it to $V$, and average it over $H$, as in (ii), to an $H$-invariant form
$\lambda$ with the same values on the $H$-fixed vectors. An $H$-invariant form
is determined by these values, as $v$ and its average over $H$ have the same
value. For each $\gamma$, $\lambda\circ T_\gamma$ is $H$-invariant and agrees
with $\lambda$ on the $H$-fixed vectors, as $T_\gamma z-z\in M$ for those $z$;
so $\lambda\circ T_\gamma=\lambda$. Then $\lambda$ vanishes on every
$T_\gamma z-z$, hence on $w$, which is absurd.
\end{proof}
% Lean: smInd_coinvariants with bump, smIndMap_bump, mem_span_bump and mem_of_bump for (i),
% IwahoriOnto.exists_fixed_map_eq for (ii), and mem_of_fixed (DualAverage.exists_invariant_form,
% dual_ext) for the kernel; the subgroups J are exists_open_normalised_subgroup_subset.

\begin{lemma}\label{asm-theta-agree-coinvariants-irreducible}
Let $k$, $n\leq d$, $t=d-n$, $\mu$, $\tau$, $\operatorname{Ind}(\tau)$ and
$R(h_1)$ be as in Lemma~\ref{asm-theta-agree-coinvariants-fibre}, with $\mu$
unitary, smooth and $\mu\neq1$, and let $\mathrm{Rel}$ be the span of the
functions $h\mapsto\tau(g,1,1)f(h)-f(h)$ with $f\in\operatorname{Ind}(\tau)$ and
$g\in GL_n(k)$. Let $N\subset\operatorname{Ind}(\tau)$ be a subspace that
contains $\mathrm{Rel}$ and is stable under every $R(h_1)$,
$h_1\in GL_d(k)$. Then $N=\operatorname{Ind}(\tau)$ or $N\subset\mathrm{Rel}$.
\end{lemma}
\leanref{Development/AutomorphicSpectrum/Proofs/BlockDecomposition/KCoinvOfNarrow.lean}{GlobalTheory.CompanionTheory.KCoNarrow.irr_gen_narrow}
\leanref{Development/AutomorphicSpectrum/Proofs/BlockDecomposition/KCoinvOfNarrow.lean}{GlobalTheory.CompanionTheory.KCoNarrow.irr_ind_narrow}
\leanref{Development/WeilBound/CoreStatements/LocallyConstantCoinvariants/TopOfMap.lean}{PositivityTheory.KCo.top_of_map}
\begin{proof}
Let $l$ be as in Lemma~\ref{asm-theta-agree-coinvariants-haar}, and
$\Lambda(f)=l\circ f$. By that lemma,
$\Lambda(f)(bh)=\mu(\det b_D)|\det b_A|^{t/2}|\det b_D|^{-n/2}\Lambda(f)(h)$ for
block upper triangular $b$, so $\Lambda$ maps $\operatorname{Ind}(\tau)$ to the
space $I(1,\mu)$ of Definition~\ref{asm-theta-agree-finite-split-models}(2), with
$(a,a')=(n,t)$, and commutes with right translation. Apply
Lemma~\ref{asm-theta-agree-coinvariants-exact} with $P$ the block upper
triangular matrices, $\sigma(b)=|\det b_A|^{t/2}|\det b_D|^{-n/2}\tau(1,b_A,b_D)$ on
$\mathcal S(GL_n)$, $\rho(b)=\mu(\det b_D)|\det b_A|^{t/2}|\det b_D|^{-n/2}$ on
$\mathbf C$, $\alpha=l$ and $T_g=\tau(g,1,1)$ for $g\in GL_n(k)$. Here
$G=PK$ by the Iwasawa decomposition; every $\phi\in\mathcal S(GL_n)$ is fixed by
$\sigma(b)$ for $b$ near $1$, as $\phi(p_1^{-1}y)=\phi(y)$ for $p_1$ near $1$ and the
characters are trivial near $1$; the $T_g$ commute with $\sigma$; and $l$ is onto,
with kernel spanned by the $T_g\phi-\phi$, by
Lemma~\ref{asm-theta-agree-coinvariants-haar}. So $\Lambda$ maps
$\operatorname{Ind}(\tau)$ onto $I(1,\mu)$, with kernel $\mathrm{Rel}$.

$\Lambda(N)$ is stable under right translation, and $1$ and $\mu$ are unitary
smooth characters with $\mu\neq1$, so by
Theorem~\ref{asm-theta-agree-finite-bernstein}, $\Lambda(N)=0$ or
$\Lambda(N)=I(1,\mu)$. In the first case $N\subset\ker\Lambda=\mathrm{Rel}$. In the
second, for $f\in\operatorname{Ind}(\tau)$ choose $f'\in N$ with
$\Lambda(f')=\Lambda(f)$; then $f-f'\in\mathrm{Rel}\subset N$, so $f\in N$.
\end{proof}
% Lean: irr_ind_narrow (tau = fibRep, l from exists_fun with xi = detChar (mu^{-1} |.|^{(n+t)/2}))
% applies irr_gen_narrow, which uses Narrow for detRep 1 and detRep mu (Theorem
% asm-theta-agree-finite-bernstein; Narrow is BernBridge.Split.NarrowBernstein),
% LeviInt.levi_integral_facts for G = PK, smInd_coinvariants and top_of_map.

\begin{lemma}\label{asm-theta-agree-dictionary-split}
Let $E$ be a CM field with maximal totally real subfield $F$ and complex
conjugation $c$, let $v$ be a finite place of $F$, let $t\colon E\to F_v$ be an
$F$-embedding, and let $\pi_t$ and $\pi_{t\circ c}$ be the maps
$\mathbf A_E\to F_v$ of Definition~\ref{asm-theta-agree-models}(1). For
$s\in F_v$ let $\bar s\in\mathbf A_F$ have component $s$ at $v$ and $0$
elsewhere.
\begin{enumerate}
\item The maps $\pi_t$ and $\pi_{t\circ c}$ are ring homomorphisms with
  $\pi_t(az)=a_v\pi_t(z)$ and $\pi_{t\circ c}(az)=a_v\pi_{t\circ c}(z)$ for
  $a\in\mathbf A_F$, with $\pi_t(x)=t(x)$ and $\pi_{t\circ c}(x)=t(c(x))$ for
  $x\in E$, and with $\pi_t\circ c=\pi_{t\circ c}$ and
  $\pi_{t\circ c}\circ c=\pi_t$. The map
  $z\mapsto(\pi_t(z),\pi_{t\circ c}(z))$ is a bijection from the elements of
  $\mathbf A_E$ at $v$ (Definition~\ref{asm-theta-agree-models}(2)) onto
  $F_v\times F_v$. If $z$ is at $v$, then so is $c(z)$, and
  $\operatorname{Tr}_{\mathbf A_E/\mathbf A_F}(z)=\overline{\pi_t(z)+\pi_{t\circ c}(z)}$.
\item For $g_1,g_2\in GL_m(F_v)$ there is $g\in GL_m(\mathbf A_E)$ with $g-1$ at
  $v$, $\pi_t(g)=g_1$ and $\pi_{t\circ c}(g)=g_2$.
\item Let $z\in\mathbf A_E^\times$ and $x\in F_v^\times$ with $z-1$ at $v$,
  $\pi_t(z)=x$ and $\pi_{t\circ c}(z)=1$. Then $\chi(z)=\chi_t(x)$ for every
  Hecke character $\chi$ of $E$
  (Definition~\ref{asm-theta-agree-coinvariants-local}(4)), and
  $|N_{E/F}(z)|_{\mathbf A_F}=|x|$, the normalized absolute value of $x$ in
  $F_v$. Every such $\chi_t$ is unitary and smooth.
\item Let $J\in M_m(\mathbf A_E)$ and $\varepsilon\in F_v$ with $\pi_t(J)$
  invertible and $\pi_{t\circ c}(J)=\varepsilon\,\pi_t(J)^{\mathsf T}$, and let
  $U_J$ be the group of the $g\in GL_m(\mathbf A_E)$ with
  $g^{\mathsf T}J\,c(g)=J$. Every $g\in U_J$ has
  $\pi_t(g)^{\mathsf T}\pi_t(J)\,\pi_{t\circ c}(g)=\pi_t(J)$, and for every
  $g_0\in GL_m(F_v)$ there is exactly one $g\in U_J$ with $g-1$ at $v$ and
  $\pi_t(g)=g_0$.
\end{enumerate}
\end{lemma}
\leanref{Development/WeilBound/CoreStatements/LocallyConstantCoinvariants/PfDictDef.lean}{PositivityTheory.KCo.IsSec}
\leanref{Development/WeilBound/CoreStatements/LocallyConstantCoinvariants/PfDictDef.lean}{PositivityTheory.KCo.SecSt}
\leanref{Development/WeilBound/CoreStatements/LocallyConstantCoinvariants/PfDictSec.lean}{PositivityTheory.KCo.section_statement}
\leanref{Development/WeilBound/CoreStatements/LocallyConstantCoinvariants/KdMulRight.lean}{PositivityTheory.KCo.IsSec.exists_gl}
\leanref{Development/WeilBound/CoreStatements/LocallyConstantCoinvariants/KdMulRight.lean}{PositivityTheory.KCo.IsSec.unit_rel}
\leanref{Development/WeilBound/CoreStatements/LocallyConstantCoinvariants/KdMulRight.lean}{PositivityTheory.KCo.IsSec.exists_ug}
\leanref{Development/WeilBound/CoreStatements/LocallyConstantCoinvariants/KdMulRight.lean}{PositivityTheory.KCo.IsSec.ug_ext}
\begin{proof}
(1) As $c\neq1$ and $t$ is injective, $t\neq t\circ c$, so the map
$E\otimes_FF_v\to F_v\times F_v$, $x\otimes s\mapsto(t(x)s,t(c(x))s)$, is an
isomorphism of $F_v$-algebras. Under it $c$ exchanges the two factors, and the
trace to $F_v$ is the sum of the two components. The elements of $\mathbf A_E$
at $v$ are the elements of $E\otimes_FF_v$ placed at $v$, and $\pi_t$ and
$\pi_{t\circ c}$ read their two components. This gives (1).

(2) By (1) there are matrices $g'$ and $g''$ at $v$ with $\pi_t(g')=g_1-1$,
$\pi_{t\circ c}(g')=g_2-1$, $\pi_t(g'')=g_1^{-1}-1$ and
$\pi_{t\circ c}(g'')=g_2^{-1}-1$. The products $(1+g')(1+g'')$ and
$(1+g'')(1+g')$ differ from $1$ by matrices at $v$ whose images under $\pi_t$
and $\pi_{t\circ c}$ are $0$, so they are $1$ by (1), and $g=1+g'$ is as
required.

(3) The idele $z$ has component $x$ at the place $w$ of $E$ defined by $t$,
component $\pi_{t\circ c}(z)=1$ at the conjugate place, and component $1$ at
every place not above $v$. So $\chi(z)$ is the value at $x$ of the component of
$\chi$ at $w$, which is $\chi_t(x)$, and $N_{E/F}(z)$ has component $x$ at $v$
and $1$ elsewhere. A Hecke character is continuous with values of absolute
value $1$, so $\chi_t$ is unitary and its kernel is open.

(4) The first claim is the image of $g^{\mathsf T}Jc(g)=J$ under $\pi_t$, by
(1). Let $A=\pi_t(J)$ and $g_0\in GL_m(F_v)$, and take $g$ as in (2) with
$\pi_t(g)=g_0$ and $\pi_{t\circ c}(g)=A^{-1}(g_0^{\mathsf T})^{-1}A$. The matrix
$g^{\mathsf T}Jc(g)-J=(g-1)^{\mathsf T}Jc(g)+J\,c(g-1)$ is at $v$. Its image
under $\pi_t$ is $g_0^{\mathsf T}A\,A^{-1}(g_0^{\mathsf T})^{-1}A-A=0$, and its
image under $\pi_{t\circ c}$ is
$A^{\mathsf T}g_0^{-1}(A^{\mathsf T})^{-1}\cdot\varepsilon A^{\mathsf T}g_0-\varepsilon A^{\mathsf T}=0$.
So $g\in U_J$ by (1). If $g,g'\in U_J$ have $g-1$ and $g'-1$ at $v$ and
$\pi_t(g)=\pi_t(g')$, then the first claim gives
$\pi_{t\circ c}(g)=A^{-1}(\pi_t(g)^{\mathsf T})^{-1}A=\pi_{t\circ c}(g')$, and
$g=g'$ by (1).
\end{proof}
% Lean: (1) and (3) are the fields at', fst, snd, ext, conj, atConj, smul, inc, trace and idele of
%   IsSec, which section_statement proves for the section tensAt E v of (splitAlgEquiv h t).symm,
%   with modu_single and eq_locIdele for idele; the last claim of (3) is usm_chrW. (2) is
%   IsSec.exists_gl; (4) is IsSec.unit_rel, IsSec.exists_ug and IsSec.ug_ext.

\begin{theorem}[Stone, von Neumann \cite{MVW87}]\label{asm-theta-agree-dictionary-stone-von-neumann}
Let $k$ be a nonarchimedean local field in which $2\neq0$, and let $\psi$ be a
continuous character of $k$ that is not $1$. For a finite set $\kappa$ and
$w=(a,b)\in k^\kappa\times k^\kappa$ let $\rho_\psi(w)$ be the operator
$(\rho_\psi(w)\phi)(x)=\psi(x\cdot b+\tfrac12a\cdot b)\,\phi(x+a)$ on functions
on $k^\kappa$, and put $[w,w']=a\cdot b'-a'\cdot b$ for $w'=(a',b')$.
\begin{enumerate}
\item Let $\phi_0\neq0$ be a Schwartz-Bruhat function on $k^\kappa$, and let
  $l$ be a linear form on the functions on $k^\kappa$ with
  $l(\rho_\psi(w)\phi_0)=0$ for every $w$. Then $l(\phi)=0$ for every
  Schwartz-Bruhat $\phi$.
\item Let $\mu$ be a second finite set and
  $e\colon k^\mu\times k^\mu\to k^\kappa\times k^\kappa$ a $k$-linear
  isomorphism with $[e(w),e(w')]=[w,w']$ for all $w,w'$. Then there is a linear
  map $\iota$ from the functions on $k^\kappa$ to the functions on $k^\mu$,
  bijective between the Schwartz-Bruhat functions, with
  $\iota(\rho_\psi(e(w))\phi)=\rho_\psi(w)\,\iota\phi$ for every $w$ and every
  Schwartz-Bruhat $\phi$.
\end{enumerate}
\end{theorem}
\leanref{Development/WeilBound/CoreStatements/LocallyConstantCoinvariants/PfDictDef.lean}{PositivityTheory.KCo.Cyc}
\leanref{Development/WeilBound/CoreStatements/LocallyConstantCoinvariants/PfDictDef.lean}{PositivityTheory.KCo.ModelEx}
\leanref{Development/WeilBound/CoreStatements/LocallyConstantCoinvariants/PfDictMod.lean}{PositivityTheory.KCo.cyclicity}
\leanref{Development/WeilBound/CoreStatements/LocallyConstantCoinvariants/PfDictMod.lean}{PositivityTheory.KCo.model_existence}

\begin{lemma}\label{asm-theta-agree-dictionary-evaluation}
Let $k$, $\psi$, $\kappa$ and $\rho_\psi$ be as in
Theorem~\ref{asm-theta-agree-dictionary-stone-von-neumann}, and let $l$ be a
linear form on the functions on $k^\kappa$ with
$l(\rho_\psi(0,y)\phi)=l(\phi)$ for every $y\in k^\kappa$ and every
Schwartz-Bruhat $\phi$. Then there is $a\in\mathbf C$ with $l(\phi)=a\,\phi(0)$
for every Schwartz-Bruhat $\phi$.
\end{lemma}
\leanref{Development/AutomorphicSpectrum/Proofs/MembershipProofs/ExistsEqMulApplyZero.lean}{GlobalTheory.CompanionTheory.WeilCoinv.exists_eq_mul_apply_zero}
\begin{proof}
Note that $(\rho_\psi(0,y)\phi)(x)=\psi(x\cdot y)\phi(x)$, and that
$x\mapsto\psi(x\cdot y)$ is locally constant, as $\psi$ is continuous with values
of absolute value $1$.

First, let $T\subset k^\kappa$ be finite and $\phi$ Schwartz-Bruhat, such that
every $x$ with $\phi(x)\neq0$ has $\psi(x\cdot y)\neq1$ for some $y\in T$. Then
$l(\phi)=0$, by induction on $T$. If $T$ is empty, then $\phi=0$. Otherwise let
$y\in T$, let $\phi_1$ be equal to $\phi$ where $\psi(x\cdot y)=1$ and to $0$
elsewhere, and let $g$ be equal to $\phi(x)/(\psi(x\cdot y)-1)$ where
$\psi(x\cdot y)\neq1$ and to $0$ elsewhere. Both are Schwartz-Bruhat, and
$\phi=\phi_1+\rho_\psi(0,y)g-g$, so $l(\phi)=l(\phi_1)$. The function $\phi_1$
satisfies the hypothesis with $T\setminus\{y\}$.

Next, let $\phi$ be Schwartz-Bruhat with $\phi(0)=0$. Its support is closed, as
$\phi$ is locally constant, and lies in a compact set, so it is compact, and it
does not contain $0$. Every $x\neq0$ has $\psi(x\cdot y)\neq1$ for some $y$:
take $s$ with $\psi(s)\neq1$ and a nonzero coordinate $x_i$, and let $y$ have
coordinate $s/x_i$ at $i$ and $0$ elsewhere. The sets
$\{x:\psi(x\cdot y)\neq1\}$ are open, so finitely many of them cover the
support, and $l(\phi)=0$ by the first step.

Finally, let $O\subset k^\kappa$ be a compact open subgroup, $\phi_0$ its
indicator function and $a=l(\phi_0)$. For Schwartz-Bruhat $\phi$, the function
$\phi-\phi(0)\phi_0$ vanishes at $0$, so $l(\phi)=a\,\phi(0)$.
\end{proof}

\begin{situation}\label{asm-theta-agree-dictionary-situation}
Let $B$ be a theta base for a CM field $E$ with character $\psi$
(Definition~\ref{asm-theta-agree-models}(1)), maximal totally real subfield $F$
and complex conjugation $c$, and put $\mathbf A=\mathbf A_F$. Let $v$ be a
finite place of $F$ and $t\colon E\to F_v$ an $F$-embedding. Let $n,d\geq1$, let
$J\in M_n(E)$ and $J'\in M_d(E)$ be matrices with $t(J)$ and $t(J')$
invertible, and suppose $t(c(J))=\varepsilon\,t(J)^{\mathsf T}$ and
$t(c(J'))=-\varepsilon\,t(J')^{\mathsf T}$ for some $\varepsilon=\pm1$. For
$s\in F_v$ let $\bar s\in\mathbf A$ have component $s$ at $v$ and $0$
elsewhere, and put $\psi_v(s)=\psi(\bar s)$. For a finite set $\kappa$ and
$w\in F_v^\kappa\times F_v^\kappa$ write $\rho_v(w)$ for the operator
$\rho_{\psi_v}(w)$ of Theorem~\ref{asm-theta-agree-dictionary-stone-von-neumann}
on functions on $F_v^\kappa$, and use the same notation on functions on
$M_{n\times d}(F_v)$, with the product $Y\cdot Z=\operatorname{tr}(Y^{\mathsf T}Z)$.
For $w\in\mathbf A^\kappa\times\mathbf A^\kappa$ let $w_v$ be its component at
$v$. Give $\mathbb W=M_{n\times d}(\mathbf A_E)$ the form
$\langle x,y\rangle=-\operatorname{Tr}_{\mathbf A_E/\mathbf A}h(x,y)$ with
$h(x,y)=\sum_{i,j,k,l}x_{ij}J_{ik}J'_{jl}\,c(y_{kl})$. Let $I$ be a finite set
and $e=(e_1,e_2)\colon\mathbb W\to\mathbf A^I\times\mathbf A^I$ an
$\mathbf A$-linear isomorphism with
$\langle u,u'\rangle=e_1(u)\cdot e_2(u')-e_1(u')\cdot e_2(u)$, and let $\rho(u)$,
$u\in\mathbb W$, be its Schrödinger operators, given by the formula of
Definition~\ref{asm-theta-agree-models}(3). For $u\in\mathbb W$ put
\[
\operatorname{co}(u)=-t(J)\,\pi_{t\circ c}(u)\,t(J')^{\mathsf T},\qquad
\operatorname{co}'(u)=-\pi_t(u),
\]
and let $U_{J'}$ be the group of the $h\in GL_d(\mathbf A_E)$ with
$h^{\mathsf T}J'c(h)=J'$.
\end{situation}

\begin{lemma}\label{asm-theta-agree-dictionary-coordinates}
In Situation~\ref{asm-theta-agree-dictionary-situation}:
\begin{enumerate}
\item The maps $\operatorname{co}$ and $\operatorname{co}'$ are additive, with
  $\operatorname{co}(au)=a_v\operatorname{co}(u)$ and
  $\operatorname{co}'(au)=a_v\operatorname{co}'(u)$ for $a\in\mathbf A$, and
  $u\mapsto(\operatorname{co}(u),\operatorname{co}'(u))$ is a bijection from the
  elements of $\mathbb W$ at $v$ onto $M_{n\times d}(F_v)\times M_{n\times d}(F_v)$.
  Moreover $\operatorname{co}(u)=0$ if and only if $\pi_{t\circ c}(u)=0$, and
  $\operatorname{co}'(u)=0$ if and only if $\pi_t(u)=0$.
\item Let $u\in\mathbb W$ be at $v$ and $u'\in\mathbb W$. Then $\langle u,u'\rangle$
  is at $v$ with component
  $[(\operatorname{co}(u),\operatorname{co}'(u)),(\operatorname{co}(u'),\operatorname{co}'(u'))]$,
  the form of Theorem~\ref{asm-theta-agree-dictionary-stone-von-neumann} for the
  product $Y\cdot Z=\operatorname{tr}(Y^{\mathsf T}Z)$.
\item Let $h\in U_{J'}$ and $h_0=\pi_t(h)$. Then
  $\operatorname{co}(uh^{\mathsf T})=\operatorname{co}(u)h_0^{-1}$ and
  $\operatorname{co}'(uh^{\mathsf T})=\operatorname{co}'(u)h_0^{\mathsf T}$ for
  every $u\in\mathbb W$.
\item There is an $F_v$-linear isomorphism $e_v$ from
  $M_{n\times d}(F_v)\times M_{n\times d}(F_v)$ onto $F_v^I\times F_v^I$ with
  $[e_v(w),e_v(w')]=[w,w']$ for all $w,w'$ and
  $e_v(\operatorname{co}(u),\operatorname{co}'(u))=e(u)_v$ for every
  $u\in\mathbb W$.
\end{enumerate}
\end{lemma}
\leanref{Development/WeilBound/CoreStatements/LocallyConstantCoinvariants/KdExistsInv.lean}{PositivityTheory.KCo.Grams}
\leanref{Development/WeilBound/CoreStatements/LocallyConstantCoinvariants/KdExistsInv.lean}{PositivityTheory.KCo.coord}
\leanref{Development/WeilBound/CoreStatements/LocallyConstantCoinvariants/KdExistsInv.lean}{PositivityTheory.KCo.symp_eq}
\leanref{Development/WeilBound/CoreStatements/LocallyConstantCoinvariants/EqZeroOfLocW.lean}{PositivityTheory.KCo.exists_coord}
\leanref{Development/WeilBound/CoreStatements/LocallyConstantCoinvariants/EqZeroOfLocW.lean}{PositivityTheory.KCo.coord_inj}
\leanref{Development/WeilBound/CoreStatements/LocallyConstantCoinvariants/EqZeroOfLocW.lean}{PositivityTheory.KCo.coord_actB}
\leanref{Development/WeilBound/CoreStatements/LocallyConstantCoinvariants/EqZeroOfLocW.lean}{PositivityTheory.KCo.exists_e}
\begin{proof}
Write $A=t(J)$ and $A'=t(J')$; by Lemma~\ref{asm-theta-agree-dictionary-split}(1),
$\pi_t(J)=A$, $\pi_{t\circ c}(J)=t(c(J))=\varepsilon A^{\mathsf T}$, and likewise
for $J'$ with $-\varepsilon$.

(1) The first claims follow from
Lemma~\ref{asm-theta-agree-dictionary-split}(1), and the last two from the
invertibility of $A$ and $A'$. Given $(Y,Z)$, the element $u$ at $v$ with
$\pi_t(u)=-Z$ and $\pi_{t\circ c}(u)=-A^{-1}Y(A'^{\mathsf T})^{-1}$, which
exists by Lemma~\ref{asm-theta-agree-dictionary-split}(1), has coordinates
$(Y,Z)$, and it is the only one at $v$ by the same lemma.

(2) As $u$ is at $v$, so is $h(u,u')$. Write $P$ and $Q$ for $\pi_t(u)$ and
$\pi_{t\circ c}(u)$, and $P'$ and $Q'$ for those of $u'$. By
Lemma~\ref{asm-theta-agree-dictionary-split}(1),
$\pi_t(h(u,u'))=\operatorname{tr}(A^{\mathsf T}PA'Q'^{\mathsf T})$ and
$\pi_{t\circ c}(h(u,u'))=\operatorname{tr}(\varepsilon A\,Q\,(-\varepsilon)A'^{\mathsf T}P'^{\mathsf T})
=-\operatorname{tr}(AQA'^{\mathsf T}P'^{\mathsf T})$. So $\langle u,u'\rangle$ is
at $v$ with component
$\operatorname{tr}(AQA'^{\mathsf T}P'^{\mathsf T})-\operatorname{tr}(A^{\mathsf T}PA'Q'^{\mathsf T})$,
which is
$\operatorname{tr}(\operatorname{co}(u)^{\mathsf T}\operatorname{co}'(u'))-\operatorname{tr}(\operatorname{co}(u')^{\mathsf T}\operatorname{co}'(u))$.

(3) By Lemma~\ref{asm-theta-agree-dictionary-split}(4) for $J'$,
$h_0^{\mathsf T}A'\pi_{t\circ c}(h)=A'$, so
$\pi_{t\circ c}(h)^{\mathsf T}A'^{\mathsf T}=A'^{\mathsf T}h_0^{-1}$. Hence
$\operatorname{co}(uh^{\mathsf T})=-A\pi_{t\circ c}(u)\pi_{t\circ c}(h)^{\mathsf T}A'^{\mathsf T}=\operatorname{co}(u)h_0^{-1}$,
and $\operatorname{co}'(uh^{\mathsf T})=-\pi_t(u)h_0^{\mathsf T}$.

(4) For $w$ in $M_{n\times d}(F_v)^2$ let $u_w$ be the element at $v$ with
coordinates $w$, given by (1), and put $e_v(w)=e(u_w)_v$. For every $u$, the
element $u_{(\operatorname{co}(u),\operatorname{co}'(u))}$ is $\bar1u$, by (1), and
$e(\bar1u)_v=e(u)_v$ as $e$ is $\mathbf A$-linear; so
$e_v(\operatorname{co}(u),\operatorname{co}'(u))=e(u)_v$. By (1),
$u_{w+w'}=u_w+u_{w'}$ and $u_{sw}=\bar s\,u_w$ for $s\in F_v$, so $e_v$ is
$F_v$-linear. If $e_v(w)=0$, then $e(u_w)$, which is at $v$, is $0$, so $u_w=0$
and $w=0$. For $(a,b)\in F_v^I\times F_v^I$ the element
$u=e^{-1}(\bar a,\bar b)$ is at $v$, and
$e_v(\operatorname{co}(u),\operatorname{co}'(u))=(a,b)$. Finally, by (2) and the
property of $e$, the component at $v$ of
$\langle u_w,u_{w'}\rangle=e_1(u_w)\cdot e_2(u_{w'})-e_1(u_{w'})\cdot e_2(u_w)$ is
both $[w,w']$ and $[e_v(w),e_v(w')]$.
\end{proof}
% Lean: the hypothesis on J and J' is Grams (with the Grams of the Cartan involutions in
%   dictX_of); co and co' are the two halves of coord. (1) is coord_add, coord_smul, exists_coord,
%   coord_inj, coord_snd and coord_fst; (2) is symp_eq; (3) is coord_actB; (4) is exists_e.

\begin{definition}\label{asm-theta-agree-dictionary-chart}
In Situation~\ref{asm-theta-agree-dictionary-situation}, a \emph{Heisenberg
chart} for $e$ at $v$ is a linear map $\iota$ from the functions on $F_v^I$ to
the functions on $M_{n\times d}(F_v)$, bijective between the Schwartz-Bruhat
functions, such that
\[
\iota\bigl(\rho_v(e(u)_v)\phi\bigr)=\rho_v\bigl(\operatorname{co}(u),\operatorname{co}'(u)\bigr)\,\iota\phi
\]
for every $u\in\mathbb W$ and every Schwartz-Bruhat function $\phi$ on $F_v^I$.
\end{definition}
\leanref{Development/WeilBound/CoreStatements/LocallyConstantCoinvariants/EqZeroOfLocW.lean}{PositivityTheory.KCo.LM}

\begin{lemma}\label{asm-theta-agree-dictionary-charts}
In Situation~\ref{asm-theta-agree-dictionary-situation}:
\begin{enumerate}
\item Let $\kappa$ be a finite set and $w\in\mathbf A^\kappa\times\mathbf A^\kappa$
  at $v$, and let $\rho(w)$ be the operator of
  Theorem~\ref{asm-theta-agree-evaluation-heisenberg} for $\psi$. Then
  $\rho_v(w_v)$ maps Schwartz-Bruhat functions to Schwartz-Bruhat functions,
  and $\rho(w)(\phi\otimes r)=(\rho_v(w_v)\phi)\otimes r$ for every function
  $\phi$ on $F_v^\kappa$ and every function $r$ on $\mathbf A^\kappa$ that does
  not depend on the $v$-component. In particular $\rho(w)$ acts at $v$ by
  $\phi\mapsto\rho_v(w_v)\phi$ (Definition~\ref{asm-theta-agree-models}(2)).
\item There is a Heisenberg chart for $e$ at $v$.
\item A linear map from the functions on $F_v^I$ to the functions on
  $M_{n\times d}(F_v)$ that is bijective between the Schwartz-Bruhat functions
  is a Heisenberg chart for $e$ at $v$ if and only if it is a linear model for
  $(\pi_{t\circ c},\pi_t,\operatorname{co})$
  (Definition~\ref{asm-theta-agree-models}(5), with the Schrödinger operators
  of $e$).
\item The matrices $J'$ and $J$, the sign $-\varepsilon$ and the isomorphism
  $e^{\mathsf T}\colon w\mapsto e(w^{\mathsf T})$ of $M_{d\times n}(\mathbf A_E)$
  also satisfy Situation~\ref{asm-theta-agree-dictionary-situation}, and their
  coordinates at $u^{\mathsf T}$ are $\operatorname{co}(u)^{\mathsf T}$ and
  $\operatorname{co}'(u)^{\mathsf T}$. If $\iota$ is a Heisenberg chart for $e$
  at $v$, then $\iota^{\mathsf T}\colon\phi\mapsto(Z\mapsto\iota\phi(Z^{\mathsf T}))$
  is a Heisenberg chart for $e^{\mathsf T}$ at $v$.
\end{enumerate}
\end{lemma}
\leanref{Development/WeilBound/CoreStatements/LocallyConstantCoinvariants/PfDictAd.lean}{PositivityTheory.KCo.ad_heis}
\leanref{Development/WeilBound/CoreStatements/LocallyConstantCoinvariants/EqZeroOfLocW.lean}{PositivityTheory.KCo.exists_lm}
\leanref{Development/WeilBound/CoreStatements/LocallyConstantCoinvariants/LmOfLinBar.lean}{PositivityTheory.KCo.linBar_of_lm}
\leanref{Development/WeilBound/CoreStatements/LocallyConstantCoinvariants/LmOfLinBar.lean}{PositivityTheory.KCo.lm_of_linBar}
\leanref{Development/WeilBound/CoreStatements/LocallyConstantCoinvariants/LmOfLinBar.lean}{PositivityTheory.KCo.Grams.swap}
\leanref{Development/WeilBound/CoreStatements/LocallyConstantCoinvariants/LmOfLinBar.lean}{PositivityTheory.KCo.LM.swap}
\begin{proof}
For $w,w'\in F_v^\kappa\times F_v^\kappa$ a direct computation from the formula
gives $\rho_v(w)\rho_v(w')=\psi_v(\tfrac12[w,w'])\,\rho_v(w+w')$.

(1) The character $\psi_v$ is continuous with values of absolute value $1$, so
it is locally constant, and $\rho_v(w_v)\phi$ is Schwartz-Bruhat when $\phi$ is.
As $w=(w_1,w_2)$ is at $v$, $x\cdot w_2=\overline{x_v\cdot w_{2,v}}$ and
$w_1\cdot w_2=\overline{w_{1,v}\cdot w_{2,v}}$, and $r(x+w_1)=r(x)$ when $r$ does
not depend on the $v$-component; this gives the formula.

(2) The character $\psi_v$ is continuous, and it is not $1$ by the first step
of the proof of Lemma~\ref{asm-theta-agree-coinvariants-eigen}(3). By
Lemma~\ref{asm-theta-agree-dictionary-coordinates}(4), $e_v$ preserves the
forms. Theorem~\ref{asm-theta-agree-dictionary-stone-von-neumann}(2), for
$k=F_v$, $\psi_v$, $\kappa=I$, the set $\mu$ of the positions of
$n\times d$ matrices and the isomorphism $e_v$, gives a linear map $\iota$,
bijective between the Schwartz-Bruhat functions, with
$\iota(\rho_v(e_v(w))\phi)=\rho_v(w)\iota\phi$. As
$e_v(\operatorname{co}(u),\operatorname{co}'(u))=e(u)_v$, $\iota$ is a
Heisenberg chart.

(3) Let $\iota$ be a Heisenberg chart. Let $u$ be at $v$, and let $\phi$ and
$\phi_1$ be Schwartz-Bruhat with $\rho(u)$ acting at $v$ by
$\phi\mapsto\phi_1$. As $\rho(u)$ is the operator of
Theorem~\ref{asm-theta-agree-evaluation-heisenberg} at $e(u)$, which is at $v$,
(1) shows that $\rho(u)$ also acts at $v$ by $\phi\mapsto\rho_v(e(u)_v)\phi$, so
$\phi_1=\rho_v(e(u)_v)\phi$ by Lemma~\ref{asm-theta-agree-finite-unique}, and
$\iota\phi_1=\rho_v(\operatorname{co}(u),\operatorname{co}'(u))\iota\phi$. If
$\pi_t(u)=0$, then $\operatorname{co}'(u)=0$ by
Lemma~\ref{asm-theta-agree-dictionary-coordinates}(1), and
$\iota\phi_1(Y)=\iota\phi(Y+\operatorname{co}(u))$. If $\pi_{t\circ c}(u)=0$, then
$\operatorname{co}(u)=0$ and
$\iota\phi_1(Y)=\psi_v(\operatorname{tr}(Y^{\mathsf T}\operatorname{co}'(u)))\,\iota\phi(Y)$.
For $x$ at $v$ with $\pi_t(x)=0$, $\operatorname{co}'(x)=0$, and part (2) of
that lemma gives
$\psi(\langle x,u\rangle)=\psi_v(\operatorname{tr}(\operatorname{co}(x)^{\mathsf T}\operatorname{co}'(u)))$,
so
$\iota\phi_1(\operatorname{co}(x))=\psi(\langle x,u\rangle)\,\iota\phi(\operatorname{co}(x))$.
So $\iota$ is a linear model for $(\pi_{t\circ c},\pi_t,\operatorname{co})$.

Conversely, let $\iota$ be such a linear model, and $\phi$ Schwartz-Bruhat. Let
$Y_0$ be a matrix and $u_1$ the element at $v$ with coordinates $(Y_0,0)$, so
that $\pi_t(u_1)=0$. By (1), $\rho(u_1)$ acts at $v$ by
$\phi\mapsto\rho_v(e(u_1)_v)\phi$, a Schwartz-Bruhat function, so
$\iota(\rho_v(e(u_1)_v)\phi)(Y)=\iota\phi(Y+Y_0)=(\rho_v(Y_0,0)\iota\phi)(Y)$. Let
$Z_0$ be a matrix and $u_2$ the element at $v$ with coordinates $(0,Z_0)$, so
that $\pi_{t\circ c}(u_2)=0$. For a matrix $Y$ let $x$ be the element at $v$ with
coordinates $(Y,0)$; then, as above,
$\iota(\rho_v(e(u_2)_v)\phi)(Y)=\psi(\langle x,u_2\rangle)\iota\phi(Y)=(\rho_v(0,Z_0)\iota\phi)(Y)$.
Now let $u$ be at $v$ with coordinates $(Y_0,Z_0)$; then $u=u_1+u_2$, and
$[e(u_1)_v,e(u_2)_v]=[(Y_0,0),(0,Z_0)]$ by
Lemma~\ref{asm-theta-agree-dictionary-coordinates}(4). By the composition rule,
\[
\psi_v(\tfrac12[(Y_0,0),(0,Z_0)])\,\iota(\rho_v(e(u)_v)\phi)
=\iota(\rho_v(e(u_1)_v)\rho_v(e(u_2)_v)\phi)
=\rho_v(Y_0,0)\rho_v(0,Z_0)\iota\phi
=\psi_v(\tfrac12[(Y_0,0),(0,Z_0)])\,\rho_v(Y_0,Z_0)\iota\phi .
\]
As $e(u)_v$ and the coordinates depend only on the component of $u$ at $v$,
$\iota$ is a Heisenberg chart.

(4) As $h(w^{\mathsf T},w'^{\mathsf T})$ for $J'$ and $J$ is $h(w,w')$ for $J$
and $J'$, the isomorphism $e^{\mathsf T}$ has the property of the situation, and
$t(c(J'))=(-\varepsilon)t(J')^{\mathsf T}$,
$t(c(J))=-(-\varepsilon)t(J)^{\mathsf T}$. The coordinates of $u^{\mathsf T}$
are $-t(J')\pi_{t\circ c}(u)^{\mathsf T}t(J)^{\mathsf T}=\operatorname{co}(u)^{\mathsf T}$
and $-\pi_t(u)^{\mathsf T}=\operatorname{co}'(u)^{\mathsf T}$. For a function $\Phi$
on $M_{n\times d}(F_v)$ put $\Phi^{\mathsf T}(Z)=\Phi(Z^{\mathsf T})$. As
$\operatorname{tr}(Y^{\mathsf T}Z)=\operatorname{tr}(YZ^{\mathsf T})$, a direct
computation gives
$(\rho_v(Y_0,Z_0)\Phi)^{\mathsf T}=\rho_v(Y_0^{\mathsf T},Z_0^{\mathsf T})\Phi^{\mathsf T}$.
So
$\iota^{\mathsf T}(\rho_v(e^{\mathsf T}(u^{\mathsf T})_v)\phi)=(\iota(\rho_v(e(u)_v)\phi))^{\mathsf T}
=\rho_v(\operatorname{co}(u)^{\mathsf T},\operatorname{co}'(u)^{\mathsf T})\,\iota^{\mathsf T}\phi$.
\end{proof}
% Lean: (1) is ad_heis (the field Heisenberg_2 of AdSt), through schr_tens; the composition rule is
%   Mod.schr_schr. (2) is exists_lm, which applies model_existence to the isomorphism of exists_e,
%   with psi_v continuous and not 1 by the field psi of AdSt. (3) is linBar_of_lm (with
%   action_uniqueness) and lm_of_linBar, stated in the Lean for the second kind for (t,1,1). (4) is
%   Grams.swap and LM.swap.

\begin{lemma}\label{asm-theta-agree-dictionary-schur}
In Situation~\ref{asm-theta-agree-dictionary-situation}, let $\kappa$ be a
finite set, let $\rho(w)$, $w\in\mathbf A^\kappa\times\mathbf A^\kappa$, be the
operators of Theorem~\ref{asm-theta-agree-evaluation-heisenberg} for $\psi$, and
call a linear form on the functions on $\mathbf A^\kappa$ admissible as there.
\begin{enumerate}
\item Let $\phi_0\neq0$ be a function on $F_v^\kappa$ and $r$ a function on
  $\mathbf A^\kappa$ that does not depend on the $v$-component, such that
  $\phi_0\otimes r$ is Schwartz-Bruhat. Then $\phi\otimes r$ is Schwartz-Bruhat
  for every Schwartz-Bruhat function $\phi$ on $F_v^\kappa$.
\item Let $T$ be a linear operator on the functions on $\mathbf A^\kappa$ that
  maps Schwartz-Bruhat functions to Schwartz-Bruhat functions, with
  $\lambda\circ T$ admissible for every admissible $\lambda$, and let $s$ be a
  map of $\mathbf A^\kappa\times\mathbf A^\kappa$ to itself with
  $T\rho(w)\Phi=\rho(s(w))T\Phi$ for every $w$ and every Schwartz-Bruhat $\Phi$.
  Let $L$ and $L'$ be linear operators on the functions on $F_v^\kappa$ that map
  Schwartz-Bruhat functions to Schwartz-Bruhat functions, with $L'L\phi=\phi$
  and $LL'\phi=\phi$ for Schwartz-Bruhat $\phi$, and let $s_v$ be a map of
  $F_v^\kappa\times F_v^\kappa$ to itself with
  $L\rho_v(w)\phi=\rho_v(s_v(w))L\phi$ for every $w$ and every Schwartz-Bruhat
  $\phi$. Suppose that for every $w$, $s(w)-w$ is at $v$ and
  $s(w)_v=s_v(w_v)$. Then there is $a\in\mathbf C$ such that $T$ acts at $v$ by
  $\phi\mapsto a\,L\phi$ for every Schwartz-Bruhat $\phi$.
\end{enumerate}
\end{lemma}
\leanref{Development/WeilBound/CoreStatements/LocallyConstantCoinvariants/PfDictDef.lean}{PositivityTheory.KCo.AdSt}
\leanref{Development/WeilBound/CoreStatements/LocallyConstantCoinvariants/PfDictAd.lean}{PositivityTheory.KCo.adjoint_statement}
\leanref{Development/WeilBound/CoreStatements/LocallyConstantCoinvariants/PfDictAd.lean}{PositivityTheory.KCo.ad_rest}
\leanref{Development/WeilBound/CoreStatements/LocallyConstantCoinvariants/PfDictAd.lean}{PositivityTheory.KCo.ad_act}
\begin{proof}
For a linear operator $S$ on the functions on $F_v^\kappa$ let $S_*$ be the
operator $(S_*\Phi)(x)=(S\operatorname{res}_x\Phi)(x_v)$ on the functions on
$\mathbf A^\kappa$, with $\operatorname{res}_x$ as in
Lemma~\ref{asm-theta-agree-finite-inputs-slice}. Directly from the
definitions:
\begin{itemize}
\item[(i)] $S_*(\phi\otimes r)=(S\phi)\otimes r$ for every $\phi$ and every $r$
  that does not depend on the $v$-component, as
  $\operatorname{res}_x(\phi\otimes r)=r(x)\phi$;
\item[(ii)] if $S$ maps Schwartz-Bruhat functions to Schwartz-Bruhat functions,
  then so does $S_*$, and $\lambda\circ S_*$ is admissible for every admissible
  $\lambda$: a Schwartz-Bruhat function on $\mathbf A^\kappa$ is a finite sum of
  functions $f\otimes\bigotimes_u\phi_u$ (proof of
  Lemma~\ref{asm-theta-agree-coinvariants-slices}), and $S_*$ replaces the
  factor $\phi_v$ by $S\phi_v$;
\item[(iii)] if $S'$ also maps Schwartz-Bruhat functions to Schwartz-Bruhat
  functions and $SS'\phi=S''\phi$ for Schwartz-Bruhat $\phi$, then
  $S_*S'_*\Phi=S''_*\Phi$ for Schwartz-Bruhat $\Phi$, by the same description;
  and $1_*$ is the identity;
\item[(iv)] if $w-w'$ is at $v$ and $S\rho_v(w_v)\phi=\rho_v(w'_v)S\phi$ for
  Schwartz-Bruhat $\phi$, then $S_*\rho(w)\Phi=\rho(w')S_*\Phi$ for
  Schwartz-Bruhat $\Phi$, since
  $\rho(w)(\phi\otimes r)=(\rho_v(w_v)\phi)\otimes\rho(w^\circ)r$, where $w^\circ$
  is $w$ with its components at $v$ replaced by $0$, and $w^\circ=w'^\circ$.
\end{itemize}

(1) Choose $y$ with $\phi_0(y)\neq0$, and let $S$ be the operator
$\phi'\mapsto\phi'(y)\phi_0(y)^{-1}\phi$. It maps Schwartz-Bruhat functions to
Schwartz-Bruhat functions, so by (i) and (ii)
$\phi\otimes r=(S\phi_0)\otimes r=S_*(\phi_0\otimes r)$ is Schwartz-Bruhat.

(2) By (ii), $L'_*T$ maps Schwartz-Bruhat functions to Schwartz-Bruhat
functions, and $\lambda\circ L'_*T$ is admissible for admissible $\lambda$. For
Schwartz-Bruhat $\phi$ and every $w$,
$L'\rho_v(s_v(w))\phi=L'\rho_v(s_v(w))LL'\phi=L'L\rho_v(w)L'\phi=\rho_v(w)L'\phi$,
using that $\rho_v(w)L'\phi$ is Schwartz-Bruhat
(Lemma~\ref{asm-theta-agree-dictionary-charts}(1)). As $s(w)-w$ is at $v$ and
$s(w)_v=s_v(w_v)$, (iv) gives
$L'_*T\rho(w)\Phi=L'_*\rho(s(w))T\Phi=\rho(w)L'_*T\Phi$ for Schwartz-Bruhat
$\Phi$. The character $\psi$ is continuous, trivial on $F$ and not $1$, so
Theorem~\ref{asm-theta-agree-evaluation-heisenberg}(5) gives $a\in\mathbf C$ with
$L'_*T\Phi=a\Phi$ for Schwartz-Bruhat $\Phi$. By (iii), as $T\Phi$ is
Schwartz-Bruhat, $T\Phi=L_*L'_*T\Phi=a\,L_*\Phi$. For Schwartz-Bruhat $\phi$ and
$r$ that does not depend on the $v$-component, with $\phi\otimes r$
Schwartz-Bruhat, (i) gives $T(\phi\otimes r)=(aL\phi)\otimes r$.
\end{proof}
% Lean: the operator S_* is lift v S (lift_tens, lift_mem, lift_adm, lift_lift, lift_one,
%   lift_schr with schr_tens); (1) is ad_rest, the field rest of AdSt; (2) is ad_act, the field act,
%   which applies GlobalTheory.CompanionTheory.AdelicSchur.schur to lift v L' * T.

\begin{lemma}\label{asm-theta-agree-dictionary-diagonal}
In Situation~\ref{asm-theta-agree-dictionary-situation}, let $\kappa$ be a
finite set, let $\rho$ be the operators of
Theorem~\ref{asm-theta-agree-evaluation-heisenberg} for $\psi$ on the functions on
$\mathbf A^{\kappa\sqcup\kappa}$, and for functions $\Phi_1$ and $\Phi_2$ on
$\mathbf A^\kappa$, or on $F_v^\kappa$, put
$(\Phi_1\boxtimes\Phi_2)(x,y)=\Phi_1(x)\Phi_2(y)$. There are an admissible
linear form $l$ on the functions on $\mathbf A^{\kappa\sqcup\kappa}$, a
Schwartz-Bruhat function $\phi$ on $F_v^\kappa$ and a function $r$ on
$\mathbf A^\kappa$ that does not depend on the $v$-component, such that:
$l(\rho((x,x),(y,-y))\Phi)=l(\Phi)$ for all $x,y\in\mathbf A^\kappa$ and every
Schwartz-Bruhat $\Phi$ on $\mathbf A^{\kappa\sqcup\kappa}$; the functions
$\phi\otimes r$ and $(\phi\boxtimes\phi)\otimes(r\boxtimes r)$ are
Schwartz-Bruhat; and $l((\phi\otimes r)\boxtimes(\phi\otimes r))\neq0$.
\end{lemma}
\leanref{Development/WeilBound/CoreStatements/LocallyConstantCoinvariants/PfDictAd.lean}{PositivityTheory.KCo.lDiag_box_ne}
\leanref{Development/AbelianVarietyTheory/Proofs/DivisorProofs/ExistsDiagonalLagrangianInvariantForm.lean}{GlobalTheory.DiagForm.lDiag}
\leanref{Development/AbelianVarietyTheory/Proofs/DivisorProofs/ExistsDiagonalLagrangianInvariantForm.lean}{GlobalTheory.DiagForm.lDiag_adm}
\leanref{Development/AbelianVarietyTheory/Proofs/DivisorProofs/ExistsDiagonalLagrangianInvariantForm.lean}{GlobalTheory.DiagForm.lDiag_schr}
\begin{proof}
Let $l(\Phi)=\int_{\mathbf A^\kappa}\Phi(x,x)\,dx$ for a Haar measure on
$\mathbf A^\kappa$. For $\Phi=f\otimes g$ it is the product of
$\int f(x,x)\,dx$ over $F_\infty^\kappa$, which is continuous in $f$, and of
$\int g(x,x)\,dx$ over $\mathbf A_f^\kappa$, so $l$ is admissible. The function
$\rho((x,x),(y,-y))\Phi$ takes the value $\Phi(z+x,z+x)$ at $(z,z)$, so the
invariance holds by that of Haar measure. Let $\phi$ be the indicator function
of $\mathcal O_v^\kappa$, let $f\geq0$ be a smooth function with compact support
on $F_\infty^\kappa$ with $f(0)=1$, and let
$r(x)=f(x_\infty)\prod_{u\neq v}\mathbf 1_{\mathcal O_u^\kappa}(x_u)$, the product
over the finite places $u\neq v$. Then $r$ does not depend on the
$v$-component, $\phi\otimes r$ and
$(\phi\boxtimes\phi)\otimes(r\boxtimes r)=(\phi\otimes r)\boxtimes(\phi\otimes r)$
are Schwartz-Bruhat, and $l((\phi\otimes r)\boxtimes(\phi\otimes r))$ is the
positive number $\int f^2$ times the volume of $\prod_u\mathcal O_u^\kappa$.
\end{proof}
% Lean: the field diag of AdSt in adjoint_statement, with l = lDiag (the integral of resD over
%   Vi x Vf), phi = stdF kappa v, r = pure f (update stdF v 1) for a bump function f, indep_pure,
%   pure_update and lDiag_box_ne.

\begin{lemma}\label{asm-theta-agree-dictionary-scalar}
In Situation~\ref{asm-theta-agree-dictionary-situation}, call a linear form on
the functions on $\mathbf A^I$ admissible as in
Theorem~\ref{asm-theta-agree-evaluation-heisenberg}. Let $T$ be a linear
operator on the functions on $\mathbf A^I$ that maps Schwartz-Bruhat functions
to Schwartz-Bruhat functions, with $\lambda\circ T$ admissible for every
admissible $\lambda$, and let $h\in U_{J'}$, with $h-1$ at $v$, be such that
$T\rho(u)\Phi=\rho(uh^{\mathsf T})T\Phi$ for every $u\in\mathbb W$ and every
Schwartz-Bruhat $\Phi$. Let $\iota$ be a Heisenberg chart for $e$ at $v$
(Definition~\ref{asm-theta-agree-dictionary-chart}), and let $j$ be a linear map
from the functions on $M_{n\times d}(F_v)$ to the functions on $F_v^I$ that maps
Schwartz-Bruhat functions to Schwartz-Bruhat functions, with $\iota(j\Phi)=\Phi$
and $j(\iota\phi)=\phi$ for Schwartz-Bruhat $\Phi$ and $\phi$. Put
$h_0=\pi_t(h)$ and $(R(h_0)\Phi)(Y)=\Phi(Yh_0)$. Then there is $a\in\mathbf C$
such that $T$ acts at $v$ by $\phi\mapsto a\,j(R(h_0)\iota\phi)$ for every
Schwartz-Bruhat $\phi$.
\end{lemma}
\leanref{Development/WeilBound/CoreStatements/LocallyConstantCoinvariants/EqZeroOfLocW.lean}{PositivityTheory.KCo.acts_scalar}
\leanref{Development/WeilBound/CoreStatements/LocallyConstantCoinvariants/KdExistsInv.lean}{PositivityTheory.KCo.rmL_schr}
\begin{proof}
Let $e_v$ be as in Lemma~\ref{asm-theta-agree-dictionary-coordinates}(4), and
define $s_v$ on $F_v^I\times F_v^I$ by
$s_v(e_v(Y,Z))=e_v(Yh_0^{-1},Zh_0^{\mathsf T})$. Put
$s(w)=e(e^{-1}(w)h^{\mathsf T})$, $L=j\circ R(h_0)\circ\iota$ and
$L'=j\circ R(h_0^{-1})\circ\iota$. These satisfy the hypotheses of
Lemma~\ref{asm-theta-agree-dictionary-schur}(2) for $\kappa=I$, as follows.

The relation $T\rho(w)=\rho(s(w))T$ on Schwartz-Bruhat functions is the
hypothesis on $T$ for $u=e^{-1}(w)$. The operators $L$ and $L'$ map
Schwartz-Bruhat functions to Schwartz-Bruhat functions and are inverse to each
other on them, as $R(h_0)R(h_0^{-1})=1$. Let $u\in\mathbb W$ have coordinates
$(Y,Z)$. For Schwartz-Bruhat $\Psi$ on $M_{n\times d}(F_v)$, the functions
$j\rho_v(Y,Z)\Psi$ and $\rho_v(e(u)_v)j\Psi$ are Schwartz-Bruhat with the same
image $\rho_v(Y,Z)\Psi$ under $\iota$, so they are equal. A direct computation,
with $\operatorname{tr}((Y'h_0)^{\mathsf T}Z)=\operatorname{tr}(Y'^{\mathsf T}Zh_0^{\mathsf T})$,
gives $R(h_0)\rho_v(Y,Z)=\rho_v(Yh_0^{-1},Zh_0^{\mathsf T})R(h_0)$. Hence, for
Schwartz-Bruhat $\phi$,
\[
L\rho_v(e(u)_v)\phi=j\,R(h_0)\rho_v(Y,Z)\iota\phi
=j\,\rho_v(Yh_0^{-1},Zh_0^{\mathsf T})R(h_0)\iota\phi
=\rho_v(e_v(Yh_0^{-1},Zh_0^{\mathsf T}))L\phi=\rho_v(s_v(e(u)_v))L\phi,
\]
using Lemma~\ref{asm-theta-agree-dictionary-coordinates}(4) and the element at
$v$ with coordinates $(Yh_0^{-1},Zh_0^{\mathsf T})$. Every element of
$F_v^I\times F_v^I$ is $e(u)_v$ for some $u$, by that lemma. Finally, for
$w=e(u)$, $s(w)-w=e(u(h-1)^{\mathsf T})$ is at $v$, as $h-1$ is, and by parts
(3) and (4) of that lemma
$s(w)_v=e_v(\operatorname{co}(uh^{\mathsf T}),\operatorname{co}'(uh^{\mathsf T}))=e_v(Yh_0^{-1},Zh_0^{\mathsf T})=s_v(w_v)$.
Lemma~\ref{asm-theta-agree-dictionary-schur}(2) gives $a$.
\end{proof}

\begin{lemma}\label{asm-theta-agree-dictionary-action}
In Situation~\ref{asm-theta-agree-dictionary-situation}, let $\chi$ be a Hecke
character of $E$, and let $\omega$ be a representation of $U_{J'}$ on the
functions on $\mathbf A^I$ whose operators satisfy condition (a) of
Definition~\ref{asm-theta-agree-models}(6) in place of the operators
$\omega(1,h)$, for the isomorphism $e$, the matrices $J$ and $J'$ in place of
$J_X$ and $J_W$, and the character $\chi$ in place of $\chi_X$. Let $\iota$ be
a Heisenberg chart for $e$ at $v$, let $h\in U_{J'}$ with $h-1$ at $v$, and put
$h_0=\pi_t(h)$. Then $\omega(h)$ reads as
$\Phi\mapsto\chi_t(\det h_0)\,|\det h_0|^{n/2}\,\Phi(\,\cdot\,h_0)$ through
$\iota$ (Definition~\ref{asm-theta-agree-finite-split-models}(3)).
\end{lemma}
\leanref{Development/WeilBound/CoreStatements/LocallyConstantCoinvariants/ActsFormula.lean}{PositivityTheory.KCo.acts_formula}
\leanref{Development/AbelianVarietyTheory/Definitions/Divisors/HilbertNinetyQuotientHomEqIff.lean}{GlobalTheory.WeilM.IsUWeilB}
\begin{proof}
Write $R(g_0)\Phi=\Phi(\,\cdot\,g_0)$, and $|z|=|N_{E/F}(z)|_{\mathbf A}$ for
$z\in\mathbf A_E^\times$. Condition (a) of
Definition~\ref{asm-theta-agree-models}(6) gives the following. Each
$\omega(h)$ maps Schwartz-Bruhat functions to Schwartz-Bruhat functions,
$\lambda\circ\omega(h)$ is admissible for admissible $\lambda$, and
$\omega(h)\rho(u)\Phi=\rho(uh^{\mathsf T})\omega(h)\Phi$ for Schwartz-Bruhat
$\Phi$. Let $J^\square=\operatorname{diag}(J',-J')$, let $G$ be the group of the
$x\in GL_{2d}(\mathbf A_E)$ with $x^{\mathsf T}J^\square c(x)=J^\square$, and let
$\Delta=\left(\begin{smallmatrix}1_d\\1_d\end{smallmatrix}\right)$. There are an
isomorphism $e^\square$ of $\mathbb W^\square=M_{n\times2d}(\mathbf A_E)$ onto
$\mathbf A^{I\sqcup I}\times\mathbf A^{I\sqcup I}$ as in the situation, for $J$
and $J^\square$, with
$e^\square(w_1\ w_2)=\bigl((e_1(w_1),e_1(w_2)),(e_2(w_1),-e_2(w_2))\bigr)$ and
Schrödinger operators $\rho^\square$, and a representation $\omega^\square$ of
$G$ on the functions on $\mathbf A^{I\sqcup I}$, such that: each
$\omega^\square(x)$ maps Schwartz-Bruhat functions to Schwartz-Bruhat
functions, $\lambda\circ\omega^\square(x)$ is admissible for admissible
$\lambda$, and $\omega^\square(x)\rho^\square(u)\Phi=\rho^\square(ux^{\mathsf T})\omega^\square(x)\Phi$
for Schwartz-Bruhat $\Phi$; for every admissible $\lambda$ with
$\lambda(\rho^\square(z\Delta^{\mathsf T})\Phi)=\lambda(\Phi)$ for all
$z\in\mathbb W$ and Schwartz-Bruhat $\Phi$, and all $p\in G$ and
$\tau\in GL_d(\mathbf A_E)$ with $p\Delta=\Delta\tau$,
$\lambda(\omega^\square(p)\Phi)=\chi(\det\tau)\,|\det\tau|^{n/2}\lambda(\Phi)$ for
Schwartz-Bruhat $\Phi$; and
$\omega^\square(\operatorname{diag}(h',1))(\Phi_1\boxtimes\Phi_2)=(\omega(h')\Phi_1)\boxtimes\Phi_2$
for $h'\in U_{J'}$ and Schwartz-Bruhat $\Phi_1$ and $\Phi_2$. The matrix
$t(J^\square)=\operatorname{diag}(t(J'),-t(J'))$ is invertible and
$t(c(J^\square))=-\varepsilon\,t(J^\square)^{\mathsf T}$, so the situation also
holds for $J$, $J^\square$, $\varepsilon$, the set $I\sqcup I$ and $e^\square$.

\emph{Step 0: the scalar.} Let $j$ be a linear map from the functions on
$M_{n\times d}(F_v)$ to the functions on $F_v^I$ that is inverse to $\iota$ on
the Schwartz-Bruhat functions; it exists as $\iota$ is bijective between them.
By Lemma~\ref{asm-theta-agree-dictionary-scalar} for $T=\omega(h)$ there is
$c_1\in\mathbf C$ such that $\omega(h)$ acts at $v$ by
$\phi\mapsto c_1\,j(R(h_0)\iota\phi)$ for Schwartz-Bruhat $\phi$. Put
$s=\chi_t(\det h_0)|\det h_0|^{n/2}$. It suffices to show $c_1=s$: then
$\phi_1=s\,j(R(h_0)\iota\phi)$ is Schwartz-Bruhat with $\iota\phi_1=sR(h_0)\iota\phi$.

\emph{Step 1: two elements of $G$.} Let
$X_0=\left(\begin{smallmatrix}1&0\\-1&1\end{smallmatrix}\right)$ and
$P_0=\left(\begin{smallmatrix}h_0&0\\h_0-1&1\end{smallmatrix}\right)$ in
$GL_{2d}(F_v)$, in blocks of size $d$. By
Lemma~\ref{asm-theta-agree-dictionary-split}(4) for $J^\square$ there are
$x,p\in G$ with $x-1$ and $p-1$ at $v$, $\pi_t(x)=X_0$ and $\pi_t(p)=P_0$. Let
$h^\square=\operatorname{diag}(h,1)$, which lies in $G$. The elements $xp$ and
$h^\square x$ of $G$ differ from $1$ by matrices at $v$, and
$\pi_t(xp)=X_0P_0=\left(\begin{smallmatrix}h_0&0\\-1&1\end{smallmatrix}\right)=\operatorname{diag}(h_0,1)X_0=\pi_t(h^\square x)$,
so $xp=h^\square x$ by the uniqueness in that lemma.

\emph{Step 2: the Siegel rule for $p$.} By
Lemma~\ref{asm-theta-agree-dictionary-split}(2) there is
$\tau\in GL_d(\mathbf A_E)$ with $\tau-1$ at $v$, $\pi_t(\tau)=h_0$ and
$\pi_{t\circ c}(\tau)=1$. The matrix $p\Delta-\Delta\tau=(p-1)\Delta-\Delta(\tau-1)$
is at $v$, and $\pi_t(p\Delta)=P_0\Delta=\Delta h_0$. With $A'=t(J')$,
$\pi_t(J^\square)\Delta=\left(\begin{smallmatrix}A'\\-A'\end{smallmatrix}\right)$
and
$P_0^{\mathsf T}\left(\begin{smallmatrix}A'\\-A'\end{smallmatrix}\right)=\left(\begin{smallmatrix}h_0^{\mathsf T}A'-(h_0^{\mathsf T}-1)A'\\-A'\end{smallmatrix}\right)=\pi_t(J^\square)\Delta$;
as $\pi_{t\circ c}(p)=\pi_t(J^\square)^{-1}(P_0^{\mathsf T})^{-1}\pi_t(J^\square)$
by Lemma~\ref{asm-theta-agree-dictionary-split}(4), $\pi_{t\circ c}(p)\Delta=\Delta$.
So $p\Delta=\Delta\tau$ by part (1) of that lemma. The element $\det\tau-1$ is
at $v$, with $\pi_t(\det\tau)=\det h_0$ and $\pi_{t\circ c}(\det\tau)=1$, so
$\chi(\det\tau)=\chi_t(\det h_0)$ and $|\det\tau|=|\det h_0|$ by part (3).

\emph{Step 3: the form.} Let $l$, $\phi_*$ and $r$ be as in
Lemma~\ref{asm-theta-agree-dictionary-diagonal} for $\kappa=I$. For
$z\in\mathbb W$, $z\Delta^{\mathsf T}=(z\ z)$ and
$e^\square(z\ z)=((e_1(z),e_1(z)),(e_2(z),-e_2(z)))$, so
$l(\rho^\square(z\Delta^{\mathsf T})\Phi)=l(\Phi)$ for Schwartz-Bruhat $\Phi$. By
Step 2, $l(\omega^\square(p)\Phi)=s\,l(\Phi)$. As $x^{-1}h^\square=px^{-1}$ by
Step 1, for Schwartz-Bruhat $\Phi$
\[
l(\omega^\square(x^{-1})\omega^\square(h^\square)\Phi)=l(\omega^\square(p)\omega^\square(x^{-1})\Phi)=s\,l(\omega^\square(x^{-1})\Phi).
\]

\emph{Step 4: an invariance.} Let $u\in\mathbb W$ be at $v$ with
$\pi_{t\circ c}(u)=0$. The element $(u\ u)x^{\mathsf T}-(u\ 0)=(u\ u)(x-1)^{\mathsf T}+(0\ u)$
is at $v$, its image under $\pi_t$ is
$(\pi_t(u)\ \pi_t(u))X_0^{\mathsf T}-(\pi_t(u)\ 0)=0$, and its image under
$\pi_{t\circ c}$ is $0$; so $(u\ u)x^{\mathsf T}=(u\ 0)$. Hence
$\rho^\square(u\ 0)\Phi=\omega^\square(x)\rho^\square(u\ u)\omega^\square(x^{-1})\Phi$
for Schwartz-Bruhat $\Phi$, and by Step 3
\[
l(\omega^\square(x^{-1})\rho^\square(u\ 0)\Phi)=l(\rho^\square(u\ u)\omega^\square(x^{-1})\Phi)=l(\omega^\square(x^{-1})\Phi).
\]

\emph{Step 5: a nonzero value.} By
Lemma~\ref{asm-theta-agree-dictionary-schur}(1), $\phi\otimes r$ is
Schwartz-Bruhat for every Schwartz-Bruhat $\phi$, as $\phi_*\neq0$; so
$(\phi_1\otimes r)\boxtimes(\phi_2\otimes r)=(\phi_1\boxtimes\phi_2)\otimes(r\boxtimes r)$
is Schwartz-Bruhat for Schwartz-Bruhat $\phi_1$ and $\phi_2$. We show that
$l(\omega^\square(x^{-1})((\phi_1\otimes r)\boxtimes(\phi_2\otimes r)))\neq0$ for
some such $\phi_1$ and $\phi_2$. Suppose not, and let
$\Lambda(\Psi)=l(\omega^\square(x^{-1})(\Psi\otimes(r\boxtimes r)))$ for functions
$\Psi$ on $F_v^{I\sqcup I}$. For $w=((a',a''),(b',b''))$,
$\rho_v(w)(\phi_*\boxtimes\phi_*)=\rho_v(a',b')\phi_*\boxtimes\rho_v(a'',b'')\phi_*$,
so $\Lambda(\rho_v(w)(\phi_*\boxtimes\phi_*))=0$ for every $w$, and
$\phi_*\boxtimes\phi_*\neq0$. By
Theorem~\ref{asm-theta-agree-dictionary-stone-von-neumann}(1), $\Lambda$
vanishes on the Schwartz-Bruhat functions. Let $\iota_2$ be a Heisenberg chart
for $e^\square$ at $v$ (Lemma~\ref{asm-theta-agree-dictionary-charts}(2)) and
$j_2$ an inverse of it on the Schwartz-Bruhat functions. By
Lemma~\ref{asm-theta-agree-dictionary-scalar} for the doubled data and
$T=\omega^\square(x)$, there is $c_2$ such that $\omega^\square(x)$ acts at $v$
by $\Psi\mapsto c_2\,j_2(R(X_0)\iota_2\Psi)$. As $r\boxtimes r$ does not depend
on the $v$-component and $(\phi_*\boxtimes\phi_*)\otimes(r\boxtimes r)$ is
Schwartz-Bruhat,
\[
l((\phi_*\otimes r)\boxtimes(\phi_*\otimes r))
=l\bigl(\omega^\square(x^{-1})\omega^\square(x)((\phi_*\boxtimes\phi_*)\otimes(r\boxtimes r))\bigr)
=\Lambda\bigl(c_2\,j_2(R(X_0)\iota_2(\phi_*\boxtimes\phi_*))\bigr)=0,
\]
which contradicts Lemma~\ref{asm-theta-agree-dictionary-diagonal}.

\emph{Step 6: the form at $0$.} Fix $\phi_1$ and $\phi_2$ as in Step 5, and put
$m(\phi)=l(\omega^\square(x^{-1})((\phi\otimes r)\boxtimes(\phi_2\otimes r)))$.
Let $y\in M_{n\times d}(F_v)$, let $u$ be the element at $v$ with coordinates
$(0,y)$, so that $\pi_{t\circ c}(u)=0$, and let $\Phi$ be Schwartz-Bruhat on
$M_{n\times d}(F_v)$. As in the proof of
Lemma~\ref{asm-theta-agree-dictionary-scalar},
$j(\rho_v(0,y)\Phi)=\rho_v(e(u)_v)j\Phi$. By
Lemma~\ref{asm-theta-agree-dictionary-charts}(1),
$(\rho_v(e(u)_v)j\Phi)\otimes r=\rho(u)(j\Phi\otimes r)$, and as
$e^\square(u\ 0)=((e_1(u),0),(e_2(u),0))$,
$(\rho(u)\Psi_1)\boxtimes\Psi_2=\rho^\square(u\ 0)(\Psi_1\boxtimes\Psi_2)$. So
Step 4 gives $m(j(\rho_v(0,y)\Phi))=m(j\Phi)$. By
Lemma~\ref{asm-theta-agree-dictionary-evaluation} for $\Phi\mapsto m(j\Phi)$
there is $a\in\mathbf C$ with $m(j\Phi)=a\,\Phi(0)$ for Schwartz-Bruhat $\Phi$;
hence $m(\phi)=a\,\iota\phi(0)$ for Schwartz-Bruhat $\phi$.

\emph{Step 7: conclusion.} Let $\Psi_1=(\phi_1\otimes r)\boxtimes(\phi_2\otimes r)$.
Then
$\omega^\square(h^\square)\Psi_1=(\omega(h)(\phi_1\otimes r))\boxtimes(\phi_2\otimes r)=(c_1j(R(h_0)\iota\phi_1)\otimes r)\boxtimes(\phi_2\otimes r)$.
By Steps 3 and 6,
\[
s\,m(\phi_1)=m\bigl(c_1\,j(R(h_0)\iota\phi_1)\bigr)=c_1a\,(R(h_0)\iota\phi_1)(0)=c_1a\,\iota\phi_1(0)=c_1m(\phi_1).
\]
As $m(\phi_1)\neq0$ by Step 5, $c_1=s$.
\end{proof}
% Lean: acts_formula, with the halves IsUWeilB.isWeilB of IsWeil. Step 0 is acts_scalar with
%   kd_exists_inv; Step 1 is X_0, P_0, h.exists_ug, exists_dbl_elem and h.ug_ext (hxp); Step 2 is
%   h.exists_gl, h.extM with h.unit_rel (hpt), hdet and h.idele; Step 3 is hA.diag, hl0, the
%   SiegelRule hR (hsie) and k1; Step 4 is k2; Step 5 is the by_contra with cyclicity (hC), schr_box,
%   tens_box and acts_scalar for x; Step 6 is exists_eq_mul_apply_zero (hm, hm'); Step 7 is f1, f2.

\begin{lemma}\label{asm-theta-agree-coinvariants-dictionary}
Let $B$ be a theta base for a CM field $E$
(Definition~\ref{asm-theta-agree-models}(1)), let $X$ be a Hermitian and $W$ a
skew-Hermitian lattice space, of ranks $n$ and $d$ and of some signatures
(Definition~\ref{asm-theta-agree-doubling}(1)), let $M$ be a model for $(X,W)$,
let $\chi_X$ and $\chi_W$ be Hecke characters of $E$, and let $\omega$ be a Weil
representation for $M$ with characters $(\chi_X,\chi_W)$
(Definition~\ref{asm-theta-agree-models}(6)). Let $v$ be a finite place of $F$
and $t\colon E\to F_v$ an $F$-embedding, and put $\nu=(\chi_W)_t$ and
$\alpha=(\chi_X)_t$ (Definition~\ref{asm-theta-agree-coinvariants-local}(4)),
which are unitary and smooth, as $\chi_W$ and $\chi_X$ are continuous with values
of absolute value $1$.
\begin{enumerate}
\item There is a linear model of the second kind for $(t,1,1)$
  (Definition~\ref{asm-theta-agree-models}(5)).
\item Every $g_0\in GL_n(F_v)$ is $\pi_t(g)$ for some $g\in U_X^{(v)}$, and every
  $h_0\in GL_d(F_v)$ is $\pi_t(h)$ for some $h\in U_W^{(v)}$.
\item Let $\iota$ be a linear model of the second kind for $(t,1,1)$, and let
  $g\in U_X^{(v)}$ and $h\in U_W^{(v)}$. Then $\omega(g,1)$ reads as
  $\omega_{\nu,\alpha}(\pi_t(g),1)$ and $\omega(1,h)$ reads as
  $\omega_{\nu,\alpha}(1,\pi_t(h))$ through $\iota$
  (Definitions~\ref{asm-theta-agree-coinvariants-local}(1)
  and~\ref{asm-theta-agree-finite-split-models}(3)).
\end{enumerate}
\end{lemma}
\leanref{Development/WeilBound/CoreStatements/LocallyConstantCoinvariants/LmOfLinBar.lean}{PositivityTheory.KCo.DictX}
\leanref{Development/WeilBound/ThetaLifts/LocalTheory/StandardValuations/VacuumImageIntertwiner.lean}{PositivityTheory.ThKG.StdDict.dictX}
\leanref{Development/WeilBound/CoreStatements/LocallyConstantCoinvariants/LmOfLinBar.lean}{PositivityTheory.KCo.dictX_of}
\begin{proof}
Let $w$ be an infinite place of $F$. A complex matrix $M$ with a Cartan
involution $\theta$ is invertible: if $My=0$ with $y\neq0$, then $x=\bar y$ has
$\operatorname{Re}((\theta x)^{\mathsf T}M\bar x)=0$. As $X$ is Hermitian and $W$
skew-Hermitian, $\tau_w(J_X)$ and $i\tau_w(J_W)$ have Cartan involutions, so
$J_X$ and $J_W$ are invertible, and so are $t(J_X)$ and $t(J_W)$. As
$c(J_X)^{\mathsf T}=J_X$ and $c(J_W)^{\mathsf T}=-J_W$,
$t(c(J_X))=t(J_X)^{\mathsf T}$ and $t(c(J_W))=-t(J_W)^{\mathsf T}$. So
Situation~\ref{asm-theta-agree-dictionary-situation} holds for $J=J_X$,
$J'=J_W$, $\varepsilon=1$, $I=\{1,\ldots,n\}\times\{1,\ldots,d\}$ and the
isomorphism $e$ of $M$, and the linear models for
$(\pi_{t\circ c},\pi_t,\operatorname{co})$ are the linear models of the second
kind for $(t,1,1)$. The characters $\nu$ and $\alpha$ are unitary and smooth by
Lemma~\ref{asm-theta-agree-dictionary-split}(3).

(1) By Lemma~\ref{asm-theta-agree-dictionary-charts}(2) there is a Heisenberg
chart for $e$ at $v$, and by part (3) of that lemma it is a linear model of the
second kind for $(t,1,1)$.

(2) Apply Lemma~\ref{asm-theta-agree-dictionary-split}(4) to $J_X$ with
$\varepsilon=1$ and to $J_W$ with $\varepsilon=-1$.

(3) By Lemma~\ref{asm-theta-agree-dictionary-charts}(3), $\iota$ is a
Heisenberg chart for $e$ at $v$. Let $h_0=\pi_t(h)$. By condition (a) of
Definition~\ref{asm-theta-agree-models}(6), the representation
$h'\mapsto\omega(1,h')$ of $U_W$ satisfies the hypotheses of
Lemma~\ref{asm-theta-agree-dictionary-action} with $\chi=\chi_X$. So
$\omega(1,h)$ reads through $\iota$ as
$\Phi\mapsto\alpha(\det h_0)|\det h_0|^{n/2}\Phi(\,\cdot\,h_0)$, which is
$\omega_{\nu,\alpha}(1,h_0)$.

Let $g_0=\pi_t(g)$. By Lemma~\ref{asm-theta-agree-dictionary-charts}(4), the
situation holds for $J_W$, $J_X$, $-1$ and the isomorphism
$e^{\mathsf T}\colon w\mapsto e(w^{\mathsf T})$ of $M_{d\times n}(\mathbf A_E)$,
and $\iota^{\mathsf T}\colon\phi\mapsto(Z\mapsto\iota\phi(Z^{\mathsf T}))$ is a
Heisenberg chart for $e^{\mathsf T}$ at $v$. By condition (b) of
Definition~\ref{asm-theta-agree-models}(6), the representation
$g'\mapsto\omega(g',1)$ of $U_X$ satisfies the hypotheses of
Lemma~\ref{asm-theta-agree-dictionary-action} for these data, with
$\chi=\chi_W$. So for Schwartz-Bruhat $\phi$, $\omega(g,1)$ acts at $v$ by
$\phi\mapsto\phi_1$ for a Schwartz-Bruhat $\phi_1$ with
$\iota^{\mathsf T}\phi_1(Z)=\nu(\det g_0)|\det g_0|^{d/2}\iota^{\mathsf T}\phi(Zg_0)$
for $Z\in M_{d\times n}(F_v)$. Taking $Z=Y^{\mathsf T}$ gives
$\iota\phi_1(Y)=\nu(\det g_0)|\det g_0|^{d/2}\iota\phi(g_0^{\mathsf T}Y)=(\omega_{\nu,\alpha}(g_0,1)\iota\phi)(Y)$.
% Lean: ThKG.StdDict.dictX = dictX_of model_existence cyclicity adjoint_statement
%   section_statement action_uniqueness (LmOfLinBar.lean). dictX_of: Grams from the Cartan
%   involutions (isUnit_det_of_cartan); (1) is linBar_of_lm of exists_lm (model_existence,
%   AdSt, a section sec from SecSt); (2) is IsSec.exists_ug; (3) passes to LM by lm_of_linBar
%   and applies acts_formula (ActsFormula.lean, PfActs.lean) to the halves hR and hL of IsWeil,
%   the latter through LM.swap. The same DictX gives the split-place dictionary of
%   Lemma asm-theta-agree-finite-inputs(1) (ThK.SplitDict.split_place_dictionary) and
%   TC.splitModel. Action uniqueness is Lemma asm-theta-agree-finite-unique.
\end{proof}

\begin{lemma}\label{asm-theta-agree-coinvariants-inputs}
Let $k$ be a nonarchimedean local field and $n,d\geq0$.
\begin{enumerate}
\item Let $\chi_1,\chi_2$ be characters of $k^\times$ with $\chi_1$ unitary and
  smooth and $\chi_1\neq1$. Then every Schwartz-Bruhat function on
  $M_{n\times d}(k)$ is $\Psi_0+z$ with $\Psi_0$ of full rank and
  $z\in\mathfrak a(\chi_1,\chi_2)$
  (Definition~\ref{asm-theta-agree-coinvariants-local}).
\item Let $n\leq d$, let $\chi_1,\chi_2$ be unitary and smooth with
  $\chi_1\neq1$, and let $N$ be a subspace of the functions on
  $M_{n\times d}(k)$ that contains $\mathfrak a(\chi_1,\chi_2)$ and contains
  $\omega_{\chi_1,\chi_2}(1,h)\Psi$ whenever $\Psi\in N$ is Schwartz-Bruhat and
  $h\in GL_d(k)$. Then either $N$ contains every function of full rank, or
  every function of full rank in $N$ lies in $\mathfrak a(\chi_1,\chi_2)$.
\item Let $B$ be a theta base for a CM field $E$
  (Definition~\ref{asm-theta-agree-models}(1)), let $X$ be a Hermitian and $W$ a
  skew-Hermitian lattice space, of ranks $n$ and $d$ and of some signatures
  (Definition~\ref{asm-theta-agree-doubling}(1)), let $M$ be a regular model for
  $(X,W)$, let $\chi_X$ and $\chi_W$ be Hecke characters of $E$, and let
  $\omega$ be a Weil representation for $M$ with characters $(\chi_X,\chi_W)$
  (Definition~\ref{asm-theta-agree-models}(6)). Let $v$ be a finite place and
  $t\colon E\to F_v$ an $F$-embedding. Then there is a linear model of the
  second kind for $(t,1,1)$
  (Definition~\ref{asm-theta-agree-models}(5)), and there are unitary smooth
  characters $\chi_1,\chi_2$ of $F_v^\times$ such that: if $n,d\geq1$ and
  $\chi_1=1$, then $(\chi_W)_t=1$ or $(\chi_W)_{t\circ c}=1$; and every linear
  model of the second kind for $(t,1,1)$ is a dictionary at $v$ with
  characters $(\chi_1,\chi_2)$.
\item Let $(B,G,\mathcal X,\omega,\omega')$ be a triple
  (Definition~\ref{asm-theta-agree-triple}). Then $\omega'(g,h)$ maps
  $\mathcal S^\circ(M')$ into itself for all $g\in U_{V'}$ and $h\in U_W$.
\end{enumerate}
\end{lemma}
\leanref{Development/WeilBound/CoreStatements/LocallyConstantCoinvariants/RightOnto.lean}{PositivityTheory.KCo.SurjSt}
\leanref{Development/AutomorphicSpectrum/Proofs/BlockDecomposition/KCoinvOfNarrow.lean}{GlobalTheory.CompanionTheory.KCoNarrow.IrrStNe}
\leanref{Development/WeilBound/CoreStatements/LocallyConstantCoinvariants/RightOnto.lean}{PositivityTheory.KCo.DictSt}
\leanref{Development/WeilBound/CoreStatements/LocallyConstantCoinvariants/RightOnto.lean}{PositivityTheory.KCo.SoSt}
\leanref{Development/WeilBound/CoreStatements/LocallyConstantCoinvariants/PfSurj.lean}{PositivityTheory.KCo.surjectivity_statement}
\leanref{Development/AutomorphicSpectrum/Proofs/BlockDecomposition/KCoinvOfNarrow.lean}{GlobalTheory.CompanionTheory.KCoNarrow.irrSt_narrow}
\leanref{Development/WeilBound/CoreStatements/LocallyConstantCoinvariants/ExistsExtZ.lean}{PositivityTheory.KCo.extZ}
\leanref{Development/WeilBound/CoreStatements/LocallyConstantCoinvariants/ExistsExtZ.lean}{PositivityTheory.KCo.irr_equiv}
\leanref{Development/WeilBound/CoreStatements/LocallyConstantCoinvariants/LmOfLinBar.lean}{PositivityTheory.KCo.dictSt_of}
\leanref{Development/WeilBound/CoreStatements/LocallyConstantCoinvariants/SoStatement.lean}{PositivityTheory.KCo.so_stable}
\leanref{Development/WeilBound/CoreStatements/LocallyConstantCoinvariants/SoStatement.lean}{PositivityTheory.KCo.so_statement}
\leanref{Development/WeilBound/CoreStatements/LocallyConstantCoinvariants/CompFinLiftE.lean}{PositivityTheory.KCo.symp_act}
\begin{proof}
(1) Let $\Psi$ be Schwartz-Bruhat, and put $\Psi^{(0)}=\Psi$. Every $Y$ is free
on $\emptyset$ (Lemma~\ref{asm-theta-agree-coinvariants-step}). For
$m=1,\ldots,n$, Lemma~\ref{asm-theta-agree-coinvariants-step} with
$S=\{1,\ldots,m-1\}$ and $i=m$ gives a Schwartz-Bruhat $\Psi^{(m)}$ that vanishes
at every $Y$ not free on $\{1,\ldots,m\}$, with
$\Psi^{(m-1)}-\Psi^{(m)}\in\mathfrak a(\chi_1,\chi_2)$. A matrix of rank different
from $n$ is not free on $\{1,\ldots,n\}$, so $\Psi_0=\Psi^{(n)}$ has full rank,
and $\Psi-\Psi_0\in\mathfrak a(\chi_1,\chi_2)$.

(2) Put $\mu=\chi_1^{-1}$, which is unitary, smooth and not $1$, and let $t=d-n$,
$R$, $\mathcal S(R)$, $\omega^R$, $\tau$, $\operatorname{Ind}(\tau)$, $R(h_1)$ and
$E$ be as in Lemma~\ref{asm-theta-agree-coinvariants-fibre} for $\mu$. For
$\phi\in\mathcal S(R)$ let $Z\phi(Y)=\phi(Y^{\mathsf T})$ if $Y\in M_{n\times d}(k)$
has rank $n$, and $Z\phi(Y)=0$ otherwise. As the matrices of rank $n$ form an
open set, $Z\phi$ is Schwartz-Bruhat, and it has full rank; $Z$ is injective; and
every function $\Phi$ of full rank is $Z\phi$, with $\phi(x)=\Phi(x^{\mathsf T})$.
As $(g^{\mathsf T}Y(h^{\mathsf T})^{-1})^{\mathsf T}=h^{-1}Y^{\mathsf T}g$,
\[
Z(\omega^R(g,h)\phi)(Y)=\chi_1(\det g)|\det g|^{d/2}\,\chi_1(\det h)^{-1}|\det h|^{-n/2}\,Z\phi\bigl(g^{\mathsf T}Y(h^{\mathsf T})^{-1}\bigr).
\]
So $Z(\omega^R(g,1)\phi)=\omega_{\chi_1,\chi_2}(g,1)Z\phi$, and
$Z(\omega^R(1,h)\phi)$ is a nonzero multiple of
$\omega_{\chi_1,\chi_2}(1,(h^{\mathsf T})^{-1})Z\phi$. Let $N'$ be the space of
the $\phi\in\mathcal S(R)$ with $Z\phi\in N$, and $\mathrm{Rel}_R$ the span of
the $\omega^R(g,1)\phi-\phi$. Then $\mathrm{Rel}_R\subset N'$, as $Z$ maps it into
$\mathfrak a(\chi_1,\chi_2)\subset N$, and $N'$ is stable under every
$\omega^R(1,h)$, by the hypothesis on $N$. By
Lemma~\ref{asm-theta-agree-coinvariants-fibre}, $E$ maps $\mathrm{Rel}_R$ onto
the space $\mathrm{Rel}$ of Lemma~\ref{asm-theta-agree-coinvariants-irreducible},
and $E(\omega^R(1,h)\phi)=R(h)E(\phi)$. So $E(N')$ contains $\mathrm{Rel}$ and is
stable under every $R(h)$, and by
Lemma~\ref{asm-theta-agree-coinvariants-irreducible}, $N'=\mathcal S(R)$ or
$N'\subset\mathrm{Rel}_R$. In the first case every function of full rank, being
some $Z\phi$, lies in $N$. In the second, a function of full rank in $N$ is
$Z\phi$ with $\phi\in N'\subset\mathrm{Rel}_R$, so it lies in
$Z(\mathrm{Rel}_R)\subset\mathfrak a(\chi_1,\chi_2)$.

(3) Put $\chi_1=(\chi_W)_t$ and $\chi_2=(\chi_X)_t$. As $\chi_W$ and $\chi_X$ are
continuous with values of absolute value $1$, $\chi_1$ and $\chi_2$ are unitary
and smooth. By Lemma~\ref{asm-theta-agree-coinvariants-dictionary}(1) there is a
linear model of the second kind for $(t,1,1)$, and the condition on $\chi_1=1$
holds as $\chi_1=(\chi_W)_t$. Let $\iota$ be such a model; it is bijective between
the Schwartz-Bruhat functions (Definition~\ref{asm-theta-agree-models}(5)). For
$g\in U_X^{(v)}$, condition (a) of a dictionary holds with $g_0=\pi_t(g)$, by
Lemma~\ref{asm-theta-agree-coinvariants-dictionary}(3). For
$g_0\in GL_n(F_v)$, Lemma~\ref{asm-theta-agree-coinvariants-dictionary}(2) gives
$g\in U_X^{(v)}$ with $\pi_t(g)=g_0$, and part (3) of that lemma applies to it.
Condition (b) follows in the same way from the statements for $W$.

(4) Let $\gamma=(g,h)\in U_{V'}\times U_W$, and write $\rho$ for the Schrödinger
operators of $M'$, $I=\{1,\ldots,a'\}\times\{1,\ldots,a+a'\}$, $v=v_1$ and
$t=\tau_1$. By conditions (a) and (b) of
Definition~\ref{asm-theta-agree-models}(6), and as
$\omega'(g,h)=\omega'(g,1)\omega'(1,h)$, $\omega'(\gamma)$ maps Schwartz-Bruhat
functions to Schwartz-Bruhat functions, and
$\omega'(\gamma)\rho(u)\Phi=\rho(\gamma u)\omega'(\gamma)\Phi$ for Schwartz-Bruhat
$\Phi$, where $\gamma u=guh^{\mathsf T}$. As $g^{\mathsf T}J_{V'}c(g)=J_{V'}$,
$h^{\mathsf T}J_Wc(h)=J_W$ and
$\langle x,y\rangle=-\operatorname{Tr}\operatorname{tr}(x^{\mathsf T}J_{V'}c(y)J_W^{\mathsf T})$,
the invariance of $\operatorname{tr}$ under cyclic permutations gives
$\langle\gamma x,\gamma y\rangle=\langle x,y\rangle$. As the triple datum is
proper, $V'$ and $W$ have signatures $(q,a'-q)$ and $(a+a'-p-q,p+q)$, so
$J_{V'}$ and $J_W$ are invertible.

By linearity it suffices to show $\omega'(\gamma)\Phi\in\mathcal S^\circ(M')$ for
$\Phi=f\otimes\bigotimes_w\phi_w$ as in
Definition~\ref{asm-theta-agree-models}(5), with $\iota$ a linear model of the
second kind for $(t,1,1)$ such that $\iota\phi_v(Y)=0$ for $Y$ of rank different
from $a'$. Let $r(x)=f(x_\infty)\prod_{w\neq v}\phi_w(x_w)$. It does not depend
on the $v$-component, $\Phi=\phi_v\otimes r$, and $\phi'\otimes r$ is
Schwartz-Bruhat for every Schwartz-Bruhat $\phi'$ on $F_v^I$. Put
$\Psi=\omega'(\gamma)\Phi$, a Schwartz-Bruhat function. By
Lemma~\ref{asm-theta-agree-coinvariants-slices} it suffices that every
$\operatorname{res}_z\Psi$ has full rank at $t$. It is Schwartz-Bruhat by
Lemma~\ref{asm-theta-agree-finite-inputs-slice}(3); we show
$\iota(\operatorname{res}_z\Psi)(Y)=0$ for $Y$ of rank different from $a'$.

By Lemma~\ref{asm-theta-agree-coinvariants-eigen}(1), $Y=\operatorname{co}(x)$
with $x$ at $v$ and $\pi_t(x)=0$. Let
$l(\phi')=\iota\bigl(\operatorname{res}_z(\omega'(\gamma)(\phi'\otimes r))\bigr)(\operatorname{co}(x))$,
a linear form, and $x_1=\gamma^{-1}x=g^{-1}x(h^{-1})^{\mathsf T}$, which is at $v$
with $\pi_t(x_1)=0$. Let $u$ be at $v$ with $\pi_{t\circ c}(u)=0$, and $\phi'$
Schwartz-Bruhat. By Lemma~\ref{asm-theta-agree-coinvariants-eigen}(2),
\[
\operatorname{res}_z\bigl(\omega'(\gamma)((\rho_v(u)\phi')\otimes r)\bigr)
=\operatorname{res}_z\bigl(\rho(\gamma u)\,\omega'(\gamma)(\phi'\otimes r)\bigr)
=\rho_v(\gamma u)\operatorname{res}_z\bigl(\omega'(\gamma)(\phi'\otimes r)\bigr).
\]
As $\gamma u$ is at $v$ with $\pi_{t\circ c}(\gamma u)=0$, $\rho(\gamma u)$ acts at
$v$ by $\rho_v(\gamma u)$, and $\iota$ is a linear model, this gives
$l(\rho_v(u)\phi')=\psi(\langle x,\gamma u\rangle)\,l(\phi')=\psi(\langle x_1,u\rangle)\,l(\phi')$.
Also
$\operatorname{co}(x_1)=-t(J_{V'})\,\pi_{t\circ c}(g)^{-1}\pi_{t\circ c}(x)\,(\pi_{t\circ c}(h)^{-1})^{\mathsf T}t(J_W)^{\mathsf T}$
has the same rank as $\pi_{t\circ c}(x)$, hence as $\operatorname{co}(x)=Y$, so
$\iota\phi_v(\operatorname{co}(x_1))=0$. By
Lemma~\ref{asm-theta-agree-coinvariants-eigen}(3), $l(\phi_v)=0$, that is,
$\iota(\operatorname{res}_z\Psi)(Y)=0$.
\end{proof}
% Lean: (1) KCo.surjectivity_statement, by induction over Finset.induction_on with step.
% (2) KCoNarrow.irrSt_narrow: extZ (extZ_mem, exists_extZ, extZ_act) and irr_equiv for the
%   equivalence omega_equiv_induced_fibre with mu = nu^{-1}; irr_ind_narrow (Lemma
%   asm-theta-agree-coinvariants-irreducible); the hypothesis USm alpha is not used.
% (3) KCo.dictSt_of ThKG.StdDict.dictX section_statement: only the second half of SecSt
%   (usm_chrW) is used, and the hypothesis that M is regular is not used.
% (4) KCo.so_statement = so_stable for G.M' with IsUWeilPair.isUSplitting of the Weil
%   representation, pure_at and pure_eq_pureA for r, slice_mem', mem_so, eigen_kill,
%   rank_coB and symp_act; det_ne_zero_of_hasSig for the invertibility.

\begin{lemma}\label{asm-theta-agree-coinvariants-split}
Let $E$ be a CM field, Galois over $\mathbf Q$, with maximal totally real
subfield $F$, and let $v$ be a finite place of $F$ such that the rational prime
below $v$ splits completely in $E$. Then $E$ has an $F$-embedding into $F_v$.
\end{lemma}
\leanref{Development/WeilBound/CoreStatements/LocallyConstantCoinvariants/ExistsFull.lean}{PositivityTheory.KCo.GoodSplit}
\leanref{Development/WeilBound/CoreStatements/LocallyConstantCoinvariants/PfSplit.lean}{PositivityTheory.KCo.good_split_of_cm}
\begin{proof}
Let $p$ be the rational prime below $v$. As $p$ splits completely in $E$, every
prime of $E$ above $p$ has ramification index and residue degree $1$ over $p$,
hence over $v$. As $E/F$ is Galois of degree $2$, the number of primes of $E$
above $v$ times their common ramification index and residue degree is $2$; so
there are two of them, and $E\otimes_FF_v\cong F_v\times F_v$. Either
projection restricted to $E$ is an $F$-embedding $E\to F_v$.
\end{proof}

\begin{lemma}\label{asm-theta-agree-coinvariants-nontrivial}
Let $(B,G,\mathcal X,\omega,\omega')$ be a triple
(Definition~\ref{asm-theta-agree-triple}), with $\nu=\alpha/\beta$, let $v$ be
a finite place clean for $\mathcal X$, and let $t\colon E\to F_v$ be an
$F$-embedding. Then $\nu_t\neq1$
(Definition~\ref{asm-theta-agree-coinvariants-local}(4)).
\end{lemma}
\leanref{Development/WeilBound/CoreStatements/LocallyConstantCoinvariants/RightOnto.lean}{PositivityTheory.KCo.NuSt}
\leanref{Development/WeilBound/CoreStatements/LocallyConstantCoinvariants/ContinuousFinId.lean}{PositivityTheory.KCo.nu_statement}
\leanref{Development/WeilBound/CoreStatements/LocallyConstantCoinvariants/ContinuousFinId.lean}{PositivityTheory.KCo.not_prod}
\begin{proof}
Suppose $\nu_t=1$, and let $w$ be the place of $E$ defined by $t$. As $v$ is
clean, $v\notin T$, so $\nu$ is unramified at $w$, and $v$ splits in $E$, so a
uniformizer $\varpi\in F$ at $v$ is one at $w$. Then $\nu_w(\varpi)=\nu_t(\varpi)=1$,
and for $x\in E^\times$, $x\varpi^{-\operatorname{ord}_w(x)}$ is a unit at $w$,
so $\nu_w(x)=1$.

The class group of $E$ is finite, so $w^k=(\pi_0)$ for some $k>0$ and some
$\pi_0\in\mathcal O_E$; then $w$ is the only prime containing $\pi_0$. For a
finite place $u\neq w$ of $E$, $\pi_0$ is a unit at $u$; as $\nu_u$ is
continuous it is trivial on $1+\varpi_u^m\mathcal O_u$ for some $m$, and as
$\mathcal O_E/u^m$ is finite, two powers of $\pi_0$ are congruent modulo $u^m$;
so $\nu_u(\pi_0)$ is a root of unity. Let $N$ be the product of the orders of
the $\nu_u(\pi_0)$ over the finitely many places $u$ above $T$, and $x=\pi_0^N$.
Then $\nu_u(x)=1$ for every finite place $u$: above $T$ by the choice of $N$,
at $w$ by the first paragraph, and elsewhere because $\nu$ is unramified there
and $x$ is a unit. As $\nu$ is trivial on $E^\times$ and its finite part is
continuous, $\nu$ is trivial on the idele with the finite components of $x$ and
archimedean components $1$; so $\nu$ is trivial on the idele with the
archimedean components of $x$ and finite components $1$. With
$m(w')=m_\alpha(w')-m_\beta(w')$ this reads
\[
\prod_{w'}\Bigl(\frac{\tau_{w'}(x)}{|\tau_{w'}(x)|}\Bigr)^{m(w')}=1,
\]
the product over the infinite places $w'$ of $F$. Here $m(w')>a+a'>0$, as $G$
is proper.

Squaring, and using $|z|^2=z\bar z$ and $\overline{\tau_{w'}(x)}=\tau_{w'}(c(x))$,
gives $\prod_{w'}\tau_{w'}(x)^{m(w')}=\prod_{w'}\tau_{w'}(c(x))^{m(w')}$. Fix an
infinite place $w_0$ of $F$. As $E$ is Galois over $\mathbf Q$, there are
$g_{w'}\in\operatorname{Gal}(E/\mathbf Q)$ with
$\tau_{w'}=\tau_{w_0}\circ g_{w'}$, and $g_{w_0}=1$. So
$\prod_{w'}g_{w'}(x)^{m(w')}=\prod_{w'}(g_{w'}c)(x)^{m(w')}$ in $\mathcal O_E$.
The left side lies in $w$, through the factor $x^{m(w_0)}$. So some
$(g_{w'}c)(x)$ lies in $w$, that is $x\in(g_{w'}c)^{-1}w$, and so
$(g_{w'}c)^{-1}w=w$, as $w$ is the only prime containing $x$. As the rational
prime below $v$ splits completely in $E$, the decomposition group of $w$ is
trivial, so $g_{w'}=c$. Then $\tau_{w_0}=\tau_{w'}\circ c$, and restricting to
$F$ gives $w'=w_0$, so $c=g_{w_0}=1$, which is absurd.
\end{proof}
% Lean: hall (the first paragraph), exists_pow_eq_one, arch_eq_one with hasExp_div and
% fin_eq_one, hT.geo.kappa for m > 0, and not_prod (the last paragraph).

\begin{lemma}\label{asm-theta-agree-coinvariants-kernel}
Let $E$ be a CM field with maximal totally real subfield $F$, let $X$ and $W$ be
lattice spaces of ranks $n\leq d$ (Definition~\ref{asm-theta-agree-doubling}(1)),
let $I=\{1,\ldots,n\}\times\{1,\ldots,d\}$, let $\omega$ be a representation of
$U_X\times U_W$ on the functions on $\mathbf A_F^I$, let $v$ be a finite place
of $F$, and let $\iota$ be a dictionary at $v$ with unitary smooth
characters $(\chi_1,\chi_2)$, $\chi_1\neq1$
(Definition~\ref{asm-theta-agree-coinvariants-local}). Let $\mathfrak a_v(\omega)$
be the span of the $\phi_1-\phi$ with $\phi,\phi_1$ Schwartz-Bruhat on $F_v^I$
and $\omega(g,1)$ acting at $v$ by $\phi\mapsto\phi_1$ for some
$g\in U_X^{(v)}$.
\begin{enumerate}
\item Every Schwartz-Bruhat $\phi$ on $F_v^I$ is $\phi_0+z$ with $\phi_0$
  Schwartz-Bruhat, $\iota\phi_0$ of full rank and $z\in\mathfrak a_v(\omega)$.
\item Let $P\neq0$ be a complex vector space and $p$ a linear map from functions
  on $F_v^I$ to $P$ that maps the Schwartz-Bruhat functions onto $P$, vanishes
  on $\mathfrak a_v(\omega)$, and has $p(\phi_1)=0$ whenever $p(\phi)=0$, the
  functions $\phi,\phi_1$ are Schwartz-Bruhat, and $\omega(1,h)$ acts at $v$ by
  $\phi\mapsto\phi_1$ for some $h\in U_W^{(v)}$. Then a Schwartz-Bruhat $\phi$
  has $p(\phi)=0$ if and only if $\phi\in\mathfrak a_v(\omega)$.
\end{enumerate}
\end{lemma}
\leanref{Development/WeilBound/CoreStatements/LocallyConstantCoinvariants/ExistsFull.lean}{PositivityTheory.KCo.exists_full}
\leanref{Development/AutomorphicSpectrum/Proofs/BlockDecomposition/KCoinvOfNarrow.lean}{GlobalTheory.CompanionTheory.KCoNarrow.ker_eq_ne}
\begin{proof}
First, $\mathfrak a(\chi_1,\chi_2)\subset\iota(\mathfrak a_v(\omega))$. Let
$\omega_{\chi_1,\chi_2}(g_0,1)\Psi-\Psi$ be a generator, and write
$\Psi=\iota\phi$ with $\phi$ Schwartz-Bruhat. By (a) of the definition of a
dictionary there are $g\in U_X^{(v)}$ and a Schwartz-Bruhat $\phi_1$ with
$\omega(g,1)$ acting at $v$ by $\phi\mapsto\phi_1$ and
$\iota\phi_1=\omega_{\chi_1,\chi_2}(g_0,1)\Psi$; so the generator is
$\iota(\phi_1-\phi)$ with $\phi_1-\phi\in\mathfrak a_v(\omega)$. As $\iota$ is
injective on Schwartz-Bruhat functions and $\mathfrak a_v(\omega)$ consists of
them, a Schwartz-Bruhat $\phi$ with $\iota\phi\in\mathfrak a(\chi_1,\chi_2)$
lies in $\mathfrak a_v(\omega)$.

(1) By Lemma~\ref{asm-theta-agree-coinvariants-inputs}(1),
$\iota\phi=\Psi_0+z'$ with $\Psi_0$ of full rank and
$z'\in\mathfrak a(\chi_1,\chi_2)$. Write $\Psi_0=\iota\phi_0$ with $\phi_0$
Schwartz-Bruhat. Then $\iota(\phi-\phi_0)=z'$, so $\phi-\phi_0\in\mathfrak a_v(\omega)$.

(2) The condition is sufficient by assumption. Let $N$ be the image under
$\iota$ of the Schwartz-Bruhat $\psi$ with $p(\psi)=0$. It contains
$\mathfrak a(\chi_1,\chi_2)$, by the first paragraph. It is stable under the
$\omega_{\chi_1,\chi_2}(1,h_0)$: for $\iota\psi\in N$, (b) of the definition of a
dictionary gives $h\in U_W^{(v)}$ and a Schwartz-Bruhat $\psi_1$ with
$\omega(1,h)$ acting at $v$ by $\psi\mapsto\psi_1$ and
$\iota\psi_1=\omega_{\chi_1,\chi_2}(1,h_0)\iota\psi$, and $p(\psi_1)=0$ by
assumption. By Lemma~\ref{asm-theta-agree-coinvariants-inputs}(2) there are two
cases.

If $N$ contains every function of full rank, take $x\neq0$ in $P$ and a
Schwartz-Bruhat $\psi$ with $p(\psi)=x$. By (1), $\psi=\psi_0+z$ with
$\iota\psi_0$ of full rank and $z\in\mathfrak a_v(\omega)$. Then
$\iota\psi_0\in N$, so by injectivity $p(\psi_0)=0$, and
$x=p(\psi_0)+p(z)=0$, which is absurd.

Otherwise every function of full rank in $N$ lies in $\mathfrak a(\chi_1,\chi_2)$.
Let $\phi$ be Schwartz-Bruhat with $p(\phi)=0$, and write $\phi=\phi_0+z$ by
(1). Then $p(\phi_0)=p(\phi)-p(z)=0$, so $\iota\phi_0\in N$ has full rank and
lies in $\mathfrak a(\chi_1,\chi_2)$. By the first paragraph
$\phi_0\in\mathfrak a_v(\omega)$, and so $\phi\in\mathfrak a_v(\omega)$.
\end{proof}
% Lean: the first paragraph is aug_le_map and mem_augAt (ExistsFull.lean), with augAt_le.

\begin{lemma}\label{asm-theta-agree-coinvariants}
Let $(B,G,\mathcal X,\omega,\omega',R,R_f)$ be a matched quadruple for a theta
base $B$ of a CM field $E$ (Definition~\ref{asm-theta-agree-matched}). Then it
describes the coinvariants (Definition~\ref{asm-theta-agree-parts}(3)).
\end{lemma}
\leanref{Development/AbelianVarietyTheory/Definitions/Divisors/MatchingOfParts.lean}{GlobalTheory.ThK.KCoinv}
\leanref{Development/WeilBound/CoreStatements/TemperedRepresentationTheory/BernsteinNarrowness/BernsteinClosure/RamifiedMultiplicityOne.lean}{PositivityTheory.BernsteinRamification.Close.kCoinv}
\leanref{Development/AutomorphicSpectrum/Proofs/BlockDecomposition/KCoinvOfNarrow.lean}{GlobalTheory.CompanionTheory.KCoNarrow.kCoinv_of_ne}
\begin{proof}
By Lemma~\ref{asm-theta-agree-coinvariants-inputs}(4), $\omega'(g,h)$ maps
$\mathcal S^\circ(M')$ into itself. Let $v$ be clean for $\mathcal X$. As $G$
is proper, $1\leq a'<a$, so both Weil pairs of the triple,
$(V,W,M,\beta,\beta/\alpha,\omega)$ and $(V',W,M',\alpha,\alpha/\beta,\omega')$,
have ranks $n\leq d$ with $n\geq1$. By
Lemma~\ref{asm-theta-agree-coinvariants-split} there is an $F$-embedding
$t_0\colon E\to F_v$.

Let $t$ be an $F$-embedding $E\to F_v$. Apply
Lemma~\ref{asm-theta-agree-coinvariants-inputs}(3) at $t$ to the Weil pair of
$\omega'$, that is, to $V'$, $W$, $M'$, the characters $(\alpha,\nu)$ with
$\nu=\alpha/\beta$, and $\omega'$: there are a
linear model $\iota$ of the second kind for $(t,1,1)$ and unitary smooth
$\chi_1,\chi_2$ such that $\iota$ is a dictionary with characters
$(\chi_1,\chi_2)$. If $\chi_1=1$, then $\nu_t=1$ or $\nu_{t\circ c}=1$,
contradicting Lemma~\ref{asm-theta-agree-coinvariants-nontrivial}; so
$\chi_1\neq1$. The same holds for the Weil pair of $\omega$, whose second
character is $\nu^{-1}$, since $(\nu^{-1})_t=1$ only if $\nu_t=1$.

\emph{The kernels.} Apply Lemma~\ref{asm-theta-agree-coinvariants-kernel}(2) at
$t_0$ to $p=p_v$ and $P=\Pi_v$. As the finite maps are matched, $\Pi_v\neq0$,
$p_v$ maps the Schwartz-Bruhat functions onto $\Pi_v$ and vanishes on
$\mathfrak a_v(\omega)$. If $p_v(\phi)=0$ and $\omega(1,h)$ acts at $v$ by
$\phi\mapsto\phi_1$ with $h\in U_W^{(v)}$, then $\omega'(1,h)$ acts at $v$ by
$0\mapsto0$ and $p_v(\phi)=p'_v(0)$, so the matching gives
$p_v(\phi_1)=p'_v(0)=0$. So $p_v(\phi)=0$ if and only if
$\phi\in\mathfrak a_v(\omega)$. The same argument applies to $p'_v$ and
$\mathfrak a_v(\omega')$.

\emph{Full rank.} Let $t$ be an $F$-embedding, $\xi\in\Pi_v$, and $\psi$
Schwartz-Bruhat with $p'_v(\psi)=\xi$. With $\iota$ for $\omega'$ at $t$ as above,
Lemma~\ref{asm-theta-agree-coinvariants-kernel}(1) gives $\psi=\psi_0+z$ with
$\iota\psi_0$ of full rank and $z\in\mathfrak a_v(\omega')$. Then $\psi_0$ has
full rank at $t$, as $\iota$ is a linear model of the second kind for
$(t,1,1)$, and $p'_v(\psi_0)=p'_v(\psi)=\xi$, as $p'_v$ vanishes on
$\mathfrak a_v(\omega')$.

\emph{Linear forms.} Let $\ell$ vanish on $\mathfrak a_v(\omega')$. As $p'_v$
maps the Schwartz-Bruhat functions onto $\Pi_v$, it has a linear section $s$
with values in them. Put $\ell_0=\ell\circ s$. For Schwartz-Bruhat $\phi$,
$p'_v(\phi-s(p'_v(\phi)))=0$, so $\phi-s(p'_v(\phi))\in\mathfrak a_v(\omega')$ by
the kernels, and $\ell(\phi)=\ell_0(p'_v(\phi))$.
\end{proof}
% Lean: Close.kCoinv hD = kCoinv_of_geo geometricLemmaFiltration hD = kCoinv_of_narrow
% (narrowBern ...) hD = kCoinv_of_ne surjectivity_statement (irrSt_narrow NB) hD
% nu_statement so_statement good_split_of_cm, with hD = KCo.dictSt_of ThKG.StdDict.dictX
% KCo.section_statement. ker_eq_ne is applied with st from hF.equiv and acts_zero; the
% contradiction for chi_1 = 1 uses the clause h1 of DictSt, NuSt and chrW_div_ne.

\subsection{Convergence and continuity of theta integrals}\label{subsection-asm-convergence-continuity}

\begin{definition}\label{asm-theta-agree-convergence-tori}
Let $E$ be a CM field with maximal totally real subfield $F$ and complex
conjugation $c$, write $\mathbf A=\mathbf A_F$ and
$\mathbf A_E=\mathbf A\otimes_FE$, and let $n,d,q,q'\geq0$. For $t\in\mathbf R$
let $u_t\in\mathbf A^\times$ have component $e^t$ at every infinite place and
$1$ at every finite place.
\begin{enumerate}
\item A matrix $J\in M_n(E)$ is \emph{nondegenerate $\epsilon$-Hermitian} if
  $\det J\neq0$ and $c(J)^{\mathsf T}=\epsilon J$ with $\epsilon=1$ or
  $\epsilon=-1$. Then $U_J$ is the group of $g\in GL_n(\mathbf A_E)$ with
  $g^{\mathsf T}J\,c(g)=J$, $U_J(F)$ is its subgroup of matrices over $E$, and
  $h(x,y)=x^{\mathsf T}J\,c(y)$ for $x,y\in E^n$.
\item A \emph{Witt torus} of size $q$ of such a $J$ consists of vectors
  $x_1,\ldots,x_q,y_1,\ldots,y_q\in E^n$ and a continuous homomorphism
  $a\colon\mathbf R^q\to U_J$ such that: $h(x_i,x_j)=h(y_i,y_j)=0$ and
  $h(x_i,y_j)=\delta_{ij}$; every $v\in E^n$ with $h(v,x_i)=h(v,y_i)=0$ for all
  $i$ and $h(v,v)=0$ is $0$; and for every $t\in\mathbf R^q$,
  $a(t)x_i=u_{t_i}x_i$, $a(t)y_i=u_{-t_i}y_i$, and $a(t)v=v$ for every
  $v\in E^n$ with $h(v,x_i)=h(v,y_i)=0$ for all $i$.
\item Put
  $\mathcal C_q=\{t\in\mathbf R^q:t_1\geq t_2\geq\cdots\geq t_q\geq0\}$ and
  $\delta_n(t)=\sum_{i=1}^q(n-2i+1)t_i$. A Witt torus $a$ of size $q$ is
  \emph{Siegel} if:
  \begin{enumerate}
  \item[(a)] there is a compact $\Omega\subset U_J$ such that every $g\in U_J$
    satisfies $\gamma g=a(t)\kappa$ for some $\gamma\in U_J(F)$,
    $t\in\mathcal C_q$ and $\kappa\in\Omega$;
  \item[(b)] for every Haar measure $\mu$ on $U_J$ and every fundamental domain
    $\mathcal D$ for $U_J(F)$ in $U_J$ with respect to $\mu$ there are
    $K<\infty$ and a compact $\Omega\subset U_J$ such that
    \[
    \int_{\mathcal D}f\,d\mu\leq K\int_{\mathcal C_q}\beta(t)\,e^{-[E:\mathbf Q]\delta_n(t)}\,dt
    \]
    for all measurable $\beta\colon\mathbf R^q\to[0,\infty]$ and all measurable
    $f\colon U_J\to[0,\infty]$ with $f(\gamma g)=f(g)$ for $\gamma\in U_J(F)$
    and $f(a(t)\kappa)\leq\beta(t)$ for $t\in\mathcal C_q$ and
    $\kappa\in\Omega$;
  \item[(c)] there are $C$ and $k$ with
    $\operatorname{ht}(a(t))\leq Ce^{k(t_1+\cdots+t_q)}$ and
    $e^{t_1+\cdots+t_q}\leq C\operatorname{ht}(\gamma a(t))^k$ for all
    $t\in\mathcal C_q$ and $\gamma\in U_J(F)$, with the height of
    Definition~\ref{asm-theta-agree-doubling}(2).
  \end{enumerate}
\item Suppose that $E$ is Galois over $\mathbf Q$, and let $B$ be a theta base
  for $E$ (Definition~\ref{asm-theta-agree-models}(1)), with character $\psi$
  and maps $\pi_w$. For $z\in M_{n\times d}(\mathbf A_E)$ put
  $\|z\|_\infty=\sum_w\sum_{i,j}|\pi_w(z_{ij})|$, the sum over the infinite
  places $w$ of $F$. For a compact set $C\subset M_{n\times d}(\mathbf A_E)$
  and an integer $N\geq0$, the \emph{majorant} $m_{C,N}(z)$ is
  $(1+\|z\|_\infty)^{-N}$ if $z-z_0$ is at $\infty$
  (Definition~\ref{asm-theta-agree-models}(2)) for some $z_0\in C$, and $0$
  otherwise. For nondegenerate $\epsilon$-Hermitian $J\in M_n(E)$ and
  $J'\in M_d(E)$, the \emph{lattice sum} is
  \[
  L_{C,N}(g,h)=\sum_{x\in M_{n\times d}(E)}m_{C,N}\bigl(g^{-1}x\,(h^{-1})^{\mathsf T}\bigr)\in[0,\infty]
  \qquad(g\in U_J,\ h\in U_{J'}).
  \]
\item For $s\in\mathbf R^q$ and $t\in\mathbf R^{q'}$ put
  \[
  P_{n,d}(s,t)=\sum_{i=1}^q\sum_{j=1}^{q'}2\max(s_i,t_j)+(d-2q')\sum_{i=1}^qs_i+(n-2q)\sum_{j=1}^{q'}t_j .
  \]
\item With $B$ as in (4), let $I$ be a finite set and $dx$ a Haar measure on
  $\mathbf A^I$. For $w=(w_1,w_2)\in\mathbf A^I\times\mathbf A^I$ and a
  function $\Phi$ on $\mathbf A^I$ put
  $(\rho(w)\Phi)(x)=\psi\bigl(x\cdot w_2+\tfrac12w_1\cdot w_2\bigr)\Phi(x+w_1)$.
  The \emph{coefficient} of Schwartz-Bruhat functions $\Phi_1$ and $\Phi_2$ at
  $w$ is
  $\langle\rho(w)\Phi_1,\Phi_2\rangle=\int_{\mathbf A^I}(\rho(w)\Phi_1)(x)\overline{\Phi_2(x)}\,dx$.
  For a model $M$ for lattice spaces $(X,W)$, with isomorphism $e$
  (Definition~\ref{asm-theta-agree-models}(3)), the Schrödinger operator of
  $z\in\mathbb W(X,W)$ is $\rho(z)=\rho(e(z))$.
\end{enumerate}
\end{definition}
\leanref{Development/WeilBound/ThetaLifts/FrameEstimates/SquareIntegrableWeilRange.lean}{PositivityTheory.ThKEst.EpsHerm}
\leanref{Development/WeilBound/ThetaLifts/FrameEstimates/SquareIntegrableWeilRange.lean}{PositivityTheory.ThKEst.rId}
\leanref{Development/WeilBound/ThetaLifts/FrameEstimates/SquareIntegrableWeilRange.lean}{PositivityTheory.ThKEst.Frame}
\leanref{Development/WeilBound/ThetaLifts/FrameEstimates/SquareIntegrableWeilRange.lean}{PositivityTheory.ThKEst.Frame.IsGood}
\leanref{Development/WeilBound/ThetaLifts/FrameEstimates/SquareIntegrableWeilRange.lean}{PositivityTheory.ThKEst.Cone}
\leanref{Development/WeilBound/ThetaLifts/FrameEstimates/SquareIntegrableWeilRange.lean}{PositivityTheory.ThKEst.rho2}
\leanref{Development/WeilBound/ThetaLifts/FrameEstimates/SquareIntegrableWeilRange.lean}{PositivityTheory.ThKEst.SiegelCover}
\leanref{Development/WeilBound/ThetaLifts/FrameEstimates/SquareIntegrableWeilRange.lean}{PositivityTheory.ThKEst.SiegelInt}
\leanref{Development/WeilBound/ThetaLifts/FrameEstimates/SquareIntegrableWeilRange.lean}{PositivityTheory.ThKEst.TorHt}
\leanref{Development/WeilBound/ThetaLifts/FrameEstimates/SquareIntegrableWeilRange.lean}{PositivityTheory.ThKEst.sz}
\leanref{Development/WeilBound/ThetaLifts/FrameEstimates/SquareIntegrableWeilRange.lean}{PositivityTheory.ThKEst.maj}
\leanref{Development/WeilBound/ThetaLifts/FrameEstimates/SquareIntegrableWeilRange.lean}{PositivityTheory.ThKEst.latSum}
\leanref{Development/WeilBound/ThetaLifts/FrameEstimates/SquareIntegrableWeilRange.lean}{PositivityTheory.ThKEst.posWt}
\leanref{Development/WeilBound/ThetaLifts/FrameEstimates/SquareIntegrableWeilRange.lean}{PositivityTheory.ThKEst.coef}

\begin{theorem}[Reduction theory \cite{Bor69,BHC62}]\label{asm-theta-agree-convergence-siegel}
Let $E$ be a CM field with maximal totally real subfield $F$, and let
$J\in M_n(E)$ be nondegenerate $\epsilon$-Hermitian
(Definition~\ref{asm-theta-agree-convergence-tori}(1)). Then:
\begin{enumerate}
\item $U_J$ is a locally compact, second countable Hausdorff group, the
  inclusion of $U_J$ in $M_n(\mathbf A_E)$ is continuous, $U_J(F)$ is a
  countable discrete subgroup, and the height
  (Definition~\ref{asm-theta-agree-doubling}(2)) is bounded on every compact
  subset of $U_J$;
\item $J$ has a Siegel Witt torus of some size $q$ with $2q\leq n$
  (Definition~\ref{asm-theta-agree-convergence-tori}(3)).
\end{enumerate}
\end{theorem}
\leanref{Development/WeilBound/ThetaLifts/FrameEstimates/SquareIntegrableWeilRange.lean}{PositivityTheory.ThKEst.GrpTop}
\leanref{Development/WeilBound/ThetaLifts/FrameEstimates/GrpTop.lean}{PositivityTheory.ThKEst.grpTop_of}
\leanref{Development/WeilBound/ThetaLifts/FrameEstimates/SquareIntegrableWeilRange.lean}{PositivityTheory.ThKEst.SiegelAll}
\leanref{Development/WeilBound/ThetaLifts/FrameEstimates/SiegelAll.lean}{PositivityTheory.ThKEst.siegelAll_of}
\leanref{Development/AbelianVarietyTheory/PublishedTheorems/Divisors₂/IsSmoothModerateOfIsAutFunSm.lean}{GlobalTheory.Published.SiegelSetUnitary}

\begin{lemma}\label{asm-theta-agree-convergence-signature}
Let $B$ be a theta base for a CM field $E$
(Definition~\ref{asm-theta-agree-models}(1)), and let $X$ be a lattice space
of rank $n$ with matrix $J_X$ that is Hermitian or skew-Hermitian of signature
$(r,s)$ (Definition~\ref{asm-theta-agree-doubling}(1)). Then $J_X$ is
nondegenerate $\epsilon$-Hermitian, $r+s=n$, and every Witt torus of $J_X$
(Definition~\ref{asm-theta-agree-convergence-tori}(2)) has size at most
$\min(r,s)$.
\end{lemma}
\leanref{Development/WeilBound/ThetaLifts/FrameEstimates/SquareIntegrableWeilRange.lean}{PositivityTheory.ThKEst.EpsOf}
\leanref{Development/WeilBound/ThetaLifts/FrameEstimates/BlockForm.lean}{PositivityTheory.ThKEst.epsOf_of}
\begin{proof}
The relation $c(J_X)^{\mathsf T}=\pm J_X$ is part of the definition. Fix an
infinite place $w$ of $F$. The matrix $M_w$ is $\tau_w(J_X)$ in the Hermitian
case and $i\,\tau_w(J_X)$ in the skew-Hermitian case. It is a Hermitian matrix
of signature $(r,s)$, so it is congruent to
$\operatorname{diag}(1_r,-1_s)$; hence $r+s=n$ and $\det J_X\neq0$. Let
$x_i,y_i$ be the vectors of a Witt torus of size $q$. As $\tau_w\circ c$ is
complex conjugation on $\tau_w(E)$,
$\tau_w(x_i)^{\mathsf T}M_w\overline{\tau_w(x_j)}=0$ and
$\tau_w(x_i)^{\mathsf T}M_w\overline{\tau_w(y_j)}=\lambda\delta_{ij}$, with
$\lambda=1$ or $\lambda=i$. So the $\tau_w(x_i)$ are linearly independent and
span a totally isotropic subspace of dimension $q$ for the Hermitian form of
$M_w$, and such a subspace has dimension at most $\min(r,s)$.
\end{proof}

\begin{lemma}\label{asm-theta-agree-convergence-theta}
Let $(X,W,M,\chi_X,\chi_W,\omega)$ be a Weil pair for a theta base $B$ of a CM
field $E$ (Definition~\ref{asm-theta-agree-triple}(1)), with $X$ and $W$ of
ranks $n$ and $d$, and put $I=\{1,\ldots,n\}\times\{1,\ldots,d\}$.
\begin{enumerate}
\item The theta sum $\theta(\Phi)=\sum_{\xi\in F^I}\Phi(\xi)$ of a
  Schwartz-Bruhat function $\Phi$ on $\mathbf A^I$ converges absolutely, and
  there is an admissible linear form $l$
  (Definition~\ref{asm-theta-agree-doubling}(6)) with $l(\Phi)=\theta(\Phi)$
  for every Schwartz-Bruhat $\Phi$.
\item For Schwartz-Bruhat functions $\Phi$ and $\Psi$ on $\mathbf A^I$ and
  $a\in\mathbf C$,
  $\theta_\omega(\Phi+a\Psi)=\theta_\omega(\Phi)+a\,\theta_\omega(\Psi)$, and
  $\theta_\omega(\Phi)$ is continuous on $U_X\times U_W$.
\item For a Schwartz-Bruhat function $\phi$ on $\mathbf A_f^I$ and
  $(g,h)\in U_X\times U_W$, the map $f\mapsto\theta_\omega(f\otimes\phi)(g,h)$
  is continuous on $\mathcal S(F_\infty^I)$.
\end{enumerate}
\end{lemma}
\leanref{Development/AutomorphicSpectrum/Proofs/MembershipProofs/ThetaSumAdmissibilityStatement.lean}{GlobalTheory.T.theta_distribution_nonzero_admissible}
\leanref{Development/AbelianVarietyTheory/Definitions/Divisors/CTSmallCone.lean}{GlobalTheory.ThK.ThetaBasic}
\leanref{Development/WeilBound/ThetaLifts/LeviDeterminants/ThetaBasicOf.lean}{PositivityTheory.ThKF0.thetaBasic_of}
\leanref{Development/WeilBound/ThetaLifts/FrameEstimates/ThickFrames.lean}{PositivityTheory.ThKEst.theta_add}
\begin{proof}
Part (1) is the continuity of Weil's theta distribution \cite{Wei64}, in the
form that asks for continuity in the archimedean component for each fixed
finite component. The operators $\omega(g,h)=\omega(g,1)\,\omega(1,h)$ preserve
the Schwartz-Bruhat functions and the admissible linear forms, by conditions (a)
and (b) of Definition~\ref{asm-theta-agree-models}(6). So
$\theta_\omega(\Phi)(g,h)=\theta(\omega(g,h)\Phi)=l(\omega(g,h)\Phi)$. This is
linear in $\Phi$, and continuous in $(g,h)$ by condition (c) there. For (3),
the linear form $l\circ\omega(g,h)$ is admissible, so
$f\mapsto l(\omega(g,h)(f\otimes\phi))$ is continuous.
\end{proof}

\begin{lemma}\label{asm-theta-agree-convergence-poisson}
Let $B$ be a theta base for a CM field $E$
(Definition~\ref{asm-theta-agree-models}(1)), with character $\psi$, let $I$ be
a finite set, and let $dx$ be a Haar measure on $\mathbf A^I$. Then there is
$c>0$ such that for all Schwartz-Bruhat functions $\Phi_1$ and $\Phi_2$ on
$\mathbf A^I$,
\[
|\theta(\Phi_1)|\,|\theta(\Phi_2)|\leq c\sum_{w\in F^I\times F^I}\bigl|\langle\rho(w)\Phi_1,\Phi_2\rangle\bigr|,
\]
with the theta sums of Lemma~\ref{asm-theta-agree-convergence-theta}(1) and the
coefficients of Definition~\ref{asm-theta-agree-convergence-tori}(6).
\end{lemma}
\leanref{Development/WeilBound/ThetaLifts/FrameEstimates/SquareIntegrableWeilRange.lean}{PositivityTheory.ThKEst.CoefSum}
\leanref{Development/WeilBound/ThetaLifts/FrameEstimates/CoefSumB.lean}{PositivityTheory.ThKEst.coefSum_of_annih}
\leanref{Development/WeilBound/ThetaLifts/FrameEstimates/CoefSumA.lean}{PositivityTheory.ThKEst.CS.Dual.poisson_le}
\begin{proof}
The theta sums converge absolutely, so
\[
\theta(\Phi_1)\overline{\theta(\Phi_2)}=\sum_{x\in F^I}\sum_{\xi\in F^I}\Phi_1(\xi+x)\overline{\Phi_2(\xi)},
\]
the double sum converging absolutely. Fix $x\in F^I$ and put
$H(u)=\Phi_1(u+x)\overline{\Phi_2(u)}$, a continuous function on
$\mathbf A^I$ whose translates by a compact set are bounded by a summable
function on $F^I$. The group $F^I$ is discrete in $\mathbf A^I$ with compact
quotient. The characters $u\mapsto\psi(u\cdot y)$, $y\in F^I$, are trivial on
$F^I$, as $\psi$ is trivial on $F$, and they separate the points of
$\mathbf A^I/F^I$, as an $a\in\mathbf A$ with $\psi(at)=1$ for every $t\in F$
lies in $F$. By the Poisson summation formula for $F^I$ in $\mathbf A^I$
\cite[\S4.2]{Tat67}, with $c$ the inverse of the volume of a fundamental domain
for $F^I$,
\[
\Bigl|\sum_{\xi\in F^I}H(\xi)\Bigr|\leq c\sum_{y\in F^I}\Bigl|\int_{\mathbf A^I}\psi(u\cdot y)H(u)\,du\Bigr| .
\]
As $\tfrac12x\cdot y\in F$, $\psi(\tfrac12x\cdot y)=1$, and the integral is
$\langle\rho(x,y)\Phi_1,\Phi_2\rangle$. Summing over $x$ gives the lemma.
\end{proof}

\begin{lemma}\label{asm-theta-agree-convergence-coefficients}
Let $(X,W,M,\chi_X,\chi_W,\omega)$ be a Weil pair for a theta base $B$ of a CM
field $E$ (Definition~\ref{asm-theta-agree-triple}(1)), with $X$ and $W$ of
ranks $n$ and $d$, put $I=\{1,\ldots,n\}\times\{1,\ldots,d\}$, and let $dx$ be
a Haar measure on $\mathbf A^I$.
\begin{enumerate}
\item For Schwartz-Bruhat functions $\Phi_1$ and $\Phi_2$ on $\mathbf A^I$,
  $g\in U_X$, $h\in U_W$ and $z\in\mathbb W(X,W)$,
  \[
  \langle\rho(z)\,\omega(g,h)\Phi_1,\omega(g,h)\Phi_2\rangle
  =\bigl\langle\rho\bigl(g^{-1}z\,(h^{-1})^{\mathsf T}\bigr)\Phi_1,\Phi_2\bigr\rangle .
  \]
\item For Schwartz-Bruhat functions $\phi_1$ and $\phi_2$ on $\mathbf A_f^I$
  and $N\geq0$ there are a continuous seminorm $p$ on $\mathcal S(F_\infty^I)$
  and a compact set $C\subset\mathbb W(X,W)$ with
  \[
  \bigl|\langle\rho(z)(f_1\otimes\phi_1),f_2\otimes\phi_2\rangle\bigr|\leq p(f_1)\,p(f_2)\,m_{C,N}(z)
  \]
  for all $f_1,f_2\in\mathcal S(F_\infty^I)$ and $z\in\mathbb W(X,W)$, with the
  majorant of Definition~\ref{asm-theta-agree-convergence-tori}(4).
\end{enumerate}
\end{lemma}
\leanref{Development/WeilBound/ThetaLifts/FrameEstimates/SquareIntegrableWeilRange.lean}{PositivityTheory.ThKEst.CoefEq}
\leanref{Development/WeilBound/ThetaLifts/FrameEstimates/CoefEq.lean}{PositivityTheory.ThKEst.coefEq_of}
\leanref{Development/WeilBound/ThetaLifts/FrameEstimates/SquareIntegrableWeilRange.lean}{PositivityTheory.ThKEst.CoefBd}
\leanref{Development/WeilBound/ThetaLifts/FrameEstimates/CoefBd.lean}{PositivityTheory.ThKEst.coefBd_of}
\leanref{Development/WeilBound/ThetaLifts/FrameEstimates/CoefBd.lean}{PositivityTheory.ThKEst.W.coef_split}
\leanref{Development/WeilBound/ThetaLifts/FrameEstimates/FinBound.lean}{PositivityTheory.ThKEst.W.fin_bound}
\leanref{Development/WeilBound/ThetaLifts/FrameEstimates/SquareIntegrableWeilRange.lean}{PositivityTheory.ThKEst.CoefReal}
\leanref{Development/WeilBound/ThetaLifts/FrameEstimates/CoefReal.lean}{PositivityTheory.ThKEst.coefReal_of}
\begin{proof}
(1) By conditions (a) and (b) of Definition~\ref{asm-theta-agree-models}(6),
$\omega(1,h)\rho(w)=\rho(wh^{\mathsf T})\omega(1,h)$ and
$\omega(g,1)\rho(w)=\rho(gw)\omega(g,1)$ on Schwartz-Bruhat functions, so
$\omega(g,h)\rho(w)=\rho(gwh^{\mathsf T})\omega(g,h)$. The operator
$\omega(g,h)$ preserves the $L^2$-norm of Schwartz-Bruhat functions, hence, by
polarization, the inner product
$\langle\Phi,\Psi\rangle=\int\Phi\overline\Psi\,dx$. With
$w=g^{-1}z(h^{-1})^{\mathsf T}$, and as $\rho(w)\Phi_1$ is Schwartz-Bruhat,
\[
\langle\rho(z)\omega(g,h)\Phi_1,\omega(g,h)\Phi_2\rangle
=\langle\omega(g,h)\rho(w)\Phi_1,\omega(g,h)\Phi_2\rangle
=\langle\rho(w)\Phi_1,\Phi_2\rangle .
\]

(2) Write a point of $\mathbf A^I$ as $(x_\infty,x_f)$ with
$x_\infty\in F_\infty^I$ and $x_f\in\mathbf A_f^I$. The measure $dx$ is a
multiple of a product of Haar measures on $F_\infty^I$ and $\mathbf A_f^I$, and
$\psi$ is the product of its archimedean part $\psi_\infty$ and its finite part
$\psi_f$. So for $w=(w_1,w_2)$ with $w_1=(x_\infty,x_f)$ and
$w_2=(y_\infty,y_f)$,
\[
\langle\rho(w)(f_1\otimes\phi_1),f_2\otimes\phi_2\rangle
=c_0\,\psi\bigl(\tfrac12w_1\cdot w_2\bigr)\,A(x_\infty,y_\infty)\,B(x_f,y_f)
\]
for a constant $c_0>0$, where
$A(x,y)=\int\psi_\infty(u\cdot y)f_1(u+x)\overline{f_2(u)}\,du$ and
$B(x,y)=\int\psi_f(u\cdot y)\phi_1(u+x)\overline{\phi_2(u)}\,du$.
\begin{enumerate}
\item[(i)] $B$ is bounded, and vanishes outside a compact set $K$ of pairs
  $(x_f,y_f)$. Indeed $B(x,y)\neq0$ forces $x$ into the compact set
  $\operatorname{supp}\phi_1-\operatorname{supp}\phi_2$; and as $\phi_1$ and
  $\phi_2$ are invariant under translation by an open subgroup $U$, the
  substitution $u\mapsto u+k$ shows that it forces $\psi_f(k\cdot y)=1$ for
  $k\in U$, which confines $y$ to a compact set.
\item[(ii)] There is a continuous seminorm $p_0$ on $\mathcal S(F_\infty^I)$
  with $|A(x,y)|(1+|x|+|y|)^N\leq p_0(f_1)p_0(f_2)$ for all $x$ and $y$: the
  Fourier-Wigner transform of two Schwartz functions is a Schwartz function,
  depending continuously on them \cite[Chapter~1]{Fol89}.
\item[(iii)] As $e^{-1}$ is $\mathbf A$-linear, the archimedean components of
  $e^{-1}(w)$ depend linearly on $(x_\infty,y_\infty)$, so
  $1+\|e^{-1}(w)\|_\infty\leq L(1+|x_\infty|+|y_\infty|)$ for a constant
  $L\geq1$.
\end{enumerate}
Let $C$ be the image under $e^{-1}$ of the pairs $((0,x_f),(0,y_f))$ with
$(x_f,y_f)\in K$, a compact set. If the coefficient at $z=e^{-1}(w)$ is
nonzero, then $(x_f,y_f)\in K$, so $z-z_0$ is at $\infty$ for the element
$z_0\in C$ with the finite components of $w$; and by (i) to (iii) the
coefficient is at most
$c_1L^Np_0(f_1)p_0(f_2)(1+\|z\|_\infty)^{-N}$, with $c_1=c_0\sup|B|$. So
$p=(1+c_1L^N)p_0$ works.
\end{proof}

\begin{lemma}\label{asm-theta-agree-convergence-square}
Let $(X,W,M,\chi_X,\chi_W,\omega)$ be a Weil pair for a theta base $B$ of a CM
field $E$ (Definition~\ref{asm-theta-agree-triple}(1)), with $X$ and $W$ of
ranks $n$ and $d$ and matrices $J_X$ and $J_W$, and put
$I=\{1,\ldots,n\}\times\{1,\ldots,d\}$. For every Schwartz-Bruhat function
$\phi$ on $\mathbf A_f^I$ and every $N\geq0$ there are a continuous seminorm $p$
on $\mathcal S(F_\infty^I)$ and a compact set
$C\subset M_{n\times d}(\mathbf A_E)$ such that
\[
|\theta_\omega(f\otimes\phi)(g,h)|^2\leq p(f)^2\,L_{C,N}(g,h)
\]
for all $f\in\mathcal S(F_\infty^I)$, $g\in U_X$ and $h\in U_W$, with the
lattice sum for $J_X$ and $J_W$
(Definition~\ref{asm-theta-agree-convergence-tori}(4)).
\end{lemma}
\leanref{Development/WeilBound/ThetaLifts/FrameEstimates/SquareIntegrableWeilRange.lean}{PositivityTheory.ThKEst.SqAll}
\leanref{Development/WeilBound/ThetaLifts/FrameEstimates/SqAllC.lean}{PositivityTheory.ThKEst.sqAll_of'}
\leanref{Development/WeilBound/ThetaLifts/FrameEstimates/SqAllC.lean}{PositivityTheory.ThKEst.theta_mul_le}
\leanref{Development/WeilBound/ThetaLifts/FrameEstimates/SqAllC.lean}{PositivityTheory.ThKEst.exists_ratEquiv}
\begin{proof}
Fix a Haar measure $dx$ on $\mathbf A^I$, and let $e$ be the isomorphism of
$M$. The map $x\mapsto e(x)$ is a bijection from $M_{n\times d}(E)$ onto
$F^I\times F^I$: it maps into $F^I\times F^I$ because $M$ is regular, it is
$F$-linear and injective there, and both sides have dimension $2nd$ over $F$,
as $[E:F]=2$.

Let $\Phi=f\otimes\phi$ and $\Psi=\omega(g,h)\Phi$, a Schwartz-Bruhat
function, so that $\theta_\omega(\Phi)(g,h)=\theta(\Psi)$. By
Lemma~\ref{asm-theta-agree-convergence-poisson} for $\Psi$ and $\Psi$, and then
Lemma~\ref{asm-theta-agree-convergence-coefficients}(1),
\[
|\theta_\omega(\Phi)(g,h)|^2\leq c\sum_{x\in M_{n\times d}(E)}|\langle\rho(x)\Psi,\Psi\rangle|
=c\sum_{x\in M_{n\times d}(E)}\bigl|\bigl\langle\rho\bigl(g^{-1}x(h^{-1})^{\mathsf T}\bigr)\Phi,\Phi\bigr\rangle\bigr| .
\]
By Lemma~\ref{asm-theta-agree-convergence-coefficients}(2) with
$\phi_1=\phi_2=\phi$ there are $p_0$ and $C$ with each term at most
$p_0(f)^2m_{C,N}(g^{-1}x(h^{-1})^{\mathsf T})$. Summing over $x$ gives the
lemma with $p=\sqrt c\,p_0$.
\end{proof}

\begin{lemma}\label{asm-theta-agree-convergence-lattice}
Let $B$ be a theta base for a CM field $E$
(Definition~\ref{asm-theta-agree-models}(1)), let $J\in M_n(E)$ and
$J'\in M_d(E)$ be nondegenerate $\epsilon$-Hermitian, and let $L_{C,N}$ be the
lattice sums of Definition~\ref{asm-theta-agree-convergence-tori}(4). Let
$C\subset M_{n\times d}(\mathbf A_E)$ be compact and $N\geq0$.
\begin{enumerate}
\item $L_{C,N}(\gamma g,\gamma'h)=L_{C,N}(g,h)$ for $\gamma\in U_J(F)$,
  $\gamma'\in U_{J'}(F)$, $g\in U_J$ and $h\in U_{J'}$.
\item For every $h\in U_{J'}$ the function $g\mapsto L_{C,N}(g,h)$ is
  measurable.
\item For compact sets $\Omega\subset U_J$ and $\Omega'\subset U_{J'}$ there
  are a compact set $C'$ and $K\geq0$ with
  $L_{C,N}(g\kappa,h\kappa')\leq K\,L_{C',N}(g,h)$ for all $g$, $h$,
  $\kappa\in\Omega$ and $\kappa'\in\Omega'$.
\item For a Witt torus $a'$ of $J'$ of size $q'$ there are a compact set $C'$
  and $K\geq0$ with
  $L_{C,N}(g,a'(t))\leq K\,e^{N(t_1+\cdots+t_{q'})}L_{C',N}(g,1)$ for all
  $g\in U_J$ and $t\in\mathcal C_{q'}$.
\item For Witt tori $a$ of $J$ of size $q$ and $a'$ of $J'$ of size $q'$
  there is $N_0$ such that for every compact $C$ there is $K\geq0$ with
  $L_{C,N_0}(a(s),a'(t))\leq K\exp\bigl([E:\mathbf Q]\,P_{n,d}(s,t)\bigr)$ for
  all $s\in\mathcal C_q$ and $t\in\mathcal C_{q'}$.
\end{enumerate}
\end{lemma}
\leanref{Development/WeilBound/ThetaLifts/FrameEstimates/SquareIntegrableWeilRange.lean}{PositivityTheory.ThKEst.LatInv}
\leanref{Development/WeilBound/ThetaLifts/FrameEstimates/LatInv.lean}{PositivityTheory.ThKEst.latInv_of}
\leanref{Development/WeilBound/ThetaLifts/FrameEstimates/SquareIntegrableWeilRange.lean}{PositivityTheory.ThKEst.LatMeas}
\leanref{Development/WeilBound/ThetaLifts/FrameEstimates/LatMeas.lean}{PositivityTheory.ThKEst.latMeas}
\leanref{Development/WeilBound/ThetaLifts/FrameEstimates/SquareIntegrableWeilRange.lean}{PositivityTheory.ThKEst.LatCpt}
\leanref{Development/WeilBound/ThetaLifts/FrameEstimates/LatCpt.lean}{PositivityTheory.ThKEst.latCpt_of}
\leanref{Development/WeilBound/ThetaLifts/FrameEstimates/SquareIntegrableWeilRange.lean}{PositivityTheory.ThKEst.LatArch}
\leanref{Development/WeilBound/ThetaLifts/FrameEstimates/LatArch.lean}{PositivityTheory.ThKEst.latArch_of}
\leanref{Development/WeilBound/ThetaLifts/FrameEstimates/SquareIntegrableWeilRange.lean}{PositivityTheory.ThKEst.LatTor}
\leanref{Development/WeilBound/ThetaLifts/FrameEstimates/LatTor.lean}{PositivityTheory.ThKEst.latTor_of}
\leanref{Development/WeilBound/ThetaLifts/FrameEstimates/SquareIntegrableWeilRange.lean}{PositivityTheory.ThKEst.FrameBasis}
\leanref{Development/WeilBound/ThetaLifts/FrameEstimates/FrameBasis.lean}{PositivityTheory.ThKEst.frameBasis_of}
\leanref{Development/WeilBound/ThetaLifts/FrameEstimates/IntMvPolynomialCoeffAbsSumMulLe.lean}{PositivityTheory.ThKEst.diag_coords}
\leanref{Development/WeilBound/ThetaLifts/FrameEstimates/SquareIntegrableWeilRange.lean}{PositivityTheory.ThKEst.OneVar}
\leanref{Development/WeilBound/ThetaLifts/FrameEstimates/OneVar.lean}{PositivityTheory.ThKEst.oneVar_of}
\begin{proof}
(1) The matrices $\gamma^{-1}$ and $\gamma'^{-1}$ have entries in $E$, so
$(\gamma g)^{-1}x((\gamma'h)^{-1})^{\mathsf T}=g^{-1}x'(h^{-1})^{\mathsf T}$
with $x'=\gamma^{-1}x(\gamma'^{-1})^{\mathsf T}$, and $x\mapsto x'$ is a
bijection of $M_{n\times d}(E)$.

(2) Each term is a measurable function of $g$: inversion and $\|\ \|_\infty$
are continuous, and the set of $z$ such that $z-z_0$ is at $\infty$ for some
$z_0\in C$ is closed, as $C$ is compact. A countable sum of measurable
functions with values in $[0,\infty]$ is measurable.

(3) For $A\in M_n(\mathbf A_E)$, $A'\in M_d(\mathbf A_E)$ and
$z\in M_{n\times d}(\mathbf A_E)$,
$\|AzA'^{\mathsf T}\|_\infty\leq nd\,\|A\|_\infty\|A'\|_\infty\|z\|_\infty$; and
if $z-z_0$ is at $\infty$, so is $AzA'^{\mathsf T}-Az_0A'^{\mathsf T}$. Let
$M_0$ and $M'_0$ bound $\|\ \|_\infty$ on $\Omega$ and $\Omega'$, put
$K_0=\max(1,ndM_0M'_0)$, and let $C'$ be the compact set of the
$\kappa z_0\kappa'^{\mathsf T}$ with $\kappa\in\Omega$, $z_0\in C$ and
$\kappa'\in\Omega'$. Then
$m_{C,N}(z)\leq K_0^Nm_{C',N}(\kappa z\kappa'^{\mathsf T})$ for
$\kappa\in\Omega$ and $\kappa'\in\Omega'$. For
$z=(g\kappa)^{-1}x((h\kappa')^{-1})^{\mathsf T}$ one has
$\kappa z\kappa'^{\mathsf T}=g^{-1}x(h^{-1})^{\mathsf T}$, and summing over $x$
gives (3) with $K=K_0^N$.

(4) The vectors $x'_i$ and $y'_i$ of the Witt torus and a basis of their
orthogonal complement form a basis of $E^d$, the columns of some
$b'\in GL_d(E)$, and
$a'(t)b'=b'\operatorname{diag}(u_{\lambda_1(t)},\ldots,u_{\lambda_d(t)})$,
where each $\lambda_k$ is one of the linear forms $t\mapsto t_i$,
$t\mapsto-t_i$ and $0$. So $a'(t)-1$ is at $\infty$, and for
$t\in\mathcal C_{q'}$, where every $\lambda_k(t)$ is at most
$t_1+\cdots+t_{q'}$, the entries satisfy
$|\pi_w(a'(t)_{ij})|\leq K_1e^{t_1+\cdots+t_{q'}}$ for a constant $K_1$. Hence
$\|z\,a'(t)^{\mathsf T}\|_\infty\leq dK_1e^{\sum_jt_j}\|z\|_\infty$, and
$za'(t)^{\mathsf T}-z_0$ is at $\infty$ when $z-z_0$ is. So
$m_{C,N}(z)\leq\bigl((1+dK_1)e^{\sum_jt_j}\bigr)^Nm_{C,N}(za'(t)^{\mathsf T})$.
Applied to $z=g^{-1}x(a'(t)^{-1})^{\mathsf T}$ and summed over $x$, this gives
(4) with $C'=C$ and $K=(1+dK_1)^N$.

(5) Take bases $b$ and $b'$ adapted to $a$ and $a'$ as in (4), with linear
forms $\lambda_k$ on $\mathbf R^q$ and $\lambda'_l$ on $\mathbf R^{q'}$. For
$x\in M_{n\times d}(E)$ put $\xi=b^{-1}x(b'^{-1})^{\mathsf T}$ and
$z=a(s)^{-1}x(a'(t)^{-1})^{\mathsf T}$; then
$b^{-1}z(b'^{-1})^{\mathsf T}$ has entries
$u_{-\lambda_k(s)-\lambda'_l(t)}\xi_{kl}$. As $C$ is compact, there are an
integer $m\geq1$ and $K_2\geq1$ such that: if $m_{C,N}(z)\neq0$, then every
$m\xi_{kl}$ is an algebraic integer, since the finite components of
$b^{-1}z(b'^{-1})^{\mathsf T}$ are those of an element of a compact set; and
$\sum_w\sum_{k,l}e^{-\lambda_k(s)-\lambda'_l(t)}|\tau_w(\xi_{kl})|\leq K_2\|z\|_\infty$.
Put $N_1=[E:\mathbf Q]+1$, $N_0=N_1nd$ and
$\nu_{kl}=\lambda_k(s)+\lambda'_l(t)$. Then
\[
m_{C,N_0}(z)\leq K_2^{N_0}\prod_{k,l}\Bigl(1+e^{-\nu_{kl}}\sum_w|\tau_w(\xi_{kl})|\Bigr)^{-N_1}
\]
when $m\xi$ is integral, and $m_{C,N_0}(z)=0$ otherwise. As $x\mapsto\xi$ is
injective, summing over $x$ bounds $L_{C,N_0}(a(s),a'(t))$ by
$K_2^{N_0}\prod_{k,l}S(\nu_{kl})$, where $S(\nu)$ is the sum of
$(1+e^{-\nu}\sum_w|\tau_w(\eta)|)^{-N_1}$ over the $\eta\in E$ with $m\eta$
integral. These $m\eta$ form a lattice in a real space of dimension
$[E:\mathbf Q]<N_1$, and comparing the sum with an integral over disjoint balls
around the lattice points, then rescaling, gives
$S(\nu)\leq K_3e^{[E:\mathbf Q]\max(0,\nu)}$ for a constant $K_3$. Finally,
for $s\in\mathcal C_q$ and $t\in\mathcal C_{q'}$,
$\sum_{k,l}\max(0,\nu_{kl})=P_{n,d}(s,t)$. Indeed the values $\lambda_k(s)$
are the $s_i$, the $-s_i$ and $n-2q$ zeros, the values $\lambda'_l(t)$ are the
$t_j$, the $-t_j$ and $d-2q'$ zeros, and for $a,b\geq0$
\[
\max(0,a+b)+\max(0,a-b)+\max(0,-a+b)+\max(0,-a-b)=2\max(a,b),\qquad
\max(0,a)+\max(0,-a)=a .
\]
\end{proof}

\begin{lemma}\label{asm-theta-agree-convergence-cone}
Let $n,d,q,q'\geq0$, with $\mathcal C_q$, $\delta_n$ and $P_{n,d}$ as in
Definition~\ref{asm-theta-agree-convergence-tori}(3) and (5).
\begin{enumerate}
\item $P_{n,d}(s,t)\leq d\sum_is_i+n\sum_jt_j$ for $s\in\mathcal C_q$ and
  $t\in\mathcal C_{q'}$.
\item If $q=0$ or $d+2q<2n$, then for every $D>0$
  \[
  \int_{\mathcal C_q}\exp\Bigl(D\Bigl(\frac d2\sum_{i=1}^qs_i-\delta_n(s)\Bigr)\Bigr)\,ds<\infty .
  \]
\end{enumerate}
\end{lemma}
\leanref{Development/WeilBound/ThetaLifts/FrameEstimates/SquareIntegrableWeilRange.lean}{PositivityTheory.ThKEst.PosWt}
\leanref{Development/WeilBound/ThetaLifts/FrameEstimates/InitialPairs/SquareIntegrableIntegralHolds.lean}{PositivityTheory.ThKEst.Sq.posWt'}
\leanref{Development/WeilBound/ThetaLifts/FrameEstimates/SquareIntegrableWeilRange.lean}{PositivityTheory.ThKEst.ConeIntegralsFinite}
\leanref{Development/WeilBound/ThetaLifts/FrameEstimates/InitialPairs/SquareIntegrableIntegralHolds.lean}{PositivityTheory.ThKEst.Sq.square_integrable_integral_holds}
\leanref{Development/WeilBound/ThetaLifts/FrameEstimates/SquareIntegrableWeilRange.lean}{PositivityTheory.ThKEst.ConeExp}
\begin{proof}
(1) As $2\max(s_i,t_j)=2s_i+2t_j-2\min(s_i,t_j)$,
$P_{n,d}(s,t)=d\sum_is_i+n\sum_jt_j-2\sum_{i,j}\min(s_i,t_j)$, and every
$\min(s_i,t_j)$ is nonnegative.

(2) For $q=0$ the integral is $1$. Otherwise the exponent is
$\sum_{i=1}^qc_is_i$ with $c_i=D(\frac d2-n+2i-1)$, and
$c_1+\cdots+c_k=Dk(\frac d2-n+k)<0$ for $1\leq k\leq q$, as $d+2q<2n$. The map
$s\mapsto v$ with $v_k=s_k-s_{k+1}$, where $s_{q+1}=0$, carries $\mathcal C_q$
onto $[0,\infty)^q$ with Jacobian $1$, and the exponent becomes
$\sum_kv_k(c_1+\cdots+c_k)$. So the integral is
$\prod_{k=1}^q|c_1+\cdots+c_k|^{-1}$.
\end{proof}

\begin{lemma}\label{asm-theta-agree-convergence-dominated}
Let $(X,W,M,\chi_X,\chi_W,\omega)$ be a Weil pair of signature $(r,s;r',s')$
for a theta base $B$ of a CM field $E$
(Definition~\ref{asm-theta-agree-triple}(1)), with $X$ and $W$ of ranks $n$ and
$d$, and suppose that $d+2\min(r,s)<2n$ or $\min(r,s)=0$. Put
$I=\{1,\ldots,n\}\times\{1,\ldots,d\}$, and let $\mu$ and $\mathcal D$ be the
measure and the domain of $M$. Then there is $k\geq0$ such that for every
Schwartz-Bruhat function $\phi$ on $\mathbf A_f^I$ there are a continuous
seminorm $p$ on $\mathcal S(F_\infty^I)$ and a measurable function
$G\colon U_X\to[0,\infty]$ with $\int_{\mathcal D}G\,d\mu<\infty$ and
\[
|\theta_\omega(f\otimes\phi)(g,h)|\leq p(f)\operatorname{ht}(h)^kG(g)
\]
for all $f\in\mathcal S(F_\infty^I)$, $g\in U_X$ and $h\in U_W$.
\end{lemma}
\leanref{Development/WeilBound/ThetaLifts/FrameEstimates/HodgeTateFromSandwich.lean}{PositivityTheory.ThKEst.DomWeilU}
\leanref{Development/WeilBound/ThetaLifts/FrameEstimates/WeilConv.lean}{PositivityTheory.ThKEst.domWeilU}
\leanref{Development/WeilBound/ThetaLifts/FrameEstimates/DomWeil.lean}{PositivityTheory.ThKEst.domWeilU_of}
\leanref{Development/WeilBound/ThetaLifts/FrameEstimates/LatSumRedH.lean}{PositivityTheory.ThKEst.latSum_redH}
\leanref{Development/WeilBound/ThetaLifts/FrameEstimates/LatSumRedH.lean}{PositivityTheory.ThKEst.exp_le_ht}
\begin{proof}
By Lemma~\ref{asm-theta-agree-convergence-signature}, $J_X$ and $J_W$ are
nondegenerate $\epsilon$-Hermitian, and by
Theorem~\ref{asm-theta-agree-convergence-siegel}(2) they have Siegel Witt tori,
$a$ of size $q$ and $a'$ of size $q'$. By
Lemma~\ref{asm-theta-agree-convergence-signature} again, $q\leq\min(r,s)$, so
$q=0$ or $d+2q<2n$. Let $N$ be the integer $N_0$ of
Lemma~\ref{asm-theta-agree-convergence-lattice}(5) for $a$ and $a'$, take the
lattice sums for $J_X$ and $J_W$, and put $D=[E:\mathbf Q]$.

\emph{Step 1: the variable $h$.} There is $k$ such that for every compact $C$
there are a compact $C'$ and $K\geq0$ with
$L_{C,N}(g,h)\leq(K\operatorname{ht}(h)^k)^2L_{C',N}(g,1)$ for all $g$ and $h$.
Let $\Omega'$ be the compact set of the Siegel property (a) of $a'$. There are
$A\geq0$ and $k_1$ such that $e^{t_1+\cdots+t_{q'}}\leq A\operatorname{ht}(h)^{k_1}$
whenever $\gamma'h=a'(t)\kappa'$ with $\gamma'\in U_W(F)$,
$t\in\mathcal C_{q'}$ and $\kappa'\in\Omega'$: then
$\gamma'^{-1}a'(t)=h\kappa'^{-1}$, so this follows from the Siegel property
(c) of $a'$, as the height is submultiplicative, unchanged by inversion and
bounded on $\Omega'$ by Theorem~\ref{asm-theta-agree-convergence-siegel}(1).
Put $k=k_1N$. Given $C$, Lemma~\ref{asm-theta-agree-convergence-lattice}(3)
for $\{1\}$ and $\Omega'$ gives $C_1$ and $K_1$, and
Lemma~\ref{asm-theta-agree-convergence-lattice}(4) for $a'$ and $C_1$ gives
$C'$ and $K_2$. For $h\in U_W$ write $\gamma'h=a'(t)\kappa'$ by the Siegel
property (a). By Lemma~\ref{asm-theta-agree-convergence-lattice}(1), (3) and
(4),
\[
L_{C,N}(g,h)=L_{C,N}(g,a'(t)\kappa')\leq K_1L_{C_1,N}(g,a'(t))
\leq K_1K_2e^{N\sum_jt_j}L_{C',N}(g,1)\leq K_1K_2A^N\operatorname{ht}(h)^kL_{C',N}(g,1),
\]
and $K_1K_2A^N\operatorname{ht}(h)^k\leq\bigl((K_1K_2A^N+1)\operatorname{ht}(h)^k\bigr)^2$,
as $\operatorname{ht}(h)\geq1$.

\emph{Step 2: the bound.} Let $\phi$ be given. Take $p_0$ and $C$ from
Lemma~\ref{asm-theta-agree-convergence-square} for $\phi$ and $N$, and then
$C'$ and $K$ from Step 1. Put $p=Kp_0$ and
$G(g)=L_{C',N}(g,1)^{1/2}$, which is measurable by
Lemma~\ref{asm-theta-agree-convergence-lattice}(2). Then
\[
|\theta_\omega(f\otimes\phi)(g,h)|^2\leq p_0(f)^2L_{C,N}(g,h)\leq\bigl(p(f)\operatorname{ht}(h)^kG(g)\bigr)^2 .
\]

\emph{Step 3: integrability.} Let $K_0$ and $\Omega$ be given by the Siegel
property (b) of $a$ for $\mu$ and $\mathcal D$.
Lemma~\ref{asm-theta-agree-convergence-lattice}(3) for $\Omega$ and $\{1\}$
gives $C_2$ and $K_3$, and Lemma~\ref{asm-theta-agree-convergence-lattice}(5)
for $C_2$ gives $K_4$. For $t\in\mathcal C_q$ and $\kappa\in\Omega$, as
$a'(0)=1$,
\[
G(a(t)\kappa)^2\leq K_3L_{C_2,N}(a(t),a'(0))\leq K_3K_4e^{DP_{n,d}(t,0)}\leq K_3K_4e^{Dd\sum_it_i},
\]
by Lemma~\ref{asm-theta-agree-convergence-cone}(1). By
Lemma~\ref{asm-theta-agree-convergence-lattice}(1), $G(\gamma g)=G(g)$ for
$\gamma\in U_X(F)$. So the Siegel property (b) applies to $G$ and the function
$t\mapsto(K_3K_4)^{1/2}e^{D\frac d2\sum_it_i}$, and
\[
\int_{\mathcal D}G\,d\mu\leq K_0(K_3K_4)^{1/2}\int_{\mathcal C_q}e^{D(\frac d2\sum_it_i-\delta_n(t))}\,dt<\infty
\]
by Lemma~\ref{asm-theta-agree-convergence-cone}(2).
\end{proof}

\begin{lemma}\label{asm-theta-agree-convergence}
Let $(X,W,M,\chi_X,\chi_W,\omega)$ be a Weil pair of signature $(r,s;r',s')$
for a theta base $B$ of a CM field $E$
(Definition~\ref{asm-theta-agree-triple}(1)), with $X$ and $W$ of ranks $n$ and
$d$, and suppose that $d+2\min(r,s)<2n$ or $\min(r,s)=0$. Then:
\begin{enumerate}
\item every Schwartz-Bruhat function $\Phi$ on
  $\mathbf A^{n\times d}$ converges, and there are $C$ and $k$ with
  $\int_{\mathcal D}|\theta_\omega(\Phi)(g,h)|\,d\mu(g)\leq C\operatorname{ht}(h)^k$
  for every $h\in U_W$;
\item for every Schwartz-Bruhat function $\phi$ on $\mathbf A_f^{n\times d}$ and
  every $h\in U_W$, the map $f\mapsto I(f\otimes\phi)(h)$ is continuous on
  $\mathcal S(F_\infty^{n\times d})$.
\end{enumerate}
\end{lemma}
\leanref{Development/AbelianVarietyTheory/Definitions/Divisors/CTSmallCone.lean}{GlobalTheory.ThK.WeilConv}
\leanref{Development/WeilBound/ThetaLifts/FrameEstimates/WeilConv.lean}{PositivityTheory.ThKEst.weilConv}
\leanref{Development/WeilBound/ThetaLifts/FrameEstimates/ThickFrames.lean}{PositivityTheory.ThKEst.weilConv_of}
\leanref{Development/WeilBound/ThetaLifts/FrameEstimates/ThickFrames.lean}{PositivityTheory.ThKEst.dom_facts}
\leanref{Development/WeilBound/ThetaLifts/FrameEstimates/ThickFrames.lean}{PositivityTheory.ThKEst.dom_sum}
\begin{proof}
Let $k$ be as in Lemma~\ref{asm-theta-agree-convergence-dominated}, and put
$I=\{1,\ldots,n\}\times\{1,\ldots,d\}$.

\emph{Pure tensors.} Let $\phi$ be a Schwartz-Bruhat function on
$\mathbf A_f^I$, and let $p$ and $G$ be as in
Lemma~\ref{asm-theta-agree-convergence-dominated} for $\phi$. Fix $h\in U_W$
and $f\in\mathcal S(F_\infty^I)$. The function
$g\mapsto\theta_\omega(f\otimes\phi)(g,h)$ is continuous by
Lemma~\ref{asm-theta-agree-convergence-theta}(2), hence measurable, and it is
bounded by $p(f)\operatorname{ht}(h)^kG$, which is integrable on $\mathcal D$.
So $f\otimes\phi$ converges, and
\[
\int_{\mathcal D}|\theta_\omega(f\otimes\phi)(g,h)|\,d\mu(g)\leq p(f)\Bigl(\int_{\mathcal D}G\,d\mu\Bigr)\operatorname{ht}(h)^k .
\]

\emph{Part (2).} Let $f_0\in\mathcal S(F_\infty^I)$. As $p$ is continuous,
$p(f)<p(f_0)+1$ for $f$ in a neighbourhood of $f_0$. For these $f$ the
functions $g\mapsto\theta_\omega(f\otimes\phi)(g,h)$ are bounded by the
integrable function $(p(f_0)+1)\operatorname{ht}(h)^kG$, and as $f\to f_0$
they converge pointwise to $g\mapsto\theta_\omega(f_0\otimes\phi)(g,h)$, by
Lemma~\ref{asm-theta-agree-convergence-theta}(3). By dominated convergence,
$f\mapsto I(f\otimes\phi)(h)$ is continuous at $f_0$.

\emph{Part (1).} A Schwartz-Bruhat function $\Phi$ on $\mathbf A^I$ is a finite
sum $\sum_if_i\otimes\phi_i$ (Definition~\ref{asm-theta-agree-models}), and
$\theta_\omega(\Phi)=\sum_i\theta_\omega(f_i\otimes\phi_i)$ by
Lemma~\ref{asm-theta-agree-convergence-theta}(2). By the first paragraph each
term is integrable on $\mathcal D$, with
$\int_{\mathcal D}|\theta_\omega(f_i\otimes\phi_i)(g,h)|\,d\mu(g)\leq C_i\operatorname{ht}(h)^k$.
So $\Phi$ converges, and (1) holds with $C=\sum_iC_i$.

The convergence in (1) is Weil's convergence criterion for theta integrals
(Theorem~\ref{lit-weil-siegel}(2)) in the form Alpöge uses. Weil's condition
is on the Witt index of $X$, which is at most $\min(r,s)$, so the hypothesis
here implies it \cite{Wei65,KR88}.
\end{proof}
% Lean: weilConv hSiegelSet thetaBasic = weilConv_of thetaBasic (domWeil_of_U (domWeilU
% hSiegelSet)) soFin_of; in Theorem K, hSiegelSet = ThKR.siegel_of_core Prw.siegelCore_holds.

\begin{definition}\label{asm-theta-agree-continuity-sums}
Let $E$ be a CM field, Galois over $\mathbf Q$, with maximal totally real
subfield $F$, and let $B$ be a theta base for $E$
(Definition~\ref{asm-theta-agree-models}(1)).
\begin{enumerate}
\item Let $J\in M_n(E)$ and $J'\in M_d(E)$ be nondegenerate
  $\epsilon$-Hermitian (Definition~\ref{asm-theta-agree-convergence-tori}(1)),
  let $C\subset M_{n\times d}(\mathbf A_E)$ be compact, and let $N\geq0$. The
  \emph{full-rank lattice sum} is
  \[
  L^{\mathrm{full}}_{C,N}(g,h)=\sum_xm_{C,N}\bigl(g^{-1}x\,(h^{-1})^{\mathsf T}\bigr)\in[0,\infty]
  \qquad(g\in U_J,\ h\in U_{J'}),
  \]
  the sum over the $x\in M_{n\times d}(E)$ of rank $n$, with the majorant of
  Definition~\ref{asm-theta-agree-convergence-tori}(4).
\item For $s\in\mathbf R^q$ and an integer $i\geq0$ let $\sigma_i(s)$ be the sum
  of the first $\min(i,q)$ coordinates of $s$. For $s\in\mathbf R^q$ and
  $t\in\mathbf R^{q'}$, the \emph{excess} is
  $\Delta(s,t)=\max_{0\leq i\leq q}\bigl(\sigma_i(s)-\sigma_i(t)\bigr)$.
\item For a finite set $I$ and a function $\Phi$ on $\mathbf A^I$, put
  $\Phi^-(x)=\Phi(-x)$.
\item Let $(X,W,M,\chi_X,\chi_W,\omega)$ be a Weil pair for $B$
  (Definition~\ref{asm-theta-agree-triple}(1)), with $X$ and $W$ of ranks $n$
  and $d$. A \emph{full-rank finite part} for $M$ is a Schwartz-Bruhat function
  $\phi$ on $\mathbf A_f^{n\times d}$ such that $f\otimes\phi\in\mathcal S^\circ(M)$
  for every $f\in\mathcal S(F_\infty^{n\times d})$.
\end{enumerate}
\end{definition}
\leanref{Development/WeilBound/ThetaLifts/FrameEstimates/SquareIntegrableWeilRange.lean}{PositivityTheory.ThKEst.fullRank}
\leanref{Development/WeilBound/ThetaLifts/FrameEstimates/SquareIntegrableWeilRange.lean}{PositivityTheory.ThKEst.psum}
\leanref{Development/WeilBound/ThetaLifts/FrameEstimates/SquareIntegrableWeilRange.lean}{PositivityTheory.ThKEst.viol}
\leanref{Development/WeilBound/ThetaLifts/FrameEstimates/SquareIntegrableWeilRange.lean}{PositivityTheory.ThKEst.soFin}

\begin{lemma}\label{asm-theta-agree-continuity-pieces}
Let $(X,W,M,\chi_X,\chi_W,\omega)$ be a Weil pair for a theta base $B$ of a CM
field $E$ (Definition~\ref{asm-theta-agree-triple}(1)), with $X$ and $W$ of
ranks $n$ and $d$. Then:
\begin{enumerate}
\item every $\Phi\in\mathcal S^\circ(M)$ is a finite sum $\sum_if_i\otimes\phi_i$
  with $f_i\in\mathcal S(F_\infty^{n\times d})$ and full-rank finite parts
  $\phi_i$ for $M$ (Definition~\ref{asm-theta-agree-continuity-sums}(4));
\item for every full-rank finite part $\phi$ for $M$ there are finitely many
  $c_i\in\mathbf C$ and, for each $i$ and each finite place $v$, a
  Schwartz-Bruhat function $\chi_{i,v}$ on $F_v^{n\times d}$, with
  $\chi_{i,v}=\mathrm{vac}_v$ for almost all $v$ and $\chi_{i,v_1}$ of full
  rank at $\tau_1$ (Definition~\ref{asm-theta-agree-models}(5)), such that
  $f\otimes\phi=\sum_ic_i\,f\otimes\bigotimes_v\chi_{i,v}$ for every
  $f\in\mathcal S(F_\infty^{n\times d})$.
\end{enumerate}
\end{lemma}
\leanref{Development/WeilBound/ThetaLifts/FrameEstimates/SquareIntegrableWeilRange.lean}{PositivityTheory.ThKEst.SoFin}
\leanref{Development/WeilBound/ThetaLifts/FrameEstimates/IsCompactIntV.lean}{PositivityTheory.ThKEst.soFin_of}
\begin{proof}
(1) By definition, $\mathcal S^\circ(M)$ is spanned by the functions
$f\otimes\bigotimes_v\phi_v$ of Definition~\ref{asm-theta-agree-models}(5).
For such a function, $\phi=\bigotimes_v\phi_v$ is a Schwartz-Bruhat function
on $\mathbf A_f^{n\times d}$, since every $\phi_v$ is one and, $M$ being
regular, $\phi_v=\mathrm{vac}_v$ is the indicator function of
$\mathcal O_v^{n\times d}$ for almost all $v$. For every
$f'\in\mathcal S(F_\infty^{n\times d})$ the function $f'\otimes\phi$ is again
one of the spanning functions. So $\phi$ is a full-rank finite part. An element
of $\mathcal S^\circ(M)$ is a finite sum $\sum_ic_i\,f_i\otimes\phi_i$ of
multiples of spanning functions, and $c_i\,f_i\otimes\phi_i=(c_if_i)\otimes\phi_i$.

(2) Choose $f_0\in\mathcal S(F_\infty^{n\times d})$ with $f_0(0)=1$. Since
$f_0\otimes\phi\in\mathcal S^\circ(M)$, it is a finite sum
$\sum_ic_i\,f_i\otimes\bigotimes_v\chi_{i,v}$ of multiples of spanning
functions. Evaluating at the points with archimedean component $0$ gives
$\phi=\sum_ic_if_i(0)\bigotimes_v\chi_{i,v}$ on $\mathbf A_f^{n\times d}$.
Hence $f\otimes\phi=\sum_ic_if_i(0)\,f\otimes\bigotimes_v\chi_{i,v}$ for every
$f$.
\end{proof}

\begin{lemma}\label{asm-theta-agree-continuity-cut-local}
Let $k$ be a nonarchimedean local field of characteristic $0$, let $\psi_k$ be
a locally constant character of $k$ with $\psi_k\neq1$, and let $n,d\geq0$.
Let $f$ be a Schwartz-Bruhat function on $M_{n\times d}(k)$ with $f(Y)=0$ for
every $Y$ of rank different from $n$. Then $f=\sum_ig_i$ for finitely many
Schwartz-Bruhat functions $g_i$ such that $g_i(Y)=0$ whenever $f(Y)=0$, and for
each $i$ there is $Y_i\in M_{n\times d}(k)$ with the following property: for
every $Z\in M_{n\times d}(k)$ of rank different from $n$ there is
$U\in M_{n\times d}(k)$ with
$\psi_k(\operatorname{tr}(YU^{\mathsf T}))=\psi_k(\operatorname{tr}(Y_iU^{\mathsf T}))$
whenever $g_i(Y)\neq0$, and
$\psi_k(\operatorname{tr}((Z-2Y_i)U^{\mathsf T}))\neq1$.
\end{lemma}
\leanref{Development/WeilBound/ThetaLifts/FrameEstimates/RankIffRight.lean}{PositivityTheory.ThKEst.W.exists_cut}
\leanref{Development/WeilBound/ThetaLifts/FrameEstimates/RankIffRight.lean}{PositivityTheory.ThKEst.W.exists_dual}
\begin{proof}
The set $S$ of the $Y$ with $f(Y)\neq0$ is compact, and its elements have rank
$n$. The matrices of rank $n$ form an open set $R$, since $Y$ has rank $n$ if
and only if $\det(YC)\neq0$ for some $C\in M_{d\times n}(k)$. With $Y$, also
$2Y$ has rank $n$. So the compact set $2S=\{2Y:Y\in S\}$ lies in $R$, and there
is a neighbourhood $V$ of $0$ with $2S+V\subset R$.

Choose $b\in k$ with $\psi_k(b)\neq1$, and $a\in k^\times$ such that $aY\in V$
for every $Y$ with integral entries; $a$ exists because the matrices with
integral entries form a compact set. Let $K$ be the set of the matrices
$(b/a)Y$ with $Y$ integral, and $\Lambda$ the set of the $Y$ with
$\psi_k(\operatorname{tr}(YU^{\mathsf T}))=1$ for every $U\in K$. Then
$\Lambda$ is a subgroup, and it is open, since $\psi_k$ is locally constant and
$K$ is compact. Moreover $\Lambda\subset V$. Indeed, let $Y\in\Lambda$, and
suppose that an entry $y=(a^{-1}Y)_{ij}$ is not integral. Then $y^{-1}$ is
integral, so $U=(b/a)y^{-1}E_{ij}$ lies in $K$, where $E_{ij}$ is the matrix
unit. But $\operatorname{tr}(YU^{\mathsf T})=b$, so
$\psi_k(\operatorname{tr}(YU^{\mathsf T}))\neq1$, which is absurd. Hence
$a^{-1}Y$ is integral and $Y=a(a^{-1}Y)\in V$.

The compact set $S$ meets finitely many cosets $Y_i+\Lambda$, and we may take
$Y_i\in S$. Let $g_i$ be $f$ times the indicator function of $Y_i+\Lambda$.
These are Schwartz-Bruhat functions with sum $f$, and $g_i$ vanishes where $f$
does. Let $Z$ have rank different from $n$. Then $Z-2Y_i\notin\Lambda$, since
otherwise $Z\in2S+V\subset R$. So there is $U\in K$ with
$\psi_k(\operatorname{tr}((Z-2Y_i)U^{\mathsf T}))\neq1$. If $g_i(Y)\neq0$, then
$Y-Y_i\in\Lambda$, so
$\psi_k(\operatorname{tr}(YU^{\mathsf T}))=\psi_k(\operatorname{tr}(Y_iU^{\mathsf T}))$.
\end{proof}

\begin{lemma}\label{asm-theta-agree-continuity-eigen}
Let $B$ be a theta base for a CM field $E$
(Definition~\ref{asm-theta-agree-models}(1)), with character $\psi$, let $I$ be
a finite set and $dx$ a Haar measure on $\mathbf A^I$, with the operators
$\rho(w)$ and the coefficients of
Definition~\ref{asm-theta-agree-convergence-tori}(6). Let $\Phi$ be a
Schwartz-Bruhat function on $\mathbf A^I$, and let
$u=(u_1,u_2)\in\mathbf A^I\times\mathbf A^I$ and $a\in\mathbf A$ be such that
$\rho(u)\Phi=\psi(a)\Phi$. Then for every
$w=(w_1,w_2)\in\mathbf A^I\times\mathbf A^I$,
\[
\langle\rho(w)\Phi,\Phi^-\rangle=\psi(u_1\cdot w_2-w_1\cdot u_2+2a)\,\langle\rho(w)\Phi,\Phi^-\rangle .
\]
\end{lemma}
\leanref{Development/WeilBound/ThetaLifts/FrameEstimates/HermitianFormDetNeZero.lean}{PositivityTheory.ThKEst.W.coef_eigen}
\begin{proof}
The substitution $x\mapsto x+u_1$, under which $dx$ is invariant, gives
\[
\langle\rho(w)\Phi,\Phi^-\rangle=\int\psi\bigl((x+u_1)\cdot w_2+\tfrac12w_1\cdot w_2\bigr)\Phi(x+u_1+w_1)\overline{\Phi(-x-u_1)}\,dx .
\]
The equation $\rho(u)\Phi=\psi(a)\Phi$ says that
$\Phi(y+u_1)=\psi(a-y\cdot u_2-\tfrac12u_1\cdot u_2)\Phi(y)$ for every $y$. At
$y=x+w_1$ it gives
$\Phi(x+u_1+w_1)=\psi(a-(x+w_1)\cdot u_2-\tfrac12u_1\cdot u_2)\Phi(x+w_1)$. At
$y=-x-u_1$ it gives $\Phi(-x)=\psi(a+x\cdot u_2+\tfrac12u_1\cdot u_2)\Phi(-x-u_1)$,
so, as $|\psi|=1$,
$\overline{\Phi(-x-u_1)}=\psi(a+x\cdot u_2+\tfrac12u_1\cdot u_2)\overline{\Phi(-x)}$.
Multiplying, the integrand is $\psi(u_1\cdot w_2-w_1\cdot u_2+2a)$ times
$\psi(x\cdot w_2+\tfrac12w_1\cdot w_2)\Phi(x+w_1)\overline{\Phi(-x)}$, the
integrand of $\langle\rho(w)\Phi,\Phi^-\rangle$.
\end{proof}

\begin{lemma}\label{asm-theta-agree-continuity-cut}
Let $(X,W,M,\chi_X,\chi_W,\omega)$ be a Weil pair for a theta base $B$ of a CM
field $E$ (Definition~\ref{asm-theta-agree-triple}(1)), with $X$ and $W$ of
ranks $n$ and $d$, and put $v=v_1$ and $t=\tau_1$. Let $\chi$ be a
Schwartz-Bruhat function on $F_v^{n\times d}$ of full rank at $t$
(Definition~\ref{asm-theta-agree-models}(5)). Then $\chi=\sum_i\chi_i$ for
finitely many functions $\chi_i$ of full rank at $t$ with the following
property. Let $r$ be a function on $\mathbf A^{n\times d}$ that does not depend
on the $v$-component and for which $\Phi=\chi_i\otimes r$ is a Schwartz-Bruhat
function, let $dx$ be a Haar measure on $\mathbf A^{n\times d}$, and let
$z\in\mathbb W(X,W)$ be such that $\pi_{t\circ c}(z)$ has rank different from
$n$. Then $\langle\rho(z)\Phi,\Phi^-\rangle=0$, with the coefficients of
Definition~\ref{asm-theta-agree-convergence-tori}(6).
\end{lemma}
\leanref{Development/WeilBound/ThetaLifts/FrameEstimates/SquareIntegrableWeilRange.lean}{PositivityTheory.ThKEst.SoCut}
\leanref{Development/WeilBound/ThetaLifts/FrameEstimates/HermitianFormDetNeZero.lean}{PositivityTheory.ThKEst.soCut_of}
\begin{proof}
Let $\iota$ be a linear model of the second kind for $(t,1,1)$ with
$\iota\chi(Y)=0$ for every $Y$ of rank different from $n$, and write
$\operatorname{co}(u)=-t(J_X)\,\pi_{t\circ c}(u)\,t(J_W)^{\mathsf T}$ for its
map (Definition~\ref{asm-theta-agree-models}(5)). The matrices $J_X$ and $J_W$
are invertible by Lemma~\ref{asm-theta-agree-convergence-signature}, so
$t(J_X)$ and $t(J_W)$ are invertible. Let $\psi_v$ be the character
$y\mapsto\psi(y)$ of $F_v$, with $y$ read as an adele at $v$. It is locally
constant, and nontrivial because $\psi$ is a standard character.

Lemma~\ref{asm-theta-agree-continuity-cut-local}, for $k=F_v$, $\psi_v$ and
$f=\iota\chi$, gives $\iota\chi=\sum_ig_i$ and matrices $Y_i$. Since $\iota$ is
a bijection between the Schwartz-Bruhat spaces, $g_i=\iota\chi_i$ for
Schwartz-Bruhat functions $\chi_i$ with $\chi=\sum_i\chi_i$. Each $\chi_i$ has
full rank at $t$, because $g_i$ vanishes wherever $\iota\chi$ does.

Fix $i$, and $r$, $\Phi$, $dx$ and $z$ as in the statement. Put
$Z=\operatorname{co}(z)$. Then $Z$ has rank different from $n$, as
$\pi_{t\circ c}(z)$ has. Lemma~\ref{asm-theta-agree-continuity-cut-local} gives
a matrix $U$ for $Z$. Let $w\in\mathbb W(X,W)$ be the element at $v$ with
$\pi_t(w)=-U$ and $\pi_{t\circ c}(w)=0$. The definition of
$\langle\ ,\ \rangle$ gives, for every $x\in\mathbb W(X,W)$, that
$\langle x,w\rangle$ is the adele at $v$ with component
$\operatorname{tr}(\operatorname{co}(x)U^{\mathsf T})$. Since $w$ is at $v$, so
is $e(w)$, and $\rho(w)$ acts at $v$ (Definition~\ref{asm-theta-agree-models}(2))
by the Schrödinger operator $\rho_v(w)$ on functions on $F_v^{n\times d}$
formed with $\psi_v$ and the $v$-components of $e(w)$; it maps Schwartz-Bruhat
functions to Schwartz-Bruhat functions. As $\pi_{t\circ c}(w)=0$, the defining
property of a linear model of the second kind gives
\[
\iota(\rho_v(w)\chi_i)(\operatorname{co}(x))=\psi_v\bigl(\operatorname{tr}(\operatorname{co}(x)U^{\mathsf T})\bigr)\,g_i(\operatorname{co}(x))
\]
for every $x$ at $v$ with $\pi_t(x)=0$. Every $Y\in M_{n\times d}(F_v)$ is
$\operatorname{co}(x)$ for such an $x$, namely the one with
$\pi_{t\circ c}(x)=-t(J_X)^{-1}Y(t(J_W)^{\mathsf T})^{-1}$. By the choice of
$U$,
$\psi_v(\operatorname{tr}(YU^{\mathsf T}))g_i(Y)=\psi_v(\operatorname{tr}(Y_iU^{\mathsf T}))g_i(Y)$
for every $Y$. So
$\iota(\rho_v(w)\chi_i)=\psi_v(\operatorname{tr}(Y_iU^{\mathsf T}))\,\iota\chi_i$,
and $\rho_v(w)\chi_i=\psi_v(\operatorname{tr}(Y_iU^{\mathsf T}))\,\chi_i$, since
$\iota$ is injective on Schwartz-Bruhat functions. Hence
$\rho(w)\Phi=\psi(a)\Phi$, where $a$ is the adele at $v$ with component
$\operatorname{tr}(Y_iU^{\mathsf T})$.

Lemma~\ref{asm-theta-agree-continuity-eigen}, with $u=e(w)$, gives
$\langle\rho(z)\Phi,\Phi^-\rangle=\psi(\langle w,z\rangle+2a)\,\langle\rho(z)\Phi,\Phi^-\rangle$,
since $e$ carries $\langle\ ,\ \rangle$ to the form
$e_1(w)\cdot e_2(z)-e_1(z)\cdot e_2(w)$. Here
$\langle w,z\rangle=-\langle z,w\rangle$, so
\[
\psi(\langle w,z\rangle+2a)=\psi_v\bigl(\operatorname{tr}((2Y_i-Z)U^{\mathsf T})\bigr)
=\psi_v\bigl(\operatorname{tr}((Z-2Y_i)U^{\mathsf T})\bigr)^{-1}\neq1 .
\]
So $\langle\rho(z)\Phi,\Phi^-\rangle=0$.
\end{proof}

\begin{lemma}\label{asm-theta-agree-continuity-fourier}
Let $B$ be a theta base for a CM field $E$
(Definition~\ref{asm-theta-agree-models}(1)), with character $\psi$, let $I$ be
a finite set and $dx$ a Haar measure on $\mathbf A^I$. Let $\Phi_1$, $\Phi_2$,
$\Psi_1$ and $\Psi_2$ be Schwartz-Bruhat functions on $\mathbf A^I$ with
$\langle\rho(w)\Phi_1,\Phi_2\rangle=\langle\rho(w)\Psi_1,\Psi_2\rangle$ for
every $w\in\mathbf A^I\times\mathbf A^I$
(Definition~\ref{asm-theta-agree-convergence-tori}(6)). Then
$\Phi_1(x)\overline{\Phi_2(u)}=\Psi_1(x)\overline{\Psi_2(u)}$ for all
$x,u\in\mathbf A^I$.
\end{lemma}
\leanref{Development/WeilBound/ThetaLifts/FrameEstimates/SquareIntegrableWeilRange.lean}{PositivityTheory.ThKEst.FourierInj}
\leanref{Development/WeilBound/ThetaLifts/FrameEstimates/FourierInj.lean}{PositivityTheory.ThKEst.fourierInj_of}
\begin{proof}
For $w_1,y\in\mathbf A^I$,
$\langle\rho(w_1,y)\Phi_1,\Phi_2\rangle=\psi(\tfrac12w_1\cdot y)\int\psi(u\cdot y)\Phi_1(u+w_1)\overline{\Phi_2(u)}\,du$,
and likewise for $\Psi_1$ and $\Psi_2$. So the function
$u\mapsto\Phi_1(u+w_1)\overline{\Phi_2(u)}-\Psi_1(u+w_1)\overline{\Psi_2(u)}$,
which is continuous and integrable, has Fourier transform $0$ for the pairing
$(u,y)\mapsto\psi(u\cdot y)$. As $\mathbf A^I$ is self-dual for this pairing,
the Fourier transform is injective on continuous integrable functions
\cite{Tat67}, and the function is $0$. Taking $w_1=x-u$ gives the claim.
\end{proof}

\begin{lemma}\label{asm-theta-agree-continuity-parity}
Let $(X,W,M,\chi_X,\chi_W,\omega)$ be a Weil pair for a theta base $B$ of a CM
field $E$ (Definition~\ref{asm-theta-agree-triple}(1)), with $X$ and $W$ of
ranks $n$ and $d$. Then:
\begin{enumerate}
\item there is $c\in\mathbf C$ with $|c|=1$ such that
  $\omega(-1,1)\Phi=c\,\Phi^-$ for every Schwartz-Bruhat function $\Phi$ on
  $\mathbf A^{n\times d}$, with $\Phi^-$ as in
  Definition~\ref{asm-theta-agree-continuity-sums}(3);
\item $\theta_\omega(\Phi^-)(g,h)=\theta_\omega(\Phi)(g,h)$ for every
  Schwartz-Bruhat function $\Phi$ on $\mathbf A^{n\times d}$, $g\in U_X$ and
  $h\in U_W$.
\end{enumerate}
\end{lemma}
\leanref{Development/WeilBound/ThetaLifts/FrameEstimates/SquareIntegrableWeilRange.lean}{PositivityTheory.ThKEst.Parity}
\leanref{Development/WeilBound/ThetaLifts/FrameEstimates/Parity.lean}{PositivityTheory.ThKEst.parity_of}
\leanref{Development/WeilBound/ThetaLifts/FrameEstimates/ThetaSumNeg.lean}{PositivityTheory.ThKEst.theta_neg}
\begin{proof}
(1) Fix a Haar measure $dx$ on $\mathbf A^{n\times d}$; it is invariant under
$x\mapsto-x$. Let $\Phi_1$ and $\Phi_2$ be Schwartz-Bruhat functions. By
Lemma~\ref{asm-theta-agree-convergence-coefficients}(1) for $g=-1$ and $h=1$,
$\langle\rho(z)\,\omega(-1,1)\Phi_1,\omega(-1,1)\Phi_2\rangle=\langle\rho(-z)\Phi_1,\Phi_2\rangle$
for every $z\in\mathbb W(X,W)$. The substitution $x\mapsto-x$ gives
$\langle\rho(-w)\Phi_1,\Phi_2\rangle=\langle\rho(w)\Phi_1^-,\Phi_2^-\rangle$.
As $e$ is onto, Lemma~\ref{asm-theta-agree-continuity-fourier} gives
\[
(\omega(-1,1)\Phi_1)(x)\,\overline{(\omega(-1,1)\Phi_2)(u)}=\Phi_1(-x)\,\overline{\Phi_2(-u)}
\]
for all $x$ and $u$. Choose a Schwartz-Bruhat function $\Phi_0$ and $u$ with
$\Phi_0(-u)\neq0$. The identity for $\Phi_1=\Phi_2=\Phi_0$ and $x=u$ gives
$|(\omega(-1,1)\Phi_0)(u)|=|\Phi_0(-u)|\neq0$. Put
$c=\overline{\Phi_0(-u)}/\overline{(\omega(-1,1)\Phi_0)(u)}$, so that $|c|=1$.
The identity for $\Phi_1=\Phi$ and $\Phi_2=\Phi_0$ gives
$(\omega(-1,1)\Phi)(x)=c\,\Phi(-x)$ for every $x$.

(2) The element $-1$ is central in $U_X$, so $\omega(g,h)$ commutes with
$\omega(-1,1)$. As $\omega(g,h)\Phi$ is a Schwartz-Bruhat function, (1) gives
$c\,\omega(g,h)\Phi^-=\omega(g,h)\omega(-1,1)\Phi=\omega(-1,1)\omega(g,h)\Phi=c\,(\omega(g,h)\Phi)^-$.
As $c\neq0$, $\omega(g,h)\Phi^-=(\omega(g,h)\Phi)^-$. The theta sum of
$\Psi^-$ is that of $\Psi$, by the substitution $\xi\mapsto-\xi$ of
$F^{n\times d}$. So $\theta_\omega(\Phi^-)(g,h)=\theta_\omega(\Phi)(g,h)$.
\end{proof}

\begin{lemma}\label{asm-theta-agree-continuity-square}
Let $(X,W,M,\chi_X,\chi_W,\omega)$ be a Weil pair for a theta base $B$ of a CM
field $E$ (Definition~\ref{asm-theta-agree-triple}(1)), with $X$ and $W$ of
ranks $n$ and $d$ and matrices $J_X$ and $J_W$, let $\phi$ be a full-rank
finite part for $M$, and let $N\geq0$. Then there are a continuous seminorm
$p$ on $\mathcal S(F_\infty^{n\times d})$ and a compact set
$C\subset M_{n\times d}(\mathbf A_E)$ such that
\[
|\theta_\omega(f\otimes\phi)(g,h)|^2\leq p(f)^2\,L^{\mathrm{full}}_{C,N}(g,h)
\]
for all $f\in\mathcal S(F_\infty^{n\times d})$, $g\in U_X$ and $h\in U_W$, with
the full-rank lattice sum for $J_X$ and $J_W$
(Definition~\ref{asm-theta-agree-continuity-sums}(1)).
\end{lemma}
\leanref{Development/WeilBound/ThetaLifts/FrameEstimates/SquareIntegrableWeilRange.lean}{PositivityTheory.ThKEst.SqSo}
\leanref{Development/WeilBound/ThetaLifts/FrameEstimates/ThetaSumNeg.lean}{PositivityTheory.ThKEst.sqSo_of'}
\leanref{Development/WeilBound/ThetaLifts/FrameEstimates/ThetaSumNeg.lean}{PositivityTheory.ThKEst.qb_piece}
\begin{proof}
Say that a map $f\mapsto\Phi_f$ from $\mathcal S(F_\infty^{n\times d})$ to the
Schwartz-Bruhat functions on $\mathbf A^{n\times d}$ is \emph{bounded by
lattice sums} if there are $p$ and $C$ as in the statement with
$|\theta_\omega(\Phi_f)(g,h)|^2\leq p(f)^2L^{\mathrm{full}}_{C,N}(g,h)$ for all
$f$, $g$ and $h$. Sums and constant multiples of such maps are bounded by
lattice sums. Indeed, if $p,C$ serve for $f\mapsto\Phi_f$ and $q,C'$ for
$f\mapsto\Psi_f$, then, as the lattice sum increases with $C$ and the theta
function is additive (Lemma~\ref{asm-theta-agree-convergence-theta}(2)),
\[
|\theta_\omega(\Phi_f+\Psi_f)(g,h)|\leq|\theta_\omega(\Phi_f)(g,h)|+|\theta_\omega(\Psi_f)(g,h)|
\leq(p(f)+q(f))\,L^{\mathrm{full}}_{C\cup C',N}(g,h)^{1/2};
\]
and $|a|p$ and $C$ serve for $f\mapsto a\Phi_f$.

By Lemma~\ref{asm-theta-agree-continuity-pieces}(2),
$f\otimes\phi=\sum_ic_i\,f\otimes\bigotimes_v\chi_{i,v}$. So it suffices to
show that $f\mapsto f\otimes\bigotimes_v\chi_v$ is bounded by lattice sums
whenever the $\chi_v$ are Schwartz-Bruhat, $\chi_v=\mathrm{vac}_v$ for almost
all $v$ and $\chi_{v_1}$ has full rank at $\tau_1$.
Lemma~\ref{asm-theta-agree-continuity-cut} writes $\chi_{v_1}=\sum_j\xi_j$. Let
$\phi_j$ be $\bigotimes_v\chi_v$ with $\chi_{v_1}$ replaced by $\xi_j$. Then
$f\otimes\bigotimes_v\chi_v=\sum_jf\otimes\phi_j$, and $f\otimes\phi_j$ is
$\xi_j\otimes r$ for the function $r$ that equals
$f\otimes\bigotimes_{v\neq v_1}\chi_v$ on the components away from $v_1$ and
does not depend on the $v_1$-component. So
Lemma~\ref{asm-theta-agree-continuity-cut} gives:
\begin{itemize}
\item[$(*)$] $\langle\rho(z)\Phi,\Phi^-\rangle=0$ for $\Phi=f\otimes\phi_j$,
  every $f\in\mathcal S(F_\infty^{n\times d})$, every Haar measure on
  $\mathbf A^{n\times d}$ and every $z\in\mathbb W(X,W)$ such that
  $\pi_{\tau_1\circ c}(z)$ has rank different from $n$.
\end{itemize}
It remains to show that $f\mapsto f\otimes\phi'$ is bounded by lattice sums for
every Schwartz-Bruhat function $\phi'$ on $\mathbf A_f^{n\times d}$ that
satisfies $(*)$.

Fix a Haar measure $dx$ on $\mathbf A^{n\times d}$, let
$f\in\mathcal S(F_\infty^{n\times d})$, and put $\Phi=f\otimes\phi'$. As in the
proof of Lemma~\ref{asm-theta-agree-convergence-square}, $x\mapsto e(x)$ is a
bijection from $M_{n\times d}(E)$ onto $F^{n\times d}\times F^{n\times d}$. So
Lemma~\ref{asm-theta-agree-convergence-poisson}, for $\omega(g,h)\Phi$ and
$\omega(g,h)\Phi^-$, Lemma~\ref{asm-theta-agree-convergence-coefficients}(1)
and Lemma~\ref{asm-theta-agree-continuity-parity}(2) give $c>0$, independent of
$f$, with
\[
|\theta_\omega(\Phi)(g,h)|^2=|\theta_\omega(\Phi)(g,h)|\,|\theta_\omega(\Phi^-)(g,h)|
\leq c\sum_{x\in M_{n\times d}(E)}\bigl|\bigl\langle\rho\bigl(g^{-1}x(h^{-1})^{\mathsf T}\bigr)\Phi,\Phi^-\bigr\rangle\bigr| .
\]
If $x$ has rank less than $n$, so has
$\pi_{\tau_1\circ c}(g^{-1}x(h^{-1})^{\mathsf T})=\pi_{\tau_1\circ c}(g)^{-1}\,(\tau_1\circ c)(x)\,\pi_{\tau_1\circ c}(h^{-1})^{\mathsf T}$,
and the term of $x$ is $0$ by $(*)$. Now $\Phi^-=(Pf)\otimes\phi'^-$, where
$(Pf)(y)=f(-y)$ defines a continuous linear map $P$ of
$\mathcal S(F_\infty^{n\times d})$, and $\phi'^-$ is Schwartz-Bruhat.
Lemma~\ref{asm-theta-agree-convergence-coefficients}(2), for $\phi'$,
$\phi'^-$ and $N$, gives a continuous seminorm $p$ and a compact set $C$ with
$|\langle\rho(z)(f_1\otimes\phi'),f_2\otimes\phi'^-\rangle|\leq p(f_1)p(f_2)\,m_{C,N}(z)$.
Hence
\[
|\theta_\omega(\Phi)(g,h)|^2\leq c\,p(f)\,p(Pf)\,L^{\mathrm{full}}_{C,N}(g,h)
\leq\bigl(\sqrt c\,q(f)\bigr)^2\,L^{\mathrm{full}}_{C,N}(g,h)
\]
for the continuous seminorm $q=p+p\circ P$.
\end{proof}

\begin{lemma}\label{asm-theta-agree-continuity-full-sums}
Let $B$ be a theta base for a CM field $E$
(Definition~\ref{asm-theta-agree-models}(1)), let $J\in M_n(E)$ and
$J'\in M_d(E)$ be nondegenerate $\epsilon$-Hermitian, let
$C\subset M_{n\times d}(\mathbf A_E)$ be compact and $N\geq0$, and let
$L^{\mathrm{full}}_{C,N}$ be the full-rank lattice sums of
Definition~\ref{asm-theta-agree-continuity-sums}(1). Then:
\begin{enumerate}
\item $L^{\mathrm{full}}_{C,N}(\gamma g,\gamma'h)=L^{\mathrm{full}}_{C,N}(g,h)$
  for $\gamma\in U_J(F)$, $\gamma'\in U_{J'}(F)$, $g\in U_J$ and
  $h\in U_{J'}$;
\item for compact sets $\Omega\subset U_J$ and $\Omega'\subset U_{J'}$ there
  are a compact set $C'$ and $K\geq0$ with
  $L^{\mathrm{full}}_{C,N}(g\kappa,h\kappa')\leq K\,L^{\mathrm{full}}_{C',N}(g,h)$
  for all $g$, $h$, $\kappa\in\Omega$ and $\kappa'\in\Omega'$.
\end{enumerate}
\end{lemma}
\leanref{Development/WeilBound/ThetaLifts/FrameEstimates/SquareIntegrableWeilRange.lean}{PositivityTheory.ThKEst.LatInv}
\leanref{Development/WeilBound/ThetaLifts/FrameEstimates/LatInv.lean}{PositivityTheory.ThKEst.latInv_of}
\leanref{Development/WeilBound/ThetaLifts/FrameEstimates/SquareIntegrableWeilRange.lean}{PositivityTheory.ThKEst.LatCpt}
\leanref{Development/WeilBound/ThetaLifts/FrameEstimates/LatCpt.lean}{PositivityTheory.ThKEst.latCpt_of}
\begin{proof}
(1) As in the proof of Lemma~\ref{asm-theta-agree-convergence-lattice}(1),
$(\gamma g)^{-1}x((\gamma'h)^{-1})^{\mathsf T}=g^{-1}x'(h^{-1})^{\mathsf T}$
with $x'=\gamma^{-1}x(\gamma'^{-1})^{\mathsf T}$, and $x\mapsto x'$ is a
bijection of $M_{n\times d}(E)$ that preserves the rank.

(2) The proof of Lemma~\ref{asm-theta-agree-convergence-lattice}(3) gives a
compact set $C'$ and $K\geq0$ with
$m_{C,N}(z)\leq K\,m_{C',N}(\kappa z\kappa'^{\mathsf T})$ for all $z$,
$\kappa\in\Omega$ and $\kappa'\in\Omega'$. For
$z=(g\kappa)^{-1}x((h\kappa')^{-1})^{\mathsf T}$ one has
$\kappa z\kappa'^{\mathsf T}=g^{-1}x(h^{-1})^{\mathsf T}$, and summing over
the $x$ of rank $n$ gives (2).
\end{proof}

\begin{lemma}\label{asm-theta-agree-continuity-minor}
Let $B$ be a theta base for a CM field $E$
(Definition~\ref{asm-theta-agree-models}(1)), let $J\in M_n(E)$ and
$J'\in M_d(E)$ be nondegenerate $\epsilon$-Hermitian, let $a$ and $a'$ be
Witt tori of $J$ and $J'$ of sizes $q$ and $q'$
(Definition~\ref{asm-theta-agree-convergence-tori}(2)), and let
$C\subset M_{n\times d}(\mathbf A_E)$ be compact. Then there is $c>0$ such that
\[
c\,e^{\Delta(s,t)/(n+1)}\leq1+\bigl\|a(s)^{-1}x\,(a'(t)^{-1})^{\mathsf T}\bigr\|_\infty
\]
for all $s\in\mathcal C_q$ and $t\in\mathcal C_{q'}$ and every
$x\in M_{n\times d}(E)$ of rank $n$ such that
$a(s)^{-1}x(a'(t)^{-1})^{\mathsf T}-z_0$ is at $\infty$ for some $z_0\in C$,
with the excess $\Delta$ of Definition~\ref{asm-theta-agree-continuity-sums}(2).
\end{lemma}
\leanref{Development/WeilBound/ThetaLifts/FrameEstimates/SquareIntegrableWeilRange.lean}{PositivityTheory.ThKEst.Minor}
\leanref{Development/WeilBound/ThetaLifts/FrameEstimates/ExistsColsDetNeZero.lean}{PositivityTheory.ThKEst.minor_of}
\leanref{Development/WeilBound/ThetaLifts/FrameEstimates/ExistsColsDetNeZero.lean}{PositivityTheory.ThKEst.sum_lam_le_psum}
\begin{proof}
Take bases $b$ of $E^n$ and $b'$ of $E^d$ adapted to $a$ and $a'$, with linear
forms $\lambda_k$ and $\lambda'_l$, and integers $Q\geq1$ and $K\geq1$, as in
the proof of Lemma~\ref{asm-theta-agree-convergence-lattice}(4) and (5). Thus
$y_i$ is the column of $b$ at a position $l_i$, with $\lambda_{l_i}(s)=-s_i$,
the $l_i$ being distinct; the values $\lambda'_l(t)$ are the $t_j$, the $-t_j$
and $0$, each $t_j$ and each $-t_j$ occurring once. Let $s$, $t$ and $x$ be as
in the statement, and put
$z=\|a(s)^{-1}x(a'(t)^{-1})^{\mathsf T}\|_\infty$ and
$\xi=b^{-1}x(b'^{-1})^{\mathsf T}$. Then the entries of $Q\xi$ are algebraic
integers, and
$\sum_w\sum_{k,l}e^{-\lambda_k(s)-\lambda'_l(t)}|\tau_w(\xi_{kl})|\leq Kz$.

Choose $i\leq q$ with $\Delta(s,t)=\sigma_i(s)-\sigma_i(t)$. The matrix $\xi$
has rank $n$, so its rows $l_1,\ldots,l_i$ are linearly independent, and there
are columns $c_1,\ldots,c_i$ such that the minor
$A=(\xi_{l_a,c_b})_{a,b\leq i}$ has $\det A\neq0$. Then $Q^i\det A$ is a
nonzero algebraic integer, so some embedding of $E$ into $\mathbf C$ maps it to
a number of absolute value at least $1$; that embedding has the same absolute
values as $\tau_w$ for an infinite place $w$ of $F$. Let $R$ be the matrix with
entries $R_{ab}=e^{s_a}e^{-\lambda'_{c_b}(t)}\tau_w(A_{ab})$. As
$\lambda_{l_a}(s)=-s_a$, every $|R_{ab}|$ is one of the terms of the sum
above, so $|R_{ab}|\leq K(1+z)$, and $|\det R|\leq i!\,(K(1+z))^i$. On the
other hand,
$|\det R|=\exp\bigl(\sigma_i(s)-\sum_b\lambda'_{c_b}(t)\bigr)|\tau_w(\det A)|$,
and $\sum_b\lambda'_{c_b}(t)\leq\sigma_i(t)$, because the $c_b$ are distinct
and $t_1\geq t_2\geq\cdots\geq0$. Hence
\[
e^{\Delta(s,t)}\leq Q^i\,|\det R|\leq i!\,(QK(1+z))^i\leq\bigl(n!\,QK(1+z)\bigr)^{n+1},
\]
as $i\leq q\leq n$ and $QK(1+z)\geq1$. So $c=(n!\,QK)^{-1}$ works.
\end{proof}

\begin{lemma}\label{asm-theta-agree-continuity-lattice}
Let $B$ be a theta base for a CM field $E$
(Definition~\ref{asm-theta-agree-models}(1)), let $J\in M_n(E)$ and
$J'\in M_d(E)$ be nondegenerate $\epsilon$-Hermitian, and let $a$ and $a'$ be
Witt tori of $J$ and $J'$ of sizes $q$ and $q'$. For every integer $\ell\geq0$
there is $N$ such that for every compact set
$C\subset M_{n\times d}(\mathbf A_E)$ there is $K\geq0$ with
\[
L^{\mathrm{full}}_{C,N}\bigl(a(s),a'(t)\bigr)\leq K\exp\bigl([E:\mathbf Q]\,P_{n,d}(s,t)-\ell\,\Delta(s,t)\bigr)
\]
for all $s\in\mathcal C_q$ and $t\in\mathcal C_{q'}$, with $P_{n,d}$ as in
Definition~\ref{asm-theta-agree-convergence-tori}(5).
\end{lemma}
\leanref{Development/WeilBound/ThetaLifts/FrameEstimates/SquareIntegrableWeilRange.lean}{PositivityTheory.ThKEst.LatFull}
\leanref{Development/WeilBound/ThetaLifts/FrameEstimates/LatFull.lean}{PositivityTheory.ThKEst.latFull_of}
\begin{proof}
Lemma~\ref{asm-theta-agree-convergence-lattice}(5) gives $N_0$ such that for
every compact $C$ there is $K_0$ with
$L_{C,N_0}(a(s),a'(t))\leq K_0\exp([E:\mathbf Q]P_{n,d}(s,t))$. Put
$N=N_0+\ell(n+1)$, let $C$ be compact, and let $c$ be as in
Lemma~\ref{asm-theta-agree-continuity-minor}. Let $x$ have rank $n$, and put
$z=a(s)^{-1}x(a'(t)^{-1})^{\mathsf T}$. If $m_{C,N}(z)\neq0$, then that lemma
gives $(1+\|z\|_\infty)^{\ell(n+1)}\geq c^{\ell(n+1)}e^{\ell\Delta(s,t)}$, so
\[
m_{C,N}(z)\leq c^{-\ell(n+1)}e^{-\ell\Delta(s,t)}\,m_{C,N_0}(z) .
\]
Summing over the $x$ of rank $n$, and bounding by the sum over all $x$, gives
the lemma with $K=c^{-\ell(n+1)}K_0$.
\end{proof}

\begin{lemma}\label{asm-theta-agree-continuity-decay}
Let $(X,W,M,\chi_X,\chi_W,\omega)$ be a Weil pair for a theta base $B$ of a CM
field $E$ (Definition~\ref{asm-theta-agree-triple}(1)), with $X$ and $W$ of
ranks $n$ and $d$, and let $N\geq0$. For $g\in U_X$ let
$\overline{\operatorname{ht}}(g)$ be the infimum of $\operatorname{ht}(\gamma g)$
over $\gamma\in U_X(F)$. Then there is $k\geq0$ such that for every full-rank
finite part $\phi$ for $M$ (Definition~\ref{asm-theta-agree-continuity-sums}(4))
there is a continuous seminorm $p$ on $\mathcal S(F_\infty^{n\times d})$ with
\[
|\theta_\omega(f\otimes\phi)(g,h)|\,\overline{\operatorname{ht}}(g)^N\leq p(f)\operatorname{ht}(h)^k
\]
for all $f\in\mathcal S(F_\infty^{n\times d})$, $g\in U_X$ and $h\in U_W$.
\end{lemma}
\leanref{Development/WeilBound/ThetaLifts/FrameEstimates/HodgeTateFromSandwich.lean}{PositivityTheory.ThKEst.SoSemiU}
\leanref{Development/WeilBound/ThetaLifts/FrameEstimates/WeilConv.lean}{PositivityTheory.ThKEst.soSemiU}
\leanref{Development/WeilBound/ThetaLifts/FrameEstimates/SoSemi.lean}{PositivityTheory.ThKEst.soSemiU_of}
\begin{proof}
Put $D=[E:\mathbf Q]$. By Lemma~\ref{asm-theta-agree-convergence-signature},
$J_X$ and $J_W$ are nondegenerate $\epsilon$-Hermitian. By
Theorem~\ref{asm-theta-agree-convergence-siegel}(2), $J_X$ has a Siegel Witt
torus $a$ of some size $q$, with a compact set $\Omega$ as in property (a) and
constants $C_0\geq0$ and $k_0$ with
$\operatorname{ht}(a(s))\leq C_0e^{k_0\sigma_q(s)}$ for $s\in\mathcal C_q$, by
property (c) (Definition~\ref{asm-theta-agree-convergence-tori}(3)); and $J_W$
has a Siegel Witt torus $a'$ of some size $q'$, with a compact set $\Omega'$
as in property (a). There are $A\geq0$ and $k_1$ with
$e^{\sigma_{q'}(t)}\leq A\operatorname{ht}(h)^{k_1}$ whenever
$\delta h=a'(t)\kappa'$ with $\delta\in U_W(F)$, $t\in\mathcal C_{q'}$ and
$\kappa'\in\Omega'$: then $\delta^{-1}a'(t)=h\kappa'^{-1}$, so this follows
from property (c) of $a'$, as the height is submultiplicative, unchanged by
inversion and bounded on $\Omega'$
(Theorem~\ref{asm-theta-agree-convergence-siegel}(1)). Let $B_0$ bound the
height on $\Omega$. Put $\ell=Dd+2k_0N$, let $N'$ be given by
Lemma~\ref{asm-theta-agree-continuity-lattice} for $a$, $a'$ and $\ell$, and
put $k=k_1(\ell+Dn)$.

Let $\phi$ be a full-rank finite part.
Lemma~\ref{asm-theta-agree-continuity-square}, for $\phi$ and $N'$, gives $p$
and $C$. Lemma~\ref{asm-theta-agree-continuity-full-sums}(2), for $\Omega$,
$\Omega'$ and $C$, gives a compact set $C_1$ and $K_1$, and
Lemma~\ref{asm-theta-agree-continuity-lattice}, for $C_1$, gives $K_2$.

Let $f$, $g$ and $h$ be given, and write $\gamma g=a(s)\kappa$ and
$\delta h=a'(t)\kappa'$ by property (a), with $\gamma\in U_X(F)$,
$\delta\in U_W(F)$, $s\in\mathcal C_q$, $t\in\mathcal C_{q'}$, $\kappa\in\Omega$
and $\kappa'\in\Omega'$. Put $S=\sigma_q(s)$, $T=\sigma_{q'}(t)$ and
$y=DP_{n,d}(s,t)-\ell\Delta(s,t)$. By
Lemma~\ref{asm-theta-agree-continuity-square},
Lemma~\ref{asm-theta-agree-continuity-full-sums}(1) and (2), and
Lemma~\ref{asm-theta-agree-continuity-lattice},
\[
|\theta_\omega(f\otimes\phi)(g,h)|^2\leq p(f)^2L^{\mathrm{full}}_{C,N'}\bigl(a(s)\kappa,a'(t)\kappa'\bigr)
\leq p(f)^2K_1L^{\mathrm{full}}_{C_1,N'}\bigl(a(s),a'(t)\bigr)\leq p(f)^2K_1K_2\,e^{y},
\]
so $|\theta_\omega(f\otimes\phi)(g,h)|\leq p(f)(K_1K_2+1)e^{y/2}$. Next,
$\overline{\operatorname{ht}}(g)\leq\operatorname{ht}(\gamma g)\leq\operatorname{ht}(a(s))\operatorname{ht}(\kappa)\leq C_0B_0e^{k_0S}$.
By Lemma~\ref{asm-theta-agree-convergence-cone}(1),
$P_{n,d}(s,t)\leq dS+nT$. By the definition of the excess,
$\Delta(s,t)\geq\sigma_q(s)-\sigma_q(t)\geq S-T$, as the $t_j$ are nonnegative.
With $\ell=Dd+2k_0N$ these give
\[
\frac y2+Nk_0S\leq\frac{DdS+DnT-\ell\Delta(s,t)}2+Nk_0S=\frac{\ell S+DnT-\ell\Delta(s,t)}2\leq(\ell+Dn)T,
\]
and $e^{(\ell+Dn)T}\leq A^{\ell+Dn}\operatorname{ht}(h)^k$. Hence
\[
|\theta_\omega(f\otimes\phi)(g,h)|\,\overline{\operatorname{ht}}(g)^N
\leq(K_1K_2+1)(C_0B_0)^NA^{\ell+Dn}\,p(f)\operatorname{ht}(h)^k,
\]
and the seminorm $(K_1K_2+1)(C_0B_0)^NA^{\ell+Dn}\,p$ works.
\end{proof}

\begin{theorem}[Borel, Harish-Chandra \cite{BHC62}]\label{asm-theta-agree-continuity-volume}
Let $E$ be a CM field with maximal totally real subfield $F$, let
$J\in M_n(E)$ be nondegenerate $\epsilon$-Hermitian
(Definition~\ref{asm-theta-agree-convergence-tori}(1)), let $\mu$ be a Haar
measure on $U_J$, and let $\mathcal D$ be a fundamental domain for $U_J(F)$ in
$U_J$ with respect to $\mu$. Then $\mu(\mathcal D)<\infty$.
\end{theorem}
\leanref{Development/WeilBound/ThetaLifts/FrameEstimates/SquareIntegrableWeilRange.lean}{PositivityTheory.ThKEst.FinVolU}
\leanref{Development/WeilBound/ThetaLifts/FrameEstimates/FundamentalDomainFiniteVolume.lean}{PositivityTheory.ThKEst.finVolU}
\leanref{CompanionTheory/SpectralGeometry/LocalGlobalProofs/AdelicAnalysis/UnitaryGroupFiniteVolume.lean}{GlobalTheory.CompanionTheory.SW.finVol_unitaryGroup}

\begin{lemma}\label{asm-theta-agree-continuity-domination}
Let $(X,W,M,\chi_X,\chi_W,\omega)$ be a Weil pair for a theta base $B$ of a CM
field $E$ (Definition~\ref{asm-theta-agree-triple}(1)), with $X$ and $W$ of
ranks $n$ and $d$, and let $\mu$ and $\mathcal D$ be the measure and the
domain of $M$. Then there is $k\geq0$ such that for every full-rank finite
part $\phi$ for $M$ (Definition~\ref{asm-theta-agree-continuity-sums}(4))
there are a continuous seminorm $p$ on $\mathcal S(F_\infty^{n\times d})$ and
a measurable function $G\colon U_X\to[0,\infty]$ with
$\int_{\mathcal D}G\,d\mu<\infty$ and
\[
|\theta_\omega(f\otimes\phi)(g,h)|\leq p(f)\operatorname{ht}(h)^kG(g)
\]
for all $f\in\mathcal S(F_\infty^{n\times d})$, $g\in U_X$ and $h\in U_W$.
\end{lemma}
\leanref{Development/WeilBound/ThetaLifts/FrameEstimates/HodgeTateFromSandwich.lean}{PositivityTheory.ThKEst.DomSoU}
\leanref{Development/WeilBound/ThetaLifts/FrameEstimates/WeilConv.lean}{PositivityTheory.ThKEst.domSoU}
\leanref{Development/WeilBound/ThetaLifts/FrameEstimates/SoSemi.lean}{PositivityTheory.ThKEst.domSoU_of}
\begin{proof}
Lemma~\ref{asm-theta-agree-continuity-decay} for $N=0$ gives $k$, and for
every full-rank finite part $\phi$ a continuous seminorm $p$ with
$|\theta_\omega(f\otimes\phi)(g,h)|\leq p(f)\operatorname{ht}(h)^k$ for all $f$,
$g$ and $h$. Take $G=1$. As $M$ is regular, $\mu$ is a Haar measure and
$\mathcal D$ a fundamental domain for $U_X(F)$ in $U_X$ with respect to $\mu$,
and $J_X$ is nondegenerate $\epsilon$-Hermitian by
Lemma~\ref{asm-theta-agree-convergence-signature}. So
$\int_{\mathcal D}G\,d\mu=\mu(\mathcal D)<\infty$ by
Theorem~\ref{asm-theta-agree-continuity-volume}.
\end{proof}

\begin{lemma}\label{asm-theta-agree-continuity-integrals}
Let $(X,W,M,\chi_X,\chi_W,\omega)$ be a Weil pair for a theta base $B$ of a CM
field $E$ (Definition~\ref{asm-theta-agree-triple}(1)), with $X$ and $W$ of
ranks $n$ and $d$, and let $\mu$ and $\mathcal D$ be the measure and the
domain of $M$. Let $\phi_1,\ldots,\phi_m$ be Schwartz-Bruhat functions on
$\mathbf A_f^{n\times d}$, and suppose that for each $i$ there are a continuous
seminorm $p_i$ on $\mathcal S(F_\infty^{n\times d})$, an integer $k_i\geq0$ and
a measurable function $G_i\colon U_X\to[0,\infty]$ with
$\int_{\mathcal D}G_i\,d\mu<\infty$ and
$|\theta_\omega(f\otimes\phi_i)(g,h)|\leq p_i(f)\operatorname{ht}(h)^{k_i}G_i(g)$
for all $f$, $g$ and $h$. Then:
\begin{enumerate}
\item for all $f_1,\ldots,f_m\in\mathcal S(F_\infty^{n\times d})$ the function
  $\Phi=\sum_if_i\otimes\phi_i$ converges
  (Definition~\ref{asm-theta-agree-models}(7)), there are $C$ and $k$ with
  $\int_{\mathcal D}|\theta_\omega(\Phi)(g,h)|\,d\mu(g)\leq C\operatorname{ht}(h)^k$
  for every $h\in U_W$, and $I(\Phi)(h)=\sum_iI(f_i\otimes\phi_i)(h)$;
\item for each $i$ and each $h\in U_W$, the map $f\mapsto I(f\otimes\phi_i)(h)$
  is continuous on $\mathcal S(F_\infty^{n\times d})$.
\end{enumerate}
\end{lemma}
\leanref{Development/WeilBound/ThetaLifts/FrameEstimates/ThickFrames.lean}{PositivityTheory.ThKEst.dom_facts}
\leanref{Development/WeilBound/ThetaLifts/FrameEstimates/ThickFrames.lean}{PositivityTheory.ThKEst.dom_sum}
\begin{proof}
(1) For each $f$, $h$ and $i$, the function
$g\mapsto\theta_\omega(f\otimes\phi_i)(g,h)$ is continuous by
Lemma~\ref{asm-theta-agree-convergence-theta}(2), hence measurable. It is
bounded by $p_i(f)\operatorname{ht}(h)^{k_i}G_i$, which is integrable over
$\mathcal D$. So $f\otimes\phi_i$ converges, and
$\int_{\mathcal D}|\theta_\omega(f\otimes\phi_i)(g,h)|\,d\mu(g)\leq C_i\operatorname{ht}(h)^{k_i}$
with $C_i=p_i(f)\int_{\mathcal D}G_i\,d\mu$. By
Lemma~\ref{asm-theta-agree-convergence-theta}(2),
$\theta_\omega(\Phi)=\sum_i\theta_\omega(f_i\otimes\phi_i)$. So $\Phi$
converges, its theta integral is the sum of those of the $f_i\otimes\phi_i$,
and, as $\operatorname{ht}\geq1$,
$\int_{\mathcal D}|\theta_\omega(\Phi)(g,h)|\,d\mu(g)\leq\sum_iC_i\operatorname{ht}(h)^{k_i}\leq\bigl(\sum_iC_i\bigr)\operatorname{ht}(h)^{k_1+\cdots+k_m}$.

(2) Fix $i$, $h$ and $f_0$. As $p_i$ is continuous, $p_i(f)<p_i(f_0)+1$ for
$f$ near $f_0$, and then
$|\theta_\omega(f\otimes\phi_i)(g,h)|\leq(p_i(f_0)+1)\operatorname{ht}(h)^{k_i}G_i(g)$,
a bound that is integrable over $\mathcal D$ and independent of $f$. For each
$g$, the map $f\mapsto\theta_\omega(f\otimes\phi_i)(g,h)$ is continuous by
Lemma~\ref{asm-theta-agree-convergence-theta}(3). By dominated convergence,
$f\mapsto\int_{\mathcal D}\theta_\omega(f\otimes\phi_i)(g,h)\,d\mu(g)$ is
continuous at $f_0$.
\end{proof}

\begin{lemma}\label{asm-theta-agree-continuity}
Let $(X,W,M,\chi_X,\chi_W,\omega)$ be a Weil pair for a theta base $B$ of a CM
field $E$ (Definition~\ref{asm-theta-agree-triple}(1)), with $X$ and $W$ of
ranks $n$ and $d$. Let $\phi_1,\ldots,\phi_m$ be Schwartz-Bruhat functions on
$\mathbf A_f^{n\times d}$, and $(u_{k,j})_k$ sequences in
$\mathcal S(F_\infty^{n\times d})$ with limits $f_j$, such that
$\Psi_k=\sum_ju_{k,j}\otimes\phi_j$ and $\Psi=\sum_jf_j\otimes\phi_j$ lie in
$\mathcal S^\circ(M)$ for every $k$. Then $I(\Psi_k)(h)\to I(\Psi)(h)$ for every
$h\in U_W$.
\end{lemma}
\leanref{Development/AbelianVarietyTheory/Definitions/Divisors/CTSmallCone.lean}{GlobalTheory.ThK.SoCont}
\leanref{Development/WeilBound/ThetaLifts/FrameEstimates/WeilConv.lean}{PositivityTheory.ThKEst.soCont}
\leanref{Development/WeilBound/ThetaLifts/FrameEstimates/ThickFrames.lean}{PositivityTheory.ThKEst.soCont_of}
\begin{proof}
By Lemma~\ref{asm-theta-agree-continuity-domination}, finitely many full-rank
finite parts for $M$ satisfy the hypothesis of
Lemma~\ref{asm-theta-agree-continuity-integrals}.

Let $V\subset\mathbf C^m$ be the subspace of the vectors $c$ such that
$f\otimes\sum_jc_j\phi_j\in\mathcal S^\circ(M)$ for every
$f\in\mathcal S(F_\infty^{n\times d})$. Choose a linear projection $P$ of
$\mathbf C^m$ onto $V$, and put $\psi_j=\sum_iP(e_j)_i\,\phi_i$, where $e_j$ is
the $j$-th standard basis vector. Each $\psi_j$ is a Schwartz-Bruhat function,
and a full-rank finite part (Definition~\ref{asm-theta-agree-continuity-sums}(4)),
as $P(e_j)\in V$.

Let $w_1,\ldots,w_m\in\mathcal S(F_\infty^{n\times d})$ with
$\sum_jw_j\otimes\phi_j\in\mathcal S^\circ(M)$. Then
$\sum_jw_j\otimes\phi_j=\sum_jw_j\otimes\psi_j$. Indeed, by
Lemma~\ref{asm-theta-agree-continuity-pieces}(1),
$\sum_jw_j\otimes\phi_j=\sum_if'_i\otimes\chi_i$ with full-rank finite parts
$\chi_i$. Fix $x\in F_\infty^{n\times d}$. Comparing values at the points with
archimedean component $x$ gives $\sum_jw_j(x)\phi_j=\sum_if'_i(x)\chi_i$. So
for every $f_0$,
$f_0\otimes\sum_jw_j(x)\phi_j=\sum_if'_i(x)\,f_0\otimes\chi_i\in\mathcal S^\circ(M)$,
and the vector $w(x)=(w_j(x))_j$ lies in $V$. Hence
$w(x)=Pw(x)=\sum_jw_j(x)P(e_j)$, and
$\sum_iw_i(x)\phi_i=\sum_jw_j(x)\psi_j$. As $x$ was arbitrary, the two
functions agree.

So $\Psi=\sum_jf_j\otimes\psi_j$ and $\Psi_k=\sum_ju_{k,j}\otimes\psi_j$. By
Lemma~\ref{asm-theta-agree-continuity-integrals}(1) for $\psi_1,\ldots,\psi_m$,
$I(\Psi)(h)=\sum_jI(f_j\otimes\psi_j)(h)$ and
$I(\Psi_k)(h)=\sum_jI(u_{k,j}\otimes\psi_j)(h)$. By
Lemma~\ref{asm-theta-agree-continuity-integrals}(2), each
$f\mapsto I(f\otimes\psi_j)(h)$ is continuous on
$\mathcal S(F_\infty^{n\times d})$, and $u_{k,j}\to f_j$. Hence
$I(\Psi_k)(h)\to I(\Psi)(h)$.
\end{proof}
% Lean: soCont hSiegelSet thetaBasic = soCont_of thetaBasic (domSoU hSiegelSet) soFin_of;
% the subspace V is soCoeff and P is Submodule.linearProjOfIsCompl.

\begin{lemma}\label{asm-theta-agree-density-moments}
Let $V$ be a finite-dimensional real vector space, and let $Q$ and $Q'$ be
$\mathbf R$-bilinear forms on $V$ with $Q(x,x)>0$ for $x\neq0$. Let
$c_0\in\mathbf C^\times$, and let $f_0(x)=c_0\exp(-Q(x,x)-iQ'(x,x))$, a
Schwartz function on $V$. Let $N$ be a subspace of $\mathcal S(V)$ that
contains $f_0$, and contains $x\mapsto l(x)f(x)$ whenever $f\in N$ and $l$ is
an $\mathbf R$-linear form on $V$. Then every continuous linear form on
$\mathcal S(V)$ that vanishes on $N$ is zero.
\end{lemma}
\leanref{Development/WeilBound/ThetaLifts/LatticeBasics/SchwartzFunctionalEqZeroOfDense.lean}{GlobalTheory.ThK.Dbl.Dense.schwartz_functional_eq_zero_of_dense}
\leanref{Development/WeilBound/ThetaLifts/SchwartzMoments/EqZeroOfMoments.lean}{PositivityTheory.ThKOwn.eq_zero_of_moments}
\leanref{Development/WeilBound/ThetaLifts/SchwartzMoments/MulLPowApply.lean}{PositivityTheory.ThKOwn.kill_exp}
\leanref{Development/WeilBound/ThetaLifts/LocalTheory/ArchimedeanBumps/DictIrr.lean}{PositivityTheory.ThKG.Loc.Bump.kill_mul}
\begin{proof}
Let $T$ be such a form, and fix a Euclidean norm on $V$.
\begin{enumerate}
\item $N$ contains $Pf_0$ for every polynomial function $P\colon V\to\mathbf C$,
  by induction on $P$: constants are clear, sums are clear, and multiplying by
  a coordinate is multiplying by a linear form.
\item Multiplication by $e(x)=c_0\exp(-iQ'(x,x))$ is a continuous linear
  bijection of $\mathcal S(V)$, with inverse the multiplication by $e^{-1}$,
  since $e$, $e^{-1}$ and all their derivatives grow at most polynomially. Put
  $g(x)=\exp(-Q(x,x))$, so that $f_0=eg$, and $T_1(f)=T(ef)$. Then $T_1$ is a
  continuous linear form, $T_1(Pg)=T(Pf_0)=0$ for every polynomial $P$ by (1),
  and $T=0$ once $T_1=0$.
\item By compactness of the unit sphere there is $w>0$ with
  $Q(x,x)\geq w\|x\|^2$. Let $L$ be an $\mathbf R$-linear form on $V$. For
  $z\in\mathbf C$ the function $G_z(x)=\exp(zL(x))g(x)$ is a Schwartz
  function, as $\operatorname{Re}(zL(x)-Q(x,x))\leq C_z-\tfrac w2\|x\|^2$. The
  derivatives of $(z,x)\mapsto\exp(zL(x)-Q(x,x))$ satisfy Schwartz bounds in
  $x$ uniformly for $z$ in discs, so $z\mapsto T_1(G_z)$ is holomorphic on
  $\mathbf C$ with derivative $z\mapsto T_1(LG_z)$. Its $m$-th derivative at
  $0$ is $T_1(L^mg)=0$. An entire function is the sum of its Taylor series at
  $0$, so $T_1(G_z)=0$ for every $z$, in particular for $z=i$.
\item Let $\phi\in\mathcal S(V)$. By Fourier inversion,
  $\phi(x)=\int_Va(\eta)\exp(2\pi i\langle\eta,x\rangle)\,d\eta$ for a Schwartz
  function $a$. The map $\eta\mapsto a(\eta)\exp(2\pi i\langle\eta,\cdot\rangle)g$
  into $\mathcal S(V)$ is continuous, and each Schwartz seminorm of it is at
  most $|a(\eta)|$ times a polynomial in $\|\eta\|$, so it is integrable. Hence
  there is $F\in\mathcal S(V)$ with
  $\lambda(F)=\int_Va(\eta)\lambda(\exp(2\pi i\langle\eta,\cdot\rangle)g)\,d\eta$
  for every continuous linear form $\lambda$ on $\mathcal S(V)$: the form
  $\lambda$ factors through finitely many of the maps
  $f\mapsto\|x\|^kD^jf(x)$ to bounded continuous functions, where the integral
  converges. Evaluation at points gives $F=\phi g$, and $\lambda=T_1$ gives
  $T_1(\phi g)=0$ by (3).
\item A smooth function $\psi$ with compact support is $(\psi/g)g$ with
  $\psi/g\in\mathcal S(V)$, so $T_1(\psi)=0$ by (4). These functions are dense
  in $\mathcal S(V)$, and $T_1$ is continuous, so $T_1=0$.
\end{enumerate}
\end{proof}
% Lean: (1) poly_mem; (2) exp_tg, quad_tg and the maps M, M', A, B of
% schwartz_functional_eq_zero_of_dense; (3) lower_bound, then eq_zero_PG and
% eq_zero_of_moments with kill_exp (exists_family, family_bound, family_deriv,
% Complex.taylorSeries_eq_of_entire'); (4) kill_mul with exists_density, semi_mul,
% cont_mod and weakInt (JetApply.lean: SchI.factor, SchI.exists_int); (5) dense_compactSupport.

\begin{lemma}\label{asm-theta-agree-density-derivative}
Let $B$ be a theta base for a CM field $E$, let $X$ and $W$ be lattice spaces
of ranks $n$ and $d$, let $M$ be a model for $(X,W)$
(Definition~\ref{asm-theta-agree-models}(3)), and put
$I=\{1,\ldots,n\}\times\{1,\ldots,d\}$ and $V=F_\infty^I$. For $x,b\in V$ let
$T(x,b)=\operatorname{Tr}_{F_\infty/\mathbf R}\bigl(\sum_{i\in I}x_ib_i\bigr)$,
and for $a,b\in V$ and $f\in\mathcal S(V)$ let
$D_{a,b}f(x)=\partial_af(x)-2\pi i\,T(x,b)f(x)$, where $\partial_a$ is the
derivative in the direction $a$. Let $u\in\mathbb W(X,W)$ be at $\infty$, and
let $a$ and $b$ be the components at $\infty$ of $e_1(u)$ and $e_2(u)$. Then for
every $f\in\mathcal S(V)$ the family $t\mapsto\rho(tu)$ acts at $\infty$ to
first order by $f\mapsto D_{a,b}f$; and if it acts at $\infty$ to first order
by $f\mapsto f_1$, then $f_1=D_{a,b}f$.
\end{lemma}
\leanref{Development/WeilBound/ThetaLifts/LatticeBasics/ExistsRealStructure.lean}{GlobalTheory.ThK.Dbl.Dense.hasDerivAt_dOp}
\leanref{Development/WeilBound/ThetaLifts/LatticeBasics/ExistsRealStructure.lean}{GlobalTheory.ThK.Dbl.Dense.actsD_unique}
\begin{proof}
As $e$ is $\mathbf A$-linear and $u$ is at $\infty$, $e_1(u)$ and $e_2(u)$ are
at $\infty$. Let $r$ be a function on $\mathbf A^I$ that does not depend on the
component at $\infty$, with $f\otimes r$ Schwartz-Bruhat. For $x\in\mathbf A^I$,
the argument $x\cdot te_2(u)+\tfrac12t^2e_1(u)\cdot e_2(u)$ of $\psi$ is at
$\infty$, where $\psi$ is $z\mapsto\exp(-2\pi i\operatorname{Tr}_{F_\infty/\mathbf R}z_\infty)$,
and $r(x+te_1(u))=r(x)$. So
\[
\rho(tu)(f\otimes r)(x)=\exp\bigl(-2\pi i\bigl(t\,T(x_\infty,b)+\tfrac12t^2\operatorname{Tr}_{F_\infty/\mathbf R}(a\cdot b)\bigr)\bigr)f(x_\infty+ta)\,r(x),
\]
whose derivative at $t=0$ is $D_{a,b}f(x_\infty)r(x)$. This is the first claim.
For the second, take for $r$ the product over the finite places $v$ of the
indicator functions of $\mathcal O_v^I$, and $x$ with finite components $0$, so
that $r(x)=1$. Uniqueness of derivatives gives $f_1(x_\infty)=D_{a,b}f(x_\infty)$.
\end{proof}

\begin{lemma}\label{asm-theta-agree-density-gaussian}
Let $B$ be a theta base for a CM field $E$, let $X$ be a Hermitian and $W$ a
skew-Hermitian lattice space, of ranks $n$ and $d$ and of some signatures, and
let $M$ be a model for $(X,W)$. Put $I=\{1,\ldots,n\}\times\{1,\ldots,d\}$.
Then there are $\mathbf R$-bilinear forms $Q$ and $Q'$ on $F_\infty^I$, with
$Q(x,x)>0$ for $x\neq0$, such that $f_0(x)=\exp(-Q(x,x)-iQ'(x,x))$ is a
Gaussian for $M$ (Definition~\ref{asm-theta-agree-models}(4)), with all signs
$\varepsilon_w=1$.
\end{lemma}
\leanref{Development/WeilBound/ThetaLifts/LatticeBasics/ExistsIsGauss.lean}{GlobalTheory.ThK.Dbl.Dense.exists_isGauss}
\leanref{Development/WeilBound/ThetaLifts/LatticeBasics/ExistsRealStructure.lean}{GlobalTheory.ThK.Dbl.Dense.exists_real_structure}
\leanref{Development/WeilBound/ThetaLifts/LatticeBasics/DenseG.lean}{GlobalTheory.ThK.Dbl.Dense.exists_vacuum}
\leanref{Development/WeilBound/ThetaLifts/LatticeBasics/DenseM.lean}{GlobalTheory.ThK.Dbl.Dense.hf_neg}
\begin{proof}
Write $V=F_\infty^I$, and use $T$ and $D_{a,b}$ of
Lemma~\ref{asm-theta-agree-density-derivative}. For $p=(p_1,p_2)$ and
$q=(q_1,q_2)$ in $V\times V$ put $D_p=D_{p_1,p_2}$ and
$\Omega(p,q)=T(p_1,q_2)-T(q_1,p_2)$.

\emph{Coordinates.} The isomorphism $e$ is $\mathbf A$-linear, so it maps the
elements of $\mathbb W(X,W)$ at $\infty$ bijectively onto the elements of
$\mathbf A^I\times\mathbf A^I$ at $\infty$, which we identify with $V\times V$;
let $s$ be the inverse map. The elements of $\mathbf A_E$ at $\infty$ form
$E\otimes_FF_\infty=\prod_wE\otimes_FF_w$, and $\tau_w$ identifies
$E\otimes_FF_w$ with $\mathbf C$. So $p\mapsto(\pi_w(s(p)))_w$ is an
$\mathbf R$-linear bijection from $V\times V$ onto $\prod_wM_{n\times d}(\mathbf C)$,
and there is a unique $\mathbf R$-linear map $J$ of $V\times V$ with
$\pi_w(s(Jp))=\kappa_w(\pi_w(s(p)))$ for every $w$, where
$\kappa_w(x)=i\,\theta_{X,w}\,x\,\theta_{W,w}^{\mathsf T}$.

\emph{The form.} For $x,y\in M_{n\times d}(\mathbf C)$ put
$H_w(x,y)=\operatorname{tr}\bigl(x^{\mathsf T}\tau_w(J_X)\,\bar y\,\tau_w(J_W)^{\mathsf T}\bigr)$.
The condition on $e$ gives
$\Omega(p,q)=\operatorname{Tr}_{F_\infty/\mathbf R}\langle s(p),s(q)\rangle$.
Since $E$ is a CM field, $\tau_w\circ c$ is the complex conjugate of $\tau_w$,
so $\pi_w(h(s(p),s(q)))=H_w(\pi_ws(p),\pi_ws(q))$; and for $z\in\mathbf A_E$ at
$\infty$, $\operatorname{Tr}_{F_\infty/\mathbf R}\operatorname{Tr}_{\mathbf A_E/\mathbf A}z
=\sum_w2\operatorname{Re}\pi_w(z)$. As
$\langle x,y\rangle=-\operatorname{Tr}_{\mathbf A_E/\mathbf A}h(x,y)$,
\[
\Omega(p,q)=-2\sum_w\operatorname{Re}H_w\bigl(\pi_ws(p),\pi_ws(q)\bigr).
\]

\emph{Cartan involutions.} Fix $w$ and write $A=\tau_w(J_X)$,
$C=\tau_w(J_W)$, $\theta=\theta_{X,w}$, $\theta'=\theta_{W,w}$, $\kappa=\kappa_w$
and $H=H_w$. The matrices $A$ and $iC$ are Hermitian, as $X$ is Hermitian and
$W$ skew-Hermitian, and $\theta$ and $\theta'$ are Cartan involutions of $A$
and $iC$. For $x,y\in M_{n\times d}(\mathbf C)$:
\begin{enumerate}
\item[(a)] $H(\kappa x,\kappa y)=H(x,y)$. Indeed
  $\overline{\kappa y}=-i\bar\theta\bar y\bar\theta'^{\mathsf T}$, so
  $H(\kappa x,\kappa y)=\operatorname{tr}(\theta'x^{\mathsf T}\theta^{\mathsf T}A\bar\theta\,\bar y\,\bar\theta'^{\mathsf T}C^{\mathsf T})$;
  now use $\theta^{\mathsf T}A\bar\theta=A$, the transpose
  $\bar\theta'^{\mathsf T}C^{\mathsf T}\theta'=C^{\mathsf T}$ of
  $\theta'^{\mathsf T}C\bar\theta'=C$, and $\operatorname{tr}(YZ)=\operatorname{tr}(ZY)$.
\item[(b)] $H(\kappa x,y)=-H(x,\kappa y)$. As $\theta^2=\theta'^2=1$,
  $\kappa(\kappa x)=-x$, so by (a)
  $H(\kappa x,y)=H(\kappa x,\kappa\kappa(-y))=H(x,\kappa(-y))=-H(x,\kappa y)$.
\item[(c)] $\operatorname{Re}H(x,\kappa x)<0$ if $x\neq0$. If $\vartheta$ is a
  Cartan involution of a Hermitian matrix $N$, then $N\bar\vartheta$ is
  Hermitian, since $\vartheta^{\mathsf T}N=N\bar\vartheta$, and positive
  definite, since $z^*N\bar\vartheta z=(\vartheta y)^{\mathsf T}N\bar y$ with
  $y=\vartheta\bar z$. So $P=A\bar\theta$ and $P'=iC\bar\theta'$ are positive
  definite Hermitian, and
  $H(x,\kappa x)=-\operatorname{tr}(x^{\mathsf T}P\bar xP'^{\mathsf T})=-\xi^*(P\otimes P')\xi$,
  where $\xi$ is the vector of entries of $\bar x$.
\end{enumerate}
Hence $\Omega(Jp,q)=-\Omega(p,Jq)$ and $\Omega(Jp,Jq)=\Omega(p,q)$ for all $p$
and $q$, and
$\Omega(p,Jp)=-2\sum_w\operatorname{Re}H_w(x_w,\kappa_wx_w)>0$ for $p\neq0$,
where $x_w=\pi_ws(p)$, as some $x_w$ is nonzero.

\emph{The Gaussian.} Let $K(b)$ be the first component of $J(0,b)$. Then
$\Omega((0,b),J(0,b))=-T(K(b),b)$, which is positive for $b\neq0$; so $K$ is
injective, hence invertible. Put $c=-K^{-1}$, so that $J(0,cx)=(-x,c'x)$ for a
linear map $c'$, and let $Q(x,y)=\pi\,T(y,cx)$ and $Q'(x,y)=\pi\,T(y,c'x)$. The
relations for $J$ applied to $(0,cx)$ and $(0,cy)$ give $T(x,cy)=T(y,cx)$,
$T(x,c'y)=T(y,c'x)$ and $T(x,cx)>0$ for $x\neq0$. So $Q$ and $Q'$ are symmetric
and $Q$ is positive definite. Then $f_0(x)=\exp(-Q(x,x)-iQ'(x,x))$ is a
Schwartz function, with $\partial_af_0(x)=-2(Q(x,a)+iQ'(x,a))f_0(x)$. The
relations for $J$ applied to $p$ and $(0,cx)$ give
$T((Jp)_1,cx)=-T(p_1,c'x)-T(x,p_2)$ and
$T((Jp)_1,c'x)+T(x,(Jp)_2)=T(p_1,cx)$. Hence
\[
\frac{D_{Jp}f_0(x)}{f_0(x)}=-2\pi T((Jp)_1,cx)-2\pi i\bigl(T((Jp)_1,c'x)+T(x,(Jp)_2)\bigr)
=2\pi\bigl(T(p_1,c'x)+T(x,p_2)\bigr)-2\pi i\,T(p_1,cx),
\]
which is $i\,D_pf_0(x)/f_0(x)$. So $D_{Jp}f_0=i\,D_pf_0$ for every $p$.

\emph{Conclusion.} Let $u$ and $u'$ be at $\infty$ with
$\pi_w(u')=i\,\theta_{X,w}\pi_w(u)\theta_{W,w}^{\mathsf T}$ for every $w$, and
let $t\mapsto\rho(tu)$ and $t\mapsto\rho(tu')$ act at $\infty$ to first order by
$f_0\mapsto f$ and $f_0\mapsto f'$. Let $p$ be the coordinates of $u$, so
$u=s(p)$. By the bijectivity in the first paragraph, $u'=s(Jp)$. By
Lemma~\ref{asm-theta-agree-density-derivative}, $f=D_pf_0$ and
$f'=D_{Jp}f_0=if$. So $f_0$ is a Gaussian for $M$, with $\varepsilon_w=1$.
\end{proof}
% Lean: the coordinates are exists_real_structure (with eq_zero_of_piInf and
% exists_adele_piInf); the form is trace_symp with hTens_tr; (a), (b), (c) are hf_cs_cs,
% hf_cs and hf_neg, with posDef_cartan; the Gaussian is exists_vacuum (with exists_gauss and
% fderiv_gauss); the conclusion uses actsD_unique twice.

\begin{lemma}\label{asm-theta-agree-density-linear}
Let $B$ be a theta base for a CM field $E$, let $X$ and $W$ be lattice spaces
of ranks $n$ and $d$, and let $M$ be a model for $(X,W)$. Put
$I=\{1,\ldots,n\}\times\{1,\ldots,d\}$. If $f\in\operatorname{Pol}(M)$ and $l$
is an $\mathbf R$-linear form on $F_\infty^I$, then $x\mapsto l(x)f(x)$ lies in
$\operatorname{Pol}(M)$.
\end{lemma}
\leanref{Development/WeilBound/ThetaLifts/LatticeBasics/ExistsRealStructure.lean}{GlobalTheory.ThK.Dbl.Dense.fock_mul_linear}
\begin{proof}
The trace form $(y,b)\mapsto\operatorname{Tr}_{F_\infty/\mathbf R}(yb)$ on
$F_\infty=F\otimes\mathbf R$ is nondegenerate, as the trace form of $F$ over
$\mathbf Q$ is. So there is $b\in F_\infty^I$ with $l(x)=T(x,b)$ for all $x$, in
the notation of Lemma~\ref{asm-theta-agree-density-derivative}. Let
$u\in\mathbb W(X,W)$ with $e(u)=(0,b)$, where $b$ is viewed as an element of
$\mathbf A^I$ at $\infty$; then $u$ is at $\infty$. By
Lemma~\ref{asm-theta-agree-density-derivative}, $t\mapsto\rho(tu)$ acts at
$\infty$ to first order by $f\mapsto D_{0,b}f=-2\pi i\,lf$. As
$\operatorname{Pol}(M)$ contains $f$ and is closed under these actions, it
contains $-2\pi i\,lf$, hence $lf$.
\end{proof}
% Lean: exists_linForm_eq (exists_trace_eq, trace_mul_eq_zero), reAd, phase_line and
% hasDerivAt_line.

\begin{lemma}\label{asm-theta-agree-density}
Let $(X,W,M,\chi_X,\chi_W,\omega)$ be a Weil pair for a theta base $B$ of a CM
field $E$ (Definition~\ref{asm-theta-agree-triple}(1)), with $X$ and $W$ of
ranks $n$ and $d$. Then $\operatorname{Pol}(M)$ is dense in
$\mathcal S(F_\infty^{n\times d})$.
\end{lemma}
\leanref{Development/WeilBound/ThetaLifts/ConeHypotheses/PolDense.lean}{PositivityTheory.ThKG.PolDense}
\leanref{Development/WeilBound/ThetaLifts/LatticeBasics/PolDense.lean}{GlobalTheory.ThK.Dbl.polDense}
\leanref{Development/WeilBound/ThetaLifts/LatticeBasics/ExistsIsGauss.lean}{GlobalTheory.ThK.Dbl.Dense.pol_dense}
\begin{proof}
Suppose not. By the Hahn-Banach theorem in the locally convex space
$\mathcal S(F_\infty^{n\times d})$, there is a nonzero continuous linear form
$\lambda$ on it that vanishes on $\operatorname{Pol}(M)$. By
Lemma~\ref{asm-theta-agree-density-gaussian}, applied to the Hermitian space
$X$ and the skew-Hermitian space $W$, there are $Q$ and $Q'$ with $Q(x,x)>0$
for $x\neq0$ such that $f_0(x)=\exp(-Q(x,x)-iQ'(x,x))$ is a Gaussian for $M$;
so $f_0\in\operatorname{Pol}(M)$. By
Lemma~\ref{asm-theta-agree-density-linear}, $\operatorname{Pol}(M)$ contains
$lf$ whenever it contains $f$ and $l$ is a real linear form. So
Lemma~\ref{asm-theta-agree-density-moments}, with $N=\operatorname{Pol}(M)$
and $c_0=1$, gives $\lambda=0$, which is absurd.
\end{proof}
% Lean: polDense, by RCLike.geometric_hahn_banach_closed_point and Dense.pol_dense, which
% is schwartz_functional_eq_zero_of_dense applied to exists_isGauss and fock_mul_linear.

\subsection{Evaluation maps along adapted planes}\label{subsection-asm-evaluation-maps}

\begin{lemma}\label{asm-theta-agree-partners}
Let $(B,G,\mathcal X,\omega,\omega',R,R_f)$ be a matched quadruple for a theta
base $B$ of a CM field $E$ (Definition~\ref{asm-theta-agree-matched}). Then:
\begin{enumerate}
\item for every $\Phi\in\operatorname{PD}(M)$ there is
  $\Phi'\in\mathcal S^\circ(M')$ corresponding to $\Phi$;
\item for every $\Phi'\in\operatorname{PD}(M')$ there is a $\Phi$
  corresponding to $\Phi'$;
\item if $\Phi'\in\mathcal S^\circ(M')$ and $\Phi$ and $\Phi'$ correspond, then
  $\Phi\in\operatorname{PD}(M)$, $\Phi'\in\operatorname{PD}(M')$, and $\Phi$ and
  $\Phi'$ match.
\end{enumerate}
\end{lemma}
\leanref{Development/AbelianVarietyTheory/Definitions/Divisors/CTSmallCone.lean}{GlobalTheory.ThK.MatchEx}
\leanref{Development/WeilBound/ThetaLifts/WeilDomination/MatchEx.lean}{PositivityTheory.ThKClose.matchEx_of}
\begin{proof}
(1) Write $\Phi=\sum_if_i\otimes\bigotimes_v\phi_{i,v}$ with
$f_i\in\operatorname{Pol}(M)$, the $\phi_{i,v}$ Schwartz-Bruhat and
$\phi_{i,v}=\mathrm{vac}_v$ outside a finite set. Let $T_1$ be the union of
$T$, $\{v_1\}$ and the places where some $\phi_{i,v}\neq\mathrm{vac}_v$. Since
$p'(\operatorname{Pol}(M'))=\Pi$, there are $f'_i\in\operatorname{Pol}(M')$
with $p'(f'_i)=p(f_i)$. Since $p'_v$ maps the Schwartz-Bruhat functions onto
$\Pi_v$, there are Schwartz-Bruhat $\chi_{i,v}$ with
$p'_v(\chi_{i,v})=p_v(\phi_{i,v})$. The place $v_1$ is clean for
$\mathcal X$, so by Lemma~\ref{asm-theta-agree-coinvariants} there are
$\zeta_i$ of full rank at $\tau_1$ with
$p'_{v_1}(\zeta_i)=p_{v_1}(\phi_{i,v_1})$. Put $\phi'_{i,v_1}=\zeta_i$,
$\phi'_{i,v}=\chi_{i,v}$ for $v\in T_1\setminus\{v_1\}$ and
$\phi'_{i,v}=\mathrm{vac}'_v$ for $v\notin T_1$, and
$\Phi'=\sum_if'_i\otimes\bigotimes_v\phi'_{i,v}$. Then
$\Phi'\in\mathcal S^\circ(M')$, and the two sums of tensors over $T_1$ agree
term by term, so $\Phi$ and $\Phi'$ correspond.

(2) Write $\Phi'=\sum_if'_i\otimes\bigotimes_v\phi'_{i,v}$ in the same way,
let $T_1$ be the union of $T$ and the places where some
$\phi'_{i,v}\neq\mathrm{vac}'_v$, choose $f_i\in\operatorname{Pol}(M)$ with
$p(f_i)=p'(f'_i)$ and Schwartz-Bruhat $\chi_{i,v}$ with
$p_v(\chi_{i,v})=p'_v(\phi'_{i,v})$, and put $\phi_{i,v}=\chi_{i,v}$ for
$v\in T_1$ and $\phi_{i,v}=\mathrm{vac}_v$ otherwise. Then
$\Phi=\sum_if_i\otimes\bigotimes_v\phi_{i,v}$ corresponds to $\Phi'$.

(3) Let $T_1$ and the sums $\Phi=\sum_if_i\otimes\bigotimes_v\phi_{i,v}$ and
$\Phi'=\sum_jf'_j\otimes\bigotimes_v\phi'_{j,v}$ be as in the definition of
correspondence. These sums show $\Phi\in\operatorname{PD}(M)$ and
$\Phi'\in\operatorname{PD}(M')$. Since $M$ and $M'$ are regular,
$\phi_{i,v}$ and $\phi'_{j,v}$ are the indicator functions of the integral
vectors for almost all $v$, so $g_i=\bigotimes_v\phi_{i,v}$ and
$g'_j=\bigotimes_v\phi'_{j,v}$ are Schwartz-Bruhat functions on the finite
adeles, and $\Phi=\sum_if_i\otimes g_i$, $\Phi'=\sum_jf'_j\otimes g'_j$. The
constant sequences $s_{k,i}=f_i$ and $s'_{k,j}=f'_j$ show that $\Phi$ and
$\Phi'$ match.
\end{proof}

\begin{lemma}\label{asm-theta-agree-levi}
Let $E$ be a CM field with a theta base $B$
(Definition~\ref{asm-theta-agree-models}), let $W$ be a skew-Hermitian lattice
space of some signature, and let $H=(x,y,\iota,\mathfrak b)$ be an adapted
plane in $W$ over a lattice space $W_1$
(Definition~\ref{asm-theta-agree-planes}(1)). Then for every
$t\in\mathbf A_E^\times$ and $h_1\in U_{W_1}$ there is a Levi element
$m\in U_W$ for $(t,h_1)$; in particular there is one for $(1,h_1)$.
\end{lemma}
\leanref{Development/WeilBound/ThetaLifts/ConeHypotheses/PolDense.lean}{PositivityTheory.ThKG.LeviEx}
\leanref{Development/WeilBound/ThetaLifts/TowerModels/LeviAd.lean}{PositivityTheory.ThKG.Tow.Levi.exists_levi}
\leanref{Development/WeilBound/ThetaLifts/TowerModels/LeviAlg.lean}{PositivityTheory.ThKG.Tow.Levi.IsFr.levi}
\begin{proof}
Let $J$ and $J_1$ be the matrices of $W$ and $W_1$, of sizes $d$ and
$d_1=d-2$, and $h(u,u')=u^{\mathsf T}J\,c(u')$. As $W$ is skew-Hermitian,
$h(u',u)=-c(h(u,u'))$. Let $A=(\iota\ x\ y)\in M_d(E)$ be the matrix with the
columns of $\iota$ followed by $x$ and $y$.

$A$ is invertible. Every $z\in L$ is $sx+\iota u+ty$, so the column space of
$A$ contains $L$, which spans $E^d$; as $A$ is square, it is invertible.

The Gram matrix of $A$ is
$A^{\mathsf T}J\,c(A)=J^\natural=\operatorname{diag}\bigl(J_1,\left(\begin{smallmatrix}0&1\\-1&0\end{smallmatrix}\right)\bigr)$.
Indeed $\iota^{\mathsf T}J\,c(\iota)=J_1$; $h(\iota u,x)=h(\iota u,y)=0$, and
hence $h(x,\iota u)=h(y,\iota u)=0$; $h(x,x)=h(y,y)=0$, $h(x,y)=1$ and
$h(y,x)=-c(1)=-1$.

Put $D=\operatorname{diag}(h_1,t,c(t)^{-1})$ and $m=ADA^{-1}$, a matrix over
$\mathbf A_E$ with inverse $AD^{-1}A^{-1}$. Reading off columns, $mx=tx$,
$my=c(t)^{-1}y$ and $m\iota=\iota h_1$. Since $h_1^{\mathsf T}J_1c(h_1)=J_1$
and
$\left(\begin{smallmatrix}t&0\\0&c(t)^{-1}\end{smallmatrix}\right)
\left(\begin{smallmatrix}0&1\\-1&0\end{smallmatrix}\right)
\left(\begin{smallmatrix}c(t)&0\\0&t^{-1}\end{smallmatrix}\right)
=\left(\begin{smallmatrix}0&1\\-1&0\end{smallmatrix}\right)$,
we have $D^{\mathsf T}J^\natural c(D)=J^\natural$. With
$J=(A^{-1})^{\mathsf T}J^\natural c(A^{-1})$ this gives
$m^{\mathsf T}J\,c(m)=(A^{-1})^{\mathsf T}D^{\mathsf T}J^\natural c(D)\,c(A^{-1})=J$.
So $m\in U_W$, and it is a Levi element for $(t,h_1)$.
\end{proof}

\begin{lemma}\label{asm-theta-agree-stabilizer}
Let $E$ be a CM field with a theta base $B$
(Definition~\ref{asm-theta-agree-models}), and let $X$ be a Hermitian lattice
space of some signature. Then the lattice stabilizer $K$ of $X$
(Definition~\ref{asm-theta-agree-planes}(6)) is a compact subgroup of $U_X$,
and there are Haar measures on $K$ and on $\mathbf A_E^\times$.
\end{lemma}
\leanref{Development/AbelianVarietyTheory/Definitions/Divisors/CTSmallCone.lean}{GlobalTheory.ThK.StabEx}
\leanref{Development/WeilBound/ThetaLifts/TowerModels/StabEx.lean}{PositivityTheory.ThKG.Tow.stabEx}
\leanref{Development/WeilBound/ThetaLifts/TowerModels/StabCpt.lean}{PositivityTheory.ThKG.Tow.Stab.isCompact_stabK}
\begin{proof}
Let $K$ be the set of $g\in U_X$ that commute with every $\theta_{X,w}$ at $w$
and satisfy $g\hat L\subset\hat L$ and $g^{-1}\hat L\subset\hat L$, where
$\hat L=\hat L_X$. It contains $1$ and is closed under products; it is closed
under inverses, since an inverse of a matrix commuting with $\theta_{X,w}$
commutes with it and the lattice condition is symmetric in $g$ and $g^{-1}$.
So $K$ is a subgroup, and it is the lattice stabilizer.

The group $\hat L$ is compact: if $s_1,\ldots,s_r$ generate $L_X$ over
$\mathbf Z$, then $\hat L$ is the image of $\hat{\mathcal O}^r$, with
$\hat{\mathcal O}=\{a\in\mathbf A:a_\infty=0,\ a_v\in\mathcal O_v\text{ for all }v\}$
compact, under the continuous map $(a_j)\mapsto\sum_ja_js_j$.

The entries of the elements of $K$ lie in compact sets. At an infinite place
$w$, let $M_w=\tau_w(J_X)$. As $\theta_{X,w}$ is a Cartan involution of $M_w$,
$Q_w(u)=\operatorname{Re}((\theta_{X,w}u)^{\mathsf T}M_w\bar u)$ is positive on
the unit sphere, so $Q_w(u)\geq\kappa_w|u|^2$ for some $\kappa_w>0$. For
$k\in K$, the matrix $k_w=\pi_w(k)$ satisfies $k_w^{\mathsf T}M_w\bar k_w=M_w$
and commutes with $\theta_{X,w}$, so $Q_w(k_wu)=Q_w(u)$; hence
$|k_we_j|^2\leq Q_w(e_j)/\kappa_w$, which bounds the entries of $k_w$. At the
finite places, let $e_f\in\mathbf A$ be $0$ at $\infty$ and $1$ at every finite
place. For $u\in E^n$ there is a compact $C_u\subset\mathbf A_E^n$ containing
$N(e_fu)$ for every matrix $N$ with $N\hat L\subset\hat L$: for $u\in L_X$ take
$C_u=\hat L$, since $e_fu\in\hat L$; and $C_{u+u'}=C_u+C_{u'}$ and
$C_{\lambda u}=\lambda C_u$ for $\lambda\in E$ cover the $E$-span of $L_X$,
which is $E^n$. With $u$ the standard basis vectors this puts $e_fk_{ij}$ in a
compact set for $k\in K$. Since $F$ is totally real, there is a continuous
$\Phi\colon\mathbf C^{\{w\}}\to\mathbf A_E$ with
$z=\Phi((\pi_w(z))_w)+e_fz$ for every $z\in\mathbf A_E$. So $K$ lies in a
compact set $S$ of matrices, and so does $K^{-1}=K$.

The conditions defining $K$ are closed, $\hat L$ being compact, and $U_X$ is
closed in $GL_n(\mathbf A_E)$. The map $g\mapsto(g,g^{-1})$ embeds
$GL_n(\mathbf A_E)$ as a closed subset of $M_n(\mathbf A_E)^2$, and it maps
$K$ into $S\times S$. So $K$ is compact. A compact group and the locally
compact group $\mathbf A_E^\times$ have Haar measures.
\end{proof}

\begin{theorem}[Schrödinger operators and the theta distribution \cite{Wei64}]\label{asm-theta-agree-evaluation-heisenberg}
Let $F$ be a number field with adele ring $\mathbf A$, finite adeles
$\mathbf A_f$ and $F_\infty=F\otimes\mathbf R$, and let $\psi$ be a continuous
character of $\mathbf A$ that is trivial on $F$ and not identically $1$. Let
$\kappa$ be a finite set, and let $\mathcal S(\mathbf A^\kappa)$ be the space of
Schwartz-Bruhat functions on $\mathbf A^\kappa$
(Definition~\ref{asm-theta-agree-models}). Call a linear form $\lambda$ on
functions on $\mathbf A^\kappa$ \emph{admissible} if $f\mapsto\lambda(f\otimes g)$
is continuous on $\mathcal S(F_\infty^\kappa)$ for every Schwartz-Bruhat
function $g$ on $\mathbf A_f^\kappa$. For $w=(a,b)\in\mathbf A^\kappa\times\mathbf A^\kappa$
let $\rho(w)$ be the \emph{Schrödinger operator}
$(\rho(w)\phi)(x)=\psi(x\cdot b+\tfrac12a\cdot b)\phi(x+a)$, and for
$\Phi\in\mathcal S(\mathbf A^\kappa)$ let $\theta(\Phi)=\sum_{\xi\in F^\kappa}\Phi(\xi)$.
Then:
\begin{enumerate}
\item the operators $\rho(w)$ map $\mathcal S(\mathbf A^\kappa)$ into itself,
  and $\theta(\rho(w)\Phi)=\theta(\Phi)$ for $w\in F^\kappa\times F^\kappa$ and
  $\Phi\in\mathcal S(\mathbf A^\kappa)$;
\item the sum $\theta(\Phi)$ converges absolutely for
  $\Phi\in\mathcal S(\mathbf A^\kappa)$, some admissible linear form agrees with
  $\theta$ on $\mathcal S(\mathbf A^\kappa)$, and $\theta(\Phi)\neq0$ for some
  $\Phi\in\mathcal S(\mathbf A^\kappa)$;
\item if $\Phi\in\mathcal S(\mathbf A^\kappa)$ and $\theta(\rho(w)\Phi)=0$ for
  every $w\in\mathbf A^\kappa\times\mathbf A^\kappa$, then $\Phi=0$;
\item if $\lambda$ is an admissible linear form with
  $\lambda(\rho(w)\Phi)=\lambda(\Phi)$ for all $w\in F^\kappa\times F^\kappa$ and
  $\Phi\in\mathcal S(\mathbf A^\kappa)$, then there is $a\in\mathbf C$ with
  $\lambda(\Phi)=a\,\theta(\Phi)$ for $\Phi\in\mathcal S(\mathbf A^\kappa)$;
\item if $T$ is a linear map on functions on $\mathbf A^\kappa$ that maps
  $\mathcal S(\mathbf A^\kappa)$ into itself, such that $\lambda\circ T$ is
  admissible for every admissible $\lambda$ and $T\rho(w)\Phi=\rho(w)T\Phi$ for
  every $w$ and $\Phi\in\mathcal S(\mathbf A^\kappa)$, then there is $a\in\mathbf C$
  with $T\Phi=a\Phi$ for $\Phi\in\mathcal S(\mathbf A^\kappa)$.
\end{enumerate}
\end{theorem}
\leanref{Development/AutomorphicSpectrum/Proofs/MembershipProofs/ThetaDistributionNonzeroAdmissibleStatement.lean}{GlobalTheory.T.theta_distribution_heisenberg_invariant}
\leanref{Development/AutomorphicSpectrum/Proofs/MembershipProofs/ThetaSumAdmissibilityStatement.lean}{GlobalTheory.T.theta_distribution_nonzero_admissible}
\leanref{Development/AutomorphicSpectrum/Proofs/MembershipProofs/HeisenbergCommutingOperatorScalarStatement.lean}{GlobalTheory.T.theta_heisenberg_translates_separate}
\leanref{Development/AutomorphicSpectrum/Proofs/MembershipProofs/AdelicCharacterRationalDualStatement.lean}{GlobalTheory.T.admissible_functional_rational_invariance}
\leanref{Development/AutomorphicSpectrum/Proofs/MembershipProofs/WeilRepresentationThetaEquivarianceStatement.lean}{GlobalTheory.T.heisenberg_commuting_operator_scalar}

\begin{theorem}[Adelic metaplectic operators \cite{Wei64}]\label{asm-theta-agree-evaluation-metaplectic}
Let $F$, $\mathbf A$, $\mathbf A_f$ and $F_\infty$ be as in
Theorem~\ref{asm-theta-agree-evaluation-heisenberg}, and let $\psi$ be a
continuous character of $\mathbf A$, trivial on $F$, with
$\psi(a)=\exp(-2\pi i\operatorname{Tr}_{F_\infty/\mathbf R}a_\infty)$ for $a$ at
$\infty$. Let $\kappa$ be a finite set, let $\rho$ be the Schrödinger operators
of Theorem~\ref{asm-theta-agree-evaluation-heisenberg}, and let $S$ be an
$\mathbf A$-linear automorphism of $\mathbf A^\kappa\times\mathbf A^\kappa$ with
$S(w)_1\cdot S(w')_2-S(w')_1\cdot S(w)_2=w_1\cdot w'_2-w'_1\cdot w_2$ for all
$w,w'$. Then there are a linear map $e$ on functions on $\mathbf A^\kappa$, a
continuous linear automorphism $e_\infty$ of $\mathcal S(F_\infty^\kappa)$, and
for each finite place $v$ a linear map $e_v$ on functions on $F_v^\kappa$, such
that:
\begin{enumerate}
\item $e$ maps $\mathcal S(\mathbf A^\kappa)$ into itself, and
  $e(\rho(w)\Phi)=\rho(Sw)e(\Phi)$ for every $w$ and
  $\Phi\in\mathcal S(\mathbf A^\kappa)$;
\item each $e_v$ maps the Schwartz-Bruhat functions on $F_v^\kappa$ bijectively
  onto themselves, and $e_v(\mathbf 1_{\mathcal O_v^\kappa})=\mathbf 1_{\mathcal O_v^\kappa}$
  for almost all $v$;
\item if $w$ is at $v$ and $\rho(w)$ acts at $v$ by $\phi\mapsto\phi_1$
  (Definition~\ref{asm-theta-agree-models}(2)) with $\phi$ Schwartz-Bruhat, then
  $\rho(Sw)$ acts at $v$ by $e_v(\phi)\mapsto e_v(\phi_1)$;
\item $e(f\otimes\bigotimes_v\phi_v)=e_\infty(f)\otimes\bigotimes_ve_v(\phi_v)$
  for $f\in\mathcal S(F_\infty^\kappa)$ and Schwartz-Bruhat $\phi_v$ with
  $\phi_v=\mathbf 1_{\mathcal O_v^\kappa}$ for almost all $v$.
\end{enumerate}
\end{theorem}
\leanref{Development/WeilBound/ThetaLifts/LeviDeterminants/MixedModels/RealSt.lean}{PositivityTheory.Mix.AdelSt}
\leanref{Development/WeilBound/ThetaLifts/LeviDeterminants/MixedModels/AdelicEstimatesHolds.lean}{PositivityTheory.Mix.adelic_estimates_holds}

\begin{lemma}\label{asm-theta-agree-evaluation-transfer}
Let $F$, $\mathbf A$, $\psi$, $\kappa$ and $\rho$ be as in
Theorem~\ref{asm-theta-agree-evaluation-metaplectic}, and let $\theta$ and
admissible linear forms be as in
Theorem~\ref{asm-theta-agree-evaluation-heisenberg}. Let $\Lambda$ be an
$\mathbf A$-module with an $\mathbf A$-bilinear form $B$, and let $d$ and $d'$ be
\emph{Darboux coordinates} on $\Lambda$: $\mathbf A$-linear isomorphisms
$\Lambda\to\mathbf A^\kappa\times\mathbf A^\kappa$ with
$B(w,w')=d(w)_1\cdot d(w')_2-d(w')_1\cdot d(w)_2$, and likewise for $d'$. A
\emph{transfer} from $d$ to $d'$ is a linear map $t$ on functions on
$\mathbf A^\kappa$ that maps $\mathcal S(\mathbf A^\kappa)$ into itself, with
$t(\rho(d(w))\Phi)=\rho(d'(w))t(\Phi)$ for $w\in\Lambda$ and
$\Phi\in\mathcal S(\mathbf A^\kappa)$, with $\lambda\circ t$ admissible for every
admissible $\lambda$, and with $t(\Phi)\neq0$ for some
$\Phi\in\mathcal S(\mathbf A^\kappa)$. Call $d$ \emph{rationally related} to $d'$
if for all $y,z\in F^\kappa$ there is $w\in\Lambda$ with $d(w)=(y,z)$ and
$d'(w)\in F^\kappa\times F^\kappa$.
\begin{enumerate}
\item Let $e$, $e_\infty$ and $e_v$ be as in
  Theorem~\ref{asm-theta-agree-evaluation-metaplectic} for $S=d'\circ d^{-1}$.
  Then $e$ is a transfer from $d$ to $d'$ and maps $\mathcal S(\mathbf A^\kappa)$
  onto itself. In particular there is a transfer from $d$ to $d'$.
\item If $d$ is rationally related to $d'$ and $t$ is a transfer from $d$ to
  $d'$, then there is $c\neq0$ with $\theta(t\Phi)=c\,\theta(\Phi)$ for
  $\Phi\in\mathcal S(\mathbf A^\kappa)$, and $t$ is injective on
  $\mathcal S(\mathbf A^\kappa)$.
\item If $d$ and $d'$ are rationally related to each other, $t$ is a transfer
  from $d$ to $d'$ and $t'$ one from $d'$ to $d$, then there is a transfer
  $t''$ from $d'$ to $d$ with $t''(t\Phi)=\Phi$ and $t(t''\Phi)=\Phi$ for
  $\Phi\in\mathcal S(\mathbf A^\kappa)$.
\end{enumerate}
\end{lemma}
\leanref{Development/WeilBound/ThetaLifts/LeviDeterminants/ExistsNeOfSurjOn.lean}{PositivityTheory.ThKF0.tw_surj}
\leanref{Development/WeilBound/ThetaLifts/LeviDeterminants/EqZeroOfTheta.lean}{PositivityTheory.ThKF0.Tw}
\leanref{Development/WeilBound/ThetaLifts/LeviDeterminants/EqZeroOfTheta.lean}{PositivityTheory.ThKF0.RatRel}
\leanref{Development/WeilBound/ThetaLifts/LeviDeterminants/EqZeroOfTheta.lean}{PositivityTheory.ThKF0.Tw.theta}
\leanref{Development/WeilBound/ThetaLifts/LeviDeterminants/EqZeroOfTheta.lean}{PositivityTheory.ThKF0.Tw.inj}
\leanref{Development/WeilBound/ThetaLifts/LeviDeterminants/EqZeroOfTheta.lean}{PositivityTheory.ThKF0.Tw.exists_inv}
\begin{proof}
(1) As $d$ and $d'$ are Darboux coordinates for the same form, $S=d'\circ d^{-1}$
satisfies the hypothesis of
Theorem~\ref{asm-theta-agree-evaluation-metaplectic}. By its part (1), $e$ maps
$\mathcal S(\mathbf A^\kappa)$ into itself and
$e(\rho(d(w))\Phi)=\rho(Sd(w))e(\Phi)=\rho(d'(w))e(\Phi)$.

Every Schwartz-Bruhat function $g$ on $\mathbf A_f^\kappa$ is a finite linear
combination of products $\bigotimes_v\phi_v$ of Schwartz-Bruhat functions with
$\phi_v=\mathbf 1_{\mathcal O_v^\kappa}$ for almost all $v$. Indeed, $g$ is
invariant under translation by $a\hat{\mathcal O}^\kappa$ for some $a\in F^\times$,
where $\hat{\mathcal O}=\prod_v\mathcal O_v$, since it is locally constant with
compact support; so $g$ is a finite linear combination of indicator functions
of cosets $r+a\hat{\mathcal O}^\kappa$, and each of these is such a product.

$e$ maps $\mathcal S(\mathbf A^\kappa)$ onto itself. By the last paragraph it
suffices to reach $f\otimes\bigotimes_v\phi_v$ with $\phi_v$ Schwartz-Bruhat and
almost all equal to $\mathbf 1_{\mathcal O_v^\kappa}$. Write $f=e_\infty(f_0)$.
For each $v$ choose a Schwartz-Bruhat $\phi^0_v$ with $e_v(\phi^0_v)=\phi_v$,
taking $\phi^0_v=\phi_v$ when $e_v(\phi_v)=\phi_v$. Then
$\phi^0_v=\mathbf 1_{\mathcal O_v^\kappa}$ for almost all $v$ by
Theorem~\ref{asm-theta-agree-evaluation-metaplectic}(2), and
$e(f_0\otimes\bigotimes_v\phi^0_v)=f\otimes\bigotimes_v\phi_v$ by part (4).

$\lambda\circ e$ is admissible when $\lambda$ is. Let $g$ be a Schwartz-Bruhat
function on $\mathbf A_f^\kappa$ and write $g=\sum_kc_k\bigotimes_v\phi_{k,v}$ as
above. For $f\in\mathcal S(F_\infty^\kappa)$, part (4) gives
$\lambda(e(f\otimes g))=\sum_kc_k\lambda(e_\infty(f)\otimes g_k)$, where
$g_k=\bigotimes_ve_v(\phi_{k,v})$ is Schwartz-Bruhat by part (2). Each term is
continuous in $f$, since $e_\infty$ is continuous and $\lambda$ is admissible.

Finally $\mathcal S(\mathbf A^\kappa)\neq0$ and $e$ maps it onto itself, so
$e(\Phi)\neq0$ for some $\Phi$. Hence $e$ is a transfer.

(2) By Theorem~\ref{asm-theta-agree-evaluation-heisenberg}(2) there is an
admissible $\ell$ equal to $\theta$ on $\mathcal S(\mathbf A^\kappa)$. Then
$\ell\circ t$ is admissible. Let $(y,z)\in F^\kappa\times F^\kappa$, and choose
$w\in\Lambda$ with $d(w)=(y,z)$ and $d'(w)$ rational. For
$\Phi\in\mathcal S(\mathbf A^\kappa)$, as $t(\rho(y,z)\Phi)=\rho(d'(w))t(\Phi)$,
\[
\ell(t(\rho(y,z)\Phi))=\theta(\rho(d'(w))t(\Phi))=\theta(t(\Phi))=\ell(t(\Phi))
\]
by Theorem~\ref{asm-theta-agree-evaluation-heisenberg}(1). By part (4) of that
theorem there is $c$ with $\theta(t\Phi)=c\,\theta(\Phi)$ on
$\mathcal S(\mathbf A^\kappa)$. Suppose $c=0$, and let $\Phi$ have $t\Phi\neq0$.
For every $w\in\Lambda$,
$\theta(\rho(d'(w))t\Phi)=\theta(t(\rho(d(w))\Phi))=c\,\theta(\rho(d(w))\Phi)=0$.
As $d'$ is onto, part (3) of that theorem gives $t\Phi=0$, a contradiction. So
$c\neq0$. If $t\Phi=0$, the same computation gives
$c\,\theta(\rho(d(w))\Phi)=\theta(\rho(d'(w))t\Phi)=0$ for every $w$; as $d$ is
onto and $c\neq0$, part (3) gives $\Phi=0$.

(3) The map $t'\circ t$ maps $\mathcal S(\mathbf A^\kappa)$ into itself, keeps
admissible forms admissible, and commutes with every $\rho(d(w))$, hence with
every $\rho(w)$. By Theorem~\ref{asm-theta-agree-evaluation-heisenberg}(5) there
is $a$ with $t'(t\Phi)=a\Phi$ on $\mathcal S(\mathbf A^\kappa)$. If $a=0$, take
$\Phi$ with $t\Phi\neq0$; then $t'(t\Phi)=0$, so $t\Phi=0$ by the injectivity of
$t'$ in (2), a contradiction. Put $t''=a^{-1}t'$, a transfer from $d'$ to $d$
with $t''(t\Phi)=\Phi$. For $\Phi\in\mathcal S(\mathbf A^\kappa)$, the function
$\Psi=t(t''\Phi)-\Phi$ lies in $\mathcal S(\mathbf A^\kappa)$ and has
$t''\Psi=t''\Phi-t''\Phi=0$, so $\Psi=0$ by the injectivity of $t''$ in (2).
\end{proof}
% Lean: (1) is tw_surj, with surjOn_of_pure, Span.exists_sum_finprod (exists_sum_cosets,
% exists_invariant, indicator_coset_scaledInt_eq_finprod), Span.comp_mem_adm and
% exists_ne_of_surjOn (exists_adelicSchwartz_ne_zero_at_origin). (2) is Tw.theta and Tw.inj,
% through eq_zero_of_theta and rho_mem; (3) is Tw.exists_inv, through schur.

\begin{lemma}\label{asm-theta-agree-evaluation-rational}
Let $B$ be a theta base for a CM field $E$, let $X$ and $W$ be lattice spaces
of ranks $n$ and $d$, put $I=\{1,\ldots,n\}\times\{1,\ldots,d\}$, and let $M$ and
$M'$ be regular models for $(X,W)$ (Definition~\ref{asm-theta-agree-models}(3))
with isomorphisms $e^M$ and $e^{M'}$. Then for all $y,z\in F^I$ there is
$x\in M_{n\times d}(E)$ with $e^M(x)=(y,z)$, and $e^{M'}(x)\in F^I\times F^I$.
In the terms of Lemma~\ref{asm-theta-agree-evaluation-transfer}, $e^M$ and
$e^{M'}$ are rationally related to each other.
\end{lemma}
\leanref{Development/WeilBound/ThetaLifts/LeviDeterminants/ModelChangeOf.lean}{PositivityTheory.ThKF0.ratRel_of_good}
\leanref{Development/WeilBound/ThetaLifts/LeviDeterminants/ModelChangeOf.lean}{PositivityTheory.ThKF0.exists_linearEquiv_of_rational}
\begin{proof}
As $M$ is regular, $e^M$ maps $M_{n\times d}(E)$ into $F^I\times F^I$. Its
restriction there is $F$-linear, since $e^M$ is $\mathbf A$-linear, and
injective, since $e^M$ is injective and $E\to\mathbf A_E$ is injective. Both
spaces have dimension $2nd$ over $F$, as $[E:F]=2$. So the restriction is a
bijection onto $F^I\times F^I$, and there is $x$ with $e^M(x)=(y,z)$. As $M'$ is
regular, $e^{M'}(x)$ is rational.
\end{proof}

\begin{lemma}\label{asm-theta-agree-evaluation-lattice}
Let $B$ be a theta base for a CM field $E$, with character $\psi$, let $X$ be a
Hermitian and $W$ a skew-Hermitian lattice space, of ranks $n$ and $d$ and of
some signatures, with lattices $L_X$ and $L_W$, and let $h$ and
$\langle\ ,\ \rangle$ be as in Definition~\ref{asm-theta-agree-models}(3). Let
$v$ be a finite place of $F$, and let $\Lambda_v\subset\mathbb W(X,W)$ be the
group generated by the matrices $a\,xy^{\mathsf T}$ with $a\in\mathbf A$ at $v$,
$a_v\in\mathcal O_v$, $x\in L_X$ and $y\in L_W$. Then:
\begin{enumerate}
\item $\psi\bigl(b\operatorname{Tr}_{\mathbf A_E/\mathbf A}(o\,h(\gamma,\gamma'))\bigr)=1$
  for $\gamma,\gamma'\in\Lambda_v$, $o\in\mathcal O_E$ and $b\in\mathbf A$ with
  $b_v\in\mathcal O_v$;
\item if $z\in\mathbb W(X,W)$ is at $v$ and $\psi(\langle\gamma,z\rangle)=1$ for
  every $\gamma\in\Lambda_v$, then $z\in\Lambda_v$.
\end{enumerate}
\end{lemma}
\leanref{Development/AbelianVarietyTheory/Proofs/DivisorProofs/StandardAddCharSymplecticLocalLatticeEqOne.lean}{GlobalTheory.LatVacArith.psi_symp_latAtV}
\leanref{Development/WeilBound/CoreStatements/CompactAveragingSplitting/EigenvectorLifting/VaceigenDual.lean}{PositivityTheory.ThK.Vac.mem_latAtV}
% Lean: both are computations with the lattices L_X and L_W; mem_latAtV goes through
% dual_family (a dual family of L_X over O_E), mem_latW, exists_tr_rep, span_incV and
% mem_int_of_psi.

\begin{lemma}\label{asm-theta-agree-evaluation-vacuum}
Let $B$ be a theta base for a CM field $E$, let $X$ be a Hermitian and $W$ a
skew-Hermitian lattice space, of some signatures, and let $M$ be a model for
$(X,W)$ (Definition~\ref{asm-theta-agree-models}(3)). Let $v$ be a finite
place, and let $\phi$ and $\phi'$ be vacua for $M$ at $v$. Then $\phi'=c\,\phi$
for some $c\neq0$.
\end{lemma}
\leanref{Development/WeilBound/ThetaLifts/LeviDeterminants/ChiIndep.lean}{PositivityTheory.ThKF0.vac_line}
\leanref{Development/WeilBound/ThetaLifts/LeviDeterminants/ChiIndep.lean}{PositivityTheory.ThKF0.chi_indep}
\leanref{Development/WeilBound/CoreStatements/CompactAveragingSplitting/EigenvectorLifting/Vaceigen.lean}{PositivityTheory.ThK.Vac.local_line}
\leanref{Development/WeilBound/CoreStatements/CompactAveragingSplitting/EigenvectorLifting/VaceigenAlg.lean}{PositivityTheory.ThK.Vac.line_of_eigen}
\begin{proof}
Let $n$ and $d$ be the ranks of $X$ and $W$, put
$I=\{1,\ldots,n\}\times\{1,\ldots,d\}$, let $\Lambda_v$ be as in
Lemma~\ref{asm-theta-agree-evaluation-lattice}, let $e^M$ be the isomorphism of
$M$ and $\rho$ its Schrödinger operators, and let $\theta,\theta'\in\mathcal O_E$
with $\theta+c(\theta)=\theta'+c(\theta')=1$ be the elements in the definition
of the two vacua. For $\gamma\in\Lambda_v$ put $T=h(\gamma,\gamma)$,
$\chi(\gamma)=\psi(\tfrac12\operatorname{Tr}(\theta T))$ and
$\chi'(\gamma)=\psi(\tfrac12\operatorname{Tr}(\theta'T))$, with
$\operatorname{Tr}=\operatorname{Tr}_{\mathbf A_E/\mathbf A}$. As $h$ is linear
in its first argument, $\rho(\gamma)$ acts at $v$ by $\phi\mapsto\chi(\gamma)\phi$
and by $\phi'\mapsto\chi'(\gamma)\phi'$.

\emph{The two characters agree.} As $X$ is Hermitian and $W$ skew-Hermitian,
$c(T)=-T$. Write $a=\theta$ and $a'=\theta'$, so $c(a)=1-a$ and $c(a')=1-a'$.
Then
$c\bigl(a(a-a')T\bigr)=(1-a)(a'-a)(-T)=(a-a')T-a(a-a')T$. The trace is
invariant under $c$, so
$\operatorname{Tr}(a(a-a')T)=\operatorname{Tr}(aT)-\operatorname{Tr}(a'T)-\operatorname{Tr}(a(a-a')T)$,
that is
$\tfrac12\operatorname{Tr}(aT)=\tfrac12\operatorname{Tr}(a'T)+\operatorname{Tr}(a(a-a')T)$.
By Lemma~\ref{asm-theta-agree-evaluation-lattice}(1) with $o=\theta(\theta-\theta')$
and $b=1$, $\psi(\operatorname{Tr}(a(a-a')T))=1$. So $\chi'=\chi$.

\emph{Local reading.} Let $\rho_v(w)$, for $w\in F_v^I\times F_v^I$, be the
Schrödinger operators on functions on $F_v^I$ for $\psi_v$, given by the formula
of Theorem~\ref{asm-theta-agree-evaluation-heisenberg}. If $\gamma$ is at $v$,
then so is $e^M(\gamma)$, and $\rho(\gamma)(\phi\otimes r)=(\rho_v(\gamma_v)\phi)\otimes r$,
where $\gamma_v$ is the component of $e^M(\gamma)$ at $v$. So
$\rho_v(\gamma_v)\phi=\chi(\gamma)\phi$ and $\rho_v(\gamma_v)\phi'=\chi(\gamma)\phi'$
for $\gamma\in\Lambda_v$. Let $\Gamma$ be the set of these $\gamma_v$, and write
$[w,w']=w_1\cdot w'_2-w'_1\cdot w_2$ on $F_v^I\times F_v^I$. A direct
computation from the formula gives
$\rho_v(w)\rho_v(w')=\psi_v([w,w'])\rho_v(w')\rho_v(w)$, and
$\langle u,u'\rangle=e^M_1(u)\cdot e^M_2(u')-e^M_1(u')\cdot e^M_2(u)$. Hence: if
$w\in F_v^I\times F_v^I$ has $\psi_v([\gamma_v,w])=1$ for all
$\gamma\in\Lambda_v$, then $w\in\Gamma$. Indeed, let $\hat w$ be the element of
$\mathbf A^I\times\mathbf A^I$ at $v$ with component $w$ there, and
$z=(e^M)^{-1}(\hat w)$, which is at $v$ as $e^M$ is $\mathbf A$-linear. Then
$\psi(\langle\gamma,z\rangle)=\psi_v([\gamma_v,w])=1$ for $\gamma\in\Lambda_v$, so
$z\in\Lambda_v$ by Lemma~\ref{asm-theta-agree-evaluation-lattice}(2), and
$w=z_v\in\Gamma$.

\emph{The line.} Let $\langle f,f'\rangle=\int\bar f f'$ for a Haar measure on
$F_v^I$. The operators $\rho_v(w)$ preserve it, and the representation $\rho_v$
on the Schwartz-Bruhat functions on $F_v^I$ is irreducible, by the Stone-von
Neumann theorem \cite[Chapitre~2]{MVW87}. Since $\phi\neq0$,
$|\chi(\gamma)|=1$. Let $\eta$ be a Schwartz-Bruhat function with
$\rho_v(\gamma_v)\eta=\chi(\gamma)\eta$ for $\gamma\in\Lambda_v$ and
$\langle\eta,\phi\rangle=0$. We show $\eta=0$. Let $w\in F_v^I\times F_v^I$. If
$w=\gamma_v$ with $\gamma\in\Lambda_v$, then
$\langle\eta,\rho_v(w)\phi\rangle=\chi(\gamma)\langle\eta,\phi\rangle=0$.
Otherwise there is $\gamma\in\Lambda_v$ with
$\zeta=\psi_v([\gamma_v,w])\neq1$. Then
$\rho_v(\gamma_v)\rho_v(w)\phi=\zeta\chi(\gamma)\rho_v(w)\phi$, so
\[
\langle\eta,\rho_v(w)\phi\rangle=\langle\rho_v(\gamma_v)\eta,\rho_v(\gamma_v)\rho_v(w)\phi\rangle
=\overline{\chi(\gamma)}\,\zeta\chi(\gamma)\langle\eta,\rho_v(w)\phi\rangle
=\zeta\langle\eta,\rho_v(w)\phi\rangle,
\]
and $\langle\eta,\rho_v(w)\phi\rangle=0$. By irreducibility the
$\rho_v(w)\phi$ span the Schwartz-Bruhat functions, so $\langle\eta,\eta\rangle=0$
and $\eta=0$. Apply this to $\eta=\phi'-c\phi$ with
$c=\langle\phi,\phi'\rangle/\langle\phi,\phi\rangle$. Then $\phi'=c\phi$, and
$c\neq0$ since $\phi'\neq0$.
\end{proof}
% Lean: the characters are chi_indep (psi_symp_latAtV, conj_hTens, trace_conj); the
% local reading is loc_eig (rho_acts_locCo, eq_of_tens_eq); the line is local_line, which
% builds the local model with SchrModel.exists_isModel (inner product and commutation
% rule), takes irreducibility from VacLoc.ModelLine.cyclic_and_line, and concludes by
% line_of_eigen.

\begin{lemma}\label{asm-theta-agree-evaluation-change}
Let $B$ be a theta base for a CM field $E$, let $X$ be a Hermitian and $W$ a
skew-Hermitian lattice space, of ranks $n$ and $d$ and of some signatures, and
let $M$ and $M'$ be regular models for $(X,W)$. Then there is an evaluation map
$e$ for $(M,M',1,1,\emptyset)$ (Definition~\ref{asm-theta-agree-planes}(5)) that
maps the Schwartz-Bruhat functions on $\mathbf A^{n\times d}$ bijectively onto
themselves.
\end{lemma}
\leanref{Development/WeilBound/ThetaLifts/LeviDeterminants/EvAdel.lean}{PositivityTheory.ThKF0.isEv_bij_of_adel}
\leanref{Development/WeilBound/ThetaLifts/LeviDeterminants/MapEqOfSmul.lean}{PositivityTheory.ThKF0.EvBij}
\leanref{Development/WeilBound/ThetaLifts/LeviDeterminants/MixedModels/DarbouxModelRationalLatticeOf.lean}{PositivityTheory.Mix.darboux_model_rational_lattice_of}
\begin{proof}
Put $I=\{1,\ldots,n\}\times\{1,\ldots,d\}$, let $e^M$ and $e^{M'}$ be the
isomorphisms of $M$ and $M'$, with Schrödinger operators $\rho$ and $\rho'$, and
write $\mathcal S$ for the Schwartz-Bruhat functions on $\mathbf A^I$. Both
$e^M$ and $e^{M'}$ are Darboux coordinates on $\mathbb W(X,W)$ for
$\langle\ ,\ \rangle$, in the sense of
Lemma~\ref{asm-theta-agree-evaluation-transfer}, and they are rationally related
to each other by Lemma~\ref{asm-theta-agree-evaluation-rational}. The character
$\psi$ of $B$ satisfies the hypotheses of
Theorem~\ref{asm-theta-agree-evaluation-metaplectic}.

\emph{The operator.} Let $t_0$, $t_\infty$ and $t_v$ be as in
Theorem~\ref{asm-theta-agree-evaluation-metaplectic} for
$S=e^{M'}\circ(e^M)^{-1}$. By Lemma~\ref{asm-theta-agree-evaluation-transfer}(1),
$t_0$ is a transfer from $e^M$ to $e^{M'}$ that maps $\mathcal S$ onto itself;
so $t_0(\rho(u)\Phi)=\rho'(u)t_0(\Phi)$ for $\Phi\in\mathcal S$.

\emph{Vacua go to vacua.} Fix $v$ and put $\phi_v=t_v(\mathrm{vac}_v)$. It is
Schwartz-Bruhat and nonzero, since $t_v$ is bijective on Schwartz-Bruhat
functions. Let $\theta$ be the element in the definition of the vacuum
$\mathrm{vac}_v$ of $M$. For $\gamma$ in the group $\Lambda_v$ of
Lemma~\ref{asm-theta-agree-evaluation-lattice}, $\gamma$ is at $v$, hence so is
$e^M(\gamma)$, and $\rho(\gamma)$ acts at $v$ by
$\mathrm{vac}_v\mapsto\chi(\gamma)\mathrm{vac}_v$ with
$\chi(\gamma)=\psi(\tfrac12\operatorname{Tr}h(\theta\gamma,\gamma))$. By
Theorem~\ref{asm-theta-agree-evaluation-metaplectic}(3) for $w=e^M(\gamma)$,
$\rho'(\gamma)$ acts at $v$ by $\phi_v\mapsto\chi(\gamma)\phi_v$. So $\phi_v$ is
a vacuum for the Schrödinger operators of $M'$ at $v$, with the same $\theta$.
For almost all $v$, $\mathrm{vac}_v=\mathbf 1_{\mathcal O_v^I}$ and
$t_v(\mathbf 1_{\mathcal O_v^I})=\mathbf 1_{\mathcal O_v^I}$, so
$\phi_v=\mathbf 1_{\mathcal O_v^I}$.

\emph{Normalisation.} By Lemma~\ref{asm-theta-agree-evaluation-vacuum}, applied
to the vacua $\mathrm{vac}'_v$ and $\phi_v$ of $M'$, there is $c_v\neq0$ with
$\phi_v=c_v\mathrm{vac}'_v$. For almost all $v$ both $\phi_v$ and
$\mathrm{vac}'_v$ are $\mathbf 1_{\mathcal O_v^I}$, as $M'$ is regular, and the
value at $0$ gives $c_v=1$. Put $c=\prod_vc_v\neq0$, $t=c^{-1}t_0$ and
$t'_v=c_v^{-1}t_v$. Then $t'_v(\mathrm{vac}_v)=\mathrm{vac}'_v$, $t'_v$ maps the
Schwartz-Bruhat functions onto themselves, and
$t(f\otimes\bigotimes_v\phi_v)=t_\infty(f)\otimes\bigotimes_vt'_v(\phi_v)$
whenever the $\phi_v$ are Schwartz-Bruhat and almost all equal $\mathrm{vac}_v$,
since $\mathrm{vac}_v=\mathbf 1_{\mathcal O_v^I}$ for almost all $v$.

\emph{Bijectivity.} The map $t$, a nonzero multiple of $t_0$, is a transfer from
$e^M$ to $e^{M'}$. By Lemma~\ref{asm-theta-agree-evaluation-transfer}(1) there is
a transfer from $e^{M'}$ to $e^M$, and by part (3) of that lemma there is $t''$
with $t''(t\Phi)=\Phi$ and $t(t''\Phi)=\Phi$ for $\Phi\in\mathcal S$, where
$t''\Phi\in\mathcal S$. So $t$ maps $\mathcal S$ bijectively onto itself.

\emph{Evaluation map.} The map $t$ maps $\mathcal S$ into itself; for
$z\in\mathbb W(X,W)$ and $\Phi\in\mathcal S$,
$t(\rho(1\cdot z\cdot1^{\mathsf T})\Phi)=\rho'(z)t(\Phi)$; the invariance
condition is empty for $\emptyset$; and its components are $t_\infty$, which is
nonzero since it is surjective and $\mathcal S(F_\infty^I)\neq0$, and the $t'_v$,
with the properties shown above. So $t$ is an evaluation map for
$(M,M',1,1,\emptyset)$.
\end{proof}
% Lean: the operator is darboux_model_rational_lattice_of (AdelSt for M.co.e.symm.trans
% co.e) with surjOn_of_pure, Span.comp_mem_adm and exists_ne_of_surjOn; the normalisation is
% vac_line and the finprod of the c v; bijectivity is twEx and Tw.exists_inv with
% ratRel_of_good both ways.

\begin{lemma}\label{asm-theta-agree-evaluation-derivative}
Let $V$ be a finite-dimensional real vector space, $a\in V$, $L$ a real linear
form on $V$ and $f\in\mathcal S(V)$. For $t\in\mathbf R$ let
\[
G_t(x)=\exp\bigl(-2\pi i\,(t\,L(x)+\tfrac12t^2L(a))\bigr)f(x+ta),
\qquad
Df(x)=\partial_af(x)-2\pi i\,L(x)f(x).
\]
Then $G_t\in\mathcal S(V)$, and for every continuous linear form $T$ on
$\mathcal S(V)$ the function $t\mapsto T(G_t)$ has derivative $T(Df)$ at $t=0$.
\end{lemma}
\leanref{Development/WeilBound/ThetaLifts/LeviDeterminants/SchwartzDq.lean}{PositivityTheory.ThKF0.schwartzDq}
\leanref{Development/WeilBound/ThetaLifts/LeviDeterminants/DqStmt.lean}{PositivityTheory.ThKF0.SchwartzDq}
\leanref{Development/WeilBound/ThetaLifts/LeviDeterminants/EvLine.lean}{PositivityTheory.ThKF0.Line.hl_schwartz}
\begin{proof}
The phase $t\,L(x)+\tfrac12t^2L(a)$ is a polynomial in $(t,x)$, so
$G_t$ is $f$ translated by $ta$ times a smooth function whose derivatives grow
at most polynomially; hence $G_t\in\mathcal S(V)$. The function
$(t,x)\mapsto G_t(x)$ is smooth, and its derivative in $t$ at $t=0$ is $Df(x)$
by the product and chain rules. Its derivatives of every order in $(t,x)$,
times any power of $|x|$, are bounded uniformly for $|t|\leq1$, as $f$ is a
Schwartz function and the phase has derivatives of polynomial growth. By
Taylor's formula in $t$, every Schwartz seminorm of $G_t-G_0-t\,Df$ is then
$O(t^2)$. A continuous linear form $T$ is bounded by finitely many Schwartz
seminorms, so $T(G_t)-T(G_0)-t\,T(Df)=O(t^2)$, which is the claim.
\end{proof}
% Lean: the joint smoothness is temperate_exp and f.smooth; the uniform bounds are
% bound_mul_translate; the O(t^2) bound on seminorms is seminorm_remainder inside
% hasFDerivAt_family, with Seminorm.bound_of_continuous for T.

\begin{lemma}\label{asm-theta-agree-evaluation-pure}
Let $B$ be a theta base for a CM field $E$, let $X$ and $W$ be lattice spaces
of ranks $n$ and $d$, and let $M$ and $M'$ be regular models for $(X,W)$. Let
$e$ be an evaluation map for $(M,M',1,1,\emptyset)$
(Definition~\ref{asm-theta-agree-planes}(5)) that maps the Schwartz-Bruhat
functions on $\mathbf A^{n\times d}$ bijectively onto themselves, and let
$e_\infty$ and $e_v$ be components of $e$ as in that definition. Then
$e_\infty$ is surjective, and
$e_\infty(\operatorname{Pol}(M))=\operatorname{Pol}(M')$.
\end{lemma}
\leanref{Development/WeilBound/ThetaLifts/LeviDeterminants/EvChange.lean}{PositivityTheory.ThKF0.Change.change_core}
\leanref{Development/WeilBound/ThetaLifts/LeviDeterminants/EvChange.lean}{PositivityTheory.ThKF0.Change.inter}
\leanref{Development/WeilBound/ThetaLifts/LeviDeterminants/EvChange.lean}{PositivityTheory.ThKF0.Change.pol_map}
\begin{proof}
Put $I=\{1,\ldots,n\}\times\{1,\ldots,d\}$, write $\mathcal S$ for the
Schwartz-Bruhat functions on $\mathbf A^I$, let $\rho$ and $\rho'$ be the
Schrödinger operators of $M$ and $M'$, and write $(x,y)$ for the point of
$\mathbf A^I$ with component $x\in F_\infty^I$ at $\infty$ and finite component
$y$. Let $g=\bigotimes_v\mathrm{vac}_v$ and $g'=\bigotimes_v\mathrm{vac}'_v$, the
products of the vacua of $M$ and $M'$. They are Schwartz-Bruhat functions on
$\mathbf A_f^I$, as the models are regular, and there are $y$ and $y'$ with
$g(y)\neq0$ and $g'(y')\neq0$: the vacua are nonzero, and at the places where
they are $\mathbf 1_{\mathcal O_v^I}$ an integral point can be chosen.

\emph{Step 1.} $e(f\otimes g)=e_\infty(f)\otimes g'$ for
$f\in\mathcal S(F_\infty^I)$, by the factorization of $e$ and
$e_v(\mathrm{vac}_v)=\mathrm{vac}'_v$.

\emph{Step 2: $e_\infty$ is injective.} If $e_\infty(f)=0$, then
$e(f\otimes g)=0=e(0)$, so $f\otimes g=0$ by injectivity on $\mathcal S$, and
$f(x)g(y)=0$ for all $x$ gives $f=0$.

\emph{Step 3.} For every $\Phi\in\mathcal S$ there is $f$ with
$e(\Phi)(x,y')=e_\infty(f)(x)$ for all $x$. By linearity, and the decomposition
of Schwartz-Bruhat functions on $\mathbf A_f^I$ in the proof of
Lemma~\ref{asm-theta-agree-evaluation-transfer}(1), it suffices to take
$\Phi=f_0\otimes\bigotimes_v\phi_v$ with $\phi_v$ Schwartz-Bruhat and almost all
equal to $\mathbf 1_{\mathcal O_v^I}=\mathrm{vac}_v$. Then
$e(\Phi)(x,y')=e_\infty(f_0)(x)\prod_ve_v(\phi_v)(y'_v)$, which is
$e_\infty(f)(x)$ for a multiple $f$ of $f_0$.

\emph{Step 4: $e_\infty$ is surjective.} Let $f_1\in\mathcal S(F_\infty^I)$. As
$e$ maps $\mathcal S$ onto itself, $e(\Phi)=f_1\otimes g'$ for some
$\Phi\in\mathcal S$. By Step 3, $f_1(x)g'(y')=e_\infty(f)(x)$ for all $x$, so
$f_1=e_\infty(g'(y')^{-1}f)$.

\emph{Step 5: the derivative operators.} Let $u\in\mathbb W(X,W)$ be at $\infty$.
Let $D_u$ and $D'_u$ be the operators $D_{a,b}$ of
Lemma~\ref{asm-theta-agree-density-derivative} for the components $(a,b)$ at
$\infty$ of $e^M(u)$ and of $e^{M'}(u)$, where $e^M$ and $e^{M'}$ are the
isomorphisms of $M$ and $M'$. We show $e_\infty(D_uf)=D'_u(e_\infty f)$. As in
the proof of Lemma~\ref{asm-theta-agree-density-derivative},
$\rho(su)(f\otimes g)=G_s\otimes g$ for $s\in\mathbf R$, where $G_s$ is the
family of Lemma~\ref{asm-theta-agree-evaluation-derivative} for $f$, the vector
$a$ and the form $x\mapsto T(x,b)$; likewise for $M'$, with a family $G'_s$ for
$e_\infty(f)$. By the condition of an evaluation map with $\iota_X=\iota_W=1$,
$e(\rho(su)(f\otimes g))=\rho'(su)e(f\otimes g)$, which by Step 1 reads
$e_\infty(G_s)\otimes g'=G'_s\otimes g'$. At $(x,y')$ this gives
$e_\infty(G_s)(x)=G'_s(x)$ for all $s$. The derivative at $s=0$ of the left side
is $e_\infty(D_uf)(x)$, by Lemma~\ref{asm-theta-agree-evaluation-derivative}
for the continuous form $f\mapsto e_\infty(f)(x)$; that of the right side is
$D'_u(e_\infty f)(x)$.

\emph{Step 6: the Fock spaces.} By
Lemma~\ref{asm-theta-agree-density-derivative}, $s\mapsto\rho(su)$ acts at
$\infty$ to first order by $f\mapsto f_1$ if and only if $f_1=D_uf$, and
likewise for $M'$. So a Gaussian for $M$
(Definition~\ref{asm-theta-agree-models}(4)) is an $f_0$ with signs
$\varepsilon_w$ such that $D_{u'}f_0=i\,D_uf_0$ for all $u,u'$ at $\infty$ with
$\pi_w(u')=\varepsilon_wi\,\theta_{X,w}\pi_w(u)\theta_{W,w}^{\mathsf T}$; and
$\operatorname{Pol}(M)$ is the smallest subspace containing the Gaussians and
stable under every $D_u$. As $e_\infty$ is bijective (Steps 2 and 4) and
$e_\infty D_u=D'_ue_\infty$ (Step 5), $f_0$ is a Gaussian for $M$ if and only if
$e_\infty(f_0)$ is one for $M'$. Hence $e_\infty^{-1}(\operatorname{Pol}(M'))$
contains the Gaussians for $M$ and is stable under the $D_u$, so it contains
$\operatorname{Pol}(M)$; and $e_\infty(\operatorname{Pol}(M))$ contains the
Gaussians for $M'$ and is stable under the $D'_u$, so it contains
$\operatorname{Pol}(M')$.
\end{proof}
% Lean: Steps 1 to 4 are hvac, hinj, hsl and hsurj of change_core (finPr_mem,
% exists_finPr_ne, mem_span_finPr); Step 5 is inter, with Lib.heisLine_tens and
% Lib.hasDerivAt_hl and the hypothesis SchwartzDq; Step 6 is pol_map with Lib.actsD_unique
% and Lib.actsD_dop, the forms in this Lean file of Lemma asm-theta-agree-density-derivative.

\begin{lemma}\label{asm-theta-agree-evaluation-compose}
Let $B$ be a theta base for a CM field $E$, let $X$, $W$, $X_1$ and $W_1$ be
lattice spaces of ranks $n$, $d$, $n_1$ and $d_1$, and put
$I=\{1,\ldots,n\}\times\{1,\ldots,d\}$ and
$I_1=\{1,\ldots,n_1\}\times\{1,\ldots,d_1\}$.
\begin{enumerate}
\item Let $M$ and $M'$ be models for $(X,W)$ and $M_1$ one for $(X_1,W_1)$, and
  let $\iota_X$, $\iota_W$ and $\mathcal J$ be as in
  Definition~\ref{asm-theta-agree-planes}(5). Let $e_1$ be an evaluation map for
  $(M,M',1,1,\emptyset)$ with components $e_{1,\infty}$ and $e_{1,v}$, where
  $e_{1,\infty}$ is surjective, and let $e_2$ be an evaluation map for
  $(M',M_1,\iota_X,\iota_W,\mathcal J)$. Then $e_2\circ e_1$ is an evaluation map
  for $(M,M_1,\iota_X,\iota_W,\mathcal J)$.
\item Let $M$ be a model for $(X,W)$ and $M_1$ one for $(X_1,W_1)$, with vacua
  $\mathrm{vac}_v$ and $\mathrm{vac}_{1,v}$ that are Schwartz-Bruhat. Let $e$ be
  a linear map from functions on $\mathbf A^I$ to functions on $\mathbf A^{I_1}$,
  $e_\infty$ a linear map from $\mathcal S(F_\infty^I)$ to $\mathcal S(F_\infty^{I_1})$
  with $e_\infty(\operatorname{Pol}(M))=\operatorname{Pol}(M_1)$, and $e_v$
  linear maps with $e_v(\mathrm{vac}_v)=\mathrm{vac}_{1,v}$ that map the
  Schwartz-Bruhat functions onto the Schwartz-Bruhat functions, such that
  $e(f\otimes\bigotimes_v\phi_v)=e_\infty(f)\otimes\bigotimes_ve_v(\phi_v)$ for
  $f\in\mathcal S(F_\infty^I)$ and Schwartz-Bruhat $\phi_v$ almost all equal to
  $\mathrm{vac}_v$. Then $e(\operatorname{PD}(M))=\operatorname{PD}(M_1)$.
\end{enumerate}
\end{lemma}
\leanref{Development/WeilBound/ThetaLifts/LeviDeterminants/MapEqOfSmul.lean}{PositivityTheory.ThKF0.Top.isEv_comp}
\leanref{Development/WeilBound/ThetaLifts/LeviDeterminants/MapEqOfSmul.lean}{PositivityTheory.ThKF0.Top.polData_map}
\begin{proof}
(1) Let $\rho$, $\rho'$ and $\rho_1$ be the Schrödinger operators of $M$, $M'$
and $M_1$. As $e_1$ is an evaluation map with $\iota_X=\iota_W=1$,
$e_1(\rho(u)\Phi)=\rho'(u)e_1(\Phi)$ for every $u$ and Schwartz-Bruhat $\Phi$.
The composite maps Schwartz-Bruhat functions to Schwartz-Bruhat functions. For
$z\in\mathbb W(X_1,W_1)$,
$e_2e_1(\rho(\iota_Xz\iota_W^{\mathsf T})\Phi)=e_2(\rho'(\iota_Xz\iota_W^{\mathsf T})e_1\Phi)=\rho_1(z)e_2e_1\Phi$,
and for $u\in\mathcal J$,
$e_2e_1(\rho(u)\Phi)=e_2(\rho'(u)e_1\Phi)=e_2e_1\Phi$. Let $e_{2,\infty}$ and
$e_{2,v}$ be components of $e_2$. Then $e_{2,\infty}e_{1,\infty}\neq0$, since
$e_{1,\infty}$ is surjective and $e_{2,\infty}\neq0$; $e_{2,v}e_{1,v}$ carries
$\mathrm{vac}_v$ to $\mathrm{vac}'_v$ to $\mathrm{vac}_{1,v}$ and the
Schwartz-Bruhat functions onto themselves. If the $\phi_v$ are Schwartz-Bruhat
and almost all equal $\mathrm{vac}_v$, then the $e_{1,v}(\phi_v)$ are
Schwartz-Bruhat and almost all equal $\mathrm{vac}'_v$, so
$e_2e_1(f\otimes\bigotimes_v\phi_v)=e_{2,\infty}e_{1,\infty}(f)\otimes\bigotimes_ve_{2,v}e_{1,v}(\phi_v)$.

(2) $\operatorname{PD}(M)$ is spanned by the $f\otimes\bigotimes_v\phi_v$ with
$f\in\operatorname{Pol}(M)$ and $\phi_v$ Schwartz-Bruhat, almost all equal to
$\mathrm{vac}_v$. The image of such a generator is
$e_\infty(f)\otimes\bigotimes_ve_v(\phi_v)$, a generator of
$\operatorname{PD}(M_1)$. Conversely, let $f_1\otimes\bigotimes_v\phi'_v$ be a
generator of $\operatorname{PD}(M_1)$. Write $f_1=e_\infty(f)$ with
$f\in\operatorname{Pol}(M)$, and for each $v$ let $\phi_v=\mathrm{vac}_v$ if
$\phi'_v=\mathrm{vac}_{1,v}$, and otherwise a Schwartz-Bruhat function with
$e_v(\phi_v)=\phi'_v$. Then $e(f\otimes\bigotimes_v\phi_v)=f_1\otimes\bigotimes_v\phi'_v$.
\end{proof}

\begin{lemma}\label{asm-theta-agree-evaluation-glue}
Let $B$ be a theta base for a CM field $E$, let $X$, $W$, $X_1$ and $W_1$ be
lattice spaces of ranks $n$, $d$, $n_1$ and $d_1$, and put
$I=\{1,\ldots,n\}\times\{1,\ldots,d\}$ and
$I_1=\{1,\ldots,n_1\}\times\{1,\ldots,d_1\}$. Let $\kappa_0$ be a finite set and
$\epsilon\colon I\to\kappa_0\sqcup I_1$ a bijection; for vectors $p$ on
$\kappa_0$ and $b$ on $I_1$, let $[p,b]$ be the vector on $I$ with components
$p$ on $\epsilon^{-1}(\kappa_0)$ and $b$ on $\epsilon^{-1}(I_1)$. Let $M_a$ be a
model for $(X,W)$ and $M_1$ one for $(X_1,W_1)$, with isomorphisms $e^a$ and
$e^1$, and let $\iota_X$, $\iota_W$ and $\mathcal J$ be as in
Definition~\ref{asm-theta-agree-planes}(5). Suppose there are maps
$p,q\colon\mathbb W(X,W)\to\mathbf A^{\kappa_0}$ and
$\pi\colon\mathbb W(X,W)\to\mathbb W(X_1,W_1)$ such that
$e^a(u)=\bigl([p(u),e^1_1(\pi u)],[q(u),e^1_2(\pi u)]\bigr)$ for every $u$;
$p(\iota_Xz\iota_W^{\mathsf T})=0$ and $\pi(\iota_Xz\iota_W^{\mathsf T})=z$ for
$z\in\mathbb W(X_1,W_1)$; $p(u)=0$ and $\pi(u)=0$ for $u\in\mathcal J$; and the
vacua satisfy $\mathrm{vac}^a_v([0,b])=\mathrm{vac}^1_v(b)$. Then
$e(\Phi)(b)=\Phi([0,b])$ defines an evaluation map for
$(M_a,M_1,\iota_X,\iota_W,\mathcal J)$, with components
$e_\infty(f)(b)=f([0,b])$ and $e_v(\phi)(b)=\phi([0,b])$.
\end{lemma}
\leanref{Development/WeilBound/ThetaLifts/LeviDeterminants/EvModel.lean}{PositivityTheory.ThKF0.exists_isEvS}
\begin{proof}
Let $\rho_a$ and $\rho_1$ be the Schrödinger operators of $M_a$ and $M_1$.
Restricting a Schwartz-Bruhat function on $\mathbf A^I$ to the vectors $[0,b]$
gives a Schwartz-Bruhat function on $\mathbf A^{I_1}$, as restriction does on
each factor $f\otimes g$. The dot products split:
$[p,b]\cdot[q,b']=p\cdot q+b\cdot b'$.

\emph{Heisenberg condition.} Let $z\in\mathbb W(X_1,W_1)$ and
$u=\iota_Xz\iota_W^{\mathsf T}$. Then $e^a(u)=([0,e^1_1(z)],[q(u),e^1_2(z)])$,
so
\[
(\rho_a(u)\Phi)([0,b])=\psi\bigl(b\cdot e^1_2(z)+\tfrac12e^1_1(z)\cdot e^1_2(z)\bigr)\Phi([0,b+e^1_1(z)])
=(\rho_1(z)e(\Phi))(b).
\]

\emph{Invariance.} For $u\in\mathcal J$, $e^a(u)=(0,[q(u),0])$, and
$[0,b]\cdot[q(u),0]=0$, so $(\rho_a(u)\Phi)([0,b])=\Phi([0,b])$.

\emph{Components.} $e_\infty$ is a continuous linear map between the Schwartz
spaces, and it is nonzero, since $e_\infty(f)(0)=f(0)$ and some Schwartz
function has $f(0)\neq0$. By hypothesis $e_v(\mathrm{vac}^a_v)=\mathrm{vac}^1_v$.
Each $e_v$ maps Schwartz-Bruhat functions to Schwartz-Bruhat functions, and onto
them: $\phi_1$ is the image of $[p,b]\mapsto\mathbf 1_{\mathcal O_v^{\kappa_0}}(p)\phi_1(b)$.
Finally
$(f\otimes\bigotimes_v\phi_v)([0,b])=f([0,b_\infty])\prod_v\phi_v([0,b_v])$, which
is the factorization.
\end{proof}
% Lean: slice_mem_schA, glue_dot and glue_add for the two conditions; exists_slice and
% exists_schwartz_ne for e_infty; slice_mem_SB, box_mem_SB and box_glue for e_v.

\begin{lemma}\label{asm-theta-agree-evaluation-adapted-frame}
Let $B$ be a theta base for a CM field $E$, and let $Y$ and $Y_1$ be lattice
spaces of ranks $m$ and $m_1$, with matrices $J$ and $J_1$ and lattices $L$ and
$L_1$, such that $c(J)^{\mathsf T}=\eta J$ for some $\eta=\pm1$. Let
$H=(x,y,\iota,\mathfrak b)$ be an adapted plane in $Y$ over $Y_1$
(Definition~\ref{asm-theta-agree-planes}(1)), and let
$\mathrm{fr}_H\in M_{m\times(m_1+2)}(E)$ be the matrix whose columns are those of
$\iota$, then $x$, then $y$. Index the coordinates of $E^{m_1+2}$ by
$1,\ldots,m_1$, $x$ and $y$, write $Q_1z$ for the first $m_1$ coordinates of a
vector $Qz$, and for integers $r_0,r_x,r_y$ put
$\Delta(r_0,r_x,r_y)=\operatorname{diag}(r_01_{m_1},r_x,r_y)$. Then
$\mathrm{fr}_H$ is invertible, and its inverse $Q$ satisfies:
\begin{enumerate}
\item $Qx$ and $Qy$ are the unit vectors at $x$ and at $y$, and $Q\iota$ has rows
  $1_{m_1}$, then $0$ and $0$;
\item $\mathrm{fr}_H^{\mathsf T}J\,c(\mathrm{fr}_H)=\operatorname{diag}\bigl(J_1,\left(\begin{smallmatrix}0&1\\\eta&0\end{smallmatrix}\right)\bigr)$,
  and $c(J_1)^{\mathsf T}=\eta J_1$;
\item for $u\in L$, $Q_1u\in L_1$, $(Qu)_x\,x\in L$ and $(Qu)_y\,y\in L$;
\item for integers $r_0,r_x,r_y$ there is $A\in M_m(E)$ with $AL\subset L$ and
  $QA=\Delta(r_0,r_x,r_y)\,Q$.
\end{enumerate}
\end{lemma}
\leanref{Development/WeilBound/ThetaLifts/LeviDeterminants/MixedModels/AdaptC.lean}{PositivityTheory.Mix.Adapt.exists_frame_inv}
\leanref{Development/WeilBound/ThetaLifts/LeviDeterminants/MixedModels/AdaptC.lean}{PositivityTheory.Mix.Adapt.frame_gram}
\leanref{Development/WeilBound/ThetaLifts/LeviDeterminants/MixedModels/AdaptC.lean}{PositivityTheory.Mix.Adapt.mask_mem}
\leanref{Development/WeilBound/ThetaLifts/LeviDeterminants/EvFrame.lean}{PositivityTheory.ThKF0.exists_frOK}
\begin{proof}
For $w\in E^{m_1+2}$, $\mathrm{fr}_Hw=w_x\,x+\iota w'+w_y\,y$, with $w'$ the
first $m_1$ coordinates of $w$. By Definition~\ref{asm-theta-agree-planes}(1),
every $z\in L$ is $sx+\iota u+ty$ with $sx\in L$, $ty\in L$ and $u\in L_1$, so
$z=\mathrm{fr}_Hw$ for a $w$ with $w'\in L_1$, $w_xx\in L$ and $w_yy\in L$. As
the lattice $L$ spans $E^m$, $\mathrm{fr}_H$ is onto, and as it is square
($m=m_1+2$) it is invertible. Then $Q\,\mathrm{fr}_H=1$ gives (1), and the $w$
above is $Qz$, which gives (3).

(2) Put $h(a,a')=a^{\mathsf T}J\,c(a')$. As $c(J)^{\mathsf T}=\eta J$ and $c$ is
an involution, $h(a',a)=\eta\,c(h(a,a'))$. The entry of
$\mathrm{fr}_H^{\mathsf T}J\,c(\mathrm{fr}_H)$ at two columns $a,a'$ of
$\mathrm{fr}_H$ is $h(a,a')$; and $\iota^{\mathsf T}J\,c(\iota)=J_1$,
$h(\iota u,x)=h(\iota u,y)=0$, hence $h(x,\iota u)=h(y,\iota u)=0$,
$h(x,x)=h(y,y)=0$, $h(x,y)=1$ and $h(y,x)=\eta$. Finally
$c(J_1)^{\mathsf T}=\iota^{\mathsf T}c(J)^{\mathsf T}c(\iota)=\eta J_1$.

(4) Put $A=\mathrm{fr}_H\Delta(r_0,r_x,r_y)Q$. Then $QA=\Delta(r_0,r_x,r_y)Q$,
and for $u\in L$,
$Au=r_x(Qu)_x\,x+\iota(r_0Q_1u)+r_y(Qu)_y\,y$ lies in $L$ by (3), as
$\iota L_1\subset L$.
\end{proof}
% Lean: FrOK collects (1)-(4) over the adeles, through incM; pq, qx, qy, qi are (1), gram,
% herm, herm1 are (2), low is (3), msk is (4) with mask.

\begin{lemma}\label{asm-theta-agree-evaluation-adapted-vacuum}
Let $B$ be a theta base for a CM field $E$, let $X$ be a Hermitian and $W$ a
skew-Hermitian lattice space, of ranks $n$ and $d$ and of some signatures, and
let $X_1$ and $W_1$ be lattice spaces of ranks $n_1$ and $d_1$. Put
$I=\{1,\ldots,n\}\times\{1,\ldots,d\}$ and
$I_1=\{1,\ldots,n_1\}\times\{1,\ldots,d_1\}$, let $\kappa_0$, $\epsilon$ and
$[p,y]$ be as in Lemma~\ref{asm-theta-agree-evaluation-glue}, let
$\operatorname{Tr}=\operatorname{Tr}_{\mathbf A_E/\mathbf A}$, and let $h_1$ be
the form $h$ of Definition~\ref{asm-theta-agree-models}(3) for $(X_1,W_1)$. Let
$e^a\colon\mathbb W(X,W)\to\mathbf A^I\times\mathbf A^I$ be an
$\mathbf A$-linear isomorphism with
$\langle w,w'\rangle=e^a_1(w)\cdot e^a_2(w')-e^a_1(w')\cdot e^a_2(w)$ that maps
$M_{n\times d}(E)$ into $F^I\times F^I$, and let $M_1$ be a model for
$(X_1,W_1)$ with isomorphism $e^1$. Suppose there are maps
$p,q\colon\mathbb W(X,W)\to\mathbf A^{\kappa_0}$ and
$\pi\colon\mathbb W(X,W)\to\mathbb W(X_1,W_1)$ with
$e^a(u)=\bigl([p(u),e^1_1(\pi u)],[q(u),e^1_2(\pi u)]\bigr)$ for every $u$. Let
$v$ be a finite place, and let $\Lambda_v$ and $\Lambda^1_v$ be the groups
generated by the matrices $a\,xy^{\mathsf T}$ with $a\in\mathbf A$ at $v$,
$a_v\in\mathcal O_v$, and $x$ and $y$ in the lattices of $X$ and $W$,
respectively of $X_1$ and $W_1$. Suppose that:
\begin{enumerate}
\item $\pi(\Lambda_v)\subset\Lambda^1_v$;
\item $\psi(p(\gamma')\cdot q(\gamma))=1$ for $\gamma,\gamma'\in\Lambda_v$;
\item for $\theta\in\mathcal O_E$ with $\theta+c(\theta)=1$ and
  $\gamma\in\Lambda_v$,
  \[
  \psi\bigl(\tfrac12\operatorname{Tr}h(\theta\gamma,\gamma)\bigr)
  =\psi\bigl(\tfrac12p(\gamma)\cdot q(\gamma)\bigr)\,
  \psi\bigl(\tfrac12\operatorname{Tr}h_1(\theta\pi\gamma,\pi\gamma)\bigr).
  \]
\end{enumerate}
Let $\phi$ be a vacuum of $M_1$ at $v$. Then there is a function $\phi_a$ on
$F_v^I$ that is a vacuum at $v$ for the models of $(X,W)$ with the isomorphism
$e^a$ (Definition~\ref{asm-theta-agree-models}(3)), with
$\phi_a([0,y])=\phi(y)$ for $y\in F_v^{I_1}$.
\end{lemma}
\leanref{Development/WeilBound/ThetaLifts/LeviDeterminants/MixedModels/AdaptC.lean}{PositivityTheory.Mix.Adapt.exists_isVac}
\leanref{Development/AbelianVarietyTheory/Proofs/DivisorProofs/HermitianTensorSkewHermitianLatticeOpenCompact.lean}{GlobalTheory.LatOpenCpt.main}
\begin{proof}
Let $\rho_v$ be the Schrödinger operators on functions on $F_v^I$, and on
$F_v^{I_1}$, for $\psi_v$, as in the proof of
Lemma~\ref{asm-theta-agree-evaluation-vacuum}; a subscript $v$ marks the
component at $v$.

\emph{A compact open subgroup.} Let $K$ be the set of the $e^a(\gamma)_v$,
$\gamma\in\Lambda_v$. By Lemma~\ref{asm-theta-agree-tower-cocycle}(1), $K$ is the
$\mathcal O_v$-span of the vectors $e^a(x_iy_j^{\mathsf T})_v$, for generators
$x_i$ and $y_j$ of the lattices of $X$ and $W$. The group generated by the
$x_iy_j^{\mathsf T}$ is an $\mathcal O_E$-module that spans $M_{n\times d}(E)$
over $E$, hence over $F$, and $e^a$ maps $M_{n\times d}(E)$ isomorphically onto
$F^I\times F^I$, as in the proof of
Lemma~\ref{asm-theta-agree-evaluation-rational}. So $K$ is a compact open
subgroup of $F_v^I\times F_v^I$. Let $\Lambda\subset F_v^{\kappa_0}$ be the set of
the $p(\gamma)_v$, $\gamma\in\Lambda_v$, the image of $K$ under the projection to
the coordinates $\kappa_0$ of the first component. It is compact, and it is open,
as it contains every $a$ with $([a,0],0)\in K$.

\emph{The function.} Put $\phi_a([a,y])=\mathbf 1_\Lambda(a)\phi(y)$. It is a
Schwartz-Bruhat function, $\phi_a([0,y])=\phi(y)$, and $\phi_a\neq0$ as
$\phi\neq0$. Let $\theta$ be the element in the definition of the vacuum $\phi$,
and $\gamma\in\Lambda_v$. As $\gamma$ is at $v$, so are $e^a(\gamma)$, $p(\gamma)$
and $q(\gamma)$, and $\pi(\gamma)\in\Lambda^1_v$ by (1). So $\rho_a(\gamma)$, the
Schrödinger operator of $e^a$, acts at $v$ by $\phi_a\mapsto\rho_v(e^a(\gamma)_v)\phi_a$.
Write $p=p(\gamma)_v$, $q=q(\gamma)_v$ and $(e_1,e_2)=e^1(\pi\gamma)_v$. As the dot
products split, $[a,y]\cdot[q,e_2]=a\cdot q+y\cdot e_2$, and
\[
\bigl(\rho_v(e^a(\gamma)_v)\phi_a\bigr)([a,y])
=\psi_v\bigl(a\cdot q+\tfrac12p\cdot q\bigr)\mathbf 1_\Lambda(a+p)\cdot
\psi_v\bigl(y\cdot e_2+\tfrac12e_1\cdot e_2\bigr)\phi(y+e_1).
\]
If $a\notin\Lambda$, then $a+p\notin\Lambda$, as $p\in\Lambda$, and both this and
$\phi_a([a,y])$ vanish. If $a\in\Lambda$, then $a=p(\gamma')_v$ for some
$\gamma'\in\Lambda_v$, so $\psi_v(a\cdot q)=\psi(p(\gamma')\cdot q(\gamma))=1$ by
(2); $\mathbf 1_\Lambda(a+p)=1$; and the last factor is
$\psi(\frac12\operatorname{Tr}h_1(\theta\pi\gamma,\pi\gamma))\phi(y)$, as $\phi$
is a vacuum of $M_1$ at $v$. By (3) the product is
$\psi(\frac12\operatorname{Tr}h(\theta\gamma,\gamma))\phi_a([a,y])$. So $\phi_a$
is a vacuum at $v$, with the same $\theta$.
\end{proof}
% Lean: hKo, hKc are LatOpenCpt.main; Lambda is K.map pr (hmem); phi_a is box eps Lambda phi
% (box_mem_SB, h0); the computation is the case split on a in Lambda, with hat
% (latAtV_at_stable) for the components away from v and hR1 for the action of M1.

\begin{lemma}\label{asm-theta-agree-evaluation-adapted-model}
Let $B$ be a theta base for a CM field $E$, let $X$ be a Hermitian and $W$ a
skew-Hermitian lattice space, of ranks $n$ and $d$ and of some signatures, let
$M$ be a regular model for $(X,W)$, let $X_1$ and $W_1$ be lattice spaces of
ranks $n_1$ and $d_1$, let $M_1$ be a regular model for $(X_1,W_1)$ with
isomorphism $e^1$, and let $\iota_X$, $\iota_W$ and $\mathcal J$ be as in
Definition~\ref{asm-theta-agree-planes}(5). Let $I$, $I_1$, $\operatorname{Tr}$,
$h_1$, $\Lambda_v$ and $\Lambda^1_v$ be as in
Lemma~\ref{asm-theta-agree-evaluation-adapted-vacuum}. Let $\kappa$ be a finite
set with $nd=2|\kappa|+n_1d_1$, and let
$\beta,\delta\colon\mathbb W(X,W)\to\mathbf A_E^\kappa$ and
$\pi\colon\mathbb W(X,W)\to\mathbb W(X_1,W_1)$ be $\mathbf A$-linear maps such
that:
\begin{enumerate}
\item $h(z,z')=\beta(z)\cdot\delta(z')-c\bigl(\beta(z')\cdot\delta(z)\bigr)+h_1(\pi z,\pi z')$
  for all $z,z'$;
\item $\beta$ and $\delta$ map $M_{n\times d}(E)$ into $E^\kappa$, and $\pi$ maps
  it into $M_{n_1\times d_1}(E)$;
\item $\beta(\iota_Xz\iota_W^{\mathsf T})=0$ and $\pi(\iota_Xz\iota_W^{\mathsf T})=z$
  for $z\in\mathbb W(X_1,W_1)$, and $\beta(u)=0$ and $\pi(u)=0$ for
  $u\in\mathcal J$;
\item for every finite place $v$ and $\gamma\in\Lambda_v$, $\pi(\gamma)\in\Lambda^1_v$,
  and there is $\gamma_0\in\Lambda_v$ with $\beta(\gamma_0)=0$, $\pi(\gamma_0)=0$
  and $\delta(\gamma_0)=\delta(\gamma)$.
\end{enumerate}
Then there are a bijection $\epsilon\colon I\to(\kappa\times\{0,1\})\sqcup I_1$,
with $[p,y]$ as in Lemma~\ref{asm-theta-agree-evaluation-glue}, maps
$p,q\colon\mathbb W(X,W)\to\mathbf A^{\kappa\times\{0,1\}}$ with $p(u)=0$ when
$\beta(u)=0$ and $q(u)=0$ when $\delta(u)=0$, and a regular model $M_a$ for
$(X,W)$, with the measure and the fundamental domain of $M$, whose isomorphism is
$e^a(u)=\bigl([p(u),e^1_1(\pi u)],[q(u),e^1_2(\pi u)]\bigr)$ and whose vacua
satisfy $\mathrm{vac}^a_v([0,y])=\mathrm{vac}^1_v(y)$. By
Lemma~\ref{asm-theta-agree-evaluation-glue}, $e(\Phi)(y)=\Phi([0,y])$ is then an
evaluation map for $(M_a,M_1,\iota_X,\iota_W,\mathcal J)$, with components
$e_\infty(f)(y)=f([0,y])$ and $e_v(\phi)(y)=\phi([0,y])$.
\end{lemma}
\leanref{Development/WeilBound/ThetaLifts/LeviDeterminants/EvModel.lean}{PositivityTheory.ThKF0.exists_modelS}
\leanref{Development/WeilBound/ThetaLifts/LeviDeterminants/MixedModels/AdaptC.lean}{PositivityTheory.Mix.Adapt.exists_model}
\leanref{Development/WeilBound/ThetaLifts/LeviDeterminants/MixedModels/AdaptC.lean}{PositivityTheory.Mix.Adapt.bijective_of_symp}
\begin{proof}
As $|I|=2|\kappa|+|I_1|$, there is a bijection $\epsilon$. Let $(o_0,o_1)$ be a
basis of $E$ over $F$; it is a basis of $\mathbf A_E$ over $\mathbf A$. For $z$
define $p(z)$ by $\beta(z)_k=p(z)_{(k,0)}o_0+p(z)_{(k,1)}o_1$, and put
$q(z)_{(k,j)}=-\operatorname{Tr}(o_j\delta(z)_k)$. These are $\mathbf A$-linear,
$p(z)=0$ when $\beta(z)=0$, $q(z)=0$ when $\delta(z)=0$, and
$p(z)\cdot q(z')=-\operatorname{Tr}(\beta(z)\cdot\delta(z'))$.

\emph{The isomorphism.} Define $e^a$ by the formula in the statement. It is
$\mathbf A$-linear. The dot products split along $\epsilon$,
$\operatorname{Tr}\circ c=\operatorname{Tr}$, and
$e^1_1(\pi z)\cdot e^1_2(\pi z')-e^1_1(\pi z')\cdot e^1_2(\pi z)=-\operatorname{Tr}h_1(\pi z,\pi z')$;
so (1) gives
\[
e^a_1(z)\cdot e^a_2(z')-e^a_1(z')\cdot e^a_2(z)
=-\operatorname{Tr}(\beta(z)\cdot\delta(z'))+\operatorname{Tr}(\beta(z')\cdot\delta(z))-\operatorname{Tr}h_1(\pi z,\pi z')
=\langle z,z'\rangle .
\]
Let $e^M$ be the isomorphism of $M$, and $T=e^a\circ(e^M)^{-1}$, an endomorphism
of $\mathbf A^I\times\mathbf A^I$ that preserves the form
$[w,w']=w_1\cdot w'_2-w'_1\cdot w_2$. With $e_k$ the unit vectors, put
$S(w)=\bigl(([w,T(0,e_k)])_k,([T(e_k,0),w])_k\bigr)$. Then $S(T(w))=w$, as $T$
preserves the form. As $ST=1$ for endomorphisms of a free module of finite rank
over a commutative ring, $T$ is bijective, and so is $e^a$. For
$z\in M_{n\times d}(E)$, $p(z)$ and
$q(z)_{(k,j)}=-\operatorname{Tr}_{E/F}(o_j\delta(z)_k)$ are rational by (2), and
so is $e^1(\pi z)$, as $M_1$ is regular; so $e^a$ maps $M_{n\times d}(E)$ into
$F^I\times F^I$.

\emph{A vanishing.} Let $\gamma,\gamma'\in\Lambda_v$ and $\theta\in\mathcal O_E$.
We show $\psi(\operatorname{Tr}(\theta\beta(\gamma')\cdot\delta(\gamma)))=1$. Let
$\gamma_0$ be as in (4). By (1),
$h(\gamma',\gamma_0)=\beta(\gamma')\cdot\delta(\gamma)$, and
Lemma~\ref{asm-theta-agree-evaluation-lattice}(1), with $o=\theta$ and $b=1$,
gives $\psi(\operatorname{Tr}(\theta h(\gamma',\gamma_0)))=1$.

\emph{Vacua.} Fix $v$. We check the hypotheses of
Lemma~\ref{asm-theta-agree-evaluation-adapted-vacuum} for $e^a$, $p$, $q$, $\pi$
and $\kappa_0=\kappa\times\{0,1\}$. Its (1) is (4). Its (2) is the vanishing with
$\theta=-1$, as $p(\gamma')\cdot q(\gamma)=\operatorname{Tr}(-\beta(\gamma')\cdot\delta(\gamma))$.
For its (3), let $\theta+c(\theta)=1$, $\gamma\in\Lambda_v$, and
$m=\beta(\gamma)\cdot\delta(\gamma)$. By (1),
$h(\gamma,\gamma)=m-c(m)+h_1(\pi\gamma,\pi\gamma)$, and
$\operatorname{Tr}(\theta c(m))=\operatorname{Tr}(c(\theta)m)=\operatorname{Tr}(m)-\operatorname{Tr}(\theta m)$,
so
\[
\tfrac12\operatorname{Tr}h(\theta\gamma,\gamma)=\operatorname{Tr}(\theta m)+\tfrac12p(\gamma)\cdot q(\gamma)+\tfrac12\operatorname{Tr}h_1(\theta\pi\gamma,\pi\gamma),
\]
and $\psi(\operatorname{Tr}(\theta m))=1$ by the vanishing. So there is a vacuum
$\phi_v$ at $v$ for the isomorphism $e^a$ with $\phi_v([0,y])=\mathrm{vac}^1_v(y)$.

\emph{The model.} Let $\mathrm{vac}^a_v$ be the indicator function of
$\mathcal O_v^I$ if it is a vacuum at $v$ for $e^a$ and $\mathrm{vac}^1_v$ is
the indicator function of $\mathcal O_v^{I_1}$, and $\phi_v$ otherwise. In both
cases $\mathrm{vac}^a_v([0,y])=\mathrm{vac}^1_v(y)$. By
Lemma~\ref{asm-theta-agree-tower-vacuum}(2) for $e^a$, and as $M_1$ is regular,
the first case holds for almost all $v$. With the measure and the fundamental
domain of $M$, this is a regular model $M_a$. By (3),
$p(\iota_Xz\iota_W^{\mathsf T})=0$, $\pi(\iota_Xz\iota_W^{\mathsf T})=z$, and
$p(u)=\pi(u)=0$ for $u\in\mathcal J$, so
Lemma~\ref{asm-theta-agree-evaluation-glue} applies.
\end{proof}
% Lean: exists_modelS and exists_model have the same proof: hpq, hPhi (with M1.co.form and
% trace_conj), hbij (bijective_of_symp, through OrzechProperty), hr, hdelta
% (LatVacArith.psi_symp_latAtV), hv (exists_isVac), then Ma with the choice G v and
% VacIndic.eventually_isVac_indicator, and exists_isEvS.

\begin{lemma}\label{asm-theta-agree-evaluation-adapted-signs}
Let $B$ be a theta base for a CM field $E$, and let $\varepsilon=(\varepsilon_w)_w$
be signs indexed by the infinite places of $F$. For lattice spaces $X$ and $Y$ of
ranks $n$ and $m$, with matrices $\theta_{X,w}$ and $\theta_{Y,w}$, and
$u,u'\in M_{n\times m}(\mathbf A_E)$, write $u\sim_\varepsilon u'$ if
$\pi_w(u')=\varepsilon_wi\,\theta_{X,w}\pi_w(u)\theta_{Y,w}^{\mathsf T}$ for every
$w$. Let $X$ be a Hermitian and $W$ a skew-Hermitian lattice space, of ranks $n$
and $d$ and of some signatures, and for $u,u'\in\mathbb W(X,W)$ put
$\omega(u,u')=\operatorname{Tr}_{F_\infty/\mathbf R}\langle u,u'\rangle_\infty$,
with $\langle\ ,\ \rangle$ as in Definition~\ref{asm-theta-agree-models}(3).
\begin{enumerate}
\item If $u\sim_\varepsilon u'$ and $u_1\sim_\varepsilon u'_1$ in $\mathbb W(X,W)$,
  then $u+u_1\sim_\varepsilon u'+u'_1$, $u'\sim_\varepsilon-u$, and
  $\omega(u',u'_1)=\omega(u,u_1)$.
\item For every $u\in M_{n\times m}(\mathbf A_E)$ there is $u'$ at $\infty$ with
  $u\sim_\varepsilon u'$; and if $u'$ and $u''$ are at $\infty$ with
  $u\sim_\varepsilon u'$ and $u\sim_\varepsilon u''$, then $u'=u''$.
\item Let $a\in\mathbf A_E^n$ and $b\in E^d$ with $b\neq0$, and let
  $ab^{\mathsf T}\sim_\varepsilon u'$. Then
  $\omega(ab^{\mathsf T},u')=\sum_w\varepsilon_wp_w$ with real numbers $p_w$ such
  that $p_w=0$ if $\pi_w(a)=0$ and $p_w>0$ otherwise.
\item Let $H'=(x',y',\iota',\mathfrak b')$ be an adapted plane in $W$ over a
  lattice space $W_1$. If $z\sim_\varepsilon z'$ in $\mathbb W(X,W_1)$, then
  $z\iota'^{\mathsf T}\sim_\varepsilon z'\iota'^{\mathsf T}$. For $a\in\mathbf A_E^n$
  there is $b$ at $\infty$ with $ax'^{\mathsf T}\sim_\varepsilon by'^{\mathsf T}$,
  and for $b\in\mathbf A_E^n$ there is $a$ at $\infty$ with
  $by'^{\mathsf T}\sim_\varepsilon ax'^{\mathsf T}$.
\end{enumerate}
\end{lemma}
\leanref{Development/WeilBound/ThetaLifts/LeviDeterminants/CartanPos.lean}{PositivityTheory.ThKF0.Rel}
\leanref{Development/WeilBound/ThetaLifts/LeviDeterminants/CartanPos.lean}{PositivityTheory.ThKF0.Rel.symm}
\leanref{Development/WeilBound/ThetaLifts/LeviDeterminants/CartanPos.lean}{PositivityTheory.ThKF0.omega_rel}
\leanref{Development/WeilBound/ThetaLifts/LeviDeterminants/CartanPos.lean}{PositivityTheory.ThKF0.exists_rel}
\leanref{Development/WeilBound/ThetaLifts/LeviDeterminants/CartanPos.lean}{PositivityTheory.ThKF0.Rel.unique}
\leanref{Development/WeilBound/ThetaLifts/LeviDeterminants/CartanPos.lean}{PositivityTheory.ThKF0.omega_pure}
\leanref{Development/WeilBound/ThetaLifts/LeviDeterminants/ExistsRelationAdaptedVector.lean}{PositivityTheory.ThKF0.rel_iota}
\leanref{Development/WeilBound/ThetaLifts/LeviDeterminants/ExistsRelationAdaptedVector.lean}{PositivityTheory.ThKF0.exists_relation_adapted_vector}
\leanref{Development/WeilBound/ThetaLifts/LeviDeterminants/ExistsRelationAdaptedVector.lean}{PositivityTheory.ThKF0.rel_y}
\begin{proof}
Write $\theta=\theta_{X,w}$ and $\theta'=\theta_{W,w}$, $A=\tau_w(J_X)$ and
$C=\tau_w(J_W)$, and let $H_w(x,y)=\operatorname{tr}(x^{\mathsf T}A\,\bar yC^{\mathsf T})$
for $x,y\in M_{n\times d}(\mathbf C)$. As in the proof of
Lemma~\ref{asm-theta-agree-density-gaussian},
$\omega(u,u')=-2\sum_w\operatorname{Re}H_w(\pi_wu,\pi_wu')$.

(1) The first claim is clear. As $\theta^2=1$ and $\theta'^2=1$, applying the
relation twice multiplies by $(\varepsilon_wi)^2=-1$, which gives the second.
For the third,
$H_w(\varepsilon_wi\,\theta x\theta'^{\mathsf T},\varepsilon_wi\,\theta y\theta'^{\mathsf T})=H_w(x,y)$,
by part (a) of the proof of Lemma~\ref{asm-theta-agree-density-gaussian} and
$(\varepsilon_wi)\overline{(\varepsilon_wi)}=1$.

(2) An element of $\mathbf A_E$ at $\infty$ lies in
$E\otimes_FF_\infty=\prod_wE\otimes_FF_w$, which the $\pi_w$ identify with
$\prod_w\mathbf C$. So there is a unique $u'$ at $\infty$ with prescribed
$\pi_w(u')$ for every $w$.

(3) Put $x_w=\pi_w(a)$ and $b_w=\tau_w(b)$, which is nonzero. Then
$\pi_w(ab^{\mathsf T})=x_wb_w^{\mathsf T}$,
$\pi_w(u')=\varepsilon_wi\,(\theta x_w)(\theta'b_w)^{\mathsf T}$, and
$H_w(xb^{\mathsf T},x'b'^{\mathsf T})=(x^{\mathsf T}A\bar x')(b^{\mathsf T}C\bar b')$.
By part (c) of the proof of Lemma~\ref{asm-theta-agree-density-gaussian},
$A\bar\theta$ and $iC\bar\theta'$ are positive definite Hermitian, so
$P_w=x_w^{\mathsf T}A\,\overline{\theta x_w}$ is real, $\geq0$, and $>0$ if
$x_w\neq0$, and $b_w^{\mathsf T}C\,\overline{\theta'b_w}=-iq_w$ with $q_w>0$.
Hence
$H_w(\pi_w(ab^{\mathsf T}),\pi_wu')=\overline{\varepsilon_wi}\,P_w(-iq_w)=-\varepsilon_wP_wq_w$,
and $\omega(ab^{\mathsf T},u')=\sum_w\varepsilon_w\,2P_wq_w$.

(4) As $\theta_{W,w}\tau_w(\iota')=\tau_w(\iota')\theta_{W_1,w}$,
$\theta_{W_1,w}^{\mathsf T}\tau_w(\iota')^{\mathsf T}=\tau_w(\iota')^{\mathsf T}\theta_{W,w}^{\mathsf T}$,
which gives the first claim. By Definition~\ref{asm-theta-agree-planes}(1),
$\tau_w(y')=\lambda_w\theta_{W,w}\tau_w(x')$ with $\lambda_w\neq0$, so
$\theta_{W,w}\tau_w(x')=\lambda_w^{-1}\tau_w(y')$ and
$\theta_{W,w}\tau_w(y')=\lambda_w\tau_w(x')$. By (2) choose $b$ at $\infty$ with
$\pi_w(b)=\varepsilon_wi\lambda_w^{-1}\theta_{X,w}\pi_w(a)$, and $a$ at $\infty$
with $\pi_w(a)=\varepsilon_wi\lambda_w\theta_{X,w}\pi_w(b)$.
\end{proof}
% Lean: Rel B X W e u u' is the relation; (1) is Rel.add, Rel.symm and omega_rel (omega_eq,
% hTens_smul, hTens_mul); (2) is exists_rel (exists_vec) and Rel.unique (at_ext); (3) is
% omega_pure with cartan_pos; (4) is rel_iota, exists_relation_adapted_vector and rel_y,
% through rel_pure.

\begin{lemma}\label{asm-theta-agree-evaluation-adapted-commutator}
Let $V$ be a finite-dimensional real vector space. For $a\in V$, a real linear
form $L$ on $V$ and $f\in\mathcal S(V)$, put
$D_{a,L}f=\partial_af-2\pi i\,Lf$. Let $f\in\mathcal S(V)$, $a,a'\in V$ and $L,L'$
real linear forms with $D_{a',L'}f=i\,D_{a,L}f$ and $L'(a)-L(a')<0$. Then $f=0$.
\end{lemma}
\leanref{Development/WeilBound/ThetaLifts/LeviDeterminants/NormIteratedFDerivIdLeOneAdd.lean}{PositivityTheory.ThKF0.Lib.eq_zero_of_neg}
\begin{proof}
Let $\langle f,g\rangle=\int_V\bar fg$ for a Haar measure, and write $D=D_{a,L}$,
$D'=D_{a',L'}$. Integration by parts, and $L$ real, give
$\langle f,Dg\rangle=-\langle Df,g\rangle$, and likewise for $D'$. The product
rule gives $DD'f-D'Df=-2\pi i(L'(a)-L(a'))f$. Put $g=Df$, so that $D'f=ig$. Then
$\langle f,DD'f\rangle=-\langle g,ig\rangle=-i\langle g,g\rangle$ and
$\langle f,D'Df\rangle=-\langle ig,g\rangle=i\langle g,g\rangle$. Subtracting,
$\langle g,g\rangle=\pi(L'(a)-L(a'))\langle f,f\rangle$. The left side is
$\geq0$ and the right side $\leq0$, so $\langle f,f\rangle=0$ and $f=0$.
\end{proof}

\begin{lemma}\label{asm-theta-agree-evaluation-adapted-fock}
Let $B$ be a theta base for a CM field $E$, let $X$ be a Hermitian and $W$ a
skew-Hermitian lattice space, of ranks $n$ and $d$ and of some signatures, let
$H'=(x',y',\iota',\mathfrak b')$ be an adapted plane in $W$ over a lattice space
$W_1$ of rank $d_1$, and let $M_a$ be a model for $(X,W)$ and $M_1$ one for
$(X,W_1)$, with isomorphisms $e^a$ and $e^1$. Put
$I=\{1,\ldots,n\}\times\{1,\ldots,d\}$ and
$I_1=\{1,\ldots,n\}\times\{1,\ldots,d_1\}$, and let $\kappa_0$, $\epsilon$ and
$[t,y]$ be as in Lemma~\ref{asm-theta-agree-evaluation-glue}. Suppose that
$e^a(z\iota'^{\mathsf T})=([0,e^1_1(z)],[0,e^1_2(z)])$ for
$z\in\mathbb W(X,W_1)$, and that for every $a\in\mathbf A_E^n$ there are
$p,q\in\mathbf A^{\kappa_0}$ with $e^a(ax'^{\mathsf T})=(0,[q,0])$ and
$e^a(ay'^{\mathsf T})=([p,0],0)$. For $f\in\mathcal S(F_\infty^I)$ put
$Rf(y)=f([0,y])$. Then $R(\operatorname{Pol}(M_a))=\operatorname{Pol}(M_1)$.
\end{lemma}
\leanref{Development/WeilBound/ThetaLifts/LeviDeterminants/ExistsRelationAdaptedVector.lean}{PositivityTheory.ThKF0.Adp.map_eq}
\leanref{Development/WeilBound/ThetaLifts/LeviDeterminants/ExistsRelationAdaptedVector.lean}{PositivityTheory.ThKF0.Adp.map_le}
\leanref{Development/WeilBound/ThetaLifts/LeviDeterminants/ExistsRelationAdaptedVector.lean}{PositivityTheory.ThKF0.Adp.gs_one}
\leanref{Development/WeilBound/ThetaLifts/LeviDeterminants/ExistsRelationAdaptedVector.lean}{PositivityTheory.ThKF0.Adp.exists_lift}
\leanref{Development/WeilBound/ThetaLifts/LeviDeterminants/ExistsGaussBox.lean}{PositivityTheory.ThKF0.exists_gauss_box}
\begin{proof}
For $u\in\mathbb W(X,W)$ at $\infty$, $e^a(u)$ is at $\infty$; let $D_u=D_{a,b}$ be
the operator of Lemma~\ref{asm-theta-agree-density-derivative} for the components
$(a,b)$ of $e^a(u)$ at $\infty$, and define $D^1_z$ likewise for $M_1$ and
$z\in\mathbb W(X,W_1)$ at $\infty$. Then $D_{u+u'}=D_u+D_{u'}$ and $D_{-u}=-D_u$.
Let $\sim_\varepsilon$ and $\omega$ be as in
Lemma~\ref{asm-theta-agree-evaluation-adapted-signs}, and
$T(x,b)=\operatorname{Tr}_{F_\infty/\mathbf R}(\sum_kx_kb_k)$ as in
Lemma~\ref{asm-theta-agree-density-derivative}, on any $F_\infty^\kappa$. As in
Step~6 of the proof of Lemma~\ref{asm-theta-agree-evaluation-pure}, a Gaussian
for $M_a$ with signs $\varepsilon$ is an $f$ with $D_{u'}f=i\,D_uf$ whenever
$u\sim_\varepsilon u'$ are at $\infty$, and $\operatorname{Pol}(M_a)$ is the
smallest subspace that contains the Gaussians and is stable under the $D_u$; the
same holds for $M_1$. By the condition on $e^a$,
$\omega(u,u')=T(a,b')-T(a',b)$ for $u,u'$ at $\infty$ with components $(a,b)$ and
$(a',b')$ of $e^a(u)$ and $e^a(u')$ at $\infty$. Points of $F_\infty^I$ are
$[t,y]$ with $t\in F_\infty^{\kappa_0}$ and $y\in F_\infty^{I_1}$.

\emph{Five facts.}
\begin{enumerate}
\item[(a)] For $z\in\mathbb W(X,W_1)$ at $\infty$,
  $(D_{z\iota'^{\mathsf T}}f)([t,y])=D^1_z(f([t,\cdot]))(y)$: the components of
  $e^a(z\iota'^{\mathsf T})$ at $\infty$ are $[0,a]$ and $[0,b]$, with $(a,b)$
  those of $e^1(z)$, and $T([t,y],[0,b])=T(y,b)$.
\item[(b)] For $a$ at $\infty$ with $e^a(ax'^{\mathsf T})=(0,[q,0])$,
  $(D_{ax'^{\mathsf T}}f)([t,y])=-2\pi i\,T(t,q)f([t,y])$.
\item[(c)] For $b$ at $\infty$ with $e^a(by'^{\mathsf T})=([p,0],0)$,
  $D_{by'^{\mathsf T}}f=\partial_{[p,0]}f$.
\item[(d)] Every $u$ at $\infty$ is $ax'^{\mathsf T}+z\iota'^{\mathsf T}+by'^{\mathsf T}$
  with $a$, $z$ and $b$ at $\infty$: with $Q'$ the inverse of $\mathrm{fr}_{H'}$
  (Lemma~\ref{asm-theta-agree-evaluation-adapted-frame}), $u=(uQ'^{\mathsf T})\mathrm{fr}_{H'}^{\mathsf T}$.
\item[(e)] For $t\in F_\infty^{\kappa_0}$ there is $b$ at $\infty$ with
  $e^a(by'^{\mathsf T})=([t,0],0)$: write the $u$ at $\infty$ with
  $e^a(u)=([t,0],0)$ as in (d), apply $e^a$ and compare the coordinates
  $\kappa_0$ of the first components.
\end{enumerate}

\emph{Step 1: $R(\operatorname{Pol}(M_a))\subset\operatorname{Pol}(M_1)$.} Let
$\mathcal T$ be the set of the $f$ for which there is $g\in\operatorname{Pol}(M_1)$
such that every $f([t,\cdot])$ is a multiple of $g$. Then $R$ maps
$\mathcal T$ into $\operatorname{Pol}(M_1)$, so it suffices to show that the span
of $\mathcal T$ contains the Gaussians for $M_a$ and is stable under the $D_u$.

Let $f$ be a Gaussian for $M_a$ with signs $\varepsilon$, and $g=Rf$. If
$z\sim_\varepsilon z'$ are at $\infty$, then
$z\iota'^{\mathsf T}\sim_\varepsilon z'\iota'^{\mathsf T}$ by
Lemma~\ref{asm-theta-agree-evaluation-adapted-signs}(4), and (a) at $t=0$ gives
$D^1_{z'}g=i\,D^1_zg$. So $g$ is a Gaussian for $M_1$, and
$g\in\operatorname{Pol}(M_1)$. Let $t\in F_\infty^{\kappa_0}$; choose $b$ by (e),
$a$ at $\infty$ with $by'^{\mathsf T}\sim_\varepsilon ax'^{\mathsf T}$ by
Lemma~\ref{asm-theta-agree-evaluation-adapted-signs}(4), and $q$ with
$e^a(ax'^{\mathsf T})=(0,[q,0])$. By (b) and (c), the Gaussian condition
$D_{ax'^{\mathsf T}}f=i\,D_{by'^{\mathsf T}}f$ at $[st,y]$, $s\in\mathbf R$,
reads $\frac d{ds}f([st,y])=-2\pi s\,T(t,q)f([st,y])$. So
$s\mapsto e^{\pi s^2T(t,q)}f([st,y])$ is constant, and
$f([t,y])=e^{-\pi T(t,q)}g(y)$. Hence $f\in\mathcal T$.

Let $f\in\mathcal T$, with $f([t,\cdot])=c_tg$, and let $u$ be at $\infty$; write
$u=ax'^{\mathsf T}+z\iota'^{\mathsf T}+by'^{\mathsf T}$ by (d). By (b),
$(D_{ax'^{\mathsf T}}f)([t,\cdot])=-2\pi i\,T(t,q)c_tg$. By (a),
$(D_{z\iota'^{\mathsf T}}f)([t,\cdot])=c_tD^1_zg$, and
$D^1_zg\in\operatorname{Pol}(M_1)$. By (c), $D_{by'^{\mathsf T}}f=\partial_{[p,0]}f$;
if $g=0$ it vanishes, and otherwise, with $g(y_0)\neq0$,
$f([t',y])=f([t',y_0])g(y)/g(y_0)$ for all $t'$ and $y$, so
$\partial_{[p,0]}f([t,y])=\partial_{[p,0]}f([t,y_0])\,g(y)/g(y_0)$. So the three
terms of $D_uf$ lie in $\mathcal T$.

\emph{Step 2: $\operatorname{Pol}(M_1)\subset R(\operatorname{Pol}(M_a))$.} By (a)
at $t=0$, $D^1_z(Rf)=R(D_{z\iota'^{\mathsf T}}f)$, so
$R(\operatorname{Pol}(M_a))$ is stable under the $D^1_z$. Let $g\neq0$ be a
Gaussian for $M_1$ with signs $\varepsilon$; it remains to find
$f\in\operatorname{Pol}(M_a)$ with $Rf=g$.

\emph{The signs are $1$.} If $nd_1=0$, then $\mathbb W(X,W_1)=0$ and $g$ is a
Gaussian with signs $1$. Otherwise, suppose $\varepsilon_{w_0}=-1$. Let $e_j$ be
a unit vector of length $d_1$, let $a$ be at $\infty$ with $\pi_{w_0}(a)\neq0$ and
$\pi_w(a)=0$ for $w\neq w_0$, put $u=ae_j^{\mathsf T}$, and choose $z'$ at $\infty$
with $u\sim_\varepsilon z'$ by Lemma~\ref{asm-theta-agree-evaluation-adapted-signs}(2).
Then $D^1_{z'}g=i\,D^1_ug$. By
Lemma~\ref{asm-theta-agree-evaluation-adapted-signs}(4),
$u\iota'^{\mathsf T}\sim_\varepsilon z'\iota'^{\mathsf T}$, and
$u\iota'^{\mathsf T}=a(\iota'e_j)^{\mathsf T}$ with $\iota'e_j\neq0$, a column of
the invertible $\mathrm{fr}_{H'}$; so by part (3) of that lemma
$\omega(u\iota'^{\mathsf T},z'\iota'^{\mathsf T})=-p_{w_0}<0$. Let $(a_u,b_u)$
and $(a_{z'},b_{z'})$ be the components at $\infty$ of $e^1(u)$ and $e^1(z')$. By
the condition on $e^a(z\iota'^{\mathsf T})$,
$\omega(u\iota'^{\mathsf T},z'\iota'^{\mathsf T})=T(a_u,b_{z'})-T(a_{z'},b_u)$.
Lemma~\ref{asm-theta-agree-evaluation-adapted-commutator}, for
$D^1_u=D_{a_u,T(\cdot,b_u)}$ and $D^1_{z'}=D_{a_{z'},T(\cdot,b_{z'})}$, gives $g=0$,
a contradiction.

\emph{A form.} Write $\sim$ for $\sim_1$. Suppose $by'^{\mathsf T}\sim ax'^{\mathsf T}$
and $b'y'^{\mathsf T}\sim a'x'^{\mathsf T}$ are at $\infty$, with
$e^a(by'^{\mathsf T})=([p,0],0)$, $e^a(b'y'^{\mathsf T})=([p',0],0)$,
$e^a(ax'^{\mathsf T})=(0,[q,0])$ and $e^a(a'x'^{\mathsf T})=(0,[q',0])$. Then
$T(p,q')=T(p',q)$: by Lemma~\ref{asm-theta-agree-evaluation-adapted-signs}(1),
$a'x'^{\mathsf T}\sim-b'y'^{\mathsf T}$ and
$\omega(by'^{\mathsf T},a'x'^{\mathsf T})=\omega(ax'^{\mathsf T},-b'y'^{\mathsf T})$,
and the two sides are $T(p,q')$ and $T(p',q)$. For each $t$ choose $b_t$ by (e),
$a_t$ at $\infty$ with $b_ty'^{\mathsf T}\sim a_tx'^{\mathsf T}$, and $q_t$ with
$e^a(a_tx'^{\mathsf T})=(0,[q_t,0])$, and put $\mathsf Q(t,t')=T(t,q_{t'})$. It is
linear in $t$ and, by the last sentence, symmetric; so it is a symmetric bilinear
form. It is positive definite: for $t\neq0$, $b_t\neq0$, so $\pi_w(b_t)\neq0$ for
some $w$, and part (3) of that lemma, with $y'\neq0$, gives
$0<\omega(b_ty'^{\mathsf T},a_tx'^{\mathsf T})=T(t,q_t)=\mathsf Q(t,t)$.

\emph{The lift.} Put $f([t,y])=e^{-\pi\mathsf Q(t,t)}g(y)$. It is a Schwartz
function, as $\mathsf Q$ is positive definite; $Rf=g$; every $f([t,\cdot])$ is a
multiple of $g$; and $\partial_{[p,0]}f([t,y])=-2\pi\mathsf Q(t,p)f([t,y])$. We
show that $f$ is a Gaussian for $M_a$ with signs $1$.
\begin{enumerate}
\item[(f)] If $b'y'^{\mathsf T}\sim a'x'^{\mathsf T}$ are at $\infty$, with $p'$ and
  $q'$ as above, then by (b) and (c)
  $(D_{a'x'^{\mathsf T}}f)([t,y])=-2\pi i\,T(t,q')f([t,y])$ and
  $(D_{b'y'^{\mathsf T}}f)([t,y])=-2\pi\mathsf Q(t,p')f([t,y])$, and
  $T(t,q')=T(p',q_t)=\mathsf Q(t,p')$ by the form paragraph for $b_t$ and $b'$. So
  $D_{a'x'^{\mathsf T}}f=i\,D_{b'y'^{\mathsf T}}f$.
\item[(g)] If $z\sim z'$ are at $\infty$, then
  $D_{z'\iota'^{\mathsf T}}f=i\,D_{z\iota'^{\mathsf T}}f$, by (a), as
  $f([t,\cdot])=e^{-\pi\mathsf Q(t,t)}g$ and $D^1_{z'}g=i\,D^1_zg$.
\end{enumerate}
Let $u\sim u'$ be at $\infty$, and write
$u=a_1x'^{\mathsf T}+z\iota'^{\mathsf T}+b_1y'^{\mathsf T}$ by (d). By
Lemma~\ref{asm-theta-agree-evaluation-adapted-signs} choose, at $\infty$, $b_2$
with $a_1x'^{\mathsf T}\sim b_2y'^{\mathsf T}$, $z'$ with $z\sim z'$, so that
$z\iota'^{\mathsf T}\sim z'\iota'^{\mathsf T}$, and $a_2$ with
$b_1y'^{\mathsf T}\sim a_2x'^{\mathsf T}$. By parts (1) and (2) of that lemma,
$u'=b_2y'^{\mathsf T}+z'\iota'^{\mathsf T}+a_2x'^{\mathsf T}$. As
$b_2y'^{\mathsf T}\sim-a_1x'^{\mathsf T}$, (f) gives
$-D_{a_1x'^{\mathsf T}}f=i\,D_{b_2y'^{\mathsf T}}f$, that is
$D_{b_2y'^{\mathsf T}}f=i\,D_{a_1x'^{\mathsf T}}f$. With (f) for $(b_1,a_2)$ and
(g), $D_{u'}f=i\,D_uf$. So $f\in\operatorname{Pol}(M_a)$ and $Rf=g$.
\end{proof}
% Lean: the operators are Dp (dop of aOf and lOf); gauss_iff, pol_gs, pol_Dp and pol_le are
% the description of Pol; (a)-(e) are Adp.slice, cx, cy, decomp and onto, and omega_co is the
% formula for omega. Step 1 is map_le, with gs_slice and slice_ode for the Gaussians and
% slice_normal for D of b y'^T; Step 2 is gs_one (eq_zero_of_neg) and exists_lift, whose sym is
% the form paragraph and exists_gauss_box the lift.

\begin{lemma}\label{asm-theta-agree-evaluation-adapted}
Let $B$ be a theta base for a CM field $E$, let $X$ be a Hermitian and $W$ a
skew-Hermitian lattice space, of ranks $n$ and $d$ and of some signatures, and
let $M$ be a regular model for $(X,W)$.
\begin{enumerate}
\item Let $H'$ be an adapted plane in $W$ over a lattice space $W_1$
  (Definition~\ref{asm-theta-agree-planes}(1)), and $M_1$ a regular model for
  $(X,W_1)$. Then there are a regular model $M_a$ for $(X,W)$ and an evaluation
  map for $(M_a,M_1,1,\iota_{H'},\mathcal J_0(H'))$ with components $e_\infty$
  and $e_v$ such that $e_\infty(\operatorname{Pol}(M_a))=\operatorname{Pol}(M_1)$.
\item Let $H$ be an adapted plane in $X$ over a lattice space $X_1$ and $H'$ one
  in $W$ over a lattice space $W_1$, and let $M_1$ be a regular model for
  $(X_1,W_1)$. Then there are a regular model $M_a$ for $(X,W)$ and an
  evaluation map for $(M_a,M_1,\iota_H,\iota_{H'},\mathcal J_\alpha(H,H'))$, and
  a regular model $M'_a$ for $(X,W)$ and an evaluation map for
  $(M'_a,M_1,\iota_H,\iota_{H'},\mathcal J_\beta(H,H'))$.
\end{enumerate}
\end{lemma}
\leanref{Development/WeilBound/ThetaLifts/LeviDeterminants/ExistsAdaptedModelSkewPlane.lean}{PositivityTheory.ThKF0.exists_adapted_model_evaluation}
\leanref{Development/WeilBound/ThetaLifts/LeviDeterminants/EvStmt.lean}{PositivityTheory.ThKF0.AdaptedModelEvaluationAtZero}
\leanref{Development/WeilBound/ThetaLifts/LeviDeterminants/ExistsAdaptedModelTwoPlanes.lean}{PositivityTheory.ThKF0.adapted_plane_model_compatibility}
\leanref{Development/WeilBound/ThetaLifts/LeviDeterminants/EigenIntertwinerTwoPlanes.lean}{PositivityTheory.ThKF0.adapted_plane_model_evaluation_eq}
\begin{proof}
In each case the model is given by
Lemma~\ref{asm-theta-agree-evaluation-adapted-model}, for maps $\beta$, $\delta$
and $\pi$ read off from the coordinates of a matrix in the frames of the planes
(Lemma~\ref{asm-theta-agree-evaluation-adapted-frame}), and that lemma gives the
evaluation map, with $e_\infty(f)(y)=f([0,y])$. A plane in the Hermitian space
$X$ has $\eta=1$, and one in the skew-Hermitian space $W$ has $\eta=-1$. For a
matrix $w$ whose rows are indexed by $\{1,\ldots,m_1\}$, $x$, $y$ and whose
columns are indexed by $\{1,\ldots,m'_1\}$, $x'$, $y'$, write $w_{r,\cdot}$ for a
row and $w_{\cdot,s}$ for a column.

\emph{(1).} Let $H'=(x',y',\iota',\mathfrak b')$, let $Q'$ be the inverse of
$\mathrm{fr}_{H'}$, and let $d_1$ be the rank of $W_1$. For $z\in\mathbb W(X,W)$
put $w(z)=zQ'^{\mathsf T}$, with columns indexed by $\{1,\ldots,d_1\}$, $x'$ and
$y'$; let $\pi(z)$ be its first $d_1$ columns, $a(z)=w(z)_{\cdot,x'}$ and
$\beta(z)=w(z)_{\cdot,y'}$. As $\mathrm{fr}_{H'}Q'=1$,
\[
z=w(z)\,\mathrm{fr}_{H'}^{\mathsf T}=a(z)x'^{\mathsf T}+\pi(z)\iota'^{\mathsf T}+\beta(z)y'^{\mathsf T}.
\]
Put $\delta(z)=-J_Xc(a(z))$ and $\kappa=\{1,\ldots,n\}$, so that
$nd=2n+nd_1$. We check the hypotheses of
Lemma~\ref{asm-theta-agree-evaluation-adapted-model} for $X_1=X$, $\iota_X=1$,
$\iota_W=\iota'$ and $\mathcal J=\mathcal J_0(H')$.
\begin{enumerate}
\item Substituting $z=w(z)\mathrm{fr}_{H'}^{\mathsf T}$, $h(z,z')$ is the form $h$
  for the matrices $J_X$ and $\mathrm{fr}_{H'}^{\mathsf T}J_Wc(\mathrm{fr}_{H'})$
  at $w(z)$ and $w(z')$, and the second matrix is
  $\operatorname{diag}(J_{W_1},\left(\begin{smallmatrix}0&1\\-1&0\end{smallmatrix}\right))$
  by Lemma~\ref{asm-theta-agree-evaluation-adapted-frame}(2). Expanding,
  $h(z,z')=h_1(\pi z,\pi z')+a(z)^{\mathsf T}J_Xc(\beta(z'))-\beta(z)^{\mathsf T}J_Xc(a(z'))$.
  Here $\beta(z)\cdot\delta(z')=-\beta(z)^{\mathsf T}J_Xc(a(z'))$, and
  $c(\beta(z')\cdot\delta(z))=-a(z)^{\mathsf T}J_Xc(\beta(z'))$ as
  $c(J_X)^{\mathsf T}=J_X$.
\item $Q'$ has entries in $E$.
\item By Lemma~\ref{asm-theta-agree-evaluation-adapted-frame}(1),
  $w(z\iota'^{\mathsf T})=z(Q'\iota')^{\mathsf T}$ has columns $z$, $0$, $0$, so
  $\beta=0$ and $\pi=z$ there; and $w(yx'^{\mathsf T})=y(Q'x')^{\mathsf T}$ vanishes
  outside the column at $x'$, so $\beta=0$ and $\pi=0$ there.
\item For a generator $\gamma=s\,x_0y_0^{\mathsf T}$ of $\Lambda_v$,
  $\pi(\gamma)=s\,x_0(Q'_1y_0)^{\mathsf T}$ with $Q'_1y_0\in L_{W_1}$ by
  Lemma~\ref{asm-theta-agree-evaluation-adapted-frame}(3); so
  $\pi(\Lambda_v)\subset\Lambda^1_v$. Part (4) of that lemma, with
  $(r_0,r_x,r_y)=(0,1,0)$, gives $C$ with $CL_W\subset L_W$ and
  $Q'C=\Delta(0,1,0)Q'$. Put $\gamma_0=\gamma C^{\mathsf T}$, which lies in
  $\Lambda_v$. Then $w(\gamma_0)=w(\gamma)\Delta(0,1,0)$ is the column of
  $w(\gamma)$ at $x'$, so $\beta(\gamma_0)=0$, $\pi(\gamma_0)=0$ and
  $\delta(\gamma_0)=\delta(\gamma)$.
\end{enumerate}
Lemma~\ref{asm-theta-agree-evaluation-adapted-model} gives $M_a$ and the evaluation
map, with $e^a(u)=([p(u),e^1_1(\pi u)],[q(u),e^1_2(\pi u)])$, where $p=0$ where
$\beta=0$ and $q=0$ where $\delta=0$. By (3),
$e^a(z\iota'^{\mathsf T})=([0,e^1_1(z)],[0,e^1_2(z)])$ and
$e^a(yx'^{\mathsf T})=(0,[q(yx'^{\mathsf T}),0])$; and $w(yy'^{\mathsf T})$
vanishes outside the column at $y'$, so $a=0$, $\delta=0$, $\pi=0$ there and
$e^a(yy'^{\mathsf T})=([p(yy'^{\mathsf T}),0],0)$. By
Lemma~\ref{asm-theta-agree-evaluation-adapted-fock},
$e_\infty(\operatorname{Pol}(M_a))=\operatorname{Pol}(M_1)$.

\emph{(2).} Let $H=(x,y,\iota,\mathfrak b)$ and $H'=(x',y',\iota',\mathfrak b')$,
let $Q$ and $Q'$ be the inverses of $\mathrm{fr}_H$ and $\mathrm{fr}_{H'}$, and
let $n_1$ and $d_1$ be the ranks of $X_1$ and $W_1$. For $z\in\mathbb W(X,W)$ put
$w(z)=Qz\,Q'^{\mathsf T}$, with rows indexed by $\{1,\ldots,n_1\}$, $x$, $y$ and
columns by $\{1,\ldots,d_1\}$, $x'$, $y'$, and let $\pi(z)$ be its block in the
first $n_1$ rows and $d_1$ columns. Then $z=\mathrm{fr}_Hw(z)\mathrm{fr}_{H'}^{\mathsf T}$,
and by Lemma~\ref{asm-theta-agree-evaluation-adapted-frame}(2), $h(z,z')$ is the
form $h$ for $G=\operatorname{diag}(J_{X_1},\left(\begin{smallmatrix}0&1\\1&0\end{smallmatrix}\right))$
and $G'=\operatorname{diag}(J_{W_1},\left(\begin{smallmatrix}0&1\\-1&0\end{smallmatrix}\right))$
at $w(z)$ and $w(z')$. Then:
\begin{itemize}
\item for $\mathcal J_\alpha(H,H')$, let $\kappa=(\{1,\ldots,n_1\}\sqcup\{x,y\})\sqcup\{1,\ldots,d_1\}$;
  let $\beta(z)$ consist of $w(z)_{\cdot,y'}$ and of the first $d_1$ entries of
  $w(z)_{y,\cdot}$, and $\delta(z)$ of $-G\,c(w(z)_{\cdot,x'})$ and of
  $J_{W_1}c(w')$, with $w'$ the first $d_1$ entries of $w(z)_{x,\cdot}$;
\item for $\mathcal J_\beta(H,H')$, let $\kappa=(\{1,\ldots,d_1\}\sqcup\{x',y'\})\sqcup\{1,\ldots,n_1\}$;
  let $\beta(z)$ consist of $w(z)_{y,\cdot}$ and of the first $n_1$ entries of
  $w(z)_{\cdot,y'}$, and $\delta(z)$ of $G'c(w(z)_{x,\cdot})$ and of
  $-J_{X_1}c(w'')$, with $w''$ the first $n_1$ entries of $w(z)_{\cdot,x'}$.
\end{itemize}
In both cases $nd=(n_1+2)(d_1+2)=2|\kappa|+n_1d_1$. We check the hypotheses of
Lemma~\ref{asm-theta-agree-evaluation-adapted-model} for $\iota_X=\iota$ and
$\iota_W=\iota'$.
\begin{enumerate}
\item Expand $h_{G,G'}(w(z),w(z'))$ as in (1), for $\mathcal J_\alpha$ first over
  the columns at $x'$, $y'$ and then over the rows at $x$, $y$, and for
  $\mathcal J_\beta$ in the other order, using $c(G)^{\mathsf T}=G$,
  $c(G')^{\mathsf T}=-G'$, $c(J_{X_1})^{\mathsf T}=J_{X_1}$ and
  $c(J_{W_1})^{\mathsf T}=-J_{W_1}$.
\item $Q$ and $Q'$ have entries in $E$.
\item $w(\iota z\iota'^{\mathsf T})=(Q\iota)z(Q'\iota')^{\mathsf T}$ is $z$ in the
  first block and $0$ elsewhere. For $\mathcal J_\alpha$,
  $w(yx'^{\mathsf T})=(Qy)(Q'x')^{\mathsf T}$ vanishes outside the column at $x'$,
  and $w(x(\iota'z')^{\mathsf T})$ outside the first $d_1$ entries of the row at
  $x$. For $\mathcal J_\beta$, $w(xz'^{\mathsf T})$ vanishes outside the row at
  $x$, and $w((\iota z)x'^{\mathsf T})$ outside the first $n_1$ entries of the
  column at $x'$. In each case $\beta=0$ and $\pi=0$ there.
\item $\pi(\Lambda_v)\subset\Lambda^1_v$ by
  Lemma~\ref{asm-theta-agree-evaluation-adapted-frame}(3) on both sides. Part (4)
  gives matrices $A_k$ and $C_k$ that preserve $L_X$ and $L_W$, with
  $QA_k=\Delta_kQ$ and $Q'C_k=\Delta'_kQ'$, and
  $\gamma_0=A_0\gamma C_0^{\mathsf T}+A_1\gamma C_1^{\mathsf T}\in\Lambda_v$ has
  $w(\gamma_0)=\Delta_0w(\gamma)\Delta'_0+\Delta_1w(\gamma)\Delta'_1$. For
  $\mathcal J_\alpha$ take $\Delta_0=1$, $\Delta'_0=\Delta(0,1,0)$,
  $\Delta_1=\Delta(0,1,0)$, $\Delta'_1=\Delta(1,0,0)$: $w(\gamma_0)$ is the column
  of $w(\gamma)$ at $x'$ plus the first $d_1$ entries of its row at $x$. For
  $\mathcal J_\beta$ take $\Delta_0=\Delta(0,1,0)$, $\Delta'_0=1$,
  $\Delta_1=\Delta(1,0,0)$, $\Delta'_1=\Delta(0,1,0)$: the row at $x$ plus the
  first $n_1$ entries of the column at $x'$. In each case $\beta(\gamma_0)=0$,
  $\pi(\gamma_0)=0$ and $\delta(\gamma_0)=\delta(\gamma)$.
\end{enumerate}
Lemma~\ref{asm-theta-agree-evaluation-adapted-model} gives $M_a$ and $M'_a$ with
the two evaluation maps.
\end{proof}
% Lean: (1) is exists_adapted_model_evaluation: exists_model0 (frOK_skew, with w z = z Q'^T,
% b = column inr 1, D = DR (J_X), pi = the inl columns, splitR for (1), mask 0 1 0 for (4),
% then exists_modelS) and Adp.map_eq. (2) is adapted_plane_model_compatibility (isoA, splitR
% then splitL, masks (1,1,1),(0,1,0) and (0,1,0),(1,0,0)) and adapted_plane_model_evaluation_eq
% (isoB, splitL then splitR, masks (0,1,0),(1,1,1) and (1,0,0),(0,1,0)), both through
% Mix.Adapt.exists_model.

\begin{lemma}\label{asm-theta-agree-evaluation-lagrangian-normal}
Let $F$ be a field, $\kappa$ a finite set, and
$L\subset F^\kappa\times F^\kappa$ an $F$-subspace such that $w\in L$ if and
only if $w_1\cdot w'_2=w'_1\cdot w_2$ for every $w'\in L$. Then there are
matrices $P,b\in M_\kappa(F)$ with $P^2=P$ and $b^{\mathsf T}=b$ such that
$(Px,bPx)\in L$ and $(0,(1-P)^{\mathsf T}x)\in L$ for every $x\in F^\kappa$.
\end{lemma}
\leanref{Development/WeilBound/ThetaLifts/LeviDeterminants/LagrangianUniqueness/LagAlg.lean}{PositivityTheory.ThKF0.LagAlg.exists_mats}
\begin{proof}
Let $L_1\subset F^\kappa$ be the image of $L$ under $w\mapsto w_1$. Choose a
linear map $\pi'\colon F^\kappa\to L_1$ that is the identity on $L_1$, and a
linear map $r\colon L_1\to L$ with $r(y)_1=y$ for $y\in L_1$. Put
$s=r\circ\pi'$, and let $P$ and $\beta$ be the matrices of $x\mapsto s(x)_1$
and $x\mapsto s(x)_2$. Then $Px=\pi'(x)$, so $P^2=P$ and $Pw_1=w_1$ for
$w\in L$. As $s(x)$ and $s(y)$ lie in $L$, $Px\cdot\beta y=Py\cdot\beta x$, so
$b=P^{\mathsf T}\beta$ is symmetric. For $w\in L$ and $y\in F^\kappa$,
$w_1\cdot(1-P)^{\mathsf T}y=((1-P)w_1)\cdot y=0$, so $(0,(1-P)^{\mathsf T}y)$
lies in $L$. Finally $s(Px)=(Px,\beta Px)$, so
$(Px,bPx)=s(Px)-(0,(1-P)^{\mathsf T}\beta Px)$ lies in $L$.
\end{proof}

\begin{lemma}\label{asm-theta-agree-evaluation-lagrangian-finite}
Let $G$ be a commutative topological group whose open subgroups form a base of
neighbourhoods of $0$, and let $U_0\subset G$ be a compact open
subgroup. Call a locally constant function with compact support on $G$ a
Schwartz-Bruhat function, and write $\mathbf 1_S$ for the indicator function of
a set $S$. Let $P\subset G$ be a subset containing $0$, and let $\lambda$ be a
linear form on the functions on $G$ such that:
\begin{enumerate}
\item $\lambda(\phi(\cdot+p))=\lambda(\phi)$ for $p\in P$ and every
  Schwartz-Bruhat function $\phi$;
\item for every $a\in G$ with $a\notin P$ there is a locally constant function
  $\chi$ on $G$ with $\chi(a)\neq1$ and $\lambda(\chi\phi)=\lambda(\phi)$ for
  every Schwartz-Bruhat function $\phi$;
\item $\lambda(\mathbf 1_{U_0})=0$.
\end{enumerate}
Then $\lambda(\phi)=0$ for every Schwartz-Bruhat function $\phi$.
\end{lemma}
\leanref{Development/WeilBound/ThetaLifts/LeviDeterminants/LagrangianUniqueness/Imports.lean}{PositivityTheory.ThKF0.FinLine.test}
\begin{proof}
\emph{Decomposition.} A Schwartz-Bruhat function $\phi$ is invariant under
translation by some open subgroup $U_1$: by compactness of its support and local
constancy there is an open subgroup $U_1$ with $\phi(z+u)=\phi(z)$ for $u\in U_1$
and $z$ in the support, and then for every $z$. If $\phi$ is invariant under an
open subgroup $U$, then $\phi=\sum_q\phi(q)\mathbf 1_{q+U}$, the finite sum over
representatives $q$ of the cosets of $U$ on which $\phi$ is nonzero. For an open
subgroup $U\subset U_0$, each $\mathbf 1_{q+U}$ is a Schwartz-Bruhat function.

\emph{Functions vanishing on $P$.} Let $\phi$ be a Schwartz-Bruhat function with
$\phi(a)=0$ for $a\in P$. We show $\lambda(\phi)=0$. For $a\notin P$ choose
$\chi_a$ as in (2) and an open subgroup $V_a$ with $\chi_a$ constant on $a+V_a$.
The support of $\phi$ is compact and contained in the complement of $P$, so it
is covered by finitely many $a+V_a$, $a\in s$. Let $U\subset U_0\cap U_1$ be an
open subgroup contained in every $V_a$, $a\in s$, and write
$\phi=\sum_q\phi(q)\mathbf 1_{q+U}$. For each $q$ with $\phi(q)\neq0$ there is
$a\in s$ with $q+U\subset a+V_a$, so
$\chi_a\mathbf 1_{q+U}=\chi_a(a)\mathbf 1_{q+U}$, and (2) gives
$\lambda(\mathbf 1_{q+U})=\chi_a(a)\lambda(\mathbf 1_{q+U})$. As
$\chi_a(a)\neq1$, $\lambda(\mathbf 1_{q+U})=0$.

\emph{Cosets.} Let $U\subset U_0$ be an open subgroup and $a\in G$. If $p-a\in U$
for some $p\in P$, then $\mathbf 1_{a+U}(\cdot+p)=\mathbf 1_U$, so
$\lambda(\mathbf 1_{a+U})=\lambda(\mathbf 1_U)$ by (1). Otherwise
$\mathbf 1_{a+U}$ vanishes on $P$, and $\lambda(\mathbf 1_{a+U})=0$ by the last
paragraph.

\emph{Conclusion.} Let $U\subset U_0$ be an open subgroup. Writing
$\mathbf 1_{U_0}$ as the sum of the $\mathbf 1_{q+U}$ over the cosets $q+U$ in
$U_0$, the last paragraph gives $\lambda(\mathbf 1_{U_0})=N\lambda(\mathbf 1_U)$,
where $N\geq1$ is the number of these cosets that meet $P$; the coset $U$ is
one of them, as $0\in P$. By (3), $\lambda(\mathbf 1_U)=0$, and then
$\lambda(\mathbf 1_{a+U})=0$ for every $a$. Finally, for a Schwartz-Bruhat
function $\phi$, invariant under $U_1$, write
$\phi=\sum_q\phi(q)\mathbf 1_{q+U}$ with $U=U_0\cap U_1$.
\end{proof}
% Lean: uniform and decomp are the decomposition; off is the second paragraph; key and zero
% inside test are the last two.

\begin{lemma}\label{asm-theta-agree-evaluation-lagrangian-real}
Let $V$ be a finite-dimensional real vector space, $P\subset V$ a subspace and
$\beta$ a symmetric bilinear form on $V$, and write $e(t)=\exp(2\pi it)$ for
$t\in\mathbf R$. Then there is $f_0\in\mathcal S(V)$ with the following
property. Let $T$ be a continuous linear form on $\mathcal S(V)$ such that:
\begin{enumerate}
\item $T(g)=T(f)$ whenever $p\in P$ and
  $g(x)=e\bigl(\beta(x,p)+\tfrac12\beta(p,p)\bigr)f(x+p)$ for all $x$;
\item $T(g)=T(f)$ whenever $l$ is a real linear form on $V$ that vanishes on $P$
  and $g(x)=e(l(x))f(x)$ for all $x$;
\item $T(f_0)=0$.
\end{enumerate}
Then $T=0$.
\end{lemma}
\leanref{Development/WeilBound/ThetaLifts/LeviDeterminants/LagrangianUniqueness/RealTest.lean}{PositivityTheory.ThKF0.RealTest}
\leanref{Development/WeilBound/ThetaLifts/LeviDeterminants/LagrangianUniqueness/HasDerivAtEc.lean}{PositivityTheory.ThKF0.realTest}
\begin{proof}
Choose an inner product $\langle\ ,\ \rangle$ on $V$, let $B$ be the
self-adjoint map with $\langle x,By\rangle=\beta(x,y)$, and let $\mathrm{pr}$ be
the orthogonal projection onto $P$. For $w=(w_1,w_2)\in V\times V$ and
$f\in\mathcal S(V)$ put
$(\rho_V(w)f)(z)=e\bigl(\langle z,w_2\rangle+\tfrac12\langle w_1,w_2\rangle\bigr)f(z+w_1)$,
a Schwartz function. Call a set $S\subset V\times V$ \emph{testing} if there is
$f_0\in\mathcal S(V)$ such that every continuous linear form $T$ on
$\mathcal S(V)$ with $T(\rho_V(w)f)=T(f)$ for $w\in S$ and $f\in\mathcal S(V)$,
and with $T(f_0)=0$, is zero.

\emph{Step 1: $\{0\}\times V$ is testing.} Identify $V$ with $\mathbf R^m$ by an
orthonormal basis; invariance under the $\rho_V(0,w_2)$ is invariance under
multiplication by $x\mapsto e(cx_i)$ for every $i$ and $c\in\mathbf R$. Let
$f_0$ be smooth with compact support and $f_0(0)=1$, and let $T$ be such a form
with $T(f_0)=0$. Let $f$ be smooth with $f(0)=0$ and with support in the set
where $|x_i|<R$ for all $i$. Then $f_1(x)=f(Rx)$ has support where $|x_i|<1$.
By Hadamard's lemma $f_1=\sum_ix_ih_i$ with smooth $h_i$. The function
$q(t)=t/(e(t)-1)$, with $q(0)=1/(2\pi i)$, is smooth on $(-1,1)$. Let $\chi$ be
smooth, equal to $1$ on the support of $f_1$, with support where $|x_i|<1$ for
all $i$, and put $g_i(x)=\chi(x)q(x_i)h_i(x)$. These are smooth with compact
support, and $f_1(x)=\sum_i(e(x_i)-1)g_i(x)$ for all $x$. Hence
$f(x)=\sum_i(e(x_i/R)-1)G_i(x)$ with $G_i(x)=g_i(x/R)$, and
$T(f)=\sum_i\bigl(T(e(x_i/R)G_i)-T(G_i)\bigr)=0$. For any smooth $f$ with
compact support, $f-f(0)f_0$ vanishes at $0$, so $T(f)=f(0)T(f_0)=0$. These
functions are dense in $\mathcal S(V)$ and $T$ is continuous, so $T=0$.

\emph{Step 2: transport.} Let $S$ be testing, with $f_0$, let $U$ be a continuous
linear map of $\mathcal S(V)$ onto itself, and let $\varphi$ be a map of
$V\times V$ with $U(\rho_V(w)f)=\rho_V(\varphi(w))U(f)$ for $w\in S$ and
$f\in\mathcal S(V)$. If $\varphi(S)\subset S'$, then $S'$ is testing, with
$U(f_0)$: if $T$ is invariant under the $\rho_V(w')$, $w'\in S'$, and
$T(Uf_0)=0$, then $T\circ U$ is invariant under the $\rho_V(w)$, $w\in S$, and
vanishes at $f_0$, so $T\circ U=0$ and $T=0$. Two such maps $U$ are the
Fourier transform $\hat f(\xi)=\int_Ve(-\langle v,\xi\rangle)f(v)\,dv$, for which
the substitution $v\mapsto v-w_1$ gives
$(\rho_V(w)f)^\wedge=\rho_V(-w_2,w_1)\hat f$; and, for a self-adjoint map $C$,
multiplication by $e(-\frac12\langle z,Cz\rangle)$, which carries $\rho_V(w)f$
to $\rho_V(w_1,w_2+Cw_1)$ applied to the product, by expanding
$\langle z+w_1,C(z+w_1)\rangle$.

\emph{Step 3.} Applying Step 2 to $\{0\}\times V$: the Fourier transform shows
that the set of $w$ with $w_2=0$ is testing; multiplication with $C=-\mathrm{pr}$
gives the set of $w$ with $w_2=-\mathrm{pr}\,w_1$; the Fourier transform gives
the set of $w$ with $w_1=\mathrm{pr}\,w_2$; and multiplication with
$C=B-\mathrm{pr}$ gives the set $S_P$ of the $(p,Bp+q)$ with $p\in P$ and
$q\perp P$. Indeed, if $w_1=\mathrm{pr}\,w_2$, then $p=w_1\in P$,
$\mathrm{pr}\,p=p$, and $w_2+(B-\mathrm{pr})w_1=Bp+(w_2-\mathrm{pr}\,w_2)$.

\emph{Step 4.} Let $f_0$ be as in the definition for $S_P$, and let $T$ satisfy
(1), (2) and (3). For $p\in P$ and $q\perp P$, $\langle p,q\rangle=0$, so
\[
(\rho_V(p,Bp+q)f)(x)=e(\langle x,q\rangle)\,e\bigl(\beta(x,p)+\tfrac12\beta(p,p)\bigr)f(x+p).
\]
The form $\langle\cdot,q\rangle$ vanishes on $P$, so (2) and then (1) give
$T(\rho_V(p,Bp+q)f)=T(f)$. As $S_P$ is testing, $T=0$.
\end{proof}
% Lean: Step 1 is RealL.test_Y with test_pi, delta_pi and DiracComb.exists_div (Hadamard's
% lemma is hadamard_lemma, q is exists_quot) and DiracComb.dense_hasCompactSupport; Step 2 is
% Test.map, Test.fourier and Test.shear; Step 3 is test_lag; Step 4 is realTest, through
% toEuclidean.

\begin{lemma}\label{asm-theta-agree-evaluation-lagrangian}
Let $F$, $\mathbf A$, $\psi$, $\kappa$ and $\rho$ be as in
Theorem~\ref{asm-theta-agree-evaluation-metaplectic}, and let admissible linear
forms be as in Theorem~\ref{asm-theta-agree-evaluation-heisenberg}. Let
$L\subset F^\kappa\times F^\kappa$ be an $F$-subspace such that $w\in L$ if and
only if $w_1\cdot w'_2=w'_1\cdot w_2$ for every $w'\in L$. Let $\lambda_1$ and
$\lambda_2$ be admissible linear forms with
$\lambda_j(\rho(tw)\Phi)=\lambda_j(\Phi)$ for $t\in\mathbf A$, $w\in L$ and
$\Phi\in\mathcal S(\mathbf A^\kappa)$, where $F^\kappa\subset\mathbf A^\kappa$,
and with $\lambda_2(\Phi)\neq0$ for some $\Phi\in\mathcal S(\mathbf A^\kappa)$.
Then there is $c\in\mathbf C$ with $\lambda_1(\Phi)=c\,\lambda_2(\Phi)$ for
$\Phi\in\mathcal S(\mathbf A^\kappa)$.
\end{lemma}
\leanref{Development/WeilBound/ThetaLifts/LeviDeterminants/LagrangianUniqueness/RealTest.lean}{PositivityTheory.ThKF0.LagLine}
\leanref{Development/WeilBound/ThetaLifts/LeviDeterminants/LagrangianUniqueness/LagLine.lean}{PositivityTheory.ThKF0.lagLine}
\leanref{Development/WeilBound/ThetaLifts/LeviDeterminants/LagrangianUniqueness/LineCore.lean}{PositivityTheory.ThKF0.LineCore.lagLine_of}
\leanref{Development/WeilBound/ThetaLifts/LeviDeterminants/LagrangianUniqueness/LineCore.lean}{PositivityTheory.ThKF0.LineCore.test}
\leanref{Development/WeilBound/ThetaLifts/LeviDeterminants/LagrangianUniqueness/LagAlg.lean}{PositivityTheory.ThKF0.LagAlg.inv_of_mem}
\begin{proof}
Write $\mathcal S=\mathcal S(\mathbf A^\kappa)$, $\operatorname{tr}=\operatorname{Tr}_{F_\infty/\mathbf R}$
and $e(t)=\exp(2\pi it)$, and identify $\mathbf A^\kappa$ with
$F_\infty^\kappa\times\mathbf A_f^\kappa$, so that $\mathbf A_f$ is the set of
adeles with component $0$ at $\infty$. Let $\psi_f$ be the restriction of $\psi$
to $\mathbf A_f$; then $\psi(a)=e(-\operatorname{tr}a_\infty)\psi_f(a_f)$. As a
continuous character, $\psi_f$ is trivial on an open subgroup, so the functions
$z\mapsto\psi_f(R(z))$, for $R$ a polynomial with coefficients in $\mathbf A_f$,
are locally constant, and their products with Schwartz-Bruhat functions on
$\mathbf A_f^\kappa$ are Schwartz-Bruhat. The character $\psi$ is nontrivial,
as it is at $\infty$, so Theorem~\ref{asm-theta-agree-invariance-tate}(4)
applies to it.

\emph{Step 1: normal form.} Lemma~\ref{asm-theta-agree-evaluation-lagrangian-normal}
gives $P,b\in M_\kappa(F)$ with $P^2=P$, $b^{\mathsf T}=b$, and $(Px,bPx)\in L$
and $(0,(1-P)^{\mathsf T}x)\in L$ for $x\in F^\kappa$. Let $\lambda$ be an
admissible form with $\lambda(\rho(tw)\Phi)=\lambda(\Phi)$ for $t\in\mathbf A$,
$w\in L$ and $\Phi\in\mathcal S$. We show:
\begin{enumerate}
\item[(a)] $\lambda(\rho(Px,bPx)\Phi)=\lambda(\Phi)$ and
  $\lambda(\rho(0,(1-P)^{\mathsf T}x)\Phi)=\lambda(\Phi)$ for $x\in\mathbf A^\kappa$
  and $\Phi\in\mathcal S$.
\end{enumerate}
Let $C,D\in M_\kappa(F)$ with $(Cy,Dy)\in L$ for $y\in F^\kappa$, and let $e_i$
be the unit vectors. For $x\in\mathbf A^\kappa$,
$(Cx,Dx)=\sum_ix_iu_i$ with $u_i=(Ce_i,De_i)\in L$. As $L$ is isotropic, the
terms $x_iu_i$ pair to zero, and $\rho(w+w')=\rho(w)\rho(w')$ for $w,w'$ with
$w_1\cdot w'_2=w'_1\cdot w_2$, by the formula for $\rho$. As the $\rho(w)$
preserve $\mathcal S$ (Theorem~\ref{asm-theta-agree-evaluation-heisenberg}(1)),
induction on the number of terms gives $\lambda(\rho(Cx,Dx)\Phi)=\lambda(\Phi)$.
Take $(C,D)=(P,bP)$ and $(C,D)=(0,(1-P)^{\mathsf T})$.

\emph{Step 2: a test function.} Let $f_0$ be given by
Lemma~\ref{asm-theta-agree-evaluation-lagrangian-real} for $V=F_\infty^\kappa$,
the subspace $P_\infty=\{p\in V:Pp=p\}$ and the form
$\beta(x,y)=-\operatorname{tr}(x\cdot by)$. For $z\in\mathbf A_f^\kappa$ put
$Q(z)=\frac12z\cdot bz$, let $U_0=\widehat{\mathcal O}^\kappa$, and put
$\Phi_0=f_0\otimes g_0$ with $g_0(z)=\psi_f(-Q(z))\mathbf 1_{U_0}(z)$. Let
$\lambda$ be an admissible form that satisfies (a) and has $\lambda(\Phi_0)=0$.
We show that $\lambda$ vanishes on $\mathcal S$.

\emph{The finite places.} For functions $g$ on $\mathbf A_f^\kappa$ put
$\mu(g)=\lambda(f_0\otimes\psi_f(-Q)g)$. We check the hypotheses of
Lemma~\ref{asm-theta-agree-evaluation-lagrangian-finite} for $G=\mathbf A_f^\kappa$,
$U_0$ and $P_f=\{p\in\mathbf A_f^\kappa:Pp=p\}$.
\begin{enumerate}
\item Let $p\in P_f$, an element of $\mathbf A^\kappa$ with component $0$ at
  $\infty$. As $Pp=p$, (a) with $x=p$ gives invariance under $\rho(p,bp)$,
  which carries $f_0\otimes\psi_f(-Q)g$ to $f_0\otimes h$ with
  \[
  h(z)=\psi_f\bigl(z\cdot bp+\tfrac12p\cdot bp-Q(z+p)\bigr)g(z+p)=\psi_f(-Q(z))g(z+p),
  \]
  as $b$ is symmetric. So $\mu(g(\cdot+p))=\mu(g)$.
\item Let $a\notin P_f$. Then $(1-P)a\neq0$; let $u=((1-P)a)_i\neq0$ be one of its
  coordinates. As $u\in\mathbf A_f$ is nonzero, it is not in $F$, so by
  Theorem~\ref{asm-theta-agree-invariance-tate}(4) there is $t\in F$ with
  $\psi_f(ut)=\psi(ut)\neq1$. Let $y\in\mathbf A_f^\kappa$ have $i$th coordinate
  $t$ and the others $0$, and put
  $\chi(z)=\psi_f(z\cdot(1-P)^{\mathsf T}y)=\psi_f(((1-P)z)_it)$. Then $\chi$ is
  locally constant, $\chi(a)\neq1$, and $\rho(0,(1-P)^{\mathsf T}y)$ carries
  $f_0\otimes h$ to $f_0\otimes\chi h$. By (a), $\mu(\chi g)=\mu(g)$.
\item $\mu(\mathbf 1_{U_0})=\lambda(\Phi_0)=0$.
\end{enumerate}
So $\mu$ vanishes on the Schwartz-Bruhat functions. Replacing $g$ by
$\psi_f(Q)g$: $\lambda(f_0\otimes g)=0$ for every Schwartz-Bruhat function $g$
on $\mathbf A_f^\kappa$.

\emph{The infinite places.} Fix such a $g$. As $\lambda$ is admissible,
$T(f)=\lambda(f\otimes g)$ is a continuous linear form on $\mathcal S(V)$. We
check the hypotheses of Lemma~\ref{asm-theta-agree-evaluation-lagrangian-real}.
\begin{enumerate}
\item Let $p\in P_\infty$, an element of $\mathbf A^\kappa$ at $\infty$. By (a)
  with $x=p$, $T$ is invariant under $\rho(p,bp)$, which carries $f\otimes g$ to
  $f'\otimes g$ with
  $f'(x)=e\bigl(-\operatorname{tr}(x\cdot bp)-\tfrac12\operatorname{tr}(p\cdot bp)\bigr)f(x+p)
  =e\bigl(\beta(x,p)+\tfrac12\beta(p,p)\bigr)f(x+p)$.
\item Let $l$ be a real linear form on $V$ that vanishes on $P_\infty$. There is
  $y\in V$ with $l(x)=-\operatorname{tr}(x\cdot y)$ for all $x$: the map from $V$
  to its dual sending $y$ to $x\mapsto-\operatorname{tr}(x\cdot y)$ is injective,
  for if $\operatorname{tr}(x\cdot y)=0$ for all $x$, then each coordinate $y_i$,
  as an adele at $\infty$, has $\psi(ay_i)=e(-\operatorname{tr}(ay_i))=1$ for
  $a\in F$, so $y_i\in F$ by Theorem~\ref{asm-theta-agree-invariance-tate}(4),
  and $y_i=0$ as its finite components vanish. As $l(Px)=0$,
  $l(x)=-\operatorname{tr}(((1-P)x)\cdot y)=-\operatorname{tr}(x\cdot(1-P)^{\mathsf T}y)$.
  So $\rho(0,(1-P)^{\mathsf T}y)$, with $y$ at $\infty$, carries $f\otimes g$ to
  $e(l)f\otimes g$, and (a) gives $T(e(l)f)=T(f)$.
\item $T(f_0)=0$ by the finite places.
\end{enumerate}
So $T=0$. As $\mathcal S$ is spanned by the $f\otimes g$, $\lambda$ vanishes on
$\mathcal S$.

\emph{Step 3: the line.} By Step 1, $\lambda_1$ and $\lambda_2$ satisfy (a). As
$\lambda_2$ does not vanish on $\mathcal S$, Step 2 gives $\lambda_2(\Phi_0)\neq0$.
Put $c=\lambda_1(\Phi_0)/\lambda_2(\Phi_0)$. The form $\lambda_1-c\lambda_2$ is
admissible, invariant under the $\rho(tw)$, and vanishes at $\Phi_0$, so by
Steps 1 and 2 it vanishes on $\mathcal S$.
\end{proof}
% Lean: Step 1 is LagAlg.exists_mats and LagAlg.inv_of_mem (with schr_schr and
% schr_mem_schA); Step 2 is LineCore.test (psi_split, schr_pure, nondeg for Theorem
% asm-theta-agree-invariance-tate(4), exists_tr, FinLine.test with levelSub F kappa 1);
% Step 3 is lagLine_of.

\begin{lemma}\label{asm-theta-agree-evaluation-unique}
Let $B$ be a theta base for a CM field $E$, let $X$ be a Hermitian and $W$ a
skew-Hermitian lattice space, of some signatures, let $M$ be a regular model for
$(X,W)$, let $H'$ be an adapted plane in $W$ over a lattice space $W_1$, and let
$M_1$ be a regular model for $(X,W_1)$. If $e$ and $e'$ are evaluation maps for
$(M,M_1,1,\iota_{H'},\mathcal J_0(H'))$, then there is $c\neq0$ with
$e'(\Phi)=c\,e(\Phi)$ for $\Phi\in\operatorname{PD}(M)$.
\end{lemma}
\leanref{Development/WeilBound/ThetaLifts/LeviDeterminants/EigenIntertwinerSkewPlaneUnique.lean}{PositivityTheory.ThKF0.evaluation_at_zero_unique}
\leanref{Development/WeilBound/ThetaLifts/LeviDeterminants/LagrangianUniqueness/GeoCore.lean}{PositivityTheory.ThKF0.Geo.map_line}
\leanref{Development/WeilBound/ThetaLifts/LeviDeterminants/LagrangianUniqueness/GeoCore.lean}{PositivityTheory.ThKF0.Geo.form_line}
\leanref{Development/WeilBound/ThetaLifts/LeviDeterminants/LagrangianUniqueness/GeoCore.lean}{PositivityTheory.ThKF0.Geo.hasLagrangian_skew_adapted}
\leanref{Development/WeilBound/ThetaLifts/LeviDeterminants/LagrangianUniqueness/GeoAlg.lean}{PositivityTheory.ThKF0.Geo.exists_lag}
\begin{proof}
Let $X$, $W$ and $W_1$ have ranks $n$, $d$ and $d_1=d-2$, put
$I=\{1,\ldots,n\}\times\{1,\ldots,d\}$ and $I_1=\{1,\ldots,n\}\times\{1,\ldots,d_1\}$,
let $e^M$ and $e^1$ be the isomorphisms of $M$ and $M_1$ with Schrödinger
operators $\rho$ and $\rho_1$, write $\mathcal S$ for the Schwartz-Bruhat
functions on $\mathbf A^I$, and let $H'=(x',y',\iota',\mathfrak b')$. Let
$\rho^\circ(w)$, for $w\in\mathbf A^I\times\mathbf A^I$, be the operators of
Theorem~\ref{asm-theta-agree-evaluation-heisenberg}, so that
$\rho(u)=\rho^\circ(e^M(u))$.

\emph{Two facts on evaluation maps.} Let $e$ be an evaluation map with
components $e_\infty$ and $e_v$. (a) $e(\Phi)\neq0$ for some $\Phi\in\mathcal S$:
take $f$ with $e_\infty(f)\neq0$ and $\Phi=f\otimes\bigotimes_v\mathrm{vac}_v$;
then $e(\Phi)=e_\infty(f)\otimes\bigotimes_v\mathrm{vac}_{1,v}$, which is nonzero
as in the proof of Lemma~\ref{asm-theta-agree-evaluation-pure}. (b) If
$\lambda$ is admissible, so is $\lambda\circ e$, by the argument of
Lemma~\ref{asm-theta-agree-evaluation-transfer}(1), since every Schwartz-Bruhat
function on $\mathbf A_f^I$ is a combination of products $\bigotimes_v\phi_v$ with
$\phi_v=\mathbf 1_{\mathcal O_v^I}=\mathrm{vac}_v$ for almost all $v$. As
evaluation at $0$ is admissible, so is $\Phi\mapsto e(\Phi)(0)$.

\emph{A rational Lagrangian.} Let $s(y)=y\,x'^{\mathsf T}$ for $y$ a vector of
length $n$, let $j(z)=z\,\iota'^{\mathsf T}$ for $z\in M_{n\times d_1}(\mathbf A_E)$,
and let $L_1$ be the set of $z\in M_{n\times d_1}(E)$ with $e^1_1(z)=0$. By
regularity, $e^M$ and $e^1$ carry the rational matrices $F$-linearly into
rational vectors, and as in the proof of
Lemma~\ref{asm-theta-agree-evaluation-rational} $e^1$ is a bijection from
$M_{n\times d_1}(E)$ onto $F^{I_1}\times F^{I_1}$; so $\dim_FL_1=nd_1$. Let
$L=e^M(s(E^n)+j(L_1))\subset F^I\times F^I$. The form vanishes on it:
$h(s(y),s(y'))=0$ as $h(x',x')=0$; $h(j(z),s(y))=0$ as $h(\iota'u,x')=0$; and
$\langle j(z),j(z')\rangle=\langle z,z'\rangle_1
=e^1_1(z)\cdot e^1_2(z')-e^1_1(z')\cdot e^1_2(z)=0$ for $z,z'\in L_1$, as
$\iota'^{\mathsf T}J_Wc(\iota')=J_{W_1}$. The map $s$ is injective, as
$x'\neq0$, so $\dim_Fs(E^n)=2n$. The map $j$ preserves the forms and
$\langle\ ,\ \rangle_1$ is nondegenerate, so $j$ is injective; and if
$s(y)=j(z)$, then $\langle z,z'\rangle_1=\langle s(y),j(z')\rangle=0$ for every
$z'$, so $z=0$. Hence $\dim_FL=2n+nd_1=nd$, the number of elements of $I$, and
an isotropic subspace of this dimension is its own orthogonal: $w\in L$ if and
only if $w_1\cdot w'_2=w'_1\cdot w_2$ for all $w'\in L$. For $t\in\mathbf A$ and
$w=e^M(s(y)+j(z_0))\in L$, the elements $ts(y)\in\mathcal J_0(H')$ and $tj(z_0)$
pair to zero, so
\[
\rho^\circ(tw)=\rho(ts(y))\,\rho(j(tz_0)),\qquad e^1_1(tz_0)=0,
\]
by the rule $\rho^\circ(w_1+w_2)=\rho^\circ(w_1)\rho^\circ(w_2)$ for $w_1,w_2$
with $w_{1,1}\cdot w_{2,2}=w_{2,1}\cdot w_{1,2}$, which follows from the formula
for $\rho^\circ$.

\emph{Uniqueness.} Put $\ell(\Phi)=e(\Phi)(0)$ and $\ell'(\Phi)=e'(\Phi)(0)$,
admissible by (b). If $e^1_1(z)=0$, then $(\rho_1(z)\Psi)(0)=\Psi(0)$, so
$\ell(\rho(j(z))\Phi)=(\rho_1(z)e(\Phi))(0)=\ell(\Phi)$ for $\Phi\in\mathcal S$;
and $\ell(\rho(u)\Phi)=\ell(\Phi)$ for $u\in\mathcal J_0(H')$. As the
$\rho(u)$ preserve $\mathcal S$
(Theorem~\ref{asm-theta-agree-evaluation-heisenberg}(1)), the display gives
$\ell(\rho^\circ(tw)\Phi)=\ell(\Phi)$ for $t\in\mathbf A$ and $w\in L$, and
likewise for $\ell'$. By (a) there are $\Phi_0\in\mathcal S$ and $y_0$ with
$e(\Phi_0)(y_0)\neq0$. For $y\in\mathbf A^{I_1}$ let $z_y$ have
$e^1(z_y)=(y,0)$; then $(\rho_1(z_y)\Psi)(0)=\Psi(y)$, so
$\ell(\rho(j(z_{y_0}))\Phi_0)=e(\Phi_0)(y_0)\neq0$. By
Lemma~\ref{asm-theta-agree-evaluation-lagrangian} there is $c$ with $\ell'=c\,\ell$
on $\mathcal S$. For $\Phi\in\mathcal S$ and every $y$,
\[
e'(\Phi)(y)=(\rho_1(z_y)e'(\Phi))(0)=\ell'(\rho(j(z_y))\Phi)=c\,\ell(\rho(j(z_y))\Phi)=c\,e(\Phi)(y).
\]
So $e'=c\,e$ on $\mathcal S$, which contains $\operatorname{PD}(M)$ as the
vacua of $M$ are almost all $\mathbf 1_{\mathcal O_v^I}$. By (a) applied to $e'$,
$c\neq0$.
\end{proof}
% Lean: (a) is EvNonzero.exists_mem_schA_ev_ne_zero (evaluation_intertwining_support);
% (b) is GeoAdm.comp_mem_adm with Geo.proj_adm (evaluation_intertwining_of_models); the
% Lagrangian is hasLagrangian_skew_adapted through hasLag_of and exists_lag (with
% lag_of_finrank and schr_add); uniqueness is map_line with form_line and lagLine.

\begin{lemma}\label{asm-theta-agree-evaluation}
Let $(X,W,M,\chi_X,\chi_W,\omega)$ be a Weil pair for a theta base $B$ of a CM
field $E$ (Definition~\ref{asm-theta-agree-triple}(1)), let $H'$ be an adapted
plane in $W$ over a lattice space $W_1$, and let $M_1$ be a regular model for
$(X,W_1)$. Then there is an evaluation map for
$(M,M_1,1,\iota_{H'},\mathcal J_0(H'))$
(Definition~\ref{asm-theta-agree-planes}(5)), and every such map carries
$\operatorname{PD}(M)$ onto $\operatorname{PD}(M_1)$.
\end{lemma}
\leanref{Development/AbelianVarietyTheory/Definitions/Divisors/CTSmallCone.lean}{GlobalTheory.ThK.EvZero}
\begin{proof}
Since $(X,W,M,\chi_X,\chi_W,\omega)$ is a Weil pair, $M$ is regular, $X$ is
Hermitian and $W$ is skew-Hermitian.

\emph{Existence.} By Lemma~\ref{asm-theta-agree-evaluation-adapted}(1) there are a
regular model $M_a$ for $(X,W)$ and an evaluation map $e_2$ for
$(M_a,M_1,1,\iota_{H'},\mathcal J_0(H'))$ with components $e_{2,\infty}$ and
$e_{2,v}$ such that $e_{2,\infty}(\operatorname{Pol}(M_a))=\operatorname{Pol}(M_1)$.
By Lemma~\ref{asm-theta-agree-evaluation-change} there is an evaluation map $e_1$
for $(M,M_a,1,1,\emptyset)$ that maps the Schwartz-Bruhat functions bijectively
onto themselves; let $e_{1,\infty}$ and $e_{1,v}$ be components of it. By
Lemma~\ref{asm-theta-agree-evaluation-pure}, $e_{1,\infty}$ is surjective and
$e_{1,\infty}(\operatorname{Pol}(M))=\operatorname{Pol}(M_a)$. By
Lemma~\ref{asm-theta-agree-evaluation-compose}(1), $e_0=e_2\circ e_1$ is an
evaluation map for $(M,M_1,1,\iota_{H'},\mathcal J_0(H'))$.

\emph{The image of $e_0$.} The vacua of the regular models $M$, $M_a$ and $M_1$
are Schwartz-Bruhat. Lemma~\ref{asm-theta-agree-evaluation-compose}(2), applied to
$e_1$ with $e_{1,\infty}$ and the $e_{1,v}$, gives
$e_1(\operatorname{PD}(M))=\operatorname{PD}(M_a)$; applied to $e_2$ with
$e_{2,\infty}$ and the $e_{2,v}$, it gives
$e_2(\operatorname{PD}(M_a))=\operatorname{PD}(M_1)$. So
$e_0(\operatorname{PD}(M))=\operatorname{PD}(M_1)$.

\emph{Every evaluation map.} Let $e$ be an evaluation map for
$(M,M_1,1,\iota_{H'},\mathcal J_0(H'))$. By
Lemma~\ref{asm-theta-agree-evaluation-unique} there is $c\neq0$ with
$e(\Phi)=c\,e_0(\Phi)$ for $\Phi\in\operatorname{PD}(M)$. As
$\operatorname{PD}(M)$ is a subspace and $c\neq0$,
$e(\operatorname{PD}(M))=c\,e_0(\operatorname{PD}(M))=\operatorname{PD}(M_1)$.
\end{proof}
% Lean: evZero = evZero_of_adapted exists_adapted_model_evaluation, that is
% evZero_of_adel adelic_estimates_holds schwartzDq, which is evZero_of_bij (evBij_of_adel)
% in MapEqOfSmul.lean: change_core, isEv_comp, polData_map twice,
% evaluation_at_zero_unique and map_eq_of_smul. What is still open lies in Lemmas
% asm-theta-agree-evaluation-adapted and -lagrangian.

\begin{lemma}\label{asm-theta-agree-evaluation-cone}
Let $(X,W,M,\chi_X,\chi_W,\omega)$ be a Weil pair for a theta base $B$ of a CM
field $E$ (Definition~\ref{asm-theta-agree-triple}(1)), let $H$ be an adapted
plane in $X$ over a lattice space $X_1$ and $H'$ one in $W$ over a lattice
space $W_1$, and let $M_1$ be a regular model for $(X_1,W_1)$. Then there are
evaluation maps for $(M,M_1,\iota_H,\iota_{H'},\mathcal J_\alpha(H,H'))$ and for
$(M,M_1,\iota_H,\iota_{H'},\mathcal J_\beta(H,H'))$
(Definition~\ref{asm-theta-agree-planes}(5)).
\end{lemma}
\leanref{Development/AbelianVarietyTheory/Definitions/Divisors/CTSmallCone.lean}{GlobalTheory.ThK.EvCone}
\begin{proof}
Let $\mathcal J$ be $\mathcal J_\alpha(H,H')$ or $\mathcal J_\beta(H,H')$. The model
$M$ is regular, $X$ is Hermitian and $W$ skew-Hermitian, since
$(X,W,M,\chi_X,\chi_W,\omega)$ is a Weil pair. By
Lemma~\ref{asm-theta-agree-evaluation-adapted}(2) there are a regular model $M_a$
for $(X,W)$ and an evaluation map $e_2$ for
$(M_a,M_1,\iota_H,\iota_{H'},\mathcal J)$. By
Lemma~\ref{asm-theta-agree-evaluation-change} there is an evaluation map $e_1$
for $(M,M_a,1,1,\emptyset)$ that maps the Schwartz-Bruhat functions bijectively
onto themselves, and by Lemma~\ref{asm-theta-agree-evaluation-pure} its
component at $\infty$ is surjective. By
Lemma~\ref{asm-theta-agree-evaluation-compose}(1), $e_2\circ e_1$ is an
evaluation map for $(M,M_1,\iota_H,\iota_{H'},\mathcal J)$.
\end{proof}
% Lean: evCone = evCone_of_adel adelic_estimates_holds, which is evCone_of_bij (evBij_of_adel)
% in MapEqOfSmul.lean: adapted_plane_model_compatibility (for isoA) and
% adapted_plane_model_evaluation_eq (for isoB), then change_core and isEv_comp. What is
% still open lies in Lemma asm-theta-agree-evaluation-adapted(2).

\subsection{Adapted planes and configurations}\label{subsection-asm-adapted-planes}

\begin{theorem}[Hasse-Minkowski \cite{OMe63}]\label{asm-theta-agree-tower-isotropic}
Let $E$ be a CM field with maximal totally real subfield $F$ and complex
conjugation $c$, and let $\tau_w\colon E\to\mathbf C$ be an embedding over
each infinite place $w$ of $F$. Let $n\geq3$, let $\varepsilon=\pm1$, let
$J\in GL_n(E)$ with $c(J)^{\mathsf T}=\varepsilon J$, put
$h(x,y)=x^{\mathsf T}J\,c(y)$, and let $\kappa=1$ if $\varepsilon=1$ and
$\kappa=i$ if $\varepsilon=-1$. Suppose that for every $w$ the Hermitian
matrix $\kappa\,\tau_w(J)$ has signature $(r,s)$ with $r,s\geq1$. Then there
is $x\in E^n$ with $x\neq0$ and $h(x,x)=0$.
\end{theorem}
\leanref{Development/WeilBound/ThetaLifts/TowerModels/PlaneIso.lean}{PositivityTheory.ThKG.Tow.Pln.exists_isotropic}
\leanref{Development/AbelianVarietyTheory/Proofs/AdelicLatticeProofs/ExistsIsotropicOfIndefinite.lean}{GlobalTheory.UTrans.exists_isotropic_of_indefinite}

\begin{lemma}\label{asm-theta-agree-tower-pair}
Let $E$ be a CM field with a theta base $B$
(Definition~\ref{asm-theta-agree-models}), with set of places $P$ and
embeddings $\tau_v\colon E\to F_v$ for $v\in P$, and let $\varepsilon=\pm1$.
Let $D=\mathcal O_E$ if $\varepsilon=1$, and let $D$ be the codifferent
$\{t\in E:\operatorname{Tr}_{E/\mathbf Q}(ot)\in\mathbf Z\text{ for all }o\in\mathcal O_E\}$
of $E$ if $\varepsilon=-1$. Let $J\in GL_n(E)$ with
$c(J)^{\mathsf T}=\varepsilon J$, put $h(x,y)=x^{\mathsf T}J\,c(y)$, and let
$L\subset E^n$ be an $\mathcal O_E$-lattice with
$L=\{z\in E^n:h(z,a)\in D\text{ for all }a\in L\}$. Let $x_0\in E^n$ with
$x_0\neq0$ and $h(x_0,x_0)=0$. Then there are $x,y\in E^n$ such that:
\begin{enumerate}
\item $h(x,x)=h(y,y)=0$ and $h(x,y)=1$;
\item $h(z,x)\,y\in L$ and $h(z,y)\,x\in L$ for every $z\in L$;
\item for every $v\in P$, $|\tau_v(s)|_v\leq1$ for every $s\in E$ with
  $sx\in L$, and $|\tau_v(s)|_v=1$ for some $s\in E$ with $sx\in L$; and the
  same holds for $y$.
\end{enumerate}
\end{lemma}
\leanref{Development/WeilBound/ThetaLifts/TowerModels/PlanePair.lean}{PositivityTheory.ThKG.Tow.Pln.exists_adapted_pair}
\leanref{Development/WeilBound/ThetaLifts/TowerModels/PlaneLat.lean}{PositivityTheory.ThKG.Tow.Pln.exists_partner}
\leanref{Development/WeilBound/ThetaLifts/TowerModels/PlaneLat.lean}{PositivityTheory.ThKG.Tow.Pln.exists_max_val}
\leanref{Development/WeilBound/ThetaLifts/TowerModels/PlaneLat.lean}{PositivityTheory.ThKG.Tow.Pln.exists_denom}
\leanref{Development/WeilBound/ThetaLifts/TowerModels/PlaneCodiff.lean}{PositivityTheory.ThKG.Tow.Pln.val_le_one_of_codE}
\begin{proof}
As $c(J)^{\mathsf T}=\varepsilon J$ and $\varepsilon^2=1$,
$h(a',a)=\varepsilon\,c(h(a,a'))$ for all $a,a'$. As $1\in D$ and $D$ is an
$\mathcal O_E$-module, $\mathcal O_E\subset D$ and $\varepsilon D=D$. For
$v\in P$ let $\mathfrak p_v$ and $\bar{\mathfrak p}_v$ be the primes of
$\mathcal O_E$ given by $\tau_v$ and $\tau_v\circ c$; they are distinct, as
$v$ splits in $E$.

\emph{Elements of $D$ are integral at $\tau_v$.} This is clear if
$D=\mathcal O_E$. Let $D$ be the codifferent, $t\in D$ and $v\in P$. For
$o\in\mathcal O_E$, $\tau_v(ot)+\tau_v(c(ot))=\operatorname{Tr}_{E/F}(ot)$ lies
in the codifferent of $F$. As $\psi$ is standard, $\psi_v$ is trivial on
$t'\mathcal O_v$ for every $t'$ in the codifferent of $F$, and as $\psi_v$ has
conductor $\mathcal O_v$, such $t'$ are integral at $v$. Choose
$p\in\bar{\mathfrak p}_v$ with $p\notin\mathfrak p_v$. Then $\tau_v(p)$ is a
unit and $\tau_v(c(p))$ lies in the maximal ideal, so for $k$ large
$\tau_v(c(p^kt))=\tau_v(c(p))^k\tau_v(c(t))$ is integral; with $o=p^k$ this makes
$\tau_v(p^kt)$ integral, and so $\tau_v(t)$ is integral.

\emph{The vector $x$.} The linear form $z\mapsto h(z,x_0)$ is not zero on $L$,
since $L$ spans $E^n$, $J$ is invertible and $x_0\neq0$. Let $S$ be the finite
set of the primes $\mathfrak p_v$ and $\bar{\mathfrak p}_v$, $v\in P$. As $L$
is finitely generated, for each $\mathfrak p\in S$ there is
$z_{\mathfrak p}\in L$ such that $h(z_{\mathfrak p},x_0)\neq0$ has the least
$\mathfrak p$-order among the $h(z,x_0)$, $z\in L$. Choose
$e_{\mathfrak p}\in\mathcal O_E$ prime to $\mathfrak p$ and lying in the other
primes of $S$, and put $z_0=\sum_{\mathfrak p\in S}e_{\mathfrak p}z_{\mathfrak p}$.
Then $c_0=h(z_0,x_0)$ has the same $\mathfrak p$-order as
$h(z_{\mathfrak p},x_0)$ for every $\mathfrak p\in S$. Put
$x=c(c_0)^{-1}x_0$. Then $h(x,x)=0$, $h(z_0,x)=c_0^{-1}h(z_0,x_0)=1$, and
$h(z,x)=c_0^{-1}h(z,x_0)$ is integral at every prime of $S$, for $z\in L$.

\emph{A partner.} The set $C=\{h(z,x):z\in L\}$ is a fractional ideal of $E$
containing $1$, so $CC^{-1}=\mathcal O_E$. Write $1=\sum_kb_kh(z_k,x)$ with
$b_k\in C^{-1}$ and $z_k\in L$, and put $y_0=\varepsilon\sum_kb_kz_k$. Then
$h(y_0,x)=\varepsilon$, so $h(x,y_0)=\varepsilon\,c(\varepsilon)=1$, and
$sy_0\in L$ for every $s\in C$, since $sb_k\in\mathcal O_E$. If $y'\in E^n$ has
$sy'\in L$ for all $s\in C$, then $h(z,y')\,x\in L$ for all $z\in L$: for
$z'\in L$,
$h(h(z,y')x,z')=\varepsilon\,h(z,y')\,c(h(z',x))=\varepsilon\,h(z,h(z',x)y')$
lies in $D$, as $h(z',x)y'\in L$, and $L$ is its own dual.

Let $q=h(y_0,y_0)$, so that $\varepsilon\,c(q)=q$. As $E/F$ is unramified at
the finite places, the trace $\mathcal O_E\to\mathcal O_F$ is onto, so there is
$\theta\in\mathcal O_E$ with $\theta+c(\theta)=1$. Put $y=y_0-\theta q\,x$. For
$s\in C$, $s\theta q\,x=\theta\,h(sy_0,y_0)\,x$ lies in $L$, by the previous
paragraph applied to $y'=y_0$ and $z=sy_0$; so $sy\in L$. Further
$h(x,y)=h(x,y_0)=1$ and
\[
h(y,y)=q-c(\theta)c(q)\varepsilon-\theta q=q\,c(\theta)-c(\theta)\,\varepsilon c(q)=0 .
\]
This gives (1). For $z\in L$, $h(z,x)\in C$ gives $h(z,x)y\in L$, and the
previous paragraph applied to $y'=y$ gives $h(z,y)x\in L$; this is (2).

\emph{The places of $P$.} Let $v\in P$. If $sx\in L$, then
$h(sx,z_0)=s\,\varepsilon\,c(h(z_0,x))=\varepsilon s$ lies in $D$, so
$|\tau_v(s)|_v\leq1$. The ideal $C$ is integral at $\bar{\mathfrak p}_v$, so
there is $d\in\mathcal O_E$ with $d\notin\bar{\mathfrak p}_v$ and
$dC\subset\mathcal O_E$. Then $c(d)x\in L$, since
$h(c(d)x,z)=\varepsilon\,c(d\,h(z,x))\in\mathcal O_E\subset D$ for $z\in L$, and
$|\tau_v(c(d))|_v=1$, since $d\notin\bar{\mathfrak p}_v$. If $sy\in L$, then
$h(sy,x)=\varepsilon s$ is integral at $\mathfrak p_v$, so
$|\tau_v(s)|_v\leq1$. Finally $y=h(z_0,x)\,y\in L$. This is (3).
\end{proof}

\begin{lemma}\label{asm-theta-agree-tower-cartan}
Let $n_1\geq0$ and $n=n_1+2$, and let $M\in M_n(\mathbf C)$ be a Hermitian
matrix, $\bar M^{\mathsf T}=M$, of signature $(r,s)$: there is an invertible
$A$ with $A^{\mathsf T}M\bar A=\operatorname{diag}(1_r,-1_s)$. Write
$\langle u,u'\rangle=u^{\mathsf T}M\bar u'$. Let $X,Y\in\mathbf C^n$ with
$\langle X,X\rangle=\langle Y,Y\rangle=0$ and $\langle X,Y\rangle=1$, and let
$\iota\in M_{n\times n_1}(\mathbf C)$ have linearly independent columns with
$\langle\iota u,X\rangle=\langle\iota u,Y\rangle=0$ for all
$u\in\mathbf C^{n_1}$. Then $M_1=\iota^{\mathsf T}M\bar\iota$ has signature
$(r-1,s-1)$, and there are Cartan involutions $\theta$ of $M$ and $\theta_1$ of
$M_1$ (Definition~\ref{asm-theta-agree-doubling}(1)) with
$\theta\iota=\iota\theta_1$ and $Y=2\theta X$.
\end{lemma}
\leanref{Development/WeilBound/ThetaLifts/TowerModels/PlaneReal.lean}{PositivityTheory.ThKG.Tow.Pln.real_plane}
\leanref{Development/WeilBound/ThetaLifts/TowerModels/PlaneReal.lean}{PositivityTheory.ThKG.Tow.Pln.cartan_of_frame}
\leanref{Development/WeilBound/ThetaLifts/TowerModels/PlaneReal.lean}{PositivityTheory.ThKG.Tow.Pln.hasSig_unique}
\begin{proof}
For a Hermitian $N$ and an invertible $A$ with $A^{\mathsf T}N\bar A=\Delta$,
$\Delta=\operatorname{diag}(1_p,-1_q)$, the matrix $\theta_A=A\Delta A^{-1}$ is
a Cartan involution of $N$: $\theta_A^2=1$; since
$N=(A^{-1})^{\mathsf T}\Delta\bar A^{-1}$,
$\theta_A^{\mathsf T}N\bar\theta_A=(A^{-1})^{\mathsf T}\Delta\Delta\Delta\bar A^{-1}=N$;
and for $x=Au\neq0$,
$(\theta_Ax)^{\mathsf T}N\bar x=u^{\mathsf T}\Delta(A^{\mathsf T}N\bar A)\bar u=\sum_k|u_k|^2>0$.

The vectors $X$, $Y$ and the columns of $\iota$ are linearly independent: if
$aX+bY+\iota u=0$, pairing with $Y$ gives $a=0$, pairing with $X$ gives $b=0$,
as $\langle Y,X\rangle=\overline{\langle X,Y\rangle}=1$, and then $u=0$. So
$\mathbf C^n=\mathbf CX\oplus\iota\mathbf C^{n_1}\oplus\mathbf CY$. The matrix
$M_1$ is Hermitian, and it is invertible: if $\langle\iota u,\iota u'\rangle=0$
for all $u'$, then $\iota u$ is orthogonal to all of $\mathbf C^n$, so
$\iota u=0$ and $u=0$. Let $A_1$ be invertible with
$A_1^{\mathsf T}M_1\bar A_1=\operatorname{diag}(1_p,-1_q)$, with columns
$a_1,\ldots,a_p,a'_1,\ldots,a'_q$. The vectors
$X+\frac12Y,\iota a_1,\ldots,\iota a_p,X-\frac12Y,\iota a'_1,\ldots,\iota a'_q$
are the columns of an invertible $A$ with
$A^{\mathsf T}M\bar A=\operatorname{diag}(1_{p+1},-1_{q+1})$, since
$\langle X\pm\frac12Y,X\pm\frac12Y\rangle=\pm1$,
$\langle X+\frac12Y,X-\frac12Y\rangle=0$, and the $\iota a$ are orthogonal to
$X$ and $Y$. By Sylvester's law of inertia, $p+1=r$ and $q+1=s$.

Put $\theta=\theta_A$ and $\theta_1=\theta_{A_1}$. Then $\theta$ fixes
$X+\frac12Y$ and the $\iota a_j$ and negates $X-\frac12Y$ and the
$\iota a'_j$, while $\theta_1$ fixes the $a_j$ and negates the $a'_j$. So
$\theta\iota$ and $\iota\theta_1$ agree on the columns of $A_1$, and
$\theta\iota=\iota\theta_1$. Finally
$\theta X=\frac12\bigl((X+\frac12Y)-(X-\frac12Y)\bigr)=\frac12Y$.
\end{proof}

\begin{lemma}\label{asm-theta-agree-tower-basis}
Let $E$ be a CM field with maximal totally real subfield $F$, let $v$ be a
finite place of $F$ and $\tau\colon E\to F_v$ an $F$-embedding, and for an
$\mathcal O_E$-lattice $N\subset E^m$ let $N_\tau\subset F_v^m$ be its
completion at the place of $E$ given by $\tau$, the $\mathcal O_v$-span of
$\tau(N)$. Let $n_1\geq0$ and $n=n_1+2$, let $L\subset E^n$ be an
$\mathcal O_E$-lattice, let $x,y\in E^n$ and $\iota\in M_{n\times n_1}(E)$, and
put $L_1=\{u\in E^{n_1}:\iota u\in L\}$. Suppose that every $z\in L$ is
$sx+\iota u+ty$ with $sx\in L$, $ty\in L$ and $u\in L_1$, and that
$|\tau(s)|_v\leq1$ for every $s\in E$ with $sx\in L$ and $|\tau(s)|_v=1$ for
some such $s$, and the same for $y$. Let $\beta_1\in GL_{n_1}(F_v)$ with
$\beta_1\mathcal O_v^{n_1}=(L_1)_\tau$, let $0\leq i_0<i_1<n$, and let $\beta$
be the $n\times n$ matrix whose column $i_0$ is $\tau(x)$, whose column $i_1$
is $\tau(y)$, and whose other columns, in order, are the columns of
$\tau(\iota)\beta_1$. Then $\beta\in GL_n(F_v)$ and
$\beta\mathcal O_v^n=L_\tau$.
\end{lemma}
\leanref{Development/WeilBound/ThetaLifts/TowerModels/PlaneBas.lean}{PositivityTheory.ThKG.Tow.Pln.exists_basis}
\begin{proof}
Let $N'=\beta\mathcal O_v^n$, the $\mathcal O_v$-span of the columns of
$\beta$. For $z\in L$ write $z=sx+\iota u+ty$ as above. Then $\tau(s)$ and
$\tau(t)$ are integral and $\tau(u)\in(L_1)_\tau=\beta_1\mathcal O_v^{n_1}$, so
$\tau(z)=\tau(s)\tau(x)+\tau(\iota)\beta_1\bigl(\beta_1^{-1}\tau(u)\bigr)+\tau(t)\tau(y)$
lies in $N'$. Hence $L_\tau\subset N'$. As $L$ spans $E^n$, $\tau(L)$ spans
$F_v^n$, so $\beta$ is onto and invertible. Conversely, if $sx\in L$ with
$|\tau(s)|_v=1$, then $\tau(x)=\tau(s)^{-1}\tau(sx)$ lies in $L_\tau$, and the
same holds for $y$; and $\tau(\iota)(L_1)_\tau\subset L_\tau$, as
$\iota L_1\subset L$, so the columns of $\tau(\iota)\beta_1$ lie in $L_\tau$.
Hence $N'=L_\tau$.
\end{proof}

\begin{proposition}\label{asm-theta-agree-tower-plane}
Let $E$ be a CM field with a theta base $B$
(Definition~\ref{asm-theta-agree-models}), let $n_1\geq1$ and $n=n_1+2$, let
$Y$ be a lattice space of rank $n$ (Definition~\ref{asm-theta-agree-doubling}(1))
that is Hermitian of signature $(r,s)$, or skew-Hermitian of signature $(r,s)$,
with $r,s\geq1$, and let $0\leq i_0<i_1<n$. For a lattice space $Y^\circ$ of
rank $n$, write $Y[Y^\circ]$ for the lattice space with the matrix and the
lattice of $Y$ and the matrices $\theta_w$ and $b_v$ of $Y^\circ$. Then there
are lattice spaces $Y^\circ$ of rank $n$ and $Y_1$ of rank $n_1$ such that
$Y[Y^\circ]$ is Hermitian, respectively skew-Hermitian, of signature $(r,s)$,
$Y_1$ is Hermitian, respectively skew-Hermitian, of signature $(r-1,s-1)$, and
there is an adapted plane in $Y[Y^\circ]$ over $Y_1$ at positions $i_0<i_1$
(Definition~\ref{asm-theta-agree-planes}(1)).
\end{proposition}
\leanref{Development/WeilBound/ThetaLifts/TowerModels/PlaneHerm.lean}{PositivityTheory.ThKG.Tow.PlaneHerm}
\leanref{Development/WeilBound/ThetaLifts/TowerModels/PlaneHerm.lean}{PositivityTheory.ThKG.Tow.PlaneSkew}
\leanref{Development/WeilBound/ThetaLifts/TowerModels/PlaneSkew.lean}{PositivityTheory.ThKG.Tow.planeHerm}
\leanref{Development/WeilBound/ThetaLifts/TowerModels/PlaneSkew.lean}{PositivityTheory.ThKG.Tow.planeSkew}
\leanref{Development/WeilBound/ThetaLifts/TowerModels/PlaneMain.lean}{PositivityTheory.ThKG.Tow.Pln.exists_plane}
\leanref{Development/WeilBound/ThetaLifts/TowerModels/PlaneComp.lean}{PositivityTheory.ThKG.Tow.Pln.exists_compl}
\begin{proof}
In the Hermitian case put $\varepsilon=1$, $\kappa=1$ and $D=\mathcal O_E$; in
the skew-Hermitian case put $\varepsilon=-1$, $\kappa=i$ and let $D$ be the
codifferent of $E$. Let $J$ and $L$ be the matrix and the lattice of $Y$, and
$h(a,a')=a^{\mathsf T}J\,c(a')$. Then $c(J)^{\mathsf T}=\varepsilon J$,
$L=\{z:h(z,a)\in D\text{ for all }a\in L\}$, and for every infinite place $w$
the matrix $\kappa\tau_w(J)$ has signature $(r,s)$, so $J$ is invertible. As
$E$ is a CM field, $\tau_w\circ c$ is the complex conjugate of $\tau_w$.

\emph{The pair.} As $n\geq3$, Theorem~\ref{asm-theta-agree-tower-isotropic}
gives $x_0\neq0$ with $h(x_0,x_0)=0$, and
Lemma~\ref{asm-theta-agree-tower-pair} gives $x$ and $y$ with its properties
(1) to (3). Then $h(y,x)=\varepsilon\,c(h(x,y))=\varepsilon$.

\emph{The complement.} The map $z\mapsto(h(z,x),h(z,y))$ from $E^n$ to $E^2$ is
onto, as $\varepsilon a\,y+b\,x\mapsto(a,b)$, so its kernel has dimension
$n_1$. Let $\iota\in M_{n\times n_1}(E)$ have as columns a basis of the kernel.
Then $h(\iota u,x)=h(\iota u,y)=0$ for all $u$, $\iota$ is injective, and every
$z$ with $h(z,x)=h(z,y)=0$ is $\iota u$ for some $u$.

\emph{The splitting.} For $z\in L$, the vector
$z-h(z,y)x-\varepsilon h(z,x)y$ is orthogonal to $x$ and $y$, so it is
$\iota u$ for some $u$. By Lemma~\ref{asm-theta-agree-tower-pair}(2),
$h(z,y)x\in L$ and $\varepsilon h(z,x)y\in L$, so $\iota u\in L$.

\emph{The smaller space.} Put $J_1=\iota^{\mathsf T}J\,c(\iota)$, so that
$h_1(u,u')=u^{\mathsf T}J_1c(u')$ equals $h(\iota u,\iota u')$ and
$c(J_1)^{\mathsf T}=\varepsilon J_1$; put
$L_1=\{u\in E^{n_1}:\iota u\in L\}$ and
$\mathfrak b=\{s\in E:sx\in L\}$. As $\iota$ is injective and $L$ is a
lattice, $L_1$ is finitely generated, and it spans $E^{n_1}$, since some
nonzero integer multiple of every $\iota u$ lies in $L$; so $L_1$ is a
lattice. As $x\neq0$, $\mathfrak b$ is a fractional ideal. The lattice $L_1$
is its own dual for $h_1$ and $D$: if $u\in L_1$, then
$h_1(u,u')=h(\iota u,\iota u')\in D$ for $u'\in L_1$; conversely, if
$h_1(u,u')\in D$ for all $u'\in L_1$, then for $z=sx+\iota u'+ty\in L$ as in
the splitting,
$h(\iota u,z)=c(s)h(\iota u,x)+h_1(u,u')+c(t)h(\iota u,y)=h_1(u,u')\in D$,
so $\iota u\in L$ and $u\in L_1$.

\emph{The infinite places.} Let $w$ be an infinite place. Apply
Lemma~\ref{asm-theta-agree-tower-cartan} to $M=\kappa\tau_w(J)$,
$X=\tau_w(x)$, $Y=\kappa\tau_w(y)$ and $\tau_w(\iota)$. The matrix $M$ is
Hermitian, as $\bar\kappa\varepsilon=\kappa$; with
$\langle u,u'\rangle=u^{\mathsf T}M\bar u'$ we have
$\langle\tau_w(a),\tau_w(a')\rangle=\kappa\tau_w(h(a,a'))$, so
$\langle X,X\rangle=0$, $\langle Y,Y\rangle=\kappa|\kappa|^2\tau_w(h(y,y))=0$,
$\langle X,Y\rangle=|\kappa|^2\tau_w(h(x,y))=1$, and the columns of
$\tau_w(\iota)$ are orthogonal to $X$ and $Y$ and linearly independent. So
$\kappa\tau_w(J_1)=\tau_w(\iota)^{\mathsf T}M\overline{\tau_w(\iota)}$ has
signature $(r-1,s-1)$, and there are Cartan involutions $\theta_w$ of
$\kappa\tau_w(J)$ and $\theta_{1,w}$ of $\kappa\tau_w(J_1)$ with
$\theta_w\tau_w(\iota)=\tau_w(\iota)\theta_{1,w}$ and
$\kappa\tau_w(y)=2\theta_w\tau_w(x)$.

\emph{The places of $P$.} Let $v\in P$. The completion $(L_1)_{\tau_v}$ of
$L_1$ at the place given by $\tau_v$ is a free $\mathcal O_v$-module of rank
$n_1$, as $\mathcal O_v$ is a principal ideal domain and $L_1$ spans
$E^{n_1}$; let $\beta_{1,v}\in GL_{n_1}(F_v)$ have a basis of it as columns.
By the splitting and Lemma~\ref{asm-theta-agree-tower-pair}(3),
Lemma~\ref{asm-theta-agree-tower-basis} applies and gives
$\beta_v\in GL_n(F_v)$ with $\beta_v\mathcal O_v^n=L_{\tau_v}$, column $i_0$
equal to $\tau_v(x)$, column $i_1$ equal to $\tau_v(y)$, and the other
columns, in order, equal to the columns of $\tau_v(\iota)\beta_{1,v}$.

\emph{The spaces and the plane.} Let $Y_1$ be the lattice space with matrix
$J_1$, lattice $L_1$, the matrices $\theta_{1,w}$, and $b_{1,v}=\beta_{1,v}$
for $v\in P$ and $b_{1,v}=1$ otherwise. Let $Y^\circ$ be the lattice space
with matrix $J$, lattice $L$, the matrices $\theta_w$, and $b_v=\beta_v$ for
$v\in P$ and $b_v=b_{Y,v}$ otherwise; then $Y[Y^\circ]=Y^\circ$. It is
Hermitian, respectively skew-Hermitian, of signature $(r,s)$: its matrix and
lattice are those of $Y$, $\theta_w$ is a Cartan involution of
$\kappa\tau_w(J)$, and $\beta_v\mathcal O_v^n=L_{\tau_v}$ for $v\in P$. The
space $Y_1$ is Hermitian, respectively skew-Hermitian, of signature
$(r-1,s-1)$ by the last three paragraphs. Finally $(x,y,\iota,\mathfrak b)$ is
an adapted plane in $Y^\circ$ over $Y_1$ at positions $i_0<i_1$:
$h(x,x)=h(y,y)=0$ and $h(x,y)=1$ by Lemma~\ref{asm-theta-agree-tower-pair}(1);
$\iota$ is orthogonal to $x$ and $y$ and $\iota^{\mathsf T}J\,c(\iota)=J_1$;
$u\in L_1$ if and only if $\iota u\in L$, and $s\in\mathfrak b$ if and only if
$sx\in L$, by definition; the splitting is the one above; at each $w$,
$\theta_w\tau_w(\iota)=\tau_w(\iota)\theta_{1,w}$ and
$\tau_w(y)=2\kappa^{-1}\theta_w\tau_w(x)$; and the columns of $\beta_v$ are as
required.
\end{proof}

\begin{lemma}\label{asm-theta-agree-tower-eigen}
Let $F$ be a number field, $v$ a finite place of $F$, $\psi_v$ a nontrivial
continuous character of $F_v$, and $I$ a finite set. For
$w=(w_1,w_2)\in F_v^I\times F_v^I$ let $\rho_v(w)$ act on functions on $F_v^I$
by $(\rho_v(w)\phi)(x)=\psi_v(x\cdot w_2+\tfrac12w_1\cdot w_2)\,\phi(x+w_1)$,
and put $\langle w,w'\rangle=w_1\cdot w'_2-w'_1\cdot w_2$. Let
$A=\mathcal O_v^m$ for some $m$, let $\lambda\colon A\to F_v^I\times F_v^I$ be
additive and continuous, and let $q\colon A\to\mathbf C$ be continuous, with
$q(a)=1$ whenever $\lambda(a)=0$, and
\[
q(a+b)=q(a)\,q(b)\,\psi_v\bigl(-\tfrac12\langle\lambda(a),\lambda(b)\rangle\bigr)
\qquad(a,b\in A).
\]
Then there is a nonzero Schwartz-Bruhat function $\phi$ on $F_v^I$ with
$\rho_v(\lambda(a))\phi=q(a)\phi$ for every $a\in A$.
\end{lemma}
\leanref{Development/WeilBound/ThetaLifts/TowerModels/TwistedSums/ExistsSchroedingerEigenvectorCompactCocycle.lean}{Tw.VacLoc.exists_schr_eigen}
\begin{proof}
From the formula,
$\rho_v(w)\rho_v(w')=\psi_v(\frac12\langle w,w'\rangle)\rho_v(w+w')$, so
$\rho_v(w)\rho_v(w')=\psi_v(\langle w,w'\rangle)\rho_v(w')\rho_v(w)$; each
$\rho_v(w)$ is invertible and preserves the Schwartz-Bruhat functions. As
$\langle w,w\rangle=0$, $1=q(0)=q(a)q(-a)$, so $q(a)\neq0$, and
$\pi(a)=q(a)^{-1}\rho_v(\lambda(a))$ is defined, with
$\pi(a+b)=\pi(a)\pi(b)$.

\emph{Step 1: $q$ is trivial near $0$.} Choose an open subgroup $A_0$ of $A$
on which $\psi_v(-\frac12\langle\lambda(a),\lambda(b)\rangle)=1$ and
$|q-1|<\frac12$. On $A_0$, $q$ is a homomorphism, so $q(a)^k=q(ka)$ is within
$\frac12$ of $1$ for every $k$, and $q(a)=1$. The subgroup $A_0+\ker\lambda$
is open, its complement $A'$ in $A$ is compact, and $\lambda(A')$ is closed and
does not contain $0$; let $U$ be its complement. If $\lambda(a)\in U$, then
$a=a_0+a_1$ with $a_0\in A_0$ and $\lambda(a_1)=0$, and
$q(a)=q(a_0)q(a_1)=1$.

\emph{Step 2: a test function.} The indicator function $\phi_1$ of
$\mathcal O_v^I$ is a nonzero Schwartz-Bruhat function, and
$\rho_v(w)\phi_1=\phi_1$ for $w$ in a neighbourhood $U'$ of $0$: it suffices
that $w_1\in\mathcal O_v^I$ and that $\psi_v$ is trivial at
$x\cdot w_2+\frac12w_1\cdot w_2$ for every $x\in\mathcal O_v^I$.

\emph{Step 3: a dual subgroup.} Let $t_0\in F_v$ with $\psi_v(t_0)\neq1$,
choose $b\in F_v^\times$ with
$b(\mathcal O_v^I\times\mathcal O_v^I)\subset U\cap U'$, and put
$Y=t_0b^{-1}(\mathcal O_v^I\times\mathcal O_v^I)$. Let $e_i$ be the unit
vectors of $F_v^I$. If $w'\notin U\cap U'$, then some coordinate of $w'$ has
absolute value larger than $|b|$: if it is the coordinate $w'_{1,i}$ of
$w'_1$, then $w=(0,(t_0/w'_{1,i})e_i)$ lies in $Y$ and
$\langle w',w\rangle=t_0$, and if it is $w'_{2,i}$, then
$w=(-(t_0/w'_{2,i})e_i,0)$ does. So every $w'$ with
$\psi_v(\langle w',w\rangle)=1$ for all $w\in Y$ lies in $U\cap U'$. As $Y$ is
compact, there is a neighbourhood $U''$ of $0$ with
$\psi_v(\langle w',w\rangle)=1$ for all $w'\in U''$ and $w\in Y$.

\emph{Step 4: a finite quotient.} The set $K$ of the $a\in A$ with
$\psi_v(\langle\lambda(a),w\rangle)=1$ for all $w\in Y$ is a subgroup of $A$
containing the neighbourhood $\lambda^{-1}(U'')$ of $0$, so it is open and of
finite index in the compact group $A$. By Step 3, $\lambda(K)\subset U\cap U'$,
so $\pi(k)\phi_1=\phi_1$ for $k\in K$, by Steps 1 and 2. For $w\in Y$,
$\beta_w(a)=\psi_v(\langle\lambda(a),w\rangle)$ is a character of $A/K$; these
characters form a group, which separates the points of $A/K$ by the
definition of $K$; and $\pi(a)\rho_v(w)=\beta_w(a)\rho_v(w)\pi(a)$.

\emph{Step 5: averaging.} Let $a_1,\ldots,a_l$ represent $A/K$, with
$a_1\in K$, and put $c_j=\pi(a_j)\phi_1$, so that $c_1=\phi_1\neq0$. Some
$w_0\in Y$ has $\sum_j\beta_{w_0}(a_j)c_j\neq0$: otherwise, summing these
vanishing sums over the group of the characters $\beta_w$, whose sum vanishes
at every $a_j\notin K$, gives $Nc_1=0$ with $N$ the order of that group. The
function $\phi_0=\rho_v(w_0)\phi_1$ is fixed by $\pi(K)$, as
$\pi(k)\rho_v(w_0)\phi_1=\beta_{w_0}(k)\rho_v(w_0)\pi(k)\phi_1=\phi_0$. Hence
$\phi=\sum_j\pi(a_j)\phi_0$ does not depend on the representatives, and
$\pi(a)\phi=\phi$ for all $a\in A$, that is,
$\rho_v(\lambda(a))\phi=q(a)\phi$. It is a Schwartz-Bruhat function, and
$\phi=\sum_j\beta_{w_0}(a_j)\rho_v(w_0)\pi(a_j)\phi_1=\rho_v(w_0)\sum_j\beta_{w_0}(a_j)c_j\neq0$.
\end{proof}

\begin{lemma}\label{asm-theta-agree-tower-cocycle}
Let $E$ be a CM field with a theta base $B$
(Definition~\ref{asm-theta-agree-models}), and let $X$ be a Hermitian lattice
space of rank $n$ and $W$ a skew-Hermitian lattice space of rank $d$, of some
signatures, with matrices $J_X$ and $J_W$ and lattices $L_X$ and $L_W$. Let $h$
and $\langle\ ,\ \rangle$ be the forms on $\mathbb W(X,W)=M_{n\times d}(\mathbf A_E)$
of Definition~\ref{asm-theta-agree-models}(3), and write
$\operatorname{Tr}=\operatorname{Tr}_{\mathbf A_E/\mathbf A}$. Let $v$ be a
finite place, let $\Gamma_v$ be the group generated by the matrices
$a\,xy^{\mathsf T}$ with $a\in\mathbf A$ at $v$, $a_v\in\mathcal O_v$,
$x\in L_X$ and $y\in L_W$, let $\theta\in\mathcal O_E$ with
$\theta+c(\theta)=1$, and put
$q(\gamma)=\psi\bigl(\tfrac12\operatorname{Tr}h(\theta\gamma,\gamma)\bigr)$.
\begin{enumerate}
\item If $x_1,\ldots,x_k$ and $y_1,\ldots,y_l$ generate $L_X$ and $L_W$ as
  groups, then $\Gamma_v$ consists of the sums
  $\sum_{i,j}a_{ij}x_iy_j^{\mathsf T}$ with $a_{ij}\in\mathbf A$ at $v$ and
  $(a_{ij})_v\in\mathcal O_v$.
\item $\psi\bigl(b\operatorname{Tr}h(\xi\gamma,\gamma')\bigr)=1$ for
  $\gamma,\gamma'\in\Gamma_v$, $\xi\in\mathcal O_E$ and $b\in\mathbf A$ with
  $b_v\in\mathcal O_v$.
\item $q(\gamma+\gamma')=q(\gamma)\,q(\gamma')\,\psi\bigl(-\frac12\langle\gamma,\gamma'\rangle\bigr)$
  for $\gamma,\gamma'\in\Gamma_v$.
\end{enumerate}
\end{lemma}
\leanref{Development/WeilBound/ThetaLifts/TowerModels/TwistedSums/TensorLatticeAtFinitePlaceEqIntegralSpanGenerato.lean}{Tw.LatParam.latAtV_param}
\leanref{Development/WeilBound/ThetaLifts/TowerModels/TwistedSums/StandardAddCharSymplecticLocalLatticeEqOne.lean}{Tw.LatVacArith.psi_symp_latAtV}
\leanref{Development/WeilBound/ThetaLifts/TowerModels/TwistedSums/AddCharHalfTraceTensorLatticeHeisenbergCocycle.lean}{Tw.VacCocycle.psi_half_symp_add}
\begin{proof}
(1) If $x=\sum_im_ix_i$ and $y=\sum_jn_jy_j$ with integers $m_i$ and $n_j$,
then $a\,xy^{\mathsf T}=\sum_{i,j}m_in_ja\,x_iy_j^{\mathsf T}$. So the
generators of $\Gamma_v$ are such sums, these sums form a group, and each
term $a_{ij}x_iy_j^{\mathsf T}$ is a generator.

(2) The argument of $\psi$ is additive in $\gamma$ and in $\gamma'$, so it
suffices to take $\gamma=a\,xy^{\mathsf T}$ and $\gamma'=a'x'y'^{\mathsf T}$.
Then $h(\xi\gamma,\gamma')=aa'\xi\,h_X(x,x')\,h_W(y,y')$, with
$h_X(x,x')=x^{\mathsf T}J_Xc(x')$ and $h_W(y,y')=y^{\mathsf T}J_Wc(y')$, so
$b\operatorname{Tr}h(\xi\gamma,\gamma')=baa'\,t$ with
$t=\operatorname{Tr}_{E/F}(\xi h_X(x,x')h_W(y,y'))\in F$. As $L_X$ is its own
dual, $h_X(x,x')\in\mathcal O_E$, and as $L_W$ is its own dual for the trace,
$\operatorname{Tr}_{F/\mathbf Q}(ot)=\operatorname{Tr}_{E/\mathbf Q}(o\xi h_X(x,x')h_W(y,y'))\in\mathbf Z$
for $o\in\mathcal O_F$. So $t$ lies in the codifferent of $F$. The adele
$baa'$ is at $v$ with integral component. A standard character is trivial on
the products of an element $t$ of the codifferent with integral finite
adeles: it is trivial at the finite part of $ot$ for $o\in\mathcal O_F$, as
$\psi$ is trivial on $F$ and
$\psi_\infty(ot)=\exp(-2\pi i\operatorname{Tr}_{F/\mathbf Q}(ot))=1$, and
$\mathcal O_F$ is dense in the integral finite adeles.

(3) As $c(J_X)^{\mathsf T}=J_X$ and $c(J_W)^{\mathsf T}=-J_W$,
$c(h(\gamma',\gamma))=-h(\gamma,\gamma')$. So
\[
\operatorname{Tr}h(\theta\gamma',\gamma)=\operatorname{Tr}c\bigl(\theta h(\gamma',\gamma)\bigr)
=-\operatorname{Tr}\bigl((1-\theta)h(\gamma,\gamma')\bigr)
=\operatorname{Tr}h(\theta\gamma,\gamma')+\langle\gamma,\gamma'\rangle ,
\]
and
$\frac12\operatorname{Tr}h(\theta(\gamma+\gamma'),\gamma+\gamma')
=\frac12\operatorname{Tr}h(\theta\gamma,\gamma)+\frac12\operatorname{Tr}h(\theta\gamma',\gamma')
+\operatorname{Tr}h(\theta\gamma,\gamma')+\frac12\langle\gamma,\gamma'\rangle$.
By (2) with $\xi=\theta$ and with $\xi=-1$,
$\psi(\operatorname{Tr}h(\theta\gamma,\gamma'))=1$ and
$\psi(\langle\gamma,\gamma'\rangle)=1$, so
$\psi(\frac12\langle\gamma,\gamma'\rangle)=\psi(-\frac12\langle\gamma,\gamma'\rangle)$,
which gives (3).
\end{proof}

\begin{lemma}\label{asm-theta-agree-tower-vacuum}
Let $E$ be a CM field with a theta base $B$
(Definition~\ref{asm-theta-agree-models}), let $X$ be a Hermitian lattice space
of rank $n$ and $W$ a skew-Hermitian lattice space of rank $d$, of some
signatures, put $I=\{1,\ldots,n\}\times\{1,\ldots,d\}$, and let
$e\colon\mathbb W(X,W)\to\mathbf A^I\times\mathbf A^I$ be an $\mathbf A$-linear
isomorphism with $\langle w,w'\rangle=e_1(w)\cdot e_2(w')-e_1(w')\cdot e_2(w)$
that maps $M_{n\times d}(E)$ into $F^I\times F^I$. Call a function on $F_v^I$ a
vacuum at $v$ if it is a vacuum for the models of $(X,W)$ with the isomorphism
$e$ (Definition~\ref{asm-theta-agree-models}(3)). Then:
\begin{enumerate}
\item for every finite place $v$ there is a vacuum at $v$;
\item for almost all $v$ the indicator function of $\mathcal O_v^I$ is a
  vacuum at $v$.
\end{enumerate}
\end{lemma}
\leanref{Development/WeilBound/ThetaLifts/TowerModels/TwistedSums/ExistsFinitePlaceVacuumAnyDarbouxCoordinates.lean}{Tw.VacAny.exists_isVac_any_darboux}
\leanref{Development/WeilBound/ThetaLifts/TowerModels/TwistedSums/EventuallySelfDualLatticeIndicatorIsVacuum.lean}{Tw.VacIndic.eventually_isVac_indicator}
\begin{proof}
As $E/F$ is unramified at the finite places, the trace
$\mathcal O_E\to\mathcal O_F$ is onto, so there is $\theta\in\mathcal O_E$
with $\theta+c(\theta)=1$. Let $\rho$ be the Schrödinger operators of $e$, let
$\Gamma_v$ and $q$ be as in Lemma~\ref{asm-theta-agree-tower-cocycle} for this
$\theta$, let $x_1,\ldots,x_k$ and $y_1,\ldots,y_l$ generate $L_X$ and $L_W$ as
groups, and put $g_{ij}=x_iy_j^{\mathsf T}\in M_{n\times d}(E)$. For an adele,
or a vector of adeles, at $v$, the index $v$ marks its component at $v$.

(1) Fix $v$ and put $A=\mathcal O_v^{k\times l}$. For $a\in A$ let
$\gamma_a=\sum_{i,j}a_{ij}g_{ij}$, with $a_{ij}$ placed at $v$; by
Lemma~\ref{asm-theta-agree-tower-cocycle}(1), $\Gamma_v$ is the set of the
$\gamma_a$. As $e(g_{ij})\in F^I\times F^I$, $e(\gamma_a)=\sum a_{ij}e(g_{ij})$
is at $v$, and $\lambda(a)=e(\gamma_a)_v$ is additive and continuous in $a$; if
$\lambda(a)=0$, then $\gamma_a=0$ and $q(\gamma_a)=1$. The function
$a\mapsto q(\gamma_a)$ is continuous, and by
Lemma~\ref{asm-theta-agree-tower-cocycle}(3), as $\psi$ is $\psi_v$ on adeles
at $v$,
\[
q(\gamma_{a+b})=q(\gamma_a)\,q(\gamma_b)\,
\psi_v\bigl(-\tfrac12(\lambda(a)_1\cdot\lambda(b)_2-\lambda(b)_1\cdot\lambda(a)_2)\bigr).
\]
Lemma~\ref{asm-theta-agree-tower-eigen}, for the nontrivial character
$\psi_v$, gives a nonzero Schwartz-Bruhat function $\phi$ on $F_v^I$ with
$\rho_v(\lambda(a))\phi=q(\gamma_a)\phi$ for all $a\in A$. Let $r$ be a
function on $\mathbf A^I$ that does not depend on the $v$-component, and
$w=e(\gamma_a)$. Then
$(\rho(\gamma_a)(\phi\otimes r))(x)=\psi_v(x_v\cdot w_{2,v}+\frac12w_{1,v}\cdot w_{2,v})\,\phi(x_v+w_{1,v})\,r(x)$,
so $\rho(\gamma_a)$ acts at $v$ by
$\phi\mapsto\rho_v(\lambda(a))\phi=q(\gamma_a)\phi$. So $\phi$ is a vacuum at
$v$.

(2) For almost all $v$: $\psi_v$ is trivial on $\mathcal O_v$;
$\frac12\in\mathcal O_v$; the coordinates of the $e(g_{ij})$, which lie in $F$,
are integral at $v$; and the finitely many elements
$\operatorname{Tr}_{E/F}h(\theta g_{ij},g_{i'j'})$ of $F$ are integral at $v$.
Let $v$ be such a place and $\gamma=\gamma_a\in\Gamma_v$. Then $w=e(\gamma)$ is
at $v$ with integral $v$-components, and
$\frac12\operatorname{Tr}h(\theta\gamma,\gamma)
=\frac12\sum a_{ij}a_{i'j'}\operatorname{Tr}_{E/F}h(\theta g_{ij},g_{i'j'})$ is
at $v$ with integral $v$-component, so $q(\gamma)=1$. For
$x_v\in\mathcal O_v^I$, $x_v+w_{1,v}\in\mathcal O_v^I$ and
$\psi_v(x_v\cdot w_{2,v}+\frac12w_{1,v}\cdot w_{2,v})=1$; for
$x_v\notin\mathcal O_v^I$, $x_v+w_{1,v}\notin\mathcal O_v^I$. So $\rho(\gamma)$
acts at $v$ on the indicator function of $\mathcal O_v^I$ by the identity,
which is multiplication by $q(\gamma)$, and this nonzero Schwartz-Bruhat
function is a vacuum at $v$.
\end{proof}

\begin{lemma}\label{asm-theta-agree-tower-model}
Let $E$ be a CM field with a theta base $B$
(Definition~\ref{asm-theta-agree-models}), and let $X$ be a Hermitian lattice
space of rank $n$ and $W$ a skew-Hermitian lattice space of rank $d$, of some
signatures. Then there is a regular model for $(X,W)$
(Definition~\ref{asm-theta-agree-models}(3)).
\end{lemma}
\leanref{Development/WeilBound/ThetaLifts/TowerModels/PlaneHerm.lean}{PositivityTheory.ThKG.Tow.ModelEx}
\leanref{Development/WeilBound/ThetaLifts/TowerModels/ModelEx.lean}{PositivityTheory.ThKG.Tow.modelEx}
\leanref{Development/WeilBound/ThetaLifts/TowerModels/TwistedSums/ExistsRationalDarbouxHermitianTensorAdeles.lean}{Tw.RatDarboux.exists_rational_darboux}
\leanref{Development/WeilBound/ThetaLifts/CartanFactorisation/Unimod.lean}{PositivityTheory.ThKF.unimodularity_holds}
\begin{proof}
Let $J_X$ and $J_W$ be the matrices of $X$ and $W$, and
$I=\{1,\ldots,n\}\times\{1,\ldots,d\}$.

\emph{Darboux coordinates.} On the $F$-vector space $M_{n\times d}(E)$, of
dimension $2|I|$, the form $\langle x,y\rangle=-\operatorname{Tr}_{E/F}h(x,y)$
is alternating, since $c(h(x,x))=-h(x,x)$ by $c(J_X)^{\mathsf T}=J_X$ and
$c(J_W)^{\mathsf T}=-J_W$. It is nondegenerate: $J_X$ and $J_W$ are
invertible, as they have signatures,
$h(x,y)=\sum_{k,l}(J_X^{\mathsf T}xJ_W)_{kl}\,c(y_{kl})$, and the trace form of
$E/F$ is nondegenerate. So it has a symplectic basis $(f_k,f'_k)_{k\in I}$,
with $\langle f_k,f_l\rangle=\langle f'_k,f'_l\rangle=0$ and
$\langle f_k,f'_l\rangle=\delta_{kl}$. The form on
$\mathbb W(X,W)=\mathbf A\otimes_FM_{n\times d}(E)$ is its $\mathbf A$-bilinear
extension, so the coordinates $e=(e_1,e_2)$ in this basis are an
$\mathbf A$-linear isomorphism $\mathbb W(X,W)\to\mathbf A^I\times\mathbf A^I$
with $\langle w,w'\rangle=e_1(w)\cdot e_2(w')-e_1(w')\cdot e_2(w)$, and they
map $M_{n\times d}(E)$ into $F^I\times F^I$.

\emph{Measure and domain.} The group $U_X$ is a closed subgroup of
$GL_n(\mathbf A_E)$, so it is locally compact and second countable and has a
Haar measure $\mu$. Its subgroup $U_X(F)$ is countable, and discrete as $E$ is
discrete in $\mathbf A_E$; so it has a measurable fundamental domain
$\mathcal D$ in $U_X$ with respect to $\mu$.

\emph{Vacua.} By Lemma~\ref{asm-theta-agree-tower-vacuum}(1) choose a vacuum
$\phi_v$ at each finite place $v$ for the isomorphism $e$, and let
$\mathrm{vac}_v$ be the indicator function of $\mathcal O_v^I$ if this
function is a vacuum at $v$, and $\phi_v$ otherwise. By
Lemma~\ref{asm-theta-agree-tower-vacuum}(2), $\mathrm{vac}_v$ is that indicator
function for almost all $v$. So $(e,(\mathrm{vac}_v)_v,\mu,\mathcal D)$ is a
regular model for $(X,W)$.
\end{proof}

\begin{theorem}[Weil representation of a unitary dual pair \cite{Kud94,HKS96}]\label{asm-theta-agree-tower-weil}
Let $E$ be a CM field with a theta base $B$
(Definition~\ref{asm-theta-agree-models}), let $X$ be a Hermitian lattice space
of rank $n$ and $W$ a skew-Hermitian lattice space of rank $d$, of some
signatures, let $M$ be a regular model for $(X,W)$, and let $\chi_X$ and
$\chi_W$ be Hecke characters of $E$ with $\chi_X|_{\mathbf A^\times}=\eta^n$
and $\chi_W|_{\mathbf A^\times}=\eta^d$. Then there is a Weil representation
for $M$ with characters $(\chi_X,\chi_W)$
(Definition~\ref{asm-theta-agree-models}(6)).
\end{theorem}
\leanref{Development/WeilBound/ThetaLifts/TowerModels/PlaneHerm.lean}{PositivityTheory.ThKG.Tow.WeilEx}
\leanref{Development/WeilBound/ThetaLifts/Reductions/WeilPairs/WeilEx.lean}{PositivityTheory.ThKW.weilEx}
\leanref{Development/AbelianVarietyTheory/PublishedTheorems/Divisors₂/WeilRepresentationExists.lean}{GlobalTheory.Published.WeilRepresentationExists}

\begin{lemma}\label{asm-theta-agree-tower-alpha}
Let $(B,G,\mathcal X,\omega,\omega')$ be a triple of ranks $(b+2,a')$ for a theta
base $B$ of a CM field $E$ (Definition~\ref{asm-theta-agree-triple}), with
$q=0$ and $p\geq1$. Then there are lattice spaces $V^\circ$, $V'^\circ$ and
$W^\circ$ such that, with $G^\circ=G[V^\circ,V'^\circ,W^\circ]$,
$(B,G^\circ,\mathcal X,\omega,\omega')$ is a triple, and there are a triple datum
$G_1$ of ranks $(b,a')$, adapted planes $H$ and $H'$ such that
$(G^\circ,G_1,H,H')$ is an $\alpha$-configuration
(Definition~\ref{asm-theta-agree-planes}(4)), and representations $\omega_1$
and $\omega'_1$ such that $(B,G_1,\mathcal X,\omega_1,\omega'_1)$ is a triple.
\end{lemma}
\leanref{Development/AbelianVarietyTheory/Definitions/Divisors/CTSmallCone.lean}{GlobalTheory.ThK.TowerA}
\leanref{Development/WeilBound/ThetaLifts/TowerModels/TowerCodimensionTwoOf.lean}{PositivityTheory.ThKG.Tow.towerA_of}
\begin{proof}
Let $G=(p,q,V,V',W,M,M')$ and $\mathcal X=(\alpha,\beta,S_0,Q)$. As $G$ is
proper with $a=b+2$ and $q=0$, $2p+1<\lfloor(b+2-a')/2\rfloor$ and
$2p+1<a'$; with $p\geq1$ this gives $b-a'\geq6$, so $b\geq1$ and
$b+2-p\geq1$.

\emph{The planes.} The space $V$ is Hermitian of signature $(b+2-p,p)$ and
$W$ is skew-Hermitian of signature $(b+2+a'-p,p)$, both of the form $(r,s)$
with $r,s\geq1$. Proposition~\ref{asm-theta-agree-tower-plane} for $V$, with
$n_1=b$ and positions $0<b+1$, gives $V^\circ$ and a Hermitian lattice space
$V_1$ of signature $(b+1-p,p-1)$ such that $V[V^\circ]$ is Hermitian of
signature $(b+2-p,p)$ and there is an adapted plane $H$ in $V[V^\circ]$ over
$V_1$ at positions $0<b+1$. For $W$, with $n_1=b+a'$ and positions $0<b+1$,
it gives $W^\circ$ and a skew-Hermitian lattice space $W_1$ of signature
$(b+1+a'-p,p-1)$ such that $W[W^\circ]$ is skew-Hermitian of signature
$(b+2+a'-p,p)$ and there is an adapted plane $H'$ in $W[W^\circ]$ over $W_1$
at positions $0<b+1$. Put $V'^\circ=V'$ and
$G^\circ=G[V^\circ,V'^\circ,W^\circ]$, whose spaces are $V[V^\circ]$, $V'$ and
$W[W^\circ]$.

\emph{The re-chosen triple.} The datum $G^\circ$ has the integers, the ranks
and the models of $G$, and its spaces have the matrices, the lattices and the
signatures of those of $G$. Regularity of a model and the conditions on a
Weil representation involve the spaces only through their matrices and
lattices. So $G^\circ$ is proper for $B$, and
$(B,G^\circ,\mathcal X,\omega,\omega')$ is a triple.

\emph{The smaller triple.} By Lemma~\ref{asm-theta-agree-tower-model} there
are regular models $M_1$ for $(V_1,W_1)$ and $M'_1$ for $(V',W_1)$. As
$\eta^2=1$, $\beta|_{\mathbf A^\times}=\eta^{b+2}=\eta^b$,
$(\beta/\alpha)|_{\mathbf A^\times}=\eta^{b+a'}$,
$\alpha|_{\mathbf A^\times}=\eta^{a'}$ and
$(\alpha/\beta)|_{\mathbf A^\times}=\eta^{a'+b}$, and $V_1$, $V'$ and $W_1$
have ranks $b$, $a'$ and $b+a'$. By Theorem~\ref{asm-theta-agree-tower-weil}
there are a Weil representation $\omega_1$ for $M_1$ with characters
$(\beta,\beta/\alpha)$ and a Weil representation $\omega'_1$ for $M'_1$ with
characters $(\alpha,\alpha/\beta)$. Let $G_1=(p-1,0,V_1,V',W_1,M_1,M'_1)$, a
triple datum of ranks $(b,a')$. It is proper for $B$: $a'<b$; $b+a'$ is even,
as $b+2+a'$ is; $b+a'<b+2+a'<m_\alpha(w)-m_\beta(w)$ for every $w$;
$2(p-1)+1<\lfloor(b-a')/2\rfloor$, as $2p+1<\lfloor(b-a')/2\rfloor+1$;
$2(p-1)+1<2p+1<a'$; the spaces $V_1$, $V'$ and $W_1$ have signatures
$(b-(p-1),p-1)$, $(0,a')$ and $(b+a'-(p-1),p-1)$; and $M_1$ and $M'_1$ are
regular. The character datum $\mathcal X$ is proper for $(B,b,a')$, since
$\beta|_{\mathbf A^\times}=\eta^b$ and its other conditions do not involve the
ranks. So $(B,G_1,\mathcal X,\omega_1,\omega'_1)$ is a triple.

\emph{The configuration.} The data $G^\circ$ and $G_1$ have ranks $(b+2,a')$
and $(b,a')$ and the same second space $V'$, with $q=q_1=0$ and $(p-1)+1=p$;
$H$ is in the first space $V[V^\circ]$ of $G^\circ$ over the first space $V_1$
of $G_1$ at positions $(0,b+1)$, and $H'$ is in $W[W^\circ]$ over $W_1$ at
positions $(0,b+1)$. So $(G^\circ,G_1,H,H')$ is an $\alpha$-configuration.
% Lean: PositivityTheory.ThKG.Tow.towerA_of Tow.planeEx Tow.modelEx ThKW.weilEx, as closeSph
% applies it; the three inputs are Proposition asm-theta-agree-tower-plane, Lemma
% asm-theta-agree-tower-model and Theorem asm-theta-agree-tower-weil.
\end{proof}

\begin{lemma}\label{asm-theta-agree-tower-beta}
Let $(B,G,\mathcal X,\omega,\omega')$ be a triple of ranks $(a,b+2)$ for a theta
base $B$ of a CM field $E$ (Definition~\ref{asm-theta-agree-triple}), with
$q\geq1$. Then there are lattice spaces $V^\circ$, $V'^\circ$ and $W^\circ$ such
that, with $G^\circ=G[V^\circ,V'^\circ,W^\circ]$,
$(B,G^\circ,\mathcal X,\omega,\omega')$ is a triple, and there are a triple datum
$G_1$ of ranks $(a,b)$, adapted planes $H$ and $H'$ such that
$(G^\circ,G_1,H,H')$ is a $\beta$-configuration
(Definition~\ref{asm-theta-agree-planes}(4)), and representations $\omega_1$
and $\omega'_1$ such that $(B,G_1,\mathcal X,\omega_1,\omega'_1)$ is a triple.
\end{lemma}
\leanref{Development/AbelianVarietyTheory/Definitions/Divisors/CTSmallCone.lean}{GlobalTheory.ThK.TowerB}
\leanref{Development/WeilBound/ThetaLifts/TowerModels/TowerCodimensionTwoOf.lean}{PositivityTheory.ThKG.Tow.towerB_of}
\begin{proof}
Let $G=(p,q,V,V',W,M,M')$ and $\mathcal X=(\alpha,\beta,S_0,Q)$. As $G$ is
proper with $a'=b+2$, $2(p+q)+1<\lfloor(a-b-2)/2\rfloor$ and
$2(p+q)+1<b+2$; with $q\geq1$ this gives $b\geq2$ and $b+2-q\geq1$.

\emph{The planes.} The space $V'$ is Hermitian of signature $(q,b+2-q)$ and
$W$ is skew-Hermitian of signature $(a+b+2-p-q,p+q)$, both of the form
$(r,s)$ with $r,s\geq1$. Proposition~\ref{asm-theta-agree-tower-plane} for
$V'$, with $n_1=b$ and positions $0<b+1$, gives $V'^\circ$ and a Hermitian
lattice space $V'_1$ of signature $(q-1,b+1-q)$ such that $V'[V'^\circ]$ is
Hermitian of signature $(q,b+2-q)$ and there is an adapted plane $H$ in
$V'[V'^\circ]$ over $V'_1$ at positions $0<b+1$. For $W$, with $n_1=a+b$ and
positions $a<a+b+1$, it gives $W^\circ$ and a skew-Hermitian lattice space
$W_1$ of signature $(a+b+1-p-q,p+q-1)$ such that $W[W^\circ]$ is
skew-Hermitian of signature $(a+b+2-p-q,p+q)$ and there is an adapted plane
$H'$ in $W[W^\circ]$ over $W_1$ at positions $a<a+b+1$. Put $V^\circ=V$ and
$G^\circ=G[V^\circ,V'^\circ,W^\circ]$, whose spaces are $V$, $V'[V'^\circ]$ and
$W[W^\circ]$. As in the proof of Lemma~\ref{asm-theta-agree-tower-alpha},
$(B,G^\circ,\mathcal X,\omega,\omega')$ is a triple.

\emph{The smaller triple.} By Lemma~\ref{asm-theta-agree-tower-model} there
are regular models $M_1$ for $(V,W_1)$ and $M'_1$ for $(V'_1,W_1)$. As
$\eta^2=1$, $\beta|_{\mathbf A^\times}=\eta^a$,
$(\beta/\alpha)|_{\mathbf A^\times}=\eta^{a+b}$,
$\alpha|_{\mathbf A^\times}=\eta^{b+2}=\eta^b$ and
$(\alpha/\beta)|_{\mathbf A^\times}=\eta^{b+a}$, and $V$, $V'_1$ and $W_1$
have ranks $a$, $b$ and $a+b$. By Theorem~\ref{asm-theta-agree-tower-weil}
there are a Weil representation $\omega_1$ for $M_1$ with characters
$(\beta,\beta/\alpha)$ and a Weil representation $\omega'_1$ for $M'_1$ with
characters $(\alpha,\alpha/\beta)$. Let $G_1=(p,q-1,V,V'_1,W_1,M_1,M'_1)$, a
triple datum of ranks $(a,b)$. It is proper for $B$: $b<b+2<a$; $a+b$ is
even, as $a+b+2$ is; $a+b<m_\alpha(w)-m_\beta(w)$ for every $w$;
$2(p+q-1)+1<\lfloor(a-b)/2\rfloor$, as
$2(p+q)+1<\lfloor(a-b)/2\rfloor-1$; $2(p+q-1)+1<b$, as $2(p+q)+1<b+2$; the
spaces $V$, $V'_1$ and $W_1$ have signatures $(a-p,p)$, $(q-1,b-(q-1))$ and
$(a+b-(p+q-1),p+q-1)$; and $M_1$ and $M'_1$ are regular. The character datum
$\mathcal X$ is proper for $(B,a,b)$, since $\alpha|_{\mathbf A^\times}=\eta^b$
and its other conditions do not involve the ranks. So
$(B,G_1,\mathcal X,\omega_1,\omega'_1)$ is a triple.

\emph{The configuration.} The data $G^\circ$ and $G_1$ have ranks $(a,b+2)$
and $(a,b)$, the same first space $V$ and the same $p$, with $(q-1)+1=q$; $H$
is in the second space $V'[V'^\circ]$ of $G^\circ$ over the second space
$V'_1$ of $G_1$ at positions $(0,b+1)$, and $H'$ is in $W[W^\circ]$ over $W_1$
at positions $(a,a+b+1)$. So $(G^\circ,G_1,H,H')$ is a
$\beta$-configuration.
% Lean: PositivityTheory.ThKG.Tow.towerB_of Tow.planeEx Tow.modelEx ThKW.weilEx, as closeSph
% applies it.
\end{proof}

\subsection{Cone integrals and the seesaw}\label{subsection-asm-cone-integrals-seesaw}

\begin{definition}\label{asm-theta-agree-split-notions}
Let $E$ be a CM field, Galois over $\mathbf Q$, with maximal totally real
subfield $F$ and complex conjugation $c$, let $B$ be a theta base for $E$
(Definition~\ref{asm-theta-agree-models}(1)), and let $X$ and $W$ be lattice
spaces (Definition~\ref{asm-theta-agree-doubling}(1)) of ranks $n$ and $d$.
Write $h_X(u,u')=u^{\mathsf T}J_X\,c(u')$ for $u,u'\in E^n$, where $J_X$ is
the matrix of $X$, and $|l|=|N_{E/F}(l)|_{\mathbf A}$ for
$l\in\mathbf A_E^\times$.
\begin{enumerate}
\item The set $\operatorname{Iso}(X)$ of \emph{isotropic vectors} of $X$
  consists of the $u\in E^n$ with $u\neq0$ and $h_X(u,u)=0$. The group
  $E^\times$ acts on it by scalar multiplication, and the action is free. For
  $x\in E^n$, a \emph{transporter family} for $x$ is a map
  $\gamma\colon\operatorname{Iso}(X)\to U_X(F)$ with $\gamma(u)u=x$ for every
  $u\in\operatorname{Iso}(X)$.
\item Let $x\in E^n$. A subgroup $K\subset U_X$ is an \emph{Iwasawa subgroup}
  for $x$ if $K$ is compact and every $g\in U_X$ is $pk$ with $k\in K$ and
  $p\in U_X$ such that $px=lx$ for some $l\in\mathbf A_E^\times$.
\item Let $\omega$ be a representation of $U_X\times U_W$ on functions on
  $\mathbf A^{n\times d}$ and $K\subset U_X$ a subgroup. A function $\Phi$ on
  $\mathbf A^{n\times d}$ is \emph{$K$-finite} for $\omega$ if the functions
  $\omega(k,1)\Phi$, $k\in K$, span a finite-dimensional space.
\item Let $M$ be a model for $(X,W)$ with Schrödinger operators $\rho$
  (Definition~\ref{asm-theta-agree-models}(3)) and $\omega$ a representation
  as in (3). Let $X_1$ and $W_1$ be lattice spaces of ranks $n_1$ and $d_1$,
  $M_0$ a model for $(X_1,W_1)$ and $\omega_0$ a representation of
  $U_{X_1}\times U_{W_1}$ on functions on $\mathbf A^{n_1\times d_1}$, and let
  $e$ be a linear map from functions on $\mathbf A^{n\times d}$ to functions
  on $\mathbf A^{n_1\times d_1}$. Let $K\subset U_X$ be a subgroup with a
  measure $\mu_K$, $\mu$ a measure on $\mathbf A_E^\times$,
  $u\colon\mathbf A_E^\times\to\mathbb W(X,W)$, $\chi$ a character of
  $\mathbf A_E^\times$ and $z\geq0$ an integer. The \emph{cone integral} for
  $(M,\omega,M_0,\omega_0,e,K,\mu_K,\mu,u,\chi,z)$ of a function $\Phi$ on
  $\mathbf A^{n\times d}$ is the function
  \[
  I^{\mathrm{cone}}(\Phi)(h_1)=\int_K\int_{\mathbf A_E^\times}I_{M_0,\omega_0}\bigl(e(\rho(u(l))\,\omega(k,1)\Phi)\bigr)(h_1)\,\chi(l)\,|l|^z\,d\mu(l)\,d\mu_K(k)
  \]
  on $U_{W_1}$, with the theta integral of
  Definition~\ref{asm-theta-agree-models}(7); here, as in the theta
  integral, the integral of a function that is not integrable is taken to be
  $0$.
\end{enumerate}
\end{definition}
\leanref{Development/WeilBound/ThetaLifts/SmallCones/ThetaUniq.lean}{PositivityTheory.ThKCT.IsoV}
\leanref{Development/WeilBound/ThetaLifts/SmallCones/IsoActFree.lean}{PositivityTheory.ThKCT.isoAct}
\leanref{Development/WeilBound/ThetaLifts/SmallCones/ThetaUniq.lean}{PositivityTheory.ThKCT.Moves}
\leanref{Development/WeilBound/ThetaLifts/SmallCones/ThetaUniq.lean}{PositivityTheory.ThKCT.Iwa}
\leanref{Development/WeilBound/ThetaLifts/SmallCones/ThetaUniq.lean}{PositivityTheory.ThKCT.KFinite}
\leanref{Development/WeilBound/ThetaLifts/SmallCones/ThetaUniq.lean}{PositivityTheory.ThKCT.coneI}

\begin{lemma}\label{asm-theta-agree-split-low}
Let $E$ be a CM field with a theta base $B$
(Definition~\ref{asm-theta-agree-models}(1)), and let $H=(x,y,\iota,\mathfrak b)$
be an adapted plane in a lattice space $Y$ over a lattice space $Y_1$ at some
positions (Definition~\ref{asm-theta-agree-planes}(1)). If $Y$ is Hermitian
of signature $(r,s)$, then $r,s\geq1$ and $Y_1$ is Hermitian of signature
$(r-1,s-1)$. The same holds with skew-Hermitian in place of Hermitian.
\end{lemma}
\leanref{Development/WeilBound/ThetaLifts/SmallCones/ThetaUniq.lean}{PositivityTheory.ThKCT.Low}
\leanref{Development/WeilBound/ThetaLifts/SmallCones/SplitLatMap.lean}{PositivityTheory.ThKCT.low_of}
\leanref{Development/WeilBound/ThetaLifts/SmallCones/SplitLatMap.lean}{PositivityTheory.ThKCT.low_eps}
\leanref{Development/WeilBound/ThetaLifts/SmallCones/SplitLatMap.lean}{PositivityTheory.ThKCT.cartan_low}
\begin{proof}
Let $m$ and $m_1=m-2$ be the ranks of $Y$ and $Y_1$, $J$ and $J_1$ their
matrices, $L$ and $L_1$ their lattices, and $h(u,u')=u^{\mathsf T}J\,c(u')$.
In the Hermitian case put $\epsilon=1$, $\kappa=1$ and $D=\mathcal O_E$; in
the skew-Hermitian case put $\epsilon=-1$, $\kappa=i$ and let $D$ be the
codifferent
$\{t\in E:\operatorname{Tr}_{E/\mathbf Q}(ot)\in\mathbf Z\text{ for all }o\in\mathcal O_E\}$.
Then $c(J)^{\mathsf T}=\epsilon J$, $L=\{u:h(u,a)\in D\text{ for all }a\in L\}$,
and the matrix $M_w$ of Definition~\ref{asm-theta-agree-doubling}(1) is
$\kappa\,\tau_w(J)$. As $\tau_w\circ c$ is complex conjugation on $\tau_w(E)$,
$\tau_w(u)^{\mathsf T}M_w\overline{\tau_w(u')}=\kappa\,\tau_w(h(u,u'))$.

\emph{A left inverse.} The vectors $x$, $y$ and the columns of $\iota$ form a
basis of $E^m$: every element of $L$ is $sx+\iota u+ty$, and $L$ spans $E^m$,
which has dimension $m_1+2$. So there is $P\in M_{m_1\times m}(E)$ with
$P(\iota u+ax+by)=u$ for all $u$, $a$ and $b$; in particular $P\iota=1$, and
$P$ maps $L$ into $L_1$, since $z=sx+\iota u+ty\in L$ has $u\in L_1$.

\emph{The signature is indefinite.} If $\min(r,s)=0$, the Hermitian matrix
$M_w$ is definite for every infinite place $w$, and
$\tau_w(x)^{\mathsf T}M_w\overline{\tau_w(x)}=\kappa\,\tau_w(h(x,x))=0$ forces
$\tau_w(x)=0$, so $x=0$, contradicting $h(x,y)=1$. So $r,s\geq1$.

\emph{The form.} By the definition of an adapted plane,
$J_1=\iota^{\mathsf T}J\,c(\iota)$, so $c(J_1)^{\mathsf T}=\epsilon J_1$, and
$h_1(u,u')=h(\iota u,\iota u')$ for the form $h_1$ of $J_1$.

\emph{The archimedean places.} Fix $w$, put $\iota_w=\tau_w(\iota)$,
$\theta=\theta_{Y,w}$ and $\theta_1=\theta_{Y_1,w}$, and write
$\langle u,u'\rangle=u^{\mathsf T}M_w\bar u'$. Then
$\langle\tau_w(x),\tau_w(x)\rangle=0$,
$\langle\kappa\tau_w(y),\kappa\tau_w(y)\rangle=0$,
$\langle\tau_w(x),\kappa\tau_w(y)\rangle=\bar\kappa\kappa=1$,
$\langle\iota_wu,\tau_w(x)\rangle=\langle\iota_wu,\kappa\tau_w(y)\rangle=0$,
and $\iota_w$ has linearly independent columns, as
$\tau_w(P)\iota_w=1$. By Lemma~\ref{asm-theta-agree-tower-cartan}, the matrix
$\iota_w^{\mathsf T}M_w\bar\iota_w=\kappa\,\tau_w(J_1)$, which is the matrix
$M_{1,w}$ of $Y_1$, has signature $(r-1,s-1)$. The adapted plane gives
$\theta\iota_w=\iota_w\theta_1$, and $\theta$ is a Cartan involution of $M_w$.
Then $\theta_1$ is one of $M_{1,w}$:
$\theta_1^2=\tau_w(P)\iota_w\theta_1^2=\tau_w(P)\theta^2\iota_w=1$;
$\theta_1^{\mathsf T}M_{1,w}\bar\theta_1=(\theta\iota_w)^{\mathsf T}M_w\overline{\theta\iota_w}=\iota_w^{\mathsf T}(\theta^{\mathsf T}M_w\bar\theta)\bar\iota_w=M_{1,w}$;
and for $u\neq0$,
$\operatorname{Re}((\theta_1u)^{\mathsf T}M_{1,w}\bar u)=\operatorname{Re}((\theta\iota_wu)^{\mathsf T}M_w\overline{\iota_wu})>0$,
as $\iota_wu\neq0$.

\emph{The lattice.} By the definition of an adapted plane,
$L_1=\{u:\iota u\in L\}$, and $L_1=P(L)$, since $u=P\iota u$. So $L_1$ is an
$\mathcal O_E$-lattice in $E^{m_1}$. If $u,u'\in L_1$, then
$h_1(u,u')=h(\iota u,\iota u')\in D$. Conversely, let $u$ satisfy
$h_1(u,u')\in D$ for every $u'\in L_1$. For $z=sx+\iota u'+ty\in L$,
\[
h(\iota u,z)=c(s)\,h(\iota u,x)+h(\iota u,\iota u')+c(t)\,h(\iota u,y)=h_1(u,u')\in D,
\]
so $\iota u\in L$ and $u\in L_1$. Hence
$L_1=\{u:h_1(u,u')\in D\text{ for all }u'\in L_1\}$.

\emph{The places in $P$.} Let $v\in P$, $\tau=\tau_v$, $\beta=b_{Y,v}$ and
$\beta_1=b_{Y_1,v}$, let $H$ be at the positions $i_0<i_1$, and let $T$ be the
$m\times m_1$ matrix whose columns are those of the identity matrix other than
the columns $i_0$ and $i_1$. The adapted plane gives
$\beta T=\tau(\iota)\beta_1$. For $u\in F_v^{m_1}$, $u$ is integral if and
only if $Tu$ is, if and only if $\beta Tu$ lies in the completion $L_\tau$ of
$L$, since $\beta\mathcal O_v^m=L_\tau$. For $u'\in F_v^{m_1}$,
$\tau(\iota)u'\in L_\tau$ if and only if $u'$ lies in the completion
$(L_1)_\tau$ of $L_1$: one direction holds as $\iota L_1\subset L$, the other
as $\tau(P)$ maps $L_\tau$ into $(L_1)_\tau$ and $\tau(P)\tau(\iota)=1$.
Together, $\beta_1u\in(L_1)_\tau$ if and only if $u$ is integral, and as
$\beta_1$ is invertible, $\beta_1\mathcal O_v^{m_1}=(L_1)_\tau$.

So $Y_1$ is Hermitian, respectively skew-Hermitian, of signature $(r-1,s-1)$.
\end{proof}

\begin{lemma}\label{asm-theta-agree-split-witt}
Let $E$ be a CM field with a theta base $B$
(Definition~\ref{asm-theta-agree-models}(1)), and let $X$ be a Hermitian
lattice space of signature $(r,s)$ (Definition~\ref{asm-theta-agree-doubling}(1)),
with isotropic vectors $\operatorname{Iso}(X)$
(Definition~\ref{asm-theta-agree-split-notions}(1)).
\begin{enumerate}
\item If $\min(r,s)=0$, then $\operatorname{Iso}(X)$ is empty.
\item For every $x\in\operatorname{Iso}(X)$ there is a transporter family for
  $x$.
\end{enumerate}
\end{lemma}
\leanref{Development/WeilBound/ThetaLifts/SmallCones/ThetaUniq.lean}{PositivityTheory.ThKCT.Witt}
\leanref{Development/WeilBound/ThetaLifts/SmallCones/WittOf.lean}{PositivityTheory.ThKCT.witt_of}
\leanref{Development/WeilBound/ThetaLifts/SmallCones/WittOf.lean}{PositivityTheory.ThKCT.aniso}
\leanref{Development/WeilBound/ThetaLifts/EigenvalueEndgame/UnipotentGroups/IsFrPlane.lean}{PositivityTheory.ThKEnd.Grp.wittVec}
\leanref{Development/WeilBound/ThetaLifts/EigenvalueEndgame/UnipotentGroups/IsFrPlane.lean}{PositivityTheory.ThKEnd.Grp.witt_step}
\begin{proof}
Let $n$, $J$ and $h(u,u')=u^{\mathsf T}J\,c(u')$ be the rank, the matrix and
the form of $X$. Then $h(u',u)=c(h(u,u'))$, and $J$ is invertible
(Lemma~\ref{asm-theta-agree-convergence-signature}).

(1) Let $w$ be an infinite place. If $\min(r,s)=0$, the Hermitian matrix
$\tau_w(J)$ is definite, and $u\in\operatorname{Iso}(X)$ would give
$\tau_w(u)^{\mathsf T}\tau_w(J)\overline{\tau_w(u)}=\tau_w(h(u,u))=0$, so
$\tau_w(u)=0$ and $u=0$.

(2) It suffices to find, for $u,x\in\operatorname{Iso}(X)$, an element
$g\in U_X(F)$ with $gu=x$; a choice of one for each $u$ is a transporter
family.

\emph{A transvection.} Suppose first that $b=h(u,x)\neq0$. Put $w=u-x$ and
$g(v)=v+b^{-1}h(v,w)\,w$. Then $h(u,w)=h(u,u)-h(u,x)=-b$, so
$gu=u-w=x$. For $v,v'\in E^n$, using $h(w,v')=c(h(v',w))$,
\[
h(gv,gv')=h(v,v')+h(v,w)\,h(w,v')\bigl(c(b)^{-1}+b^{-1}+b^{-1}c(b)^{-1}h(w,w)\bigr),
\]
and $h(w,w)=-h(u,x)-h(x,u)=-b-c(b)$, so the bracket is $0$. Hence
$g^{\mathsf T}J\,c(g)=J$; as $J$ is invertible, so is $g$, and $g\in U_X(F)$.

\emph{Two transvections.} Suppose that $h(u,x)=0$. As $J$ is invertible there
are $a,a'$ with $h(u,a)\neq0$ and $h(x,a')\neq0$; one of $a$, $a'$ and $a+a'$
is a vector $v$ with $h(u,v)\neq0$ and $h(x,v)\neq0$. Put
$y=v+\lambda u$ with $\lambda=-h(v,v)/(2h(u,v))$. As $h(v,v)\in F$,
\[
h(y,y)=h(v,v)+\lambda h(u,v)+c(\lambda h(u,v))=0,
\]
and $h(u,y)=h(u,v)\neq0$, $h(y,x)=h(v,x)\neq0$; in particular $y\neq0$. Two
transvections as above take $u$ to $y$ and $y$ to $x$, and their product is
the element $g$.
\end{proof}

\begin{lemma}\label{asm-theta-agree-split-inputs-characters}
Let $A$ be an abelian group, $T$ a set and $(e_i)_{i\in T}$ functions
$A\to\mathbf C$ with $e_i(a+b)=e_i(a)e_i(b)$ and $|e_i(a)|\leq1$ for all $i$,
$a$ and $b$. Let $i_0\in T$ with $e_{i_0}=1$, and suppose that $e_i\neq1$ for
$i\neq i_0$. Let $(c_i)_{i\in T}$ be complex numbers with
$\sum_i|c_i|<\infty$ and $\sum_ic_ie_i(a)=0$ for every $a\in A$. Then
$c_{i_0}=0$.
\end{lemma}
\leanref{Development/AutomorphicSpectrum/Proofs/MembershipProofs/HeisenbergCommutingOperatorScalar.lean}{GlobalTheory.CompanionTheory.ThetaInj.charSeries_coeff_eq_zero}
\leanref{Development/AutomorphicSpectrum/Proofs/MembershipProofs/HeisenbergCommutingOperatorScalar.lean}{GlobalTheory.CompanionTheory.ThetaInj.charSeries_restrict}
\leanref{Development/AutomorphicSpectrum/Proofs/MembershipProofs/HeisenbergCommutingOperatorScalar.lean}{GlobalTheory.CompanionTheory.ThetaInj.charSeries_remove}
\begin{proof}
\emph{Restriction.} Let $t\in A$. Then $\sum_{e_i(t)=1}c_ie_i(a)=0$ for every
$a$, and the coefficients $c_i$ for $e_i(t)=1$ and $0$ otherwise again satisfy
the hypotheses. Put $q_i=(1+e_i(t))/2$, so that $|q_i|\leq1$, $q_i=1$ if
$e_i(t)=1$, and $|q_i|<1$ otherwise, as $|e_i(t)|\leq1$. By induction on $k$,
$\sum_ic_iq_i^ke_i(a)=0$ for every $a$: the case $k=0$ is the hypothesis, and
$c_iq_i^{k+1}e_i(a)=\frac12\bigl(c_iq_i^ke_i(a)+c_iq_i^ke_i(a+t)\bigr)$. The
terms are bounded by $|c_i|$, and $q_i^k$ tends to $1$ if $e_i(t)=1$ and to $0$
otherwise. By dominated convergence the sums tend to
$\sum_{e_i(t)=1}c_ie_i(a)$, which is therefore $0$.

\emph{Conclusion.} Suppose $c_{i_0}\neq0$, and choose a finite set $S\ni i_0$
with $\sum_{i\notin S}|c_i|<|c_{i_0}|$. For each $i\in S$ other than $i_0$
choose $t_i$ with $e_i(t_i)\neq1$. Applying the restriction for each $t_i$ in
turn gives coefficients $c'_i$, each equal to $c_i$ or $0$, with
$c'_{i_0}=c_{i_0}$, as $e_{i_0}=1$, with $c'_i=0$ for $i\in S\setminus\{i_0\}$,
and with $\sum_ic'_ie_i(a)=0$ for every $a$. At $a=0$ this gives
$c_{i_0}=-\sum_{i\notin S}c'_ie_i(0)$, so $|c_{i_0}|\leq\sum_{i\notin S}|c_i|$,
a contradiction.
\end{proof}

\begin{definition}\label{asm-theta-agree-split-inputs-polar}
Let $B$ be a theta base for a CM field $E$
(Definition~\ref{asm-theta-agree-models}(1)), let $X$, $W$, $X_1$ and $W_1$ be
lattice spaces of ranks $n$, $d$, $n_1$ and $d_1$, with the forms
$\langle\,,\rangle$ on $\mathbb W(X,W)$ and $\langle\,,\rangle_1$ on
$\mathbb W(X_1,W_1)$ of Definition~\ref{asm-theta-agree-models}(3), let
$\iota_X\in M_{n\times n_1}(E)$ and $\iota_W\in M_{d\times d_1}(E)$, and let
$\mathcal J\subset\mathbb W(X,W)$. For $w\in\mathbb W(X_1,W_1)$ write
$\bar w=\iota_Xw\iota_W^{\mathsf T}$. A \emph{polar system} for
$(\iota_X,\iota_W,\mathcal J)$ consists of an integer $m\geq0$ and matrices
$u_1,\ldots,u_m,z_1,\ldots,z_m\in M_{n\times d}(E)$ such that:
\begin{enumerate}
\item $az_j\in\mathcal J$ for every $a\in\mathbf A$ and every $j$, and
  $\langle z_i,u_j\rangle=\delta_{ij}$;
\item $\langle u_i,u_j\rangle=\langle z_i,z_j\rangle=0$,
  $\langle u_i,\bar w\rangle=\langle z_i,\bar w\rangle=0$ and
  $\langle\bar w,\bar w'\rangle=\langle w,w'\rangle_1$ for all $i$, $j$ and
  $w,w'\in\mathbb W(X_1,W_1)$;
\item every $x\in M_{n\times d}(E)$ is
  $\sum_ja_ju_j+\sum_jb_jz_j+\iota_Xx_1\iota_W^{\mathsf T}$ for some
  $a_j,b_j\in F$ and $x_1\in M_{n_1\times d_1}(E)$.
\end{enumerate}
For $y\in\mathbf A^m$ put $u(y)=\sum_jy_ju_j$ and $z(y)=\sum_jy_jz_j$.
\end{definition}
\leanref{Development/WeilBound/ThetaLifts/SmallCones/ThetaUniq.lean}{PositivityTheory.ThKCT.Dual}
\leanref{Development/WeilBound/ThetaLifts/SmallCones/ThetaUniq.lean}{PositivityTheory.ThKCT.Polar}
\leanref{Development/WeilBound/ThetaLifts/SmallCones/ThetaUniq.lean}{PositivityTheory.ThKCT.lin}

\begin{lemma}\label{asm-theta-agree-split-inputs-expansion}
Let $B$ be a theta base for a CM field $E$ with character $\psi$, let $M$ and
$M_1$ be regular models for lattice spaces $(X,W)$ and $(X_1,W_1)$, of ranks
$(n,d)$ and $(n_1,d_1)$, with Schrödinger operators $\rho$ and $\rho_1$
(Definition~\ref{asm-theta-agree-models}(3)), let $e$ be an evaluation map
for $(M,M_1,\iota_X,\iota_W,\mathcal J)$
(Definition~\ref{asm-theta-agree-planes}(5)), and let $(u_j,z_j)_{j\leq m}$ be
a polar system for $(\iota_X,\iota_W,\mathcal J)$
(Definition~\ref{asm-theta-agree-split-inputs-polar}). Write $\theta$ for theta
sums, and call linear forms admissible, as in
Theorem~\ref{asm-theta-agree-evaluation-heisenberg}.
\begin{enumerate}
\item For every admissible linear form $\lambda$ on the functions on
  $\mathbf A^{n_1\times d_1}$ and every $u\in\mathbb W(X,W)$, the form
  $\Phi\mapsto\lambda(e(\rho(u)\Phi))$ is admissible.
\item There is $C\neq0$ such that for every Schwartz-Bruhat function $\Phi$ on
  $\mathbf A^{n\times d}$ the sum
  $\sum_{y\in F^m}C\,\theta\bigl(e(\rho(u(y))\Phi)\bigr)$ converges absolutely
  to $\theta(\Phi)$.
\item If $e'$ is also an evaluation map for
  $(M,M_1,\iota_X,\iota_W,\mathcal J)$, there is $c\neq0$ with
  $e'(\Phi)=c\,e(\Phi)$ for every Schwartz-Bruhat function $\Phi$.
\item For every Schwartz-Bruhat function $\Phi$ there is
  $b\colon F^m\to[0,\infty)$ with $\sum_yb(y)<\infty$ and
  $|\theta(\rho_1(w)e(\rho(u(y))\Phi))|\leq b(y)$ for all $y\in F^m$ and
  $w\in\mathbb W(X_1,W_1)$.
\item Let $M_0$ be a regular model for lattice spaces $(X_0,W_0)$, $e_0$ an
  evaluation map for $(M_0,M,\kappa_X,\kappa_W,\mathcal J_0)$, and
  $\mathcal K\subset\mathbb W(X_0,W_0)$ a set each of whose elements lies in
  $\mathcal J_0$ or is $\kappa_Xw\kappa_W^{\mathsf T}$ with $w\in\mathcal J$.
  Then $e\circ e_0$ is an evaluation map for
  $(M_0,M_1,\kappa_X\iota_X,\kappa_W\iota_W,\mathcal K)$.
\end{enumerate}
\end{lemma}
\leanref{Development/WeilBound/ThetaLifts/SmallCones/ThetaUniq.lean}{PositivityTheory.ThKCT.EvAdm}
\leanref{Development/WeilBound/ThetaLifts/SmallCones/NormLeSum.lean}{PositivityTheory.ThKCT.evAdm_of}
\leanref{Development/WeilBound/ThetaLifts/SmallCones/ThetaUniq.lean}{PositivityTheory.ThKCT.ThetaExp}
\leanref{Development/WeilBound/ThetaLifts/SmallCones/SliceSummable.lean}{PositivityTheory.ThKCT.thetaExp_of}
\leanref{Development/WeilBound/ThetaLifts/SmallCones/ThetaUniq.lean}{PositivityTheory.ThKCT.EvUniq}
\leanref{Development/WeilBound/ThetaLifts/SmallCones/SliceSummable.lean}{PositivityTheory.ThKCT.evUniq_of}
\leanref{Development/WeilBound/ThetaLifts/SmallCones/ContinuousEvI.lean}{PositivityTheory.ThKCT.Ker.slice_unif}
\leanref{Development/WeilBound/ThetaLifts/SmallCones/KerPol.lean}{PositivityTheory.ThKCT.Ker.isEv_comp}
\leanref{Development/WeilBound/ThetaLifts/SmallCones/KerPol.lean}{PositivityTheory.ThKCT.Ker.comp_ne}
\begin{proof}
Put $I=\{1,\ldots,n\}\times\{1,\ldots,d\}$ and
$I_1=\{1,\ldots,n_1\}\times\{1,\ldots,d_1\}$, let $e^M$ and $e^{M_1}$ be the
isomorphisms of $M$ and $M_1$, and write $\mathcal S$ and $\mathcal S_1$ for
the Schwartz-Bruhat functions on $\mathbf A^I$ and $\mathbf A^{I_1}$. Then
$\rho(w)$ is the Schrödinger operator of
Theorem~\ref{asm-theta-agree-evaluation-heisenberg} at $e^M(w)$, so it maps
$\mathcal S$ into itself, and its formula and
$\langle w,w'\rangle=e^M_1(w)\cdot e^M_2(w')-e^M_1(w')\cdot e^M_2(w)$ give
\[
\rho(w)\rho(w')=\psi\bigl(\tfrac12\langle w,w'\rangle\bigr)\rho(w+w')
=\psi(\langle w,w'\rangle)\,\rho(w')\rho(w).\tag{R}
\]
By Lemma~\ref{asm-theta-agree-invariance-theta}(1), $|\psi|=1$; and $\psi$ is
trivial on $F$. For $y\in\mathbf A^m$ put $s(y)=\sum_j\sum_w|y_{j,w}|$, the
inner sum over the infinite places $w$ of $F$, as in
Lemma~\ref{asm-theta-agree-local-bump-slices}.

\emph{Step 1: the polar system.} Let $y,t\in\mathbf A^m$ and
$w\in\mathbb W(X_1,W_1)$, and write $\bar w=\iota_Xw\iota_W^{\mathsf T}$. As the
form is alternating and $\mathbf A$-bilinear, the conditions on a polar system
give $\langle u(y),z(t)\rangle=-y\cdot t$, $\langle u(y),u(t)\rangle=0$,
$\langle z(y),z(t)\rangle=0$ and
$\langle u(y),\bar w\rangle=\langle z(y),\bar w\rangle=0$. By (R), $\rho(z(t))$
is the product of the $\rho(t_jz_j)$; as each $t_jz_j$ lies in $\mathcal J$,
$e(\rho(z(t))\Phi)=e(\Phi)$ for $\Phi\in\mathcal S$. With (R) and the definition
of an evaluation map this gives, for $\Phi\in\mathcal S$,
\[
e\bigl(\rho(u(y))\rho(z(t))\Phi\bigr)=\psi(-y\cdot t)\,e\bigl(\rho(u(y))\Phi\bigr),\qquad
e\bigl(\rho(u(y))\rho(\bar w)\Phi\bigr)=\rho_1(w)\,e\bigl(\rho(u(y))\Phi\bigr).\tag{E}
\]
Let $x\in M_{n\times d}(E)$, written as $x=u(a)+z(b)+\bar x_1$ with
$a,b\in F^m$ and $x_1\in M_{n_1\times d_1}(E)$. By (R),
$\rho(u(a))\rho(z(b))\rho(\bar x_1)=\psi(-\tfrac12a\cdot b)\rho(x)=\rho(x)$. By
Lemma~\ref{asm-theta-agree-evaluation-rational}, every element of
$F^I\times F^I$ is $e^M(x)$ for such an $x$.

\emph{Step 2: a Fourier argument.} Let $a\colon F^m\to\mathbf C$ satisfy
$\sum_y|a(y)|<\infty$ and $\sum_ya(y)\psi(y\cdot t)=0$ for every
$t\in\mathbf A^m$. Then $a=0$. Indeed, fix $y_0$ and apply
Lemma~\ref{asm-theta-agree-split-inputs-characters} to $A=\mathbf A^m$,
$e_y(t)=\psi(y\cdot t)$, $c_y=a(y+y_0)$ and $i_0=0$: the sum
$\sum_yc_ye_y(t)$ is $\psi(-y_0\cdot t)\sum_ya(y)\psi(y\cdot t)=0$; and if
$y_i\neq0$ and $\psi(r)\neq1$, then the vector $t$ with coordinate $y_i^{-1}r$
at $i$ and $0$ elsewhere has $e_y(t)=\psi(r)\neq1$. So $a(y_0)=c_0=0$.
Consequently, if linear forms $a_y$ on $\mathcal S$, $y\in F^m$, satisfy
$a_y(\rho(z(t))\Phi)=\psi(-y\cdot t)\,a_y(\Phi)$ for $t\in\mathbf A^m$, and
$\sum_ya_y(\Phi)$ converges absolutely to $0$ for every $\Phi\in\mathcal S$,
then every $a_y$ vanishes on $\mathcal S$: apply the above to
$y\mapsto a_y(\rho(z(-t))\Phi)=\psi(y\cdot t)\,a_y(\Phi)$.

\emph{Step 3: a weight.} Let $C\subset\mathbf A_f^m$ be compact. Then there is
$N$ with $\sum(1+|\xi_\infty|)^{-N}<\infty$, the sum over the $\xi\in F^m$ whose
finite component lies in $C$, where $|\cdot|$ is the largest absolute value of
the components at the infinite places; as $1+|\xi_\infty|\leq1+s(\xi)$, the
same holds for $\sum(1+s(\xi))^{-N}$. Indeed, there is an integer $D\geq1$ with
$DC\subset\prod_v\mathcal O_v^m$; such $\xi$ have $D\xi_j\in\mathcal O_F$, so
they lie in the lattice $(D^{-1}\mathcal O_F)^m$ of $F_\infty^m$, and the sum of
$(1+|\lambda|)^{-N}$ over a lattice converges when $N$ exceeds its rank.

(1) Let $\Phi=f\otimes g$ with $f\in\mathcal S(F_\infty^I)$ and $g$
Schwartz-Bruhat on $\mathbf A_f^I$. By
Lemma~\ref{asm-theta-agree-local-bump-slices-adelic}(2) and (3),
$\rho(u)\Phi=S(w)f\otimes g'$, where $w$ is the archimedean component of
$e^M(u)$ read in $V\times V^*$ as there and $g'$ is Schwartz-Bruhat. By Step 1 of
the proof of Lemma~\ref{asm-theta-agree-local-bump-slices},
$e(S(w)f\otimes g')=e_\infty(S(w)f)\otimes\epsilon(g')$ with $\epsilon(g')$
Schwartz-Bruhat. So $f\mapsto\lambda(e(\rho(u)(f\otimes g)))$ is the composite
of the continuous maps $S(w)$ and $e_\infty$ with the form
$f_1\mapsto\lambda(f_1\otimes\epsilon(g'))$, which is continuous as $\lambda$
is admissible.

(2) Let $\lambda_1$ be an admissible linear form that agrees with $\theta$ on
$\mathcal S_1$ (Theorem~\ref{asm-theta-agree-evaluation-heisenberg}(2)). Let
$\Phi=f\otimes\phi$ with $f\in\mathcal S(F_\infty^I)$ and $\phi$
Schwartz-Bruhat on $\mathbf A_f^I$. Lemma~\ref{asm-theta-agree-local-bump-slices},
whose hypotheses on the $u_j$ and $z_j$ are condition (1) of a polar system,
gives Schwartz-Bruhat functions $\phi_{1,i}$, a compact set $C$ and maps
$A_i(y)$, $i\leq r$. Let $N$ be as in Step 3 for $C$, and
$T_i(f_1)=\lambda_1(f_1\otimes\phi_{1,i})$, a continuous linear form. By
Lemma~\ref{asm-theta-agree-local-bump-slices}(4) for the seminorm $|T_i|$ there
is a continuous seminorm $p_i$ with $|T_i(A_i(y)f)|(1+s(y))^N\leq p_i(f)$. By
Lemma~\ref{asm-theta-agree-local-bump-slices}(2) and (3), for $y\in F^m$,
\[
|\theta(e(\rho(u(y))\Phi))|=\Bigl|\sum_iT_i(A_i(y)f)\Bigr|\leq P(f)\,\varpi(y),
\]
with $P=\sum_ip_i$, $\varpi(y)=(1+s(y))^{-N}$ if the finite component of $y$
lies in $C$ and $\varpi(y)=0$ otherwise; and $\sum_{y\in F^m}\varpi(y)<\infty$.
By linearity, $L(\Phi)=\sum_{y\in F^m}\theta(e(\rho(u(y))\Phi))$ converges
absolutely for every $\Phi\in\mathcal S$. It is linear on $\mathcal S$; extend
it to a linear form on all functions on $\mathbf A^I$. It is admissible, as
$|L(f\otimes\phi)|\leq P(f)\sum_y\varpi(y)$ with $P$ continuous.

The form $L$ is invariant under the rational Schrödinger operators. Let
$x\in M_{n\times d}(E)$, and write $\rho(x)=\rho(u(a))\rho(z(b))\rho(\bar x_1)$
as in Step 1; it suffices to treat the three factors. Let $\Phi\in\mathcal S$.
As $\rho(u(y))\rho(u(a))=\rho(u(y+a))$ by (R) and Step 1, the terms of
$L(\rho(u(a))\Phi)$ are those of $L(\Phi)$, reindexed by $y\mapsto y+a$. By (E),
the terms of $L(\rho(z(b))\Phi)$ are those of $L(\Phi)$ times
$\psi(-y\cdot b)=1$, and the terms of $L(\rho(\bar x_1)\Phi)$ are
$\theta(\rho_1(x_1)e(\rho(u(y))\Phi))=\theta(e(\rho(u(y))\Phi))$, as $M_1$ is
regular, so that $e^{M_1}(x_1)$ is rational, and by
Theorem~\ref{asm-theta-agree-evaluation-heisenberg}(1). As every rational
element of $\mathbf A^I\times\mathbf A^I$ is some $e^M(x)$ (Step 1),
Theorem~\ref{asm-theta-agree-evaluation-heisenberg}(4) gives $c_0$ with
$L=c_0\,\theta$ on $\mathcal S$.

Suppose $c_0=0$. By (E), the forms $a_y(\Phi)=\theta(e(\rho(u(y))\Phi))$ satisfy
the hypotheses of Step 2, so $\theta(e(\Phi))=a_0(\Phi)=0$ for every
$\Phi\in\mathcal S$. For $w\in\mathbb W(X_1,W_1)$ and $\Phi\in\mathcal S$ then
$\theta(\rho_1(w)e(\Phi))=\theta(e(\rho(\bar w)\Phi))=0$; as $e^{M_1}$ is onto
$\mathbf A^{I_1}\times\mathbf A^{I_1}$,
Theorem~\ref{asm-theta-agree-evaluation-heisenberg}(3) gives $e(\Phi)=0$ for
every $\Phi\in\mathcal S$. This contradicts fact (a) in the proof of
Lemma~\ref{asm-theta-agree-evaluation-unique}. So $c_0\neq0$, and $C=c_0^{-1}$
works.

(3) Let $C$ and $C'$ be the constants of (2) for $e$ and $e'$, and put
$a_y(\Phi)=C\theta(e(\rho(u(y))\Phi))-C'\theta(e'(\rho(u(y))\Phi))$ for
$\Phi\in\mathcal S$. By (2), $\sum_ya_y(\Phi)$ converges absolutely to
$\theta(\Phi)-\theta(\Phi)=0$, and by (E) for $e$ and $e'$ the hypotheses of
Step 2 hold. So $\theta(Ce(\Phi)-C'e'(\Phi))=a_0(\Phi)=0$ for every
$\Phi\in\mathcal S$. Applied to $\rho(\bar w)\Phi$, this gives
$\theta(\rho_1(w)(Ce(\Phi)-C'e'(\Phi)))=0$ for every $w$, and as in (2),
$Ce(\Phi)=C'e'(\Phi)$. Put $c=C/C'$.

(4) \emph{A uniform bound for theta sums.} Let $g$ be Schwartz-Bruhat on
$\mathbf A_f^{I_1}$. There is a continuous seminorm $q$ on
$\mathcal S(F_\infty^{I_1})$ with $|\theta(\rho^\circ(w)(f_1\otimes g))|\leq q(f_1)$
for every $f_1$ and every $w\in\mathbf A^{I_1}\times\mathbf A^{I_1}$, where
$\rho^\circ$ are the operators of
Theorem~\ref{asm-theta-agree-evaluation-heisenberg}. Indeed, by
Theorem~\ref{asm-theta-agree-invariance-tate}(1) there is a compact set
$K\subset\mathbf A$ with $\mathbf A=F+K$. Write the first coordinate of $w$ as
$w_1=\alpha+k$ with $\alpha\in F^{I_1}$ and $k\in K^{I_1}$. As $|\psi|=1$,
reindexing by $\xi\mapsto\xi-\alpha$ gives
\[
|\theta(\rho^\circ(w)(f_1\otimes g))|\leq\sum_{\xi\in F^{I_1}}|f_1(\xi_\infty+k_\infty)|\,|g(\xi_f+k_f)|.
\]
The term vanishes unless $\xi_f$ lies in the compact set
$C_1=\operatorname{supp}g-K_f^{I_1}$, where $K_f$ is the image of $K$ in
$\mathbf A_f$. Let $N$ be as in Step 3 for $C_1$, let $G$ bound $|g|$ and $R$
bound the archimedean components of $K^{I_1}$, and let $q_N$ be a continuous
seminorm with $(1+|x|)^N|f_1(x)|\leq q_N(f_1)$. As
$1+|\xi_\infty|\leq(1+R)(1+|\xi_\infty+k_\infty|)$, the sum is at most
$G(1+R)^Nq_N(f_1)\sum(1+|\xi_\infty|)^{-N}$, the sum over the $\xi$ with
$\xi_f\in C_1$, which is finite.

\emph{The bound.} By linearity it suffices to treat $\Phi=f\otimes\phi$. With
$\phi_{1,i}$, $C$, $A_i(y)$, $N$ and $\varpi$ as in (2), let $q_i$ be the
seminorm just found for $g=\phi_{1,i}$, and $p_i$ that of
Lemma~\ref{asm-theta-agree-local-bump-slices}(4) for $q_i$ and $N$. For $y\in F^m$
and $w\in\mathbb W(X_1,W_1)$,
$\rho_1(w)e(\rho(u(y))\Phi)=\sum_i\rho_1(w)(A_i(y)f\otimes\phi_{1,i})$, so
\[
|\theta(\rho_1(w)e(\rho(u(y))\Phi))|\leq\sum_iq_i(A_i(y)f)\leq\Bigl(\sum_ip_i(f)\Bigr)\varpi(y),
\]
and $b=(\sum_ip_i(f))\,\varpi$ works.

(5) \emph{The composite is nonzero.} Let $\rho_0$ be the Schrödinger operators
of $M_0$, and $\mathcal S_0$ the Schwartz-Bruhat functions of its index set.
There is $\Phi\in\mathcal S_0$ with $e(e_0(\Phi))\neq0$. Otherwise, for every
$\Psi\in\mathcal S_0$, (2) applied to $e_0(\Psi)$ gives
$\theta(e_0(\Psi))=\sum_yC\theta(e(\rho(u(y))e_0(\Psi)))$, and
$e(\rho(u(y))e_0(\Psi))=e(e_0(\rho_0(\kappa_Xu(y)\kappa_W^{\mathsf T})\Psi))=0$;
so $\theta(e_0(\Psi))=0$. Then
$\theta(\rho(w)e_0(\Phi))=\theta(e_0(\rho_0(\kappa_Xw\kappa_W^{\mathsf T})\Phi))=0$
for every $w\in\mathbb W(X,W)$ and $\Phi\in\mathcal S_0$, and
Theorem~\ref{asm-theta-agree-evaluation-heisenberg}(3) gives $e_0=0$ on
$\mathcal S_0$, contradicting fact (a) in the proof of
Lemma~\ref{asm-theta-agree-evaluation-unique}.

\emph{The conditions.} The composite maps Schwartz-Bruhat functions to
Schwartz-Bruhat functions. For $w\in\mathbb W(X_1,W_1)$ and
$\Phi\in\mathcal S_0$,
$ee_0(\rho_0(\kappa_X\iota_Xw\iota_W^{\mathsf T}\kappa_W^{\mathsf T})\Phi)=e(\rho(\iota_Xw\iota_W^{\mathsf T})e_0\Phi)=\rho_1(w)ee_0(\Phi)$.
If $k\in\mathcal K$ lies in $\mathcal J_0$, then $e_0(\rho_0(k)\Phi)=e_0(\Phi)$;
if $k=\kappa_Xw\kappa_W^{\mathsf T}$ with $w\in\mathcal J$, then
$ee_0(\rho_0(k)\Phi)=e(\rho(w)e_0\Phi)=ee_0(\Phi)$. Let $e_\infty,e_v$ and
$e_{0,\infty},e_{0,v}$ be components of $e$ and $e_0$. The maps $e_ve_{0,v}$
carry the vacua of $M_0$ to those of $M_1$ and the Schwartz-Bruhat functions
onto the Schwartz-Bruhat functions, and
$ee_0(f\otimes\bigotimes_v\phi_v)=e_\infty e_{0,\infty}(f)\otimes\bigotimes_ve_ve_{0,v}(\phi_v)$
when the $\phi_v$ are Schwartz-Bruhat and almost all equal to the vacua, since
then so are the $e_{0,v}(\phi_v)$. Finally $e_\infty e_{0,\infty}\neq0$:
otherwise $ee_0$ vanishes on these products, which span $\mathcal S_0$ as $M_0$
is regular (Step 1 of the proof of Lemma~\ref{asm-theta-agree-local-bump-slices}),
contradicting the first paragraph.
\end{proof}
% Lean: (1) is evAdm_of (NormLeSum.lean), with ev_pure and schr_pure. (2) is thetaExp_of
% (SliceSummable.lean): Step 3 is summable_wt (ExpSum.lean: exists_nat_mul_integral_of_isCompact,
% denomLattice, summable_lat_wt); the bound and the admissible form L are slice_bound and
% slice_adm (LinearMap.exists_extend); the invariance is hu, hz (ev_rho_uz) and hm (ev_rho_um,
% theta_schr_rat), with rat_surj and rho_rat (ExpAlg.lean); L = c theta is thetaUniq_of; c /= 0 is
% hc0 with term_zero, fourier_uniq (Step 2, through charSeries_coeff_eq_zero), thetaInj_of and
% ev_ne_zero. (3) is evUniq_of. (4) is Ker.slice_unif with theta_unif (ContinuousEvI.lean:
% exists_compact_adele, summable_wt, sz_le). (5) is Ker.isEv_comp with Ker.comp_ne (KerPol.lean)
% and Span.exists_sum_finprod for the spanning products; the Lean takes the nonvanishing of the
% composite as a hypothesis of isEv_comp and proves it in comp_ne from PolT.exp.

\begin{lemma}\label{asm-theta-agree-split-inputs-systems}
Let $B$ be a theta base for a CM field $E$, let $X$ be a Hermitian and $W$ a
skew-Hermitian lattice space, of ranks $n$ and $d$ and of some signatures
(Definition~\ref{asm-theta-agree-doubling}(1)), and write
$h_X(v,v')=v^{\mathsf T}J_Xc(v')$ and $h_W(v,v')=v^{\mathsf T}J_Wc(v')$. For an
adapted plane $H=(x,y,\iota,\mathfrak b)$ in $X$ over a lattice space $X_1$ of
rank $n_1$ put $\mathcal J_x(H)=\{xz^{\mathsf T}\}$, with $z$ running over the
vectors over $\mathbf A_E$ of length $d$; for adapted planes $H$ and
$H'=(x',y',\iota',\mathfrak b')$ in $W$ over a lattice space $W_1$ of rank
$d_1$ put
\[
\mathcal J_\alpha^+(H,H')=\{vx'^{\mathsf T}+x(\iota'z)^{\mathsf T}\},\qquad
\mathcal J_\beta^+(H,H')=\{xv^{\mathsf T}+(\iota z)x'^{\mathsf T}\},
\]
with $v$ and $z$ running over vectors over $\mathbf A_E$ of the appropriate
lengths (Definition~\ref{asm-theta-agree-planes}). Polar systems are taken for
the forms of the pairs of lattice spaces indicated
(Definition~\ref{asm-theta-agree-split-inputs-polar}).
\begin{enumerate}
\item For every $F$-basis $(b_j)$ of $E^n$ there are $z_j$ such that
  $(b_jy'^{\mathsf T},z_j)$ is a polar system for $(1,\iota',\mathcal J_0(H'))$,
  for $(X,W)$ and $(X,W_1)$.
\item For every $F$-basis $(b_j)$ of $E^d$ there are $z_j$ such that
  $(yb_j^{\mathsf T},z_j)$ is a polar system for $(\iota,1,\mathcal J_x(H))$, for
  $(X,W)$ and $(X_1,W)$.
\item For every $F$-basis $(b_j,b'_j)$ of $E^n\times E^{d_1}$ there are $z_j$
  such that $(b_jy'^{\mathsf T}+y(\iota'b'_j)^{\mathsf T},z_j)$ is a polar
  system for $(\iota,\iota',\mathcal J_\alpha^+(H,H'))$, for $(X,W)$ and
  $(X_1,W_1)$.
\item For every $F$-basis $(b_j,b'_j)$ of $E^d\times E^{n_1}$ there are $z_j$
  such that $(yb_j^{\mathsf T}+(\iota b'_j)y'^{\mathsf T},z_j)$ is a polar
  system for $(\iota,\iota',\mathcal J_\beta^+(H,H'))$, for $(X,W)$ and
  $(X_1,W_1)$.
\item For models $M$ and $M_1$ for $(X,W)$ and $(X_1,W_1)$, every evaluation map
  for $(M,M_1,\iota,\iota',\mathcal J_\alpha(H,H'))$ is one for
  $(M,M_1,\iota,\iota',\mathcal J_\alpha^+(H,H'))$, and every evaluation map for
  $(M,M_1,\iota,\iota',\mathcal J_\beta(H,H'))$ is one for
  $(M,M_1,\iota,\iota',\mathcal J_\beta^+(H,H'))$.
\end{enumerate}
\end{lemma}
\leanref{Development/WeilBound/ThetaLifts/SmallCones/KerGeo.lean}{PositivityTheory.ThKCT.Ker.polar_of}
\leanref{Development/WeilBound/ThetaLifts/SmallCones/KerPol.lean}{PositivityTheory.ThKCT.Ker.pol0}
\leanref{Development/WeilBound/ThetaLifts/SmallCones/KerPol.lean}{PositivityTheory.ThKCT.Ker.polX}
\leanref{Development/WeilBound/ThetaLifts/SmallCones/KerPol.lean}{PositivityTheory.ThKCT.Ker.polA}
\leanref{Development/WeilBound/ThetaLifts/SmallCones/KerPol.lean}{PositivityTheory.ThKCT.Ker.polB}
\leanref{Development/WeilBound/ThetaLifts/SmallCones/KerPol.lean}{PositivityTheory.ThKCT.Ker.isEv_sumA}
\leanref{Development/WeilBound/ThetaLifts/SmallCones/KerPol.lean}{PositivityTheory.ThKCT.Ker.isEv_sumB}
\begin{proof}
Write $h$ for the form of Definition~\ref{asm-theta-agree-models}(3) on
$\mathbb W(X,W)$, so that $\langle p,q\rangle=-\operatorname{Tr}h(p,q)$,
$h(\lambda p,q)=\lambda h(p,q)$, $h(p,\lambda q)=c(\lambda)h(p,q)$ and
$h(ab^{\mathsf T},a'b'^{\mathsf T})=h_X(a,a')h_W(b,b')$; for rational $p,q$,
$\langle p,q\rangle=-\operatorname{Tr}_{E/F}h(p,q)$. The matrices of $X$, $W$,
$X_1$ and $W_1$ are invertible (Lemmas~\ref{asm-theta-agree-split-low}
and~\ref{asm-theta-agree-convergence-signature}), $h_X(v',v)=c(h_X(v,v'))$ and
$h_W(v',v)=-c(h_W(v,v'))$. By the definition of an adapted plane and these
rules, $h_X(x,x)=h_X(y,y)=0$, $h_X(x,y)=h_X(y,x)=1$, $h_X$ of $x$ or $y$ with a
vector $\iota z$ vanishes in either order, and
$h_X(\iota z,\iota z')=h_{X_1}(z,z')$; the same holds for $H'$, except that
$h_W(y',x')=-1$. As $x$, $y$ and the columns of $\iota$ form a basis of $E^n$
(proof of Lemma~\ref{asm-theta-agree-split-low}), every
$p\in M_{n\times d}(E)$ is $xs^{\mathsf T}+\iota q+yt^{\mathsf T}$ with
$s,t\in E^d$ and $q$ rational, applying this to the columns; likewise $p$ is
$sx'^{\mathsf T}+q\iota'^{\mathsf T}+ty'^{\mathsf T}$, applying the same for
$H'$ to the rows.

\emph{A criterion.} Let $X'$ and $W'$ be lattice spaces of ranks $n'$ and $d'$,
let $\iota_X\in M_{n\times n'}(E)$ and $\iota_W\in M_{d\times d'}(E)$ satisfy
$\iota_X^{\mathsf T}J_Xc(\iota_X)=J_{X'}$ and
$\iota_W^{\mathsf T}J_Wc(\iota_W)=J_{W'}$, and let
$\mathcal J\subset\mathbb W(X,W)$.
Let $T$ be a finite-dimensional $E$-vector space and $U,Z\colon T\to M_{n\times d}(E)$
$E$-linear maps such that: $aZ(t)\in\mathcal J$ for $a\in\mathbf A$; $t=0$
whenever $h(Z(t'),U(t))=0$ for every $t'$; $h(U(t),U(t'))=h(Z(t),Z(t'))=0$;
$h(U(t),\iota_Xw\iota_W^{\mathsf T})=h(Z(t),\iota_Xw\iota_W^{\mathsf T})=0$ for
$w\in\mathbb W(X',W')$; and every $p\in M_{n\times d}(E)$ is
$U(t)+Z(t')+\iota_Xp_1\iota_W^{\mathsf T}$ with $p_1$ rational. Then for every
$F$-basis $(b_j)$ of $T$ there are $z_j$ such that $(U(b_j),z_j)$ is a polar
system for $(\iota_X,\iota_W,\mathcal J)$. Indeed,
$\beta(t',t)=-\operatorname{Tr}_{E/F}h(Z(t'),U(t))$ is an $F$-bilinear form on
$T$. It is nondegenerate: if $\beta(t',t)=0$ for every $t'$, then, replacing
$t'$ by $\lambda t'$, $\operatorname{Tr}_{E/F}(\lambda h(Z(t'),U(t)))=0$ for every
$\lambda\in E$, so $h(Z(t'),U(t))=0$, as the trace form of $E/F$ is
nondegenerate, and $t=0$. Let $(b'_j)$ be the basis with
$\beta(b'_i,b_j)=\delta_{ij}$ and put $z_j=Z(b'_j)$. Then
$\langle z_i,U(b_j)\rangle=\delta_{ij}$, and the other conditions of a polar
system follow from the hypotheses: for
$\langle\iota_Xw\iota_W^{\mathsf T},\iota_Xw'\iota_W^{\mathsf T}\rangle=\langle w,w'\rangle'$,
write $h(p,q)=\operatorname{tr}(p^{\mathsf T}J_Xc(q)J_W^{\mathsf T})$ and use
the relations between the matrices; for condition (3), write
$p=U(t)+Z(t')+\iota_Xp_1\iota_W^{\mathsf T}$ and expand $t$ in the basis
$(b_j)$ and $t'$ in $(b'_j)$.

(1) Take $T=E^n$, $U(t)=ty'^{\mathsf T}$, $Z(t)=tx'^{\mathsf T}$, $\iota_X=1$,
$\iota_W=\iota'$ and $\mathcal J=\mathcal J_0(H')$. Then
$aZ(t)=(at)x'^{\mathsf T}\in\mathcal J_0(H')$; $h(Z(t'),U(t))=h_X(t',t)h_W(x',y')=h_X(t',t)$,
which vanishes for every $t'$ only if $t=0$, as $J_X$ is invertible;
$h(U(t),U(t'))=h_X(t,t')h_W(y',y')=0$, and likewise for $Z$. The matrix
$w\iota'^{\mathsf T}$ is a sum of matrices $a'(\iota'z)^{\mathsf T}$, and
$h(ab^{\mathsf T},a'(\iota'z)^{\mathsf T})$ contains the factor
$h_W(b,\iota'z)$, which is $0$ for $b=y'$ and $b=x'$. Finally
$p=sx'^{\mathsf T}+q\iota'^{\mathsf T}+ty'^{\mathsf T}=U(t)+Z(s)+1\cdot q\cdot\iota'^{\mathsf T}$.

(2) is (1) with the roles of $X$ and $W$ exchanged: $T=E^d$,
$U(t)=yt^{\mathsf T}$, $Z(t)=xt^{\mathsf T}$,
$h(Z(t'),U(t))=h_X(x,y)h_W(t',t)=h_W(t',t)$, and
$p=xs^{\mathsf T}+\iota q+yt^{\mathsf T}=U(t)+Z(s)+\iota\cdot q\cdot1$.

(3) Take $T=E^n\times E^{d_1}$, $U(t,t')=ty'^{\mathsf T}+y(\iota't')^{\mathsf T}$
and $Z(t,t')=tx'^{\mathsf T}+x(\iota't')^{\mathsf T}$. Clearly
$aZ(t,t')\in\mathcal J_\alpha^+(H,H')$. Expanding,
$h(Z(s,s'),U(t,t'))=h_X(s,t)+h_{W_1}(s',t')$, the two cross terms containing
$h_W(x',\iota't')$ and $h_W(\iota's',y')$. If this vanishes for all $(s,s')$,
then $t=0$ and $t'=0$, as the matrices of $X$ and $W_1$ are invertible. Each of
the four terms of $h(U(t,t'),U(s,s'))$ contains one of $h_W(y',y')$,
$h_W(y',\iota's')$, $h_W(\iota't',y')$ and $h_X(y,y)$, so it is $0$; likewise
for $Z$, with $x$ and $x'$. Each of the two terms of $U$ and of $Z$ is
orthogonal to every $\iota w\iota'^{\mathsf T}$, as $h_W(y',\iota'z)=h_W(x',\iota'z)=0$
and $h_X(y,\iota z)=h_X(x,\iota z)=0$. Finally, write
$p=sx'^{\mathsf T}+q\iota'^{\mathsf T}+ty'^{\mathsf T}$ and
$q=xs'^{\mathsf T}+\iota q_1+yt'^{\mathsf T}$ with $s',t'\in E^{d_1}$; then
$q\iota'^{\mathsf T}=x(\iota's')^{\mathsf T}+\iota q_1\iota'^{\mathsf T}+y(\iota't')^{\mathsf T}$
and $p=U(t,t')+Z(s,s')+\iota q_1\iota'^{\mathsf T}$.

(4) is (3) with the roles of $X$ and $W$ exchanged: $T=E^d\times E^{n_1}$,
$U(t,t')=yt^{\mathsf T}+(\iota t')y'^{\mathsf T}$,
$Z(t,t')=xt^{\mathsf T}+(\iota t')x'^{\mathsf T}$,
$h(Z(s,s'),U(t,t'))=h_W(s,t)+h_{X_1}(s',t')$, and with
$p=xs^{\mathsf T}+\iota q+yt^{\mathsf T}$ and
$q=s'x'^{\mathsf T}+q_1\iota'^{\mathsf T}+t'y'^{\mathsf T}$,
$p=U(t,t')+Z(s,s')+\iota q_1\iota'^{\mathsf T}$.

(5) Let $e$ be an evaluation map for $\mathcal J_\alpha(H,H')$, and let
$a=vx'^{\mathsf T}+x(\iota'z)^{\mathsf T}$. Then
$\langle vx'^{\mathsf T},x(\iota'z)^{\mathsf T}\rangle=-\operatorname{Tr}h_X(v,x)h_W(x',\iota'z)=0$,
so by the formula for the Schrödinger operators $\rho(a)=\rho(vx'^{\mathsf T})\rho(x(\iota'z)^{\mathsf T})$,
and $e$ is invariant under both factors on the Schwartz-Bruhat functions. For
$\mathcal J_\beta(H,H')$, use
$\langle xv^{\mathsf T},(\iota z)x'^{\mathsf T}\rangle=-\operatorname{Tr}h_X(x,\iota z)h_W(v,x')=0$.
\end{proof}
% Lean: the criterion is Ker.polar_of (KerGeo.lean), with the dual basis of the form beta
% (LinearMap.BilinForm.dualBasis, traceForm_nondegenerate), sy_inc, hT_mm and incM_gram; the
% plane identities are PG (pg_of, pg_herm, pg_skew), with span_l and span_r. (1)-(4) are pol0,
% polX, polA and polB (KerPol.lean), the sets being iso0, isoX, sumA and sumB; (5) is isEv_sumA
% and isEv_sumB. Through PolT.exp and PolT.uniq, Lemma asm-theta-agree-split-inputs-expansion(2)
% and (3) are applied to these systems, with the sum over F^m read as the sum over T by the basis.

\begin{lemma}\label{asm-theta-agree-split-inputs-kfinite}
Let $(X,W,M,\chi_X,\chi_W,\omega)$ be a Weil pair for a theta base $B$ of a CM
field $E$ (Definition~\ref{asm-theta-agree-triple}(1)), and let $K$ be the
lattice stabilizer of $X$ (Definition~\ref{asm-theta-agree-planes}(6)). Then
every $\Phi\in\operatorname{PD}(M)$ is $K$-finite for $\omega$
(Definition~\ref{asm-theta-agree-split-notions}(3)).
\end{lemma}
\leanref{Development/WeilBound/ThetaLifts/SmallCones/KFinDataOf.lean}{PositivityTheory.ThKCT.kFinData_of}
\leanref{Development/WeilBound/ThetaLifts/SmallCones/KFinDataOf.lean}{PositivityTheory.ThKCT.So.orbit_fin}
\leanref{Development/WeilBound/ThetaLifts/SmallCones/KFinAn.lean}{PositivityTheory.ThKCT.So.level_fin}
\leanref{Development/WeilBound/ThetaLifts/SmallCones/KFinAn.lean}{PositivityTheory.ThKCT.So.vacs_fin}
\leanref{Development/WeilBound/ThetaLifts/SmallCones/KFinArch.lean}{PositivityTheory.ThKCT.So.polM_stab}
\begin{proof}
Let $X$ and $W$ have ranks $n$ and $d$, put $I=\{1,\ldots,n\}\times\{1,\ldots,d\}$
and $V=F_\infty^I$, write $(\xi,y)$ for the point of $\mathbf A^I$ with
component $\xi\in V$ at $\infty$ and finite component $y\in\mathbf A_f^I$, and
for a function $\Psi$ on $\mathbf A^I$ put $\Psi_y(\xi)=\Psi(\xi,y)$. Use $T$
and $D_{a,b}$ of Lemma~\ref{asm-theta-agree-density-derivative}, and write
$D_p=D_{a,b}$ for $p=(a,b)\in V\times V$; it is $\mathbf R$-linear in $p$. For
$u\in\mathbb W(X,W)$ at $\infty$ let $p(u)\in V\times V$ be the components at
$\infty$ of $e(u)$, where $e$ is the isomorphism of $M$; as in the proof of
Lemma~\ref{asm-theta-agree-density-gaussian}, $u\mapsto p(u)$ is a bijection
from the elements at $\infty$ onto $V\times V$, and $u\mapsto(\pi_w(u))_w$ one
onto $\prod_wM_{n\times d}(\mathbf C)$.

\emph{Step 1: Gaussians.} By Lemma~\ref{asm-theta-agree-density-derivative}, as
in Step 6 of the proof of Lemma~\ref{asm-theta-agree-evaluation-pure}, a
Gaussian with signs $\varepsilon=(\varepsilon_w)$ is an $f\in\mathcal S(V)$ with
$D_{p(u')}f=i\,D_{p(u)}f$ for all $u,u'$ at $\infty$ with
$\pi_w(u')=\varepsilon_wi\,\theta_{X,w}\pi_w(u)\theta_{W,w}^{\mathsf T}$ for every
$w$; call such $u,u'$ related. Let $G_\varepsilon$ be the space of these $f$. If
$f\in G_\varepsilon$ and $f(0)=0$, then $f=0$; so $\dim G_\varepsilon\leq1$.
Indeed, let $\Gamma\subset(V\times V)\times(V\times V)$ be the real span of the
pairs $(p(u),p(u'))$ with $u,u'$ related. As $D_p$ is linear in $p$,
$D_{q'}f=i\,D_qf$ for every $(q,q')\in\Gamma$, and every $q\in V\times V$ is the
first component of an element of $\Gamma$. Let $\pi\colon\Gamma\to V\times V$
send $((a,b),(a',b'))$ to $(a,a')$. If $\pi$ is injective, then
$\dim\Gamma\leq\dim(V\times V)\leq\dim\Gamma$, so $\pi$ is bijective; for
$x\in V$ take $((0,b),(x,b'))\in\Gamma$. Then $D_{x,b'}f=i\,D_{0,b}f$ at the
points $sx$, $s\in\mathbf R$, gives
$\frac{d}{ds}f(sx)=s\,c\,f(sx)$ with $c=2\pi i(T(x,b')-iT(x,b))$, so
$f(sx)=f(0)e^{cs^2/2}$ and $f(x)=0$. If $\pi$ is not injective, there is
$((0,b),(0,b'))\neq0$ in $\Gamma$, and
$f(\xi)\bigl(T(\xi,b')-iT(\xi,b)\bigr)=0$ for every $\xi$. If $f\neq0$, the real
linear forms $T(\cdot,b)$ and $T(\cdot,b')$ vanish on the nonempty open set where
$f\neq0$, hence everywhere, and $b=b'=0$, as $T$ is nondegenerate; a
contradiction.

\emph{Step 2: a filtration.} Fix a real basis $(q_i)$ of $V\times V$. Put
$P_0=\sum_\varepsilon G_\varepsilon$, over the finitely many sign vectors, and
$P_{m+1}=P_m+\sum_iD_{q_i}(P_m)$. By Step 1 and induction every $P_m$ is
finite-dimensional, $P_m\subset P_{m+1}$, and $D_q(P_m)\subset P_{m+1}$ for every
$q$, as $D_q$ is a combination of the $D_{q_i}$. By
Lemma~\ref{asm-theta-agree-density-derivative}, $\operatorname{Pol}(M)$ is the
smallest subspace of $\mathcal S(V)$ that contains the Gaussians and is stable
under every $D_{p(u)}$. The union of the $P_m$ is such a subspace, so every
$f\in\operatorname{Pol}(M)$ lies in some $P_m$.

\emph{Step 3: slices.} Let $k\in K$, $g$ a Schwartz-Bruhat function on
$\mathbf A_f^I$ and $y\in\mathbf A_f^I$. As $\omega(k,1)(f\otimes g)$ is
Schwartz-Bruhat, its slice at $y$ lies in $\mathcal S(V)$; let $A$ be the
linear map $f\mapsto(\omega(k,1)(f\otimes g))_y$ of $\mathcal S(V)$. Then
$A(P_m)\subset P_m$ for every $m$.

(i) $AD_{p(u)}=D_{p(ku)}A$ for every $u$ at $\infty$. Suppose first that the
first component $a$ of $p(u)=(a,b)$ is nonzero. Fix $\xi\in V$. Evaluation at
$(\xi,y)$ is admissible, and $\omega(k,1)$ preserves admissible forms
(Definition~\ref{asm-theta-agree-models}(6)), so $f\mapsto(Af)(\xi)$ is a
continuous linear form. As in Step 5 of the proof of
Lemma~\ref{asm-theta-agree-evaluation-pure}, $\rho(su)(f\otimes g)=G_s\otimes g$
with $G_s$ the family of Lemma~\ref{asm-theta-agree-evaluation-derivative} for
$f$, $a$ and $T(\cdot,b)$; and as $ku$ is at $\infty$, the slice at $y$ of
$\rho(s\,ku)\Psi$ is the family $G'_s$ of that lemma for $\Psi_y$ and $p(ku)$.
By Definition~\ref{asm-theta-agree-models}(6)(b),
$\omega(k,1)\rho(su)=\rho(s\,ku)\omega(k,1)$ on Schwartz-Bruhat functions, so
with $\Psi=\omega(k,1)(f\otimes g)$, $A(G_s)(\xi)=G'_s(\xi)$ for every $s$.
Differentiating at $s=0$, by Lemma~\ref{asm-theta-agree-evaluation-derivative}
on the left, gives $(AD_{p(u)}f)(\xi)=(D_{p(ku)}Af)(\xi)$. If $a=0$, choose $v$
at $\infty$ such that $p(v)$ has nonzero first component; the claim for $u+v$
and for $v$, and the linearity of $D_p$ in $p$, give it for $u$.

(ii) $A(G_\varepsilon)\subset G_\varepsilon$. Let $f\in G_\varepsilon$ and $u,u'$
related. As $k\in K$, $\pi_w(k)$ commutes with $\theta_{X,w}$, so $k^{-1}u$ and
$k^{-1}u'$ are related. By (i),
$D_{p(u')}Af=AD_{p(k^{-1}u')}f=A(i\,D_{p(k^{-1}u)}f)=i\,D_{p(u)}Af$.

(iii) By (ii), $A(P_0)\subset P_0$. If $A(P_m)\subset P_m$, then for $u_i$ with
$p(u_i)=q_i$, (i) gives $A(D_{q_i}P_m)=D_{p(ku_i)}A(P_m)\subset P_{m+1}$, so
$A(P_{m+1})\subset P_{m+1}$.

\emph{Step 4: the orbit.} Let $\Phi\in\operatorname{PD}(M)$. As $M$ is regular,
$\Phi$ is a finite sum of functions $f\otimes g$ with $f\in\operatorname{Pol}(M)$
and $g$ Schwartz-Bruhat on $\mathbf A_f^I$ (the products $\bigotimes_v\phi_v$
with $\phi_v=\mathrm{vac}_v$ for almost all $v$ are Schwartz-Bruhat). By Steps 2
and 3 there is $m$ with $(\omega(k,1)\Phi)_y\in P_m$ for all $k\in K$ and $y$.

For a compact set $C\subset\mathbf A_f^I$ and a neighbourhood $U$ of $0$ in
$\mathbf A_f^I$, let $N(C,U)$ be the space of functions $\Psi$ on $\mathbf A^I$
with $\Psi_y\in P_m$ for every $y$, $\Psi(\xi,y)=0$ for $y\notin C$, and
$\Psi(\xi,y+u)=\Psi(\xi,y)$ for $u\in U$. It is finite-dimensional: choose
$y_1,\ldots,y_r$ with $C\subset\bigcup_j(y_j+U)$; then
$\Psi\mapsto(\Psi_{y_j})_j$ is injective into $P_m^r$. Every Schwartz-Bruhat
function $\Psi$ has $\Psi(\xi,y)=0$ for $y$ outside a compact set $C$ and is
invariant under translation of $y$ by a neighbourhood $U$ of $0$, as a locally
constant function of compact support on $\mathbf A_f^I$ is.

Let $C_1\subset C_2\subset\cdots$ be compact sets such that every compact
subset of $\mathbf A_f^I$ lies in some $C_a$, and $U_1\supset U_2\supset\cdots$
a basis of neighbourhoods of $0$. For each point $(\xi,y)$,
$k\mapsto(\omega(k,1)\Phi)(\xi,y)$ is continuous on $K$, as evaluation at a point
is admissible (Definition~\ref{asm-theta-agree-models}(6)(c)). So the sets
$Z_{a,b}$ of the $k\in K$ for which $\omega(k,1)\Phi$ vanishes for $y\notin C_a$
and is invariant under translation of $y$ by $U_b$ are closed, and they cover
$K$, as every $\omega(k,1)\Phi$ is Schwartz-Bruhat. The group $K$ is compact
(Lemma~\ref{asm-theta-agree-stabilizer}), so by Baire's theorem some $Z_{a,b}$
has an interior point $k_0$, and by compactness there are $k_1,\ldots,k_r\in K$
such that every $k\in K$ has $k_0k_j^{-1}k\in Z_{a,b}$ for some $j$. Then
$\omega(k_0k_j^{-1}k,1)\Phi\in N(C_a,U_b)$, and
$\omega(k,1)\Phi=\omega(k_jk_0^{-1},1)\,\omega(k_0k_j^{-1}k,1)\Phi$. So the
functions $\omega(k,1)\Phi$, $k\in K$, lie in the finite-dimensional space
$\sum_j\omega(k_jk_0^{-1},1)N(C_a,U_b)$.
\end{proof}
% Lean: kFinData_of hS (KFinDataOf.lean). Step 1 is vacs_fin (KFinAn.lean) for vacE (KFinArch.lean),
% with ray_zero for the ordinary differential equation and lamF_inj for the nondegeneracy of T;
% the pairs come from exists_ph and exists_mat_inf. Step 2 is fil, fil_fin, polM, polM_fin and
% exists_polM (actsD_dop, actsD_unique, dopL_mem). Step 3 is exists_slice and polM_stab with
% fil_stab; (i) is inter' and inter_ne, with KPol.hl_line_deriv (the derivative of Lemma
% asm-theta-agree-evaluation-derivative, which the Lean proves for a nonzero, hence the case
% split), line_slice and the admissibility of point evaluations (hs.adm); (ii) uses hk from
% IsStab. Step 4 is pure_eq_pureA, then orbit_fin with exists_lev, sb_unif, level_fin,
% CompactExhaustion.choice, nonempty_interior_of_iUnion_of_closed and a finite subcover; the
% compactness of K is StabEx (Lemma asm-theta-agree-stabilizer), and the continuity is IsUSplitting.cont.

\begin{lemma}\label{asm-theta-agree-split-inputs-kill}
Let $B$ be a theta base for a CM field $E$ with character $\psi$, place $v=v_1$
and embedding $t=\tau_1$ (Definition~\ref{asm-theta-agree-models}(1)), let $X$
be a Hermitian and $W$ a skew-Hermitian lattice space, of ranks $n$ and $d$ and
of some signatures, and let $M$ be a regular model for $(X,W)$, with
Schrödinger operators $\rho$ and form $\langle\,,\rangle$. Put
$I=\{1,\ldots,n\}\times\{1,\ldots,d\}$, and for $c\in F_v$ let $\bar c\in\mathbf A$
have component $c$ at $v$ and $0$ elsewhere.
\begin{enumerate}
\item Let $Z$ be a topological space, call a function on $Z$ Schwartz-Bruhat if
  it is locally constant with compact support, let $N$ be a subspace of the
  functions on $Z$, and let $Z_0\subset Z$. Suppose that for every
  $Y\notin Z_0$ there are a locally constant function $m$ on $Z$ and
  $c\in\mathbf C$ with $m(Y)\neq c$ and $(m-c)\Phi\in N$ for every
  Schwartz-Bruhat $\Phi$. Then every Schwartz-Bruhat $\Phi$ that vanishes on
  $Z_0$ lies in $N$.
\item Let $x,u\in\mathbb W(X,W)$ with $x$ at $v$ and $\pi_t(x)=0$, and put
  $C=t(J_X)\,\pi_{t\circ c}(x)\,t(J_W)^{\mathsf T}$. Then the component at $v$ of
  $\langle x,u\rangle$ is $\operatorname{tr}(C^{\mathsf T}\pi_t(u))$.
\item Let $L$ be a linear form on the functions on $\mathbf A^I$, and $S$ a set
  of triples $(p,q,\beta)$ with $p\in E^n$, $q\in E^d$ and $\beta\in F_v$.
  Suppose that $L(\rho(\bar c\,\varepsilon\,pq^{\mathsf T})\Phi)=\psi(\overline{\beta c})L(\Phi)$
  for every $(p,q,\beta)\in S$, $c\in F_v$, Schwartz-Bruhat $\Phi$ and
  $\varepsilon\in\mathbf A_E$ at $v$ with $\pi_t(\varepsilon)=1$ and
  $\pi_{t\circ c}(\varepsilon)=0$; and that every $C\in M_{n\times d}(F_v)$ with
  $t(p)^{\mathsf T}C\,t(q)=\beta$ for all $(p,q,\beta)\in S$ has a nonzero
  $r\in F_v^n$ with $r^{\mathsf T}C=0$. Then $L$ vanishes on
  $\mathcal S^\circ(M)$ (Definition~\ref{asm-theta-agree-models}(5)).
\item Let $e$ be a linear map from functions on $\mathbf A^I$ to functions on
  some $\mathbf A^{I'}$, and $x\in E^n$ nonzero, such that
  $e(\rho(xz^{\mathsf T})\Phi)=e(\Phi)$ for every vector $z$ over $\mathbf A_E$
  of length $d$ and every Schwartz-Bruhat $\Phi$. Then $e$ vanishes on
  $\mathcal S^\circ(M)$.
\end{enumerate}
\end{lemma}
\leanref{Development/WeilBound/ThetaLifts/SmallCones/SoKill.lean}{PositivityTheory.ThKCT.So.kill_set}
\leanref{Development/WeilBound/ThetaLifts/SmallCones/SoDefs.lean}{PositivityTheory.ThKCT.LinEx}
\leanref{Development/WeilBound/ThetaLifts/SmallCones/LinExOf.lean}{PositivityTheory.ThKCT.linEx_of}
\leanref{Development/WeilBound/ThetaLifts/SmallCones/SoKill.lean}{PositivityTheory.ThKCT.So.eigen_kill_set}
\leanref{Development/WeilBound/ThetaLifts/SmallCones/SoKill.lean}{PositivityTheory.ThKCT.So.so_kill}
\leanref{Development/WeilBound/ThetaLifts/SmallCones/FormSwapZero.lean}{PositivityTheory.ThKCT.So.ev_zero}
\begin{proof}
(1) For $Y\notin Z_0$ let $m_Y$ and $c_Y$ be as in the hypothesis, and
$U_Y=\{y:m_Y(y)=m_Y(Y)\}$, an open and closed set containing $Y$. Let $\Phi$
vanish on $Z_0$. Its support is compact and does not meet $Z_0$, so it is
covered by finitely many $U_{Y^1},\ldots,U_{Y^r}$. We argue by induction on $r$;
for $r=0$, $\Phi=0$. Let $U=U_{Y^1}$, $m=m_{Y^1}$ and $c=c_{Y^1}$. Then
$\mathbf 1_U\Phi=(m-c)\,(m(Y^1)-c)^{-1}\mathbf 1_U\Phi$, as $m=m(Y^1)$ on $U$, and
$\mathbf 1_U\Phi$ is Schwartz-Bruhat, so $\mathbf 1_U\Phi\in N$; and
$\mathbf 1_{Z\setminus U}\Phi$ is Schwartz-Bruhat, vanishes on $Z_0$ and has its
support covered by $U_{Y^2},\ldots,U_{Y^r}$, so it lies in $N$ by induction.

(2) The matrices $J_X$ and $J_W$ are invertible: for an infinite place $w$,
$\tau_w(J_X)$ and $i\,\tau_w(J_W)$ have the Cartan involutions $\theta_{X,w}$
and $\theta_{W,w}$, and a matrix $A$ with a Cartan involution $\theta$ has
$A\bar u\neq0$ for $u\neq0$, as $\operatorname{Re}((\theta u)^{\mathsf T}A\bar u)>0$.
So $t(J_X)$ and $t(J_W)$ are invertible. As $c(J_X)^{\mathsf T}=J_X$ and
$c(J_W)^{\mathsf T}=-J_W$, $t(c(J_X))=t(J_X)^{\mathsf T}$ and
$t(c(J_W))=-t(J_W)^{\mathsf T}$. So
Situation~\ref{asm-theta-agree-dictionary-situation} holds for $J=J_X$,
$J'=J_W$, $\varepsilon=1$ and the isomorphism of $M$. By
Lemma~\ref{asm-theta-agree-dictionary-coordinates}(2), as $x$ is at $v$, the
component at $v$ of $\langle x,u\rangle$ is
$\operatorname{tr}(\operatorname{co}(x)^{\mathsf T}\operatorname{co}'(u))-\operatorname{tr}(\operatorname{co}(u)^{\mathsf T}\operatorname{co}'(x))$.
By part (1) of that lemma, $\operatorname{co}'(x)=0$, as $\pi_t(x)=0$. As
$\operatorname{co}(x)=-C$ and $\operatorname{co}'(u)=-\pi_t(u)$, the component is
$\operatorname{tr}(C^{\mathsf T}\pi_t(u))$.

(3) The matrices of $X$ and $W$ are invertible
(Lemma~\ref{asm-theta-agree-convergence-signature}). Use $\rho_v$, $\operatorname{co}$
and the vectors $x_Y$ of Lemma~\ref{asm-theta-agree-coinvariants-eigen}: for
$Y\in M_{n\times d}(F_v)$, $x_Y$ is the element at $v$ with $\pi_t(x_Y)=0$ and
$\operatorname{co}(x_Y)=Y$, and $\operatorname{co}(x_Y)=-t(J_X)\pi_{t\circ c}(x_Y)t(J_W)^{\mathsf T}$.
Fix $\varepsilon$ as in (3), which exists by the identification in (2).

\emph{An eigenvalue argument.} Let $\iota$ be a linear model of the second kind
for $(t,1,1)$, $l$ a linear form on the functions on $F_v^I$, and $A$ a set of
pairs $(u,a)$, $u\in\mathbb W(X,W)$ at $v$ with $\pi_{t\circ c}(u)=0$ and
$a\in\mathbf C$, with $l(\rho_v(u)\phi)=a\,l(\phi)$ for Schwartz-Bruhat $\phi$.
Let $\phi$ be Schwartz-Bruhat with $\iota\phi(Y)=0$ for every $Y$ with
$\psi(\langle x_Y,u\rangle)=a$ for all $(u,a)\in A$. Then $l(\phi)=0$. Indeed,
apply (1) on $M_{n\times d}(F_v)$ to the image $N$ under $\iota$ of the
Schwartz-Bruhat functions in the kernel of $l$, and to the set $Z_0$ of these
$Y$. For $Y\notin Z_0$ choose $(u,a)\in A$ with $\psi(\langle x_Y,u\rangle)\neq a$
and put $m(y)=\psi(\langle x_y,u\rangle)$. By the definition of a linear model of
the second kind, $\iota(\rho_v(u)\phi')=m\cdot\iota\phi'$ for Schwartz-Bruhat
$\phi'$, and as in the proof of
Lemma~\ref{asm-theta-agree-coinvariants-eigen}(3), $m$ is locally constant.
Then $(m-a)\iota\phi'=\iota(\rho_v(u)\phi'-a\phi')\in N$. By (1),
$\iota\phi\in N$, and as $\iota$ is injective on Schwartz-Bruhat functions,
$l(\phi)=0$.

\emph{The conclusion.} $\mathcal S^\circ(M)$ is spanned by functions
$\Phi=f\otimes\bigotimes_w\phi_w$ with $\phi_v$ of full rank at $t$, so that
$\iota\phi_v(Y)=0$ for $Y$ of rank different from $n$, for a linear model $\iota$
of the second kind for $(t,1,1)$. Write $\Phi=\phi_v\otimes r$ with $r$ not
depending on the component at $v$, so that $\phi'\otimes r$ is Schwartz-Bruhat
for Schwartz-Bruhat $\phi'$, and put $l(\phi')=L(\phi'\otimes r)$. Let $A$ be the
set of the pairs $(\bar c\,\varepsilon\,pq^{\mathsf T},\psi(\overline{\beta c}))$ with
$(p,q,\beta)\in S$ and $c\in F_v$. Each such $u$ is at $v$ with
$\pi_{t\circ c}(u)=0$, and $\rho(u)(\phi'\otimes r)=(\rho_v(u)\phi')\otimes r$
(Lemma~\ref{asm-theta-agree-coinvariants-eigen}(2)), so the hypothesis gives
$l(\rho_v(u)\phi')=\psi(\overline{\beta c})\,l(\phi')$. Let $Y$ satisfy
$\psi(\langle x_Y,u\rangle)=a$ for all $(u,a)\in A$, and put
$C=t(J_X)\pi_{t\circ c}(x_Y)t(J_W)^{\mathsf T}=-Y$. For $(p,q,\beta)\in S$ let $b$ be
the component at $v$ of $\langle x_Y,\varepsilon pq^{\mathsf T}\rangle$; by (2),
$b=\operatorname{tr}(C^{\mathsf T}t(p)t(q)^{\mathsf T})=t(p)^{\mathsf T}C\,t(q)$.
Suppose $b\neq\beta$. By the proof of
Lemma~\ref{asm-theta-agree-coinvariants-eigen}(3) there is $s_0\in F_v$ with
$\psi(\bar s_0)\neq1$; with $c=s_0/(b-\beta)$,
$\langle x_Y,\bar c\varepsilon pq^{\mathsf T}\rangle=\bar c\langle x_Y,\varepsilon pq^{\mathsf T}\rangle=\overline{cb}=\bar s_0+\overline{\beta c}$,
and the condition on $Y$ gives $\psi(\bar s_0)\psi(\overline{\beta c})=\psi(\overline{\beta c})$,
which is false. So $t(p)^{\mathsf T}Ct(q)=\beta$ for all $(p,q,\beta)\in S$, and
by hypothesis $r^{\mathsf T}C=0$ for some $r\neq0$; so $Y=-C$ has rank different
from $n$, and $\iota\phi_v(Y)=0$. The eigenvalue argument gives
$L(\Phi)=l(\phi_v)=0$.

(4) Fix a point $y$ of $\mathbf A^{I'}$, and apply (3) to $L(\Phi)=e(\Phi)(y)$
and $S=\{(x,q,0):q\in E^d\}$. As
$\bar c\varepsilon xq^{\mathsf T}=x(\bar c\varepsilon q)^{\mathsf T}$, the
hypothesis on $e$ gives the first condition. If
$t(x)^{\mathsf T}C\,t(q)=0$ for every $q\in E^d$, then, taking for $q$ the unit
vectors, $t(x)^{\mathsf T}C=0$, and $t(x)\neq0$ as $t$ is injective.
\end{proof}
% Lean: (1) is kill_set (SoKill.lean), with the fibres {m_Y = m_Y(Y)}; it is
% Lemma asm-theta-agree-coinvariants-kill with a set in place of a point. (2) is the second part
% of LinEx, proved by linEx_of: Grams (isUnit_det_of_cartan for the invertibility, hX.herm and
% hW.skew with epsilon = 1), then sy, KCo.symp_eq (Lemma asm-theta-agree-dictionary-coordinates(2),
% with section_statement), ev_sgl, coord_snd (part (1) of that lemma) and dot_ofM for the
% trace; the first part of LinEx, a linear model of the second kind,
% is Lemma asm-theta-agree-coinvariants-dictionary(1). (3) is so_kill: the eigenvalue argument is
% eigen_kill_set (exists_coB, mult', lc_of_mul), the decomposition is pure_at and rho_tens, the
% element epsilon is exists_pair, s0 is GlobalTheory.PsiNontriv.psiV_ne_one, and
% Y has rank different from n by rank_ne_iff and coB_eq. (4) is ev_zero with vec_left and map_sgl.

\begin{lemma}\label{asm-theta-agree-split-inputs-cone-haar}
Let $G$ be a locally compact, second countable Hausdorff group with a left
Haar measure $\nu$, let $K\subset G$ be a compact subgroup, and let $P\subset G$
be a subset with $G=PK$. Let $m$ be a Borel measure on $G$ that is finite on
compact sets and is carried to itself by the left translations $g\mapsto pg$,
$p\in P$, and by the right translations $g\mapsto gk$, $k\in K$. Then
$m=c\,\nu$ for some $c\in[0,\infty)$.
\end{lemma}
\leanref{Development/WeilBound/ThetaLifts/SmallCones/RightOfCompact.lean}{PositivityTheory.ThKCT.Unf.eq_smul}
\leanref{Development/WeilBound/ThetaLifts/SmallCones/RightOfCompact.lean}{PositivityTheory.ThKCT.Unf.conv_eq}
\leanref{Development/WeilBound/ThetaLifts/SmallCones/RightOfCompact.lean}{PositivityTheory.ThKCT.Unf.cross_le}
\begin{proof}
For a measure $\lambda$ and a measurable $f\geq0$ write
$\lambda(f)=\int f\,d\lambda$. As $G$ is second countable and $m$ is finite on
compact sets, $m$ is $\sigma$-finite.

\emph{Step 1.} Let $f,b\colon G\to[0,\infty]$ be measurable with $b(gk)=b(g)$
for $g\in G$ and $k\in K$. Then
\[
\int b(g)\int f(xg)\,dm(x)\,d\nu(g)=\Bigl(\int b(x^{-1})\,dm(x)\Bigr)\,\nu(f).
\]
Indeed, write $g=pk$ with $p\in P$ and $k\in K$. The invariance of $m$ under
$x\mapsto px$ and that of $b$ under $k$ give
$\int b(x^{-1}g)\,dm(x)=\int b(x^{-1}k)\,dm(x)=\int b(x^{-1})\,dm(x)$. By
Tonelli's theorem and the left invariance of $\nu$, the left side is
$\int\int b(x^{-1}g)f(g)\,d\nu(g)\,dm(x)=\int f(g)\int b(x^{-1}g)\,dm(x)\,d\nu(g)$,
which is the right side.

\emph{Step 2.} Let $f,f_0\geq0$ be continuous with compact support. Then
$\nu(f)\,m(f_0)\leq\nu(f_0)\,m(f)$. Suppose not. The functions
$u(g)=\nu(f_0)\int f(xg)\,dm(x)$ and $v(g)=\nu(f)\int f_0(xg)\,dm(x)$ are
finite and continuous, and for $k\in K$ the invariance of $m$ under
$x\mapsto xk$ gives $u(k)=\nu(f_0)m(f)<\nu(f)m(f_0)=v(k)$. The open set
$\{u<v\}$ contains the compact set $K$, so it contains $VK$ for an open
neighbourhood $V$ of $1$. Let $C$ be a compact neighbourhood of $1$ and
$W=\operatorname{int}(V\cap C)\,K$. It is open, contains $1$, satisfies $Wk=W$
for $k\in K$, lies in $\{u<v\}$, and $m(W^{-1})\leq m((CK)^{-1})<\infty$. Step 1
with $b=\mathbf 1_W$, applied to $f$ and to $f_0$, gives
\[
\int_Wu\,d\nu=\nu(f_0)\,m(W^{-1})\,\nu(f)=\int_Wv\,d\nu ,
\]
a finite number. As $u<v$ on $W$ and $\nu(W)>0$, $\int_Wu\,d\nu<\int_Wv\,d\nu$,
a contradiction.

\emph{Step 3.} Choose $f_0$ as in Step 2 with $f_0(1)>0$. Then
$0<\nu(f_0)<\infty$; put $c=m(f_0)/\nu(f_0)$, which is finite. By Step 2 in both
orders, $m(f)=c\,\nu(f)$ for every continuous $f\geq0$ with compact support.
Two Borel measures finite on compact sets that agree on these functions are
equal, so $m=c\,\nu$.
\end{proof}
% Lean: Step 1 is conv_eq; the continuity of u and v is cont_M (continuous_integral_apply_inv_mul
% for the measure m^{-1}); Step 2 is cross_le, with compact_open_separated_mul_left for V and
% setLIntegral_strict_mono; Step 3 is eq_smul, with exists_continuous_nonneg_pos and
% Measure.ext_of_integral_eq_on_compactlySupported_nnreal.

\begin{lemma}\label{asm-theta-agree-split-inputs-cone-parabolic}
Let $E$ be a CM field with a theta base $B$
(Definition~\ref{asm-theta-agree-models}(1)), let $X$ be a Hermitian lattice
space of rank $n$ and some signature, and let $H$ be an adapted plane in $X$
over a lattice space $X_1$ (Definition~\ref{asm-theta-agree-planes}(1)), with
unipotent group $N_H$. Let
$\lambda\colon\mathbf A_E^\times\times U_{X_1}\to U_X$ be the continuous
homomorphism of Lemma~\ref{asm-theta-agree-split-inputs}(1)(c), so that
$\lambda(t,g_1)$ is the Levi element for $(t,g_1)$, and let $\nu_N$ be a Haar
measure on the closed subgroup $N_H$ of $U_X$. Then:
\begin{enumerate}
\item the map $(u,g_1,t)\mapsto u\,\lambda(t,g_1)$ from
  $N_H\times U_{X_1}\times\mathbf A_E^\times$ to $U_X$ is a homeomorphism onto a
  closed subset;
\item for $g_1\in U_{X_1}$ the map $u\mapsto\lambda(1,g_1)\,u\,\lambda(1,g_1)^{-1}$
  carries $\nu_N$ to itself, and for $t\in\mathbf A_E^\times$ the map
  $u\mapsto\lambda(t,1)\,u\,\lambda(t,1)^{-1}$ carries $\nu_N$ to
  $|t|^{1-n}\nu_N$.
\end{enumerate}
\end{lemma}
\leanref{Development/WeilBound/ThetaLifts/SmallCones/FcCols.lean}{PositivityTheory.ThKCT.Unf.parEx}
\leanref{Development/WeilBound/ThetaLifts/SmallCones/UnfAbs.lean}{PositivityTheory.ThKCT.Unf.Par}
\leanref{Development/WeilBound/ThetaLifts/SmallCones/UnfMod.lean}{PositivityTheory.ThKCT.Unf.modulus}
\begin{proof}
This is a computation in the coordinates of the frame of $H$
(Lemma~\ref{asm-theta-agree-evaluation-adapted-frame}), whose result we state.
In these coordinates there is a closed embedding of
$\mathbf A_E^{n-2}\times\mathbf A$ onto $N_H$ that carries an additive Haar
measure to a Haar measure on $N_H$, and the map of (1) is a closed embedding. Conjugation by $\lambda(1,g_1)$ acts on the first coordinate by a
map that preserves the Haar measure and fixes the second. Conjugation by
$\lambda(t,1)$ acts on the first coordinate by a scalar and on the second by
$N_{E/F}(t)$, and this substitution multiplies the Haar measure by
$|t|^{-(n-1)}$. Any two Haar measures on $N_H$ are proportional.
\end{proof}
% Lean: parEx (FcCols.lean) gives N = (Fin 1 -> Fin (n-2) -> A_E) x A_F with a measure muN, the
% map nn onto N_H, iota = lambda(1, .) and mm = lambda(., 1), and the record Par: emb is (1),
% h1 is the left invariance of muN, h2 and h3 are the two halves of (2) with chi(t) |t|^(n-1) = 1
% (modulus, UnfMod.lean), h4 says that iota and mm commute. The coordinates come from par2
% (UnfPar.lean), transported by Par.transport and parH (UnfParH.lean) through the frame fc of H.

\begin{lemma}\label{asm-theta-agree-split-inputs-cone-unfold}
Let $E$ be a CM field with a theta base $B$, let $X$ be a Hermitian lattice
space of rank $n$ and some signature, and let $H=(x,y,\iota,\mathfrak b)$ be an
adapted plane in $X$ over a lattice space $X_1$
(Definition~\ref{asm-theta-agree-planes}(1)). Let $K\subset U_X$ be an Iwasawa
subgroup for $x$ (Definition~\ref{asm-theta-agree-split-notions}(2)) with a
Haar measure $\mu_K$, let $\mu$ be a Haar measure on $\mathbf A_E^\times$, let
$\nu$ be a Haar measure on $U_X$ and $\mathcal D$ a fundamental domain for
$U_X(F)$ in $U_X$ with respect to $\nu$, let $\nu_1$ and $\mathcal D_1$ be the
same for $X_1$, and let $\lambda(t,g_1)$ be the Levi element for $(t,g_1)$
(Lemma~\ref{asm-theta-agree-split-inputs}(1)(c)). Then there is $\kappa>0$ such
that the following holds for every transporter family $\gamma$ for $x$ and every
continuous function $f$ on $U_X$ with $f(ug)=f(g)$ for $u\in N_H$,
$f(\delta g)=f(g)$ for $\delta\in U_X(F)$ with $\delta x=x$, and
\[
\int_K\int_{\mathbf A_E^\times}\Bigl(\int_{\mathcal D_1}|f(\lambda(l,g_1)k)|\,d\nu_1(g_1)\Bigr)|l|^{1-n}\,d\mu(l)\,d\mu_K(k)<\infty .
\]
For almost every $g\in\mathcal D$ the sum
$\Sigma(g)=\sum_{v\in\operatorname{Iso}(X)}f(\gamma(v)g)$ converges absolutely,
$\Sigma$ is integrable on $\mathcal D$, and
\[
\int_{\mathcal D}\Sigma\,d\nu=\kappa\int_K\int_{\mathbf A_E^\times}\Bigl(\int_{\mathcal D_1}f(\lambda(l,g_1)k)\,d\nu_1(g_1)\Bigr)|l|^{1-n}\,d\mu(l)\,d\mu_K(k).
\]
\end{lemma}
\leanref{Development/WeilBound/ThetaLifts/SmallCones/ThetaUniq.lean}{PositivityTheory.ThKCT.Unfold}
\leanref{Development/WeilBound/ThetaLifts/SmallCones/Unfold.lean}{PositivityTheory.ThKCT.unfold_of}
\leanref{Development/WeilBound/ThetaLifts/SmallCones/UnfMain.lean}{PositivityTheory.ThKCT.Unf.unfold_abs}
\leanref{Development/WeilBound/ThetaLifts/SmallCones/UnfAbs.lean}{PositivityTheory.ThKCT.Unf.haar_eq}
\leanref{Development/WeilBound/ThetaLifts/SmallCones/UnfAbs.lean}{PositivityTheory.ThKCT.Unf.fold}
\leanref{Development/WeilBound/CoreStatements/ThetaDilations/Nuabs.lean}{PositivityTheory.PlacePermOperatorBicommutantSpan.unf}
\leanref{Development/WeilBound/CoreStatements/ThetaDilations/Nuabs.lean}{PositivityTheory.PlacePermOperatorBicommutantSpan.nuabs}
\begin{proof}
Put $\Gamma=U_X(F)$, $Q=\{\delta\in\Gamma:\delta x=x\}$,
$\Gamma_N=\Gamma\cap N_H$ and $\Gamma_1=U_{X_1}(F)$. By
Lemma~\ref{asm-theta-agree-split-low}, $X_1$ is Hermitian of some signature, so
by Lemma~\ref{asm-theta-agree-convergence-signature} and
Theorem~\ref{asm-theta-agree-convergence-siegel}(1), $U_X$ and $U_{X_1}$ are
locally compact, second countable and Hausdorff, and $\Gamma$ and $\Gamma_1$
are countable and discrete. By Lemma~\ref{asm-theta-agree-split-inputs}(1)(c),
$\lambda(t,g_1)=\lambda(1,g_1)\lambda(t,1)=\lambda(t,1)\lambda(1,g_1)$. Let
$\nu_N$ be a Haar measure on $N_H$.

\emph{Step 1: decompositions.} Every $g\in U_X$ is $u\lambda(l,g_1)k$ with
$u\in N_H$, $g_1\in U_{X_1}$, $l\in\mathbf A_E^\times$ and $k\in K$: as $K$ is
an Iwasawa subgroup, $g=pk$ with $px=lx$, and by
Lemma~\ref{asm-theta-agree-split-inputs}(1)(e), $p=u\lambda(l,g_1)$. Every
$q\in Q$ is $\lambda(1,\gamma_1)\delta$ for exactly one pair
$(\gamma_1,\delta)\in\Gamma_1\times\Gamma_N$. Indeed, by
Lemma~\ref{asm-theta-agree-split-inputs}(1)(e) and (f), $q=u\lambda(1,\gamma_1)$
with $u\in N_H$, $\gamma_1\in\Gamma_1$ and $\lambda(1,\gamma_1)\in\Gamma$; then
$\delta=\lambda(1,\gamma_1)^{-1}u\lambda(1,\gamma_1)$ lies in $N_H$ by part (g)
of that lemma, and in $\Gamma$ as a product of rational elements. If
$\lambda(1,\gamma_1)\delta=\lambda(1,\gamma'_1)\delta'$, then
$(\lambda(1,\gamma_1)\delta\lambda(1,\gamma_1)^{-1})\lambda(1,\gamma_1)=(\lambda(1,\gamma'_1)\delta'\lambda(1,\gamma'_1)^{-1})\lambda(1,\gamma'_1)$,
the conjugates lying in $N_H$, and
Lemma~\ref{asm-theta-agree-split-inputs-cone-parabolic}(1) gives
$\gamma_1=\gamma'_1$ and then $\delta=\delta'$. Conversely
$\Gamma_N\subset Q$ and $\lambda(1,\Gamma_1)\subset Q$, as elements of $N_H$
and Levi elements for $(1,g_1)$ fix $x$.

\emph{Step 2: Haar measure.} Let $m$ be the image of
$\nu_N\times\nu_1\times|t|^{1-n}\mu\times\mu_K$ under
$(u,g_1,t,k)\mapsto u\lambda(1,g_1)\lambda(t,1)k$. By
Lemma~\ref{asm-theta-agree-split-inputs-cone-parabolic}(1) and the compactness
of $K$, $m$ is finite on compact sets. It is carried to itself by the left
translations by $u_0\in N_H$, by $\lambda(1,g_0)$ and by $\lambda(t_0,1)$: for
instance
$\lambda(t_0,1)u\lambda(1,g_1)\lambda(t,1)=u'\lambda(1,g_1)\lambda(t_0t,1)$ with
$u'=\lambda(t_0,1)u\lambda(t_0,1)^{-1}$; by
Lemma~\ref{asm-theta-agree-split-inputs-cone-parabolic}(2) the substitution
$u\mapsto u'$ multiplies $\nu_N$ by $|t_0|^{1-n}$, and $t\mapsto t_0t$ carries
$|t|^{1-n}\mu$ to $|t_0|^{n-1}|t|^{1-n}\mu$. Similarly for $\lambda(1,g_0)$,
using part (2) of that lemma and the invariance of $\nu_1$. Haar measure on the
compact group $K$ is invariant under right translations, so $m$ is carried to
itself by the right translations by $K$. By Step 1, $U_X=PK$ for the set $P$ of
the $u\lambda(t,g_1)$. Lemma~\ref{asm-theta-agree-split-inputs-cone-haar} gives
$c$ with $m=c\,\nu$, and $c>0$ as $m\neq0$. So for measurable
$\varphi\geq0$,
\[
c\int\varphi\,d\nu=\int_K\int_{\mathbf A_E^\times}\Bigl(\int_{U_{X_1}}\int_{N_H}\varphi(u\lambda(1,g_1)\lambda(t,1)k)\,d\nu_N(u)\,d\nu_1(g_1)\Bigr)|t|^{1-n}\,d\mu(t)\,d\mu_K(k).\tag{M}
\]

\emph{Step 3: folding.} Let $\mathcal F\subset U_X$ be a measurable set meeting
every orbit $\Gamma_Ng$ in exactly one point, and put
$a=\nu_N(\mathcal F\cap N_H)$. Then $0<a<\infty$. Indeed,
$\mathcal F\cap N_H$ is a fundamental domain for the countable group
$\Gamma_N$ acting on $N_H$, so $a>0$. By
Lemma~\ref{asm-theta-agree-split-inputs}(1)(b) there is a compact
$C\subset N_H$ with $N_H=\Gamma_NC$; the sets $\delta^{-1}\mathcal F$,
$\delta\in\Gamma_N$, are disjoint, so
$a\leq\sum_\delta\nu_N(\mathcal F\cap\delta C)=\sum_\delta\nu_N(\delta^{-1}\mathcal F\cap C)\leq\nu_N(C)<\infty$.
Let $\beta\colon U_X\to[0,\infty]$ be measurable with
$\sum_{q\in Q}\beta(qg)=1$ for every $g$, and let $\varphi\geq0$ be measurable
with $\varphi(ug)=\varphi(g)$ for $u\in N_H$ and $\varphi(qg)=\varphi(g)$ for
$q\in Q$. We show, for $y\in U_X$,
\[
\int_{U_{X_1}}\int_{N_H}\beta(u\lambda(1,g_1)y)\varphi(u\lambda(1,g_1)y)\,d\nu_N(u)\,d\nu_1(g_1)=a\int_{\mathcal D_1}\varphi(\lambda(1,g_1)y)\,d\nu_1(g_1).\tag{F}
\]
First, for $g\in U_X$,
$\sum_{\gamma_1\in\Gamma_1}\int_{N_H}\beta(u\lambda(1,\gamma_1)g)\,d\nu_N(u)=a$:
by Lemma~\ref{asm-theta-agree-split-inputs-cone-parabolic}(2) the term at
$\gamma_1$ is
$\int_{N_H}\beta(\lambda(1,\gamma_1)ug)\,d\nu_N(u)=\sum_{\delta\in\Gamma_N}\int_{\mathcal F\cap N_H}\beta(\lambda(1,\gamma_1)\delta ug)\,d\nu_N(u)$,
and by Step 1 the sum over $(\gamma_1,\delta)$ is
$\int_{\mathcal F\cap N_H}\sum_{q\in Q}\beta(qug)\,d\nu_N(u)=a$. Now
$\varphi(u\lambda(1,g_1)y)=\varphi(\lambda(1,g_1)y)$, and as $\mathcal D_1$ is a
fundamental domain for $\Gamma_1$ and $\varphi(\lambda(1,\gamma_1)g)=\varphi(g)$,
the left side of (F) is
\[
\int_{\mathcal D_1}\sum_{\gamma_1\in\Gamma_1}\Bigl(\int_{N_H}\beta(u\lambda(1,\gamma_1)\lambda(1,g_1)y)\,d\nu_N(u)\Bigr)\varphi(\lambda(1,g_1)y)\,d\nu_1(g_1),
\]
which is the right side of (F).

\emph{Step 4: unfolding.} For $\delta\in\Gamma$ there is exactly one
$v\in\operatorname{Iso}(X)$ with $\delta\gamma(v)^{-1}\in Q$, namely
$v=\delta^{-1}x$: this vector is isotropic and nonzero, as $\delta$ is a
rational unitary matrix, and $\delta\gamma(v)^{-1}x=\delta v$. So the
$\gamma(v)$ represent the cosets $Q\backslash\Gamma$, and the set
$\mathcal F'=\bigcup_v\gamma(v)\mathcal D$ is a fundamental domain for $Q$ in
$U_X$ with respect to $\nu$, with
$\int_{\mathcal D}\sum_v\varphi(\gamma(v)g)\,d\nu(g)=\int_{\mathcal F'}\varphi\,d\nu$
for measurable $\varphi\geq0$, by the left invariance of $\nu$. Let
$\mathcal G\subset U_X$ be a measurable set meeting every orbit $Qg$ in exactly
one point; it is also a fundamental domain for $Q$, and
$\sum_{q\in Q}\mathbf 1_{\mathcal G}(qg)=1$. For $\varphi$ as in Step 3,
$\int_{\mathcal F'}\varphi\,d\nu=\int_{\mathcal G}\varphi\,d\nu$, and (M) for
$\mathbf 1_{\mathcal G}\varphi$ and (F) with $\beta=\mathbf 1_{\mathcal G}$ and
$y=\lambda(t,1)k$ give
\[
\int_{\mathcal D}\sum_v\varphi(\gamma(v)g)\,d\nu(g)=\frac ac\int_K\int_{\mathbf A_E^\times}\Bigl(\int_{\mathcal D_1}\varphi(\lambda(t,g_1)k)\,d\nu_1(g_1)\Bigr)|t|^{1-n}\,d\mu(t)\,d\mu_K(k).\tag{K}
\]

\emph{Step 5: the function $f$.} Put $\kappa=a/c$. By (K) for $\varphi=|f|$, the
integral of $\sum_v|f(\gamma(v)g)|$ over $\mathcal D$ is finite. So the sum
converges absolutely for almost every $g\in\mathcal D$, and $\Sigma$ is
integrable on $\mathcal D$. For every measurable
$\psi\colon\mathbf C\to[0,\infty]$, (K) for $\varphi=\psi\circ f$ says that the
image under $f$ of the sum over $v$ of the images of $\nu|_{\mathcal D}$ under
$g\mapsto\gamma(v)g$ is the image under $f$ of $\kappa$ times the image of
$\mu_K\times|t|^{1-n}\mu\times\nu_1|_{\mathcal D_1}$ under
$(k,t,g_1)\mapsto\lambda(t,g_1)k$. So $f$ has the same integral against these
two measures, against both of which it is integrable by the case $\psi=|\cdot|$.
By Fubini's theorem these integrals are the two sides of the claimed identity.
\end{proof}
% Lean: unfold_of (Unfold.lean) applies unfold_abs (UnfMain.lean) with Gamma = ratU, Q = fixQ,
% Gamma_N = ratU inf unip and Gamma_1 = ratU of X_1, the data of parEx, levi_uniq for lambda,
% stabDec_of for Step 1 (with levi_conj and the injectivity from hP.emb), and exists_iso for the
% vector delta^{-1} x. Step 2 is haar_eq (eq_smul, nuabs for the left invariance,
% right_of_compact for K); Step 3 is fold, with exists_cut and the bound l 𝓕 <= l (nn '' C);
% Step 4 is unf (Nuabs.lean), f1.setLIntegral_eq and key; Step 5 is the end of unfold_abs, with
% integral_eq_of for the images of the two measures.

\begin{lemma}\label{asm-theta-agree-split-inputs-cone-tate}
Let $E$ be a CM field with a theta base $B$, let $\mu$ be a Haar measure on
$\mathbf A_E^\times$ and $z\geq2$ an integer. Then there is $N\geq0$ such that for
every compact set $C\subset\mathbf A_E$
\[
\int_{A(C)}|l|^z\Bigl(1+\sum_w|\pi_w(l)|\Bigr)^{-N}d\mu(l)<\infty ,
\]
where the sum is over the infinite places $w$ of $F$, $\pi_w$ is as in
Definition~\ref{asm-theta-agree-models}(1), and $A(C)$ is the set of
$l\in\mathbf A_E^\times$ such that $l-c$ is at $\infty$ for some $c\in C$.
\end{lemma}
\leanref{Development/WeilBound/ThetaLifts/SmallCones/ThetaUniq.lean}{PositivityTheory.ThKCT.TateAbs}
\leanref{Development/WeilBound/ThetaLifts/SmallCones/TateAbs.lean}{PositivityTheory.ThKCT.tateAbs}
\leanref{GroundTheory/GlobalFields/TateThesis/GlobalZeta.lean}{GroundTheory.Global.integrable_zetaBound}
\begin{proof}
Identify $\mathbf A_E$ with the adele ring of $E$, so that $\mu$ is a Haar
measure on the idele group of $E$ and $|l|$ is the idele norm,
$|l|=\prod_W|l_W|_W\cdot|l_f|$ with $W$ running over the infinite places of $E$,
$|\cdot|_W$ the normalized absolute value and $|l_f|$ the norm of the finite
part. For $l\in A(C)$ the finite part of $l$ is that of some $c\in C$. The
finite parts of the elements of $C$ form a compact set, contained in the set
$\hat L$ of finite adeles lying in $L$ at every place, for a nonzero fractional
ideal $L$ of $E$. Let $r$ be the number of infinite places of $E$, $k=2z+3$ and
$N=kr$. Each $|l_W|$ is $|\pi_w(l)|$ for the place $w$ of $F$ below $W$, so
$(1+\sum_w|\pi_w(l)|)^N\geq\prod_W(1+|l_W|)^k$, and on $A(C)$
\[
|l|^z\Bigl(1+\sum_w|\pi_w(l)|\Bigr)^{-N}\leq\prod_W\frac{|l_W|_W^{\,z}}{(1+|l_W|)^k}\cdot\mathbf 1_{\hat L}(l_f)\,|l_f|^z .
\]
The right side is a product of local functions, and it is integrable on the
ideles of $E$: its integral is the product of the local integrals, which
converge at the infinite places as $k>2z$, and whose product over the finite
places is bounded by a multiple of $\zeta_E(z)<\infty$, as $z>1$. This is the
absolute convergence of global zeta integrals for $\operatorname{Re}s>1$ in
Tate's thesis \cite[\S4.4]{Tat67}, here at $s=z$ with $k>2z+2$.
\end{proof}
% Lean: tateAbs (TateAbs.lean): L from exists_subset_idealFiniteAdeles, N = (2z+3) times the
% number of infinite places of E, Tate.absE_eq, Tate.norm_le and Tate.snd_eq for the bound, and
% integrable_zetaBound (TateThesis/GlobalZeta.lean) for the integrability, transported to mu by
% Tate.transport.

\begin{lemma}\label{asm-theta-agree-split-inputs-cone-line}
Let $B$ be a theta base for a CM field $E$, let $M$ and $M_1$ be regular
models for lattice spaces $(X,W)$ and $(X_1,W_1)$, of ranks $(n,d)$ and
$(n_1,d_1)$, with Schrödinger operators $\rho$ for $M$
(Definition~\ref{asm-theta-agree-models}(3)), and let $e$ be an evaluation map
for $(M,M_1,\iota_X,\iota_W,\mathcal J)$
(Definition~\ref{asm-theta-agree-planes}(5)). Put
$I_1=\{1,\ldots,n_1\}\times\{1,\ldots,d_1\}$, and let $\mu_1$ and $\mathcal D_1$
be the measure and the domain of $M_1$. Let $v_0,z_0\in M_{n\times d}(E)$ with
$az_0\in\mathcal J$ for every $a\in\mathbf A_E$ and $h(z_0,v_0)\neq0$, for the
form $h$ of Definition~\ref{asm-theta-agree-models}(3). Let $\tau$ be the
identity or $c$, $\mu$ a Haar measure on $\mathbf A_E^\times$, $z\geq2$ an
integer, $\nu$ a Hecke character of $E$, and $\Psi$ a Schwartz-Bruhat function
on $\mathbf A^{n\times d}$. For $l\in\mathbf A_E^\times$ put
$\Psi_l=e(\rho(\tau(l)v_0)\Psi)$.
\begin{enumerate}
\item Let $\omega_1$ be a representation of $U_{X_1}\times U_{W_1}$ on
  functions on $\mathbf A^{I_1}$ that maps Schwartz-Bruhat functions to
  Schwartz-Bruhat functions, and let $\mathcal F$ be a set of Schwartz-Bruhat
  functions on $\mathbf A_f^{I_1}$ such that for every $\phi\in\mathcal F$
  there are $k\geq0$, a continuous seminorm $p$ on $\mathcal S(F_\infty^{I_1})$
  and a measurable $G\colon U_{X_1}\to[0,\infty]$ with
  $\int_{\mathcal D_1}G\,d\mu_1<\infty$ and
  $|\theta_{\omega_1}(f\otimes\phi)(g,h)|\leq p(f)\operatorname{ht}(h)^kG(g)$ for
  all $f$, $g$ and $h$. Suppose that $\mathcal F$ contains every
  Schwartz-Bruhat function, or that $\mathcal F$ contains every Schwartz-Bruhat
  $\phi$ with $f\otimes\phi\in\mathcal S^\circ(M_1)$ for all $f$ and
  $e(\rho(av_0)\Psi)\in\mathcal S^\circ(M_1)$ for every $a\in\mathbf A_E$. Then
  for every $h_1\in U_{W_1}$
  \[
  \int_{\mathbf A_E^\times}\Bigl(\int_{\mathcal D_1}|\theta_{\omega_1}(\Psi_l)(g,h_1)|\,d\mu_1(g)\Bigr)|l|^z\,d\mu(l)<\infty .
  \]
\item There is a Schwartz-Bruhat function $\Psi_1$ on $\mathbf A^{I_1}$ such
  that for every admissible linear form $\lambda$ on the functions on
  $\mathbf A^{I_1}$ (Theorem~\ref{asm-theta-agree-evaluation-heisenberg}) the
  function $l\mapsto\lambda(\Psi_l)\nu(l)|l|^z$ is $\mu$-integrable, with
  integral $\lambda(\Psi_1)$.
\end{enumerate}
\end{lemma}
\leanref{Development/WeilBound/ThetaLifts/SmallCones/ThetaUniq.lean}{PositivityTheory.ThKCT.LineInt}
\leanref{Development/WeilBound/ThetaLifts/SmallCones/LineInt.lean}{PositivityTheory.ThKCT.lineInt_of}
\leanref{Development/WeilBound/ThetaLifts/SmallCones/NormPiInfSmul.lean}{PositivityTheory.ThKCT.Line.nice}
\leanref{Development/WeilBound/ThetaLifts/SmallCones/NormPiInfSmul.lean}{PositivityTheory.ThKCT.Line.exists_dual}
\leanref{Development/WeilBound/ThetaLifts/SmallCones/LineInt.lean}{PositivityTheory.ThKCT.Line.dom_int}
\leanref{Development/WeilBound/ThetaLifts/SmallCones/LineInt.lean}{PositivityTheory.ThKCT.Line.exists_so}
\begin{proof}
For $y\in\mathbf A^m$ put $s(y)=\sum_j\sum_w|y_{j,w}|$, as in
Lemma~\ref{asm-theta-agree-local-bump-slices}, and write
$\langle\,,\rangle$ for the form of Definition~\ref{asm-theta-agree-models}(3).

\emph{Step 1: a decomposition.} We show that there are Schwartz-Bruhat
functions $\phi_1,\ldots,\phi_r$ on $\mathbf A_f^{I_1}$ and continuous maps
$g_i\colon\mathbf A_E\to\mathcal S(F_\infty^{I_1})$ with
$e(\rho(\tau(a)v_0)\Psi)=\sum_ig_i(a)\otimes\phi_i$ for every $a\in\mathbf A_E$,
such that $l\mapsto q(g_i(l))|l|^z$ is $\mu$-integrable on $\mathbf A_E^\times$
for every continuous seminorm $q$. The set of $\Psi$ with this property is a
subspace, so we may take $\Psi=f\otimes\phi$. Let $b_1,b_2$ be a basis of $E$
over $F$, which is a basis of $\mathbf A_E$ over $\mathbf A$, and $b_1^\vee,b_2^\vee$
the dual basis for $\operatorname{Tr}_{E/F}$. Put $h_0=h(z_0,v_0)$,
$u_j=b_jv_0$ and $z_j=-c(b_j^\vee)h_0^{-1}z_0$. Then $az_j\in\mathcal J$ for
$a\in\mathbf A$, and as $h(\alpha u,\beta u')=\alpha c(\beta)h(u,u')$,
$\langle z_i,u_j\rangle=\operatorname{Tr}_{E/F}c(b_jb_i^\vee)=\delta_{ij}$. Let
$\phi_{1,i}$, $C$ and $A_i(y)$ be given by
Lemma~\ref{asm-theta-agree-local-bump-slices} for these vectors and $\phi$. For
$a\in\mathbf A_E$ let $y(a)\in\mathbf A^2$ be the coordinates of $\tau(a)$, so
that $\tau(a)=\sum_jy_j(a)b_j$; they are continuous in $a$, and
$\sum_jy_j(a)u_j=\tau(a)v_0$. By parts (1) and (2) of that lemma,
$g_i(a)=A_i(y(a))f$ and $\phi_i=\phi_{1,i}$ give the decomposition. Let $q$ be a
continuous seminorm, $N$ as in Lemma~\ref{asm-theta-agree-split-inputs-cone-tate}
for $z$, and $p$ as in Lemma~\ref{asm-theta-agree-local-bump-slices}(4) for $q$
and $N$. As $\tau$ is an involution, $a=\sum_jy_j(a)\tau(b_j)$, so
$\sum_w|\pi_w(a)|\leq s(y(a))\,K_0$ with
$K_0=\sum_j\sum_w|\pi_w(\tau(b_j))|$. Hence, by parts (3) and (4) of that lemma,
$q(g_i(l))|l|^z$ is $0$ unless the finite component of $y(l)$ lies in $C$, and
is at most $p(f)(1+K_0)^N|l|^z(1+\sum_w|\pi_w(l)|)^{-N}$. If the finite
component of $y(l)$ lies in $C$, then $l-c$ is at $\infty$ for
$c=\sum_j\bar y_j\tau(b_j)$, where $\bar y_j$ has component $0$ at $\infty$ and
the finite component of $y_j(l)$ elsewhere; these $c$ form a compact set. So
Lemma~\ref{asm-theta-agree-split-inputs-cone-tate} gives the integrability.

\emph{Step 2: proof of (1).} In the second case, let $V$ be the space of the
Schwartz-Bruhat $\phi'$ with $f\otimes\phi'\in\mathcal S^\circ(M_1)$ for all $f$,
choose a linear projection $\varpi$ from the functions on $\mathbf A_f^{I_1}$
onto $V$, and replace each $\phi_i$ by $\varpi(\phi_i)\in V$. This does not
change $\sum_ig_i(a)\otimes\phi_i$ when the sum lies in $\mathcal S^\circ(M_1)$,
which holds for $a=\tau(l)$: for every point $x$ of $F_\infty^{I_1}$ the
function $\sum_ig_i(a)(x)\phi_i$ is the restriction of the sum to the points
with archimedean component $x$, which lies in $V$ as $\mathcal S^\circ(M_1)$ is
spanned by products with full-rank factor at $v_1$; so $\varpi$ fixes it. In
both cases $\Psi_l=\sum_ig_i(l)\otimes\phi_i$ with $\phi_i\in\mathcal F$. As
$\omega_1(g,h_1)$ preserves Schwartz-Bruhat functions and the theta sum is
linear on them,
$|\theta_{\omega_1}(\Psi_l)(g,h_1)|\leq\sum_ip_i(g_i(l))\operatorname{ht}(h_1)^{k_i}G_i(g)$,
with the data $k_i,p_i,G_i$ of $\phi_i$. The integral of (1) is therefore at
most
$\sum_i\operatorname{ht}(h_1)^{k_i}\int_{\mathcal D_1}G_i\,d\mu_1\int p_i(g_i(l))|l|^z\,d\mu(l)$,
which is finite by Step 1.

\emph{Step 3: proof of (2).} Put $\varsigma(l)=\nu(l)|l|^z$, a continuous function
with $|\varsigma(l)|=|l|^z$. The maps $l\mapsto\varsigma(l)g_i(l)$ are continuous,
and every Schwartz seminorm of them is integrable by Step 1. By
Lemma~\ref{asm-theta-agree-local-bump-schwartz}(2) there are
$f_i\in\mathcal S(F_\infty^{I_1})$ with
$T(f_i)=\int T(\varsigma(l)g_i(l))\,d\mu(l)$ for every continuous linear form
$T$. Put $\Psi_1=\sum_if_i\otimes\phi_i$. For admissible $\lambda$ the forms
$T_i(f')=\lambda(f'\otimes\phi_i)$ are continuous, and
$\lambda(\Psi_l)\varsigma(l)=\sum_iT_i(\varsigma(l)g_i(l))$ is integrable with integral
$\sum_iT_i(f_i)=\lambda(\Psi_1)$.
\end{proof}
% Lean: lineInt_of (LineInt.lean) from Slice, SchInt and TateAbs. Step 1 is nice
% (NormPiInfSmul.lean), by span induction, with exists_dual for the z_j (traceDual) and the
% Slice data of Lemma asm-theta-agree-local-bump-slices, the coordinates y of Bt.equivFun, the
% bound through norm_piInf_smul and the constant K, and TateAbs for the image of C. Step 2 is
% dom_int, with exists_so and Line.contract for the second case (soSub, a left inverse of its
% inclusion). Step 3 is the second half of lineInt_of, through the property SchInt.

\begin{lemma}\label{asm-theta-agree-split-inputs-cone-exchange}
Let $(X,W,M,\chi_X,\chi_W,\omega)$ and $(X_1,W_1,M_0,\chi_X,\chi_W,\omega_0)$ be
Weil pairs of some signatures for a theta base $B$ of a CM field $E$
(Definition~\ref{asm-theta-agree-triple}(1)), with $X$ and $W$ of ranks $n$ and
$d$, with Schrödinger operators $\rho$ for $M$ and with measure $\mu_0$ and
domain $\mathcal D_0$ of $M_0$. Let $e$ be an evaluation map for
$(M,M_0,\iota_X,\iota_W,\mathcal J)$, $K\subset U_X$ a compact subgroup with a
Haar measure $\mu_K$, $\mu$ a Haar measure on $\mathbf A_E^\times$,
$u\colon\mathbf A_E^\times\to\mathbb W(X,W)$, $\nu$ a Hecke character and
$z\geq0$ an integer. Let $\mathcal S$ be a space of Schwartz-Bruhat functions on
$\mathbf A^{n\times d}$ with $\omega(g,h)\mathcal S\subset\mathcal S$ for all
$(g,h)$, and suppose:
\begin{enumerate}
\item[(i)] for $\Psi\in\mathcal S$ and $h_1\in U_{W_1}$, the function
  $(l,g)\mapsto\theta_{\omega_0}(e(\rho(u(l))\Psi))(g,h_1)\,\nu(l)|l|^z$ is
  integrable on $\mathbf A_E^\times\times\mathcal D_0$ for $\mu\times\mu_0$;
\item[(ii)] for every Schwartz-Bruhat $\Psi$ there is a Schwartz-Bruhat
  function $\Psi'$ with
  $\lambda(\Psi')=\int\lambda(e(\rho(u(l))\Psi))\,\nu(l)|l|^z\,d\mu(l)$, the
  integrand being integrable, for every admissible linear form $\lambda$.
\end{enumerate}
Let $\Phi\in\mathcal S$ be $K$-finite for $\omega$, and let $I^{\mathrm{cone}}$
be the cone integral for $(M,\omega,M_0,\omega_0,e,K,\mu_K,\mu,u,\nu,z)$
(Definition~\ref{asm-theta-agree-split-notions}(4)) and $\operatorname{cone}$
the cone for $(M,\omega,e,K,\mu_K,\mu,u,\nu,z)$
(Definition~\ref{asm-theta-agree-planes}(7)). Then for every
$h_1\in U_{W_1}$:
\begin{enumerate}
\item $\int_K\int_{\mathbf A_E^\times}\int_{\mathcal D_0}|\theta_{\omega_0}(e(\rho(u(l))\omega(k,1)\Phi))(g,h_1)|\,|l|^z\,d\mu_0(g)\,d\mu(l)\,d\mu_K(k)<\infty$;
\item $I^{\mathrm{cone}}(\Phi)(h_1)=I_{M_0,\omega_0}(\operatorname{cone}(\Phi))(h_1)$.
\end{enumerate}
\end{lemma}
\leanref{Development/WeilBound/ThetaLifts/SmallCones/ConeEx.lean}{PositivityTheory.ThKCT.LineOK}
\leanref{Development/WeilBound/ThetaLifts/SmallCones/ConeEx.lean}{PositivityTheory.ThKCT.LineOK.coneI_eq}
\leanref{Development/WeilBound/ThetaLifts/SmallCones/ConeUnf.lean}{PositivityTheory.ThKCT.LineOK.fin}
\leanref{Development/WeilBound/ThetaLifts/SmallCones/ConeBasic.lean}{PositivityTheory.ThKCT.kfin_dec}
\begin{proof}
\emph{Step 1: $K$-finiteness.} The span $V$ of the $\omega(k,1)\Phi$, $k\in K$,
is finite-dimensional and lies in $\mathcal S$. Choose a basis
$\Psi_1,\ldots,\Psi_m$ of $V$ and write $\omega(k,1)\Phi=\sum_ic_i(k)\Psi_i$.
Evaluation at a point of $\mathbf A^{n\times d}$ is an admissible form, so
$k\mapsto(\omega(k,1)\Phi)(y)$ is continuous for every $y$
(Definition~\ref{asm-theta-agree-models}(6)(c)); as $V$ is
finite-dimensional, the topology of pointwise convergence on it is its usual
topology, so the $c_i$ are continuous, and they are $\mu_K$-integrable on the
compact group $K$.

\emph{Step 2: two identities.} The functions $e(\rho(u(l))\Psi)$ with $\Psi$
Schwartz-Bruhat are Schwartz-Bruhat
(Theorem~\ref{asm-theta-agree-evaluation-heisenberg}(1)). For $g$ and $h_1$, the
form $\Psi'\mapsto\theta_{\omega_0}(\Psi')(g,h_1)=\theta(\omega_0(g,h_1)\Psi')$
agrees on the Schwartz-Bruhat functions with an admissible form, by
Theorem~\ref{asm-theta-agree-evaluation-heisenberg}(2) and
Definition~\ref{asm-theta-agree-models}(6). Let $\Lambda_i$ be the function
$\Psi'$ of (ii) for $\Psi_i$. Then (ii) for the evaluation at a point $y$ and for
this form give
\[
\Lambda_i(y)=\int e(\rho(u(l))\Psi_i)(y)\,\nu(l)|l|^z\,d\mu(l),\qquad
\theta_{\omega_0}(\Lambda_i)(g,h_1)=\int\theta_{\omega_0}(e(\rho(u(l))\Psi_i))(g,h_1)\,\nu(l)|l|^z\,d\mu(l).
\]

\emph{Step 3: proof of (1).} Let
$J_i(l,g)=\theta_{\omega_0}(e(\rho(u(l))\Psi_i))(g,h_1)\,\nu(l)|l|^z$, which is
integrable by (i). By linearity, the integrand of (1) is
$|\sum_ic_i(k)J_i(l,g)|$, so the integral is at most
$\sum_i\int_K|c_i|\,d\mu_K\int|J_i|\,d(\mu\times\mu_0)<\infty$.

\emph{Step 4: the cone.} By the first identity of Step 2 and linearity,
$\int e(\rho(u(l))\omega(k,1)\Phi)(y)\,\nu(l)|l|^z\,d\mu(l)=\sum_ic_i(k)\Lambda_i(y)$,
so $\operatorname{cone}(\Phi)=\sum_i\bigl(\int_Kc_i\,d\mu_K\bigr)\Lambda_i$.

\emph{Step 5: proof of (2).} By the second identity of Step 2 and Fubini's
theorem for $J_i$, $g\mapsto\theta_{\omega_0}(\Lambda_i)(g,h_1)$ is integrable on
$\mathcal D_0$, and
$T_i=I_{M_0,\omega_0}(\Lambda_i)(h_1)$ is the integral of the integrable
function $l\mapsto I_{M_0,\omega_0}(e(\rho(u(l))\Psi_i))(h_1)\,\nu(l)|l|^z$. For
almost every $l$ all the functions $g\mapsto J_i(l,g)$ are integrable on
$\mathcal D_0$; for such $l$, linearity gives
$I_{M_0,\omega_0}(e(\rho(u(l))\omega(k,1)\Phi))(h_1)=\sum_ic_i(k)I_{M_0,\omega_0}(e(\rho(u(l))\Psi_i))(h_1)$.
Integrating over $l$, the inner integral of $I^{\mathrm{cone}}(\Phi)(h_1)$ at
$k$ is $\sum_ic_i(k)T_i$. So
$I^{\mathrm{cone}}(\Phi)(h_1)=\sum_i\bigl(\int_Kc_i\,d\mu_K\bigr)T_i$, which by
Step 4 and the linearity of the theta integral on the integrable functions
$g\mapsto\theta_{\omega_0}(\Lambda_i)(g,h_1)$ is
$I_{M_0,\omega_0}(\operatorname{cone}(\Phi))(h_1)$.
\end{proof}
% Lean: Step 1 is kfin_dec (ConeBasic.lean), with ev_adm and pair_cont. The hypotheses are the
% fields of LineOK (ConeEx.lean): le, stab, sch, th (pair_adm with the admissible theta sum), int
% for (i) and lam for (ii); cone_core builds LineOK from the Weil pairs. (1) is LineOK.fin
% (ConeUnf.lean), (2) is LineOK.coneI_eq: Steps 2 to 5 are its h2, h1, h3 to h6 and the final
% computation.

\begin{lemma}\label{asm-theta-agree-split-inputs-xunique}
Let $B$ be a theta base for a CM field $E$
(Definition~\ref{asm-theta-agree-models}(1)), let $X$ and $X_1$ be Hermitian and
$W$ a skew-Hermitian lattice space, of ranks $n$, $n_1$ and $d$ and of some
signatures, let $H=(x,y,\iota,\mathfrak b)$ be an adapted plane in $X$ over
$X_1$, and let $M$ and $M_1$ be regular models for $(X,W)$ and $(X_1,W)$, with
Schrödinger operators $\rho$ and $\rho_1$. Put
$\mathcal J_x(H)=\{xq^{\mathsf T}\}$, with $q$ running over the vectors over
$\mathbf A_E$ of length $d$. Let $e$ be an evaluation map for
$(M,M_1,\iota,1,\mathcal J_x(H))$ (Definition~\ref{asm-theta-agree-planes}(5)),
and let $e'$ be a linear map from functions on $\mathbf A^{n\times d}$ to
functions on $\mathbf A^{n_1\times d}$ such that, for every Schwartz-Bruhat
function $\Phi$, $e'(\rho(\iota z)\Phi)=\rho_1(z)e'(\Phi)$ for
$z\in\mathbb W(X_1,W)$ and $e'(\rho(u)\Phi)=e'(\Phi)$ for $u\in\mathcal J_x(H)$,
and such that $\Phi\mapsto e'(\Phi)(0)$ is admissible
(Theorem~\ref{asm-theta-agree-evaluation-heisenberg}). Then there is
$c\in\mathbf C$ with $e'(\Phi)=c\,e(\Phi)$ for every Schwartz-Bruhat function
$\Phi$.
\end{lemma}
\leanref{Development/WeilBound/ThetaLifts/LeviDeterminants/LagrangianUniqueness/GeoLeaves.lean}{PositivityTheory.ThKF0.hermitian_plane_theta_restriction}
\leanref{Development/WeilBound/ThetaLifts/LeviDeterminants/LagrangianUniqueness/GeoLeaves.lean}{PositivityTheory.ThKF0.evUniq}
\leanref{Development/WeilBound/ThetaLifts/LeviDeterminants/LagrangianUniqueness/GeoCore.lean}{PositivityTheory.ThKF0.Geo.hasLag_X}
\leanref{Development/WeilBound/ThetaLifts/LeviDeterminants/LagrangianUniqueness/GeoCore.lean}{PositivityTheory.ThKF0.Geo.map_line}
\begin{proof}
The proof of Lemma~\ref{asm-theta-agree-evaluation-unique} applies with the
roles of $X$ and $W$ exchanged; we give the changes. Put
$I=\{1,\ldots,n\}\times\{1,\ldots,d\}$ and $I_1=\{1,\ldots,n_1\}\times\{1,\ldots,d\}$,
let $e^M$ and $e^1$ be the isomorphisms of $M$ and $M_1$, $\rho^\circ$ as
there, and $\mathcal S$ the Schwartz-Bruhat functions on $\mathbf A^I$. By the
facts (a) and (b) of that proof, $e(\Phi_0)(y_0)\neq0$ for some
$\Phi_0\in\mathcal S$ and some $y_0$, and $\ell(\Phi)=e(\Phi)(0)$ is admissible;
$\ell'(\Phi)=e'(\Phi)(0)$ is admissible by hypothesis.

\emph{A rational Lagrangian.} Let $s(q)=x\,q^{\mathsf T}$ for $q\in E^d$,
$j(z)=\iota z$ for $z\in M_{n_1\times d}(\mathbf A_E)$, and let $L_1$ be the set
of $z\in M_{n_1\times d}(E)$ with $e^1_1(z)=0$, of dimension $n_1d$ over $F$.
Put $L=e^M(s(E^d)+j(L_1))\subset F^I\times F^I$. The form vanishes on it:
$h(s(q),s(q'))=0$ as $h_X(x,x)=0$; $h(j(z),s(q))=0$ as $h_X(\iota u,x)=0$; and
$\langle j(z),j(z')\rangle=\langle z,z'\rangle_1=0$ for $z,z'\in L_1$, as
$\iota^{\mathsf T}J_Xc(\iota)=J_{X_1}$. As $x\neq0$, $s$ is injective, so
$\dim_Fs(E^d)=2d$; $j$ preserves the forms, so it is injective, and
$s(E^d)\cap j(L_1)=0$ as in that proof. Hence $\dim_FL=2d+n_1d=nd$, and $L$ is
its own orthogonal. For $t\in\mathbf A$ and $w=e^M(s(q)+j(z_0))\in L$,
$\rho^\circ(tw)=\rho(ts(q))\rho(j(tz_0))$, with $ts(q)=x(tq)^{\mathsf T}\in\mathcal J_x(H)$
and $e^1_1(tz_0)=0$.

\emph{Uniqueness.} As in that proof, $\ell(\rho^\circ(tw)\Phi)=\ell(\Phi)$ and
$\ell'(\rho^\circ(tw)\Phi)=\ell'(\Phi)$ for $t\in\mathbf A$, $w\in L$ and
$\Phi\in\mathcal S$, and $\ell(\rho(j(z_{y_0}))\Phi_0)=e(\Phi_0)(y_0)\neq0$,
where $e^1(z_y)=(y,0)$. Lemma~\ref{asm-theta-agree-evaluation-lagrangian} gives
$c$ with $\ell'=c\,\ell$ on $\mathcal S$, and for $\Phi\in\mathcal S$ and every
$y$, $e'(\Phi)(y)=\ell'(\rho(j(z_y))\Phi)=c\,\ell(\rho(j(z_y))\Phi)=c\,e(\Phi)(y)$.
\end{proof}
% Lean: hermitian_plane_theta_restriction = evUniq (GeoLeaves.lean) with hasLag_X (GeoCore.lean),
% which is hasLagrangian_skew_adapted (Lemma asm-theta-agree-evaluation-unique) with s y = x y^T
% and iota on the left (hasLag_of, x_ne, orth_alt, adapted_plane_isotropic_vector), and
% Geo.map_line with lagLine, GeoAdm.comp_mem_adm and exists_mem_schA_ev_ne_zero.

\begin{lemma}\label{asm-theta-agree-split-inputs-xcommutator}
Let $B$ be a theta base for a CM field $E$, let $X$ and $X_1$ be Hermitian
lattice spaces of some signatures, and let $H=(x,y,\iota,\mathfrak b)$ be an
adapted plane in $X$ over $X_1$. Let $u\in U_X$ with $ux=x$ and
$u\iota z\in\iota z+\mathbf A_Ex$ for every vector $z$ over $\mathbf A_E$, and
let $m,m'\in U_X$ and an integer $t\geq1$ satisfy $mx=2x$, $2\,my=y$,
$m\iota=\iota$, $m'x=tx$, $t\,m'y=y$ and $m'\iota=\iota$. Put
$v=u^{-1}u^{-1}mum^{-1}$. Then $m'vm'^{-1}=v^{t^2}$.
\end{lemma}
\leanref{Development/WeilBound/ThetaLifts/LeviDeterminants/MixedModels/MixedModelPlanes/UnitaryFixingVectorPreserves.lean}{PositivityTheory.Mix.T.unitary_fixing_vector_preserves}
\leanref{Development/WeilBound/ThetaLifts/LeviDeterminants/MixedModels/UnitaryFixingIsotropicComplement.lean}{PositivityTheory.Mix.unitary_fixing_isotropic_complement}
\begin{proof}
Write $u\iota z=\iota z+a(z)x$ with $a(z)\in\mathbf A_E$. The vectors $x$, $y$
and the columns of $\iota$ form a basis (proof of
Lemma~\ref{asm-theta-agree-split-low}), so $uy=s_0x+r_0y+\iota z_0$ for some
$s_0,r_0\in\mathbf A_E$ and $z_0$; as $h_X(uy,ux)=h_X(y,x)=1$ and
$h_X(x,x)=h_X(\iota z_0,x)=0$, $r_0=1$. Hence $u^{-1}x=x$,
$u^{-1}\iota z=\iota z-a(z)x$ and $u^{-1}y=y-s_0x-\iota z_0+a(z_0)x$. Applying
$m^{-1}$, $u$, $m$, $u^{-1}$ and $u^{-1}$ in turn gives
\[
vx=x,\qquad v\iota z=\iota z,\qquad vy=y+c_0x\quad\text{with }c_0=2s_0-a(z_0).
\]
By induction, $v^kx=x$, $v^k\iota z=\iota z$ and $v^ky=y+kc_0x$. We compare
$m'v$ and $v^{t^2}m'$ on the basis. Both send $x$ to $tx$ and $\iota z$ to
$\iota z$. Further $t\,m'vy=t\,m'(y+c_0x)=y+t^2c_0x$ and
$t\,v^{t^2}m'y=v^{t^2}y=y+t^2c_0x$; as $t$ is invertible, $m'vy=v^{t^2}m'y$. So
$m'v=v^{t^2}m'$.
\end{proof}
% Lean: unitary_fixing_isotropic_complement (UnitaryFixingIsotropicComplement.lean): planeOK_herm,
% span and hpres for r0 = 1 (hr), hUix, hUiι and hUiy for u^{-1}, hAx, hAι and hAy for v with
% c0 = 2 s0 - a z0, then the comparison with the powers on x, the iota z and y.

\begin{definition}\label{asm-theta-agree-split-inputs-heis}
Let $B$ be a theta base for a CM field $E$
(Definition~\ref{asm-theta-agree-models}(1)), with maximal totally real
subfield $F$, $\mathbf A=\mathbf A_F$ and character $\psi$. For a finite set
$\kappa$ write $\mathcal S(\kappa)$ for the space of Schwartz-Bruhat functions
on $\mathbf A^\kappa$; call linear forms on the functions on $\mathbf A^\kappa$
admissible, and let $\rho(w)$, for $w=(w_1,w_2)\in\mathbf A^\kappa\times\mathbf A^\kappa$,
be the Schrödinger operators, as in
Theorem~\ref{asm-theta-agree-evaluation-heisenberg}. Put
$[w,w']=w_1\cdot w'_2-w'_1\cdot w_2$ and $\operatorname{ev}_0(\phi)=\phi(0)$.
\begin{enumerate}
\item A \emph{symplectic automorphism} of $\mathbf A^\kappa\times\mathbf A^\kappa$
  is an $\mathbf A$-linear bijection $S$ with $[Sw,Sw']=[w,w']$ for all $w,w'$.
  A linear map $T$ on the functions on $\mathbf A^\kappa$ \emph{implements} $S$
  if it maps $\mathcal S(\kappa)$ into itself and is injective on
  $\mathcal S(\kappa)$, $\lambda\circ T$ is admissible for every admissible
  $\lambda$, and $T(\rho(w)\phi)=\rho(Sw)T(\phi)$ for all $w$ and
  $\phi\in\mathcal S(\kappa)$.
\item Let $j\colon\kappa_1\to\kappa$ be injective. For a vector $a$ indexed by
  $\kappa_1$, let $j_!a$ be the vector indexed by $\kappa$ with entry $a_k$ at
  $j(k)$ and $0$ off $j(\kappa_1)$, and for a function $\phi$ on
  $\mathbf A^\kappa$ put $\operatorname{res}_j(\phi)(y)=\phi(j_!y)$. Let $S$
  be a symplectic automorphism. An element $w$ of
  $\mathbf A^\kappa\times\mathbf A^\kappa$ \emph{lies over} an element $w'$ of
  $\mathbf A^{\kappa_1}\times\mathbf A^{\kappa_1}$ for $(S,j)$ if
  $(Sw)_1=j_!w'_1$ and $(Sw)_2\circ j=w'_2$. A linear map $e$ from the
  functions on $\mathbf A^\kappa$ to those on $\mathbf A^{\kappa_1}$
  \emph{satisfies the law} of $(S,j)$ if $e(\rho(w)\phi)=\rho(w')e(\phi)$
  whenever $w$ lies over $w'$ and $\phi\in\mathcal S(\kappa)$. It is an
  \emph{$(S,j)$-evaluation} if it satisfies the law of $(S,j)$, maps
  $\mathcal S(\kappa)$ into $\mathcal S(\kappa_1)$, has $\lambda\circ e$
  admissible for every admissible $\lambda$, and has $e(\phi)\neq0$ for some
  $\phi\in\mathcal S(\kappa)$.
\item An admissible linear form $\ell$ on the functions on $\mathbf A^\kappa$ is
  \emph{cyclic} if for every nonzero $\phi\in\mathcal S(\kappa)$ there is $w$
  with $\ell(\rho(w)\phi)\neq0$.
\item For functions $\phi_1$ on $\mathbf A^\kappa$ and $\phi_2$ on
  $\mathbf A^{\kappa'}$, let $\phi_1\boxtimes\phi_2$ be the function
  $(x,x')\mapsto\phi_1(x)\phi_2(x')$ on $\mathbf A^{\kappa\sqcup\kappa'}$. For
  $w,w'\in\mathbf A^\kappa\times\mathbf A^\kappa$ put
  $w\boxminus w'=\bigl((w_1,w'_1),(w_2,-w'_2)\bigr)$, an element of
  $\mathbf A^{\kappa\sqcup\kappa}\times\mathbf A^{\kappa\sqcup\kappa}$. Let
  $S^\square$ be the map with $S^\square(w\boxminus w')=Sw\boxminus Sw'$, and
  $j\sqcup j\colon\kappa_1\sqcup\kappa_1\to\kappa\sqcup\kappa$ the map that is
  $j$ on each copy.
\end{enumerate}
\end{definition}
\leanref{Development/WeilBound/ThetaLifts/LeviDeterminants/CharacterAdmissibility/ResAdm.lean}{PositivityTheory.ThKF0.Chr.IsSp}
\leanref{Development/WeilBound/ThetaLifts/LeviDeterminants/CharacterAdmissibility/SchrodingerApplyZeroOfFst.lean}{PositivityTheory.ThKF0.Chr.IsImp}
\leanref{Development/WeilBound/ThetaLifts/LeviDeterminants/CharacterAdmissibility/ResAdm.lean}{PositivityTheory.ThKF0.Chr.Rel}
\leanref{Development/WeilBound/ThetaLifts/LeviDeterminants/CharacterAdmissibility/SchrodingerApplyZeroOfFst.lean}{PositivityTheory.ThKF0.Chr.Law}
\leanref{Development/WeilBound/ThetaLifts/LeviDeterminants/CharacterAdmissibility/SchrodingerApplyZeroOfFst.lean}{PositivityTheory.ThKF0.Chr.IsE}
\leanref{Development/WeilBound/ThetaLifts/LeviDeterminants/CharacterAdmissibility/SchrodingerApplyZeroOfFst.lean}{PositivityTheory.ThKF0.Chr.Cyc}
\leanref{Development/WeilBound/ThetaLifts/LeviDeterminants/CharacterAdmissibility/BoxCancel.lean}{PositivityTheory.ThKF0.Chr.dv}
\leanref{Development/WeilBound/ThetaLifts/LeviDeterminants/CharacterAdmissibility/BoxCancel.lean}{PositivityTheory.ThKF0.Chr.dbS}

\begin{lemma}\label{asm-theta-agree-split-inputs-heis-axioms}
Let $B$ be a theta base for a CM field $E$, with the notation of
Definition~\ref{asm-theta-agree-split-inputs-heis}, and let $\kappa$ and
$\kappa'$ be finite sets.
\begin{enumerate}
\item The operators $\rho(w)$ map $\mathcal S(\kappa)$ into itself;
  $\mathcal S(\kappa)\neq0$; $\operatorname{ev}_0$ is admissible; for an
  injective $j\colon\kappa_1\to\kappa$, $\operatorname{res}_j$ maps
  $\mathcal S(\kappa)$ into $\mathcal S(\kappa_1)$ and
  $\lambda\circ\operatorname{res}_j$ is admissible for every admissible
  $\lambda$; and for $\phi_1\in\mathcal S(\kappa)$ and
  $\phi_2\in\mathcal S(\kappa')$, $\phi_1\boxtimes\phi_2$ lies in
  $\mathcal S(\kappa\sqcup\kappa')$, and $\phi\mapsto\lambda(\phi\boxtimes\phi_2)$
  is admissible for every admissible $\lambda$.
\item If $\lambda$ is admissible and $\lambda(\rho(0,y)\phi)=\lambda(\phi)$ for
  all $y\in\mathbf A^\kappa$ and $\phi\in\mathcal S(\kappa)$, then there is
  $a\in\mathbf C$ with $\lambda=a\operatorname{ev}_0$ on $\mathcal S(\kappa)$.
\item Every symplectic automorphism of $\mathbf A^\kappa\times\mathbf A^\kappa$
  is implemented by some linear map.
\item There is an admissible linear form $I$ on the functions on
  $\mathbf A^\kappa$, not zero on $\mathcal S(\kappa)$, with
  $I(\phi(\cdot+a))=I(\phi)$ for $a\in\mathbf A^\kappa$ and
  $\phi\in\mathcal S(\kappa)$, and with
  $I(\phi\circ g^{-1})=|\det g|_{\mathbf A}\,I(\phi)$ for every
  $g\in GL_\kappa(\mathbf A)$ and every $\phi\in\mathcal S(\kappa)$ with
  $\phi\circ g^{-1}\in\mathcal S(\kappa)$.
\end{enumerate}
\end{lemma}
\leanref{Development/WeilBound/ThetaLifts/LeviDeterminants/CharacterAdmissibility/LawWLevi.lean}{PositivityTheory.ThKF0.Chr.ax_of}
\leanref{Development/WeilBound/ThetaLifts/LeviDeterminants/CharacterAdmissibility/AdelicCoordinatesEvalZeroFormMemAdmissible.lean}{PositivityTheory.ThKF0.Chr.axB}
\leanref{Development/WeilBound/ThetaLifts/LeviDeterminants/CharacterAdmissibility/AdelicCoordinatesEvalZeroFormMemAdmissible.lean}{PositivityTheory.ThKF0.Chr.axDelta_of}
\leanref{Development/WeilBound/ThetaLifts/LeviDeterminants/CharacterAdmissibility/AdelicCoordinatesEvalZeroFormMemAdmissible.lean}{PositivityTheory.ThKF0.Chr.axImp_of}
\leanref{Development/WeilBound/ThetaLifts/LeviDeterminants/CharacterAdmissibility/AdelInt.lean}{PositivityTheory.ThKF0.Chr.axInt}
\begin{proof}
(1) The first claim is Theorem~\ref{asm-theta-agree-evaluation-heisenberg}(1).
A product $f\otimes g$ of a nonzero Schwartz function on $F_\infty^\kappa$ and
a nonzero Schwartz-Bruhat function on $\mathbf A_f^\kappa$ is nonzero, and
$f\mapsto f(0)g(0)$ is continuous; so $\mathcal S(\kappa)\neq0$ and
$\operatorname{ev}_0$ is admissible. Restriction to the vectors $j_!y$ carries
$f\otimes g$ to $(f\circ j_!)\otimes(g\circ j_!)$, a Schwartz function times a
Schwartz-Bruhat function, and $f\mapsto f\circ j_!$ is continuous between the
Schwartz spaces; this gives the claims on $\operatorname{res}_j$. Likewise
$(f_1\otimes g_1)\boxtimes(f_2\otimes g_2)=(f_1\boxtimes f_2)\otimes(g_1\boxtimes g_2)$,
and $f_1\mapsto f_1\boxtimes f_2$ is continuous between the Schwartz spaces.

(2) The subspace $L=\{0\}\times F^\kappa$ of $F^\kappa\times F^\kappa$ is its own
orthogonal for $[\ ,\ ]$. For $t\in\mathbf A$ and $(0,y)\in L$,
$\lambda(\rho(t(0,y))\phi)=\lambda(\phi)$ by hypothesis, and
$\operatorname{ev}_0(\rho(0,ty)\phi)=\phi(0)$ by the formula for $\rho$. As
$\operatorname{ev}_0$ is admissible and not zero on $\mathcal S(\kappa)$,
Lemma~\ref{asm-theta-agree-evaluation-lagrangian} with $\lambda_1=\lambda$ and
$\lambda_2=\operatorname{ev}_0$ gives $a$.

(3) Let $S$ be a symplectic automorphism. The identity and $S$ are Darboux
coordinates on $\Lambda=\mathbf A^\kappa\times\mathbf A^\kappa$ for the form
$[\ ,\ ]$ in the sense of Lemma~\ref{asm-theta-agree-evaluation-transfer}, and
the character $\psi$ of $B$ satisfies the hypotheses of
Theorem~\ref{asm-theta-agree-evaluation-metaplectic}. By
Lemma~\ref{asm-theta-agree-evaluation-transfer}(1) there are a transfer $T$
from the identity to $S$ and a transfer $T'$ from $S$ to the identity, each
mapping $\mathcal S(\kappa)$ onto itself. So $T(\rho(w)\phi)=\rho(Sw)T(\phi)$
for $\phi\in\mathcal S(\kappa)$, and $T'\circ T$ maps $\mathcal S(\kappa)$ into
itself, keeps admissible forms admissible, and commutes with every $\rho(w)$ on
$\mathcal S(\kappa)$. By Theorem~\ref{asm-theta-agree-evaluation-heisenberg}(5),
$T'T=a$ on $\mathcal S(\kappa)$ for some $a\in\mathbf C$. Choose $\Psi$ with
$T'\Psi\neq0$ and $\phi\in\mathcal S(\kappa)$ with $T\phi=\Psi$; then
$a\phi=T'\Psi\neq0$. So $a\neq0$, $T$ is injective on $\mathcal S(\kappa)$, and
$T$ implements $S$.

(4) Let $\mu$ be a Haar measure on $\mathbf A^\kappa=F_\infty^\kappa\times\mathbf A_f^\kappa$,
the product of Haar measures $\mu_\infty$ and $\mu_f$, and let $I$ be a linear
form on the functions on $\mathbf A^\kappa$ with $I(\phi)=\int\phi\,d\mu$ for
$\phi\in\mathcal S(\kappa)$. By Fubini's theorem
$I(f\otimes g)=\int f\,d\mu_\infty\int g\,d\mu_f$, which is continuous in $f$;
so $I$ is admissible. If $f\geq0$ is a smooth bump function with $f(0)>0$ and
$g$ is the indicator function of $\prod_v\mathcal O_v^\kappa$, then
$I(f\otimes g)>0$. Haar measures are translation invariant. At each place an
invertible linear map multiplies the Haar measures by the absolute value of its
determinant \cite[Chapter~I, \S2]{Wei67}, and at almost all places $g$ and
$g^{-1}$ preserve $\mathcal O_v^\kappa$; so $\mu(gB)=|\det g|_{\mathbf A}\,\mu(B)$
for Borel sets $B$, and the substitution $x=gx'$ gives
$\int\phi(g^{-1}x)\,d\mu(x)=|\det g|_{\mathbf A}\int\phi\,d\mu$.
\end{proof}
% Lean: ax_of hA hpsi = { axB hpsi with delta := axDelta_of lagLine hpsi, imp := axImp_of hA hpsi }
% (LawWLevi.lean). (1) is axB: theta_distribution_heisenberg_invariant, exists_mem_schA_ne_zero,
% adelicCoordinatesEvalZeroForm_mem_admissible, RestrSchA.mem_schA_comp_extend and restrS,
% SchwartzBox.boxFun_mem_schA and box_adm (exists_boxCLM). (2) is axDelta_of with lagLine
% (Lemma asm-theta-agree-evaluation-lagrangian) for the Lagrangian ker fst. (3) is axImp_of through
% tw_bij (ExistsNeOfSurjOn.lean): tw_surj both ways and schur. (4) is axInt: Lext of the integral
% for the Haar measure mu on Vi x Vf (integral_prod_mul), a bump times the indicator of the
% integral finite adeles, integral_add_right_eq_self, and chi_eq for the modulus of mulVecCE.

\begin{lemma}\label{asm-theta-agree-split-inputs-heis-lines}
Let $B$ be a theta base for a CM field $E$, with the notation of
Definition~\ref{asm-theta-agree-split-inputs-heis}. Let $S$ be a symplectic
automorphism of $\mathbf A^\kappa\times\mathbf A^\kappa$ and
$j\colon\kappa_1\to\kappa$ an injective map.
\begin{enumerate}
\item If $T$ is a linear map on the functions on $\mathbf A^\kappa$ with
  $\operatorname{ev}_0\circ T$ admissible and $T(\rho(w)\phi)=\rho(w)T(\phi)$
  for all $w$ and $\phi\in\mathcal S(\kappa)$, then there is $a\in\mathbf C$
  with $T(\phi)=a\phi$ for $\phi\in\mathcal S(\kappa)$.
\item Let $T$ implement $S$, and let $\ell$ be admissible with
  $\ell(\rho(w)\phi)=\ell(\phi)$ whenever $(Sw)_1=0$ and
  $\phi\in\mathcal S(\kappa)$. Then there is $a\in\mathbf C$ with
  $\ell(\phi)=a\,T(\phi)(0)$ for $\phi\in\mathcal S(\kappa)$. If moreover
  $\ell$ is not zero on $\mathcal S(\kappa)$, then $\ell$ is cyclic.
\item There is a cyclic form $\ell$ with $\ell(\rho(w)\phi)=\ell(\phi)$ whenever
  $(Sw)_1=0$ and $\phi\in\mathcal S(\kappa)$.
\item If $T$ implements $S$, then $\operatorname{res}_j\circ T$ is an
  $(S,j)$-evaluation. In particular there is an $(S,j)$-evaluation.
\item Let $e$ be an $(S,j)$-evaluation. Then $\operatorname{ev}_0\circ e$ is
  cyclic. If $\ell$ is a cyclic form on the functions on $\mathbf A^{\kappa_1}$
  and $\phi\in\mathcal S(\kappa)$ is nonzero, then there are $w$ and $w'$ with
  $\ell(e(\rho(w')\rho(w)\phi))\neq0$; in particular $\ell(e(\phi))\neq0$ for
  some $\phi\in\mathcal S(\kappa)$.
\item Let $e$ be an $(S,j)$-evaluation, and let $e'$ be a linear map from the
  functions on $\mathbf A^\kappa$ to those on $\mathbf A^{\kappa_1}$ that
  satisfies the law of $(S,j)$, with $\operatorname{ev}_0\circ e'$ admissible.
  Then there is $c\in\mathbf C$ with $e'(\phi)=c\,e(\phi)$ for
  $\phi\in\mathcal S(\kappa)$.
\end{enumerate}
\end{lemma}
\leanref{Development/WeilBound/ThetaLifts/LeviDeterminants/CharacterAdmissibility/SchrodingerApplyZeroOfFst.lean}{PositivityTheory.ThKF0.Chr.Ax.schur}
\leanref{Development/WeilBound/ThetaLifts/LeviDeterminants/CharacterAdmissibility/SchrodingerApplyZeroOfFst.lean}{PositivityTheory.ThKF0.Chr.Ax.line}
\leanref{Development/WeilBound/ThetaLifts/LeviDeterminants/CharacterAdmissibility/SchrodingerApplyZeroOfFst.lean}{PositivityTheory.ThKF0.Chr.Ax.cyc}
\leanref{Development/WeilBound/ThetaLifts/LeviDeterminants/CharacterAdmissibility/SchrodingerApplyZeroOfFst.lean}{PositivityTheory.ThKF0.Chr.Ax.exists_form}
\leanref{Development/WeilBound/ThetaLifts/LeviDeterminants/CharacterAdmissibility/SchrodingerApplyZeroOfFst.lean}{PositivityTheory.ThKF0.Chr.Ax.isE_res}
\leanref{Development/WeilBound/ThetaLifts/LeviDeterminants/CharacterAdmissibility/SchrodingerApplyZeroOfFst.lean}{PositivityTheory.ThKF0.Chr.Ax.exists_E}
\leanref{Development/WeilBound/ThetaLifts/LeviDeterminants/CharacterAdmissibility/SchrodingerApplyZeroOfFst.lean}{PositivityTheory.ThKF0.Chr.Ax.E_cyc}
\leanref{Development/WeilBound/ThetaLifts/LeviDeterminants/CharacterAdmissibility/SchrodingerApplyZeroOfFst.lean}{PositivityTheory.ThKF0.Chr.Ax.comp_cyc}
\leanref{Development/WeilBound/ThetaLifts/LeviDeterminants/CharacterAdmissibility/SchrodingerApplyZeroOfFst.lean}{PositivityTheory.ThKF0.Chr.Ax.comp_ne}
\leanref{Development/WeilBound/ThetaLifts/LeviDeterminants/CharacterAdmissibility/SchrodingerApplyZeroOfFst.lean}{PositivityTheory.ThKF0.Chr.Ax.uniq}
\begin{proof}
By the formula for $\rho$, $(\rho(x,0)\phi)(0)=\phi(x)$, and
$(\rho(w)\phi)(0)=\phi(0)$ if $w_1=0$. For $y\in\mathbf A^{\kappa_1}$ put
$w_y=S^{-1}(j_!y,0)$; it lies over $(y,0)$. More generally
$S^{-1}(j_!w'_1,j_!w'_2)$ lies over $w'$, and if $(Sw)_1=0$, then $w$ lies over
$(0,(Sw)_2\circ j)$. So, if $e$ satisfies the law of $(S,j)$, then for
$\phi\in\mathcal S(\kappa)$:
\begin{enumerate}
\item[(L1)] $e(\phi)(y)=(\rho(y,0)e(\phi))(0)=e(\rho(w_y)\phi)(0)$ for every $y$;
\item[(L2)] $e(\rho(w)\phi)(0)=(\rho(0,(Sw)_2\circ j)e(\phi))(0)=e(\phi)(0)$
  whenever $(Sw)_1=0$.
\end{enumerate}

(1) The form $\operatorname{ev}_0\circ T$ is admissible, and
$T(\rho(0,y)\phi)(0)=(\rho(0,y)T\phi)(0)=T(\phi)(0)$. By
Lemma~\ref{asm-theta-agree-split-inputs-heis-axioms}(2),
$\operatorname{ev}_0\circ T=a\operatorname{ev}_0$ on $\mathcal S(\kappa)$. As
$\rho(x,0)\phi\in\mathcal S(\kappa)$,
$T(\phi)(x)=(\rho(x,0)T\phi)(0)=T(\rho(x,0)\phi)(0)=a\,(\rho(x,0)\phi)(0)=a\phi(x)$.

(2) By Lemma~\ref{asm-theta-agree-split-inputs-heis-axioms}(3) let $T'$
implement $S^{-1}$. The form $\ell\circ T'$ is admissible, and
$\ell(T'\rho(0,y)\phi)=\ell(\rho(S^{-1}(0,y))T'\phi)=\ell(T'\phi)$, as
$SS^{-1}(0,y)$ has first component $0$. By
Lemma~\ref{asm-theta-agree-split-inputs-heis-axioms}(2),
$\ell\circ T'=a'\operatorname{ev}_0$ on $\mathcal S(\kappa)$. The map $T'T$ maps
$\mathcal S(\kappa)$ into itself, commutes with every $\rho(w)$ there, as
$T'T\rho(w)=T'\rho(Sw)T=\rho(w)T'T$, and $\operatorname{ev}_0\circ T'\circ T$ is
admissible; by (1), $T'T=b$ on $\mathcal S(\kappa)$. If $\phi\in\mathcal S(\kappa)$
is nonzero, then $T'T\phi\neq0$ by the injectivity of $T$ and $T'$; so $b\neq0$,
and $\ell(\phi)=b^{-1}\ell(T'T\phi)=b^{-1}a'\,T(\phi)(0)$. Put $a=b^{-1}a'$.
If $\ell$ is not zero on $\mathcal S(\kappa)$, then $a\neq0$. Let $\phi$ be
nonzero; then $T\phi\neq0$, so $T(\phi)(x)\neq0$ for some $x$, and
$w=S^{-1}(x,0)$ gives
$\ell(\rho(w)\phi)=a\,(\rho(x,0)T\phi)(0)=a\,T(\phi)(x)\neq0$. So $\ell$ is
cyclic.

(3) Let $T$ implement $S$ and $\ell=\operatorname{ev}_0\circ T$. It is
admissible; if $(Sw)_1=0$, then
$\ell(\rho(w)\phi)=(\rho(Sw)T\phi)(0)=T(\phi)(0)$; and for nonzero $\phi$ the
$w$ of (2) gives $\ell(\rho(w)\phi)=T(\phi)(x)\neq0$.

(4) The map $e=\operatorname{res}_j\circ T$ maps $\mathcal S(\kappa)$ into
$\mathcal S(\kappa_1)$ and keeps admissible forms admissible, by
Lemma~\ref{asm-theta-agree-split-inputs-heis-axioms}(1). Let $w$ lie over $w'$
and $\phi\in\mathcal S(\kappa)$. For $x\in\mathbf A^{\kappa_1}$,
\[
e(\rho(w)\phi)(x)=(\rho(Sw)T\phi)(j_!x)
=\psi\bigl(j_!x\cdot(Sw)_2+\tfrac12(Sw)_1\cdot(Sw)_2\bigr)\,T(\phi)\bigl(j_!x+(Sw)_1\bigr).
\]
As $(Sw)_1=j_!w'_1$ and $(Sw)_2\circ j=w'_2$, $j_!x\cdot(Sw)_2=x\cdot w'_2$,
$(Sw)_1\cdot(Sw)_2=w'_1\cdot w'_2$ and $j_!x+(Sw)_1=j_!(x+w'_1)$; so this is
$(\rho(w')e(\phi))(x)$. For nonzero $\phi\in\mathcal S(\kappa)$
(Lemma~\ref{asm-theta-agree-split-inputs-heis-axioms}(1)), the $w$ of (2) gives
$e(\rho(w)\phi)(0)=T(\phi)(x)\neq0$. The last claim follows from
Lemma~\ref{asm-theta-agree-split-inputs-heis-axioms}(3).

(5) The form $\operatorname{ev}_0\circ e$ is admissible and, by (L2), invariant
under the $\rho(w)$ with $(Sw)_1=0$. Choose $\phi$ and $y$ with
$e(\phi)(y)\neq0$; by (L1), $e(\rho(w_y)\phi)(0)\neq0$. By (2) applied to any
$T$ that implements $S$, $\operatorname{ev}_0\circ e$ is cyclic. Let $\ell$ be
cyclic and $\phi\neq0$. There is $w$ with $e(\rho(w)\phi)(0)\neq0$, so the
function $e(\rho(w)\phi)\in\mathcal S(\kappa_1)$ is nonzero, and there is $v$
with $\ell(\rho(v)e(\rho(w)\phi))\neq0$. With $w'=S^{-1}(j_!v_1,j_!v_2)$, which
lies over $v$, $\rho(v)e(\rho(w)\phi)=e(\rho(w')\rho(w)\phi)$.

(6) Let $T$ implement $S$. The forms $\operatorname{ev}_0\circ e'$ and
$\operatorname{ev}_0\circ e$ are admissible and, by (L2), invariant under the
$\rho(w)$ with $(Sw)_1=0$. By (2) there are $a$ and $b$ with
$e'(\phi)(0)=a\,T(\phi)(0)$ and $e(\phi)(0)=b\,T(\phi)(0)$ on
$\mathcal S(\kappa)$, and $b\neq0$, as $\operatorname{ev}_0\circ e$ is cyclic by
(5) and $\mathcal S(\kappa)\neq0$. For $\phi\in\mathcal S(\kappa)$ and every $y$,
by (L1) for $e'$ and $e$,
\[
e'(\phi)(y)=e'(\rho(w_y)\phi)(0)=a\,T(\rho(w_y)\phi)(0)=\tfrac ab\,e(\rho(w_y)\phi)(0)=\tfrac ab\,e(\phi)(y).
\]
\end{proof}
% Lean: (L1) and (L2) are Law.apply and Law.inv (with rel_lift and schrodinger_translate_apply,
% schrodinger_apply_zero_of_fst). (1) Ax.schur; (2) Ax.line and, for cyclicity, IsImp.cyc and
% Ax.cyc; (3) Ax.exists_form; (4) IsImp.law_res, Ax.isE_res (extend_dot, extend_add,
% extend_zero') and Ax.exists_E; (5) Ax.E_cyc, Ax.comp_cyc and Ax.comp_ne; (6) Ax.uniq.

\begin{lemma}\label{asm-theta-agree-split-inputs-heis-double}
Let $B$ be a theta base for a CM field $E$, with the notation of
Definition~\ref{asm-theta-agree-split-inputs-heis}. Let $S$ be a symplectic
automorphism of $\mathbf A^\kappa\times\mathbf A^\kappa$,
$j\colon\kappa_1\to\kappa$ an injective map, $e$ an $(S,j)$-evaluation and $\mathcal E$
an $(S^\square,j\sqcup j)$-evaluation.
\begin{enumerate}
\item $S^\square$ is a symplectic automorphism; every element of
  $\mathbf A^{\kappa\sqcup\kappa}\times\mathbf A^{\kappa\sqcup\kappa}$ is
  $w\boxminus w'$ for some $w,w'$; $w\boxminus w'$ lies over
  $v\boxminus v'$ for $(S^\square,j\sqcup j)$ if and only if $w$ lies over
  $v$ and $w'$ over $v'$ for $(S,j)$; and
  $\rho(w\boxminus w')(\phi_1\boxtimes\phi_2)=\rho(w)\phi_1\boxtimes\rho(w'_1,-w'_2)\phi_2$.
\item For every $\phi_2\in\mathcal S(\kappa)$ there is $B\in\mathcal S(\kappa_1)$
  with $\mathcal E(\phi_1\boxtimes\phi_2)=e(\phi_1)\boxtimes B$ for all
  $\phi_1\in\mathcal S(\kappa)$.
\item If $\ell$ is a cyclic form on the functions on
  $\mathbf A^{\kappa_1\sqcup\kappa_1}$, there are $\phi_1,\phi_2\in\mathcal S(\kappa)$
  with $\ell(\mathcal E(\phi_1\boxtimes\phi_2))\neq0$.
\item For every finite set $\kappa_2$ there is a cyclic form $\ell$ on the
  functions on $\mathbf A^{\kappa_2\sqcup\kappa_2}$ with
  $\ell(\rho(w\boxminus w)\phi)=\ell(\phi)$ for all
  $w\in\mathbf A^{\kappa_2}\times\mathbf A^{\kappa_2}$ and
  $\phi\in\mathcal S(\kappa_2\sqcup\kappa_2)$.
\item Let $\ell$ be a cyclic form on the functions on
  $\mathbf A^{\kappa_1\sqcup\kappa_1}$ with $\ell(\rho(w\boxminus w)\phi)=\ell(\phi)$
  for all $w$ and $\phi$. Let $\kappa'$ be a finite set and
  $i\colon\kappa'\to\kappa$ an injective map whose image is the complement of
  $j(\kappa_1)$, and for $u\in\mathbf A^{\kappa'}$ put $v_u=S^{-1}(i_!u,0)$.
  Then there is an admissible form $L$ on the functions on
  $\mathbf A^{\kappa\sqcup\kappa}$, not zero on $\mathcal S(\kappa\sqcup\kappa)$,
  with $L(\rho(w\boxminus w)\phi)=L(\phi)$ for all $w$ and $\phi$, such that
  the following holds. Let $O$ be a linear map that maps
  $\mathcal S(\kappa\sqcup\kappa)$ into itself, $a\neq0$ a complex number and
  $g\in GL_{\kappa'}(\mathbf A)$, such that for all $u$ and
  $\phi\in\mathcal S(\kappa\sqcup\kappa)$,
  $O(\rho(v_u\boxminus v_u)\phi)=\rho(v_{gu}\boxminus v_{gu})O(\phi)$ and
  $\ell(\mathcal E(O\phi))=a\,\ell(\mathcal E(\phi))$. Then $L(O\phi)=a\,|\det g|_{\mathbf A}L(\phi)$
  for $\phi\in\mathcal S(\kappa\sqcup\kappa)$.
\end{enumerate}
\end{lemma}
\leanref{Development/WeilBound/ThetaLifts/LeviDeterminants/CharacterAdmissibility/BoxCancel.lean}{PositivityTheory.ThKF0.Chr.IsSp.dbS}
\leanref{Development/WeilBound/ThetaLifts/LeviDeterminants/CharacterAdmissibility/BoxCancel.lean}{PositivityTheory.ThKF0.Chr.rel_dv}
\leanref{Development/WeilBound/ThetaLifts/LeviDeterminants/CharacterAdmissibility/BoxCancel.lean}{PositivityTheory.ThKF0.Chr.schr_dv}
\leanref{Development/WeilBound/ThetaLifts/LeviDeterminants/CharacterAdmissibility/BoxCancel.lean}{PositivityTheory.ThKF0.Chr.Ax.box_E}
\leanref{Development/WeilBound/ThetaLifts/LeviDeterminants/CharacterAdmissibility/BoxCancel.lean}{PositivityTheory.ThKF0.Chr.Ax.box_ne}
\leanref{Development/WeilBound/ThetaLifts/LeviDeterminants/CharacterAdmissibility/BoxCancel.lean}{PositivityTheory.ThKF0.Chr.Ax.exists_diag}
\leanref{Development/WeilBound/ThetaLifts/LeviDeterminants/CharacterAdmissibility/BoxCancel.lean}{PositivityTheory.ThKF0.Chr.Ax.exists_avg}
\leanref{Development/WeilBound/ThetaLifts/LeviDeterminants/CharacterAdmissibility/BoxCancel.lean}{PositivityTheory.ThKF0.Chr.isSp_kE}
\begin{proof}
(1) The dot products split over the two copies of $\kappa$, so
$[w\boxminus w',u\boxminus u']=[w,u]-[w',u']$. Hence
$[S^\square(w\boxminus w'),S^\square(u\boxminus u')]=[Sw,Su]-[Sw',Su']=[w\boxminus w',u\boxminus u']$;
$S^\square$ is $\mathbf A$-linear and bijective, as $S$ is, and
$((a,a'),(b,b'))=(a,b)\boxminus(a',-b')$. The first component of
$S^\square(w\boxminus w')$ is $((Sw)_1,(Sw')_1)$ and the second composed with
$j\sqcup j$ is $((Sw)_2\circ j,-(Sw')_2\circ j)$; comparing them with
$(j_!v_1,j_!v'_1)$ and $(v_2,-v'_2)$ gives the equivalence. The
formula for $\rho(w\boxminus w')$ follows from the formula for $\rho$, as the
dot products split.

(2) Fix $\phi_2$. For $y\in\mathbf A^{\kappa_1}$ let $w=S^{-1}(j_!y,0)$, which lies
over $(y,0)$, put $\phi_2^y=\rho(w_1,-w_2)\phi_2\in\mathcal S(\kappa)$, and
define $e_y(\phi_1)(x)=\mathcal E(\phi_1\boxtimes\phi_2^y)(x,0)$. The map $e_y$ satisfies
the law of $(S,j)$: if $u$ lies over $u_1$, then $u\boxminus0$ lies over
$u_1\boxminus0$ by (1), $\rho(u\boxminus0)(\phi_1\boxtimes\phi_2^y)=\rho(u)\phi_1\boxtimes\phi_2^y$,
and so $\mathcal E(\rho(u)\phi_1\boxtimes\phi_2^y)=\rho(u_1\boxminus0)\mathcal E(\phi_1\boxtimes\phi_2^y)$,
whose value at $(x,0)$ is $(\rho(u_1)e_y(\phi_1))(x)$. The form
$\operatorname{ev}_0\circ e_y$ is admissible, as $\operatorname{ev}_0\circ\mathcal E$ is,
by Lemma~\ref{asm-theta-agree-split-inputs-heis-axioms}(1). By
Lemma~\ref{asm-theta-agree-split-inputs-heis-lines}(6) there is $c_y$ with
$e_y=c_ye$ on $\mathcal S(\kappa)$. As $0$ lies over $0$, $0\boxminus w$ lies over
$0\boxminus(y,0)$, and $\rho(0\boxminus w)(\phi_1\boxtimes\phi_2)=\phi_1\boxtimes\phi_2^y$;
so $\mathcal E(\phi_1\boxtimes\phi_2^y)=\rho(0\boxminus(y,0))\mathcal E(\phi_1\boxtimes\phi_2)$, whose
value at $(x,0)$ is $\mathcal E(\phi_1\boxtimes\phi_2)(x,y)$. Hence
$\mathcal E(\phi_1\boxtimes\phi_2)(x,y)=c_y\,e(\phi_1)(x)$, and $B(y)=c_y$ works. To see
$B\in\mathcal S(\kappa_1)$, choose $\phi$ and $x$ with $e(\phi)(x)\neq0$. Then
$B(y)=e(\phi)(x)^{-1}\mathcal E(\phi\boxtimes\phi_2)(x,y)$, which is
$e(\phi)(x)^{-1}$ times the restriction of $\rho(((x,0),0))\mathcal E(\phi\boxtimes\phi_2)$
to the vectors $(0,y)$; it lies in $\mathcal S(\kappa_1)$ by
Lemma~\ref{asm-theta-agree-split-inputs-heis-axioms}(1).

(3) Let $\phi\in\mathcal S(\kappa)$ be nonzero; then $\phi\boxtimes\phi$ is a
nonzero element of $\mathcal S(\kappa\sqcup\kappa)$. By
Lemma~\ref{asm-theta-agree-split-inputs-heis-lines}(5) for $\mathcal E$ there are $u$ and
$u'$ with $\ell(\mathcal E(\rho(u')\rho(u)(\phi\boxtimes\phi)))\neq0$. Write
$u=a\boxminus b$ and $u'=a'\boxminus b'$ by (1); then
$\rho(u')\rho(u)(\phi\boxtimes\phi)=\phi_1\boxtimes\phi_2$ with
$\phi_1=\rho(a')\rho(a)\phi$ and $\phi_2=\rho(b'_1,-b'_2)\rho(b_1,-b_2)\phi$.

(4) Let $K$ be the map with
$K(w\boxminus w')=\bigl((w_1-w'_1,w_2-w'_2),(w_2,-w'_1)\bigr)$. It is
$\mathbf A$-linear and bijective, and a direct expansion gives
$[K(w\boxminus w'),K(u\boxminus u')]=[w,u]-[w',u']$, so it is symplectic; and
$K(w\boxminus w)$ has first component $0$. By
Lemma~\ref{asm-theta-agree-split-inputs-heis-lines}(3) for $K$ there is a cyclic
$\ell$ invariant under the $\rho(z)$ with $(Kz)_1=0$, in particular under the
$\rho(w\boxminus w)$.

(5) For $k\in\kappa$ write $k\in j$ if $k\in j(\kappa_1)$. Let $K'$ be the map
with $K'(w\boxminus w')=((w_1-w'_1,s),(w_2,s'))$, where $s_k=w_{2,k}-w'_{2,k}$
and $s'_k=-w'_{1,k}$ if $k\in j$, and $s_k=w'_{1,k}$ and
$s'_k=w_{2,k}-w'_{2,k}$ otherwise. It is $\mathbf A$-linear, bijective and
symplectic, by the computation of (4) at the $k\in j$ and a similar one at the
other $k$. Put $K=K'\circ S^\square$, a symplectic automorphism, let $T$
implement $K$ (Lemma~\ref{asm-theta-agree-split-inputs-heis-axioms}(3)), let
$i_2\colon\kappa'\to\kappa\sqcup\kappa$ be $i$ followed by the inclusion of the
second copy, and put $\mathcal E'=\operatorname{res}_{i_2}\circ T$, a
$(K,i_2)$-evaluation by Lemma~\ref{asm-theta-agree-split-inputs-heis-lines}(4).

(a) For every $w$, $w\boxminus w$ lies over $((Sw)_1\circ i,0)$ for $(K,i_2)$.
Indeed $K(w\boxminus w)=((0,s),((Sw)_2,s'))$ with $s_k=0$ and $s'_k=-(Sw)_{1,k}$
for $k\in j$, and $s_k=(Sw)_{1,k}$ and $s'_k=0$ otherwise. As the image of $i$
is the complement of $j(\kappa_1)$, $(0,s)=(i_2)_!((Sw)_1\circ i)$ and
$s'\circ i=0$.

(b) The form $\ell\circ\mathcal E$ is admissible, and it is invariant under the
$\rho(z)$ with $(Kz)_1=0$. Indeed, write $z=w\boxminus w'$; then $(Kz)_1=0$ says
$(Sw)_1=(Sw')_1$, $(Sw)_{2,k}=(Sw')_{2,k}$ for $k\in j$, and $(Sw')_{1,k}=0$ for
$k\notin j$. So $w$ and $w'$ both lie over
$v=((Sw')_1\circ j,(Sw')_2\circ j)$, $z$ lies over $v\boxminus v$ by (1), and
$\ell(\mathcal E(\rho(z)\phi))=\ell(\rho(v\boxminus v)\mathcal E(\phi))=\ell(\mathcal E(\phi))$. By
Lemma~\ref{asm-theta-agree-split-inputs-heis-lines}(2) there is $c$ with
$\ell(\mathcal E(\phi))=c\,T(\phi)(0)$ on $\mathcal S(\kappa\sqcup\kappa)$, and $c\neq0$,
as $\ell\circ\mathcal E$ is not zero on $\mathcal S(\kappa\sqcup\kappa)$ by
Lemma~\ref{asm-theta-agree-split-inputs-heis-lines}(5).

(c) For $u\in\mathbf A^{\kappa'}$, $(Sv_u)_1\circ i=u$, so by (a) and the law of
$(K,i_2)$, $\mathcal E'(\rho(v_u\boxminus v_u)\phi)=\rho(u,0)\mathcal E'(\phi)$, whose value at $0$
is $\mathcal E'(\phi)(u)$. As $\mathcal E'(\psi)(0)=T(\psi)(0)$, this gives
$\mathcal E'(\phi)(u)=T(\rho(v_u\boxminus v_u)\phi)(0)$, and by (b)
\[
c\,\mathcal E'(\phi)(u)=\ell\bigl(\mathcal E(\rho(v_u\boxminus v_u)\phi)\bigr)\qquad(\phi\in\mathcal S(\kappa\sqcup\kappa)).\tag{F}
\]

(d) Let $I$ be the form of
Lemma~\ref{asm-theta-agree-split-inputs-heis-axioms}(4) for $\kappa'$. It is
invariant under the $\rho(w)$ with $w_2=0$, which are the translations, and these
are the $w$ with $(S'w)_1=0$ for the symplectic automorphism $S'(w)=(w_2,-w_1)$;
as $I$ is not zero on $\mathcal S(\kappa')$, it is cyclic by
Lemma~\ref{asm-theta-agree-split-inputs-heis-lines}(2). Put $L=I\circ\mathcal E'$. It is
admissible; it is not zero on $\mathcal S(\kappa\sqcup\kappa)$ by
Lemma~\ref{asm-theta-agree-split-inputs-heis-lines}(5) for $\mathcal E'$ and $I$; and by
(a), $\mathcal E'(\rho(w\boxminus w)\phi)=\rho((Sw)_1\circ i,0)\mathcal E'(\phi)$ is a translate of
$\mathcal E'(\phi)$, so $L(\rho(w\boxminus w)\phi)=L(\phi)$.

(e) Let $O$, $a$ and $g$ be as in (5), and $\phi\in\mathcal S(\kappa\sqcup\kappa)$.
By (F), the hypothesis on $O$ with $g^{-1}u$ in place of $u$, the hypothesis on
$\ell$, and (F) again,
\[
c\,\mathcal E'(O\phi)(u)=\ell\bigl(\mathcal E(O\rho(v_{g^{-1}u}\boxminus v_{g^{-1}u})\phi)\bigr)
=a\,\ell\bigl(\mathcal E(\rho(v_{g^{-1}u}\boxminus v_{g^{-1}u})\phi)\bigr)=a\,c\,\mathcal E'(\phi)(g^{-1}u).
\]
So $\mathcal E'(O\phi)=a\,\mathcal E'(\phi)\circ g^{-1}$; both $\mathcal E'(\phi)$ and
$\mathcal E'(\phi)\circ g^{-1}=a^{-1}\mathcal E'(O\phi)$ lie in $\mathcal S(\kappa')$, and
Lemma~\ref{asm-theta-agree-split-inputs-heis-axioms}(4) gives
$L(O\phi)=a\,I(\mathcal E'(\phi)\circ g^{-1})=a\,|\det g|_{\mathbf A}L(\phi)$.
\end{proof}
% Lean: (1) is dE, exists_dv, dbS_dv, IsSp.dbS, rel_dv and schr_box/schr_dv. (2) is Ax.box_E
% (e_y is res inl . E . boxL with schr (w.1,-w.2) phi2, made proportional to e by Ax.uniq; the
% coefficient function A is then a restriction of a translate). (3) is Ax.box_ne. (4) is
% Ax.exists_diag with kE (fun _ => True) and isSp_kE. (5) is Ax.exists_avg: K = (dbS S).trans
% (kE (fun k => exists a, j a = k)), E' = res (inr . i) . T, hrel for (a), Ax.line for (b) (with
% hc0 from Ax.comp_ne), hF for (F), I from AxInt with Ax.cyc for isSp_fE, and the final
% computation e1 with hIg.

\begin{lemma}\label{asm-theta-agree-split-inputs-heis-scalars}
Let $B$ be a theta base for a CM field $E$, with the notation of
Definition~\ref{asm-theta-agree-split-inputs-heis}. Let $S$ be a symplectic
automorphism of $\mathbf A^\kappa\times\mathbf A^\kappa$,
$j\colon\kappa_1\to\kappa$ an injective map and $e$ an $(S,j)$-evaluation.
\begin{enumerate}
\item Let $T$ be a linear map on the functions on $\mathbf A^\kappa$ that maps
  $\mathcal S(\kappa)$ into itself and keeps admissible forms admissible, and
  let $T_1$ and $T'_1$ be linear maps on the functions on $\mathbf A^{\kappa_1}$
  with $T_1T'_1$ the identity and $\lambda\circ T'_1$ admissible for every
  admissible $\lambda$. Suppose that for every $w$ lying over $v$ there are
  $w'$ lying over $v'$ with $T(\rho(w)\phi)=\rho(w')T(\phi)$ for
  $\phi\in\mathcal S(\kappa)$ and $T'_1(\rho(v')\psi)=\rho(v)T'_1(\psi)$ for
  $\psi\in\mathcal S(\kappa_1)$. Then there is $a\in\mathbf C$ with
  $e(T\phi)=a\,T_1(e(\phi))$ for $\phi\in\mathcal S(\kappa)$.
\item Let $G$ be a topological group, and $\omega$ and $\omega_1$ representations
  of $G$ on the functions on $\mathbf A^\kappa$ and on $\mathbf A^{\kappa_1}$, such
  that every $\omega(g)$ maps $\mathcal S(\kappa)$ into itself, every
  $\omega(g)$ and $\omega_1(g)$ keeps admissible forms admissible, the functions
  $g\mapsto\lambda(\omega(g)\phi)$ and $g\mapsto\lambda(\omega_1(g)\psi)$ are
  continuous for admissible $\lambda$, $\phi\in\mathcal S(\kappa)$ and
  $\psi\in\mathcal S(\kappa_1)$, and (1) applies to $T=\omega(g)$,
  $T_1=\omega_1(g)$ and $T'_1=\omega_1(g^{-1})$ for every $g$. Then there is a
  continuous homomorphism $A\colon G\to\mathbf C^\times$ with
  $e(\omega(g)\phi)=A(g)\,\omega_1(g)e(\phi)$ for $g\in G$ and
  $\phi\in\mathcal S(\kappa)$.
\end{enumerate}
\end{lemma}
\leanref{Development/WeilBound/ThetaLifts/LeviDeterminants/CharacterAdmissibility/ExistsDoubledUnitaryBlock.lean}{PositivityTheory.ThKF0.Chr.lev_scalar}
\leanref{Development/WeilBound/ThetaLifts/LeviDeterminants/CharacterAdmissibility/LevScalar2.lean}{PositivityTheory.ThKF0.Chr.lev_scalar2}
\leanref{Development/WeilBound/ThetaLifts/LeviDeterminants/CharacterAdmissibility/DetConjMulVec.lean}{PositivityTheory.ThKF0.Chr.Ax.exists_char}
\begin{proof}
(1) Put $e'=T'_1\circ e\circ T$. If $w$ lies over $v$, with $w'$ and $v'$ as in
the hypothesis, then for $\phi\in\mathcal S(\kappa)$, as $T\phi\in\mathcal S(\kappa)$
and $e(T\phi)\in\mathcal S(\kappa_1)$,
$e'(\rho(w)\phi)=T'_1(e(\rho(w')T\phi))=T'_1(\rho(v')e(T\phi))=\rho(v)e'(\phi)$.
So $e'$ satisfies the law of $(S,j)$, and
$\operatorname{ev}_0\circ e'=((\operatorname{ev}_0\circ T'_1)\circ e)\circ T$ is
admissible. By Lemma~\ref{asm-theta-agree-split-inputs-heis-lines}(6),
$e'=a\,e$ on $\mathcal S(\kappa)$, and applying $T_1$ gives the claim.

(2) For each $g$, (1) gives $b(g)$ with $e(\omega(g)\phi)=b(g)\,\omega_1(g)e(\phi)$.
By Lemma~\ref{asm-theta-agree-split-inputs-heis-lines}(5) choose
$\phi_0\in\mathcal S(\kappa)$ with $e(\phi_0)(0)\neq0$. The number $b(g)$ is
determined by $g$: if $x\,\omega_1(g)e(\phi_0)=b(g)\,\omega_1(g)e(\phi_0)$, apply
$\omega_1(g^{-1})$ and evaluate at $0$. Hence $b(1)=1$, and
$e(\omega(gg')\phi)=b(g)\,\omega_1(g)e(\omega(g')\phi)=b(g)b(g')\,\omega_1(gg')e(\phi)$
gives $b(gg')=b(g)b(g')$; so $b$ is a homomorphism into $\mathbf C^\times$. Let
$g_0\in G$ and $\phi=\omega(g_0^{-1})\phi_0\in\mathcal S(\kappa)$. Then
$b(g_0)\,(\omega_1(g_0)e(\phi))(0)=e(\phi_0)(0)\neq0$. The functions
$g\mapsto e(\omega(g)\phi)(0)$ and $g\mapsto(\omega_1(g)e(\phi))(0)$ are
continuous, the second is nonzero near $g_0$, and $b$ is their quotient there.
So $A=b$ is continuous.
\end{proof}
% Lean: (1) is the common core of lev_scalar and lev_scalar2 (and of the first step of
% Ax.exists_char): e' = om1 g^-1 . e . om m, made proportional to e by Ax.uniq. (2) is
% Ax.exists_char: key, uq, the homomorphism B and its continuity through hs.cont (c1, c2) at
% phi = om g0^-1 phi0.

\begin{definition}\label{asm-theta-agree-split-inputs-side}
Let $B$ be a theta base for a CM field $E$, with the notation of
Definition~\ref{asm-theta-agree-split-inputs-heis}, and let $A$ and $Y$ be
lattice spaces (Definition~\ref{asm-theta-agree-doubling}(1)) of ranks $a$ and
$b$ with matrices $J_A$ and $J_Y$, let $\kappa$ be a finite set, the
\emph{index set}, give $M_{a\times b}(\mathbf A_E)$ the form
$\langle u,u'\rangle=-\operatorname{Tr}_{\mathbf A_E/\mathbf A}\sum_{i,j,k,l}u_{ij}(J_A)_{ik}(J_Y)_{jl}\,c(u'_{kl})$,
and let $G_Y$, $P_Y$, $\Delta$ and $t(p)$ be as in
Definition~\ref{asm-theta-agree-doubling}(2). A \emph{Weil side}
$(A,Y,d,\chi,\omega)$ consists of an $\mathbf A$-linear isomorphism $d$ from
$M_{a\times b}(\mathbf A_E)$ onto $\mathbf A^\kappa\times\mathbf A^\kappa$ with
$\langle u,u'\rangle=[d(u),d(u')]$, a Hecke character $\chi$ of $E$, a
representation $\omega$ of $U_Y$ on the functions on $\mathbf A^\kappa$ and a
representation $\omega^\square$ of $G_Y$ on the functions on
$\mathbf A^{\kappa\sqcup\kappa}$ such that, with
$d^\square(u_1\ u_2)=d(u_1)\boxminus d(u_2)$ for
$u_1,u_2\in M_{a\times b}(\mathbf A_E)$:
\begin{enumerate}
\item[(a)] the operators $\omega(h)$ map $\mathcal S(\kappa)$ into itself and
  keep admissible forms admissible, $\omega(h)\rho(d(u))\phi=\rho(d(uh^{\mathsf T}))\omega(h)\phi$
  for $u\in M_{a\times b}(\mathbf A_E)$ and $\phi\in\mathcal S(\kappa)$, and
  $h\mapsto\lambda(\omega(h)\phi)$ is continuous for admissible $\lambda$ and
  $\phi\in\mathcal S(\kappa)$;
\item[(b)] the operators $\omega^\square(x)$, $x\in G_Y$, satisfy (a) with
  $d^\square$ in place of $d$, and for every admissible $\lambda$ with
  $\lambda(\rho(d^\square(z\Delta^{\mathsf T}))\phi)=\lambda(\phi)$ for all
  $z\in M_{a\times b}(\mathbf A_E)$ and $\phi\in\mathcal S(\kappa\sqcup\kappa)$,
  \[
  \lambda(\omega^\square(p)\phi)=\chi(\det t(p))\,|\det t(p)|^{a/2}\lambda(\phi)
  \qquad(p\in P_Y,\ \phi\in\mathcal S(\kappa\sqcup\kappa));
  \]
\item[(c)] $\omega^\square(\operatorname{diag}(h,1))(\phi_1\boxtimes\phi_2)=(\omega(h)\phi_1)\boxtimes\phi_2$
  for $h\in U_Y$ and $\phi_1,\phi_2\in\mathcal S(\kappa)$.
\end{enumerate}
Let $(X,W,M,\chi_X,\chi_W,\omega)$ be a Weil pair
(Definition~\ref{asm-theta-agree-triple}(1)) and $e^M$ the isomorphism of $M$,
and put $\tau(u)=u^{\mathsf T}$. By (a) and (b) of
Definition~\ref{asm-theta-agree-models}(6), with $\omega^\square$ the operators
there, $(X,W,e^M,\chi_X,\omega(1,\cdot))$ and
$(W,X,e^M\circ\tau,\chi_W,\omega(\cdot,1))$ are Weil sides, both with index set
$\{1,\ldots,n\}\times\{1,\ldots,d\}$ for the ranks $n$ and $d$ of $X$ and $W$;
the action
$u\mapsto gu$ of $g\in U_X$ on $\mathbb W(X,W)$ is $U\mapsto Ug^{\mathsf T}$ on
$U=u^{\mathsf T}$.
\end{definition}
\leanref{Development/AbelianVarietyTheory/Definitions/Divisors/HilbertNinetyQuotientHomEqIff.lean}{GlobalTheory.WeilM.IsUWeilB}
\leanref{Development/AbelianVarietyTheory/Definitions/Divisors/HilbertNinetyQuotientHomEqIff.lean}{GlobalTheory.WeilM.IsDblCoord}
\leanref{Development/AbelianVarietyTheory/Definitions/Divisors/HilbertNinetyQuotientHomEqIff.lean}{GlobalTheory.WeilM.SiegelRule}

\begin{situation}\label{asm-theta-agree-split-inputs-siegel-situation}
Let $E$ be a commutative ring with a map $c\colon E\to E$, and let $R$ be a ring
with the structure of an $E$-algebra. Let $\sigma\colon R\to R$ be an additive
map with $\sigma(1)=1$, let $G$ be a group with a homomorphism of monoids
$\varphi\colon G\to R$, let $\lambda\colon G\to\mathbf C^\times$ be a
homomorphism, and let $e,f\in R$ satisfy $e+f=1$, $e^2=e$, $f^2=f$ and
$ef=fe=0$.
\end{situation}

\begin{lemma}\label{asm-theta-agree-split-inputs-siegel-unipotent}
In Situation~\ref{asm-theta-agree-split-inputs-siegel-situation}, suppose that
$\sigma(ax)=c(a)\sigma(x)$ for $a\in E$ and $x\in R$, that $\varphi$ is
injective and its image contains every $x\in R$ with $\sigma(x)x=1$, that
$\sigma(e)=f$ and $\sigma(f)=e$, and that $2$ and $3$ are invertible in $E$ with
$c(2)=2$, $c(1/2)=1/2$ and $c(1/3)=1/3$. Let $s\in R$ satisfy $es=0$, $sf=0$ and
$\sigma(s)=-s$. Then there is $n\in G$ with $\varphi(n)=1+s$ and $\lambda(n)=1$.
\end{lemma}
\leanref{Development/AutomorphicSpectrum/Proofs/MembershipProofs/SplitQuadraticAlgebraOrder.lean}{GlobalTheory.CompanionTheory.SiegelChar.lower_triv}
\begin{proof}
Let $t\in R$ satisfy $et=0$ and $tf=0$. Then $ft=(e+f)t=t$ and $te=t(e+f)=t$,
so $t^2=(te)(ft)=t(ef)t=0$. Hence $(1+t)^k=1+kt$ for every $k\ge0$, by induction
on $k$. If moreover $\sigma(t)=-t$, then $\sigma(1+t)(1+t)=(1-t)(1+t)=1-t^2=1$,
so $1+t=\varphi(n)$ for some $n\in G$.

Let $t$ be as above with $\sigma(t)=-t$, and let $n\in G$ satisfy
$\varphi(n)=1+t$. Then $\lambda(n)^3=1$. Indeed, put $d=2e+\frac12f$ and
$d'=\frac12e+2f$. As $e$ and $f$ are orthogonal idempotents with sum $1$,
$dd'=d'd=e+f=1$, and $\sigma(d)d=(2f+\frac12e)(2e+\frac12f)=e+f=1$ by
$\sigma(e)=f$, $\sigma(f)=e$, $c(2)=2$ and $c(1/2)=1/2$. So $d=\varphi(m)$ for
some $m\in G$, and $\varphi(m^{-1})=d'd\,\varphi(m^{-1})=d'$. From $et=0$ and
$ft=t$ we get $dt=\frac12t$, and from $tf=0$ and $te=t$ we get
$td'=\frac12t$. Hence $\varphi(mnm^{-1})=d(1+t)d'=1+\frac14t$, and as
$(\frac14t)^2=0$, $\varphi((mnm^{-1})^4)=(1+\frac14t)^4=1+t=\varphi(n)$. So
$(mnm^{-1})^4=n$, as $\varphi$ is injective. As $\mathbf C^\times$ is
commutative, $\lambda(mnm^{-1})=\lambda(n)$, so $\lambda(n)^4=\lambda(n)$ and
$\lambda(n)^3=1$.

Now let $s$ be as in the statement. By the first paragraph there is $n\in G$
with $\varphi(n)=1+s$. The element $\frac13s$ satisfies $e(\frac13s)=0$,
$(\frac13s)f=0$ and $\sigma(\frac13s)=c(1/3)\sigma(s)=-\frac13s$, so there is
$n'\in G$ with $\varphi(n')=1+\frac13s$. Then
$\varphi(n'^3)=(1+\frac13s)^3=1+s=\varphi(n)$, so $n'^3=n$. By the second
paragraph, applied to $\frac13s$ and $n'$, $\lambda(n)=\lambda(n')^3=1$.
\end{proof}
% Lean: lower_triv: hsfe, htee, hsq (t^2 = 0), hpow, hunit; key is the second paragraph (d, d',
% td from hsurj, hm, hm4 with hinj, hlm); the last paragraph is n, n' and hn3.

\begin{lemma}\label{asm-theta-agree-split-inputs-siegel-cell}
In Situation~\ref{asm-theta-agree-split-inputs-siegel-situation}, suppose that
$\sigma(xy)=\sigma(y)\sigma(x)$ for $x,y\in R$, that
$\sigma(\varphi(g))\varphi(g)=1$ for $g\in G$, and that $\sigma(e)=f$ and
$\sigma(f)=e$. Suppose also:
\begin{enumerate}
\item for every $s\in R$ with $es=0$, $sf=0$ and $\sigma(s)=-s$ there is $n\in G$
  with $\varphi(n)=1+s$ and $\lambda(n)=1$;
\item $\lambda(g)=1$ for every $g\in G$ with $f\varphi(g)e=0$.
\end{enumerate}
Then $\lambda(g)=1$ for every $g\in G$ such that $\varphi(g)e+f$ is a unit of
$R$.
\end{lemma}
\leanref{Development/AutomorphicSpectrum/Proofs/MembershipProofs/SplitQuadraticAlgebraOrder.lean}{GlobalTheory.CompanionTheory.SiegelChar.cell_triv}
\begin{proof}
Let $x=\varphi(g)$, so that $\sigma(x)x=1$, and let $w$ be the inverse of
$xe+f$. From $(xe+f)f=f$ and $(xe+f)e=xe$ we get $wf=f$ and $wxe=e$. As
$xe=exe+fxe$, this gives $wfxe=fxe$ and $wexe=e-fxe$. Put $s=-fwe$. Then $es=0$,
$sf=0$, $se=s$ and $fs=s$, and $s\cdot exe=-fw\cdot exe=-f(e-fxe)=fxe$.
Multiplying $(xe+f)w=1$ by $e$ on both sides gives $exe\cdot ewe=e$, so
$s=se=s\cdot exe\cdot ewe=fxe\cdot ewe$.

As $\sigma$ reverses products, $\sigma(fxe)=f\sigma(x)e$ and
$\sigma(exe)=f\sigma(x)f$, so
\[
 \sigma(fxe)\,exe+\sigma(exe)\,fxe=f\sigma(x)(e+f)xe=f\sigma(x)xe=fe=0.
\]
Hence
$\sigma(fxe)=\sigma(fxe)e=\sigma(fxe)\,exe\cdot ewe=-\sigma(exe)\,fxe\cdot ewe
=-\sigma(exe)s$, and
\[
 \sigma(s)=\sigma(ewe)\sigma(fxe)=-\sigma(ewe)\sigma(exe)s
 =-\sigma(exe\cdot ewe)s=-\sigma(e)s=-fs=-s.
\]
By (1) there is $n\in G$ with $\varphi(n)=1+s$ and $\lambda(n)=1$. As
$s^2=(se)(fs)=s(ef)s=0$, $(1-s)(1+s)=1$ and $\varphi(n^{-1})=1-s$. Since $sf=0$,
$sxe=s\cdot exe+s\cdot fxe=fxe$, so, using $fs=s$,
$f\varphi(n^{-1}g)e=f(1-s)xe=fxe-sxe=0$. By (2), $\lambda(n^{-1}g)=1$, and
$\lambda(g)=\lambda(n)\lambda(n^{-1}g)=1$.
\end{proof}
% Lean: cell_triv: hwf, hwxe, hxe, hwfxe, hweaxe; s = -(f w e) with hsa, haa', hs_eq; hrel,
% hσc2 and hσs; then hlow, hninv, hsxe, hpar and hP.

\begin{lemma}\label{asm-theta-agree-split-inputs-siegel-finite}
Let $F$ be a field, $E$ a commutative Artinian $F$-algebra, $\beta$ a finite set
and $A,B,D\in M_\beta(E)$. If $A+\tau B+\tau^2D$ is invertible for some
$\tau\in E$, then $A+tB+t^2D$ is invertible for all but finitely many $t\in F$.
\end{lemma}
\leanref{Development/AutomorphicSpectrum/Proofs/DivisibilityProofs/FiniteNonunitMatrix.lean}{GlobalTheory.CompanionTheory.Ar.RT.Siegel.finite_nonunit_matrix}
\leanref{Development/AutomorphicSpectrum/Proofs/DivisibilityProofs/FiniteNonunitMatrix.lean}{GlobalTheory.CompanionTheory.Ar.RT.Siegel.finite_nonunit}
\begin{proof}
A square matrix over $E$ is invertible exactly when its determinant is a unit.
Let $p=\det(A+XB+X^2D)\in E[X]$; then $p(z)=\det(A+zB+z^2D)$ for $z\in E$, and
$p(\tau)$ is a unit. The ring $E$ has finitely many maximal ideals. For each of
them, $\mathfrak m$, the image of $p$ in $(E/\mathfrak m)[X]$ is nonzero, since
its value at the class of $\tau$ is nonzero; so it has finitely many roots, and
as $F\to E/\mathfrak m$ is injective, $p(t)\in\mathfrak m$ for only finitely
many $t\in F$. A $p(t)$ that is not a unit lies in some maximal ideal.
\end{proof}

\begin{lemma}\label{asm-theta-agree-split-inputs-siegel-generic}
In Situation~\ref{asm-theta-agree-split-inputs-siegel-situation}, suppose that
$E$ is an Artinian algebra over an infinite field $F$, that $c$ is an
endomorphism of $F$-algebras, and that $R=M_\beta(E)$ for a finite set $\beta$,
with $E$ acting by scalar multiplication. Suppose that $\sigma(ax)=c(a)\sigma(x)$
and $\sigma(xy)=\sigma(y)\sigma(x)$ for $a\in E$ and $x,y\in R$, that the image
of $\varphi$ is the set of $x\in R$ with $\sigma(x)x=1$, that $2$ is a unit of
$E$, and that $E$ has a unit $\xi$ with $c(\xi)=-\xi$. If $\lambda(g)=1$ for
every $g\in G$ such that $\varphi(g)e+f$ is invertible, then $\lambda$ is
trivial.
\end{lemma}
\leanref{Development/AutomorphicSpectrum/Proofs/DivisibilityProofs/FiniteNonunitMatrix.lean}{GlobalTheory.CompanionTheory.Ar.RT.Siegel.triv_of_cell}
\begin{proof}
Let $g\in G$, $h=\varphi(g)$ and $k=\varphi(g^{-1})$. Then $hk=kh=1$,
$\sigma(h)=\sigma(h)hk=k$ and likewise $\sigma(k)=h$. For units $a,b$ of $E$,
$ae+bf$ is invertible, with inverse $a^{-1}e+b^{-1}f$. For $t\in E$ put
\[
 y(t)=(t^2+t\xi)h+(t-\xi),\qquad y'(t)=(t^2-t\xi)k+(t+\xi).
\]
Then $y'(t)=\xi+t(1-\xi k)+t^2k$, and for $z\in R$
\[
 y(t)e+y'(t)z=(-\xi e+\xi z)+t(\xi he+e-\xi kz+z)+t^2(he+kz).
\]
By Lemma~\ref{asm-theta-agree-split-inputs-siegel-finite}, applied three times,
each of the following matrices is invertible for all but finitely many $t\in F$:
\begin{enumerate}
\item[(a)] $y'(t)$, since $y'(0)=\xi=\xi e+\xi f$ is invertible;
\item[(b)] $y(t)e+y'(t)f$, which is $-\xi e+\xi f$ at $t=0$;
\item[(c)] $y(t)e+y'(t)hf$, which is $h(2\xi^2e+2\xi f)$ at $t=\xi$, since
  $y(\xi)=2\xi^2h$ and $y'(\xi)=2\xi$; this is invertible, as $h$, $2\xi^2$ and
  $2\xi$ are.
\end{enumerate}
As $F$ is infinite, there is $t\in F$ at which all three are invertible. Write
$Y=y(t)$, $Y'=y'(t)$ and $w=Y'^{-1}$.

As $c$ fixes $t\in F$ and $c(\xi)=-\xi$,
$\sigma(Y)=(t^2-t\xi)k+(t+\xi)=Y'$, and likewise $\sigma(Y')=Y$. As $Y$ and $Y'$
are polynomials in $h$ and $k=h^{-1}$, $YY'=Y'Y$, so $wY=wYY'w=wY'Yw=Yw$. Also
$\sigma(w)Y=\sigma(w)\sigma(Y')=\sigma(Y'w)=1$, so
\[
 \sigma(wY)\,wY=Y'\sigma(w)wY=Y'\sigma(w)Yw=Y'w=1,
\]
and $wY=\varphi(x)$ for some $x\in G$. Now $\varphi(x)e+f=w(Ye+Y'f)$ is
invertible by (b), so $\lambda(x)=1$. Since $kwY'hf=khf=f$,
$\varphi(g^{-1}x)e+f=kw(Ye+Y'hf)$ is invertible by (c), so $\lambda(g^{-1}x)=1$.
Hence $\lambda(g)^{-1}=\lambda(g^{-1}x)\lambda(x)^{-1}=1$.
\end{proof}
% Lean: triv_of_cell (FiniteNonunitMatrix.lean): hσh, hσk, hunit; y, y', e1, e2; the finite sets
% f1, f2, f3; then s1, s2, cm, cw, x from hsurj, and l1, l2 from hcell.

\begin{lemma}\label{asm-theta-agree-split-inputs-siegel-artinian}
Let $F$ be a field of characteristic $0$, $E$ a commutative Artinian
$F$-algebra, $c$ an endomorphism of $F$-algebras of $E$ with $c(c(a))=a$ for
$a\in E$, and $\xi$ a unit of $E$ with $c(\xi)=-\xi$; apply $c$ to matrices
entrywise. Let $\beta$ and $\nu$ be finite sets, let $J\in M_\beta(E)$ have a
unit as determinant, and let $C_y\in M_{\beta\times\nu}(E)$ satisfy
$C_y^{\mathsf T}Jc(C_y)=0$. Suppose that there are
$C_x\in M_{\beta\times\nu}(E)$ and $M\in M_{(\nu\sqcup\nu)\times\beta}(E)$ with
$C_x^{\mathsf T}Jc(C_x)=0$, $(C_y\;C_x)M=1$ and $M(C_y\;C_x)=1$. Let $U$ be the
group of $g\in GL_\beta(E)$ with $g^{\mathsf T}Jc(g)=J$. Then every homomorphism
$\lambda\colon U\to\mathbf C^\times$ with $\lambda(p)=1$ for every $p\in U$ such
that $pC_y=C_yt$ for some $t\in GL_\nu(E)$ is trivial.
\end{lemma}
\leanref{Development/AutomorphicSpectrum/Proofs/DivisibilityProofs/FiniteNonunitMatrix.lean}{GlobalTheory.CompanionTheory.Ar.RT.Siegel.char_trivial_artinian}
\begin{proof}
For $x\in M_\beta(E)$ put $\sigma(x)=c(J^{-1}x^{\mathsf T}J)$. Then $\sigma$ is
additive, $\sigma(1)=1$, $\sigma(ax)=c(a)\sigma(x)$ for $a\in E$, and
$\sigma(xy)=\sigma(y)\sigma(x)$, as $c$ is a ring homomorphism and
$J^{-1}y^{\mathsf T}x^{\mathsf T}J=(J^{-1}y^{\mathsf T}J)(J^{-1}x^{\mathsf T}J)$.
As $c$ is an involution, $\sigma(x)y=c(J^{-1}x^{\mathsf T}Jc(y))$. So
$\sigma(x)x=1$ exactly when $x^{\mathsf T}Jc(x)=J$, and then $x$ is invertible,
with inverse $\sigma(x)$. Hence
Situation~\ref{asm-theta-agree-split-inputs-siegel-situation} holds with
$R=M_\beta(E)$, the group $U$, its inclusion $\varphi$ into $M_\beta(E)$, which
is injective with image the set of $x$ with $\sigma(x)x=1$, and the idempotents
$e$ and $f$ below.

Write $M=\left(\begin{smallmatrix}M_1\\M_2\end{smallmatrix}\right)$ with
$M_1,M_2\in M_{\nu\times\beta}(E)$. Then $(C_y\;C_x)M=1$ says
$C_yM_1+C_xM_2=1$, and $M(C_y\;C_x)=1$ says $M_1C_y=1$, $M_1C_x=0$, $M_2C_y=0$
and $M_2C_x=1$. Put $e=C_yM_1$ and $f=C_xM_2$. Then $e+f=1$, $e^2=e$, $f^2=f$
and $ef=fe=0$. If $C\in M_{\beta\times\nu}(E)$ satisfies
$C^{\mathsf T}Jc(C)=0$, then for every $N\in M_{\nu\times\beta}(E)$
\[
 \sigma(CN)\,CN=c\big(J^{-1}N^{\mathsf T}(C^{\mathsf T}Jc(C))c(N)\big)=0.
\]
So $\sigma(e)e=0$ and $\sigma(f)f=0$. As $\sigma(e)+\sigma(f)=\sigma(1)=1$,
\[
 \sigma(e)=\sigma(e)(e+f)+\sigma(f)f=(\sigma(e)+\sigma(f))f=f,
\]
and then $\sigma(f)=1-f=e$.

Let $g\in U$ with $fge=0$. Multiplying by $C_y$ on the right and using
$M_1C_y=1$ gives $fgC_y=0$, so $gC_y=(e+f)gC_y=C_y(M_1gC_y)$. The matrix
$t=M_1gC_y\in M_\nu(E)$ satisfies
$(M_1g^{-1}C_y)t=M_1g^{-1}C_yt=M_1g^{-1}gC_y=M_1C_y=1$, so it
is invertible. Hence $\lambda(g)=1$.

As $F$ has characteristic $0$, it is infinite, and $2$ and $3$ are invertible
in $E$, with $c(2)=2$, $c(1/2)=1/2$ and $c(1/3)=1/3$ as $c$ is $F$-linear. So
Lemma~\ref{asm-theta-agree-split-inputs-siegel-unipotent} gives hypothesis (1)
of Lemma~\ref{asm-theta-agree-split-inputs-siegel-cell}, and the previous
paragraph gives its hypothesis (2). That lemma shows that $\lambda(g)=1$ for
every $g\in U$ such that $ge+f$ is invertible, and
Lemma~\ref{asm-theta-agree-split-inputs-siegel-generic} shows that $\lambda$ is
trivial.
\end{proof}
% Lean: char_trivial_artinian: σ, hsm, hmul, h1, hform, hu, hsurj; M split into M1, M2 (m11, m12,
% m21, m22), e = Cy M1 and f = Cx M2 (hee, hff, hef0, hfe0); iso, ie, if'; hP; hsum, hse, hsf;
% h2, h3; then triv_of_cell (cell_triv (lower_triv ...) hP).

\begin{lemma}\label{asm-theta-agree-split-inputs-siegel}
Let $E$ be a CM field with a theta base $B$
(Definition~\ref{asm-theta-agree-models}(1)), and let $Y$ be a lattice space of
rank $n$ whose matrix $J$ satisfies $c(J)^{\mathsf T}=\epsilon J$ for some
$\epsilon=\pm1$ and $\det J\neq0$. Let $G_Y$ and $P_Y$ be as in
Definition~\ref{asm-theta-agree-doubling}(2). Then every continuous
homomorphism $\lambda\colon G_Y\to\mathbf C^\times$ with $\lambda(p)=1$ for
$p\in P_Y$ is trivial.
\end{lemma}
\leanref{Development/WeilBound/ThetaLifts/LeviDeterminants/DoubledUnitaryCharacterTrivialOfSiegel.lean}{PositivityTheory.ThKF0.parabolic_character_trivial_on_diagonal}
\leanref{Development/WeilBound/ThetaLifts/LeviDeterminants/DoubledUnitaryCharacterTrivialOfSiegel.lean}{PositivityTheory.ThKF0.local_global_character_interpolation}
\leanref{Development/WeilBound/ThetaLifts/LeviDeterminants/DoubledUnitaryCharacterTrivialOfSiegel.lean}{PositivityTheory.ThKF0.Sg.triv_piece}
\leanref{Development/AutomorphicSpectrum/Proofs/DivisibilityProofs/HermitianDarbouxSplitProduct.lean}{GlobalTheory.CompanionTheory.Ar.RT.CellSplit.char_trivial_first}
\begin{proof}
For a finite set $S$ of finite places of $F$, let $\varepsilon_S\in\mathbf A$ have
components $1$ at the infinite places and at the places in $S$, and $0$
elsewhere. It is an idempotent, and $c(\varepsilon_S)=\varepsilon_S$ in
$\mathbf A_E$. Write $J^\square=\operatorname{diag}(J,-J)$ and
$\Delta=\left(\begin{smallmatrix}1_n\\1_n\end{smallmatrix}\right)$.

\emph{Step 1: elements trivial away from $S$.} Let $R_S=F_\infty\times\prod_{v\in S}F_v$
and $E_S=R_S\otimes_FE$, a finite product of fields, hence an Artinian ring, with
the involution $c$. The map sending $a$ to its components at $\infty$ and in $S$
and to its class modulo $\varepsilon_S$ is a ring isomorphism
$\mathbf A\cong R_S\times\mathbf A/\varepsilon_S\mathbf A$; tensoring with $E$
gives $\mathbf A_E\cong E_S\times(\mathbf A/\varepsilon_S\mathbf A)\otimes_FE$,
compatibly with $c$. Let $G_S$ be the group of $x\in GL_{2n}(E_S)$ with
$x^{\mathsf T}J^\square c(x)=J^\square$, and $P_S$ its subgroup of the $x$ with
$x\Delta=\Delta t$ for some $t\in GL_n(E_S)$. The map $\iota_S\colon G_S\to G_Y$
that is $x$ in the first factor and $1$ in the second is a homomorphism; its
image is the set of $g\in G_Y$ with $g-1\in\varepsilon_SM_{2n}(\mathbf A_E)$, and
it maps $P_S$ into $P_Y$, with $t$ extended by $1$ in the second factor. So
$\lambda\circ\iota_S$ is a homomorphism $G_S\to\mathbf C^\times$, trivial on $P_S$.

Lemma~\ref{asm-theta-agree-split-inputs-siegel-artinian} applies over $F$, which
has characteristic $0$, to the Artinian $F$-algebra $E_S$ with the involution
$c$, which is $F$-linear, to the image of $J^\square$ in $M_{2n}(E_S)$ and to
$C_y=\Delta$. Its hypotheses hold. As $E$ is a CM field, $c(a)\neq a$ for some
$a\in E$, and $\xi=1\otimes(a-c(a))$ is a unit of $E_S$ with $c(\xi)=-\xi$. The
determinant $(-1)^n(\det J)^2$ of $J^\square$ is a unit, as $\det J\in E$ is
nonzero. With $C_x=\left(\begin{smallmatrix}1_n\\-1_n\end{smallmatrix}\right)$
and $M=\frac12\left(\begin{smallmatrix}1_n&1_n\\1_n&-1_n\end{smallmatrix}\right)$,
$\Delta^{\mathsf T}J^\square c(\Delta)=J-J=0$,
$C_x^{\mathsf T}J^\square c(C_x)=J-J=0$ and $(\Delta\;C_x)M=M(\Delta\;C_x)=1$. The
group $U$ of that lemma is $G_S$, and its parabolic subgroup is $P_S$. So
$\lambda\circ\iota_S$ is trivial, and $\lambda(g)=1$ for every $g\in G_Y$ with
$g-1\in\varepsilon_SM_{2n}(\mathbf A_E)$.
% Lean: char_trivial_first (HermitianDarbouxSplitProduct.lean) applies char_trivial_artinian
% to E1 = R_S tensor E (IsArtinianRing.of_finite over Ainf F, isArtinian_inf), c1 = id tensor
% sigma, xi = 1 tensor (x - sigma x) from exists_moved, dbl (J.map pr1) with isUnit_det_dbl, and
% diagCols with isLagr_diagCols (Cx = fromRows 1 (-1), M = (1/2) fromBlocks 1 1 1 (-1)). The rest
% of Step 1 is uLift, glPr, pr_mem, lift_diag and snd_lift_sub_one_apply_eq_zero, with splitS
% and split (DoubledUnitaryCharacterTrivialOfSiegel.lean).

\emph{Step 2: approximation.} Let $g\in G_Y$ and put
$q_S=\varepsilon_Sg+(1-\varepsilon_S)1$. As $\varepsilon_S$ is an idempotent fixed by
$c$, $q_S$ has inverse $\varepsilon_Sg^{-1}+(1-\varepsilon_S)1$ and
$q_S^{\mathsf T}J^\square c(q_S)=\varepsilon_Sg^{\mathsf T}J^\square c(g)+(1-\varepsilon_S)J^\square=J^\square$.
So $q_S\in G_Y$, and $q_S-1=\varepsilon_S(g-1)$, so $\lambda(q_S)=1$ by Step 1. For
$z\in\mathbf A_E$, $\varepsilon_Sz$ agrees with $z$ at the infinite places and in
$S$ and is $0$ elsewhere; as $z$ is integral at almost all places,
$\varepsilon_Sz\to z$ as $S$ runs through the finite sets of finite places, ordered
by inclusion. So $q_S\to g$ and $q_S^{-1}\to g^{-1}$ entrywise, and by
continuity $\lambda(g)=\lim\lambda(q_S)=1$.
\end{proof}
% Lean: parabolic_character_trivial_on_diagonal: hJ from det J /= 0; Step 2 is
% local_global_character_interpolation (piece, mix_mul, mix_self, tendsto_mix' with
% tendsto_eps_tens) and the limit h2; Step 1 is triv_piece through hermitian_darboux_split_product
% = char_trivial_first, for the splitting split E (splitS S) of A_E.

\begin{proposition}\label{asm-theta-agree-split-inputs-seesaw}
Let $B$ be a theta base for a CM field $E$, with the notation of
Definitions~\ref{asm-theta-agree-split-inputs-heis}
and~\ref{asm-theta-agree-split-inputs-side}.
\begin{enumerate}
\item Let $(A,Y,d,\chi,\omega)$ and $(A_1,Y,d_1,\chi,\omega_1)$ be Weil sides,
  with $A$, $A_1$ and $Y$ of ranks $a$, $a_1=a-2$ and $b$ and with index sets
  $\kappa$ and $\kappa_1$, where the matrix $J_Y$ of $Y$ satisfies
  $c(J_Y)^{\mathsf T}=\epsilon J_Y$ for some $\epsilon=\pm1$ and
  $\det J_Y\neq0$. Let $S$ be a symplectic automorphism of
  $\mathbf A^\kappa\times\mathbf A^\kappa$, $j\colon\kappa_1\to\kappa$ an
  injective map, $e$ an $(S,j)$-evaluation, $\iota\in M_{a\times a_1}(\mathbf A_E)$
  and $x,y\in\mathbf A_E^a$, such that:
  \begin{enumerate}
  \item[(i)] $d(u)$ lies over $d_1(z)$ for $(S,j)$ if and only if
    $u=\iota z+xq^{\mathsf T}$ for some $q\in\mathbf A_E^b$;
  \item[(ii)] $(S\,d(yq^{\mathsf T}))_2=0$ and $(S\,d(yq^{\mathsf T}))_1\circ j=0$
    for every $q\in\mathbf A_E^b$;
  \item[(iii)] $yq^{\mathsf T}=0$ only for $q=0$;
  \item[(iv)] every $u\in M_{a\times b}(\mathbf A_E)$ is
    $\iota z+xq^{\mathsf T}+yq'^{\mathsf T}$ for some $z$, $q$ and $q'$.
  \end{enumerate}
  Then $e(\omega(h)\phi)=\omega_1(h)e(\phi)$ for $h\in U_Y$ and
  $\phi\in\mathcal S(\kappa)$.
\item Let $(A,Y,d,\chi,\omega)$ and $(A,Y_1,d_1,\chi,\omega_1)$ be Weil sides,
  with $A$, $Y$ and $Y_1$ of ranks $a$, $b$ and $b_1$ and with index sets
  $\kappa$ and $\kappa_1$, and let $H=(x,y,\iota,\mathfrak b)$ be an adapted
  plane in $Y$ over $Y_1$ (Definition~\ref{asm-theta-agree-planes}(1)), where
  $c(J_Y)^{\mathsf T}=\epsilon J_Y$ for some $\epsilon=\pm1$ and
  $\det J_{Y_1}\neq0$. Let $S$, $j$ and $e$ be as in (1), such that $d(u)$ lies
  over $d_1(z)$ if and only if $u=z\iota^{\mathsf T}+qx^{\mathsf T}$ for some
  $q\in\mathbf A_E^a$. Let $m\in U_Y$, $g\in U_{Y_1}$, $t\in\mathbf A_E^\times$
  and $r\in M_{1\times b_1}(\mathbf A_E)$ satisfy $m\iota=\iota g+xr$ and
  $mx=tx$. Then $e(\omega(m)\phi)=\chi(t)|t|^{a/2}\omega_1(g)e(\phi)$ for
  $\phi\in\mathcal S(\kappa)$.
\end{enumerate}
\end{proposition}
\leanref{Development/WeilBound/ThetaLifts/LeviDeterminants/CharacterAdmissibility/DetConjMulVec.lean}{PositivityTheory.ThKF0.Chr.TA}
\leanref{Development/WeilBound/ThetaLifts/LeviDeterminants/CharacterAdmissibility/DetConjMulVec.lean}{PositivityTheory.ThKF0.Chr.TA'}
\leanref{Development/WeilBound/ThetaLifts/LeviDeterminants/CharacterAdmissibility/DetConjMulVec.lean}{PositivityTheory.ThKF0.Chr.det_conj_mulVec}
\leanref{Development/WeilBound/ThetaLifts/LeviDeterminants/CharacterAdmissibility/ExistsDoubledUnitaryBlock.lean}{PositivityTheory.ThKF0.Chr.TB}
\leanref{Development/WeilBound/ThetaLifts/LeviDeterminants/CharacterAdmissibility/LevScalar2.lean}{PositivityTheory.ThKF0.Chr.TU}
\leanref{Development/WeilBound/ThetaLifts/LeviDeterminants/CharacterAdmissibility/DoublingGramFromBlocks.lean}{PositivityTheory.ThKF0.Chr.Fr.emb}
\leanref{Development/WeilBound/ThetaLifts/LeviDeterminants/CharacterAdmissibility/DoublingGramFromBlocks.lean}{PositivityTheory.ThKF0.Chr.Fr.exists_gamma}
\leanref{Development/WeilBound/CoreStatements/CompactAveragingSplitting/EigenvectorLifting/SiegelTransportAlg.lean}{PositivityTheory.ThK.SiegelT.rule_of_conj}
\begin{proof}
For a matrix $N$ write $N^{(2)}=\operatorname{diag}(N,N)$, and let
$\Delta=\left(\begin{smallmatrix}1_b\\1_b\end{smallmatrix}\right)$. For a Weil
side, $d^\square(u\Delta^{\mathsf T})=d(u)\boxminus d(u)$, as
$u\Delta^{\mathsf T}=(u\ u)$; and $d^\square$ is bijective, as $d$ is. Let
$\omega^\square$ and $\omega_1^\square$ be the doubled operators of the two
sides.

(1) \emph{Step 0: the complement of $j$.} Let $\kappa'$ be the complement of
$j(\kappa_1)$ in $\kappa$ and $i$ its inclusion. For $q\in\mathbf A_E^b$ put
$\theta(q)=(S\,d(yq^{\mathsf T}))_1\circ i$; by (ii),
$S\,d(yq^{\mathsf T})=(i_!\theta(q),0)$. The $\mathbf A$-linear map
$\theta\colon\mathbf A_E^b\to\mathbf A^{\kappa'}$ is bijective. If $\theta(q)=0$,
then $S\,d(yq^{\mathsf T})=0$, so $q=0$ by (iii). If $v\in\mathbf A^{\kappa'}$,
write $d^{-1}S^{-1}(i_!v,0)=\iota z+xq'^{\mathsf T}+yq^{\mathsf T}$ by (iv); by
(i), $d(\iota z+xq'^{\mathsf T})$ lies over $d_1(z)$, so the first component of
its image under $S$ vanishes on $\kappa'$, and $v=\theta(q)$. So, for
$u\in\mathbf A^{\kappa'}$,
\[
v_u=S^{-1}(i_!u,0)=d\bigl(y\,\theta^{-1}(u)^{\mathsf T}\bigr).\tag{V}
\]

\emph{Step 1: the doubled objects.} By (i) and
Lemma~\ref{asm-theta-agree-split-inputs-heis-double}(1), $d^\square(u_1\ u_2)$
lies over $d_1^\square(z_1\ z_2)$ for $(S^\square,j\sqcup j)$ if and only if
$(u_1\ u_2)=\iota(z_1\ z_2)+xq^{\mathsf T}$ for some $q\in\mathbf A_E^{2b}$. By
Lemma~\ref{asm-theta-agree-split-inputs-heis-lines}(4) let $\mathcal E$ be an
$(S^\square,j\sqcup j)$-evaluation, and by
Lemma~\ref{asm-theta-agree-split-inputs-heis-double}(4) let $\ell$ be a cyclic
form on the functions on $\mathbf A^{\kappa_1\sqcup\kappa_1}$ invariant under the
$\rho(w\boxminus w)$.

\emph{Step 2: a character of $G_Y$.} Let $p\in G_Y$, and let $d^\square(u)$ lie
over $d_1^\square(z)$, so $u=\iota z+xq^{\mathsf T}$. Then
$up^{\mathsf T}=\iota(zp^{\mathsf T})+x(pq)^{\mathsf T}$, so $d^\square(up^{\mathsf T})$
lies over $d_1^\square(zp^{\mathsf T})$; by (b) of the two sides,
$\omega^\square(p)\rho(d^\square(u))=\rho(d^\square(up^{\mathsf T}))\omega^\square(p)$
and
$\omega_1^\square(p^{-1})\rho(d_1^\square(zp^{\mathsf T}))=\rho(d_1^\square(z))\omega_1^\square(p^{-1})$
on the Schwartz-Bruhat functions. With (a) and (b) of the two sides,
Lemma~\ref{asm-theta-agree-split-inputs-heis-scalars}(2) gives a continuous
homomorphism $\Lambda\colon G_Y\to\mathbf C^\times$ with
\[
\mathcal E(\omega^\square(p)\Psi)=\Lambda(p)\,\omega_1^\square(p)\mathcal E(\Psi)\qquad(p\in G_Y,\ \Psi\in\mathcal S(\kappa\sqcup\kappa)).\tag{A}
\]

\emph{Step 3: $\Lambda$ is trivial on $P_Y$.} Let $L$ be the form of
Lemma~\ref{asm-theta-agree-split-inputs-heis-double}(5) for $\kappa'$ and $i$.
Let $p\in P_Y$ with $p\Delta=\Delta t$, and let $g=\theta\circ t\circ\theta^{-1}$,
an $\mathbf A$-linear automorphism of $\mathbf A^{\kappa'}$. Its determinant is
that of $t$ as an $\mathbf A$-linear map of $\mathbf A_E^b$, which is
$N_{\mathbf A_E/\mathbf A}(\det t)$; so $|\det g|_{\mathbf A}=|\det t|$. We check
the hypotheses of Lemma~\ref{asm-theta-agree-split-inputs-heis-double}(5) for
$O=\omega^\square(p)$, $g$ and the scalar
$\Lambda(p)\chi(\det t)|\det t|^{a_1/2}$ in place of $a$. Let
$u\in\mathbf A^{\kappa'}$ and $q=\theta^{-1}(u)$. By (V),
$v_u\boxminus v_u=d^\square(yq^{\mathsf T}\Delta^{\mathsf T})$, and
$yq^{\mathsf T}\Delta^{\mathsf T}p^{\mathsf T}=yq^{\mathsf T}t^{\mathsf T}\Delta^{\mathsf T}=y(tq)^{\mathsf T}\Delta^{\mathsf T}$
with $\theta(tq)=gu$; so
$\omega^\square(p)\rho(v_u\boxminus v_u)=\rho(v_{gu}\boxminus v_{gu})\omega^\square(p)$.
The form $\ell$ is invariant under
$\rho(d_1^\square(z\Delta^{\mathsf T}))=\rho(d_1(z)\boxminus d_1(z))$, so the
rule (b) of the side $(A_1,Y,d_1,\chi,\omega_1)$ and (A) give
$\ell(\mathcal E(\omega^\square(p)\Psi))=\Lambda(p)\chi(\det t)|\det t|^{a_1/2}\ell(\mathcal E(\Psi))$.
So $L(\omega^\square(p)\Psi)=\Lambda(p)\chi(\det t)|\det t|^{a_1/2+1}L(\Psi)$. On
the other hand $L$ is invariant under
$\rho(d^\square(z\Delta^{\mathsf T}))=\rho(d(z)\boxminus d(z))$, so the rule (b)
of the side $(A,Y,d,\chi,\omega)$ gives
$L(\omega^\square(p)\Psi)=\chi(\det t)|\det t|^{a/2}L(\Psi)$. Taking $\Psi$ with
$L(\Psi)\neq0$, and as $a=a_1+2$, $\Lambda(p)=1$.

\emph{Step 4.} By Lemma~\ref{asm-theta-agree-split-inputs-siegel} for $Y$,
$\Lambda=1$.

\emph{Step 5: cancellation.} By
Lemma~\ref{asm-theta-agree-split-inputs-heis-double}(3) choose
$\phi_1,\phi_2\in\mathcal S(\kappa)$ with $\ell(\mathcal E(\phi_1\boxtimes\phi_2))\neq0$,
and by part (2) of that lemma $B\in\mathcal S(\kappa_1)$ with
$\mathcal E(\phi\boxtimes\phi_2)=e(\phi)\boxtimes B$ for $\phi\in\mathcal S(\kappa)$; then
$B\neq0$. Let $h\in U_Y$; then $\bar h=\operatorname{diag}(h,1)\in G_Y$. For
$\phi\in\mathcal S(\kappa)$, by (c) of the two sides and Step 4,
\[
e(\omega(h)\phi)\boxtimes B=\mathcal E(\omega^\square(\bar h)(\phi\boxtimes\phi_2))
=\omega_1^\square(\bar h)(e(\phi)\boxtimes B)=\omega_1(h)e(\phi)\boxtimes B.
\]
As $B\neq0$, $e(\omega(h)\phi)=\omega_1(h)e(\phi)$.

(2) \emph{Step 0: the frame.} Let $Q$ be the inverse of the matrix
$\mathrm{fr}_H$ of Lemma~\ref{asm-theta-agree-evaluation-adapted-frame}, let
$\pi$ be its first $b_1$ rows and $\rho_H$ its two rows at $x$ and $y$, and put
$Z=(x\ y)\in M_{b\times2}(E)$. By that lemma,
$\iota\pi+Z\rho_H=1$, $\pi\iota=1$, $\pi Z=0$, $\rho_H\iota=0$, $\rho_HZ=1$,
$\iota^{\mathsf T}J_Yc(\iota)=J_{Y_1}$, $\iota^{\mathsf T}J_Yc(Z)=0$,
$Z^{\mathsf T}J_Yc(Z)=H_\epsilon=\left(\begin{smallmatrix}0&1\\\epsilon&0\end{smallmatrix}\right)$
and $c(J_{Y_1})^{\mathsf T}=\epsilon J_{Y_1}$. Let $G_\epsilon$ be the group of
$k\in GL_4(\mathbf A_E)$ with
$k^{\mathsf T}\operatorname{diag}(H_\epsilon,-H_\epsilon)c(k)=\operatorname{diag}(H_\epsilon,-H_\epsilon)$,
the doubled group of $H_\epsilon$.
For $p\in M_{2b_1}(\mathbf A_E)$ and $k\in M_4(\mathbf A_E)$ put
$\operatorname{em}(p,k)=\iota^{(2)} p\pi^{(2)}+Z^{(2)} k\rho_H^{(2)}$.
The relations above, for the doubled matrices, give
$\operatorname{em}(p,k)\operatorname{em}(p',k')=\operatorname{em}(pp',kk')$,
$\operatorname{em}(1,1)=1$, $\operatorname{em}(p,k)\iota^{(2)}=\iota^{(2)} p$,
$\operatorname{em}(p,k)Z^{(2)}=Z^{(2)} k$, and
\[
\operatorname{em}(p,k)^{\mathsf T}J_Y^\square c(\operatorname{em}(p,k))
=(\pi^{(2)})^{\mathsf T}\bigl(p^{\mathsf T}J_{Y_1}^\square c(p)\bigr)c(\pi^{(2)})
+(\rho_H^{(2)})^{\mathsf T}\bigl(k^{\mathsf T}H_\epsilon^\square c(k)\bigr)c(\rho_H^{(2)}),
\]
where, as in Definition~\ref{asm-theta-agree-doubling}(2),
$N^\square=\operatorname{diag}(N,-N)$. For $p\in G_{Y_1}$ and
$k\in G_\epsilon$ the right side is the value at $(1,1)$, which is $J_Y^\square$;
so $\operatorname{em}$ maps $G_{Y_1}\times G_\epsilon$ into $G_Y$, and
$p\mapsto\operatorname{em}(p,1)$ is a continuous homomorphism.

\emph{Step 1: a transported rule.} Let $\gamma_0$ be the matrix, in the basis
of the columns of $Z^{(2)}$ (that is $(x,0)$, $(y,0)$, $(0,x)$, $(0,y)$),
\[
\gamma_0=\begin{pmatrix}1&0&0&0\\0&1&0&-1\\0&0&0&\epsilon\\-1&0&1&0\end{pmatrix}.
\]
A direct computation shows $\gamma_0\in G_\epsilon$ and
$\gamma_0\left(\begin{smallmatrix}1_2\\1_2\end{smallmatrix}\right)=(e_1\ \ \epsilon e_3)$,
with $e_1,\ldots,e_4$ the unit vectors.
Put $\gamma=\operatorname{em}(1,\gamma_0)\in G_Y$,
$\Delta_1=\left(\begin{smallmatrix}1_{b_1}\\1_{b_1}\end{smallmatrix}\right)$ and
$C=(\iota^{(2)}\Delta_1\ \ x^{(2)})$, a $2b\times(b_1+2)$ matrix. As
$\iota\pi+Z\rho_H=1$,
$\Delta=\iota^{(2)}\Delta_1\pi+Z^{(2)}\left(\begin{smallmatrix}1_2\\1_2\end{smallmatrix}\right)\rho_H$,
so
\[
\gamma\Delta=\iota^{(2)}\Delta_1\pi+Z^{(2)}(e_1\ \ \epsilon e_3)\rho_H
=\iota^{(2)}\Delta_1\pi+x^{(2)}\operatorname{diag}(1,\epsilon)\rho_H=Cs,
\]
with $s$ the matrix with rows $\pi$ and $\operatorname{diag}(1,\epsilon)\rho_H$,
which has inverse $s'=(\iota\ \ Z\operatorname{diag}(1,\epsilon))$. Then:
\begin{enumerate}
\item[(R)] if $\lambda$ is admissible with
  $\lambda(\rho(d^\square(zC^{\mathsf T}))\Psi)=\lambda(\Psi)$ for all
  $z\in M_{a\times(b_1+2)}(\mathbf A_E)$ and $\Psi\in\mathcal S(\kappa\sqcup\kappa)$,
  and $p\in G_Y$ and $t'\in GL_{b_1+2}(\mathbf A_E)$ satisfy $pC=Ct'$, then
  $\lambda(\omega^\square(p)\Psi)=\chi(\det t')|\det t'|^{a/2}\lambda(\Psi)$.
\end{enumerate}
Indeed, $\lambda\circ\omega^\square(\gamma)$ is admissible, and as
$z\Delta^{\mathsf T}\gamma^{\mathsf T}=z(\gamma\Delta)^{\mathsf T}=(zs^{\mathsf T})C^{\mathsf T}$,
it is invariant under the $\rho(d^\square(z\Delta^{\mathsf T}))$. Further
$\gamma^{-1}p\gamma\Delta=\gamma^{-1}pCs=\gamma^{-1}Ct's=\Delta s't's$, so
$\gamma^{-1}p\gamma\in P_Y$ with $\det(s't's)=\det t'$. The rule (b) of the side
$(A,Y,d,\chi,\omega)$, for $\lambda\circ\omega^\square(\gamma)$, $\gamma^{-1}p\gamma$
and $\omega^\square(\gamma^{-1})\Psi$, gives (R).

\emph{Step 2: the doubled objects.} By the hypothesis on $(S,j)$ and
Lemma~\ref{asm-theta-agree-split-inputs-heis-double}(1),
$d^\square(u_1\ u_2)$ lies over $d_1^\square(z)$ for $(S^\square,j\sqcup j)$ if
and only if $(u_1\ u_2)=z(\iota^{(2)})^{\mathsf T}+q(x^{(2)})^{\mathsf T}$
for some $q\in M_{a\times2}(\mathbf A_E)$. Let $\mathcal E$ and $\ell$ be as in Step 1 of
(1), for $\kappa_1$ the index set of $d_1$. For
$z=(z_1\ z_2)$ with $z_1\in M_{a\times b_1}(\mathbf A_E)$,
$zC^{\mathsf T}=z_1\Delta_1^{\mathsf T}(\iota^{(2)})^{\mathsf T}+z_2(x^{(2)})^{\mathsf T}$,
so $d^\square(zC^{\mathsf T})$ lies over $d_1^\square(z_1\Delta_1^{\mathsf T})=d_1(z_1)\boxminus d_1(z_1)$
and
$\ell(\mathcal E(\rho(d^\square(zC^{\mathsf T}))\Psi))=\ell(\mathcal E(\Psi))$. So (R) applies to
$\ell\circ\mathcal E$. The form $\ell$ is invariant under the
$\rho(d_1^\square(z\Delta_1^{\mathsf T}))$, so the rule (b) of the side
$(A,Y_1,d_1,\chi,\omega_1)$ applies to it. By
Lemma~\ref{asm-theta-agree-split-inputs-heis-lines}(5) choose $\Psi_0$ with
$\ell(\mathcal E(\Psi_0))\neq0$.

\emph{Step 3: a character of $G_{Y_1}$.} For $p\in G_{Y_1}$ put
$\operatorname{em}(p)=\operatorname{em}(p,1)$, so that
$\operatorname{em}(p)\iota^{(2)}=\iota^{(2)} p$ and
$\operatorname{em}(p)x^{(2)}=x^{(2)}$. If $u=z(\iota^{(2)})^{\mathsf T}+q(x^{(2)})^{\mathsf T}$,
then
$u\operatorname{em}(p)^{\mathsf T}=(zp^{\mathsf T})(\iota^{(2)})^{\mathsf T}+q(x^{(2)})^{\mathsf T}$,
so $d^\square(u\operatorname{em}(p)^{\mathsf T})$ lies over $d_1^\square(zp^{\mathsf T})$.
As in Step 2 of (1),
Lemma~\ref{asm-theta-agree-split-inputs-heis-scalars}(1) for
$T=\omega^\square(\operatorname{em}(p))$, $T_1=\omega_1^\square(p)$ and
$T'_1=\omega_1^\square(p^{-1})$ gives $\Lambda(p)$ with
$\mathcal E(\omega^\square(\operatorname{em}(p))\Psi)=\Lambda(p)\,\omega_1^\square(p)\mathcal E(\Psi)$.
As $\mathcal E(\Psi_0)\neq0$ and the $\omega_1^\square(p)$ are invertible, $\Lambda$ is a
homomorphism $G_{Y_1}\to\mathbf C^\times$, as in the proof of
Lemma~\ref{asm-theta-agree-split-inputs-heis-scalars}(2). It is continuous: let
$p_0\in G_{Y_1}$; as $\ell$ is cyclic, there is $z$ with
$\ell(\rho(d_1^\square(z))\omega_1^\square(p_0)\mathcal E(\Psi_0))\neq0$; with
$z'=z(p_0^{-1})^{\mathsf T}$ and
$\Psi=\rho(d^\square(z'(\iota^{(2)})^{\mathsf T}))\Psi_0$,
$\rho(d_1^\square(z))\omega_1^\square(p_0)\mathcal E(\Psi_0)=\omega_1^\square(p_0)\rho(d_1^\square(z'))\mathcal E(\Psi_0)=\omega_1^\square(p_0)\mathcal E(\Psi)$;
so $\ell(\omega_1^\square(p)\mathcal E(\Psi))\neq0$ near $p_0$, and there
$\Lambda(p)=\ell(\mathcal E(\omega^\square(\operatorname{em}(p))\Psi))/\ell(\omega_1^\square(p)\mathcal E(\Psi))$,
a quotient of continuous functions. It is trivial on $P_{Y_1}$: if
$p\Delta_1=\Delta_1t'$, then $\operatorname{em}(p)C=C\operatorname{diag}(t',1_2)$,
and (R) and the rule (b) for $\ell$ give
\[
\chi(\det t')|\det t'|^{a/2}\ell(\mathcal E(\Psi_0))=\ell(\mathcal E(\omega^\square(\operatorname{em}(p))\Psi_0))
=\Lambda(p)\,\ell(\omega_1^\square(p)\mathcal E(\Psi_0))=\Lambda(p)\chi(\det t')|\det t'|^{a/2}\ell(\mathcal E(\Psi_0)).
\]
By Lemma~\ref{asm-theta-agree-split-inputs-siegel} for $Y_1$, $\Lambda=1$:
\[
\mathcal E(\omega^\square(\operatorname{em}(p))\Psi)=\omega_1^\square(p)\mathcal E(\Psi)\qquad(p\in G_{Y_1}).\tag{M}
\]

\emph{Step 4: the scalar.} Put $\bar m=\operatorname{diag}(m,1)\in G_Y$,
$\bar g=\operatorname{diag}(g,1)\in G_{Y_1}$, $\bar r=\operatorname{diag}(r,0)$
and $p=\bar m\operatorname{em}(\bar g)^{-1}$. As $\bar m\iota^{(2)}=\iota^{(2)} \bar g+x^{(2)} \bar r$
and $\bar m x^{(2)}=x^{(2)}\operatorname{diag}(t,1)$,
\[
p\iota^{(2)}=\iota^{(2)}+x^{(2)} \bar r(\bar g)^{-1},\qquad
px^{(2)}=x^{(2)}\operatorname{diag}(t,1).
\]
So if $u=z(\iota^{(2)})^{\mathsf T}+q(x^{(2)})^{\mathsf T}$, then
$up^{\mathsf T}=z(\iota^{(2)})^{\mathsf T}+q'(x^{(2)})^{\mathsf T}$ for some
$q'$, and $d^\square(up^{\mathsf T})$ lies over $d_1^\square(z)$.
Lemma~\ref{asm-theta-agree-split-inputs-heis-scalars}(1) for
$T=\omega^\square(p)$ and $T_1=T'_1=1$ gives $\beta$ with
$\mathcal E(\omega^\square(p)\Psi)=\beta \mathcal E(\Psi)$. Further $pC=CN$ with $N$ the block
lower triangular matrix with diagonal blocks $1_{b_1}$ and $\operatorname{diag}(t,1)$
and lower left block $\bar r(\bar g)^{-1}\Delta_1$, of determinant $t$. By
(R) for $\ell\circ\mathcal E$,
$\beta\,\ell(\mathcal E(\Psi_0))=\ell(\mathcal E(\omega^\square(p)\Psi_0))=\chi(t)|t|^{a/2}\ell(\mathcal E(\Psi_0))$,
so $\beta=\chi(t)|t|^{a/2}$.

\emph{Step 5: cancellation.} Choose $\phi_2$ and $B\neq0$ as in Step 5 of (1).
For $\phi\in\mathcal S(\kappa)$, as
$\omega^\square(\bar m)=\omega^\square(p)\omega^\square(\operatorname{em}(\bar g))$,
(c) of the two sides and (M) give
\[
e(\omega(m)\phi)\boxtimes B=\mathcal E(\omega^\square(\bar m)(\phi\boxtimes\phi_2))
=\beta\,\omega_1^\square(\bar g)(e(\phi)\boxtimes B)
=\beta\,\omega_1(g)e(\phi)\boxtimes B,
\]
and as $B\neq0$, $e(\omega(m)\phi)=\chi(t)|t|^{a/2}\omega_1(g)e(\phi)$.
\end{proof}
% Lean: (1) is TA' (the complement subtype, theta0 and its bijectivity from hy, hy0, hsp and hrel)
% and TA: hrel2, E2 from Ax.exists_E, l from Ax.exists_diag, hdiag (D), A from Ax.exists_char,
% L from Ax.exists_avg, hA1 with det_conj_mulVec (g, |det g| = modu of the norm of det t), hr
% and hr1 (the Siegel rules of the doubled sides) and e4 for the exponents, hch, then
% exists_unitaryDoubled_blockDiag_one, Ax.box_ne, Ax.box_E and box_cancel. (2) is TU (and TB, its
% case r = 0, which the Lean proves by the same argument with lev_scalar and blk): the frame is
% Fr (fr_plane), em is Fr.emb (Fr.dbl, Fr.mul, Fr.gram, emb_i, embeddingAtLocalFactor,
% continuous_emb), gamma0 is exists_gamma0, gamma, s and s' are Fr.exists_gamma, (R) is
% rule_of_conj, Step 3 is the homomorphism A with hA (continuity through hcyc and hF.dbl
% continuous_emb, triviality on the parabolic through hL with blk t' 1 and hL1), Step 4 is k3, k4,
% lev_scalar2 and hbv with low, Step 5 is e5, e6 and box_cancel.

\begin{lemma}\label{asm-theta-agree-split-inputs-charts}
Let $B$ be a theta base for a CM field $E$, with the notation of
Definition~\ref{asm-theta-agree-split-inputs-heis}. Let $X$ be a Hermitian and
$W$ a skew-Hermitian lattice space, of ranks $n$ and $d$ and of some
signatures, let $M$ be a regular model for $(X,W)$ with isomorphism $e^M$ and
Schrödinger operators $\rho_M$, put $I=\{1,\ldots,n\}\times\{1,\ldots,d\}$, and
let $h$ and $\langle\ ,\ \rangle$ be as in
Definition~\ref{asm-theta-agree-models}(3).
\begin{enumerate}
\item Let $H=(x,y,\iota,\mathfrak b)$ be an adapted plane in $X$ over a
  Hermitian lattice space $X_1$ of rank $n_1$ and some signature, let $M_1$ be a
  regular model for $(X_1,W)$ with isomorphism $e^1$, and put
  $I_1=\{1,\ldots,n_1\}\times\{1,\ldots,d\}$. Then there are an injective map
  $j\colon I_1\to I$ and an $\mathbf A$-linear isomorphism $e^a$ from
  $\mathbb W(X,W)$ onto $\mathbf A^I\times\mathbf A^I$ with
  $\langle u,u'\rangle=[e^a(u),e^a(u')]$, mapping $M_{n\times d}(E)$ into
  $F^I\times F^I$, such that for $z\in\mathbb W(X_1,W)$ and $q\in\mathbf A_E^d$:
  $e^a(\iota z)=(j_!e^1_1(z),j_!e^1_2(z))$; $e^a(xq^{\mathsf T})_1=0$ and
  $e^a(xq^{\mathsf T})_2\circ j=0$; and $e^a(yq^{\mathsf T})_2=0$ and
  $e^a(yq^{\mathsf T})_1\circ j=0$.
\item In (1), let $e$ be an evaluation map for $(M,M_1,\iota,1,\mathcal J_x(H))$,
  with $\mathcal J_x(H)=\{xq^{\mathsf T}\}$ as in
  Lemma~\ref{asm-theta-agree-split-inputs-xunique}, and put
  $S=e^a\circ(e^M)^{-1}$. Then $S$ is a symplectic automorphism, $e$ is an
  $(S,j)$-evaluation, $e^M(u)$ lies over $e^1(z)$ for $(S,j)$ if and only if
  $u=\iota z+xq^{\mathsf T}$ for some $q\in\mathbf A_E^d$, and
  $(S\,e^M(yq^{\mathsf T}))_2=0$ and $(S\,e^M(yq^{\mathsf T}))_1\circ j=0$ for
  every $q$.
\item Let $H'=(x',y',\iota',\mathfrak b')$ be an adapted plane in $W$ over a
  lattice space $W_1$ of rank $d_1$, let $M_1$ be a regular model for $(X,W_1)$
  with isomorphism $e^1$, put $I_1=\{1,\ldots,n\}\times\{1,\ldots,d_1\}$, and
  let $e$ be an evaluation map for $(M,M_1,1,\iota',\mathcal J_0(H'))$. Then
  there are a symplectic automorphism $S$ of $\mathbf A^I\times\mathbf A^I$ and
  an injective map $j\colon I_1\to I$ such that $e$ is an $(S,j)$-evaluation,
  $e^M(u)$ lies over $e^1(z)$ for $(S,j)$ if and only if
  $u=z\iota'^{\mathsf T}+qx'^{\mathsf T}$ for some $q\in\mathbf A_E^n$, and
  $(S\,e^M(qy'^{\mathsf T}))_2=0$ and $(S\,e^M(qy'^{\mathsf T}))_1\circ j=0$ for
  every $q\in\mathbf A_E^n$.
\end{enumerate}
\end{lemma}
\leanref{Development/WeilBound/ThetaLifts/LeviDeterminants/MixedModels/ExistsAdaptedIndexHermitian.lean}{PositivityTheory.Mix.exists_adapted_index_hermitian}
\leanref{Development/WeilBound/ThetaLifts/LeviDeterminants/CharacterAdmissibility/DetNeHerm.lean}{PositivityTheory.ThKF0.Chr.setX}
\leanref{Development/WeilBound/ThetaLifts/LeviDeterminants/CharacterAdmissibility/FrOfFrame.lean}{PositivityTheory.ThKF0.Chr.chartW}
\leanref{Development/WeilBound/ThetaLifts/LeviDeterminants/CharacterAdmissibility/FrOfFrame.lean}{PositivityTheory.ThKF0.Chr.setW}
\leanref{Development/WeilBound/ThetaLifts/LeviDeterminants/CharacterAdmissibility/FrOfFrame.lean}{PositivityTheory.ThKF0.Chr.isEv_adm}
\begin{proof}
(1) Let $Q$ be the inverse of the matrix $\mathrm{fr}_H$ of
Lemma~\ref{asm-theta-agree-evaluation-adapted-frame}, index the rows of $Q$ by
$\{1,\ldots,n_1\}$, $x$ and $y$, and for $z\in\mathbb W(X,W)$ put $w(z)=Qz$.
Let $\beta(z)\in\mathbf A_E^d$ be the row of $w(z)$ at $y$,
$\delta(z)=J_Wc(w(z)_x)$ with $w(z)_x$ the row of $w(z)$ at $x$ read as a
column, and $\pi(z)$ the block of the first $n_1$ rows of $w(z)$. As $z=\mathrm{fr}_Hw(z)$ and, by part (2) of that
lemma with $\eta=1$,
$\mathrm{fr}_H^{\mathsf T}J_Xc(\mathrm{fr}_H)=\operatorname{diag}\bigl(J_{X_1},\left(\begin{smallmatrix}0&1\\1&0\end{smallmatrix}\right)\bigr)$,
expanding over the rows at $x$ and $y$ gives
\[
h(z,z')=\sum_{k,l}w_{x,k}(J_W)_{kl}\,c(w'_{y,l})+\sum_{k,l}w_{y,k}(J_W)_{kl}\,c(w'_{x,l})+h_1(\pi z,\pi z'),
\]
with $w=w(z)$, $w'=w(z')$ and $h_1$ the form $h$ for $(X_1,W)$. The second sum
is $\beta(z)\cdot\delta(z')$, and as $c(J_W)^{\mathsf T}=-J_W$ the first is
$-c(\beta(z')\cdot\delta(z))$. So
$h(z,z')=\beta(z)\cdot\delta(z')-c(\beta(z')\cdot\delta(z))+h_1(\pi z,\pi z')$,
which is hypothesis (1) of
Lemma~\ref{asm-theta-agree-evaluation-adapted-model} for $\kappa=\{1,\ldots,d\}$,
$X_1$ and $W_1=W$; and its hypothesis (2) holds, as $Q$ has entries in $E$. The
first two paragraphs of the proof of that lemma use only these two hypotheses
and $nd=2d+n_1d$; they give a bijection
$\epsilon\colon I\to(\kappa\times\{0,1\})\sqcup I_1$, maps $p$ and $q$ with
$p(u)=0$ when $\beta(u)=0$ and $q(u)=0$ when $\delta(u)=0$, and the isomorphism
$e^a(u)=([p(u),e^1_1(\pi u)],[q(u),e^1_2(\pi u)])$, with
$\langle u,u'\rangle=[e^a(u),e^a(u')]$, which maps $M_{n\times d}(E)$ into
$F^I\times F^I$. Let $j$ be $\epsilon^{-1}$ on $I_1$, so that $[0,b]=j_!b$ and
$[a,0]\circ j=0$. By part (1) of
Lemma~\ref{asm-theta-agree-evaluation-adapted-frame}, $Q\iota$ has rows $1$, $0$,
$0$, and $Qx$ and $Qy$ are the unit vectors at $x$ and $y$. So $w(\iota z)$ is
$z$ in the first $n_1$ rows and $0$ elsewhere, and $\beta=\delta=0$,
$\pi=z$ there; $w(xq^{\mathsf T})$ vanishes outside the row at $x$, so
$\beta=0$ and $\pi=0$ there; and $w(yq^{\mathsf T})$ vanishes outside the row at
$y$, so $\delta=0$ and $\pi=0$ there. This gives the three formulas.

(2) As $e^a$ and $e^M$ carry $\langle\ ,\ \rangle$ to $[\ ,\ ]$, $S$ is a
symplectic automorphism. Let $u\in\mathbb W(X,W)$ and $z\in\mathbb W(X_1,W)$.
As the columns of $\mathrm{fr}_H$ form a basis,
$u=\iota z'+xa^{\mathsf T}+yb^{\mathsf T}$ for some $z'$, $a$ and $b$, and by (1)
\[
S\,e^M(u)=e^a(u)=\bigl(j_!e^1_1(z')+e^a(yb^{\mathsf T})_1,\ j_!e^1_2(z')+e^a(xa^{\mathsf T})_2\bigr).
\]
If $e^M(u)$ lies over $e^1(z)$, then the first component is $j_!e^1_1(z)$;
reading it on $j(I_1)$ gives $e^1_1(z')=e^1_1(z)$, and off $j(I_1)$ it gives
$e^a(yb^{\mathsf T})_1=0$ there, hence everywhere. The second component composed
with $j$ is $e^1_2(z)$, which gives $e^1_2(z')=e^1_2(z)$. So $z'=z$,
$e^a(yb^{\mathsf T})=0$, $yb^{\mathsf T}=0$, and $u=\iota z+xa^{\mathsf T}$.
Conversely, $e^a(\iota z+xq^{\mathsf T})=(j_!e^1_1(z),j_!e^1_2(z)+e^a(xq^{\mathsf T})_2)$
lies over $e^1(z)$ by (1). The claims on $S\,e^M(yq^{\mathsf T})$ are (1).

The map $e$ maps Schwartz-Bruhat functions to Schwartz-Bruhat functions, by the
definition of an evaluation map; it keeps admissible forms admissible, and
$e(\Phi)\neq0$ for some Schwartz-Bruhat $\Phi$, by the facts (b) and (a) in the
proof of Lemma~\ref{asm-theta-agree-evaluation-unique}, whose proofs use only
the factorization of an evaluation map and the regularity of the models. For
the law, let $e^M(u)$ lie over $e^1(z)$, so $u=\iota z+xq^{\mathsf T}$. As
$h_X(\iota s,x)=0$ for every $s$, $\langle\iota z,xq^{\mathsf T}\rangle=0$, and by
the rule (R) of the proof of
Lemma~\ref{asm-theta-agree-split-inputs-expansion},
$\rho_M(u)=\rho_M(xq^{\mathsf T})\rho_M(\iota z)$. So for Schwartz-Bruhat $\Phi$,
$e(\rho_M(u)\Phi)=e(\rho_M(\iota z)\Phi)=\rho_1(z)e(\Phi)$, by the invariance
under $\mathcal J_x(H)$ and the definition of an evaluation map, where $\rho_1$
are the Schrödinger operators of $M_1$.

(3) The proof of Lemma~\ref{asm-theta-agree-evaluation-adapted}(1) checks
hypotheses (1) and (2) of Lemma~\ref{asm-theta-agree-evaluation-adapted-model}
for $\kappa=\{1,\ldots,n\}$, $X_1=X$, $\iota_W=\iota'$ and maps $\beta$,
$\delta$ and $\pi$ read off from $zQ'^{\mathsf T}$, where $Q'$ is the inverse of
$\mathrm{fr}_{H'}$; with the first two paragraphs of the proof of that lemma it
gives a bijection $\epsilon$ and an isomorphism $e^a$ with
$\langle u,u'\rangle=[e^a(u),e^a(u')]$,
$e^a(z\iota'^{\mathsf T})=([0,e^1_1(z)],[0,e^1_2(z)])$,
$e^a(qx'^{\mathsf T})=(0,[q_0,0])$ and $e^a(qy'^{\mathsf T})=([p_0,0],0)$ for some
$p_0$ and $q_0$ depending on $q$. With $j=\epsilon^{-1}$ on $I_1$ and
$S=e^a\circ(e^M)^{-1}$, the argument of (2), with
$u=z'\iota'^{\mathsf T}+ax'^{\mathsf T}+by'^{\mathsf T}$ and the invariance under
$\mathcal J_0(H')$, gives the claims; here
$\langle z\iota'^{\mathsf T},qx'^{\mathsf T}\rangle=0$ as $h_W(\iota's,x')=0$.
\end{proof}
% Lean: (1) is exists_adapted_index_hermitian (ExistsAdaptedIndexHermitian.lean): b, D and pi from
% incM Q * z, hH' (hTens_mul, frame_gram with eta = 1, hTens_plane_left, conj_pair), p and q
% through the basis of E over F, Phi, hPhi, hbij (bijective_of_symp, with the rational chart c0 of
% exists_rational_darboux_of_isGood), rationality through hM1.rat, and mul_iota,
% adapted_frame_mul_vecMulVec, mul_y for the three formulas. (2) is setX (DetNeHerm.lean) with
% fr_plane and fr_span for the basis, isEv_adm and EvNonzero.exists_mem_schA_ev_ne_zero for (b)
% and (a), and schr_add for (R). (3) is chartW and setW (FrOfFrame.lean): t z = z * Q'^T, b the
% column inr 1, D = DR (J_X) the column inr 0, pi the inl columns, splitR, and hrel through
% hF.one.

\begin{lemma}\label{asm-theta-agree-split-inputs-xvacuum}
In the situation of Lemma~\ref{asm-theta-agree-split-inputs-charts}(1), let $v$
be a finite place of $F$, let $\mathrm{vac}^1_v$ be the vacuum of $M_1$ at $v$,
and let $\phi$ be a vacuum at $v$ for the models of $(X,W)$ with the isomorphism
$e^a$ (Definition~\ref{asm-theta-agree-models}(3)). Then there is $c\neq0$ with
$\phi(j_!y)=c\,\mathrm{vac}^1_v(y)$ for every $y\in F_v^{I_1}$.
\end{lemma}
\leanref{Development/WeilBound/ThetaLifts/LeviDeterminants/MixedModels/AdaptedIndexVectorCompatibility.lean}{PositivityTheory.Mix.adapted_index_vector_compatibility}
\leanref{Development/WeilBound/ThetaLifts/LeviDeterminants/MixedModels/AdaptedIndexVectorCompatibility.lean}{PositivityTheory.Mix.hTens_frame}
\leanref{Development/WeilBound/ThetaLifts/LeviDeterminants/MixedModels/AdaptedIndexVectorCompatibility.lean}{PositivityTheory.Mix.chi_indep}
\begin{proof}
Let $\Lambda_v$ and $\Lambda^1_v$ be the groups of
Lemma~\ref{asm-theta-agree-evaluation-lattice} for $(X,W)$ and for $(X_1,W)$,
let $\rho_v$ and $\rho_{1,v}$ be the Schrödinger operators on the functions on
$F_v^I$ and $F_v^{I_1}$ for $\psi_v$, as in the proof of
Lemma~\ref{asm-theta-agree-evaluation-vacuum}, let $\theta$ and $\theta_1$ be the
elements in the definitions of the vacua $\phi$ and $\mathrm{vac}^1_v$, and
$\operatorname{Tr}=\operatorname{Tr}_{\mathbf A_E/\mathbf A}$. A subscript $v$
marks the component at $v$. The function $\phi_1(y)=\phi(j_!y)$ is
Schwartz-Bruhat on $F_v^{I_1}$, as a restriction to a coordinate subspace.

\emph{Eigenvalues.} Let $\gamma\in\Lambda^1_v$. As $\iota$ maps the lattice of
$X_1$ into that of $X$ (Definition~\ref{asm-theta-agree-planes}(1)),
$\iota\gamma\in\Lambda_v$; and as $\iota^{\mathsf T}J_Xc(\iota)=J_{X_1}$,
$h(\iota\gamma,\iota\gamma')=h_1(\gamma,\gamma')$, with $h_1$ the form $h$ for
$(X_1,W)$. So $\rho(e^a(\iota\gamma))$ acts at $v$ by
$\phi\mapsto\psi(\tfrac12\operatorname{Tr}h_1(\theta\gamma,\gamma))\phi$, and by
the first step of the proof of Lemma~\ref{asm-theta-agree-evaluation-vacuum}, for
$(X_1,W)$, this factor is
$\chi_1(\gamma)=\psi(\tfrac12\operatorname{Tr}h_1(\theta_1\gamma,\gamma))$, the
factor of $\mathrm{vac}^1_v$. By
Lemma~\ref{asm-theta-agree-split-inputs-charts}(1),
$e^a(\iota\gamma)=(j_!e^1_1(\gamma),j_!e^1_2(\gamma))$; with
$(a,b)=e^1(\gamma)_v$, for $y\in F_v^{I_1}$,
\[
(\rho_v(j_!a,j_!b)\phi)(j_!y)=\psi_v\bigl(j_!y\cdot j_!b+\tfrac12j_!a\cdot j_!b\bigr)\phi(j_!(y+a))
=(\rho_{1,v}(a,b)\phi_1)(y).
\]
So $\rho_{1,v}(e^1(\gamma)_v)\phi_1=\chi_1(\gamma)\phi_1$ for every
$\gamma\in\Lambda^1_v$.

\emph{The line.} The function $\mathrm{vac}^1_v$ satisfies the same equations and
is nonzero. The last step of the proof of
Lemma~\ref{asm-theta-agree-evaluation-vacuum}, for $M_1$, uses of its second
function only that it is Schwartz-Bruhat and satisfies these equations; it gives
$c$ with $\phi_1=c\,\mathrm{vac}^1_v$.

\emph{$c\neq0$.} Suppose $c=0$, so that $\phi$ vanishes at the vectors $j_!y$.
Choose $x_0$ with $\phi(x_0)\neq0$, and put $a_0=x_0-j_!(x_0\circ j)$, which
vanishes on $j(I_1)$. Let $u\in\mathbb W(X,W)$ have $e^a(u)=(\hat a_0,0)$, where
$\hat a_0$ is the vector at $v$ with component $a_0$; as $e^a$ is
$\mathbf A$-linear, $u$ is at $v$. We show $u\in\Lambda_v$, by
Lemma~\ref{asm-theta-agree-evaluation-lattice}(2). Let $\gamma\in\Lambda_v$ and
$w=Q\gamma$, with $Q$ as in the proof of
Lemma~\ref{asm-theta-agree-split-inputs-charts}(1). By
Lemma~\ref{asm-theta-agree-evaluation-adapted-frame}(4) for $(0,1,0)$ there is
$A$ with $AL_X\subset L_X$ and $QA=\Delta(0,1,0)Q$; so $\gamma_0=A\gamma$ lies in
$\Lambda_v$, and $\gamma_0=xw_x^{\mathsf T}$ with $w_x$ the row of $w$ at $x$. As
$\gamma=\iota w_1+\gamma_0+yw_y^{\mathsf T}$, with $w_1$ the first $n_1$ rows and
$w_y$ the row at $y$, Lemma~\ref{asm-theta-agree-split-inputs-charts}(1) gives
$e^a(\gamma)_2=j_!e^1_2(w_1)+e^a(\gamma_0)_2$; as $a_0$ vanishes on $j(I_1)$,
$\hat a_0\cdot e^a(\gamma)_2=\hat a_0\cdot e^a(\gamma_0)_2$. Now
$h(\gamma_0,\gamma_0)=0$, as $h_X(x,x)=0$, and $e^a(\gamma_0)_1=0$; so the vacuum
property of $\phi$ at $\gamma_0$ reads $\psi_v(x'\cdot b_0)\phi(x')=\phi(x')$ for
all $x'$, with $b_0=e^a(\gamma_0)_{2,v}$, and at $x'=x_0$ it gives
$\psi_v(x_0\cdot b_0)=1$. As $b_0$ vanishes on $j(I_1)$,
$x_0\cdot b_0=a_0\cdot b_0$. Hence
$\psi(\langle u,\gamma\rangle)=\psi(\hat a_0\cdot e^a(\gamma)_2)=\psi_v(a_0\cdot b_0)=1$,
and $\psi(\langle\gamma,u\rangle)=\psi(\langle u,\gamma\rangle)^{-1}=1$. So
$u\in\Lambda_v$, and $-u\in\Lambda_v$. The vacuum property of $\phi$ at $-u$,
with $e^a(-u)_v=(-a_0,0)$, gives at $x_0$
\[
\psi\bigl(\tfrac12\operatorname{Tr}h(\theta u,u)\bigr)\phi(x_0)=(\rho_v(-a_0,0)\phi)(x_0)=\phi(x_0-a_0)=\phi(j_!(x_0\circ j))=0,
\]
which contradicts $\phi(x_0)\neq0$, as $|\psi|=1$.
\end{proof}
% Lean: adapted_index_vector_compatibility (AdaptedIndexVectorCompatibility.lean): hs' for the
% restriction; the eigenvalues are lat_map (hH.low), e1 (hTens_frame, frame_iota), chi_indep and
% loc_eig with l1, l2, l3; the line is Vac.local_line for M1.co (the argument of local_line, as
% in Lemma asm-theta-agree-evaluation-vacuum); c /= 0 is the final block: a = x - extend (x o j),
% u = co.e.symm (liftV a, 0), huL through Vac.mem_latAtV with mask_mem 0 1 0, frame_alt, dec, h2,
% h3, h4, h5, then h6 at -u.

\begin{lemma}\label{asm-theta-agree-split-inputs-xexists}
Let $(X,W,M,\chi_X,\chi_W,\omega)$ be a Weil pair for a theta base $B$ of a CM
field $E$ (Definition~\ref{asm-theta-agree-triple}(1)), let $H$ be an adapted
plane in $X$ over a lattice space $X_1$, and let
$(X_1,W,M_1,\chi_X,\chi_W,\omega_1)$ be a Weil pair. Then there is an
evaluation map for $(M,M_1,\iota_H,1,\mathcal J_x(H))$
(Definition~\ref{asm-theta-agree-planes}(5)), with
$\mathcal J_x(H)=\{x_Hq^{\mathsf T}\}$ as in
Lemma~\ref{asm-theta-agree-split-inputs-xunique}.
\end{lemma}
\leanref{Development/WeilBound/ThetaLifts/LeviDeterminants/MixedModels/MixedModelPlanes/AdaptedPlanePairRestrictionStatement.lean}{PositivityTheory.Mix.T.adapted_plane_pair_restriction}
\leanref{Development/WeilBound/ThetaLifts/LeviDeterminants/MixedModels/MixedModelPlanes/AdaptedPlanePairRestriction.lean}{GlobalTheory.ThKMixXEx.exists_isEv_mixedX_m}
\leanref{Development/WeilBound/ThetaLifts/LeviDeterminants/MixedModels/DarbouxModelRationalLattice.lean}{PositivityTheory.Mix.darboux_model_rational_lattice}
\begin{proof}
As the two are Weil pairs, $X$ and $X_1$ are Hermitian, $W$ is skew-Hermitian,
and $M$ and $M_1$ are regular. Let $H=(x,y,\iota,\mathfrak b)$, let $n$, $n_1$
and $d$ be the ranks of $X$, $X_1$ and $W$, let $I$, $I_1$, $j$ and $e^a$ be as
in Lemma~\ref{asm-theta-agree-split-inputs-charts}(1), let $e^M$ and $e^1$ be the
isomorphisms of $M$ and $M_1$, and let $\rho_a(u)$ be the Schrödinger operator
for $e^a$ at $u$, as in Definition~\ref{asm-theta-agree-models}(3).

\emph{The transfer.} The map $e^a\circ(e^M)^{-1}$ is a symplectic automorphism.
The steps ``The operator'' and ``Vacua go to vacua'' of the proof of
Lemma~\ref{asm-theta-agree-evaluation-change} use of the second model only its
isomorphism; for $e^a$ in its place they give a linear map $t$ with
components $t_\infty$, a continuous automorphism of $\mathcal S(F_\infty^I)$, and
$t_v$, bijections of the Schwartz-Bruhat functions on $F_v^I$, such that $t$ maps
the Schwartz-Bruhat functions into themselves, $t(\rho_M(u)\Phi)=\rho_a(u)t(\Phi)$,
$t(f\otimes\bigotimes_v\phi_v)=t_\infty(f)\otimes\bigotimes_vt_v(\phi_v)$ when the
$\phi_v$ are Schwartz-Bruhat and almost all equal $\mathrm{vac}_v$, each
$\phi_v^a=t_v(\mathrm{vac}_v)$ is a vacuum at $v$ for $e^a$, and
$\phi_v^a=\mathbf 1_{\mathcal O_v^I}$ for almost all $v$.

\emph{Normalisation.} By Lemma~\ref{asm-theta-agree-split-inputs-xvacuum} there
is $c_v\neq0$ with $\phi_v^a(j_!y)=c_v\mathrm{vac}^1_v(y)$. For almost all $v$,
$\phi_v^a$ and $\mathrm{vac}^1_v$ are the indicator functions of
$\mathcal O_v^I$ and $\mathcal O_v^{I_1}$, as $M_1$ is regular, and the value at
$y=0$ gives $c_v=1$. Put $c=\prod_vc_v\neq0$.

\emph{The map.} Put $e(\Phi)=c^{-1}\operatorname{res}_j(t(\Phi))$, that is
$e(\Phi)(y)=c^{-1}t(\Phi)(j_!y)$. It maps Schwartz-Bruhat functions to
Schwartz-Bruhat functions by
Lemma~\ref{asm-theta-agree-split-inputs-heis-axioms}(1). For
$z\in\mathbb W(X_1,W)$, $e^a(\iota z)=(j_!e^1_1(z),j_!e^1_2(z))$, so the
computation of the proof of
Lemma~\ref{asm-theta-agree-split-inputs-heis-lines}(4) gives
$\operatorname{res}_j(\rho_a(\iota z)\Psi)=\rho_1(z)\operatorname{res}_j(\Psi)$, and
$e(\rho_M(\iota z)\Phi)=\rho_1(z)e(\Phi)$, where $\rho_1$ are the Schrödinger
operators of $M_1$. For $q\in\mathbf A_E^d$, $e^a(xq^{\mathsf T})=(0,b)$ with
$b\circ j=0$, so $(\rho_a(xq^{\mathsf T})\Psi)(j_!y)=\psi(j_!y\cdot b)\Psi(j_!y)=\Psi(j_!y)$,
and $e(\rho_M(xq^{\mathsf T})\Phi)=e(\Phi)$. Put
$e_\infty(f)=\operatorname{res}_j(t_\infty(f))$ and
$e_v=c_v^{-1}\operatorname{res}_j\circ t_v$. The map $e_\infty$ is continuous, and
it is nonzero: as $t_\infty$ is surjective, some $f$ has $t_\infty(f)=f_0$ with
$f_0(0)\neq0$, and $e_\infty(f)(0)=f_0(0)$. Each $e_v$ carries $\mathrm{vac}_v$
to $c_v^{-1}\phi^a_v\circ j_!=\mathrm{vac}^1_v$, and maps the Schwartz-Bruhat
functions onto those on $F_v^{I_1}$, as $t_v$ is bijective and every
Schwartz-Bruhat $\phi_1$ on $F_v^{I_1}$ is the restriction of
$[a,y]\mapsto\mathbf 1_{\mathcal O_v}(a)\phi_1(y)$, in the coordinates of
Lemma~\ref{asm-theta-agree-split-inputs-charts}(1). Finally, as only finitely
many $c_v$ differ from $1$ and restriction commutes with products,
$e(f\otimes\bigotimes_v\phi_v)=e_\infty(f)\otimes\bigotimes_ve_v(\phi_v)$ when
the $\phi_v$ are Schwartz-Bruhat and almost all equal $\mathrm{vac}_v$. So $e$ is
an evaluation map for $(M,M_1,\iota,1,\mathcal J_x(H))$.
\end{proof}
% Lean: exists_isEv_mixedX_m (AdaptedPlanePairRestriction.lean): j and co from
% exists_adapted_basis_hermitian = exists_adapted_index_hermitian (Lemma
% asm-theta-agree-split-inputs-charts(1)); e1, ei1, ev1 from darboux_coordinates_good_model =
% darboux_model_rational_lattice = darboux_model_rational_lattice_of adelic_estimates_holds
% (Theorem asm-theta-agree-evaluation-metaplectic for M.co.e.symm.trans co.e); c v from
% adapted_basis_hermitian_unique = adapted_index_vector_compatibility, hc1 with std_extend_m;
% e = (finprod c)^-1 . funLeft (extend j) . e1, with extend_by_zero_mem_schA, extend_dot_m,
% extend_add_m, extend_dot_zero_m, exists_restriction_extend_by_zero (R), a bump b with b 0 = 1,
% extend_by_zero_schwartzBruhat_map and finprod_mul_distrib.

\begin{lemma}\label{asm-theta-agree-split-inputs}
Let $E$ be a CM field, Galois over $\mathbf Q$, with a theta base $B$
(Definition~\ref{asm-theta-agree-models}(1)). Lattice
spaces are those of Definition~\ref{asm-theta-agree-doubling}(1); adapted
planes, Levi elements, unipotent groups, evaluation maps, lattice stabilizers
and cones those of Definition~\ref{asm-theta-agree-planes}; and isotropic
vectors, transporter families, Iwasawa subgroups, $K$-finiteness and cone
integrals those of Definition~\ref{asm-theta-agree-split-notions}. For an
adapted plane $H$ write $H=(x_H,y_H,\iota_H,\mathfrak b_H)$.
\begin{enumerate}
\item \emph{Unipotent groups.} Let $Y$ be a lattice space of rank $m$ that is
  Hermitian or skew-Hermitian of some signature, and $H$ an adapted plane in
  $Y$ over a lattice space $Y_1$, with unipotent group $N_H$. Then:
  \begin{enumerate}
  \item[(a)] $N_H$ is closed in $U_Y$, every Haar measure on $N_H$ is
    invariant under right translations, and there are a Haar measure on
    $N_H$ and a fundamental domain for $N_H\cap U_Y(F)$ in $N_H$ of measure
    $1$;
  \item[(b)] there is a compact $C\subset N_H$ such that every $u\in N_H$ has
    some $\gamma\in N_H\cap U_Y(F)$ with $\gamma u\in C$;
  \item[(c)] there is a continuous homomorphism
    $\lambda\colon\mathbf A_E^\times\times U_{Y_1}\to U_Y$ such that, for all
    $t$ and $g_1$, $\lambda(t,g_1)$ is the only Levi element for $(t,g_1)$;
  \item[(d)] every $u\in N_H$ has $ux_H=x_H$ and
    $uy_H=y_H+\iota_Ha+bx_H$ for some $a\in\mathbf A_E^{m-2}$ and
    $b\in\mathbf A_E$;
  \item[(e)] if $p\in U_Y$ and $t\in\mathbf A_E^\times$ have $px_H=tx_H$,
    then $p=um$ with $u\in N_H$ and $m$ a Levi element for $(t,g_1)$ for some
    $g_1\in U_{Y_1}$, and $g_1\in U_{Y_1}(F)$ if $p\in U_Y(F)$;
  \item[(f)] every Levi element for $(\xi,g_1)$ with $\xi\in E^\times$ and
    $g_1\in U_{Y_1}(F)$ lies in $U_Y(F)$;
  \item[(g)] $m^{-1}um\in N_H$ for every Levi element $m$ and every $u\in N_H$.
  \end{enumerate}
\item \emph{Finiteness.} Let $(X,W,M,\chi_X,\chi_W,\omega)$ be a Weil pair for
  $B$ (Definition~\ref{asm-theta-agree-triple}(1)) and $K$ the lattice
  stabilizer of $X$. Then every $\Phi\in\operatorname{PD}(M)$ is a
  Schwartz-Bruhat function and is $K$-finite for $\omega$.
\item \emph{Iwasawa decomposition.} Let $X$ be a Hermitian lattice space of some
  signature and $x\in\operatorname{Iso}(X)$. Then the lattice stabilizer of $X$
  is an Iwasawa subgroup for $x$.
\item \emph{Evaluation at a line of $W$.} Let $(X,W,M,\chi_X,\chi_W,\omega)$ be
  a Weil pair of signature $(r,s;r',s')$, with $X$ of rank $n$, let $H'$ be an
  adapted plane in $W$ over a lattice space $W_1$, let
  $(X,W_1,M_1,\chi_X,\chi_W,\omega_1)$ be a Weil pair of signature
  $(r,s;r_1,s_1)$, and let $e$ be an evaluation map for
  $(M,M_1,1,\iota_{H'},\mathcal J_0(H'))$. Then for $g\in U_X$, a Levi
  element $m\in U_W$ for $(t,h_1)$, $u\in N_{H'}$ and a Schwartz-Bruhat
  function $\Phi$,
  \[
  e(\omega(g,um)\Phi)=\chi_X(t)|t|^{n/2}\,\omega_1(g,h_1)\,e(\Phi).
  \]
\item \emph{Evaluation at two lines.} Let $H$ be an adapted plane in a lattice
  space $X$ over $X_1$ and $H'$ one in $W$ over $W_1$, with ranks $n$, $n_1$,
  $d$ and $d_1$, and let $(X,W,M,\chi_X,\chi_W,\omega)$ and
  $(X_1,W_1,M_0,\chi_X,\chi_W,\omega_0)$ be Weil pairs of signatures
  $(r+1,s+1;r'+1,s'+1)$ and $(r,s;r',s')$. Let $p\in U_X$ be a Levi element for
  $(t,g_1)$, $u\in N_H$, $p'\in U_W$ a Levi element for $(t',h_1)$,
  $u'\in N_{H'}$, and $\Phi$ a Schwartz-Bruhat function. If $e$ is an
  evaluation map for $(M,M_0,\iota_H,\iota_{H'},\mathcal J_\alpha(H,H'))$, then
  \[
  e(\omega(up,u'p')\Phi)=\chi_W(t)|t|^{d_1/2}\chi_X(t')|t'|^{n/2}\,\omega_0(g_1,h_1)\,e(\Phi);
  \]
  if $e$ is one for $(M,M_0,\iota_H,\iota_{H'},\mathcal J_\beta(H,H'))$, the
  same holds with the exponents $d/2$ and $n_1/2$ in place of $d_1/2$ and
  $n/2$.
\item \emph{Stability of full rank.} For every Weil pair
  $(X,W,M,\chi_X,\chi_W,\omega)$ for $B$, the operators $\omega(g,h)$ map
  $\mathcal S^\circ(M)$ (Definition~\ref{asm-theta-agree-models}(5)) into
  itself.
\item \emph{Full rank and evaluation.} Let $H$ be an adapted plane in a
  Hermitian lattice space $X$ of signature $(r,s)$ over $X_1$, $H'$ one in a
  skew-Hermitian lattice space $W$ of signature $(r',s')$ over $W_1$, and $M$
  and $M_0$ regular models for $(X,W)$ and $(X_1,W_1)$, with Schrödinger
  operators $\rho$ for $M$.
  \begin{enumerate}
  \item[(a)] If $e$ is an evaluation map for
    $(M,M_0,\iota_H,\iota_{H'},\mathcal J_\beta(H,H'))$ and
    $\Psi\in\mathcal S^\circ(M)$, then $e(\Psi)=0$ and
    $e(\rho(a\,y_Hx_{H'}^{\mathsf T})\Psi)\in\mathcal S^\circ(M_0)$ for every
    $a\in\mathbf A_E$.
  \item[(b)] Let $M_1$ be a regular model for $(X,W_1)$, and suppose that
    there are Weil pairs $(X,W_1,M_1,\chi_X,\chi_W,\omega_1)$ and
    $(X_1,W_1,M_0,\chi_X,\chi_W,\omega_0)$ in which $W_1$ has the same
    signature. If $e_0$ is an evaluation map for
    $(M,M_1,1,\iota_{H'},\mathcal J_0(H'))$, $e$ one for
    $(M,M_0,\iota_H,\iota_{H'},\mathcal J_\alpha(H,H'))$, and $\Psi$ is a
    Schwartz-Bruhat function with $e_0(\Psi)\in\mathcal S^\circ(M_1)$, then
    $e(\Psi)=0$.
  \end{enumerate}
\item \emph{Two expansions.} Let $H$ be an adapted plane in a Hermitian
  lattice space $X$ of signature $(r,s)$ over $X_1$, $H'$ one in a
  skew-Hermitian lattice space $W$ of signature $(r',s')$ over $W_1$, $M$ and
  $M_0$ regular models for $(X,W)$ and $(X_1,W_1)$, with Schrödinger operators
  $\rho$ for $M$, and $e_\alpha$ and $e_\beta$ evaluation maps for
  $(M,M_0,\iota_H,\iota_{H'},\mathcal J_\alpha(H,H'))$ and
  $(M,M_0,\iota_H,\iota_{H'},\mathcal J_\beta(H,H'))$. Then there is $c\neq0$
  such that for every Schwartz-Bruhat function $\Psi$ the sums
  \[
  \sum_{\xi\in E}\theta\bigl(e_\alpha(\rho(\xi\,x_Hy_{H'}^{\mathsf T})\Psi)\bigr)
  \quad\text{and}\quad
  \sum_{\xi\in E}\theta\bigl(e_\beta(\rho(\xi\,y_Hx_{H'}^{\mathsf T})\Psi)\bigr)
  \]
  of theta sums converge absolutely, and the first is $c$ times the second.
\item \emph{Sorting the theta kernel.} Let $X$, $W$ and $W_1$ be lattice spaces
  of ranks $n$, $d$ and $d_1$, $M$ a model for $(X,W)$ with Schrödinger
  operators $\rho$, $M_1$ a regular model for $(X,W_1)$, $H'$ an adapted plane
  in $W$ over $W_1$, and $e_0$ an evaluation map for
  $(M,M_1,1,\iota_{H'},\mathcal J_0(H'))$. In (a), (b) and (c), $\Psi$ is a
  Schwartz-Bruhat function on $\mathbf A^{n\times d}$, $\nu$ a Haar measure on
  $N_{H'}$ and $\mathcal F$ a fundamental domain for $N_{H'}\cap U_W(F)$ in
  $N_{H'}$ with $\nu(\mathcal F)=1$, and the Weil pairs
  $(X,W,M,\chi_X,\chi_W,\omega)$ have any signature. For $x\in E^n$ put
  \[
  T_\omega(x,\Psi)=\int_{\mathcal F}\theta\bigl(e_0(\rho(xy_{H'}^{\mathsf T})\,\omega(1,u)\Psi)\bigr)\,d\nu(u).
  \]
  \begin{enumerate}
  \item[(a)] There is $C_0\neq0$ such that for every Weil pair, every $\Psi$,
    $\nu$ and $\mathcal F$, the sum $\sum_{x\in E^n}C_0\,T_\omega(x,\Psi)$
    converges absolutely to $\int_{\mathcal F}\theta(\omega(1,u)\Psi)\,d\nu(u)$.
  \item[(b)] For every Weil pair, every $\Psi$, $\nu$ and $\mathcal F$, and
    every $x\in E^n$ with $h_X(x,x)\neq0$, $T_\omega(x,\Psi)=0$.
  \item[(c)] For every adapted plane $H$ in $X$ over a lattice space $X_1$,
    every regular model $M_2$ for $(X_1,W_1)$ and every evaluation map $e$
    for $(M,M_2,\iota_H,\iota_{H'},\mathcal J_\alpha(H,H'))$, there is
    $C\neq0$ such that
    $T_\omega(x_H,\Psi)=C\,\theta(e(\rho(x_Hy_{H'}^{\mathsf T})\Psi))$ for every
    Weil pair, every $\Psi$, $\nu$ and $\mathcal F$.
  \end{enumerate}
\item \emph{Unfolding along $\mathcal J_\alpha$.} Let $H$ be an adapted plane in
  a lattice space $X$ over $X_1$ and $H'$ one in $W$ over $W_1$, with ranks $n$,
  $n_1$, $d$ and $d_1$, let $M$ be a model for $(X,W)$ with Schrödinger
  operators $\rho$, measure $\mu_M$ and domain $\mathcal D$, let $M_0$ be a
  model for $(X_1,W_1)$ and $e$ an evaluation map for
  $(M,M_0,\iota_H,\iota_{H'},\mathcal J_\alpha(H,H'))$, let $r,s,r',s'\geq0$
  and $z\geq0$ be integers with $2n=d+2z$ and $\min(r,s)+2\leq z$, let
  $K\subset U_X$ be an Iwasawa subgroup for $x_H$, and let $\mu_K$ and $\mu$
  be Haar measures on $K$ and $\mathbf A_E^\times$. Then there is
  $\kappa\neq0$ such that the following holds for all Hecke characters
  $\chi_X,\chi_W$ and representations $\omega,\omega_0$ for which
  $(X,W,M,\chi_X,\chi_W,\omega)$ and $(X_1,W_1,M_0,\chi_X,\chi_W,\omega_0)$
  are Weil pairs of signatures $(r+1,s+1;r'+1,s'+1)$ and $(r,s;r',s')$, and
  every Schwartz-Bruhat function $\Phi$ that is $K$-finite for $\omega$. Let
  $I^{\mathrm{cone}}$ be the cone integral for
  $(M,\omega,M_0,\omega_0,e,K,\mu_K,\mu,u_\alpha,\chi_W^{-1},z)$ and
  $\operatorname{cone}$ the cone for
  $(M,\omega,e,K,\mu_K,\mu,u_\alpha,\chi_W^{-1},z)$.
  \begin{enumerate}
  \item[(a)] For every transporter family $\gamma$ for $x_H$ and every
    $h\in U_W$: for almost every $g\in\mathcal D$ the sum
    $S(g)=\sum_{v\in\operatorname{Iso}(X)}\theta(e(\rho(x_Hy_{H'}^{\mathsf T})\,\omega(\gamma(v)g,h)\Phi))$
    converges absolutely, $S$ is integrable on $\mathcal D$, and
    $\int_{\mathcal D}S\,d\mu_M=\kappa\,I^{\mathrm{cone}}(\omega(1,h)\Phi)(1)$.
  \item[(b)] For every Levi element $p\in U_W$ for $(t,h_1)$ and every
    $h\in U_W$,
    $I^{\mathrm{cone}}(\omega(1,ph)\Phi)(1)=(\chi_X/\chi_W)(t)\,|t|^{(d-n)/2}\,I^{\mathrm{cone}}(\omega(1,h)\Phi)(h_1)$.
  \item[(c)] $I^{\mathrm{cone}}(\omega(1,h)\Phi)(h_1)=I_{M_0,\omega_0}(\operatorname{cone}(\omega(1,h)\Phi))(h_1)$
    for $h\in U_W$ and $h_1\in U_{W_1}$.
  \end{enumerate}
\item \emph{Unfolding along $\mathcal J_\beta$.} The statement of (10) holds with
  $\mathcal J_\beta(H,H')$ in place of $\mathcal J_\alpha(H,H')$, the
  conditions $d+2=2n+2z$ and $2\leq z$ on $z$, functions
  $\Phi\in\mathcal S^\circ(M)$ that are $K$-finite for $\omega$, the map
  $u_\beta$ and the character $\chi_W$ in place of $u_\alpha$ and
  $\chi_W^{-1}$, and the sums
  $\sum_{v\in\operatorname{Iso}(X)}\theta(e(\rho(y_Hx_{H'}^{\mathsf T})\,\omega(\gamma(v)g,h)\Phi))$.
\item \emph{Evaluation at a line of $X$.} Let $(X,W,M,\chi_X,\chi_W,\omega)$ be
  a Weil pair of signature $(r,s;r',s')$, with $W$ of rank $d$, let $H$ be an
  adapted plane in $X$ over a lattice space $X_1$, and let
  $(X_1,W,M_1,\chi_X,\chi_W,\omega_1)$ be a Weil pair of signature
  $(r_1,s_1;r',s')$. Put $\mathcal J_x(H)=\{x_Hz^{\mathsf T}\}$, with $z$
  running over the vectors over $\mathbf A_E$ of length $d$. Then there is an
  evaluation map for $(M,M_1,\iota_H,1,\mathcal J_x(H))$, and for every such
  map $e$, every Levi element $p\in U_X$ for $(t,g_1)$, $u\in N_H$, $h\in U_W$
  and Schwartz-Bruhat function $\Phi$,
  \[
  e(\omega(up,h)\Phi)=\chi_W(t)|t|^{d/2}\,\omega_1(g_1,h)\,e(\Phi).
  \]
\end{enumerate}
\end{lemma}
\leanref{Development/WeilBound/ThetaLifts/CTKernel.lean}{PositivityTheory.ThKG.MixX}
\leanref{Development/WeilBound/ThetaLifts/LeviDeterminants/MixedModels/MixXClosed.lean}{PositivityTheory.Mix.mixX_closed}
\leanref{Development/WeilBound/ThetaLifts/SmallCones/ThetaUniq.lean}{PositivityTheory.ThKCT.UnipEx}
\leanref{Development/WeilBound/ThetaLifts/SmallCones/UnipEx.lean}{PositivityTheory.ThKCT.unipEx_of}
\leanref{Development/WeilBound/ThetaLifts/SmallCones/HodgeTateFromSandwich.lean}{PositivityTheory.ThKCT.StabDec}
\leanref{Development/WeilBound/ThetaLifts/SmallCones/StabDec.lean}{PositivityTheory.ThKCT.stabDec_of}
\leanref{Development/WeilBound/ThetaLifts/SmallCones/ThetaUniq.lean}{PositivityTheory.ThKCT.KFinData}
\leanref{Development/WeilBound/ThetaLifts/SmallCones/KFinDataOf.lean}{PositivityTheory.ThKCT.kFinData_of}
\leanref{Development/WeilBound/ThetaLifts/SmallCones/ThetaUniq.lean}{PositivityTheory.ThKCT.IwaEx}
\leanref{Development/WeilBound/ThetaLifts/SmallCones/IwaExOf.lean}{PositivityTheory.ThKCT.iwaEx_of}
\leanref{Development/WeilBound/ThetaLifts/CTKernel.lean}{PositivityTheory.ThKG.MixW}
\leanref{Development/WeilBound/ThetaLifts/LeviDeterminants/MixedModels/MixWClosed.lean}{PositivityTheory.Mix.mixW_closed}
\leanref{Development/WeilBound/ThetaLifts/SmallCones/ThetaUniq.lean}{PositivityTheory.ThKCT.MixAB}
\leanref{Development/WeilBound/ThetaLifts/SmallCones/ThetaUniq.lean}{PositivityTheory.ThKCT.Mix2}
\leanref{Development/WeilBound/ThetaLifts/SmallCones/MixABOf.lean}{PositivityTheory.ThKCT.Ker.mixAB_of}
\leanref{Development/WeilBound/ThetaLifts/SmallCones/ThetaUniq.lean}{PositivityTheory.ThKCT.SoStab}
\leanref{Development/WeilBound/ThetaLifts/SmallCones/SoStab.lean}{PositivityTheory.ThKCT.soStab_of}
\leanref{Development/WeilBound/ThetaLifts/SmallCones/ThetaUniq.lean}{PositivityTheory.ThKCT.SoEv}
\leanref{Development/WeilBound/ThetaLifts/SmallCones/FormSwapZero.lean}{PositivityTheory.ThKCT.soEv_of}
\leanref{Development/WeilBound/ThetaLifts/SmallCones/SoDefs.lean}{PositivityTheory.ThKCT.SoZeroA'}
\leanref{Development/WeilBound/ThetaLifts/SmallCones/SoZero.lean}{PositivityTheory.ThKCT.soZeroA'_of}
\leanref{Development/WeilBound/ThetaLifts/SmallCones/ThetaUniq.lean}{PositivityTheory.ThKCT.TwoExp}
\leanref{Development/WeilBound/ThetaLifts/SmallCones/TwoExp.lean}{PositivityTheory.ThKCT.Ker.twoExp_of}
\leanref{Development/WeilBound/ThetaLifts/SmallCones/Kernel.lean}{PositivityTheory.ThKCT.Ker.sum0}
\leanref{Development/WeilBound/ThetaLifts/SmallCones/ConstantTermRationalUnipotentInvariant.lean}{PositivityTheory.ThKCT.Ker.constant_term_anisotropic_integral_zero}
\leanref{Development/WeilBound/ThetaLifts/SmallCones/Kernel.lean}{PositivityTheory.ThKCT.Ker.constant_term_cone_integral_eq}
\leanref{Development/WeilBound/ThetaLifts/SmallCones/HodgeTateFromSandwich.lean}{PositivityTheory.ThKCT.ConeA}
\leanref{Development/WeilBound/ThetaLifts/SmallCones/ConeAB.lean}{PositivityTheory.ThKCT.coneA_of}
\leanref{Development/WeilBound/ThetaLifts/SmallCones/HodgeTateFromSandwich.lean}{PositivityTheory.ThKCT.ConeB}
\leanref{Development/WeilBound/ThetaLifts/SmallCones/ConeAB.lean}{PositivityTheory.ThKCT.coneB_of}
\begin{proof}
(1) Let $J$ and $J_1$ be the matrices of $Y$ and $Y_1$, write
$h(u,u')=u^{\mathsf T}J\,c(u')$, $x=x_H$, $y=y_H$ and $\iota=\iota_H$, and let
$\epsilon=1$ if $Y$ is Hermitian and $\epsilon=-1$ if it is skew-Hermitian, so
that $h(u',u)=\epsilon\,c(h(u,u'))$ and $h(y,x)=\epsilon$. Let $A=(\iota\ x\ y)$.
As in the proof of Lemma~\ref{asm-theta-agree-levi}, $A\in GL_m(E)$ and
$A^{\mathsf T}J\,c(A)=J^\natural=\operatorname{diag}\bigl(J_1,\left(\begin{smallmatrix}0&1\\\epsilon&0\end{smallmatrix}\right)\bigr)$.
As $\det J\neq0$ (Lemma~\ref{asm-theta-agree-convergence-signature}),
$\det J_1\neq0$, and $(x,y,\iota)$ is a plane in $Y$ over $Y_1$
(Definition~\ref{asm-theta-agree-end-notions}(1)). For $p\in U_Y$ put
$P=A^{-1}pA$, with rows and columns indexed by the columns of $A$; then
$P^{\mathsf T}J^\natural c(P)=J^\natural$. Let $e_x$ and $e_y$ be the unit
vectors at $x$ and $y$.

(a), (b) The proofs of Lemma~\ref{asm-theta-agree-end-unipotent}(1) and (3) use
of $W$ only that its matrix is invertible and $\epsilon$-Hermitian
(Lemma~\ref{asm-theta-agree-convergence-signature} and
Theorem~\ref{asm-theta-agree-convergence-siegel}(1)) and that the Gram matrix of
$A$ is $J^\natural$. With $h(y,x)=\epsilon$ in place of $-1$ they apply to $Y$
and $(x,y,\iota)$, and give (a), (b) and the description of $N_H$ in
Lemma~\ref{asm-theta-agree-end-unipotent}(3): $u\in U_Y$ lies in $N_H$ if and
only if $ux=x$ and $uz-z\in\mathbf A_Ex$ whenever $h(z,x)=0$.

(c) The computation in the proof of Lemma~\ref{asm-theta-agree-levi}, with
$\left(\begin{smallmatrix}0&1\\\epsilon&0\end{smallmatrix}\right)$ in place of
$\left(\begin{smallmatrix}0&1\\-1&0\end{smallmatrix}\right)$, shows that
$\lambda(t,g_1)=A\operatorname{diag}(g_1,t,c(t)^{-1})A^{-1}$ is a Levi element
for $(t,g_1)$. Two Levi elements for $(t,g_1)$ agree on the columns of $A$, so
they are equal. The product of Levi elements for $(t,g_1)$ and $(t',g'_1)$ is
one for $(tt',g_1g'_1)$, so $\lambda$ is a homomorphism. It is continuous, and
so is $(t,g_1)\mapsto\lambda(t,g_1)^{-1}=\lambda(t^{-1},g_1^{-1})$, as the
entries of $\lambda(t,g_1)$ are linear in those of $g_1$, $t$ and $c(t)^{-1}$.

(d) Let $u\in N_H$. Then $ux=x$, as every generator of $N_H$ fixes $x$, so the
column of $P=A^{-1}uA$ at $x$ is $e_x$. As the only nonzero entry of the row of
$J^\natural$ at $x$ is $1$ at $y$, the entry of $P^{\mathsf T}J^\natural c(P)$ at
$(x,y)$ is $c(P_{yy})$. It equals $J^\natural_{xy}=1$, so $P_{yy}=1$, and
$uy=APe_y=y+\iota a+bx$, with $a$ and $b$ the other entries of the column of $P$
at $y$.

(e) Let $px=tx$, so that the column of $P$ at $x$ is $te_x$. For an index $j$ of
the block of $\iota$, the entry of $P^{\mathsf T}J^\natural c(P)$ at $(j,x)$ is
$\epsilon P_{yj}c(t)$, which must be $J^\natural_{jx}=0$; so $P_{yj}=0$, and
$p\iota=\iota g_1+x\alpha^{\mathsf T}$, where $g_1$ is the block of $P$ on the
indices of $\iota$ and $\alpha_j=P_{xj}$. The block of
$P^{\mathsf T}J^\natural c(P)$ on these indices is then $g_1^{\mathsf T}J_1c(g_1)$,
so $g_1^{\mathsf T}J_1c(g_1)=J_1$; as $\det J_1\neq0$, $g_1$ is invertible and
$g_1\in U_{Y_1}$. Let $m=\lambda(t,g_1)$ and $u=pm^{-1}$. Then
$ux=t^{-1}px=x$ and $u\iota z=p\iota g_1^{-1}z=\iota z+(\alpha^{\mathsf T}g_1^{-1}z)\,x$
for every $z$, so $u$ is one of the generators of $N_H$
(Definition~\ref{asm-theta-agree-planes}(2)), and $p=um$. If $p\in U_Y(F)$, then
$P$ and its block $g_1$ have entries in $E$. This is the argument of Steps 5
and 6 of the proof of Lemma~\ref{asm-theta-agree-end-iwasawa}.

(f) By the uniqueness in (c), the Levi element for $(\xi,g_1)$ is
$A\operatorname{diag}(g_1,\xi,c(\xi)^{-1})A^{-1}$, which has entries in $E$ when
$\xi\in E^\times$ and $g_1\in U_{Y_1}(F)$.

(g) Let $m$ be a Levi element for $(t,g_1)$ and $u\in N_H$. Then
$m^{-1}x=t^{-1}x$, so $m^{-1}umx=t\,m^{-1}x=x$. For $z\in\mathbf A_E^{m-2}$,
$h(\iota g_1z,x)=0$, so $u\iota g_1z=\iota g_1z+bx$ for some $b$, by the
description of $N_H$ in (a); hence
$m^{-1}um\iota z=m^{-1}u\iota g_1z=\iota z+bt^{-1}x$. So $m^{-1}um$ is one of the
generators of $N_H$.

(2) As in the first paragraph of the proof of
Lemma~\ref{asm-theta-agree-end-automorphic}, every element of
$\operatorname{PD}(M)$ is a Schwartz-Bruhat function, as $M$ is regular. It is
$K$-finite by Lemma~\ref{asm-theta-agree-split-inputs-kfinite}.

(3) Let $J$ be the matrix of $X$, $n$ its rank, $K$ its lattice stabilizer and
$\hat L=\hat L_X$. By Lemma~\ref{asm-theta-agree-stabilizer}, $K$ is a compact
subgroup of $U_X$. Let $g\in U_X$.

\emph{The infinite places.} For an infinite place $w$, $M_w=\tau_w(J)$ is a
Hermitian matrix with Cartan involution $\theta_{X,w}$, and $v=\tau_w(x)$ is
nonzero with $v^{\mathsf T}M_w\bar v=\tau_w(x^{\mathsf T}J\,c(x))=0$. Steps 1 to 4
of the proof of Lemma~\ref{asm-theta-agree-end-iwasawa} use of $M_w$ only that it
is Hermitian with Cartan involution $\theta_{X,w}$, and of $v$ only that it is
nonzero and isotropic. So they give $q\in U_X$ with components $1$ at the
finite places and $t\in\mathbf A_E^\times$ with $qx=tx$, such that
$\pi_w(q^{-1}g)$ commutes with $\theta_{X,w}$ for every $w$.

\emph{Integrality.} Call $h\in U_X$ \emph{integral at} a finite place $v$ if
$h$ and $h^{-1}$ map every element of $\hat L$ at $v$ into $\hat L$. As $h$ acts
componentwise, it maps vectors at $v$ to vectors at $v$, and integrality at $v$
depends only on the component of $h$ at $v$: if $p\in U_X^{(v')}$ with
$v'\neq v$, then $ph$ is integral at $v$ if and only if $h$ is; and $h$ is
integral at $v$ if and only if $h_{(v)}$ is, where $h_{(v)}\in U_X^{(v)}$ has the
component of $h$ at $v$ and $1$ elsewhere. Every element of $K_X(v)$
(Definition~\ref{asm-theta-agree-finite-local-objects}(1)) is integral at $v$,
as the elements of $\hat L$ at $v$ lie in $L_v$, and $L_v\subset\hat L$. Let
$\widehat{\mathcal O}$ be the set of $a\in\mathbf A$ with $a_\infty=0$ and
$a_v\in\mathcal O_v$ for every $v$, and let $s_1,\ldots,s_r$ generate $L_X$
over $\mathbf Z$, so that $\hat L$ is the set of the $\sum_ja_js_j$ with
$a_j\in\widehat{\mathcal O}$ (proof of
Lemma~\ref{asm-theta-agree-stabilizer}).

Every $h$ is integral at almost all $v$. Call $w\in\mathbf A_E^n$
\emph{almost integral} if for almost all $v$, $aw\in\hat L$ for every
$a\in\widehat{\mathcal O}$ at $v$. The almost integral vectors form an
$\mathbf A$-submodule: if $c\in\mathbf A$, then $c_v\in\mathcal O_v$ for almost
all $v$, and for such $v$ and $a\in\widehat{\mathcal O}$ at $v$, $ac$ lies in
$\widehat{\mathcal O}$ and is at $v$. This submodule contains $L_X$, by the
definition of $\hat L$. It contains $tx$ for $t\in E$ and $x\in L_X$: choose an
integer $N\neq0$ with $Nt\in\mathcal O_E$; then $(Nt)x\in L_X$ and
$tx=N^{-1}\cdot(Nt)x$. As $L_X$ spans $E^n$ over $E$, and $E^n$ spans
$\mathbf A_E^n$ over $\mathbf A$, every vector is almost integral. Apply this to
the $2r$ vectors $hs_j$ and $h^{-1}s_j$: for almost all $v$, $a\,hs_j$ and
$a\,h^{-1}s_j$ lie in $\hat L$ for every $j$ and every $a\in\widehat{\mathcal O}$
at $v$. For such $v$, let $z\in\hat L$ be at $v$, and let $\epsilon_v\in\mathbf A$
be $1$ at $v$ and $0$ at the other places. Then $z=\sum_ja_js_j$ with
$a_j\in\widehat{\mathcal O}$, and $z=\epsilon_vz=\sum_j(\epsilon_va_j)s_j$ with
$\epsilon_va_j\in\widehat{\mathcal O}$ at $v$; so
$hz=\sum_j(\epsilon_va_j)\,hs_j\in\hat L$, and likewise $h^{-1}z\in\hat L$.

A vector $z$ with $z_\infty=0$ whose
part at each $v$ lies in $\hat L$ lies in $\hat L$: if the part at $v$ is
$\sum_ja_{v,j}s_j$, then $z=\sum_ja_js_j$ with $a_{j,v}$ the component of
$a_{v,j}$ at $v$. So if $h$ is integral at every finite place, then
$h\hat L\subset\hat L$ and $h^{-1}\hat L\subset\hat L$, since the part at $v$ of
$hz$ is $h$ applied to the part at $v$ of $z$, which lies in $\hat L$.

\emph{The finite places.} We show, by induction on the size of a finite set $S$
of finite places, that every $h\in U_X$ integral at every $v\notin S$ is $p\,h'$
with $p\in U_X$ having components $1$ at the infinite places, $px=t'x$ for some
$t'\in\mathbf A_E^\times$, and $h'$ integral at every finite place. If $S$ is
empty, take $p=1$. Otherwise let $v\in S$. By
Theorem~\ref{asm-theta-agree-local-mult-unramified-iwasawa},
$h_{(v)}=p_1k$ with $p_1\in U_X^{(v)}$, $p_1x=t_1x$ and $k\in K_X(v)$. Then
$p_1^{-1}h$ is integral at $v$, as its component at $v$ is that of $k$, and at
every $v'\notin S$, as $p_1$ is at $v$. The induction hypothesis for
$S\setminus\{v\}$ gives $p_1^{-1}h=p_2h'$, and $p=p_1p_2$ works.

\emph{Conclusion.} Apply this to $q^{-1}g$ and the finite set of places at
which it is not integral: $q^{-1}g=pk$ with $p$ as above and $k$ integral at
every finite place. Then $\pi_w(k)=\pi_w(q^{-1}g)$ commutes with $\theta_{X,w}$
for every $w$, and $k\hat L\subset\hat L$ and $k^{-1}\hat L\subset\hat L$; so
$k\in K$. Finally $g=(qp)k$ and $qpx=t't\,x$.

(4) As $\omega$ is a representation,
$\omega(g,um)=\omega(g,1)\omega(1,u)\omega(1,m)$, and these operators preserve
the Schwartz-Bruhat functions. So (4) follows from three laws, for
Schwartz-Bruhat $\Phi$:
\begin{enumerate}
\item[(4a)] $e(\omega(g,1)\Phi)=\omega_1(g,1)e(\Phi)$ for $g\in U_X$;
\item[(4b)] $e(\omega(1,u)\Phi)=e(\Phi)$ for $u\in N_{H'}$;
\item[(4c)] $e(\omega(1,m)\Phi)=\chi_X(t)|t|^{n/2}\,\omega_1(1,h_1)e(\Phi)$ for a
  Levi element $m\in U_W$ for $(t,h_1)$.
\end{enumerate}
Indeed, they give
$e(\omega(g,um)\Phi)=\omega_1(g,1)e(\omega(1,m)\Phi)=\chi_X(t)|t|^{n/2}\,\omega_1(g,1)\omega_1(1,h_1)e(\Phi)$,
which is the formula of (4).

\emph{The laws.} As the two are Weil pairs, $X$ is Hermitian, $W$ and $W_1$
are skew-Hermitian, $M$ and $M_1$ are regular, and the matrices of $X$ and
$W_1$ are invertible (Lemma~\ref{asm-theta-agree-convergence-signature}). Write
$H'=(x',y',\iota',\mathfrak b')$, and let $S$ and $j$ be as in
Lemma~\ref{asm-theta-agree-split-inputs-charts}(3) for $e$, so that $e$ is an
$(S,j)$-evaluation. By Definition~\ref{asm-theta-agree-split-inputs-side},
$\Sigma=(X,W,e^M,\chi_X,\omega(1,\cdot))$ and
$\Sigma_1=(X,W_1,e^{M_1},\chi_X,\omega_1(1,\cdot))$ are Weil sides, and so are
$\Sigma'=(W,X,e^M\circ\tau,\chi_W,\omega(\cdot,1))$ and
$\Sigma'_1=(W_1,X,e^{M_1}\circ\tau,\chi_W,\omega_1(\cdot,1))$, with
$\tau(u)=u^{\mathsf T}$ and $e^M$, $e^{M_1}$ the isomorphisms of $M$ and $M_1$.

(4a) Apply Proposition~\ref{asm-theta-agree-split-inputs-seesaw}(1) to
$\Sigma'$ and $\Sigma'_1$, with $Y=X$ and $\epsilon=1$, to $S$, $j$ and $e$, and
to $\iota'$, $x'$ and $y'$. Its hypothesis (i) holds: $e^M(U^{\mathsf T})$ lies
over $e^{M_1}(Z^{\mathsf T})$ if and only if
$U^{\mathsf T}=Z^{\mathsf T}\iota'^{\mathsf T}+qx'^{\mathsf T}$ for some $q$,
that is $U=\iota'Z+x'q^{\mathsf T}$. Hypothesis (ii) is the last claim of
Lemma~\ref{asm-theta-agree-split-inputs-charts}(3), as
$e^M(\tau(y'q^{\mathsf T}))=e^M(qy'^{\mathsf T})$; (iii) holds as $y'\neq0$; and
(iv) holds as the columns of $\mathrm{fr}_{H'}$ form a basis
(Lemma~\ref{asm-theta-agree-evaluation-adapted-frame}). The proposition gives
(4a).

(4b) Let $u$ be one of the generators of $N_{H'}$, so $ux'=x'$ and
$u\iota's=\iota's+a(s)x'$ with $a(s)\in\mathbf A_E$ for every vector $s$. With
$r$ the row with entries $a(s_1),\ldots,a(s_{d_1})$ at the unit vectors $s_k$,
$u\iota'=\iota'+x'r$. Proposition~\ref{asm-theta-agree-split-inputs-seesaw}(2),
for $\Sigma$, $\Sigma_1$, $H'$ and $\epsilon=-1$, with the relation of
Lemma~\ref{asm-theta-agree-split-inputs-charts}(3), $m=u$, $g=1$ and $t=1$,
gives $e(\omega(1,u)\Phi)=e(\Phi)$. The set of $u\in U_W$ with
$e(\omega(1,u)\Phi)=e(\Phi)$ for every Schwartz-Bruhat $\Phi$ is a subgroup, as
the $\omega(1,u)$ preserve the Schwartz-Bruhat functions; it contains the
generators, hence $N_{H'}$.

(4c) A Levi element $m$ for $(t,h_1)$ has $mx'=tx'$ and $m\iota'=\iota'h_1$.
Proposition~\ref{asm-theta-agree-split-inputs-seesaw}(2), as in (4b) with
$g=h_1$ and $r=0$, gives (4c).

(5) Let $H=(x,y,\iota,\mathfrak b)$ and $H'=(x',y',\iota',\mathfrak b')$. The
space $X$ is Hermitian of signature $(r+1,s+1)$ and $W$ skew-Hermitian of
signature $(r'+1,s'+1)$, and $n_1=n-2$, $d_1=d-2$. Fix $p,u,p',u'$ and $\Phi$
as in (5); then $\omega(up,u'p')\Phi$ is Schwartz-Bruhat.

\emph{The case $\mathcal J_\alpha$.} By Lemma~\ref{asm-theta-agree-split-low},
$W_1$ is skew-Hermitian of signature $(r',s')$. By
Lemma~\ref{asm-theta-agree-tower-model} there is a regular model $M_1$ for
$(X,W_1)$, and by Theorem~\ref{asm-theta-agree-tower-weil} a Weil
representation $\omega_1$ for it with characters $(\chi_X,\chi_W)$, as
$\chi_W|_{\mathbf A^\times}=\eta^d=\eta^{d_1}$; so
$(X,W_1,M_1,\chi_X,\chi_W,\omega_1)$ is a Weil pair of signature
$(r+1,s+1;r',s')$. By Lemma~\ref{asm-theta-agree-evaluation} there is an
evaluation map $e_0$ for $(M,M_1,1,\iota',\mathcal J_0(H'))$, and by (12), for
the Weil pairs of $\omega_1$ and $\omega_0$, there is an evaluation map $e_1$
for $(M_1,M_0,\iota,1,\mathcal J_x(H))$. Every element of
$\mathcal J_\alpha(H,H')$ lies in $\mathcal J_0(H')$ or is
$x(\iota'z)^{\mathsf T}=1\cdot(xz^{\mathsf T})\cdot\iota'^{\mathsf T}$ with
$xz^{\mathsf T}\in\mathcal J_x(H)$. By
Lemma~\ref{asm-theta-agree-split-inputs-expansion}(5), with the polar system of
Lemma~\ref{asm-theta-agree-split-inputs-systems}(2) for $(X,W_1)$ and
$(X_1,W_1)$, $e_1\circ e_0$ is an evaluation map for
$(M,M_0,\iota,\iota',\mathcal J_\alpha(H,H'))$. By
Lemma~\ref{asm-theta-agree-split-inputs-systems}(5), $e$ and $e_1\circ e_0$ are
evaluation maps for $(M,M_0,\iota,\iota',\mathcal J_\alpha^+(H,H'))$, and by
Lemma~\ref{asm-theta-agree-split-inputs-expansion}(3), with the polar system of
Lemma~\ref{asm-theta-agree-split-inputs-systems}(3), there is $c\neq0$ with
$e=c\,e_1e_0$ on the Schwartz-Bruhat functions. By (4) with $g=up$, and by (12)
for $\omega_1$, whose second space $W_1$ has rank $d_1$, with $h=h_1$,
\[
e(\omega(up,u'p')\Phi)=c\,\chi_X(t')|t'|^{n/2}\,e_1(\omega_1(up,h_1)e_0\Phi)
=\chi_W(t)|t|^{d_1/2}\chi_X(t')|t'|^{n/2}\,\omega_0(g_1,h_1)\,e(\Phi).
\]

\emph{The case $\mathcal J_\beta$.} By Lemma~\ref{asm-theta-agree-split-low},
$X_1$ is Hermitian of signature $(r,s)$, and as above there is a Weil pair
$(X_1,W,M_1,\chi_X,\chi_W,\omega_1)$ of signature $(r,s;r'+1,s'+1)$, as
$\chi_X|_{\mathbf A^\times}=\eta^n=\eta^{n_1}$. By (12), for the Weil pairs of
$\omega$ and $\omega_1$, there is an evaluation map $e_0$ for
$(M,M_1,\iota,1,\mathcal J_x(H))$, and by Lemma~\ref{asm-theta-agree-evaluation},
for the Weil pair of $\omega_1$, one $e_1$ for
$(M_1,M_0,1,\iota',\mathcal J_0(H'))$. Every element of
$\mathcal J_\beta(H,H')$ lies in $\mathcal J_x(H)$ or is
$(\iota z)x'^{\mathsf T}=\iota\cdot(zx'^{\mathsf T})\cdot1$ with
$zx'^{\mathsf T}\in\mathcal J_0(H')$. By
Lemma~\ref{asm-theta-agree-split-inputs-expansion}(5), with the polar system of
Lemma~\ref{asm-theta-agree-split-inputs-systems}(1) for $(X_1,W)$ and
$(X_1,W_1)$, $e_1\circ e_0$ is an evaluation map for
$(M,M_0,\iota,\iota',\mathcal J_\beta(H,H'))$, and as above, with
Lemma~\ref{asm-theta-agree-split-inputs-systems}(4) and (5), $e=c\,e_1e_0$ for
some $c\neq0$. By (12) with $h=u'p'$, and by (4) for $\omega_1$, whose first
space $X_1$ has rank $n_1$, with $g=g_1$,
\[
e(\omega(up,u'p')\Phi)=c\,\chi_W(t)|t|^{d/2}\,e_1(\omega_1(g_1,u'p')e_0\Phi)
=\chi_W(t)|t|^{d/2}\chi_X(t')|t'|^{n_1/2}\,\omega_0(g_1,h_1)\,e(\Phi).
\]

(6) The proof of Lemma~\ref{asm-theta-agree-coinvariants-inputs}(4) uses of
the triple only that $M'$ is a regular model for a Hermitian and a
skew-Hermitian lattice space whose matrices are invertible, and that $\omega'$
is a Weil representation for $M'$. For a Weil pair
$(X,W,M,\chi_X,\chi_W,\omega)$ these hold for $M$ and $\omega$, the matrices
being invertible by Lemma~\ref{asm-theta-agree-convergence-signature}; the same
argument gives (6).

(7) Let $H=(x,y,\iota,\mathfrak b)$ and $H'=(x',y',\iota',\mathfrak b')$, and
let $\rho$ and $\rho_0$ be the Schrödinger operators of $M$ and $M_0$. Write
$v=v_1$, $t=\tau_1$, $\bar c$ and $\varepsilon$ as in
Lemma~\ref{asm-theta-agree-split-inputs-kill}, and use $h$, $h_X$ and $h_W$ as
in the proof of Lemma~\ref{asm-theta-agree-split-inputs-systems}. By
Lemma~\ref{asm-theta-agree-split-low}, $X_1$ is Hermitian and $W_1$
skew-Hermitian.

(a) Let $\Psi\in\mathcal S^\circ(M)$. As $xz^{\mathsf T}\in\mathcal J_\beta(H,H')$
for every $z$, Lemma~\ref{asm-theta-agree-split-inputs-kill}(4) gives
$e(\Psi)=0$.

Let $\lambda\in\mathbf A_E$ and $s_0=\lambda\,yx'^{\mathsf T}$; the function
$e(\rho(s_0)\Psi)$ is Schwartz-Bruhat. By
Lemma~\ref{asm-theta-agree-coinvariants-slices} it suffices to show that for
every $z$ the function $\operatorname{res}_z(e(\rho(s_0)\Psi))$ has full rank at
$t$. Let $\iota_0$ be a linear model of the second kind for $(t,1,1)$ for $M_0$
(Lemma~\ref{asm-theta-agree-coinvariants-dictionary}(1)), let
$Y\in M_{n_1\times d_1}(F_v)$ have rank different from $n_1$, and let $x_0$ be the
element of $\mathbb W(X_1,W_1)$ at $v$ with $\pi_t(x_0)=0$ and
$\operatorname{co}(x_0)=Y$ (Lemma~\ref{asm-theta-agree-coinvariants-eigen}(1)).
Put $C_0=t(J_{X_1})\pi_{t\circ c}(x_0)t(J_{W_1})^{\mathsf T}=-Y$, and choose
$c_0\in F_v^{n_1}$ nonzero with $c_0^{\mathsf T}C_0=0$. We apply
Lemma~\ref{asm-theta-agree-split-inputs-kill}(3) to
$L(\Phi)=\iota_0\bigl(\operatorname{res}_z(e(\rho(s_0)\Phi))\bigr)(Y)$ and the
set $S$ of the triples $(x,x',0)$, $(x,\iota'q,0)$, $(\iota p,x',0)$ and
$(\iota p,\iota'q,t(p)^{\mathsf T}C_0t(q))$, for $p\in E^{n_1}$ and
$q\in E^{d_1}$.

\emph{The first condition.} Let $(p',q',\beta)\in S$, $c\in F_v$ and
$u=\bar c\varepsilon p'q'^{\mathsf T}$. In each case $h_W(x',q')=0$, so
$h(s_0,u)=\lambda\,c(\bar c\varepsilon)h_X(y,p')h_W(x',q')=0$, and
$\rho(s_0)\rho(u)=\rho(u)\rho(s_0)$ by (R) of the proof of
Lemma~\ref{asm-theta-agree-split-inputs-expansion}. If $p'=x$, then
$u=x(\bar c\varepsilon q')^{\mathsf T}\in\mathcal J_\beta(H,H')$; if
$(p',q')=(\iota p,x')$, then
$u=(\iota\,\bar c\varepsilon p)x'^{\mathsf T}\in\mathcal J_\beta(H,H')$; in both
cases $\beta=0$ and $L(\rho(u)\Phi)=L(\Phi)$. If $(p',q')=(\iota p,\iota'q)$,
then $u=\iota w\iota'^{\mathsf T}$ with $w=\bar c\varepsilon pq^{\mathsf T}$,
which is at $v$ with $\pi_{t\circ c}(w)=0$. So
$e(\rho(u)\rho(s_0)\Phi)=\rho_0(w)e(\rho(s_0)\Phi)$, and by
Lemma~\ref{asm-theta-agree-coinvariants-eigen}(2) and the definition of a
linear model of the second kind,
$L(\rho(u)\Phi)=\psi(\langle x_0,w\rangle_0)L(\Phi)$. By
Lemma~\ref{asm-theta-agree-split-inputs-kill}(2) for $(X_1,W_1)$,
$\langle x_0,w\rangle_0=\bar c\,\langle x_0,\varepsilon pq^{\mathsf T}\rangle_0=\overline{c\,t(p)^{\mathsf T}C_0t(q)}=\overline{\beta c}$.

\emph{The second condition.} Let $C\in M_{n\times d}(F_v)$ satisfy
$t(p')^{\mathsf T}C\,t(q')=\beta$ for $(p',q',\beta)\in S$, and put $r_1=t(x)$,
$r_2=t(\iota)c_0$ and $w_0=t(y')$. First, $r_2^{\mathsf T}Cw=0$ whenever
$w\in F_v^d$ and $w_1\in F_v^{d_1}$ satisfy
$(t(\iota)t(p))^{\mathsf T}Cw=t(p)^{\mathsf T}C_0w_1$ for all $p\in E^{n_1}$:
taking for $p$ the unit vectors, $t(\iota)^{\mathsf T}Cw=C_0w_1$, and
$r_2^{\mathsf T}Cw=c_0^{\mathsf T}C_0w_1=0$. Every unit vector of $E^d$ is
$\alpha x'+\iota'q+\beta'y'$ with $\alpha,\beta'\in E$ and $q\in E^{d_1}$
(proof of Lemma~\ref{asm-theta-agree-split-low}); with
$w=t(\alpha)t(x')+t(\iota'q)$, the corresponding unit vector of $F_v^d$ is
$w+t(\beta')w_0$. The triples $(x,x',0)$ and $(x,\iota'q,0)$ give
$r_1^{\mathsf T}Cw=0$, and the triples $(\iota p,x',0)$ and
$(\iota p,\iota'q,t(p)^{\mathsf T}C_0t(q))$ give, by the first remark with
$w_1=0$ and with $w_1=t(q)$, $r_2^{\mathsf T}Cw=0$. Hence, for
$a_1,a_2\in F_v$ with $a_1r_1^{\mathsf T}Cw_0+a_2r_2^{\mathsf T}Cw_0=0$, the row
$(a_1r_1+a_2r_2)^{\mathsf T}C$ vanishes on every unit vector, so it is $0$. If
$r_1^{\mathsf T}Cw_0=0$, take $(a_1,a_2)=(1,0)$, as $r_1\neq0$. Otherwise take
$a_1=r_2^{\mathsf T}Cw_0$ and $a_2=-r_1^{\mathsf T}Cw_0\neq0$; then
$a_1r_1+a_2r_2\neq0$: multiplying $a_1r_1+a_2r_2$ on the right by
$t(J_Xc(\iota))$, as $t(x)^{\mathsf T}t(J_Xc(\iota))=0$ and
$t(\iota)^{\mathsf T}t(J_Xc(\iota))=t(J_{X_1})$, gives
$a_2c_0^{\mathsf T}t(J_{X_1})$, which is nonzero as $c_0\neq0$ and $J_{X_1}$ is
invertible.

So Lemma~\ref{asm-theta-agree-split-inputs-kill}(3) gives $L(\Psi)=0$, that is,
$\iota_0(\operatorname{res}_z(e(\rho(s_0)\Psi)))(Y)=0$ for every $Y$ of rank
different from $n_1$.

(b) By (12), for the two Weil pairs of (b), there is an evaluation map $e_X$ for
$(M_1,M_0,\iota,1,\mathcal J_x(H))$, with $x=x_H$ and vectors of length $d_1$. As
in (9)(c), by Lemma~\ref{asm-theta-agree-split-inputs-expansion}(5), with the
polar system of Lemma~\ref{asm-theta-agree-split-inputs-systems}(2) for $(X,W_1)$
and $(X_1,W_1)$, $e_X\circ e_0$ is an evaluation map for
$(M,M_0,\iota,\iota',\mathcal J_\alpha(H,H'))$, and
Lemma~\ref{asm-theta-agree-split-inputs-systems}(3) and (5) and
Lemma~\ref{asm-theta-agree-split-inputs-expansion}(3) give $c\neq0$ with
$e=c\,e_Xe_0$ on the Schwartz-Bruhat functions. As $e_X$ is invariant under the
$\rho_1(xz^{\mathsf T})$, where $\rho_1$ are the Schrödinger operators of $M_1$,
Lemma~\ref{asm-theta-agree-split-inputs-kill}(4) for $M_1$ shows that $e_X$
vanishes on $\mathcal S^\circ(M_1)$. So $e(\Psi)=c\,e_X(e_0(\Psi))=0$.

(8) Write $H=(x,y,\iota,\mathfrak b)$ and $H'=(x',y',\iota',\mathfrak b')$, let
$\rho$ be the Schrödinger operators of $M$, and use $h$, $h_X$ and $h_W$ as in
the proof of Lemma~\ref{asm-theta-agree-split-inputs-systems}. Put
$U_\alpha(v,w)=vy'^{\mathsf T}+y(\iota'w)^{\mathsf T}$ for
$(v,w)\in E^n\times E^{d_1}$ and
$U_\beta(v,w)=yv^{\mathsf T}+(\iota w)y'^{\mathsf T}$ for
$(v,w)\in E^d\times E^{n_1}$. For an $F$-basis $(b_j,b'_j)$ of
$E^n\times E^{d_1}$, $\xi\mapsto(\sum_j\xi_jb_j,\sum_j\xi_jb'_j)$ is a bijection
from $F^m$ onto $E^n\times E^{d_1}$, and $u(\xi)$ is $U_\alpha$ of its image. So
Lemma~\ref{asm-theta-agree-split-inputs-systems}(3) and (5) and
Lemma~\ref{asm-theta-agree-split-inputs-expansion}(2) give $C\neq0$ with
\[
\sum_{t\in E^n\times E^{d_1}}C\,\theta\bigl(e_\alpha(\rho(U_\alpha(t))\Psi)\bigr)=\theta(\Psi)
\]
for every Schwartz-Bruhat $\Psi$, the sum converging absolutely; in the same
way, with Lemma~\ref{asm-theta-agree-split-inputs-systems}(4), there is
$C'\neq0$ with the same identity for $e_\beta$, $U_\beta$ and
$t\in E^d\times E^{n_1}$. As $x$ and $x'$ are nonzero, $\xi\mapsto(\xi x,0)$ and
$\xi\mapsto(\xi x',0)$ are injective on $E$, and
$U_\alpha(\xi x,0)=\xi\,xy'^{\mathsf T}$ and $U_\beta(\xi x',0)=\xi\,yx'^{\mathsf T}$.
So the two sums of (8) are parts of absolutely convergent sums, and converge
absolutely.

\emph{Characters.} Let
$V=\mathcal J_\alpha^+(H,H')\cap\mathcal J_\beta^+(H,H')$, an
$\mathbf A_E$-submodule of $\mathbb W(X,W)$. By
Lemma~\ref{asm-theta-agree-split-inputs-systems}(5), $e_\alpha$ and $e_\beta$
are invariant under $\rho(a)$ for $a\in V$. By the rule
$\rho(U)\rho(a)=\psi(\langle U,a\rangle)\rho(a)\rho(U)$ (proof of
Lemma~\ref{asm-theta-agree-split-inputs-expansion}), the two identities applied
to $\rho(a)\Psi$ give, for every $a\in V$,
\[
\theta(\rho(a)\Psi)=\sum_tc_t\,\chi_t(a)=\sum_sc'_s\,\chi'_s(a),\tag{T}
\]
with $c_t=C\theta(e_\alpha(\rho(U_\alpha(t))\Psi))$,
$\chi_t(a)=\psi(\langle U_\alpha(t),a\rangle)$,
$c'_s=C'\theta(e_\beta(\rho(U_\beta(s))\Psi))$ and
$\chi'_s(a)=\psi(\langle U_\beta(s),a\rangle)$. These are characters of $V$ with
values of absolute value $1$. We show that $\chi_t=1$ if and only if
$t=(\xi x,0)$ for some $\xi\in E$, and $\chi'_s=1$ if and only if
$s=(\xi x',0)$.

If $t=(\xi x,0)$, then $U_\alpha(t)=\xi xy'^{\mathsf T}$, and every $a\in V$ is
$xv^{\mathsf T}+(\iota z)x'^{\mathsf T}$, as $a\in\mathcal J_\beta^+(H,H')$; so
$h(U_\alpha(t),a)=\xi\bigl(h_X(x,x)h_W(y',v)+h_X(x,\iota z)h_W(y',x')\bigr)=0$
and $\chi_t=1$. Likewise $\chi'_{(\xi x',0)}=1$, writing
$a=vx'^{\mathsf T}+x(\iota'z)^{\mathsf T}$. Conversely, let $U$ and $q$ be
rational with $q\in V$ and $h(U,q)\neq0$. Choose $\lambda\in E$ with
$\tau=-\operatorname{Tr}_{E/F}(\lambda h(U,q))\neq0$, as the trace form is
nondegenerate, and $b\in\mathbf A$ with $\psi(b)\neq1$, and put
$a=b\tau^{-1}c(\lambda)q\in V$. Then $h(U,a)=b\tau^{-1}\lambda h(U,q)$ and
$\langle U,a\rangle=b$, so $\psi(\langle U,a\rangle)\neq1$. The rational
matrices $xx'^{\mathsf T}$, $(\iota z)x'^{\mathsf T}$ and
$x(\iota'z')^{\mathsf T}$ lie in $V$. So if $\chi_t=1$ for $t=(v,w)$, then
$U=U_\alpha(t)$ has
\[
0=h(U,xx'^{\mathsf T})=-h_X(v,x),\quad
0=h(U,(\iota z)x'^{\mathsf T})=-h_X(v,\iota z),\quad
0=h(U,x(\iota'z')^{\mathsf T})=h_{W_1}(w,z')
\]
for all $z$ and $z'$. So $w=0$; and writing $v=sx+\iota v_1+s'y$,
$h_X(v,x)=s'$ and $h_X(v,\iota z)=h_{X_1}(v_1,z)$, so $s'=0$, $v_1=0$ and
$v=sx$. In the same way, if $\chi'_s=1$ for $s=(v,w)$, then
$h(U_\beta(s),xx'^{\mathsf T})=h_W(v,x')$,
$h(U_\beta(s),x(\iota'z')^{\mathsf T})=h_W(v,\iota'z')$ and
$h(U_\beta(s),(\iota z)x'^{\mathsf T})=-h_{X_1}(w,z)$ vanish, so $w=0$ and
$v\in Ex'$.

\emph{Comparison.} Put $P=\sum_{\xi\in E}c_{(\xi x,0)}$ and
$P'=\sum_{\xi\in E}c'_{(\xi x',0)}$. By (T), for $a\in V$ the sums of
$c_t\chi_t(a)$ over the other $t$ and of $c'_s\chi'_s(a)$ over the other $s$
converge absolutely to $\theta(\rho(a)\Psi)-P$ and $\theta(\rho(a)\Psi)-P'$.
Apply Lemma~\ref{asm-theta-agree-split-inputs-characters} to $A=V$, the index
set consisting of a point $i_0$, the other $t$ and the other $s$, the
characters $1$, $\chi_t$ and $\chi'_s$, and the coefficients $P-P'$, $c_t$ and
$-c'_s$: the sum is
$(P-P')+(\theta(\rho(a)\Psi)-P)-(\theta(\rho(a)\Psi)-P')=0$. So $P=P'$, that
is,
$C\sum_\xi\theta(e_\alpha(\rho(\xi xy'^{\mathsf T})\Psi))=C'\sum_\xi\theta(e_\beta(\rho(\xi yx'^{\mathsf T})\Psi))$,
and $c=C'/C$ works.

(9) Write $H'=(x',y',\iota',\mathfrak b')$, let $\rho_1$ be the Schrödinger
operators of $M_1$, and use $h$, $h_X$ and $h_W$ as in the proof of
Lemma~\ref{asm-theta-agree-split-inputs-systems}, and the rule (R) of the proof
of Lemma~\ref{asm-theta-agree-split-inputs-expansion}. Let
$(X,W,M,\chi_X,\chi_W,\omega)$ be a Weil pair, $\Psi$ a Schwartz-Bruhat
function and $\nu$, $\mathcal F$ as in (9). For $\xi\in E^n$ and $u\in N_{H'}$
put $f_\xi(u)=\theta(e_0(\rho(\xi y'^{\mathsf T})\omega(1,u)\Psi))$, so that
$T_\omega(\xi,\Psi)=\int_{\mathcal F}f_\xi\,d\nu$.

\emph{Common facts.} As in the proof of
Lemma~\ref{asm-theta-agree-split-kernel}, there is a Weil pair
$(X,W_1,M_1,\chi_X,\chi_W,\omega_1)$, and (4) with $g=1$ and the Levi element
$1$ for $(1,1)$ gives $e_0(\omega(1,u)\Phi)=e_0(\Phi)$ for $u\in N_{H'}$ and
Schwartz-Bruhat $\Phi$. As $\omega(1,u)\rho(w)=\rho(wu^{\mathsf T})\omega(1,u)$
(Definition~\ref{asm-theta-agree-models}(6)), this gives
\[
e_0\bigl(\rho(p)\,\omega(1,u)\Phi\bigr)=e_0\bigl(\rho(p(u^{-1})^{\mathsf T})\Phi\bigr)
\qquad(p\in\mathbb W(X,W)).\tag{U}
\]
By (1)(d) for $u^{-1}$, $u^{-1}y'=y'+\iota'b+ax'$ with $b$ and $a$ adelic, and
rational if $u$ is, as the matrix $P$ of the proof of (1)(d) is then rational.
By (R) and the definition of an evaluation map, for $j\in\mathcal J_0(H')$,
$q\in\mathbb W(X,W_1)$ and $p\in\mathbb W(X,W)$,
\[
e_0(\rho(j+p)\Phi)=\psi(-\tfrac12\langle j,p\rangle)\,e_0(\rho(p)\Phi),\qquad
e_0(\rho(q\iota'^{\mathsf T}+p)\Phi)=\psi(-\tfrac12\langle q\iota'^{\mathsf T},p\rangle)\,\rho_1(q)\,e_0(\rho(p)\Phi).\tag{J}
\]
As
$\xi(u^{-1}y')^{\mathsf T}=(a\xi)x'^{\mathsf T}+(\xi b^{\mathsf T})\iota'^{\mathsf T}+\xi y'^{\mathsf T}$,
(U) and (J) give
\[
f_\xi(u)=\psi(r)\,\theta\bigl(\rho_1(\xi b^{\mathsf T})\,e_0(\rho(\xi y'^{\mathsf T})\Psi)\bigr)\tag{F}
\]
for some $r\in\mathbf A$, with $r\in F$ if $u$ is rational.

\emph{Eichler elements.} Let $z\in\mathbf A_E^{d_1}$ and $\alpha\in\mathbf A_E$
with $h_W(\iota'z,\iota'z)+\alpha-c(\alpha)=0$, and let $\mathcal E_{z,\alpha}$ be the
linear map $w\mapsto w-h_W(w,x')(\iota'z+\alpha x')-h_W(w,\iota'z)x'$ of
$\mathbf A_E^d$. Then $\mathcal E_{z,\alpha}\in N_{H'}$, and
\[
\mathcal E_{z,\alpha}x'=x',\qquad \mathcal E_{z,\alpha}y'=y'+\iota'z+\alpha x',\qquad
\mathcal E_{z,\alpha}\iota'v=\iota'v-h_W(\iota'v,\iota'z)x'.
\]
The three values follow from $h_W(x',x')=h_W(x',\iota'z)=0$,
$h_W(y',x')=-1$, $h_W(y',\iota'z)=0$ and $h_W(\iota'v,x')=0$. Expanding with
$h_W(w',w)=-c(h_W(w,w'))$,
\[
h_W(\mathcal E_{z,\alpha}w,\mathcal E_{z,\alpha}w')=h_W(w,w')+h_W(w,x')\,c(h_W(w',x'))\bigl(\alpha-c(\alpha)+h_W(\iota'z,\iota'z)\bigr)=h_W(w,w'),
\]
and $\mathcal E_{z,\alpha}\mathcal E_{-z,\alpha'}=1$ for $\alpha'=-\alpha-h_W(\iota'z,\iota'z)$, as
both sides agree on $x'$, $y'$ and the $\iota'v$. So $\mathcal E_{z,\alpha}\in U_W$. As
it fixes $x'$ and moves each $\iota'v$ by a multiple of $x'$, it is one of the
generators of $N_{H'}$ (Definition~\ref{asm-theta-agree-planes}(2)).

\emph{(a)} Let $(b_j)$ be an $F$-basis of $E^n$. By
Lemma~\ref{asm-theta-agree-split-inputs-systems}(1) and
Lemma~\ref{asm-theta-agree-split-inputs-expansion}(2), as
$\xi\mapsto\sum_j\xi_jb_j$ is a bijection from $F^m$ onto $E^n$, there is
$C_0\neq0$, depending only on $e_0$, such that
$\sum_{\xi\in E^n}C_0\,\theta(e_0(\rho(\xi y'^{\mathsf T})\Phi))$ converges
absolutely to $\theta(\Phi)$ for every Schwartz-Bruhat $\Phi$. By (1)(a), $\nu$
is invariant under right translations, and $\nu(\mathcal F)=1$. Each $f_\xi$ is
continuous: by Lemma~\ref{asm-theta-agree-split-inputs-expansion}(1),
$\Phi\mapsto\theta(e_0(\rho(\xi y'^{\mathsf T})\Phi))$ agrees on the
Schwartz-Bruhat functions with an admissible form $\lambda$, and
$u\mapsto\lambda(\omega(1,u)\Psi)$ is continuous by
Definition~\ref{asm-theta-agree-models}(6)(c). By (F) and
Lemma~\ref{asm-theta-agree-split-inputs-expansion}(4), there is
$\beta\colon E^n\to[0,\infty)$ with $\sum_\xi\beta(\xi)<\infty$ and
$|f_\xi(u)|\leq\beta(\xi)$ for all $\xi$ and $u$. For each $u$,
$\sum_\xi C_0f_\xi(u)=\theta(\omega(1,u)\Psi)$, absolutely. As
$\nu(\mathcal F)=1$, the $f_\xi$ are integrable on $\mathcal F$,
$\sum_\xi\int_{\mathcal F}|C_0f_\xi|\,d\nu\leq|C_0|\sum_\xi\beta(\xi)$, and by
dominated convergence $\sum_\xi C_0T_\omega(\xi,\Psi)$ converges absolutely to
$\int_{\mathcal F}\theta(\omega(1,u)\Psi)\,d\nu(u)$.

\emph{(b)} Let $\xi\in E^n$ with $h_X(\xi,\xi)\neq0$. By (F), $f_\xi$ is
bounded and continuous, so integrable on $\mathcal F$. For
$\gamma\in N_{H'}\cap U_W(F)$ and $u\in N_{H'}$, $f_\xi(\gamma u)=f_\xi(u)$:
apply (F) with $\omega(1,u)\Psi$ in place of $\Psi$ and $\gamma$ in place of
$u$; then $\psi(r)=1$, and $b$ is rational, so
$\theta(\rho_1(\xi b^{\mathsf T})\Psi_1)=\theta(\Psi_1)$ by
Theorem~\ref{asm-theta-agree-evaluation-heisenberg}(1), as $M_1$ is regular.
Choose $a_0\in\mathbf A$ with $\psi(a_0)\neq1$ and put
$\alpha=a_0h_X(\xi,\xi)^{-1}$; as $X$ is Hermitian, $h_X(\xi,\xi)\in F$, so
$\alpha\in\mathbf A$, $\alpha=c(\alpha)$, and $\mathcal E_{0,\alpha}$ is defined. Put
$g=\mathcal E_{0,\alpha}^{-1}\in N_{H'}$. For $u\in N_{H'}$,
$(ug)^{-1}y'=\mathcal E_{0,\alpha}u^{-1}y'=u^{-1}y'+\alpha x'$, as $\mathcal E_{0,\alpha}$ fixes
$x'$ and the $\iota'v$. By (U) and (J), with $p=\xi(u^{-1}y')^{\mathsf T}$,
\[
f_\xi(ug)=\theta\bigl(e_0(\rho(\alpha\xi x'^{\mathsf T}+p)\Psi)\bigr)=\psi(-\tfrac12\langle\alpha\xi x'^{\mathsf T},p\rangle)\,f_\xi(u),
\]
and
$\langle\alpha\xi x'^{\mathsf T},p\rangle=-\operatorname{Tr}\bigl(\alpha h_X(\xi,\xi)h_W(x',u^{-1}y')\bigr)=-\operatorname{Tr}(a_0)=-2a_0$,
as $h_W(x',u^{-1}y')=h_W(x',y')=1$. So $f_\xi(ug)=\psi(a_0)f_\xi(u)$. As $\nu$
is invariant under right translations,
$\int_{\mathcal F}f_\xi(ug)\,d\nu(u)=\int_{\mathcal Fg}f_\xi\,d\nu$, and
$\mathcal Fg$ is again a fundamental domain for the group
$\Gamma=N_{H'}\cap U_W(F)$, which is countable
(Theorem~\ref{asm-theta-agree-convergence-siegel}(1)); as $f_\xi$ is invariant
under $\Gamma$ on the left, $\int_{\mathcal Fg}f_\xi\,d\nu=\int_{\mathcal F}f_\xi\,d\nu$.
So $\psi(a_0)T_\omega(\xi,\Psi)=T_\omega(\xi,\Psi)$, and $T_\omega(\xi,\Psi)=0$.

\emph{(c)} Let $H=(x,y,\iota,\mathfrak b)$, $M_2$ and $e$ be given. If there is
no Weil pair $(X,W,M,\ldots)$ there is nothing to prove; otherwise fix one, with
characters $(\chi^0_X,\chi^0_W)$. By Lemma~\ref{asm-theta-agree-split-low},
$X_1$ is Hermitian and $W_1$ skew-Hermitian, and as above there are Weil pairs
$(X,W_1,M_1,\chi^0_X,\chi^0_W,\omega^0_1)$ and, by
Theorem~\ref{asm-theta-agree-tower-weil},
$(X_1,W_1,M_2,\chi^0_X,\chi^0_W,\omega^0_2)$. By (12) for these two Weil pairs
there is an evaluation map $e_X$ for $(M_1,M_2,\iota,1,\mathcal J_x(H))$. By
Lemma~\ref{asm-theta-agree-split-inputs-systems}(2), for $(X,W_1)$ and
$(X_1,W_1)$, and Lemma~\ref{asm-theta-agree-split-inputs-expansion}(2), there is
$C\neq0$ such that
$\sum_{v\in E^{d_1}}C\,\theta(e_X(\rho_1(yv^{\mathsf T})\Psi_1))$ converges
absolutely to $\theta(\Psi_1)$ for every Schwartz-Bruhat function $\Psi_1$ on
$\mathbf A^{n\times d_1}$. Every element of $\mathcal J_\alpha(H,H')$ lies in
$\mathcal J_0(H')$ or is $1\cdot(xz^{\mathsf T})\cdot\iota'^{\mathsf T}$ with
$xz^{\mathsf T}\in\mathcal J_x(H)$. By
Lemma~\ref{asm-theta-agree-split-inputs-expansion}(5), $e_Xe_0$ is an evaluation
map for $(M,M_2,\iota,\iota',\mathcal J_\alpha(H,H'))$, and by
Lemma~\ref{asm-theta-agree-split-inputs-systems}(3) and (5) and
Lemma~\ref{asm-theta-agree-split-inputs-expansion}(3), there is $c\neq0$ with
$e_Xe_0=c\,e$ on the Schwartz-Bruhat functions. The constants $C$ and $c$ do
not depend on the Weil pair; we show that (c) holds with $Cc$.

Let $(X,W,M,\chi_X,\chi_W,\omega)$ be a Weil pair. By the common facts,
$e(\omega(1,u)\Phi)=c^{-1}e_Xe_0(\omega(1,u)\Phi)=e(\Phi)$ for $u\in N_{H'}$,
so (U) holds with $e$ in place of $e_0$. For $v\in E^{d_1}$ put
$U_v=xy'^{\mathsf T}+y(\iota'v)^{\mathsf T}$ and
$g_v(u)=\theta(e(\rho(U_v)\omega(1,u)\Psi))$. For $u\in N_{H'}$ let
$\Psi'=\rho(xy'^{\mathsf T})\omega(1,u)\Psi$. As
$\rho_1(yv^{\mathsf T})e_0(\Psi')=e_0(\rho(y(\iota'v)^{\mathsf T})\Psi')$ and
$\langle y(\iota'v)^{\mathsf T},xy'^{\mathsf T}\rangle=-\operatorname{Tr}h_X(y,x)h_W(\iota'v,y')=0$,
the expansion of $e_X$ applied to $e_0(\Psi')$ gives
\[
f_x(u)=\sum_{v\in E^{d_1}}Cc\,g_v(u),\tag{S}
\]
absolutely. By (U) for $e$, $g_v(u)=\theta(e(\rho(U_v(u^{-1})^{\mathsf T})\Psi))$.
Write $u^{-1}y'=y'+\iota'b+ax'$ as above, and $u^{-1}\iota'v=\iota'v+a'x'$, by
the description of $N_{H'}$ in (1)(a), as $h_W(\iota'v,x')=0$. Then
$U_v(u^{-1})^{\mathsf T}=U_v+j_u$ with
$j_u=(ax+a'y)x'^{\mathsf T}+x(\iota'b)^{\mathsf T}\in\mathcal J_\alpha^+(H,H')$,
so by Lemma~\ref{asm-theta-agree-split-inputs-systems}(5) and (R),
\[
g_v(u)=\psi(-\tfrac12\langle j_u,U_v\rangle)\,g_v(1).\tag{G}
\]
In particular $|g_v(u)|=|g_v(1)|$, and $\sum_v|Cc\,g_v(1)|<\infty$ by (S) at
$u=1$. Each $g_v$ is continuous, by
Lemma~\ref{asm-theta-agree-split-inputs-expansion}(1) as in (a). By dominated
convergence, $T_\omega(x,\Psi)=\sum_vCc\int_{\mathcal F}g_v\,d\nu$.

For $v=0$, $a'=0$ and $j_u=x(ax'+\iota'b)^{\mathsf T}$, so
$\langle j_u,U_0\rangle=-\operatorname{Tr}h_X(x,x)h_W(ax'+\iota'b,y')=0$; so
$g_0$ is constant and $\int_{\mathcal F}g_0\,d\nu=\theta(e(\rho(xy'^{\mathsf T})\Psi))$.

Let $v\neq0$. For rational $\gamma\in N_{H'}$ and $u\in N_{H'}$, (G) with
$\omega(1,u)\Psi$ in place of $\Psi$ gives $g_v(\gamma u)=g_v(u)$, as
$j_\gamma$ and $U_v$ are rational. As the matrix of $W_1$ is invertible, there
is $z_0\in E^{d_1}$ with $h_W(\iota'z_0,\iota'v)\neq0$. Choose
$a_0\in\mathbf A$ with $\psi(a_0)\neq1$, and put
$z=\frac12a_0h_W(\iota'z_0,\iota'v)^{-1}z_0$, so that
$h_W(\iota'z,\iota'v)=\frac12a_0$, and $\alpha=-\frac12h_W(\iota'z,\iota'z)$. As
$c(h_W(\iota'z,\iota'z))=-h_W(\iota'z,\iota'z)$, $\mathcal E_{z,\alpha}$ is defined; put
$g=\mathcal E_{z,\alpha}^{-1}\in N_{H'}$. Then
$j_g=x(\iota'z)^{\mathsf T}+\alpha\,xx'^{\mathsf T}-h_W(\iota'v,\iota'z)\,yx'^{\mathsf T}$,
and
\[
h(j_g,U_v)=h_X(x,y)h_W(\iota'z,\iota'v)-h_W(\iota'v,\iota'z)h_X(y,x)h_W(x',y')
=\tfrac12a_0+c(\tfrac12a_0)=a_0,
\]
the other terms containing $h_X(x,x)$, $h_W(x',\iota'v)$ or $h_X(y,y)$. So
$\psi(-\frac12\langle j_g,U_v\rangle)=\psi(a_0)$, and (G) for $\omega(1,g)\Psi$
and for $\Psi$ gives
$g_v(ug)=\psi(-\frac12\langle j_u,U_v\rangle)g_v(g)=\psi(a_0)g_v(u)$. As in (b),
$\int_{\mathcal F}g_v\,d\nu=0$. Hence
$T_\omega(x,\Psi)=Cc\,\theta(e(\rho(xy'^{\mathsf T})\Psi))$.

(10) Let $H=(x,y,\iota,\mathfrak b)$ and $H'=(x',y',\iota',\mathfrak b')$, so
that $n_1=n-2$ and $d_1=d-2$, let $\rho$ be the Schrödinger operators of $M$,
and let $\lambda$ and $\lambda'$ be the homomorphisms of (1)(c) for $H$ and
$H'$. If no characters and representations give the two Weil pairs of (10),
put $\kappa=1$. Otherwise fix such data. Then $X$ is Hermitian and $M$ and
$M_0$ are regular, so their measures $\mu_M$ and $\mu_0$ are Haar measures with
fundamental domains $\mathcal D$ and $\mathcal D_0$, and
Lemma~\ref{asm-theta-agree-split-inputs-cone-unfold} for $X$, $H$, $K$,
$\mu_K$, $\mu$, $(\mu_M,\mathcal D)$ and $(\mu_0,\mathcal D_0)$ gives $\kappa>0$,
which depends only on these data. Now let $\chi_X,\chi_W,\omega,\omega_0$ give
the two Weil pairs, and let $\Phi$ be Schwartz-Bruhat and $K$-finite for
$\omega$. For a Schwartz-Bruhat function $\Psi$:
\begin{enumerate}
\item[(L)] if $p\in U_W$ is a Levi element for $(t,h_1)$ and
  $l\in\mathbf A_E^\times$, then
  $e(\rho(l\,xy'^{\mathsf T})\omega(1,p)\Psi)=\chi_X(t)|t|^{n/2}\,\omega_0(1,h_1)\,e(\rho(l\,c(t)\,xy'^{\mathsf T})\Psi)$;
  indeed $\rho(l\,xy'^{\mathsf T})\omega(1,p)\Psi=\omega(1,p)\rho(l\,x(p^{-1}y')^{\mathsf T})\Psi$
  (Definition~\ref{asm-theta-agree-models}(6)) and $p^{-1}y'=c(t)y'$, and (5)
  with the Levi element $1\in U_X$ for $(1,1)$ gives the factor;
\item[(T)] if $p\in U_X$ is a Levi element for $(l^{-1},g_1)$ and $p'\in U_W$
  one for $(1,h_1)$, then
  $|l^{-1}|^{1-n}\,\theta(e(\rho(xy'^{\mathsf T})\omega(p,p')\Psi))=\theta_{\omega_0}(e(\rho(u_\alpha(l))\Psi))(g_1,h_1)\,\chi_W(l)^{-1}|l|^z$;
  indeed (A) of the proof of Lemma~\ref{asm-theta-agree-split-shift} gives
  $e(\rho(xy'^{\mathsf T})\omega(p,p')\Psi)=\chi_W(l)^{-1}|l|^{-d_1/2}\omega_0(g_1,h_1)e(\rho(u_\alpha(l))\Psi)$,
  and $n-1-d_1/2=z$ as $2n=d+2z$;
\item[(N)] $\chi_W(c(t))=\chi_W(t)^{-1}$ and $|c(t)|=|t|$ for
  $t\in\mathbf A_E^\times$, as $t\,c(t)=N_{E/F}(t)$ and
  $\chi_W(N_{E/F}(t))=\eta^d(N_{E/F}(t))=1$.
\end{enumerate}

\emph{Integrability.} By Lemma~\ref{asm-theta-agree-split-inputs-expansion}(1)
the form $\Psi\mapsto\theta(e(\rho(xy'^{\mathsf T})\Psi))$ agrees on the
Schwartz-Bruhat functions with an admissible form, so
$(g,h)\mapsto\theta(e(\rho(xy'^{\mathsf T})\omega(g,h)\Psi))$ is continuous
(Definition~\ref{asm-theta-agree-models}(6)(c)). By (T) with
$p=\lambda(l^{-1},g)$ and $p'=\lambda'(1,h_1)$, the function
$(l,g)\mapsto\theta_{\omega_0}(e(\rho(u_\alpha(l))\Psi))(g,h_1)\,\chi_W(l)^{-1}|l|^z$
is continuous on $\mathbf A_E^\times\times U_{X_1}$. Apply
Lemma~\ref{asm-theta-agree-split-inputs-cone-line} to $e$, $v_0=xy'^{\mathsf T}$,
$z_0=yx'^{\mathsf T}$, $\tau=1$ and $\chi_W^{-1}$, so that
$\tau(l)v_0=u_\alpha(l)$: for $a\in\mathbf A_E$,
$az_0=(ay)x'^{\mathsf T}\in\mathcal J_0(H')\subset\mathcal J_\alpha(H,H')$, and
$h(z_0,v_0)=h_X(y,x)h_W(x',y')=1$. As $2n=d+2z\geq d+2\min(r,s)+4$, the second
Weil pair has $d_1+2\min(r,s)<2n_1$, and
Lemma~\ref{asm-theta-agree-convergence-dominated} for it gives the hypothesis of
Lemma~\ref{asm-theta-agree-split-inputs-cone-line}(1) with $\mathcal F$ the set
of all Schwartz-Bruhat functions, for $\omega_0$. With the continuity and
Tonelli's theorem, part (1) of that lemma shows that the function above is
integrable on $\mathbf A_E^\times\times\mathcal D_0$; and part (2) holds. So
Lemma~\ref{asm-theta-agree-split-inputs-cone-exchange} applies to the space
$\mathcal S$ of Schwartz-Bruhat functions, $u_\alpha$, $\chi_W^{-1}$ and $z$.

\emph{(c)} For $h\in U_W$, the function $\omega(1,h)\Phi$ is Schwartz-Bruhat and
$K$-finite, as $\omega(k,1)\omega(1,h)\Phi=\omega(1,h)\omega(k,1)\Phi$. So (c) is
Lemma~\ref{asm-theta-agree-split-inputs-cone-exchange}(2).

\emph{(b)} Let $p$ be a Levi element for $(t,h_1)$ and $h\in U_W$, and put
$\Psi_k=\omega(k,h)\Phi$, so that $\omega(k,1)\omega(1,ph)\Phi=\omega(1,p)\Psi_k$.
As $I_{M_0,\omega_0}(\omega_0(1,h_1)\Psi')(1)=I_{M_0,\omega_0}(\Psi')(h_1)$, (L)
gives
\[
I_{M_0,\omega_0}\bigl(e(\rho(u_\alpha(l))\omega(1,p)\Psi_k)\bigr)(1)=\chi_X(t)|t|^{n/2}\,I_{M_0,\omega_0}\bigl(e(\rho(u_\alpha(l\,c(t)))\Psi_k)\bigr)(h_1).
\]
Multiply by $\chi_W(l)^{-1}|l|^z$ and integrate over $l$; the substitution
$l\mapsto l\,c(t)^{-1}$, which preserves $\mu$, and (N) turn the right side into
$\chi_X(t)\chi_W(t)^{-1}|t|^{n/2-z}$ times the corresponding integral for
$h_1$. As $n/2-z=(d-n)/2$, integrating over $K$ gives (b).

\emph{(a)} Let $\gamma$ be a transporter family for $x$ and $h\in U_W$. Put
$\Phi'=\omega(1,h)\Phi$ and
$f(g)=\theta(e(\rho(xy'^{\mathsf T})\omega(g,h)\Phi))$, so that
$S(g)=\sum_{v\in\operatorname{Iso}(X)}f(\gamma(v)g)$. The function $f$ is
continuous, $f(ug)=f(g)$ for $u\in N_H$ by the invariance under $N_H$ in the
proof of Lemma~\ref{asm-theta-agree-split-shift}, and $f(\delta g)=f(g)$ for
$\delta\in U_X(F)$ with $\delta x=x$ by part (1) of that lemma. For $k\in K$,
$l\in\mathbf A_E^\times$ and $g_1\in U_{X_1}$, (T) for $\Psi=\omega(k,1)\Phi'$,
$p=\lambda(l^{-1},g_1)$ and $p'=1$ gives
\[
f(\lambda(l^{-1},g_1)k)\,|l^{-1}|^{1-n}=\theta_{\omega_0}\bigl(e(\rho(u_\alpha(l))\omega(k,1)\Phi')\bigr)(g_1,1)\,\chi_W(l)^{-1}|l|^z.\tag{H}
\]
As $\mu$ is invariant under $l\mapsto l^{-1}$, (H) and
Lemma~\ref{asm-theta-agree-split-inputs-cone-exchange}(1) for $\Phi'$ and
$h_1=1$ show that the integral in the hypothesis of
Lemma~\ref{asm-theta-agree-split-inputs-cone-unfold} is finite. That lemma gives
the convergence of $S(g)$ for almost every $g\in\mathcal D$ and the
integrability of $S$, and, after the substitution $l\mapsto l^{-1}$ and (H),
\[
\int_{\mathcal D}S\,d\mu_M=\kappa\int_K\int_{\mathbf A_E^\times}\Bigl(\int_{\mathcal D_0}\theta_{\omega_0}\bigl(e(\rho(u_\alpha(l))\omega(k,1)\Phi')\bigr)(g_1,1)\,d\mu_0(g_1)\Bigr)\chi_W(l)^{-1}|l|^z\,d\mu(l)\,d\mu_K(k),
\]
which is $\kappa\,I^{\mathrm{cone}}(\omega(1,h)\Phi)(1)$.

(11) The proof of (10) applies with the following changes. Now $e$ is an
evaluation map for $\mathcal J_\beta(H,H')$, $u=u_\beta$, the character is
$\chi_W$, $d+2=2n+2z$, no substitution $l\mapsto l^{-1}$ is made, and
$\mathcal S=\mathcal S^\circ(M)$, which $\omega$ preserves by (6) and which
consists of Schwartz-Bruhat functions
(Lemma~\ref{asm-theta-agree-continuity-pieces}(1)).
\begin{enumerate}
\item[(L)] If $p\in U_W$ is a Levi element for $(t,h_1)$, then $p^{-1}x'=t^{-1}x'$,
  and (5) for $\mathcal J_\beta$ gives
  $e(\rho(c(l)\,yx'^{\mathsf T})\omega(1,p)\Psi)=\chi_X(t)|t|^{n_1/2}\,\omega_0(1,h_1)\,e(\rho(u_\beta(l\,c(t)^{-1}))\Psi)$,
  as $c(l)t^{-1}=c(l\,c(t)^{-1})$.
\item[(T)] If $p\in U_X$ is a Levi element for $(l,g_1)$ and $p'\in U_W$ one for
  $(1,h_1)$, then (B) of the proof of Lemma~\ref{asm-theta-agree-split-shift}
  and $d/2+1-n=z$ give
  $|l|^{1-n}\,\theta(e(\rho(yx'^{\mathsf T})\omega(p,p')\Psi))=\theta_{\omega_0}(e(\rho(u_\beta(l))\Psi))(g_1,h_1)\,\chi_W(l)|l|^z$.
\end{enumerate}
Lemma~\ref{asm-theta-agree-split-inputs-cone-line} applies with
$v_0=yx'^{\mathsf T}$, $z_0=xy'^{\mathsf T}$, $\tau=c$ and $\chi_W$: for
$a\in\mathbf A_E$, $az_0=x(ay')^{\mathsf T}\in\mathcal J_\beta(H,H')$, and
$h(z_0,v_0)=h_X(x,y)h_W(y',x')=-1$. Its part (1) is used in the second form:
for $\Psi\in\mathcal S^\circ(M)$, $e(\rho(a\,yx'^{\mathsf T})\Psi)\in\mathcal S^\circ(M_0)$
by (7)(a), and Lemma~\ref{asm-theta-agree-continuity-domination} for the
second Weil pair gives the bound for the Schwartz-Bruhat $\phi$ with
$f\otimes\phi\in\mathcal S^\circ(M_0)$ for all $f$. In (b), the substitution
$l\mapsto l\,c(t)$ and (N) give the factor
$\chi_X(t)|t|^{n_1/2}\chi_W(t)^{-1}|t|^z$, and $n_1/2+z=(d-n)/2$. In (a),
$f(g)=\theta(e(\rho(yx'^{\mathsf T})\omega(g,h)\Phi))$, its invariances are those
of $f_\beta$ in Lemma~\ref{asm-theta-agree-split-shift}, and (H) reads
$f(\lambda(l,g_1)k)|l|^{1-n}=\theta_{\omega_0}(e(\rho(u_\beta(l))\omega(k,1)\Phi'))(g_1,1)\chi_W(l)|l|^z$.

(12) Let $\rho$ and $\rho_1$ be the Schrödinger operators of $M$ and $M_1$. As
$\omega(up,h)=\omega(u,1)\omega(p,1)\omega(1,h)$, the formula of (12) follows
from three laws, for Schwartz-Bruhat $\Phi$:
\begin{enumerate}
\item[(12a)] $e(\omega(1,h)\Phi)=\omega_1(1,h)e(\Phi)$ for $h\in U_W$;
\item[(12b)] $e(\omega(u,1)\Phi)=e(\Phi)$ for $u\in N_H$;
\item[(12c)] $e(\omega(p,1)\Phi)=\chi_W(t)|t|^{d/2}\,\omega_1(g_1,1)e(\Phi)$ for a
  Levi element $p\in U_X$ for $(t,g_1)$.
\end{enumerate}
Indeed, they give
$e(\omega(up,h)\Phi)=e(\omega(p,1)\omega(1,h)\Phi)=\chi_W(t)|t|^{d/2}\,\omega_1(g_1,1)\omega_1(1,h)e(\Phi)$.

\emph{Proof of (12b).} As the two are Weil pairs, $X$ and $X_1$ are Hermitian,
$W$ is skew-Hermitian and $M$ and $M_1$ are regular.

\emph{A twist.} Let $g\in U_X$ with $gx=ax$ for some $a\in\mathbf A_E$ and
$g\iota z\in\iota z+\mathbf A_Ex$ for every $z$, and write
$g\iota=\iota+x\beta^{\mathsf T}$. Then there is $c_g\in\mathbf C$ with
$e(\omega(g,1)\Phi)=c_g\,e(\Phi)$ for every Schwartz-Bruhat $\Phi$. Indeed, put
$e'(\Phi)=e(\omega(g,1)\Phi)$. By condition (b) of
Definition~\ref{asm-theta-agree-models}(6),
$\omega(g,1)\rho(w)=\rho(gw)\omega(g,1)$ on Schwartz-Bruhat functions. For
$z\in\mathbb W(X_1,W)$, $g\iota z=\iota z+xb^{\mathsf T}$ with
$b^{\mathsf T}=\beta^{\mathsf T}z$, and $\langle\iota z,xb^{\mathsf T}\rangle=0$ as
$h_X(\iota u,x)=0$; so by (R) of the proof of
Lemma~\ref{asm-theta-agree-split-inputs-expansion},
$\rho(g\iota z)=\rho(\iota z)\rho(xb^{\mathsf T})$. As $e$ is an evaluation map,
$e'(\rho(\iota z)\Phi)=e(\rho(\iota z)\rho(xb^{\mathsf T})\omega(g,1)\Phi)=\rho_1(z)e'(\Phi)$.
For a vector $q$, $g\,xq^{\mathsf T}=x(aq)^{\mathsf T}\in\mathcal J_x(H)$, so
$e'(\rho(xq^{\mathsf T})\Phi)=e'(\Phi)$. The form $\Phi\mapsto e'(\Phi)(0)$ is
admissible, by fact (b) of the proof of
Lemma~\ref{asm-theta-agree-evaluation-unique} and as $\omega(g,1)$ preserves
admissible forms (Definition~\ref{asm-theta-agree-models}(6)). So
Lemma~\ref{asm-theta-agree-split-inputs-xunique} gives $c_g$.

\emph{The constants.} If $e$ vanishes on the Schwartz-Bruhat functions, (12b)
holds; let $e(\Phi_0)\neq0$. For $g\in U_X$ and $c\in\mathbf C$ write $Q(g,c)$ if
$e(\omega(g,1)\Phi)=c\,e(\Phi)$ for every Schwartz-Bruhat $\Phi$. Then $c$ is
determined by $g$; $Q(g,c)$ and $Q(g',c')$ give $Q(gg',cc')$; $Q(1,1)$ holds;
and $Q(g,c)$ gives $c\neq0$ and $Q(g^{-1},c^{-1})$, as
$e(\Phi)=e(\omega(g,1)\omega(g^{-1},1)\Phi)=c\,e(\omega(g^{-1},1)\Phi)$.

\emph{Generators.} Let $u\in U_X$ with $ux=x$ and $u\iota z\in\iota z+\mathbf A_Ex$
for every $z$. By the twist, $Q(u,c_u)$. For an integer $t\geq1$ the Levi element
$m_t$ for $(t,1)$ (part (1)(c)) has $m_tx=tx$, $t\,m_ty=y$ and $m_t\iota=\iota$,
and the twist gives $Q(m_t,c_t)$. Put $v=u^{-1}u^{-1}m_2um_2^{-1}$. Then
$Q(v,c_v)$ with $c_v=c_u^{-2}c_2c_uc_2^{-1}=c_u^{-1}$. By
Lemma~\ref{asm-theta-agree-split-inputs-xcommutator},
$m_tvm_t^{-1}=v^{t^2}$, so $Q(v^{t^2},c_v^{t^2})$ and $Q(m_tvm_t^{-1},c_v)$ give
$c_v^{t^2}=c_v$. With $t=2$ and $t=3$, $c_v^4=c_v$ and $c_v^9=c_v$; as $c_v\neq0$,
$c_v^3=1$ and $c_v=(c_v^3)^3=1$. So $c_u=1$.

\emph{Conclusion.} The group $N_H$ is generated by these $u$, and $Q$ is
compatible with products and inverses, so $Q(u,1)$ for every $u\in N_H$, which
is (12b).

\emph{Existence, and the laws (12a) and (12c).} Write $H=(x,y,\iota,\mathfrak b)$
and let $n$, $n_1$ and $d$ be the ranks of $X$, $X_1$ and $W$. An evaluation map
for $(M,M_1,\iota,1,\mathcal J_x(H))$ exists by
Lemma~\ref{asm-theta-agree-split-inputs-xexists}. Let $e$ be one, and let $S$ and
$j$ be as in Lemma~\ref{asm-theta-agree-split-inputs-charts}(2), so that $e$ is
an $(S,j)$-evaluation. The matrices of $W$ and $X_1$ are invertible
(Lemma~\ref{asm-theta-agree-convergence-signature}). By
Definition~\ref{asm-theta-agree-split-inputs-side},
$\Sigma=(X,W,e^M,\chi_X,\omega(1,\cdot))$ and
$\Sigma_1=(X_1,W,e^{M_1},\chi_X,\omega_1(1,\cdot))$ are Weil sides, and so are
$\Sigma'=(W,X,e^M\circ\tau,\chi_W,\omega(\cdot,1))$ and
$\Sigma'_1=(W,X_1,e^{M_1}\circ\tau,\chi_W,\omega_1(\cdot,1))$, with
$\tau(u)=u^{\mathsf T}$ and $e^M$, $e^{M_1}$ the isomorphisms of $M$ and $M_1$.

(12a) Apply Proposition~\ref{asm-theta-agree-split-inputs-seesaw}(1) to
$\Sigma$ and $\Sigma_1$, with $Y=W$ and $\epsilon=-1$, to $S$, $j$ and $e$, and
to $\iota$, $x$ and $y$. Its hypotheses (i) and (ii) are
Lemma~\ref{asm-theta-agree-split-inputs-charts}(2); (iii) holds as $y\neq0$; and
(iv) holds as the columns of $\mathrm{fr}_H$ form a basis
(Lemma~\ref{asm-theta-agree-evaluation-adapted-frame}). The proposition gives
(12a).

(12c) For $U\in M_{d\times n}(\mathbf A_E)$ and $Z\in M_{d\times n_1}(\mathbf A_E)$,
$e^M(U^{\mathsf T})$ lies over $e^{M_1}(Z^{\mathsf T})$ if and only if
$U^{\mathsf T}=\iota Z^{\mathsf T}+xq^{\mathsf T}$ for some $q$, that is
$U=Z\iota^{\mathsf T}+qx^{\mathsf T}$. A Levi element $p\in U_X$ for $(t,g_1)$
has $px=tx$ and $p\iota=\iota g_1$.
Proposition~\ref{asm-theta-agree-split-inputs-seesaw}(2), for $\Sigma'$,
$\Sigma'_1$, $H$ and $\epsilon=1$, with $m=p$, $g=g_1$ and $r=0$, gives
$e(\omega(p,1)\Phi)=\chi_W(t)|t|^{d/2}\omega_1(g_1,1)e(\Phi)$, which is (12c).
% Lean, with Ext the record of ct_all (ThetaExp.lean): (1) unipEx_of and stabDec_of
% (UnipEx.lean, StabDec.lean). The frame is isFr (MatCancel.lean), with IsFr.gram, inv1 and inv2;
% (a) and (b) are ThKEnd.Grp.unipGood (Lemma asm-theta-agree-end-unipotent), applied to J
% Hermitian or skew-Hermitian with det_ne; (c) is exists_levi (IsFr.levi), levi_uniq (IsFr.ext),
% levi_mul and the continuity of z |-> mat (lev z) and its inverse; (d) is IsFr.fix_y; (e) is
% stab_dec (IsFr.stab, exists_UA, levi_inv, levi_vec) with the rational block of A' p fr; (f) is
% the second clause of stabDec_of (exists_levi, levi_uniq); (g) is levi_conj (mem_unip). The
% last clause of UnipEx, an idele of absolute value > 1, is not used here. (2): the first claim is
% the first component of kFinData_of (pure_eq_pureA); the K-finiteness is kFinData_of h.stab
% (polM, polM_stab, exists_slice, polM_fin and orbit_fin; KFinDataOf.lean, KFinAn.lean,
% KFinArch.lean). (3) iwaEx_of h.stab (IwaExOf.lean) = IwaG.iwa_glob with Iwa.iwaLoc (Theorem
% asm-theta-agree-local-mult-unramified-iwasawa) (IwaGlob.lean): the infinite places are
% exists_arch (exists_loc, which is exists_move, exists_lift and exists_inf); integrality is
% IntAt, with int_mul_ne, int_prj, ae_mulVec (EulerB.lean: lat_gen, aeInt, mem_aeInt through
% mem_of_lat, at_iff and idem_sq for the idempotent epsilon_v), applied in iwa_glob to mat (q^-1 g)
% and its inverse, int_of_KAt (idem_latW, latW_le), glue
% (lat_gen) and stab_of_int; the finite places are exists_fin by Finset induction; the
% conclusion is the end of iwa_glob. (4) Ext.mixW = Mix.mixW_closed = ThKF0.mixW_of_adel adelic_estimates_holds
% (Chr.lawW_left, lawW_unip, lawW_levi; CharacterAdmissibility/, about 2,250 lines); (12)
% Ext.mixX = Mix.mixX_closed = ThKMixX.mixed_model_plane_X_m (MixedModels/, about 2,500 lines).
% (5) Ext.mixAB = Ker.mixAB_of mixW mixX evZero pair thetaExp evUniq (MixABOf.lean): PairEx and
% EvZero give M1, om1 and e0, MixX gives e1 and the X-side formula, isEv_comp with comp_ne (polX,
% resp. pol0) makes e1 e0 an evaluation map for isoA (resp. isoB), (polA ..).uniq (resp. polB)
% with isEv_sumA (resp. isEv_sumB) gives c, and the formula is hW then h2 (resp. h2 then hW).
% (6) soStab_of (KCo.so_stable with det_herm, det_skew; SoStab.lean). (7) soEv_of linEx_of
% low_of (FormSwapZero.lean): (a) first part ev_zero (Lemma asm-theta-agree-split-inputs-kill(4)),
% second part so_kill with L built from mem_so, slice_mem', the linear model hι0 of LinEx and
% exists_coB, the set S of the three kinds of triples, hcomm (hTens_vec, hww), vec_left,
% vec_right, vec_mid, slice_rho, mult', the trace formula hf (LinEx) with trace_vec, and the
% second condition two_rows with hrow, plane_span, g1, g2 and unit_map; (b) soZeroA'_of linEx_of
% mixX thetaExp evUniq (SoZero.lean): MixX gives eX, isEv_comp with comp_ne (polX),
% (polA ..).uniq, then ev_zero for M1. (8) Ker.twoExp_of thetaExp (TwoExp.lean): (polA ..).exp and (polB ..).exp,
% capAB for V, the trivial characters by sy_of_hT, p.ad.xx, p.ad.xi, the nontrivial ones by
% exists_ne with cap_xx, cap_ix, cap_xi and PG.line, PG.nd_low', and the comparison triv_part
% (split_triv, charSeries_coeff_eq_zero). (9) is ctKernel_of (Kernel.lean): the common facts are
% hi0 (pair_low, mix0), tm_eq (rho_om), unip_y and unip_y_rat (KerUni.lean, coef_y), ev_I and
% ev_mid; (F) is constant_term_norm_eq_theta (up to the factor psi(r)) and
% constant_term_rational_unipotent_invariant; Eichler elements are eich, eich_apply, eich_unit,
% eich_mul and exists_eich (Tow.Levi.exists_UA, Fr.mem_unip). (9)(a) Ker.sum0 with (pol0 ..).exp
% thetaExp, slice_unif slice_of, tm_cont (evAdm_of, IsUSplitting.cont) and MeanOK.sum
% (mean_hasSum), the right invariance from UnipEx and the countability from GrpTop (meanOK);
% (9)(b) Ker.constant_term_anisotropic_integral_zero (exists_eich with z = 0, MeanOK.zero =
% mean_zero, half_sy); (9)(c) Ker.constant_term_cone_integral_eq with (polX ..).exp for eX from
% MixX (pair_low, PairEx), isEv_comp with comp_ne, and (polA ..).uniq evUniq giving c; (S) is hs,
% (G) is tA_eig, the term v = 0 is the hasSum_single computation with tm_eq and ev_I, and the
% terms v /= 0 are tA_zero (exists_eich with z and -h(z,z)/2, hb, hb', hgc), all as ctKernel_of
% applies them, with the constants C0 and C * c taken from one fixed Weil pair hp0; (10)
% and (11) Ext.coneA and Ext.coneB = coneA_of, coneB_of (ConeAB.lean) through cone_core
% (ConeCore.lean). For (10): u = shiftA, iota = inverse, the record Laws from lawsA (ConeLawAB.lean)
% with the alpha half of MixAB, that is (5): its field tor is (T) (lawsTA: rho_om, isLevi_inv',
% levi_inv_y, Mix2.levi; the same computation as (A) of Lemma asm-theta-agree-split-shift), uni is
% the invariance under N_H, torW is (L) with s = conjUnit t, chi_conj and absE_conjUnit for (N);
% LawsT.rat is Lemma asm-theta-agree-split-shift(1). key is LineInt (lineInt_of, Lemma
% asm-theta-agree-split-inputs-cone-line) with v0 = x y'^T, z0 = y x'^T, tau = id, and hW = DomWeil
% (Lemma asm-theta-agree-convergence-dominated) for the second pair. In cone_core: kappa from Unfold
% (Lemma asm-theta-agree-split-inputs-cone-unfold) for the fixed pair hq; LineOK with int from
% int_prod and the continuity c1 (pair_cont, LawsT.theta); kfin_dec; (b) is Laws.coneI_law with
% integral_mul_right_eq_self; (c) is LineOK.coneI_eq; (a) is LineOK.unf with lintegral_inv_eq_self
% and integral_inv_eq_self, then hUnf, with LawsT.rat and hL.uni. For (11): u = shiftB, iota = id,
% lawsB, S = M.So, hle from SoFin, SoStab (6), the second case of LineInt with domSo (SoSemi and
% FinVolU, Lemma asm-theta-agree-continuity-domination) and hso = SoEv (7)(a). thetaExp and evUniq
% are thetaExp_of slice_of and evUniq_of thetaExp (SliceSummable.lean), written out as Lemma
% asm-theta-agree-split-inputs-expansion; the K-finiteness is Lemma
% asm-theta-agree-split-inputs-kfinite. (4) is MixW = mixW_of_adel adelic_estimates_holds
% (MixWOfAdel.lean): h2 is the factorization, and the three laws are (4a) Chr.lawW_left (TA'),
% (4b) Chr.lawW_unip (TU on the generators of unip, closure induction) and (4c) Chr.lawW_levi
% (TB) (LawWLevi.lean). (12) is MixX = ThKMixX.mixed_model_plane_X_m
% (AdaptedPlaneFirstVariable.lean): hsplit is the factorization; the existence is
% exists_isEv_mixedX_m (AdaptedPlanePairRestriction.lean); (12a) is
% adapted_plane_pair_local_factor = Chr.e_om_one_eq_om1_one_e (TA'), (12c) is
% adapted_plane_pair_seesaw_identity = Chr.mixX_levi (TB) (DetNeHerm.lean, LeavesX.lean); (12b)
% is adapted_plane_pair_theta_integral = MixXUnip.unip_trivial_m: the twist is
% hermitian_adapted_pair_of_pair (twist_heis_inv, with hcomp, hperp and hgι) and
% hermitian_adapted_plane_local (Lemma asm-theta-agree-split-inputs-xunique), with
% hermitian_adapted_pair_converse (GeoAdm.comp_mem_adm) and hs.adm for the admissibility; the
% constants are Q, Quniq, Qmul, Qone, Qinv and Qpow; scl is exists_unitary_scaling_isotropic
% (exists_scale), which the book takes as the Levi element for (t,1); step is
% unitary_fixing_vector_preserves (Lemma asm-theta-agree-split-inputs-xcommutator) with t = 2, 3;
% hcl is the closure induction. The laws (4a), (4b), (4c), (12a), (12c) and the existence in
% (12) are written out through Lemmas asm-theta-agree-split-inputs-heis-* , -side, -charts,
% -xvacuum, -xexists and Proposition asm-theta-agree-split-inputs-seesaw; (4a) is lawW_left =
% TA' with setW, (4b) lawW_unip = TU, (4c) lawW_levi = TB, (12a) e_om_one_eq_om1_one_e = TA'
% with setX, (12c) mixX_levi = TB with the transposed charts d', d1' of hP.weil. What is still
% open lies in Lemma asm-theta-agree-split-inputs-siegel. The abstract theorems TA, TA', TB, TU (DetConjMulVec.lean,
% ExistsDoubledUnitaryBlock.lean, LevScalar2.lean) rest on the axioms Ax and AxInt of a test class
% (ResAdm.lean, SchrodingerApplyZeroOfFst.lean, BoxCancel.lean), proved for the adeles by ax_of
% (axB, axDelta_of lagLine, axImp_of AdelSt) and axInt (AdelInt.lean), on setW/setX and fr_plane
% (FrOfFrame.lean), and on parabolic_character_trivial_on_diagonal
% (DoubledUnitaryCharacterTrivialOfSiegel.lean); about 4,800 lines with the MixedModels files.
\end{proof}

\subsection{Splitting theta integrals along a plane}\label{subsection-asm-splitting}

\begin{lemma}\label{asm-theta-agree-split-kernel}
Let $E$ be a CM field with a theta base $B$
(Definition~\ref{asm-theta-agree-models}(1)). Let $X$, $W$ and $W_1$ be lattice
spaces of ranks $n$, $d$ and $d_1$, $M$ a model for $(X,W)$ with Schrödinger
operators $\rho$, $M_1$ a regular model for $(X,W_1)$, $H'$ an adapted plane in
$W$ over $W_1$, and $e_0$ an evaluation map for
$(M,M_1,1,\iota_{H'},\mathcal J_0(H'))$ (Definition~\ref{asm-theta-agree-planes}).
Then there is $C_0\neq0$ such that:
\begin{enumerate}
\item if $\operatorname{Iso}(X)$ is empty
  (Definition~\ref{asm-theta-agree-split-notions}(1)), then for every Weil pair
  $(X,W,M,\chi_X,\chi_W,\omega)$ of any signature
  (Definition~\ref{asm-theta-agree-triple}(1)), every Schwartz-Bruhat function
  $\Phi$ on $\mathbf A^{n\times d}$ and every $g\in U_X$, the function
  $h\mapsto C_0\,\theta(e_0(\omega(g,h)\Phi))$ is the constant term of
  $\theta_\omega(\Phi)(g,\cdot)$ along $H'$;
\item for every adapted plane $H$ in $X$ over a lattice space $X_1$, every
  regular model $M_2$ for $(X_1,W_1)$ and every evaluation map $e$ for
  $(M,M_2,\iota_H,\iota_{H'},\mathcal J_\alpha(H,H'))$, there is $C_1\neq0$
  such that the following holds for every transporter family $\gamma$ for
  $x_H$, every Weil pair $(X,W,M,\chi_X,\chi_W,\omega)$ of any signature,
  every Schwartz-Bruhat function $\Phi$ and every $g\in U_X$: for every
  $h\in U_W$ the sum
  \[
  S(g,h)=\sum_{v\in\operatorname{Iso}(X)}\theta\bigl(e(\rho(x_Hy_{H'}^{\mathsf T})\,\omega(\gamma(v)g,h)\Phi)\bigr)
  \]
  converges absolutely, and $h\mapsto C_0\,\theta(e_0(\omega(g,h)\Phi))+C_1S(g,h)$
  is the constant term of $\theta_\omega(\Phi)(g,\cdot)$ along $H'$.
\end{enumerate}
\end{lemma}
\leanref{Development/WeilBound/ThetaLifts/CTKernel.lean}{PositivityTheory.ThKG.CTKernel}
\leanref{Development/WeilBound/ThetaLifts/SmallCones/Kernel.lean}{PositivityTheory.ThKCT.ctKernel_of}
\leanref{Development/WeilBound/ThetaLifts/SmallCones/Kernel.lean}{PositivityTheory.ThKCT.Ker.constant_term_integral_at_zero}
\leanref{Development/WeilBound/ThetaLifts/SmallCones/Kernel.lean}{PositivityTheory.ThKCT.Ker.constant_term_rational_translate}
\leanref{Development/WeilBound/ThetaLifts/SmallCones/Kernel.lean}{PositivityTheory.ThKCT.Ker.split_sum}
\begin{proof}
Write $y'=y_{H'}$, and use the integrals $T_\omega(x,\Psi)$ of
Lemma~\ref{asm-theta-agree-split-inputs}(9), for Weil pairs
$(X,W,M,\chi_X,\chi_W,\omega)$, Schwartz-Bruhat functions $\Psi$, and $\nu$ and
$\mathcal F$ as there. Let $C_0$ be the constant of
Lemma~\ref{asm-theta-agree-split-inputs}(9)(a).

\emph{An auxiliary Weil pair.} Given a Weil pair
$(X,W,M,\chi_X,\chi_W,\omega)$, the space $W_1$ is skew-Hermitian
(Lemma~\ref{asm-theta-agree-split-low}), and
Theorem~\ref{asm-theta-agree-tower-weil} gives $\omega_1$ such that
$(X,W_1,M_1,\chi_X,\chi_W,\omega_1)$ is a Weil pair, as
$\chi_W|_{\mathbf A^\times}=\eta^d=\eta^{d_1}$.

\emph{The zero vector.} By Lemma~\ref{asm-theta-agree-split-inputs}(4), for
these two Weil pairs, with $g=1$ and the Levi element $1$ for $(1,1)$,
$e_0(\omega(1,u)\Psi)=e_0(\Psi)$ for $u\in N_{H'}$. As $\nu(\mathcal F)=1$,
$T_\omega(0,\Psi)=\theta(e_0(\Psi))$.

\emph{Rational translates.} Let $\gamma\in U_X(F)$, $x\in E^n$ and
$u\in N_{H'}$, and put $\Psi'=\omega(1,u)\Psi$. The function
$e_0(\rho(xy'^{\mathsf T})\Psi')$ is Schwartz-Bruhat.
Lemma~\ref{asm-theta-agree-anisotropic-automorphy} for the Weil pair of
$\omega_1$, Lemma~\ref{asm-theta-agree-split-inputs}(4) with $g=\gamma$ and
$m=u=1$, and the rule $\omega(g,1)\rho(w)=\rho(gw)\omega(g,1)$ of
Definition~\ref{asm-theta-agree-models}(6)(b) give
\[
\theta(e_0(\rho(xy'^{\mathsf T})\Psi'))=\theta(\omega_1(\gamma,1)e_0(\rho(xy'^{\mathsf T})\Psi'))
=\theta(e_0(\omega(\gamma,1)\rho(xy'^{\mathsf T})\Psi'))
=\theta(e_0(\rho((\gamma x)y'^{\mathsf T})\omega(\gamma,1)\Psi')).
\]
As $\omega(\gamma,1)\omega(1,u)=\omega(1,u)\omega(\gamma,1)$, integrating over
$\mathcal F$ gives $T_\omega(x,\Psi)=T_\omega(\gamma x,\omega(\gamma,1)\Psi)$.

\emph{The constant term.} Let $\Phi$ be Schwartz-Bruhat, $g\in U_X$ and
$h\in U_W$, and put $\Psi_h=\omega(g,h)\Phi$. As
$\omega(g,uh)\Phi=\omega(1,u)\Psi_h$, the integral
$\int_{\mathcal F}\theta_\omega(\Phi)(g,uh)\,d\nu(u)$ equals
$\int_{\mathcal F}\theta(\omega(1,u)\Psi_h)\,d\nu(u)$, which by
Lemma~\ref{asm-theta-agree-split-inputs}(9)(a) is the sum of the absolutely
convergent series $\sum_{x\in E^n}C_0T_\omega(x,\Psi_h)$. By (9)(b) its terms
vanish unless $x=0$ or $x\in\operatorname{Iso}(X)$, and the term at $0$ is
$C_0\theta(e_0(\omega(g,h)\Phi))$. Hence
\[
\sum_{x\in\operatorname{Iso}(X)}C_0T_\omega(x,\Psi_h)=\int_{\mathcal F}\theta_\omega(\Phi)(g,uh)\,d\nu(u)-C_0\,\theta(e_0(\omega(g,h)\Phi)),\tag{K}
\]
the sum converging absolutely.

(1) If $\operatorname{Iso}(X)$ is empty, the left side of (K) is $0$, for
every $\nu$ and $\mathcal F$.

(2) Let $C$ be the constant of Lemma~\ref{asm-theta-agree-split-inputs}(9)(c)
for $H$, $M_2$ and $e$, and put $C_1=C_0C$. For $v\in\operatorname{Iso}(X)$,
the rational translates by $\gamma(v)$ and (9)(c) give
\[
T_\omega(v,\Psi_h)=T_\omega(x_H,\omega(\gamma(v)g,h)\Phi)=C\,\theta\bigl(e(\rho(x_Hy'^{\mathsf T})\,\omega(\gamma(v)g,h)\Phi)\bigr).
\]
So by (K), $\sum_vC_1\theta(e(\rho(x_Hy'^{\mathsf T})\omega(\gamma(v)g,h)\Phi))$
converges absolutely to
$\int_{\mathcal F}\theta_\omega(\Phi)(g,uh)\,d\nu(u)-C_0\theta(e_0(\omega(g,h)\Phi))$.
As $C_1\neq0$ and there is such a pair $(\nu,\mathcal F)$
(Lemma~\ref{asm-theta-agree-split-inputs}(1)(a)), $S(g,h)$ converges
absolutely, and for every $\nu$ and $\mathcal F$,
$\int_{\mathcal F}\theta_\omega(\Phi)(g,uh)\,d\nu(u)=C_0\theta(e_0(\omega(g,h)\Phi))+C_1S(g,h)$.
\end{proof}

\begin{lemma}\label{asm-theta-agree-split-unique}
Let $(B,G,\mathcal X,\omega,\omega')$ be a triple of ranks $(a,a')$ for a theta
base $B$ of a CM field $E$ (Definition~\ref{asm-theta-agree-triple}), with
skew-Hermitian space $W$, and let $H'$ be an adapted plane in $W$ over a
lattice space $W_1$ at some positions. If $(f_\alpha,f_\beta)$ and
$(f'_\alpha,f'_\beta)$ both split a function $f$ on $U_W$ along $H'$
(Definition~\ref{asm-theta-agree-planes}(3)), then $f_\alpha=f'_\alpha$ and
$f_\beta=f'_\beta$.
\end{lemma}
\leanref{Development/WeilBound/ThetaLifts/SmallCones/ThetaKernelGeneralIdentity.lean}{PositivityTheory.ThKCT.split2_unique}
\leanref{Development/WeilBound/ThetaLifts/SmallCones/ThetaKernelGeneralIdentity.lean}{PositivityTheory.ThKCT.split_uniq}
\begin{proof}
Let $\mathcal X=(\alpha,\beta,S_0,Q)$ and let $p,q$ be the integers of $G$.
As $G$ is proper, $a'<a$ and $W$ is skew-Hermitian of signature
$(a+a'-p-q,p+q)$. By
Lemma~\ref{asm-theta-agree-split-inputs}(1)(a) there are a Haar measure $\nu$
on $N_{H'}$ and a fundamental domain $\mathcal F$ for $N_{H'}\cap U_W(F)$ in
$N_{H'}$ with $\nu(\mathcal F)=1$. As $f_\alpha+f_\beta$ and
$f'_\alpha+f'_\beta$ are both the constant term of $f$,
\[
f_\alpha(h)+f_\beta(h)=\int_{\mathcal F}f(uh)\,d\nu(u)=f'_\alpha(h)+f'_\beta(h)
\qquad(h\in U_W).
\]
Let $t\in\mathbf A_E^\times$ be $2$ at one infinite place of $F$ and $1$ at the
other places; then $|t|=|N_{E/F}(t)|_{\mathbf A}=4>1$. By
Lemma~\ref{asm-theta-agree-split-inputs}(1)(c) there is a Levi element $m$ for
$(t,1)$. Put $\delta=f_\alpha-f'_\alpha$; by the display,
$\delta=f'_\beta-f_\beta$. The laws of $f_\alpha,f'_\alpha$ and of
$f_\beta,f'_\beta$ give
\[
\delta(mh)=\alpha(t)|t|^{a'/2}\delta(h)\quad\text{and}\quad\delta(mh)=\beta(t)|t|^{a/2}\delta(h).
\]
If $\delta(h)\neq0$, then $\alpha(t)|t|^{a'/2}=\beta(t)|t|^{a/2}$, and taking
absolute values, as $|\alpha(t)|=|\beta(t)|=1$, gives $|t|^{a'/2}=|t|^{a/2}$,
which is false since $|t|>1$ and $a'<a$. So $\delta=0$.
\end{proof}

\begin{lemma}\label{asm-theta-agree-split-fubini}
Let $(X,W,M,\chi_X,\chi_W,\omega)$ be a Weil pair of signature $(r,s;r',s')$
for a theta base $B$ of a CM field $E$
(Definition~\ref{asm-theta-agree-triple}(1)), with $X$ and $W$ of ranks $n$
and $d$, and with measure $\mu_M$ and domain $\mathcal D$ of $M$. Let $H'$ be
an adapted plane in $W$ over a lattice space $W_1$, and let $\Phi$ be a
Schwartz-Bruhat function on $\mathbf A^{n\times d}$ for which there are $C$
and $k$ with
$\int_{\mathcal D}|\theta_\omega(\Phi)(g,h)|\,d\mu_M(g)\leq C\operatorname{ht}(h)^k$
for every $h\in U_W$. For each $g\in U_X$ let $c_g$ be the constant term of
$\theta_\omega(\Phi)(g,\cdot)$ along $H'$
(Definition~\ref{asm-theta-agree-planes}(2)). Then for every $h\in U_W$ the
function $g\mapsto c_g(h)$ is integrable on $\mathcal D$, and
$h\mapsto\int_{\mathcal D}c_g(h)\,d\mu_M(g)$ is the constant term of
$I_{M,\omega}(\Phi)$ along $H'$.
\end{lemma}
\leanref{Development/WeilBound/ThetaLifts/SmallCones/ThetaUniq.lean}{PositivityTheory.ThKCT.CTFub}
\leanref{Development/WeilBound/ThetaLifts/SmallCones/IntegrableProdOfBound.lean}{PositivityTheory.ThKCT.ctFub_of}
\leanref{Development/WeilBound/ThetaLifts/SmallCones/IntegrableProdOfBound.lean}{PositivityTheory.ThKCT.integrable_prod_of_bound}
\begin{proof}
Let $\nu$ be a Haar measure on $N_{H'}$ and $\mathcal F$ a fundamental domain
for $N_{H'}\cap U_W(F)$ in $N_{H'}$ with $\nu(\mathcal F)=1$; by
Lemma~\ref{asm-theta-agree-split-inputs}(1)(a) there is one. Fix $h\in U_W$ and
put $f(g,u)=\theta_\omega(\Phi)(g,uh)$ on $U_X\times N_{H'}$. It is continuous
by Lemma~\ref{asm-theta-agree-convergence-theta}(2). Let $C_N\subset N_{H'}$ be
the compact set of Lemma~\ref{asm-theta-agree-split-inputs}(1)(b). The height
is bounded, by $D_0$ say, on the compact set $C_Nh$
(Theorem~\ref{asm-theta-agree-convergence-siegel}(1), $J_W$ being
nondegenerate by Lemma~\ref{asm-theta-agree-convergence-signature}). For
$u\in N_{H'}$ choose $\gamma\in N_{H'}\cap U_W(F)$ with $\gamma u\in C_N$. By
Lemma~\ref{asm-theta-agree-anisotropic-automorphy},
$f(g,u)=\theta_\omega(\Phi)(g,\gamma uh)$, so
\[
\int_{\mathcal D}|f(g,u)|\,d\mu_M(g)\leq C\operatorname{ht}(\gamma uh)^k\leq|C|D_0^k .
\]
As $\nu(\mathcal F)=1$, Tonelli's theorem shows that $f$ is integrable on
$\mathcal D\times\mathcal F$ for $\mu_M\times\nu$. By Fubini's theorem,
$g\mapsto\int_{\mathcal F}f(g,u)\,d\nu(u)=c_g(h)$ is integrable on $\mathcal D$,
and
\[
\int_{\mathcal D}c_g(h)\,d\mu_M(g)=\int_{\mathcal F}\int_{\mathcal D}\theta_\omega(\Phi)(g,uh)\,d\mu_M(g)\,d\nu(u)=\int_{\mathcal F}I_{M,\omega}(\Phi)(uh)\,d\nu(u).
\]
This holds for every such $\nu$ and $\mathcal F$, which is the second claim.
\end{proof}

\begin{lemma}\label{asm-theta-agree-split-domains}
Let $E$ be a CM field with a theta base $B$
(Definition~\ref{asm-theta-agree-models}(1)), let $X$ be a Hermitian lattice
space of some signature, let $\mu$ and $\mu'$ be Haar measures on $U_X$, and
let $\mathcal D$ and $\mathcal D'$ be fundamental domains for $U_X(F)$ in $U_X$
with respect to $\mu$ and $\mu'$. Take the integral of a function that is not
integrable to be $0$. Then there is $c>0$ such that
$\int_{\mathcal D}\varphi\,d\mu=c\int_{\mathcal D'}\varphi\,d\mu'$ for every
function $\varphi$ on $U_X$ with $\varphi(\gamma g)=\varphi(g)$ for all
$\gamma\in U_X(F)$ and $g\in U_X$.
\end{lemma}
\leanref{Development/WeilBound/ThetaLifts/SmallCones/GenAnisoOf.lean}{PositivityTheory.ThKCT.dom_change}
\begin{proof}
The group $U_X$ is locally compact, second countable and Hausdorff, and
$U_X(F)$ is a countable discrete subgroup
(Theorem~\ref{asm-theta-agree-convergence-siegel}(1) and
Lemma~\ref{asm-theta-agree-convergence-signature}). By the uniqueness of Haar
measure, $\mu=c\mu'$ with $c>0$. Then $\mathcal D'$ is also a fundamental
domain with respect to $\mu$. A function invariant under $U_X(F)$ on the left
has the same integral over two fundamental domains, and is integrable on one
if and only if it is on the other, as in the proof of
Lemma~\ref{asm-theta-agree-anisotropic-invariance}. Hence
$\int_{\mathcal D}\varphi\,d\mu=\int_{\mathcal D'}\varphi\,d\mu=c\int_{\mathcal D'}\varphi\,d\mu'$.
\end{proof}

\begin{lemma}\label{asm-theta-agree-split-zero}
Let $(X,W,M,\chi_X,\chi_W,\omega)$ be a Weil pair of signature $(r,s;r',s')$
for a theta base $B$ of a CM field $E$
(Definition~\ref{asm-theta-agree-triple}(1)), with $X$ and $W$ of ranks $n$
and $d$, and with measure $\mu_M$ and domain $\mathcal D$ of $M$. Let $H'$ be
an adapted plane in $W$ over a lattice space $W_1$, let
$(X,W_1,M_1,\chi_X,\chi_W,\omega_1)$ be a Weil pair of signature
$(r,s;r_1,s_1)$, with measure $\mu_1$ and domain $\mathcal D_1$ of $M_1$, and
let $e$ be an evaluation map for $(M,M_1,1,\iota_{H'},\mathcal J_0(H'))$
(Definition~\ref{asm-theta-agree-planes}(5)). Take the integral of a function
that is not integrable to be $0$. Let $c\in\mathbf C$ be such that
$\int_{\mathcal D}\varphi\,d\mu_M=c\int_{\mathcal D_1}\varphi\,d\mu_1$ for every
function $\varphi$ on $U_X$ invariant under $U_X(F)$ on the left, let
$C_0\in\mathbf C$ and let $\Phi$ be a Schwartz-Bruhat function on
$\mathbf A^{n\times d}$. Then the function
\[
A(h)=\int_{\mathcal D}C_0\,\theta(e(\omega(g,h)\Phi))\,d\mu_M(g)
\]
on $U_W$ has the law $(\chi_X,n/2)$, and
$A(ph)=C_0c\,I_{M_1,\omega_1}(e(\omega(1,h)\Phi))(h_1)$ for every Levi element
$p$ for $(1,h_1)$ and every $h\in U_W$.
\end{lemma}
\leanref{Development/WeilBound/ThetaLifts/SmallCones/GenAnisoOf.lean}{PositivityTheory.ThKCT.zero_term}
\begin{proof}
Let $g\in U_X$, $h\in U_W$ and let $p$ be a Levi element for $(t,h_1)$. As
$\omega$ is a representation, $\omega(g,ph)\Phi=\omega(g,p)\Psi$ with
$\Psi=\omega(1,h)\Phi$, a Schwartz-Bruhat function. By
Lemma~\ref{asm-theta-agree-split-inputs}(4) with $u=1$,
\[
\theta(e(\omega(g,ph)\Phi))=\chi_X(t)|t|^{n/2}\,\theta(\omega_1(g,h_1)e(\Psi))=\chi_X(t)|t|^{n/2}\,\theta_{\omega_1}(e(\Psi))(g,h_1).\tag{Z}
\]
With $t=1$ and $p=1$, a Levi element for $(1,1)$, this gives
$\theta(e(\omega(g,h)\Phi))=\theta_{\omega_1}(e(\Psi))(g,1)$. So for a Levi
element $m$ for $(t,1)$, (Z) gives
$\theta(e(\omega(g,mh)\Phi))=\chi_X(t)|t|^{n/2}\theta(e(\omega(g,h)\Phi))$ for
every $g$, and integrating over $\mathcal D$ gives
$A(mh)=\chi_X(t)|t|^{n/2}A(h)$.

Let $p$ be a Levi element for $(1,h_1)$. By (Z),
$A(ph)=C_0\int_{\mathcal D}\theta_{\omega_1}(e(\Psi))(g,h_1)\,d\mu_M(g)$. The
function $e(\Psi)$ is Schwartz-Bruhat, so
$g\mapsto\theta_{\omega_1}(e(\Psi))(g,h_1)$ is invariant under $U_X(F)$ on the
left (Lemma~\ref{asm-theta-agree-anisotropic-automorphy} for the second Weil
pair). Hence
$A(ph)=C_0c\int_{\mathcal D_1}\theta_{\omega_1}(e(\Psi))(g,h_1)\,d\mu_1(g)=C_0c\,I_{M_1,\omega_1}(e(\Psi))(h_1)$.
\end{proof}

\begin{lemma}\label{asm-theta-agree-split-shift}
Let $B$ be a theta base of a CM field $E$ with complex conjugation $c$
(Definition~\ref{asm-theta-agree-models}(1)). Let $H$ be an adapted plane in a
lattice space $X$ over $X_1$ and $H'$ one in $W$ over $W_1$
(Definition~\ref{asm-theta-agree-planes}(1)), with ranks $n$, $n_1$, $d$ and
$d_1$, and write $x=x_H$, $y=y_H$, $x'=x_{H'}$ and $y'=y_{H'}$. Let
$(X,W,M,\chi_X,\chi_W,\omega)$ and $(X_1,W_1,M_0,\chi_X,\chi_W,\omega_0)$ be
Weil pairs of signatures $(r+1,s+1;r'+1,s'+1)$ and $(r,s;r',s')$
(Definition~\ref{asm-theta-agree-triple}(1)), with Schrödinger operators
$\rho$ for $M$, and let $e_\alpha$ and $e_\beta$ be evaluation maps for
$(M,M_0,\iota_H,\iota_{H'},\mathcal J_\alpha(H,H'))$ and
$(M,M_0,\iota_H,\iota_{H'},\mathcal J_\beta(H,H'))$. Let $\Phi$ be a
Schwartz-Bruhat function on $\mathbf A^{n\times d}$ and $h\in U_W$, and put
\[
f_\alpha(g)=\theta\bigl(e_\alpha(\rho(xy'^{\mathsf T})\,\omega(g,h)\Phi)\bigr),\qquad
f_\beta(g)=\theta\bigl(e_\beta(\rho(yx'^{\mathsf T})\,\omega(g,h)\Phi)\bigr)\qquad(g\in U_X).
\]
\begin{enumerate}
\item If $\delta\in U_X(F)$ and $\delta x=x$, then $f_\alpha(\delta g)=f_\alpha(g)$
  and $f_\beta(\delta g)=f_\beta(g)$ for every $g\in U_X$.
\item Let $\gamma$ be a transporter family for $x$
  (Definition~\ref{asm-theta-agree-split-notions}(1)), $v\in\operatorname{Iso}(X)$,
  $\xi\in E^\times$, $g\in U_X$, and $\Psi=\omega(\gamma(v)g,h)\Phi$. Then
  \[
  f_\alpha(\gamma(\xi v)g)=\theta\bigl(e_\alpha(\rho(\xi\,xy'^{\mathsf T})\Psi)\bigr),\qquad
  f_\beta(\gamma(\xi v)g)=\theta\bigl(e_\beta(\rho(c(\xi)^{-1}yx'^{\mathsf T})\Psi)\bigr).
  \]
\end{enumerate}
\end{lemma}
\leanref{Development/WeilBound/ThetaLifts/SmallCones/ConeLaw.lean}{PositivityTheory.ThKCT.LawsT}
\leanref{Development/WeilBound/ThetaLifts/SmallCones/ConeLaw.lean}{PositivityTheory.ThKCT.LawsT.rat}
\leanref{Development/WeilBound/ThetaLifts/SmallCones/IsoActFree.lean}{PositivityTheory.ThKCT.LawsT.mult}
\leanref{Development/WeilBound/ThetaLifts/SmallCones/ConeLawAB.lean}{PositivityTheory.ThKCT.lawsTA}
\leanref{Development/WeilBound/ThetaLifts/SmallCones/ConeLawAB.lean}{PositivityTheory.ThKCT.lawsB}
\leanref{Development/WeilBound/ThetaLifts/SmallCones/ConeBasic.lean}{PositivityTheory.ThKCT.rho_om}
\begin{proof}
\emph{Two rules.} By conditions (a) and (b) of
Definition~\ref{asm-theta-agree-models}(6),
$\omega(g,h)\rho(w)=\rho(gwh^{\mathsf T})\omega(g,h)$ on Schwartz-Bruhat
functions, so for $a\in\mathbf A_E^n$ and $b\in\mathbf A_E^d$
\[
\rho(ab^{\mathsf T})\,\omega(g,h)\Psi=\omega(g,h)\,\rho\bigl((g^{-1}a)(h^{-1}b)^{\mathsf T}\bigr)\Psi .\tag{R}
\]
The operators $\omega(g,h)$ and $\rho(w)$ and the maps $e_\alpha,e_\beta$
preserve the Schwartz-Bruhat functions. By
Lemma~\ref{asm-theta-agree-split-inputs}(5), with all Levi elements equal to
$1$, $e_\alpha(\omega(u,1)\Psi)=e_\alpha(\Psi)$ and
$e_\beta(\omega(u,1)\Psi)=e_\beta(\Psi)$ for $u\in N_H$.

\emph{Invariance under $N_H$.} Let $u\in N_H$ and $\Psi$ be Schwartz-Bruhat.
By Lemma~\ref{asm-theta-agree-split-inputs}(1)(d), $u^{-1}x=x$ and
$u^{-1}y=y+\iota_Ha+bx$ for some $a$ and $b$. By (R),
$e_\alpha(\rho(xy'^{\mathsf T})\omega(u,1)\Psi)=e_\alpha(\omega(u,1)\rho(xy'^{\mathsf T})\Psi)=e_\alpha(\rho(xy'^{\mathsf T})\Psi)$.
For $e_\beta$, (R) gives
$\rho(yx'^{\mathsf T})\omega(u,1)\Psi=\omega(u,1)\rho(w_1+w_2+w_3)\Psi$ with
$w_1=yx'^{\mathsf T}$, $w_2=(\iota_Ha)x'^{\mathsf T}$ and
$w_3=x(bx')^{\mathsf T}$. The symplectic form of $\mathbb W(X,W)$ vanishes on
pairs of elements $a_1x'^{\mathsf T}$ and $a_2x'^{\mathsf T}$, as
$x'^{\mathsf T}J_W\,c(x')=0$ for the matrix $J_W$ of $W$, so
$\rho(w_1+w_2+w_3)=\rho(w_3)\rho(w_2)\rho(w_1)$. As $w_2$ and $w_3$ lie in
$\mathcal J_\beta(H,H')$, $e_\beta$ is invariant under $\rho(w_2)$ and
$\rho(w_3)$, and so
$e_\beta(\rho(yx'^{\mathsf T})\omega(u,1)\Psi)=e_\beta(\rho(yx'^{\mathsf T})\Psi)$.

\emph{Levi elements.} Let $l\in\mathbf A_E^\times$, $p'\in U_W$ a Levi element
for $(1,h_1)$, and $\Psi$ Schwartz-Bruhat. Then $p'^{-1}x'=x'$ and
$p'^{-1}y'=y'$. If $p\in U_X$ is a Levi element for $(l^{-1},g_1)$, then
$p^{-1}x=lx$, and (R) and Lemma~\ref{asm-theta-agree-split-inputs}(5) give
\[
e_\alpha(\rho(xy'^{\mathsf T})\omega(p,p')\Psi)=e_\alpha(\omega(p,p')\rho(l\,xy'^{\mathsf T})\Psi)=\chi_W(l)^{-1}|l|^{-d_1/2}\,\omega_0(g_1,h_1)\,e_\alpha(\rho(l\,xy'^{\mathsf T})\Psi).\tag{A}
\]
If $p$ is a Levi element for $(l,g_1)$, then $p^{-1}y=c(l)y$, and in the same
way
\[
e_\beta(\rho(yx'^{\mathsf T})\omega(p,p')\Psi)=\chi_W(l)|l|^{d/2}\,\omega_0(g_1,h_1)\,e_\beta(\rho(c(l)\,yx'^{\mathsf T})\Psi).\tag{B}
\]

(1) By Lemma~\ref{asm-theta-agree-split-inputs}(1)(e) with $t=1$,
$\delta=um$ with $u\in N_H$ and $m$ a Levi element for $(1,g_1)$, with
$g_1\in U_{X_1}(F)$. Put $\Psi=\omega(g,h)\Phi$, so that
$\omega(\delta g,h)\Phi=\omega(u,1)\omega(m,1)\Psi$. By the invariance under
$N_H$ and by (A) with $l=1$ and $p'=1$,
\[
f_\alpha(\delta g)=\theta\bigl(e_\alpha(\rho(xy'^{\mathsf T})\omega(m,1)\Psi)\bigr)=\theta\bigl(\omega_0(g_1,1)e_\alpha(\rho(xy'^{\mathsf T})\Psi)\bigr)=\theta_{\omega_0}\bigl(e_\alpha(\rho(xy'^{\mathsf T})\Psi)\bigr)(g_1,1),
\]
which is $\theta(e_\alpha(\rho(xy'^{\mathsf T})\Psi))=f_\alpha(g)$ by
Lemma~\ref{asm-theta-agree-anisotropic-automorphy} for the second Weil pair,
as $g_1$ is rational. The same argument with (B) applies to $f_\beta$.

(2) Let $p\in U_X$ be the Levi element for $(\xi^{-1},1)$
(Lemma~\ref{asm-theta-agree-split-inputs}(1)(c)); it lies in $U_X(F)$ by
Lemma~\ref{asm-theta-agree-split-inputs}(1)(f). Then
$p\gamma(v)(\xi v)=\xi\,px=x$, so
$\delta=\gamma(\xi v)(p\gamma(v))^{-1}$ lies in $U_X(F)$ and
$\delta x=\gamma(\xi v)(\xi v)=x$. By (1),
$f_\alpha(\gamma(\xi v)g)=f_\alpha(p\gamma(v)g)$, and
$\omega(p\gamma(v)g,h)\Phi=\omega(p,1)\Psi$. By (A) with $l=\xi$, $g_1=1$ and
$p'=1$,
\[
f_\alpha(\gamma(\xi v)g)=\chi_W(\xi)^{-1}|\xi|^{-d_1/2}\,\theta\bigl(e_\alpha(\rho(\xi\,xy'^{\mathsf T})\Psi)\bigr)=\theta\bigl(e_\alpha(\rho(\xi\,xy'^{\mathsf T})\Psi)\bigr),
\]
since $\chi_W$ is trivial on $E^\times$ and $|\xi|=1$ by the product formula.
In the same way, $f_\beta(\gamma(\xi v)g)=f_\beta(p\gamma(v)g)$, and (B) with
$l=\xi^{-1}$ gives
$f_\beta(\gamma(\xi v)g)=\theta(e_\beta(\rho(c(\xi)^{-1}yx'^{\mathsf T})\Psi))$.
\end{proof}

\begin{lemma}\label{asm-theta-agree-split-aniso}
Let $E$ be a CM field with a theta base $B$
(Definition~\ref{asm-theta-agree-models}(1)). Let $X$, $W$ and $W_1$ be
lattice spaces of ranks $n$, $d$ and $d_1$, $H'$ an adapted plane in $W$ over
$W_1$, $M$ a model for $(X,W)$, $M_1$ a model for $(X,W_1)$, and $e$ an
evaluation map for $(M,M_1,1,\iota_{H'},\mathcal J_0(H'))$
(Definition~\ref{asm-theta-agree-planes}). Let $r,s,r',s',r_1,s_1\geq0$ with
$\min(r,s)=0$. Take the integral of a function that is not integrable to be
$0$. Then there is $C\neq0$ such that the following holds for all Hecke
characters $\chi_X,\chi_W$ and representations $\omega,\omega_1$ for which
$(X,W,M,\chi_X,\chi_W,\omega)$ and $(X,W_1,M_1,\chi_X,\chi_W,\omega_1)$ are Weil
pairs of signatures $(r,s;r',s')$ and $(r,s;r_1,s_1)$
(Definition~\ref{asm-theta-agree-triple}(1)). For every Schwartz-Bruhat
function $\Phi$ on $\mathbf A^{n\times d}$ there is a function $A_0$ on $U_W$
that is the constant term of $I_{M,\omega}(\Phi)$ along $H'$, has the law
$(\chi_X,n/2)$, and satisfies
$A_0(ph)=C\,I_{M_1,\omega_1}(e(\omega(1,h)\Phi))(h_1)$ for every Levi element
$p$ for $(1,h_1)$ and every $h\in U_W$.
\end{lemma}
\leanref{Development/WeilBound/ThetaLifts/SmallCones/ThetaUniq.lean}{PositivityTheory.ThKCT.GenAniso}
\leanref{Development/WeilBound/ThetaLifts/SmallCones/GenAnisoOf.lean}{PositivityTheory.ThKCT.genAniso_of}
\begin{proof}
If there are no characters and representations with the two Weil pairs, put
$C=1$. Otherwise fix such data. Then $X$ is Hermitian of signature $(r,s)$, and
$M$ and $M_1$ are regular. Let $C_0\neq0$ be given by
Lemma~\ref{asm-theta-agree-split-kernel} for $(X,W,W_1,M,M_1,H',e)$, and
$c\neq0$ by Lemma~\ref{asm-theta-agree-split-domains} for the measures and
domains of $M$ and $M_1$. Put $C=C_0c$.

Let $\chi_X,\chi_W,\omega,\omega_1$ give the two Weil pairs, and $\Phi$ be
Schwartz-Bruhat. Put
$A_0(h)=\int_{\mathcal D}C_0\,\theta(e(\omega(g,h)\Phi))\,d\mu_M(g)$, with the
measure $\mu_M$ and domain $\mathcal D$ of $M$. By
Lemma~\ref{asm-theta-agree-split-zero}, $A_0$ has the law $(\chi_X,n/2)$ and
$A_0(ph)=C_0c\,I_{M_1,\omega_1}(e(\omega(1,h)\Phi))(h_1)$. As $\min(r,s)=0$,
$\operatorname{Iso}(X)$ is empty (Lemma~\ref{asm-theta-agree-split-witt}(1)),
so by Lemma~\ref{asm-theta-agree-split-kernel}(1) the constant term of
$\theta_\omega(\Phi)(g,\cdot)$ along $H'$ is
$h\mapsto C_0\,\theta(e(\omega(g,h)\Phi))$ for every $g$. By
Lemma~\ref{asm-theta-agree-convergence}(1), which applies as $\min(r,s)=0$,
there are $C'$ and $k$ with
$\int_{\mathcal D}|\theta_\omega(\Phi)(g,h)|\,d\mu_M(g)\leq C'\operatorname{ht}(h)^k$.
By Lemma~\ref{asm-theta-agree-split-fubini}, $A_0$ is the constant term of
$I_{M,\omega}(\Phi)$ along $H'$.
\end{proof}

\begin{lemma}\label{asm-theta-agree-split-big}
Let $E$ be a CM field with a theta base $B$
(Definition~\ref{asm-theta-agree-models}(1)). Let $H$ be an adapted plane in a
lattice space $X$ over $X_1$ and $H'$ one in $W$ over $W_1$, with ranks $n$,
$n_1$, $d$ and $d_1$. Let $M$ be a model for $(X,W)$, $M_0$ and $M_1$ regular
models for $(X_1,W_1)$ and $(X,W_1)$, $e_0$ an evaluation map for
$(M,M_1,1,\iota_{H'},\mathcal J_0(H'))$ and $e$ one for
$(M,M_0,\iota_H,\iota_{H'},\mathcal J_\alpha(H,H'))$
(Definition~\ref{asm-theta-agree-planes}). Let $r,s,r',s'\geq0$ and $z\geq0$ be
integers with $2n=d+2z$ and $\min(r,s)+2\leq z$, let $K\subset U_X$ be an
Iwasawa subgroup for $x_H$ (Definition~\ref{asm-theta-agree-split-notions}(2)),
and let $\mu_K$ and $\mu$ be Haar measures on $K$ and $\mathbf A_E^\times$.
Take the integral of a function that is not integrable to be $0$. Then there
are $C_0,C_1\neq0$ such that the following holds for all Hecke characters
$\chi_X,\chi_W$ and representations $\omega,\omega_1,\omega_0$ for which
$(X,W,M,\chi_X,\chi_W,\omega)$, $(X,W_1,M_1,\chi_X,\chi_W,\omega_1)$ and
$(X_1,W_1,M_0,\chi_X,\chi_W,\omega_0)$ are Weil pairs of signatures
$(r+1,s+1;r'+1,s'+1)$, $(r+1,s+1;r',s')$ and $(r,s;r',s')$
(Definition~\ref{asm-theta-agree-triple}(1)). Let $\operatorname{cone}$ be the
cone for $(M,\omega,e,K,\mu_K,\mu,u_\alpha,\chi_W^{-1},z)$
(Definition~\ref{asm-theta-agree-planes}(7)). For every Schwartz-Bruhat
function $\Phi$ on $\mathbf A^{n\times d}$ that is $K$-finite for $\omega$,
there are functions $A_0$ and $A_1$ on $U_W$ such that $A_0+A_1$ is the
constant term of $I_{M,\omega}(\Phi)$ along $H'$, $A_0$ has the law
$(\chi_X,n/2)$, $A_1$ has the law $(\chi_X/\chi_W,(d-n)/2)$, and
\[
A_1(ph)=C_1\,I_{M_0,\omega_0}(\operatorname{cone}(\omega(1,h)\Phi))(h_1),\qquad
A_0(ph)=C_0\,I_{M_1,\omega_1}(e_0(\omega(1,h)\Phi))(h_1)
\]
for every Levi element $p$ for $(1,h_1)$ and every $h\in U_W$.
\end{lemma}
\leanref{Development/WeilBound/ThetaLifts/SmallCones/ThetaUniq.lean}{PositivityTheory.ThKCT.GenBig}
\leanref{Development/WeilBound/ThetaLifts/SmallCones/GenAnisoOf.lean}{PositivityTheory.ThKCT.genBig_of}
\begin{proof}
If there are no $\chi_X,\chi_W,\omega$ for which
$(X,W,M,\chi_X,\chi_W,\omega)$ is a Weil pair of signature
$(r+1,s+1;r'+1,s'+1)$, put $C_0=C_1=1$. Otherwise fix such data; then $X$ is
Hermitian of signature $(r+1,s+1)$ and $M$ is regular. Let $C_0'\neq0$ be given
by Lemma~\ref{asm-theta-agree-split-kernel} for $(X,W,W_1,M,M_1,H',e_0)$,
and $C_1'\neq0$ by its part (2) for $(H,M_0,e)$. Let $\kappa\neq0$ be given by
Lemma~\ref{asm-theta-agree-split-inputs}(10) for $e$, $z$, $K$, $\mu_K$ and
$\mu$, and $c\neq0$ by Lemma~\ref{asm-theta-agree-split-domains} for the
measures and domains of $M$ and $M_1$. As $x_H\neq0$ and $h_X(x_H,x_H)=0$,
Lemma~\ref{asm-theta-agree-split-witt}(2) gives a transporter family $\gamma$
for $x_H$. Put $C_0=C_0'c$ and $C_1=C_1'\kappa$.

Let the characters and representations give the three Weil pairs, and let
$\Phi$ be Schwartz-Bruhat and $K$-finite. Write $\mu_M$ and $\mathcal D$ for
the measure and domain of $M$, $\rho$ for its Schrödinger operators, and put
\[
S(g,h)=\sum_{v\in\operatorname{Iso}(X)}\theta\bigl(e(\rho(x_Hy_{H'}^{\mathsf T})\,\omega(\gamma(v)g,h)\Phi)\bigr).
\]
Let $I^{\mathrm{cone}}$ be the cone integral for
$(M,\omega,M_0,\omega_0,e,K,\mu_K,\mu,u_\alpha,\chi_W^{-1},z)$, and put
\[
A_0(h)=\int_{\mathcal D}C_0'\,\theta(e_0(\omega(g,h)\Phi))\,d\mu_M(g),\qquad
A_1(h)=C_1'\kappa\,I^{\mathrm{cone}}(\omega(1,h)\Phi)(1).
\]

By Lemma~\ref{asm-theta-agree-split-zero}, $A_0$ has the law $(\chi_X,n/2)$ and
$A_0(ph)=C_0'c\,I_{M_1,\omega_1}(e_0(\omega(1,h)\Phi))(h_1)$. By
Lemma~\ref{asm-theta-agree-split-inputs}(10)(b) with $h_1=1$, $A_1$ has the law
$(\chi_X/\chi_W,(d-n)/2)$; by (10)(b) with $t=1$ and (10)(c),
$A_1(ph)=C_1'\kappa\,I_{M_0,\omega_0}(\operatorname{cone}(\omega(1,h)\Phi))(h_1)$.

It remains to compute the constant term. As $2n=d+2z$ and
$\min(r+1,s+1)=\min(r,s)+1<z$, we have $d+2\min(r+1,s+1)<2n$, and
Lemma~\ref{asm-theta-agree-convergence}(1) gives $C$ and $k$ with
$\int_{\mathcal D}|\theta_\omega(\Phi)(g,h)|\,d\mu_M(g)\leq C\operatorname{ht}(h)^k$.
By Lemma~\ref{asm-theta-agree-split-kernel}(2), for every $g$ the constant
term of $\theta_\omega(\Phi)(g,\cdot)$ along $H'$ is
$h\mapsto C_0'\theta(e_0(\omega(g,h)\Phi))+C_1'S(g,h)$. By
Lemma~\ref{asm-theta-agree-split-fubini}, this is integrable in $g$ on
$\mathcal D$ for every $h$, and the constant term of $I_{M,\omega}(\Phi)$ is
$h\mapsto\int_{\mathcal D}(C_0'\theta(e_0(\omega(g,h)\Phi))+C_1'S(g,h))\,d\mu_M(g)$.
By Lemma~\ref{asm-theta-agree-split-inputs}(10)(a), $g\mapsto S(g,h)$ is
integrable on $\mathcal D$ with integral $\kappa I^{\mathrm{cone}}(\omega(1,h)\Phi)(1)$.
So $g\mapsto C_0'\theta(e_0(\omega(g,h)\Phi))$ is integrable, and the integral
splits as $A_0(h)+A_1(h)$.
\end{proof}

\begin{lemma}\label{asm-theta-agree-split-small}
Let $E$ be a CM field with complex conjugation $c$ and a theta base $B$
(Definition~\ref{asm-theta-agree-models}(1)). Let $H$ be an adapted plane in a
lattice space $X$ over $X_1$ and $H'$ one in $W$ over $W_1$, with ranks $n$,
$n_1$, $d$ and $d_1$. Let $M$ be a model for $(X,W)$, $M_0$ and $M_1$ regular
models for $(X_1,W_1)$ and $(X,W_1)$, $e_0$ an evaluation map for
$(M,M_1,1,\iota_{H'},\mathcal J_0(H'))$ and $e$ one for
$(M,M_0,\iota_H,\iota_{H'},\mathcal J_\beta(H,H'))$
(Definition~\ref{asm-theta-agree-planes}). Let $r,s,r',s'\geq0$ and $z\geq0$ be
integers with $d+2=2n+2z$ and $2\leq z$, let $K\subset U_X$ be an Iwasawa
subgroup for $x_H$ (Definition~\ref{asm-theta-agree-split-notions}(2)), and
let $\mu_K$ and $\mu$ be Haar measures on $K$ and $\mathbf A_E^\times$. Take
the integral of a function that is not integrable to be $0$. Then there are
$C_0,C_1\neq0$ such that the following holds for all Hecke characters
$\chi_X,\chi_W$ and representations $\omega,\omega_1,\omega_0$ for which
$(X,W,M,\chi_X,\chi_W,\omega)$, $(X,W_1,M_1,\chi_X,\chi_W,\omega_1)$ and
$(X_1,W_1,M_0,\chi_X,\chi_W,\omega_0)$ are Weil pairs of signatures
$(r+1,s+1;r'+1,s'+1)$, $(r+1,s+1;r',s')$ and $(r,s;r',s')$
(Definition~\ref{asm-theta-agree-triple}(1)). Let $\operatorname{cone}$ be the
cone for $(M,\omega,e,K,\mu_K,\mu,u_\beta,\chi_W,z)$
(Definition~\ref{asm-theta-agree-planes}(7)). For every
$\Phi\in\mathcal S^\circ(M)$ (Definition~\ref{asm-theta-agree-models}(5)) that
is $K$-finite for $\omega$, there are functions $A_0$ and $A_1$ on $U_W$ such
that $A_0+A_1$ is the constant term of $I_{M,\omega}(\Phi)$ along $H'$, $A_0$
has the law $(\chi_X,n/2)$, $A_1$ has the law $(\chi_X/\chi_W,(d-n)/2)$, and
for every Levi element $p$ for $(1,h_1)$ and every $h\in U_W$,
\[
A_1(ph)=C_1\,I_{M_0,\omega_0}(\operatorname{cone}(\omega(1,h)\Phi))(h_1),
\]
and, if $e_0(\omega(1,h)\Phi)\in\mathcal S^\circ(M_1)$,
$A_0(ph)=C_0\,I_{M_1,\omega_1}(e_0(\omega(1,h)\Phi))(h_1)$.
\end{lemma}
\leanref{Development/WeilBound/ThetaLifts/SmallCones/ThetaUniq.lean}{PositivityTheory.ThKCT.GenSmall}
\leanref{Development/WeilBound/ThetaLifts/SmallCones/GenSmallOf.lean}{PositivityTheory.ThKCT.genSmall_of}
\leanref{Development/WeilBound/ThetaLifts/SmallCones/IsoActFree.lean}{PositivityTheory.ThKCT.orbit_sum}
\leanref{Development/WeilBound/ThetaLifts/SmallCones/IsoActFree.lean}{PositivityTheory.ThKCT.units_sum}
\begin{proof}
If there are no $\chi_X,\chi_W,\omega$ for which
$(X,W,M,\chi_X,\chi_W,\omega)$ is a Weil pair of signature
$(r+1,s+1;r'+1,s'+1)$, put $C_0=C_1=1$. Otherwise fix such data; then $X$ is
Hermitian of signature $(r+1,s+1)$, $W$ skew-Hermitian of signature
$(r'+1,s'+1)$, and $M$ regular. By Lemma~\ref{asm-theta-agree-evaluation-cone}
there is an evaluation map $e_\alpha$ for
$(M,M_0,\iota_H,\iota_{H'},\mathcal J_\alpha(H,H'))$. Let $C_0'\neq0$ be given
by Lemma~\ref{asm-theta-agree-split-kernel} for $(X,W,W_1,M,M_1,H',e_0)$,
$C_1'\neq0$ by its part (2) for $(H,M_0,e_\alpha)$, $c_2\neq0$ by
Lemma~\ref{asm-theta-agree-split-inputs}(8) for $e_\alpha$ and $e_\beta=e$,
$\kappa\neq0$ by Lemma~\ref{asm-theta-agree-split-inputs}(11) for $e$, $z$,
$K$, $\mu_K$ and $\mu$, and $\kappa_1\neq0$ by
Lemma~\ref{asm-theta-agree-split-domains} for the measures and domains of $M$
and $M_1$. Lemma~\ref{asm-theta-agree-split-witt}(2) gives a transporter family
$\gamma$ for $x_H$. Put $C_0=C_0'\kappa_1$ and $C_1=C_1'c_2\kappa$.

Let the characters and representations give the three Weil pairs, and let
$\Phi\in\mathcal S^\circ(M)$ be $K$-finite; it is a Schwartz-Bruhat function.
Write $x=x_H$, $y=y_H$, $x'=x_{H'}$, $y'=y_{H'}$, $\mu_M$ and $\mathcal D$ for
the measure and domain of $M$ and $\rho$ for its Schrödinger operators, and put
\[
a(g,h)=\theta\bigl(e_\alpha(\rho(xy'^{\mathsf T})\omega(g,h)\Phi)\bigr),\qquad
b(g,h)=\theta\bigl(e(\rho(yx'^{\mathsf T})\omega(g,h)\Phi)\bigr).
\]
The group $E^\times$ acts freely on $\operatorname{Iso}(X)$; choose a
representative $v_q$ of each orbit $q$, and put
\[
L(g,h)=-\sum_q\theta\bigl(e_\alpha(\omega(\gamma(v_q)g,h)\Phi)\bigr),
\]
the sum over the orbits. By Lemma~\ref{asm-theta-agree-split-inputs}(6),
$\omega(g,h)\Phi\in\mathcal S^\circ(M)$ for all $g$ and $h$.

\emph{The two sums.} Let $g$ and $h$ be such that
$\sum_{v\in\operatorname{Iso}(X)}b(\gamma(v)g,h)$ converges absolutely. Then
\[
\sum_{v\in\operatorname{Iso}(X)}a(\gamma(v)g,h)-c_2\sum_{v\in\operatorname{Iso}(X)}b(\gamma(v)g,h)=L(g,h),\tag{L}
\]
the sum of the $a$ converging absolutely by
Lemma~\ref{asm-theta-agree-split-kernel}(2). Indeed, $\operatorname{Iso}(X)$
is the disjoint union of the orbits $E^\times v_q$, so the absolutely
convergent left side is the sum over $q$ of the sums over $\xi\in E^\times$ of
the terms at $\xi v_q$. Fix $q$ and put $\Psi=\omega(\gamma(v_q)g,h)\Phi$. By
Lemma~\ref{asm-theta-agree-split-shift}(2),
\[
a(\gamma(\xi v_q)g,h)=\theta\bigl(e_\alpha(\rho(\xi\,xy'^{\mathsf T})\Psi)\bigr),\qquad
b(\gamma(\xi v_q)g,h)=\theta\bigl(e(\rho(c(\xi)^{-1}yx'^{\mathsf T})\Psi)\bigr).
\]
As $\xi\mapsto c(\xi)^{-1}$ is a bijection of $E^\times$, the sum over the
orbit is
$\sum_{\xi\in E^\times}\theta(e_\alpha(\rho(\xi xy'^{\mathsf T})\Psi))-c_2\sum_{\xi\in E^\times}\theta(e(\rho(\xi yx'^{\mathsf T})\Psi))$.
By Lemma~\ref{asm-theta-agree-split-inputs}(8) the sums over $\xi\in E$
converge absolutely and the first is $c_2$ times the second. Removing the terms
at $\xi=0$, where $\rho(0)=1$, the sum over the orbit is
$-\theta(e_\alpha(\Psi))+c_2\,\theta(e(\Psi))$. As $\Psi\in\mathcal S^\circ(M)$,
$e(\Psi)=0$ by Lemma~\ref{asm-theta-agree-split-inputs}(7)(a). Summing over
$q$ gives (L).

\emph{The constant term.} By Lemma~\ref{asm-theta-agree-decay-full} with $N=0$
there are $C$ and $k$ with
$|\theta_\omega(\Phi)(g,h)|\leq C\operatorname{ht}(h)^k$, and
$\mu_M(\mathcal D)<\infty$ by Theorem~\ref{asm-theta-agree-continuity-volume},
so $\int_{\mathcal D}|\theta_\omega(\Phi)(g,h)|\,d\mu_M(g)\leq C\mu_M(\mathcal D)\operatorname{ht}(h)^k$.
By Lemma~\ref{asm-theta-agree-split-kernel}(2), for every $g$ the constant
term of $\theta_\omega(\Phi)(g,\cdot)$ along $H'$ is
$h\mapsto C_0'\theta(e_0(\omega(g,h)\Phi))+C_1'\sum_va(\gamma(v)g,h)$. By
Lemma~\ref{asm-theta-agree-split-fubini}, this is integrable in $g$ on
$\mathcal D$, and its integral over $\mathcal D$ is the constant term of
$I_{M,\omega}(\Phi)$. Let $I^{\mathrm{cone}}$ be the cone integral for
$(M,\omega,M_0,\omega_0,e,K,\mu_K,\mu,u_\beta,\chi_W,z)$. By
Lemma~\ref{asm-theta-agree-split-inputs}(11)(a), for almost every
$g\in\mathcal D$ the sum $\sum_vb(\gamma(v)g,h)$ converges absolutely; it is
integrable on $\mathcal D$, with integral $\kappa I^{\mathrm{cone}}(\omega(1,h)\Phi)(1)$.
Put
\[
A_0(h)=\int_{\mathcal D}\bigl(C_0'\theta(e_0(\omega(g,h)\Phi))+C_1'L(g,h)\bigr)\,d\mu_M(g),\qquad
A_1(h)=C_1'c_2\kappa\,I^{\mathrm{cone}}(\omega(1,h)\Phi)(1).
\]
By (L), the integrand of $A_0$ agrees almost everywhere on $\mathcal D$ with the
difference of the integrable functions
$C_0'\theta(e_0(\omega(g,h)\Phi))+C_1'\sum_va(\gamma(v)g,h)$ and
$C_1'c_2\sum_vb(\gamma(v)g,h)$. So it is integrable, and $A_0+A_1$ is the
constant term of $I_{M,\omega}(\Phi)$ along $H'$.

\emph{The law of $A_0$.} Let $m$ be a Levi element for $(t,1)$, and $g,h$
given. By Lemma~\ref{asm-theta-agree-split-inputs}(4) with $u=1$, applied to
$\omega(g,h)\Phi$ at the element $1$ of $U_X$,
$e_0(\omega(g,mh)\Phi)=\chi_X(t)|t|^{n/2}e_0(\omega(g,h)\Phi)$. By
Lemma~\ref{asm-theta-agree-split-inputs}(5), with the Levi element $1$ of $U_X$
for $(1,1)$, $e_\alpha(\omega(1,m)\Psi)=\chi_X(t)|t|^{n/2}e_\alpha(\Psi)$ for
Schwartz-Bruhat $\Psi$, so $L(g,mh)=\chi_X(t)|t|^{n/2}L(g,h)$. Hence
$A_0(mh)=\chi_X(t)|t|^{n/2}A_0(h)$.

\emph{The law and the values of $A_1$.} By
Lemma~\ref{asm-theta-agree-split-inputs}(11)(b) with $h_1=1$, $A_1$ has the law
$(\chi_X/\chi_W,(d-n)/2)$. For a Levi element $p$ for $(1,h_1)$, (11)(b) with
$t=1$ and (11)(c) give
$A_1(ph)=C_1'c_2\kappa\,I_{M_0,\omega_0}(\operatorname{cone}(\omega(1,h)\Phi))(h_1)$.

\emph{The values of $A_0$.} Let $p$ be a Levi element for $(1,h_1)$, and
suppose $e_0(\omega(1,h)\Phi)\in\mathcal S^\circ(M_1)$. For $g'\in U_X$,
Lemma~\ref{asm-theta-agree-split-inputs}(4) gives
$e_0(\omega(g',ph)\Phi)=\omega_1(g',h_1)e_0(\omega(1,h)\Phi)$, which lies in
$\mathcal S^\circ(M_1)$ by Lemma~\ref{asm-theta-agree-split-inputs}(6) for the
Weil pair of $\omega_1$. By Lemma~\ref{asm-theta-agree-split-inputs}(7)(b), for
the Weil pairs of $\omega_1$ and $\omega_0$, $e_\alpha(\omega(g',ph)\Phi)=0$.
So $L(g,ph)=0$ for every $g$, and by Lemma~\ref{asm-theta-agree-split-zero}
\[
A_0(ph)=\int_{\mathcal D}C_0'\theta(e_0(\omega(g,ph)\Phi))\,d\mu_M(g)=C_0'\kappa_1\,I_{M_1,\omega_1}(e_0(\omega(1,h)\Phi))(h_1).
\]
\end{proof}

\begin{lemma}\label{asm-theta-agree-split-any}
Let $E$ be a CM field with a theta base $B$
(Definition~\ref{asm-theta-agree-models}(1)). Let $X$, $W$ and $W_1$ be lattice
spaces of ranks $n$, $d$ and $d_1$ with $n\geq3$, $H'$ an adapted plane in $W$
over $W_1$, $M$ a model for $(X,W)$, $M_1$ a model for $(X,W_1)$, and $e_0$ an
evaluation map for $(M,M_1,1,\iota_{H'},\mathcal J_0(H'))$
(Definition~\ref{asm-theta-agree-planes}). Let $r,s,r',s'\geq0$, and suppose
that one of the following holds for some integer $z\geq0$:
\begin{enumerate}
\item[(big)] $2n=d+2z$ and $\min(r,s)+1\leq z$;
\item[(small)] $d+2=2n+2z$ and $2\leq z$.
\end{enumerate}
Take the integral of a function that is not integrable to be $0$. Then there
is $C_0\neq0$ such that the following holds for all Hecke characters
$\chi_X,\chi_W$ and representations $\omega,\omega_1$ for which
$(X,W,M,\chi_X,\chi_W,\omega)$ and $(X,W_1,M_1,\chi_X,\chi_W,\omega_1)$ are Weil
pairs of signatures $(r,s;r'+1,s'+1)$ and $(r,s;r',s')$
(Definition~\ref{asm-theta-agree-triple}(1)). Let
$\Phi\in\operatorname{PD}(M)$, with $\Phi\in\mathcal S^\circ(M)$ in the case
(small) (Definition~\ref{asm-theta-agree-models}). Then there are functions
$A_0$ and $A_1$ on $U_W$ such that $A_0+A_1$ is the constant term of
$I_{M,\omega}(\Phi)$ along $H'$, $A_0$ has the law $(\chi_X,n/2)$, $A_1$ has
the law $(\chi_X/\chi_W,(d-n)/2)$, and
$A_0(ph)=C_0\,I_{M_1,\omega_1}(e_0(\omega(1,h)\Phi))(h_1)$ for every Levi
element $p$ for $(1,h_1)$ and every $h\in U_W$, in the case (small) whenever
$e_0(\omega(1,h)\Phi)\in\mathcal S^\circ(M_1)$.
\end{lemma}
\leanref{Development/WeilBound/ThetaLifts/SmallCones/ThetaKernelGeneralIdentity.lean}{PositivityTheory.ThKCT.AnySplit}
\leanref{Development/WeilBound/ThetaLifts/SmallCones/ThetaKernelGeneralIdentity.lean}{PositivityTheory.ThKCT.theta_kernel_general_identity}
\leanref{Development/WeilBound/ThetaLifts/SmallCones/ThetaKernelGeneralIdentity.lean}{PositivityTheory.ThKCT.anyBig}
\leanref{Development/WeilBound/ThetaLifts/SmallCones/ThetaKernelGeneralIdentity.lean}{PositivityTheory.ThKCT.anySmall}
\begin{proof}
\emph{The case $\min(r,s)=0$.} Let $C$ be given by
Lemma~\ref{asm-theta-agree-split-aniso} for $(X,W,W_1,H',M,M_1,e_0)$ and the
signatures $(r,s;r'+1,s'+1)$ and $(r,s;r',s')$, and put $C_0=C$. Let the
characters and representations give the two Weil pairs, and let
$\Phi\in\operatorname{PD}(M)$. By
Lemma~\ref{asm-theta-agree-split-inputs}(2), $\Phi$ is a Schwartz-Bruhat
function. Lemma~\ref{asm-theta-agree-split-aniso} gives $A_0$ with the
required properties, and $A_1=0$ has every law.

\emph{The case $r,s\geq1$.} If there are no $\chi_X,\chi_W,\omega,\omega_1$
with the two Weil pairs, put $C_0=1$. Otherwise fix such data; then $X$ is
Hermitian of signature $(r,s)$, $W_1$ is skew-Hermitian of signature $(r',s')$,
and $M$ and $M_1$ are regular. Write $r=r_0+1$ and $s=s_0+1$, and $J_X$ for the
matrix of $X$. By Proposition~\ref{asm-theta-agree-tower-plane} for $X$, with
$n_1=n-2\geq1$ and the positions $0<1$, there are lattice spaces $X^\circ$ and
$X_1$ such that $X[X^\circ]$ is Hermitian of signature $(r,s)$, $X_1$ is
Hermitian of signature $(r_0,s_0)$, and there is an adapted plane $H$ in
$X[X^\circ]$ over $X_1$ at the positions $0<1$. Let $M^\circ$ and $M_1^\circ$ be
$M$ and $M_1$ regarded as models for $(X[X^\circ],W)$ and
$(X[X^\circ],W_1)$, with the same isomorphisms, vacua, measures and domains.
They have the same Schrödinger operators and theta integrals as $M$ and $M_1$,
and the same spaces of full-rank data, which depend only on the matrices, the
isomorphisms and the vacua; $e_0$ is an evaluation map for
$(M^\circ,M_1^\circ,1,\iota_{H'},\mathcal J_0(H'))$; and a Weil pair
$(X,W,M,\chi_X,\chi_W,\omega)$ of signature $(r,s;\cdot)$ gives a Weil pair
$(X[X^\circ],W,M^\circ,\chi_X,\chi_W,\omega)$ of the same signature, and
likewise for $M_1$, as $X[X^\circ]$ is Hermitian of signature $(r,s)$.

By Lemma~\ref{asm-theta-agree-tower-model} there is a regular model $M_0$ for
$(X_1,W_1)$. By Lemma~\ref{asm-theta-agree-evaluation-cone}, for the fixed Weil
pair of $M^\circ$, there are evaluation maps $e_\alpha$ and $e_\beta$ for
$(M^\circ,M_0,\iota_H,\iota_{H'},\mathcal J_\alpha(H,H'))$ and
$(M^\circ,M_0,\iota_H,\iota_{H'},\mathcal J_\beta(H,H'))$. Let $K$ be the
lattice stabilizer of $X$, with Haar measures $\mu_K$ and $\mu$ on $K$ and
$\mathbf A_E^\times$ (Lemma~\ref{asm-theta-agree-stabilizer}). As $x_H\neq0$
and $x_H^{\mathsf T}J_Xc(x_H)=0$, $K$ is an Iwasawa subgroup for $x_H$
(Lemma~\ref{asm-theta-agree-split-inputs}(3)). In the case (big), apply
Lemma~\ref{asm-theta-agree-split-big} to
$(H,H',M^\circ,M_0,M_1^\circ,e_0,e_\alpha)$ with the integers
$r_0,s_0,r',s'$ and $z$, which satisfy $\min(r_0,s_0)+2=\min(r,s)+1\leq z$; in
the case (small), apply Lemma~\ref{asm-theta-agree-split-small} to
$(H,H',M^\circ,M_0,M_1^\circ,e_0,e_\beta)$ with the same integers, and $K$,
$\mu_K$ and $\mu$ in both cases. Let $C_0$ be the constant it gives.

Let $\chi_X,\chi_W,\omega,\omega_1$ give the two Weil pairs, and $\Phi$ be as
in the statement. As $\chi_X|_{\mathbf A^\times}=\eta^n=\eta^{n-2}$ and
$\chi_W|_{\mathbf A^\times}=\eta^{d_1}$, Theorem~\ref{asm-theta-agree-tower-weil}
gives $\omega_0$ with $(X_1,W_1,M_0,\chi_X,\chi_W,\omega_0)$ a Weil pair of
signature $(r_0,s_0;r',s')$. By Lemma~\ref{asm-theta-agree-split-inputs}(2),
$\Phi$ is Schwartz-Bruhat and $K$-finite for $\omega$. The three Weil pairs of
$M^\circ$, $M_1^\circ$ and $M_0$ have the signatures required in
Lemma~\ref{asm-theta-agree-split-big} or~\ref{asm-theta-agree-split-small}, and
$\Phi\in\mathcal S^\circ(M^\circ)$ in the case (small). That lemma gives
$A_0$ and $A_1$, and, as the theta integrals and the full-rank data of
$M^\circ$ and $M_1^\circ$ are those of $M$ and $M_1$, they have the required
properties.
\end{proof}

\begin{lemma}\label{asm-theta-agree-split}
Let $(B,G,\mathcal X,\omega,\omega')$ be a triple for a theta base $B$ of a CM
field $E$ (Definition~\ref{asm-theta-agree-triple}), and let $H'$ be an adapted
plane in $W$ over a lattice space $W_1$. Then for every
$\Phi\in\operatorname{PD}(M)$ there is exactly one pair $(f_\alpha,f_\beta)$
that splits $I(\Phi)$ along $H'$
(Definition~\ref{asm-theta-agree-planes}(3)), and for every
$\Phi'\in\operatorname{PD}(M')\cap\mathcal S^\circ(M')$ there is exactly one pair
that splits $I'(\Phi')$ along $H'$.
\end{lemma}
\leanref{Development/AbelianVarietyTheory/Definitions/Divisors/CTSmallCone.lean}{GlobalTheory.ThK.SplitEx}
\leanref{Development/WeilBound/ThetaLifts/SmallCones/Leaves.lean}{PositivityTheory.ThKCT.splitEx_of}
\leanref{Development/WeilBound/ThetaLifts/SmallCones/ThetaExp.lean}{PositivityTheory.ThKCT.ct_all}
\begin{proof}
Let the triple have ranks $(a,a')$, integers $p,q$, spaces $V,V',W$, models
$M,M'$ and character datum $\mathcal X=(\alpha,\beta,S_0,Q)$. By
Definition~\ref{asm-theta-agree-triple}(4), $(V,W,M,\beta,\beta/\alpha,\omega)$
is a Weil pair of signature $(a-p,p;a+a'-p-q,p+q)$ and
$(V',W,M',\alpha,\alpha/\beta,\omega')$ one of signature
$(q,a'-q;a+a'-p-q,p+q)$. As $G$ is proper, $a'<a$, $a+a'$ is even,
$2(p+q)+1<\lfloor(a-a')/2\rfloor$ and $2(p+q)+1<a'$. By
Lemma~\ref{asm-theta-agree-split-low} for $H'$, $p+q\geq1$ and $W_1$ is
skew-Hermitian of signature $(a+a'-p-q-1,p+q-1)$; so $a'\geq4$ and
$a>a'$. By Lemma~\ref{asm-theta-agree-split-unique}, a function on $U_W$
has at most one splitting along $H'$. It remains to find one.

\emph{The first theta integral.} By Lemma~\ref{asm-theta-agree-tower-model}
there is a regular model $M_1$ for $(V,W_1)$, and by
Theorem~\ref{asm-theta-agree-tower-weil} a Weil representation $\omega_1$ for
$M_1$ with characters $(\beta,\beta/\alpha)$, as
$(\beta/\alpha)|_{\mathbf A^\times}=\eta^{a+a'}=\eta^{a+a'-2}$. So
$(V,W_1,M_1,\beta,\beta/\alpha,\omega_1)$ is a Weil pair of signature
$(a-p,p;a+a'-p-q-1,p+q-1)$. By Lemma~\ref{asm-theta-agree-evaluation} there is
an evaluation map $e_0$ for $(M,M_1,1,\iota_{H'},\mathcal J_0(H'))$. Put
$z=(a-a')/2$. Then $2a=(a+a')+2z$, and $\min(a-p,p)+1\leq p+1\leq z$, as
$2p+1<z$. Lemma~\ref{asm-theta-agree-split-any}, in the case (big), with
$n=a$, $d=a+a'$, $X=V$, $(r,s)=(a-p,p)$ and $(r'+1,s'+1)=(a+a'-p-q,p+q)$,
gives for $\Phi\in\operatorname{PD}(M)$ functions $A_0$ and $A_1$ with
$A_0+A_1$ the constant term of $I(\Phi)$ along $H'$, $A_0$ of law
$(\beta,a/2)$ and $A_1$ of law $(\beta/(\beta/\alpha),((a+a')-a)/2)=(\alpha,a'/2)$.
So $(A_1,A_0)$ splits $I(\Phi)$ along $H'$.

\emph{The second theta integral.} In the same way there are a regular model
$M'_1$ for $(V',W_1)$, a Weil representation $\omega'_1$ for it with
characters $(\alpha,\alpha/\beta)$, and an evaluation map $e'_0$ for
$(M',M'_1,1,\iota_{H'},\mathcal J_0(H'))$. Put $z=(a-a')/2+1$. Then
$(a+a')+2=2a'+2z$ and $z\geq2$. Lemma~\ref{asm-theta-agree-split-any}, in the
case (small), with $n=a'$, $d=a+a'$, $X=V'$ and $(r,s)=(q,a'-q)$, gives for
$\Phi'\in\operatorname{PD}(M')\cap\mathcal S^\circ(M')$ functions $A_0$ and
$A_1$ with $A_0+A_1$ the constant term of $I'(\Phi')$, $A_0$ of law
$(\alpha,a'/2)$ and $A_1$ of law
$(\alpha/(\alpha/\beta),((a+a')-a')/2)=(\beta,a/2)$. So $(A_0,A_1)$ splits
$I'(\Phi')$ along $H'$.
% Lean: (PositivityTheory.ThKCT.ct_all Ext).2.1, where Ext is the record ThKCT.Ext of
% mixW_closed, mixX_closed, thetaBasic, EvZero, EvCone, StabEx, Tow.pairEx_of,
% Tow.planeHerm, grpTop_of, epsOf_of, triplePair_holds, WeilConv, domWeil_of, soSemi_of,
% soDecay, finVolU, soFin_of; the component is splitEx_of (Leaves.lean). What is not yet
% written out of the 12,800 lines shared with the four Lemmas asm-theta-agree-ct-* is in
% Lemma asm-theta-agree-split-inputs.
\end{proof}

\begin{lemma}\label{asm-theta-agree-ct-beta-zero}
Let $B$ be a theta base of a CM field $E$, let $(G,G_1,H,H')$ be a
$\beta$-configuration (Definition~\ref{asm-theta-agree-planes}(4)) with $G$ of
ranks $(a,b+2)$, and let $e$ be an evaluation map for
$(M,M_1,1,\iota_{H'},\mathcal J_0(H'))$. Then there is $C\neq0$ such that the
following holds for all triples $(B,G,\mathcal X,\omega,\omega')$ and
$(B,G_1,\mathcal X,\omega_1,\omega'_1)$ with theta integrals $I$ and $I_1$
(Definition~\ref{asm-theta-agree-triple}). If $\Phi\in\operatorname{PD}(M)$,
$(f_\alpha,f_\beta)$ splits $I(\Phi)$ along $H'$, $h_1\in U_{W_1}$, $m$ is a
Levi element for $(1,h_1)$ and $h\in U_W$, then
$f_\beta(mh)=C\,I_1\bigl(e(\omega(1,h)\Phi)\bigr)(h_1)$.
\end{lemma}
\leanref{Development/AbelianVarietyTheory/Definitions/Divisors/CTSmallCone.lean}{GlobalTheory.ThK.CTBigZero}
\leanref{Development/WeilBound/ThetaLifts/SmallCones/Leaves.lean}{PositivityTheory.ThKCT.ctBigZero_of}
\begin{proof}
Let $G=(p,q,V,V',W,M,M')$ and $G_1=(p_1,q_1,V_1,V'_1,W_1,M_1,M'_1)$; as the
configuration is a $\beta$-configuration, $V_1=V$, $p_1=p$ and $q=q_1+1$. If
there is no triple $(B,G,\mathcal X,\omega,\omega')$, put $C=1$. Otherwise fix
one. As $G$ is proper, $b+2<a$, $a+b$ is even,
$2(p+q)+1<\lfloor(a-b-2)/2\rfloor$ and $2(p+q)+1<b+2$. Put
$z=(a-b-2)/2$; then $2a=(a+b+2)+2z$, $\min(a-p,p)+1\leq z$ and $a\geq3$. Let
$C$ be the constant $C_0$ of Lemma~\ref{asm-theta-agree-split-any}, in the case
(big), for $n=a$, $d=a+b+2$, $d_1=a+b$, $X=V$, $W$, $W_1$, $H'$, $M$, $M_1$,
$e$ and the integers $(r,s)=(a-p,p)$ and $(r',s')=(a+b-p-q_1,p+q_1)$.

Let the two triples, $\Phi$, $(f_\alpha,f_\beta)$, $h_1$, $m$ and $h$ be as in
the statement. By Definition~\ref{asm-theta-agree-triple}(4),
$(V,W,M,\beta,\beta/\alpha,\omega)$ is a Weil pair of signature
$(a-p,p;a+b+2-p-q,p+q)=(r,s;r'+1,s'+1)$, and
$(V,W_1,M_1,\beta,\beta/\alpha,\omega_1)$ one of signature $(r,s;r',s')$.
Lemma~\ref{asm-theta-agree-split-any} gives $A_0$ and $A_1$ with $A_0+A_1$ the
constant term of $I(\Phi)$ along $H'$, $A_0$ of law $(\beta,a/2)$ and $A_1$ of
law $(\alpha,(b+2)/2)$, and $A_0(mh)=C\,I_1(e(\omega(1,h)\Phi))(h_1)$. So
$(A_1,A_0)$ splits $I(\Phi)$, and by Lemma~\ref{asm-theta-agree-split-unique},
$f_\beta=A_0$.
% Lean: PositivityTheory.ThKCT.ctBigZero_of GenBig GenAniso KFinData StabEx IwaEx PairEx
% PlaneHerm EvCone UnipEx TriplePair, through anyBig and split_uniq.
\end{proof}

\begin{lemma}\label{asm-theta-agree-ct-beta-cone}
Let $B$ be a theta base of a CM field $E$, let $(G,G_1,H,H')$ be a
$\beta$-configuration (Definition~\ref{asm-theta-agree-planes}(4)) with $G$ of
ranks $(a,b+2)$, let $e'$ be an evaluation map for
$(M',M'_1,\iota_H,\iota_{H'},\mathcal J_\beta(H,H'))$, let $K\subset U_{V'}$ be the
lattice stabilizer of $V'$, and let $\mu_K$ and $\mu$ be Haar measures on $K$ and
$\mathbf A_E^\times$. Then there is $C\neq0$ such that the following holds for
all triples $(B,G,\mathcal X,\omega,\omega')$ and
$(B,G_1,\mathcal X,\omega_1,\omega'_1)$ with theta integrals $I'$ and $I'_1$ of
the second spaces (Definition~\ref{asm-theta-agree-triple}). Let
$\operatorname{cone}_\beta$ be the cone for $(M',\omega',e',K,\mu_K,\mu,u_\beta,\nu)$
and $z=k_\beta+1$, and take the integral of a function that is not integrable
to be $0$. If $\Phi'\in\operatorname{PD}(M')\cap\mathcal S^\circ(M')$,
$(f'_\alpha,f'_\beta)$ splits $I'(\Phi')$ along $H'$, $h_1\in U_{W_1}$, $m$ is a
Levi element for $(1,h_1)$ and $h\in U_W$, then
$f'_\beta(mh)=C\,I'_1\bigl(\operatorname{cone}_\beta(\omega'(1,h)\Phi')\bigr)(h_1)$.
\end{lemma}
\leanref{Development/AbelianVarietyTheory/Definitions/Divisors/CTSmallCone.lean}{GlobalTheory.ThK.CTSmallCone}
\leanref{Development/WeilBound/ThetaLifts/SmallCones/Leaves.lean}{PositivityTheory.ThKCT.ctSmallCone_of}
\begin{proof}
Let $G=(p,q,V,V',W,M,M')$ and $G_1=(p_1,q_1,V_1,V'_1,W_1,M_1,M'_1)$, with
$p_1=p$ and $q=q_1+1$, and let $H$ be in $V'$ over $V'_1$. If there are no
triples $(B,G,\mathcal X,\omega,\omega')$ and
$(B,G_1,\mathcal X,\omega_1,\omega'_1)$ with a common character datum, put
$C=1$. Otherwise fix such a pair. As $G$ is proper, $b+2<a$, $a+b$ is even,
$2(p+q)+1<\lfloor(a-b-2)/2\rfloor$ and $2(p+q)+1<b+2$; as $G_1$ is proper,
$M'_1$ is regular. By Lemma~\ref{asm-theta-agree-tower-model} there is a
regular model $M_2$ for $(V',W_1)$, and by Lemma~\ref{asm-theta-agree-evaluation}, for the Weil pair
$(V',W,M',\alpha,\alpha/\beta,\omega')$ of the fixed triple, an evaluation map
$e_0$ for $(M',M_2,1,\iota_{H'},\mathcal J_0(H'))$. As $x_H\neq0$ and $x_H$ is
isotropic, $K$ is an Iwasawa subgroup for $x_H$
(Lemma~\ref{asm-theta-agree-split-inputs}(3)). Put $z=k_\beta+1=(a-b-2)/2+1$;
then $(a+b+2)+2=2(b+2)+2z$ and $z\geq2$. Let $C$ be the constant $C_1$ of
Lemma~\ref{asm-theta-agree-split-small} for $n=b+2$, $d=a+b+2$, $n_1=b$,
$d_1=a+b$, $X=V'$, $W$, $X_1=V'_1$, $W_1$, $H$, $H'$, $M'$, $M_0=M'_1$,
$M_1=M_2$, $e_0$, $e'$, the integers $(r,s)=(q_1,b-q_1)$ and
$(r',s')=(a+b-p-q_1,p+q_1)$, $z$, $K$, $\mu_K$ and $\mu$.

Let the two triples, $\Phi'$, $(f'_\alpha,f'_\beta)$, $h_1$, $m$ and $h$ be as
in the statement. By Definition~\ref{asm-theta-agree-triple}(4),
$(V',W,M',\alpha,\alpha/\beta,\omega')$ is a Weil pair of signature
$(q,b+2-q;a+b+2-p-q,p+q)=(r+1,s+1;r'+1,s'+1)$ and
$(V'_1,W_1,M'_1,\alpha,\alpha/\beta,\omega'_1)$ one of signature
$(r,s;r',s')$. Theorem~\ref{asm-theta-agree-tower-weil} gives a Weil
representation $\omega_2$ for $M_2$ with characters $(\alpha,\alpha/\beta)$,
so that $(V',W_1,M_2,\alpha,\alpha/\beta,\omega_2)$ is a Weil pair of signature
$(r+1,s+1;r',s')$. By Lemma~\ref{asm-theta-agree-split-inputs}(2), $\Phi'$ is
$K$-finite for $\omega'$. Lemma~\ref{asm-theta-agree-split-small} gives $A_0$
and $A_1$ with $A_0+A_1$ the constant term of $I'(\Phi')$ along $H'$, $A_0$ of
law $(\alpha,(b+2)/2)$, $A_1$ of law
$(\alpha/(\alpha/\beta),((a+b+2)-(b+2))/2)=(\beta,a/2)$, and
$A_1(mh)=C\,I'_1(\operatorname{cone}_\beta(\omega'(1,h)\Phi'))(h_1)$, the cone
being for the character $\alpha/\beta=\nu$. So $(A_0,A_1)$ splits $I'(\Phi')$,
and by Lemma~\ref{asm-theta-agree-split-unique}, $f'_\beta=A_1$.
% Lean: (PositivityTheory.ThKCT.ct_all Ext).2.2.2.2.2, which is ctSmallCone_of GenSmall
% KFinData IwaEx PairEx EvZero UnipEx TriplePair.
\end{proof}

\begin{lemma}\label{asm-theta-agree-ct-alpha-cone}
Let $B$ be a theta base of a CM field $E$, let $(G,G_1,H,H')$ be an
$\alpha$-configuration (Definition~\ref{asm-theta-agree-planes}(4)) with $G$ of
ranks $(b+2,a')$, let $e$ be an evaluation map for
$(M,M_1,\iota_H,\iota_{H'},\mathcal J_\alpha(H,H'))$, let $K\subset U_V$ be the
lattice stabilizer of $V$, and let $\mu_K$ and $\mu$ be Haar measures on $K$ and
$\mathbf A_E^\times$. Then there is $C\neq0$ such that the following holds for
all triples $(B,G,\mathcal X,\omega,\omega')$ and
$(B,G_1,\mathcal X,\omega_1,\omega'_1)$ with theta integrals $I$ and $I_1$
(Definition~\ref{asm-theta-agree-triple}). Let $\operatorname{cone}_\alpha$ be
the cone for $(M,\omega,e,K,\mu_K,\mu,u_\alpha,\nu)$ and $z=k_\alpha$, and take
the integral of a function that is not integrable to be $0$. If
$\Phi\in\operatorname{PD}(M)$, $(f_\alpha,f_\beta)$ splits $I(\Phi)$ along $H'$,
$h_1\in U_{W_1}$, $m$ is a Levi element for $(1,h_1)$ and $h\in U_W$, then
$f_\alpha(mh)=C\,I_1\bigl(\operatorname{cone}_\alpha(\omega(1,h)\Phi)\bigr)(h_1)$.
\end{lemma}
\leanref{Development/AbelianVarietyTheory/Definitions/Divisors/CTSmallCone.lean}{GlobalTheory.ThK.CTBigCone}
\leanref{Development/WeilBound/ThetaLifts/SmallCones/Leaves.lean}{PositivityTheory.ThKCT.ctBigCone_of}
\begin{proof}
Let $G=(p,q,V,V',W,M,M')$ and $G_1=(p_1,q_1,V_1,V'_1,W_1,M_1,M'_1)$; as the
configuration is an $\alpha$-configuration, $V'_1=V'$, $q=q_1=0$ and
$p=p_1+1$, and $H$ is in $V$ over $V_1$. If there are no triples
$(B,G,\mathcal X,\omega,\omega')$ and $(B,G_1,\mathcal X,\omega_1,\omega'_1)$
with a common character datum, put $C=1$. Otherwise fix such a pair. As $G$ is
proper, $a'<b+2$, $b+2+a'$ is even, $2p+1<\lfloor(b+2-a')/2\rfloor$ and
$2p+1<a'$; as $G_1$ is proper, $M_1$ is regular. By
Lemma~\ref{asm-theta-agree-tower-model} there is a regular model $M_2$ for
$(V,W_1)$, and by Lemma~\ref{asm-theta-agree-evaluation}, for the
Weil pair $(V,W,M,\beta,\beta/\alpha,\omega)$ of the fixed triple, an
evaluation map $e_0$ for $(M,M_2,1,\iota_{H'},\mathcal J_0(H'))$. By
Lemma~\ref{asm-theta-agree-split-inputs}(3), $K$ is an Iwasawa subgroup for
$x_H$. Put $z=k_\alpha=(b+2-a')/2$; then $2(b+2)=(b+2+a')+2z$ and
$\min(b-p_1,p_1)+2\leq p_1+2\leq z$. Let $C$ be the constant $C_1$ of
Lemma~\ref{asm-theta-agree-split-big} for $n=b+2$, $d=b+2+a'$, $n_1=b$,
$d_1=b+a'$, $X=V$, $W$, $X_1=V_1$, $W_1$, $H$, $H'$, $M$, $M_0=M_1$,
$M_1=M_2$, $e_0$, $e$, the integers $(r,s)=(b-p_1,p_1)$ and
$(r',s')=(b+a'-p_1,p_1)$, $z$, $K$, $\mu_K$ and $\mu$.

Let the two triples, $\Phi$, $(f_\alpha,f_\beta)$, $h_1$, $m$ and $h$ be as in
the statement. By Definition~\ref{asm-theta-agree-triple}(4),
$(V,W,M,\beta,\beta/\alpha,\omega)$ is a Weil pair of signature
$(b+2-p,p;b+2+a'-p,p)=(r+1,s+1;r'+1,s'+1)$ and
$(V_1,W_1,M_1,\beta,\beta/\alpha,\omega_1)$ one of signature $(r,s;r',s')$.
Theorem~\ref{asm-theta-agree-tower-weil} gives a Weil representation
$\omega_2$ for $M_2$ with characters $(\beta,\beta/\alpha)$, so that
$(V,W_1,M_2,\beta,\beta/\alpha,\omega_2)$ is a Weil pair of signature
$(r+1,s+1;r',s')$. By Lemma~\ref{asm-theta-agree-split-inputs}(2), $\Phi$ is
Schwartz-Bruhat and $K$-finite for $\omega$.
Lemma~\ref{asm-theta-agree-split-big} gives $A_0$ and $A_1$ with $A_0+A_1$ the
constant term of $I(\Phi)$ along $H'$, $A_0$ of law $(\beta,(b+2)/2)$, $A_1$ of
law $(\beta/(\beta/\alpha),((b+2+a')-(b+2))/2)=(\alpha,a'/2)$, and
$A_1(mh)=C\,I_1(\operatorname{cone}_\alpha(\omega(1,h)\Phi))(h_1)$, the cone
being for the character $(\beta/\alpha)^{-1}=\nu$. So $(A_1,A_0)$ splits
$I(\Phi)$, and by Lemma~\ref{asm-theta-agree-split-unique}, $f_\alpha=A_1$.
% Lean: PositivityTheory.ThKCT.ctBigCone_of GenBig KFinData IwaEx PairEx EvZero UnipEx
% TriplePair; the character (beta/alpha)^{-1} is rewritten as nu by inv_div.
\end{proof}

\begin{lemma}\label{asm-theta-agree-ct-alpha-zero}
Let $B$ be a theta base of a CM field $E$, let $(G,G_1,H,H')$ be an
$\alpha$-configuration (Definition~\ref{asm-theta-agree-planes}(4)) with $G$ of
ranks $(b+2,a')$, and let $e'$ be an evaluation map for
$(M',M'_1,1,\iota_{H'},\mathcal J_0(H'))$. Then there is $C\neq0$ such that the
following holds for all triples $(B,G,\mathcal X,\omega,\omega')$ and
$(B,G_1,\mathcal X,\omega_1,\omega'_1)$ with theta integrals $I'$ and $I'_1$ of
the second spaces (Definition~\ref{asm-theta-agree-triple}). For every
$\Phi'\in\operatorname{PD}(M')$ there is $f'_\alpha$ such that
$(f'_\alpha,0)$ splits $I'(\Phi')$ along $H'$ and
$f'_\alpha(mh)=C\,I'_1\bigl(e'(\omega'(1,h)\Phi')\bigr)(h_1)$ for every
$h_1\in U_{W_1}$, every Levi element $m$ for $(1,h_1)$ and every $h\in U_W$.
\end{lemma}
\leanref{Development/AbelianVarietyTheory/Definitions/Divisors/CTSmallCone.lean}{GlobalTheory.ThK.CTSmallZero}
\leanref{Development/WeilBound/ThetaLifts/SmallCones/Leaves.lean}{PositivityTheory.ThKCT.ctSmallZero_of}
\begin{proof}
Let $G=(p,q,V,V',W,M,M')$ and $G_1=(p_1,q_1,V_1,V'_1,W_1,M_1,M'_1)$, with
$V'_1=V'$ and $q=q_1=0$. Let $C$ be given by
Lemma~\ref{asm-theta-agree-split-aniso} for $n=a'$, $X=V'$, $W$, $W_1$, $H'$,
$M'$, $M'_1$, $e'$ and the signatures $(r,s)=(0,a')$,
$(r',s')=(b+2+a'-p,p)$ and $(r_1,s_1)=(b+a'-p_1,p_1)$; here $\min(r,s)=0$.

Let the two triples and $\Phi'\in\operatorname{PD}(M')$ be given. By
Definition~\ref{asm-theta-agree-triple}(4),
$(V',W,M',\alpha,\alpha/\beta,\omega')$ is a Weil pair of signature
$(0,a';r',s')$ and $(V',W_1,M'_1,\alpha,\alpha/\beta,\omega'_1)$ one of
signature $(0,a';r_1,s_1)$. Let $K$ be the lattice stabilizer of $V'$
(Lemma~\ref{asm-theta-agree-stabilizer}); by
Lemma~\ref{asm-theta-agree-split-inputs}(2), $\Phi'$ is Schwartz-Bruhat.
Lemma~\ref{asm-theta-agree-split-aniso} gives a function $A_0$ that is the
constant term of $I'(\Phi')$ along $H'$, has the law $(\alpha,a'/2)$, and
satisfies $A_0(mh)=C\,I'_1(e'(\omega'(1,h)\Phi'))(h_1)$. Put $f'_\alpha=A_0$.
As $0$ has the law $(\beta,a/2)$, $(f'_\alpha,0)$ splits $I'(\Phi')$ along
$H'$.
% Lean: PositivityTheory.ThKCT.ctSmallZero_of GenAniso KFinData StabEx TriplePair.
\end{proof}

\begin{lemma}\label{asm-theta-agree-factors}
Let $(B,G,\mathcal X,\omega,\omega')$ be a triple for a theta base $B$ of a CM
field $E$ (Definition~\ref{asm-theta-agree-triple}). Then
$\sigma_{e_1,e_2,s}(\mathfrak b,\mathfrak c;z)\neq0$ for all signs
$e_1,e_2,s$, fractional ideals $\mathfrak b,\mathfrak c$ of $E$ and
$z\in\mathbf Z$, and $L^{S_0}(\nu,k)\neq0$ for every integer $k\geq2$
(Definition~\ref{asm-theta-agree-parts}(4)).
\end{lemma}
\leanref{Development/AbelianVarietyTheory/Definitions/Divisors/CTSmallCone.lean}{GlobalTheory.ThK.FacNe}
\leanref{Development/WeilBound/ThetaLifts/WeilDomination/ExistsUnitChrWNe.lean}{PositivityTheory.ThK.Close.facNe}
\leanref{Development/WeilBound/ThetaLifts/WeilDomination/FacNeGauss.lean}{PositivityTheory.ThK.GaussNe.gauss_ne_zero}
\begin{proof}
The factors $\nu(\mathfrak b)^{e_1}$, $\nu(\mathfrak c)^{e_2}$ and
$\nu_v(s)$ are values of characters, so they are nonzero, and
$\epsilon(\nu_v,z)$ is $\mathfrak g(\nu_v)$ times nonzero factors. It remains
to show $\mathfrak g(\nu_v)\neq0$ for $v\in P\cap T$.

Let $v\in P\cap T$, and write $\varpi=\varpi_v$, $c=c_v$ and $q=N(v)$. As
$v\in T$, $\nu$ is ramified at the place of $E$ given by $\tau_v$, so $\nu_v$
is nontrivial on $\mathcal O_v^\times$ and $c\geq1$. As $v\in P$, $\psi_v$ is
trivial on $\mathcal O_v$, and for every $x\notin\mathcal O_v$ there is
$y\in\mathcal O_v$ with $\psi_v(xy)\neq1$. For $x\in\mathcal O_v$ put
$G(x)=\sum_y\nu_v(y)^{-1}\psi_v(xy\varpi^{-c})$, the sum over
$y\in(\mathcal O_v/\varpi^c)^\times$; it does not depend on the representatives,
and $G(1)=\mathfrak g(\nu_v)$.
\begin{enumerate}
\item For a unit $u$, $G(xu)=\nu_v(u)G(x)$, by the substitution $y\mapsto yu^{-1}$.
\item If $x\in\varpi\mathcal O_v$ then $G(x)=0$. By the minimality of $c$ there
  is a unit $u\in1+\varpi^{c-1}\mathcal O_v$ with $\nu_v(u)\neq1$. Then
  $x(u-1)y\varpi^{-c}\in\mathcal O_v$, so $G(xu)=G(x)$, and (1) gives
  $(\nu_v(u)-1)G(x)=0$.
\item Suppose $G(1)=0$. By (1) and (2), $G$ vanishes on $\mathcal O_v$. Then,
  summing over $x\in\mathcal O_v/\varpi^c$,
  \[
  0=\sum_xG(x)\psi_v(-x\varpi^{-c})=\sum_y\nu_v(y)^{-1}\sum_x\psi_v\bigl(x(y-1)\varpi^{-c}\bigr).
  \]
  If $y\not\equiv1$ modulo $\varpi^c$, then $(y-1)\varpi^{-c}\notin\mathcal O_v$,
  so there is $t\in\mathcal O_v$ with $\psi_v(t(y-1)\varpi^{-c})\neq1$, and the
  substitution $x\mapsto x+t$ shows that the inner sum is $0$. For $y=1$ the
  inner sum is $q^c$. So $0=q^c$, which is absurd.
\end{enumerate}
Hence $\mathfrak g(\nu_v)\neq0$.

For $k\geq2$, the factors of $L^{S_0}(\nu,k)$ are nonzero and the product
converges absolutely, since $|\nu_w(\varpi_w)|=1$ and $\sum_wN(w)^{-k}$
converges. So $L^{S_0}(\nu,k)\neq0$.
\end{proof}

\subsection{Automorphy and invariance of theta integrals}\label{subsection-asm-invariance}

\begin{definition}\label{asm-theta-agree-invariance-heisenberg}
Let $F$ be a number field, with adeles $\mathbf A=\mathbf A_F$, finite adeles
$\mathbf A_f$, $F_\infty=F\otimes\mathbf R$ and
$\widehat{\mathcal O}=\prod_v\mathcal O_v\subset\mathbf A_f$, and let $\psi$ be a
nontrivial continuous character of $\mathbf A$ that is trivial on $F$. Let $I$
be a finite set, and let $\mathcal S=\mathcal S(\mathbf A^I)$ be the space of
Schwartz-Bruhat functions on $\mathbf A^I$
(Definition~\ref{asm-theta-agree-models}).
\begin{enumerate}
\item The \emph{Schrödinger operator} of $u=(u_1,u_2)\in\mathbf A^I\times\mathbf A^I$
  is $(\rho(u)\phi)(x)=\psi\bigl(x\cdot u_2+\tfrac12u_1\cdot u_2\bigr)\phi(x+u_1)$.
  The element $u$ is \emph{rational} if $u\in F^I\times F^I$, and
  $\langle u,u'\rangle=u_1\cdot u'_2-u'_1\cdot u_2$.
\item A linear form $\lambda$ on the functions on $\mathbf A^I$ is
  \emph{admissible} if, for every Schwartz-Bruhat function $g$ on
  $\mathbf A_f^I$, the form $f\mapsto\lambda(f\otimes g)$ is continuous on the
  Schwartz space $\mathcal S(F_\infty^I)$.
\item The \emph{theta sum} of a function $\Phi$ on $\mathbf A^I$ is
  $\theta(\Phi)=\sum_{\xi\in F^I}\Phi(\xi)$, when the sum converges absolutely.
\item For functions $\phi_1$ on $\mathbf A^I$ and $\phi_2$ on $\mathbf A^{I'}$,
  $\phi_1\boxtimes\phi_2$ is the function $(x,y)\mapsto\phi_1(x)\phi_2(y)$ on
  $\mathbf A^{I\sqcup I'}$.
\item A subspace $L\subset F^I\times F^I$ is \emph{Lagrangian} if it consists of
  the $w$ with $\langle w,w'\rangle=0$ for all $w'\in L$.
\end{enumerate}
\end{definition}
\leanref{Development/AbelianVarietyTheory/Definitions/Divisors/HilbertNinetyQuotientHomEqIff.lean}{GlobalTheory.WeilM.schr}
\leanref{Development/AbelianVarietyTheory/Definitions/Divisors/HilbertNinetyQuotientHomEqIff.lean}{GlobalTheory.WeilM.adeleClass}
\leanref{Development/AbelianVarietyTheory/Definitions/Divisors/HilbertNinetyQuotientHomEqIff.lean}{GlobalTheory.WeilM.thetaSum}
\leanref{Development/AbelianVarietyTheory/Definitions/Divisors/HilbertNinetyQuotientHomEqIff.lean}{GlobalTheory.WeilM.boxFun}

\begin{theorem}[Tate \cite{Tat67}]\label{asm-theta-agree-invariance-tate}
Let $F$ be a number field, with $\mathbf A$, $\mathbf A_f$, $F_\infty$ and
$\widehat{\mathcal O}$ as in Definition~\ref{asm-theta-agree-invariance-heisenberg},
let $\mathfrak d^{-1}$ be the inverse different of $F$, and let
$\operatorname{Tr}\colon F_\infty\to\mathbf R$ be the trace.
\begin{enumerate}
\item There is a compact set $C\subset F_\infty$ such that every $x\in\mathbf A$
  is $k+y$ with $k\in F$ and $y\in C\times\widehat{\mathcal O}$. The subgroup $F$
  of $\mathbf A$ is closed.
\item There is a continuous character $\psi_0$ of $\mathbf A$, trivial on $F$,
  with $\psi_0(y)=e^{-2\pi i\operatorname{Tr}(y_\infty)}$ for every $y$ whose
  finite components lie in the completions of $\mathfrak d^{-1}$. Every
  continuous character of $\mathbf A$ with values of absolute value $1$ that is
  trivial on $F$ is $x\mapsto\psi_0(cx)$ for a unique $c\in F$.
\item Every continuous character of $F_\infty$ with values of absolute value $1$
  that is trivial on $\mathcal O_F$ is $m\mapsto e^{-2\pi i\operatorname{Tr}(\xi m)}$
  for some $\xi\in\mathfrak d^{-1}$.
\item Let $\psi$ be a nontrivial continuous character of $\mathbf A$ that is
  trivial on $F$. If $x\in\mathbf A$ and $\psi(xa)=1$ for every $a\in F$, then
  $x\in F$.
\item For $a\in F^\times$, multiplication by $a$ preserves the Haar measures of
  $\mathbf A$.
\end{enumerate}
\end{theorem}
\leanref{GroundTheory/GlobalFields/Adele/Discrete.lean}{GroundTheory.Global.exists_isCompact_forall_exists_sub_mem}
\leanref{GroundTheory/GlobalFields/Adele/Discrete.lean}{GroundTheory.Global.isClosed_principalSubgroup}
\leanref{GroundTheory/GlobalFields/TateThesis/AdelicChar.lean}{GroundTheory.Global.adelicChar_of_mem_invDifferentAdeles}
\leanref{GroundTheory/GlobalFields/TateThesis/AdelicCharDual.lean}{GroundTheory.Global.existsUnique_eq_adelicChar_mul}
\leanref{GroundTheory/GlobalFields/TateThesis/AdelicCharDual.lean}{GroundTheory.Global.exists_mixedSpace_char_eq}
\leanref{Development/AutomorphicSpectrum/Proofs/MembershipProofs/ThetaDistributionHeisenbergInvariant.lean}{GlobalTheory.CompanionTheory.AdeleAnn.unitary_and_annihilator_eq}
\leanref{Development/AutomorphicSpectrum/Proofs/MembershipProofs/ThetaHeisenbergTranslatesSeparateStatement.lean}{GlobalTheory.T.rational_units_haar_character_one}
\leanref{Development/AutomorphicSpectrum/Proofs/MembershipProofs/ThetaHeisenbergTranslatesSeparate.lean}{GlobalTheory.CompanionTheory.AdeleModule.distribHaarChar_principal}

\begin{lemma}\label{asm-theta-agree-invariance-theta}
Let $F$, $\psi$, $I$ and $\mathcal S$ be as in
Definition~\ref{asm-theta-agree-invariance-heisenberg}.
\begin{enumerate}
\item $|\psi(x)|=1$ for every $x\in\mathbf A$.
\item For $\Phi\in\mathcal S$ the theta sum converges absolutely, and $\theta$
  agrees on $\mathcal S$ with an admissible linear form. There is
  $\Phi_0\in\mathcal S$ with $\theta(\Phi_0)\neq0$.
\item For $\Phi\in\mathcal S$ and $u\in\mathbf A^I\times\mathbf A^I$,
  $\rho(u)\Phi\in\mathcal S$. If $u$ is rational, then
  $\theta(\rho(u)\Phi)=\theta(\Phi)$.
\item For $\phi_1\in\mathcal S(\mathbf A^I)$ and $\phi_2\in\mathcal S(\mathbf A^{I'})$,
  $\phi_1\boxtimes\phi_2\in\mathcal S(\mathbf A^{I\sqcup I'})$ and
  $\theta(\phi_1\boxtimes\phi_2)=\theta(\phi_1)\,\theta(\phi_2)$.
\end{enumerate}
\end{lemma}
\leanref{Development/AutomorphicSpectrum/Proofs/MembershipProofs/ThetaDistributionNonzeroAdmissible.lean}{GlobalTheory.CompanionTheory.AdelicHeis.norm_eq_one}
\leanref{Development/AutomorphicSpectrum/Proofs/MembershipProofs/ThetaSumAdmissibilityStatement.lean}{GlobalTheory.T.theta_distribution_nonzero_admissible}
\leanref{Development/AutomorphicSpectrum/Proofs/MembershipProofs/ThetaSumAdmissibility.lean}{GlobalTheory.CompanionTheory.ThetaAdm.theta_adm}
\leanref{Development/AutomorphicSpectrum/Proofs/MembershipProofs/ThetaDistributionNonzeroAdmissibleStatement.lean}{GlobalTheory.T.theta_distribution_heisenberg_invariant}
\leanref{Development/AutomorphicSpectrum/Proofs/MembershipProofs/ThetaDistributionNonzeroAdmissible.lean}{GlobalTheory.CompanionTheory.AdelicHeis.schr_mem_and_thetaSum}
\leanref{Development/AutomorphicSpectrum/Proofs/MembershipProofs/RemovableLogGrowthSingularityStatement.lean}{GlobalTheory.T.theta_sum_tensor_product}
\leanref{Development/AutomorphicSpectrum/Proofs/MembershipProofs/RemovableLogGrowthSingularity.lean}{GlobalTheory.CompanionTheory.ThetaBox.theta_box}
\begin{proof}
(1) By Theorem~\ref{asm-theta-agree-invariance-tate}(1) there is a compact set
$D\subset\mathbf A$ with $\mathbf A=F+D$. As $\psi$ is trivial on $F$, $|\psi|$
is bounded on $\mathbf A$ by its maximum $M$ on $D$. If $|\psi(x)|>1$, then
$|\psi(nx)|=|\psi(x)|^n$ exceeds $M$ for large $n$. So $|\psi|\leq1$, and
$|\psi(x)|\,|\psi(-x)|=1$ gives $|\psi(x)|=1$.

(2) It suffices to treat $\Phi=f\otimes g$. The support of $g$ is compact, so
there is a nonzero $d\in\mathcal O_F$ with $d\,\operatorname{supp}(g)\subset\widehat{\mathcal O}^I$.
Hence $g(\xi)\neq0$ for $\xi\in F^I$ only if $\xi$ lies in the lattice
$\Lambda=(d^{-1}\mathcal O_F)^I$ of $F_\infty^I$. Let $C$ bound $|g|$, and let
$k$ exceed the rank of $\Lambda$, so that
$\sum_{\xi\in\Lambda}(1+|\xi|)^{-k}$ converges. There is a Schwartz seminorm
$p_k$ with $|f(x)|\leq2^kp_k(f)(1+|x|)^{-k}$. So
\[
\sum_{\xi\in F^I}|f(\xi)g(\xi)|\leq2^kC\,p_k(f)\sum_{\xi\in\Lambda}(1+|\xi|)^{-k}.
\]
This gives the absolute convergence and the continuity of $f\mapsto\theta(f\otimes g)$.
Any linear extension of $\theta$ from $\mathcal S$ to all functions is then
admissible. For $\Phi_0$ take $f$ a smooth function of compact support on
$F_\infty^I$ with values in $[0,1]$ and $f(0)=1$, and $g$ the indicator function
of $\widehat{\mathcal O}^I$. Every term of $\theta(\Phi_0)$ is $\geq0$, and the
term at $\xi=0$ is $1$.

(3) Let $\Phi=f\otimes g$. Then $\rho(u)\Phi$ is the constant
$\psi(\frac12u_1\cdot u_2)$ times the product of
$x_\infty\mapsto\psi(x_\infty\cdot u_{2,\infty})f(x_\infty+u_{1,\infty})$ and
$x_f\mapsto\psi(x_f\cdot u_{2,f})g(x_f+u_{1,f})$. The function
$x_\infty\mapsto\psi(x_\infty\cdot u_{2,\infty})$ is a continuous character of
$F_\infty^I$ of absolute value $1$, so it is $e^{i\ell(x_\infty)}$ for a linear
form $\ell$; it is smooth with all derivatives bounded, so the first factor is a
Schwartz function. The restriction of $\psi$ to $\mathbf A_f$ is trivial on an
open subgroup, so the second factor is locally constant with compact support.
If $u$ is rational, then $\psi(\xi\cdot u_2+\frac12u_1\cdot u_2)=1$ for
$\xi\in F^I$, as $\psi$ is trivial on $F$, and $\xi\mapsto\xi+u_1$ permutes
$F^I$. Hence $\theta(\rho(u)\Phi)=\sum_\xi\Phi(\xi+u_1)=\theta(\Phi)$.

(4) For $\phi_k=f_k\otimes g_k$, $\phi_1\boxtimes\phi_2=(f_1\boxtimes f_2)\otimes(g_1\boxtimes g_2)$,
and these span. The theta sum over $F^{I\sqcup I'}=F^I\times F^{I'}$ is the
product of two absolutely convergent sums.
\end{proof}

\begin{theorem}[Poisson summation \cite{Hor90}]\label{asm-theta-agree-invariance-poisson}
Let $V$ be a finite-dimensional real vector space with a basis
$(b_p)_{p\in P}$, and write $x=\sum_px_pb_p$ for $x\in V$. Let $T$ be a
continuous linear form on the Schwartz space $\mathcal S(V)$ such that
$T(g)=T(f)$ whenever, for some $p\in P$, $g(x)=e^{2\pi ix_p}f(x)$ for all $x$,
and whenever, for some $m\in\mathbf Z^P$, $g(x)=f(x+\sum_pm_pb_p)$ for all $x$.
Then there is $a\in\mathbf C$ with
$T(f)=a\sum_{m\in\mathbf Z^P}f\bigl(\sum_pm_pb_p\bigr)$ for every
$f\in\mathcal S(V)$.
\end{theorem}
\leanref{Development/AutomorphicSpectrum/Proofs/MembershipProofs/AdelicCharacterRationalDual.lean}{GlobalTheory.CompanionTheory.ThetaUniq.comb_of_basis}
\leanref{Development/AutomorphicSpectrum/Proofs/MembershipProofs/DenseHasCompactSupport.lean}{GlobalTheory.CompanionTheory.DiracComb.eq_mul_comb}

\begin{lemma}\label{asm-theta-agree-invariance-modulation}
Let $F$, $\psi$ and $\widehat{\mathcal O}$ be as in
Definition~\ref{asm-theta-agree-invariance-heisenberg}, let $N\geq1$ be an
integer, and let $\varphi\colon F_\infty\to\mathbf R$ be an $\mathbf R$-linear
map with $\varphi(Na)\in\mathbf Z$ for every $a\in\mathcal O_F$. Then there is
$y\in F$ with $\psi(xy)=e^{2\pi i\varphi(x_\infty)}$ for every $x\in\mathbf A$
whose finite part lies in $N\widehat{\mathcal O}$.
\end{lemma}
\leanref{Development/AutomorphicSpectrum/Proofs/MembershipProofs/AdelicCharacterRationalDual.lean}{GlobalTheory.CompanionTheory.ThetaUniq.exists_rat_modulation}
\begin{proof}
By Lemma~\ref{asm-theta-agree-invariance-theta}(1) and
Theorem~\ref{asm-theta-agree-invariance-tate}(2), $\psi(x)=\psi_0(cx)$ for some
$c\in F$, and $c\neq0$ as $\psi$ is nontrivial. The character
$m\mapsto e^{2\pi i\varphi(Nm)}$ of $F_\infty$ is trivial on $\mathcal O_F$. By
Theorem~\ref{asm-theta-agree-invariance-tate}(3) it is
$m\mapsto e^{-2\pi i\operatorname{Tr}(\xi m)}$ for some $\xi\in\mathfrak d^{-1}$.
Put $y=\xi/(Nc)$. Let $x\in\mathbf A$ have finite part $Nu$ with
$u\in\widehat{\mathcal O}$, and put $x'=x/N$. Then $cxy=\xi x'$, and the finite
components $\xi u_v$ of $\xi x'$ lie in the completions of $\mathfrak d^{-1}$. By
Theorem~\ref{asm-theta-agree-invariance-tate}(2),
\[
\psi(xy)=\psi_0(\xi x')=e^{-2\pi i\operatorname{Tr}(\xi x'_\infty)}
=e^{2\pi i\varphi(Nx'_\infty)}=e^{2\pi i\varphi(x_\infty)}.\qedhere
\]
\end{proof}

\begin{lemma}\label{asm-theta-agree-invariance-unique}
Let $F$, $\psi$, $I$ and $\mathcal S$ be as in
Definition~\ref{asm-theta-agree-invariance-heisenberg}, and let $\lambda$ be an
admissible linear form with $\lambda(\rho(w)\Phi)=\lambda(\Phi)$ for every
rational $w$ and every $\Phi\in\mathcal S$. Then there is $a\in\mathbf C$ with
$\lambda(\Phi)=a\,\theta(\Phi)$ for every $\Phi\in\mathcal S$.
\end{lemma}
\leanref{Development/AutomorphicSpectrum/Proofs/MembershipProofs/AdelicCharacterRationalDualStatement.lean}{GlobalTheory.T.admissible_functional_rational_invariance}
\leanref{Development/AutomorphicSpectrum/Proofs/MembershipProofs/AdelicCharacterRationalDual.lean}{GlobalTheory.CompanionTheory.ThetaUniq.theta_unique}
\leanref{Development/AutomorphicSpectrum/Proofs/MembershipProofs/AdelicCharacterRationalDual.lean}{GlobalTheory.CompanionTheory.ThetaUniq.level_zero}
\leanref{Development/AutomorphicSpectrum/Proofs/MembershipProofs/AdelicCharacterRationalDual.lean}{GlobalTheory.CompanionTheory.ThetaUniq.level_span}
\leanref{Development/AutomorphicSpectrum/Proofs/MembershipProofs/AdelicCharacterRationalDual.lean}{GlobalTheory.CompanionTheory.ThetaUniq.mem_span_cosetInd}
\begin{proof}
For $d\in F^I$, $y\in F$ and $i\in I$, the hypothesis for $w=(d,0)$ and for
$w=(0,ye_i)$, with $e_i$ the $i$th unit vector, gives
\[
\lambda\bigl(\Phi(\cdot+d)\bigr)=\lambda(\Phi),\qquad
\lambda\bigl(\psi(x_iy)\Phi(x)\bigr)=\lambda(\Phi)\qquad(\Phi\in\mathcal S).
\]
For an integer $N\geq1$ and $d\in F^I$, let $1_{N,d}$ be the indicator function
of $d+(N\widehat{\mathcal O})^I$ in $\mathbf A_f^I$. It is a Schwartz-Bruhat
function.

\emph{Step 1.} Fix $N$. Then there is $c_N\in\mathbf C$ with
$\lambda(f\otimes1_{N,0})=c_N\,\theta(f\otimes1_{N,0})$ for every
$f\in\mathcal S(F_\infty^I)$. Let $\omega_1,\ldots,\omega_r$ be a
$\mathbf Z$-basis of $\mathcal O_F$, so that $N\omega_1,\ldots,N\omega_r$ is an
$\mathbf R$-basis of $F_\infty$, and let $B$ be the basis of $F_\infty^I$
whose element $b_{ij}$ is $N\omega_j$ in the coordinate $i$ and $0$ elsewhere,
with coordinates $x_{ij}$. The form $T(f)=\lambda(f\otimes1_{N,0})$ is
continuous on $\mathcal S(F_\infty^I)$, as $\lambda$ is admissible. For
$m\in\mathbf Z^{I\times\{1,\ldots,r\}}$, $\sum m_{ij}b_{ij}$ is the image of
$d=(N\sum_jm_{ij}\omega_j)_i\in(N\mathcal O_F)^I$, and $1_{N,0}(x+d)=1_{N,0}(x)$;
so the first invariance gives $T(f(\cdot+\sum m_{ij}b_{ij}))=T(f)$. The
coordinate $\varphi_j$ of $F_\infty$ for the basis $(N\omega_j)$ takes
$N\mathcal O_F$ into $\mathbf Z$. By
Lemma~\ref{asm-theta-agree-invariance-modulation} there is $y\in F$ with
$\psi(ty)=e^{2\pi i\varphi_j(t_\infty)}$ for every $t\in\mathbf A$ with finite
part in $N\widehat{\mathcal O}$. Hence
$e^{2\pi ix_{ij}}(f\otimes1_{N,0})(x)=\psi(x_iy)(f\otimes1_{N,0})(x)$ for every
$x\in\mathbf A^I$, both sides being $0$ unless $x_f\in(N\widehat{\mathcal O})^I$.
The second invariance gives $T(e^{2\pi ix_{ij}}f)=T(f)$. By
Theorem~\ref{asm-theta-agree-invariance-poisson},
$T(f)=c_N\sum_{a\in\mathcal O_F^I}f(Na)$ for some $c_N$. An element of $F^I$ lies
in $(N\widehat{\mathcal O})^I$ exactly when it lies in $(N\mathcal O_F)^I$, so the
sum is $\theta(f\otimes1_{N,0})$.

\emph{Step 2.} For $d\in F^I$, $(f\otimes1_{N,d})(x+d)=f(x_\infty+d)1_{N,0}(x_f)$.
The first invariance, Step 1 and
Lemma~\ref{asm-theta-agree-invariance-theta}(3) give
\[
\lambda(f\otimes1_{N,d})=\lambda\bigl(f(\cdot+d)\otimes1_{N,0}\bigr)
=c_N\,\theta\bigl(f(\cdot+d)\otimes1_{N,0}\bigr)=c_N\,\theta(f\otimes1_{N,d}).
\]
By linearity, $\lambda(f\otimes g)=c_N\theta(f\otimes g)$ for every $g$ in the
span of the $1_{N,d}$, $d\in F^I$.

\emph{Step 3.} Let $g$ be a Schwartz-Bruhat function on $\mathbf A_f^I$. As $g$
is locally constant with compact support, it is invariant under translation by
$(N_0\widehat{\mathcal O})^I$ for some $N_0\geq1$. If $N_0\mid N$, then $g$ is a
finite linear combination of indicator functions of cosets of
$(N\widehat{\mathcal O})^I$. Every such coset contains a point of $F^I$, since
$\mathbf A_f=F+\widehat{\mathcal O}$ by
Theorem~\ref{asm-theta-agree-invariance-tate}(1); applied to $y/N$ for
$y\in\mathbf A_f$, this gives $\mathbf A_f=F+N\widehat{\mathcal O}$. So $g$ lies in the span of the
$1_{N,d}$. As $\mathcal S$ is spanned by the $f\otimes g$, Step 2 gives: for
every $\Phi\in\mathcal S$ there is $N_0\geq1$ with
$\lambda(\Phi)=c_N\theta(\Phi)$ for every $N$ divisible by $N_0$.

\emph{Step 4.} If $\theta$ vanishes on $\mathcal S$, then so does $\lambda$ by
Step 3, and $a=0$ works. Otherwise choose $\Phi_0\in\mathcal S$ with
$\theta(\Phi_0)\neq0$ and put $a=\lambda(\Phi_0)/\theta(\Phi_0)$. For
$\Phi\in\mathcal S$, choose $N$ divisible by the integers $N_0$ of Step 3 for
$\Phi$ and for $\Phi_0$. Then $\lambda(\Phi_0)=c_N\theta(\Phi_0)$ gives $c_N=a$,
and $\lambda(\Phi)=c_N\theta(\Phi)=a\,\theta(\Phi)$.
\end{proof}

\begin{lemma}\label{asm-theta-agree-invariance-average}
Let $F$, $\psi$, $I$ and $\mathcal S$ be as in
Definition~\ref{asm-theta-agree-invariance-heisenberg}, and let
$L\subset F^I\times F^I$ be a Lagrangian subspace. Then there is an admissible
linear form $l$ with the following properties.
\begin{enumerate}
\item $l(\Phi_1)\neq0$ for some $\Phi_1\in\mathcal S$.
\item $l(\rho(y)\Phi)=l(\Phi)$ for every $y$ in the $\mathbf A$-span of $L$ in
  $\mathbf A^I\times\mathbf A^I$ and every $\Phi\in\mathcal S$.
\item Let $s$ be an $F$-linear automorphism of $F^I\times F^I$ with $s(L)=L$, and
  let $s_{\mathbf A}$ be an $\mathbf A$-linear endomorphism of
  $\mathbf A^I\times\mathbf A^I$ that agrees with $s$ on $F^I\times F^I$. Let $T$
  be a linear operator on the functions on $\mathbf A^I$ that preserves
  $\mathcal S$, with $T\rho(u)\Phi=\rho(s_{\mathbf A}u)T\Phi$ for every
  $u\in\mathbf A^I\times\mathbf A^I$ and $\Phi\in\mathcal S$. If $a\in\mathbf C$
  and $\theta(T\Phi)=a\,\theta(\Phi)$ for every $\Phi\in\mathcal S$, then
  $l(T\Phi)=a\,l(\Phi)$ for every $\Phi\in\mathcal S$.
\end{enumerate}
\end{lemma}
\leanref{Development/AutomorphicSpectrum/Proofs/MembershipProofs/RationalUnitsHaarCharacterOneStatement.lean}{GlobalTheory.T.lagrangian_invariant_functional_exists}
\leanref{Development/AutomorphicSpectrum/Proofs/MembershipProofs/RationalUnitsHaarCharacterOne.lean}{GlobalTheory.CompanionTheory.ThetaAvg.theta_mean}
\leanref{Development/AutomorphicSpectrum/Proofs/MembershipProofs/RationalUnitsHaarCharacterOne.lean}{GlobalTheory.CompanionTheory.ThetaAvg.exists_avg}
\leanref{Development/AutomorphicSpectrum/Proofs/MembershipProofs/RationalUnitsHaarCharacterOne.lean}{GlobalTheory.CompanionTheory.ThetaAvg.exists_mean}
\begin{proof}
A direct computation gives
$\rho(u)\rho(u')=\psi(\frac12\langle u,u'\rangle)\rho(u+u')$, and hence
$\rho(u)\rho(u')=\psi(\langle u,u'\rangle)\rho(u')\rho(u)$.

Let $b_1,\ldots,b_n$ be a basis of $L$, and let
$\beta(t)=\sum_it_ib_i$ for $t\in\mathbf A^n$. The image of $\beta$ is the
$\mathbf A$-span of $L$. As $L$ is Lagrangian, $\langle\beta t,\beta t'\rangle=0$,
so $\rho(\beta t)\rho(\beta t')=\rho(\beta(t+t'))$. For $\Phi\in\mathcal S$ put
$F_\Phi(t)=\theta(\rho(\beta t)\Phi)$, which is defined by
Lemma~\ref{asm-theta-agree-invariance-theta}(2) and~(3). It is continuous, by
dominated convergence: for $\Phi=f\otimes g$ and $t$ in a compact set, the bound
in the proof of Lemma~\ref{asm-theta-agree-invariance-theta}(2) holds for the
translates $\Phi(\cdot+(\beta t)_1)$ with one summable majorant, and
$|\rho(u)\Phi(x)|=|\Phi(x+u_1)|$. For $c\in F^n$, $\beta c$ is rational and
$F_\Phi(t+c)=\theta(\rho(\beta c)\rho(\beta t)\Phi)=F_\Phi(t)$ by
Lemma~\ref{asm-theta-agree-invariance-theta}(3). So $F_\Phi$ is a continuous
function on $(\mathbf A/F)^n$, which is a compact group by
Theorem~\ref{asm-theta-agree-invariance-tate}(1). Let $dt$ be its Haar
probability measure, put $l(\Phi)=\int_{(\mathbf A/F)^n}F_\Phi(t)\,dt$ for
$\Phi\in\mathcal S$, and extend $l$ linearly to all functions on $\mathbf A^I$.

\emph{Admissibility.} Let $D\subset\mathbf A^n$ be compact with
$\mathbf A^n=F^n+D$ (Theorem~\ref{asm-theta-agree-invariance-tate}(1)). Then
$|l(\Phi)|\leq\sup_{t\in D}|F_\Phi(t)|$. For $\Phi=f\otimes g$, the uniform bound
above gives a Schwartz seminorm $p$ and a constant $C$ with
$|F_\Phi(t)|\leq C\,p(f)$ for $t\in D$. So $f\mapsto l(f\otimes g)$ is
continuous.

(2) Let $y=\beta t_0$. Then $F_{\rho(y)\Phi}(t)=F_\Phi(t+t_0)$, and $dt$ is
invariant under translation.

(1) Suppose that $l$ vanishes on $\mathcal S$, and let $\Phi_0\in\mathcal S$.
For $c\in F^n$ let $e_c(t)=\psi(\sum_it_ic_i)$, a continuous function on
$(\mathbf A/F)^n$. The form $\langle\ ,\ \rangle$ is nondegenerate on
$F^I\times F^I$, so there is a rational $w$ with $\langle b_i,w\rangle=c_i$ for
every $i$. Then $\langle\beta t,w\rangle=\sum_it_ic_i$, and the commutation
rule with Lemma~\ref{asm-theta-agree-invariance-theta}(3) for $w$ gives
$F_{\rho(w)\Phi_0}=e_cF_{\Phi_0}$. So
$\int e_cF_{\Phi_0}\,dt=l(\rho(w)\Phi_0)=0$ for every $c$. The functions $e_c$
contain $1$, are closed under products and, by
Lemma~\ref{asm-theta-agree-invariance-theta}(1), under complex conjugation.
They separate the points of $(\mathbf A/F)^n$: if $e_c(t)=e_c(t')$ for all $c$,
then $\psi((t'_i-t_i)a)=1$ for all $a\in F$, so $t'_i-t_i\in F$ by
Theorem~\ref{asm-theta-agree-invariance-tate}(4). By the Stone-Weierstrass
theorem $F_{\Phi_0}=0$, and in particular $\theta(\Phi_0)=F_{\Phi_0}(0)=0$. This
contradicts Lemma~\ref{asm-theta-agree-invariance-theta}(2) for a suitable
$\Phi_0$.

(3) As $s(L)=L$, the restriction of $s$ to $L$ has a matrix $S\in GL_n(F)$ in the
basis $(b_i)$, and $s_{\mathbf A}(\beta t)=\beta(St)$ since $s_{\mathbf A}$ is
$\mathbf A$-linear. For $\Phi\in\mathcal S$,
\[
F_{T\Phi}(St)=\theta\bigl(\rho(s_{\mathbf A}\beta t)T\Phi\bigr)
=\theta\bigl(T\rho(\beta t)\Phi\bigr)=a\,F_\Phi(t),
\]
using $\rho(\beta t)\Phi\in\mathcal S$. The map $t\mapsto St$ induces a
continuous surjective endomorphism of the compact group $(\mathbf A/F)^n$, which
preserves its Haar probability measure. Hence
$l(T\Phi)=\int F_{T\Phi}(St)\,dt=a\,l(\Phi)$.
\end{proof}

\begin{lemma}\label{asm-theta-agree-invariance-character}
Let $F$ be a field of characteristic $0$, and let $E$ be a commutative
$F$-algebra with an $F$-algebra automorphism $c$ and an invertible $\xi\in E$
such that $c(\xi)=-\xi$ and $1,\xi$ is an $F$-basis of $E$. Let $\beta$ and $\nu$
be finite sets, and let $J\in M_\beta(E)$ with $c(J)^{\mathsf T}=J$ and $\det J$
invertible; let $U_J(E)$ be the group of $g\in GL_\beta(E)$ with
$g^{\mathsf T}Jc(g)=J$. Let $C_y,C_x\in M_{\beta\times\nu}(E)$ with
$C_y^{\mathsf T}Jc(C_y)=0$ and $C_x^{\mathsf T}Jc(C_x)=0$, such that the matrix
$(C_y\ C_x)$ has a two-sided inverse. Let $\lambda\colon U_J(E)\to\mathbf C^\times$
be a homomorphism with $\lambda(p)=1$ whenever $pC_y=C_yt$ for some
$t\in GL_\nu(E)$. Then $\lambda=1$.
\end{lemma}
\leanref{Development/AutomorphicSpectrum/Proofs/MembershipProofs/SplitQuadraticAlgebraOrderStatement.lean}{GlobalTheory.T.lagrangian_stabilizer_character_trivial}
\leanref{Development/AutomorphicSpectrum/Proofs/MembershipProofs/SplitQuadraticAlgebraOrder.lean}{GlobalTheory.CompanionTheory.SiegelChar.char_trivial}
\leanref{Development/AutomorphicSpectrum/Proofs/MembershipProofs/SplitQuadraticAlgebraOrder.lean}{GlobalTheory.CompanionTheory.SiegelChar.lower_triv}
\leanref{Development/AutomorphicSpectrum/Proofs/MembershipProofs/SplitQuadraticAlgebraOrder.lean}{GlobalTheory.CompanionTheory.SiegelChar.cell_triv}
\leanref{Development/AutomorphicSpectrum/Proofs/MembershipProofs/SplitQuadraticAlgebraOrder.lean}{GlobalTheory.CompanionTheory.SiegelChar.triv_of_cell}
\begin{proof}
For $x\in M_\beta(E)$ put $\sigma(x)=(Jc(x)J^{-1})^{\mathsf T}$. Then $\sigma$ is
additive, $\sigma(ax)=c(a)\sigma(x)$ for $a\in E$, $\sigma(xy)=\sigma(y)\sigma(x)$,
$\sigma(1)=1$, and
$\sigma(x)x=(J^{-1})^{\mathsf T}(x^{\mathsf T}Jc(x))^{\mathsf T}$. So
$x\in U_J(E)$ if and only if $\sigma(x)x=1$. Let $\left(\begin{smallmatrix}M_t\\M_b\end{smallmatrix}\right)$ be the
inverse of $(C_y\ C_x)$, and put $e=C_yM_t$ and $f=C_xM_b$. Then $e+f=1$,
$e^2=e$, $f^2=f$ and $ef=fe=0$, and the isotropy of $C_y$ and $C_x$ gives
$Jc(e)J^{-1}=M_b^{\mathsf T}C_x^{\mathsf T}$, so $\sigma(e)=f$ and
$\sigma(f)=e$. All matrices below are in $M_\beta(E)$.

(a) \emph{If $g\in U_J(E)$ and $fge=0$, then $\lambda(g)=1$.} Indeed $ge=(e+f)ge=ege$,
and $eC_y=C_y$ as $M_tC_y=1$. So $gC_y=geC_y=egeC_y=C_y\,(M_tgC_y)$, and
$M_tgC_y$ has the inverse $M_tg^{-1}C_y$, because $egC_y=gC_y$.

(b) \emph{If $es=0$, $sf=0$ and $\sigma(s)=-s$, then $1+s\in U_J(E)$ and
$\lambda(1+s)=1$.} Here $s=fs=se$, so $s^2=s(ef)s=0$, and
$\sigma(1+s)(1+s)=(1-s)(1+s)=1$. The element $d=2e+\frac12f$ has
$\sigma(d)d=(2f+\frac12e)(2e+\frac12f)=1$, so $d\in U_J(E)$, and
$d(1+s)d^{-1}=1+\frac14s$, so $\lambda(1+\frac14s)=\lambda(1+s)$. As $s^2=0$,
$(1+\frac14s)^4=1+s$. So
$\lambda(1+s)=\lambda(1+\frac14s)^4=\lambda(1+s)^4$, and $\lambda(1+s)^3=1$.
The same holds for $\frac13s$, and $(1+\frac13s)^3=1+s$; hence
$\lambda(1+s)=\lambda(1+\frac13s)^3=1$.

(c) \emph{If $g\in U_J(E)$ and $ge+f$ is invertible, then $\lambda(g)=1$.} Let
$w=(ge+f)^{-1}$ and $s=-fwe$. Then $es=0$ and $sf=0$. From $w(ge+f)=1$ follow
$wf=f$ and $wge=e$, hence $w(ege)=e-fge$ and $s(ege)=fge$; also
$(ege)(ewe)=e$, so $s=(fge)(ewe)$. From $\sigma(g)g=1$,
$\sigma(fge)(ege)+\sigma(ege)(fge)=f\sigma(g)ge=0$. Therefore
$\sigma(fge)=-\sigma(ege)s$, and
$\sigma(s)=\sigma(ewe)\sigma(fge)=-\sigma\bigl((ege)(ewe)\bigr)s=-\sigma(e)s=-s$.
By (b), $n=1+s$ lies in $U_J(E)$ with $\lambda(n)=1$. Since $n^{-1}=1-s$ and
$s(ge)=fge$, $fn^{-1}ge=0$; so $\lambda(n^{-1}g)=1$ by (a), and $\lambda(g)=1$.

(d) \emph{If $g\in U_J(E)$ and $g+1$ is invertible, then $\lambda(g)=1$.} Let
$A,B\in M_\beta(E)$ with $A+\tau B$ invertible for some $\tau\in E$. Then
$A+tB$ is invertible for all but finitely many $t\in F$: $E$ is finite over $F$
and has finitely many maximal ideals, and modulo each of them $\det(A+XB)$ is a
nonzero polynomial. As $F$ is infinite, there is $t\in F$ such that
\[
N=(g+1)+t(g-1),\quad D=(g+1)-t(g-1),\quad (g+1)+t(g-1)(e-f),\quad
(g+1)(ge+f)+t(g-1)(f-ge)
\]
are invertible: the first three are invertible at $t=0$, and the last is $2g$ at
$t=1$. The matrices $N$ and $D$ commute, and $\sigma(g)g=1$ gives
$\sigma(N)=\sigma(g)D$ and $\sigma(D)=\sigma(g)N$. So $g_a=ND^{-1}$ has
$\sigma(g_a)g_a=\sigma(D)^{-1}\sigma(g)DND^{-1}=\sigma(D)^{-1}\sigma(g)N=1$, and
$g_a\in U_J(E)$. Now $g_ae+f=D^{-1}(Ne+Df)$ with
$Ne+Df=(g+1)+t(g-1)(e-f)$, which is invertible; so $\lambda(g_a)=1$ by (c).
Also $g_a^{-1}ge+f=N^{-1}(Dge+Nf)$ with
$Dge+Nf=(g+1)(ge+f)+t(g-1)(f-ge)$, which is invertible; so
$\lambda(g_a^{-1}g)=1$ by (c). Hence $\lambda(g)=1$.

(e) Let $g\in U_J(E)$. As in (d), there is $t\in F$ such that
$(g+1)+t\xi(g-1)$, $1-\xi t$ and $1+\xi t$ are invertible: the first is $2g$ at
$t=\xi^{-1}$, the others are $1$ at $t=0$. Put $w=1-\xi t$, so that
$c(w)=1+\xi t$, and $z=c(w)w^{-1}$. Then $c(z)z=1$, so the scalar matrix $z$ lies
in $U_J(E)$. The matrix $ze+f$ has the inverse $c(z)e+f$, so $\lambda(z)=1$ by
(c). And $w(zg+1)=c(w)g+w=(g+1)+t\xi(g-1)$ is invertible, so $gz+1$ is
invertible and $\lambda(gz)=1$ by (d). Hence $\lambda(g)=1$.
\end{proof}

\begin{definition}\label{asm-theta-agree-invariance-datum}
Let $F$, $\mathbf A$, $\psi$ and the Schwartz-Bruhat spaces $\mathcal S$ be as
in Definition~\ref{asm-theta-agree-invariance-heisenberg}, let $E\supset F$ be a
field with an $F$-automorphism $c$, put $\mathbf A_E=\mathbf A\otimes_FE$ with
$c$ acting on the second factor, and for $a\in\mathbf A^\times$ let $|a|$ be the
factor by which multiplication by $a$ scales the Haar measures of $\mathbf A$.
For a square matrix $J$ over $E$ of size $\gamma$, let $U_J$ be the group of
$g\in GL_\gamma(\mathbf A_E)$ with $g^{\mathsf T}Jc(g)=J$, and $U_J(E)$ its
subgroup of matrices over $E$. Let $\alpha$, $\beta$, $\kappa$ and $\nu$ be
finite sets, $\xi\in E$, $J_A\in M_\alpha(E)$ and $J_B\in M_\beta(E)$.
\begin{enumerate}
\item The space $M_{\alpha\times\beta}(\mathbf A_E)$ carries the form
  $\langle x,y\rangle=\operatorname{Tr}_{\mathbf A_E/\mathbf A}\bigl(\xi\,h(x,y)\bigr)$,
  where $h(x,y)=\sum_{i,j,k,l}x_{ij}(J_A)_{ik}(J_B)_{jl}c(y_{kl})$. A
  \emph{Darboux isomorphism} for $(\xi,J_A,J_B)$ is an $\mathbf A$-linear
  isomorphism $d=(d_1,d_2)$ from $M_{\alpha\times\beta}(\mathbf A_E)$ to
  $\mathbf A^\kappa\times\mathbf A^\kappa$ with
  $\langle x,y\rangle=d_1(x)\cdot d_2(y)-d_1(y)\cdot d_2(x)$. It is
  \emph{rational} if it maps $M_{\alpha\times\beta}(E)$ bijectively onto
  $F^\kappa\times F^\kappa$. Write $\rho_d(x)=\rho(d(x))$.
\item A \emph{splitting} for $(d,J_B)$ is a homomorphism $\omega$ from $U_{J_B}$
  to the linear operators on the functions on $\mathbf A^\kappa$ such that every
  $\omega(h)$ preserves $\mathcal S$, $\lambda\circ\omega(h)$ is admissible for
  every admissible $\lambda$, and
  $\omega(h)\rho_d(x)\Phi=\rho_d(xh^{\mathsf T})\omega(h)\Phi$ for
  $x\in M_{\alpha\times\beta}(\mathbf A_E)$ and $\Phi\in\mathcal S$.
\item Let $C_y\in M_{\beta\times\nu}(E)$, and let $\chi$ be a character of
  $\mathbf A_E^\times$ with values of absolute value $1$. A splitting $\omega$
  satisfies the \emph{Siegel rule} for $(C_y,\chi)$ if every admissible
  $\lambda$ with $\lambda(\rho_d(zC_y^{\mathsf T})\Phi)=\lambda(\Phi)$ for all
  $z\in M_{\alpha\times\nu}(\mathbf A_E)$ and $\Phi\in\mathcal S$ satisfies
  \[
  \lambda(\omega(p)\Phi)=\chi(\det t)\,\bigl|N_{\mathbf A_E/\mathbf A}(\det t)\bigr|^{|\alpha|/2}\lambda(\Phi)
  \]
  whenever $\Phi\in\mathcal S$, $p\in U_{J_B}$, $t\in GL_\nu(\mathbf A_E)$ and
  $pC_y=C_yt$.
\item Let $\chi$ be as in (3). A \emph{Weil datum} for $(\xi,J_A,J_B,d,\chi)$
  is a splitting $\omega$ for $(d,J_B)$ for which there are a Darboux
  isomorphism $d^\square$ for $(\xi,J_A,J_B^\square)$, where
  $J_B^\square=\operatorname{diag}(J_B,-J_B)$, with
  $d^\square(x_1\ x_2)=\bigl((d_1(x_1),d_1(x_2)),(d_2(x_1),-d_2(x_2))\bigr)$,
  and a splitting $\omega^\square$ for $(d^\square,J_B^\square)$ that satisfies
  the Siegel rule for
  $(\Delta,\chi)$, $\Delta=\left(\begin{smallmatrix}1\\1\end{smallmatrix}\right)$,
  and has $\omega^\square(\operatorname{diag}(h,1))(\phi_1\boxtimes\phi_2)=(\omega(h)\phi_1)\boxtimes\phi_2$
  for $h\in U_{J_B}$ and $\phi_1,\phi_2\in\mathcal S$.
\end{enumerate}
\end{definition}
\leanref{Development/AbelianVarietyTheory/Definitions/Divisors/HilbertNinetyQuotientHomEqIff.lean}{GlobalTheory.WeilM.symp}
\leanref{Development/AbelianVarietyTheory/Definitions/Divisors/HilbertNinetyQuotientHomEqIff.lean}{GlobalTheory.WeilM.Darboux}
\leanref{Development/AbelianVarietyTheory/Definitions/Divisors/HilbertNinetyQuotientHomEqIff.lean}{GlobalTheory.WeilM.IsSplitting}
\leanref{Development/AbelianVarietyTheory/Definitions/Divisors/HilbertNinetyQuotientHomEqIff.lean}{GlobalTheory.WeilM.SiegelRule}
\leanref{Development/AbelianVarietyTheory/Definitions/Divisors/HilbertNinetyQuotientHomEqIff.lean}{GlobalTheory.WeilM.IsDblCoord}
\leanref{Development/AbelianVarietyTheory/Definitions/Divisors/HilbertNinetyQuotientHomEqIff.lean}{GlobalTheory.WeilM.IsWeilB}

\begin{lemma}\label{asm-theta-agree-invariance-split}
Let $F$, $\psi$, $E$, $c$, $\xi$, $J_A$, $J_B$, $\kappa$ and $\nu$ be as in
Definition~\ref{asm-theta-agree-invariance-datum}, with $\xi$ invertible,
$c(\xi)=-\xi$ and $1,\xi$ an $F$-basis of $E$, and with
$c(J_A)^{\mathsf T}=J_A$, $c(J_B)^{\mathsf T}=J_B$ and $\det J_A,\det J_B\neq0$.
Let $d$ be a rational Darboux isomorphism for $(\xi,J_A,J_B)$. Let
$C_y,C_x\in M_{\beta\times\nu}(E)$ with $C_y^{\mathsf T}J_Bc(C_y)=0$ and
$C_x^{\mathsf T}J_Bc(C_x)=0$, such that $(C_y\ C_x)$ has a two-sided inverse.
Let $\chi$ be a character of $\mathbf A_E^\times$ with values of absolute value
$1$ that is trivial on $E^\times$, and let $\omega$ be a splitting for
$(d,J_B)$ that satisfies the Siegel rule for $(C_y,\chi)$. Then
$\theta(\omega(\gamma)\Phi)=\theta(\Phi)$ for every $\gamma\in U_{J_B}(E)$ and
every $\Phi\in\mathcal S(\mathbf A^\kappa)$.
\end{lemma}
\leanref{Development/AutomorphicSpectrum/Proofs/MembershipProofs/SeesawThetaDualPairStatement.lean}{GlobalTheory.T.weil_representation_theta_equivariance}
\leanref{Development/AutomorphicSpectrum/Proofs/MembershipProofs/SeesawThetaDualPair.lean}{GlobalTheory.CompanionTheory.ThetaSplit.theta_split}
\leanref{Development/AutomorphicSpectrum/Proofs/MembershipProofs/SeesawThetaDualPair.lean}{GlobalTheory.CompanionTheory.ThetaSplit.theta_split_inc}
\leanref{Development/AutomorphicSpectrum/Proofs/MembershipProofs/SeesawThetaDualPair.lean}{GlobalTheory.CompanionTheory.ThetaSplit.lagr_iff}
\begin{proof}
Let $d_0$ be the restriction of $d$ to $M_{\alpha\times\beta}(E)$, a bijection
onto $F^\kappa\times F^\kappa$. For $x,y\in M_{\alpha\times\beta}(E)$ put
$\langle x,y\rangle_F=\operatorname{Tr}_{E/F}(\xi h(x,y))$. As
$\operatorname{Tr}_{\mathbf A_E/\mathbf A}$ restricts to $\operatorname{Tr}_{E/F}$,
$\langle x,y\rangle_F=\langle d_0x,d_0y\rangle$ in the notation of
Definition~\ref{asm-theta-agree-invariance-heisenberg}(1).

\emph{Step 1.} For $p\in U_{J_B}(E)$ there is $a(p)\in\mathbf C$ with
$\theta(\omega(p)\Phi)=a(p)\theta(\Phi)$ for $\Phi\in\mathcal S$. By
Lemma~\ref{asm-theta-agree-invariance-theta}(2), $\theta$ agrees on $\mathcal S$
with an admissible form $\lambda_0$, and $\lambda_0\circ\omega(p)$ is admissible.
A rational element is $d(x)$ with $x\in M_{\alpha\times\beta}(E)$. Then
$\omega(p)\rho_d(x)\Phi=\rho_d(xp^{\mathsf T})\omega(p)\Phi$, and $d(xp^{\mathsf T})$
is rational. By Lemma~\ref{asm-theta-agree-invariance-theta}(3),
$\theta(\omega(p)\rho_d(x)\Phi)=\theta(\omega(p)\Phi)$. Lemma~\ref{asm-theta-agree-invariance-unique}
gives $a(p)$.

\emph{Step 2.} $a$ is a homomorphism $U_{J_B}(E)\to\mathbf C^\times$. Let
$\Phi_0\in\mathcal S$ with $\theta(\Phi_0)\neq0$. Then $a(1)=1$, and
$a(pq)\theta(\Phi_0)=\theta(\omega(p)\omega(q)\Phi_0)=a(p)a(q)\theta(\Phi_0)$.

\emph{Step 3.} The subspace $L=d_0\bigl(M_{\alpha\times\nu}(E)C_y^{\mathsf T}\bigr)$
of $F^\kappa\times F^\kappa$ is Lagrangian. Note that
$h(x,y)=\operatorname{tr}(J_A^{\mathsf T}xJ_Bc(y)^{\mathsf T})$. It suffices to
show that $x\in M_{\alpha\times\nu}(E)C_y^{\mathsf T}$ if and only if
$\langle x,zC_y^{\mathsf T}\rangle_F=0$ for every $z\in M_{\alpha\times\nu}(E)$.
If $x=uC_y^{\mathsf T}$, then
$h(x,zC_y^{\mathsf T})=\operatorname{tr}(J_A^{\mathsf T}uC_y^{\mathsf T}J_Bc(C_y)c(z)^{\mathsf T})=0$.
Conversely, suppose the condition holds, and put
$Y=J_A^{\mathsf T}xJ_Bc(C_y)$, so that
$h(x,zC_y^{\mathsf T})=\operatorname{tr}(Yc(z)^{\mathsf T})$. If $Y_{ij}\neq0$,
let $z$ have the single nonzero entry $c((\xi Y_{ij})^{-1})$ at $(i,j)$; then
$\xi h(x,zC_y^{\mathsf T})=1$, and $\operatorname{Tr}_{E/F}(1)=2\neq0$, a
contradiction. So $Y=0$, and $xJ_Bc(C_y)=0$ as $J_A$ is invertible. Let
$\left(\begin{smallmatrix}M_t\\M_b\end{smallmatrix}\right)$ be the inverse of $(C_y\ C_x)$, and write
$x=uC_y^{\mathsf T}+vC_x^{\mathsf T}$ with $u=xM_t^{\mathsf T}$ and
$v=xM_b^{\mathsf T}$. Then $vC_x^{\mathsf T}J_Bc(C_y)=0$ and
$vC_x^{\mathsf T}J_Bc(C_x)=0$, so $vC_x^{\mathsf T}J_Bc\bigl((C_y\ C_x)\bigr)=0$.
As $J_B$ and $(C_y\ C_x)$ are invertible, $vC_x^{\mathsf T}=0$ and
$x=uC_y^{\mathsf T}$.

\emph{Step 4.} Let $l$ and $\Phi_1$ be as in
Lemma~\ref{asm-theta-agree-invariance-average} for $L$. For
$z\in M_{\alpha\times\nu}(\mathbf A_E)$, $d(zC_y^{\mathsf T})$ lies in the
$\mathbf A$-span of $L$, because $M_{\alpha\times\nu}(\mathbf A_E)$ is spanned
over $\mathbf A$ by $M_{\alpha\times\nu}(E)$. So
$l(\rho_d(zC_y^{\mathsf T})\Phi)=l(\Phi)$ for $\Phi\in\mathcal S$.

\emph{Step 5.} Let $p\in U_{J_B}(E)$ and $t\in GL_\nu(E)$ with $pC_y=C_yt$. By
the Siegel rule and Step 4,
$l(\omega(p)\Phi_1)=\chi(\det t)|N(\det t)|^{|\alpha|/2}l(\Phi_1)=l(\Phi_1)$:
$\det t\in E^\times$, on which $\chi$ is trivial, and $N(\det t)\in F^\times$
has $|N(\det t)|=1$ by Theorem~\ref{asm-theta-agree-invariance-tate}(5). On the
other hand, apply Lemma~\ref{asm-theta-agree-invariance-average}(3) with
$s=d_0\circ(x\mapsto xp^{\mathsf T})\circ d_0^{-1}$,
$s_{\mathbf A}=d\circ(x\mapsto xp^{\mathsf T})\circ d^{-1}$, $T=\omega(p)$ and
$a=a(p)$. Here $s(L)=L$, since
$zC_y^{\mathsf T}p^{\mathsf T}=z(pC_y)^{\mathsf T}=(zt^{\mathsf T})C_y^{\mathsf T}$
and $t$ is invertible; $T$ preserves $\mathcal S$ and
$T\rho(u)=\rho(s_{\mathbf A}u)T$ on $\mathcal S$, as $\omega$ is a splitting;
and $\theta\circ T=a(p)\theta$ on $\mathcal S$ by Step 1. So
$l(\omega(p)\Phi_1)=a(p)l(\Phi_1)$. As $l(\Phi_1)\neq0$, $a(p)=1$.

\emph{Step 6.} By Step 5 and Lemma~\ref{asm-theta-agree-invariance-character},
applied to $J_B$, $C_y$ and $C_x$, the homomorphism $a$ is trivial. By Step 1,
$\theta(\omega(\gamma)\Phi)=\theta(\Phi)$.
\end{proof}

\begin{lemma}\label{asm-theta-agree-invariance-doubling}
Let $F$, $\psi$, $E$, $c$, $\xi$, $J_A$, $J_B$ and $\kappa$ be as in
Definition~\ref{asm-theta-agree-invariance-datum}, with $\xi$ invertible,
$c(\xi)=-\xi$ and $1,\xi$ an $F$-basis of $E$, and with
$c(J_A)^{\mathsf T}=J_A$, $c(J_B)^{\mathsf T}=J_B$ and $\det J_A,\det J_B\neq0$.
Let $d$ be a rational Darboux isomorphism for $(\xi,J_A,J_B)$, let $\chi$ be a
character of $\mathbf A_E^\times$ with values of absolute value $1$ that is
trivial on $E^\times$, and let $\omega$ be a Weil datum for
$(\xi,J_A,J_B,d,\chi)$. Then $\theta(\omega(\gamma)\Phi)=\theta(\Phi)$ for every
$\gamma\in U_{J_B}(E)$ and every $\Phi\in\mathcal S(\mathbf A^\kappa)$.
\end{lemma}
\leanref{Development/AutomorphicSpectrum/Proofs/MembershipProofs/ThetaSumTensorProductStatement.lean}{GlobalTheory.T.seesaw_theta_dual_pair}
\leanref{Development/AutomorphicSpectrum/Proofs/MembershipProofs/ThetaSumTensorProduct.lean}{GlobalTheory.CompanionTheory.ThetaMember.theta_om_rational}
\leanref{Development/AutomorphicSpectrum/Proofs/MembershipProofs/ThetaSumTensorProduct.lean}{GlobalTheory.CompanionTheory.ThetaMember.split_of_eq}
\leanref{Development/AbelianVarietyTheory/Definitions/Divisors/HilbertNinetyQuotientHomEqIff.lean}{GlobalTheory.WeilM.isLagr_diagCols}
\begin{proof}
Let $d^\square$ and $\omega^\square$ be as in
Definition~\ref{asm-theta-agree-invariance-datum}(4). The matrix
$J_B^\square=\operatorname{diag}(J_B,-J_B)$ has
$c(J_B^\square)^{\mathsf T}=J_B^\square$ and $\det J_B^\square\neq0$. The
matrices $\Delta=\left(\begin{smallmatrix}1\\1\end{smallmatrix}\right)$ and
$\Delta'=\left(\begin{smallmatrix}1\\-1\end{smallmatrix}\right)$ have
$\Delta^{\mathsf T}J_B^\square c(\Delta)=J_B-J_B=0$ and
$\Delta'^{\mathsf T}J_B^\square c(\Delta')=0$, and $(\Delta\ \Delta')$ has the
inverse $\frac12\left(\begin{smallmatrix}1&1\\1&-1\end{smallmatrix}\right)$. The
isomorphism $d^\square$ is rational, by its formula and the rationality of $d$.
So Lemma~\ref{asm-theta-agree-invariance-split} applies to $\omega^\square$ with
$C_y=\Delta$ and $C_x=\Delta'$: $\theta(\omega^\square(\gamma^\square)\Psi)=\theta(\Psi)$
for $\gamma^\square\in U_{J_B^\square}(E)$ and $\Psi\in\mathcal S(\mathbf A^{\kappa\sqcup\kappa})$.

Let $\gamma\in U_{J_B}(E)$, so that $\gamma^\square=\operatorname{diag}(\gamma,1)$
lies in $U_{J_B^\square}(E)$. By Lemma~\ref{asm-theta-agree-invariance-theta}(2)
there is $\Phi_0\in\mathcal S$ with $\theta(\Phi_0)\neq0$. For $\Phi\in\mathcal S$,
$\Phi\boxtimes\Phi_0\in\mathcal S(\mathbf A^{\kappa\sqcup\kappa})$ and
$\omega^\square(\gamma^\square)(\Phi\boxtimes\Phi_0)=(\omega(\gamma)\Phi)\boxtimes\Phi_0$.
By Lemma~\ref{asm-theta-agree-invariance-theta}(4),
\[
\theta(\omega(\gamma)\Phi)\,\theta(\Phi_0)
=\theta\bigl(\omega^\square(\gamma^\square)(\Phi\boxtimes\Phi_0)\bigr)
=\theta(\Phi\boxtimes\Phi_0)=\theta(\Phi)\,\theta(\Phi_0).
\]
Dividing by $\theta(\Phi_0)$ gives the claim.
\end{proof}

\begin{proposition}[Weil \cite{Wei64}]\label{asm-theta-agree-invariance-rational}
Let $(X,W,M,\chi_X,\chi_W,\omega)$ be a Weil pair for a theta base $B$ of a CM
field $E$ with maximal totally real subfield $F$
(Definition~\ref{asm-theta-agree-triple}(1)), with $X$ and $W$ of ranks $n$ and
$d$, and let $\Phi$ be a Schwartz-Bruhat function on $\mathbf A^{n\times d}$.
Then
\[
\sum_{\xi\in F^{n\times d}}\bigl(\omega(\gamma,\gamma')\Phi\bigr)(\xi)=\sum_{\xi\in F^{n\times d}}\Phi(\xi)
\]
for all $\gamma\in U_X(F)$ and $\gamma'\in U_W(F)$.
\end{proposition}
\leanref{Development/WeilBound/ThetaLifts/LeviDeterminants/ThetaSumInvariantRationalUnitary.lean}{PositivityTheory.ThKF0.theta_sum_rational_invariance}
\leanref{Development/WeilBound/ThetaLifts/LeviDeterminants/WeilRepRationalDarbouxFixesThetaDistribution.lean}{PositivityTheory.ThKF0.darboux_coordinates_rational_exist}
\leanref{Development/WeilBound/ThetaLifts/LeviDeterminants/WeilRepRationalDarbouxFixesThetaDistribution.lean}{PositivityTheory.ThKF0.Resc.member}
\leanref{Development/WeilBound/ThetaLifts/LeviDeterminants/WeilRepRationalDarbouxFixesThetaDistribution.lean}{PositivityTheory.ThKF0.Resc.isWeilB_resc}
\leanref{Development/WeilBound/ThetaLifts/LeviDeterminants/WeilRepRationalDarbouxFixesThetaDistribution.lean}{PositivityTheory.ThKF0.Resc.exists_ratCoord}
\leanref{Development/WeilBound/ThetaLifts/LeviDeterminants/WeilRepRationalDarbouxFixesThetaDistribution.lean}{PositivityTheory.ThKF0.Resc.exists_isQuad}
\begin{proof}
Let $J_X$ and $J_W$ be the matrices of $X$ and $W$, $e$ the isomorphism of the model $M$, and
$I=\{1,\ldots,n\}\times\{1,\ldots,d\}$. Both sides are theta sums
(Definition~\ref{asm-theta-agree-invariance-heisenberg}(3)) for the character
$\psi$ of $B$, which is continuous, trivial on $F$ and nontrivial on
$F\otimes\mathbf R$.

\emph{Step 1.} $\det J_X\neq0$ and $\det J_W\neq0$: at an infinite place $w$, the
matrices $\tau_w(J_X)$ and $i\,\tau_w(J_W)$ have signatures $(r,s)$ and
$(r',s')$ with $r+s=n$ and $r'+s'=d$, so they are invertible.

\emph{Step 2.} As $E\neq F$, there is $x\in E$ with $c(x)\neq x$. Then
$\delta=x-c(x)$ is nonzero, $c(\delta)=-\delta$, and $1,\delta$ is an $F$-basis
of $E$. The matrix $-\delta^{-1}J_W$ satisfies
$c(-\delta^{-1}J_W)^{\mathsf T}=-\delta^{-1}J_W$, as $c(J_W)^{\mathsf T}=-J_W$.

\emph{Step 3.} The isomorphism $e$ is rational in the sense of
Definition~\ref{asm-theta-agree-invariance-datum}(1): as $M$ is regular, $e$
maps $M_{n\times d}(E)$ into $F^I\times F^I$, injectively, and both are
$F$-vector spaces of dimension $2nd$. So is $x\mapsto e(x^{\mathsf T})$ on
$M_{d\times n}(\mathbf A_E)$.

\emph{Step 4.} By Definition~\ref{asm-theta-agree-models}(6)(a), the map
$h\mapsto\omega(1,h)$ is a Weil datum for $(-1,J_X,J_W,e,\chi_X)$
(Definition~\ref{asm-theta-agree-invariance-datum}(4)); the absolute value
$|x|=|N_{E/F}(x)|_{\mathbf A_F}$ of Definition~\ref{asm-theta-agree-doubling} is
the factor by which multiplication by $N_{E/F}(x)$ scales the Haar measures of
$\mathbf A$. Since
$-1=\delta\cdot1\cdot(-\delta^{-1})$, the form of
Definition~\ref{asm-theta-agree-invariance-datum}(1) for $(-1,J_X,J_W)$ is the
form for $(\delta,J_X,-\delta^{-1}J_W)$. Also $U_{-\delta^{-1}J_W}=U_{J_W}$, and
likewise for the doubled matrices. So $h\mapsto\omega(1,h)$ is a Weil datum for
$(\delta,J_X,-\delta^{-1}J_W,e,\chi_X)$. The Hecke character $\chi_X$ is trivial
on $E^\times$. By Step 3 and Lemma~\ref{asm-theta-agree-invariance-doubling},
$\theta(\omega(1,\gamma')\Psi)=\theta(\Psi)$ for every Schwartz-Bruhat
function $\Psi$.

\emph{Step 5.} By Definition~\ref{asm-theta-agree-models}(6)(b), in the same way,
$g\mapsto\omega(g,1)$ is a Weil datum for
$(-1,J_W,J_X,x\mapsto e(x^{\mathsf T}),\chi_W)$, hence for
$(\delta,-\delta^{-1}J_W,J_X,x\mapsto e(x^{\mathsf T}),\chi_W)$, since
$-1=\delta\cdot(-\delta^{-1})\cdot1$. By
Lemma~\ref{asm-theta-agree-invariance-doubling},
$\theta(\omega(\gamma,1)\Psi)=\theta(\Psi)$ for every Schwartz-Bruhat function
$\Psi$.

\emph{Step 6.} As $\omega$ is a representation,
$\omega(\gamma,\gamma')\Phi=\omega(\gamma,1)\bigl(\omega(1,\gamma')\Phi\bigr)$, and
$\omega(1,\gamma')\Phi$ is a Schwartz-Bruhat function. Steps 5 and 4 give
$\theta(\omega(\gamma,\gamma')\Phi)=\theta(\omega(1,\gamma')\Phi)=\theta(\Phi)$.
\end{proof}

\begin{lemma}\label{asm-theta-agree-anisotropic-automorphy}
Let $(X,W,M,\chi_X,\chi_W,\omega)$ be a Weil pair for a theta base $B$ of a CM
field $E$ (Definition~\ref{asm-theta-agree-triple}(1)), with $X$ and $W$ of
ranks $n$ and $d$, and let $\Phi$ be a Schwartz-Bruhat function on
$\mathbf A^{n\times d}$. Then
$\theta_\omega(\Phi)(\gamma g,\gamma'h)=\theta_\omega(\Phi)(g,h)$ for all
$\gamma\in U_X(F)$, $\gamma'\in U_W(F)$, $g\in U_X$ and $h\in U_W$.
\end{lemma}
\leanref{Development/AbelianVarietyTheory/Definitions/Divisors/CTSmallCone.lean}{GlobalTheory.ThK.ThetaBasic}
\leanref{Development/WeilBound/ThetaLifts/LeviDeterminants/ThetaSumInvariantRationalUnitary.lean}{PositivityTheory.ThKF0.theta_sum_rational_invariance}
\begin{proof}
As $\omega$ is a representation, $\theta_\omega(\Phi)(\gamma g,\gamma'h)$ is the
theta sum $\sum_{\xi}(\omega(\gamma,\gamma')\Psi)(\xi)$, $\xi\in F^{n\times d}$,
of $\Psi=\omega(g,h)\Phi$, which is a Schwartz-Bruhat function. So it suffices
to show $\sum_\xi(\omega(\gamma,\gamma')\Psi)(\xi)=\sum_\xi\Psi(\xi)$ for every
Schwartz-Bruhat function $\Psi$, and this is
Proposition~\ref{asm-theta-agree-invariance-rational}.
% Lean: theta_sum_rational_invariance (Proposition asm-theta-agree-invariance-rational);
% thetaBasic_of (ThetaBasicOf.lean) proves this clause of ThetaBasic from it.
% ThetaBasic is also used by Lemmas asm-theta-agree-decay, -convergence and -continuity.
\end{proof}

\begin{lemma}\label{asm-theta-agree-anisotropic-invariance}
Let $(X,W,M,\chi_X,\chi_W,\omega)$ be a Weil pair for a theta base $B$ of a CM
field $E$ (Definition~\ref{asm-theta-agree-triple}(1)), with $X$ and $W$ of
ranks $n$ and $d$, and let $\Phi$ be a Schwartz-Bruhat function on
$\mathbf A^{n\times d}$. Take the integral of a function that is not integrable
to be $0$ in the theta integral $I$ (Definition~\ref{asm-theta-agree-models}(7)).
Then $I(\omega(g_0,1)\Phi)=I(\Phi)$ for every $g_0\in U_X$.
\end{lemma}
\leanref{Development/WeilBound/ThetaLifts/CartanFactorisation/Factor.lean}{PositivityTheory.ThKF.ILeft}
\leanref{Development/WeilBound/ThetaLifts/CartanFactorisation/ExistsBeta.lean}{PositivityTheory.ThKF.I_left}
\leanref{Development/WeilBound/ThetaLifts/CartanFactorisation/IntegrableOnRight.lean}{PositivityTheory.ThKF.L2.integral_right}
\leanref{Development/WeilBound/ThetaLifts/CartanFactorisation/IntegrableOnRight.lean}{PositivityTheory.ThKF.L2.integrableOn_right}
\leanref{Development/WeilBound/ThetaLifts/CartanFactorisation/Unimod.lean}{PositivityTheory.ThKF.unimodularity_holds}
\begin{proof}
Fix $h\in U_W$ and put $f(g)=\theta_\omega(\Phi)(g,h)$. As $\omega$ is a
representation, $\theta_\omega(\omega(g_0,1)\Phi)(g,h)=f(gg_0)$. By
Lemma~\ref{asm-theta-agree-anisotropic-automorphy}, $f(\gamma g)=f(g)$ for
$\gamma\in U_X(F)$.

The group $U_X(F)$ is countable, and $U_X$ is unimodular, being the group of
adelic points of a unitary group. So the Haar measure $\mu$ of the regular
model $M$ is invariant under right translations. Right translation by $g_0$
commutes with left translation by $U_X(F)$, so it carries the fundamental
domain $\mathcal D$ to a fundamental domain $\mathcal Dg_0$, and it preserves
$\mu$. Hence
\[
\int_{\mathcal D}f(gg_0)\,d\mu(g)=\int_{\mathcal Dg_0}f\,d\mu=\int_{\mathcal D}f\,d\mu .
\]
The second equality holds because a function invariant under $U_X(F)$ on the
left has the same integral over any two fundamental domains $\mathcal D$ and
$\mathcal D'$:
$\int_{\mathcal D'}f\,d\mu=\sum_\gamma\int_{\mathcal D'\cap\gamma\mathcal D}f\,d\mu
=\sum_\gamma\int_{\gamma^{-1}\mathcal D'\cap\mathcal D}f\,d\mu=\int_{\mathcal D}f\,d\mu$,
the sums over $\gamma\in U_X(F)$. These computations hold for $f$ integrable on
$\mathcal D$. The same computations for $|f|$, with values in $[0,\infty]$, show
that $g\mapsto f(gg_0)$ is integrable on $\mathcal D$ if and only if $f$ is. If
they are, the display is $I(\omega(g_0,1)\Phi)(h)=I(\Phi)(h)$. If they are not,
both sides are $0$.
\end{proof}

\begin{lemma}\label{asm-theta-agree-anisotropic-augmentation}
Let $(X,W,M,\chi_X,\chi_W,\omega)$ be a Weil pair for a theta base $B$ of a CM
field $E$ (Definition~\ref{asm-theta-agree-triple}(1)), with $X$ and $W$ of
ranks $n$ and $d$, such that every Schwartz-Bruhat function on $\mathbf A^I$,
$I=\{1,\ldots,n\}\times\{1,\ldots,d\}$, converges. Let $v$ be a finite place,
and let $\mathfrak a_v(\omega)$ be the span of the $\phi_1-\phi$ with $\phi$
and $\phi_1$ Schwartz-Bruhat on $F_v^I$ and $\omega(g,1)$ acting at $v$ by
$\phi\mapsto\phi_1$ for some $g\in U_X^{(v)}$
(Definition~\ref{asm-theta-agree-models}(2)). Let $f\in\mathcal S(F_\infty^I)$,
and for each finite place $u\neq v$ let $\phi_u$ be a Schwartz-Bruhat function
on $F_u^I$, with $\phi_u=\mathrm{vac}_u$ for almost all $u$. For a function $z$
on $F_v^I$, let $\Phi_z=f\otimes\bigotimes_u\phi_u$ with $\phi_v=z$. Then
$I(\Phi_z)=0$ for every $z\in\mathfrak a_v(\omega)$.
\end{lemma}
\leanref{Development/WeilBound/ThetaLifts/CartanFactorisation/ExistsEnlarge.lean}{PositivityTheory.ThKF.kill_aug}
\begin{proof}
Let $r(x)=f(x_\infty)\prod_{u\neq v}\phi_u(x_u)$. It is a function on
$\mathbf A^I$ that does not depend on the $v$-component, and
$\Phi_z(x)=z(x_v)r(x)$. For Schwartz-Bruhat $z$ the function $\Phi_z$ is
Schwartz-Bruhat, since $M$ is regular and so $\mathrm{vac}_u$ is the indicator
function of $\mathcal O_u^I$ for almost all $u$; hence $\Phi_z$ converges. The
theta sum is additive, and so is the integral of integrable functions, so
$z\mapsto I(\Phi_z)$ is linear on the Schwartz-Bruhat functions on $F_v^I$. It
therefore suffices to show $I(\Phi_{\phi_1})=I(\Phi_\phi)$ when $\omega(g,1)$
acts at $v$ by $\phi\mapsto\phi_1$ for some $g\in U_X^{(v)}$. By the definition
of acting at $v$, $\omega(g,1)\Phi_\phi=\Phi_{\phi_1}$, and
Lemma~\ref{asm-theta-agree-anisotropic-invariance} gives
$I(\Phi_{\phi_1})=I(\omega(g,1)\Phi_\phi)=I(\Phi_\phi)$.
\end{proof}
% Lean: the linear extension J of exists_Ilin (Factor.lean) carries the linearity of I on
% Schwartz-Bruhat functions; r is the function of PureFactor.exists_rest_factor.

\begin{lemma}\label{asm-theta-agree-anisotropic}
Let $(B,G,\mathcal X,\omega,\omega',R,R_f)$ be a matched quadruple for a theta
base $B$ of a CM field $E$ (Definition~\ref{asm-theta-agree-matched}), with
$q=0$, and let $c$ be a polynomial constant of it. Then $I(\Phi)=c\,I'(\Phi')$
whenever $\Phi$ and $\Phi'$ correspond.
\end{lemma}
\leanref{Development/AbelianVarietyTheory/Definitions/Divisors/CTSmallCone.lean}{GlobalTheory.ThK.AnAll}
\leanref{Development/WeilBound/ThetaLifts/CartanFactorisation/ExistsEnlarge.lean}{PositivityTheory.ThKF.anAll_of}
\begin{proof}
Let $T_1\supset T$, $\Phi=\sum_if_i\otimes\bigotimes_v\phi_{i,v}$ and
$\Phi'=\sum_jf'_j\otimes\bigotimes_v\phi'_{j,v}$ be as in the definition of
corresponding data (Definition~\ref{asm-theta-agree-matched}(4)).

The space $V'$ has signature $(q,a'-q)=(0,a')$. So
Lemma~\ref{asm-theta-agree-convergence}(1), applied to the Weil pair
$(V',W,M',\alpha,\alpha/\beta,\omega')$, shows that every Schwartz-Bruhat
function converges for $(M',\omega')$. In particular $I'$ is linear on the
Schwartz-Bruhat functions.

The place $v_1$ is clean for $\mathcal X$, as $\mathcal X$ is proper. By
Lemma~\ref{asm-theta-agree-coinvariants}, for each $j$ there is a
Schwartz-Bruhat function $y_j$ on $F_{v_1}^{I'_0}$, of full rank at $\tau_1$,
with $p'_{v_1}(y_j)=p'_{v_1}(\phi'_{j,v_1})$; by the same lemma
$y_j-\phi'_{j,v_1}\in\mathfrak a_{v_1}(\omega')$, since $p'_{v_1}$ vanishes on
it. Let $\phi^0_{j,v}=\phi'_{j,v}$ for $v\neq v_1$ and $\phi^0_{j,v_1}=y_j$, and
put $\Phi'_0=\sum_jf'_j\otimes\bigotimes_v\phi^0_{j,v}$. As
$f'_j\in\operatorname{Pol}(M')$, $\phi^0_{j,v}=\mathrm{vac}'_v$ for
$v\notin T_1\cup\{v_1\}$, and $\phi^0_{j,v_1}$ has full rank at $\tau_1$, we have
$\Phi'_0\in\mathcal S^\circ(M')$.

Next, $\Phi$ and $\Phi'_0$ correspond, with the finite set
$T_2=T_1\cup\{v_1\}$. The components of both are vacua outside $T_2$. If
$v_1\in T_1$, the identity in $\Pi\otimes\bigotimes_{v\in T_1}\Pi_v$ for $\Phi$
and $\Phi'$ is the one required, since $p'_{v_1}(y_j)=p'_{v_1}(\phi'_{j,v_1})$.
If $v_1\notin T_1$, apply to that identity the linear map
$x\mapsto x\otimes p_{v_1}(\mathrm{vac}_{v_1})$ from
$\Pi\otimes\bigotimes_{v\in T_1}\Pi_v$ to $\Pi\otimes\bigotimes_{v\in T_2}\Pi_v$.
On the left, $p_{v_1}(\phi_{i,v_1})=p_{v_1}(\mathrm{vac}_{v_1})$. On the right,
$p'_{v_1}(\phi^0_{j,v_1})=p'_{v_1}(\phi'_{j,v_1})=p'_{v_1}(\mathrm{vac}'_{v_1})
=p_{v_1}(\mathrm{vac}_{v_1})$, the last because $v_1\notin T$ and the finite
maps are matched. This gives the identity for $T_2$.

As $c$ is a polynomial constant and $\Phi'_0\in\mathcal S^\circ(M')$,
$I(\Phi)=c\,I'(\Phi'_0)$. Since $f\otimes\bigotimes_v\phi_v$ is linear in
$\phi_{v_1}$, the difference $\Phi'_0-\Phi'$ is
$\sum_jf'_j\otimes\bigotimes_v\phi''_{j,v}$ with $\phi''_{j,v}=\phi'_{j,v}$ for
$v\neq v_1$ and $\phi''_{j,v_1}=y_j-\phi'_{j,v_1}\in\mathfrak a_{v_1}(\omega')$.
Each of its terms has $I'=0$ by
Lemma~\ref{asm-theta-agree-anisotropic-augmentation} for
$(V',W,M',\alpha,\alpha/\beta,\omega')$ at $v_1$. So $I'(\Phi'_0)=I'(\Phi')$, and
$I(\Phi)=c\,I'(\Phi')$.
\end{proof}
% Lean: anAll_of WeilConv KCoinv (ThKF.iLeft thetaBasic). The second and third clauses
% of Coinv at v_1 (c2, c3) give y_j; exists_enlarge builds the map to T_2 = insert v_1 T;
% pure_mem_So gives Phi'_0 in So; pureAct is the linearity in the v_1-component.

\subsection{Changing the lattice spaces}\label{subsection-asm-changing-lattice-spaces}

\begin{definition}\label{asm-theta-agree-rechoice-terms}
Let $(X,W,M,\chi_X,\chi_W,\omega)$ be a Weil pair for a theta base $B$ of a CM
field $E$ (Definition~\ref{asm-theta-agree-triple}(1)), with $X$ and $W$ of
ranks $n$ and $d$, and put $I=\{1,\ldots,n\}\times\{1,\ldots,d\}$.
\begin{enumerate}
\item Let $g\in U_X$ and $h\in U_W$ be at $\infty$
  (Definition~\ref{asm-theta-agree-models}(2)). An \emph{archimedean operator}
  of $\omega$ at $(g,h)$ is a continuous linear operator $T$ on
  $\mathcal S(F_\infty^I)$ such that $\omega(g,h)(f\otimes\phi)=T(f)\otimes\phi$
  for every $f\in\mathcal S(F_\infty^I)$ and every Schwartz-Bruhat function
  $\phi$ on $\mathbf A_f^I$, and $\omega(g,h)$ acts at $\infty$ by
  $f\mapsto T(f)$ for every $f$.
\item A \emph{theta subspace} of $\operatorname{Pol}(M)$ is a subspace
  $N\subset\operatorname{Pol}(M)$ with the following properties. If
  $f\in\operatorname{Pol}(M)$, $\gamma$ is a line in $U_X$
  (Definition~\ref{asm-theta-agree-matched}(1)) and $t\mapsto\omega(\gamma(t),1)$
  acts at $\infty$ to first order by $f\mapsto f_1$, then $f_1\in N$. If
  $f\in\operatorname{Pol}(M)$, $k\in K_\infty(X)$ and $\omega(k,1)$ acts at
  $\infty$ by $f\mapsto f_1$, then $f_1-f\in N$. If $f\in N$ and $f\mapsto f_1$
  is the action at $\infty$ to first order of $t\mapsto\omega(1,\delta(t))$ for
  a line $\delta$ in $U_W$, or the action at $\infty$ of $\omega(1,k)$ for some
  $k\in K_\infty(W)$, then $f_1\in N$. These are the theta subspaces of
  Theorem~\ref{lit-howe-duality-real} for $\operatorname{Pol}(M)$.
\item Let $v$ be a finite place, and let $\mathfrak a_v(\omega)$ be the span of
  the $\phi_1-\phi$ with $\phi$ and $\phi_1$ Schwartz-Bruhat on $F_v^I$ and
  $\omega(g,1)$ acting at $v$ by $\phi\mapsto\phi_1$ for some $g\in U_X^{(v)}$.
  A \emph{theta subspace at $v$} of $\omega$ is a subspace $N$ with
  $\mathfrak a_v(\omega)\subset N\subset\mathcal S(F_v^I)$ such that
  $\phi_1\in N$ whenever $\phi\in N$, $\phi_1$ is Schwartz-Bruhat,
  $h\in U_W^{(v)}$ and $\omega(1,h)$ acts at $v$ by $\phi\mapsto\phi_1$.
\item Let $Y$ and $Y^\circ$ be lattice spaces of the same rank with the same
  matrix (Definition~\ref{asm-theta-agree-doubling}(1)). An element $g\in U_Y$
  at $\infty$ \emph{moves} $Y^\circ$ to $Y$ if
  $\pi_w(g)\theta_{Y^\circ,w}=\theta_{Y,w}\pi_w(g)$ for every infinite place
  $w$ of $F$.
\end{enumerate}
\end{definition}
\leanref{Development/WeilBound/ThetaLifts/CartanFactorisation/CartanConj.lean}{PositivityTheory.ThKF.IsArchOp}
\leanref{Development/AbelianVarietyTheory/PublishedTheorems/Divisors₂/HoweDualityReal.lean}{GlobalTheory.Published.ThetaSubInf}
\leanref{Development/AbelianVarietyTheory/PublishedTheorems/Divisors₂/HoweDualityUnitary.lean}{GlobalTheory.Published.ThetaSub}
\leanref{Development/WeilBound/ThetaLifts/CartanFactorisation/CommuteConj.lean}{PositivityTheory.ThKF.Moves}

\begin{lemma}\label{asm-theta-agree-rechoice-schur}
Let $Q$, $Q'$, $P$ and $P_2$ be complex vector spaces with $P_2\neq0$, let
$S\subset Q$ and $S'\subset Q'$ be subspaces, and let $p\colon Q\to P$,
$p'\colon Q'\to P$, $p_2\colon Q\to P_2$ and $p'_2\colon Q'\to P_2$ be linear
maps with $p(S)=p'(S')=P$ and $p_2(S)=p'_2(S')=P_2$. Let $\iota$ be a set, and
for $i\in\iota$ let $A_i\colon S\to S$ be a linear map and $A'_i\colon S'\to S'$
a map. Assume:
\begin{enumerate}
\item[(a)] for $x\in S$, $p(x)=0$ implies $p_2(x)=0$, and for $x'\in S'$,
  $p'(x')=0$ if and only if $p'_2(x')=0$;
\item[(b)] for all $i\in\iota$, $x\in S$ and $x'\in S'$, $p(x)=p'(x')$ implies
  $p(A_ix)=p'(A'_ix')$, and $p_2(x)=p'_2(x')$ implies
  $p_2(A_ix)=p'_2(A'_ix')$;
\item[(c)] $0$ and $P$ are the only subspaces $N\subset P$ with $p(A_ix)\in N$
  whenever $i\in\iota$, $x\in S$ and $p(x)\in N$;
\item[(d)] $S$ has countable dimension.
\end{enumerate}
Then there are a linear map $e\colon P\to P_2$ and $c\in\mathbf C$ with
$c\neq0$, $e(p(x))=p_2(x)$ for $x\in S$, and $e(p'(x'))=c\,p'_2(x')$ for
$x'\in S'$.
\end{lemma}
\leanref{Development/WeilBound/ThetaLifts/CartanFactorisation/CountableCompactOpen.lean}{PositivityTheory.ThKF.schur_match}
\leanref{Development/WeilBound/ThetaLifts/CartanFactorisation/ExistsPureML.lean}{PositivityTheory.ThKF.exists_section}
\begin{proof}
As $p(S)=P$ and $p'_2(S')=P_2$, there are linear maps $s\colon P\to S$ and
$t\colon P_2\to S'$ with $p(s(y))=y$ and $p'_2(t(y))=y$. For $x\in S$, the
element $s(p(x))-x$ of $S$ lies in the kernel of $p$, so
$p_2(s(p(x)))=p_2(x)$ by (a). For $x'\in S'$, the element $t(p'_2(x'))-x'$ of
$S'$ lies in the kernel of $p'_2$, so $p'(t(p'_2(x')))=p'(x')$ by (a).

For $i\in\iota$ put $\rho_i=p\circ A_i\circ s$, an operator on $P$. Then
$\rho_i(p(x))=p(A_ix)$ for $x\in S$. Indeed, choose $x'\in S'$ with
$p'(x')=p(x)$. By (b), $p(A_ix)=p'(A'_ix')$, and
$p(A_is(p(x)))=p'(A'_ix')$ since $p(s(p(x)))=p(x)=p'(x')$.

Put $U=p'\circ t\circ p_2\circ s$, an operator on $P$. If $x\in S$ and
$x'\in S'$ have $p_2(x)=p'_2(x')$, then
$U(p(x))=p'(t(p_2(x)))=p'(t(p'_2(x')))=p'(x')$. The operator $U$ commutes with
every $\rho_i$. Let $y=p(x)$ with $x\in S$, and choose $x'\in S'$ with
$p'_2(x')=p_2(x)$ and $x_3\in S$ with $p(x_3)=p'(x')$. Then
$U(y)=p'(x')=p(x_3)$, so $\rho_i(U(y))=p(A_ix_3)=p'(A'_ix')$ by (b). Also
$\rho_i(y)=p(A_ix)$ and $p_2(A_ix)=p'_2(A'_ix')$ by (b), so
$U(\rho_i(y))=p'(A'_ix')$.

The space $P=p(S)$ has countable dimension by (d). By (c) and the formula for
$\rho_i(p(x))$, $0$ and $P$ are the only subspaces of $P$ stable under every
$\rho_i$. By Lemma~\ref{th-first-lift-dixmier}, $U$ is multiplication by some
$c\in\mathbf C$.

Let $x\in S$ and $x'\in S'$ with $p(x)=p'(x')$, and choose $y\in S'$ with
$p'_2(y)=p_2(x)$. Then $p'(y)=U(p(x))=c\,p(x)=c\,p'(x')$, so
$p'(y-cx')=0$, and $p'_2(y-cx')=0$ by (a). Hence $p_2(x)=p'_2(y)=c\,p'_2(x')$.

Put $e=p_2\circ s$. Then $e(p(x))=p_2(x)$ for $x\in S$. For $x'\in S'$, choose
$x\in S$ with $p(x)=p'(x')$; then $e(p'(x'))=p_2(x)=c\,p'_2(x')$. Finally,
choose $x\in S$ with $p_2(x)\neq0$, possible as $P_2\neq0$, and $x'\in S'$ with
$p'(x')=p(x)$. Then $p_2(x)=c\,p'_2(x')$, so $c\neq0$.
\end{proof}

\begin{lemma}\label{asm-theta-agree-rechoice-operators}
Let $(X,W,M,\chi_X,\chi_W,\omega)$ be a Weil pair for a theta base $B$ of a CM
field $E$ (Definition~\ref{asm-theta-agree-triple}(1)), with $X$ and $W$ of
ranks $n$ and $d$, and put $I=\{1,\ldots,n\}\times\{1,\ldots,d\}$.
\begin{enumerate}
\item For $g\in U_X$ and $h\in U_W$ at $\infty$ there is an archimedean
  operator of $\omega$ at $(g,h)$
  (Definition~\ref{asm-theta-agree-rechoice-terms}(1)).
\item Let $x,y\in U_X\times U_W$ and $f,f_1,f_2\in\mathcal S(F_\infty^I)$. If
  $\omega(y)$ acts at $\infty$ by $f\mapsto f_1$ and $\omega(x)$ acts at
  $\infty$ by $f_1\mapsto f_2$, then $\omega(xy)$ acts at $\infty$ by
  $f\mapsto f_2$.
\item Let $T$ be an archimedean operator of $\omega$ at $z=(g,h)$, let $x$ be
  an element of $U_X\times U_W$ whose two components are at $\infty$, and let
  $y\in U_X\times U_W$ with $xz=zy$. If $\omega(y)$ acts at $\infty$ by
  $f\mapsto f_1$, then
  $\omega(x)$ acts at $\infty$ by $T(f)\mapsto T(f_1)$.
\end{enumerate}
\end{lemma}
\leanref{Development/WeilBound/ThetaLifts/CartanFactorisation/CartanConj.lean}{PositivityTheory.ThKF.ArchOp}
\leanref{Development/WeilBound/ThetaLifts/CartanFactorisation/ArchimedeanEvaluationVector.lean}{PositivityTheory.ThKF.A.archOp}
\leanref{Development/WeilBound/ThetaLifts/CartanFactorisation/ArchimedeanEvaluationVector.lean}{PositivityTheory.ThKF.A.exists_op}
\leanref{Development/WeilBound/ThetaLifts/CartanFactorisation/ExistsCharNe.lean}{PositivityTheory.ThKF.A.exists_lin}
\leanref{Development/WeilBound/ThetaLifts/CartanFactorisation/ExistsCharNe.lean}{PositivityTheory.ThKF.A.apply_pure}
\leanref{Development/WeilBound/ThetaLifts/CartanFactorisation/EqOfInner.lean}{PositivityTheory.ThKF.E.Impl.continuous}
\leanref{Development/WeilBound/ThetaLifts/CartanFactorisation/ActsComp.lean}{PositivityTheory.ThKF.acts_comp}
\leanref{Development/WeilBound/ThetaLifts/CartanFactorisation/CommuteConj.lean}{PositivityTheory.ThKF.transK}
\begin{proof}
Let $\rho$ be the Schrödinger operators of $M$, with isomorphism $e$, and call
$m\in\mathbf A^I\times\mathbf A^I$ \emph{finite} if its components at $\infty$
are $0$.

(1) Put $S=\omega(g,h)$.

\emph{Step 1.} $S$ maps Schwartz-Bruhat functions to Schwartz-Bruhat functions,
by conditions (a) and (b) of a Weil representation, and $S\rho(m)=\rho(m)S$ on
them for every finite $m$. Indeed, write $m=e(u)$. By the same conditions,
$S\rho(u)=\rho(guh^{\mathsf T})S$ on Schwartz-Bruhat functions. The components
of $u$ at $\infty$ are $0$, and $g-1$ and $h-1$ are at $\infty$, so
$(g-1)u=0$ and $u(h-1)^{\mathsf T}=0$; hence $guh^{\mathsf T}=u$.

\emph{Step 2.} Let $g_0$ be a Schwartz-Bruhat function on $\mathbf A_f^I$ with
$g_0(0)=1$, the indicator function of a compact open subgroup. Let
$f\in\mathcal S(F_\infty^I)$, and let $\phi$ be Schwartz-Bruhat on
$\mathbf A_f^I$ with $\phi(0)=0$. For $z\neq0$ in $\mathbf A_f^I$ there is
$w\in\mathbf A_f^I$ with $\psi(z\cdot w)\neq1$, as $\psi$ is not trivial on any
$F_v$. The support of $\phi$ is compact and does not contain $0$, so there are
$w_1,\ldots,w_k$ such that $G(x)=\sum_j|\psi(x\cdot w_j)-1|^2$ is positive on
it. Put $\phi_j(x)=\overline{(\psi(x\cdot w_j)-1)}\,G(x)^{-1}\phi(x)$ on the
support of $\phi$ and $\phi_j=0$ elsewhere; these are Schwartz-Bruhat, and
$\phi=\sum_j(\psi(x\cdot w_j)-1)\phi_j$. View elements of $\mathbf A_f^I$ as
elements of $\mathbf A^I$ with component $0$ at $\infty$. In terms of the
finite elements $m_j=(0,w_j)$, this says
$f\otimes\phi=\sum_j(\rho(m_j)-1)(f\otimes\phi_j)$.
By Step 1, $S(f\otimes\phi)=\sum_j(\rho(m_j)-1)S(f\otimes\phi_j)$, which
vanishes at every point with finite component $0$, since $\psi(0)=1$.

Now let $\phi$ be any Schwartz-Bruhat function on $\mathbf A_f^I$, and let
$y\in F_\infty^I$ and $z\in\mathbf A_f^I$. The operator $\rho(z,0)$ is the
translation by $z$, which changes only the finite component, so Step 1 gives
$S(f\otimes\phi)(y,z)=S(f\otimes\phi(\,\cdot+z))(y,0)$. As
$\phi(\,\cdot+z)-\phi(z)g_0$ vanishes at $0$, the previous paragraph gives
$S(f\otimes\phi)(y,z)=\phi(z)\,S(f\otimes g_0)(y,0)$. Put
$T(f)(y)=S(f\otimes g_0)(y,0)$, a Schwartz function, as the restriction of a
Schwartz-Bruhat function to $F_\infty^I\times\{0\}$. Then
$S(f\otimes\phi)=T(f)\otimes\phi$, and $T$ is linear.

\emph{Step 3.} $T$ is continuous. By condition (b) of a Weil representation,
$S$ preserves $L^2$-norms, so $\|T(f)\|_2=\|f\|_2$. Let $s$ be the
$\mathbf R$-linear automorphism of $F_\infty^I\times F_\infty^I$ with
$s(e(u)_\infty)=e(guh^{\mathsf T})_\infty$ for $u$ at $\infty$, and $\rho_\infty$
the Schrödinger operators on $\mathcal S(F_\infty^I)$. Applying
$S\rho(u)=\rho(guh^{\mathsf T})S$ to $f\otimes g_0$ gives
$T(\rho_\infty(a)f)=\rho_\infty(s(a))T(f)$ for every $a$. The operator $T$ is
onto, since the operator built in the same way from $(g^{-1},h^{-1})$ is
inverse to it. Differentiating
$\langle T\varphi,T\rho_\infty(ta)f\rangle=\langle\varphi,\rho_\infty(ta)f\rangle$
at $t=0$, with $T\rho_\infty(ta)f=\rho_\infty(ts(a))Tf$ and $T$ onto, gives
$T\circ D_a=D_{s(a)}\circ T$, where $D_a=D_{a_1,a_2}$ for $a=(a_1,a_2)$ is the
operator of Lemma~\ref{asm-theta-agree-real-derivative}. So every element $A$ of the
algebra generated by the $D_a$ has $A\circ T=T\circ A'$ for some $A'$ in it.
Every Schwartz seminorm $q$ is bounded by a finite sum
$q(f)\le\sum_r\|A_rf\|_2$ with $A_r$ in this algebra, by Sobolev's inequality.
Then $q(Tf)\le\sum_r\|TA'_rf\|_2=\sum_r\|A'_rf\|_2$, which is continuous in
$f$.

\emph{Step 4.} $S(0)=0$, and $S$ acts weakly at $\infty$ by $f\mapsto T(f)$
(Definition~\ref{asm-theta-agree-real-fock}(7)) by Step 2, so it acts at
$\infty$ by $f\mapsto T(f)$ by Lemma~\ref{asm-theta-agree-real-weak}(2). Thus
$T$ is an archimedean operator at $(g,h)$.

(2) Let $r$ be a function on $\mathbf A^I$ that does not depend on the
component at $\infty$, with $f\otimes r$ Schwartz-Bruhat. Then
$\omega(y)(f\otimes r)=f_1\otimes r$ is Schwartz-Bruhat, and
$\omega(xy)(f\otimes r)=\omega(x)(f_1\otimes r)=f_2\otimes r$.

(3) Let $T_1$ be an archimedean operator of $\omega$ at $x$, by (1). By (2),
$\omega(xz)$ acts at $\infty$ by $f\mapsto T_1(T(f))$ and $\omega(zy)$ by
$f\mapsto T(f_1)$. As $xz=zy$, $T_1(T(f))=T(f_1)$ by
Lemma~\ref{asm-theta-agree-real-weak}(3). So $\omega(x)$ acts at $\infty$ by
$T(f)\mapsto T(f_1)$.
\end{proof}
% Lean: (1) exists_op: exists_strictHenselian_absolutely_closed_local_ring (Step 2, through
% exists_lin, apply_pure, exists_char_ne, exists_schwartz_slice and phB_fin), impl_of and
% surj_of with Impl.continuous (Step 3: Impl.dRho, weylAlg_induction, exists_seminorm_le), and
% acts_of (Step 4). The clause of ArchOp for K_inf is map_Pol; see
% Lemma asm-theta-agree-rechoice-fock.

\begin{lemma}\label{asm-theta-agree-rechoice-intertwine}
Let $(X,W,M,\chi_X,\chi_W,\omega)$ be a Weil pair for a theta base $B$ of a CM
field $E$, with $X$ and $W$ of ranks $n$ and $d$ and
$I=\{1,\ldots,n\}\times\{1,\ldots,d\}$, let $\rho$ be the Schrödinger operators
of $M$, and let $D_u$, for $u\in\mathbb W(X,W)$ at $\infty$, be the operators
of Lemma~\ref{asm-theta-agree-real-derivative} for the isomorphism of $M$.
\begin{enumerate}
\item For $u$ at $\infty$ and $f,f_1\in\mathcal S(F_\infty^I)$, the family
  $t\mapsto\rho(tu)$ acts at $\infty$ to first order by $f\mapsto f_1$ if and
  only if $f_1=D_uf$.
\item Let $T$ be an archimedean operator of $\omega$ at $(g,h)$. Then
  $T(D_uf)=D_{guh^{\mathsf T}}(T(f))$ for every $u$ at $\infty$ and
  $f\in\mathcal S(F_\infty^I)$.
\end{enumerate}
\end{lemma}
\leanref{Development/WeilBound/ThetaLifts/CartanFactorisation/Cartan.lean}{PositivityTheory.ThKF.Cart.intertw}
\begin{proof}
(1) By Lemma~\ref{asm-theta-agree-real-derivative}, $t\mapsto\rho(tu)$ acts
weakly at $\infty$ to first order by $f\mapsto D_uf$, and $\rho(tu)(0)=0$. By
Lemma~\ref{asm-theta-agree-real-weak}(1) it acts at $\infty$ to first order by
$f\mapsto D_uf$, and by part (3) of that lemma by no other map.

(2) Let $\phi$ be Schwartz-Bruhat on $\mathbf A_f^I$ and $\lambda$ an admissible
linear form. By conditions (a) and (b) of a Weil representation,
$\omega(g,h)\rho(tu)=\rho(t\,guh^{\mathsf T})\omega(g,h)$ on Schwartz-Bruhat
functions, and $\lambda\circ\omega(g,h)$ is admissible. So
\[
\lambda\bigl(\rho(t\,guh^{\mathsf T})(T(f)\otimes\phi)\bigr)
=(\lambda\circ\omega(g,h))\bigl(\rho(tu)(f\otimes\phi)\bigr),
\]
which by Lemma~\ref{asm-theta-agree-real-derivative} has derivative
$(\lambda\circ\omega(g,h))(D_uf\otimes\phi)=\lambda(T(D_uf)\otimes\phi)$ at
$t=0$. So $t\mapsto\rho(t\,guh^{\mathsf T})$ acts weakly at $\infty$ to first
order by $T(f)\mapsto T(D_uf)$. As $guh^{\mathsf T}$ is at $\infty$, it also
does so by $T(f)\mapsto D_{guh^{\mathsf T}}(T(f))$, and
Lemma~\ref{asm-theta-agree-real-weak}(3) gives the identity.
\end{proof}
% Lean: intertw uses actsLaws (fifth clause), hs.impl and hs.adm; (1) is hD1 in cartanMove_of.

\begin{lemma}\label{asm-theta-agree-rechoice-fock}
Let $(X,W,M,\chi_X,\chi_W,\omega)$ and $(X^\circ,W^\circ,M^\circ,\chi_X,\chi_W,\omega)$
be Weil pairs for a theta base $B$ of a CM field $E$, such that $X^\circ$ and
$W^\circ$ have the matrices and lattices of $X$ and $W$ and $M^\circ$ has the
isomorphism of $M$. Let $g\in U_X$ move $X^\circ$ to $X$ and $h\in U_W$ move
$W^\circ$ to $W$ (Definition~\ref{asm-theta-agree-rechoice-terms}(4)), and let
$T$ be an archimedean operator of $\omega$ at $(g,h)$. Then
$T(\operatorname{Pol}(M^\circ))\subset\operatorname{Pol}(M)$. In particular, if
$g\in K_\infty(X)$ and $h\in K_\infty(W)$, then $T$ maps
$\operatorname{Pol}(M)$ into itself.
\end{lemma}
\leanref{Development/WeilBound/ThetaLifts/CartanFactorisation/Cartan.lean}{PositivityTheory.ThKF.Cart.pol_le}
\leanref{Development/WeilBound/ThetaLifts/CartanFactorisation/ArchimedeanEvaluationVector.lean}{PositivityTheory.ThKF.A.map_Pol}
\begin{proof}
The two models have the same isomorphism, so the same Schrödinger operators
and the same operators $D_u$; by
Lemma~\ref{asm-theta-agree-rechoice-intertwine}(1), for both the action at
$\infty$ to first order of $t\mapsto\rho(tu)$ is $f\mapsto D_uf$. Let $N$ be the
set of $f\in\mathcal S(F_\infty^I)$ with $T(f)\in\operatorname{Pol}(M)$, a
subspace. It suffices to show that $N$ contains the Gaussians of $M^\circ$ and
contains $D_uf$ whenever it contains $f$ and $u$ is at $\infty$.

From $\pi_w(g)\theta_{X^\circ,w}=\theta_{X,w}\pi_w(g)$ we get
$\pi_w(g^{-1})\theta_{X,w}=\theta_{X^\circ,w}\pi_w(g^{-1})$, and likewise for
$h$. Let $f_0$ be a Gaussian of $M^\circ$, with signs $\varepsilon_w$. Let $u$
and $u'$ be at $\infty$ with
$\pi_w(u')=\varepsilon_wi\,\theta_{X,w}\pi_w(u)\theta_{W,w}^{\mathsf T}$ for
every $w$, and put $u_0=g^{-1}u(h^{-1})^{\mathsf T}$ and
$u'_0=g^{-1}u'(h^{-1})^{\mathsf T}$, which are at $\infty$. Then
$\pi_w(u'_0)=\varepsilon_wi\,\theta_{X^\circ,w}\pi_w(u_0)\theta_{W^\circ,w}^{\mathsf T}$,
so $D_{u'_0}f_0=iD_{u_0}f_0$. By
Lemma~\ref{asm-theta-agree-rechoice-intertwine}(2),
$D_u(T(f_0))=T(D_{u_0}f_0)$ and $D_{u'}(T(f_0))=T(D_{u'_0}f_0)=iD_u(T(f_0))$. So
$T(f_0)$ is a Gaussian of $M$, with the same signs, and $f_0\in N$.

If $f\in N$ and $u$ is at $\infty$, then
$T(D_uf)=D_{guh^{\mathsf T}}(T(f))$ by
Lemma~\ref{asm-theta-agree-rechoice-intertwine}(2), and this lies in
$\operatorname{Pol}(M)$, since $T(f)$ does and $guh^{\mathsf T}$ is at
$\infty$. So $D_uf\in N$.

If $g\in K_\infty(X)$ and $h\in K_\infty(W)$, then $g$ moves $X$ to $X$ and
$h$ moves $W$ to $W$, and the first part applies with $M^\circ=M$.
\end{proof}
% Lean: the first part is pol_le applied in cartanMove_of with a = mat g^-1, b = mat h^-1 and
% the identity of intertw. The second part is the clause of ArchOp for K_inf, which archOp
% proves as map_Pol with the same identity obtained from T_D (Impl.dRho).

\begin{lemma}\label{asm-theta-agree-rechoice-lie}
Let $(X,W,M,\chi_X,\chi_W,\omega)$ be a Weil pair for a theta base $B$ of a CM
field $E$, with $X$ and $W$ of ranks $n$ and $d$ and
$I=\{1,\ldots,n\}\times\{1,\ldots,d\}$. Let $\gamma$ and $\delta$ be lines in
$U_X$ and $U_W$, for each $t\in\mathbf R$ let $T_t$ be an archimedean operator
of $\omega$ at $(\gamma(t),\delta(t))$, and let $f\in\operatorname{Pol}(M)$. Then
there is $f_1\in\operatorname{Pol}(M)$ such that
$t\mapsto\omega(\gamma(t),\delta(t))$ acts at $\infty$ to first order by
$f\mapsto f_1$, and $t\mapsto L(T_tf)$ has derivative $L(f_1)$ at $t=0$ for
every continuous linear form $L$ on $\mathcal S(F_\infty^I)$.
\end{lemma}
\leanref{Development/WeilBound/ThetaLifts/CartanFactorisation/CartanConj.lean}{PositivityTheory.ThKF.ArchLie}
\leanref{Development/WeilBound/ThetaLifts/Reductions/ArchLie.lean}{PositivityTheory.ThKR.archLie'}
\leanref{Development/WeilBound/ThetaLifts/Reductions/ExistsSmoothFun.lean}{PositivityTheory.ThKR.lie'}
\leanref{Development/WeilBound/ThetaLifts/Reductions/ExistsSmoothFun.lean}{PositivityTheory.ThKR.exists_sp}
\leanref{Development/WeilBound/ThetaLifts/Reductions/ExistsSmoothFun.lean}{PositivityTheory.ThKR.exists_smooth_fun}
\leanref{Development/WeilBound/ThetaLifts/CartanFactorisation/EqOfInner.lean}{PositivityTheory.ThKF.E.Impl.scalar}
\leanref{Development/WeilBound/ThetaLifts/Reductions/MetaDef.lean}{PositivityTheory.ThKR.LocalLift}
\leanref{Development/WeilBound/ThetaLifts/Reductions/LocalLift.lean}{PositivityTheory.ThKR.localLift_holds}
\begin{proof}
Put $x_t=(\gamma(t),\delta(t))$, $V=F_\infty^I$, and let $e$, $\rho_\infty$
and $D_a$, for $a\in V\times V$, be as in the proof of
Lemma~\ref{asm-theta-agree-rechoice-operators}. Let $s_t$ be the
$\mathbf R$-linear automorphism of $V\times V$ with
$s_t(e(u)_\infty)=e(\gamma(t)u\delta(t)^{\mathsf T})_\infty$ for $u$ at
$\infty$.

\emph{Step 1.} The family $(s_t)$ is a continuous one-parameter group of
automorphisms of $V\times V$ preserving the symplectic form, since $e$ is
symplectic and $u\mapsto\gamma(t)u\delta(t)^{\mathsf T}$ preserves the form of
$\mathbb W(X,W)$. Let $g_0$ be a Schwartz-Bruhat function on $\mathbf A_f^I$
with $g_0(0)=1$. An archimedean operator $T$ at a point $x$ is determined by
$T(f)\otimes g_0=\omega(x)(f\otimes g_0)$, so $T_t$ is the operator of Step 2
of the proof of Lemma~\ref{asm-theta-agree-rechoice-operators} for $x_t$. By
Step 3 there, $T_t$ preserves $L^2$-norms, is onto, and has
$T_t\rho_\infty(a)=\rho_\infty(s_ta)T_t$. We have $T_0=1$ and
$T_{t+t'}=T_tT_{t'}$, since $\omega$ is a representation and archimedean
operators are determined in this way. For a continuous linear form $L$ there
is an admissible linear form $\lambda$ with $\lambda(f'\otimes g_0)=L(f')$, so
$t\mapsto L(T_tf)=\lambda(\omega(x_t)(f\otimes g_0))$ is continuous by
condition (c) of a Weil representation.

\emph{Step 2.} A continuous one-parameter group of linear maps is an
exponential, so $s_t=\exp(tY)$ for a matrix $Y$; differentiating the
invariance of the symplectic form at $t=0$ shows that $Y$ lies in its
symplectic Lie algebra.

\emph{Step 3.} Choose linear coordinates $V\cong\mathbf R^l$ and
$V\times V\cong\mathbf R^{2l}$ under which $\rho_\infty$ and the $D_b$ become the
operators $\rho_S$ and $d\rho$ of Lemma~\ref{asm-theta-agree-real-metaplectic}
and the symplectic form becomes $\operatorname{sympl}$. By part (2) of that
lemma, the metaplectic representation near the identity, there are
$\varepsilon>0$, $a\in\mathbf C$ and continuous invertible operators $U_t$ on
$\mathcal S(V)$ for $|t|<\varepsilon$, preserving $L^2$-norms, with $U_0=1$,
$U_t\rho_\infty(b)=\rho_\infty(\exp(tY)b)U_t$, and $t\mapsto L(U_tf)$
differentiable on $(-\varepsilon,\varepsilon)$ with derivative
$L(q_Yf+af)$ at $0$ for every continuous linear form $L$, where
$q_Y=\operatorname{quad}(Y)$ lies in the algebra generated by the operators
$D_b$.

\emph{Step 4.} For $|t|<\varepsilon$, the operator $T_{-t}U_t$ preserves
$L^2$-norms, is onto and commutes with every $\rho_\infty(b)$. By the theorem of
Stone and von Neumann \cite[Chapitre~2]{MVW87} it is a scalar. So
$T_t=c(t)U_t$ with $c(t)\in\mathbf C$,
and $c(0)=1$.

\emph{Step 5.} $c$ is differentiable at $0$. Let $\varphi\neq0$ in
$\mathcal S(V)$. As $t\mapsto\langle\varphi,T_t\varphi\rangle$ is continuous
and equal to $\|\varphi\|_2^2\neq0$ at $0$, there is $\delta_0$ with
$\mu(\varphi)\neq0$ for the continuous linear form
$\mu(f')=\int_0^{\delta_0}\langle\varphi,T_sf'\rangle\,ds$. Then
$\mu(T_t\varphi)=\int_t^{t+\delta_0}\langle\varphi,T_s\varphi\rangle\,ds$ is
differentiable in $t$, and $c(t)=\mu(T_t\varphi)/\mu(U_t\varphi)$ near $0$.

\emph{Step 6.} By Steps 3 to 5, $t\mapsto L(T_tf)$ has derivative $L(Qf)$ at
$0$ for every continuous linear form $L$, where $Q=c'(0)+a+q_Y$ lies in the
algebra generated by the $D_b$. Every $b\in V\times V$ is $e(u)_\infty$ for
some $u$ at $\infty$, and $D_u$ maps $\operatorname{Pol}(M)$ into itself by the
definition of $\operatorname{Pol}(M)$ and
Lemma~\ref{asm-theta-agree-rechoice-intertwine}(1). So
$f_1=Qf\in\operatorname{Pol}(M)$. For a Schwartz-Bruhat $\phi$ on
$\mathbf A_f^I$ and an admissible $\lambda$,
$\lambda(\omega(x_t)(f\otimes\phi))=\lambda(T_tf\otimes\phi)$, and
$f'\mapsto\lambda(f'\otimes\phi)$ is a continuous linear form. So
$t\mapsto\omega(x_t)$ acts weakly at $\infty$ to first order by $f\mapsto f_1$,
and by Lemma~\ref{asm-theta-agree-real-weak}(1) it acts at $\infty$ to first
order by $f\mapsto f_1$.
\end{proof}
% Lean: archLie' localLift_holds, through lie' (exists_sp, Impl.scalar, exists_smooth_fun),
% impl_of, exists_adm, weylAlg_induction with D_mem_Pol, and actsD_of. Step 3 is LocalLift,
% which the Lean proves as localLift_holds; the closing of Theorem K passes archLie'
% localLift_holds for ArchLie. Its other proof, archLie_of, assumes
% Published.RealMetaplecticRepresentation.

\begin{lemma}\label{asm-theta-agree-rechoice-count}
Let $(X,W,M,\chi_X,\chi_W,\omega)$ be a Weil pair for a theta base $B$ of a CM
field $E$. Then $\operatorname{Pol}(M)$ has countable dimension.
\end{lemma}
\leanref{Development/WeilBound/ThetaLifts/CartanFactorisation/PolCountDef.lean}{PositivityTheory.ThKF.PolCount}
\leanref{Development/WeilBound/ThetaLifts/CartanFactorisation/PolCountFinal.lean}{PositivityTheory.ThKF.polCount_of_derEx}
\leanref{Development/WeilBound/ThetaLifts/CartanFactorisation/PolCount.lean}{PositivityTheory.ThKF.PC.polCount_of}
\begin{proof}
Let $(r,s;r',s')$ be the signature of the pair. For every infinite place $w$,
Lemma~\ref{asm-theta-agree-real-frames} gives a frame $A_w$ of $\tau_w(J_X)$ of
type $(r,s)$ with $\theta_{X,w}A_w=A_wJ_{r,s}$ and a frame $C_w$ of
$i\tau_w(J_W)$ of type $(r',s')$ with $\theta_{W,w}C_w=C_wJ_{r',s'}$. By
Lemma~\ref{asm-theta-agree-real-fock-map}, with $a_0=1$ and $b_0=i$, there is
a Fock map $\Phi$ for $(M,(A_w),(C_w))$, and its image is
$\operatorname{Pol}(M)$. Its source $\bigotimes_w\mathcal P(r,s;r',s')$ is a
finite tensor product of polynomial rings in finitely many variables, which has
a countable basis of monomials.
\end{proof}
% Lean: polCount_of_derEx is PC.polCount_of (RL.fockEx derEx) (RL.vacLine derEx); exists_fock
% takes the frames from RL.frameEx, pol_le gives Pol(M) inside the range of the Fock map
% (VacLine), and rank_tens bounds the rank of the tensor product of polynomial rings.

\begin{lemma}\label{asm-theta-agree-rechoice-linear}
Let $(X,W,M,\chi_X,\chi_W,\omega)$ be a Weil pair for a theta base $B$ of a CM
field $E$, and let $f\in\operatorname{Pol}(M)$.
\begin{enumerate}
\item If $\gamma$ and $\delta$ are lines in $U_X$ and $U_W$, there is exactly
  one $f_1$ such that $t\mapsto\omega(\gamma(t),\delta(t))$ acts at $\infty$ to
  first order by $f\mapsto f_1$, and $f_1\in\operatorname{Pol}(M)$.
\item If $k\in K_\infty(X)$ and $k'\in K_\infty(W)$, there is exactly one $f_1$
  such that $\omega(k,k')$ acts at $\infty$ by $f\mapsto f_1$, and
  $f_1\in\operatorname{Pol}(M)$.
\item For a line $\delta$ in $U_W$, the map $f\mapsto f_1$ of (1) with
  $\gamma=1$ is linear on $\operatorname{Pol}(M)$; for $k'\in K_\infty(W)$, the
  map $f\mapsto f_1$ of (2) with $k=1$ is linear on $\operatorname{Pol}(M)$.
\end{enumerate}
\end{lemma}
\leanref{Development/WeilBound/ThetaLifts/CartanFactorisation/RinfZero.lean}{PositivityTheory.ThKF.pol_lie}
\leanref{Development/WeilBound/ThetaLifts/CartanFactorisation/RinfZero.lean}{PositivityTheory.ThKF.pol_k}
\leanref{Development/WeilBound/ThetaLifts/CartanFactorisation/RinfZero.lean}{PositivityTheory.ThKF.rinf_ex}
\leanref{Development/WeilBound/ThetaLifts/CartanFactorisation/CommuteConj.lean}{PositivityTheory.ThKF.exists_linop}
\leanref{Development/WeilBound/ThetaLifts/CartanFactorisation/CartanConj.lean}{PositivityTheory.ThKF.ArchUniq}
\leanref{Development/WeilBound/ThetaLifts/CartanFactorisation/ArchimedeanEvaluationVector.lean}{PositivityTheory.ThKF.A.archUniq}
\begin{proof}
Uniqueness in (1) and (2) is Lemma~\ref{asm-theta-agree-real-weak}(3).

(1) Choose archimedean operators $T_t$ of $\omega$ at
$(\gamma(t),\delta(t))$ by Lemma~\ref{asm-theta-agree-rechoice-operators}(1).
Lemma~\ref{asm-theta-agree-rechoice-lie} gives $f_1$.

(2) Let $T$ be an archimedean operator of $\omega$ at $(k,k')$. Then $\omega(k,k')$
acts at $\infty$ by $f\mapsto T(f)$, and $T(f)\in\operatorname{Pol}(M)$ by
Lemma~\ref{asm-theta-agree-rechoice-fock}.

(3) For $k'$, the map is the restriction of an archimedean operator at
$(1,k')$, which is linear. For $\delta$, choose archimedean operators $T_t$ at
$(1,\delta(t))$, and write $D(f)$ for the $f_1$ of (1). By
Lemma~\ref{asm-theta-agree-rechoice-lie}, $L(D(f))$ is the derivative at $0$ of
$t\mapsto L(T_tf)$, for every continuous linear form $L$. The derivative is
linear in $f$, as each $T_t$ is. So $L(D(f+f'))=L(D(f)+D(f'))$ and
$L(D(cf))=L(cD(f))$ for every $L$. The evaluations at points are continuous
linear forms on $\mathcal S(F_\infty^I)$, so $D(f+f')=D(f)+D(f')$ and
$D(cf)=cD(f)$.
\end{proof}

\begin{lemma}\label{asm-theta-agree-rechoice-howe}
Let $(B,G,\mathcal X,\omega,\omega')$ be a triple for a theta base $B$ of a CM
field $E$ (Definition~\ref{asm-theta-agree-triple}), and let
$(Y,W,M_0,\chi_Y,\chi_W,\omega_0)$ be one of its two Weil pairs, with $Y$ and
$W$ of ranks $n$ and $d$ and $I=\{1,\ldots,n\}\times\{1,\ldots,d\}$.
\begin{enumerate}
\item Any two maximal elements of the set of theta subspaces of
  $\operatorname{Pol}(M_0)$ other than $\operatorname{Pol}(M_0)$
  (Definition~\ref{asm-theta-agree-rechoice-terms}(2)) are equal.
\item Let $v$ be a finite place. Any two maximal elements of the set of theta
  subspaces at $v$ of $\omega_0$ other than $\mathcal S(F_v^I)$
  (Definition~\ref{asm-theta-agree-rechoice-terms}(3)) are equal.
\end{enumerate}
\end{lemma}
\leanref{Development/WeilBound/ThetaLifts/Reductions/JoinDef.lean}{PositivityTheory.ThKR.JoinTri}
\leanref{Development/WeilBound/ThetaLifts/Reductions/JoinDef.lean}{PositivityTheory.ThKR.JoinAt}
\leanref{Development/WeilBound/ThetaLifts/Reductions/JoinTriOf.lean}{PositivityTheory.ThKR.Join.joinTri_of}
\leanref{Development/WeilBound/ThetaLifts/Reductions/ArchimedeanMatchedMapsCompare.lean}{PositivityTheory.ThKF.inf_uniq2}
\leanref{Development/WeilBound/ThetaLifts/CartanFactorisation/RealTrans.lean}{PositivityTheory.ThKF.max_uniq}
\leanref{Development/WeilBound/CoreStatements/TemperedRepresentationTheory/HoweMaximalTori/MaxSplit.lean}{PositivityTheory.HoweDualityUnitaryAlt.MaxFinT}
\leanref{Development/WeilBound/CoreStatements/TemperedRepresentationTheory/HoweMaximalTori/MaxSplit.lean}{PositivityTheory.HoweDualityUnitaryAlt.MaxAt}
\leanref{Development/WeilBound/ThetaLifts/SphericalLattices/ThetaSubSup.lean}{PositivityTheory.Sph.maxFinT_of}
\leanref{Development/WeilBound/ThetaLifts/SphericalLattices/ThetaSubSup.lean}{PositivityTheory.Sph.maxAt_of}
\leanref{Development/WeilBound/ThetaLifts/SphericalLattices/ThetaSubSup.lean}{PositivityTheory.Sph.Stab.sup_ne}
\begin{proof}
(1) Let $N$ and $N'$ be such maximal elements. By
Theorem~\ref{lit-howe-duality-real}(2), for the pair $(Y,W)$ with the character
$\psi$ of $B$, the characters $\chi_Y$ and $\chi_W$, the representation
$\omega_0$ and the Cartan involutions $\theta_{Y,w}$ and $\theta_{W,w}$, whose
Fock space is $\operatorname{Pol}(M_0)$, we have
$N+N'\neq\operatorname{Pol}(M_0)$. The sum $N+N'$ is a theta subspace. It lies
in $\operatorname{Pol}(M_0)$ and contains the elements required by the first
two conditions, as $N$ does. For the third, let $f\mapsto f_1$ be the action at
$\infty$ to first order of $t\mapsto\omega_0(1,\delta(t))$ for a line $\delta$
in $U_W$, or the action at $\infty$ of $\omega_0(1,k)$ for $k\in K_\infty(W)$. By
Lemma~\ref{asm-theta-agree-rechoice-linear}(3) there is a linear map $A$ on
$\operatorname{Pol}(M_0)$ such that this action sends $f$ to $A(f)$, and to no
other function. For $f=u+u'$ with $u\in N$ and $u'\in N'$, this gives
$f_1=A(u)+A(u')$ with $A(u)\in N$ and $A(u')\in N'$. So $N+N'$ is a theta
subspace other than $\operatorname{Pol}(M_0)$ containing $N$ and $N'$, and
$N=N+N'=N'$ by maximality.

(2) Let $(\omega_v,\rho_v)$ be local operators at $v$ for the pair
(Definition~\ref{asm-theta-agree-finite-local-objects}(4)), which exist by
Lemma~\ref{asm-theta-agree-finite-inputs}(2). For $h\in U_W^{(v)}$ and a
Schwartz-Bruhat $\phi$, $\omega_0(1,h)$ acts at $v$ by
$\phi\mapsto\omega_v(1,h)\phi$, and by no other map
(Lemma~\ref{asm-theta-agree-finite-unique}). So a subspace $N$ with
$\mathfrak a_v(\omega_0)\subset N\subset\mathcal S(F_v^I)$ is a theta subspace
at $v$ if and only if it is stable under every $\omega_v(1,h)$. Since $G$ is
proper, the models are regular, and
Theorem~\ref{th-howe-finite} applies, with $\varpi_v(1,h)=\omega_v(1,h)$ and
$\mathcal R_v=\mathfrak a_v(\omega_0)$.

Let $N$ and $N'$ be maximal theta subspaces at $v$ other than
$\mathcal S(F_v^I)$. Then $N+N'$ is a theta subspace at $v$, and it suffices to
show $N+N'\neq\mathcal S(F_v^I)$. In the first alternative of
Theorem~\ref{th-howe-finite}, $N=N'=\mathfrak a_v(\omega_0)$. In the second,
let $s$ be as there. No theta subspace $P\neq\mathcal S(F_v^I)$ contains $s$,
since it would contain $\mathfrak a_v(\omega_0)$ and every $\omega_v(1,h)s$,
which span $\mathcal S(F_v^I)$. Suppose $s=x+y$ with $x\in N$ and $y\in N'$.
For $\phi=x$, take $h_0,\ldots,h_m$, $w_0,\ldots,w_m$ and $c$ as in the
theorem, and put $x'=\sum_iw_i\omega_v(1,h_i)x\in N$ and
$y'=\sum_iw_i\omega_v(1,h_i)y\in N'$. Then $x'+y'=(m+1)s$ and
$x'-cs\in\mathfrak a_v(\omega_0)\subset N\cap N'$. So $cs\in N$ and
$(m+1-c)s=y'+(x'-cs)\in N'$. If $c\neq0$ then $s\in N$, and if $c=0$ then
$s\in N'$; both are impossible. So $s\notin N+N'$.
\end{proof}
% Lean: (1) is inf_uniq2 with JoinAt from JoinTri, which the closing of Theorem K proves as
% Join.joinTri_of from FockEx, PolRange, LieFirst and LieSecond; JoinAt is part (2) of
% Theorem lit-howe-duality-real for the pair. The older proof polRe_of uses
% Published.HoweDualityReal through inf_uniq. (2) is MaxFinT = Sph.maxFinT_of SphTri, through
% locRep_tri, maxAt_of, thetaSub_sup and Stab.sup_ne.

\begin{lemma}\label{asm-theta-agree-rechoice-kernels}
Let $(B,G,\mathcal X,\omega,\omega')$ be a triple of ranks $(a,a')$ for a theta
base $B$ of a CM field $E$ (Definition~\ref{asm-theta-agree-triple}), with
models $M$ and $M'$, and put $I_0=\{1,\ldots,a\}\times\{1,\ldots,a+a'\}$ and
$I'_0=\{1,\ldots,a'\}\times\{1,\ldots,a+a'\}$.
\begin{enumerate}
\item Let $(\Pi,p,p')$ be real maps matched to the triple
  (Definition~\ref{asm-theta-agree-matched}(2)). Then
  $\ker p\cap\operatorname{Pol}(M)$ is a maximal element of the set of theta
  subspaces of $\operatorname{Pol}(M)$ other than $\operatorname{Pol}(M)$, for
  the Weil pair of $\omega$; likewise for $\ker p'\cap\operatorname{Pol}(M')$
  and $\omega'$.
\item Let $(\Pi_v,p_v,p'_v)_v$ be finite maps matched to the triple
  (Definition~\ref{asm-theta-agree-matched}(3)), and let $v$ be a finite place.
  Then $\ker p_v\cap\mathcal S(F_v^{I_0})$ is a maximal element of the set of
  theta subspaces at $v$ of $\omega$ other than $\mathcal S(F_v^{I_0})$;
  likewise for $p'_v$ and $\omega'$.
\end{enumerate}
\end{lemma}
\leanref{Development/WeilBound/ThetaLifts/CartanFactorisation/RinfZero.lean}{PositivityTheory.ThKF.maxInfP}
\leanref{Development/WeilBound/ThetaLifts/CartanFactorisation/RinfZero.lean}{PositivityTheory.ThKF.maxFinP}
\leanref{Development/WeilBound/ThetaLifts/CartanFactorisation/RinfZero.lean}{PositivityTheory.ThKF.absMax}
\leanref{Development/WeilBound/ThetaLifts/CartanFactorisation/RinfZero.lean}{PositivityTheory.ThKF.max_inf}
\leanref{Development/WeilBound/ThetaLifts/CartanFactorisation/RinfZero.lean}{PositivityTheory.ThKF.max_fin}
\leanref{Development/WeilBound/ThetaLifts/CartanFactorisation/ExistsBeta.lean}{PositivityTheory.ThKF.maximal_of}
\begin{proof}
(1) Let $\iota$ be the set of lines in $U_W$ together with the elements of
$K_\infty(W)$. For $i\in\iota$ and $f\in\operatorname{Pol}(M)$ let $A_if$ be the
image of $f$ under the action at $\infty$ to first order of
$t\mapsto\omega(1,\delta(t))$ if $i=\delta$ is a line, and under the action at
$\infty$ of $\omega(1,k)$ if $i=k$; by
Lemma~\ref{asm-theta-agree-rechoice-linear} it exists, is unique, lies in
$\operatorname{Pol}(M)$, and $A_i0=0$. Define $A'_i$ on $\operatorname{Pol}(M')$
in the same way with $\omega'$. Condition (b) of matched real maps says that
$p(f)=p'(f')$ implies $p(A_if)=p'(A'_if')$.

\emph{Step 1.} If $f\in\operatorname{Pol}(M)$ and $p(f)=0$, then $p(A_if)=0$:
apply (b) with $f'=0$. Likewise for $p'$.

\emph{Step 2.} If $N\subset\operatorname{Pol}(M)$ is stable under every $A_i$,
then $p(N)$ is $0$ or $\Pi$. Indeed, let $f\in\operatorname{Pol}(M)$ with
$p(f)=p(n)$ for some $n\in N$. As $p'$ maps $\operatorname{Pol}(M')$ onto $\Pi$,
there is $f'$ with $p'(f')=p(f)$, and (b) gives
$p(A_if)=p'(A'_if')=p(A_in)\in p(N)$. By condition (c), $p(N)=0$ or
$p(N)=\Pi$. If $N\subset\operatorname{Pol}(M')$ is stable under every $A'_i$,
then $p'(N)$ is $0$ or $\Pi$ in the same way: if $p(f)=p'(n)$ with $n\in N$,
then $p(A_if)=p'(A'_in)\in p'(N)$.

\emph{Step 3.} Put $K=\ker p\cap\operatorname{Pol}(M)$. It is a theta subspace.
If $\gamma$ is a line in $U_V$ and $t\mapsto\omega(\gamma(t),1)$ acts at
$\infty$ to first order by $f\mapsto f_1$ with $f\in\operatorname{Pol}(M)$, then
$f_1\in\operatorname{Pol}(M)$ by Lemma~\ref{asm-theta-agree-rechoice-linear}(1)
and $p(f_1)=0$ by condition (a); likewise $f_1-f\in K$ for the action of
$\omega(k,1)$, $k\in K_\infty(V)$. By Step 1, $K$ is stable under the $A_i$. As
$p(\operatorname{Pol}(M))=\Pi\neq0$, $K\neq\operatorname{Pol}(M)$. Let $N\supset K$
be a theta subspace other than $\operatorname{Pol}(M)$. It is stable under the
$A_i$, so by Step 2 $p(N)=0$ or $p(N)=\Pi$. In the second case, every
$f\in\operatorname{Pol}(M)$ has $p(f)=p(n)$ for some $n\in N$, so
$f-n\in K\subset N$ and $f\in N$, which is excluded. So $p(N)=0$ and $N=K$. The
same argument applies to $p'$.

(2) Let $\iota=U_W^{(v)}$. By Lemma~\ref{asm-theta-agree-finite-inputs}(2) and
Lemma~\ref{asm-theta-agree-finite-unique}, for $h\in\iota$ and a Schwartz-Bruhat
$\phi$ there is exactly one function $A_h\phi$ such that $\omega(1,h)$ acts at
$v$ by $\phi\mapsto A_h\phi$, and it is Schwartz-Bruhat; define $A'_h$ with
$\omega'$ in the same way. The conditions on matched finite maps at $v$ give the
analogues of conditions (a) to (c): $p_v$ and $p'_v$ map the Schwartz-Bruhat
functions onto $\Pi_v\neq0$ and vanish on $\mathfrak a_v(\omega)$ and
$\mathfrak a_v(\omega')$, $p_v(\phi)=p'_v(\phi')$ implies
$p_v(A_h\phi)=p'_v(A'_h\phi')$, and the condition on subspaces of $\Pi_v$ holds.
Steps 1 and 2 then hold as in (1). The space
$K_v=\ker p_v\cap\mathcal S(F_v^{I_0})$ contains $\mathfrak a_v(\omega)$ and is
stable under the $A_h$ by Step 1, so it is a theta subspace at $v$, and the
argument of Step 3 shows that it is maximal. Likewise for $p'_v$.
\end{proof}

\begin{lemma}\label{asm-theta-agree-rechoice-compare}
Let $B$ be a theta base of a CM field $E$.
\begin{enumerate}
\item Let $(B,G,\mathcal X,\omega,\omega')$ be a triple
  (Definition~\ref{asm-theta-agree-triple}), with models $M$ and $M'$, and let
  $(\Pi,p,p')$ and $(\Pi_2,p_2,p'_2)$ be real maps matched to it
  (Definition~\ref{asm-theta-agree-matched}(2)). Then there are a linear map
  $e\colon\Pi\to\Pi_2$ and $c\in\mathbf C$ with $c\neq0$, $e(p(f))=p_2(f)$ for
  $f\in\operatorname{Pol}(M)$, and $e(p'(f'))=c\,p'_2(f')$ for
  $f'\in\operatorname{Pol}(M')$.
\item Let $(B,G,\mathcal X,\omega,\omega',R,R_f)$ and
  $(B,G^\circ,\mathcal X,\omega,\omega',R_2,R_{f,2})$ be matched quadruples, where
  $G^\circ=G[V^\circ,V'^\circ,W^\circ]$ for some lattice spaces $V^\circ$,
  $V'^\circ$ and $W^\circ$ (Definition~\ref{asm-theta-agree-triple}(2)), with
  finite maps $R_f=(\Pi_v,p_v,p'_v)_v$ and $R_{f,2}=(\Pi_{2,v},p_{2,v},p'_{2,v})_v$.
  Then for every finite place $v$ there are a linear map
  $e_v\colon\Pi_v\to\Pi_{2,v}$ and $c_v\in\mathbf C$ with $c_v\neq0$,
  $e_v(p_v(\phi))=p_{2,v}(\phi)$ and $e_v(p'_v(\phi'))=c_v\,p'_{2,v}(\phi')$ for
  Schwartz-Bruhat $\phi$ and $\phi'$. If $v\notin T$, then $c_v=1$.
\end{enumerate}
\end{lemma}
\leanref{Development/WeilBound/ThetaLifts/Reductions/ArchimedeanMatchedMapsCompare.lean}{PositivityTheory.ThKF.archimedean_matched_maps_compare}
\leanref{Development/AutomorphicSpectrum/Proofs/BlockDecomposition/PolReOf2T.lean}{GlobalTheory.CompanionTheory.ThetaK.cmp_finT}
\leanref{Development/WeilBound/ThetaLifts/CartanFactorisation/CountableCompactOpen.lean}{PositivityTheory.ThKF.rank_SB}
\begin{proof}
(1) Apply Lemma~\ref{asm-theta-agree-rechoice-schur} with
$S=\operatorname{Pol}(M)$, $S'=\operatorname{Pol}(M')$, the maps $p$, $p'$,
$p_2$ and $p'_2$, and the set $\iota$ and the maps $A_i$ and $A'_i$ of the proof
of Lemma~\ref{asm-theta-agree-rechoice-kernels}(1); the $A_i$ are linear by
Lemma~\ref{asm-theta-agree-rechoice-linear}(3). By condition (a) of matched real
maps, the four maps are onto, and $\Pi_2\neq0$ by condition (c). For hypothesis
(a), Lemma~\ref{asm-theta-agree-rechoice-kernels}(1) applied to both real maps
shows that $\ker p\cap\operatorname{Pol}(M)$ and
$\ker p_2\cap\operatorname{Pol}(M)$ are maximal elements of the same set, so
they are equal by Lemma~\ref{asm-theta-agree-rechoice-howe}(1); likewise
$\ker p'\cap\operatorname{Pol}(M')=\ker p'_2\cap\operatorname{Pol}(M')$.
Hypothesis (b) is condition (b) for both real maps, (c) is condition (c) for
$(\Pi,p,p')$, and (d) is Lemma~\ref{asm-theta-agree-rechoice-count}.

(2) The two triples have the same representations $\omega$ and $\omega'$, and
their models have the same isomorphisms and vacua. So the augmentation spaces,
the theta subspaces at $v$ and the actions at $v$ are the same for both, and
$R_{f,2}$ satisfies every condition of finite maps matched to
$(B,G,\mathcal X,\omega,\omega')$ at $v$ other than the zeta condition, which
the following does not use. Apply Lemma~\ref{asm-theta-agree-rechoice-schur}
with $S=\mathcal S(F_v^{I_0})$, $S'=\mathcal S(F_v^{I'_0})$, the maps $p_v$,
$p'_v$, $p_{2,v}$ and $p'_{2,v}$, $\iota=U_W^{(v)}$, $A_h=\omega_v(1,h)$ for local
operators $(\omega_v,\rho_v)$ at $v$ for the pair of $\omega$, and $A'_h$ as in
the proof of Lemma~\ref{asm-theta-agree-rechoice-kernels}(2). Hypothesis (a)
follows from Lemma~\ref{asm-theta-agree-rechoice-kernels}(2), applied to both
finite maps, and Lemma~\ref{asm-theta-agree-rechoice-howe}(2). Hypothesis (b) is
the condition on $U_W^{(v)}$ of matched finite maps for both, and (c) is the
condition on subspaces of $\Pi_v$. For (d), the Schwartz-Bruhat functions on
$F_v^{I_0}$ are spanned by the indicator functions of compact open sets. Every
compact open set is a finite union of members of a countable basis of the
topology, so there are countably many of them.

Let $v\notin T$. Both finite maps are matched, and the vacua are the same, so
$p_v(\mathrm{vac}_v)=p'_v(\mathrm{vac}'_v)$ and
$p_{2,v}(\mathrm{vac}_v)=p'_{2,v}(\mathrm{vac}'_v)\neq0$. Then
$p_{2,v}(\mathrm{vac}_v)=e_v(p_v(\mathrm{vac}_v))=e_v(p'_v(\mathrm{vac}'_v))
=c_v\,p'_{2,v}(\mathrm{vac}'_v)=c_v\,p_{2,v}(\mathrm{vac}_v)$, so $c_v=1$.
\end{proof}
% Lean: (1) is archimedean_matched_maps_compare, with exists_linop for the A_i, inf_uniq2 and
% maxInfP for (a), and PolCount for (d). (2) is cmp_finT with locRepEx, ex_fin, maxFinP and
% MaxFinT for (a), and rank_SB (countable_compactOpen); c_v = 1 is cf1 in polRe_of2T.

\begin{lemma}\label{asm-theta-agree-rechoice-cartan}
Let $B$ be a theta base of a CM field $E$, and let $Y$ and $Y^\circ$ be lattice
spaces of rank $m$ with the same matrix $J$ that are both Hermitian, or both
skew-Hermitian, of the same signature $(r,s)$
(Definition~\ref{asm-theta-agree-doubling}(1)). Then there is $g\in U_Y$ at
$\infty$ that moves $Y^\circ$ to $Y$
(Definition~\ref{asm-theta-agree-rechoice-terms}(4)).
\end{lemma}
\leanref{Development/WeilBound/ThetaLifts/CartanFactorisation/CartanConj.lean}{PositivityTheory.ThKF.CartanConj}
\leanref{Development/WeilBound/ThetaLifts/CartanFactorisation/Cartan.lean}{PositivityTheory.ThKF.cartan_conjugacy_holds}
\leanref{Development/WeilBound/ThetaLifts/CartanFactorisation/CartanA.lean}{PositivityTheory.ThKF.Cart.exists_conj}
\leanref{Development/WeilBound/ThetaLifts/CartanFactorisation/CartanA.lean}{PositivityTheory.ThKF.Cart.conj_mat}
\leanref{Development/WeilBound/ThetaLifts/CartanFactorisation/CartanA.lean}{PositivityTheory.ThKF.Cart.exists_point}
\begin{proof}
Fix an infinite place $w$, and put $H=\tau_w(J)$ in the Hermitian case and
$H=i\tau_w(J)$ in the skew-Hermitian case. Then $H$ has signature $(r,s)$, and
$\theta=\theta_{Y,w}$ and $\theta^\circ=\theta_{Y^\circ,w}$ are Cartan
involutions of $H$. By Lemma~\ref{asm-theta-agree-real-frames} there are frames
$A$ and $C$ of $H$ of type $(r,s)$ with $\theta A=AJ_{r,s}$ and
$\theta^\circ C=CJ_{r,s}$; they are square, so invertible. Put $k_w=AC^{-1}$.
From $A^{\mathsf T}H\bar A=J_{r,s}=C^{\mathsf T}H\bar C$ we get
$k_w^{\mathsf T}H\bar k_w=H$, so $k_w^{\mathsf T}\tau_w(J)\bar k_w=\tau_w(J)$.
From $\theta^\circ=CJ_{r,s}C^{-1}$ we get
$k_w\theta^\circ=AJ_{r,s}C^{-1}=\theta k_w$.

As $F$ is totally real and $E$ is a CM field, $\pi_w$ identifies the component
of $\mathbf A_E$ at $w$ with $\mathbf C$, carrying $c$ to complex conjugation. So
there is $g\in GL_m(\mathbf A_E)$ with $g-1$ at $\infty$ and $\pi_w(g)=k_w$ for
every $w$. Then $g^{\mathsf T}Jc(g)=J$: at the infinite places this is the
identity for $k_w$, and at the finite places $g$ is $1$. So $g\in U_Y$ is at
$\infty$ and moves $Y^\circ$ to $Y$.
\end{proof}
% Lean: cartan_conjugacy_holds is Cart.exists_conj with c = i; cartanMove_of uses it with
% c = 1. conj_mat takes the frames from RL.frameEx and frame_inv; exists_point builds g
% through exists_mat, mat_ext, map_form and piInf_cA.

\begin{lemma}\label{asm-theta-agree-rechoice-move}
Let $(X,W,M,\chi_X,\chi_W,\omega)$ and $(X^\circ,W^\circ,M^\circ,\chi_X,\chi_W,\omega)$
be Weil pairs of the same signature for a theta base $B$ of a CM field $E$
(Definition~\ref{asm-theta-agree-triple}(1)), such that $X^\circ$ and $W^\circ$
have the matrices and lattices of $X$ and $W$ and $M^\circ$ has the isomorphism
of $M$. Let $h\in U_W$ move $W^\circ$ to $W$
(Definition~\ref{asm-theta-agree-rechoice-terms}(4)). Then there is $g\in U_X$
that moves $X^\circ$ to $X$ such that
$T(\operatorname{Pol}(M^\circ))=\operatorname{Pol}(M)$ for every archimedean
operator $T$ of $\omega$ at $(g,h)$.
\end{lemma}
\leanref{Development/WeilBound/ThetaLifts/CartanFactorisation/CartanConj.lean}{PositivityTheory.ThKF.CartanMove}
\leanref{Development/WeilBound/ThetaLifts/CartanFactorisation/Cartan.lean}{PositivityTheory.ThKF.cartanMove_of}
\begin{proof}
The spaces $X$ and $X^\circ$ are Hermitian of the same signature with the same
matrix, so Lemma~\ref{asm-theta-agree-rechoice-cartan} gives $g\in U_X$ at
$\infty$ moving $X^\circ$ to $X$. Let $T$ be an archimedean operator of
$\omega$ at $(g,h)$. By Lemma~\ref{asm-theta-agree-rechoice-fock},
$T(\operatorname{Pol}(M^\circ))\subset\operatorname{Pol}(M)$.

From $\pi_w(g)\theta_{X^\circ,w}=\theta_{X,w}\pi_w(g)$ we get
$\pi_w(g^{-1})\theta_{X,w}=\theta_{X^\circ,w}\pi_w(g^{-1})$, so $g^{-1}$ moves $X$
to $X^\circ$; likewise $h^{-1}$ moves $W$ to $W^\circ$. Let $T'$ be an
archimedean operator of $\omega$ at $(g^{-1},h^{-1})$
(Lemma~\ref{asm-theta-agree-rechoice-operators}(1)). By
Lemma~\ref{asm-theta-agree-rechoice-fock} with the two pairs exchanged,
$T'(\operatorname{Pol}(M))\subset\operatorname{Pol}(M^\circ)$. For
$f\in\mathcal S(F_\infty^I)$ and Schwartz-Bruhat $\phi$ on $\mathbf A_f^I$,
\[
T(T'(f))\otimes\phi=\omega(g,h)\omega(g^{-1},h^{-1})(f\otimes\phi)=f\otimes\phi,
\]
so $T(T'(f))=f$, taking $\phi$ with $\phi(0)=1$. Hence every
$f\in\operatorname{Pol}(M)$ is $T(T'(f))$ with $T'(f)\in\operatorname{Pol}(M^\circ)$.
\end{proof}
% Lean: cartanMove_of hO: Cart.exists_conj with c = 1, derEx and actsLaws for hD1 and hD2,
% intertw for T and T', pol_le twice (l1, l2), and actsLaws for T (T' f) = f.

\begin{lemma}\label{asm-theta-agree-rechoice-transport}
Let $(B,G,\mathcal X,\omega,\omega')$ be a triple for a theta base $B$ of a CM
field $E$ (Definition~\ref{asm-theta-agree-triple}), with models $M$ and $M'$,
and let $R=(\Pi,p,p')$ be real maps matched to it
(Definition~\ref{asm-theta-agree-matched}(2)). Let $V^\circ$, $V'^\circ$ and $W^\circ$
be lattice spaces such that $(B,G^\circ,\mathcal X,\omega,\omega')$ is a triple,
where $G^\circ=G[V^\circ,V'^\circ,W^\circ]$, with lattice spaces $\tilde V$,
$\tilde V'$ and $\tilde W$ and models $M^\circ$ and $M'^\circ$. Let $g\in U_V$
move $\tilde V$ to $V$, $g'\in U_{V'}$ move $\tilde V'$ to $V'$, and $h\in U_W$
move $\tilde W$ to $W$ (Definition~\ref{asm-theta-agree-rechoice-terms}(4)), and
let $T$ and $T'$ be archimedean operators of $\omega$ at $(g,h)$ and of
$\omega'$ at $(g',h)$ with $T(\operatorname{Pol}(M^\circ))=\operatorname{Pol}(M)$
and $T'(\operatorname{Pol}(M'^\circ))=\operatorname{Pol}(M')$. Then
$(\Pi,p\circ T,p'\circ T')$ are real maps matched to
$(B,G^\circ,\mathcal X,\omega,\omega')$.
\end{lemma}
\leanref{Development/WeilBound/ThetaLifts/CartanFactorisation/RealTrans.lean}{PositivityTheory.ThKF.real_trans}
\leanref{Development/WeilBound/ThetaLifts/CartanFactorisation/CommuteConj.lean}{PositivityTheory.ThKF.quot_trans}
\leanref{Development/WeilBound/ThetaLifts/CartanFactorisation/CommuteConj.lean}{PositivityTheory.ThKF.transL}
\leanref{Development/WeilBound/ThetaLifts/CartanFactorisation/CommuteConj.lean}{PositivityTheory.ThKF.Moves.up}
\leanref{Development/WeilBound/ThetaLifts/CartanFactorisation/CommuteConj.lean}{PositivityTheory.ThKF.Moves.down}
\begin{proof}
\emph{Step 1.} Let $(X,W,M_0,\ldots,\omega_0)$ be the Weil pair of $\omega$ or of
$\omega'$ for $G$, with $(X,z,T_0)$ equal to $(V,g,T)$ or $(V',g',T')$, and
$M_0^\circ$ the corresponding model of $G^\circ$. Let $\gamma,\gamma_2$ be lines
in $U_X$ and $\delta,\delta_2$ lines in $U_W$ with
$(\gamma(t),\delta(t))(z,h)=(z,h)(\gamma_2(t),\delta_2(t))$ for all $t$. If
$f\in\operatorname{Pol}(M_0^\circ)$ and $t\mapsto\omega_0(\gamma_2(t),\delta_2(t))$
acts at $\infty$ to first order by $f\mapsto f_1$, then
$t\mapsto\omega_0(\gamma(t),\delta(t))$ acts at $\infty$ to first order by
$T_0(f)\mapsto T_0(f_1)$. Indeed, choose archimedean operators $S_t$ of
$\omega_0$ at $(\gamma_2(t),\delta_2(t))$ and $S'_t$ at $(\gamma(t),\delta(t))$.
Lemma~\ref{asm-theta-agree-rechoice-lie}, for the pair of $M_0^\circ$, gives $z_1$
such that the first family acts by $f\mapsto z_1$ and $L(z_1)$ is the derivative
of $t\mapsto L(S_tf)$ at $0$; then $z_1=f_1$ by
Lemma~\ref{asm-theta-agree-real-weak}(3). The same lemma, for the pair of $M_0$
and $T_0(f)\in\operatorname{Pol}(M_0)$, gives $y$ such that the second family
acts by $T_0(f)\mapsto y$ and $L(y)$ is the derivative of
$t\mapsto L(S'_tT_0(f))$ at $0$. By
Lemma~\ref{asm-theta-agree-rechoice-operators}(3), $\omega_0(\gamma(t),\delta(t))$
acts at $\infty$ by $T_0(f)\mapsto T_0(S_tf)$, and also by
$T_0(f)\mapsto S'_tT_0(f)$, so $S'_tT_0(f)=T_0(S_tf)$. For a continuous linear
form $L$, $L\circ T_0$ is one too, so $L(T_0(f_1))$ is the derivative of
$t\mapsto L(T_0(S_tf))=L(S'_tT_0(f))$ at $0$, which is $L(y)$. Hence
$y=T_0(f_1)$.

\emph{Step 2: condition (a).} The map $p\circ T$ sends
$\operatorname{Pol}(M^\circ)$ onto $p(\operatorname{Pol}(M))=\Pi$. Let $\gamma$
be a line in $U_V$, $f\in\operatorname{Pol}(M^\circ)$, and let
$t\mapsto\omega(\gamma(t),1)$ act at $\infty$ to first order by $f\mapsto f_1$.
As $g-1$ is at $\infty$, $g\gamma g^{-1}$ is a line in $U_V$, and
$(g\gamma(t)g^{-1},1)(g,h)=(g,h)(\gamma(t),1)$. By Step 1,
$t\mapsto\omega(g\gamma(t)g^{-1},1)$ acts at $\infty$ to first order by
$T(f)\mapsto T(f_1)$, with $T(f)\in\operatorname{Pol}(M)$, so $p(T(f_1))=0$ by
condition (a) for $R$. Let $k\in K_\infty(\tilde V)$ and let $\omega(k,1)$ act at
$\infty$ by $f\mapsto f_1$. As $g$ moves $\tilde V$ to $V$, $gkg^{-1}\in K_\infty(V)$,
and $(gkg^{-1},1)(g,h)=(g,h)(k,1)$. By
Lemma~\ref{asm-theta-agree-rechoice-operators}(3), $\omega(gkg^{-1},1)$ acts at
$\infty$ by $T(f)\mapsto T(f_1)$, so $p(T(f_1))=p(T(f))$. The same argument
applies to $p'\circ T'$, with $g'$.

\emph{Step 3: condition (b).} Let $\gamma$ be a line in $U_W$,
$f\in\operatorname{Pol}(M^\circ)$ and $f'\in\operatorname{Pol}(M'^\circ)$ with
$p(T(f))=p'(T'(f'))$, and let $f\mapsto f_1$ and $f'\mapsto f'_1$ be the actions
at $\infty$ to first order of $t\mapsto\omega(1,\gamma(t))$ and
$t\mapsto\omega'(1,\gamma(t))$. Then $h\gamma h^{-1}$ is a line in $U_W$, and
$(1,h\gamma(t)h^{-1})(z,h)=(z,h)(1,\gamma(t))$ for $z=g,g'$. By Step 1, the
actions along $h\gamma h^{-1}$ send $T(f)$ to $T(f_1)$ and $T'(f')$ to
$T'(f'_1)$, and condition (b) for $R$ gives $p(T(f_1))=p'(T'(f'_1))$. For
$k\in K_\infty(\tilde W)$, use $hkh^{-1}\in K_\infty(W)$ and
Lemma~\ref{asm-theta-agree-rechoice-operators}(3) in the same way.

\emph{Step 4: condition (c).} $\Pi\neq0$. Let $N\subset\Pi$ be a subspace with
$p(T(f_1))\in N$ whenever $f\in\operatorname{Pol}(M^\circ)$, $p(T(f))\in N$ and
$f\mapsto f_1$ is one of the actions of $\omega$ in (b) for $G^\circ$. We check
the hypothesis of condition (c) for $R$. Let $f\in\operatorname{Pol}(M)$ with
$p(f)\in N$, and write $f=T(u)$ with $u\in\operatorname{Pol}(M^\circ)$. Let
$\gamma$ be a line in $U_W$, and let $t\mapsto\omega(1,\gamma(t))$ act at
$\infty$ to first order by $f\mapsto y$. Then $h^{-1}\gamma h$ is a line in
$U_W$, and $(1,\gamma(t))(g,h)=(g,h)(1,h^{-1}\gamma(t)h)$. By
Lemma~\ref{asm-theta-agree-rechoice-linear}(1) for the pair of $M^\circ$, there is
$z_1\in\operatorname{Pol}(M^\circ)$ such that $t\mapsto\omega(1,h^{-1}\gamma(t)h)$
acts at $\infty$ to first order by $u\mapsto z_1$. By Step 1 and
Lemma~\ref{asm-theta-agree-real-weak}(3), $y=T(z_1)$, so $p(y)\in N$. For
$k\in K_\infty(W)$, $h^{-1}kh\in K_\infty(\tilde W)$, and
Lemma~\ref{asm-theta-agree-rechoice-linear}(2) and
Lemma~\ref{asm-theta-agree-rechoice-operators}(3) give the same conclusion. By
condition (c) for $R$, $N=0$ or $N=\Pi$.
\end{proof}
% Lean: real_trans (quot_trans for (a); r.lie and r.cpt with transL, transK, conj_line and
% cj2 for (b); pol_lie, pol_k, cj3, Moves.down and ArchUniq for (c)). Step 1 is transL.

\begin{lemma}\label{asm-theta-agree-rechoice}
Let $(B,G,\mathcal X,\omega,\omega',R,R_f)$ be a matched quadruple for a theta
base $B$ of a CM field $E$ (Definition~\ref{asm-theta-agree-matched}), let
$V^\circ$, $V'^\circ$ and $W^\circ$ be lattice spaces, $G^\circ=G[V^\circ,V'^\circ,W^\circ]$,
and $R_2$ and $R_{f,2}$ maps such that
$(B,G^\circ,\mathcal X,\omega,\omega',R_2,R_{f,2})$ is a matched quadruple with a
nonzero polynomial constant. Then $(B,G,\mathcal X,\omega,\omega',R,R_f)$ has a
nonzero polynomial constant.
\end{lemma}
\leanref{Development/AbelianVarietyTheory/Definitions/Divisors/CTSmallCone.lean}{GlobalTheory.ThK.PolRe}
\leanref{Development/AutomorphicSpectrum/Proofs/BlockDecomposition/PolReOf2T.lean}{GlobalTheory.CompanionTheory.ThetaK.polRe_of2T}
\leanref{Development/WeilBound/CoreStatements/TemperedRepresentationTheory/HoweMaximalTori/MaxFinTOfPublished.lean}{PositivityTheory.HoweDualityUnitaryAlt.polRe_of2T}
\begin{proof}
Write $R=(\Pi,p,p')$, $R_f=(\Pi_v,p_v,p'_v)_v$, $R_2=(\Pi_2,p_2,p'_2)$ and
$R_{f,2}=(\Pi_{2,v},p_{2,v},p'_{2,v})_v$, and let $c_2\neq0$ be a polynomial
constant of the second quadruple. Let $(a,a')$ be the ranks, $M$ and $M'$ the
models of $G$, and $\tilde V$, $\tilde V'$, $\tilde W$, $M^\circ$ and $M'^\circ$
the lattice spaces and models of $G^\circ$. The models of $G^\circ$ have the
isomorphisms, vacua, measures and domains of those of $G$, so the two
quadruples have the same theta integrals $I$ and $I'$. Also
$\mathcal S^\circ(M'^\circ)=\mathcal S^\circ(M')$, since full rank at $\tau_1$
involves only the matrices of $V'$ and $W$.

\emph{Step 1: moving the Cartan involutions.} Both triples are proper and have
the same integers $p$ and $q$, so $W$ and $\tilde W$ are skew-Hermitian of
signature $(a+a'-p-q,p+q)$ with the matrix of $W$. By
Lemma~\ref{asm-theta-agree-rechoice-cartan} there is $h\in U_W$ that moves
$\tilde W$ to $W$. The Weil pairs $(V,W,M,\beta,\beta/\alpha,\omega)$ and
$(\tilde V,\tilde W,M^\circ,\beta,\beta/\alpha,\omega)$ have the same signature,
so Lemma~\ref{asm-theta-agree-rechoice-move} gives $g\in U_V$ that moves
$\tilde V$ to $V$ with $T(\operatorname{Pol}(M^\circ))=\operatorname{Pol}(M)$ for
every archimedean operator $T$ of $\omega$ at $(g,h)$. In the same way there is
$g'\in U_{V'}$ that moves $\tilde V'$ to $V'$ with
$T'(\operatorname{Pol}(M'^\circ))=\operatorname{Pol}(M')$ for every archimedean
operator $T'$ of $\omega'$ at $(g',h)$. Fix such $T$ and $T'$
(Lemma~\ref{asm-theta-agree-rechoice-operators}(1)).

\emph{Step 2: the archimedean factor.} By
Lemma~\ref{asm-theta-agree-rechoice-transport}, $(\Pi,p\circ T,p'\circ T')$ are
real maps matched to $(B,G^\circ,\mathcal X,\omega,\omega')$, and so are
$(\Pi_2,p_2,p'_2)$. Lemma~\ref{asm-theta-agree-rechoice-compare}(1) gives a
linear map $e\colon\Pi\to\Pi_2$ and $c\neq0$ with $e(p(T(u)))=p_2(u)$ for
$u\in\operatorname{Pol}(M^\circ)$ and $e(p'(T'(u')))=c\,p'_2(u')$ for
$u'\in\operatorname{Pol}(M'^\circ)$.

\emph{Step 3: the finite factors.}
Lemma~\ref{asm-theta-agree-rechoice-compare}(2) gives, for every finite $v$, a
linear map $e_v\colon\Pi_v\to\Pi_{2,v}$ and $c_v\neq0$ with
$e_v(p_v(\phi))=p_{2,v}(\phi)$ and $e_v(p'_v(\phi'))=c_v\,p'_{2,v}(\phi')$ for
Schwartz-Bruhat $\phi$ and $\phi'$, and $c_v=1$ for $v\notin T$. Put
$\Lambda=c\prod_{v\in T}c_v$, which is not $0$. We show that $c_2\Lambda$ is a
polynomial constant of the first quadruple.

\emph{Step 4: moving the data.} Let $\Phi'\in\mathcal S^\circ(M')$, and let
$\Phi$ and $\Phi'$ correspond, with $T_1\supset T$,
$\Phi=\sum_if_i\otimes\bigotimes_v\phi_{i,v}$ and
$\Phi'=\sum_jf'_j\otimes\bigotimes_v\phi'_{j,v}$ as in
Definition~\ref{asm-theta-agree-matched}(4), so that
\[
\sum_ip(f_i)\otimes\bigotimes_{v\in T_1}p_v(\phi_{i,v})
=\sum_jp'(f'_j)\otimes\bigotimes_{v\in T_1}p'_v(\phi'_{j,v}).
\tag{$*$}
\]
By Step 1, $f_i=T(u_i)$ with $u_i\in\operatorname{Pol}(M^\circ)$ and
$f'_j=T'(u'_j)$ with $u'_j\in\operatorname{Pol}(M'^\circ)$. Put
$\Psi=\sum_iu_i\otimes\bigotimes_v\phi_{i,v}$ and
$\Psi'=\sum_ju'_j\otimes\bigotimes_v\phi'_{j,v}$. As $\phi_{i,v}=\mathrm{vac}_v$
for $v\notin T_1$ and $\mathrm{vac}_v$ is the indicator function of
$\mathcal O_v^{I_0}$ for almost all $v$, $\bigotimes_v\phi_{i,v}$ is a
Schwartz-Bruhat function on $\mathbf A_f^{I_0}$. So $\Psi$ is Schwartz-Bruhat,
and $\omega(g,h)\Psi=\Phi$, as $T$ is an archimedean operator. Likewise
$\omega'(g',h)\Psi'=\Phi'$. By Lemma~\ref{asm-theta-agree-coinvariants} the
first quadruple describes the coinvariants, so $\omega'((g',h)^{-1})$ maps
$\mathcal S^\circ(M')$ into itself, and
$\Psi'=\omega'((g',h)^{-1})\Phi'\in\mathcal S^\circ(M')=\mathcal S^\circ(M'^\circ)$.

\emph{Step 5: the moved data correspond.} We have
$\Lambda\Psi'=\sum_j(\Lambda u'_j)\otimes\bigotimes_v\phi'_{j,v}$ with
$\Lambda u'_j\in\operatorname{Pol}(M'^\circ)$. Apply to $(*)$ the linear map
$e\otimes\bigotimes_{v\in T_1}e_v$. The left side becomes
$\sum_ip_2(u_i)\otimes\bigotimes_{v\in T_1}p_{2,v}(\phi_{i,v})$. The right side
becomes
$\sum_jc\,p'_2(u'_j)\otimes\bigotimes_{v\in T_1}c_v\,p'_{2,v}(\phi'_{j,v})
=\sum_jp'_2(\Lambda u'_j)\otimes\bigotimes_{v\in T_1}p'_{2,v}(\phi'_{j,v})$,
since $c_v=1$ for $v\in T_1\setminus T$. So $\Psi$ and $\Lambda\Psi'$ correspond
for the second quadruple, with the same set $T_1$.

\emph{Step 6: the theta integrals.} As $c_2$ is a polynomial constant of the
second quadruple and $\Lambda\Psi'\in\mathcal S^\circ(M'^\circ)$,
$I(\Psi)=c_2I'(\Lambda\Psi')=c_2\Lambda I'(\Psi')$. Let $x\in U_W$. As
$\omega(g,h)=\omega(g,1)\omega(1,h)$ and $\omega(1,h)\Psi$ is Schwartz-Bruhat,
Lemma~\ref{asm-theta-agree-anisotropic-invariance} gives
$I(\Phi)(x)=I(\omega(1,h)\Psi)(x)$, with the convention there. As $\omega$ is a
representation,
$\theta_\omega(\omega(1,h)\Psi)(g_1,x)=\theta_\omega(\Psi)(g_1,xh)$, so
$I(\omega(1,h)\Psi)(x)=I(\Psi)(xh)$. In the same way
$I'(\Phi')(x)=I'(\Psi')(xh)$, since $\Psi'$ is Schwartz-Bruhat. Hence
$I(\Phi)(x)=c_2\Lambda I'(\Psi')(xh)=c_2\Lambda I'(\Phi')(x)$.
\end{proof}
% Lean: polRe_of2T hH hM hO hL hQ hC hJ hV hK hI. Step 1 is hJ (CartanConj) and hV
% (CartanMove) twice, with hO (ArchOp) for T and T'; Step 2 is real_trans and
% archimedean_matched_maps_compare; Step 3 is cmp_finT and cf1; Step 4 is move_sum, pure_mem_schA
% and the first clause of KCoinv with so_re; Step 5 is the map TensorProduct.map e
% (PiTensorProduct.map ef) and Finset.prod_subset; Step 6 is I_move (ILeft and I_right) and
% I_smul. In the closing of Theorem K, closeSph passes JoinTri = Join.joinTri_of, MaxFinT =
% Sph.maxFinT_of SphTri, ArchOp = A.archOp, ArchLie = archLie' localLift_holds, ArchUniq =
% A.archUniq, PolCount = polCount_of_derEx, CartanConj = cartan_conjugacy_holds and CartanMove =
% cartanMove_of.

\subsection{Nonvanishing through the doubling method}\label{subsection-asm-nonvanishing-doubling}

\begin{lemma}\label{asm-theta-agree-nonvanishing-schwartz}
Let $(X,W,M,\chi_X,\chi_W,\omega)$ be a Weil pair of signature $(r,s;r',s')$
for a theta base $B$ of a CM field $E$
(Definition~\ref{asm-theta-agree-triple}(1)), with $X$ and $W$ of ranks $n$ and
$d$. Then every element of $\operatorname{PD}(M)$ is a Schwartz-Bruhat function
on $\mathbf A^{n\times d}$. If moreover $\min(r,s)=0$ and $I(\Phi)=0$ for every
$\Phi\in\operatorname{PD}(M)$, then $I(\Phi)=0$ for every Schwartz-Bruhat
function $\Phi$ on $\mathbf A^{n\times d}$.
\end{lemma}
\leanref{Development/WeilBound/ThetaLifts/LatticeBasics/SchToPol.lean}{GlobalTheory.ThK.Dbl.PolSch}
\leanref{Development/WeilBound/ThetaLifts/LatticeBasics/ExistsIsGauss.lean}{GlobalTheory.ThK.Dbl.polSch}
\leanref{Development/WeilBound/ThetaLifts/LatticeBasics/SchToPol.lean}{GlobalTheory.ThK.Dbl.SchToPol}
\leanref{Development/WeilBound/ThetaLifts/LatticeBasics/ExistsIsGauss.lean}{GlobalTheory.ThK.Dbl.schToPol_of}
\leanref{Development/WeilBound/ThetaLifts/LatticeBasics/IsCompactIntV.lean}{GlobalTheory.ThK.Dbl.Dense.schA_le}
\begin{proof}
\emph{Polynomial data.} Let $f\in\operatorname{Pol}(M)$, and for each finite
place $v$ let $\phi_v$ be a Schwartz-Bruhat function on $F_v^{n\times d}$ with
$\phi_v=\mathrm{vac}_v$ for almost all $v$. As $M$ is regular,
$\mathrm{vac}_v$ is the indicator function of $\mathcal O_v^{n\times d}$ for
almost all $v$, and hence so is $\phi_v$. So $\phi(x)=\prod_v\phi_v(x_v)$ is a
Schwartz-Bruhat function on $\mathbf A_f^{n\times d}$, and
$f\otimes\bigotimes_v\phi_v=f\otimes\phi$ is a Schwartz-Bruhat function on
$\mathbf A^{n\times d}$. These functions span $\operatorname{PD}(M)$, which
proves the first assertion.

\emph{Linearity.} Now let $\min(r,s)=0$. By
Lemma~\ref{asm-theta-agree-convergence}(1), every Schwartz-Bruhat function
converges. So for Schwartz-Bruhat $\Phi$ and $\Psi$, $c\in\mathbf C$ and
$h\in U_W$, the functions $g\mapsto\theta_\omega(\Phi)(g,h)$ and
$g\mapsto\theta_\omega(\Psi)(g,h)$ are integrable on $\mathcal D$, and
Lemma~\ref{asm-theta-agree-convergence-theta}(2) gives
$I(\Phi+c\Psi)=I(\Phi)+c\,I(\Psi)$.

\emph{Pure tensors.} Let $\phi=\prod_v\phi_v$ with every $\phi_v$
Schwartz-Bruhat on $F_v^{n\times d}$ and $\phi_v$ the indicator function of
$\mathcal O_v^{n\times d}$ for almost all $v$, and fix $h\in U_W$. As $M$ is
regular, $\phi_v=\mathrm{vac}_v$ for almost all $v$, so
$f\otimes\phi\in\operatorname{PD}(M)$ for every $f\in\operatorname{Pol}(M)$. The
map $\lambda(f)=I(f\otimes\phi)(h)$ on $\mathcal S(F_\infty^{n\times d})$ is
linear by the last paragraph, continuous by
Lemma~\ref{asm-theta-agree-convergence}(2), and zero on $\operatorname{Pol}(M)$
by hypothesis. As $\operatorname{Pol}(M)$ is dense
(Lemma~\ref{asm-theta-agree-density}), $\lambda=0$.

\emph{All Schwartz-Bruhat functions.} A Schwartz-Bruhat function on
$\mathbf A_f^{n\times d}$ is locally constant with compact support, so it is a
finite linear combination of indicator functions of sets
$x+\prod_v\varpi_v^{k_v}\mathcal O_v^{n\times d}$ with $k_v=0$ and
$x_v\in\mathcal O_v^{n\times d}$ for almost all $v$, and each of these is a
product $\phi$ as in the last paragraph. So every Schwartz-Bruhat function on
$\mathbf A^{n\times d}$ is a finite linear combination of functions
$f\otimes\phi$ as there, and $I$ vanishes on it by linearity.
\end{proof}
% Lean: polSch; schToPol_of WeilConv, with PositivityTheory.ThKG.I_add for the linearity,
% Dense.schA_le (through mem_span_pf) for the span of the pure tensors with Dense.Good
% components, the second clause of WeilConv for the continuity, and Dense.pol_dense.

\begin{lemma}\label{asm-theta-agree-nonvanishing-full-rank}
Let $(B,G,\mathcal X,\omega,\omega')$ be a triple for a theta base $B$ of a CM
field $E$ (Definition~\ref{asm-theta-agree-triple}), with $q=0$. Then for every
$\Phi'\in\operatorname{PD}(M')$ there is
$\Psi\in\operatorname{PD}(M')\cap\mathcal S^\circ(M')$ with
$I'(\Psi)=I'(\Phi')$.
\end{lemma}
\leanref{Development/WeilBound/ThetaLifts/LatticeBasics/SchToPol.lean}{GlobalTheory.ThK.Dbl.ToSo}
\leanref{Development/WeilBound/ThetaLifts/LatticeBasics/ConvOfCpt.lean}{GlobalTheory.ThK.Dbl.toSo_of}
\leanref{Development/WeilBound/ThetaLifts/LatticeBasics/ConvOfCpt.lean}{GlobalTheory.ThK.Dbl.conv_of_cpt}
\leanref{Development/WeilBound/ThetaLifts/LatticeBasics/ConvOfCpt.lean}{GlobalTheory.ThK.Dbl.I_aug}
\begin{proof}
Let $G$ have ranks $(a,a')$, and put
$I'_0=\{1,\ldots,a'\}\times\{1,\ldots,a+a'\}$. As $q=0$, the space $V'$ is
Hermitian of signature $(0,a')$, and $(V',W,M',\alpha,\alpha/\beta,\omega')$ is
a Weil pair (Definition~\ref{asm-theta-agree-triple}(4)). Let $\mu'$ and
$\mathcal D'$ be the Haar measure and the fundamental domain of $M'$.

\emph{Convergence.} Let $\Phi$ be a Schwartz-Bruhat function on
$\mathbf A^{I'_0}$ and $h\in U_W$. The function
$g\mapsto\theta_{\omega'}(\Phi)(g,h)$ is continuous
(Lemma~\ref{asm-theta-agree-convergence-theta}(2)) and invariant under
$U_{V'}(F)$ on the left (Lemma~\ref{asm-theta-agree-anisotropic-automorphy}),
so it is bounded by Lemma~\ref{asm-theta-agree-definite-compact}(4), for the
definite space $V'$. As $\mu'(\mathcal D')<\infty$ by
Lemma~\ref{asm-theta-agree-definite-compact}(3), it is integrable on
$\mathcal D'$. So every Schwartz-Bruhat function converges for
$(M',\omega')$, and, as in the proof of
Lemma~\ref{asm-theta-agree-nonvanishing-schwartz}, $I'$ is linear on the
Schwartz-Bruhat functions.

\emph{The place $v_1$.} By Lemmas~\ref{asm-theta-agree-real}
and~\ref{asm-theta-agree-finite} there are real maps $R$ and finite maps
$R_f=(\Pi_v,p_v,p'_v)_v$ matched to the triple, and by
Lemma~\ref{asm-theta-agree-coinvariants} the matched quadruple
$(B,G,\mathcal X,\omega,\omega',R,R_f)$ describes the coinvariants
(Definition~\ref{asm-theta-agree-parts}(3)). The place $v_1$ is clean for
$\mathcal X$, as $\mathcal X$ is proper. So a Schwartz-Bruhat function $\phi$ on
$F_{v_1}^{I'_0}$ has $p'_{v_1}(\phi)=0$ if and only if
$\phi\in\mathfrak a_{v_1}(\omega')$, and every element of $\Pi_{v_1}$ is
$p'_{v_1}(y)$ for some $y$ of full rank at $\tau_1$.

\emph{Generators.} Let $\Phi'=f\otimes\bigotimes_v\phi_v$ with
$f\in\operatorname{Pol}(M')$, every $\phi_v$ Schwartz-Bruhat and
$\phi_v=\mathrm{vac}'_v$ for almost all $v$. Choose $y$ of full rank at $\tau_1$
with $p'_{v_1}(y)=p'_{v_1}(\phi_{v_1})$, and let $\Psi$ and $\Phi''$ be
$\Phi'$ with its component at $v_1$ replaced by $y$ and by $\phi_{v_1}-y$. Then
$\Psi\in\operatorname{PD}(M')\cap\mathcal S^\circ(M')$ and $\Phi'=\Psi+\Phi''$.
As $p'_{v_1}(\phi_{v_1}-y)=0$, $\phi_{v_1}-y\in\mathfrak a_{v_1}(\omega')$, so
$I'(\Phi'')=0$ by Lemma~\ref{asm-theta-agree-anisotropic-augmentation}, applied
to $(V',W,M',\alpha,\alpha/\beta,\omega')$ at $v_1$. By linearity,
$I'(\Phi')=I'(\Psi)$.

\emph{Sums.} The set of $\Phi'\in\operatorname{PD}(M')$ for which such a $\Psi$
exists contains $0$ and these generators. It is closed under sums and scalar
multiples: if $I'(\Psi_1)=I'(\Phi'_1)$ and $I'(\Psi_2)=I'(\Phi'_2)$, then
$I'(\Psi_1+c\Psi_2)=I'(\Phi'_1+c\Phi'_2)$ by linearity, since all four
functions are Schwartz-Bruhat (Lemma~\ref{asm-theta-agree-nonvanishing-schwartz}).
So it is all of $\operatorname{PD}(M')$.
\end{proof}
% Lean: toSo_of KReal KFin KCoinv ThetaBasic CptAll. CptAll (the record Cpt) is Lemma
% asm-theta-agree-definite-compact; conv_of_cpt is the convergence, I_lin the linearity, pol_sch the
% first part of Lemma asm-theta-agree-nonvanishing-schwartz, and I_aug (from I_lin and
% I_act) the replacement at v_1, as in Lemma asm-theta-agree-anisotropic-augmentation.

\begin{definition}\label{asm-theta-agree-nonvanishing-doubling}
Let $E$ be a CM field with a theta base $B$
(Definition~\ref{asm-theta-agree-models}), and let $X$ be a lattice space of
rank $n$ (Definition~\ref{asm-theta-agree-doubling}(1)), with doubled
group $G_X$, Siegel parabolic subgroup $P_X$ and $t(p)$ for $p\in P_X$
(Definition~\ref{asm-theta-agree-doubling}(2)), and with $\hat L_X$ defined from
the lattice $L_X$ as in Definition~\ref{asm-theta-agree-planes}(6).
\begin{enumerate}
\item The \emph{lattice stabilizer of the doubled space} of $X$ is the subgroup
  $K$ of the $k\in G_X$ such that $\pi_w(k)$ commutes with
  $\operatorname{diag}(\theta_{X,w},-\theta_{X,w})$ for every infinite place
  $w$, and $k$ and $k^{-1}$ map $\hat L_X\times\hat L_X$ into itself. The
  \emph{doubling embedding} $\iota\colon U_X\times U_X\to G_X$ is
  $\iota(h,h')=\operatorname{diag}(h,h')$.
\item Let $v$ be a finite place of $F$. An element $g\in U_X$ is \emph{integral
  at $v$} if $g$ and $g^{-1}$ map every element of $\hat L_X$ at $v$ into
  $\hat L_X$, and \emph{trivial at $v$} if its component at $v$ is $1$. For a
  finite set $T$ of finite places, $\Omega_T\subset U_X$ is the set of the $g$
  that are integral at every finite place $v\notin T$. For a Hecke character
  $\chi$ of $E$ and $z\in\mathbf C$, a function $f$ on $G_X$ is \emph{spherical
  at $v$} for $(\chi,z)$ if
  $f(pk)=\chi(\det t(p))\,|\det t(p)|^z$ for every $p\in P_X$ with $p-1$ at
  $v$ and every $k\in K$ with $k-1$ at $v$.
\item For a finite place $v$ put $\epsilon_v=1$ if there is an $F$-embedding
  $E\to F_v$ and $\epsilon_v=-1$ otherwise, and
  \[
  d_v(s)=\prod_{i=1}^n\bigl(1-\epsilon_v^{\,n-i}N(v)^{-(2s+i)}\bigr)^{-1}.
  \]
  For a Hecke character $\chi$ of $E$, $L_v(s,\chi)$ is the product of
  $(1-\chi_w(\varpi_w)N(w)^{-s})^{-1}$ over the places $w$ of $E$ above $v$ at
  which $\chi$ is unramified, for uniformizers $\varpi_w$, and
  $L(s,\chi)=\prod_vL_v(s,\chi)$, the product over the finite places of $F$,
  for $\operatorname{Re}s>1$.
\item Let $(X,W,M,\chi_X,\chi_W,\omega)$ be a Weil pair for $B$
  (Definition~\ref{asm-theta-agree-triple}(1)), with $W$ of rank $d$, let $\nu$
  be a Haar measure on $\mathbf A^{n\times d}$, and let $a$ be a function on
  $G_X$ with positive values. The \emph{matrix coefficient} of a
  Schwartz-Bruhat function $\Phi$ on $\mathbf A^{n\times d}$ is
  $c_\Phi(g)=\int(\omega(g,1)\Phi)(u)\,\overline{\Phi(u)}\,d\nu(u)$, and its
  \emph{zeta integrand} is
  $Z_\Phi(s,g)=a(\iota(g,1))^{s-(d-n)/2}\,c_\Phi(g)$, for $s\in\mathbf C$ and
  $g\in U_X$.
\end{enumerate}
\end{definition}
\leanref{Development/WeilBound/ThetaLifts/LatticeBasics/Hypotheses2.lean}{GlobalTheory.ThK.Dbl.IsStab2}
\leanref{Development/WeilBound/ThetaLifts/LatticeBasics/Hypotheses2.lean}{GlobalTheory.ThK.Dbl.IsEmb}
\leanref{Development/WeilBound/ThetaLifts/LatticeBasics/HodgeTateFromSandwich.lean}{GlobalTheory.ThK.Dbl.locGrp}
\leanref{Development/WeilBound/ThetaLifts/LatticeBasics/HodgeTateFromSandwich.lean}{GlobalTheory.ThK.Dbl.IntAt}
\leanref{Development/WeilBound/ThetaLifts/LatticeBasics/HodgeTateFromSandwich.lean}{GlobalTheory.ThK.Dbl.OffAt}
\leanref{Development/WeilBound/ThetaLifts/LatticeBasics/HodgeTateFromSandwich.lean}{GlobalTheory.ThK.Dbl.IsSphAt}
\leanref{Development/WeilBound/ThetaLifts/LatticeBasics/HodgeTateFromSandwich.lean}{GlobalTheory.ThK.Dbl.locD}
\leanref{Development/AbelianVarietyTheory/Definitions/Foundations/NormEulerPolySubOneLe.lean}{GlobalTheory.LFun.Lv}
\leanref{Development/AbelianVarietyTheory/Definitions/Foundations/NotIsSplitOfIsInertAt.lean}{GlobalTheory.LFun.charL}
\leanref{Development/WeilBound/ThetaLifts/LatticeBasics/HodgeTateFromSandwich.lean}{GlobalTheory.ThK.Dbl.coef}
\leanref{Development/WeilBound/ThetaLifts/LatticeBasics/HodgeTateFromSandwich.lean}{GlobalTheory.ThK.Dbl.zint}

\begin{theorem}[Iwasawa decomposition \cite{BT72,Car79,MW95}]\label{asm-theta-agree-doubling-iwasawa}
Let $B$ be a theta base for a CM field $E$
(Definition~\ref{asm-theta-agree-models}(1)), and let $Y$ be a lattice space of
rank $n$, with matrix $J$, lattice $L$ and matrices $\theta_w$, that is
Hermitian or skew-Hermitian of some signature
(Definition~\ref{asm-theta-agree-doubling}(1)). For a place $v$ of $F$, finite
or infinite, let $G_Y(F_v)$ be the group of the
$x\in GL_{2n}(E\otimes_FF_v)$ with $x^{\mathsf T}J^\square c(x)=J^\square$, and
$P_Y(F_v)$ its subgroup of the $p$ with $p\Delta=\Delta t$ for some
$t\in GL_n(E\otimes_FF_v)$ (Definition~\ref{asm-theta-agree-doubling}(2)).
\begin{enumerate}
\item Let $v$ be finite, and let $L_v\subset(E\otimes_FF_v)^n$ be the closure
  of $L$. Then every $x\in G_Y(F_v)$ is $pk$ with $p\in P_Y(F_v)$ and
  $k\in G_Y(F_v)$ such that $k(L_v\oplus L_v)=L_v\oplus L_v$.
\item Let $w$ be infinite, and let $\tau_w$ also denote the isomorphism
  $E\otimes_FF_w\to\mathbf C$ it defines. Then every $x\in G_Y(F_w)$ is $pk$
  with $p\in P_Y(F_w)$ and $k\in G_Y(F_w)$ such that $\tau_w(k)$ commutes with
  $\operatorname{diag}(\theta_w,-\theta_w)$.
\end{enumerate}
\end{theorem}
\leanref{Development/AbelianVarietyTheory/Definitions/Foundations/T2SpaceAdeleRing.lean}{GlobalTheory.Doubling.HasIwasawa}
\leanref{Development/WeilBound/ThetaLifts/LatticeBasics/HgtFin.lean}{GlobalTheory.ThK.Dbl.Hgt.hasIwasawa_fin'}
\leanref{Development/WeilBound/ThetaLifts/LatticeBasics/HgtArch.lean}{GlobalTheory.ThK.Dbl.Hgt.exists_arch}

\begin{lemma}\label{asm-theta-agree-doubling-stabilizer}
Let $B$ be a theta base for a CM field $E$
(Definition~\ref{asm-theta-agree-models}(1)), let $Y$ be a lattice space of
rank $n$ that is Hermitian or skew-Hermitian of some signature
(Definition~\ref{asm-theta-agree-doubling}(1)), with doubled group $G_Y$ and
Siegel parabolic subgroup $P_Y$ (Definition~\ref{asm-theta-agree-doubling}(2)),
and let $K$ be the lattice stabilizer of the doubled space of $Y$ and $\iota$
the doubling embedding
(Definition~\ref{asm-theta-agree-nonvanishing-doubling}(1)). Then:
\begin{enumerate}
\item $K$ is a compact subgroup of $G_Y$;
\item every $x\in G_Y$ is $pk$ with $p\in P_Y$ and $k\in K$;
\item there is exactly one height function $a$ for $K$
  (Definition~\ref{asm-theta-agree-doubling}(4)), and it is continuous with
  positive values;
\item $\iota$ is a continuous homomorphism from $U_Y\times U_Y$ to $G_Y$.
\end{enumerate}
\end{lemma}
\leanref{Development/WeilBound/ThetaLifts/LatticeBasics/PairsStable.lean}{GlobalTheory.ThK.Dbl.HgtEx}
\leanref{Development/WeilBound/ThetaLifts/LatticeBasics/HgtIntOf.lean}{GlobalTheory.ThK.Dbl.hgtEx_of}
\leanref{Development/WeilBound/ThetaLifts/LatticeBasics/ContinuousFlat.lean}{GlobalTheory.ThK.Dbl.Stab2Cpt}
\leanref{Development/WeilBound/ThetaLifts/LatticeBasics/CartanBound2.lean}{GlobalTheory.ThK.Dbl.stab2Cpt}
\leanref{Development/WeilBound/ThetaLifts/LatticeBasics/CartanBound2.lean}{GlobalTheory.ThK.Dbl.cartan_bound2}
\leanref{Development/WeilBound/ThetaLifts/LatticeBasics/HgtExA.lean}{GlobalTheory.ThK.Dbl.Hgt.stab2}
\leanref{Development/WeilBound/ThetaLifts/LatticeBasics/HgtExA.lean}{GlobalTheory.ThK.Dbl.Hgt.iwasawa2}
\leanref{Development/WeilBound/ThetaLifts/LatticeBasics/HgtExA.lean}{GlobalTheory.ThK.Dbl.Hgt.hgtEx_core}
\leanref{Development/WeilBound/ThetaLifts/LatticeBasics/HgtExA.lean}{GlobalTheory.ThK.Dbl.Hgt.emb2}
\leanref{Development/AbelianVarietyTheory/Proofs/FoundationProofs/ExistsSiegelHProjArchEq.lean}{GlobalTheory.Doubling.HasIwasawa.adelic}
\leanref{Development/AbelianVarietyTheory/Definitions/Foundations/IsIwasawaOfIsCompact.lean}{GlobalTheory.Doubling.isIwasawa_of_isCompact}
\leanref{Development/AbelianVarietyTheory/Definitions/Foundations/IsIwasawaOfIsCompact.lean}{GlobalTheory.Doubling.IsIwasawa.continuous_iwAbs}
\begin{proof}
Let $J$, $L$ and $\theta_w$ be the matrix, the lattice and the matrices at
infinity of $Y$, and put $\Theta_w=\operatorname{diag}(\theta_w,-\theta_w)$.

(1) The conditions defining $K$ are preserved by products and inverses, so $K$
is a subgroup. Each condition is closed: $\pi_w$ is continuous, and
$\hat L_Y\times\hat L_Y$ is compact, hence closed. Let $c=1$ if $Y$ is
Hermitian and $c=i$ if $Y$ is skew-Hermitian, and fix an infinite place $w$.
Then $M_w=c\,\tau_w(J)$ has the Cartan involution $\theta_w$, and $\Theta_w$ is
a Cartan involution of $M_w^\square=\operatorname{diag}(M_w,-M_w)$: it squares
to $1$, preserves the form, and for $(y,z)\neq0$ the real part of
$(\Theta_w(y,z))^{\mathsf T}M_w^\square\overline{(y,z)}$ is the sum of those of
$(\theta_wy)^{\mathsf T}M_w\bar y$ and $(\theta_wz)^{\mathsf T}M_w\bar z$, which
is positive. For $k\in K$, the matrix $A=\pi_w(k)$ satisfies
$A^{\mathsf T}M_w^\square\bar A=M_w^\square$ and commutes with $\Theta_w$, so it
preserves the form $(x,y)\mapsto(\Theta_wx)^{\mathsf T}M_w^\square\bar y$, whose
real part on the diagonal is positive definite. Hence the entries of $\pi_w(k)$
are bounded by a constant $R_w$ depending only on $M_w$ and $\theta_w$. At the
finite places, $k$ and $k^{-1}$ map $\hat L_Y\times\hat L_Y$ into itself, so the
finite components of the entries of each of the four blocks of $k$ lie in a
compact set depending only on $L$. So the entries of $k$ and of $k^{-1}$ lie in
fixed compact subsets of $\mathbf A_E$. The pairs $(k,k^{-1})$, $k\in K$, form
a closed subset of a product of compact sets, so $K$ is compact in $G_Y$.

(2) There is $\xi\in E^\times$ with $c(\xi J^{\mathsf T})^{\mathsf T}=\xi J^{\mathsf T}$:
$\xi=1$ if $Y$ is Hermitian, and $\xi=y-c(y)$ for some $y\in E$ with
$c(y)\neq y$ if $Y$ is skew-Hermitian. Up to transposition, $G_Y$ is the unitary
group of the Hermitian form $\operatorname{diag}(\xi J^{\mathsf T},-\xi J^{\mathsf T})$,
with the same Siegel parabolic subgroup, so
Theorem~\ref{asm-theta-agree-doubling-iwasawa} applies at every place. Let
$x\in G_Y$. For almost all finite places $v$ the component $x_v$ maps
$L_v\oplus L_v$ onto itself; let $S$ be the finite set of the other finite
places. For $v\in S$ write $x_v=p_vk_v$ by
Theorem~\ref{asm-theta-agree-doubling-iwasawa}(1), and for every infinite place
$w$ write $x_w=p_wk_w$ by Theorem~\ref{asm-theta-agree-doubling-iwasawa}(2).
Let $p\in P_Y$ have the components $p_v$ for $v\in S$, $p_w$ at the infinite
places and $1$ at the other finite places. Then $k=p^{-1}x$ has finite
components that map $L_v\oplus L_v$ onto itself, and so do those of $k^{-1}$;
as an element of $\mathbf A_E^n$ lies in $\hat L_Y$ if and only if its
components at $\infty$ are $0$ and its component at every finite $v$ lies in
$L_v$, $k$ and $k^{-1}$ map $\hat L_Y\times\hat L_Y$ into itself. Its components
at infinity are the $k_w$, so $\pi_w(k)$ commutes with $\Theta_w$. Hence
$k\in K$ and $x=pk$.

(3) If $a$ is a height function for $K$, then $a(x)=|\det t(p)|$ for any
decomposition $x=pk$ of (2), so there is at most one. The map
$p\mapsto|\det t(p)|$ is a continuous homomorphism from $P_Y$ to
$\mathbf R_{>0}$, as $t(pp')=t(p)t(p')$. The group $P_Y\cap K$ is compact, as
$P_Y$ is closed, so its image is a compact subgroup of $\mathbf R_{>0}$, which
is $\{1\}$. If $pk=p'k'$ with $p,p'\in P_Y$ and $k,k'\in K$, then
$p'^{-1}p=k'k^{-1}\in P_Y\cap K$, so $|\det t(p)|=|\det t(p')|$. Hence
$a(pk)=|\det t(p)|$ defines a height function for $K$, with positive values.
The multiplication map $m\colon P_Y\times K\to G_Y$ is continuous, surjective by
(2), and closed, since $P_Y$ is closed and $K$ is compact; so it is a quotient
map. As $a\circ m$ is the continuous map $(p,k)\mapsto|\det t(p)|$, $a$ is
continuous.

(4) For $h,h'\in U_Y$,
$\operatorname{diag}(h,h')^{\mathsf T}J^\square c(\operatorname{diag}(h,h'))=\operatorname{diag}(h^{\mathsf T}Jc(h),-h'^{\mathsf T}Jc(h'))=J^\square$,
so $\iota(h,h')\in G_Y$; and $\iota$ is a continuous homomorphism.
\end{proof}
% Lean: (1) stab2Cpt: cartan_bound2 (through cartan_bound) at infinity, fin_bound at the finite
% places, exists_section and isCompact_of_units. (2) Hgt.iwasawa2: exists_scal gives xi, mem_HR
% and back the transposed group, Frames.exists_isLatticeFrame_of_arch the adelic frame a,
% HasIwasawa.adelic (eventually_proj_mem_stabLoc, exists_siegelH_proj_eq,
% exists_siegelH_projArch_eq) with hasIwasawa_fin' and exists_arch, then keeps_back
% (mem_latHat) and back_comm. (3) hgtEx_core: isIwasawa_of_isCompact (triv_of_isCompact) and
% IsIwasawa.continuous_iwAbs (continuous_of_prod); uniqueness as in hgtInt_of. (4) emb2.

\begin{lemma}\label{asm-theta-agree-doubling-count}
Let $E$ be a CM field with maximal totally real subfield $F$, write $\mathbf A$
for the adeles of $F$, $N=N_{E/F}$ and $|u|=|N(u)|_{\mathbf A}$ for
$u\in\mathbf A_E^\times$, let $H$ be a closed subgroup of $GL_n(\mathbf A_E)$
with a Haar measure $\mu$, and let $\Omega\subset M_n(\mathbf A_E)$ and
$\Omega_0\subset\mathbf A$ be compact. For $j\geq0$ let $H_j$ be the set of the
$g\in H$ for which some $u\in\mathbf A_E^\times$ has $N(u)\in\Omega_0$,
$|u|\geq2^{-j}$, $N(u)g\in\Omega$ and $N(u)g^{-1}\in\Omega$. Then there are
$R\geq1$, $C<\infty$ and $b\geq0$ such that $|u|\leq R$ for every
$u\in\mathbf A_E^\times$ with $N(u)\in\Omega_0$, and $\mu(H_j)\leq Cb^j$ for
every $j\geq0$.
\end{lemma}
\leanref{Development/WeilBound/ThetaLifts/LatticeBasics/CompactBound.lean}{GlobalTheory.ThK.Dbl.Hgt.count_ad}
\leanref{Development/WeilBound/ThetaLifts/LatticeBasics/CompactBound.lean}{GlobalTheory.ThK.Dbl.Hgt.scale_ad}
\leanref{Development/WeilBound/ThetaLifts/LatticeBasics/HgtCount.lean}{GlobalTheory.ThK.Dbl.Hgt.count_grp}
\leanref{Development/WeilBound/ThetaLifts/LatticeBasics/HgtCount.lean}{GlobalTheory.ThK.Dbl.Hgt.exists_cover}
\begin{proof}
This is a volume computation. Let $\rho\in\mathbf A^\times$ have the
components $\tfrac12$ at the infinite places and $1$ at the finite places, and
let $\delta$ be the factor by which multiplication by $\rho^{-1}$ scales a Haar
measure of $\mathbf A$. There are $i_0\geq0$ and a compact set
$C_0\subset\mathbf A$ with $m\rho^{2i}\in C_0$ for all integers
$0\leq m\leq2^i$, such that every $u$ with $N(u)\in\Omega_0$ has
$|u|\leq2^{i_0}$, and, if moreover $|u|\geq2^{-j}$, then
$m\rho^{2(j+i_0)}=x\,N(u)$ for some integer $0<m\leq2^{j+i_0}$ and some
$x\in C_0$. So $H_j$ is contained in the union over these $m$ of the sets of the
$g$ with $\lambda g,\lambda g^{-1}\in C_0\Omega$, $\lambda=m\rho^{2(j+i_0)}$.
Covering a compact subset of $\mathbf A$ by translates of a neighbourhood of
$0$ scaled by $\lambda^2$ bounds the measure of each of these sets by
$(A\delta^{4(j+i_0)})^{q}\mu(U)$, with $q$ the number of entries, a fixed
compact set $U\subset H$ and a constant $A$. Summing over $m\leq2^{j+i_0}$
gives the bound with $b=2\delta^{4q}$ and $R=2^{i_0}$.
\end{proof}

\begin{lemma}\label{asm-theta-agree-doubling-integrable}
Let $B$ be a theta base for a CM field $E$
(Definition~\ref{asm-theta-agree-models}(1)), let $Y$ be a lattice space that
is Hermitian or skew-Hermitian of some signature
(Definition~\ref{asm-theta-agree-doubling}(1)), and let $K$, $a$ and $\iota$ be
the lattice stabilizer of the doubled space of $Y$, its height function and the
doubling embedding (Lemma~\ref{asm-theta-agree-doubling-stabilizer}). Then for
every Haar measure $\mu$ on $U_Y$ there is $t_0\in\mathbf R$ such that
$h\mapsto a(\iota(1,h))^t$ is $\mu$-integrable on $U_Y$ for every $t\geq t_0$.
\end{lemma}
\leanref{Development/WeilBound/ThetaLifts/LatticeBasics/PairsStable.lean}{GlobalTheory.ThK.Dbl.HgtInt}
\leanref{Development/WeilBound/ThetaLifts/LatticeBasics/HgtIntOf.lean}{GlobalTheory.ThK.Dbl.hgtInt_of}
\leanref{Development/WeilBound/ThetaLifts/LatticeBasics/HgtIntOf.lean}{GlobalTheory.ThK.Dbl.Hgt.graph_alg}
\leanref{Development/WeilBound/ThetaLifts/LatticeBasics/HgtCount.lean}{GlobalTheory.ThK.Dbl.Hgt.integrable_of_count}
\begin{proof}
Let $n$ be the rank of $Y$ and $J$ its matrix, and $N=N_{E/F}$. For $k\in G_Y$
write $k^{-1}\Delta=\left(\begin{smallmatrix}T(k)\\T'(k)\end{smallmatrix}\right)$
with $n\times n$ blocks $T(k)$ and $T'(k)$; these are continuous in $k$.

\emph{The key identity.} Let $g\in U_Y$, and write $\iota(1,g)=pk$ with
$p\in P_Y$ and $k\in K$ (Lemma~\ref{asm-theta-agree-doubling-stabilizer}(2)),
and $t=t(p)$. Then $k^{-1}=\iota(1,g^{-1})p$, so
$k^{-1}\Delta=\iota(1,g^{-1})\Delta t$, that is, $T(k)=t$ and
$T'(k)=g^{-1}t$; and $a(\iota(1,g))=|\det t|$. Taking determinants in
$g^{-1\,\mathsf T}Jc(g^{-1})=J$ gives $\det g^{-1}\,c(\det g^{-1})=1$, as
$\det J\neq0$. With $u=\det t$, so that $N(u)=u\,c(u)$, Cramer's rule gives
\[
N(u)\,g=c(\det T'(k))\,T(k)\operatorname{adj}T'(k),\qquad
N(u)\,g^{-1}=c(\det T(k))\,T'(k)\operatorname{adj}T(k).
\]

\emph{The bound.} Let $\Omega\subset M_n(\mathbf A_E)$ be the union of the
images of the compact group $K$
(Lemma~\ref{asm-theta-agree-doubling-stabilizer}(1)) under the two continuous
maps on the right, and $\Omega_0\subset\mathbf A$ the image of $K$ under
$k\mapsto N(\det T(k))$. By the identity, every $g\in U_Y$ has $u$ with
$N(u)\in\Omega_0$, $N(u)g\in\Omega$, $N(u)g^{-1}\in\Omega$ and
$a(\iota(1,g))=|u|$. As $U_Y$ is closed in $GL_n(\mathbf A_E)$,
Lemma~\ref{asm-theta-agree-doubling-count} gives $R\geq1$, $C$ and $b$ with
$a(\iota(1,g))\leq R$ for all $g$ and
$\mu\{g:a(\iota(1,g))\geq2^{-j}\}\leq Cb^j$ for all $j$.

\emph{Integrability.} The function $f(g)=a(\iota(1,g))$ is continuous and
positive. Choose an integer $t_0\geq0$ with $2^{-t_0}(b+1)<1$, and let
$t\geq t_0$. The sets $P_j=\{g:2^{-j}\leq f(g)\leq2R\,2^{-j}\}$, $j\geq0$, cover
$U_Y$, and $f^t\leq(2R)^t2^{-jt}$ on $P_j$, so
\[
\int_{U_Y}f^t\,d\mu\leq\sum_{j\geq0}(2R)^t\,C\,(2^{-t}b)^j<\infty,
\]
as $2^{-t}b<1$.
\end{proof}
% Lean: hgtInt_of: key (with e1, e2), Hgt.graph_alg, count_ad, integrable_of_count.

\begin{lemma}\label{asm-theta-agree-doubling-frame}
Let $B$ be a theta base for a CM field $E$
(Definition~\ref{asm-theta-agree-models}(1)), let $Y$ be a lattice space that
is Hermitian or skew-Hermitian of some signature
(Definition~\ref{asm-theta-agree-doubling}(1)), and let $K$ be the lattice
stabilizer of the doubled space of $Y$
(Definition~\ref{asm-theta-agree-nonvanishing-doubling}(1)). Then there are a
frame $b$ of $Y$, with standard compact subgroup $K_b$
(Definition~\ref{asm-theta-agree-doubling}(3)), and $p\in P_Y$ such that
$K=pK_bp^{-1}$.
\end{lemma}
\leanref{Development/WeilBound/ThetaLifts/LatticeBasics/HypothesesE.lean}{GlobalTheory.ThK.Dbl.FrameConj}
\leanref{Development/WeilBound/ThetaLifts/LatticeBasics/FrameConj.lean}{GlobalTheory.ThK.Dbl.frameConj}
\leanref{Development/WeilBound/ThetaLifts/LatticeBasics/FrameA.lean}{GlobalTheory.ThK.Dbl.Frame.frame_gram}
\leanref{Development/WeilBound/ThetaLifts/LatticeBasics/FrameA.lean}{GlobalTheory.ThK.Dbl.Frame.exists_arch}
\leanref{Development/WeilBound/ThetaLifts/LatticeBasics/FrameB.lean}{GlobalTheory.ThK.Dbl.Frame.exists_ad}
\leanref{Development/WeilBound/ThetaLifts/LatticeBasics/FrameConj.lean}{GlobalTheory.ThK.Dbl.Frame.keeps_iff}
\begin{proof}
Let $n$, $J$, $L$ and $\theta_w$ be the rank, the matrix, the lattice and the
matrices at infinity of $Y$, and let $\epsilon=1$ if $Y$ is Hermitian and
$\epsilon=-1$ if it is skew-Hermitian, so that $c(J)^{\mathsf T}=\epsilon J$.
For a ring $S$ with an involution $c$, $g,T\in GL_n(S)$ and $\alpha\in S$, put
\[
b(g,\alpha,T)=\begin{pmatrix}g&0\\0&g\end{pmatrix}\begin{pmatrix}1&\alpha T\\1&(\alpha-1)T\end{pmatrix}.
\]
If $\alpha+c(\alpha)=1$ and $g^{\mathsf T}Jc(g)c(T)=1$, then
$b^{\mathsf T}J^\square c(b)=\left(\begin{smallmatrix}0&1\\\epsilon&0\end{smallmatrix}\right)$
for $b=b(g,\alpha,T)$, and the first $n$ columns of $b$ are $\Delta g$; both
are direct computations with $2\times2$ block matrices.

\emph{The choices.} There is $g\in GL_n(\mathbf A_E)$ such that: some
$a\in GL_n(\mathbf A_E)$ has $y\in\mathcal O_{E,v}^n$ if and only if $a_vy\in L_v$, for
every finite place $v$, with $a_v$ a unit multiple of $g_v$; at every finite
$v$, the Gram matrix $g_v^{\mathsf T}Jc(g_v)$ and its inverse have integral
entries; and at every infinite $w$, $\tau_w(g_w)$ is a basis $c_w$ chosen for
$\theta_w$ as below. Let $T\in GL_n(\mathbf A_E)$ with
$g^{\mathsf T}Jc(g)c(T)=1$, and $T_0\in GL_n(E)$ with $Jc(T_0)=1$. Let $t_1$ be
an element of $\mathcal O_E$ with $t_1+c(t_1)=1$, which exists by the
conditions on a good base, and let $\alpha\in\mathbf A_E$ have the components
$\tfrac12$ at the infinite places and $t_1$ at the finite places; then
$\alpha+c(\alpha)=1$. Put $b=b(1,\tfrac12,T_0)$, a matrix over $E$, and
$b'=b(g,\alpha,T)$. By the computation, $b$ is a frame of $Y$, and $b$ and $b'$
have the same Gram matrix for $J^\square$, so $p=b'b^{-1}\in G_Y$; and
$p\Delta=b'\left(\begin{smallmatrix}1\\0\end{smallmatrix}\right)=\Delta g$, so
$p\in P_Y$ with $t(p)=g$.

\emph{$K=pK_bp^{-1}$.} For $k\in G_Y$, $b^{-1}(p^{-1}kp)b=b'^{-1}kb'$. So
$p^{-1}kp\in K_b$ if and only if $\tau_w(b'^{-1}kb')$ is unitary for every
infinite $w$ and the finite components of $b'^{-1}kb'$ and $b'^{-1}k^{-1}b'$
are integral. At a finite place $v$: $k$ and $k^{-1}$ map
$\hat L_Y\times\hat L_Y$ into itself if and only if
$\operatorname{diag}(a,a)^{-1}k^{\pm1}\operatorname{diag}(a,a)$ have integral
finite components; $a_v$ is a unit multiple of $g_v$; and the second factor of
$b'$ and its inverse are integral at $v$, because $\alpha_v$ is integral and
$T_v$ and $T_v^{-1}$ are, $T_v$ being the inverse of the conjugate of an
integral Gram matrix with integral inverse. So the two integrality conditions
agree. At an infinite place $w$: let $c=1$ if $Y$ is Hermitian and $c=i$
otherwise, so that $M_w=c\,\tau_w(J)$ is Hermitian with the Cartan involution
$\theta_w$. Choose $c_w$ so that, in the basis of its columns, $M_w$ is
diagonal with entries $\pm1$ and a nonzero multiple of the form
$(x,y)\mapsto(\theta_wx)^{\mathsf T}M_w\bar y$ is the standard Hermitian form;
then for every $x$ with
$x^{\mathsf T}\tau_w(J)^\square\bar x=\tau_w(J)^\square$,
$\tau_w(b')^{-1}x\,\tau_w(b')$ is unitary if and only if $x$ commutes with
$\operatorname{diag}(\theta_w,-\theta_w)$. Hence $p^{-1}kp\in K_b$ if and only if
$k\in K$.
\end{proof}
% Lean: frameConj, with Frame.two g = diag(g, g) and Frame.wit alpha T the second factor;
% frame_gram and frame_diag are the computation, exists_dual gives T and T_0,
% exists_int_add_conj_eq_one_of_base_good gives t_1, exists_ad (with exists_loc) gives a and g
% (Uni is the integrality of the Gram matrix and its inverse), par_of gives p in P_Y. The
% finite places are keeps_iff, int_dual, int_wit and int_conj; the infinite places are
% Frame.exists_arch, through Hgt.exists_cartan_frame and Hgt.unitary_iff_commute.

\begin{lemma}\label{asm-theta-agree-doubling-continuation}
Assume the meromorphic continuation of Siegel Eisenstein series in the form of
Theorem~\ref{asm-theta-agree-eisenstein}. Let $B$ be a theta base for a CM
field $E$ (Definition~\ref{asm-theta-agree-models}(1)), let $Y$ be a lattice
space of rank $n\geq1$ that is Hermitian or skew-Hermitian of some signature
(Definition~\ref{asm-theta-agree-doubling}(1)), and let $K$ and $a$ be the
lattice stabilizer of the doubled space of $Y$ and its height function
(Lemma~\ref{asm-theta-agree-doubling-stabilizer}). Let $\chi$ be a Hecke
character of $E$ with $\chi|_{\mathbf A^\times}=\eta^c$ for some $c\geq0$, let
$s_1\in\mathbf C$, and let $f$ be a section of degree $s_1$ for $(K,\chi)$, with
flat family $(f_s)$ (Definition~\ref{asm-theta-agree-doubling}(4)). Then there
are an open set $U\subset\mathbf C$ and a function $\hat E$ on
$\mathbf C\times G_Y$ such that:
\begin{enumerate}
\item $U$ is connected, and every $z\in\mathbf C$ has a punctured neighbourhood
  contained in $U$;
\item there is $t\in\mathbf R$ such that every $s$ with $\operatorname{Re}s>t$
  lies in $U$ and, for every $x\in G_Y$, the Eisenstein sum of $f_s$ at $x$
  converges unconditionally to $\hat E(s,x)$;
\item $s\mapsto\hat E(s,x)$ is holomorphic on $U$ for every $x$, and $\hat E$ is
  continuous on $U\times G_Y$;
\item every $s\in U$ has a neighbourhood $D$ such that, for some $C$ and $N$,
  $|\hat E(s',x)|\leq C\operatorname{ht}(x)^N$ for all $s'\in D$ and
  $x\in G_Y$;
\item for every $s_0\in\mathbf C$ there are $k\geq0$, $r>0$ and a function $H$ on
  $\{|s-s_0|<r\}\times G_Y$ with $H(s,x)=(s-s_0)^k\hat E(s,x)$ for
  $0<|s-s_0|<r$, such that $s\mapsto H(s,x)$ is holomorphic for every $x$, $H$
  is continuous, and $|H(s,x)|\leq C\operatorname{ht}(x)^N$ for some $C$ and
  $N$.
\end{enumerate}
\end{lemma}
\leanref{Development/WeilBound/ThetaLifts/LatticeBasics/Hypotheses2.lean}{GlobalTheory.ThK.Dbl.EisContinues}
\leanref{Development/WeilBound/ThetaLifts/LatticeBasics/Hypotheses2.lean}{GlobalTheory.ThK.Dbl.IsEisCont}
\leanref{Development/WeilBound/ThetaLifts/LatticeBasics/Hypotheses2.lean}{GlobalTheory.ThK.Dbl.IsRegAt}
\leanref{Development/WeilBound/ThetaLifts/LatticeBasics/EisCont.lean}{GlobalTheory.ThK.Dbl.eisContinues_of}
\leanref{Development/WeilBound/ThetaLifts/LatticeBasics/EisCont.lean}{GlobalTheory.ThK.Dbl.reg_of_weak}
\leanref{Development/WeilBound/ThetaLifts/LatticeBasics/EisCont.lean}{GlobalTheory.ThK.Dbl.eisCont_of_reg}
\leanref{Development/WeilBound/ThetaLifts/LatticeBasics/RegOfBound.lean}{GlobalTheory.ThK.Dbl.An.unif_bound}
\leanref{Development/WeilBound/ThetaLifts/LatticeBasics/RegOfBound.lean}{GlobalTheory.ThK.Dbl.An.reg_of_bound}
\begin{proof}
Call a function $G$ on $\mathbf C\times G_Y$ \emph{regular} at $s_0$ if on some
neighbourhood $D$ of $s_0$ the functions $s\mapsto G(s,x)$ are holomorphic, $G$
is continuous on $D\times G_Y$, and $|G(s,x)|\leq C\operatorname{ht}(x)^N$ for
$s\in D$ and some $C$ and $N$. Heights are at least $1$ and continuous.

\emph{Step 1: conjugation.} By Lemma~\ref{asm-theta-agree-doubling-frame} there
are a frame $b$ of $Y$ and $p\in P_Y$ with $K=pK_bp^{-1}$; put
$d=|\det t(p)|>0$. Then $a_0(x)=d\,a(xp^{-1})$ is a height function for $K_b$:
if $xp^{-1}=qk$ with $q\in P_Y$ and $k\in K$, then
$x=(qp)(p^{-1}kp)$; and for $q\in P_Y$ and $k_0\in K_b$,
$a_0(qk_0)=d\,a(qp^{-1}\cdot pk_0p^{-1})=d\,|\det t(qp^{-1})|=|\det t(q)|$. The
function $f^0(x)=f(xp^{-1})$ is a section of degree $s_1$ for $(K_b,\chi)$: it
is continuous, has the same law under $P_Y$ on the left, and its right
translates by $K_b$ are the right translates of $f$ by $K$ composed with
$x\mapsto xp^{-1}$. Its flat family for $a_0$ is
$a_0(x)^{s-s_1}f^0(x)=d^{s-s_1}f_s(xp^{-1})$.

\emph{Step 2: the half-plane.} Theorem~\ref{asm-theta-agree-eisenstein}, for $b$,
$K_b$, $a_0$, $\chi$ and $f^0$, gives $\hat E_0$ and $t$. Put
$\hat E(s,x)=d^{s_1-s}\hat E_0(s,xp)$. For $\operatorname{Re}s>t$ the
Eisenstein sum of $f_s$ at $x$ is $d^{s_1-s}$ times the Eisenstein sum of the
flat family of $f^0$ at $xp$, so it converges unconditionally to
$\hat E(s,x)$, and $s\mapsto\hat E(s,x)$ is holomorphic there.

\emph{Step 3: near a point.} Let $s_0\in\mathbf C$, and let $k$, a
neighbourhood $U_0$ of $s_0$ and the functions $H_s$ of moderate growth be
given by Theorem~\ref{asm-theta-agree-eisenstein}(2) for $\hat E_0$. Let
$D\subset U_0$ be a compact neighbourhood of $s_0$. For every growth functional
$\mu$, $s\mapsto\mu(H_s)$ is holomorphic, hence bounded, on $D$. This gives
$C$ and $N$ with $|H_s(x)|\leq C\operatorname{ht}(x)^N$ for $s\in D$ and all
$x$, by a uniform boundedness argument: otherwise there are $s_m\in D$ and
$x_m$ with $|H_{s_m}(x_m)|>m2^m\operatorname{ht}(x_m)^m$; for a bounded
sequence $(c_m)$ the form $f\mapsto\sum_mc_m2^{-m}\operatorname{ht}(x_m)^{-m}f(x_m)$
is a growth functional, so the Banach-Steinhaus theorem on the bounded
sequences bounds the norms of $(c_m)\mapsto\sum_mc_m2^{-m}\operatorname{ht}(x_m)^{-m}H_{s_l}(x_m)$
uniformly in $l$, which fails at the sequence with a single $1$ in place $l$.
Evaluation at $x$ is a growth functional, so $s\mapsto H_s(x)$ is holomorphic
on $D$, and each $H_s$ is continuous. With the uniform bound, Cauchy's
estimate gives $|H_z(x)-H_y(x)|\leq2C\operatorname{ht}(x)^Nr^{-1}|z-y|$ on a
disc of radius $r$ in $D$, so $(s,x)\mapsto H_s(x)$ is continuous near $s_0$.
So $(s,x)\mapsto H_s(x)$ is regular at $s_0$, and so is
$(s,x)\mapsto d^{s_1-s}H_s(xp)$, as
$\operatorname{ht}(xp)\leq\operatorname{ht}(x)\operatorname{ht}(p)$; and it
equals $(s-s_0)^k\hat E(s,x)$ for $s\neq s_0$ near $s_0$. This gives (5).

\emph{Step 4: the set $U$.} Let $U$ be the set of the $s$ at which $\hat E$ is
regular; it is open. Let $s_0\in\mathbf C$, and $H$ as in Step 3. Near $s_0$
and away from it, $\hat E=(s-s_0)^{-k}H$ with $(s-s_0)^{-k}$ holomorphic, so
$\hat E$ is regular at every $s\neq s_0$ near $s_0$. So every point has a
punctured neighbourhood in $U$; hence $\mathbf C\setminus U$ is discrete,
hence countable, and $U$ is connected, as the complement of a countable set in
the plane. If $\operatorname{Re}s_0>t$, then $\hat E(\cdot,x)$ is holomorphic
near $s_0$, and on a small circle around $s_0$ it is bounded by
$Cr^{-k}\operatorname{ht}(x)^N$; by the maximum principle the same bound holds
inside, and continuity at $(s_0,x)$ follows in the same way, so $s_0\in U$.
This gives (1) and (2), and (3) and (4) hold because $\hat E$ is regular at
every point of $U$.
\end{proof}
% Lean: eisContinues_of EisBL FrameConj (ha0 and hf0 are Step 1); reg_of_weak (with
% An.unif_bound, which applies Banach-Steinhaus to the bilinear form B, and An.reg_of_bound)
% and Reg.mul are Step 3; eisCont_of_reg is Step 4, with An.Reg.punct, An.preconn and
% An.Reg.remove. The set U is {s | An.Reg ht E-hat s}.

\begin{lemma}\label{asm-theta-agree-doubling-datum}
Let $(X,W,M,\chi_X,\chi_W,\omega)$ be a Weil pair for a theta base $B$ of a CM
field $E$ (Definition~\ref{asm-theta-agree-triple}(1)), with $X$ and $W$ of
ranks $n$ and $d$, let $e$ be the isomorphism of $M$, put
$I=\{1,\ldots,n\}\times\{1,\ldots,d\}$, let $I\sqcup I$ be the disjoint union
of two copies of $I$, and write
$(\phi_1\boxtimes\phi_2)(x,y)=\phi_1(x)\phi_2(y)$ for functions on
$\mathbf A^I$. Let $\iota$ denote the doubling embeddings of $X$ and of $W$
(Definition~\ref{asm-theta-agree-nonvanishing-doubling}(1)).
\begin{enumerate}
\item There is a doubling datum $(I\sqcup I,e^\square,\omega_2,S,l)$ for $(X,W)$
  with character $\chi_X$ (Definition~\ref{asm-theta-agree-doubling}(6)), with
  $e^\square(w_1\ w_2)=\bigl((e_1(w_1),e_1(w_2)),(e_2(w_1),-e_2(w_2))\bigr)$,
  such that for all Schwartz-Bruhat functions $\Phi_1,\Phi_2$ on
  $\mathbf A^I$, all $g\in U_X$, all $h,h'\in U_W$ and $x=\iota(h,h')$:
  $\Phi_1\boxtimes\bar\Phi_2$ is a Schwartz-Bruhat function,
  \[
  \chi_X(\det h')^{-1}\,\omega_2(x)(\Phi_1\boxtimes\bar\Phi_2)=(\omega(1,h)\Phi_1)\boxtimes\overline{\omega(1,h')\Phi_2},
  \]
  and
  \[
  \chi_X(\det h')^{-1}\,\theta\bigl(\omega_2(x)S(g)(\Phi_1\boxtimes\bar\Phi_2)\bigr)=\theta_\omega(\Phi_1)(g,h)\,\overline{\theta_\omega(\Phi_2)(g,h')}.
  \]
\item There is a doubling datum $(I\sqcup I,e^\square,\omega_2,S,l)$ for $(W,X)$
  with character $\chi_W$, with
  $e^\square(w_1\ w_2)=\bigl((e_1(w_1^{\mathsf T}),e_1(w_2^{\mathsf T})),(e_2(w_1^{\mathsf T}),-e_2(w_2^{\mathsf T}))\bigr)$,
  such that for all Schwartz-Bruhat $\Phi_1,\Phi_2$, all $h\in U_W$, all
  $g,g'\in U_X$ and $x=\iota(g,g')$:
  \[
  \chi_W(\det g')^{-1}\,\omega_2(x)(\Phi_1\boxtimes\bar\Phi_2)=(\omega(g,1)\Phi_1)\boxtimes\overline{\omega(g',1)\Phi_2},
  \]
  \[
  \chi_W(\det g')^{-1}\,\theta\bigl(\omega_2(x)S(h)(\Phi_1\boxtimes\bar\Phi_2)\bigr)=\theta_\omega(\Phi_1)(g,h)\,\overline{\theta_\omega(\Phi_2)(g',h)}.
  \]
\item In (1) and (2), for every Haar measure $\nu$ on $\mathbf A^I$ there is
  $c'>0$ with $l(\Psi_1\boxtimes\bar\Psi_2)=c'\int\Psi_1\bar\Psi_2\,d\nu$ for all
  Schwartz-Bruhat $\Psi_1,\Psi_2$ on $\mathbf A^I$.
\end{enumerate}
\end{lemma}
\leanref{Development/WeilBound/ThetaLifts/LatticeBasics/PairsStable.lean}{GlobalTheory.ThK.Dbl.DblW}
\leanref{Development/WeilBound/ThetaLifts/LatticeBasics/BoxSecX.lean}{GlobalTheory.ThK.Dbl.DblX}
\leanref{Development/WeilBound/ThetaLifts/LatticeBasics/RatEquiv.lean}{GlobalTheory.ThK.Dbl.dblW}
\leanref{Development/WeilBound/ThetaLifts/LatticeBasics/RatEquiv.lean}{GlobalTheory.ThK.Dbl.dblX}
\leanref{Development/WeilBound/ThetaLifts/LatticeBasics/RatEquiv.lean}{GlobalTheory.ThK.Dbl.dbl_main}
\leanref{Development/WeilBound/ThetaLifts/LatticeBasics/RatEquiv.lean}{GlobalTheory.ThK.Dbl.dbl_box}
\leanref{Development/WeilBound/ThetaLifts/LatticeBasics/RatEquiv.lean}{GlobalTheory.ThK.Dbl.dbl_emb}
\leanref{Development/WeilBound/ThetaLifts/LatticeBasics/ThetaBox.lean}{GlobalTheory.ThK.Dbl.lD}
\leanref{Development/WeilBound/ThetaLifts/LatticeBasics/ThetaBox.lean}{GlobalTheory.ThK.Dbl.lD_inj}
\leanref{Development/WeilBound/ThetaLifts/LatticeBasics/ThetaBox.lean}{GlobalTheory.ThK.Dbl.lD_haar}
\begin{proof}
We prove (1) and (3); (2) is the same argument with $X$ and $W$ exchanged,
using condition (b) of Definition~\ref{asm-theta-agree-models}(6) in the place
of (a). Let $\rho$ be the Schrödinger operators of $M$ and $\rho^\square$ those
of $e^\square$, and let $\omega^\square$ and $\omega'^\square$ be the operators
on the doubled groups $G_W$ and $G_X$ given by conditions (a) and (b) of
Definition~\ref{asm-theta-agree-models}(6). For $\Phi$ Schwartz-Bruhat on
$\mathbf A^{I\sqcup I}$ let $l(\Phi)=\int\Phi(u,u)\,du$, for the self-dual
measure of $\mathbf A^I$, and extend $l$ linearly to all functions. Put
$\omega_2=\omega^\square$ and
$S(g)=\chi_W(\det g)^{-1}\omega'^\square(\iota(g,g))$.

\emph{Properties of $l$.} It is admissible. Writing $\psi$ as
$x\mapsto\psi_0(\xi x)$ with $\psi_0$ the standard character and
$\xi\in F^\times$, a direct computation gives, for $a,a',b,b'\in\mathbf A^I$,
\[
l\bigl(\rho^\square((a,a'),(b,b'))\Phi\bigr)=\psi\bigl(\tfrac12(a\cdot b+a'\cdot b')-a'\cdot(b+b')\bigr)\,\hat\Phi_{a-a'}\bigl(\xi(b+b')\bigr),
\]
where $\hat\Phi_c$ is the Fourier transform of $u\mapsto\Phi(u+c,u)$. Three
consequences: (i) $l$ is invariant under the $\rho^\square((a,a),(b,-b))$, which
are the operators $\rho^\square(z\Delta^{\mathsf T})$; (ii) if
$l(\rho^\square((a,0),(b,0))\Phi)=0$ for all $a$ and $b$, then all the
$\hat\Phi_a$ vanish, so $\Phi=0$ by Fourier inversion; (iii) the sum of the
$l(\rho^\square((y,-y))\Phi)$ over $y\in M_{n\times d}(E)$ is, after the change
of variables given by $e^\square$ on rational points (which maps
$M_{n\times 2d}(E)$ onto $F^{I\sqcup I}\times F^{I\sqcup I}$, as $M$ is
regular and $\dim_FE=2$), Poisson summation for the slices, with sum
$\theta(\Phi)$. Finally
$l(\Psi_1\boxtimes\bar\Psi_2)=\int\Psi_1\bar\Psi_2$ for the self-dual measure,
which gives (3), and the unitarity of $\omega(1,h)$ gives, by polarization,
$l(\omega(1,h)\phi\boxtimes\overline{\omega(1,h)\phi'})=l(\phi\boxtimes\bar\phi')$.

\emph{The doubling datum.} Condition (a) of
Definition~\ref{asm-theta-agree-doubling}(6) is that of $e^\square$, with the
rationality just used. Condition (b) for $\omega_2$ is condition (a) of the
Weil representation. For (c): $S(g)$ preserves the Schwartz-Bruhat functions
and the admissible forms, it is unitary, it is continuous in $g$ against
admissible forms, and
$S(g)\rho^\square(w)=\rho^\square(gw)S(g)$, because $\rho^\square$ of a pair of
columns is $\rho'^\square$ of the transposed pair for the doubled coordinates
of condition (b). For (d): $l$ is admissible and invariant under the
$\rho^\square(z\Delta^{\mathsf T})$ by (i); the displayed rule of condition (b)
for $\omega'^\square$ at $\iota(g,g)\in P_X$, with $t=g$ and $|\det g|=1$, gives
$l(\omega'^\square(\iota(g,g))\Phi)=\chi_W(\det g)\,l(\Phi)$, that is,
$l(S(g)\Phi)=l(\Phi)$; and the theta condition is (iii).

\emph{The box identity.} Let $h\in U_W$ and $\phi,\phi'$ Schwartz-Bruhat. Then
$\omega^\square(\iota(h,h))(\phi\boxtimes\bar\phi')=\chi_X(\det h)\,\omega(1,h)\phi\boxtimes\overline{\omega(1,h)\phi'}$.
Both sides are Schwartz-Bruhat, so by (ii) it suffices that $l$ agrees on
$\rho^\square((a,0),(b,0))$ of both. Write $((a,0),(b,0))=e^\square(wh^{\mathsf T},0)$.
On box functions $\rho^\square((u,0))$ acts on the first factor by $\rho(u)$.
Using $\omega^\square(x)\rho^\square(w)=\rho^\square(wx^{\mathsf T})\omega^\square(x)$,
the rule $l(\omega^\square(\iota(h,h))\Phi)=\chi_X(\det h)l(\Phi)$ (from
condition (b), as above), $\omega(1,h)\rho(w)=\rho(wh^{\mathsf T})\omega(1,h)$
and the unitarity of $\omega(1,h)$ under $l$, both sides become
$\chi_X(\det h)\,l(\rho(w)\phi\boxtimes\bar\phi')$. For
$x=\iota(h,h')=\iota(hh'^{-1},1)\iota(h',h')$, condition (a) gives
$\omega^\square(\iota(k,1))(\phi_1\boxtimes\phi_2)=\omega(1,k)\phi_1\boxtimes\phi_2$,
so
$\chi_X(\det h')^{-1}\omega_2(x)(\phi\boxtimes\bar\phi')=\omega(1,h)\phi\boxtimes\overline{\omega(1,h')\phi'}$,
the first display. In the same way, with condition (b),
$S(g)(\phi\boxtimes\bar\phi')=\omega(g,1)\phi\boxtimes\overline{\omega(g,1)\phi'}$.

\emph{The seesaw identity.} With $\phi_i=\omega(g,1)\Phi_i$, the first display
and $\omega(1,h)\omega(g,1)=\omega(g,h)$ give
\[
\chi_X(\det h')^{-1}\theta\bigl(\omega_2(x)S(g)(\Phi_1\boxtimes\bar\Phi_2)\bigr)=\theta\bigl(\omega(g,h)\Phi_1\boxtimes\overline{\omega(g,h')\Phi_2}\bigr)=\theta_\omega(\Phi_1)(g,h)\,\overline{\theta_\omega(\Phi_2)(g,h')},
\]
since the theta sum of a box function over $F^{I\sqcup I}=F^I\times F^I$ is the
product of the theta sums.
\end{proof}
% Lean: dblW and dblX from dbl_main (rat_equiv for the rational points, lD_kill, dbl_fix
% with modu_det, lD_theta, lD_adm, and S g = detInv chi_W g * omd' (dgU g)), dbl_box (lD_inj,
% lD_unit, lD_schr) and dbl_emb, and theta_box. (3) is lD_haar (the self-dual measure is
% adeleMeasurePi, c' its Haar scalar factor against nu). DblX is (2), with IsDblCoordX the
% transposed coordinates. About 250 lines with ThetaBox.lean.

\begin{lemma}\label{asm-theta-agree-doubling-level-gaussian}
Let $V$ be a finite-dimensional real vector space, and for $a\in V$ let
$\partial_a$ be the derivative in the direction $a$.
\begin{enumerate}
\item Let $f\neq0$ be a Schwartz function on $V$, and
  $m\colon V\times V\to\mathbf C$ an $\mathbf R$-bilinear map with
  $\partial_af(y)=m(a,y)f(y)$ for all $a,y\in V$. Then
  $f(y)=f(0)\exp\bigl(m(y,y)/2\bigr)$ for every $y$, $f$ vanishes nowhere, $m$
  is symmetric, and $\operatorname{Re}m(a,a)<0$ for every $a\neq0$.
\item Let $f$ and $g$ be as in (1), with maps $m$ and $m'$. Let $a,b\in V$, and
  let $L$ and $L'$ be $\mathbf R$-linear forms on $V$ with
  \[
  \partial_bf-2\pi i\,L'f=i\bigl(\partial_af-2\pi i\,Lf\bigr),\qquad
  -\partial_bg+2\pi i\,L'g=i\bigl(\partial_ag-2\pi i\,Lg\bigr).
  \]
  Then $a=b=0$ and $L=0$.
\item Let $T$ be an $\mathbf R$-bilinear form on $V$ such that $b=0$ whenever
  $T(y,b)=0$ for all $y$. For $p=(p_1,p_2)\in V\times V$ and a Schwartz
  function $f$ put $D_pf=\partial_{p_1}f-2\pi i\,T(\cdot,p_2)f$, and let $J$ be
  an $\mathbf R$-linear map of $V\times V$. If $f\neq0$ is a Schwartz function
  with $D_{Jp}f=i\,D_pf$ for every $p$, then there is an $\mathbf R$-bilinear
  map $m\colon V\times V\to\mathbf C$ with $\partial_af(y)=m(a,y)f(y)$ for all
  $a,y\in V$.
\end{enumerate}
\end{lemma}
\leanref{Development/AutomorphicSpectrum/Proofs/TwistedProofs/SchwartzLogDerivLinearGaussian.lean}{GlobalTheory.CompanionTheory.Twist.schwartz_logDeriv_linear_formula}
\leanref{Development/AutomorphicSpectrum/Proofs/TwistedProofs/SchwartzLogDerivLinearGaussian.lean}{GlobalTheory.CompanionTheory.Twist.schwartz_logDeriv_linear_gaussian}
\leanref{Development/WeilBound/ThetaLifts/LatticeBasics/BoxFinA.lean}{GlobalTheory.ThK.Dbl.Box.Gau}
\leanref{Development/WeilBound/ThetaLifts/LatticeBasics/BoxFinA.lean}{GlobalTheory.ThK.Dbl.Box.gau_of}
\leanref{Development/WeilBound/ThetaLifts/LatticeBasics/BoxFinA.lean}{GlobalTheory.ThK.Dbl.Box.opp}
\leanref{Development/AutomorphicSpectrum/Proofs/TwistedProofs/EqZeroOfForallNhds.lean}{GlobalTheory.CompanionTheory.Twist.FockLogDeriv.main}
\leanref{Development/AutomorphicSpectrum/Proofs/TwistedProofs/EqZeroOfForallNhds.lean}{GlobalTheory.CompanionTheory.Twist.FockLogDeriv.eq_zero_of_forall_nhds}
\begin{proof}
(1) Let $y\in V$ and $\varphi(t)=f(ty)$ for $t\in\mathbf R$. Then
$\varphi'(t)=\partial_yf(ty)=m(y,ty)f(ty)=t\,m(y,y)\varphi(t)$, so
$t\mapsto\varphi(t)\exp(-t^2m(y,y)/2)$ has derivative $0$, and
$f(y)=\varphi(1)=f(0)\exp(m(y,y)/2)$. As $f\neq0$, $f(0)\neq0$, so $f$ vanishes
nowhere. Let $a\neq0$. If $\operatorname{Re}m(a,a)\geq0$, then
$|f(ta)|=|f(0)|\exp\bigl(t^2\operatorname{Re}m(a,a)/2\bigr)\geq|f(0)|>0$ for
every $t\in\mathbf R$, while $|t|\,\|a\|\,|f(ta)|$ is bounded in $t$, as $f$ is
a Schwartz function, for a norm $\|\cdot\|$ on $V$. So
$\operatorname{Re}m(a,a)<0$. Differentiating the formula for $f$ in
the direction $a$ gives $\partial_af(y)=\tfrac12\bigl(m(a,y)+m(y,a)\bigr)f(y)$;
comparing with $\partial_af(y)=m(a,y)f(y)$ and dividing by $f(y)\neq0$ gives
$m(a,y)=m(y,a)$.

(2) By (1), $f$ and $g$ vanish nowhere. Dividing the two relations at $y$ by
$f(y)$ and by $g(y)$ gives
$m(b,y)-2\pi i\,L'(y)=i\,m(a,y)+2\pi L(y)$ and
$-m'(b,y)+2\pi i\,L'(y)=i\,m'(a,y)+2\pi L(y)$. Their real and imaginary parts
are
\[
\begin{aligned}
&\text{(i)}\ \operatorname{Re}m(b,y)=-\operatorname{Im}m(a,y)+2\pi L(y),&\qquad
&\text{(ii)}\ \operatorname{Im}m(b,y)-2\pi L'(y)=\operatorname{Re}m(a,y),\\
&\text{(iii)}\ -\operatorname{Re}m'(b,y)=-\operatorname{Im}m'(a,y)+2\pi L(y),&\qquad
&\text{(iv)}\ -\operatorname{Im}m'(b,y)+2\pi L'(y)=\operatorname{Re}m'(a,y).
\end{aligned}
\]
Put $\delta=\operatorname{Im}m(a,b)-\operatorname{Im}m'(a,b)$, which equals
$\operatorname{Im}m(b,a)-\operatorname{Im}m'(b,a)$ as $m$ and $m'$ are
symmetric. Subtracting (iii) from (i) at $y=b$ gives
$\operatorname{Re}m(b,b)+\operatorname{Re}m'(b,b)=-\delta$, and adding (ii)
and (iv) at $y=a$ gives $\operatorname{Re}m(a,a)+\operatorname{Re}m'(a,a)=\delta$.
So the four numbers $\operatorname{Re}m(a,a)$, $\operatorname{Re}m'(a,a)$,
$\operatorname{Re}m(b,b)$ and $\operatorname{Re}m'(b,b)$ add up to $0$. Each is
at most $0$, and negative unless its argument is $0$, by (1). Hence $a=b=0$,
and then (i) gives $L(y)=0$ for every $y$.

(3) The map $\Phi(p)=D_pf$ from $V\times V$ to the functions on $V$ is
$\mathbf R$-linear, with $\Phi(Jp)=i\,\Phi(p)$. Put
$\Theta(w_1,w_2)=\Phi(0,w_1)+\Phi(J(0,w_2))=\Phi(0,w_1)+i\,\Phi(0,w_2)$, so that
$\Theta(w_1,w_2)(y)=-2\pi i\bigl(T(y,w_1)+i\,T(y,w_2)\bigr)f(y)$. The map
$\Theta$ is injective: if $\Theta(w_1,w_2)=0$, then $T(y,w_1)=T(y,w_2)=0$ for
every $y$ in the set where $f\neq0$, which is open and nonempty; a linear form
that vanishes on a nonempty open set is $0$, so $T(\cdot,w_1)=T(\cdot,w_2)=0$,
and $w_1=w_2=0$. So the image of $\Theta$ has dimension $2\dim V$. It is
contained in the image of $\Phi$, which has dimension at most $2\dim V$, so the
two images are equal. Hence there is an $\mathbf R$-linear map
$W\colon V\to V\times V$ with $\Theta(W(a))=\Phi(a,0)=\partial_af$ for every
$a$, and
$m(a,y)=-2\pi i\bigl(T(y,W(a)_1)+i\,T(y,W(a)_2)\bigr)$ has the required
property.
\end{proof}
% Lean: (1) is gau_of (BoxFinA.lean) through schwartz_logDeriv_linear_gaussian and
% schwartz_logDeriv_linear_formula; the Lean's Gau structure is the conclusion of (1). (2) is
% opp, whose linarith steps use (i) and (iii) at b and (ii) and (iv) at a, as above. (3) is
% Twist.FockLogDeriv.main, with eq_zero_of_forall_nhds for the open set.

\begin{lemma}\label{asm-theta-agree-doubling-level}
Let $B$ be a theta base for a CM field $E$
(Definition~\ref{asm-theta-agree-models}(1)), and let $\kappa$ be a finite set.
Call an element of $\mathbf A^\kappa$, or of $\mathbf A^\kappa\times\mathbf A^\kappa$,
\emph{finite} if its components at $\infty$ are $0$.
\begin{enumerate}
\item Let $A$ and $Y$ be lattice spaces of ranks $m$ and $n$, let
  $d\colon M_{m\times2n}(\mathbf A_E)\to\mathbf A^\kappa\times\mathbf A^\kappa$
  be an isomorphism as in condition (a) of
  Definition~\ref{asm-theta-agree-doubling}(6), with Schrödinger operators
  $\rho$, let $K\subset G_Y$ be compact, and let $\Phi$ be a Schwartz-Bruhat
  function on $\mathbf A^\kappa$. Then there are a set $\Lambda$ of
  $u\in M_{m\times2n}(\mathbf A_E)$ with $d(u)$ finite, such that
  $uk^{\mathsf T}\in\Lambda$ for $u\in\Lambda$ and $k\in K$ and
  $\rho(u)\Phi=\Phi$ for $u\in\Lambda$, and a finite set $s$ of finite elements
  of $\mathbf A^\kappa$, with the following property. If $\Psi$ is a function
  on $\mathbf A^\kappa$ with $\rho(u)\Psi=\Psi$ for every $u\in\Lambda$, and
  $x\in\mathbf A^\kappa$ is at $\infty$ with $\Psi(x+y)=0$ for every $y\in s$,
  then $\Psi(x+y)=0$ for every finite $y$.
\item Let $\beta$ be a finite set, let
  $e\colon M_{m\times\beta}(\mathbf A_E)\to\mathbf A^\kappa\times\mathbf A^\kappa$
  be an $\mathbf A$-linear isomorphism, let $D_u$, for $u$ at $\infty$, be the
  operators of Lemma~\ref{asm-theta-agree-real-derivative} for $e$, and let
  $\theta_w\in M_m(\mathbf C)$ and $\theta'_w\in M_\beta(\mathbf C)$ for the
  infinite places $w$. For signs $\varepsilon=(\varepsilon_w)$ let
  $G_\varepsilon$ be the space of the $f\in\mathcal S(F_\infty^\kappa)$ with
  $D_{u'}f=i\,D_uf$ whenever $u$ and $u'$ are at $\infty$ and
  $\pi_w(u')=\varepsilon_wi\,\theta_w\pi_w(u)\theta_w'^{\mathsf T}$ for every
  $w$. Then a nonzero $f\in G_\varepsilon$ vanishes nowhere; and if
  $f\in G_\varepsilon$ and $f'\in G_{\varepsilon'}$ are nonzero, then
  $f'\in G_\varepsilon$.
\end{enumerate}
\end{lemma}
\leanref{Development/WeilBound/ThetaLifts/LatticeBasics/SecondFromFirstStatement.lean}{GlobalTheory.ThK.Dbl.FinLevel}
\leanref{Development/WeilBound/ThetaLifts/LatticeBasics/FinLevel.lean}{GlobalTheory.ThK.Dbl.finLevel}
\leanref{Development/WeilBound/ThetaLifts/LatticeBasics/BoxFinV.lean}{GlobalTheory.ThK.Dbl.Box.Gs}
\leanref{Development/WeilBound/ThetaLifts/LatticeBasics/BoxFinV.lean}{GlobalTheory.ThK.Dbl.Box.Gs.gau}
\leanref{Development/WeilBound/ThetaLifts/LatticeBasics/BoxFinV.lean}{GlobalTheory.ThK.Dbl.Box.Gs.same}
\begin{proof}
(1) Let $\epsilon\in\mathbf A$ have the components $0$ at $\infty$ and $1$ at
the finite places. Then $\epsilon^2=\epsilon$, $\epsilon x$ is finite for every
$x$, and $x$ is finite if and only if $\epsilon x=x$. For
$w=(w_1,w_2)\in\mathbf A^\kappa\times\mathbf A^\kappa$ put
$(\rho_0(w)\phi)(x)=\psi\bigl(x\cdot w_2+\tfrac12w_1\cdot w_2\bigr)\phi(x+w_1)$,
so that $\rho(u)=\rho_0(d(u))$.

\emph{Step 1: a level.} For every Schwartz-Bruhat function $\Phi$ on
$\mathbf A^\kappa$ there is a neighbourhood $N$ of $0$ with
$\rho_0(\epsilon w)\Phi=\Phi$ for all $w\in N$. The set of the $\Phi$ with this
property contains $0$ and is closed under sums and scalar multiples, as two
neighbourhoods meet in a neighbourhood. A Schwartz-Bruhat function is a finite
sum of functions
$G(x)=f(x_\infty)\prod_{i\in\kappa}1_{b_i+\hat J}(x_{f,i})$, with $f$ a Schwartz
function on $F_\infty^\kappa$, $J$ a fractional ideal of $F$ with closure
$\hat J$ in $\mathbf A_f$, and $b_i\in\mathbf A_f$. Fix such a $G$, and let $C$
be the support of its finite factor. As $\hat J$ is open, $C$ is compact and
$\psi$, being continuous, is locally constant on the finite elements, there is
a neighbourhood $N$ of $0$ such that for $w\in N$ every finite
component of $w_1$ lies in $\hat J$ and
$\psi\bigl(z\cdot\epsilon w_2+\tfrac12\epsilon w_1\cdot\epsilon w_2\bigr)=1$
for every finite $z$ with finite part in $C$. Let $w\in N$. Then
$x+\epsilon w_1$ has the component of $x$ at $\infty$, and its finite
components lie in the same cosets of $\hat J$ as those of $x$, so
$G(x+\epsilon w_1)=G(x)$. If $G(x)=0$, both sides of
$\rho_0(\epsilon w)G(x)=G(x)$ vanish. Otherwise the finite part of $x$ lies in
$C$, and as $x\cdot\epsilon w_2=\epsilon x\cdot\epsilon w_2$ the factor
$\psi(x\cdot\epsilon w_2+\tfrac12\epsilon w_1\cdot\epsilon w_2)$ is $1$. So
$\rho_0(\epsilon w)G=G$.

\emph{Step 2: finitely many values.} Let $N$ be a neighbourhood of $0$ in
$\mathbf A^\kappa\times\mathbf A^\kappa$. Then there is a finite set $s$ of
finite elements of $\mathbf A^\kappa$ such that: if $\Psi$ is a function on
$\mathbf A^\kappa$ with $\rho_0(\epsilon w)\Psi=\Psi$ for every $w\in N$, and
$x$ is at $\infty$ with $\Psi(x+t)=0$ for every $t\in s$, then $\Psi(x+y)=0$
for every finite $y$. Indeed, $N$ contains $N_1\times N_2$ for neighbourhoods
$N_1$ and $N_2$ of $0$. As $B$ is a theta base, $\psi(z)=\psi_0(\xi z)$ for the
standard character $\psi_0$ of $\mathbf A/F$ and some $\xi\in F^\times$. There
is a fractional ideal $J$ such that the vector $a\,e_i$, with component
$a\in\mathbf A$ at $i$ and $0$ elsewhere, lies in $N_2$ whenever $a$ is finite
with finite part in $\hat J$ and $i\in\kappa$. Let $J^\vee$ be the trace dual
of $J$, and $C$ the set of the finite $y\in\mathbf A^\kappa$ whose finite
components lie in $\xi^{-1}\hat J^\vee$; it is compact. Cover it by finitely
many sets $t+N_1$ with $t\in C$, and let $s$ be the set of these $t$. Now let
$\Psi$ and $x$ be as stated, let $y$ be finite, and suppose that
$\Psi(x+y)\neq0$. For $a$ as above and $i\in\kappa$, invariance at
$w=(0,a\,e_i)$, with $\epsilon x=0$ and $\epsilon a=a$, gives
$\Psi(x+y)=\psi(a\,y_i)\Psi(x+y)$, so $\psi_0(\xi\,a\,y_i)=1$. For a finite
$c\in\mathbf A$, $\psi_0(ac)=1$ for all finite $a$ with finite part in $\hat J$
if and only if the finite part of $c$ lies in $\hat J^\vee$ \cite{Tat67}. So
$\xi y_i\in\hat J^\vee$ for every $i$, and $y\in C$. Choose $t\in s$ with
$y-t\in N_1$. As $y-t$ is finite, invariance at $w=(y-t,0)$ gives
$\Psi(x+t)=\Psi(x+t+(y-t))=\Psi(x+y)\neq0$, against $t\in s$.

\emph{Step 3: the set $\Lambda$.} The map
$(w,k)\mapsto d\bigl(d^{-1}(w)k^{\mathsf T}\bigr)$ on
$(\mathbf A^\kappa\times\mathbf A^\kappa)\times G_Y$ is continuous and takes
the value $0$ at $(0,k)$. By Step 1, for every $k\in K$ there are
neighbourhoods of $0$ and of $k$ on whose product
$\rho_0\bigl(\epsilon\,d(d^{-1}(w)k'^{\mathsf T})\bigr)\Phi=\Phi$. As $K$ is
compact, there is a neighbourhood $N$ of $0$ with
$\rho_0\bigl(\epsilon\,d(d^{-1}(w)k^{\mathsf T})\bigr)\Phi=\Phi$ for all
$w\in N$ and $k\in K$. Let $s$ be given by Step 2 for $N$, and let $\Lambda$ be
the set of the $u$ with $\epsilon u=u$ and $\rho(uk^{\mathsf T})\Phi=\Phi$ for
every $k\in K$. For $u\in\Lambda$: $d(u)=d(\epsilon u)=\epsilon\,d(u)$ is
finite, as $d$ is $\mathbf A$-linear; $\rho(u)\Phi=\Phi$, taking $k=1$; and
for $k\in K$, $uk^{\mathsf T}\in\Lambda$, since
$(uk^{\mathsf T})k'^{\mathsf T}=u(k'k)^{\mathsf T}$ with $k'k\in K$. Let $\Psi$
be fixed by every $\rho(u)$, $u\in\Lambda$, and let $w\in N$. Then
$u=\epsilon\,d^{-1}(w)$ lies in $\Lambda$, since $\epsilon u=u$ and
$\rho(uk^{\mathsf T})\Phi=\rho_0\bigl(\epsilon\,d(d^{-1}(w)k^{\mathsf T})\bigr)\Phi=\Phi$
for $k\in K$; and $\rho(u)=\rho_0(\epsilon w)$. So $\rho_0(\epsilon w)\Psi=\Psi$
for every $w\in N$, and Step 2 gives the property of $s$.

(2) Put $V=F_\infty^\kappa$ and
$T(x,b)=\operatorname{Tr}_{F_\infty/\mathbf R}\bigl(\sum_{i\in\kappa}x_ib_i\bigr)$;
if $T(x,b)=0$ for every $x$, then $b=0$, as the trace form of $F$ over
$\mathbf Q$ is nondegenerate. For $p=(p_1,p_2)\in V\times V$ let $D_p$ be the
operator $D_{p_1,p_2}$ of Lemma~\ref{asm-theta-agree-density-derivative}, so that
$D_u=D_p$ when $p$ is the pair of components at $\infty$ of $e_1(u)$ and
$e_2(u)$, and $D_{p+q}=D_p+D_q$.

\emph{Coordinates.} As $e$ is $\mathbf A$-linear, it identifies the elements of
$M_{m\times\beta}(\mathbf A_E)$ at $\infty$ with $V\times V$; write $s(p)$ for
the element with coordinates $p$. As in the proof of
Lemma~\ref{asm-theta-agree-density-gaussian}, $p\mapsto(\pi_w(s(p)))_w$ is an
$\mathbf R$-linear bijection from $V\times V$ onto
$\prod_wM_{m\times\beta}(\mathbf C)$, since $E$ is a CM field. So for every
family of signs $\varepsilon$ there is a unique $\mathbf R$-linear map
$J_\varepsilon$ of $V\times V$ with
$\pi_w(s(J_\varepsilon p))=\varepsilon_wi\,\theta_w\pi_w(s(p))\theta_w'^{\mathsf T}$
for every $w$, and $f\in G_\varepsilon$ if and only if
$D_{J_\varepsilon p}f=i\,D_pf$ for every $p$.

\emph{Nowhere vanishing.} Let $f\in G_\varepsilon$ be nonzero. By
Lemma~\ref{asm-theta-agree-doubling-level-gaussian}(3), for $T$ and
$J=J_\varepsilon$, there is an $\mathbf R$-bilinear $m$ with
$\partial_af(y)=m(a,y)f(y)$, so $f$ vanishes nowhere by part (1) of that lemma.

\emph{The signs.} Let $f\in G_\varepsilon$ and $f'\in G_{\varepsilon'}$ be
nonzero; by the last paragraph, both satisfy the hypotheses of
Lemma~\ref{asm-theta-agree-doubling-level-gaussian}(1). Let $p\in V\times V$,
and put $q=J_\varepsilon p-J_{\varepsilon'}p$. Then
$\pi_w(s(q))=(\varepsilon_w-\varepsilon'_w)\,i\,\theta_w\pi_w(s(p))\theta_w'^{\mathsf T}$,
which is $0$ when $\varepsilon_w=\varepsilon'_w$. Hence
$\pi_w(s(J_{\varepsilon'}q))=-\pi_w(s(J_\varepsilon q))$ for every $w$: both
sides are $0$ where $\varepsilon_w=\varepsilon'_w$, and
$\varepsilon'_w=-\varepsilon_w$ elsewhere. So $J_{\varepsilon'}q=-J_\varepsilon q$.
As $f\in G_\varepsilon$ and $f'\in G_{\varepsilon'}$,
$D_{J_\varepsilon q}f=i\,D_qf$ and $D_{-J_\varepsilon q}f'=i\,D_qf'$. With
$q=(q_1,q_2)$ and $J_\varepsilon q=(r_1,r_2)$, these are the relations of
Lemma~\ref{asm-theta-agree-doubling-level-gaussian}(2) for $f$, $g=f'$,
$a=q_1$, $b=r_1$, $L=T(\cdot,q_2)$ and $L'=T(\cdot,r_2)$. So $q_1=0$ and
$T(\cdot,q_2)=0$, that is, $D_q=0$. Hence
$D_{J_\varepsilon p}f'=D_{J_{\varepsilon'}p}f'+D_qf'=i\,D_pf'$, and
$f'\in G_\varepsilon$.
\end{proof}
% Lean: (1) finLevel (FinLevel.lean): Step 1 is level, by induction over cosetSpanPi
% (schwartzBruhatPi_eq_cosetSpanPi), with finChar_isLocallyConstant; Step 2 is
% finite_values, with std_eq for xi, exists_idealFiniteAdeles_subset for J and
% forall_finiteAdelicChar_mul_eq_one_iff for the duality; Step 3 is finLevel itself, with
% Lambda = {u | epsF u = u and rho(actB k u) Phi = Phi for k in K}, where actB k u = u k^T.
% (2) Gs.gau and Gs.same (BoxFinV.lean): the coordinates s and the map J are
% exists_real_structure', the linear form T is linForm (nondegenerate by
% trace_mul_eq_zero), and Gs.same uses exists_rel for J_eps' p and J_eps q, then opp.

\begin{lemma}\label{asm-theta-agree-doubling-section}
Let $(X,W,M,\chi_X,\chi_W,\omega)$ be a Weil pair of signature $(r,s;r',s')$
for a theta base $B$ of a CM field $E$
(Definition~\ref{asm-theta-agree-triple}(1)), with $X$ and $W$ of ranks
$n,d\geq1$, put $I=\{1,\ldots,n\}\times\{1,\ldots,d\}$, and let
$\Phi_1,\Phi_2\in\operatorname{PD}(M)$.
\begin{enumerate}
\item Suppose that $\min(r,s)=0$. Let $(I\sqcup I,e^\square,\omega_2,S,l)$ be a
  doubling datum for $(X,W)$ with character $\chi_X$ and the isomorphism
  $e^\square$ of Lemma~\ref{asm-theta-agree-doubling-datum}(1), and let $K$ be
  the lattice stabilizer of the doubled space of $W$. Then the functions
  $\omega_2(k)(\Phi_1\boxtimes\bar\Phi_2)$, $k\in K$, span a finite-dimensional
  space, and $f(x)=l(\omega_2(x)(\Phi_1\boxtimes\bar\Phi_2))$ is a section of
  degree $(n-d)/2$ for $(K,\chi_X)$ (Definition~\ref{asm-theta-agree-doubling}(4)).
\item Suppose that $\min(r',s')=0$. Let $(I\sqcup I,e^\square,\omega_2,S,l)$ be
  a doubling datum for $(W,X)$ with character $\chi_W$ and the isomorphism
  $e^\square$ of Lemma~\ref{asm-theta-agree-doubling-datum}(2), and let $K$ be
  the lattice stabilizer of the doubled space of $X$. Then the functions
  $\omega_2(k)(\Phi_1\boxtimes\bar\Phi_2)$, $k\in K$, span a finite-dimensional
  space, and $f(x)=l(\omega_2(x)(\Phi_1\boxtimes\bar\Phi_2))$ is a section of
  degree $(d-n)/2$ for $(K,\chi_W)$.
\end{enumerate}
\end{lemma}
\leanref{Development/WeilBound/ThetaLifts/LatticeBasics/PairsStable.lean}{GlobalTheory.ThK.Dbl.BoxSec}
\leanref{Development/WeilBound/ThetaLifts/LatticeBasics/BoxSecX.lean}{GlobalTheory.ThK.Dbl.BoxSecX}
\leanref{Development/WeilBound/ThetaLifts/LatticeBasics/BoxSec.lean}{GlobalTheory.ThK.Dbl.boxSec_of}
\leanref{Development/WeilBound/ThetaLifts/LatticeBasics/BoxSec.lean}{GlobalTheory.ThK.Dbl.boxSecX_of}
\leanref{Development/WeilBound/ThetaLifts/LatticeBasics/BoxSec.lean}{GlobalTheory.ThK.Dbl.Box.isSec_of}
\leanref{Development/WeilBound/ThetaLifts/LatticeBasics/BoxSec.lean}{GlobalTheory.ThK.Dbl.Box.polData_kfin}
\leanref{Development/WeilBound/ThetaLifts/LatticeBasics/BoxFinC.lean}{GlobalTheory.ThK.Dbl.Box.box_kfin}
\leanref{Development/WeilBound/ThetaLifts/LatticeBasics/BoxFinC.lean}{GlobalTheory.ThK.Dbl.Box.vac_kfin}
\leanref{Development/WeilBound/ThetaLifts/LatticeBasics/BoxFinC.lean}{GlobalTheory.ThK.Dbl.Box.D_kfin}
\leanref{Development/WeilBound/ThetaLifts/LatticeBasics/BoxFinC.lean}{GlobalTheory.ThK.Dbl.Box.D_box}
\begin{proof}
We prove (1); (2) is the same with $X$ and $W$ exchanged. Let $\rho^\square$ be
the Schrödinger operators of $e^\square$, and for $u\in M_{n\times2d}(\mathbf A_E)$
at $\infty$ let $D_u$ be the operator of
Lemma~\ref{asm-theta-agree-real-derivative} for $e^\square$; for
$u\in M_{n\times d}(\mathbf A_E)$ at $\infty$ let $D_u$ also denote the one for
the isomorphism $e$ of $M$. By Lemma~\ref{asm-theta-agree-doubling-stabilizer},
$K$ is compact. Let $\mathcal K$ be the set of Schwartz-Bruhat functions $\Phi$
on $\mathbf A^{I\sqcup I}$ that lie in a finite-dimensional space of
Schwartz-Bruhat functions stable under every $\omega_2(k)$, $k\in K$; it is a
subspace.

\emph{The section.} By condition (b) of
Definition~\ref{asm-theta-agree-doubling}(6), $x\mapsto l(\omega_2(x)\Phi)$ is
continuous for Schwartz-Bruhat $\Phi$, as $l$ is admissible, and, as $l$ is
invariant under the $\rho^\square(z\Delta^{\mathsf T})$ by condition (d), it has
the law $l(\omega_2(p)\Psi)=\chi_X(\det t(p))|\det t(p)|^{n/2}l(\Psi)$ with
$n/2=(n-d)/2+d/2$. If $\Phi\in\mathcal K$, lying in a $K$-stable
finite-dimensional space $N$, the right translates of
$x\mapsto l(\omega_2(x)\Phi)$ by $K$ lie in the image of $N$ under
$\Psi\mapsto(x\mapsto l(\omega_2(x)\Psi))$. So it suffices to show that
$\Phi_1\boxtimes\bar\Phi_2\in\mathcal K$, and, as $\boxtimes$ is linear in the
first factor and conjugate linear in the second, that
$\Phi_1\boxtimes\bar\Phi_2\in\mathcal K$ for $\Phi_i=f_i\otimes\phi_i$ with
$f_i\in\operatorname{Pol}(M)$ and $\phi_i$ Schwartz-Bruhat on
$\mathbf A_f^I$. Fix $\phi_1$ and $\phi_2$, and write
$b(f_1,f_2)=(f_1\otimes\phi_1)\boxtimes\overline{f_2\otimes\phi_2}$.

\emph{Derivatives.} If $\Phi\in\mathcal K$ and $u$ is at $\infty$, then
$D_u\Phi\in\mathcal K$. Indeed, let $N$ be as above. The $D_v\Psi$ with $v$ at
$\infty$ and $\Psi\in N$ span a finite-dimensional space, since $D_v$ depends
linearly on the components at $\infty$ of $e^\square(v)$, which range over a
finite-dimensional real space; and it is stable under $K$, since
$\omega_2(k)D_v=D_{vk^{\mathsf T}}\omega_2(k)$, by the rule
$\omega_2(k)\rho^\square(w)=\rho^\square(wk^{\mathsf T})\omega_2(k)$ and
Lemma~\ref{asm-theta-agree-real-derivative}. On box functions,
$D_{(u_1\ u_2)}b(f_1,f_2)=b(D_{u_1}f_1,f_2)+b(f_1,D_{u_2}f_2)$, by the form of
$e^\square$. With $u_2=0$, respectively $u_1=0$, the set of $f_1$ with
$b(f_1,f_2)\in\mathcal K$ is stable under the $D_u$ for fixed $f_2$, and
likewise in $f_2$.

\emph{Gaussians.} Let $f_1$ and $f_2$ be Gaussians for $M$
(Definition~\ref{asm-theta-agree-models}(4)), with signs $\varepsilon$ and
$\varepsilon'$, both nonzero. In the notation of
Lemma~\ref{asm-theta-agree-doubling-level}(2) for $e$ and the matrices
$\theta_{X,w}$ and $\theta_{W,w}$, $f_1\in G_\varepsilon$ and
$f_2\in G_{\varepsilon'}$, by Lemma~\ref{asm-theta-agree-density-derivative};
so $f_2\in G_\varepsilon$ by Lemma~\ref{asm-theta-agree-doubling-level}(2). By
the formula for $D$ on box functions, $\Phi=b(f_1,f_2)$ then satisfies
$D_{u'}\Phi=i\,D_u\Phi$ whenever $u,u'$ are at $\infty$ with
$\pi_w(u')=\varepsilon_wi\,\theta_{X,w}\pi_w(u)\Theta_w^{\mathsf T}$,
$\Theta_w=\operatorname{diag}(\theta_{W,w},-\theta_{W,w})$; the sign of the
second block matches the conjugation of $f_2$. We show $\Phi\in\mathcal K$.
Let $\Lambda$ and $s$ be given by Lemma~\ref{asm-theta-agree-doubling-level}(1)
for $e^\square$, $K$ and $\Phi$, and let $N$ be the space of Schwartz-Bruhat
$\Psi$ with $\rho^\square(u)\Psi=\Psi$ for $u\in\Lambda$ and with
$D_{u'}\Psi=i\,D_u\Psi$ for all such $u,u'$. Then $\Phi\in N$. The space $N$ is
stable under $K$: for $k\in K$ and $u\in\Lambda$,
$\rho^\square(u)\omega_2(k)\Psi=\omega_2(k)\rho^\square(u(k^{-1})^{\mathsf T})\Psi=\omega_2(k)\Psi$,
as $u(k^{-1})^{\mathsf T}\in\Lambda$; and $\omega_2(k)D_u=D_{uk^{\mathsf T}}\omega_2(k)$,
where $u\mapsto uk^{\mathsf T}$ preserves the relation between $u$ and $u'$
because $\pi_w(k)$ commutes with $\Theta_w$. The space $N$ is
finite-dimensional, because $\Psi\mapsto(\Psi(y))_{y\in s}$ is injective on
$N$: let $\Psi\in N$ vanish on $s$, and let $y\in s$. The function
$x\mapsto\Psi(x+y)$ on the elements at $\infty$ is a Schwartz function
satisfying the Gaussian relations of
Lemma~\ref{asm-theta-agree-doubling-level}(2) for $e^\square$ and the matrices
$\theta_{X,w}$ and $\Theta_w$, and it vanishes at $0$, so it is $0$ by that
lemma. So $\Psi(x+y)=0$ for $x$ at $\infty$ and $y\in s$, and
Lemma~\ref{asm-theta-agree-doubling-level}(1) gives $\Psi(x+y)=0$ for all
finite $y$, that is, $\Psi=0$.

\emph{Conclusion.} If $f_1$ or $f_2$ is $0$ then $b(f_1,f_2)=0$. So for each
Gaussian $f_2$ the set of $f_1\in\operatorname{Pol}(M)$ with
$b(f_1,f_2)\in\mathcal K$ is a subspace containing every Gaussian and stable
under the $D_u$; by Lemma~\ref{asm-theta-agree-density-derivative} and the
definition of $\operatorname{Pol}(M)$ it is all of $\operatorname{Pol}(M)$.
Then for each $f_1\in\operatorname{Pol}(M)$ the set of $f_2$ with
$b(f_1,f_2)\in\mathcal K$ is a subspace containing every Gaussian and stable
under the $D_u$, so it is $\operatorname{Pol}(M)$.
\end{proof}
% Lean: boxSec_of finLevel stab2Cpt and boxSecX_of: isSec_of (sec_cont_law for continuity
% and the law), polData_kfin (the reduction to pure tensors), box_kfin: D_kfin (exists_D_fin,
% om_D), D_box, base with Gs.same and vac_kfin (N, its K-stability from hs.impl and om_D, and
% finiteness from finLevel and Gs.gau). kfin om K is the set K-fin. In BoxSecX the transposed
% coordinates IsDblCoordX are used.

\begin{lemma}\label{asm-theta-agree-doubling-siegel-weil}
Assume the Siegel-Weil formula in the form of
Theorem~\ref{asm-theta-agree-siegel-weil}. Let $(X,W,M,\chi_X,\chi_W,\omega)$
be a Weil pair of signature $(r,s;r',s')$ for a theta base $B$ of a CM field
$E$ (Definition~\ref{asm-theta-agree-triple}(1)), with $X$ and $W$ of ranks
$n,d\geq1$, let $\Phi_1,\Phi_2\in\operatorname{PD}(M)$ and
$\Phi=\Phi_1\boxtimes\bar\Phi_2$.
\begin{enumerate}
\item Suppose that $\min(r,s)=0$. Let $(I\sqcup I,e^\square,\omega_2,S,l)$ and
  $K$ be as in Lemma~\ref{asm-theta-agree-doubling-section}(1), let $a$ be the
  height function for $K$, let $\mu_X$ and $\mathcal D_X$ be the measure and
  the domain of $M$, and put $s_0=(n-d)/2$, $\kappa=1$ if $d<n$ and $\kappa=2$ if
  $d\geq n$, $f(x)=l(\omega_2(x)\Phi)$ and $f_s=a^{s-s_0}f$. Then for every
  $x\in G_W$ there is a meromorphic function $e_x$ on $\mathbf C$ such that for
  $\operatorname{Re}s$ large enough the Eisenstein sum of $f_s$ at $x$
  converges unconditionally to $e_x(s)$, $e_x$ is holomorphic at $s_0$, and
  \[
  e_x(s_0)=\frac{\kappa}{\mu_X(\mathcal D_X)}\int_{\mathcal D_X}\theta\bigl(\omega_2(x)S(g)\Phi\bigr)\,d\mu_X(g).
  \]
\item Suppose that $\min(r',s')=0$. Let $(I\sqcup I,e^\square,\omega_2,S,l)$ and
  $K$ be as in Lemma~\ref{asm-theta-agree-doubling-section}(2), let $a$ be the
  height function for $K$, let $\mu$ be a Haar measure on $U_W$ and
  $\mathcal D$ a fundamental domain for $U_W(F)$ with respect to $\mu$, and put
  $s_0=(d-n)/2$, $\kappa=1$ if $n<d$ and $\kappa=2$ if $n\geq d$,
  $f(x)=l(\omega_2(x)\Phi)$ and $f_s=a^{s-s_0}f$. Then the same holds for
  $x\in G_X$, with
  $e_x(s_0)=\kappa\,\mu(\mathcal D)^{-1}\int_{\mathcal D}\theta(\omega_2(x)S(h)\Phi)\,d\mu(h)$.
\end{enumerate}
\end{lemma}
\leanref{Development/WeilBound/ThetaLifts/LatticeBasics/Hypotheses2.lean}{GlobalTheory.ThK.Dbl.SWAniso}
\leanref{Development/WeilBound/ThetaLifts/LatticeBasics/Hypotheses2.lean}{GlobalTheory.ThK.Dbl.EisValueAt}
\leanref{Development/WeilBound/ThetaLifts/LatticeBasics/SiegelWeilAnisoOf.lean}{GlobalTheory.ThK.Dbl.siegelWeilAniso_of}
\leanref{Development/WeilBound/ThetaLifts/LatticeBasics/SiegelWeilAnisoOf.lean}{GlobalTheory.ThK.Dbl.Box.sw_transl}
\leanref{Development/WeilBound/ThetaLifts/LatticeBasics/SiegelWeilAnisoOf.lean}{GlobalTheory.ThK.Dbl.Box.dbl_comm}
\leanref{Development/WeilBound/ThetaLifts/LatticeBasics/SiegelWeilAnisoOf.lean}{GlobalTheory.ThK.Dbl.Box.dbl_inj}
\leanref{Development/WeilBound/ThetaLifts/LatticeBasics/SiegelWeilAnisoOf.lean}{GlobalTheory.ThK.Dbl.Box.theta_inj}
\begin{proof}
We prove (1); (2) is the same with the roles of $X$ and $W$ exchanged. Write
$\rho$ for the Schrödinger operators of $e^\square$.

\emph{Commutation.} For $p\in P_W$, $g\in U_X$ and Schwartz-Bruhat $\Psi$,
$\omega_2(p)S(g)\Psi=S(g)\omega_2(p)\Psi$. First, a Schwartz-Bruhat $\Psi'$ with
$l(\rho(w)\Psi')=0$ for every $w$ is $0$: by condition (d) of
Definition~\ref{asm-theta-agree-doubling}(6), $\theta(\rho(w)\Psi')$ is the sum
of the $l(\rho((y,-y))\rho(w)\Psi')$, and $\rho((y,-y))\rho(w)$ is a scalar
multiple of $\rho((y,-y)+w)$, so $\theta(\rho(w)\Psi')=0$ for every $w$; for
fixed $x$, the numbers $\theta(\rho((x,y))\Psi')$, as $y$ varies, are the values
of a Fourier series with coefficients $\psi(\cdot)\Psi'(\xi+x)$, $\xi\in F^{I\sqcup I}$,
so $\Psi'(x)=0$. Now let $w$ be arbitrary and write
$w=g\,w'p^{\mathsf T}$. By the rules
$\omega_2(p)\rho(w'')=\rho(w''p^{\mathsf T})\omega_2(p)$ and
$S(g)\rho(w'')=\rho(gw'')S(g)$, the law
$l(\omega_2(p)\Psi'')=\chi_X(\det t(p))|\det t(p)|^{n/2}l(\Psi'')$ and
$l(S(g)\Psi'')=l(\Psi'')$, both
$l(\rho(w)\omega_2(p)S(g)\Psi)$ and $l(\rho(w)S(g)\omega_2(p)\Psi)$ equal
$\chi_X(\det t(p))|\det t(p)|^{n/2}l(\rho(w')\Psi)$.

\emph{Conjugation.} By Lemma~\ref{asm-theta-agree-doubling-frame} there are a
frame $b$ of $W$ and $p\in P_W$ with $K=pK_bp^{-1}$; put $d_p=|\det t(p)|$.
As in Step 1 of the proof of
Lemma~\ref{asm-theta-agree-doubling-continuation}, $a_0(x)=d_p\,a(xp^{-1})$ is
a height function for $K_b$. The function $\Psi=\omega_2(p^{-1})\Phi$ is
Schwartz-Bruhat, and its translates $\omega_2(k_0)\Psi=\omega_2(p^{-1})\omega_2(pk_0p^{-1})\Phi$,
$k_0\in K_b$, span a finite-dimensional space, as $pk_0p^{-1}\in K$ and by
Lemma~\ref{asm-theta-agree-doubling-section}(1).

\emph{The value.} Theorem~\ref{asm-theta-agree-siegel-weil}, for $A=X$, $Y=W$,
the character $\chi_X$, the doubling datum, $b$, $a_0$, $\mu_X$,
$\mathcal D_X$ and $\Psi$, at the point $xp$, gives a meromorphic $e$. Put
$e_x(s)=d_p^{\,s_0-s}e(s)$. For $\gamma\in G_W(F)$,
$a_0(\gamma xp)^{s-s_0}l(\omega_2(\gamma xp)\Psi)=d_p^{\,s-s_0}f_s(\gamma x)$, so for
$\operatorname{Re}s$ large the Eisenstein sum of $f_s$ at $x$ converges
unconditionally to $e_x(s)$; $e_x$ is holomorphic at $s_0$, and by the
commutation
\[
e_x(s_0)=e(s_0)=\frac{\kappa}{\mu_X(\mathcal D_X)}\int_{\mathcal D_X}\theta\bigl(\omega_2(xp)S(g)\omega_2(p^{-1})\Phi\bigr)\,d\mu_X(g)=\frac{\kappa}{\mu_X(\mathcal D_X)}\int_{\mathcal D_X}\theta\bigl(\omega_2(x)S(g)\Phi\bigr)\,d\mu_X(g).\qedhere
\]
\end{proof}
% Lean: siegelWeilAniso_of SiegelWeilStd FrameConj FinLevel Stab2Cpt; its two halves are
% SWAniso (1) and SWAnisoX (2), both sw_transl with polData_kfin and box_kfin. The
% commutation is dbl_comm (through dbl_inj and theta_inj, the latter by
% SchrodingerUniq.coeff_zero_eq_zero).

\begin{lemma}\label{asm-theta-agree-doubling-value}
Assume the meromorphic continuation of Siegel Eisenstein series in the form of
Theorem~\ref{asm-theta-agree-eisenstein} and the Siegel-Weil formula in the form
of Theorem~\ref{asm-theta-agree-siegel-weil}. In the situation of
Lemma~\ref{asm-theta-agree-doubling-siegel-weil}(1), respectively (2), there
are an open set $U\subset\mathbf C$ containing $s_0$ and a function $\hat E$ on
$\mathbf C\times G_W$, respectively $\mathbf C\times G_X$, with properties (1)
to (4) of Lemma~\ref{asm-theta-agree-doubling-continuation} for the flat family
$(f_s)$, such that $\hat E(s_0,x)$ is the value $e_x(s_0)$ of
Lemma~\ref{asm-theta-agree-doubling-siegel-weil} for every $x$.
\end{lemma}
\leanref{Development/WeilBound/ThetaLifts/LatticeBasics/PairsStable.lean}{GlobalTheory.ThK.Dbl.SWFull}
\leanref{Development/WeilBound/ThetaLifts/LatticeBasics/BoxSecX.lean}{GlobalTheory.ThK.Dbl.SWFullX}
\leanref{Development/WeilBound/ThetaLifts/LatticeBasics/SWFull.lean}{GlobalTheory.ThK.Dbl.swFull_of}
\leanref{Development/WeilBound/ThetaLifts/LatticeBasics/NormZeta.lean}{GlobalTheory.ThK.Dbl.swFullX_of}
\leanref{Development/WeilBound/ThetaLifts/LatticeBasics/SWFull.lean}{GlobalTheory.ThK.Dbl.IsEisCont.update}
\begin{proof}
We treat (1). By Lemma~\ref{asm-theta-agree-doubling-section}(1), $f$ is a
section of degree $s_0$ for $(K,\chi_X)$, and $\chi_X|_{\mathbf A^\times}=\eta^n$.
Let $U_1$ and $\hat E_1$ be given by
Lemma~\ref{asm-theta-agree-doubling-continuation}, let $k$, $r$ and $H$ be given
by its part (5) at $s_0$, and let $e_x$ be given by
Lemma~\ref{asm-theta-agree-doubling-siegel-weil}(1). Put
$U=U_1\cup\{s_0\}$, and let $\hat E$ equal $\hat E_1$ except that
$\hat E(s_0,x)=e_x(s_0)$.

The set $U$ is open, as $s_0$ has a punctured neighbourhood in $U_1$; it is
connected, as $U_1$ is connected and $s_0$ lies in its closure; and every point
has a punctured neighbourhood in $U_1\subset U$. Property (2) holds with
$\max(t,\operatorname{Re}s_0)$ for $t$. Away from $s_0$, properties (3) and (4)
are those of $\hat E_1$. At $s_0$: the function $e_x-\hat E_1(\cdot,x)$ is
meromorphic on the connected open set $U_1$ and vanishes on a half-plane, so
it vanishes on $U_1$; as $e_x$ is holomorphic at $s_0$, $\hat E(\cdot,x)$
agrees with $e_x$ near $s_0$ and is holomorphic there. Choose $\rho<r$ with the
closed disc of radius $\rho$ around $s_0$ in $U$. On its boundary circle,
$|\hat E(s,x)|=\rho^{-k}|H(s,x)|\leq C\rho^{-k}\operatorname{ht}(x)^N$, so by
the maximum principle the same bound holds on the disc, which gives (4) at
$s_0$. For continuity at $(s_0,x_0)$: the circle is compact and $\hat E$ is
continuous on the circle times $G_W$, so
$\sup_{|z-s_0|=\rho}|\hat E(z,x)-\hat E(z,x_0)|$ tends to $0$ as $x\to x_0$; by
the maximum principle so does $\sup_{|s-s_0|\leq\rho}|\hat E(s,x)-\hat E(s,x_0)|$,
and $\hat E(\cdot,x_0)$ is continuous at $s_0$.
\end{proof}
% Lean: swFull_of SWAniso EisContinues BoxSec and swFullX_of SWAnisoX EisContinues BoxSecX,
% through IsEisCont.update (mero_eq for the identity theorem, contAt_of_sphere for the
% continuity, Complex.norm_le_of_forall_mem_frontier_norm_le for the bound).

\begin{lemma}\label{asm-theta-agree-doubling-kernel}
Let $B$ be a theta base for a CM field $E$
(Definition~\ref{asm-theta-agree-models}(1)), let $Y$ be a lattice space of
rank $n\geq1$ that is Hermitian or skew-Hermitian of signature $(r,s)$ with
$\min(r,s)=0$ (Definition~\ref{asm-theta-agree-doubling}(1)), and let $K$, $a$
and $\iota$ be as in Lemma~\ref{asm-theta-agree-doubling-stabilizer}. Let
$\chi$ be a Hecke character of $E$, let $z,s_0\in\mathbf C$, let $f$ be a
continuous function on $G_Y$ with
$f(px)=\chi(\det t(p))|\det t(p)|^zf(x)$ for $p\in P_Y$ and $x\in G_Y$, put
$f_s=a^{s-s_0}f$, and let $U$ and $\hat E$ have properties (1) to (4) of
Lemma~\ref{asm-theta-agree-doubling-continuation} for $(f_s)$. Let $\mu$ be a
Haar measure on $U_Y$ and $\mathcal D$ a fundamental domain for $U_Y(F)$ with
respect to $\mu$. For $s\in U$ and a rapidly decreasing function $\phi$ on
$U_Y$ (Definition~\ref{asm-theta-agree-definite-products}(1)) put
\[
P(s,\phi)=\int_{\mathcal D\times\mathcal D}\hat E(s,\iota(h,h'))\,\chi(\det h')^{-1}\phi(h')\,\overline{\phi(h)}\,d\mu(h)\,d\mu(h').
\]
Then $s\mapsto P(s,\phi)$ is holomorphic on $U$ for every rapidly decreasing
$\phi$, and there is $t_1\in\mathbf R$ such that for $\operatorname{Re}s>t_1$:
$s\in U$, the function $w_s(g)=f_s(\iota(1,g))\chi(\det g)^{-1}$ is
$\mu$-integrable on $U_Y$, and
$P(s,\phi)=\int_{U_Y}w_s(g)\,\langle R_g\phi,\phi\rangle_{\mathcal D}\,d\mu(g)$
for every rapidly decreasing $\phi$.
\end{lemma}
\leanref{Development/WeilBound/ThetaLifts/LatticeBasics/ContinuousFlat.lean}{GlobalTheory.ThK.Dbl.kpair}
\leanref{Development/WeilBound/ThetaLifts/LatticeBasics/ContinuousFlat.lean}{GlobalTheory.ThK.Dbl.kpair_cont}
\leanref{Development/WeilBound/ThetaLifts/LatticeBasics/PairsStable.lean}{GlobalTheory.ThK.Dbl.Unfold0}
\leanref{Development/WeilBound/ThetaLifts/LatticeBasics/Unfold0.lean}{GlobalTheory.ThK.Dbl.unfold0}
\leanref{Development/WeilBound/ThetaLifts/LatticeBasics/UnfoldB.lean}{GlobalTheory.ThK.Dbl.orbit}
\leanref{Development/WeilBound/ThetaLifts/LatticeBasics/Unfold0.lean}{GlobalTheory.ThK.Dbl.Law.eis_inv}
\leanref{Development/WeilBound/ThetaLifts/LatticeBasics/ContinuousFlat.lean}{GlobalTheory.ThK.Dbl.Law.bound}
\leanref{Development/WeilBound/ThetaLifts/LatticeBasics/FarRight.lean}{GlobalTheory.ThK.Dbl.An.far_right}
\begin{proof}
Let $J$ be the matrix of $Y$ and $h_Y(x,y)=x^{\mathsf T}Jc(y)$.

\emph{Facts.} By Lemma~\ref{asm-theta-agree-definite-compact}, $U_Y(F)$ is
countable, $\mu(\mathcal D)<\infty$, there is a compact $\Omega\subset U_Y$ with
$U_Y(F)\Omega=U_Y$, and rapidly decreasing functions are bounded. The form
$h_Y$ is anisotropic: $h_Y(x,x)=0$ with $x\neq0$ is impossible, as $h_Y$ is
definite at an infinite place. As $a(px)=|\det t(p)|a(x)$,
$f_s(px)=\chi(\det t(p))|\det t(p)|^{s-s_0+z}f_s(x)$ for $p\in P_Y$. For
$p\in P_Y\cap G_Y(F)$, $t(p)$ has entries in
$E$, so $\chi(\det t(p))=1$ and $|\det t(p)|=1$; for $h\in U_Y$,
$\iota(h,h)\in P_Y$ with $t=h$ and $|\det h|=1$, so
$f_s(\iota(h,h)x)=\chi(\det h)f_s(x)$; and $\chi(\det\delta)=1$ for
$\delta\in U_Y(F)$.

\emph{The cosets.} Every coset of $P_Y\cap G_Y(F)$ in $G_Y(F)$ contains
$\iota(1,\delta)$ for exactly one $\delta\in U_Y(F)$. Indeed, let
$\gamma\in G_Y(F)$ and $\gamma^{-1}\Delta=\left(\begin{smallmatrix}A\\B\end{smallmatrix}\right)$.
As $\Delta$ is isotropic for $J^\square$,
$A^{\mathsf T}Jc(A)=B^{\mathsf T}Jc(B)$. If $Av=0$, then $h_Y(Bv,Bv)=h_Y(Av,Av)=0$,
so $Bv=0$, and $v=0$ as $\gamma$ is invertible; so $A$ is invertible, and so is
$B$ in the same way. Then $\delta=AB^{-1}\in U_Y(F)$ and
$\iota(1,\delta)\gamma^{-1}\Delta=\Delta A$, so $\iota(1,\delta)\gamma^{-1}\in P_Y$,
and its inverse $\gamma\,\iota(1,\delta)^{-1}$ lies in $P_Y\cap G_Y(F)$: the
coset of $\gamma$ contains $\iota(1,\delta)$. If
$\iota(1,\delta)\iota(1,\delta')^{-1}\in P_Y$, then
$\operatorname{diag}(1,\delta\delta'^{-1})\Delta=\Delta t$ forces $t=1$ and
$\delta=\delta'$. Hence, for $\operatorname{Re}s>t$, with $t$ from (2),
$\hat E(s,x)=\sum_{\delta\in U_Y(F)}f_s(\iota(1,\delta)x)$.

\emph{Invariance.} For $s\in U$, $\gamma,\delta\in U_Y(F)$ and $x\in G_Y$,
$\hat E(s,\iota(\gamma,\delta)x)=\hat E(s,x)$. For $\operatorname{Re}s>t$,
$f_s(\iota(1,\eta)\iota(\gamma,\delta)x)=f_s(\iota(1,\gamma^{-1}\eta\delta)x)$ by
the facts, and $\eta\mapsto\gamma^{-1}\eta\delta$ permutes $U_Y(F)$; both sides
are holomorphic on the connected set $U$, so they agree on $U$.

\emph{Holomorphy.} Let $S\subset U$ be compact. For $h,h'\in U_Y$ choose
$\gamma,\delta\in U_Y(F)$ with $\gamma h,\delta h'\in\Omega$; by the invariance,
$|\hat E(s,\iota(h,h'))|$ is at most the maximum of $|\hat E|$ on the compact set
$S\times\iota(\Omega\times\Omega)$, finite by (3). So the integrand of $P$ is
bounded on $S\times\mathcal D\times\mathcal D$, holomorphic in $s$, and
$\mathcal D\times\mathcal D$ has finite measure; differentiation under the
integral sign shows that $P(\cdot,\phi)$ is holomorphic on $U$.

\emph{Integrability.} Let $t_0$ be given by
Lemma~\ref{asm-theta-agree-doubling-integrable} for $\mu$, and put
$t_1=\max(t,t_0+\operatorname{Re}s_0-\operatorname{Re}z)$. Let
$\operatorname{Re}s>t_1$. Then $s\in U$. For $x=pk$ with $p\in P_Y$ and $k\in K$,
$|f_s(x)|=|\det t(p)|^{\operatorname{Re}(s-s_0+z)}|f_s(k)|\leq C\,a(x)^{\operatorname{Re}(s-s_0+z)}$,
with $C$ the maximum of $|f_s|$ on the compact group $K$. As
$\operatorname{Re}(s-s_0+z)\geq t_0$, $w_s$ is $\mu$-integrable.

\emph{Unfolding.} Let $h\in U_Y$, and
$k(g)=w_s(g)\phi(hg)$, integrable as $\phi$ is bounded. By the invariance of
$\mu$ and as $\mathcal D$ is a fundamental domain,
\[
\int_{U_Y}k\,d\mu=\int_{U_Y}k(h^{-1}g)\,d\mu(g)=\int_{\mathcal D}\sum_{\delta\in U_Y(F)}k(h^{-1}\delta h')\,d\mu(h').
\]
By the facts, with $\iota(h,h)\iota(1,h^{-1}\delta h')=\iota(1,\delta)\iota(h,h')$
and $\phi(\delta h')=\phi(h')$,
$k(h^{-1}\delta h')=f_s(\iota(1,\delta)\iota(h,h'))\chi(\det h')^{-1}\phi(h')$.
By the cosets, the sum over $\delta$ is
$\hat E(s,\iota(h,h'))\chi(\det h')^{-1}\phi(h')$. So
\[
\int_{\mathcal D}\hat E(s,\iota(h,h'))\chi(\det h')^{-1}\phi(h')\,d\mu(h')=\int_{U_Y}w_s(g)\phi(hg)\,d\mu(g).
\]
Multiplying by $\overline{\phi(h)}$, integrating over $h\in\mathcal D$ and
exchanging the integrals, which is allowed as $w_s$ is integrable and $\phi$ is
bounded, gives the formula for $P(s,\phi)$.
\end{proof}
% Lean: kpair_cont (ContinuousFlat.lean): iv (An.ident with Law.eis_inv through Law.hasSum,
% orbit and Law.diag), hol (An.diffOn_integral with c.cov), Law.bound and the integrability
% i1, i2 from HgtInt, and An.far_right with unfold0 (An.unfold, Law.hasSum). The facts are
% aniso_of, Law.rat, Law.diag, detInv_rat; the Cpt record is CptAll, Lemma
% asm-theta-agree-definite-compact.

\begin{lemma}\label{asm-theta-agree-nonvanishing-zeta-local-dichotomy}
Let $k$ be a nonarchimedean local field with ring of integers $\mathcal O$ and
maximal ideal $\mathfrak m$, and let $\psi$ be a character of $k$ such that
$\psi(x\mathcal O)=1$ if and only if $x\in\mathcal O$. Let $n\leq d$, let
$a\in\mathfrak m$ be nonzero, and write $\Lambda_{m,m'}=a\,M_{m\times m'}(\mathcal O)$.
Let $y\in M_{n\times d}(\mathcal O)$ and $r\in M_{d\times n}(\mathcal O)$ with
$yr=1$.
\begin{enumerate}
\item If $h\in M_n(k)$ has $h-1\in\Lambda_{n,n}$ and $Y\in M_{n\times d}(k)$ has
  $Y-y\in\Lambda_{n,d}$, then $hY-y\in\Lambda_{n,d}$.
\item Let $h\in GL_n(k)$. Then either $h-1$ and $h^{-1}-1$ lie in
  $\Lambda_{n,n}$, or there is a $k$-linear form $l$ on $M_{n\times d}(k)$ with
  $\psi(l(Y))=\psi(l(hY))=1$ for every $Y\in\Lambda_{n,d}$ and
  $\psi(l(hy-y))\neq1$.
\end{enumerate}
\end{lemma}
\leanref{Development/WeilBound/ThetaLifts/LatticeBasics/MemLam.lean}{GlobalTheory.ThK.Dbl.Bmp.keep}
\leanref{Development/WeilBound/ThetaLifts/LatticeBasics/MemLam.lean}{GlobalTheory.ThK.Dbl.Bmp.dich}
\leanref{Development/WeilBound/ThetaLifts/LatticeBasics/MemLam.lean}{GlobalTheory.ThK.Dbl.Bmp.sep}
\begin{proof}
(1) Write $h=1+a\alpha$ and $Y=y+a\beta$ with $\alpha$ and $\beta$ integral.
Then $hY-y=a(\beta+\alpha y+a\alpha\beta)$.

(2) The set $L=\Lambda_{n,d}+h\Lambda_{n,d}$ is an $\mathcal O$-lattice in
$M_{n\times d}(k)$, as $a\neq0$ and $h$ is invertible. Suppose first that
$hy-y\in L$, say $hy-y=aA_1+haA_2$ with $A_1,A_2\in M_{n\times d}(\mathcal O)$.
Multiplying on the right by $r$ gives $h-1=aA_1r+h\,aA_2r$, that is, $hP=Q$ with
$P=1-aA_2r$ and $Q=1+aA_1r$. These matrices have integral entries and
determinant $1$ modulo $\mathfrak m$, as $a\in\mathfrak m$, so they lie in
$GL_n(\mathcal O)$. Hence $h=QP^{-1}$,
$h-1=(Q-P)P^{-1}=a(A_1+A_2)rP^{-1}$ and
$h^{-1}-1=(P-Q)Q^{-1}=-a(A_1+A_2)rQ^{-1}$, both in $\Lambda_{n,n}$. Suppose now
that $hy-y\notin L$. An $\mathcal O$-basis of $L$ is a $k$-basis of
$M_{n\times d}(k)$, and some coordinate $x$ of $hy-y$ in it is not in
$\mathcal O$; choose $s\in\mathcal O$ with $\psi(xs)\neq1$, and let $l$ be $s$
times this coordinate. Then $l(L)\subset\mathcal O$, so $\psi(l(L))=1$; this
gives the claim, as $\Lambda_{n,d}$ and $h\Lambda_{n,d}$ lie in $L$.
\end{proof}
% Lean: keep; dich, with sep for the second case (Module.Free.chooseBasis of the lattice
% lam + map G lam, Submodule.IsLattice.sup) and isUnit_det_lemma_3 for P and Q.

\begin{lemma}\label{asm-theta-agree-nonvanishing-zeta-local-vacuum-skew}
Let $A$ be a commutative ring with an involution $c$, and $r\colon A\to\kappa$ a
surjective ring homomorphism onto a field, such that every $z\in A$ with
$r(z)\neq0$ and $r(c(z))\neq0$ is a unit, and such that one of the following
holds:
\begin{enumerate}
\item[(i)] $r(c(z))=0$ whenever $r(z)=0$, and $r(u_0-c(u_0))\neq0$ for some
  $u_0\in A$;
\item[(ii)] $r(\eta)=1$ and $r(c(\eta))=0$ for some $\eta\in A$.
\end{enumerate}
Let $G\in GL_n(A)$ with $c(G)^{\mathsf T}=G$, and call $S\in M_n(A)$
\emph{skew} if $S^{\mathsf T}G+G\,c(S)=0$. Let $M,m,P,Q\in M_n(A)$ with
$M^{\mathsf T}G\,c(m)+m^{\mathsf T}G\,c(M)=0$ and $PM+Qm=1$. Then there is a
skew $S$ such that $\det(M+Sm)$ is a unit.
\end{lemma}
\leanref{Development/WeilBound/ThetaLifts/LatticeBasics/SkewTwo.lean}{GlobalTheory.ThK.Dbl.Vf2.exists_skew}
\leanref{Development/WeilBound/ThetaLifts/LatticeBasics/SkewTwo.lean}{GlobalTheory.ThK.Dbl.Vf2.step}
\leanref{Development/WeilBound/ThetaLifts/LatticeBasics/SkewTwo.lean}{GlobalTheory.ThK.Dbl.Vf2.skew_two}
\leanref{Development/WeilBound/ThetaLifts/LatticeBasics/SkewTwo.lean}{GlobalTheory.ThK.Dbl.Vf2.rank_lt}
\begin{proof}
For a matrix or vector $N$ over $A$ write $\bar N$ for its image over $\kappa$
under $r$, and $\tilde N$ for its image under $r\circ c$. Let $e_i$ be the unit
vectors.

\emph{Two facts.} (a) If $S$ is skew, then $M_S=M+Sm$ satisfies
$M_S^{\mathsf T}G\,c(m)+m^{\mathsf T}G\,c(M_S)=0$, the difference with the
relation for $M$ being $m^{\mathsf T}(S^{\mathsf T}G+G\,c(S))c(m)$, and
$PM_S+(Q-PS)m=1$. (b) For $u,w\in A^n$ the matrix
$u\,(G\,c(w))^{\mathsf T}-w\,(G\,c(u))^{\mathsf T}$ is skew, by a direct
computation with $c(G)^{\mathsf T}=G$.

\emph{Step.} Let $M'$ satisfy $M'^{\mathsf T}G\,c(m)+m^{\mathsf T}G\,c(M')=0$
and $P'M'+Q'm=1$ for some $P',Q'$, and suppose that $\det\bar M'=0$. Then there
is a skew $S'$ with $\operatorname{rank}\overline{M'+S'm}>\operatorname{rank}\bar M'$.
Indeed, choose $y\in A^n$ with $\bar y\neq0$ and $\bar M'\bar y=0$. Some entry
$i$ has $r((my)_i)\neq0$: otherwise $\bar m\bar y=0$ and
$\bar y=\bar P'\bar M'\bar y+\bar Q'\bar m\bar y=0$. Put
$w=c(G^{-1}e_i)$, so that $G\,c(w)=e_i$, and for $u\in A^n$ let
$S_u=u\,e_i^{\mathsf T}-w\,(G\,c(u))^{\mathsf T}$, which is skew by (b).
Suppose that $\bar S_u=z\,e_i^{\mathsf T}$ for some $z\in\kappa^n$ outside the
image of $\bar M'$. Then $\overline{M'+S_um}=\bar M'+z\varphi^{\mathsf T}$,
where $\varphi^{\mathsf T}$ is the $i$-th row of $\bar m$, and
$\varphi^{\mathsf T}\bar y=r((my)_i)\neq0$. The image of
$\bar M'+z\varphi^{\mathsf T}$ contains $z$, the image of
$(\varphi^{\mathsf T}\bar y)^{-1}\bar y$, and hence the image of $\bar M'$; as
$z$ is not in the latter, the rank grows. It remains to find $u$.

In case (i), take $u=u_0w$. Then
$\bar S_u=r(u_0-c(u_0))\,\bar w\,e_i^{\mathsf T}$, and
$z=r(u_0-c(u_0))\bar w$ is not in the image of $\bar M'$. Otherwise
$\bar w=\bar M'\bar x$ for some $x\in A^n$, so $r$ kills the entries of
$w-M'x$, and by (i) so does $r\circ c$. Multiplying the relation for $M'$ by
$y^{\mathsf T}$ on the left and $c(x)$ on the right gives
$(M'y)^{\mathsf T}G\,c(mx)+(my)^{\mathsf T}G\,c(M'x)=0$, so
\[
(my)_i=(my)^{\mathsf T}G\,c(w)=(my)^{\mathsf T}G\,c(w-M'x)-(M'y)^{\mathsf T}G\,c(mx),
\]
and $r((my)_i)=0$, as $\overline{M'y}=0$; a contradiction. In case (ii), choose
$z_0\in\kappa^n$ outside the image of $\bar M'$, which is not onto, lift it to
$z\in A^n$, and take $u=\eta z$: as $r(\eta)=1$ and $r(c(\eta))=0$,
$\bar S_u=\bar z\,e_i^{\mathsf T}$.

\emph{Conclusion.} Starting from $S=0$ and applying the step to $M_S$, with
$P'=P$ and $Q'=Q-PS$ by (a), at most $n+1$ times, as ranks are at most $n$,
gives a skew $S$ with $\det\overline{M_S}\neq0$, that is,
$r(\det M_S)\neq0$. Also $r(c(\det M_S))=\det\tilde M_S\neq0$. Otherwise there is
$y\in\kappa^n$, $y\neq0$, with $\tilde M_Sy=0$. Applying $r$ to the relation in
(a) and multiplying by $y$ gives $\overline{M_S}^{\mathsf T}\bar G\,\tilde my=0$, so
$\tilde my=0$, as $\overline{M_S}$ and $\bar G$ are invertible; and applying
$r\circ c$ to $PM_S+(Q-PS)m=1$ gives $y=\tilde P\tilde M_Sy+\widetilde{(Q-PS)}\tilde my=0$,
a contradiction. So $\det M_S$ is a unit.
\end{proof}
% Lean: exists_skew (SkewTwo.lean), by induction on k with S' from step (ind); (a) is l2 and
% i2, (b) is skew_two, the rank increase is rank_lt, and the two cases of step are the two
% branches of hcase (fin, with e for case (i)). The last paragraph is the end of exists_skew.

\begin{lemma}\label{asm-theta-agree-nonvanishing-zeta-local-vacuum-columns}
Let $B$ be a theta base of a CM field $E$ with complex conjugation $c$, let $X$
be a Hermitian lattice space of rank $n$ and some signature, with matrix $J$
and lattice $L$, let $G_X$ and $\Delta$ be as in
Definition~\ref{asm-theta-agree-doubling}(2), and let $v$ be a finite place.
Let $L_v$ be as in Lemma~\ref{asm-theta-agree-local-vacuum-lattice}, and say
that an $n\times n$ matrix $N$ over $\mathbf A_E$ \emph{preserves} $L_v$ if
$NL_v\subset L_v$. Let $g\in G_X$ with
$g\Delta=\left(\begin{smallmatrix}x\\y\end{smallmatrix}\right)$, and let $a_1$
and $a_2$ be $n\times n$ matrices over $\mathbf A_E$, such that $x$, $y$, $a_1$
and $a_2$ preserve $L_v$, $a_1x+a_2y=1$, and $x-1$ and $y-1$ are at $v$. Then
there are $n\times n$ matrices $b$, $b'$ and $t'$ over $\mathbf A_E$ with
$b^{\mathsf T}J=J\,c(b')$, $b+b'=-1$, $b$ and $t'$ preserving $L_v$, and
$(y+b(y-x))t'=1$.
\end{lemma}
\leanref{Development/WeilBound/ThetaLifts/LatticeBasics/CompCA.lean}{GlobalTheory.ThK.Dbl.Vf2.exists_beta}
\leanref{Development/WeilBound/ThetaLifts/LatticeBasics/CompCA.lean}{GlobalTheory.ThK.Dbl.Vf2.exists_res}
\leanref{Development/WeilBound/ThetaLifts/LatticeBasics/CompCA.lean}{GlobalTheory.ThK.Dbl.Vf2.fr_iff}
\leanref{Development/WeilBound/ThetaLifts/LatticeBasics/CompCA.lean}{GlobalTheory.ThK.Dbl.Vf2.gram_fr}
\leanref{Development/WeilBound/ThetaLifts/LatticeBasics/SkewTwo.lean}{GlobalTheory.ThK.Dbl.Vf2.lag_of}
\begin{proof}
Put $E_v=E\otimes_FF_v$ and $\mathcal O_{E,v}=\mathcal O_E\otimes\mathcal O_v$,
on which $c$ acts, and for a matrix $N$ over $\mathbf A_E$ let $N_v$ be its
component over $E_v$.

\emph{The lattice.} There is $a\in GL_n(E)$ whose columns form a basis of the
$\mathcal O_{E,v}$-lattice $L\otimes\mathcal O_v\subset E_v^n$. As $L_v$
consists of the vectors at $v$ with component in $L\otimes\mathcal O_v$, a
matrix $N$ over $\mathbf A_E$ preserves $L_v$ if and only if $a^{-1}N_va$ has
entries in $\mathcal O_{E,v}$. As $X$ is Hermitian, $L$ is its own
dual for $h(y,x)=y^{\mathsf T}J\,c(x)$, so the Gram matrix
$G=a^{\mathsf T}J\,c(a)$ and its inverse have entries in $\mathcal O_{E,v}$; and
$c(G)^{\mathsf T}=G$. For matrices $N$ and $N'$ over $\mathbf A_E$,
$a^{\mathsf T}(N^{\mathsf T}J\,c(N'))_vc(a)=(a^{-1}N_va)^{\mathsf T}G\,c(a^{-1}N'_va)$.

\emph{The residue map.} There are a field $\kappa$ and a surjective ring
homomorphism $r\colon\mathcal O_{E,v}\to\kappa$ that satisfy the hypotheses of
Lemma~\ref{asm-theta-agree-nonvanishing-zeta-local-vacuum-skew} for the
involution $c$. If there is an $F$-embedding $\tau\colon E\to F_v$, then
$\mathcal O_{E,v}\cong\mathcal O_v\times\mathcal O_v$ through
$(\tau,\tau\circ c)$, and $c$ exchanges the factors; let $r$ be the residue map
of the first factor. Then $\eta=(1,0)$ gives (ii), and an element whose two
components have nonzero residues is a unit. Otherwise, as $E/F$ is unramified at
the finite places, $v$ is inert in $E$, $\mathcal O_{E,v}=\mathcal O_w$ for the
place $w$ above $v$, and $c$ induces the nontrivial automorphism of the residue
field $\kappa$ of $w$; let $r$ be the residue map. Then (i) holds, with $u_0$
whose residue is not fixed, and every $z$ with $r(z)\neq0$ is a unit.

\emph{A first choice.} As $E/F$ is unramified at the finite places, there is
$\theta\in\mathcal O_E$ with $\theta+c(\theta)=1$. The scalar matrices
$b_0=-\theta$ and $b'_0=-c(\theta)$ have $b_0^{\mathsf T}J=J\,c(b'_0)$ and
$b_0+b'_0=-1$, and $b_0$ preserves $L_v$
(Lemma~\ref{asm-theta-agree-local-vacuum-lattice}(3)). Put $m=y-x$ and
$M=y+b_0m$. As $g\in G_X$, $\Delta^{\mathsf T}J^\square c(\Delta)=0$ gives
$x^{\mathsf T}J\,c(x)=y^{\mathsf T}J\,c(y)$, with $J^\square=\operatorname{diag}(J,-J)$;
with $b_0^{\mathsf T}J=-J-J\,c(b_0)$, a direct computation gives
$M^{\mathsf T}J\,c(m)+m^{\mathsf T}J\,c(M)=y^{\mathsf T}J\,c(y)-x^{\mathsf T}J\,c(x)=0$.
Also $(a_1+a_2)M+(-a_1-(a_1+a_2)b_0)m=a_1x+a_2y=1$.

\emph{The skew correction.} The matrices $N_M=a^{-1}M_va$, $N_m=a^{-1}m_va$,
$N_P=a^{-1}(a_1+a_2)_va$ and $N_Q=a^{-1}(-a_1-(a_1+a_2)b_0)_va$ have entries in
$\mathcal O_{E,v}$, as $M$, $m$, $a_1+a_2$ and $-a_1-(a_1+a_2)b_0$ preserve
$L_v$. By the last formula of the first paragraph,
$N_M^{\mathsf T}G\,c(N_m)+N_m^{\mathsf T}G\,c(N_M)=0$, and $N_PN_M+N_QN_m=1$.
Lemma~\ref{asm-theta-agree-nonvanishing-zeta-local-vacuum-skew} gives $S$ over
$\mathcal O_{E,v}$ with $S^{\mathsf T}G+G\,c(S)=0$ and $\det(N_M+SN_m)$ a unit.
Let $S_A$ be the matrix over $\mathbf A_E$ that is at $v$ with component
$aSa^{-1}$ there. Then $S_A$ preserves $L_v$, and
$S_A^{\mathsf T}J+J\,c(S_A)=0$, as $a^{\mathsf T}$ and $c(a)$ turn its component
at $v$ into $S^{\mathsf T}G+G\,c(S)$. Put $b=b_0+S_A$ and $b'=b'_0-S_A$. Then
$b^{\mathsf T}J=J\,c(b')$, $b+b'=-1$, and $b$ preserves $L_v$.

\emph{The inverse.} The matrix $f=y+bm=M+S_Am$ has $f-1$ at $v$, as $y-1$ and
$m$ are, and $a^{-1}f_va=N_M+SN_m$ has a unit determinant. So $f_v$ is
invertible, and the matrix $t'$ that is $1$ away from $v$ and $f_v^{-1}$ at $v$
has $ft'=1$; it preserves $L_v$, as $a^{-1}t'_va=(N_M+SN_m)^{-1}$ has entries in
$\mathcal O_{E,v}$.
\end{proof}
% Lean: exists_beta (CompCA.lean). The lattice: Frames.exists_rational_frame for a = bc iota P,
% fr_iff, Hgt.gram_loc (with hX.dual) for G and H, and gram_fr; the residue map is exists_res
% (KFinLoc.split_or_inert, splitIdem, existsUnique_isConjAt, exists_isUnramInv); theta is
% exists_int_add_conj_eq_one_of_base_good; the relation is lag_of (SkewTwo.lean); S is
% exists_skew; S_A is SA, with sk; the inverse is exists_inv (MemIntLocOfIdem.lean).

\begin{lemma}\label{asm-theta-agree-nonvanishing-zeta-local-vacuum-decomposition}
Let $E$ be a CM field with complex conjugation $c$, let $X$ be a lattice space
of rank $n$ with invertible matrix $J$, and let $G_X$, $P_X$ and $\Delta$ be as
in Definition~\ref{asm-theta-agree-doubling}(2). Let $g\in G_X$ with
$g\Delta=\left(\begin{smallmatrix}x\\y\end{smallmatrix}\right)$, put $m=y-x$,
let $b$, $b'$ and $t'$ be $n\times n$ matrices over $\mathbf A_E$ with
$b^{\mathsf T}J=J\,c(b')$, $b+b'=-1$ and $(y+bm)t'=1$, and let
$\varepsilon\in\mathbf A_E$ with $\varepsilon^2=\varepsilon$,
$c(\varepsilon)=\varepsilon$ and $\varepsilon m=m$. Let $C$ and $S$ be as in
Lemma~\ref{asm-theta-agree-local-vacuum-decomposition}, and put
\[
T=-mt',\qquad t=y+bm,\qquad s=CSC^{-1},\qquad n_0=C\begin{pmatrix}1&T\\0&1\end{pmatrix}C^{-1}.
\]
Then $t$ is invertible with $t^{-1}=t'$; $s$ and $n_0$ lie in $G_X$; $s^2=1$;
$n_0^{-1}=C\left(\begin{smallmatrix}1&-T\\0&1\end{smallmatrix}\right)C^{-1}$ and
$n_0\Delta=\Delta$; and $p=(sn_0s)^{-1}g$ lies in $G_X$ and has $p\Delta=\Delta t$.
\end{lemma}
\leanref{Development/WeilBound/ThetaLifts/LatticeBasics/SkewTwo.lean}{GlobalTheory.ThK.Dbl.Vf2.exists_decomp}
\begin{proof}
As $t$ and $t'$ are square matrices over a commutative ring with $tt'=1$, also
$t't=1$. As in the proof of
Lemma~\ref{asm-theta-agree-local-vacuum-decomposition}, with
$J'=\left(\begin{smallmatrix}0&J\\J&0\end{smallmatrix}\right)$ and
$J^\square=\operatorname{diag}(J,-J)$: $C^{\mathsf T}J^\square c(C)=J'$, so a
matrix lies in $G_X$ when its conjugate by $C^{-1}$ preserves $J'$; $S$
preserves $J'$ and $S^2=1$, so $s\in G_X$ and $s^2=1$.

As $g\in G_X$, $x^{\mathsf T}J\,c(x)=y^{\mathsf T}J\,c(y)$, and as in the proof
of Lemma~\ref{asm-theta-agree-nonvanishing-zeta-local-vacuum-columns},
$t^{\mathsf T}J\,c(m)+m^{\mathsf T}J\,c(t)=0$. As $Tt=-m$,
\[
t^{\mathsf T}\bigl(T^{\mathsf T}J+J\,c(T)\bigr)c(t)=-m^{\mathsf T}J\,c(t)-t^{\mathsf T}J\,c(m)=0,
\]
and $t$ is invertible, so $T^{\mathsf T}J+J\,c(T)=0$. Hence
$\left(\begin{smallmatrix}1&T\\0&1\end{smallmatrix}\right)$ preserves $J'$, and
$n_0\in G_X$. The formula for $n_0^{-1}$ is clear, and $n_0\Delta=\Delta$ as
$C^{-1}\Delta=\left(\begin{smallmatrix}1\\0\end{smallmatrix}\right)$ and
$C\left(\begin{smallmatrix}1\\0\end{smallmatrix}\right)=\Delta$.

As $\varepsilon T=T$,
$S\left(\begin{smallmatrix}1&T\\0&1\end{smallmatrix}\right)S=\left(\begin{smallmatrix}1&0\\T&1\end{smallmatrix}\right)$,
so $sn_0s\Delta=C\left(\begin{smallmatrix}1\\T\end{smallmatrix}\right)$, with
blocks $1+(1+b)T$ and $1+bT$. As $Tt=-m$, $(1+bT)t=t-bm=y$ and
$(1+(1+b)T)t=t-(1+b)m=x$. So $sn_0s\Delta t=g\Delta$, that is, $p\Delta=\Delta t$,
and $p\in G_X$.
\end{proof}
% Lean: Vf2.exists_decomp (SkewTwo.lean), with mkU (conjugation by CC b = C), hS and hSu for
% SS eps = S, sk for T (from lag_of), hU and hUu for n_0, hL (S n S), hD and hD', and e2 for
% p Delta = Delta t.

\begin{lemma}\label{asm-theta-agree-nonvanishing-zeta-local}
Let $(X,W,M,\chi_X,\chi_W,\omega)$ be a Weil pair for a theta base $B$ of a CM
field $E$ (Definition~\ref{asm-theta-agree-triple}(1)), with $X$ and $W$ of
ranks $n$ and $d$, let $\rho$ be the Schrödinger operators of $M$, and put
$I=\{1,\ldots,n\}\times\{1,\ldots,d\}$.
\begin{enumerate}
\item Suppose that $n\leq d$. Then for every $v\in P$ there is a nonzero
  Schwartz-Bruhat function $\beta_v$ on $F_v^I$ such that every $g\in U_X$ at
  $v$ (Definition~\ref{asm-theta-agree-models}(2)) satisfies one of the
  following: $g$ is integral at $v$
  (Definition~\ref{asm-theta-agree-nonvanishing-doubling}(2)) and $\omega(g,1)$
  acts at $v$ by $\beta_v\mapsto\beta_v$; or there are $u\in M_{n\times d}(\mathbf A_E)$
  at $v$ and $c\neq c'$ in $\mathbf C$ such that $\rho(u)$ acts at $v$ by
  $\beta_v\mapsto c\beta_v$ and $\rho(gu)$ acts at $v$ by
  $\beta_v\mapsto c'\beta_v$.
\item Let $(I\sqcup I,e^\square,\omega_2,S,l)$ be a doubling datum for $(W,X)$
  with character $\chi_W$ and the isomorphism $e^\square$ of
  Lemma~\ref{asm-theta-agree-doubling-datum}(2), let $K$ be the lattice
  stabilizer of the doubled space of $X$, and let $v$ be a finite place such
  that $\chi_W$ is unramified at every place of $E$ above $v$. Then for every
  $k\in K$ with $k-1$ at $v$, $\omega_2(k)$ acts at $v$ by
  $\mathrm{vac}_v\boxtimes\overline{\mathrm{vac}_v}\mapsto\mathrm{vac}_v\boxtimes\overline{\mathrm{vac}_v}$.
\end{enumerate}
\end{lemma}
\leanref{Development/WeilBound/ThetaLifts/LatticeBasics/BadStmt.lean}{GlobalTheory.ThK.Dbl.Bump}
\leanref{Development/WeilBound/ThetaLifts/LatticeBasics/NearUnits.lean}{GlobalTheory.ThK.Dbl.bump}
\leanref{Development/WeilBound/ThetaLifts/LatticeBasics/BadStmt.lean}{GlobalTheory.ThK.Dbl.VacFix2}
\leanref{Development/WeilBound/ThetaLifts/LatticeBasics/PresDOf.lean}{GlobalTheory.ThK.Dbl.vacFix2}
\leanref{Development/WeilBound/ThetaLifts/LatticeBasics/PresDOf.lean}{GlobalTheory.ThK.Dbl.Vf2.presD_of}
\leanref{Development/WeilBound/ThetaLifts/LatticeBasics/SchrUnit.lean}{GlobalTheory.ThK.Dbl.Bad.norm_eig}
\begin{proof}
(1) Let $v\in P$, let $t\colon E\to F_v$ be the $F$-embedding that $B$ gives
at $v$, and write $\mathcal O$ for the ring of integers of $F_v$, $\mathfrak m$
for its maximal ideal, $|\cdot|$ for its normalized absolute value,
$\nu=(\chi_W)_t$ (Definition~\ref{asm-theta-agree-coinvariants-local}(4)), and
$\psi_v(s)=\psi(\bar s)$ for $s\in F_v$, where $\bar s\in\mathbf A$ has
component $s$ at $v$ and $0$ elsewhere. As $B$ is a theta base and $v\in P$,
$\psi_v(x\mathcal O)=1$ if and only if $x\in\mathcal O$. Let $\iota$ be a
linear model of the second kind for $(t,1,1)$
(Lemma~\ref{asm-theta-agree-coinvariants-dictionary}(1)), and put
$A=t(J_X)$ and $C=t(J_W)$, which are invertible, as in the proof of that
lemma.

\emph{Step 1: a level.} There is a nonzero $a\in\mathfrak m$ such that every
$g\in U_X^{(v)}$ for which $\pi_t(g)^{\mathsf T}-1$ and
$\pi_t(g^{-1})^{\mathsf T}-1$ lie in $a\,M_n(\mathcal O)$ is integral at $v$
and has $\nu(\det\pi_t(g))=1$ and $|\det\pi_t(g)|=1$. Indeed, for
$g\in U_X^{(v)}$, Lemma~\ref{asm-theta-agree-dictionary-split}(4) gives
$\pi_{t\circ c}(g)=A^{-1}\pi_t(g^{-1})^{\mathsf T}A$. By part (1) of that lemma
there is a continuous map $S$ from $M_n(F_v)^2$ to $M_n(\mathbf A_E)$, with
$S(1,1)=1$, such that $S(p,q)-1$ is at $v$, $\pi_t(S(p,q))=p$ and
$\pi_{t\circ c}(S(p,q))=A^{-1}q^{\mathsf T}A$; so
$g=S(\pi_t(g),\pi_t(g^{-1}))$ and $g^{-1}=S(\pi_t(g^{-1}),\pi_t(g))$. The
elements of $U_X$ integral at $v$ form an open set containing $1$, as the
condition is that the component at $v$ and its inverse preserve the closure
of $L_X$ at $v$, a compact open subgroup; and the $g_0\in GL_n(F_v)$ with
$\nu(\det g_0)=1=|\det g_0|$ form an open subgroup, as $\nu$ is smooth. As
$g\mapsto(g,g^{-1})$ identifies $U_X$ and $GL_n(F_v)$ with subspaces of
$M_n(\mathbf A_E)^2$ and $M_n(F_v)^2$, there is a neighbourhood $N$ of $(1,1)$
in $M_n(F_v)^2$ such that every $g\in U_X^{(v)}$ with
$(\pi_t(g),\pi_t(g^{-1}))\in N$ has the three properties. As
$M_n(\mathcal O)^2$ is compact and $a_0^j\to0$ for a nonzero $a_0\in\mathfrak m$,
there is $j\geq1$ such that $N$ contains the $(p,q)$ with $p-1$ and $q-1$ in
$a_0^jM_n(\mathcal O)$. Take $a=a_0^j$; transposition preserves
$a\,M_n(\mathcal O)$.

\emph{Step 2: the function.} Let $y\in M_{n\times d}(\mathcal O)$ have entries
$1$ at the places $(i,i)$ and $0$ elsewhere, so that $yy^{\mathsf T}=1$, as
$n\leq d$. Let $\beta_v$ be the Schwartz-Bruhat function with $\iota\beta_v$
the indicator function of the compact open set $y+a\,M_{n\times d}(\mathcal O)$.
It is nonzero, as $\iota\beta_v(y)=1$.

\emph{Step 3: Schrödinger operators.} Let $u\in M_{n\times d}(\mathbf A_E)$ be at
$v$ with $\pi_{t\circ c}(u)=0$, and let $c\in\mathbf C$ with
$\psi_v\bigl(-\operatorname{tr}(Y^{\mathsf T}\pi_t(u))\bigr)=c$ for every
$Y\in y+a\,M_{n\times d}(\mathcal O)$. Then $\rho(u)$ acts at $v$ by
$\beta_v\mapsto c\beta_v$. Indeed, as $u$ is at $v$, $\rho(u)$ acts at $v$ by
$\beta_v\mapsto\beta'$ for a Schwartz-Bruhat $\beta'$. By
Lemma~\ref{asm-theta-agree-dictionary-split}(1), every
$Y\in M_{n\times d}(F_v)$ is $-A\,\pi_{t\circ c}(x)\,C^{\mathsf T}$ for some $x$
at $v$ with $\pi_t(x)=0$, and this is the map $\operatorname{co}(x)$ of a
linear model of the second kind for $(t,1,1)$
(Definition~\ref{asm-theta-agree-models}(5)). For such $x$, with $h$ the form
of Definition~\ref{asm-theta-agree-models}(3),
$\pi_t(h(x,u))=0$ and, as $t(c(J_X))=A^{\mathsf T}$, $t(c(J_W))=-C^{\mathsf T}$
and $\pi_{t\circ c}(c(u))=\pi_t(u)$,
$\pi_{t\circ c}(h(x,u))=-\operatorname{tr}\bigl(\pi_{t\circ c}(x)^{\mathsf T}A^{\mathsf T}\pi_t(u)\,C\bigr)$;
as the trace of an element $z$ at $v$ is $\pi_t(z)+\pi_{t\circ c}(z)$ placed at
$v$, $\psi(\langle x,u\rangle)=\psi_v\bigl(-\operatorname{tr}(Y^{\mathsf T}\pi_t(u))\bigr)$.
So the definition of a linear model gives
$\iota\beta'(Y)=\psi_v\bigl(-\operatorname{tr}(Y^{\mathsf T}\pi_t(u))\bigr)\iota\beta_v(Y)$
for every $Y$, which is $c\,\iota\beta_v(Y)$, as $\iota\beta_v$ vanishes off
$y+a\,M_{n\times d}(\mathcal O)$. As $\iota$ is injective on the
Schwartz-Bruhat functions, $\beta'=c\beta_v$.

\emph{Step 4: the dichotomy.} Let $g\in U_X$ be at $v$, and put
$g_0=\pi_t(g)$. Apply
Lemma~\ref{asm-theta-agree-nonvanishing-zeta-local-dichotomy}(2) to $F_v$,
$\psi_v$, $a$, $y$, $r=y^{\mathsf T}$ and $h=g_0^{\mathsf T}$.

In the first case $g_0^{\mathsf T}-1$ and
$(g_0^{\mathsf T})^{-1}-1=\pi_t(g^{-1})^{\mathsf T}-1$ lie in
$a\,M_n(\mathcal O)$, so by Step 1, $g$ is integral at $v$,
$\nu(\det g_0)=1$ and $|\det g_0|=1$. By
Lemma~\ref{asm-theta-agree-coinvariants-dictionary}(3), $\omega(g,1)$ acts at
$v$ by $\beta_v\mapsto\beta'$ with
$\iota\beta'(Y)=\nu(\det g_0)|\det g_0|^{d/2}\iota\beta_v(g_0^{\mathsf T}Y)=\iota\beta_v(g_0^{\mathsf T}Y)$.
By part (1) of the dichotomy lemma, for $h=g_0^{\mathsf T}$ and for its inverse,
$g_0^{\mathsf T}Y-y\in a\,M_{n\times d}(\mathcal O)$ if and only if
$Y-y\in a\,M_{n\times d}(\mathcal O)$. So $\iota\beta'=\iota\beta_v$, and
$\beta'=\beta_v$.

In the second case there is a linear form $l$ with
$\psi_v(l(Y))=\psi_v(l(g_0^{\mathsf T}Y))=1$ for
$Y\in a\,M_{n\times d}(\mathcal O)$ and $\psi_v(l(g_0^{\mathsf T}y-y))\neq1$.
Choose $u_0\in M_{n\times d}(F_v)$ with
$l(Y)=-\operatorname{tr}(Y^{\mathsf T}u_0)$ for all $Y$, and, by
Lemma~\ref{asm-theta-agree-dictionary-split}(1), $u$ at $v$ with $\pi_t(u)=u_0$
and $\pi_{t\circ c}(u)=0$. Put $c=\psi_v(l(y))$ and
$c'=\psi_v(l(g_0^{\mathsf T}y))$; then $c'/c=\psi_v(l(g_0^{\mathsf T}y-y))\neq1$.
For $Y\in y+a\,M_{n\times d}(\mathcal O)$,
$\psi_v(-\operatorname{tr}(Y^{\mathsf T}u_0))=\psi_v(l(y))\,\psi_v(l(Y-y))=c$,
so $\rho(u)$ acts at $v$ by $\beta_v\mapsto c\beta_v$ by Step 3. The element
$gu$ is at $v$, with $\pi_{t\circ c}(gu)=0$ and $\pi_t(gu)=g_0u_0$, and
$-\operatorname{tr}(Y^{\mathsf T}g_0u_0)=l(g_0^{\mathsf T}Y)$, so for the same
$Y$,
$\psi_v(-\operatorname{tr}(Y^{\mathsf T}g_0u_0))=\psi_v(l(g_0^{\mathsf T}y))\,\psi_v(l(g_0^{\mathsf T}(Y-y)))=c'$,
and $\rho(gu)$ acts at $v$ by $\beta_v\mapsto c'\beta_v$ by Step 3.

(2) Let $J=J_X$, which is invertible as $\tau_w(J)$ has a signature, let
$L_v$, $\Lambda_v$ and the matrices that preserve $L_v$ be as in
Lemma~\ref{asm-theta-agree-local-vacuum-lattice}, let $1_v\in\mathbf A$ be $1$
at $v$ and $0$ elsewhere, and let $\varepsilon$ be its image in $\mathbf A_E$.
Let $x_0$ and $r_0$ be as in
Lemma~\ref{asm-theta-agree-finite-inputs-slice}(4), for $I$, and put
$\Phi=\mathrm{vac}_v\otimes r_0$ and $\Phi^\square=\Phi\boxtimes\overline\Phi$.
As in Step 3 of the proof of Lemma~\ref{asm-theta-agree-local-factor-vacuum},
$\Phi^\square$ is a nonzero Schwartz-Bruhat function on $\mathbf A^{I\sqcup I}$,
equal to
$(\mathrm{vac}_v\boxtimes\overline{\mathrm{vac}_v})\otimes(r_0\boxtimes\overline{r_0})$,
and $r_0\boxtimes\overline{r_0}$ does not depend on the $v$-components. Let
$k\in K$ with $k-1$ at $v$.

\emph{Step 1: the blocks.} An $n\times n$ matrix $N$ over $\mathbf A_E$ that maps
$\hat L_X$ into itself preserves $L_v$: as in the proof of
Lemma~\ref{asm-theta-agree-local-vacuum-lattice}(1), $L_v\subset\hat L_X$ and
$1_v\hat L_X\subset L_v$, and for $y\in L_v$ the vector $Ny$ is at $v$, so
$Ny=1_vNy\in L_v$. As $k$ and $k^{-1}$ map $\hat L_X\times\hat L_X$ into itself,
their four $n\times n$ blocks map $\hat L_X$ into itself, and so preserve $L_v$.

\emph{Step 2: eigenvectors.} The operators $\omega_2(x)$, $x\in G_X$, satisfy
condition (b) of Definition~\ref{asm-theta-agree-doubling}(6) for the
isomorphism $e^\square$ of Lemma~\ref{asm-theta-agree-doubling-datum}(2), which
is what the proof of Lemma~\ref{asm-theta-agree-local-vacuum-doubled} uses of
its operators $\omega^\square$. So for every $g\in G_X$ with $g-1$ at $v$ whose
inverse has four blocks that preserve $L_v$, $\omega_2(g)$ acts at $v$ by
$\mathrm{vac}_v\boxtimes\overline{\mathrm{vac}_v}\mapsto a_g\,\mathrm{vac}_v\boxtimes\overline{\mathrm{vac}_v}$
for some $a_g\in\mathbf C$, and then $\omega_2(g)\Phi^\square=a_g\Phi^\square$.
By Step 1 this applies to $g=k$; let $c=a_k$. It remains to show $c=1$.

\emph{Step 3: a decomposition.} Write $k$ and $k^{-1}$ in $n\times n$ blocks
$k_{ij}$ and $k'_{ij}$, and put $x=k_{11}+k_{12}$ and $y=k_{21}+k_{22}$, so that
$k\Delta=\left(\begin{smallmatrix}x\\y\end{smallmatrix}\right)$. From
$k^{-1}k\Delta=\Delta$, $k'_{11}x+k'_{12}y=1$; $x-1$ and $y-1$ are at $v$, as
$k-1$ is; and $x$, $y$, $k'_{11}$ and $k'_{12}$ preserve $L_v$ by Step 1. Let
$b$, $b'$ and $t'$ be given by
Lemma~\ref{asm-theta-agree-nonvanishing-zeta-local-vacuum-columns} for $k$,
$a_1=k'_{11}$ and $a_2=k'_{12}$. As $y-x$ is at $v$, $\varepsilon(y-x)=y-x$, and
$\varepsilon^2=\varepsilon=c(\varepsilon)$. Let $C$, $S$, $T$, $t$, $s$, $n_0$
and $p$ be as in
Lemma~\ref{asm-theta-agree-nonvanishing-zeta-local-vacuum-decomposition}, for
$g=k$; so $k=sn_0s\,p$.

\emph{Step 4: $s$ and $n_0$.} The blocks $1$, $1+b$, $b$, $-b$ and $-1$ of $C$
and $C^{-1}$, the blocks $\varepsilon$ and $1-\varepsilon$ of $S$, and
$(y-x)t'$ preserve $L_v$ (Lemma~\ref{asm-theta-agree-local-vacuum-lattice}(3)
for $\varepsilon$ and $1-\varepsilon$), and so do the blocks of a product of
block matrices with such blocks. Hence the blocks of $s^{-1}=s$ and of
$n_0^{-1}=C\left(\begin{smallmatrix}1&(y-x)t'\\0&1\end{smallmatrix}\right)C^{-1}$
preserve $L_v$. Also $s-1=C(S-1)C^{-1}$ and
$n_0-1=C\left(\begin{smallmatrix}0&T\\0&0\end{smallmatrix}\right)C^{-1}$ are at
$v$, as $S-1$ has entries $\pm\varepsilon$ and $T=-(y-x)t'$. By Step 2 there
are $a,a'\in\mathbf C$ with $\omega_2(s)\Phi^\square=a\Phi^\square$ and
$\omega_2(n_0)\Phi^\square=a'\Phi^\square$. As $s^2=1$ and $\Phi^\square\neq0$,
$a^2=1$.

\emph{Step 5: the Siegel rule.} As in Step 5 of the proof of
Lemma~\ref{asm-theta-agree-local-factor-vacuum}, with the rule for $P_X$ in
condition (b) of Definition~\ref{asm-theta-agree-doubling}(6) for $\omega_2$,
the character $\chi_W$ and $m=d$: if $q\in P_X$ has
$\omega_2(q)\Phi^\square=a_q\Phi^\square$, then
$a_q=\chi_W(\det t(q))|\det t(q)|^{d/2}$. For $q=n_0$, with $t(n_0)=1$, this
gives $a'=1$. So $\omega_2(sn_0s)\Phi^\square=a^2\Phi^\square=\Phi^\square$, and
\[
\omega_2(p)\Phi^\square=\omega_2(sn_0s)^{-1}\omega_2(k)\Phi^\square=c\,\Phi^\square .
\]
As $p\in P_X$ with $t(p)=t$, $c=\chi_W(\det t)|\det t|^{d/2}$.

\emph{Step 6: $c=1$.} The matrices $t=y+b(y-x)$ and $t^{-1}=t'$ preserve $L_v$,
and $t-1=(y-1)+b(y-x)$ is at $v$. As in Step 6 of the proof of
Lemma~\ref{asm-theta-agree-local-factor-vacuum}, $\chi_W(\det t)=1$, as $\chi_W$
is unramified at the places above $v$; so $c=|\det t|^{d/2}$ is a positive real
number. The operator $\omega_2(k)$ preserves the $L^2$-norm of Schwartz-Bruhat
functions (condition (b) of Definition~\ref{asm-theta-agree-doubling}(6)), so
$\|\Phi^\square\|=\|\omega_2(k)\Phi^\square\|=|c|\,\|\Phi^\square\|$, and
$\|\Phi^\square\|\neq0$ as $\Phi^\square$ is continuous and nonzero. So $|c|=1$,
$c=1$, and $\omega_2(k)$ acts at $v$ by
$\mathrm{vac}_v\boxtimes\overline{\mathrm{vac}_v}\mapsto\mathrm{vac}_v\boxtimes\overline{\mathrm{vac}_v}$.
% Lean: (1) bump (NearUnits.lean), followed above: Step 1 is near_UA (with IsSec.unit_rel and
% Euler.isOpen_int), near_units and exists_level; Step 2 is ind_mem and the bijectivity
% hiota.bij of the linear model of dictX (Lemma asm-theta-agree-coinvariants-dictionary, through
% ThK.SplitDict.dictX = KCo.dictX_of ...); Step 3 is mult, with Heisenberg2 of
% adjoint_statement, exists_coord, co_eq and symp_eq; Step 4 is dich (Lemma
% asm-theta-agree-nonvanishing-zeta-local-dichotomy) with keep, the reading (k iota hiota).1 of
% omega(g, 1) through omB, and secM sec u0 0 for u. (2) VacFix2 = vacFix2 (PresDOf.lean),
% followed above: Phi = tens vac r0 (exists_base) and bt, rd for Phi^box; Step 1 is presD_of
% (q1, q2, with pres_of_hat); Step 2 is ev and hc (Vf.dbl_eigen, Lemma
% asm-theta-agree-local-vacuum-doubled, with hsp.toIsSplitting); Step 3 is hD, hI, ax, ay,
% Vf2.exists_beta and Vf2.exists_decomp with eps v (eps_mul, eps_smul); Step 4 is pC, pCi, pe,
% pe', the eigenvalues a and a' and aa; Step 5 is rule_diag (lDiag, lDiag_ne) for a1 and hpos,
% with e2 and e3; Step 6 is chi_det and Bad.norm_eig (hsp.unit, through lD_unit and lD_haar).
% Both parts are inputs of badFactor_of (BadFin.lean).
\end{proof}

\begin{lemma}\label{asm-theta-agree-nonvanishing-zeta}
Let $(X,W,M,\chi_X,\chi_W,\omega)$ be a Weil pair of signature $(r,s;r',s')$
for a theta base $B$ of a CM field $E$
(Definition~\ref{asm-theta-agree-triple}(1)), with $X$ and $W$ of ranks $n$ and
$d$, Haar measure $\mu$ of $M$, and $\min(r,s)=0$. Let $K$ be the lattice
stabilizer of the doubled space of $X$
(Definition~\ref{asm-theta-agree-nonvanishing-doubling}(1)), and $\nu$ a Haar
measure on $\mathbf A^{n\times d}$.
\begin{enumerate}
\item There is a continuous height function $a>0$ for $K$
  (Definition~\ref{asm-theta-agree-doubling}(4)).
\item Let $\Pi\neq0$ be a complex vector space and
  $p\colon\mathcal S(F_\infty^{n\times d})\to\Pi$ a linear map with
  $p(\operatorname{Pol}(M))=\Pi$, such that $p(f_1)=0$ whenever
  $f\in\operatorname{Pol}(M)$, $\gamma$ is a line in $U_X$ and
  $t\mapsto\omega(\gamma(t),1)$ acts at $\infty$ to first order by
  $f\mapsto f_1$, and $p(f_1)=p(f)$ whenever $f\in\operatorname{Pol}(M)$,
  $k\in K_\infty(X)$ and $\omega(k,1)$ acts at $\infty$ by $f\mapsto f_1$
  (Definition~\ref{asm-theta-agree-matched}(1)). Then there is a nonzero
  $f_0\in\operatorname{Pol}(M)$ such that $\omega(k,1)$ acts at $\infty$ by
  $f_0\mapsto f_0$ for every $k\in K_\infty(X)$.
\item Assume the meromorphic continuation of Siegel Eisenstein series in the
  form of Theorem~\ref{asm-theta-agree-eisenstein} and the Siegel-Weil formula
  in the form of Theorem~\ref{asm-theta-agree-siegel-weil}. Suppose that
  $\min(r',s')=0$ and $n,d\geq1$, let $a$ be a height function for $K$, let
  $Z_\Phi$ be the zeta integrands for $\nu$ and $a$
  (Definition~\ref{asm-theta-agree-nonvanishing-doubling}(4)), and let
  $\Phi\in\operatorname{PD}(M)$. Then there are an open set $U\subset\mathbf C$
  such that every $z\in\mathbf C$ has a punctured neighbourhood contained in
  $U$, a function $Z$ holomorphic on $U$ and $t_0\in\mathbf R$ such that
  $(d-n)/2\in U$; for $\operatorname{Re}z>t_0$, $z\in U$, the function
  $g\mapsto Z_\Phi(z,g)$ is $\mu$-integrable on $U_X$ and
  $Z(z)=\int_{U_X}Z_\Phi(z,g)\,d\mu(g)$; and $Z((d-n)/2)=0$ if $I(\Phi)=0$.
\item Suppose that $\min(r',s')=0$, $n\leq d$ and $d$ is even, let $T\subset P$
  be a finite set such that $\chi_X$ and $\chi_W$ are unramified at every place
  of $E$ not above $T$, let $f_0\neq0$ be as in (2), let $a$ be a height
  function for $K$, and let $Z_\Phi$ be the zeta integrands for $\nu$ and $a$.
  Then there is $\Phi\in\operatorname{PD}(M)$ such that: for
  every $z\in\mathbf C$ and every finite place $v\notin T$ there is a function
  $f$ on $G_X$, spherical at $v$ for $(\chi_W,z+n/2)$, with
  $f(\iota(u,1))=1$ for every $u\in U_X^{(v)}$ integral at $v$ and
  $Z_\Phi(z,ug)=f(\iota(u,1))\,Z_\Phi(z,g)$ for every $u\in U_X^{(v)}$ and every
  $g$ trivial at $v$; the function
  $H(z)=\int_{\Omega_T}Z_\Phi(z,g)\,d\mu(g)$ is entire; and
  $H((d-n)/2)\neq0$.
\end{enumerate}
\end{lemma}
\leanref{Development/WeilBound/ThetaLifts/LatticeBasics/PairsStable.lean}{GlobalTheory.ThK.Dbl.HgtEx}
\leanref{Development/WeilBound/ThetaLifts/LatticeBasics/HodgeTateFromSandwich.lean}{GlobalTheory.ThK.Dbl.RealInv}
\leanref{Development/WeilBound/ThetaLifts/LatticeBasics/HodgeTateFromSandwich.lean}{GlobalTheory.ThK.Dbl.NormZeta}
\leanref{Development/WeilBound/ThetaLifts/LatticeBasics/HodgeTateFromSandwich.lean}{GlobalTheory.ThK.Dbl.BadFactor}
\leanref{Development/WeilBound/ThetaLifts/LatticeBasics/RealInv.lean}{GlobalTheory.ThK.Dbl.realInv_of}
\leanref{Development/WeilBound/ThetaLifts/LatticeBasics/RealInv.lean}{GlobalTheory.ThK.Dbl.Dense.pol_finite}
\leanref{Development/WeilBound/ThetaLifts/LatticeBasics/RealInv.lean}{GlobalTheory.ThK.Dbl.Dense.exists_fixed}
\leanref{Development/WeilBound/ThetaLifts/LatticeBasics/NormZeta.lean}{GlobalTheory.ThK.Dbl.normZeta_of}
\leanref{Development/WeilBound/ThetaLifts/LatticeBasics/BadFactor.lean}{GlobalTheory.ThK.Dbl.badFactor_of}
\leanref{Development/WeilBound/ThetaLifts/LatticeBasics/BadFactor.lean}{GlobalTheory.ThK.Dbl.Bad.orth}
\leanref{Development/WeilBound/ThetaLifts/LatticeBasics/BadFactor.lean}{GlobalTheory.ThK.Dbl.Bad.fix_of_parts}
\begin{proof}
Put $I=\{1,\ldots,n\}\times\{1,\ldots,d\}$, let $\rho$ be the Schrödinger
operators of $M$, and for $u\in M_{n\times d}(\mathbf A_E)$ at $\infty$ let
$D_u$ be the operator of Lemma~\ref{asm-theta-agree-real-derivative} for the
isomorphism of $M$.

(1) This is Lemma~\ref{asm-theta-agree-doubling-stabilizer}(3), for the
Hermitian lattice space $X$.

(2) \emph{Finite-dimensional stable spaces.} Every $f\in\operatorname{Pol}(M)$
lies in a finite-dimensional subspace $N\subset\mathcal S(F_\infty^I)$ with
$T(N)\subset N$ for every $k\in K_\infty(X)$ and every archimedean operator $T$
of $\omega$ at $(k,1)$ (Definition~\ref{asm-theta-agree-rechoice-terms}(1)).
The set of such $f$ is a subspace. For such $k$ and $T$ and $u$ at $\infty$,
$T(D_uf)=D_{ku}(T(f))$ by Lemma~\ref{asm-theta-agree-rechoice-intertwine}(2).
A Gaussian with signs $\varepsilon$ lies in the space $G_\varepsilon$ of
Lemma~\ref{asm-theta-agree-doubling-level}(2) for the isomorphism of $M$ and
the matrices $\theta_{X,w}$ and $\theta_{W,w}$, by
Lemma~\ref{asm-theta-agree-density-derivative}. A nonzero element of
$G_\varepsilon$ vanishes nowhere, so $f\mapsto f(0)$ is injective on
$G_\varepsilon$, which therefore has dimension at most $1$. It is stable under
$T$: if $u$ and $u'$ satisfy the relation defining $G_\varepsilon$, so do $ku$
and $ku'$, as $\pi_w(k)$ commutes with $\theta_{X,w}$. So the subspace contains
every Gaussian. If $f$ lies in such an $N$ and $u$ is at $\infty$, then $D_uf$
lies in the span $N'$ of the $D_vg$ with $v$ at $\infty$ and $g\in N$, which is
finite-dimensional, as $D_v$ depends linearly on the components at $\infty$ of
the image of $v$ under the isomorphism of $M$, and stable under every such $T$,
as $T(D_vg)=D_{kv}(T(g))$. By Lemma~\ref{asm-theta-agree-density-derivative}
and the definition of $\operatorname{Pol}(M)$, every $f\in\operatorname{Pol}(M)$
has such an $N$.

\emph{The fixed vector.} As $\Pi\neq0$ and $p(\operatorname{Pol}(M))=\Pi$,
there is $f\in\operatorname{Pol}(M)$ with $p(f)\neq0$; let $N$ be as above for
$f$, and $N_0=N\cap\operatorname{Pol}(M)$, finite-dimensional. Give
$\mathcal S(F_\infty^I)$ the $L^2$-norm of a Haar measure. Let $z$ be the
orthogonal projection of $f$ to the finite-dimensional space
$N_0\cap\ker p$, and $f_0=f-z\in N_0$. Then $p(f_0)=p(f)\neq0$, so $f_0\neq0$;
and $f_0$ is orthogonal to $N_0\cap\ker p$, so if $g\in N_0$ has
$p(g)=p(f_0)$ and $\|g\|=\|f_0\|$, then
$\|g\|^2=\|f_0\|^2+\|g-f_0\|^2$, and $g=f_0$. Let $k\in K_\infty(X)$. By
Lemma~\ref{asm-theta-agree-rechoice-operators}(1) there is an archimedean
operator $T$ of $\omega$ at $(k,1)$, and $T$ maps $\operatorname{Pol}(M)$ into
itself by Lemma~\ref{asm-theta-agree-rechoice-fock}, as $1\in K_\infty(W)$. So
$g=T(f_0)\in N_0$. As $\omega(k,1)$ acts at $\infty$ by $f_0\mapsto T(f_0)$,
the hypothesis on $p$ gives $p(g)=p(f_0)$. As $\omega(k,1)$ preserves the
$L^2$-norm and $\omega(k,1)(f_0\otimes\phi)=T(f_0)\otimes\phi$ for a nonzero
Schwartz-Bruhat $\phi$ on $\mathbf A_f^I$, $\|g\|=\|f_0\|$. Hence $T(f_0)=f_0$:
$\omega(k,1)$ acts at $\infty$ by $f_0\mapsto f_0$.

(3) Let $\mu_X$ and $\mathcal D_X$ be the Haar measure $\mu$ and the domain of
$M$, $s_0=(d-n)/2$, $\Phi'=\Phi\boxtimes\bar\Phi$, and let $\iota$ be the
doubling embedding of $X$.

\emph{The data.} By Lemma~\ref{asm-theta-agree-definite-compact} for $X$ and
for $W$, both definite, $U_W$ has a Haar measure $\mu_W$ with a fundamental
domain $\mathcal D_W$, the measures $\mu_X(\mathcal D_X)$ and
$\mu_W(\mathcal D_W)$ are finite and $\mu_X(\mathcal D_X)\neq0$, the constant
function $1$ on $U_X$ is rapidly decreasing, every Haar measure on $U_X$ is
invariant under right translations, and every continuous function on
$U_X\times U_W$ invariant under $U_X(F)\times U_W(F)$ is bounded. By
Lemma~\ref{asm-theta-agree-nonvanishing-schwartz}, $\Phi$ is Schwartz-Bruhat,
so $\theta_\omega(\Phi)$ is continuous
(Lemma~\ref{asm-theta-agree-convergence-theta}(2)), invariant under
$U_X(F)\times U_W(F)$ (Lemma~\ref{asm-theta-agree-anisotropic-automorphy}), and
bounded. Let $(I\sqcup I,e^\square,\omega_2,S,l)$ be the doubling datum for
$(W,X)$ of Lemma~\ref{asm-theta-agree-doubling-datum}(2), with character
$\chi_W$, and $c'>0$ given by its part (3) for $\nu$. Let $f$ be the section
$x\mapsto l(\omega_2(x)\Phi')$, of degree $s_0$ for $(K,\chi_W)$ by
Lemma~\ref{asm-theta-agree-doubling-section}(2), with flat family $f_s$, and
let $U$ and $\hat E$ be given by Lemma~\ref{asm-theta-agree-doubling-value}
for $\mu_W$ and $\mathcal D_W$, so that $s_0\in U$ and
$\hat E(s_0,x)=\kappa\mu_W(\mathcal D_W)^{-1}\int_{\mathcal D_W}\theta(\omega_2(x)S(h)\Phi')\,d\mu_W(h)$,
with $\kappa\in\{1,2\}$ as in
Lemma~\ref{asm-theta-agree-doubling-siegel-weil}(2).
Let $P(s,\phi)$ and $t_1$ be given by
Lemma~\ref{asm-theta-agree-doubling-kernel} for $Y=X$, $\chi=\chi_W$,
$z=d/2$, $f$, $s_0$, $U$, $\hat E$, $\mu_X$ and $\mathcal D_X$.

\emph{The identity.} For $s\in\mathbf C$ and $g\in U_X$,
$f_s(\iota(1,g))\chi_W(\det g)^{-1}=c'Z_\Phi(s,g^{-1})$. Indeed,
$\iota(g,g)\in P_X$ with $t=g$ and $|\det g|=1$, so
$f_s(\iota(1,g))\chi_W(\det g)^{-1}=f_s(\iota(g,g)\iota(g^{-1},1))\chi_W(\det g)^{-1}=f_s(\iota(g^{-1},1))$.
By Lemma~\ref{asm-theta-agree-doubling-datum}(2) with $g'=1$,
$\omega_2(\iota(g^{-1},1))\Phi'=\omega(g^{-1},1)\Phi\boxtimes\bar\Phi$, so
$f(\iota(g^{-1},1))=c'\int(\omega(g^{-1},1)\Phi)\bar\Phi\,d\nu=c'c_\Phi(g^{-1})$,
and $f_s(\iota(g^{-1},1))=a(\iota(g^{-1},1))^{s-s_0}c'c_\Phi(g^{-1})=c'Z_\Phi(s,g^{-1})$.

\emph{The continuation.} As $\mu_X$ is invariant under right translations,
the measure $A\mapsto\mu_X(A^{-1})$ is again a Haar measure, so there is
$c\neq0$ with $\int F(g^{-1})\,d\mu_X(g)=c\int F\,d\mu_X$ for every integrable
$F$, and $F$ is integrable if $g\mapsto F(g^{-1})$ is. Put
$Z(s)=(c'c\,\mu_X(\mathcal D_X))^{-1}P(s,1)$, holomorphic on $U$; every point
of $\mathbf C$ has a punctured neighbourhood in $U$, $U$ is open, and
$s_0\in U$. Let $\operatorname{Re}s>t_1$, so that $s\in U$. As
$\langle R_g1,1\rangle_{\mathcal D_X}=\mu_X(\mathcal D_X)$, the identity and
Lemma~\ref{asm-theta-agree-doubling-kernel} show that
$g\mapsto Z_\Phi(s,g^{-1})$, hence $Z_\Phi(s,\cdot)$, is integrable, and
\[
P(s,1)=\mu_X(\mathcal D_X)\int_{U_X}c'Z_\Phi(s,g^{-1})\,d\mu_X(g)=c'c\,\mu_X(\mathcal D_X)\int_{U_X}Z_\Phi(s,g)\,d\mu_X(g),
\]
that is, $Z(s)=\int_{U_X}Z_\Phi(s,g)\,d\mu_X(g)$.

\emph{The value.} Suppose that $I(\Phi)=0$. By the value of $\hat E$ and
Lemma~\ref{asm-theta-agree-doubling-datum}(2),
$\chi_W(\det g')^{-1}\hat E(s_0,\iota(g,g'))=\kappa\mu_W(\mathcal D_W)^{-1}\int_{\mathcal D_W}\theta_\omega(\Phi)(g,h)\overline{\theta_\omega(\Phi)(g',h)}\,d\mu_W(h)$.
As $\theta_\omega(\Phi)$ is bounded and continuous and the domains have finite
measure, Fubini's theorem gives
\[
P(s_0,1)=\frac{\kappa}{\mu_W(\mathcal D_W)}\int_{\mathcal D_W}\Bigl(\int_{\mathcal D_X}\theta_\omega(\Phi)(g,h)\,d\mu_X(g)\Bigr)\overline{\Bigl(\int_{\mathcal D_X}\theta_\omega(\Phi)(g',h)\,d\mu_X(g')\Bigr)}\,d\mu_W(h)=\frac{\kappa}{\mu_W(\mathcal D_W)}\int_{\mathcal D_W}|I(\Phi)(h)|^2\,d\mu_W(h)=0.
\]
So $Z(s_0)=0$.

(4) Let $\mu_X$ be the Haar measure $\mu$ of $M$, $s_0=(d-n)/2$, $\iota$ the
doubling embedding of $X$, and $\langle\cdot,\cdot\rangle$ the $L^2$-pairing
for $\nu$. As $T\subset P$ and $n\leq d$, Lemma~\ref{asm-theta-agree-nonvanishing-zeta-local}(1)
gives $\beta_v$ for $v\in T$. Let $\Phi=f_0\otimes\bigotimes_v\phi_v$ with
$\phi_v=\beta_v$ for $v\in T$ and $\phi_v=\mathrm{vac}_v$ otherwise. Then
$\Phi\in\operatorname{PD}(M)$; it is Schwartz-Bruhat
(Lemma~\ref{asm-theta-agree-nonvanishing-schwartz}) and nonzero, as $f_0$ and
all $\phi_v$ are. Let $(I\sqcup I,e^\square,\omega_2,S,l)$ be the doubling datum
of Lemma~\ref{asm-theta-agree-doubling-datum}(2), $c'>0$ given by its part (3)
for $\nu$, and $\Phi'=\Phi\boxtimes\bar\Phi$.

\emph{Local facts.} (i) An operator that acts at a finite place $v$ by
$\phi_v\mapsto c\phi_v$ maps $\Phi$ to $c\Phi$; one that acts at $\infty$ by
$f_0\mapsto f_0$ fixes $\Phi$; and one that acts at $v$ by
$\phi_v\boxtimes\bar\phi_v\mapsto\phi_v\boxtimes\bar\phi_v$ fixes $\Phi'$. This
is Definition~\ref{asm-theta-agree-models}(2), applied to the product of the
other components. (ii) For $v\notin T$ and $k\in K$ with $k-1$ at $v$,
$\omega_2(k)\Phi'=\Phi'$, by
Lemma~\ref{asm-theta-agree-nonvanishing-zeta-local}(2), as $\chi_W$ is
unramified above $v$, and (i). (iii) By
Lemma~\ref{asm-theta-agree-doubling-datum}(2) with $g'=1$,
$\omega_2(\iota(g,1))\Phi'=\omega(g,1)\Phi\boxtimes\bar\Phi$, so
$l(\omega_2(\iota(g,1))\Phi')=c'c_\Phi(g)$; in particular $c_\Phi$ is
continuous, by condition (b) of
Definition~\ref{asm-theta-agree-doubling}(6) and
Lemma~\ref{asm-theta-agree-doubling-stabilizer}(4). (iv)
$c_\Phi(1)=\langle\Phi,\Phi\rangle\neq0$, as $\Phi$ is continuous and nonzero.
(v) $a(xk)=a(x)$ and $a(k)=1$ for $k\in K$. (vi) As $X$ is definite, each
$\theta_{X,w}$ is $\pm1$, so every element of $U_X$ at $\infty$ lies in
$K_\infty(X)$, and if $g\in U_X$ is integral at every finite place, then
$\iota(g,1)\in K$.

\emph{The spherical functions.} Let $z\in\mathbf C$ and $v\notin T$ finite.
Let $G$ be the value at $z$ of the flat family of the section
$x\mapsto l(\omega_2(x)\Phi')$, of degree $s_0$ for $(K,\chi_W)$
(Lemma~\ref{asm-theta-agree-doubling-section}(2)); so
$G(px)=\chi_W(\det t(p))|\det t(p)|^{z+n/2}G(x)$ for $p\in P_X$. By (ii) and (v),
$G(yk)=G(y)$ for $k\in K$ with $k-1$ at $v$. By (iii),
$G(\iota(g,1))=c'Z_\Phi(z,g)$, and $G(1)=c'c_\Phi(1)\neq0$ by (iv) and (v). Put
$f=G/G(1)$. For $p\in P_X$ with $p-1$ at $v$ and $k\in K$ with $k-1$ at $v$,
$f(pk)=\chi_W(\det t(p))|\det t(p)|^{z+n/2}$, so $f$ is spherical at $v$ for
$(\chi_W,z+n/2)$. Let $u\in U_X^{(v)}$; then $\iota(u,1)-1$ is at $v$. If $u$ is
integral at $v$, it is integral at every finite place, as its other components
are $1$, so $\iota(u,1)\in K$ by (vi) and $f(\iota(u,1))=1$. Let $g$ be trivial
at $v$. Write $\iota(u,1)=pk$ with $p\in P_X$ and $k\in K$
(Lemma~\ref{asm-theta-agree-doubling-stabilizer}(2)); replacing $p$ and $k$ by
the elements with the same components at $v$ and $1$ elsewhere, which lie in
$P_X$ and $K$, we may assume that $p-1$ and $k-1$ are at $v$. As $k$ and
$\iota(g,1)$ differ from $1$ at disjoint sets of places, they commute, so
$G(\iota(ug,1))=G(p\,\iota(g,1)\,k)=\chi_W(\det t(p))|\det t(p)|^{z+n/2}G(\iota(g,1))=f(\iota(u,1))\,G(\iota(g,1))$,
that is, $Z_\Phi(z,ug)=f(\iota(u,1))Z_\Phi(z,g)$.

\emph{The key property.} Let $g\in\Omega_T$ with $c_\Phi(g)\neq0$. Then $g$ is
integral at every finite place and $c_\Phi(g)=c_\Phi(1)$. For a finite place $v$
let $g_v\in U_X^{(v)}$ have the component of $g$ at $v$ and $1$ elsewhere, and
let $g_\infty$ have the components of $g$ at $\infty$ and $1$ elsewhere. Let
$v\in T$. Suppose the second alternative of
Lemma~\ref{asm-theta-agree-nonvanishing-zeta-local}(1) held for $g_v$, with $u$,
$c$ and $c'$. As $u$ is at $v$, $g_vu=gu$, so $\rho(u)\Phi=c\Phi$ and
$\rho(gu)\Phi=c'\Phi$ by (i). As $\omega(g,1)\rho(u)=\rho(gu)\omega(g,1)$ and
$\rho(gu)$ preserves the pairing,
$c_\Phi(g)=\langle\rho(gu)\omega(g,1)\Phi,\rho(gu)\Phi\rangle=c\bar c'c_\Phi(g)$
and $\langle\Phi,\Phi\rangle=|c'|^2\langle\Phi,\Phi\rangle$; so
$c\bar c'=1=c'\bar c'$ and $c=c'$, a contradiction. So $g_v$, hence $g$, is
integral at $v$, and $\omega(g_v,1)\Phi=\Phi$ by (i). For $v\notin T$, $g$ is
integral at $v$ as $g\in\Omega_T$, so $\iota(g_v,1)\in K$ by (vi), with
$\iota(g_v,1)-1$ at $v$; by (ii) and (iii),
$\omega(g_v,1)\Phi\boxtimes\bar\Phi=\Phi\boxtimes\bar\Phi$, so
$\omega(g_v,1)\Phi=\Phi$. By (vi), $g_\infty\in K_\infty(X)$, so
$\omega(g_\infty,1)\Phi=\Phi$ by the choice of $f_0$ and (i). As $g$ is the limit
of the products of $g_\infty$ and the $g_v$, $v\in S$, along the finite sets
$S$ of finite places, and $g'\mapsto(\omega(g',1)\Phi)(x)$ is continuous,
evaluation at $x$ being admissible, $\omega(g,1)\Phi=\Phi$; so
$c_\Phi(g)=c_\Phi(1)$.

\emph{The integral.} For $g\in\Omega_T$ with $c_\Phi(g)\neq0$, $\iota(g,1)\in K$
by the key property and (vi), so $a(\iota(g,1))=1$ and
$Z_\Phi(z,g)=c_\Phi(1)$; where $c_\Phi(g)=0$, $Z_\Phi(z,g)=0$. Hence
\[
H(z)=\int_{\Omega_T}Z_\Phi(z,g)\,d\mu_X(g)=\mu_X\bigl(\Omega_T\cap\{c_\Phi\neq0\}\bigr)\,c_\Phi(1)
\]
for every $z$, so $H$ is constant, hence entire. The set $\{c_\Phi\neq0\}$ is
open and contains $1$, and $\Omega_T$ contains the set of elements integral at
every finite place, which is a neighbourhood of $1$; so the measure is
positive. By the key property, $\Omega_T\cap\{c_\Phi\neq0\}$ lies in the set of
elements integral at every finite place. This set is the lattice stabilizer of
$X$ (Definition~\ref{asm-theta-agree-planes}(6)): the condition at infinity
holds by (vi), and an element of $\mathbf A_E^n$ with components $0$ at
$\infty$ lies in $\hat L_X$ if and only if its component at every finite $v$
lies in the closure $L_v$ of $L_X$. It is compact by
Lemma~\ref{asm-theta-agree-stabilizer}, so the measure is finite. Hence
$H((d-n)/2)\neq0$.
% Lean: (1) HgtEx = hgtEx_of stab2Cpt; it also gives the embedding emb2 = iota.
% (2) RealInv = realInv_of ArchOp (RealInv.lean): Dense.pol_finite (vacSp, vacSp_fin through
%   the Gaussian structure, vacSp_map and exists_dop_fin, with inter = the intertwining of an
%   archimedean operator with the derivatives), Dense.exists_fixed (orthogonal projection in
%   L^2 through toLp), and the clause map_Pol of ArchOp; the L^2 norm is preserved by
%   Dense.lintegral_eq. The book uses Lemma asm-theta-agree-doubling-level(2) for vacSp_fin.
% (3) NormZeta = normZeta_of SWFullX BoxSecX DblX HgtEx HgtInt Stab2Cpt CptAll ThetaBasic
%   PolSch: id is the identity, inv_scale the constant c, kpair_cont the kernel lemma, and
%   An.at_value the Fubini step. The shared parts are Lemmas asm-theta-agree-doubling-*.
% (4) BadFactor = badFactor = badFactor_of dblX stabEx vacFix2 bump (BadFin.lean,
%   BadFactor.lean), followed above: Phi = pure f_0 phi; (i) is eig, fix_real and fix_box
%   (SchrUnit.lean); (ii) is fx; (iii) is bx and cf, with ct for continuity; (iv) is sq_ne;
%   (v) is hgt_mul and hgt_one; (vi) is commute_theta and emb_mem. The spherical functions are
%   the first refine (law, R, Z, G1, loc_iwa, comm_off, at_emb, int_all); the key property is
%   key (orth through schr_unit, fix_of_parts, prj_eq, Euler.int_prj, partU_real_mem); the
%   integral is val, with int_nhds, stab_iff and StabEx. Bump and VacFix2 are Lemma
%   asm-theta-agree-nonvanishing-zeta-local, proved above.
% NormZeta and BadFactor hold for any K, a and e with IsStab2, IsHgt and IsEmb.
\end{proof}

\begin{lemma}\label{asm-theta-agree-nonvanishing-euler}
Let $E$ be a CM field with a theta base $B$
(Definition~\ref{asm-theta-agree-models}), let $X$ be a Hermitian lattice space
of some signature (Definition~\ref{asm-theta-agree-doubling}(1)), and let $\mu$
be a Haar measure on $U_X$.
\begin{enumerate}
\item For every finite place $v$ there is a Haar measure $\mu_v$ on
  $U_X^{(v)}$ for which the set of elements integral at $v$
  (Definition~\ref{asm-theta-agree-nonvanishing-doubling}(2)) has measure $1$.
\item Let $T$ be a finite set of finite places, let $\mu_v$ be as in (1) for
  every finite $v$, let $f$ be a $\mu$-integrable function on $U_X$, and for
  every finite place $v\notin T$ let $f_v$ be a $\mu_v$-integrable function on
  $U_X^{(v)}$ such that $f_v(u)=1$ for every $u$ integral at $v$ and
  $f(ug)=f_v(u)f(g)$ for every $u\in U_X^{(v)}$ and every $g$ trivial at $v$.
  Then, along the finite sets $S$ of finite places outside $T$ ordered by
  inclusion,
  \[
  \Bigl(\prod_{v\in S}\int f_v\,d\mu_v\Bigr)\int_{\Omega_T}f\,d\mu\ \longrightarrow\ \int_{U_X}f\,d\mu .
  \]
\end{enumerate}
\end{lemma}
\leanref{Development/WeilBound/ThetaLifts/LatticeBasics/HodgeTateFromSandwich.lean}{GlobalTheory.ThK.Dbl.LocMeas}
\leanref{Development/WeilBound/ThetaLifts/LatticeBasics/LocMeas.lean}{GlobalTheory.ThK.Dbl.locMeas}
\leanref{Development/WeilBound/ThetaLifts/LatticeBasics/HodgeTateFromSandwich.lean}{GlobalTheory.ThK.Dbl.EulerProd}
\leanref{Development/WeilBound/ThetaLifts/LatticeBasics/EulerProd.lean}{GlobalTheory.ThK.Dbl.eulerProd}
\leanref{Development/WeilBound/ThetaLifts/LatticeBasics/EulerProd.lean}{GlobalTheory.ThK.Dbl.Euler.step}
\leanref{Development/WeilBound/ThetaLifts/LatticeBasics/EulerA.lean}{GlobalTheory.ThK.Dbl.Euler.factor}
\begin{proof}
(1) The group $U_X^{(v)}$ is locally compact, and its elements integral at $v$
form a compact open subgroup $K_v$. So there is a Haar measure with
$\mu_v(K_v)=1$.

(2) For $u\in U_X^{(v)}$ and $g\in U_X$, $ug$ is integral at a place $w\neq v$
if and only if $g$ is. For a finite set $S$ of finite places outside $T$, let
$A_S$ be the set of $g\in U_X$ that are integral at every finite place outside
$T\cup S$, a closed set with $A_\emptyset=\Omega_T$.

\emph{One place.} Let $v\notin T\cup S$, put $A=A_{S\cup\{v\}}$, and let
$\pi_v\colon U_X\to U_X^{(v)}$ be the continuous map that keeps the component
at $v$ and puts $1$ elsewhere, so that $\pi_v(ug)=u\,\pi_v(g)$ for
$u\in U_X^{(v)}$, and $g$ is integral at $v$ if and only if $\pi_v(g)\in K_v$.
For every $g$, $\mu_v\{u:u^{-1}\pi_v(g)\in K_v\}=\mu_v(K_v)=1$, by
invariance. By Fubini's theorem and the invariance of $\mu$ under left
translation by $u$,
\[
\int_Af\,d\mu
=\int_{U_X^{(v)}}\int_{U_X}\mathbf 1_A(g)\,\mathbf 1_{K_v}(u^{-1}\pi_v(g))\,f(g)\,d\mu(g)\,d\mu_v(u)
=\int_{U_X}\mathbf 1_A(g)\,\mathbf 1_{K_v}(\pi_v(g))\int_{U_X^{(v)}}f(ug)\,d\mu_v(u)\,d\mu(g),
\]
since $ug\in A$ if and only if $g\in A$. Let $g\in A$ with $\pi_v(g)\in K_v$,
and put $g'=\pi_v(g)^{-1}g$, which is trivial at $v$. Then
$f(ug)=f_v(u\pi_v(g))f(g')$, and in particular $f(g)=f_v(\pi_v(g))f(g')=f(g')$.
As $\mu_v$ is invariant under right translation by elements of the compact
subgroup $K_v$, $\int f(ug)\,d\mu_v(u)=f(g)\int f_v\,d\mu_v$. So
$\int_Af\,d\mu=\bigl(\int f_v\,d\mu_v\bigr)\int_{A_S}f\,d\mu$, as
$A\cap\pi_v^{-1}(K_v)=A_S$.

\emph{All places.} By induction on $S$,
$\int_{A_S}f\,d\mu=\bigl(\prod_{v\in S}\int f_v\,d\mu_v\bigr)\int_{\Omega_T}f\,d\mu$.
The sets $A_S$ increase with $S$, and their union is $U_X$, since $g$ and
$g^{-1}$ have entries integral at almost all places, so $g$ is integral at
almost all $v$. As $f$ is integrable, $\int_{A_S}f\,d\mu\to\int_{U_X}f\,d\mu$ by
dominated convergence.
\end{proof}
% Lean: (1) locMeas, the Haar measure normalized on the compact open subgroup intGrp
% (isCompact_int, isOpen_int). (2) eulerProd: Euler.step (from Euler.factor with K the
% integral elements, i = the inclusion, pi = prj v; int_mul_ne and int_prj), induction on
% S, ae_mulVec for the union, and tendsto_setIntegral_of_monotone.

\begin{lemma}\label{asm-theta-agree-nonvanishing-normalizer}
Let $E$ be a CM field with maximal totally real subfield $F$, let $n\geq0$, and
let $T$ be a finite set of finite places of $F$. Then there is a function
$d_T$, holomorphic on $\operatorname{Re}z>0$, such that for
$\operatorname{Re}z>0$, $d_T(z)\neq0$ and the product of the $d_v(z)$
(Definition~\ref{asm-theta-agree-nonvanishing-doubling}(3)) over the finite
places $v\notin T$ converges unconditionally to $d_T(z)$.
\end{lemma}
\leanref{Development/WeilBound/ThetaLifts/LatticeBasics/HodgeTateFromSandwich.lean}{GlobalTheory.ThK.Dbl.DProd}
\leanref{Development/WeilBound/ThetaLifts/LatticeBasics/EulerHalf.lean}{GlobalTheory.ThK.Dbl.dProd}
\leanref{Development/WeilBound/ThetaLifts/LatticeBasics/EulerHalf.lean}{GlobalTheory.ThK.Dbl.euler_half}
\begin{proof}
Fix $1\leq i\leq n$, and for $v\notin T$ put
$x_v(z)=-\epsilon_v^{\,n-i}N(v)^{-(2z+i)}$. Then
$|x_v(z)|\leq N(v)^{-t}$ for $t\leq2\operatorname{Re}z+1$, and
$|x_v(z)|\leq\frac12$ for $\operatorname{Re}z>0$, as $N(v)\geq2$. As
$|\log(1+x)|\leq\frac32|x|$ for $|x|\leq\frac12$ and $\sum_vN(v)^{-t}$
converges for $t>1$, the series $\sum_{v\notin T}\log(1+x_v(z))$ converges
absolutely and locally uniformly on $\operatorname{Re}z>0$; its sum $\ell_i(z)$
is holomorphic there. So
$D_i(z)=\exp(-\ell_i(z))$ is holomorphic and nonzero on $\operatorname{Re}z>0$,
and $\prod_{v\notin T}(1+x_v(z))^{-1}$ converges unconditionally to $D_i(z)$.
Put $d_T=D_1\cdots D_n$.
\end{proof}

\begin{theorem}[Unramified doubling zeta integral \cite{GPSR87,Li92}]\label{asm-theta-agree-nonvanishing-unramified}
Let $E$ be a CM field with a theta base $B$
(Definition~\ref{asm-theta-agree-models}), let $X$ be a Hermitian lattice space
of rank $n$ and some signature (Definition~\ref{asm-theta-agree-doubling}(1)),
and let $\chi$ be a Hecke character of $E$ that is trivial on
$\mathbf A^\times$. Let $v$ be a finite place such that $\chi$ is unramified at
every place of $E$ above $v$, let $\mu_v$ be a Haar measure on $U_X^{(v)}$ for
which the set of elements integral at $v$ has measure $1$, let $z\in\mathbf C$
with $\operatorname{Re}z>n/2$, and let $f$ be a function on $G_X$ that is
spherical at $v$ for $(\chi,z+n/2)$
(Definition~\ref{asm-theta-agree-nonvanishing-doubling}). Then
$u\mapsto f(\iota(u,1))$ is $\mu_v$-integrable on $U_X^{(v)}$, and
\[
d_v(z)\int f(\iota(u,1))\,d\mu_v(u)=\prod_{i=0}^{n-1}L_v(z+n/2-i,\chi).
\]
\end{theorem}
\leanref{Development/WeilBound/ThetaLifts/LatticeBasics/HodgeTateFromSandwich.lean}{GlobalTheory.ThK.Dbl.LocUnr}
\leanref{Development/WeilBound/ThetaLifts/LatticeBasics/LocUnr.lean}{GlobalTheory.ThK.Dbl.locUnr_of}

\begin{lemma}\label{asm-theta-agree-nonvanishing-pair}
Assume the meromorphic continuation of Siegel Eisenstein series in the form of
Theorem~\ref{asm-theta-agree-eisenstein} and the Siegel-Weil formula in the
form of Theorem~\ref{asm-theta-agree-siegel-weil}. Let
$(X,W,M,\chi_X,\chi_W,\omega)$ be a Weil pair of signature $(r,s;r',s')$ for a
theta base $B$ of a CM field $E$ (Definition~\ref{asm-theta-agree-triple}(1)),
with $X$ and $W$ of ranks $n$ and $d$, such that $\min(r,s)=\min(r',s')=0$,
$1\leq n<d$ and $d$ is even. Let $T\subset P$ be a finite set such that
$\chi_X$ and $\chi_W$ are unramified at every place of $E$ not above $T$, and
$\chi_W$ is ramified at every place of $E$ above $T$. Let $\Pi\neq0$ and $p$
be as in Lemma~\ref{asm-theta-agree-nonvanishing-zeta}(2), and let $\Lambda$ be
an entire function with $\Lambda(z)=L(z,\chi_W)$ for $\operatorname{Re}z>1$
(Definition~\ref{asm-theta-agree-nonvanishing-doubling}(3)) and
$\Lambda(z)\neq0$ for every integer $z$ with $2z\geq-d$. Then $I(\Phi)\neq0$
for some $\Phi\in\operatorname{PD}(M)$.
\end{lemma}
\leanref{Development/WeilBound/ThetaLifts/LatticeBasics/BotAsm.lean}{GlobalTheory.ThK.Dbl.pair_ne_zero}
\leanref{Development/WeilBound/ThetaLifts/LatticeBasics/BotAsm.lean}{GlobalTheory.ThK.Dbl.eq_zero_of_right}
\leanref{Development/WeilBound/ThetaLifts/LatticeBasics/BotAsm.lean}{GlobalTheory.ThK.Dbl.hasProd_off}
\begin{proof}
Let $\mu$ be the Haar measure of $M$, $K$ the lattice stabilizer of the doubled
space of $X$, $\nu$ a Haar measure on $\mathbf A^{n\times d}$, and
$s_0=(d-n)/2$; then $\operatorname{Re}s_0>0$, as $n<d$.

\emph{The data.} By Lemma~\ref{asm-theta-agree-nonvanishing-zeta}(1) there is
a height function $a$ for $K$; let $Z_\Phi$ be the zeta integrands for $\nu$
and $a$. By Lemma~\ref{asm-theta-agree-nonvanishing-zeta}(2) there is a nonzero
$f_0\in\operatorname{Pol}(M)$ fixed at $\infty$ by the $\omega(k,1)$,
$k\in K_\infty(X)$. Lemma~\ref{asm-theta-agree-nonvanishing-zeta}(4), for $T$
and $f_0$, gives $\Phi\in\operatorname{PD}(M)$, functions $f_{z,v}$ for
$z\in\mathbf C$ and $v\notin T$, and the entire function $H$ with
$H(s_0)\neq0$. Lemma~\ref{asm-theta-agree-nonvanishing-zeta}(3), for $\Phi$,
gives $U$, $Z$ and $t_0$. Let $d_T$ be as in
Lemma~\ref{asm-theta-agree-nonvanishing-normalizer}, and for every finite place
$v$ let $\mu_v$ be as in Lemma~\ref{asm-theta-agree-nonvanishing-euler}(1).

\emph{The identity.} Let $\operatorname{Re}z>\max(t_0,n/2)$. Then $z\in U$,
$Z_\Phi(z,\cdot)$ is $\mu$-integrable and $Z(z)=\int_{U_X}Z_\Phi(z,g)\,d\mu(g)$.
As $d$ is even, $\chi_W$ is trivial on $\mathbf A^\times$. For $v\notin T$,
$\chi_W$ is unramified at the places above $v$, so
Theorem~\ref{asm-theta-agree-nonvanishing-unramified}, applied to $f_{z,v}$,
shows that $f_v(u)=f_{z,v}(\iota(u,1))$ is $\mu_v$-integrable with
$d_v(z)\int f_v\,d\mu_v=\prod_{i=0}^{n-1}L_v(z+n/2-i,\chi_W)$. By
Lemma~\ref{asm-theta-agree-nonvanishing-euler}(2), applied to
$Z_\Phi(z,\cdot)$ and the $f_v$, and by
Lemma~\ref{asm-theta-agree-nonvanishing-normalizer},
\[
\prod_{v\in S}\Bigl(\prod_{i=0}^{n-1}L_v(z+n/2-i,\chi_W)\Bigr)H(z)
=\Bigl(\prod_{v\in S}d_v(z)\Bigr)\Bigl(\prod_{v\in S}\int f_v\,d\mu_v\Bigr)H(z)
\longrightarrow d_T(z)Z(z)
\]
along the finite sets $S$ of finite places outside $T$. For $0\leq i<n$,
$\operatorname{Re}(z+n/2-i)>n-i\geq1$, so the product of the
$L_v(z+n/2-i,\chi_W)$ over all finite $v$ converges unconditionally to
$\Lambda(z+n/2-i)$; and $L_v(z+n/2-i,\chi_W)=1$ for $v\in T$, as $\chi_W$ is
ramified at every place above $v$. So the left side tends to
$\prod_{i=0}^{n-1}\Lambda(z+n/2-i)\,H(z)$, and
\[
Z(z)\,d_T(z)=\prod_{i=0}^{n-1}\Lambda(z+n/2-i)\,H(z)\qquad(\operatorname{Re}z>\max(t_0,n/2)).
\]

\emph{Continuation.} Let $D(z)$ be the left side minus the right side; it is
holomorphic on $V=U\cap\{\operatorname{Re}z>0\}$. Every point of $\mathbf C$
has a punctured neighbourhood in $U$, so $\mathbf C\setminus U$ is discrete,
hence countable. The map $\psi(w)=e^{\operatorname{Re}w}+i\operatorname{Im}w$
is a homeomorphism of $\mathbf C$ onto the half-plane
$\operatorname{Re}z>0$, and $V$ is the image under $\psi$ of the complement of
the countable set $\psi^{-1}(\mathbf C\setminus U)$, which is connected. So $V$
is connected. As $D$ vanishes on a nonempty open subset of $V$, $D=0$ on $V$
by the identity theorem. As $s_0\in U$ and $\operatorname{Re}s_0>0$,
$D(s_0)=0$.

\emph{Conclusion.} Suppose that $I(\Phi)=0$. Then $Z(s_0)=0$, so
$\prod_{i=0}^{n-1}\Lambda(s_0+n/2-i)\,H(s_0)=0$. For $0\leq i<n$,
$s_0+n/2-i=d/2-i$ is an integer, as $d$ is even, and $2(d/2-i)=d-2i>-d$.
So every $\Lambda(s_0+n/2-i)$ is nonzero, and $H(s_0)\neq0$, a contradiction.
Hence $I(\Phi)\neq0$.
\end{proof}
% Lean: pair_ne_zero HgtEx NormZeta EulerProd LocMeas LocUnr BadFactor RealInv DProd; the
% identity is key (tendsto_nhds_unique of EulerProd times DProd against hasProd_prod with
% hasProd_off), the continuation is eq_zero_of_right, and the conclusion uses l3 at the
% integers k - i, where d = 2k.

\begin{lemma}\label{asm-theta-agree-nonvanishing-hecke}
Let $(B,G,\mathcal X,\omega,\omega')$ be a triple of ranks $(a,a')$ for a theta
base $B$ of a CM field $E$ (Definition~\ref{asm-theta-agree-triple}), with
$\nu=\alpha/\beta$, and let $\chi=\nu$ or $\chi=\nu^{-1}$. Then there is an
entire function $\Lambda$ with $\Lambda(z)=L(z,\chi)$ for
$\operatorname{Re}z>1$ (Definition~\ref{asm-theta-agree-nonvanishing-doubling}(3))
and $\Lambda(z)\neq0$ for every integer $z$ with $2z\geq-(a+a')$.
\end{lemma}
\leanref{Development/WeilBound/ThetaLifts/LatticeBasics/ExistsIsHeckeLOfCompleted.lean}{GlobalTheory.ThK.Dbl.LInt2}
\leanref{Development/WeilBound/ThetaLifts/Reductions/WeightConstants/ConstOfGood.lean}{GlobalTheory.ThK.Dbl.lInt2}
\leanref{Development/WeilBound/ThetaLifts/Reductions/WeightConstants/Bridge.lean}{PositivityTheory.ThKRed.HK.heckeL_ne_zero}
\begin{proof}
As $\mathcal X$ is proper, $\alpha$ and $\beta$ have the infinity types
$m_\alpha$ and $m_\beta$, so $\chi$ is a unitary Hecke character of infinity
type $\pm m$ with $m(w)=m_\alpha(w)-m_\beta(w)$, and
$|m(w)|>a+a'\geq0$ for every infinite place $w$ of $F$, as $G$ is proper. In
particular $\chi$ is not a power of the norm. By Hecke's theorem
\cite{Hec20,Tat67}, its completed $L$-function, the product of $L(s,\chi)$, a
power of the norm of the conductor and the gamma factors
$\Gamma_{\mathbf C}(s+|m(w)|/2)$, is entire and satisfies a functional equation
relating its values at $s$ and at $1-s$ to those of the completed $L$-function
of $\chi^{-1}$, with a nonzero constant; and $L(s,\chi)$ and $L(s,\chi^{-1})$
have no zeros on $\operatorname{Re}s\geq1$. Let $\Lambda$ be the completed
$L$-function divided by the conductor term and the gamma factors; as
$1/\Gamma$ is entire, $\Lambda$ is entire, and it equals $L(z,\chi)$ for
$\operatorname{Re}z>1$. Let $z$ be an integer with $2z\geq-(a+a')$. If
$z\geq1$, $\Lambda(z)=L(z,\chi)\neq0$. If $z\leq0$, the functional equation
expresses $\Lambda(z)$ as a nonzero constant times the inverse gamma factors
at $z$ times the completed $L$-function of $\chi^{-1}$ at $1-z\geq1$, which is
nonzero; and $1/\Gamma(z+|m(w)|/2)\neq0$, since
$2z+|m(w)|>2z+a+a'\geq0$.
\end{proof}
% Lean: lInt2, through exists_hchar_unitary (the unitary Hecke character H of nu, with
% infinity type m_alpha - m_beta by hasExp_nu) and ThKRed.HK.heckeL_ne_zero, whose cases
% z >= 1 and z <= 0 are continuation_ne_zero_of_one_le_re for mu and mu^{-1}; the bound
% |m(w)| > a + a' is h.geo.kappa.

\begin{lemma}\label{asm-theta-agree-nonvanishing-bottom}
Assume the meromorphic continuation of Siegel Eisenstein series in the form of
Theorem~\ref{asm-theta-agree-eisenstein} and the Siegel-Weil formula in the
form of Theorem~\ref{asm-theta-agree-siegel-weil}. Let
$(B,G,\mathcal X,\omega,\omega')$ be a triple for a theta base $B$ of a CM field
$E$ (Definition~\ref{asm-theta-agree-triple}), with $p=q=0$. Then
$I(\Phi)\neq0$ for some $\Phi\in\operatorname{PD}(M)$, and $I'(\Phi')\neq0$ for
some $\Phi'\in\operatorname{PD}(M')$.
\end{lemma}
\leanref{Development/WeilBound/ThetaLifts/LatticeBasics/PairsStable.lean}{GlobalTheory.ThK.Dbl.NonZeroBot}
\leanref{Development/WeilBound/ThetaLifts/LatticeBasics/BotAsm.lean}{GlobalTheory.ThK.Dbl.nonZeroBottom_of_doubling_integral}
\begin{proof}
Let $G$ have ranks $(a,a')$, $\mathcal X=(\alpha,\beta,S_0,Q)$,
$\nu=\alpha/\beta$ and $T=S_0\cup Q$. As $G$ is proper, $a'<a$, $a+a'$ is even
and $1=2(p+q)+1<a'$, so $1\leq a<a+a'$ and $1\leq a'<a+a'$. As $\mathcal X$ is
proper, $T\subset P$, $\nu$ and $\nu^{-1}$ are ramified at every place of $E$
above $T$, and $\alpha$, $\beta$, $\nu$ and $\nu^{-1}$ are unramified at every
place of $E$ not above $T$. By Lemma~\ref{asm-theta-agree-real} there are real
maps $(\Pi,p_\infty,p'_\infty)$ matched to the triple; by conditions (a) and
(c) of Definition~\ref{asm-theta-agree-matched}(2), $\Pi\neq0$, and
$p_\infty$, for $M$ and $\omega$, and $p'_\infty$, for $M'$ and $\omega'$,
satisfy the hypotheses of Lemma~\ref{asm-theta-agree-nonvanishing-zeta}(2). By
Lemma~\ref{asm-theta-agree-nonvanishing-hecke} there are entire functions
$\Lambda$ and $\Lambda'$ for $\nu$ and $\nu^{-1}$.

The Weil pair $(V,W,M,\beta,\nu^{-1},\omega)$, as $\beta/\alpha=\nu^{-1}$, has
signature $(a,0;a+a',0)$. Lemma~\ref{asm-theta-agree-nonvanishing-pair},
applied to it with $T$, $p_\infty$ and $\Lambda'$, gives
$\Phi\in\operatorname{PD}(M)$ with $I(\Phi)\neq0$. The Weil pair
$(V',W,M',\alpha,\nu,\omega')$ has signature $(0,a';a+a',0)$.
Lemma~\ref{asm-theta-agree-nonvanishing-pair}, applied to it with $T$,
$p'_\infty$ and $\Lambda$, gives $\Phi'\in\operatorname{PD}(M')$ with
$I'(\Phi')\neq0$.
\end{proof}
% Lean: nonZeroBottom_of_doubling_integral KReal LInt2 HgtEx NormZeta EulerProd LocMeas
% LocUnr BadFactor RealInv DProd, applying pair_ne_zero to pairV with chi_W = nu^{-1} (hr.quot,
% Lambda' from LInt2 at nu^{-1}) and to pairV' with chi_W = nu (hr.quot', Lambda at nu).

\begin{lemma}\label{asm-theta-agree-nonvanishing}
Assume the meromorphic continuation of Siegel Eisenstein series in the form of
Theorem~\ref{asm-theta-agree-eisenstein} and the Siegel-Weil formula in the
form of Theorem~\ref{asm-theta-agree-siegel-weil}. Let
$(B,G,\mathcal X,\omega,\omega')$ be a triple for a theta base $B$ of a CM field
$E$ (Definition~\ref{asm-theta-agree-triple}). If $p=0$, then $I(\Phi)\neq0$ for
some $\Phi\in\operatorname{PD}(M)$. If $q=0$, then $I'(\Phi')\neq0$ for some
$\Phi'\in\operatorname{PD}(M')\cap\mathcal S^\circ(M')$.
\end{lemma}
\leanref{Development/AbelianVarietyTheory/Definitions/Divisors/CTSmallCone.lean}{GlobalTheory.ThK.NonZeroAn}
\leanref{Development/WeilBound/ThetaLifts/LatticeBasics/NonZero.lean}{GlobalTheory.ThK.Dbl.nonZeroAn_of}
\leanref{Development/WeilBound/ThetaLifts/LatticeBasics/NonZero.lean}{GlobalTheory.ThK.Dbl.big_of}
\leanref{Development/WeilBound/ThetaLifts/LatticeBasics/NonZero.lean}{GlobalTheory.ThK.Dbl.small_of}
\begin{proof}
Fix $B$ and $\mathcal X=(\alpha,\beta,S_0,Q)$. The first Weil pair
$(V,W,M,\beta,\beta/\alpha,\omega)$ of a triple of ranks $(a,a')$ has signature
$(a-p,p;a+a'-p-q,p+q)$, and the second, $(V',W,M',\alpha,\alpha/\beta,\omega')$,
has signature $(q,a'-q;a+a'-p-q,p+q)$
(Definition~\ref{asm-theta-agree-triple}(4)). For a re-chosen datum
$G^\circ=G[V^\circ,V'^\circ,W^\circ]$, write $\tilde V$, $\tilde V'$ and
$\tilde W$ for its spaces, and $\operatorname{PD}^\circ(M)$ and
$\operatorname{PD}^\circ(M')$ for the polynomial data of $M$ and $M'$ as models
for $(\tilde V,\tilde W)$ and $(\tilde V',\tilde W)$. The theta integrals of
$G^\circ$ are $I$ and $I'$, and $U_{\tilde W}=U_W$.

\emph{Part 1: $p=0$.} Fix $a$. We show by induction on $q\geq0$ that for every
triple $(B,G,\mathcal X,\omega,\omega')$ of ranks $(a,a')$, for any $a'$, with
$p=0$ and the given $q$, some $\Phi\in\operatorname{PD}(M)$ has $I(\Phi)\neq0$.
For $q=0$ this is Lemma~\ref{asm-theta-agree-nonvanishing-bottom}. Let
$q\geq1$. As $G$ is proper, $2q+1<a'$, so $a'=b+2$ with $b\geq2$.
Lemma~\ref{asm-theta-agree-tower-beta} gives lattice spaces $V^\circ$,
$V'^\circ$ and $W^\circ$ such that $(B,G^\circ,\mathcal X,\omega,\omega')$ is a
triple, a triple $(B,G_1,\mathcal X,\omega_1,\omega'_1)$ of ranks $(a,b)$, with
first model $M_1$, third space $W_1$ and theta integral $I_1$, and adapted
planes $H$ and $H'$ such that $(G^\circ,G_1,H,H')$ is a $\beta$-configuration.
Then $p_1=p=0$ and $q_1=q-1$, so by induction some
$\Phi_1\in\operatorname{PD}(M_1)$ has $I_1(\Phi_1)\neq0$.

The first Weil pair of $G^\circ$ is
$(\tilde V,\tilde W,M,\beta,\beta/\alpha,\omega)$; $H'$ is an adapted plane in
$\tilde W$ over $W_1$; and $M_1$ is a regular model for $(\tilde V,W_1)$, since
the first space of $G_1$ is $\tilde V$. By
Lemma~\ref{asm-theta-agree-evaluation} there is an evaluation map $e$ for
$(M,M_1,1,\iota_{H'},\mathcal J_0(H'))$, and
$e(\operatorname{PD}^\circ(M))=\operatorname{PD}(M_1)$; choose
$\Phi\in\operatorname{PD}^\circ(M)$ with $e(\Phi)=\Phi_1$. Let $C\neq0$ be as in
Lemma~\ref{asm-theta-agree-ct-beta-zero} for $(G^\circ,G_1,H,H')$ and $e$.

Suppose that $I(\Psi)=0$ for every $\Psi\in\operatorname{PD}(M)$. As $p=0$, the
first Weil pair of $G$ has signature $(a,0;\ldots)$, so $I$ vanishes on every
Schwartz-Bruhat function by Lemma~\ref{asm-theta-agree-nonvanishing-schwartz}.
By the same lemma for the first Weil pair of $G^\circ$, $\Phi$ is a
Schwartz-Bruhat function, so $I(\Phi)=0$. Let $h_1\in U_{W_1}$. By
Lemma~\ref{asm-theta-agree-levi}, for the skew-Hermitian space $\tilde W$ and
$H'$, there is a Levi element $m$ for $(1,h_1)$. The pair $(0,0)$ splits the
function $I(\Phi)=0$ along $H'$ (Definition~\ref{asm-theta-agree-planes}(3)),
so Lemma~\ref{asm-theta-agree-ct-beta-zero}, for the triples of $G^\circ$ and
$G_1$, with $f_\beta=0$ and $h=1$, gives
$0=C\,I_1(e(\omega(1,1)\Phi))(h_1)=C\,I_1(\Phi_1)(h_1)$. As $C\neq0$,
$I_1(\Phi_1)=0$, a contradiction.

\emph{Part 2: $q=0$.} Fix $a'$. We show by induction on $p\geq0$ that for every
triple $(B,G,\mathcal X,\omega,\omega')$ of ranks $(a,a')$, for any $a$, with
$q=0$ and the given $p$, some $\Phi'\in\operatorname{PD}(M')$ has
$I'(\Phi')\neq0$. For $p=0$ this is
Lemma~\ref{asm-theta-agree-nonvanishing-bottom}. Let $p\geq1$. As $G$ is
proper, $a'<a$ and $2p+1<a'$, so $a=b+2$ with $b\geq2$.
Lemma~\ref{asm-theta-agree-tower-alpha} gives lattice spaces $V^\circ$,
$V'^\circ$ and $W^\circ$ such that $(B,G^\circ,\mathcal X,\omega,\omega')$ is a
triple, a triple $(B,G_1,\mathcal X,\omega_1,\omega'_1)$ of ranks $(b,a')$, with
second model $M'_1$, third space $W_1$ and theta integral $I'_1$ of its second
space, and adapted planes $H$ and $H'$ such that $(G^\circ,G_1,H,H')$ is an
$\alpha$-configuration. Then $p_1=p-1$ and $q_1=0$, so by induction some
$\Phi_1\in\operatorname{PD}(M'_1)$ has $I'_1(\Phi_1)\neq0$.

The second Weil pair of $G^\circ$ is
$(\tilde V',\tilde W,M',\alpha,\alpha/\beta,\omega')$; $H'$ is an adapted plane
in $\tilde W$ over $W_1$; and $M'_1$ is a regular model for $(\tilde V',W_1)$,
since the second space of $G_1$ is $\tilde V'$. By
Lemma~\ref{asm-theta-agree-evaluation} there is an evaluation map $e'$ for
$(M',M'_1,1,\iota_{H'},\mathcal J_0(H'))$, and
$e'(\operatorname{PD}^\circ(M'))=\operatorname{PD}(M'_1)$; choose
$\Phi\in\operatorname{PD}^\circ(M')$ with $e'(\Phi)=\Phi_1$. Let $C\neq0$ be as
in Lemma~\ref{asm-theta-agree-ct-alpha-zero} for $(G^\circ,G_1,H,H')$ and $e'$.

Suppose that $I'(\Psi)=0$ for every $\Psi\in\operatorname{PD}(M')$. As $q=0$,
the second Weil pair of $G$ has signature $(0,a';\ldots)$, so by
Lemma~\ref{asm-theta-agree-nonvanishing-schwartz}, applied to it and to the
second Weil pair of $G^\circ$, $I'(\Phi)=0$.
Lemma~\ref{asm-theta-agree-ct-alpha-zero}, for the triples of $G^\circ$ and
$G_1$, gives $f'_\alpha$ such that $(f'_\alpha,0)$ splits $I'(\Phi)=0$ along
$H'$ and $f'_\alpha(mh)=C\,I'_1(e'(\omega'(1,h)\Phi))(h_1)$ for every
$h_1\in U_{W_1}$, every Levi element $m$ for $(1,h_1)$ and every $h\in U_W$. By
Lemma~\ref{asm-theta-agree-split}, for the triple of $G^\circ$ and
$0\in\operatorname{PD}^\circ(M)$, exactly one pair splits $I(0)=0$ along $H'$.
As $(0,0)$ splits $0$, $f'_\alpha=0$. Let $h_1\in U_{W_1}$, and let $m$ be a
Levi element for $(1,h_1)$, by Lemma~\ref{asm-theta-agree-levi}. With $h=1$,
$0=C\,I'_1(\Phi_1)(h_1)$. As $C\neq0$, $I'_1(\Phi_1)=0$, a contradiction.

\emph{Conclusion.} If $p=0$, Part 1 gives the first assertion. If $q=0$, Part 2
gives $\Phi''\in\operatorname{PD}(M')$ with $I'(\Phi'')\neq0$, and
Lemma~\ref{asm-theta-agree-nonvanishing-full-rank} gives
$\Phi'\in\operatorname{PD}(M')\cap\mathcal S^\circ(M')$ with
$I'(\Phi')=I'(\Phi'')\neq0$.
\end{proof}
% Lean: GlobalTheory.ThK.Dbl.nonZeroAn_of NonZeroBot TowerA TowerB EvZero CTBigZero
% CTSmallZero SchToPol PolSch LeviEx SplitEx ToSo: Part 1 is big_of (induction on q), Part 2
% is small_of (induction on p; the uniqueness is SplitEx at Phi = 0 with i0), and ToSo
% concludes. NonZeroBot is Lemma asm-theta-agree-nonvanishing-bottom, whose inputs, Lemma
% asm-theta-agree-nonvanishing-zeta among them (with -zeta-local), are proved.

\subsection{The definite case}\label{subsection-asm-definite-case}

\begin{theorem}[Finite volume and unimodularity \cite{BHC62,Bor69}]\label{asm-theta-agree-definite-volume}
Let $E$ be a CM field with maximal totally real subfield $F$, and let
$J\in M_n(E)$ be nondegenerate $\epsilon$-Hermitian
(Definition~\ref{asm-theta-agree-convergence-tori}(1)). Then $U_J(F)$ is
countable; $U_J$ has a Haar measure and a fundamental domain for $U_J(F)$ with
respect to it; every Haar measure on $U_J$ is invariant under right
translations; and every fundamental domain for $U_J(F)$ with respect to a Haar
measure on $U_J$ has finite measure.
\end{theorem}
\leanref{Development/WeilBound/ThetaLifts/FrameEstimates/SquareIntegrableWeilRange.lean}{PositivityTheory.ThKEst.FinVolU}
\leanref{Development/WeilBound/ThetaLifts/FrameEstimates/FundamentalDomainFiniteVolume.lean}{PositivityTheory.ThKEst.finVolU}
\leanref{Development/WeilBound/ThetaLifts/CartanFactorisation/CartanConj.lean}{PositivityTheory.ThKF.Unimod}
\leanref{Development/WeilBound/ThetaLifts/CartanFactorisation/Unimod.lean}{PositivityTheory.ThKF.unimodularity_holds}

\begin{definition}\label{asm-theta-agree-definite-products}
Let $B$ be a theta base for a CM field $E$
(Definition~\ref{asm-theta-agree-models}(1)), let $Y$ be a lattice space
(Definition~\ref{asm-theta-agree-doubling}(1)), let $\mu$ be a Haar measure on
$U_Y$ and $\mathcal D$ a fundamental domain for $U_Y(F)$ with respect to $\mu$.
\begin{enumerate}
\item For functions $f$ and $f'$ on $U_Y$ put
  $\langle f,f'\rangle_{\mathcal D}=\int_{\mathcal D}f(h)\overline{f'(h)}\,d\mu(h)$
  when the integral converges absolutely, and for $g\in U_Y$ let $R_gf$ be the
  function $h\mapsto f(hg)$. For $h\in U_Y$ let
  $\overline{\operatorname{ht}}(h)$ be the infimum of
  $\operatorname{ht}(\gamma h)$ over $\gamma\in U_Y(F)$, with the height of
  Definition~\ref{asm-theta-agree-doubling}(2). A function $\phi$ on $U_Y$ is
  \emph{rapidly decreasing} if it is continuous, $\phi(\gamma h)=\phi(h)$ for
  all $\gamma\in U_Y(F)$ and $h\in U_Y$, and for every $N\geq0$ there is $C$
  with $|\phi(h)|\,\overline{\operatorname{ht}}(h)^N\leq C$ for all $h$.
\item Let $(X,W,M,\chi_X,\chi_W,\omega)$ be a Weil pair for $B$
  (Definition~\ref{asm-theta-agree-triple}(1)), with $X$ and $W$ of ranks $n$
  and $d$, and let $Y=W$. For a function $\phi$ on $U_W$ and a function $\Psi$
  on $\mathbf A^{n\times d}$, the \emph{theta lift} of $\phi$ against $\Psi$ is
  the function
  $\theta^*_\omega(\phi,\Psi)(g)=\langle\phi,\theta_\omega(\Psi)(g,\cdot)\rangle_{\mathcal D}$
  of $g\in U_X$.
\item Let $(B,G,\mathcal X,\omega,\omega',R,R_f)$ be a matched quadruple
  (Definition~\ref{asm-theta-agree-matched}(4)), with spaces $V$, $V'$ and $W$
  and models $M$ and $M'$, and let $Y=W$. Its \emph{pairs} are the
  $(\Phi,\Phi')$ with $\Phi'\in\mathcal S^\circ(M')$ such that $\Phi$ and
  $\Phi'$ correspond; $\mathcal P$ is the set of them. For $s,t\in\mathbf C$
  and $y=(\Phi,\Phi')\in\mathcal P$, put
  $\mathcal I_{s,t}(y)=s\,I(\Phi)+t\,I'(\Phi')$, a function on $U_W$.
\end{enumerate}
\end{definition}
\leanref{Development/WeilBound/ThetaLifts/LatticeBasics/PairsStable.lean}{GlobalTheory.ThK.Dbl.pet}
\leanref{Development/WeilBound/ThetaLifts/LatticeBasics/PairsStable.lean}{GlobalTheory.ThK.Dbl.IsRapid}
\leanref{Development/WeilBound/ThetaLifts/LatticeBasics/PairsStable.lean}{GlobalTheory.ThK.Dbl.th}
\leanref{Development/WeilBound/ThetaLifts/LatticeBasics/PairsStable.lean}{GlobalTheory.ThK.Dbl.pairs}
\leanref{Development/WeilBound/ThetaLifts/LatticeBasics/PairsStable.lean}{GlobalTheory.ThK.Dbl.comb}

\begin{lemma}\label{asm-theta-agree-definite-compact}
Let $B$ be a theta base for a CM field $E$
(Definition~\ref{asm-theta-agree-models}(1)), and let $Y$ be a lattice space
that is Hermitian or skew-Hermitian of signature $(r,s)$ with $\min(r,s)=0$
(Definition~\ref{asm-theta-agree-doubling}(1)). Then:
\begin{enumerate}
\item $U_Y$ is locally compact, second countable and Hausdorff, and $U_Y(F)$ is
  countable;
\item there is a compact set $\Omega\subset U_Y$ such that every $g\in U_Y$ has
  $\gamma g\in\Omega$ for some $\gamma\in U_Y(F)$, and
  $\overline{\operatorname{ht}}$
  (Definition~\ref{asm-theta-agree-definite-products}(1)) is bounded on $U_Y$;
\item $U_Y$ has a Haar measure and a fundamental domain for $U_Y(F)$ with
  respect to it; every Haar measure $\mu$ on $U_Y$ is invariant under right
  translations; and every fundamental domain $\mathcal D$ for $U_Y(F)$ with
  respect to $\mu$ has $0<\mu(\mathcal D)<\infty$;
\item every continuous function $f$ on $U_Y$ with $f(\gamma g)=f(g)$ for all
  $\gamma\in U_Y(F)$ and $g\in U_Y$ is bounded, and if moreover
  $\langle f,f\rangle_{\mathcal D}=0$ for $\mu$ and $\mathcal D$ as in (3),
  then $f=0$; and if $Y'$ is a second lattice space as in the lemma, every
  continuous function on $U_Y\times U_{Y'}$ invariant under left translation by
  $U_Y(F)\times U_{Y'}(F)$ is bounded;
\item every $f$ as in (4) is rapidly decreasing
  (Definition~\ref{asm-theta-agree-definite-products}(1)).
\end{enumerate}
\end{lemma}
\leanref{Development/WeilBound/ThetaLifts/LatticeBasics/FdEqZero.lean}{GlobalTheory.ThK.Dbl.Cpt}
\leanref{Development/WeilBound/ThetaLifts/LatticeBasics/FdEqZero.lean}{GlobalTheory.ThK.Dbl.CptAll}
\leanref{Development/WeilBound/ThetaLifts/LatticeBasics/FdEqZero.lean}{GlobalTheory.ThK.Dbl.cptAll_of}
\leanref{Development/WeilBound/ThetaLifts/LatticeBasics/FdEqZero.lean}{GlobalTheory.ThK.Dbl.Cpt.bdd}
\leanref{Development/WeilBound/ThetaLifts/LatticeBasics/FdEqZero.lean}{GlobalTheory.ThK.Dbl.Cpt.bdd₂}
\leanref{Development/WeilBound/ThetaLifts/LatticeBasics/FdEqZero.lean}{GlobalTheory.ThK.Dbl.Cpt.vol_ne}
\leanref{Development/WeilBound/ThetaLifts/LatticeBasics/FdEqZero.lean}{GlobalTheory.ThK.Dbl.Cpt.eq_zero}
\leanref{Development/WeilBound/ThetaLifts/LatticeBasics/FdEqZero.lean}{GlobalTheory.ThK.Dbl.Cpt.rapid}
\begin{proof}
By Lemma~\ref{asm-theta-agree-convergence-signature}, the matrix $J$ of $Y$ is
nondegenerate $\epsilon$-Hermitian, and every Witt torus of $J$ has size at
most $\min(r,s)=0$. Theorem~\ref{asm-theta-agree-convergence-siegel}(1) gives
(1), and that the height is bounded on every compact subset of $U_Y$.

(2) By Theorem~\ref{asm-theta-agree-convergence-siegel}(2), $J$ has a Siegel
Witt torus $a$, of size $0$. Then $a(t)=1$ for the only element $t$ of
$\mathbf R^0$, and the Siegel property (a) gives a compact set $\Omega$ such
that every $g$ has $\gamma g\in\Omega$ for some $\gamma\in U_Y(F)$. If the
height is bounded by $C$ on $\Omega$, then
$\overline{\operatorname{ht}}(g)\leq\operatorname{ht}(\gamma g)\leq C$.

(3) Theorem~\ref{asm-theta-agree-definite-volume} gives all of (3) except
$\mu(\mathcal D)>0$. As $U_Y(F)$ is countable and the translates
$\gamma\mathcal D$ cover $U_Y$ up to a null set, $\mu(\mathcal D)=0$ would give
$\mu=0$.

(4) Let $f$ be as in (4). It is bounded on the compact set $\Omega$ of (2), say
by $C$, and $f(g)=f(\gamma g)$ with $\gamma g\in\Omega$; so $|f|\leq C$. The
same argument with $\Omega\times\Omega'$ applies on $U_Y\times U_{Y'}$. Suppose
now that $\langle f,f\rangle_{\mathcal D}=\int_{\mathcal D}|f|^2\,d\mu=0$;
the integral converges, as $f$ is bounded and $\mu(\mathcal D)<\infty$. The set
$S=\{g:f(g)\neq0\}$ is invariant under $U_Y(F)$ and $S\cap\mathcal D$ is null,
so
$\mu(S)\leq\sum_\gamma\mu(S\cap\gamma\mathcal D)=\sum_\gamma\mu(\gamma(S\cap\mathcal D))=0$,
the sums over $\gamma\in U_Y(F)$. As $f$ is continuous, $S$ is open, and a
nonempty open set has positive Haar measure. So $S$ is empty and $f=0$.

(5) If $|f|\leq C$ and $\overline{\operatorname{ht}}\leq C'$, then
$|f|\,\overline{\operatorname{ht}}^N\leq CC'^N$.
\end{proof}

\begin{lemma}\label{asm-theta-agree-definite-continuous}
Let $(X,W,M,\chi_X,\chi_W,\omega)$ be a Weil pair of signature $(r,s;r',s')$
for a theta base $B$ of a CM field $E$
(Definition~\ref{asm-theta-agree-triple}(1)), with $X$ and $W$ of ranks $n$ and
$d$, and suppose that $d+2\min(r,s)<2n$ or $\min(r,s)=0$. Then for every
Schwartz-Bruhat function $\Phi$ on $\mathbf A^{n\times d}$, the theta integral
$I(\Phi)$ is a continuous function on $U_W$.
\end{lemma}
\leanref{Development/WeilBound/ThetaLifts/CartanFactorisation/CartanConj.lean}{PositivityTheory.ThKF.ICont}
\leanref{Development/WeilBound/ThetaLifts/FrameEstimates/IContOf.lean}{PositivityTheory.ThKEst.iCont_of}
\leanref{Development/WeilBound/ThetaLifts/FrameEstimates/DomContinuous.lean}{PositivityTheory.ThKEst.dom_continuous}
\begin{proof}
Write $\Phi=\sum_if_i\otimes\phi_i$ with $f_i\in\mathcal S(F_\infty^{n\times d})$
and $\phi_i$ Schwartz-Bruhat on $\mathbf A_f^{n\times d}$. By
Lemma~\ref{asm-theta-agree-convergence}(1) each $f_i\otimes\phi_i$ converges,
and $\theta_\omega(\Phi)=\sum_i\theta_\omega(f_i\otimes\phi_i)$ by
Lemma~\ref{asm-theta-agree-convergence-theta}(2), so
$I(\Phi)=\sum_iI(f_i\otimes\phi_i)$. It suffices to show that
$I(f\otimes\phi)$ is continuous for such $f$ and $\phi$.

Let $k$ be as in Lemma~\ref{asm-theta-agree-convergence-dominated}, and $p$ and
$G$ as there for $\phi$. Let $h_0\in U_W$. By
Lemma~\ref{asm-theta-agree-convergence-signature} and
Theorem~\ref{asm-theta-agree-convergence-siegel}(1), $U_W$ is locally compact
and second countable, so $h_0$ has a compact neighbourhood $N$, on which the
height is bounded by some $C$. For $h\in N$ the function
$g\mapsto\theta_\omega(f\otimes\phi)(g,h)$ is bounded by $p(f)C^kG(g)$, which
is integrable on $\mathcal D$, and for each $g$ it depends continuously on $h$,
by Lemma~\ref{asm-theta-agree-convergence-theta}(2). By dominated convergence
along sequences,
$h\mapsto I(f\otimes\phi)(h)=\int_{\mathcal D}\theta_\omega(f\otimes\phi)(g,h)\,d\mu(g)$
is continuous at $h_0$.
\end{proof}

\begin{lemma}\label{asm-theta-agree-definite-dist-tools}
Call a complex function on a topological space \emph{Schwartz-Bruhat} if it is
locally constant with compact support. Let $V$ be a Hausdorff abelian
topological group in which every neighbourhood of $0$ contains an open
subgroup.
\begin{enumerate}
\item Let $l$ be a linear form on the complex functions on $V$, and $\phi$ a
  Schwartz-Bruhat function on $V$ such that every $x$ with $\phi(x)\neq0$ has
  a neighbourhood $U$ with $l(\psi)=0$ for every Schwartz-Bruhat $\psi$
  vanishing outside $U$. Then $l(\phi)=0$.
\item Suppose that $V$ is locally compact, let $Z$ be a Hausdorff space and
  $S\subset Z$, and let $l$ be a linear form on the complex functions on
  $Z\times V$ with $l(f\otimes g)=0$ for all Schwartz-Bruhat functions $f$ on
  $Z$ vanishing outside $S$ and $g$ on $V$, where
  $(f\otimes g)(z,x)=f(z)g(x)$. Then $l(F)=0$ for every Schwartz-Bruhat $F$
  on $Z\times V$ vanishing outside $S\times V$.
\item Suppose that $V$ is locally compact, and call a Schwartz-Bruhat function
  on $V$ with values in $\{0,1\}$ and value $1$ at $0$ a \emph{bump}. Then $V$
  has a bump, and for bumps $g_0$ and $g_1$ there is a real number $s>0$ such
  that every linear form $m$ on the complex functions on $V$ with
  $m(g(\cdot+a))=m(g)$ for all Schwartz-Bruhat $g$ and all $a\in V$ has
  $m(g_0)=s\,m(g_1)$, and has $m(g)=0$ for every Schwartz-Bruhat $g$ if
  $m(g_0)=0$.
\end{enumerate}
\end{lemma}
\leanref{Development/WeilBound/ThetaLifts/LatticeBasics/LocDistA.lean}{GlobalTheory.ThK.Dbl.LD.kill_loc}
\leanref{Development/WeilBound/ThetaLifts/LatticeBasics/LocDistA.lean}{GlobalTheory.ThK.Dbl.LD.kill_tens}
\leanref{Development/WeilBound/ThetaLifts/LatticeBasics/LocDistA.lean}{GlobalTheory.ThK.Dbl.LD.Bump}
\leanref{Development/WeilBound/ThetaLifts/LatticeBasics/LocDistA.lean}{GlobalTheory.ThK.Dbl.LD.exists_bump}
\leanref{Development/WeilBound/ThetaLifts/LatticeBasics/LocDistA.lean}{GlobalTheory.ThK.Dbl.LD.bump_ratio}
\begin{proof}
Write $1_A$ for the indicator function of a set $A$. The support of a
Schwartz-Bruhat function is compact, being closed, as the function is locally
constant, and contained in a compact set.

(1) Each $x$ in the support of $\phi$ has a neighbourhood $U_x$ as in the
hypothesis, and an open subgroup $K_x$ with $x+K_x\subset U_x$; the set
$x+K_x$ is open and closed. Finitely many of these sets, $A_1,\ldots,A_r$, cover
the support. The sets $B_i=A_i\setminus(A_1\cup\cdots\cup A_{i-1})$ are open
and closed, so the functions $1_{B_i}\phi$ are Schwartz-Bruhat, and
$\phi=\sum_i1_{B_i}\phi$. As $1_{B_i}\phi$ vanishes outside $U_{x_i}$ for the
centre $x_i$ of $A_i$, $l(1_{B_i}\phi)=0$. So $l(\phi)=0$.

(2) As $V$ is locally compact, a compact neighbourhood of $0$ contains an open
subgroup $K_0$, which is closed and so compact. Let $C$ be the support of $F$.
As $F$ is locally constant and $C$ is compact, there is an open subgroup $K_1$
of $V$ with $F(z,x+a)=F(z,x)$ for $(z,x)\in C$ and $a\in K_1$; this then holds
for all $(z,x)$, since if $(z,x+a)\in C$ it follows from the case of $(z,x+a)$
and $-a$. Put $K=K_0\cap K_1$, a compact open subgroup. The image of $C$ in the
discrete group $V/K$ is compact, hence finite; let $Q$ be a set of
representatives of its elements. Then
\[
F(z,x)=\sum_{q\in Q}F(z,q)\,1_{q+K}(x)
\]
for all $(z,x)$. Each function $z\mapsto F(z,q)$ is Schwartz-Bruhat and
vanishes outside $S$, and $1_{q+K}$ is Schwartz-Bruhat; so $l(F)=0$.

(3) The indicator function of a compact open subgroup, which exists as in (2),
is a bump. Let $\mu$ be a Haar measure on $V$ and $m$ a linear form as in (3).
Let $g$ be Schwartz-Bruhat. As in (2), there is a compact open subgroup $K$
such that $g$ is constant on the cosets of $K$, so $g$ is a finite linear
combination of functions $1_{a+K}$, and $m(1_{a+K})=m(1_K)$. If $K'\subset K$
are compact open subgroups, then $1_K$ is the sum of the $[K:K']$ functions
$1_{a+K'}$ for the cosets of $K'$ in $K$, so $m(1_K)=[K:K']\,m(1_{K'})$, and
likewise $\mu(K)=[K:K']\,\mu(K')$. As two compact open subgroups contain their
intersection, $m(1_K)=c\,\mu(K)$ for a constant $c$ independent of $K$, and so
$m(g)=c\int g\,d\mu$. A bump is the indicator function of a compact open set
containing $0$, so its integral is positive. Put
$s=\int g_0\,d\mu/\int g_1\,d\mu$. Then $m(g_0)=s\,m(g_1)$, and if $m(g_0)=0$
then $c=0$.
\end{proof}

\begin{lemma}\label{asm-theta-agree-definite-dist-group}
Let $G$ be a Hausdorff topological group, $\chi$ a homomorphism from $G$ to the
multiplicative monoid of $\mathbf C$, $\theta$ a continuous endomorphism of $G$
with $\theta(\theta(x))=x$ for all $x$, and $g\in G$ with $\theta(g)=g^{-1}$.
Let $l$ be a linear form on the complex functions on $G$ such that
$l(\phi(\cdot\,h))=\chi(h)\,l(\phi)$ for every Schwartz-Bruhat $\phi$
(Lemma~\ref{asm-theta-agree-definite-dist-tools}) and every $h\in G$. Then
$l(\phi(g\,\cdot))=\chi(g)\,l(\phi)$ for every Schwartz-Bruhat $\phi$.
\end{lemma}
\leanref{Development/WeilBound/ThetaLifts/LatticeBasics/LocDistB.lean}{GlobalTheory.ThK.Dbl.LD.left_of_right}
\leanref{Development/WeilBound/ThetaLifts/LatticeBasics/LocDistB.lean}{GlobalTheory.ThK.Dbl.LD.ind_index}
\leanref{Development/WeilBound/ThetaLifts/LatticeBasics/LocDistB.lean}{GlobalTheory.ThK.Dbl.LD.ind_zero}
\begin{proof}
Write $1_K$ for the indicator function of a subgroup $K$. We may assume
$\phi\neq0$; let $\phi(x_0)\neq0$.

\emph{A subgroup.} The set $K_0$ of the $k\in G$ with $\phi(kx)=\phi(x)$ for
all $x$ is a subgroup. It is open, as $\phi$ is locally constant with compact
support; so it is closed, and it lies in $Cx_0^{-1}$, with $C$ the support of
$\phi$, so it is compact. As $\theta$ is an automorphism with $\theta^2=1$,
$K=K_0\cap\theta^{-1}(K_0)$ is a compact open subgroup with $\theta(K)=K$, and
so is $K'=g^{-1}Kg$. The automorphism $\alpha(x)=\theta(gxg^{-1})$ has
$\alpha(K)=\theta(g)K\theta(g)^{-1}=K'$ and $\alpha(K')=\theta(K)=K$.

\emph{The two indicator functions.} We show $l(1_{K'})=l(1_K)$. Suppose first
that $\chi=1$ on $K$; then $\chi=1$ on $K'$, as $\chi(g^{-1}kg)=\chi(k)$. Put
$H=K\cap K'$, an open subgroup of finite index in the compact groups $K$ and
$K'$. If $Q$ is a set of representatives of $K/H$, then
$1_K=\sum_{q\in Q}1_H(\cdot\,q)$, so $l(1_K)=[K:H]\,l(1_H)$, and likewise
$l(1_{K'})=[K':H]\,l(1_H)$. As $\alpha$ maps $K$ onto $K'$ and $H$ onto
$K'\cap K=H$, $[K:H]=[K':H]$. Suppose next that $\chi(k)\neq1$ for some
$k\in K$. As $1_K(\cdot\,k)=1_K$, $l(1_K)=\chi(k)\,l(1_K)$, so $l(1_K)=0$; and
likewise $l(1_{K'})=0$, with $g^{-1}kg\in K'$.

\emph{Conclusion.} As $\phi$ is invariant under left translation by $K$ and has
compact support, there is a finite set $Q\subset G$ with
$\phi(x)=\sum_{q\in Q}\phi(q^{-1})\,1_K(xq)$ for all $x$: take representatives
of the cosets $x^{-1}K$ with $x$ in the support. Since $gxq\in K$ if and only if
$xqg\in K'$,
\[
l(\phi(g\,\cdot))=\sum_{q\in Q}\phi(q^{-1})\,l\bigl(1_{K'}(\cdot\,qg)\bigr)
=\chi(g)\sum_{q\in Q}\phi(q^{-1})\chi(q)\,l(1_K)
=\chi(g)\sum_{q\in Q}\phi(q^{-1})\,l\bigl(1_K(\cdot\,q)\bigr)=\chi(g)\,l(\phi).
\]
\end{proof}

\begin{lemma}\label{asm-theta-agree-definite-dist-degenerate}
Let $k$ be a nonarchimedean local field, let $\varpi\in k$ with
$\varpi\neq0,1$, and let $c\colon k\to\mathbf C$ be a function with
$c(a)\neq0$ for $a\neq0$ and $c(\varpi)\,t\neq1$ for every real $t>0$. Let
$n,d\geq0$; call $X\in M_{n\times d}(k)$ \emph{degenerate} if $Xy=0$ for some
nonzero $y\in k^d$, and write $\phi^h$ for the function $X\mapsto\phi(Xh)$.
Let $l$ be a linear form on the complex functions on $M_{n\times d}(k)$ such
that:
\begin{enumerate}
\item[(a)] for every $h\in GL_d(k)$ there is a real $t>0$ with
  $l(\phi^h)=c(\det h)\,t\,l(\phi)$ for every Schwartz-Bruhat $\phi$
  (Lemma~\ref{asm-theta-agree-definite-dist-tools});
\item[(b)] $l(\phi)=0$ for every Schwartz-Bruhat $\phi$ that vanishes at every
  degenerate matrix.
\end{enumerate}
Then $l(\phi)=0$ for every Schwartz-Bruhat $\phi$.
\end{lemma}
\leanref{Development/WeilBound/ThetaLifts/LatticeBasics/LocDistD.lean}{GlobalTheory.ThK.Dbl.LD.Claim}
\leanref{Development/WeilBound/ThetaLifts/LatticeBasics/LocDistD.lean}{GlobalTheory.ThK.Dbl.LD.claim}
\leanref{Development/WeilBound/ThetaLifts/LatticeBasics/LocDistC.lean}{GlobalTheory.ThK.Dbl.LD.Eqv}
\leanref{Development/WeilBound/ThetaLifts/LatticeBasics/LocDistC.lean}{GlobalTheory.ThK.Dbl.LD.Eqv.unip}
\leanref{Development/WeilBound/ThetaLifts/LatticeBasics/LocDistD.lean}{GlobalTheory.ThK.Dbl.LD.step_chart}
\leanref{Development/WeilBound/ThetaLifts/LatticeBasics/LocDistD.lean}{GlobalTheory.ThK.Dbl.LD.step_row}
\leanref{Development/WeilBound/ThetaLifts/LatticeBasics/LocDistD.lean}{GlobalTheory.ThK.Dbl.LD.step_closed}
\begin{proof}
The spaces $k^m$ and $M_{m\times m'}(k)$ have the properties required of $V$
in Lemma~\ref{asm-theta-agree-definite-dist-tools}. We argue by induction on
$n$, for all $d$ at once. If $d=0$, no matrix is degenerate, so $l(\phi)=0$ by
(b). If $n=0$ and $d\geq1$, then $M_{0\times d}(k)$ is a point and
$\phi^h=\phi$; for $h=\operatorname{diag}(\varpi,1,\ldots,1)$, (a) gives
$l(\phi)=c(\varpi)\,t\,l(\phi)$ with $c(\varpi)t\neq1$, so $l(\phi)=0$.

Now let $n,d\geq0$, assume the lemma for $n$ rows and every number of columns,
and let $l$ be as in the lemma for $M_{(n+1)\times(d+1)}(k)$. Index rows and
columns from $0$. For a row vector $b\in k^d$ and $h\in GL_d(k)$ put
\[
u_b=\begin{pmatrix}1&b\\0&1_d\end{pmatrix},\qquad
m_h=\begin{pmatrix}1&0\\0&h\end{pmatrix},\qquad
D=\begin{pmatrix}\varpi&0\\0&1_d\end{pmatrix},
\]
so that $\det u_b=1$, $\det m_h=\det h$ and $\det D=\varpi$. As
$\phi^{gh}=(\phi^h)^g$, (a) shows that $l(\phi^{gh})$ is $c(\det g)c(\det h)$
times a positive number times $l(\phi)$.

\emph{Step 1: $l(\phi^{u_b})=l(\phi)$.} Put $b'=(\varpi-1)^{-1}b$, so that
$\varpi b'=b+b'$ and $Du_{b'}=u_bu_{b'}D$. Let $t_1$, $t_2$ and $t_3$ be as in
(a) for $u_b$, $u_{b'}$ and $D$. Computing $l(\phi^{Du_{b'}})$ and
$l(\phi^{u_bu_{b'}D})$ by (a) gives
\[
c(\varpi)t_3\,c(1)t_2\,l(\phi)=c(1)t_1\,c(1)t_2\,c(\varpi)t_3\,l(\phi).
\]
Cancelling, $l(\phi)=c(1)t_1\,l(\phi)=l(\phi^{u_b})$.

\emph{Step 2: $l(\phi)=0$ if $\phi$ vanishes at every $X$ with $X_{00}=0$.}
Let $\Omega$ be the set of $(x,b,Y)\in k^{n+1}\times k^d\times M_{n\times d}(k)$
with $x_0\neq0$, indexing the entries of $b$ and $Y$ from $1$. For $X$ with
$X_{00}\neq0$ put $E(X)=(x,b,Y)$ with
\[
x_i=X_{i0},\qquad b_j=X_{0j}/X_{00},\qquad Y_{ij}=X_{ij}-X_{i0}X_{0j}/X_{00}\quad(i,j\geq1).
\]
Then $E$ is a homeomorphism onto $\Omega$, with inverse the restriction of the
continuous map $R$ that sends $(x,b,Y)$ to the matrix with column $0$ equal to
$x$, entries $x_0b_j$ in row $0$, and $x_ib_j+Y_{ij}$ elsewhere. For
$X_{00}\neq0$ a direct computation gives $E(Xu_{b'})=(x,b+b',Y)$ and
$E(Xm_h)=(x,bh,Yh)$; and if $X$ is degenerate, so is $Y$: if $Xy=0$ with
$y\neq0$, then $y_0=-X_{00}^{-1}\sum_{j\geq1}X_{0j}y_j$, so
$(y_1,\ldots,y_d)\neq0$, and it is killed by $Y$. For a Schwartz-Bruhat
function $F$ on $k^{n+1}\times k^d\times M_{n\times d}(k)$ vanishing where
$x_0=0$, let $\tilde F$ be equal to $F\circ E$ where $X_{00}\neq0$ and to $0$
elsewhere; it is Schwartz-Bruhat, as the support of $F$ is a compact subset of
$\Omega$. For any function $F$ define $\tilde F$ in the same way, and put
$\Lambda(F)=l(\tilde F)$.

Let $\phi$ be as in Step 2. Then $F=\phi\circ R$ is Schwartz-Bruhat, vanishes
where $x_0=0$, and has $\tilde F=\phi$. By
Lemma~\ref{asm-theta-agree-definite-dist-tools}(2), applied to $\Lambda$ with
$Z=k^{n+1}\times k^d$, $V=M_{n\times d}(k)$ and $S$ given by $x_0\neq0$, and
then, for fixed $\psi$, to $F_1\mapsto\Lambda(F_1\otimes\psi)$ with
$Z=k^{n+1}$, $V=k^d$ and $S$ given by $x_0\neq0$, it suffices to show that
$\Lambda(f\otimes g\otimes\psi)=0$ for Schwartz-Bruhat functions $f$ on
$k^{n+1}$ vanishing where $x_0=0$, $g$ on $k^d$ and $\psi$ on
$M_{n\times d}(k)$.

Fix $f$ and $\psi$, and put $m(g)=\Lambda(f\otimes g\otimes\psi)$. By Step 1
and the formula for $E(Xu_{b'})$, $m(g(\cdot+b'))=m(g)$ for Schwartz-Bruhat
$g$ and all $b'$. Let $g_0$ be a bump on $k^d$. By
Lemma~\ref{asm-theta-agree-definite-dist-tools}(3), it suffices that
$m(g_0)=0$. Put $l_1(\psi')=\Lambda(f\otimes g_0\otimes\psi')$ for functions
$\psi'$ on $M_{n\times d}(k)$. It satisfies (a) and (b) for $n$ rows and $d$
columns. For (a), let $h\in GL_d(k)$ and $t$ be as in (a) for $m_h$; the
function $g_1(b)=g_0(bh)$ is a bump, and let $s$ be as in
Lemma~\ref{asm-theta-agree-definite-dist-tools}(3) for $g_0$ and $g_1$. That
lemma, applied to the form $g\mapsto\Lambda(f\otimes g\otimes\psi'^h)$, which is
invariant under translation as above, the formula for $E(Xm_h)$, and (a) for
$m_h$ give, with $F_1=f\otimes g_0\otimes\psi'$,
\[
l_1(\psi'^h)=s\,\Lambda(f\otimes g_1\otimes\psi'^h)=s\,l\bigl((\tilde F_1)^{m_h}\bigr)=s\,c(\det h)\,t\,l_1(\psi').
\]
For (b), if $\psi'$ is Schwartz-Bruhat and vanishes at every degenerate
matrix, then $\tilde F_1$ vanishes at every degenerate $X$, as the last
component of $E(X)$ is then degenerate; so
$l_1(\psi')=0$ by (b) for $l$. By the induction hypothesis, $l_1$ vanishes on
the Schwartz-Bruhat functions, and $m(g_0)=0$.

\emph{Step 3: $l(\phi)=0$ if $\phi$ vanishes at every $X$ whose row $0$ is
zero.} By Lemma~\ref{asm-theta-agree-definite-dist-tools}(1), it suffices that
every $X$ whose row $0$ is nonzero has a neighbourhood $U$ such that $l$ kills
every Schwartz-Bruhat $\psi$ vanishing outside $U$. Choose $j$ with
$X_{0j}\neq0$, let $U$ be the set of $X'$ with $X'_{0j}\neq0$, and let $s$ be
the permutation matrix that exchanges the columns $0$ and $j$. If $\psi$
vanishes outside $U$, then $\psi^s$ vanishes at every $X'$ with $X'_{00}=0$, so
$l(\psi^s)=0$ by Step 2; and $l(\psi^s)=c(\det s)\,t\,l(\psi)$ with
$c(\det s)t\neq0$. So $l(\psi)=0$.

\emph{Step 4: conclusion.} By Step 3, $l(\phi)=l(\phi')$ whenever $\phi$ and
$\phi'$ are Schwartz-Bruhat and agree at every $X$ whose row $0$ is zero. For
$X\in M_{(n+1)\times(d+1)}(k)$ let $X_0\in k^{d+1}$ be its row $0$ and
$X'\in M_{n\times(d+1)}(k)$ the matrix of its other rows. Let $g_0$ be a bump
on $k^{d+1}$, and put $l_2(\psi)=l\bigl(X\mapsto g_0(X_0)\psi(X')\bigr)$ for
functions $\psi$ on $M_{n\times(d+1)}(k)$. Then $l_2$ satisfies (a) and (b) for
$n$ rows and $d+1$ columns. For (a), let $h\in GL_{d+1}(k)$; the functions
$X\mapsto g_0(X_0)\psi(X'h)$ and $X\mapsto g_0((Xh)_0)\psi((Xh)')$ agree where
$X_0=0$, as $(Xh)_0=X_0h$, $(Xh)'=X'h$ and $g_0(0)=1$, so
$l_2(\psi^h)=c(\det h)\,t\,l_2(\psi)$ with $t$ as in (a) for $l$. For (b), if
$Xy=0$ then $X'y=0$. By the induction hypothesis, $l_2$ vanishes on the
Schwartz-Bruhat functions. Every Schwartz-Bruhat $\phi$ agrees, where $X_0=0$,
with $X\mapsto g_0(X_0)\phi_1(X')$, where $\phi_1(Y)$ is the value of $\phi$ at
the matrix with row $0$ zero and other rows $Y$. So $l(\phi)=l_2(\phi_1)=0$.
\end{proof}

\begin{lemma}\label{asm-theta-agree-definite-dist}
Let $v$ be a finite place of a number field $F$, $\varpi$ a uniformizer of
$F_v$, and $\xi\colon F_v^\times\to\mathbf C^\times$ a homomorphism such that
$\xi(\varpi)$ is not a positive real number. Let $n,d\geq0$, and let $l$ be a
linear form on the complex functions on $M_{n\times d}(F_v)$ with
$l(\phi(\cdot\,h))=\xi(\det h)\,l(\phi)$ for every Schwartz-Bruhat $\phi$
(Lemma~\ref{asm-theta-agree-definite-dist-tools}) and every
$h\in GL_d(F_v)$.
\begin{enumerate}
\item If $n<d$, then $l(\phi)=0$ for every Schwartz-Bruhat $\phi$.
\item If $n=d$, then $l(\phi(g\,\cdot))=\xi(\det g)\,l(\phi)$ for every
  Schwartz-Bruhat $\phi$ and every $g\in GL_n(F_v)$.
\end{enumerate}
\end{lemma}
\leanref{Development/WeilBound/ThetaLifts/LatticeBasics/PairsStable.lean}{GlobalTheory.ThK.Dbl.LocDist}
\leanref{Development/WeilBound/ThetaLifts/LatticeBasics/LocDist.lean}{GlobalTheory.ThK.Dbl.locDist}
\leanref{Development/WeilBound/ThetaLifts/LatticeBasics/LocDistE.lean}{GlobalTheory.ThK.Dbl.LD.vanish}
\leanref{Development/WeilBound/ThetaLifts/LatticeBasics/LocDistE.lean}{GlobalTheory.ThK.Dbl.LD.lt_case}
\leanref{Development/WeilBound/ThetaLifts/LatticeBasics/LocDistE.lean}{GlobalTheory.ThK.Dbl.LD.top_cell}
\leanref{Development/WeilBound/ThetaLifts/LatticeBasics/LocDistE.lean}{GlobalTheory.ThK.Dbl.LD.eq_case}
\begin{proof}
Put $k=F_v$, and let $c(a)=\xi(a)$ for $a\neq0$ and $c(0)=0$. Then $c(a)\neq0$
for $a\neq0$; $\varpi\neq1$, as $\xi(1)=1$; and $c(\varpi)t\neq1$ for real
$t>0$, since otherwise $\xi(\varpi)=t^{-1}$. So $l$ satisfies (a) of
Lemma~\ref{asm-theta-agree-definite-dist-degenerate} with $t=1$.

(1) If $n<d$, every $X\in M_{n\times d}(k)$ is degenerate, so a function
vanishing at every degenerate matrix is $0$, and (b) of
Lemma~\ref{asm-theta-agree-definite-dist-degenerate} holds. That lemma gives
the claim.

(2) Let $n=d$; then $X$ is degenerate if and only if $\det X=0$. Write
${}^g\phi$ for the function $\phi(g\,\cdot)$.

\emph{Symmetric $g$, and $\phi$ vanishing on singular matrices.} Let
$g^{\mathsf T}=g$, and let $\phi$ vanish at every singular matrix. The group
$G=GL_d(k)$ is an open subset of $M_d(k)$. For a function $\psi$ on $G$ let
$\tilde\psi$ be the function on $M_d(k)$ equal to $\psi$ on $G$ and to $0$
elsewhere; if $\psi$ is Schwartz-Bruhat, so is $\tilde\psi$, as its support is
a compact subset of the open set $G$. So $m(\psi)=l(\tilde\psi)$ is a linear
form on the functions on $G$ with $m(\psi(\cdot\,h))=\xi(\det h)\,m(\psi)$ for
Schwartz-Bruhat $\psi$ and $h\in G$. The map $\theta(x)=(x^{\mathsf T})^{-1}$ is
a continuous endomorphism of $G$ with $\theta^2=1$, and $\theta(g)=g^{-1}$. By
Lemma~\ref{asm-theta-agree-definite-dist-group} for $\chi=\xi\circ\det$,
$m(\psi(g\,\cdot))=\xi(\det g)\,m(\psi)$. For $\psi$ the restriction of $\phi$
to $G$, $\tilde\psi=\phi$ and $\widetilde{\psi(g\,\cdot)}={}^g\phi$; so
$l({}^g\phi)=\xi(\det g)\,l(\phi)$.

\emph{Symmetric $g$.} Let $g^{\mathsf T}=g$, and put
$l'(\phi)=l({}^g\phi)-\xi(\det g)\,l(\phi)$. As ${}^g(\phi^h)=({}^g\phi)^h$,
$l'$ satisfies (a) of Lemma~\ref{asm-theta-agree-definite-dist-degenerate} with
$t=1$, and by the last paragraph it satisfies (b). By that lemma, $l'(\phi)=0$
for every Schwartz-Bruhat $\phi$.

\emph{All $g$.} Let $P$ be the set of $g\in GL_d(k)$ with
$l({}^g\phi)=\xi(\det g)\,l(\phi)$ for every Schwartz-Bruhat $\phi$. If
$g,g'\in P$, then $gg'\in P$, as ${}^{gg'}\phi={}^{g'}({}^g\phi)$. By the last
paragraph, $P$ contains the invertible symmetric matrices: the invertible
diagonal matrices, and, for $i\neq j$ and $a\in k$, the permutation matrix $s$
of the transposition of $i$ and $j$ and the matrix $s(1+aE_{ij})=s+aE_{jj}$,
where $E_{ij}$ is a matrix unit. So $P$ contains the transvections
$1+aE_{ij}=s\cdot s(1+aE_{ij})$. As $GL_d(k)$ is generated by its invertible
diagonal matrices and its transvections, $P=GL_d(k)$.
\end{proof}

\begin{lemma}\label{asm-theta-agree-definite-dist-bilinear}
Let $v$, $\varpi$ and $\xi$ be as in Lemma~\ref{asm-theta-agree-definite-dist},
let $n,n'\geq0$ and $d=n+n'$, and let $\beta(a,b)\in\mathbf C$ be given for
Schwartz-Bruhat functions $a$ on $M_{n\times d}(F_v)$ and $b$ on
$M_{n'\times d}(F_v)$, bilinear in $a$ and $b$, with
$\beta(a(\cdot\,h),b(\cdot\,h))=\xi(\det h)\,\beta(a,b)$ for every
$h\in GL_d(F_v)$. Then $\beta(a(g\,\cdot),b)=\xi(\det g)\,\beta(a,b)$ for every
$g\in GL_n(F_v)$.
\end{lemma}
\leanref{Development/WeilBound/ThetaLifts/LatticeBasics/ThConstLoc.lean}{GlobalTheory.ThK.Dbl.TC.locBil}
\leanref{Development/WeilBound/ThetaLifts/LatticeBasics/ThConstLoc.lean}{GlobalTheory.ThK.Dbl.TC.boxL_ker}
\leanref{Development/WeilBound/ThetaLifts/LatticeBasics/ThConstLoc.lean}{GlobalTheory.ThK.Dbl.TC.mem_range_boxL}
\begin{proof}
For $Y\in M_d(F_v)$ let $Y_1$ and $Y_2$ be the matrices of its first $n$ and
its last $n'$ rows, and for functions $a$ and $b$ put
$(a\boxtimes b)(Y)=a(Y_1)\,b(Y_2)$. The linear map $a\otimes b\mapsto a\boxtimes b$,
from the tensor product of the spaces of Schwartz-Bruhat functions on
$M_{n\times d}(F_v)$ and on $M_{n'\times d}(F_v)$ to the functions on
$M_d(F_v)$, is injective: if $\sum_ia_i\boxtimes b_i=0$ with linearly
independent $a_i$, then $\sum_ib_i(Y_2)\,a_i=0$ for every $Y_2$, so every
$b_i=0$. Its image is the space of Schwartz-Bruhat functions on $M_d(F_v)$.
Indeed, if $\phi$ is one, the map $Y_1\mapsto\phi(Y_1,\cdot)$ is locally
constant, so it takes finitely many nonzero values $\phi_1,\ldots,\phi_r$ on
the compact projection of the support of $\phi$, each Schwartz-Bruhat, on open
and compact sets $A_1,\ldots,A_r$; and $\phi=\sum_i1_{A_i}\boxtimes\phi_i$.
Hence there is a linear form $l$ on the functions on $M_d(F_v)$ with
$l(a\boxtimes b)=\beta(a,b)$ for all Schwartz-Bruhat $a$ and $b$.

For $h\in GL_d(F_v)$, $(a\boxtimes b)(Yh)=a(Y_1h)\,b(Y_2h)$, so
$l(\phi(\cdot\,h))=\xi(\det h)\,l(\phi)$ for $\phi=a\boxtimes b$, and hence for
every Schwartz-Bruhat $\phi$. By Lemma~\ref{asm-theta-agree-definite-dist}(2),
$l(\phi(G\,\cdot))=\xi(\det G)\,l(\phi)$ for $G\in GL_d(F_v)$. For
$G=\operatorname{diag}(g,1_{n'})$, $(a\boxtimes b)(GY)=a(gY_1)\,b(Y_2)$ and
$\det G=\det g$, so
$\beta(a(g\,\cdot),b)=l\bigl((a\boxtimes b)(G\,\cdot)\bigr)=\xi(\det g)\,\beta(a,b)$.
\end{proof}

\begin{lemma}\label{asm-theta-agree-definite-hecke}
Let $E$ be a CM field, Galois over $\mathbf Q$, with maximal totally real
subfield $F$ and complex conjugation $c$, and for each infinite place $w$ of
$F$ let $\tau_w\colon E\to\mathbf C$ be an embedding over $w$. Let $H$ be a
Hecke character of $E$ (Definition~\ref{asm-theta-agree-doubling}) and $n(w)$
nonzero integers such that
$H(x_\infty)=\prod_w\bigl(\tau_w(x)/|\tau_w(x)|\bigr)^{n(w)}$ for every
$x\in E^\times$, where $x_\infty\in\mathbf A_E^\times$ has components $x$ at
the infinite places and $1$ at the finite places. Let $\mathfrak P$ be a prime
of $E$ whose absolute norm is a prime number $p$, with $p\notin\mathfrak P^2$,
and let $\varpi\in E^\times$ have $\operatorname{ord}_{\mathfrak P}(\varpi)=1$.
Then $H_{\mathfrak P}(\varpi)$ is not a root of unity, where $H_{\mathfrak P}$
is the component of $H$ at $\mathfrak P$.
\end{lemma}
\leanref{Development/AbelianVarietyTheory/Proofs/DivisorProofs/HeckeValueNotFiniteOrder.lean}{GlobalTheory.CMSplitValue.not_isOfFinOrder_val}
\leanref{Development/AbelianVarietyTheory/Proofs/DivisorProofs/HeckeValueNotFiniteOrder.lean}{GlobalTheory.CMSplitValue.not_isOfFinOrder_val_of_tame}
\leanref{Development/AbelianVarietyTheory/Proofs/DecoratedLatticeProofs/ExistsPowHeckeCharUnramifiedEverywhere.lean}{GlobalTheory.Corev.UnrPow.exists_pow_unrAt}
\leanref{Development/AbelianVarietyTheory/Proofs/GlobalInterpolation/ExistsPowLocKMulInf.lean}{GlobalTheory.TamePow.exists_pow_locK_mul_inf}
\leanref{Development/AbelianVarietyTheory/Proofs/TwistedOperators/HeckeCharacterInfiniteOrderAtSplit.lean}{GlobalTheory.SplitValue.val_map_eq_one}
\begin{proof}
Suppose that $H_{\mathfrak P}(\varpi)^a=1$ with $a\geq1$.

\emph{Reduction to unramified $H$.} The restriction of $H$ to the product $U$
of the groups $\mathcal O_{\mathfrak Q}^\times$ over the finite places
$\mathfrak Q$ of $E$ is continuous, so it is trivial on an open subgroup, which
contains $\mathcal O_{\mathfrak Q}^\times$ for $\mathfrak Q$ outside a finite
set $S$ and $1+\mathfrak Q^e\mathcal O_{\mathfrak Q}$ for $\mathfrak Q\in S$,
for some $e$. As the groups
$\mathcal O_{\mathfrak Q}^\times/(1+\mathfrak Q^e\mathcal O_{\mathfrak Q})$ are
finite, $H^N$ is trivial on $U$ for some $N\geq1$. The character $H^N$
satisfies the hypotheses with $Nn(w)$ in place of $n(w)$, and
$H^N_{\mathfrak P}(\varpi)=H_{\mathfrak P}(\varpi)^N$ is a root of unity. So we
may assume that $H$ is unramified at every finite place.

\emph{An element of $E$.} As the class group of $E$ is finite, there are
$M\geq1$ and $z\in E^\times$ with $z\mathcal O_E=\mathfrak P^M$. Then
$z=\varpi^Mu$ with $u$ a unit at $\mathfrak P$, and $z$ is a unit at every
other finite place. As $H$ is trivial on $E^\times$ and unramified,
$1=H(z)=H(z_\infty)\,H_{\mathfrak P}(\varpi)^M$, so $H(z_\infty)^a=1$.

\emph{Galois conjugates.} Fix an infinite place $w_0$ of $F$ and put
$\tau_0=\tau_{w_0}$. As $E/\mathbf Q$ is Galois, there are
$g_w\in\operatorname{Gal}(E/\mathbf Q)$ with $\tau_w=\tau_0\circ g_w$, and
$g_{w_0}=1$. As $E$ is a CM field, $\overline{\tau_w(x)}=\tau_w(c(x))$, so
$|\tau_w(z)|^2=\tau_w(z)\,\tau_w(c(z))$ and
\[
H(z_\infty)^2=\prod_w\Bigl(\frac{\tau_w(z)}{\tau_w(c(z))}\Bigr)^{n(w)}=\tau_0(y),\qquad
y=\prod_w\Bigl(\frac{g_w(z)}{g_w(c(z))}\Bigr)^{n(w)}\in E^\times .
\]
Then $\tau_0(y^a)=H(z_\infty)^{2a}=1$, so $y^a=1$ and
$\operatorname{ord}_{\mathfrak P}(y)=0$.

\emph{Valuations.} For $g\in\operatorname{Gal}(E/\mathbf Q)$,
$g(z)\mathcal O_E=g(\mathfrak P)^M$. As $\mathfrak P$ has ramification index and
residue degree $1$ over $p$, its decomposition group in
$\operatorname{Gal}(E/\mathbf Q)$ is trivial, so
$\operatorname{ord}_{\mathfrak P}(g(z))=0$ for $g\neq1$. Now $g_w\neq1$ for
$w\neq w_0$, as $\tau_w$ and $\tau_0$ differ on $F$; and $g_wc\neq1$ for every
$w$, since otherwise $g_w=c$ is the identity on $F$, so that $w=w_0$ and
$c=g_{w_0}=1$. Hence
$\operatorname{ord}_{\mathfrak P}(y)=n(w_0)\operatorname{ord}_{\mathfrak P}(z)=n(w_0)M\neq0$,
a contradiction.
\end{proof}

\begin{lemma}\label{asm-theta-agree-definite-inputs-tensor}
\begin{enumerate}
\item Let $F$ be a number field, $\kappa$ a finite set and $T$ a finite set of
  finite places of $F$, and for every finite place $v$ let $w_v\neq0$ be a
  function on $F_v^\kappa$, the indicator function of $\mathcal O_v^\kappa$
  for almost all $v$. Let $f_i$ and $f'_j$ be functions on $F_\infty^\kappa$
  and $\phi_{i,v}$ and $\phi'_{j,v}$ functions on $F_v^\kappa$, for $i$ and $j$
  in finite sets and $v$ finite, with $\phi_{i,v}=\phi'_{j,v}=w_v$ for
  $v\notin T$. If
  $\sum_if_i\otimes\bigotimes_v\phi_{i,v}=\sum_jf'_j\otimes\bigotimes_v\phi'_{j,v}$
  as functions on $\mathbf A_F^\kappa$, then
  $\sum_if_i\otimes\bigotimes_{v\in T}\phi_{i,v}=\sum_jf'_j\otimes\bigotimes_{v\in T}\phi'_{j,v}$
  in the tensor product of the spaces of functions on $F_\infty^\kappa$ and on
  the $F_v^\kappa$, $v\in T$.
\item Let $A$, $P$, $Z$, and $Y_v$ and $Q_v$ for $v$ in a finite set $T$, be
  complex vector spaces, let $p\colon A\to P$ and $q_v\colon Y_v\to Q_v$ be
  linear, let $A_0\subset A$ and $S_v\subset Y_v$ be subspaces, and let
  $\Lambda\colon A\times\prod_{v\in T}Y_v\to Z$ be multilinear, with
  $\Lambda(f,\psi)=0$ whenever $f\in A_0$, $\psi_v\in S_v$ for every $v$, and
  $p(f)=0$ or $q_v(\psi_v)=0$ for some $v$. Then there is a linear map
  $L\colon P\otimes\bigotimes_{v\in T}Q_v\to Z$ with
  $L\bigl(p(f)\otimes\bigotimes_vq_v(\psi_v)\bigr)=\Lambda(f,\psi)$ for all
  $f\in A_0$ and $\psi$ with $\psi_v\in S_v$.
\end{enumerate}
\end{lemma}
\leanref{Development/WeilBound/ThetaLifts/CartanFactorisation/ExistsPureML.lean}{PositivityTheory.ThKF.tens_eq}
\leanref{Development/WeilBound/ThetaLifts/CartanFactorisation/ExistsPureML.lean}{PositivityTheory.ThKF.tens_eq_map}
\leanref{Development/WeilBound/ThetaLifts/CartanFactorisation/ExistsPureML.lean}{PositivityTheory.ThKF.desc}
\begin{proof}
(1) Choose $z_v\in F_v^\kappa$ with $w_v(z_v)\neq0$ for $v\notin T$, with
$z_v=0$ where $w_v$ is the indicator function of $\mathcal O_v^\kappa$. At the
points of $\mathbf A_F^\kappa$ with components $z_v$ outside $T\cup\{\infty\}$
both sides equal $\prod_{v\notin T}w_v(z_v)\neq0$ times the functions
$\sum_if_i\otimes\bigotimes_{v\in T}\phi_{i,v}$ and
$\sum_jf'_j\otimes\bigotimes_{v\in T}\phi'_{j,v}$ on
$F_\infty^\kappa\times\prod_{v\in T}F_v^\kappa$. So these functions agree. The
map from a tensor product of spaces of functions on sets to the functions on
the product of the sets is injective, which gives the claim.

(2) Choose linear maps $s_0\colon P\to A_0$ and $s_v\colon Q_v\to S_v$ with
$p(s_0(p(f)))=p(f)$ for $f\in A_0$ and $q_v(s_v(q_v(\psi)))=q_v(\psi)$ for
$\psi\in S_v$: on $p(A_0)$ choose a section of $p$ through $A_0$, and extend it
by $0$ on a complement. Let $L$ be the linear map with
$L(x\otimes\bigotimes_vy_v)=\Lambda(s_0(x),(s_v(y_v))_v)$. For $f\in A_0$ and
$\psi_v\in S_v$, $s_0(p(f))-f$ lies in $A_0\cap\ker p$ and
$s_v(q_v(\psi_v))-\psi_v$ in $S_v\cap\ker q_v$, so expanding by multilinearity,
$\Lambda(s_0(p(f)),(s_v(q_v(\psi_v)))_v)=\Lambda(f,\psi)$.
\end{proof}
% Lean: (1) is tens_eq, through tensInj, and tens_eq_map is its image under p and the q_v.
% (2) is desc, with exists_section for s_0 and s_v and ml_congr for the expansion.

\begin{lemma}\label{asm-theta-agree-definite-inputs-constancy-hecke}
Let $K$ be a number field, $\chi$ a continuous homomorphism from
$\mathbf A_K^\times$ to $\mathbf C^\times$ that is trivial on $K^\times$ and
has values of absolute value $1$, and $S$ a finite set of finite places of
$K$. Suppose that $\chi(x_w)=1$ for every finite place $w\notin S$ whose
absolute norm $N(w)$ is a prime number and every $x\in K_w^\times$, where
$x_w\in\mathbf A_K^\times$ has component $x$ at $w$ and $1$ elsewhere. Then
$\chi=1$.
\end{lemma}
\leanref{Development/WeilBound/ThetaLifts/LatticeBasics/ConstHecke.lean}{GlobalTheory.ThK.Dbl.hchar_eq_one}
\leanref{ReciprocityTheory/NeutralLocalFields/LocalFieldComparison/ExistsPrimeAbsNormOfHasDirichletDensity.lean}{ReciprocityTheory.Neutral.Bridge.exists_bound_primeIdealZetaSum_not_prime}
\leanref{ReciprocityTheory/NeutralLocalFields/LocalFieldComparison/TendstoPrimeIdealZetaSumUniv.lean}{ReciprocityTheory.Neutral.Bridge.tendsto_primeIdealZetaSum_univ}
\begin{proof}
Suppose that $\chi\neq1$. For a finite place $w$ and real $s>1$ let
$e_w(s)=(1-\chi(\varpi_w)N(w)^{-s})^{-1}$ if $\chi$ is unramified at $w$, where
$\varpi_w\in\mathbf A_K^\times$ has a uniformizer of $K_w$ as its component at
$w$ and $1$ elsewhere, and $e_w(s)=1$ otherwise. By Hecke's theorem
\cite{Hec20,Tat67}, there is a function $\Lambda$, holomorphic in a
neighbourhood of $1$, with $\Lambda(1)\neq0$ and $\Lambda(s)=\prod_we_w(s)$
for real $s>1$, the product converging. The proof below does not use
$\Lambda(1)\neq0$.

Let $s>1$ be real, and put $x_w=N(w)^{-s}$, so that $0<x_w<1$. Let $A$ be the
set of finite places whose absolute norm is not a prime number. There is $C$,
independent of $s$, with $\sum_{w\in A}x_w\le C$: if $w\in A$ lies above the
prime $p$, then $p^2\le N(w)$, so $x_w\le p^{-2}$, and at most $[K:\mathbf Q]$
places lie above $p$; so $C=[K:\mathbf Q]\sum_{m\ge1}m^{-2}$ will do.

For $w\in A\cup S$, $|e_w(s)|\ge e^{-x_w}$: if $\chi$ is unramified at $w$,
then $|1-\chi(\varpi_w)x_w|\le1+x_w\le e^{x_w}$, as $|\chi|=1$. For $w\notin A\cup S$,
$\chi$ is trivial on $K_w^\times$, so it is unramified at $w$ with
$\chi(\varpi_w)=1$, and $e_w(s)=(1-x_w)^{-1}\ge e^{x_w}$. So for every $w$,
$|e_w(s)|\ge\exp\bigl(x_w-2x_w[w\in A]-2x_w[w\in S]\bigr)$, where $[\,\cdot\,]$
is $1$ if the condition holds and $0$ otherwise. Taking products over finite
sets of places and passing to the limit,
\[
|\Lambda(s)|\ge\exp\Bigl(\sum_wN(w)^{-s}-2C-2|S|\Bigr),
\]
as $x_w\le1$ for $w\in S$. The Dedekind zeta function
$\zeta_K(s)=\prod_w(1-N(w)^{-s})^{-1}$ tends to $\infty$ as $s\to1^+$
\cite{Lan94}, and $\log\zeta_K(s)\le2\sum_wN(w)^{-s}$, as
$-\log(1-x)\le2x$ for $0\le x\le\tfrac12$ and $N(w)\ge2$. So
$\sum_wN(w)^{-s}\to\infty$ and $|\Lambda(s)|\to\infty$ as $s\to1^+$, which
contradicts the continuity of $\Lambda$ at $1$.
\end{proof}
% Lean: hchar_eq_one, from (H K chi hu).2 (Published.HeckeLEntireAndNonzeroAtOne, whose
% IsContAtOne gives AnalyticAt Lambda 1 and the Euler product for Re s > 1; Lambda 1 /= 0 is
% not used). The bound for A is exists_bound_primeIdealZetaSum_not_prime
% (rpow_neg_le_inv_ratPrimeBelow_sq, card_fiber_ratPrimeBelow_le); the factors are
% exp_neg_le_norm_inv and exp_le_inv; the divergence is tendsto_primeIdealZetaSum_univ.

\begin{lemma}\label{asm-theta-agree-definite-inputs-constancy-ideles}
Let $K$ be a totally complex number field, let $t_K\in\mathbf A_K^\times$ have
components $\exp(1)$ at the infinite places and $1$ at the finite places, and
let $S$ be a finite set of finite places of $K$. Let $M$ be a closed subgroup
of $\mathbf A_K^\times$ that contains $K^\times$ and $t_K$, and that contains,
for every finite place $w\notin S$ whose absolute norm is a prime number, every
element of $\mathbf A_K^\times$ whose components away from $w$ are $1$. Then
$M=\mathbf A_K^\times$.
\end{lemma}
\leanref{Development/WeilBound/ThetaLifts/LatticeBasics/ConstHecke.lean}{GlobalTheory.ThK.Dbl.closed_eq_top}
\leanref{Development/AbelianVarietyTheory/Proofs/InertiaEigenvalues/ExistsCompactLatticeGL.lean}{GlobalTheory.IdeleCovolume.exists_compact_idele}
\begin{proof}
There is a compact set $C\subset\mathbf A_K^\times$ such that for every
$y\in\mathbf A_K^\times$ there are $k\in K^\times$ and $m\in\mathbf Z$ with
$k\,t_K^my\in C$; this is the compactness of the group of idele classes of
idelic norm $1$ \cite{Wei67}. As $M$ contains $K^\times$ and $t_K$, the class
of $y$ in $\mathbf A_K^\times/M$ is that of $k\,t_K^my$. So
$\mathbf A_K^\times/M$ is the image of $C$, and it is a compact abelian group,
Hausdorff as $M$ is closed. Suppose that some $z\in\mathbf A_K^\times$ is not in
$M$. As the characters of a compact abelian group separate its points
\cite{Fol95}, there is a continuous homomorphism $\chi$ from
$\mathbf A_K^\times/M$ to the unit circle with $\chi(zM)\neq1$. Its composite
with the quotient map is a continuous homomorphism on $\mathbf A_K^\times$ with
values of absolute value $1$ that is trivial on $M$, hence on $K^\times$ and on
the elements of $\mathbf A_K^\times$ whose components away from $w$ are $1$, for
$w\notin S$ of prime absolute norm. By
Lemma~\ref{asm-theta-agree-definite-inputs-constancy-hecke} it is trivial,
against $\chi(zM)\neq1$.
\end{proof}
% Lean: closed_eq_top, with the compact set of exists_compact_idele (eIdele is t_K) and the
% characters of CompactAbelianFourier.exists_char_ne_one; the Hecke character is X, and
% hchar_eq_one gives X = 1.

\begin{lemma}\label{asm-theta-agree-definite-inputs-constancy-norm}
Let $E$ be a CM field, Galois over $\mathbf Q$, with maximal totally real
subfield $F$ and complex conjugation $c$, put
$\mathbf A_E=\mathbf A_F\otimes_FE$, and let $\mathbf A_E^1$ be the group of
$z\in\mathbf A_E^\times$ with $z\,c(z)=1$. Let $T$ be a finite set of finite
places of $F$. Let $N$ be a closed subgroup of $\mathbf A_E^1$ that contains
$E^\times\cap\mathbf A_E^1$ and that contains, for every finite place
$v\notin T$ of $F$ such that the rational prime below $v$ splits completely in
$E$, every $z\in\mathbf A_E^1$ whose components away from $v$ are $1$. Then
$N=\mathbf A_E^1$.
\end{lemma}
\leanref{Development/WeilBound/ThetaLifts/LatticeBasics/ConstIdele.lean}{GlobalTheory.ThK.Dbl.normOne_mem}
\leanref{Development/WeilBound/ThetaLifts/LatticeBasics/ConstIdele.lean}{GlobalTheory.ThK.Dbl.exists_h90_eq_of_normOne}
\leanref{Development/AbelianVarietyTheory/Proofs/UnitaryExistence/Hilbert90AdeleRingTensorQuadratic.lean}{GlobalTheory.AdLat.hilbert90_adeleRing_tensor_quadratic}
\leanref{Development/WeilBound/ThetaLifts/LatticeBasics/ConstIdele.lean}{GlobalTheory.ThK.Dbl.splitsFully_of}
\begin{proof}
\emph{Hilbert's Theorem 90.} The map $q(z)=z/c(z)$ is a continuous
homomorphism from $\mathbf A_E^\times$ to $\mathbf A_E^1$, as $c$ is a
continuous involution. It is onto. Indeed, let $u\in\mathbf A_E^1$, and put
$s=(u+c(u))/2\in\mathbf A_F$. There is an idempotent $e\in\mathbf A_F$ such that
$e(1+s)+(1-e)(1-s)$ is a unit: at a finite place $v$ take $e_v=1$ if
$|1+s_v|_v\ge|1-s_v|_v$ and $e_v=0$ otherwise, so that the component at $v$ has
absolute value at least $|2|_v$, and at an infinite place take $e_v=1$ if
$1+s_v\neq0$ and $e_v=0$ otherwise. Let $\delta\in E^\times$ with
$c(\delta)=-\delta$, and put $x=e(1+u)+(1-e)\delta(1-u)$. Then $u\,c(x)=x$,
and $x\,c(x)=2\bigl(e(1+s)+(1-e)(1-s)\bigr)\bigl(e-(1-e)\delta^2\bigr)$ is a
unit; so $x\in\mathbf A_E^\times$ and $u=q(x)$.

\emph{The subgroup $M=q^{-1}(N)$.} It is a closed subgroup of
$\mathbf A_E^\times$, which is the idele group of $E$. For $k\in E^\times$,
$q(k)=k/c(k)$ lies in $E^\times\cap\mathbf A_E^1$, so $k\in M$. The element
$t_E$ of Lemma~\ref{asm-theta-agree-definite-inputs-constancy-ideles} has real
components, so $c(t_E)=t_E$, $q(t_E)=1$ and $t_E\in M$. Let $S$ be the finite
set of finite places of $E$ that lie above $T$ or are ramified over
$\mathbf Q$. Let $w\notin S$ be a finite place of $E$ whose absolute norm is a
prime number $p$, let $v$ be the place of $F$ below $w$, and let
$y\in\mathbf A_E^\times$ have components $1$ away from $w$. Then $y-1$ has
components $0$ away from $v$, hence so do $c(y)-1$ and
$q(y)-1=\bigl((y-1)-(c(y)-1)\bigr)c(y)^{-1}$. As $w$ is unramified over
$\mathbf Q$ with residue degree $1$ over $p$ and $E$ is Galois over
$\mathbf Q$, $p$ splits completely in $E$; and $v\notin T$. So $q(y)\in N$ and
$y\in M$.

As $E$ is totally complex,
Lemma~\ref{asm-theta-agree-definite-inputs-constancy-ideles} gives
$M=\mathbf A_E^\times$. So every $u=q(x)$ in $\mathbf A_E^1$ lies in $N$.
\end{proof}
% Lean: normOne_mem, with q = h90 (cA sigma) composed with adeleEquiv, onto by
% exists_h90_eq_of_normOne (hilbert90_adeleRing_tensor_quadratic: the idempotent is
% exists_idempotent_isUnit_one_add_or_sub_adele = AdIdem.exists_idem_isUnit, through local_val
% and local_field, and x is hilbert90_involution_of_idempotent_unit). The three conditions of
% closed_eq_top are k / sigma k, cA_eIdele and at_localIdele, with badSet, finite_badSet and
% splitsFully_of.

\begin{lemma}\label{asm-theta-agree-definite-inputs-constancy-cayley}
Let $E$ be a CM field with maximal totally real subfield $F$ and complex
conjugation $c$, let $v$ be a place of $F$, and put $E_v=F_v\otimes_FE$, on
which $c$ acts through the second factor. Let $J\in GL_n(E)$ with
$c(J)^{\mathsf T}=J$. Then every $g\in M_n(E_v)$ with
$g^{\mathsf T}J\,c(g)=J$ is a limit of matrices $\gamma\in M_n(E)$ with
$\gamma^{\mathsf T}J\,c(\gamma)=J$.
\end{lemma}
\leanref{Development/WeilBound/CoreStatements/AbelianHeightsDuality/UnitaryDensity/Cayley.lean}{PositivityTheory.SA.dense_unitary}
\leanref{Development/WeilBound/CoreStatements/AbelianHeightsDuality/UnitaryDensity/Cayley.lean}{PositivityTheory.SA.units_fd}
\leanref{Development/AbelianVarietyTheory/Proofs/AdelicLatticeProofs/CayleyTransformUnitarySkewInvolutive.lean}{GlobalTheory.CayleyInv.cayley_transform}
\leanref{Development/AbelianVarietyTheory/Proofs/FoundationProofs/ExistsUnitaryClassRepTrivialAtCm.lean}{GlobalTheory.NormSum.RepAt.exists_unit_det}
\begin{proof}
The ring $E_v$ is a finite-dimensional algebra over the field $F_v$. So it has
finitely many maximal ideals, its units form an open set on which inversion is
continuous, and $E$ is dense in it, as $F$ is dense in $F_v$.

\emph{Cayley transforms.} Let $R$ be a commutative ring in which $2$ is a
unit, with an involution $\sigma$, and $G\in GL_n(R)$ with
$\sigma(G)^{\mathsf T}=G$. Call $S\in M_n(R)$ \emph{skew} if
$S^{\mathsf T}G+G\sigma(S)=0$. The following are direct computations. If
$h^{\mathsf T}G\sigma(h)=G$ and $1+h$ is invertible, then
$S=(1-h)(1+h)^{-1}$ is skew, $1+S$ is invertible and $(1-S)(1+S)^{-1}=h$. If
$S$ is skew and $1+S$ is invertible, then $h=(1-S)(1+S)^{-1}$ has
$h^{\mathsf T}G\sigma(h)=G$. For every $A\in M_n(R)$,
$P(A)=\tfrac12\bigl(A-(G\sigma(A)G^{-1})^{\mathsf T}\bigr)$ is skew.

\emph{A unit.} The group of $\lambda\in E^\times$ with $\lambda c(\lambda)=1$
is infinite, as $c\neq1$. For each maximal ideal $\mathfrak m$ of $E_v$, the
image of $\det(1+\lambda g)$ in the field $E_v/\mathfrak m$ is the value at
$\lambda$ of a nonzero polynomial, the reversed characteristic polynomial
$\det(1+Xg)$ of $-g$ modulo $\mathfrak m$, so it vanishes for finitely many
$\lambda$. So there is $\lambda\in E$ with $\lambda c(\lambda)=1$ and
$\det(1+\lambda g)$ a unit.

\emph{Approximation.} Apply the computations over $E_v$, with $G=J$ and
$\sigma=c$, to $h=\lambda g$, which has $h^{\mathsf T}J\,c(h)=J$ as
$\lambda c(\lambda)=1$. Then $S=(1-h)(1+h)^{-1}$ is skew, so
$J\,c(S)J^{-1}=-S^{\mathsf T}$ and $P(S)=S$. The map
$A\mapsto\lambda^{-1}(1-P(A))(1+P(A))^{-1}$ is defined and continuous near
$S$, as $1+P(S)=1+S$ is invertible, and its value at $S$ is
$\lambda^{-1}h=g$. As $M_n(E)$ is dense in $M_n(E_v)$, there are
$A\in M_n(E)$ arbitrarily close to $S$ with $1+P(A)$ invertible. For such $A$,
$T=P(A)$ lies in $M_n(E)$, is skew and has $1+T$ invertible, so
$\gamma=\lambda^{-1}(1-T)(1+T)^{-1}\in M_n(E)$ has
$\gamma^{\mathsf T}J\,c(\gamma)=J$, as $\lambda^{-1}c(\lambda^{-1})=1$. These
$\gamma$ tend to $g$.
\end{proof}
% Lean: dense_unitary, applied in cut_place to R = k (x)_F E with units_fd for the units;
% exists_unit_det with normOne_infinite for lambda; cayley_transform (its four clauses c1, c2,
% d3, d4) for the computations; P, PS, hf and hT as in the text.

\begin{lemma}\label{asm-theta-agree-definite-inputs-constancy-lines}
Let $E$ be a CM field with maximal totally real subfield $F$ and complex
conjugation $c$, and let $\mathbf A_E$ and $\mathbf A_E^1$ be as in
Lemma~\ref{asm-theta-agree-definite-inputs-constancy-norm}. Let
$J\in GL_n(E)$ with $c(J)^{\mathsf T}=J$, let $U$ be the group of
$g\in GL_n(\mathbf A_E)$ with $g^{\mathsf T}J\,c(g)=J$, and $U(F)$ its
subgroup of matrices over $E$.
\begin{enumerate}
\item Let $R\in M_n(E)$ with $R^2=R$ and $R^{\mathsf T}J=J\,c(R)$. Then
  $z\mapsto1+(z-1)R$ is a continuous homomorphism from $\mathbf A_E^1$ to $U$,
  which maps $E^\times\cap\mathbf A_E^1$ into $U(F)$.
\item Let $L$ be a closed subgroup of $U$ that contains $1+(z-1)R$ for every
  $z\in\mathbf A_E^1$ and every $R$ as in (1). Then $L=U$.
\end{enumerate}
\end{lemma}
\leanref{Development/WeilBound/CoreStatements/AbelianHeightsDuality/UnitaryDensity/EqTopOfLines.lean}{PositivityTheory.SA.exists_line}
\leanref{Development/WeilBound/CoreStatements/AbelianHeightsDuality/UnitaryDensity/EqTopOfLines.lean}{PositivityTheory.SA.eq_top_of_lines}
\leanref{Development/WeilBound/CoreStatements/AbelianHeightsDuality/UnitaryDensity/EqTopOfLines.lean}{PositivityTheory.SA.cut_rat}
\leanref{Development/WeilBound/CoreStatements/AbelianHeightsDuality/UnitaryDensity/EqTopOfLines.lean}{PositivityTheory.SA.cut_place}
\leanref{Development/WeilBound/CoreStatements/AbelianHeightsDuality/UnitaryDensity/EqTopOfLines.lean}{PositivityTheory.SA.isInducing_mat}
\begin{proof}
(1) Put $S(z)=1+(z-1)R$. As $R^2=R$,
$S(z)S(z')=1+(z-1+z'-1+(z-1)(z'-1))R=S(zz')$. As $R^{\mathsf T}J=J\,c(R)$ and
$R^{\mathsf T}J\,c(R)=J\,c(R^2)=J\,c(R)$,
\[
S(z)^{\mathsf T}J\,c(S(z))=J+\bigl((z-1)+(c(z)-1)+(z-1)(c(z)-1)\bigr)J\,c(R)=J+(z\,c(z)-1)J\,c(R),
\]
which is $J$ for $z\in\mathbf A_E^1$. So $S$ is a homomorphism to $U$, and it is
continuous. For $z\in E^\times$, $S(z)$ is a matrix over $E$.

(2) For $a\in\mathbf A_F$ and $M\in M_n(\mathbf A_E)$ put $M_a=1+a(M-1)$. If
$a^2=a$, then $M_aN_a=(MN)_a$; and if $ab=0$, then $M_aM_b=M_{a+b}$.

\emph{Step 1: rational elements.} Let $a\in\mathbf A_F$ with $a^2=a$, and
$\gamma\in U(F)$. Then $\gamma_a\in L$. Indeed, put $h(x,y)=x^{\mathsf T}J\,c(y)$
on $E^n$; then $h(\gamma x,\gamma y)=h(x,y)$. By \cite{Die71}, $\gamma$ is a
product of quasi-reflections $r=1+(\lambda-1)R$ with
$R=h(p,p)^{-1}p\,(J\,c(p))^{\mathsf T}$, where $p\in E^n$ has $h(p,p)\neq0$ and
$\lambda\in E$ has $\lambda c(\lambda)=1$. Such an $R$ has $R^2=R$, as
$(J\,c(p))^{\mathsf T}p=h(p,p)$, and $R^{\mathsf T}J=J\,c(R)$, as
$c(J)^{\mathsf T}=J$ and $c(h(p,p))=h(p,p)$. The element $z=1+a(\lambda-1)$ has
$z\,c(z)=1+a(\lambda c(\lambda)-1)=1$, so $z\in\mathbf A_E^1$, and
$r_a=1+(z-1)R$ lies in $L$. So $\gamma_a$, the product of the elements $r_a$,
lies in $L$.

\emph{Step 2: one place.} Let $C$ be the closure of $L$ in
$M_n(\mathbf A_E)$. Let $v$ be a place of $F$, $e_v\in\mathbf A_F$ the element
with component $1$ at $v$ and $0$ elsewhere, $p\colon\mathbf A_E\to E_v=F_v\otimes_FE$
the projection and $q\colon E_v\to\mathbf A_E$ the inclusion at $v$; they are
continuous, commute with $c$, and $q(p(z))=e_vz$. Let $g\in U$. Then
$p(g)^{\mathsf T}J\,c(p(g))=J$, so by
Lemma~\ref{asm-theta-agree-definite-inputs-constancy-cayley}, $p(g)$ is a limit
of matrices $\gamma\in U(F)$. The map $\Phi(M)=1+q(M-1)$ from $M_n(E_v)$ to
$M_n(\mathbf A_E)$ is continuous, and $\Phi(p(M))=M_{e_v}$ for
$M\in M_n(\mathbf A_E)$. So $g_{e_v}=\Phi(p(g))$ is a limit of the matrices
$\Phi(\gamma)=\gamma_{e_v}$, which lie in $L$ by Step 1, and $g_{e_v}\in C$.

\emph{Step 3: all places.} The set $C$ is closed under multiplication, as $L$
is. For a finite set $\Sigma$ of finite places put
$\eta_\Sigma=\sum_we_w+\sum_{v\in\Sigma}e_v$, the first sum over the infinite
places. These $e_w$ and $e_v$ are orthogonal idempotents, so $g_{\eta_\Sigma}$
is the product of the $g_{e_w}$ and the $g_{e_v}$, and lies in $C$. As
$\Sigma$ grows, $\eta_\Sigma x\to x$ for every $x\in\mathbf A_F$, hence for
every $x\in\mathbf A_E$; so $g_{\eta_\Sigma}=1+\eta_\Sigma(g-1)\to g$, and
$g\in C$. On $U$,
$g^{-1}=(J\,c(g)J^{-1})^{\mathsf T}$ is a continuous function of $g$, so the
topology of $U$ is induced from $M_n(\mathbf A_E)$. Hence $g$ lies in the
closure of $L$ in $U$, which is $L$.
\end{proof}
% Lean: (1) is exists_line (sec, with the unitarity computed as in the text). (2) is
% eq_top_of_lines: M_a is cutM a M, with cutM_mul and cutM_add; Step 1 is cut_rat, through
% T.unitaryMatrix_eq_prod_quasiReflections (cited to Dieudonne as in Lemma
% is-eigenvalue-conjugate-split-correction) and the induction on the list; Step 2 is cut_place
% with p = map pi id, q = map j id (qp), Phi and hPhi, and dense_unitary for key; Step 3 is
% hsum and h1 (etaS, jI, jF, with jI_ne and jF_ne for orthogonality), h2 and h3
% (tendsto_etaS), and isInducing_mat.

\begin{lemma}\label{asm-theta-agree-definite-inputs-pairs-density}
Let $A$ be a complex vector space and $T$ a finite set, and for $v\in T$ let
$X_v$ be a nonzero complex vector space of countable dimension and $O_v$ a set
of linear operators on $X_v$ such that $0$ and $X_v$ are the only subspaces of
$X_v$ stable under every element of $O_v$. For $v\in T$ and an operator
$\mathrm{Op}$ on $X_v$, let $\mathrm{Op}^{(v)}$ be the operator on
$A\otimes\bigotimes_{u\in T}X_u$ that is $\mathrm{Op}$ on the factor $X_v$ and
the identity on the other factors. Let $S\neq0$ be a subspace of
$A\otimes\bigotimes_{u\in T}X_u$ with $\mathrm{Op}^{(v)}S\subset S$ for every
$v\in T$ and $\mathrm{Op}\in O_v$. Then there is $a_0\neq0$ in $A$ with
$a_0\otimes y\in S$ for every $y\in\bigotimes_{u\in T}X_u$.
\end{lemma}
\leanref{Development/WeilBound/ThetaLifts/LatticeBasics/PairsD.lean}{GlobalTheory.ThK.Dbl.Dens.exists_tmul_mem}
\leanref{Development/WeilBound/ThetaLifts/LatticeBasics/PairsD.lean}{GlobalTheory.ThK.Dbl.Dens.exists_mem_adjoin_eq}
\leanref{Development/WeilBound/ThetaLifts/LatticeBasics/PairsD.lean}{GlobalTheory.ThK.Dbl.Dens.exists_smul_of_rank_le_aleph0}
\begin{proof}
For $v\in T$ let $\mathcal A_v$ be the set of operators $\mathrm{Op}$ on $X_v$
with $\mathrm{Op}^{(v)}S\subset S$. It is a subalgebra containing $O_v$, as
$(\mathrm{Op}\,\mathrm{Op}')^{(v)}=\mathrm{Op}^{(v)}\mathrm{Op}'^{(v)}$. By
Lemma~\ref{th-first-lift-dixmier}, every operator on $X_v$ that commutes with
the elements of $O_v$ is a scalar. So by the Jacobson density theorem
\cite[\S11]{Lam01}, for every operator $f$ on $X_v$ and every finite set of
vectors of $X_v$, some element of the algebra generated by $O_v$, and so of
$\mathcal A_v$, agrees with $f$ on these vectors. If $\mathrm{Op}_v\in\mathcal A_v$
for every $v$, then the operator $\bigotimes_v\mathrm{Op}_v$, the identity on
$A$, is the composite of the $\mathrm{Op}_v^{(v)}$ and maps $S$ into $S$.

Choose bases $(b_{v,j})_j$ of the $X_v$ and $x\neq0$ in $S$, and write
$x=\sum_qa_q\otimes\bigotimes_vb_{v,q_v}$, a finite sum over families
$q=(q_v)_{v\in T}$, with $a_p\neq0$ for some $p$. Let $y_v\in X_v$ for $v\in T$,
and let $f_v$ be the operator on $X_v$ with $f_v(b_{v,p_v})=y_v$ and
$f_v(b_{v,j})=0$ for $j\neq p_v$. Choose $\mathrm{Op}_v\in\mathcal A_v$ that agrees
with $f_v$ on the finitely many $b_{v,q_v}$ occurring in $x$. Then
$(\bigotimes_v\mathrm{Op}_v)(x)=(\bigotimes_vf_v)(x)=a_p\otimes\bigotimes_vy_v$
lies in $S$. As the pure tensors span $\bigotimes_vX_v$, $a_0=a_p$ will do.
\end{proof}
% Lean: exists_tmul_mem (PairsD.lean): Asub is A_v, h3 the composite, exists_mem_adjoin_eq the
% density (jacobson_density, with exists_smul_of_rank_le_aleph0 for the scalar endomorphisms),
% and hpure the pure tensors, through the basis Basis.piTensorProduct and
% equivFinsuppOfBasisRight.

\begin{lemma}\label{asm-theta-agree-definite-inputs-pairs-finite}
Let $(X,W,M,\chi_X,\chi_W,\omega)$ be a Weil pair for a theta base $B$
(Definition~\ref{asm-theta-agree-triple}(1)), with $X$ and $W$ of ranks $n$ and
$d$, put $I=\{1,\ldots,n\}\times\{1,\ldots,d\}$, and let $K_\infty(W)$ be as in
Definition~\ref{asm-theta-agree-matched}(1). Then for every
$f\in\operatorname{Pol}(M)$ there is a finite-dimensional subspace
$U\subset\mathcal S(F_\infty^I)$ that contains $f_1$ whenever
$k\in K_\infty(W)$ and $\omega(1,k)$ acts at $\infty$ by $f\mapsto f_1$.
\end{lemma}
\leanref{Development/WeilBound/ThetaLifts/LatticeBasics/PolFin.lean}{GlobalTheory.ThK.Dbl.polFin}
\leanref{Development/WeilBound/ThetaLifts/LatticeBasics/PairsIrrOfArch.lean}{GlobalTheory.ThK.Dbl.polFinTri}
\leanref{Development/WeilBound/ThetaLifts/SchwartzMoments/PolynomialFunctionalEvalContinuous.lean}{PositivityTheory.ThKOwn.KPol.gauss_impl}
\leanref{Development/WeilBound/ThetaLifts/SchwartzMoments/PolynomialFunctionalEvalContinuous.lean}{PositivityTheory.ThKOwn.KPol.inter}
\leanref{Development/WeilBound/ThetaLifts/SchwartzMoments/PolynomialFunctionalEvalContinuous.lean}{PositivityTheory.ThKOwn.KPol.acts_unique}
\begin{proof}
Let $D_u$, for $u\in\mathbb W(X,W)$ at $\infty$, be the operators of
Lemma~\ref{asm-theta-agree-rechoice-intertwine}, and for $k\in K_\infty(W)$ let
$T_k$ be an archimedean operator of $\omega$ at $(1,k)$
(Lemma~\ref{asm-theta-agree-rechoice-operators}(1)). If $\omega(1,k)$ acts at
$\infty$ by $f\mapsto f_1$, then $f_1=T_k(f)$: for a Schwartz-Bruhat $g\neq0$
on $\mathbf A_f^I$, $\omega(1,k)(f\otimes g)$ is both $f_1\otimes g$ and
$T_k(f)\otimes g$. Let $\mathcal U$ be the set of $f\in\mathcal S(F_\infty^I)$
for which some finite-dimensional subspace contains $T_k(f)$ for every
$k\in K_\infty(W)$; it is a subspace.

\emph{Gaussians.} Let $f_0$ be a Gaussian with signs $\varepsilon$
(Definition~\ref{asm-theta-agree-models}(4)), and $k\in K_\infty(W)$. Then
$T_k(f_0)$ is a Gaussian with signs $\varepsilon$. Indeed, let $u$ and $u'$ be
at $\infty$ with
$\pi_w(u')=\varepsilon_wi\,\theta_{X,w}\pi_w(u)\theta_{W,w}^{\mathsf T}$ for
every $w$, and put $m=(k^{-1})^{\mathsf T}$. As $\pi_w(k^{-1})$ commutes with
$\theta_{W,w}$, the elements $um$ and $u'm$ satisfy the same relation, so
$D_{u'm}f_0=i\,D_{um}f_0$, by
Lemma~\ref{asm-theta-agree-rechoice-intertwine}(1). By part (2) of that lemma,
as $umk^{\mathsf T}=u$, $T_k(D_{um}f_0)=D_uT_k(f_0)$, and likewise for $u'$. So
$D_{u'}T_k(f_0)=i\,D_uT_k(f_0)$. Hence $T_k(f_0)$ lies in the space
$G_\varepsilon$ of the proof of
Lemma~\ref{asm-theta-agree-definite-inputs}(1), which has dimension at most
$1$, and $f_0\in\mathcal U$.

\emph{Derivatives.} Let $f\in\mathcal U$, with a finite-dimensional space $U$
containing every $T_k(f)$, and let $u$ be at $\infty$. As $D_u$ is
$\mathbf R$-linear in $u$, the operators $D_u$ are combinations of finitely
many $D_{u_1},\ldots,D_{u_m}$, and $U'=\sum_jD_{u_j}U$ is finite-dimensional.
By Lemma~\ref{asm-theta-agree-rechoice-intertwine}(2),
$T_k(D_uf)=D_{uk^{\mathsf T}}T_k(f)\in U'$. So $D_uf\in\mathcal U$.

By Lemma~\ref{asm-theta-agree-rechoice-intertwine}(1),
$\operatorname{Pol}(M)$ is the smallest subspace of $\mathcal S(F_\infty^I)$
that contains the Gaussians and contains $D_uf$ whenever it contains $f$ and
$u$ is at $\infty$. So $\operatorname{Pol}(M)\subset\mathcal U$.
\end{proof}
% Lean: polFin (PolFin.lean), with S the set of f with a finite-dimensional U containing A f
% for every k in KInf and every A with KPol.Impl om k A; the Gaussians are gauss_impl and
% Dense.vacSp_fin (the space G_eps), the derivatives Dense.exists_dop_fin and KPol.inter, and
% the conclusion KPol.exists_impl, exists_arch_operator and acts_unique.

\begin{lemma}\label{asm-theta-agree-definite-inputs-pairs-operators}
Let $(B,G,\mathcal X,\omega,\omega',R,R_f)$ be a matched quadruple of ranks
$(a,a')$ (Definition~\ref{asm-theta-agree-matched}), with spaces $V$, $V'$ and
$W$, models $M$ and $M'$, real maps $R=(\Pi,p,p')$ and finite maps
$R_f=(\Pi_v,p_v,p'_v)_v$, and let $I_0$ and $I'_0$ be as there.
\begin{enumerate}
\item Let $v$ be a finite place and $h\in U_W^{(v)}$. There is an operator
  $\mathrm{Op}$ on $\Pi_v$ such that: for every Schwartz-Bruhat $\phi$ on
  $F_v^{I_0}$ there is a Schwartz-Bruhat $\phi_1$ with $\omega(1,h)$ acting at
  $v$ by $\phi\mapsto\phi_1$; and $p_v(\phi_1)=\mathrm{Op}(p_v(\phi))$ whenever
  $\phi$ and $\phi_1$ are Schwartz-Bruhat and $\omega(1,h)$ acts at $v$ by
  $\phi\mapsto\phi_1$. The same holds for $\omega'$ and $p'_v$, with the same
  $\mathrm{Op}$.
\item Let $k\in K_\infty(W)$. There is an operator $\mathrm{Op}$ on $\Pi$ such
  that: for every $f\in\operatorname{Pol}(M)$ there is
  $f_1\in\operatorname{Pol}(M)$ with $\omega(1,k)$ acting at $\infty$ by
  $f\mapsto f_1$; and $p(f_1)=\mathrm{Op}(p(f))$ whenever
  $f\in\operatorname{Pol}(M)$ and $\omega(1,k)$ acts at $\infty$ by
  $f\mapsto f_1$. The same holds for $\omega'$, $\operatorname{Pol}(M')$ and
  $p'$, with the same $\mathrm{Op}$.
\end{enumerate}
\end{lemma}
\leanref{Development/WeilBound/ThetaLifts/LatticeBasics/PairsL.lean}{GlobalTheory.ThK.Dbl.Prs.fin_ops}
\leanref{Development/WeilBound/ThetaLifts/LatticeBasics/PairsL.lean}{GlobalTheory.ThK.Dbl.Prs.inf_ops}
\leanref{Development/WeilBound/ThetaLifts/LatticeBasics/PairsL.lean}{GlobalTheory.ThK.Dbl.Prs.common_op}
\leanref{Development/WeilBound/ThetaLifts/LatticeBasics/PairsL.lean}{GlobalTheory.ThK.Dbl.Prs.loc_op}
\leanref{Development/WeilBound/ThetaLifts/LatticeBasics/PairsL.lean}{GlobalTheory.ThK.Dbl.Prs.inf_op}
\begin{proof}
(1) By Lemma~\ref{asm-theta-agree-finite-inputs}(2), the Weil pairs of $M$ and
$M'$ have local operators at $v$
(Definition~\ref{asm-theta-agree-finite-local-objects}(4)). Let $A$ and $A'$ be
their operators at $(1,h)$: they map Schwartz-Bruhat functions to
Schwartz-Bruhat functions, and $\omega(1,h)$ acts at $v$ by
$\phi\mapsto A\phi$ and $\omega'(1,h)$ by $\phi'\mapsto A'\phi'$. By the
condition on $h$ in Definition~\ref{asm-theta-agree-matched}(3),
$p_v(A\phi)=p'_v(A'\phi')$ whenever $p_v(\phi)=p'_v(\phi')$; so, as in the proof
of Lemma~\ref{asm-theta-agree-definite-inputs}(6), there is an operator
$\mathrm{Op}$ on $\Pi_v$ with $\mathrm{Op}(p_v(\phi))=p_v(A\phi)$ and
$\mathrm{Op}(p'_v(\phi'))=p'_v(A'\phi')$ for all Schwartz-Bruhat $\phi$ and
$\phi'$. Let $\omega(1,h)$ act at $v$ by $\phi\mapsto\phi_1$, with $\phi$ and
$\phi_1$ Schwartz-Bruhat, and choose $\phi'$ with $p'_v(\phi')=p_v(\phi)$,
which exists as $p'_v$ is onto. By the same condition,
$p_v(\phi_1)=p'_v(A'\phi')=\mathrm{Op}(p'_v(\phi'))=\mathrm{Op}(p_v(\phi))$. The
argument for $\omega'$ is the same.

(2) By Lemmas~\ref{asm-theta-agree-rechoice-operators}(1)
and~\ref{asm-theta-agree-rechoice-fock}, there are archimedean operators $T$ of
$\omega$ and $T'$ of $\omega'$ at $(1,k)$ that map $\operatorname{Pol}(M)$ and
$\operatorname{Pol}(M')$ into themselves. As $p(\operatorname{Pol}(M))=\Pi$ and
$p'(\operatorname{Pol}(M'))=\Pi$, condition (b) of
Definition~\ref{asm-theta-agree-matched}(2) gives, as in (1), an operator
$\mathrm{Op}$ on $\Pi$ with $\mathrm{Op}(p(f))=p(T(f))$ and
$\mathrm{Op}(p'(f'))=p'(T'(f'))$, and $p(f_1)=\mathrm{Op}(p(f))$ whenever
$f\in\operatorname{Pol}(M)$ and $\omega(1,k)$ acts at $\infty$ by $f\mapsto f_1$;
and likewise for $\omega'$.
\end{proof}
% Lean: fin_ops and inf_ops (PairsL.lean), through common_op (the operator p o A o s for a
% section s of p, exists_rightInverse_of_surjective), loc_op (LocRepEx, Lemma
% asm-theta-agree-finite-inputs(2)) and inf_op (ArchOp, with the clause for KInf); the
% second clauses use hm.fin.equiv with hm.fin.onto' and hm.real.cpt with hm.real.quot'.onto.

\begin{lemma}\label{asm-theta-agree-definite-inputs-pairs-fill}
Let $(B,G,\mathcal X,\omega,\omega',R,R_f)$ be a matched quadruple of ranks
$(a,a')$ with $p=q=0$, with spaces $V$, $V'$ and $W$, models $M$ and $M'$, set
$T$, real maps $(\Pi,p,p')$ and finite maps $(\Pi_v,p_v,p'_v)_v$
(Definition~\ref{asm-theta-agree-matched}). For a finite set $T_1\supset T$ of
finite places put $\Pi_{T_1}=\Pi\otimes\bigotimes_{v\in T_1}\Pi_v$, and let
$C_{T_1}$ be the span of the pairs
\[
\Bigl(f\otimes\bigotimes_v\phi_v,\ p(f)\otimes\bigotimes_{v\in T_1}p_v(\phi_v)\Bigr),
\]
in the product of the space of functions on $\mathbf A^{I_0}$ and $\Pi_{T_1}$,
with $f\in\operatorname{Pol}(M)$, every $\phi_v$ Schwartz-Bruhat and
$\phi_v=\mathrm{vac}_v$ for $v\notin T_1$. Let $N_1$ be a space of functions on
$\mathbf A^{I_0}$ with $\omega(1,h)N_1\subset N_1$ for every $h\in U_W$, and
let $(\Phi_0,\xi_0)\in C_{T_1}$ with $\Phi_0\in N_1$ and $\xi_0\neq0$. Then for
every $\xi\in\Pi_{T_1}$ there is $\Phi\in N_1$ with $(\Phi,\xi)\in C_{T_1}$.
\end{lemma}
\leanref{Development/WeilBound/ThetaLifts/LatticeBasics/PairsB.lean}{GlobalTheory.ThK.Dbl.Prs.fill}
\leanref{Development/WeilBound/ThetaLifts/LatticeBasics/PairsB.lean}{GlobalTheory.ThK.Dbl.Prs.irr_inf}
\leanref{Development/WeilBound/ThetaLifts/LatticeBasics/PairsB.lean}{GlobalTheory.ThK.Dbl.Prs.deriv_mem}
\leanref{Development/WeilBound/ThetaLifts/LatticeBasics/PairsB.lean}{GlobalTheory.ThK.Dbl.Prs.rank_pi}
\leanref{Development/WeilBound/ThetaLifts/LatticeBasics/PairsL.lean}{GlobalTheory.ThK.Dbl.Prs.cl_fin}
\leanref{Development/WeilBound/ThetaLifts/LatticeBasics/PairsL.lean}{GlobalTheory.ThK.Dbl.Prs.cl_inf}
\begin{proof}
Let $S$ be the set of $\eta\in\Pi_{T_1}$ with $(\Phi,\eta)\in C_{T_1}$ for some
$\Phi\in N_1$, a subspace containing $\xi_0$.

\emph{The finite places.} For $v\in T_1$ let $O_v$ be the set of the operators
$\mathrm{Op}$ of Lemma~\ref{asm-theta-agree-definite-inputs-pairs-operators}(1),
for all $h\in U_W^{(v)}$, and for $\mathrm{Op}\in O_v$ let $\mathrm{Op}^{(v)}$ be
the operator on $\Pi_{T_1}$ that is $\mathrm{Op}$ on the factor $\Pi_v$. Let
$\mathrm{Op}\in O_v$ belong to $h$. A generator of $C_{T_1}$ has first component
$\phi_v\otimes r$, with $r$ not depending on the $v$-component; so, with
$\phi_1$ Schwartz-Bruhat, $\omega(1,h)$ acting at $v$ by $\phi_v\mapsto\phi_1$
and $p_v(\phi_1)=\mathrm{Op}(p_v(\phi_v))$, $\omega(1,h)$ maps it to the
function with $\phi_1$ in the place of $\phi_v$. Hence
$(\omega(1,h)\Phi,\mathrm{Op}^{(v)}\eta)\in C_{T_1}$ whenever
$(\Phi,\eta)\in C_{T_1}$, and, as $N_1$ is stable under $\omega(1,h)$,
$\mathrm{Op}^{(v)}S\subset S$. The space $\Pi_v$ is nonzero; it has countable
dimension, as $p_v$ maps onto it the Schwartz-Bruhat functions on $F_v^{I_0}$,
which are spanned by the countably many indicator functions of compact open
sets; and by the irreducibility in Definition~\ref{asm-theta-agree-matched}(3)
and the second property of the operators in $O_v$, $0$ and $\Pi_v$ are its only
subspaces stable under $O_v$. By
Lemma~\ref{asm-theta-agree-definite-inputs-pairs-density}, there is $a_0\neq0$
in $\Pi$ with $a_0\otimes y\in S$ for every $y\in\bigotimes_{v\in T_1}\Pi_v$.

\emph{The infinite places.} Let $N_\infty$ be the set of $u\in\Pi$ with
$u\otimes y\in S$ for every $y$, a subspace containing $a_0$. Let $f\in\operatorname{Pol}(M)$ with
$p(f)\in N_\infty$. (a) Let $k\in K_\infty(W)$, and let $\mathrm{Op}$ be as in
Lemma~\ref{asm-theta-agree-definite-inputs-pairs-operators}(2). A generator of
$C_{T_1}$ is $f'\otimes r$ with $r$ not depending on the $\infty$-component, so
$\omega(1,k)$ maps it to $f'_1\otimes r$, with $f'_1\in\operatorname{Pol}(M)$ and
$p(f'_1)=\mathrm{Op}(p(f'))$; as above, $(\mathrm{Op}\otimes1)S\subset S$. So if
$\omega(1,k)$ acts at $\infty$ by $f\mapsto f_1$, then
$p(f_1)=\mathrm{Op}(p(f))\in N_\infty$. (b) Let $\gamma$ be a line in $U_W$ and
let $t\mapsto\omega(1,\gamma(t))$ act at $\infty$ to first order by
$f\mapsto f_1$. As $p=q=0$, $\theta_{W,w}=1$ for every infinite place $w$, as
in the proof of Lemma~\ref{asm-theta-agree-definite-inputs}(8); so
$\gamma(t)\in K_\infty(W)$ for every $t$. Let $f_t\in\operatorname{Pol}(M)$ be such
that $\omega(1,\gamma(t))$ acts at $\infty$ by $f\mapsto f_t$
(Lemma~\ref{asm-theta-agree-definite-inputs-pairs-operators}(2)). By (a) and
Lemma~\ref{asm-theta-agree-definite-inputs-pairs-finite}, every $f_t$ lies in
the finite-dimensional space $U_0=U\cap p^{-1}(N_\infty)$, with $U$ given by
that lemma for $f$. For every finite place $v$ choose $b_v$ with
$\mathrm{vac}_v(b_v)\neq0$, in $\mathcal O_v^{I_0}$ when $\mathrm{vac}_v$ is the
indicator function of $\mathcal O_v^{I_0}$, so that $y=(b_v)_v$ is a finite
adele and $c=\prod_v\mathrm{vac}_v(b_v)\neq0$. The function
$r(z)=\prod_v\mathrm{vac}_v(z_v)$ does not depend on the $\infty$-component, and
$f\otimes r$ is Schwartz-Bruhat, so $\omega(1,\gamma(t))(f\otimes r)=f_t\otimes r$,
and at the point with components $x$ at $\infty$ and $y$ at the finite places,
the action to first order says that $t\mapsto c\,f_t(x)$ has derivative
$c\,f_1(x)$ at $t=0$. So $f_1$ is the pointwise limit of the functions
$(f_t-f_0)/t\in U_0$. A finite-dimensional space of functions is closed under
pointwise limits, so $f_1\in U_0$ and $p(f_1)\in N_\infty$.

By (a), (b) and condition (c) of Definition~\ref{asm-theta-agree-matched}(2),
$N_\infty$ is $0$ or $\Pi$; as $a_0\in N_\infty$, $N_\infty=\Pi$. So $S$ contains
every $u\otimes y$, and $S=\Pi_{T_1}$.
\end{proof}
% Lean: fill (PairsB.lean), with S the image under snd of Cl B G.M R.p Rf.p T meet N x top; the
% finite places are exists_tmul_mem with O_v from fin_ops, cl_fin (PureAct), rank_pi
% (rank_SB_Fv, countable_compactOpen, mem_span_indicator), hm.fin.nontriv and hm.fin.irred;
% the infinite places are Ni = N_infty with inf_ops and cl_inf, and irr_inf: hm.real.irred,
% comm_theta (cartan_one), PolFin (Lemma asm-theta-agree-definite-inputs-pairs-finite),
% Dat.rest with T = empty, the point z from b and y, and deriv_mem.

\begin{lemma}\label{asm-theta-agree-definite-inputs-growth-lattice}
Let $F$ be a number field and $\kappa$ a finite set, and for $x\in F_\infty^\kappa$
let $|x|$ be the largest absolute value of its components at the infinite
places; for $\xi\in F^\kappa$ let $|\xi|$ be that of its image in $F_\infty^\kappa$.
There are integers $N_0,k\geq0$ and $C\geq0$ such that the following holds. Let
$A\in\mathbf R$, let $m\geq1$ be an integer, and let $\Psi$ be a function on
$\mathbf A_F^\kappa$ with $|\Psi(\xi)|(1+|\xi|)^{N_0}\leq A$ for every
$\xi\in F^\kappa$, and with $m\xi\in\mathcal O_F^\kappa$ whenever
$\xi\in F^\kappa$ and $\Psi(\xi)\neq0$. Then $\sum_{\xi\in F^\kappa}\Psi(\xi)$
converges absolutely, and its absolute value is at most $C\,A\,m^k$.
\end{lemma}
\leanref{Development/WeilBound/ThetaLifts/LeviDeterminants/AdelicThetaSumNormLePolynomialScale.lean}{PositivityTheory.ThKF0.theta_sum_norm_bound_denominator}
\leanref{Development/WeilBound/ThetaLifts/LeviDeterminants/AdelicThetaSumNormLePolynomialScale.lean}{PositivityTheory.ThKF0.LatSum.thetaSum_norm_le_pow}
\leanref{Development/WeilBound/ThetaLifts/LeviDeterminants/AdelicThetaSumNormLePolynomialScale.lean}{PositivityTheory.ThKF0.LatSum.summable_inv_one_add_pow}
\begin{proof}
The image $L$ of $\mathcal O_F^\kappa$ in $F_\infty^\kappa$ is a lattice of some
rank $r$. Put $N_0=k=r+1$ and $C=\sum_{z\in L}(1+|z|)^{-N_0}$, which is finite
as $N_0>r$. Let $\xi\in F^\kappa$ with $\Psi(\xi)\neq0$. Then $m\xi\in L$, and as
$m\geq1$, $1+|m\xi|\leq m(1+|\xi|)$, so
\[
|\Psi(\xi)|\leq A(1+|\xi|)^{-N_0}\leq A\,m^{N_0}(1+|m\xi|)^{-N_0}.
\]
As $\xi\mapsto m\xi$ is injective, the sum of the $|\Psi(\xi)|$ is at most
$A\,m^{N_0}C$.
\end{proof}
% Lean: LatSum.thetaSum_norm_le_pow, with L = latPi (the span of latticeBasis), N_0 the rank plus
% one, summable_inv_one_add_pow (ZLattice.summable_norm_rpow) for C, and key for the
% inequality, through tsum_comp_le_tsum_of_inj.

\begin{lemma}\label{asm-theta-agree-definite-inputs-growth-inputs-denominator}
Let $F$ be a number field, and for every finite place $v$ of $F$ let
$|\cdot|_v$ be the normalized absolute value of $F_v$ and $H_v\geq1$ a real
number, with $H_v=1$ for almost all $v$. Then there is an integer $d$ with
$1\leq d\leq\prod_vH_v$ such that $|dy|_v\leq1$ for every finite place $v$ and
every $y\in F_v$ with $|y|_v\leq H_v$.
\end{lemma}
\leanref{Development/AbelianVarietyTheory/Proofs/HeightProofs/UnitaryLatticeDenominatorHeightBoundProof.lean}{GlobalTheory.SmoothForms.denom_global}
\leanref{Development/AbelianVarietyTheory/Proofs/HeightProofs/NormNatCastLe.lean}{GlobalTheory.AdicNormZpowEstimates.exists_denominator}
\begin{proof}
Let $N(v)$ be the absolute norm of $v$, and $n_v\geq0$ the integer with
$N(v)^{n_v}\leq H_v<N(v)^{n_v+1}$; then $n_v=0$ when $H_v=1$. Put
$d=\prod_vN(v)^{n_v}$, so that $1\leq d\leq\prod_vH_v$. Let $|y|_v\leq H_v$. The
nonzero values of $|\cdot|_v$ are the powers of $N(v)$, so
$|y|_v\leq N(v)^{n_v}$. As $N(v)$ lies in the maximal ideal of $\mathcal O_v$,
$|N(v)^{n_v}|_v\leq N(v)^{-n_v}$, and the other factors of $d$ are integers, of
absolute value at most $1$. So $|dy|_v\leq1$.
\end{proof}
% Lean: denom_global = exists_denominator (NormNatCastLe.lean), with n_v from
% exists_nat_pow_near, norm_le_of_lt (through norm_eq_zpow) and norm_absNorm_le.

\begin{lemma}\label{asm-theta-agree-definite-inputs-growth-inputs-entries}
Let $E$ be a finite extension of a number field $F$, and put
$\mathbf A=\mathbf A_F$ and $\mathbf A_E=\mathbf A\otimes_FE$. For
$g\in GL_n(\mathbf A_E)$ let
$\operatorname{ht}(g)=\prod_v\max(1,\|g_v\|,\|g_v^{-1}\|)$, the product over
the places $v$ of $F$, where $\|\cdot\|$ is the largest operator norm under the
complex embeddings of $E$ over $v$ if $v$ is infinite, and the largest absolute
value of an entry if $v$ is finite; there the absolute value of
$y\in F_v\otimes_FE$ is the largest $|\tau(y)|_v$ over the $F$-embeddings $\tau$
of $E$ into an algebraic closure of $F_v$, with $|\cdot|_v$ extending the
normalized absolute value of $F_v$. When $E$ is a CM field and $F$ its maximal
totally real subfield, this is the height of
Definition~\ref{asm-theta-agree-doubling}(2). For $z\in\mathbf A$ let
$|z_\infty|$ be the largest absolute value of its components at the infinite
places. Then there are a basis $\beta_1,\ldots,\beta_r$ of $\mathbf A_E$ over
$\mathbf A$ and a real number $C\geq0$ with the following property. For every
$n\geq1$ and $g\in GL_n(\mathbf A_E)$ there is an integer $m$ with
$1\leq m\leq\operatorname{ht}(g)$ such that every coordinate $z\in\mathbf A$ of
an entry of $g$ in this basis has $|z_\infty|\leq C\operatorname{ht}(g)$ and
$(mz)_v\in\mathcal O_v$ for every finite place $v$.
\end{lemma}
\leanref{Development/WeilBound/ThetaLifts/LeviDeterminants/GrowHt.lean}{PositivityTheory.ThKF0.Grow.Ht.exists_entries}
\leanref{Development/WeilBound/ThetaLifts/LeviDeterminants/GrowHt.lean}{PositivityTheory.ThKF0.Grow.Ht.exists_norm_repr_le}
\leanref{Development/AbelianVarietyTheory/Proofs/FoundationProofs/ModerateGrowthOfHasCompactSupport.lean}{GlobalTheory.AdelicU.HeightProper.nnnorm_repr_le_absFin}
\begin{proof}
Choose a basis $b_1,\ldots,b_r$ of $E$ over $F$ of algebraic integers, let
$b^*_1,\ldots,b^*_r$ be the dual basis for the trace form, so that
$\operatorname{Tr}_{E/F}(b^*_kb_l)$ is $1$ if $k=l$ and $0$ otherwise, and put
$\beta_k=1\otimes b^*_k$. The $k$-th coordinate of $y\in\mathbf A_E$ is
$\operatorname{Tr}_{\mathbf A_E/\mathbf A}(yb_k)$. Its component at a place $v$
is the $k$-th coordinate of the component $y_v\in F_v\otimes_FE$ in the basis
$(1\otimes b^*_k)_k$. The image of this coordinate in $\mathbf C$ if $v$ is
infinite, and in an algebraic closure of $F_v$ if $v$ is finite, is
$\sum_\tau\tau(y_v)\tau(b_k)$, the sum over the $F$-embeddings $\tau$ of $E$
into $\mathbf C$ over $v$, or into the algebraic closure, extended to
$F_v\otimes_FE$: both sides are additive in $y_v$ and agree on the elements
$x\otimes e$. So at an infinite place $w$ the $k$-th coordinate has absolute
value at most $B_{k,w}\max_\tau|\tau(y_w)|$, with $B_{k,w}=\sum_\tau|\tau(b_k)|$.
At a finite place $v$ it has absolute value at most $\max_\tau|\tau(y_v)|_v$, as
$|\cdot|_v$ is ultrametric and $|\tau(b_k)|_v\leq1$ for the algebraic integer
$b_k$. Put $C=\sum_{k,w}B_{k,w}$.

Let $g\in GL_n(\mathbf A_E)$, and put $H_v=\max(1,\|g_v\|,\|g_v^{-1}\|)$ for
finite $v$; then $H_v=1$ for almost all $v$.
Lemma~\ref{asm-theta-agree-definite-inputs-growth-inputs-denominator} gives an
integer $m$ with $1\leq m\leq\prod_vH_v\leq\operatorname{ht}(g)$, the factors of
$\operatorname{ht}(g)$ at the infinite places being at least $1$. Let $y$ be an
entry of $g$ and $z$ its $k$-th coordinate. At a finite place $v$,
$|z_v|_v\leq\max_\tau|\tau(y_v)|_v\leq\|g_v\|\leq H_v$, so $|mz_v|_v\leq1$. At an
infinite place $w$, an entry of a complex matrix is bounded by its operator
norm, so $|\tau(y_w)|\leq\|g_w\|\leq\operatorname{ht}(g)$ for every $\tau$, and
$|z_w|\leq B_{k,w}\operatorname{ht}(g)\leq C\operatorname{ht}(g)$.
\end{proof}
% Lean: exists_entries (GrowHt.lean), with b from exists_is_basis_integral and the basis
% Algebra.TensorProduct.basis (AF F) b.traceDual; the coordinate at w is
% HeightProper.repr_map, its bound exists_norm_repr_le (key: the coordinate is the sum of
% infEmb tau y * tau (b i)) with norm_entry_le_opNorm, normInf_le_iff and
% htInf_le_heightGL; the finite bound is nnnorm_repr_le_absFin (algebraMap_basis_repr,
% absFin_tmul_le) with absFin_le_normFin and le_htOf; m is the d of denom_global for
% H_v = htFin, and norm_evH_le gives |z_infty|.

\begin{lemma}\label{asm-theta-agree-definite-inputs-growth-inputs-bilinear}
Let $E$ be a finite extension of a number field $F$, put $\mathbf A=\mathbf A_F$
and $\mathbf A_E=\mathbf A\otimes_FE$, and let $\beta_1,\ldots,\beta_r$ be a
basis of $\mathbf A_E$ over $\mathbf A$. Let $\kappa$ be a finite set and
$e\colon M_{n\times d}(\mathbf A_E)\to\mathbf A^\kappa\times\mathbf A^\kappa$ an
$\mathbf A$-linear isomorphism. Call $w\in\mathbf A^\kappa\times\mathbf A^\kappa$
\emph{integral} if its components at $\infty$ are $0$ and its components at
every finite place $v$ have entries in $\mathcal O_v$, and \emph{at $\infty$} if
its components at the finite places are $0$. For $z\in\mathbf A$ let
$|z_\infty|$ be the largest absolute value of its components at the infinite
places, and for $w$ let $|w|$ be the largest $|z_\infty|$ over its entries $z$.
There are an integer $m_0\geq1$ and a real number $K\geq0$ with the following
property. Let $a,b\geq1$ be integers and $x,y\geq0$ real numbers, and let
$M\in M_n(\mathbf A_E)$ and $N\in M_d(\mathbf A_E)$ be such that every
coordinate $z$, in the basis $(\beta_k)$, of an entry of $M$ has
$|z_\infty|\leq x$ and $(az)_v\in\mathcal O_v$ for every finite place $v$, and
every such coordinate $z$ of an entry of $N$ has $|z_\infty|\leq y$ and
$(bz)_v\in\mathcal O_v$ for every finite place $v$. Let $s$ be the
$\mathbf A$-linear map $w\mapsto e\bigl(M\,e^{-1}(w)\,N^{\mathsf T}\bigr)$. Then
$abm_0\,s(w)$ is integral for every integral $w$, and $s$ maps every $w$ at
$\infty$ to an element at $\infty$ with $|s(w)|\leq Kxy\,|w|$.
\end{lemma}
\leanref{Development/WeilBound/ThetaLifts/LeviDeterminants/GrowDef.lean}{PositivityTheory.ThKF0.Grow.Bd}
\leanref{Development/WeilBound/ThetaLifts/LeviDeterminants/GrowHt.lean}{PositivityTheory.ThKF0.Grow.Ht.bd_bilin}
\leanref{Development/WeilBound/ThetaLifts/LeviDeterminants/GrowHt.lean}{PositivityTheory.ThKF0.Grow.Ht.bd_of_sz}
\leanref{Development/WeilBound/ThetaLifts/LeviDeterminants/GrowHt.lean}{PositivityTheory.ThKF0.Grow.Ht.exists_sz}
\begin{proof}
Let $(\varepsilon_j)_{j\in\kappa\sqcup\kappa}$ be the standard basis of
$\mathbf A^\kappa\times\mathbf A^\kappa$ over $\mathbf A$, put
$u_j=e^{-1}(\varepsilon_j)$, and let $l_i(Y)$ be the $i$-th coordinate of
$e(Y)$, an $\mathbf A$-linear form. For a matrix $Y$ over $\mathbf A_E$ let
$Y^{(k)}_{pq}$ be the $k$-th coordinate of its entry $Y_{pq}$, so that
$Y=\sum_{p,q,k}Y^{(k)}_{pq}\beta_kE_{pq}$, with $E_{pq}$ the matrix units. The
$i$-th coordinate of $s(\varepsilon_j)$ is
\[
s_{ij}=l_i\bigl(Mu_jN^{\mathsf T}\bigr)=\sum M^{(k)}_{pq}N^{(k')}_{p'q'}\,c_{ij}(p,q,k;p',q',k'),\qquad
c_{ij}(p,q,k;p',q',k')=l_i\bigl(\beta_kE_{pq}\,u_j\,(\beta_{k'}E_{p'q'})^{\mathsf T}\bigr),
\]
a sum of $P=n^2d^2r^2$ terms. The finitely many adeles $c_{ij}(\ldots)$ do not
depend on $M$ and $N$. Each of them, say $c$, has components in $\mathcal O_v$
for almost all $v$, so
Lemma~\ref{asm-theta-agree-definite-inputs-growth-inputs-denominator}, for
$H_v=\max(1,|c_v|_v)$, gives an integer $m_c\geq1$ with
$(m_cc)_v\in\mathcal O_v$ for every finite $v$. Let $m_0$ be the product of the
$m_c$ and $c_0$ the sum of the $|c_\infty|$. Then every term of $s_{ij}$ has
absolute value at most $xyc_0$ at every infinite place, and its product with
$abm_0$ has components in $\mathcal O_v$ at every finite $v$. So
$|(s_{ij})_\infty|\leq Pxyc_0$ and $(abm_0s_{ij})_v\in\mathcal O_v$.

Now $s(w)=\sum_i\bigl(\sum_jw_js_{ij}\bigr)\varepsilon_i$, where the $w_j$ are
the coordinates of $w$. If $w$ is integral, each $abm_0w_js_{ij}$ has component
$0$ at $\infty$ and components in $\mathcal O_v$, so $abm_0s(w)$ is integral. If
$w$ is at $\infty$, so is each $w_js_{ij}$, and
$|\sum_jw_js_{ij}|\leq2|\kappa|Pxyc_0|w|$. So $K=2|\kappa|Pc_0$ will do.
\end{proof}
% Lean: bd_bilin (GrowHt.lean), with l and u as in the text, bilin and expand for the
% expansion, exists_sz (exists_den, through denom_global) for m_0 and c_0, Sz.mul and Sz.sum
% for the terms, and bd_of_sz (crd_apply, intF_mul, intF_sum, inf_mul, inf_sum) for the two
% clauses of Bd.

\begin{definition}\label{asm-theta-agree-definite-inputs-growth-inputs-derivatives}
Let $F$ be a number field, put $\mathbf A=\mathbf A_F$, let $\kappa$ be a
finite set, and let $\psi$ be a continuous character of $\mathbf A$, trivial on
$F$, with $\psi(z)=\exp(-2\pi i\operatorname{Tr}_{F_\infty/\mathbf R}z_\infty)$
for $z$ at $\infty$. Put $V=F_\infty^\kappa$, and regard $V$ and $V\times V$ as
the elements of $\mathbf A^\kappa$ and of $\mathbf A^\kappa\times\mathbf A^\kappa$
at $\infty$, whose components at the finite places are $0$. For
$w=(w_1,w_2)\in\mathbf A^\kappa\times\mathbf A^\kappa$ and a function $\phi$ on
$\mathbf A^\kappa$ put
$(\rho(w)\phi)(x)=\psi\bigl(x\cdot w_2+\tfrac12w_1\cdot w_2\bigr)\phi(x+w_1)$.
Call $u\in\mathbf A^\kappa$, or $w\in\mathbf A^\kappa\times\mathbf A^\kappa$,
\emph{integral} if its components at $\infty$ are $0$ and its components at every
finite place $v$ have entries in $\mathcal O_v$. Let $|\cdot|$ be the largest
absolute value of the components at the infinite places, on $V$ and on
$V\times V$, and put
$T(x,b)=\operatorname{Tr}_{F_\infty/\mathbf R}\bigl(\sum_{i\in\kappa}x_ib_i\bigr)$
for $x,b\in V$.
\begin{enumerate}
\item For $p\in V\times V$ and a function $\Phi$ on $\mathbf A^\kappa$, the
  \emph{derivative} $D_p\Phi(x)$ is the derivative at $t=0$ of
  $t\mapsto(\rho(tp)\Phi)(x)$, where it exists. For a finite sequence
  $l=(p_1,\ldots,p_k)$ in $V\times V$, of \emph{length} $k$, put
  $D_l=D_{p_1}\circ\cdots\circ D_{p_k}$, the identity if $k=0$.
\item Let $m\geq1$ be an integer. A function $\Theta$ on $\mathbf A^\kappa$ is
  \emph{$m$-periodic} if $\Theta(x+mu)=\Theta(x)$ for every $x$ and every
  integral $u\in\mathbf A^\kappa$.
\item For a Haar measure $\mu$ on $\mathbf A^\kappa$ and a function $\Theta$ on
  $\mathbf A^\kappa$ put $\|\Theta\|_\mu=\bigl(\int|\Theta|^2\,d\mu\bigr)^{1/2}$,
  which may be infinite.
\end{enumerate}
\end{definition}
\leanref{Development/WeilBound/ThetaLifts/LeviDeterminants/GrowDef.lean}{PositivityTheory.ThKF0.Grow.IntP}
\leanref{Development/WeilBound/ThetaLifts/LeviDeterminants/GrowDef.lean}{PositivityTheory.ThKF0.Grow.dR}
\leanref{Development/WeilBound/ThetaLifts/LeviDeterminants/GrowDef.lean}{PositivityTheory.ThKF0.Grow.word}
\leanref{Development/WeilBound/ThetaLifts/LeviDeterminants/GrowDef.lean}{PositivityTheory.ThKF0.Grow.Lev}
\leanref{Development/WeilBound/ThetaLifts/LeviDeterminants/GrowDef.lean}{PositivityTheory.ThKF0.Grow.nsq}

\begin{lemma}\label{asm-theta-agree-definite-inputs-growth-inputs-derivatives-bounds}
Let $F$, $\kappa$, $\psi$, $V$, $\rho$, $|\cdot|$, $T$, the derivatives and the
norms $\|\cdot\|_\mu$ be as in
Definition~\ref{asm-theta-agree-definite-inputs-growth-inputs-derivatives}, let
$\Phi$ be a Schwartz-Bruhat function on $\mathbf A^\kappa$, and let
$p=(a,b)\in V\times V$.
\begin{enumerate}
\item For $f\in\mathcal S(V)$ and a Schwartz-Bruhat function $g$ on
  $\mathbf A_f^\kappa$, $D_p(f\otimes g)=D_{a,b}f\otimes g$, where
  $D_{a,b}f(x)=\partial_af(x)-2\pi i\,T(x,b)f(x)$ and $\partial_a$ is the
  derivative in the direction $a$. For every $x$ the function
  $t\mapsto(\rho(tp)\Phi)(x)$ is differentiable at $0$; $D_p\Phi$ is a
  Schwartz-Bruhat function, linear in $\Phi$ and $\mathbf R$-linear in $p$; and
  if $\Phi$ is $m$-periodic, so is $D_p\Phi$. For every function $\Theta$ on
  $\mathbf A^\kappa$, $D_{(0,b)}\Theta(x)=-2\pi i\,T(x_\infty,b)\,\Theta(x)$.
\item There is a continuous function $G$ on $\mathbf A^\kappa$ with
  $\|G\|_\mu<\infty$ for every Haar measure $\mu$, such that
  $|(\rho(tp)\Phi)(x)-\Phi(x)|\leq|t|\,|G(x)|$ for every $x$ and every real $t$
  with $|t|\leq1$.
\item Let $\mu$ be a Haar measure on $\mathbf A^\kappa$ and $k\geq0$ an integer.
  There is $C$ such that $\|D_l\Phi\|_\mu^2\leq C\,c^{2k}$ for every $c\geq0$ and
  every sequence $l$ in $V\times V$ of length $k$ whose terms have
  $|\cdot|\leq c$.
\end{enumerate}
\end{lemma}
\leanref{Development/WeilBound/ThetaLifts/LeviDeterminants/GrowDef.lean}{PositivityTheory.ThKF0.Grow.DerivOK}
\leanref{Development/WeilBound/ThetaLifts/LeviDeterminants/ContinuousOfMem.lean}{PositivityTheory.ThKF0.Grow.derivOK}
\leanref{Development/WeilBound/ThetaLifts/LeviDeterminants/ContinuousOfMem.lean}{PositivityTheory.ThKF0.Grow.dR_pure}
\leanref{Development/WeilBound/ThetaLifts/LeviDeterminants/GrowInt.lean}{PositivityTheory.ThKF0.Grow.dR_sum}
\leanref{Development/WeilBound/ThetaLifts/LeviDeterminants/GrowFinal.lean}{PositivityTheory.ThKF0.Grow.dR_mult}
\leanref{Development/WeilBound/ThetaLifts/LeviDeterminants/GrowInt.lean}{PositivityTheory.ThKF0.Grow.exists_dom}
\leanref{Development/WeilBound/ThetaLifts/LeviDeterminants/GrowLine.lean}{PositivityTheory.ThKF0.Grow.lineD_bound}
\leanref{Development/WeilBound/ThetaLifts/LeviDeterminants/GrowInt.lean}{PositivityTheory.ThKF0.Grow.word_bound}
\begin{proof}
(1) As $p$ is at $\infty$, for $x\in\mathbf A^\kappa$ and real $t$ the argument
$tx\cdot b+\tfrac12t^2a\cdot b$ of $\psi$ is at $\infty$, and $x+ta$ has the
finite components of $x$. So
\[
\rho(tp)(f\otimes g)(x)=\exp\bigl(-2\pi i\bigl(t\,T(x_\infty,b)+t^2c_p\bigr)\bigr)f(x_\infty+ta)\,g(x_f),\qquad
c_p=\tfrac12\operatorname{Tr}_{F_\infty/\mathbf R}(a\cdot b),
\]
whose derivative at $t=0$ is $D_{a,b}f(x_\infty)g(x_f)$, as in the proof of
Lemma~\ref{asm-theta-agree-density-derivative}; and $D_{a,b}f$ is a Schwartz
function. A Schwartz-Bruhat function is a finite sum of such products, which
gives the differentiability, and that $D_p\Phi$ is Schwartz-Bruhat and linear in
$\Phi$; and $D_{a,b}f$ is $\mathbf R$-linear in $(a,b)$. Let $u$ be integral.
Then $mu\cdot b=0$, as $u$ has components $0$ at $\infty$ and $b$ has components
$0$ at the finite places. So if $\Phi$ is $m$-periodic, then
$(\rho(tp)\Phi)(x+mu)=(\rho(tp)\Phi)(x)$ for all $t$, and differentiating at
$t=0$ gives $D_p\Phi(x+mu)=D_p\Phi(x)$. Finally,
$(\rho(t(0,b))\Theta)(x)=\psi(tx\cdot b)\Theta(x)=\exp(-2\pi i\,t\,T(x_\infty,b))\Theta(x)$.

(2) Let first $\Phi=f\otimes g$, and put $k=\dim_{\mathbf R}V+1$. For fixed $x$
the function $\ell(t)=\exp(-2\pi i(tT(x_\infty,b)+t^2c_p))f(x_\infty+ta)$ has
derivative
\[
\ell'(t)=\exp\bigl(-2\pi i(tT(x_\infty,b)+t^2c_p)\bigr)\bigl(-2\pi i\bigl(T(x_\infty,b)+2tc_p\bigr)f(q)+\partial_af(q)\bigr),\qquad q=x_\infty+ta.
\]
Let $|t|\leq1$. Then $|x_\infty|\leq|q|+|a|$, so
$|T(x_\infty,b)+2tc_p|\leq A(1+|q|)$ for a constant $A$, and
$1+|x_\infty|\leq(1+|a|)(1+|q|)$. As $f$ is a Schwartz function,
$(1+|q|)^{k+1}|f(q)|$ and $(1+|q|)^k|\partial_af(q)|$ are bounded. So
$|\ell'(t)|\leq K(1+|x_\infty|)^{-k}$, with $K$ independent of $x$ and $t$, and
by the mean value inequality
$|(\rho(tp)\Phi)(x)-\Phi(x)|=|\ell(t)-\ell(0)|\,|g(x_f)|\leq|t|K(1+|x_\infty|)^{-k}|g(x_f)|$.
The function $G(x)=K(1+|x_\infty|)^{-k}g(x_f)$ is continuous with
$\|G\|_\mu<\infty$: a Haar measure on $\mathbf A^\kappa=V\times\mathbf A_f^\kappa$
is a multiple of a product of Haar measures, $(1+|z|)^{-2k}$ is integrable on $V$
as $2k>\dim_{\mathbf R}V$, and $g$ is bounded with compact support. If the claim
holds for $\Phi_1$ and $\Phi_2$, with $G_1$ and $G_2$, it holds for
$\Phi_1+\Phi_2$ with $|G_1|+|G_2|$, and for $z\Phi_1$ with $zG_1$.

(3) We argue by induction on $k$, for all $\Phi$ at once. For $k=0$ take
$C=\|\Phi\|_\mu^2$, which is finite, as Schwartz functions are square
integrable and Schwartz-Bruhat functions on $\mathbf A_f^\kappa$ are bounded with
compact support. Let $k\geq1$, let $p_1,\ldots,p_J$ be an $\mathbf R$-basis of
$V\times V$, and $\lambda_j$ the coordinate functions, with
$|\lambda_j(p')|\leq K_j|p'|$. By induction, applied to the Schwartz-Bruhat
functions $D_{p_j}\Phi$, there are $C_j$ with
$\|D_{l'}D_{p_j}\Phi\|_\mu^2\leq C_jc^{2(k-1)}$ for every sequence $l'$ of length
$k-1$ with terms of $|\cdot|\leq c$. A sequence of length $k$ is $l=(l',p')$,
with last term $p'$, and by (1),
$D_l\Phi=D_{l'}D_{p'}\Phi=\sum_j\lambda_j(p')D_{l'}D_{p_j}\Phi$. As
$|\sum_{j=1}^J\Theta_j|^2\leq2^J\sum_j|\Theta_j|^2$ pointwise,
$\|D_l\Phi\|_\mu^2\leq2^J\sum_jK_j^2c^2C_jc^{2(k-1)}$.
\end{proof}
% Lean: (1) is dR_pure (ContinuousOfMem.lean: hasDerivAt_pure through schr_emb, phase_line
% and lineFn), hasDerivAt_dR, dR_mem, dR_add, dR_smul, dR_sum (GrowInt.lean), the third
% clause of derivOK, and dR_mult (GrowFinal.lean). (2) is exists_dom (GrowInt.lean), with
% lineD_bound (GrowLine.lean) for K, weight_lt_top for k, nsq_pure_lt_top, and exists_major
% for sums. (3) is word_bound, with Module.finBasis, dR_sum, nsq_sum_le and nsq_lt_top.

\begin{lemma}\label{asm-theta-agree-definite-inputs-growth-inputs-transport}
Let $F$, $\kappa$, $V$, $\rho$, $|\cdot|$, the derivatives and the norms
$\|\cdot\|_\mu$ be as in
Definition~\ref{asm-theta-agree-definite-inputs-growth-inputs-derivatives}. Let
$s$ be an $\mathbf A$-linear map of $\mathbf A^\kappa\times\mathbf A^\kappa$, and
$\Omega$ a linear operator on functions on $\mathbf A^\kappa$ that maps
Schwartz-Bruhat functions to Schwartz-Bruhat functions, has
$\|\Omega\Theta\|_\mu=\|\Theta\|_\mu$ for every Schwartz-Bruhat $\Theta$ and
every Haar measure $\mu$, and has $\Omega\rho(s(w))\Theta=\rho(w)\Omega\Theta$
for every $w$ and every Schwartz-Bruhat $\Theta$. Let $\Phi$ be a
Schwartz-Bruhat function.
\begin{enumerate}
\item If $p,p'\in V\times V$ and $s(p)=p'$, then
  $\Omega(D_{p'}\Phi)=D_p(\Omega\Phi)$.
\item Let $c\geq0$, and suppose that $s$ maps every $p\in V\times V$ to some
  $p'\in V\times V$ with $|p'|\leq c|p|$. Let $R\geq0$, and let $l$ be a finite
  sequence in $V\times V$ whose terms have $|\cdot|\leq R$. Then there is a
  sequence $l'$ of the same length, whose terms have $|\cdot|\leq cR$, with
  $D_l(\Omega\Phi)=\Omega(D_{l'}\Phi)$.
\end{enumerate}
\end{lemma}
\leanref{Development/WeilBound/ThetaLifts/LeviDeterminants/GrowDef.lean}{PositivityTheory.ThKF0.Grow.Impl}
\leanref{Development/WeilBound/ThetaLifts/LeviDeterminants/GrowInt.lean}{PositivityTheory.ThKF0.Grow.Impl.dR}
\leanref{Development/WeilBound/ThetaLifts/LeviDeterminants/GrowInt.lean}{PositivityTheory.ThKF0.Grow.Impl.word}
\begin{proof}
(1) Fix a Haar measure $\mu$. For $n\geq1$ put
$q_n=n\bigl(\rho(n^{-1}p')\Phi-\Phi\bigr)-D_{p'}\Phi$. By the formula in the
proof of
Lemma~\ref{asm-theta-agree-definite-inputs-growth-inputs-derivatives-bounds}(1),
$\rho(tp')$ maps Schwartz-Bruhat functions to Schwartz-Bruhat functions, so $q_n$
is Schwartz-Bruhat. By (1) of that lemma, $q_n(x)\to0$ for every $x$; by its
(2) for $p'$ and $t=n^{-1}$, $|q_n(x)|\leq|G(x)|+|D_{p'}\Phi(x)|$, where $G$ and
$D_{p'}\Phi$ have finite norm. By dominated convergence, $\|q_n\|_\mu\to0$. As $s$ is $\mathbf A$-linear,
$s(n^{-1}p)=n^{-1}p'$, so $\Omega\rho(n^{-1}p')\Phi=\rho(n^{-1}p)\Omega\Phi$ and
\[
\Omega q_n=n\bigl(\rho(n^{-1}p)\Omega\Phi-\Omega\Phi\bigr)-\Omega(D_{p'}\Phi),
\]
which tends pointwise to $\Delta=D_p(\Omega\Phi)-\Omega(D_{p'}\Phi)$, by (1) of
that lemma for the Schwartz-Bruhat function $\Omega\Phi$. By Fatou's lemma,
$\|\Delta\|_\mu\leq\liminf_n\|\Omega q_n\|_\mu=\liminf_n\|q_n\|_\mu=0$. The
function $\Delta$ is Schwartz-Bruhat, hence continuous, and it vanishes almost
everywhere for $\mu$, which is positive on nonempty open sets. So $\Delta=0$.

(2) We argue by induction on the length; for length $0$ take $l'=l$. Let
$l=(p,l_1)$, so that $D_l=D_pD_{l_1}$. By induction there is a sequence $l'_1$,
of the length of $l_1$ and with terms of $|\cdot|\leq cR$, with
$D_{l_1}(\Omega\Phi)=\Omega(D_{l'_1}\Phi)$. Let $p'=s(p)$, which lies in
$V\times V$ with $|p'|\leq c|p|\leq cR$. By (1), for the Schwartz-Bruhat function
$D_{l'_1}\Phi$,
$D_l(\Omega\Phi)=D_p\Omega(D_{l'_1}\Phi)=\Omega(D_{p'}D_{l'_1}\Phi)$; so
$l'=(p',l'_1)$ will do.
\end{proof}
% Lean: Impl psi T s is the hypothesis on Omega (pres, impl, unit). (1) is Impl.dR
% (GrowInt.lean): q n uses 1/(n+1); t0 is the dominated convergence (exists_dom,
% exists_major, slope_tendsto), t1 the pointwise limit through hT.impl, t2 Fatou
% (lintegral_liminf_le) with hT.unit, and the end is Continuous.ae_eq_iff_eq. (2) is
% Impl.word, with the clause inf of Bd for p'.

\begin{lemma}\label{asm-theta-agree-definite-inputs-growth-inputs-sobolev-approximation}
Let $F$ be a number field with ring of finite adeles $\mathbf A_f$, put
$\widehat{\mathcal O}=\prod_v\mathcal O_v$, the product over the finite places $v$
of $F$, and let $N\geq1$ be an integer.
\begin{enumerate}
\item Every $y\in\mathbf A_f$ is $x+Na$ with $x\in F$ and $a\in\widehat{\mathcal O}$.
\item Every $y\in\widehat{\mathcal O}$ is $x+Na$ with $x\in\mathcal O_F$ and
  $a\in\widehat{\mathcal O}$.
\end{enumerate}
\end{lemma}
\leanref{Development/AbelianVarietyTheory/Proofs/UnitaryExistence/FiniteAdeleStrongApproximationModNat.lean}{GlobalTheory.AdLat.strong_approximation_mod_nat}
\leanref{Development/AbelianVarietyTheory/Proofs/UnitaryExistence/FiniteAdeleStrongApproximationModNat.lean}{GlobalTheory.AdLat.exists_eq_algebraMap_add_nsmul}
\begin{proof}
For a finite place $v$ let $\mathfrak p_v$ be its prime ideal and $v(\cdot)$ its
valuation.

\emph{One element.} Let $T$ be a finite set of finite places and $c_v\in F$ for
$v\in T$. Then there is $w\in F$ with $w-c_v\in\mathcal O_v$ for $v\in T$ and
$w\in\mathcal O_v$ for $v\notin T$. Indeed, choose a nonzero $d\in\mathcal O_F$
with $dc_v\in\mathcal O_F$ for every $v\in T$, put $g_v=dc_v$ for $v\in T$ and
$g_v=0$ for $v\notin T$, and let $S$ be the union of $T$ and the finitely many
places $v$ with $v(d)>0$. By the Chinese remainder theorem in $\mathcal O_F$
there is $r\in\mathcal O_F$ with $r-g_v\in\mathfrak p_v^{v(d)}$ for every
$v\in S$. Put $w=r/d$. For $v\in S$, $w-g_v/d=(r-g_v)/d$ has valuation at least
$v(d)-v(d)=0$; at $v\in T$ this is the claim, as $g_v/d=c_v$, and at
$v\in S\setminus T$ it gives $w\in\mathcal O_v$, as $g_v=0$. For $v\notin S$,
$v(d)=0$, so $w=r/d\in\mathcal O_v$.

\emph{An adele.} Let $z\in\mathbf A_f$. Then there is $w\in F$ with
$z_v-w\in\mathcal O_v$ for every $v$. Indeed, the set $T$ of the places $v$ with
$z_v\notin\mathcal O_v$ is finite. For $v\in T$, the set $z_v+\mathcal O_v$ is
open in $F_v$ and $F$ is dense in $F_v$, so there is $c_v\in F$ with
$z_v-c_v\in\mathcal O_v$. Let $w$ be given by the first step. For $v\in T$,
$z_v-w=(z_v-c_v)+(c_v-w)\in\mathcal O_v$; for $v\notin T$, both $z_v$ and $w$
lie in $\mathcal O_v$.

(1) The second step, for $z=N^{-1}y$, gives $w\in F$ with
$a=N^{-1}y-w\in\widehat{\mathcal O}$. Then $y=Nw+Na$, and $x=Nw$ will do.

(2) Write $y=x+Na$ as in (1). As $y$ and $Na$ lie in $\widehat{\mathcal O}$, the
element $x=y-Na$ of $F$ lies in $\mathcal O_v$ for every finite place $v$, so
$x\in\mathcal O_F$.
\end{proof}
% Lean: the first step is exists_valuation_sub_lt_of_injective (private; d from
% IsLocalization.exist_integer_multiples, r from IsDedekindDomain.exists_forall_sub_mem_ideal
% with the exponents count v (span {d})), the second exists_sub_mem_adicCompletionIntegers'
% (private; HeightOneSpectrum.denseRange_algebraMap), (1) exists_eq_algebraMap_add_nsmul, and
% (2) the second clause of strong_approximation_mod_nat, through
% HeightOneSpectrum.mem_integers_of_valuation_le_one.

\begin{lemma}\label{asm-theta-agree-definite-inputs-growth-inputs-sobolev}
Let $F$, $\kappa$, $V$, the derivatives, $m$-periodicity and the norms
$\|\cdot\|_\mu$ be as in
Definition~\ref{asm-theta-agree-definite-inputs-growth-inputs-derivatives}. There
are a Haar measure $\mu$ on $\mathbf A^\kappa$, a finite set $S$ of finite
sequences in $V\times V$, and real numbers $C\geq0$ and $k\geq0$ such that
\[
|\Theta(x)|^2\leq C\,m^k\sum_{l\in S}\|D_l\Theta\|_\mu^2
\]
for every integer $m\geq1$, every $m$-periodic Schwartz-Bruhat function
$\Theta$ on $\mathbf A^\kappa$ and every $x\in\mathbf A^\kappa$.
\end{lemma}
\leanref{Development/WeilBound/ThetaLifts/LeviDeterminants/GrowDef.lean}{PositivityTheory.ThKF0.Grow.Sob}
\leanref{Development/WeilBound/ThetaLifts/LeviDeterminants/GrowSob.lean}{PositivityTheory.ThKF0.Grow.sob_of}
\leanref{Development/WeilBound/ThetaLifts/LeviDeterminants/GrowSob.lean}{PositivityTheory.ThKF0.Grow.slice}
\leanref{Development/WeilBound/ThetaLifts/LeviDeterminants/GrowSobA.lean}{PositivityTheory.ThKF0.Grow.oneD}
\leanref{Development/WeilBound/ThetaLifts/LeviDeterminants/GrowSobA.lean}{PositivityTheory.ThKF0.Grow.sobAbs}
\leanref{Development/WeilBound/ThetaLifts/LeviDeterminants/GrowSobB.lean}{PositivityTheory.ThKF0.Grow.measure_Om_le}
\leanref{Development/WeilBound/ThetaLifts/LeviDeterminants/GrowSobB.lean}{PositivityTheory.ThKF0.Grow.exists_reps}
\begin{proof}
\emph{One variable.} Let $u\colon\mathbf R\to\mathbf C$ be differentiable with
continuous derivative $u'$. Then
$|u(0)|^2\leq2\int_{\mathbf R}(|u|^2+|u'|^2)$. Indeed, $(t-1)u(t)^2$ has
derivative $u(t)^2+2(t-1)u(t)u'(t)$ and vanishes at $t=1$, so
$u(0)^2=\int_0^1\bigl(u(t)^2+2(t-1)u(t)u'(t)\bigr)\,dt$; and on $[0,1]$,
$|t-1|\leq1$ and $2|u||u'|\leq|u|^2+|u'|^2$, so the integrand has absolute value
at most $2(|u|^2+|u'|^2)$.

\emph{Several variables.} For $c\in V$ and a Schwartz-Bruhat $\Theta$,
$(\rho(t(c,0))\Theta)(x)=\Theta(x+tc)$; so by
Lemma~\ref{asm-theta-agree-definite-inputs-growth-inputs-derivatives-bounds}(1)
the function $t\mapsto\Theta(x+tc)$ has derivative $D_{(c,0)}\Theta(x+tc)$ at
every $t$, which is continuous in $t$, as Schwartz-Bruhat functions are
continuous. We show: for distinct $c_1,\ldots,c_j\in V$ there is a finite set
$L$ of finite sequences in $\{(c_1,0),\ldots,(c_j,0)\}$ such that
\[
|\Theta(x)|^2\leq2^j\int_{\mathbf R^j}\sum_{l\in L}\Bigl|D_l\Theta\Bigl(x+\sum_it_ic_i\Bigr)\Bigr|^2dt
\]
for every Schwartz-Bruhat $\Theta$ and every $x$. For $j=0$ take $L=\{\varnothing\}$. If $L$
works for $c_2,\ldots,c_j$, then $L\cup\{(l,(c_1,0)):l\in L\}$, a disjoint
union, works for $c_1,\ldots,c_j$, where $(l,(c_1,0))$ is $l$ followed by
$(c_1,0)$: apply the one-variable
inequality to $u(s)=\Theta(x+sc_1)$, whose derivative is
$D_{(c_1,0)}\Theta(x+sc_1)$, then the inequality for $c_2,\ldots,c_j$ to the
Schwartz-Bruhat functions $\Theta$ and $D_{(c_1,0)}\Theta$ at $x+sc_1$, and
Tonelli's theorem.

\emph{Periodic integrals.} Let $b_1,\ldots,b_D$ be an $\mathbf R$-basis of $V$,
$\nu$ a Haar measure on $\mathbf A_f^\kappa$,
$\widehat{\mathcal O}=\prod_v\mathcal O_v$, and $\mu$ the image of the product of
Lebesgue measure on $\mathbf R^D$ and $\nu$ under $J(t,y)=(\sum_it_ib_i,y)$, a
Haar measure on $\mathbf A^\kappa=V\times\mathbf A_f^\kappa$. Put
$k=[F:\mathbf Q]\,|\kappa|$. For every integer $m\geq1$,
$\nu(\widehat{\mathcal O}^\kappa)\leq m^k\nu(m\widehat{\mathcal O}^\kappa)$.
Indeed, let $e_1,\ldots,e_n$ be a $\mathbf Z$-basis of $\mathcal O_F$, with
$n=[F:\mathbf Q]$. By
Lemma~\ref{asm-theta-agree-definite-inputs-growth-inputs-sobolev-approximation}(2),
every $y\in\widehat{\mathcal O}$ is $z+ma$ with $z\in\mathcal O_F$ and
$a\in\widehat{\mathcal O}$. Write $z=\sum_jz_je_j$ and $z_j=r_j+mq_j$ with
integers $0\leq r_j<m$; then
$y=\sum_jr_je_j+m\bigl(\sum_jq_je_j+a\bigr)$, and
$\sum_jq_je_j+a\in\widehat{\mathcal O}$. So $\widehat{\mathcal O}$ is covered by
the $m^n$ translates of $m\widehat{\mathcal O}$ by the $\sum_jr_je_j$, and
$\widehat{\mathcal O}^\kappa$ by $m^k$ translates of
$m\widehat{\mathcal O}^\kappa$. Now let
$G\geq0$ be a continuous function on $\mathbf A^\kappa$ with $G(x+mu)=G(x)$ for
every $x$ and every integral $u$, and let $x=J(t_0,y)$. By the translation
invariance of Lebesgue measure,
$\int_{\mathbf R^D}G(x+\sum_it_ib_i)\,dt=\int_{\mathbf R^D}G(J(t,y))\,dt$. For
$y'\in y+m\widehat{\mathcal O}^\kappa$, $G(J(t,y'))=G(J(t,y))$. Integrating over
these $y'$,
\[
\nu(m\widehat{\mathcal O}^\kappa)\int_{\mathbf R^D}G(J(t,y))\,dt\leq\int_{\mathbf R^D}\int G(J(t,y'))\,d\nu(y')\,dt=\int G\,d\mu .
\]
As $\widehat{\mathcal O}^\kappa$ is compact and open,
$0<\nu(\widehat{\mathcal O}^\kappa)<\infty$. So
$\int_{\mathbf R^D}G(x+\sum_it_ib_i)\,dt\leq\nu(\widehat{\mathcal O}^\kappa)^{-1}m^k\int G\,d\mu$.

\emph{Conclusion.} Let $S$ be the set $L$ of the second step for
$b_1,\ldots,b_D$. Let $\Theta$ be an $m$-periodic Schwartz-Bruhat function, and
put $G=\sum_{l\in S}|D_l\Theta|^2$. It is continuous, and $G(x+mu)=G(x)$ for
integral $u$, as every $D_l\Theta$ is $m$-periodic by
Lemma~\ref{asm-theta-agree-definite-inputs-growth-inputs-derivatives-bounds}(1).
So
\[
|\Theta(x)|^2\leq2^D\int_{\mathbf R^D}G\Bigl(x+\sum_it_ib_i\Bigr)dt\leq2^D\nu(\widehat{\mathcal O}^\kappa)^{-1}m^k\sum_{l\in S}\|D_l\Theta\|_\mu^2 .
\]
\end{proof}
% Lean: Sob, proved by sob_of derivOK (GrowSob.lean): the second step is sobAbs (GrowSobA.lean)
% with oneD and lintegral_fin_succ, for the lines lineA (b i) (schr_lineA); the third is
% slice, with J built from splitE, measure_Om_le and exists_reps (GrowSobB.lean, through
% AdLat.strong_approximation_mod_nat), lintegral_add_left_eq_self and lintegral_prod; the
% conclusion uses continuous_sumSq and lintegral_sumSq, with C = 2^D nu(Om)^-1.

\begin{lemma}\label{asm-theta-agree-definite-inputs-growth-inputs-trace}
Let $F$ be a number field and $\kappa$ a finite set, put $V=F_\infty^\kappa$,
let $|x|$ be the largest absolute value of the components of $x\in V$ at the
infinite places, and put
$T(x,b)=\operatorname{Tr}_{F_\infty/\mathbf R}\bigl(\sum_{i\in\kappa}x_ib_i\bigr)$
for $x,b\in V$. There is $c\geq0$ such that for every $x\in V$ there is $b\in V$
with $|b|\leq1$ and $|x|\leq c\,|T(x,b)|$.
\end{lemma}
\leanref{Development/WeilBound/ThetaLifts/LeviDeterminants/GrowDef.lean}{PositivityTheory.ThKF0.Grow.TrBd}
\leanref{Development/WeilBound/ThetaLifts/LeviDeterminants/GrowTr.lean}{PositivityTheory.ThKF0.Grow.trBd}
\leanref{Development/WeilBound/ThetaLifts/LeviDeterminants/GrowTr.lean}{PositivityTheory.ThKF0.Grow.exists_trace_ne}
\begin{proof}
The form $T$ is nondegenerate. Indeed, let $x\in V$ with $x_i\neq0$ for some
$i$; let $y\in F_\infty$ have the component $x_{i,w}^{-1}$ at every infinite
place $w$ with $x_{i,w}\neq0$ and $0$ at the others, and let $b$ have $b_i=y$
and $b_j=0$ for $j\neq i$. Then $x_iy$ is an idempotent of $F_\infty$ with
$x_iyx_i=x_i\neq0$, so $T(x,b)=\operatorname{Tr}_{F_\infty/\mathbf R}(x_iy)$,
the rank of multiplication by $x_iy$, is positive. So $x\mapsto T(x,\cdot)$ is
an injective linear map from $V$ to its dual, and there is $K$ with
$|x|\leq K\|T(x,\cdot)\|$ for every $x$, where $\|\cdot\|$ is the operator norm
for $|\cdot|$. For $x=0$ take $b=0$. Otherwise $T(x,\cdot)\neq0$, and there is
$b$ with $|b|\leq1$ and $|T(x,b)|>\|T(x,\cdot)\|/2$; then $|x|\leq2K|T(x,b)|$.
So $c=2K$ will do.
\end{proof}
% Lean: trBd (GrowTr.lean): B is the form, A its curried map, exists_trace_ne the
% nondegeneracy (through LinearMap.IsIdempotentElem.isProj_range and its trace), K from
% exists_antilipschitzWith, and b from exists_lt_apply_of_lt_opNorm.

\begin{lemma}\label{asm-theta-agree-definite-inputs-growth-inputs}
Let $(X,W,M,\chi_X,\chi_W,\omega)$ be a Weil pair for a theta base $B$ with
character $\psi$ (Definition~\ref{asm-theta-agree-triple}(1)), with $X$ and $W$
of ranks $n$ and $d$, put $I=\{1,\ldots,n\}\times\{1,\ldots,d\}$ and
$\mathbf A=\mathbf A_F$, and let $e$ be the isomorphism of $M$. For
$w=(w_1,w_2)\in\mathbf A^I\times\mathbf A^I$ put
$(\rho(w)\phi)(x)=\psi(x\cdot w_2+\tfrac12w_1\cdot w_2)\phi(x+w_1)$, so that the
Schrödinger operator of $u\in\mathbb W(X,W)$ is $\rho(e(u))$. For $g\in U_X$ and
$h\in U_W$ let $s_{g,h}$ be the $\mathbf A$-linear map of
$\mathbf A^I\times\mathbf A^I$ with $s_{g,h}(e(u))=e(g^{-1}u(h^{-1})^{\mathsf T})$.
Call $w$ \emph{integral} if its components at $\infty$ are $0$ and its
components at every finite place $v$ have entries in $\mathcal O_v$. Let
$|\cdot|$ be the largest absolute value of the components at the infinite
places, on $F_\infty^I$ and on $F_\infty^I\times F_\infty^I$, and let $\Phi$ be a
Schwartz-Bruhat function on $\mathbf A^I$.
\begin{enumerate}
\item There is $C\geq0$ such that for all $g\in U_X$ and $h\in U_W$ there is an
  integer $m$ with $1\leq m\leq C\operatorname{ht}(g)\operatorname{ht}(h)$, such
  that $m\,s_{g,h}(w)$ is integral for every integral $w$, and $s_{g,h}$ maps
  every $w$ at $\infty$ to an element at $\infty$ with
  $|s_{g,h}(w)|\leq C\operatorname{ht}(g)\operatorname{ht}(h)\,|w|$.
\item There is an integer $m_0\geq1$ such that
  $\rho(m_0mw)\,\omega(g,h)\Phi=\omega(g,h)\Phi$ for every integral $w$, whenever
  $g$, $h$ and $m$ are as in (1).
\item For every $N\geq0$ there are $C$ and $k$ with
  $|(\omega(g,h)\Phi)(x)|(1+|x_\infty|)^N\leq C(\operatorname{ht}(g)\operatorname{ht}(h))^k$
  for all $g\in U_X$, $h\in U_W$ and $x\in\mathbf A^I$.
\end{enumerate}
\end{lemma}
\leanref{Development/WeilBound/ThetaLifts/LeviDeterminants/GrowDef.lean}{PositivityTheory.ThKF0.Grow.HtBd}
\leanref{Development/WeilBound/ThetaLifts/LeviDeterminants/GrowHt.lean}{PositivityTheory.ThKF0.Grow.htBd}
\leanref{Development/WeilBound/ThetaLifts/LeviDeterminants/GrowDef.lean}{PositivityTheory.ThKF0.Grow.FinInv}
\leanref{Development/WeilBound/ThetaLifts/LeviDeterminants/GrowFin.lean}{PositivityTheory.ThKF0.Grow.finInv}
\leanref{Development/WeilBound/ThetaLifts/LeviDeterminants/GrowFin.lean}{PositivityTheory.ThKF0.Grow.Fi.fix}
\leanref{Development/WeilBound/ThetaLifts/LeviDeterminants/GrowDef.lean}{PositivityTheory.ThKF0.Grow.impl_sOf}
\leanref{Development/WeilBound/ThetaLifts/LeviDeterminants/ThetaGrow.lean}{PositivityTheory.ThKF0.ArchBd}
\leanref{Development/WeilBound/ThetaLifts/LeviDeterminants/WeilRepresentationPointwiseArchimedeanDecayBound.lean}{PositivityTheory.ThKF0.weil_rep_rapid_decay_height}
\leanref{Development/WeilBound/ThetaLifts/LeviDeterminants/GrowFinal.lean}{PositivityTheory.ThKF0.Grow.weil_rep_rapid_decay_of_growth}
\leanref{Development/WeilBound/ThetaLifts/LeviDeterminants/GrowFinal.lean}{PositivityTheory.ThKF0.Grow.weighted}
\leanref{Development/WeilBound/ThetaLifts/LeviDeterminants/GrowFinal.lean}{PositivityTheory.ThKF0.Grow.core}
\leanref{Development/WeilBound/ThetaLifts/LeviDeterminants/GrowFinal.lean}{PositivityTheory.ThKF0.Grow.list_bound}
\begin{proof}
By condition (a) of a Weil representation
(Definition~\ref{asm-theta-agree-models}(6)),
$\omega(g,h)\rho(e(u'))\Theta=\rho(e(gu'h^{\mathsf T}))\omega(g,h)\Theta$ for
Schwartz-Bruhat $\Theta$; with $u'=g^{-1}u(h^{-1})^{\mathsf T}$ this reads
$\omega(g,h)\rho(s_{g,h}(w))\Theta=\rho(w)\,\omega(g,h)\Theta$.

(1) Let $\beta_1,\ldots,\beta_r$ and $C_1$ be given by
Lemma~\ref{asm-theta-agree-definite-inputs-growth-inputs-entries} for $E/F$,
and $m_0$ and $K$ by
Lemma~\ref{asm-theta-agree-definite-inputs-growth-inputs-bilinear} for this
basis, $\kappa=I$ and the isomorphism $e$; put $C=m_0+KC_1^2$. As $E$ is a CM
field, the height of the first lemma is $\operatorname{ht}$, and the integral
elements and $|\cdot|$ of the second are those above. Let $g\in U_X$ and
$h\in U_W$. As
$\operatorname{ht}(x^{-1})=\operatorname{ht}(x)$,
Lemma~\ref{asm-theta-agree-definite-inputs-growth-inputs-entries} for $g^{-1}$
and $h^{-1}$ gives integers $a$ and $b$ with $1\leq a\leq\operatorname{ht}(g)$
and $1\leq b\leq\operatorname{ht}(h)$ such that the coordinates of the entries
of $g^{-1}$ and of $h^{-1}$ satisfy the hypotheses of
Lemma~\ref{asm-theta-agree-definite-inputs-growth-inputs-bilinear} with
$x=C_1\operatorname{ht}(g)$ and $y=C_1\operatorname{ht}(h)$. As
$s_{g,h}(w)=e\bigl(g^{-1}e^{-1}(w)(h^{-1})^{\mathsf T}\bigr)$, that lemma gives
(1) with $m=abm_0$: $1\leq m\leq m_0\operatorname{ht}(g)\operatorname{ht}(h)$,
$m\,s_{g,h}(w)$ is integral for integral $w$, and
$|s_{g,h}(w)|\leq KC_1^2\operatorname{ht}(g)\operatorname{ht}(h)\,|w|$ for $w$
at $\infty$.

(2) First, for every Schwartz-Bruhat $\Theta$ on $\mathbf A^I$ there is
$m_0\geq1$ with $\rho(m_0w)\Theta=\Theta$ for every integral $w$. If this holds
for $\Theta_1$ and $\Theta_2$ with $m_1$ and $m_2$, it holds for
$\Theta_1+\Theta_2$ with $m_1m_2$, as $m_2w$ and $m_1w$ are integral; so we may
assume $\Theta=f\otimes\phi$ with $f\in\mathcal S(F_\infty^I)$ and $\phi$
Schwartz-Bruhat on $\mathbf A_f^I$. For integral $w$, with finite parts
$w'_1,w'_2$,
$(\rho(w)\Theta)(x)=f(x_\infty)\,\psi_f(x_f\cdot w'_2+\tfrac12w'_1\cdot w'_2)\,\phi(x_f+w'_1)$,
where $\psi_f$ is the restriction of $\psi$ to the finite adeles, a locally
constant character. As $\phi$ is locally constant with compact support, there
is a neighbourhood $U$ of $0$ in $\mathbf A_f^I\times\mathbf A_f^I$ such that
$\psi_f(y\cdot w'_2+\tfrac12w'_1\cdot w'_2)\phi(y+w'_1)=\phi(y)$ for all
$(w'_1,w'_2)\in U$ and all $y$: both sides vanish for $y$ outside a compact set,
on which local constancy gives a uniform neighbourhood. As the integral vectors
form a compact group and every neighbourhood of $0$ contains the multiples by
$N$ of all of them for some $N\geq1$, there is $m_0$ as required.

Now let $m_0$ be this integer for $\Phi$, and let $g$, $h$ and $m$ be as in (1).
For integral $w$, $m\,s_{g,h}(w)$ is integral, so by the first paragraph
\[
\rho(m_0mw)\,\omega(g,h)\Phi=\omega(g,h)\,\rho\bigl(m_0\,m\,s_{g,h}(w)\bigr)\Phi=\omega(g,h)\Phi .
\]

(3) Fix $N\geq0$. Let $C$ be as in (1), and $m_0$ as in (2) for $\Phi$. Use
Definition~\ref{asm-theta-agree-definite-inputs-growth-inputs-derivatives} for
$\kappa=I$ and the character $\psi$ of $B$, which is trivial on $F$ with
$\psi(z)=\exp(-2\pi i\operatorname{Tr}_{F_\infty/\mathbf R}z_\infty)$ for $z$ at
$\infty$ (Definition~\ref{gal-good-base}); its operators $\rho$, integral
elements and $|\cdot|$ are those above, and $V=F_\infty^I$. Let $\mu$, $S$,
$C_S$ and $k_S$ be given by
Lemma~\ref{asm-theta-agree-definite-inputs-growth-inputs-sobolev}, and $c$ by
Lemma~\ref{asm-theta-agree-definite-inputs-growth-inputs-trace} for
$\kappa=I$. For $l\in S$ let $|l|$ be its length and $R_l\geq1$ a bound for
the $|\cdot|$ of its terms.

Let $g\in U_X$ and $h\in U_W$, let $m$ be as in (1), and put
$y=C\operatorname{ht}(g)\operatorname{ht}(h)$, so that $1\leq m\leq y$, and
$\Theta=\omega(g,h)\Phi$. The operators $\omega(g,1)$ and $\omega(1,h)$ satisfy
condition (b) of Definition~\ref{asm-theta-agree-doubling}(6) before the
displayed rule (Definition~\ref{asm-theta-agree-models}(6)), so
$\Omega=\omega(g,h)=\omega(g,1)\omega(1,h)$ maps Schwartz-Bruhat functions to
Schwartz-Bruhat functions and preserves their norms $\|\cdot\|_{\mu'}$ for every
Haar measure $\mu'$. By the first paragraph,
$\Omega\rho(s_{g,h}(w))\Theta'=\rho(w)\Omega\Theta'$ for every Schwartz-Bruhat
$\Theta'$, and by (1), $s_{g,h}$ maps every $p\in V\times V$ to some
$p'\in V\times V$ with $|p'|\leq y|p|$. So
Lemma~\ref{asm-theta-agree-definite-inputs-growth-inputs-transport} applies to
$\Omega$ and $s=s_{g,h}$, with $y$ in place of $c$.

\emph{Step 1.} For every $j\geq0$ there are $C_j$ and $k_j$, independent of $g$,
$h$, $x$ and $b$, with $|D_{(0,b)}^j\Theta(x)|\leq C_jy^{k_j}$ for all $g$, $h$
and $x$ and every $b\in V$ with $|b|\leq1$; here $D_{(0,b)}^j$ is the $j$-fold
composite of $D_{(0,b)}$. Indeed, let $u\in\mathbf A^I$ be integral. Then
$w=(u,0)$ is integral, and by (2),
$\Theta(x+m_0mu)=(\rho(m_0mw)\Theta)(x)=\Theta(x)$. So $\Theta$ is
$m_0m$-periodic, and so is the Schwartz-Bruhat function
$\Theta_j=D^j_{(0,b)}\Theta$, by
Lemma~\ref{asm-theta-agree-definite-inputs-growth-inputs-derivatives-bounds}(1).
By Lemma~\ref{asm-theta-agree-definite-inputs-growth-inputs-sobolev},
$|\Theta_j(x)|^2\leq C_S(m_0m)^{k_S}\sum_{l\in S}\|D_l\Theta_j\|_\mu^2$. Let
$l\in S$. The sequence $l$ followed by $j$ terms $(0,b)$ has terms of
$|\cdot|\leq R_l$, so by
Lemma~\ref{asm-theta-agree-definite-inputs-growth-inputs-transport}(2),
$D_l\Theta_j=\Omega(D_{l'}\Phi)$ for a sequence $l'$ of length $|l|+j$ with
terms of $|\cdot|\leq yR_l$. As $\Omega$ preserves norms,
Lemma~\ref{asm-theta-agree-definite-inputs-growth-inputs-derivatives-bounds}(3),
for $\Phi$, $\mu$ and the length $|l|+j$, gives $A_l$, independent of $g$, $h$
and $b$, with
$\|D_l\Theta_j\|_\mu^2=\|D_{l'}\Phi\|_\mu^2\leq A_l(yR_l)^{2(|l|+j)}$. As
$m\leq y$ and $y\geq1$, $|\Theta_j(x)|^2\leq Ay^{k'}$, with
$A=C_Sm_0^{k_S}\sum_{l\in S}A_lR_l^{2(|l|+j)}$ and
$k'=k_S+2\max_{l\in S}(|l|+j)$. So
$|\Theta_j(x)|\leq1+|\Theta_j(x)|^2\leq(1+A)y^{k'}$.

\emph{Step 2.} Let $x\in\mathbf A^I$, and choose $b\in V$ with $|b|\leq1$ and
$|x_\infty|\leq c\,|T(x_\infty,b)|$. By
Lemma~\ref{asm-theta-agree-definite-inputs-growth-inputs-derivatives-bounds}(1),
$D^N_{(0,b)}\Theta(x)=(-2\pi i\,T(x_\infty,b))^N\Theta(x)$. As $2\pi\geq1$ and by
Step 1,
$|\Theta(x)|\,|x_\infty|^N\leq c^N|D^N_{(0,b)}\Theta(x)|\leq c^NC_Ny^{k_N}$. Put
$e_N=2^{N-1}$ if $N\geq1$ and $e_0=1$. Then
$(1+|x_\infty|)^N\leq e_N(1+|x_\infty|^N)$: for $N\geq1$ by the convexity of
$t\mapsto t^N$ on $[0,\infty)$, and for $N=0$ as $1\leq2$. With Step 1 for
$j=0$, and as $y^{k_0}$ and $y^{k_N}$ are at most $y^{k_0+k_N}$,
\[
|\Theta(x)|\,(1+|x_\infty|)^N\leq e_N\bigl(C_0+c^NC_N\bigr)y^{k_0+k_N}
=e_N\bigl(C_0+c^NC_N\bigr)C^{k_0+k_N}\bigl(\operatorname{ht}(g)\operatorname{ht}(h)\bigr)^{k_0+k_N},
\]
which is (3).
\end{proof}
% Lean: the first paragraph is impl_sOf. (1) is Grow.htBd (GrowHt.lean): exists_entries for
% g^-1 and h^-1 (heightGL_inv), bd_bilin for M.co.e, and the constant m_0 + K C C. (2) is
% finInv with Fi.fix (fix_pure, eventually_fix, exists_pi, integralPoint_natCast_smul). (3) is
% weil_rep_rapid_decay_height = Grow.weil_rep_rapid_decay_of_growth htBd finInv (sob_of
% derivOK) trBd (GrowFinal.lean), for T = om (g, h) and s = sOf M g h with Bd s m y: Step 1
% is core (FinInv for the periodicity, word_lev, Sob, and list_bound through Impl.word,
% hT.unit and word_bound), Step 2 is weighted (TrBd, word_mult, add_pow_le, whose constant
% 2^(N-1) uses truncated subtraction, so e_0 = 1, with the constant 2^(N-1) (C_0 + c^N C_1) and
% exponent k_0 + k_1). Lemma asm-theta-agree-decay-full gives the bound only for Phi in So(M).

\begin{lemma}\label{asm-theta-agree-definite-inputs}
Let $B$ be a theta base for a CM field $E$
(Definition~\ref{asm-theta-agree-models}(1)), and fix a uniformizer
$\varpi_v$ of $F_v$ for every finite place $v$.
\begin{enumerate}
\item Let $X$ be a Hermitian and $W$ a skew-Hermitian lattice space for $B$, of
  some signatures, and $M$ a model for $(X,W)$. Then the polynomial Fock space
  $\operatorname{Pol}(M)$ (Definition~\ref{asm-theta-agree-models}(4)) has
  countable dimension.
\item Let $(X,W,M,\chi_X,\chi_W,\omega)$ be a Weil pair for $B$
  (Definition~\ref{asm-theta-agree-triple}(1)), with $X$ and $W$ of ranks $n$
  and $d$. For every Schwartz-Bruhat function $\Phi$ on $\mathbf A^{n\times d}$
  there are $C$ and $k$ with
  $|\theta_\omega(\Phi)(g,h)|\leq C\bigl(\operatorname{ht}(g)\operatorname{ht}(h)\bigr)^k$
  for all $g\in U_X$ and $h\in U_W$.
\item Let $(X,W,M,\chi_X,\chi_W,\omega)$ and $(X',W,M',\chi'_X,\chi'_W,\omega')$
  be Weil pairs for $B$ with the same space $W$, with $X$, $X'$ and $W$ of
  ranks $n$, $n'$ and $d=n+n'$, and with $\chi_W=\chi_X/\chi'_X$. Let $v$ be a
  finite place and $\tau\colon E\to F_v$ an $F$-embedding with
  $(\chi_W)_\tau(\varpi_v)\neq1$
  (Definition~\ref{asm-theta-agree-coinvariants-local}(4)). Then every
  $h_0\in GL_d(F_v)$ is $\pi_\tau(h)$ for some $h\in U_W^{(v)}$
  (Definition~\ref{asm-theta-agree-models}(1) and~(2)). Let
  $L(\Phi',\Psi)\in\mathbf C$ be given for all functions $\Phi'$ on
  $\mathbf A^{n'\times d}$ and $\Psi$ on $\mathbf A^{n\times d}$, linear in
  $\Phi'$ and conjugate-linear in $\Psi$, such that for every
  $h_0\in GL_d(F_v)$ some $h\in U_W^{(v)}$ with $\pi_\tau(h)=h_0$ has
  $L(\omega'(1,h)\Phi',\omega(1,h)\Psi)=L(\Phi',\Psi)$ for all
  $\Phi'\in\operatorname{PD}(M')$ and Schwartz-Bruhat $\Psi$. Then
  $L(\Phi',\omega(u,1)\Psi)=L(\Phi',\Psi)$ for all $u\in U_X^{(v)}$,
  $\Phi'\in\operatorname{PD}(M')$ and Schwartz-Bruhat $\Psi$.
\item Let $X$ be a Hermitian lattice space for $B$ of rank $n\geq2$ and some
  signature, and $T$ a finite set of finite places. Let $f$ be a continuous
  function on $U_X$ with $f(\gamma g)=f(g)$ for all $\gamma\in U_X(F)$ and
  $g\in U_X$, and with $f(gu)=f(g)$ for all $g\in U_X$, every finite place
  $v\notin T$ such that the rational prime below $v$ splits completely in $E$,
  and every $u\in U_X^{(v)}$. Then $f$ is constant. The hypothesis $n\geq2$ is
  not used in the proof.
\end{enumerate}
Let now $(B,G,\mathcal X,\omega,\omega',R,R_f)$ be a matched quadruple of ranks
$(a,a')$ (Definition~\ref{asm-theta-agree-matched}), with spaces $V$, $V'$ and
$W$, models $M$ and $M'$, $\nu=\alpha/\beta$, and pairs $\mathcal P$
(Definition~\ref{asm-theta-agree-definite-products}(3)).
\begin{enumerate}
\setcounter{enumi}{4}
\item If $\Phi$ and $0$ correspond, then $I(\Phi)=0$. If
  $\Phi'\in\mathcal S^\circ(M')$ and $0$ and $\Phi'$ correspond, then
  $I'(\Phi')=0$.
\item Let $v$ be a finite place such that $E$ has an $F$-embedding into $F_v$,
  and let $h\in U_W^{(v)}$. If $\Phi$ and $\Phi'$ correspond, then so do
  $\omega(1,h)\Phi$ and $\omega'(1,h)\Phi'$.
\item If $v$ is clean for $\mathcal X$, then $E$ has an $F$-embedding into
  $F_v$, and $\nu_\tau(\varpi_v)\neq1$ for every $F$-embedding
  $\tau\colon E\to F_v$.
\end{enumerate}
Suppose further that $p=q=0$.
\begin{enumerate}
\setcounter{enumi}{7}
\item If $y,y'\in\mathcal P$ and $z\in\mathbf C$, then $y+zy'\in\mathcal P$;
  and if $(\Phi,\Phi')\in\mathcal P$ and $h\in U_W$, then
  $(\omega(1,h)\Phi,\omega'(1,h)\Phi')\in\mathcal P$.
\item Let $N\subset\mathcal P$ be a linear subspace that contains every
  $(\Phi,\Phi')\in\mathcal P$ such that $\Phi$ and $0$ correspond, and contains
  $(\omega(1,h)\Phi,\omega'(1,h)\Phi')$ whenever $(\Phi,\Phi')\in N$ and
  $h\in U_W$. Then either $\Phi$ and $0$ correspond for every
  $(\Phi,\Phi')\in N$, or $N=\mathcal P$.
\item Assume the meromorphic continuation of Siegel Eisenstein series in the
  form of Theorem~\ref{asm-theta-agree-eisenstein} and the Siegel-Weil formula
  in the form of Theorem~\ref{asm-theta-agree-siegel-weil}. Let $\mu$ be a Haar
  measure on $U_W$ and $\mathcal D$ a fundamental domain for $U_W(F)$ with
  respect to $\mu$, let $\mu'$ and $\mathcal D'$ be the measure and the domain
  of $M'$, and let $\Psi\in\operatorname{PD}(M')$. Then there are a connected
  open set $U\subset\mathbf C$ containing $z_0=-a/2$, a real number $t_0$,
  functions $F_z$ on $U_W$ for $\operatorname{Re}z>t_0$, and for every rapidly
  decreasing function $\phi$ on $U_W$ a holomorphic function $Z_\phi$ on $U$,
  such that: for $\operatorname{Re}z>t_0$, $z\in U$, $F_z$ is integrable on
  $U_W$ with respect to $\mu$, and
  \[
  Z_\phi(z)=\int_{U_W}F_z(g)\,\langle R_g\phi,\phi\rangle_{\mathcal D}\,d\mu(g);
  \]
  and
  \[
  Z_\phi(z_0)=\frac{2}{\mu'(\mathcal D')}\,\langle\theta^*_{\omega'}(\phi,\Psi),\theta^*_{\omega'}(\phi,\Psi)\rangle_{\mathcal D'},
  \]
  the product on $U_{V'}$ taken with $\mu'$ and $\mathcal D'$
  (Definition~\ref{asm-theta-agree-definite-products}).
\end{enumerate}
\end{lemma}
\leanref{Development/WeilBound/ThetaLifts/LatticeBasics/RankPolData.lean}{GlobalTheory.ThK.Dbl.PolRank}
\leanref{Development/AbelianVarietyTheory/Definitions/Divisors/CTSmallCone.lean}{GlobalTheory.ThK.ThetaBasic}
\leanref{Development/WeilBound/ThetaLifts/LeviDeterminants/ThetaGrow.lean}{PositivityTheory.ThKF0.theta_moderate_growth_of_decay}
\leanref{Development/WeilBound/ThetaLifts/LeviDeterminants/WeilRepRationalSupportCommonDenominator.lean}{PositivityTheory.ThKF0.weil_rep_bounded_denominator}
\leanref{Development/WeilBound/ThetaLifts/LeviDeterminants/WeilRepRationalSupportCommonDenominator.lean}{PositivityTheory.ThKF0.exists_dual_den}
\leanref{Development/WeilBound/ThetaLifts/LatticeBasics/ThConstInv.lean}{GlobalTheory.ThK.Dbl.TC.inv_loc}
\leanref{Development/WeilBound/ThetaLifts/LatticeBasics/PairsStable.lean}{GlobalTheory.ThK.Dbl.ConstOfGood}
\leanref{Development/WeilBound/CoreStatements/AbelianHeightsDuality/UnitaryDensity/EqTopOfLines.lean}{PositivityTheory.SA.constOfGood}
\leanref{Development/WeilBound/CoreStatements/AbelianHeightsDuality/UnitaryDensity/EqTopOfLines.lean}{PositivityTheory.SA.mem_of_good}
\leanref{Development/WeilBound/ThetaLifts/LatticeBasics/ConstOfGood.lean}{GlobalTheory.ThK.Dbl.conj_mem_UAt}
\leanref{Development/AbelianVarietyTheory/Definitions/Divisors/CTSmallCone.lean}{GlobalTheory.ThK.Factor}
\leanref{Development/WeilBound/ThetaLifts/LatticeBasics/ThConstMatch.lean}{GlobalTheory.ThK.Dbl.TC.pm_fin}
\leanref{Development/WeilBound/ThetaLifts/LatticeBasics/PairsStable.lean}{GlobalTheory.ThK.Dbl.GoodNe}
\leanref{Development/WeilBound/ThetaLifts/LatticeBasics/PairsStable.lean}{GlobalTheory.ThK.Dbl.PairsStable}
\leanref{Development/WeilBound/ThetaLifts/LatticeBasics/PairsStable.lean}{GlobalTheory.ThK.Dbl.PairsIrr}
\leanref{Development/WeilBound/ThetaLifts/LatticeBasics/PairsB.lean}{GlobalTheory.ThK.Dbl.pairsIrr_of}
\leanref{Development/WeilBound/ThetaLifts/LatticeBasics/PairsL.lean}{GlobalTheory.ThK.Dbl.Prs.polMatch_iff}
\leanref{Development/WeilBound/ThetaLifts/LatticeBasics/ContinuousFlat.lean}{GlobalTheory.ThK.Dbl.kpair_cont}
\leanref{Development/WeilBound/ThetaLifts/LatticeBasics/PairsStable.lean}{GlobalTheory.ThK.Dbl.SWFull}
\leanref{Development/WeilBound/ThetaLifts/LatticeBasics/HermRank.lean}{GlobalTheory.ThK.Dbl.polRank}
\leanref{Development/WeilBound/ThetaLifts/LatticeBasics/HermRank.lean}{GlobalTheory.ThK.Dbl.Herm.exists_vac}
\leanref{Development/WeilBound/ThetaLifts/LatticeBasics/HermRank.lean}{GlobalTheory.ThK.Dbl.Herm.exists_stable}
\leanref{Development/WeilBound/ThetaLifts/LatticeBasics/ThConstInv.lean}{GlobalTheory.ThK.Dbl.TC.SplitModel}
\leanref{Development/WeilBound/ThetaLifts/LatticeBasics/ThConstSplit.lean}{GlobalTheory.ThK.Dbl.TC.splitModel}
\leanref{Development/WeilBound/ThetaLifts/LatticeBasics/ThConstMatch.lean}{GlobalTheory.ThK.Dbl.TC.loc_op}
\leanref{Development/WeilBound/ThetaLifts/LatticeBasics/GoodNe.lean}{GlobalTheory.ThK.Dbl.goodNe}
\leanref{Development/WeilBound/ThetaLifts/LatticeBasics/GoodNe.lean}{GlobalTheory.ThK.Dbl.nu_pow_ne_one}
\leanref{Development/WeilBound/ThetaLifts/LatticeBasics/GoodNe.lean}{GlobalTheory.ThK.Dbl.nonempty_emb_of_splitsFully}
\leanref{Development/WeilBound/ThetaLifts/LatticeBasics/PairsFin.lean}{GlobalTheory.ThK.Dbl.pairsStable}
\leanref{Development/WeilBound/ThetaLifts/LatticeBasics/PairsB.lean}{GlobalTheory.ThK.Dbl.pairsStable_of}
\leanref{Development/WeilBound/ThetaLifts/LatticeBasics/PairsB.lean}{GlobalTheory.ThK.Dbl.Prs.cartan_one}
\begin{proof}
(1) Let $n$ and $d$ be the ranks of $X$ and $W$, put
$I=\{1,\ldots,n\}\times\{1,\ldots,d\}$ and $V=F_\infty^I$, let $T$ and
$D_{a,b}$ be as in Lemma~\ref{asm-theta-agree-density-derivative}, put
$D_p=D_{p_1,p_2}$ for $p=(p_1,p_2)\in V\times V$, and let $s$, $\Omega$, $H_w$
and $\kappa_w$ be as in the proof of
Lemma~\ref{asm-theta-agree-density-gaussian}. For a family
$\varepsilon=(\varepsilon_w)$ of signs, indexed by the infinite places of $F$,
let $J_\varepsilon$ be the $\mathbf R$-linear map of $V\times V$ with
$\pi_w(s(J_\varepsilon p))=\varepsilon_w\kappa_w(\pi_w(s(p)))$ for every $w$;
it exists because $p\mapsto(\pi_w(s(p)))_w$ is an $\mathbf R$-linear bijection
onto $\prod_wM_{n\times d}(\mathbf C)$. Let $e_\varepsilon\in F_\infty$ have
components $\varepsilon_w$, and let $G_\varepsilon$ be the space of
$f\in\mathcal S(V)$ with $D_{J_\varepsilon p}f=i\,D_pf$ for every
$p\in V\times V$.

\emph{Gaussians.} A Gaussian $f_0$ with signs $\varepsilon$
(Definition~\ref{asm-theta-agree-models}(4)) lies in $G_\varepsilon$. Indeed,
for $p\in V\times V$ the elements $u=s(p)$ and $u'=s(J_\varepsilon p)$ are at
$\infty$, with
$\pi_w(u')=\varepsilon_wi\,\theta_{X,w}\pi_w(u)\theta_{W,w}^{\mathsf T}$, and by
Lemma~\ref{asm-theta-agree-density-derivative} the families $t\mapsto\rho(tu)$
and $t\mapsto\rho(tu')$ act at $\infty$ to first order by $f_0\mapsto D_pf_0$
and $f_0\mapsto D_{J_\varepsilon p}f_0$. So $D_{J_\varepsilon p}f_0=i\,D_pf_0$.

\emph{Each $G_\varepsilon$ has dimension at most $1$.} For $b\in V$ let $K(b)$
be the first component of $J_\varepsilon(0,b)$. Let $p\neq0$ and
$x_w=\pi_w(s(p))$. As $e$ is $\mathbf A$-linear, $s(e_\varepsilon p)=e_\varepsilon s(p)$,
and by (c) in the proof of Lemma~\ref{asm-theta-agree-density-gaussian},
\[
\Omega(e_\varepsilon p,J_\varepsilon p)=-2\sum_w\operatorname{Re}H_w(\varepsilon_wx_w,\varepsilon_w\kappa_wx_w)=-2\sum_w\operatorname{Re}H_w(x_w,\kappa_wx_w)>0 .
\]
For $p=(0,b)$ the left side is $-T(K(b),e_\varepsilon b)$; so $K$ is
injective, hence bijective. Let $f\in G_\varepsilon$ with $f(0)=0$, and let
$x\in V$. Choose $b$ with $K(b)=x$, and write $J_\varepsilon(0,b)=(x,q)$. The
relation $D_{(x,q)}f=i\,D_{(0,b)}f$ says that $\partial_xf(y)=c(y)f(y)$ with
$c(y)=2\pi i\,T(y,q)+2\pi\,T(y,b)$, which is linear in $y$. So
$\varphi(t)=f(tx)$ has $\varphi'(t)=t\,c(x)\varphi(t)$, the function
$\varphi(t)e^{-c(x)t^2/2}$ is constant, and $f(x)=\varphi(0)e^{c(x)/2}=0$. So
$f\mapsto f(0)$ is injective on $G_\varepsilon$.

\emph{Countable dimension.} There are finitely many families of signs, so the
sum $G$ of the spaces $G_\varepsilon$ is finite-dimensional. Let
$p_1,\ldots,p_m$ be an $\mathbf R$-basis of $V\times V$, and let $N$ be the
span of the functions $D_{p_{i_1}}\cdots D_{p_{i_k}}f$ with $k\geq0$,
$i_1,\ldots,i_k\in\{1,\ldots,m\}$ and $f\in G$. It has countable dimension, and
$D_pN\subset N$ for every $p$, as $D_p$ is $\mathbf R$-linear in $p$. It
contains every Gaussian; and if $f\in N$ and $t\mapsto\rho(tu)$, for some $u$
at $\infty$, acts at $\infty$ to first order by $f\mapsto f_1$, then
$f_1=D_pf\in N$ by Lemma~\ref{asm-theta-agree-density-derivative}, with $p$ the
components at $\infty$ of $e_1(u)$ and $e_2(u)$. So
$\operatorname{Pol}(M)\subset N$, and $\operatorname{Pol}(M)$ has countable
dimension.

(2) Let $\psi$ be the character of $B$, let $\rho$, the integral elements and
$|\cdot|$ be as in Lemma~\ref{asm-theta-agree-definite-inputs-growth-inputs},
let $N_0$, $k_2$ and $C_2$ be given by
Lemma~\ref{asm-theta-agree-definite-inputs-growth-lattice} for $\kappa=I$, let
$C_1$ and $k_1$ be given by part (3) of
Lemma~\ref{asm-theta-agree-definite-inputs-growth-inputs} for $N=N_0$, and let
$C$ and $m_0$ be given by its parts (1) and (2). Let $D=|\operatorname{disc}(F)|$.
Then $D\xi\in\mathcal O_F$ for every $\xi\in F$ with
$\operatorname{Tr}_{F/\mathbf Q}(\xi\mathcal O_F)\subset\mathbf Z$: such $\xi$ form
the inverse of the different of $F$, and $D$, the absolute norm of the
different, lies in it.

Let $g\in U_X$ and $h\in U_W$, put $H=\operatorname{ht}(g)\operatorname{ht}(h)$
and $\Theta=\omega(g,h)\Phi$, and let $m\leq CH$ be as in part (1) of that
lemma. Let $\xi\in F^I$ with $\Theta(\xi)\neq0$, let $i\in I$ and
$y\in\mathcal O_F$, and let $y_f\in\mathbf A$ have components $0$ at the
infinite places and $y$ at the finite places. The element $w$ with first
component $0$ and second component $y_f$ at $i$ and $0$ elsewhere is integral,
so by part (2), $\rho(m_0mw)\Theta=\Theta$; at $\xi$ this reads
$\psi(m_0m\,\xi_iy_f)\Theta(\xi)=\Theta(\xi)$. As $\psi$ is trivial on $F$ and
$\psi(x)=\exp(-2\pi i\operatorname{Tr}x)$ on $F_\infty$
(Definition~\ref{gal-good-base}), $\psi(a_f)=\exp(2\pi i\operatorname{Tr}_{F/\mathbf Q}(a))$
for $a\in F$, with $a_f$ defined like $y_f$. So
$\operatorname{Tr}_{F/\mathbf Q}(m_0m\,\xi_iy)\in\mathbf Z$ for every
$y\in\mathcal O_F$, and $Dm_0m\,\xi_i\in\mathcal O_F$. Hence $m'=Dm_0m$ has
$m'\xi\in\mathcal O_F^I$ whenever $\xi\in F^I$ and $\Theta(\xi)\neq0$, and
$m'\leq Dm_0CH$.

By part (3) of Lemma~\ref{asm-theta-agree-definite-inputs-growth-inputs},
$|\Theta(\xi)|(1+|\xi|)^{N_0}\leq C_1H^{k_1}$ for $\xi\in F^I$. So
Lemma~\ref{asm-theta-agree-definite-inputs-growth-lattice}, with
$A=C_1H^{k_1}$, gives
\[
|\theta_\omega(\Phi)(g,h)|=\Bigl|\sum_{\xi\in F^I}\Theta(\xi)\Bigr|\leq C_2C_1H^{k_1}(Dm_0CH)^{k_2},
\]
which is a bound of the required form.

(3) The first claim is part (2) of
Lemma~\ref{asm-theta-agree-coinvariants-dictionary}. Put $k=F_v$ and
$\varpi=\varpi_v$, let $|\cdot|$ be the normalized absolute value of $k$, and
put $\nu=(\chi_W)_\tau$, $\alpha=(\chi_X)_\tau$, $\nu'=(\chi'_W)_\tau$ and
$\alpha'=(\chi'_X)_\tau$. These characters have values of absolute value $1$,
and $\nu=\alpha/\alpha'$. By parts (1) and (3) of
Lemma~\ref{asm-theta-agree-coinvariants-dictionary}, there are linear models
$\iota$ and $\iota'$ of the second kind for $(\tau,1,1)$, for $M$ and for $M'$,
through which $\omega(g,1)$ and $\omega(1,h)$ read as
$\omega_{\nu,\alpha}(\pi_\tau(g),1)$ and $\omega_{\nu,\alpha}(1,\pi_\tau(h))$,
and $\omega'(1,h)$ reads as $\omega_{\nu',\alpha'}(1,\pi_\tau(h))$, for
$g\in U_X^{(v)}$ and $h\in U_W^{(v)}$
(Definitions~\ref{asm-theta-agree-coinvariants-local}(1)
and~\ref{asm-theta-agree-finite-split-models}(3)). Let $j$ and $j'$ be the
inverses of $\iota$ and $\iota'$ on the Schwartz-Bruhat functions. For a
function $a$ on a space of matrices over $k$ write $a^h(Y)=a(Yh)$ and
${}^ga(Y)=a(gY)$. Then
\[
\omega_{\nu,\alpha}(1,h_0)a=\alpha(\det h_0)\,|\det h_0|^{n/2}a^{h_0},\qquad
\omega_{\nu,\alpha}(g_0,1)a=\nu(\det g_0)\,|\det g_0|^{d/2}\,{}^{g_0^{\mathsf T}}a,
\]
and $\omega_{\nu',\alpha'}(1,h_0)b=\alpha'(\det h_0)\,|\det h_0|^{n'/2}b^{h_0}$.

\emph{Pure tensors.} Both sides of the claim are linear in $\Phi'$ and
conjugate-linear in $\Psi$, and $\operatorname{PD}(M')$ and the Schwartz-Bruhat
functions are spanned by functions of the following forms. So we may assume
that $\Phi'=f'\otimes\bigotimes_w\phi'_w$ with $f'\in\operatorname{Pol}(M')$,
every $\phi'_w$ Schwartz-Bruhat and $\phi'_w=\mathrm{vac}'_w$ for almost all
$w$, and that $\Psi=f\otimes\bigotimes_w\phi_w$ with
$f\in\mathcal S(F_\infty^{n\times d})$, every $\phi_w$ Schwartz-Bruhat and
$\phi_w$ the indicator function of $\mathcal O_w^{n\times d}$ for almost all
$w$. For a function $\psi$ on $F_v^{n\times d}$ let $\Psi[\psi]$ be $\Psi$ with
$\phi_v$ replaced by $\psi$, and define $\Phi'[\psi']$ from $\Phi'$ in the same
way. Then $\Psi[\psi]=\psi\otimes r$ with $r$ not depending on the
$v$-component. So, for Schwartz-Bruhat $\psi$ and $\psi'$, $\Psi[\psi]$ is
Schwartz-Bruhat, and $\Phi'[\psi']$ lies in $\operatorname{PD}(M')$ and is
Schwartz-Bruhat, as $M'$ is regular; and an operator $S$ that reads as $S_0$ through $\iota$ has
$S(\Psi[\psi])=\Psi[jS_0\iota\psi]$, by
Definition~\ref{asm-theta-agree-models}(2), and likewise for $\Phi'$ and
$\iota'$.

\emph{A bilinear form.} For Schwartz-Bruhat $a$ on $M_{n\times d}(k)$ and $b$
on $M_{n'\times d}(k)$ put $Q(a,b)=L(\Phi'[j'b],\Psi[j\bar a])$. It is
bilinear, as $L$ is conjugate-linear in its second argument. Let
$\xi(x)=\nu(x)|x|^{-d/2}$. Let $h_0\in GL_d(k)$, and let $h$ be as in the
hypothesis. As $\omega(1,h)$ and $\omega'(1,h)$ read as
$\omega_{\nu,\alpha}(1,h_0)$ and $\omega_{\nu',\alpha'}(1,h_0)$,
\[
Q(a,b)=L\bigl(\omega'(1,h)\Phi'[j'b],\omega(1,h)\Psi[j\bar a]\bigr)
=\alpha'(\det h_0)\,\overline{\alpha(\det h_0)}\,|\det h_0|^{(n+n')/2}\,Q(a^{h_0},b^{h_0}),
\]
and as $|\alpha|=1$ and $\nu=\alpha/\alpha'$,
$Q(a^{h_0},b^{h_0})=\xi(\det h_0)\,Q(a,b)$. If $\xi(\varpi)$ were a real number
$t>0$, then $\nu(\varpi)=t|\varpi|^{d/2}$ would be a positive real number of
absolute value $1$, so $\nu(\varpi)=1$, against the hypothesis. By
Lemma~\ref{asm-theta-agree-definite-dist-bilinear},
$Q({}^ga,b)=\xi(\det g)\,Q(a,b)$ for $g\in GL_n(k)$.

\emph{Conclusion.} Let $u\in U_X^{(v)}$, and put $g_0=\pi_\tau(u)$,
$a=\overline{\iota\phi_v}$ and $b=\iota'\phi'_v$, so that $\Psi=\Psi[j\bar a]$
and $\Phi'=\Phi'[j'b]$. Then
\[
\omega(u,1)\Psi=\Psi\bigl[j\,\omega_{\nu,\alpha}(g_0,1)\iota\phi_v\bigr]=\nu(\det g_0)\,|\det g_0|^{d/2}\,\Psi\bigl[j\,{}^{g_0^{\mathsf T}}\bar a\bigr],
\]
so
\[
L(\Phi',\omega(u,1)\Psi)=\overline{\nu(\det g_0)}\,|\det g_0|^{d/2}\,Q({}^{g_0^{\mathsf T}}a,b)
=\nu(\det g_0)^{-1}|\det g_0|^{d/2}\,\xi(\det g_0)\,Q(a,b)=L(\Phi',\Psi).
\]

(4) Let $c$ be the complex conjugation of $E$ and $J=J_X$; then
$c(J)^{\mathsf T}=J$, and $J$ is invertible, as $\tau_w(J)$ has a signature.
Let $\mathbf A_E^1$ be as in
Lemma~\ref{asm-theta-agree-definite-inputs-constancy-norm}. Let $L$ be the set
of $x\in U_X$ with $f(xg)=f(g)$ for every $g\in U_X$. It is a subgroup of
$U_X$, closed as $f$ is continuous, and it contains $U_X(F)$ by the invariance
of $f$ on the left. Let $v\notin T$ be a finite place such that the rational
prime below $v$ splits completely in $E$, and let $u\in U_X^{(v)}$. For
$g\in U_X$, $g^{-1}ug-1=g^{-1}(u-1)g$ is at $v$, so $g^{-1}ug\in U_X^{(v)}$ and
$f(ug)=f(g\cdot g^{-1}ug)=f(g)$. So $U_X^{(v)}\subset L$.

Let $R\in M_n(E)$ with $R^2=R$ and $R^{\mathsf T}J=J\,c(R)$, and let $S_R$ be
the continuous homomorphism $z\mapsto1+(z-1)R$ from $\mathbf A_E^1$ to $U_X$ of
Lemma~\ref{asm-theta-agree-definite-inputs-constancy-lines}(1). The closed
subgroup $S_R^{-1}(L)$ of $\mathbf A_E^1$ contains $E^\times\cap\mathbf A_E^1$,
which $S_R$ maps into $U_X(F)$; and, for $v$ as above, it contains every
$z\in\mathbf A_E^1$ with components $1$ away from $v$, as then $S_R(z)-1=(z-1)R$
is at $v$. As $E$ is Galois over $\mathbf Q$,
Lemma~\ref{asm-theta-agree-definite-inputs-constancy-norm} gives
$S_R^{-1}(L)=\mathbf A_E^1$. So $L$ contains $1+(z-1)R$ for every such $R$ and
every $z\in\mathbf A_E^1$, and $L=U_X$ by
Lemma~\ref{asm-theta-agree-definite-inputs-constancy-lines}(2). Hence
$f(g)=f(g\cdot1)=f(1)$ for every $g\in U_X$.

(5) By Lemma~\ref{asm-theta-agree-decay}, every Schwartz-Bruhat function
converges for $(M,\omega)$ and every element of $\mathcal S^\circ(M')$
converges for $(M',\omega')$. So, as in the proof of
Lemma~\ref{asm-theta-agree-nonvanishing-schwartz}, $I$ is linear on the
Schwartz-Bruhat functions and $I'$ on $\mathcal S^\circ(M')$, by
Lemma~\ref{asm-theta-agree-convergence-theta}(2). Extend them to linear maps
$J$ and $J'$ on all functions. For a function $f$ at $\infty$ and a family
$\phi=(\phi_v)$ of functions at the finite places write
$f\otimes\phi=f\otimes\bigotimes_v\phi_v$.

\emph{The first claim.} Let $\Phi$ and $0$ correspond, with a set $T_1$ and
sums $\Phi=\sum_if_i\otimes\phi_i$ and $0=\sum_jf'_j\otimes\phi'_j$ as in
Definition~\ref{asm-theta-agree-matched}(4). By
Lemma~\ref{asm-theta-agree-definite-inputs-tensor}(1), for $w_v=\mathrm{vac}'_v$
and the empty sum on the other side,
$\sum_jf'_j\otimes\bigotimes_{v\in T_1}\phi'_{j,v}=0$. Applying $p'$ and the
$p'_v$, the right side of the identity of
Definition~\ref{asm-theta-agree-matched}(4) vanishes, so
$\sum_ip(f_i)\otimes\bigotimes_{v\in T_1}p_v(\phi_{i,v})=0$. For
$f\in\operatorname{Pol}(M)$ and Schwartz-Bruhat $\psi_v$, $v\in T_1$, put
$\Lambda(f,\psi)=J\bigl(f\otimes\bigotimes_{v\in T_1}\psi_v\otimes\bigotimes_{v\notin T_1}\mathrm{vac}_v\bigr)$.
This function is Schwartz-Bruhat, as $M$ is regular, so $\Lambda(f,\psi)$ is
its theta integral, which is $0$ if $p(f)=0$ or $p_v(\psi_v)=0$ for some $v$,
by Lemma~\ref{asm-theta-agree-local-injective-kills}(1).
Lemma~\ref{asm-theta-agree-definite-inputs-tensor}(2) gives a linear map $L$
with
$I(\Phi)=\sum_i\Lambda(f_i,(\phi_{i,v})_{v\in T_1})=L\bigl(\sum_ip(f_i)\otimes\bigotimes_{v\in T_1}p_v(\phi_{i,v})\bigr)=0$.

\emph{The second claim.} Let $\Phi'\in\mathcal S^\circ(M')$, and let $0$ and
$\Phi'$ correspond, with $T_1$ and sums $0=\sum_if_i\otimes\phi_i$ and
$\Phi'=\sum_jf'_j\otimes\phi'_j$. As in the first claim, now with
$w_v=\mathrm{vac}_v$,
$\sum_jp'(f'_j)\otimes\bigotimes_{v\in T_1}p'_v(\phi'_{j,v})=0$. Let $N$ be the
span of the Schwartz-Bruhat functions on $F_{v_1}^{I'_0}$ of full rank at
$\tau_1$, and $\pi$ a linear map of the functions on $F_{v_1}^{I'_0}$ into $N$
that is the identity on $N$. As in the proof of
Lemma~\ref{asm-theta-agree-local-augmentation-slice}, there is a linear
operator $O$ on the functions on $\mathbf A^{I'_0}$ with
$O(f\otimes\phi)=f\otimes\tilde\phi$, where $\tilde\phi$ is $\phi$ with
$\phi_{v_1}$ replaced by $\pi(\phi_{v_1})$, whenever almost all $\phi_v$ are
the indicator functions of $\mathcal O_v^{I'_0}$. It is the identity on the
generators of $\mathcal S^\circ(M')$, hence on $\mathcal S^\circ(M')$, so
$\Phi'=O(\Phi')=\sum_jf'_j\otimes\tilde\phi'_j$. Moreover
$\sum_jp'(f'_j)\otimes\bigotimes_{v\in T_1}p'_v(\tilde\phi'_{j,v})=0$: if
$v_1\notin T_1$ the components at the places of $T_1$ are unchanged, and if
$v_1\in T_1$ both sums for $\Phi'$ have the components $\mathrm{vac}'_v$
outside $T_1$, and Lemma~\ref{asm-theta-agree-definite-inputs-tensor}(1)
applies. Put $w_v=\mathrm{vac}'_v$ for $v\neq v_1$ and
$w_{v_1}=\pi(\mathrm{vac}'_{v_1})$, and let $S_{v_1}=N$ and $S_v$ be the
Schwartz-Bruhat functions for $v\neq v_1$. For $f\in\operatorname{Pol}(M')$
and $\psi_v\in S_v$, $v\in T_1$, put
$\Lambda'(f,\psi)=J'\bigl(f\otimes\bigotimes_{v\in T_1}\psi_v\otimes\bigotimes_{v\notin T_1}w_v\bigr)$.
This function lies in $\mathcal S^\circ(M')$, as its component at $v_1$ lies
in $N$, so $\Lambda'(f,\psi)$ is its theta integral, which is $0$ if
$p'(f)=0$ or $p'_v(\psi_v)=0$ for some $v$, by
Lemma~\ref{asm-theta-agree-local-injective-kills}(2). The families
$\tilde\phi'_j$ have components in $S_v$ at $v\in T_1$ and $w_v$ outside $T_1$,
so Lemma~\ref{asm-theta-agree-definite-inputs-tensor}(2) gives a linear map
$L'$ with
$I'(\Phi')=\sum_j\Lambda'(f'_j,(\tilde\phi'_{j,v})_{v\in T_1})=L'\bigl(\sum_jp'(f'_j)\otimes\bigotimes_{v\in T_1}p'_v(\tilde\phi'_{j,v})\bigr)=0$.

(6) Let $\tau\colon E\to F_v$ be an $F$-embedding, and let $\Phi$ and $\Phi'$
correspond, with a set $T_1$ and sums
$\Phi=\sum_if_i\otimes\bigotimes_u\phi_{i,u}$ and
$\Phi'=\sum_jf'_j\otimes\bigotimes_u\phi'_{j,u}$ as in
Definition~\ref{asm-theta-agree-matched}(4). If $v\notin T_1$, then
$v\notin T$, so $p_v(\mathrm{vac}_v)=p'_v(\mathrm{vac}'_v)$, and tensoring the
identity there with this element shows that $T_1\cup\{v\}$ serves as well; so
we may assume $v\in T_1$. By parts (1) and (3) of
Lemma~\ref{asm-theta-agree-coinvariants-dictionary}, for the Weil pairs
$(V,W,M,\beta,\beta/\alpha,\omega)$ and $(V',W,M',\alpha,\alpha/\beta,\omega')$
and the embedding $\tau$, there are linear maps $\iota$ and $\iota'$, bijective between the
Schwartz-Bruhat functions, through which $\omega(1,h)$ and $\omega'(1,h)$ read
as operators $S$ and $S'$. For Schwartz-Bruhat $\phi$ and $\phi'$ put
$O\phi=\iota^{-1}S\iota\phi$ and $O'\phi'=\iota'^{-1}S'\iota'\phi'$; then
$O\phi$ is Schwartz-Bruhat and $\omega(1,h)$ acts at $v$ by $\phi\mapsto O\phi$,
and likewise for $O'$. As in the proof of (3),
$\omega(1,h)\Phi=\sum_if_i\otimes\bigotimes_u\tilde\phi_{i,u}$ with
$\tilde\phi_{i,v}=O\phi_{i,v}$ and $\tilde\phi_{i,u}=\phi_{i,u}$ for $u\neq v$,
and $\omega'(1,h)\Phi'=\sum_jf'_j\otimes\bigotimes_u\tilde\phi'_{j,u}$ with
$\tilde\phi'_{j,v}=O'\phi'_{j,v}$ and $\tilde\phi'_{j,u}=\phi'_{j,u}$ for
$u\neq v$.

By the conditions on matched finite maps
(Definition~\ref{asm-theta-agree-matched}(3)), $p_v(O\phi)=p'_v(O'\phi')$
whenever $\phi$ and $\phi'$ are Schwartz-Bruhat with $p_v(\phi)=p'_v(\phi')$.
There is a linear map $A$ of $\Pi_v$ with $A(p_v(\phi))=p_v(O\phi)$ and
$A(p'_v(\phi'))=p'_v(O'\phi')$ for all Schwartz-Bruhat $\phi$ and $\phi'$. Indeed,
as $p_v$ maps the Schwartz-Bruhat functions onto $\Pi_v$, there is a linear map
$\sigma$ from $\Pi_v$ to them with $p_v\circ\sigma=1$; put
$A=p_v\circ O\circ\sigma$. Then $A(p'_v(\phi'))=p'_v(O'\phi')$, as
$p_v(\sigma(p'_v(\phi')))=p'_v(\phi')$; and given $\phi$, choose $\phi'$ with
$p'_v(\phi')=p_v(\phi)$, which exists as $p'_v$ is onto, so that
$A(p_v(\phi))=p'_v(O'\phi')=p_v(O\phi)$. Applying to the identity of
Definition~\ref{asm-theta-agree-matched}(4) the map that is $A$ on the factor
$\Pi_v$ and the identity on the other factors gives that identity for the sums
for $\omega(1,h)\Phi$ and $\omega'(1,h)\Phi'$. Their local factors are
Schwartz-Bruhat and equal to $\mathrm{vac}_u$ and $\mathrm{vac}'_u$ for
$u\notin T_1$; so $\omega(1,h)\Phi$ and $\omega'(1,h)\Phi'$ correspond.

(7) Let $v$ be clean for $\mathcal X$, and let $p$ be the rational prime below
$v$. As $p$ splits completely in $E$, every prime of $E$ above $p$ has
ramification index and residue degree $1$ over $p$: its absolute norm is $p$,
and $p$ does not lie in its square. As $B$ is a theta base, $E/F$ is
unramified at $v$, so $v$ is split or inert in $E$. If $v$ were inert, the
prime $v\mathcal O_E$ of $E$ would lie above $p$ and have absolute norm
$N(v)^2$, which is not the prime number $p$. So $v$ splits, and either of the
two places of $E$ above $v$ gives an $F$-embedding of $E$ into $F_v$.

Let $\tau\colon E\to F_v$ be an $F$-embedding, and $\mathfrak P$ the prime of
$E$ it defines. As $v\notin T$, $\mathfrak P$ is not above $T$, so $\alpha$
and $\beta$, hence $\nu$, are unramified at $\mathfrak P$
(Definition~\ref{asm-theta-agree-triple}(3)). As $\tau$ identifies
$E_{\mathfrak P}$ with $F_v$ and $\nu$ is unramified at $\mathfrak P$,
$\nu_\tau(\varpi_v)=\nu_{\mathfrak P}(\varpi)$
for every $\varpi\in E^\times$ with $\operatorname{ord}_{\mathfrak P}(\varpi)=1$.
For $z\in\mathbf A_E^\times$ whose finite components are $1$,
$\nu(z)=\prod_w(\pi_w(z)/|\pi_w(z)|)^{n(w)}$ with
$n(w)=m_\alpha(w)-m_\beta(w)>a+a'$, as the character datum and $G$ are proper
(Definition~\ref{asm-theta-agree-triple}(2) and~(3)); so $n(w)\neq0$, and for
$z=x_\infty$ with $x\in E^\times$, $\pi_w(z)=\tau_w(x)$. By
Lemma~\ref{asm-theta-agree-definite-hecke}, applied to $\nu$ and $\mathfrak P$,
$\nu_\tau(\varpi_v)$ is not a root of unity; in particular it is not $1$.

(8) The space $\mathcal S^\circ(M')$ is a linear subspace. If $\Phi_1$ and
$\Phi'_1$ correspond and $\Phi_2$ and $\Phi'_2$ correspond, then so do
$\Phi_1+z\Phi_2$ and $\Phi'_1+z\Phi'_2$: after enlarging the sets $T_1$ as in
the proof of (6), both pairs are written with the same set, and the identity of
Definition~\ref{asm-theta-agree-matched}(4) for the first pair plus $z$ times
that for the second is the identity for the combination. This gives the first
claim. As $p=q=0$, $W$ is skew-Hermitian of signature $(a+a',0)$, so for every
infinite place $w$ the matrix $N_w=i\,\tau_w(J_W)$ has $x^{\mathsf T}N_w\bar x>0$
for $x\neq0$. Its Cartan involution $\theta=\theta_{W,w}$ is $1$: if
$\theta u\neq u$, then $x=\theta u-u$ is nonzero and $\theta x=-x$, so
$\operatorname{Re}\bigl((\theta x)^{\mathsf T}N_w\bar x\bigr)=-x^{\mathsf T}N_w\bar x<0$,
against Definition~\ref{asm-theta-agree-doubling}(1). So every $h\in U_W$ lies
in $K^\theta_W$ (Definition~\ref{asm-theta-agree-end-notions}(2)), and
Lemma~\ref{asm-theta-agree-end-stable} gives the second claim.

(9) Let $C_{T_1}$ and $\Pi_{T_1}$ be as in
Lemma~\ref{asm-theta-agree-definite-inputs-pairs-fill}, and let $C'_{T_1}$ be
defined in the same way from $M'$, $p'$, the $p'_v$ and the
$\mathrm{vac}'_v$. By Definition~\ref{asm-theta-agree-matched}(4), $\Phi$ and
$\Phi'$ correspond if and only if there are a finite set $T_1\supset T$ and
$\xi\in\Pi_{T_1}$ with $(\Phi,\xi)\in C_{T_1}$ and $(\Phi',\xi)\in C'_{T_1}$; and
then, as in the proof of (6), for every finite $T_2\supset T_1$ the element
$\xi\otimes\bigotimes_{v\in T_2\setminus T_1}p_v(\mathrm{vac}_v)$ serves for
$T_2$, as $p_v(\mathrm{vac}_v)=p'_v(\mathrm{vac}'_v)$ for $v\notin T$.

Suppose that some $(\Phi_0,\Phi'_0)\in N$ is such that $\Phi_0$ and $0$ do not
correspond, and let $y=(\Phi,\Phi')\in\mathcal P$. Choose $T_1$ that serves for
both pairs: $(\Phi_0,\xi_0)$ and $(\Phi,\xi)$ lie in $C_{T_1}$, and
$(\Phi'_0,\xi_0)$ and $(\Phi',\xi)$ in $C'_{T_1}$. Then $\xi_0\neq0$, since
otherwise $(0,\xi_0)\in C'_{T_1}$ would make $\Phi_0$ and $0$ correspond. The
set $N_1$ of first components of elements of $N$ is a space with
$\omega(1,h)N_1\subset N_1$ for $h\in U_W$, as $N$ is stable under the pairs
$(\omega(1,h),\omega'(1,h))$. By
Lemma~\ref{asm-theta-agree-definite-inputs-pairs-fill} there is
$(\Psi,\Psi')\in N$ with $(\Psi,\xi)\in C_{T_1}$. By (8),
$y-(\Psi,\Psi')\in\mathcal P$; and $\Phi-\Psi$ and $0$ correspond, as
$(\Phi-\Psi,0)\in C_{T_1}$ and $(0,0)\in C'_{T_1}$. So $y-(\Psi,\Psi')\in N$ by
the hypothesis on $N$, and $y\in N$. Hence $N=\mathcal P$.

(10) As $p=q=0$, $V'$ is Hermitian of signature $(0,a')$ and $W$ is
skew-Hermitian of signature $(a+a',0)$, so the Weil pair
$(V',W,M',\alpha,\alpha/\beta,\omega')$ has signature $(0,a';a+a',0)$, with
ranks $n=a'$ and $d=a+a'$, both at least $1$ as $a'>2(p+q)+1$
(Definition~\ref{asm-theta-agree-triple}(2)). Let $(I\sqcup I,e^\square,\omega_2,S,l)$ be the
doubling datum for $(V',W)$ of
Lemma~\ref{asm-theta-agree-doubling-datum}(1), with character $\alpha$, let $K$
be the lattice stabilizer of the doubled space of $W$, $a$ its height function
and $\iota$ the doubling embedding of $W$
(Lemma~\ref{asm-theta-agree-doubling-stabilizer}), and put
$\Phi'=\Psi\boxtimes\bar\Psi$, $f(x)=l(\omega_2(x)\Phi')$ and $f_s=a^{s-z_0}f$.
Then $z_0=(n-d)/2$.

\emph{The data.} By Lemma~\ref{asm-theta-agree-doubling-section}(1), $f$ is a
section of degree $z_0$ for $(K,\alpha)$, so
$f(px)=\alpha(\det t(p))|\det t(p)|^{a'/2}f(x)$ for $p$ in the Siegel parabolic
subgroup of $G_W$, as $z_0+(a+a')/2=a'/2$. Let $U$ and $\hat E$ be given by
Lemma~\ref{asm-theta-agree-doubling-value} for
Lemma~\ref{asm-theta-agree-doubling-siegel-weil}(1), with the measure $\mu'$
and the domain $\mathcal D'$ of $M'$: $U$ is a connected open set containing
$z_0$, $\hat E$ has properties (1) to (4) of
Lemma~\ref{asm-theta-agree-doubling-continuation} for $(f_s)$, and, as $d\geq n$,
\[
\hat E(z_0,x)=\frac{2}{\mu'(\mathcal D')}\int_{\mathcal D'}\theta\bigl(\omega_2(x)S(g)\Phi'\bigr)\,d\mu'(g)
\]
for $x\in G_W$. Let $P(s,\phi)$ and $t_1$ be given by
Lemma~\ref{asm-theta-agree-doubling-kernel} for $Y=W$, $\chi=\alpha$,
$z=a'/2$, $f$, $s_0=z_0$, $U$, $\hat E$, $\mu$ and $\mathcal D$. Put
$t_0=t_1$, $Z_\phi(z)=P(z,\phi)$ for $z\in U$, and
$F_z(g)=f_z(\iota(1,g))\,\alpha(\det g)^{-1}$. By that lemma, $Z_\phi$ is
holomorphic on $U$, and for $\operatorname{Re}z>t_0$, $z\in U$, $F_z$ is
$\mu$-integrable and
$Z_\phi(z)=\int_{U_W}F_z(g)\langle R_g\phi,\phi\rangle_{\mathcal D}\,d\mu(g)$.

\emph{The value at $z_0$.} Let $\theta=\theta_{\omega'}(\Psi)$. The function
$\Psi$ is Schwartz-Bruhat (Lemma~\ref{asm-theta-agree-nonvanishing-schwartz}),
so by the seesaw identity of Lemma~\ref{asm-theta-agree-doubling-datum}(1),
for $h,h'\in U_W$ and $g\in U_{V'}$,
$\alpha(\det h')^{-1}\theta\bigl(\omega_2(\iota(h,h'))S(g)\Phi'\bigr)=\theta(g,h)\overline{\theta(g,h')}$.
Hence
\[
\alpha(\det h')^{-1}\hat E(z_0,\iota(h,h'))=\frac{2}{\mu'(\mathcal D')}\int_{\mathcal D'}\theta(g,h)\overline{\theta(g,h')}\,d\mu'(g).
\]
The function $\theta$ is continuous
(Lemma~\ref{asm-theta-agree-convergence-theta}(2)) and invariant under
$U_{V'}(F)\times U_W(F)$
(Lemma~\ref{asm-theta-agree-anisotropic-automorphy}), so it is bounded by
Lemma~\ref{asm-theta-agree-definite-compact}(4), for the definite spaces $V'$
and $W$; a rapidly decreasing $\phi$ is bounded by the same lemma; and
$\mu(\mathcal D)$ and $\mu'(\mathcal D')$ are finite and nonzero by part (3)
of that lemma. So Fubini's theorem applies to
$(h,h',g)\mapsto\theta(g,h)\overline{\theta(g,h')}\phi(h')\overline{\phi(h)}$
on $\mathcal D\times\mathcal D\times\mathcal D'$, and
\[
Z_\phi(z_0)=\frac{2}{\mu'(\mathcal D')}\int_{\mathcal D'}\Bigl(\int_{\mathcal D}\phi(h')\overline{\theta(g,h')}\,d\mu(h')\Bigr)\overline{\Bigl(\int_{\mathcal D}\phi(h)\overline{\theta(g,h)}\,d\mu(h)\Bigr)}\,d\mu'(g).
\]
The inner integrals are $\theta^*_{\omega'}(\phi,\Psi)(g)$
(Definition~\ref{asm-theta-agree-definite-products}(2)), which gives the
formula for $Z_\phi(z_0)$.

% Lean, done above: (1) polRank (HermRank.lean): G eps is the infimum over p of the kernels
% of dL (J p) - I dL p in exists_vac, with J from Dense.exists_real_structure for the signs
% eps and the injectivity of f -> f 0 from ode_zero; N is exists_stable for the dL of a
% basis, and actsD_unique identifies the first-order actions. (3) is TC.inv_loc
% (ThConstInv.lean) with TC.SplitModel (TC.splitModel = splitModel_of
% KCo.dictionary_extended, Lemma asm-theta-agree-coinvariants-dictionary), as TC.const_of
% uses them: the pure tensors are eqOn_span' and schA_le, Psi[psi] is
% intersection_update_fsHerm_neg_one, Q is beta, and the bilinear step is locBil with
% LocDist (Lemmas asm-theta-agree-definite-dist-bilinear and -dist). (6) is TC.pm_fin
% (ThConstMatch.lean): On.mono for v in T_1, loc_op for O and O', and A from
% exists_rightInverse_of_surjective and hm.fin.equiv. (7) is goodNe: the embedding is
% nonempty_emb_of_splitsFully (KFinLoc.split_or_inert, SplitFullyNorm, absNorm_algebraMap);
% the value is nu_pow_ne_one with k = 1 (chrW_unifAt, HeckeBridge.exists_hchar, hasExp_nu,
% Geo kappa and CMSplitValue.not_isOfFinOrder_val, Lemma asm-theta-agree-definite-hecke).
% (8) is pairsStable = pairsStable_of (ThKEnd.Grp.matchStab_of ...) (PairsB.lean): pairs_lin
% (polMatch_add_smul), and MatchStab (Lemma asm-theta-agree-end-stable, whose proof is still
% open) for k = h, which lies in FinK by comm_theta and cartan_one.
% Lean, (4), done above: ConstOfGood, given by PositivityTheory.SA.constOfGood
% (EqTopOfLines.lean) from Published.HeckeLEntireAndNonzeroAtOne_holds: L is the subgroup
% with carrier the intersection of {x | f (x g) = f g}, the conjugates are conj_mem_UAt, and
% mem_of_good is eq_top_of_lines with HasLines from exists_line and normOne_mem (Lemmas
% asm-theta-agree-definite-inputs-constancy-lines and -norm); hermND_transpose for J. The
% hypothesis 2 <= n is not used.
% Lean, (5), done above: Factor = ThKF.factor_of KDecay Kills (Factor.lean), a property used
% by ThKR.polynomial_vanishing_of_compactness; KDecay is Lemma asm-theta-agree-decay and Kills
% is Lemma asm-theta-agree-local-injective-kills. The extensions J and J' are exists_Ilin,
% the descent is vanish (desc and exists_pureML, Lemma
% asm-theta-agree-definite-inputs-tensor(2)), the operator O and pi are so_proj
% (exists_factorOp, exists_section), pure_mem_So puts the pure functions in So(M'), and the
% vanishing tensors come from tens_eq_map (Lemma asm-theta-agree-definite-inputs-tensor(1)).
% Lean, (2), done above: the third clause of the first part of ThetaBasic (thetaBasic_of
% weil_rep_rapid_decay_height), ThKF0.theta_moderate_growth_of_decay (ThetaGrow.lean): hlat is
% theta_sum_norm_bound_denominator (Lemma asm-theta-agree-definite-inputs-growth-lattice), harch
% is ArchBd and hfin is weil_rep_bounded_denominator, whose proof is the second paragraph, with
% htBd, finInv, exists_dual_den (D = |disc F| in differentIdeal) and trace_int
% (PsiNontriv.isStd_fin_rat). htBd and ArchBd are parts (1) and (3) of Lemma
% asm-theta-agree-definite-inputs-growth-inputs, proved above.
% Lean, (9), done above: pairsIrr = pairsIrr_of_arch A.archOp = pairsIrr_of locRepEx pureAct
% archOp polFinTri (PairsFin.lean, PairsB.lean, PairsL.lean): C_T1 is Cl B G.M R.p Rf.p T1,
% the description of correspondence is polMatch_iff with polMatch_big and enlarge, the pair
% (Psi, Psi') is fill (Lemma asm-theta-agree-definite-inputs-pairs-fill), and y - (Psi, Psi')
% is pairs_lin with hnull.
% Lean, (10), done above: coefLine_of (CoefLine.lean) up to its term val. The data are
% DblW (dblW) for pairV' T, HgtEx (hgtEx_of stab2Cpt) for the lattice stabilizer K of the
% doubled space of W, HgtInt and Stab2Cpt, the section of BoxSec (boxSec_of finLevel
% stab2Cpt) at Psi and Psi, SWFull (swFull_of (siegelWeilAniso_of SiegelWeilStd frameConj
% finLevel stab2Cpt).1 (eisContinues_of EisBL frameConj) boxSec), and kpair_cont, which
% gives hol, t_1 and far: Z_phi(z) = kpair alpha e mu D (E^ z) phi and
% F_z(g) = flat a f z0 z (e(1,g)) alpha(det g)^{-1}. The value is val: hv with
% if_neg (not a + a' < a') for the factor 2, the seesaw identity d2 of DblW, and
% An.at_value (FarRight.lean, Fubini through integral_integral_swap and integral_prod_mul)
% with the bounds (b2 _ s3) and cV.bdd2 cW for theta and cW.bdd for phi.
\end{proof}

\begin{lemma}\label{asm-theta-agree-definite-orbits}
Let $G$ be a topological group that is a Baire space, let $\rho$ be a
representation of $G$ on a complex vector space $P$ of countable dimension, and
let $u$ be an injective linear map from $P$ to the continuous functions on $G$
with $u(\rho(g)x)(h)=u(x)(hg)$ for all $g,h\in G$ and $x\in P$. Then for every
$x\in P$ and every compact set $C\subset G$ there is a finite-dimensional
subspace $V\subset P$ with $\rho(g)x\in V$ for all $g\in C$.
\end{lemma}
\leanref{Development/WeilBound/ThetaLifts/LatticeBasics/FinOrbit.lean}{GlobalTheory.ThK.Dbl.Herm.fin_orbit}
\begin{proof}
Let $(b_i)_{i\in\mathcal B}$ be a basis of $P$; the set $\mathcal B$ is
countable. For a finite set $\beta\subset\mathcal B$ let $V_\beta$ be the span
of the $b_i$ with $i\in\beta$, and $A_\beta=\{g\in G:\rho(g)x\in V_\beta\}$.
Give the space of functions on $G$ the topology of pointwise convergence. The
map $g\mapsto u(\rho(g)x)$, that is $g\mapsto(h\mapsto u(x)(hg))$, is
continuous, as $u(x)$ is. The finite-dimensional subspace $u(V_\beta)$ is
closed, and $A_\beta$ is its inverse image, as $u$ is injective; so $A_\beta$
is closed. Every $\rho(g)x$ lies in some $V_\beta$, so the countably many
closed sets $A_\beta$ cover $G$. As $G$ is a Baire space, some $A_\beta$ has an
interior point $g_0$.

Let $k\in G$. The set $N_k=\{g:g_0k^{-1}g\in A_\beta\}$ is a neighbourhood of
$k$, and for $g\in N_k$,
$\rho(g)x=\rho(kg_0^{-1})\rho(g_0k^{-1}g)x\in\rho(kg_0^{-1})V_\beta$, a
finite-dimensional subspace. Finitely many of the $N_k$ cover $C$, and the sum
$V$ of the corresponding subspaces $\rho(kg_0^{-1})V_\beta$ has the required
property.
\end{proof}

\begin{lemma}\label{asm-theta-agree-definite-eigen}
Let $U$ be a complex inner product space, with inner product $\langle x,y\rangle$
conjugate-linear in $x$. Let $K$ be a group acting on $U$ by linear maps
$\kappa(k)$ with $\langle\kappa(k)x,\kappa(k)y\rangle=\langle x,y\rangle$, such
that the span $O_x$ of the orbit $\{\kappa(k)x:k\in K\}$ of every $x\in U$ is
finite-dimensional. Let $l_1,\ldots,l_m$ be linear forms on $U$ such that $x=0$
whenever $l_i(\kappa(k)x)=0$ for all $i$ and $k$. Let $b$ be a form on $U$,
conjugate-linear in the first variable and linear in the second, with
$b(\kappa(k)x,\kappa(k)y)=b(x,y)$ for all $k$, $x$ and $y$, and let $f\in U$ be
nonzero. Then there are a nonzero $x\in U$ and $c\in\mathbf C$ with
$b(x,y)=c\langle x,y\rangle$ for all $y\in U$.
\end{lemma}
\leanref{Development/WeilBound/ThetaLifts/LatticeBasics/HermAbs.lean}{GlobalTheory.ThK.Dbl.Herm.exists_eigen}
\begin{proof}
Each $O_x$ is stable under $K$. Let $H$ be the space of linear maps
$L\colon O_f\to U$ with $L(\kappa(k)w)=\kappa(k)L(w)$ for all $k\in K$ and
$w\in O_f$. It is finite-dimensional: the linear map sending $L$ to the tuple
$(l_i\circ L)_i$ of linear forms on the finite-dimensional space $O_f$ is
injective, since if $l_i(L(w))=0$ for all $i$ and $w$, then
$l_i(\kappa(k)L(w))=l_i(L(\kappa(k)w))=0$ for all $i$ and $k$, so $L(w)=0$. Let
$Z$ be the span of the $L(w)$ with $L\in H$ and $w\in O_f$. It is
finite-dimensional, it is stable under $K$, since $\kappa(k)L(w)=L(\kappa(k)w)$,
and it contains $f$, the image of $f$ under the inclusion of $O_f$. So
$Z\neq0$.

\emph{Step 1: $b(z,y)=0$ for $z\in Z$ and $y\in Z^\perp$.} As $Z$ is stable
under $K$ and the $\kappa(k)$ preserve the inner product, $Z^\perp$ is stable
under $K$, so $O_y\subset Z^\perp$. As $O_y$ is finite-dimensional, there is a
linear map $S\colon U\to O_y$ with $\langle S(x),w\rangle=b(x,w)$ for $x\in U$
and $w\in O_y$. It commutes with the $\kappa(k)$: for $w\in O_y$,
\[
\langle S(\kappa(k)x),w\rangle=b(\kappa(k)x,w)=b(x,\kappa(k)^{-1}w)
=\langle S(x),\kappa(k)^{-1}w\rangle=\langle\kappa(k)S(x),w\rangle .
\]
So for $L\in H$ the map $S\circ L$, with values in $O_y\subset U$, lies in $H$,
and $S(L(w))\in Z$ for $w\in O_f$. Hence
$b(L(w),y)=\langle S(L(w)),y\rangle=0$, as $y\in Z^\perp$. As the $L(w)$ span
$Z$, Step 1 follows.

\emph{Step 2.} As $Z$ is finite-dimensional, there is a linear map
$S\colon U\to Z$ with $\langle S(x),z\rangle=b(x,z)$ for $x\in U$ and $z\in Z$.
Its restriction to $Z$ is an endomorphism of the nonzero finite-dimensional
complex space $Z$, so it has an eigenvector $x\neq0$, with $S(x)=e\,x$. Let
$y\in U$, and write $y=y_1+y_2$ with $y_1\in Z$ and $y_2\in Z^\perp$. By Step 1,
$b(x,y_2)=0$, and $\langle x,y_2\rangle=0$. So
$b(x,y)=b(x,y_1)=\langle e\,x,y_1\rangle=\bar e\langle x,y_1\rangle=\bar e\langle x,y\rangle$,
and $c=\bar e$ has the required property.
\end{proof}

\begin{lemma}\label{asm-theta-agree-definite-schur}
Let $G$ be a locally compact Hausdorff group, $K\subset G$ a compact open
subgroup and $\Omega\subset G$ a compact set. Let $\rho$ be a representation of
$G$ on a complex vector space $P$, and $u$ a linear map from $P$ to the
continuous functions on $G$ such that $u(\rho(g)x)(h)=u(x)(hg)$ for all
$g,h\in G$ and $x\in P$, the image of $u$ has countable dimension, and for
every $g\in G$ there is $g'\in\Omega$ with $u(x)(g)=u(x)(g')$ for all $x\in P$.
Let $b$ and $c$ be forms on $P$, conjugate-linear in the first variable and
linear in the second, with $b(\rho(g)x,\rho(g)y)=b(x,y)$ and
$c(\rho(g)x,\rho(g)y)=c(x,y)$ for all $g$, $x$ and $y$. Suppose that
$c(y,x)=\overline{c(x,y)}$ and $\operatorname{Re}c(x,x)\geq0$ for all $x$ and
$y$; that $u(x)=0$ whenever $c(x,x)=0$; that $b(x,y)=b(y,x)=c(y,x)=0$ for all
$y$ whenever $u(x)=0$; and that every subspace $N\subset P$ containing
$\ker u$ with $\rho(g)N\subset N$ for all $g$ satisfies $N\subset\ker u$ or
$N=P$. Then there is $\lambda\in\mathbf C$ with $b(x,y)=\lambda\,c(x,y)$ for
all $x,y\in P$.
\end{lemma}
\leanref{Development/WeilBound/ThetaLifts/LatticeBasics/FinOrbit.lean}{GlobalTheory.ThK.Dbl.Herm.line}
\begin{proof}
If $u=0$, then $b=0$, and $\lambda=0$ will do. Otherwise let $f\in P$ with
$u(f)\neq0$. The kernel of $u$ is stable under $\rho$, as
$u(\rho(g)x)=R_gu(x)$. Let $\bar P=P/\ker u$, with the induced representation
$\bar\rho$, the induced injective map $\bar u$ and the induced forms $\bar b$
and $\bar c$; these are defined, since $b(x,y)=b(y,x)=0$ and
$c(y,x)=\overline{c(x,y)}=0$ for all $y$ when $u(x)=0$. By the conditions on
$c$, $\bar c$ is an inner product on $\bar P$, conjugate-linear in the first
variable. The space $\bar P$ is isomorphic to the image of $u$, so it has
countable dimension.

The open sets $gK$ with $g\in\Omega$ cover $\Omega$; let $g_1,\ldots,g_m$ be
such that $\Omega\subset g_1K\cup\cdots\cup g_mK$, and put
$l_i(z)=\bar u(z)(g_i)$. Apply Lemma~\ref{asm-theta-agree-definite-eigen} to
$\bar P$ with the inner product $\bar c$, the group $K$ acting by $\bar\rho$,
the forms $l_i$ and $\bar b$, and the image of $f$. Its hypotheses hold. The
$\bar\rho(k)$ preserve $\bar c$ and $\bar b$. As $G$ is locally compact and
Hausdorff, it is a Baire space, so by
Lemma~\ref{asm-theta-agree-definite-orbits} for $\bar\rho$, $\bar u$ and the
compact set $K$, the orbit of every $z\in\bar P$ under $K$ spans a
finite-dimensional space. If $l_i(\bar\rho(k)z)=0$ for all $i$ and $k\in K$,
let $g\in G$; take $g'\in\Omega$ with $\bar u(z)(g)=\bar u(z)(g')$, and $i$
and $k\in K$ with $g'=g_ik$. Then
$\bar u(z)(g)=\bar u(z)(g_ik)=\bar u(\bar\rho(k)z)(g_i)=0$; so $\bar u(z)=0$
and $z=0$. The lemma gives $x_0\in P$ with $u(x_0)\neq0$ and
$\lambda\in\mathbf C$ with $b(x_0,y)=\lambda\,c(x_0,y)$ for all $y\in P$.

Let $N$ be the set of $x\in P$ with $b(x,y)=\lambda\,c(x,y)$ for all $y$. It is
a subspace containing $\ker u$, and it is stable under $\rho$: for $x\in N$,
$b(\rho(g)x,y)=b(x,\rho(g^{-1})y)=\lambda\,c(x,\rho(g^{-1})y)=\lambda\,c(\rho(g)x,y)$.
As $x_0\in N$ and $x_0\notin\ker u$, $N\not\subset\ker u$, so $N=P$.
\end{proof}

\begin{lemma}\label{asm-theta-agree-definite-open}
Let $B$ be a theta base for a CM field $E$
(Definition~\ref{asm-theta-agree-models}(1)), and let $Y$ be a lattice space
of rank $n$ that is Hermitian or skew-Hermitian of signature $(r,s)$ with
$\min(r,s)=0$ (Definition~\ref{asm-theta-agree-doubling}(1)). Then $U_Y$ has a
compact open subgroup.
\end{lemma}
\leanref{Development/WeilBound/ThetaLifts/LatticeBasics/RankPolData.lean}{GlobalTheory.ThK.Dbl.CptOpen}
\leanref{Development/WeilBound/ThetaLifts/LatticeBasics/HermCpt.lean}{GlobalTheory.ThK.Dbl.cptOpen}
\begin{proof}
Let $\widehat{\mathcal O}_F\subset\mathbf A_f$ be the product of the
$\mathcal O_v$, let $(b_i)$ be a $\mathbf Z$-basis of $\mathcal O_E$, and let
$\Lambda\subset\mathbf A_E$ be the set of sums $\sum_ia_ib_i$ with
$a_i\in\mathbf A$ of archimedean component $0$ and finite component in
$\widehat{\mathcal O}_F$. It is a compact additive subgroup of $\mathbf A_E$,
closed under multiplication, and the finite component $a_f$ of every
$a\in\mathbf A_E$ near $0$ lies in $\Lambda$. Let $K$ be the set of $g\in U_Y$
such that the finite components of $g-1$ and $g^{-1}-1$ have entries in
$\Lambda$. As $gh-1=(g-1)(h-1)+(g-1)+(h-1)$ and $(gh)^{-1}=h^{-1}g^{-1}$, $K$
is a subgroup. It contains a neighbourhood of $1$, as $g\mapsto g$ and
$g\mapsto g^{-1}$ are continuous; so $K$ is open, hence also closed.

Let $w$ be an infinite place of $F$ and $M_w$ the Hermitian matrix of
signature $(r,s)$ of Definition~\ref{asm-theta-agree-doubling}(1). Every
$g\in U_Y$ has $\pi_w(g)^{\mathsf T}M_w\overline{\pi_w(g)}=M_w$, as $c$ induces
complex conjugation under $\tau_w$. As $\min(r,s)=0$, $M_w$ or $-M_w$ is
positive definite, so these matrices $\pi_w(g)$ lie in a compact set $C_w$, a
conjugate of the unitary group. There is a continuous map $S$ from
$\prod_w\mathbf C$ to $\mathbf A_E$ with $a=S((\pi_w(a))_w)+a_f$ for every
$a\in\mathbf A_E$. So for $g\in K$ the matrices $g$ and $g^{-1}$ lie in the
image of the compact set $\prod_wC_w\times M_n(\Lambda)$ under the continuous
map $((m_w)_w,l)\mapsto S((m_w)_w)+l+1_f$, applied entrywise, where $1_f$ is
the identity matrix at the finite places. The topology of $U_Y$ is that of a
closed subgroup of $GL_n(\mathbf A_E)$, so a closed subset of $U_Y$ on which
$g$ and $g^{-1}$ range over compact sets of matrices is compact. Hence $K$ is
compact.
\end{proof}

\begin{lemma}\label{asm-theta-agree-definite-line}
Assume the meromorphic continuation of Siegel Eisenstein series in the form of
Theorem~\ref{asm-theta-agree-eisenstein} and the Siegel-Weil formula in the
form of Theorem~\ref{asm-theta-agree-siegel-weil}. Let
$(B,G,\mathcal X,\omega,\omega',R,R_f)$ be a matched quadruple for a theta base
$B$ of a CM field $E$ (Definition~\ref{asm-theta-agree-matched}) with $p=q=0$,
let $\mu$ be a Haar measure on $U_W$ and $\mathcal D$ a fundamental domain for
$U_W(F)$ with respect to $\mu$, and let $\mathcal P$ and $\mathcal I_{s,t}$ be
as in Definition~\ref{asm-theta-agree-definite-products}. Then for all
$s,t,s',t'\in\mathbf C$ there is $\lambda\in\mathbf C$ with
\[
\langle\mathcal I_{s,t}(y),\mathcal I_{s',t'}(y')\rangle_{\mathcal D}
=\lambda\,\langle\mathcal I_{0,1}(y),\mathcal I_{0,1}(y')\rangle_{\mathcal D}
\qquad(y,y'\in\mathcal P).
\]
\end{lemma}
\leanref{Development/WeilBound/ThetaLifts/LatticeBasics/PairsStable.lean}{GlobalTheory.ThK.Dbl.HermLine}
\leanref{Development/WeilBound/ThetaLifts/LatticeBasics/HermCpt.lean}{GlobalTheory.ThK.Dbl.hermLine_of'}
\leanref{Development/WeilBound/ThetaLifts/LatticeBasics/RankPolData.lean}{GlobalTheory.ThK.Dbl.hermLine_of}
\leanref{Development/WeilBound/ThetaLifts/LatticeBasics/RankPolData.lean}{GlobalTheory.ThK.Dbl.Herm.line_cpt}
\leanref{Development/WeilBound/ThetaLifts/LatticeBasics/RankPolData.lean}{GlobalTheory.ThK.Dbl.Herm.exists_form}
\begin{proof}
The space $W$ is skew-Hermitian of signature $(a+a',0)$, and $V$ and $V'$ are
Hermitian of signatures $(a,0)$ and $(0,a')$.

\emph{The space and the maps.} The pair $(0,0)$ lies in $\mathcal P$, as
$0\in\mathcal S^\circ(M')$ and $0$ corresponds to $0$ with $T_1=T$ and empty
sums. By Lemma~\ref{asm-theta-agree-definite-inputs}(8), $\mathcal P$ is a
complex vector space, and
$\rho(h)(\Phi,\Phi')=(\omega(1,h)\Phi,\omega'(1,h)\Phi')$ defines a
representation $\rho$ of $U_W$ on it. Let $(\Phi,\Phi')\in\mathcal P$. By
Lemma~\ref{asm-theta-agree-partners}(3), $\Phi\in\operatorname{PD}(M)$ and
$\Phi'\in\operatorname{PD}(M')$, and both are Schwartz-Bruhat, since the models
are regular. As $\min(a,0)=\min(0,a')=0$,
Lemma~\ref{asm-theta-agree-convergence}(1), applied to the Weil pairs
$(V,W,M,\beta,\beta/\alpha,\omega)$ and $(V',W,M',\alpha,\alpha/\beta,\omega')$,
shows that every Schwartz-Bruhat function converges for $(M,\omega)$ and for
$(M',\omega')$. As theta functions are linear
(Lemma~\ref{asm-theta-agree-convergence-theta}(2)), $I$ and $I'$ are linear on
the Schwartz-Bruhat functions, and $\mathcal I_{s,t}$ is linear on
$\mathcal P$. By Lemma~\ref{asm-theta-agree-definite-continuous}, every
$\mathcal I_{s,t}(y)$ is continuous; by
Lemma~\ref{asm-theta-agree-anisotropic-automorphy}, it is invariant under left
translation by $U_W(F)$; and as $\omega$ and $\omega'$ are representations,
$\theta_\omega(\omega(1,g)\Phi)(g',h)=\theta_\omega(\Phi)(g',hg)$ and likewise
for $\omega'$, so $\mathcal I_{s,t}(\rho(g)y)=R_g\mathcal I_{s,t}(y)$.

\emph{The kernel.} Put $u=\mathcal I_{0,1}$, so that $u(\Phi,\Phi')=I'(\Phi')$.
Let $y=(\Phi,\Phi')\in\mathcal P$ be such that $\Phi$ and $0$ correspond. Then
$(\Phi,0)\in\mathcal P$, so $(0,\Phi')=y-(\Phi,0)\in\mathcal P$, that is, $0$
and $\Phi'$ correspond. By Lemma~\ref{asm-theta-agree-definite-inputs}(5),
$I(\Phi)=0$ and $I'(\Phi')=0$, so $\mathcal I_{s,t}(y)=0$ for all $s$ and $t$.

Let $N\subset\mathcal P$ be a subspace containing $\ker u$ and stable under
$\rho$. By the last paragraph, $N$ contains every $(\Phi,\Phi')\in\mathcal P$
with $\Phi$ and $0$ corresponding. By
Lemma~\ref{asm-theta-agree-definite-inputs}(9), either $\Phi$ and $0$
correspond for every $(\Phi,\Phi')\in N$, so that $N\subset\ker u$ by the last
paragraph, or $N=\mathcal P$. Apply this to $N=\ker u$, which is stable under
$\rho$ as $u(\rho(g)y)=R_gu(y)$. If $\ker u=\mathcal P$, then $I'(\Phi')=0$
for all $(\Phi,\Phi')\in\mathcal P$; but by
Lemma~\ref{asm-theta-agree-nonvanishing} there is
$\Phi'\in\operatorname{PD}(M')\cap\mathcal S^\circ(M')$ with $I'(\Phi')\neq0$,
as $q=0$, and by Lemma~\ref{asm-theta-agree-partners}(2) some $\Phi$
corresponds to it, so that $(\Phi,\Phi')\in\mathcal P$. Hence $\Phi$ and $0$
correspond for every $(\Phi,\Phi')\in\ker u$, and every $\mathcal I_{s,t}$
vanishes on $\ker u$.

\emph{Countable dimension.} By Lemma~\ref{asm-theta-agree-partners}(3), the
image of $u$ lies in $I'(\operatorname{PD}(M'))$. A spanning function
$f\otimes\bigotimes_v\phi_v$ of $\operatorname{PD}(M')$ is $f\otimes\phi$ with
$\phi(x)=\prod_v\phi_v(x_v)$, which is Schwartz-Bruhat on $\mathbf A_f^{I'_0}$,
since $\phi_v=\mathrm{vac}'_v$ is the indicator function of
$\mathcal O_v^{I'_0}$ for almost all $v$. The space $\operatorname{Pol}(M')$
has countable dimension by Lemma~\ref{asm-theta-agree-definite-inputs}(1), and
so does the space of Schwartz-Bruhat functions on $\mathbf A_f^{I'_0}$, which
is spanned by the indicator functions of the countably many compact open
subsets. Hence $\operatorname{PD}(M')$ and the image of $u$ have countable
dimension.

\emph{The forms.} Let $w$ and $w'$ be among the maps $\mathcal I_{s,t}$, and
put $\beta_{w,w'}(x,y)=\langle w(y),w'(x)\rangle_{\mathcal D}$ for
$x,y\in\mathcal P$. The integral converges, since the integrand is continuous
and bounded by Lemma~\ref{asm-theta-agree-definite-compact}(4), and
$\mu(\mathcal D)<\infty$. The form is linear in $y$ and conjugate-linear in
$x$, and $\beta_{w,w'}(\rho(g)x,\rho(g)y)=\beta_{w,w'}(x,y)$: the function
$h\mapsto w(y)(h)\overline{w'(x)(h)}$ is invariant under left translation by
$U_W(F)$, right translation by $g$ preserves $\mu$
(Lemma~\ref{asm-theta-agree-definite-compact}(3)) and carries $\mathcal D$ to
the fundamental domain $\mathcal Dg$, and the integrals of such a function over
two fundamental domains agree, as in the proof of
Lemma~\ref{asm-theta-agree-anisotropic-invariance}.

\emph{Conclusion.} Apply Lemma~\ref{asm-theta-agree-definite-schur} to the
group $U_W$, which is locally compact and Hausdorff by
Lemma~\ref{asm-theta-agree-definite-compact}(1), a compact open subgroup $K$
(Lemma~\ref{asm-theta-agree-definite-open}), the compact set $\Omega$ of
Lemma~\ref{asm-theta-agree-definite-compact}(2), the representation $\rho$,
the map $u$, and the forms $b=\beta_{\mathcal I_{s,t},\mathcal I_{s',t'}}$ and
$c=\beta_{u,u}$. Its hypotheses hold. For every $g$ there is $\gamma$ with
$\gamma g\in\Omega$, and $u(x)(g)=u(x)(\gamma g)$. The form $c$ has
$c(y,x)=\overline{c(x,y)}$ and $c(x,x)=\int_{\mathcal D}|u(x)|^2\,d\mu\geq0$,
and $c(x,x)=0$ implies $u(x)=0$ by
Lemma~\ref{asm-theta-agree-definite-compact}(4). If $u(x)=0$, then
$\mathcal I_{s,t}(x)=\mathcal I_{s',t'}(x)=0$ by the kernel paragraph, so
$b(x,y)=b(y,x)=c(y,x)=0$ for all $y$. The condition on subspaces was shown
above. So there is $\lambda$ with $b=\lambda c$, that is,
$\langle\mathcal I_{s,t}(y),\mathcal I_{s',t'}(x)\rangle_{\mathcal D}=\lambda\,\langle\mathcal I_{0,1}(y),\mathcal I_{0,1}(x)\rangle_{\mathcal D}$
for all $x,y\in\mathcal P$.
\end{proof}

\begin{lemma}\label{asm-theta-agree-definite-norms}
Assume the meromorphic continuation of Siegel Eisenstein series in the form of
Theorem~\ref{asm-theta-agree-eisenstein} and the Siegel-Weil formula in the
form of Theorem~\ref{asm-theta-agree-siegel-weil}. Let
$(B,G,\mathcal X,\omega,\omega',R,R_f)$ be a matched quadruple for a theta base
$B$ of a CM field $E$ (Definition~\ref{asm-theta-agree-matched}) with $p=q=0$,
let $\mu$ be a Haar measure on $U_W$ and $\mathcal D$ a fundamental domain for
$U_W(F)$ with respect to $\mu$, let $\mu'$ and $\mathcal D'$ be the measure and
the domain of $M'$, and let $\mathcal P$, $\mathcal I_{s,t}$, the products and
the theta lifts be as in Definition~\ref{asm-theta-agree-definite-products},
with $\mu$ and $\mathcal D$ on $U_W$ and $\mu'$ and $\mathcal D'$ on $U_{V'}$.
Let $s,t,\lambda\in\mathbf C$ be such that
$\langle\mathcal I_{s,t}(y),\mathcal I_{s,t}(y')\rangle_{\mathcal D}=\lambda\,\langle\mathcal I_{0,1}(y),\mathcal I_{0,1}(y')\rangle_{\mathcal D}$
for all $y,y'\in\mathcal P$. Then for every $y\in\mathcal P$ and every
$\Psi\in\operatorname{PD}(M')$,
\[
\langle\theta^*_{\omega'}(\mathcal I_{s,t}(y),\Psi),\theta^*_{\omega'}(\mathcal I_{s,t}(y),\Psi)\rangle_{\mathcal D'}
=\lambda\,\langle\theta^*_{\omega'}(\mathcal I_{0,1}(y),\Psi),\theta^*_{\omega'}(\mathcal I_{0,1}(y),\Psi)\rangle_{\mathcal D'} .
\]
\end{lemma}
\leanref{Development/WeilBound/ThetaLifts/LatticeBasics/PairsStable.lean}{GlobalTheory.ThK.Dbl.CoefLine}
\leanref{Development/WeilBound/ThetaLifts/LatticeBasics/CoefLine.lean}{GlobalTheory.ThK.Dbl.coefLine_of}
\leanref{Development/WeilBound/ThetaLifts/LatticeBasics/FarRight.lean}{GlobalTheory.ThK.Dbl.An.ident}
\begin{proof}
Let $y=(\Phi,\Phi')\in\mathcal P$ and $\Psi\in\operatorname{PD}(M')$. As in the
proof of Lemma~\ref{asm-theta-agree-definite-line}, for all
$s_1,t_1\in\mathbf C$ the function $\mathcal I_{s_1,t_1}(y)$ is continuous and
invariant under left translation by $U_W(F)$, so it is rapidly decreasing by
Lemma~\ref{asm-theta-agree-definite-compact}(5); and
$R_g\mathcal I_{s_1,t_1}(y)=\mathcal I_{s_1,t_1}(y_g)$ with
$y_g=(\omega(1,g)\Phi,\omega'(1,g)\Phi')$, which lies in $\mathcal P$ by
Lemma~\ref{asm-theta-agree-definite-inputs}(8).

Let $U$, $z_0$, $t_0$, $F_z$ and $Z_\phi$ be as in
Lemma~\ref{asm-theta-agree-definite-inputs}(10) for $\Psi$. For
$\operatorname{Re}z>t_0$, the hypothesis for the pairs $y_g$ and $y$ gives
\[
Z_{\mathcal I_{s,t}(y)}(z)=\int_{U_W}F_z(g)\,\langle\mathcal I_{s,t}(y_g),\mathcal I_{s,t}(y)\rangle_{\mathcal D}\,d\mu(g)
=\lambda\int_{U_W}F_z(g)\,\langle\mathcal I_{0,1}(y_g),\mathcal I_{0,1}(y)\rangle_{\mathcal D}\,d\mu(g)
=\lambda\,Z_{\mathcal I_{0,1}(y)}(z).
\]
The functions $Z_{\mathcal I_{s,t}(y)}$ and $\lambda Z_{\mathcal I_{0,1}(y)}$
are holomorphic on the connected open set $U$ and agree on the nonempty open
set of $z$ with $\operatorname{Re}z>t_0$, which lies in $U$. So they agree on
$U$, and in particular at $z_0$. By the value at $z_0$ in
Lemma~\ref{asm-theta-agree-definite-inputs}(10),
\[
\frac{2}{\mu'(\mathcal D')}\langle\theta^*_{\omega'}(\mathcal I_{s,t}(y),\Psi),\theta^*_{\omega'}(\mathcal I_{s,t}(y),\Psi)\rangle_{\mathcal D'}
=\lambda\,\frac{2}{\mu'(\mathcal D')}\langle\theta^*_{\omega'}(\mathcal I_{0,1}(y),\Psi),\theta^*_{\omega'}(\mathcal I_{0,1}(y),\Psi)\rangle_{\mathcal D'} .
\]
As $V'$ has signature $(0,a')$,
Lemma~\ref{asm-theta-agree-definite-compact}(3) gives
$0<\mu'(\mathcal D')<\infty$, and the factor $2/\mu'(\mathcal D')$ cancels.
\end{proof}

\begin{lemma}\label{asm-theta-agree-definite-constant}
Let $(B,G,\mathcal X,\omega,\omega',R,R_f)$ be a matched quadruple of ranks
$(a,a')$ for a theta base $B$ of a CM field $E$
(Definition~\ref{asm-theta-agree-matched}), put
$I_0=\{1,\ldots,a\}\times\{1,\ldots,a+a'\}$ and
$I'_0=\{1,\ldots,a'\}\times\{1,\ldots,a+a'\}$, let $\mu$ be a Haar measure on
$U_W$ and $\mathcal D$ a fundamental domain for $U_W(F)$ with respect to $\mu$,
and let $s,t\in\mathbf C$ be such that $\mathcal I_{s,t}(y)$ is rapidly
decreasing for every $y\in\mathcal P$
(Definition~\ref{asm-theta-agree-definite-products}). Let $y\in\mathcal P$.
Then $\theta^*_\omega(\mathcal I_{s,t}(y),\Psi)$ is constant on $U_V$ for
every Schwartz-Bruhat function $\Psi$ on $\mathbf A^{I_0}$, and
$\theta^*_{\omega'}(\mathcal I_{s,t}(y),\Psi)$ is constant on $U_{V'}$ for
every Schwartz-Bruhat function $\Psi$ on $\mathbf A^{I'_0}$.
\end{lemma}
\leanref{Development/WeilBound/ThetaLifts/LatticeBasics/PairsStable.lean}{GlobalTheory.ThK.Dbl.ThConst}
\leanref{Development/WeilBound/ThetaLifts/LatticeBasics/ThConstFin.lean}{GlobalTheory.ThK.Dbl.thConst_of'}
\leanref{Development/WeilBound/ThetaLifts/LatticeBasics/ThConst.lean}{GlobalTheory.ThK.Dbl.thConst_of}
\leanref{Development/WeilBound/ThetaLifts/LatticeBasics/ThConst.lean}{GlobalTheory.ThK.Dbl.TC.exists_maps}
\leanref{Development/WeilBound/ThetaLifts/LatticeBasics/ThConst.lean}{GlobalTheory.ThK.Dbl.TC.const_of}
\begin{proof}
\emph{Step 1: two linear maps.} The pairs $(\Phi,\Phi')$ that correspond form
a linear subspace $Q$: two corresponding pairs can be written with the same
finite set $T_1$, since $p_v(\mathrm{vac}_v)=p'_v(\mathrm{vac}'_v)$ for
$v\notin T$ allows enlarging $T_1$, and the sum of the identities in
Definition~\ref{asm-theta-agree-matched}(4), with the first multiplied by a
scalar, is the identity for the linear combination. So $\mathcal P$ is the
subspace of the $(\Phi,\Phi')\in Q$ with $\Phi'\in\mathcal S^\circ(M')$. By
Lemma~\ref{asm-theta-agree-decay}, every Schwartz-Bruhat function converges
for $(M,\omega)$ and every element of $\mathcal S^\circ(M')$ converges for
$(M',\omega')$; as theta functions are linear
(Lemma~\ref{asm-theta-agree-convergence-theta}(2)),
$y\mapsto\mathcal I_{s,t}(y)$ is linear on $\mathcal P$. If
$(0,\Phi')\in\mathcal P$, then $0$ and $\Phi'$ correspond, so $I'(\Phi')=0$ by
Lemma~\ref{asm-theta-agree-definite-inputs}(5), and
$\mathcal I_{s,t}(0,\Phi')=0$. Hence there is a linear map $C$ from the
functions on $\mathbf A^{I_0}$ to the functions on $U_W$ with
$C(\Phi)=\mathcal I_{s,t}(\Phi,\Phi')$ for every $(\Phi,\Phi')\in\mathcal P$:
define it on the image of the projection $(\Phi,\Phi')\mapsto\Phi$ of
$\mathcal P$ and extend it linearly. If $(\Phi,0)\in Q$, then
$(\Phi,0)\in\mathcal P$ and $C(\Phi)=s\,I(\Phi)=0$, by
Lemma~\ref{asm-theta-agree-definite-inputs}(5). So in the same way there is a
linear map $C'$ from the functions on $\mathbf A^{I'_0}$ to the functions on
$U_W$ with $C'(\Phi')=C(\Phi)$ whenever $\Phi$ and $\Phi'$ correspond.

\emph{Step 2: properties of $C$ and $C'$.} (a) For
$\Phi\in\operatorname{PD}(M)$ the function $C(\Phi)$ is rapidly decreasing: by
Lemma~\ref{asm-theta-agree-partners}(1) some $\Phi'\in\mathcal S^\circ(M')$
corresponds to $\Phi$, and $C(\Phi)=\mathcal I_{s,t}(\Phi,\Phi')$. For
$\Phi'\in\operatorname{PD}(M')$, some $\Phi$ corresponds to $\Phi'$ by
Lemma~\ref{asm-theta-agree-partners}(2); then $\Phi\in\operatorname{PD}(M)$ by
Definition~\ref{asm-theta-agree-matched}(4), and $C'(\Phi')=C(\Phi)$ is
rapidly decreasing.

(b) Let $v$ be a finite place such that $E$ has an $F$-embedding into $F_v$,
and $h\in U_W^{(v)}$. For $\Phi\in\operatorname{PD}(M)$, take
$\Phi'\in\mathcal S^\circ(M')$ corresponding to $\Phi$. By
Lemma~\ref{asm-theta-agree-definite-inputs}(6), $\omega(1,h)\Phi$ and
$\omega'(1,h)\Phi'$ correspond, so $\omega(1,h)\Phi\in\operatorname{PD}(M)$;
and $\omega'(1,h)\Phi'\in\mathcal S^\circ(M')$ by
Lemma~\ref{asm-theta-agree-coinvariants-inputs}(4). As $\omega$ and $\omega'$
are representations,
$C(\omega(1,h)\Phi)=\mathcal I_{s,t}(\omega(1,h)\Phi,\omega'(1,h)\Phi')=R_h\mathcal I_{s,t}(\Phi,\Phi')=R_hC(\Phi)$.
Likewise, for $\Phi'\in\operatorname{PD}(M')$ and $\Phi$ corresponding to it,
$\omega'(1,h)\Phi'\in\operatorname{PD}(M')$ and
$C'(\omega'(1,h)\Phi')=C(\omega(1,h)\Phi)=R_hC(\Phi)=R_hC'(\Phi')$.

\emph{Step 3: the constancy.} Let $(X,X',C_X)$ be $(V,V',C')$ or $(V',V,C)$,
and let $(X,W,M_X,\chi_X,\chi_W,\omega_X)$ and
$(X',W,M_{X'},\chi_{X'},\chi'_W,\omega_{X'})$ be the corresponding Weil pairs
of Definition~\ref{asm-theta-agree-triple}(4), with $X$ and $X'$ of ranks $n$
and $n'$. Then $n+n'=a+a'$, $\chi_W=\chi_X/\chi_{X'}$, which is $\beta/\alpha$
or $\alpha/\beta$, and $n\geq2$, as $a>a'>2(p+q)+1\geq1$. Let
$\Phi\in\operatorname{PD}(M_{X'})$ and let $\Psi$ be a Schwartz-Bruhat function
on $\mathbf A^{n\times(a+a')}$. We show that
$f(g)=\theta^*_{\omega_X}(C_X(\Phi),\Psi)(g)$ is constant on $U_X$. Write
$y=(\Phi_y,\Phi'_y)$; then $\Phi_y\in\operatorname{PD}(M)$ and
$\Phi'_y\in\operatorname{PD}(M')$ by
Definition~\ref{asm-theta-agree-matched}(4), and
$C(\Phi_y)=C'(\Phi'_y)=\mathcal I_{s,t}(y)$ by Step 1. So the case
$(X,X')=(V,V')$, $\Phi=\Phi'_y$ gives the first claim, and the case
$(X,X')=(V',V)$, $\Phi=\Phi_y$ gives the second.

Let $A_0$ and $k$ be as in Lemma~\ref{asm-theta-agree-definite-inputs}(2) for
$\Psi$. By Lemma~\ref{asm-theta-agree-anisotropic-automorphy},
$\theta_{\omega_X}(\Psi)(g,h)=\theta_{\omega_X}(\Psi)(g,\gamma h)$ for
$\gamma\in U_W(F)$, so
$|\theta_{\omega_X}(\Psi)(g,h)|\leq A_0\operatorname{ht}(g)^k\,\overline{\operatorname{ht}}(h)^k$.
As $C_X(\Phi)$ is rapidly decreasing by Step 2(a), there is $K_0$ with
$|C_X(\Phi)(h)\,\theta_{\omega_X}(\Psi)(g,h)|\leq K_0\operatorname{ht}(g)^k$.
By Theorem~\ref{asm-theta-agree-convergence-siegel}(1) and
Lemma~\ref{asm-theta-agree-convergence-signature}, $U_X$ is locally compact and
second countable and the height is bounded on its compact subsets, and
$\mu(\mathcal D)<\infty$ by Theorem~\ref{asm-theta-agree-definite-volume}. As
$\theta_{\omega_X}(\Psi)$ is continuous
(Lemma~\ref{asm-theta-agree-convergence-theta}(2)), dominated convergence shows
that $f$ is continuous, and the same bound shows that
$L(\Phi_1,\Psi_1)=\theta^*_{\omega_X}(C_X(\Phi_1),\Psi_1)(g)$ converges for
$\Phi_1\in\operatorname{PD}(M_{X'})$, Schwartz-Bruhat $\Psi_1$ and $g\in U_X$.
By Lemma~\ref{asm-theta-agree-anisotropic-automorphy}, $f(\gamma g)=f(g)$ for
$\gamma\in U_X(F)$.

Let $v\notin T$ be a finite place such that the rational prime below $v$
splits completely in $E$, that is, $v$ is clean for $\mathcal X$, and let
$u\in U_X^{(v)}$ and $g\in U_X$. By
Lemma~\ref{asm-theta-agree-definite-inputs}(7) there is an $F$-embedding
$\tau\colon E\to F_v$ with $\nu_\tau(\varpi_v)\neq1$, and then also
$(\nu^{-1})_\tau(\varpi_v)=\nu_\tau(\varpi_v)^{-1}\neq1$; so
$(\chi_W)_\tau(\varpi_v)\neq1$. For this $g$, the form $L$ above is linear in
$\Phi_1$ and conjugate-linear in $\Psi_1$; extend it to all functions by
composing with linear projections onto $\operatorname{PD}(M_{X'})$ and onto the
Schwartz-Bruhat functions. Let $h\in U_W^{(v)}$,
$\Phi_1\in\operatorname{PD}(M_{X'})$ and $\Psi_1$ Schwartz-Bruhat. By
Step 2(b), $\omega_{X'}(1,h)\Phi_1\in\operatorname{PD}(M_{X'})$ and
$C_X(\omega_{X'}(1,h)\Phi_1)=R_hC_X(\Phi_1)$; the function
$\omega_X(1,h)\Psi_1$ is Schwartz-Bruhat, and
$\theta_{\omega_X}(\omega_X(1,h)\Psi_1)(g,x)=\theta_{\omega_X}(\Psi_1)(g,xh)$.
So
\[
L(\omega_{X'}(1,h)\Phi_1,\omega_X(1,h)\Psi_1)=\int_{\mathcal D}C_X(\Phi_1)(xh)\,\overline{\theta_{\omega_X}(\Psi_1)(g,xh)}\,d\mu(x)=L(\Phi_1,\Psi_1),
\]
as the integrand is invariant under left translation by $U_W(F)$ and
integrable, $U_W$ is unimodular (Theorem~\ref{asm-theta-agree-definite-volume}),
and integrals of such functions over the fundamental domains $\mathcal D$ and
$\mathcal Dh$ agree, as in the proof of
Lemma~\ref{asm-theta-agree-anisotropic-invariance}. As every
$h_0\in GL_{a+a'}(F_v)$ is $\pi_\tau(h)$ for some $h\in U_W^{(v)}$, by
Lemma~\ref{asm-theta-agree-definite-inputs}(3),
$L(\Phi,\omega_X(u,1)\Psi)=L(\Phi,\Psi)$. As
$\theta_{\omega_X}(\omega_X(u,1)\Psi)(g,x)=\theta_{\omega_X}(\Psi)(gu,x)$, this
says $f(gu)=f(g)$.

By Lemma~\ref{asm-theta-agree-definite-inputs}(4) for $X$, of rank $n\geq2$,
and the set $T$, the function $f$ is constant.
\end{proof}
% Lean: Step 1 is TC.exists_maps (KDecay, Factor, exists_fac, pm_add, pm_zero); Step 2 is
% r1 and e1 in thConst_of (MatchEx, pm_fin, KCo.so_statement, pol_mem); Step 3 is
% TC.const_of, with Ker.cont, Ker.bd, TC.exists_form, th_right, theta_act, GoodNe and
% ConstOfGood. The first claim is const_of for (pairV, pairV', C'), the second for
% (pairV', pairV, C).

\begin{lemma}\label{asm-theta-agree-definite}
Assume the meromorphic continuation of Siegel Eisenstein series in the form of
Theorem~\ref{asm-theta-agree-eisenstein} and the Siegel-Weil formula in the
form of Theorem~\ref{asm-theta-agree-siegel-weil}. Let
$(B,G,\mathcal X,\omega,\omega',R,R_f)$ be a matched quadruple for a theta base
$B$ of a CM field $E$ (Definition~\ref{asm-theta-agree-matched}) with $p=q=0$.
Then it has a nonzero polynomial constant.
\end{lemma}
\leanref{Development/AbelianVarietyTheory/Definitions/Divisors/CTSmallCone.lean}{GlobalTheory.ThK.PolZero}
\leanref{Development/WeilBound/ThetaLifts/Reductions/TopPolarZero/PolZero.lean}{PositivityTheory.ThKR.polynomial_vanishing_of_compactness}
\leanref{Development/WeilBound/ThetaLifts/LatticeBasics/PolZero.lean}{GlobalTheory.ThK.Dbl.polZero_of}
\begin{proof}
The space $W$ is skew-Hermitian of signature $(a+a',0)$, and $V$ and $V'$ are
Hermitian of signatures $(a,0)$ and $(0,a')$. So
Lemma~\ref{asm-theta-agree-definite-compact} applies to $W$ and to $V'$. Let
$\mu$ be a Haar measure on $U_W$ and $\mathcal D$ a fundamental domain for
$U_W(F)$ with respect to $\mu$, which exist by
Lemma~\ref{asm-theta-agree-definite-compact}(3), and let $\mu'$ and
$\mathcal D'$ be the measure and the domain of $M'$; as $M'$ is regular, they
are a Haar measure and a fundamental domain, and
$0<\mu'(\mathcal D')<\infty$ by
Lemma~\ref{asm-theta-agree-definite-compact}(3). Let $\mathcal P$,
$\mathcal I_{s,t}$, the products $\langle\,,\rangle_{\mathcal D}$ on $U_W$ and
$\langle\,,\rangle_{\mathcal D'}$ on $U_{V'}$, and the theta lifts be as in
Definition~\ref{asm-theta-agree-definite-products}.

\emph{Step 1: the functions $\mathcal I_{s,t}(y)$.} Let
$y=(\Phi,\Phi')\in\mathcal P$. By Lemma~\ref{asm-theta-agree-partners}(3),
$\Phi\in\operatorname{PD}(M)$ and $\Phi'\in\operatorname{PD}(M')$, and both are
Schwartz-Bruhat, since the models are regular. By
Lemma~\ref{asm-theta-agree-definite-continuous}, applied to the Weil pairs
$(V,W,M,\beta,\beta/\alpha,\omega)$ and $(V',W,M',\alpha,\alpha/\beta,\omega')$,
whose first spaces have $\min(a,0)=\min(0,a')=0$, the functions $I(\Phi)$ and
$I'(\Phi')$ are continuous, and by
Lemma~\ref{asm-theta-agree-anisotropic-automorphy} they are invariant under left
translation by $U_W(F)$. Hence so is $\mathcal I_{s,t}(y)$ for all
$s,t\in\mathbf C$, and it is bounded and rapidly decreasing by
Lemma~\ref{asm-theta-agree-definite-compact}(4) and~(5). In particular all the
products $\langle\mathcal I_{s,t}(y),\mathcal I_{s',t'}(y')\rangle_{\mathcal D}$
converge.

\emph{Step 2: the candidate constant.} By
Lemma~\ref{asm-theta-agree-definite-line} for $(s,t,s',t')=(1,0,0,1)$ there is
$c\in\mathbf C$ with
$\langle I(\Phi_1),I'(\Phi'_2)\rangle_{\mathcal D}=c\,\langle I'(\Phi'_1),I'(\Phi'_2)\rangle_{\mathcal D}$
for all $(\Phi_1,\Phi'_1),(\Phi_2,\Phi'_2)\in\mathcal P$. Taking both pairs
equal to $y$ gives
\[
\langle\mathcal I_{1,-c}(y),\mathcal I_{0,1}(y)\rangle_{\mathcal D}
=\langle I(\Phi),I'(\Phi')\rangle_{\mathcal D}-c\,\langle I'(\Phi'),I'(\Phi')\rangle_{\mathcal D}=0
\qquad(y=(\Phi,\Phi')\in\mathcal P).
\]
By Lemma~\ref{asm-theta-agree-definite-line} for $(1,-c,1,-c)$ there is
$\lambda\in\mathbf C$ with
$\langle\mathcal I_{1,-c}(y),\mathcal I_{1,-c}(y')\rangle_{\mathcal D}=\lambda\,\langle\mathcal I_{0,1}(y),\mathcal I_{0,1}(y')\rangle_{\mathcal D}$
for all $y,y'\in\mathcal P$.

\emph{Step 3: a test pair.} By Lemma~\ref{asm-theta-agree-nonvanishing} there
are $\Phi_0\in\operatorname{PD}(M)$ with $I(\Phi_0)\neq0$, as $p=0$, and
$\Phi'\in\operatorname{PD}(M')\cap\mathcal S^\circ(M')$ with $I'(\Phi')\neq0$,
as $q=0$. By Lemma~\ref{asm-theta-agree-partners}(2) some $\Phi$ corresponds to
$\Phi'$, so $y_0=(\Phi,\Phi')\in\mathcal P$. Put
$x=\langle I'(\Phi'),I'(\Phi')\rangle_{\mathcal D}$. Then $x\neq0$ by
Lemma~\ref{asm-theta-agree-definite-compact}(4), as $I'(\Phi')\neq0$.

\emph{Step 4: theta lifts against $\Phi'$.} Let $s,t\in\mathbf C$ and
$\phi=\mathcal I_{s,t}(y_0)$. By Step 1 and
Lemma~\ref{asm-theta-agree-definite-constant}, the function
$\theta^*_{\omega'}(\phi,\Phi')$ on $U_{V'}$ is constant, say equal to $k$. The
function $(g,h)\mapsto\phi(h)\overline{\theta_{\omega'}(\Phi')(g,h)}$ on
$U_{V'}\times U_W$ is continuous and invariant under left translation by
$U_{V'}(F)\times U_W(F)$, by
Lemmas~\ref{asm-theta-agree-convergence-theta}(2)
and~\ref{asm-theta-agree-anisotropic-automorphy}, so it is bounded by
Lemma~\ref{asm-theta-agree-definite-compact}(4); and $\mathcal D'$ and
$\mathcal D$ have finite measure. By Fubini's theorem,
\[
\mu'(\mathcal D')\,k=\int_{\mathcal D'}\theta^*_{\omega'}(\phi,\Phi')(g)\,d\mu'(g)
=\int_{\mathcal D}\phi(h)\,\overline{\int_{\mathcal D'}\theta_{\omega'}(\Phi')(g,h)\,d\mu'(g)}\,d\mu(h)
=\langle\phi,I'(\Phi')\rangle_{\mathcal D}.
\]
So $\theta^*_{\omega'}(\mathcal I_{s,t}(y_0),\Phi')$ is the constant
$\mu'(\mathcal D')^{-1}\langle\mathcal I_{s,t}(y_0),\mathcal I_{0,1}(y_0)\rangle_{\mathcal D}$.

\emph{Step 5: $\lambda=0$.} By Step 2, Lemma~\ref{asm-theta-agree-definite-norms}
applies to $(s,t)=(1,-c)$ and $\lambda$; take $y_0$ and $\Psi=\Phi'$. By
Steps 4 and 2, $\theta^*_{\omega'}(\mathcal I_{1,-c}(y_0),\Phi')=0$, and by
Step 4, $\theta^*_{\omega'}(\mathcal I_{0,1}(y_0),\Phi')$ is the constant
$x/\mu'(\mathcal D')$. So
\[
0=\lambda\,\mu'(\mathcal D')\,\Bigl|\frac{x}{\mu'(\mathcal D')}\Bigr|^2,
\]
and $\lambda=0$, as $x\neq0$ and $0<\mu'(\mathcal D')<\infty$.

\emph{Step 6: conclusion.} Let $\Phi'_1\in\mathcal S^\circ(M')$ and $\Phi_1$
correspond, so that $z=(\Phi_1,\Phi'_1)\in\mathcal P$. By Steps 2 and 5,
$\langle\mathcal I_{1,-c}(z),\mathcal I_{1,-c}(z)\rangle_{\mathcal D}=0$, so
$\mathcal I_{1,-c}(z)=0$ by Step 1 and
Lemma~\ref{asm-theta-agree-definite-compact}(4). That is,
$I(\Phi_1)=c\,I'(\Phi'_1)$, and $c$ is a polynomial constant. By
Lemma~\ref{asm-theta-agree-partners}(1) some $\Psi'\in\mathcal S^\circ(M')$
corresponds to $\Phi_0$, so $I(\Phi_0)=c\,I'(\Psi')$. As $I(\Phi_0)\neq0$,
$c\neq0$.
\end{proof}
% Lean: ThKR.polynomial_vanishing_of_compactness is polZero_of (LatticeBasics/PolZero.lean)
% with CptAll = cptAll_of GrpTop EpsOf SiegelAll FinVolU Unimod, HermLine = hermLine_of'
% (cptOpen, polRank, pairsStable, pairsIrr), CoefLine = coefLine_of (swFull_of, boxSec_of,
% dblW, hgtEx_of, hgtInt_of, stab2Cpt) and ThConst = thConst_of' (locDist, goodNe,
% SA.constOfGood HeckeLEntireAndNonzeroAtOne_holds). Step 4 is key with int_th; Step 5 is l0;
% Step 6 is main. PolSch (polSch) gives that polynomial data are Schwartz-Bruhat. Lemma
% asm-theta-agree-definite-inputs, which supplies the inputs, is proved.

\subsection{Local cones and local factors}\label{subsection-asm-local-factors}

\begin{definition}\label{asm-theta-agree-local-cones}
Let $E$ and $F$ be as in Definition~\ref{asm-theta-agree-models}. For a model
$M$ for lattice spaces $(X,W)$ of ranks $(n,d)$ put
$I_M=\{1,\ldots,n\}\times\{1,\ldots,d\}$. For a finite place $v$ and a finite
set $I$ let $\mathcal S(F_v^I)$ be the space of Schwartz-Bruhat functions on
$F_v^I$. Let $1_v\in\mathbf A$ be $1$ at $v$ and $0$ at the other places, and
$1_\infty\in\mathbf A$ be $1$ at the infinite places and $0$ at the finite
ones. Let $M$ be a model for $(X,W)$ with Schrödinger operators $\rho$, and
$\omega$ a representation of $U_X\times U_W$ on functions on $\mathbf A^{I_M}$.
\begin{enumerate}
\item Let $T$ be an operator on functions on $\mathbf A^{I_M}$. For a finite
  place $v$ and a function $\phi$ on $F_v^{I_M}$, let $T_v\phi$ be a
  $\phi_1\in\mathcal S(F_v^{I_M})$ such that $T$ acts at $v$ by
  $\phi\mapsto\phi_1$ (Definition~\ref{asm-theta-agree-models}(2)), if there is
  one, and $0$ otherwise. For $f\in\mathcal S(F_\infty^{I_M})$ let $T_\infty f$
  be an $f_1\in\mathcal S(F_\infty^{I_M})$ such that $T$ acts at $\infty$ by
  $f\mapsto f_1$, if there is one, and $0$ otherwise. In both cases the choice
  depends only on the set of such $\phi_1$, respectively $f_1$.
\item Let $K\subset U_X$ be a subgroup. For a finite place $v$, or
  $v=\infty$, let $K_v$ be the intersection of $K$ with the subgroup generated
  by $U_X^{(v)}$, and $A_v$ the subgroup of $\mathbf A_E^\times$ generated by
  the $l$ with $l-1$ at $v$. A \emph{local measure system} $m$ for $K$
  consists of measures $m_{K,v}$ on $K_v$ and $m_{A,v}$ on $A_v$, for every
  finite place $v$ and for $v=\infty$. It is \emph{normalized} if for every
  finite $v$ the measure $m_{K,v}$ is a Haar measure of total mass $1$ and
  $m_{A,v}$ is a Haar measure that gives mass $1$ to the set of $l\in A_v$
  whose powers $l^k$, $k\in\mathbf Z$, have compact closure, and if
  $m_{K,\infty}$ and $m_{A,\infty}$ are Haar measures.
\item Let $m$ be a local measure system for a subgroup $K\subset U_X$, let
  $u\colon\mathbf A_E^\times\to\mathbb W(X,W)$ be a map, $\chi$ a character of
  $\mathbf A_E^\times$, $z\geq0$ an integer and $I_1$ a finite set. For
  $k\in U_X$, $l\in\mathbf A_E^\times$ and $v$ a finite place or $\infty$, put
  $T^v_{k,l}=\rho(1_vu(l))\,\omega(k,1)$. For a finite place $v$ and a linear
  map $e_v$ from functions on $F_v^{I_M}$ to functions on $F_v^{I_1}$, the
  \emph{local cone} for $(M,\omega,m,e_v,u,\chi,z)$ of a function $\phi$ on
  $F_v^{I_M}$ is
  \[
  \operatorname{cone}_v(\phi)(y)=\int_{K_v}\int_{A_v}e_v\bigl((T^v_{k,l})_v\phi\bigr)(y)\,\chi(l)|l|^z\,dm_{A,v}(l)\,dm_{K,v}(k).
  \]
  For a continuous linear map $e_\infty\colon\mathcal S(F_\infty^{I_M})\to\mathcal S(F_\infty^{I_1})$,
  a continuous linear map $L$ between the same spaces is an \emph{archimedean
  cone} for $(M,\omega,m,e_\infty,u,\chi,z)$ if for all $f$ and $y$
  \[
  L(f)(y)=\int_{K_\infty}\int_{A_\infty}e_\infty\bigl((T^\infty_{k,l})_\infty f\bigr)(y)\,\chi(l)|l|^z\,dm_{A,\infty}(l)\,dm_{K,\infty}(k).
  \]
\end{enumerate}
\end{definition}
\leanref{Development/WeilBound/ThetaLifts/LocalTheory/IsArchimedeanLocalEquivariantMap.lean}{PositivityTheory.ThKG.Loc.opF}
\leanref{Development/WeilBound/ThetaLifts/LocalTheory/IsArchimedeanLocalEquivariantMap.lean}{PositivityTheory.ThKG.Loc.opI}
\leanref{Development/WeilBound/ThetaLifts/LocalTheory/IsArchimedeanLocalEquivariantMap.lean}{PositivityTheory.ThKG.Loc.kAt}
\leanref{Development/WeilBound/ThetaLifts/LocalTheory/IsArchimedeanLocalEquivariantMap.lean}{PositivityTheory.ThKG.Loc.idAt}
\leanref{Development/WeilBound/ThetaLifts/LocalTheory/IsArchimedeanLocalEquivariantMap.lean}{PositivityTheory.ThKG.Loc.LocMeas}
\leanref{Development/WeilBound/ThetaLifts/LocalTheory/IsArchimedeanLocalEquivariantMap.lean}{PositivityTheory.ThKG.Loc.LocMeas.Is}
\leanref{Development/WeilBound/ThetaLifts/LocalTheory/IsArchimedeanLocalEquivariantMap.lean}{PositivityTheory.ThKG.Loc.coneF}
\leanref{Development/WeilBound/ThetaLifts/LocalTheory/IsArchimedeanLocalEquivariantMap.lean}{PositivityTheory.ThKG.Loc.IsConeI}

\begin{lemma}\label{asm-theta-agree-local-measures}
Let $B$ be a theta base of a CM field $E$, let $X$ be a Hermitian lattice
space of some signature, and let $K\subset U_X$ be the lattice stabilizer of
$X$ (Definition~\ref{asm-theta-agree-planes}(6)). Then $K$ has a normalized
local measure system (Definition~\ref{asm-theta-agree-local-cones}(2)).
\end{lemma}
\leanref{Development/WeilBound/ThetaLifts/LocalTheory/UnramifiedDeterminedBySpherical.lean}{PositivityTheory.ThKG.Loc.MeasEx}
\leanref{Development/WeilBound/ThetaLifts/LocalTheory/ConeShifts/ContinuousFc.lean}{PositivityTheory.ThKG.Loc.Cone.measEx_of}
\begin{proof}
Let $v$ be a finite place or $v=\infty$. An element $z$ of $\mathbf A_E$, or a
matrix over $\mathbf A_E$, is at $v$ if and only if $(1-1_v)z=0$. So
$U_X^{(v)}$ is the set of $g\in U_X$ with $(1-1_v)(g-1)=0$, which is closed, and
it is a subgroup, since $gh-1=(g-1)h+(h-1)$ and $g^{-1}-1=-(g-1)g^{-1}$. Hence
$K_v=K\cap U_X^{(v)}$. In the same way the $l\in\mathbf A_E^\times$ with $l-1$ at
$v$ form a closed subgroup, which is $A_v$.

\emph{The measures on $K_v$.} By Lemma~\ref{asm-theta-agree-stabilizer}, $K$ is
compact, and $U_X$ is Hausdorff by
Theorem~\ref{asm-theta-agree-convergence-siegel}(1). So $K_v$, a closed subgroup
of $K$, is a compact group, and it has a Haar measure $m_{K,v}$ of total mass
$1$. This holds for $v=\infty$ as well.

\emph{The measures on $A_v$, $v$ finite.} The group $\mathbf A_E^\times$ is
locally compact, so its closed subgroup $A_v$ is too. Let $U_v$ be the set of
$l\in A_v$ with $|l_w|_w=1$ for every finite place $w$ of $E$. It is a subgroup
of $A_v$. It is open, as the $l\in\mathbf A_E^\times$ such that $l$ and $l^{-1}$
have integral components at every finite place form an open set. It is
compact: an element of $U_v$ is $1$ at the places of $E$ not above $v$ and a
unit at those above $v$, so $U_v$ is the image of the compact group
$\prod_{w\mid v}\mathcal O_w^\times$. A Haar measure on $A_v$ gives the compact
open set $U_v$ a finite positive mass; rescaling, let $m_{A,v}$ be the Haar
measure with $m_{A,v}(U_v)=1$.

It remains to show that $U_v$ is the set of $l\in A_v$ whose powers $l^k$,
$k\in\mathbf Z$, have compact closure. If $l\in U_v$, its powers lie in $U_v$,
which is compact and closed. Conversely, let the closure $C$ of the powers of
$l$ be compact. The cosets of the open subgroup $U_v$ form an open cover of
$C$, so the powers $l^n$, $n\geq0$, meet only finitely many of them, and
$l^iU_v=l^jU_v$ for some $i<j$. Then $l^k\in U_v$ for $k=j-i>0$, so
$|l_w|_w^k=1$ and $|l_w|_w=1$ for every finite place $w$, and $l\in U_v$.

\emph{The measure on $A_\infty$.} It is a closed subgroup of
$\mathbf A_E^\times$, hence locally compact; let $m_{A,\infty}$ be a Haar
measure on it.
\end{proof}
% Lean: measEx_of StabEx (ConeShifts/ContinuousFc.lean). K is compact by stab_compact;
% kAt_eq and idAt_eq describe K_v and A_v as the fixed points of the projections to v
% (closed by sysU and sysI); co_U and exists_haar give m_{K,v}; isOpen_UI, cpt_I, co_I and
% exists_haar give m_{A,v}; units_eq, with pow_mem_of_compact, is the last paragraph but one.

\begin{lemma}\label{asm-theta-agree-local-real-unique-diagonal}
Let $F$, $\mathbf A$ and $\psi$ be as in
Theorem~\ref{asm-theta-agree-evaluation-heisenberg}, let $\kappa$ be a finite
set, and let $\rho$ be the Schrödinger operators of that theorem on functions
on $\mathbf A^\kappa\times\mathbf A^\kappa=\mathbf A^{\kappa\sqcup\kappa}$. Then
there are an admissible linear form $l$ on functions on
$\mathbf A^{\kappa\sqcup\kappa}$ and a Schwartz-Bruhat function $\phi_0$ with
$l(\phi_0)\neq0$ such that $l(\rho(w)\phi)=l(\phi)$ for every Schwartz-Bruhat
$\phi$ and every $w=((x,x),(y,-y))$ with $x,y\in\mathbf A^\kappa$.
\end{lemma}
\leanref{Development/AbelianVarietyTheory/Proofs/DivisorProofs/ExistsDiagonalLagrangianInvariantForm.lean}{GlobalTheory.diagLagrangian_form}
\begin{proof}
Let $dy$ be a Haar measure on $\mathbf A^\kappa$. For a Schwartz-Bruhat $\phi$
on $\mathbf A^{\kappa\sqcup\kappa}$ the function $y\mapsto\phi(y,y)$ is a
Schwartz-Bruhat function on $\mathbf A^\kappa$, hence integrable. Let $l$ be a
linear form on all functions on $\mathbf A^{\kappa\sqcup\kappa}$ with
$l(\phi)=\int\phi(y,y)\,dy$ whenever $y\mapsto\phi(y,y)$ is integrable; it
exists since these $\phi$ form a subspace on which the integral is linear. For
$\phi=f\otimes g$, $l(\phi)=\int_{F_\infty^\kappa}f(u,u)\,du\cdot\int_{\mathbf A_f^\kappa}g(v,v)\,dv$,
and $f\mapsto\int f(u,u)\,du$ is continuous on the Schwartz space; so $l$ is
admissible. For $w=((x,x),(y,-y))$, the dot products $(u,u)\cdot(y,-y)$ and
$(x,x)\cdot(y,-y)$ vanish, so $(\rho(w)\phi)(u,u)=\phi(u+x,u+x)$, and
$l(\rho(w)\phi)=l(\phi)$ by the translation invariance of $dy$. Let $f$ be a
smooth function with compact support on $F_\infty^{\kappa\sqcup\kappa}$, with
$f\geq0$ and $f(0)=1$, and $g$ the indicator function of the integral vectors
in $\mathbf A_f^{\kappa\sqcup\kappa}$. Then $\phi_0=f\otimes g$ has
$l(\phi_0)>0$: the first factor is the integral of a nonnegative continuous
function that is positive at $0$, and the second is the volume of the compact
open set of integral vectors in $\mathbf A_f^\kappa$.
\end{proof}
% Lean: diagLagrangian_form with DiagForm.lDiag = Lext . resD (the restriction to the
% diagonal and an extension of the integral, LinearMap.exists_extend); lDiag_adm, lDiag_schr
% and exists_lDiag_ne (a ContDiffBump times the indicator of the integral finite adeles).

\begin{lemma}\label{asm-theta-agree-local-real-unique-cayley}
Let $E$ be a CM field with complex conjugation $c$, let $Y$ be a lattice space
of rank $n$ with matrix $J$ (Definition~\ref{asm-theta-agree-doubling}(1)), and
let $G_Y$ and $\Delta$ be as in Definition~\ref{asm-theta-agree-doubling}(2).
Let $g\in U_Y$ be such that $t=\tfrac12(g+1)$ is invertible, and put
$T=\tfrac14(g-1)t^{-1}$,
\[
n=\begin{pmatrix}1+T&-T\\T&1-T\end{pmatrix},\qquad
w=\begin{pmatrix}1&0\\0&-1\end{pmatrix}.
\]
Then $\operatorname{diag}(g,1)$, $n$ and $w$ lie in $G_Y$, $n\Delta=\Delta$,
and $\operatorname{diag}(g,1)=wnw^{-1}p$ with $p\in G_Y$ and $p\Delta=\Delta t$.
\end{lemma}
\leanref{Development/WeilBound/ThetaLifts/LocalTheory/ToralCohomology/DarbouxExt.lean}{PositivityTheory.ThKG.CayleyTransformDarboux.cayley}
\begin{proof}
Put $J^\square=\operatorname{diag}(J,-J)$. As $g^{\mathsf T}Jc(g)=J$,
$\operatorname{diag}(g,1)^{\mathsf T}J^\square c(\operatorname{diag}(g,1))=J^\square$,
and clearly $w^{\mathsf T}J^\square c(w)=J^\square$. Next,
$T^{\mathsf T}J+Jc(T)=0$: as $t$ is invertible it suffices that
$t^{\mathsf T}(T^{\mathsf T}J+Jc(T))c(t)=(Tt)^{\mathsf T}Jc(t)+t^{\mathsf T}Jc(Tt)$
vanishes, and with $Tt=\tfrac14(g-1)$, $t=\tfrac12(g+1)$ and $c(\tfrac12)=\tfrac12$
this is
$\tfrac18\bigl((g-1)^{\mathsf T}Jc(g+1)+(g+1)^{\mathsf T}Jc(g-1)\bigr)=\tfrac18(2g^{\mathsf T}Jc(g)-2J)=0$.
Using this relation, the four blocks of $n^{\mathsf T}J^\square c(n)$ are $J$,
$0$, $0$ and $-J$; for instance the upper left block is
$(1+T)^{\mathsf T}J(1+c(T))-T^{\mathsf T}Jc(T)=J+T^{\mathsf T}J+Jc(T)=J$. The
matrix $n$ has inverse
$\left(\begin{smallmatrix}1-T&T\\-T&1+T\end{smallmatrix}\right)$, so
$n\in G_Y$, and $n\Delta=\Delta$.

Put $p=(wnw^{-1})^{-1}\operatorname{diag}(g,1)\in G_Y$. As
$wnw^{-1}=\left(\begin{smallmatrix}1+T&T\\-T&1-T\end{smallmatrix}\right)$,
$wnw^{-1}\Delta t$ has blocks $(1+2T)t=t+\tfrac12(g-1)=g$ and
$(1-2T)t=t-\tfrac12(g-1)=1$, which are the blocks of
$\operatorname{diag}(g,1)\Delta$. So $p\Delta=\Delta t$.
\end{proof}

\begin{lemma}\label{asm-theta-agree-local-real-unique}
Let $B$ be a theta base of a CM field $E$, and let
$(X,W,M,\chi_X,\chi_W,\omega)$ and $(X,W,M,\chi'_X,\chi'_W,\omega')$ be Weil
pairs of the same signature for $B$ (Definition~\ref{asm-theta-agree-triple}(1)).
Suppose that there are integers $k(w)$, for the infinite places $w$ of $F$,
with $\chi_W(z)=\chi'_W(z)=\prod_w(\pi_w(z)/|\pi_w(z)|)^{k(w)}$ for every
$z\in\mathbf A_E^\times$ whose finite components are $1$. Then
$\omega(g,1)\Phi=\omega'(g,1)\Phi$ for every $g\in U_X^{(\infty)}$ and every
Schwartz-Bruhat function $\Phi$.
\end{lemma}
\leanref{Development/WeilBound/ThetaLifts/LocalTheory/UnramifiedDeterminedBySpherical.lean}{PositivityTheory.ThKG.Loc.RealUniq}
\leanref{Development/WeilBound/ThetaLifts/LocalTheory/ToralCohomology/DarbouxExt.lean}{PositivityTheory.ThKG.CayleyTransformDarboux.realUniq}
\begin{proof}
Let $n$ and $d$ be the ranks of $X$ and $W$, $I=\{1,\ldots,n\}\times\{1,\ldots,d\}$,
and $e$ the isomorphism of $M$. By condition (b) of
Definition~\ref{asm-theta-agree-models}(6) for $\omega$, there are operators
$\omega^\square(x)$, $x\in G_X$, on functions on
$\mathbf A^I\times\mathbf A^I=\mathbf A^{I\sqcup I}$ that satisfy condition (b) of
Definition~\ref{asm-theta-agree-doubling}(6) with the character $\chi_W$ and
$m=d$, for the isomorphism $e^\square$ of $M_{d\times2n}(\mathbf A_E)$ given by
$e^\square(w_1\ w_2)=\bigl((e_1(w_1^{\mathsf T}),e_1(w_2^{\mathsf T})),(e_2(w_1^{\mathsf T}),-e_2(w_2^{\mathsf T}))\bigr)$,
and that satisfy
$\omega^\square(\operatorname{diag}(g,1))(\phi_1\boxtimes\phi_2)=(\omega(g,1)\phi_1)\boxtimes\phi_2$
for $g\in U_X$ and Schwartz-Bruhat $\phi_1,\phi_2$. The same holds for $\omega'$,
with operators $\omega'^\square(x)$ and the character $\chi'_W$; the
isomorphism $e^\square$ is the same, as it is determined by $e$. Let
$\rho^\square(w)=\rho(e^\square(w))$ be its Schrödinger operators, and let $l$
and $\phi_0$ be as in Lemma~\ref{asm-theta-agree-local-real-unique-diagonal}
for $\kappa=I$.

\emph{Step 1: a scalar for each $x\in G_X$.} Let $x\in G_X$ and
$S=\omega^\square(x^{-1})\omega'^\square(x)$. It maps Schwartz-Bruhat
functions to Schwartz-Bruhat functions, $\lambda\circ S$ is admissible for
every admissible $\lambda$, and for every $w$ and Schwartz-Bruhat $\phi$,
\[
S\rho^\square(w)\phi=\omega^\square(x^{-1})\rho^\square(wx^{\mathsf T})\omega'^\square(x)\phi
=\rho^\square(w)S\phi.
\]
As $e^\square$ is onto, Theorem~\ref{asm-theta-agree-evaluation-heisenberg}(5)
gives $s(x)\in\mathbf C$ with $S\phi=s(x)\phi$, so
$\omega'^\square(x)\phi=s(x)\,\omega^\square(x)\phi$ for Schwartz-Bruhat $\phi$.

\emph{Step 2: $s$ is multiplicative.} As $l(\phi_0)\neq0$, $\phi_0\neq0$, and so
$\omega^\square(x)\phi_0\neq0$ for every $x$, since
$\omega^\square(x^{-1})\omega^\square(x)\phi_0=\phi_0$. For $x,y\in G_X$,
$\omega'^\square(xy)\phi_0=\omega'^\square(x)\omega'^\square(y)\phi_0=s(x)s(y)\,\omega^\square(xy)\phi_0$,
and it is also $s(xy)\,\omega^\square(xy)\phi_0$. So $s(xy)=s(x)s(y)$, and
likewise $s(1)=1$.

\emph{Step 3: $s(p)=1$ for $p\in P_X$ with
$\chi_W(\det t(p))=\chi'_W(\det t(p))$.} For $z\in M_{d\times n}(\mathbf A_E)$,
$e^\square(z\Delta^{\mathsf T})=e^\square(z\ z)$ has the form $((x,x),(y,-y))$,
so $l(\rho^\square(z\Delta^{\mathsf T})\phi)=l(\phi)$ for Schwartz-Bruhat
$\phi$. By the rule for $P_X$ in Definition~\ref{asm-theta-agree-doubling}(6)(b),
with $u=t(p)$,
\[
s(p)\,l(\omega^\square(p)\phi_0)=l(\omega'^\square(p)\phi_0)=\chi'_W(\det u)|\det u|^{d/2}l(\phi_0)
=\chi_W(\det u)|\det u|^{d/2}l(\phi_0)=l(\omega^\square(p)\phi_0),
\]
and this is nonzero. So $s(p)=1$.

\emph{Step 4: the Cayley transform.} Let $g\in U_X$ be such that
$t=\tfrac12(g+1)$ is invertible and $\chi_W(\det t)=\chi'_W(\det t)$. We show
$\omega(g,1)\Phi=\omega'(g,1)\Phi$ for Schwartz-Bruhat $\Phi$. Let $w$, $n$ and
$p$ be as in Lemma~\ref{asm-theta-agree-local-real-unique-cayley}. Then $n$
and $p$ lie in $P_X$, with $t(n)=1$ and $t(p)=t$, so $s(n)=s(p)=1$ by Step 3,
and by Step 2
$s(\operatorname{diag}(g,1))=s(w)s(n)s(w)^{-1}s(p)=1$. Let $\phi_2$ be a
Schwartz-Bruhat function on $\mathbf A^I$ with $\phi_2(0)\neq0$. Then
$\Phi\boxtimes\phi_2$ is Schwartz-Bruhat, and by Step 1
\[
(\omega'(g,1)\Phi)\boxtimes\phi_2=\omega'^\square(\operatorname{diag}(g,1))(\Phi\boxtimes\phi_2)
=\omega^\square(\operatorname{diag}(g,1))(\Phi\boxtimes\phi_2)=(\omega(g,1)\Phi)\boxtimes\phi_2.
\]
Evaluating at $(x,0)$ and dividing by $\phi_2(0)$ gives the claim.

\emph{Step 5: elements in general position.} Let $g\in U_X$ be at $\infty$
with $\det(\pi_w(g)+1)\neq0$ for every infinite place $w$. Then Step 4 applies
to $g$. Indeed, $t=\tfrac12(g+1)=1+\tfrac12(g-1)$, and $g-1$ is at $\infty$, so
$\det t-1$ is at $\infty$, and $\pi_w(\det t)=2^{-n}\det(\pi_w(g)+1)\neq0$ for
every $w$. An element of $\mathbf A_E$ that is $1$ at the finite places and has
nonzero components at the infinite places is a unit, so $t$ is invertible;
and as the finite components of $\det t$ are $1$, the hypothesis gives
$\chi_W(\det t)=\prod_w(\pi_w(\det t)/|\pi_w(\det t)|)^{k(w)}=\chi'_W(\det t)$.

\emph{Step 6: conclusion.} Let $g\in U_X^{(\infty)}$. For each $w$, the
polynomial $q_w(X)=\det\bigl((1+iX)\pi_w(g)+(1-iX)\bigr)$ is nonzero, as
$q_w(i)=2^n$. Choose an integer $k$ that is a root of no $q_w$, and put
$z=(1+ik)/(1-ik)$. Then $|z|=1$, $\bar z+1=2/(1+ik)\neq0$, and
$\det(z\pi_w(g)+1)=(1-ik)^{-n}q_w(k)\neq0$ for every $w$. As in fact (iii) in
the proof of Lemma~\ref{asm-theta-agree-real-lines}, there is
$u\in\mathbf A_E$ with $u-1$ at $\infty$ and $\pi_w(u)=z$ for every $w$;
then $uc(u)=1$, since $uc(u)-1$ is at $\infty$ with components
$z\bar z-1=0$ (fact (i) there). So $u$ is a unit, and $\zeta=u\cdot1_n$ lies in
$U_X$. The elements $\zeta g$ and
$\zeta^{-1}$ are at $\infty$, as $\zeta g-1=u(g-1)+(u-1)$, and
$\det(\pi_w(\zeta g)+1)=\det(z\pi_w(g)+1)\neq0$ and
$\det(\pi_w(\zeta^{-1})+1)=(\bar z+1)^n\neq0$. By Step 5 for both, and as
$\omega'(\zeta g,1)\Phi$ is Schwartz-Bruhat,
\[
\omega(g,1)\Phi=\omega(\zeta^{-1},1)\,\omega(\zeta g,1)\Phi
=\omega(\zeta^{-1},1)\,\omega'(\zeta g,1)\Phi
=\omega'(\zeta^{-1},1)\,\omega'(\zeta g,1)\Phi=\omega'(g,1)\Phi.
\]
\end{proof}
% Lean: CayleyTransformDarboux.realUniq (ToralCohomology/DarbouxExt.lean). darboux_ext
% identifies the two doubled coordinates; Steps 1 to 4 are weilB_eq (with ThKF0.schur for
% Step 1, diagLagrangian_form for l, par for Step 3, unitary_diag_one_mem_dbl and cayley
% for Step 4, exists_mem_schA_ne_zero and boxFun_mem_schA for phi_2); Step 5 is good and
% good_unit (with exists_unit); Step 6 is exists_rot and the unit u, with g1 and g2.

\begin{lemma}\label{asm-theta-agree-local-cone-unique}
Let $E$ and $F$ be as in Definition~\ref{asm-theta-agree-models}, let $M$ be a
model for lattice spaces $(X,W)$, and let $\omega$ and $\omega'$ be
representations of $U_X\times U_W$ on functions on $\mathbf A^{I_M}$ with
$\omega(g,1)\Phi=\omega'(g,1)\Phi$ for $g\in U_X^{(\infty)}$ and Schwartz-Bruhat
$\Phi$. Let $\chi$ and $\chi'$ be characters of $\mathbf A_E^\times$ with
$\chi(l)=\chi'(l)$ whenever $l-1$ is at $\infty$. Let $m$ be a local measure
system for a subgroup $K\subset U_X$, $e_\infty$ a continuous linear map from
$\mathcal S(F_\infty^{I_M})$ to a Schwartz space, $u\colon\mathbf A_E^\times\to\mathbb W(X,W)$
a map and $z\geq0$ an integer. If $L$ is an archimedean cone for
$(M,\omega,m,e_\infty,u,\chi,z)$ and $L'$ one for
$(M,\omega',m,e_\infty,u,\chi',z)$
(Definition~\ref{asm-theta-agree-local-cones}(3)), then $L=L'$.
\end{lemma}
\leanref{Development/WeilBound/ThetaLifts/LocalTheory/MemPolData.lean}{PositivityTheory.ThKG.Loc.coneI_eq}
\begin{proof}
The elements of $U_X$ at $\infty$ form a subgroup, since
$gh-1=(g-1)h+(h-1)$ and $g^{-1}-1=-(g-1)g^{-1}$; so $K_\infty\subset U_X^{(\infty)}$.
In the same way every $l\in A_\infty$ has $l-1$ at $\infty$. Let $k\in K_\infty$
and $l\in A_\infty$. The operators $T^\infty_{k,l}$ for $\omega$ and for
$\omega'$ agree on Schwartz-Bruhat functions, and whether an operator acts at
$\infty$ by $f\mapsto f_1$ depends only on its values on Schwartz-Bruhat
functions. So the two operators at $\infty$ have the same values, and
$\chi(l)=\chi'(l)$. Hence the integrands defining $L(f)(y)$ and $L'(f)(y)$
agree, and $L=L'$.
\end{proof}

\begin{lemma}\label{asm-theta-agree-local-factor-sums}
Let $(\Omega,\mu)$ be a measure space, $Y$ a set and $G$ a map from $\Omega$ to
the complex functions on $Y$ that takes finitely many values, each on a
measurable set. Let $D\subset\Omega$ be a set with $G(a)=0$ for $a\notin D$, and
$w$ a complex function on $\Omega$ that is integrable on $D$. Then for every
$y\in Y$ the function $a\mapsto G(a)(y)w(a)$ is integrable, the function
\[
\textstyle\int Gw\colon y\mapsto\int_\Omega G(a)(y)\,w(a)\,d\mu(a)
\]
lies in every subspace that contains all the $G(a)$, and
$l(\int Gw)=\int_\Omega l(G(a))\,w(a)\,d\mu(a)$ for every linear form $l$ on
the functions on $Y$.
\end{lemma}
\leanref{Development/WeilBound/ThetaLifts/LocalTheory/ConeShifts/PureCancel.lean}{PositivityTheory.ThKG.Loc.Cone.Fin.int_fin}
\leanref{Development/WeilBound/ThetaLifts/LocalTheory/UnramifiedCohomology/SpanSoLoc.lean}{PositivityTheory.ThKG.Loc.FL.int_lin}
\begin{proof}
Let $g_1,\ldots,g_r$ be the nonzero values of $G$ and $\Omega_i=G^{-1}(g_i)$.
These sets are measurable and contained in $D$, so the functions
$\mathbf 1_{\Omega_i}w$ are integrable. As
$G(a)(y)w(a)=\sum_ig_i(y)\mathbf 1_{\Omega_i}(a)w(a)$, the first claim holds,
and $\int Gw=\sum_ic_ig_i$ with $c_i=\int_{\Omega_i}w\,d\mu$. This lies in every
subspace containing the $g_i$. For the last claim, let $Y'=Y\sqcup\{*\}$ and
apply the second claim to the functions $G'(a)$ on $Y'$ equal to $G(a)$ on $Y$
and to $l(G(a))$ at $*$, and to the subspace of functions $f$ on $Y'$ with
$f(*)=l(f|_Y)$. It gives
$\int_\Omega l(G(a))w(a)\,d\mu(a)=l(\int Gw)$.
\end{proof}

\begin{lemma}\label{asm-theta-agree-local-factor-evaluation}
Let $B$ be a theta base of a CM field $E$, with character $\psi$, and $M$ a
model for lattice spaces $(X,W)$, with isomorphism $e^M=(e^M_1,e^M_2)$,
Schrödinger operators $\rho$ and form $\langle\,,\rangle$ on $\mathbb W(X,W)$
(Definition~\ref{asm-theta-agree-models}(3)). For a finite place $v$ and
$w\in\mathbb W(X,W)$ at $v$, let $\rho_v(w)$ be the operator
\[
(\rho_v(w)\phi)(\xi)=\psi\bigl(\bar\xi\cdot e^M_2(w)+\tfrac12e^M_1(w)\cdot e^M_2(w)\bigr)\,\phi\bigl(\xi+e^M_1(w)_v\bigr)
\]
on functions on $F_v^{I_M}$, where $\bar\xi\in\mathbf A^{I_M}$ has component
$\xi$ at $v$ and $0$ elsewhere.
\begin{enumerate}
\item Let $w,w'\in\mathbb W(X,W)$ be at $v$. Then $\rho_v(w)$ maps
  Schwartz-Bruhat functions to Schwartz-Bruhat functions, $\rho(w)$ acts at
  $v$ by $\phi\mapsto\rho_v(w)\phi$, and
  \[
  \rho_v(w)\rho_v(w')=\psi(\tfrac12\langle w,w'\rangle)\,\rho_v(w+w')=\psi(\langle w,w'\rangle)\,\rho_v(w')\rho_v(w).
  \]
\item Let $M$ be regular, let $M_1$ be a regular model for lattice spaces
  $(X_1,W_1)$, and let $e$ be an evaluation map for
  $(M,M_1,\iota_X,\iota_W,\mathcal J)$, for some $\iota_X$, $\iota_W$ and
  $\mathcal J$ (Definition~\ref{asm-theta-agree-planes}(5)), with components
  $e_\infty$ and $e_v$, let $S$ and $S_1$ be operators on functions on
  $\mathbf A^{I_M}$ and on $\mathbf A^{I_{M_1}}$, and $c\in\mathbf C$, with
  $e(S\Phi)=c\,S_1(e(\Phi))$ for every Schwartz-Bruhat $\Phi$. Let $\phi$ and
  $\phi'$ be Schwartz-Bruhat functions on $F_v^{I_M}$ such that $S$ acts at $v$
  by $\phi\mapsto\phi'$, and suppose that $S_1$ acts at $v$ by
  $e_v(\phi)\mapsto\chi$. Then $e_v(\phi')=c\,\chi$.
\item Let $M$, $M_1$ and $e$ be as in (2), and let $w\in\mathbb W(X,W)$ be at $v$
  with $e(\rho(w)\Phi)=e(\Phi)$ for every Schwartz-Bruhat $\Phi$. Then
  $e_v(\rho_v(w)\phi)=e_v(\phi)$ for every Schwartz-Bruhat function $\phi$ on
  $F_v^{I_M}$.
\end{enumerate}
\end{lemma}
\leanref{Development/WeilBound/ThetaLifts/LocalTheory/ConeShifts/CoefficientSpans.lean}{PositivityTheory.ThKG.Loc.Cone.Fin.rv}
\leanref{Development/WeilBound/ThetaLifts/LocalTheory/ConeShifts/CoefficientSpans.lean}{PositivityTheory.ThKG.Loc.Cone.Fin.rv_sb}
\leanref{Development/WeilBound/ThetaLifts/LocalTheory/ConeShifts/CoefficientSpans.lean}{PositivityTheory.ThKG.Loc.Cone.Fin.rv_acts}
\leanref{Development/WeilBound/ThetaLifts/LocalTheory/ConeShifts/PureCancel.lean}{PositivityTheory.ThKG.Loc.Cone.Fin.schr_mul}
\leanref{Development/WeilBound/ThetaLifts/LocalTheory/UnramifiedCohomology/SpanSoLoc.lean}{PositivityTheory.ThKG.Loc.FL.ev_loc}
\leanref{Development/WeilBound/ThetaLifts/LocalTheory/ConeShifts/CoefficientSpans.lean}{PositivityTheory.ThKG.Loc.Cone.Fin.ev_inv}
\begin{proof}
(1) Write $(a_1,a_2)=e^M(w)$ and $(a'_1,a'_2)=e^M(w')$; as $e^M$ is
$\mathbf A$-linear, these are at $v$. The function $s\mapsto\psi(\bar s)$ on
$F_v$ is a continuous character, hence locally constant, so $\rho_v(w)\phi$ is
locally constant with compact support when $\phi$ is. For $\xi'\in\mathbf
A^{I_M}$, $\xi'\cdot a_2=\overline{\xi'_v}\cdot a_2$, and $r(\xi'+a_1)=r(\xi')$
when $r$ does not depend on the $v$-component; so
$\rho(w)(\phi\otimes r)=(\rho_v(w)\phi)\otimes r$, and $\rho(w)$ acts at $v$ by
$\phi\mapsto\rho_v(w)\phi$. As $a_1$ and $a'_2$ are at $v$,
$\overline{\xi+a_{1,v}}\cdot a'_2=\bar\xi\cdot a'_2+a_1\cdot a'_2$. Hence
$(\rho_v(w)\rho_v(w')\phi)(\xi)$ is
$\phi(\xi+a_{1,v}+a'_{1,v})$ times $\psi$ of
\[
\bar\xi\cdot(a_2+a'_2)+\tfrac12a_1\cdot a_2+\tfrac12a'_1\cdot a'_2+a_1\cdot a'_2
=\bar\xi\cdot(a_2+a'_2)+\tfrac12(a_1+a'_1)\cdot(a_2+a'_2)+\tfrac12(a_1\cdot a'_2-a'_1\cdot a_2).
\]
As $\langle w,w'\rangle=a_1\cdot a'_2-a'_1\cdot a_2$, this is the first
identity. Exchanging $w$ and $w'$ and using
$\langle w',w\rangle=-\langle w,w'\rangle$ gives the second.

(2) As $e_\infty\neq0$, choose $f$ with $e_\infty(f)\neq0$. For a
Schwartz-Bruhat function $\chi'$ on $F_v^{I_M}$ the function
$f\otimes\chi'\otimes\bigotimes_{v'\neq v}\mathrm{vac}_{v'}$ is
Schwartz-Bruhat, and it is $\chi'\otimes r$ with
$r=f\otimes\bigotimes_{v'\neq v}\mathrm{vac}_{v'}$, which does not depend on the
$v$-component. So $S(\phi\otimes r)=\phi'\otimes r$. The components of $e$ give
$e(\chi'\otimes r)=e_v(\chi')\otimes r_1$ with
$r_1=e_\infty(f)\otimes\bigotimes_{v'\neq v}\mathrm{vac}_{1,v'}$. Hence
\[
e_v(\phi')\otimes r_1=e(S(\phi\otimes r))=c\,S_1(e_v(\phi)\otimes r_1)=c\,\chi\otimes r_1 .
\]
The function $r_1$ is not $0$: $e_\infty(f)$ and every $\mathrm{vac}_{1,v'}$ are
nonzero, and $\mathrm{vac}_{1,v'}(0)=1$ for almost all $v'$. So
$e_v(\phi')=c\,\chi$.

(3) Apply (2) with $S=\rho(w)$, $S_1$ the identity, $c=1$ and
$\phi'=\rho_v(w)\phi$, which is Schwartz-Bruhat by (1).
\end{proof}
% Lean: (1) is rv_sb, rv_acts (VacBridge.rho_acts_locCo), schr_mul and schr_comm, with loc_form
% for the form; (2) is ev_loc, through pureAct and pure_cancel; (3) is ev_inv, the same argument.

\begin{definition}\label{asm-theta-agree-local-factor-setting}
Let $B$ be a theta base of a CM field $E$. A \emph{cone setting}
$(M,M_1,H,H',e,u)$ for $B$ consists of: lattice spaces $X$, $W$, $X_1$ and
$W_1$, of ranks $n$, $d$, $n_1$ and $d_1$, with $X$ and $X_1$ Hermitian and $W$
and $W_1$ skew-Hermitian, of some signatures; regular models $M$ for $(X,W)$
and $M_1$ for $(X_1,W_1)$ (Definition~\ref{asm-theta-agree-models}(3)), with
Schrödinger operators $\rho$ and $\rho_1$ and vacua $\mathrm{vac}_v$ and
$\mathrm{vac}_{1,v}$; adapted planes $H=(x,y,\iota,\mathfrak b)$ in $X$ over $X_1$
and $H'=(x',y',\iota',\mathfrak b')$ in $W$ over $W_1$, at some positions
(Definition~\ref{asm-theta-agree-planes}(1)); and either an evaluation map $e$
for $(M,M_1,\iota,\iota',\mathcal J_\alpha(H,H'))$ and $u=u_\alpha$, or an
evaluation map $e$ for $(M,M_1,\iota,\iota',\mathcal J_\beta(H,H'))$ and
$u=u_\beta$ (Definition~\ref{asm-theta-agree-planes}(5) and~(7)), with
components $e_\infty$ and $e_v$. For a finite place $v$, $\Gamma_v$ is the
group $\Lambda_v\subset\mathbb W(X,W)$ of
Lemma~\ref{asm-theta-agree-evaluation-lattice}, and $\rho_v(w)$, for $w$ at
$v$, is the operator of Lemma~\ref{asm-theta-agree-local-factor-evaluation}.
\end{definition}
\leanref{Development/WeilBound/ThetaLifts/LocalTheory/ConeShifts/ArchCone.lean}{PositivityTheory.ThKG.Loc.Cone.Cfg}

\begin{lemma}\label{asm-theta-agree-local-factor-directions}
Let $(M,M_1,H,H',e,u)$ be a cone setting for a theta base $B$
(Definition~\ref{asm-theta-agree-local-factor-setting}), and let $h$ be the form
on $\mathbb W(X,W)$ of Definition~\ref{asm-theta-agree-models}(3). Put
$(V_0,V_1,\sigma)=(xy'^{\mathsf T},yx'^{\mathsf T},1)$ if $u=u_\alpha$ and
$(V_0,V_1,\sigma)=(yx'^{\mathsf T},xy'^{\mathsf T},c)$ if $u=u_\beta$, and let $v$
be a finite place.
\begin{enumerate}
\item $u(l)=\sigma(l)V_0$ for $l\in\mathbf A_E^\times$;
  $h(V_0,V_0)=h(V_1,V_1)=0$; $e(\rho(tV_1)\Phi)=e(\Phi)$ for $t\in\mathbf A_E$
  and Schwartz-Bruhat $\Phi$; and for every $\gamma\in\Gamma_v$ there is
  $t\in\mathbf A_E$ with $tV_1\in\Gamma_v$ and $h(\gamma,V_0)=t\,h(V_1,V_0)$.
\item Let $\phi$ be a Schwartz-Bruhat function on $F_v^{I_M}$ and $a\in\mathbf A$
  with $\rho_v(atV_1)\phi=\phi$ whenever $tV_1\in\Gamma_v$, and let
  $b\in\mathbf A_E$ with $bV_0$ at $v$. If $e_v(\rho_v(bV_0)\phi)\neq0$, then
  $abV_0\in\Gamma_v$.
\item If $\gamma\in\Gamma_v$ and $h(\gamma,\gamma)=0$, then $\rho(\gamma)$ acts at
  $v$ by $\mathrm{vac}_v\mapsto\mathrm{vac}_v$, and
  $\rho_v(\gamma)\mathrm{vac}_v=\mathrm{vac}_v$.
\end{enumerate}
\end{lemma}
\leanref{Development/WeilBound/ThetaLifts/LocalTheory/ConeShifts/Plane.lean}{PositivityTheory.ThKG.Loc.Cone.Fin.dirs}
\leanref{Development/WeilBound/ThetaLifts/LocalTheory/ConeShifts/Plane.lean}{PositivityTheory.ThKG.Loc.Cone.Fin.comp_lat}
\leanref{Development/WeilBound/ThetaLifts/LocalTheory/ConeShifts/CoefficientSpans.lean}{PositivityTheory.ThKG.Loc.Cone.Fin.key}
\leanref{Development/WeilBound/ThetaLifts/LocalTheory/ConeShifts/CoefficientSpans.lean}{PositivityTheory.ThKG.Loc.Cone.Fin.vac_fix}
\begin{proof}
For $p,p'\in E^n$ and $q,q'\in E^d$, the definition of $h$ gives
$h(pq^{\mathsf T},p'q'^{\mathsf T})=h_X(p,p')\,h_W(q,q')$, with
$h_X(p,p')=p^{\mathsf T}J_Xc(p')$ and $h_W$ likewise; and $h(sw,tw')=s\,c(t)\,h(w,w')$
for $s,t\in\mathbf A_E$.

(1) The formula for $u$ is the definition of $u_\alpha$ and $u_\beta$. As $x$,
$y$, $x'$ and $y'$ are isotropic, $h(V_0,V_0)=h(V_1,V_1)=0$. If $u=u_\alpha$,
then $tV_1=(ty)x'^{\mathsf T}\in\mathcal J_0(H')\subset\mathcal J_\alpha(H,H')$; if
$u=u_\beta$, then $tV_1=x(ty')^{\mathsf T}\in\mathcal J_\beta(H,H')$. In both cases
$e(\rho(tV_1)\Phi)=e(\Phi)$. For the last claim, let $u=u_\alpha$; the case
$u=u_\beta$ is the same with $x$ and $y$ exchanged and $x'$ and $y'$ exchanged.
Every $p\in L_X$ is $sx+\iota z_1+ty$ with $ty\in L_X$, so $h_X(p,x)=t\,h_X(y,x)$,
as $h_X(x,x)=h_X(\iota z_1,x)=0$. Likewise every $q\in L_W$ is
$s'x'+\iota'z'_1+t'y'$ with $s'x'\in L_W$, and $h_W(q,y')=s'h_W(x',y')$. So for a
generator $\gamma=a\,pq^{\mathsf T}$ of $\Gamma_v$,
$h(\gamma,V_0)=a\,t\,s'\,h(V_1,V_0)$, and $ats'V_1=a\,(ty)(s'x')^{\mathsf T}\in\Gamma_v$.
The set of $\gamma$ with the property is a group, as $h$ is additive in its
first argument, so it is $\Gamma_v$.

(2) By Lemma~\ref{asm-theta-agree-evaluation-lattice}(2) it suffices to show
$\psi(\langle\gamma,abV_0\rangle)=1$ for $\gamma\in\Gamma_v$, since $abV_0$ is at
$v$. Let $\gamma\in\Gamma_v$, and let $t$ be as in (1). Then $atV_1$ is at $v$, as
$\Gamma_v$ consists of elements at $v$. Put $\chi=\rho_v(bV_0)\phi$, a
Schwartz-Bruhat function. By (1) and
Lemma~\ref{asm-theta-agree-local-factor-evaluation}(3),
$e_v(\rho_v(atV_1)\chi)=e_v(\chi)$. By
Lemma~\ref{asm-theta-agree-local-factor-evaluation}(1) and the hypothesis on
$\phi$,
\[
\rho_v(atV_1)\chi=\psi(\langle atV_1,bV_0\rangle)\,\rho_v(bV_0)\rho_v(atV_1)\phi=\psi(\langle atV_1,bV_0\rangle)\,\chi .
\]
As $e_v(\chi)\neq0$, $\psi(\langle atV_1,bV_0\rangle)=1$. Finally
$h(atV_1,bV_0)=a\,c(b)\,t\,h(V_1,V_0)=a\,c(b)\,h(\gamma,V_0)=h(\gamma,abV_0)$, as
$a\in\mathbf A$, so $\langle atV_1,bV_0\rangle=\langle\gamma,abV_0\rangle$.

(3) Let $\theta$ be the element in the definition of the vacuum
$\mathrm{vac}_v$. As $\Gamma_v$ is the group generated by the matrices of that
definition, $\rho(\gamma)$ acts at $v$ by
$\mathrm{vac}_v\mapsto\psi(\tfrac12\operatorname{Tr}h(\theta\gamma,\gamma))\mathrm{vac}_v=\mathrm{vac}_v$,
as $h(\theta\gamma,\gamma)=\theta\,h(\gamma,\gamma)=0$. It also acts at $v$ by
$\mathrm{vac}_v\mapsto\rho_v(\gamma)\mathrm{vac}_v$, so the two functions agree by
Lemma~\ref{asm-theta-agree-finite-inputs-schur}(1).
\end{proof}
% Lean: dirs (Plane.lean), with hh_rk, hh_smul, comp_lat, exists_smul_y_mem_and_form_eq_of_isAdapted
% and plane_y; key (CoefficientSpans.lean), with Vac.mem_latAtV, ev_inv, schr_comm and
% DualLoc.psi_form; vac_fix, with acts_uniq for the last claim.

\begin{lemma}\label{asm-theta-agree-local-ideal-lines}
Let $B$ be a theta base of a CM field $E$, with set $P$, and let $Y$ be a lattice
space of rank $m$ with lattice $L$, Hermitian or skew-Hermitian of some
signature. Put $D=\mathcal O_E$ if $Y$ is Hermitian, and let $D$ be the
codifferent $\{t\in E:\operatorname{Tr}_{E/\mathbf Q}(ot)\in\mathbf Z\text{ for all }o\in\mathcal O_E\}$
if $Y$ is skew-Hermitian. Call a pair $(u,\mathfrak p)$ of a vector $u\in E^m$ and a
nonzero fractional ideal $\mathfrak p$ a \emph{line} of $L$ if $su\in L$ for
$s\in\mathfrak p$ and there is $p\in E^m$ with $u\cdot p=1$ and $z\cdot p\in\mathfrak p$
for every $z\in L$. Let $H=(x,y,\iota,\mathfrak b)$ be an adapted plane in $Y$
over some lattice space, at positions $(i,j)$
(Definition~\ref{asm-theta-agree-planes}(1)).
\begin{enumerate}
\item $(x,\mathfrak b)$ and $(y,D\,c(\mathfrak b)^{-1})$ are lines of $L$.
\item $\operatorname{ord}_wD=0$ for every place $w$ of $E$ above $P$; and if
  $i,j<m$, then $\operatorname{ord}_w\mathfrak b=0$ for every place $w$ of $E$ above
  $P$.
\end{enumerate}
\end{lemma}
\leanref{Development/WeilBound/ThetaLifts/LocalTheory/ConeShifts/IdealB.lean}{PositivityTheory.ThKG.Loc.Cone.Id.Line}
\leanref{Development/WeilBound/ThetaLifts/LocalTheory/ConeShifts/IdealB.lean}{PositivityTheory.ThKG.Loc.Cone.Id.plane_lines}
\leanref{Development/WeilBound/ThetaLifts/LocalTheory/ConeShifts/IdealB.lean}{PositivityTheory.ThKG.Loc.Cone.Id.plane_pr}
\leanref{Development/WeilBound/ThetaLifts/LocalTheory/ConeShifts/IdealB.lean}{PositivityTheory.ThKG.Loc.Cone.Id.pr_codiff}
\begin{proof}
Let $J$ be the matrix of $Y$, $h(a,b)=a^{\mathsf T}J\,c(b)$, and $\varepsilon=1$ or
$-1$ in the two cases, so that $h(b,a)=\varepsilon\,c(h(a,b))$ and
$L=\{a:h(a,z)\in D\text{ for all }z\in L\}$
(Definition~\ref{asm-theta-agree-doubling}(1)); multiplying by $\varepsilon$
preserves $D$.

(1) Every $z\in L$ is $sx+\iota u+ty$ with $sx\in L$ and $ty\in L$, so $s\in\mathfrak b$,
and $h(z,y)=s$ and $h(z,x)=t\varepsilon$, as $h(x,x)=h(y,y)=h(\iota u,x)=h(\iota u,y)=0$
and $h(x,y)=1$. With $p=J\,c(y)$, $x\cdot p=h(x,y)=1$ and $z\cdot p=h(z,y)\in\mathfrak b$;
and $sx\in L$ for $s\in\mathfrak b$ by the definition of an adapted plane. So
$(x,\mathfrak b)$ is a line, and $\{h(z,y):z\in L\}=\mathfrak b$. Now
$h(ty,z)=t\varepsilon\,c(h(z,y))$, so $ty\in L$ if and only if
$t\,c(\mathfrak b)\subset D$, that is $t\in D\,c(\mathfrak b)^{-1}$. With
$p=\varepsilon J\,c(x)$, $y\cdot p=\varepsilon h(y,x)=1$ and
$z\cdot p=\varepsilon h(z,x)=t$ for $z=sx+\iota u+ty$, which lies in
$D\,c(\mathfrak b)^{-1}$ as $ty\in L$.

(2) As $\psi$ is trivial on $F$ and equals $\exp(-2\pi i\operatorname{Tr}(x))$ on
$F_\infty$, it is the standard character $\psi_{\mathbf Q}\circ\operatorname{Tr}_{F/\mathbf Q}$,
whose component at $v$ has conductor the codifferent of $F$ at $v$
\cite{Tat67}. As $\psi_v$ has conductor $\mathcal O_v$ for $v\in P$, the
codifferent of $F$ has order $0$ at $v\in P$; as $E/F$ is unramified, the
codifferent of $E$ is that of $F$ times $\mathcal O_E$, and has order $0$ above
$P$. So $\operatorname{ord}_wD=0$ for $w$ above $P$ in both cases.

Let $v\in P$, let $w$ be the place of $E$ given by $\tau_v$, and identify $E_w$ with
$F_v$ through $\tau_v$. We show: if $(u,\mathfrak p)$ is a line of $L$ and $\tau_v(u)$
is a column of $b_{Y,v}$, then $\operatorname{ord}_w\mathfrak p=0$. The columns of
$b_{Y,v}$ are a basis of the completion $L_w$ of $L$. Let $p$ be as in the
definition of a line. For $z\in L$, $\operatorname{ord}_w(z\cdot p)\geq\operatorname{ord}_w\mathfrak p$,
so the same holds for $z\in L_w$; for $z=\tau_v(u)$ this gives
$0\geq\operatorname{ord}_w\mathfrak p$. Choose $s\in\mathfrak p$ with
$\operatorname{ord}_ws=\operatorname{ord}_w\mathfrak p$. Then $su\in L$, so the
coordinate $\tau_v(s)$ of $\tau_v(su)$ in the basis lies in $\mathcal O_v$, and
$\operatorname{ord}_w\mathfrak p\geq0$.

If $i,j<m$, the columns $i$ and $j$ of $b_{Y,v}$ are $\tau_v(x)$ and $\tau_v(y)$. By
(1), $\operatorname{ord}_w\mathfrak b=0$ and
$\operatorname{ord}_w(D\,c(\mathfrak b)^{-1})=0$, so
$\operatorname{ord}_{c(w)}\mathfrak b=\operatorname{ord}_wc(\mathfrak b)=\operatorname{ord}_wD=0$.
As $v$ splits in $E$, $w$ and $c(w)$ are the places above $v$.
\end{proof}
% Lean: Line (IdealB.lean) with fields ne, mem and ex; plane_lines for (1) (with D = 1 and
% epsilon = 1 for Hermitian spaces, D = codiff E and epsilon = -1 for skew-Hermitian ones, as in
% lines of IdealId0.lean); (2) is pr_codiff (count_codiff_F with PsiUnr, codiff_up with
% hB.unram) and plane_pr (Line.count_eq for column i of b_{Y,v} at tau_v, and for column j with
% count_bar at tau_v o c).

\begin{lemma}\label{asm-theta-agree-local-ideal-lattice}
Let $B$ be a theta base of a CM field $E$, let $X$ be a Hermitian and $W$ a
skew-Hermitian lattice space, with lattices $L_X$ and $L_W$, let $(x,\mathfrak p)$
be a line of $L_X$ and $(y,\mathfrak q)$ one of $L_W$
(Lemma~\ref{asm-theta-agree-local-ideal-lines}), and let $v$ be a finite place
and $\Lambda_v$ as in Lemma~\ref{asm-theta-agree-evaluation-lattice}. Then for
$c\in\mathbf A_E$, $1_v\,c\,xy^{\mathsf T}\in\Lambda_v$ if and only if
$c_w\in\mathfrak p\mathfrak q\,\mathcal O_w$ for every place $w$ of $E$ above $v$.
\end{lemma}
\leanref{Development/WeilBound/ThetaLifts/LocalTheory/ConeShifts/IdealB.lean}{PositivityTheory.ThKG.Loc.Cone.Id.line_lat}
\leanref{Development/WeilBound/ThetaLifts/LocalTheory/ConeShifts/IdealA.lean}{PositivityTheory.ThKG.Loc.Cone.Id.lat_iff}
\begin{proof}
For a nonzero fractional ideal $\mathfrak c$ let $\Lambda_v(\mathfrak c)\subset\mathbf A_E$
be the group generated by the $as$ with $a\in\mathbf A$ at $v$, $a_v\in\mathcal O_v$
and $s\in\mathfrak c$. It consists of the elements at $v$ whose components at the
places $w\mid v$ lie in $\mathfrak c\mathcal O_w$, as $\mathfrak c\otimes\mathcal O_v$
maps onto $\prod_{w\mid v}\mathfrak c\mathcal O_w$. So $1_vc\in\Lambda_v(\mathfrak p\mathfrak q)$
if and only if $c_w\in\mathfrak p\mathfrak q\,\mathcal O_w$ for every $w\mid v$.

If $1_vc\in\Lambda_v(\mathfrak p\mathfrak q)$, then $1_vc$ is a sum of terms
$\pm a\,st$ with $s\in\mathfrak p$ and $t\in\mathfrak q$, and
$a\,st\,xy^{\mathsf T}=a\,(sx)(ty)^{\mathsf T}\in\Lambda_v$. Conversely let $p$ and
$q$ be as in the definition of the two lines. The map $Z\mapsto p^{\mathsf T}Zq$
sends a generator $a\,x'y'^{\mathsf T}$ of $\Lambda_v$ to
$a\,(x'\cdot p)(y'\cdot q)\in\Lambda_v(\mathfrak p\mathfrak q)$, so it maps
$\Lambda_v$ into $\Lambda_v(\mathfrak p\mathfrak q)$; and it sends
$1_vc\,xy^{\mathsf T}$ to $1_vc$.
\end{proof}
% Lean: line_lat (IdealB.lean), with the map key (Z |-> p . Z q) for one direction and
% FractionalIdeal.mul_induction_on for the other; Lat a v is the group Lambda_v(c), and lat_iff
% (IdealA.lean, through lat_of_inAt and inAt_of_lat) is the first paragraph.

\begin{lemma}\label{asm-theta-agree-local-ideal-values}
Let $B$ be a theta base of a CM field $E$, with set $P$, and let $\nu$ be a
continuous character of $\mathbf A_E^\times$ with values of absolute value $1$,
trivial on $\mathbf A^\times$ and unramified at the places of $E$ not above $P$.
Let $\nu(\mathfrak a)$ be as in Definition~\ref{asm-theta-agree-parts}(4).
\begin{enumerate}
\item $\nu(c(l))=\nu(l)^{-1}$ for $l\in\mathbf A_E^\times$; and
  $\nu(\mathfrak a\mathfrak a')=\nu(\mathfrak a)\nu(\mathfrak a')$ and
  $\nu(\mathfrak a^{-1})=\nu(\mathfrak a)^{-1}$.
\item If $\operatorname{ord}_w\mathfrak a=0$ for every $w$ above $P$, then
  $\nu(c(\mathfrak a))=\nu(\mathfrak a)^{-1}$.
\item $\nu(D)=1$ for the codifferent $D$ of $E$.
\end{enumerate}
\end{lemma}
\leanref{Development/WeilBound/ThetaLifts/LocalTheory/ConeShifts/IdealB.lean}{PositivityTheory.ThKG.Loc.Cone.Id.nu_conj}
\leanref{Development/WeilBound/ThetaLifts/LocalTheory/ConeShifts/IdealA.lean}{PositivityTheory.ThKG.Loc.Cone.Id.idealVal_cj}
\leanref{Development/WeilBound/ThetaLifts/LocalTheory/ConeShifts/IdealB.lean}{PositivityTheory.ThKG.Loc.Cone.Id.idealVal_codiff}
\begin{proof}
(1) The idele $l\,c(l)$ is fixed by $c$, so it lies in $\mathbf A^\times$, as the
elements of $\mathbf A_E=\mathbf A\otimes_FE$ fixed by $c$ are those of
$\mathbf A$. So $\nu(l)\nu(c(l))=1$. The other two identities hold as
$\operatorname{ord}_w$ is additive.

(2) Let $l$ be the idele with component $\varpi_w^{\operatorname{ord}_w\mathfrak a}$ at
the finitely many $w$ with $\operatorname{ord}_w\mathfrak a\neq0$ and $1$ elsewhere.
These $w$ are not above $P$, so $\nu$ is unramified there, and
$\nu(l)=\nu(\mathfrak a)$. The idele $c(l)$ has component a power of a uniformizer
of exponent $\operatorname{ord}_w\mathfrak a=\operatorname{ord}_{c(w)}c(\mathfrak a)$
at each $c(w)$ and $1$ elsewhere; so likewise $\nu(c(l))=\nu(c(\mathfrak a))$. By
(1), $\nu(c(\mathfrak a))=\nu(l)^{-1}=\nu(\mathfrak a)^{-1}$.

(3) As in the proof of Lemma~\ref{asm-theta-agree-local-ideal-lines}(2), $D$ is the
codifferent $D_F$ of $F$ times $\mathcal O_E$, and $D_F$ has order $0$ at $P$. Let
$a\in\mathbf A^\times$ have component $\varpi_v^{\operatorname{ord}_vD_F}$ at the
finitely many $v$ where this order is not $0$, and $1$ elsewhere. As $E/F$ is
unramified, $a$ has order $\operatorname{ord}_wD$ at every place $w$ of $E$; and
$\nu$ is unramified at the places above these $v$. So $\nu(D)=\nu(a)=1$.
\end{proof}
% Lean: nu_conj (IdealB.lean, with fixed_cA); idealVal_mul and idealVal_inv; idealVal_cj (IdealA.lean)
% through exists_gen, Gen.conj and Gen.eq (the idele l of (2)); idealVal_codiff with codiff_up and
% count_codiff_F.

\begin{lemma}\label{asm-theta-agree-local-factor-ideal}
Let $B$ be a theta base of a CM field $E$, let $X$ be a Hermitian and $W$ a
skew-Hermitian lattice space, of ranks $n$ and $d$ and of some signatures, and
let $H=(x,y,\iota,\mathfrak b)$ and $H'=(x',y',\iota',\mathfrak b')$ be adapted
planes in $X$ and in $W$ over some lattice spaces, at positions $(i,j)$ and
$(i',j')$ (Definition~\ref{asm-theta-agree-planes}(1)). Let $u$ be $u_\alpha$ or
$u_\beta$, and let $\Gamma_v\subset\mathbb W(X,W)$ be the group $\Lambda_v$ of
Lemma~\ref{asm-theta-agree-evaluation-lattice}. For a fractional ideal
$\mathfrak a$ of $E$, a finite place $v$ and $l\in\mathbf A_E^\times$, say that $l$
\emph{lies in $\mathfrak a$ at $v$} if $l_w\in\mathfrak a\mathcal O_w$ for every
place $w$ of $E$ above $v$.
\begin{enumerate}
\item There is a fractional ideal $\mathfrak a\neq0$ of $E$ such that for every
  finite place $v$ and every $l\in\mathbf A_E^\times$, $1_vu(l)\in\Gamma_v$ if and
  only if $l$ lies in $\mathfrak a$ at $v$.
\item There are signs $a_1,a_2,b_1,b_2\in\{\pm1\}$, independent of all the data,
  such that the following holds if $i,j<n$ and $i',j'<d$. Let
  $(e_1,e_2)=(a_1,a_2)$ if $u=u_\alpha$ and $(e_1,e_2)=(b_1,b_2)$ if $u=u_\beta$.
  Then $\mathfrak a$ in (1) can be chosen with $\operatorname{ord}_w\mathfrak a=0$
  for every place $w$ of $E$ above $P$ and
  $\nu(\mathfrak a)=\nu(\mathfrak b)^{e_1}\nu(\mathfrak b')^{e_2}$
  (Definition~\ref{asm-theta-agree-parts}(4)) for every continuous character
  $\nu$ of $\mathbf A_E^\times$ that is trivial on $\mathbf A^\times$ and
  unramified at the places of $E$ not above $P$.
\end{enumerate}
\end{lemma}
\leanref{Development/WeilBound/ThetaLifts/LocalTheory/ConeShifts/IdealStmt.lean}{PositivityTheory.ThKG.Loc.Cone.IdealId0}
\leanref{Development/WeilBound/ThetaLifts/LocalTheory/ConeShifts/IdealId0.lean}{PositivityTheory.ThKG.Loc.Cone.Id.idealId0}
\leanref{Development/WeilBound/ThetaLifts/LocalTheory/ConeShifts/IdealStmt.lean}{PositivityTheory.ThKG.Loc.Cone.IdealId'}
\leanref{Development/WeilBound/ThetaLifts/LocalTheory/ConeShifts/IdealId0.lean}{PositivityTheory.ThKG.Loc.Cone.Id.idealId'}
\begin{proof}
Let $D$ be the codifferent of $E$. By
Lemma~\ref{asm-theta-agree-local-ideal-lines}(1), $(x,\mathfrak b)$ and
$(y,c(\mathfrak b)^{-1})$ are lines of $L_X$, and $(x',\mathfrak b')$ and
$(y',D\,c(\mathfrak b')^{-1})$ are lines of $L_W$.

(1) If $u=u_\alpha$, put $\mathfrak a=\mathfrak b\,D\,c(\mathfrak b')^{-1}$. As
$1_vu_\alpha(l)=1_vl\,xy'^{\mathsf T}$,
Lemma~\ref{asm-theta-agree-local-ideal-lattice} gives the claim. If $u=u_\beta$, put
$\mathfrak a=c\bigl(c(\mathfrak b)^{-1}\mathfrak b'\bigr)$. As
$1_vu_\beta(l)=1_vc(l)\,yx'^{\mathsf T}$, the same lemma shows that
$1_vu_\beta(l)\in\Gamma_v$ if and only if $c(l)$ lies in $c(\mathfrak b)^{-1}\mathfrak b'$
at $v$; as $c$ carries the component at $w$ to the component at $c(w)$ and
$c(\mathfrak c\mathcal O_w)=c(\mathfrak c)\mathcal O_{c(w)}$, this holds if and only
if $l$ lies in $\mathfrak a$ at $v$.

(2) Take $(a_1,a_2,b_1,b_2)=(1,1,-1,-1)$, and let $i,j<n$ and $i',j'<d$. By
Lemma~\ref{asm-theta-agree-local-ideal-lines}(2), $\mathfrak b$, $\mathfrak b'$ and
$D$ have order $0$ at every place above $P$, and so do their conjugates, as $c$
permutes the places above each $v$; hence so does $\mathfrak a$. Let $\nu$ be as in
(2). By Lemma~\ref{asm-theta-agree-local-ideal-values}, if $u=u_\alpha$,
$\nu(\mathfrak a)=\nu(\mathfrak b)\,\nu(D)\,\nu(c(\mathfrak b'))^{-1}=\nu(\mathfrak b)\nu(\mathfrak b')$;
and if $u=u_\beta$,
$\nu(\mathfrak a)=\nu(c(\mathfrak b)^{-1}\mathfrak b')^{-1}=\nu(c(\mathfrak b))\,\nu(\mathfrak b')^{-1}=\nu(\mathfrak b)^{-1}\nu(\mathfrak b')^{-1}$.
\end{proof}
% Lean: Id.idealId0 and Id.idealId' (ConeShifts/IdealId0.lean): lines (plane_lines for X and W),
% memA (line_lat and lat_iff) and memB (with inAt_cj_iff) for (1); for (2), plane_pr and
% pr_codiff for the orders above P, and idealVal_mul, idealVal_inv, idealVal_cj (with nu_conj)
% and idealVal_codiff for the values. The Lean takes (a1, a2, b1, b2) = (1, 1, -1, -1).

\begin{lemma}\label{asm-theta-agree-local-ideles-parts}
Let $E$ be a CM field with maximal totally real subfield $F$, and let
$U\subset\mathbf A_E^\times$ be the group of ideles whose components at all finite
places of $E$ are units. For an idempotent $e\in\mathbf A_E$ and
$l\in\mathbf A_E^\times$ put $l^{(e)}=el+(1-e)$, and let $A_e$ be the set of $l$
with $l^{(e)}=l$.
\begin{enumerate}
\item $l\mapsto l^{(e)}$ is a continuous homomorphism of $\mathbf A_E^\times$ onto
  the closed subgroup $A_e$; $(l^{(f)})^{(e)}=l^{(ef)}$; and if $ef=0$, then
  $l^{(e+f)}=l^{(e)}l^{(f)}$ and $(x,y)\mapsto xy$ is an isomorphism of
  topological groups $A_e\times A_f\to A_{e+f}$, with inverse
  $l\mapsto(l^{(e)},l^{(f)})$.
\item If $e$ is $0$ at the infinite places, $A_e\cap U$ is a compact open subgroup
  of $A_e$; let $\beta_e$ be the Haar measure on $A_e$ with $\beta_e(A_e\cap U)=1$.
  If $e$ and $f$ are two such idempotents with $ef=0$, then $\beta_{e+f}$ is the
  image of $\beta_e\times\beta_f$ under the isomorphism of (1).
\item For a finite place $w$ of $E$ let $e_w$ be the idempotent that is $1$ at $w$
  and $0$ elsewhere, so that $A_{e_w}=E_w^\times$; for a finite place $v$ of $F$,
  $1_v=\sum_{w\mid v}e_w$, $A_{1_v}=A_v$ (Definition~\ref{asm-theta-agree-local-cones}(2)),
  and $\beta_{1_v}=m_{A,v}$ for every normalized local measure system $m$.
\end{enumerate}
\end{lemma}
\leanref{Development/WeilBound/ThetaLifts/LocalTheory/ConeShifts/ContinuousFc.lean}{PositivityTheory.ThKG.Loc.Cone.sysI}
\leanref{Development/WeilBound/ThetaLifts/LocalTheory/ConeShifts/FubLoc.lean}{PositivityTheory.ThKG.Loc.Cone.PI_sum}
\leanref{Development/WeilBound/ThetaLifts/LocalTheory/ConeShifts/ContinuousFc.lean}{PositivityTheory.ThKG.Loc.Cone.units_eq}
\begin{proof}
(1) As $e(1-e)=0$, $(ex+(1-e))(ey+(1-e))=exy+(1-e)$, and $l^{(e)}$ has inverse
$(l^{-1})^{(e)}$. The other identities are expansions of the same kind; $A_e$ is
closed as the fixed set of a continuous map, and $(l^{(e)})^{(e)}=l^{(e)}$. If
$ef=0$, $x\in A_e$ and $y\in A_f$, then $(xy)^{(e)}=x\,y^{(ef)}=x$ and
$(xy)^{(f)}=y$, and $l=l^{(e+f)}=l^{(e)}l^{(f)}$ for $l\in A_{e+f}$.

(2) The components of $e$ are $0$ or $1$. If $e$ is $0$ at the infinite places,
$A_e\cap U$ is the product of the unit groups $\mathcal O_w^\times$ over the finite
$w$ where $e$ is $1$, embedded with components $1$ elsewhere; it is compact, and
open in $A_e$ as $U$ is open in $\mathbf A_E^\times$. If $ef=0$, then $l\in A_{e+f}$
lies in $U$ if and only if $l^{(e)}$ and $l^{(f)}$ do, as their components at $w$
are $l_w$ or $1$. By the uniqueness of Haar measures, the image of
$\beta_e\times\beta_f$ is $c\,\beta_{e+f}$ for some $c>0$, and evaluating on
$A_{e+f}\cap U$ gives $c=1$.

(3) The first claims are clear. The proof of
Lemma~\ref{asm-theta-agree-local-measures} shows that $A_v\cap U$ is the set of
$l\in A_v$ whose powers have compact closure, so the normalization of $m_{A,v}$
is $m_{A,v}(A_v\cap U)=1$.
\end{proof}
% Lean: sysI (ContinuousFc.lean) is the system of maps PI e on the ideles; Sys.split and
% Split.fub give (1) and the comparison in (2), with the constant 1 from the masses as in
% Sys.prod; co_I and cpt_I give the compactness; idAt_eq and units_eq give (3); PI_sum is the
% product formula for several orthogonal idempotents.

\begin{lemma}\label{asm-theta-agree-local-ideles-tate}
Let $E$ be a CM field, $w$ a finite place of $E$ with uniformizer $\varpi_w$ and
norm $N(w)$, and $\gamma$ the Haar measure on $E_w^\times$ with
$\gamma(\mathcal O_w^\times)=1$. Let $\nu$ be a character of $\mathbf A_E^\times$ with
values of absolute value $1$, trivial on $\mathcal O_w^\times$, let $z\geq1$ and
$m$ be integers, and put $t_w=\nu(\varpi_w)N(w)^{-z}$, where $\varpi_w$ is read as
an idele that is $1$ away from $w$. Then $|t_w|<1$, the function
$l\mapsto\mathbf 1[\operatorname{ord}_w(l)\geq m]\,\nu(l)|l|^z$ on $E_w^\times$ is
$\gamma$-integrable, and its integral is $t_w^m(1-t_w)^{-1}\neq0$.
\end{lemma}
\leanref{Development/WeilBound/ThetaLifts/LocalTheory/ConeShifts/FubTate.lean}{PositivityTheory.ThKG.Fub.tate}
\leanref{Development/WeilBound/ThetaLifts/LocalTheory/ConeShifts/FubLoc.lean}{PositivityTheory.ThKG.Loc.Cone.locW}
\begin{proof}
For $l\in E_w^\times$, $|l|=N(w)^{-\operatorname{ord}_w(l)}$, so $|t_w|=N(w)^{-z}<1$.
The group $E_w^\times$ is the disjoint union of the cosets
$\varpi_w^k\mathcal O_w^\times$, $k\in\mathbf Z$, each of $\gamma$-measure $1$, and on
$\varpi_w^k\mathcal O_w^\times$ the function $\nu(l)|l|^z$ is $t_w^k$. So the integral
is $\sum_{k\geq m}t_w^k=t_w^m(1-t_w)^{-1}$, and the series of absolute values
converges.
\end{proof}
% Lean: Fub.tate (FubTate.lean), for a commutative group with a homomorphism o to Z, open
% kernel of measure 1 and a character trivial on the kernel; locW (FubLoc.lean) applies it with
% o = ord_w, pw for the uniformizer (vw_pw, abs_pw, wt_pw), nu_ew and abs_UI.

\begin{lemma}\label{asm-theta-agree-local-ideles-tail}
Let $E$ be a CM field with maximal totally real subfield $F$, let $\mathfrak a\neq0$
be a fractional ideal of $E$, $z\geq2$ an integer, $T$ a finite set of finite
places of $F$, and $\nu$ a continuous character of $\mathbf A_E^\times$ with values
of absolute value $1$, unramified at the places of $E$ not above $T$. Put
$m_w=\operatorname{ord}_w\mathfrak a$ and $t_w$ as in
Lemma~\ref{asm-theta-agree-local-ideles-tate}, let
$r_T=1-1_\infty-\sum_{v\in T}1_v$, and let $\beta_{r_T}$ be as in
Lemma~\ref{asm-theta-agree-local-ideles-parts}(2). Let $T_0\supset T$ be a finite
set of finite places containing the places below every $w$ with $m_w\neq0$. Then
the function $\mathcal T(l)$, equal to $\nu(l)|l|^z$ if $l$ lies in $\mathfrak a$
at every finite $v\notin T$ (Lemma~\ref{asm-theta-agree-local-factor-ideal}) and
to $0$ otherwise, is $\beta_{r_T}$-integrable on $A_{r_T}$, and
\[
\int_{A_{r_T}}\mathcal T\,d\beta_{r_T}=\lim_{T'}\prod_{w}t_w^{m_w}(1-t_w)^{-1},
\]
the product over the places $w$ of $E$ above $(T'\cup T_0)\setminus T$ and the
limit over the finite sets $T'$ of finite places, ordered by inclusion.
\end{lemma}
\leanref{Development/WeilBound/ThetaLifts/LocalTheory/ConeShifts/FubMain.lean}{PositivityTheory.ThKG.Loc.Cone.tail}
\leanref{Development/WeilBound/ThetaLifts/LocalTheory/ConeShifts/FubLoc.lean}{PositivityTheory.ThKG.Loc.Cone.nu_rest}
\leanref{Development/WeilBound/ThetaLifts/LocalTheory/ConeShifts/FubLoc.lean}{PositivityTheory.ThKG.Loc.Cone.prodW}
\begin{proof}
\emph{Step 1: $\nu=1$ on $A_{r_T}\cap U$.} For finite $A\supset T$ the sets
$V_A=A_{r_A}\cap U$ are compact and decrease to $\{1\}$, so some $V_{A_0}$ lies in
the neighbourhood $\{l:|\nu(l)-1|<1\}$ of $1$. Let $l\in A_{r_T}\cap U$, and
$A=T\cup A_0$. By Lemma~\ref{asm-theta-agree-local-ideles-parts}(1), $l$ is the
product of its components $l^{(e_w)}\in\mathcal O_w^\times$ for $w$ above
$A\setminus T$ and of $l^{(r_A)}\in V_A$. As $\nu$ is unramified at these $w$,
$\nu(l)=\nu(l^{(r_A)})$, and all powers of $\nu(l^{(r_A)})$ lie in
$\{s:|s-1|<1\}$. A number of absolute value $1$ with this property is $1$.

\emph{Step 2: the integral over $D_{T'}$.} Let $T'\supset T_0$ be finite, and let
$D_{T'}$ be the set of $l\in A_{r_T}$ whose components at the places not above $T'$
are units. Let $W$ be the set of places of $E$ above $T'\setminus T$. By
Lemma~\ref{asm-theta-agree-local-ideles-parts}, $A_{r_T}\cong\prod_{w\in W}E_w^\times\times A_{r_{T'}}$,
with $\beta_{r_T}$ corresponding to $\prod_w\gamma_w\times\beta_{r_{T'}}$, where
$\gamma_w=\beta_{e_w}$; and $D_{T'}$ corresponds to
$\prod_wE_w^\times\times(A_{r_{T'}}\cap U)$. On $A_{r_{T'}}\cap U$, $\nu=1$ by
Step 1 and $|\cdot|=1$; and for $v\notin T'$, the condition that $l$ lie in
$\mathfrak a$ at $v$ holds on $D_{T'}$, as $m_w=0$ above $v$. So on $D_{T'}$,
$\mathcal T(l)=\prod_{w\in W}\mathbf 1[\operatorname{ord}_w(l_w)\geq m_w]\nu(l^{(e_w)})|l^{(e_w)}|^z$,
and by Fubini's theorem and Lemma~\ref{asm-theta-agree-local-ideles-tate},
$\mathcal T$ is integrable on $D_{T'}$ with integral
$\prod_{w\in W}t_w^{m_w}(1-t_w)^{-1}$.

\emph{Step 3: the limit.} Every idele has unit components at almost all places,
so the $D_{T'}$ increase to $A_{r_T}$. The function $|\mathcal T|$ is $\mathcal T$
for the trivial character, so by Step 2 its integral over $D_{T'}$ is
$\prod_{w\in W}N(w)^{-zm_w}(1-N(w)^{-z})^{-1}$, which is bounded independently of
$T'$ as $\prod_w(1-N(w)^{-z})^{-1}=\zeta_E(z)<\infty$ for $z\geq2$. By monotone
convergence $\mathcal T$ is integrable, and by dominated convergence its integral
is the limit of the integrals over $D_{T'}$.
\end{proof}
% Lean: tail (FubMain.lean): the cover D (AECover by the sets where PI (io (rest (T' u T0))) x is
% a unit, with fc for the cofiniteness), key (prodW with rest_split, ew_rest and nu_rest), the
% bound through lim_tw for the trivial character, and integral_tendsto_of_countably_generated;
% nu_rest is Step 1, through exists_subset_nhds_of_isCompact' and circle_one.

\begin{lemma}\label{asm-theta-agree-local-factor-ideles}
Let $B$ be a theta base of a CM field $E$, let $m$ be a normalized local measure
system for a subgroup $K$ of $U_X$, for a lattice space $X$
(Definition~\ref{asm-theta-agree-local-cones}(2)), $\mu$ a Haar measure on
$\mathbf A_E^\times$, $\mathfrak a\neq0$ a fractional ideal of $E$ and $z\geq2$
an integer. For $l\in\mathbf A_E^\times$ and $v$ a finite place or $v=\infty$
put $l^{(v)}=1_vl+(1-1_v)$, and for a character $\nu$ of $\mathbf A_E^\times$
put $\nu_z(l)=\nu(l)|l|^z$. Then $l^{(v)}\in A_v$, and there is $\kappa\neq0$
such that the following holds for every continuous character $\nu$ of
$\mathbf A_E^\times$ and all finite sets $S_0\subset S$ of finite places such
that $\nu$ is unramified at the places of $E$ not above $S$, and every place $w$
of $E$ above $S$ has $\nu$ ramified at $w$ and $\operatorname{ord}_w\mathfrak a=0$.
\begin{enumerate}
\item For every finite place $v\notin S$, the function equal to $\nu_z(l)$ if
  $l$ lies in $\mathfrak a$ at $v$
  (Lemma~\ref{asm-theta-agree-local-factor-ideal}) and to $0$ otherwise is
  $m_{A,v}$-integrable on $A_v$, and its integral $u_v(\nu)$ is not $0$.
\item Let $T\supset S$ be a finite set of finite places, and let $g_\infty$ on
  $A_\infty$ and $g_v$ on $A_v$, for $v\in T$, be functions such that
  $g_\infty\nu_z$ is $m_{A,\infty}$-integrable and every $g_v\nu_z$ is
  $m_{A,v}$-integrable. Let $G(l)=g_\infty(l^{(\infty)})\prod_{v\in T}g_v(l^{(v)})$
  if $l$ lies in $\mathfrak a$ at every finite place $v\notin T$, and $G(l)=0$
  otherwise. Then $G\nu_z$ is $\mu$-integrable and
  \[
  \Bigl(\prod_{v\in T\setminus S}u_v(\nu)\Bigr)\int G\nu_z\,d\mu
  =\kappa\,\nu(\mathfrak a)\,L^{S_0}(\nu,z)\int_{A_\infty}g_\infty\nu_z\,dm_{A,\infty}\prod_{v\in T}\int_{A_v}g_v\nu_z\,dm_{A,v},
  \]
  with $\nu(\mathfrak a)$ and $L^{S_0}(\nu,z)$ as in
  Definition~\ref{asm-theta-agree-parts}(4).
\end{enumerate}
\end{lemma}
\leanref{Development/WeilBound/ThetaLifts/LocalTheory/ConeShifts/ArchCone.lean}{PositivityTheory.ThKG.Loc.Cone.IdeleFub}
\leanref{Development/WeilBound/ThetaLifts/LocalTheory/ConeShifts/FubMain.lean}{PositivityTheory.ThKG.Loc.Cone.ideleFub}
\begin{proof}
That $l^{(v)}\in A_v$ is Lemma~\ref{asm-theta-agree-local-ideles-parts}(1) and (3),
as $l^{(v)}=l^{(1_v)}$; the same holds at $v=\infty$ with the idempotent
$1_\infty$. Let $m_w$, $t_w$, $r_T$ and $\beta_{r_T}$ be as in
Lemma~\ref{asm-theta-agree-local-ideles-tail}, and put
$\kappa_{\mathfrak a}=\prod_wN(w)^{-zm_w}$, a finite product.

\emph{The constant.} As $1_\infty+r_\emptyset=1$,
Lemma~\ref{asm-theta-agree-local-ideles-parts}(1) gives
$\mathbf A_E^\times\cong A_\infty\times A_{r_\emptyset}$, and by the uniqueness of
Haar measures $\mu$ corresponds to $c\,(m_{A,\infty}\times\beta_{r_\emptyset})$ for
some $c>0$. Put $\kappa=c\,\kappa_{\mathfrak a}\neq0$. Now let $\nu$, $S_0$ and $S$
be as in the statement, and put $\nu_z(l)=\nu(l)|l|^z$, a homomorphism.

\emph{Part (1).} Let $v\notin S$. Then $\nu$ is unramified at the places $w$ above
$v$, and $A_v\cong\prod_{w\mid v}E_w^\times$ with $m_{A,v}$ corresponding to the
product of the $\gamma_w=\beta_{e_w}$
(Lemma~\ref{asm-theta-agree-local-ideles-parts}). As $l$ lies in $\mathfrak a$ at
$v$ exactly when $\operatorname{ord}_w(l_w)\geq m_w$ for every $w\mid v$, Fubini's
theorem and Lemma~\ref{asm-theta-agree-local-ideles-tate} give the integrability
and $u_v(\nu)=\prod_{w\mid v}t_w^{m_w}(1-t_w)^{-1}\neq0$.

\emph{Part (2): Fubini.} Let $T\supset S$ and the functions $g_\infty$ and $g_v$ be
as in (2). By Lemma~\ref{asm-theta-agree-local-ideles-parts},
$A_{r_\emptyset}\cong\prod_{v\in T}A_v\times A_{r_T}$ with $\beta_{r_\emptyset}$
corresponding to $\prod_{v\in T}m_{A,v}\times\beta_{r_T}$, and every $l$ is
$l^{(\infty)}\prod_{v\in T}l^{(v)}\,l^{(r_T)}$. Whether $l$ lies in $\mathfrak a$ at
$v\notin T$ depends only on $l^{(r_T)}$, so
\[
G(l)\nu_z(l)=g_\infty(l^{(\infty)})\nu_z(l^{(\infty)})\prod_{v\in T}g_v(l^{(v)})\nu_z(l^{(v)})\cdot\mathcal T(l^{(r_T)}),
\]
with $\mathcal T$ as in Lemma~\ref{asm-theta-agree-local-ideles-tail}, which is
$\beta_{r_T}$-integrable by that lemma, applied with $T$ and
$T_0=T\cup\{v:m_w\neq0\text{ for some }w\mid v\}$. By Fubini's theorem $G\nu_z$ is
$\mu$-integrable and
\[
\int G\nu_z\,d\mu=c\int_{A_\infty}g_\infty\nu_z\,dm_{A,\infty}\prod_{v\in T}\int_{A_v}g_v\nu_z\,dm_{A,v}\cdot\int_{A_{r_T}}\mathcal T\,d\beta_{r_T}.
\]

\emph{Part (2): the Euler product.} By part (1) and
Lemma~\ref{asm-theta-agree-local-ideles-tail}, $\prod_{v\in T\setminus S}u_v(\nu)$
times the last integral is the limit over $T'$ of
$\prod_wt_w^{m_w}(1-t_w)^{-1}$, over the places $w$ above
$(T'\cup T_0)\setminus S$. As $m_w=0$ above $S$ and outside $T_0$, this product is
\[
\prod_wN(w)^{-zm_w}\nu_w(\varpi_w)^{m_w}\cdot\prod_{w}(1-t_w)^{-1}
=\kappa_{\mathfrak a}\,\nu(\mathfrak a)\prod_{w}(1-t_w)^{-1},
\]
the last product over the $w$ above $(T'\cup T_0)\setminus S$. As $T'$ grows it
tends to the Euler product over all finite $w$ not above $S$, which converges
absolutely as $z\geq2$; as $\nu$ is unramified at these $w$ and ramified at the
places above $S\setminus S_0$, this Euler product is $L^{S_0}(\nu,z)$
(Definition~\ref{asm-theta-agree-parts}(4)). Multiplying the Fubini identity by
$\prod_{v\in T\setminus S}u_v(\nu)$ gives the claim.
\end{proof}
% Lean: Cone.ideleFub (ConeShifts/FubMain.lean). The constant is c * kappa_a with c from sp.fub
% (realAd F 1 and rest empty) and kappa_a the product of npow w (-z)^count; part (1) is uLocP
% (prodW with c = 0); part (2) is sy.prod for the idempotents io (idem F v), v in T, and
% io (rest T), c1 for the infinite factor, tail for the remaining integral, and euler (lim_tw,
% charL, idealVal) for the Euler product; ef is the factorization of G nu_z.

\begin{lemma}\label{asm-theta-agree-local-factor-tate}
Let $(M,M_1,H,H',e,u)$ be a cone setting for a theta base $B$, $v$ a finite
place, $m$ a normalized local measure system for a subgroup $K\subset U_X$,
$\nu$ a continuous character of $\mathbf A_E^\times$, $z\geq2$ an integer and
$\phi$ a Schwartz-Bruhat function on $F_v^{I_M}$. Then the function
$G_\phi(l)=e_v(\rho_v(1_vu(l))\phi)$ on $A_v$ takes finitely many values and is
locally constant, and $l\mapsto\nu(l)|l|^z$ is $m_{A,v}$-integrable on the set
of $l\in A_v$ with $G_\phi(l)\neq0$.
\end{lemma}
\leanref{Development/WeilBound/ThetaLifts/LocalTheory/UnramifiedCohomology/SpanSoLoc.lean}{PositivityTheory.ThKG.Loc.FL.tate'}
\leanref{Development/WeilBound/ThetaLifts/LocalTheory/ConeShifts/KAtCompact.lean}{PositivityTheory.ThKG.Loc.Cone.Fin.tate}
\begin{proof}
Let $V_0$, $V_1$ and $\sigma$ be as in
Lemma~\ref{asm-theta-agree-local-factor-directions}, so that
$1_vu(l)=(1_v\sigma(l))V_0$. For $w$ at $v$, $\rho_v(w)$ depends only on the
$v$-components $w^\flat\in F_v^{I_M}\times F_v^{I_M}$ of $e^M(w)$; write
$\rho_v(w)=\rho^\flat(w^\flat)$, with $\rho^\flat$ given by the formula of
Lemma~\ref{asm-theta-agree-local-factor-evaluation}.

\emph{Step 1: a level.} As $L_X$ and $L_W$ are finitely generated, $\Gamma_v$ is
the set of sums $\sum_ja_jg_j$ with $a_j\in\mathbf A$ at $v$, $a_{j,v}\in\mathcal
O_v$, for finitely many fixed $g_j\in M_{n\times d}(E)$, so the $\gamma^\flat$,
$\gamma\in\Gamma_v$, lie in a compact set $C$. As $\phi$ is locally constant with
compact support and the character $\psi_v$ of $F_v$ is locally constant, there is
a neighbourhood $V$ of $0$ in $F_v^{I_M}\times F_v^{I_M}$ with
$\rho^\flat(w^\flat)\phi=\phi$ for $w^\flat\in V$. As $C$ is compact, there is a
neighbourhood $U$ of $0$ in $F_v$ with $sw^\flat\in V$ for $s\in U$ and $w^\flat\in
C$. It contains $N=q^j$ for some $j\geq1$, where $q$ is the norm of $v$. So
$\rho_v(N\gamma)\phi=\phi$ for $\gamma\in\Gamma_v$.

\emph{Step 2: the support.} The function
$g(w^\flat)=e_v(\rho^\flat(w^\flat)\phi)$ is locally constant, since $\phi$ and
$\psi_v$ are. As $l\mapsto(1_vu(l))^\flat$ is continuous, $G_\phi$ is locally
constant. By Step 1, Lemma~\ref{asm-theta-agree-local-factor-directions}(2), with
$a=N$ and $b=1_v\sigma(l)$, shows: if $G_\phi(l)\neq0$, then $N1_vu(l)\in\Gamma_v$.

\emph{Step 3: finitely many values.} If $G_\phi(l)\neq0$, then
$G_\phi(l)=g(N^{-1}(N1_vu(l))^\flat)$ with $(N1_vu(l))^\flat\in C$. The locally
constant function $w^\flat\mapsto g(N^{-1}w^\flat)$ takes finitely many values on
the compact set $C$. So $G_\phi$ takes finitely many values.

\emph{Step 4: integrability.} By
Lemma~\ref{asm-theta-agree-local-factor-ideal}(1) there is $\mathfrak a\neq0$
with $1_vu(l)\in\Gamma_v$ exactly when $l$ lies in $\mathfrak a$ at $v$. Let $\mu$
be a Haar measure on $\mathbf A_E^\times$
(Lemma~\ref{asm-theta-agree-stabilizer}). By
Lemma~\ref{asm-theta-agree-local-factor-ideles}(1), for the trivial character
and $S_0=S=\emptyset$, the function $l\mapsto\mathbf 1[l\in\mathfrak a\text{ at
}v]\,|l|^z$ is $m_{A,v}$-integrable on $A_v$. Let $l_0=N^{(v)}\in A_v$, equal to $N$
at $v$ and $1$ elsewhere. Then $1_vu(l_0l)=N1_vu(l)$, as $\sigma(N)=N$. If
$G_\phi(l)\neq0$, then by Step 2 $l_0l$ lies in $\mathfrak a$ at $v$, and
\[
|\nu(l)|l|^z|=|l|^z=|l_0|^{-z}\,\mathbf 1[l_0l\in\mathfrak a\text{ at }v]\,|l_0l|^z .
\]
The right side is integrable on $A_v$, by the invariance of the Haar measure
$m_{A,v}$ under translation by $l_0$. So $\nu(l)|l|^z$ is integrable on the set
where $G_\phi\neq0$, which is open and closed as $G_\phi$ is locally constant.
\end{proof}
% Lean: tate' (UnramifiedCohomology/SpanSoLoc.lean) = tate (ConeShifts/KAtCompact.lean) without
% the last step: dirs, hJ (IdealId0), hS (StabEx, for mu), hI (IdeleFub at nu = 1, S = empty),
% exists_level (lat_compact, schr_fix, exists_nat), schr_lc, cont_loc, key; l0 is pi v N.

\begin{lemma}\label{asm-theta-agree-local-factor-finite}
Let $(M,M_1,H,H',e,u)$ be a cone setting for a theta base $B$ and $v$ a finite
place.
\begin{enumerate}
\item For $l\in\mathbf A_E^\times$, $e_v(\rho_v(1_vu(l))\mathrm{vac}_v)$ is
  $\mathrm{vac}_{1,v}$ if $1_vu(l)\in\Gamma_v$ and $0$ otherwise.
\item Let $(X,W,M,\chi_X,\chi_W,\omega)$ be a Weil pair for $B$
  (Definition~\ref{asm-theta-agree-triple}(1)), $K\subset U_X$ the lattice
  stabilizer of $X$, $m$ a normalized local measure system for $K$, $\nu$ a
  continuous character of $\mathbf A_E^\times$ and $z\geq2$ an integer. Then there
  is a linear map $\Lambda$ from functions on $F_v^{I_M}$ to functions on
  $F_v^{I_{M_1}}$ such that for every Schwartz-Bruhat $\phi$: $\Lambda\phi$ is the
  local cone $\operatorname{cone}_v(\phi)$ for $(M,\omega,m,e_v,u,\nu,z)$
  (Definition~\ref{asm-theta-agree-local-cones}(3)); it is Schwartz-Bruhat; and for
  every $y$ the integrand in $\operatorname{cone}_v(\phi)(y)$ is
  $m_{A,v}$-integrable in $l$ for every $k\in K_v$, and its integral over $A_v$ is
  $m_{K,v}$-integrable in $k$.
\end{enumerate}
\end{lemma}
\leanref{Development/WeilBound/ThetaLifts/LocalTheory/ConeShifts/FinStmt.lean}{PositivityTheory.ThKG.Loc.Cone.FinCone2}
\leanref{Development/WeilBound/ThetaLifts/LocalTheory/ConeShifts/KAtCompact.lean}{PositivityTheory.ThKG.Loc.Cone.Fin.finCone2_of}
\leanref{Development/WeilBound/ThetaLifts/LocalTheory/ConeShifts/ActsRhoFin.lean}{PositivityTheory.ThKG.Loc.Cone.VacShift}
\leanref{Development/WeilBound/ThetaLifts/LocalTheory/ConeShifts/PureIte.lean}{PositivityTheory.ThKG.Loc.Cone.FinM}
\begin{proof}
Let $V_0$, $V_1$ and $\sigma$ be as in
Lemma~\ref{asm-theta-agree-local-factor-directions}.

(1) Put $\gamma=1_vu(l)=(1_v\sigma(l))V_0$; then $h(\gamma,\gamma)=0$. If
$\gamma\in\Gamma_v$, then $\rho_v(\gamma)\mathrm{vac}_v=\mathrm{vac}_v$ by
Lemma~\ref{asm-theta-agree-local-factor-directions}(3), and
$e_v(\mathrm{vac}_v)=\mathrm{vac}_{1,v}$. Otherwise, suppose
$e_v(\rho_v(\gamma)\mathrm{vac}_v)\neq0$. If $tV_1\in\Gamma_v$, then
$h(tV_1,tV_1)=0$, so $\rho_v(tV_1)\mathrm{vac}_v=\mathrm{vac}_v$ by
Lemma~\ref{asm-theta-agree-local-factor-directions}(3). Then
Lemma~\ref{asm-theta-agree-local-factor-directions}(2), with $\phi=\mathrm{vac}_v$,
$a=1$ and $b=1_v\sigma(l)$, gives $\gamma\in\Gamma_v$, a contradiction.

(2) Let $(\omega_v,\rho'_v)$ be local operators at $v$ for the Weil pair
(Lemma~\ref{asm-theta-agree-finite-inputs}(2)). For $w$ at $v$ and
Schwartz-Bruhat $\phi$, $\rho(w)$ acts at $v$ by $\phi\mapsto\rho'_v(w)\phi$ and by
$\phi\mapsto\rho_v(w)\phi$, so these agree
(Lemma~\ref{asm-theta-agree-finite-inputs-schur}(1)). The group $K$ is compact
(Lemma~\ref{asm-theta-agree-stabilizer}), so $K_v=K\cap U_X^{(v)}$ is compact.

\emph{The integrand.} Let $k\in K_v$, $l\in A_v$ and $\phi$ Schwartz-Bruhat. The
operator $\omega(k,1)$ acts at $v$ by $\phi\mapsto\omega_v(k,1)\phi$ and maps
Schwartz-Bruhat functions to Schwartz-Bruhat functions, and $\rho(1_vu(l))$ acts
at $v$ by $\chi\mapsto\rho_v(1_vu(l))\chi$. So
$T^v_{k,l}=\rho(1_vu(l))\,\omega(k,1)$ acts at $v$ by
$\phi\mapsto\rho_v(1_vu(l))\omega_v(k,1)\phi$, and by uniqueness
$(T^v_{k,l})_v\phi=\rho_v(1_vu(l))\omega_v(k,1)\phi$.

\emph{The integral over $A_v$.} For Schwartz-Bruhat $\chi$ put
\[
A(\chi)(y)=\int_{A_v}e_v(\rho_v(1_vu(l))\chi)(y)\,\nu(l)|l|^z\,dm_{A,v}(l).
\]
The values of $G_\chi(l)=e_v(\rho_v(1_vu(l))\chi)$ are Schwartz-Bruhat. By
Lemma~\ref{asm-theta-agree-local-factor-tate} and
Lemma~\ref{asm-theta-agree-local-factor-sums}, with $w=\nu(l)|l|^z$ and $D$ the set
where $G_\chi\neq0$, the integrand is integrable for every $y$ and $A(\chi)$ is
Schwartz-Bruhat. The map $A$ is linear on the Schwartz-Bruhat functions.

\emph{The integral over $K_v$.} Let $\phi$ be Schwartz-Bruhat. The map
$k\mapsto\omega_v(k,1)\phi$ is locally constant on $K_v$: given $k_0$,
Lemma~\ref{asm-theta-agree-finite-inputs}(2) gives a neighbourhood $U$ of $1$ in
$U_X$ whose elements $g$ at $v$ have $\omega_v(g,1)\omega_v(k_0,1)\phi=\omega_v(k_0,1)\phi$,
and $\omega_v(k,1)\phi=\omega_v(kk_0^{-1},1)\omega_v(k_0,1)\phi$ for $k\in K_v$. So
$k\mapsto A(\omega_v(k,1)\phi)$ is locally constant on the compact group $K_v$,
takes finitely many values, each on an open set, and these values are
Schwartz-Bruhat. By Lemma~\ref{asm-theta-agree-local-factor-sums}, with $w=1$ and
$m_{K,v}(K_v)=1$, the function $k\mapsto A(\omega_v(k,1)\phi)(y)$ is integrable and
\[
\operatorname{cone}_v(\phi)=\int_{K_v}A(\omega_v(k,1)\phi)\,dm_{K,v}(k)
\]
is Schwartz-Bruhat; the equality is the definition of the local cone, by the
description of the integrand.

\emph{The map $\Lambda$.} As $A$ and every $\omega_v(k,1)$ are linear on the
Schwartz-Bruhat functions, so is $\phi\mapsto\operatorname{cone}_v(\phi)$. Extend
this linear map from the subspace of Schwartz-Bruhat functions to all functions
on $F_v^{I_M}$, choosing a complementary subspace; this is $\Lambda$.
\end{proof}
% Lean: finCone2_of (ConeShifts/KAtCompact.lean): the first part with vac_fix, key and acts_uniq;
% the second with locRepEx, StabEx and kAt_compact, opF_eq and acts_comp (the integrand), tate
% = tate' with int_fin (A), lc from ActsEx (local constancy), int_fin (cone), and
% LinearMap.exists_extend. vacShift_of and finM_of read off VacShift and FinM.

\begin{lemma}\label{asm-theta-agree-local-stabilizer-parts}
Let $B$ be a theta base of a CM field $E$, let $X$ be a Hermitian lattice space of
some signature, and $K\subset U_X$ its lattice stabilizer. For $a\in\mathbf A$
with $a^2=a$ and $g\in U_X$ put $g^{(a)}=a\,g+(1-a)\cdot1$.
\begin{enumerate}
\item $g^{(a)}\in U_X$, with inverse $(g^{-1})^{(a)}$, and $g^{(a)}-1=a(g-1)$.
  The map $g\mapsto g^{(a)}$ is a continuous homomorphism of $U_X$. For
  idempotents $a,b$, $(g^{(b)})^{(a)}=g^{(ab)}$, $g^{(0)}=1$ and $g^{(1)}=g$; if
  $ab=0$, then $a+b$ is idempotent and $g^{(a+b)}=g^{(a)}g^{(b)}=g^{(b)}g^{(a)}$.
\item If the components of $a$ at the infinite places are all $0$ or all $1$, then
  $k^{(a)}\in K$ for $k\in K$.
\item Let $v$ be a finite place or $v=\infty$, and $1_v$ as in
  Definition~\ref{asm-theta-agree-local-cones}. Then an element $g\in U_X$ is at
  $v$ if and only if $g^{(1_v)}=g$, so that $K_v=\{k\in K:k^{(1_v)}=k\}$; and
  $k^{(1_v)}\in K_v$ for $k\in K$.
\item For finite sets $S$ of finite places put $a_S=1_\infty+\sum_{v\in S}1_v$.
  Then $g^{(a_S)}\to g$ in $U_X$ as $S$ increases.
\end{enumerate}
\end{lemma}
\leanref{Development/WeilBound/ThetaLifts/LocalTheory/ConeShifts/EvFinIdemSelf.lean}{PositivityTheory.ThKG.Loc.Cone.partU}
\leanref{Development/WeilBound/ThetaLifts/LocalTheory/ConeShifts/EvFinIdemSelf.lean}{PositivityTheory.ThKG.Loc.Cone.partU_add}
\leanref{Development/WeilBound/ThetaLifts/LocalTheory/ConeShifts/EvFinIdemSelf.lean}{PositivityTheory.ThKG.Loc.Cone.partU_mem_kAt}
\leanref{Development/WeilBound/ThetaLifts/LocalTheory/ConeShifts/EvFinIdemSelf.lean}{PositivityTheory.ThKG.Loc.Cone.tendsto_partU}
\leanref{Development/WeilBound/ThetaLifts/LocalTheory/ConeShifts/ExtendByZeroOnSubsetApplyCoe.lean}{PositivityTheory.ThKG.Loc.Cone.PU_mem}
\leanref{Development/WeilBound/ThetaLifts/LocalTheory/ConeShifts/ExtendByZeroOnSubsetApplyCoe.lean}{PositivityTheory.ThKG.Loc.Cone.kAt_eq}
\begin{proof}
(1) As $a(1-a)=0$, $(ax+(1-a))(ay+(1-a))=a\,xy+(1-a)$ for matrices $x,y$; with
$x=g$, $y=g^{-1}$ this gives the inverse, and in general the homomorphism property.
As $c(a)=a$ and $g^{\mathsf T}J\,c(g)=J$, for the matrix $J$ of $X$, the same expansion gives
$(g^{(a)})^{\mathsf T}J\,c(g^{(a)})=aJ+(1-a)J=J$. The other identities are
expansions of the same kind, and continuity holds as the entries of $g^{(a)}$ and
of its inverse are affine in those of $g$ and $g^{-1}$.

(2) If $a$ is $0$ at the infinite places, $\pi_w(k^{(a)})=1$; if it is $1$ there,
$\pi_w(k^{(a)})=\pi_w(k)$. Either way $\pi_w(k^{(a)})$ commutes with
$\theta_{X,w}$. The group $\hat L_X$ is stable under multiplication by $a$ and by
$1-a$, as their finite components are $0$ or $1$; so
$k^{(a)}z=a\,kz+(1-a)z\in\hat L_X$ for $z\in\hat L_X$, and likewise for
$(k^{(a)})^{-1}=(k^{-1})^{(a)}$.

(3) An element $z$ is at $v$ if and only if $1_vz=z$, and $g^{(1_v)}-1=1_v(g-1)$.
By the proof of Lemma~\ref{asm-theta-agree-local-measures}, $K_v=K\cap U_X^{(v)}$.
For $k\in K$, $k^{(1_v)}\in K$ by (2), and it is at $v$ by (1).

(4) The idempotent $1-a_S$ is $0$ at the infinite places and at $S$, and $1$
elsewhere, so it lies in every neighbourhood of $0$ in $\mathbf A$ of the form
$\prod_{v\in S_0\cup\{\infty\}}U_v\times\prod_{v\notin S_0}\mathcal O_v$ once
$S\supset S_0$. So $a_S\to1$, and the entries of $g^{(a_S)}$ and of its inverse
tend to those of $g$ and $g^{-1}$.
\end{proof}
% Lean: partU (EvFinIdemSelf.lean), mix_mul, mat_partU, partU_add, partU_sub_one; PU and PU_mem
% (ExtendByZeroOnSubsetApplyCoe.lean, with smul_latHat and piInf_alg; Gd is the condition on the
% infinite components); kAt_eq with at_fin_iff and at_inf_iff; partU_mem_kAt; tendsto_partU with
% tendsto_aS.

\begin{lemma}\label{asm-theta-agree-local-factor-stabilizer}
Let $B$ be a theta base of a CM field $E$, $X$ a Hermitian lattice space of some
signature, $K\subset U_X$ its lattice stabilizer, $m$ a normalized local measure
system for $K$ and $\mu_K$ a Haar measure on $K$. For $k\in K$ and $v$ a finite
place or $v=\infty$ put $k^{(v)}=1_vk+(1-1_v)1$. Then $k^{(v)}\in K_v$, and there
is $\kappa\neq0$ such that for every finite set $T$ of finite places, every
$m_{K,\infty}$-integrable $h_\infty$ on $K_\infty$ and all $m_{K,v}$-integrable
$h_v$ on $K_v$, $v\in T$, the function
$k\mapsto h_\infty(k^{(\infty)})\prod_{v\in T}h_v(k^{(v)})$ is $\mu_K$-integrable
and
\[
\int_Kh_\infty(k^{(\infty)})\prod_{v\in T}h_v(k^{(v)})\,d\mu_K(k)
=\kappa\int_{K_\infty}h_\infty\,dm_{K,\infty}\prod_{v\in T}\int_{K_v}h_v\,dm_{K,v}.
\]
\end{lemma}
\leanref{Development/WeilBound/ThetaLifts/LocalTheory/ConeShifts/ArchCone.lean}{PositivityTheory.ThKG.Loc.Cone.KFub}
\leanref{Development/WeilBound/ThetaLifts/LocalTheory/ConeShifts/ExtendByZeroOnSubsetApplyCoe.lean}{PositivityTheory.ThKG.Loc.Cone.kFub_of}
\begin{proof}
By Lemma~\ref{asm-theta-agree-local-stabilizer-parts}(3), $k^{(v)}=k^{(1_v)}\in K_v$.
For an idempotent $a\in\mathbf A$ whose infinite components are all $0$ or all
$1$, let $K_a=\{k\in K:k^{(a)}=k\}$, a closed, hence compact, subgroup of $K$
(Lemma~\ref{asm-theta-agree-stabilizer}); by
Lemma~\ref{asm-theta-agree-local-stabilizer-parts}(2) and (1) it is the image of
$K$ under $k\mapsto k^{(a)}$, and $K_{1_v}=K_v$ by part (3) of that lemma. For a
finite set $S$ of finite places put $r_S=1-1_\infty-\sum_{v\in S}1_v$, an
idempotent that is $0$ at the infinite places, and let $\beta_S$ be the Haar
measure of total mass $1$ on the compact group $K_{r_S}$.

\emph{Step 1: splitting.} Let $a$ and $b$ be idempotents as above with $ab=0$.
Then $(x,y)\mapsto xy$ is an isomorphism of topological groups from
$K_a\times K_b$ onto $K_{a+b}$, with inverse $z\mapsto(z^{(a)},z^{(b)})$: for
$x\in K_a$ and $y\in K_b$, $x$ and $y$ commute,
$(xy)^{(a)}=x\,(y^{(b)})^{(a)}=x\,y^{(0)}=x$ and likewise $(xy)^{(b)}=y$, and
$z=z^{(a)}z^{(b)}$ for $z\in K_{a+b}$ by
Lemma~\ref{asm-theta-agree-local-stabilizer-parts}(1). By the uniqueness of Haar
measures, the image of a Haar measure of $K_{a+b}$ is a positive multiple of the
product of Haar measures of $K_a$ and $K_b$.

\emph{Step 2: the constant.} As $1_\infty+r_\emptyset=1$ and $K_1=K$, Step 1 gives
$\kappa>0$ such that $\mu_K$ corresponds to $\kappa\,(m_{K,\infty}\times\beta_\emptyset)$
under $K_\infty\times K_{r_\emptyset}\cong K$.

\emph{Step 3: the finite places.} Let $T$ be a finite set of finite places. Then
$r_\emptyset=\sum_{v\in T}1_v+r_T$, and the idempotents $1_v$, $v\in T$, and $r_T$
are pairwise orthogonal. By induction on $T$, Step 1 gives an isomorphism
$\prod_{v\in T}K_v\times K_{r_T}\to K_{r_\emptyset}$,
$((x_v),y)\mapsto\prod_vx_v\cdot y$, with inverse
$z\mapsto((z^{(1_v)}),z^{(r_T)})$, under which $\beta_\emptyset$ corresponds to
$\prod_{v\in T}m_{K,v}\times\beta_T$: the constant of each step is $1$, as all
these measures have total mass $1$.

\emph{Step 4: the integral.} Let $h_\infty$ and $h_v$ be as in the statement. For
$k\in K$ write $k=k^{(\infty)}z$ with $z=k^{(r_\emptyset)}$; then
$k^{(v)}=z^{(1_v)}$ for $v\in T$, as $1_vr_\emptyset=1_v$. By Steps 2 and 3 and
Fubini's theorem for the finite product measures, the function
$k\mapsto h_\infty(k^{(\infty)})\prod_{v\in T}h_v(k^{(v)})$ is $\mu_K$-integrable,
and its integral is
\[
\kappa\int_{K_\infty}h_\infty\,dm_{K,\infty}\prod_{v\in T}\int_{K_v}h_v\,dm_{K,v}\cdot\beta_T(K_{r_T}),
\]
with $\beta_T(K_{r_T})=1$.
\end{proof}
% Lean: kFub_of StabEx (ConeShifts/ExtendByZeroOnSubsetApplyCoe.lean), with the framework of
% FubGen.lean: sysU (the system of maps PU a on K for the idempotents with Gd a), Sys.split and
% Split.fub (Step 1, with haarScalarFactor), sp.fub for realAd F 1 and rest empty (Step 2),
% Sys.prod (Step 3, the constants being 1 by the total masses), exists_haar for beta_T, and
% kAt_eq for K_v; ext0 extends the h_v by zero to U_X.

\begin{lemma}\label{asm-theta-agree-local-vacuum-lattice}
Let $B$ be a theta base of a CM field $E$, let $X$ be a Hermitian and $W$ a
skew-Hermitian lattice space, of ranks $n$ and $d$ and of some signatures, let
$K\subset U_X$ be the lattice stabilizer of $X$, and let $v$ be a finite place.
Let $L_v\subset\mathbf A_E^n$ be the group of
Definition~\ref{asm-theta-agree-finite-local-objects}(1), $\Lambda_v$ that of
Lemma~\ref{asm-theta-agree-evaluation-lattice}, and $K_v$ that of
Definition~\ref{asm-theta-agree-local-cones}(2). Say that an $n\times n$ matrix
$N$ over $\mathbf A_E$ \emph{preserves} $L_v$ if $NL_v\subset L_v$.
\begin{enumerate}
\item Every $k\in K_v$ is at $v$, and $k$ and $k^{-1}$ preserve $L_v$.
\item If $N$ preserves $L_v$, then $N\gamma\in\Lambda_v$ for every
  $\gamma\in\Lambda_v$.
\item If $z\in\mathbf A_E$ has its component in $E\otimes_FF_v$ in
  $\mathcal O_E\otimes\mathcal O_v$, then the scalar matrix $z$ preserves $L_v$.
\end{enumerate}
\end{lemma}
\leanref{Development/WeilBound/ThetaLifts/LocalTheory/ConeShifts/PresOfHat.lean}{PositivityTheory.ThKG.Loc.Cone.Vf.pres_of_hat}
\leanref{Development/WeilBound/ThetaLifts/LocalTheory/ConeShifts/StarMemSchA.lean}{PositivityTheory.ThKG.Loc.Cone.Vf.mul_latAtV}
\leanref{Development/WeilBound/ThetaLifts/LocalTheory/ConeShifts/MemIntLocOfIdem.lean}{PositivityTheory.ThKG.Loc.Cone.Vf.smul_mem_Pres}
\begin{proof}
Let $1_v$ be as in Definition~\ref{asm-theta-agree-local-cones}, and $\hat L$
the group $\hat L_X$ of Definition~\ref{asm-theta-agree-planes}(6).

(1) By the proof of Lemma~\ref{asm-theta-agree-local-measures},
$K_v=K\cap U_X^{(v)}$, so $k$ and $k^{-1}$ are at $v$ and in $K$. Let $N$ be
$k$ or $k^{-1}$, so that $N\hat L\subset\hat L$ and $N-1$ is at $v$. Every
generator $a\,y$ of $L_v$ has $a_\infty=0$ and $a_u\in\mathcal O_u$ for every
finite $u$, so $L_v\subset\hat L$; and $1_v$ maps the generators $a\,y$ of $\hat L$
to the generators $(1_va)y$ of $L_v$, so $1_v\hat L\subset L_v$. Let $y\in L_v$.
Then $Ny\in\hat L$, and $Ny=y+(N-1)y$ is at $v$, so $Ny=1_vNy\in L_v$.

(2) For a generator $a\,xy^{\mathsf T}$ of $\Lambda_v$, with $x\in L_X$ and
$y\in L_W$, $N(a\,xy^{\mathsf T})=(N(ax))y^{\mathsf T}$, and $N(ax)\in L_v$ is a
sum of terms $\pm a'x'$ with $a'$ at $v$, $a'_v\in\mathcal O_v$ and $x'\in L_X$.
So $N(a\,xy^{\mathsf T})$ is a sum of terms $\pm a'x'y^{\mathsf T}\in\Lambda_v$. The
elements $\gamma$ with $N\gamma\in\Lambda_v$ form a group, so they include
$\Lambda_v$.

(3) For a generator $a\,y$ of $L_v$, $za=(1_vz)a$, and $1_vz$ is a sum of
products $a_io_i$ with $a_i\in\mathbf A$ at $v$, $a_{i,v}\in\mathcal O_v$ and
$o_i\in\mathcal O_E$. So $z(ay)=\sum_i(a_ia)(o_iy)$, and $o_iy\in L_X$ as $L_X$ is
an $\mathcal O_E$-module.
\end{proof}
% Lean: (1) is pres_of_hat (PresOfHat.lean, with up and dn), with hU in vacFixX_of; (2) is
% mul_latAtV (StarMemSchA.lean); (3) is smul_mem_Pres through exists_fr (MemIntLocOfIdem.lean).

\begin{lemma}\label{asm-theta-agree-local-vacuum-residue}
Let $\kappa$ be a finite field, $N,N'\in M_n(\kappa)$ with $N'N=1$, and
$h,\delta\in\kappa$ with $2h=1$ and $\delta\neq0$. Put $Y=2+N+N'$. Then there is
an integer $m\geq1$ such that $1+(1-N)\bigl(-h+\delta(1-Y^m)\bigr)$ is
invertible.
\end{lemma}
\leanref{Development/WeilBound/ThetaLifts/LocalTheory/ConeShifts/ExistsPowIdem.lean}{PositivityTheory.ThKG.Loc.Cone.Vf.exists_det_ne}
\begin{proof}
As $M_n(\kappa)$ is finite, the powers of $Y$ repeat, and there is $m\geq1$ with
$e=Y^m$ idempotent. Put $A=1+(1-N)(-h+\delta(1-e))$ and $U=1+N$. As $NN'=1$,
$N$ commutes with $Y$ and $e$, and $Y=N'U^2$. Suppose $Ay=0$ with $y\neq0$.

From $e(1-e)=0$ and $2h=1$, $eA=e-h(1-N)e=hUe$. So $Uey=0$. As
$e=e^2=Y^{m-1}N'U^2e$, $ey=Y^{m-1}N'U\,(Uey)=0$. As
$A=2\delta-(\delta-h)U-\delta(1-N)e$, $Ay=0$ gives $2\delta y=(\delta-h)Uy$, so
$y=\lambda^jU^jy$ for every $j$, with $\lambda=(\delta-h)/(2\delta)$. But
$U^{2m}=(NY)^m=N^me$, so $U^{2m}y=N^mey=0$ and $y=0$, a contradiction.
\end{proof}
% Lean: exists_det_ne with exists_pow_idem (ExistsPowIdem.lean); f1 is eA = hUe, y2 is ey = 0,
% y3 and y4 the last step.

\begin{lemma}\label{asm-theta-agree-local-vacuum-units}
Let $B$ be a theta base of a CM field $E$, let $X$ be a Hermitian lattice space of
rank $n$ and some signature, with matrix $J$, and let $v$ be a finite place. Let
$k\in U_X$ be at $v$, with $k$ and $k^{-1}$ preserving $L_v$
(Lemma~\ref{asm-theta-agree-local-vacuum-lattice}). Then there are $n\times n$
matrices $b$, $b'$ and $f'$ over $\mathbf A_E$ such that
$b^{\mathsf T}J=J\,c(b')$, $b'^{\mathsf T}J=J\,c(b)$, $b+b'=-1$, $b$ and $f'$
preserve $L_v$, and $(1+(1-k)b)f'=1$.
\end{lemma}
\leanref{Development/WeilBound/ThetaLifts/LocalTheory/ConeShifts/MemIntLocOfIdem.lean}{PositivityTheory.ThKG.Loc.Cone.Vf.exists_beta}
\begin{proof}
Put $E_v=E\otimes_FF_v$ and $\mathcal O_{E,v}=\mathcal O_E\otimes\mathcal O_v$, and
for a matrix $N$ over $\mathbf A_E$ let $N_v$ be its component over $E_v$. The
group $L_v$ consists of the vectors at $v$ whose component lies in the
$\mathcal O_{E,v}$-lattice $L\otimes\mathcal O_v\subset E_v^n$, which is free of
rank $n$. Let $a\in GL_n(E_v)$ have a basis of it as columns. Then $N$ preserves
$L_v$ if and only if $a^{-1}N_va$ has entries in $\mathcal O_{E,v}$. Let $1_v$ be as
in Definition~\ref{asm-theta-agree-local-cones}; as $k-1$ is at $v$,
$1_v(1-k)=1-k$.

\emph{Split case.} Suppose there is an $F$-embedding $\tau\colon E\to F_v$. Then
$E_v\cong F_v\times F_v$ through $(\tau,\tau\circ c)$, and $c$ exchanges the two
factors. Let $\eta\in\mathbf A_E$ be at $v$ with component $(1,0)$. Then
$\eta^2=\eta$, $\eta+c(\eta)=1_v$, and $\eta_v\in\mathcal O_{E,v}$. Put
$\vartheta=\eta+\tfrac12(1-1_v)$, $b=-\vartheta$, $b'=-c(\vartheta)$ (scalar
matrices) and $f'=1-\eta(1-k^{-1})$. Then
$b^{\mathsf T}J=-\vartheta J=J\,c(b')$ and $b'^{\mathsf T}J=J\,c(b)$, as $c$ is an
involution, and $b+b'=-(1_v+(1-1_v))=-1$. By
Lemma~\ref{asm-theta-agree-local-vacuum-lattice}(3), $b$ and $f'$ preserve
$L_v$, since $\vartheta_v=\eta_v$. As $(1-1_v)(1-k)=0$,
$1+(1-k)b=1-\eta(1-k)$, and
\[
(1-\eta(1-k))(1-\eta(1-k^{-1}))=1-\eta(2-k-k^{-1})+\eta^2(1-k)(1-k^{-1})=1 .
\]

\emph{Inert case.} Otherwise, as $E/F$ is unramified at the finite places, $v$
is inert in $E$, and by the conditions on a good base $v\nmid2$. Let $w$ be the
place of $E$ above $v$; then $\mathcal O_{E,v}=\mathcal O_w$, and $c$ induces the
nontrivial automorphism of the residue field $\kappa_w$, which is finite and
quadratic over that of $v$. Choose $x\in\mathcal O_w$ whose residue is not fixed
by it, and let $\delta\in\mathbf A_E$ be at $v$ with $\delta_v=c(x)-x$, a unit of
$\mathcal O_w$; then $c(\delta)=-\delta$. Put $h=\tfrac12$, whose component at
$v$ is integral. Let $N=a^{-1}k_va$ and $N'=a^{-1}k_v^{-1}a$, matrices over
$\mathcal O_w$ with $N'N=1$. By Lemma~\ref{asm-theta-agree-local-vacuum-residue},
applied over $\kappa_w$ to the residues of $N$, $N'$, $h$ and $\delta_v$, there is
$m\geq1$ such that the residue of
$1+(1-N)(-h+\delta_v(1-(2+N+N')^m))$ is invertible. Put $Y=2+k+k^{-1}$,
\[
b=-h+\delta(1-Y^m),\qquad b'=-h-\delta(1-Y^m).
\]
As $k\in U_X$, $k^{\mathsf T}J=J\,c(k^{-1})$ and $(k^{-1})^{\mathsf T}J=J\,c(k)$, so
$(Y^m)^{\mathsf T}J=J\,c(Y^m)$; with $c(\delta)=-\delta$ and $c(h)=h$ this gives
$b^{\mathsf T}J=J\,c(b')$ and $b'^{\mathsf T}J=J\,c(b)$, and $b+b'=-2h=-1$. By
Lemma~\ref{asm-theta-agree-local-vacuum-lattice}(3), $b$ preserves $L_v$. The
matrix $A=1+(1-k)b$ has $A-1$ at $v$, and $a^{-1}A_va$ is the matrix above, with
entries in $\mathcal O_w$ and a unit determinant. So $A_v$ is invertible and
$a^{-1}A_v^{-1}a$ has entries in $\mathcal O_w$. Let $f'$ be $1$ away from $v$ and
$A_v^{-1}$ at $v$. Then $Af'=1$ and $f'$ preserves $L_v$.
\end{proof}
% Lean: exists_beta (MemIntLocOfIdem.lean). The split case is the second branch, with exists_eta
% (eta) and mem_intLoc_of_idem; the inert case is the first, with inert_of_local, exists_isUnramInv
% (x with |c x - x| = 1), Formalism.UG.P.tjcosets_finite, exists_det_ne and exists_inv; exists_fr
% gives the matrix a.

\begin{lemma}\label{asm-theta-agree-local-vacuum-decomposition}
Let $E$ be a CM field with complex conjugation $c$, let $X$ be a lattice space of
rank $n$ with invertible matrix $J$, and let $G_X$, $P_X$ and $\Delta$ be as in
Definition~\ref{asm-theta-agree-doubling}(2). Let $k\in U_X$, let $b$, $b'$ and
$f'$ be $n\times n$ matrices over $\mathbf A_E$ with $b^{\mathsf T}J=J\,c(b')$,
$b'^{\mathsf T}J=J\,c(b)$, $b+b'=-1$ and $(1+(1-k)b)f'=1$, and let
$\varepsilon\in\mathbf A_E$ with $\varepsilon^2=\varepsilon$,
$c(\varepsilon)=\varepsilon$ and $\varepsilon(1-k)=1-k$. Put
\[
C=\begin{pmatrix}1&1+b\\1&b\end{pmatrix},\quad
S=\begin{pmatrix}1-\varepsilon&\varepsilon\\\varepsilon&1-\varepsilon\end{pmatrix},\quad
T=-f'(1-k),\quad
s=CSC^{-1},\quad n_0=C\begin{pmatrix}1&T\\0&1\end{pmatrix}C^{-1},
\]
and $t=1+b(1-k)$. Then
$C^{-1}=\left(\begin{smallmatrix}-b&1+b\\1&-1\end{smallmatrix}\right)$; $s$ and
$n_0$ lie in $G_X$; $s^2=1$;
$n_0^{-1}=C\left(\begin{smallmatrix}1&-T\\0&1\end{smallmatrix}\right)C^{-1}$ and
$n_0\Delta=\Delta$; $t$ is invertible with $t^{-1}=1-bf'(1-k)$; and
$p=(sn_0s)^{-1}\operatorname{diag}(k,1)$ lies in $G_X$ and has $p\Delta=\Delta t$.
\end{lemma}
\leanref{Development/WeilBound/ThetaLifts/LocalTheory/ConeShifts/ExistsPowIdem.lean}{PositivityTheory.ThKG.Loc.Cone.Vf.exists_decomp}
\leanref{Development/WeilBound/ThetaLifts/LocalTheory/ConeShifts/ExistsPowIdem.lean}{PositivityTheory.ThKG.Loc.Cone.Vf.skew_alt}
\leanref{Development/WeilBound/ThetaLifts/LocalTheory/ConeShifts/ExistsPowIdem.lean}{PositivityTheory.ThKG.Loc.Cone.Vf.gram_CC}
\begin{proof}
Put $J'=\left(\begin{smallmatrix}0&J\\J&0\end{smallmatrix}\right)$ and
$J^\square=\operatorname{diag}(J,-J)$. Multiplying out, $C^{-1}C=1$. As
$c(b)=-1-c(b')$, $b^{\mathsf T}J=J\,c(b')=-J-J\,c(b)$, and with this,
$C^{\mathsf T}J^\square c(C)=J'$. So a matrix $x$ lies in $G_X$ when
$C^{-1}xC$ preserves $J'$. As $c(\varepsilon)=\varepsilon$, $S$ preserves $J'$, and
$S^2=1$ as $\varepsilon(1-\varepsilon)=0$; so $s\in G_X$ and $s^2=1$. The matrix
$\left(\begin{smallmatrix}1&T\\0&1\end{smallmatrix}\right)$ preserves $J'$ if
$T^{\mathsf T}J+J\,c(T)=0$. Put $m=1-k$, $m'=1-k^{-1}$, $f=1+mb$ and
$g=1+b'm'$. As $k\in U_X$, $m^{\mathsf T}J=J\,c(m')$ and
$m'^{\mathsf T}J=J\,c(m)$, so $f^{\mathsf T}J=J\,c(g)$ and $g^{\mathsf T}J=J\,c(f)$;
$f'f=1$, $fT=-m$, and $mm'=m+m'$ gives $fm'+mg=0$. Hence
\[
g^{\mathsf T}\bigl(T^{\mathsf T}J+J\,c(T)\bigr)c(g)=-J\,c(fm')-J\,c(mg)=0,
\]
and $g$ is invertible, as $g^{\mathsf T}=J\,c(f)J^{-1}$. So $n_0\in G_X$. The
formula for $n_0^{-1}$ is clear. As $C^{-1}\Delta=\left(\begin{smallmatrix}1\\0\end{smallmatrix}\right)$
and $C\left(\begin{smallmatrix}1\\0\end{smallmatrix}\right)=\Delta$, $n_0\Delta=\Delta$.

As $\varepsilon T=T$, $S\left(\begin{smallmatrix}1&T\\0&1\end{smallmatrix}\right)S=\left(\begin{smallmatrix}1&0\\T&1\end{smallmatrix}\right)$,
so $sn_0s\Delta=C\left(\begin{smallmatrix}1\\T\end{smallmatrix}\right)$, with blocks
$1+(1+b)T$ and $1+bT$. Now $Tt=-f'(1+mb)m=-m$, so $(1+bT)t=1$ and
$(1+(1+b)T)t=k$. Hence $sn_0s\Delta t$ has blocks $k$ and $1$, those of
$\operatorname{diag}(k,1)\Delta$; that is, $p\Delta=\Delta t$; and $t^{-1}=1+bT$. Finally
$\operatorname{diag}(k,1)\in G_X$, so $p\in G_X$.
\end{proof}
% Lean: exists_decomp (ExistsPowIdem.lean), with gram_CC, mkU (the conjugation by CC b = C),
% skew_alt (T^T J + J c(T) = 0), SS (S), hL (S n S), and e2 for p Delta = Delta t; diag(k,1) in
% G_X is GlobalTheory.unitary_diag_one_mem_dbl.

\begin{lemma}\label{asm-theta-agree-local-vacuum-doubled}
Let $(X,W,M,\chi_X,\chi_W,\omega)$ be a Weil pair for a theta base $B$
(Definition~\ref{asm-theta-agree-triple}(1)), with $X$ of rank $n$, $I=I_M$ and
vacua $\mathrm{vac}_v$, and let $\omega^\square(x)$, $x\in G_X$, be operators on
functions on $\mathbf A^I\times\mathbf A^I=\mathbf A^{I\sqcup I}$ as in condition
(b) of Definition~\ref{asm-theta-agree-models}(6). Let $v$ be a finite place and
$g\in G_X$ with $g-1$ at $v$, such that the four $n\times n$ blocks of $g^{-1}$
preserve $L_v$ (Lemma~\ref{asm-theta-agree-local-vacuum-lattice}). Then there is
$a\in\mathbf C$ such that $\omega^\square(g)$ acts at $v$ by
$\mathrm{vac}_v\boxtimes\overline{\mathrm{vac}_v}\mapsto a\,\mathrm{vac}_v\boxtimes\overline{\mathrm{vac}_v}$.
\end{lemma}
\leanref{Development/WeilBound/ThetaLifts/LocalTheory/ConeShifts/StarMemSchA.lean}{PositivityTheory.ThKG.Loc.Cone.Vf.dbl_eigen}
\leanref{Development/WeilBound/ThetaLifts/LocalTheory/ConeShifts/StarMemSchA.lean}{PositivityTheory.ThKG.Loc.Cone.Vf.eigen}
\begin{proof}
Let $e=(e_1,e_2)$ be the isomorphism of $M$, $\rho$ its Schrödinger operators,
and $e^\square$ and $\rho^\square(w)=\rho(e^\square(w))$, for
$w\in M_{d\times2n}(\mathbf A_E)$, the isomorphism and Schrödinger operators of
condition (b), as in the proof of Lemma~\ref{asm-theta-agree-local-real-unique}.
Put $\phi_0=\mathrm{vac}_v\boxtimes\overline{\mathrm{vac}_v}$, a nonzero
Schwartz-Bruhat function on $F_v^{I\sqcup I}$, and let $x_0$ and $r_0$ be as in
Lemma~\ref{asm-theta-agree-finite-inputs-slice}(4), for $I\sqcup I$. For $w$ at
$v$, $\rho^\square(w)$ acts at $v$ by $\phi\mapsto\rho^\square_v(w)\phi$, where
$\rho^\square_v(w)$ is given by the formula of
Lemma~\ref{asm-theta-agree-local-factor-evaluation} with $e^\square$; the same
holds for $\rho$ and $\rho_v$.

\emph{The eigencharacter.} Let $\theta$ be the element in the definition of the
vacuum $\mathrm{vac}_v$, and
$\chi(\gamma)=\psi(\tfrac12\operatorname{Tr}_{\mathbf A_E/\mathbf A}h(\theta\gamma,\gamma))$
for $\gamma\in\Lambda_v$, so that $\rho_v(\gamma)\mathrm{vac}_v=\chi(\gamma)\mathrm{vac}_v$
(Lemma~\ref{asm-theta-agree-finite-inputs-schur}(1)) and $|\chi(\gamma)|=1$. Let
$\Gamma$ be the set of $z=(\gamma_1^{\mathsf T}\ \gamma_2^{\mathsf T})$ with
$\gamma_1,\gamma_2\in\Lambda_v$, and put
$\chi^\square(z)=\chi(\gamma_1)\overline{\chi(\gamma_2)}$; it is $\psi$ of half the
trace of the form of $\theta z^{\mathsf T}$ and $z^{\mathsf T}$ for the matrices
$J^\square=\operatorname{diag}(J_X,-J_X)$ and $J_W$. As
$e^\square(z)=((e_1(\gamma_1),e_1(\gamma_2)),(e_2(\gamma_1),-e_2(\gamma_2)))$ and
$\psi(-s)=\overline{\psi(s)}$, the formula gives
$\rho^\square_v(z)(\phi\boxtimes\phi')=\rho_v(\gamma_1)\phi\boxtimes\overline{\rho_v(\gamma_2)\overline{\phi'}}$,
so $\rho^\square_v(z)\phi_0=\chi^\square(z)\phi_0$.

\emph{Schur.} The operator $\omega^\square(g)$ maps Schwartz-Bruhat functions to
Schwartz-Bruhat functions, preserves the admissible forms, and has
$\omega^\square(g)\rho^\square(w)=\rho^\square(wg^{\mathsf T})\omega^\square(g)$
on Schwartz-Bruhat functions. If $1_vw=0$, then $w(g-1)^{\mathsf T}=0$ as $g-1$
is at $v$, so $wg^{\mathsf T}=w$. As $e^\square$ is $\mathbf A$-linear and onto,
Lemma~\ref{asm-theta-agree-finite-inputs-schur}(2) shows that $\omega^\square(g)$
acts at $v$ by $\phi_0\mapsto\phi_1$, with
$\phi_1=\operatorname{res}_{x_0}(\omega^\square(g)(\phi_0\otimes r_0))$
Schwartz-Bruhat.

\emph{$\phi_1$ is an eigenvector.} Let $z\in\Gamma$ and $z'=z(g^{-1})^{\mathsf T}$.
If $g^{-1}$ has blocks $g_{11},g_{12},g_{21},g_{22}$, then $z'^{\mathsf T}$ has
blocks $g_{11}\gamma_1+g_{12}\gamma_2$ and $g_{21}\gamma_1+g_{22}\gamma_2$, which lie in
$\Lambda_v$ by Lemma~\ref{asm-theta-agree-local-vacuum-lattice}(2); so $z'\in\Gamma$.
As $g^{-1}$ preserves $J^\square$, $\chi^\square(z')=\chi^\square(z)$. Since
$\rho^\square(z)$ acts at $v$ by $\phi_0\mapsto\chi^\square(z)\phi_0$,
\[
\rho^\square(z)\,\omega^\square(g)(\phi_0\otimes r_0)=\omega^\square(g)\,\rho^\square(z')(\phi_0\otimes r_0)
=\chi^\square(z)\,\omega^\square(g)(\phi_0\otimes r_0)=\chi^\square(z)\,\phi_1\otimes r_0,
\]
and the left side is $(\rho^\square_v(z)\phi_1)\otimes r_0$. Evaluating at the
points $x_0\langle y\rangle$ gives $\rho^\square_v(z)\phi_1=\chi^\square(z)\phi_1$.

\emph{The line.} Let $w$ be a pair of vectors over $F_v^{I\sqcup I}$ with
$\psi_v([z_v,w])=1$ for every $z\in\Gamma$, where $z_v$ is the component of
$e^\square(z)$ at $v$ and $[\ ,\ ]$ is as in the proof of
Lemma~\ref{asm-theta-agree-evaluation-vacuum}. Let $z''=(\zeta_1^{\mathsf T}\ \zeta_2^{\mathsf T})$
be the element at $v$ with $e^\square(z'')$ equal to $w$ at $v$. From the
formula for $e^\square$ and
$\langle u,u'\rangle=e_1(u)\cdot e_2(u')-e_1(u')\cdot e_2(u)$,
$[z_v,w]$ is the component at $v$ of
$\langle\gamma_1,\zeta_1\rangle-\langle\gamma_2,\zeta_2\rangle$. Taking $\gamma_2=0$,
then $\gamma_1=0$, Lemma~\ref{asm-theta-agree-evaluation-lattice}(2) gives
$\zeta_1,\zeta_2\in\Lambda_v$, so $w=z''_v$ with $z''\in\Gamma$. This is the
property of $\Lambda_v$ that the last paragraph of the proof of
Lemma~\ref{asm-theta-agree-evaluation-vacuum} uses. That paragraph, for the set
$I\sqcup I$, the group $\Gamma$, the character $\chi^\square$ and the nonzero
function $\phi_0$, shows that every Schwartz-Bruhat $\eta$ with
$\rho^\square_v(z)\eta=\chi^\square(z)\eta$ for $z\in\Gamma$ is a multiple of
$\phi_0$. So $\phi_1=a\phi_0$.
\end{proof}
% Lean: dbl_eigen (StarMemSchA.lean) applies eigen to the doubled coordinates dd: the set
% {fromCols g1^T g2^T}, the self-duality by mem_latAtV and symp_dbl, the eigencharacter by
% schr_box and star_psi, A = actB g with g^-1 in PresD and hTens_mul. eigen is acts_res (Schur)
% followed by local_line, the line of Lemma asm-theta-agree-evaluation-vacuum.

\begin{lemma}\label{asm-theta-agree-local-factor-vacuum}
Let $(X,W,M,\chi_X,\chi_W,\omega)$ be a Weil pair for a theta base $B$, with
vacua $\mathrm{vac}_v$, let $K\subset U_X$ be the lattice stabilizer of $X$, and
let $v$ be a finite place such that $\chi_X$ and $\chi_W$ are unramified at every
place of $E$ above $v$. Then $\omega(k,1)$ acts at $v$ by
$\mathrm{vac}_v\mapsto\mathrm{vac}_v$ for every $k\in K_v$.
\end{lemma}
\leanref{Development/WeilBound/ThetaLifts/LocalTheory/ConeShifts/ArchCone.lean}{PositivityTheory.ThKG.Loc.Cone.VacFixX}
\leanref{Development/WeilBound/ThetaLifts/LocalTheory/ConeShifts/PresOfHat.lean}{PositivityTheory.ThKG.Loc.Cone.vacFixX_of}
\begin{proof}
Let $n$ and $d$ be the ranks of $X$ and $W$, $I=I_M$, $\Lambda_v$ as in
Lemma~\ref{asm-theta-agree-evaluation-lattice}, and $x_0$ and $r_0$ as in
Lemma~\ref{asm-theta-agree-finite-inputs-slice}(4). Let $k\in K_v$. By
Lemma~\ref{asm-theta-agree-local-vacuum-lattice}(1), $k$ and $k^{-1}$ are at $v$
and preserve $L_v$, and so does every element of $K_v$.

\emph{Step 1: an eigenvalue.} Let $\theta$ and $\chi$ be as in the proof of
Lemma~\ref{asm-theta-agree-local-vacuum-doubled}. The operator $\omega(k,1)$ maps
Schwartz-Bruhat functions to Schwartz-Bruhat functions, preserves the admissible
forms, and $\omega(k,1)\rho(w)=\rho(kw)\omega(k,1)$ on Schwartz-Bruhat functions;
if $1_vw=0$, then $kw=w$. By
Lemma~\ref{asm-theta-agree-finite-inputs-schur}(2), $\omega(k,1)$ acts at $v$ by
$\mathrm{vac}_v\mapsto\phi_1$ for a Schwartz-Bruhat $\phi_1$. Let
$\gamma\in\Lambda_v$. By Lemma~\ref{asm-theta-agree-local-vacuum-lattice}(2),
$\gamma'=k^{-1}\gamma\in\Lambda_v$, and $\chi(\gamma')=\chi(\gamma)$ as $k\in U_X$.
As in the proof of Lemma~\ref{asm-theta-agree-local-vacuum-doubled}, the identity
$\rho(\gamma)\omega(k,1)(\mathrm{vac}_v\otimes r_0)=\omega(k,1)\rho(\gamma')(\mathrm{vac}_v\otimes r_0)$
gives $\rho_v(\gamma)\phi_1=\chi(\gamma)\phi_1$. As $\mathrm{vac}_v\neq0$ has the
same property, the last paragraph of the proof of
Lemma~\ref{asm-theta-agree-evaluation-vacuum} gives $\phi_1=c\,\mathrm{vac}_v$ for
some $c\in\mathbf C$. In the same way, $\omega(k^{-1},1)$ acts at $v$ by
$\mathrm{vac}_v\mapsto c'\mathrm{vac}_v$ for some $c'$. Put
$\Phi=\mathrm{vac}_v\otimes r_0$, a Schwartz-Bruhat function with
$\Phi(x_0\langle y\rangle)=\mathrm{vac}_v(y)$. Then $\omega(k,1)\Phi=c\Phi$ and
$\omega(k^{-1},1)\Phi=c'\Phi$, so $cc'\Phi=\Phi$ and $cc'=1$, as $\Phi\neq0$.

\emph{Step 2: $|c|=1$.} Choose $y$ with $\mathrm{vac}_v(y)\neq0$. The linear form
$\Psi\mapsto\Psi(x_0\langle y\rangle)$ is admissible
(Lemma~\ref{asm-theta-agree-finite-inputs-slice}(3)), so
$f(g)=(\omega(g,1)\Phi)(x_0\langle y\rangle)$ is continuous on $U_X$ by condition
(c) of Definition~\ref{asm-theta-agree-models}(6). As $K$ is compact
(Lemma~\ref{asm-theta-agree-stabilizer}), $f$ is bounded on $K$. But
$f(k^j)=c^j\mathrm{vac}_v(y)$ for $j\geq0$. So $|c|\leq1$, likewise
$|c'|\leq1$, and $|c|=1$ as $cc'=1$.

\emph{Step 3: the doubled operators.} Let $\omega^\square(x)$, $x\in G_X$, be the
operators of condition (b) of Definition~\ref{asm-theta-agree-models}(6) for
$\omega$, with the character $\chi_W$ and $m=d$, and put
$\Phi^\square=\Phi\boxtimes\overline\Phi$, a nonzero Schwartz-Bruhat function on
$\mathbf A^{I\sqcup I}$. It is
$(\mathrm{vac}_v\boxtimes\overline{\mathrm{vac}_v})\otimes(r_0\boxtimes\overline{r_0})$,
and $r_0\boxtimes\overline{r_0}$ does not depend on the $v$-components. So if
$g\in G_X$ is as in Lemma~\ref{asm-theta-agree-local-vacuum-doubled}, then
$\omega^\square(g)\Phi^\square=a\Phi^\square$ for some $a\in\mathbf C$. As
$\operatorname{diag}(k,1)\in G_X$,
$\omega^\square(\operatorname{diag}(k,1))\Phi^\square=(\omega(k,1)\Phi)\boxtimes\overline\Phi=c\Phi^\square$.

\emph{Step 4: a decomposition.} Let $J$ be the matrix of $X$, invertible as
$\tau_w(J)$ has a Cartan involution, and let $b$, $b'$ and $f'$ be as in
Lemma~\ref{asm-theta-agree-local-vacuum-units} for $k$. Let $\varepsilon$ be the
image of $1_v$ in $\mathbf A_E$; then $\varepsilon^2=\varepsilon$,
$c(\varepsilon)=\varepsilon$ and $\varepsilon(1-k)=1-k$. Let $C$, $S$, $T$, $s$,
$n_0$, $t$ and $p$ be as in Lemma~\ref{asm-theta-agree-local-vacuum-decomposition}.
The blocks $1$, $1+b$, $b$, $-b$ and $-1$ of $C$ and $C^{-1}$, the blocks
$\varepsilon$ and $1-\varepsilon$ of $S$ and the matrix $f'(1-k)$ preserve $L_v$
(Lemma~\ref{asm-theta-agree-local-vacuum-lattice}(3)), and so do the blocks of a
product of block matrices with such blocks. Hence the blocks of $s^{-1}=s$ and of
$n_0^{-1}$ preserve $L_v$. Also $s-1=C(S-1)C^{-1}$ and
$n_0-1=C\left(\begin{smallmatrix}0&T\\0&0\end{smallmatrix}\right)C^{-1}$ are at
$v$, as $S-1$ has entries $\pm\varepsilon$ and $T=\varepsilon T$. So by Step 3
there are $a,a'$ with $\omega^\square(s)\Phi^\square=a\Phi^\square$ and
$\omega^\square(n_0)\Phi^\square=a'\Phi^\square$. As $s^2=1$, $a^2=1$.

\emph{Step 5: the Siegel rule.} Let $l$ be the linear form of
Lemma~\ref{asm-theta-agree-local-real-unique-diagonal}, for $\kappa=I$. As in
Step 3 of the proof of Lemma~\ref{asm-theta-agree-local-real-unique},
$l(\rho^\square(z\Delta^{\mathsf T})\Psi)=l(\Psi)$ for Schwartz-Bruhat $\Psi$, so the
rule for $P_X$ in Definition~\ref{asm-theta-agree-doubling}(6)(b) gives
$l(\omega^\square(q)\Psi)=\chi_W(\det t(q))|\det t(q)|^{d/2}l(\Psi)$ for $q\in P_X$.
Moreover $l(\Phi^\square)=\int|\Phi(y)|^2\,dy>0$, as $\Phi$ is continuous and
nonzero. If $q\in P_X$ has $\omega^\square(q)\Phi^\square=a_q\Phi^\square$, it
follows that $a_q=\chi_W(\det t(q))|\det t(q)|^{d/2}$. For $q=n_0$, with
$t(n_0)=1$, this gives $a'=1$. So
$\omega^\square(sn_0s)\Phi^\square=a^2\Phi^\square=\Phi^\square$, and
\[
\omega^\square(p)\Phi^\square=\omega^\square(sn_0s)^{-1}\omega^\square(\operatorname{diag}(k,1))\Phi^\square=c\,\Phi^\square .
\]
As $p\in P_X$ with $t(p)=t$, $c=\chi_W(\det t)|\det t|^{d/2}$.

\emph{Step 6: $c=1$.} The matrices $t=1+b(1-k)$ and $t^{-1}=1-bf'(1-k)$ preserve
$L_v$, and $t-1$ is at $v$. So $\det t-1$ is at $v$, and in the basis of the
proof of Lemma~\ref{asm-theta-agree-local-vacuum-units} the component of $t$ at
$v$ and its inverse have integral entries; hence the component of $\det t$ at
every place of $E$ above $v$ is a unit, and $\det t$ is $1$ at the other places.
As $\chi_W$ is unramified at the places above $v$, $\chi_W(\det t)=1$. So
$c=|\det t|^{d/2}$ is a positive real number, and $c=1$ by Step 2. Hence
$\omega(k,1)$ acts at $v$ by $\mathrm{vac}_v\mapsto\mathrm{vac}_v$.
\end{proof}
% Lean: vacFixX_of StabEx (ConeShifts/PresOfHat.lean). Step 1 is eig (eigen, mem_latAtV, hP from
% pres_of_hat) and cc; Step 2 is bd (norm_eq_one, res_adm); Step 3 is ev (dbl_eigen with rd and
% bt) and e1 (the restriction rule hres, unitary_diag_one_mem_dbl); Step 4 is exists_beta,
% exists_decomp, pC, pCi, pe and pe'; Step 5 is rule_diag (lDiag, lDiag_ne, lDiag_schr) for a1
% and hpos, with aa; Step 6 is chi_det (at_det, exists_fr) and n4.

\begin{lemma}\label{asm-theta-agree-local-archimedean-orthogonal}
Let $B$ be a theta base of a CM field $E$, let $X$ be a Hermitian and $W$ a
skew-Hermitian lattice space, of some signatures and of ranks $n$ and $d$, and
let $H=(x,y,\iota,\mathfrak b)$ be an adapted plane in $X$ over a lattice space
$X_1$ of rank $n_1$ and $H'=(x',y',\iota',\mathfrak b')$ one in $W$ over a
lattice space $W_1$ of rank $d_1$ (Definition~\ref{asm-theta-agree-planes}(1)).
Put $h_X(a,a')=a^{\mathsf T}J_X\,c(a')$ for $a,a'\in\mathbf A_E^n$ and
$h_W(b,b')=b^{\mathsf T}J_W\,c(b')$ for $b,b'\in\mathbf A_E^d$, let $h$ and
$\langle\,,\rangle$ be the forms on $\mathbb W(X,W)$ of
Definition~\ref{asm-theta-agree-models}(3), and let $e_k$ be the unit vectors.
\begin{enumerate}
\item For $z\in\mathbf A_E^{n_1}$, $h_X(x,\iota z)=h_X(y,\iota z)=0$, and every
  $a\in\mathbf A_E^n$ is $sx+\iota z+ty$ with $s,t\in\mathbf A_E$ and
  $z\in\mathbf A_E^{n_1}$. The same holds for $H'$ in $W$.
\item For $a\in\mathbf A_E^n$, $b\in\mathbf A_E^d$, $\gamma\in\mathbf A_E$ and
  $w\in\mathbb W(X_1,W_1)$,
  \[
  h(\gamma\,ab^{\mathsf T},\iota w\iota'^{\mathsf T})=\gamma\sum_{k,l}c(w_{kl})\,h_X(a,\iota e_k)\,h_W(b,\iota'e_l).
  \]
  In particular $\langle\gamma\,ab^{\mathsf T},\iota w\iota'^{\mathsf T}\rangle=0$
  if $a\in\{x,y\}$ or $b\in\{x',y'\}$.
\item For every $u\in\mathbb W(X,W)$ at $\infty$ there is
  $z\in\mathbb W(X_1,W_1)$ at $\infty$ with
  $\langle u-\iota z\iota'^{\mathsf T},\iota w\iota'^{\mathsf T}\rangle=0$ for
  every $w\in\mathbb W(X_1,W_1)$.
\end{enumerate}
\end{lemma}
\leanref{Development/WeilBound/ThetaLifts/LeviDeterminants/MixedModels/PlaneAd.lean}{PositivityTheory.Mix.planeOK}
\leanref{Development/WeilBound/ThetaLifts/LeviDeterminants/MixedModels/PlaneGen.lean}{PositivityTheory.Mix.PlaneOK.xi}
\leanref{Development/WeilBound/ThetaLifts/LocalTheory/ArchimedeanBumps/SlicePol.lean}{PositivityTheory.ThKG.Loc.Bump.sy_vec}
\leanref{Development/WeilBound/ThetaLifts/LocalTheory/ArchimedeanBumps/SlicePol.lean}{PositivityTheory.ThKG.Loc.Bump.one_eq}
\leanref{Development/WeilBound/ThetaLifts/LocalTheory/ArchimedeanBumps/SlicePol.lean}{PositivityTheory.ThKG.Loc.Bump.decompS}
\begin{proof}
(1) As $c(J_X)^{\mathsf T}=J_X$, $h_X(a,a')=c(h_X(a',a))$ for all $a,a'$. For
$u\in E^{n_1}$, $h_X(\iota u,x)=h_X(\iota u,y)=0$ by the definition of an
adapted plane, and $h_X$ is $\mathbf A_E$-linear in its first argument; so
$h_X(\iota z,x)=h_X(\iota z,y)=0$ for $z=\sum_kz_ke_k\in\mathbf A_E^{n_1}$, and
$h_X(x,\iota z)=c(h_X(\iota z,x))=0$, and likewise for $y$. The lattice $L$ of
$X$ spans $E^n$, and every element of $L$ is $sx+\iota u+ty$ with $s,t\in E$ and
$u\in E^{n_1}$; so every element of $E^n$ has this form. For
$a=\sum_ja_je_j\in\mathbf A_E^n$, write $e_j=s_jx+\iota u_j+t_jy$; then
$a=(\sum_ja_js_j)x+\iota(\sum_ja_ju_j)+(\sum_ja_jt_j)y$. For $W$ the same
argument applies, with $c(J_W)^{\mathsf T}=-J_W$ and
$h_W(b,b')=-c(h_W(b',b))$.

(2) The definition of $h$ gives
$h(\gamma\,ab^{\mathsf T},\gamma'a'b'^{\mathsf T})=\gamma\,c(\gamma')\,h_X(a,a')\,h_W(b,b')$
for vectors $a,a',b,b'$ and $\gamma,\gamma'\in\mathbf A_E$, and $h$ is additive in
its second argument. As
$\iota w\iota'^{\mathsf T}=\sum_{k,l}w_{kl}(\iota e_k)(\iota'e_l)^{\mathsf T}$, this
is the formula. If $a\in\{x,y\}$, every $h_X(a,\iota e_k)$ is $0$ by (1); if
$b\in\{x',y'\}$, every $h_W(b,\iota'e_l)$ is. Then $h$ vanishes, and so does
$\langle\,,\rangle=-\operatorname{Tr}_{\mathbf A_E/\mathbf A}h$.

(3) By (1), applied to the unit vectors, there are $s,t\in\mathbf A_E^n$ and
$Z\in M_{n_1\times n}(\mathbf A_E)$ with
$1_n=\iota Z+xs^{\mathsf T}+yt^{\mathsf T}$, and likewise
$s',t'\in\mathbf A_E^d$ and $Z'$ with
$1_d=\iota'Z'+x's'^{\mathsf T}+y't'^{\mathsf T}$. Put $z=ZuZ'^{\mathsf T}$, which
is at $\infty$ as $u$ is. Multiplying $u$ on the left by the first identity, and
$\iota Zu$ on the right by the transpose of the second, gives
\[
u-\iota z\iota'^{\mathsf T}=x(u^{\mathsf T}s)^{\mathsf T}+y(u^{\mathsf T}t)^{\mathsf T}+(\iota Zus')x'^{\mathsf T}+(\iota Zut')y'^{\mathsf T}.
\]
Each term is $ab^{\mathsf T}$ with $a\in\{x,y\}$ or $b\in\{x',y'\}$, so (2) with
$\gamma=1$ gives the claim.
\end{proof}
% Lean: (1) is planeOK (PlaneAd.lean; the span from hH.split and the lattice spanning E^n, the
% orthogonality orth) with PlaneOK.xi and yi (PlaneGen.lean, from symm); (2) is sy_vec, with mid_sum
% and hTens_vec; (3) is decompS with one_eq. The Lean states (1) through PlaneOK, for planeOK_herm
% and planeOK_skew.

\begin{lemma}\label{asm-theta-agree-local-archimedean-gaussians}
Let $B$ be a theta base of a CM field $E$, let $X$ be a Hermitian and $W$ a
skew-Hermitian lattice space, of some signatures, and let $M$ be a model for
$(X,W)$. Then the Gaussians for $M$ (Definition~\ref{asm-theta-agree-models}(4))
span a subspace of $\mathcal S(F_\infty^{I_M})$ of dimension at most $1$.
\end{lemma}
\leanref{Development/WeilBound/ThetaLifts/LocalTheory/ArchimedeanBumps/SlicePol.lean}{PositivityTheory.ThKG.Loc.Bump.gauss_fin}
\begin{proof}
Let $(r,s)$ and $(r',s')$ be the signatures of $X$ and $W$. By
Lemma~\ref{asm-theta-agree-real-frames}, applied to $\tau_w(J_X)$ and
$i\tau_w(J_W)$ with their Cartan involutions, there are frames $A_w$ of
$\tau_w(J_X)$ of type $(r,s)$ and $C_w$ of $i\tau_w(J_W)$ of type $(r',s')$, with
$\theta_{X,w}A_w=A_wJ_{r,s}$ and $\theta_{W,w}C_w=C_wJ_{r',s'}$. By
Lemma~\ref{asm-theta-agree-real-fock-map}(1), with $a_0=1$ and $b_0=i$, there is
a Fock map $\Phi$ for $(M,(A_w),(C_w))$. For these data, the paragraph of the
proof of Lemma~\ref{asm-theta-agree-real-fock-map}(2) that shows
$\operatorname{Pol}(M)\subset\Phi(\mathcal T)$ proves that every Gaussian for $M$
is a multiple of $\Phi(\mathbf 1)$, where $\mathbf 1=1\otimes\cdots\otimes1$. So the
Gaussians span a subspace of $\mathbf C\,\Phi(\mathbf 1)$.
\end{proof}
% Lean: gauss_fin (SlicePol.lean): FrameEx for A and C, FockEx for Phi (Framed with a = 1 and
% b = I), and VacLine, whose hypothesis is the Gaussian condition rewritten with hA2 and hC2.

\begin{lemma}\label{asm-theta-agree-local-archimedean-slices}
Let $(M,M_1,H,H',e,u)$ be a cone setting for a theta base $B$
(Definition~\ref{asm-theta-agree-local-factor-setting}), with
$H=(x,y,\iota,\mathfrak b)$, $H'=(x',y',\iota',\mathfrak b')$ and Schrödinger
operators $\rho$ of $M$. For $s\in\mathbb W(X,W)$ at $\infty$ and
$f\in\mathcal S(F_\infty^{I_M})$ put $A_sf=e_\infty\bigl((\rho(s))_\infty f\bigr)$,
with the notation of Definition~\ref{asm-theta-agree-local-cones}(1). Let
$S_\infty$ be the set of infinite places of $F$, and for
$c\in\mathbf C^{S_\infty}$ let $\hat c\in\mathbf A_E$ be the element at $\infty$
with $\pi_w(\hat c)=c_w$ for every $w$. Then for every finite-dimensional
subspace $V$ of $\operatorname{Pol}(M)$ there is a finite-dimensional subspace
$V_1$ of $\operatorname{Pol}(M_1)$ with $A_{\hat cu(1)}g\in V_1$ for every $g\in V$
and every $c\in\mathbf C^{S_\infty}$.
\end{lemma}
\leanref{Development/WeilBound/ThetaLifts/LocalTheory/ArchimedeanBumps/SlicePol.lean}{PositivityTheory.ThKG.Loc.Bump.pol_fd}
\begin{proof}
Write $I=I_M$ and $I_1=I_{M_1}$, let $\rho_1$ be the Schrödinger operators of
$M_1$ and $e^M$ the isomorphism of $M$. Let $\psi_0$ and $\psi_1$ be the finite
parts of the vacua of $M$ and $M_1$, nonzero Schwartz-Bruhat functions on
$\mathbf A_f^I$ and $\mathbf A_f^{I_1}$ (proof of
Lemma~\ref{asm-theta-agree-local-bump-cone}), and choose $a$ with
$\psi_1(a)\neq0$. For $v\in\mathbb W(X,W)$ at $\infty$ let $D_v$ be the operator
of Lemma~\ref{asm-theta-agree-real-derivative} for $M$, and for
$z\in\mathbb W(X_1,W_1)$ at $\infty$ let $D^1_z$ be that for $M_1$. Write $(y,b)$
for the point of $\mathbf A^{I_1}$ with component $y$ at $\infty$ and finite
component $b$.

\emph{Step 1: five facts.} (a) Let $s$ be at $\infty$, and let
$S_s=S(e^M(s)_\infty^\flat)$, with the notation of
Lemma~\ref{asm-theta-agree-local-bump-slices-adelic}. A function on $\mathbf A^I$
that does not depend on the component at $\infty$ is $x\mapsto\phi(x_f)$ for a
function $\phi$ on $\mathbf A_f^I$, and as the finite component of $e^M(s)$ is
$0$, part (2) of that lemma gives $\rho(s)(f\otimes\phi)=S_sf\otimes\phi$. So
$\rho(s)$ acts at $\infty$ by $f\mapsto S_sf$, and by no other map on $f$, as
one sees on $f\otimes\psi_0$; hence $A_s=e_\infty\circ S_s$, a continuous linear
map. As $f\otimes\psi_0=f\otimes\bigotimes_v\mathrm{vac}_v$ and
$e_v(\mathrm{vac}_v)=\mathrm{vac}_{1,v}$, the factorization of $e$ gives
$e(\rho(s)(f\otimes\psi_0))=A_sf\otimes\psi_1$.

(b) For every continuous linear form $L$ on $\mathcal S(F_\infty^{I_1})$ there is
an admissible linear form $l$ on the functions on $\mathbf A^{I_1}$ with
$l(f\otimes\phi)=\phi(a)L(f)$ for $f\in\mathcal S(F_\infty^{I_1})$ and
Schwartz-Bruhat $\phi$ on $\mathbf A_f^{I_1}$: put $l(\Phi)=L(y\mapsto\Phi(y,a))$
on the Schwartz-Bruhat functions, as in Step 3 of the proof of
Lemma~\ref{asm-theta-agree-local-bump-cone}, and extend it linearly.

(c) Let $s$ and $v$ be at $\infty$, $f\in\mathcal S(F_\infty^I)$, and $l$ an
admissible linear form on the functions on $\mathbf A^{I_1}$. Then
$t\mapsto l\bigl(e(\rho(s)\rho(tv)(f\otimes\psi_0))\bigr)$ has derivative
$l(A_s(D_vf)\otimes\psi_1)$ at $t=0$. Indeed, $l'(\Phi)=l(e(\rho(s)\Phi))$ is
admissible: for Schwartz-Bruhat $\phi$ on $\mathbf A_f^I$,
$\rho(s)(f'\otimes\phi)=S_sf'\otimes\phi$ by (a), and Step 1 of the proof of
Lemma~\ref{asm-theta-agree-local-bump-slices} gives a Schwartz-Bruhat
$\epsilon(\phi)$ with $e(f''\otimes\phi)=e_\infty(f'')\otimes\epsilon(\phi)$ for
all $f''$; so $l'(f'\otimes\phi)=l(e_\infty(S_sf')\otimes\epsilon(\phi))$ is
continuous in $f'$. By Lemma~\ref{asm-theta-agree-real-derivative},
$t\mapsto l'(\rho(tv)(f\otimes\psi_0))$ has derivative
$l'(D_vf\otimes\psi_0)$ at $0$, which is $l(A_s(D_vf)\otimes\psi_1)$ by (a).

(d) For $v$ at $\infty$, $t\mapsto\rho(tv)$ acts at $\infty$ to first order by
$f\mapsto D_vf$, and by no other map on $f$, by
Lemmas~\ref{asm-theta-agree-real-derivative}
and~\ref{asm-theta-agree-real-weak}(1) and (3), as $\rho(tv)0=0$; likewise for
$M_1$. Hence $\operatorname{Pol}(M)$ is the smallest subspace of
$\mathcal S(F_\infty^I)$ that contains the Gaussians for $M$ and is stable under
every $D_v$; a Gaussian $f_0$ for $M$ with signs $\varepsilon_w$ has
$D_{v'}f_0=iD_vf_0$ whenever $v,v'$ are at $\infty$ with
$\pi_w(v')=\varepsilon_wi\,\theta_{X,w}\pi_w(v)\theta_{W,w}^{\mathsf T}$ for
every $w$; and every $D^1_z$ maps $\operatorname{Pol}(M_1)$ into itself.

(e) $D_{v+v'}=D_v+D_{v'}$ for $v,v'$ at $\infty$: $D_v$ is the operator
$D_{p,q}$ of Lemma~\ref{asm-theta-agree-density-derivative} with $(p,q)$ the
components at $\infty$ of $e^M(v)$, and $D_{p,q}$ is additive in $(p,q)$.

Let $\Sigma$ be the set of $s\in\mathbb W(X,W)$ at $\infty$ with
$\langle s,\iota w\iota'^{\mathsf T}\rangle=0$ for every
$w\in\mathbb W(X_1,W_1)$. It is closed under sums and under multiplication by
real numbers.

\emph{Step 2: the slices intertwine.} Let $s\in\Sigma$, let
$z\in\mathbb W(X_1,W_1)$ be at $\infty$ and $f\in\mathcal S(F_\infty^I)$; then
$A_s(D_{\iota z\iota'^{\mathsf T}}f)=D^1_z(A_sf)$. Let $L$ be a continuous linear
form on $\mathcal S(F_\infty^{I_1})$ and $l$ as in (b). As $\iota z\iota'^{\mathsf T}$
is at $\infty$, (c) shows that
$t\mapsto l\bigl(e(\rho(s)\rho(t\iota z\iota'^{\mathsf T})(f\otimes\psi_0))\bigr)$ has
derivative $\psi_1(a)L(A_s(D_{\iota z\iota'^{\mathsf T}}f))$ at $0$. As
$\langle s,t\iota z\iota'^{\mathsf T}\rangle=0$,
Lemma~\ref{asm-theta-agree-local-bump-slices-adelic}(1) gives
$\rho(s)\rho(t\iota z\iota'^{\mathsf T})=\rho(t\iota z\iota'^{\mathsf T})\rho(s)$, and
as $e$ is an evaluation map,
$e(\rho(\iota(tz)\iota'^{\mathsf T})\Phi)=\rho_1(tz)e(\Phi)$ for Schwartz-Bruhat
$\Phi$. With (a), the function is
$t\mapsto l(\rho_1(tz)(A_sf\otimes\psi_1))$, whose derivative at $0$ is
$\psi_1(a)L(D^1_z(A_sf))$, by Lemma~\ref{asm-theta-agree-real-derivative} for
$M_1$. As $\psi_1(a)\neq0$, $L$ takes the same value at the two functions. This
holds for every $L$, in particular for the evaluations at points.

\emph{Step 3: orthogonal directions.} Let $V_1$ be a finite-dimensional subspace
of $\mathcal S(F_\infty^{I_1})$ and $f$ a function with $A_sf\in V_1$ for every
$s\in\Sigma$. Then $A_s(D_vf)\in V_1$ for all $s,v\in\Sigma$. Let $L$ be a
continuous linear form that vanishes on $V_1$, and $l$ as in (b). For real $t$,
$s+tv\in\Sigma$, and the definition of $\rho$ gives
$\rho(s)\rho(tv)=\psi(\tfrac12\langle s,tv\rangle)\rho(s+tv)$, as in the proof of
Lemma~\ref{asm-theta-agree-local-factor-evaluation}(1). So by (a),
\[
l\bigl(e(\rho(s)\rho(tv)(f\otimes\psi_0))\bigr)=\psi(\tfrac12\langle s,tv\rangle)\,\psi_1(a)\,L(A_{s+tv}f)=0
\]
for every $t$. By (c) its derivative at $0$ is $\psi_1(a)L(A_s(D_vf))$, so
$L(A_s(D_vf))=0$. As $V_1$ is finite-dimensional, it is closed and is the
intersection of the kernels of the continuous linear forms that vanish on it,
by the Hahn-Banach theorem; so $A_s(D_vf)\in V_1$.

\emph{Step 4: the polynomial Fock space.} Let $N$ be the set of
$f\in\mathcal S(F_\infty^I)$ for which there is a finite-dimensional subspace
$V_1\subset\operatorname{Pol}(M_1)$ with $A_sf\in V_1$ for every $s\in\Sigma$. As
every $A_s$ is linear, $N$ is a subspace. We show
$\operatorname{Pol}(M)\subset N$; by (d) it suffices to show that $N$ contains
the Gaussians for $M$ and contains $D_vf$ whenever it contains $f$ and $v$ is at
$\infty$.

Let $V_G$ be the span of the Gaussians for $M_1$. It is finite-dimensional by
Lemma~\ref{asm-theta-agree-local-archimedean-gaussians}, as $X_1$ is Hermitian
and $W_1$ skew-Hermitian, and it lies in $\operatorname{Pol}(M_1)$. Let $f_0$ be a
Gaussian for $M$ with signs $\varepsilon_w$, and $s\in\Sigma$; we show that
$A_sf_0$ is a Gaussian for $M_1$ with the same signs. Let $z,z'$ be at $\infty$
with
$\pi_w(z')=\varepsilon_wi\,\theta_{X_1,w}\pi_w(z)\theta_{W_1,w}^{\mathsf T}$ for every
$w$, and suppose that $t\mapsto\rho_1(tz)$ and $t\mapsto\rho_1(tz')$ act at
$\infty$ to first order by $A_sf_0\mapsto g$ and $A_sf_0\mapsto g'$. By (d) and
Step 2, $g=D^1_z(A_sf_0)=A_s(D_{\iota z\iota'^{\mathsf T}}f_0)$ and
$g'=A_s(D_{\iota z'\iota'^{\mathsf T}}f_0)$. As $\iota$ and $\iota'$ have entries in
$E$, $\pi_w(\iota z\iota'^{\mathsf T})=\tau_w(\iota)\pi_w(z)\tau_w(\iota')^{\mathsf T}$,
and as $\theta_{X,w}\tau_w(\iota)=\tau_w(\iota)\theta_{X_1,w}$ and
$\theta_{W,w}\tau_w(\iota')=\tau_w(\iota')\theta_{W_1,w}$ (adapted planes),
\[
\pi_w(\iota z'\iota'^{\mathsf T})=\varepsilon_wi\,\theta_{X,w}\,\pi_w(\iota z\iota'^{\mathsf T})\,\theta_{W,w}^{\mathsf T}.
\]
So $D_{\iota z'\iota'^{\mathsf T}}f_0=iD_{\iota z\iota'^{\mathsf T}}f_0$ by (d), and
$g'=ig$. Hence $A_sf_0\in V_G$ for every $s\in\Sigma$, and $f_0\in N$.

Let $f\in N$, with a space $V_1$, and let $v$ be at $\infty$. By
Lemma~\ref{asm-theta-agree-local-archimedean-orthogonal}(3) there is $z$ at
$\infty$ with $v_0=v-\iota z\iota'^{\mathsf T}\in\Sigma$. By (e),
$D_vf=D_{\iota z\iota'^{\mathsf T}}f+D_{v_0}f$. For $s\in\Sigma$, Steps 2 and 3
give
\[
A_s(D_vf)=D^1_z(A_sf)+A_s(D_{v_0}f)\in D^1_z(V_1)+V_1 .
\]
The space $D^1_z(V_1)+V_1$ is finite-dimensional and lies in
$\operatorname{Pol}(M_1)$ by (d). So $D_vf\in N$.

\emph{Step 5: conclusion.} Let $c\in\mathbf C^{S_\infty}$. The element
$\hat cu(1)$ is at $\infty$, as $\hat c$ is. If $u=u_\alpha$, then
$u(1)=xy'^{\mathsf T}$, and if $u=u_\beta$, then $u(1)=yx'^{\mathsf T}$; in both cases
Lemma~\ref{asm-theta-agree-local-archimedean-orthogonal}(2), with $\gamma=\hat c$,
gives $\hat cu(1)\in\Sigma$. Let $g_1,\ldots,g_r$ be a basis of $V$. By Step 4
each $g_i$ lies in $N$, with a space $V_{1,i}$; put $V_1=\sum_iV_{1,i}$. As
$A_{\hat cu(1)}$ is linear, $A_{\hat cu(1)}g\in V_1$ for every $g\in V$.
\end{proof}
% Lean: pol_fd (SlicePol.lean). Step 1 is (a) ev_op (ArchHeis, ArchUniq; SliceLine.lean) with
% hval, hadd and hsm, (b) exists_admissible_functional_adelic (SliceMain.lean), (c) hder with
% EvAdm (evAdm_of in SmallCones/NormLeSum.lean: schr_pure and ev_pure) and DerEx, (d) uq, uq1 and
% polD through ActsLaws (actsD_of), (e) the computation hs. The set Sigma is S. Step 2 is hmid
% (rho_comm, IsEv.Heisenberg2, eq_of_weak); Step 3 is hort (rho_rho, mem_of_weak); Step 4 is the
% submodule N and hN, with gauss_fin (Lemma asm-theta-agree-local-archimedean-gaussians), the
% Cartan conditions hp.cartan and hp'.cartan with map_inc, and decompS; Step 5 is hS (at_smul,
% sy_vec with oX.xi or oX.yi) and the span induction.

\begin{lemma}\label{asm-theta-agree-local-factor-archimedean}
Let $(M,M_1,H,H',e,u)$ be a cone setting for a theta base $B$, $K\subset U_X$ the
lattice stabilizer of $X$, $m$ a normalized local measure system for $K$,
$(X,W,M,\chi_X,\chi_W,\omega)$ a Weil pair, $\nu$ a continuous character of
$\mathbf A_E^\times$ and $z\geq1$ an integer. Then there is an archimedean cone
$L$ for $(M,\omega,m,e_\infty,u,\nu,z)$
(Definition~\ref{asm-theta-agree-local-cones}(3)) with
$L(\operatorname{Pol}(M))\subset\operatorname{Pol}(M_1)$, such that for every
$f\in\mathcal S(F_\infty^{I_M})$ and every $y$ the integrand defining $L(f)(y)$ is
$m_{A,\infty}$-integrable in $l$ for every $k\in K_\infty$, and its integral over
$A_\infty$ is $m_{K,\infty}$-integrable in $k$.
\end{lemma}
\leanref{Development/WeilBound/ThetaLifts/LocalTheory/ConeShifts/ArchCone.lean}{PositivityTheory.ThKG.Loc.Cone.ArchCone}
\leanref{Development/WeilBound/ThetaLifts/LocalTheory/ArchimedeanBumps/SlicePol.lean}{PositivityTheory.ThKG.Loc.Bump.archCone_of}
\begin{proof}
Let $T(k)$, $k\in K_\infty$, and $A(c)$, $c\in\mathbf C^{S_\infty}$, be the
continuous linear maps of Steps 3 and 2 of the proof of
Lemma~\ref{asm-theta-agree-local-bump-cone}, for the same data; that lemma's
hypotheses are those here. The construction there does not use an archimedean
cone.

\emph{Step 1: the map $L$.} Step 5 of that proof, for $\chi=\nu$ and $z$, gives a
continuous linear map $L$ from $\mathcal S(F_\infty^{I_M})$ to
$\mathcal S(F_\infty^{I_{M_1}})$ such that for every continuous linear form $T'$
the function
$(k,t)\mapsto T'\bigl(e_\infty((\rho(1_\infty u(t))\,\omega(k,1))_\infty f)\bigr)\nu(t)|t|^z$
is $m_{K,\infty}\times m_{A,\infty}$-integrable, with integral $T'(L(f))$. Taking
for $T'$ the evaluation at $y$, Fubini's theorem gives the two integrability
claims and the formula of Definition~\ref{asm-theta-agree-local-cones}(3), so $L$
is an archimedean cone for $(M,\omega,m,e_\infty,u,\nu,z)$. By Step 4 of that
proof the integrand is $T'(A(c(t))T(k)f)\,\nu(t)|t|^z$, with $c(t)$ as in its
Step 5.

\emph{Step 2: a finite-dimensional space.} Let $f\in\operatorname{Pol}(M)$. By
Step 1 of that proof $K_\infty\subset K_\infty(X)$, so $T(k)$ maps
$\operatorname{Pol}(M)$ into itself
(Lemma~\ref{asm-theta-agree-rechoice-fock}, with $h=1$). Let $\psi_0$ be the
finite part of the vacua of $M$, a nonzero Schwartz-Bruhat function on
$\mathbf A_f^{I_M}$, so that $f\otimes\psi_0\in\operatorname{PD}(M)$. By
Lemma~\ref{asm-theta-agree-split-inputs}(2) the functions
$\omega(k,1)(f\otimes\psi_0)=T(k)f\otimes\psi_0$, $k\in K_\infty$, span a
finite-dimensional space. As $g\mapsto g\otimes\psi_0$ is injective, the span $V$
of the $T(k)f$ is finite-dimensional, and $V\subset\operatorname{Pol}(M)$.

\emph{Step 3: the slices.} By Step 2 of the proof of
Lemma~\ref{asm-theta-agree-local-bump-cone},
$A(c)g=e_\infty\bigl((\rho(\hat c\,u(1)))_\infty g\bigr)$, with $\hat c$ as there.
So Lemma~\ref{asm-theta-agree-local-archimedean-slices}, applied to the space $V$
of Step 2, gives a finite-dimensional subspace $V_1$ of
$\operatorname{Pol}(M_1)$ with $A(c)g\in V_1$ for every $g\in V$ and every
$c\in\mathbf C^{S_\infty}$.
% Lean: Bump.pol_fd (ArchimedeanBumps/SlicePol.lean), applied to V = span of the T k f with
% hTp; Lemma asm-theta-agree-local-archimedean-slices.

\emph{Step 4: $L(f)\in\operatorname{Pol}(M_1)$.} Let $V_1$ be as in Step 3 for the
space $V$ of Step 2. Every value $A(c(t))T(k)f$ lies in $V_1$, so by Step 1,
$l(L(f))=0$ for every continuous linear form $l$ that vanishes on $V_1$. As $V_1$
is finite-dimensional, it is closed and is the intersection of the kernels of
these forms, by the Hahn-Banach theorem. So $L(f)\in V_1\subset\operatorname{Pol}(M_1)$.
\end{proof}
% Lean: Bump.archCone_of ArchDict StabEx ArchOp Slice ArchHeis ArchUniq SchInt ActsLaws DerEx
% VacLine FockEx FrameEx EvAdm KFinData (ArchimedeanBumps/SlicePol.lean). Step 1 is cone_core (the
% same map as in archSlice_of) with evAt; Step 2 is k_ops (the clause hTp of ArchOp for K_inf),
% vac_fin, KFinData and the injective map j; Step 4 is op_comp, shift_eq and mem_of_weak.

\begin{lemma}\label{asm-theta-agree-local-pure-schrodinger}
Let $F$ be a number field, $I$ a finite set, and $\psi$ a continuous character of
$\mathbf A$, trivial on $F$, with $\psi(x)=\exp(-2\pi i\operatorname{Tr}(x))$ on
$F_\infty$, and Schrödinger operators $\rho(w)$ as in
Theorem~\ref{asm-theta-agree-evaluation-heisenberg}. Let
$w=(w_1,w_2)\in\mathbf A^I\times\mathbf A^I$. For $f$ on $F_\infty^I$ and for
$\phi$ on $F_v^I$, $v$ finite, put
\[
(\rho_\infty(w)f)(x)=\psi_\infty\bigl(x\cdot w_{2,\infty}+\tfrac12w_{1,\infty}\cdot w_{2,\infty}\bigr)f(x+w_{1,\infty}),\qquad
(\rho_v(w)\phi)(x)=\psi_v\bigl(x\cdot w_{2,v}+\tfrac12w_{1,v}\cdot w_{2,v}\bigr)\phi(x+w_{1,v}).
\]
Let $f\in\mathcal S(F_\infty^I)$ and let $\phi_v$ be Schwartz-Bruhat functions on
$F_v^I$, equal to the indicator function of $\mathcal O_v^I$ for almost all $v$.
Then $\rho_\infty(w)f\in\mathcal S(F_\infty^I)$; $\rho_v(w)\phi_v$ is
Schwartz-Bruhat for every $v$ and is the indicator function of $\mathcal O_v^I$
for almost all $v$; and
\[
\rho(w)\Bigl(f\otimes\bigotimes_v\phi_v\Bigr)=\rho_\infty(w)f\otimes\bigotimes_v\rho_v(w)\phi_v .
\]
\end{lemma}
\leanref{Development/WeilBound/ThetaLifts/LeviDeterminants/MixedModels/PureSt.lean}{PositivityTheory.Mix.PureSt}
\leanref{Development/WeilBound/ThetaLifts/LeviDeterminants/MixedModels/MCPure.lean}{PositivityTheory.Mix.pure_tensor_estimates_holds}
\begin{proof}
The function $\rho_\infty(w)f$ is a translate of $f$ times a smooth character of
$F_\infty^I$ of absolute value $1$ whose derivatives are bounded, so it is a
Schwartz function; and $\rho_v(w)\phi_v$ is locally constant with compact support,
as $\psi_v$ is locally constant. As $\psi=\prod_v\psi_v$, the value of
$\rho(w)(f\otimes\bigotimes_v\phi_v)$ at $x$ is the product of the values of the
factors at the components of $x$. For almost all $v$, $\phi_v$ is the indicator
function of $\mathcal O_v^I$, the components $w_{1,v}$ and $w_{2,v}$ lie in
$\mathcal O_v^I$, $\tfrac12\in\mathcal O_v$, and $\psi_v$ is trivial on
$\mathcal O_v$. For such $v$ and $x\in\mathcal O_v^I$, $x+w_{1,v}\in\mathcal O_v^I$ and
the argument of $\psi_v$ lies in $\mathcal O_v$, so $(\rho_v(w)\phi_v)(x)=1$; for
$x\notin\mathcal O_v^I$, $x+w_{1,v}\notin\mathcal O_v^I$ and the value is $0$.
\end{proof}
% Lean: pure_tensor_estimates_holds (MCPure.lean): exists_isTr for the archimedean factor, the
% filter argument with psiV_int for almost all v, and the factorization of psi (e1, e2).

\begin{lemma}\label{asm-theta-agree-local-pure-partial}
Let $(X,W,M,\chi_X,\chi_W,\omega)$ be a Weil pair for a theta base $B$, with
$I=I_M$ and vacua $\mathrm{vac}_v$, let $K\subset U_X$ be the lattice stabilizer
of $X$, and let $T$ be a finite set of finite places such that $\chi_X$ and
$\chi_W$ are unramified at the places of $E$ not above $T$. Let
$f,f_1\in\mathcal S(F_\infty^I)$, let $\phi_v$ be Schwartz-Bruhat functions on
$F_v^I$ with $\phi_v=\mathrm{vac}_v$ for $v\notin T$, and let $k\in K$, with
$k^{(v)}=k^{(1_v)}$ as in Lemma~\ref{asm-theta-agree-local-stabilizer-parts}.
Suppose that $\omega(k^{(\infty)},1)$ acts at $\infty$ by $f\mapsto f_1$. Then
there are Schwartz-Bruhat functions $\phi'_v$ on $F_v^I$ such that
$\omega(k^{(v)},1)$ acts at $v$ by $\phi_v\mapsto\phi'_v$, $\phi'_v=\mathrm{vac}_v$
for $v\notin T$, and
$\omega(k,1)(f\otimes\bigotimes_v\phi_v)=f_1\otimes\bigotimes_v\phi'_v$.
\end{lemma}
\leanref{Development/WeilBound/ThetaLifts/LocalTheory/ConeShifts/ActsRhoFin.lean}{PositivityTheory.ThKG.Loc.Cone.om_pure}
\begin{proof}
Let $(\omega_v,\rho_v)$ be local operators at $v$ for the pair
(Lemma~\ref{asm-theta-agree-finite-inputs}(2)), and put
$\phi'_v=\omega_v(k^{(v)},1)\phi_v$, a Schwartz-Bruhat function; $k^{(v)}$ is at
$v$, so $\omega(k^{(v)},1)$ acts at $v$ by $\phi_v\mapsto\phi'_v$. If $v\notin T$,
then $k^{(v)}\in K_v$
(Lemma~\ref{asm-theta-agree-local-stabilizer-parts}(3)), so $\omega(k^{(v)},1)$
also acts at $v$ by $\mathrm{vac}_v\mapsto\mathrm{vac}_v$
(Lemma~\ref{asm-theta-agree-local-factor-vacuum}), and
$\phi'_v=\mathrm{vac}_v$ by Lemma~\ref{asm-theta-agree-finite-inputs-schur}(1).
As $M$ is regular, $\mathrm{vac}_v$ is the indicator function of $\mathcal O_v^I$
for almost all $v$, so all the pure tensors below are Schwartz-Bruhat. Put
$\Phi=f\otimes\bigotimes_v\phi_v$.

\emph{Step 1: finitely many places.} Let $a_S$ be as in
Lemma~\ref{asm-theta-agree-local-stabilizer-parts}(4), and let $\phi^S_v$ be
$\phi'_v$ for $v\in S$ and $\phi_v$ for $v\notin S$. We show
$\omega(k^{(a_S)},1)\Phi=f_1\otimes\bigotimes_v\phi^S_v$ by induction on $S$. For
$S=\emptyset$, $a_S=1_\infty$, and this is the action at $\infty$, with
$r=\bigotimes_v\phi_v$, which does not depend on the infinite component. If
$v\notin S$, then $1_va_S=0$, so
$k^{(a_{S\cup\{v\}})}=k^{(v)}k^{(a_S)}$
(Lemma~\ref{asm-theta-agree-local-stabilizer-parts}(1)), and
$\omega(k^{(v)},1)$ maps $f_1\otimes\bigotimes_u\phi^S_u$, whose factor at $v$ is
$\phi_v$, to the same tensor with $\phi'_v$ at $v$, by the action at $v$.

\emph{Step 2: the limit.} Let $x\in\mathbf A^I$. The linear form
$\Psi\mapsto\Psi(x)$ is admissible, so $g\mapsto(\omega(g,1)\Phi)(x)$ is continuous
on $U_X$ (condition (c) of Definition~\ref{asm-theta-agree-models}(6)). For
$S\supset T$, $\phi^S_v=\phi'_v$ for every $v$, as $\phi'_v=\mathrm{vac}_v=\phi_v$
for $v\notin T$. So $(\omega(k^{(a_S)},1)\Phi)(x)=(f_1\otimes\bigotimes_v\phi'_v)(x)$
for $S\supset T$, and as $k^{(a_S)}\to k$
(Lemma~\ref{asm-theta-agree-local-stabilizer-parts}(4)),
$(\omega(k,1)\Phi)(x)=(f_1\otimes\bigotimes_v\phi'_v)(x)$.
\end{proof}
% Lean: om_pure (ActsRhoFin.lean) with ActsEx for phi'_v, VacFixX (Lemma
% asm-theta-agree-local-factor-vacuum) and partU_mem_kAt off T, step (induction with
% pure_real, partU_add and ThKF.pureAct), and the limit through proj_adm, the continuity field of
% the Weil pair and tendsto_partU.

\begin{lemma}\label{asm-theta-agree-local-factor-pure}
Let $(M,M_1,H,H',e,u)$ be a cone setting for a theta base $B$, $K\subset U_X$
the lattice stabilizer of $X$, $(X,W,M,\chi_X,\chi_W,\omega)$ a Weil pair, and
$T$ a finite set of finite places such that $\chi_X$ and $\chi_W$ are unramified
at the places of $E$ not above $T$. Let $f\in\mathcal S(F_\infty^{I_M})$, let
$\phi_v$ be Schwartz-Bruhat functions on $F_v^{I_M}$ with $\phi_v=\mathrm{vac}_v$
for $v\notin T$, and let $k\in K$ and $l\in\mathbf A_E^\times$, with $k^{(v)}$ and
$l^{(v)}$ as in Lemmas~\ref{asm-theta-agree-local-factor-stabilizer}
and~\ref{asm-theta-agree-local-factor-ideles}. Let $T^v_{k,l}$ be the operators
of Definition~\ref{asm-theta-agree-local-cones}(3), for $u$. Then
$e\bigl(\rho(u(l))\,\omega(k,1)\,(f\otimes\bigotimes_v\phi_v)\bigr)$ is $0$ if
$1_vu(l)\notin\Gamma_v$ for some finite $v\notin T$, and otherwise it is
\[
e_\infty\bigl((T^\infty_{k^{(\infty)},l^{(\infty)}})_\infty f\bigr)\otimes
\bigotimes_{v\in T}e_v\bigl((T^v_{k^{(v)},l^{(v)}})_v\phi_v\bigr)\otimes
\bigotimes_{v\notin T}\mathrm{vac}_{1,v}.
\]
\end{lemma}
\leanref{Development/WeilBound/ThetaLifts/LocalTheory/ConeShifts/ArchCone.lean}{PositivityTheory.ThKG.Loc.Cone.PureFac}
\leanref{Development/WeilBound/ThetaLifts/LocalTheory/ConeShifts/ActsRhoFin.lean}{PositivityTheory.ThKG.Loc.Cone.pureFac_of}
\begin{proof}
Let $e^M$ be the isomorphism of $M$, $\rho$ its Schrödinger operators and
$I=I_M$, and put $w=e^M(u(l))$ and $\Phi=f\otimes\bigotimes_v\phi_v$. By
Lemma~\ref{asm-theta-agree-rechoice-operators}(1) there is $f_1$ such that
$\omega(k^{(\infty)},1)$ acts at $\infty$ by $f\mapsto f_1$. Let $\phi'_v$ be as
in Lemma~\ref{asm-theta-agree-local-pure-partial}, so that
$\omega(k,1)\Phi=f_1\otimes\bigotimes_v\phi'_v$, with $\phi'_v=\mathrm{vac}_v$ for
$v\notin T$. With $\rho_\infty(w)$ and $\rho_v(w)$ as in
Lemma~\ref{asm-theta-agree-local-pure-schrodinger}, put $g=\rho_\infty(w)f_1$ and
$\phi''_v=\rho_v(w)\phi'_v$. By that lemma,
\[
\rho(u(l))\,\omega(k,1)\Phi=g\otimes\bigotimes_v\phi''_v,
\]
with $\phi''_v$ Schwartz-Bruhat and equal to the indicator function of
$\mathcal O_v^I$, hence to $\mathrm{vac}_v$, for almost all $v$. By the components
of the evaluation map $e$ (Definition~\ref{asm-theta-agree-planes}(5)),
\[
e\bigl(\rho(u(l))\,\omega(k,1)\Phi\bigr)=e_\infty(g)\otimes\bigotimes_ve_v(\phi''_v).
\]

\emph{The finite factors.} Let $v$ be a finite place. The $v$-components of
$e^M(1_vu(l))=1_vw$ are those of $w$, so $\phi''_v=\rho_v(1_vu(l))\phi'_v$ with the
operator of Lemma~\ref{asm-theta-agree-local-factor-evaluation}, and $\rho(1_vu(l))$
acts at $v$ by $\phi'_v\mapsto\phi''_v$ by part (1) of that lemma. If $v\notin T$,
then $\phi'_v=\mathrm{vac}_v$, and
Lemma~\ref{asm-theta-agree-local-factor-finite}(1) gives
$e_v(\phi''_v)=\mathrm{vac}_{1,v}$ if $1_vu(l)\in\Gamma_v$ and $0$ otherwise. In
particular the whole tensor is $0$ if $1_vu(l)\notin\Gamma_v$ for some finite
$v\notin T$. If $v\in T$: by
Lemma~\ref{asm-theta-agree-local-factor-directions}(1), $u(l)=\sigma(l)V_0$, and
$1_vl^{(v)}=1_vl$ with $\sigma(1_v)=1_v$, so $1_vu(l^{(v)})=1_vu(l)$. So
$T^v_{k^{(v)},l^{(v)}}=\rho(1_vu(l))\,\omega(k^{(v)},1)$. As $\omega(k^{(v)},1)$ maps
Schwartz-Bruhat functions to Schwartz-Bruhat functions and acts at $v$ by
$\phi_v\mapsto\phi'_v$, this operator acts at $v$ by $\phi_v\mapsto\phi''_v$; so
$(T^v_{k^{(v)},l^{(v)}})_v\phi_v=\phi''_v$, by
Lemma~\ref{asm-theta-agree-finite-inputs-schur}(1).

\emph{The infinite factor.} In the same way
$T^\infty_{k^{(\infty)},l^{(\infty)}}=\rho(1_\infty u(l))\,\omega(k^{(\infty)},1)$.
For a function $r$ that does not depend on the infinite component, the formula
for $\rho$ and $\psi(1_\infty s)=\psi_\infty(s_\infty)$ give
$\rho(1_\infty u(l))(f_1\otimes r)=g\otimes r$. So
$T^\infty_{k^{(\infty)},l^{(\infty)}}$ acts at $\infty$ by $f\mapsto g$, and
$(T^\infty_{k^{(\infty)},l^{(\infty)}})_\infty f=g$ by the uniqueness of
Lemma~\ref{asm-theta-agree-real-weak}(3). Putting the factors together gives the
claim.
\end{proof}
% Lean: pureFac_of ActsEx ArchOp ArchUniq VacFixX VacShift PureSt (ConeShifts/ActsRhoFin.lean):
% ArchOp gives f1 (Ti), om_pure the first display, PureSt (pure_tensor_estimates_holds) the
% second, IsEvLoc.pure the third; off is the case v not in T (VacShift = vacShift_of, the first
% part of Lemma asm-theta-agree-local-factor-finite), part_shift the identity 1_v u(l^(v)) =
% 1_v u(l), acts_comp, acts_rho_fin, acts_rho_inf, opF_eq and opI_eq (with ArchUniq) the
% identification of the factors, and pure_zero the vanishing.

\begin{lemma}\label{asm-theta-agree-local-factor-span}
Let $F$ be a number field, $I$ a finite set, and for every finite place $v$ let
$\mathrm{vac}_v$ be a Schwartz-Bruhat function on $F_v^I$, equal to the indicator
function of $\mathcal O_v^I$ for almost all $v$. Then the Schwartz-Bruhat
functions on $\mathbf A^I$ are the finite sums of functions
$f\otimes\bigotimes_v\phi_v$ with $f\in\mathcal S(F_\infty^I)$, every $\phi_v$
Schwartz-Bruhat and $\phi_v=\mathrm{vac}_v$ for $v$ outside a finite set.
\end{lemma}
\leanref{Development/WeilBound/ThetaLifts/LocalTheory/ConeShifts/ArchCone.lean}{PositivityTheory.ThKG.Loc.Cone.SchSpan}
\leanref{Development/WeilBound/ThetaLifts/LocalTheory/ConeShifts/FinSpan.lean}{PositivityTheory.ThKG.Loc.Cone.Fin.schSpan}
\begin{proof}
Such a function is Schwartz-Bruhat: the product $\bigotimes_v\phi_v$ is locally
constant with compact support on $\mathbf A_f^I$, as almost all factors are the
indicator function of $\mathcal O_v^I$. Conversely, it suffices to write
$f\otimes g$, with $g$ Schwartz-Bruhat on $\mathbf A_f^I$, as such a sum. Let
$\hat{\mathcal O}=\prod_v\mathcal O_v$. The function $g$ is invariant under
translation by an open subgroup of $\mathbf A_f^I$, which contains
$a\hat{\mathcal O}^I$ for some $a\in F^\times$, and its support is compact, so it
meets finitely many cosets $r+a\hat{\mathcal O}^I$, $r\in R$. Hence
$g=\sum_{r\in R}g(r)\mathbf 1_{r+a\hat{\mathcal O}^I}$. Each indicator function is
$\bigotimes_v\mathbf 1_{r_v+a\mathcal O_v^I}$, and for almost all $v$ both
$a\in\mathcal O_v^\times$ and $r_v\in\mathcal O_v^I$, so that the factor is the
indicator function of $\mathcal O_v^I$, which is $\mathrm{vac}_v$ for almost all
$v$. So $f\otimes g=\sum_{r\in R}g(r)\,f\otimes\bigotimes_v\mathbf 1_{r_v+a\mathcal
O_v^I}$, with factors $\mathrm{vac}_v$ outside a finite set.
\end{proof}
% Lean: schSpan (ConeShifts/FinSpan.lean): exists_invariant and
% exists_scaled_integral_lattice_subset (the subgroup a O^I), exists_sum_cosets,
% indicator_coset_scaledInt_eq_finprod and eventually_cosetBoxV_eq (the product), and
% TensUniq.pure_mem_schA for the converse.

\begin{lemma}\label{asm-theta-agree-local-factor}
Let $B$ be a theta base of a CM field $E$.
\begin{enumerate}
\item There are signs $e_1,e_2\in\{\pm1\}$ such that the following holds. Let
  $G$ and $G_1$ be triple data of ranks $(b+2,a')$ and $(b,a')$
  (Definition~\ref{asm-theta-agree-triple}(2)) with $V'_1=V'$, $q_1=q$ and
  $p_1+1=p$, let $H$ be an adapted plane in $V$ over $V_1$ and $H'$ one in $W$
  over $W_1$, both at positions $(0,b+1)$
  (Definition~\ref{asm-theta-agree-planes}(1)); that is, $(G,G_1,H,H')$ is a
  plane-removal configuration (Definition~\ref{asm-theta-agree-end-notions}(7)),
  and for $q=0$ an $\alpha$-configuration
  (Definition~\ref{asm-theta-agree-planes}(4)). Let
  $e$ be an evaluation map for $(M,M_1,\iota_H,\iota_{H'},\mathcal J_\alpha(H,H'))$
  with components $e_\infty$ and $e_v$
  (Definition~\ref{asm-theta-agree-planes}(5)), $K\subset U_V$ the lattice
  stabilizer of $V$, $\mu_K$ and $\mu$ Haar measures on $K$ and
  $\mathbf A_E^\times$, and $m$ a normalized local measure system for $K$
  (Definition~\ref{asm-theta-agree-local-cones}). Then there is $\kappa\neq0$
  such that for all triples $(B,G,\mathcal X,\omega,\omega')$ and
  $(B,G_1,\mathcal X,\omega_1,\omega'_1)$ there are a continuous linear map
  $L$, linear maps $\Lambda_v$ and numbers $u_v$, for the finite places $v$,
  with the following properties. Let $\operatorname{cone}_\alpha$ be the cone
  for $(M,\omega,e,K,\mu_K,\mu,u_\alpha,\nu)$ and $z=k_\alpha$, let
  $\mathrm{vac}_{1,v}$ be the vacua of $M_1$, and put
  $\lambda_0=\kappa\,\nu(\mathfrak b_H)^{e_1}\nu(\mathfrak b_{H'})^{e_2}L^{S_0}(\nu,k_\alpha)$.
  \begin{enumerate}
  \item[(a)] $L$ is an archimedean cone for $(M,\omega,m,e_\infty,u_\alpha,\nu,k_\alpha)$,
    and $L(\operatorname{Pol}(M))\subset\operatorname{Pol}(M_1)$.
  \item[(b)] For every $v$ and $\phi\in\mathcal S(F_v^{I_M})$, $\Lambda_v\phi$
    is the local cone of $\phi$ for $(M,\omega,m,e_v,u_\alpha,\nu,k_\alpha)$,
    and it is Schwartz-Bruhat.
  \item[(c)] For $v\notin T$, $u_v\neq0$ and $\Lambda_v(\mathrm{vac}_v)=u_v\,\mathrm{vac}_{1,v}$.
  \item[(d)] $\operatorname{cone}_\alpha(\Phi+c\Psi)=\operatorname{cone}_\alpha(\Phi)+c\operatorname{cone}_\alpha(\Psi)$
    for Schwartz-Bruhat $\Phi$ and $\Psi$ and $c\in\mathbf C$.
  \item[(e)] Let $T'\supset T$ be a finite set of finite places,
    $f\in\mathcal S(F_\infty^{I_M})$, and $\phi_v\in\mathcal S(F_v^{I_M})$
    with $\phi_v=\mathrm{vac}_v$ for $v\notin T'$. Put $\psi_v=\Lambda_v\phi_v$
    for $v\in T'$ and $\psi_v=\mathrm{vac}_{1,v}$ for $v\notin T'$. Then
    \[
    \Bigl(\prod_{v\in T'\setminus T}u_v\Bigr)\operatorname{cone}_\alpha\Bigl(f\otimes\bigotimes_v\phi_v\Bigr)
    =\lambda_0\,L(f)\otimes\bigotimes_v\psi_v.
    \]
  \end{enumerate}
\item There are signs $e_1,e_2$ such that the same holds for a
  $\beta$-configuration $(G,G_1,H,H')$ with $G$ of ranks $(a,b+2)$, an
  evaluation map $e'$ for $(M',M'_1,\iota_H,\iota_{H'},\mathcal J_\beta(H,H'))$
  with components $e'_\infty$ and $e'_v$, and the lattice stabilizer
  $K\subset U_{V'}$ of $V'$, with $M'$, $\omega'$, $e'$, $M'_1$, $u_\beta$,
  $k_\beta+1$ and the vacua $\mathrm{vac}'_v$ and $\mathrm{vac}'_{1,v}$ of $M'$
  and $M'_1$ in place of $M$, $\omega$, $e$, $M_1$, $u_\alpha$, $k_\alpha$,
  $\mathrm{vac}_v$ and $\mathrm{vac}_{1,v}$; the cone is then
  $\operatorname{cone}_\beta$, and
  $\lambda_0=\kappa\,\nu(\mathfrak b_H)^{e_1}\nu(\mathfrak b_{H'})^{e_2}L^{S_0}(\nu,k_\beta+1)$.
\end{enumerate}
\end{lemma}
\leanref{Development/WeilBound/ThetaLifts/EigenvalueEndgame/EigenKilling/ConeIntegralFactorisation.lean}{PositivityTheory.ThKEnd.Comp.ConeIntegralFactorisation}
\leanref{Development/WeilBound/ThetaLifts/LocalTheory/ConeShifts/ConeFactorisationOfCodimensionTwo.lean}{PositivityTheory.ThKG.Loc.Cone.coneIntegralFactorisation_of}
\leanref{Development/WeilBound/ThetaLifts/LocalTheory/UnramifiedDeterminedBySpherical.lean}{PositivityTheory.ThKG.Loc.ConeFacA}
\leanref{Development/WeilBound/ThetaLifts/LocalTheory/UnramifiedDeterminedBySpherical.lean}{PositivityTheory.ThKG.Loc.ConeFacB}
\leanref{Development/WeilBound/ThetaLifts/LocalTheory/ConeShifts/VacShiftOf.lean}{PositivityTheory.ThKG.Loc.Cone.coneFac_all}
\begin{proof}
The argument below proves (1). It proves (2) with the Weil pair
$(V',W,M',\alpha,\alpha/\beta,\omega')$, the map $u_\beta$, $z=k_\beta+1$ and the
signs $(b_1,b_2)$ of Lemma~\ref{asm-theta-agree-local-factor-ideal}(2) in place
of $(V,W,M,\beta,\beta/\alpha,\omega)$, $u_\alpha$, $z=k_\alpha$ and $(a_1,a_2)$.
Let $(e_1,e_2)=(a_1,a_2)$, and fix the data of (1) that precede $\kappa$. If
there are no triples $(B,G,\mathcal X,\omega,\omega')$ and
$(B,G_1,\mathcal X,\omega_1,\omega'_1)$, every $\kappa$ will do. Otherwise $G$ and
$G_1$ are proper for $B$. Then $(M,M_1,H,H',e,u_\alpha)$ is a cone setting for $B$
(Definition~\ref{asm-theta-agree-local-factor-setting}), the positions of $H$ and
$H'$ are less than the ranks $b+2$ of $V$ and $b+2+a'$ of $W$, and
$z=k_\alpha\geq2$, as $\lfloor(b+2-a')/2\rfloor>2(p+q)+1$.

\emph{Step 1: the constant.} By
Lemma~\ref{asm-theta-agree-local-factor-ideal}(2) there is a fractional ideal
$\mathfrak a\neq0$ such that $1_vu_\alpha(l)\in\Gamma_v$ exactly when $l$ lies in
$\mathfrak a$ at $v$; $\operatorname{ord}_w\mathfrak a=0$ for the places $w$ above
$P$; and
$\nu'(\mathfrak a)=\nu'(\mathfrak b_H)^{e_1}\nu'(\mathfrak b_{H'})^{e_2}$ for every
continuous character $\nu'$ trivial on $\mathbf A^\times$ and unramified at the
places not above $P$. Let $\kappa_1$ be the constant of
Lemma~\ref{asm-theta-agree-local-factor-ideles} for $(B,V,m,\mu,\mathfrak a,z)$,
and $\kappa_2$ that of Lemma~\ref{asm-theta-agree-local-factor-stabilizer} for
$(B,V,K,m,\mu_K)$. Put $\kappa=\kappa_1\kappa_2\neq0$.

Now let the triples be given, with $\nu=\alpha/\beta$. The character $\nu$ is
continuous; it is trivial on $\mathbf A^\times$, as
$\alpha|_{\mathbf A^\times}=\eta^{a'}$, $\beta|_{\mathbf A^\times}=\eta^{b+2}$, $a'+b+2$
is even and $\eta^2=1$; $\alpha$ and $\beta$, hence $\nu$, $\beta$ and $\beta/\alpha$,
are unramified at the places not above $T$; $\nu$ is ramified at the places above
$T$; and $T\subset P$, so $\operatorname{ord}_w\mathfrak a=0$ for $w$ above $T$.
Hence Lemma~\ref{asm-theta-agree-local-factor-ideles} applies to $\nu$ and
$S_0\subset T$, and $\nu(\mathfrak a)=\nu(\mathfrak b_H)^{e_1}\nu(\mathfrak b_{H'})^{e_2}$.
For $v\notin T$ put $u_v=u_v(\nu)\neq0$
(Lemma~\ref{asm-theta-agree-local-factor-ideles}(1)). Let $L$ be the archimedean
cone of Lemma~\ref{asm-theta-agree-local-factor-archimedean} and $\Lambda_v$ the
map of Lemma~\ref{asm-theta-agree-local-factor-finite}(2), both for the Weil pair
$(V,W,M,\beta,\beta/\alpha,\omega)$, $K$, $m$, $\nu$ and $z$. Then (a) and (b) hold,
and $\lambda_0=\kappa\,\nu(\mathfrak a)\,L^{S_0}(\nu,k_\alpha)$.

\emph{Step 2: the integrand of a pure tensor.} Let $T'$, $f$ and $\phi_v$ be as in
(e), $\Phi=f\otimes\bigotimes_v\phi_v$, and $y$ a point of
$\mathbf A^{I_{M_1}}$. For $k\in K$ and $l\in\mathbf A_E^\times$ put
$I(k,l)=e(\rho(u_\alpha(l))\,\omega(k,1)\Phi)(y)\,\nu(l)|l|^z$. For $k'$ in
$K_\infty$ or $K_v$ and $l'$ in $A_\infty$ or $A_v$ put
\[
g_\infty(k',l')=e_\infty\bigl((T^\infty_{k',l'})_\infty f\bigr)(y_\infty),\qquad
g_v(k',l')=e_v\bigl((T^v_{k',l'})_v\phi_v\bigr)(y_v),
\]
and $t=\prod_{v\notin T'}\mathrm{vac}_{1,v}(y_v)$. As $\beta$ and $\beta/\alpha$
are unramified at the places not above $T'$,
Lemma~\ref{asm-theta-agree-local-factor-pure} and the choice of $\mathfrak a$ give
\[
I(k,l)=t\,g_\infty(k^{(\infty)},l^{(\infty)})\prod_{v\in T'}g_v(k^{(v)},l^{(v)})\,\nu(l)|l|^z
\]
if $l$ lies in $\mathfrak a$ at every finite $v\notin T'$, and $I(k,l)=0$
otherwise.

\emph{Step 3: the integral over $\mathbf A_E^\times$.} Fix $k\in K$. By
Lemmas~\ref{asm-theta-agree-local-factor-archimedean}
and~\ref{asm-theta-agree-local-factor-finite}(2), the functions
$l'\mapsto g_\infty(k^{(\infty)},l')\nu(l')|l'|^z$ and
$l'\mapsto g_v(k^{(v)},l')\nu(l')|l'|^z$ are integrable on $A_\infty$ and $A_v$. By
Lemma~\ref{asm-theta-agree-local-factor-ideles}(2), with $T'\supset T$, the
function $l\mapsto I(k,l)$ is $\mu$-integrable and
\[
\Bigl(\prod_{v\in T'\setminus T}u_v\Bigr)\int I(k,l)\,d\mu(l)
=t\,\kappa_1\,\nu(\mathfrak a)\,L^{S_0}(\nu,z)\,h_\infty(k^{(\infty)})\prod_{v\in T'}h_v(k^{(v)}),
\]
where $h_\infty(k')=\int_{A_\infty}g_\infty(k',l')\nu(l')|l'|^z\,dm_{A,\infty}(l')$
and $h_v(k')=\int_{A_v}g_v(k',l')\nu(l')|l'|^z\,dm_{A,v}(l')$.

\emph{Step 4: the integral over $K$, (e).} By the same two lemmas, $h_\infty$ and
the $h_v$ are integrable on $K_\infty$ and $K_v$, with
$\int_{K_\infty}h_\infty\,dm_{K,\infty}=L(f)(y_\infty)$ and
$\int_{K_v}h_v\,dm_{K,v}=\Lambda_v(\phi_v)(y_v)$, by the definitions of the
archimedean and local cones. By
Lemma~\ref{asm-theta-agree-local-factor-stabilizer}, $k\mapsto\int I(k,l)\,d\mu(l)$
is $\mu_K$-integrable and, as
$\operatorname{cone}_\alpha(\Phi)(y)=\int_K\int I(k,l)\,d\mu(l)\,d\mu_K(k)$,
\[
\Bigl(\prod_{v\in T'\setminus T}u_v\Bigr)\operatorname{cone}_\alpha(\Phi)(y)
=t\,\kappa_1\kappa_2\,\nu(\mathfrak a)\,L^{S_0}(\nu,z)\,L(f)(y_\infty)\prod_{v\in T'}\Lambda_v(\phi_v)(y_v).
\]
The right side is $\lambda_0$ times the value at $y$ of
$L(f)\otimes\bigotimes_v\psi_v$, which is (e).

\emph{Step 5: the vacuum, (c).} Let $v\notin T$, $k\in K_v$ and $l\in A_v$. By
Lemma~\ref{asm-theta-agree-local-factor-vacuum}, as $\beta$ and $\beta/\alpha$ are
unramified above $v$, $\omega(k,1)$ acts at $v$ by
$\mathrm{vac}_v\mapsto\mathrm{vac}_v$; and $\rho(1_vu_\alpha(l))$ acts at $v$ by
$\mathrm{vac}_v\mapsto\rho_v(1_vu_\alpha(l))\mathrm{vac}_v$
(Lemma~\ref{asm-theta-agree-local-factor-evaluation}(1)). So
$(T^v_{k,l})_v\mathrm{vac}_v=\rho_v(1_vu_\alpha(l))\mathrm{vac}_v$
(Lemma~\ref{asm-theta-agree-finite-inputs-schur}(1)), and by
Lemma~\ref{asm-theta-agree-local-factor-finite}(1) and the choice of $\mathfrak a$,
$e_v((T^v_{k,l})_v\mathrm{vac}_v)=\mathrm{vac}_{1,v}$ if $l$ lies in $\mathfrak a$
at $v$ and $0$ otherwise. As $m_{K,v}(K_v)=1$,
\[
\Lambda_v(\mathrm{vac}_v)=\operatorname{cone}_v(\mathrm{vac}_v)
=\Bigl(\int_{A_v}\mathbf 1[l\in\mathfrak a\text{ at }v]\,\nu(l)|l|^z\,dm_{A,v}(l)\Bigr)\mathrm{vac}_{1,v}
=u_v\,\mathrm{vac}_{1,v},
\]
and $u_v\neq0$.

\emph{Step 6: linearity, (d).} Let $Q$ be the set of Schwartz-Bruhat $\Phi$ such
that for every $y$ the function
$l\mapsto e(\rho(u_\alpha(l))\omega(k,1)\Phi)(y)\nu(l)|l|^z$ is $\mu$-integrable
for every $k\in K$, and its integral is $\mu_K$-integrable in $k$. The integrand is
linear in $\Phi$, as $\rho(u_\alpha(l))$, $\omega(k,1)$ and $e$ are linear. So $Q$
is closed under linear combinations, and $\operatorname{cone}_\alpha$ is linear on
$Q$. By Steps 3 and 4, applied with $T'$ the union of $T$ and the set of places
where $\phi_v\neq\mathrm{vac}_v$, $Q$ contains every $f\otimes\bigotimes_v\phi_v$
with $\phi_v$ Schwartz-Bruhat and $\phi_v=\mathrm{vac}_v$ outside a finite set. By
Lemma~\ref{asm-theta-agree-local-factor-span}, these span the Schwartz-Bruhat
functions, so $Q$ contains all of them, and (d) holds.
\end{proof}
% Lean: Cone.coneFac_all (ConeShifts/VacShiftOf.lean): coneIntegralFactorisation_of (part (1),
% read as ConeFacA by coneFactorisation_of_codimension_two) and coneFacB_of (part (2))
% (ConeFactorisationOfCodimensionTwo.lean) apply gen (ConeShifts/PureIte.lean) to pairV t,
% resp. pairV' t, after IdealM (Id.idealId'), unr_div, nu_cont and nu_res. In gen, Step 1 is
% hI, hKF, hAr and hF; Steps 2 to 4 are main (e1 from PureFac, HT from IdeleFub, HK from
% KFub); Step 5 is the vacuum clause (ActsEx, VacFixX, opF_eq, acts_comp, VacShift); Step 6
% is Iadd, Qadd and QA with SchSpan. The inputs are ideleFub, kFub_of StabEx, vacFixX_of StabEx,
% finCone2_of (VacShift, FinM), archCone_of, pureFac_of and schSpan. What remains open is in
% the sublemmas asm-theta-agree-local-factor-ideal, -ideles, -stabilizer, -vacuum,
% -archimedean and -pure.

\subsection{Values and invariance of local cones}\label{subsection-asm-local-cone-values}

\begin{lemma}\label{asm-theta-agree-local-ramified-values-integrals}
Let $k$ be a nonarchimedean local field with ring of integers $\mathcal O$,
uniformizer $\varpi$ and residue field of $q$ elements, let $dr$ be the Haar
measure on $k$ with $\int_{\mathcal O}dr=1$, and let $\psi$ be a continuous
character of $k$ with values of absolute value $1$ such that, for $x\in k$,
$\psi(xy)=1$ for every $y\in\mathcal O$ if and only if $x\in\mathcal O$. Let
$\nu$ be a unitary character of $k^\times$ and $c\geq1$ an integer such that
$\nu$ is trivial on the units in $1+\varpi^c\mathcal O$ and not on the units in
$1+\varpi^{c-1}\mathcal O$. Put $\delta(x)=\nu(x)^{-1}$ for
$x\in\mathcal O^\times$, $\delta(x)=0$ for $x\in k\setminus\mathcal O^\times$,
and
\[
G(y)=\int_k\delta(r)\,\psi(yr)\,dr\qquad(y\in k).
\]
Rows and columns of matrices are counted from $0$, and $E_{i,j}$ is the matrix
with entry $1$ at $(i,j)$ and $0$ elsewhere. For an injective map $e$ from
$\{0,\ldots,n-1\}$ to $\{0,\ldots,d-1\}$, let $s_e$ be the function on
$M_{n\times d}(k)$ with $s_e(Y)=\delta(\det Y_e)$, where
$Y_e=(Y_{i,e(l)})_{i,l}$, if the columns of $Y$ with index in the image of $e$
have entries in $\mathcal O$ and the other columns have entries in
$\varpi^c\mathcal O$, and $s_e(Y)=0$ otherwise. For a set $P$ of positions,
read $Z\in k^P$ as the matrix with entry $Z_p$ at $p\in P$ and $0$ elsewhere,
and let $dZ$ be the product of the measures $dr$. For injective maps
$j_r\colon\{0,\ldots,n_1-1\}\to\{0,\ldots,n-1\}$ and
$j_c\colon\{0,\ldots,d_1-1\}\to\{0,\ldots,d-1\}$ and
$y\in M_{n_1\times d_1}(k)$, let $\operatorname{ext}(y)$ be the $n\times d$
matrix with entry $y_{il}$ at $(j_r(i),j_c(l))$ and $0$ elsewhere.
\begin{enumerate}
\item $G(uy)=\nu(u)G(y)$ for $u\in\mathcal O^\times$; $G(y)=0$ unless
  $y\varpi^c\in\mathcal O^\times$; and $G(\varpi^{-c})=q^{-c}\mathfrak g(\nu)$,
  where
  $\mathfrak g(\nu)=\sum_{y\in(\mathcal O/\varpi^c)^\times}\nu(y)^{-1}\psi(y\varpi^{-c})$
  for fixed representatives $y\in\mathcal O$.
\item Each $s_e$ is a Schwartz-Bruhat function.
\item Let $n_1=n$ and $j_r$ the identity, let $e_1$ be a map with
  $j_c\circ e_1=e$, let $j_0$ be a column outside the images of $e$ and $j_c$,
  and let $P$ be the set of positions in column $j_0$. Then
  \[
  \int_{k^P}s_e\bigl(Z+\operatorname{ext}(y)\bigr)\,dZ=q^{-cn}s_{e_1}(y).
  \]
\item Let $n=b+2$, $n_1=b$ and $j_r(i)=i+1$, let $e_1$ be a map with
  $j_c(e_1(i))=e(i+1)$ for $i<b$, and suppose that $e(0)$ and $e(b+1)$ are not
  in the image of $j_c$. Let $P$ be a set of positions that contains
  $(b+1,e(0))$ and contains neither $(0,e(b+1))$ nor any $(i+1,j_c(l))$, such
  that either $P$ contains no $(0,e(i))$ with $i\neq b+1$ and no $(i+1,e(0))$,
  or $P$ contains no $(i,e(b+1))$ with $i\neq0$ and no $(b+1,e(i+1))$. Let $r$
  be the number of positions in $P$ whose column is not in the image of $e$.
  Then for $s,t\in k$
  \[
  \int_{k^P}\psi(-tZ_{b+1,e(0)})\,s_e\bigl(Z+\operatorname{ext}(y)+sE_{0,e(b+1)}\bigr)\,dZ
  =q^{-cr}\,\delta(s)\,G(t)\,s_{e_1}(y).
  \]
\item Let $E$ and $F$ be as in Definition~\ref{asm-theta-agree-models}, let
  $k=F_v$ for a finite place $v$ of $F$, and let $t$ and $\bar t$ be the two
  $F$-embeddings of $E$ into $F_v$. Let $\chi$ be a Hecke character of $E$ with
  $\chi_{\bar t}=\chi_t^{-1}$ (Definition~\ref{asm-theta-agree-coinvariants-local}(4))
  such that $\nu=\chi_t$, let $A_v$ be as in
  Definition~\ref{asm-theta-agree-local-cones}(2), and let $\mu$ be a Haar
  measure on $A_v$ that gives mass $1$ to the set of $l\in A_v$ whose powers
  have compact closure. Then for $C\in\mathbf C$ and an integer $z\geq0$
  \[
  \int_{A_v}\delta(\pi_t(l))\,G(\pi_{\bar t}(l))\,C\,\chi(l)|l|^z\,d\mu(l)
  =G(\varpi^{-c})\,\chi_t(\varpi)^c\,q^{cz}\,C.
  \]
\end{enumerate}
\end{lemma}
\leanref{Development/WeilBound/ThetaLifts/LocalTheory/StandardValuations/IntIndPow.lean}{PositivityTheory.ThKG.Std.dl}
\leanref{Development/WeilBound/ThetaLifts/LocalTheory/StandardValuations/IntIndPow.lean}{PositivityTheory.ThKG.Std.gI}
\leanref{Development/WeilBound/ThetaLifts/LocalTheory/StandardValuations/IntIndPow.lean}{PositivityTheory.ThKG.Std.gI_unit}
\leanref{Development/WeilBound/ThetaLifts/LocalTheory/StandardValuations/IntIndPow.lean}{PositivityTheory.ThKG.Std.gI_zero}
\leanref{Development/WeilBound/ThetaLifts/LocalTheory/StandardValuations/IntIndPow.lean}{PositivityTheory.ThKG.Std.gI_gauss}
\leanref{Development/WeilBound/ThetaLifts/LocalTheory/StandardValuations/SplitProd.lean}{PositivityTheory.ThKG.Std.std}
\leanref{Development/WeilBound/ThetaLifts/LocalTheory/StandardValuations/SplitProd.lean}{PositivityTheory.ThKG.Std.std_mem}
\leanref{Development/WeilBound/ThetaLifts/LocalTheory/StandardValuations/SplitProd.lean}{PositivityTheory.ThKG.Std.int_ev}
\leanref{Development/WeilBound/ThetaLifts/LocalTheory/StandardValuations/SplitProd.lean}{PositivityTheory.ThKG.Std.int_cone}
\leanref{Development/WeilBound/ThetaLifts/LocalTheory/StandardValuations/IdeleInt.lean}{PositivityTheory.ThKG.Std.idele_int}
\begin{proof}
(1) The substitution $r\mapsto ur$ preserves $dr$ and gives $G(uy)=\nu(u)G(y)$.
Suppose $y\varpi^c\in\varpi\mathcal O$. As $\nu$ is not trivial on the units in
$1+\varpi^{c-1}\mathcal O$, there is such a unit $u$ with $\nu(u)\neq1$. For
$r\in\mathcal O^\times$, $y(u-1)r\in\mathcal O$, so $\psi(yur)=\psi(yr)$ and
$G(uy)=G(y)$. Hence $(\nu(u)-1)G(y)=0$ and $G(y)=0$. Suppose
$y\varpi^c\notin\mathcal O$, and choose $s\in\mathcal O$ with
$\psi(y\varpi^cs)\neq1$. The translation $r\mapsto r+\varpi^cs$ preserves $dr$
and $\delta$, as $\nu$ is trivial on the units in $1+\varpi^c\mathcal O$, and
multiplies $\psi(yr)$ by $\psi(y\varpi^cs)$. So $G(y)=\psi(y\varpi^cs)G(y)$, and
$G(y)=0$. For $y=\varpi^{-c}$ the integrand vanishes off $\mathcal O$ and is
invariant under translation by $\varpi^c\mathcal O$, a set of measure $q^{-c}$;
summing over the classes modulo $\varpi^c$, on which $\delta$ vanishes unless
the class is a unit, gives $G(\varpi^{-c})=q^{-c}\mathfrak g(\nu)$.

(2) The function $s_e$ is locally constant, as $\delta$ is (by the triviality of
$\nu$ on $1+\varpi^c\mathcal O$), and it vanishes outside a compact set.

(3) For $Z$ in column $j_0$, the columns of $Z+\operatorname{ext}(y)$ in the image
of $e$ are those of $y$ in the image of $e_1$; the other columns are the other
columns of $y$, the column $j_0$, which is $Z$, and zero columns. So the
integrand is $s_{e_1}(y)$ times the indicator function of
$(\varpi^c\mathcal O)^n$ at $Z$, whose integral is $q^{-cn}$.

(4) Let $Y=Z+\operatorname{ext}(y)+sE_{0,e(b+1)}$. The entry of $Y$ at
$(0,e(b+1))$ is $s$, that at $(b+1,e(0))$ is $Z_{b+1,e(0)}$, and the minor of
$Y_e$ on the rows and columns $1,\ldots,b$ is $y_{e_1}$. Under the first
alternative for $P$, the entries of the row $0$ of $Y_e$ other than the last and
the entries of its column $0$ in the rows $1,\ldots,b$ vanish; under the second,
the entries of its column $b+1$ other than the first and the entries of its row
$b+1$ in the columns $1,\ldots,b$ vanish. Expanding the determinant gives
$\det Y_e=-s\,Z_{b+1,e(0)}\det y_{e_1}$ in both cases. When the integrand is not
$0$, these three factors lie in $\mathcal O$, and $\delta$ is multiplicative on
$\mathcal O$; the condition on $s$ is that $s\in\mathcal O$, which $\delta(s)$
already records. So the integrand is the product of $\delta(s)\,s_{e_1}(y)$ and
of one factor for each $p\in P$: the factor $\delta(-Z_p)\psi(-tZ_p)$ for
$p=(b+1,e(0))$, the indicator function of $\mathcal O$ for the other $p$ with
column in the image of $e$, and that of $\varpi^c\mathcal O$ for the $r$
positions with column outside it. Their integrals are $G(t)$, after
$r\mapsto-r$, then $1$ and $q^{-c}$.

(5) Through $l\mapsto(\pi_t(l),\pi_{\bar t}(l))$ the group $A_v$ is
$k^\times\times k^\times$, with $\chi(l)=\chi_t(\pi_t(l))\chi_t(\pi_{\bar t}(l))^{-1}$
and $|l|=|\pi_t(l)||\pi_{\bar t}(l)|$, and the set of $l$ whose powers have
compact closure is $\mathcal O^\times\times\mathcal O^\times$. Translate the
integral by the element with components $(1,\varpi^{-c})$, which preserves
$\mu$. At $l$ with components $(x,y)$ the new integrand is
$\delta(x)G(\varpi^{-c}y)C\nu(x)\nu(\varpi^{-c}y)^{-1}|x|^z|\varpi^{-c}y|^z$. By
(1) it vanishes unless $x,y\in\mathcal O^\times$, and there it equals
$G(\varpi^{-c})\nu(\varpi)^cq^{cz}C$.
\end{proof}
% Lean: (1) gI_unit, gI_zero (gI_small, gI_big), gI_gauss (int_cosets, int_ind_pow); the
% conductor exponent is condExp with cond_ex. (2) std_mem. (3) int_ev (split_prod, prod_vol,
% prod_mid). (4) int_cone (detA, detB, split_prod, gI_neg, prod_vol, prod_mid). (5) idele_int
% (uI, cmp_uI, chr_id, abs_id, units_eq, gI_zero).

\begin{lemma}\label{asm-theta-agree-local-ramified-values-zeta}
Let $k$, $\mathcal O$ and $\varpi$ be as in
Lemma~\ref{asm-theta-agree-local-ramified-values-integrals}, let $a>a'\geq0$ with
$d=a+a'$ even, let $\alpha$ and $\beta$ be unitary smooth characters of
$k^\times$ (Definition~\ref{asm-theta-agree-finite-split-models}) with
$\nu=\alpha/\beta$ ramified, and let $c\geq1$ be such that $\nu$ is trivial on
the units in $1+\varpi^c\mathcal O$. Let $Z$, $Z'$, $s_A$ and $s_B$ be as in
Lemma~\ref{th-bottom-split-zeta} for these data. Then $Z(s_A)=Z'(s_B)$ and
$Z(s_A)\neq0$.
\end{lemma}
\leanref{Development/WeilBound/ThetaLifts/LocalTheory/StandardValuations/SmallFSing.lean}{PositivityTheory.ThKG.Std.std_pair}
\leanref{Development/WeilBound/ThetaLifts/LocalTheory/StandardValuations/SmallFSing.lean}{PositivityTheory.ThKG.Std.pRam_ne}
\begin{proof}
Let $h\in GL_d(k)$, and write $h$ and $h^{-1}$ in blocks of sizes $a$ and
$a'$, with lower right blocks $N_{22}$ and $N'_{22}$ and upper left blocks
$N_{11}$ and $N'_{11}$. Suppose first that $\det N_{22}=0$. Then
$\det N'_{11}=0$: otherwise the identity $hh^{-1}=1$ in blocks gives
$N_{22}(N'_{22}-N'_{21}N'^{-1}_{11}N'_{12})=1$. The first $a$ columns of
$gX_0(h^{-1})^{\mathsf T}$ form $g\,N'^{\mathsf T}_{11}$, and the last $a'$
columns of $gY_0h$ form $gN_{22}$; both are singular, so $s_A$ and $s_B$ vanish
there, and $Z(s_A)(h)=0=Z'(s_B)(h)$.

Suppose now that $\det N_{22}\neq0$, put $y=N_{22}^{-1}N_{21}$ and let
$\ell(y)\in GL_d(k)$ have blocks $1_a,0,y,1_{a'}$. Then $h\ell(y)^{-1}$ is block
upper triangular. By Lemma~\ref{asm-theta-agree-finite-zeta}(1), $Z(s_A)$ and
$Z'(s_B)$ lie in $I(\alpha,\beta)$, so it suffices to compare their values at
$\ell(y)$. There $\det\ell(y)=1$, $X_0(\ell(y)^{-1})^{\mathsf T}=(1_a\ \,-y^{\mathsf T})$
and $Y_0\ell(y)=(y\ \,1_{a'})$. For $g\in GL_a(k)$, $s_A(g\,(1_a\ \,-y^{\mathsf T}))$
vanishes unless $g\in GL_a(\mathcal O)$, and then equals $\nu(\det g)^{-1}$ if
$y$ has entries in $\varpi^c\mathcal O$ and $0$ otherwise. So only the double
coset of $\varpi^0$ contributes, and $Z(s_A)(\ell(y))$ is $1$ or $0$ according as
$y$ has entries in $\varpi^c\mathcal O$ or not. The same computation gives the
same value for $Z'(s_B)(\ell(y))$. At $h=1$ both values are $1$, so
$Z(s_A)\neq0$.
\end{proof}
% Lean: std_pair (with low, pRam_low, pRam'_low, bigF_std, smallF_std, std_mul, bigF_sing,
% smallF_sing, and pRam_mem, pRam'_mem for the law), pRam_ne. Lemma th-bottom-split-zeta cites
% another Lean statement, which allows a nonzero multiple of s_B; here the multiple is 1.

\begin{definition}\label{asm-theta-agree-local-ramified-values-coordinates}
Let $B$ be a theta base of a CM field $E$, let $M$ be a model for lattice spaces
$(X,W)$ of ranks $(n,d)$, with Schrödinger operators $\rho$, put
$I=\{0,\ldots,n-1\}\times\{0,\ldots,d-1\}$, and let $v$ be a finite place of
$F$. For $n\times d$ matrices $Y$ and $Z$ put $Y\cdot Z=\sum_{i,j}Y_{ij}Z_{ij}$.
Let $c_1$ and $c_2$ be maps from $\mathbb W(X,W)$ to $M_{n\times d}(F_v)$. A
\emph{coordinate model} at $v$ with coordinates $(c_1,c_2)$ is a linear map
$\iota$ from functions on $F_v^I$ to functions on $M_{n\times d}(F_v)$,
bijective between the Schwartz-Bruhat functions, such that: $\iota(\mathrm{vac}_v)$
is the indicator function of $M_{n\times d}(\mathcal O_v)$; for all
$Y_1,Y_2\in M_{n\times d}(F_v)$ there is $u$ at $v$ with $c_1(u)=Y_1$ and
$c_2(u)=Y_2$; and for every $u$ at $v$ with $c_1(u)\cdot c_2(u)=0$ and every
Schwartz-Bruhat $\phi$, $\rho(u)$ acts at $v$
(Definition~\ref{asm-theta-agree-models}(2)) by $\phi\mapsto\phi_1$ for a
Schwartz-Bruhat $\phi_1$ with
\[
\iota\phi_1(Y)=\psi_v\bigl(Y\cdot c_2(u)\bigr)\,\iota\phi\bigl(Y+c_1(u)\bigr).
\]
\end{definition}
\leanref{Development/WeilBound/ThetaLifts/LocalTheory/StandardValuations/LocalSplittingLinearSkew.lean}{PositivityTheory.ThKG.Std.Lin}
\leanref{Development/WeilBound/ThetaLifts/LocalTheory/StandardValuations/LocalSplittingLinearSkew.lean}{PositivityTheory.ThKG.Std.Hv}

\begin{lemma}\label{asm-theta-agree-local-ramified-values-linear}
Let $B$ be a theta base of a CM field $E$, let $X$ be a Hermitian and $W$ a
skew-Hermitian lattice space, of ranks $n$ and $d$ and some signatures, with
matrices $J_X$ and $J_W$, let $M$ be a model for $(X,W)$, let $v$ be a finite
place of $F$ and $t$ and $\bar t$ the two $F$-embeddings of $E$ into $F_v$, and
let $b\in GL_n(F_v)$ and $b'\in GL_d(F_v)$. Put
\[
c_W(u)=b^{-1}\pi_t(u)(b'^{-1})^{\mathsf T},\qquad
c_B(u)=-b^{\mathsf T}t(J_X)\,\pi_{\bar t}(u)\,t(J_W)^{\mathsf T}b'.
\]
A linear model of the first kind for $(t,b,b')$
(Definition~\ref{asm-theta-agree-models}(5)) that carries $\mathrm{vac}_v$ to
the indicator function of $M_{n\times d}(\mathcal O_v)$ is a coordinate model at
$v$ with coordinates $(c_W,c_B)$
(Definition~\ref{asm-theta-agree-local-ramified-values-coordinates}), and a
linear model of the second kind for $(t,b,b')$ with the same property is one
with coordinates $(c_B,-c_W)$.
\end{lemma}
\leanref{Development/WeilBound/ThetaLifts/LocalTheory/StandardValuations/LocalSplittingLinearSkew.lean}{PositivityTheory.ThKG.Std.local_splitting_linear_skew}
\leanref{Development/WeilBound/ThetaLifts/LocalTheory/StandardValuations/LocalSplittingLinearSkew.lean}{PositivityTheory.ThKG.Std.lin_Bar}
\begin{proof}
The maps $\pi_t$ and $\pi_{\bar t}$ identify the elements of
$\mathbf A_E$ at $v$ with $F_v\times F_v$, and $t(J_X)$ and $t(J_W)$ are
invertible. So for all $Y_1,Y_2$ there is $u$ at $v$ with $c_W(u)=Y_1$ and
$c_B(u)=Y_2$, and every $u$ at $v$ is $u_1+u_0$ with $u_1,u_0$ at $v$,
$\pi_{\bar t}(u_1)=0$ and $\pi_t(u_0)=0$; then $c_W(u_1)=c_W(u)$,
$c_B(u_0)=c_B(u)$ and $c_B(u_1)=0=c_W(u_0)$. A direct computation from the
definition of the form of $\mathbb W(X,W)$, with $c(J_X)^{\mathsf T}=J_X$ and
$c(J_W)^{\mathsf T}=-J_W$, shows that the form of two elements $x,u$ at $v$ is
$1_v$ times $c_W(x)\cdot c_B(u)-c_W(u)\cdot c_B(x)$. For $u$ at $v$, the
definition of $\rho$ shows that $\rho(u)$ acts at $v$ by a map that preserves
Schwartz-Bruhat functions, and that
$\rho(u_1)\rho(u_0)=\psi(\tfrac12\langle u_1,u_0\rangle)\rho(u)$.

Let $\iota$ be of the first kind, let $u$ be at $v$ with
$c_W(u)\cdot c_B(u)=0$, decomposed as above, and let $\phi$ be Schwartz-Bruhat.
Then $\langle u_1,u_0\rangle$ is $1_v$ times $c_W(u)\cdot c_B(u)=0$, so $\rho(u)$
acts at $v$ by $\phi\mapsto\phi_1$, where $\rho(u_0)$ acts by
$\phi\mapsto\phi'$ and $\rho(u_1)$ by $\phi'\mapsto\phi_1$. Let $Y$ be a matrix,
and choose $x$ at $v$ with $\pi_{\bar t}(x)=0$ and $c_W(x)=Y+c_W(u)$. As
$\pi_{\bar t}(u_1)=0$, the first property of a linear model gives
$\iota\phi_1(Y)=\iota\phi'(c_W(x))$. As $\pi_t(u_0)=0$, the second gives
$\iota\phi'(c_W(x))=\psi(\langle x,u_0\rangle)\iota\phi(c_W(x))$, and
$\langle x,u_0\rangle$ is $1_v$ times $c_W(x)\cdot c_B(u)=Y\cdot c_B(u)$. So
$\iota\phi_1(Y)=\psi_v(Y\cdot c_B(u))\,\iota\phi(Y+c_W(u))$. For the second
kind the same argument applies with $t$ and $\bar t$ exchanged, $c_B$ in place of
$c_W$ and $-c_W$ in place of $c_B$, since the form is then
$c_B(x)\cdot(-c_W(u))-c_B(u)\cdot(-c_W(x))$.
\end{proof}
% Lean: lin_of (the decomposition dec, KCo.schr_schr, IsLin.trans, IsLin.mult), with symp_c,
% dec_c and sur_c, and grams for the invertibility of t(J_X) and t(J_W).

\begin{lemma}\label{asm-theta-agree-local-ramified-values-evaluation}
Let $B$ be a theta base of a CM field $E$, let $v\in P$, and let $M$ and $M_1$ be
regular models for lattice spaces $(X,W)$ and $(X_1,W_1)$ of ranks $(n,d)$ and
$(n_1,d_1)$. Read matrices over $F_v$, and the matrices
$\operatorname{ext}(y)$ and the measures $dZ$, as in
Lemma~\ref{asm-theta-agree-local-ramified-values-integrals} for $k=F_v$.
\begin{enumerate}
\item Let $e$ be an evaluation map for $(M,M_1,\iota_X,\iota_W,\mathcal J)$ with
  components $e_v$ (Definition~\ref{asm-theta-agree-planes}(5)), let $\iota$ and
  $\iota_1$ be coordinate models at $v$ for $M$ and $M_1$ with coordinates
  $(c_1,c_2)$ and $(c'_1,c'_2)$
  (Definition~\ref{asm-theta-agree-local-ramified-values-coordinates}), let
  $j_r$ and $j_c$ be injective maps from the rows and columns of size $n_1$ and
  $d_1$ to those of size $n$ and $d$, and let $P$ be a set of positions in
  $n\times d$ matrices. Suppose that: (a) $c_i(\iota_Xz\iota_W^{\mathsf T})=\operatorname{ext}(c'_i(z))$
  for $i=1,2$ and every $z\in\mathbb W(X_1,W_1)$; (b) for $p\in P$ and
  $s\in F_v$ there is $u\in\mathcal J$ at $v$ with $c_1(u)=sE_p$ and
  $c_2(u)=0$; (c) every $p\notin P$ is $(j_r(i),j_c(l))$ for some $i,l$, or for
  every $s\in F_v$ there is $u\in\mathcal J$ at $v$ with $c_1(u)=0$ and
  $c_2(u)=sE_p$. Then for Schwartz-Bruhat $\phi$ and
  $y\in M_{n_1\times d_1}(F_v)$
  \[
  \iota_1(e_v\phi)(y)=\int_{F_v^P}\iota\phi\bigl(Z+\operatorname{ext}(y)\bigr)\,dZ.
  \]
\item Two coordinate models at $v$ for $M$ with the same coordinates agree on
  the Schwartz-Bruhat functions.
\end{enumerate}
In (3) and (4), let $(G,G_1,H,H')$ be a configuration of triple data for $B$
(Definition~\ref{asm-theta-agree-planes}(4)) whose models are regular, with
Hermitian $V$, $V'$, $V_1$, $V'_1$ and skew-Hermitian $W$, $W_1$, and let
$\iota$, $\iota_1$, $\iota'$ and $\iota'_1$ be coordinate models at $v$ for $M$,
$M_1$, $M'$ and $M'_1$, the first two with the coordinates $(c_W,c_B)$ and the
last two with $(c_B,-c_W)$ of
Lemma~\ref{asm-theta-agree-local-ramified-values-linear} for $t=\tau_v$ and the
matrices $b_v$ of the spaces involved. Let $j$ be the increasing map from
$\{0,\ldots,d_1-1\}$ to $\{0,\ldots,d-1\}$ that omits the two positions of $H'$,
and let $\phi$ and $\phi'$ be Schwartz-Bruhat.
\begin{enumerate}
\setcounter{enumi}{2}
\item If $(G,G_1,H,H')$ is an $\alpha$-configuration with $G$ of ranks
  $(b+2,a')$, $e$ is an evaluation map for
  $(M,M_1,\iota_H,\iota_{H'},\mathcal J_\alpha(H,H'))$ and $e'$ one for
  $(M',M'_1,1,\iota_{H'},\mathcal J_0(H'))$, with components $e_v$ and $e'_v$,
  then $\iota_1(e_v\phi)(y)$ is given by the formula of (1) with $P$ the set of
  positions in the column $0$ or in the row $0$ outside the column $b+1$,
  $j_r(i)=i+1$ and $j_c=j$, and $\iota'_1(e'_v\phi')(y)$ by the formula with $P$
  the set of positions in the column $b+1$, $j_r$ the identity and $j_c=j$.
\item If $(G,G_1,H,H')$ is a $\beta$-configuration with $G$ of ranks
  $(a,b+2)$, $e$ is an evaluation map for
  $(M,M_1,1,\iota_{H'},\mathcal J_0(H'))$ and $e'$ one for
  $(M',M'_1,\iota_H,\iota_{H'},\mathcal J_\beta(H,H'))$, then
  $\iota_1(e_v\phi)(y)$ is given by the formula of (1) with $P$ the set of
  positions in the column $a$, $j_r$ the identity and $j_c=j$, and
  $\iota'_1(e'_v\phi')(y)$ by the formula with $P$ the set of positions in the
  row $b+1$ or in the column $a+b+1$ outside the row $0$, $j_r(i)=i+1$ and
  $j_c=j$.
\end{enumerate}
\end{lemma}
\leanref{Development/WeilBound/ThetaLifts/LocalTheory/StandardValuations/LocalSplittingLinearSkew.lean}{PositivityTheory.ThKG.Std.ev_model}
\leanref{Development/WeilBound/ThetaLifts/LocalTheory/StandardValuations/ContinuousCut.lean}{PositivityTheory.ThKG.Std.evInt}
\leanref{Development/WeilBound/ThetaLifts/LocalTheory/StandardValuations/ConeVal.lean}{PositivityTheory.ThKG.Std.lin_uniq}
\leanref{Development/WeilBound/ThetaLifts/LocalTheory/StandardValuations/ConeVal.lean}{PositivityTheory.ThKG.Std.evA}
\leanref{Development/WeilBound/ThetaLifts/LocalTheory/StandardValuations/ConeVal.lean}{PositivityTheory.ThKG.Std.ev0B}
\leanref{Development/WeilBound/ThetaLifts/LocalTheory/StandardValuations/ConeVal.lean}{PositivityTheory.ThKG.Std.ev0W}
\leanref{Development/WeilBound/ThetaLifts/LocalTheory/StandardValuations/ConeVal.lean}{PositivityTheory.ThKG.Std.evB}
\begin{proof}
(1) Let $\rho$ and $\rho_1$ be the Schrödinger operators of $M$ and $M_1$, and
for $u$ at $v$ let $\phi\mapsto\rho_v(u)\phi$ be the map by which $\rho(u)$ acts
at $v$. If $u$ and $u_1$ are at $v$ and $e(\rho(u)\Phi)=\rho_1(u_1)e(\Phi)$ for
every Schwartz-Bruhat $\Phi$, then $e_v(\rho_v(u)\phi)=\rho_{1,v}(u_1)e_v(\phi)$
for Schwartz-Bruhat $\phi$: apply both sides to $f\otimes\phi\otimes\bigotimes_{w\neq v}\mathrm{vac}_w$
with $e_\infty(f)\neq0$, use the factorization of $e$, and compare the
components at $v$ of the two products, whose other components agree and are
nonzero. This applies to $u\in\mathcal J$ with $u_1=0$, and to
$u=\iota_Xz\iota_W^{\mathsf T}$ with $u_1=z$.

For Schwartz-Bruhat $\Phi$ on $M_{n\times d}(F_v)$ put
$l(\Phi)=\iota_1(e_v\phi)(0)$, where $\iota\phi=\Phi$, and extend $l$ linearly.
Let $p\in P$ and $s\in F_v$, and take $u$ as in (b). Then
$e_v(\rho_v(u)\phi)=e_v(\phi)$ and, as $c_1(u)\cdot c_2(u)=0$,
$\iota\rho_v(u)\phi(Y)=\iota\phi(Y+sE_p)$. So $l$ is invariant under the
translations by the matrices supported on $P$. Let $p\notin P$. If
$p=(j_r(i),j_c(l))$, choose $z$ at $v$ with $c'_1(z)=0$ and $c'_2(z)=sE_{i,l}$;
by (a), $u=\iota_Xz\iota_W^{\mathsf T}$ has $c_1(u)=0$ and $c_2(u)=sE_p$, and
$\iota_1(\rho_{1,v}(z)e_v\phi)(0)=\iota_1(e_v\phi)(0)$. Otherwise take $u$ as in
(c). In both cases $l(Y\mapsto\psi_v(Y\cdot sE_p)\Phi(Y))=l(\Phi)$. Finally,
$l$ takes the value $\iota_1(\mathrm{vac}_{1,v})(0)=1$ at the indicator
function $\mathbf 1$ of $M_{n\times d}(\mathcal O_v)$.

The linear form $\Phi\mapsto\int_{F_v^P}\Phi(Z)\,dZ$ has the same two
invariances and the same value at $\mathbf 1$, so its difference $l_0$ with $l$
is invariant under the Schrödinger operators of the pairs $(Z,Z')$ with $Z$
supported on $P$ and $Z'$ off $P$, and vanishes at $\mathbf 1$. Every pair $w$
is $w_0+w_1$ with $w_0$ of this kind and $w_1$ supported in the other
coordinates. If $w_1$ is integral, its Schrödinger operator fixes $\mathbf 1$;
otherwise some integral pair $w'$ of the first kind fixes $\mathbf 1$ and has
$\psi_v(\langle w',w_1\rangle)\neq1$, and the commutation relation shows that
$l_0$ vanishes at the translate of $\mathbf 1$ by $w_1$. So $l_0$ vanishes at
all Schrödinger translates of $\mathbf 1$. These span the Schwartz-Bruhat
functions, by the irreducibility of the Schrödinger representation
\cite{Wei64}. So $l(\Phi)=\int_{F_v^P}\Phi(Z)\,dZ$.

For $y\in M_{n_1\times d_1}(F_v)$ choose $z$ at $v$ with $c'_1(z)=y$ and
$c'_2(z)=0$, and put $u=\iota_Xz\iota_W^{\mathsf T}$, so that
$c_1(u)=\operatorname{ext}(y)$ and $c_2(u)=0$. Then
\[
\iota_1(e_v\phi)(y)=\iota_1\bigl(\rho_{1,v}(z)e_v\phi\bigr)(0)
=\iota_1\bigl(e_v\rho_v(u)\phi\bigr)(0)=l\bigl(\iota\rho_v(u)\phi\bigr)
=\int_{F_v^P}\iota\phi\bigl(Z+\operatorname{ext}(y)\bigr)\,dZ.
\]

(2) Let $\iota$ and $\iota'$ be the two coordinate models. The identity is an
evaluation map for $(M,M,1,1,\emptyset)$, with components the identity, and (1)
applies to it with $\iota_1=\iota'$, $P=\emptyset$ and $j_r$, $j_c$ the
identities. It gives $\iota'\phi(y)=\iota\phi(y)$.

(3) and (4) Let $H=(x,y,\iota_H,\mathfrak b)$ be in $V$ (in (4), in $V'$) at
positions $(i_0,i_1)=(0,b+1)$, with matrices $b=b_{V,v}$ and $b_1=b_{V_1,v}$,
and $H'=(x',y',\iota_{H'},\mathfrak b')$ in $W$ at positions $(j_0,j_1)$, which
are $(0,b+1)$ in (3) and $(a,a+b+1)$ in (4), with $b'=b_{W,v}$ and
$b'_1=b_{W_1,v}$. By the conditions on an adapted plane at $v\in P$, $b$ has
columns $\tau_v(x)$ and $\tau_v(y)$ at $i_0$ and $i_1$ and the columns of
$\tau_v(\iota_H)b_1$ elsewhere; with $h(x,y)=1$, $h(x,x)=h(y,y)=0$,
$h(\iota_Hu,x)=h(\iota_Hu,y)=0$ and $\iota_H^{\mathsf T}J\,c(\iota_H)=J_1$
this gives $b^{-1}\tau_v(x)=e_{i_0}$, $b^{-1}\tau_v(y)=e_{i_1}$,
$b^{\mathsf T}\tau_v(J)\bar\tau_v(x)=e_{i_1}$ and
$b^{\mathsf T}\tau_v(J)\bar\tau_v(y)=e_{i_0}$, where $\bar\tau_v$ is the other
embedding, and $b^{-1}\tau_v(\iota_H)$ and
$b^{\mathsf T}\tau_v(J)\bar\tau_v(\iota_H)$ are the $0,1$-matrix of
$i\mapsto i+1$ times $b_1^{-1}$ and $b_1^{\mathsf T}\tau_v(J_1)$; the same
holds for $H'$, with $-e_{j_1}$ for $x'$ as $W$ is skew-Hermitian. Hence (a)
holds for $\iota_H$ and $\iota_{H'}$, and also for $1$ and $\iota_{H'}$. For a
row $i$ and $s,s'\in F_v$, choose $p$ at $v$ with $b^{-1}\pi_{\tau_v}(p)=se_i$
and $b^{\mathsf T}\tau_v(J)\pi_{\bar\tau_v}(p)=s'e_i$; then
$px'^{\mathsf T}\in\mathcal J_0(H')$ has $c_W=sE_{i,j_0}$ and $c_B=s'E_{i,j_1}$.
For a column $j(l)$ and $s,s'$, choose $q$ at $v$ in the same way for
$\iota_{H'}q$; then $x(\iota_{H'}q)^{\mathsf T}\in\mathcal J_\alpha(H,H')$ has
$c_W=sE_{i_0,j(l)}$ and $c_B=-s'E_{i_1,j(l)}$. In (3) these elements give (b)
for the column $j_0$ and for the row $i_0$ outside $j_1$, and (c) for the column
$j_1$ and the row $i_1$; the other positions are $(i+1,j(l))$. For $e'$ only
the first elements are in $\mathcal J_0(H')$, and in the coordinates
$(c_B,-c_W)$ they give (b) for the column $j_1$ and (c) for the column $j_0$.
Case (4) is the same with the elements $xq^{\mathsf T}$ and
$(\iota_Hp)x'^{\mathsf T}$ of $\mathcal J_\beta(H,H')$ and rows and columns
exchanged. So (1) gives (3) and (4).
\end{proof}
% Lean: (1) ev_model, through ev_loc (pure_cancel, ThKF.pureAct), evInt (ContinuousCut.lean;
% the uniqueness is ThKH.Evl.core, whose cyclicity step ThKH.Evl.cyc is the irreducibility of
% the Schrödinger representation), Lin.act and Lin.sur. (2) lin_uniq. (3) evA and ev0B,
% (4) ev0W and evB, through plane (PE, PX), emb_c, vec_ex, vec_ex', col_el and row_el.

\begin{lemma}\label{asm-theta-agree-local-ramified-values-dictionary}
Let $(X,W,M,\chi_X,\chi_W,\omega)$ be a Weil pair for a theta base $B$ of a CM
field $E$ (Definition~\ref{asm-theta-agree-triple}(1)), with $X$ and $W$ of ranks
$n$ and $d$.
\begin{enumerate}
\item Let $v$ be a finite place of $F$ and $t\colon E\to F_v$ an $F$-embedding,
  and put $\nu'=(\chi_W)_t$ and $\alpha'=(\chi_X)_t$. There is a linear model of
  the second kind for $(t,1,1)$; $\pi_t$ maps $U_X^{(v)}$ onto $GL_n(F_v)$ and
  $U_W^{(v)}$ onto $GL_d(F_v)$; and through every linear model of the second
  kind for $(t,1,1)$, $\omega(g,1)$ reads as
  $\omega_B^{\nu',\alpha'}(\pi_t(g),1)$ and $\omega(1,h)$ as
  $\omega_B^{\nu',\alpha'}(1,\pi_t(h))$ for $g\in U_X^{(v)}$ and
  $h\in U_W^{(v)}$ (Definition~\ref{asm-theta-agree-finite-split-models}).
\end{enumerate}
In (2) and (3), let $v\in P$, $t=\tau_v$ and $\bar t$ the other $F$-embedding of
$E$ into $F_v$, and put $b=b_{X,v}$ and $b'=b_{W,v}$. For $g\in U_X^{(v)}$ and
$h\in U_W^{(v)}$ put $\gamma(g)=b^{-1}\pi_t(g)b$ and $\eta(h)=b'^{-1}\pi_t(h)b'$.
\begin{enumerate}
\setcounter{enumi}{1}
\item There is a linear model $\iota$ of the first kind for $(t,b,b')$ that
  carries $\mathrm{vac}_v$ to the indicator function of
  $M_{n\times d}(\mathcal O_v)$, through which $\omega(g,1)$ reads as
  $\omega_A^{\nu,\beta}(\gamma(g),1)$ and $\omega(1,h)$ as
  $\omega_A^{\nu,\beta}(1,\eta(h))$ for all such $g$ and $h$
  (Definition~\ref{asm-theta-agree-finite-split-models}), where
  $\nu=(\chi_W)_{\bar t}$ and $\beta=(\chi_X)_{\bar t}^{-1}$.
\item There is a linear model $\iota'$ of the second kind for $(t,b,b')$ with
  the same property of the vacuum, through which $\omega(g,1)$ reads as
  $\omega_B^{\nu',\alpha'}(\gamma(g),1)$ and $\omega(1,h)$ as
  $\omega_B^{\nu',\alpha'}(1,\eta(h))$, where $\nu'=(\chi_W)_t$ and
  $\alpha'=(\chi_X)_t$.
\end{enumerate}
\end{lemma}
\leanref{Development/WeilBound/CoreStatements/LocallyConstantCoinvariants/LmOfLinBar.lean}{PositivityTheory.KCo.DictX}
\leanref{Development/WeilBound/CoreStatements/LocallyConstantCoinvariants/PfDict.lean}{PositivityTheory.KCo.dictionary_extended}
\leanref{Development/WeilBound/ThetaLifts/LocalTheory/StandardValuations/VacuumImageIntertwiner.lean}{PositivityTheory.ThKG.StdDict.model_alt}
\leanref{Development/WeilBound/ThetaLifts/LocalTheory/StandardValuations/VacuumImageIntertwiner.lean}{PositivityTheory.ThKG.StdDict.model_Bar}
\leanref{Development/WeilBound/ThetaLifts/LocalTheory/StandardValuations/VacuumImageIntertwiner.lean}{PositivityTheory.ThKG.StdDict.inB}
\begin{proof}
(1) This is Lemma~\ref{asm-theta-agree-coinvariants-dictionary}, as
$\omega_B^{\nu',\alpha'}=\omega_{\nu',\alpha'}$ for the unitary smooth characters
$\nu'$ and $\alpha'$ (Definition~\ref{asm-theta-agree-finite-split-models}(1)).

(2) and (3). As $M$ is regular, $\mathrm{vac}_v$ is a vacuum. By
Definition~\ref{asm-theta-agree-doubling}(1), $b\,\mathcal O_v^n$ and
$b'\,\mathcal O_v^d$ are the completions of the lattices of $X$ and $W$ at the
place of $E$ given by $t=\tau_v$, and $\psi_v$ has conductor $\mathcal O_v$ as
$v\in P$ (Definition~\ref{gal-good-base}), so $\delta=1$ has the property of
Lemma~\ref{asm-theta-agree-dictionary-vacuum}. With $\bar t=t\circ c$, parts (1)
and (2) of Lemma~\ref{asm-theta-agree-dictionary-models}, for the matrices $b$
and $b'$ and $\delta=1$, are (2) and (3).
% Lean: (1) is KCo.DictX, proved as KCo.dictionary_extended = ThKG.StdDict.dictX =
% dictX_of model_existence cyclicity adjoint_statement section_statement action_uniqueness
% (LmOfLinBar.lean); compare Lemma asm-theta-agree-coinvariants-inputs(3), which follows from it.
% (2) and (3): model_alt and model_Bar (VacuumImageIntertwiner.lean; they are stated and proved as
% ThK.SplitDict.model_alt and model_Bar in PfDict.lean, Lemma asm-theta-agree-dictionary-models)
% apply KCo.DictX (LmOfLinBar.lean, proved as
% ThKG.StdDict.dictX = dictX_of model_existence cyclicity adjoint_statement section_statement
% action_uniqueness) to a linear model of the second kind for (t,1,1) (for model_alt, at the
% embedding t o c, with the Gram matrices A and D of grams), change coordinates by lin_chg and
% reads_chg (A_left, A_right, B_left, B_right, gdet_rel, gdet_conj, h.unit_rel), and rescale
% the model by the nonzero constant of vacuum_image_intertwiner, respectively vac_Bar with
% delta = 1 (both from vac_ind with lat_idem, lat_to, lat_from and hTens_zero). DictX is
% stronger than Lemma asm-theta-agree-coinvariants-inputs(3) (KCo.DictSt), which does not
% name the matrices g_0 and h_0.
\end{proof}

\begin{lemma}\label{asm-theta-agree-local-ramified-values-standard}
Let $(B,G,\mathcal X,\omega,\omega',R,R_f)$ be a matched quadruple of ranks
$(a,a')$ (Definition~\ref{asm-theta-agree-matched}), with finite maps
$(\Pi_v,p_v,p'_v)_v$, let $v\in T$, let $\nu_v$, $c_v$ and $\varpi_v$ be as in
Definition~\ref{asm-theta-agree-parts}(4), and let $\bar\tau_v$ be the
$F$-embedding of $E$ into $F_v$ other than $\tau_v$. Then $c_v\geq1$, and
$\nu_{\bar\tau_v}=\nu_v^{-1}$. Let $s_A$ and $s_B$ be the functions of
Lemma~\ref{th-bottom-split-zeta} for $k=F_v$, $\varpi_v$, $\nu_v$ and $c=c_v$.
There are linear models $\iota$ and $\iota'$ as in the zeta condition at $v$
(Definition~\ref{asm-theta-agree-matched}(3)) and Schwartz-Bruhat functions
$\phi$ and $\phi'$ with $\iota\phi=s_A$, $\iota'\phi'=s_B$ and
$p_v(\phi)=p'_v(\phi')\neq0$. Moreover, let $K\subset U_V$ and $K'\subset U_{V'}$
be the lattice stabilizers of $V$ and $V'$
(Definition~\ref{asm-theta-agree-planes}(6)), with the subgroups $K_v$ and $K'_v$
of Definition~\ref{asm-theta-agree-local-cones}(2). Then $\omega(k,1)$ acts at
$v$ by $\phi\mapsto\phi$ for every $k\in K_v$, and $\omega'(k',1)$ by
$\phi'\mapsto\phi'$ for every $k'\in K'_v$.
\end{lemma}
\leanref{Development/WeilBound/ThetaLifts/LocalTheory/StandardValuations/ExistsVol.lean}{PositivityTheory.ThKG.Std.data}
\leanref{Development/WeilBound/ThetaLifts/LocalTheory/StandardValuations/ExistsVol.lean}{PositivityTheory.ThKG.Std.nu_data}
\leanref{Development/WeilBound/ThetaLifts/LocalTheory/StandardValuations/AdelicInclusionMapPlaceSection.lean}{PositivityTheory.ThKG.Std.inB_mem}
\leanref{Development/WeilBound/ThetaLifts/LocalTheory/StandardValuations/SmallFSing.lean}{PositivityTheory.ThKG.Std.omA_std}
\leanref{Development/WeilBound/ThetaLifts/LocalTheory/StandardValuations/SmallFSing.lean}{PositivityTheory.ThKG.Std.omB_std}
\begin{proof}
By Lemma~\ref{asm-theta-agree-finite-characters}, $\nu_v$ is ramified, so
$c_v\geq1$, and $\alpha_v$ and $\beta_v$ are unitary and smooth with
$\nu_v=\alpha_v/\beta_v$. As $\alpha$ and $\beta$ restrict to powers of $\eta$
on $\mathbf A^\times$ and $\eta$ is trivial at the split place $v$, the
characters of $F_v^\times$ given by $\alpha$ and $\beta$ through $\bar\tau_v$
are $\alpha_v^{-1}$ and $\beta_v^{-1}$; so $\nu_{\bar\tau_v}=\nu_v^{-1}$.

Let $\iota$ and $\iota'$ be given by the zeta condition at $v$. By
Lemma~\ref{asm-theta-agree-local-ramified-values-integrals}(2), $s_A$ and $s_B$
are Schwartz-Bruhat, so there are Schwartz-Bruhat $\phi$ and $\phi'$ with
$\iota\phi=s_A$ and $\iota'\phi'=s_B$. The triple datum is proper, so $a'<a$ and
$a+a'$ is even, and Lemma~\ref{asm-theta-agree-local-ramified-values-zeta}
gives $Z(s_A)=Z'(s_B)$; by the zeta condition, $p_v(\phi)=p'_v(\phi')$. If
$p_v(\phi)=0=p'_v(0)$, the zeta condition would give $Z(s_A)=Z'(0)=0$, which
that lemma excludes.

Let $k\in K_v$. As the elements of $U_V$ at $v$ form a subgroup, $k$ is at $v$.
Lemma~\ref{asm-theta-agree-local-ramified-values-dictionary}(2), for the Weil
pair $(V,W,M,\beta,\beta/\alpha,\omega)$, gives a linear model $\iota_0$ of the
first kind for $(\tau_v,b_{V,v},b_{W,v})$ through which $\omega(k,1)$ reads as
$\omega_A^{\nu_v,\beta_v}(g_0,1)$ with $g_0=b_{V,v}^{-1}\pi_{\tau_v}(k)b_{V,v}$,
by the first paragraph. By
Lemma~\ref{asm-theta-agree-local-ramified-values-linear}, $\iota_0$ and $\iota$
are coordinate models at $v$ with the same coordinates, so they agree on the
Schwartz-Bruhat functions by
Lemma~\ref{asm-theta-agree-local-ramified-values-evaluation}(2). As $k$ and
$k^{-1}$ preserve $\hat L_V$ and $b_{V,v}\mathcal O_v^a$ is the completion of
$L_V$ at the place given by $\tau_v$, $g_0\in GL_a(\mathcal O_v)$. Then
$|\det g_0|=1$ and $s_A(g_0^{-1}X)=\nu_v(\det g_0)s_A(X)$, so
$\omega_A^{\nu_v,\beta_v}(g_0,1)s_A=s_A$. Reading through $\iota_0$ gives a
Schwartz-Bruhat $\phi_1$ with $\omega(k,1)$ acting at $v$ by $\phi\mapsto\phi_1$
and $\iota\phi_1=\iota_0\phi_1=s_A=\iota\phi$; so $\phi_1=\phi$. The same
argument with Lemma~\ref{asm-theta-agree-local-ramified-values-dictionary}(3),
for $(V',W,M',\alpha,\alpha/\beta,\omega')$, and
$s_B(g_0^{\mathsf T}Y)=\nu_v(\det g_0)^{-1}s_B(Y)$ for
$g_0\in GL_{a'}(\mathcal O_v)$, gives the statement for $\phi'$.
\end{proof}
% Lean: data (with T.fin.zeta, local_splitting_linear_skew, lin_Bar, std_mem, std_pair,
% pRam_ne, model_alt, model_Bar, lin_uniq, uat_closure, inB_mem, omA_std, omB_std) and nu_data
% (split_character_facts, chrW_conj, cond_ex).

\begin{lemma}\label{asm-theta-agree-local-ramified-values-cone}
Let $B$ be a theta base of a CM field $E$, let $v\in P$, $t=\tau_v$ and $\bar t$
the other $F$-embedding of $E$ into $F_v$, let $\varpi$ be a uniformizer of
$F_v$ and $q$ the number of elements of its residue field. Let $M$ and $M_1$ be
regular models for $(X,W)$ and $(X_1,W_1)$, of ranks $(b+2,d)$ and $(b,d_1)$,
with coordinate models $\iota$ and $\iota_1$ at $v$ with coordinates
$(c_1,c_2)$ and $(c'_1,c'_2)$
(Definition~\ref{asm-theta-agree-local-ramified-values-coordinates}). Let
$\chi$ be a Hecke character of $E$ with $\chi_{\bar t}=\chi_t^{-1}$ such that
$\chi_t$ is trivial on the units in $1+\varpi^c\mathcal O_v$ and not on those in
$1+\varpi^{c-1}\mathcal O_v$, for an integer $c\geq1$, and let $j_c$, $e$, $e_1$,
$P$ and $r$ be as in
Lemma~\ref{asm-theta-agree-local-ramified-values-integrals}(4) for $k=F_v$,
$\nu=\chi_t$ and $\psi=\psi_v$. Let $e_v$ be a linear map from functions on
$F_v^{I_M}$ to functions on $F_v^{I_{M_1}}$ that maps the Schwartz-Bruhat
functions onto the Schwartz-Bruhat functions and has
\[
\iota_1(e_v\phi)(y)=\int_{F_v^P}\iota\phi\bigl(Z+\operatorname{ext}(y)\bigr)\,dZ
\]
for Schwartz-Bruhat $\phi$, with $j_r(i)=i+1$. Let $\omega$ be a representation
of $U_X\times U_W$ on functions on $\mathbf A^{I_M}$, $K\subset U_X$ a subgroup
with a normalized local measure system $m$
(Definition~\ref{asm-theta-agree-local-cones}(2)),
$u\colon\mathbf A_E^\times\to\mathbb W(X,W)$ a map with
$c_1(1_vu(l))=\pi_t(l)E_{0,e(b+1)}$ and $c_2(1_vu(l))=-\pi_{\bar t}(l)E_{b+1,e(0)}$
for every $l$, and $z\geq0$ an integer. Let $\phi$ and $\phi_1$ be
Schwartz-Bruhat with $\iota\phi=s_e$ and $\iota_1\phi_1=s_{e_1}$, such that
$\omega(k,1)$ acts at $v$ by $\phi\mapsto\phi$ for every $k\in K_v$. Then the
local cone of $\phi$ for $(M,\omega,m,e_v,u,\chi,z)$
(Definition~\ref{asm-theta-agree-local-cones}(3)) is
\[
\operatorname{cone}_v(\phi)=G(\varpi^{-c})\,\chi_t(\varpi)^c\,q^{cz}\,q^{-cr}\,\phi_1,
\]
with $G$ as in Lemma~\ref{asm-theta-agree-local-ramified-values-integrals} for
$\nu=\chi_t$ and $\psi=\psi_v$.
\end{lemma}
\leanref{Development/WeilBound/ThetaLifts/LocalTheory/StandardValuations/ConeVal.lean}{PositivityTheory.ThKG.Std.cone_val}
\begin{proof}
The character $\psi_v$ has conductor $\mathcal O_v$, as $v\in P$, so
Lemma~\ref{asm-theta-agree-local-ramified-values-integrals} applies. Let
$k\in K_v$ and $l\in A_v$, and put $u_l=1_vu(l)$, which is at $v$. As
$\omega(k,1)$ acts at $v$ by $\phi\mapsto\phi$, the operator
$T^v_{k,l}=\rho(u_l)\omega(k,1)$ acts at $v$ by $\phi\mapsto\phi_l$, where
$\rho(u_l)$ acts by $\phi\mapsto\phi_l$; so $(T^v_{k,l})_v\phi=\phi_l$ by
Lemma~\ref{asm-theta-agree-finite-unique}. As $c_1(u_l)$ and $c_2(u_l)$ are
supported at the positions $(0,e(b+1))$ and $(b+1,e(0))$, which differ,
$c_1(u_l)\cdot c_2(u_l)=0$, and
\[
\iota\phi_l(Y)=\psi_v\bigl(-\pi_{\bar t}(l)Y_{b+1,e(0)}\bigr)\,s_e\bigl(Y+\pi_t(l)E_{0,e(b+1)}\bigr).
\]
The matrix $\operatorname{ext}(y)$ vanishes at $(b+1,e(0))$, so
\[
\iota_1(e_v\phi_l)(y)=\int_{F_v^P}\psi_v\bigl(-\pi_{\bar t}(l)Z_{b+1,e(0)}\bigr)
s_e\bigl(Z+\operatorname{ext}(y)+\pi_t(l)E_{0,e(b+1)}\bigr)\,dZ
=q^{-cr}\delta(\pi_t(l))\,G(\pi_{\bar t}(l))\,s_{e_1}(y)
\]
by Lemma~\ref{asm-theta-agree-local-ramified-values-integrals}(4). As
$\iota_1$ is injective on the Schwartz-Bruhat functions,
$e_v\phi_l=q^{-cr}\delta(\pi_t(l))G(\pi_{\bar t}(l))\phi_1$. So the integrand
of the local cone at $(k,l)$ is
$q^{-cr}\delta(\pi_t(l))G(\pi_{\bar t}(l))\phi_1(y)\chi(l)|l|^z$. Its integral
over $A_v$ is $G(\varpi^{-c})\chi_t(\varpi)^cq^{cz}q^{-cr}\phi_1(y)$ by
Lemma~\ref{asm-theta-agree-local-ramified-values-integrals}(5), and $K_v$ has
total mass $1$.
\end{proof}
% Lean: cone_val (key: opF_eq, acts_Hv, Hv_mem, Lin.act, int_cone; then idele_int and the
% mass of m.kf v). The map u enters through the hypothesis hs on its coordinates.

\begin{lemma}\label{asm-theta-agree-local-ramified-values}
Let $B$ be a theta base of a CM field $E$.
\begin{enumerate}
\item There is a sign $s\in\{\pm1\}$ such that the following holds. Let
  $(G,G_1,H,H')$ be an $\alpha$-configuration with $G$ of ranks $(b+2,a')$,
  $e$ an evaluation map for $(M,M_1,\iota_H,\iota_{H'},\mathcal J_\alpha(H,H'))$
  with components $e_\infty$ and $e_v$, $e'$ one for
  $(M',M'_1,1,\iota_{H'},\mathcal J_0(H'))$ with components $e'_\infty$ and
  $e'_v$, $K\subset U_V$ the lattice stabilizer of $V$, and $m$ a normalized
  local measure system for $K$. Let $(B,G,\mathcal X,\omega,\omega',R,R_f)$ and
  $(B,G_1,\mathcal X,\omega_1,\omega'_1,R_1,R_{f,1})$ be matched quadruples
  (Definition~\ref{asm-theta-agree-matched}), with finite maps
  $(\Pi_v,p_v,p'_v)_v$ and $(\Pi_{1,v},p_{1,v},p'_{1,v})_v$, and let $v\in T$.
  Then there are $\phi\in\mathcal S(F_v^{I_M})$ and
  $\phi'\in\mathcal S(F_v^{I_{M'}})$ with $p_v(\phi)=p'_v(\phi')$,
  $p'_{1,v}(e'_v\phi')\neq0$ and
  \[
  p_{1,v}(\operatorname{cone}_v(\phi))=\nu_v(s)\,\epsilon(\nu_v,1-k_\alpha)\,p'_{1,v}(e'_v\phi'),
  \]
  where $\operatorname{cone}_v$ is the local cone for
  $(M,\omega,m,e_v,u_\alpha,\nu,k_\alpha)$ and $\nu_v$ and
  $\epsilon(\nu_v,z)$ are as in Definition~\ref{asm-theta-agree-parts}(4).
\item There is a sign $s$ such that the following holds. Let $(G,G_1,H,H')$
  be a $\beta$-configuration with $G$ of ranks $(a,b+2)$, $e$ an evaluation
  map for $(M,M_1,1,\iota_{H'},\mathcal J_0(H'))$ with components $e_\infty$
  and $e_v$, $e'$ one for $(M',M'_1,\iota_H,\iota_{H'},\mathcal J_\beta(H,H'))$
  with components $e'_\infty$ and $e'_v$, $K\subset U_{V'}$ the lattice
  stabilizer of $V'$, and $m$ a normalized local measure system for $K$. Let
  the matched quadruples and $v\in T$ be as in (1). Then there are
  $\phi\in\mathcal S(F_v^{I_M})$ and $\phi'\in\mathcal S(F_v^{I_{M'}})$ with
  $p_v(\phi)=p'_v(\phi')$, $p_{1,v}(e_v\phi)\neq0$ and
  \[
  p'_{1,v}(\operatorname{cone}_v(\phi'))=\nu_v(s)\,\epsilon(\nu_v,-k_\beta)\,p_{1,v}(e_v\phi),
  \]
  where $\operatorname{cone}_v$ is the local cone for
  $(M',\omega',m,e'_v,u_\beta,\nu,k_\beta+1)$.
\end{enumerate}
\end{lemma}
\leanref{Development/WeilBound/ThetaLifts/LocalTheory/UnramifiedDeterminedBySpherical.lean}{PositivityTheory.ThKG.Loc.StdA}
\leanref{Development/WeilBound/ThetaLifts/LocalTheory/UnramifiedDeterminedBySpherical.lean}{PositivityTheory.ThKG.Loc.StdB}
\leanref{Development/WeilBound/ThetaLifts/LocalTheory/StandardValuations/ExistsVol.lean}{PositivityTheory.ThKG.Std.stdA}
\leanref{Development/WeilBound/ThetaLifts/LocalTheory/StandardValuations/ExistsVol.lean}{PositivityTheory.ThKG.Std.stdB}
\begin{proof}
In both parts take $s=1$, so that $\nu_v(s)=1$. Let $v\in T$, write $Q$ and $Q_1$
for the two matched quadruples, $t=\tau_v$, $\bar t$ for the other
$F$-embedding of $E$ into $F_v$, $c=c_v$, $\varpi=\varpi_v$ and $q$ for the
number of elements of the residue field of $F_v$. Read the functions $s_e$,
the matrices $\operatorname{ext}(y)$ and the measures $dZ$ as in
Lemma~\ref{asm-theta-agree-local-ramified-values-integrals} for $k=F_v$,
$\nu=\nu_v$ and $\psi=\psi_v$, and let $e_A$ and $e_B$ be the maps onto the
first and the last columns, so that the functions $s_A$ and $s_B$ of
Lemma~\ref{th-bottom-split-zeta} are $s_{e_A}$ and $s_{e_B}$.

(1) By Lemma~\ref{asm-theta-agree-local-ramified-values-standard} for $Q$ there
are linear models $\iota$ and $\iota'$ as in the zeta condition and
Schwartz-Bruhat $\phi$ and $\phi'$ with $\iota\phi=s_{e_A}$,
$\iota'\phi'=s_{e_B}$, $p_v(\phi)=p'_v(\phi')$ and $\phi$ fixed by $\omega(k,1)$
at $v$ for $k\in K_v$; moreover $c\geq1$ and $\nu_{\bar t}=\nu_v^{-1}$. The same
lemma for $Q_1$ gives $\iota_1$, $\iota'_1$, $\phi_1$ and $\phi'_1$ with
$\iota_1\phi_1=s_{e_A}$, $\iota'_1\phi'_1=s_{e_B}$ (now for matrices with $b$,
respectively $a'$, rows and $b+a'$ columns) and
$p_{1,v}(\phi_1)=p'_{1,v}(\phi'_1)\neq0$. By
Lemma~\ref{asm-theta-agree-local-ramified-values-linear}, $\iota$ and $\iota_1$
are coordinate models at $v$ with the coordinates $(c_W,c_B)$ for
$(t,b_{V,v},b_{W,v})$ and $(t,b_{V_1,v},b_{W_1,v})$, and $\iota'$ and
$\iota'_1$ are coordinate models with the coordinates $(c_B,-c_W)$ for
$(t,b_{V',v},b_{W,v})$ and $(t,b_{V',v},b_{W_1,v})$, as $V'_1=V'$. Let $j$ be
the increasing map from $\{0,\ldots,b+a'-1\}$ to $\{0,\ldots,b+a'+1\}$ that
omits $0$ and $b+1$.

\emph{Step 1: the evaluation of $\phi'$.} By
Lemma~\ref{asm-theta-agree-local-ramified-values-evaluation}(3),
$\iota'_1(e'_v\phi')(y)$ is the integral of $s_{e_B}(Z+\operatorname{ext}(y))$
over the column $b+1$, with $j_r$ the identity and $j_c=j$. The column $b+1$ lies
outside the image of $e_B$, which consists of the columns $b+2,\ldots,b+a'+1$,
and outside that of $j$, and $j$ carries the last $a'$ columns of $b+a'$ to
those of $b+a'+2$. So Lemma~\ref{asm-theta-agree-local-ramified-values-integrals}(3)
gives $\iota'_1(e'_v\phi')=q^{-ca'}s_{e_B}=\iota'_1(q^{-ca'}\phi'_1)$. As $e'_v$
preserves Schwartz-Bruhat functions and $\iota'_1$ is injective on them,
$e'_v\phi'=q^{-ca'}\phi'_1$. Hence
\[
p'_{1,v}(e'_v\phi')=q^{-ca'}p'_{1,v}(\phi'_1)=q^{-ca'}p_{1,v}(\phi_1)\neq0.
\]

\emph{Step 2: the cone.} The plane $H=(x,y,\ldots)$ is at the positions $0$
and $b+1$ of $V$ and $H'=(x',y',\ldots)$ at the positions $0$ and $b+1$ of $W$.
The identities for $x$, $y$, $x'$ and $y'$ in the proof of
Lemma~\ref{asm-theta-agree-local-ramified-values-evaluation}(3) give, for
$u_\alpha(l)=l\,xy'^{\mathsf T}$,
\[
c_W(1_vu_\alpha(l))=\pi_t(l)E_{0,b+1},\qquad c_B(1_vu_\alpha(l))=-\pi_{\bar t}(l)E_{b+1,0}.
\]
Apply Lemma~\ref{asm-theta-agree-local-ramified-values-cone} with $\chi=\nu$,
$e=e_A$, $e_1$ the map onto the first $b$ columns, $j_c=j$, $z=k_\alpha$ and $P$
the set of positions in the column $0$ or in the row $0$ outside the column
$b+1$. Its hypotheses hold: $e_A(0)=0$ and $e_A(b+1)=b+1$ are not in the image of
$j$, and $j(e_1(i))=e_A(i+1)$; $P$ contains $(b+1,0)$ and neither $(0,b+1)$ nor
any $(i+1,j(l))$; $P$ contains no $(i,b+1)$ with $i\neq0$ and no
$(b+1,i+1)$, which is the second alternative of
Lemma~\ref{asm-theta-agree-local-ramified-values-integrals}(4); the evaluation
$e_v$ has the required form by
Lemma~\ref{asm-theta-agree-local-ramified-values-evaluation}(3); and $\phi$ is
fixed by $\omega(k,1)$ for $k\in K_v$. The positions of $P$ whose column is not
in the image of $e_A$ are those of the row $0$ in the columns
$b+2,\ldots,b+a'+1$, so $r=a'$. Hence
\[
\operatorname{cone}_v(\phi)=G(\varpi^{-c})\,\nu_v(\varpi)^c\,q^{ck_\alpha}\,q^{-ca'}\,\phi_1.
\]

\emph{Step 3: the epsilon factor.} By
Lemma~\ref{asm-theta-agree-local-ramified-values-integrals}(1) and
Definition~\ref{asm-theta-agree-parts}(4),
$G(\varpi^{-c})\nu_v(\varpi)^cq^{ck_\alpha}=q^{-c(1-k_\alpha)}\mathfrak g(\nu_v)\nu_v(\varpi)^c=\epsilon(\nu_v,1-k_\alpha)$.
With Step 1,
\[
p_{1,v}(\operatorname{cone}_v(\phi))=\epsilon(\nu_v,1-k_\alpha)\,q^{-ca'}p_{1,v}(\phi_1)
=\epsilon(\nu_v,1-k_\alpha)\,p'_{1,v}(e'_v\phi'),
\]
and $p_v(\phi)=p'_v(\phi')$ and $p'_{1,v}(e'_v\phi')\neq0$.

(2) The proof is that of (1) with the roles of the two sides exchanged.
Lemma~\ref{asm-theta-agree-local-ramified-values-standard} gives $\iota$,
$\iota'$, $\phi$, $\phi'$ for $Q$, now with $\phi'$ fixed by $\omega'(k,1)$ for
$k\in K_v$, where $K$ is the lattice stabilizer of $V'$, and $\iota_1$,
$\iota'_1$, $\phi_1$, $\phi'_1$ for $Q_1$, with $V_1=V$. Here $H$ is at the
positions $0$ and $b+1$ of $V'$, $H'$ at the positions $a$ and $a+b+1$ of $W$,
and $j$ omits $a$ and $a+b+1$. By
Lemma~\ref{asm-theta-agree-local-ramified-values-evaluation}(4), $e_v$ is the
integral over the column $a$, which lies outside the images of $e_A$ and $j$; so
$e_v\phi=q^{-ca}\phi_1$ by
Lemma~\ref{asm-theta-agree-local-ramified-values-integrals}(3), and
$p_{1,v}(e_v\phi)=q^{-ca}p_{1,v}(\phi_1)\neq0$. In the coordinates
$(c_B,-c_W)$ of $\iota'$, $u_\beta(l)=c(l)\,yx'^{\mathsf T}$ has
$c_B(1_vu_\beta(l))=\pi_t(l)E_{0,a+b+1}$ and
$-c_W(1_vu_\beta(l))=-\pi_{\bar t}(l)E_{b+1,a}$. Apply
Lemma~\ref{asm-theta-agree-local-ramified-values-cone} to $M'$, $M'_1$,
$\iota'$, $\iota'_1$, $\omega'$, $K$, $e=e_B$ (so $e(0)=a$ and
$e(b+1)=a+b+1$), $e_1$ the map onto the last $b$ columns, $z=k_\beta+1$, and $P$
the set of positions in the row $b+1$ or in the column $a+b+1$ outside the row
$0$; now the first alternative of
Lemma~\ref{asm-theta-agree-local-ramified-values-integrals}(4) holds, and the
positions of $P$ outside the image of $e_B$ are those of the row $b+1$ in the
columns $0,\ldots,a-1$, so $r=a$. This gives
$\operatorname{cone}_v(\phi')=G(\varpi^{-c})\nu_v(\varpi)^cq^{c(k_\beta+1)}q^{-ca}\phi'_1$,
and as in Step 3,
$G(\varpi^{-c})\nu_v(\varpi)^cq^{c(k_\beta+1)}=\epsilon(\nu_v,-k_\beta)$ and
\[
p'_{1,v}(\operatorname{cone}_v(\phi'))=\epsilon(\nu_v,-k_\beta)\,q^{-ca}p'_{1,v}(\phi'_1)
=\epsilon(\nu_v,-k_\beta)\,p_{1,v}(e_v\phi).
\]
\end{proof}
% Lean: StdA is Std.stdA = stdA_gen with s = 1, and StdB is Std.stdB with s = 1
% (StandardValuations/ExistsVol.lean). Step 1 is e1 (ev0B, int_ev); Step 2 is e2 (cone_val with
% evA, shiftA_c and the second alternative); Step 3 is eps_eq with mt1 and ne1. In stdB, e1 is
% ev0W with int_ev, and e2 is cone_val with evB, shiftB_c and the first alternative. What
% remains open is Lemma asm-theta-agree-local-ramified-values-dictionary.

\begin{lemma}\label{asm-theta-agree-local-bump-schwartz}
Let $V$ and $V_1$ be finite-dimensional real inner product spaces with Lebesgue
measure, let $\mathcal S(V)$ be the Schwartz space of $V$, with the seminorms
$\|f\|_{k,j}=\sup_x|x|^k|D^jf(x)|$, and put $e(s)=e^{2\pi is}$ for
$s\in\mathbf R$.
\begin{enumerate}
\item For every $\varphi\in\mathcal S(V)$ there is $b\in\mathcal S(V)$ with
  $\varphi(x)=\int_Ve(\langle\eta,x\rangle)b(\eta)\,d\eta$ for every $x$.
\item Let $\Omega$ be a second countable space with a Borel measure $\mu$, and
  $G\colon\Omega\to\mathcal S(V)$ a continuous map such that every function
  $a\mapsto\|G(a)\|_{k,j}$ is $\mu$-integrable. Then there is
  $f\in\mathcal S(V)$ such that for every continuous linear form $T$ on
  $\mathcal S(V)$ the function $T\circ G$ is integrable and
  $T(f)=\int T(G(a))\,d\mu(a)$.
\item Let $(T_i)_{i\in I}$ be continuous linear maps from $\mathcal S(V)$ to
  $\mathcal S(V_1)$ and $q$ a continuous seminorm on $\mathcal S(V_1)$ such
  that $q(T_if)$ is bounded in $i$ for every $f$. Then there is a continuous
  seminorm $p$ on $\mathcal S(V)$ with $q(T_if)\leq p(f)$ for all $i$ and $f$.
  If $(g_i)_{i\in I}$ are elements of $\mathcal S(V)$ and $T(g_i)$ is bounded in
  $i$ for every continuous linear form $T$, then $q(g_i)$ is bounded in $i$ for
  every continuous seminorm $q$.
\item Let $K$ be a topological space and $g\colon K\to\mathcal S(V)$ a map such
  that $q\circ g$ is bounded for every continuous seminorm $q$ and $T\circ g$
  is continuous for every continuous linear form $T$. Then $g$ is continuous.
\item Let $\Omega$ be a measure space with measure $\mu$, and $G$ a map from
  $\Omega$ to the continuous linear maps from $\mathcal S(V)$ to
  $\mathcal S(V_1)$. Suppose that for every $f$ there is $g\in\mathcal S(V_1)$
  such that, for every continuous linear form $T$, $a\mapsto T(G(a)f)$ is
  integrable with integral $T(g)$; and that for all $k,j$ there are a
  continuous seminorm $p$ and an integrable function $h$ with
  $\|G(a)f\|_{k,j}\leq p(f)h(a)$ for all $a$ and $f$. Then there is a
  continuous linear map $L\colon\mathcal S(V)\to\mathcal S(V_1)$ with
  $T(L(f))=\int T(G(a)f)\,d\mu(a)$ for all $f$ and $T$.
\end{enumerate}
\end{lemma}
\leanref{Development/WeilBound/ThetaLifts/LocalTheory/ArchimedeanBumps/ExistsDensity.lean}{PositivityTheory.ThKG.Loc.Bump.exists_density}
\leanref{Development/WeilBound/ThetaLifts/SmallCones/ThetaUniq.lean}{PositivityTheory.ThKCT.SchInt}
\leanref{Development/WeilBound/ThetaLifts/SmallCones/SchInt.lean}{PositivityTheory.ThKCT.schInt}
\leanref{Development/WeilBound/ThetaLifts/LocalTheory/ArchimedeanBumps/SliceAn.lean}{PositivityTheory.ThKG.Loc.Bump.equicont}
\leanref{Development/WeilBound/ThetaLifts/LocalTheory/ArchimedeanBumps/SliceAn.lean}{PositivityTheory.ThKG.Loc.Bump.weak_bounded}
\leanref{Development/WeilBound/ThetaLifts/LocalTheory/ArchimedeanBumps/SliceAn.lean}{PositivityTheory.ThKG.Loc.Bump.strong_of_weak}
\leanref{Development/WeilBound/ThetaLifts/LocalTheory/ArchimedeanBumps/SliceAn.lean}{PositivityTheory.ThKG.Loc.Bump.weak_op}
\begin{proof}
(1) is the Fourier inversion formula on $\mathcal S(V)$
\cite[Section~7.1]{Hor90}, with $b$ the Fourier transform of $\varphi$.

(2) Put $f(x)=\int G(a)(x)\,d\mu(a)$. As $|D^jG(a)(x)|\leq\|G(a)\|_{0,j}$ and
these bounds are integrable, differentiation under the integral sign shows
that $f$ is smooth with $D^jf(x)=\int D^jG(a)(x)\,d\mu(a)$, and
$\|f\|_{k,j}\leq\int\|G(a)\|_{k,j}\,d\mu(a)$; so $f\in\mathcal S(V)$. A
continuous linear form $T$ has $|T(g)|\leq C\max_{(k,j)\in s}\|g\|_{k,j}$ for a
finite set $s$. So, by the Hahn-Banach theorem, $T=T'\circ J$, where $J$ maps
$g$ to $(x\mapsto|x|^kD^jg(x))_{(k,j)\in s}$, a vector in a finite product of
Banach spaces of bounded continuous functions, and $T'$ is continuous and
linear. The map $a\mapsto J(G(a))$ is Bochner integrable, its integral is
$J(f)$ by the formula for $D^jf$, and $T'$ commutes with the Bochner integral.

(3) The first claim is the Banach-Steinhaus theorem for the Fréchet space
$\mathcal S(V)$ \cite[Theorem~2.6]{Rud91}, applied to the maps $T_i$ into the
seminormed space $(\mathcal S(V_1),q)$. The second claim is the boundedness of
weakly bounded sets in a locally convex space \cite[Theorem~3.18]{Rud91}.

(4) For $x\in V$ and $n\geq0$, $k\mapsto D^ng(k)(x)$ is continuous, since its
coordinates are continuous linear forms of $g(k)$. Let $k_0\in K$, $j,n\geq0$ and
$\varepsilon>0$, and let $C_1$ and $C_2$ bound $\|g(k)\|_{j+1,n}$ and
$\|g(k)\|_{0,n+1}$. For $|x|\geq R=4C_1/\varepsilon+1$,
$|x|^j|D^n(g(k)-g(k_0))(x)|\leq2C_1/|x|<\varepsilon/2$. On the ball of radius
$R$, $D^n(g(k)-g(k_0))$ is $2C_2$-Lipschitz; so its smallness at the finitely
many points of a fine net of the ball, which holds for $k$ near $k_0$, gives
$|x|^j|D^n(g(k)-g(k_0))(x)|<\varepsilon/2$ on the ball. So
$\|g(k)-g(k_0)\|_{j,n}<\varepsilon$ for $k$ near $k_0$.

(5) The $g$ of the hypothesis is unique, as continuous linear forms separate
the points of $\mathcal S(V_1)$; so $L(f)=g$ defines a linear map. For $f$ and
$(k,j)$, the Hahn-Banach theorem gives a continuous linear form $T$ with
$|T(L(f))|=\|L(f)\|_{k,j}$ and $|T(g')|\leq\|g'\|_{k,j}$ for all $g'$. Then
$\|L(f)\|_{k,j}\leq\int\|G(a)f\|_{k,j}\,d\mu(a)\leq p(f)\int h\,d\mu$, so $L$ is
continuous.
\end{proof}
% Lean: (1) exists_density takes the Fourier transform with fourierInv_fourier_eq. (2) is the
% property SchInt, which schInt (SmallCones/SchInt.lean) proves by exists_int (taylor, the
% derivatives under the integral) and factor (the Hahn-Banach factorization); WeakInt
% (DictIrr.lean) has the same statement. (3) is equicont (Baire property of the Schwartz space
% and Seminorm.continuous_iSup) and weak_bounded (banach_steinhaus with the norming forms of
% norming). (4) is strong_of_weak, with the net of finite_cover_balls_of_compact. (5) is weak_op,
% with eq_of_weak and norming'.

\begin{lemma}\label{asm-theta-agree-local-bump-positivity}
Let $S$ be a finite set. For $x,t\in\mathbf C^S$ put
$[x,t]=\sum_w2\operatorname{Re}(x_wt_w)$ and $N(t)=\prod_w|t_w|^2$, write $|t|$
for the largest $|t_w|$, and for an integrable function $h$ on $\mathbf C^S$ put
$\hat h(t)=\int e(-[x,t])h(x)\,dx$, with $e(s)=e^{2\pi is}$ and $dx$ the
Lebesgue measure. Let $\kappa\colon S\to\mathbf Z_{\geq0}$ and $z\geq0$, and put
$r(x)=\prod_w(x_w/|x_w|)^{\kappa(w)}|x_w|^{2z}$ for $x$ with nonzero
components.
\begin{enumerate}
\item If $g$ is a continuous integrable function on $\mathbf C^S$ with
  $g(0)\neq0$, then $\hat g(t)\neq0$ for some $t$.
\item If $\kappa(w)>2z$ for every $w$, there is $Z\in\mathbf C$ with
  $\int\hat H(t)r(t)^{-1}N(t)^{-1}\,dt=Z\int H(x)r(x)\,dx$ for every integrable
  $H$ such that $Hr$ and $\hat Hr^{-1}N^{-1}$ are integrable.
\end{enumerate}
Let $K_0$ be a compact Hausdorff space with a finite Borel measure $\mu_0$
that is positive on nonempty open sets, let $\xi\colon K_0\to(\mathbf C^n)^S$ be
continuous with $\xi(k)_w\neq0$ for all $k$ and $w$, and put
$\pi(k,c)=(2c_w\xi(k)_w)_w$ for $c\in\mathbf C^S$, in the real inner product
space $(\mathbf C^n)^S$ with
$\langle\eta,y\rangle=\operatorname{Re}\sum_{w,i}\overline{\eta_{w,i}}y_{w,i}$.
Let $a$ be a continuous function on $K_0\times\mathbf C^S$ such that
$|a(k,c)|(1+|c|)^N$ is bounded for every $N$, put $a_k=a(k,\cdot)$, and let
$k_1\in K_0$ and $c_1\in\mathbf C$.
\begin{enumerate}
\item[(3)] If $\operatorname{Re}(a(k,c)r(c)c_1)>0$ for every $(k,c)$ with
  $\pi(k,c)=\pi(k_1,1)$, then
  \[
  \int_{K_0}\int e(\langle\eta,\pi(k,c)\rangle)\,a(k,c)\,r(c)\,dc\,d\mu_0(k)\neq0
  \]
  for some $\eta\in(\mathbf C^n)^S$.
\item[(4)] Let $K_0$ be first countable, let $\kappa(w)>2z$ for every $w$, and
  put $\beta(k,\eta)_w=\sum_i\xi(k)_{w,i}\overline{\eta_{w,i}}$. If
  $\operatorname{Re}(\hat a_k(c)r(c)^{-1}N(c)^{-1}c_1)>0$ for every $(k,c)$ with
  $\pi(k,c)=\pi(k_1,1)$, then
  \[
  \int_{K_0}\int a(k,c+\beta(k,\eta))\,r(c)\,dc\,d\mu_0(k)\neq0
  \]
  for some $\eta\in(\mathbf C^n)^S$.
\end{enumerate}
\end{lemma}
\leanref{Development/WeilBound/ThetaLifts/LocalTheory/ArchimedeanBumps/RationalWeightSuccMulInvNorm.lean}{PositivityTheory.ThKG.Loc.Bump.exists_FT_ne}
\leanref{Development/WeilBound/ThetaLifts/LocalTheory/ArchimedeanBumps/IntegrableConjPowGaussian.lean}{PositivityTheory.ThKG.Loc.Bump.tate}
\leanref{Development/WeilBound/ThetaLifts/LocalTheory/ArchimedeanBumps/ContinuousOnRationalWeight.lean}{PositivityTheory.ThKG.Loc.Bump.coreA}
\leanref{Development/WeilBound/ThetaLifts/LocalTheory/ArchimedeanBumps/ContinuousOnRationalWeight.lean}{PositivityTheory.ThKG.Loc.Bump.coreB}
\leanref{Development/WeilBound/ThetaLifts/LocalTheory/ArchimedeanBumps/RationalWeightSuccMulInvNorm.lean}{PositivityTheory.ThKG.Loc.Bump.prp_pim}
\begin{proof}
Let $m_0(x)=(1+|x|)^{-2|S|-1}$, an integrable function on $\mathbf C^S$. As
$1+|x|\leq(1+|x+b|)(1+|b|)$ and $|r(x)|\leq(1+|x|)^{2z|S|}$, the bounds on $a$
give, for every $M$, constants $C$ and $N'$ with
$|a(k,x+b)|(1+|x|)^M\leq C(1+|b|)^{N'}m_0(x)$ for all $k$, $x$ and $b$.

\emph{A bump.} Let $\Omega$ be a Hausdorff space with a Borel measure $\mu$
that is finite on compact sets and positive on nonempty open sets, $Y$ a
finite-dimensional real inner product space, $\pi_0\colon\Omega\to Y$
continuous, $y_0\in Y$, $U\subset\Omega$ open, $G$ continuous on $U$,
$S_1\subset U$ compact and $\varepsilon>0$, such that every $p$ with
$|\pi_0(p)-y_0|\leq\varepsilon$ lies in $S_1$, $\operatorname{Re}G(p)>0$
whenever $\pi_0(p)=y_0$, and $\pi_0(p_0)=y_0$ for some $p_0$. Then there is
$\varphi\in\mathcal S(Y)$, vanishing at $\pi_0(p)$ for $p\notin S_1$, such that
$p\mapsto\varphi(\pi_0(p))G(p)$ is integrable with nonzero integral. Indeed, the
$p\in S_1$ with $\operatorname{Re}G(p)\leq0$ form a compact set, whose image
under $\pi_0$ is compact and does not contain $y_0$; so it misses the ball of
some radius $\delta>0$ around $y_0$. Let $\varphi$ be smooth with values in
$[0,1]$, equal to $1$ at $y_0$ and supported in the ball of radius
$\min(\delta,\varepsilon)$ around $y_0$. Then $\varphi(\pi_0(p))G(p)$ is
continuous with compact support in $S_1$, its real part is $\geq0$ and is
positive at $p_0$, so the real part of its integral is positive.

\emph{Fourier vanishing.} Let $\Omega$, $\mu$, $Y$ and $\pi_0$ be as above and
$g$ integrable on $\Omega$ with $\int e(\langle\eta,\pi_0(p)\rangle)g(p)\,d\mu(p)=0$
for every $\eta\in Y$. Then $\int\varphi(\pi_0(p))g(p)\,d\mu(p)=0$ for every
$\varphi\in\mathcal S(Y)$: write
$\varphi(y)=\int e(\langle\eta,y\rangle)b(\eta)\,d\eta$ by
Lemma~\ref{asm-theta-agree-local-bump-schwartz}(1), and exchange the integrals,
as $g$ and $b$ are integrable.

(1) Suppose $\hat g=0$. Take $\Omega=Y=\mathbf C^S$, $\pi_0(x)=2x$, $y_0=0$,
$U=\Omega$, $G=g/g(0)$, $S_1$ the set of $x$ with $|x|\leq1$, $\varepsilon=1$
and $p_0=0$; then $\pi_0(x)=0$ only for $x=0$. As
$\langle\eta,2x\rangle=-[x,-\bar\eta]$, every
$\int e(\langle\eta,2x\rangle)g(x)\,dx=\hat g(-\bar\eta)$ is $0$, so Fourier
vanishing contradicts the bump.

(2) Put $g_\kappa(y)=\prod_w\bar y_w^{\kappa(w)}e^{-2\pi|y_w|^2}$. The Fourier
transform of $\bar y^ke^{-2\pi|y|^2}$ on $\mathbf C$ gives
$\hat g_\kappa(x)=\prod_w\tfrac12(-i)^{\kappa(w)}x_w^{\kappa(w)}e^{-2\pi|x_w|^2}$.
Put $j_1=g_\kappa r$ and $j_2(y)=\hat g_\kappa(y^{-1})r(y)N(y)^{-1}$, where
$y^{-1}=(y_w^{-1})_w$. Then
$j_1(y)=\prod_w|y_w|^{\kappa(w)+2z}e^{-2\pi|y_w|^2}$ is integrable with positive
integral, and $j_2$ is integrable, as $\kappa(w)>2z$ makes
$|y_w|^{2z-\kappa(w)-2}e^{-2\pi|y_w|^{-2}}$ integrable. Put
$Z=\int j_2/\int j_1$. For $y$ with nonzero components, $x\mapsto yx$ multiplies
the Lebesgue measure by $N(y)$, and the Fourier transform of $t\mapsto m(yt)$ is
$x\mapsto N(y)^{-1}\hat m(xy^{-1})$. Also $r$ and $N$ are multiplicative, and
$\int\hat hm=\int h\hat m$ for integrable $h$ and $m$. Compute
$I=\int\int\hat H(t)r(t)^{-1}j_1(ty)\,dt\,dy$ in two ways. Integrating over $y$
first gives $I=\int j_1\cdot\int\hat Hr^{-1}N^{-1}$. For fixed $y$,
$r(t)^{-1}j_1(ty)=r(y)g_\kappa(yt)$, so the inner integral is
\[
r(y)\int H(x)\,N(y)^{-1}\hat g_\kappa(xy^{-1})\,dx=\int H(x)r(x)N(x)^{-1}j_2(x^{-1}y)\,dx,
\]
and integrating over $y$ first gives $I=\int j_2\cdot\int Hr$. Fubini's theorem
applies in both cases, by the integrability of $\hat Hr^{-1}N^{-1}$, $Hr$, $j_1$
and $j_2$.

(3) Let $U$ be the open set of the $(k,c)$ with every $c_w\neq0$, and
$G(k,c)=a(k,c)r(c)c_1$, which is continuous on $U$. On $\mathbf C^n$ use the
largest absolute value of the coordinates. Let $m>0$ be a lower bound of the
$|\xi(k)_w|$, which exists as $K_0$ is compact, and $S_1$ the set of $(k,c)$
with $|c_w||\xi(k)_w|\geq m/2$ for every $w$ and $|c|\leq(|\xi(k_1)|+m)/m$; it is
compact and contained in $U$. If $|\pi(k,c)-\pi(k_1,1)|\leq m$, then
$|c_w\xi(k)_w-\xi(k_1)_w|\leq m/2$, so $|c_w||\xi(k)_w|\geq m/2$ and
$m|c_w|\leq|\xi(k_1)_w|+m/2$, for every $w$; so $(k,c)\in S_1$. The bump, for
$\Omega=K_0\times\mathbf C^S$ with $\mu_0\times dc$, $Y=(\mathbf C^n)^S$,
$y_0=\pi(k_1,1)$, $\varepsilon=m$ and $p_0=(k_1,1)$, gives $\varphi$ with
$\int\int\varphi(\pi(k,c))a(k,c)r(c)c_1\,dc\,d\mu_0\neq0$. The function $ar$
is integrable on $K_0\times\mathbf C^S$, by the bound above with $M=2z|S|$ and
$b=0$. If the integrals in (3) were $0$ for every $\eta$, Fourier vanishing for
$g=arc_1$ would make this integral $0$.

(4) The function $(k,t)\mapsto\hat a_k(t)$ is continuous, by dominated
convergence with the bound $|a(k,x)|\leq Cm_0(x)$. Let $U$, $m$ and $S_1$ be as
in (3); the bump for $G(k,t)=\hat a_k(t)r(t)^{-1}N(t)^{-1}c_1$, which is
continuous on $U$, gives $\varphi$ with
$\int\int\varphi(\pi(k,t))G(k,t)\,dt\,d\mu_0\neq0$, the integrand being
integrable. By Lemma~\ref{asm-theta-agree-local-bump-schwartz}(1),
$\varphi(y)=\int e(\langle\eta,y\rangle)b(\eta)\,d\eta$ with $b$ Schwartz. As
$K_0$ is compact, $|\beta(k,\eta)|\leq D|\eta|$ for some $D$, so
$|b(\eta)a(k,x+\beta(k,\eta))r(x)|\leq C(1+D)^{N'}(1+|\eta|)^{N'}|b(\eta)|m_0(x)$
for suitable $C$ and $N'$, an integrable bound. Put
$H_k(x)=\int b(\eta)a(k,x+\beta(k,\eta))\,d\eta$. As
$[\beta(k,\eta),t]=\langle\eta,\pi(k,t)\rangle$ and
$\int e(-[x,t])a_k(x+b')\,dx=e([b',t])\hat a_k(t)$, Fubini's theorem gives
$\hat H_k(t)=\varphi(\pi(k,t))\hat a_k(t)$. For almost every $k$, (2) applies
to $H_k$, and
\[
\int\int\varphi(\pi(k,t))G(k,t)\,dt\,d\mu_0(k)=Zc_1\int\int H_k(x)r(x)\,dx\,d\mu_0(k)
=Zc_1\int b(\eta)\int_{K_0}\int a(k,x+\beta(k,\eta))r(x)\,dx\,d\mu_0(k)\,d\eta.
\]
If the integrals in (4) were $0$ for every $\eta$, this would be $0$.
\end{proof}
% Lean: the bound is Dec.bound with mj; the bump is core_bump and Fourier vanishing is
% fourier_vanish with exists_density (ExistsDensity.lean). (1) is exists_FT_ne. (2) is tate
% (IntegrableConjPowGaussian.lean), with gm, FT_gm (from complexFourier_conj_pow_mul_gaussian),
% kernel, integral_scale, FT_scale, FT_mul and rationalWeight_mul. (3) is coreA with
% prp_pim (exists_lb, the set S_1). (4) is coreB with prp_pim, bet, bet_le, pr_bet, FT_shift,
% hFH (the transform of H_k) and h0.

\begin{lemma}\label{asm-theta-agree-local-bump-levi}
Let $B$ be a theta base of a CM field $E$, let
$(X,W,M,\chi_X,\chi_W,\omega)$ and $(X_1,W_1,M_1,\chi_X,\chi_W,\omega_1)$ be
Weil pairs for $B$ of some signatures
(Definition~\ref{asm-theta-agree-triple}(1)), let $H=(x,y,\iota,\mathfrak b)$
be an adapted plane in $X$ over $X_1$ and $H'$ one in $W$ over $W_1$, at some
positions (Definition~\ref{asm-theta-agree-planes}), and let $e$ be an
evaluation map for $(M,M_1,\iota,\iota_{H'},\mathcal J)$, for some $\mathcal J$,
with components $e_\infty$ and $e_v$. Let $r_X,r_W\in\mathbf R$ be such that
\[
e(\omega(up,u'p')\Phi)=\chi_W(t)|t|^{r_X}\chi_X(t')|t'|^{r_W}\,\omega_1(g_1,h_1)\,e(\Phi)
\]
for every Levi element $p\in U_X$ for $(t,g_1)$, $u\in N_H$, Levi element
$p'\in U_W$ for $(t',h_1)$, $u'\in N_{H'}$ and Schwartz-Bruhat $\Phi$. Let $l$
be a continuous linear form on $\mathcal S(F_\infty^{I_{M_1}})$ with
$l(f_1)=l(f)$ whenever $\omega_1(g,1)$ acts at $\infty$ by $f\mapsto f_1$ for
some $g\in U_{X_1}^{(\infty)}$. Let $k\in U_X^{(\infty)}$ and
$t\in\mathbf A_E^\times$ with $kx=tx$.
\begin{enumerate}
\item If $\omega(k,1)$ acts at $\infty$ by $f\mapsto f_1$, then
  $l(e_\infty(f_1))=\chi_W(t)|t|^{r_X}\,l(e_\infty(f))$.
\item $k^{-1}y=\iota a+bx+c(t)y$ for some vector $a$ and some
  $b\in\mathbf A_E$.
\end{enumerate}
\end{lemma}
\leanref{Development/WeilBound/ThetaLifts/LocalTheory/ArchimedeanBumps/HeisenbergActionAddOfOrthogonal.lean}{PositivityTheory.ThKG.Loc.Bump.mix_un}
\begin{proof}
(1) By Lemma~\ref{asm-theta-agree-split-inputs}(1)(e), $k=um$ with $u\in N_H$
and $m$ a Levi element for $(t,g_1)$, for some $g_1\in U_{X_1}$. Let $P$ be the
matrix with the columns of $\iota$ followed by $x$ and $y$. As $mx=tx$,
$my=c(t)^{-1}y$ and $m\iota=\iota g_1$, the matrix $P^{-1}mP$ is block diagonal
with first block $g_1$; and as $u\iota z\in\iota z+\mathbf A_Ex$ for every $z$,
which holds for the generators of $N_H$ and so for all of $N_H$, the first
diagonal block of $P^{-1}uP$ is $1$. So the first diagonal block of $P^{-1}kP$
is $g_1$, and $g_1-1$, a block of
$P^{-1}(k-1)P$, is at $\infty$; so $g_1\in U_{X_1}^{(\infty)}$. By
Lemma~\ref{asm-theta-agree-rechoice-operators}(1) there is an archimedean
operator $T$ of $\omega_1$ at $(g_1,1)$.

Let $\mathrm{vac}=\bigotimes_v\mathrm{vac}_v$ and
$\mathrm{vac}_1=\bigotimes_v\mathrm{vac}_{1,v}$ be the vacua of $M$ and $M_1$,
which are nonzero. The hypothesis with $u'=p'=1$, a Levi element for $(1,1)$,
gives $e(\omega(k,1)\Phi)=\chi_W(t)|t|^{r_X}\omega_1(g_1,1)e(\Phi)$. For
$\Phi=f\otimes\mathrm{vac}$, $\omega(k,1)\Phi=f_1\otimes\mathrm{vac}$ and
$e(\Phi)=e_\infty(f)\otimes\mathrm{vac}_1$, so
\[
e_\infty(f_1)\otimes\mathrm{vac}_1=\chi_W(t)|t|^{r_X}\,T(e_\infty(f))\otimes\mathrm{vac}_1,
\]
and $e_\infty(f_1)=\chi_W(t)|t|^{r_X}T(e_\infty(f))$. As $\omega_1(g_1,1)$ acts
at $\infty$ by $e_\infty(f)\mapsto T(e_\infty(f))$, $l(T(e_\infty(f)))=l(e_\infty(f))$.

(2) With $u$ and $m$ as in (1), $u^{-1}\in N_H$, so
$u^{-1}y=y+\iota a'+b'x$ by Lemma~\ref{asm-theta-agree-split-inputs}(1)(d), and
$m^{-1}$ is a Levi element for $(t^{-1},g_1^{-1})$. Hence
$k^{-1}y=m^{-1}u^{-1}y=c(t)y+\iota g_1^{-1}a'+b't^{-1}x$.
\end{proof}
% Lean: mix_un. (1): isFr gives the frame P, stab_dec the decomposition k = u m (MatCancel.lean),
% the block identity e1 shows that g1 is at infinity, ArchOp gives T, and Mix2 with u' = p' = 1
% (levi_one), e_pv, pv_act, un_pv (pure_inj) and the invariance hl give the equation. (2):
% IsFr.ad.fix_y for u^-1, levi_inv and levi_vec.

\begin{lemma}\label{asm-theta-agree-local-bump-dense}
Let $(M,M_1,H,H',e,u)$ be a cone setting for a theta base $B$
(Definition~\ref{asm-theta-agree-local-factor-setting}), and let $l$ be a
nonzero continuous linear form on $\mathcal S(F_\infty^{I_{M_1}})$. Then
$l(e_\infty(f))\neq0$ for some $f\in\mathcal S(F_\infty^{I_M})$.
\end{lemma}
\leanref{Development/WeilBound/ThetaLifts/LocalTheory/ArchimedeanBumps/ArchSlice.lean}{PositivityTheory.ThKG.Loc.Bump.EiDense}
\leanref{Development/WeilBound/ThetaLifts/LocalTheory/ArchimedeanBumps/ArchHeisOf.lean}{PositivityTheory.ThKG.Loc.Bump.eiDense_of}
\leanref{Development/WeilBound/ThetaLifts/LocalTheory/ArchimedeanBumps/DictIrr.lean}{PositivityTheory.ThKG.Loc.Bump.irr}
\begin{proof}
Write $V_1=F_\infty^{I_{M_1}}$, a Euclidean space, and suppose that
$l(e_\infty(f))=0$ for every $f$. As $e_\infty\neq0$, $g=e_\infty(f_0)\neq0$ for
some $f_0$.

\emph{Step 1: $l$ vanishes on the functions $x\mapsto e^{iL(x)}g(x-a)$, for
$a\in V_1$ and $L$ a real linear form on $V_1$.} As $F$ is totally real,
$L(x)=-2\pi\operatorname{Tr}_{F_\infty/\mathbf R}(x\cdot b)$ for some
$b\in V_1$, by the nondegeneracy of the trace form. Let $z\in\mathbb W(X_1,W_1)$
be at $\infty$ with $e^{M_1}(z)=(-a,b)$, with $a$ and $b$ read as elements of
$\mathbf A^{I_{M_1}}$ at $\infty$. As $\psi$ is a standard character,
$\psi(s)=e(-\operatorname{Tr}_{F_\infty/\mathbf R}(s_\infty))$ for $s$ at
$\infty$, so $\rho_1(z)$ acts at $\infty$ by $g'\mapsto\rho_{1,\infty}(z)g'$
with
$\rho_{1,\infty}(z)g'(x)=\psi(-\tfrac12a\cdot b)\,e^{iL(x)}g'(x-a)$; and in the
same way $\rho(\iota z\iota'^{\mathsf T})$ acts at $\infty$ by some
$f_0\mapsto f_1$. With $\mathrm{vac}$ and $\mathrm{vac}_1$ the vacua of $M$ and
$M_1$, the property
$e(\rho(\iota z\iota'^{\mathsf T})\Phi)=\rho_1(z)e(\Phi)$ of evaluation maps
for $\Phi=f_0\otimes\mathrm{vac}$ gives
$e_\infty(f_1)\otimes\mathrm{vac}_1=\rho_{1,\infty}(z)g\otimes\mathrm{vac}_1$.
As $\mathrm{vac}_1\neq0$, $e_\infty(f_1)=\rho_{1,\infty}(z)g$, and
$l(\rho_{1,\infty}(z)g)=0$.

\emph{Step 2: $l(\varphi\cdot g(\,\cdot-a))=0$ for $\varphi\in\mathcal S(V_1)$
and $a\in V_1$.} Put $h=g(\,\cdot-a)$, and write
$\varphi(x)=\int e(\langle\eta,x\rangle)b(\eta)\,d\eta$ by
Lemma~\ref{asm-theta-agree-local-bump-schwartz}(1). The map
$\eta\mapsto b(\eta)\,e(\langle\eta,\cdot\rangle)h$ into $\mathcal S(V_1)$ is
continuous, and its seminorms $\|\cdot\|_{k,j}$ are bounded by a constant times
$(1+|\eta|^2)^j|b(\eta)|$, which is integrable. By
Lemma~\ref{asm-theta-agree-local-bump-schwartz}(2) it has a weak integral $f$;
evaluating at points shows $f=\varphi h$. So
$l(\varphi h)=\int b(\eta)\,l(e(\langle\eta,\cdot\rangle)h)\,d\eta=0$ by Step 1.

\emph{Step 3: $l=0$.} The smooth functions with compact support are dense in
$\mathcal S(V_1)$ \cite[Section~7.1]{Hor90}, so it suffices that $l(\varphi)=0$
for such $\varphi$. As $g(y)\neq0$ for some $y$, the open sets
$\{x:g(x-a)\neq0\}$, $a\in V_1$, cover $V_1$, and finitely many of them, for
$a$ in a finite set $A$, cover the support of $\varphi$. Then
$w(x)=\sum_{a\in A}|g(x-a)|^2$ does not vanish on the support of $\varphi$, and
$\varphi=\sum_{a\in A}\varphi_a\,g(\,\cdot-a)$ with
$\varphi_a=\varphi\,\overline{g(\,\cdot-a)}/w$, smooth with compact support.
By Step 2, $l(\varphi)=0$. This contradicts $l\neq0$.
\end{proof}
% Lean: eiDense_of (WeakInt, which is the property SchInt): Step 1 is hh (Heisenberg_2 of the
% evaluation map), lf_surj, schr_acts and archHeis_of, with rest_pure and pure_inj for the
% vacua; Steps 2 and 3 are irr, irr_inner and kill_mul (DictIrr.lean), with semi_mul and cont_mod
% for the bounds and the continuity, and ThKOwn.dense_compactSupport for the density.

\begin{lemma}\label{asm-theta-agree-local-bump-slices-heisenberg}
Let $V$ be a finite-dimensional real vector space with a norm $|\cdot|$, let
$V^*$ be its dual with the operator norm, and let $\mathcal S(V)$ be the
Schwartz space of $V$ with the seminorms
$\|f\|_{k,j}=\sup_x|x|^k|D^jf(x)|$. For $w=(a,\lambda)\in V\times V^*$ put
$\|w\|=\max(|a|,\|\lambda\|)$ and
\[
(S(w)f)(x)=e^{i(\lambda(x)+\lambda(a)/2)}f(x+a).
\]
\begin{enumerate}
\item Each $S(w)$ is a continuous linear operator on $\mathcal S(V)$,
  $S(0)=1$, $S(w)S(-w)=1$, and $w\mapsto S(w)f$ is continuous for every
  $f\in\mathcal S(V)$.
\item For every continuous seminorm $q$ on $\mathcal S(V)$ there are a
  continuous seminorm $p$ and an integer $M\geq0$ with
  $q(S(w)f)\leq(1+\|w\|)^Mp(f)$ for all $w$ and $f$.
\item For every continuous seminorm $q$ on $\mathcal S(V)$ there is a
  continuous seminorm $p$ with $q(S(w)f-f)\leq\|w\|\,p(f)$ for all $f$ and all
  $w$ with $\|w\|\leq1$.
\item Let $Y$ be a set, $Z$ a complex vector space with a seminorm $q$,
  $A_y\colon\mathcal S(V)\to Z$ a linear map for each $y\in Y$,
  $s\colon Y\to[0,\infty)$ a function and $K\geq0$, such that for every $y$
  with $s(y)>0$ there is $w$ with $\|w\|s(y)\leq K$ and
  $A_y(S(w)f)=-A_y(f)$ for every $f$. Let $p$ be a continuous seminorm on
  $\mathcal S(V)$ and $M\geq0$ an integer with
  $q(A_yf)\leq(1+s(y))^Mp(f)$ for all $y$ and $f$. Then for every $N\geq0$
  there is a continuous seminorm $p'$ on $\mathcal S(V)$ with
  $q(A_yf)(1+s(y))^N\leq p'(f)$ for all $y$ and $f$.
\end{enumerate}
\end{lemma}
\leanref{Development/WeilBound/ThetaLifts/SmallCones/SliceArch.lean}{PositivityTheory.ThKCT.Arch.S}
\leanref{Development/WeilBound/ThetaLifts/SmallCones/SliceArch.lean}{PositivityTheory.ThKCT.Arch.S_add}
\leanref{Development/WeilBound/ThetaLifts/SmallCones/SliceArch.lean}{PositivityTheory.ThKCT.Arch.continuous_S}
\leanref{Development/WeilBound/ThetaLifts/SmallCones/SliceArch.lean}{PositivityTheory.ThKCT.Arch.grow}
\leanref{Development/WeilBound/ThetaLifts/SmallCones/SliceArch.lean}{PositivityTheory.ThKCT.Arch.lip}
\leanref{Development/WeilBound/ThetaLifts/SmallCones/SliceArch.lean}{PositivityTheory.ThKCT.Arch.decay}
\begin{proof}
For $m\geq0$ let $P_m$ be the largest of the $\|\cdot\|_{k,j}$ with
$k,j\leq m$, and for a smooth function $g$ on $V$ let $B_m(g)$ be the least $C$
with $(1+|x|)^m|D^jg(x)|\leq C$ for all $x$ and all $j\leq m$. Then
$P_m\leq B_m\leq2^mP_m$ on $\mathcal S(V)$. As the $\|\cdot\|_{k,j}$ define
the topology of $\mathcal S(V)$, every continuous seminorm $q$ has $q\leq CP_m$
for some $C$ and $m$. For $w=(a,\lambda)$ put
$\phi_w(x)=e^{i(\lambda(x)+\lambda(a)/2)}$; then
$D^j\phi_w(x)=i^j\lambda^{\otimes j}\phi_w(x)$, so $|D^j\phi_w(x)|\leq\|\lambda\|^j$.
Two estimates are used. By the Leibniz rule, if $|D^i\varphi(x)|\leq A(1+|x|)^r$
for all $x$ and $i\leq m$, then $B_m(\varphi g)\leq2^mA\,B_{m+r}(g)$. By
Peetre's inequality $1+|x|\leq(1+|a|)(1+|x+a|)$,
$B_m(g(\cdot+a))\leq(1+|a|)^mB_m(g)$.

(2) The two estimates, with $A=(1+\|w\|)^m$ and $r=0$, give
$P_m(S(w)f)\leq B_m(S(w)f)\leq2^m(1+\|w\|)^{2m}B_m(f)\leq4^m(1+\|w\|)^{2m}P_m(f)$.
If $q\leq CP_m$, take $p=4^mCP_m$ and $M=2m$.

(3) Let $q\leq CP_m$ and $\|w\|\leq1$, and write
$S(w)f-f=(\phi_w-1)f(\cdot+a)+(f(\cdot+a)-f)$. As
$|\lambda(x)+\lambda(a)/2|\leq\|w\|(1+|x|)$ and $|D^i\phi_w|\leq\|w\|^i\leq\|w\|$
for $i\geq1$, every derivative of $\phi_w-1$ of order $i$ has
$|D^i(\phi_w-1)(x)|\leq\|w\|(1+|x|)$. The two estimates, with $r=1$, give
$B_m((\phi_w-1)f(\cdot+a))\leq2^m\|w\|2^{m+1}B_{m+1}(f)$. The mean value theorem
for $D^jf$ on the segment from $x$ to $x+a$ and Peetre's inequality give
$B_m(f(\cdot+a)-f)\leq2^m|a|B_{m+1}(f)$. So $q(S(w)f-f)\leq\|w\|p(f)$ with
$p=C(2^{2m+1}+2^m)2^{m+1}P_{m+1}$.

(1) $S(w)$ is linear, and by the estimate in (2) it maps $\mathcal S(V)$
continuously into itself. Clearly $S(0)=1$. A direct computation gives
$S(w+w')=e^{i(\lambda(a')-\lambda'(a))/2}S(w)S(w')$ for $w'=(a',\lambda')$, and
$w'=-w$ gives $S(w)S(-w)=S(0)=1$. By (3), $S(u)f\to f$ as $u\to0$. For fixed $w$,
the same rule gives $S(u)f=e^{i\vartheta(u)}S(w)S(u-w)f$ with real
$\vartheta(u)\to0$ as $u\to w$; as $S(w)$ is continuous, $S(u)f\to S(w)f$.

(4) Applying (3) $n$ times gives, for every continuous seminorm $q'$ and every
$n$, a continuous seminorm $p_n$ with $q'((S(w)-1)^nf)\leq\|w\|^np_n(f)$ for
$\|w\|\leq1$. Let $p_1$ be this seminorm for $q'=p$ and $n=M+N$, put
$R=\max(K,1)$ and $p'=(1+R)^{M+N}p+(2R)^{M+N}p_1$. If $s(y)\leq R$, then
$q(A_yf)(1+s(y))^N\leq(1+s(y))^{M+N}p(f)\leq(1+R)^{M+N}p(f)$. If $s(y)>R$, let
$w$ be as in the hypothesis, so that $\|w\|\leq R/s(y)<1$. As
$A_y((S(w)-1)g)=-2A_y(g)$, $q(A_yg)\leq q(A_y((S(w)-1)g))$ for every $g$, and by
induction
\[
q(A_yf)\leq q\bigl(A_y((S(w)-1)^{M+N}f)\bigr)\leq(1+s(y))^M\|w\|^{M+N}p_1(f).
\]
As $\|w\|(1+s(y))\leq1+R\leq2R$, $q(A_yf)(1+s(y))^N\leq(2R)^{M+N}p_1(f)$.
\end{proof}
% Lean: SmallCones/SliceArch.lean (namespace Arch). B_m is Bd and P_m is P, with exists_P, bd_P
% and P_le; the two estimates are Bd.mul and Bd.shift (peetre), with ph_bd for the derivatives of
% the phase. (2) is grow. (3) is lip, with ph_sub_bd and Bd.shift_sub. (1) is S (smulLeftCLM
% composed with compSubConstCLM), S_zero, S_add, S_neg, tendsto_S and continuous_S. (4) is decay,
% with lip_pow for the iterate.

\begin{lemma}\label{asm-theta-agree-local-bump-slices-adelic}
Let $B$ be a theta base of a CM field $E$
(Definition~\ref{asm-theta-agree-models}(1)), with character $\psi$ of
$\mathbf A=F_\infty\times\mathbf A_f$, put $\psi_f(b)=\psi(0,b)$ for
$b\in\mathbf A_f$, and let $I$ be a finite set. Let $\rho$ be the Schrödinger
operators of Theorem~\ref{asm-theta-agree-evaluation-heisenberg} on functions on
$\mathbf A^I$, and for $w=(a,b)\in\mathbf A_f^I\times\mathbf A_f^I$ let
$(\rho_f(w)\phi)(x)=\psi_f(x\cdot b+\tfrac12a\cdot b)\phi(x+a)$ on functions on
$\mathbf A_f^I$. Give $F_\infty$ the norm $|t|=\max_w|t_w|$ over the infinite
places $w$ of $F$, and $V=F_\infty^I$ the largest norm of the coordinates. For
$(a,b)\in V\times V$ let $(a,b)^\flat=(a,\lambda_b)\in V\times V^*$ with
$\lambda_b(x)=-2\pi\operatorname{Tr}_{F_\infty/\mathbf R}(x\cdot b)$, and let $S$
be the operators of Lemma~\ref{asm-theta-agree-local-bump-slices-heisenberg} on
$\mathcal S(V)$. Write $w_\infty$ and $w_f$ for the archimedean and finite
components of $w\in\mathbf A^I\times\mathbf A^I$.
\begin{enumerate}
\item For $w=(a,b)$ and $w'=(a',b')$ in $\mathbf A^I\times\mathbf A^I$,
  $\rho(w)\rho(w')=\psi(a\cdot b'-a'\cdot b)\,\rho(w')\rho(w)$.
\item For $w\in\mathbf A^I\times\mathbf A^I$, $f\in\mathcal S(V)$ and a function
  $\phi$ on $\mathbf A_f^I$,
  $\rho(w)(f\otimes\phi)=S(w_\infty^\flat)f\otimes\rho_f(w_f)\phi$.
\item If $\phi$ is a Schwartz-Bruhat function on $\mathbf A_f^I$, then so is
  every $\rho_f(w)\phi$, and $w\mapsto\rho_f(w)\phi$ is locally constant on
  $\mathbf A_f^I\times\mathbf A_f^I$.
\item For every neighbourhood $U$ of $0$ in $\mathbf A_f$ there is a compact set
  $C_0\subset\mathbf A_f$ that contains every $b$ with $\psi_f(ab)=1$ for all
  $a\in U$.
\item There is $K\geq0$ such that every nonzero $t\in F_\infty$ has some
  $\alpha\in F_\infty$ with $|\alpha||t|\leq K$ and
  $\operatorname{Tr}_{F_\infty/\mathbf R}(t\alpha)=\frac12$.
\end{enumerate}
\end{lemma}
\leanref{Development/WeilBound/ThetaLifts/SmallCones/SliceAd.lean}{PositivityTheory.ThKCT.Ad.schr_comm}
\leanref{Development/WeilBound/ThetaLifts/SmallCones/SliceAd.lean}{PositivityTheory.ThKCT.Ad.schr_pure}
\leanref{Development/WeilBound/ThetaLifts/SmallCones/SliceAd.lean}{PositivityTheory.ThKCT.Ad.fw_lc}
\leanref{Development/WeilBound/ThetaLifts/SmallCones/SliceAd.lean}{PositivityTheory.ThKCT.Ad.annih}
\leanref{Development/WeilBound/ThetaLifts/SmallCones/SliceAd.lean}{PositivityTheory.ThKCT.Ad.exists_half_trace}
\begin{proof}
(1) Applied to $\phi$ at $x$, the two sides are
$\psi(x\cdot b+\tfrac12a\cdot b)\,\psi((x+a)\cdot b'+\tfrac12a'\cdot b')\,\phi(x+a+a')$
and $\psi(a\cdot b'-a'\cdot b)$ times
$\psi(x\cdot b'+\tfrac12a'\cdot b')\,\psi((x+a')\cdot b+\tfrac12a\cdot b)\,\phi(x+a+a')$,
which agree.

(2) By the definition of a theta base,
$\psi(s)=e^{-2\pi i\operatorname{Tr}(s_\infty)}\psi_f(s_f)$ for $s\in\mathbf A$.
So for $w=(a,b)$ and $x\in\mathbf A^I$, $(\rho(w)(f\otimes\phi))(x)$ is the
product of
$e^{-2\pi i\operatorname{Tr}(x_\infty\cdot b_\infty+\frac12a_\infty\cdot b_\infty)}f(x_\infty+a_\infty)$,
which is $(S(w_\infty^\flat)f)(x_\infty)$, and
$\psi_f(x_f\cdot b_f+\tfrac12a_f\cdot b_f)\phi(x_f+a_f)=(\rho_f(w_f)\phi)(x_f)$.

(3) The character $\psi_f$ is continuous on the totally disconnected group
$\mathbf A_f$, so it is trivial on an open subgroup, and
$(w,x)\mapsto\psi_f(x\cdot b+\frac12a\cdot b)$ is locally constant. So
$\rho_f(w)\phi$ is locally constant, with support in the compact set
$\operatorname{supp}\phi-a$. As $\phi$ is locally constant with compact support,
there is an open subgroup $U_0$ of $\mathbf A_f^I$ with $\phi(x+u)=\phi(x)$ for
all $x$ and $u\in U_0$. Let $w_0=(a_0,b_0)$, and let $\Sigma$ be the compact
support of $\phi(\cdot+a_0)$. By compactness of $\Sigma$, for $w$ near $w_0$ the
factor $\psi_f(x\cdot b+\frac12a\cdot b)$ equals its value at $w_0$ for every
$x\in\Sigma$; if moreover $a-a_0\in U_0$, then $\phi(x+a)=\phi(x+a_0)$ for every
$x$, and $\rho_f(w)\phi=\rho_f(w_0)\phi$.

(4) By Theorem~\ref{asm-theta-agree-invariance-tate}(2), $\psi(s)=\psi_0(\xi s)$
for some $\xi\in F$, and $\xi\neq0$ as $\psi$ is not trivial on $F_\infty$. The
neighbourhood $U$ contains $N\widehat{\mathcal O}$ for some integer $N\geq1$,
where $\widehat{\mathcal O}=\prod_v\mathcal O_v$, as these sets form a basis of
neighbourhoods of $0$ in $\mathbf A_f$. Let $b$ have $\psi_f(ab)=1$ for all
$a\in U$, and let $v$ be a finite place. Then the component $\psi_{0,v}$ of
$\psi_0$ at $v$ is trivial on $\mathcal O_v\,\xi Nb_v$. The character $\psi_0$ is
Tate's standard character, and $\psi_{0,v}$ is trivial on $\mathcal O_vt$ only
for $t$ in the inverse different $\mathfrak d_v^{-1}$ of $F_v$ \cite[\S2.2]{Tat67}.
So $\xi Nb$ lies in the compact set $\prod_v\mathfrak d_v^{-1}$, and $b$ lies in
its image $C_0$ under multiplication by $(\xi N)^{-1}$.

(5) The trace form $(t,t')\mapsto\operatorname{Tr}(tt')$ on
$F_\infty=F\otimes_{\mathbf Q}\mathbf R$ is the extension of scalars of the trace
form of $F/\mathbf Q$, which is nondegenerate. So $\beta(t)=\operatorname{Tr}(t\,\cdot)$
is an injective linear map from $F_\infty$ to its dual, and $|t|\leq K\|\beta(t)\|$
for some $K$ and all $t$. For $t\neq0$ choose $\alpha'$ with $|\alpha'|\leq1$ and
$|\beta(t)(\alpha')|\geq\|\beta(t)\|/2>0$, and put
$\alpha=\alpha'/(2\beta(t)(\alpha'))$. Then $\operatorname{Tr}(t\alpha)=\frac12$
and $|\alpha||t|\leq|t|/\|\beta(t)\|\leq K$.
\end{proof}
% Lean: SmallCones/SliceAd.lean (namespace Ad). (1) is schr_comm. (2) is schr_pure with psi_eq;
% w^flat is aw0 F κ (ev2 w), and rho_f(w_f) is fw ψ (fn2 w) (finS). (3) is fw_mem (finS_mem) and
% fw_lc (HaarUniq.exists_invariant, finChar_isLocallyConstant). (4) is annih:
% existsUnique_eq_adelicChar_mul, exists_forall_mul_mem_of_mem_nhds, and the local conductor
% statement forall_finiteAdelicChar_mul_eq_one_iff with isCompact_idealFiniteAdeles. (5) is
% exists_half_trace = Arch.exists_half (exists_antilipschitzWith, exists_lt_apply_of_lt_opNorm)
% with trace_nondeg.

\begin{lemma}\label{asm-theta-agree-local-bump-slices}
Let $B$ be a theta base of a CM field $E$, let $M$ and $M_1$ be regular models
for lattice spaces $(X,W)$ and $(X_1,W_1)$, of ranks $(n,d)$ and $(n_1,d_1)$,
with Schrödinger operators $\rho$ and form $\langle\,,\rangle$ on
$\mathbb W(X,W)$ (Definition~\ref{asm-theta-agree-models}(3)), and let $e$ be an
evaluation map for $(M,M_1,\iota_X,\iota_W,\mathcal J)$
(Definition~\ref{asm-theta-agree-planes}(5)). Let
$u_1,\ldots,u_m,z_1,\ldots,z_m\in M_{n\times d}(E)$ with $az_j\in\mathcal J$
for every $a\in\mathbf A$ and $\langle z_i,u_j\rangle=\delta_{ij}$, and let
$\phi$ be a Schwartz-Bruhat function on $\mathbf A_f^{I_M}$. Then there are
Schwartz-Bruhat functions $\phi_{1,1},\ldots,\phi_{1,r}$ on
$\mathbf A_f^{I_{M_1}}$, a compact set $C\subset\mathbf A_f^m$ and continuous
linear maps $A_i(y)\colon\mathcal S(F_\infty^{I_M})\to\mathcal S(F_\infty^{I_{M_1}})$,
for $y\in\mathbf A^m$ and $i\leq r$, such that:
\begin{enumerate}
\item $y\mapsto A_i(y)f$ is continuous for every $f$;
\item
  $e\bigl(\rho(\sum_jy_ju_j)(f\otimes\phi)\bigr)=\sum_iA_i(y)f\otimes\phi_{1,i}$
  for all $y$ and $f$;
\item $A_i(y)=0$ unless the finite component of $y$ lies in $C$;
\item for every continuous seminorm $q$ on $\mathcal S(F_\infty^{I_{M_1}})$ and
  every $N$ there is a continuous seminorm $p$ on $\mathcal S(F_\infty^{I_M})$
  with
  $q(A_i(y)f)\bigl(1+\sum_j\sum_w|y_{j,w}|\bigr)^N\leq p(f)$ for all $i$, $y$
  and $f$, where $y_{j,w}$ is the component of $y_j$ at the infinite place $w$.
\end{enumerate}
\end{lemma}
\leanref{Development/WeilBound/ThetaLifts/SmallCones/ThetaUniq.lean}{PositivityTheory.ThKCT.Slice}
\leanref{Development/WeilBound/ThetaLifts/SmallCones/NormLeSum.lean}{PositivityTheory.ThKCT.slice_of}
\begin{proof}
Write $I=I_M$ and $I_1=I_{M_1}$, let $\psi$ be the character of $B$, let $e^M$
be the isomorphism of $M$, and let $e_\infty$ and $e_v$ be components of $e$.
Then $\rho(w)$ is the Schrödinger operator of
Theorem~\ref{asm-theta-agree-evaluation-heisenberg} at $e^M(w)$. Use $\psi_f$,
$\rho_f$, $S$, $w\mapsto w^\flat$ and the norms of
Lemma~\ref{asm-theta-agree-local-bump-slices-adelic}, and for $w\in V\times V$
let $\|w\|$ be the larger of the norms of its two coordinates. Put
$c_j=e^M(u_j)$ and $\zeta_j=e^M(z_j)$ in $\mathbf A^I\times\mathbf A^I$. For
$y\in\mathbf A^m$ write $y_\infty$ and $y_f$ for its archimedean and finite
components, and put $s(y)=\sum_j\sum_w|y_{j,w}|$.

\emph{Step 1: pure tensors.} For every Schwartz-Bruhat function $\phi'$ on
$\mathbf A_f^I$ there is a Schwartz-Bruhat function $\epsilon(\phi')$ on
$\mathbf A_f^{I_1}$ with $e(f\otimes\phi')=e_\infty(f)\otimes\epsilon(\phi')$
for every $f\in\mathcal S(F_\infty^I)$. Indeed, as in the proof of
Lemma~\ref{asm-theta-agree-evaluation-transfer}(1), $\phi'$ is a finite linear
combination of products $\bigotimes_v\phi_v$ of Schwartz-Bruhat functions with
$\phi_v=\mathbf 1_{\mathcal O_v^I}$ for almost all $v$, and
$\mathbf 1_{\mathcal O_v^I}=\mathrm{vac}_v$ for almost all $v$, as $M$ is
regular. For such a product,
$e(f\otimes\bigotimes_v\phi_v)=e_\infty(f)\otimes\bigotimes_ve_v(\phi_v)$, and
$\bigotimes_ve_v(\phi_v)$ is Schwartz-Bruhat: each $e_v(\phi_v)$ is, and
$e_v(\phi_v)=\mathrm{vac}_{1,v}=\mathbf 1_{\mathcal O_v^{I_1}}$ for almost all
$v$, as $M_1$ is regular. The claim follows by linearity. Fix such an
$\epsilon(\phi')$ for each $\phi'$.

\emph{Step 2: the twist.} For $y\in\mathbf A^m$, $j\leq m$, $a\in\mathbf A$ and
Schwartz-Bruhat $\Phi$,
\[
e\bigl(\rho(\textstyle\sum_iy_iu_i)\,\rho(az_j)\Phi\bigr)=\psi(-ay_j)\,e\bigl(\rho(\textstyle\sum_iy_iu_i)\Phi\bigr).
\]
Indeed, the form of $\mathbb W(X,W)$ is
$\langle w,w'\rangle=e_1(w)\cdot e_2(w')-e_1(w')\cdot e_2(w)$, which is
$\mathbf A$-bilinear and alternating, so
$\langle\sum_iy_iu_i,az_j\rangle=-a\sum_iy_i\langle z_j,u_i\rangle=-ay_j$. By
Lemma~\ref{asm-theta-agree-local-bump-slices-adelic}(1),
$\rho(\sum_iy_iu_i)\rho(az_j)=\psi(-ay_j)\rho(az_j)\rho(\sum_iy_iu_i)$. The
function $\rho(\sum_iy_iu_i)\Phi$ is Schwartz-Bruhat
(Theorem~\ref{asm-theta-agree-evaluation-heisenberg}(1)) and $az_j\in\mathcal J$,
so $e$ does not see $\rho(az_j)$.

\emph{Step 3: the factorization.} For $y\in\mathbf A^m$ let
$T(y)=e_\infty\circ S\bigl((\sum_jy_{j,\infty}c_{j,\infty})^\flat\bigr)$, a
continuous linear map, and for $b\in\mathbf A_f^m$ let
$G(b)=\epsilon\bigl(\rho_f(\sum_jb_jc_{j,f})\phi\bigr)$. As
$e^M(\sum_jy_ju_j)=\sum_jy_jc_j$, Lemma~\ref{asm-theta-agree-local-bump-slices-adelic}(2)
and Step 1 give
\[
e\bigl(\rho(\textstyle\sum_jy_ju_j)(f\otimes\phi)\bigr)=T(y)f\otimes G(y_f)\tag{a}
\]
for all $y$ and $f$. By Lemma~\ref{asm-theta-agree-local-bump-slices-adelic}(3),
$b\mapsto\rho_f(\sum_jb_jc_{j,f})\phi$ is locally constant, so $G$ is locally
constant, and every $G(b)$ is Schwartz-Bruhat.

\emph{Step 4: a compact set.} As $e_\infty\neq0$, there is $f_0$ with
$e_\infty(f_0)\neq0$. By Lemma~\ref{asm-theta-agree-local-bump-slices-adelic}(3)
at $0$ there is a neighbourhood $U$ of $0$ in $\mathbf A_f$ with
$\rho_f(a\zeta_{j,f})\phi=\phi$ for all $a\in U$ and all $j$. Let $C_0$ be the
compact set of Lemma~\ref{asm-theta-agree-local-bump-slices-adelic}(4) for $U$,
and $C=(-C_0)^m$. We show that $G(b)\neq0$ implies $b\in C$. Fix $j$ and $a\in U$,
and let $y\in\mathbf A^m$ and $a'\in\mathbf A$ have archimedean components $0$ and
finite components $b$ and $a$. By
Lemma~\ref{asm-theta-agree-local-bump-slices-adelic}(2),
$\rho(a'z_j)(f_0\otimes\phi)=f_0\otimes\rho_f(a\zeta_{j,f})\phi=f_0\otimes\phi$,
and $T(y)=e_\infty$. So Step 2 and (a) give
$e_\infty(f_0)\otimes G(b)=\psi(-a'y_j)\,e_\infty(f_0)\otimes G(b)$. Evaluating
at a point where both factors are nonzero gives $\psi(-a'y_j)=1$, that is,
$\psi_f(a\cdot(-b_j))=1$, as $a'y_j$ has archimedean component $0$. As $a\in U$
was arbitrary, $-b_j\in C_0$.

\emph{Step 5: the maps $A_i$.} As $G$ is locally constant and $C$ is compact,
$G(C)$ is finite. Let $\phi_{1,1},\ldots,\phi_{1,r}$ be its nonzero elements,
which are Schwartz-Bruhat, and put $A_i(y)=T(y)$ if $G(y_f)=\phi_{1,i}$ and
$A_i(y)=0$ otherwise. For (1), $y\mapsto A_i(y)f$ is the product of the locally
constant indicator function of $\{y:G(y_f)=\phi_{1,i}\}$ and the function
$y\mapsto e_\infty\bigl(S((\sum_jy_{j,\infty}c_{j,\infty})^\flat)f\bigr)$, which
is continuous by Lemma~\ref{asm-theta-agree-local-bump-slices-heisenberg}(1).
For (2), if $G(y_f)=0$, both sides vanish by (a), as no $\phi_{1,i}$ is $0$;
otherwise $y_f\in C$ by Step 4, so $G(y_f)=\phi_{1,i_0}$ for exactly one $i_0$,
and the sum is $T(y)f\otimes G(y_f)$, which is (a). For (3), if $A_i(y)\neq0$,
then $G(y_f)=\phi_{1,i}\neq0$, so $y_f\in C$.

\emph{Step 6: the bound (4).} Let $q$ and $N$ be given. By
Lemma~\ref{asm-theta-agree-local-bump-slices-heisenberg}(2) for the seminorm
$q\circ e_\infty$ there are $p$ and $M_0$ with
$q(e_\infty(S(w)f))\leq(1+\|w\|)^{M_0}p(f)$. As $w\mapsto w^\flat$ is linear on
a finite-dimensional space, $\|w^\flat\|\leq K_1\|w\|$ for some $K_1$. With
$C_c=\sum_j\|c_{j,\infty}\|$ and $|y_{j,\infty}|\leq\sum_w|y_{j,w}|$,
$\|(\sum_jy_{j,\infty}c_{j,\infty})^\flat\|\leq K_1s(y)C_c$, so
\[
q(A_i(y)f)\leq(1+s(y))^{M_0}(1+K_1C_c)^{M_0}p(f).
\]
Apply Lemma~\ref{asm-theta-agree-local-bump-slices-heisenberg}(4) to the pairs
$(i,y)$, the maps $A_i(y)$, the function $s$ and $K=K_1Z\,m\,|S_\infty|\,K_2$,
where $Z=\sum_j\|\zeta_{j,\infty}\|$, $S_\infty$ is the set of infinite places of
$F$ and $K_2$ is the constant of
Lemma~\ref{asm-theta-agree-local-bump-slices-adelic}(5). Its hypothesis holds.
Let $s(y)>0$ and choose $j$ with $|y_{j,\infty}|$ largest. Then
$s(y)\leq m|S_\infty||y_{j,\infty}|$, so $y_{j,\infty}\neq0$, and
Lemma~\ref{asm-theta-agree-local-bump-slices-adelic}(5) gives $\alpha\in F_\infty$
with $|\alpha||y_{j,\infty}|\leq K_2$ and
$\operatorname{Tr}(y_{j,\infty}\alpha)=\frac12$. Put
$w=(\alpha\zeta_{j,\infty})^\flat$. Then $\|w\|\leq K_1|\alpha|Z$, so
$\|w\|s(y)\leq K$. Let $a'\in\mathbf A$ have archimedean component $\alpha$ and
finite component $0$. If $G(y_f)=\phi_{1,i}$, then
Lemma~\ref{asm-theta-agree-local-bump-slices-adelic}(2) gives
$\rho(a'z_j)(f\otimes\phi)=S(w)f\otimes\phi$, and Step 2 and (a) give
\[
T(y)S(w)f\otimes\phi_{1,i}=\psi(-a'y_j)\,T(y)f\otimes\phi_{1,i},
\]
with $\psi(-a'y_j)=e^{2\pi i\operatorname{Tr}(\alpha y_{j,\infty})}=-1$, as the
finite component of $a'y_j$ is $0$. As $\phi_{1,i}\neq0$,
$A_i(y)S(w)f=-A_i(y)f$; if $G(y_f)\neq\phi_{1,i}$, both sides are $0$.
Lemma~\ref{asm-theta-agree-local-bump-slices-heisenberg}(4) gives $p'$ with
$q(A_i(y)f)(1+s(y))^N\leq p'(f)$.
\end{proof}
% Lean: slice_of (SmallCones/NormLeSum.lean) = Ad.slice_core with ev_pure (SliceAd.lean). Step 1
% is ev_pure (IsEv.loc, Dense.mem_span_pf, hM.std, hM1.std). Step 2 is key (schr_comm, hinv from
% IsEv.inv and Dual.mem, hp from Dual.pair and M.co.form). Step 3 is T, G and main (schr_pure,
% pairEvaluation_linear_combination), with hG (fw_lc) and hGm. Step 4 is the set C (annih with the
% neighbourhood hU, pureA_cancel and psi_eq; the Lean applies key to S(-wy y) f0, which is f0 here).
% Step 5 is hfin, phi1 and A, with the continuity from continuous_S and the decomposition by
% cases. Step 6 is grow, exists_half_trace, norm_le_sum, sum_le_norm and decay over the pairs
% (i, y), with K = K1 Z (m |S_infty|) K2. Slice is one of the properties that Bump.bumpA_fin and
% bumpB_fin take; it is also used, through thetaExp_of, slice_unif and lineInt_of, in the open
% proof of Lemma asm-theta-agree-split-inputs.

\begin{lemma}\label{asm-theta-agree-local-bump-cone}
Let $(M,M_1,H,H',e,u)$ be a cone setting for a theta base $B$
(Definition~\ref{asm-theta-agree-local-factor-setting}), with $M$ a model for
$(X,W)$, let $(X,W,M,\chi_X,\chi_W,\omega)$ be a Weil pair for $B$ of some
signature (Definition~\ref{asm-theta-agree-triple}(1)), let $K\subset U_X$ be
the lattice stabilizer of $X$ and $m$ a normalized local measure system for
$K$, with $K_\infty$, $A_\infty$, $m_{K,\infty}$ and $m_{A,\infty}$ as in
Definition~\ref{asm-theta-agree-local-cones}(2). Let $S_\infty$ be the set of
infinite places of $F$, and for $c\in\mathbf C^{S_\infty}$ let
$\hat c\in\mathbf A_E$ be the element at $\infty$ with $\pi_w(\hat c)=c_w$ for
every $w$. For $f\in\mathcal S(F_\infty^{I_M})$, a continuous linear form $l$ on
$\mathcal S(F_\infty^{I_{M_1}})$, $k\in K_\infty$ and $c\in\mathbf C^{S_\infty}$
put
\[
a_{f,l}(k,c)=l\Bigl(e_\infty\bigl(\bigl(\rho(\hat c\,u(1))\,\omega(k,1)\bigr)_\infty f\bigr)\Bigr),
\]
with the notation of Definition~\ref{asm-theta-agree-local-cones}(1). Then
$K_\infty$ is compact and contained in $K_\infty(X)$, and for all $f$ and $l$:
$a_{f,l}$ is continuous; $|a_{f,l}(k,c)|(1+|c|)^N$ is bounded for every $N$;
and for every continuous character $\chi$ of $\mathbf A_E^\times$ with values
of absolute value $1$, every integer $z\geq1$ and every archimedean cone $L$
for $(M,\omega,m,e_\infty,u,\chi,z)$,
\[
l(L(f))=\int_{K_\infty}\int_{A_\infty}l\Bigl(e_\infty\bigl(\bigl(\rho(1_\infty u(t))\,\omega(k,1)\bigr)_\infty f\bigr)\Bigr)\chi(t)|t|^z\,dm_{A,\infty}(t)\,dm_{K,\infty}(k).
\]
\end{lemma}
\leanref{Development/WeilBound/ThetaLifts/LocalTheory/ArchimedeanBumps/ArchSlice.lean}{PositivityTheory.ThKG.Loc.Bump.ArchSlice}
\leanref{Development/WeilBound/ThetaLifts/LocalTheory/ArchimedeanBumps/SliceMain.lean}{PositivityTheory.ThKG.Loc.Bump.archSlice_of}
\leanref{Development/WeilBound/ThetaLifts/LocalTheory/ArchimedeanBumps/SliceLine.lean}{PositivityTheory.ThKG.Loc.Bump.line_ex}
\leanref{Development/WeilBound/ThetaLifts/LocalTheory/ArchimedeanBumps/SliceMain.lean}{PositivityTheory.ThKG.Loc.Bump.k_ops}
\leanref{Development/WeilBound/ThetaLifts/LocalTheory/ArchimedeanBumps/SliceMain.lean}{PositivityTheory.ThKG.Loc.Bump.cone_core}
\begin{proof}
Write $H=(x,y,\iota,\mathfrak b)$ and $H'=(x',y',\iota',\mathfrak b')$, and let
$\mathrm{vac}$ and $\mathrm{vac}_1$ be the vacua of $M$ and $M_1$. Their finite
parts $\psi_0(x_f)=\prod_v\mathrm{vac}_v(x_v)$ and
$\psi_1(x_f)=\prod_v\mathrm{vac}_{1,v}(x_v)$ are nonzero Schwartz-Bruhat
functions on $\mathbf A_f^{I_M}$ and $\mathbf A_f^{I_{M_1}}$, with
$f\otimes\mathrm{vac}=f\otimes\psi_0$. For an operator $T$ on functions on
$\mathbf A^{I_M}$, the function $f_1$ with $T$ acting at $\infty$ by $f\mapsto f_1$
is unique if it exists, as one sees on $f\otimes\psi_0$. Through
$t\mapsto(\pi_w(t))_w$, $A_\infty$ is $(\mathbf C^\times)^{S_\infty}$, its Haar
measures are the multiples of $\prod_w|c_w|^{-2}\,dc$, and
$|t|=\prod_w|c_w|^2$.

\emph{Step 1: $K_\infty$.} By Lemma~\ref{asm-theta-agree-stabilizer}, $K$ is
compact, and $K_\infty=K\cap U_X^{(\infty)}$ is closed in $K$, as in the proof of
Lemma~\ref{asm-theta-agree-local-measures}; so $K_\infty$ is compact. Its elements
are at $\infty$ and commute with every $\theta_{X,w}$, so $K_\infty\subset K_\infty(X)$.

\emph{Step 2: slices.} There are continuous linear maps
$A(c)\colon\mathcal S(F_\infty^{I_M})\to\mathcal S(F_\infty^{I_{M_1}})$, for
$c\in\mathbf C^{S_\infty}$, such that $c\mapsto A(c)f$ is continuous,
$A(c)f=e_\infty\bigl((\rho(\hat c\,u(1)))_\infty f\bigr)$, and for every
continuous seminorm $q$ and every $N$ there is a continuous seminorm $p$ with
$q(A(c)f)(1+|c|)^N\leq p(f)$. Indeed, let $u(1)=x_0y_0^{\mathsf T}$ and put
$(x_1,y_1)=(y,x')$ if $u=u_\alpha$, where $(x_0,y_0)=(x,y')$, and
$(x_1,y_1)=(x,y')$ if $u=u_\beta$, where $(x_0,y_0)=(y,x')$. Then
$ax_1y_1^{\mathsf T}\in\mathcal J$ for every $a\in\mathbf A_E$, and
$s=-h(x_1,x_0)h_W(y_1,y_0)\neq0$. Let $(b'_j)$ be a basis of $E$ over $F$ with
dual basis $(b'^\vee_j)$ for the trace form, and put $b_j=c(s^{-1}b'_j)$,
$u_j=b_ju(1)$, $z_j=b'^\vee_jx_1y_1^{\mathsf T}$ and
$Y_j(a)=\operatorname{Tr}_{\mathbf A_E/\mathbf A}(b'^\vee_jsc(a))$ for
$a\in\mathbf A_E$. Then $az_j\in\mathcal J$ for $a\in\mathbf A$,
$\langle z_i,u_j\rangle=\delta_{ij}$, $\sum_jY_j(a)b_j=a$ and
$\sum_jY_j(a)u_j=a\,u(1)$. Lemma~\ref{asm-theta-agree-local-bump-slices} for
$\phi=\psi_0$ gives $\phi_{1,i}$ and $A_i(y)$. Choose $a_0$ with
$\psi_1(a_0)\neq0$ and put
$A(c)=\sum_i\psi_1(a_0)^{-1}\phi_{1,i}(a_0)\,A_i\bigl((Y_j(\hat c))_j\bigr)$. As
$\rho(\hat c\,u(1))$ acts at $\infty$ by some $f\mapsto f'$,
$e(\rho(\hat c\,u(1))(f\otimes\psi_0))=e_\infty(f')\otimes\psi_1$, and comparing
with Lemma~\ref{asm-theta-agree-local-bump-slices}(2) at the points with finite
component $a_0$ gives $A(c)f=e_\infty(f')$. The continuity follows from
Lemma~\ref{asm-theta-agree-local-bump-slices}(1). For the bound,
$|c_w|=|\sum_jY_j(\hat c)_w\tau_w(b_j)|\leq K_1\sum_j\sum_{w'}|Y_j(\hat c)_{w'}|$ with
$K_1=1+\sum_j\sum_w|\tau_w(b_j)|$, so
$(1+|c|)^N\leq K_1^N(1+\sum_j\sum_w|Y_j(\hat c)_w|)^N$, and
Lemma~\ref{asm-theta-agree-local-bump-slices}(4) applies.

\emph{Step 3: the operators of $K_\infty$.} For $k\in K_\infty$ let $T(k)$ be an
archimedean operator of $\omega$ at $(k,1)$
(Lemma~\ref{asm-theta-agree-rechoice-operators}(1)). Then $k\mapsto T(k)f$ is
continuous, and for every continuous seminorm $q$ there is a continuous
seminorm $p$ with $q(T(k)f)\leq p(f)$ for all $k$ and $f$. Indeed, let $l'$ be a
continuous linear form on $\mathcal S(F_\infty^{I_M})$ and $a_1$ a point with
$\psi_0(a_1)\neq0$. The linear form
$\lambda(\Phi)=\psi_0(a_1)^{-1}l'(y\mapsto\Phi(y,a_1))$ is admissible, and
$\lambda(\omega(k,1)(f\otimes\psi_0))=\lambda(T(k)f\otimes\psi_0)=l'(T(k)f)$, which
is continuous in $k$ by Definition~\ref{asm-theta-agree-models}(6)(c), hence
bounded on the compact group $K_\infty$. By
Lemma~\ref{asm-theta-agree-local-bump-schwartz}(3), (4) and (3), in this order,
$q(T(k)f)$ is bounded in $k$, $k\mapsto T(k)f$ is continuous, and the $T(k)$ are
equicontinuous.

\emph{Step 4: $a_{f,l}$.} For $k\in K_\infty$ and $w$ at $\infty$,
$\omega(k,1)$ acts at $\infty$ by $f\mapsto T(k)f$ and $\rho(w)$ by
$T(k)f\mapsto(\rho(w))_\infty T(k)f$, so
$(\rho(w)\omega(k,1))_\infty f=(\rho(w))_\infty T(k)f$ by the uniqueness above.
Hence $a_{f,l}(k,c)=l(A(c)T(k)f)$. As
$A(c)T(k)f-A(c_0)T(k_0)f=A(c)(T(k)f-T(k_0)f)+(A(c)-A(c_0))T(k_0)f$ and the
$A(c)$ are equicontinuous by Step 2 with $N=0$, $a_{f,l}$ is continuous; and
$|a_{f,l}(k,c)|(1+|c|)^N\leq p(T(k)f)\leq p'(f)$ by Steps 2 and 3.

\emph{Step 5: the integral.} For $t\in A_\infty$ let $c(t)=(\pi_w(t))_w$ if
$u=u_\alpha$ and $c(t)=(\overline{\pi_w(t)})_w$ if $u=u_\beta$; then
$1_\infty u(t)=\widehat{c(t)}\,u(1)$ and $|c(t)_w|=|\pi_w(t)|$. Put
$G(k,t)=\chi(t)|t|^zA(c(t))T(k)$. As $z\geq1$, the function
$h(t)=|t|^z(1+|c(t)|)^{-N}$ is $m_{A,\infty}$-integrable for $N$ large, by the
description of the Haar measures above. By Steps 2 and 3, for all $k',j$ there
is a continuous seminorm $p$ with $\|G(k,t)f\|_{k',j}\leq p(f)h(t)$, and by the
argument of Step 4, $(k,t)\mapsto G(k,t)f$ is continuous. So
Lemma~\ref{asm-theta-agree-local-bump-schwartz}(2) and (5), for the measure
$m_{K,\infty}\times m_{A,\infty}$ on the second countable space
$K_\infty\times A_\infty$, give a continuous linear map $L'$ with
$T'(L'(f))=\int\int T'(G(k,t)f)\,dm_{A,\infty}(t)\,dm_{K,\infty}(k)$ for every
continuous linear form $T'$, by Fubini's theorem. By Step 4,
$T'(G(k,t)f)=T'\bigl(e_\infty((\rho(1_\infty u(t))\omega(k,1))_\infty f)\bigr)\chi(t)|t|^z$.
For $T'$ the evaluation at a point, this shows $L'(f)=L(f)$; for $T'=l$ it is
the formula.
\end{proof}
% Lean: archSlice_of StabEx ArchOp Slice archHeis_of ArchUniq SchInt (SliceMain.lean), with the
% identification lift of ArchDict (archDict_of, TracePiReal.lean, and IsLift.haar, DictHaar.lean).
% Step 1 is kAt_cpt and mem_kAt with IsStab. Step 2 is line_ex (SliceLine.lean) with
% exists_dual, vac_fin, ev_op and pureA_ptA. Step 3 is k_ops with
% exists_admissible_functional_adelic, IsUSplitting.cont, weak_bounded, strong_of_weak and
% equicont. Step 4 is op_comp (opI_eq, ArchUniq) and joint_cont. Step 5 is cone_core, with
% shift_eq, xi_int and dom_int for the integrability and weak_op; the last sentence is the end of
% archSlice_of (evAt).

\begin{lemma}\label{asm-theta-agree-local-bump}
Let $B$ be a theta base of a CM field $E$.
\begin{enumerate}
\item Let $(G,G_1,H,H')$ be an $\alpha$-configuration with $G$ of ranks
  $(b+2,a')$, $e$ an evaluation map for
  $(M,M_1,\iota_H,\iota_{H'},\mathcal J_\alpha(H,H'))$ with components
  $e_\infty$ and $e_v$, $K\subset U_V$ the lattice stabilizer of $V$, and $m$ a
  normalized local measure system for $K$. Let $(B,G,\mathcal X,\omega,\omega')$
  and $(B,G_1,\mathcal X,\omega_1,\omega'_1)$ be triples, $L$ an archimedean
  cone for $(M,\omega,m,e_\infty,u_\alpha,\nu,k_\alpha)$, and $l$ a nonzero
  continuous linear form on $\mathcal S(F_\infty^{I_{M_1}})$ with $l(f_1)=l(f)$
  whenever $\omega_1(g,1)$ acts at $\infty$ by $f\mapsto f_1$ for some
  $g\in U_{V_1}^{(\infty)}$. Then $l(L(f))\neq0$ for some
  $f\in\mathcal S(F_\infty^{I_M})$.
\item The same holds for a $\beta$-configuration with $G$ of ranks
  $(a,b+2)$, an evaluation map $e'$ for
  $(M',M'_1,\iota_H,\iota_{H'},\mathcal J_\beta(H,H'))$ with components
  $e'_\infty$ and $e'_v$, the lattice stabilizer $K\subset U_{V'}$ of $V'$, an
  archimedean cone $L$ for $(M',\omega',m,e'_\infty,u_\beta,\nu,k_\beta+1)$, and
  forms $l$ on $\mathcal S(F_\infty^{I_{M'_1}})$ invariant in the same sense
  under the $\omega'_1(g,1)$, $g\in U_{V'_1}^{(\infty)}$.
\end{enumerate}
\end{lemma}
\leanref{Development/WeilBound/ThetaLifts/LocalTheory/UnramifiedDeterminedBySpherical.lean}{PositivityTheory.ThKG.Loc.BumpA}
\leanref{Development/WeilBound/ThetaLifts/LocalTheory/UnramifiedDeterminedBySpherical.lean}{PositivityTheory.ThKG.Loc.BumpB}
\leanref{Development/WeilBound/ThetaLifts/LocalTheory/ArchimedeanBumps/BumpAFin.lean}{PositivityTheory.ThKG.Loc.Bump.bumpA_fin}
\leanref{Development/WeilBound/ThetaLifts/LocalTheory/ArchimedeanBumps/BumpAFin.lean}{PositivityTheory.ThKG.Loc.Bump.bumpB_fin}
\begin{proof}
(1) Let $S_\infty$ be the set of infinite places of $F$, and for
$c\in\mathbf C^{S_\infty}$ let $\hat c\in\mathbf A_E$ be the element at $\infty$
with $\pi_w(\hat c)=c_w$ for every $w$. Write $H=(x,y,\iota_H,\mathfrak b_H)$ and
$H'=(x',y',\iota_{H'},\mathfrak b_{H'})$, so that $u_\alpha(l)=l\,xy'^{\mathsf T}$,
and let $K_\infty$, $m_{K,\infty}$ and $m_{A,\infty}$ be as in
Definition~\ref{asm-theta-agree-local-cones}(2). For
$f\in\mathcal S(F_\infty^{I_M})$, $k\in K_\infty$ and $c\in\mathbf C^{S_\infty}$
put
\[
a_f(k,c)=l\Bigl(e_\infty\bigl(\bigl(\rho(\hat c\,xy'^{\mathsf T})\,\omega(k,1)\bigr)_\infty f\bigr)\Bigr),
\]
with the notation of Definition~\ref{asm-theta-agree-local-cones}(1). The
argument uses the following facts.
\begin{enumerate}
\item[(a)] $K_\infty$ is compact and contained in $K_\infty(V)$; for every $f$,
  $a_f$ is continuous and $|a_f(k,c)|(1+|c|)^N$ is bounded for every $N$; and
  for every continuous Hecke character $\chi$, every integer $z\geq1$ and every
  archimedean cone $L_0$ for $(M,\omega,m,e_\infty,u_\alpha,\chi,z)$,
  \[
  l(L_0(f))=\int_{K_\infty}\int_{A_\infty}l\Bigl(e_\infty\bigl(\bigl(\rho(1_\infty u_\alpha(t))\,\omega(k,1)\bigr)_\infty f\bigr)\Bigr)\chi(t)|t|^z\,dm_{A,\infty}(t)\,dm_{K,\infty}(k).
  \]
\item[(b)] $l(e_\infty(f_0))\neq0$ for some $f_0$.
\item[(c)] There is $r\in\mathbf R$ such that, if $k\in K_\infty$ and
  $t\in\mathbf A_E^\times$ with $t-1$ at $\infty$ have $kx=tx$, then
  $l(e_\infty(\omega(k,1)_\infty f))=(\beta/\alpha)(t)|t|^r\,l(e_\infty(f))$ for
  every $f$.
\item[(d)] Let $K_0$ be a compact space with a finite measure $\mu_0$ that is
  positive on nonempty open sets, let $\xi$ be a continuous map from $K_0$ to
  $(\mathbf C^n)^{S_\infty}$ with nonzero components, put
  $\pi(k,c)=(2c_w\xi(k)_w)_w$, viewed in the real inner product space
  $(\mathbf C^n)^{S_\infty}$, and let $\kappa\colon S_\infty\to\mathbf Z_{\geq0}$,
  $z\geq0$ and
  $r_0(c)=\prod_w(c_w/|c_w|)^{\kappa(w)}|c_w|^{2z}$. Let $a$ be continuous on
  $K_0\times\mathbf C^{S_\infty}$ with $|a(k,c)|(1+|c|)^N$ bounded for every
  $N$, and let $k_1\in K_0$ and $c_1\in\mathbf C$. If
  $\operatorname{Re}(a(p)r_0(p)c_1)>0$ for every $p=(k,c)$ with
  $\pi(p)=\pi(k_1,1)$, writing $r_0(p)=r_0(c)$, then
  $\int_{K_0}\int e^{2\pi i\langle\eta,\pi(k,c)\rangle}a(k,c)r_0(c)\,dc\,d\mu_0(k)\neq0$
  for some $\eta$.
\end{enumerate}
Fact (a) is Lemma~\ref{asm-theta-agree-local-bump-cone}, for the cone setting
$(M,M_1,H,H',e,u_\alpha)$ and the Weil pair $(V,W,M,\beta,\beta/\alpha,\omega)$
of the triple. Fact (b) is Lemma~\ref{asm-theta-agree-local-bump-dense}. Fact
(c) is Lemma~\ref{asm-theta-agree-local-bump-levi}(1), with
$\chi_W=\beta/\alpha$; its hypothesis on $e$ holds, for some $r$, by
Lemma~\ref{asm-theta-agree-split-inputs}(5) for the Weil pairs of the two
triples, and $\omega(k,1)$ acts at $\infty$ by $f\mapsto\omega(k,1)_\infty f$ by
Lemma~\ref{asm-theta-agree-rechoice-operators}(1). Fact (d) is
Lemma~\ref{asm-theta-agree-local-bump-positivity}(3).
% Lean: (a) is ArchSlice = archSlice_of StabEx ArchOp Slice archHeis_of ArchUniq SchInt, with
% the identification lift of ArchDict (archDict_of, IsLift); (b) is EiDense = eiDense_of SchInt;
% (c) is mix_un with Mix2 from MixAB (bumpA_of); (d) is coreA with prp_pim.

Suppose that $l(L(f))=0$ for every $f$. As the triple datum $G$ is proper,
$k_\alpha=\lfloor(b+2-a')/2\rfloor>2(p+q)+1\geq1$; write $k_\alpha=z+1$. Put
$\kappa(w)=m_\alpha(w)-m_\beta(w)$, so that
$\nu(t)=\prod_w(\pi_w(t)/|\pi_w(t)|)^{\kappa(w)}$ for $t\in A_\infty$, by the
properness of $\mathcal X$.

\emph{Step 1: the cone as an integral over $K_\infty\times\mathbf C^{S_\infty}$.}
Through $t\mapsto(\pi_w(t))_w$, the group $A_\infty$ is
$(\mathbf C^\times)^{S_\infty}$, its Haar measures are the multiples of
$\prod_w|c_w|^{-2}\,dc$, and $|t|=\prod_w|c_w|^2$; moreover
$1_\infty u_\alpha(t)=\hat c\,xy'^{\mathsf T}$ for $c=(\pi_w(t))_w$. So (a), for
$L_0=L$, $\chi=\nu$ and $z+1$, gives $c_0\neq0$ with
\[
0=l(L(f))=c_0\int_{K_\infty}\int a_f(k,c)\,r_0(c)\,dc\,dm_{K,\infty}(k),\qquad
r_0(c)=\prod_w(c_w/|c_w|)^{\kappa(w)}|c_w|^{2z},
\]
for every $f$.

\emph{Step 2: a function with $a_f(1,1)\neq0$.} Take $f_0$ as in (b), and let
$\rho(-\hat1\,xy'^{\mathsf T})$ act at $\infty$ by $f_0\mapsto f$. As
$h(x,x)=0$, the form of $\mathbb W(V,W)$ vanishes on $\hat1\,xy'^{\mathsf T}$ and
$-\hat1\,xy'^{\mathsf T}$, so $\rho(\hat1\,xy'^{\mathsf T})\rho(-\hat1\,xy'^{\mathsf T})=1$
and $a_f(1,1)=l(e_\infty(f_0))\neq0$.

\emph{Step 3: positivity on an orbit.} For $k\in K_\infty$ let
$\xi(k)_w=\tau_w(k^{-1}x)\in\mathbf C^{b+2}$, which is continuous in $k$ and
nonzero, and put $\pi(k,c)=(2c_w\xi(k)_w)_w$. Let $p=(k,c)$ with
$\pi(p)=\pi(1,1)$. Then $c_w\tau_w(k^{-1}x)=\tau_w(x)$ for every $w$, so every
$c_w\neq0$, and the element $t\in\mathbf A_E^\times$ with $t-1$ at $\infty$ and
$\pi_w(t)=c_w$ has $kx=tx=\hat c x$. Hence
$\rho(\hat c\,xy'^{\mathsf T})\omega(k,1)=\omega(k,1)\rho(\hat1\,xy'^{\mathsf T})$,
as $\omega(k,1)\rho(w)\omega(k,1)^{-1}=\rho(kw)$, and (c) gives
$a_f(p)=(\beta/\alpha)(t)|t|^ra_f(1,1)$. As $\nu(t)=\prod_w(c_w/|c_w|)^{\kappa(w)}$,
$|t|=\prod_w|c_w|^2$ and $(\beta/\alpha)(t)\nu(t)=1$,
\[
a_f(p)\,r_0(c)\,a_f(1,1)^{-1}=|t|^{r+z}>0.
\]

\emph{Step 4: twists.} For $\eta\in(\mathbf C^{b+2})^{S_\infty}$ there is
$Y_\eta\in\mathbf A_E^{b+2}$ at $\infty$ with
$\pi_w(h(x,kY_\eta))=\sum_i\xi(k)_{w,i}\overline{\eta_{w,i}}$ for every $k$ and
$w$. Let $\rho(Y_\eta x'^{\mathsf T})$ act at $\infty$ by $f\mapsto f_\eta$. As
$\omega(k,1)\rho(Y_\eta x'^{\mathsf T})=\rho(kY_\eta x'^{\mathsf T})\omega(k,1)$,
the element $kY_\eta x'^{\mathsf T}$ lies in $\mathcal J_0(H')\subset\mathcal J_\alpha(H,H')$,
so $e$ does not see it, and the two Schrödinger operators
$\rho(\hat c\,xy'^{\mathsf T})$ and $\rho(kY_\eta x'^{\mathsf T})$ commute up to
$\psi$ of the form of the two elements, which is computed from
$h_W(y',x')=-1$ and the choice of $Y_\eta$. This gives
$a_{f_\eta}(k,c)=e^{2\pi i\langle\eta,\pi(k,c)\rangle}a_f(k,c)$ after replacing
$\eta$ by $-\eta$, and Step 1 for $f_\eta$ gives
$\int_{K_\infty}\int e^{2\pi i\langle\eta,\pi(k,c)\rangle}a_f(k,c)r_0(c)\,dc\,dm_{K,\infty}(k)=0$
for every $\eta$.

Steps 3 and 4 contradict (d), applied to $K_0=K_\infty$ with $m_{K,\infty}$,
$\xi$, $\kappa$, $z$, $a=a_f$, $k_1=1$ and $c_1=a_f(1,1)^{-1}$; the hypotheses on
$a_f$ hold by (a). So $l(L(f))\neq0$ for some $f$.

(2) Now $H=(x,y,\iota_H,\mathfrak b_H)$ is a plane in $V'$, so that
$u_\beta(t)=c(t)\,yx'^{\mathsf T}$; let $\rho'$ be the Schrödinger operators of
$M'$, and let $K_\infty$, $m_{K,\infty}$, $\hat c$, $\xi$, $\pi$ and $Y_\eta$ be as
in (1), for $V'$, which has rank $b+2$. Suppose that $l(L(f))=0$ for every $f$.
Put $z=k_\beta$ and $\kappa(w)=m_\alpha(w)-m_\beta(w)$, so that
$\nu(t)=\prod_w(\pi_w(t)/|\pi_w(t)|)^{\kappa(w)}$ for $t\in A_\infty$; by the
properness of $\mathcal X$, $\kappa(w)>a+b+2>2z$. For $f$, $k\in K_\infty$ and
$c\in\mathbf C^{S_\infty}$ put
\[
a_f(k,c)=l\Bigl(e'_\infty\bigl(\bigl(\rho'(\widehat{\bar c}\,yx'^{\mathsf T})\,\omega'(k,1)\bigr)_\infty f\bigr)\Bigr),
\]
with $\bar c=(\bar c_w)_w$, and let $\hat a_{f,k}$ be the transform of
$a_f(k,\cdot)$, with $[\,,]$, $N$ and $r$ for $\kappa$ and $z$ as in
Lemma~\ref{asm-theta-agree-local-bump-positivity}. By
Lemma~\ref{asm-theta-agree-local-bump-cone}, for the cone setting
$(M',M'_1,H,H',e',u_\beta)$ and the Weil pair
$(V',W,M',\alpha,\alpha/\beta,\omega')$, $a_f$ is continuous and
$|a_f(k,c)|(1+|c|)^N$ is bounded for every $N$.

\emph{Step 1.} For $t\in A_\infty$ and $c=(\pi_w(t))_w$,
$1_\infty u_\beta(t)=\widehat{\bar c}\,yx'^{\mathsf T}$, and
$\nu(t)|t|^{z+1}\prod_w|c_w|^{-2}=r(c)$. So
Lemma~\ref{asm-theta-agree-local-bump-cone}, for $\chi=\nu$ and $z+1$, gives
$c_0\neq0$ with
$0=l(L(f))=c_0\int_{K_\infty}\int a_f(k,c)r(c)\,dc\,dm_{K,\infty}(k)$ for every
$f$.

\emph{Step 2: a function with $\hat a_{f,1}(1)\neq0$.} By
Lemma~\ref{asm-theta-agree-local-bump-dense} there is $f_0$ with
$l(e'_\infty(f_0))\neq0$, so $a_{f_0}(1,0)\neq0$, and by
Lemma~\ref{asm-theta-agree-local-bump-positivity}(1) there is $t_0$ with
$\hat a_{f_0,1}(t_0)\neq0$. Let $\rho'(\widehat{1-t_0}\,xy'^{\mathsf T})$ act at
$\infty$ by $f_0\mapsto f$. This element lies in $\mathcal J_\beta(H,H')$, so
$e'$ is invariant under its Schrödinger operator, and it commutes with
$\rho'(\widehat{\bar c}\,yx'^{\mathsf T})$ up to $\psi$ of the form of the two
elements, which is
$-\operatorname{Tr}(\widehat{\bar c}\,c(\widehat{1-t_0}))$ as
$h(y,x)=h_W(x',y')=1$. As $\psi(s)=e(-\operatorname{Tr}_{F_\infty/\mathbf R}s)$
for $s$ at $\infty$, this gives $a_f(1,c)=e(-[c,t_0-1])\,a_{f_0}(1,c)$, so
$\hat a_{f,1}(1)=\hat a_{f_0,1}(t_0)\neq0$.

\emph{Step 3: positivity.} Let $p=(k,c)$ with $\pi(p)=\pi(1,1)$. As in Step 3
of (1), every $c_w\neq0$, and the $t\in\mathbf A_E^\times$ with $t-1$ at $\infty$
and $\pi_w(t)=c_w$ has $kx=tx$; also $\nu(t)=\prod_w(c_w/|c_w|)^{\kappa(w)}$ and
$|t|=N(c)$. By Lemma~\ref{asm-theta-agree-split-inputs}(5), the hypothesis of
Lemma~\ref{asm-theta-agree-local-bump-levi} holds for $e'$, with
$\chi_W=\alpha/\beta=\nu$ and some $r_X$. By part (2) of that lemma,
$k^{-1}y=\iota_Ha'+b'x+c(t)y$. As $\rho'(w)\omega'(k,1)=\omega'(k,1)\rho'(k^{-1}w)$,
as $e'$ is invariant under the Schrödinger operators of $(\iota_Ha')x'^{\mathsf T}$
and $(b'x)x'^{\mathsf T}$, which lie in $\mathcal J_\beta(H,H')$ and are
orthogonal to every $vx'^{\mathsf T}$ since $h_W(x',x')=0$, and by part (1) of
that lemma,
\[
a_f(k,c')=\nu(t)|t|^{r_X}a_f(1,cc')\qquad(c'\in\mathbf C^{S_\infty}).
\]
As $x\mapsto cx$ multiplies the Lebesgue measure by $N(c)$, this gives
$\hat a_{f,k}(c)=\nu(t)|t|^{r_X}N(c)^{-1}\hat a_{f,1}(1)$. As
$r(c)=\nu(t)N(c)^z$, for $c_1=\hat a_{f,1}(1)^{-1}$,
\[
\hat a_{f,k}(c)\,r(c)^{-1}N(c)^{-1}c_1=N(c)^{r_X-z-2}>0.
\]

\emph{Step 4: twists.} For $\eta\in(\mathbf C^{b+2})^{S_\infty}$ let
$\rho'(Y_\eta x'^{\mathsf T})$ act at $\infty$ by $f\mapsto g$, and put
$\beta(k,\eta)_w=\sum_i\xi(k)_{w,i}\overline{\eta_{w,i}}$. As
$\omega'(k,1)\rho'(Y_\eta x'^{\mathsf T})=\rho'(kY_\eta x'^{\mathsf T})\omega'(k,1)$,
the elements $\widehat{\bar c}\,yx'^{\mathsf T}$ and $kY_\eta x'^{\mathsf T}$ are
orthogonal, and
$kY_\eta=\iota_Ha''+b''x+c(h(x,kY_\eta))y$ with
$\pi_w(c(h(x,kY_\eta)))=\overline{\beta(k,\eta)_w}$, the invariance of $e'$ as in
Step 3 gives $a_g(k,c)=a_f(k,c+\beta(k,\eta))$. So Step 1 for $g$ gives
$\int_{K_\infty}\int a_f(k,c+\beta(k,\eta))r(c)\,dc\,dm_{K,\infty}(k)=0$ for every
$\eta$.

Steps 3 and 4 contradict
Lemma~\ref{asm-theta-agree-local-bump-positivity}(4), applied to $K_0=K_\infty$,
which is first countable as a subspace of a metrizable group, with
$m_{K,\infty}$, $\xi$, $\kappa$, $z$, $a=a_f$, $k_1=1$ and
$c_1=\hat a_{f,1}(1)^{-1}$. So $l(L(f))\neq0$ for some $f$.
\end{proof}
% Lean: Bump.bumpA_fin and bumpB_fin StabEx ArchOp Slice ArchUniq SchInt MixAB
% (ArchimedeanBumps/BumpAFin.lean) = bumpA_of (bumpB_of) with eiDense_of SchInt and
% archSlice_of StabEx ArchOp Slice archHeis_of ArchUniq SchInt. Part (1) is bumpA_of = genA
% (BumpAOf.lean) with kap_ex: Step 1 is zero (with conv, FormYx.lean, for the Haar measure of
% IsLift.haar); Step 2 is h11 (pv_rho, heisenberg_action_add_of_orthogonal); Step 3 is the
% positivity branch of coreA (sig_pt, nu_eq, mix_un, rationalWeight_eq_mul_norm_pow); Step 4
% is the twist branch (Yv, at_Yv, pi_beta, pr_bet, IsLift.ch, om_rho, rho_comm, IsEv.inv).
% Part (2) is bumpB_of = genB (BumpAOf.lean) with kap_ex: a_f is a with star c; Step 1 is zero
% (conv); Step 2 is exists_FT_ne and h11 (rho_comm, IsEv.inv with Or.inl, sy_vv, form_yx,
% IsLift.ch, ch_add); Step 3 is the positivity branch of coreB (sig_pt, nu_eq, mix_un, k2 with
% mat_cancel', om_rho and kill, k4 with FT_scale, k3); Step 4 is the twist branch (decomp,
% pi_beta, kill, heisenberg_action_add_of_orthogonal with sy0). The remaining open step of this
% lemma is the proof of Lemma asm-theta-agree-local-bump-slices (Slice).

\begin{lemma}\label{asm-theta-agree-local-augmentation-point}
Let $(M,M_1,H,H',e,u_\beta)$ be a cone setting for a theta base $B$
(Definition~\ref{asm-theta-agree-local-factor-setting}), with $e$ an evaluation
map for $\mathcal J_\beta(H,H')$, and let $(X,W,M,\chi_X,\chi_W,\omega)$ and
$(X_1,W_1,M_1,\chi_X,\chi_W,\omega_1)$ be Weil pairs, with local operators
$(\omega_v,\rho_v)$ and $(\omega_{1,v},\rho_{1,v})$ at a finite place $v$
(Definition~\ref{asm-theta-agree-finite-local-objects}(4)). Let $r_X,r_W\in\mathbf R$
be such that
\[
e(\omega(up,u'p')\Phi)=\chi_W(t)|t|^{r_X}\chi_X(t')|t'|^{r_W}\,\omega_1(g_1,h_1)\,e(\Phi)
\]
for all Levi elements $p\in U_X$ for $(t,g_1)$ and $p'\in U_W$ for $(t',h_1)$,
$u\in N_H$, $u'\in N_{H'}$ and Schwartz-Bruhat $\Phi$
(Definition~\ref{asm-theta-agree-planes}(2)). Let $l$ be a linear form on the
functions on $F_v^{I_{M_1}}$ that vanishes on every $\phi_1-\phi$ with $\phi,\phi_1$
Schwartz-Bruhat and $\omega_1(g,1)$ acting at $v$ by $\phi\mapsto\phi_1$ for some
$g\in U_{X_1}^{(v)}$. Let $u\in N_H$ and let $p$ be a Levi element for $(t,g_1)$,
with $u$, $p$ and $g_1$ at $v$. Then for every Schwartz-Bruhat $\psi'$ on
$F_v^{I_M}$ and every $a\in\mathbf A_E^\times$,
\[
l\bigl(e_v(\rho_v(1_vu_\beta(a))\,\omega_v(up,1)\psi')\bigr)
=\chi_W(t)|t|^{r_X}\,l\bigl(e_v(\rho_v(1_vu_\beta(ta))\psi')\bigr).
\]
\end{lemma}
\leanref{Development/WeilBound/ThetaLifts/LocalTheory/UnramifiedCohomology/FLocBOf.lean}{PositivityTheory.ThKG.Loc.FL.point}
\begin{proof}
On Schwartz-Bruhat functions the operators $\rho_v(w)$ of the local operators
agree with those of Lemma~\ref{asm-theta-agree-local-factor-evaluation}, as both
give the action of $\rho(w)$ at $v$
(Lemma~\ref{asm-theta-agree-finite-inputs-schur}(1)). Write
$H=(x,y,\iota,\mathfrak b)$ and $H'=(x',y',\iota',\mathfrak b')$, so that
$u_\beta(a)=c(a)\,yx'^{\mathsf T}$. By
Lemma~\ref{asm-theta-agree-split-inputs}(1)(d), applied to $u^{-1}\in N_H$,
$u^{-1}y=y+\iota z_1+bx$ for some $z_1\in\mathbf A_E^{n_1}$ and $b\in\mathbf A_E$.
As $p$ is a Levi element for $(t,g_1)$, $p^{-1}x=t^{-1}x$, $p^{-1}y=c(t)y$ and
$p^{-1}\iota=\iota g_1^{-1}$. Hence
\[
(up)^{-1}\bigl(1_vc(a)y\bigr)=P_1+P_2+P_3,\qquad
P_1=1_vc(a)\,b\,t^{-1}x,\quad P_2=\iota\bigl(1_vc(a)\,g_1^{-1}z_1\bigr),\quad P_3=1_vc(ta)\,y .
\]
Put $w=(P_1+P_2+P_3)x'^{\mathsf T}$, which is at $v$, and
$\psi''=\rho_v(w)\psi'$, a Schwartz-Bruhat function. As
$(up)w=1_vc(a)yx'^{\mathsf T}=1_vu_\beta(a)$, the commutation rule of local
operators gives
\[
\omega_v(up,1)\psi''=\rho_v((up)w)\,\omega_v(up,1)\psi'=\rho_v(1_vu_\beta(a))\,\omega_v(up,1)\psi' .
\]
So the left side of the claim is $l(e_v(\omega_v(up,1)\psi''))$.

By the hypothesis with $u'=p'=1$, a Levi element for $(1,1)$,
$e(\omega(up,1)\Phi)=\chi_W(t)|t|^{r_X}\,\omega_1(g_1,1)e(\Phi)$ for Schwartz-Bruhat
$\Phi$. As $\omega(up,1)$ acts at $v$ by
$\psi''\mapsto\omega_v(up,1)\psi''$ and $\omega_1(g_1,1)$ by
$e_v(\psi'')\mapsto\omega_{1,v}(g_1,1)e_v(\psi'')$,
Lemma~\ref{asm-theta-agree-local-factor-evaluation}(2) gives
$e_v(\omega_v(up,1)\psi'')=\chi_W(t)|t|^{r_X}\,\omega_{1,v}(g_1,1)e_v(\psi'')$. As $l$
vanishes on $\omega_{1,v}(g_1,1)e_v(\psi'')-e_v(\psi'')$,
\[
l\bigl(e_v(\omega_v(up,1)\psi'')\bigr)=\chi_W(t)|t|^{r_X}\,l(e_v(\psi'')).
\]

Finally, as $h(Px'^{\mathsf T},Qx'^{\mathsf T})=h_X(P,Q)\,h_W(x',x')=0$, the form
$\langle\,,\rangle$ vanishes on pairs of matrices $Px'^{\mathsf T}$, so by
Lemma~\ref{asm-theta-agree-local-factor-evaluation}(1),
$\rho_v(w)=\rho_v(P_1x'^{\mathsf T})\rho_v(P_2x'^{\mathsf T})\rho_v(P_3x'^{\mathsf T})$.
Here $P_1x'^{\mathsf T}=x(1_vc(a)bt^{-1}x')^{\mathsf T}$ and
$P_2x'^{\mathsf T}=(\iota\,1_vc(a)g_1^{-1}z_1)x'^{\mathsf T}$ lie in
$\mathcal J_\beta(H,H')$ and are at $v$, so by
Lemma~\ref{asm-theta-agree-local-factor-evaluation}(3),
$e_v(\psi'')=e_v(\rho_v(P_3x'^{\mathsf T})\psi')=e_v(\rho_v(1_vu_\beta(ta))\psi')$.
\end{proof}
% Lean: point (UnramifiedCohomology/FLocBOf.lean): fix_y of Fr.isFr for u^{-1}, levi_inv, e1 for
% the decomposition, L.impl (e2), ev_loc with Mix2 at u' = p' = 1 (e3), hl (e4), and rv_add
% with ev_inv twice (e5). The Lean states it for any representations om, om1 with IsLocRep.

\begin{lemma}\label{asm-theta-agree-local-augmentation-slice}
Let $M$ be a regular model for lattice spaces $(X,W)$, for a theta base $B$,
with vacua $\mathrm{vac}_v$, let $f\in\mathcal S(F_\infty^{I_M})$ be nonzero,
and let $\phi$ be a function on $F_{v_1}^{I_M}$ such that
$f\otimes\phi\otimes\bigotimes_{v\neq v_1}\mathrm{vac}_v$ lies in
$\mathcal S^\circ(M)$ (Definition~\ref{asm-theta-agree-models}(5)). Then $\phi$
lies in the span of the functions of full rank at $\tau_1$.
\end{lemma}
\leanref{Development/WeilBound/ThetaLifts/LocalTheory/AdmissiblePureTensorSchwartz.lean}{PositivityTheory.ThKG.Loc.so_slice}
\begin{proof}
Let $N$ be that span, $N'$ a complementary subspace, and $\pi$ the projection
onto $N'$ with kernel $N$. For a function $\Phi$ on $\mathbf A^{I_M}$ put
$O(\Phi)(\xi)=\pi(\Phi_\xi)(\xi_{v_1})$, where $\Phi_\xi(\zeta)=\Phi(\xi\langle\zeta\rangle)$
and $\xi\langle\zeta\rangle$ has component $\zeta$ at $v_1$ and the components of
$\xi$ elsewhere. For $f'\otimes\bigotimes_v\phi_v$ with almost all $\phi_v$ the
indicator function of $\mathcal O_v^{I_M}$, $\Phi_\xi$ is a multiple of
$\phi_{v_1}$, and so
$O(f'\otimes\bigotimes_v\phi_v)=f'\otimes\pi(\phi_{v_1})\otimes\bigotimes_{v\neq v_1}\phi_v$.
On the generators of $\mathcal S^\circ(M)$, $\phi_{v_1}$ has full rank at $\tau_1$
and $\pi(\phi_{v_1})=0$; so $O$ vanishes on $\mathcal S^\circ(M)$. Hence
$f\otimes\pi(\phi)\otimes\bigotimes_{v\neq v_1}\mathrm{vac}_v=0$. As $f\neq0$, every
$\mathrm{vac}_v$ is nonzero and almost all are the indicator function of
$\mathcal O_v^{I_M}$, $\pi(\phi)=0$, that is, $\phi\in N$.
\end{proof}
% Lean: so_slice (AdmissiblePureTensorSchwartz.lean), with ThKF.exists_factorOp for O,
% pure_upd_zero and pure_ne.

\begin{lemma}\label{asm-theta-agree-local-augmentation-subspace}
Let $(X,W,M,\chi_X,\chi_W,\omega)$ be a Weil pair for a theta base $B$, with $X$ of
rank $n$, let $v$ be a finite place and $t\colon E\to F_v$ an $F$-embedding.
\begin{enumerate}
\item If $\phi$ has full rank at $t$ (Definition~\ref{asm-theta-agree-models}(5))
  and $\iota$ is a linear model of the second kind for $(t,1,1)$ for $M$, then
  $\iota\phi(Y)=0$ for every $Y$ of rank different from $n$.
\item Every function in the span of the functions of full rank at $t$ has full
  rank at $t$.
\end{enumerate}
\end{lemma}
\leanref{Development/WeilBound/CoreStatements/LocallyConstantCoinvariants/FullOfSoLoc.lean}{PositivityTheory.KCo.full_of_soLoc}
\leanref{Development/WeilBound/ThetaLifts/SmallCones/SoDefs.lean}{PositivityTheory.ThKCT.LinEx}
\leanref{Development/WeilBound/ThetaLifts/LocalTheory/UnramifiedCohomology/SpanSoLoc.lean}{PositivityTheory.ThKG.Loc.FL.span_soLoc}
\begin{proof}
The matrices of $X$ and $W$ are invertible, as their images under the $\tau_w$
have Cartan involutions. Let
$\operatorname{co}(u)=-t(J_X)\pi_{t\circ c}(u)t(J_W)^{\mathsf T}$.

(1) Let $\iota'$ be a linear model of the second kind for $(t,1,1)$ with
$\iota'\phi(Y)=0$ for $Y$ of rank different from $n$. Let $Y$ have rank
different from $n$. By Lemma~\ref{asm-theta-agree-coinvariants-eigen}(1),
$Y=\operatorname{co}(x_1)$ for some $x_1$ at $v$ with $\pi_t(x_1)=0$. Let
$l(\chi)=\iota\chi(Y)$. For $u$ at $v$ with $\pi_{t\circ c}(u)=0$ and
Schwartz-Bruhat $\chi$, $\rho(u)$ acts at $v$ by $\chi\mapsto\rho_v(u)\chi$, so the
definition of a linear model of the second kind gives
$l(\rho_v(u)\chi)=\psi(\langle x_1,u\rangle)\,l(\chi)$. By
Lemma~\ref{asm-theta-agree-coinvariants-eigen}(3), for $\iota'$ and $x_1$,
$l(\phi)=0$, as $\iota'\phi(\operatorname{co}(x_1))=\iota'\phi(Y)=0$.

(2) By Lemma~\ref{asm-theta-agree-coinvariants-dictionary}(1) there is a linear
model $\iota$ of the second kind for $(t,1,1)$. By (1) and the definition, the
functions of full rank at $t$ are the Schwartz-Bruhat $\phi$ with $\iota\phi(Y)=0$
for every $Y$ of rank different from $n$. As $\iota$ is linear, these form a
subspace.
\end{proof}
% Lean: full_of_soLoc (exists_coB, eigen_kill and mult'); span_soLoc takes the linear model from
% the first part of LinEx (linEx_of: linBar_of_lm of exists_lm, the construction of
% Lemma asm-theta-agree-coinvariants-dictionary(1)).

\begin{lemma}\label{asm-theta-agree-local-invariance-parabolic}
Let $B$ be a theta base of a CM field $E$, let $Y$ be a lattice space of rank $m$,
Hermitian or skew-Hermitian of some signature, and let
$H=(x,y,\iota,\mathfrak b)$ be an adapted plane in $Y$ over a lattice space $Y_1$
of rank $n_1$, with unipotent group $N_H$ and Levi elements $\lambda(t,g_1)$
(Lemma~\ref{asm-theta-agree-split-inputs}(1)(c)). An element $p\in U_Y$ is
\emph{parabolic with multiplier} $t\in\mathbf A_E^\times$ if $px=tx$; let $P$ be
the set of parabolic elements, and put $|t|=|N_{E/F}(t)|_{\mathbf A}$. Then
$P$ is a subgroup of $U_Y$, and there is a nonzero measure $\mu_N$ on the Borel
sets of $N_H$, finite on compact sets and invariant under left and right
translations by $N_H$, such that:
\begin{enumerate}
\item the map $\Psi(u,g_1,t)=u\lambda(t,g_1)$ is a homeomorphism from
  $N_H\times U_{Y_1}\times\mathbf A_E^\times$ onto $P$, which is closed in $U_Y$,
  and $u\lambda(t,g_1)$ is parabolic with multiplier $t$;
\item for $g_1\in U_{Y_1}$ and $t\in\mathbf A_E^\times$, conjugation
  $u\mapsto\lambda u\lambda^{-1}$ by $\lambda=\lambda(t,g_1)$ maps $N_H$ onto
  itself, and $\mu_N(\lambda S\lambda^{-1})=|t|^{n_1+1}\mu_N(S)$ for every Borel
  set $S\subset N_H$.
\end{enumerate}
\end{lemma}
\leanref{Development/WeilBound/ThetaLifts/SmallCones/FcCols.lean}{PositivityTheory.ThKCT.Unf.parEx}
\leanref{Development/WeilBound/ThetaLifts/SmallCones/FcCols.lean}{PositivityTheory.ThKCT.Unf.conjEx}
\leanref{Development/WeilBound/ThetaLifts/SmallCones/UnfParH.lean}{PositivityTheory.ThKCT.Unf.parH}
\leanref{Development/WeilBound/ThetaLifts/SmallCones/UnfPar.lean}{PositivityTheory.ThKCT.Unf.par2}
\leanref{Development/WeilBound/ThetaLifts/SmallCones/UnfMod.lean}{PositivityTheory.ThKCT.Unf.modulus}
\leanref{Development/WeilBound/ThetaLifts/SmallCones/UnfAbs.lean}{PositivityTheory.ThKCT.Unf.Par}
\begin{proof}
If $px=tx$ and $p'x=t'x$, then $pp'x=tt'x$ and $p^{-1}x=t^{-1}x$; so $P$ is a
subgroup, and multipliers multiply. Let $J$ and $J_1$ be the matrices of $Y$
and $Y_1$, and $h(a,a')=a^{\mathsf T}J\,c(a')$. Let $\varepsilon=1$ if $Y$ is
Hermitian and $\varepsilon=-1$ if it is skew-Hermitian, so that
$h(a,a')=\varepsilon\,c(h(a',a))$. As $E$ is a CM field, there are $\xi\in E^\times$
with $c(\xi)=\varepsilon\xi$ and $\delta\in E^\times$ with $c(\delta)=-\delta$.

\emph{Step 1: a hyperbolic basis.} Put $y'=c(\xi)^{-1}y$, and let $A\in M_m(E)$
have the columns $x$, $y'$, $\iota e_1,\ldots,\iota e_{n_1}$. The conditions on an
adapted plane, with $h(a,a')=\varepsilon\,c(h(a',a))$ and $c(\xi)=\varepsilon\xi$,
give $\xi h(y',x)=\xi h(x,y')=1$, $h(x,x)=h(y',y')=0$, $h(\iota e_k,\iota e_l)=(J_1)_{kl}$,
and $h(a,\iota e_k)=h(\iota e_k,a)=0$ for $a\in\{x,y'\}$; that is,
\[
\xi\bigl(A^{\mathsf T}J\,c(A)\bigr)^{\mathsf T}=J'=\begin{pmatrix}0&1&0\\1&0&0\\0&0&J'_1\end{pmatrix},\qquad J'_1=\xi J_1^{\mathsf T},
\]
in blocks of sizes $1$, $1$ and $n_1$. The lattice of $Y$ spans $E^m$, and each
of its elements is $sx+\iota u+ty$ with $s,t\in E$
(Definition~\ref{asm-theta-agree-planes}(1)); so the $m=n_1+2$ columns of $A$
span $E^m$, and $A$ is invertible. As $J$ is invertible, its images $\tau_w(J)$
having Cartan involutions, so are $J'$ and $J_1$. For
$g\in GL_m(\mathbf A_E)$ put $g'=A^{-1}gA$; a computation from the
displayed identity shows that $g\in U_Y$ if and only if
$c(g')^{\mathsf T}J'g'=J'$, and $g\mapsto g'$ is a homeomorphism. As $Ae_0=x$, with
$e_0$ the first unit vector, $g$ is parabolic with multiplier $t$ if and only if
$g'e_0=te_0$. By the definition of a Levi element, $\lambda(t,g_1)'$ is the block
diagonal matrix $\operatorname{diag}(t,c(t)^{-1},g_1)$.

\emph{Step 2: coordinates.} For $w\in\mathbf A_E^{n_1}$ and $s\in\mathbf A$ put
$z(w,s)=-\tfrac12c(w)^{\mathsf T}J'_1w+s\delta$, so that
$z+c(z)+c(w)^{\mathsf T}J'_1w=0$, and
\[
n(w,s)=\begin{pmatrix}1&z(w,s)&-c(w)^{\mathsf T}J'_1\\0&1&0\\0&w&1\end{pmatrix}.
\]
A direct computation with these block matrices gives:
\begin{enumerate}
\item[(a)] $c(n)^{\mathsf T}J'n=J'$ for $n=n(w,s)$, and
  $n(w,s)n(w',s')=n(w+w',s+s'+c_0(w,w'))$ for a continuous function $c_0$ with
  values in $\mathbf A$;
\item[(b)] for $l=\operatorname{diag}(t,c(t)^{-1},g_1)$ with $g_1\in U_{Y_1}$,
  $l\,n(w,s)\,l^{-1}=n\bigl(c(t)g_1w,N_{E/F}(t)s\bigr)$;
\item[(c)] the map $(w,s,g_1,t)\mapsto n(w,s)\operatorname{diag}(t,c(t)^{-1},g_1)$
  is a homeomorphism from $\mathbf A_E^{n_1}\times\mathbf A\times U_{Y_1}\times\mathbf A_E^\times$
  onto a closed subset, which contains every $g'$ with
  $c(g')^{\mathsf T}J'g'=J'$ and $g'e_0\in\mathbf A_E^\times e_0$;
\item[(d)] the matrices $An(w,s)A^{-1}$ are the elements of $N_H$.
\end{enumerate}
For (d): $n(w,s)$ fixes $e_0$ and maps $e_{k+2}$ into $e_{k+2}+\mathbf A_Ee_0$, so
$An(w,s)A^{-1}$ fixes $x$ and maps every $\iota z$ into $\iota z+\mathbf A_Ex$;
conversely an element of $N_H$ has these two properties, and by (c) it is
$An(w,s)\operatorname{diag}(1,1,g_1)A^{-1}$, where the second property forces
$g_1=1$.

\emph{Step 3: the measure.} Let $\mu_W$ and $\mu_A$ be Haar measures on
$\mathbf A_E^{n_1}$ and $\mathbf A$, and let $\mu_N$ be the image of
$\mu_W\times\mu_A$ under $(w,s)\mapsto An(w,s)A^{-1}$, a homeomorphism onto $N_H$
by (c) and (d). By (a), left translation by $An(w_0,s_0)A^{-1}$ acts on the
coordinates by $(w,s)\mapsto(w_0+w,s_0+s+c_0(w_0,w))$, a translation in $w$
followed by a translation in $s$ that depends on $w$, which preserves
$\mu_W\times\mu_A$ by Fubini's theorem; likewise for right translations. By (b),
conjugation by $\lambda(t,g_1)$ acts by
$(w,s)\mapsto(c(t)g_1w,N_{E/F}(t)s)$. At each place, an invertible linear map of
a finite-dimensional vector space multiplies its Haar measures by the absolute
value of its determinant \cite[Chapter~I, \S2]{Wei67}; at almost all places the
map preserves the integral points, which have measure $1$. Hence multiplication
by $N_{E/F}(t)$ multiplies the Haar measures of $\mathbf A$ by
$|N_{E/F}(t)|_{\mathbf A}=|t|$, multiplication by $c(t)$ multiplies those of
$\mathbf A_E^{n_1}$ by $|N_{E/F}(c(t))|_{\mathbf A}^{n_1}=|t|^{n_1}$, and $g_1$
multiplies them by $|N_{E/F}(\det g_1)|_{\mathbf A}=1$, as
$\det(g_1)\,c(\det g_1)=1$ by $g_1^{\mathsf T}J_1c(g_1)=J_1$. So
$\mu_N(\lambda(t,g_1)S\lambda(t,g_1)^{-1})=|t|^{n_1}|t|\,\mu_N(S)$.

\emph{Step 4: the claims.} By Steps 1 and 2(c), $\Psi$ is a homeomorphism onto a
closed subset of $U_Y$, which contains $P$; and $u\lambda(t,g_1)x=u(tx)=tx$, as $u$
fixes $x$, so the image is $P$. This is (1), and (2) is Step 3, with (b) for the
first claim.
\end{proof}
% Lean: parEx (SmallCones/FcCols.lean). The first sentence is used in htEx (act, inv). Step 1 is
% exists_eps, exists_scale, the matrix fc x y' iota with hA (the identity, entry by entry) and hu,
% and conjEx with conj_iff; Step 2 is parH (UnfParH.lean) from par2 (UnfPar.lean: nn, iota, mm,
% hval, hmn, hin, hmi, hdec, the closed embedding hPsi), with exists_anti for delta, fc_cols and
% mem_unip for (d); Step 3 is the skew products for h1 and hR, adelic_haar_invariant_unitary_action
% for g_1, and modulus (UnfMod.lean) with chi1 and chi2: chi(a) |a|^(M+1) = 1, where chi(a) is the
% factor by which (w, s) -> (c(a) w, N(a) s) pushes the measure forward. The Lean proves chi2 in its
% own way: chi is trivial on E^x by a lattice argument (lattice_one), equals the right power of
% exp on a real flow, and a character bounded on a compact set of representatives is trivial.

\begin{lemma}\label{asm-theta-agree-local-invariance-height}
Let $B$ be a theta base of a CM field $E$, $Y$ a Hermitian lattice space of some
signature, $H$ an adapted plane in $Y$ over a lattice space $Y_1$ of rank $n_1$,
with Levi elements $\lambda(t,g_1)$, and $K\subset U_Y$ an Iwasawa subgroup for
$x_H$ (Definition~\ref{asm-theta-agree-split-notions}(2)). Call $p\in U_Y$
parabolic with multiplier $t$ if $px_H=tx_H$, and put
$|t|=|N_{E/F}(t)|_{\mathbf A}$.
\begin{enumerate}
\item Every $g\in U_Y$ is $u\lambda(t,g_1)k$ with $u\in N_H$, $g_1\in U_{Y_1}$,
  $t\in\mathbf A_E^\times$ and $k\in K$.
\item If $k\in K$ is parabolic with multiplier $t$, then $|t|=1$.
\item There is a continuous function $q\colon U_Y\to\mathbf R_{>0}$ with
  $q(gk)=q(g)$ for $g\in U_Y$ and $k\in K$, and $q(pg)=|t|^{n_1+1}q(g)$ for
  $g\in U_Y$ and $p$ parabolic with multiplier $t$.
\end{enumerate}
\end{lemma}
\leanref{Development/WeilBound/ThetaLifts/LocalTheory/UnramifiedCohomology/KAtCompact.lean}{PositivityTheory.ThKG.Loc.KInv.htEx}
\begin{proof}
Let $\Psi$ be the map of Lemma~\ref{asm-theta-agree-local-invariance-parabolic}(1),
and $P$ the group of parabolic elements, in which multipliers multiply.

(1) As $K$ is an Iwasawa subgroup, $g=pk$ with $k\in K$ and $p$ parabolic, with
multiplier $t$ say. By Lemma~\ref{asm-theta-agree-split-inputs}(1)(e) and (c),
$p=u\lambda(t,g_1)$ with $u\in N_H$ and $g_1\in U_{Y_1}$.

(2) As $K$ is compact and $\Psi$ is a homeomorphism onto the closed set $P$,
$\Psi^{-1}(K\cap P)$ is compact, so $|t|\leq C$ on it for some $C$. Let $k\in K$
be parabolic with multiplier $t$. For $j\geq0$, $k^j\in K$ is parabolic with
multiplier $t^j$, so $k^j=\Psi(u_j,g_j,t^j)$ and $|t|^j\leq C$. Hence
$|t|\leq1$. Applied to $k^{-1}$, with multiplier $t^{-1}$, this gives
$|t|^{-1}\leq1$; so $|t|=1$.

(3) By (1), choose for every $g$ a decomposition
$g=u_g\lambda(t_g,h_g)k_g$, and put $q(g)=|t_g|^{n_1+1}>0$. Let
$g=u\lambda(t,g_1)k$ be any decomposition with $k\in K$. Then
$(u\lambda(t,g_1))^{-1}u_g\lambda(t_g,h_g)=kk_g^{-1}$ lies in $K$ and is parabolic
with multiplier $t^{-1}t_g$, so $|t_g|=|t|$ by (2). Hence
$q(u\lambda(t,g_1)k)=|t|^{n_1+1}$ for every such decomposition. As
$gk=u_g\lambda(t_g,h_g)(k_gk)$, $q(gk)=q(g)$ for $k\in K$. If $p$ is parabolic
with multiplier $t$, then $pu_g\lambda(t_g,h_g)$ is parabolic with multiplier
$tt_g$, so it is $u\lambda(tt_g,g_1)$ for some $u$ and $g_1$, by
Lemma~\ref{asm-theta-agree-local-invariance-parabolic}(1); hence
$q(pg)=|tt_g|^{n_1+1}=|t|^{n_1+1}q(g)$. For the continuity, let
$C\subset\mathbf R$ be closed. Then $q^{-1}(C)=\Psi(Z)K$, with $Z$ the set of
$(u,g_1,t)$ with $|t|^{n_1+1}\in C$, which is closed as $t\mapsto|t|$ is
continuous. As $\Psi$ is a homeomorphism onto a closed set, $\Psi(Z)$ is closed
in $U_Y$, and in a topological group the product of a closed set and a compact
set is closed. So $q^{-1}(C)$ is closed, and $q$ is continuous.
\end{proof}
% Lean: htEx (UnramifiedCohomology/KAtCompact.lean), with the data of parEx (hemb, hu, hL). (1) is
% hG, with dec (Fr.stab_dec and Fr.levi_uniq) and act; (2) is le1 (the bound C from
% hemb.isCompact_preimage and pow_unbounded_of_one_lt) and one; (3) is ht with spec, htK, htP and
% htc (hemb.isClosedMap and IsClosed.mul_right_of_isCompact).

\begin{lemma}\label{asm-theta-agree-local-invariance-integral}
Let $B$ be a theta base of a CM field $E$, $Y$ a Hermitian lattice space of some
signature, $H$ an adapted plane in $Y$ over a lattice space $Y_1$ of rank $n_1$,
and $K\subset U_Y$ an Iwasawa subgroup for $x_H$. Call $p\in U_Y$ parabolic with
multiplier $t$ if $px_H=tx_H$, and put $|t|=|N_{E/F}(t)|_{\mathbf A}$. Let $f_0$
be a continuous function on $U_Y$ with $f_0(pg)=|t|^{n_1+1}f_0(g)$ for every
$g\in U_Y$ and every $p$ parabolic with multiplier $t$. Then for every Haar
measure $\mu_K$ on $K$ and every $g\in U_Y$,
\[
\int_Kf_0(kg)\,d\mu_K(k)=\int_Kf_0(k)\,d\mu_K(k).
\]
\end{lemma}
\leanref{Development/WeilBound/ThetaLifts/LocalTheory/UnramifiedCohomology/Lint.lean}{PositivityTheory.ThKG.Loc.KInv.int}
\leanref{Development/WeilBound/ThetaLifts/LocalTheory/UnramifiedCohomology/Lint.lean}{PositivityTheory.ThKG.Loc.KInv.lint}
\leanref{Development/WeilBound/ThetaLifts/SmallCones/UnfAbs.lean}{PositivityTheory.ThKCT.Unf.haar_eq}
\begin{proof}
Let $P$ be the group of parabolic elements, closed in $U_Y$ by
Lemma~\ref{asm-theta-agree-local-invariance-parabolic}(1). The group $U_Y$ is
unimodular (Theorem~\ref{asm-theta-agree-definite-volume}), $K$ is compact, and
$U_Y=PK$, as $K$ is an Iwasawa subgroup for $x_H$. For such data,
\cite[Lemma~2.6.1]{Bum97} shows that $\int_Kf(kg)\,dk=\int_Kf(k)\,dk$ for every
continuous $f$ on $U_Y$ with $f(pg)=\delta_P(p)f(g)$, where $\delta_P$ is the
modular quasicharacter of $P$:
$\int_P\varphi(pp_0)\,dp=\delta_P(p_0)\int_P\varphi(p)\,dp$ for a left Haar measure
$dp$ on $P$ and every continuous $\varphi$ with compact support. For the Borel
subgroup of $GL_2$ this normalization gives
$\delta\bigl(\begin{smallmatrix}y_1&x\\0&y_2\end{smallmatrix}\bigr)=|y_1/y_2|$.
So it suffices to show that $\delta_P(u\lambda(t,g_1))=|t|^{n_1+1}$.

Let $\mu_N$ and $\Psi$ be as in
Lemma~\ref{asm-theta-agree-local-invariance-parabolic}, and let $\mu_1$ and $\mu_A$
be Haar measures on $U_{Y_1}$ and $\mathbf A_E^\times$. Through the homeomorphism
$\Psi$ define a measure on $P$ by
\[
\int_P\varphi(p)\,dp=\int_{\mathbf A_E^\times}\int_{U_{Y_1}}\int_{N_H}\varphi(u\lambda(t,g_1))\,|t|^{-(n_1+1)}\,d\mu_N(u)\,d\mu_1(g_1)\,d\mu_A(t).
\]
It is invariant under left translations. Let $p_0=u_0\lambda_0$ with
$\lambda_0=\lambda(t_0,h_0)$. As $\lambda$ is a homomorphism,
$p_0u\lambda(t,g_1)=u_0(\lambda_0u\lambda_0^{-1})\lambda(t_0t,h_0g_1)$. The
substitution $u\mapsto u_0\lambda_0u\lambda_0^{-1}$ multiplies the integral over
$N_H$ by $|t_0|^{-(n_1+1)}$
(Lemma~\ref{asm-theta-agree-local-invariance-parabolic}(2) and the left
invariance of $\mu_N$), the substitutions $t\mapsto t_0t$ and
$g_1\mapsto h_0g_1$ preserve $\mu_A$ and $\mu_1$, and
$|t_0t|^{-(n_1+1)}=|t_0|^{-(n_1+1)}|t|^{-(n_1+1)}$; so the integrals of
$\varphi(p_0p)$ and $\varphi(p)$ agree. For right translations,
$u\lambda(t,g_1)p_0=u(\lambda u_0\lambda^{-1})\lambda(tt_0,g_1h_0)$ with
$\lambda=\lambda(t,g_1)$. The substitution $u\mapsto u\lambda u_0\lambda^{-1}$
preserves $\mu_N$, $g_1\mapsto g_1h_0$ preserves $\mu_1$, as $U_{Y_1}$ is
unimodular (Theorem~\ref{asm-theta-agree-definite-volume}), and $t\mapsto tt_0$
preserves $\mu_A$; the weight gives
$\int_P\varphi(pp_0)\,dp=|t_0|^{n_1+1}\int_P\varphi(p)\,dp$. Hence
$\delta_P(u_0\lambda(t_0,h_0))=|t_0|^{n_1+1}$, and the hypothesis on $f_0$ is
$f_0(pg)=\delta_P(p)f_0(g)$.
\end{proof}
% Lean: KInv.int (Lint.lean), applied in kInv to the data of parEx (hpar, hR), with the Haar
% measures nu on U_Y and muH on U_{Y_1} from ThKF.unimodularity_holds. The Lean proves the
% invariance in its own way for a group G with the data Par of a closed embedding of
% N x H x T (Lemma asm-theta-agree-local-invariance-parabolic): Unf.haar_eq writes a Haar measure
% of G as the integral over K of the measure above, and lint compares integrals of f against a
% right K-invariant bump with the right invariance of nu. Its law is f(nn p iota u mm t x) chi t =
% f x, with chi t = |t|^(-(n1+1)).

\begin{lemma}\label{asm-theta-agree-local-augmentation-invariance}
Let $B$ be a theta base of a CM field $E$, $Y$ a Hermitian lattice space of some
signature, $H$ an adapted plane in $Y$ over a lattice space $Y_1$ of rank $n_1$,
and $K\subset U_Y$ the lattice stabilizer of $Y$; suppose that $K$ is an Iwasawa
subgroup for the first vector $x_H$ of $H$
(Definition~\ref{asm-theta-agree-split-notions}(2)).
Let $v$ be a finite place, $\mu$ a Haar measure on $K_v$, and $f$ a function on
$U_Y$ that is continuous on $U_Y^{(v)}$ and satisfies
$f(upx)=|t|^{n_1+1}f(x)$ whenever $u\in N_H$ is at $v$, $p$ is a Levi element for
$(t,g_1)$ with $p$ and $g_1$ at $v$ and $t-1$ at $v$, and $x\in U_Y^{(v)}$. Then
\[
\int_{K_v}f(kg)\,d\mu(k)=\int_{K_v}f(k)\,d\mu(k)\qquad(g\in U_Y^{(v)}).
\]
\end{lemma}
\leanref{Development/WeilBound/ThetaLifts/LocalTheory/UnramifiedCohomology/KAtCompact.lean}{PositivityTheory.ThKG.Loc.KInv.kInv}
\begin{proof}
Write $x_0=x_H$, and say that $p\in U_Y$ is \emph{parabolic with multiplier} $t$
if $px_0=tx_0$. For an idempotent $a\in\mathbf A$ and $g$ in $U_Y$ or $U_{Y_1}$ let
$g^{(a)}$ be as in Lemma~\ref{asm-theta-agree-local-stabilizer-parts}, whose part
(1) holds for every lattice space with the same proof, and put $g^{(v)}=g^{(1_v)}$
and $t^{(v)}=1_vt+(1-1_v)$ for $t\in\mathbf A_E^\times$.

\emph{Step 1: the adelic group.} Let $\operatorname{ht}\colon U_Y\to\mathbf R_{>0}$
be the continuous function $q$ of
Lemma~\ref{asm-theta-agree-local-invariance-height}(3): $\operatorname{ht}(gk)=\operatorname{ht}(g)$
for $k\in K$, and $\operatorname{ht}(pg)=|t|^{n_1+1}\operatorname{ht}(g)$ for $p$
parabolic with multiplier $t$. By
Lemma~\ref{asm-theta-agree-local-invariance-integral}, for every continuous
$f_0$ on $U_Y$ with $f_0(pg)=|t|^{n_1+1}f_0(g)$ for such $p$, and every Haar
measure $\mu_K$ on $K$,
$\int_Kf_0(kg)\,d\mu_K(k)=\int_Kf_0(k)\,d\mu_K(k)$ for every $g\in U_Y$.
% Lean: htEx (UnramifiedCohomology/KAtCompact.lean), Lemma
% asm-theta-agree-local-invariance-height, with the data of Unf.parEx (Lemma
% asm-theta-agree-local-invariance-parabolic); and KInv.int (UnramifiedCohomology/Lint.lean),
% Lemma asm-theta-agree-local-invariance-integral, applied in kInv to f' (main).

\emph{Step 2: parts at $v$.} Let $p$ be parabolic with multiplier $t$. By
Lemma~\ref{asm-theta-agree-split-inputs}(1)(e), $p=um$ with $u\in N_H$ and $m$ a
Levi element for $(t,g_1)$. Then $p^{(v)}=u^{(v)}m^{(v)}$; $m^{(v)}$ is a Levi
element for $(t^{(v)},g_1^{(v)})$, as $1_v$ commutes with $c$; $u^{(v)}\in N_H$, as
$u\mapsto u^{(v)}$ is a homomorphism that maps the generators of $N_H$ to
generators; and $u^{(v)}$, $m^{(v)}$, $g_1^{(v)}$ and $t^{(v)}-1=1_v(t-1)$ are at
$v$. By hypothesis,
$f(p^{(v)}x)=|t^{(v)}|^{n_1+1}f(x)$ for $x\in U_Y^{(v)}$.

\emph{Step 3: a global function.} Let $\operatorname{ht}$ be as in Step 1, and put
$f'(x)=f(x^{(v)})\operatorname{ht}(x)/\operatorname{ht}(x^{(v)})$. It is continuous
on $U_Y$, as $x\mapsto x^{(v)}$ is continuous with values in $U_Y^{(v)}$. For $p$
parabolic with multiplier $t$, $(px)^{(v)}=p^{(v)}x^{(v)}$ and $p^{(v)}$ is parabolic
with multiplier $t^{(v)}$, so by Step 2
\[
f'(px)=\frac{|t^{(v)}|^{n_1+1}f(x^{(v)})\,|t|^{n_1+1}\operatorname{ht}(x)}{|t^{(v)}|^{n_1+1}\operatorname{ht}(x^{(v)})}=|t|^{n_1+1}f'(x).
\]
Let $g\in U_Y^{(v)}$ and $k\in K$. Then $k=k^{(v)}k^{(1-1_v)}$, the two factors lie
in $K$ (Lemma~\ref{asm-theta-agree-local-stabilizer-parts}(2)), and $k^{(1-1_v)}$
commutes with $g$, so $kg=k^{(v)}g\,k^{(1-1_v)}$; as $(kg)^{(v)}=k^{(v)}g$ and
$\operatorname{ht}$ is right $K$-invariant, $f'(kg)=f(k^{(v)}g)$.

\emph{Step 4: the projection.} The map $\pi\colon K\to K_v$, $k\mapsto k^{(v)}$, is a
continuous homomorphism onto $K_v$
(Lemma~\ref{asm-theta-agree-local-stabilizer-parts}), and $K$ is compact
(Lemma~\ref{asm-theta-agree-stabilizer}). The image under $\pi$ of a Haar measure
$\mu_K$ of $K$ is a Haar measure of $K_v$, so $\mu=c\,\pi_*\mu_K$ for some $c>0$.
For $g\in U_Y^{(v)}$, $k\mapsto f(kg)$ is continuous on $K_v$, and by Step 3 and
Step 1, applied to $f_0=f'$,
\[
\int_{K_v}f(kg)\,d\mu(k)=c\int_Kf(k^{(v)}g)\,d\mu_K(k)=c\int_Kf'(kg)\,d\mu_K(k)
=c\int_Kf'(k)\,d\mu_K(k)=\int_{K_v}f(k)\,d\mu(k).
\]
\end{proof}
% Lean: KInv.kInv (UnramifiedCohomology/KAtCompact.lean). Step 2 is part_levi and part_unip with
% the decomposition of htEx; Step 3 is f' (with cont_part), the identity e1 to e6 in the
% application of KInv.int, and val (part_mul, part_at, part_comm, htK); Step 4 is push and tr.

\begin{lemma}\label{asm-theta-agree-local-augmentation}
Let $B$ be a theta base of a CM field $E$, let $(G,G_1,H,H')$ be a
$\beta$-configuration with $G$ of ranks $(a,b+2)$, $e'$ an evaluation map for
$(M',M'_1,\iota_H,\iota_{H'},\mathcal J_\beta(H,H'))$ with components
$e'_\infty$ and $e'_v$, $K\subset U_{V'}$ the lattice stabilizer of $V'$, and
$m$ a normalized local measure system for $K$. Let
$(B,G,\mathcal X,\omega,\omega')$ and $(B,G_1,\mathcal X,\omega_1,\omega'_1)$ be
triples, and let $\operatorname{cone}_{v_1}$ be the local cone for
$(M',\omega',m,e'_{v_1},u_\beta,\nu,k_\beta+1)$ at the place $v_1$ of $B$.
\begin{enumerate}
\item If $g\in U_{V'}^{(v_1)}$, $\phi$ and $\phi_1$ are Schwartz-Bruhat and
  $\omega'(g,1)$ acts at $v_1$ by $\phi\mapsto\phi_1$, then
  $l(\operatorname{cone}_{v_1}(\phi_1))=l(\operatorname{cone}_{v_1}(\phi))$ for
  every linear form $l$ on functions on $F_{v_1}^{I_{M'_1}}$ that vanishes on
  the augmentation space $\mathfrak a_{v_1}(\omega'_1)$
  (Definition~\ref{asm-theta-agree-matched}(1)).
\item $\operatorname{cone}_{v_1}$ maps every Schwartz-Bruhat function on
  $F_{v_1}^{I_{M'}}$ of full rank at $\tau_1$ for $M'$ to one of full rank at
  $\tau_1$ for $M'_1$ (Definition~\ref{asm-theta-agree-models}(5)).
\end{enumerate}
\end{lemma}
\leanref{Development/WeilBound/ThetaLifts/LocalTheory/UnramifiedDeterminedBySpherical.lean}{PositivityTheory.ThKG.Loc.FLocB}
\leanref{Development/WeilBound/ThetaLifts/LocalTheory/UnramifiedCohomology/FLocBOf.lean}{PositivityTheory.ThKG.Loc.FL.fLocB_of}
\begin{proof}
Put $v=v_1$, $\nu=\alpha/\beta$ and $z=k_\beta+1$, so that $\operatorname{cone}_v$
is the local cone for $(M',\omega',m,e'_v,u_\beta,\nu,z)$. As $G$ is proper,
$a+b+2$ is even and $k_\beta=(a-b-2)/2>2(p+q)+1$, so $z\geq2$ and
$(a+b+2)/2=(b+1)+z$. As $G$ and $G_1$ are proper, $(M',M'_1,H,H',e',u_\beta)$ is a
cone setting for $B$ (Definition~\ref{asm-theta-agree-local-factor-setting}).

\emph{Step 1: the data.} The Weil pairs $(V',W,M',\alpha,\nu,\omega')$ and
$(V'_1,W_1,M'_1,\alpha,\nu,\omega'_1)$ of the two triples have signatures
$(q_1+1,b-q_1+1;r'+1,s'+1)$ and $(q_1,b-q_1;r',s')$, with $r'=a+b-p-q_1$ and
$s'=p+q_1$, as $q=q_1+1$ and $p_1=p$. So
Lemma~\ref{asm-theta-agree-split-inputs}(5), for $\mathcal J_\beta(H,H')$, gives
the hypothesis of Lemma~\ref{asm-theta-agree-local-augmentation-point} for $e'$,
$\omega'$ and $\omega'_1$, with $\chi_X=\alpha$, $\chi_W=\nu$, $r_X=(a+b+2)/2$ and
$r_W=b/2$. Let $(\omega_v,\rho_v)$ and $(\omega_{1,v},\rho_{1,v})$ be local
operators at $v$ for the two pairs
(Lemma~\ref{asm-theta-agree-finite-inputs}(2)). By
Lemma~\ref{asm-theta-agree-split-inputs}(3), $K$ is an Iwasawa subgroup for
$x_H$, which is isotropic and nonzero; in particular $K$ and $K_v$ are compact.

\emph{Step 2: the cone at $v$.} As in the proof of
Lemma~\ref{asm-theta-agree-local-factor-finite}(2), for Schwartz-Bruhat $\phi$,
\[
\operatorname{cone}_v(\phi)=\int_{K_v}A(\omega_v(k,1)\phi)\,dm_{K,v}(k),\qquad
A(\chi)=\int_{A_v}G_\chi(s)\,\nu(s)|s|^z\,dm_{A,v}(s),
\]
with $G_\chi(s)=e'_v(\rho_v(1_vu_\beta(s))\chi)$; the map
$k\mapsto\omega_v(k,1)\phi$ is locally constant on $U_{V'}^{(v)}$, so
$k\mapsto A(\omega_v(k,1)\phi)$ takes finitely many values on $K_v$, each on an
open set; and by Lemma~\ref{asm-theta-agree-local-factor-tate}, for
Schwartz-Bruhat $\chi$, $G_\chi$ takes finitely many values, is locally
constant, and $\nu(s)|s|^z$ is integrable on the set where $G_\chi\neq0$.

\emph{Step 3: part (1).} Let $g$, $\phi$, $\phi_1$ and $l$ be as in (1). By
Lemma~\ref{asm-theta-agree-finite-inputs-schur}(1), $\phi_1=\omega_v(g,1)\phi$. By
the last claim of Lemma~\ref{asm-theta-agree-local-factor-sums}, applied to the
two integrals of Step 2,
\[
l(\operatorname{cone}_v(\phi))=\int_{K_v}l(A(\omega_v(k,1)\phi))\,dm_{K,v}(k),\qquad
l(A(\chi))=\int_{A_v}l(G_\chi(s))\,\nu(s)|s|^z\,dm_{A,v}(s).
\]
Let $f(x)=l(A(\omega_v(x,1)\phi))$ for $x\in U_{V'}^{(v)}$, and $f(x)=0$ otherwise.
It is locally constant, hence continuous, on $U_{V'}^{(v)}$. Let $u\in N_H$ be at
$v$, $p$ a Levi element for $(t,g_1)$ with $p$ and $g_1$ at $v$ and $t-1$ at $v$,
and $x\in U_{V'}^{(v)}$; put $\chi=\omega_v(x,1)\phi$, so that
$\omega_v(upx,1)\phi=\omega_v(up,1)\chi$. By
Lemma~\ref{asm-theta-agree-local-augmentation-point}, for every $s$,
$l(G_{\omega_v(up,1)\chi}(s))=\nu(t)|t|^{(a+b+2)/2}\,l(G_\chi(ts))$. As $t\in A_v$
and $m_{A,v}$ is a Haar measure,
$\int l(G_\chi(s))\nu(s)|s|^z\,dm_{A,v}=\int l(G_\chi(ts))\nu(ts)|ts|^z\,dm_{A,v}(s)$.
Hence
\[
f(upx)=\nu(t)|t|^{(a+b+2)/2}\int l(G_\chi(ts))\,\nu(s)|s|^z\,dm_{A,v}(s)
=|t|^{(a+b+2)/2-z}f(x)=|t|^{b+1}f(x).
\]
The space $V'_1$ has rank $b$, so
Lemma~\ref{asm-theta-agree-local-augmentation-invariance}, for $Y=V'$ and the
Haar measure $m_{K,v}$, gives $\int_{K_v}f(kg)\,dm_{K,v}=\int_{K_v}f(k)\,dm_{K,v}$. As
$f(kg)=l(A(\omega_v(k,1)\phi_1))$, this says
$l(\operatorname{cone}_v(\phi_1))=l(\operatorname{cone}_v(\phi))$.

\emph{Step 4: part (2).} Let $\phi$ be Schwartz-Bruhat of full rank at $\tau_1$
for $M'$, and let $N_1$ be the span of the functions of full rank at $\tau_1$ for
$M'_1$. Choose $f$ with $e'_\infty(f)\neq0$. Let $k\in K_v$ and
$s\in\mathbf A_E^\times$. The function
$\Phi=f\otimes\phi\otimes\bigotimes_{w\neq v}\mathrm{vac}'_w$ lies in
$\mathcal S^\circ(M')$, so $\omega'(k,1)\Phi$ does, by
Lemma~\ref{asm-theta-agree-split-inputs}(6); as $\omega'(k,1)$ acts at $v$ by
$\phi\mapsto\omega_v(k,1)\phi$, it is
$f\otimes\omega_v(k,1)\phi\otimes\bigotimes_{w\neq v}\mathrm{vac}'_w$. By
Lemma~\ref{asm-theta-agree-split-inputs}(7)(a), with $a=1_vc(s)$, so that
$a\,yx'^{\mathsf T}=1_vu_\beta(s)$,
\[
e'\bigl(\rho(1_vu_\beta(s))\,\omega'(k,1)\Phi\bigr)
=e'_\infty(f)\otimes G_{\omega_v(k,1)\phi}(s)\otimes\bigotimes_{w\neq v}\mathrm{vac}'_{1,w}
\]
lies in $\mathcal S^\circ(M'_1)$; the equality holds as $\rho(1_vu_\beta(s))$ acts at
$v$ by $\rho_v(1_vu_\beta(s))$ and by the components of $e'$. By
Lemma~\ref{asm-theta-agree-local-augmentation-slice},
$G_{\omega_v(k,1)\phi}(s)\in N_1$. By the second claim of
Lemma~\ref{asm-theta-agree-local-factor-sums}, applied to the two integrals of
Step 2, $A(\omega_v(k,1)\phi)\in N_1$ for every $k$, and then
$\operatorname{cone}_v(\phi)\in N_1$. By
Lemma~\ref{asm-theta-agree-local-augmentation-subspace}(2), for the Weil pair
$(V'_1,W_1,M'_1,\alpha,\nu,\omega'_1)$, $\operatorname{cone}_v(\phi)$ has full rank
at $\tau_1$ for $M'_1$.
\end{proof}
% Lean: FL.fLocB_of (UnramifiedCohomology/FLocBOf.lean). Step 1 is hc, TriplePair (hP, hP1),
% the arithmetic hz, hzr, MixAB (hx), locRepEx, IwaEx (hIK) and kAt_compact; Step 2 is op, A,
% tate' (hA), lc, lcK and ce; Step 3 is cl, al (int_lin), point and KInv.kInv; Step 4 is gen, up,
% pt (SoStab, SoEv, so_slice), int_fin twice and span_soLoc (LinEx). What remains open is
% Lemma asm-theta-agree-local-augmentation-invariance (KInv.kInv, about 2,800 lines) and the
% sublemmas of Lemma asm-theta-agree-local-factor that Lemma asm-theta-agree-local-factor-tate
% uses (ideal, ideles).

\subsection{Local multiplicity one}\label{subsection-asm-local-multiplicity-one}

\begin{definition}\label{asm-theta-agree-local-equivariant}
Let $B$ be a theta base of a CM field $E$, and let
$Q=(B,G,\mathcal X,\omega,\omega',R,R_f)$ and
$Q_1=(B,G_1,\mathcal X,\omega_1,\omega'_1,R_1,R_{f,1})$ be matched quadruples
(Definition~\ref{asm-theta-agree-matched}) of ranks $(a,a')$ and $(a_1,a'_1)$,
with real maps $(\Pi,p,p')$ and $(\Pi_1,p_1,p'_1)$ and finite maps
$(\Pi_v,p_v,p'_v)_v$ and $(\Pi_{1,v},p_{1,v},p'_{1,v})_v$. Let $H'$ be an
adapted plane in $W$ over $W_1$, the skew-Hermitian spaces of $G$ and $G_1$,
let $\chi$ be a Hecke character of $E$ and $r\in\mathbf R$.
\begin{enumerate}
\item $(Q,Q_1)$ has \emph{step type} $(\chi,r)$ if either $a_1=a$,
  $a'_1+2=a'$, $p_1=p$, $q_1+1=q$, $\chi=\beta$ and $r=a/2$, or $a_1+2=a$,
  $a'_1=a'$, $p_1+1=p$, $q_1=q$, $\chi=\alpha$ and $r=a'/2$, where $p,q$ and
  $p_1,q_1$ are the integers of $G$ and $G_1$.
\item Let $v$ be a finite place. A linear map $\lambda\colon\Pi_v\to\Pi_{1,v}$
  is \emph{equivariant at $v$} for $(H',\chi,r)$ if the following holds. Let
  $m$ be a Levi element for $(t,h_1)$ and $u\in N_{H'}$ with $u$ and $m$ at
  $v$, and let $\phi,\phi',\phi_1,\phi'_1$ be Schwartz-Bruhat functions such
  that $\omega(1,um)$ acts at $v$ by $\phi\mapsto\phi'$, $\omega_1(1,h_1)$
  acts at $v$ by $\phi_1\mapsto\phi'_1$, and $\lambda(p_v(\phi))=p_{1,v}(\phi_1)$.
  Then $\lambda(p_v(\phi'))=\chi(t)|t|^r\,p_{1,v}(\phi'_1)$.
\item A linear map $\lambda\colon\Pi\to\Pi_1$ is \emph{equivariant at $\infty$}
  for $(H',\chi,r)$ if the following hold.
  \begin{enumerate}
  \item[(a)] If $\gamma$ is a line in $U_W$ with values in $N_{H'}$,
    $f\in\operatorname{Pol}(M)$ and $t\mapsto\omega(1,\gamma(t))$ acts at
    $\infty$ to first order by $f\mapsto f'$, then $\lambda(p(f'))=0$.
  \item[(b)] Let $\gamma$ and $\gamma_1$ be lines in $U_W$ and $U_{W_1}$ and
    $\tau\colon\mathbf R\to\mathbf A_E^\times$ a map such that $\gamma(s)$ is
    a Levi element for $(\tau(s),\gamma_1(s))$ for every $s$, and let $c$ be
    the derivative at $0$ of $s\mapsto\chi(\tau(s))|\tau(s)|^r$. If
    $f\in\operatorname{Pol}(M)$, $f_1\in\operatorname{Pol}(M_1)$,
    $t\mapsto\omega(1,\gamma(t))$ and $t\mapsto\omega_1(1,\gamma_1(t))$ act at
    $\infty$ to first order by $f\mapsto f'$ and $f_1\mapsto f'_1$, and
    $\lambda(p(f))=p_1(f_1)$, then $\lambda(p(f'))=c\,p_1(f_1)+p_1(f'_1)$.
  \item[(c)] Let $k\in K_\infty(W)$ and $k_1\in K_\infty(W_1)$ with $k$ a Levi
    element for $(t,k_1)$. If $f\in\operatorname{Pol}(M)$,
    $f_1\in\operatorname{Pol}(M_1)$, $\omega(1,k)$ and $\omega_1(1,k_1)$ act at
    $\infty$ by $f\mapsto f'$ and $f_1\mapsto f'_1$, and $\lambda(p(f))=p_1(f_1)$,
    then $\lambda(p(f'))=\chi(t)|t|^r\,p_1(f'_1)$.
  \end{enumerate}
\end{enumerate}
\end{definition}
\leanref{Development/WeilBound/ThetaLifts/LocalTheory/UnramifiedDeterminedBySpherical.lean}{PositivityTheory.ThKG.Loc.StepTy}
\leanref{Development/WeilBound/ThetaLifts/LocalTheory/IsArchimedeanLocalEquivariantMap.lean}{PositivityTheory.ThKG.Loc.IsFiniteLocalEquivariantMap}
\leanref{Development/WeilBound/ThetaLifts/LocalTheory/IsArchimedeanLocalEquivariantMap.lean}{PositivityTheory.ThKG.Loc.IsArchimedeanLocalEquivariantMap}

\begin{theorem}[Iwasawa decomposition \cite{BT72,Car79}]\label{asm-theta-agree-local-mult-unramified-iwasawa}
Let $B$ be a theta base of a CM field $E$
(Definition~\ref{asm-theta-agree-models}(1)), let $Y$ be a lattice space of
rank $n$ with matrix $J$ that is Hermitian or skew-Hermitian of some signature
(Definition~\ref{asm-theta-agree-doubling}(1)), let $x\in E^n$ be nonzero with
$x^{\mathsf T}J\,c(x)=0$, and let $v$ be a finite place of $F$. Then every
$g\in U_Y^{(v)}$ is $pk$ with $k\in K_Y(v)$
(Definition~\ref{asm-theta-agree-finite-local-objects}(1)) and
$p\in U_Y^{(v)}$ such that $px=tx$ for some $t\in\mathbf A_E^\times$.
\end{theorem}
\leanref{Development/WeilBound/ThetaLifts/LocalTheory/UnramifiedCohomology/MemLatWIff.lean}{PositivityTheory.ThKG.Loc.Iwa.iwaLoc}

\begin{lemma}\label{asm-theta-agree-local-mult-unramified-levi}
Let $B$ be a theta base of a CM field $E$, let $W$ be a skew-Hermitian lattice
space of rank $d$ and some signature, let $H=(x,y,\iota,\mathfrak b)$ be an
adapted plane in $W$ over a lattice space $W_1$
(Definition~\ref{asm-theta-agree-planes}(1)), and let $v$ be a finite place. If
$p\in U_W^{(v)}$ and $t\in\mathbf A_E^\times$ have $px=tx$, then $p=um$ with
$u\in N_H\cap U_W^{(v)}$ and $m\in U_W^{(v)}$ a Levi element for $(t,h_1)$
(Definition~\ref{asm-theta-agree-planes}(2)) for some $h_1\in U_{W_1}^{(v)}$.
\end{lemma}
\leanref{Development/WeilBound/ThetaLifts/LocalTheory/UnramifiedCohomology/UnramifiedDeterminedBySphericalHolds.lean}{PositivityTheory.ThKG.Loc.parDec}
\leanref{Development/WeilBound/ThetaLifts/TowerModels/LeviAlg.lean}{PositivityTheory.ThKG.Tow.Levi.IsFr.stab}
\begin{proof}
Let $J$ and $J_1$ be the matrices of $W$ and $W_1$, and let $A=(\iota\ x\ y)$
and $J^\natural$ be as in the proof of Lemma~\ref{asm-theta-agree-levi}:
$A\in GL_d(E)$ and
$A^{\mathsf T}J\,c(A)=J^\natural=\operatorname{diag}\bigl(J_1,\left(\begin{smallmatrix}0&1\\-1&0\end{smallmatrix}\right)\bigr)$.
Index the rows and columns of $J^\natural$ by the columns of $A$, and put
$P=A^{-1}pA$. Then $P^{\mathsf T}J^\natural c(P)=J^\natural$, and the column
of $P$ at $x$ is $t$ times the unit vector at $x$, as $px=tx$.

\emph{The block of $\iota$.} For an index $j$ of the block of $\iota$, the entry
of $P^{\mathsf T}J^\natural c(P)$ in row $j$ and column $x$ is $-P_{yj}c(t)$,
and that of $J^\natural$ is $0$. So $P_{yj}=0$, and with $h_1$ the block of $P$
on the indices of $\iota$ and $\alpha$ the row of $P$ at $x$ on these indices,
\[
p\iota z=\iota h_1z+(\alpha z)\,x\qquad(z\in\mathbf A_E^{d-2}).
\]
On the block of $\iota$, the equation $P^{\mathsf T}J^\natural c(P)=J^\natural$
reads $h_1^{\mathsf T}J_1c(h_1)=J_1$, since $P_{yj}=0$ and the entries of
$J^\natural$ at $(x,x)$ and $(y,y)$ vanish. As $J$ is invertible
(Lemma~\ref{asm-theta-agree-convergence-signature}), so are $J^\natural$ and
$J_1$; taking determinants, $\det h_1$ is a unit, and $h_1\in U_{W_1}$. Let $m$ be a Levi element for $(t,h_1)$
(Lemma~\ref{asm-theta-agree-levi}), and put $u=pm^{-1}$.

\emph{At $v$.} As $p-1$ is at $v$, $(1-1_v)(p-1)=0$. Applied to $x$ and to
$\iota z$ this gives $(1-1_v)(t-1)x=0$ and
$(1-1_v)\bigl(\iota(h_1-1)z+(\alpha z)x\bigr)=0$. The columns of $A$ are
linearly independent over $\mathbf A_E$, as $A$ is invertible, so $t-1$ is at
$v$, and so is $(h_1-1)z$ for every $z$: $h_1\in U_{W_1}^{(v)}$. Then
$c(t)^{-1}-1=c(t)^{-1}c(1-t)$ is at $v$. The matrix $m-1$ maps $\iota z$, $x$
and $y$ to $\iota(h_1-1)z$, $(t-1)x$ and $(c(t)^{-1}-1)y$, so $(1-1_v)(m-1)A=0$
and $m\in U_W^{(v)}$. Hence $u=pm^{-1}\in U_W^{(v)}$.

\emph{$u\in N_H$.} As $m^{-1}x=t^{-1}x$, $ux=p(t^{-1}x)=x$. As
$m^{-1}\iota h_1z=\iota z$, $u\iota h_1z=p\iota z=\iota h_1z+(\alpha z)x$; as
$h_1$ is invertible, $u\iota z'\in\iota z'+\mathbf A_Ex$ for every $z'$. So $u$
is one of the generators of $N_H$.
\end{proof}

\begin{lemma}\label{asm-theta-agree-local-mult-unramified-vacuum}
Let $(X,W,M,\chi_X,\chi_W,\omega)$ be a Weil pair for a theta base $B$ of a CM
field $E$ (Definition~\ref{asm-theta-agree-triple}(1)), let $v$ be a finite
place, let $(\omega_v,\rho_v)$ be local operators at $v$ for the pair
(Definition~\ref{asm-theta-agree-finite-local-objects}(4)), and let
$\mathrm{vac}_v$ be the vacuum of $M$ at $v$. Then for every $k\in K_W(v)$
there is $c\in\mathbf C$ with $\omega_v(1,k)\mathrm{vac}_v=c\,\mathrm{vac}_v$.
\end{lemma}
\leanref{Development/WeilBound/CoreStatements/CompactAveragingSplitting/EigenvectorLifting/InertDef.lean}{PositivityTheory.ThK.VacEigen}
\leanref{Development/WeilBound/CoreStatements/CompactAveragingSplitting/EigenvectorLifting/Vaceigen.lean}{PositivityTheory.ThK.vacuum_eigenline_of_local}
\leanref{Development/WeilBound/CoreStatements/CompactAveragingSplitting/EigenvectorLifting/Vaceigen.lean}{PositivityTheory.ThK.Vac.latAtV_mul}
\begin{proof}
Let $\Lambda_v$ be as in Lemma~\ref{asm-theta-agree-evaluation-lattice}, let
$\theta\in\mathcal O_E$ be the element in the definition of the vacuum
$\mathrm{vac}_v$ (Definition~\ref{asm-theta-agree-models}(3)), and put
$\chi(\gamma)=\psi(\tfrac12\operatorname{Tr}_{\mathbf A_E/\mathbf A}h(\theta\gamma,\gamma))$
for $\gamma\in\Lambda_v$. Then $\rho(\gamma)$ acts at $v$ by
$\mathrm{vac}_v\mapsto\chi(\gamma)\mathrm{vac}_v$, and by
$\mathrm{vac}_v\mapsto\rho_v(\gamma)\mathrm{vac}_v$. By
Lemma~\ref{asm-theta-agree-finite-inputs-slice}(4) such a function is unique, so
$\rho_v(\gamma)\mathrm{vac}_v=\chi(\gamma)\mathrm{vac}_v$.

Let $k\in K_W(v)$ and put $\phi'=\omega_v(1,k)\mathrm{vac}_v$, a
Schwartz-Bruhat function. Let $\gamma\in\Lambda_v$ and
$\gamma'=\gamma(k^{-1})^{\mathsf T}$. Then $\gamma'\in\Lambda_v$: for a
generator $a\,x'y^{\mathsf T}$ of $\Lambda_v$,
$a\,x'y^{\mathsf T}(k^{-1})^{\mathsf T}=x'(k^{-1}ay)^{\mathsf T}$, and
$k^{-1}ay$ lies in the group $L_v$ of
Definition~\ref{asm-theta-agree-finite-local-objects}(1), whose elements are
sums of vectors $a'y'$ with $a'$ at $v$, $a'_v\in\mathcal O_v$ and $y'$ in the
lattice of $W$. As $k^{-1}\in U_W$, $h(\theta\gamma',\gamma')=h(\theta\gamma,\gamma)$,
so $\chi(\gamma')=\chi(\gamma)$. By the commutation rule of local operators,
\[
\rho_v(\gamma)\phi'=\rho_v(\gamma'k^{\mathsf T})\,\omega_v(1,k)\mathrm{vac}_v
=\omega_v(1,k)\rho_v(\gamma')\mathrm{vac}_v=\chi(\gamma)\phi'.
\]
So for every $\gamma\in\Lambda_v$, $\rho(\gamma)$ acts at $v$ by
$\mathrm{vac}_v\mapsto\chi(\gamma)\mathrm{vac}_v$ and by
$\phi'\mapsto\chi(\gamma)\phi'$, and $\mathrm{vac}_v\neq0$. The last two
paragraphs of the proof of Lemma~\ref{asm-theta-agree-evaluation-vacuum} use
only these properties of the two functions there, and they give
$\phi'=c\,\mathrm{vac}_v$ for some $c$.
\end{proof}
% Lean: vacuum_eigenline_of_local (Vaceigen.lean): hv and hrho give the first paragraph,
% latAtV_mul and hTens_mul the invariance of Lambda_v and chi, L.impl the display, and
% local_line (the last paragraph of the proof of Lemma asm-theta-agree-evaluation-vacuum)
% the conclusion.

\begin{lemma}\label{asm-theta-agree-local-mult-unramified}
Let $B$ be a theta base of a CM field $E$, let $Q$ and $Q_1$ be matched
quadruples of step type $(\chi,r)$, and let $H'$ be an adapted plane in $W$
over $W_1$ at some positions
(Definition~\ref{asm-theta-agree-local-equivariant}). Let $v\notin T$ be a
finite place, let $\lambda,\lambda'\colon\Pi_v\to\Pi_{1,v}$ be equivariant at
$v$ for $(H',\chi,r)$, and let $c\in\mathbf C$ with
$\lambda'(p_v(\mathrm{vac}_v))=c\,\lambda(p_v(\mathrm{vac}_v))$. Then
$\lambda'=c\lambda$.
\end{lemma}
\leanref{Development/WeilBound/ThetaLifts/LocalTheory/UnramifiedMultiplicityOneProportional.lean}{PositivityTheory.ThKG.Loc.UnramifiedMultiplicityOneProportional}
\leanref{Development/WeilBound/ThetaLifts/LocalTheory/UnramifiedCohomology/UnramifiedDeterminedBySphericalHolds.lean}{PositivityTheory.ThKG.Loc.unramifiedProportional_holds}
\begin{proof}
Let the triples of $Q$ and $Q_1$ be $(B,G,\mathcal X,\omega,\omega')$ and
$(B,G_1,\mathcal X,\omega_1,\omega'_1)$, with $G$ of ranks $(a,a')$, spaces
$V,V',W$ and models $M,M'$, and $G_1$ with spaces $V_1,V'_1,W_1$ and model
$M_1$. The step type is not used. Put $\mu=\lambda'-c\lambda$; we show $\mu=0$.
By Definition~\ref{asm-theta-agree-triple}(4) and
Lemma~\ref{asm-theta-agree-finite-inputs}(2), the Weil pairs
$(V,W,M,\ldots,\omega)$, $(V',W,M',\ldots,\omega')$ and
$(V_1,W_1,M_1,\ldots,\omega_1)$ have local operators $\omega_v$, $\omega'_v$ and
$\omega_{1,v}$ at $v$ (Definition~\ref{asm-theta-agree-finite-local-objects}(4)).
By Lemma~\ref{asm-theta-agree-finite-inputs-slice}(4), a function $\phi_1$ with
an operator acting at $v$ by $\phi\mapsto\phi_1$ is unique, so for
Schwartz-Bruhat $\phi$ and $\phi_1$, $\omega(1,h)$ acts at $v$ by
$\phi\mapsto\phi_1$ if and only if $\phi_1=\omega_v(1,h)\phi$.

\emph{Step 1: kernels.} Let $h\in U_W^{(v)}$ and $\phi$ Schwartz-Bruhat with
$p_v(\phi)=0$. Then $p_v(\omega_v(1,h)\phi)=0$. Indeed $p_v(\phi)=0=p'_v(0)$,
$\omega(1,h)$ acts at $v$ by $\phi\mapsto\omega_v(1,h)\phi$ and $\omega'(1,h)$ by
$0\mapsto0$, so by Definition~\ref{asm-theta-agree-matched}(3),
$p_v(\omega_v(1,h)\phi)=p'_v(0)=0$.

\emph{Step 2: the space $S$.} Let $S$ be the space of Schwartz-Bruhat $\phi$ on
$F_v^{I_0}$ with $\mu(p_v(\omega_v(1,h)\phi))=0$ for every $h\in U_W^{(v)}$.

\emph{Step 3: $\mathrm{vac}_v\in S$.} Let $h\in U_W^{(v)}$, and let
$H'=(x,y,\iota,\mathfrak b)$. As $h(x,x)=0$ and $h(x,y)=1$, $x$ is a nonzero
isotropic vector. By
Theorem~\ref{asm-theta-agree-local-mult-unramified-iwasawa}, $h=pk$ with
$k\in K_W(v)$, $p\in U_W^{(v)}$ and $px=tx$ for some $t\in\mathbf A_E^\times$.
By Lemma~\ref{asm-theta-agree-local-mult-unramified-levi}, $p=um$ with $u\in N_{H'}$
and $m\in U_W^{(v)}$ a Levi element for $(t,h_1)$, where $u$ and $h_1$ are at
$v$. By Lemma~\ref{asm-theta-agree-local-mult-unramified-vacuum},
$\omega_v(1,k)\mathrm{vac}_v=c'\mathrm{vac}_v$ for some $c'$, so
$\omega_v(1,h)\mathrm{vac}_v=c'\,\omega_v(1,um)\mathrm{vac}_v$. As $p_{1,v}$ maps
the Schwartz-Bruhat functions onto $\Pi_{1,v}$, there is a Schwartz-Bruhat $b$
with $p_{1,v}(b)=\lambda(p_v(\mathrm{vac}_v))$; then
$p_{1,v}(cb)=\lambda'(p_v(\mathrm{vac}_v))$. The operators $\omega(1,um)$ and
$\omega_1(1,h_1)$ act at $v$ by $\mathrm{vac}_v\mapsto\omega_v(1,um)\mathrm{vac}_v$,
$b\mapsto\omega_{1,v}(1,h_1)b$ and $cb\mapsto c\,\omega_{1,v}(1,h_1)b$. The
equivariance of $\lambda$ and of $\lambda'$
(Definition~\ref{asm-theta-agree-local-equivariant}(2)) gives
\[
\lambda\bigl(p_v(\omega_v(1,um)\mathrm{vac}_v)\bigr)=\chi(t)|t|^r\,p_{1,v}(\omega_{1,v}(1,h_1)b),\qquad
\lambda'\bigl(p_v(\omega_v(1,um)\mathrm{vac}_v)\bigr)=c\,\chi(t)|t|^r\,p_{1,v}(\omega_{1,v}(1,h_1)b).
\]
So $\mu(p_v(\omega_v(1,h)\mathrm{vac}_v))=c'\cdot0=0$.

\emph{Step 4: $p_v(S)$ is stable.} Let $h\in U_W^{(v)}$, let $\phi,\phi_1$ be
Schwartz-Bruhat with $\omega(1,h)$ acting at $v$ by $\phi\mapsto\phi_1$, and
suppose $p_v(\phi)=p_v(\psi)$ with $\psi\in S$. Then
$\phi_1=\omega_v(1,h)\phi$. The function $\omega_v(1,h)\psi$ lies in $S$, since
$\omega_v(1,h')\omega_v(1,h)\psi=\omega_v(1,h'h)\psi$ and $h'h\in U_W^{(v)}$ for
$h'\in U_W^{(v)}$. By Step 1 for $\psi-\phi$,
$p_v(\omega_v(1,h)\psi)=p_v(\phi_1)$. So $p_v(\phi_1)\in p_v(S)$.

\emph{Step 5: conclusion.} By Step 4 and the irreducibility condition of
Definition~\ref{asm-theta-agree-matched}(3), $p_v(S)$ is $0$ or $\Pi_v$. It
contains $p_v(\mathrm{vac}_v)$ by Step 3, which is nonzero as $v\notin T$. So
$p_v(S)=\Pi_v$. Every $\xi\in\Pi_v$ is then $p_v(\psi)$ with $\psi\in S$, and
$\mu(\xi)=\mu(p_v(\omega_v(1,1)\psi))=0$.
\end{proof}
% Lean: unramifiedProportional_holds (UnramifiedCohomology/UnramifiedDeterminedBySphericalHolds.lean),
% with lam = l2 - c0 l1 and the local representations of ThKH.Locrep.locRepEx for the three
% pairs (PairOK). Step 1 is ker (hM.fin.equiv with phi' = 0); Step 3 is hS, with Iwa.iwaLoc,
% parDec and vacuum_eigenline_of_local; Steps 4 and 5 are hN (hM.fin.irred) and hM.fin.sph.

\begin{situation}\label{asm-theta-agree-local-mult-ramified-situation}
Let $k$ be a nonarchimedean local field with ring of integers $\mathcal O$, let
$n\leq d_1$, put $d=d_1+2$, and index the columns of $n\times d$ matrices by
$\{0,\ldots,d_1-1\}\sqcup\{x,y\}$ through a fixed bijection with
$\{0,\ldots,d-1\}$. For $Y\in M_{n\times d}(k)$ let $Y_{\mathrm m}\in M_{n\times d_1}(k)$
be its columns in $\{0,\ldots,d_1-1\}$, and $Y_x,Y_y\in k^n$ its columns $x$ and
$y$. Let $\nu$ be a unitary smooth character of $k^\times$
(Definition~\ref{asm-theta-agree-coinvariants-local}) with $\nu(l)\neq1$ for some
$l\in\mathcal O^\times$, and write $\omega=\omega_{\nu,1}$, so that
$\omega(g,1)\Phi(Y)=\nu(\det g)|\det g|^{d/2}\Phi(g^{\mathsf T}Y)$. Let
$P^\circ$ be the set of $h\in M_d(k)$ whose column $x$ is $h_{xx}e_x$ with
$h_{xx}\in\mathcal O^\times$, with $h_{yy}\in\mathcal O^\times$, whose block on
the rows and columns $\{0,\ldots,d_1-1\}$ is the identity, and whose row $y$
vanishes in the columns $\{0,\ldots,d_1-1\}$. For $g\in GL_{d_1}(k)$ let
$\operatorname{blk}(g)\in GL_d(k)$ act as $g$ on the columns
$\{0,\ldots,d_1-1\}$ and fix $e_x$ and $e_y$. Let $\mathbf 1_{\mathcal O}$ be the
indicator function of $\mathcal O^n$, and for a function $\psi$ on
$M_{n\times d_1}(k)$ put
$j(\psi)(Y)=\psi(Y_{\mathrm m})\mathbf 1_{\mathcal O}(Y_x)\mathbf 1_{\mathcal O}(Y_y)$.
Let $\mathfrak a^\circ$ be the span of the $\omega_{\nu,1}(g,1)\psi-\psi$ on
$M_{n\times d_1}(k)$, with $g\in GL_n(k)$ and $\psi$ of full rank. Consider a
linear map $L$ from the functions on $M_{n\times d}(k)$ to a complex vector
space with:
\begin{enumerate}
\item[(a)] $L(\omega(g,1)\Phi)=L(\Phi)$ for $g\in GL_n(k)$ and Schwartz-Bruhat
  $\Phi$;
\item[(b)] $L(\Phi(\,\cdot\,h))=L(\Phi)$ for $h\in P^\circ$ and Schwartz-Bruhat
  $\Phi$.
\end{enumerate}
\end{situation}
\leanref{Development/WeilBound/ThetaLifts/LocalTheory/DixmierCommutants/KilEmpty.lean}{PositivityTheory.ThKG.LocalRamifiedLeviRestriction.InP}
\leanref{Development/WeilBound/ThetaLifts/LocalTheory/DixmierCommutants/KilEmpty.lean}{PositivityTheory.ThKG.LocalRamifiedLeviRestriction.Inv}
\leanref{Development/WeilBound/ThetaLifts/LocalTheory/DixmierCommutants/LocalMultiplicityOneIrreducibleGL.lean}{PositivityTheory.ThKG.LocalRamifiedLeviRestriction.tn}
\leanref{Development/WeilBound/ThetaLifts/LocalTheory/DixmierCommutants/LocalMultiplicityOneIrreducibleGL.lean}{PositivityTheory.ThKG.LocalRamifiedLeviRestriction.blk}
\leanref{Development/WeilBound/CoreStatements/LocallyConstantCoinvariants/FullOfSoLoc.lean}{PositivityTheory.KCo.augF}

\begin{lemma}\label{asm-theta-agree-local-mult-ramified-reduction}
In Situation~\ref{asm-theta-agree-local-mult-ramified-situation}, let $L$
satisfy (a) and (b).
\begin{enumerate}
\item If $L(j(\psi))=0$ for every $\psi$ of full rank on $M_{n\times d_1}(k)$,
  then $L(\Phi)=0$ for every Schwartz-Bruhat $\Phi$ on $M_{n\times d}(k)$.
\item $L\circ j$ vanishes on $\mathfrak a^\circ$.
\end{enumerate}
\end{lemma}
\leanref{Development/WeilBound/ThetaLifts/LocalTheory/DixmierCommutants/LocalMultiplicityOneIrreducibleGL.lean}{PositivityTheory.ThKG.LocalRamifiedLeviRestriction.red}
\leanref{Development/WeilBound/ThetaLifts/LocalTheory/DixmierCommutants/LocalMultiplicityOneIrreducibleGL.lean}{PositivityTheory.ThKG.LocalRamifiedLeviRestriction.aug_ker}
\begin{proof}
For a subspace $Z$ of the functions on $M_{n\times d}(k)$, say that $Z$
\emph{kills} a set $U$ if it contains every Schwartz-Bruhat function that
vanishes outside $U$. If every point of an open set $U$ has a neighbourhood that
$Z$ kills, then $Z$ kills $U$: a Schwartz-Bruhat function vanishing outside $U$
is a finite sum of indicator functions of small cosets of a compact open
subgroup, each inside such a neighbourhood; call this the first principle. For
the second principle, let $Z$ kill a set $U_0$, let $X_0$ be a matrix, let
$P(X)\in M_n(k)$ and $Q(X)\in M_d(k)$ depend continuously on $X$ near $X_0$, with
$P(X_0)$ and $Q(X_0)$ idempotent, and let $c\neq1$. Put
$T_X(S)=(1+(l^{-1}-1)P(X))\,S\,(1+(l-1)Q(X))$. Suppose that every $X$ near $X_0$
lies in $U_0$ or has $P(X)$ and $Q(X)$ idempotent, $P(X)X=XQ(X)$, and
$\Phi\circ T_X-c\Phi\in Z$ for every Schwartz-Bruhat $\Phi$. Then $Z$ kills a
neighbourhood of $X_0$. Indeed, near $X_0$ a Schwartz-Bruhat function is a
combination of indicator functions of cosets $X+\Lambda$ of a small compact open
subgroup $\Lambda$ that every $T_X$ preserves. Such a coset lies in $U_0$, or
contains an $X\notin U_0$; then $T_X$ fixes $X$, as $P(X)X=XQ(X)$ and $P(X)$ and
$Q(X)$ are idempotent, so the indicator function $\varphi$ of $X+\Lambda$ has
$\varphi\circ T_X=\varphi$ and $(1-c)\varphi\in Z$.

For $x,c\in k^n$ with $c\cdot x=1$ and $h\in P^\circ$, the matrix
$g^{\mathsf T}=1+(l^{-1}-1)xc^{\mathsf T}$ has $\det g=l^{-1}$, and (a) for
$\Phi(\,\cdot\,h)$ together with (b) gives
$L\bigl(Y\mapsto\Phi(g^{\mathsf T}Yh)\bigr)=\nu(l)L(\Phi)$. Apply the second
principle with $Z=\ker L$ and $c=\nu(l)$.

\emph{Step 1: $\ker L$ kills the set where $Y_x\neq0$.} Let
$(Y_0)_{x,i}\neq0$, and for $Y$ near $Y_0$ put $c(Y)=e_i/Y_{x,i}$,
$P(Y)=Y_xc(Y)^{\mathsf T}$ and $Q(Y)=e_x\,(c(Y)^{\mathsf T}Y)$, with $U_0$ empty.
Then $P(Y)Y=YQ(Y)$, and $1+(l-1)Q(Y)$ has column $x$ equal to $le_x$, entry $1$
at $(y,y)$, identity middle block and row $y$ zero, so it lies in $P^\circ$. So
$\ker L$ kills a neighbourhood of $Y_0$. In particular, for Schwartz-Bruhat
$\psi$ on $M_{n\times d_1}(k)$ and $f,f_1,f'$ on $k^n$ with $f(0)=cf_1(0)$, the
function $Y\mapsto\psi(Y_{\mathrm m})(f(Y_x)-cf_1(Y_x))f'(Y_y)$ vanishes where
$Y_x=0$, so $L$ takes the same value on $\psi(Y_{\mathrm m})f(Y_x)f'(Y_y)$ and
on $c\,\psi(Y_{\mathrm m})f_1(Y_x)f'(Y_y)$.

\emph{Step 2: the column $y$.} Let $\psi$ be of full rank, $f$ Schwartz-Bruhat
on $k^n$, and write $L_\psi(f')=L(Y\mapsto\psi(Y_{\mathrm m})f(Y_x)f'(Y_y))$.
Then $L_\psi(f'(\,\cdot\,+w))=L_\psi(f')$ for $w\in k^n$: near a point $Z_0$ with
$Z_0R=1$, the element $h=1+\tilde w e_y^{\mathsf T}$ of $P^\circ$, where $\tilde w$
has middle part $Rw$ and entries $0$ at $x$ and $y$, adds $Y_{\mathrm m}Rw$ to
the column $y$, and $Y_{\mathrm m}Rw$ is close to $w$; so (b) and the first
principle give the claim. A translation-invariant linear form on the
Schwartz-Bruhat functions on $k^n$ is a multiple of the Haar integral, so
$|\det g|\,L_\psi(f'(g\,\cdot\,))=L_\psi(f')$ for $g\in GL_n(k)$.

\emph{Step 3: $\ker L$ kills the set where $Y_{\mathrm m}$ has rank $n$.} Near a
point where $Y_{\mathrm m}$ has rank $n$, a Schwartz-Bruhat function is a
combination of indicator functions of small cosets, each of the form
$\psi(Y_{\mathrm m})f(Y_x)f'(Y_y)$ with $\psi$ of full rank and $f$ and $f'$
indicator functions of cosets $z+\varpi^m\mathcal O^n$. By Step 1 the factor $f$
can be replaced by $f(0)\mathbf 1_{\mathcal O}$, and by Step 2 the factor $f'$
by a multiple of $\mathbf 1_{\mathcal O}$. The resulting value is a multiple of
$L(j(\psi))=0$.

\emph{Step 4: $\ker L$ kills the matrices of rank $n$.} By Steps 1 and 3 and the
first principle, $\ker L$ kills the set $U_0$ where $Y_x\neq0$ or
$Y_{\mathrm m}$ has rank $n$. Let $Y_0R=1$, and for $Y$ near $Y_0$ put
$c(Y)^{\mathsf T}=R_y^{\mathsf T}(YR)^{-1}$, with $R_y$ the row $y$ of $R$,
$P(Y)=Y_yc(Y)^{\mathsf T}$ and $Q(Y)=e_y(c(Y)^{\mathsf T}Y)$. If $Y\notin U_0$,
then $Y_x=0$ and some $c'\neq0$ has $c'^{\mathsf T}Y_{\mathrm m}=0$, so
$c'^{\mathsf T}Y$ is a nonzero multiple of $e_y^{\mathsf T}$; this gives
$c(Y)^{\mathsf T}Y=e_y^{\mathsf T}$, so $Q(Y)=e_ye_y^{\mathsf T}$,
$P(Y)Y=YQ(Y)$ and $1+(l-1)Q(Y)\in P^\circ$. So $\ker L$ kills a neighbourhood of
$Y_0$, and by the first principle it contains every function of full rank.

\emph{Step 5: proof of (1).} By
Lemma~\ref{asm-theta-agree-coinvariants-inputs}(1), with $\chi_1=\nu\neq1$ and
$\chi_2=1$, every Schwartz-Bruhat $\Phi$ is $\Phi_0+z$ with $\Phi_0$ of full
rank and $z\in\mathfrak a(\nu,1)$. By (a), $L(z)=0$, and $L(\Phi_0)=0$ by
Step 4.

(2) Let $\psi$ be of full rank and $g\in GL_n(k)$, and put $g'=g^{\mathsf T}$.
The function $j(\psi)$ is Schwartz-Bruhat, and
\[
\omega(g,1)j(\psi)(Y)=\nu(\det g)|\det g|^{d/2}\psi(g'Y_{\mathrm m})\mathbf 1_{\mathcal O}(g'Y_x)\mathbf 1_{\mathcal O}(g'Y_y).
\]
By Step 1 the factor $\mathbf 1_{\mathcal O}(g'Y_x)$ can be replaced by
$\mathbf 1_{\mathcal O}(Y_x)$, and by Step 2 the factor
$\mathbf 1_{\mathcal O}(g'Y_y)$ by $|\det g|^{-1}\mathbf 1_{\mathcal O}(Y_y)$, as
$Z\mapsto\psi(g'Z)$ has full rank. As $d=d_1+2$, this gives
$L(\omega(g,1)j(\psi))=L(j(\omega_{\nu,1}(g,1)\psi))$, and the left side is
$L(j(\psi))$ by (a).
\end{proof}
% Lean: (1) red, through kil_loc (the first principle, ExistsInvar.lean), fixk (the second
% principle, with lam and idp), Inv.op (the identity before Step 1), kil_one (Step 1),
% swap_mid, trans_inv and mod_inv (with invariant_comp_mulVec, Step 2), h2 in red (Step 3),
% kil_full (Step 4) and KCo.surjectivity_statement (Step 5). (2) aug_ker, through full_omB,
% swap_mid and mod_inv.

\begin{lemma}\label{asm-theta-agree-local-mult-ramified-irreducible}
Let $k$ be a nonarchimedean local field, let $n\leq d$, let $\nu$ be a unitary
smooth character of $k^\times$
(Definition~\ref{asm-theta-agree-coinvariants-local}) with $\nu\neq1$, and let
$\alpha$ be a character of $k^\times$. Let
$\mathfrak a^\circ(\nu,\alpha)$ be the span of the
$\omega_{\nu,\alpha}(g,1)\Psi-\Psi$ with $g\in GL_n(k)$ and $\Psi$ of full rank
on $M_{n\times d}(k)$, and let $N$ be a subspace of the functions on
$M_{n\times d}(k)$ that contains $\mathfrak a^\circ(\nu,\alpha)$ and contains
$\omega_{\nu,\alpha}(1,h)\Psi$ whenever $\Psi\in N$ is Schwartz-Bruhat and
$h\in GL_d(k)$. Then $N$ contains every function of full rank, or every function
of full rank in $N$ lies in $\mathfrak a^\circ(\nu,\alpha)$.
\end{lemma}
\leanref{Development/AutomorphicSpectrum/Proofs/BlockDecomposition/LocalMultiplicityOneIrreducibleGL.lean}{GlobalTheory.CompanionTheory.BernNarrow.irrF'}
\begin{proof}
Write $d=n+t$ and put $\mu=\nu^{-1}$, which is unitary, smooth and not $1$. Let
$R$, $\mathcal S(R)$, $\omega^R$, $\tau$, $\operatorname{Ind}(\tau)$, $R(h_1)$ and
$E$ be as in Lemma~\ref{asm-theta-agree-coinvariants-fibre} for $\mu$. For
$\phi\in\mathcal S(R)$ let $Z\phi(Y)=\phi(Y^{\mathsf T})$ if
$Y\in M_{n\times d}(k)$ has rank $n$, and $Z\phi(Y)=0$ otherwise. As the
matrices of rank $n$ form an open set, $Z\phi$ is Schwartz-Bruhat, so it has full
rank; and every function $\Psi$ of full rank is $Z\phi$, with
$\phi(x)=\Psi(x^{\mathsf T})$. As $g^{\mathsf T}Y(h^{\mathsf T})^{-1}$ has the
rank of $Y$ and its transpose is $h^{-1}Y^{\mathsf T}g$,
\[
Z(\omega^R(g,h)\phi)(Y)=\nu(\det g)|\det g|^{d/2}\,\nu(\det h)^{-1}|\det h|^{-n/2}\,Z\phi\bigl(g^{\mathsf T}Y(h^{\mathsf T})^{-1}\bigr).
\]
So $Z(\omega^R(g,1)\phi)=\omega_{\nu,\alpha}(g,1)Z\phi$, and
$Z(\omega^R(1,h)\phi)=\alpha(\det h)\nu(\det h)^{-1}\,\omega_{\nu,\alpha}(1,(h^{\mathsf T})^{-1})Z\phi$.

Let $N'$ be the space of the $\phi\in\mathcal S(R)$ with $Z\phi\in N$, and
$\mathrm{Rel}_R$ the span of the $\omega^R(g,1)\phi-\phi$ with $g\in GL_n(k)$.
As $Z\phi$ has full rank,
$Z(\omega^R(g,1)\phi-\phi)=\omega_{\nu,\alpha}(g,1)Z\phi-Z\phi$ lies in
$\mathfrak a^\circ(\nu,\alpha)$; so
$Z(\mathrm{Rel}_R)\subset\mathfrak a^\circ(\nu,\alpha)\subset N$, and
$\mathrm{Rel}_R\subset N'$. If $\phi\in N'$ and $h\in GL_d(k)$, then $Z\phi\in N$
is Schwartz-Bruhat, so $N$ contains
$\omega_{\nu,\alpha}(1,(h^{\mathsf T})^{-1})Z\phi$ and its multiple
$Z(\omega^R(1,h)\phi)$. So $N'$ is stable under every $\omega^R(1,h)$.

By Lemma~\ref{asm-theta-agree-coinvariants-fibre},
$E(\omega^R(g,1)\phi)(x)=\tau(g,1,1)\bigl(E(\phi)(x)\bigr)$ and
$E(\omega^R(1,h)\phi)=R(h)E(\phi)$. As $E$ is bijective, it maps
$\mathrm{Rel}_R$ onto the space $\mathrm{Rel}$ of
Lemma~\ref{asm-theta-agree-coinvariants-irreducible}, and $E(N')$ contains
$\mathrm{Rel}$ and is stable under every $R(h)$. By
Lemma~\ref{asm-theta-agree-coinvariants-irreducible}, $E(N')=\operatorname{Ind}(\tau)$
or $E(N')\subset\mathrm{Rel}$; as $E$ is bijective, $N'=\mathcal S(R)$ or
$N'\subset\mathrm{Rel}_R$. In the first case every function of full rank, being
some $Z\phi$, lies in $N$. In the second, a function of full rank in $N$ is
$Z\phi$ with $\phi\in N'\subset\mathrm{Rel}_R$, so it lies in
$Z(\mathrm{Rel}_R)\subset\mathfrak a^\circ(\nu,\alpha)$.
\end{proof}
% Lean: BernNarrow.irrF' (BlockDecomposition/LocalMultiplicityOneIrreducibleGL.lean). Z is
% KCo.extZ (ExistsExtZ.lean), with extZ_mem, extZ_neg (full rank), exists_extZ and extZ_act
% (the displayed formula); E is T.omega_equiv_induced_fibre for mu = nu^-1; hR is the inclusion
% Rel_R in N' and the convert step with invT h the stability of N'; irr_equiv carries the
% dichotomy through E. BernNarrow.irr_ind and irr_gen, which irrF' applies, are verbatim copies
% of KCoNarrow.irr_ind_narrow and irr_gen_narrow (Lemma
% asm-theta-agree-coinvariants-irreducible). The proof of Lemma
% asm-theta-agree-coinvariants-inputs(2), KCoNarrow.irrSt_narrow, has the same steps for the
% augmentation of all Schwartz-Bruhat functions.

\begin{lemma}\label{asm-theta-agree-local-mult-ramified-gl}
In Situation~\ref{asm-theta-agree-local-mult-ramified-situation}, let $P$ be a
complex vector space of countable dimension and $\rho$ a map from
$GL_{d_1}(k)$ to the linear operators on $P$ such that $0$ and $P$ are the only
subspaces stable under every $\rho(g)$. Let $L$ and $L'$ be linear maps from the
functions on $M_{n\times d}(k)$ to $P$ that satisfy (a) and (b) and have
$L(\Phi(\,\cdot\,\operatorname{blk}(g)))=\rho(g)L(\Phi)$ and
$L'(\Phi(\,\cdot\,\operatorname{blk}(g)))=\rho(g)L'(\Phi)$ for Schwartz-Bruhat
$\Phi$ and $g\in GL_{d_1}(k)$. If $L(\Phi)\neq0$ for some Schwartz-Bruhat $\Phi$,
then there is $c\in\mathbf C$ with $L'(\Phi)=cL(\Phi)$ for every
Schwartz-Bruhat $\Phi$.
\end{lemma}
\leanref{Development/AutomorphicSpectrum/Proofs/BlockDecomposition/LocalMultiplicityOneIrreducibleGL.lean}{GlobalTheory.CompanionTheory.BernNarrow.local_multiplicity_one_irreducible_GL}
\begin{proof}
Let $S$ be the space of functions of full rank on $M_{n\times d_1}(k)$, and put
$\mu=L\circ j$ and $\mu'=L'\circ j$ on $S$. For $g\in GL_{d_1}(k)$ and $\psi\in S$
the function $\psi(\,\cdot\,g)$ lies in $S$, and
$j(\psi(\,\cdot\,g))=j(\psi)(\,\cdot\,\operatorname{blk}(g))$; so
$\mu(\psi(\,\cdot\,g))=\rho(g)\mu(\psi)$, and likewise for $\mu'$. By
Lemma~\ref{asm-theta-agree-local-mult-ramified-reduction}(1), $\mu\neq0$.

Let $N$ be the space of $\psi\in S$ with $\mu(\psi)=0$. By
Lemma~\ref{asm-theta-agree-local-mult-ramified-reduction}(2) it contains
$\mathfrak a^\circ$, whose elements have full rank. As
$\omega_{\nu,1}(1,g)\psi=|\det g|^{n/2}\psi(\,\cdot\,g)$, it contains
$\omega_{\nu,1}(1,g)\psi$ for $\psi\in N$. By
Lemma~\ref{asm-theta-agree-local-mult-ramified-irreducible}, for $\nu$, $\alpha=1$
and $n\leq d_1$, either $N=S$, which $\mu\neq0$ excludes, or $N\subset\mathfrak a^\circ$.
By Lemma~\ref{asm-theta-agree-local-mult-ramified-reduction}(2) for $L'$, $\mu'$
vanishes on $\mathfrak a^\circ$, so $\mu'$ vanishes on the kernel of $\mu$.

The image of $\mu$ is a nonzero subspace of $P$ stable under every $\rho(g)$, so
it is $P$. Let $s$ be a linear map from $P$ to $S$ with $\mu\circ s=1$, and put
$T=\mu'\circ s$. As $\mu'$ vanishes on the kernel of $\mu$, $T\circ\mu=\mu'$, and
then $T\rho(g)\mu(\psi)=\mu'(\psi(\,\cdot\,g))=\rho(g)T\mu(\psi)$. So $T$
commutes with every $\rho(g)$, and $T=c$ for some $c\in\mathbf C$ by
Lemma~\ref{th-first-lift-dixmier}. Hence $L'-cL$, which satisfies (a) and (b),
vanishes on $j(\psi)$ for every $\psi\in S$, and
Lemma~\ref{asm-theta-agree-local-mult-ramified-reduction}(1) gives $L'=cL$ on the
Schwartz-Bruhat functions.
\end{proof}
% Lean: local_multiplicity_one_irreducible_GL (BernNarrow), with red, aug_ker, full_omB,
% tn_mem, irrF' (for alpha = 1) and dixmier_commutant_scalar.

\begin{lemma}\label{asm-theta-agree-local-mult-ramified-pair}
Let $(X,W,M,\chi_X,\chi_W,\omega)$ be a Weil pair for a theta base $B$ of a CM
field $E$ (Definition~\ref{asm-theta-agree-triple}(1)), with $X$ and $W$ of
ranks $n$ and $d$, put $I=\{0,\ldots,n-1\}\times\{0,\ldots,d-1\}$, and let $H'$ be
an adapted plane in $W$ over a skew-Hermitian lattice space $W_1$ of rank $d_1$
at some positions, with $n\leq d_1$. Let $v$ be a finite place of $F$ and
$t\colon E\to F_v$ an $F$-embedding such that $(\chi_W)_t$ is ramified and
$\pi_t$ maps $U_{W_1}^{(v)}$ onto $GL_{d_1}(F_v)$. Let $P$ be a complex vector
space of countable dimension and $\tau$ a map from $U_{W_1}$ to the linear
operators on $P$ with $\tau(1)=1$, such that $0$ and $P$ are the only subspaces
stable under every $\tau(h)$ with $h\in U_{W_1}^{(v)}$, and let $r\in\mathbf R$.
Consider linear maps $\Lambda$ from the functions on $F_v^I$ to $P$ such that:
$\Lambda$ vanishes on the span $\mathfrak a_v(\omega)$ of the $\phi_1-\phi$ with
$\phi,\phi_1$ Schwartz-Bruhat and $\omega(g,1)$ acting at $v$ by
$\phi\mapsto\phi_1$ for some $g\in U_X^{(v)}$; and
$\Lambda(\phi')=\chi_X(z)|z|^r\tau(h_1)\Lambda(\phi)$ whenever $m\in U_W^{(v)}$
is a Levi element for $(z,h_1)$ with $h_1\in U_{W_1}^{(v)}$
(Definition~\ref{asm-theta-agree-planes}(2)), $u\in N_{H'}$ is at $v$, $\phi$ and
$\phi'$ are Schwartz-Bruhat and $\omega(1,um)$ acts at $v$ by $\phi\mapsto\phi'$.
If $\Lambda$ and $\Lambda'$ are two such maps and $\Lambda(\phi)\neq0$ for some
Schwartz-Bruhat $\phi$, then there is $c\in\mathbf C$ with
$\Lambda'(\phi)=c\Lambda(\phi)$ for every Schwartz-Bruhat $\phi$.
\end{lemma}
\leanref{Development/AutomorphicSpectrum/Proofs/BlockDecomposition/LocalMultiplicityOneIrreducibleGL.lean}{GlobalTheory.CompanionTheory.BernNarrow.pair_multiplicity_one_ramified}
\leanref{Development/WeilBound/ThetaLifts/LocalTheory/DixmierCommutants/RamifiedMultiplicityOneOfBernstein.lean}{PositivityTheory.ThKG.LocalRamifiedLeviRestriction.Eqv}
\begin{proof}
Let $\nu=(\chi_W)_t$ and $\alpha'=(\chi_X)_t$, unitary smooth characters of
$F_v^\times$, and $\bar t$ the other $F$-embedding of $E$ into $F_v$. Write
$H'=(x',y',\iota_{H'},\mathfrak b')$, index $\{0,\ldots,d-1\}$ by
$\{0,\ldots,d_1-1\}\sqcup\{x,y\}$ through a bijection, and let $f\in M_d(E)$ have
the columns of $\iota_{H'}$ in the places $\{0,\ldots,d_1-1\}$ and the columns
$x'$ and $y'$ in the places $x$ and $y$. By the conditions on an adapted plane,
$f^{\mathsf T}J_Wc(f)$ is block diagonal with blocks $J_{W_1}$ and
$\left(\begin{smallmatrix}0&1\\-1&0\end{smallmatrix}\right)$; so $f$ and
$C=t(f)\in GL_d(F_v)$ are invertible.

By Lemma~\ref{asm-theta-agree-local-ramified-values-dictionary}(1) there is a
linear model $\iota_0$ of the second kind for $(t,1,1)$. Put
$\iota\phi(Y)=\iota_0\phi(YC^{-1})$; it is bijective between the
Schwartz-Bruhat functions. By the same lemma, every $g_0\in GL_n(F_v)$ is
$\pi_t(g)$ for some $g\in U_X^{(v)}$, and $\omega(g,1)$ reads as
$\omega_B^{\nu,\alpha'}(g_0,1)$ through $\iota$; and if $h\in U_W^{(v)}$ and
$p\in GL_d(F_v)$ have $\pi_t(h)=CpC^{-1}$, then $\omega(1,h)$ reads as
$\omega_B^{\nu,\alpha'}(1,p)$ through $\iota$, since
$\omega_B^{\nu,\alpha'}(1,CpC^{-1})$ and right multiplication by $C^{-1}$ give
$\alpha'(\det p)|\det p|^{n/2}\iota_0\phi(YpC^{-1})$. Let $L$ and $L'$ be linear
maps from the functions on $M_{n\times d}(F_v)$ to $P$ with
$L(\iota\phi)=\Lambda(\phi)$ and $L'(\iota\phi)=\Lambda'(\phi)$ for
Schwartz-Bruhat $\phi$. For $g\in GL_{d_1}(F_v)$ choose
$h_1(g)\in U_{W_1}^{(v)}$ with $\pi_t(h_1(g))=g$, and put
$\rho(g)=(\alpha'(\det g)|\det g|^{n/2})^{-1}\tau(h_1(g))$. We check the
hypotheses of Lemma~\ref{asm-theta-agree-local-mult-ramified-gl} for $k=F_v$,
$\nu$ and $\rho$; the argument is the same for $L$ and $L'$.

(a) Let $g_0\in GL_n(F_v)$ and $\Phi=\iota\phi$. With $g$ as above, $\omega(g,1)$
acts at $v$ by $\phi\mapsto\phi_1$ with $\iota\phi_1=\omega_B^{\nu,\alpha'}(g_0,1)\Phi=\omega_{\nu,1}(g_0,1)\Phi$,
and $\Lambda(\phi_1-\phi)=0$.

(b) Let $h\in P^\circ$, put $x_0=h_{xx}$ and $y_0=h_{yy}$, and let
$p_0\in GL_d(F_v)$ be diagonal with entries $x_0$ at $x$, $y_0$ at $y$ and $1$
elsewhere. Choose $u,m\in U_W^{(v)}$ with $\pi_t(u)=Chp_0^{-1}C^{-1}$ and
$\pi_t(m)=Cp_0C^{-1}$. The matrix $hp_0^{-1}$ has column $x$ equal to $e_x$,
entry $1$ at $(y,y)$, identity middle block and row $y$ zero in the middle
columns; so $u$ fixes $x'$ and moves each $\iota_{H'}z$ within
$\iota_{H'}z+\mathbf A_Ex'$, the components at $\bar t$ being determined by
those at $t$ through the unitarity of $u$, and $u\in N_{H'}$. Let
$z\in\mathbf A_E^\times$ have $z-1$ at $v$, $\pi_t(z)=x_0$ and
$\pi_{\bar t}(z)=y_0^{-1}$; then $|z|=1$, and $\chi_X(z)=\alpha'(x_0y_0)$, since
$\chi_X$ restricts to a power of $\eta$ on $\mathbf A^\times$ and $\eta$ is
trivial at $v$. In the same way $m$ is a Levi element for $(z,1)$. As
$\pi_t(um)=ChC^{-1}$, $\omega(1,um)$ reads as $\omega_B^{\nu,\alpha'}(1,h)$, which
maps $\Phi$ to $\alpha'(x_0y_0)\Phi(\,\cdot\,h)$ since $\det h=x_0y_0$ is a unit.
The equivariance of $\Lambda$ gives
$\alpha'(x_0y_0)L(\Phi(\,\cdot\,h))=\chi_X(z)|z|^rL(\Phi)$, so
$L(\Phi(\,\cdot\,h))=L(\Phi)$.

(c) Let $g\in GL_{d_1}(F_v)$, and choose $m\in U_W^{(v)}$ with
$\pi_t(m)=C\operatorname{blk}(g)C^{-1}$; as in (b), $m$ is a Levi element for
$(1,h_1(g))$. Through $\iota$, $\omega(1,m)$ reads as
$\omega_B^{\nu,\alpha'}(1,\operatorname{blk}(g))$, which maps $\Phi$ to
$\alpha'(\det g)|\det g|^{n/2}\Phi(\,\cdot\,\operatorname{blk}(g))$, and the
equivariance with $u=1$ gives $L(\Phi(\,\cdot\,\operatorname{blk}(g)))=\rho(g)L(\Phi)$.

An element of $U_{W_1}^{(v)}$ is determined by its component at $t$, so
$h=h_1(\pi_t(h))$ for $h\in U_{W_1}^{(v)}$, and a subspace stable under every
$\rho(g)$ is stable under every such $\tau(h)$; so it is $0$ or $P$. As $\nu$ is
ramified, $\nu(l)\neq1$ for some $l\in\mathcal O_v^\times$. So
Lemma~\ref{asm-theta-agree-local-mult-ramified-gl} gives $c$ with $L'=cL$ on the
Schwartz-Bruhat functions, and $\Lambda'=c\Lambda$ on them as $\iota$ is
bijective.
\end{proof}
% Lean: pair_multiplicity_one_ramified (BernNarrow), with dictionary_extended (dictionary
% (1)), fr, fr_gram, gm_det and Fr for C, chg 1 C^-1 1 for iota with reads_chg, B_left and
% B_right, inP_det, mem_unip (through col_eq), idele_two, is_levi, ua_ext and the hypothesis
% hS. The Lean states it for a good base, X Hermitian, W and W1 skew-Hermitian, a Weil
% representation, a Hecke character chi_W and chi_X with ResEtaPow, which a Weil pair provides.

\begin{lemma}\label{asm-theta-agree-local-mult-ramified}
Let $B$ be a theta base of a CM field $E$, let $Q$ and $Q_1$ be matched
quadruples of step type $(\chi,r)$, and let $H'$ be an adapted plane in $W$
over $W_1$ at some positions
(Definition~\ref{asm-theta-agree-local-equivariant}). Let $v\in T$, and let
$\lambda,\lambda'\colon\Pi_v\to\Pi_{1,v}$ be equivariant at $v$ for
$(H',\chi,r)$ with $\lambda\neq0$. Then $\lambda'=c\lambda$ for some
$c\in\mathbf C$.
\end{lemma}
\leanref{Development/WeilBound/ThetaLifts/LocalTheory/UnramifiedDeterminedBySpherical.lean}{PositivityTheory.ThKG.Loc.RamifiedMultiplicityOne}
\leanref{Development/WeilBound/CoreStatements/TemperedRepresentationTheory/BernsteinNarrowness/BernsteinClosure/RamifiedMultiplicityOne.lean}{PositivityTheory.BernsteinRamification.Close.ramifiedMultiplicityOne}
\begin{proof}
Let $(a,a')$ and $(a_1,a'_1)$ be the ranks of $Q$ and $Q_1$, $d_1=a_1+a'_1$ the
rank of $W_1$, and $t=\tau_v$.

\emph{Step 1: the action on $\Pi_{1,v}$.} By
Lemma~\ref{asm-theta-agree-local-ramified-values-dictionary}(1), for the Weil pair
$(V_1,W_1,M_1,\beta,\beta/\alpha,\omega_1)$ of $Q_1$, $\pi_t$ maps
$U_{W_1}^{(v)}$ onto $GL_{d_1}(F_v)$, and for $h\in U_{W_1}^{(v)}$ and
Schwartz-Bruhat $\phi$ there is a Schwartz-Bruhat $\phi_1$ with
$\omega_1(1,h)$ acting at $v$ by $\phi\mapsto\phi_1$. If $\omega_1(1,h)$ acts at
$v$ by $\phi\mapsto\phi_1$ and by $\psi\mapsto\psi_1$, and
$p_{1,v}(\phi)=p_{1,v}(\psi)$, then $\omega_1(1,h)$ acts by
$\phi-\psi\mapsto\phi_1-\psi_1$ and $\omega'_1(1,h)$ by $0\mapsto0$, and
$p_{1,v}(\phi-\psi)=0=p'_{1,v}(0)$; so the matching of the finite maps gives
$p_{1,v}(\phi_1-\psi_1)=0$. As $p_{1,v}$ maps the Schwartz-Bruhat functions onto
$\Pi_{1,v}$, there are linear operators $\tau(h)$ on $\Pi_{1,v}$, for
$h\in U_{W_1}^{(v)}$, with $\tau(h)p_{1,v}(\phi)=p_{1,v}(\phi_1)$; put
$\tau(h)=0$ for the other $h\in U_{W_1}$. Then $\tau(1)=1$, and $0$ and
$\Pi_{1,v}$ are the only subspaces stable under every $\tau(h)$,
$h\in U_{W_1}^{(v)}$, by the matching. The space $\Pi_{1,v}$ has countable
dimension, as the image of the Schwartz-Bruhat functions on a power of $F_v$,
which are spanned by the countably many indicator functions of cosets of the
subgroups $\varpi_v^m\mathcal O_v^I$.

\emph{Step 2: the equivariance.} Let $\lambda$ be equivariant at $v$ for
$(H',\chi,r)$. Then $\Lambda=\lambda\circ p_v$ satisfies the hypotheses of
Lemma~\ref{asm-theta-agree-local-mult-ramified-pair} for the Weil pair
$(V,W,M,\beta,\beta/\alpha,\omega)$, with $\chi$ in place of $\chi_X$, $\tau$ and
$r$. It vanishes on $\mathfrak a_v(\omega)$ as $p_v$ does. Let $m$, $u$, $z$,
$h_1$, $\phi$ and $\phi'$ be as there. Choose a Schwartz-Bruhat $\phi_1$ with
$p_{1,v}(\phi_1)=\lambda(p_v(\phi))$ and $\phi'_1$ with $\omega_1(1,h_1)$ acting
at $v$ by $\phi_1\mapsto\phi'_1$. By
Definition~\ref{asm-theta-agree-local-equivariant}(2),
\[
\lambda(p_v(\phi'))=\chi(z)|z|^rp_{1,v}(\phi'_1)=\chi(z)|z|^r\tau(h_1)\lambda(p_v(\phi)).
\]
Likewise $\Lambda=\lambda\circ p'_v$ satisfies them for
$(V',W,M',\alpha,\alpha/\beta,\omega')$, with $\chi$ in place of $\chi_X$: if
$\omega'(1,um)$ acts at $v$ by $\phi\mapsto\phi'$, choose $\psi$ with
$p_v(\psi)=p'_v(\phi)$ and $\psi'$ with $\omega(1,um)$ acting by $\psi\mapsto\psi'$;
the matching gives $p_v(\psi')=p'_v(\phi')$, and the first case applies to $\psi$.

\emph{Step 3.} If the step type is $(\beta,a/2)$, then $\chi=\beta$, $a_1=a$ and
$a'=a'_1+2$, so $a\leq d_1$; and $(\beta/\alpha)_t=\nu_t^{-1}$ is ramified by
Lemma~\ref{asm-theta-agree-finite-characters}. Apply
Lemma~\ref{asm-theta-agree-local-mult-ramified-pair} to
$(V,W,M,\beta,\beta/\alpha,\omega)$, $H'$, $t$, $P=\Pi_{1,v}$, $\tau$ and the
maps $\lambda\circ p_v$ and $\lambda'\circ p_v$ of Step 2; the first is not $0$
on the Schwartz-Bruhat functions since $\lambda\neq0$ and $p_v$ maps them onto
$\Pi_v$. It gives $c$ with $\lambda'\circ p_v=c\,\lambda\circ p_v$ on them, so
$\lambda'=c\lambda$. If the step type is $(\alpha,a'/2)$, then $\chi=\alpha$,
$a=a_1+2$ and $a'_1=a'$, so $a'\leq d_1$, and $(\alpha/\beta)_t=\nu_t$ is
ramified; the same argument with $(V',W,M',\alpha,\alpha/\beta,\omega')$ and
$p'_v$ gives $\lambda'=c\lambda$.
\end{proof}
% Lean: BernsteinRamification.Close.ramifiedMultiplicityOne =
% ramifiedMultiplicityOne_of_geometricFiltration Geo.geometricLemmaFiltration =
% BernNarrow.first_cohomology_ramified_narrow (narrowBern hG) =
% ramifiedMultiplicityOne_of_bernstein (LocalMultiplicityOneIrreducibleGL.lean). Step 1 is
% acts_ex, wd, exTau, tau1, hr (rank_SB) and hi; Step 2 is eq and eq'; Step 3 is the case
% split on StepTy with pair_multiplicity_one_ramified. What remains open is Lemma
% asm-theta-agree-local-mult-ramified-irreducible, and the dictionary (1) of Lemma
% asm-theta-agree-local-ramified-values-dictionary.

\begin{definition}\label{asm-theta-agree-local-mult-archimedean-modules}
Let $\Sigma$ be a finite set and $R,S\geq0$, and use the notation of
Definitions~\ref{asm-theta-agree-real-fock}
and~\ref{asm-theta-agree-real-quotient-lie-module}. For $w\in\Sigma$ and
$Y\in M_{R+S}(\mathbf C)$ let $\delta_wY\in M_{R+S}(\mathbf C)^\Sigma$ be the
family with $Y$ at $w$ and $0$ elsewhere, and put $E^w_{ab}=\delta_wE_{ab}$.
Families of matrices are multiplied placewise.
\begin{enumerate}
\item A \emph{Lie module of type $(R,S)$ over $\Sigma$} is a complex vector
  space $M$ with a $\mathbf C$-linear map $Z\mapsto M(Z)$ from
  $M_{R+S}(\mathbf C)^\Sigma$ to $\operatorname{End}M$ such that
  $M(ZZ'-Z'Z)=M(Z)M(Z')-M(Z')M(Z)$ for all $Z$ and $Z'$. A subspace of $M$ is
  \emph{stable} if it is stable under every $M(Z)$.
\item Let $M$ be a Lie module of type $(R,S)$ over $\Sigma$, let
  $\Lambda\colon\Sigma\times\mathrm{Ix}(R,S)\to\mathbf C$, written
  $(w,a)\mapsto\Lambda_w(a)$, and let $f\in M$. Then $(M,f)$ is a
  \emph{highest-weight module of weight $\Lambda$} if:
  \begin{enumerate}
  \item[(a)] $M(E^w_{aa})f=\Lambda_w(a)f$ for all $w$ and $a$;
  \item[(b)] $M(E^w_{ab})f=0$ for all $w$ and all $a\prec b$;
  \item[(c)] $M$ is the only stable subspace that contains $f$;
  \item[(d)] every $u\in M$ that satisfies (a) and (b) in place of $f$ is a
    multiple of $f$.
  \end{enumerate}
\end{enumerate}
\end{definition}
\leanref{Development/WeilBound/ThetaLifts/LocalTheory/ToralCohomology/ShiftLow.lean}{PositivityTheory.ThKG.CayleyTransformDarboux.Lie}
\leanref{Development/WeilBound/ThetaLifts/LocalTheory/ToralCohomology/ShiftLow.lean}{PositivityTheory.ThKG.CayleyTransformDarboux.St}
\leanref{Development/WeilBound/ThetaLifts/LocalTheory/ToralCohomology/ShiftLow.lean}{PositivityTheory.ThKG.CayleyTransformDarboux.un}
\leanref{Development/WeilBound/ThetaLifts/LocalTheory/ToralCohomology/SumPlSingle.lean}{PositivityTheory.ThKG.CayleyTransformDarboux.HWMod}

\begin{lemma}\label{asm-theta-agree-local-mult-archimedean-gen}
Let $A$ be an associative $\mathbf C$-algebra, $M$ a complex vector space and
$Z\mapsto M(Z)$ a $\mathbf C$-linear map from $A$ to $\operatorname{End}M$ with
$M(ZZ'-Z'Z)=M(Z)M(Z')-M(Z')M(Z)$ for all $Z,Z'\in A$. Let $\mathfrak g$,
$\mathfrak q$ and $\mathfrak l$ be subspaces of $A$ such that
$ZZ'-Z'Z\in\mathfrak g$ for all $Z,Z'\in\mathfrak g$, and
$\mathfrak q,\mathfrak l\subset\mathfrak g\subset\mathfrak q+\mathfrak l$. Let
$V\subset N$ be subspaces of $M$ such that $M(Z)V\subset V$ for
$Z\in\mathfrak q$ and $M(Z)N\subset N$ for $Z\in\mathfrak l$. Then there is a
subspace $W$ with $V\subset W\subset N$ and $M(Z)W\subset W$ for every
$Z\in\mathfrak g$.
\end{lemma}
\leanref{Development/WeilBound/ThetaLifts/LocalTheory/ToralCohomology/ShiftLow.lean}{PositivityTheory.ThKG.CayleyTransformDarboux.gen_le}
\begin{proof}
Put $F_0=V$ and $F_{n+1}=F_n+\sum_{Y\in\mathfrak l}M(Y)F_n$, and let $W$ be the
union of the $F_n$. Then $V\subset W$, and $F_n\subset N$ by induction on $n$,
as $N$ is stable under the $M(Y)$ with $Y\in\mathfrak l$. By induction on $n$,
$F_n$ is stable under the $M(Z)$ with $Z\in\mathfrak q$. For $n=0$ this is the
hypothesis on $V$. Suppose it holds for $F_n$. Then $M(Z)F_n\subset F_{n+1}$
for every $Z\in\mathfrak g$, by writing $Z=Z_1+Z_2$ with $Z_1\in\mathfrak q$
and $Z_2\in\mathfrak l$. Let $Z\in\mathfrak q$, $Y\in\mathfrak l$ and
$x\in F_n$. Then $M(Z)M(Y)x=M(Y)M(Z)x+M(ZY-YZ)x$. The first term lies in
$F_{n+1}$ as $M(Z)x\in F_n$, and the second because $ZY-YZ\in\mathfrak g$. So
$F_{n+1}$ is stable under the $M(Z)$ with $Z\in\mathfrak q$. Hence
$M(Z)F_n\subset F_{n+1}$ for all $n$ and all $Z\in\mathfrak g$, and $W$ is
stable under the $M(Z)$ with $Z\in\mathfrak g$.
\end{proof}
% Lean: gen_le (ShiftLow.lean, with filt; about 35 lines). up and mono are the inclusions
% M(Y)F_n in F_{n+1} and F_n in F_{n+1}, Q is M(Z)F_n in F_{n+1} for Z in g, P the stability
% under q, le the inclusion in N; W is the supremum of the filt act l V n.

\begin{lemma}\label{asm-theta-agree-local-mult-archimedean-low}
Let $M$ be a complex vector space, $T$ an operator on $M$ and
$\lambda\in\mathbf C$. Let $M_\lambda(T)$ be the space of $x\in M$ with
$Tx=\lambda x$, and $M_{<\lambda}(T)$ the sum of the spaces $M_{\lambda-t}(T)$
over the real numbers $t>0$.
\begin{enumerate}
\item $M_{<\lambda}(T)\cap M_\lambda(T)=0$.
\item If $D\colon M\to M'$ is linear with $DT=\lambda D$, then $D$ vanishes on
  $M_{<\lambda}(T)$.
\item If $\varphi\colon M\to M'$ is linear and $T'$ is an operator on $M'$ with
  $\varphi T=T'\varphi$, then $\varphi$ maps $M_{<\lambda}(T)$ into
  $M'_{<\lambda}(T')$.
\item Let $Z$ be an operator on $M$ and $k\in\mathbf C$ with $TZ-ZT=kZ$. Then
  $Z$ maps $M_\mu(T)$ into $M_{\mu+k}(T)$ for every $\mu\in\mathbf C$. If $k$ is
  real and negative and $V\subset M_\lambda(T)$ is a subspace, then
  $V+M_{<\lambda}(T)$ is stable under $Z$.
\end{enumerate}
\end{lemma}
\leanref{Development/WeilBound/ThetaLifts/LocalTheory/ToralCohomology/ShiftLow.lean}{PositivityTheory.ThKG.CayleyTransformDarboux.low}
\leanref{Development/WeilBound/ThetaLifts/LocalTheory/ToralCohomology/ShiftLow.lean}{PositivityTheory.ThKG.CayleyTransformDarboux.low_eig}
\leanref{Development/WeilBound/ThetaLifts/LocalTheory/ToralCohomology/ShiftLow.lean}{PositivityTheory.ThKG.CayleyTransformDarboux.low_ker}
\leanref{Development/WeilBound/ThetaLifts/LocalTheory/ToralCohomology/ShiftLow.lean}{PositivityTheory.ThKG.CayleyTransformDarboux.low_map}
\leanref{Development/WeilBound/ThetaLifts/LocalTheory/ToralCohomology/ShiftLow.lean}{PositivityTheory.ThKG.CayleyTransformDarboux.shift_eig}
\leanref{Development/WeilBound/ThetaLifts/LocalTheory/ToralCohomology/ShiftLow.lean}{PositivityTheory.ThKG.CayleyTransformDarboux.shift_low}
\begin{proof}
(1) Eigenvectors of $T$ for distinct eigenvalues are linearly independent, and
$\lambda-t\neq\lambda$ for $t>0$. (2) If $Tx=(\lambda-t)x$ with $t>0$, then
$\lambda Dx=DTx=(\lambda-t)Dx$, so $tDx=0$ and $Dx=0$. (3) If
$Tx=(\lambda-t)x$, then $T'\varphi(x)=(\lambda-t)\varphi(x)$. (4) If $Tx=\mu x$,
then $TZx=ZTx+kZx=(\mu+k)Zx$. So $Z$ maps $M_\lambda(T)$ into
$M_{\lambda+k}(T)$ and $M_{\lambda-t}(T)$ into $M_{\lambda-t+k}(T)$, and both
lie in $M_{<\lambda}(T)$ when $k<0$.
\end{proof}
% Lean: low, low_eig (from Module.End.eigenspaces_iSupIndep), low_ker, low_map, shift_eig
% and shift_low (ShiftLow.lean, about 45 lines). In shift_eig and shift_low the operators are
% act H and act Z for a map act with the bracket relation (Lie) and H * Z - Z * H = k Z.

\begin{lemma}\label{asm-theta-agree-local-mult-archimedean-cyclic}
Let $\Sigma$ and $I$ be finite sets, $M$ a complex vector space and
$Z\mapsto M(Z)$ a $\mathbf C$-linear map from $M_I(\mathbf C)^\Sigma$ to
$\operatorname{End}M$ with $M(ZZ'-Z'Z)=M(Z)M(Z')-M(Z')M(Z)$, families being
multiplied placewise. For $w\in\Sigma$ and $a,b\in I$ let $E^w_{ab}$ be the
family with the matrix unit $E_{ab}$ at $w$ and $0$ elsewhere, and for
$\mathcal R\subset I\times I$ let $\mathfrak a(\mathcal R)$ be the span of the
$E^w_{ab}$ with $(a,b)\in\mathcal R$. Let $\mathcal K,\mathcal N,\mathcal Z\subset I\times I$
and $h\colon I\to\mathbf R$ satisfy:
\begin{enumerate}
\item[(a)] $\mathcal K$ is reflexive and transitive, and $\mathcal N\subset\mathcal K$;
\item[(b)] $h(b)\leq h(a)$ for $(a,b)\in\mathcal N$, and $(a,b)\in\mathcal K$
  and $h(a)=h(b)$ for $(a,b)\in\mathcal Z$;
\item[(c)] if $(a,b)\in\mathcal K$ and $h(b)\leq h(a)$, then $a=b$,
  $(a,b)\in\mathcal N$ or $(a,b)\in\mathcal Z$.
\end{enumerate}
Let $H$ be the family with $\operatorname{diag}(h)$ at every place, let
$\Lambda\colon\Sigma\times I\to\mathbf C$, written $(w,a)\mapsto\Lambda_w(a)$,
and put $\lambda=\sum_w\sum_ah(a)\Lambda_w(a)$. Let $f\in M$ with
$M(E^w_{aa})f=\Lambda_w(a)f$ for all $w$ and $a$ and $M(E^w_{ab})f=0$ for all
$w$ and $(a,b)\in\mathcal N$, and let $N\ni f$ be a subspace stable under the
$M(Z)$ with $Z\in\mathfrak a(\mathcal Z)$. Then there is a subspace $W\ni f$,
stable under the $M(Z)$ with $Z\in\mathfrak a(\mathcal K)$, with
$W\subset N+M_{<\lambda}(M(H))$
(Lemma~\ref{asm-theta-agree-local-mult-archimedean-low}).
\end{lemma}
\leanref{Development/WeilBound/ThetaLifts/LocalTheory/ToralCohomology/ShiftLow.lean}{PositivityTheory.ThKG.CayleyTransformDarboux.cyc_le}
\leanref{Development/WeilBound/ThetaLifts/LocalTheory/ToralCohomology/ShiftLow.lean}{PositivityTheory.ThKG.CayleyTransformDarboux.alg}
\leanref{Development/WeilBound/ThetaLifts/LocalTheory/ToralCohomology/ShiftLow.lean}{PositivityTheory.ThKG.CayleyTransformDarboux.dg}
\begin{proof}
Let $\mathcal K_\leq$ be the set of $(a,b)\in\mathcal K$ with $h(b)\leq h(a)$
and $\mathcal K_<$ the set of $(a,b)\in\mathcal K$ with $h(a)<h(b)$. The product
$E^w_{ab}E^{w'}_{cd}$ is $E^w_{ad}$ if $w=w'$ and $b=c$, and $0$ otherwise; as
$\mathcal K$ is transitive, $\mathfrak a(\mathcal K)$ and
$\mathfrak a(\mathcal K_\leq)$ are closed under commutators. Also
$HE^w_{ab}-E^w_{ab}H=(h(a)-h(b))E^w_{ab}$, and $M(H)f=\lambda f$, as
$H=\sum_w\sum_ah(a)E^w_{aa}$.

\emph{Step 1.} Apply Lemma~\ref{asm-theta-agree-local-mult-archimedean-gen}
with $\mathfrak g=\mathfrak a(\mathcal K_\leq)$, $\mathfrak q$ the span of the
$E^w_{aa}$ and of the $E^w_{ab}$ with $(a,b)\in\mathcal N$,
$\mathfrak l=\mathfrak a(\mathcal Z)$, $V=\mathbf Cf$ and
$N\cap M_\lambda(M(H))$ in place of $N$. By (a), (b) and (c),
$\mathfrak q,\mathfrak l\subset\mathfrak g\subset\mathfrak q+\mathfrak l$. The
line $\mathbf Cf$ is stable under the $M(Z)$ with $Z\in\mathfrak q$, and it lies
in $N\cap M_\lambda(M(H))$. For $(a,b)\in\mathcal Z$, $h(a)=h(b)$, so
$M(E^w_{ab})$ commutes with $M(H)$ and preserves $M_\lambda(M(H))$; hence
$N\cap M_\lambda(M(H))$ is stable under the $M(Z)$ with $Z\in\mathfrak l$. This
gives a subspace $V'$ with $f\in V'\subset N\cap M_\lambda(M(H))$, stable under
the $M(Z)$ with $Z\in\mathfrak a(\mathcal K_\leq)$.

\emph{Step 2.} Apply Lemma~\ref{asm-theta-agree-local-mult-archimedean-gen}
with $\mathfrak g=\mathfrak a(\mathcal K)$,
$\mathfrak q=\mathfrak a(\mathcal K_\leq)$, $\mathfrak l=\mathfrak a(\mathcal K_<)$,
$V'$ and $V'+M_{<\lambda}(M(H))$. For $(a,b)\in\mathcal K_<$,
$M(H)M(E^w_{ab})-M(E^w_{ab})M(H)=kM(E^w_{ab})$ with $k=h(a)-h(b)<0$, so
$V'+M_{<\lambda}(M(H))$ is stable under $M(E^w_{ab})$ by
Lemma~\ref{asm-theta-agree-local-mult-archimedean-low}(4), as
$V'\subset M_\lambda(M(H))$. This gives a subspace $W$ with
$V'\subset W\subset V'+M_{<\lambda}(M(H))\subset N+M_{<\lambda}(M(H))$, stable
under the $M(Z)$ with $Z\in\mathfrak a(\mathcal K)$, and $f\in W$.
\end{proof}
% Lean: cyc_le (ShiftLow.lean, about 30 lines, with alg, alg_le, alg_br, alg_mono, alg_sup,
% st_alg, dg_br, dg_eq and dg_hw, about 85 lines). The relations sb, pos and z are K, N and Z.
% Step 1 is the first gen_le, with g = alg (sb and h b <= h a), q = alg (a = b or pos), l =
% alg z and N inf eigenspace; Step 2 the second, with l = alg (sb and h a < h b) and shift_low.

\begin{lemma}\label{asm-theta-agree-local-mult-archimedean-generator}
Let $R,S\geq0$ and $\Lambda\colon\mathrm{Ix}(R,S)\to\mathbf Q$, and let $T$ be a
Lie module of type $(R,S)$ generated with Casimir gap in weight $\Lambda$
(Definition~\ref{asm-theta-agree-real-quotient-lie-module}). Then there is
$f\in T$ with the following properties.
\begin{enumerate}
\item $T(E_{aa})f=\Lambda_af$ for every $a$, and $T(E_{ab})f=0$ whenever
  $a\prec b$.
\item $T$ is the stable subspace generated by $f$.
\item Let $N$ be a Lie module of type $(R,S)$, $J$ a finite set, and for
  $j\in J$ let $\varphi_j\colon T\to N$ be a linear map with
  $\varphi_jT(Y)=N(Y)\varphi_j$ for every $Y$, and $x_j\in T$. If
  $u=\sum_j\varphi_j(x_j)$ satisfies $N(E_{aa})u=\Lambda_au$ for every $a$ and
  $N(E_{ab})u=0$ whenever $a\prec b$, then $u=\sum_jc_j\varphi_j(f)$ for some
  $c_j\in\mathbf C$.
\end{enumerate}
\end{lemma}
\leanref{Development/WeilBound/ThetaLifts/LocalTheory/ToralCohomology/SumPlSingle.lean}{PositivityTheory.ThKG.CayleyTransformDarboux.side_one}
\leanref{Development/WeilBound/ThetaLifts/LocalTheory/ToralCohomology/SumPlSingle.lean}{PositivityTheory.ThKG.CayleyTransformDarboux.poly_step}
\begin{proof}
Let $(A,\rho,\iota,f_A)$ be a highest-weight datum of $T$ as in
Definition~\ref{asm-theta-agree-real-quotient-lie-module}(3), put
$f=\iota(f_A)$, $\Omega=\Omega_T$ and $c=c(\Lambda)$, and let $B\subset A$ be the
smallest subspace containing $f_A$ and stable under every $\rho(a,b)$. Parts (1)
and (2) hold as $\iota\rho(a,b)=T(E_{ab})\iota$, and by the definition.

\emph{Step 1.} By step (i) of the proof of
Lemma~\ref{asm-theta-agree-real-quotient-local}, $B$ is the subspace called $B$
there, and by steps (iv) and (v) of that proof,
Lemma~\ref{asm-theta-agree-real-quotient-filtration} applies to $B$, with the
restrictions of the $\rho(a,b)$ and of $\iota$ and with $c_0=c$; and the stable
subspace generated by $\iota(B)$ is $T$, as it contains $f$. So:
($\alpha$) if $x\in T$ and $\Omega x-cx\in\iota(B)$, then $x\in\iota(B)$;
($\beta$) every $x\in T$ lies in a finite-dimensional subspace stable under the
$T(E_{ab})$ with $a,b$ in the same block, hence under $\Omega$. Also $\iota(B)$
is stable under these $T(E_{ab})$, as $B$ is stable under the $\rho(a,b)$, and so
under $\Omega$.

\emph{Step 2: a polynomial.} For each $j$ there is a polynomial $g_j$ with
$g_j(c)\neq0$ and $g_j(\Omega)x_j\in\iota(B)$. Indeed, let $F$ be as in
($\beta$) for $x_j$, and $g$ a monic polynomial with $g(\Omega)=0$ on $F$; then
$g(\Omega)x_j=0\in\iota(B)$. If $g(c)=0$, write $g=(X-c)g'$; then
$(\Omega-c)g'(\Omega)x_j\in\iota(B)$, so $g'(\Omega)x_j\in\iota(B)$ by
($\alpha$), and $g'$ has smaller degree. Repeating this gives $g_j$. Put
$g=\prod_jg_j$. Then $g(c)\neq0$, and $g(\Omega)x_j\in\iota(B)$ for every $j$,
as $\iota(B)$ is stable under $\Omega$.

\emph{Step 3: lower weights.} Let $h\colon\mathrm{Ix}(R,S)\to\mathbf R$ take
the value $-k$ at the element of each block that is preceded by $k$ elements of
that block, let $H=\operatorname{diag}(h)$ and $\lambda=\sum_ah(a)\Lambda_a$.
Then $\iota(B)\subset\mathbf Cf+T_{<\lambda}(T(H))$
(Lemma~\ref{asm-theta-agree-local-mult-archimedean-low}). Indeed, apply
Lemma~\ref{asm-theta-agree-local-mult-archimedean-cyclic} with $\Sigma$ a
point, $I=\mathrm{Ix}(R,S)$, $\mathcal K$ the set of pairs in the same block,
$\mathcal N$ the set of pairs with $a\prec b$, $\mathcal Z$ empty, $h$, $f$ and
$N=\mathbf Cf$: if $a\prec b$, then $h(b)<h(a)$; and if $a$ and $b$ lie in the
same block and $h(b)\leq h(a)$, then $a=b$ or $a\prec b$. It gives a subspace
$W\ni f$, stable under the $T(E_{ab})$ with $a,b$ in the same block, with
$W\subset\mathbf Cf+T_{<\lambda}(T(H))$. The $v\in B$ with $\iota(v)\in W$ form
a subspace that contains $f_A$ and is stable under every $\rho(a,b)$, so it is
$B$. Hence $g(\Omega)x_j=c_jf+z_j$ with $c_j\in\mathbf C$ and
$z_j\in T_{<\lambda}(T(H))$.

\emph{Step 4.} The computation in step (iv) of the proof of
Lemma~\ref{asm-theta-agree-real-quotient-local}, made with the operators
$N(E_{ab})$, which satisfy the bracket relations of the matrix units, and the
vector $u$, gives $\Omega_Nu=cu$. As $\varphi_j\Omega=\Omega_N\varphi_j$,
\[
g(c)u=g(\Omega_N)u=\sum_j\varphi_j(g(\Omega)x_j)=\sum_jc_j\varphi_j(f)+z,
\qquad z=\sum_j\varphi_j(z_j).
\]
By Lemma~\ref{asm-theta-agree-local-mult-archimedean-low}(3),
$z\in N_{<\lambda}(N(H))$. As $N(H)=\sum_ah(a)N(E_{aa})$, both $u$ and the
$\varphi_j(f)$ lie in $N_\lambda(N(H))$, and so does $z$. So $z=0$ by
Lemma~\ref{asm-theta-agree-local-mult-archimedean-low}(1), and
$u=\sum_jg(c)^{-1}c_j\varphi_j(f)$.
\end{proof}
% Lean: side_one (SumPlSingle.lean, about 105 lines, with poly_step, aeval_int and aeval_mem,
% about 30). B is cycK rho f with rhoV; hcV is cas_cyc and hsV the star condition for range
% iotaV; cas_bound_gl (Lemma asm-theta-agree-real-quotient-filtration) gives hinj and hfin
% (Step 1); hp and hP are Step 2 (LinearMap.exists_monic_and_aeval_eq_zero and poly_step); cyc_le
% for one act with sb = sameBlock, pos = posC, z = False and h = hD, hD (inl i) = -i, is Step 3
% (hVW); hcu (cas_hw with glK_of_lie), e1, e2 and e3 (low_eig, low_map, dg_hw) are Step 4.

\begin{lemma}\label{asm-theta-agree-local-mult-archimedean-tensor}
Let $\Sigma$ be a finite set and $R,S\geq0$, and for each $w\in\Sigma$ let
$\Lambda_w\colon\mathrm{Ix}(R,S)\to\mathbf Q$ and let $T_w$ be a Lie module of
type $(R,S)$ generated with Casimir gap in weight $\Lambda_w$
(Definition~\ref{asm-theta-agree-real-quotient-lie-module}). Let $P$ be a
complex vector space with, for each $w$, a structure $Y\mapsto P_w(Y)$ of Lie
module of type $(R,S)$ and a surjective linear map $\mu_w\colon P\to T_w$ with
$\mu_wP_w(Y)=T_w(Y)\mu_w$ for every $Y$. Let $M$ be a complex vector space and
$q\colon\bigotimes_{w\in\Sigma}P\to M$ a surjective linear map with
$q(\bigotimes_vy_v)=0$ whenever $\mu_w(y_w)=0$ for some $w$, and let
$Z\mapsto M(Z)$ be a linear map from $M_{R+S}(\mathbf C)^\Sigma$ to
$\operatorname{End}M$ with $M(Z)\,q=q\circ\sum_wP_w(Z_w)^{(w)}$ for every $Z$
(Definition~\ref{asm-theta-agree-real-fock}(4)). Then $M$ is a Lie module of
type $(R,S)$ over $\Sigma$, and there is $f\in M$ such that $(M,f)$ is a
highest-weight module of weight $(w,a)\mapsto\Lambda_w(a)$
(Definition~\ref{asm-theta-agree-local-mult-archimedean-modules}).
\end{lemma}
\leanref{Development/WeilBound/ThetaLifts/LocalTheory/ToralCohomology/SumPlSingle.lean}{PositivityTheory.ThKG.CayleyTransformDarboux.tens_hw}
\leanref{Development/WeilBound/ThetaLifts/LocalTheory/ToralCohomology/SumPlSingle.lean}{PositivityTheory.ThKG.CayleyTransformDarboux.sumPl_single}
\begin{proof}
For $y=(y_v)_v\in P^\Sigma$ write $\otimes y=\bigotimes_vy_v$, and for
$x\in P$ let $y[w\mapsto x]$ be $y$ with its entry at $w$ replaced by $x$. By
the hypothesis on $M(Z)$,
\[
M(\delta_wY)\,q(\otimes y)=q\bigl(\otimes y[w\mapsto P_w(Y)y_w]\bigr).\tag{$*$}
\]
For each $w$ let $f^0_w\in T_w$ be as in
Lemma~\ref{asm-theta-agree-local-mult-archimedean-generator}, choose
$f_w\in P$ with $\mu_w(f_w)=f^0_w$, and put $f=q(\otimes(f_v)_v)$. For $w$ and
$y$, the map $x\mapsto q(\otimes y[w\mapsto x])$ vanishes on the kernel of the
surjective map $\mu_w$, so there is a linear map $\varphi_{w,y}\colon T_w\to M$
with $\varphi_{w,y}(\mu_w(x))=q(\otimes y[w\mapsto x])$. Then
$q(\otimes y)=\varphi_{w,y}(\mu_w(y_w))$, and by ($*$)
$\varphi_{w,y}T_w(Y)=M(\delta_wY)\varphi_{w,y}$.

\emph{Bracket relation.} As $Z=\sum_w\delta_wZ_w$, it suffices to show
$M(\delta_wX)M(\delta_{w'}Y)-M(\delta_{w'}Y)M(\delta_wX)=M(\delta_wX\,\delta_{w'}Y-\delta_{w'}Y\,\delta_wX)$
on the vectors $q(\otimes y)$, which span $M$. If $w\neq w'$, both sides vanish
by ($*$), as the two operators change different entries of $y$ and
$\delta_wX\,\delta_{w'}Y=0$. If $w=w'$, write $q(\otimes y)=\varphi_{w,y}(\mu_w(y_w))$;
the relation follows from that of $T_w$ through $\varphi_{w,y}$.

\emph{Conditions (a) and (b).} As $f=\varphi_{w,(f_v)}(f^0_w)$,
$M(E^w_{aa})f=\varphi_{w,(f_v)}(T_w(E_{aa})f^0_w)=\Lambda_w(a)f$, and in the
same way $M(E^w_{ab})f=0$ for $a\prec b$.

\emph{Condition (c).} Let $W\ni f$ be stable. For a set
$\Sigma_0\subset\Sigma$, every $q(\otimes y)$ with $y_v=f_v$ for
$v\notin\Sigma_0$ lies in $W$; this is proved by induction on $\Sigma_0$. For
$\Sigma_0$ empty it is $f\in W$. Let $\Sigma_1=\Sigma_0\cup\{w\}$ with
$w\notin\Sigma_0$, and $y_v=f_v$ for $v\notin\Sigma_1$. The set of $x\in T_w$ with
$\varphi_{w,y}(x)\in W$ is stable under the $T_w(Y)$ and contains
$f^0_w$, as $\varphi_{w,y}(f^0_w)=q(\otimes y[w\mapsto f_w])\in W$ by
induction. By Lemma~\ref{asm-theta-agree-local-mult-archimedean-generator}(2)
it is $T_w$, so $q(\otimes y)=\varphi_{w,y}(\mu_w(y_w))\in W$. For
$\Sigma_0=\Sigma$ this gives $W=M$, as $q$ is surjective.

\emph{Condition (d).} Let $u\in M$ satisfy (a) and (b). For a set
$\Sigma_0\subset\Sigma$, $u$ lies in the span of the $q(\otimes y)$ with
$y_v=f_v$ for $v\in\Sigma_0$; this is proved by induction on $\Sigma_0$. For
$\Sigma_0$ empty it holds as $q$ is surjective. Let
$\Sigma_1=\Sigma_0\cup\{w\}$ with $w\notin\Sigma_0$, and
$u=\sum_kc_kq(\otimes y_k)$ with $y_{k,v}=f_v$ for $v\in\Sigma_0$. Then
$u=\sum_k\varphi_{w,y_k}(c_k\mu_w(y_{k,w}))$. By the bracket relation,
$Y\mapsto M(\delta_wY)$ makes $M$ a Lie module of type $(R,S)$; the
$\varphi_{w,y_k}$ intertwine it with $T_w$, and $u$ satisfies the conditions of
Lemma~\ref{asm-theta-agree-local-mult-archimedean-generator}(3) for
$\Lambda_w$. So $u=\sum_kc'_k\varphi_{w,y_k}(f^0_w)=\sum_kc'_kq(\otimes y_k[w\mapsto f_w])$,
which lies in the span for $\Sigma_1$. For $\Sigma_0=\Sigma$ the span is
$\mathbf Cf$.
\end{proof}
% Lean: tens_hw (SumPlSingle.lean, about 120 lines, with atPl_tprod and sumPl_single). f0, f
% and phi are from side_one, hmk and exists_factor; a1 is (*), hphia the intertwining; L1 and hL
% the bracket relation; the four fields of HWMod are proved in the order of the text, (c) by
% Finset.induction_on with genG_le', (d) by Finset.induction_on with hf3 (side_one (3)) for
% act composed with LinearMap.single w.

\begin{lemma}\label{asm-theta-agree-local-mult-archimedean-action}
Let $(B,G,\mathcal X,\omega,\omega')$ be a triple of ranks $(a,a')$ for a theta
base $B$ of a CM field $E$ (Definition~\ref{asm-theta-agree-triple}), with real
maps $(\Pi,p,p')$ matched to it
(Definition~\ref{asm-theta-agree-matched}(2)). Let $(P,Q)$ and $(R,S)$ be the
signatures of $V$ and $W$, and for every infinite place $w$ of $F$ put
$\kappa_w=m_\alpha(w)-m_\beta(w)$ and let $C_w$ be a frame of
$i\tau_w(J_W)$ of type $(R,S)$ (Definition~\ref{asm-theta-agree-real-fock}(5))
with $\theta_{W,w}C_w=C_wJ_{R,S}$. With the notation $c_{R,S}=\tfrac12(R-S)$ of
Definition~\ref{asm-theta-agree-real-quotient-shapes}(1), let
$\Lambda_w\colon\mathrm{Ix}(R,S)\to\mathbf Q$ take the value
$\tfrac12\kappa_w-c_{R,S}+c_{P,Q}+\tfrac12m_\beta(w)$ at the first $P$
elements of the first block, $c_{P,Q}+\tfrac12m_\beta(w)$ at its other
elements, $\tfrac12\kappa_w+c_{R,S}-c_{P,Q}+\tfrac12m_\beta(w)$ at the first
$Q$ elements of the second block, and $-c_{P,Q}+\tfrac12m_\beta(w)$ at its
other elements. Then there are a structure $Z\mapsto\varrho(Z)$ on $\Pi$ of Lie
module of type $(R,S)$ over the set of infinite places of $F$ and a vector
$f_0\in\Pi$ such that $(\Pi,f_0)$ is a highest-weight module of weight
$(w,b)\mapsto\Lambda_w(b)$
(Definition~\ref{asm-theta-agree-local-mult-archimedean-modules}), and such
that the following holds: if $\gamma$ is a line in $U_W$ generated by
$Y=(Y_w)_w$ relative to $(C_w)_w$ (Definition~\ref{asm-theta-agree-real-fock}(6))
and $f\in\operatorname{Pol}(M)$, there is $f_1$ such that
$t\mapsto\omega(1,\gamma(t))$ acts at $\infty$ to first order by $f\mapsto f_1$
and $p(f_1)=\varrho(Y)p(f)$.
\end{lemma}
\leanref{Development/WeilBound/ThetaLifts/LocalTheory/ToralCohomology/PiInfUnitNe.lean}{PositivityTheory.ThKG.CayleyTransformDarboux.exists_mod}
\leanref{Development/WeilBound/ThetaLifts/LocalTheory/ToralCohomology/PiInfUnitNe.lean}{PositivityTheory.ThKG.CayleyTransformDarboux.lamW}
\leanref{Development/WeilBound/ThetaLifts/LocalTheory/ToralCohomology/PiInfUnitNe.lean}{PositivityTheory.ThKG.CayleyTransformDarboux.Dic}
\leanref{Development/WeilBound/ThetaLifts/LocalTheory/ToralCohomology/PiInfUnitNe.lean}{PositivityTheory.ThKG.CayleyTransformDarboux.exists_act}
\leanref{Development/WeilBound/ThetaLifts/LocalTheory/ToralCohomology/PiInfUnitNe.lean}{PositivityTheory.ThKG.CayleyTransformDarboux.ker_stable}
\leanref{Development/WeilBound/ThetaLifts/LocalTheory/ToralCohomology/PiInfUnitNe.lean}{PositivityTheory.ThKG.CayleyTransformDarboux.upq_split}
\begin{proof}
Let $p_0,q_0$ be the integers of $G$, so that $(P,Q)=(a-p_0,p_0)$ and
$(R,S)=(a+a'-p_0-q_0,p_0+q_0)$, and let $\mathcal T=\bigotimes_w\mathcal P(P,Q;R,S)$.
For an infinite place $w$ and $Y\in M_{R+S}(\mathbf C)$ put
$\tilde d_w(Y)=d_2(Y)+\tfrac12m_\beta(w)\operatorname{tr}(Y)$, and for a family
$Y=(Y_w)$ put $L_Y=\sum_w\tilde d_w(Y_w)^{(w)}$.

\emph{Step 1: the Fock model.} By Lemma~\ref{asm-theta-agree-real-frames}
there are frames $A_w$ of $\tau_w(J_V)$ of type $(P,Q)$ with
$\theta_{V,w}A_w=A_wJ_{P,Q}$, and by
Lemma~\ref{asm-theta-agree-real-fock-map} a Fock map $\Phi$ for
$(M,(A_w),(C_w))$, with $a_0=1$ and $b_0=i$, whose image is
$\operatorname{Pol}(M)$. As in the paragraph \emph{Derivatives} of the proof of
Lemma~\ref{asm-theta-agree-real}, which uses only that $\Phi$ is such a Fock
map, for $T\in\mathcal T$:
\begin{enumerate}
\item[(i)] if $\gamma$ is a line in $U_V$ generated by $(X_w)$ relative to
  $(A_w)$, then $t\mapsto\omega(\gamma(t),1)$ acts at $\infty$ to first order
  by $\Phi(T)\mapsto\Phi(\sum_w(d_1(X_w)-\tfrac12\kappa_w\operatorname{tr}X_w)^{(w)}T)$;
\item[(ii)] if $\gamma$ is a line in $U_W$ generated by $Y=(Y_w)$ relative to
  $(C_w)$, then $t\mapsto\omega(1,\gamma(t))$ acts at $\infty$ to first order
  by $\Phi(T)\mapsto\Phi(L_YT)$.
\end{enumerate}
Put $q=p\circ\Phi\colon\mathcal T\to\Pi$. It is surjective, as
$\Phi(\mathcal T)=\operatorname{Pol}(M)$ and $p(\operatorname{Pol}(M))=\Pi$.

\emph{Step 2: $q$ kills the first action.} For every $w$,
$X\in M_{P+Q}(\mathbf C)$ and $T\in\mathcal T$,
$q((d_1(X)-\tfrac12\kappa_w\operatorname{tr}X)^{(w)}T)=0$. For
$X\in\mathfrak u(P,Q)$, the family with $X$ at $w$ and $0$ elsewhere generates a
line in $U_V$ relative to $(A_w)$
(Lemma~\ref{asm-theta-agree-real-lines}(2)), and the claim follows from (i) and
condition (a) of Definition~\ref{asm-theta-agree-matched}(2). Every $X$ is
$X_1+iX_2$ with $X_1,X_2\in\mathfrak u(P,Q)$ (proof of
Lemma~\ref{asm-theta-agree-real-tensor}), and the expression is linear in $X$.

\emph{Step 3: the action.} The kernel of $q$ is stable under every $L_Y$.
Suppose first that $Y_w\in\mathfrak u(R,S)$ for every $w$, and let $\gamma$ be
the line in $U_W$ that $Y$ generates relative to $(C_w)$
(Lemma~\ref{asm-theta-agree-real-lines}(2)). Let $q(T)=0$. Then
$p(\Phi(T))=0=p'(0)$ with $0\in\operatorname{Pol}(M')$, and
$t\mapsto\omega'(1,\gamma(t))$ acts at $\infty$ to first order by $0\mapsto0$, as
these operators are linear. By (ii) and condition (b) of
Definition~\ref{asm-theta-agree-matched}(2), $q(L_YT)=p'(0)=0$. A general family
is $Y=Y'+iY''$ with all $Y'_w,Y''_w\in\mathfrak u(R,S)$, and $L_Y$ is linear
in $Y$. As $q$ is surjective, there are operators $\varrho(Y)$ on $\Pi$, linear
in $Y$, with $\varrho(Y)q=qL_Y$.

\emph{Step 4: the quotient at one place.} Fix $w$. By
Lemma~\ref{asm-theta-agree-real-parity}, $\kappa_w=2k$ for an integer $k$, and
$a+a'=2d$ for an integer $d$ as $G$ is proper; so $c_{R,S}=d-S$. Put
$n=k-c_{R,S}$ and $m=k+c_{R,S}$, which are integers. As $G$ is proper,
$a+a'<\kappa_w$, so $k>d$ and $n>S\geq0$; and $2(p_0+q_0)+1<a'<a$, so $S<d$
and $m>0$. Also $P\leq R$ and $Q\leq S$, as $q_0\leq a'$; $n+2c_{P,Q}=n+a-2p_0>0$,
as $2p_0<a$; and $m-2c_{P,Q}=k+\tfrac12(a'-a)+p_0-q_0>a'+p_0-q_0>0$, as $k>d$
and $2q_0<a'$. By Lemma~\ref{asm-theta-agree-real-quotient-weights}(1), its map
$\Lambda^+$ for $(P,Q,R,S)$, $n$ and $m$ is a gap weight for
$(P,Q,k,\tfrac12m_\beta(w))$, and $\Lambda^++\tfrac12m_\beta(w)=\Lambda_w$. Let
$\Theta_w$ be the quotient of $\mathcal P(P,Q;R,S)$ by the span of the
$d_1(X)x-k\operatorname{tr}(X)x$, with quotient map $\mu_w$ and the operators
$\vartheta_w(Y)$ induced by $\tilde d_w(Y)$
(Lemma~\ref{asm-theta-agree-real-fock-relations}(1)). By
Lemma~\ref{asm-theta-agree-real-quotient-nonzero}(1), $\Theta_w\neq0$, and by
Proposition~\ref{asm-theta-agree-real-quotient-side}, with $c_0=k$ and
$\tau_0=\tfrac12m_\beta(w)$, $\Theta_w$ with the $\vartheta_w(Y)$ is a Lie
module of type $(R,S)$ generated with Casimir gap in weight $\Lambda_w$.

\emph{Step 5: the highest-weight vector.} Apply
Lemma~\ref{asm-theta-agree-local-mult-archimedean-tensor} with $\Sigma$ the set
of infinite places, $\mathcal P(P,Q;R,S)$ with the structures $\tilde d_w$, which are
Lie modules by Lemma~\ref{asm-theta-agree-real-fock-relations}(1) as
$\operatorname{tr}(YY'-Y'Y)=0$, the modules $\Theta_w$ and the maps $\mu_w$,
$M=\Pi$, $q$ and $\varrho$. The condition on $q$ holds: the kernel of $\mu_w$ is
spanned by the $d_1(X)x-k\operatorname{tr}(X)x$, and a pure tensor whose entry at
$w$ is of this form is
$(d_1(X)-\tfrac12\kappa_w\operatorname{tr}X)^{(w)}$ applied to a pure tensor,
which $q$ kills by Step 2. This gives $f_0$.

\emph{Step 6: compatibility.} Let $\gamma$ be generated by $Y$ relative to
$(C_w)$, and $f\in\operatorname{Pol}(M)$; write $f=\Phi(T)$. By (ii),
$t\mapsto\omega(1,\gamma(t))$ acts at $\infty$ to first order by $f\mapsto f_1$
with $f_1=\Phi(L_YT)$, and $p(f_1)=q(L_YT)=\varrho(Y)q(T)=\varrho(Y)p(f)$.
\end{proof}
% Lean: exists_mod (PiInfUnitNe.lean, about 90 lines, with sumL, sumL_eq, ker_stable, exists_act,
% hasExp_div and actsD_zero, about 50). Dic bundles the Fock-model lemmas: laws (Lemma
% asm-theta-agree-real-weak), fock and range (Lemma asm-theta-agree-real-fock-map), fst and snd
% (Lemma asm-theta-agree-real-lie), line (Lemma asm-theta-agree-real-lines), frame (Lemma
% asm-theta-agree-real-frames), even_im_hermitian_sub (Lemma asm-theta-agree-real-parity). Step 1
% is A, fr, Phi, hr, hW2 (snd with beta and m_beta) and hW1 (fst with beta/alpha and exponents
% m_beta - m_alpha by hasExp_div), with s2; Step 2 is s1 (hR.quot.lie, ker_stable, upq_split);
% Step 3 is exists_act (hR.lie with f' = 0 and actsD_zero); Step 4 is side, from side_of with
% comb_pos and nontrivial_of_pair (pair_pos), the weight lamW = laPos + m_beta/2; Step 5 is
% tens_hw with aP = twAct (m_beta/2) and mk = thMkC; Step 6 is the final refine.

\begin{definition}\label{asm-theta-agree-local-mult-archimedean-embeddings}
Let $R,S,R_1,S_1\geq0$, and use the notation of
Definitions~\ref{asm-theta-agree-real-fock}
and~\ref{asm-theta-agree-real-quotient-lie-module}.
\begin{enumerate}
\item An \emph{embedding} of $\mathrm{Ix}(R_1,S_1)$ into $\mathrm{Ix}(R,S)$ is
  a quadruple $(e,o_1,o_2,\sigma)$ of an injective map
  $e\colon\mathrm{Ix}(R_1,S_1)\to\mathrm{Ix}(R,S)$, an element $o_1$ of the
  first block and an element $o_2$ of the second block of $\mathrm{Ix}(R,S)$,
  and a real number $\sigma\neq0$, such that: every element of
  $\mathrm{Ix}(R,S)$ is $o_1$, $o_2$ or in the image of $e$, and $o_1$ and $o_2$
  are not in the image; $e(a)$ and $e(b)$ lie in the same block if and only if
  $a$ and $b$ do; $e(a)\prec e(b)$ whenever $a\prec b$; and if
  $a\in\{o_1,o_2\}$ and $a$ and $e(b)$ lie in the same block, then
  $a\prec e(b)$ if $\sigma>0$ and $e(b)\prec a$ if $\sigma<0$. Its \emph{matrix}
  is $P_e\in M_{\mathrm{Ix}(R,S)\times\mathrm{Ix}(R_1,S_1)}(\mathbf C)$, with
  entry $1$ at $(e(b),b)$ for every $b$ and $0$ elsewhere. With
  $\mathbf 1_a\in\mathbf C^{\mathrm{Ix}(R,S)}$ the vector with entry $1$ at $a$
  and $0$ elsewhere, put $x_e=\mathbf 1_{o_1}+\mathbf 1_{o_2}$ and
  $y_e=\mathbf 1_{o_1}-\mathbf 1_{o_2}$.
\item Let $R=R_1+1$ and $S=S_1+1$. The \emph{first-index embedding} maps the
  $k$-th element of each block of $\mathrm{Ix}(R_1,S_1)$ to the $(k+1)$-st
  element of the block of the same name, with $o_1$ and $o_2$ the first
  elements of the two blocks and $\sigma=1$. The \emph{last-index embedding}
  maps the $k$-th element of each block to the $k$-th element of the block of
  the same name, with $o_1$ and $o_2$ the last elements of the two blocks and
  $\sigma=-1$. Both are embeddings, and both map each block into the block of
  the same name.
\end{enumerate}
\end{definition}
\leanref{Development/WeilBound/ThetaLifts/LocalTheory/ToralCohomology/ParZero.lean}{PositivityTheory.ThKG.CayleyTransformDarboux.Emb}
\leanref{Development/WeilBound/ThetaLifts/LocalTheory/ToralCohomology/EmbeddingTransposeVectorZero.lean}{PositivityTheory.ThKG.CayleyTransformDarboux.Emb.P}
\leanref{Development/WeilBound/ThetaLifts/LocalTheory/ToralCohomology/EmbeddingTransposeVectorZero.lean}{PositivityTheory.ThKG.CayleyTransformDarboux.Emb.x}
\leanref{Development/WeilBound/ThetaLifts/LocalTheory/ToralCohomology/EmbeddingTransposeVectorZero.lean}{PositivityTheory.ThKG.CayleyTransformDarboux.Emb.y}
\leanref{Development/WeilBound/ThetaLifts/LocalTheory/ToralCohomology/EmbeddingTransposeVectorZero.lean}{PositivityTheory.ThKG.CayleyTransformDarboux.embA}
\leanref{Development/WeilBound/ThetaLifts/LocalTheory/ToralCohomology/EmbeddingTransposeVectorZero.lean}{PositivityTheory.ThKG.CayleyTransformDarboux.embB}

\begin{definition}\label{asm-theta-agree-local-mult-archimedean-relations}
Let $\Sigma$ be a finite set, $(e,o_1,o_2,\sigma)$ an embedding of
$\mathrm{Ix}(R_1,S_1)$ into $\mathrm{Ix}(R,S)$
(Definition~\ref{asm-theta-agree-local-mult-archimedean-embeddings}), $M$ and
$M_1$ Lie modules over $\Sigma$ of types $(R,S)$ and $(R_1,S_1)$
(Definition~\ref{asm-theta-agree-local-mult-archimedean-modules}), and
$c,c'\colon\Sigma\to\mathbf C$. A linear map $D\colon M\to M_1$ satisfies the
\emph{parabolic relations} for $(e,c,c')$ if for all $w\in\Sigma$, all
$a,b\in\mathrm{Ix}(R_1,S_1)$ and all $x\in M$:
\begin{enumerate}
\item $D\,M(E^w_{e(a)e(b)})x=M_1(E^w_{ab})Dx$;
\item $D\,M(E^w_{o_1o_1}+E^w_{o_2o_2})x=c_wDx$;
\item $D\,M(E^w_{o_1o_2}+E^w_{o_2o_1})x=c'_wDx$;
\item $D\,M(E^w_{o_1o_1}-E^w_{o_1o_2}+E^w_{o_2o_1}-E^w_{o_2o_2})x=0$;
\item $D\,M(E^w_{e(b)o_1}-E^w_{e(b)o_2})x=0$;
\item $D\,M(E^w_{o_1e(b)}+E^w_{o_2e(b)})x=0$.
\end{enumerate}
\end{definition}
\leanref{Development/WeilBound/ThetaLifts/LocalTheory/ToralCohomology/ParZero.lean}{PositivityTheory.ThKG.CayleyTransformDarboux.IsPar}

\begin{lemma}\label{asm-theta-agree-local-mult-archimedean-zero}
Let $\Sigma$ be a finite set, $(e,o_1,o_2,\sigma)$ an embedding of
$\mathrm{Ix}(R_1,S_1)$ into $\mathrm{Ix}(R,S)$
(Definition~\ref{asm-theta-agree-local-mult-archimedean-embeddings}), $(M,f)$
a highest-weight module of type $(R,S)$ over $\Sigma$ and of weight $\Lambda$,
and $M_1$ a Lie module of type $(R_1,S_1)$ over $\Sigma$
(Definition~\ref{asm-theta-agree-local-mult-archimedean-modules}). Let
$c,c'\colon\Sigma\to\mathbf C$ with $c_w=\Lambda_w(o_1)+\Lambda_w(o_2)$ for
every $w$, and let $D\colon M\to M_1$ satisfy the parabolic relations for
$(e,c,c')$ (Definition~\ref{asm-theta-agree-local-mult-archimedean-relations}).
If $Df=0$, then $D=0$.
\end{lemma}
\leanref{Development/WeilBound/ThetaLifts/LocalTheory/ToralCohomology/ParZero.lean}{PositivityTheory.ThKG.CayleyTransformDarboux.par_zero}
\begin{proof}
Let $h\colon\mathrm{Ix}(R,S)\to\mathbf R$ be $\sigma$ at $o_1$ and $o_2$ and
$0$ on the image of $e$, let $H$ be the family with $\operatorname{diag}(h)$ at
every place, and put $\lambda=\sum_w\sum_ah(a)\Lambda_w(a)=\sigma\sum_wc_w$.

\emph{Step 1.} $D$ vanishes on $M_{<\lambda}(M(H))$
(Lemma~\ref{asm-theta-agree-local-mult-archimedean-low}). Indeed
$H=\sigma\sum_w(E^w_{o_1o_1}+E^w_{o_2o_2})$, so $DM(H)=\sigma\sum_wc_wD=\lambda D$
by relation (2), and part (2) of that lemma applies.

\emph{Step 2.} Apply Lemma~\ref{asm-theta-agree-local-mult-archimedean-cyclic}
with $I=\mathrm{Ix}(R,S)$, $\mathcal K$ the set of pairs in the same block,
$\mathcal N$ the set of pairs with $a\prec b$, $\mathcal Z$ the set of pairs
$(e(a),e(b))$ with $a$ and $b$ in the same block, $h$, $f$ and $N=\ker D$.
Condition (a) holds. For (b), let $a\prec b$. If $a\in\{o_1,o_2\}$, then $b$ is
in the image of $e$, as $o_1$ and $o_2$ lie in different blocks; $\sigma<0$
would give $b\prec a$; so $\sigma>0$ and $h(b)=0<h(a)$. If $b\in\{o_1,o_2\}$,
then in the same way $\sigma<0$ and $h(b)<0=h(a)$. Otherwise $h(a)=h(b)=0$. On
the pairs in $\mathcal Z$, $h$ vanishes. For (c), let $a$ and $b$ lie in the
same block with $h(b)\leq h(a)$. If both lie in $\{o_1,o_2\}$, then $a=b$. If
only $a$ does, then $\sigma=h(a)\geq h(b)=0$, so $\sigma>0$ and $a\prec b$; if
only $b$ does, then $\sigma<0$ and $a\prec b$; if neither does,
$(a,b)\in\mathcal Z$. Finally, $f\in\ker D$, and $\ker D$ is stable under the
$M(E^w_{e(a)e(b)})$ by relation (1). So there is a subspace $V\ni f$, stable
under $M(\mathfrak k)$, where $\mathfrak k$ is the span of the $E^w_{ab}$ with
$a$ and $b$ in the same block, with $V\subset\ker D+M_{<\lambda}(M(H))$, which is
$\ker D$ by Step 1.

\emph{Step 3.} Let $\mathfrak p$ be the space of families $Z$ with
$M(Z)\ker D\subset\ker D$. It contains every $Z$ with $DM(Z)=tD$ for some
$t\in\mathbf C$. So by the parabolic relations it contains $E^w_{e(a)e(b)}$,
$Z_2=E^w_{o_1o_1}+E^w_{o_2o_2}$, $Z_3=E^w_{o_1o_2}+E^w_{o_2o_1}$,
$Z_4=E^w_{o_1o_1}-E^w_{o_1o_2}+E^w_{o_2o_1}-E^w_{o_2o_2}$,
$Z_5=E^w_{e(b)o_1}-E^w_{e(b)o_2}$ and $Z_6=E^w_{o_1e(b)}+E^w_{o_2e(b)}$, for
all $w$, $a$ and $b$. Every $E^w_{ab}$ lies in $\mathfrak k+\mathfrak p$. This
is clear if $a$ and $b$ lie in the same block or both lie in the image of $e$.
Next, $E^w_{o_1o_2}=\tfrac12(Z_3-Z_4+E^w_{o_1o_1}-E^w_{o_2o_2})$ and
$E^w_{o_2o_1}=\tfrac12(Z_3+Z_4-E^w_{o_1o_1}+E^w_{o_2o_2})$. Let $b'=e(b)$. If
$b'$ lies in the block of $o_1$, then $E^w_{b'o_1},E^w_{o_1b'}\in\mathfrak k$,
$E^w_{b'o_2}=E^w_{b'o_1}-Z_5$ and $E^w_{o_2b'}=Z_6-E^w_{o_1b'}$; if it lies in
the block of $o_2$, then $E^w_{b'o_1}=Z_5+E^w_{b'o_2}$ and
$E^w_{o_1b'}=Z_6-E^w_{o_2b'}$. As every element of $\mathrm{Ix}(R,S)$ is $o_1$,
$o_2$ or in the image of $e$, this covers every pair.

\emph{Step 4.} By Lemma~\ref{asm-theta-agree-local-mult-archimedean-gen}, with
$A=M_{R+S}(\mathbf C)^\Sigma$, $\mathfrak g=A$, $\mathfrak q=\mathfrak k$,
$\mathfrak l=\mathfrak p$, $V$ and $N=\ker D$, there is a stable subspace $W$
with $V\subset W\subset\ker D$. As $f\in W$, $W=M$ by condition (c) of
Definition~\ref{asm-theta-agree-local-mult-archimedean-modules}(2), so $D=0$.
\end{proof}
% Lean: par_zero (ParZero.lean, about 75 lines, with Emb.h, Emb.cases, Emb.eq_of and Emb.sum).
% hs and hlow are Step 1 (low_ker); cyc_le with z a b = sameBlock a b and both in the image of
% m.e is Step 2; K = alg sameBlock and P = stab act (ker D), with mP, mK, hij, tri and hb, are
% Step 3, and top is K + P = everything; gen_le with g = top and hM.gen are Step 4.

\begin{lemma}\label{asm-theta-agree-local-mult-archimedean-proportional}
Let $\Sigma$ be a finite set, $(e,o_1,o_2,\sigma)$ an embedding of
$\mathrm{Ix}(R_1,S_1)$ into $\mathrm{Ix}(R,S)$
(Definition~\ref{asm-theta-agree-local-mult-archimedean-embeddings}), and
$(M,f)$ and $(M_1,f_1)$ highest-weight modules over $\Sigma$ of types $(R,S)$
and $(R_1,S_1)$ and of weights $\Lambda$ and $\Lambda_1$
(Definition~\ref{asm-theta-agree-local-mult-archimedean-modules}), with
$\Lambda_w(e(b))=\Lambda_{1,w}(b)$ for all $w$ and $b$. Let
$c,c'\colon\Sigma\to\mathbf C$, and let $D,D'\colon M\to M_1$ satisfy the
parabolic relations for $(e,c,c')$
(Definition~\ref{asm-theta-agree-local-mult-archimedean-relations}), with
$D\neq0$. Then $D'=tD$ for some $t\in\mathbf C$.
\end{lemma}
\leanref{Development/WeilBound/ThetaLifts/LocalTheory/ToralCohomology/ParZero.lean}{PositivityTheory.ThKG.CayleyTransformDarboux.par_line}
\leanref{Development/WeilBound/ThetaLifts/LocalTheory/ToralCohomology/ParZero.lean}{PositivityTheory.ThKG.CayleyTransformDarboux.HWMod.centre}
\leanref{Development/WeilBound/ThetaLifts/LocalTheory/ToralCohomology/ParZero.lean}{PositivityTheory.ThKG.CayleyTransformDarboux.IsPar.sub}
\begin{proof}
(i) For every $w$, $\sum_aM(E^w_{aa})$ is $\sum_a\Lambda_w(a)$ times the
identity of $M$, and likewise for $M_1$. Indeed $\sum_aE^w_{aa}=\delta_w1$
commutes with every family, so $M(\delta_w1)$ commutes with every $M(Z)$ and its
eigenspace for $\sum_a\Lambda_w(a)$ is stable. It contains $f$ by condition (a)
of Definition~\ref{asm-theta-agree-local-mult-archimedean-modules}(2), so it is
$M$ by condition (c).

(ii) $c_w=\Lambda_w(o_1)+\Lambda_w(o_2)$ for every $w$. Choose $x$ with
$Dx\neq0$, apply $D$ to $\sum_aM(E^w_{aa})x=\sum_a\Lambda_w(a)x$, and split
each sum over $a\in\{o_1,o_2\}$ and over $a=e(b)$. By relations (2) and (1) and
by (i) for $M_1$, the left side becomes $c_wDx+\sum_b\Lambda_{1,w}(b)Dx$. As
$\Lambda_w(e(b))=\Lambda_{1,w}(b)$, this gives
$c_wDx=(\Lambda_w(o_1)+\Lambda_w(o_2))Dx$.

(iii) If $D''$ satisfies the parabolic relations for $(e,c,c')$, then
$D''f\in\mathbf Cf_1$. By relation (1),
$M_1(E^w_{aa})D''f=D''M(E^w_{e(a)e(a)})f=\Lambda_{1,w}(a)D''f$, and for
$a\prec b$, $M_1(E^w_{ab})D''f=D''M(E^w_{e(a)e(b)})f=0$, as $e(a)\prec e(b)$.
Condition (d) for $(M_1,f_1)$ applies.

(iv) Write $Df=tf_1$ and $D'f=t'f_1$. If $t=0$, then $D=0$ by (ii) and
Lemma~\ref{asm-theta-agree-local-mult-archimedean-zero}, which is excluded. The
map $D'-(t'/t)D$ satisfies the parabolic relations for $(e,c,c')$, as they are
linear in the map, and it kills $f$; so $D'=(t'/t)D$ by
Lemma~\ref{asm-theta-agree-local-mult-archimedean-zero}.
\end{proof}
% Lean: par_line (ParZero.lean, about 22 lines), with HWMod.centre (about 25 lines, (i)), hc
% ((ii)), val ((iii)), t0 and IsPar.sub with par_zero ((iv)).

\begin{lemma}\label{asm-theta-agree-local-mult-archimedean-frame}
Let $(e,o_1,o_2,\sigma)$ be an embedding of $\mathrm{Ix}(R_1,S_1)$ into
$\mathrm{Ix}(R,S)$ that maps each block into the block of the same name, with
matrix $P_e$ and vectors $x_e$ and $y_e$
(Definition~\ref{asm-theta-agree-local-mult-archimedean-embeddings}). Let
$d=R+S$ and $d_1\geq0$, let $M\in M_d(\mathbf C)$ with $\bar M^{\mathsf T}=M$
and a Cartan involution $\theta$ of $M$
(Definition~\ref{asm-theta-agree-doubling}(1)), and write
$\langle u,v\rangle=u^{\mathsf T}M\bar v$ for $u,v\in\mathbf C^d$. Let
$M_1,\theta_1\in M_{d_1}(\mathbf C)$ and let $C_1$ be a frame of $M_1$ of type
$(R_1,S_1)$ (Definition~\ref{asm-theta-agree-real-fock}(5)) with
$\theta_1C_1=C_1J_{R_1,S_1}$. Let $x,y\in\mathbf C^d$, $c\in\mathbf C$ and
$\iota\in M_{d\times d_1}(\mathbf C)$ with $x\neq0$, $\langle x,x\rangle=0$,
$c\neq0$, $y=c\,\theta x$, $\langle\iota u,x\rangle=\langle\iota u,y\rangle=0$
for every $u\in\mathbf C^{d_1}$, $\iota^{\mathsf T}M\bar\iota=M_1$ and
$\theta\iota=\iota\theta_1$. Then there are a frame $C$ of $M$ of type $(R,S)$
and $n,n'\in\mathbf C$ with $n\neq0$, $\theta C=CJ_{R,S}$, $x=n\,Cx_e$,
$y=n'\,Cy_e$ and $CP_e=\iota C_1$.
\end{lemma}
\leanref{Development/WeilBound/ThetaLifts/LocalTheory/ToralCohomology/EmbeddingTransposeVectorZero.lean}{PositivityTheory.ThKG.CayleyTransformDarboux.exists_frame}
\begin{proof}
As $M$ is Hermitian, $\langle v,u\rangle=\overline{\langle u,v\rangle}$, and as
$\theta^{\mathsf T}M\bar\theta=M$, $\langle\theta u,\theta v\rangle=\langle u,v\rangle$.
Put $z=\theta x$. Then $\langle z,z\rangle=\langle x,x\rangle=0$, and
$\langle x,z\rangle=\langle\theta x,\theta z\rangle=\langle z,x\rangle$ as
$\theta^2=1$; so $\langle z,x\rangle$ is real, and it is positive by the
positivity condition on $\theta$, as $x\neq0$. Put $k=\langle z,x\rangle$ and
$N=(2k)^{-1/2}$. For every column $\iota_l$ of $\iota$,
$\langle\iota_l,x\rangle=0$, and $\langle\iota_l,z\rangle=0$ as
$\langle\iota_l,y\rangle=\bar c\langle\iota_l,z\rangle$ and $c\neq0$; so the
columns of $\iota C_1$ are orthogonal to $x$ and $z$.

Let $C\in M_{d\times\mathrm{Ix}(R,S)}(\mathbf C)$ have the columns
$C_{o_1}=N(x+z)$, $C_{o_2}=N(x-z)$ and $C_{e(b)}$ equal to the $b$-th column of
$\iota C_1$. Then $\langle C_{o_1},C_{o_1}\rangle=N^2(\langle x,z\rangle+\langle z,x\rangle)=1$,
$\langle C_{o_2},C_{o_2}\rangle=-1$,
$\langle C_{o_1},C_{o_2}\rangle=N^2(\langle z,x\rangle-\langle x,z\rangle)=0$,
the columns $C_{e(b)}$ are orthogonal to $C_{o_1}$ and $C_{o_2}$, and
$\langle C_{e(a)},C_{e(b)}\rangle$ is the entry at $(a,b)$ of
$C_1^{\mathsf T}\iota^{\mathsf T}M\bar\iota\bar C_1=C_1^{\mathsf T}M_1\bar C_1=J_{R_1,S_1}$.
As $e$ is injective and maps each block into the block of the same name,
$C^{\mathsf T}M\bar C=J_{R,S}$. Then $(J_{R,S}C^{\mathsf T}M)\bar C=1$; as
$\bar C$ is square, it is invertible, and so is $C$. So $C$ is a frame of $M$ of
type $(R,S)$. Next, $\theta C_{o_1}=N(z+x)=C_{o_1}$,
$\theta C_{o_2}=N(z-x)=-C_{o_2}$ and
$\theta\iota C_1=\iota\theta_1C_1=\iota C_1J_{R_1,S_1}$; so
$\theta C=CJ_{R,S}$. Finally, put $n=(2N)^{-1}$ and $n'=c(2N)^{-1}$. Then
$Cx_e=C_{o_1}+C_{o_2}=2Nx$, so $x=nCx_e$; $Cy_e=2Nz$, so $y=cz=n'Cy_e$; and
$CP_e=\iota C_1$ by the choice of the columns $C_{e(b)}$.
\end{proof}
% Lean: exists_frame (EmbeddingTransposeVectorZero.lean, about 110 lines), with hf for the form,
% herm, iso, k0 (the positivity field of IsCartan), zk, NN, oz, col, the columns cl, gram, the
% right inverse from Matrix.mul_eq_one_comm_of_card_eq, and the identities for theta C, x, y
% and C P.

\begin{lemma}\label{asm-theta-agree-local-mult-archimedean-torus}
Let $B$ be a theta base of a CM field $E$, let $\chi$ be a homomorphism from
$\mathbf A_E^\times$ to the complex numbers of absolute value $1$ with exponents
$m$ (Lemma~\ref{asm-theta-agree-real-lie}), let $r\in\mathbf R$, and let
$l_w\in\mathbf C$ for every infinite place $w$ of $F$. Then there are
$t(s)\in\mathbf A_E^\times$, for $s\in\mathbf R$, with $t(s)-1$ at $\infty$
(Definition~\ref{asm-theta-agree-models}(2)) and
$\pi_w(t(s))=\exp(sl_w)$ for all $s$ and $w$, such that the function
$s\mapsto\chi(t(s))|t(s)|^r$ is differentiable at $s=0$.
\end{lemma}
\leanref{Development/WeilBound/ThetaLifts/LocalTheory/ToralCohomology/ExistsTorus.lean}{PositivityTheory.ThKG.CayleyTransformDarboux.exists_torus}
\leanref{Development/WeilBound/ThetaLifts/LocalTheory/ToralCohomology/ExistsTorus.lean}{PositivityTheory.ThKG.CayleyTransformDarboux.exists_unit}
\leanref{Development/WeilBound/ThetaLifts/LocalTheory/ToralCohomology/ExistsTorus.lean}{PositivityTheory.ThKG.CayleyTransformDarboux.exists_inf}
\leanref{Development/WeilBound/ThetaLifts/LocalTheory/ToralCohomology/ExistsTorus.lean}{PositivityTheory.ThKG.CayleyTransformDarboux.eq_of_inf}
\leanref{Development/WeilBound/ThetaLifts/LocalTheory/ToralCohomology/ExistsTorus.lean}{PositivityTheory.ThKG.CayleyTransformDarboux.absE_inf}
\begin{proof}
Fix $s$. By fact (iii) of the proof of Lemma~\ref{asm-theta-agree-real-lines},
there are $u,u'\in\mathbf A_E$ at $\infty$ with $\pi_w(u)=\exp(sl_w)-1$ and
$\pi_w(u')=\exp(-sl_w)-1$ for every $w$. Then $(1+u)(1+u')-1$ is at $\infty$ and
all its $\pi_w$-components vanish, so $(1+u)(1+u')=1$ by the same fact; put
$t(s)=1+u$. The finite components of $t(s)$ are $1$, so
$|t(s)|=\prod_w|\pi_w(t(s))|^2$, as at an infinite place of $F$ the absolute
value of the norm from $\mathbf C$ to $\mathbf R$ is the square of the usual
absolute value; and
$\chi(t(s))=\prod_w(\pi_w(t(s))/|\pi_w(t(s))|)^{m(w)}$. As
$\pi_w(t(s))=\exp(sl_w)$,
\[
\chi(t(s))|t(s)|^r=\exp\Bigl(s\sum_w\bigl(i\,m(w)\operatorname{Im}l_w+2r\operatorname{Re}l_w\bigr)\Bigr),
\]
which is differentiable at $s=0$.
\end{proof}
% Lean: exists_torus (ExistsTorus.lean, about 20 lines), with exists_unit, exists_inf (from
% exists_moved and exists_re, as in fact (iii)), eq_of_inf (eq_zero_of_piInf), absE_inf
% (absEw_eq_sq) and the exponent law HasExp; ExistsTorus.lean has about 120 lines.

\begin{lemma}\label{asm-theta-agree-local-mult-archimedean-levi}
Let $B$ be a theta base of a CM field $E$, let $W$ and $W_1$ be lattice spaces
of ranks $d$ and $d_1$, and let $H'=(x,y,\iota,\mathfrak b)$ be an adapted plane
in $W$ over $W_1$ at some positions
(Definition~\ref{asm-theta-agree-planes}(1)). Let $(e,o_1,o_2,\sigma)$ be an
embedding of $\mathrm{Ix}(R_1,S_1)$ into $\mathrm{Ix}(R,S)$, with matrix $P_e$
and vectors $x_e$ and $y_e$
(Definition~\ref{asm-theta-agree-local-mult-archimedean-embeddings}). For every
infinite place $w$ of $F$ let $C_w\in M_{d\times\mathrm{Ix}(R,S)}(\mathbf C)$,
let $C_{1,w}\in M_{d_1\times\mathrm{Ix}(R_1,S_1)}(\mathbf C)$ with
$C_{1,w}A_w=1$ for some matrix $A_w$, and let $\xi_w,\xi'_w\in\mathbf C$ with
$\xi_w\neq0$, such that $\tau_w(x)=\xi_wC_wx_e$, $\tau_w(y)=\xi'_wC_wy_e$ and
$C_wP_e=\tau_w(\iota)C_{1,w}$.
\begin{enumerate}
\item Let $\gamma$ and $\gamma_1$ be lines in $U_W$ and $U_{W_1}$
  (Definition~\ref{asm-theta-agree-matched}(1)) generated by $(Y_w)$ and
  $(Y_{1,w})$ relative to $(C_w)$ and $(C_{1,w})$
  (Definition~\ref{asm-theta-agree-real-fock}(6)), and let $l_w\in\mathbf C$
  with $Y_wx_e=l_wx_e$, $Y_wy_e=-\bar l_wy_e$ and $Y_wP_e=P_eY_{1,w}$ for every
  $w$. Let $s\in\mathbf R$ and $t\in\mathbf A_E^\times$ with $t-1$ at $\infty$
  and $\pi_w(t)=\exp(sl_w)$ for every $w$. Then $\gamma(s)$ is a Levi element
  for $(t,\gamma_1(s))$ (Definition~\ref{asm-theta-agree-planes}(2)).
\item Let $\gamma$ be a line in $U_W$ generated by $(Y_w)$ relative to $(C_w)$,
  with $Y_wx_e=0$ and $Y_wP_eu\in\mathbf Cx_e$ for every $w$ and every
  $u\in\mathbf C^{\mathrm{Ix}(R_1,S_1)}$. Then $\gamma(s)\in N_{H'}$ for every
  $s\in\mathbf R$.
\end{enumerate}
\end{lemma}
\leanref{Development/WeilBound/ThetaLifts/LocalTheory/ToralCohomology/ArchimedeanMultiplicityOneOf.lean}{PositivityTheory.ThKG.CayleyTransformDarboux.FD}
\leanref{Development/WeilBound/ThetaLifts/LocalTheory/ToralCohomology/ArchimedeanMultiplicityOneOf.lean}{PositivityTheory.ThKG.CayleyTransformDarboux.FD.levi}
\leanref{Development/WeilBound/ThetaLifts/LocalTheory/ToralCohomology/ArchimedeanMultiplicityOneOf.lean}{PositivityTheory.ThKG.CayleyTransformDarboux.FD.unip}
\leanref{Development/WeilBound/ThetaLifts/LocalTheory/ToralCohomology/ArchimedeanMultiplicityOneOf.lean}{PositivityTheory.ThKG.CayleyTransformDarboux.eig_vec}
\begin{proof}
Elements of $\mathbf A_E$, vectors and matrices over it that are at $\infty$ and
have all $\pi_w$-components $0$ are $0$ (fact (iii) of the proof of
Lemma~\ref{asm-theta-agree-real-lines}); and $\gamma(s)-1$, $\gamma_1(s)-1$ and
$t-1$ are at $\infty$. For every $w$ and $z\in\mathbf A_E^{d_1}$, put
$u_w=A_w\pi_w(z)$, so that $\pi_w(z)=C_{1,w}u_w$ and
$\tau_w(\iota)\pi_w(z)=C_wP_eu_w$.

(1) The vector $\gamma(s)x-tx=(\gamma(s)-1)x-(t-1)x$ is at $\infty$, and
\[
\pi_w(\gamma(s))\tau_w(x)=\xi_wC_w\exp(sY_w)x_e=\xi_we^{sl_w}C_wx_e=\pi_w(t)\tau_w(x),
\]
so $\gamma(s)x=tx$. In the same way $\gamma(s)y=c(t)^{-1}y$: $c(t)^{-1}-1$ is
at $\infty$, and $\pi_w(c(t)^{-1})=\overline{\pi_w(t)}^{-1}=\exp(-s\bar l_w)$, as
$E$ is a CM field. For $z\in\mathbf A_E^{d_1}$, the vector
$\gamma(s)\iota z-\iota\gamma_1(s)z$ is at $\infty$, and as
$\exp(sY_w)P_e=P_e\exp(sY_{1,w})$,
\[
\pi_w(\gamma(s))\tau_w(\iota)\pi_w(z)=C_w\exp(sY_w)P_eu_w=C_wP_e\exp(sY_{1,w})u_w
=\tau_w(\iota)C_{1,w}\exp(sY_{1,w})u_w=\tau_w(\iota)\pi_w(\gamma_1(s))\pi_w(z).
\]
So $\gamma(s)\iota=\iota\gamma_1(s)$, and $\gamma(s)$ is a Levi element for
$(t,\gamma_1(s))$.

(2) As in (1) with $l_w=0$ and $t=1$, $\gamma(s)x=x$. Let
$z\in\mathbf A_E^{d_1}$, and write $Y_wP_eu_w=a_wx_e$. By fact (iii) there is
$g\in\mathbf A_E$ at $\infty$ with $\pi_w(g)=sa_w/\xi_w$ for every $w$. As
$Y_w(Y_wP_eu_w)=a_wY_wx_e=0$, $\exp(sY_w)P_eu_w=P_eu_w+sa_wx_e$, so
\[
\pi_w(\gamma(s))\tau_w(\iota)\pi_w(z)=C_wP_eu_w+sa_wC_wx_e
=\tau_w(\iota)\pi_w(z)+\pi_w(g)\tau_w(x).
\]
As $\gamma(s)\iota z-\iota z-gx$ is at $\infty$, $\gamma(s)\iota z=\iota z+gx$.
So $\gamma(s)$ is one of the generators of $N_{H'}$.
\end{proof}
% Lean: FD (the hypotheses on C, C1, n = xi, n' = xi'), eig_vec, FD.iota, FD.levi and FD.unip
% (ArchimedeanMultiplicityOneOf.lean, about 95 lines), with exp_int, exp_eig and exp_nil
% (EmbeddingTransposeVectorZero.lean, about 40 lines), at_mulVec, at_cA, inv_at, piInf_cA,
% exists_inf and eq_of_inf. Neither FD.levi nor FD.unip uses the adaptedness of H'.

\begin{lemma}\label{asm-theta-agree-local-mult-archimedean}
Let $B$ be a theta base of a CM field $E$, let $Q$ and $Q_1$ be matched
quadruples of step type $(\chi,r)$, and let $H'$ be an adapted plane in $W$
over $W_1$ at some positions
(Definition~\ref{asm-theta-agree-local-equivariant}). Let
$\lambda,\lambda'\colon\Pi\to\Pi_1$ be equivariant at $\infty$ for
$(H',\chi,r)$ with $\lambda\neq0$. Then $\lambda'=c\lambda$ for some
$c\in\mathbf C$.
\end{lemma}
\leanref{Development/WeilBound/ThetaLifts/LocalTheory/UnramifiedDeterminedBySpherical.lean}{PositivityTheory.ThKG.Loc.ArchimedeanMultiplicityOne}
\leanref{Development/WeilBound/ThetaLifts/LocalTheory/ToralCohomology/ArchimedeanMultiplicityOneOfLie.lean}{PositivityTheory.ThKG.CayleyTransformDarboux.archimedean_multiplicity_one_of_lie}
\leanref{Development/WeilBound/ThetaLifts/LocalTheory/ToralCohomology/ArchimedeanMultiplicityOneOf.lean}{PositivityTheory.ThKG.CayleyTransformDarboux.archimedean_multiplicity_one_of}
\leanref{Development/WeilBound/ThetaLifts/LocalTheory/ToralCohomology/ArchimedeanMultiplicityOneOf.lean}{PositivityTheory.ThKG.CayleyTransformDarboux.core}
\begin{proof}
Let $(a,a')$ and $(a_1,a'_1)$ be the ranks of $Q$ and $Q_1$, $(\Pi,p,p')$ and
$(\Pi_1,p_1,p'_1)$ their real maps, $M$ and $M_1$ the models of their first
Weil pairs, with Weil representations $\omega$ and $\omega_1$, and $(R,S)$ and
$(R_1,S_1)$ the signatures of $W$ and $W_1$. Write $H'=(x,y,\iota,\mathfrak b)$.
For an infinite place $w$ and a complex number or matrix $Y$, let $\delta_wY$ be
the family with $Y$ at $w$ and $0$ at the other infinite places.

\emph{Step 1: the embedding.} By the step type, either $a_1=a$,
$a'_1+2=a'$, $\chi=\beta$ and $r=a/2$, and the second integer of $G_1$ is one
less than that of $G$ while the first is the same; or $a_1+2=a$, $a'_1=a'$,
$\chi=\alpha$ and $r=a'/2$, and the first integer of $G_1$ is one less than that
of $G$ while the second is the same. In both cases $R=R_1+1$ and $S=S_1+1$. In
the first case let $(e,o_1,o_2,\sigma)$ be the last-index embedding of
$\mathrm{Ix}(R_1,S_1)$ into $\mathrm{Ix}(R,S)$, and in the second the
first-index embedding
(Definition~\ref{asm-theta-agree-local-mult-archimedean-embeddings}(2)). Let
$\Lambda$ and $\Lambda_1$ be the weights of
Lemma~\ref{asm-theta-agree-local-mult-archimedean-action} for $Q$ and $Q_1$.
Then $\Lambda_w(e(b))=\Lambda_{1,w}(b)$ for all $w$ and $b$. Indeed $\kappa_w$
and $m_\beta(w)$ are those of $B$ in both, $c_{R,S}=c_{R_1,S_1}$, and with
$(n_+,n_-)$ and $(n_{1,+},n_{1,-})$ the signatures of $V$ and $V_1$: in the first
case $(n_+,n_-)=(n_{1,+},n_{1,-})$ and $e$ keeps the position of each index in
its block; in the second $(n_+,n_-)=(n_{1,+}+1,n_{1,-}+1)$ and $e$ moves each
index one position on. In both cases $c_{n_+,n_-}=c_{n_{1,+},n_{1,-}}$, and
$e(b)$ is among the first $n_+$ (or $n_-$) elements of its block if and only if
$b$ is among the first $n_{1,+}$ (or $n_{1,-}$) elements of its block.

\emph{Step 2: frames adapted to $H'$.} By
Lemma~\ref{asm-theta-agree-real-frames}, for every infinite place $w$ there is a
frame $C_{1,w}$ of $i\tau_w(J_{W_1})$ of type $(R_1,S_1)$ with
$\theta_{W_1,w}C_{1,w}=C_{1,w}J_{R_1,S_1}$. Fix $w$. As $E$ is a CM field,
$\tau_w(c(z))=\overline{\tau_w(z)}$; so $i\tau_w(J_W)$ is Hermitian, as $J_W$ is
skew-Hermitian, and
$\tau_w(u)^{\mathsf T}i\tau_w(J_W)\overline{\tau_w(v)}=i\tau_w(h(u,v))$ for
$u,v\in E^{a+a'}$, with $h$ as in Definition~\ref{asm-theta-agree-planes}. By
Definition~\ref{asm-theta-agree-planes}(1): $\tau_w(x)\neq0$, as $h(x,y)=1$;
$h(x,x)=0$; $h(\iota u,x)=h(\iota u,y)=0$ for every $u$;
$\iota^{\mathsf T}J_Wc(\iota)=J_{W_1}$; $\theta_{W,w}\tau_w(\iota)=\tau_w(\iota)\theta_{W_1,w}$;
and $\tau_w(y)=\lambda_w\theta_{W,w}\tau_w(x)$ with $\lambda_w\neq0$. Also
$a+a'=R+S$, and $\theta_{W,w}$ is a Cartan involution of $i\tau_w(J_W)$. So
Lemma~\ref{asm-theta-agree-local-mult-archimedean-frame} applies to
$i\tau_w(J_W)$, $\theta_{W,w}$, $i\tau_w(J_{W_1})$, $\theta_{W_1,w}$,
$C_{1,w}$, $\tau_w(x)$, $\tau_w(y)$, $\lambda_w$ and $\tau_w(\iota)$, as $e$
maps each block into the block of the same name. It gives a frame $C_w$ of
$i\tau_w(J_W)$ of type $(R,S)$ and $\xi_w,\xi'_w\in\mathbf C$ with
$\xi_w\neq0$, $\theta_{W,w}C_w=C_wJ_{R,S}$, $\tau_w(x)=\xi_wC_wx_e$,
$\tau_w(y)=\xi'_wC_wy_e$ and $C_wP_e=\tau_w(\iota)C_{1,w}$. As frames have right
inverses, Lemma~\ref{asm-theta-agree-local-mult-archimedean-levi} applies to
these data.

\emph{Step 3: the actions.}
Lemma~\ref{asm-theta-agree-local-mult-archimedean-action}, for the triple of
$Q$ with the frames $C_w$ and for that of $Q_1$ with the frames $C_{1,w}$, gives
structures $\varrho$ and $\varrho_1$ on $\Pi$ and $\Pi_1$ of Lie modules over
the set of infinite places of types $(R,S)$ and $(R_1,S_1)$, and vectors $f_0$
and $f_{1,0}$, such that $(\Pi,f_0)$ and $(\Pi_1,f_{1,0})$ are highest-weight
modules of weights $\Lambda$ and $\Lambda_1$ and the compatibility of that lemma
with lines holds.

\emph{Step 4: tori.} As $\mathcal X$ is proper, $\chi$ has exponents $m_\beta$
or $m_\alpha$ (Definition~\ref{asm-theta-agree-triple}(3)). For a family
$l=(l_w)_w$ of complex numbers, let $t_l(s)$ be the elements given by
Lemma~\ref{asm-theta-agree-local-mult-archimedean-torus} for $\chi$, these
exponents, $r$ and $l$, and $c(l)$ the derivative at $s=0$ of
$s\mapsto\chi(t_l(s))|t_l(s)|^r$. For $l=0$, $t_0(s)-1$ is at $\infty$ with all
$\pi_w$-components $0$, so $t_0(s)=1$ by fact (iii) of the proof of
Lemma~\ref{asm-theta-agree-real-lines}, and $c(0)=0$.

\emph{Step 5: three relations.} Let $D\colon\Pi\to\Pi_1$ be equivariant at
$\infty$ for $(H',\chi,r)$. By condition (a) of
Definition~\ref{asm-theta-agree-matched}(2), every $v\in\Pi$ is $p(g)$ for some
$g\in\operatorname{Pol}(M)$, and every element of $\Pi_1$ is $p_1(g_1)$ for
some $g_1\in\operatorname{Pol}(M_1)$.
\begin{enumerate}
\item[(L)] Let $Y_w\in\mathfrak u(R,S)$, $Y_{1,w}\in\mathfrak u(R_1,S_1)$ and
  $l_w\in\mathbf C$ for every $w$, with $Y_wx_e=l_wx_e$, $Y_wy_e=-\bar l_wy_e$
  and $Y_wP_e=P_eY_{1,w}$. Then $D\varrho(Y)v=c(l)Dv+\varrho_1(Y_1)Dv$ for every
  $v\in\Pi$. Indeed, by Lemma~\ref{asm-theta-agree-real-lines}(2) there are
  lines $\gamma$ in $U_W$ and $\gamma_1$ in $U_{W_1}$ generated by $(Y_w)$ and
  $(Y_{1,w})$ relative to $(C_w)$ and $(C_{1,w})$, and by
  Lemma~\ref{asm-theta-agree-local-mult-archimedean-levi}(1), $\gamma(s)$ is a
  Levi element for $(t_l(s),\gamma_1(s))$ for every $s$. Write $v=p(g)$ and
  $Dv=p_1(g_1)$. By Lemma~\ref{asm-theta-agree-local-mult-archimedean-action},
  $t\mapsto\omega(1,\gamma(t))$ and $t\mapsto\omega_1(1,\gamma_1(t))$ act at
  $\infty$ to first order by $g\mapsto g'$ and $g_1\mapsto g'_1$ with
  $p(g')=\varrho(Y)v$ and $p_1(g'_1)=\varrho_1(Y_1)Dv$. Condition (b) of
  Definition~\ref{asm-theta-agree-local-equivariant}(3) gives
  $D(p(g'))=c(l)p_1(g_1)+p_1(g'_1)$.
\item[(U)] Let $w$ be an infinite place and $Y\in\mathfrak u(R,S)$ with
  $Yx_e=0$ and $YP_eu\in\mathbf Cx_e$ for every $u$. Then
  $D\varrho(\delta_wY)v=0$ for every $v$. Indeed, $\delta_wY$ generates a line
  $\gamma$ in $U_W$ relative to $(C_w)$
  (Lemma~\ref{asm-theta-agree-real-lines}(2)), with values in $N_{H'}$ by
  Lemma~\ref{asm-theta-agree-local-mult-archimedean-levi}(2). With $v=p(g)$ and
  $g'$ as in Lemma~\ref{asm-theta-agree-local-mult-archimedean-action},
  condition (a) of Definition~\ref{asm-theta-agree-local-equivariant}(3) gives
  $D\varrho(\delta_wY)v=D(p(g'))=0$.
\item[(T)] Let $w$ be an infinite place, $Y\in\mathfrak u(R,S)$ and
  $l\in\mathbf C$ with $Yx_e=lx_e$, $Yy_e=-\bar ly_e$ and $YP_e=0$. Then
  $D\varrho(\delta_wY)v=c(\delta_wl)Dv$ for every $v$: apply (L) to $\delta_wY$,
  the zero family and $\delta_wl$.
\end{enumerate}

\emph{Step 6: the parabolic relations.} Every $D$ as in Step 5 satisfies the
parabolic relations for $(e,c_*,c'_*)$, with $c_{*,w}=-i\,c(\delta_wi)$ and
$c'_{*,w}=c(\delta_w1)$
(Definition~\ref{asm-theta-agree-local-mult-archimedean-relations}). Write
$\epsilon(a)$ and $E_{ab}$ as in Definition~\ref{asm-theta-agree-real-fock}.
\begin{enumerate}
\item Let $Y\in\mathfrak u(R_1,S_1)$. As $e$ maps each block into the block of
  the same name, $J_{R,S}P_e=P_eJ_{R_1,S_1}$, and as $P_e$ is real,
  $P_eYP_e^{\mathsf T}\in\mathfrak u(R,S)$. As $o_1$ and $o_2$ are not in the
  image of $e$, $P_e^{\mathsf T}x_e=P_e^{\mathsf T}y_e=0$, and
  $P_e^{\mathsf T}P_e=1$. So (L), applied to $\delta_w(P_eYP_e^{\mathsf T})$,
  $\delta_wY$ and $l=0$, gives
  $D\varrho(\delta_w(P_eYP_e^{\mathsf T}))v=\varrho_1(\delta_wY)Dv$, as
  $c(0)=0$. Both sides are linear in $Y$, and every matrix is $Y'+iY''$ with
  $Y',Y''\in\mathfrak u(R_1,S_1)$ (proof of
  Lemma~\ref{asm-theta-agree-real-tensor}); so this holds for every $Y$. For
  $Y=E_{ab}$, $P_eE_{ab}P_e^{\mathsf T}=E_{e(a)e(b)}$, which gives relation (1).
\item The matrix $Y=i(E_{o_1o_1}+E_{o_2o_2})$ lies in $\mathfrak u(R,S)$, with
  $Yx_e=ix_e$, $Yy_e=iy_e$ and $YP_e=0$. By (T) with $l=i$,
  $iD\varrho(E^w_{o_1o_1}+E^w_{o_2o_2})v=c(\delta_wi)Dv$.
\item The matrix $Y=E_{o_1o_2}+E_{o_2o_1}$ lies in $\mathfrak u(R,S)$, with
  $Yx_e=x_e$, $Yy_e=-y_e$ and $YP_e=0$; apply (T) with $l=1$.
\item The matrix $Y=i(E_{o_1o_1}-E_{o_1o_2}+E_{o_2o_1}-E_{o_2o_2})$ lies in
  $\mathfrak u(R,S)$, with $Yx_e=0$ and $YP_e=0$; apply (U).
\item[(5), (6)] Let $b\in\mathrm{Ix}(R_1,S_1)$, $b'=e(b)$ and
  $\epsilon'=\epsilon(b')$, and for $s\in\mathbf C$ put
  $Y_s=s(E_{b'o_1}-E_{b'o_2})-\bar s\epsilon'(E_{o_1b'}+E_{o_2b'})$. Then
  $Y_s\in\mathfrak u(R,S)$, $Y_sx_e=0$, and $Y_sP_eu=-\bar s\epsilon'u_bx_e$
  for every $u$. By (U), $sA-\bar s\epsilon'B=0$, where
  $A=D\varrho(E^w_{b'o_1}-E^w_{b'o_2})v$ and
  $B=D\varrho(E^w_{o_1b'}+E^w_{o_2b'})v$. For $s=1$ and $s=i$ this gives
  $A-\epsilon'B=0$ and $A+\epsilon'B=0$, so $A=B=0$.
\end{enumerate}

\emph{Step 7: conclusion.} By Step 6, $\lambda$ and $\lambda'$ satisfy the
parabolic relations for $(e,c_*,c'_*)$, and $\lambda\neq0$. By Steps 1 and 3,
Lemma~\ref{asm-theta-agree-local-mult-archimedean-proportional} applies to
$(\Pi,f_0)$, $(\Pi_1,f_{1,0})$, $\lambda$ and $\lambda'$, and gives
$\lambda'=c\lambda$ for some $c\in\mathbf C$.
\end{proof}
% Lean: archimedean_multiplicity_one_of_lie LieFirst LieSecond
% (ToralCohomology/ArchimedeanMultiplicityOneOfLie.lean) = archimedean_multiplicity_one_of
% applied to the bundle Dic of RL.actsLaws, fockEx, polRange_of, LieFirst, LieSecond, lineEx,
% frameEx and even_im_hermitian_sub (Lemmas asm-theta-agree-real-weak, -fock-map, -lie, -lines,
% -frames and -parity). archimedean_multiplicity_one_of (ArchimedeanMultiplicityOneOf.lean,
% about 10 lines) is the case split of Step 1 on StepTy, with core applied to embB and the
% exponents of beta (beta-step) or embA and those of alpha (alpha-step), and the check of lamW
% through the embedding (the hypothesis named hLambda). core (about 180 lines):
% C1 from hD.frame and C from exists_frame with fm (Step 2), fd : FD; exists_mod twice (Step
% 3); exists_torus and c0 (Step 4); onto and onto1; par, with L, U, T, xi, N and the six fields
% of IsPar (Steps 5 and 6; key with upq_split and m.conj_single for (1)); par_line (Step 7).
% The ToralCohomology files other than DarbouxExt.lean have about 1,700 lines.

\subsection{Theta classes and reduction of rank}\label{subsection-asm-theta-classes}

\begin{definition}\label{asm-theta-agree-kills-class}
Let $(X,W,M,\chi_X,\chi_W,\omega)$ be a Weil pair for a theta base $B$ of a CM
field $E$ (Definition~\ref{asm-theta-agree-triple}(1)), with $X$ and $W$ of
ranks $n$ and $d$ and theta integral $I$. A \emph{theta class} for the pair is
a space $\mathcal C$ of Schwartz-Bruhat functions on $\mathbf A^{n\times d}$
with the following properties.
\begin{enumerate}
\item Every $\Phi\in\mathcal C$ converges
  (Definition~\ref{asm-theta-agree-models}(7)).
\item The function $I(\Phi)$ on $U_W$ is continuous for every
  $\Phi\in\mathcal C$.
\item Let $\xi$ be a Schwartz-Bruhat function on $\mathbf A_f^{n\times d}$ such
  that $g\otimes\xi\in\mathcal C$ for every $g\in\mathcal S(F_\infty^{n\times d})$.
  Then for every $h\in U_W$ the map $g\mapsto I(g\otimes\xi)(h)$ is continuous
  on $\mathcal S(F_\infty^{n\times d})$.
\item Let $\mu$ be a Haar measure on $U_W$ and $\mathcal D$ a fundamental domain
  for $U_W(F)$ in $U_W$ with respect to $\mu$. Then $|I(\Phi)|^2$ is integrable
  on $\mathcal D$ for every $\Phi\in\mathcal C\cap\operatorname{PD}(M)$
  (Definition~\ref{asm-theta-agree-models}(4)).
\end{enumerate}
\end{definition}
\leanref{Development/WeilBound/ThetaLifts/CartanFactorisation/ExistsBeta.lean}{PositivityTheory.ThKF.Cls}

\begin{lemma}\label{asm-theta-agree-kills-class-triple}
Let $(B,G,\mathcal X,\omega,\omega')$ be a triple of ranks $(a,a')$ for a theta
base $B$ of a CM field $E$ (Definition~\ref{asm-theta-agree-triple}), with
models $M$ and $M'$, and put $I_0=\{1,\ldots,a\}\times\{1,\ldots,a+a'\}$. Then
the Schwartz-Bruhat functions on $\mathbf A^{I_0}$ form a theta class for the
Weil pair $(V,W,M,\beta,\beta/\alpha,\omega)$, and $\mathcal S^\circ(M')$ is a
theta class for the Weil pair $(V',W,M',\alpha,\alpha/\beta,\omega')$
(Definitions~\ref{asm-theta-agree-models}(5)
and~\ref{asm-theta-agree-kills-class}).
\end{lemma}
\leanref{Development/WeilBound/ThetaLifts/CartanFactorisation/Kills0Of.lean}{PositivityTheory.ThKF.clsV}
\leanref{Development/WeilBound/ThetaLifts/CartanFactorisation/Kills0Of.lean}{PositivityTheory.ThKF.clsV'}
\leanref{Development/WeilBound/ThetaLifts/CartanFactorisation/Kills0Of.lean}{PositivityTheory.ThKF.so_leK}
\begin{proof}
By Definition~\ref{asm-theta-agree-triple}(4), these are Weil pairs, of
signatures $(a-p,p;a+a'-p-q,p+q)$ and $(q,a'-q;a+a'-p-q,p+q)$.

\emph{The first pair.} As $G$ is proper,
$2(p+q)+1<\lfloor(a-a')/2\rfloor$, so $4p<a-a'$ and
$a+a'+2\min(a-p,p)\leq a+a'+2p<2a$, as in the proof of
Lemma~\ref{asm-theta-agree-decay}. So
Lemma~\ref{asm-theta-agree-convergence} applies, with $n=a$ and $d=a+a'$. Its
part (1) gives condition (1), and its part (2) gives condition (3).
Lemma~\ref{asm-theta-agree-definite-continuous} gives condition (2), and
Lemma~\ref{asm-theta-agree-end-square} gives condition (4).

\emph{The second pair.} Every element of $\mathcal S^\circ(M')$ is a
Schwartz-Bruhat function: for a spanning element $f\otimes\bigotimes_v\phi_v$,
the function $\bigotimes_v\phi_v$ is a Schwartz-Bruhat function on the finite
adeles, as $M'$ is regular and so $\mathrm{vac}'_v$ is the indicator function
of the integral vectors for almost all $v$. By
Lemma~\ref{asm-theta-agree-decay}, every element of $\mathcal S^\circ(M')$
converges for $(M',\omega')$. For condition (2), let
$\Phi'\in\mathcal S^\circ(M')$. The argument for the continuity of $I'(\Phi')$
in the proof of Lemma~\ref{asm-theta-agree-end-automorphic} uses only that
$\Phi'\in\mathcal S^\circ(M')$, so $I'(\Phi')$ is continuous. For condition (3), let $\xi$ be as
there, $h\in U_W$, and $u_k\to g$ a convergent sequence in
$\mathcal S(F_\infty^{I'_0})$, with $I'_0=\{1,\ldots,a'\}\times\{1,\ldots,a+a'\}$.
Lemma~\ref{asm-theta-agree-continuity}, with $m=1$, $\phi_1=\xi$ and
$u_{k,1}=u_k$, gives $I'(u_k\otimes\xi)(h)\to I'(g\otimes\xi)(h)$. So the map
is sequentially continuous, hence continuous, as the Schwartz space is
metrizable. Condition (4) is Lemma~\ref{asm-theta-agree-end-square}.
\end{proof}
% Lean: clsV WeilConv ICont SqInt (with weilCondV) and clsV' KDecay KCoinv SoCont ICont SqInt
% (Kills0Of.lean); so_leK is the second paragraph's first sentence, and contf of clsV' is
% SoCont with SeqContinuous.continuous. The field cont of clsV' is the second clause of ICont,
% which iCont_of proves without using its hypothesis on So (soFin, dom_of_domSoU, dom_sum,
% dom_continuous), as the continuity paragraph of Lemma asm-theta-agree-end-automorphic does.

\begin{lemma}\label{asm-theta-agree-kills-forms}
Let $(X,W,M,\chi_X,\chi_W,\omega)$ be a Weil pair for a theta base $B$ of a CM
field $E$ (Definition~\ref{asm-theta-agree-triple}(1)), let $\mu$ be a Haar
measure on $U_W$, and let $\mathcal D$ be a fundamental domain for $U_W(F)$ in
$U_W$ with respect to $\mu$. Let $Q$ be a complex vector space and $u$ a linear
map from $Q$ to the functions on $U_W$ such that for every $x\in Q$ the
function $u(x)$ is continuous, $u(x)(\gamma h)=u(x)(h)$ for all
$\gamma\in U_W(F)$ and $h\in U_W$, and $|u(x)|^2$ is integrable on $\mathcal D$.
For $x,y\in Q$ put
\[
\beta(x,y)=\int_{\mathcal D}\overline{u(x)(h)}\,u(y)(h)\,d\mu(h).
\]
\begin{enumerate}
\item The integral converges absolutely; $\beta$ is conjugate-linear in $x$ and
  linear in $y$; $\beta(y,x)=\overline{\beta(x,y)}$; and $\beta(x,x)\geq0$.
\item If $\beta(x,x)=0$, then $\beta(x,y)=0$ for every $y\in Q$, and $u(x)=0$.
\item Let $h_0\in U_W$, and let $A$ be a linear map of $Q$ with
  $u(Ax)(h)=u(x)(hh_0)$ for all $x\in Q$ and $h\in U_W$. Then
  $\beta(Ax,Ay)=\beta(x,y)$ for all $x,y\in Q$.
\item Let $\delta$ be a line in $U_W$
  (Definition~\ref{asm-theta-agree-matched}(1)), and $D$ a linear map of $Q$
  such that for all $x\in Q$ and $h\in U_W$ the function
  $t\mapsto u(x)(h\delta(t))$ has derivative $u(Dx)(h)$ at $t=0$. Then
  $\beta(Dx,y)+\beta(x,Dy)=0$ for all $x,y\in Q$.
\end{enumerate}
\end{lemma}
\leanref{Development/WeilBound/ThetaLifts/CartanFactorisation/ExistsBeta.lean}{PositivityTheory.ThKF.exists_beta}
\leanref{Development/WeilBound/ThetaLifts/CartanFactorisation/RadOfZero.lean}{PositivityTheory.ThKF.rad_of_zero}
\leanref{Development/WeilBound/ThetaLifts/CartanFactorisation/IntegrableOnRight.lean}{PositivityTheory.ThKF.L2.eq_zero}
\leanref{Development/WeilBound/ThetaLifts/CartanFactorisation/IntegrableOnRight.lean}{PositivityTheory.ThKF.L2.form_right}
\leanref{Development/WeilBound/ThetaLifts/CartanFactorisation/IntegrableOnRight.lean}{PositivityTheory.ThKF.L2.integral_deriv}
\leanref{Development/WeilBound/ThetaLifts/CartanFactorisation/IntegrableOnRight.lean}{PositivityTheory.ThKF.L2.skew}
\begin{proof}
(1) As $|\bar ab|\leq|a|^2+|b|^2$ and $u(x)$ and $u(y)$ are continuous, the
integrand is integrable on $\mathcal D$. The other claims follow from the
linearity of $u$ and of the integral, by conjugating the integrand, and from
$\beta(x,x)=\int_{\mathcal D}|u(x)|^2\,d\mu$.

(2) By (1), $\beta$ is a positive semidefinite Hermitian form on $Q$, so the
Cauchy-Schwarz inequality $|\beta(x,y)|^2\leq\beta(x,x)\beta(y,y)$ holds, and
$\beta(x,x)=0$ gives $\beta(x,y)=0$. As $\int_{\mathcal D}|u(x)|^2\,d\mu=0$,
the argument of the last paragraph of the proof of
Lemma~\ref{asm-theta-agree-end-perp-zero} applies to the continuous function
$u(x)$, which is invariant under left translation by $U_W(F)$, and gives
$u(x)=0$.

(3) The matrix of $W$ is nondegenerate $\epsilon$-Hermitian
(Lemma~\ref{asm-theta-agree-convergence-signature}), so by
Theorem~\ref{asm-theta-agree-definite-volume}, $U_W(F)$ is countable and $\mu$
is invariant under right translations. Let $G$ be a function on $U_W$,
integrable on $\mathcal D$, with $G(\gamma h)=G(h)$ for $\gamma\in U_W(F)$.
Then $\int_{\mathcal D}G(hh_0)\,d\mu(h)=\int_{\mathcal D}G\,d\mu$, by the
argument of the proof of Lemma~\ref{asm-theta-agree-anisotropic-invariance}:
right translation by $h_0$ carries $\mathcal D$ to a fundamental domain
$\mathcal Dh_0$ and preserves $\mu$, and the integral of $G$ is the same over
any two fundamental domains. Apply this to
$G(h)=\overline{u(x)(h)}u(y)(h)$, for which
$G(hh_0)=\overline{u(Ax)(h)}u(Ay)(h)$.

(4) Put $G(h)=\overline{u(x)(h)}\,u(y)(h)$ and
$G_1(h)=\overline{u(Dx)(h)}\,u(y)(h)+\overline{u(x)(h)}\,u(Dy)(h)$. Both are
continuous, invariant under left translation by $U_W(F)$, and integrable on
$\mathcal D$, by (1). By the product rule, $t\mapsto G(h\delta(t))$ has
derivative $G_1(h)$ at $t=0$; as $\delta$ is a homomorphism, it has derivative
$G_1(h\delta(s))$ at $t=s$, for every $s$. So
$G(h\delta(1))-G(h)=\int_0^1G_1(h\delta(s))\,ds$. The function
$(s,h)\mapsto G_1(h\delta(s))$ is continuous, and by the argument of (3)
$\int_{\mathcal D}|G_1(h\delta(s))|\,d\mu(h)=\int_{\mathcal D}|G_1|\,d\mu$ for
every $s$; so it is integrable on $[0,1]\times\mathcal D$. By Fubini's theorem
and the argument of (3), applied to $G_1$ and to $G$,
\[
\int_{\mathcal D}G_1\,d\mu=\int_0^1\!\!\int_{\mathcal D}G_1(h\delta(s))\,d\mu(h)\,ds
=\int_{\mathcal D}\bigl(G(h\delta(1))-G(h)\bigr)\,d\mu(h)=0,
\]
and the left side is $\beta(Dx,y)+\beta(x,Dy)$.
\end{proof}
% Lean: (1) exists_beta (with L2.integrableOn_mul); (2) rad_of_zero and L2.eq_zero; (3)
% L2.form_right = L2.integral_right with fd_right; (4) L2.skew = L2.integral_deriv. The
% right invariance of mu is uniW (Unimod, Theorem asm-theta-agree-definite-volume).

\begin{lemma}\label{asm-theta-agree-kills-finite}
Let $(Y,W,M_0,\chi_Y,\chi_W,\omega_0)$ be a Weil pair for a theta base $B$ of a
CM field $E$ (Definition~\ref{asm-theta-agree-triple}(1)), with $Y$ and
$W$ of ranks $n$ and $d$, vacua $\mathrm{vac}_u$ and theta integral $I$. Let
$\mathcal C$ be a theta class for the pair
(Definition~\ref{asm-theta-agree-kills-class}), let $v$ be a finite place, and
let $\mathcal S_v$ be the space of Schwartz-Bruhat functions on
$F_v^{n\times d}$. Suppose that one of the following holds, with
$\mathfrak a_v(\omega_0)$ and the theta subspaces at $v$ of
Definition~\ref{asm-theta-agree-rechoice-terms}(3):
\begin{enumerate}
\item[(i)] every theta subspace at $v$ of $\omega_0$ is
  $\mathfrak a_v(\omega_0)$ or $\mathcal S_v$;
\item[(ii)] for all local operators $(\omega_v,\rho_v)$ at $v$ for the pair
  (Definition~\ref{asm-theta-agree-finite-local-objects}(4)) there is
  $s\in\mathcal S_v$ such that $\mathcal S_v$ is spanned by
  $\mathfrak a_v(\omega_0)$ and the vectors $\omega_v(1,h)s$ with
  $h\in U_W^{(v)}$, and such that for every $x\in\mathcal S_v$ there are
  $h_0,\ldots,h_m\in U_W^{(v)}$ and $w_0,\ldots,w_m,c\in\mathbf C$ with
  $w_i\omega_v(1,h_i)s=s$ for every $i$ and
  $\sum_iw_i\omega_v(1,h_i)x-c\,s\in\mathfrak a_v(\omega_0)$.
\end{enumerate}
Let $N$ be a theta subspace at $v$ of $\omega_0$ other than $\mathcal S_v$. Let
$f\in\operatorname{Pol}(M_0)$, and let $\phi=(\phi_u)_u$ be a family of
Schwartz-Bruhat functions $\phi_u$ on $F_u^{n\times d}$, for the finite places
$u$, with $\phi_u=\mathrm{vac}_u$ for almost all $u$. Write
$f\otimes\phi=f\otimes\bigotimes_u\phi_u$, and for a function $x$ on
$F_v^{n\times d}$ let $\phi[v\mapsto x]$ be $\phi$ with $\phi_v$ replaced by
$x$. Suppose that $f\otimes\phi[v\mapsto x]\in\mathcal C$ for every
$x\in\mathcal S_v$. If $\phi_v\in N$, then $I(f\otimes\phi)=0$.
\end{lemma}
\leanref{Development/WeilBound/ThetaLifts/SphericalLattices/ThetaKillsCuspFormsSpherical.lean}{PositivityTheory.Sph.killFin3}
\leanref{Development/WeilBound/ThetaLifts/CartanFactorisation/UpdSelf.lean}{PositivityTheory.ThKF.intersection_update_fsHerm_neg_one}
\leanref{Development/WeilBound/ThetaLifts/SphericalLattices/ThetaSubSup.lean}{PositivityTheory.Sph.stab_of}
\begin{proof}
\emph{Step 1: the map $j$.} For a function $x$ on $F_v^{n\times d}$ put
$j(x)=f\otimes\phi[v\mapsto x]$. It is linear in $x$, and
$j(\phi_v)=f\otimes\phi$. Let $r(z)=f(z_\infty)\prod_{u\neq v}\phi_u(z_u)$; it
does not depend on the $v$-component, and $j(x)$ is the function
$z\mapsto x(z_v)r(z)$. For $x\in\mathcal S_v$, $j(x)$ is a Schwartz-Bruhat
function, as $M_0$ is regular, and it lies in
$\mathcal C\cap\operatorname{PD}(M_0)$. So, by
Definition~\ref{asm-theta-agree-models}(2), if an operator $T$ acts at $v$ by
$x\mapsto x_1$ with $x\in\mathcal S_v$, then $T(j(x))=j(x_1)$. As the functions
in $\mathcal C$ converge and theta functions are additive
(Lemma~\ref{asm-theta-agree-convergence-theta}(2)), $I$ is linear on
$\mathcal C$, and $x\mapsto I(j(x))$ is linear on $\mathcal S_v$.

\emph{Step 2: local operators.} Let $(\omega_v,\rho_v)$ be local operators at
$v$ for the pair (Lemma~\ref{asm-theta-agree-finite-inputs}(2)). For
$h\in U_W^{(v)}$ and $x\in\mathcal S_v$, $\omega_0(1,h)$ acts at $v$ by
$x\mapsto\omega_v(1,h)x$, which lies in $\mathcal S_v$. So
$j(\omega_v(1,h)x)=\omega_0(1,h)j(x)$, and, as $\omega_0$ is a representation,
$I(j(\omega_v(1,h)x))(g)=I(j(x))(gh)$ for $g\in U_W$. As $N$ is a theta
subspace at $v$, $\omega_v(1,h)N\subset N$.

\emph{Step 3: the augmentation.} $I(j(z))=0$ for every
$z\in\mathfrak a_v(\omega_0)$. By Step 1, it suffices to take $z=x_1-x$ with
$x,x_1\in\mathcal S_v$ and $\omega_0(g,1)$ acting at $v$ by $x\mapsto x_1$ for
some $g\in U_Y^{(v)}$. Then $j(x_1)=\omega_0(g,1)j(x)$, and
$I(j(x_1))=I(j(x))$ by Lemma~\ref{asm-theta-agree-anisotropic-invariance}.

\emph{Step 4: the two cases.} If (i) holds, then $N=\mathfrak a_v(\omega_0)$,
as $N\neq\mathcal S_v$; then $\phi_v\in\mathfrak a_v(\omega_0)$, and
$I(f\otimes\phi)=I(j(\phi_v))=0$ by Step 3. Suppose now that (ii) holds, and
let $s\in\mathcal S_v$ be as there, for the local operators of Step 2.

\emph{Step 5: the form.} The matrix of $W$ is nondegenerate
$\epsilon$-Hermitian (Lemma~\ref{asm-theta-agree-convergence-signature}), so by
Theorem~\ref{asm-theta-agree-definite-volume} there are a Haar measure $\mu$ on
$U_W$ and a fundamental domain $\mathcal D$ for $U_W(F)$ with respect to $\mu$.
For $x\in\mathcal S_v$ put $u(x)=I(j(x))$. By Step 1, $u$ is linear; $u(x)$ is
continuous, by condition (2) of a theta class; $u(x)(\gamma h)=u(x)(h)$ for
$\gamma\in U_W(F)$, by Lemma~\ref{asm-theta-agree-anisotropic-automorphy}, as
$j(x)$ is Schwartz-Bruhat; and $|u(x)|^2$ is integrable on $\mathcal D$, by
condition (4). Let $\beta$ be the form of
Lemma~\ref{asm-theta-agree-kills-forms} on $Q=\mathcal S_v$. By Step 2 and part
(3) of that lemma, $\beta(\omega_v(1,h)x,\omega_v(1,h)y)=\beta(x,y)$ for
$h\in U_W^{(v)}$. By Step 3, $u(z)=0$ for $z\in\mathfrak a_v(\omega_0)$, so
$\beta(x,y)=0$ whenever $x$ or $y$ lies in $\mathfrak a_v(\omega_0)$.

\emph{Step 6: $s\notin N$ and $s\neq0$.} If $s\in N$, then $N$ contains
$\mathfrak a_v(\omega_0)$ and every $\omega_v(1,h)s$, which span
$\mathcal S_v$; so $N=\mathcal S_v$, which is excluded. As $0\in N$, $s\neq0$.

\emph{Step 7: $\beta(s,y)=0$ for $y\in N$.} Let $y\in N$, and take
$h_0,\ldots,h_m\in U_W^{(v)}$ and $w_0,\ldots,w_m,c\in\mathbf C$ as in (ii),
for $x=y$. Put $z=\sum_iw_i\omega_v(1,h_i)y$. Then $z\in N$ by Step 2, and
$z-cs\in\mathfrak a_v(\omega_0)\subset N$; so $cs\in N$, and $c=0$ by Step 6.
So $z\in\mathfrak a_v(\omega_0)$, and $\beta(s,z)=0$ by Step 5. As
$s=w_i\omega_v(1,h_i)s$,
\[
\beta(s,\omega_v(1,h_i)y)=\beta\bigl(w_i\omega_v(1,h_i)s,\omega_v(1,h_i)y\bigr)=\bar w_i\,\beta(s,y),
\]
so $0=\beta(s,z)=\bigl(\sum_i|w_i|^2\bigr)\beta(s,y)$. As $s\neq0$ and
$w_0\omega_v(1,h_0)s=s$, $w_0\neq0$; so $\beta(s,y)=0$.

\emph{Step 8: $\beta(x,y)=0$ for $x\in\mathcal S_v$ and $y\in N$.} Let
$h\in U_W^{(v)}$. As $\omega_v(1,h)\omega_v(1,h^{-1})y=y$, Step 5 gives
$\beta(\omega_v(1,h)s,y)=\beta(s,\omega_v(1,h^{-1})y)$, which is $0$ by Step 7,
as $\omega_v(1,h^{-1})y\in N$. Also $\beta(z,y)=0$ for
$z\in\mathfrak a_v(\omega_0)$, by Step 5. As $\mathcal S_v$ is spanned by
$\mathfrak a_v(\omega_0)$ and the $\omega_v(1,h)s$, and $\beta$ is
conjugate-linear in the first variable, $\beta(x,y)=0$ for every
$x\in\mathcal S_v$.

\emph{Step 9: conclusion.} As $\phi_v\in N$, Step 8 gives
$\beta(\phi_v,\phi_v)=0$. By Lemma~\ref{asm-theta-agree-kills-forms}(2),
$u(\phi_v)=0$, that is, $I(f\otimes\phi)=0$.
\end{proof}
% Lean: killFin3 (ThetaKillsCuspFormsSpherical.lean). Step 1 is
% intersection_update_fsHerm_neg_one with pureAct and I_add'; Step 2 locRepEx, uA and stab_of
% (with LocCpt for the group U_W^(v)); Step 3 uaug (I_left); Step 4 the alternative
% hsa : IrrAt om v (hypothesis (i)) or SphAt M psi om v (hypothesis (ii), with IsSph: gen
% and line), which theta_kills_cusp_forms_spherical takes from SphTri (Theorem
% th-howe-finite); Step 5 uniW, exists_beta, iso (L2.form_right), bz and bz';
% Step 6 sN (Stab.not_mem with hs.gen) and s0; Step 7 kA (hs.line); Step 8 kB and kC; Step 9
% key with L2.eq_zero.

\begin{definition}\label{asm-theta-agree-kills-radical-maximal}
Let $(Y,W,M_0,\chi_Y,\chi_W,\omega_0)$ be a Weil pair for a theta base $B$ of a
CM field $E$ (Definition~\ref{asm-theta-agree-triple}(1)). The pair has
\emph{maximal radicals} if the following holds. Let $\beta$ be
a map from $\operatorname{Pol}(M_0)\times\operatorname{Pol}(M_0)$ to
$\mathbf C$, conjugate-linear in the first variable and linear in the second,
with $\beta(y,x)=\overline{\beta(x,y)}$, and let $R$ be the space of
$x\in\operatorname{Pol}(M_0)$ with $\beta(x,y)=0$ for every $y$. Suppose that:
\begin{enumerate}
\item[(a)] if $x,y\in\operatorname{Pol}(M_0)$, $\gamma$ is a line in $U_Y$
  (Definition~\ref{asm-theta-agree-matched}(1)), and
  $t\mapsto\omega_0(\gamma(t),1)$ acts at $\infty$ to first order by
  $x\mapsto y$ (Definition~\ref{asm-theta-agree-models}(2)), then $y\in R$;
\item[(b)] if $x,x_1,y,y_1\in\operatorname{Pol}(M_0)$, $\delta$ is a line in
  $U_W$, and $t\mapsto\omega_0(1,\delta(t))$ acts at $\infty$ to first order by
  $x\mapsto x_1$ and by $y\mapsto y_1$, then $\beta(x_1,y)+\beta(x,y_1)=0$.
\end{enumerate}
Let $P$ be a subspace of $\operatorname{Pol}(M_0)$ that contains $R$, such that
$y\in P$ whenever $x\in P$, $y\in\operatorname{Pol}(M_0)$, $\delta$ is a line in
$U_W$ and $t\mapsto\omega_0(1,\delta(t))$ acts at $\infty$ to first order by
$x\mapsto y$. Then $P\subset R$ or $P=\operatorname{Pol}(M_0)$.
\end{definition}
\leanref{Development/WeilBound/ThetaLifts/CartanFactorisation/UnitDef.lean}{PositivityTheory.ThKF.RadMaxAt}

\begin{lemma}\label{asm-theta-agree-kills-radical-greatest}
Let $R,S\geq0$ and $\Lambda\colon\mathrm{Ix}(R,S)\to\mathbf Q$, and let $T$ be a
Lie module of type $(R,S)$ generated with Casimir gap in weight $\Lambda$
(Definition~\ref{asm-theta-agree-real-quotient-lie-module}). Let $K$ be the sum
of the stable subspaces of $T$ other than $T$. Then $K$ is stable and
$K\neq T$. Let $\beta\colon T\times T\to\mathbf C$ be conjugate-linear in the
first variable and linear in the second, with
$\beta(T(E_{ab})x,y)=\epsilon_{ab}\beta(x,T(E_{ba})y)$ for all
$a,b\in\mathrm{Ix}(R,S)$ and $x,y\in T$, where $\epsilon_{ab}=1$ if $a$ and $b$
lie in the same block and $\epsilon_{ab}=-1$ otherwise. Then $\beta(x,y)=0$
for all $x\in K$ and $y\in T$.
\end{lemma}
\leanref{Development/WeilBound/ThetaLifts/CartanFactorisation/ExistsGreatestProperInvariantSubmoduleOfSide.lean}{PositivityTheory.ThKF.Unit.exists_greatest_proper_invariant_submodule_of_side}
\leanref{Development/WeilBound/ThetaLifts/CartanFactorisation/ExistsGreatestProperInvariantSubmoduleOfSide.lean}{PositivityTheory.ThKF.Unit.casimir_eigenvalue_and_sub_smul_onto_proper_submodule}
\begin{proof}
Let $(A,\rho,\iota,f)$ be a highest-weight datum as in
Definition~\ref{asm-theta-agree-real-quotient-lie-module}(3), let $B\subset A$
be the smallest subspace containing $f$ and stable under every $\rho(a,b)$, and
put $\Omega=\Omega_T$ and $c=c(\Lambda)$.

(i) As $T\neq0$ is generated by $\iota(f)$, $\iota(f)\neq0$, and
Lemma~\ref{asm-theta-agree-real-quotient-local} gives a subspace $Z$ with
$\iota(f)\notin Z$ that contains every stable subspace other than $T$. So
$K\subset Z$ and $K\neq T$; and $K$ is stable, as a sum of stable subspaces.

(ii) $\Omega\iota(f)=c\,\iota(f)$, and every $x\in K$ is $\Omega y-cy$ for some
$y\in T$. By steps (i) and (iv) of the proof of
Lemma~\ref{asm-theta-agree-real-quotient-local}, $B$ is the subspace called $B$
there, and $\sum\rho(a,b)\rho(b,a)=c$ on $B$. So
Lemma~\ref{asm-theta-agree-real-quotient-filtration} applies to $B$, with the
restrictions of the $\rho(a,b)$ and of $\iota$ and with $c_0=c$, and the stable
subspace generated by $\iota(B)$ is $T$. For $u\in B$,
$\Omega\iota(u)=\iota(\sum\rho(a,b)\rho(b,a)u)=c\,\iota(u)$. The $u\in B$ with
$\iota(u)\in K$ form a subspace stable under every $\rho(a,b)$; by step (iii) of
the proof of Lemma~\ref{asm-theta-agree-real-quotient-local} it is $0$ or
contains $f$, and $\iota(f)\notin K$; so $K\cap\iota(B)=0$. Let $x\in K$. By
part (2) of Lemma~\ref{asm-theta-agree-real-quotient-filtration}, $x$ lies in a
finite-dimensional subspace $F$ stable under the $T(E_{ab})$ with $a,b$ in the
same block. Then $C=\Omega-c$ maps $K\cap F$ into itself. If $y\in K\cap F$ and
$Cy=0$, then $y\in\iota(B)$ by part (1) of that lemma, so $y=0$. Hence $C$ is
injective on $K\cap F$, so surjective, and $x=Cy$ with $y\in K\cap F$.

(iii) For $a,b$ in the same block $\epsilon_{ab}=1$, so
$\beta(\Omega u,v)=\beta(u,\Omega v)$ for all $u,v$. Let $O$ be the set of
$z\in T$ with $\beta(x,z)=0$ for every $x\in K$. For $x\in K$, write
$x=\Omega y-cy$; as $c$ is real,
$\beta(x,\iota(f))=\beta(y,\Omega\iota(f)-c\,\iota(f))=0$. So $\iota(f)\in O$.
If $z\in O$, $x\in K$ and $a,b\in\mathrm{Ix}(R,S)$, then $T(E_{ba})x\in K$ and
$0=\beta(T(E_{ba})x,z)=\epsilon_{ba}\beta(x,T(E_{ab})z)$; so $O$ is stable
under every $T(E_{ab})$, hence under every $T(Y)$. As $T$ is generated by
$\iota(f)$, $O=T$.
\end{proof}

\begin{lemma}\label{asm-theta-agree-kills-radical-place}
Let $a,a',p,q,k$, $(P,Q)$, $(P',Q')$ and $(R,S)$ be as in
Lemma~\ref{asm-theta-agree-real-quotient}. Put $\mathcal P=\mathcal P(P,Q;R,S)$
and $c=\tfrac k2$, or $\mathcal P=\mathcal P(P',Q';R,S)$ and $c=-\tfrac k2$
(Definition~\ref{asm-theta-agree-real-fock}(1)). Then there is a subspace
$K\neq\mathcal P$ of $\mathcal P$, stable under every $d_2(Y)$, such that every
subspace $N\supset K$ stable under every $d_2(Y)$ is $K$ or $\mathcal P$, and
with the following property. Let $\beta\colon\mathcal P\times\mathcal P\to\mathbf C$
be conjugate-linear in the first variable and linear in the second, with
$\beta(g,y)=\beta(y,g)=0$ for every $y$ and every $g=d_1(X)x-c\operatorname{tr}(X)x$,
and with $\beta(d_2(E_{ab})x,y)=\epsilon_{ab}\beta(x,d_2(E_{ba})y)$ for all
$a,b\in\mathrm{Ix}(R,S)$, $\epsilon_{ab}$ as in
Lemma~\ref{asm-theta-agree-kills-radical-greatest}. Then $\beta(x,y)=0$ for all
$x\in K$ and $y\in\mathcal P$.
\end{lemma}
\leanref{Development/WeilBound/ThetaLifts/CartanFactorisation/ExistsGreatestProperInvariantSubmoduleOfSide.lean}{PositivityTheory.ThKF.Unit.One}
\leanref{Development/WeilBound/ThetaLifts/CartanFactorisation/ExistsGreatestProperInvariantSubmoduleOfSide.lean}{PositivityTheory.ThKF.Unit.one_of}
\leanref{Development/WeilBound/ThetaLifts/CartanFactorisation/ExistsGreatestProperInvariantSubmoduleOfSide.lean}{PositivityTheory.ThKF.Unit.one_data}
\begin{proof}
Let $\Theta$ be the quotient of $\mathcal P$ by the span of the polynomials
$g$, with quotient map $\pi$, and let $\vartheta(Y)$ be the operator induced by
$d_2(Y)+\tau_0\operatorname{tr}(Y)$, with $\tau_0=0$ in the first case and
$\tau_0=\tfrac k2$ in the second. These are the modules $\Theta$ and $\Theta'$
of Lemma~\ref{asm-theta-agree-local-mult-archimedean-sides}, which are
generated with Casimir gap. Let $K_\Theta$ be the subspace $K$ of
Lemma~\ref{asm-theta-agree-kills-radical-greatest} for $\Theta$, and
$K=\pi^{-1}(K_\Theta)$. As $\pi$ is onto, $K\neq\mathcal P$. As
$\pi(d_2(Y)x)=\vartheta(Y)\pi(x)-\tau_0\operatorname{tr}(Y)\pi(x)$, $K$ is
stable, and for a stable $N\supset K$ the space $\pi(N)$ is stable under the
$\vartheta(Y)$. If $\pi(N)\neq\Theta$, then $\pi(N)\subset K_\Theta$ and
$N=K$. If $\pi(N)=\Theta$, then every $x\in\mathcal P$ has $\pi(x)=\pi(y)$
with $y\in N$, and $x-y\in\ker\pi\subset K\subset N$; so $N=\mathcal P$.

Let $\beta$ be as in the statement. It vanishes when one of its arguments lies
in $\ker\pi$, so $\bar\beta(\pi(x),\pi(y))=\beta(x,y)$ defines a map on
$\Theta\times\Theta$, conjugate-linear in the first variable and linear in the
second. As $\tau_0$ is real and $\epsilon_{aa}=1$,
\[
\bar\beta(\vartheta(E_{ab})\pi(x),\pi(y))=\beta(d_2(E_{ab})x,y)+\delta_{ab}\tau_0\beta(x,y)
=\epsilon_{ab}\bar\beta(\pi(x),\vartheta(E_{ba})\pi(y)).
\]
By Lemma~\ref{asm-theta-agree-kills-radical-greatest}, $\bar\beta(u,v)=0$ for
$u\in K_\Theta$; so $\beta(x,y)=0$ for $x\in K$.
\end{proof}

\begin{lemma}\label{asm-theta-agree-kills-radical-tensor}
Let $\iota$ be a finite set, and for $i\in\iota$ let $V_i$ be a complex vector
space of countable dimension, $\mathcal O_i$ a set of operators on $V_i$, and
$K_i\neq V_i$ a subspace of $V_i$ stable under $\mathcal O_i$ such that every
subspace $N\supset K_i$ stable under $\mathcal O_i$ is $K_i$ or $V_i$. Let
$R\subset\bigotimes_iV_i$ be a subspace that contains every pure tensor
$\bigotimes_if_i$ with $f_i\in K_i$ for some $i$, and let $P\supset R$ be a
subspace stable under $L^{(i)}$ for every $i$ and every $L\in\mathcal O_i$
(Definition~\ref{asm-theta-agree-real-fock}(4)). Then $P\subset R$ or
$P=\bigotimes_iV_i$.
\end{lemma}
\leanref{Development/WeilBound/ThetaLifts/CartanFactorisation/TensMax.lean}{PositivityTheory.ThKF.Unit.tens_max}
\leanref{Development/WeilBound/ThetaLifts/CartanFactorisation/TensMax.lean}{PositivityTheory.ThKF.Unit.ker_le}
\leanref{Development/WeilBound/CoreStatements/GaussianSeminormBounds/RealSeminormFiltration/Tensor.lean}{PositivityTheory.KReal.tensorIrr}
\begin{proof}
Let $\pi_i\colon V_i\to V_i/K_i$ be the quotient maps and
$\pi=\bigotimes_i\pi_i$. First, $\ker\pi\subset R$. Choose linear maps $s_i$
with $\pi_is_i=1$. In a pure tensor $\bigotimes_if_i$, replace the factors
$f_i$ by $s_i\pi_i(f_i)$ one at a time; each step changes the tensor by a pure
tensor with a factor $f_i-s_i\pi_i(f_i)\in K_i$, which lies in $R$. So
$x-(\bigotimes_is_i)\pi(x)\in R$ for every $x$, and $x\in R$ if $\pi(x)=0$.

Each $V_i/K_i$ is nonzero, of countable dimension, and its only subspaces stable
under the operators induced by $\mathcal O_i$ are $0$ and itself. By
Lemma~\ref{th-first-lift-dixmier}, every operator on $V_i/K_i$ commuting with
them is a scalar, and the argument of the proof of
Lemma~\ref{asm-theta-agree-real-tensor} shows that $0$ and
$\bigotimes_i(V_i/K_i)$ are its only subspaces stable under the induced
operators $L^{(i)}$. The subspace $\pi(P)$ is stable under them. If
$\pi(P)=0$, then $P\subset\ker\pi\subset R$. Otherwise, every $x$ has
$\pi(x)=\pi(y)$ with $y\in P$, so $x-y\in\ker\pi\subset R\subset P$, and
$x\in P$.
\end{proof}

\begin{lemma}\label{asm-theta-agree-kills-radical}
Let $(B,G,\mathcal X,\omega,\omega')$ be a triple for a theta base $B$ of a CM
field $E$ (Definition~\ref{asm-theta-agree-triple}). Then its two Weil pairs
have maximal radicals
(Definition~\ref{asm-theta-agree-kills-radical-maximal}).
\end{lemma}
\leanref{Development/WeilBound/ThetaLifts/CartanFactorisation/UnitDef.lean}{PositivityTheory.ThKF.RadMax}
\leanref{Development/WeilBound/ThetaLifts/CartanFactorisation/FockSecondActionMatrixUnitEq.lean}{PositivityTheory.ThKF.Unit.radMax_of}
\begin{proof}
Let $(a,a')$ be the ranks of the triple and $p,q$ the integers of $G$, put
$(P,Q)=(a-p,p)$, $(P',Q')=(q,a'-q)$ and $(R,S)=(a+a'-p-q,p+q)$, and
$\kappa_w=m_\alpha(w)-m_\beta(w)$ for each infinite place $w$ of $F$. We treat
the first Weil pair $(V,W,M,\beta,\beta/\alpha,\omega)$ and put
$c_w=\tfrac12\kappa_w$ and $m(w)=m_\beta(w)$; for the second pair replace $V$,
$M$, $\omega$ and $(P,Q)$ by $V'$, $M'$, $\omega'$ and $(P',Q')$, and put
$c_w=-\tfrac12\kappa_w$ and $m(w)=m_\alpha(w)$. Let $I$ be the index set of the
pair, $J=J_{R,S}$, and $\epsilon_{ab}$ as in
Lemma~\ref{asm-theta-agree-kills-radical-greatest}.

\emph{Step 1: the Fock model.} As in the proof of
Lemma~\ref{asm-theta-agree-real}, there are frames $A_w$ of $\tau_w(J_V)$ and
$C_w$ of $i\tau_w(J_W)$, of types $(P,Q)$ and $(R,S)$, and a Fock map
$\Phi\colon\mathcal T\to\mathcal S(F_\infty^I)$ for $(M,(A_w),(C_w))$, with
$\mathcal T=\bigotimes_w\mathcal P(P,Q;R,S)$, which is injective with image
$\operatorname{Pol}(M)$; and for $T\in\mathcal T$:
\begin{enumerate}
\item[(i)] if $\gamma$ is a line in $U_V$ generated by $(Y_w)$ relative to
  $(A_w)$, then $t\mapsto\omega(\gamma(t),1)$ acts at $\infty$ to first order
  by $\Phi(T)\mapsto\Phi(\sum_w(d_1(Y_w)-c_w\operatorname{tr}Y_w)^{(w)}T)$;
\item[(ii)] if $\gamma$ is a line in $U_W$ generated by $(Y_w)$ relative to
  $(C_w)$, then $t\mapsto\omega(1,\gamma(t))$ acts at $\infty$ to first order
  by $\Phi(T)\mapsto\Phi(\sum_w(d_2(Y_w)+\tfrac12m(w)\operatorname{tr}Y_w)^{(w)}T)$.
\end{enumerate}
Every $Z\in M_{P+Q}(\mathbf C)$ is $Z_1+iZ_2$ with $Z_1,Z_2\in\mathfrak u(P,Q)$,
and then $J_{P,Q}Z^*J_{P,Q}=-Z_1+iZ_2$ (proof of
Lemma~\ref{asm-theta-agree-real-tensor}); the same holds for $(R,S)$. For
$Y\in\mathfrak u(R,S)$, $\overline{\operatorname{tr}Y}=-\operatorname{tr}Y$.

\emph{Step 2: the form on $\mathcal T$.} Let $b$ be a map as in
Definition~\ref{asm-theta-agree-kills-radical-maximal} for the pair, with
radical $R_b$ (the space called $R$ there), and let $P_0$ be a subspace as
there (called $P$ there). Put $b'(T,T')=b(\Phi(T),\Phi(T'))$.
\begin{enumerate}
\item[(iii)] $b'((d_1(Z)-c_w\operatorname{tr}Z)^{(w)}T,T')=0$ for all $w$,
  $Z\in M_{P+Q}(\mathbf C)$ and $T,T'$. For $Z\in\mathfrak u(P,Q)$, the family
  with $Z$ at $w$ and $0$ at the other places generates a line in $U_V$
  (Lemma~\ref{asm-theta-agree-real-lines}(2)), and by (i) and hypothesis (a) of
  the definition, $\Phi((d_1(Z)-c_w\operatorname{tr}Z)^{(w)}T)\in R_b$. The
  general case follows from $Z=Z_1+iZ_2$, as $b'$ is conjugate-linear in the
  first variable.
\item[(iv)] $b'(d_2(Y)^{(w)}T,T')=b'(T,d_2(JY^*J)^{(w)}T')$ for all $w$,
  $Y\in M_{R+S}(\mathbf C)$ and $T,T'$; in particular
  $b'(d_2(E_{ab})^{(w)}T,T')=\epsilon_{ab}b'(T,d_2(E_{ba})^{(w)}T')$, as
  $JE_{ab}^*J=\epsilon_{ab}E_{ba}$. For $Y\in\mathfrak u(R,S)$, the family with
  $Y$ at $w$ generates a line in $U_W$, and by (ii) and hypothesis (b),
  $b'(d_2(Y)^{(w)}T+\tfrac12m(w)\operatorname{tr}(Y)T,T')
  +b'(T,d_2(Y)^{(w)}T'+\tfrac12m(w)\operatorname{tr}(Y)T')=0$. As
  $\tfrac12m(w)$ is real and $\operatorname{tr}Y$ is imaginary, the terms with
  $\operatorname{tr}Y$ cancel, so
  $b'(d_2(Y)^{(w)}T,T')=-b'(T,d_2(Y)^{(w)}T')=b'(T,d_2(JY^*J)^{(w)}T')$. The
  general case follows from $Y=Y_1+iY_2$.
\end{enumerate}

\emph{Step 3: the subspaces $K_w$.} Fix $w$. By
Lemma~\ref{asm-theta-agree-real-parity} $\kappa_w$ is even, and as $G$ is
proper, the hypotheses of Lemma~\ref{asm-theta-agree-real-quotient} hold for
$(a,a',p,q,\kappa_w)$, as in the proof of
Lemma~\ref{asm-theta-agree-real-one-place}. Let $K_w\subset\mathcal P(P,Q;R,S)$
be the subspace of Lemma~\ref{asm-theta-agree-kills-radical-place} for
$c=c_w$. Let $\mathcal R_0$ be the set of $T\in\mathcal T$ with $b'(T,T')=0$
for every $T'$. Let $f=\bigotimes_uf_u$ be a pure tensor with $f_w\in K_w$,
and $t=\bigotimes_ut_u$ any pure tensor. For $x,y\in\mathcal P(P,Q;R,S)$ let
$b_0(x,y)=b'(f[w\mapsto x],t[w\mapsto y])$, where $f[w\mapsto x]$ is $f$ with
the factor at $w$ replaced by $x$. By (iii), by
$b'(T,T')=\overline{b'(T',T)}$ and by (iv), $b_0$ satisfies the hypotheses of
Lemma~\ref{asm-theta-agree-kills-radical-place} for $c=c_w$, so
$b'(f,t)=b_0(f_w,t_w)=0$. As the pure tensors span $\mathcal T$,
$f\in\mathcal R_0$.

\emph{Step 4: conclusion.} Let $\mathcal T_0=\Phi^{-1}(P_0)$. As $\Phi$ is a
bijection onto $\operatorname{Pol}(M)$,
$\mathcal T_0\supset\mathcal R_0=\Phi^{-1}(R_b)$. For $Y\in\mathfrak u(R,S)$
and $T\in\mathcal T_0$, by Lemma~\ref{asm-theta-agree-real-lines}(2), (ii) and
the hypothesis on $P_0$,
$d_2(Y)^{(w)}T+\tfrac12m(w)\operatorname{tr}(Y)T\in\mathcal T_0$, so
$d_2(Y)^{(w)}T\in\mathcal T_0$; with $Y=Y_1+iY_2$, $\mathcal T_0$ is stable
under every $d_2(Y)^{(w)}$. Lemma~\ref{asm-theta-agree-kills-radical-tensor},
with $V_w=\mathcal P(P,Q;R,S)$, which has countable dimension, the operators
$d_2(Y)$, the subspaces $K_w$, $\mathcal R_0$ and $\mathcal T_0$, gives
$\mathcal T_0\subset\mathcal R_0$ or $\mathcal T_0=\mathcal T$; that is,
$P_0\subset R_b$ or $P_0=\operatorname{Pol}(M)$.
\end{proof}
% Lean: Unit.radMax_of FockEx PolRange LieFirst LieSecond
% (CartanFactorisation/FockSecondActionMatrixUnitEq.lean). For each Weil pair of the triple,
% PC.exists_fock gives the Fock model Phi of the tensor product over the infinite places of
% the spaces Pol(P,Q,R,S) of Definition asm-theta-agree-real-fock, injective with range
% Pol(M_0) (PolRange), on which the lines act by fockAct1 and fockAct2 up to the characters
% c_1 and c_2 (LieFirst, LieSecond, as in Lemma asm-theta-agree-real-lie). radMax_pair
% transports beta to the tensor product, gets the relations sA and sB from line_at
% (RL.lineEx, Lemma asm-theta-agree-real-lines) and decomp (Z = Z_1 + i Z_2 with Z_1, Z_2 in
% u(p,q)) and star_E, and applies tens_max (TensMax.lean, through KReal.tensorIrr) with the
% property One(P,Q,R,S,-c_1(w)) at each infinite place w. One is one_data (with
% even_im_hermitian_sub for the parity, Lemma asm-theta-agree-real-parity), through one_of:
% exists_greatest_proper_invariant_submodule_of_side
% (ExistsGreatestProperInvariantSubmoduleOfSide.lean) applied to side_of, combV, combV' and
% nontrivial_of_pair. In the text: Step 1 is exists_fock, PolRange, LieFirst and LieSecond, as
% in Lemma asm-theta-agree-real (c_1(w) = -c_w, from hasExp_div; c_2(w) = m(w)/2, real); Step 2
% is beta', sA and sB; Step 3 is od (one_data) and the argument hR of tens_max (k4 applied to
% the form (beta'.compl2 (j t)).comp (j f)); Step 4 is the stability argument for N.comap phi
% and tens_max (Lemma asm-theta-agree-kills-radical-tensor).

\begin{lemma}\label{asm-theta-agree-kills-real}
Let $(Y,W,M_0,\chi_Y,\chi_W,\omega_0)$ be a Weil pair for a theta base $B$ of a
CM field $E$ (Definition~\ref{asm-theta-agree-triple}(1)), with $Y$ and
$W$ of ranks $n$ and $d$, vacua $\mathrm{vac}_v$ and theta integral $I$.
Suppose that the pair has maximal radicals
(Definition~\ref{asm-theta-agree-kills-radical-maximal}), and that
$N_1+N_2\neq\operatorname{Pol}(M_0)$ whenever $N_1$ and $N_2$ are theta
subspaces of $\operatorname{Pol}(M_0)$ other than $\operatorname{Pol}(M_0)$
(Definition~\ref{asm-theta-agree-rechoice-terms}(2)). Let
$\mathcal C$ be a theta class for the pair
(Definition~\ref{asm-theta-agree-kills-class}), and let $N$ be a maximal
element of the set of theta subspaces of $\operatorname{Pol}(M_0)$ other than
$\operatorname{Pol}(M_0)$.
Let $\phi_v$ be Schwartz-Bruhat functions on $F_v^{n\times d}$, for the finite
places $v$, with $\phi_v=\mathrm{vac}_v$ for almost all $v$, such that
$g\otimes\bigotimes_v\phi_v\in\mathcal C$ for every
$g\in\mathcal S(F_\infty^{n\times d})$. Then
$I(f\otimes\bigotimes_v\phi_v)=0$ for every $f\in N$.
\end{lemma}
\leanref{Development/WeilBound/ThetaLifts/CartanFactorisation/KillInf.lean}{PositivityTheory.ThKF.killInf'}
\leanref{Development/WeilBound/ThetaLifts/CartanFactorisation/ActsComp.lean}{PositivityTheory.ThKF.exists_L}
\begin{proof}
Put $\xi=\bigotimes_v\phi_v$, a Schwartz-Bruhat function on
$\mathbf A_f^{n\times d}$ as $M_0$ is regular, so that
$g\otimes\bigotimes_v\phi_v=g\otimes\xi$. As the functions in $\mathcal C$
converge and theta functions are additive
(Lemma~\ref{asm-theta-agree-convergence-theta}(2)), $I$ is linear on
$\mathcal C$. For $h\in U_W$ let $L_h(g)=I(g\otimes\xi)(h)$; it is a linear form
on $\mathcal S(F_\infty^{n\times d})$, continuous by condition (3) of a theta
class. Write $Q=\operatorname{Pol}(M_0)$.

\emph{Step 1: translations.} Let $g_0\in U_Y$ and $h_0\in U_W$ be at $\infty$,
and $T$ an archimedean operator of $\omega_0$ at $(g_0,h_0)$
(Definition~\ref{asm-theta-agree-rechoice-terms}(1)), so that
$\omega_0(g_0,h_0)(g\otimes\xi)=T(g)\otimes\xi$. If $g_0=1$, then
$L_h(T(g))=L_{hh_0}(g)$ for all $h$, as
$\theta_{\omega_0}(\omega_0(1,h_0)\Phi)(g',h)=\theta_{\omega_0}(\Phi)(g',hh_0)$.
If $h_0=1$, then $L_h(T(g))=L_h(g)$, by
Lemma~\ref{asm-theta-agree-anisotropic-invariance}.

\emph{Step 2: the form.} As in Step 5 of the proof of
Lemma~\ref{asm-theta-agree-kills-finite}, there are a Haar measure $\mu$ on
$U_W$ and a fundamental domain $\mathcal D$ for $U_W(F)$ with respect to $\mu$.
For $x\in Q$ put $u(x)=I(x\otimes\xi)$, so that $u(x)(h)=L_h(x)$. It is linear
in $x$; $u(x)$ is continuous, by condition (2), as $x\otimes\xi\in\mathcal C$;
$u(x)(\gamma h)=u(x)(h)$
for $\gamma\in U_W(F)$, by Lemma~\ref{asm-theta-agree-anisotropic-automorphy};
and $|u(x)|^2$ is integrable on $\mathcal D$, by condition (4), as
$x\otimes\xi\in\mathcal C\cap\operatorname{PD}(M_0)$. Let $\beta$ be
the form of Lemma~\ref{asm-theta-agree-kills-forms} on $Q$, and $R$ the space of
$x\in Q$ with $\beta(x,y)=0$ for every $y\in Q$.

\emph{Step 3: the operators of $U_W$.} Let $k\in K_\infty(W)$. By
Lemma~\ref{asm-theta-agree-rechoice-linear}(2) and (3) there is a linear map
$A_k$ of $Q$ such that $\omega_0(1,k)$ acts at $\infty$ by $x\mapsto A_kx$, and
by no other map. So $A_kx=T(x)$ for an archimedean operator $T$ at $(1,k)$
(Lemma~\ref{asm-theta-agree-rechoice-operators}(1)), and by Step 1,
$u(A_kx)(h)=u(x)(hk)$. Lemma~\ref{asm-theta-agree-kills-forms}(3) gives
$\beta(A_kx,A_ky)=\beta(x,y)$. Let $\delta$ be a line in $U_W$. By
Lemma~\ref{asm-theta-agree-rechoice-linear}(1) and (3), with the constant line
in $U_Y$, there is a linear map $D_\delta$ of $Q$ such that
$t\mapsto\omega_0(1,\delta(t))$ acts at $\infty$ to first order by
$x\mapsto D_\delta x$, and by no other map. Let $T_t$ be archimedean operators
of $\omega_0$ at $(1,\delta(t))$
(Lemma~\ref{asm-theta-agree-rechoice-operators}(1)). By
Lemma~\ref{asm-theta-agree-rechoice-lie} and this uniqueness,
$t\mapsto L(T_tx)$ has derivative $L(D_\delta x)$ at $t=0$ for every continuous
linear form $L$. For $L=L_h$, Step 1 gives $L_h(T_tx)=u(x)(h\delta(t))$; so
$t\mapsto u(x)(h\delta(t))$ has derivative $u(D_\delta x)(h)$ at $0$, and
Lemma~\ref{asm-theta-agree-kills-forms}(4) gives
$\beta(D_\delta x,y)+\beta(x,D_\delta y)=0$.

\emph{Step 4: the augmentation.} Let $A_0$ be the span of the functions $y$
such that $t\mapsto\omega_0(\gamma(t),1)$ acts at $\infty$ to first order by
$x\mapsto y$ for some $x\in Q$ and some line $\gamma$ in $U_Y$, and of the
functions $y-x$ such that $\omega_0(k,1)$ acts at $\infty$ by $x\mapsto y$ for
some $x\in Q$ and $k\in K_\infty(Y)$. Then $A_0\subset Q$, and every $L_h$
vanishes on $A_0$. For the first kind, $y\in Q$ by
Lemma~\ref{asm-theta-agree-rechoice-linear}(1); with archimedean operators
$T_t$ at $(\gamma(t),1)$, Lemma~\ref{asm-theta-agree-rechoice-lie} and
Lemma~\ref{asm-theta-agree-real-weak}(3) show that $t\mapsto L_h(T_tx)$ has
derivative $L_h(y)$ at $0$, and $L_h(T_tx)=L_h(x)$ for all $t$ by Step 1; so
$L_h(y)=0$. For the second kind, $y\in Q$ by
Lemma~\ref{asm-theta-agree-rechoice-linear}(2), $y=T(x)$ for an archimedean
operator $T$ at $(k,1)$, by the uniqueness there, and $L_h(y)=L_h(x)$ by
Step 1. So $u(z)=0$ for
$z\in A_0$, hence $A_0\subset R$.

\emph{Step 5: theta subspaces.} By the uniqueness in Step 3, a subspace
$P\subset Q$ is a theta subspace if and only if $A_0\subset P$ and
$A_kP\subset P$ and $D_\delta P\subset P$ for every $k\in K_\infty(W)$ and every
line $\delta$ in $U_W$: the first two conditions of
Definition~\ref{asm-theta-agree-rechoice-terms}(2) say that the generators of
$A_0$ lie in $P$, and the third says that $P$ is stable under the $A_k$ and the
$D_\delta$.

\emph{Step 6: $R$ is a theta subspace.} It contains $A_0$ by Step 4. If
$x\in R$, then $\beta(A_kx,A_kx)=\beta(x,x)=0$ by Step 3, so $A_kx\in R$ by
Lemma~\ref{asm-theta-agree-kills-forms}(2); and
$\beta(D_\delta x,y)=-\beta(x,D_\delta y)=0$ for every $y$, so $D_\delta x\in R$.
By Step 5, $R$ is a theta subspace.

\emph{Step 7: $N\subset R$.} We may suppose $R\neq Q$. The form $\beta$
satisfies the hypotheses (a) and (b) of
Definition~\ref{asm-theta-agree-kills-radical-maximal}: it is
Hermitian by Lemma~\ref{asm-theta-agree-kills-forms}(1); (a) holds as the
functions $y$ there lie in $A_0\subset R$; and (b) holds by Step 3, as
$x_1=D_\delta x$ and $y_1=D_\delta y$. The subspace $N+R$ contains $R$ and is
stable under the $D_\delta$, since $N$ and $R$ are, by Steps 5 and 6. As the
pair has maximal radicals, $N+R\subset R$ or $N+R=Q$. In the second case, $N$
and $R$ are theta subspaces other than $\operatorname{Pol}(M_0)$ whose sum is
$\operatorname{Pol}(M_0)$, which the hypothesis excludes. So $N\subset R$.

\emph{Step 8: conclusion.} Let $f\in N$. Then $f\in R$ by Step 7, so
$\beta(f,f)=0$, and $u(f)=0$ by Lemma~\ref{asm-theta-agree-kills-forms}(2).
That is, $I(f\otimes\bigotimes_v\phi_v)=0$.
\end{proof}
% Lean: killInf' (KillInf.lean), with hypotheses hAd = RadMaxAt (Definition
% asm-theta-agree-kills-radical-maximal) and hsum = JoinAt, which
% theta_kills_cusp_forms_spherical takes from RadMax (Lemma asm-theta-agree-kills-radical) and
% JoinTri (Theorem lit-howe-duality-real(2)). L is exists_L
% (Cls.contf and I_add'); Step 1 is LW (I_right) and LX (I_left); Step 2 uniW and exists_beta;
% Step 3 kappa (op from ArchOp, hA from ArchUniq, acts_comp), iso (L2.form_right), exD
% (ArchLie, with hD and hDu) and sk (L2.skew); Step 4 gen, augle and uaug; Step 5 hiff;
% Step 6 b0 and a2 (rad_of_zero); Step 7 hk; Step 8 key with L2.eq_zero. The maximality of N
% is not used.

\begin{lemma}\label{asm-theta-agree-kills-full-rank}
Let $(B,G,\mathcal X,\omega,\omega')$ be a triple of ranks $(a,a')$ for a theta
base $B$ of a CM field $E$ (Definition~\ref{asm-theta-agree-triple}), with
model $M'$, and put $I'_0=\{1,\ldots,a'\}\times\{1,\ldots,a+a'\}$. Let
$\mathfrak a^\circ_{v_1}(\omega')$ be the span of the functions $\phi_1-\phi$
with $\phi$ a Schwartz-Bruhat function on $F_{v_1}^{I'_0}$ of full rank at
$\tau_1$ (Definition~\ref{asm-theta-agree-models}(5)), $\phi_1$ Schwartz-Bruhat,
and $\omega'(g,1)$ acting at $v_1$ by $\phi\mapsto\phi_1$ for some
$g\in U_{V'}^{(v_1)}$. If $\phi$ lies in the span of the Schwartz-Bruhat
functions of full rank at $\tau_1$ and in $\mathfrak a_{v_1}(\omega')$
(Definition~\ref{asm-theta-agree-matched}(1)), then
$\phi\in\mathfrak a^\circ_{v_1}(\omega')$.
\end{lemma}
\leanref{Development/WeilBound/ThetaLifts/CartanFactorisation/CartanConj.lean}{PositivityTheory.ThKF.SoCoinv}
\leanref{Development/WeilBound/ThetaLifts/CartanFactorisation/CartanConj.lean}{PositivityTheory.ThKF.augOn}
\leanref{Development/WeilBound/CoreStatements/TemperedRepresentationTheory/BernsteinNarrowness/BernsteinClosure/RamifiedMultiplicityOne.lean}{PositivityTheory.BernsteinRamification.Close.soCoinv}
\leanref{Development/AutomorphicSpectrum/Proofs/BlockDecomposition/SoCoinvOfNarrow.lean}{GlobalTheory.CompanionTheory.SoCoinvNarrow.soCoinv_of}
\leanref{Development/WeilBound/CoreStatements/LocallyConstantCoinvariants/FullOfSoLoc.lean}{PositivityTheory.KCo.full_of_soLoc}
\begin{proof}
Let $S$ be the span of the Schwartz-Bruhat functions on $F_{v_1}^{I'_0}$ of
full rank at $\tau_1$. By Lemma~\ref{asm-theta-agree-finite} there are finite
maps $(\Pi_v,p_v,p'_v)_v$ matched to the triple; write $p'=p'_{v_1}$. By
Definition~\ref{asm-theta-agree-matched}(3), $\Pi_{v_1}\neq0$, $p'$ maps the
Schwartz-Bruhat functions onto $\Pi_{v_1}$ and vanishes on
$\mathfrak a_{v_1}(\omega')$. As $G$ is proper, $a'\geq1$.

\emph{A dictionary.} Apply
Lemma~\ref{asm-theta-agree-coinvariants-inputs}(3) to the Weil pair
$(V',W,M',\alpha,\alpha/\beta,\omega')$ at $v_1$ and $\tau_1$. It gives a
linear model $\iota$ of the second kind for $(\tau_1,1,1)$ for $M'$ and unitary
smooth characters $\chi_1,\chi_2$ of $F_{v_1}^\times$ such that $\iota$ is a
dictionary at $v_1$ with characters $(\chi_1,\chi_2)$
(Definition~\ref{asm-theta-agree-coinvariants-local}(3)), and such that
$\chi_1=1$ only if $\nu_{\tau_1}=1$ or $\nu_{\tau_1\circ c}=1$, where
$\nu=\alpha/\beta$. As $v_1$ is clean for $\mathcal X$
(Definition~\ref{asm-theta-agree-triple}(3)),
Lemma~\ref{asm-theta-agree-coinvariants-nontrivial} excludes both; so
$\chi_1\neq1$. Write $\omega_\circ=\omega_{\chi_1,\chi_2}$, and let
$\mathfrak a^\circ$ be the span of the $\omega_\circ(g_0,1)\Psi-\Psi$ with
$g_0\in GL_{a'}(F_{v_1})$ and $\Psi$ of full rank on
$M_{a'\times(a+a')}(F_{v_1})$.

\emph{Full rank under $\iota$.} If $\kappa$ is a Schwartz-Bruhat function of
full rank at $\tau_1$, then $\iota\kappa$ has full rank. Indeed, let $\iota'$ be
a linear model of the second kind for $(\tau_1,1,1)$ with $\iota'\kappa(Y)=0$
for every $Y$ of rank different from $a'$, and let $Y$ be such a matrix. The
matrices of $V'$ and $W$ are invertible
(Lemma~\ref{asm-theta-agree-convergence-signature}), so by
Lemma~\ref{asm-theta-agree-coinvariants-eigen}(1) there is $x$ at $v_1$ with
$\pi_{\tau_1}(x)=0$ and $\operatorname{co}(x)=Y$. Let $l(\kappa')=\iota\kappa'(Y)$.
For $u$ at $v_1$ with $\pi_{\tau_1\circ c}(u)=0$ and Schwartz-Bruhat $\kappa'$,
$\rho(u)$ acts at $v_1$ by $\kappa'\mapsto\rho_{v_1}(u)\kappa'$
(Lemma~\ref{asm-theta-agree-coinvariants-eigen}(2)), so, as $\iota$ is a linear
model of the second kind,
$l(\rho_{v_1}(u)\kappa')=\psi(\langle x,u\rangle)l(\kappa')$, with $\psi$ the
character of $B$. By Lemma~\ref{asm-theta-agree-coinvariants-eigen}(3), for
$\iota'$ and $x_1=x$, $l(\kappa)=0$, as
$\iota'\kappa(\operatorname{co}(x))=\iota'\kappa(Y)=0$. So $\iota\kappa(Y)=0$.
Hence $\iota$ maps $S$ into the functions of full rank. Conversely, a
Schwartz-Bruhat $\kappa$ with $\iota\kappa$ of full rank has full rank at
$\tau_1$, with $\iota$ in the definition.

\emph{The space $N$.} Let $K$ be the space of Schwartz-Bruhat $y$ on
$F_{v_1}^{I'_0}$ with $p'(y)=0$, and $N=\iota(K)$. First,
$\omega_\circ(g_0,1)\Psi-\Psi\in N$ for every $g_0\in GL_{a'}(F_{v_1})$ and
every Schwartz-Bruhat $\Psi$; so $N$ contains $\mathfrak a^\circ$. Indeed,
$\Psi=\iota\kappa$ with $\kappa$ Schwartz-Bruhat, as $\iota$ is bijective
between the Schwartz-Bruhat functions; by (a) of the definition of a dictionary
there is $g\in U_{V'}^{(v_1)}$ such that $\omega'(g,1)$ acts at $v_1$ by
$\kappa\mapsto\kappa_1$ with $\kappa_1$ Schwartz-Bruhat and
$\iota\kappa_1=\omega_\circ(g_0,1)\Psi$; so
$\omega_\circ(g_0,1)\Psi-\Psi=\iota(\kappa_1-\kappa)$, and
$\kappa_1-\kappa\in\mathfrak a_{v_1}(\omega')\subset K$. Second, let $\Psi=\iota y$
with $y\in K$, and $h_0\in GL_{a+a'}(F_{v_1})$. By (b) of that definition
there is $h\in U_W^{(v_1)}$ such that $\omega'(1,h)$ acts at $v_1$ by
$y\mapsto y_1$ with $y_1$ Schwartz-Bruhat and
$\iota y_1=\omega_\circ(1,h_0)\Psi$. As $\omega(1,h)$ acts at $v_1$ by
$0\mapsto0$ and $p_{v_1}(0)=0=p'(y)$, the matching
(Definition~\ref{asm-theta-agree-matched}(3)) gives $p'(y_1)=p_{v_1}(0)=0$.
So $\omega_\circ(1,h_0)\Psi\in N$. By
Lemma~\ref{asm-theta-agree-local-mult-ramified-irreducible}, with $n=a'$,
$d=a+a'$, $\nu=\chi_1$ and $\alpha=\chi_2$, either $N$ contains every function
of full rank, or every function of full rank in $N$ lies in $\mathfrak a^\circ$.

\emph{The first case does not occur.} As $\Pi_{v_1}\neq0$, there is a
Schwartz-Bruhat $y$ with $p'(y)\neq0$. By
Lemma~\ref{asm-theta-agree-coinvariants-kernel}(1), $y=y_0+z$ with $y_0$
Schwartz-Bruhat, $\iota y_0$ of full rank and $z\in\mathfrak a_{v_1}(\omega')$.
In the first case $\iota y_0\in N$, so $\iota y_0=\iota y'$ with $y'\in K$, and
$y_0=y'$, as $\iota$ is injective on the Schwartz-Bruhat functions. Then
$p'(y)=p'(y_0)+p'(z)=0$, a contradiction.

\emph{The second case.} As $\phi\in\mathfrak a_{v_1}(\omega')\subset K$, the
function $\iota\phi$ lies in $N$, and it has full rank, as $\phi\in S$. So
$\iota\phi\in\mathfrak a^\circ$. Let $\omega_\circ(g_0,1)\Psi-\Psi$ be a
generator of $\mathfrak a^\circ$, with $\Psi$ of full rank, and write
$\Psi=\iota\kappa$, $g$ and $\kappa_1$ as above. Then $\kappa$ has full rank at
$\tau_1$, so $\kappa_1-\kappa\in\mathfrak a^\circ_{v_1}(\omega')$, and the
generator is $\iota(\kappa_1-\kappa)$. So
$\mathfrak a^\circ\subset\iota(\mathfrak a^\circ_{v_1}(\omega'))$, and
$\iota\phi=\iota\kappa'$ with $\kappa'\in\mathfrak a^\circ_{v_1}(\omega')$, a
Schwartz-Bruhat function. As $\iota$ is injective on the Schwartz-Bruhat
functions, $\phi=\kappa'$.
\end{proof}
% Lean: Close.soCoinv = soCoinv_of_geo geometricLemmaFiltration = SoCoinvNarrow.soCoinv_of
% (narrowBern hG) dictionary_statement hK (SoCoinvOfNarrow.lean). hK (KFin) gives Rf; hD
% (DictSt, Lemma asm-theta-agree-coinvariants-inputs(3)) and nu_statement give iota and
% nu != 1; hf is full_of_soLoc (exists_coB, eigen_kill and mult'); N is
% (SB inf ker p').map iota, with aug_le_map and rightOnto; irrF is applied with the
% hypothesis aug <= N, and the book's Lemma asm-theta-agree-local-mult-ramified-irreducible
% (irrF', the same argument) needs only augF <= N (augF_le); the first case uses
% exists_full (Lemma asm-theta-agree-coinvariants-kernel(1)) and hF.nontriv, hF.onto',
% hF.inv'; the second is key (leftOnto) with hb.2.1.

\begin{lemma}\label{asm-theta-agree-local-injective-kills}
Let $(B,G,\mathcal X,\omega,\omega',R,R_f)$ be a matched quadruple for a theta
base $B$ of a CM field $E$ (Definition~\ref{asm-theta-agree-matched}), with
real maps $(\Pi,p,p')$, finite maps $(\Pi_v,p_v,p'_v)_v$ and theta integrals
$I$ and $I'$.
\begin{enumerate}
\item Let $f\in\operatorname{Pol}(M)$, and let $\phi_v$ be Schwartz-Bruhat
  functions on $F_v^{I_0}$ with $\phi_v=\mathrm{vac}_v$ for almost all $v$. If
  $p(f)=0$, or $p_v(\phi_v)=0$ for some $v$, then
  $I(f\otimes\bigotimes_v\phi_v)=0$.
\item Let $f\in\operatorname{Pol}(M')$, and let $\phi_v$ be Schwartz-Bruhat
  functions on $F_v^{I'_0}$ with $\phi_v=\mathrm{vac}'_v$ for almost all $v$
  and $\phi_{v_1}$ in the span of the Schwartz-Bruhat functions of full rank
  at $\tau_1$. If $p'(f)=0$, or $p'_v(\phi_v)=0$ for some $v$, then
  $I'(f\otimes\bigotimes_v\phi_v)=0$.
\end{enumerate}
\end{lemma}
\leanref{Development/WeilBound/ThetaLifts/CartanFactorisation/CartanConj.lean}{PositivityTheory.ThKF.Kills}
\leanref{Development/WeilBound/ThetaLifts/SphericalLattices/ThetaKillsCuspFormsSpherical.lean}{PositivityTheory.Sph.theta_kills_cusp_forms_spherical}
\begin{proof}
By Definition~\ref{asm-theta-agree-triple}(4),
$(V,W,M,\beta,\beta/\alpha,\omega)$ and $(V',W,M',\alpha,\alpha/\beta,\omega')$
are Weil pairs, the two Weil pairs of the triple. By
Lemma~\ref{asm-theta-agree-kills-class-triple}, the Schwartz-Bruhat functions
on $\mathbf A^{I_0}$ form a theta class $\mathcal C$ for the first, and
$\mathcal C'=\mathcal S^\circ(M')$ is a theta class for the second. For a family
$\phi=(\phi_u)_u$ of functions $\phi_u$ on $F_u^{I_0}$ or on $F_u^{I'_0}$, over
the finite places $u$, and for $g$ on $F_\infty^{I_0}$ or $F_\infty^{I'_0}$,
write $g\otimes\phi=g\otimes\bigotimes_u\phi_u$, and for a finite place $v$ and
a function $x$ let $\phi[v\mapsto x]$ be $\phi$ with $\phi_v$ replaced by $x$.
As $M$ and $M'$ are regular, $g\otimes\phi$ is a Schwartz-Bruhat function when
$g$ is a Schwartz function, every $\phi_u$ is Schwartz-Bruhat and $\phi_u$ is
the vacuum for almost all $u$.

Both Weil pairs have maximal radicals, by
Lemma~\ref{asm-theta-agree-kills-radical}. For each of them, two theta
subspaces of its Fock space other than the whole space do not have the whole
space as their sum, by Theorem~\ref{lit-howe-duality-real}(2), applied as in
the proof of Lemma~\ref{asm-theta-agree-rechoice-howe}(1). Let $v$ be a finite
place and $(\omega_v,\rho_v)$ local operators at $v$ for one of the pairs. As in
the proof of Lemma~\ref{asm-theta-agree-rechoice-howe}(2),
Theorem~\ref{th-howe-finite} applies, with $\varpi_v(1,h)=\omega_v(1,h)$ and
$\mathcal R_v$ the space $\mathfrak a_v(\omega)$ or $\mathfrak a_v(\omega')$ of
the pair (Definition~\ref{asm-theta-agree-rechoice-terms}(3)), and the theta
subspaces at $v$ are the subspaces between $\mathcal R_v$ and the
Schwartz-Bruhat functions that are stable under every $\omega_v(1,h)$. Its first alternative is hypothesis (i) of
Lemma~\ref{asm-theta-agree-kills-finite}, and its second is (ii), since by
Lemma~\ref{asm-theta-agree-finite-unique} any two choices of local operators at
$v$ agree on the Schwartz-Bruhat functions. So the hypotheses of
Lemmas~\ref{asm-theta-agree-kills-finite} and~\ref{asm-theta-agree-kills-real}
on the pair hold for both pairs and every finite place $v$.

(1) Let $f$ and $\phi$ be as in (1). Then $g\otimes\phi$ and
$f\otimes\phi[v\mapsto x]$ lie in $\mathcal C$ for every Schwartz function $g$,
every finite place $v$ and every Schwartz-Bruhat $x$. If $p(f)=0$, then $f$ lies
in $\ker p\cap\operatorname{Pol}(M)$, a maximal element of the set of theta
subspaces of $\operatorname{Pol}(M)$ other than $\operatorname{Pol}(M)$, by
Lemma~\ref{asm-theta-agree-rechoice-kernels}(1); and
Lemma~\ref{asm-theta-agree-kills-real} gives $I(f\otimes\phi)=0$. If
$p_v(\phi_v)=0$, then $\phi_v$ lies in $\ker p_v\cap\mathcal S(F_v^{I_0})$, a
theta subspace at $v$ of $\omega$ other than $\mathcal S(F_v^{I_0})$, by
Lemma~\ref{asm-theta-agree-rechoice-kernels}(2); and
Lemma~\ref{asm-theta-agree-kills-finite} gives $I(f\otimes\phi)=0$.

(2) Let $f$ and $\phi$ be as in (2), and let $S$ be the span of the
Schwartz-Bruhat functions on $F_{v_1}^{I'_0}$ of full rank at $\tau_1$. If
$\chi$ is a family of Schwartz-Bruhat functions, almost all equal to the vacua,
with $\chi_{v_1}\in S$, then $g\otimes\chi\in\mathcal S^\circ(M')$ for every
Schwartz function $g$: the function $g\otimes\chi[v_1\mapsto x]$ is linear in
$x$, and lies in $\mathcal S^\circ(M')$ when $x$ has full rank at $\tau_1$. So
$g\otimes\phi\in\mathcal C'$ for every $g$, and
$f\otimes\phi[v\mapsto x]\in\mathcal C'$ for every finite place $v\neq v_1$ and
every Schwartz-Bruhat $x$. If $p'(f)=0$, then
Lemmas~\ref{asm-theta-agree-rechoice-kernels}(1)
and~\ref{asm-theta-agree-kills-real} give $I'(f\otimes\phi)=0$, as in (1). If
$p'_v(\phi_v)=0$ for some $v\neq v_1$, then
Lemmas~\ref{asm-theta-agree-rechoice-kernels}(2)
and~\ref{asm-theta-agree-kills-finite} give $I'(f\otimes\phi)=0$.

Let now $p'_{v_1}(\phi_{v_1})=0$. The place $v_1$ is clean for $\mathcal X$
(Definition~\ref{asm-theta-agree-triple}(3)), and the quadruple describes the
coinvariants (Lemma~\ref{asm-theta-agree-coinvariants}), so
$\phi_{v_1}\in\mathfrak a_{v_1}(\omega')$. By
Lemma~\ref{asm-theta-agree-kills-full-rank}, $\phi_{v_1}$ lies in the span of
the functions $x_1-x$ with $x$ of full rank at $\tau_1$, $x_1$
Schwartz-Bruhat, and $\omega'(g,1)$ acting at $v_1$ by $x\mapsto x_1$ for some
$g\in U_{V'}^{(v_1)}$. Put $j(y)=f\otimes\phi[v_1\mapsto y]$; it is linear in
$y$, and $j(\phi_{v_1})=f\otimes\phi$. The elements of $\mathcal C'$ converge,
so $I'$ is linear on $\mathcal C'$, as theta functions are additive
(Lemma~\ref{asm-theta-agree-convergence-theta}(2)). Let $x$, $x_1$ and $g$ be
as above. Then $j(x)\in\mathcal C'$. By the definition of acting at $v_1$,
$j(x_1)=\omega'(g,1)j(x)$, so $j(x_1)\in\mathcal C'$, as $\omega'(g,1)$ maps
$\mathcal S^\circ(M')$ into itself (Lemma~\ref{asm-theta-agree-coinvariants});
and $I'(j(x_1))=I'(j(x))$ by
Lemma~\ref{asm-theta-agree-anisotropic-invariance}. So
$j(x_1-x)\in\mathcal C'$ and $I'(j(x_1-x))=0$. By linearity, $j$ maps the span
of these functions into $\mathcal C'$, and $I'$ vanishes on its image. So
$I'(f\otimes\phi)=I'(j(\phi_{v_1}))=0$.
\end{proof}
% Lean: Sph.theta_kills_cusp_forms_spherical ThetaBasic WeilConv SoCont SqInt KDecay KCoinv
% ICont ArchOp ArchLie ArchUniq RadMax SphTri JoinTri SoCoinv
% (SphericalLattices/ThetaKillsCuspFormsSpherical.lean). c1 = clsV and c2 = clsV' (Lemma
% asm-theta-agree-kills-class-triple), sch from pure_eq_pureA and pureA_mem. killInf' (Lemma
% asm-theta-agree-kills-real) with maxInfP (Lemma asm-theta-agree-rechoice-kernels(1)), RadMax
% (Lemma asm-theta-agree-kills-radical) and JoinTri for p(f) = 0 and p'(f) = 0, the latter
% with pure_soK; killFin3 (Lemma asm-theta-agree-kills-finite) with maxFinP and SphTri for
% p_v(phi_v) = 0 and, in the second space, for v != v_1. For v = v_1: the first clause of
% KCoinv gives the augmentation, SoCoinv (Lemma asm-theta-agree-kills-full-rank) the span,
% and key is the span induction with pure_soK, the second clause of KCoinv, I_add' and I_left.
% Kills is also the input of Factor = ThKF.factor_of KDecay Kills (Lemma
% asm-theta-agree-definite-inputs (5)). The hypotheses of killFin3 and killInf' on the pair
% are (hSph ...).1 and .2, (hAi ...).1 and .2 and (hH ...).1 and .2, the paragraph after the
% first one.

\begin{lemma}\label{asm-theta-agree-local-injective-kills-zero}
Let $(B,G,\mathcal X,\omega,\omega',R,R_f)$ be a matched quadruple for a theta
base $B$ of a CM field $E$ with $q=0$, with real maps $(\Pi,p,p')$, finite
maps $(\Pi_v,p_v,p'_v)_v$ and theta integral $I'$ of the second space. Let
$f\in\operatorname{Pol}(M')$, and let $\phi_v$ be Schwartz-Bruhat functions on
$F_v^{I'_0}$ with $\phi_v=\mathrm{vac}'_v$ for almost all $v$. If $p'(f)=0$, or
$p'_v(\phi_v)=0$ for some $v$, then $I'(f\otimes\bigotimes_v\phi_v)=0$.
\end{lemma}
\leanref{Development/WeilBound/ThetaLifts/LocalTheory/SplitInv.lean}{PositivityTheory.ThKG.Loc.Kills0}
\leanref{Development/WeilBound/ThetaLifts/CartanFactorisation/Kills0Of.lean}{PositivityTheory.ThKF.kills0_of}
\begin{proof}
By Definition~\ref{asm-theta-agree-triple}(4),
$(V',W,M',\alpha,\alpha/\beta,\omega')$ is a Weil pair of signature
$(q,a'-q;\ldots)=(0,a';\ldots)$, so by Lemma~\ref{asm-theta-agree-convergence}
every Schwartz-Bruhat function converges for $(M',\omega')$, and $I'$ is
additive on Schwartz-Bruhat functions. For a function $y$ on $F_{v_1}^{I'_0}$
let $j(y)=f\otimes\bigotimes_v\phi_v$ with $\phi_{v_1}$ replaced by $y$; it
depends linearly on $y$, and it is Schwartz-Bruhat for Schwartz-Bruhat $y$, as
$M'$ is regular.

The place $v_1$ is clean for $\mathcal X$
(Definition~\ref{asm-theta-agree-triple}(3)), and the quadruple describes the
coinvariants (Lemma~\ref{asm-theta-agree-coinvariants}). So there is a
Schwartz-Bruhat $x$ of full rank at $\tau_1$ with
$p'_{v_1}(x)=p'_{v_1}(\phi_{v_1})$, and then $\phi_{v_1}-x\in\mathfrak a_{v_1}(\omega')$,
as $p'_{v_1}(\phi_{v_1}-x)=0$. By
Lemma~\ref{asm-theta-agree-anisotropic-augmentation},
$I'(j(\phi_{v_1}-x))=0$, so
$I'(f\otimes\bigotimes_v\phi_v)=I'(j(\phi_{v_1}))=I'(j(x))$. The family
defining $j(x)$ has the component $x$ of full rank at $\tau_1$ at $v_1$, the
other components of $\phi$, and $p'_{v_1}(x)=p'_{v_1}(\phi_{v_1})$. So if
$p'(f)=0$ or $p'_v(\phi_v)=0$ for some $v$,
Lemma~\ref{asm-theta-agree-local-injective-kills}(2) gives $I'(j(x))=0$.
\end{proof}
% Lean: kills0_of ThetaBasic WeilConv KCoinv Kills (Kills0Of.lean): x and the augmentation
% come from the first clause of KCoinv (k2, k3); key is the vanishing on the augmentation
% space (I_left), which is Lemma asm-theta-agree-anisotropic-augmentation.

\begin{lemma}\label{asm-theta-agree-local-injective-general}
Let $(X,W,M,\chi_X,\chi_W,\omega)$ be a Weil pair for a theta base $B$ of a CM
field $E$ (Definition~\ref{asm-theta-agree-triple}(1)), with $X$ and $W$ of
ranks $n$ and $d$, vacua $\mathrm{vac}_v$ and theta integral $I$. Let
$\mathcal C$ be a space of Schwartz-Bruhat functions on $\mathbf A^{n\times d}$
such that every $\Phi\in\mathcal C$ converges
(Definition~\ref{asm-theta-agree-models}(7)), and such that
$f\mapsto I(f\otimes\xi)(h)$ is continuous on $\mathcal S(F_\infty^{n\times d})$
for every $h\in U_W$ and every Schwartz-Bruhat $\xi$ on
$\mathbf A_f^{n\times d}$ with $f\otimes\xi\in\mathcal C$ for all $f$. For each
finite place $v$ let $S_v$ be a space of Schwartz-Bruhat functions on
$F_v^{n\times d}$. Call a family $\phi=(\phi_v)_v$ \emph{in $S$} if
$\phi_v\in S_v$ for every $v$ and $\phi_v=\mathrm{vac}_v$ for almost all $v$,
and suppose that $f\otimes\phi=f\otimes\bigotimes_v\phi_v$ lies in $\mathcal C$
for every $f\in\mathcal S(F_\infty^{n\times d})$ and every family $\phi$ in $S$.
Let $p$ be a linear map from $\mathcal S(F_\infty^{n\times d})$ to a complex
vector space $\Pi$, and $p_v$ linear maps from the functions on
$F_v^{n\times d}$ to complex vector spaces $\Pi_v$, such that:
\begin{enumerate}
\item[(a)] if $f\in\operatorname{Pol}(M)$, $\phi$ is a family in $S$, and
  $p(f)=0$ or $p_v(\phi_v)=0$ for some $v$, then $I(f\otimes\phi)=0$;
\item[(b)] for every $v$, $0$ and $\Pi_v$ are the only subspaces $N\subset\Pi_v$
  with $p_v(x_1)\in N$ whenever $x,x_1\in S_v$, $p_v(x)\in N$ and $\omega(1,h)$
  acts at $v$ by $x\mapsto x_1$ for some $h\in U_W^{(v)}$;
\item[(c)] $0$ and $\Pi$ are the only subspaces $N\subset\Pi$ with
  $p(f_1)\in N$ whenever $f\in\operatorname{Pol}(M)$, $p(f)\in N$, and either
  $t\mapsto\omega(1,\gamma(t))$ acts at $\infty$ to first order by
  $f\mapsto f_1$ for a line $\gamma$ in $U_W$, or $\omega(1,k)$ acts at
  $\infty$ by $f\mapsto f_1$ for some $k\in K_\infty(W)$
  (Definition~\ref{asm-theta-agree-matched}(1)).
\end{enumerate}
If $f\in\operatorname{Pol}(M)$ and $\phi$ is a family in $S$ with $p(f)\neq0$,
$p_v(\phi_v)\neq0$ for every $v$ and $I(f\otimes\phi)=0$, then
$I(g\otimes\psi)=0$ for every $g\in\operatorname{Pol}(M)$ and every family
$\psi$ in $S$.
\end{lemma}
\leanref{Development/WeilBound/ThetaLifts/CartanFactorisation/CuspidalClassOfWeilConvergence.lean}{PositivityTheory.ThKF.Cls0}
\leanref{Development/WeilBound/ThetaLifts/CartanFactorisation/CuspidalClassOfWeilConvergence.lean}{PositivityTheory.ThKF.injGen}
\begin{proof}
As the functions in $\mathcal C$ converge and theta functions are additive
(Lemma~\ref{asm-theta-agree-convergence-theta}(2)), $I$ is linear on
$\mathcal C$. As $\omega$ is a representation,
$\theta_\omega(\omega(1,h)\Phi)(g,h')=\theta_\omega(\Phi)(g,h'h)$, so
$I(\omega(1,h)\Phi)(h')=I(\Phi)(h'h)$. For a family $\chi$, a finite place $v$
and a function $x$ on $F_v^{n\times d}$, write $\chi[v\mapsto x]$ for $\chi$
with $\chi_v$ replaced by $x$. Then $f\otimes\chi[v\mapsto x]$ depends linearly
on $x$, and if $\omega(1,h)$ acts at $v$ by $x\mapsto x_1$, then
$\omega(1,h)(f\otimes\chi[v\mapsto x])=f\otimes\chi[v\mapsto x_1]$, by the
definition of acting at $v$ (Definition~\ref{asm-theta-agree-models}(2)) with
$r(y)=f(y_\infty)\prod_{u\neq v}\chi_u(y_u)$.

\emph{Step 1: one place.} Let $\chi$ be a family in $S$ with
$I(f\otimes\chi)=0$, and $v$ a place with $p_v(\chi_v)\neq0$. We show
$I(f\otimes\chi[v\mapsto x])=0$ for every $x\in S_v$. Let $K$ be the set of
$x\in S_v$ with $I(f\otimes\chi[v\mapsto x])=0$, a subspace by the linearity of
$I$ on $\mathcal C$. By (a), $K$ contains every $x\in S_v$ with $p_v(x)=0$;
hence for $x\in S_v$, $p_v(x)\in p_v(K)$ if and only if $x\in K$. The subspace
$N=p_v(K)$ has the property in (b): if $x,x_1\in S_v$, $p_v(x)\in N$ and
$\omega(1,h)$ acts at $v$ by $x\mapsto x_1$, then $x\in K$ and
$I(f\otimes\chi[v\mapsto x_1])(h')=I(f\otimes\chi[v\mapsto x])(h'h)=0$, so
$x_1\in K$. As $\chi_v\in K$ and $p_v(\chi_v)\neq0$, (b) gives $N=\Pi_v$, so
$K=S_v$.

\emph{Step 2: all finite places.} Let $\psi$ be a family in $S$. For a finite
set $T$ of finite places let $\chi_T$ be the family equal to $\psi$ on $T$ and
to $\phi$ off $T$. By induction on $T$, $I(f\otimes\chi_T)=0$: for $v\notin T$,
$\chi_{T\cup\{v\}}=\chi_T[v\mapsto\psi_v]$, and $\chi_T$ has the component
$\phi_v$ at $v$, with $p_v(\phi_v)\neq0$, so Step 1 applies. As $\psi$ and
$\phi$ differ at finitely many places, $I(f\otimes\psi)=0$.

\emph{Step 3: the infinite places.} As $M$ is regular,
$\xi=\bigotimes_v\psi_v$ is a Schwartz-Bruhat function on the finite adeles,
and $g\otimes\psi=g\otimes\xi$ for every $g$. For $h\in U_W$, the map
$L_h(x)=I(x\otimes\xi)(h)$ is a continuous linear form on
$\mathcal S(F_\infty^{n\times d})$. Let $K$ be the set of $x$ with $L_h(x)=0$
for every $h$. By (a), $K$ contains every $x\in\operatorname{Pol}(M)$ with
$p(x)=0$, so for $x\in\operatorname{Pol}(M)$, $p(x)\in p(K\cap\operatorname{Pol}(M))$
if and only if $x\in K$. The subspace $N=p(K\cap\operatorname{Pol}(M))$ has the
property in (c). Let $x\in\operatorname{Pol}(M)$ with $p(x)\in N$, so that
$x\in K$. If $k\in K_\infty(W)$ and $\omega(1,k)$ acts at $\infty$ by
$x\mapsto x_1$, then $\omega(1,k)(x\otimes\xi)=x_1\otimes\xi$, so
$L_h(x_1)=L_{hk}(x)=0$, and $x_1\in\operatorname{Pol}(M)$ by
Lemma~\ref{asm-theta-agree-rechoice-linear}(2). If $\gamma$ is a line in $U_W$
and $t\mapsto\omega(1,\gamma(t))$ acts at $\infty$ to first order by
$x\mapsto x_1$, let $T_t$ be archimedean operators of $\omega$ at
$(1,\gamma(t))$ (Lemma~\ref{asm-theta-agree-rechoice-operators}(1)). By
Lemma~\ref{asm-theta-agree-rechoice-lie}, with the constant line in $U_X$, there
is $y\in\operatorname{Pol}(M)$ such that $t\mapsto\omega(1,\gamma(t))$ acts at
$\infty$ to first order by $x\mapsto y$ and $t\mapsto L(T_tx)$ has derivative
$L(y)$ at $0$ for every continuous linear form $L$; by
Lemma~\ref{asm-theta-agree-real-weak}(3), $y=x_1$. As
$T_tx\otimes\xi=\omega(1,\gamma(t))(x\otimes\xi)$,
$L_h(T_tx)=L_{h\gamma(t)}(x)=0$ for all $t$, so $L_h(x_1)=0$ for every $h$. In
both cases $x_1\in K\cap\operatorname{Pol}(M)$. As $f\in K$ by Step 2 and
$p(f)\neq0$, (c) gives $N=\Pi$, so every $g\in\operatorname{Pol}(M)$ lies in $K$:
$I(g\otimes\psi)=0$.
\end{proof}
% Lean: injGen (CartanFactorisation/CuspidalClassOfWeilConvergence.lean), with the class
% Cls0 for C. Step 1 is stepF (kerIn, mem_kerIn, pureAct, inj_right), Step 2 is fin and
% hpsi0, Step 3 is the rest, with ArchOp (op), ArchLie and ArchUniq.

\begin{lemma}\label{asm-theta-agree-local-injective}
Let $(B,G,\mathcal X,\omega,\omega',R,R_f)$ be a matched quadruple for a theta
base $B$ of a CM field $E$, with real maps $(\Pi,p,p')$ and finite maps
$(\Pi_v,p_v,p'_v)_v$, such that $I(\Phi)\neq0$ for some
$\Phi\in\operatorname{PD}(M)$ and $I'(\Phi')\neq0$ for some
$\Phi'\in\operatorname{PD}(M')\cap\mathcal S^\circ(M')$.
\begin{enumerate}
\item Let $f\in\operatorname{Pol}(M)$ and $\phi_v\in\mathcal S(F_v^{I_M})$ with
  $\phi_v=\mathrm{vac}_v$ for almost all $v$. If $p(f)\neq0$ and
  $p_v(\phi_v)\neq0$ for every $v$, then $I(f\otimes\bigotimes_v\phi_v)\neq0$.
\item Let $f\in\operatorname{Pol}(M')$ and $\phi_v\in\mathcal S(F_v^{I_{M'}})$
  with $\phi_v=\mathrm{vac}'_v$ for almost all $v$, and suppose that $q=0$ or
  that $\phi_{v_1}$ is in the span of the Schwartz-Bruhat functions of full
  rank at $\tau_1$. If $p'(f)\neq0$ and $p'_v(\phi_v)\neq0$ for every $v$, then
  $I'(f\otimes\bigotimes_v\phi_v)\neq0$.
\end{enumerate}
\end{lemma}
\leanref{Development/WeilBound/ThetaLifts/LocalTheory/UnramifiedDeterminedBySpherical.lean}{PositivityTheory.ThKG.Loc.InjPure}
\leanref{Development/WeilBound/ThetaLifts/CartanFactorisation/CuspidalClassOfWeilConvergence.lean}{PositivityTheory.ThKF.injPure_of}
\begin{proof}
By Definition~\ref{asm-theta-agree-triple}(4), $(V,W,M,\beta,\beta/\alpha,\omega)$
and $(V',W,M',\alpha,\alpha/\beta,\omega')$ are Weil pairs; as $M$ and $M'$ are
regular, the pure tensors $f\otimes\bigotimes_v\phi_v$ with Schwartz-Bruhat
$\phi_v$, almost all equal to the vacua, are Schwartz-Bruhat functions. In each
part we suppose that the theta integral of the given pure tensor vanishes, and
apply Lemma~\ref{asm-theta-agree-local-injective-general}; it shows that the
theta integral vanishes on a space that contains a function with nonzero theta
integral, a contradiction. Two remarks give conditions (b) and (c) of that
lemma for the second space.

\emph{(i)} Let $v$ be a finite place and $S'_v$ a space of Schwartz-Bruhat
functions on $F_v^{I'_0}$ with $p'_v(S'_v)=\Pi_v$, such that $x_1\in S'_v$
whenever $x\in S'_v$, $x_1$ is Schwartz-Bruhat and $\omega'(1,h)$ acts at $v$ by
$x\mapsto x_1$ for some $h\in U_W^{(v)}$. Then condition (b) holds for
$\omega'$, $S'_v$ and $p'_v$. Indeed, let $N\subset\Pi_v$ have the property in
(b). Let $h\in U_W^{(v)}$ and $\phi,\phi_1$ Schwartz-Bruhat with $\omega(1,h)$
acting at $v$ by $\phi\mapsto\phi_1$ and $p_v(\phi)\in N$. Choose $y\in S'_v$
with $p'_v(y)=p_v(\phi)$; with local operators $\omega'_v$ at $v$ for the Weil
pair of $M'$ (Lemma~\ref{asm-theta-agree-finite-inputs}(2)), $\omega'(1,h)$ acts
at $v$ by $y\mapsto y_1=\omega'_v(1,h)y$, and $y_1\in S'_v$. By
Definition~\ref{asm-theta-agree-matched}(3), $p_v(\phi_1)=p'_v(y_1)\in N$. So
$N$ satisfies the irreducibility condition there, and $N$ is $0$ or $\Pi_v$.

\emph{(ii)} Condition (c) holds for $\omega'$, $M'$ and $p'$. Indeed, let
$N\subset\Pi$ have the property in (c). Let $f\in\operatorname{Pol}(M)$ with
$p(f)\in N$, and $f\mapsto f_1$ one of the actions of $\omega$ in
Definition~\ref{asm-theta-agree-matched}(2)(b). Choose
$f'\in\operatorname{Pol}(M')$ with $p'(f')=p(f)$, by (2)(a) there, and let
$f'\mapsto f'_1$ be the action of $\omega'$ for the same line or the same $k$,
which exists by Lemma~\ref{asm-theta-agree-rechoice-linear}(1) and (2). By
(2)(b) there, $p(f_1)=p'(f'_1)\in N$. So $N$ satisfies (2)(c) there, and $N$
is $0$ or $\Pi$.

\emph{Part (1).} Suppose $I(f\otimes\bigotimes_v\phi_v)=0$. The Weil pair of
$M$ has signature $(a-p,p;\ldots)$ and $a+a'+2p<2a$, as in the proof of
Lemma~\ref{asm-theta-agree-constants}; so by
Lemma~\ref{asm-theta-agree-convergence} every Schwartz-Bruhat function
converges for $(M,\omega)$, with the continuity required in
Lemma~\ref{asm-theta-agree-local-injective-general}. Apply that lemma with
$\mathcal C$ the Schwartz-Bruhat functions, $S_v$ the Schwartz-Bruhat
functions, $p$ and $p_v$: condition (a) is
Lemma~\ref{asm-theta-agree-local-injective-kills}(1), (b) is the irreducibility
condition of Definition~\ref{asm-theta-agree-matched}(3), and (c) is
Definition~\ref{asm-theta-agree-matched}(2)(c). So $I(g\otimes\psi)=0$ for all
$g\in\operatorname{Pol}(M)$ and all families $\psi$ of Schwartz-Bruhat functions
almost all equal to the vacua. As $I$ is linear on Schwartz-Bruhat functions,
$I$ vanishes on $\operatorname{PD}(M)$, a contradiction.

\emph{Part (2), $q=0$.} Suppose $I'(f\otimes\bigotimes_v\phi_v)=0$. The Weil pair
of $M'$ has signature $(0,a';\ldots)$, so
Lemma~\ref{asm-theta-agree-convergence} applies to it as in Part (1). Apply
Lemma~\ref{asm-theta-agree-local-injective-general} with all $S_v$ the
Schwartz-Bruhat functions, $p'$ and $p'_v$: (a) is
Lemma~\ref{asm-theta-agree-local-injective-kills-zero}, (b) is (i) with $S'_v$
the Schwartz-Bruhat functions, which $p'_v$ maps onto $\Pi_v$, and (c) is (ii).
As in Part (1), $I'$ vanishes on $\operatorname{PD}(M')$, a contradiction.

\emph{Part (2), $\phi_{v_1}$ in the span $S_{v_1}$ of the functions of full rank
at $\tau_1$.} Suppose $I'(f\otimes\bigotimes_v\phi_v)=0$. Let $S_v$ be the
Schwartz-Bruhat functions for $v\neq v_1$; a family in $S$ gives a function in
$\mathcal S^\circ(M')$, and let $\mathcal C=\mathcal S^\circ(M')$. By
Lemma~\ref{asm-theta-agree-decay}, every function in $\mathcal S^\circ(M')$
converges for $(M',\omega')$. By Lemma~\ref{asm-theta-agree-continuity},
$f\mapsto I'(f\otimes\xi)(h)$ is sequentially continuous, hence continuous, on
the Schwartz space, which is metrizable, whenever all $f\otimes\xi$ lie in
$\mathcal S^\circ(M')$. Condition (a) is
Lemma~\ref{asm-theta-agree-local-injective-kills}(2), and (c) is (ii). For (b)
we use (i), with $S'_v$ the Schwartz-Bruhat functions for $v\neq v_1$ and
$S'_{v_1}=S_{v_1}$. The place $v_1$ is clean, so by
Lemma~\ref{asm-theta-agree-coinvariants} every element of $\Pi_{v_1}$ is
$p'_{v_1}(x)$ for some $x$ of full rank at $\tau_1$, and
$p'_{v_1}(S_{v_1})=\Pi_{v_1}$. For the stability of $S_{v_1}$, let $\pi$ be a
linear projection of the space of functions on $F_{v_1}^{I'_0}$ onto $S_{v_1}$,
and define an operator $O$ on functions on $\mathbf A^{I'_0}$ by
$(O\Phi)(y)=\pi\bigl(z\mapsto\Phi(y\langle z\rangle)\bigr)(y_{v_1})$, where
$y\langle z\rangle$ is $y$ with its component at $v_1$ replaced by $z$. Then
$O(g\otimes\bigotimes_v\chi_v)=g\otimes\bigotimes_v\chi'_v$ with
$\chi'_{v_1}=\pi(\chi_{v_1})$ and $\chi'_v=\chi_v$ otherwise, so $O$ fixes every
element of $\mathcal S^\circ(M')$. Let $x\in S_{v_1}$, $h\in U_W^{(v_1)}$ and
$x_1$ Schwartz-Bruhat with $\omega'(1,h)$ acting at $v_1$ by $x\mapsto x_1$, and
let $g\neq0$ be a Schwartz function. The function $\Psi=g\otimes\bigotimes_v\chi_v$
with $\chi_{v_1}=x$ and $\chi_v=\mathrm{vac}'_v$ otherwise lies in
$\mathcal S^\circ(M')$, as it is linear in $x$, and so does
$\omega'(1,h)\Psi$, as the quadruple describes the coinvariants; this is the
same tensor with $x_1$ in place of $x$. As $O$ fixes it, it equals the tensor
with $\pi(x_1)$ in place of $x$, and evaluating at points where $g$ and the
other vacua do not vanish gives $x_1=\pi(x_1)\in S_{v_1}$. So (i) applies.

Lemma~\ref{asm-theta-agree-local-injective-general} gives $I'(g\otimes\psi)=0$
for all $g\in\operatorname{Pol}(M')$ and families $\psi$ in $S$. Let
$\Phi'\in\operatorname{PD}(M')\cap\mathcal S^\circ(M')$ with $I'(\Phi')\neq0$,
and write $\Phi'=\sum_ig_i\otimes\psi_i$ with $g_i\in\operatorname{Pol}(M')$ and
families $\psi_i$ of Schwartz-Bruhat functions almost all equal to the vacua.
Then $\Phi'=O\Phi'=\sum_ig_i\otimes\psi'_i$, where $\psi'_i$ is $\psi_i$ with its
component at $v_1$ replaced by its image under $\pi$, a family in $S$. As $I'$
is linear on $\mathcal S^\circ(M')$, $I'(\Phi')=0$, a contradiction.
\end{proof}
% Lean: ThKF.injPure_of ArchOp ArchLie ArchUniq ActsEx WeilConv SoCont KDecay KCoinv Kills
% Kills0 (CartanFactorisation/CuspidalClassOfWeilConvergence.lean). (i) is irrF (with
% ActsEx and hm.fin.equiv), (ii) is irrI (hLie, hm.real.lie and hm.real.cpt); Part (1)
% is injGen with cuspidal_class_of_weil_convergence and I_span; Part (2) with q = 0 uses
% Kills0 = kills0_of; the last case uses the class S-circ (KDecay, SoCont), exists_proj
% (exists_factorOp), pure_so, co2 (KCoinv) and pure_inj.

\begin{lemma}\label{asm-theta-agree-local-split-invariant}
Let $(B,G,\mathcal X,\omega,\omega')$ be a triple of ranks $(a,a')$ for a
theta base $B$ of a CM field $E$, let $H'$ be an adapted plane in $W$ over a
lattice space $W_1$ at some positions, and let $f$ be a function on $U_W$
that is $I(\Phi)$ for some $\Phi\in\operatorname{PD}(M)$, or $I'(\Phi')$ for
some $\Phi'\in\operatorname{PD}(M')$ with $\Phi'\in\mathcal S^\circ(M')$ or
$q=0$. If $(f_\alpha,f_\beta)$ splits $f$ along $H'$
(Definition~\ref{asm-theta-agree-planes}(3)), then $f_\alpha(uh)=f_\alpha(h)$
and $f_\beta(uh)=f_\beta(h)$ for all $u\in N_{H'}$ and $h\in U_W$, and $f=0$
implies $f_\alpha=f_\beta=0$.
\end{lemma}
\leanref{Development/WeilBound/ThetaLifts/LocalTheory/SplitInv2.lean}{PositivityTheory.ThKG.Loc.SplitInv2}
\leanref{Development/WeilBound/ThetaLifts/SmallCones/SplitInv.lean}{PositivityTheory.ThKCT.splitInv2_of}
\begin{proof}
Let $\mathcal X=(\alpha,\beta,S_0,Q)$. By
Definition~\ref{asm-theta-agree-triple}(4),
$(V,W,M,\beta,\beta/\alpha,\omega)$ and $(V',W,M',\alpha,\alpha/\beta,\omega')$
are Weil pairs. Let $(M_0,\omega_0,\Phi_0)$ be $(M,\omega,\Phi)$ or
$(M',\omega',\Phi')$, so that $f=I_{M_0,\omega_0}(\Phi_0)$ with
$\Phi_0\in\operatorname{PD}(M_0)$; the condition on $\Phi'$ is not used.

\emph{Step 1: $f(\gamma h)=f(h)$ for $\gamma\in U_W(F)$ and $h\in U_W$.} By
Lemma~\ref{asm-theta-agree-split-inputs}(2), $\Phi_0$ is a Schwartz-Bruhat
function. By Lemma~\ref{asm-theta-agree-anisotropic-automorphy}, applied with
$1\in U_{V}(F)$, respectively $U_{V'}(F)$, and $\gamma$,
$\theta_{\omega_0}(\Phi_0)(g,\gamma h)=\theta_{\omega_0}(\Phi_0)(g,h)$ for every
$g$. The two theta integrals at $\gamma h$ and at $h$ integrate the same
function over the domain of $M_0$, so they are equal.

\emph{Step 2: vanishing.} If $f=0$, then $(0,0)$ splits $f$
(Definition~\ref{asm-theta-agree-planes}(3)), and
Lemma~\ref{asm-theta-agree-split-unique} gives $f_\alpha=f_\beta=0$.

\emph{Step 3: the sum of the parts.} The space $W$ is skew-Hermitian of some
signature. By Lemma~\ref{asm-theta-agree-split-inputs}(1)(a) there are a Haar
measure $\nu$ on $N_{H'}$ and a fundamental domain $\mathcal F$ for
$N_{H'}\cap U_W(F)$ in $N_{H'}$ with $\nu(\mathcal F)=1$, and $\nu$ is
invariant under right translations. As $f_\alpha+f_\beta$ is the constant term
of $f$ along $H'$,
\[
f_\alpha(h)+f_\beta(h)=\int_{\mathcal F}f(u'h)\,d\nu(u')\qquad(h\in U_W).
\]
Let $u\in N_{H'}$ and $h\in U_W$. By Step 1 the function $u'\mapsto f(u'h)$ on
$N_{H'}$ is invariant under left translation by $N_{H'}\cap U_W(F)$, which is
countable by Lemma~\ref{asm-theta-agree-convergence-signature} and
Theorem~\ref{asm-theta-agree-convergence-siegel}(1). The computation in the
proof of Lemma~\ref{asm-theta-agree-anisotropic-invariance}, with $N_{H'}$,
$N_{H'}\cap U_W(F)$, $\nu$ and $\mathcal F$ in place of $U_X$, $U_X(F)$, $\mu$
and $\mathcal D$, shows that the integrals of $u'\mapsto f(u'uh)$ and of
$u'\mapsto f(u'h)$ over $\mathcal F$ agree. Hence
\[
f_\alpha(uh)+f_\beta(uh)=f_\alpha(h)+f_\beta(h)\qquad(u\in N_{H'},\ h\in U_W).
\tag{$*$}
\]

\emph{Step 4: separating the parts.} Let $t\in\mathbf A_E^\times$ have
$|t|>1$, for example $t$ equal to $2$ at one infinite place of $F$ and to $1$
at the other places, as in the proof of
Lemma~\ref{asm-theta-agree-split-unique}. Let $\lambda$ be the homomorphism of
Lemma~\ref{asm-theta-agree-split-inputs}(1)(c) for $H'$, and put
$m=\lambda(t,1)$, a Levi element for $(t,1)$. Then $m^{-1}=\lambda(t^{-1},1)$
is a Levi element for $(t^{-1},1)$, so $mum^{-1}\in N_{H'}$ for every
$u\in N_{H'}$ by Lemma~\ref{asm-theta-agree-split-inputs}(1)(g). Let
$u\in N_{H'}$ and $h\in U_W$. By ($*$) for $mum^{-1}$ and $mh$,
\[
f_\alpha(muh)+f_\beta(muh)=f_\alpha(mh)+f_\beta(mh).
\]
Put $x=\alpha(t)|t|^{a'/2}$ and $y=\beta(t)|t|^{a/2}$. By the laws of
$f_\alpha$ and $f_\beta$, this reads
$x\,f_\alpha(uh)+y\,f_\beta(uh)=x\,f_\alpha(h)+y\,f_\beta(h)$. Subtracting $y$
times ($*$) gives
\[
(x-y)\bigl(f_\alpha(uh)-f_\alpha(h)\bigr)=0.
\]
As $|\alpha(t)|=|\beta(t)|=1$, $|x|=|t|^{a'/2}$ and $|y|=|t|^{a/2}$, and these
differ since $|t|>1$ and $a'<a$, $G$ being proper. So $x\neq y$, and
$f_\alpha(uh)=f_\alpha(h)$. Then $f_\beta(uh)=f_\beta(h)$ by ($*$).
\end{proof}
% Lean: ThKCT.splitInv2_of UnipEx StabDec ThetaBasic TriplePair GrpTop EpsOf
% (SmallCones/SplitInv.lean). Step 1 is I_rat (ThetaBasic, polData_le); Step 2 is
% split_uniq; Step 3 is hc, with L2.integral_right and the countability from GrpTop;
% Step 4 is the rest, with the Levi homomorphism and the t with 1 < |t| from UnipEx, the
% conjugation from StabDec, and law_ne for x /= y.

\begin{definition}\label{asm-theta-agree-local-transfer}
Let $E$ and $F$ be as in Definition~\ref{asm-theta-agree-models}, let $B$ be a
theta base for $E$, and use the notation of
Definition~\ref{asm-theta-agree-local-cones}.
\begin{enumerate}
\item Let $M$ be a model for $(X,W)$ with vacua $\mathrm{vac}_v$, and let
  $S=(S_v)_v$ consist of subspaces $S_v\subset\mathcal S(F_v^{I_M})$, one for
  each finite place $v$. A family $\phi=(\phi_v)_v$ of functions on the
  $F_v^{I_M}$ is \emph{in $S$} if $\phi_v\in S_v$ for every $v$ and
  $\phi_v=\mathrm{vac}_v$ for almost all $v$. Write
  $f\otimes\phi=f\otimes\bigotimes_v\phi_v$, and $\phi[v_0\mapsto x]$ for the
  family $\phi$ with $\phi_{v_0}$ replaced by $x$. Let
  $S^{\mathrm{SB}}=(\mathcal S(F_v^{I_M}))_v$, and let $S^\circ(M)$ have
  $S^\circ_v(M)=\mathcal S(F_v^{I_M})$ for $v\neq v_1$ and
  $S^\circ_{v_1}(M)$ the span of the Schwartz-Bruhat functions of full rank
  at $\tau_1$ (Definition~\ref{asm-theta-agree-models}(4)).
\item A \emph{side} $(M,\omega,\pi,p,S)$ consists of a model $M$ for some
  $(X,W)$, a representation $\omega$ of $U_X\times U_W$ on functions on
  $\mathbf A^{I_M}$, linear maps $\pi$ from $\mathcal S(F_\infty^{I_M})$ and
  $p_v$ from the functions on $F_v^{I_M}$ to complex vector spaces, and a
  family $S$ as in (1); its theta integral is $I_{M,\omega}$. A matched
  quadruple $(B,G,\mathcal X,\omega,\omega',R,R_f)$ with real maps
  $(\Pi,p,p')$ and finite maps $(\Pi_v,p_v,p'_v)_v$ has the \emph{first side}
  $(M,\omega,p,(p_v)_v,S^{\mathrm{SB}})$, the \emph{full-rank side}
  $(M',\omega',p',(p'_v)_v,S^\circ(M'))$ and, if $q=0$, the \emph{second
  side} $(M',\omega',p',(p'_v)_v,S^{\mathrm{SB}})$.
\item Let $(M,\omega,\pi,p,S)$ and $(M_1,\omega_1,\pi_1,p_1,S_1)$ be sides,
  with $M$ for $(X,W)$ and $M_1$ for $(X_1,W_1)$ and theta integrals $I$ and
  $I_1$, let $H'$ be an adapted plane in $W$ over $W_1$, $\chi$ a Hecke
  character of $E$ and $r\in\mathbf R$. A \emph{transfer} along $(H',\chi,r)$
  from the first side to the second is a triple $(\mathrm{Ev},L,\Lambda)$ of
  a map $\mathrm{Ev}$ from functions on $\mathbf A^{I_M}$ to functions on
  $\mathbf A^{I_{M_1}}$, a continuous linear map
  $L\colon\mathcal S(F_\infty^{I_M})\to\mathcal S(F_\infty^{I_{M_1}})$ and
  linear maps $\Lambda_v$ from functions on $F_v^{I_M}$ to functions on
  $F_v^{I_{M_1}}$, such that:
  \begin{enumerate}
  \item[(a)] $L(\operatorname{Pol}(M))\subset\operatorname{Pol}(M_1)$, and
    $\Lambda_v(S_v)\subset S_{1,v}$ for every $v$;
  \item[(b)] for every family $\phi$ in $S$ and every finite place $v_0$
    there are $c\neq0$ and a family $\psi$ in $S_1$ such that every
    $\Lambda_v\phi_v$ is a multiple of $\psi_v$, and
    $\mathrm{Ev}(f\otimes\phi[v_0\mapsto x])=c\,L(f)\otimes\psi[v_0\mapsto\Lambda_{v_0}x]$
    for all $f\in\mathcal S(F_\infty^{I_M})$ and $x\in S_{v_0}$;
  \item[(c)] if $f\in\operatorname{Pol}(M)$, $\phi$ is in $S$ and
    $I(f\otimes\phi)=0$, then $I_1(\mathrm{Ev}(f\otimes\phi))=0$;
  \item[(d)] if $f\in\operatorname{Pol}(M)$, $\phi$ is in $S$, $m$ is a Levi
    element for $(t,h_1)$, $u\in N_{H'}$ and $h'\in U_{W_1}$, then
    $I_1\bigl(\mathrm{Ev}(\omega(1,um)(f\otimes\phi))\bigr)(h')=\chi(t)|t|^r\,I_1(\mathrm{Ev}(f\otimes\phi))(h'h_1)$.
  \end{enumerate}
\item The side $(M_1,\omega_1,\pi_1,p_1,S_1)$ is \emph{injective} if
  $I_1(f\otimes\phi)\neq0$ whenever $f\in\operatorname{Pol}(M_1)$, $\phi$ is
  in $S_1$, $\pi_1(f)\neq0$ and $p_{1,v}(\phi_v)\neq0$ for every $v$. With
  $L$ and $\Lambda$ as in (3), a \emph{test vector} for $(L,\Lambda)$
  relative to $S$, $\pi_1$ and $p_1$ is a pair of an $f\in\operatorname{Pol}(M)$
  and a family $\phi$ in $S$ with $\pi_1(L(f))\neq0$ and
  $p_{1,v}(\Lambda_v\phi_v)\neq0$ for every $v$.
\end{enumerate}
\end{definition}
\leanref{Development/WeilBound/ThetaLifts/LocalTheory/ISmulEqZero.lean}{PositivityTheory.ThKG.Loc.admitted}
\leanref{Development/WeilBound/ThetaLifts/LocalTheory/AdmissiblePureTensorSchwartz.lean}{PositivityTheory.ThKG.Loc.SS}
\leanref{Development/WeilBound/ThetaLifts/LocalTheory/ISmulEqZero.lean}{PositivityTheory.ThKG.Loc.Br}
\leanref{Development/WeilBound/ThetaLifts/LocalTheory/ISmulEqZero.lean}{PositivityTheory.ThKG.Loc.Inj}
\leanref{Development/WeilBound/ThetaLifts/LocalTheory/ISmulEqZero.lean}{PositivityTheory.ThKG.Loc.Pt}

\begin{lemma}\label{asm-theta-agree-local-transfer-criteria}
Let $(M,\omega,\pi,p,S)$ and $(M_1,\omega_1,\pi_1,p_1,S_1)$ be sides for a
theta base $B$ of a CM field $E$ (Definition~\ref{asm-theta-agree-local-transfer}),
with $M$ for $(X,W)$ and $M_1$ for $(X_1,W_1)$, theta integrals $I$ and $I_1$
and vacua $\mathrm{vac}_v$ and $\mathrm{vac}_{1,v}$. Let $H'$ be an adapted
plane in $W$ over $W_1$, $\chi$ a Hecke character of $E$, $r\in\mathbf R$, and
let $\mathrm{Ev}$, $L$ and $\Lambda=(\Lambda_v)_v$ be maps as in
Definition~\ref{asm-theta-agree-local-transfer}(3).
\begin{enumerate}
\item Suppose that there is a Levi element for $(1,h_1)$ for every
  $h_1\in U_{W_1}$, that conditions (a) and (b) of a transfer hold, and that
  there is $C\neq0$ with the following property: for every
  $f\in\operatorname{Pol}(M)$ and every family $\phi$ in $S$ there is a
  function $T$ on $U_W$ with the law $(\chi,r)$ along $H'$
  (Definition~\ref{asm-theta-agree-planes}(2)), with $T(uh)=T(h)$ for
  $u\in N_{H'}$ and $h\in U_W$, with $T=0$ if $I(f\otimes\phi)=0$, and with
  $T(mh)=C\,I_1\bigl(\mathrm{Ev}(\omega(1,h)(f\otimes\phi))\bigr)(h_1)$ for every
  Levi element $m$ for $(1,h_1)$ and every $h\in U_W$. Then
  $(\mathrm{Ev},L,\Lambda)$ is a transfer along $(H',\chi,r)$.
\item Suppose that $S_v\subset\mathcal S(F_v^{I_M})$ and
  $\Lambda_v(S_v)\subset S_{1,v}$ for every $v$, and that $\mathrm{Ev}$ is an
  evaluation map (Definition~\ref{asm-theta-agree-planes}(5)) with components
  $L$ and $\Lambda_v$. Then condition (b) of a transfer holds.
\item Let $T_0$ be a finite set of finite places, $k\neq0$, and $u_v\neq0$
  numbers with $\Lambda_v(\mathrm{vac}_v)=u_v\,\mathrm{vac}_{1,v}$ for
  $v\notin T_0$. Suppose that $\Lambda_v(S_v)\subset S_{1,v}$ for every $v$,
  that $\mathrm{vac}_{1,v}\in S_{1,v}$ whenever $\mathrm{vac}_v\in S_v$, and that
  \[
  \Bigl(\prod_{v\in T\setminus T_0}u_v\Bigr)\mathrm{Ev}\Bigl(f\otimes\bigotimes_v\phi_v\Bigr)=k\,L(f)\otimes\bigotimes_v\psi_v
  \]
  whenever $T\supset T_0$ is a finite set of finite places,
  $f\in\mathcal S(F_\infty^{I_M})$ and $\phi_v\in\mathcal S(F_v^{I_M})$ with
  $\phi_v=\mathrm{vac}_v$ for $v\notin T$, where $\psi_v=\Lambda_v\phi_v$ for
  $v\in T$ and $\psi_v=\mathrm{vac}_{1,v}$ for $v\notin T$. Then condition (b)
  of a transfer holds.
\end{enumerate}
\end{lemma}
\leanref{Development/WeilBound/ThetaLifts/LocalTheory/MemPolData.lean}{PositivityTheory.ThKG.Loc.br_of}
\leanref{Development/WeilBound/ThetaLifts/LocalTheory/MemPolData.lean}{PositivityTheory.ThKG.Loc.fac_ev}
\leanref{Development/WeilBound/ThetaLifts/LocalTheory/MemPolData.lean}{PositivityTheory.ThKG.Loc.fac_cone}
\leanref{Development/WeilBound/ThetaLifts/LocalTheory/UatClosure.lean}{PositivityTheory.ThKG.Loc.law_par}
\begin{proof}
(1) Conditions (a) and (b) hold by assumption. For (c), let
$f\in\operatorname{Pol}(M)$ and $\phi$ in $S$ with $I(f\otimes\phi)=0$, and let
$T$ be as in the assumption; then $T=0$. Let $h_1\in U_{W_1}$ and let $m$ be a
Levi element for $(1,h_1)$. As $\omega(1,1)$ is the identity,
$0=T(m)=C\,I_1(\mathrm{Ev}(f\otimes\phi))(h_1)$, and $C\neq0$.

For (d), let $f$, $\phi$ and $T$ be as before, $m$ a Levi element for
$(t,h_1)$, $u\in N_{H'}$ and $h'\in U_{W_1}$. Choose Levi elements $m'$ for
$(1,h')$ and $l$ for $(1,h'h_1)$. Then
$m'um=(m'um'^{-1})(m'ml^{-1})l$. Conjugation by $m'$ preserves $N_{H'}$, as
$m'$ fixes $x$ and maps $\iota z$ to $\iota h'z$; and $m'ml^{-1}$ is a Levi
element for $(t,1)$, by the defining relations. By the invariance of $T$ and
its law,
\[
T(m'um)=T\bigl((m'ml^{-1})l\bigr)=\chi(t)|t|^r\,T(l).
\]
By the last property of $T$, with $h=um$ and with $h=1$,
$T(m'\cdot um)=C\,I_1(\mathrm{Ev}(\omega(1,um)(f\otimes\phi)))(h')$ and
$T(l)=C\,I_1(\mathrm{Ev}(f\otimes\phi))(h'h_1)$. Dividing by $C$ gives (d).

(2) Let $\phi$ be in $S$ and $v_0$ a finite place. Take $c=1$ and
$\psi_v=\Lambda_v\phi_v$. Then $\psi$ is in $S_1$, as $\Lambda_v$ maps $S_v$
into $S_{1,v}$ and $\Lambda_v(\mathrm{vac}_v)=\mathrm{vac}_{1,v}$. For
$x\in S_{v_0}$ the family $\phi[v_0\mapsto x]$ consists of Schwartz-Bruhat
functions, almost all equal to the vacua, so by the factorization of an
evaluation map
$\mathrm{Ev}(f\otimes\phi[v_0\mapsto x])=L(f)\otimes\psi[v_0\mapsto\Lambda_{v_0}x]$.

(3) Let $\phi$ be in $S$ and $v_0$ a finite place. Let $T$ be the union of
$T_0$, $\{v_0\}$ and the finite set of $v$ with $\phi_v\neq\mathrm{vac}_v$, put
$u=\prod_{v\in T\setminus T_0}u_v\neq0$ and $c=u^{-1}k\neq0$, and let
$\psi_v=\Lambda_v\phi_v$ for $v\in T$ and $\psi_v=\mathrm{vac}_{1,v}$ for
$v\notin T$. Then $\psi$ is in $S_1$: for $v\in T$,
$\Lambda_v\phi_v\in S_{1,v}$, and for $v\notin T$, $\phi_v=\mathrm{vac}_v\in S_v$,
so $\mathrm{vac}_{1,v}\in S_{1,v}$. Each $\Lambda_v\phi_v$ is a multiple of
$\psi_v$: with factor $1$ for $v\in T$, and $u_v$ for $v\notin T$. For
$x\in S_{v_0}$ the family $\phi[v_0\mapsto x]$ equals the vacua off $T$, as
$v_0\in T$, and the assumption gives
$u\,\mathrm{Ev}(f\otimes\phi[v_0\mapsto x])=k\,L(f)\otimes\psi[v_0\mapsto\Lambda_{v_0}x]$.
\end{proof}

\begin{lemma}\label{asm-theta-agree-local-evaluation-pol}
Let $E$ and $F$ be as in Definition~\ref{asm-theta-agree-models}, let $M$ and
$M_1$ be regular models (Definition~\ref{asm-theta-agree-models}(3)), and let
$e$ be an evaluation map for some data $(M,M_1,\iota_X,\iota_W,\mathcal J)$
(Definition~\ref{asm-theta-agree-planes}(5)), with components $e_\infty$ and
$e_v$, that carries $\operatorname{PD}(M)$ onto $\operatorname{PD}(M_1)$. Then
$e_\infty(\operatorname{Pol}(M))=\operatorname{Pol}(M_1)$.
\end{lemma}
\leanref{Development/WeilBound/ThetaLifts/LocalTheory/MemPolData.lean}{PositivityTheory.ThKG.Loc.ev_pol}
\leanref{Development/WeilBound/ThetaLifts/LocalTheory/AdmissiblePureTensorSchwartz.lean}{PositivityTheory.ThKG.Loc.arch_slice}
\begin{proof}
Write $\mathrm{vac}_v$ and $\mathrm{vac}_{1,v}$ for the vacua of $M$ and
$M_1$, and let $I_1=I_{M_1}$. We first show: if $N\subset\mathcal S(F_\infty^{I_1})$
is a subspace and $g\otimes\bigotimes_v\mathrm{vac}_{1,v}$ lies in the span
of the functions $f\otimes\bigotimes_v\psi_v$ with $f\in N$, all $\psi_v$
Schwartz-Bruhat and $\psi_v=\mathrm{vac}_{1,v}$ for almost all $v$, then
$g\in N$. Each $\mathrm{vac}_{1,v}$ is nonzero, and it is the indicator
function of $\mathcal O_v^{I_1}$ for almost all $v$, as $M_1$ is regular. So
there is a finite adele $x_0$ with
$c_0=\prod_v\mathrm{vac}_{1,v}(x_{0,v})\neq0$, where $x_{0,v}=0$ whenever
$\mathrm{vac}_{1,v}$ is that indicator function. For every $\Phi$ in the span,
the function $z\mapsto\Phi(z,x_0)$ on $F_\infty^{I_1}$ lies in $N$: for
$\Phi=f\otimes\bigotimes_v\psi_v$ it is $(\prod_v\psi_v(x_{0,v}))f$, a finite
product. Applied to $\Phi=g\otimes\bigotimes_v\mathrm{vac}_{1,v}$ this gives
$c_0g\in N$, so $g\in N$.

Let $f\in\operatorname{Pol}(M)$. Then $f\otimes\bigotimes_v\mathrm{vac}_v\in\operatorname{PD}(M)$,
so $e_\infty(f)\otimes\bigotimes_v\mathrm{vac}_{1,v}=e(f\otimes\bigotimes_v\mathrm{vac}_v)$
lies in $\operatorname{PD}(M_1)$, which is spanned by the
$f_1\otimes\bigotimes_v\psi_v$ with $f_1\in\operatorname{Pol}(M_1)$. With
$N=\operatorname{Pol}(M_1)$ this gives $e_\infty(f)\in\operatorname{Pol}(M_1)$.
Let $f_1\in\operatorname{Pol}(M_1)$. Then $f_1\otimes\bigotimes_v\mathrm{vac}_{1,v}$
lies in $\operatorname{PD}(M_1)=e(\operatorname{PD}(M))$, the span of the
$e(f\otimes\bigotimes_v\psi_v)=e_\infty(f)\otimes\bigotimes_ve_v\psi_v$ with
$f\in\operatorname{Pol}(M)$, and $e_v\psi_v=\mathrm{vac}_{1,v}$ for almost all
$v$. With $N=e_\infty(\operatorname{Pol}(M))$ this gives
$f_1\in e_\infty(\operatorname{Pol}(M))$.
\end{proof}

\begin{lemma}\label{asm-theta-agree-local-maps-sides}
Let $(B,G,\mathcal X,\omega,\omega',R,R_f)$ be a matched quadruple for a theta
base $B$ of a CM field $E$, with real maps $(\Pi,p,p')$ and finite maps
$(\Pi_v,p_v,p'_v)_v$. Let $(M_\Sigma,\omega_\Sigma,\pi,\varpi,S)$ be its first
side, its full-rank side or, if $q=0$, its second side
(Definition~\ref{asm-theta-agree-local-transfer}(2)), with theta integral
$I_\Sigma$, and with families in $S$ as in
Definition~\ref{asm-theta-agree-local-transfer}(1).
\begin{enumerate}
\item Every $S_v$ consists of Schwartz-Bruhat functions, $\varpi_v(S_v)=\Pi_v$,
  and $\pi(\operatorname{Pol}(M_\Sigma))=\Pi$.
\item For every finite place $v$, $h\in U_W^{(v)}$ and $x\in S_v$ there is a
  Schwartz-Bruhat $x_1$ with $\omega_\Sigma(1,h)$ acting at $v$ by
  $x\mapsto x_1$, and every such $x_1$ lies in $S_v$.
\item There is a space of Schwartz-Bruhat functions that converge for
  $(M_\Sigma,\omega_\Sigma)$ and that contains $f\otimes\phi$ for every
  Schwartz function $f$ and every family $\phi$ in $S$; and for every such
  $\phi$ and every $h\in U_W$, $f\mapsto I_\Sigma(f\otimes\phi)(h)$ is
  continuous.
\item If $f\in\operatorname{Pol}(M_\Sigma)$, $\phi$ is a family in $S$, and
  $\pi(f)=0$ or $\varpi_v(\phi_v)=0$ for some $v$, then $I_\Sigma(f\otimes\phi)=0$.
\item Let $v$ be a finite place, $h\in U_W^{(v)}$, $\phi,\phi_1$ Schwartz-Bruhat
  functions on $F_v^{I_0}$ and $x,x_1$ Schwartz-Bruhat functions of the side
  at $v$, such that $\omega(1,h)$ acts at $v$ by $\phi\mapsto\phi_1$,
  $\omega_\Sigma(1,h)$ acts at $v$ by $x\mapsto x_1$, and
  $p_v(\phi)=\varpi_v(x)$. Then $p_v(\phi_1)=\varpi_v(x_1)$. Likewise, let
  $f\in\operatorname{Pol}(M)$ and $x\in\operatorname{Pol}(M_\Sigma)$ with
  $p(f)=\pi(x)$, and let $f\mapsto f_1$ and $x\mapsto x_1$ be the actions at
  $\infty$ to first order of $t\mapsto\omega(1,\gamma(t))$ and
  $t\mapsto\omega_\Sigma(1,\gamma(t))$ for a line $\gamma$ in $U_W$, or the
  actions at $\infty$ of $\omega(1,k)$ and $\omega_\Sigma(1,k)$ for some
  $k\in K_\infty(W)$. Then $p(f_1)=\pi(x_1)$.
\item The statements of (5) hold with $(M_\Sigma,\omega_\Sigma,\pi,\varpi)$ in
  place of $(M,\omega,p,(p_v)_v)$.
\end{enumerate}
\end{lemma}
\leanref{Development/WeilBound/ThetaLifts/LocalTheory/ISmulEqZero.lean}{PositivityTheory.ThKG.Loc.Side}
\leanref{Development/WeilBound/ThetaLifts/LocalTheory/Cross.lean}{PositivityTheory.ThKG.Loc.Cross}
\leanref{Development/WeilBound/ThetaLifts/LocalTheory/AdmissiblePureTensorSchwartz.lean}{PositivityTheory.ThKG.Loc.side_of}
\leanref{Development/WeilBound/ThetaLifts/LocalTheory/AdmissiblePureTensorSchwartz.lean}{PositivityTheory.ThKG.Loc.sideV}
\leanref{Development/WeilBound/ThetaLifts/LocalTheory/AdmissiblePureTensorSchwartz.lean}{PositivityTheory.ThKG.Loc.sideSo}
\leanref{Development/WeilBound/ThetaLifts/LocalTheory/AdmissiblePureTensorSchwartz.lean}{PositivityTheory.ThKG.Loc.side_condition_vanishing}
\begin{proof}
For the first side, $(M_\Sigma,\omega_\Sigma,\pi,\varpi)=(M,\omega,p,(p_v)_v)$ and
all $S_v$ are the Schwartz-Bruhat functions; call the second space, with
$(M',\omega',p',(p'_v)_v)$, the \emph{other space}. For the full-rank and
second sides, $(M_\Sigma,\omega_\Sigma,\pi,\varpi)=(M',\omega',p',(p'_v)_v)$, and
the other space is the first one. Write $(M_o,\omega_o,\pi_o,\varpi_o)$ for the
other space. Pure tensors with Schwartz-Bruhat components almost all equal to
the vacua are Schwartz-Bruhat functions, as the models are regular.

(1) This is Definition~\ref{asm-theta-agree-matched}(2)(a) and~(3), except for
$S_{v_1}$ in the full-rank side, where $\varpi_{v_1}(S_{v_1})=\Pi_{v_1}$ was shown
in the proof of Lemma~\ref{asm-theta-agree-local-injective} (Part (2), last
case).

(2) A function $x_1$ exists by local operators at $v$
(Lemma~\ref{asm-theta-agree-finite-inputs}(2)). It lies in $S_v$ trivially
unless $S_v$ is the span $S_{v_1}$ of the functions of full rank at $\tau_1$,
and for that case see the proof of Lemma~\ref{asm-theta-agree-local-injective}
(Part (2), last case).

(3) For the first and second sides take all Schwartz-Bruhat functions, and for
the full-rank side $\mathcal S^\circ(M')$. Convergence and continuity follow, as
in the proof of Lemma~\ref{asm-theta-agree-local-injective}, from
Lemma~\ref{asm-theta-agree-convergence} for the first and second sides and from
Lemmas~\ref{asm-theta-agree-decay} and~\ref{asm-theta-agree-continuity} for the
full-rank side.

(4) This is Lemma~\ref{asm-theta-agree-local-injective-kills}(1) for the first
side, Lemma~\ref{asm-theta-agree-local-injective-kills-zero} for the second side
and Lemma~\ref{asm-theta-agree-local-injective-kills}(2) for the full-rank side.

(6) Let $\omega_\Sigma(1,h)$ act at $v$ by $a\mapsto a'$ and by $b\mapsto b'$,
all Schwartz-Bruhat, with $\varpi_v(a)=\varpi_v(b)$. Choose a Schwartz-Bruhat
$z$ with $\varpi_{o,v}(z)=\varpi_v(a)$, and by local operators a
Schwartz-Bruhat $z'$ with $\omega_o(1,h)$ acting at $v$ by $z\mapsto z'$. The
compatibility of the finite maps with $\omega(1,h)$ and $\omega'(1,h)$ in
Definition~\ref{asm-theta-agree-matched}(3), which is symmetric in the two
spaces, gives $\varpi_v(a')=\varpi_{o,v}(z')=\varpi_v(b')$. At $\infty$ the
argument is the same, with $z\in\operatorname{Pol}(M_o)$ and $\pi_o(z)=\pi(a)$
from Definition~\ref{asm-theta-agree-matched}(2)(a), the action $z\mapsto z'$
from Lemma~\ref{asm-theta-agree-rechoice-linear}(1) and~(2), and the
compatibility (2)(b) there.

(5) For the first side this is (6). For the other sides it is the
compatibility of Definition~\ref{asm-theta-agree-matched}(3), respectively
(2)(b), for the first and second spaces.
\end{proof}
% Lean: Side (ISmulEqZero.lean) records (1) (le, onto, ontoI), (2) (stab), (3) (cls, cont),
% (4) (kill) and (6) (wdF, wdL, wdK), with op, lie and uniq; (5) is the Cross crossM
% (fin, lie, cpt) and its symm, or Side.cross for the first side. side_of builds a Side from
% these inputs (wdF through c.fin and ActsEx); sideV, sideSo and side_condition_vanishing
% supply them with Kills, Kills0, WeilConv, KDecay, SoCont and KCoinv (so_slice for (2)).

\begin{lemma}\label{asm-theta-agree-local-maps}
Let $B$ be a theta base of a CM field $E$, let
$Q=(B,G,\mathcal X,\omega,\omega',R,R_f)$ and
$Q_1=(B,G_1,\mathcal X,\omega_1,\omega'_1,R_1,R_{f,1})$ be matched quadruples
with real maps $(\Pi,p,p')$ and $(\Pi_1,p_1,p'_1)$ and finite maps
$(\Pi_v,p_v,p'_v)_v$ and $(\Pi_{1,v},p_{1,v},p'_{1,v})_v$, let $H'$ be an
adapted plane in $W$ over $W_1$, the skew-Hermitian spaces of $G$ and $G_1$, at
some positions, let $\chi$ be a Hecke character of $E$ and $r\in\mathbf R$.
Let $(M_\Sigma,\omega_\Sigma,\pi,\varpi,S)$ be the first side, the full-rank
side or the second side of $Q$, and $(M_{\Sigma,1},\omega_{\Sigma,1},\pi_1,\varpi_1,S_1)$
the side of the same kind of $Q_1$
(Definition~\ref{asm-theta-agree-local-transfer}), with theta integral
$I_{\Sigma,1}$; here $M_{\Sigma,1}$ is a model for $(Y_1,W_1)$.
\begin{enumerate}
\item Let $L\colon\mathcal S(F_\infty^{I_{M_\Sigma}})\to\mathcal S(F_\infty^{I_{M_{\Sigma,1}}})$
  be a continuous linear map with
  $L(\operatorname{Pol}(M_\Sigma))\subset\operatorname{Pol}(M_{\Sigma,1})$,
  such that $l\circ L\neq0$ for every nonzero continuous linear form $l$ on
  $\mathcal S(F_\infty^{I_{M_{\Sigma,1}}})$ with $l(f_1)=l(f)$ whenever
  $\omega_{\Sigma,1}(g,1)$ acts at $\infty$ by $f\mapsto f_1$ for some
  $g\in U_{Y_1}^{(\infty)}$. If $I_{\Sigma,1}(\Phi)\neq0$ for some $\Phi$ in
  the span of the $f\otimes\phi$ with $\phi$ in $S_1$, then $\pi_1(L(f))\neq0$
  for some $f\in\operatorname{Pol}(M_\Sigma)$.
\item Let $(\mathrm{Ev},L,\Lambda)$ be a transfer along $(H',\chi,r)$ from
  the side of $Q$ to that of $Q_1$, such that the side of $Q_1$ is injective
  and $(L,\Lambda)$ has a test vector relative to $S$, $\pi_1$ and $\varpi_1$.
  Then there are linear maps $l_\infty\colon\Pi\to\Pi_1$ and
  $l_v\colon\Pi_v\to\Pi_{1,v}$ with $l_\infty(\pi(f))=\pi_1(L(f))$ for
  $f\in\operatorname{Pol}(M_\Sigma)$ and $l_v(\varpi_v(x))=\varpi_{1,v}(\Lambda_vx)$
  for $x\in S_v$, such that $l_\infty$ is equivariant at $\infty$ and every
  $l_v$ is equivariant at $v$ for $(H',\chi,r)$
  (Definition~\ref{asm-theta-agree-local-equivariant}).
\end{enumerate}
\end{lemma}
\leanref{Development/WeilBound/ThetaLifts/LocalTheory/Cross.lean}{PositivityTheory.ThKG.Loc.lams}
\leanref{Development/WeilBound/ThetaLifts/LocalTheory/MemPolData.lean}{PositivityTheory.ThKG.Loc.bump_pt}
\leanref{Development/WeilBound/ThetaLifts/LocalTheory/MemPolData.lean}{PositivityTheory.ThKG.Loc.nz_pure}
\leanref{Development/WeilBound/ThetaLifts/LocalTheory/AdmissiblePureTensorSchwartz.lean}{PositivityTheory.ThKG.Loc.sideV}
\leanref{Development/WeilBound/ThetaLifts/LocalTheory/AdmissiblePureTensorSchwartz.lean}{PositivityTheory.ThKG.Loc.sideSo}
\leanref{Development/WeilBound/ThetaLifts/LocalTheory/AdmissiblePureTensorSchwartz.lean}{PositivityTheory.ThKG.Loc.side_condition_vanishing}
\begin{proof}
Write $I_1=I_{\Sigma,1}$. Lemma~\ref{asm-theta-agree-local-maps-sides} applies
to the side of $Q$ and to that of $Q_1$; we cite it as ``sides''. For a family
$\chi$, a finite place $v$ and a function $x$ at $v$, write $\chi[v\mapsto x]$ for
$\chi$ with $\chi_v$ replaced by $x$. Recall that
$I(\omega(1,h)\Phi)(h')=I(\Phi)(h'h)$ for every theta integral.

\emph{Part (1).} By sides~(3) for $Q_1$, $I_1$ is linear on a space containing
the $f\otimes\phi$ with $\phi$ in $S_1$. If all these had theta integral $0$,
so would every $\Phi$ in their span. So there are a Schwartz function $f_0$, a
family $\phi$ in $S_1$ and $h\in U_{W_1}$ with $I_1(f_0\otimes\phi)(h)\neq0$. By
sides~(3), $l(f)=I_1(f\otimes\phi)(h)$ is a nonzero continuous linear form. It
is invariant: let $g\in U_{Y_1}^{(\infty)}$ and $\omega_{\Sigma,1}(g,1)$ act at
$\infty$ by $f\mapsto f_1$, and let $T$ be an archimedean operator of
$\omega_{\Sigma,1}$ at $(g,1)$
(Lemma~\ref{asm-theta-agree-rechoice-operators}(1)). By
Lemma~\ref{asm-theta-agree-real-weak}(3), $f_1=T(f)$, so
$\omega_{\Sigma,1}(g,1)(f\otimes\phi)=f_1\otimes\phi$, and as $f\otimes\phi$ is a
Schwartz-Bruhat function, Lemma~\ref{asm-theta-agree-anisotropic-invariance}
gives $l(f_1)=l(f)$. By hypothesis $l(L(f))\neq0$ for some $f$. Suppose that
$\pi_1(L(g))=0$ for every $g\in\operatorname{Pol}(M_\Sigma)$. As
$L(g)\in\operatorname{Pol}(M_{\Sigma,1})$, sides~(4) for $Q_1$ gives
$l(L(g))=0$. So the continuous linear form $l\circ L$ vanishes on
$\operatorname{Pol}(M_\Sigma)$, which is dense by
Lemma~\ref{asm-theta-agree-density}, and $l\circ L=0$, a contradiction.

\emph{Part (2).} Let $(f_0,\phi)$ be the test vector: $f_0\in\operatorname{Pol}(M_\Sigma)$,
$\phi$ is a family in $S$, $\pi_1(L(f_0))\neq0$ and
$\varpi_{1,v}(\Lambda_v\phi_v)\neq0$ for every $v$.

\emph{Step 1: test vectors.} Let $v$ be a finite place. Condition (b) of a
transfer, for $\phi$ and $v_0=v$, gives $c\neq0$ and a family $\psi$ in $S_1$
such that every $\Lambda_u\phi_u$ is a multiple of $\psi_u$ and
$\mathrm{Ev}(f\otimes\phi[v\mapsto x])=c\,L(f)\otimes\psi[v\mapsto\Lambda_vx]$
for all $f$ and $x\in S_v$. For $u\neq v$, $\varpi_{1,u}(\psi_u)\neq0$, since a
multiple of $\psi_u$ has nonzero image. As the side of $Q_1$ is injective and
$\pi_1(L(f_0))\neq0$:
\begin{enumerate}
\item[(T$_v$)] if $y\in S_{1,v}$ and $I_1(L(f_0)\otimes\psi[v\mapsto y])=0$, then
  $\varpi_{1,v}(y)=0$.
\end{enumerate}
For $v=v_1$ put $\psi'=\psi[v_1\mapsto\Lambda_{v_1}\phi_{v_1}]$, a family in
$S_1$ as $\Lambda_{v_1}(S_{v_1})\subset S_{1,v_1}$. Then
$\mathrm{Ev}(f\otimes\phi)=c'\,L(f)\otimes\psi'$ for every $f$, where $c'\neq0$
is the constant $c$ for $v_1$, and all $\varpi_{1,u}(\psi'_u)\neq0$; so by
injectivity:
\begin{enumerate}
\item[(T$_\infty$)] if $g\in\operatorname{Pol}(M_{\Sigma,1})$ and $I_1(g\otimes\psi')=0$,
  then $\pi_1(g)=0$.
\end{enumerate}

\emph{Step 2: kernels.} If $f\in\operatorname{Pol}(M_\Sigma)$ and $\pi(f)=0$, then
$I_\Sigma(f\otimes\phi)=0$ by sides~(4); by condition (c) of a transfer,
$0=I_1(\mathrm{Ev}(f\otimes\phi))=c'I_1(L(f)\otimes\psi')$, and (T$_\infty$) gives
$\pi_1(L(f))=0$. If $x\in S_v$ and $\varpi_v(x)=0$, then
$I_\Sigma(f_0\otimes\phi[v\mapsto x])=0$ by sides~(4), so
$0=c\,I_1(L(f_0)\otimes\psi[v\mapsto\Lambda_vx])$, and (T$_v$) gives
$\varpi_{1,v}(\Lambda_vx)=0$.

\emph{Step 3: the maps.} By sides~(1) there is a linear map
$s\colon\Pi\to\operatorname{Pol}(M_\Sigma)$ with $\pi\circ s=1$; put
$l_\infty=\pi_1\circ L\circ s$. For $f\in\operatorname{Pol}(M_\Sigma)$,
$\pi(s(\pi(f))-f)=0$, so $l_\infty(\pi(f))=\pi_1(L(f))$ by Step 2. In the same
way, with $s_v\colon\Pi_v\to S_v$ and $\varpi_v\circ s_v=1$,
$l_v=\varpi_{1,v}\circ\Lambda_v\circ s_v$ has
$l_v(\varpi_v(x))=\varpi_{1,v}(\Lambda_vx)$ for $x\in S_v$.

\emph{Step 4: $l_v$ is equivariant at $v$.} Let $m$ be a Levi element for
$(t,h_1)$ and $u\in N_{H'}$, both at $v$, and let $\eta,\eta',\phi_1,\phi'_1$
be Schwartz-Bruhat with $\omega(1,um)$ acting at $v$ by
$\eta\mapsto\eta'$, $\omega_1(1,h_1)$ by $\phi_1\mapsto\phi'_1$,
and $l_v(p_v(\eta))=p_{1,v}(\phi_1)$; put $a=\chi(t)|t|^r$. As $m-1$ is at
$v$ and $m\iota=\iota h_1$, $\iota(h_1-1)$ is at $v$, so $h_1$ is at $v$, the
columns of $\iota$ being linearly independent. By sides~(1) and~(2), choose
$x\in S_v$ with $\varpi_v(x)=p_v(\eta)$ and $x'\in S_v$ with
$\omega_\Sigma(1,um)$ acting at $v$ by $x\mapsto x'$; $y\in S_{1,v}$ with
$\varpi_{1,v}(y)=p_{1,v}(\phi_1)$ and $y'\in S_{1,v}$ with $\omega_{\Sigma,1}(1,h_1)$
acting at $v$ by $y\mapsto y'$; and $z\in S_{1,v}$ with $\omega_{\Sigma,1}(1,h_1)$
acting at $v$ by $\Lambda_vx\mapsto z$. By sides~(5),
$p_v(\eta')=\varpi_v(x')$ and $p_{1,v}(\phi'_1)=\varpi_{1,v}(y')$, so
$l_v(p_v(\eta'))=\varpi_{1,v}(\Lambda_vx')$. As
$\varpi_{1,v}(\Lambda_vx)=l_v(p_v(\eta))=\varpi_{1,v}(y)$, sides~(6) gives
$\varpi_{1,v}(z)=\varpi_{1,v}(y')$. With $f_0$, $\psi$ and $c$ of Step 1 for $v$,
and $\omega_\Sigma(1,um)(f_0\otimes\phi[v\mapsto x])=f_0\otimes\phi[v\mapsto x']$,
condition (d) of a transfer gives, for $h'\in U_{W_1}$,
\[
c\,I_1(L(f_0)\otimes\psi[v\mapsto\Lambda_vx'])(h')
=a\,c\,I_1(L(f_0)\otimes\psi[v\mapsto\Lambda_vx])(h'h_1)
=a\,c\,I_1(L(f_0)\otimes\psi[v\mapsto z])(h').
\]
So $I_1(L(f_0)\otimes\psi[v\mapsto\Lambda_vx'-az])=0$, and by (T$_v$),
$l_v(p_v(\eta'))=\varpi_{1,v}(\Lambda_vx')=a\,\varpi_{1,v}(z)=a\,p_{1,v}(\phi'_1)$.

\emph{Step 5: $l_\infty$ is equivariant at $\infty$.} For $h'\in U_{W_1}$ let
$l_{h'}(g)=I_1(g\otimes\psi')(h')$, a continuous linear form by sides~(3). Let
$x\in\operatorname{Pol}(M_\Sigma)$, let $m$ be a Levi element for $(t,h_1)$ and
$u\in N_{H'}$ with $um$ at $\infty$, and let $T$ be an archimedean operator of
$\omega_\Sigma$ at $(1,um)$. Then
$\omega_\Sigma(1,um)(x\otimes\phi)=T(x)\otimes\phi$, and condition (d) of a
transfer gives, after division by $c'$,
\[
l_{h'}(L(T(x)))=\chi(t)|t|^r\,l_{h'h_1}(L(x)).\tag{$*$}
\]
(a) Let $\gamma$ be a line with values in $N_{H'}$, $f\in\operatorname{Pol}(M)$,
and $t\mapsto\omega(1,\gamma(t))$ act at $\infty$ to first order by $f\mapsto f'$.
Choose $x\in\operatorname{Pol}(M_\Sigma)$ with $\pi(x)=p(f)$ and archimedean
operators $T_t$ of $\omega_\Sigma$ at $(1,\gamma(t))$. By
Lemma~\ref{asm-theta-agree-rechoice-lie} there is
$x'\in\operatorname{Pol}(M_\Sigma)$ such that $t\mapsto\omega_\Sigma(1,\gamma(t))$
acts at $\infty$ to first order by $x\mapsto x'$ and $t\mapsto\lambda(T_tx)$ has
derivative $\lambda(x')$ at $0$ for every continuous linear form $\lambda$. By
sides~(5), $l_\infty(p(f'))=\pi_1(L(x'))$. By ($*$) with $m=1$ and
$u=\gamma(t)$, $t\mapsto l_{h'}(L(T_tx))$ is constant, so $l_{h'}(L(x'))=0$ for
every $h'$, and $\pi_1(L(x'))=0$ by (T$_\infty$).

(b) Let $\gamma,\gamma_1,\tau$ and $c_0$ be as in
Definition~\ref{asm-theta-agree-local-equivariant}(3)(b), with
$c_0$ the derivative of $s\mapsto\chi(\tau(s))|\tau(s)|^r$ at $0$, and let
$f\mapsto f'$ and $f_1\mapsto f'_1$ be as there, with $l_\infty(p(f))=p_1(f_1)$.
As $\gamma(0)=1$ is a Levi element for $(\tau(0),1)$, $\tau(0)=1$. Choose $x$,
$T_t$ and $x'$ as in (a), so that $l_\infty(p(f'))=\pi_1(L(x'))$; choose
$y\in\operatorname{Pol}(M_{\Sigma,1})$ with $\pi_1(y)=p_1(f_1)$ and $y'$ with
$t\mapsto\omega_{\Sigma,1}(1,\gamma_1(t))$ acting at $\infty$ to first order by
$y\mapsto y'$ (Lemma~\ref{asm-theta-agree-rechoice-linear}(1)), so that
$p_1(f'_1)=\pi_1(y')$ by sides~(5) for $Q_1$; and archimedean operators
$T_{1,t}$ of $\omega_{\Sigma,1}$ at $(1,\gamma_1(t))$, with $g_1$ given by
Lemma~\ref{asm-theta-agree-rechoice-lie} for $L(x)$. By ($*$) with $u=1$ and
$m=\gamma(t)$, and as
$l_{h'\gamma_1(t)}(L(x))=l_{h'}(T_{1,t}L(x))$,
\[
l_{h'}(L(T_tx))=\chi(\tau(t))|\tau(t)|^r\,l_{h'}(T_{1,t}L(x)).
\]
At $t=0$ the last factor is $l_{h'}(L(x))$, as $\omega_{\Sigma,1}(1,1)$ is the
identity. Differentiating at $t=0$ gives
$l_{h'}(L(x'))=c_0\,l_{h'}(L(x))+l_{h'}(g_1)$ for every $h'$, and (T$_\infty$) gives
$\pi_1(L(x'))=c_0\pi_1(L(x))+\pi_1(g_1)$. Here
$\pi_1(L(x))=l_\infty(p(f))=\pi_1(y)$, and then $\pi_1(g_1)=\pi_1(y')$ by
sides~(6). So $l_\infty(p(f'))=c_0\,p_1(f_1)+p_1(f'_1)$.

(c) Let $k$, $k_1$, $t$, $f\mapsto f'$ and $f_1\mapsto f'_1$ be as in
Definition~\ref{asm-theta-agree-local-equivariant}(3)(c), with
$l_\infty(p(f))=p_1(f_1)$. Choose $x$ with $\pi(x)=p(f)$, an archimedean operator
$T$ of $\omega_\Sigma$ at $(1,k)$, $y$ with $\pi_1(y)=p_1(f_1)$ and an archimedean
operator $T_1$ of $\omega_{\Sigma,1}$ at $(1,k_1)$. By
Lemma~\ref{asm-theta-agree-rechoice-linear}(2), $T(x)$ and $T_1(L(x))$ are
polynomial, and by sides~(5), $l_\infty(p(f'))=\pi_1(L(T(x)))$ and
$p_1(f'_1)=\pi_1(T_1(y))$. By ($*$) with $u=1$ and $m=k$,
$l_{h'}(L(T(x)))=\chi(t)|t|^r\,l_{h'}(T_1(L(x)))$ for every $h'$, so
$\pi_1(L(T(x)))=\chi(t)|t|^r\,\pi_1(T_1(L(x)))$ by (T$_\infty$), and
$\pi_1(T_1(L(x)))=\pi_1(T_1(y))$ by sides~(6), as $\pi_1(L(x))=\pi_1(y)$. So
$l_\infty(p(f'))=\chi(t)|t|^r\,p_1(f'_1)$.
\end{proof}
% Lean: (1) is bump_pt with nz_pure (MemPolData.lean), with PolDense, ArchOp and Hyp.unim
% (I_left). (2) is lams (Cross.lean): Step 1 is exists_test_vectors_at_place and
% exists_test_vectors, Step 2 kerF and kerI, Step 3 exists_section, Step 4 eqF with
% levi_at, Step 5 (a), (b), (c) eqN, eqL (with levi_t_one) and eqK (ISmulEqZero.lean),
% each combined with the Cross maps c and c1 (sides (5)) as in lams.

\begin{lemma}\label{asm-theta-agree-local-alpha}
There are signs $e_1,e_2,s\in\{\pm1\}$ such that the following holds. Let $B$ be
a theta base of a CM field $E$, $(G,G_1,H,H')$ an $\alpha$-configuration
(Definition~\ref{asm-theta-agree-planes}(4)) with $G$ of ranks $(b+2,a')$, $R$
and $R_1$ real maps for $G$ and $G_1$, $e$ an evaluation map for
$(M,M_1,\iota_H,\iota_{H'},\mathcal J_\alpha(H,H'))$, $e'$ one for
$(M',M'_1,1,\iota_{H'},\mathcal J_0(H'))$, $K\subset U_V$ the lattice stabilizer
of $V$, and $\mu_K$ and $\mu$ Haar measures on $K$ and $\mathbf A_E^\times$. Then
there is $C\neq0$ such that for all matched quadruples
$(B,G,\mathcal X,\omega,\omega',R,R_f)$ and
$(B,G_1,\mathcal X,\omega_1,\omega'_1,R_1,R_{f,1})$ such that nondegeneracy holds
below $p+q$ (Definition~\ref{asm-theta-agree-matched}(6)), and all
corresponding $\Phi$ and $\Phi'$, the functions $\operatorname{cone}_\alpha(\Phi)$
and $\lambda\,e'(\Phi')$ correspond for the second quadruple. Here
$\operatorname{cone}_\alpha$ is the cone for
$(M,\omega,e,K,\mu_K,\mu,u_\alpha,\nu)$ and $z=k_\alpha$, and
$\lambda=C\sigma_{e_1,e_2,s}(\mathfrak b_H,\mathfrak b_{H'};1-k_\alpha)L^{S_0}(\nu,k_\alpha)$.
\end{lemma}
\leanref{Development/WeilBound/ThetaLifts/ConeHypotheses/PolDense.lean}{PositivityTheory.ThKG.LocA'}
\begin{proof}
Let $e_1,e_2$ be the signs of
Lemma~\ref{asm-theta-agree-local-factor}(1) and $s$ the sign of
Lemma~\ref{asm-theta-agree-local-ramified-values}(1). Let $B$,
$(G,G_1,H,H')$, $R=(\Pi,p,p')$, $R_1=(\Pi_1,p_1,p'_1)$, $e$, $e'$, $K$,
$\mu_K$ and $\mu$ be given, let $e_\infty,e_v$ and $e'_\infty,e'_v$ be the
components of $e$ and $e'$, and let $\mathrm{vac}_{1,v}$ and
$\mathrm{vac}'_{1,v}$ be the vacua of $M_1$ and $M'_1$. If there is no triple
$(B,G,\mathcal X,\omega,\omega')$, there is no matched quadruple for $G$, and
$C=1$ will do. Otherwise $V$ is Hermitian, and
Lemma~\ref{asm-theta-agree-local-measures} gives a normalized local measure
system $m$ for $K$. Lemma~\ref{asm-theta-agree-local-factor}(1) gives
$\kappa\neq0$, Lemma~\ref{asm-theta-agree-ct-alpha-zero} gives $C_1\neq0$ for
$e'$, and Lemma~\ref{asm-theta-agree-ct-alpha-cone} gives $C_2\neq0$ for $e$,
$K$, $\mu_K$ and $\mu$.

Call a pair of matched quadruples $Q=(B,G,\mathcal X,\omega,\omega',R,R_f)$
and $Q_1=(B,G_1,\mathcal X,\omega_1,\omega'_1,R_1,R_{f,1})$ \emph{eligible} if
nondegeneracy holds below $p+q$, and for it put
$\sigma=\sigma_{e_1,e_2,s}(\mathfrak b_H,\mathfrak b_{H'};1-k_\alpha)$ and
$L_0=L^{S_0}(\nu,k_\alpha)$. Steps 2 to 8 show that for every eligible pair
there are an archimedean cone $L$ for $(M,\omega,m,e_\infty,u_\alpha,\nu,k_\alpha)$
(Definition~\ref{asm-theta-agree-local-cones}) and $c\neq0$ such that:
\begin{enumerate}
\item[(i)] $p_1(L(f))=c\,p'_1(e'_\infty(f'))$ whenever $f\in\operatorname{Pol}(M)$,
  $f'\in\operatorname{Pol}(M')$ and $p(f)=p'(f')$;
\item[(ii)] for all corresponding $\Phi$ and $\Phi'$ for $Q$, the functions
  $\operatorname{cone}_\alpha(\Phi)$ and $\kappa c\sigma L_0\,e'(\Phi')$
  correspond for $Q_1$.
\end{enumerate}

\emph{Step 1: one constant.} Let $(L,c)$ and $(\tilde L,\tilde c)$ be as in
(i) and (ii) for eligible pairs $(Q,Q_1)$ and $(\tilde Q,\tilde Q_1)$. The Weil
pairs $(V,W,M,\beta,\beta/\alpha,\omega)$ and
$(V,W,M,\tilde\beta,\tilde\beta/\tilde\alpha,\tilde\omega)$ of $Q$ and $\tilde Q$
have the same signature. As both character data are proper, $\beta/\alpha$ and
$\tilde\beta/\tilde\alpha$ have the same infinity type $m_\beta-m_\alpha$, and
$\nu(l)=\tilde\nu(l)$ whenever $l-1$ is at $\infty$. So
Lemma~\ref{asm-theta-agree-local-real-unique} gives
$\omega(g,1)\Phi=\tilde\omega(g,1)\Phi$ for $g\in U_V^{(\infty)}$ and
Schwartz-Bruhat $\Phi$, and Lemma~\ref{asm-theta-agree-local-cone-unique}
gives $L=\tilde L$. Now (i) holds for $L$ with $c$ and with $\tilde c$. As
$R_1$ is matched to the triple of $Q_1$, $p'_1$ maps $\operatorname{Pol}(M'_1)$
onto $\Pi_1\neq0$; choose $f_1\in\operatorname{Pol}(M'_1)$ with
$p'_1(f_1)\neq0$. By Lemma~\ref{asm-theta-agree-evaluation}, $e'$ carries
$\operatorname{PD}(M')$ onto $\operatorname{PD}(M'_1)$, so
Lemma~\ref{asm-theta-agree-local-evaluation-pol} gives
$f'\in\operatorname{Pol}(M')$ with $e'_\infty(f')=f_1$; and as $p$ maps
$\operatorname{Pol}(M)$ onto $\Pi$, there is $f\in\operatorname{Pol}(M)$ with
$p(f)=p'(f')$. Then $c\,p'_1(f_1)=p_1(L(f))=\tilde c\,p'_1(f_1)$, so
$c=\tilde c$. Let $c$ be this common value, or $c=1$ if there is no eligible
pair, and put $C=\kappa c\neq0$. By (ii), the lemma holds with $C$.

\emph{Step 2: the smaller quadruple.} Fix an eligible pair $(Q,Q_1)$. As
$p_1+q_1=p+q-1$, $Q_1$ is nondegenerate: $I_1(\Phi_1)\neq0$ for some
$\Phi_1\in\operatorname{PD}(M_1)$, and $I'_1$ does not vanish on
$\operatorname{PD}(M'_1)\cap\mathcal S^\circ(M'_1)$. By
Lemma~\ref{asm-theta-agree-local-injective}, as $q_1=0$, the first side and the
second side of $Q_1$ are injective
(Definition~\ref{asm-theta-agree-local-transfer}). As in Step 1,
$e'_\infty(\operatorname{Pol}(M'))=\operatorname{Pol}(M'_1)$.

\emph{Step 3: the transfer by $e'$.} We apply
Lemma~\ref{asm-theta-agree-local-transfer-criteria}(1) and (2), with $C_1$, to
$(e',e'_\infty,(e'_v)_v)$ and the second sides of $Q$ and $Q_1$, which exist
as $q=q_1=0$. Levi elements for $(1,h_1)$ exist by
Lemma~\ref{asm-theta-agree-levi}; $e'_\infty$ maps $\operatorname{Pol}(M')$
into $\operatorname{Pol}(M'_1)$, and $e'_v$ maps Schwartz-Bruhat functions to
Schwartz-Bruhat functions. Let $f\in\operatorname{Pol}(M')$ and $\phi$ be a
family in $S^{\mathrm{SB}}$. Then $f\otimes\phi\in\operatorname{PD}(M')$, and
Lemma~\ref{asm-theta-agree-ct-alpha-zero} gives $f'_\alpha$ such that
$(f'_\alpha,0)$ splits $I'(f\otimes\phi)$ along $H'$ and
$f'_\alpha(mh)=C_1I'_1(e'(\omega'(1,h)(f\otimes\phi)))(h_1)$ for every Levi
element $m$ for $(1,h_1)$. So $f'_\alpha$ has the law $(\alpha,a'/2)$, and by
Lemma~\ref{asm-theta-agree-local-split-invariant}, as $q=0$, it is invariant
under left translation by $N_{H'}$ and vanishes if $I'(f\otimes\phi)=0$. With
$T=f'_\alpha$, the lemma shows that $(e',e'_\infty,(e'_v)_v)$ is a transfer
along $(H',\alpha,a'/2)$ from the second side of $Q$ to that of $Q_1$.

\emph{Step 4: the transfer by the cone.} Let $L$, $\Lambda_v$ and $u_v$ be
given by Lemma~\ref{asm-theta-agree-local-factor}(1) for the triples of $Q$
and $Q_1$, and put
$\lambda_0=\kappa\,\nu(\mathfrak b_H)^{e_1}\nu(\mathfrak b_{H'})^{e_2}L_0$.
Then $\lambda_0\neq0$: the values of $\nu$ on ideals are nonzero, and
$L_0\neq0$ by Lemma~\ref{asm-theta-agree-factors}, since $k_\alpha\geq2$ as
$G$ is proper. Parts (a), (b), (c) and (e) of
Lemma~\ref{asm-theta-agree-local-factor}(1) give the hypotheses of
Lemma~\ref{asm-theta-agree-local-transfer-criteria}(3) for $T_0=T$,
$k=\lambda_0$ and $S=S_1=S^{\mathrm{SB}}$, and those of (1) other than the last
one. For that, let $f\in\operatorname{Pol}(M)$ and $\phi$ be a family in
$S^{\mathrm{SB}}$. Lemma~\ref{asm-theta-agree-split} gives a pair
$(f_\alpha,f_\beta)$ that splits $I(f\otimes\phi)$ along $H'$; $f_\alpha$ has
the law $(\alpha,a'/2)$, it is invariant under $N_{H'}$ and vanishes if
$I(f\otimes\phi)=0$ by Lemma~\ref{asm-theta-agree-local-split-invariant}, and
$f_\alpha(mh)=C_2I_1(\operatorname{cone}_\alpha(\omega(1,h)(f\otimes\phi)))(h_1)$
by Lemma~\ref{asm-theta-agree-ct-alpha-cone}. With $T=f_\alpha$ and $C_2$,
$(\operatorname{cone}_\alpha,L,\Lambda)$ is a transfer along $(H',\alpha,a'/2)$
from the first side of $Q$ to that of $Q_1$.

\emph{Step 5: test vectors.} For $v\in T$ let
$w_v=\nu_v(s)\,\epsilon(\nu_v,1-k_\alpha)$, and let $\phi_v$ and $\phi'_v$ be
the functions of Lemma~\ref{asm-theta-agree-local-ramified-values}(1); by
Lemma~\ref{asm-theta-agree-local-factor}(1)(b),
$p_{1,v}(\Lambda_v\phi_v)=w_v\,p'_{1,v}(e'_v\phi'_v)$. As $T\subset P$,
$\sigma=\nu(\mathfrak b_H)^{e_1}\nu(\mathfrak b_{H'})^{e_2}\prod_{v\in T}w_v$,
and $\sigma\neq0$ by Lemma~\ref{asm-theta-agree-factors}; so every
$w_v\neq0$. For $v\notin T$, as $R_{f,1}$ is matched,
$p'_{1,v}(e'_v\mathrm{vac}'_v)=p'_{1,v}(\mathrm{vac}'_{1,v})=p_{1,v}(\mathrm{vac}_{1,v})\neq0$,
and $p_{1,v}(\Lambda_v\mathrm{vac}_v)=u_v\,p_{1,v}(\mathrm{vac}_{1,v})\neq0$.
Choose $f_1\in\operatorname{Pol}(M'_1)$ with $p'_1(f_1)\neq0$ and
$f\in\operatorname{Pol}(M')$ with $e'_\infty(f)=f_1$ (Step 2). Then $f$ and the
family equal to $\phi'_v$ for $v\in T$ and to $\mathrm{vac}'_v$ for $v\notin T$
form a test vector for $(e'_\infty,(e'_v)_v)$ relative to
$S^{\mathrm{SB}}$, $p'_1$ and $(p'_{1,v})_v$. Next,
Lemma~\ref{asm-theta-agree-local-maps}(1), for the first sides of $Q$ and
$Q_1$, applies to $L$: $L(\operatorname{Pol}(M))\subset\operatorname{Pol}(M_1)$,
the condition on linear forms is Lemma~\ref{asm-theta-agree-local-bump}(1),
and $I_1(\Phi_1)\neq0$, where $\operatorname{PD}(M_1)$ is spanned by the
$f\otimes\phi$ with $\phi$ in $S^{\mathrm{SB}}$. So $p_1(L(f))\neq0$ for some
$f\in\operatorname{Pol}(M)$, and $f$ with the family equal to $\phi_v$ for
$v\in T$ and to $\mathrm{vac}_v$ for $v\notin T$ is a test vector for
$(L,\Lambda)$ relative to $S^{\mathrm{SB}}$, $p_1$ and $(p_{1,v})_v$.

\emph{Step 6: local maps.} Lemma~\ref{asm-theta-agree-local-maps}(2), for the
second sides and Steps 2, 3 and 5, gives linear maps
$l_\infty\colon\Pi\to\Pi_1$ and $l_v\colon\Pi_v\to\Pi_{1,v}$ with
$l_\infty(p'(f))=p'_1(e'_\infty(f))$ for $f\in\operatorname{Pol}(M')$ and
$l_v(p'_v(x))=p'_{1,v}(e'_vx)$ for Schwartz-Bruhat $x$. For the first sides and
Steps 2, 4 and 5 it gives $l'_\infty$ and $l'_v$ with
$l'_\infty(p(f))=p_1(L(f))$ for $f\in\operatorname{Pol}(M)$ and
$l'_v(p_v(x))=p_{1,v}(\Lambda_vx)$ for Schwartz-Bruhat $x$. All these maps are
equivariant for $(H',\alpha,a'/2)$, and $(Q,Q_1)$ has step type
$(\alpha,a'/2)$ (Definition~\ref{asm-theta-agree-local-equivariant}), as
$(G,G_1,H,H')$ is an $\alpha$-configuration.

\emph{Step 7: comparison of the local maps.} The test vector of Step 5 for
$e'$ gives $l_\infty(p'(f))\neq0$, so $l_\infty\neq0$, and
Lemma~\ref{asm-theta-agree-local-mult-archimedean} gives $c$ with
$l'_\infty=c\,l_\infty$. Then $c\neq0$, as $l'_\infty(p(f))=p_1(L(f))\neq0$ for
the test vector of $L$. Let $v\in T$. Then $l_v(p'_v\phi'_v)\neq0$, and
Lemma~\ref{asm-theta-agree-local-mult-ramified} gives $k$ with $l'_v=k\,l_v$.
As $p_v(\phi_v)=p'_v(\phi'_v)$,
\[
k\,p'_{1,v}(e'_v\phi'_v)=k\,l_v(p_v\phi_v)=l'_v(p_v\phi_v)=p_{1,v}(\Lambda_v\phi_v)=w_v\,p'_{1,v}(e'_v\phi'_v),
\]
so $l'_v=w_vl_v$. Let $v\notin T$. As $p_v(\mathrm{vac}_v)=p'_v(\mathrm{vac}'_v)$,
\[
l'_v(p_v\mathrm{vac}_v)=p_{1,v}(\Lambda_v\mathrm{vac}_v)=u_v\,p_{1,v}(\mathrm{vac}_{1,v})=u_v\,p'_{1,v}(e'_v\mathrm{vac}'_v)=u_v\,l_v(p'_v\mathrm{vac}'_v),
\]
and Lemma~\ref{asm-theta-agree-local-mult-unramified} gives $l'_v=u_vl_v$; put
$w_v=u_v$ for $v\notin T$. Now (i) holds: if $p(f)=p'(f')$, then
$p_1(L(f))=l'_\infty(p(f))=c\,l_\infty(p'(f'))=c\,p'_1(e'_\infty(f'))$.

\emph{Step 8: correspondence.} Let $\Phi$ and $\Phi'$ correspond for $Q$,
through a finite set $T'\supset T$ and sums $\Phi=\sum_jf_j\otimes\phi_j$ and
$\Phi'=\sum_if'_i\otimes\phi'_i$ as in
Definition~\ref{asm-theta-agree-matched}(4), and put
$u=\prod_{v\in T'\setminus T}u_v\neq0$. By (d) and (e) of
Lemma~\ref{asm-theta-agree-local-factor}(1),
$\operatorname{cone}_\alpha(\Phi)=\sum_ju^{-1}\lambda_0\,L(f_j)\otimes\psi_j$,
with $\psi_{j,v}=\Lambda_v\phi_{j,v}$ for $v\in T'$ and
$\psi_{j,v}=\mathrm{vac}_{1,v}$ for $v\notin T'$; and
$\kappa c\sigma L_0\,e'(\Phi')=\sum_i\kappa c\sigma L_0\,e'_\infty(f'_i)\otimes\bigotimes_ve'_v\phi'_{i,v}$,
with $e'_v\phi'_{i,v}=\mathrm{vac}'_{1,v}$ for $v\notin T'$. These sums have
the form required in the definition of correspondence for $Q_1$, with the
same set $T'$. Apply the linear map
$l'_\infty\otimes\bigotimes_{v\in T'}l'_v=c\prod_{v\in T'}w_v\,(l_\infty\otimes\bigotimes_{v\in T'}l_v)$
to the identity in $\Pi\otimes\bigotimes_{v\in T'}\Pi_v$ that expresses the
correspondence of $\Phi$ and $\Phi'$. By Step 6 this gives
\[
\sum_jp_1(L(f_j))\otimes\bigotimes_{v\in T'}p_{1,v}(\Lambda_v\phi_{j,v})
=c\prod_{v\in T'}w_v\sum_ip'_1(e'_\infty(f'_i))\otimes\bigotimes_{v\in T'}p'_{1,v}(e'_v\phi'_{i,v}).
\]
As
$u^{-1}\lambda_0c\prod_{v\in T'}w_v=\kappa c\,\nu(\mathfrak b_H)^{e_1}\nu(\mathfrak b_{H'})^{e_2}\prod_{v\in T}w_v\,L_0=\kappa c\sigma L_0$,
multiplying by $u^{-1}\lambda_0$ gives the identity that makes
$\operatorname{cone}_\alpha(\Phi)$ and $\kappa c\sigma L_0\,e'(\Phi')$
correspond for $Q_1$. This is (ii).
\end{proof}
% Lean: PositivityTheory.ThKG.Loc.locA_of Hyp Hyp2 ConeFacA StdA (LocalTheory/LocAOf.lean).
% Step 1 is one_const (StepFacEq.lean) with coneI_eq (pres from Hyp2.runiq, nu_eq,
% hasExp_div) and the uniqueness argument through ex_p_ne, ev_pol and sideV.ontoI; the
% first paragraph uses Hyp2.meas, ConeFacA, Hyp2.ctSZ and Hyp2.ctBC. Step 2 is nz (from Below), Hyp.inj and
% ep; Step 3 is b1 = br_of lev ep.1 evS (fac_ev ...) hC1 with ct1 and Hyp.sinv; Step 4 is
% b2 = br_of ... (fac_cone ...) hC2 with Hyp.split, Hyp.sinv, ct2 and Hyp.fac; Step 5 is
% std, wne (stepFac_eq), pt1 and pt2 (pt_of, bump_pt with Hyp2.bumpA); Step 6 is the two
% calls of lams with side_condition_vanishing and sideV; Step 7 is Hyp2.inf, Hyp2.ram and
% Hyp2.unr (hw); Step 8 is cone_sum, cfac, tens_step and prod_w. Size of locA_of and the
% generic lemmas it calls (StepFacEq.lean, MemPolData.lean): about 940 lines of closure,
% most of them shared with Lemma asm-theta-agree-local-beta and the sublemmas above.

\begin{lemma}\label{asm-theta-agree-local-beta}
There are signs $e_1,e_2,s\in\{\pm1\}$ such that the following holds. Let $B$ be
a theta base of a CM field $E$, $(G,G_1,H,H')$ a $\beta$-configuration
(Definition~\ref{asm-theta-agree-planes}(4)) with $G$ of ranks $(a,b+2)$, $R$
and $R_1$ real maps for $G$ and $G_1$, $e$ an evaluation map for
$(M,M_1,1,\iota_{H'},\mathcal J_0(H'))$, $e'$ one for
$(M',M'_1,\iota_H,\iota_{H'},\mathcal J_\beta(H,H'))$, $K\subset U_{V'}$ the
lattice stabilizer of $V'$, and $\mu_K$ and $\mu$ Haar measures on $K$ and
$\mathbf A_E^\times$. Then there is $C\neq0$ such that for all matched
quadruples $(B,G,\mathcal X,\omega,\omega',R,R_f)$ and
$(B,G_1,\mathcal X,\omega_1,\omega'_1,R_1,R_{f,1})$ such that nondegeneracy holds
below $p+q$ (Definition~\ref{asm-theta-agree-matched}(6)), and all corresponding
$\Phi$ and $\Phi'$ with $\Phi'\in\mathcal S^\circ(M')$: the function
$\operatorname{cone}_\beta(\Phi')$ lies in $\mathcal S^\circ(M'_1)$, and
$\lambda\,e(\Phi)$ and $\operatorname{cone}_\beta(\Phi')$ correspond for the
second quadruple. Here $\operatorname{cone}_\beta$ is the cone for
$(M',\omega',e',K,\mu_K,\mu,u_\beta,\nu)$ and $z=k_\beta+1$, and
$\lambda=C\sigma_{e_1,e_2,s}(\mathfrak b_H,\mathfrak b_{H'};-k_\beta)L^{S_0}(\nu,k_\beta+1)$.
\end{lemma}
\leanref{Development/WeilBound/ThetaLifts/ConeHypotheses/PolDense.lean}{PositivityTheory.ThKG.LocB'}
\begin{proof}
The proof follows that of Lemma~\ref{asm-theta-agree-local-alpha}, with the
roles of the two spaces exchanged and with full-rank data on the second
space. Let $e_1,e_2$ and $s$ be the signs of
Lemmas~\ref{asm-theta-agree-local-factor}(2)
and~\ref{asm-theta-agree-local-ramified-values}(2). Let the data be given,
with $R=(\Pi,p,p')$, $R_1=(\Pi_1,p_1,p'_1)$, components $e_\infty,e_v$ and
$e'_\infty,e'_v$ of $e$ and $e'$, and vacua $\mathrm{vac}_{1,v}$ and
$\mathrm{vac}'_{1,v}$ of $M_1$ and $M'_1$. If there is no triple for $G$, take
$C=1$. Otherwise $V'$ is Hermitian, Lemma~\ref{asm-theta-agree-local-measures}
gives a normalized local measure system $m$ for $K$,
Lemma~\ref{asm-theta-agree-local-factor}(2) gives $\kappa\neq0$, and
Lemmas~\ref{asm-theta-agree-ct-beta-zero} and~\ref{asm-theta-agree-ct-beta-cone}
give $C_1\neq0$ for $e$ and $C_2\neq0$ for $e'$, $K$, $\mu_K$ and $\mu$. For an
eligible pair $(Q,Q_1)$, defined as before, put
$\sigma=\sigma_{e_1,e_2,s}(\mathfrak b_H,\mathfrak b_{H'};-k_\beta)$ and
$L_0=L^{S_0}(\nu,k_\beta+1)$. Steps 2 to 9 show that for every eligible pair
there are an archimedean cone $L$ for
$(M',\omega',m,e'_\infty,u_\beta,\nu,k_\beta+1)$ and $c\neq0$ such that:
\begin{enumerate}
\item[(i)] $p'_1(L(f'))=c\,p_1(e_\infty(f))$ whenever $f\in\operatorname{Pol}(M)$,
  $f'\in\operatorname{Pol}(M')$ and $p(f)=p'(f')$;
\item[(ii)] $\operatorname{cone}_\beta$ maps $\mathcal S^\circ(M')$ into
  $\mathcal S^\circ(M'_1)$, and for corresponding $\Phi$ and $\Phi'$ for $Q$
  the functions $\kappa c\sigma L_0\,e(\Phi)$ and
  $\operatorname{cone}_\beta(\Phi')$ correspond for $Q_1$.
\end{enumerate}

\emph{Step 1: one constant.} As in Step 1 of the proof of
Lemma~\ref{asm-theta-agree-local-alpha}, the archimedean cones of two eligible
pairs are equal, by Lemmas~\ref{asm-theta-agree-local-real-unique}
and~\ref{asm-theta-agree-local-cone-unique} applied to the Weil pairs
$(V',W,M',\alpha,\alpha/\beta,\omega')$, whose characters $\alpha/\beta$ have
the infinity type $m_\alpha-m_\beta$. For the uniqueness of $c$, choose
$f_1\in\operatorname{Pol}(M_1)$ with $p_1(f_1)\neq0$, $f\in\operatorname{Pol}(M)$
with $e_\infty(f)=f_1$ (Lemmas~\ref{asm-theta-agree-evaluation}
and~\ref{asm-theta-agree-local-evaluation-pol}), and
$f'\in\operatorname{Pol}(M')$ with $p'(f')=p(f)$; then (i) determines $c$.
With $c$ the common value, or $c=1$ if there is no eligible pair, $C=\kappa c$
works by (ii).

\emph{Step 2: the smaller quadruple.} Fix an eligible pair. As
$p_1+q_1=p+q-1$, $Q_1$ is nondegenerate. By
Lemma~\ref{asm-theta-agree-local-injective}, the first side and the full-rank
side of $Q_1$ are injective, and as in Step 1,
$e_\infty(\operatorname{Pol}(M))=\operatorname{Pol}(M_1)$.

\emph{Step 3: the transfer by $e$.} As in Step 3 of the proof of
Lemma~\ref{asm-theta-agree-local-alpha}, with
Lemma~\ref{asm-theta-agree-split}, the part $f_\beta$ of the pair that splits
$I(f\otimes\phi)$, Lemma~\ref{asm-theta-agree-local-split-invariant} and
Lemma~\ref{asm-theta-agree-ct-beta-zero} with $C_1$,
Lemma~\ref{asm-theta-agree-local-transfer-criteria}(1) and (2) show that
$(e,e_\infty,(e_v)_v)$ is a transfer along $(H',\beta,a/2)$ from the first side
of $Q$ to that of $Q_1$.

\emph{Step 4: the transfer by the cone.} Let $L$, $\Lambda_v$ and $u_v$ be
given by Lemma~\ref{asm-theta-agree-local-factor}(2), and put
$\lambda_0=\kappa\,\nu(\mathfrak b_H)^{e_1}\nu(\mathfrak b_{H'})^{e_2}L_0\neq0$,
using Lemma~\ref{asm-theta-agree-factors} and $k_\beta+1\geq2$. The map
$\Lambda_v$ sends $S^\circ_v(M')$ into $S^\circ_v(M'_1)$: for $v\neq v_1$ both
are the Schwartz-Bruhat functions, and for $v=v_1$ this follows from
Lemma~\ref{asm-theta-agree-local-augmentation}(2) and part (b) of
Lemma~\ref{asm-theta-agree-local-factor}, by linearity. The place $v_1$ is
clean, so $v_1\notin T$; hence if $\mathrm{vac}'_{v_1}\in S^\circ_{v_1}(M')$,
then $u_{v_1}\mathrm{vac}'_{1,v_1}=\Lambda_{v_1}\mathrm{vac}'_{v_1}$ lies in
$S^\circ_{v_1}(M'_1)$ with $u_{v_1}\neq0$, and so does $\mathrm{vac}'_{1,v_1}$.
So Lemma~\ref{asm-theta-agree-local-transfer-criteria}(3) applies with
$T_0=T$, $k=\lambda_0$, $S=S^\circ(M')$ and $S_1=S^\circ(M'_1)$. For (1), let
$f\in\operatorname{Pol}(M')$ and $\phi$ be a family in $S^\circ(M')$. Then
$f\otimes\phi\in\operatorname{PD}(M')\cap\mathcal S^\circ(M')$, as it is
linear in $\phi_{v_1}$. Lemma~\ref{asm-theta-agree-split} gives a pair
$(f'_\alpha,f'_\beta)$ that splits $I'(f\otimes\phi)$ along $H'$, and $T=f'_\beta$
has the law $(\beta,a/2)$, is invariant under $N_{H'}$ and vanishes with
$I'(f\otimes\phi)$ by Lemma~\ref{asm-theta-agree-local-split-invariant}, and
satisfies
$T(mh)=C_2I'_1(\operatorname{cone}_\beta(\omega'(1,h)(f\otimes\phi)))(h_1)$ by
Lemma~\ref{asm-theta-agree-ct-beta-cone}. So
$(\operatorname{cone}_\beta,L,\Lambda)$ is a transfer along $(H',\beta,a/2)$
from the full-rank side of $Q$ to that of $Q_1$.

\emph{Step 5: augmentation at $v_1$.} We show: if $x$ and $y$ are
Schwartz-Bruhat functions on $F_{v_1}^{I_{M'}}$ with
$p'_{v_1}(x)=p'_{v_1}(y)$, then
$p'_{1,v_1}(\Lambda_{v_1}x)=p'_{1,v_1}(\Lambda_{v_1}y)$. Let $\ell$ be a linear
form on $\Pi_{1,v_1}$. As $R_{f,1}$ is matched, $\ell\circ p'_{1,v_1}$
vanishes on $\mathfrak a_{v_1}(\omega'_1)$, so by
Lemma~\ref{asm-theta-agree-local-augmentation}(1) it takes the same value on
$\Lambda_{v_1}\phi_1$ and $\Lambda_{v_1}\phi$ whenever $\omega'(g,1)$ acts at
$v_1$ by $\phi\mapsto\phi_1$ with $g\in U_{V'}^{(v_1)}$. As $\ell$ is
arbitrary, $p'_{1,v_1}\circ\Lambda_{v_1}$ vanishes on
$\mathfrak a_{v_1}(\omega')$. By Lemma~\ref{asm-theta-agree-coinvariants}, $Q$
describes the coinvariants, and $v_1$ is clean, so $p'_{v_1}(x-y)=0$ gives
$x-y\in\mathfrak a_{v_1}(\omega')$, and the claim follows.

\emph{Step 6: test vectors.} For $v\in T$ let
$w_v=\nu_v(s)\,\epsilon(\nu_v,-k_\beta)$, and let $\phi_v$ and $\phi'_v$ be the
functions of Lemma~\ref{asm-theta-agree-local-ramified-values}(2), so that
$p'_{1,v}(\Lambda_v\phi'_v)=w_v\,p_{1,v}(e_v\phi_v)$; as before, $w_v\neq0$.
Choose $f\in\operatorname{Pol}(M)$ with $p_1(e_\infty(f))\neq0$ (Step 2). Then
$f$ and the family equal to $\phi_v$ for $v\in T$ and to $\mathrm{vac}_v$ for
$v\notin T$ form a test vector for $(e_\infty,(e_v)_v)$ relative to
$S^{\mathrm{SB}}$, $p_1$ and $(p_{1,v})_v$, as
$p_{1,v}(e_v\mathrm{vac}_v)=p_{1,v}(\mathrm{vac}_{1,v})\neq0$ for $v\notin T$.
Next, Lemma~\ref{asm-theta-agree-local-maps}(1), for the full-rank sides,
applies to $L$, by part (a) of Lemma~\ref{asm-theta-agree-local-factor}(2),
Lemma~\ref{asm-theta-agree-local-bump}(2) and Step 2, as
$\mathcal S^\circ(M'_1)$ is spanned by the $f\otimes\phi$ with $\phi$ in
$S^\circ(M'_1)$; it gives $f'\in\operatorname{Pol}(M')$ with
$p'_1(L(f'))\neq0$. As $Q$ describes the coinvariants, there is a function $x$
of full rank at $\tau_1$ with $p'_{v_1}(x)=p'_{v_1}(\mathrm{vac}'_{v_1})$, and
by Step 5,
$p'_{1,v_1}(\Lambda_{v_1}x)=p'_{1,v_1}(\Lambda_{v_1}\mathrm{vac}'_{v_1})=u_{v_1}p'_{1,v_1}(\mathrm{vac}'_{1,v_1})\neq0$.
So $f'$ and the family equal to $x$ at $v_1$, to $\phi'_v$ for $v\in T$ and to
$\mathrm{vac}'_v$ elsewhere form a test vector for $(L,\Lambda)$ relative to
$S^\circ(M')$, $p'_1$ and $(p'_{1,v})_v$; at $v\notin T\cup\{v_1\}$,
$p'_{1,v}(\Lambda_v\mathrm{vac}'_v)=u_v\,p'_{1,v}(\mathrm{vac}'_{1,v})\neq0$.

\emph{Step 7: local maps.} Lemma~\ref{asm-theta-agree-local-maps}(2), for the
first sides and Steps 2, 3 and 6, gives $l_\infty$ and $l_v$ with
$l_\infty(p(f))=p_1(e_\infty(f))$ for $f\in\operatorname{Pol}(M)$ and
$l_v(p_v(x))=p_{1,v}(e_vx)$ for Schwartz-Bruhat $x$. For the full-rank sides
and Steps 2, 4 and 6 it gives $l'_\infty$ and $l'_v$ with
$l'_\infty(p'(f'))=p'_1(L(f'))$ for $f'\in\operatorname{Pol}(M')$ and
$l'_v(p'_v(x))=p'_{1,v}(\Lambda_vx)$ for $x\in S^\circ_v(M')$. The last
identity holds for every Schwartz-Bruhat $x$: for $v\neq v_1$ this is
immediate, and for $v=v_1$, as $Q$ describes the coinvariants, there is
$y\in S^\circ_{v_1}(M')$ with $p'_{v_1}(y)=p'_{v_1}(x)$, and Step 5 gives
$l'_{v_1}(p'_{v_1}(x))=p'_{1,v_1}(\Lambda_{v_1}y)=p'_{1,v_1}(\Lambda_{v_1}x)$.
All these maps are equivariant for $(H',\beta,a/2)$, and $(Q,Q_1)$ has step
type $(\beta,a/2)$, as $(G,G_1,H,H')$ is a $\beta$-configuration.

\emph{Step 8: comparison of the local maps.} As in Step 7 of the proof of
Lemma~\ref{asm-theta-agree-local-alpha}, the test vectors give $l_\infty\neq0$
and $l'_\infty=c\,l_\infty$ with $c\neq0$, by
Lemma~\ref{asm-theta-agree-local-mult-archimedean}. For $v\in T$,
Lemma~\ref{asm-theta-agree-local-mult-ramified} gives $l'_v=k\,l_v$, and
\[
k\,p_{1,v}(e_v\phi_v)=k\,l_v(p'_v\phi'_v)=l'_v(p'_v\phi'_v)=p'_{1,v}(\Lambda_v\phi'_v)=w_v\,p_{1,v}(e_v\phi_v),
\]
so $l'_v=w_vl_v$. For $v\notin T$,
$l'_v(p_v\mathrm{vac}_v)=l'_v(p'_v\mathrm{vac}'_v)=p'_{1,v}(\Lambda_v\mathrm{vac}'_v)=u_v\,p'_{1,v}(\mathrm{vac}'_{1,v})=u_v\,p_{1,v}(e_v\mathrm{vac}_v)=u_v\,l_v(p_v\mathrm{vac}_v)$,
so $l'_v=u_vl_v$ by Lemma~\ref{asm-theta-agree-local-mult-unramified}; put
$w_v=u_v$ for $v\notin T$. Now (i) holds: if $p(f)=p'(f')$, then
$p'_1(L(f'))=l'_\infty(p'(f'))=c\,l_\infty(p(f))=c\,p_1(e_\infty(f))$.

\emph{Step 9: full rank and correspondence.} The space $\mathcal S^\circ(M')$
is spanned by the $f\otimes\phi$ with $f\in\mathcal S(F_\infty^{I_{M'}})$ and
$\phi$ in $S^\circ(M')$. For such $f\otimes\phi$, condition (b) of the
transfer of Step 4, with $v_0=v_1$ and $x=\phi_{v_1}$, writes
$\operatorname{cone}_\beta(f\otimes\phi)$ as a multiple of
$L(f)\otimes\psi[v_1\mapsto\Lambda_{v_1}\phi_{v_1}]$ with $\psi$ in
$S^\circ(M'_1)$; this family is in $S^\circ(M'_1)$, as
$\Lambda_{v_1}\phi_{v_1}\in S^\circ_{v_1}(M'_1)$, so the function lies in
$\mathcal S^\circ(M'_1)$, being linear in its component at $v_1$. By part (d)
of Lemma~\ref{asm-theta-agree-local-factor}(2), $\operatorname{cone}_\beta$ is
linear on Schwartz-Bruhat functions, so it maps $\mathcal S^\circ(M')$ into
$\mathcal S^\circ(M'_1)$. Now let $\Phi=\sum_if_i\otimes\phi_i$ and
$\Phi'=\sum_jf'_j\otimes\phi'_j$ correspond for $Q$ through $T'\supset T$, and
put $u=\prod_{v\in T'\setminus T}u_v$. Then
$\kappa c\sigma L_0\,e(\Phi)=\sum_i\kappa c\sigma L_0\,e_\infty(f_i)\otimes\bigotimes_ve_v\phi_{i,v}$
and
$\operatorname{cone}_\beta(\Phi')=\sum_ju^{-1}\lambda_0\,L(f'_j)\otimes\psi'_j$,
with $\psi'_{j,v}=\Lambda_v\phi'_{j,v}$ for $v\in T'$ and
$\psi'_{j,v}=\mathrm{vac}'_{1,v}$ otherwise. Applying
$l'_\infty\otimes\bigotimes_{v\in T'}l'_v=c\prod_{v\in T'}w_v\,(l_\infty\otimes\bigotimes_{v\in T'}l_v)$
to the identity expressing the correspondence of $\Phi$ and $\Phi'$ gives, by
Step 7,
\[
\sum_jp'_1(L(f'_j))\otimes\bigotimes_{v\in T'}p'_{1,v}(\Lambda_v\phi'_{j,v})
=c\prod_{v\in T'}w_v\sum_ip_1(e_\infty(f_i))\otimes\bigotimes_{v\in T'}p_{1,v}(e_v\phi_{i,v}),
\]
and as $u^{-1}\lambda_0c\prod_{v\in T'}w_v=\kappa c\sigma L_0$, the functions
$\kappa c\sigma L_0\,e(\Phi)$ and $\operatorname{cone}_\beta(\Phi')$ correspond
for $Q_1$. This is (ii).
\end{proof}
% Lean: PositivityTheory.ThKG.Loc.locB_of Hyp Hyp2 ConeFacB StdB (LocalTheory/LocBOf.lean).
% Step 1 is one_const with coneI_eq and sideSo.ontoI; Step 2 is nz, j1 and j2 (Hyp.inj, with
% SS); Step 3 is b1 (br_of, fac_ev, ct1 = Hyp2.ctBZ, Hyp.split, Hyp.sinv); Step 4 is fl
% (Hyp2.floc), LambdaS, hvac, cf = fac_cone and b2 (ct2 = Hyp2.ctSC); Step 5 is aug, with co
% = Hyp.coinv; Step 6 is std (StdB), pt1 and pt2 (bump_pt with Hyp2.bumpB and nz_pure); Step 7
% is the two calls of lams with sideV and sideSo, and e3; Step 8 is Hyp2.inf, Hyp2.ram,
% Hyp2.unr (hw); Step 9 is the span induction on So with pure_mem_so, and cone_sum, cfac,
% tens_step and prod_w. Size of locB_of and the generic lemmas it calls: about 1,100 lines of
% closure, most of them shared with Lemma asm-theta-agree-local-alpha and the sublemmas
% above. The two lemmas together have about 14,100 lines beyond the other properties of
% Theorem K, counted by the declarations reachable from them and from no other property.

\subsection{Sets of differences and their lower parts}\label{subsection-asm-differences}

\begin{definition}\label{asm-theta-agree-end-notions}
Let $E$ be a CM field with maximal totally real subfield $F$ and complex
conjugation $c$, let $B$ be a theta base for $E$
(Definition~\ref{asm-theta-agree-models}(1)), and let $Y$, $Y_1$ and $W$ be
lattice spaces (Definition~\ref{asm-theta-agree-doubling}(1)) of ranks $m$,
$m_1$ and $d$, with matrices $J$, $J_1$ and $J_W$. Write
$h(u,u')=u^{\mathsf T}J\,c(u')$ and $h_W(u,u')=u^{\mathsf T}J_W\,c(u')$ for
vectors $u,u'$ over $\mathbf A_E$ of the appropriate sizes, and
$|t|=|N_{E/F}(t)|_{\mathbf A}$ for $t\in\mathbf A_E^\times$. The integral of a
function that is not integrable is taken to be $0$.
\begin{enumerate}
\item A \emph{plane} $H$ in $Y$ over $Y_1$ consists of vectors $x,y\in E^m$ and
  a matrix $\iota\in M_{m\times m_1}(E)$ such that $m=m_1+2$,
  $h(x,x)=h(y,y)=0$, $h(x,y)=1$, $h(\iota u,x)=h(\iota u,y)=0$ for every
  $u\in E^{m_1}$, $\iota^{\mathsf T}J\,c(\iota)=J_1$ and $\det J_1\neq0$. Levi
  elements, the unipotent group $N_H$, constant terms along $H$ and laws are
  defined for a plane by the formulas of
  Definition~\ref{asm-theta-agree-planes}(2).
\item $K^\theta_W$ is the set of $k\in U_W$ such that $\pi_w(k)$ commutes with
  $\theta_{W,w}$ for every infinite place $w$ of $F$.
\item A function $f$ on $U_W$ is \emph{automorphic} if it is continuous and
  $f(\gamma h)=f(h)$ for all $\gamma\in U_W(F)$ and $h\in U_W$. Let $\chi$ be a
  homomorphism from $\mathbf A_E^\times$ to the complex numbers of absolute
  value $1$, and let $z\in\mathbf R$. The function $f$ is \emph{pure of law}
  $(\chi,z)$ if there are a lattice space $W_1$, a plane $H'$ in $W$ over
  $W_1$ and a function on $U_W$ that is the constant term of $f$ along $H'$ and
  has the law $(\chi,z)$.
\item For a nonzero $E$-subspace $U\subset E^d$ with $h_W(u,u')=0$ for all
  $u,u'\in U$, let $N_U$ be the subgroup of $U_W$ generated by the $g\in U_W$
  with $gu=u$ for every $u\in U$ and with $gz-z$ in the $\mathbf A_E$-span of
  $U$ for every $z\in E^d$ such that $h_W(z,u)=0$ for all $u\in U$. A function
  $f$ on $U_W$ is \emph{cuspidal} if $\int_{\mathcal F}f(gh)\,d\nu(g)=0$ for
  every such $U$, every Haar measure $\nu$ on $N_U$, every fundamental domain
  $\mathcal F$ for $N_U\cap U_W(F)$ in $N_U$ with respect to $\nu$, and every
  $h\in U_W$. A function $g$ on $U_W$ is \emph{orthogonal} to $f$ if
  $\int_{\mathcal D}g(h)\overline{f(h)}\,d\mu(h)=0$ for every Haar measure
  $\mu$ on $U_W$ and every fundamental domain $\mathcal D$ for $U_W(F)$ in
  $U_W$ with respect to $\mu$.
\item Let $(B,G,\mathcal X,\omega,\omega')$ be a triple
  (Definition~\ref{asm-theta-agree-triple}(4)), with models $M$ and $M'$ and
  theta integrals $I$ and $I'$. Its \emph{set of differences} $\Sigma$ is the
  set of the functions $I(\Phi)-I'(\Phi')$ on $U_W$ with
  $\Phi\in\operatorname{PD}(M)$ and
  $\Phi'\in\operatorname{PD}(M')\cap\mathcal S^\circ(M')$.
\item Let $H'$ be a plane in $W$ over $W_1$, let $T$ be a function on $U_W$
  and let $k\in U_W$. A function $f_1$ on $U_{W_1}$ is the \emph{lower part}
  of $T$ at $k$ if $f_1(h_1)=T(lk)$ for every $h_1\in U_{W_1}$ and every Levi
  element $l$ for $(1,h_1)$.
\item Let $G$ and $G_1$ be triple data (Definition~\ref{asm-theta-agree-triple}(2))
  of ranks $(b+2,a')$ and $(b,a')$, let $H$ be an adapted plane in the first
  space $V$ of $G$ over the first space $V_1$ of $G_1$, and let $H'$ be one in
  $W$ over $W_1$ (Definition~\ref{asm-theta-agree-planes}(1)). Then
  $(G,G_1,H,H')$ is a \emph{plane-removal configuration} if $V'_1=V'$,
  $q_1=q$, $p_1+1=p$, and $H$ and $H'$ are at positions $(0,b+1)$. The
  $\alpha$-configurations of Definition~\ref{asm-theta-agree-planes}(4) are the
  plane-removal configurations with $q=0$.
\end{enumerate}
\end{definition}
\leanref{Development/WeilBound/ThetaLifts/EigenvalueEndgame/PlaneRemovalConfiguration.lean}{PositivityTheory.ThKEnd.IsPlane}
\leanref{Development/WeilBound/ThetaLifts/EigenvalueEndgame/PlaneRemovalConfiguration.lean}{PositivityTheory.ThKEnd.FinK}
\leanref{Development/WeilBound/ThetaLifts/EigenvalueEndgame/PlaneRemovalConfiguration.lean}{PositivityTheory.ThKEnd.IsAut}
\leanref{Development/WeilBound/ThetaLifts/EigenvalueEndgame/PlaneRemovalConfiguration.lean}{PositivityTheory.ThKEnd.Pure}
\leanref{Development/WeilBound/ThetaLifts/EigenvalueEndgame/PlaneRemovalConfiguration.lean}{PositivityTheory.ThKEnd.unipU}
\leanref{Development/WeilBound/ThetaLifts/EigenvalueEndgame/PlaneRemovalConfiguration.lean}{PositivityTheory.ThKEnd.IsCusp}
\leanref{Development/WeilBound/ThetaLifts/EigenvalueEndgame/PlaneRemovalConfiguration.lean}{PositivityTheory.ThKEnd.Perp}
\leanref{Development/WeilBound/ThetaLifts/EigenvalueEndgame/PlaneRemovalConfiguration.lean}{PositivityTheory.ThKEnd.Sig}
\leanref{Development/WeilBound/ThetaLifts/EigenvalueEndgame/PlaneRemovalConfiguration.lean}{PositivityTheory.ThKEnd.IsLow}
\leanref{Development/WeilBound/ThetaLifts/EigenvalueEndgame/PlaneRemovalConfiguration.lean}{PositivityTheory.ThKEnd.PlaneRemovalConfiguration}

\begin{lemma}\label{asm-theta-agree-end-unipotent}
Let $E$ be a CM field with a theta base $B$
(Definition~\ref{asm-theta-agree-models}(1)), let $W$ be a skew-Hermitian
lattice space of rank $d$ and some signature
(Definition~\ref{asm-theta-agree-doubling}(1)), with matrix $J_W$ and
$h_W(u,u')=u^{\mathsf T}J_W\,c(u')$, and let $H=(x,y,\iota)$ be a plane in $W$
over a lattice space $W_1$ (Definition~\ref{asm-theta-agree-end-notions}(1)).
Then:
\begin{enumerate}
\item $N_H$ is a closed subgroup of $U_W$, it is locally compact and second
  countable, and $N_H\cap U_W(F)$ is countable; every Haar measure on $N_H$ is
  invariant under right translations; for every Haar measure $\nu$ on $N_H$
  there is a measurable fundamental domain for $N_H\cap U_W(F)$ in $N_H$, and
  every such fundamental domain $\mathcal F$ has $0<\nu(\mathcal F)<\infty$;
  some Haar measure on $N_H$ has a fundamental domain of measure $1$; and there
  is a compact set $C\subset N_H$ such that every $u\in N_H$ has $\gamma u\in C$
  for some $\gamma\in N_H\cap U_W(F)$;
\item for all $t\in\mathbf A_E^\times$ and $h_1\in U_{W_1}$ there is a Levi
  element for $(t,h_1)$, and $mum^{-1}\in N_H$ for every Levi element $m$ and
  every $u\in N_H$;
\item an element $u\in U_W$ lies in $N_H$ if and only if $ux=x$ and
  $uz-z\in\mathbf A_Ex$ for every $z\in\mathbf A_E^d$ with $h_W(z,x)=0$;
\item for all nonzero $x',x''\in E^d$ with $h_W(x',x')=h_W(x'',x'')=0$ there is
  $\gamma\in U_W(F)$ with $\gamma x'=x''$.
\end{enumerate}
\end{lemma}
\leanref{Development/WeilBound/ThetaLifts/EigenvalueEndgame/UnipotentGroups/UnipUGood.lean}{PositivityTheory.ThKEnd.Grp.UnipGood}
\leanref{Development/WeilBound/ThetaLifts/EigenvalueEndgame/UnipotentGroups/IsFrPlane.lean}{PositivityTheory.ThKEnd.Grp.unipGood}
\leanref{Development/WeilBound/ThetaLifts/EigenvalueEndgame/UnipotentGroups/UnipUGood.lean}{PositivityTheory.ThKEnd.Grp.LeviPlane}
\leanref{Development/WeilBound/ThetaLifts/EigenvalueEndgame/UnipotentGroups/IsFrPlane.lean}{PositivityTheory.ThKEnd.Grp.leviPlane}
\leanref{Development/WeilBound/ThetaLifts/EigenvalueEndgame/UnipotentGroups/IsFrPlane.lean}{PositivityTheory.ThKEnd.Grp.mem_unip}
\leanref{Development/WeilBound/ThetaLifts/EigenvalueEndgame/UnipotentGroups/UnipUGood.lean}{PositivityTheory.ThKEnd.Grp.WittVec}
\leanref{Development/WeilBound/ThetaLifts/EigenvalueEndgame/UnipotentGroups/IsFrPlane.lean}{PositivityTheory.ThKEnd.Grp.wittVec}
\begin{proof}
Write $h=h_W$, let $J_1$ be the matrix of $W_1$, and let $A=(\iota\ x\ y)$ be
the matrix with the columns of $\iota$ followed by $x$ and $y$. As $W$ is
skew-Hermitian, $h(u',u)=-c(h(u,u'))$; so $h(y,x)=-1$ and
$h(x,\iota u)=h(y,\iota u)=0$, and the Gram matrix $A^{\mathsf T}J_W\,c(A)$ is
$\operatorname{diag}\bigl(J_1,\left(\begin{smallmatrix}0&1\\-1&0\end{smallmatrix}\right)\bigr)$.
It is invertible, as $\det J_1\neq0$, so $A$ is invertible and the columns of
$A$ are a basis of $\mathbf A_E^d$ over $\mathbf A_E$. If
$z=\iota u+sx+s'y$ has $h(z,x)=0$, then $0=h(z,x)=-s'$; so
\begin{enumerate}
\item[(a)] every $z\in\mathbf A_E^d$ with $h(z,x)=0$ is $\iota u+sx$ with
  $u\in\mathbf A_E^{d_1}$ and $s\in\mathbf A_E$.
\end{enumerate}

(3) Let $S$ be the set of $u\in U_W$ with $ux=x$ and $uz-z\in\mathbf A_Ex$ for
every $z\in\mathbf A_E^d$ with $h(z,x)=0$. Then $S$ is the generating set of
$N_H$ (Definition~\ref{asm-theta-agree-planes}(2)): if $u\in S$, then
$u\iota z\in\iota z+\mathbf A_Ex$, as $h(\iota z,x)=0$; conversely, if $u$ is a
generator and $z=\iota u'+sx$ as in (a), then
$uz-z=(u\iota u'-\iota u')+s(ux-x)\in\mathbf A_Ex$. And $S$ is a subgroup. Let
$u,u'\in S$ and $h(z,x)=0$, with $uz=z+bx$ and $u'z=z+b'x$. Then
$uu'x=x$ and $uu'z=uz+b'ux=z+(b+b')x$; and $u^{-1}x=x$ and
$z=u^{-1}(z+bx)=u^{-1}z+bx$, so $u^{-1}z=z-bx$. Hence $N_H=S$, which is (3).

(1) By (a) and (3), $u\in U_W$ lies in $N_H$ if and only if $ux=x$ and
$u\iota_j-\iota_j\in\mathbf A_Ex$ for the columns $\iota_j$ of $\iota$. A vector of $\mathbf A_E^d$ lies in $\mathbf A_Ex$ if and
only if its coordinates along the columns of $\iota$ and along $y$ in the basis
$A$ vanish, a closed condition. So $N_H$ is closed in $U_W$. By
Lemma~\ref{asm-theta-agree-convergence-signature}, $J_W$ is nondegenerate
$\epsilon$-Hermitian, so by Theorem~\ref{asm-theta-agree-convergence-siegel}(1)
$U_W$ is locally compact, second countable and Hausdorff and $U_W(F)$ is a
countable discrete subgroup. Hence $N_H$ is locally compact and second
countable, and $\Gamma=N_H\cap U_W(F)$ is countable and discrete in $N_H$.

\emph{A compact set $C$.} We show that there is a compact set $C\subset N_H$
such that every $u\in N_H$ has $\gamma u\in C$ for some $\gamma\in\Gamma$. Let
$x^\perp=\{z\in E^d:h(z,x)=0\}$, and let $e,f\in M_d(E)$ be projections onto
$Ex$ and onto $x^\perp$. Call $X\in M_d(\mathbf A_E)$ \emph{triangular} if $Xe=0$,
$eXf=Xf$ and $fX=X$, that is, if $X$ kills $x$, maps
$x^\perp_{\mathbf A}=\mathbf A_E\otimes_Ex^\perp$ into $\mathbf A_Ex$ and has image in
$x^\perp_{\mathbf A}$.
\begin{enumerate}
\item[(b)] An element $u\in U_W$ lies in $N_H$ if and only if $u-1$ is
  triangular. Indeed, if $ux=x$, then $h((u-1)z,x)=h(uz,ux)-h(z,x)=0$ for every
  $z$, so the image of $u-1$ lies in $x^\perp_{\mathbf A}$; and the other two
  conditions are those of (3).
\item[(c)] Sums, scalar multiples and products of triangular matrices are
  triangular. If $X$ and $Y$ are triangular, then $XYf=X(eYf)=0$, so
  $XYZ=XYfZ=0$ for every triangular $Z$; and if $Xf=0$, then $XY=XfY=0$.
\item[(d)] For triangular $X$ put $\exp X=1+X+\frac12X^2$, and for $u$ with
  $u-1$ triangular put $\log u=(u-1)-\frac12(u-1)^2$. As products of three
  triangular matrices vanish, $\exp$ and $\log$ are inverse to each other,
  $\exp X-1$ and $\log u$ are triangular, $(\exp X-1)f=Xf$ and
  $(\log u)f=(u-1)f$, and $(gg'-1)f=(g-1)f+(g'-1)f$ when $g-1$ and $g'-1$ are
  triangular. The same identities show that $\exp X\in U_W$ if
  $X^{\mathsf T}J_W+J_Wc(X)=0$: then $(\exp X)^{\mathsf T}J_W=J_W\exp(-c(X))$
  and $\exp(-c(X))\exp(c(X))=1$; and that $X=\log u$ satisfies
  $X^{\mathsf T}J_W+J_Wc(X)=0$ if $u\in N_H$.
\end{enumerate}
Let $\mathfrak n$ be the set of triangular $X$ with
$X^{\mathsf T}J_W+J_Wc(X)=0$, and $\mathfrak z$ the set of $X\in\mathfrak n$
with $Xf=0$; let $\mathfrak n(F)$ and $\mathfrak z(F)$ be their subsets of
matrices over $E$. Both $\mathfrak n$ and $\mathfrak z$ are kernels of
$\mathbf A_F$-linear maps on $M_d(\mathbf A_E)=M_d(E)\otimes_F\mathbf A_F$
obtained by base change from $F$-linear maps on $M_d(E)$, whose kernels are
$\mathfrak n(F)$ and $\mathfrak z(F)$. For an $F$-linear map $T\colon V\to V'$
of finite-dimensional $F$-vector spaces, with base change $T_{\mathbf A}$,
there is a compact set $C_T\subset\ker T_{\mathbf A}$ with
$\ker T_{\mathbf A}=\ker T+C_T$: if $b_1,\ldots,b_k$ is a basis of $\ker T$,
then $\ker T_{\mathbf A}=\sum_j\mathbf A_Fb_j$, as $\mathbf A_F$ is flat over
$F$, and $\mathbf A_F=F+K$ for a compact set $K$, as $\mathbf A_F/F$ is compact;
take $C_T=\sum_jKb_j$. So there are compact sets $C_1\subset\mathfrak n$ and
$C_2\subset\mathfrak z$ with $\mathfrak n=\mathfrak n(F)+C_1$ and
$\mathfrak z=\mathfrak z(F)+C_2$.

By (b) and (d), $\exp$ maps $\mathfrak n$ into $N_H$ and $\mathfrak n(F)$ into
$\Gamma$, and it is continuous. Let $C=\exp(C_2)\exp(C_1)$, a compact subset of
$N_H$, and let $u\in N_H$. By (d), $\log u\in\mathfrak n$, so
$\log u=v+k$ with $v\in\mathfrak n(F)$ and $k\in C_1$. Put $\gamma=\exp v$,
$\kappa=\exp k$ and $z=\gamma^{-1}u\kappa^{-1}\in N_H$. By (d),
\[
vf+kf=(\log u)f=(u-1)f=(\gamma z\kappa-1)f=vf+(z-1)f+kf,
\]
so $(z-1)f=0$. By (c), $(z-1)^2=0$, so $\log z=z-1$ lies in $\mathfrak z$, and
$z-1=w+k'$ with $w\in\mathfrak z(F)$ and $k'\in C_2$. By (c),
$w^2=k'^2=wk'=0$, so $\exp w\,\exp k'=(1+w)(1+k')=1+w+k'=z$. Hence
$u=\gamma\exp(w)\exp(k')\kappa$, with $\gamma\exp(w)\in\Gamma$ and
$\exp(k')\kappa\in C$.

Given $C$, the remaining assertions of (1) follow. Let $\nu$ be a Haar measure
on $N_H$. As $\Gamma$ is a countable discrete subgroup of the locally compact,
second countable group $N_H$, it has a measurable fundamental domain with
respect to $\nu$. Let $\mathcal F$ be any fundamental domain. Its translates
$\gamma\mathcal F$, $\gamma\in\Gamma$, cover $N_H$ and are countably many, so
$\nu(\mathcal F)>0$. Every $u\in\mathcal F$ lies in $\gamma C$ for some
$\gamma\in\Gamma$, so
$\nu(\mathcal F)\leq\sum_{\gamma\in\Gamma}\nu(\gamma C\cap\mathcal F)=\nu(C)<\infty$,
the last equality because the sets $C\cap\gamma^{-1}\mathcal F$ partition $C$
up to a null set and $\nu$ is invariant under left translations. Rescaling
$\nu$ gives a Haar measure with a fundamental domain of measure $1$. For
$g\in N_H$, the measure $B\mapsto\nu(Bg)$ is again a Haar measure, so
$\nu(Bg)=\Delta(g)\nu(B)$ for some $\Delta(g)>0$. The set $\mathcal Fg$ is a
fundamental domain for $\Gamma$, as $\gamma(\mathcal Fg)=(\gamma\mathcal F)g$,
and two fundamental domains have the same measure, by the computation in the
proof of Lemma~\ref{asm-theta-agree-anisotropic-invariance} for the constant
function $1$. So $\Delta(g)\nu(\mathcal F)=\nu(\mathcal F)$, and $\Delta(g)=1$
as $0<\nu(\mathcal F)<\infty$: $\nu$ is invariant under right translations.

(2) By the computation above, the proof of Lemma~\ref{asm-theta-agree-levi},
with the invertibility of $A$ shown here in place of the argument through the
lattice, gives for $t\in\mathbf A_E^\times$ and $h_1\in U_{W_1}$ the Levi
element $m=ADA^{-1}$ for $(t,h_1)$, with
$D=\operatorname{diag}(h_1,t,c(t)^{-1})$. Let $m$ be a Levi element for
$(t,h_1)$ and $u\in N_H$. Then $m^{-1}x=t^{-1}x$, so
$mum^{-1}x=t^{-1}mux=t^{-1}mx=x$. Let $h(z,x)=0$. As $m$ preserves $h$,
$h(m^{-1}z,x)=h(z,mx)=c(t)h(z,x)=0$, so by (3) $um^{-1}z=m^{-1}z+ax$ for some
$a$, and $mum^{-1}z=z+atx$. By (3), $mum^{-1}\in N_H$.

(4) This is the argument of Lemma~\ref{asm-theta-agree-split-witt}(2), for the
skew-Hermitian form $h=h_W$ on $E^d$, which is nondegenerate by
Lemma~\ref{asm-theta-agree-convergence-signature}. It suffices to find
$\gamma\in U_W(F)$ with $\gamma x'=x''$.

\emph{A transvection.} Suppose first that $b=h(x',x'')\neq0$. Put $w=x'-x''$
and $g(v)=v+b^{-1}h(v,w)\,w$. Then $h(x',w)=-b$, so $gx'=x'-w=x''$. For
$v,v'\in E^d$, using $h(w,v')=-c(h(v',w))$,
\[
h(gv,gv')=h(v,v')+h(v,w)\,c(h(v',w))\bigl(c(b)^{-1}-b^{-1}+b^{-1}c(b)^{-1}h(w,w)\bigr),
\]
and $h(w,w)=-h(x',x'')-h(x'',x')=c(b)-b$, so the bracket is $0$. Hence $g$
preserves $h$; as $h$ is nondegenerate, $g$ is invertible, and $g\in U_W(F)$.

\emph{Two transvections.} Suppose that $h(x',x'')=0$, so also
$h(x'',x')=0$. As $h$ is nondegenerate there are $a',a''$ with
$h(x',a')\neq0$ and $h(x'',a'')\neq0$; one of $a'$, $a''$ and $a'+a''$ is a
vector $v$ with $h(x',v)\neq0$ and $h(x'',v)\neq0$. Put $y'=v+\lambda x'$ with
$\lambda=-h(v,v)/(2h(x',v))$. As $c(h(v,v))=-h(v,v)$ and
$h(v,x')=-c(h(x',v))$,
\[
h(y',y')=h(v,v)+\lambda h(x',v)-c\bigl(\lambda h(x',v)\bigr)=h(v,v)-\tfrac12h(v,v)-\tfrac12h(v,v)=0,
\]
and $h(x',y')=h(x',v)\neq0$, $h(y',x'')=h(v,x'')=-c(h(x'',v))\neq0$; in
particular $y'\neq0$. Two transvections as above take $x'$ to $y'$ and $y'$ to
$x''$, and their product is $\gamma$.
\end{proof}
% Lean: (a) is fr_inv and exists_decomp_of_fa_incV_x_eq_zero (IsFrPlane.lean), with form_yx;
% (3) is mem_unip: unipLine and line_iff identify the generating set of N_H with that of the
% group unipU of the line E x, and unipU_eq shows that this generating set is a subgroup,
% through gen_iff and Tri.grp, by describing it with projection matrices (the book checks
% closure under products and inverses directly); par_of and Par give the form of (3). (1) is
% unipGood = unipUGood: triG_closed, triG_cov and unipGood_of (lcsc_UA, countable_ratU,
% isDiscreteWithDomains_of_discrete, fd_ne, fd_fin, right_of_fd). In the compact set C, (b) is
% gen_iff, (c) is Tri.mul, Tri.mf, Tri.three and Tri.zero_mul, (d) is ex, ex_neg, Tri.ex,
% Tri.ex_f, Tri.log, Tri.grp, ex_unit and log_skew; the kernels are eqs and lie, with
% lie_cocompact for p = 0 and p = f from ker_cocompact (KerCocompact.lean); the rest is
% triG_cov and expG_rat. Lemma asm-theta-agree-split-inputs(1)(b) rests on the same triG_cov. (2) is leviPlane, with
% isFr_plane, IsFr.levi and exists_UA, and Par.conj. (4) is wittVec and witt_step.

\begin{lemma}\label{asm-theta-agree-end-invariant}
Let $E$ be a CM field with a theta base $B$
(Definition~\ref{asm-theta-agree-models}(1)), let $W$ be a skew-Hermitian
lattice space of some signature, let $H$ be a plane in $W$ over a lattice space
$W_1$ (Definition~\ref{asm-theta-agree-end-notions}(1)), and let $f$ be an
automorphic function on $U_W$ (Definition~\ref{asm-theta-agree-end-notions}(3)).
\begin{enumerate}
\item For every Haar measure $\nu$ on $N_H$, every fundamental domain
  $\mathcal F$ for $N_H\cap U_W(F)$ in $N_H$ with respect to $\nu$ and every
  $h\in U_W$, the function $u\mapsto f(uh)$ is integrable on $\mathcal F$.
\item If $c$ is the constant term of $f$ along $H$, then $c(uh)=c(h)$ for all
  $u\in N_H$ and $h\in U_W$.
\end{enumerate}
\end{lemma}
\leanref{Development/WeilBound/ThetaLifts/EigenvalueEndgame/UnipotentGroups/ExistsAbsGt.lean}{PositivityTheory.ThKEnd.Grp.UnipGood.integrable}
\leanref{Development/WeilBound/ThetaLifts/EigenvalueEndgame/UnipotentGroups/ExistsAbsGt.lean}{PositivityTheory.ThKEnd.Grp.ct_inv}
\begin{proof}
(1) The function $\varphi(u)=f(uh)$ is continuous on $N_H$. Let $C$ be the
compact set of Lemma~\ref{asm-theta-agree-end-unipotent}(1). Then $\varphi$ is
bounded on $C$, say by $b$. For $u\in N_H$ choose
$\gamma\in N_H\cap U_W(F)$ with $\gamma u\in C$. As $f$ is automorphic,
$\varphi(u)=f(\gamma uh)=\varphi(\gamma u)$, so $|\varphi(u)|\leq b$. By
Lemma~\ref{asm-theta-agree-end-unipotent}(1), $\nu(\mathcal F)<\infty$, so
$\varphi$ is integrable on $\mathcal F$.

(2) By Lemma~\ref{asm-theta-agree-end-unipotent}(1) there are a Haar measure
$\nu$ on $N_H$ and a fundamental domain $\mathcal F$ for $N_H\cap U_W(F)$ of
measure $1$. Fix $h$ and put $\varphi(v)=f(vh)$. It is invariant under left
translation by $N_H\cap U_W(F)$, as $f$ is automorphic, and integrable on every
fundamental domain, by (1). Let $u\in N_H$. As $\nu$ is invariant under right
translations, right translation by $u$ carries $\mathcal F$ to a fundamental
domain $\mathcal Fu$ and preserves $\nu$. So
\[
c(uh)=\int_{\mathcal F}\varphi(vu)\,d\nu(v)=\int_{\mathcal Fu}\varphi\,d\nu
=\int_{\mathcal F}\varphi\,d\nu=c(h).
\]
The third equality holds because a function invariant on the left under
$N_H\cap U_W(F)$ has the same integral over any two fundamental domains, by the
computation in the proof of Lemma~\ref{asm-theta-agree-anisotropic-invariance}.
\end{proof}

\begin{lemma}\label{asm-theta-agree-end-constant-terms}
Let $E$ be a CM field with a theta base $B$
(Definition~\ref{asm-theta-agree-models}(1)), let $W$ be a skew-Hermitian
lattice space of some signature, and let $H$ be a plane in $W$ over a lattice
space $W_1$ (Definition~\ref{asm-theta-agree-end-notions}(1)).
\begin{enumerate}
\item Let $f$ and $g$ be automorphic functions on $U_W$
  (Definition~\ref{asm-theta-agree-end-notions}(3)), let $c_f$ and $c_g$ be the
  constant terms of $f$ and $g$ along $H$, and let $c\in\mathbf C$. Then
  $c_f-c\,c_g$ is the constant term of $f-c\,g$ along $H$.
\item Let $f$ be an automorphic function on $U_W$, let $\chi$ and $\chi'$ be
  homomorphisms from $\mathbf A_E^\times$ to the complex numbers of absolute
  value $1$, and let $z\neq z'$ be real numbers. If $T+T'$ is the constant term
  of $f$ along $H$, $T$ has the law $(\chi,z)$ and $T'$ has the law
  $(\chi',z')$, then $T(uh)=T(h)$ and $T'(uh)=T'(h)$ for all $u\in N_H$ and
  $h\in U_W$.
\end{enumerate}
\end{lemma}
\leanref{Development/WeilBound/ThetaLifts/EigenvalueEndgame/PlaneRemovalConfiguration.lean}{PositivityTheory.ThKEnd.CTBasic}
\leanref{Development/WeilBound/ThetaLifts/EigenvalueEndgame/UnipotentGroups/ExistsAbsGt.lean}{PositivityTheory.ThKEnd.Grp.ctBasic}
\leanref{Development/WeilBound/ThetaLifts/EigenvalueEndgame/UnipotentGroups/ExistsAbsGt.lean}{PositivityTheory.ThKEnd.Grp.exists_abs_gt}
\begin{proof}
(1) Let $\nu$ be a Haar measure on $N_H$ and $\mathcal F$ a fundamental domain
for $N_H\cap U_W(F)$ in $N_H$ of measure $1$, and let $h\in U_W$. By
Lemma~\ref{asm-theta-agree-end-invariant}(1), $u\mapsto f(uh)$ and
$u\mapsto g(uh)$ are integrable on $\mathcal F$. So
\[
\int_{\mathcal F}(f-c\,g)(uh)\,d\nu(u)=\int_{\mathcal F}f(uh)\,d\nu(u)-c\int_{\mathcal F}g(uh)\,d\nu(u)=c_f(h)-c\,c_g(h).
\]

(2) Let $u\in N_H$ and $h\in U_W$. By Lemma~\ref{asm-theta-agree-end-invariant}(2),
the constant term $T+T'$ is invariant under left translation by $N_H$, so
\[
T(uh)+T'(uh)=T(h)+T'(h).\tag{a}
\]
Let $w_0$ be an infinite place of $F$, and let $t\in\mathbf A_E^\times$ have
component $2$ at $w_0$ and $1$ at every other place. Then
$|t|=|N_{E/F}(2)|_{w_0}=4>1$. By Lemma~\ref{asm-theta-agree-end-unipotent}(2)
there is a Levi element $m$ for $(t,1)$, and $mum^{-1}\in N_H$. Applying (a)
to $mum^{-1}$ and $mh$, and then the two laws, gives
\[
\chi(t)|t|^zT(uh)+\chi'(t)|t|^{z'}T'(uh)=\chi(t)|t|^zT(h)+\chi'(t)|t|^{z'}T'(h).\tag{b}
\]
Put $\delta=T(uh)-T(h)$ and $\delta'=T'(uh)-T'(h)$. By (a), $\delta'=-\delta$,
and (b) becomes $(\chi(t)|t|^z-\chi'(t)|t|^{z'})\delta=0$. The two terms in
the bracket have absolute values $|t|^z$ and $|t|^{z'}$, which differ because
$|t|>1$ and $z\neq z'$. So $\delta=0$, and then $\delta'=0$.
\end{proof}

\begin{lemma}\label{asm-theta-agree-end-move}
Let $E$ be a CM field with a theta base $B$
(Definition~\ref{asm-theta-agree-models}(1)), let $W$ be a skew-Hermitian
lattice space of some signature, let $H$ be a plane in $W$ over a lattice space
$W_1$ (Definition~\ref{asm-theta-agree-end-notions}(1)), let $\chi$ be a
homomorphism from $\mathbf A_E^\times$ to the complex numbers of absolute value
$1$, and let $z\in\mathbf R$. If $f$ is an automorphic function on $U_W$ that
is pure of law $(\chi,z)$ (Definition~\ref{asm-theta-agree-end-notions}(3)),
then there is a function $c$ on $U_W$ that is the constant term of $f$ along
$H$ and has the law $(\chi,z)$.
\end{lemma}
\leanref{Development/WeilBound/ThetaLifts/EigenvalueEndgame/PlaneRemovalConfiguration.lean}{PositivityTheory.ThKEnd.LineMove}
\leanref{Development/WeilBound/ThetaLifts/EigenvalueEndgame/UnipotentGroups/ExistsAbsGt.lean}{PositivityTheory.ThKEnd.Grp.lineMove}
\leanref{Development/WeilBound/ThetaLifts/EigenvalueEndgame/UnipotentGroups/ExistsAbsGt.lean}{PositivityTheory.ThKEnd.Grp.ct_conj}
\begin{proof}
Write $H=(x,y,\iota)$, $d$ for the rank of $W$ and $h_W$ for its form. As $f$ is
pure, there are a lattice space $W_2$, a plane $H_2=(x_2,y_2,\iota_2)$ in $W$
over $W_2$ and a function $c_2$ that is the constant term of $f$ along $H_2$
and has the law $(\chi,z)$.

\emph{A rational element.} The vectors $x$ and $x_2$ are isotropic, and
nonzero because $h_W(x,y)=h_W(x_2,y_2)=1$. By
Lemma~\ref{asm-theta-agree-end-unipotent}(4) there is $\gamma\in U_W(F)$ with
$\gamma x_2=x$.

\emph{The unipotent groups.} For $u\in U_W$, $u\in N_{H_2}$ if and only if
$\gamma u\gamma^{-1}\in N_H$. Indeed, $\gamma u\gamma^{-1}x=\gamma ux_2$, which
is $x$ if and only if $ux_2=x_2$. As $\gamma$ preserves $h_W$,
$h_W(z,x)=0$ if and only if $h_W(\gamma^{-1}z,x_2)=0$, and
$\gamma u\gamma^{-1}z-z=\gamma(u\gamma^{-1}z-\gamma^{-1}z)$ lies in
$\mathbf A_Ex$ if and only if $u\gamma^{-1}z-\gamma^{-1}z$ lies in
$\mathbf A_Ex_2$. Lemma~\ref{asm-theta-agree-end-unipotent}(3), for $H$ and for
$H_2$, gives the claim.

\emph{The constant term.} Put $c(h)=c_2(\gamma^{-1}h)$. Let $\nu$ be a Haar
measure on $N_H$ and $\mathcal F$ a fundamental domain for $N_H\cap U_W(F)$ in
$N_H$ of measure $1$. The map $v\mapsto\gamma v\gamma^{-1}$ is an isomorphism
of topological groups from $N_{H_2}$ onto $N_H$, and it maps
$N_{H_2}\cap U_W(F)$ onto $N_H\cap U_W(F)$, as $\gamma\in U_W(F)$. So it
carries a Haar measure $\nu_2$ on $N_{H_2}$ to $\nu$, and
$\mathcal F_2=\gamma^{-1}\mathcal F\gamma$ is a fundamental domain for
$N_{H_2}\cap U_W(F)$ in $N_{H_2}$ with $\nu_2(\mathcal F_2)=1$. As $f$ is
automorphic, for $h\in U_W$
\[
\int_{\mathcal F}f(uh)\,d\nu(u)=\int_{\mathcal F_2}f(\gamma v\gamma^{-1}h)\,d\nu_2(v)
=\int_{\mathcal F_2}f(v\gamma^{-1}h)\,d\nu_2(v)=c_2(\gamma^{-1}h)=c(h).
\]
So $c$ is the constant term of $f$ along $H$.

\emph{Levi elements.} Let $m$ be a Levi element for $(t,1)$ along $H$. Then
$mx=tx$, and $mz-z\in\mathbf A_Ex$ for every $z\in\mathbf A_E^d$ with
$h_W(z,x)=0$. Indeed, the Gram matrix of $x$, $y$ and the columns of $\iota$ is
$\operatorname{diag}\bigl(J_1,\left(\begin{smallmatrix}0&1\\-1&0\end{smallmatrix}\right)\bigr)$,
since $h_W(y,x)=-c(h_W(x,y))=-1$ as $W$ is skew-Hermitian; it is invertible,
as $\det J_1\neq0$. So these vectors form a basis, and
$z=sx+\iota u+s'y$ with $s,s'\in\mathbf A_E$. Then $0=h_W(z,x)=-s'$, and
$mz=stx+\iota u=z+s(t-1)x$. The same holds for a Levi element for $(t,1)$
along $H_2$, with $x_2$ in place of $x$.

\emph{The law.} Let $m$ be a Levi element for $(t,1)$ along $H$, and $m_0$ one
along $H_2$, which exists by Lemma~\ref{asm-theta-agree-end-unipotent}(2). Put
$n=\gamma^{-1}m\gamma m_0^{-1}$. Then $nx_2=\gamma^{-1}m\gamma t^{-1}x_2=x_2$.
Let $z\in\mathbf A_E^d$ with $h_W(z,x_2)=0$. Then $m_0z=z+bx_2$ for some $b$,
so $m_0^{-1}z=z-bt^{-1}x_2$; and $h_W(\gamma z,x)=0$, so $m\gamma z=\gamma z+ax$
for some $a$, and $\gamma^{-1}m\gamma z=z+ax_2$. Hence
$nz=z+ax_2-bt^{-1}tx_2=z+(a-b)x_2$. By
Lemma~\ref{asm-theta-agree-end-unipotent}(3), $n\in N_{H_2}$. By
Lemma~\ref{asm-theta-agree-end-invariant}(2) for $H_2$, and the law of $c_2$,
\[
c(mh)=c_2(\gamma^{-1}mh)=c_2(n\,m_0\gamma^{-1}h)=c_2(m_0\gamma^{-1}h)
=\chi(t)|t|^zc_2(\gamma^{-1}h)=\chi(t)|t|^zc(h).
\]
So $c$ has the law $(\chi,z)$.
\end{proof}

\begin{lemma}\label{asm-theta-agree-end-iwasawa}
Let $E$ be a CM field with a theta base $B$
(Definition~\ref{asm-theta-agree-models}(1)), let $W$ be a skew-Hermitian
lattice space of signature $(r,s)$, and let $H'=(x,y,\iota,\mathfrak b)$ be an
adapted plane in $W$ over a lattice space $W_1$ at some positions
(Definition~\ref{asm-theta-agree-planes}(1)). Then $(x,y,\iota)$ is a plane in
$W$ over $W_1$ (Definition~\ref{asm-theta-agree-end-notions}(1)), and every
$g\in U_W$ is $umlk$ with $u\in N_{H'}$, $m$ a Levi element for $(t,1)$ for
some $t\in\mathbf A_E^\times$, $l$ a Levi element for $(1,h_1)$ for some
$h_1\in U_{W_1}$, and $k\in K^\theta_W$.
\end{lemma}
\leanref{Development/WeilBound/ThetaLifts/EigenvalueEndgame/PlaneRemovalConfiguration.lean}{PositivityTheory.ThKEnd.Iwa}
\leanref{Development/WeilBound/ThetaLifts/TowerModels/LeviEx.lean}{PositivityTheory.ThKG.Tow.iwa}
\begin{proof}
The conditions on $x$, $y$ and $\iota$ in the definition of a plane, other than
$\det J_1\neq0$, are among those of an adapted plane. Let $J$ and $J_1$ be the
matrices of $W$ and $W_1$ and $A=(\iota\ x\ y)$. As in the proof of
Lemma~\ref{asm-theta-agree-levi}, $A$ is invertible and
$A^{\mathsf T}J\,c(A)=\operatorname{diag}\bigl(J_1,\left(\begin{smallmatrix}0&1\\-1&0\end{smallmatrix}\right)\bigr)$.
Taking determinants, $\det A\cdot\det J\cdot c(\det A)=\det J_1$, and
$\det J\neq0$ by Lemma~\ref{asm-theta-agree-convergence-signature}. So
$\det J_1\neq0$, and $(x,y,\iota)$ is a plane.

Let $g\in U_W$. For an infinite place $w$ of $F$ let $M_w=i\,\tau_w(J)$,
which is Hermitian, $\theta=\theta_{W,w}$, a Cartan involution of $M_w$
(Definition~\ref{asm-theta-agree-doubling}(1)), and
$\langle u,u'\rangle=u^{\mathsf T}M_w\bar u'$ on $\mathbf C^d$; and let
$g_w=\pi_w(g)$, which preserves $\langle\ ,\ \rangle$, as $\tau_w\circ c$ is
complex conjugation on $\tau_w(E)$.

\emph{Step 1: reflections.} As $\theta$ preserves the form and $\theta^2=1$,
$\langle u,\theta u\rangle=\langle\theta u,u\rangle$; as the form is
Hermitian, $\langle\theta u,u\rangle$ is real, and it is positive for
$u\neq0$. Let $\epsilon=\pm1$ and let $u,u'\in\mathbf C^d$ have
$\theta u=\epsilon u$, $\theta u'=\epsilon u'$ and
$\langle u,u\rangle=\langle u',u'\rangle$. We show that some $\kappa$ that
preserves the form and commutes with $\theta$ has $\kappa u=u'$ and $\kappa z=z$
whenever $\theta z=-\epsilon z$. Put $a=u-u'$ and $b=\langle u,a\rangle$; then
$\theta a=\epsilon a$ and $\langle a,a\rangle=b+\bar b$. If $b=0$, then
$\langle\theta a,a\rangle=\epsilon\langle a,a\rangle=0$, so $a=0$ and
$\kappa=1$ works. Otherwise put $\kappa z=z-b^{-1}\langle z,a\rangle a$. Then
\[
\langle\kappa z,\kappa z'\rangle=\langle z,z'\rangle-\langle z,a\rangle\overline{\langle z',a\rangle}\bigl(\bar b^{-1}+b^{-1}-b^{-1}\bar b^{-1}\langle a,a\rangle\bigr)=\langle z,z'\rangle ;
\]
$\langle\theta z,a\rangle=\epsilon\langle\theta z,\theta a\rangle=\epsilon\langle z,a\rangle$,
so $\kappa\theta z=\theta z-b^{-1}\langle z,a\rangle\theta a=\theta\kappa z$;
$\kappa u=u-a=u'$; and if $\theta z=-\epsilon z$, then
$\langle z,a\rangle=\langle\theta z,\theta a\rangle=-\langle z,a\rangle$, so
$\kappa z=z$.

\emph{Step 2: isotropic lines.} Let $v,v'\in\mathbf C^d$ be nonzero with
$\langle v,v\rangle=\langle v',v'\rangle=0$. We show that some $\kappa$ that
preserves the form and commutes with $\theta$ has $\kappa v=sv'$ for some real
$s>0$. Put
$s=(\langle\theta v,v\rangle/\langle\theta v',v'\rangle)^{1/2}$ and $v''=sv'$,
so that $\langle v'',v''\rangle=0$ and
$\langle\theta v'',v''\rangle=\langle\theta v,v\rangle$. Put
$v_\pm=v\pm\theta v$ and $v''_\pm=v''\pm\theta v''$. Then
$\theta v_\pm=\pm v_\pm$, $\theta v''_\pm=\pm v''_\pm$, and
$\langle v_\pm,v_\pm\rangle=\pm2\langle\theta v,v\rangle=\langle v''_\pm,v''_\pm\rangle$,
as $\langle\theta v,\theta v\rangle=\langle v,v\rangle=0$. Step 1 gives
$\kappa_1$ with $\kappa_1v_+=v''_+$ that fixes $v_-$, and $\kappa_2$ with
$\kappa_2v_-=v''_-$ that fixes $v''_+$. Then $\kappa=\kappa_2\kappa_1$ has
$\kappa v=\tfrac12\kappa_2(v''_++v_-)=\tfrac12(v''_++v''_-)=sv'$.

\emph{Step 3: the infinite places.} The vector $v=\tau_w(x)$ is nonzero, as
$h(x,y)=1$, and $\langle v,v\rangle=i\,\tau_w(h(x,x))=0$, with
$h(u,u')=u^{\mathsf T}J\,c(u')$. So Step 2, for $v$ and $v'=g_w^{-1}v$, gives
$\kappa_w$ and $s_w>0$ with $\kappa_wv=s_wg_w^{-1}v$. Put $p_w=g_w\kappa_w$.
Then $p_w$ preserves the form, $p_w\tau_w(x)=s_w\tau_w(x)$, and
$p_w^{-1}g_w=\kappa_w^{-1}$ commutes with $\theta_{W,w}$.

\emph{Step 4: an adelic element.} Let $p\in GL_d(\mathbf A_E)$ have
$\pi_w(p)=p_w$ at every infinite place $w$ and component $1$ at every finite
place, and let $t\in\mathbf A_E^\times$ have $\pi_w(t)=s_w$ and component $1$ at
every finite place. As each $p_w$ preserves $M_w$, $p\in U_W$; and $px=tx$,
and $k=p^{-1}g\in K^\theta_W$, as $\pi_w(k)=\kappa_w^{-1}$.

\emph{Step 5: the stabilizer of the line.} Put
$J^\natural=\operatorname{diag}\bigl(J_1,\left(\begin{smallmatrix}0&1\\-1&0\end{smallmatrix}\right)\bigr)$
and $P=A^{-1}pA$. Then $P^{\mathsf T}J^\natural c(P)=J^\natural$, and the
column of $P$ for $x$ is $t$ times the corresponding standard basis vector. For
an index $j$ of the block of $\iota$, the entry $(j,x)$ of
$P^{\mathsf T}J^\natural c(P)$ is $-P_{yj}\,c(t)$, which must be $0$; so
$P_{yj}=0$. Hence $p\iota=\iota h_1+x\alpha^{\mathsf T}$, where $h_1$ is the
block of $P$ on the indices of $\iota$ and $\alpha_j=P_{xj}$. The block of
$P^{\mathsf T}J^\natural c(P)$ on these indices is then
$h_1^{\mathsf T}J_1c(h_1)$, so $h_1^{\mathsf T}J_1c(h_1)=J_1$; as
$\det J_1\neq0$, $h_1$ is invertible, and $h_1\in U_{W_1}$.

\emph{Step 6: the decomposition.} By Lemma~\ref{asm-theta-agree-levi} there
are a Levi element $m$ for $(t,1)$ and a Levi element $l$ for $(1,h_1)$. Put
$u=p(ml)^{-1}$, so that $g=pk=umlk$. As $mlx=tx$, $ux=t^{-1}px=x$. For
$z\in\mathbf A_E^{d-2}$, $ml\,\iota h_1^{-1}z=m\iota z=\iota z$, so
$u\iota z=p\iota h_1^{-1}z=\iota z+(\alpha^{\mathsf T}h_1^{-1}z)\,x$. So $u$
is one of the generators of $N_{H'}$
(Definition~\ref{asm-theta-agree-planes}(2)).
\end{proof}
% Lean: the plane is Levi.isFr with IsFr.isUnit_det. The decomposition is iwa through
% Levi.exists_stab (LeviEx.lean). Step 1 is exists_refl and Step 2 exists_move
% (ExistsRefl.lean); Step 3 is exists_loc; Step 4 is exists_lift and exists_inf (LeviAd.lean),
% assembled in exists_stab; Step 5 is IsFr.stab for isFrA, with exists_UA; Step 6 is
% exists_levi twice and the closing computation of iwa. Steps 1 to 4 are shared with iwaEx_of
% (Lemma asm-theta-agree-split-inputs(3)).

\begin{lemma}\label{asm-theta-agree-end-automorphic}
Let $(B,G,\mathcal X,\omega,\omega')$ be a triple for a theta base $B$ of a CM
field $E$ (Definition~\ref{asm-theta-agree-triple}), with models $M$ and $M'$
and theta integrals $I$ and $I'$. Then $I(\Phi)$ is automorphic
(Definition~\ref{asm-theta-agree-end-notions}(3)) for every
$\Phi\in\operatorname{PD}(M)$, and $I'(\Phi')$ is automorphic for every
$\Phi'\in\operatorname{PD}(M')\cap\mathcal S^\circ(M')$.
\end{lemma}
\leanref{Development/WeilBound/ThetaLifts/EigenvalueEndgame/PlaneRemovalConfiguration.lean}{PositivityTheory.ThKEnd.IAut}
\leanref{Development/WeilBound/ThetaLifts/EigenvalueEndgame/UnipotentGroups/PerpZeroOf.lean}{PositivityTheory.ThKEnd.Grp.iAut_of}
\leanref{Development/WeilBound/ThetaLifts/EigenvalueEndgame/UnipotentGroups/PerpZeroOf.lean}{PositivityTheory.ThKEnd.Grp.I_inv}
\begin{proof}
Let $G=(p,q,V,V',W,M,M')$, of ranks $(a,a')$. Every element of
$\operatorname{PD}(M)$ or of $\operatorname{PD}(M')$ is a Schwartz-Bruhat
function: for a spanning element $f\otimes\bigotimes_v\phi_v$, the function
$\bigotimes_v\phi_v$ is a Schwartz-Bruhat function on the finite adeles, since
every $\phi_v$ is one and, the models being regular, $\phi_v=\mathrm{vac}_v$ is
the indicator function of the integral vectors for almost all $v$.

\emph{Invariance.} Let $\Psi$ be a Schwartz-Bruhat function on
$\mathbf A^{a\times(a+a')}$, and let
$\gamma\in U_W(F)$ and $h\in U_W$. By
Lemma~\ref{asm-theta-agree-anisotropic-automorphy},
$\theta_\omega(\Psi)(g,\gamma h)=\theta_\omega(\Psi)(g,h)$ for every $g\in U_V$.
So the integrands defining $I(\Psi)(\gamma h)$ and $I(\Psi)(h)$ agree, and
$I(\Psi)(\gamma h)=I(\Psi)(h)$. The same holds for $M'$, $\omega'$ and $I'$.

\emph{Continuity of $I(\Phi)$.} The Weil pair $(V,W,M,\beta,\beta/\alpha,\omega)$
has signature $(a-p,p;a+a'-p-q,p+q)$, and
$a+a'+2\min(a-p,p)\leq a+a'+2p<2a$, because $2(p+q)+1<\lfloor(a-a')/2\rfloor$.
So $I(\Phi)$ is continuous by Lemma~\ref{asm-theta-agree-definite-continuous}.

\emph{Continuity of $I'(\Phi')$.} Let $\mu'$ and $\mathcal D'$ be the measure
and the domain of $M'$. By Lemma~\ref{asm-theta-agree-continuity-pieces}(1) for
the Weil pair $(V',W,M',\alpha,\alpha/\beta,\omega')$,
$\Phi'=\sum_if_i\otimes\phi_i$ with full-rank finite parts $\phi_i$.
Lemma~\ref{asm-theta-agree-continuity-domination} gives $k$, and for each $i$ a
continuous seminorm $p_i$ and a function $G_i\geq0$ with
$\int_{\mathcal D'}G_i\,d\mu'<\infty$ and
$|\theta_{\omega'}(f\otimes\phi_i)(g,h)|\leq p_i(f)\operatorname{ht}(h)^kG_i(g)$.
By Lemma~\ref{asm-theta-agree-continuity-integrals}(1),
$I'(\Phi')=\sum_iI'(f_i\otimes\phi_i)$. Each term is continuous, by the
argument of the proof of Lemma~\ref{asm-theta-agree-definite-continuous}: a
point $h_0\in U_W$ has a compact neighbourhood on which the height is bounded
by some $C$; there the integrand $g\mapsto\theta_{\omega'}(f_i\otimes\phi_i)(g,h)$
is bounded by $p_i(f_i)C^kG_i(g)$, which is integrable on $\mathcal D'$, and it
depends continuously on $h$ by
Lemma~\ref{asm-theta-agree-convergence-theta}(2); dominated convergence gives
continuity at $h_0$.
\end{proof}
% Lean: iAut_of thetaBasic iCont soStab: the first clause of ICont with weilCondV, and its
% second clause (iCont_of with soFin, dom_of_domSoU, dom_sum and dom_continuous), whose
% hypothesis on So is SoStab; I_inv is the invariance; polData_le gives PD in schA.

\begin{lemma}\label{asm-theta-agree-end-stable-level}
Let $(X,W,M,\chi_X,\chi_W,\omega)$ be a Weil pair for a theta base $B$ of a CM
field $E$ (Definition~\ref{asm-theta-agree-triple}(1)), with $X$ and $W$ of
ranks $n$ and $d$, and put $I=\{1,\ldots,n\}\times\{1,\ldots,d\}$. Then for
every Schwartz-Bruhat function $\phi$ on $\mathbf A_f^I$ there is a
neighbourhood $U$ of $1$ in $U_W$ such that $\omega(1,u)(f\otimes\phi)=f\otimes\phi$
for every $f\in\mathcal S(F_\infty^I)$ and every $u\in U$ with $\pi_w(u)=1$ for
every infinite place $w$ of $F$.
\end{lemma}
\leanref{Development/WeilBound/ThetaLifts/FrameEstimates/TowerRallisLemma.lean}{PositivityTheory.ThKEst.Level}
\leanref{Development/WeilBound/ThetaLifts/FrameEstimates/SchrAddPt.lean}{PositivityTheory.ThKEst.level_of}
\begin{proof}
Let $e=(e_1,e_2)$ be the isomorphism of $M$ and $\rho$ its Schrödinger
operators, and call $u\in U_W$ \emph{finite} if $\pi_w(u)=1$ for every $w$,
that is, if $u-1$ vanishes at the infinite places. If $f\otimes\phi=0$ for
every $f$, there is nothing to prove. Otherwise choose $f_0$ and
$y\in\mathbf A^I$ with $\Psi_0(y)\neq0$, where $\Psi_0=f_0\otimes\phi$, and fix
a Haar measure on $\mathbf A^I$ for the coefficients of
Definition~\ref{asm-theta-agree-convergence-tori}(6).

(i) Lemma~\ref{asm-theta-agree-convergence-coefficients}(2), for
$\phi_1=\phi_2=\phi$ and $N=0$, gives a compact set
$C\subset M_{n\times d}(\mathbf A_E)$ such that
$\langle\rho(z)(f\otimes\phi),f'\otimes\phi\rangle=0$ unless $z-c$ is at
$\infty$ for some $c\in C$.

(ii) There is a neighbourhood $U_1$ of $1$ such that
$\rho(zu^{\mathsf T})(f\otimes\phi)=\rho(z)(f\otimes\phi)$ for every finite
$u\in U_1$, every $f$ and every $z$ with $z-c$ at $\infty$ for some $c\in C$.
Indeed, for such $u$ and $z$, $(z-c)(u-1)^{\mathsf T}=0$, so
$zu^{\mathsf T}=z+\delta'$ with $\delta'=c(u-1)^{\mathsf T}$; put $w=e(z)$ and
$\delta=e(\delta')$, which vanishes at $\infty$. By the definition of $\rho$,
\[
\rho(w+\delta)=\psi\bigl(\tfrac12(\delta_1\cdot w_2-w_1\cdot\delta_2)\bigr)\rho(w)\rho(\delta),
\]
and $\delta_1\cdot w_2-w_1\cdot\delta_2=\delta_1\cdot e_2(c)-e_1(c)\cdot\delta_2$,
as $\delta$ vanishes at $\infty$ and $w-e(c)$ is at $\infty$. On the finite
adeles $\psi$ is locally constant, and $\phi$ is locally constant with compact
support. So there is a neighbourhood $D$ of $0$ such that, for $\delta\in D$
vanishing at $\infty$, $\rho(\delta)(f\otimes\phi)=f\otimes\phi$ for every $f$,
and, $C$ being compact, the phase above is $1$ for every $c\in C$. As $(u,c)\mapsto e(c(u-1)^{\mathsf T})$
is continuous and vanishes for $u=1$, and $C$ is compact, there is a
neighbourhood $U_1$ of $1$ with $e(c(u-1)^{\mathsf T})\in D$ for all $u\in U_1$
and $c\in C$.

(iii) The set $U_0$ of the $u$ with
$|(\omega(1,u)\Psi_0)(y)/\Psi_0(y)-1|<\frac12$ is an open neighbourhood of $1$,
as evaluation at $y$ is admissible and condition (c) of
Definition~\ref{asm-theta-agree-models}(6) makes $u\mapsto(\omega(1,u)\Psi_0)(y)$
continuous. There is a neighbourhood $U_2$ of $1$ such that $u^p\in U_0$ for
every finite $u\in U_2$ and every $p\geq0$. Indeed, as $U_W$ carries the
topology of the pairs $(g,g^{-1})$, there is a neighbourhood $V$ of $0$ in
$\mathbf A_E$ such that every $g\in U_W$ for which the entries of $g-1$ and
$g^{-1}-1$ lie in $V$ belongs to $U_0$. Fix a basis of $E$ over $F$. For an
integer $N\geq1$, let $\Lambda_N$ be the set of the $z\in\mathbf A_E$ that
vanish at the infinite places and whose coordinates in this basis are $N$ times
integral finite adeles. It is an additive subgroup, and for a suitable $N$ it is
closed under products and contained in $V$; put $\Lambda=\Lambda_N$. There is a
neighbourhood $V'$ of $0$ in $\mathbf A_E$ such that every $z\in V'$ that
vanishes at the infinite places lies in $\Lambda$, as the coordinates of $z/N$
are integral for $z$ near $0$. Let $U_2$ be the set of the $u$ for which the
entries of $u-1$ and $u^{-1}-1$ lie in $V'$. Let $u\in U_2$ be finite. Then
$u^{-1}$ is finite, so the entries of $u-1$ and $u^{-1}-1$ lie in $\Lambda$. By
induction on $p$, using
$u^{p+1}-1=(u^p-1)(u-1)+(u^p-1)+(u-1)$, the entries of $u^p-1$ lie in
$\Lambda$, and in the same way those of $(u^p)^{-1}-1=(u^{-1})^p-1$. As
$\Lambda\subset V$, $u^p\in U_0$.

(iv) Let $u$ be finite with $u\in U_1$, $u^{-1}\in U_1$ and $u\in U_2$, let
$f\in\mathcal S(F_\infty^I)$, and put $\Phi_1=\omega(1,u)(f\otimes\phi)$ and
$\Phi_2=\omega(1,u)\Psi_0$. For $z\in M_{n\times d}(\mathbf A_E)$,
Lemma~\ref{asm-theta-agree-convergence-coefficients}(1) gives
$\langle\rho(z)\Phi_1,\Phi_2\rangle=\langle\rho(z(u^{-1})^{\mathsf T})(f\otimes\phi),\Psi_0\rangle$,
and this is $\langle\rho(z)(f\otimes\phi),\Psi_0\rangle$: if $z-c$ is at $\infty$
for some $c\in C$, by (ii) for $u^{-1}$; if $z(u^{-1})^{\mathsf T}-c$ is at
$\infty$ for some $c\in C$, by (ii) for $u$ and $z(u^{-1})^{\mathsf T}$; and
otherwise both coefficients vanish by (i). As $e$ is onto,
Lemma~\ref{asm-theta-agree-continuity-fourier} gives
$\Phi_1(x)\overline{\Phi_2(y)}=(f\otimes\phi)(x)\overline{\Psi_0(y)}$ for all
$x$. For $f=f_0$ and $x=y$ this reads $|\Phi_2(y)|^2=|\Psi_0(y)|^2$, which is
not $0$, so $\Phi_2(y)\neq0$. So
$\omega(1,u)(f\otimes\phi)=c(u)\,f\otimes\phi$ for every $f$, with
$c(u)=\overline{\Psi_0(y)}/\overline{\Phi_2(y)}$.

(v) For such $u$, as $\omega$ is a representation,
$\omega(1,u^p)\Psi_0=c(u)^p\Psi_0$ for every $p\geq0$, and $u^p\in U_0$ by
(iii). So $|c(u)^p-1|<\frac12$ for every $p$, and $c(u)=1$ by
Lemma~\ref{asm-theta-agree-finite-inputs-powers}. Hence
$U=U_1\cap U_1^{-1}\cap U_2$ has the property of the lemma.
\end{proof}
% Lean: level_of coefEq_of fourierInj_of (coefBd_of coefReal_of) (SchrAddPt.lean): (i) is
% CoefBd, (ii) is coef_near through schr_near, fix_near (eventually_fix, from lc_fin) and
% schr_add_pt, (iii) is pow_near (the neighbourhood V of 0, Lambda = exists_lattice with
% N = N_1 N_2, the induction key, then key for u^{-1}) (LevelA.lean) and ev_adm, (iv) is the
% key identity through CoefEq and FourierInj, with the case split a, a' of the Lean and h4,
% and (v) is h5 and eq_one_of_pow_near. Checked against level_of line by line. The theorem level (CuspRapid.lean) is this
% with the three inputs.

\begin{lemma}\label{asm-theta-agree-end-stable}
Let $(B,G,\mathcal X,\omega,\omega',R,R_f)$ be a matched quadruple for a theta
base $B$ of a CM field $E$ (Definition~\ref{asm-theta-agree-matched}), and let
$k\in K^\theta_W$ (Definition~\ref{asm-theta-agree-end-notions}(2)). If
$\Phi'\in\mathcal S^\circ(M')$ and $\Phi$ and $\Phi'$ correspond, then
$\omega'(1,k)\Phi'\in\mathcal S^\circ(M')$, and $\omega(1,k)\Phi$ and
$\omega'(1,k)\Phi'$ correspond.
\end{lemma}
\leanref{Development/WeilBound/ThetaLifts/EigenvalueEndgame/PlaneRemovalConfiguration.lean}{PositivityTheory.ThKEnd.MatchStab}
\leanref{Development/WeilBound/ThetaLifts/EigenvalueEndgame/UnipotentGroups/TendstoIndMul.lean}{PositivityTheory.ThKEnd.Grp.matchStab_of}
\begin{proof}
By Lemma~\ref{asm-theta-agree-split-inputs}(6) for the Weil pair
$(V',W,M',\alpha,\alpha/\beta,\omega')$, the operator $\omega'(1,k)$ maps
$\mathcal S^\circ(M')$ into itself.

Let $V$, $V'$ and $W$ be the spaces of $G$ and $M$, $M'$ its models, and let
$\mathcal X=(\alpha,\beta,S_0,Q)$, $R=(\Pi,p,p')$ and
$R_f=(\Pi_v,p_v,p'_v)_v$. As $\Phi$ and $\Phi'$ correspond, $\Phi\in\operatorname{PD}(M)$ and
$\Phi'\in\operatorname{PD}(M')$, so both are Schwartz-Bruhat functions (proof
of Lemma~\ref{asm-theta-agree-end-automorphic}).

\emph{Step 1: a common operator.} Let $Z$, $Z'$ and $N$ be complex vector
spaces, $\pi\colon Z\to N$ and $\pi'\colon Z'\to N$ linear maps with $\pi$
onto, and $A$, $A'$ operators on $Z$ and $Z'$ such that $\pi(x)=\pi'(x')$
implies $\pi(Ax)=\pi'(A'x')$. Then some operator $O$ on $N$ has $O\pi=\pi A$
and $O\pi'=\pi'A'$. Indeed, let $s$ be a linear section of
$\pi$ and $O=\pi As$. If $\pi(x)=0$, the hypothesis for $x$ and $x'=0$ gives
$\pi(Ax)=0$; so $\pi As\pi(x)=\pi(Ax)$, as $s\pi(x)-x$ lies in the kernel of
$\pi$. For $x'\in Z'$, the hypothesis for $s\pi'(x')$ and $x'$ gives
$O\pi'(x')=\pi'(A'x')$.

\emph{Step 2: a finite place.} Let $v$ be a finite place, $h\in U_W^{(v)}$,
and let $\Psi$ and $\Psi'$ correspond. We show that $\omega(1,h)\Psi$ and
$\omega'(1,h)\Psi'$ correspond. Let $(\omega_v,\rho_v)$ and
$(\omega'_v,\rho'_v)$ be local operators at $v$ for the two Weil pairs of the
triple (Lemma~\ref{asm-theta-agree-finite-inputs}(2) and
Definition~\ref{asm-theta-agree-finite-local-objects}(4)), so that
$\omega(1,h)$ acts at $v$ by $\phi\mapsto\omega_v(1,h)\phi$, and likewise for
$\omega'$. As $R_f$ is matched, $p_v$ maps the Schwartz-Bruhat functions onto
$\Pi_v$, and $p_v(\phi)=p'_v(\phi')$ implies
$p_v(\omega_v(1,h)\phi)=p'_v(\omega'_v(1,h)\phi')$. Step 1 gives an operator
$O_v$ on $\Pi_v$ with $O_vp_v(\phi)=p_v(\omega_v(1,h)\phi)$ and
$O_vp'_v(\phi')=p'_v(\omega'_v(1,h)\phi')$ for Schwartz-Bruhat $\phi$ and
$\phi'$. Write $\Psi=\sum_if_i\otimes\bigotimes_u\phi_{i,u}$ and
$\Psi'=\sum_jf'_j\otimes\bigotimes_u\phi'_{j,u}$ with a set $T_1\supset T$ as
in Definition~\ref{asm-theta-agree-matched}(4). As $\omega(1,h)$ acts at $v$
by $\phi_{i,v}\mapsto\omega_v(1,h)\phi_{i,v}$, it maps
$f_i\otimes\bigotimes_u\phi_{i,u}$ to the same tensor with
$\omega_v(1,h)\phi_{i,v}$ in place of $\phi_{i,v}$; likewise for $\omega'$.
Put $T_2=T_1\cup\{v\}$. If $v\in T_1$, apply $O_v$ to the factor at $v$ of the
identity of the correspondence. If $v\notin T_1$, then $\phi_{i,v}=\mathrm{vac}_v$
and $\phi'_{j,v}=\mathrm{vac}'_v$ for all $i$ and $j$, and
$p_v(\mathrm{vac}_v)=p'_v(\mathrm{vac}'_v)$ as $v\notin T$; tensor the identity
with $O_vp_v(\mathrm{vac}_v)=O_vp'_v(\mathrm{vac}'_v)$ at $v$. In both cases
the result is the identity, over $T_2$, that makes $\omega(1,h)\Psi$ and
$\omega'(1,h)\Psi'$ correspond.

\emph{Step 3: the infinite places.} Let $k_0\in K_\infty(W)$, and let $\Psi$
and $\Psi'$ correspond. We show that $\omega(1,k_0)\Psi$ and
$\omega'(1,k_0)\Psi'$ correspond. By
Lemma~\ref{asm-theta-agree-rechoice-operators}(1) there are archimedean
operators $A$ of $\omega$ and $A'$ of $\omega'$ at $(1,k_0)$, and by
Lemma~\ref{asm-theta-agree-rechoice-fock} they map $\operatorname{Pol}(M)$ and
$\operatorname{Pol}(M')$ into themselves. As $R$ is matched,
$p(\operatorname{Pol}(M))=\Pi$, and $p(f)=p'(f')$ implies $p(Af)=p'(A'f')$ for
$f\in\operatorname{Pol}(M)$ and $f'\in\operatorname{Pol}(M')$ (condition (b)
of Definition~\ref{asm-theta-agree-matched}(2)). Step 1, for the restrictions
of $p$ and $p'$ to the polynomial Fock spaces, gives an operator $O$ on $\Pi$
with $Op(f)=p(Af)$ and $Op'(f')=p'(A'f')$. With $\Psi$ and $\Psi'$ written as
in Step 2, $\omega(1,k_0)$ maps $f_i\otimes\bigotimes_u\phi_{i,u}$ to
$A(f_i)\otimes\bigotimes_u\phi_{i,u}$, and likewise for $\omega'$. Applying
$O\otimes1$ to the identity of the correspondence gives the identity that makes
$\omega(1,k_0)\Psi$ and $\omega'(1,k_0)\Psi'$ correspond.

\emph{Step 4: pieces of $k$.} For $e\in\mathbf A_F$ with $e^2=e$ and
$g\in U_W$, put $g_e=1+e(g-1)$. As $e^2=e$ and $c(e)=e$,
$g_e^{\mathsf T}J_Wc(g_e)=J_W+e\bigl(g^{\mathsf T}J_Wc(g)-J_W\bigr)=J_W$ and
$g_e(g^{-1})_e=(gg^{-1})_e=1$, so $g_e\in U_W$; and if $e=e_1+e_2$ with
$e_1e_2=0$ and $e_i^2=e_i$, then $g_e=g_{e_1}g_{e_2}$. For a finite set $S$ of
finite places let $\varepsilon_S\in\mathbf A_F$ be $1$ at the infinite places
and at the places of $S$ and $0$ at the other places, and for a finite place
$v$ let $\iota_v$ be $1$ at $v$ and $0$ elsewhere. Then
$\varepsilon_S=\varepsilon_\emptyset+\sum_{v\in S}\iota_v$, a sum of pairwise
orthogonal idempotents, so
$k_{\varepsilon_S}=k_{\varepsilon_\emptyset}\prod_{v\in S}k_{\iota_v}$ in some
order. Here $k_0=k_{\varepsilon_\emptyset}$ is at $\infty$ with
$\pi_w(k_0)=\pi_w(k)$, so $k_0\in K_\infty(W)$ as $k\in K^\theta_W$, and each
$k_{\iota_v}$ lies in $U_W^{(v)}$. As $S$ grows, $\varepsilon_Sx\to x$ for every
$x\in\mathbf A_E$, so $k_{\varepsilon_S}\to k$ and
$r_S=k_{\varepsilon_S}^{-1}k\to1$ in $U_W$; and $\pi_w(r_S)=1$ for every $w$,
as $k_{\varepsilon_S}$ has the archimedean components of $k$.

\emph{Step 5: a neighbourhood fixing $\Phi$ and $\Phi'$.} Write
$\Phi=\sum_if_i\otimes\phi_i$ and $\Phi'=\sum_jf'_j\otimes\phi'_j$. Applying
Lemma~\ref{asm-theta-agree-end-stable-level} to each $\phi_i$, for the Weil
pair $(V,W,M,\beta,\beta/\alpha,\omega)$, and to each $\phi'_j$, for
$(V',W,M',\alpha,\alpha/\beta,\omega')$, and intersecting the finitely many
neighbourhoods so obtained gives a neighbourhood
$U$ of $1$ in $U_W$ with $\omega(1,u)\Phi=\Phi$ and $\omega'(1,u)\Phi'=\Phi'$
for every $u\in U$ with $\pi_w(u)=1$ for all $w$.

\emph{Step 6: conclusion.} By Step 4 choose $S$ with $r=r_S\in U$. Then
$k=k_0\,l_1\cdots l_m\,r$, where $l_1,\ldots,l_m$ are the $k_{\iota_v}$,
$v\in S$, and by Step 5 $\omega(1,r)\Phi=\Phi$ and $\omega'(1,r)\Phi'=\Phi'$.
So
$\omega(1,k)\Phi=\omega(1,k_0)\omega(1,l_1)\cdots\omega(1,l_m)\Phi$, and
likewise for $\omega'$ and $\Phi'$. By Step 2, applied $m$ times starting with
$l_m$, and then Step 3, $\omega(1,k)\Phi$ and $\omega'(1,k)\Phi'$ correspond.
\end{proof}
% Lean: matchStab_of ArchOp PureAct Level SoStab (TendstoIndMul.lean). Step 1 is common_op;
% Step 2 is match_fin, with Locrep.locRepEx, PureAct and exists_ins; Step 3 is match_inf,
% with ArchOp (its clause for K_infinity is Lemma asm-theta-agree-rechoice-fock); Step 4 is
% piece, mix_mul, piece_add, piece_prod, piece_inf, piInf_piece, tendsto_eps and
% tendsto_piece; Steps 5 and 6 are level_sch and decomp_of with the closing induction.

\begin{lemma}\label{asm-theta-agree-end-rechoose}
Let $(B,G,\mathcal X,\omega,\omega')$ be a triple of ranks $(a,a')$ for a theta
base $B$ of a CM field $E$ (Definition~\ref{asm-theta-agree-triple}), and let
$V^\circ$, $V'^\circ$ and $W^\circ$ be lattice spaces such that, with
$G^\circ=G[V^\circ,V'^\circ,W^\circ]$, $(B,G^\circ,\mathcal X,\omega,\omega')$ is
a triple. Then there is $h_0\in U_W$ such that for every $f$ in the set of
differences of $(B,G,\mathcal X,\omega,\omega')$
(Definition~\ref{asm-theta-agree-end-notions}(5)), the function
$h\mapsto f(hh_0)$ lies in the set of differences of
$(B,G^\circ,\mathcal X,\omega,\omega')$.
\end{lemma}
\leanref{Development/WeilBound/ThetaLifts/EigenvalueEndgame/PlaneRemovalConfiguration.lean}{PositivityTheory.ThKEnd.SigRe}
\leanref{Development/WeilBound/ThetaLifts/EigenvalueEndgame/UnipotentGroups/SigReOf.lean}{PositivityTheory.ThKEnd.Grp.sigRe_of}
\leanref{Development/WeilBound/ThetaLifts/EigenvalueEndgame/UnipotentGroups/SigReOf.lean}{PositivityTheory.ThKEnd.Grp.move}
\leanref{Development/WeilBound/ThetaLifts/EigenvalueEndgame/UnipotentGroups/SigReOf.lean}{PositivityTheory.ThKEnd.Grp.I_both}
\leanref{Development/WeilBound/ThetaLifts/EigenvalueEndgame/UnipotentGroups/SigReOf.lean}{PositivityTheory.ThKEnd.Grp.so_re}
\begin{proof}
Let $G=(p,q,V,V',W,M,M')$, let $\tilde V$, $\tilde V'$ and $\tilde W$ be the
spaces and $M^\circ$ and $M'^\circ$ the models of $G^\circ$, and let $\Sigma$
and $\Sigma^\circ$ be the two sets of differences. The two triples have the
same theta integrals $I$ and $I'$ (Definition~\ref{asm-theta-agree-triple}(2)),
and $U_{\tilde W}=U_W$, since a unitary group depends only on the matrix.

\emph{Moving the spaces.} The spaces $W$ and $\tilde W$ have the same matrix
and are skew-Hermitian of the same signature $(a+a'-p-q,p+q)$, as both triples
are proper with the same $p$ and $q$. By
Lemma~\ref{asm-theta-agree-rechoice-cartan} there is $h\in U_W$ that moves
$\tilde W$ to $W$ (Definition~\ref{asm-theta-agree-rechoice-terms}(4)). The
Weil pairs $(V,W,M,\beta,\beta/\alpha,\omega)$ and
$(\tilde V,\tilde W,M^\circ,\beta,\beta/\alpha,\omega)$ have the same signature,
$\tilde V$ and $\tilde W$ have the matrices and lattices of $V$ and $W$, and
$M^\circ$ has the isomorphism of $M$. By
Lemma~\ref{asm-theta-agree-rechoice-move} there is $g\in U_V$ that moves
$\tilde V$ to $V$ such that $T(\operatorname{Pol}(M^\circ))=\operatorname{Pol}(M)$
for every archimedean operator $T$ of $\omega$ at $(g,h)$. In the same way
there is $g'\in U_{V'}$ that moves $\tilde V'$ to $V'$ such that
$T'(\operatorname{Pol}(M'^\circ))=\operatorname{Pol}(M')$ for every archimedean
operator $T'$ of $\omega'$ at $(g',h)$.

\emph{Polynomial data.} The operator $\omega(g^{-1},h^{-1})$ maps
$\operatorname{PD}(M)$ into $\operatorname{PD}(M^\circ)$. By
Lemma~\ref{asm-theta-agree-rechoice-operators}(1) there is an archimedean
operator $T$ of $\omega$ at $(g,h)$, as $g$ and $h$ are at $\infty$. It
suffices to treat a spanning element $\Phi=f\otimes\bigotimes_v\phi_v$ with
$f\in\operatorname{Pol}(M)$. The function $\phi=\bigotimes_v\phi_v$ is a
Schwartz-Bruhat function on the finite adeles, as $M$ is regular, and
$f=T(f_2)$ for some $f_2\in\operatorname{Pol}(M^\circ)$. As $\omega$ is a
representation and $T$ an archimedean operator,
\[
\omega(g^{-1},h^{-1})\Phi=\omega(g^{-1},h^{-1})\,\omega(g,h)(f_2\otimes\phi)=f_2\otimes\phi,
\]
which lies in $\operatorname{PD}(M^\circ)$, since $M^\circ$ has the vacua of $M$.
In the same way $\omega'(g'^{-1},h^{-1})$ maps $\operatorname{PD}(M')$ into
$\operatorname{PD}(M'^\circ)$. By Lemma~\ref{asm-theta-agree-split-inputs}(6),
$\omega'(g'^{-1},h^{-1})$ maps $\mathcal S^\circ(M')$ into itself, and
$\mathcal S^\circ(M')\subset\mathcal S^\circ(M'^\circ)$: full rank at $\tau_1$
is defined through linear models of the second kind for $(\tau_1,1,1)$, which
involve the spaces only through their matrices, and $M'^\circ$ has the
isomorphism and the vacua of $M'$.

\emph{Translation.} Put $h_0=h^{-1}$. Let $\Psi$ be a Schwartz-Bruhat function
for $M$ and $x\in U_W$. As
$\omega(g^{-1},h^{-1})=\omega(g^{-1},1)\,\omega(1,h^{-1})$,
Lemma~\ref{asm-theta-agree-anisotropic-invariance} gives
$I(\omega(g^{-1},h^{-1})\Psi)=I(\omega(1,h^{-1})\Psi)$, and
$I(\omega(1,h^{-1})\Psi)(x)=I(\Psi)(xh^{-1})$, since
$\theta_\omega(\omega(1,h^{-1})\Psi)(y,x)=\theta_\omega(\Psi)(y,xh^{-1})$. The
same holds for $M'$, $\omega'$ and $g'$. Let $f=I(\Phi)-I'(\Phi')\in\Sigma$.
Elements of $\operatorname{PD}(M)$ and $\operatorname{PD}(M')$ are
Schwartz-Bruhat, as in the proof of
Lemma~\ref{asm-theta-agree-end-automorphic}. So
\[
f(xh_0)=I(\omega(g^{-1},h^{-1})\Phi)(x)-I'(\omega'(g'^{-1},h^{-1})\Phi')(x),
\]
with $\omega(g^{-1},h^{-1})\Phi\in\operatorname{PD}(M^\circ)$ and
$\omega'(g'^{-1},h^{-1})\Phi'\in\operatorname{PD}(M'^\circ)\cap\mathcal S^\circ(M'^\circ)$.
Hence $x\mapsto f(xh_0)$ lies in $\Sigma^\circ$.
\end{proof}

\begin{lemma}\label{asm-theta-agree-end-tower}
Let $(B,G,\mathcal X,\omega,\omega')$ be a triple of ranks $(b+2,a')$ for a theta
base $B$ of a CM field $E$ (Definition~\ref{asm-theta-agree-triple}), with
$p\geq1$. Then there are lattice spaces $V^\circ$, $V'^\circ$ and $W^\circ$ such
that, with $G^\circ=G[V^\circ,V'^\circ,W^\circ]$,
$(B,G^\circ,\mathcal X,\omega,\omega')$ is a triple, and there are a triple datum
$G_1$ of ranks $(b,a')$, adapted planes $H$ and $H'$ such that
$(G^\circ,G_1,H,H')$ is a plane-removal configuration
(Definition~\ref{asm-theta-agree-end-notions}(7)), and representations
$\omega_1$ and $\omega'_1$ such that $(B,G_1,\mathcal X,\omega_1,\omega'_1)$ is a
triple.
\end{lemma}
\leanref{Development/WeilBound/ThetaLifts/EigenvalueEndgame/PlaneRemovalConfiguration.lean}{PositivityTheory.ThKEnd.TowerCodimensionTwo}
\leanref{Development/WeilBound/ThetaLifts/TowerModels/TowerCodimensionTwoOf.lean}{PositivityTheory.ThKG.Tow.towerCodimensionTwo_of}
\begin{proof}
Let $G=(p,q,V,V',W,M,M')$ and $\mathcal X=(\alpha,\beta,S_0,Q)$. As $G$ is
proper with $a=b+2$, $2(p+q)+1<\lfloor(b+2-a')/2\rfloor$ and $2(p+q)+1<a'$.
With $p\geq1$ this gives $b+2-a'\geq8$, so $b\geq1$, and $b+2-p\geq1$.

\emph{The planes.} The space $V$ is Hermitian of signature $(b+2-p,p)$ and $W$
is skew-Hermitian of signature $(b+2+a'-p-q,p+q)$, both of the form $(r,s)$
with $r,s\geq1$. Proposition~\ref{asm-theta-agree-tower-plane} for $V$, with
$n_1=b$ and positions $0<b+1$, gives $V^\circ$ and a Hermitian lattice space
$V_1$ of signature $(b+1-p,p-1)$ such that $V[V^\circ]$ is Hermitian of
signature $(b+2-p,p)$ and there is an adapted plane $H$ in $V[V^\circ]$ over
$V_1$ at positions $0<b+1$. For $W$, with $n_1=b+a'$ and positions $0<b+1$, it
gives $W^\circ$ and a skew-Hermitian lattice space $W_1$ of signature
$(b+1+a'-p-q,p+q-1)$ such that $W[W^\circ]$ is skew-Hermitian of signature
$(b+2+a'-p-q,p+q)$ and there is an adapted plane $H'$ in $W[W^\circ]$ over
$W_1$ at positions $0<b+1$. Put $V'^\circ=V'$ and
$G^\circ=G[V^\circ,V'^\circ,W^\circ]$, whose spaces are $V[V^\circ]$, $V'$ and
$W[W^\circ]$.

\emph{The re-chosen triple.} The datum $G^\circ$ has the integers, the ranks
and the models of $G$, and its spaces have the matrices, the lattices and the
signatures of those of $G$. Regularity of a model and the conditions on a Weil
representation involve the spaces only through their matrices and lattices. So
$G^\circ$ is proper for $B$, and $(B,G^\circ,\mathcal X,\omega,\omega')$ is a
triple.

\emph{The smaller triple.} By Lemma~\ref{asm-theta-agree-tower-model} there are
regular models $M_1$ for $(V_1,W_1)$ and $M'_1$ for $(V',W_1)$. As $\eta^2=1$,
$\beta|_{\mathbf A^\times}=\eta^{b+2}=\eta^b$,
$(\beta/\alpha)|_{\mathbf A^\times}=\eta^{b+a'}$,
$\alpha|_{\mathbf A^\times}=\eta^{a'}$ and
$(\alpha/\beta)|_{\mathbf A^\times}=\eta^{a'+b}$, and $V_1$, $V'$ and $W_1$ have
ranks $b$, $a'$ and $b+a'$. By Theorem~\ref{asm-theta-agree-tower-weil} there
are a Weil representation $\omega_1$ for $M_1$ with characters
$(\beta,\beta/\alpha)$ and a Weil representation $\omega'_1$ for $M'_1$ with
characters $(\alpha,\alpha/\beta)$. Let $G_1=(p-1,q,V_1,V',W_1,M_1,M'_1)$, a
triple datum of ranks $(b,a')$. It is proper for $B$: $a'<b$, as
$b-a'\geq6$; $b+a'$ is even, as $b+2+a'$ is; $b+a'<b+2+a'<m_\alpha(w)-m_\beta(w)$
for every $w$; $2(p-1+q)+1<\lfloor(b-a')/2\rfloor$, as
$\lfloor(b-a')/2\rfloor=\lfloor(b+2-a')/2\rfloor-1$; $2(p-1+q)+1<a'$; the spaces
$V_1$, $V'$ and $W_1$ have signatures $(b-(p-1),p-1)$, $(q,a'-q)$ and
$(b+a'-(p-1+q),p-1+q)$; and $M_1$ and $M'_1$ are regular. The character datum
$\mathcal X$ is proper for $(B,b,a')$, since $\beta|_{\mathbf A^\times}=\eta^b$
and its other conditions do not involve the ranks. So
$(B,G_1,\mathcal X,\omega_1,\omega'_1)$ is a triple.

\emph{The configuration.} The data $G^\circ$ and $G_1$ have ranks $(b+2,a')$
and $(b,a')$ and the same second space $V'$, with $q_1=q$ and $(p-1)+1=p$; $H$
is in the first space $V[V^\circ]$ of $G^\circ$ over the first space $V_1$ of
$G_1$ at positions $(0,b+1)$, and $H'$ is in $W[W^\circ]$ over $W_1$ at
positions $(0,b+1)$. So $(G^\circ,G_1,H,H')$ is a plane-removal configuration.
\end{proof}
% Lean: towerCodimensionTwo_of planeEx modelEx WeilEx (TowerCodimensionTwoOf.lean), the
% proof of towerA_of without q = 0; G_1 is (p - 1, q, V_1, reS V' V', W_1, M_1, M'_1).

\begin{lemma}\label{asm-theta-agree-end-alpha-zero-stable}
Let $(B,G,\mathcal X,\omega,\omega')$ be a triple for a theta base $B$ of a CM
field $E$ (Definition~\ref{asm-theta-agree-triple}), with models $M$ and $M'$,
and let $k\in K^\theta_W$ (Definition~\ref{asm-theta-agree-end-notions}(2)).
Then $\omega(1,k)\Phi\in\operatorname{PD}(M)$ for every
$\Phi\in\operatorname{PD}(M)$, and
$\omega'(1,k)\Phi'\in\operatorname{PD}(M')\cap\mathcal S^\circ(M')$ for every
$\Phi'\in\operatorname{PD}(M')\cap\mathcal S^\circ(M')$.
\end{lemma}
\leanref{Development/WeilBound/ThetaLifts/EigenvalueEndgame/EigenKilling/ConeIntegralFactorisation.lean}{PositivityTheory.ThKEnd.Comp.PolStab}
\leanref{Development/WeilBound/ThetaLifts/EigenvalueEndgame/EigenKilling/AlphaZeroOf.lean}{PositivityTheory.ThKEnd.Comp.polStab_of}
\begin{proof}
By Lemmas~\ref{asm-theta-agree-real} and~\ref{asm-theta-agree-finite} there
are real maps $R$ and finite maps $R_f$ matched to the triple, so
$(B,G,\mathcal X,\omega,\omega',R,R_f)$ is a matched quadruple
(Definition~\ref{asm-theta-agree-matched}(4)).

Let $\Phi\in\operatorname{PD}(M)$. By
Lemma~\ref{asm-theta-agree-partners}(1) some $\Phi'\in\mathcal S^\circ(M')$
corresponds to $\Phi$. By Lemma~\ref{asm-theta-agree-end-stable},
$\omega(1,k)\Phi$ and $\omega'(1,k)\Phi'$ correspond, so
$\omega(1,k)\Phi\in\operatorname{PD}(M)$
(Definition~\ref{asm-theta-agree-matched}(4)).

Let $\Phi'\in\operatorname{PD}(M')\cap\mathcal S^\circ(M')$. By
Lemma~\ref{asm-theta-agree-partners}(2) some $\Phi$ corresponds to $\Phi'$. By
Lemma~\ref{asm-theta-agree-end-stable},
$\omega'(1,k)\Phi'\in\mathcal S^\circ(M')$, and $\omega(1,k)\Phi$ and
$\omega'(1,k)\Phi'$ correspond, so $\omega'(1,k)\Phi'\in\operatorname{PD}(M')$.
\end{proof}
% Lean: polStab_of KReal KFin MatchEx MatchStab (AlphaZeroOf.lean); the memberships in PD
% are polMatch_mem.

\begin{lemma}\label{asm-theta-agree-end-lower-kill}
Let $Z$ be a topological space, call a function on $Z$ Schwartz-Bruhat if it is
locally constant with compact support, let $N$ be a subspace of the functions
on $Z$, and let $Z_0\subset Z$. Suppose that for every $Y\in Z\setminus Z_0$
there are a locally constant function $m$ on $Z$ and $c\in\mathbf C$ with
$m(Y)\neq c$ and $(m-c)\Phi\in N$ for every Schwartz-Bruhat $\Phi$. Then every
Schwartz-Bruhat function that vanishes on $Z_0$ lies in $N$.
\end{lemma}
\leanref{Development/WeilBound/ThetaLifts/EigenvalueEndgame/EigenKilling/EigenKillSet.lean}{PositivityTheory.ThKEnd.Comp.SoL.kill_set}
\begin{proof}
For $Y\notin Z_0$ choose such $m_Y$ and $c_Y$, and let $U_Y$ be the set of the
$y$ with $m_Y(y)=m_Y(Y)$. It is open and closed, as $m_Y$ is locally constant,
and it contains $Y$. For a Schwartz-Bruhat $\Phi$ the function
$\mathbf 1_{U_Y}\Phi$ is Schwartz-Bruhat, and on $U_Y$ the function $m_Y-c_Y$ is
the nonzero constant $m_Y(Y)-c_Y$. So the hypothesis, applied to
$(m_Y(Y)-c_Y)^{-1}\mathbf 1_{U_Y}\Phi$, gives
\[
\mathbf 1_{U_Y}\Phi=(m_Y-c_Y)\,(m_Y(Y)-c_Y)^{-1}\mathbf 1_{U_Y}\Phi\in N.
\]

Let $T\subset Z\setminus Z_0$ be finite. Every Schwartz-Bruhat $\Phi$ such that
each $y$ with $\Phi(y)\neq0$ lies in some $U_Y$, $Y\in T$, lies in $N$. This
holds by induction on $T$. If $T$ is empty, then $\Phi=0$. Otherwise let $Y\in T$
and write $\Phi=\mathbf 1_{U_Y}\Phi+\mathbf 1_{Z\setminus U_Y}\Phi$. The first
term lies in $N$, and the second is Schwartz-Bruhat and satisfies the condition
for $T\setminus\{Y\}$.

Let $\Phi$ be Schwartz-Bruhat and vanish on $Z_0$. The set of the $y$ with
$\Phi(y)\neq0$ is closed, as $\Phi$ is locally constant, and lies in a compact
set, so it is compact; it does not meet $Z_0$. It is covered by the open sets
$U_Y$ with $Y$ in it, so by finitely many of them, and $\Phi\in N$.
\end{proof}
% Lean: kill_set (EigenKillSet.lean, 40 lines). The induction is key, over
% Finset.induction_on, with indicator_mem and isClopen_fiber; the cover is
% elim_nhds_subcover. It is the set form of KCo.kill (Lemma
% asm-theta-agree-coinvariants-kill), with the fibre of m_Y through Y in place of the
% complement of the fibre through Y_1.

\begin{lemma}\label{asm-theta-agree-end-lower-eigen}
Let $B$ be a theta base for a CM field $E$ with character $\psi$
(Definition~\ref{asm-theta-agree-models}(1)), let $X$ and $W$ be lattice spaces
of ranks $n$ and $d$ whose matrices $J_X$ and $J_W$ are invertible, let
$I=\{1,\ldots,n\}\times\{1,\ldots,d\}$, and let $M$ be a model for $(X,W)$, with
form $\langle\,,\rangle$ on $\mathbb W(X,W)$
(Definition~\ref{asm-theta-agree-models}(3)). Let $v$ be a finite place of $F$,
$t\colon E\to F_v$ an $F$-embedding, $\iota$ a linear model of the second kind
for $(t,1,1)$, with
$\operatorname{co}(u)=-t(J_X)\,\pi_{t\circ c}(u)\,t(J_W)^{\mathsf T}$
(Definition~\ref{asm-theta-agree-models}(5)), and $\rho_v(u)$, for $u$ at $v$,
the operators of Lemma~\ref{asm-theta-agree-coinvariants-eigen}(2). Let $l$ be
a linear form on the functions on $F_v^I$, and let $A$ be a set of pairs
$(u,\lambda)$, with $u\in\mathbb W(X,W)$ at $v$, $\pi_{t\circ c}(u)=0$ and
$\lambda\in\mathbf C$, such that $l(\rho_v(u)\phi)=\lambda\,l(\phi)$ for every
$(u,\lambda)\in A$ and every Schwartz-Bruhat $\phi$. Let $\phi$ be a
Schwartz-Bruhat function on $F_v^I$ such that $\iota\phi(\operatorname{co}(x))=0$
for every $x\in\mathbb W(X,W)$ at $v$ with $\pi_t(x)=0$ and
$\psi(\langle x,u\rangle)=\lambda$ for all $(u,\lambda)\in A$. Then $l(\phi)=0$.
\end{lemma}
\leanref{Development/WeilBound/ThetaLifts/EigenvalueEndgame/EigenKilling/EigenKillSet.lean}{PositivityTheory.ThKEnd.Comp.SoL.eigen_kill_set}
\begin{proof}
For $Y\in M_{n\times d}(F_v)$ let $x_Y$ be the element of
Lemma~\ref{asm-theta-agree-coinvariants-eigen}(1), at $v$ with $\pi_t(x_Y)=0$
and $\operatorname{co}(x_Y)=Y$. Let $N$ be the space of the $\iota\phi'$ with
$\phi'$ Schwartz-Bruhat and $l(\phi')=0$, and $Z_0$ the set of the $Y$ with
$\psi(\langle x_Y,u\rangle)=\lambda$ for all $(u,\lambda)\in A$. We apply
Lemma~\ref{asm-theta-agree-end-lower-kill} on $M_{n\times d}(F_v)$.

Let $Y\notin Z_0$, and choose $(u,\lambda)\in A$ with
$\psi(\langle x_Y,u\rangle)\neq\lambda$. Put $m(y)=\psi(\langle x_y,u\rangle)$.
For Schwartz-Bruhat $\phi'$, $\rho(u)$ acts at $v$ by
$\phi'\mapsto\rho_v(u)\phi'$ (Lemma~\ref{asm-theta-agree-coinvariants-eigen}(2)),
so, as $\pi_{t\circ c}(u)=0$, the definition of a linear model of the second
kind, applied at every $x_y$, gives $\iota(\rho_v(u)\phi')=m\cdot\iota\phi'$. As
in the proof of Lemma~\ref{asm-theta-agree-coinvariants-eigen}(3), $m$ is
locally constant. Let $\Phi=\iota\phi'$ be Schwartz-Bruhat. Then
$(m-\lambda)\Phi=\iota(\rho_v(u)\phi'-\lambda\phi')$ and
$l(\rho_v(u)\phi'-\lambda\phi')=\lambda\,l(\phi')-\lambda\,l(\phi')=0$, so
$(m-\lambda)\Phi\in N$. And $m(Y)\neq\lambda$.

The function $\iota\phi$ is Schwartz-Bruhat and vanishes on $Z_0$, as
$\iota\phi(Y)=\iota\phi(\operatorname{co}(x_Y))=0$ for $Y\in Z_0$. By
Lemma~\ref{asm-theta-agree-end-lower-kill}, $\iota\phi=\iota\phi'$ for some
Schwartz-Bruhat $\phi'$ with $l(\phi')=0$. As $\iota$ is injective on
Schwartz-Bruhat functions, $\phi=\phi'$, and $l(\phi)=0$.
\end{proof}
% Lean: eigen_kill_set (EigenKillSet.lean, 37 lines), with exists_coB for x_Y, mult' for the
% eigenvalue m, lc_of_mul and locRho_mem. It is the set form of KCo.eigen_kill (Lemma
% asm-theta-agree-coinvariants-eigen(3)).

\begin{lemma}\label{asm-theta-agree-end-lower-step}
Let $k$ be a nonarchimedean local field, let $n,d\geq0$, let $\chi_1$ and
$\chi_2$ be characters of $k^\times$ with $\chi_1$ unitary and smooth and
$\chi_1\neq1$ (Definition~\ref{asm-theta-agree-coinvariants-local}), and call
$Y\in M_{n\times d}(k)$ free on $S\subset\{1,\ldots,n\}$ if its rows $Y_s$,
$s\in S$, are linearly independent. Let $S\subset\{1,\ldots,n\}$ and $i\notin S$,
and let $N$ be a subspace of the functions on $M_{n\times d}(k)$ that contains
$\omega_{\chi_1,\chi_2}(g,1)\Phi'-\Phi'$ for every $g\in GL_n(k)$ and every
Schwartz-Bruhat $\Phi'$ that vanishes at every matrix not free on $S$. Let
$\Phi$ be a Schwartz-Bruhat function that vanishes at every matrix not free on
$S$. Then there is a Schwartz-Bruhat function $\Psi$ that vanishes at every
matrix not free on $S\cup\{i\}$, with $\Phi-\Psi\in N$.
\end{lemma}
\leanref{Development/WeilBound/ThetaLifts/EigenvalueEndgame/EigenKilling/EigenKillSet.lean}{PositivityTheory.ThKEnd.Comp.SoL.stepN}
\begin{proof}
This is the proof of Lemma~\ref{asm-theta-agree-coinvariants-step}, with $N$ in
place of $\mathfrak a(\chi_1,\chi_2)$. There the moves are applied only to
$\Phi$ and to $\Phi_1=\sum_q\Phi\circ m_{1,a_q}$. Both are Schwartz-Bruhat and
vanish at every matrix not free on $S$, since each $m_{x,a}$ keeps the rows
$Y_s$, $s\in S$. So $\Phi_1-N_1\Phi$, $\Phi_2-N_2\Phi_1$ and $\Phi_3-\Phi_1$ lie
in $N$, and the rest of that proof applies unchanged.
\end{proof}
% Lean: stepN (EigenKillSet.lean, 82 lines), with mv_N and sum_N. It is KCo.step with mv_aug
% and sum_aug replaced by mv_N and sum_N, which use the hypothesis on N only for functions
% with Sup S; checked against step line by line.

\begin{lemma}\label{asm-theta-agree-end-lower-corank}
Let $k$ be a nonarchimedean local field, let $n,d\geq0$, let $\chi_1$ and
$\chi_2$ be characters of $k^\times$ with $\chi_1$ unitary and smooth and
$\chi_1\neq1$ (Definition~\ref{asm-theta-agree-coinvariants-local}), and call
$Y\in M_{n\times d}(k)$ free on $S\subset\{1,\ldots,n\}$ if its rows $Y_s$,
$s\in S$, are linearly independent. Put $S_i=\{1,\ldots,n\}\setminus\{i\}$. Let
$N$ be a subspace of the functions on $M_{n\times d}(k)$ that contains
$\omega_{\chi_1,\chi_2}(g,1)\Phi'-\Phi'$ for every $i$, every $g\in GL_n(k)$ and
every Schwartz-Bruhat $\Phi'$ that vanishes at every matrix not free on $S_i$.
Let $\Phi$ be a Schwartz-Bruhat function that vanishes at every matrix that is
free on none of the sets $S_i$. Then $\Phi=\Psi+\eta$ with $\Psi$ of full rank
and $\eta\in N$.
\end{lemma}
\leanref{Development/WeilBound/ThetaLifts/EigenvalueEndgame/EigenKilling/EigenKillSet.lean}{PositivityTheory.ThKEnd.Comp.SoL.corank_one}
\leanref{Development/WeilBound/ThetaLifts/EigenvalueEndgame/EigenKilling/EigenKillSet.lean}{PositivityTheory.ThKEnd.Comp.SoL.indep_nhds}
\begin{proof}
The functions of full rank form a subspace $\mathcal F$. We apply
Lemma~\ref{asm-theta-agree-end-lower-kill} on $M_{n\times d}(k)$ to the subspace
$\mathcal F+N$ and the set $Z_0$ of the matrices free on none of the $S_i$.

Let $Y\notin Z_0$, so that $Y$ is free on $S_i$ for some $i$. The matrices free
on $S_i$ form a neighbourhood of $Y$: let $R_{Y'}$ be the matrix of the rows
$Y'_s$, $s\in S_i$, and $C$ a matrix with $R_YC=1$; then $\det(R_{Y'}C)\neq0$
for $Y'$ near $Y$, and then the rows of $R_{Y'}$ are independent. As
$M_{n\times d}(k)$ is locally compact and totally disconnected, this
neighbourhood contains an open compact neighbourhood $V$ of $Y$. Put
$m=\mathbf 1_V$ and $c=0$, so that $m(Y)\neq c$. Let $\Psi'$ be
Schwartz-Bruhat. Then $\mathbf 1_V\Psi'$ is Schwartz-Bruhat and vanishes at
every matrix not free on $S_i$, so Lemma~\ref{asm-theta-agree-end-lower-step},
for $S_i$ and $i$, gives a Schwartz-Bruhat $\Psi''$ that vanishes at every
matrix not free on $\{1,\ldots,n\}$, that is, of full rank, with
$\mathbf 1_V\Psi'-\Psi''\in N$. So $(m-c)\Psi'=\mathbf 1_V\Psi'\in\mathcal F+N$.

As $\Phi$ vanishes on $Z_0$, Lemma~\ref{asm-theta-agree-end-lower-kill} gives
$\Phi\in\mathcal F+N$.
\end{proof}
% Lean: corank_one (EigenKillSet.lean, 30 lines), through kill_set for Fu ⊔ N, indep_nhds
% (20 lines) for the neighbourhood, loc_compact_Haus_tot_disc_of_zero_dim for V, and stepN.

\begin{lemma}\label{asm-theta-agree-end-lower-rank}
Let $k$ be a field, let $n\geq1$ and $d,d_1\geq0$, and let
$Y_1\in M_{n\times d_1}(k)$ be such that for every $i$ the rows of $Y_1$ other
than the $i$-th are linearly dependent. Then $Y_1\kappa+yr^{\mathsf T}$ has rank
less than $n$ for all $\kappa\in M_{d_1\times d}(k)$, $y\in k^n$ and $r\in k^d$.
\end{lemma}
\leanref{Development/WeilBound/ThetaLifts/EigenvalueEndgame/EigenKilling/EigenKillSet.lean}{PositivityTheory.ThKEnd.Comp.SoL.rank_ne}
\begin{proof}
A matrix $C$ with $n$ rows has rank less than $n$ if and only if
$p^{\mathsf T}C=0$ for some $p\neq0$. If $p^{\mathsf T}Y_1=0$, then
$p^{\mathsf T}(Y_1\kappa+yr^{\mathsf T})=(p\cdot y)\,r^{\mathsf T}$. By the
hypothesis for $i=1$ there is $p_0\neq0$ with $p_{0,1}=0$ and
$p_0^{\mathsf T}Y_1=0$; choose $i$ with $p_{0,i}\neq0$. By the hypothesis for
$i$ there is $p_1\neq0$ with $p_{1,i}=0$ and $p_1^{\mathsf T}Y_1=0$. If
$p_1\cdot y=0$, then $p=p_1$ will do. Otherwise
$p=(p_1\cdot y)p_0-(p_0\cdot y)p_1$ will do: its $i$-th entry is
$(p_1\cdot y)p_{0,i}\neq0$, $p^{\mathsf T}Y_1=0$ and $p\cdot y=0$.
\end{proof}
% Lean: rank_ne (EigenKillSet.lean, 18 lines), with rank_ne_iff and indep_iff.

\begin{lemma}\label{asm-theta-agree-end-lower-lift}
Let $k$ be a nonarchimedean local field with normalized absolute value
$|\cdot|$, let $n,d,d_1\geq0$ and $i\in\{1,\ldots,n\}$, call a function on a
space of matrices over $k$ Schwartz-Bruhat if it is locally constant with compact
support, and call a matrix with $n$ rows free on $S_i=\{1,\ldots,n\}\setminus\{i\}$
if its rows other than the $i$-th are linearly independent. Let $\mathcal A$ be a
complex vector space with a subspace $\mathcal A_0$, let $\iota$ be a linear map
from $\mathcal A$ to the functions on $M_{n\times d}(k)$ that maps $\mathcal A_0$
bijectively onto the Schwartz-Bruhat functions, and let $\tau$ be a linear map
from $\mathcal A$ to the functions on $M_{n\times d_1}(k)$. Let
$\kappa\in M_{d_1\times d}(k)$, $P\in M_{d\times d_1}(k)$ and $r,s\in k^d$ with
$\kappa P=1$, $r^{\mathsf T}P=0$, $\kappa s=0$ and $r\cdot s=1$. Suppose:
\begin{enumerate}
\item[(a)] if $\phi\in\mathcal A_0$ and $Y_1\in M_{n\times d_1}(k)$ have
  $\iota\phi(Y_1\kappa+yr^{\mathsf T})=0$ for every $y\in k^n$, then
  $\tau\phi(Y_1)=0$;
\item[(b)] if $\phi,\phi'\in\mathcal A_0$, $Z_1\in M_{n\times d_1}(k)$ and
  $y\in k^n$ have $\iota\phi'(Y)=\iota\phi(Y+Z_1\kappa+yr^{\mathsf T})$ for every
  $Y$, then $\tau\phi'(Y_1)=\tau\phi(Y_1+Z_1)$ for every $Y_1$;
\item[(c)] $\tau\phi\neq0$ for some $\phi\in\mathcal A_0$.
\end{enumerate}
Then for every Schwartz-Bruhat $\Psi$ on $M_{n\times d_1}(k)$ that vanishes at
every matrix not free on $S_i$ there is $\phi\in\mathcal A_0$ with $\iota\phi$
of full rank (Definition~\ref{asm-theta-agree-coinvariants-local}(2)) and
$\tau\phi=\Psi$.
\end{lemma}
\leanref{Development/WeilBound/ThetaLifts/EigenvalueEndgame/EigenKilling/EigenKillSet.lean}{PositivityTheory.ThKEnd.Comp.SoL.lift}
\begin{proof}
Let $e_i$ be the $i$-th standard basis vector of $k^n$. For a Schwartz-Bruhat
$\Phi$ and a matrix $Z$, the function $Y\mapsto\Phi(Y+Z)$ is Schwartz-Bruhat, so
it is $\iota\phi'$ for some $\phi'\in\mathcal A_0$.

\emph{Step 1: a normalized element.} By (c) there are $\phi_0\in\mathcal A_0$
and $Y_0$ with $\tau\phi_0(Y_0)\neq0$. Let $\phi_1\in\mathcal A_0$ have
$\iota\phi_1(Y)=\iota\phi_0(Y+Y_0\kappa)$. By (b) with $Z_1=Y_0$ and $y=0$,
$\tau\phi_1(0)=\tau\phi_0(Y_0)\neq0$.

\emph{Step 2: constants.} The support of $\Psi$ is compact and its elements are
free on $S_i$. As in the proof of Lemma~\ref{asm-theta-agree-coinvariants-step}
(boxes), there is $c_0\in k$ such that $|b_s|\leq|c_0|$ for every $s\in S_i$
whenever $\Psi(Y_1)\neq0$ and $b\in k^{S_i}$ has $Y_{1,i}=\sum_{s}b_sY_{1,s}$,
where $Y_{1,j}$ are the rows of $Y_1$. Let $c_2=1$ if $|c_0|\leq1$ and $c_2=c_0$
otherwise, so that $|c_2|\geq1$ and $|c_2|\geq|c_0|$. The image of the closed
support of $\iota\phi_1$ under $W\mapsto Ws$ is compact, so there is $c_1\neq0$
with $|(Ws)_j|\leq|c_1|$ for every $j$ and every $W$ in it. Choose $r_0\in k$ with
$0<|r_0|<1$, and put $N=r_0^{-1}c_2c_1$, so that $|N|>|c_2c_1|$.

\emph{Step 3: the function.} Put $T(Y)=Y-YP\kappa-Ne_ir^{\mathsf T}$ and
\[
\Phi(Y)=\tau\phi_1(0)^{-1}\,\Psi(YP)\,\iota\phi_1(T(Y)).
\]
It is locally constant. If $\Phi(Y)\neq0$, then $YP$ lies in the support of
$\Psi$, $T(Y)$ in that of $\iota\phi_1$, and $Y=T(Y)+(YP)\kappa+Ne_ir^{\mathsf T}$;
so $\Phi$ has compact support. Let $\phi\in\mathcal A_0$ have $\iota\phi=\Phi$.

\emph{Step 4: full rank.} Let $Y$ have rank less than $n$, and suppose that
$\Phi(Y)\neq0$, so that $\Psi(YP)\neq0$ and $\iota\phi_1(T(Y))\neq0$. Choose
$p\neq0$ with $p^{\mathsf T}Y=0$; then $p^{\mathsf T}YP=0$. As $YP$ is free on
$S_i$, $p_i\neq0$, and $(YP)_i=\sum_{s\in S_i}(-p_s/p_i)(YP)_s$; so
$|p_s/p_i|\leq|c_0|$ for $s\in S_i$. As $\kappa s=0$ and $r\cdot s=1$,
$T(Y)s=Ys-Ne_i$, and so $p\cdot T(Y)s=-p_iN$. Hence
\[
N=-(T(Y)s)_i+\sum_{s'\in S_i}(-p_{s'}/p_i)\,(T(Y)s)_{s'}.
\]
The entries of $T(Y)s$ have absolute value at most $|c_1|$, as $T(Y)$ lies in
the support of $\iota\phi_1$. By the ultrametric inequality
$|N|\leq\max(|c_1|,|c_0||c_1|)\leq|c_2c_1|$, which is false. So $\Phi(Y)=0$, and
$\iota\phi$ has full rank.

\emph{Step 5: the value.} Let $Y_1\in M_{n\times d_1}(k)$, and let
$\phi'\in\mathcal A_0$ have
$\iota\phi'(Y)=\iota\phi_1(Y-Y_1\kappa-Ne_ir^{\mathsf T})$. By (b) with $Z_1=-Y_1$
and $y=-Ne_i$, $\tau\phi'(Y_1)=\tau\phi_1(0)$. Put
$c=\tau\phi_1(0)^{-1}\Psi(Y_1)$. For $y\in k^n$,
$(Y_1\kappa+yr^{\mathsf T})P=Y_1$, as $\kappa P=1$ and $r^{\mathsf T}P=0$, and so
$T(Y_1\kappa+yr^{\mathsf T})=(y-Ne_i)r^{\mathsf T}$. Hence
\[
\iota(\phi-c\phi')(Y_1\kappa+yr^{\mathsf T})
=c\,\iota\phi_1((y-Ne_i)r^{\mathsf T})-c\,\iota\phi_1((y-Ne_i)r^{\mathsf T})=0
\]
for every $y$. By (a), $\tau\phi(Y_1)=c\,\tau\phi'(Y_1)=\Psi(Y_1)$.
\end{proof}
% Lean: lift (EigenKillSet.lean, 102 lines), with sb_shift. Step 1 is phi_1 and hc; Step 2
% is exists_bound and bx_large for c_0, c_2, bx_large for c_1 on the image under
% W -> W y', and Valuation.IsNontrivial.exists_lt_one for r_0; Step 3 is T, Phi and hPhi;
% Step 4 is the full-rank branch (rank_ne_iff, indep_iff for hpi, hb', hw, ew, z, eN);
% Step 5 is t1 (htr) and t2 (hkill). The matrices P, r and s are the Lean's iota', rho
% and y'.

\begin{lemma}\label{asm-theta-agree-end-lower-local}
Let $(X,W,M,\chi_X,\chi_W,\omega)$ and $(X,W_1,M_1,\chi_X,\chi_W,\omega_1)$ be
Weil pairs for a theta base $B$ of a CM field $E$ with character $\psi$
(Definition~\ref{asm-theta-agree-triple}(1)), with $X$, $W$ and $W_1$ of ranks
$n\geq1$, $d$ and $d_1\geq1$ and matrices $J_X$, $J_W$ and $J_{W_1}$, and put
$I=\{1,\ldots,n\}\times\{1,\ldots,d\}$. Let $H'=(x',y',\iota',\mathfrak b')$ be
an adapted plane in $W$ over $W_1$ at some positions
(Definition~\ref{asm-theta-agree-planes}(1)), and $e$ an evaluation map for
$(M,M_1,1,\iota',\mathcal J_0(H'))$ with components $e_\infty$ and $e_w$ at the
finite places $w$ (Definition~\ref{asm-theta-agree-planes}(5)). Write $v=v_1$
and $t=\tau_1$, and
$t(A)$ for the matrix of the images under $t$ of the entries of a matrix $A$
over $E$. Put $J=t(J_W)$, $J_1=t(J_{W_1})$, $P=t(\iota')$ and
\[
\kappa=(J_1^{\mathsf T})^{-1}\,t(c(\iota'))^{\mathsf T}J^{\mathsf T},\qquad
r=J\,t(c(x')),\qquad s=-t(y').
\]
Let $\iota$ and $\iota_1$ be linear models of the second kind for $(t,1,1)$ for
$M$ and $M_1$ (Definition~\ref{asm-theta-agree-models}(5)). Then:
\begin{enumerate}
\item $\kappa P=1$, $r^{\mathsf T}P=0$, $\kappa s=0$ and $r\cdot s=1$;
\item if $\phi$ and $\phi'$ are Schwartz-Bruhat functions on $F_v^I$,
  $Z_1\in M_{n\times d_1}(F_v)$, $y_0\in F_v^n$ and
  $\iota\phi'(Y)=\iota\phi(Y+Z_1\kappa+y_0r^{\mathsf T})$ for every $Y$, then
  $\iota_1e_v(\phi')(Y_1)=\iota_1e_v(\phi)(Y_1+Z_1)$ for every $Y_1$;
\item if $\phi$ is a Schwartz-Bruhat function on $F_v^I$ and
  $Y_1\in M_{n\times d_1}(F_v)$ has $\iota\phi(Y_1\kappa+y_0r^{\mathsf T})=0$ for
  every $y_0\in F_v^n$, then $\iota_1e_v(\phi)(Y_1)=0$.
\end{enumerate}
\end{lemma}
\leanref{Development/WeilBound/ThetaLifts/EigenvalueEndgame/EigenKilling/EigenKillSet.lean}{PositivityTheory.ThKEnd.Comp.SoL.soLower_of}
\begin{proof}
Write $h_W(u,u')=u^{\mathsf T}J_Wc(u')$, and $\operatorname{co}$ and
$\operatorname{co}_1$ for the maps of a linear model of the second kind for
$(t,1,1)$ for $M$ and $M_1$:
$\operatorname{co}(u)=-t(J_X)\pi_{t\circ c}(u)J^{\mathsf T}$ and
$\operatorname{co}_1(z)=-t(J_X)\pi_{t\circ c}(z)J_1^{\mathsf T}$. The matrices
$J_X$, $J_W$ and $J_{W_1}$ are invertible
(Lemma~\ref{asm-theta-agree-convergence-signature}), so $t(J_X)$, $J$ and $J_1$
are. As $W$ is skew-Hermitian, $h_W(u',u)=-c(h_W(u,u'))$. For $w\in E^d$,
\[
r\cdot t(w)=t(w)^{\mathsf T}J\,t(c(x'))=t(h_W(w,x')),\tag{a}
\]
and the $j$-th entry of $t(c(\iota'))^{\mathsf T}J^{\mathsf T}t(w)$ is
$t(h_W(w,\iota'e_j))=-t(c(h_W(\iota'e_j,w)))$; so
\[
\kappa\,t(w)=0\quad\text{if }h_W(\iota'u,w)=0\text{ for every }u\in E^{d_1}.\tag{b}
\]

(1) As $\iota'^{\mathsf T}J_Wc(\iota')=J_{W_1}$, $P^{\mathsf T}J\,t(c(\iota'))=J_1$;
transposing, $t(c(\iota'))^{\mathsf T}J^{\mathsf T}P=J_1^{\mathsf T}$, and
$\kappa P=1$. By (a), $r\cdot Pt(u)=t(h_W(\iota'u,x'))=0$ for $u\in E^{d_1}$, so
$r^{\mathsf T}P=0$. By (b), as $h_W(\iota'u,y')=0$, $\kappa s=0$. By (a),
$r\cdot s=-t(h_W(y',x'))=t(c(h_W(x',y')))=1$.

\emph{Evaluation at $v$.} Let $\rho_v$ and $\rho_{1,v}$ be the operators of
Lemma~\ref{asm-theta-agree-coinvariants-eigen}(2) for $M$ and $M_1$, and let
$\phi$ be Schwartz-Bruhat. The function $e_v(\phi)$ is Schwartz-Bruhat. For
$z\in\mathbb W(X,W_1)$ at $v$, $e(\rho(z\iota'^{\mathsf T})\Phi)=\rho_1(z)e(\Phi)$
for Schwartz-Bruhat $\Phi$, and $\rho(z\iota'^{\mathsf T})$ and $\rho_1(z)$ act at
$v$ by $\rho_v(z\iota'^{\mathsf T})$ and $\rho_{1,v}(z)$; so
Lemma~\ref{asm-theta-agree-local-factor-evaluation}(1) and~(2) give
\[
e_v(\rho_v(z\iota'^{\mathsf T})\phi)=\rho_{1,v}(z)\,e_v(\phi).\tag{c}
\]
For $a\in\mathbf A_E^n$ at $v$, $ax'^{\mathsf T}\in\mathcal J_0(H')$ is at $v$,
so Lemma~\ref{asm-theta-agree-local-factor-evaluation}(3) gives
\[
e_v(\rho_v(ax'^{\mathsf T})\phi)=e_v(\phi).\tag{d}
\]
As $\pi_{t\circ c}(\iota')=t(c(\iota'))$ and $\pi_{t\circ c}(x')=t(c(x'))$,
\[
\operatorname{co}(z\iota'^{\mathsf T})=\operatorname{co}_1(z)\,\kappa,\qquad
\operatorname{co}(ax'^{\mathsf T})=\bigl(-t(J_X)\pi_{t\circ c}(a)\bigr)r^{\mathsf T}.\tag{e}
\]

(2) By Lemma~\ref{asm-theta-agree-coinvariants-eigen}(1) for $M_1$,
$Z_1=\operatorname{co}_1(z)$ for some $z$ at $v$ with $\pi_t(z)=0$. By
Lemma~\ref{asm-theta-agree-dictionary-split}(1) there is $a\in\mathbf A_E^n$ at
$v$ with $\pi_t(a)=0$ and $\pi_{t\circ c}(a)=-t(J_X)^{-1}y_0$. Put
$\phi''=\rho_v(z\iota'^{\mathsf T})\rho_v(ax'^{\mathsf T})\phi$. The elements
$z\iota'^{\mathsf T}$ and $ax'^{\mathsf T}$ are at $v$, and their images under
$\pi_t$ are $0$. So the definition of a linear model of the second kind, with
Lemma~\ref{asm-theta-agree-coinvariants-eigen}(2), and (e) give
\[
\iota\phi''(Y)=\iota\phi\bigl(Y+\operatorname{co}(z\iota'^{\mathsf T})+\operatorname{co}(ax'^{\mathsf T})\bigr)
=\iota\phi(Y+Z_1\kappa+y_0r^{\mathsf T})=\iota\phi'(Y).
\]
As $\iota$ is injective on Schwartz-Bruhat functions, $\phi'=\phi''$. By (c)
and (d), $e_v(\phi')=\rho_{1,v}(z)e_v(\phi)$. As $\pi_t(z)=0$, the definition of
a linear model of the second kind for $\iota_1$ gives
$\iota_1e_v(\phi')(Y_1)=\iota_1e_v(\phi)(Y_1+\operatorname{co}_1(z))$.

(3) \emph{The form.} By Lemma~\ref{asm-theta-agree-dictionary-coordinates}(2),
in Situation~\ref{asm-theta-agree-dictionary-situation} for $J_X$, $J_W$ and
$\varepsilon=1$ (proof of Lemma~\ref{asm-theta-agree-coinvariants-dictionary}),
for $x\in\mathbb W(X,W)$ at $v$ with $\pi_t(x)=0$ and every $u$, the element
$\langle x,u\rangle$ is at $v$ with component
$-\operatorname{tr}(\operatorname{co}(x)^{\mathsf T}\pi_t(u))$. For $\sigma\in F_v$
let $\bar\sigma\in\mathbf A$ have component $\sigma$ at $v$ and $0$ elsewhere, and
let $\varepsilon\in\mathbf A_E$ be at $v$ with $\pi_t(\varepsilon)=1$ and
$\pi_{t\circ c}(\varepsilon)=0$ (Lemma~\ref{asm-theta-agree-dictionary-split}(1)).
For $p\in E^n$ and $q\in E^d$, $\pi_t(\bar\sigma\varepsilon pq^{\mathsf T})=\sigma\,t(p)t(q)^{\mathsf T}$,
so
\[
\langle x,\bar\sigma\varepsilon\,pq^{\mathsf T}\rangle=\overline{-\sigma\,t(p)^{\mathsf T}\operatorname{co}(x)\,t(q)},\tag{f}
\]
and the same holds for $M_1$, $\operatorname{co}_1$ and $q\in E^{d_1}$. Finally,
there is $\sigma_0\in F_v$ with $\psi(\bar\sigma_0)\neq1$ (proof of
Lemma~\ref{asm-theta-agree-coinvariants-eigen}(3)); so if $b,b'\in F_v$ have
$\psi(\overline{b\sigma})=\psi(\overline{b'\sigma})$ for every $\sigma$, then
$b=b'$, as otherwise $\sigma=\sigma_0/(b-b')$ gives a contradiction.

\emph{The eigenvalues.} By Lemma~\ref{asm-theta-agree-coinvariants-eigen}(1) for
$M_1$, $Y_1=\operatorname{co}_1(x_1)$ for some $x_1$ at $v$ with $\pi_t(x_1)=0$.
Let $l(\phi')=\iota_1e_v(\phi')(Y_1)$, and let $A$ be the set of the pairs
\[
\bigl(\bar\sigma\varepsilon\,p(\iota'q)^{\mathsf T},\ \psi(\langle x_1,\bar\sigma\varepsilon\,pq^{\mathsf T}\rangle_1)\bigr)
\quad\text{and}\quad
\bigl(\bar\sigma\varepsilon\,px'^{\mathsf T},\ 1\bigr)
\]
with $\sigma\in F_v$, $p\in E^n$ and $q\in E^{d_1}$, where $\langle\,,\rangle_1$
is the form on $\mathbb W(X,W_1)$. Their first entries $u$ are at $v$ with
$\pi_{t\circ c}(u)=0$. Let $\phi'$ be Schwartz-Bruhat. For a pair of the first
kind, $u=z\iota'^{\mathsf T}$ with $z=\bar\sigma\varepsilon pq^{\mathsf T}$, which
has $\pi_{t\circ c}(z)=0$; so by (c) and the definition of a linear model of the
second kind for $\iota_1$ at $x_1$,
$l(\rho_v(u)\phi')=\iota_1(\rho_{1,v}(z)e_v(\phi'))(\operatorname{co}_1(x_1))=\psi(\langle x_1,z\rangle_1)\,l(\phi')$.
For a pair of the second kind, $l(\rho_v(u)\phi')=l(\phi')$ by (d).

\emph{The condition.} Let $x$ be at $v$ with $\pi_t(x)=0$ and
$\psi(\langle x,u\rangle)=\lambda$ for every $(u,\lambda)\in A$, and put
$C=\operatorname{co}(x)$. For the standard basis vectors $p=e_i$ and $q=e_j$,
the pairs of the first kind and (f) give
$\psi(\overline{-\sigma\,e_i^{\mathsf T}CPe_j})=\psi(\overline{-\sigma\,e_i^{\mathsf T}Y_1e_j})$
for every $\sigma$, as $t(\iota'e_j)=Pe_j$; so $CP=Y_1$. The pairs of the second
kind with $p=e_i$ give $\psi(\overline{-\sigma\,e_i^{\mathsf T}C\,t(x')})=1$ for
every $\sigma$, so $C\,t(x')=0$. We show that
$C=Y_1\kappa+(Cs)r^{\mathsf T}$. Let $w\in E^d$. By the definition of an adapted
plane, every vector of the lattice of $W$ is $ax'+\iota'u+by'$ with $a,b\in E$
and $u\in E^{d_1}$, and the lattice spans $E^d$; so $w$ is of this form, and
$t(w)=t(a)t(x')+Pt(u)-t(b)s$. Then $Ct(w)=Y_1t(u)-t(b)\,Cs$. By (b),
$\kappa t(x')=0$, as $h_W(\iota'u,x')=0$; with (1), $\kappa t(w)=t(u)$. By (a),
$r\cdot t(x')=t(h_W(x',x'))=0$; with (1), $r\cdot t(w)=-t(b)$. So
$(Y_1\kappa+(Cs)r^{\mathsf T})t(w)=Y_1t(u)-t(b)\,Cs=Ct(w)$. Taking for $w$ the
standard basis vectors gives the claim. Hence
$\iota\phi(\operatorname{co}(x))=\iota\phi(Y_1\kappa+(Cs)r^{\mathsf T})=0$, by the
hypothesis for $y_0=Cs$.

By Lemma~\ref{asm-theta-agree-end-lower-eigen}, for $M$, $\iota$, $l$, $A$ and
$\phi$, $l(\phi)=0$, that is, $\iota_1e_v(\phi)(Y_1)=0$.
\end{proof}
% Lean: the facts proved inside soLower_of (EigenKillSet.lean). The matrices are J, J_1,
% iota' = H.iota.map t, iota-bar = H.iota.map (t o sigma), kappa, rho and y'; (a) is hrho
% (map_form), (b) is hkappa (with form_swap through sw); (1) is gram, e1, e2', e2, e3 and
% e4; (c) is lmid and (d) is linv, both through SoL.ev_loc, which proves the statement of
% Lemma asm-theta-agree-local-factor-evaluation(2) for c = 1 again, by the same argument
% (exists_continuousLinearMap_tens_update and position_inj); (e) is cmid and cvv; (2) is
% htr, with exists_coB, exists_pair, hiota.trans twice and hiota_1.trans; (3) is hkill,
% through eigen_kill_set with the set A of the Lean, ev1 (from linEx_of, the second part of
% LinEx, through symp_eq) for (f), sep for the last fact of the form, c1, c2 and c3 (with
% plane_span for the decomposition of w).

\begin{lemma}\label{asm-theta-agree-end-alpha-zero-lower}
Let $(X,W,M,\chi_X,\chi_W,\omega)$ and $(X,W_1,M_1,\chi_X,\chi_W,\omega_1)$ be
Weil pairs for a theta base $B$ of a CM field $E$
(Definition~\ref{asm-theta-agree-triple}(1)), with $X$, $W$ and $W_1$ of ranks
$n\geq1$, $d$ and $d_1\geq1$, and put $I=\{1,\ldots,n\}\times\{1,\ldots,d\}$.
Let $H'$ be an adapted plane in $W$ over $W_1$ at some positions
(Definition~\ref{asm-theta-agree-planes}(1)), and suppose that
$(\chi_W)_t\neq1$ for every $F$-embedding $t\colon E\to F_{v_1}$
(Definition~\ref{asm-theta-agree-coinvariants-local}(4)). Let $e$ be an
evaluation map for $(M,M_1,1,\iota_{H'},\mathcal J_0(H'))$ with components
$e_\infty$ and $e_v$ (Definition~\ref{asm-theta-agree-planes}(5)). Let $S_1$ be
the set of Schwartz-Bruhat functions on $F_{v_1}^I$ of full rank at $\tau_1$
for $M$ (Definition~\ref{asm-theta-agree-models}(5)), and $\mathfrak a$ the
span of the $\phi_1-\phi$ with $\phi\in S_1$, $\phi_1$ Schwartz-Bruhat and
$\omega(g,1)$ acting at $v_1$ by $\phi\mapsto\phi_1$ for some
$g\in U_X^{(v_1)}$. Then for every $\phi\in S_1$ there is $\phi'\in S_1$ with
$\phi-\phi'\in\mathfrak a$ and $e_{v_1}(\phi')$ of full rank at $\tau_1$ for
$M_1$.
\end{lemma}
\leanref{Development/WeilBound/ThetaLifts/EigenvalueEndgame/EigenKilling/ConeIntegralFactorisation.lean}{PositivityTheory.ThKEnd.Comp.SoLower}
\leanref{Development/WeilBound/ThetaLifts/EigenvalueEndgame/EigenKilling/EigenKillSet.lean}{PositivityTheory.ThKEnd.Comp.SoL.soLower_of}
\begin{proof}
Write $v=v_1$, $t=\tau_1$ and $k=F_v$, let $H'=(x',y',\iota',\mathfrak b')$, and
let $\kappa$, $P$, $r$ and $s$ be the matrices and vectors of
Lemma~\ref{asm-theta-agree-end-lower-local}.

\emph{Step 1: dictionaries.} By
Lemma~\ref{asm-theta-agree-coinvariants-inputs}(3) for the Weil pair
$(X,W,M,\chi_X,\chi_W,\omega)$ at $v$ and $t$, there are a linear model $\iota$
of the second kind for $(t,1,1)$ for $M$ and unitary smooth characters $\nu$ and
$\alpha$ of $k^\times$ such that $\iota$ is a dictionary at $v$ with characters
$(\nu,\alpha)$ (Definition~\ref{asm-theta-agree-coinvariants-local}(3)). In the
same way, for $(X,W_1,M_1,\chi_X,\chi_W,\omega_1)$ there are $\iota_1$ for
$M_1$ and unitary smooth $\nu_1$ and $\alpha_1$ such that $\iota_1$ is a
dictionary with characters $(\nu_1,\alpha_1)$, and, as $n,d_1\geq1$, if
$\nu_1=1$ then $(\chi_W)_t=1$ or $(\chi_W)_{t\circ c}=1$. Both are excluded by
hypothesis, as $t\circ c$ is also an $F$-embedding; so $\nu_1\neq1$.

\emph{Step 2: full rank.} By
Lemma~\ref{asm-theta-agree-local-augmentation-subspace}(1) and the definition,
a Schwartz-Bruhat function $\phi$ on $F_v^I$ lies in $S_1$ if and only if
$\iota\phi$ has full rank. So $S_1$ is a subspace. A Schwartz-Bruhat $\phi_1$ on
$F_v^{I_1}$, with $I_1=\{1,\ldots,n\}\times\{1,\ldots,d_1\}$, has full rank at
$t$ for $M_1$ if $\iota_1\phi_1$ has full rank. For $i\in\{1,\ldots,n\}$ put
$S_i=\{1,\ldots,n\}\setminus\{i\}$, and say that a matrix is free on $S_i$ if
its rows other than the $i$-th are linearly independent.

\emph{Step 3: a space of relations.} Let $N$ be the space of the
$\iota_1e_v(\eta)$ with $\eta\in S_1\cap\mathfrak a$. We show that $N$ contains
$\omega_{\nu_1,\alpha_1}(g_0,1)\Psi-\Psi$ for every $i$, every
$g_0\in GL_n(k)$ and every Schwartz-Bruhat $\Psi$ on $M_{n\times d_1}(k)$ that
vanishes at every matrix not free on $S_i$.

Apply Lemma~\ref{asm-theta-agree-end-lower-lift} to $\mathcal A$ the functions
on $F_v^I$, $\mathcal A_0$ the Schwartz-Bruhat functions, $\iota$,
$\tau=\iota_1\circ e_v$, $\kappa$, $P$, $r$ and $s$. Its conditions on the
matrices are Lemma~\ref{asm-theta-agree-end-lower-local}(1), and (a) and (b) are
parts (3) and (2) of that lemma. For (c), $e_v(\mathrm{vac}_v)=\mathrm{vac}_{1,v}$
is a nonzero Schwartz-Bruhat function, so $\iota_1(\mathrm{vac}_{1,v})\neq0$. The
lemma gives a Schwartz-Bruhat $\chi$ with $\iota\chi$ of full rank, so that
$\chi\in S_1$, and $\iota_1e_v(\chi)=\Psi$.

By condition (a) of a dictionary for $\iota_1$, there is $g\in U_X^{(v)}$ such
that for every Schwartz-Bruhat $\phi$ there is a Schwartz-Bruhat $\phi_1$ with
$\omega_1(g,1)$ acting at $v$ by $\phi\mapsto\phi_1$ and
$\iota_1\phi_1=\omega_{\nu_1,\alpha_1}(g_0,1)\iota_1\phi$. By condition (a) for
$\iota$, there is $g_1\in GL_n(k)$ such that for every Schwartz-Bruhat $\phi$
there is a Schwartz-Bruhat $\phi_1$ with $\omega(g,1)$ acting at $v$ by
$\phi\mapsto\phi_1$ and $\iota\phi_1=\omega_{\nu,\alpha}(g_1,1)\iota\phi$. Let
$\chi_1$ be given by the second statement for $\chi$, and $\chi_2$ by the first
for $e_v(\chi)$.

By Lemma~\ref{asm-theta-agree-split-inputs}(4), for the Levi element $1$ for
$(1,1)$ and $1\in N_{H'}$, $e(\omega(g,1)\Phi)=\omega_1(g,1)e(\Phi)$ for every
Schwartz-Bruhat $\Phi$. So
Lemma~\ref{asm-theta-agree-local-factor-evaluation}(2), with $S=\omega(g,1)$,
$S_1=\omega_1(g,1)$ and $c=1$, gives $e_v(\chi_1)=\chi_2$. The function
$\iota\chi_1(Y)=\nu(\det g_1)|\det g_1|^{d/2}\,\iota\chi(g_1^{\mathsf T}Y)$ has
full rank, as $g_1^{\mathsf T}Y$ has the rank of $Y$; so $\chi_1\in S_1$, and
$\eta=\chi_1-\chi\in S_1$. By the definition of $\mathfrak a$, $\eta\in\mathfrak a$.
Finally
\[
\iota_1e_v(\eta)=\iota_1\chi_2-\iota_1e_v(\chi)=\omega_{\nu_1,\alpha_1}(g_0,1)\Psi-\Psi.
\]

\emph{Step 4: conclusion.} Let $\phi\in S_1$, so that $\iota\phi$ has full rank,
and put $\Phi=\iota_1e_v(\phi)$, a Schwartz-Bruhat function. Let $Y_1$ be free on
none of the $S_i$. By Lemma~\ref{asm-theta-agree-end-lower-rank}, as $n\geq1$,
every $Y_1\kappa+y_0r^{\mathsf T}$ has rank less than $n$, so
$\iota\phi(Y_1\kappa+y_0r^{\mathsf T})=0$ for every $y_0\in k^n$, and
$\Phi(Y_1)=0$ by Lemma~\ref{asm-theta-agree-end-lower-local}(3). By Step 3 and
Lemma~\ref{asm-theta-agree-end-lower-corank} for $\nu_1\neq1$,
$\Phi=\Psi+\iota_1e_v(\eta)$ with $\Psi$ of full rank and $\eta\in S_1\cap\mathfrak a$.
Put $\phi'=\phi-\eta$. Then $\phi'\in S_1$ and $\phi-\phi'=\eta\in\mathfrak a$.
And $e_v(\phi')$ is Schwartz-Bruhat with $\iota_1e_v(\phi')=\Phi-(\Phi-\Psi)=\Psi$
of full rank, so $e_v(\phi')$ has full rank at $t$ for $M_1$.
\end{proof}
% Lean: SoL.soLower_of MixW (EigenKilling/EigenKillSet.lean, 205 lines). Step 1 is
% dictionary_statement (DictSt) twice, with hD.left, hD_1.leftOnto and hnu1 from its clause
% h1; Step 2 is so_iff (eigen_kill, as in full_of_soLoc) and Fu; Step 3 is the hypothesis of
% corank_one, through lift (with hkill, htr, e1 to e4 and hne), hD_1.leftOnto, hD.left,
% e2 (ev_loc with MixW at g, 1, 1, 1), b1 (omB and rank_mul_eq_right_of_isUnit_det) and
% the generator of augOn; Step 4 is corank_one with rank_ne, hkill and hfull, and the final
% refine. Size: about 800 lines in all, all in EigenKillSet.lean.

\begin{lemma}\label{asm-theta-agree-end-alpha-zero-move}
Let $(B,G,\mathcal X,\omega,\omega')$ be a triple of ranks $(a,a')$ for a theta
base $B$ of a CM field $E$ (Definition~\ref{asm-theta-agree-triple}), with
models $M$ and $M'$ and theta integrals $I$ and $I'$, let $H'$ be an adapted
plane in $W$ over a lattice space $W_1$ at some positions
(Definition~\ref{asm-theta-agree-planes}(1)), let
$(V',W_1,M_1,\alpha,\alpha/\beta,\omega_1)$ be a Weil pair of signature
$(q,a'-q;r_1,s_1)$ for some $r_1,s_1$, and let $e$ be an evaluation map for
$(M',M_1,1,\iota_{H'},\mathcal J_0(H'))$
(Definition~\ref{asm-theta-agree-planes}(5)). Then for every
$\Phi'\in\operatorname{PD}(M')\cap\mathcal S^\circ(M')$ there is
$\Psi\in\operatorname{PD}(M')\cap\mathcal S^\circ(M')$ with $I'(\Psi)=I'(\Phi')$
and $e(\Psi)\in\mathcal S^\circ(M_1)$.
\end{lemma}
\leanref{Development/WeilBound/ThetaLifts/EigenvalueEndgame/EigenKilling/ConeIntegralFactorisation.lean}{PositivityTheory.ThKEnd.Comp.SoMove}
\leanref{Development/WeilBound/ThetaLifts/EigenvalueEndgame/EigenKilling/SubMemKer.lean}{PositivityTheory.ThKEnd.Comp.soMove_of}
\begin{proof}
Let $e_\infty$ and $e_v$ be components of $e$, let $\mathrm{vac}'_v$ and
$\mathrm{vac}_{1,v}$ be the vacua of $M'$ and $M_1$, and put
$I_0=\{1,\ldots,a'\}\times\{1,\ldots,a+a'\}$. The Weil pair
$(V',W,M',\alpha,\alpha/\beta,\omega')$ of the triple has signature
$(q,a'-q;a+a'-p-q,p+q)$. Let $S_1$ and $S_1^{(1)}$ be the sets of
Schwartz-Bruhat functions at $v_1$ of full rank at $\tau_1$ for $M'$ and for
$M_1$, and $\mathfrak a$ the span of the $\phi_1-\phi$ with $\phi\in S_1$,
$\phi_1$ Schwartz-Bruhat and $\omega'(g,1)$ acting at $v_1$ by
$\phi\mapsto\phi_1$ for some $g\in U_{V'}^{(v_1)}$. A function
$f\otimes\bigotimes_v\phi_v$, with $\phi_v$ Schwartz-Bruhat and
$\phi_v=\mathrm{vac}'_v$ for almost all $v$, lies in $\mathcal S^\circ(M')$
whenever $\phi_{v_1}$ lies in the span of $S_1$, as it is linear in
$\phi_{v_1}$; likewise for $M_1$.

\emph{Step 1: a projection.} Let $\mathcal G$ be the space of the $\phi$ in the
span of $S_1$ with $e_{v_1}(\phi)$ in the span of $S_1^{(1)}$. There is a
linear map $\pi$ on the functions on $F_{v_1}^{I_0}$ with values in $\mathcal G$
such that $\phi-\pi(\phi)\in\mathfrak a$ for every $\phi\in\mathcal G+\mathfrak a$:
compose the quotient map modulo $\mathfrak a$, a linear projection of the
quotient onto the image of $\mathcal G$, and a linear section over this image
of the map from $\mathcal G$. Every $\pi(\phi)$ lies in the span of $S_1$, so
it is Schwartz-Bruhat.

\emph{Step 2: an operator.} For a function $\Psi$ on $\mathbf A^{I_0}$ and
$x\in\mathbf A^{I_0}$, let $\Psi_x$ be the function on $F_{v_1}^{I_0}$ that
maps $y$ to the value of $\Psi$ at $x$ with its component at $v_1$ replaced by
$y$, and put $(P\Psi)(x)=\pi(\Psi_x)(x_{v_1})$. The operator $P$ is linear. If
$\Psi=f\otimes\bigotimes_v\phi_v$ with $\phi_v=\mathrm{vac}'_v$ for almost all
$v$, then $\Psi_x=c_x\phi_{v_1}$ for a number $c_x$, so $P\Psi$ is the same
tensor with $\pi(\phi_{v_1})$ in place of $\phi_{v_1}$.

\emph{Step 3: the evaluation.} We show that
$P\Phi'\in\operatorname{PD}(M')\cap\mathcal S^\circ(M')$ and
$e(P\Phi')\in\mathcal S^\circ(M_1)$. As $P$ and $e$ are linear, it suffices to
treat a spanning element $\Psi=f\otimes\bigotimes_v\phi_v$ of
$\operatorname{PD}(M')$, with $f\in\operatorname{Pol}(M')$, every $\phi_v$
Schwartz-Bruhat and $\phi_v=\mathrm{vac}'_v$ for almost all $v$. By Step 2,
$P\Psi$ is this tensor with $\pi(\phi_{v_1})$ at $v_1$, which is
Schwartz-Bruhat and in the span of $S_1$; so
$P\Psi\in\operatorname{PD}(M')\cap\mathcal S^\circ(M')$. Further
$e(P\Psi)=e_\infty(f)\otimes\bigotimes_v\psi_v$ with
$\psi_{v_1}=e_{v_1}(\pi(\phi_{v_1}))$, which is in the span of $S_1^{(1)}$ as
$\pi(\phi_{v_1})\in\mathcal G$, and $\psi_v=e_v(\phi_v)$ for $v\neq v_1$,
which is Schwartz-Bruhat and equals $\mathrm{vac}_{1,v}$ when
$\phi_v=\mathrm{vac}'_v$. So $e(P\Psi)\in\mathcal S^\circ(M_1)$.

\emph{Step 4: the theta integral.} By
Lemmas~\ref{asm-theta-agree-continuity-pieces}(1)
and~\ref{asm-theta-agree-continuity-domination}, every element of
$\mathcal S^\circ(M')$ converges (proof of
Lemma~\ref{asm-theta-agree-end-beta-zero-vanish}), and $I'$ is additive on the
functions that converge. Let $\mathcal K$ be the space of the converging
Schwartz-Bruhat functions $\Psi$ with $I'(\Psi)=0$. As $1-P$ is linear and
$\Phi'\in\mathcal S^\circ(M')$, to show $\Phi'-P\Phi'\in\mathcal K$ it suffices
to show $(1-P)\Psi\in\mathcal K$ for a spanning element
$\Psi=f\otimes\bigotimes_v\phi_v$ of $\mathcal S^\circ(M')$, with
$\phi_{v_1}\in S_1$. The place $v_1$ is clean for $\mathcal X$, so
$(\alpha/\beta)_t\neq1$ for every $F$-embedding $t\colon E\to F_{v_1}$
(Lemma~\ref{asm-theta-agree-coinvariants-nontrivial}). By
Lemma~\ref{asm-theta-agree-end-alpha-zero-lower}, for the two Weil pairs
$(V',W,M',\alpha,\alpha/\beta,\omega')$ and
$(V',W_1,M_1,\alpha,\alpha/\beta,\omega_1)$, with $a'\geq1$ and
$a+a'-2\geq1$, there is $\phi'\in S_1$ with $\phi_{v_1}-\phi'\in\mathfrak a$ and
$e_{v_1}(\phi')\in S_1^{(1)}$. Then $\phi'\in\mathcal G$, so
$\phi_{v_1}\in\mathcal G+\mathfrak a$ and
$\phi_{v_1}-\pi(\phi_{v_1})\in\mathfrak a$. By Step 2, $(1-P)\Psi$ is the tensor
with $\phi_{v_1}-\pi(\phi_{v_1})$ at $v_1$. This is linear in the component at
$v_1$, so it suffices to show that the tensor $\Psi_1-\Psi_0$ lies in
$\mathcal K$, where $\Psi_0$ and $\Psi_1$ are $\Psi$ with $\phi_0$ and with
$\phi_1$ at $v_1$, for $\phi_0\in S_1$, $\phi_1$ Schwartz-Bruhat and
$\omega'(g,1)$ acting at $v_1$ by $\phi_0\mapsto\phi_1$ for some
$g\in U_{V'}^{(v_1)}$. As $\Psi_0$ is $\phi_0$ at $v_1$ times a function that
does not depend on the component at $v_1$, $\Psi_1=\omega'(g,1)\Psi_0$. Now
$\Psi_0\in\mathcal S^\circ(M')$, and so is $\Psi_1$ by
Lemma~\ref{asm-theta-agree-split-inputs}(6). Both converge, and
$I'(\Psi_1)=I'(\Psi_0)$ by Lemma~\ref{asm-theta-agree-anisotropic-invariance}.
So $\Psi_1-\Psi_0\in\mathcal K$.

\emph{Step 5: conclusion.} Put $\Psi=P\Phi'$. By Step 3,
$\Psi\in\operatorname{PD}(M')\cap\mathcal S^\circ(M')$ and
$e(\Psi)\in\mathcal S^\circ(M_1)$. By Step 4, $\Phi'-\Psi\in\mathcal K$; as
$\Phi'$ and $\Psi$ converge, $I'(\Phi')=I'(\Psi)$.
\end{proof}
% Lean: soMove_of thetaBasic unimodularity_holds DomSoU SoFin TriplePair SoLower
% (SubMemKer.lean). Step 1 is exists_proj; Step 2 is opL (Ar.RT.Move.opL) with opL_pure and
% upd_so; Step 3 is the claim A (with the IsEvLoc clause pure); Step 4 is so_conv, kerS,
% sub_mem_ker, the span induction on hPhi'.2 with Inj.pure_at, esub, the key SoLower (with
% KCo.nu_statement), pure_update_mem, KCo.so_stable and I_left.

\begin{lemma}\label{asm-theta-agree-end-alpha-zero-second}
Let $(B,G,\mathcal X,\omega,\omega')$ be a triple of ranks $(a,a')$ for a theta
base $B$ of a CM field $E$ (Definition~\ref{asm-theta-agree-triple}), with
$p+q\geq1$, models $M$ and $M'$ and theta integral $I'$, let $H'$ be an adapted
plane in $W$ over a lattice space $W_1$ at some positions
(Definition~\ref{asm-theta-agree-planes}(1)), and let
$(V',W_1,M_1,\alpha,\alpha/\beta,\omega_1)$ be a Weil pair of signature
$(q,a'-q;a+a'-p-q-1,p+q-1)$ (Definition~\ref{asm-theta-agree-triple}(1)). Then
for every $\Phi'\in\operatorname{PD}(M')\cap\mathcal S^\circ(M')$ there are
functions $y_\alpha$ and $y_\beta$ that split $I'(\Phi')$ along $H'$
(Definition~\ref{asm-theta-agree-planes}(3)) and
$\Psi_1\in\operatorname{PD}(M_1)\cap\mathcal S^\circ(M_1)$ such that
$y_\alpha(l)=I_{M_1,\omega_1}(\Psi_1)(h_1)$ for every $h_1\in U_{W_1}$ and every
Levi element $l$ for $(1,h_1)$.
\end{lemma}
\leanref{Development/WeilBound/ThetaLifts/EigenvalueEndgame/EigenKilling/AlphaZeroOf.lean}{PositivityTheory.ThKEnd.Comp.small}
\begin{proof}
As $G$ is proper, $2(p+q)+1<a'$ and $2(p+q)+1<\lfloor(a-a')/2\rfloor$, so
$a'\geq4$ and $a-a'\geq2$; and $a+a'$ is even. The Weil pair
$(V',W,M',\alpha,\alpha/\beta,\omega')$ of the triple has signature
$(q,a'-q;a+a'-p-q,p+q)$. By Lemma~\ref{asm-theta-agree-evaluation} there is an
evaluation map $e_0$ for $(M',M_1,1,\iota_{H'},\mathcal J_0(H'))$, and it
carries $\operatorname{PD}(M')$ onto $\operatorname{PD}(M_1)$. By
Lemma~\ref{asm-theta-agree-end-alpha-zero-move} there is
$\Psi\in\operatorname{PD}(M')\cap\mathcal S^\circ(M')$ with $I'(\Psi)=I'(\Phi')$
and $e_0(\Psi)\in\mathcal S^\circ(M_1)$.

Apply Lemma~\ref{asm-theta-agree-split-any} in the case (small) to $V'$, $W$,
$W_1$, $H'$, $M'$, $M_1$ and $e_0$, with $n=a'\geq3$, $d=a+a'$,
$(r,s)=(q,a'-q)$, $(r',s')=(a+a'-p-q-1,p+q-1)$ and $z=(a-a')/2+1$, so that
$d+2=2n+2z$ and $z\geq2$. For the two Weil pairs above and $\Psi$ it gives
$C_0$ and functions $A_0$ and $A_1$ such that $A_0+A_1$ is the constant term of
$I'(\Psi)=I'(\Phi')$ along $H'$, $A_0$ has the law $(\alpha,a'/2)$, $A_1$ has the
law $(\alpha/(\alpha/\beta),(a+a'-a')/2)=(\beta,a/2)$, and
$A_0(lh)=C_0\,I_{M_1,\omega_1}(e_0(\omega'(1,h)\Psi))(h_1)$ for every Levi
element $l$ for $(1,h_1)$ and every $h$ with
$e_0(\omega'(1,h)\Psi)\in\mathcal S^\circ(M_1)$. So $(A_0,A_1)$ splits
$I'(\Phi')$ along $H'$. Put $\Psi_1=C_0\,e_0(\Psi)$, which lies in
$\operatorname{PD}(M_1)\cap\mathcal S^\circ(M_1)$. For $h=1$,
$\omega'(1,1)\Psi=\Psi$ and $e_0(\Psi)\in\mathcal S^\circ(M_1)$, so
$A_0(l)=C_0\,I_{M_1,\omega_1}(e_0(\Psi))(h_1)=I_{M_1,\omega_1}(\Psi_1)(h_1)$.
Take $(y_\alpha,y_\beta)=(A_0,A_1)$.
\end{proof}
% Lean: small SmallAny SoMove EvZero TriplePair (AlphaZeroOf.lean); SmallAny is
% Comp.smallAny_of = ThKCT.anySmall, Lemma asm-theta-agree-split-any (small); I_smul for the
% last equality.

\begin{lemma}\label{asm-theta-agree-end-alpha-zero}
Let $(B,G,\mathcal X,\omega,\omega')$ be a triple of ranks $(a,a')$ for a theta
base $B$ of a CM field $E$ (Definition~\ref{asm-theta-agree-triple}), with
$p=0$ and $q\geq1$, and let $H'$ be an adapted plane in $W$ over a lattice
space $W_1$ that is skew-Hermitian of some signature. If $f$ lies in the set of
differences of the triple (Definition~\ref{asm-theta-agree-end-notions}(5)) and
$(f_\alpha,f_\beta)$ splits $f$ along $H'$
(Definition~\ref{asm-theta-agree-planes}(3)), then $f_\alpha=0$.
\end{lemma}
\leanref{Development/WeilBound/ThetaLifts/EigenvalueEndgame/PlaneRemovalConfiguration.lean}{PositivityTheory.ThKEnd.AlphaZero}
\leanref{Development/WeilBound/ThetaLifts/EigenvalueEndgame/EigenKilling/AlphaZeroOf.lean}{PositivityTheory.ThKEnd.Comp.alphaZero_of}
\begin{proof}
Let $G=(0,q,V,V',W,M,M')$, $\mathcal X=(\alpha,\beta,S_0,Q)$, and write
$f=I(\Phi)-I'(\Phi')$ with $\Phi\in\operatorname{PD}(M)$ and
$\Phi'\in\operatorname{PD}(M')\cap\mathcal S^\circ(M')$. As $G$ is proper with
$p=0$: $V$ and $V'$ are Hermitian of signatures $(a,0)$ and $(q,a'-q)$, $W$ is
skew-Hermitian of signature $(a+a'-q,q)$, $a'<a$, $2q+1<a'$ and
$a+a'<m_\alpha(w)-m_\beta(w)$ for every $w$.

\emph{Step 1: the plane.} By Lemma~\ref{asm-theta-agree-split-low}, $W_1$ is
skew-Hermitian of signature $(a+a'-q-1,q-1)$. By
Lemma~\ref{asm-theta-agree-end-iwasawa}, $H'$ gives a plane in $W$ over $W_1$.
By Lemma~\ref{asm-theta-agree-end-automorphic}, $I(\Phi)$ and $I'(\Phi')$ are
automorphic, hence so is $f$. As $a'/2\neq a/2$,
Lemma~\ref{asm-theta-agree-end-constant-terms}(2), for the constant term
$f_\alpha+f_\beta$ of $f$, gives $f_\alpha(uh)=f_\alpha(h)$ for all
$u\in N_{H'}$ and $h\in U_W$.

\emph{Step 2: reduction to $K^\theta_W$.} It suffices to show $f_\alpha(lk)=0$
for every $h_1\in U_{W_1}$, every Levi element $l$ for $(1,h_1)$ and every
$k\in K^\theta_W$. Indeed, by Lemma~\ref{asm-theta-agree-end-iwasawa} every
$g\in U_W$ is $umlk$ with $u\in N_{H'}$, $m$ a Levi element for $(t,1)$, $l$ one
for $(1,h_1)$ and $k\in K^\theta_W$; by Step 1 and the law of $f_\alpha$,
$f_\alpha(g)=f_\alpha(mlk)=\alpha(t)|t|^{a'/2}f_\alpha(lk)$.

Fix $h_1$, $l$ and $k$. By Lemma~\ref{asm-theta-agree-end-alpha-zero-stable},
$\omega(1,k)\Phi\in\operatorname{PD}(M)$ and
$\omega'(1,k)\Phi'\in\operatorname{PD}(M')\cap\mathcal S^\circ(M')$.

\emph{Step 3: the first space.} By Lemma~\ref{asm-theta-agree-tower-model}
there is a regular model $M_1$ for $(V,W_1)$, and, as
$\beta|_{\mathbf A^\times}=\eta^a$ and
$(\beta/\alpha)|_{\mathbf A^\times}=\eta^{a+a'}=\eta^{a+a'-2}$,
Theorem~\ref{asm-theta-agree-tower-weil} gives a Weil representation
$\omega_1$ for $M_1$ with characters $(\beta,\beta/\alpha)$. So
$(V,W_1,M_1,\beta,\beta/\alpha,\omega_1)$ is a Weil pair of signature
$(a,0;a+a'-q-1,q-1)$. By Lemma~\ref{asm-theta-agree-evaluation} there is an
evaluation map $e_0$ for $(M,M_1,1,\iota_{H'},\mathcal J_0(H'))$. As
$\min(a,0)=0$, Lemma~\ref{asm-theta-agree-split-aniso}, for $V$, $W$, $W_1$,
$H'$, $M$, $M_1$ and $e_0$, applied to the Schwartz-Bruhat function
$\omega(1,k)\Phi$ (Lemma~\ref{asm-theta-agree-split-inputs}(2)), gives a
function $A_0$ that is the constant term of $I(\omega(1,k)\Phi)$ along $H'$ and
has the law $(\beta,a/2)$. As $0$ has every law, $(0,A_0)$ splits
$I(\omega(1,k)\Phi)$ along $H'$.

\emph{Step 4: the second space.} In the same way there are a regular model
$M_2$ for $(V',W_1)$ and a Weil representation $\omega_2$ for $M_2$ with
characters $(\alpha,\alpha/\beta)$, so that
$(V',W_1,M_2,\alpha,\alpha/\beta,\omega_2)$ is a Weil pair of signature
$(q,a'-q;a+a'-q-1,q-1)$. Lemma~\ref{asm-theta-agree-end-alpha-zero-second},
applied to $\omega'(1,k)\Phi'$, gives functions $y_\alpha$ and $y_\beta$ that
split $I'(\omega'(1,k)\Phi')$ along $H'$ and
$\Psi_1\in\operatorname{PD}(M_2)\cap\mathcal S^\circ(M_2)$ with
$y_\alpha(l')=I_{M_2,\omega_2}(\Psi_1)(h'_1)$ for every Levi element $l'$ for
$(1,h'_1)$.

\emph{Step 5: comparison.} As
$\theta_\omega(\omega(1,k)\Phi)(g,h)=\theta_\omega(\Phi)(g,hk)$,
$I(\omega(1,k)\Phi)(h)=I(\Phi)(hk)$, and likewise for $I'$. So
$h\mapsto f(hk)$ is $I(\omega(1,k)\Phi)-I'(\omega'(1,k)\Phi')$. Right
translation by $k$ preserves constant terms along $H'$ and laws, so
$h\mapsto f_\alpha(hk)$ and $h\mapsto f_\beta(hk)$ split it. The functions
$I(\omega(1,k)\Phi)$ and $I'(\omega'(1,k)\Phi')$ are automorphic
(Lemma~\ref{asm-theta-agree-end-automorphic}), so by
Lemma~\ref{asm-theta-agree-end-constant-terms}(1) the constant term of their
difference is $(0+A_0)-(y_\alpha+y_\beta)$, and, as laws are preserved by linear
combinations, $(-y_\alpha,A_0-y_\beta)$ splits it too. By
Lemma~\ref{asm-theta-agree-split-unique}, $f_\alpha(hk)=-y_\alpha(h)$ for every
$h$. In particular $f_\alpha(lk)=-I_{M_2,\omega_2}(\Psi_1)(h_1)$.

\emph{Step 6: vanishing.} The characters $\alpha$ and $\beta$ have exponents
$m_\alpha$ and $m_\beta$ (Definition~\ref{asm-theta-agree-triple}(3)), so
$\alpha/\beta$ has exponents $m_\alpha-m_\beta$. For the Weil pair
$(V',W_1,M_2,\alpha,\alpha/\beta,\omega_2)$: $a+a'-2<|m_\alpha(w)-m_\beta(w)|$
for every $w$; and $\min(a+a'-q-1,q-1)=q-1<\min(q,a'-q)$, as $2q+1<a'$. So
Lemma~\ref{asm-theta-agree-end-beta-zero-vanish} gives
$I_{M_2,\omega_2}(\Psi_1)=0$, as $\Psi_1\in\mathcal S^\circ(M_2)$. Hence
$f_\alpha(lk)=0$.
\end{proof}
% Lean: Comp.alphaZero_of Iwa CTBasic IAut UnipEx Low PairEx TriplePair EvZero GenAniso
% SmallAny SoMove PolStab IZero (EigenKilling/AlphaZeroOf.lean). Step 1 is Low, Iwa.1,
% sig_isAut, half_ne and CTBasic.2; Step 2 is comp_zero (TowerCodimensionTwo.lean) and PolStab
% (Lemma asm-theta-agree-end-alpha-zero-stable); Step 3 is PairEx, EvZero and GenAniso
% (polData_le); Step 4 is PairEx and small (Lemma asm-theta-agree-end-alpha-zero-second);
% Step 5 is split_right (I_trans, isCT_right, hasLaw_right) and comp_eq (CTBasic.1,
% hasLaw_sub, split_uniq); Step 6 is hasExp_div, geo.kappa and the second part of IZero.
% Lemma asm-theta-agree-end-alpha-zero-lower (SoLower, about 800 lines) is now written out.

\begin{lemma}\label{asm-theta-agree-end-beta-zero-forms}
Let $P,Q,R,S\geq0$ and $m\in\mathbf Z$ with $R+S<|m|$ and
$\min(R,S)<\min(P,Q)$, and use the notation of
Definition~\ref{asm-theta-agree-real-fock}(1). Let $\ell$ be a linear form on
$\mathcal P(P,Q;R,S)$ with
$\ell(d_1(Y)f)=-\tfrac12m\operatorname{tr}(Y)\,\ell(f)$ for every
block-diagonal $Y\in M_{P+Q}(\mathbf C)$ and every $f$. Then $\ell=0$.
\end{lemma}
\leanref{Development/WeilBound/ThetaLifts/EigenvalueEndgame/EigenKilling/FockActionFirstEq.lean}{PositivityTheory.ThKEnd.Comp.noForm}
\leanref{Development/WeilBound/ThetaLifts/EigenvalueEndgame/EigenKilling/FockActionFirstEq.lean}{PositivityTheory.ThKEnd.Comp.core_pos}
\leanref{Development/WeilBound/ThetaLifts/EigenvalueEndgame/EigenKilling/FockActionFirstEq.lean}{PositivityTheory.ThKEnd.Comp.fockAction_first_eq}
\begin{proof}
Write $\mathcal P=\mathcal P(P,Q;R,S)$, with variables $x_{(a,c)}$ for
$a\in\mathrm{Ix}(P,Q)$ and $c\in\mathrm{Ix}(R,S)$, and $\partial_{(a,c)}$ for
$\partial/\partial x_{(a,c)}$. For $u\in\mathbf C^{\mathrm{Ix}(P,Q)\times\mathrm{Ix}(R,S)}$
put $\lambda_u=\sum_lu_lx_l$. For a set $\Pi\subset\mathrm{Ix}(R,S)$ and
$v,w\in\mathbf C^{\mathrm{Ix}(P,Q)}$ let $D^\Pi_{v,w}$ be the derivation
\[
D^\Pi_{v,w}=\sum_{a,b}v_aw_b\Bigl(\sum_{c\in\Pi}x_{(a,c)}\partial_{(b,c)}-\sum_{c\notin\Pi}x_{(b,c)}\partial_{(a,c)}\Bigr)
\]
of $\mathcal P$. Exchanging $a$ and $b$ shows
$D^{\Pi'}_{v,w}=-D^\Pi_{w,v}$, where $\Pi'$ is the complement of $\Pi$.

\emph{Step 1: the derivations.} Let $\beta$ be a block of $\mathrm{Ix}(P,Q)$,
let $\Pi_\beta$ be the block of $\mathrm{Ix}(R,S)$ of the same kind (the first
if $\beta$ is the first), and let $v,w$ vanish outside $\beta$ with
$\sum_av_aw_a=1$. For $a,b\in\beta$ and $c\in\mathrm{Ix}(R,S)$,
$\omega_{(a,c),(b,c)}$ is $x_{(a,c)}\partial_{(b,c)}+\frac12\delta_{ab}$ if
$c\in\Pi_\beta$, as then $s_{(a,c)}=s_{(b,c)}=1$, and
$-\partial_{(a,c)}x_{(b,c)}+\frac12\delta_{ab}=-x_{(b,c)}\partial_{(a,c)}-\frac12\delta_{ab}$
otherwise. Summing with the weights $v_aw_b$ gives
$d_1(vw^{\mathsf T})=D^{\Pi_\beta}_{v,w}+\kappa_\beta$, with
$\kappa_\beta=\frac{R-S}2$ for the first block and $\frac{S-R}2$ for the
second. As $vw^{\mathsf T}$ is block-diagonal with trace $1$, the hypothesis
gives
\[
\ell(D^{\Pi_\beta}_{v,w}f)=k_\beta\,\ell(f),\qquad k_\beta=-\tfrac m2-\kappa_\beta .\tag{$*$}
\]

\emph{Step 2.} Let $\beta$ be a block of $\mathrm{Ix}(P,Q)$ and
$\Pi\subset\mathrm{Ix}(R,S)$ with $|\Pi|<|\beta|$, and let $k>0$. Suppose that
$\ell(D^\Pi_{v,w}f)=k\,\ell(f)$ for all $f$ and all $v,w$ vanishing outside
$\beta$ with $\sum_av_aw_a=1$. Then $\ell=0$. As the powers of linear forms
span $\mathcal P$, it suffices to show $\ell(\lambda_u^N)=0$ for every $u$ and
$N\geq0$. The linear map
$\mathbf C^\beta\to\mathbf C^\Pi$, $w\mapsto(\sum_{b\in\beta}w_bu_{(b,c)})_c$,
has a nonzero kernel. Let $w$ be a nonzero element of it, extended by $0$, let
$w_{a_0}\neq0$, and let $v=w_{a_0}^{-1}e_{a_0}$; then $\sum_av_aw_a=1$ and
$\sum_bw_bu_{(b,c)}=0$ for $c\in\Pi$. Let $u'_{(a,c)}=0$ for $c\in\Pi$ and
$u'_{(a,c)}=w_a\sum_{a'}v_{a'}u_{(a',c)}$ for $c\notin\Pi$. By the definition,
$D^\Pi_{v,w}\lambda_u=\lambda_{u''}$ with
$u''_{(a,c)}=v_a\sum_bw_bu_{(b,c)}$ for $c\in\Pi$ and
$u''_{(a,c)}=-w_a\sum_{a'}v_{a'}u_{(a',c)}$ for $c\notin\Pi$. So
$D^\Pi_{v,w}\lambda_u=-\lambda_{u'}$, and, as $\sum_av_aw_a=1$,
$D^\Pi_{v,w}\lambda_{u'}=-\lambda_{u'}$. Hence $D^\Pi_{v,w}$ kills
$\lambda_u-\lambda_{u'}$, and, being a derivation, it maps
$g_{ij}=(\lambda_u-\lambda_{u'})^i\lambda_{u'}^j$ to $-jg_{ij}$. The hypothesis
gives $(k+j)\ell(g_{ij})=0$, so $\ell(g_{ij})=0$, and
$\ell(\lambda_u^N)=\sum_i\binom Ni\ell(g_{i,N-i})=0$.

\emph{Step 3.} Let $\beta$ and $\Pi$ be as in Step 2 and let $k<0$. If
$\ell(D^{\Pi'}_{v,w}f)=k\,\ell(f)$ for all $f$ and all $v,w$ as there, then
$\ell(D^\Pi_{v,w}f)=-\ell(D^{\Pi'}_{w,v}f)=-k\,\ell(f)$, and Step 2 gives
$\ell=0$.

\emph{Step 4: conclusion.} Suppose first $m>R+S$. Then $k_\beta<-R$ for the
first block and $k_\beta<-S$ for the second, so both are negative. If
$R\leq S$, then $R<Q$; Step 3 for the second block $\beta$, with $\Pi$ the
first block of $\mathrm{Ix}(R,S)$, of size $R$, and ($*$) give $\ell=0$. If
$S\leq R$, then $S<P$, and Step 3 applies to the first block with $\Pi$ the
second block of $\mathrm{Ix}(R,S)$. Suppose now $m<-(R+S)$. Then $k_\beta>S$
for the first block and $k_\beta>R$ for the second, so both are positive. If
$R\leq S$, then $R<P$, and Step 2 applies to the first block with
$\Pi=\Pi_\beta$, of size $R$; if $S\leq R$, then $S<Q$, and Step 2 applies to
the second block with $\Pi=\Pi_\beta$.
\end{proof}
% Lean: noForm, with fockAction_first_eq (Step 1), core_pos (Step 2, with der_lin and
% eq_zero_of_pow for the powers of linear forms), core_neg and der_neg (Step 3).

\begin{lemma}\label{asm-theta-agree-end-beta-zero-real}
Let $(X,W,M,\chi_X,\chi_W,\omega)$ be a Weil pair of signature $(r,s;r',s')$
for a theta base $B$ of a CM field $E$ (Definition~\ref{asm-theta-agree-triple}(1)),
with $X$ and $W$ of ranks $n$ and $d$, and put
$I=\{1,\ldots,n\}\times\{1,\ldots,d\}$. Suppose that $\chi_W$ has exponents $m$
(Lemma~\ref{asm-theta-agree-real-lie}) with $d<|m(w)|$ for every infinite
place $w$ of $F$, and that $\min(r',s')<\min(r,s)$. Let $L$ be a continuous
linear form on $\mathcal S(F_\infty^I)$ with $L(T(f))=L(f)$ for every
$g\in U_X$ at $\infty$, every archimedean operator $T$ of $\omega$ at $(g,1)$
(Definition~\ref{asm-theta-agree-rechoice-terms}(1)) and every $f$. Then
$L=0$.
\end{lemma}
\leanref{Development/WeilBound/ThetaLifts/EigenvalueEndgame/EigenKilling/ConeIntegralFactorisation.lean}{PositivityTheory.ThKEnd.Comp.RealVanish}
\leanref{Development/WeilBound/ThetaLifts/EigenvalueEndgame/EigenKilling/RealVanishOf.lean}{PositivityTheory.ThKEnd.Comp.realVanish_of}
\begin{proof}
For every infinite place $w$, the matrices $\tau_w(J_X)$ and $i\,\tau_w(J_W)$
have signatures $(r,s)$ and $(r',s')$, with Cartan involutions $\theta_{X,w}$
and $\theta_{W,w}$ (Definition~\ref{asm-theta-agree-doubling}(1)); in
particular $r'+s'=d$. By Lemma~\ref{asm-theta-agree-real-frames} there are
frames $A_w$ of $\tau_w(J_X)$ of type $(r,s)$ and $C_w$ of $i\,\tau_w(J_W)$ of
type $(r',s')$ with $\theta_{X,w}A_w=A_wJ_{r,s}$ and
$\theta_{W,w}C_w=C_wJ_{r',s'}$. Let $\Phi$ be a Fock map for
$(M,(A_w),(C_w))$ (Lemma~\ref{asm-theta-agree-real-fock-map}(1), with $a_0=1$
and $b_0=i$), defined on $\mathcal T=\bigotimes_w\mathcal P(r,s;r',s')$; by
Lemma~\ref{asm-theta-agree-real-fock-map}(2) its image is
$\operatorname{Pol}(M)$.

Choose a Schwartz-Bruhat function $g_0$ on $\mathbf A_f^I$ with $g_0(0)=1$.
The linear form $l$ on functions on $\mathbf A^I$ that maps a Schwartz-Bruhat
function $\Psi$ to $L(\Psi(\cdot,0))$, the value of $L$ on the restriction to
finite component $0$, and is extended linearly to all functions, satisfies
$l(f\otimes g)=g(0)L(f)$; so it is admissible
(Definition~\ref{asm-theta-agree-doubling}(6)).

Fix $w$ and a family $(f_v)_v$ of elements of $\mathcal P(r,s;r',s')$. For
$f\in\mathcal P(r,s;r',s')$ let $T_f\in\mathcal T$ be the tensor of the family
with $f$ in place of $f_w$, and put $\ell(f)=L(\Phi(T_f))$. Let
$Y\in\mathfrak u(r,s)$. By Lemma~\ref{asm-theta-agree-real-lines}(2), the
family equal to $Y$ at $w$ and $0$ at the other places generates a line
$\gamma$ in $U_X$ relative to $(A_w)$. For each $t$, $\gamma(t)$ is at
$\infty$, so $\omega(\gamma(t),1)$ has an archimedean operator $T_t$
(Lemma~\ref{asm-theta-agree-rechoice-operators}(1)), and
\[
l\bigl(\omega(\gamma(t),1)(\Phi(T_f)\otimes g_0)\bigr)=l\bigl(T_t\Phi(T_f)\otimes g_0\bigr)=L(T_t\Phi(T_f))=L(\Phi(T_f))
\]
does not depend on $t$. The operators $\omega(g,1)$ satisfy condition (b) of
Definition~\ref{asm-theta-agree-models}(6) with the character $\chi_W$, which
has exponents $m$. So by Lemma~\ref{asm-theta-agree-real-lie}(2) the
derivative at $t=0$ is $L(\Phi(T'))$, where $T'=T_{f'}$ with
$f'=d_1(Y)f+\frac12m(w)\operatorname{tr}(Y)f$. Hence
$\ell(d_1(Y)f)=-\frac12m(w)\operatorname{tr}(Y)\,\ell(f)$ for
$Y\in\mathfrak u(r,s)$. Every $Y\in M_{r+s}(\mathbf C)$ is $Y_1+iY_2$ with
$Y_1=\frac12(Y-J_{r,s}Y^*J_{r,s})$ and $Y_2=-\frac i2(Y+J_{r,s}Y^*J_{r,s})$ in
$\mathfrak u(r,s)$, and $d_1$ is $\mathbf C$-linear
(Lemma~\ref{asm-theta-agree-real-fock-relations}(1)); so the identity holds
for every $Y$. As $r'+s'=d<|m(w)|$ and $\min(r',s')<\min(r,s)$,
Lemma~\ref{asm-theta-agree-end-beta-zero-forms} gives $\ell=0$.

So $L\circ\Phi$ vanishes on the pure tensors, hence on $\mathcal T$, and $L$
vanishes on $\operatorname{Pol}(M)$. This space is dense in
$\mathcal S(F_\infty^I)$ by Lemma~\ref{asm-theta-agree-density}, and $L$ is
continuous, so $L=0$.
\end{proof}
% Lean: realVanish_of ArchOp PolDense FrameEx FockEx LineEx PolRange LieFirst
% (RealVanishOf.lean); l and g_0 are ThKF.A.exists_adm and exists_sb_one, the decomposition of
% Y is upq_span, and sumPl_single isolates the place w.

\begin{lemma}\label{asm-theta-agree-end-beta-zero-vanish}
Let $(X,W,M,\chi_X,\chi_W,\omega)$ be a Weil pair of signature $(r,s;r',s')$
for a theta base $B$ of a CM field $E$ (Definition~\ref{asm-theta-agree-triple}(1)),
with $X$ and $W$ of ranks $n$ and $d$ and theta integral $I$, taking the
integral of a function that is not integrable to be $0$. Suppose that $\chi_W$
has exponents $m$ (Lemma~\ref{asm-theta-agree-real-lie}) with $d<|m(w)|$ for
every infinite place $w$ of $F$, and that $\min(r',s')<\min(r,s)$. Then
$I(\Phi)=0$ for every $\Phi\in\mathcal S^\circ(M)$ and, if
$d+2\min(r,s)<2n$, for every Schwartz-Bruhat function $\Phi$ on
$\mathbf A^{n\times d}$.
\end{lemma}
\leanref{Development/WeilBound/ThetaLifts/EigenvalueEndgame/EigenKilling/ConeIntegralFactorisation.lean}{PositivityTheory.ThKEnd.Comp.IZero}
\leanref{Development/WeilBound/ThetaLifts/EigenvalueEndgame/EigenKilling/SubMemKer.lean}{PositivityTheory.ThKEnd.Comp.iZero_of}
\begin{proof}
Let $\mu$ and $\mathcal D$ be the measure and the domain of $M$, and put
$I_0=\{1,\ldots,n\}\times\{1,\ldots,d\}$. Call a Schwartz-Bruhat function
$\phi$ on $\mathbf A_f^{I_0}$ \emph{dominated} if there are $k\geq0$, a
continuous seminorm $p$ on $\mathcal S(F_\infty^{I_0})$ and a measurable
$G\colon U_X\to[0,\infty]$ with $\int_{\mathcal D}G\,d\mu<\infty$ and
$|\theta_\omega(f\otimes\phi)(g,h)|\leq p(f)\operatorname{ht}(h)^kG(g)$ for
all $f$, $g$ and $h$. Every full-rank finite part for $M$ is dominated, by
Lemma~\ref{asm-theta-agree-continuity-domination}, and, if
$d+2\min(r,s)<2n$, every Schwartz-Bruhat function on $\mathbf A_f^{I_0}$ is,
by Lemma~\ref{asm-theta-agree-convergence-dominated}.

Let $\phi$ be dominated. Then every $f\otimes\phi$ converges, as
$g\mapsto\theta_\omega(f\otimes\phi)(g,h)$ is continuous
(Lemma~\ref{asm-theta-agree-convergence-theta}(2)) and bounded by an
integrable function. Fix $h\in U_W$ and put $L(f)=I(f\otimes\phi)(h)$. It is
linear, as $\theta_\omega$ is (Lemma~\ref{asm-theta-agree-convergence-theta}(2))
and the integrals converge, and
$|L(f)|\leq p(f)\operatorname{ht}(h)^k\int_{\mathcal D}G\,d\mu$, so $L$ is
continuous. Let $g\in U_X$ be at $\infty$ and $T$ an archimedean operator of
$\omega$ at $(g,1)$. Then $\omega(g,1)(f\otimes\phi)=T(f)\otimes\phi$, so
$L(T(f))=I(\omega(g,1)(f\otimes\phi))(h)=L(f)$ by
Lemma~\ref{asm-theta-agree-anisotropic-invariance}. By
Lemma~\ref{asm-theta-agree-end-beta-zero-real}, $L=0$. So
$I(f\otimes\phi)=0$ for every $f$.

Every $\Phi\in\mathcal S^\circ(M)$ is a finite sum $\sum_if_i\otimes\phi_i$
with full-rank finite parts $\phi_i$
(Lemma~\ref{asm-theta-agree-continuity-pieces}(1)), and every Schwartz-Bruhat
$\Phi$ is such a sum with Schwartz-Bruhat $\phi_i$. In the two cases of the
lemma the $\phi_i$ are dominated, so every $f_i\otimes\phi_i$ converges and
$I(\Phi)=\sum_iI(f_i\otimes\phi_i)=0$.
\end{proof}
% Lean: iZero_of thetaBasic unimodularity_holds DomWeil DomSoU SoFin RealVanish
% (SubMemKer.lean): dom_conv is the second paragraph (with I_left), kerS and I_add' the third.
% DomWeil and DomSoU are Lemmas asm-theta-agree-convergence-dominated and
% -continuity-domination. Lemma asm-theta-agree-end-alpha-zero uses the second part.

\begin{lemma}\label{asm-theta-agree-end-beta-zero}
Let $(B,G,\mathcal X,\omega,\omega')$ be a triple of ranks $(a,a')$ for a theta
base $B$ of a CM field $E$ (Definition~\ref{asm-theta-agree-triple}), with
$q=0$ and $p\geq1$, and let $H'$ be an adapted plane in $W$ over a lattice
space $W_1$. If $f$ lies in the set of differences of the triple
(Definition~\ref{asm-theta-agree-end-notions}(5)) and $(f_\alpha,f_\beta)$ splits
$f$ along $H'$ (Definition~\ref{asm-theta-agree-planes}(3)), then $f_\beta=0$.
\end{lemma}
\leanref{Development/WeilBound/ThetaLifts/EigenvalueEndgame/PlaneRemovalConfiguration.lean}{PositivityTheory.ThKEnd.BetaZero}
\leanref{Development/WeilBound/ThetaLifts/EigenvalueEndgame/EigenKilling/AlphaZeroOf.lean}{PositivityTheory.ThKEnd.Comp.betaZero_of}
\begin{proof}
Let $G=(p,0,V,V',W,M,M')$ and $\mathcal X=(\alpha,\beta,S_0,Q)$, and write
$f=I(\Phi)-I'(\Phi')$ with $\Phi\in\operatorname{PD}(M)$ and
$\Phi'\in\operatorname{PD}(M')\cap\mathcal S^\circ(M')$. As $G$ is proper with
$q=0$: $V$ and $V'$ are Hermitian of signatures $(a-p,p)$ and $(0,a')$, $W$ is
skew-Hermitian of signature $(a+a'-p,p)$, $2p+1<\lfloor(a-a')/2\rfloor$,
$2p+1<a'$ and $a+a'<m_\alpha(w)-m_\beta(w)$ for every $w$. In particular
$4p+2<a-a'$, so $a-p>p$, and $a>a'>2p+1\geq3$.

\emph{Step 1: the smaller spaces.} By Lemma~\ref{asm-theta-agree-split-low},
$W_1$ is skew-Hermitian of signature $(a+a'-p-1,p-1)$. By
Lemma~\ref{asm-theta-agree-tower-model} there are regular models $M_1$ for
$(V,W_1)$ and $M_2$ for $(V',W_1)$. As $\eta^2=1$,
$(\beta/\alpha)|_{\mathbf A^\times}=\eta^{a+a'}=\eta^{a+a'-2}$, and likewise
for $\alpha/\beta$; so by Theorem~\ref{asm-theta-agree-tower-weil} there are
Weil representations $\omega_1$ for $M_1$ with characters $(\beta,\beta/\alpha)$
and $\omega_2$ for $M_2$ with characters $(\alpha,\alpha/\beta)$. Thus
$(V,W_1,M_1,\beta,\beta/\alpha,\omega_1)$ and
$(V',W_1,M_2,\alpha,\alpha/\beta,\omega_2)$ are Weil pairs of signatures
$(a-p,p;a+a'-p-1,p-1)$ and $(0,a';a+a'-p-1,p-1)$. By
Lemma~\ref{asm-theta-agree-evaluation} there are evaluation maps $e_0$ for
$(M,M_1,1,\iota_{H'},\mathcal J_0(H'))$ and $e_2$ for
$(M',M_2,1,\iota_{H'},\mathcal J_0(H'))$.

\emph{Step 2: the first space.} Apply Lemma~\ref{asm-theta-agree-split-any} in
the case (big) to $V$, $W$, $W_1$, $H'$, $M$, $M_1$ and $e_0$, with $n=a\geq3$,
$d=a+a'$, $(r,s)=(a-p,p)$, $(r',s')=(a+a'-p-1,p-1)$ and $z=(a-a')/2$, an
integer as $a+a'$ is even: $2n=d+2z$ and $\min(r,s)+1=p+1\leq z$. It gives
$C_0\neq0$ and functions $A_0$ and $A_1$ such that $A_0+A_1$ is the constant
term of $I(\Phi)$ along $H'$, $A_0$ has the law $(\beta,a/2)$, $A_1$ has the
law $(\beta/(\beta/\alpha),a'/2)=(\alpha,a'/2)$, and
$A_0(lh)=C_0\,I_{M_1,\omega_1}(e_0(\omega(1,h)\Phi))(h_1)$ for every Levi
element $l$ for $(1,h_1)$ and every $h$. So $(A_1,A_0)$ splits $I(\Phi)$
along $H'$.

\emph{Step 3: the second space.} As $\min(0,a')=0$,
Lemma~\ref{asm-theta-agree-split-aniso}, applied to $V'$, $W$, $W_1$, $H'$,
$M'$, $M_2$ and $e_2$, gives a function $A$ that is the constant term of
$I'(\Phi')$ along $H'$ and has the law $(\alpha,a'/2)$. As $0$ has every law,
$(A,0)$ splits $I'(\Phi')$ along $H'$.

\emph{Step 4: comparison.} By Lemma~\ref{asm-theta-agree-end-automorphic},
$I(\Phi)$ and $I'(\Phi')$ are automorphic, and by
Lemma~\ref{asm-theta-agree-end-iwasawa} the adapted plane $H'$ gives a plane.
By Lemma~\ref{asm-theta-agree-end-constant-terms}(1), $(A_0+A_1)-A$ is the
constant term of $f$ along $H'$. The function $A_1-A$ has the law
$(\alpha,a'/2)$, as laws are preserved by linear combinations, and $A_0$ has
the law $(\beta,a/2)$. So $(A_1-A,A_0)$ splits $f$ along $H'$, and
Lemma~\ref{asm-theta-agree-split-unique} gives $f_\beta=A_0$.

\emph{Step 5: vanishing.} The identity of $U_W$ is a Levi element for
$(1,1)$, so $A_0(h)=C_0\,I_{M_1,\omega_1}(e_0(\omega(1,h)\Phi))(1)$ for
$h\in U_W$, and $e_0(\omega(1,h)\Phi)$ is a Schwartz-Bruhat function, as
$\Phi$ is one (proof of Lemma~\ref{asm-theta-agree-end-automorphic}) and
$\omega(1,h)$ and $e_0$ preserve these functions. The character $\beta/\alpha$
has exponents $m_\beta-m_\alpha$, and for the Weil pair
$(V,W_1,M_1,\beta,\beta/\alpha,\omega_1)$ of Step 1: $a+a'-2<|m_\beta(w)-m_\alpha(w)|$
for every $w$; $\min(a+a'-p-1,p-1)=p-1<p=\min(a-p,p)$; and
$(a+a'-2)+2p<2a$. So Lemma~\ref{asm-theta-agree-end-beta-zero-vanish} gives
$I_{M_1,\omega_1}(\Psi)=0$ for every Schwartz-Bruhat $\Psi$. Hence $A_0=0$, and
$f_\beta=0$.
\end{proof}
% Lean: Comp.betaZero_of Iwa CTBasic IAut UnipEx Low PairEx TriplePair EvZero GenAniso BigAny
% IZero (EigenKilling/AlphaZeroOf.lean). Step 1 is Low, PairEx (ModelEx and WeilEx) and EvZero;
% Step 2 is BigAny = Comp.bigAny_of = ThKCT.anyBig (Lemma asm-theta-agree-split-any (big));
% Step 3 is GenAniso; Step 4 is comp_eq (CTBasic, hasLaw_sub and split_uniq); Step 5 is
% isLevi_one, hasExp_div and the first part of IZero.

\begin{lemma}\label{asm-theta-agree-end-alpha-low-cone}
Let $B$ be a theta base of a CM field $E$, let $(G,G_1,H,H')$ be a
plane-removal configuration (Definition~\ref{asm-theta-agree-end-notions}(7)),
with $G$ of ranks $(b+2,a')$, first space $V$ and models $M$ and $M'$, and $G_1$
with models $M_1$ and $M'_1$, and let $e$ be an evaluation map for
$(M,M_1,\iota_H,\iota_{H'},\mathcal J_\alpha(H,H'))$
(Definition~\ref{asm-theta-agree-planes}(5)). Let $K\subset U_V$ be the lattice
stabilizer of $V$ (Definition~\ref{asm-theta-agree-planes}(6)), and $\mu_K$ and
$\mu$ Haar measures on $K$ and $\mathbf A_E^\times$. Let
$(B,G,\mathcal X,\omega,\omega')$ and $(B,G_1,\mathcal X,\omega_1,\omega'_1)$ be
triples (Definition~\ref{asm-theta-agree-triple}), put $\nu=\alpha/\beta$ and
$k_\alpha=\lfloor(b+2-a')/2\rfloor$, and let $\operatorname{cone}$ be the cone
for $(M,\omega,e,K,\mu_K,\mu,u_\alpha,\nu,k_\alpha)$
(Definition~\ref{asm-theta-agree-planes}(7)). Then
$\operatorname{cone}(\Phi)\in\operatorname{PD}(M_1)$ for every
$\Phi\in\operatorname{PD}(M)$.
\end{lemma}
\leanref{Development/WeilBound/ThetaLifts/EigenvalueEndgame/EigenKilling/ConeIntegralFactorisation.lean}{PositivityTheory.ThKEnd.Comp.ConePol}
\leanref{Development/WeilBound/ThetaLifts/EigenvalueEndgame/EigenKilling/SubMemKer.lean}{PositivityTheory.ThKEnd.Comp.conePol_of}
\leanref{Development/WeilBound/ThetaLifts/EigenvalueEndgame/EigenKilling/ConeIntegralFactorisation.lean}{PositivityTheory.ThKEnd.Comp.ConeIntegralFactorisation}
\leanref{Development/WeilBound/ThetaLifts/LocalTheory/ConeShifts/ConeFactorisationOfCodimensionTwo.lean}{PositivityTheory.ThKG.Loc.Cone.coneIntegralFactorisation_of}
\begin{proof}
By Lemma~\ref{asm-theta-agree-local-measures}, $K$ has a normalized local
measure system $m$. Lemma~\ref{asm-theta-agree-local-factor}(1) holds for every
plane-removal configuration, so it applies to $(G,G_1,H,H')$, $e$, $K$,
$\mu_K$, $\mu$ and $m$: there is $\kappa\neq0$, and for the two triples there
are $L$, $\Lambda_v$ and $u_v$ with the properties (a) to (e) of that lemma, for
the cone above.

By (d) with $\Phi=\Psi=0$ and $c=-1$, $\operatorname{cone}(0)=0$, so
$\operatorname{cone}$ is linear on the Schwartz-Bruhat functions, and every
element of $\operatorname{PD}(M)$ is Schwartz-Bruhat
(Lemma~\ref{asm-theta-agree-split-inputs}(2)). So it suffices to treat a
spanning element $\Phi=f\otimes\bigotimes_v\phi_v$ of $\operatorname{PD}(M)$,
with $f\in\operatorname{Pol}(M)$, every $\phi_v$ Schwartz-Bruhat and
$\phi_v=\mathrm{vac}_v$ outside a finite set. Let $T'$ be the union of $T$ and
the set of $v$ with $\phi_v\neq\mathrm{vac}_v$. By (e),
\[
\Bigl(\prod_{v\in T'\setminus T}u_v\Bigr)\operatorname{cone}(\Phi)=\lambda_0\,L(f)\otimes\bigotimes_v\psi_v,
\]
with $\psi_v=\Lambda_v(\phi_v)$ for $v\in T'$ and $\psi_v=\mathrm{vac}_{1,v}$
for $v\notin T'$. The product is nonzero by (c), $L(f)\in\operatorname{Pol}(M_1)$
by (a), and every $\psi_v$ is Schwartz-Bruhat, by (b) and as $M_1$ is regular.
So $\operatorname{cone}(\Phi)\in\operatorname{PD}(M_1)$.
\end{proof}
% Lean: conePol_of MeasEx ConeIntegralFactorisation (SubMemKer.lean): the span induction on
% Phi with z (cone 0 = 0), hlin and hT'. ConeIntegralFactorisation is the form of
% Lemma asm-theta-agree-local-factor(1) for plane-removal configurations, proved by
% Cone.coneIntegralFactorisation_of; ConeFacA, the form for alpha-configurations, is read off
% it by coneFactorisation_of_codimension_two.

\begin{lemma}\label{asm-theta-agree-end-alpha-low}
Let $B$ be a theta base of a CM field $E$, and let
$(B,G,\mathcal X,\omega,\omega')$ and $(B,G_1,\mathcal X,\omega_1,\omega'_1)$ be
triples (Definition~\ref{asm-theta-agree-triple}), with $G$ of ranks $(b+2,a')$
and $G_1$ of ranks $(b,a')$, such that $(G,G_1,H,H')$ is a plane-removal
configuration (Definition~\ref{asm-theta-agree-end-notions}(7)) and $q\geq1$.
Let $\Sigma$ and $\Sigma_1$ be the sets of differences of the two triples. If
$f\in\Sigma$, $(f_\alpha,f_\beta)$ splits $f$ along $H'$
(Definition~\ref{asm-theta-agree-planes}(3)) and $k\in K^\theta_W$, then some
$f_1\in\Sigma_1$ is the lower part of $f_\alpha$ at $k$
(Definition~\ref{asm-theta-agree-end-notions}(6)).
\end{lemma}
\leanref{Development/WeilBound/ThetaLifts/EigenvalueEndgame/PlaneRemovalConfiguration.lean}{PositivityTheory.ThKEnd.AlphaLow}
\leanref{Development/WeilBound/ThetaLifts/EigenvalueEndgame/EigenKilling/AlphaZeroOf.lean}{PositivityTheory.ThKEnd.Comp.alphaLow_of}
\begin{proof}
Let $G=(p,q,V,V',W,M,M')$ and $\mathcal X=(\alpha,\beta,S_0,Q)$. As
$(G,G_1,H,H')$ is a plane-removal configuration,
$G_1=(p_1,q,V_1,V',W_1,M_1,M'_1)$ with $p_1+1=p$, $H$ is an adapted plane in
$V$ over $V_1$ and $H'$ one in $W$ over $W_1$, both at positions $(0,b+1)$.
Write $f=I(\Phi)-I'(\Phi')$ with $\Phi\in\operatorname{PD}(M)$ and
$\Phi'\in\operatorname{PD}(M')\cap\mathcal S^\circ(M')$, and let $I_1$ and
$I'_1$ be the theta integrals of $G_1$. As $G$ and $G_1$ are proper: $V$, $W$,
$V_1$ and $W_1$ have signatures $(b+2-p,p)$, $(b+2+a'-p-q,p+q)$, $(b-p_1,p_1)$
and $(b+a'-p_1-q,p_1+q)$; $b+2+a'$ is even; and
$2(p+q)+1<k_\alpha=\lfloor(b+2-a')/2\rfloor$.

\emph{Step 1: stability.} By Lemma~\ref{asm-theta-agree-end-alpha-zero-stable},
$\omega(1,k)\Phi\in\operatorname{PD}(M)$ and
$\omega'(1,k)\Phi'\in\operatorname{PD}(M')\cap\mathcal S^\circ(M')$.

\emph{Step 2: the cone.} By Lemma~\ref{asm-theta-agree-evaluation-cone}, for the
Weil pair $(V,W,M,\beta,\beta/\alpha,\omega)$, $H$, $H'$ and the regular model
$M_1$, there is an evaluation map $e$ for
$(M,M_1,\iota_H,\iota_{H'},\mathcal J_\alpha(H,H'))$. By
Lemma~\ref{asm-theta-agree-stabilizer} there are Haar measures $\mu_K$ on the
lattice stabilizer $K$ of $V$ and $\mu$ on $\mathbf A_E^\times$. Let
$\operatorname{cone}$ be the cone for
$(M,\omega,e,K,\mu_K,\mu,u_\alpha,\alpha/\beta,k_\alpha)$. By
Lemma~\ref{asm-theta-agree-end-alpha-low-cone},
$\operatorname{cone}(\omega(1,k)\Phi)\in\operatorname{PD}(M_1)$.

\emph{Step 3: the first space.} By Lemma~\ref{asm-theta-agree-tower-model} there
is a regular model $M_2$ for $(V,W_1)$, and, as
$\beta|_{\mathbf A^\times}=\eta^{b+2}$ and
$(\beta/\alpha)|_{\mathbf A^\times}=\eta^{b+2+a'}=\eta^{b+a'}$,
Theorem~\ref{asm-theta-agree-tower-weil} gives a Weil representation $\omega_2$
for $M_2$ with characters $(\beta,\beta/\alpha)$. By
Lemma~\ref{asm-theta-agree-evaluation} there is an evaluation map $e_0$ for
$(M,M_2,1,\iota_{H'},\mathcal J_0(H'))$. The vector $x_H$ is nonzero, as
$h(x_H,y_H)=1$, and isotropic, so $K$ is an Iwasawa subgroup for $x_H$
(Lemma~\ref{asm-theta-agree-split-inputs}(3)), and $\omega(1,k)\Phi$ is
Schwartz-Bruhat and $K$-finite for $\omega$
(Lemma~\ref{asm-theta-agree-split-inputs}(2)). Apply
Lemma~\ref{asm-theta-agree-split-big} to $H$, $H'$, $M$, $M_0=M_1$, $M_2$, $e_0$
and $e$, with $n=b+2$, $d=b+2+a'$, $(r,s)=(b-p_1,p_1)$,
$(r',s')=(b+a'-p_1-q,p_1+q)$, $z=k_\alpha$, $K$, $\mu_K$ and $\mu$: here
$2n=d+2z$, as $b+2-a'$ is even, and $\min(r,s)+2\leq p_1+2=p+1<z$. The three
Weil pairs $(V,W,M,\beta,\beta/\alpha,\omega)$,
$(V,W_1,M_2,\beta,\beta/\alpha,\omega_2)$ and
$(V_1,W_1,M_1,\beta,\beta/\alpha,\omega_1)$ have the signatures
$(r+1,s+1;r'+1,s'+1)$, $(r+1,s+1;r',s')$ and $(r,s;r',s')$ required there. The
lemma gives $C_1\neq0$ and functions $A_0$ and $A_1$ such that $A_0+A_1$ is the
constant term of $I(\omega(1,k)\Phi)$ along $H'$, $A_0$ has the law
$(\beta,(b+2)/2)$, $A_1$ has the law
$(\beta/(\beta/\alpha),(d-n)/2)=(\alpha,a'/2)$, and
$A_1(lh)=C_1\,I_1(\operatorname{cone}(\omega(1,h)\omega(1,k)\Phi))(h_1)$ for every
Levi element $l$ for $(1,h_1)$ and every $h$. So $(A_1,A_0)$ splits
$I(\omega(1,k)\Phi)$ along $H'$.

\emph{Step 4: the second space.} The Weil pair
$(V',W_1,M'_1,\alpha,\alpha/\beta,\omega'_1)$ of the triple of $G_1$ has
signature $(q,a'-q;b+a'-p_1-q,p_1+q)$, which is
$(q,a'-q;(b+2+a')-p-q-1,p+q-1)$. Lemma~\ref{asm-theta-agree-end-alpha-zero-second},
for the triple of $G$, $H'$ and this Weil pair, applied to
$\omega'(1,k)\Phi'$, gives functions $y_\alpha$ and $y_\beta$ that split
$I'(\omega'(1,k)\Phi')$ along $H'$ and
$\Psi_1\in\operatorname{PD}(M'_1)\cap\mathcal S^\circ(M'_1)$ with
$y_\alpha(l)=I'_1(\Psi_1)(h_1)$ for every Levi element $l$ for $(1,h_1)$.

\emph{Step 5: comparison.} As in Step 5 of the proof of
Lemma~\ref{asm-theta-agree-end-alpha-zero}, $h\mapsto f_\alpha(hk)$ and
$h\mapsto f_\beta(hk)$ split
$h\mapsto f(hk)=I(\omega(1,k)\Phi)(h)-I'(\omega'(1,k)\Phi')(h)$ along $H'$, and
so does $(A_1-y_\alpha,A_0-y_\beta)$, by
Lemmas~\ref{asm-theta-agree-end-automorphic}
and~\ref{asm-theta-agree-end-constant-terms}(1). By
Lemma~\ref{asm-theta-agree-split-unique}, $f_\alpha(hk)=A_1(h)-y_\alpha(h)$ for
every $h$.

\emph{Step 6: the lower part.} Put
$f_1=I_1(C_1\operatorname{cone}(\omega(1,k)\Phi))-I'_1(\Psi_1)$. By Steps 2
and 4, $C_1\operatorname{cone}(\omega(1,k)\Phi)\in\operatorname{PD}(M_1)$ and
$\Psi_1\in\operatorname{PD}(M'_1)\cap\mathcal S^\circ(M'_1)$, so
$f_1\in\Sigma_1$. Let $h_1\in U_{W_1}$ and $l$ be a Levi element for $(1,h_1)$.
By Steps 3 to 5, with $h=1$ and $\omega(1,1)=1$,
\[
f_\alpha(lk)=A_1(l)-y_\alpha(l)=C_1\,I_1(\operatorname{cone}(\omega(1,k)\Phi))(h_1)-I'_1(\Psi_1)(h_1)=f_1(h_1).
\]
So $f_1$ is the lower part of $f_\alpha$ at $k$.
\end{proof}
% Lean: Comp.alphaLow_of Iwa CTBasic IAut UnipEx PairEx TriplePair EvZero EvCone StabEx IwaEx
% KFinData GenBig SmallAny SoMove PolStab ConePol (EigenKilling/AlphaZeroOf.lean). Step 1 is
% PolStab; Step 2 is EvCone, StabEx and ConePol (Lemma asm-theta-agree-end-alpha-low-cone);
% Step 3 is PairEx, EvZero, IwaEx, KFinData and GenBig; Step 4 is small (Lemma
% asm-theta-agree-end-alpha-zero-second, through SmallAny and SoMove); Step 5 is split_right
% and comp_eq; Step 6 is c4 at h = 1 with inv_div and ThKG.I_smul. Lemma
% asm-theta-agree-end-alpha-zero-lower (SoLower), used through Lemma
% asm-theta-agree-end-alpha-zero-second, is now written out.

\subsection{Purity of lower parts}\label{subsection-asm-purity}

\begin{lemma}\label{asm-theta-agree-end-two-step-moderate}
Let $(B,G,\mathcal X,\omega,\omega')$ be a triple for a theta base $B$ of a CM
field $E$ (Definition~\ref{asm-theta-agree-triple}), and let $\operatorname{ht}$
be the height of Definition~\ref{asm-theta-agree-parts}. Then for every $f$ in
the set of differences of the triple
(Definition~\ref{asm-theta-agree-end-notions}(5)) there are $C$ and $k$ with
$|f(h)|\leq C\operatorname{ht}(h)^k$ for every $h\in U_W$.
\end{lemma}
\leanref{Development/WeilBound/ThetaLifts/EigenvalueEndgame/SignatureSplitting/SigSplit.lean}{PositivityTheory.ThKEnd.Two.SigMod}
\leanref{Development/WeilBound/ThetaLifts/EigenvalueEndgame/SignatureSplitting/SigSplitOf.lean}{PositivityTheory.ThKEnd.Two.sigMod_of}
\begin{proof}
Let the triple have ranks $(a,a')$, integers $p,q$ and models $M$ and $M'$, and
write $f=I(\Phi)-I'(\Phi')$ with $\Phi\in\operatorname{PD}(M)$ and
$\Phi'\in\operatorname{PD}(M')\cap\mathcal S^\circ(M')$. The Weil pair
$(V,W,M,\beta,\beta/\alpha,\omega)$ has signature
$(a-p,p;a+a'-p-q,p+q)$, and $a+a'+2\min(a-p,p)\leq a+a'+2p<2a$, as
$2(p+q)+1<\lfloor(a-a')/2\rfloor$. The function $\Phi$ is Schwartz-Bruhat
(Lemma~\ref{asm-theta-agree-split-inputs}(2)), so by
Lemma~\ref{asm-theta-agree-convergence}(1) there are $C_1$ and $k_1$ with
$|I(\Phi)(h)|\leq\int_{\mathcal D}|\theta_\omega(\Phi)(g,h)|\,d\mu(g)\leq C_1\operatorname{ht}(h)^{k_1}$.
By Lemma~\ref{asm-theta-agree-decay}, there are $C_2$ and $k_2$ with
$|I'(\Phi')(h)|\leq C_2\operatorname{ht}(h)^{k_2}$. As
$\operatorname{ht}(h)\geq1$, $C=\max(C_1,0)+\max(C_2,0)$ and
$k=\max(k_1,k_2)$ will do.
\end{proof}
% Lean: sigMod_of WeilConv KDecay (SigSplitOf.lean), with pairV and weilCondV for the first
% part of WeilConv, the third clause of Decay for Phi', and one_le_heightGL.

\begin{lemma}\label{asm-theta-agree-end-two-step-split}
Let $(B,G,\mathcal X,\omega,\omega')$ be a triple of ranks $(a,a')$ for a theta
base $B$ of a CM field $E$ (Definition~\ref{asm-theta-agree-triple}), with
$p+q\geq1$, and let $\Sigma$ be its set of differences
(Definition~\ref{asm-theta-agree-end-notions}(5)). Then there are a lattice
space $W_1$ and a plane $H$ in $W$ over $W_1$
(Definition~\ref{asm-theta-agree-end-notions}(1)) such that every $f\in\Sigma$ is
split along $H$: there are functions $f_\alpha$ and $f_\beta$ on $U_W$ such that
$f_\alpha+f_\beta$ is the constant term of $f$ along $H$, $f_\alpha$ has the law
$(\alpha,a'/2)$ and $f_\beta$ has the law $(\beta,a/2)$.
\end{lemma}
\leanref{Development/WeilBound/ThetaLifts/EigenvalueEndgame/SignatureSplitting/SigSplit.lean}{PositivityTheory.ThKEnd.Two.SigSplit}
\leanref{Development/WeilBound/ThetaLifts/EigenvalueEndgame/SignatureSplitting/SigSplitOf.lean}{PositivityTheory.ThKEnd.Two.sigSplit_of}
\leanref{Development/WeilBound/ThetaLifts/EigenvalueEndgame/SignatureSplitting/SigSplitOf.lean}{PositivityTheory.ThKEnd.Two.split_sub}
\begin{proof}
Let $G=(p,q,V,V',W,M,M')$. As $G$ is proper, $a'<a$, and $W$ is
skew-Hermitian of rank $a+a'\geq3$ and signature $(a+a'-p-q,p+q)$, with
$a+a'-p-q\geq1$ and $p+q\geq1$. By
Proposition~\ref{asm-theta-agree-tower-plane} for $W$, with $n_1=a+a'-2$ and the
positions $0<1$, there are lattice spaces $Y$ and $W_1$ such that $W[Y]$ is
skew-Hermitian of signature $(a+a'-p-q,p+q)$ and there is an adapted plane $H$
in $W[Y]$ over $W_1$ at positions $0<1$. The datum $G^\circ=G[V,V',Y]$ has the
integers, the ranks and the models of $G$, and its spaces have the matrices,
the lattices and the signatures of those of $G$; so, as in the proof of
Lemma~\ref{asm-theta-agree-end-tower}, $(B,G^\circ,\mathcal X,\omega,\omega')$
is a triple. Its set of differences $\Sigma^\circ$ consists of functions on
$U_{W[Y]}=U_W$. By Lemma~\ref{asm-theta-agree-end-iwasawa}, $H$ gives a plane
in $W[Y]$ over $W_1$, which is a plane in $W$ over $W_1$, as $W$ and $W[Y]$ have
the same matrix.

By Lemma~\ref{asm-theta-agree-end-rechoose} there is $h_0\in U_W$ such that
$x\mapsto f(xh_0)$ lies in $\Sigma^\circ$ for every $f\in\Sigma$. Let
$f\in\Sigma$, and write $f(xh_0)=I(\Phi)(x)-I'(\Phi')(x)$ with $\Phi$ and
$\Phi'$ polynomial data for $G^\circ$, $\Phi'$ of full rank. By
Lemma~\ref{asm-theta-agree-split} for $G^\circ$ and $H$ there are pairs
$(g_\alpha,g_\beta)$ and $(g'_\alpha,g'_\beta)$ that split $I(\Phi)$ and
$I'(\Phi')$ along $H$. By Lemma~\ref{asm-theta-agree-end-automorphic}, $I(\Phi)$
and $I'(\Phi')$ are automorphic, so by
Lemma~\ref{asm-theta-agree-end-constant-terms}(1) the constant term of
$I(\Phi)-I'(\Phi')$ along $H$ is
$(g_\alpha-g'_\alpha)+(g_\beta-g'_\beta)$; and $g_\alpha-g'_\alpha$ and
$g_\beta-g'_\beta$ have the laws $(\alpha,a'/2)$ and $(\beta,a/2)$. Right
translation by $h_0^{-1}$ preserves constant terms along $H$ and laws, and
$f(x)=(I(\Phi)-I'(\Phi'))(xh_0^{-1})$. So
$f_\alpha(x)=(g_\alpha-g'_\alpha)(xh_0^{-1})$ and
$f_\beta(x)=(g_\beta-g'_\beta)(xh_0^{-1})$ split $f$ along $H$.
\end{proof}
% Lean: sigSplit_of SplitEx PlaneSkew Iwa CTBasic IAut SigRe (SigSplitOf.lean): PlaneSkew
% for G.W at positions 0 < 1, the triple t' for reG G G.V G.V' Y, Iwa.1, SigRe, SplitEx and
% split_sub (CTBasic.1, isCT_right, hasLaw_sub, hasLaw_right) at y = h_0^{-1}.

\begin{definition}\label{asm-theta-agree-end-kill-notions}
Let $G$ be a topological group.
\begin{enumerate}
\item A subgroup $\Gamma\subset G$ is a \emph{lattice} in $G$ if it is
  countable and discrete for the subspace topology and there is a compact
  $C\subset G$ such that every $g\in G$ has some $\gamma\in\Gamma$ with
  $\gamma g\in C$.
\item Let $G$ be locally compact, $\Gamma\subset G$ a subgroup,
  $\iota\colon G\to H$ a map to a group $H$ and $f$ a function on $H$. A
  function $c$ on $H$ is the \emph{average} of $f$ over $\Gamma$ along $\iota$
  if $c(h)=\int_{\mathcal F}f(\iota(u)h)\,d\mu(u)$ for every $h\in H$, every Haar
  measure $\mu$ on $G$ and every fundamental domain $\mathcal F$ for $\Gamma$ in
  $G$ with respect to $\mu$ with $\mu(\mathcal F)=1$. When $G$ is a subgroup of
  $H$ and $\iota$ is the inclusion, $c$ is the \emph{average} of $f$ over
  $\Gamma$. The constant term of $f$ along a plane $H'$ in a lattice space $W$
  is its average over $N_{H'}\cap U_W(F)$
  (Definition~\ref{asm-theta-agree-planes}(2)).
\item Let $a\in G$ and let $m\geq0$ be an integer. A function $\varphi$ on $G$
  has \emph{growth of order $m$ along $a$} if for every compact $\Omega\subset G$
  there is $C$ with $|\varphi(a^k\kappa)|\leq C\,2^{km}$ for every integer
  $k\geq0$ and every $\kappa\in\Omega$. A \emph{bump} in a set $V\subset G$ is a
  continuous function $\eta\colon G\to[0,\infty)$ with compact support
  contained in $V$ and $\eta(1)\neq0$.
\item Let $\Gamma\subset G$ be a subgroup, $X$ a topological space,
  $n\colon X\to G$ a map and $a\in G$. Then $(\Gamma,n,a)$ has the \emph{Siegel
  property} if there is a compact $\Omega\subset G$ such that for every
  $g\in G$ there are $\gamma\in\Gamma$, an integer $k\geq0$ and $\kappa\in\Omega$
  with $\gamma g=a^k\kappa$. It has the \emph{Lipschitz property} if every
  neighbourhood $V$ of $1$ contains a bump $\eta$ with the following property:
  for every compact $\Omega\subset G$ there are $D\in\mathbf R$ and a compact
  $\Omega'\subset G$ such that for all $x\in X$ and $k\geq0$ there are
  $\gamma\in\Gamma$ and $w\in\Omega'$ with $\gamma\,n(x)\,a^k=a^kw$ and
  $|\eta(\kappa^{-1}w^{-1}y)-\eta(\kappa^{-1}y)|\leq D\,2^{-k}$ for all
  $\kappa\in\Omega$ and $y\in G$.
\item In (4), let $\lambda$ be a Borel probability measure on $X$ and
  $l\in\mathbf C$. A pair $(\varphi,c)$ of functions on $G$ is a
  \emph{recursion pair} for $(\Gamma,n,a,\lambda,l)$ if $\varphi$ and $c$ are
  continuous and, for all $g\in G$ and $\gamma\in\Gamma$,
  $\varphi(\gamma g)=\varphi(g)$, $c(g)=\int_X\varphi(n(x)g)\,d\lambda(x)$ and
  $c(a^2g)-c(ag)=l\,(c(ag)-c(g))$.
\end{enumerate}
\end{definition}
\leanref{Development/WeilBound/ThetaLifts/EigenvalueEndgame/SignatureSplitting/RightOfFd.lean}{PositivityTheory.ThKEnd.Two.Lat}
\leanref{Development/WeilBound/ThetaLifts/EigenvalueEndgame/SignatureSplitting/RightOfFd.lean}{PositivityTheory.ThKEnd.Two.IsAv}
\leanref{Development/WeilBound/ThetaLifts/EigenvalueEndgame/SignatureSplitting/ConvCont.lean}{PositivityTheory.ThKEnd.Two.Gr}
\leanref{Development/WeilBound/ThetaLifts/EigenvalueEndgame/SignatureSplitting/ConvCont.lean}{PositivityTheory.ThKEnd.Two.Bump}
\leanref{Development/WeilBound/ThetaLifts/EigenvalueEndgame/SignatureSplitting/ConvCont.lean}{PositivityTheory.ThKEnd.Two.Cusp}
\leanref{Development/WeilBound/ThetaLifts/EigenvalueEndgame/SignatureSplitting/ConvCont.lean}{PositivityTheory.ThKEnd.Two.Dat}

\begin{lemma}\label{asm-theta-agree-end-kill-averages}
Let $G$ be a locally compact, second countable Hausdorff group.
\begin{enumerate}
\item Let $\Gamma$ be a lattice in $G$
  (Definition~\ref{asm-theta-agree-end-kill-notions}(1)). Then every Haar
  measure on $G$ is invariant under right translations, and there are a Haar
  measure $\mu$ on $G$ and a measurable fundamental domain $\mathcal F$ for
  $\Gamma$ with respect to $\mu$ with $\mu(\mathcal F)=1$.
\item Let $\Gamma$ be a lattice in $G$, let $\iota\colon G\to H$ be a
  homomorphism to a group $H$, and let $f$ be a function on $H$ with
  $f(\iota(\gamma)h)=f(h)$ for all $\gamma\in\Gamma$ and $h\in H$. Then $f$ has
  an average over $\Gamma$ along $\iota$
  (Definition~\ref{asm-theta-agree-end-kill-notions}(2)). Moreover, for every
  map $\iota'\colon G\to H$, a function on $H$ has at most one average over
  $\Gamma$ along $\iota'$, and there is a probability measure $\lambda$ on $G$,
  depending only on $\Gamma$, with $c(h)=\int_Gf'(\iota'(u)h)\,d\lambda(u)$ for
  every $h$ whenever $c$ is the average of a function $f'$ along a map
  $\iota'$.
\item In (2), let $H$ be a locally compact, first countable topological group,
  and let $\iota$ and $f$ be continuous. Then for every compact $K\subset H$
  there is $M$ with $|f(\iota(u)h)|\leq M$ for all $u\in G$ and $h\in K$, and
  the average of $f$ over $\Gamma$ along $\iota$ is continuous.
\item Let $\Gamma\subset G$ be a subgroup, $A\subset G$ a closed subgroup such
  that $\Gamma\cap A$ is a lattice in $A$, $f$ a function on $G$, and $c$ the
  average of $f$ over $\Gamma\cap A$. If $\gamma\in\Gamma$ has
  $\gamma^{-1}A\gamma=A$ and $f(\gamma g)=f(g)$ for every $g\in G$, then
  $c(\gamma g)=c(g)$ for every $g\in G$.
\end{enumerate}
\end{lemma}
\leanref{Development/WeilBound/ThetaLifts/EigenvalueEndgame/SignatureSplitting/RightOfFd.lean}{PositivityTheory.ThKEnd.Two.Lat.right}
\leanref{Development/WeilBound/ThetaLifts/EigenvalueEndgame/SignatureSplitting/RightOfFd.lean}{PositivityTheory.ThKEnd.Two.Lat.pair}
\leanref{Development/WeilBound/ThetaLifts/EigenvalueEndgame/SignatureSplitting/RightOfFd.lean}{PositivityTheory.ThKEnd.Two.Lat.isAv_ex}
\leanref{Development/WeilBound/ThetaLifts/EigenvalueEndgame/SignatureSplitting/RightOfFd.lean}{PositivityTheory.ThKEnd.Two.Lat.isAv_uniq}
\leanref{Development/WeilBound/ThetaLifts/EigenvalueEndgame/SignatureSplitting/RightOfFd.lean}{PositivityTheory.ThKEnd.Two.Lat.isAv_prob}
\leanref{Development/WeilBound/ThetaLifts/EigenvalueEndgame/SignatureSplitting/RightOfFd.lean}{PositivityTheory.ThKEnd.Two.Lat.bound}
\leanref{Development/WeilBound/ThetaLifts/EigenvalueEndgame/SignatureSplitting/RightOfFd.lean}{PositivityTheory.ThKEnd.Two.IsAv.cont}
\leanref{Development/WeilBound/ThetaLifts/EigenvalueEndgame/SignatureSplitting/StageSub.lean}{PositivityTheory.ThKEnd.Two.IsAv.conj_sub}
\begin{proof}
(1) As $\Gamma$ is a countable discrete subgroup of the locally compact, second
countable group $G$, it has a measurable fundamental domain with respect to
every Haar measure. With the compact set $C$ of
Definition~\ref{asm-theta-agree-end-kill-notions}(1), the argument of the proof
of Lemma~\ref{asm-theta-agree-end-unipotent}(1), given there for $N_H$, shows
that every fundamental domain has finite positive measure and that every Haar
measure is invariant under right translations. Rescaling a Haar measure gives
$\mu$ and $\mathcal F$.

(2) Let $\mu$ and $\mathcal F$ be as in (1), and put
$c(h)=\int_{\mathcal F}f(\iota(u)h)\,d\mu(u)$. Let $\mu'$ be a Haar measure and
$\mathcal F'$ a fundamental domain for $\Gamma$ with respect to $\mu'$ with
$\mu'(\mathcal F')=1$. Then $\mu=t\mu'$ for some $t>0$, and $\mathcal F'$ is a
fundamental domain with respect to $\mu$. Two fundamental domains have the same
measure, by the computation in the proof of
Lemma~\ref{asm-theta-agree-anisotropic-invariance} for the constant function
$1$; so $t=1$ and $\mu'=\mu$. The function $u\mapsto f(\iota(u)h)$ is invariant
under left translation by $\Gamma$, as $\iota$ is a homomorphism, so the same
computation shows that its integrals over $\mathcal F$ and $\mathcal F'$ agree.
So $c$ is an average. An average $c$ of $f'$ along $\iota'$ has
$c(h)=\int_{\mathcal F}f'(\iota'(u)h)\,d\mu(u)$, which gives the uniqueness and,
with $\lambda$ the restriction of $\mu$ to $\mathcal F$, the last claim.

(3) Let $C$ be as in (1). The set $\iota(C)K$ is compact, so $f$ is bounded on it,
by $M$ say. For $u\in G$ choose $\gamma\in\Gamma$ with $\gamma u\in C$; then
$f(\iota(u)h)=f(\iota(\gamma u)h)$, so $|f(\iota(u)h)|\leq M$ for $h\in K$. By
(2) the average is $h\mapsto\int_Gf(\iota(u)h)\,d\lambda(u)$. Every $h_0$ has a
compact neighbourhood $K$, on which the integrand is bounded by $M$ and
continuous in $h$; dominated convergence gives continuity at $h_0$, as $H$ is
first countable.

(4) The map $e(u)=\gamma^{-1}u\gamma$ is an automorphism of the topological
group $A$, and it maps $\Gamma\cap A$ onto itself, as $\gamma\in\Gamma$. Let
$\mu$ be a Haar measure on $A$ and $\mathcal F$ a fundamental domain for
$\Gamma\cap A$ with $\mu(\mathcal F)=1$, as in (1). Then $e^{-1}(\mathcal F)$ is
a fundamental domain for the Haar measure $\mu'=\mu\circ e$, with
$\mu'(e^{-1}(\mathcal F))=1$. As $u\gamma=\gamma e(u)$ and $f(\gamma g)=f(g)$,
\[
c(\gamma g)=\int_{e^{-1}(\mathcal F)}f(u\gamma g)\,d\mu'(u)
=\int_{e^{-1}(\mathcal F)}f(e(u)g)\,d\mu'(u)=\int_{\mathcal F}f(ug)\,d\mu(u)=c(g).
\]
\end{proof}
% Lean: (1) is Lat.right (right_of_fd, fd_fin) and Lat.pair, with
% FundDomDiscrete.exists_measurable_section; (2) is int_indep, Lat.isAv_ex, Lat.isAv_uniq and
% Lat.isAv_prob; (3) is Lat.bound and IsAv.cont; (4) is IsAv.conj_sub =
% ConjugateSubgroup.conj_sub, through IsAv.conj for the automorphism e. Size: about 150
% lines (RightOfFd.lean, StageSub.lean, ConjugateSubgroup.lean).

\begin{lemma}\label{asm-theta-agree-end-kill-stages}
Let $G$ and $B$ be locally compact, second countable Hausdorff groups, let
$\Gamma\subset G$ be a countable discrete subgroup, let $A\subset S\subset G$ be
closed subgroups such that $\Gamma\cap A$ is a lattice in $A$, and let
$\Gamma_B$ be a lattice in $B$
(Definition~\ref{asm-theta-agree-end-kill-notions}(1)). Let
$\iota\colon B\to G$ and $\pi\colon S\to B$ be continuous homomorphisms with
$\iota(B)\subset S$, $\pi(\iota(b))=b$ for every $b\in B$, kernel of $\pi$ equal
to $A$, $\iota(\Gamma_B)\subset\Gamma$ and $\pi(\Gamma\cap S)\subset\Gamma_B$.
Then $\Gamma\cap S$ is a lattice in $S$. Let further $f$ be a continuous
function on $G$ with $f(\gamma g)=f(g)$ for $\gamma\in\Gamma$, let $c_A$ be the
average of $f$ over $\Gamma\cap A$ and $c$ the average of $c_A$ over
$\Gamma_B$ along $\iota$ (Definition~\ref{asm-theta-agree-end-kill-notions}(2)).
Then $c$ is the average of $f$ over $\Gamma\cap S$.
\end{lemma}
\leanref{Development/WeilBound/ThetaLifts/EigenvalueEndgame/SignatureSplitting/StageSub.lean}{PositivityTheory.ThKEnd.Two.stage_sub}
\leanref{Development/WeilBound/ThetaLifts/EigenvalueEndgame/SignatureSplitting/RightOfFd.lean}{PositivityTheory.ThKEnd.Two.Semi.lat}
\leanref{Development/WeilBound/ThetaLifts/EigenvalueEndgame/SignatureSplitting/RightOfFd.lean}{PositivityTheory.ThKEnd.Two.Semi.stage}
\begin{proof}
For $s\in S$ put $p(s)=s\,\iota(\pi(s))^{-1}$. Then $\pi(p(s))=1$, so
$p(s)\in A$, and $s=p(s)\,\iota(\pi(s))$. For $a\in A$, $p(as)=a\,p(s)$, as
$\pi(as)=\pi(s)$; and $p(\iota(b))=1$. So $e(b,a)=a\,\iota(b)$ is a
homeomorphism from $B\times A$ onto $S$, with inverse $s\mapsto(\pi(s),p(s))$.
The groups $A$ and $S$ are locally compact and second countable, being closed in
$G$, and $\Gamma\cap S$ is countable and discrete.

\emph{Lattice.} Let $C_A$ and $C_B$ be compact sets for the lattices $\Gamma\cap A$
and $\Gamma_B$, and $s\in S$. Choose $\beta\in\Gamma_B$ with
$\beta\pi(s)\in C_B$ and $\alpha\in\Gamma\cap A$ with
$\alpha\,p(\iota(\beta)s)\in C_A$, and put $\gamma=\alpha\iota(\beta)$, which
lies in $\Gamma\cap S$. Then $\pi(\gamma s)=\beta\pi(s)\in C_B$ and
$p(\gamma s)=\alpha\,p(\iota(\beta)s)\in C_A$, so $\gamma s$ lies in the compact
set $C_A\iota(C_B)$.

\emph{A Haar measure.} By Lemma~\ref{asm-theta-agree-end-kill-averages}(1) there
are Haar measures $\mu_A$ and $\mu_B$, invariant under right translations, and
measurable fundamental domains $\mathcal A$ for $\Gamma\cap A$ and $\mathcal B$
for $\Gamma_B$ of measure $1$. Let $\nu$ be the image of $\mu_B\times\mu_A$
under $e$. For $s\in S$ and $(b,a)\in B\times A$,
\[
e(b,a)\,s=a\,\iota(b)\,p(s)\,\iota(\pi(s))
=a\,p\bigl(\iota(b)p(s)\iota(b)^{-1}\bigr)\,\iota(b\pi(s))
=e\bigl(b\pi(s),\ a\,p(\iota(b)p(s)\iota(b)^{-1})\bigr),
\]
as $\iota(b)p(s)\iota(b)^{-1}\in A$. The map
$(b,a)\mapsto(b\pi(s),a\,p(\iota(b)p(s)\iota(b)^{-1}))$ preserves
$\mu_B\times\mu_A$, as $\mu_B$ and $\mu_A$ are invariant under right
translations; so $\nu$ is invariant under right translations by $S$. It is
finite on compact sets and positive on open sets, so the image of $\nu$ under
$s\mapsto s^{-1}$ is a Haar measure. By
Lemma~\ref{asm-theta-agree-end-kill-averages}(1) for the lattice $\Gamma\cap S$
it is invariant under right translations, so $\nu$ is invariant under left
translations, and $\nu$ is a Haar measure on $S$.

\emph{A fundamental domain.} The set $\mathcal F_S=e(\mathcal B\times\mathcal A)$
is a fundamental domain for $\Gamma\cap S$, with $\nu(\mathcal F_S)=1$. Indeed,
let $s\in S$, let $\beta\in\Gamma_B$ be the element with $\beta\pi(s)\in\mathcal B$
and $\alpha\in\Gamma\cap A$ the one with $\alpha\,p(\iota(\beta)s)\in\mathcal A$;
as above, $\gamma=\alpha\iota(\beta)$ has $\gamma s\in\mathcal F_S$. If
$\gamma'\in\Gamma\cap S$ also has $\gamma's\in\mathcal F_S$, then
$\pi(\gamma')\pi(s)\in\mathcal B$ with $\pi(\gamma')\in\Gamma_B$, so
$\pi(\gamma')=\beta$; then $\gamma'=p(\gamma')\iota(\beta)$ with
$p(\gamma')=\gamma'\iota(\beta)^{-1}\in\Gamma\cap A$ and
$p(\gamma's)=p(\gamma')\,p(\iota(\beta)s)\in\mathcal A$, so $p(\gamma')=\alpha$
and $\gamma'=\gamma$.

\emph{Integration.} Let $\mu$ be a Haar measure on $S$ and $\mathcal F$ a
fundamental domain for $\Gamma\cap S$ with $\mu(\mathcal F)=1$, and let $g\in G$.
The function $u\mapsto f(ug)$ on $S$ is invariant under left translation by
$\Gamma\cap S$, and bounded, by
Lemma~\ref{asm-theta-agree-end-kill-averages}(3) for $\{g\}$. As in the proof of
Lemma~\ref{asm-theta-agree-end-kill-averages}(2), $\mu=\nu$ and the integral of
$u\mapsto f(ug)$ over $\mathcal F$ equals its integral over $\mathcal F_S$. By
Fubini's theorem,
\[
\int_{\mathcal F_S}f(ug)\,d\nu(u)=\int_{\mathcal B}\int_{\mathcal A}f(a\,\iota(b)g)\,d\mu_A(a)\,d\mu_B(b)
=\int_{\mathcal B}c_A(\iota(b)g)\,d\mu_B(b)=c(g).
\]
\end{proof}
% Lean: stage_sub (StageSub.lean, 21 lines) builds the record Semi for j the inclusion of A,
% iota, p and pi; Semi.lat (18 lines) is the lattice, with p_mul and p_iota; Semi.stage (84
% lines) is the rest: the homeomorphism e, the measure nu (map of the product), its right
% invariance through the skew product m1, the Haar property through nu.inv and Lat.right,
% the fundamental domain hfd, hnu1, and int_indep, setIntegral_prod and the two averages.

\begin{lemma}\label{asm-theta-agree-end-kill-induction}
Let $G$ be a locally compact, second countable Hausdorff group with a Haar
measure $\mu$, $\Gamma\subset G$ a subgroup, $X$ a topological space with a
Borel probability measure $\lambda$, $n\colon X\to G$ a continuous map, $a\in G$
and $l\in\mathbf C$ with $|l|<1$, and suppose that $(\Gamma,n,a)$ has the Siegel
and Lipschitz properties (Definition~\ref{asm-theta-agree-end-kill-notions}(4)).
For a function $\varphi$ on $G$ and a bump $\eta$ put
$\operatorname{conv}_\eta\varphi(g)=\int_G\varphi(gy)\eta(y)\,d\mu(y)$. Recursion
pairs are for $(\Gamma,n,a,\lambda,l)$
(Definition~\ref{asm-theta-agree-end-kill-notions}(5)).
\begin{enumerate}
\item If $(\varphi,c)$ is a recursion pair and $K\subset G$ is compact, there is
  $M$ with $|\varphi(n(x)g)|\leq M$ for all $x\in X$ and $g\in K$.
\item If $(\varphi,c)$ is a recursion pair and $\eta$ is a bump, then
  $(\operatorname{conv}_\eta\varphi,\operatorname{conv}_\eta c)$ is a recursion
  pair.
\item Let $P$ be a property of pairs of functions on $G$. Suppose that
  $P(\varphi,c)$ holds for every recursion pair $(\varphi,c)$ with $\varphi$
  bounded, and for every recursion pair $(\varphi,c)$ such that every
  neighbourhood of $1$ contains a bump $\eta$ with
  $P(\operatorname{conv}_\eta\varphi,\operatorname{conv}_\eta c)$. Then
  $P(\varphi,c)$ holds for every recursion pair $(\varphi,c)$ such that
  $\varphi$ has growth of some order along $a$.
\item If $(\varphi,c)$ is a recursion pair and $\varphi$ has growth of some
  order along $a$, then $c(ag)=c(g)$ for every $g\in G$.
\end{enumerate}
\end{lemma}
\leanref{Development/WeilBound/ThetaLifts/EigenvalueEndgame/SignatureSplitting/ConvCont.lean}{PositivityTheory.ThKEnd.Two.conv}
\leanref{Development/WeilBound/ThetaLifts/EigenvalueEndgame/SignatureSplitting/ConvCont.lean}{PositivityTheory.ThKEnd.Two.Dat.bound}
\leanref{Development/WeilBound/ThetaLifts/EigenvalueEndgame/SignatureSplitting/ConvCont.lean}{PositivityTheory.ThKEnd.Two.Dat.conv}
\leanref{Development/WeilBound/ThetaLifts/EigenvalueEndgame/SignatureSplitting/ConvCont.lean}{PositivityTheory.ThKEnd.Two.kill_ind}
\leanref{Development/WeilBound/ThetaLifts/EigenvalueEndgame/SignatureSplitting/ConvCont.lean}{PositivityTheory.ThKEnd.Two.kill}
\begin{proof}
For a continuous $\varphi$ and a bump $\eta$, the function
$y\mapsto\varphi(gy)\eta(y)$ is continuous with compact support, and
$\operatorname{conv}_\eta\varphi$ is continuous, as an integral over the compact
support of $\eta$ of a continuous function of $(g,y)$.

(1) The Lipschitz property for the neighbourhood $G$ and the compact set
$\emptyset$ gives a compact $\Omega'$ such that for every $x$ there are
$\gamma\in\Gamma$ and $w\in\Omega'$ with $\gamma n(x)=w$ (take $k=0$). Let $M$
bound $|\varphi|$ on the compact set $\Omega'K$. For $g\in K$,
$\varphi(n(x)g)=\varphi(\gamma n(x)g)=\varphi(wg)$, so $|\varphi(n(x)g)|\leq M$.

(2) Both functions are continuous. For $\gamma\in\Gamma$,
$\operatorname{conv}_\eta\varphi(\gamma g)=\int\varphi(\gamma gy)\eta(y)\,d\mu(y)=\operatorname{conv}_\eta\varphi(g)$.
Let $g\in G$. By (1) for the compact set $g\operatorname{supp}\eta$, the
function $(x,y)\mapsto\varphi(n(x)gy)\eta(y)$ is bounded by $M\eta(y)$, which
is integrable for $\lambda\times\mu$. So Fubini's theorem gives
\[
\operatorname{conv}_\eta c(g)=\int_G\int_X\varphi(n(x)gy)\,d\lambda(x)\,\eta(y)\,d\mu(y)
=\int_X\operatorname{conv}_\eta\varphi(n(x)g)\,d\lambda(x).
\]
The recursion for $\operatorname{conv}_\eta c$ follows from that for $c$ at
$gy$, multiplied by $\eta(y)$ and integrated.

\emph{Growth of $c$.} Let $(\varphi,c)$ be a recursion pair. By induction on
$k$, $c(a^{k+1}g)-c(a^kg)=l^k(c(ag)-c(g))$, so
$c(a^k\kappa)=c(\kappa)+\bigl(\sum_{i<k}l^i\bigr)(c(a\kappa)-c(\kappa))$, and
$|\sum_{i<k}l^i|\leq(1-|l|)^{-1}$. So for compact $\Omega$, with $M$ and $M'$
bounds of $|c|$ and $|c(a\,\cdot)-c|$ on $\Omega$,
$|c(a^k\kappa)|\leq M+(1-|l|)^{-1}M'$ for $k\geq0$ and $\kappa\in\Omega$.

\emph{One step.} Let $(\varphi,c)$ be a recursion pair with $\varphi$ of growth
of order $m+1$ along $a$, and $V$ a neighbourhood of $1$. Let $\eta$ be the bump
in $V$ of the Lipschitz property. We show that $\operatorname{conv}_\eta\varphi$
has growth of order $m$. Let $\Omega$ be compact, and $D$ and $\Omega'$ as in
the Lipschitz property. By (2) and the growth of $c$,
$|\operatorname{conv}_\eta c(a^k\kappa)|\leq C_0$ for some $C_0$, all $k\geq0$
and $\kappa\in\Omega$. The set
$K=\Omega'\Omega\operatorname{supp}\eta\cup\Omega\operatorname{supp}\eta$ is
compact; let $C$ be such that $|\varphi(a^kz)|\leq C2^{k(m+1)}$ for $z\in K$. Let
$x\in X$, $k\geq0$ and $\kappa\in\Omega$, and let $\gamma$ and $w$ be as in the
Lipschitz property. As $\mu$ is invariant under left translations,
$\operatorname{conv}_\eta\varphi(a^kz)=\int\varphi(a^ky)\eta(z^{-1}y)\,d\mu(y)$,
and the integrand vanishes for $y\notin K$ when $z=\kappa$ or $z=w\kappa$. As
$\operatorname{conv}_\eta\varphi$ is invariant under $\Gamma$ and
$\gamma n(x)a^k=a^kw$,
$\operatorname{conv}_\eta\varphi(n(x)a^k\kappa)=\operatorname{conv}_\eta\varphi(a^kw\kappa)$,
so
\[
\bigl|\operatorname{conv}_\eta\varphi(a^k\kappa)-\operatorname{conv}_\eta\varphi(n(x)a^k\kappa)\bigr|
\leq\int_K|\varphi(a^ky)|\,|\eta(\kappa^{-1}y)-\eta(\kappa^{-1}w^{-1}y)|\,d\mu(y)
\leq CD\mu(K)\,2^{km}.
\]
As $\lambda$ is a probability measure, (2) gives
$\operatorname{conv}_\eta\varphi(a^k\kappa)=\operatorname{conv}_\eta c(a^k\kappa)+\int_X(\operatorname{conv}_\eta\varphi(a^k\kappa)-\operatorname{conv}_\eta\varphi(n(x)a^k\kappa))\,d\lambda(x)$,
so $|\operatorname{conv}_\eta\varphi(a^k\kappa)|\leq(C_0+CD\mu(K))2^{km}$.

(3) By induction on the order $m$. If $\varphi$ has growth of order $0$, let
$\Omega$ be the compact set of the Siegel property and $C$ the constant of the
growth for $\Omega$. Every $g$ has $\gamma g=a^k\kappa$ with $\gamma\in\Gamma$ and
$\kappa\in\Omega$, so $|\varphi(g)|=|\varphi(a^k\kappa)|\leq C$, and $P(\varphi,c)$
holds. If $\varphi$ has growth of order $m+1$, then by the previous paragraph and
(2) every neighbourhood of $1$ contains a bump $\eta$ such that
$(\operatorname{conv}_\eta\varphi,\operatorname{conv}_\eta c)$ is a recursion
pair whose first function has growth of order $m$; by induction
$P(\operatorname{conv}_\eta\varphi,\operatorname{conv}_\eta c)$ holds, and so
does $P(\varphi,c)$.

(4) Apply (3) to the property that $c(ag)=c(g)$ for every $g$. If $|\varphi|$ is
bounded by $M$, so is $|c|$, as $\lambda$ is a probability measure. For $g$ and
$k\geq0$, the identity $c(a^{k+1}g')-c(a^kg')=l^k(c(ag')-c(g'))$ for
$g'=a^{-k}g$ gives $|c(ag)-c(g)|\leq2M|l|^k$, which tends to $0$. For the second
hypothesis, put $\Delta(g)=c(ag)-c(g)$, a continuous function. For every bump
$\eta$ with $\operatorname{conv}_\eta c(ag)=\operatorname{conv}_\eta c(g)$ for
all $g$, $\int\Delta(gy)\eta(y)\,d\mu(y)=0$. Suppose that $\Delta(g)\neq0$, and
let $V$ be the set of the $y$ with $\operatorname{Re}(\Delta(gy)/\Delta(g))>0$,
an open neighbourhood of $1$. For a bump $\eta$ in $V$ with this property, the
function $y\mapsto\operatorname{Re}(\Delta(gy)/\Delta(g))\eta(y)$ is continuous,
nonnegative, with compact support and positive at $1$; so its integral is
positive, but it is $\operatorname{Re}(0/\Delta(g))=0$. So $\Delta(g)=0$.
\end{proof}
% Lean: ConvCont.lean (212 lines): the first paragraph is Bump.int and Bump.conv_cont; (1) is
% Dat.bound; (2) is Dat.conv; the growth of c is Dat.pow and Dat.gr; one step is Dat.step;
% (3) is kill_ind; (4) is kill, with Bump.zero for the last part. Part (3) is shared with
% Lemma asm-theta-agree-end-cusp, whose kill_coef (KillCoef.lean) applies kill_ind with
% l = 0.

\begin{theorem}[Reduction theory for $SL_2$ \cite{God63,Bor63}]\label{asm-theta-agree-end-kill-siegel-reduction}
Let $E$ be a number field. For a real $t>0$ let $d(t)\in SL_2(\mathbf A_E)$ have
component $\operatorname{diag}(t,t^{-1})$ at every infinite place and $1$ at
every finite place. There are compact sets $\Omega_0$ and $K_0$ in
$SL_2(\mathbf A_E)$, with every element of $\Omega_0$ upper triangular, and a
real $t_0>0$, such that every $g\in SL_2(\mathbf A_E)$ has
$\gamma g=p\,d(t)\,k$ for some $\gamma\in SL_2(E)$, $p\in\Omega_0$, real
$t\geq t_0$ and $k\in K_0$.
\end{theorem}
\leanref{Formalism/UnipotentAdelicQuotients/SiegelSetsSL2Adelic.lean}{Formalism.D.siegelSetsSL2Adelic}
\leanref{Development/AutomorphicSpectrum/Proofs/AdelicProofs/SiegelSetsSL2AdelicTie.lean}{GlobalTheory.CompanionTheory.Classical.siegelSetsSL2Adelic_tie}
% Lean: Formalism.D.siegelSetsSL2Adelic, supplied by Classical.siegelSetsSL2Adelic_tie =
% GroundTheory.BlockD.siegelSetsSL2Adelic_holds (ExistsPosBoundsOfIsCompact.lean), which
% derives it from the Lean's Siegel sets for GL_n over the adeles (exists_siegelSet).

\begin{lemma}\label{asm-theta-agree-end-kill-siegel-adeles}
Let $E$ be a number field with ring of integers $\mathcal O_E$, and write
$\mathbf A_E=\mathbf A_{E,\infty}\times\mathbf A_{E,f}$ with
$\mathbf A_{E,\infty}=\prod_{w\mid\infty}E_w$. For $x\in\mathbf A_E$ let
$x_\infty$ and $x_f$ be its two components, and put
$\|x_\infty\|=\max_w|x_w|$, with the usual absolute value on $E_w$, which is
$\mathbf R$ or $\mathbf C$.
\begin{enumerate}
\item For $r\in\mathbf R$ and a compact $C\subset\mathbf A_{E,f}$, the set of
  the $x\in\mathbf A_E$ with $\|x_\infty\|\leq r$ and $x_f\in C$ is compact.
\item Every neighbourhood of $0$ in $\mathbf A_E$ contains the set of the $x$
  with $\|x_\infty\|\leq\varepsilon$ and $x_f\in U$, for some $\varepsilon>0$
  and some neighbourhood $U$ of $0$ in $\mathbf A_{E,f}$.
\item For every neighbourhood $U$ of $0$ in $\mathbf A_{E,f}$ there are an open
  compact subgroup $L\subset U$ of $\mathbf A_{E,f}$ and $D\in\mathbf R$ such
  that every $x\in\mathbf A_E$ has some $\xi\in E$ with
  $\|(\xi+x)_\infty\|\leq D$ and $(\xi+x)_f\in L$.
\end{enumerate}
\end{lemma}
\leanref{Development/WeilBound/ThetaLifts/EigenvalueEndgame/SignatureSplitting/NonarchFiniteAdeleRing.lean}{PositivityTheory.ThKEnd.Two.ball_cpt}
\leanref{Development/WeilBound/ThetaLifts/EigenvalueEndgame/SignatureSplitting/NonarchFiniteAdeleRing.lean}{PositivityTheory.ThKEnd.Two.nhds_ad}
\leanref{Development/WeilBound/ThetaLifts/EigenvalueEndgame/SignatureSplitting/NonarchFiniteAdeleRing.lean}{PositivityTheory.ThKEnd.Two.small_ad}
\leanref{Development/WeilBound/ThetaLifts/EigenvalueEndgame/SignatureSplitting/NonarchFiniteAdeleRing.lean}{PositivityTheory.ThKEnd.Two.nonarch_finiteAdeleRing}
\begin{proof}
(1) The set is the product of a closed ball in the finite-dimensional real
vector space $\mathbf A_{E,\infty}$, which is compact, with $C$.

(2) The topology of $\mathbf A_E$ is the product topology, and the balls
$\|z\|\leq\varepsilon$ with $\varepsilon>0$ form a basis of neighbourhoods of
$0$ in $\mathbf A_{E,\infty}$.

(3) The ring $\widehat{\mathcal O}_E=\prod_v\mathcal O_{E,v}$ of integral finite
adeles is an open compact subring of $\mathbf A_{E,f}$, with the product
topology. The subgroups of $\widehat{\mathcal O}_E$ that are products of open
subgroups of the $\mathcal O_{E,v}$, equal to $\mathcal O_{E,v}$ for almost all
$v$, are open and form a basis of neighbourhoods of $0$; so $U$ contains an
open subgroup $V$ of $\mathbf A_{E,f}$. Put $L=V\cap\widehat{\mathcal O}_E$. It is
an open subgroup, hence closed, so it is compact. As $\widehat{\mathcal O}_E$ is
compact and $L$ is open in it, $M=[\widehat{\mathcal O}_E:L]$ is finite, and
$Mz\in L$ for every $z\in\widehat{\mathcal O}_E$. Every finite adele differs
from an element of $E$ by an element of $\widehat{\mathcal O}_E$, by the
approximation theorem, and $\mathcal O_E$ is a lattice in $\mathbf A_{E,\infty}$
\cite[Chapter~I, \S5]{Neu99}. So there is a compact $C\subset\mathbf A_{E,\infty}$
such that every $x\in\mathbf A_E$ has some $k\in E$ with $(x-k)_\infty\in C$ and
$(x-k)_f\in\widehat{\mathcal O}_E$. Let $D_0$ bound $\|z\|$ on $C$, and put
$D=MD_0$. For $x\in\mathbf A_E$ choose such a $k$ for $M^{-1}x$, and put
$\xi=-Mk$. Then $\xi+x=M(M^{-1}x-k)$, so $\|(\xi+x)_\infty\|\leq MD_0$ and
$(\xi+x)_f\in L$.
\end{proof}
% Lean: ball_cpt (8 lines), nhds_ad (5 lines) and small_ad (28 lines), with
% nonarch_finiteAdeleRing, nonarch_pi and nonarch_of_isInducing (about 35 lines) for the open
% subgroup V, isCompact_integralFiniteAdeles, quotient_finite_of_isOpen and
% nsmul_index_mem for M, and GroundTheory.Global.exists_isCompact_forall_exists_sub_mem
% (exists_sub_mem_integralFiniteAdeles and the fundamental set of O_E) for C.

\begin{lemma}\label{asm-theta-agree-end-kill-siegel-bump}
Let $E$ be a number field, with the notation of
Lemma~\ref{asm-theta-agree-end-kill-siegel-adeles}, and let $W$ be a
neighbourhood of $1$ in $SL_2(\mathbf A_E)$. There are a bump $\eta$ in $W$
(Definition~\ref{asm-theta-agree-end-kill-notions}(3)) and an open subgroup $L$
of $\mathbf A_{E,f}$ with the following property: if
$y,y'\in SL_2(\mathbf A_E)$ and $b\in\mathbf R$ are such that every entry $z$ of
$y-y'$ has $z_f\in L$ and $\|z_\infty\|\leq b$, then
$|\eta(y')-\eta(y)|\leq4b$.
\end{lemma}
\leanref{Development/WeilBound/ThetaLifts/EigenvalueEndgame/SignatureSplitting/ClosedEmb.lean}{PositivityTheory.ThKEnd.Two.bump_ex}
\leanref{Development/WeilBound/ThetaLifts/EigenvalueEndgame/SignatureSplitting/ClosedEmb.lean}{PositivityTheory.ThKEnd.Two.gauge}
\leanref{Development/WeilBound/ThetaLifts/EigenvalueEndgame/SignatureSplitting/ClosedEmb.lean}{PositivityTheory.ThKEnd.Two.nhds_box}
\begin{proof}
As $SL_2(\mathbf A_E)$ carries the topology of a subspace of
$M_2(\mathbf A_E)=\mathbf A_E^4$, there is a neighbourhood $O$ of $0$ in
$\mathbf A_E$ such that $y\in W$ whenever every entry of $y-1$ lies in $O$. By
Lemma~\ref{asm-theta-agree-end-kill-siegel-adeles}(2) there are
$\varepsilon_0>0$ and a neighbourhood $U$ of $0$ in $\mathbf A_{E,f}$ such that
$x\in O$ whenever $\|x_\infty\|\leq\varepsilon_0$ and $x_f\in U$, and by
part~(3) of that lemma there is an open compact subgroup $L\subset U$. Put
$\varepsilon=\min(\varepsilon_0,1/2)$.

For $x\in\mathbf A_E$ put $\delta(x)=\|x_\infty\|$ if $x_f\in L$ and
$\delta(x)=\|x_\infty\|+1$ otherwise. As $L$ is open and closed, $\delta$ is
continuous. Also $\delta(0)=0$ and $\delta\geq0$. If $\delta(x)\leq\varepsilon$,
then $x_f\in L$, as $\varepsilon<1$, and $\|x_\infty\|\leq\varepsilon$. If
$(x-x')_f\in L$, then $x_f\in L$ if and only if $x'_f\in L$, so
$|\delta(x)-\delta(x')|\leq\|(x-x')_\infty\|$.

For $y\in SL_2(\mathbf A_E)$ put $e(y)=\sum_{i,j}\delta((y-1)_{ij})$ and
$\eta(y)=\max(\varepsilon-e(y),0)$. Then $\eta$ is continuous and nonnegative,
and $\eta(1)=\varepsilon\neq0$. On the closure of the set where $\eta\neq0$,
$e(y)\leq\varepsilon$, so every entry $x$ of $y-1$ has $\delta(x)\leq\varepsilon$,
hence $\|x_\infty\|\leq\varepsilon$ and $x_f\in L$. The set of the $y$ with this
property is compact, by Lemma~\ref{asm-theta-agree-end-kill-siegel-adeles}(1)
and as $SL_2(\mathbf A_E)$ is closed in $M_2(\mathbf A_E)$; and it lies in $W$,
as $\varepsilon\leq\varepsilon_0$ and $L\subset U$. So $\eta$ is a bump in $W$.
Finally, let $y$, $y'$ and $b$ be as in the statement. The entries of
$(y-1)-(y'-1)$ are those of $y-y'$, so
\[
|\eta(y')-\eta(y)|\leq|e(y)-e(y')|\leq\sum_{i,j}\bigl|\delta((y-1)_{ij})-\delta((y'-1)_{ij})\bigr|\leq4b.
\]
\end{proof}
% Lean: bump_ex (ClosedEmb.lean, 50 lines), with nhds_box (15 lines) for O, the
% hypotheses hn and hL of lip_of (nhds_ad and small_ad in cusp_of), and gauge (30 lines)
% for delta. The Lean states it for a topological ring R with continuous ring maps p to a
% normed ring and q to a topological ring satisfying (1) to (3) of Lemma
% asm-theta-agree-end-kill-siegel-adeles; cusp_of takes the two projections of the adele
% ring of E.

\begin{lemma}\label{asm-theta-agree-end-kill-siegel}
Let $E$ be a CM field with maximal totally real subfield $F$, let
$u\in\mathbf A_E^\times$ have component $2$ at the infinite places and $1$ at
the finite places, and put $a_0=\operatorname{diag}(u,u^{-1})$ and
$n(\xi)=\left(\begin{smallmatrix}1&\xi\\0&1\end{smallmatrix}\right)$ in
$SL_2(\mathbf A_E)$ for $\xi\in\mathbf A_E$.
\begin{enumerate}
\item $\mathbf A_E$ is locally compact and second countable, and $E$ is a
  lattice in the additive group of $\mathbf A_E$
  (Definition~\ref{asm-theta-agree-end-kill-notions}(1)).
\item $(SL_2(E),n,a_0)$ has the Siegel property and the Lipschitz property in
  $SL_2(\mathbf A_E)$ (Definition~\ref{asm-theta-agree-end-kill-notions}(4)).
\end{enumerate}
\end{lemma}
\leanref{Development/WeilBound/ThetaLifts/EigenvalueEndgame/SignatureSplitting/AeFacts.lean}{PositivityTheory.ThKEnd.Two.ae_facts}
\leanref{Development/WeilBound/ThetaLifts/EigenvalueEndgame/SignatureSplitting/NonarchFiniteAdeleRing.lean}{PositivityTheory.ThKEnd.Two.lat_of}
\leanref{Development/WeilBound/ThetaLifts/EigenvalueEndgame/SignatureSplitting/NonarchFiniteAdeleRing.lean}{PositivityTheory.ThKEnd.Two.cusp_of}
\begin{proof}
(1) The ring $\mathbf A_E=\mathbf A_F\otimes_FE$ is the ring of adeles of $E$,
which is locally compact and second countable. The field $E$ is countable, it is
discrete in $\mathbf A_E$, and $\mathbf A_E/E$ is compact
\cite[Chapter~VII, \S2]{Lan94}. As $\mathbf A_E$ is locally compact, finitely
many compact neighbourhoods of points cover the image of $\mathbf A_E/E$, and
their union is a compact set $C$ with $E+C=\mathbf A_E$.

(2) As in (1), identify $\mathbf A_E$ with the ring of adeles of $E$, and use
the notation of Lemma~\ref{asm-theta-agree-end-kill-siegel-adeles}. The adele
$u^{-1}$ has components $1/2$ at the infinite places and $1$ at the finite
places, so $\|(u^{-1}z)_\infty\|\leq\|z_\infty\|/2$ and $(u^{-1}z)_f=z_f$ for
$z\in\mathbf A_E$. For $z\in\mathbf A_E$ and $k\geq0$,
$a_0^{-k}n(z)a_0^k=n(u^{-2k}z)$; and $n(\xi)\in SL_2(E)$ for $\xi\in E$.

\emph{The Siegel property.} Let $\Omega_0$, $K_0$ and $t_0$ be as in
Theorem~\ref{asm-theta-agree-end-kill-siegel-reduction}. Let $C$ be the union of
the four compact sets $\{p_{ij}:p\in\Omega_0\}$, let $r$ bound $\|x_\infty\|$ on
$C$, and put $C_f=\{x_f:x\in C\}$. Let $\Omega_1$ be the set of the
$g\in SL_2(\mathbf A_E)$ all of whose entries $x$ have $\|x_\infty\|\leq r$ and
$x_f\in C_f$, and $\Omega_2$ the set of those all of whose entries $x$ have
$\|x_\infty\|\leq2t_0+t_0^{-1}$ and $x_f\in\{0,1\}$. By
Lemma~\ref{asm-theta-agree-end-kill-siegel-adeles}(1), and as
$SL_2(\mathbf A_E)$ is closed in $M_2(\mathbf A_E)$, $\Omega_1$ and $\Omega_2$
are compact, and so is $\Omega=\Omega_1\Omega_2K_0$. Let
$g\in SL_2(\mathbf A_E)$, and let $\gamma g=p\,d(t)\,k$ as in
Theorem~\ref{asm-theta-agree-end-kill-siegel-reduction}. As $t/t_0\geq1$, there
is an integer $j\geq0$ with $2^jt_0\leq t<2^{j+1}t_0$. Then
\[
\gamma g=a_0^j\cdot\bigl(a_0^{-j}pa_0^j\bigr)\cdot\bigl(a_0^{-j}d(t)\bigr)\cdot k .
\]
As $p$ is upper triangular, $a_0^{-j}pa_0^j$ has the entries of $p$, except
the upper right one, which is $z=u^{-2j}p_{12}$; and
$\|z_\infty\|=4^{-j}\|(p_{12})_\infty\|\leq r$ and $z_f=(p_{12})_f$. So
$a_0^{-j}pa_0^j\in\Omega_1$. The element $a_0^{-j}d(t)$ has components
$\operatorname{diag}(2^{-j}t,2^jt^{-1})$ at the infinite places and $1$ at the
finite places, and $t_0\leq2^{-j}t<2t_0$, so it lies in $\Omega_2$. Hence
$\gamma g=a_0^j\kappa$ with $\kappa\in\Omega$.

\emph{The Lipschitz property.} Let $V$ be a neighbourhood of $1$, and let $\eta$
and $L$ be as in Lemma~\ref{asm-theta-agree-end-kill-siegel-bump} for $W=V$. Let
$\Omega\subset SL_2(\mathbf A_E)$ be compact, and put
$N=\left(\begin{smallmatrix}0&1\\0&0\end{smallmatrix}\right)$. The map
$(\kappa,z)\mapsto\kappa^{-1}N\kappa z$ from $SL_2(\mathbf A_E)^2$ to
$M_2(\mathbf A_E)$ is continuous, so the set $C$ of the entries of the matrices
$\kappa^{-1}N\kappa z$ with $\kappa\in\Omega$ and $z$ in the support of $\eta$
is compact. Let $D_1\geq0$ bound $\|c_\infty\|$ for $c\in C$. As $C$ is compact,
$L$ is open and $(s,c)\mapsto sc_f$ is continuous, the set $U$ of the
$s\in\mathbf A_{E,f}$ with $sc_f\in L$ for every $c\in C$ is a neighbourhood of
$0$. By Lemma~\ref{asm-theta-agree-end-kill-siegel-adeles}(3) there are an open
compact subgroup $L'\subset U$ and $D$ such that every $x\in\mathbf A_E$ has
some $\xi\in E$ with $\|(\xi+x)_\infty\|\leq D$ and $(\xi+x)_f\in L'$. Let
$\Omega'$ be the set of the $n(w)$ with $\|w_\infty\|\leq D$ and $w_f\in L'$,
which is compact by Lemma~\ref{asm-theta-agree-end-kill-siegel-adeles}(1). We
show that $D'=4DD_1$ and $\Omega'$ have the property required.

Let $x\in\mathbf A_E$ and $k\geq0$. Choose $\xi$ as above, so that $D\geq0$, and
put $\gamma=n(\xi)\in SL_2(E)$ and $w=u^{-2k}(\xi+x)$. Then
$\|w_\infty\|\leq4^{-k}D\leq2^{-k}D$ and $w_f=(\xi+x)_f\in L'$, so
$n(w)\in\Omega'$, and
$\gamma\,n(x)\,a_0^k=n(\xi+x)\,a_0^k=a_0^k\,n(w)$. Let $\kappa\in\Omega$ and
$y\in SL_2(\mathbf A_E)$. As $n(w)^{-1}=1-wN$ and $N^2=0$,
\[
\kappa^{-1}y-\kappa^{-1}n(w)^{-1}y=w\,\kappa^{-1}Ny,\qquad
N\,n(w)^{-1}y=Ny .
\]
Suppose that one of $\kappa^{-1}y$ and $\kappa^{-1}n(w)^{-1}y$ lies in the
support of $\eta$, and call it $z$. In both cases $\kappa^{-1}N\kappa z=\kappa^{-1}Ny$,
so the entries of $\kappa^{-1}Ny$ lie in $C$, and every entry of
$\kappa^{-1}y-\kappa^{-1}n(w)^{-1}y$ is $wc$ with $c\in C$. Then
$(wc)_f=w_fc_f\in L$, as $w_f\in U$, and
$\|(wc)_\infty\|\leq2^{-k}DD_1$. By
Lemma~\ref{asm-theta-agree-end-kill-siegel-bump},
$|\eta(\kappa^{-1}n(w)^{-1}y)-\eta(\kappa^{-1}y)|\leq D'2^{-k}$. If neither
lies in the support of $\eta$, both values of $\eta$ are $0$.
% Lean: (2) is cusp_of (NonarchFiniteAdeleRing.lean, 30 lines), which ae_facts transports
% along exists_tensor_adele_equiv from the adele ring of E to A_F tensor E, with a_0 = dia
% (halfE)^{-1} and norm_half. The Siegel property is siegel_ad (73 lines): Omega_1 and
% Omega_2 are the two box sets (box, ball_cpt), j comes from exists_nat_pow_near, and the
% two memberships are c and n1 to n7. The Lipschitz property is lip_of (ClosedEmb.lean, 67
% lines): bump_ex for eta and L, T for the map (kappa, z) -> kappa^{-1} N kappa z, D_1 from
% hC, L' and D from the hypothesis hL (small_ad) for the neighbourhood U, w and w1, w2 for
% its bounds, uni_dia, and diff_uni and nil_uni for the two identities. Size: about 415
% lines in SignatureSplitting, beyond the Siegel sets of Theorem
% asm-theta-agree-end-kill-siegel-reduction. Part (1) is lat_of (with small_ad and
% discreteTopology_principalSubgroup) and sl_facts. This lemma is shared with Lemma
% asm-theta-agree-end-cusp, whose wh_two uses ae_facts.
\end{proof}

\begin{lemma}\label{asm-theta-agree-end-kill-growth}
Let $E$ be a CM field with a theta base $B$
(Definition~\ref{asm-theta-agree-models}(1)), let $W$ be a skew-Hermitian
lattice space of some signature, and let $\operatorname{ht}$ be the height of
Definition~\ref{asm-theta-agree-parts}. Call a function $f$ on $U_W$
\emph{moderate} if there are $C$ and $k$ with $|f(h)|\leq C\operatorname{ht}(h)^k$
for every $h\in U_W$.
\begin{enumerate}
\item Let $f$ be moderate, $G'$ a topological group, $m\colon G'\to U_W$ a
  continuous homomorphism, $a\in G'$ and $h\in U_W$. Then
  $g\mapsto f(m(g)h)$ has growth of some order along $a$
  (Definition~\ref{asm-theta-agree-end-kill-notions}(3)).
\item Let $N\subset U_W$ be a closed subgroup such that $N\cap U_W(F)$ is a
  lattice in $N$ (Definition~\ref{asm-theta-agree-end-kill-notions}(1)), and
  let $f$ be moderate with $f(\gamma h)=f(h)$ for $\gamma\in U_W(F)$. Then the
  average of $f$ over $N\cap U_W(F)$ is moderate.
\end{enumerate}
\end{lemma}
\leanref{Development/WeilBound/ThetaLifts/EigenvalueEndgame/SignatureSplitting/AeFacts.lean}{PositivityTheory.ThKEnd.Two.Mod}
\leanref{Development/WeilBound/ThetaLifts/EigenvalueEndgame/SignatureSplitting/AeFacts.lean}{PositivityTheory.ThKEnd.Two.Mod.gr}
\leanref{Development/WeilBound/ThetaLifts/EigenvalueEndgame/SignatureSplitting/AeFacts.lean}{PositivityTheory.ThKEnd.Two.Mod.av}
\begin{proof}
The height is continuous on $U_W$ and $\operatorname{ht}(gh)\leq\operatorname{ht}(g)\operatorname{ht}(h)$,
as operator norms and absolute values are submultiplicative. Let $C$ and $k$ be
as in the definition, and put $C^+=\max(C,0)$.

(1) Choose $n$ with $\operatorname{ht}(m(a))<2^n$. Let $\Omega$ be compact and
$M$ a bound of $\kappa\mapsto\operatorname{ht}(m(\kappa)h)$ on $\Omega$. For
$i\geq0$ and $\kappa\in\Omega$,
$\operatorname{ht}(m(a^i\kappa)h)\leq\operatorname{ht}(m(a))^i\operatorname{ht}(m(\kappa)h)\leq M2^{in}$,
so $|f(m(a^i\kappa)h)|\leq C^+M^k2^{ink}$. So the order $nk$ will do.

(2) Let $C_N$ be a compact set for the lattice $N\cap U_W(F)$, and $M$ a bound of
the height on $C_N$. By Lemma~\ref{asm-theta-agree-end-kill-averages}(1) there
are a Haar measure $\mu$ on $N$ and a fundamental domain $\mathcal F$ with
$\mu(\mathcal F)=1$, so the average is
$c(h)=\int_{\mathcal F}f(uh)\,d\mu(u)$. For $u\in N$ choose
$\gamma\in N\cap U_W(F)$ with $\gamma u\in C_N$; then
$|f(uh)|=|f(\gamma uh)|\leq C^+(M\operatorname{ht}(h))^k$. So
$|c(h)|\leq C^+M^k\operatorname{ht}(h)^k$.
\end{proof}
% Lean: Mod.gr and Mod.av (AeFacts.lean, 36 lines), with ht_mul (heightGL_mul_le) and
% ht_cont; Mod.av takes the compact set and the pair from UnipGood. Shared with Lemma
% asm-theta-agree-end-cusp (wh_two).

\begin{lemma}\label{asm-theta-agree-end-kill-unipotent}
Let $E$ be a CM field with a theta base $B$
(Definition~\ref{asm-theta-agree-models}(1)), let $W$ be a skew-Hermitian
lattice space of rank $d$ and some signature, with
$h_W(u,u')=u^{\mathsf T}J_Wc(u')$, and let $U\subset E^d$ be an $E$-subspace.
Let $N_U$ be the subgroup of $U_W$ generated by the $g\in U_W$ with $gu=u$ for
every $u\in U$ and $gz-z\in\mathbf A_EU$ for every $z\in E^d$ with $h_W(z,u)=0$
for all $u\in U$.
\begin{enumerate}
\item An element $g\in U_W$ lies in $N_U$ if and only if $gu=u$ for every
  $u\in U$ and $gz-z\in\mathbf A_EU$ for every $z\in\mathbf A_E^d$ with
  $h_W(z,u)=0$ for all $u\in U$, and if and only if this holds for the
  $z\in E^d$.
\item The group $N_U$ is closed in $U_W$, locally compact and second
  countable, $N_U\cap U_W(F)$ is a lattice in $N_U$
  (Definition~\ref{asm-theta-agree-end-kill-notions}(1)), every Haar measure
  on $N_U$ is invariant under right translations, and some Haar measure on
  $N_U$ has a fundamental domain for $N_U\cap U_W(F)$ of measure $1$.
\end{enumerate}
For $U\neq0$ with $h_W(u,u')=0$ for all $u,u'\in U$, $N_U$ is the group of
Definition~\ref{asm-theta-agree-end-notions}(4); for a plane $H=(x,y,\iota)$ in
$W$, $N_{Ex}=N_H$, by (1) and Lemma~\ref{asm-theta-agree-end-unipotent}(3).
\end{lemma}
\leanref{Development/WeilBound/ThetaLifts/EigenvalueEndgame/UnipotentGroups/UnipUGood.lean}{PositivityTheory.ThKEnd.Grp.UnipUGood}
\leanref{Development/WeilBound/ThetaLifts/EigenvalueEndgame/UnipotentGroups/ContinuousExpG.lean}{PositivityTheory.ThKEnd.Grp.unipUGood}
\leanref{Development/WeilBound/ThetaLifts/EigenvalueEndgame/UnipotentGroups/ContinuousExpG.lean}{PositivityTheory.ThKEnd.Grp.gen_iff}
\leanref{Development/WeilBound/ThetaLifts/EigenvalueEndgame/SignatureSplitting/StabClosed.lean}{PositivityTheory.ThKEnd.Two.lat_of_good}
\begin{proof}
Let $U^\perp$ be the subspace of the $z\in E^d$ with $h_W(z,u)=0$ for every
$u\in U$, and $e,f\in M_d(E)$ projections onto $U$ and onto $U^\perp$. As
$\mathbf A_E$ is flat over $E$, the $z\in\mathbf A_E^d$ with $h_W(z,u)=0$ for
every $u\in U$ form $\mathbf A_EU^\perp$, on which $f$ is the identity. Call
$X\in M_d(\mathbf A_E)$ \emph{triangular} if $Xe=0$, $eXf=Xf$ and $fX=X$.

(1) Let $g\in U_W$ fix $U$ and map every $z\in U^\perp$ into
$z+\mathbf A_EU$. Then $(g-1)e=0$. For $w\in\mathbf A_E^d$ and $u\in U$,
$h_W((g-1)w,u)=h_W(gw,gu)-h_W(w,u)=0$, so the image of $g-1$ lies in
$\mathbf A_EU^\perp$ and $f(g-1)=g-1$. The columns of $(g-1)f$ are the images
under $g-1$ of the columns of $f$, which lie in $U^\perp$, so they lie in
$\mathbf A_EU$, and $e(g-1)f=(g-1)f$. So $g-1$ is triangular. Conversely, if
$g-1$ is triangular, then $g$ fixes $U$, and for $z\in\mathbf A_EU^\perp$,
$(g-1)z=(g-1)fz=e(g-1)fz\in\mathbf A_EU$. Sums and products of triangular
matrices are triangular, and products of three of them vanish, as in (c) of the
proof of Lemma~\ref{asm-theta-agree-end-unipotent}. So if $g-1$ and $g'-1$ are
triangular, then so are $gg'-1=(g-1)+(g'-1)+(g-1)(g'-1)$ and
$g^{-1}-1=-(g-1)+(g-1)^2$. Hence the $g\in U_W$ with $g-1$ triangular form a
group, which contains the generators of $N_U$ and consists of generators. So it
is $N_U$, and (1) follows.

(2) The conditions of triangularity are closed, so $N_U$ is closed in $U_W$,
which is locally compact, second countable and Hausdorff
(Theorem~\ref{asm-theta-agree-convergence-siegel}(1)); and $N_U\cap U_W(F)$ is
countable and discrete, as $U_W(F)$ is. The construction of a compact set $C$ in
the proof of Lemma~\ref{asm-theta-agree-end-unipotent}(1) uses of $e$ and $f$
only the definition of triangular matrices and (c) and (d) there, so it gives a
compact $C\subset N_U$ such that every $u\in N_U$ has $\gamma u\in C$ for some
$\gamma\in N_U\cap U_W(F)$. So $N_U\cap U_W(F)$ is a lattice in $N_U$, and
Lemma~\ref{asm-theta-agree-end-kill-averages}(1) gives the rest.
\end{proof}
% Lean: unipUGood (ContinuousExpG.lean): unipU_eq with exists_proj and perp, gen_iff (59 lines)
% for (1), and unipGood_of triG_closed triG_cov for (2); lat_of_good gives the lattice. The
% A_E-form of (1) is gen_iff read through Tri; the Lean states (1) for z in E^d. Lemma
% asm-theta-agree-end-unipotent(1) is the case U = E x (unipLine).

\begin{lemma}\label{asm-theta-agree-end-kill-levi}
Let $E$ be a CM field with complex conjugation $c$ and a theta base $B$
(Definition~\ref{asm-theta-agree-models}(1)), let $W$ be a skew-Hermitian
lattice space of rank $d$ and some signature, let $H=(x,y,\iota)$ be a plane in
$W$ over a lattice space $W_1$ of rank $d_1$
(Definition~\ref{asm-theta-agree-end-notions}(1)), and let $S_x$ be the group of
the $g\in U_W$ with $gx=x$.
\begin{enumerate}
\item For $t\in\mathbf A_E^\times$ and $h_1\in U_{W_1}$ there is at most one Levi
  element for $(t,h_1)$. If $m$ and $m'$ are Levi elements for $(t,h_1)$ and
  $(t',h'_1)$, then $mm'$ is a Levi element for $(tt',h_1h'_1)$.
\item There are continuous homomorphisms $\lambda_H\colon U_{W_1}\to U_W$ and
  $r_H\colon S_x\to U_{W_1}$ such that $\lambda_H(h_1)$ is a Levi element for
  $(1,h_1)$ and $r_H(\lambda_H(h_1))=h_1$ for every $h_1\in U_{W_1}$;
  $\lambda_H(U_{W_1}(F))\subset U_W(F)$ and
  $r_H(S_x\cap U_W(F))\subset U_{W_1}(F)$; and
  $g\iota z-\iota r_H(g)z\in\mathbf A_Ex$ for all $g\in S_x$ and
  $z\in\mathbf A_E^{d_1}$.
\end{enumerate}
\end{lemma}
\leanref{Development/WeilBound/ThetaLifts/EigenvalueEndgame/SignatureSplitting/StabClosed.lean}{PositivityTheory.ThKEnd.Two.levi_uniq}
\leanref{Development/WeilBound/ThetaLifts/EigenvalueEndgame/SignatureSplitting/StabClosed.lean}{PositivityTheory.ThKEnd.Two.levi_mul}
\leanref{Development/WeilBound/ThetaLifts/EigenvalueEndgame/SignatureSplitting/StabClosed.lean}{PositivityTheory.ThKEnd.Two.levi_data}
\begin{proof}
Let $A=(\iota\ x\ y)$; as in the proof of
Lemma~\ref{asm-theta-agree-end-unipotent}, $A$ is invertible, with
$A^{\mathsf T}J_Wc(A)=\operatorname{diag}\bigl(J_1,\left(\begin{smallmatrix}0&1\\-1&0\end{smallmatrix}\right)\bigr)$
for the matrix $J_1$ of $W_1$. Let $Q\in M_{d_1\times d}(E)$ consist of the rows
of $A^{-1}$ for the block of $\iota$, so that $Q\iota=1$, $Qx=0$ and $Qy=0$. Two
matrices $M$ and $M'$ over $\mathbf A_E$ with $d$ columns and $M\iota=M'\iota$,
$Mx=M'x$ and $My=M'y$ are equal, as $MA=M'A$.

(1) Two Levi elements for $(t,h_1)$ agree on $\iota$, $x$ and $y$, so they are
equal. And $mm'x=tt'x$, $mm'y=c(t)^{-1}c(t')^{-1}y=c(tt')^{-1}y$ and
$mm'\iota=m\iota h'_1=\iota h_1h'_1$.

(2) For $h_1\in U_{W_1}$ let $\lambda_H(h_1)$ be the Levi element for $(1,h_1)$,
which exists by Lemma~\ref{asm-theta-agree-end-unipotent}(2) and is unique by
(1); by (1), $\lambda_H$ is a homomorphism. The matrix $\iota h_1Q+1-\iota Q$
agrees with $\lambda_H(h_1)$ on $\iota$, $x$ and $y$, so it is $\lambda_H(h_1)$.
So $h_1\mapsto\lambda_H(h_1)$ and $h_1\mapsto\lambda_H(h_1)^{-1}=\lambda_H(h_1^{-1})$
are continuous, and $\lambda_H(h_1)$ has entries in $E$ when $h_1$ has.

Let $g\in S_x$. The computation of Step 5 of the proof of
Lemma~\ref{asm-theta-agree-end-iwasawa}, for $p=g$ and $t=1$, uses only the
invertibility of $A$ and the Gram matrix above; it gives
$g\iota=\iota g_1+x\alpha^{\mathsf T}$ with $g_1^{\mathsf T}J_1c(g_1)=J_1$.
Then $Qg\iota=g_1$; put $r_H(g)=Qg\iota$. For $g,g'\in S_x$,
$gg'\iota=g(\iota r_H(g')+x\alpha'^{\mathsf T})=g\iota\,r_H(g')+x\alpha'^{\mathsf T}$,
so $r_H(gg')=r_H(g)r_H(g')$. As $r_H(1)=Q\iota=1$, $r_H(g)$ is invertible, so it
lies in $U_{W_1}$. The map $r_H$ is continuous, $r_H(g)$ has entries in $E$ when
$g$ has, and $g\iota z-\iota r_H(g)z=(\alpha\cdot z)\,x$. Finally
$r_H(\lambda_H(h_1))=Q\iota h_1=h_1$.
\end{proof}
% Lean: levi_uniq and levi_mul for (1); levi_data (StabClosed.lean, 47 lines) for (2), with
% exists_Q (fr_inv), fr_ext, leviPlane for the Levi elements, the formula e, cont_hom,
% isFr_plane and IsFr.stab for g iota = iota g_1 + x alpha^T, mu for the multiplicativity,
% and rat_iff; r_H(lambda_H(h_1)) = h_1 is l-rho in the proof of two_orders. Lemma
% asm-theta-agree-end-two-step takes its Levi map from (2), as the Lean does.

\begin{lemma}\label{asm-theta-agree-end-two-step-kill}
Let $E$ be a CM field with complex conjugation $c$ and a theta base $B$
(Definition~\ref{asm-theta-agree-models}(1)), let $W$ and $W_1$ be
skew-Hermitian lattice spaces of some signatures
(Definition~\ref{asm-theta-agree-doubling}(1)), with $W$ of rank $d$, matrix
$J_W$ and $h_W(u,u')=u^{\mathsf T}J_W\,c(u')$, let $H=(x,y,\iota)$ be a plane in
$W$ over $W_1$, and $H_1=(x_1,y_1,\iota_1)$ a plane in $W_1$ over a lattice
space $W_2$ (Definition~\ref{asm-theta-agree-end-notions}(1)). Let
$X=(x\ \ \iota x_1)$ and $Y=(y\ \ \iota y_1)$ be the $d\times2$ matrices with
these columns, and for $g\in SL_2(\mathbf A_E)$ and $z\in\mathbf A_E^d$ put
$g'=c(g^{-1})^{\mathsf T}$ and
\[
\varepsilon(g)z=z+X(g-1)\begin{pmatrix}h_W(z,y)\\ h_W(z,\iota y_1)\end{pmatrix}-Y(g'-1)\begin{pmatrix}h_W(z,x)\\ h_W(z,\iota x_1)\end{pmatrix}.
\]
\begin{enumerate}
\item The map $\varepsilon$ is a continuous homomorphism from
  $SL_2(\mathbf A_E)$ to $U_W$, with $\varepsilon(g)X=Xg$ and
  $\varepsilon(g)Y=Yg'$. For $t\in\mathbf A_E^\times$ and a Levi element
  $m\in U_{W_1}$ along $H_1$ for $(t^{-1},1)$, $\varepsilon(\operatorname{diag}(t,t^{-1}))$
  is a Levi element along $H$ for $(t,m)$.
\item Let $f$ be an automorphic function on $U_W$
  (Definition~\ref{asm-theta-agree-end-notions}(3)) with
  $|f(h)|\leq C\operatorname{ht}(h)^k$ for some $C$, $k$ and every $h$
  (Definition~\ref{asm-theta-agree-parts}), let $c_0$ be the constant term of
  $f$ along $H$, and let $c_1$ be a function on $U_W$ such that
  $c_1(g)=\int_{\mathcal F}c_0(\lambda(n)g)\,d\nu(n)$ for every $g\in U_W$, every
  map $\lambda\colon U_{W_1}\to U_W$ with $\lambda(y')$ a Levi element along
  $H$ for $(1,y')$ for every $y'$, every Haar measure $\nu$ on $N_{H_1}$ and
  every fundamental domain $\mathcal F$ for $N_{H_1}\cap U_{W_1}(F)$ in
  $N_{H_1}$ with $\nu(\mathcal F)=1$. Let $u\in\mathbf A_E^\times$ be $2$ at
  the infinite places and $1$ at the finite places, put
  $a=\varepsilon(\operatorname{diag}(u,u^{-1}))$, and let $l\in\mathbf C$ with
  $|l|<1$. If $c_1(a^2g)-c_1(ag)=l\,(c_1(ag)-c_1(g))$ for every $g\in U_W$,
  then $c_1(ag)=c_1(g)$ for every $g\in U_W$.
\end{enumerate}
\end{lemma}
\leanref{Development/WeilBound/ThetaLifts/EigenvalueEndgame/SignatureSplitting/Emb2DiaMulVecIncVY.lean}{PositivityTheory.ThKEnd.Two.emb2}
\leanref{Development/WeilBound/ThetaLifts/EigenvalueEndgame/SignatureSplitting/EmbMMul.lean}{PositivityTheory.ThKEnd.Two.Dual.emb}
\leanref{Development/WeilBound/ThetaLifts/EigenvalueEndgame/SignatureSplitting/Emb2DiaMulVecIncVY.lean}{PositivityTheory.ThKEnd.Two.emb2_dia_isLevi}
\leanref{Development/WeilBound/ThetaLifts/EigenvalueEndgame/SignatureSplitting/SpanPairIff.lean}{PositivityTheory.ThKEnd.Two.two_kill}
\begin{proof}
(1) As $W$ and $W_1$ are skew-Hermitian, $h_W(u',u)=-c(h_W(u,u'))$, and the
form of $W_1$ is $h_1(v,v')=h_W(\iota v,\iota v')$. The definition of a plane
gives $X^{\mathsf T}J_Wc(X)=Y^{\mathsf T}J_Wc(Y)=0$ and $X^{\mathsf T}J_Wc(Y)=1$,
hence $Y^{\mathsf T}J_Wc(X)=-1$. With $\varepsilon(g)=1+X(g-1)(J_Wc(Y))^{\mathsf T}-Y(g'-1)(J_Wc(X))^{\mathsf T}$,
these relations give $\varepsilon(g)X=Xg$ and $\varepsilon(g)Y=Yg'$, and a
direct computation from them gives $\varepsilon(gg_2)=\varepsilon(g)\varepsilon(g_2)$
and $\varepsilon(g)^{\mathsf T}J_Wc(\varepsilon(g))=J_W$; as $\varepsilon(1)=1$,
$\varepsilon(g)$ is invertible, and it lies in $U_W$. It is continuous, as its
entries are polynomials in those of $g$, $g^{-1}$ and their conjugates.

Let $g=\operatorname{diag}(t,t^{-1})$, so $g'=\operatorname{diag}(c(t)^{-1},c(t))$.
Then $\varepsilon(g)x=tx$ and $\varepsilon(g)y=c(t)^{-1}y$. For
$v\in\mathbf A_E^{d-2}$, $h_W(\iota v,x)=h_W(\iota v,y)=0$, so
\[
\varepsilon(g)\iota v=\iota\bigl(v+(t^{-1}-1)h_1(v,y_1)x_1-(c(t)-1)h_1(v,x_1)y_1\bigr).
\]
Write $v=h_1(v,y_1)x_1-h_1(v,x_1)y_1+\iota_1w$, which is possible as
$h_1(x_1,y_1)=1$, $h_1(y_1,x_1)=-1$ and $x_1$, $y_1$ and the columns of $\iota_1$
form a basis (proof of Lemma~\ref{asm-theta-agree-end-unipotent}). As
$mx_1=t^{-1}x_1$, $my_1=c(t)y_1$ and $m\iota_1=\iota_1$, the bracket is $mv$.
So $\varepsilon(g)\iota=\iota m$, and $\varepsilon(g)$ is a Levi element for
$(t,m)$.

(2) Write $\Gamma=U_W(F)$, $U=Ex+E\iota x_1\subset E^d$, and $N_U$ for the group
of Lemma~\ref{asm-theta-agree-end-kill-unipotent}. Let $a_0$ and $n(\xi)$ be as
in Lemma~\ref{asm-theta-agree-end-kill-siegel}, so that $a=\varepsilon(a_0)$,
and let $S_x$, $\lambda_H$ and $r_H$ be as in
Lemma~\ref{asm-theta-agree-end-kill-levi} for $H$, with the matrix $Q$ of its
proof ($Q\iota=1$, $Qx=0$). The groups $U_W$, $U_{W_1}$, $SL_2(\mathbf A_E)$ and
$\mathbf A_E$ are locally compact, second countable and Hausdorff
(Theorem~\ref{asm-theta-agree-convergence-siegel}(1) and
Lemma~\ref{asm-theta-agree-end-kill-siegel}(1)), and $\Gamma$ is countable and
discrete in $U_W$.

\emph{Step 1: the embedding.} For $\xi\in\mathbf A_E$, the only nonzero entry
of $n(\xi)-1$ is $\xi$ at $(1,2)$, and that of $n(\xi)'-1$ is $-c(\xi)$ at
$(2,1)$. So the formula for $\varepsilon$ reads
\[
\varepsilon(n(\xi))w=w+\xi\,h_W(w,\iota y_1)\,x+c(\xi)\,h_W(w,x)\,\iota y_1 .
\]
Hence $\varepsilon(n(\xi))x=x$, $\varepsilon(n(\xi))\iota x_1=\iota x_1+\xi x$
and $\varepsilon(n(\xi))\iota z=\iota z+\xi h_1(z,y_1)x$ for
$z\in\mathbf A_E^{d-2}$. For $g\in SL_2(E)$ the entries of $\varepsilon(g)$ lie
in $E$, so $\varepsilon(g)\in\Gamma$; and $\varepsilon(g)X=Xg$ and
$\varepsilon(g)^{-1}X=\varepsilon(g^{-1})X=Xg^{-1}$, so $\varepsilon(g)$ and
$\varepsilon(g)^{-1}$ map $U$ onto $U$.

\emph{Step 2: the average along $U$.} By
Lemma~\ref{asm-theta-agree-end-kill-unipotent}(2), $N_U$ is closed and
$\Gamma\cap N_U$ is a lattice in it, so by
Lemma~\ref{asm-theta-agree-end-kill-averages}(2) and~(3), $f$ has an average
$\varphi$ over $\Gamma\cap N_U$, and $\varphi$ is continuous. Let
$g\in SL_2(E)$ and $\gamma=\varepsilon(g)$. Then $\gamma^{-1}N_U\gamma=N_U$. Indeed,
let $m\in N_U$. Then $\gamma^{-1}m\gamma$ fixes $U$, as $\gamma$ maps $U$ onto
$U$. If $z\in E^d$ has $h_W(z,u)=0$ for $u\in U$, then $\gamma z\in E^d$ has
$h_W(\gamma z,u)=h_W(z,\gamma^{-1}u)=0$, so
$(\gamma^{-1}m\gamma-1)z=\gamma^{-1}(m-1)\gamma z\in\gamma^{-1}\mathbf A_EU=\mathbf A_EU$.
By Lemma~\ref{asm-theta-agree-end-kill-unipotent}(1),
$\gamma^{-1}m\gamma\in N_U$, and in the same way $\gamma m\gamma^{-1}\in N_U$. As
$f$ is automorphic, Lemma~\ref{asm-theta-agree-end-kill-averages}(4) gives
$\varphi(\gamma h)=\varphi(h)$ for every $h$. In particular
$\varphi(\varepsilon(n(\xi))h)=\varphi(h)$ for $\xi\in E$. By
Lemma~\ref{asm-theta-agree-end-kill-siegel}(1), $E$ is a lattice in
$\mathbf A_E$, and $\xi\mapsto\varepsilon(n(\xi))$ is a continuous
homomorphism; so by Lemma~\ref{asm-theta-agree-end-kill-averages}(2) and~(3),
$\varphi$ has an average $c'$ over $E$ along it, and $c'$ is continuous.

\emph{Step 3: two orders of integration.} We show that $c_1=c'$. Let $S$ be the
group of the $g\in S_x$ with $r_H(g)\in N_{H_1}$. It is closed in $U_W$, as
$S_x$ is closed, $r_H$ is continuous and $N_{H_1}$ is closed in $U_{W_1}$
(Lemma~\ref{asm-theta-agree-end-unipotent}(1)).

(i) For $g\in S_x$, $r_H(g)=1$ if and only if $g\in N_H$. If $r_H(g)=1$, then
$g\iota z\in\iota z+\mathbf A_Ex$ for every $z$, so $g$ is one of the generators
of $N_H$. If $g\in N_H$, then $gx=x$ and $g\iota z-\iota z\in\mathbf A_Ex$
(Lemma~\ref{asm-theta-agree-end-unipotent}(3)), so
$\iota(r_H(g)z-z)\in\mathbf A_Ex$, and applying $Q$ gives $r_H(g)z=z$. In
particular $N_H\subset S$.

(ii) For $g\in S_x$, $g\in N_U$ if and only if $r_H(g)\in N_{H_1}$ and
$g\iota x_1=\iota x_1$. Write $g\iota z=\iota r_H(g)z+\alpha(z)x$. Let
$g\in N_U$. Then $g\iota x_1=\iota x_1$, so $\iota(r_H(g)x_1-x_1)\in\mathbf A_Ex$,
and $r_H(g)x_1=x_1$ after applying $Q$. Let $z\in\mathbf A_E^{d-2}$ have
$h_1(z,x_1)=0$. Then $h_W(\iota z,x)=0$ and $h_W(\iota z,\iota x_1)=0$, so
$g\iota z-\iota z=\beta x+\beta'\iota x_1$ for some $\beta,\beta'$, by
Lemma~\ref{asm-theta-agree-end-kill-unipotent}(1); applying $Q$ gives
$r_H(g)z-z=\beta'x_1$. By Lemma~\ref{asm-theta-agree-end-unipotent}(3) for
$H_1$, $r_H(g)\in N_{H_1}$. Conversely, let $r_H(g)\in N_{H_1}$ and
$g\iota x_1=\iota x_1$. Then $g$ fixes $x$ and $\iota x_1$, hence $U$. Let
$z\in\mathbf A_E^d$ have $h_W(z,x)=h_W(z,\iota x_1)=0$. By (a) of the proof of
Lemma~\ref{asm-theta-agree-end-unipotent}, $z=\iota w+sx$, and
$h_1(w,x_1)=h_W(z,\iota x_1)=0$. By Lemma~\ref{asm-theta-agree-end-unipotent}(3)
for $H_1$, $r_H(g)w=w+\beta x_1$, so
$gz=\iota r_H(g)w+\alpha(w)x+sx=z+\beta\iota x_1+\alpha(w)x$. By
Lemma~\ref{asm-theta-agree-end-kill-unipotent}(1), $g\in N_U$. In particular
$N_U\subset S$.

(iii) For $g\in S$, $g\iota x_1=\iota x_1+\pi(g)x$ with
$\pi(g)=h_W(g\iota x_1,y)$, as $r_H(g)x_1=x_1$, $h_W(\iota x_1,y)=0$ and
$h_W(x,y)=1$. As $gx=x$, $\pi(gg')=h_W(g(\iota x_1+\pi(g')x),y)=\pi(g)+\pi(g')$:
$\pi$ is a continuous homomorphism from $S$ to $\mathbf A_E$, and by (ii) its
kernel is $N_U$. By Step 1, $\varepsilon(n(\xi))$ fixes $x$ and maps every
$\iota z$ into $\iota z+\mathbf A_Ex$, so it lies in $N_H$; by (i) it lies in
$S$, and $\pi(\varepsilon(n(\xi)))=\xi$. For $\xi\in E$, $\varepsilon(n(\xi))\in\Gamma$,
and $\pi(g)\in E$ for $g\in S\cap\Gamma$.

(iv) Apply Lemma~\ref{asm-theta-agree-end-kill-stages} to $G=U_W$, $\Gamma$, $S$,
$A=N_H$, $B=N_{H_1}$ with the lattice $\Gamma_B=N_{H_1}\cap U_{W_1}(F)$, the
restriction of $\lambda_H$ to $N_{H_1}$, with values in $S$ as
$\lambda_H(b)x=x$ and $r_H(\lambda_H(b))=b$, and $\pi=r_H$ on $S$. The kernel is
$N_H$ by (i), $\lambda_H(\Gamma_B)\subset\Gamma$ and $r_H(\Gamma\cap S)\subset\Gamma_B$
by Lemma~\ref{asm-theta-agree-end-kill-levi}(2), and $\Gamma\cap N_H$ and
$\Gamma_B$ are lattices (Lemma~\ref{asm-theta-agree-end-unipotent}(1) for $H$
and $H_1$). So $\Gamma\cap S$ is a lattice in $S$. The function $f$ is
continuous and invariant under $\Gamma$, $c_0$ is its average over $\Gamma\cap N_H$
(Definition~\ref{asm-theta-agree-end-kill-notions}(2)), and $c_1$ is the
average of $c_0$ over $\Gamma_B$ along $\lambda_H$, by hypothesis. So $c_1$ is
the average of $f$ over $\Gamma\cap S$.

(v) Apply Lemma~\ref{asm-theta-agree-end-kill-stages} again, to $A=N_U$,
$B=\mathbf A_E$ with the lattice $E$, $\xi\mapsto\varepsilon(n(\xi))$ and $\pi$,
using (iii) and Step 2. So $c'$ is the average of $f$ over $\Gamma\cap S$. By
Lemma~\ref{asm-theta-agree-end-kill-averages}(2), $c_1=c'$.

\emph{Step 4: growth.} The function $f$ is moderate in the sense of
Lemma~\ref{asm-theta-agree-end-kill-growth}, so by part (2) of that lemma
$\varphi$ is moderate, and by part (1), for the continuous homomorphism
$\varepsilon$, for every $h\in U_W$ the function $g\mapsto\varphi(\varepsilon(g)h)$
has growth of some order along $a_0$.

\emph{Step 5: conclusion.} Let $\lambda$ be the probability measure on
$\mathbf A_E$ of Lemma~\ref{asm-theta-agree-end-kill-averages}(2) for the
lattice $E$, and fix $h\in U_W$. The pair of functions
$g\mapsto\varphi(\varepsilon(g)h)$ and $g\mapsto c_1(\varepsilon(g)h)$ on
$SL_2(\mathbf A_E)$ is a recursion pair for $(SL_2(E),n,a_0,\lambda,l)$
(Definition~\ref{asm-theta-agree-end-kill-notions}(5)): both are continuous, by
Step 2 and Step 3; the first is invariant under $SL_2(E)$ by Step 2; by Step 3
and Lemma~\ref{asm-theta-agree-end-kill-averages}(2),
$c_1(\varepsilon(g)h)=c'(\varepsilon(g)h)=\int\varphi(\varepsilon(n(\xi)g)h)\,d\lambda(\xi)$;
and the recursion is the hypothesis at $\varepsilon(g)h$, as
$\varepsilon(a_0^2g)=a^2\varepsilon(g)$. By
Lemma~\ref{asm-theta-agree-end-kill-siegel}(2) the triple $(SL_2(E),n,a_0)$ has
the Siegel and Lipschitz properties, $n$ is continuous and $|l|<1$. With Step 4,
Lemma~\ref{asm-theta-agree-end-kill-induction}(4) gives
$c_1(\varepsilon(a_0g)h)=c_1(\varepsilon(g)h)$ for every $g$; for $g=1$ this is
$c_1(ah)=c_1(h)$.
\end{proof}
% Lean: two_kill (SpanPairIff.lean, 51 lines). The setup is lcsc_UA, ae_facts and sl_facts;
% Step 1 is emb2_mulVec_incV_x_and_complement_x, emb2_mulVec_complement_eq,
% emb2_map_inclR and emb2_norm; Step 2 is phi (isAv_ex with lat_of_good G2, G2 =
% unipUGood), IsAv.cont, phi_i (IsAv.conj_sub with conj_unipU) and c' (LR.isAv_ex); Step 3
% is two_orders (92 lines): S, k0 (i), mem_pair (51 lines) (ii), act and pi (iii), the first
% stage_sub with levi_data (iv), the second (v) and uniq_sub; Step 4 is Mod.av and Mod.gr;
% Step 5 is isAv_prob, the record D of Dat, and kill.

\begin{lemma}\label{asm-theta-agree-end-two-step}
Let $B$ be a theta base of a CM field $E$, and let
$(B,G,\mathcal X,\omega,\omega')$ and $(B,G_1,\mathcal X,\omega_1,\omega'_1)$ be
triples (Definition~\ref{asm-theta-agree-triple}), with $G$ of ranks $(b+2,a')$
and $G_1$ of ranks $(b,a')$, such that $(G,G_1,H,H')$ is a plane-removal
configuration (Definition~\ref{asm-theta-agree-end-notions}(7)) and $q\geq1$.
Let $\Sigma$ and $\Sigma_1$ be the sets of differences of the two triples. Let
$f\in\Sigma$ and let $f_\alpha$ be such that $(f_\alpha,0)$ splits $f$ along
$H'$ (Definition~\ref{asm-theta-agree-planes}(3)), and suppose that for every
$k\in K^\theta_W$ some element of $\Sigma_1$ is the lower part of $f_\alpha$ at
$k$ (Definition~\ref{asm-theta-agree-end-notions}(6)). Then every $f_1\in\Sigma_1$
that is the lower part of $f_\alpha$ at some $k\in K^\theta_W$ is pure of law
$(\alpha,a'/2)$ (Definition~\ref{asm-theta-agree-end-notions}(3)).
\end{lemma}
\leanref{Development/WeilBound/ThetaLifts/EigenvalueEndgame/PlaneRemovalConfiguration.lean}{PositivityTheory.ThKEnd.TwoStep}
\leanref{Development/WeilBound/ThetaLifts/EigenvalueEndgame/SignatureSplitting/TwoStepOf.lean}{PositivityTheory.ThKEnd.Two.twoStep_of'}
\begin{proof}
Let $G_1=(p_1,q,V_1,V',W_1,M_1,M'_1)$, with $p_1+1=p$, and let $k\in K^\theta_W$
and $f_1\in\Sigma_1$ be the lower part of $f_\alpha$ at $k$. Write $|t|$ for
$|N_{E/F}(t)|_{\mathbf A}$.

\emph{Step 1: the data.} By Lemma~\ref{asm-theta-agree-end-iwasawa}, $H'$
gives a plane in $W$ over $W_1$, and every $g\in U_W$ is $u\,m\,l\,\kappa$ with
$u\in N_{H'}$, $m$ a Levi element for $(s,1)$, $l$ one for $(1,y)$ and
$\kappa\in K^\theta_W$. As $p_1+q\geq1$,
Lemma~\ref{asm-theta-agree-end-two-step-split} for the triple of $G_1$ gives a
lattice space $W_2$ and a plane $H_1$ in $W_1$ over $W_2$ such that every element
of $\Sigma_1$ is split along $H_1$, with the laws $(\alpha,a'/2)$ and
$(\beta,b/2)$. By Lemma~\ref{asm-theta-agree-end-automorphic}, $f$ is
automorphic, and $f_\alpha=f_\alpha+0$ is its constant term along $H'$. By
Lemma~\ref{asm-theta-agree-end-kill-levi}(2) for the plane $H'$ there is a
continuous homomorphism $\lambda_1\colon U_{W_1}\to U_W$ such that
$\lambda_1(y)$ is a Levi element for $(1,y)$, and by part (1) of that lemma a
Levi element for $(t,y)$ is unique and products of Levi elements are Levi
elements. If $m$ is a Levi element for $(s,1)$, then $\lambda_1(y)m$ and
$m\lambda_1(y)$ are Levi elements for $(s,y)$, so
$\lambda_1(y)m=m\lambda_1(y)$. If
$g_1\in\Sigma_1$ is the lower part of $f_\alpha$ at $\kappa$, then
$f_\alpha(\lambda_1(z)\lambda_1(y)\kappa)=f_\alpha(\lambda_1(zy)\kappa)=g_1(zy)$.

\emph{Step 2: a splitting for every $g$.} For $g\in U_W$ let $F_g$ be the
function $w\mapsto f_\alpha(\lambda_1(w)g)$ on $U_{W_1}$. We show that $F_g$ is
split along $H_1$: that there are functions $A$ and $C$ on $U_{W_1}$ with $A+C$
the constant term of $F_g$ along $H_1$, $A$ of law $(\alpha,a'/2)$ and $C$ of
law $(\beta,b/2)$. Write $g=um\lambda_1(y)\kappa$ as in Step 1, with
$l=\lambda_1(y)$ by uniqueness. By hypothesis some $g_1\in\Sigma_1$ is the
lower part of $f_\alpha$ at $\kappa$; let $(A,C)$ split $g_1$ along $H_1$. Let
$z\in U_{W_1}$. By Lemma~\ref{asm-theta-agree-end-unipotent}(2),
$\lambda_1(z)u\lambda_1(z)^{-1}\in N_{H'}$, and $f_\alpha$ is invariant under
left translation by $N_{H'}$ (Lemma~\ref{asm-theta-agree-end-invariant}(2)). So,
by Step 1 and the law of $f_\alpha$,
\[
F_g(z)=f_\alpha\bigl(\lambda_1(z)u\lambda_1(z)^{-1}\cdot\lambda_1(z)m\lambda_1(y)\kappa\bigr)
=f_\alpha(m\lambda_1(z)\lambda_1(y)\kappa)=\alpha(s)|s|^{a'/2}g_1(zy).
\]
Right translation by $y$ and multiplication by a constant preserve constant
terms along $H_1$ and laws, so
$z\mapsto\alpha(s)|s|^{a'/2}A(zy)$ and $z\mapsto\alpha(s)|s|^{a'/2}C(zy)$ split
$F_g$ along $H_1$.

For every $g$ choose such a pair $(A_g,C_g)$, and put $c_1(g)=A_g(1)+C_g(1)$. By
Lemma~\ref{asm-theta-agree-end-unipotent}(1) there are a Haar measure $\nu$ on
$N_{H_1}$ and a fundamental domain $\mathcal F$ for $N_{H_1}\cap U_{W_1}(F)$ of
measure $1$. If $(A',C')$ is any pair that splits $F_g$ along $H_1$, then
$A'(1)+C'(1)=\int_{\mathcal F}F_g(n)\,d\nu(n)=c_1(g)$.

\emph{Step 3: the element $a$.} Let $u_0\in\mathbf A_E^\times$ be $2$ at the
infinite places and $1$ at the finite places, so that $|u_0|=4^{[F:\mathbf Q]}>1$.
By Lemma~\ref{asm-theta-agree-end-unipotent}(2) there are a Levi element
$m\in U_{W_1}$ along $H_1$ for $(u_0^{-1},1)$ and a Levi element $m_0$ along
$H'$ for $(u_0,1)$. Then $a=m_0\lambda_1(m)$ is a Levi element along $H'$ for
$(u_0,m)$, and so is $\varepsilon(\operatorname{diag}(u_0,u_0^{-1}))$, with the
embedding $\varepsilon$ of Lemma~\ref{asm-theta-agree-end-two-step-kill}(1) for
$H'$ and $H_1$. By uniqueness,
$a=\varepsilon(\operatorname{diag}(u_0,u_0^{-1}))$. Put
\[
\ell=\alpha(u_0)|u_0|^{a'/2}\,\beta(u_0^{-1})|u_0|^{-b/2}.
\]
Then $|\ell|=|u_0|^{(a'-b)/2}<1$, as $a'<b$.

\emph{Step 4: one step.} Let $(A',C')$ split $F_g$ along $H_1$. Then $F_{ag}$
has a splitting $(A'',C'')$ with $A''(1)=A'(1)$ and $C''(1)=\ell\,C'(1)$.
Indeed, by Step 1 and the law of $f_\alpha$,
$F_{ag}(z)=f_\alpha(m_0\lambda_1(zm)g)=\alpha(u_0)|u_0|^{a'/2}F_g(zm)$, so
$A''(z)=\alpha(u_0)|u_0|^{a'/2}A'(zm)$ and
$C''(z)=\alpha(u_0)|u_0|^{a'/2}C'(zm)$ split $F_{ag}$. As $m$ is a Levi element
for $(u_0^{-1},1)$, the laws of $A'$ and $C'$ give
$A''(1)=\alpha(u_0)|u_0|^{a'/2}\alpha(u_0^{-1})|u_0|^{-a'/2}A'(1)=A'(1)$ and
$C''(1)=\alpha(u_0)|u_0|^{a'/2}\beta(u_0^{-1})|u_0|^{-b/2}C'(1)=\ell\,C'(1)$.
Applying this twice to $(A_g,C_g)$ and using Step 2,
$c_1(ag)=A_g(1)+\ell C_g(1)$ and $c_1(a^2g)=A_g(1)+\ell^2C_g(1)$, so
\[
c_1(a^2g)-c_1(ag)=\ell\,(c_1(ag)-c_1(g)).
\]

\emph{Step 5: invariance.} By Lemma~\ref{asm-theta-agree-end-two-step-moderate},
$|f(h)|\leq C_f\operatorname{ht}(h)^{k_f}$ for some $C_f$ and $k_f$. A map
$\lambda'\colon U_{W_1}\to U_W$ with $\lambda'(y)$ a Levi element for $(1,y)$
for every $y$ is $\lambda_1$, and by Step 2
$c_1(g)=\int_{\mathcal F}f_\alpha(\lambda_1(n)g)\,d\nu(n)$ for every Haar
measure $\nu$ and fundamental domain $\mathcal F$ of measure $1$. So
Lemma~\ref{asm-theta-agree-end-two-step-kill}(2), for $f$, $c_0=f_\alpha$,
$c_1$, $\ell$ and $a$, gives $c_1(ag)=c_1(g)$ for every $g$.

\emph{Step 6: conclusion.} Let $(A',C')$ split $f_1$ along $H_1$ (Step 1), and
let $y\in U_{W_1}$. As $f_1$ is the lower part of $f_\alpha$ at $k$,
$F_{\lambda_1(y)k}(z)=f_1(zy)$ by Step 1, so $z\mapsto A'(zy)$ and
$z\mapsto C'(zy)$ split $F_{\lambda_1(y)k}$, and
$c_1(\lambda_1(y)k)=A'(y)+C'(y)$ by Step 2. By Step 4,
$c_1(a\lambda_1(y)k)=A'(y)+\ell C'(y)$. By Step 5 these are equal, so
$(\ell-1)C'(y)=0$, and $C'(y)=0$ as $\ell\neq1$. Hence $C'=0$, so $A'$ is the
constant term of $f_1$ along $H_1$, and it has the law $(\alpha,a'/2)$: $f_1$
is pure of law $(\alpha,a'/2)$.
\end{proof}
% Lean: Two.twoStep_of' SigMod SigSplit Iwa IAut (SignatureSplitting/TwoStepOf.lean). Step 1
% is Iwa, SigSplit (Lemma asm-theta-agree-end-two-step-split), sig_isAut, levi_data (lev, v1),
% levi_uniq, com and low; Step 2 is ex (ct_inv, leviPlane, split_move), choose and val
% (unipGood.ex); Step 3 is m, m_0, ea (levi_mul, emb2_dia_isLevi), r1 (absE_halfE) and l1;
% Step 4 is step and the law clause of key; Step 5 is key = two_kill with SigMod (Lemma
% asm-theta-agree-end-two-step-moderate); Step 6 is z. The Levi map is lev from levi_data
% for the plane H' (Lemma asm-theta-agree-end-kill-levi), as in the Lean.

\subsection{Cuspidality of differences}\label{subsection-asm-cuspidality}

\begin{definition}\label{asm-theta-agree-end-cusp-notions}
Let $E$ be a CM field with maximal totally real subfield $F$ and complex
conjugation $c$, and identify $\mathbf A_E=\mathbf A_F\otimes_FE$ with the ring
of adeles of $E$. Let $W$ be a skew-Hermitian lattice space of rank $d$
(Definition~\ref{asm-theta-agree-doubling}(1)) with matrix $J_W$, and put
$h_W(u,u')=u^{\mathsf T}J_W\,c(u')$. For a finite set $I$ fix a Haar measure
$\mu_I$ on $\mathbf A_E^I$ and a fundamental domain $D_I$ for $E^I$ with
$\mu_I(D_I)=1$, and write $\langle\Phi\rangle_I=\int_{D_I}\Phi\,d\mu_I$ for a
function $\Phi$ on $\mathbf A_E^I$, the integral of a function that is not
integrable being $0$. Let $a,j\geq1$.
\begin{enumerate}
\item A \emph{frame} of rank $j$ in $W$ is a pair $(u,u')$ of matrices in
  $M_{d\times j}(E)$ with $u^{\mathsf T}J_W\,c(u)=0$,
  $u'^{\mathsf T}J_W\,c(u')=0$ and $u^{\mathsf T}J_W\,c(u')=1$. Its
  \emph{span} $U_u\subset E^d$ is the $E$-span of the columns of $u$. For
  $1\leq r\leq j$, $u_{\leq r}$ is the matrix of the first $r$ columns of $u$.
\item The \emph{Levi map} of a frame $(u,u')$ of rank $j$ is the map
  $\lambda_u$ from $GL_j(\mathbf A_E)$ to $M_d(\mathbf A_E)$ with
  \[
  \lambda_u(g)=1+u(g-1)\bigl(J_W\,c(u')\bigr)^{\mathsf T}-u'\bigl(c(g^{-1})^{\mathsf T}-1\bigr)\bigl(J_W\,c(u)\bigr)^{\mathsf T}.
  \]
\item Let $U\subset E^d$ be a nonzero $E$-subspace with $h_W(U,U)=0$, and fix a
  Haar measure $\nu_U$ on $N_U$ and a fundamental domain $D_U$ for
  $N_U\cap U_W(F)$ with $\nu_U(D_U)=1$
  (Lemma~\ref{asm-theta-agree-end-kill-unipotent}(2)). The \emph{constant term}
  along $U$ of a function $f$ on $U_W$ is the function
  $c_Uf(h)=\int_{D_U}f(nh)\,d\nu_U(n)$.
\item For $1\leq r\leq m\leq a-1$, let $I_{r,m}$ be the set of pairs $(i,k)$
  with $i\leq r<k\leq m+1$, and for $y\in\mathbf A_E^{I_{r,m}}$ let
  $n_{r,m}(y)\in GL_a(\mathbf A_E)$ be the unipotent matrix with the entries
  $y_{ik}$ at the places $(i,k)\in I_{r,m}$ and $0$ at the other places off the
  diagonal. So $n_{k,k}(x)$, for $x\in\mathbf A_E^{I_{k,k}}$, has the entries
  $x_{1,k+1},\ldots,x_{k,k+1}$ above the diagonal in column $k+1$.
\item Let $\psi$ be a character of $\mathbf A_E$ and $\varphi$ a function on
  $GL_a(\mathbf A_E)$. For $1\leq k\leq a-1$ put
  $T_k\varphi(g)=\bigl\langle x\mapsto\psi(-x_{k,k+1})\varphi(n_{k,k}(x)g)\bigr\rangle_{I_{k,k}}$,
  and for $0\leq m\leq a-1$ put $\mathcal W^m_\psi\varphi=T_1T_2\cdots T_m\varphi$,
  so that $\mathcal W^0_\psi\varphi=\varphi$.
\item A \emph{standard character} of $\mathbf A_E$ is a continuous map $\psi$
  from $\mathbf A_E$ to the complex numbers of absolute value $1$ with
  $\psi(x+y)=\psi(x)\psi(y)$ and $\psi(E)=1$, such that every $x\in\mathbf A_E$
  with $\psi(xe)=1$ for all $e\in E$ lies in $E$, and such that for every
  finite place $w$ of $E$ some $y\in E_w\subset\mathbf A_E$ has $\psi(y)\neq1$.
\item For a finite place $v$ of $F$ and an $F$-embedding $t\colon E\to F_v$,
  let $\pi_t\colon\mathbf A_E\to F_v$ be the projection to $E\otimes_FF_v$
  followed by the map induced by $t$. A character $\chi$ of $\mathbf A_E$ (a
  homomorphism to $\mathbf C^\times$) is \emph{nondegenerate} if it is
  continuous and for all such $v$ and $t$ there is $z\in\mathbf A_E$ with
  $\chi(z)\neq1$, $\pi_{t\circ c}(z)=0$, and components $0$ at the places of $E$
  not above $v$.
\item Let $(u,u')$ be a frame of rank $j$ and $\chi$ a character of
  $\mathbf A_E$. A \emph{Whittaker functional} for $(u,\chi)$ is a map $\Lambda$
  from the functions on $U_W$ to $\mathbf C$ such that
  $\Lambda(f+f')=\Lambda(f)+\Lambda(f')$ for automorphic $f$ and $f'$
  (Definition~\ref{asm-theta-agree-end-notions}(3));
  $\Lambda(zf)=z\Lambda(f)$ for every $z\in\mathbf C$ and every function $f$;
  and $\Lambda(f(\cdot\,h))=\chi\bigl(\sum_{i=1}^{j-1}g_{i,i+1}\bigr)\Lambda(f)$
  for every automorphic $f$, every $h\in U_W$ and every upper triangular
  $g\in M_j(\mathbf A_E)$ with diagonal entries $1$ such that $hu=ug$ and
  $hz-z\in\mathbf A_EU_u$ for every $z\in E^d$ with $h_W(z,U_u)=0$. Here
  $f(\cdot\,h)$ is the function $x\mapsto f(xh)$.
\item Let $(u,u')$ be a frame of rank $j$ and $\psi$ a character of
  $\mathbf A_E$. The \emph{Whittaker coefficient} along $(u,\psi)$ of a
  function $f$ on $U_W$ is
  $\mathcal W_{u,\psi}(f)=\mathcal W^{j-1}_\psi\varphi_f(1)$, where
  $\varphi_f(g)=c_{U_u}f(\lambda_u(g))$ for $g\in GL_j(\mathbf A_E)$.
\end{enumerate}
\end{definition}
\leanref{Development/WeilBound/ThetaLifts/EigenvalueEndgame/CuspidalAverages/KeyThree.lean}{PositivityTheory.ThKEnd.Cusp.IsFrame}
\leanref{Development/WeilBound/ThetaLifts/EigenvalueEndgame/CuspidalAverages/FirstIdempotentAdjointThird.lean}{PositivityTheory.ThKEnd.Cusp.Fr}
\leanref{Development/WeilBound/ThetaLifts/EigenvalueEndgame/CuspidalAverages/KeyThree.lean}{PositivityTheory.ThKEnd.Cusp.colSp}
\leanref{Development/WeilBound/ThetaLifts/EigenvalueEndgame/CuspidalAverages/KeyThree.lean}{PositivityTheory.ThKEnd.Cusp.levM}
\leanref{Development/WeilBound/ThetaLifts/EigenvalueEndgame/CuspidalAverages/KeyThree.lean}{PositivityTheory.ThKEnd.Cusp.IsLev}
\leanref{Development/WeilBound/ThetaLifts/EigenvalueEndgame/CuspidalAverages/IsAvLeft.lean}{PositivityTheory.ThKEnd.Cusp.IsCt}
\leanref{Development/WeilBound/ThetaLifts/EigenvalueEndgame/CuspidalAverages/IsAvLeft.lean}{PositivityTheory.ThKEnd.Cusp.Ctx.ct}
\leanref{Development/AbelianVarietyTheory/Proofs/InertiaEigenvalues/UnipotentBlockColumnIndex.lean}{GlobalTheory.CuspGeneric.avg}
\leanref{Development/AbelianVarietyTheory/Proofs/InertiaEigenvalues/UnipotentBlockColumnIndex.lean}{GlobalTheory.CuspGeneric.blkU}
\leanref{Development/AbelianVarietyTheory/Proofs/InertiaEigenvalues/UnipotentBlockColumnIndex.lean}{GlobalTheory.CuspGeneric.colT}
\leanref{Development/AbelianVarietyTheory/Proofs/InertiaEigenvalues/UnipotentBlockColumnIndex.lean}{GlobalTheory.CuspGeneric.whit}
\leanref{Development/AbelianVarietyTheory/Proofs/InertiaEigenvalues/UnipotentBlockColumnIndex.lean}{GlobalTheory.CuspGeneric.StdPsi}
\leanref{Development/WeilBound/ThetaLifts/EigenvalueEndgame/CuspidalAverages/KeyThree.lean}{PositivityTheory.ThKEnd.Cusp.IsChar}
\leanref{Development/WeilBound/ThetaLifts/EigenvalueEndgame/CuspidalAverages/KeyThree.lean}{PositivityTheory.ThKEnd.Cusp.IsWh}
\leanref{Development/WeilBound/ThetaLifts/EigenvalueEndgame/CuspidalAverages/CuspOfLineOfKeys.lean}{PositivityTheory.ThKEnd.Cusp.wh}

\begin{lemma}\label{asm-theta-agree-end-cusp-frames}
Let $E$ be a CM field with a theta base $B$
(Definition~\ref{asm-theta-agree-models}(1)), let $W$ be a skew-Hermitian
lattice space of rank $d$ and some signature, with the notation of
Definition~\ref{asm-theta-agree-end-cusp-notions}, and let $(u,u')$ be a frame
of rank $j$ in $W$ with span $U$.
\begin{enumerate}
\item The Levi map $\lambda_u$ is a continuous homomorphism from
  $GL_j(\mathbf A_E)$ to $U_W$, with $\lambda_u(g)u=ug$, and
  $\lambda_u(GL_j(E))\subset U_W(F)$.
\item For $g\in GL_j(\mathbf A_E)$ and $n\in N_U$,
  $\lambda_u(g)\,n\,\lambda_u(g)^{-1}\in N_U$.
\item Let $h\in U_W$ and $g\in GL_j(\mathbf A_E)$ be such that $hu=ug$ and
  $hz-z\in\mathbf A_EU$ for every $z\in E^d$ with $h_W(z,U)=0$. Then
  $h\lambda_u(g)^{-1}\in N_U$.
\item For $1\leq r\leq j$, $(u_{\leq r},u'_{\leq r})$ is a frame of rank $r$.
  Every nonzero $E$-subspace $U'\subset E^d$ with $h_W(U',U')=0$ is the span of
  a frame of rank $\dim U'$.
\end{enumerate}
\end{lemma}
\leanref{Development/WeilBound/ThetaLifts/EigenvalueEndgame/CuspidalAverages/AdelicInclusionThirdIdempotent.lean}{PositivityTheory.ThKEnd.Cusp.lev}
\leanref{Development/WeilBound/ThetaLifts/EigenvalueEndgame/CuspidalAverages/AdelicInclusionThirdIdempotent.lean}{PositivityTheory.ThKEnd.Cusp.isLev}
\leanref{Development/WeilBound/ThetaLifts/EigenvalueEndgame/CuspidalAverages/AdelicInclusionThirdIdempotent.lean}{PositivityTheory.ThKEnd.Cusp.lev_rat}
\leanref{Development/WeilBound/ThetaLifts/EigenvalueEndgame/CuspidalAverages/AdelicInclusionThirdIdempotent.lean}{PositivityTheory.ThKEnd.Cusp.lev_norm}
\leanref{Development/WeilBound/ThetaLifts/EigenvalueEndgame/CuspidalAverages/AdelicInclusionThirdIdempotent.lean}{PositivityTheory.ThKEnd.Cusp.mem_N}
\leanref{Development/WeilBound/ThetaLifts/EigenvalueEndgame/CuspidalAverages/AdelicInclusionThirdIdempotent.lean}{PositivityTheory.ThKEnd.Cusp.mid_of_move}
\leanref{Development/WeilBound/ThetaLifts/EigenvalueEndgame/CuspidalAverages/AdelicInclusionThirdIdempotent.lean}{PositivityTheory.ThKEnd.Cusp.par_mem}
\leanref{Development/WeilBound/ThetaLifts/EigenvalueEndgame/CuspidalAverages/IsAvLeft.lean}{PositivityTheory.ThKEnd.Cusp.fr_fst}
\leanref{Development/WeilBound/ThetaLifts/EigenvalueEndgame/CuspidalAverages/AvLine.lean}{PositivityTheory.ThKEnd.Cusp.Ctx.frame}
\leanref{Development/WeilBound/ThetaLifts/EigenvalueEndgame/CuspidalAverages/FirstIdempotentAdjointThird.lean}{PositivityTheory.ThKEnd.Cusp.exists_dual}
\begin{proof}
Write $J=J_W$, and put $e=(J\,c(u'))^{\mathsf T}$ and $e'=(J\,c(u))^{\mathsf T}$,
matrices with $j$ rows. As $c(J)^{\mathsf T}=-J$, the relations of a frame give
$eu=1$, $eu'=0$, $e'u=0$ and $e'u'=-1$. For $g,g'\in M_j(\mathbf A_E)$ put
$L(g,g')=1+u(g-1)e-u'(g'-1)e'$, so that
$\lambda_u(g)=L(g,c(g^{-1})^{\mathsf T})$. Expanding with these relations gives
$L(g,g')u=ug$, $L(g,g')u'=u'g'$, $L(g,g')L(k,k')=L(gk,g'k')$ and
$L(1,1)=1$, and $L(g,g')^{\mathsf T}J=J\,c\bigl(L(c(g')^{\mathsf T},c(g)^{\mathsf T})\bigr)$.

(1) So $\lambda_u(g)\lambda_u(k)=\lambda_u(gk)$, $\lambda_u(1)=1$,
$\lambda_u(g)u=ug$, and
$\lambda_u(g)^{\mathsf T}J=J\,c(\lambda_u(g^{-1}))$, hence
$\lambda_u(g)^{\mathsf T}J\,c(\lambda_u(g))=J$ and $\lambda_u(g)\in U_W$. The
entries of $\lambda_u(g)$ are polynomials in those of $g$, $g^{-1}$ and their
conjugates, so $\lambda_u$ is continuous, and they lie in $E$ when $g\in GL_j(E)$.

\emph{The unipotent group.} Put $e_1=ue$, $e_3=-u'e'$ and $q=1-e_1-e_3$. The
relations give $qu=qu'=0$, $eq=e'q=0$ and $q^2=q$, and
$qL(g,g')=L(g,g')q=q$. We show that an element $n\in U_W$ lies in $N_U$ if and
only if $nu=u$ and $q(n-1)q=0$. By
Lemma~\ref{asm-theta-agree-end-kill-unipotent}(1), $n\in N_U$ if and only if
$nu=u$ and $nz-z\in\mathbf A_EU$ for every $z\in E^d$ with $h_W(z,U)=0$, that
is, with $e'z=0$. If this holds, apply it to the columns of $q$, which satisfy
$e'q=0$: the columns of $(n-1)q$ lie in $\mathbf A_EU=u\mathbf A_E^j$, which $q$
kills, so $q(n-1)q=0$. Conversely, let $nu=u$ and $q(n-1)q=0$. Then
$(n-1)e_1=0$, and $e_3(n-1)=0$ as $n^{\mathsf T}J\,c(n)=J$ and $n^{-1}u=u$; with
$q(n-1)q=0$ this gives $(n-1)(1-e_3)=u\,e(n-1)(1-e_3)$. If $e'z=0$, then
$(1-e_3)z=z$, so $nz-z=u\,e(n-1)z\in\mathbf A_EU$.

(2) Let $n\in N_U$ and $m=\lambda_u(g)n\lambda_u(g)^{-1}$. Then
$mu=\lambda_u(g)nug^{-1}=\lambda_u(g)ug^{-1}=u$, and
$q(m-1)q=q\lambda_u(g)(n-1)\lambda_u(g)^{-1}q=q(n-1)q=0$. So $m\in N_U$.

(3) Put $m=h\lambda_u(g)^{-1}$. Then $mu=hug^{-1}=u$. The argument of the first
direction above, which uses only the condition on the $z$ with $e'z=0$, gives
$q(h-1)q=0$, and $q(m-1)q=qh\lambda_u(g)^{-1}q-q=q(h-1)q=0$. So $m\in N_U$.

(4) The relations of a frame for the first $r$ columns are submatrices of those
for $u$ and $u'$. Let $U'$ be as in the statement, of dimension $r$, and let
$v\in M_{d\times r}(E)$ have as columns a basis of $U'$, so that
$v^{\mathsf T}J\,c(v)=0$. Let $L$ be a left inverse of $v$, put
$w=c(J^{-1}L^{\mathsf T})$, so that $v^{\mathsf T}J\,c(w)=1$, and
$S=w^{\mathsf T}J\,c(w)$. Then $v'=w-\tfrac12vS^{\mathsf T}$ satisfies
$v'^{\mathsf T}J\,c(v')=0$ and $v^{\mathsf T}J\,c(v')=1$, as a computation with
$c(J)^{\mathsf T}=-J$ shows; so $(v,v')$ is a frame of rank $r$ with span $U'$.
\end{proof}
% Lean: lev (AdelicInclusionThirdIdempotent.lean) is mkU of lv (FirstIdempotentAdjointThird.lean)
% at (g, gs g), with Fr.lv_u, lv_v, lv_mul, lv_adj and lv_one for L; isLev and lev_rat for
% (1). The characterization of N_U is mem_N, through unipU_eq, mid_of_move and Fr.move, with
% mid s = e_1 + q s q + e_3 and mid_iff; (2) is lev_norm and (3) is par_mem with
% mid_of_move, both through mid_lv. (4) is fr_fst and Ctx.frame with exists_dual. Size:
% about 490 lines (FirstIdempotentAdjointThird.lean, AdelicInclusionThirdIdempotent.lean).

\begin{lemma}\label{asm-theta-agree-end-cusp-constant}
Let $E$ be a CM field with a theta base $B$
(Definition~\ref{asm-theta-agree-models}(1)), let $W$ be a skew-Hermitian
lattice space of rank $d$ and some signature, with the notation of
Definition~\ref{asm-theta-agree-end-cusp-notions}, let $(u,u')$ be a frame of
rank $j$ in $W$ with span $U$, and let $f$ be an automorphic function on $U_W$.
\begin{enumerate}
\item The constant term $c_Uf$ is the average of $f$ over $N_U\cap U_W(F)$
  (Definition~\ref{asm-theta-agree-end-kill-notions}(2)), and it is continuous.
  For $h\in U_W$, $n\in N_U$ and $\gamma\in GL_j(E)$,
  $c_Uf(nh)=c_Uf(h)$ and $c_Uf(\lambda_u(\gamma)h)=c_Uf(h)$. For every
  function $f'$ on $U_W$, $k\in U_W$ and $z\in\mathbf C$,
  $c_U(f'(\cdot\,k))=(c_Uf')(\cdot\,k)$ and $c_U(zf')=z\,c_Uf'$; and
  $c_U(f+f'')=c_Uf+c_Uf''$ for every automorphic $f''$.
\item Let $1\leq r\leq j-1$, let $U_r$ be the span of $u_{\leq r}$, and suppose
  that $0$ is the average of $f$ over $N_{U_r}\cap U_W(F)$. Then
  $\bigl\langle y\mapsto c_Uf(\lambda_u(n_{r,j-1}(y))h)\bigr\rangle_{I_{r,j-1}}=0$
  for every $h\in U_W$, where $n_{r,j-1}$ is taken in $GL_j(\mathbf A_E)$.
\end{enumerate}
\end{lemma}
\leanref{Development/WeilBound/ThetaLifts/EigenvalueEndgame/CuspidalAverages/IsAvLeft.lean}{PositivityTheory.ThKEnd.Cusp.isAv_left}
\leanref{Development/WeilBound/ThetaLifts/EigenvalueEndgame/CuspidalAverages/IsAvLeft.lean}{PositivityTheory.ThKEnd.Cusp.ct_cont}
\leanref{Development/WeilBound/ThetaLifts/EigenvalueEndgame/CuspidalAverages/IsAvLeft.lean}{PositivityTheory.ThKEnd.Cusp.ct_left}
\leanref{Development/WeilBound/ThetaLifts/EigenvalueEndgame/CuspidalAverages/IsAvLeft.lean}{PositivityTheory.ThKEnd.Cusp.ct_rat}
\leanref{Development/WeilBound/ThetaLifts/EigenvalueEndgame/CuspidalAverages/IsAvLeft.lean}{PositivityTheory.ThKEnd.Cusp.ct_van}
\leanref{Development/WeilBound/ThetaLifts/EigenvalueEndgame/CuspidalAverages/StageB.lean}{PositivityTheory.ThKEnd.Cusp.prod_stage}
\leanref{Development/WeilBound/ThetaLifts/EigenvalueEndgame/CuspidalAverages/StageB.lean}{PositivityTheory.ThKEnd.Cusp.zero_of_retract}
\begin{proof}
(1) By Lemma~\ref{asm-theta-agree-end-kill-unipotent}(2), $N_U$ is closed and
$\Gamma_U=N_U\cap U_W(F)$ is a lattice in it, so by
Lemma~\ref{asm-theta-agree-end-kill-averages}(2) and~(3) $f$ has an average over
$\Gamma_U$, which is $h\mapsto\int_{D_U}f(nh)\,d\nu_U(n)=c_Uf(h)$, and it is
continuous. Let $n\in N_U$. As $\nu_U$ is invariant under right translations,
$D_Un$ is a fundamental domain of measure $1$, so
$c_Uf(nh)=\int_{D_Un}f(mh)\,d\nu_U(m)=c_Uf(h)$, the average not depending on
the fundamental domain. Let $\gamma\in GL_j(E)$. By
Lemma~\ref{asm-theta-agree-end-cusp-frames}(1) and~(2), applied to $\gamma$ and
$\gamma^{-1}$, $\lambda_u(\gamma)$ lies in $U_W(F)$ and conjugation by it maps
$N_U$ onto $N_U$; so Lemma~\ref{asm-theta-agree-end-kill-averages}(4) gives
$c_Uf(\lambda_u(\gamma)h)=c_Uf(h)$. The last claims follow from the definition
of $c_U$; for the additivity, $f$ and $f''$ are bounded on $D_Uh$ by
Lemma~\ref{asm-theta-agree-end-kill-averages}(3), so both are integrable there.

(2) Write $\Gamma=U_W(F)$, $I=I_{r,j-1}$, $n(y)=n_{r,j-1}(y)$ and
$\iota(y)=\lambda_u(n(y))$ for $y\in\mathbf A_E^I$, and let $e$ be as in the
proof of Lemma~\ref{asm-theta-agree-end-cusp-frames}. By that lemma, $\iota$ is a
continuous homomorphism, $\iota(E^I)\subset\Gamma$, and $\iota(y)$ normalizes
$N_U$.

\emph{The group $S$.} For $s\in U_W$ let $\pi(s)\in\mathbf A_E^I$ consist of the
entries of $e\,s\,u$ at the places in $I$. If $n\in N_U$, then $nu=u$, so
$e\,n\,\iota(y)u=e\,u\,n(y)=n(y)$ and $\pi(n\iota(y))=y$. So
$S=N_U\,\iota(\mathbf A_E^I)$ is a subgroup of $U_W$, closed as it is the set of
the $s$ with $s\iota(\pi(s))^{-1}\in N_U$, and $\pi$ restricts to a continuous
homomorphism from $S$ onto $\mathbf A_E^I$ with kernel $N_U$, with
$\pi(\iota(y))=y$ and $\pi(S\cap\Gamma)\subset E^I$. The group $E^I$ is a
lattice in $\mathbf A_E^I$, by Lemma~\ref{asm-theta-agree-end-kill-siegel}(1). By
Lemma~\ref{asm-theta-agree-end-kill-stages}, $S\cap\Gamma$ is a lattice in $S$,
and the average $c$ of $c_Uf$ over $E^I$ along $\iota$ is the average of $f$ over
$S\cap\Gamma$.

\emph{A retraction.} As $n(y)$ fixes the first $r$ standard basis vectors, every
$s\in S$ fixes the columns of $u_{\leq r}$. Let $e^{(r)}$, $e_1^{(r)}$,
$e_3^{(r)}$ and $q^{(r)}$ be the matrices $e$, $e_1$, $e_3$ and $q$ of the proof
of Lemma~\ref{asm-theta-agree-end-cusp-frames} for the frame
$(u_{\leq r},u'_{\leq r})$, and for $s\in S$ put
$\rho(s)=e^{(r)}_1+q^{(r)}sq^{(r)}+e^{(r)}_3$. A computation with the relations
of that proof, and with $su_{\leq r}=u_{\leq r}$, shows that $\rho(s)\in U_W$,
$\rho(ss')=\rho(s)\rho(s')$, $\rho(\rho(s))=\rho(s)$, $\rho(s)u_{\leq r}=u_{\leq r}$
and $q^{(r)}(s\rho(s)^{-1}-1)q^{(r)}=0$; and $\rho(s)\in\Gamma$ when
$s\in\Gamma$. By the characterization of $N_{U_r}$ in that proof,
$s\in N_{U_r}$ if and only if $\rho(s)=1$, and $s\rho(s)^{-1}\in N_{U_r}$.
Moreover $N_{U_r}\subset S$: let $n\in N_{U_r}$. By the characterization of
$N_{U_r}$, $nu_{\leq r}=u_{\leq r}$ and $q^{(r)}(n-1)q^{(r)}=0$, and the
argument of its converse direction, for the frame $(u_{\leq r},u'_{\leq r})$,
gives
$(n-1)(1-e^{(r)}_3)=u_{\leq r}\,e^{(r)}(n-1)(1-e^{(r)}_3)$. As
$u^{\mathsf T}J_W\,c(u)=0$, $e^{(r)}_3u=0$; so multiplying on the right by $u$
gives $(n-1)u=u_{\leq r}Y$ with $Y=e^{(r)}(n-1)u$, whose first $r$ columns
vanish as $n$ fixes $u_{\leq r}$. So $nu=u\,n(y)$ for the entries $y$ of $Y$ at
the places in $I$. For $z\in E^d$ with $h_W(z,U)=0$, also $h_W(z,U_r)=0$, so
$nz-z\in\mathbf A_EU_r\subset\mathbf A_EU$ by
Lemma~\ref{asm-theta-agree-end-kill-unipotent}(1); and
Lemma~\ref{asm-theta-agree-end-cusp-frames}(3) gives $n\iota(y)^{-1}\in N_U$,
so $n\in S$. Hence $\rho(s)\in S$ for $s\in S$, and
$\rho$ is a continuous homomorphism from $S$ onto the closed subgroup
$R=\rho(S)$ of its fixed points, with kernel $N_{U_r}$ and $\rho(t)=t$ for
$t\in R$. The group $R\cap\Gamma$ is countable and discrete, and if $C\subset S$
is compact with $(S\cap\Gamma)C=S$, then $(R\cap\Gamma)\rho(C)=R$, as
$\rho(S\cap\Gamma)\subset R\cap\Gamma$. So $R\cap\Gamma$ is a lattice in $R$.

\emph{Conclusion.} By Lemma~\ref{asm-theta-agree-end-kill-stages} for $S$, the
normal subgroup $N_{U_r}$, $R$ with the lattice $R\cap\Gamma$, the inclusion of
$R$ and $\rho$, the average of $f$ over $S\cap\Gamma$ is the average over
$R\cap\Gamma$ of the average of $f$ over $N_{U_r}\cap\Gamma$, which is $0$. So
$c=0$ by Lemma~\ref{asm-theta-agree-end-kill-averages}(2). The function
$y\mapsto c_Uf(\iota(y)h)$ is invariant under $E^I$, by (1), so its integral
over $D_I$ is $c(h)=0$.
\end{proof}
% Lean: (1) is Ctx.ct (with Lat.isAv_prob and Lat.bound), IsCt.av, ct_cont (IsAv.cont),
% ct_left (isAv_left) and ct_rat (IsAv.conj with conjE and lev_rat, lev_norm). (2) is ct_van
% (IsAvLeft.lean, 100 lines): beta, iota and pi, prod_stage (StageB.lean) for S = prodG and
% the first average (s1, s2, s3), fix, le (with Fr.move for the small frame, par_mem and
% mid_of_move), rho_0 = mkU of mid for the small frame (mid_mat, mid_mid, mid_adj, mem_N for
% the small frame, incM_mid), zero_of_retract (StageB.lean, through Semi.stage), uniqueness
% isAv_uniq, and avg_eq_of_isAv. In le, Y comes from Fr.move for the small frame (mv),
% e_3 u = 0 is e3u (with (frA h).uu'), the vanishing of the first r columns is sup, and
% nz - z in A_E U_r is unipU_eq for the small frame.

\begin{lemma}\label{asm-theta-agree-end-cusp-expansion}
Let $E$ be a number field, let $\psi$ be a standard character of $\mathbf A_E$,
let $a\geq1$, and use the notation of
Definition~\ref{asm-theta-agree-end-cusp-notions}(4) to~(6). For $0\leq m\leq a-1$,
call a function $\varphi$ on $GL_a(\mathbf A_E)$ \emph{adapted at stage $m$} if
it is continuous, $\varphi(\gamma g)=\varphi(g)$ for every $g$ and every
$\gamma\in GL_a(E)$ whose rows $m+1,\ldots,a$ and columns $m+2,\ldots,a$ are
those of the identity matrix, and
$\varphi(n_{k,k}(x)g)=\psi(x_{k,k+1})\varphi(g)$ for $m<k\leq a-1$,
$x\in\mathbf A_E^{I_{k,k}}$ and every $g$. Call it \emph{cuspidal at stage
$m$} if moreover
$\langle y\mapsto\varphi(n_{r,m}(y)g)\rangle_{I_{r,m}}=0$ for $1\leq r\leq m$
and every $g$.
\begin{enumerate}
\item If $\varphi$ is adapted at stage $m$, then $\mathcal W^m_\psi\varphi$ is
  adapted at stage $0$.
\item If $\varphi$ and $\varphi'$ are adapted at stage $m$, then
  $\mathcal W^m_\psi(\varphi+\varphi')=\mathcal W^m_\psi\varphi+\mathcal W^m_\psi\varphi'$.
  For every function $\varphi$, $z\in\mathbf C$ and $h,g\in GL_a(\mathbf A_E)$,
  $\mathcal W^m_\psi(z\varphi)=z\,\mathcal W^m_\psi\varphi$ and
  $\mathcal W^m_\psi(\varphi(\cdot\,h))(g)=\mathcal W^m_\psi\varphi(gh)$.
\item If $\varphi$ is adapted at stage $0$, then
  $\varphi(ng)=\psi\bigl(\sum_{i=1}^{a-1}n_{i,i+1}\bigr)\varphi(g)$ for every
  $g$ and every upper triangular $n\in GL_a(\mathbf A_E)$ with diagonal entries
  $1$.
\item If $\varphi$ is cuspidal at stage $m$, then $\mathcal W^m_\psi\varphi$ is
  cuspidal at stage $0$, and $\mathcal W^m_\psi\varphi\neq0$ if
  $\varphi\neq0$.
\end{enumerate}
\end{lemma}
\leanref{Development/WeilBound/ThetaLifts/EigenvalueEndgame/CuspidalAverages/AvgEqOfIsAv.lean}{PositivityTheory.ThKEnd.Cusp.Pre}
\leanref{Development/AbelianVarietyTheory/Proofs/InertiaEigenvalues/UnipotentBlockColumnIndex.lean}{GlobalTheory.CuspGeneric.Good}
\leanref{Development/WeilBound/ThetaLifts/EigenvalueEndgame/CuspidalAverages/AvgEqOfIsAv.lean}{PositivityTheory.ThKEnd.Cusp.whit_pre}
\leanref{Development/WeilBound/ThetaLifts/EigenvalueEndgame/CuspidalAverages/AvgEqOfIsAv.lean}{PositivityTheory.ThKEnd.Cusp.whit_add'}
\leanref{Development/AbelianVarietyTheory/Proofs/InertiaEigenvalues/IsGenericOfIsLocComp.lean}{GlobalTheory.CuspRep.Generic.whit_smul}
\leanref{Development/AbelianVarietyTheory/Proofs/InertiaEigenvalues/IsGenericOfIsLocComp.lean}{GlobalTheory.CuspRep.Generic.whit_rt}
\leanref{Development/AbelianVarietyTheory/Proofs/InertiaEigenvalues/WhittakerLawUpperTriangularUnipotentAdelicStatement.lean}{GlobalTheory.T.whittakerLaw_upperTriangularUnipotent_adelic}
\leanref{Development/AbelianVarietyTheory/Proofs/InertiaEigenvalues/IsGenericOfIsLocComp.lean}{GlobalTheory.CuspRep.Generic.whit_good}
\leanref{Development/WeilBound/ThetaLifts/EigenvalueEndgame/CuspidalAverages/AvgEqOfIsAv.lean}{PositivityTheory.ThKEnd.Cusp.pre_colT}
\leanref{Development/AbelianVarietyTheory/Proofs/InertiaEigenvalues/IsGenericOfIsLocComp.lean}{GlobalTheory.CuspRep.Generic.good_colT}
\leanref{Development/AbelianVarietyTheory/Proofs/InertiaEigenvalues/UnipotentColumnWhittakerTransformNeZero.lean}{GlobalTheory.CuspGeneric.colT_ne_zero_of_good}
\leanref{Development/AbelianVarietyTheory/Proofs/InertiaEigenvalues/WhittakerCoefficientRationalMirabolicMulStatement.lean}{GlobalTheory.T.whittakerCoefficient_rationalMirabolic_mul}
\leanref{Development/AbelianVarietyTheory/Proofs/InertiaEigenvalues/PeriodicEqZeroOfFourierCoefficientsZeroStatement.lean}{GlobalTheory.T.periodic_eq_zero_of_fourier_coefficients_zero}
\leanref{Development/AbelianVarietyTheory/Proofs/InertiaEigenvalues/ExistsRationalMirabolicMovingRowToLastStatement.lean}{GlobalTheory.T.exists_rationalMirabolic_moving_row_to_last}
\leanref{Development/AbelianVarietyTheory/Proofs/InertiaEigenvalues/QuotientAverageTranslateInvariantStatement.lean}{GlobalTheory.T.quotientAverage_translate_invariant}
\leanref{Development/AbelianVarietyTheory/Proofs/InertiaEigenvalues/QuotientAverageAddOfPeriodicStatement.lean}{GlobalTheory.T.quotientAverage_add_of_periodic}
\leanref{Development/AbelianVarietyTheory/Proofs/InertiaEigenvalues/ContinuousQuotientAveragePeriodicFamilyStatement.lean}{GlobalTheory.T.continuous_quotientAverage_periodic_family}
\leanref{Development/AbelianVarietyTheory/Proofs/InertiaEigenvalues/QuotientAverageSumIndexFubiniStatement.lean}{GlobalTheory.T.quotientAverage_sum_index_fubini}
\leanref{Development/AbelianVarietyTheory/Proofs/InertiaEigenvalues/QuotientAverageAdelicRationalPointsCommStatement.lean}{GlobalTheory.T.quotientAverage_adelicRationalPoints_comm}
\leanref{Development/AbelianVarietyTheory/Proofs/InertiaEigenvalues/QuotientAverageUnipotentBlockColumnTransformEqZeStatement.lean}{GlobalTheory.T.quotientAverage_unipotentBlock_columnTransform_eq_zero}
\begin{proof}
Write $P_m$ for the set of the $\gamma\in GL_a(E)$ whose rows $m+1,\ldots,a$
and columns $m+2,\ldots,a$ are those of the identity matrix. For
$1\leq k\leq a-1$, $x\in\mathbf A_E^{I_{k,k}}$ and $\xi\in E^k$ put
$\xi\cdot x=\sum_{i\leq k}\xi_ix_{i,k+1}$ and
\[
C^k_\xi\varphi(g)=\bigl\langle x\mapsto\psi(-\xi\cdot x)\varphi(n_{k,k}(x)g)\bigr\rangle_{I_{k,k}},
\]
so that $T_k=C^k_{e_k}$ for $e_k=(0,\ldots,0,1)\in E^k$; and for
$\gamma\in GL_a(E)$ let $\xi\gamma\in E^k$ have the entries
$\sum_{s\leq k}\xi_s\gamma_{st}$ for $t\leq k$. Five facts are used.
\begin{enumerate}
\item[(A)] Let $I$ be a finite set, and let $\Phi$ and $\Phi'$ be continuous
  functions on $\mathbf A_E^I$ invariant under translation by $E^I$. Then
  $\langle\Phi(\cdot+y)\rangle_I=\langle\Phi\rangle_I$ for
  $y\in\mathbf A_E^I$, and
  $\langle\Phi+\Phi'\rangle_I=\langle\Phi\rangle_I+\langle\Phi'\rangle_I$. If
  $F$ is a continuous function on $X\times\mathbf A_E^I$, for a topological
  space $X$, invariant under $E^I$ in the second variable, then
  $p\mapsto\langle F(p,\cdot)\rangle_I$ is continuous.
\item[(B)] Let $\varphi$ be continuous, with $\varphi(\gamma g)=\varphi(g)$ for
  $\gamma\in P_k$ and every $g$. Then
  $C^k_\xi\varphi(\gamma g)=C^k_{\xi\gamma}\varphi(g)$ for $\gamma\in P_k$,
  $\xi\in E^k$ and every $g$.
\item[(C)] If $\Phi$ is as in (A) and
  $\langle x\mapsto\psi(-\sum_{t\in I}\xi_tx_t)\Phi(x)\rangle_I=0$ for every
  $\xi\in E^I$, then $\Phi=0$.
\item[(D)] For every nonzero $\xi\in E^k$ there is $\gamma\in P_k$ with
  $e_k\gamma=\xi$.
\item[(E)] Let $I$ and $I'$ be disjoint finite sets, and let $\Phi$ be a
  continuous function on
  $\mathbf A_E^{I\cup I'}=\mathbf A_E^I\times\mathbf A_E^{I'}$ invariant under
  translation by $E^{I\cup I'}$. Then
  $\langle\Phi\rangle_{I\cup I'}=\bigl\langle x\mapsto\langle y\mapsto\Phi(x,y)\rangle_{I'}\bigr\rangle_I
  =\bigl\langle y\mapsto\langle x\mapsto\Phi(x,y)\rangle_I\bigr\rangle_{I'}$.
\end{enumerate}
In (A), $\langle\Phi\rangle_I$ is the integral of the function that $\Phi$
induces on the compact group $\mathbf A_E^I/E^I$ against its invariant
probability measure. (B) is a change of variables in the integral. (C) is
Fourier analysis on this compact group, whose characters are the
$x\mapsto\psi(\sum_t\xi_tx_t)$ with $\xi\in E^I$ \cite[\S4.1]{Tat67}. For (D),
let $i_0\leq k$ have $\xi_{i_0}\neq0$, and let $\gamma$ have row $k$ equal to
$(\xi_1,\ldots,\xi_k,0,\ldots,0)$, row $i_0$ equal to the $k$-th standard basis
vector if $i_0\neq k$, and its other rows those of the identity; its
determinant is $\pm\xi_{i_0}$, and $\gamma\in P_k$. (E) is Fubini's theorem on
$\mathbf A_E^{I\cup I'}/E^{I\cup I'}$, the product of the compact groups
$\mathbf A_E^I/E^I$ and $\mathbf A_E^{I'}/E^{I'}$, whose invariant probability
measure is the product measure.

\emph{One column.} Let $1\leq k\leq a-1$ and let $\varphi$ be adapted at stage
$k$. For every $g$ and $\xi$ the function
$x\mapsto\psi(-\xi\cdot x)\varphi(n_{k,k}(x)g)$ is continuous and invariant
under $E^{I_{k,k}}$, as $n_{k,k}(x+x')=n_{k,k}(x')n_{k,k}(x)$, $n_{k,k}(x')\in P_k$
and $\psi(\xi\cdot x')=1$ for $x'\in E^{I_{k,k}}$.
\begin{enumerate}
\item[(S1)] \emph{$T_k\varphi$ is adapted at stage $k-1$.} It is continuous, by
  (A). Let $\gamma\in P_{k-1}$. Then $\gamma\in P_k$ and $e_k\gamma=e_k$, as row
  $k$ of $\gamma$ is that of the identity; so (B) gives
  $T_k\varphi(\gamma g)=T_k\varphi(g)$. Let $k-1<k'\leq a-1$ and
  $y\in\mathbf A_E^{I_{k',k'}}$. If $k'=k$, then
  $n_{k,k}(x)n_{k,k}(y)=n_{k,k}(x+y)$, and (A), for the translation by $y$,
  gives $T_k\varphi(n_{k,k}(y)g)=\psi(y_{k,k+1})T_k\varphi(g)$. If $k'>k$, then
  $n_{k,k}(x)n_{k',k'}(y)=n_{k',k'}(y')n_{k,k}(x)$, where $y'$ differs from $y$
  only in the rows $i\leq k$; so $y'_{k',k'+1}=y_{k',k'+1}$, and the law of
  $\varphi$ at $k'$ gives
  $T_k\varphi(n_{k',k'}(y)g)=\psi(y_{k',k'+1})T_k\varphi(g)$.
\item[(S2)] \emph{If $\varphi$ is cuspidal at stage $k$ and $\varphi\neq0$, then
  $T_k\varphi\neq0$.} Let $\varphi(g_0)\neq0$. The function
  $\Phi(x)=\varphi(n_{k,k}(x)g_0)$ is continuous, invariant under
  $E^{I_{k,k}}$, and $\Phi(0)\neq0$; so by (C) there is $\xi$ with
  $C^k_\xi\varphi(g_0)\neq0$. As $\varphi$ is cuspidal at stage $k$, for $r=k$,
  $C^k_0\varphi(g_0)=0$, so $\xi\neq0$. Let $\gamma$ be as in (D). By (B),
  $T_k\varphi(\gamma g_0)=C^k_{e_k\gamma}\varphi(g_0)=C^k_\xi\varphi(g_0)\neq0$.
\item[(S3)] \emph{If $\varphi$ is cuspidal at stage $k$, then $T_k\varphi$ is
  cuspidal at stage $k-1$.} By (S1) it remains to show that
  $\langle y\mapsto T_k\varphi(n_{r,k-1}(y)g)\rangle_{I_{r,k-1}}=0$ for
  $1\leq r\leq k-1$ and every $g$. Fix $r$ and $g$, and split $I_{k,k}$ into
  $H=\{(i,k+1):r<i\leq k\}$ and $L=\{(i,k+1):i\leq r\}$; then $I_{r,k}$ is the
  disjoint union of $I_{r,k-1}$ and $L$. For $z\in\mathbf A_E^H$ and
  $w\in\mathbf A_E^L$ write $(z,w)$ for the point of $\mathbf A_E^{I_{k,k}}$
  with these entries, and likewise $(y,w)\in\mathbf A_E^{I_{r,k}}$ for
  $y\in\mathbf A_E^{I_{r,k-1}}$. As $n_{r,k}(x+x')=n_{r,k}(x)n_{r,k}(x')$ and
  $n_{r,k}(\gamma)\in P_k$ for $\gamma\in E^{I_{r,k}}$, the function
  $x\mapsto\varphi(n_{r,k}(x)h)$ is continuous and invariant under
  $E^{I_{r,k}}$ for every $h$.

  \emph{A product formula.} For $y\in\mathbf A_E^{I_{r,k-1}}$ and
  $z\in\mathbf A_E^H$ let $\kappa(y,z)\in\mathbf A_E^L$ have the entries
  $\kappa(y,z)_{i,k+1}=\sum_{r<l\leq k}y_{il}z_{l,k+1}$. Then
  \[
  n_{k,k}(z,w)\,n_{r,k-1}(y)=n_{r,k}\bigl(y,w-\kappa(y,z)\bigr)\,n_{k,k}(z,0).
  \]
  Indeed, let $P$, $Q$ and $N$ be the matrices with the entries of $z$, $w$
  and $y$ at their places and $0$ elsewhere. Then $PN=QN=QP=0$ and $P^2=0$,
  and $NP$ has the entries of $\kappa(y,z)$ at the places in $L$ and $0$
  elsewhere; so both sides are $1+P+Q+N$.

  \emph{The inner average.} As $r<k$, the place $(k,k+1)$ lies in $H$, so
  $e_k\cdot(z,w)=z_{k,k+1}$. By the product formula and (A), for the
  translation by $-\kappa(y,z)$,
  \[
  \bigl\langle w\mapsto\psi(-e_k\cdot(z,w))\,\varphi\bigl(n_{k,k}(z,w)\,n_{r,k-1}(y)g\bigr)\bigr\rangle_L
  =\Phi(y,z),
  \]
  where $\Phi(y,z)=\psi(-z_{k,k+1})\Psi(y,z)$ and
  $\Psi(y,z)=\langle w\mapsto\varphi(n_{r,k}(y,w)\,n_{k,k}(z,0)g)\rangle_L$.
  By (E) for $I_{k,k}=H\cup L$,
  $T_k\varphi(n_{r,k-1}(y)g)=\langle z\mapsto\Phi(y,z)\rangle_H$.

  \emph{Invariance.} The function $\Phi$ on
  $\mathbf A_E^{I_{r,k-1}}\times\mathbf A_E^H$ is continuous, by (A), and it
  is invariant under $E^{I_{r,k-1}}\times E^H$. For the translations of $y$ by
  $\gamma\in E^{I_{r,k-1}}$ this holds as
  $n_{r,k}(y+\gamma,w)=n_{r,k}(\gamma,0)\,n_{r,k}(y,w)$ with
  $n_{r,k}(\gamma,0)\in P_k$. For the translations of $z$ by $\gamma'\in E^H$
  it holds for the left side of the last display, as
  $n_{k,k}(z+\gamma',w)=n_{k,k}(\gamma',0)\,n_{k,k}(z,w)$ with
  $n_{k,k}(\gamma',0)\in P_k$ and $\psi(e_k\cdot(\gamma',0))=1$.

  \emph{Conclusion.} By (E), the average to be computed is
  \[
  \bigl\langle y\mapsto\langle z\mapsto\Phi(y,z)\rangle_H\bigr\rangle_{I_{r,k-1}}
  =\bigl\langle z\mapsto\psi(-z_{k,k+1})\,\langle y\mapsto\Psi(y,z)\rangle_{I_{r,k-1}}\bigr\rangle_H.
  \]
  By (E) for $I_{r,k}=I_{r,k-1}\cup L$,
  $\langle y\mapsto\Psi(y,z)\rangle_{I_{r,k-1}}=\langle x\mapsto\varphi(n_{r,k}(x)\,n_{k,k}(z,0)g)\rangle_{I_{r,k}}$,
  which is $0$ as $\varphi$ is cuspidal at stage $k$ and $r\leq k$. So the
  average is $0$.
  % Lean: the van field of good_colT is
  % GlobalTheory.T.quotientAverage_unipotentBlock_columnTransform_eq_zero =
  % CuspGeneric.avg_colT_blk_eq_zero, proved by ColSplit.colT_avg_blk_eq_zero
  % (QuotientAverageUnipotentBlockColumnTransformEqZe.lean, 244 lines). H and L are HiIx and
  % LoIx (PosMApplyNeZero.lean), the product formula is blkU_col_mul_blkU with kappa(y,z) = corr;
  % hA is (E) for I_{k,k} (avg_sum_reix with colEquiv, quotientAverage_sum_index_fubini), hB
  % the inner average (dot_col, quotientAverage_translate_invariant, per_position); Phi and
  % Psi are Phi_2 and H, hHc and hPer the continuity and invariance (blk_rat, map_dot), hD
  % the vanishing (avg_sum_reix with blkEquiv and the van field at stage k), and the
  % exchange of averages is quotientAverage_adelicRationalPoints_comm.
\end{enumerate}

(1) For $m\geq1$, $\mathcal W^m_\psi\varphi=\mathcal W^{m-1}_\psi(T_m\varphi)$,
so (1) follows from (S1) by induction on $m$.

(2) The identities for $z\varphi$ and for $\varphi(\cdot\,h)$ hold for each
$T_k$, by the definition. The additivity follows by induction from (A) and
(S1).

(3) For $1\leq p\leq a$ let (3)$_p$ be the claim for the matrices whose
columns $p+1,\ldots,a$ are those of the identity; (3)$_a$ is (3). We show
(3)$_p$ by induction on $p$. In (3)$_1$ the matrix is $1$. Let $p\geq2$, let $u$
be upper triangular with diagonal entries $1$ and columns $p+1,\ldots,a$ those
of the identity, let $y\in\mathbf A_E^{I_{p-1,p-1}}$ consist of the entries of
column $p$ of $u$ above the diagonal, and put $n'=n_{p-1,p-1}(y)$. Row $p$ of
$u$ is that of the identity, so $n'^{-1}u=u-(n'-1)$: it has the entries of $u$,
except $0$ above the diagonal in column $p$, and its columns $p,\ldots,a$ are
those of the identity. By the law of $\varphi$ at $p-1$ and (3)$_{p-1}$,
$\varphi(ug)=\psi(u_{p-1,p})\varphi(n'^{-1}ug)=\psi\bigl(\sum_iu_{i,i+1}\bigr)\varphi(g)$.

(4) By (S2) and (S3), by induction on $m$, as in (1).
% Lean: P_m is IsMir m. (A) is GlobalTheory.T.quotientAverage_translate_invariant,
% quotientAverage_add_of_periodic and continuous_quotientAverage_periodic_family, (B) is
% whittakerCoefficient_rationalMirabolic_mul (= CuspGeneric.coefF_ratGL_mul, 226 lines), (C) is
% periodic_eq_zero_of_fourier_coefficients_zero (= eq_zero_of_avg_character_eq_zero, through
% SeparatingCharacters.eq_zero_of_coefficients_eq_zero and StdPsi.forall_dot_eq_one_iff) and
% (D) is exists_rationalMirabolic_moving_row_to_last (= mir_exists_row), and (E) is
% quotientAverage_sum_index_fubini and quotientAverage_adelicRationalPoints_comm (CuspGeneric.avg_sum_index
% and avg_comm). (S1) is pre_colT, (S2)
% is colT_ne_zero_of_good and (S3) is good_colT. C^k_xi is coefF psi k xi.
% Lean: (1) is whit_pre, by induction from pre_colT (AvgEqOfIsAv.lean), which uses
% GlobalTheory.T.continuous_quotientAverage_periodic_family,
% whittakerCoefficient_rationalMirabolic_mul and quotientAverage_translate_invariant. (2) is
% whit_add' (with quotientAverage_add_of_periodic), whit_smul and whit_rt. (3) is
% GlobalTheory.T.whittakerLaw_upperTriangularUnipotent_adelic = CuspGeneric.tri_law. (4) is
% CuspRep.Generic.whit_good, by induction from good_colT
% (continuous_unipotentColumnWhittakerTransform, unipotentColumnWhittakerTransform_mirabolic_invariant,
% columnWhittakerTransform_keeps_later_law, unipotentColumnWhittakerTransform_column_law,
% quotientAverage_unipotentBlock_columnTransform_eq_zero) and colT_ne_zero_of_good
% (UnipotentColumnWhittakerTransformNeZero.lean: Fourier inversion on (E\A_E)^(m+1) by
% periodic_eq_zero_of_fourier_coefficients_zero, the vanishing of the coefficient at 0 by the
% cuspidality at r = m+1, and exists_rationalMirabolic_moving_row_to_last). This is the
% inductive step of the Fourier expansion of cusp forms on GL_a of Shalika and
% Piatetski-Shapiro, under the weaker hypothesis of cuspidality at stage m. The lemmas are in
% Development/AbelianVarietyTheory/Proofs/InertiaEigenvalues and are shared with other
% chapters' genericity arguments (isGeneric_of_isLocComp).
\end{proof}

\begin{lemma}\label{asm-theta-agree-end-cusp-coefficient}
Let $E$ be a CM field with a theta base $B$
(Definition~\ref{asm-theta-agree-models}(1)), let $W$ be a skew-Hermitian
lattice space of rank $d$ and some signature, with the notation of
Definition~\ref{asm-theta-agree-end-cusp-notions}, let $(u,u')$ be a frame of
rank $j$ in $W$ with span $U$, and let $\psi$ be a standard character of
$\mathbf A_E$. For an automorphic function $f$ on $U_W$ and $x\in U_W$ put
$\varphi_{f,x}(g)=c_Uf(\lambda_u(g)x)$ for $g\in GL_j(\mathbf A_E)$.
\begin{enumerate}
\item The character $\psi$ is nondegenerate, and $\mathcal W_{u,\psi}$ is a
  Whittaker functional for $(u,\psi)$.
\item For automorphic $f$, $x\in U_W$ and $g\in GL_j(\mathbf A_E)$,
  $\mathcal W_{u,\psi}(f(\cdot\,\lambda_u(g)x))=\mathcal W^{j-1}_\psi\varphi_{f,x}(g)$.
\item Let $f$ be automorphic. Suppose that for $1\leq r<j$, $0$ is the average
  of $f$ over $N_{U_r}\cap U_W(F)$, where $U_r$ is the span of $u_{\leq r}$, and
  that $\mathcal W_{u,\psi}(f(\cdot\,x))=0$ for every $x\in U_W$. Then
  $c_Uf=0$, and $0$ is the average of $f$ over $N_U\cap U_W(F)$.
\end{enumerate}
\end{lemma}
\leanref{Development/WeilBound/ThetaLifts/EigenvalueEndgame/CuspidalAverages/CuspOfLineOfKeys.lean}{PositivityTheory.ThKEnd.Cusp.chi}
\leanref{Development/WeilBound/ThetaLifts/EigenvalueEndgame/CuspidalAverages/CuspOfLineOfKeys.lean}{PositivityTheory.ThKEnd.Cusp.isChar}
\leanref{Development/WeilBound/ThetaLifts/EigenvalueEndgame/CuspidalAverages/CuspOfLineOfKeys.lean}{PositivityTheory.ThKEnd.Cusp.isWh}
\leanref{Development/WeilBound/ThetaLifts/EigenvalueEndgame/CuspidalAverages/CuspOfLineOfKeys.lean}{PositivityTheory.ThKEnd.Cusp.wh_right}
\leanref{Development/WeilBound/ThetaLifts/EigenvalueEndgame/CuspidalAverages/CuspOfLineOfKeys.lean}{PositivityTheory.ThKEnd.Cusp.wh_ne}
\leanref{Development/WeilBound/ThetaLifts/EigenvalueEndgame/CuspidalAverages/IsAvLeft.lean}{PositivityTheory.ThKEnd.Cusp.ct_pre}
\leanref{Development/WeilBound/ThetaLifts/EigenvalueEndgame/CuspidalAverages/AvgEqOfIsAv.lean}{PositivityTheory.ThKEnd.Cusp.pre_top}
\begin{proof}
Let $f$ be automorphic and $x\in U_W$. By
Lemma~\ref{asm-theta-agree-end-cusp-constant}(1) and
Lemma~\ref{asm-theta-agree-end-cusp-frames}(1), $\varphi_{f,x}$ is continuous
and invariant under left translation by $GL_j(E)$. So it is adapted at stage
$j-1$ (Lemma~\ref{asm-theta-agree-end-cusp-expansion}), the condition on the
$n_{k,k}$ with $k>j-1$ being empty. By
Lemma~\ref{asm-theta-agree-end-cusp-constant}(1),
$c_U(f(\cdot\,k))=(c_Uf)(\cdot\,k)$ for $k\in U_W$.

(2) By the definition and as $\lambda_u$ is a homomorphism,
$\mathcal W_{u,\psi}(f(\cdot\,\lambda_u(g)x))$ is
$\mathcal W^{j-1}_\psi$ of the function
$g'\mapsto c_Uf(\lambda_u(g'g)x)=\varphi_{f,x}(g'g)$, evaluated at $1$. By
Lemma~\ref{asm-theta-agree-end-cusp-expansion}(2) this is
$\mathcal W^{j-1}_\psi\varphi_{f,x}(g)$.

(1) \emph{Nondegeneracy.} The map $\psi$ is continuous. Let $v$ be a finite
place of $F$ and $t\colon E\to F_v$ an $F$-embedding. Through $\pi_t$ and
$\pi_{t\circ c}$, $E\otimes_FF_v$ is $F_v\times F_v$; let $e\in\mathbf A_E$ have
the component $(1,0)$ there and $0$ at the other places, so that $\pi_t(e)=1$
and $\pi_{t\circ c}(e)=0$. Suppose that $\psi(z)=1$ for every $z\in\mathbf A_E$
with $\pi_{t\circ c}(z)=0$ and components $0$ at the places not above $v$. Every
$ae$ with $a\in E$ is such a $z$, so $\psi(ae)=1$ for all $a\in E$, and $e\in E$
as $\psi$ is standard. Then $t(c(e))=\pi_{t\circ c}(e)=0$, so $e=0$, which
contradicts $\pi_t(e)=1$.

\emph{Whittaker functional.} For automorphic $f$ and $f'$,
$c_U(f+f')=c_Uf+c_Uf'$, so $\varphi_{f+f',1}=\varphi_{f,1}+\varphi_{f',1}$,
and Lemma~\ref{asm-theta-agree-end-cusp-expansion}(2) gives
$\mathcal W_{u,\psi}(f+f')=\mathcal W_{u,\psi}(f)+\mathcal W_{u,\psi}(f')$. In
the same way $c_U(zf)=z\,c_Uf$ and that lemma give
$\mathcal W_{u,\psi}(zf)=z\,\mathcal W_{u,\psi}(f)$ for every function $f$. Let
$f$ be automorphic, and let $h$ and $g$ be as in
Definition~\ref{asm-theta-agree-end-cusp-notions}(8). Then $\det g=1$, and by
Lemma~\ref{asm-theta-agree-end-cusp-frames}(3), $n_0=h\lambda_u(g)^{-1}\in N_U$.
For $g'\in GL_j(\mathbf A_E)$,
$\lambda_u(g')h=\bigl(\lambda_u(g')n_0\lambda_u(g')^{-1}\bigr)\lambda_u(g'g)$,
and the first factor lies in $N_U$
(Lemma~\ref{asm-theta-agree-end-cusp-frames}(2)); so
$c_Uf(\lambda_u(g')h)=\varphi_{f,1}(g'g)$ by
Lemma~\ref{asm-theta-agree-end-cusp-constant}(1). Hence, by
Lemma~\ref{asm-theta-agree-end-cusp-expansion}(2),
$\mathcal W_{u,\psi}(f(\cdot\,h))=\mathcal W^{j-1}_\psi\varphi_{f,1}(g)$. By
parts (1) and (3) of that lemma, $\mathcal W^{j-1}_\psi\varphi_{f,1}$ is adapted
at stage $0$ and
$\mathcal W^{j-1}_\psi\varphi_{f,1}(g)=\psi\bigl(\sum_ig_{i,i+1}\bigr)\mathcal W^{j-1}_\psi\varphi_{f,1}(1)=\psi\bigl(\sum_ig_{i,i+1}\bigr)\mathcal W_{u,\psi}(f)$.

(3) Suppose that $c_Uf(x)\neq0$ for some $x\in U_W$, and put
$\varphi=\varphi_{f,x}$. For $1\leq r\leq j-1$ and $g\in GL_j(\mathbf A_E)$,
Lemma~\ref{asm-theta-agree-end-cusp-constant}(2) for $h=\lambda_u(g)x$ gives
$\langle y\mapsto\varphi(n_{r,j-1}(y)g)\rangle_{I_{r,j-1}}=0$. So $\varphi$ is
cuspidal at stage $j-1$, and $\varphi(1)=c_Uf(x)\neq0$. By
Lemma~\ref{asm-theta-agree-end-cusp-expansion}(4),
$\mathcal W^{j-1}_\psi\varphi(g)\neq0$ for some $g$. By (2) this is
$\mathcal W_{u,\psi}(f(\cdot\,\lambda_u(g)x))$, which is $0$ by hypothesis. So
$c_Uf=0$, and it is the average of $f$ over $N_U\cap U_W(F)$ by
Lemma~\ref{asm-theta-agree-end-cusp-constant}(1).
\end{proof}
% Lean: the first paragraph is ct_pre (pre_top with ct_cont, isLev, th_cont and ct_rat).
% (2) is wh_right (IsCt.right, whit_rt). (1) is isChar, with the element e = s (1, 0) from
% PositivityTheory.KCo.section_statement (PfDictSec.lean) and StdPsi.ann; and isWh: add is
% IsCt.add and whit_add', smul is IsCt.smul and whit_smul, law is par_mem with mid_of_move,
% e2 through ct_left and lev_norm, whit_rt, and
% GlobalTheory.T.whittakerLaw_upperTriangularUnipotent_adelic with whit_pre. (3) is wh_ne,
% with whit_good, ct_van for the cuspidality and IsCt.av. The Lean's functional is chi C psi =
% psi composed with the isomorphism C.ad of A_F tensor E with the adele ring of E.

\begin{lemma}\label{asm-theta-agree-end-cusp-line}
Let $E$ be a CM field with a theta base $B$
(Definition~\ref{asm-theta-agree-models}(1)), let $W$ be a skew-Hermitian
lattice space of rank $d$ and some signature, with matrix $J_W$ and
$h_W(u,u')=u^{\mathsf T}J_W\,c(u')$, let $H'$ be a plane in $W$ over a lattice
space $W_1$ (Definition~\ref{asm-theta-agree-end-notions}(1)), and let $f$ be an
automorphic function on $U_W$ such that $0$ is the constant term of $f$ along
$H'$. Then for every nonzero $x\in E^d$ with $h_W(x,x)=0$, $0$ is the average of
$f$ over $N_{Ex}\cap U_W(F)$ (Definition~\ref{asm-theta-agree-end-kill-notions}(2)
and Lemma~\ref{asm-theta-agree-end-kill-unipotent}).
\end{lemma}
\leanref{Development/WeilBound/ThetaLifts/EigenvalueEndgame/CuspidalAverages/AvLine.lean}{PositivityTheory.ThKEnd.Cusp.av_line}
\leanref{Development/WeilBound/ThetaLifts/EigenvalueEndgame/CuspidalAverages/AvLine.lean}{PositivityTheory.ThKEnd.Cusp.conj_N}
\leanref{Development/WeilBound/ThetaLifts/EigenvalueEndgame/CuspidalAverages/StageB.lean}{PositivityTheory.ThKEnd.Cusp.isAv_transport}
\begin{proof}
Let $H'=(x',y',\iota)$. By Lemma~\ref{asm-theta-agree-end-kill-unipotent},
$N_{Ex'}=N_{H'}$, so $0$ is the average of $f$ over $N_{Ex'}\cap U_W(F)$. As
$h_W(x',y')=1$, $x'\neq0$, and $h_W(x',x')=0$; so by
Lemma~\ref{asm-theta-agree-end-unipotent}(4) there is $\gamma\in U_W(F)$ with
$\gamma x=x'$. Let $s\in N_{Ex}$ and $s'=\gamma s\gamma^{-1}$. Then
$s'x'=\gamma sx=x'$, and if $z\in E^d$ has $h_W(z,x')=0$, then
$h_W(\gamma^{-1}z,x)=0$ and $\gamma^{-1}z\in E^d$, so
$s'z-z=\gamma(s\gamma^{-1}z-\gamma^{-1}z)\in\gamma\mathbf A_Ex=\mathbf A_Ex'$. By
Lemma~\ref{asm-theta-agree-end-kill-unipotent}(1), $s'\in N_{Ex'}$; and in the
same way $\gamma^{-1}N_{Ex'}\gamma\subset N_{Ex}$. So $\epsilon(s)=\gamma s\gamma^{-1}$
is an isomorphism of topological groups from $N_{Ex}$ onto $N_{Ex'}$, and it
maps $N_{Ex}\cap U_W(F)$ onto $N_{Ex'}\cap U_W(F)$, as $\gamma\in U_W(F)$.

Let $\mu$ be a Haar measure on $N_{Ex}$, $\mathcal F$ a fundamental domain for
$N_{Ex}\cap U_W(F)$ with $\mu(\mathcal F)=1$, and $h\in U_W$. The image of $\mu$
under $\epsilon$ is a Haar measure on $N_{Ex'}$, for which $\epsilon(\mathcal F)$
is a fundamental domain for $N_{Ex'}\cap U_W(F)$ of measure $1$. So
$\int_{\mathcal F}f(\epsilon(s)\gamma h)\,d\mu(s)=0$. As
$\epsilon(s)\gamma h=\gamma sh$ and $f$ is automorphic,
$\int_{\mathcal F}f(sh)\,d\mu(s)=0$.
\end{proof}
% Lean: av_line (AvLine.lean, 30 lines): h1 is unipLine, gamma comes from wittVec (Lemma
% asm-theta-agree-end-unipotent(4)), conj_N for the conjugation, and isAv_transport
% (StageB.lean) for the transport of Haar measures and fundamental domains.

\begin{lemma}\label{asm-theta-agree-end-cusp-kill}
Let $G$ be a locally compact, second countable Hausdorff group, $\Gamma\subset G$
a subgroup, $X$ a topological space with a Borel probability measure $\varrho$,
$n\colon X\to G$ a continuous map and $a\in G$, and suppose that
$(\Gamma,n,a)$ has the Siegel and Lipschitz properties
(Definition~\ref{asm-theta-agree-end-kill-notions}(4)). Let $e\colon X\to\mathbf C$
be continuous with $|e|\leq1$, let $b_i\in G$ and $y_i\in\mathbf C$ for integers
$i\geq0$ with $|y_i|\to\infty$, let $(\varphi,c)$ be a recursion pair for
$(\Gamma,n,a,\varrho,0)$ (Definition~\ref{asm-theta-agree-end-kill-notions}(5))
such that $\varphi$ has growth of some order along $a$, and let $V$ be a
neighbourhood of $1$ with
\[
y_0\int_Xe(x)\varphi(n(x)b_iz)\,d\varrho(x)=y_i\int_Xe(x)\varphi(n(x)z)\,d\varrho(x)
\]
for all $z\in V$ and $i\geq0$. Then $\int_Xe(x)\varphi(n(x))\,d\varrho(x)=0$.
\end{lemma}
\leanref{Development/WeilBound/ThetaLifts/EigenvalueEndgame/CuspidalAverages/KillCoef.lean}{PositivityTheory.ThKEnd.Cusp.kill_coef}
\begin{proof}
Let $\mu$ be a Haar measure on $G$, and use $\operatorname{conv}_\eta$ as in
Lemma~\ref{asm-theta-agree-end-kill-induction}, with $l=0$. Apply part (3) of
that lemma to the property $P(\varphi,c)$: for every neighbourhood $V$ of $1$
with the displayed identity for all $z\in V$ and $i\geq0$,
$\int_Xe(x)\varphi(n(x))\,d\varrho(x)=0$.

\emph{Bounded functions.} Let $(\varphi,c)$ be a recursion pair with
$|\varphi|\leq M$, let $V$ be as in $P$, and suppose that
$w=\int_Xe(x)\varphi(n(x))\,d\varrho(x)\neq0$. As $|y_i|\to\infty$, some $i$
has $|y_i|\,|w|>|y_0|M$. But the identity for $z=1$ gives
$|y_i|\,|w|=|y_0|\,\bigl|\int_Xe(x)\varphi(n(x)b_i)\,d\varrho(x)\bigr|\leq|y_0|M$.

\emph{Smoothing.} Let $(\varphi,c)$ be a recursion pair such that every
neighbourhood of $1$ contains a bump $\eta$ with
$P(\operatorname{conv}_\eta\varphi,\operatorname{conv}_\eta c)$, and let $V$ be
as in $P$. Put $w(g)=\int_Xe(x)\varphi(n(x)g)\,d\varrho(x)$. By
Lemma~\ref{asm-theta-agree-end-kill-induction}(1) the integrand is bounded for
$g$ in a compact neighbourhood of a given point, so $w$ is continuous, by
dominated convergence. Choose a neighbourhood $V'$ of $1$ with $V'V'\subset V$.
Let $U$ be a neighbourhood of $1$, and let $\eta$ be a bump in $U\cap V'$ with
$P(\operatorname{conv}_\eta\varphi,\operatorname{conv}_\eta c)$. By
Lemma~\ref{asm-theta-agree-end-kill-induction}(1) for the compact set
$g\operatorname{supp}\eta$ and Fubini's theorem,
$\int_Xe(x)\operatorname{conv}_\eta\varphi(n(x)g)\,d\varrho(x)=\operatorname{conv}_\eta w(g)$
for every $g$. For $z\in V'$ and $y$ in the support of $\eta$, $zy\in V$, so
$y_0w(b_izy)=y_iw(zy)$; multiplying by $\eta(y)$ and integrating gives
$y_0\operatorname{conv}_\eta w(b_iz)=y_i\operatorname{conv}_\eta w(z)$. So the
identity of $P$ holds for $\operatorname{conv}_\eta\varphi$ and $V'$, and
$P(\operatorname{conv}_\eta\varphi,\operatorname{conv}_\eta c)$ gives
$\operatorname{conv}_\eta w(1)=0$. As $U$ was arbitrary, the argument at the
end of the proof of Lemma~\ref{asm-theta-agree-end-kill-induction}(4), for the
continuous function $w$ at $1$, gives $w(1)=0$.

By Lemma~\ref{asm-theta-agree-end-kill-induction}(3), $P(\varphi,c)$ holds for
the given pair, and the given $V$ gives the claim.
\end{proof}
% Lean: kill_coef (KillCoef.lean, 50 lines): kill_ind (Measure.haar) for the property in the
% statement; the bounded case is the first bullet, the smoothing case the second, with w, cw
% (Dat.bound), exists_nhds_one_split for V', fub (integral_integral_swap) and Bump.zero.

\begin{lemma}\label{asm-theta-agree-end-cusp-two}
Let $E$ be a CM field with a theta base $B$
(Definition~\ref{asm-theta-agree-models}(1)), let $W$ be a skew-Hermitian
lattice space of rank $d$ and some signature, with the notation of
Definition~\ref{asm-theta-agree-end-cusp-notions}, let $(u,u')$ be a frame of
rank $2$ in $W$ with span $U$, and let $\psi$ be a standard character of
$\mathbf A_E$. Let $f$ be an automorphic function on $U_W$ with
$|f(h)|\leq C\operatorname{ht}(h)^k$ for some $C$, $k$ and every $h$
(Definition~\ref{asm-theta-agree-parts}), such that $0$ is the average of $f$
over $N_{U_1}\cap U_W(F)$ for the span $U_1$ of the first column of $u$. Let
$x_0\in U_W$, let $V$ be a neighbourhood of $x_0$, and for integers $i\geq0$ let
$b_i\in GL_2(\mathbf A_E)$ and $y_i\in\mathbf C$, with $b_0=1$, $\det b_i=1$ and
$|y_i|\to\infty$, such that every $x\in V$ has some $z\in\mathbf C$ with
$\mathcal W_{u,\psi}(f(\cdot\,\lambda_u(b_i)x))=zy_i$ for every $i$. Then
$\mathcal W_{u,\psi}(f(\cdot\,x_0))=0$.
\end{lemma}
\leanref{Development/WeilBound/ThetaLifts/EigenvalueEndgame/CuspidalAverages/CuspOfLineOfKeys.lean}{PositivityTheory.ThKEnd.Cusp.wh_two}
\begin{proof}
Let $m$ be the restriction of $\lambda_u$ to $SL_2(\mathbf A_E)$, a continuous
homomorphism to $U_W$ with $m(SL_2(E))\subset U_W(F)$
(Lemma~\ref{asm-theta-agree-end-cusp-frames}(1)). Let $a_0$ and $n(\xi)$ be as in
Lemma~\ref{asm-theta-agree-end-kill-siegel}, so that $n(\xi)=n_{1,1}(\xi)$ in
the notation of Definition~\ref{asm-theta-agree-end-cusp-notions}(4), with
$I_{1,1}$ the single place $(1,2)$, and let $\varrho$ be the restriction of
$\mu_{I_{1,1}}$ to $D_{I_{1,1}}$, a probability measure on
$\mathbf A_E^{I_{1,1}}=\mathbf A_E$. By the definition and
Lemma~\ref{asm-theta-agree-end-cusp-constant}(1), for $z\in U_W$,
\[
\mathcal W_{u,\psi}(f(\cdot\,z))=\int_{\mathbf A_E}\psi(-\xi)\,c_Uf(m(n(\xi))z)\,d\varrho(\xi).\tag{$*$}
\]
Put $\varphi(g)=c_Uf(m(g)x_0)$ for $g\in SL_2(\mathbf A_E)$.

\emph{A recursion pair.} The pair $(\varphi,0)$ is a recursion pair for
$(SL_2(E),n,a_0,\varrho,0)$
(Definition~\ref{asm-theta-agree-end-kill-notions}(5)). Indeed, $\varphi$ is
continuous and invariant under $SL_2(E)$, by
Lemma~\ref{asm-theta-agree-end-cusp-constant}(1); for $g\in SL_2(\mathbf A_E)$,
$\int\varphi(n(\xi)g)\,d\varrho(\xi)=\bigl\langle\xi\mapsto c_Uf(\lambda_u(n_{1,1}(\xi))m(g)x_0)\bigr\rangle_{I_{1,1}}$,
which is $0$ by Lemma~\ref{asm-theta-agree-end-cusp-constant}(2) for $r=1$ and
$h=m(g)x_0$; and the recursion for $c=0$ holds.

\emph{Growth.} By Lemma~\ref{asm-theta-agree-end-kill-unipotent}(2) and
Lemma~\ref{asm-theta-agree-end-kill-growth}(2) for $N=N_U$, the average $c_Uf$ of
$f$ satisfies a bound of the same kind, and by part (1) of that lemma, for the
continuous homomorphism $m$, $\varphi$ has growth of some order along $a_0$.

\emph{Conclusion.} The group $SL_2(\mathbf A_E)$ is locally compact, second
countable and Hausdorff, as $\mathbf A_E$ is
(Lemma~\ref{asm-theta-agree-end-kill-siegel}(1)) and $SL_2(\mathbf A_E)$ is
closed in $M_2(\mathbf A_E)$; $n$ is continuous; and by
Lemma~\ref{asm-theta-agree-end-kill-siegel}(2), $(SL_2(E),n,a_0)$ has the
Siegel and Lipschitz properties. Regard the $b_i$ as elements of
$SL_2(\mathbf A_E)$, and let $V'$ be the neighbourhood of $1$ of the $g$ with
$m(g)x_0\in V$. Let $g\in V'$ and $i\geq0$, and let $z$ be as in the statement
for $x=m(g)x_0$. By $(*)$ for $\lambda_u(b_i)x$ and for $x$, as $b_0=1$,
\[
\int\psi(-\xi)\varphi(n(\xi)b_ig)\,d\varrho(\xi)=zy_i,\qquad
\int\psi(-\xi)\varphi(n(\xi)g)\,d\varrho(\xi)=zy_0,
\]
so $y_0$ times the first integral is $y_i$ times the second. By
Lemma~\ref{asm-theta-agree-end-cusp-kill}, with $e(\xi)=\psi(-\xi)$,
$\int\psi(-\xi)\varphi(n(\xi))\,d\varrho(\xi)=0$, and by $(*)$ for $z=x_0$ this
is $\mathcal W_{u,\psi}(f(\cdot\,x_0))$.
\end{proof}
% Lean: wh_two (CuspOfLineOfKeys.lean, 60 lines): ae_facts and sl_facts for the setup, mu and
% F with m1, m2, m3 and av for rho (the Lean shows that avg equals the integral against
% mu.restrict F), n, m, en, er and b', hC', ew for (*), hk from Mod.av and Mod.gr, the record
% D of Dat (ct_cont, ct_rat, per_blk and ct_van for r = 1), and kill_coef.

\begin{lemma}\label{asm-theta-agree-end-cusp-keys-places}
Let $E$ be a CM field with maximal totally real subfield $F$, and let $\pi_t$ be
as in Definition~\ref{asm-theta-agree-models}(1).
\begin{enumerate}
\item For every finite place $v$ of $F$ and every $F$-embedding
  $t\colon E\to F_v$, the map $\pi_t\colon\mathbf A_E\to F_v$ is continuous.
\item For every $z\in\mathbf A_E$, $\pi_t(z)\in\mathcal O_v$ for almost all
  finite places $v$ of $F$ and all $F$-embeddings $t\colon E\to F_v$.
\item The finite places $v$ of $F$ for which there is an $F$-embedding
  $E\to F_v$ form an infinite set.
\end{enumerate}
\end{lemma}
\leanref{Development/WeilBound/ThetaLifts/EigenvalueEndgame/CuspidalAverages/KeyThree.lean}{PositivityTheory.ThKEnd.Cusp.Places}
\leanref{Development/WeilBound/ThetaLifts/EigenvalueEndgame/CuspidalAverages/PlaContinuous.lean}{PositivityTheory.ThKEnd.Cusp.places}
\begin{proof}
(1) The map $\pi_t$ is the projection to the component $E\otimes_FF_v$ at $v$,
which is continuous, followed by the $F_v$-linear map
$x\otimes s\mapsto t(x)s$.

(2) If the claim holds for $z$ and $z'$, it holds for $z+z'$, as a union of two
finite sets is finite. So let $z=a\otimes x$ with $a\in\mathbf A_F$ and
$x\in E$, and let $N\neq0$ be an integer with $Nx$ integral over $\mathbf Z$. For
almost all $v$, $a_v\in\mathcal O_v$ and $N^{-1}\in\mathcal O_v$. For these $v$
and every $t$, $t(Nx)$ is integral over $\mathbf Z$, so it lies in
$\mathcal O_v$, and $\pi_t(z)=a_vN^{-1}t(Nx)\in\mathcal O_v$.

(3) These are the places of $F$ that split in the quadratic extension $E$, and
there are infinitely many of them by the Chebotarev density theorem
\cite[Chapter~VII, \S13]{Neu99}.
\end{proof}
% Lean: Cusp.places (PlaContinuous.lean, 60 lines): (1) is pla_continuous, (2) is pla_int
% (induction on tensors, with exists_integral_multiple) and pla_le_one_iff, and (3) is
% GlobalTheory.Orb.exists_split_place_notMem, which the Lean proves its own way. The third
% clause of Places, on unramified characters, is used in the proof of Lemma
% asm-theta-agree-end-cusp-keys(2), where it is the definition of being unramified.

\begin{definition}\label{asm-theta-agree-end-cusp-keys-family}
Let $E$ be a CM field with maximal totally real subfield $F$, let $W$ be a
lattice space of rank $d$ (Definition~\ref{asm-theta-agree-doubling}(1)), let
$v$ be a finite place of $F$, $t\colon E\to F_v$ an $F$-embedding and $n\geq0$,
and let $\xi$ and $\alpha$ be homomorphisms from $F_v^\times$ to
$\mathbf C^\times$. A \emph{local family} for $(t,n)$ with laws $(\xi,\alpha)$
is a map $S$ from the functions on $M_{n\times d}(F_v)$ to the functions on
$U_W$ with the following properties, for all Schwartz-Bruhat functions $\phi$
and $\phi'$ on $M_{n\times d}(F_v)$.
\begin{enumerate}
\item $S(\phi)$ is automorphic (Definition~\ref{asm-theta-agree-end-notions}(3)),
  $S(\phi+\phi')=S(\phi)+S(\phi')$, and $S(z\phi'')=z\,S(\phi'')$ for every
  function $\phi''$ and every $z\in\mathbf C$.
\item $S\bigl(Y\mapsto\phi(\gamma Y)\bigr)=\xi(\det\gamma)\,S(\phi)$ for
  $\gamma\in GL_n(F_v)$.
\item $S(\phi)(xh)=\alpha(\det h_0)\,S\bigl(Y\mapsto\phi(Yh_0)\bigr)(x)$ for
  $x\in U_W$ and $h\in U_W^{(v)}$ (Definition~\ref{asm-theta-agree-models}(2))
  with $\pi_t(h)=h_0$.
\end{enumerate}
\end{definition}
\leanref{Development/WeilBound/ThetaLifts/EigenvalueEndgame/CuspidalAverages/ActsOfReads.lean}{PositivityTheory.ThKEnd.Cusp.Key.IsLoc}

\begin{lemma}\label{asm-theta-agree-end-cusp-keys-local}
Let $(X,W,M,\chi_X,\chi_W,\omega)$ be a Weil pair for a theta base $B$ of a CM
field $E$ (Definition~\ref{asm-theta-agree-triple}(1)), with $X$ and $W$ of
ranks $n$ and $d$ and theta integral $I$, and put
$I_0=\{1,\ldots,n\}\times\{1,\ldots,d\}$. For a function $\Phi$ on
$\mathbf A^{I_0}$, a finite place $v$ of $F$ and a function $\phi$ on
$F_v^{I_0}$, let $\Phi[v\mapsto\phi](x)=\phi(x_v)\Phi(x^v)$, where $x^v$ is $x$
with its component at $v$ replaced by $0$. Let $\mathcal C$ be a space of
Schwartz-Bruhat functions on $\mathbf A^{I_0}$ such that every
$\Phi\in\mathcal C$ converges and $I(\Phi)$ is automorphic
(Definition~\ref{asm-theta-agree-end-notions}(3)). Let $\Phi\in\mathcal C$ be
such that for almost all finite places $v$, $\Phi[v\mapsto\mathrm{vac}_v]=\Phi$
and $\Phi[v\mapsto\phi]\in\mathcal C$ for every Schwartz-Bruhat function $\phi$
on $F_v^{I_0}$. Then for almost all finite places $v$ and every $F$-embedding
$t\colon E\to F_v$, with $\nu=(\chi_W)_t$ and $\alpha=(\chi_X)_t$
(Definition~\ref{asm-theta-agree-coinvariants-local}(4)), there are a linear
model $\iota$ of the second kind for $(t,1,1)$
(Definition~\ref{asm-theta-agree-models}(5)) and a local family $S$ for $(t,n)$
with laws $\bigl((\nu|\cdot|^{d/2})^{-1},\alpha|\cdot|^{n/2}\bigr)$
(Definition~\ref{asm-theta-agree-end-cusp-keys-family}) such that
$\iota(\mathrm{vac}_v)$ is a Schwartz-Bruhat function and
$S(\iota(\mathrm{vac}_v))=I(\Phi)$.
\end{lemma}
\leanref{Development/WeilBound/ThetaLifts/EigenvalueEndgame/CuspidalAverages/ActsOfReads.lean}{PositivityTheory.ThKEnd.Cusp.Key.exists_loc}
\leanref{Development/WeilBound/ThetaLifts/EigenvalueEndgame/CuspidalAverages/ActsOfReads.lean}{PositivityTheory.ThKEnd.Cusp.Key.insertLocal}
\begin{proof}
Let $v$ be a finite place with $\Phi[v\mapsto\mathrm{vac}_v]=\Phi$ and
$\Phi[v\mapsto\phi]\in\mathcal C$ for every Schwartz-Bruhat $\phi$, and let
$t\colon E\to F_v$ be an $F$-embedding. By
Lemma~\ref{asm-theta-agree-coinvariants-dictionary}(1) there is a linear model
$\iota$ of the second kind for $(t,1,1)$. It is bijective between the
Schwartz-Bruhat functions, so $\iota(\mathrm{vac}_v)$ is one, and there is a
linear map $j$ from the functions on $M_{n\times d}(F_v)$ to the functions on
$F_v^{I_0}$ such that $j(\phi)$ is Schwartz-Bruhat and $\iota(j(\phi))=\phi$ for
Schwartz-Bruhat $\phi$, and $j(\iota(\phi_0))=\phi_0$ for Schwartz-Bruhat
$\phi_0$: extend the inverse of $\iota$ on the Schwartz-Bruhat functions
linearly. Put $S(\phi)=I(\Phi[v\mapsto j(\phi)])$. Then
$S(\iota(\mathrm{vac}_v))=I(\Phi[v\mapsto\mathrm{vac}_v])=I(\Phi)$.

\emph{Property (1).} For Schwartz-Bruhat $\phi$,
$\Phi[v\mapsto j(\phi)]\in\mathcal C$, so $S(\phi)$ is automorphic. The map
$\phi_0\mapsto\Phi[v\mapsto\phi_0]$ is linear, and so is $j$. The theta
function is additive on Schwartz-Bruhat functions
(Lemma~\ref{asm-theta-agree-convergence-theta}(2)), so the theta integral of a
sum of two functions that converge is the sum of their theta integrals; and the
theta integral of $z\Psi$ is $z$ times that of $\Psi$, also when one of the
integrals is not defined, as both are then $0$. This gives the additivity and
the homogeneity of $S$.

\emph{Property (2).} Let $\phi$ be Schwartz-Bruhat, $\gamma\in GL_n(F_v)$ and
$g_0=\gamma^{\mathsf T}$. By Lemma~\ref{asm-theta-agree-coinvariants-dictionary}(2)
and~(3) there is $g\in U_X^{(v)}$ with $\pi_t(g)=g_0$, and $\omega(g,1)$ reads as
$\omega_{\nu,\alpha}(g_0,1)$ through $\iota$. So $\omega(g,1)$ acts at $v$ by
$j(\phi)\mapsto\phi_1$ for a Schwartz-Bruhat $\phi_1$ with
$\iota(\phi_1)=\omega_{\nu,\alpha}(g_0,1)\phi=\nu(\det\gamma)|\det\gamma|^{d/2}\phi(\gamma\,\cdot\,)$,
and $\phi_1=j(\iota(\phi_1))=\nu(\det\gamma)|\det\gamma|^{d/2}j(\phi(\gamma\,\cdot\,))$.
The function $\Phi[v\mapsto j(\phi)]$ is $x\mapsto j(\phi)(x_v)r(x)$ with
$r(x)=\Phi(x^v)$, which does not depend on the component at $v$, and it is
Schwartz-Bruhat; so $\omega(g,1)\Phi[v\mapsto j(\phi)]=\Phi[v\mapsto\phi_1]$, by
the definition of acting at $v$. By
Lemma~\ref{asm-theta-agree-anisotropic-invariance}, the theta integrals of these
two functions agree, so
$\nu(\det\gamma)|\det\gamma|^{d/2}S(\phi(\gamma\,\cdot\,))=S(\phi)$.

\emph{Property (3).} Let $h\in U_W^{(v)}$ with $\pi_t(h)=h_0$. As $\omega$ is a
representation, $\theta_\omega(\Psi)(g,xh)=\theta_\omega(\omega(1,h)\Psi)(g,x)$,
so $I(\Psi)(xh)=I(\omega(1,h)\Psi)(x)$ for every $\Psi$. By
Lemma~\ref{asm-theta-agree-coinvariants-dictionary}(3), $\omega(1,h)$ reads as
$\omega_{\nu,\alpha}(1,h_0)$ through $\iota$, and the argument of (2) gives
$\omega(1,h)\Phi[v\mapsto j(\phi)]=\alpha(\det h_0)|\det h_0|^{n/2}\Phi[v\mapsto j(\phi(\cdot\,h_0))]$.
So $S(\phi)(xh)=\alpha(\det h_0)|\det h_0|^{n/2}S(\phi(\cdot\,h_0))(x)$.
\end{proof}
% Lean: exists_loc (ActsOfReads.lean): Phi[v -> phi] is insertLocal v Phi phi, the hypothesis
% on Phi is Phi in slots M.vac C (mem_placewise); iota and the readings come from
% dictionary_extended (DictX, Lemma asm-theta-agree-coinvariants-dictionary), j from
% kd_exists_inv, and S phi = M.I om (insertLocal v Phi (j phi)). The additivity is thetaSum_add
% and integral_add with the convergence hc, the homogeneity I_smul, the left law
% acts_of_reads, insertLocal_acts, omB_left and ThKF.I_left (Lemma
% asm-theta-agree-anisotropic-invariance), and the right law I_trans, acts_of_reads,
% insertLocal_acts and omB_right.

\begin{lemma}\label{asm-theta-agree-end-cusp-keys-complement}
Let $K$ be a field, and let $u\in M_{d\times j}(K)$ and $p\in M_{j\times d}(K)$
with $pu=1$. Then there are an injective map
$\kappa\colon\{1,\ldots,j\}\to\{1,\ldots,d\}$ and $C\in GL_d(K)$ such that
column $\kappa(i)$ of $C$ is column $i$ of $u$ for every $i$, and $pC$ has the
entry $1$ at the places $(i,\kappa(i))$ and $0$ elsewhere.
\end{lemma}
\leanref{Development/WeilBound/ThetaLifts/EigenvalueEndgame/CuspidalAverages/ExistsCond.lean}{PositivityTheory.ThKEnd.Cusp.Key.exists_compl}
\leanref{Development/WeilBound/ThetaLifts/EigenvalueEndgame/CuspidalAverages/ActsOfReads.lean}{PositivityTheory.ThKEnd.Cusp.Key.frame_compl}
\begin{proof}
As $pu=1$, $p$ maps $K^d$ onto $K^j$, so its kernel has dimension $d-j$. Let
$w\in M_{d\times(d-j)}(K)$ have as columns a basis of the kernel, let $C$ be the
matrix whose columns are those of $u$ followed by those of $w$, and put
$\kappa(i)=i$. If $ux+wy=0$, then applying $p$ gives $x=0$, and then $y=0$; so
$C$ is invertible. Finally $pC$ is $pu=1$ followed by $pw=0$.
\end{proof}
% Lean: exists_compl (ExistsCond.lean), with a basis of ker p (Module.finBasis) and the
% reindexing finSumFinEquiv; frame_compl applies it to p = (J c(u'))^T for a frame.

\begin{lemma}\label{asm-theta-agree-end-cusp-keys-form}
Let $E$ be a CM field with a theta base $B$
(Definition~\ref{asm-theta-agree-models}(1)), let $W$ be a skew-Hermitian
lattice space of rank $d$ and some signature, with the notation of
Definition~\ref{asm-theta-agree-end-cusp-notions}, let $(u,u')$ be a frame of
rank $j$ in $W$ with columns $u_1,\ldots,u_j$ of $u$, let $\chi$ be a character
of $\mathbf A_E$, and let $\Lambda$ be a Whittaker functional for $(u,\chi)$.
Let $v$ be a finite place of $F$ and $t\colon E\to F_v$ an $F$-embedding. For
$s\in F_v$ let $\tilde s\in\mathbf A_E$ be the element at $v$ with
$\pi_t(\tilde s)=s$ and $\pi_{t\circ c}(\tilde s)=0$
(Lemma~\ref{asm-theta-agree-dictionary-split}(1)), and put
$\psi_t(s)=\chi(\tilde s)$. Let $n\geq0$, let $S$ be a local family for $(t,n)$
with laws $(\xi,\alpha)$ (Definition~\ref{asm-theta-agree-end-cusp-keys-family}),
let $\kappa\colon\{1,\ldots,j\}\to\{1,\ldots,d\}$ be a map and $C\in GL_d(F_v)$ a
matrix whose column $\kappa(i)$ is $t(u_i)$ for every $i$, and let $g\in U_W$.
For $Y\in M_{n\times d}(F_v)$ write $Y_b$ for its column $b$. Then there is a
linear form $l$ on the functions on $M_{n\times d}(F_v)$ such that:
\begin{enumerate}
\item $l\bigl(Y\mapsto\phi(\gamma Y)\bigr)=\xi(\det\gamma)\,l(\phi)$ for every
  Schwartz-Bruhat $\phi$ and every $\gamma\in GL_n(F_v)$;
\item for $1\leq i<j$, every set $T$ of columns that contains
  $\kappa(1),\ldots,\kappa(i)$, every Schwartz-Bruhat $\phi$ and every
  $\rho\in F_v^d$ with $\rho_s=0$ for $s\in T$,
  $l\bigl(Y\mapsto\phi(Y+Y_{\kappa(i)}\rho^{\mathsf T})\bigr)=\psi_t(\rho_{\kappa(i+1)})\,l(\phi)$;
\item for every $k\in U_W^{(v)}$ and every Schwartz-Bruhat $\phi$,
  \[
  \Lambda\bigl(S(\phi)(\cdot\,kg)\bigr)=\alpha\bigl(\det(\pi_t(k)\pi_t(g))\bigr)\,l\bigl(Y\mapsto\phi(YC^{-1}\pi_t(k)\pi_t(g))\bigr).
  \]
\end{enumerate}
\end{lemma}
\leanref{Development/WeilBound/ThetaLifts/EigenvalueEndgame/CuspidalAverages/ActsOfReads.lean}{PositivityTheory.ThKEnd.Cusp.Key.exists_l}
\leanref{Development/WeilBound/ThetaLifts/EigenvalueEndgame/CuspidalAverages/ActsOfReads.lean}{PositivityTheory.ThKEnd.Cusp.Key.ex_at}
\leanref{Development/WeilBound/ThetaLifts/EigenvalueEndgame/CuspidalAverages/ActsOfReads.lean}{PositivityTheory.ThKEnd.Cusp.Key.ex_split}
\leanref{Development/WeilBound/ThetaLifts/EigenvalueEndgame/CuspidalAverages/ActsOfReads.lean}{PositivityTheory.ThKEnd.Cusp.Key.transv}
\leanref{Development/WeilBound/ThetaLifts/EigenvalueEndgame/CuspidalAverages/ActsOfReads.lean}{PositivityTheory.ThKEnd.Cusp.Key.wh_transv}
\begin{proof}
Write $J=J_W$ and $A=t(J)$. As in the proof of
Lemma~\ref{asm-theta-agree-coinvariants-dictionary}, $J$ is invertible, and
$t(c(J))=-A^{\mathsf T}$ as $c(J)^{\mathsf T}=-J$; so
Lemma~\ref{asm-theta-agree-dictionary-split}(4) applies to $J$ with
$\varepsilon=-1$. Four facts are used.
\begin{enumerate}
\item[(i)] For every $h_0\in GL_d(F_v)$ there is $k\in U_W^{(v)}$ with
  $\pi_t(k)=h_0$, by Lemma~\ref{asm-theta-agree-dictionary-split}(4).
\item[(ii)] $g=rg_v$ with $g_v\in U_W^{(v)}$, $\pi_t(g_v)=\pi_t(g)$, and
  $r\in U_W$ commuting with every element of $U_W^{(v)}$. Indeed, take $g_v$ by
  (i) and put $r=gg_v^{-1}$. Then $\pi_t(r)=1$, and as
  $\pi_t(r)^{\mathsf T}A\,\pi_{t\circ c}(r)=A$ by
  Lemma~\ref{asm-theta-agree-dictionary-split}(4), $\pi_{t\circ c}(r)=1$. For
  $k\in U_W^{(v)}$, $rk-kr=r(k-1)-(k-1)r$ is at $v$, and its images under
  $\pi_t$ and $\pi_{t\circ c}$ are $0$; so $rk=kr$ by
  Lemma~\ref{asm-theta-agree-dictionary-split}(1).
\item[(iii)] Let $k\in U_W^{(v)}$ with $\pi_t(k)=1+t(x)w^{\mathsf T}$ for some
  $x\in E^d$ and $w\in F_v^d$, and let $z\in E^d$ with $h_W(x,z)=0$. Then
  $kz=z+\widetilde{w\cdot t(z)}\,x$. Indeed, both sides minus $z$ are at $v$,
  and their images under $\pi_t$ are $t(x)\,(w\cdot t(z))$. Let
  $y=\pi_{t\circ c}(z)=t(c(z))$. As $t(x)^{\mathsf T}Ay=t(h_W(x,z))=0$,
  $\pi_t(k)^{\mathsf T}Ay=Ay$; with
  $\pi_t(k)^{\mathsf T}A\,\pi_{t\circ c}(k)=A$ and the invertibility of
  $\pi_t(k)^{\mathsf T}A$ this gives $\pi_{t\circ c}(k)y=y$, so both images
  under $\pi_{t\circ c}$ are $0$. By
  Lemma~\ref{asm-theta-agree-dictionary-split}(1) the two sides agree.
\item[(iv)] Let $f$ be automorphic, $1\leq i<j$, and $k\in U_W^{(v)}$ with
  $\pi_t(k)=1+t(u_i)w^{\mathsf T}$, where $w\cdot t(u_b)=0$ for $b\leq i$. Then
  $\Lambda(f(\cdot\,k))=\psi_t(w\cdot t(u_{i+1}))\,\Lambda(f)$. Indeed,
  $h_W(u_i,u_b)=0$ for every $b$, as $u^{\mathsf T}J\,c(u)=0$, so by (iii)
  $ku=uG$, where $G\in M_j(\mathbf A_E)$ is $1$ plus the matrix whose row $i$ has
  the entries $\widetilde{w\cdot t(u_b)}$ and whose other rows are $0$. It is
  upper triangular with diagonal entries $1$, as $w\cdot t(u_b)=0$ for
  $b\leq i$. If $z\in E^d$ has $h_W(z,U_u)=0$, then
  $h_W(u_i,z)=-c(h_W(z,u_i))=0$, as $c(J)^{\mathsf T}=-J$, so
  $kz-z\in\mathbf A_EU_u$ by (iii). The law of $\Lambda$ gives
  $\Lambda(f(\cdot\,k))=\chi(G_{i,i+1})\Lambda(f)$, and
  $G_{i,i+1}=\widetilde{w\cdot t(u_{i+1})}$.
\end{enumerate}

\emph{The form.} Write $g=rg_v$ as in (ii). For a Schwartz-Bruhat $\phi$ put
$L(\phi)=\Lambda\bigl(S(\phi(\cdot\,C))(\cdot\,r)\bigr)$, where
$\phi(\cdot\,C)$ is $Y\mapsto\phi(YC)$. Right translates of automorphic
functions are automorphic, so by the properties of $S$ and of $\Lambda$, $L$ is
linear on the space of Schwartz-Bruhat functions. Let $l$ be a linear form on
the functions on $M_{n\times d}(F_v)$ that extends $L$.

(1) For $\gamma\in GL_n(F_v)$,
$S\bigl(Y\mapsto\phi(\gamma YC)\bigr)=\xi(\det\gamma)S(\phi(\cdot\,C))$, so
$l(\phi(\gamma\,\cdot\,))=\xi(\det\gamma)\,l(\phi)$ by the homogeneity of
$\Lambda$.

(2) Let $\epsilon_{\kappa(i)}\in F_v^d$ be the standard basis vector and
$N=1+\epsilon_{\kappa(i)}\rho^{\mathsf T}$. Then
$Y+Y_{\kappa(i)}\rho^{\mathsf T}=YN$ and $\det N=1+\rho_{\kappa(i)}=1$. By (i)
let $k\in U_W^{(v)}$ have $\pi_t(k)=CNC^{-1}$. As
$C\epsilon_{\kappa(b)}=t(u_b)$, $\pi_t(k)=1+t(u_i)w^{\mathsf T}$ with
$w^{\mathsf T}=\rho^{\mathsf T}C^{-1}$, and $w\cdot t(u_b)=\rho_{\kappa(b)}$,
which is $0$ for $b\leq i$. The function $f=S(\phi(\cdot\,C))(\cdot\,r)$ is
automorphic, and by (iv),
$\Lambda(f(\cdot\,k))=\psi_t(\rho_{\kappa(i+1)})\,l(\phi)$. On the other hand,
$r$ and $k$ commute, and property (3) of $S$ for $k$, with
$\det(CNC^{-1})=1$, gives
$f(xk)=S(\phi(\cdot\,C))(xrk)=S(\phi(\cdot\,CN))(xr)$; so
$\Lambda(f(\cdot\,k))=l(\phi(\cdot\,N))$.

(3) Let $k\in U_W^{(v)}$. Then $kg_v\in U_W^{(v)}$,
$\pi_t(kg_v)=\pi_t(k)\pi_t(g)$ and $xkg=xrkg_v$. Property (3) of $S$ for $kg_v$
gives
$S(\phi)(xkg)=\alpha\bigl(\det(\pi_t(k)\pi_t(g))\bigr)S(\phi(\cdot\,\pi_t(k)\pi_t(g)))(xr)$,
and $\phi(Y\pi_t(k)\pi_t(g))=\phi'(YC)$ for
$\phi'(Y)=\phi(YC^{-1}\pi_t(k)\pi_t(g))$. The homogeneity of $\Lambda$ gives
(3).
\end{proof}
% Lean: exists_l (ActsOfReads.lean): (i) is ex_at (IsSec.exists_ug), (ii) is ex_split with
% snd_fix and IsSec.extM, (iii) is transv with snd_fix, (iv) is wh_transv with form_swap and
% the law field of IsWh. L is a linear map on the subspace SB, extended by
% LinearMap.exists_extend; (2) uses colOp, d1 (the determinant of N), ex_at for c N c^-1, hw
% and the right law of IsLoc; (3) uses uAt_mul and the right law. Here s -> s~ is sA v sec,
% and psi_t is chi.compAddMonoidHom (sA v sec).

\begin{lemma}\label{asm-theta-agree-end-cusp-keys-three}
Let $k$ be a nonarchimedean local field with ring of integers $\mathcal O$, let
$n,d\geq0$, put $M=M_{n\times d}(k)$, and write $Y_b$ for column $b$ of
$Y\in M$. Let $\psi\neq1$ be a homomorphism from $k$ to $\mathbf C^\times$,
$\xi$ a homomorphism from $k^\times$ to $\mathbf C^\times$, and $l$ a linear
form on the functions on $M$ with
$l\bigl(Y\mapsto\phi(\gamma Y)\bigr)=\xi(\det\gamma)\,l(\phi)$ for every
Schwartz-Bruhat $\phi$ and every $\gamma\in GL_n(k)$. For distinct columns $i$
and $i'$ and a set $T$ of columns, say that $l$ has the \emph{column law}
$(T,i,i')$ if
$l\bigl(Y\mapsto\phi(Y+Y_i\rho^{\mathsf T})\bigr)=\psi(\rho_{i'})\,l(\phi)$ for
every Schwartz-Bruhat $\phi$ and every $\rho\in k^d$ with $\rho_s=0$ for
$s\in T$. Let $i_0$ and $i_1$ be distinct columns such that $l$ has the column
law $(\{i_0\},i_0,i_1)$.
\begin{enumerate}
\item If $\phi$ is a Schwartz-Bruhat function with $\phi(Y)=0$ for every $Y$
  whose columns $Y_{i_0}$ and $Y_{i_1}$ are linearly dependent, then
  $l(\phi)=0$.
\item If $i_2$ is a column different from $i_0$ and $i_1$ and $l$ also has the
  column law $(\{i_0,i_1\},i_1,i_2)$, then $l(\phi)=0$ for every
  Schwartz-Bruhat function $\phi$.
\end{enumerate}
\end{lemma}
\leanref{Development/WeilBound/ThetaLifts/EigenvalueEndgame/CuspidalAverages/KeyThree.lean}{PositivityTheory.ThKEnd.Cusp.LeftLaw}
\leanref{Development/WeilBound/ThetaLifts/EigenvalueEndgame/CuspidalAverages/KeyThree.lean}{PositivityTheory.ThKEnd.Cusp.ColLaw}
\leanref{Development/WeilBound/ThetaLifts/EigenvalueEndgame/CuspidalAverages/KeyThree.lean}{PositivityTheory.ThKEnd.Cusp.LocSupp}
\leanref{Development/WeilBound/ThetaLifts/EigenvalueEndgame/CuspidalAverages/KeyThree.lean}{PositivityTheory.ThKEnd.Cusp.LocThree}
\leanref{Development/WeilBound/ThetaLifts/EigenvalueEndgame/CuspidalAverages/LocalFieldT2Space.lean}{PositivityTheory.ThKEnd.Cusp.locSupp}
\leanref{Development/WeilBound/ThetaLifts/EigenvalueEndgame/CuspidalAverages/LocalFieldT2Space.lean}{PositivityTheory.ThKEnd.Cusp.locThree}
\leanref{Development/WeilBound/ThetaLifts/EigenvalueEndgame/CuspidalAverages/ZariskiLocallyAffineLine.lean}{PositivityTheory.ThKEnd.Cusp.small}
\leanref{Development/WeilBound/ThetaLifts/EigenvalueEndgame/CuspidalAverages/ZariskiLocallyAffineLine.lean}{PositivityTheory.ThKEnd.Cusp.kills_of_local}
\leanref{Development/WeilBound/ThetaLifts/EigenvalueEndgame/CuspidalAverages/ZariskiLocallyAffineLine.lean}{PositivityTheory.ThKEnd.Cusp.kills_interior}
\leanref{Development/WeilBound/ThetaLifts/EigenvalueEndgame/CuspidalAverages/ZariskiLocallyAffineLine.lean}{PositivityTheory.ThKEnd.Cusp.dead}
\begin{proof}
Say that $l$ \emph{vanishes on} a set $U\subset M$ if $l(\phi)=0$ for every
Schwartz-Bruhat $\phi$ with $\phi(Y)=0$ for $Y\notin U$. For a compact open
subgroup $\Lambda\subset M$, $a\in k^\times$ and $x\in M$, the indicator
function $\mathbf 1_{x+a\Lambda}$ of $x+a\Lambda$ is Schwartz-Bruhat. Put
$\Lambda_0=M_{n\times d}(\mathcal O)$, and fix $\tau\in k$ with
$\psi(\tau)\neq1$.

\emph{(a) Small cosets.} Let $\Lambda\subset M$ be a compact open subgroup and
$\phi$ a Schwartz-Bruhat function such that every $Y_0$ with $\phi(Y_0)\neq0$ has
a neighbourhood $\Omega$ with $l(\mathbf 1_{x+a\Lambda})=0$ whenever
$a\in k^\times$ and $x+a\Lambda\subset\Omega$. Then $l(\phi)=0$. Indeed, the set
$K$ where $\phi\neq0$ is compact, and so is $\Lambda$; so there is
$a\in k^\times$ such that for every $x\in K$, $\phi$ is constant on $x+a\Lambda$
and $x+a\Lambda$ lies in one of these neighbourhoods $\Omega$. Finitely many
cosets of $a\Lambda$ cover $K$, and $\phi$ is the sum of the
$\phi(x)\mathbf 1_{x+a\Lambda}$ over representatives $x\in K$ of the distinct
cosets that meet $K$; so $l(\phi)=0$.

\emph{(b) Local vanishing.} If every point of $U$ has a neighbourhood on which
$l$ vanishes, then $l$ vanishes on $U$. This is (a) for $\Lambda=\Lambda_0$, as
$\mathbf 1_{x+a\Lambda_0}$ vanishes outside $\Omega$ when
$x+a\Lambda_0\subset\Omega$.

\emph{(c) Symmetries.} Let $l$ vanish on $U$, let $\Lambda\subset M$ be a
compact open subgroup, and let $\Omega\subset M$ be a set such that every
$Y\in\Omega$ outside $U$ has a $k$-linear map $s\colon M\to M$ with $s(Y)=Y$,
$s^{-1}(\Lambda)=\Lambda$ and $l(\phi\circ s)=c\,l(\phi)$ for every
Schwartz-Bruhat $\phi$, for some $c\neq1$. Then $l$ vanishes on the interior of
$\Omega$. Indeed, by (a) it suffices that $l(\mathbf 1_{x+a\Lambda})=0$ when
$x+a\Lambda\subset\Omega$. If $x+a\Lambda\subset U$, this holds as $l$ vanishes
on $U$. Otherwise let $y\in x+a\Lambda$ lie outside $U$, and let $s$ be as above
for $y$. As $s$ is linear, fixes $y$ and has $s^{-1}(a\Lambda)=a\Lambda$,
$s(z)\in y+a\Lambda$ if and only if $z\in y+a\Lambda$, and $y+a\Lambda=x+a\Lambda$.
So $\mathbf 1_{x+a\Lambda}\circ s=\mathbf 1_{x+a\Lambda}$, and
$l(\mathbf 1_{x+a\Lambda})=c\,l(\mathbf 1_{x+a\Lambda})$ gives
$l(\mathbf 1_{x+a\Lambda})=0$.

\emph{(d) Unipotent symmetries.} Let $l$ vanish on $U$, and let $Y_0\in M$. For
$Y$ near $Y_0$ let $A(Y)$ and $B(Y)$ be $k$-linear endomorphisms of $M$ such
that $(Y,Z)\mapsto A(Y)Z$ and $(Y,Z)\mapsto B(Y)Z$ are continuous at $(Y_0,Z)$
for every $Z$, with $A(Y_0)^2=B(Y_0)^2=0$ and $A(Y_0)B(Y_0)=B(Y_0)A(Y_0)$.
Suppose that every $Y$ near $Y_0$ lies in $U$ or has $A(Y)^2=B(Y)^2=0$,
$s_Y(Y)=Y$ for $s_Y=(1+A(Y))(1+B(Y))$, and $l(\phi\circ s_Y)=c_Y\,l(\phi)$ for
every Schwartz-Bruhat $\phi$, for some $c_Y\neq1$. Then $l$ vanishes on a
neighbourhood of $Y_0$. Indeed, put $A_0=A(Y_0)$ and $B_0=B(Y_0)$, and let
$\Lambda$ be the set of the $Z\in\Lambda_0$ with $A_0Z$, $B_0Z$ and $A_0B_0Z$ in
$\Lambda_0$, a compact open subgroup. The relations give $A_0\Lambda\subset\Lambda$
and $B_0\Lambda\subset\Lambda$. As $\Lambda$ is compact and open, the continuity
gives $A(Y)\Lambda\subset\Lambda$ and $B(Y)\Lambda\subset\Lambda$ for all $Y$ in a
neighbourhood $\Omega$ of $Y_0$, which can be taken inside the neighbourhood of
the hypothesis. If $D$ is linear with $D\Lambda\subset\Lambda$ and $D^2=0$, then
$Z+DZ\in\Lambda$ if and only if $Z\in\Lambda$, as $Z=(Z+DZ)-D(Z+DZ)$. So
$s_Y^{-1}(\Lambda)=\Lambda$ for $Y\in\Omega$ outside $U$, and by (c) $l$
vanishes on the interior of $\Omega$.

\emph{(e) Independent columns.} Let $U_1$ be the set of the $Y$ with $Y_{i_0}$
and $Y_{i_1}$ linearly independent. By (b), (1) follows once every $Y_0\in U_1$
has a neighbourhood on which $l$ vanishes. Choose $\mu,\nu\in k^n$ with
$\mu\cdot Y_{0,i_0}=0$, $\mu\cdot Y_{0,i_1}=1$ and $\nu\cdot Y_{0,i_0}=1$. Near
$Y_0$,
$D(Y)=(\nu\cdot Y_{i_0})(\mu\cdot Y_{i_1})-(\mu\cdot Y_{i_0})(\nu\cdot Y_{i_1})$
is nonzero, and
$\lambda(Y)=D(Y)^{-1}\bigl((\nu\cdot Y_{i_0})\mu-(\mu\cdot Y_{i_0})\nu\bigr)$
is continuous, with $\lambda(Y)\cdot Y_{i_0}=0$ and $\lambda(Y)\cdot Y_{i_1}=1$.
Put $N(Y)=-\tau\,Y_{i_0}\lambda(Y)^{\mathsf T}\in M_n(k)$, $A(Y)Z=N(Y)Z$,
$\rho(Y)=\tau\,Y^{\mathsf T}\lambda(Y)\in k^d$ and
$B(Y)Z=Z_{i_0}\rho(Y)^{\mathsf T}$. Then $N(Y)^2=0$, $\rho(Y)_{i_0}=0$ and
$\rho(Y)_{i_1}=\tau$, so $A(Y)^2=B(Y)^2=0$ and $A(Y)B(Y)=B(Y)A(Y)$, and
\[
s_Y(Y)=Y+Y_{i_0}\rho(Y)^{\mathsf T}+N(Y)\bigl(Y+Y_{i_0}\rho(Y)^{\mathsf T}\bigr)=Y+Y_{i_0}\rho(Y)^{\mathsf T}-\tau\,Y_{i_0}\lambda(Y)^{\mathsf T}Y=Y.
\]
As $\det(1+N(Y))=1-\tau\,\lambda(Y)\cdot Y_{i_0}=1$, the law of $l$ for
$\gamma=1+N(Y)$ gives $l(\phi\circ(1+A(Y)))=l(\phi)$, and the column law
$(\{i_0\},i_0,i_1)$ for $\phi\circ(1+A(Y))$ and $\rho(Y)$ gives
$l(\phi\circ s_Y)=\psi(\tau)\,l(\phi)$. So (d), with $U=\emptyset$, applies.

\emph{(f) A nonzero column.} Under the hypotheses of (2), $l$ vanishes on the set
$U_2$ of the $Y$ with $Y_{i_0}\neq0$. Let $Y_0\in U_2$, let $a_0$ be a row with
$(Y_0)_{a_0i_0}\neq0$, and let $\epsilon_{i_2}\in k^d$ be the standard basis
vector. For $Y$ near $Y_0$, with $Y_{a_0i_0}\neq0$, put
$A(Y)Z=\tau\,Z_{i_1}\epsilon_{i_2}^{\mathsf T}$ and
$B(Y)Z=Z_{i_0}\rho(Y)^{\mathsf T}$ with
$\rho(Y)=-\tau\,Y_{a_0i_1}Y_{a_0i_0}^{-1}\epsilon_{i_2}$. As $i_0$ and $i_1$
differ from $i_2$, $A(Y)^2=B(Y)^2=0$ and $A(Y)B(Y)=B(Y)A(Y)=0$. Let $Y$ be such a
point outside $U_1$. As $Y_{i_0}\neq0$, $Y_{i_1}=sY_{i_0}$ for some $s\in k$, so
$\rho(Y)=-\tau s\,\epsilon_{i_2}$ and
$s_Y(Y)=Y-\tau sY_{i_0}\epsilon_{i_2}^{\mathsf T}+\tau Y_{i_1}\epsilon_{i_2}^{\mathsf T}=Y$.
The column law $(\{i_0,i_1\},i_1,i_2)$ for $\tau\epsilon_{i_2}$ gives
$l(\phi\circ(1+A(Y)))=\psi(\tau)\,l(\phi)$, and the column law
$(\{i_0\},i_0,i_1)$ for $\rho(Y)$, whose entry at $i_1$ is $0$, gives
$l(\phi\circ s_Y)=l(\phi\circ(1+A(Y)))$. By (e), $l$ vanishes on $U_1$; so (d),
with $U=U_1$, and (b) give the claim.

\emph{(g) Part (2).} Let $Y_0\in M$, and put $A(Y)=0$ and
$B(Y)Z=\tau\,Z_{i_0}\epsilon_{i_1}^{\mathsf T}$. Then $B(Y)^2=0$ as
$i_0\neq i_1$. If $Y$ lies outside $U_2$, then $Y_{i_0}=0$, so $s_Y(Y)=Y$, and the
column law $(\{i_0\},i_0,i_1)$ gives $l(\phi\circ s_Y)=\psi(\tau)\,l(\phi)$. By
(f), (d) with $U=U_2$, and (b), $l$ vanishes on $M$.
\end{proof}
% Lean: (a) is small (with decomp, cosI, near), (b) kills_of_local, (c) kills_interior, (d) dead
% (ZariskiLocallyAffineLine.lean, 207 lines); (e) is kills_indep (exists_rows, Dt, lam, N, rho,
% with LeftLaw for 1 + N and ColLaw for rho, through sb_sq) and locSupp; (f) is kills_ne; (g)
% is locThree (LocalFieldT2Space.lean, 260 lines). "Vanishes on U" is Kills l U.

\begin{lemma}\label{asm-theta-agree-end-cusp-keys-two}
Let $k$, $\mathcal O$, $n$, $d$, $M$, $\psi$, $\xi$ and $l$ be as in
Lemma~\ref{asm-theta-agree-end-cusp-keys-three}, and let $i_0$ and $i_1$ be
distinct columns such that $l$ has the column law $(\{i_0\},i_0,i_1)$. Let
$\varpi$ be a uniformizer of $k$, $q$ the order of its residue field, and $c$ an
integer such that $\psi(x)=1$ for every $x$ with $|x|\leq|\varpi|^c$ and
$\psi(x)\neq1$ for some $x$ with $|x|\leq|\varpi|^{c-1}$. Suppose that $\xi$ is
trivial on $\mathcal O^\times$, and put $r=\xi(\varpi)q^{n-d}$. For $m\geq0$ let
$t_m\in M_d(k)$ be the diagonal matrix with the entries $\varpi^m$ at $i_0$,
$\varpi^{-m}$ at $i_1$ and $1$ elsewhere. Then there is $x\in\mathbf C$ such that
for every integer $m\geq0$
\[
l\bigl(Y\mapsto\mathbf 1_{\Lambda_0}(Yt_m)\bigr)=x\,(qr)^{-m}\sum_{s=0}^{2m-c}r^s,
\]
where $\mathbf 1_{\Lambda_0}$ is the indicator function of
$\Lambda_0=M_{n\times d}(\mathcal O)$ and the sum is $0$ when $2m-c<0$.
\end{lemma}
\leanref{Development/WeilBound/ThetaLifts/EigenvalueEndgame/CuspidalAverages/KeyThree.lean}{PositivityTheory.ThKEnd.Cusp.LocTwo}
\leanref{Development/WeilBound/ThetaLifts/EigenvalueEndgame/CuspidalAverages/ColOpSingle.lean}{PositivityTheory.ThKEnd.Cusp.LocB.stepA}
\leanref{Development/WeilBound/ThetaLifts/EigenvalueEndgame/CuspidalAverages/ColOpSingle.lean}{PositivityTheory.ThKEnd.Cusp.LocB.stepB}
\leanref{Development/WeilBound/ThetaLifts/EigenvalueEndgame/CuspidalAverages/ColOpSingle.lean}{PositivityTheory.ThKEnd.Cusp.LocB.bx_zero}
\leanref{Development/WeilBound/ThetaLifts/EigenvalueEndgame/CuspidalAverages/ColOpSingle.lean}{PositivityTheory.ThKEnd.Cusp.LocB.solve}
\begin{proof}
This is a computation with indicator functions, and its result is stated. For
integers $p$ and $\mu$ let $\Lambda(p,\mu)$ be the set of the $Y\in M$ with
$|Y_{ai_0}|\leq|\varpi|^p$ and $|Y_{ai_1}|\leq|\varpi|^\mu$ for every row $a$ and
all other entries in $\mathcal O$, and put $\beta(p,\mu)=l(\mathbf 1_{\Lambda(p,\mu)})$.
As $Yt_m\in\Lambda_0$ if and only if $Y\in\Lambda(-m,m)$, the left side of the
claim is $\beta(-m,m)$. The computation gives three relations.
\begin{enumerate}
\item[(A)] If $c\leq\mu-1-p$, then
  $\beta(p,\mu)-\beta(p+1,\mu)=r\bigl(\beta(p+1,\mu)-\beta(p+2,\mu)\bigr)$.
\item[(B)] If $c\leq\mu-p$, then
  $\beta(p,\mu)-\beta(p+1,\mu)=q\bigl(\beta(p,\mu+1)-\beta(p+1,\mu+1)\bigr)$.
\item[(C)] If $\mu-p<c$, then $\beta(p,\mu)=0$.
\end{enumerate}
For (C), choose $\tau$ with $|\tau|\leq|\varpi|^{c-1}$ and $\psi(\tau)\neq1$.
As $|\tau Y_{ai_0}|\leq|\varpi|^{c-1+p}\leq|\varpi|^\mu$ for
$Y\in\Lambda(p,\mu)$, this set is stable under adding $\pm\tau Y_{i_0}$ to
column $i_1$, and the column law gives $\beta(p,\mu)=\psi(\tau)\beta(p,\mu)$. For
(A) and (B), $\beta(p,\mu)-\beta(p+1,\mu)$ is a combination, with coefficients
independent of $p$ and $\mu$, of the values $l_a(p,\mu)$ of $l$ at the indicator
functions of the sets of the $Y\in\Lambda(p,\mu)$ with $|Y_{ai_0}|=|\varpi|^p$
and $|Y_{a'i_0}|\leq|\varpi|^{p+1}$ for the rows $a'\neq a$. The law of $l$ for
the diagonal matrix with $\varpi$ in row $a$ and for the transvections between
rows, and the column law for adding multiples of column $i_0$ to the other
columns, give $l_a(p,\mu)=\xi(\varpi)q^{n-d}\,l_a(p+1,\mu)$ when
$c\leq\mu-1-p$, which gives (A). The column law for adding multiples of column
$i_0$ to column $i_1$, together with
Lemma~\ref{asm-theta-agree-end-cusp-keys-three}(1), gives
$l_a(p,\mu)=q\,l_a(p,\mu+1)$ when $c\leq\mu-p$, which gives (B).

These relations determine $\beta(-m,m)$. By (C) and (A), induction on $j\geq0$
gives $\beta(\mu-c-j,\mu)-\beta(\mu-c-j+1,\mu)=r^j\beta(\mu-c,\mu)$, and so
$\beta(\mu-c-j+1,\mu)=\beta(\mu-c,\mu)\sum_{s<j}r^s$. By (B) for $p=m-c$ and
$\mu=m$, by (C), and by the first of these identities for $\mu=m+1$ and $j=1$,
$\beta(m-c,m)=qr\,\beta(m+1-c,m+1)$; as $qr\neq0$, induction gives
$\beta(m-c,m)=(qr)^{-m}\beta(-c,0)$. For $j=2m-c+1\geq0$ this gives the claim
with $x=\beta(-c,0)$; when $2m-c<0$, both sides are $0$ by (C).
\end{proof}
% Lean: LocTwo is proved as locTwo_of locSupp (ColOpSingle.lean and LeSuccIff.lean, about 520
% lines). beta(p, mu) is bx with T = univ, l_a(p, mu) is val (indicators of Piv), the
% decomposition is bx_weights and bx_empty, (A) is stepA (scale, what2, cols), (B) is stepB
% (cols, trace, LocSupp), (C) is bx_zero, and the recursion is solved by solve. The moves of
% the exponents are LeSuccIff.move and count. qr is nonzero, as q > 0 and xi takes values in
% C^x.

\begin{lemma}\label{asm-theta-agree-end-cusp-keys-vacuum}
Let $B$ be a theta base for a CM field $E$ (Definition~\ref{asm-theta-agree-models}(1)),
let $X$ be a Hermitian and $W$ a skew-Hermitian lattice space, of ranks $n$ and
$d$ and of some signatures, and let $M$ be a regular model for $(X,W)$. Then for
almost all finite places $v$ of $F$, every $F$-embedding $t\colon E\to F_v$ and
every linear model $\iota$ of the second kind for $(t,1,1)$
(Definition~\ref{asm-theta-agree-models}(5)), $\iota(\mathrm{vac}_v)$ is a
multiple of the indicator function of $M_{n\times d}(\mathcal O_v)$.
\end{lemma}
\leanref{Development/WeilBound/ThetaLifts/EigenvalueEndgame/CuspidalAverages/KeyThree.lean}{PositivityTheory.ThKEnd.Cusp.VacBar}
\leanref{Development/WeilBound/ThetaLifts/EigenvalueEndgame/CuspidalAverages/VacIntZ.lean}{PositivityTheory.ThKEnd.Cusp.vacBar}
\begin{proof}
Put $I_0=\{1,\ldots,n\}\times\{1,\ldots,d\}$ and let $e$ be the isomorphism of
$M$. As in the proof of Lemma~\ref{asm-theta-agree-coinvariants-dictionary},
Situation~\ref{asm-theta-agree-dictionary-situation} holds for $J=J_X$,
$J'=J_W$, $\varepsilon=1$, $I=I_0$ and $e$, and by
Lemma~\ref{asm-theta-agree-dictionary-charts}(3) every linear model of the
second kind for $(t,1,1)$ is a Heisenberg chart for $e$ at $v$
(Definition~\ref{asm-theta-agree-dictionary-chart}). Let $\delta\in E$ with
$\delta\neq c(\delta)$, and let $E_{ab}\in M_{n\times d}(E)$ be the matrix units.
For almost all $v$ the following hold, by
Lemma~\ref{asm-theta-agree-end-cusp-keys-places}(2) and as $M$ is regular: the
$x\in F_v$ with $\psi_v(x\mathcal O_v)=1$ form $\mathcal O_v$, as they form the
completion at $v$ of the inverse different of $F$ (proof of
Lemma~\ref{asm-theta-agree-dictionary-models}(3)); $2$ is a unit in
$\mathcal O_v$; $\mathrm{vac}_v$ is the indicator function of
$\mathcal O_v^{I_0}$; for every $t$ the elements $t(\delta)$ and
$t((\delta-c(\delta))^{-1})$ and the entries of $t(J_X^{-1})$ and $t(J_W^{-1})$
lie in $\mathcal O_v$; and the components at $v$ of $e(E_{ab})$ and of
$e(\delta E_{ab})$ have entries in $\mathcal O_v$. Fix such $v$, an
$F$-embedding $t$ and a linear model $\iota$, and for $s\in F_v$ let
$\bar s\in\mathbf A_F$ have component $s$ at $v$ and $0$ elsewhere.

(i) For $P,Q\in M_{n\times d}(\mathcal O_v)$ there is $u\in M_{n\times d}(\mathbf A_E)$
at $v$ with $\pi_t(u)=P$, $\pi_{t\circ c}(u)=Q$ and $e(u)_v$ integral. Indeed,
the matrices $B_1=(P-Q)\,t((\delta-c(\delta))^{-1})$ and $A_1=P-t(\delta)B_1$
are integral; put
$u=\sum_{a,b}\bigl(\overline{(A_1)_{ab}}\,E_{ab}+\overline{(B_1)_{ab}}\,\delta E_{ab}\bigr)$.
Then $\pi_t(u)=A_1+t(\delta)B_1=P$,
$\pi_{t\circ c}(u)=A_1+t(c(\delta))B_1=P-(P-Q)=Q$, and $e(u)_v$ is integral as
$e$ is $\mathbf A$-linear.

(ii) Let $y=(y_1,y_2)$ with $y_1,y_2\in M_{n\times d}(\mathcal O_v)$, and take
$u$ by (i) for $P=-y_2$ and
$Q=-t(J_X)^{-1}y_1\bigl(t(J_W)^{\mathsf T}\bigr)^{-1}$. Then
$(\operatorname{co}(u),\operatorname{co}'(u))=y$. Let $(a_1,b_1)=e(u)_v$, which
is integral. For integral $x$, $x\cdot b_1+\tfrac12a_1\cdot b_1\in\mathcal O_v$,
as $2$ is a unit, so $\psi_v$ of it is $1$, and $x+a_1$ is integral; for $x$ not
integral, neither is $x+a_1$. So $\rho_v(e(u)_v)$ fixes the indicator function
of $\mathcal O_v^{I_0}$, which is $\mathrm{vac}_v$. As
$\iota$ is a Heisenberg chart,
$\rho_v(y)\iota(\mathrm{vac}_v)=\iota(\rho_v(e(u)_v)\mathrm{vac}_v)=\iota(\mathrm{vac}_v)$.

(iii) Let $\Phi=\iota(\mathrm{vac}_v)$. For $Y\in M_{n\times d}(\mathcal O_v)$,
(ii) for $y=(Y,0)$ gives $\Phi(Y)=(\rho_v(Y,0)\Phi)(0)=\Phi(0)$. If $Y$ has an
entry $Y_{ab}\notin\mathcal O_v$, choose $s\in\mathcal O_v$ with
$\psi_v(Y_{ab}s)\neq1$, and let $y_2$ have the entry $s$ at $(a,b)$ and $0$
elsewhere; then (ii) gives $\Phi(Y)=(\rho_v(0,y_2)\Phi)(Y)=\psi_v(Y_{ab}s)\Phi(Y)$,
so $\Phi(Y)=0$. Hence $\Phi$ is $\Phi(0)$ times the indicator function of
$M_{n\times d}(\mathcal O_v)$.
\end{proof}
% Lean: vacBar (VacIntZ.lean, 204 lines): the Cartan involutions give dX and dW, delta comes
% from the CM structure, the conditions on v are hd (psiUnr_of_not_le, through the different),
% vac_int_F, hM.std, vac_int_E and vac_locW; lm_of_linBar makes iota a Heisenberg chart (LM),
% (i) is key with vac_sum, (ii) is hfix with coord, vac_schr and LM.Heisenberg2, and (iii) is
% vac_eq and vac_one.

\begin{lemma}\label{asm-theta-agree-end-cusp-keys-torus}
Let $E$ be a CM field with a theta base $B$
(Definition~\ref{asm-theta-agree-models}(1)), let $W$ be a skew-Hermitian
lattice space of rank $d$ and some signature, with the notation of
Definition~\ref{asm-theta-agree-end-cusp-notions}, and let $(u,u')$ be a frame
of rank $2$ in $W$ with columns $u_1,u_2$ of $u$. Let
$\kappa\colon\{1,2\}\to\{1,\ldots,d\}$ be injective and $C\in GL_d(E)$ such that
column $\kappa(i)$ of $C$ is $u_i$ and $(J_W\,c(u'))^{\mathsf T}C$ has the entry
$1$ at $(i,\kappa(i))$ and $0$ elsewhere. Let $v$ be a finite place of $F$,
$t\colon E\to F_v$ an $F$-embedding and $x\in F_v^\times$, and let
$t_x\in M_d(F_v)$ be the diagonal matrix with the entries $x$ at $\kappa(1)$,
$x^{-1}$ at $\kappa(2)$ and $1$ elsewhere. Then there is $b\in GL_2(\mathbf A_E)$
with $\det b=1$, with $b=1$ if $x=1$, with $\lambda_u(b)\in U_W^{(v)}$
(Definition~\ref{asm-theta-agree-models}(2)) and with
$\pi_t(\lambda_u(b))\,t(C)=t(C)\,t_x$.
\end{lemma}
\leanref{Development/WeilBound/ThetaLifts/EigenvalueEndgame/CuspidalAverages/ActsOfReads.lean}{PositivityTheory.ThKEnd.Cusp.Key.exists_torus}
\begin{proof}
By Lemma~\ref{asm-theta-agree-dictionary-split}(2) there is $b\in GL_2(\mathbf A_E)$
with $b-1$ at $v$, $\pi_t(b)=\operatorname{diag}(x,x^{-1})$ and
$\pi_{t\circ c}(b)=1$; if $x=1$, then $b=1$ by part (1) of that lemma. Then
$b^{-1}-1=-b^{-1}(b-1)$ is at $v$, and $\pi_{t\circ c}(b^{-1})=1$. The element
$\det b-1$ is at $v$, as $b-1$ is, and its images under $\pi_t$ and
$\pi_{t\circ c}$ are $0$; so $\det b=1$. Let $e=(J_W\,c(u'))^{\mathsf T}$ and
$e'=(J_W\,c(u))^{\mathsf T}$ be as in the proof of
Lemma~\ref{asm-theta-agree-end-cusp-frames}. Then
$\lambda_u(b)-1=u(b-1)e-u'\bigl(c(b^{-1})^{\mathsf T}-1\bigr)e'$ is at $v$, as
$c(b^{-1})-1$ is (Lemma~\ref{asm-theta-agree-dictionary-split}(1)). Finally
$\pi_t(c(b^{-1}))=\pi_{t\circ c}(b^{-1})=1$, so
$\pi_t(\lambda_u(b))=1+t(u)\bigl(\operatorname{diag}(x,x^{-1})-1\bigr)t(e)$. Let
$P\in M_{d\times2}(F_v)$ have the entry $1$ at $(\kappa(i),i)$ and $0$
elsewhere. Then $t(e)\,t(C)=P^{\mathsf T}$ and $t(u)=t(C)P$, so
$\pi_t(\lambda_u(b))\,t(C)=t(C)\bigl(1+P(\operatorname{diag}(x,x^{-1})-1)P^{\mathsf T}\bigr)=t(C)\,t_x$.
\end{proof}
% Lean: exists_torus (ActsOfReads.lean): b from IsSec.exists_gl, b_0 = 1 by IsSec.extM, the
% determinant by IsSec.ext, the support at v from IsLev and levM with IsSec.atConj, and the
% last identity from torus_eq with hp and hu.

\begin{lemma}\label{asm-theta-agree-end-cusp-keys}
Let $(B,G,\mathcal X,\omega,\omega')$ be a triple for a theta base $B$ of a CM
field $E$ (Definition~\ref{asm-theta-agree-triple}), with set of differences
$\Sigma$ (Definition~\ref{asm-theta-agree-end-notions}(5)), use the notation of
Definition~\ref{asm-theta-agree-end-cusp-notions} for its second space $W$, and
let $\chi$ be a nondegenerate character of $\mathbf A_E$.
\begin{enumerate}
\item Let $(u,u')$ be a frame of rank $j\geq3$ in $W$, and $\Lambda$ a Whittaker
  functional for $(u,\chi)$. Then $\Lambda(f(\cdot\,h))=0$ for every
  $f\in\Sigma$ and $h\in U_W$.
\item Let $(u,u')$ be a frame of rank $2$ in $W$, and $\Lambda$ a Whittaker
  functional for $(u,\chi)$. For every $f\in\Sigma$ and $h_0\in U_W$ there are
  a neighbourhood $V$ of $h_0$, and for integers $i\geq0$ elements
  $b_i\in GL_2(\mathbf A_E)$ and $y_i\in\mathbf C$, with $b_0=1$,
  $\det b_i=1$ and $|y_i|\to\infty$, such that every $h\in V$ has some
  $z\in\mathbf C$ with $\Lambda(f(\cdot\,\lambda_u(b_i)h))=zy_i$ for every $i$.
\end{enumerate}
\end{lemma}
\leanref{Development/WeilBound/ThetaLifts/EigenvalueEndgame/CuspidalAverages/KeyThree.lean}{PositivityTheory.ThKEnd.Cusp.KeyThree}
\leanref{Development/WeilBound/ThetaLifts/EigenvalueEndgame/CuspidalAverages/KeyThree.lean}{PositivityTheory.ThKEnd.Cusp.KeyTwo}
\leanref{Development/WeilBound/ThetaLifts/EigenvalueEndgame/CuspidalAverages/ActsOfReads.lean}{PositivityTheory.ThKEnd.Cusp.Key.keyThree_of_locThree}
\leanref{Development/WeilBound/ThetaLifts/EigenvalueEndgame/CuspidalAverages/ActsOfReads.lean}{PositivityTheory.ThKEnd.Cusp.Key.keyTwo_of_locTwo_vacBar}
\leanref{Development/WeilBound/ThetaLifts/EigenvalueEndgame/CuspidalAverages/ActsOfReads.lean}{PositivityTheory.ThKEnd.Cusp.Key.three}
\leanref{Development/WeilBound/ThetaLifts/EigenvalueEndgame/CuspidalAverages/ActsOfReads.lean}{PositivityTheory.ThKEnd.Cusp.Key.two}
\leanref{Development/WeilBound/ThetaLifts/EigenvalueEndgame/CuspidalAverages/ExistsCond.lean}{PositivityTheory.ThKEnd.Cusp.Key.exists_cond}
\leanref{Development/WeilBound/ThetaLifts/EigenvalueEndgame/CuspidalAverages/ExistsCond.lean}{PositivityTheory.ThKEnd.Cusp.Key.yy_inv}
\leanref{Development/WeilBound/ThetaLifts/EigenvalueEndgame/CuspidalAverages/ExistsCond.lean}{PositivityTheory.ThKEnd.Cusp.Key.yy_tendsto}
\begin{proof}
Let $G=(p,q,V,V',W,M,M')$, of ranks $(a,a')$, so that $W$ has rank $d=a+a'$,
put $I_0=\{1,\ldots,a\}\times\{1,\ldots,d\}$, and let $\mathrm{vac}'_v$ be the
vacua of $M'$. Let $f=I(\Phi)-I'(\Phi')\in\Sigma$, with
$\Phi\in\operatorname{PD}(M)$ and
$\Phi'\in\operatorname{PD}(M')\cap\mathcal S^\circ(M')$.

\emph{Local families.} By Lemma~\ref{asm-theta-agree-decay}, every
Schwartz-Bruhat function converges for $(M,\omega)$, and every element of
$\mathcal S^\circ(M')$ converges for $(M',\omega')$. By
Lemma~\ref{asm-theta-agree-end-automorphic} and its proof, $I(\Psi)$ is
automorphic for $\Psi\in\operatorname{PD}(M)$ and $I'(\Psi')$ is automorphic for
$\Psi'\in\operatorname{PD}(M')\cap\mathcal S^\circ(M')$, and these spaces
consist of Schwartz-Bruhat functions. The map $\Psi\mapsto\Psi[v\mapsto\phi]$
of Lemma~\ref{asm-theta-agree-end-cusp-keys-local} is linear. For a spanning
element $\Psi=g_\infty\otimes\bigotimes_w\phi_w$ of $\operatorname{PD}(M)$,
$\phi_v=\mathrm{vac}_v$ is the indicator function of $\mathcal O_v^{I_0}$ for
almost all $v$, as $M$ is regular. At such $v$, $\phi_v(0)=1$, so
$\Psi[v\mapsto\phi]$ is $\Psi$ with $\phi_v$ replaced by $\phi$, which lies in
$\operatorname{PD}(M)$ when $\phi$ is Schwartz-Bruhat, and
$\Psi[v\mapsto\mathrm{vac}_v]=\Psi$. So $\Phi$, a finite sum of such elements,
satisfies the hypothesis of Lemma~\ref{asm-theta-agree-end-cusp-keys-local} for
$\mathcal C=\operatorname{PD}(M)$. In the same way, using also $v\neq v_1$ for
$\mathcal S^\circ(M')$, $\Phi'$ satisfies it for
$\mathcal C=\operatorname{PD}(M')\cap\mathcal S^\circ(M')$. So that lemma, for
the Weil pairs $(V,W,M,\beta,\beta/\alpha,\omega)$ and
$(V',W,M',\alpha,\alpha/\beta,\omega')$ of
Definition~\ref{asm-theta-agree-triple}(4), gives a finite set $Z_0$ of finite
places such that for every finite place $v\notin Z_0$ and every $F$-embedding
$t\colon E\to F_v$ there are linear models $\iota$ and $\iota'$ of the second
kind for $(t,1,1)$, with $\iota(\mathrm{vac}_v)$ and
$\iota'(\mathrm{vac}'_v)$ Schwartz-Bruhat, and local families $S$ for $(t,a)$
with laws $(\xi,\beta_t|\cdot|^{a/2})$ and $S'$ for $(t,a')$ with laws
$(\xi',\alpha_t|\cdot|^{a'/2})$, where
$\xi=\bigl((\beta/\alpha)_t|\cdot|^{d/2}\bigr)^{-1}$ and
$\xi'=\bigl((\alpha/\beta)_t|\cdot|^{d/2}\bigr)^{-1}$, such that
$S(\iota(\mathrm{vac}_v))=I(\Phi)$ and $S'(\iota'(\mathrm{vac}'_v))=I'(\Phi')$.

\emph{The character $\psi_t$.} For such $v$ and $t$ let $\psi_t$ be as in
Lemma~\ref{asm-theta-agree-end-cusp-keys-form}. It is not $1$: as $\chi$ is
nondegenerate, there is $z\in\mathbf A_E$ at $v$ with $\pi_{t\circ c}(z)=0$ and
$\chi(z)\neq1$, and $z=\widetilde{\pi_t(z)}$ by
Lemma~\ref{asm-theta-agree-dictionary-split}(1), so $\psi_t(\pi_t(z))\neq1$. It
is continuous, as $\chi$ is and $\tilde s=\bar s\,\tilde1$, where
$\bar s\in\mathbf A_F$ has component $s$ at $v$ and $0$ elsewhere.

\emph{Additivity.} As $I(\Phi)$ and $I'(\Phi')$ are automorphic, so are their
right translates, and the additivity of $\Lambda$ gives
$\Lambda(f(\cdot\,h))=\Lambda(I(\Phi)(\cdot\,h))-\Lambda(I'(\Phi')(\cdot\,h))$
for every $h\in U_W$.

\emph{A completion of $u$.} By
Lemma~\ref{asm-theta-agree-end-cusp-keys-complement} for
$p=(J_W\,c(u'))^{\mathsf T}$, which has $pu=(u^{\mathsf T}J_W\,c(u'))^{\mathsf T}=1$,
there are an injective map $\kappa\colon\{1,\ldots,j\}\to\{1,\ldots,d\}$ and
$C_E\in GL_d(E)$ whose column $\kappa(i)$ is the column $u_i$ of $u$ for every
$i$, and such that $pC_E$ has the entry $1$ at $(i,\kappa(i))$ and $0$
elsewhere.

(1) Let $h\in U_W$. By Lemma~\ref{asm-theta-agree-end-cusp-keys-places}(3)
there are a finite place $v\notin Z_0$ and an $F$-embedding $t\colon E\to F_v$.
By the additivity, it suffices to show that
$\Lambda(S_0(\phi)(\cdot\,g))=0$ for every local family $S_0$ for $(t,n)$, with
some laws $(\xi_0,\alpha_0)$, every Schwartz-Bruhat $\phi$ and every
$g\in U_W$: apply it to $S$ and $\iota(\mathrm{vac}_v)$, and to $S'$ and
$\iota'(\mathrm{vac}'_v)$. Lemma~\ref{asm-theta-agree-end-cusp-keys-form}, for
$S_0$, $\kappa$, $C=t(C_E)$ and $g$, gives a linear form $l$, and its part (3)
for $k=1$ gives $\Lambda(S_0(\phi)(\cdot\,g))=\alpha_0(\det\pi_t(g))\,l(\phi_g)$
with $\phi_g(Y)=\phi(YC^{-1}\pi_t(g))$, a Schwartz-Bruhat function. By parts (1)
and (2) of that lemma, for $i=1$ with $T=\{\kappa(1)\}$ and for $i=2$ with
$T=\{\kappa(1),\kappa(2)\}$, $l$ satisfies the hypotheses of
Lemma~\ref{asm-theta-agree-end-cusp-keys-three}(2) for $\psi_t$, $\xi_0$ and the
distinct columns $\kappa(1)$, $\kappa(2)$ and $\kappa(3)$, which exist as
$j\geq3$. So $l(\phi_g)=0$.
% Lean: keyThree_of_locThree (ActsOfReads.lean, 710 lines with exists_loc, exists_l, three,
% exists_torus, two and keyTwo_of_locTwo_vacBar): v and t from Places (the infinitely many
% split places, frequently_cofinite_iff_infinite) and the eventual statements locV and locV'
% (exists_loc with polData_le_placewise, so_positions and polData_inf_so_le_placewise; the
% convergence from KDecay and the automorphy from IAut); the additivity is wh_sub; three
% applies frame_compl, exists_l at k = 1 and LocThree with the columns e 0, e 1, e 2 and the
% character chi.compAddMonoidHom (sA v sec), which is not 1 by psi_ne.

(2) Let $j=2$ and $h_0\in U_W$. By
Lemma~\ref{asm-theta-agree-end-cusp-keys-places}(2) and~(3),
Lemma~\ref{asm-theta-agree-end-cusp-keys-vacuum} for $M$ and for $M'$, and as
the set $T$ of the character datum is finite, there are a finite place
$v\notin Z_0\cup T$ and an $F$-embedding $t\colon E\to F_v$ such that the
conclusion of Lemma~\ref{asm-theta-agree-end-cusp-keys-vacuum} holds at $v$ for
$M$ and for $M'$, and the entries of $\pi_t(h_0)$, $\pi_t(h_0^{-1})$, $t(C_E)$
and $t(C_E^{-1})$ lie in $\mathcal O_v$. Put $C=t(C_E)$,
$\Lambda_0=M_{a\times d}(\mathcal O_v)$ and
$\Lambda'_0=M_{a'\times d}(\mathcal O_v)$, let $\varpi$ be a uniformizer of
$F_v$ and $q$ the order of its residue field, and let $t_i$, for $i\geq0$, be
the diagonal matrix with the entries $\varpi^i$ at $\kappa(1)$, $\varpi^{-i}$ at
$\kappa(2)$ and $1$ elsewhere. By
Lemma~\ref{asm-theta-agree-end-cusp-keys-vacuum},
$\iota(\mathrm{vac}_v)=z_0\mathbf 1_{\Lambda_0}$ and
$\iota'(\mathrm{vac}'_v)=z'_0\mathbf 1_{\Lambda'_0}$ for some
$z_0,z'_0\in\mathbf C$.

\emph{The conductor.} As $\psi_t$ is continuous, the image of $\varpi^c\mathcal O_v$
under it lies in $\{w\in\mathbf C:|w-1|<1/2\}$ for $c$ large, and the only
subgroup of $\mathbf C^\times$ in this set is $\{1\}$; so $\psi_t$ is trivial on
$\varpi^c\mathcal O_v$ for $c$ large. As $\psi_t\neq1$, it is not trivial on
$\varpi^c\mathcal O_v$ for $c$ small. Let $c_0$ be the least integer such that
$\psi_t$ is trivial on $\varpi^{c_0}\mathcal O_v$; then $\psi_t(x)\neq1$ for some
$x$ with $|x|\leq|\varpi|^{c_0-1}$.

\emph{The elements $b_i$.} For $i\geq0$, let $b_i$ be given by
Lemma~\ref{asm-theta-agree-end-cusp-keys-torus} for $\kappa$, $C_E$ and
$x=\varpi^i$. Then $b_0=1$, $\det b_i=1$, $\lambda_u(b_i)\in U_W^{(v)}$ and
$\pi_t(\lambda_u(b_i))=Ct_iC^{-1}$, which has determinant $1$.

\emph{The sequence $y_i$.} As $a+a'$ is even and $a'<a$, $\ell=(a-a')/2$ is an
integer, and $\ell\geq2$ as $2(p+q)+1<\lfloor(a-a')/2\rfloor$. Put
$r=\alpha_t(\varpi)\beta_t(\varpi)^{-1}q^\ell$ and
$y_i=(qr)^{-i}\sum_{s=0}^{2i-c_0}r^s$. As $\alpha_t$ and $\beta_t$ are unitary
(Lemma~\ref{asm-theta-agree-dictionary-split}(3)), $|r|=q^\ell>q>1$. For
$i\geq c_0$, put $N=2i-c_0+1$; then $|r^N-1|\leq(|r|+1)\bigl|\sum_{s<N}r^s\bigr|$,
so
$|y_i|\geq(q|r|)^{-i}(|r|^N-1)/(|r|+1)\geq\bigl(|r|^{1-c_0}(|r|/q)^i-1\bigr)/(|r|+1)$,
which tends to infinity.

\emph{The laws at $v$.} As $v\notin T$, $\alpha$ and $\beta$ are unramified at
the place of $E$ given by $t$, so $\alpha_t$ and $\beta_t$ are trivial on
$\mathcal O_v^\times$, and so are $\xi$ and $\xi'$. As $|\varpi|=q^{-1}$,
$\xi(\varpi)=\alpha_t(\varpi)\beta_t(\varpi)^{-1}q^{d/2}$, so
$\xi(\varpi)q^{a-d}=r$; and in the same way $\xi'(\varpi)q^{a'-d}=r^{-1}$.

\emph{The neighbourhood.} Let $V$ be the set of the $h\in U_W$ such that
$\pi_t(h)$ and $\pi_t(h^{-1})$ have their entries in $\mathcal O_v$. It contains
$h_0$, and it is open by Lemma~\ref{asm-theta-agree-end-cusp-keys-places}(1), as
$\mathcal O_v$ is open in $F_v$.

\emph{The values.} Let $h\in V$. Lemma~\ref{asm-theta-agree-end-cusp-keys-form},
for $S$, $\kappa$, $C$ and $g=h$, gives a linear form $l$, which by its parts
(1) and (2), for $i=1$ and $T=\{\kappa(1)\}$, satisfies the hypotheses of
Lemma~\ref{asm-theta-agree-end-cusp-keys-two} for $\psi_t$, $\xi$, the columns
$\kappa(1)$ and $\kappa(2)$, $\varpi$ and $c_0$, with the number
$\xi(\varpi)q^{a-d}=r$. So there is $x\in\mathbf C$ with
$l\bigl(Y\mapsto\mathbf 1_{\Lambda_0}(Yt_i)\bigr)=x\,y_i$ for every $i\geq0$. Let
$i\geq0$ and $k=\lambda_u(b_i)$. Then
$C^{-1}\pi_t(k)\pi_t(h)=t_iC^{-1}\pi_t(h)$, and $C^{-1}\pi_t(h)$ and its inverse
have their entries in $\mathcal O_v$, so
$\mathbf 1_{\Lambda_0}(YC^{-1}\pi_t(k)\pi_t(h))=\mathbf 1_{\Lambda_0}(Yt_i)$.
Part (3) of Lemma~\ref{asm-theta-agree-end-cusp-keys-form} gives
\[
\Lambda\bigl(I(\Phi)(\cdot\,kh)\bigr)=\Lambda\bigl(S(z_0\mathbf 1_{\Lambda_0})(\cdot\,kh)\bigr)
=z_0\,\beta_t(\det\pi_t(h))\,|\det\pi_t(h)|^{a/2}\,x\,y_i.
\]
In the same way, with $S'$, $\xi'$ and the number $r^{-1}$,
$\Lambda(I'(\Phi')(\cdot\,kh))=z_1(qr^{-1})^{-i}\sum_{s=0}^{2i-c_0}r^{-s}$ for
some $z_1\in\mathbf C$ independent of $i$, and this is $z_1r^{c_0}y_i$, as
$\sum_{s=0}^{N}r^{-s}=r^{-N}\sum_{s=0}^{N}r^s$ for $N=2i-c_0$ and
$(qr^{-1})^{-i}r^{-N}=r^{c_0}(qr)^{-i}$. By the additivity,
$\Lambda(f(\cdot\,\lambda_u(b_i)h))=zy_i$ for every $i$, with
$z=z_0\,\beta_t(\det\pi_t(h))\,|\det\pi_t(h)|^{a/2}\,x-z_1r^{c_0}$.
\end{proof}
% Lean: keyTwo_of_locTwo_vacBar (ActsOfReads.lean): p1 to p4 from Places, frame_compl, the
% integrality im, the choice of v with locV, locV', VacBar for G.M and G.M', i1 to i4 and
% X.T.eventually_cofinite_notMem; kappa and kappa' are z_0 and z'_0, c_0 is exists_cond, b is
% exists_torus at varpi^i, ell is epsilon (Geo.even, Geo.eBig), r1 and r2 the identities for
% xi and xi', rn the bound |r| > q (with the unitarity us from section_statement), un the
% unramified laws; V is {g | IntM (cp t g) and IntM (cp t g^-1)}; two (LocTwo, exists_l,
% oneO_tor, det_torus) gives the values, yy_inv the factor r^(c_0) and yy_tendsto the growth.

\begin{lemma}\label{asm-theta-agree-end-cusp}
Let $(B,G,\mathcal X,\omega,\omega')$ be a triple for a theta base $B$ of a CM
field $E$ (Definition~\ref{asm-theta-agree-triple}), and let $H'$ be a plane in
$W$ over a lattice space $W_1$ (Definition~\ref{asm-theta-agree-end-notions}(1)).
If $f$ lies in the set of differences of the triple and $0$ is the constant term
of $f$ along $H'$, then $f$ is cuspidal
(Definition~\ref{asm-theta-agree-end-notions}(4)).
\end{lemma}
\leanref{Development/WeilBound/ThetaLifts/EigenvalueEndgame/PlaneRemovalConfiguration.lean}{PositivityTheory.ThKEnd.CuspOfLine}
\leanref{Development/WeilBound/ThetaLifts/EigenvalueEndgame/CuspidalAverages/CuspOfLineOfKeys.lean}{PositivityTheory.ThKEnd.Cusp.cuspOfLine_of_keys}
\begin{proof}
Let $W$ be the second space of the triple, of rank $d$, with the notation of
Definition~\ref{asm-theta-agree-end-cusp-notions}, and let $f$ be as in the
statement. By Lemma~\ref{asm-theta-agree-end-automorphic}, $f$ is automorphic,
and by Lemma~\ref{asm-theta-agree-end-two-step-moderate} there are $C$ and $k$
with $|f(h)|\leq C\operatorname{ht}(h)^k$ for every $h\in U_W$. Let $\psi$ be a
standard character of $\mathbf A_E$, which exists \cite[\S4.1]{Tat67}; by
Lemma~\ref{asm-theta-agree-end-cusp-coefficient}(1) it is nondegenerate.

\emph{Induction.} We show by induction on $j\geq1$ that for every frame
$(u,u')$ of rank $j$ in $W$, $0$ is the average of $f$ over
$N_{U_u}\cap U_W(F)$. Let $(u,u')$ be a frame of rank $j$, and suppose that this
holds for the ranks $r<j$. By Lemma~\ref{asm-theta-agree-end-cusp-frames}(4),
$(u_{\leq r},u'_{\leq r})$ is a frame of rank $r$, so $0$ is the average of $f$
over $N_{U_r}\cap U_W(F)$ for the span $U_r$ of $u_{\leq r}$ and $1\leq r<j$. By
Lemma~\ref{asm-theta-agree-end-cusp-coefficient}(1), $\mathcal W_{u,\psi}$ is a
Whittaker functional for $(u,\psi)$.
\begin{itemize}
\item If $j=1$, then $U_u=Ex$ for the column $x$ of $u$, with $x\neq0$, as
  $u^{\mathsf T}J_W\,c(u')=1$, and $h_W(x,x)=0$. Lemma~\ref{asm-theta-agree-end-cusp-line}
  gives the claim.
\item If $j=2$, let $x_0\in U_W$. Lemma~\ref{asm-theta-agree-end-cusp-keys}(2)
  for $\Lambda=\mathcal W_{u,\psi}$ and $h_0=x_0$ gives a neighbourhood $V$ of
  $x_0$ and elements $b_i$ and $y_i$ as in
  Lemma~\ref{asm-theta-agree-end-cusp-two}, which gives
  $\mathcal W_{u,\psi}(f(\cdot\,x_0))=0$. So
  Lemma~\ref{asm-theta-agree-end-cusp-coefficient}(3) gives the claim.
\item If $j\geq3$, Lemma~\ref{asm-theta-agree-end-cusp-keys}(1) for
  $\Lambda=\mathcal W_{u,\psi}$ gives $\mathcal W_{u,\psi}(f(\cdot\,x))=0$ for
  every $x\in U_W$, and Lemma~\ref{asm-theta-agree-end-cusp-coefficient}(3)
  gives the claim.
\end{itemize}

\emph{Conclusion.} Let $U\subset E^d$ be a nonzero $E$-subspace with
$h_W(U,U)=0$, $\nu$ a Haar measure on $N_U$, $\mathcal F$ a fundamental domain
for $N_U\cap U_W(F)$ with respect to $\nu$, and $h\in U_W$. By
Lemma~\ref{asm-theta-agree-end-cusp-frames}(4), $U$ is the span of a frame of
rank $\dim U\geq1$, so $0$ is the average of $f$ over $N_U\cap U_W(F)$. This
group is a lattice in $N_U$ (Lemma~\ref{asm-theta-agree-end-kill-unipotent}(2)),
so $0<\nu(\mathcal F)<\infty$, by the proof of
Lemma~\ref{asm-theta-agree-end-kill-averages}(1). For the Haar measure
$\nu'=\nu(\mathcal F)^{-1}\nu$, $\nu'(\mathcal F)=1$, so
$\int_{\mathcal F}f(nh)\,d\nu(n)=\nu(\mathcal F)\int_{\mathcal F}f(nh)\,d\nu'(n)=0$.
So $f$ is cuspidal.
% Lean: cuspOfLine_of_keys (CuspOfLineOfKeys.lean, 30 lines): C (the record Ctx), sig_isAut
% (IAut), psi from GlobalTheory.T.exists_standardAdelicAddChar, key by strong induction on j
% with C.ct for the constant term and low from fr_fst; j = 1 is av_line, j = 2 is wh_ne with
% KeyTwo and wh_two (with sig_mod, which is the proof of Lemma
% asm-theta-agree-end-two-step-moderate), and j + 3 is wh_ne with KeyThree; the conclusion
% is C.frame and C.int_zero (UnipGood.fin). In the closure, cuspOfLine_of_keys is applied as
% Cusp.cuspOfLine_of_keys WeilConv KDecay IAut KeyThree KeyTwo, with
% KeyThree = Cusp.Key.keyThree_of_locThree
% thetaBasic KDecay IAut LocThree Places, KeyTwo = Cusp.Key.keyTwo_of_locTwo_vacBar thetaBasic
% KDecay IAut LocTwo VacBar Places, LocThree = Cusp.locThree, LocTwo = Cusp.locTwo_of
% Cusp.locSupp, VacBar = Cusp.vacBar and Places = Cusp.places. Size: about 5,750 lines
% beyond the lemmas stated elsewhere and Lemma asm-theta-agree-end-automorphic, of which
% about 4,700 are used by no other lemma of this group.
\end{proof}

\subsection{Orthogonality of cuspidal differences}\label{subsection-asm-orthogonality}

\begin{definition}\label{asm-theta-agree-end-orth-notions}
Let $B$ be a theta base for a CM field $E$
(Definition~\ref{asm-theta-agree-models}(1)), with maps $\pi_w$, and let $W$
be a lattice space of rank $d$ (Definition~\ref{asm-theta-agree-doubling}(1)).
\begin{enumerate}
\item A function $f$ on $U_W$ has \emph{uniform moderate growth} if there are
  an integer $k\geq0$ and a set $S$ of functions on $U_W$ with $f\in S$ such
  that every $f_1\in S$ is continuous, satisfies
  $|f_1(h)|\leq C\operatorname{ht}(h)^k$ for some $C$ and all $h\in U_W$, and
  for every line $\gamma$ in $U_W$
  (Definition~\ref{asm-theta-agree-matched}(1)) there is $f_2\in S$ such that
  $t\mapsto f_1(h\gamma(t))$ has derivative $f_2(h)$ at $t=0$ for every
  $h\in U_W$.
\item A function $f$ on $U_W$ is \emph{smooth at the finite places} if there is
  a neighbourhood $U$ of $1$ in $U_W$ such that $f(hu)=f(h)$ for every
  $h\in U_W$ and every $u\in U$ with $\pi_w(u)=1$ for every infinite place $w$
  of $F$.
\item Let $(X,W,M,\chi_X,\chi_W,\omega)$ be a Weil pair for $B$
  (Definition~\ref{asm-theta-agree-triple}(1)), with $X$ of rank $n$, put
  $I=\{1,\ldots,n\}\times\{1,\ldots,d\}$, and let $v$ be a finite place of
  $F$. A \emph{local map} at $v$ for $\omega$ is a linear map $T$ from the
  functions on $F_v^I$ to the functions on $U_W$ such that $T(\phi)$ is rapidly
  decreasing (Definition~\ref{asm-theta-agree-definite-products}(1)) for every
  Schwartz-Bruhat $\phi$, and $T(\phi_1)(x)=T(\phi)(xh)$ for every $x\in U_W$
  whenever $h\in U_W^{(v)}$, the functions $\phi$ and $\phi_1$ are
  Schwartz-Bruhat and $\omega(1,h)$ acts at $v$ by $\phi\mapsto\phi_1$
  (Definition~\ref{asm-theta-agree-models}(2)). A function $f$ on $U_W$
  \emph{comes from $v$} for $\omega$ if $f=\sum_{i=1}^mT_i(\phi_i)$ for some
  local maps $T_1,\ldots,T_m$ at $v$ for $\omega$ and Schwartz-Bruhat functions
  $\phi_1,\ldots,\phi_m$ on $F_v^I$.
\end{enumerate}
\end{definition}
\leanref{Development/WeilBound/ThetaLifts/FrameEstimates/TowerRallisLemma.lean}{PositivityTheory.ThKEst.UMG}
\leanref{Development/WeilBound/ThetaLifts/FrameEstimates/TowerRallisLemma.lean}{PositivityTheory.ThKEst.SmoothFin}
\leanref{Development/WeilBound/ThetaLifts/EigenvalueEndgame/OrthogonalDecomposition/HodgeTateFromSandwich.lean}{PositivityTheory.ThKEnd.Orth.LocMap}
\leanref{Development/WeilBound/ThetaLifts/EigenvalueEndgame/OrthogonalDecomposition/HodgeTateFromSandwich.lean}{PositivityTheory.ThKEnd.Orth.Comes}

\begin{lemma}\label{asm-theta-agree-end-orth-lift}
Let $(X,W,M,\chi_X,\chi_W,\omega)$ be a Weil pair for a theta base $B$ of a CM
field $E$ (Definition~\ref{asm-theta-agree-triple}(1)), with $X$ and $W$ of
ranks $n$ and $d$, let $\mu$ be a Haar measure on $U_W$ and $\mathcal D$ a
fundamental domain for $U_W(F)$ with respect to $\mu$, and let $f$ be a rapidly
decreasing function on $U_W$
(Definition~\ref{asm-theta-agree-definite-products}(1)).
\begin{enumerate}
\item For every integer $k\geq0$ there is $K\geq0$ such that
  $|f(h)|\,T(h)\leq\max(A,0)\,K$ for every $h\in U_W$, every $A\in\mathbf R$ and
  every function $T\colon U_W\to\mathbf R$ with $T(\gamma h)=T(h)$ for
  $\gamma\in U_W(F)$ and $T(h)\leq A\operatorname{ht}(h)^k$ for all $h$.
\item $\mu(\mathcal D)<\infty$. For every Schwartz-Bruhat function $\Psi$ on
  $\mathbf A^{n\times d}$ and every $g\in U_X$, the function
  $h\mapsto f(h)\overline{\theta_\omega(\Psi)(g,h)}$ is $\mu$-integrable on
  $\mathcal D$, so the theta lift $\theta^*_\omega(f,\Psi)$
  (Definition~\ref{asm-theta-agree-definite-products}(2)) is defined. It is
  conjugate-linear in $\Psi$, and
  $\theta^*_\omega(f,\omega(u,1)\Psi)(g)=\theta^*_\omega(f,\Psi)(gu)$ for
  $u,g\in U_X$.
\item For every Schwartz-Bruhat $\Psi$, the function $\theta^*_\omega(f,\Psi)$
  is continuous on $U_X$, and
  $\theta^*_\omega(f,\Psi)(\gamma g)=\theta^*_\omega(f,\Psi)(g)$ for
  $\gamma\in U_X(F)$ and $g\in U_X$.
\end{enumerate}
\end{lemma}
\leanref{Development/WeilBound/ThetaLifts/EigenvalueEndgame/OrthogonalDecomposition/RapidBound.lean}{PositivityTheory.ThKEnd.Orth.rapid_bound}
\leanref{Development/WeilBound/ThetaLifts/EigenvalueEndgame/OrthogonalDecomposition/RapidBound.lean}{PositivityTheory.ThKEnd.Orth.finW}
\leanref{Development/WeilBound/ThetaLifts/EigenvalueEndgame/OrthogonalDecomposition/RapidBound.lean}{PositivityTheory.ThKEnd.Orth.th_bd}
\leanref{Development/WeilBound/ThetaLifts/EigenvalueEndgame/OrthogonalDecomposition/RapidBound.lean}{PositivityTheory.ThKEnd.Orth.th_int}
\leanref{Development/WeilBound/ThetaLifts/EigenvalueEndgame/OrthogonalDecomposition/QhtUnbdd.lean}{PositivityTheory.ThKEnd.Orth.th_add}
\leanref{Development/WeilBound/ThetaLifts/EigenvalueEndgame/OrthogonalDecomposition/QhtUnbdd.lean}{PositivityTheory.ThKEnd.Orth.th_smul}
\leanref{Development/WeilBound/ThetaLifts/EigenvalueEndgame/OrthogonalDecomposition/QhtUnbdd.lean}{PositivityTheory.ThKEnd.Orth.th_left}
\leanref{Development/WeilBound/ThetaLifts/EigenvalueEndgame/OrthogonalDecomposition/HodgeTateFromSandwich.lean}{PositivityTheory.ThKEnd.Orth.ThAut}
\leanref{Development/WeilBound/ThetaLifts/EigenvalueEndgame/OrthogonalDecomposition/RapidBound.lean}{PositivityTheory.ThKEnd.Orth.thAut_of}
\begin{proof}
(1) Let $C$ be such that $|f(h)|\,\overline{\operatorname{ht}}(h)^k\leq C$ for
all $h$, and put $K=2^k\max(C,0)$. Let $h\in U_W$. Heights are at least $1$, so
$\overline{\operatorname{ht}}(h)\geq1$, and there is $\gamma\in U_W(F)$ with
$\operatorname{ht}(\gamma h)\leq2\,\overline{\operatorname{ht}}(h)$. Then
$T(h)=T(\gamma h)\leq\max(A,0)\operatorname{ht}(\gamma h)^k\leq\max(A,0)\,2^k\,\overline{\operatorname{ht}}(h)^k$,
and
$|f(h)|\,T(h)\leq\max(A,0)\,2^k\,|f(h)|\,\overline{\operatorname{ht}}(h)^k\leq\max(A,0)\,K$.

(2) By Lemma~\ref{asm-theta-agree-convergence-signature}, the matrix of $W$ is
nondegenerate $\epsilon$-Hermitian, so $\mu(\mathcal D)<\infty$ by
Theorem~\ref{asm-theta-agree-definite-volume}. By
Lemma~\ref{asm-theta-agree-definite-inputs}(2) there are $C_0$ and $k$ with
$|\theta_\omega(\Psi)(g,h)|\leq C_0\bigl(\operatorname{ht}(g)\operatorname{ht}(h)\bigr)^k$,
and $\theta_\omega(\Psi)(g,\gamma h)=\theta_\omega(\Psi)(g,h)$ for
$\gamma\in U_W(F)$ by Lemma~\ref{asm-theta-agree-anisotropic-automorphy}. Let
$K$ be given by (1) for $k$. For $b\geq0$ and $g$ with
$\operatorname{ht}(g)\leq b$, (1) applied to
$T(h)=|\theta_\omega(\Psi)(g,h)|$ and $A=\max(C_0,0)\,b^k$ gives
\[
\bigl|f(h)\,\theta_\omega(\Psi)(g,h)\bigr|\leq\max(C_0,0)\,b^kK\qquad(h\in U_W).\tag{a}
\]
The function $f$ is continuous, and $\theta_\omega(\Psi)$ is continuous on
$U_X\times U_W$ by Lemma~\ref{asm-theta-agree-convergence-theta}(2). So for
each $g$ the integrand is continuous and bounded, hence integrable on
$\mathcal D$. By Lemma~\ref{asm-theta-agree-convergence-theta}(2),
$\theta_\omega(\Psi+a\Psi')=\theta_\omega(\Psi)+a\,\theta_\omega(\Psi')$, and
the integral is linear on integrable functions; this gives the
conjugate-linearity. As $\omega$ is a representation,
$\theta_\omega(\omega(u,1)\Psi)(g,h)=\theta_\omega(\Psi)(gu,h)$, which gives
the last formula.

(3) By Lemma~\ref{asm-theta-agree-convergence-signature} and
Theorem~\ref{asm-theta-agree-convergence-siegel}(1), $U_X$ is locally compact
and the height is bounded on its compact subsets. Let $g_0\in U_X$, let
$\Omega$ be a compact neighbourhood of $g_0$, and let $b$ bound the height on
$\Omega$. By (a), the integrand $f(h)\overline{\theta_\omega(\Psi)(g,h)}$ is
bounded by a constant for $g\in\Omega$ and $h\in U_W$; it is continuous in
$g$; and $\mu(\mathcal D)<\infty$. By dominated convergence,
$\theta^*_\omega(f,\Psi)$ is continuous at $g_0$. The invariance holds because
$\theta_\omega(\Psi)(\gamma g,h)=\theta_\omega(\Psi)(g,h)$ for
$\gamma\in U_X(F)$, by Lemma~\ref{asm-theta-agree-anisotropic-automorphy}.
\end{proof}

\begin{lemma}\label{asm-theta-agree-end-orth-perp}
Let $(X,W,M,\chi_X,\chi_W,\omega)$ be a Weil pair for a theta base $B$ of a CM
field $E$ (Definition~\ref{asm-theta-agree-triple}(1)), with $X$ and $W$ of
ranks $n$ and $d$ and with measure $\mu_X$ and domain $\mathcal D_X$ of $M$,
and let $f$ be a rapidly decreasing function on $U_W$
(Definition~\ref{asm-theta-agree-definite-products}(1)). Let $\Psi$ be a
Schwartz-Bruhat function on $\mathbf A^{n\times d}$ that converges
(Definition~\ref{asm-theta-agree-models}(7)), and suppose that there are $C$
and $k$ with
$\int_{\mathcal D_X}|\theta_\omega(\Psi)(g,h)|\,d\mu_X(g)\leq C\operatorname{ht}(h)^k$
for every $h\in U_W$. If $\theta^*_\omega(f,\Psi)=0$
(Definition~\ref{asm-theta-agree-definite-products}(2)) for every Haar measure
$\mu$ on $U_W$ and every fundamental domain $\mathcal D$ for $U_W(F)$ with
respect to $\mu$, then $I(\Psi)$ is orthogonal to $f$
(Definition~\ref{asm-theta-agree-end-notions}(4)).
\end{lemma}
\leanref{Development/WeilBound/ThetaLifts/EigenvalueEndgame/OrthogonalDecomposition/HodgeTateFromSandwich.lean}{PositivityTheory.ThKEnd.Orth.ThPerp}
\leanref{Development/WeilBound/ThetaLifts/EigenvalueEndgame/OrthogonalDecomposition/RapidBound.lean}{PositivityTheory.ThKEnd.Orth.thPerp_of}
\leanref{Development/WeilBound/ThetaLifts/EigenvalueEndgame/OrthogonalDecomposition/RapidBound.lean}{PositivityTheory.ThKEnd.Orth.gen_perp}
\begin{proof}
Let $\mu$ be a Haar measure on $U_W$ and $\mathcal D$ a fundamental domain for
$U_W(F)$ with respect to $\mu$. By Lemma~\ref{asm-theta-agree-end-orth-lift}(2),
$\mu(\mathcal D)<\infty$. As $M$ is regular, $\mu_X$ is a Haar measure and
$\mathcal D_X$ a fundamental domain for $U_X(F)$, so
$\mu_X(\mathcal D_X)<\infty$ by Lemma~\ref{asm-theta-agree-convergence-signature}
and Theorem~\ref{asm-theta-agree-definite-volume}; and $U_X$ and $U_W$ are
second countable by Theorem~\ref{asm-theta-agree-convergence-siegel}(1). Put
$k(g,h)=\theta_\omega(\Psi)(g,h)$, a continuous function on $U_X\times U_W$
(Lemma~\ref{asm-theta-agree-convergence-theta}(2)). The function
$T(h)=\int_{\mathcal D_X}|k(g,h)|\,d\mu_X(g)$ is invariant under left
translation by $U_W(F)$ (Lemma~\ref{asm-theta-agree-anisotropic-automorphy})
and at most $C\operatorname{ht}(h)^k$, so by
Lemma~\ref{asm-theta-agree-end-orth-lift}(1) there is $K'$ with
$|f(h)|\,T(h)\leq K'$ for every $h$. By Tonelli's theorem,
\[
\int_{\mathcal D_X\times\mathcal D}|f(h)\,k(g,h)|\,d(\mu_X\times\mu)=\int_{\mathcal D}|f(h)|\,T(h)\,d\mu(h)\leq K'\mu(\mathcal D)<\infty,
\]
so $(g,h)\mapsto f(h)\overline{k(g,h)}$ is integrable on
$\mathcal D_X\times\mathcal D$. As $\Psi$ converges,
$I(\Psi)(h)=\int_{\mathcal D_X}k(g,h)\,d\mu_X(g)$. By Fubini's theorem,
\[
\int_{\mathcal D}f(h)\overline{I(\Psi)(h)}\,d\mu(h)=\int_{\mathcal D_X}\int_{\mathcal D}f(h)\overline{k(g,h)}\,d\mu(h)\,d\mu_X(g)=\int_{\mathcal D_X}\theta^*_\omega(f,\Psi)(g)\,d\mu_X(g)=0 .
\]
Taking complex conjugates,
$\int_{\mathcal D}I(\Psi)(h)\overline{f(h)}\,d\mu(h)=0$.
\end{proof}

\begin{lemma}\label{asm-theta-agree-end-orth-unbounded}
Let $B$ be a theta base for a CM field $E$
(Definition~\ref{asm-theta-agree-models}(1)), and let $X$ be a lattice space of
rank $n\geq3$ that is Hermitian of signature $(r,s)$ with $r,s\geq1$
(Definition~\ref{asm-theta-agree-doubling}(1)). Then for every $b\in\mathbf R$
there is $g\in U_X$ with $\overline{\operatorname{ht}}(g)>b$
(Definition~\ref{asm-theta-agree-definite-products}(1)).
\end{lemma}
\leanref{Development/WeilBound/ThetaLifts/EigenvalueEndgame/OrthogonalDecomposition/QhtUnbdd.lean}{PositivityTheory.ThKEnd.Orth.qht_unbdd}
\begin{proof}
Suppose that $\overline{\operatorname{ht}}(g)\leq b$ for every $g\in U_X$; then
$b\geq\overline{\operatorname{ht}}(1)\geq0$. Let $J$ be the matrix of $X$ and
$h(u,u')=u^{\mathsf T}J\,c(u')$. By
Lemma~\ref{asm-theta-agree-convergence-signature}, $J$ is nondegenerate
$\epsilon$-Hermitian, so by Theorem~\ref{asm-theta-agree-convergence-siegel}(2)
it has a Siegel Witt torus $a$ of some size $q$, with vectors $x_i,y_i$
(Definition~\ref{asm-theta-agree-convergence-tori}(3)); let $C$ and $k$ be as
in (c) there.

We have $q\geq1$. Indeed, by Proposition~\ref{asm-theta-agree-tower-plane},
with $n_1=n-2\geq1$ and positions $0<1$, there is an adapted plane
$(x,y,\iota,\mathfrak b)$ in a lattice space with the matrix $J$; so
$h(x,x)=0$ and $h(x,y)=1$. If $q=0$, the condition on a Witt torus that a
vector $v$ with $h(v,v)=0$ orthogonal to all $x_i$ and $y_i$ be $0$ applies to
$x$, and gives $x=0$, against $h(x,y)=1$.

Put $A=\max(C,0)(b+1)^k$ and $t=(|A|,\ldots,|A|)\in\mathcal C_q$. There is
$\gamma\in U_X(F)$ with
$\operatorname{ht}(\gamma a(t))<\overline{\operatorname{ht}}(a(t))+1\leq b+1$.
By (c),
\[
e^{q|A|}=e^{t_1+\cdots+t_q}\leq C\operatorname{ht}(\gamma a(t))^k\leq\max(C,0)(b+1)^k=A .
\]
As $q\geq1$, this gives $1+|A|\leq e^{|A|}\leq e^{q|A|}\leq A\leq|A|$, which is
absurd.
\end{proof}

\begin{lemma}\label{asm-theta-agree-end-orth-zero}
Let $(X,W,M,\chi_X,\chi_W,\omega)$ be a Weil pair of signature $(r,s;r',s')$
for a theta base $B$ of a CM field $E$
(Definition~\ref{asm-theta-agree-triple}(1)), with $X$ and $W$ of ranks $n$ and
$d$, such that $r,s\geq1$, $n\geq3$, $d\geq1$, and $(\chi_W)_\tau\neq1$ for
every $F$-embedding $\tau\colon E\to F_{v_1}$
(Definition~\ref{asm-theta-agree-coinvariants-local}(4)). Let $\mu$ be a Haar
measure on $U_W$, $\mathcal D$ a fundamental domain for $U_W(F)$ with respect
to $\mu$, and $f$ a rapidly decreasing function on $U_W$
(Definition~\ref{asm-theta-agree-definite-products}(1)). If the theta lift
$\theta^*_\omega(f,\Psi)$ (Definition~\ref{asm-theta-agree-definite-products}(2))
is constant on $U_X$ for every Schwartz-Bruhat function $\Psi$ on
$\mathbf A^{n\times d}$, then $\theta^*_\omega(f,\Psi)=0$ for every such $\Psi$.
\end{lemma}
\leanref{Development/WeilBound/ThetaLifts/EigenvalueEndgame/OrthogonalDecomposition/HodgeTateFromSandwich.lean}{PositivityTheory.ThKEnd.Orth.ThZero}
\leanref{Development/WeilBound/ThetaLifts/EigenvalueEndgame/OrthogonalDecomposition/QhtUnbdd.lean}{PositivityTheory.ThKEnd.Orth.thZero_of}
\leanref{Development/WeilBound/ThetaLifts/EigenvalueEndgame/OrthogonalDecomposition/ActsOfReads.lean}{PositivityTheory.ThKEnd.Orth.intersection_update_fsHerm_neg_one}
\begin{proof}
Put $I=\{1,\ldots,n\}\times\{1,\ldots,d\}$. By
Lemma~\ref{asm-theta-agree-end-orth-lift}(2), $\mu(\mathcal D)<\infty$, and for
Schwartz-Bruhat $\Psi$ the theta lift $\theta^*_\omega(f,\Psi)$ is defined,
conjugate-linear in $\Psi$, and
$\theta^*_\omega(f,\omega(u,1)\Psi)(1)=\theta^*_\omega(f,\Psi)(u)$ for
$u\in U_X$.

\emph{Step 1: full-rank data.} Let $\Psi\in\mathcal S^\circ(M)$
(Definition~\ref{asm-theta-agree-models}(5)). It is Schwartz-Bruhat, as the
$\mathrm{vac}_v$ are and equal the indicator function of $\mathcal O_v^I$ for
almost all $v$. We show $\theta^*_\omega(f,\Psi)(1)=0$. By
Lemma~\ref{asm-theta-agree-decay-full} with $N=1$ there are $C$ and $k$ with
$|\theta_\omega(\Psi)(g,h)|\,\overline{\operatorname{ht}}(g)\leq C\operatorname{ht}(h)^k$.
Let $K$ be given by Lemma~\ref{asm-theta-agree-end-orth-lift}(1) for $k$. For
fixed $g$, the function
$T(h)=|\theta_\omega(\Psi)(g,h)|\,\overline{\operatorname{ht}}(g)$ is invariant
under left translation by $U_W(F)$
(Lemma~\ref{asm-theta-agree-anisotropic-automorphy}), so
$|f(h)\,\theta_\omega(\Psi)(g,h)|\,\overline{\operatorname{ht}}(g)\leq\max(C,0)K$
for all $h$, and
\[
|\theta^*_\omega(f,\Psi)(g)|\,\overline{\operatorname{ht}}(g)\leq\max(C,0)\,K\mu(\mathcal D).
\]
As $\theta^*_\omega(f,\Psi)(g)=\theta^*_\omega(f,\Psi)(1)$ for every $g$, if
$\theta^*_\omega(f,\Psi)(1)\neq0$ then $\overline{\operatorname{ht}}$ is
bounded on $U_X$ by $\max(C,0)K\mu(\mathcal D)/|\theta^*_\omega(f,\Psi)(1)|$.
This contradicts Lemma~\ref{asm-theta-agree-end-orth-unbounded}, as $X$ is
Hermitian of signature $(r,s)$ with $r,s\geq1$ and $n\geq3$.

\emph{Step 2: reduction to products.} Let $K_0$ be the set of the
Schwartz-Bruhat $\Psi$ with $\theta^*_\omega(f,\Psi)(1)=0$. By the
conjugate-linearity it is a linear subspace, and it contains
$\mathcal S^\circ(M)$ by Step 1. Every Schwartz-Bruhat function on
$\mathbf A^I$ is a finite linear combination of products
$\Phi=f_0\otimes\bigotimes_v\phi_v$ with $f_0\in\mathcal S(F_\infty^I)$, every
$\phi_v$ Schwartz-Bruhat on $F_v^I$, and $\phi_v$ the indicator function of
$\mathcal O_v^I$ for almost all $v$. So it suffices to show that such a $\Phi$
lies in $K_0$.

\emph{Step 3: the place $v_1$.} For a function $\psi$ on $F_{v_1}^I$ let
$T(\psi)$ be the product $\Phi$ with its factor at $v_1$ replaced by $\psi$.
Then $T$ is linear, $T(\phi_{v_1})=\Phi$, $T(\psi)$ is Schwartz-Bruhat when
$\psi$ is, and if an operator $S$ acts at $v_1$ by $\psi\mapsto\psi_1$
(Definition~\ref{asm-theta-agree-models}(2)), with $\psi$ Schwartz-Bruhat, then
$S(T(\psi))=T(\psi_1)$, since $T(\psi)$ is $\psi$ times a function that does
not depend on the component at $v_1$.

By Lemma~\ref{asm-theta-agree-coinvariants-inputs}(3) for $v_1$ and $\tau_1$
there are a linear model $\iota$ of the second kind for $(\tau_1,1,1)$
(Definition~\ref{asm-theta-agree-models}(5)) and unitary smooth characters
$\chi_1,\chi_2$ of $F_{v_1}^\times$ such that $\iota$ is a dictionary at $v_1$
with characters $(\chi_1,\chi_2)$
(Definition~\ref{asm-theta-agree-coinvariants-local}(3)); and $\chi_1\neq1$,
since otherwise, as $n,d\geq1$, $(\chi_W)_{\tau_1}=1$ or
$(\chi_W)_{\tau_1\circ c}=1$, against the hypothesis. By
Lemma~\ref{asm-theta-agree-coinvariants-inputs}(1),
$\iota\phi_{v_1}=\Psi_0+z'$ with $\Psi_0$ of full rank and
$z'\in\mathfrak a(\chi_1,\chi_2)$. Write $\Psi_0=\iota\phi_0$ with $\phi_0$
Schwartz-Bruhat. Each generator
$\omega_{\chi_1,\chi_2}(g_0,1)\iota\psi-\iota\psi$ of $\mathfrak a(\chi_1,\chi_2)$,
with $\psi$ Schwartz-Bruhat, is $\iota(\psi_1-\psi)$ for some $g\in U_X^{(v_1)}$
and Schwartz-Bruhat $\psi_1$ with $\omega(g,1)$ acting at $v_1$ by
$\psi\mapsto\psi_1$, by (a) of the definition of a dictionary. As $\iota$ is
linear and injective on Schwartz-Bruhat functions, $\phi_{v_1}-\phi_0$ lies in
the augmentation space $\mathfrak a_{v_1}(\omega)$
(Definition~\ref{asm-theta-agree-rechoice-terms}(3)). The function $\phi_0$ has
full rank at $\tau_1$, as $\iota\phi_0=\Psi_0$ has full rank.

So $T(\phi_0)\in\mathcal S^\circ(M)$: its factors are Schwartz-Bruhat, equal to
$\mathrm{vac}_v$ for almost all $v$ (as $M$ is regular), and its factor at
$v_1$ has full rank at $\tau_1$. By Step 1, $T(\phi_0)\in K_0$. Next, $T$ maps
$\mathfrak a_{v_1}(\omega)$ into $K_0$. Indeed, let $\psi$ and $\psi_1$ be
Schwartz-Bruhat with $\omega(g,1)$ acting at $v_1$ by $\psi\mapsto\psi_1$ for
some $g\in U_X^{(v_1)}$. Then $T(\psi_1)=\omega(g,1)T(\psi)$, so
$\theta^*_\omega(f,T(\psi_1))(1)=\theta^*_\omega(f,T(\psi))(g)=\theta^*_\omega(f,T(\psi))(1)$
by the hypothesis; hence $T(\psi_1)-T(\psi)\in K_0$, and $T$ is linear. So
$\Phi=T(\phi_{v_1}-\phi_0)+T(\phi_0)\in K_0$.

\emph{Step 4: conclusion.} By Steps 2 and 3, $\theta^*_\omega(f,\Psi)(1)=0$ for
every Schwartz-Bruhat $\Psi$, and by the hypothesis
$\theta^*_\omega(f,\Psi)(g)=\theta^*_\omega(f,\Psi)(1)=0$ for every $g$.
\end{proof}
% Lean: thZero_of (QhtUnbdd.lean). Step 1 is so, with SoFin (its third clause, that the
% elements of So are Schwartz-Bruhat), SoDecay, rapid_bound, qht_unbdd and
% norm_integral_le_of_norm_le_const; Step 2 is the submodule K (th_add, th_smul) and
% KCo.schA_le; Step 3 is DictSt at v1 and tau1, KCo.exists_full with surjectivity_statement
% (the argument of Lemma asm-theta-agree-coinvariants-kernel(1), which does not use n <= d),
% the map T of intersection_update_fsHerm_neg_one, k1 (th_left) and k2.

\begin{theorem}[Rapid decay of cusp forms \cite{HC68,MW95}]\label{asm-theta-agree-end-orth-cusp-rapid}
Let $B$ be a theta base for a CM field $E$
(Definition~\ref{asm-theta-agree-models}(1)), and let $W$ be a lattice space of
rank $d$ (Definition~\ref{asm-theta-agree-doubling}(1)) whose matrix is
nondegenerate $\epsilon$-Hermitian
(Definition~\ref{asm-theta-agree-convergence-tori}(1)). Let $f$ be a function
on $U_W$ with $f(\gamma h)=f(h)$ for all $\gamma\in U_W(F)$ and $h\in U_W$ that
has uniform moderate growth, is smooth at the finite places
(Definition~\ref{asm-theta-agree-end-orth-notions}), and is cuspidal
(Definition~\ref{asm-theta-agree-end-notions}(4)). Then for every $N\geq0$
there is $C$ with $|f(h)|\,\overline{\operatorname{ht}}(h)^N\leq C$ for all
$h\in U_W$ (Definition~\ref{asm-theta-agree-definite-products}(1)).
\end{theorem}
\leanref{Development/WeilBound/ThetaLifts/FrameEstimates/TowerRallisLemma.lean}{PositivityTheory.ThKEst.CuspRapid}
\leanref{Development/WeilBound/ThetaLifts/FrameEstimates/TowerRallisLemma.lean}{PositivityTheory.ThKEst.Rapid}
\leanref{Development/WeilBound/ThetaLifts/FrameEstimates/CuspRapid.lean}{PositivityTheory.ThKEst.cuspRapid}
\leanref{Development/WeilBound/ThetaLifts/FrameEstimates/CuspRapidOf.lean}{PositivityTheory.ThKEst.cuspRapid_of}

\begin{lemma}\label{asm-theta-agree-end-orth-smooth-pure}
Let $(X,W,M,\chi_X,\chi_W,\omega)$ be a Weil pair for a theta base $B$ of a CM
field $E$ (Definition~\ref{asm-theta-agree-triple}(1)), with $X$ and $W$ of
ranks $n$ and $d$, put $I=\{1,\ldots,n\}\times\{1,\ldots,d\}$, and let $\mu$
and $\mathcal D$ be the measure and the domain of $M$. Let $\phi$ be a
Schwartz-Bruhat function on $\mathbf A_f^I$, and suppose that there are a
continuous seminorm $p$ on $\mathcal S(F_\infty^I)$, an integer $k\geq0$ and a
measurable function $G\colon U_X\to[0,\infty]$ with
$\int_{\mathcal D}G\,d\mu<\infty$ and
$|\theta_\omega(f\otimes\phi)(g,h)|\leq p(f)\operatorname{ht}(h)^kG(g)$ for all
$f\in\mathcal S(F_\infty^I)$, $g\in U_X$ and $h\in U_W$. Then for every
$f\in\operatorname{Pol}(M)$ (Definition~\ref{asm-theta-agree-models}(4)) the
function $I(f\otimes\phi)$ on $U_W$ is invariant under left translation by
$U_W(F)$, has uniform moderate growth and is smooth at the finite places
(Definition~\ref{asm-theta-agree-end-orth-notions}).
\end{lemma}
\leanref{Development/WeilBound/ThetaLifts/FrameEstimates/DomContinuous.lean}{PositivityTheory.ThKEst.iUMG_of}
\leanref{Development/WeilBound/ThetaLifts/FrameEstimates/DomContinuous.lean}{PositivityTheory.ThKEst.dom_continuous}
\leanref{Development/WeilBound/ThetaLifts/FrameEstimates/DomContinuous.lean}{PositivityTheory.ThKEst.dom_bound}
\leanref{Development/WeilBound/ThetaLifts/FrameEstimates/OneMemLines.lean}{PositivityTheory.ThKEst.iDer_of}
\leanref{Development/WeilBound/ThetaLifts/FrameEstimates/IInvOfParts.lean}{PositivityTheory.ThKEst.iInv_of_parts}
\begin{proof}
Let $f\in\operatorname{Pol}(M)$, and let $S$ be the set of the functions
$I(f_1\otimes\phi)$ with $f_1\in\operatorname{Pol}(M)$; it contains
$I(f\otimes\phi)$. By Lemma~\ref{asm-theta-agree-continuity-integrals}, for
$m=1$ and $\phi_1=\phi$, every $f_1\otimes\phi$ converges, and
$f_1\mapsto I(f_1\otimes\phi)(h)$ is a continuous linear form on
$\mathcal S(F_\infty^I)$ for every $h\in U_W$; it is linear by
Lemma~\ref{asm-theta-agree-convergence-theta}(2).

\emph{Invariance.} By Lemma~\ref{asm-theta-agree-anisotropic-automorphy},
$\theta_\omega(\Phi)(g,\gamma h)=\theta_\omega(\Phi)(g,h)$ for
$\gamma\in U_W(F)$, so $I(\Phi)(\gamma h)=I(\Phi)(h)$ for every
Schwartz-Bruhat $\Phi$.

\emph{Growth.} For $f_1\in\operatorname{Pol}(M)$ and $h\in U_W$,
\[
|I(f_1\otimes\phi)(h)|\leq\int_{\mathcal D}|\theta_\omega(f_1\otimes\phi)(g,h)|\,d\mu(g)\leq\Bigl(p(f_1)\int_{\mathcal D}G\,d\mu\Bigr)\operatorname{ht}(h)^k,
\]
with the exponent $k$ common to all elements of $S$.

\emph{Continuity.} Let $f_1\in\operatorname{Pol}(M)$ and $h_0\in U_W$. By
Lemma~\ref{asm-theta-agree-convergence-signature}, the matrix of $W$ is
nondegenerate $\epsilon$-Hermitian, so by
Theorem~\ref{asm-theta-agree-convergence-siegel}(1) $U_W$ is locally compact and
second countable, and the height is bounded, by $C$ say, on a compact
neighbourhood $\Omega$ of $h_0$. For $h\in\Omega$ the function
$g\mapsto\theta_\omega(f_1\otimes\phi)(g,h)$ is bounded by
$p(f_1)C^kG(g)$, which is integrable on $\mathcal D$, and it depends
continuously on $h$ by Lemma~\ref{asm-theta-agree-convergence-theta}(2). By
dominated convergence, $I(f_1\otimes\phi)$ is continuous at $h_0$.

\emph{Derivatives.} Let $f_1\in\operatorname{Pol}(M)$ and let $\gamma$ be a line
in $U_W$ (Definition~\ref{asm-theta-agree-matched}(1)). For $t\in\mathbf R$ the
element $\gamma(t)$ is at $\infty$, and by
Lemma~\ref{asm-theta-agree-rechoice-operators}(1) there is an archimedean
operator $T_t$ of $\omega$ at $(1,\gamma(t))$
(Definition~\ref{asm-theta-agree-rechoice-terms}(1)). The constant map $1$ is a
line in $U_X$. By Lemma~\ref{asm-theta-agree-rechoice-lie}, for the lines $1$
and $\gamma$, there is $f_2\in\operatorname{Pol}(M)$ such that
$t\mapsto L(T_tf_1)$ has derivative $L(f_2)$ at $t=0$ for every continuous
linear form $L$ on $\mathcal S(F_\infty^I)$. Let $h\in U_W$, and take
$L(f')=I(f'\otimes\phi)(h)$. As $T_t$ is an archimedean operator,
$\omega(1,\gamma(t))(f_1\otimes\phi)=T_t(f_1)\otimes\phi$; and as $\omega$ is a
representation, $\theta_\omega(\omega(1,\gamma(t))\Phi)(g,h)=\theta_\omega(\Phi)(g,h\gamma(t))$.
So $L(T_tf_1)=I(f_1\otimes\phi)(h\gamma(t))$, and
$t\mapsto I(f_1\otimes\phi)(h\gamma(t))$ has derivative $I(f_2\otimes\phi)(h)$
at $t=0$, for every $h$. As $I(f_2\otimes\phi)\in S$, $I(f\otimes\phi)$ has
uniform moderate growth.

\emph{Smoothness at the finite places.} By
Lemma~\ref{asm-theta-agree-end-stable-level} there is a neighbourhood $U$ of
$1$ in $U_W$ with $\omega(1,u)(f\otimes\phi)=f\otimes\phi$ for every $u\in U$
such that $\pi_w(u)=1$ for every infinite place $w$. For such $u$ and every
$h\in U_W$, as $\omega$ is a representation,
$I(f\otimes\phi)(hu)=I(\omega(1,u)(f\otimes\phi))(h)=I(f\otimes\phi)(h)$.
\end{proof}
% Lean: iUMG_of ThetaBasic GrpTop EpsOf IDer Level IInv (DomContinuous.lean): S and k are
% dom_bound; the continuity is dom_continuous (continuousAt_of_dominated on a compact
% neighbourhood, GrpTop for the height bound); the derivatives are IDer = iDer_of ThetaBasic
% ArchOp ArchLie IInv (OneMemLines.lean), with exists_clm (dom_facts) for L, one_mem_Lines and
% one_mem_UAt, ArchOp (Lemma asm-theta-agree-rechoice-operators(1)) and ArchLie (Lemma
% asm-theta-agree-rechoice-lie); the smoothness is Level (Lemma
% asm-theta-agree-end-stable-level) with the third clause of IInv; the invariance is the first
% clause of IInv (iInv_of_parts, from ThetaBasic). The domination is the Lean's Dom M om F.

\begin{lemma}\label{asm-theta-agree-end-orth-smooth}
Let $(B,G,\mathcal X,\omega,\omega')$ be a triple for a theta base $B$ of a CM
field $E$ (Definition~\ref{asm-theta-agree-triple}), and let $f$ be in its set
of differences (Definition~\ref{asm-theta-agree-end-notions}(5)). Then
$f(\gamma h)=f(h)$ for all $\gamma\in U_W(F)$ and $h\in U_W$, and $f$ has
uniform moderate growth and is smooth at the finite places
(Definition~\ref{asm-theta-agree-end-orth-notions}).
\end{lemma}
\leanref{Development/WeilBound/ThetaLifts/FrameEstimates/SigRapid.lean}{PositivityTheory.ThKEst.Nice}
\leanref{Development/WeilBound/ThetaLifts/FrameEstimates/SigRapid.lean}{PositivityTheory.ThKEst.sig_nice}
\leanref{Development/WeilBound/ThetaLifts/FrameEstimates/TowerRallisLemma.lean}{PositivityTheory.ThKEst.IUMG}
\leanref{Development/WeilBound/ThetaLifts/FrameEstimates/CuspRapid.lean}{PositivityTheory.ThKEst.iUMG}
\leanref{Development/WeilBound/ThetaLifts/FrameEstimates/TowerRallisLemma.lean}{PositivityTheory.ThKEst.IInv}
\leanref{Development/WeilBound/ThetaLifts/FrameEstimates/TowerRallisLemma.lean}{PositivityTheory.ThKEst.UMGLin}
\begin{proof}
Let $G$ have ranks $(a,a')$ and integers $p$ and $q$, put
$I_0=\{1,\ldots,a\}\times\{1,\ldots,a+a'\}$ and
$I'_0=\{1,\ldots,a'\}\times\{1,\ldots,a+a'\}$, and write $f=I(\Phi)-I'(\Phi')$
with $\Phi\in\operatorname{PD}(M)$ and
$\Phi'\in\operatorname{PD}(M')\cap\mathcal S^\circ(M')$. By
Definition~\ref{asm-theta-agree-triple}(4), $(V,W,M,\beta,\beta/\alpha,\omega)$
is a Weil pair of signature $(a-p,p;a+a'-p-q,p+q)$ and
$(V',W,M',\alpha,\alpha/\beta,\omega')$ one of signature
$(q,a'-q;a+a'-p-q,p+q)$. Call a function on $U_W$ \emph{nice} if it is
invariant under left translation by $U_W(F)$, has uniform moderate growth and
is smooth at the finite places. We show that $I(\Phi)$ and $I'(\Phi')$ are
nice.

\emph{Linear combinations.} The nice functions form a linear subspace. The
function $0$ is nice, with the set $\{0\}$. Let $f_1$ and $f'_1$ be nice,
with exponents $k$ and $k'$ and sets $S$ and $S'$ as in
Definition~\ref{asm-theta-agree-end-orth-notions}(1), and neighbourhoods $U$
and $U'$ as in (2) there, and let $a_1,b_1\in\mathbf C$. Then
$a_1f_1+b_1f'_1$ is invariant under $U_W(F)$, and smooth at the finite places
with the neighbourhood $U\cap U'$. It has uniform moderate growth with the
exponent $\max(k,k')$ and the set $S''$ of the functions $a_1g+b_1g'$ with
$g\in S$ and $g'\in S'$: these are continuous; if $|g|\leq C\operatorname{ht}^k$
and $|g'|\leq C'\operatorname{ht}^{k'}$, then
$|a_1g+b_1g'|\leq(|a_1||C|+|b_1||C'|)\operatorname{ht}^{\max(k,k')}$, as heights
are at least $1$; and for a line $\gamma$ with derivatives $g_2\in S$ and
$g'_2\in S'$ of $g$ and $g'$ along $\gamma$, the derivative of $a_1g+b_1g'$ is
$a_1g_2+b_1g'_2\in S''$.

\emph{The first side.} As $2(p+q)+1<\lfloor(a-a')/2\rfloor$,
$a+a'+2\min(a-p,p)\leq a+a'+2p<2a$. So by
Lemma~\ref{asm-theta-agree-convergence-dominated} for the first Weil pair, every
Schwartz-Bruhat function $\phi$ on $\mathbf A_f^{I_0}$ satisfies the hypothesis
of Lemma~\ref{asm-theta-agree-end-orth-smooth-pure}. The space
$\operatorname{PD}(M)$ is spanned by the
$f_\infty\otimes\bigotimes_u\phi_u$ with $f_\infty\in\operatorname{Pol}(M)$,
every $\phi_u$ Schwartz-Bruhat and $\phi_u=\mathrm{vac}_u$ for almost all $u$;
as $M$ is regular, $\bigotimes_u\phi_u$ is then a Schwartz-Bruhat function on
$\mathbf A_f^{I_0}$. So $\Phi=\sum_if_i\otimes\phi_i$ with
$f_i\in\operatorname{Pol}(M)$ and $\phi_i$ Schwartz-Bruhat. By
Lemma~\ref{asm-theta-agree-continuity-integrals}(1),
$I(\Phi)=\sum_iI(f_i\otimes\phi_i)$, and each term is nice by
Lemma~\ref{asm-theta-agree-end-orth-smooth-pure}. So $I(\Phi)$ is nice.

\emph{The second side.} In the same way,
$\Phi'=\sum_{j=1}^mf'_j\otimes\phi'_j$ with $f'_j\in\operatorname{Pol}(M')$ and
$\phi'_j$ Schwartz-Bruhat on $\mathbf A_f^{I'_0}$. We replace the $\phi'_j$ by
full-rank finite parts for $M'$
(Definition~\ref{asm-theta-agree-continuity-sums}(4)). Let $\Lambda\subset\mathbf C^m$
be the subspace of the $c$ with
$f_0\otimes\sum_jc_j\phi'_j\in\mathcal S^\circ(M')$ for every
$f_0\in\mathcal S(F_\infty^{I'_0})$, let $P$ be a linear projection of
$\mathbf C^m$ onto $\Lambda$, let $e_1,\ldots,e_m$ be the standard basis, and put
$\psi_j=\sum_iP(e_j)_i\phi'_i$. As $P(e_j)\in\Lambda$, each $\psi_j$ is a full-rank
finite part. We show $\Phi'=\sum_jf'_j\otimes\psi_j$. By
Lemma~\ref{asm-theta-agree-continuity-pieces}(1),
$\Phi'=\sum_lg_l\otimes\chi_l$ with Schwartz functions $g_l$ and full-rank
finite parts $\chi_l$. Let $x\in F_\infty^{I'_0}$ and $w=(f'_j(x))_j$.
Evaluating the two expressions for $\Phi'$ at the points with archimedean
component $x$ gives $\sum_jw_j\phi'_j=\sum_lg_l(x)\chi_l$, so
$f_0\otimes\sum_jw_j\phi'_j=\sum_lg_l(x)\,f_0\otimes\chi_l\in\mathcal S^\circ(M')$
for every $f_0$, and $w\in\Lambda$. Hence $w=P(w)=\sum_jw_jP(e_j)$, and
\[
\sum_jw_j\psi_j=\sum_i\Bigl(\sum_jw_jP(e_j)_i\Bigr)\phi'_i=\sum_iw_i\phi'_i .
\]
So $\Phi'(x,y)=\sum_jf'_j(x)\psi_j(y)$ for all $x$ and $y$. By
Lemma~\ref{asm-theta-agree-continuity-domination} for the second Weil pair, each
$\psi_j$ satisfies the hypothesis of
Lemma~\ref{asm-theta-agree-end-orth-smooth-pure}. By
Lemma~\ref{asm-theta-agree-continuity-integrals}(1),
$I'(\Phi')=\sum_jI'(f'_j\otimes\psi_j)$, and each term is nice by
Lemma~\ref{asm-theta-agree-end-orth-smooth-pure}. So $I'(\Phi')$ is nice,
and so is $f=I(\Phi)-I'(\Phi')$.
\end{proof}
% Lean: sig_nice ThetaBasic umgLin_of epsOf_of (iUMG ThetaBasic ArchOp ArchLie)
% (iInv_of ThetaBasic) TriplePair DomWeil DomSoU SoFin (SigRapid.lean, with the files of Lemma
% asm-theta-agree-end-orth-smooth-pure about 280 lines). Nice is the Lean's Nice B G.W;
% nice_zero and Nice.lin (umgLin_of, UMGLin.lean, with umgLin_pow_le) are the linear
% combinations; polData_sum writes Phi and Phi'; exists_soFin (with soCoeff, the subspace C,
% and the first clause of SoFin) replaces the phi'_j; nice_sum uses dom_sum (Lemma
% asm-theta-agree-continuity-integrals(1)), IInv and IUMG; the domination is DomWeil (Lemma
% asm-theta-agree-convergence-dominated, by Or.inl from TriplePair) for the first side and
% dom_of_domSoU (Lemma asm-theta-agree-continuity-domination) for the second.

\begin{lemma}\label{asm-theta-agree-end-orth-rapid}
Let $(B,G,\mathcal X,\omega,\omega')$ be a triple for a theta base $B$ of a CM
field $E$ (Definition~\ref{asm-theta-agree-triple}), and let $f$ be a cuspidal
function (Definition~\ref{asm-theta-agree-end-notions}(4)) in its set of
differences (Definition~\ref{asm-theta-agree-end-notions}(5)). Then $f$ is
rapidly decreasing (Definition~\ref{asm-theta-agree-definite-products}(1)).
\end{lemma}
\leanref{Development/WeilBound/ThetaLifts/EigenvalueEndgame/OrthogonalDecomposition/SigThConst.lean}{PositivityTheory.ThKEnd.Orth.SigRapid}
\leanref{Development/WeilBound/ThetaLifts/FrameEstimates/CuspRapid.lean}{PositivityTheory.ThKEst.sigRapid}
\leanref{Development/WeilBound/ThetaLifts/FrameEstimates/SigRapid.lean}{PositivityTheory.ThKEst.sigRapid_of}
\begin{proof}
By Lemma~\ref{asm-theta-agree-end-orth-smooth}, $f$ is invariant under left
translation by $U_W(F)$, has uniform moderate growth and is smooth at the
finite places. In particular $f$ is continuous, being an element of the set $S$
of Definition~\ref{asm-theta-agree-end-orth-notions}(1). The space $W$ is
skew-Hermitian (Definition~\ref{asm-theta-agree-triple}(2)), so its matrix is
nondegenerate $\epsilon$-Hermitian by
Lemma~\ref{asm-theta-agree-convergence-signature}. By
Theorem~\ref{asm-theta-agree-end-orth-cusp-rapid}, for every $N$ there is $C$
with $|f(h)|\,\overline{\operatorname{ht}}(h)^N\leq C$ for all $h$.
\end{proof}

\begin{lemma}\label{asm-theta-agree-end-orth-local-decomp}
Let $H$ be a set, let $\Pi\neq0$ and $Q$ be complex vector spaces, and for
$h\in H$ let $\rho(h)$ and $c(h)$ be linear maps of $\Pi$ and of $Q$. Suppose
that $0$ and $\Pi$ are the only subspaces of $\Pi$ stable under every
$\rho(h)$, and call a linear map $U\colon\Pi\to Q$ an \emph{intertwiner} if
$U\rho(h)=c(h)U$ for every $h\in H$. Let $R\subset Q$ be a subspace stable under
every $c(h)$, and let $f\in R$ be a finite sum of elements $U(x)$ with $U$ an
intertwiner and $x\in\Pi$. Then $f=\sum_{i=1}^mU_i(\xi_i)$ for some
intertwiners $U_i$ with values in $R$ and some $\xi_i\in\Pi$.
\end{lemma}
\leanref{Development/WeilBound/ThetaLifts/EigenvalueEndgame/OrthogonalDecomposition/Imports.lean}{PositivityTheory.ThKEnd.Orth.decomp}
\begin{proof}
Let $A$ be the free associative $\mathbf C$-algebra on the set $H$. The maps
$\rho(h)$ and $c(h)$ make $\Pi$ and $Q$ into $A$-modules; the intertwiners are
the $A$-linear maps $\Pi\to Q$, $R$ is an $A$-submodule of $Q$, and $\Pi$ is a
simple $A$-module. Write $f=\sum_{i=1}^mU'_i(x_i)$ with intertwiners $U'_i$
and $x_i\in\Pi$, and let $T\colon\Pi^m\to Q$ be the $A$-linear map
$T(y)=\sum_iU'_i(y_i)$. The inverse image $N=T^{-1}(R)$ is an $A$-submodule of
$\Pi^m$ that contains $x=(x_1,\ldots,x_m)$. As $\Pi^m$ is a direct sum of simple
$A$-modules, $N$ has a complementary $A$-submodule
\cite[Chapter~XVII, \S2]{Lan02}; let $\pi\colon\Pi^m\to N$ be the projection
along it, which is $A$-linear. For $y\in\Pi$ let $e_i(y)\in\Pi^m$ have $i$-th
component $y$ and the others $0$, and put $U_i=T\circ\pi\circ e_i$. These are
intertwiners with values in $T(N)\subset R$, and
$\sum_iU_i(x_i)=T(\pi(x))=T(x)=f$.
\end{proof}
% Lean: decomp (Imports.lean): A is FreeAlgebra C H, Pi is IsSimpleModule A P by irr, R' is R
% as an A-submodule, T is the sum of the V i o proj i, and the complement is exists_isCompl for
% R'.comap T; the Lean keeps the coefficients z_i of the span apart, with xi_i = z_i x_i.

\begin{lemma}\label{asm-theta-agree-end-orth-local}
Let $(B,G,\mathcal X,\omega,\omega')$ be a triple for a theta base $B$ of a CM
field $E$ (Definition~\ref{asm-theta-agree-triple}), let $f$ be a rapidly
decreasing function (Definition~\ref{asm-theta-agree-definite-products}(1)) in
its set of differences (Definition~\ref{asm-theta-agree-end-notions}(5)), and
let $v\neq v_1$ be a finite place that is clean for $\mathcal X$. Then $f$
comes from $v$ for $\omega$ and for $\omega'$
(Definition~\ref{asm-theta-agree-end-orth-notions}(3)).
\end{lemma}
\leanref{Development/WeilBound/ThetaLifts/EigenvalueEndgame/OrthogonalDecomposition/HodgeTateFromSandwich.lean}{PositivityTheory.ThKEnd.Orth.LocDecomp}
\leanref{Development/WeilBound/ThetaLifts/EigenvalueEndgame/OrthogonalDecomposition/LocDecompOf.lean}{PositivityTheory.ThKEnd.Orth.locDecomp_of}
\begin{proof}
Let $G$ have ranks $(a,a')$, put $I_0=\{1,\ldots,a\}\times\{1,\ldots,a+a'\}$ and
$I'_0=\{1,\ldots,a'\}\times\{1,\ldots,a+a'\}$, and write $f=I(\Phi)-I'(\Phi')$
with $\Phi\in\operatorname{PD}(M)$ and
$\Phi'\in\operatorname{PD}(M')\cap\mathcal S^\circ(M')$. For $h\in U_W$ let $R_h$
be right translation by $h$ on the functions on $U_W$. By
Definition~\ref{asm-theta-agree-triple}(4), $(V,W,M,\beta,\beta/\alpha,\omega)$
and $(V',W,M',\alpha,\alpha/\beta,\omega')$ are Weil pairs.

\emph{The data at $v$.} By Lemmas~\ref{asm-theta-agree-real}
and~\ref{asm-theta-agree-finite} there are real maps $R$ and finite maps
$R_f=(\Pi_u,p_u,p'_u)_u$ matched to the triple, and by
Lemma~\ref{asm-theta-agree-coinvariants} the matched quadruple
$(B,G,\mathcal X,\omega,\omega',R,R_f)$ describes the coinvariants
(Definition~\ref{asm-theta-agree-parts}(3)). As $v$ is clean, a Schwartz-Bruhat
$\phi$ on $F_v^{I_0}$ with $p_v(\phi)=0$ lies in $\mathfrak a_v(\omega)$, and
likewise for $p'_v$ and $\mathfrak a_v(\omega')$. By
Lemma~\ref{asm-theta-agree-decay}, every Schwartz-Bruhat function on
$\mathbf A^{I_0}$ converges for $(M,\omega)$, and every element of
$\mathcal S^\circ(M')$ converges for $(M',\omega')$. By
Lemma~\ref{asm-theta-agree-definite-inputs}(7) there is an $F$-embedding
$t\colon E\to F_v$. By Lemma~\ref{asm-theta-agree-coinvariants-inputs}(3) for
the first Weil pair, a linear model of the second kind for $(t,1,1)$ exists and
is a dictionary at $v$ (Definition~\ref{asm-theta-agree-coinvariants-local}(3)).
So for every $h\in U_W^{(v)}$ there is a linear map $A(h)$ of the functions on
$F_v^{I_0}$, carrying Schwartz-Bruhat functions to Schwartz-Bruhat functions,
such that $\omega(1,h)$ acts at $v$ by $\phi\mapsto A(h)\phi$ for every
Schwartz-Bruhat $\phi$: namely $A(h)=j\circ\omega_{\chi_1,\chi_2}(1,h_0)\circ\iota$,
for such a model $\iota$ with characters $(\chi_1,\chi_2)$, an element $h_0$
given by condition (b) of a dictionary, and a linear map $j$ inverse to $\iota$
on the Schwartz-Bruhat functions. In the same way there are linear maps $A'(h)$
for $\omega'$. By Lemma~\ref{asm-theta-agree-finite-unique}, if $\omega(1,h)$
acts at $v$ by $\phi\mapsto\phi_1$ with $\phi$ Schwartz-Bruhat, then
$\phi_1=A(h)\phi$; likewise for $\omega'$ and $A'(h)$.

\emph{An action on $\Pi_v$.} As $p_v$ maps the Schwartz-Bruhat functions onto
$\Pi_v$, there is a linear map $\sigma$ from $\Pi_v$ to the Schwartz-Bruhat
functions with $p_v\circ\sigma=1$. For $h\in U_W^{(v)}$ put
$\rho(h)=p_v\circ A(h)\circ\sigma$, a linear map of $\Pi_v$. Then
\[
\rho(h)p_v(\phi)=p_v(A(h)\phi),\qquad\rho(h)p'_v(\phi')=p'_v(A'(h)\phi')\tag{b}
\]
for Schwartz-Bruhat $\phi$ on $F_v^{I_0}$ and $\phi'$ on $F_v^{I'_0}$. Indeed,
$p_v(\sigma(p'_v(\phi')))=p'_v(\phi')$, so the condition on matched finite maps
for $h$ (Definition~\ref{asm-theta-agree-matched}(3)), applied to
$\sigma(p'_v(\phi'))$ and $\phi'$, gives the second identity. For the first,
choose a Schwartz-Bruhat $\phi'$ with $p'_v(\phi')=p_v(\phi)$; then
$\rho(h)p_v(\phi)=p'_v(A'(h)\phi')=p_v(A(h)\phi)$, by the second identity and the
same condition for $\phi$ and $\phi'$. The space $\Pi_v$ is not $0$, and $0$ and
$\Pi_v$ are its only subspaces $N$ stable under every $\rho(h)$: for such $N$,
if $p_v(\phi)\in N$ and $\omega(1,h)$ acts at $v$ by $\phi\mapsto\phi_1$, then
$p_v(\phi_1)=p_v(A(h)\phi)=\rho(h)p_v(\phi)\in N$, so the irreducibility
condition on matched finite maps applies to $N$.

\emph{Intertwiners.} Call a linear map $U$ from $\Pi_v$ to the functions on
$U_W$ an \emph{intertwiner} if $U\rho(h)=R_hU$ for every $h\in U_W^{(v)}$. We
show that $I(\Phi)$ is a finite sum of values of intertwiners. The elements of
$\operatorname{PD}(M)$ are Schwartz-Bruhat
(Lemma~\ref{asm-theta-agree-nonvanishing-schwartz}), so they converge, and $I$
is linear on them. So we may assume that
$\Phi=f_\infty\otimes\bigotimes_u\phi_u$ with $f_\infty\in\operatorname{Pol}(M)$,
every $\phi_u$ Schwartz-Bruhat and $\phi_u=\mathrm{vac}_u$ for almost all $u$.
For a function $\psi$ on $F_v^{I_0}$ let $\Phi[\psi]$ be $\Phi$ with $\phi_v$
replaced by $\psi$; for Schwartz-Bruhat $\psi$ it lies in
$\operatorname{PD}(M)$. Put $J(\psi)=I(\Phi[\psi])$ for Schwartz-Bruhat $\psi$;
it is linear in $\psi$. If $g\in U_V^{(v)}$ and $\omega(g,1)$ acts at $v$ by
$\psi\mapsto\psi_1$, then $\Phi[\psi_1]=\omega(g,1)\Phi[\psi]$
(Definition~\ref{asm-theta-agree-models}(2)), so $J(\psi_1)=J(\psi)$ by
Lemma~\ref{asm-theta-agree-anisotropic-invariance}. Hence $J$ vanishes on
$\mathfrak a_v(\omega)$, and $J(\psi)=J(\psi')$ whenever $p_v(\psi)=p_v(\psi')$,
as then $\psi-\psi'\in\mathfrak a_v(\omega)$. Put $U=J\circ\sigma$; then
$U(p_v(\psi))=J(\psi)$ for Schwartz-Bruhat $\psi$. It is an intertwiner: for
$h\in U_W^{(v)}$, $\omega(1,h)\Phi[\psi]=\Phi[A(h)\psi]$, and
$\theta_\omega(\omega(1,h)\Phi_1)(g,x)=\theta_\omega(\Phi_1)(g,xh)$ as $\omega$ is a
representation, so by (b)
\[
U(\rho(h)p_v(\psi))=J(A(h)\psi)=I(\omega(1,h)\Phi[\psi])=R_hI(\Phi[\psi])=R_hU(p_v(\psi)).
\]
And $I(\Phi)=J(\phi_v)=U(p_v(\phi_v))$.

The same argument applies to $I'(\Phi')$, with $p'_v$, $A'$, $\omega'$, $V'$
and $\mathcal S^\circ(M')$ in place of $p_v$, $A$, $\omega$, $V$ and
$\operatorname{PD}(M)$, and a linear map $\sigma'$ with $p'_v\circ\sigma'=1$ in
place of $\sigma$. Here $\mathcal S^\circ(M')$ is spanned by products
$f_\infty\otimes\bigotimes_u\phi'_u$ with $f_\infty$ a Schwartz function, every
$\phi'_u$ Schwartz-Bruhat, $\phi'_u=\mathrm{vac}'_u$ for almost all $u$ and
$\phi'_{v_1}$ of full rank at $\tau_1$; replacing $\phi'_v$ by a Schwartz-Bruhat
function keeps such a product in $\mathcal S^\circ(M')$, as $v\neq v_1$; and the
elements of $\mathcal S^\circ(M')$ are Schwartz-Bruhat, as $M'$ is regular, so
they converge. So $f$ is a finite sum of values of intertwiners.

\emph{Decomposition.} The rapidly decreasing functions form a subspace
$\mathcal R$ of the functions on $U_W$ that contains $f$ and is stable under
every $R_h$,
$h\in U_W$. Indeed, let $\varphi$ be rapidly decreasing. Then $R_h\varphi$ is
continuous and invariant under left translation by $U_W(F)$. As the height is
submultiplicative,
$\operatorname{ht}(\gamma x)\leq\operatorname{ht}(\gamma xh)\operatorname{ht}(h^{-1})$
for $\gamma\in U_W(F)$, so
$\overline{\operatorname{ht}}(x)\leq\overline{\operatorname{ht}}(xh)\operatorname{ht}(h^{-1})$;
and if $|\varphi|\,\overline{\operatorname{ht}}^N\leq C$, then
$|\varphi(xh)|\,\overline{\operatorname{ht}}(x)^N\leq C\operatorname{ht}(h^{-1})^N$.
By Lemma~\ref{asm-theta-agree-end-orth-local-decomp}, for the set $U_W^{(v)}$,
the maps $\rho(h)$ and $R_h$ and the subspace $\mathcal R$, there are
intertwiners $U_1,\ldots,U_m$ with values in $\mathcal R$ and
$\xi_1,\ldots,\xi_m\in\Pi_v$ with
$f=\sum_iU_i(\xi_i)$.

\emph{Local maps.} Choose Schwartz-Bruhat $\phi_i$ on $F_v^{I_0}$ with
$p_v(\phi_i)=\xi_i$, and put $T_i=U_i\circ p_v$. Then $T_i$ is a local map at
$v$ for $\omega$: $T_i(\phi)$ is rapidly decreasing; and if $h\in U_W^{(v)}$
and $\omega(1,h)$ acts at $v$ by $\phi\mapsto\phi_1$, with $\phi$ and $\phi_1$
Schwartz-Bruhat, then $\phi_1=A(h)\phi$, so
$T_i(\phi_1)=U_i(\rho(h)p_v(\phi))=R_hT_i(\phi)$ by (b). As
$f=\sum_iT_i(\phi_i)$, $f$ comes from $v$ for $\omega$. In the same way, with
Schwartz-Bruhat $\phi'_i$ on $F_v^{I'_0}$ such that $p'_v(\phi'_i)=\xi_i$, the
maps $T'_i=U_i\circ p'_v$ and (b) for $p'_v$, $f$ comes from $v$ for
$\omega'$.
\end{proof}
% Lean: locDecomp_of ThetaBasic EpsOf Unimod TriplePair KDecay SoFin PolSch DictSt GoodNe
% KReal KFin KCoinv (LocDecompOf.lean, about 250 lines). The data at v: p and p' (TriplePair),
% R (KReal), Rf (KFin), k and k' (the first clause of KCoinv at v), d1 and d2 (KDecay), t
% (GoodNe), A and A' (exists_op, from DictSt with exists_inv' and acts_of_reads). (b) is r'
% and r0 (hf.equiv, hf.onto, hf.onto'), with q for sigma; the irreducibility is hf.irred with
% action_uniqueness. The intertwiners are mem_vals (intersection_update_fsHerm_neg_one for
% Phi[psi], I_add, I_smul', and I_left, the case of Lemma
% asm-theta-agree-anisotropic-invariance), with S the spanning sets of PolData (PolSch, d1)
% and of So (SoFin, d2, and v /= v1). R is rapid with rapid_rt (qht_le, heightGL_mul_le); the
% decomposition is decomp (Lemma asm-theta-agree-end-orth-local-decomp), and comes_of (with
% action_uniqueness) gives the local maps.

\begin{lemma}\label{asm-theta-agree-end-orth-local-pair}
Let $B$ be a theta base for a CM field $E$
(Definition~\ref{asm-theta-agree-models}(1)), and fix a uniformizer
$\varpi_v$ of $F_v$ for every finite place $v$. Let
$(X,W,M,\chi_X,\chi_W,\omega)$ and $(X',W,M',\chi_{X'},\chi'_W,\omega')$ be
Weil pairs for $B$ of signatures $(r,s;r',s')$ and $(r_1,s_1;r',s')$
(Definition~\ref{asm-theta-agree-triple}(1)), with the same space $W$, with
$X$, $X'$ and $W$ of ranks $n$, $n'$ and $d=n+n'$, and with
$\chi_W=\chi_X/\chi_{X'}$. Let $v$ be a finite place and $\tau\colon E\to F_v$
an $F$-embedding with $(\chi_W)_\tau(\varpi_v)\neq1$
(Definition~\ref{asm-theta-agree-coinvariants-local}(4)). Let $\mu$ be a Haar
measure on $U_W$, $\mathcal D$ a fundamental domain for $U_W(F)$ with respect
to $\mu$, and $f$ a function on $U_W$ that comes from $v$ for $\omega'$
(Definition~\ref{asm-theta-agree-end-orth-notions}(3)). Then
$\theta^*_\omega(f,\Psi)(gu)=\theta^*_\omega(f,\Psi)(g)$
(Definition~\ref{asm-theta-agree-definite-products}(2)) for every
Schwartz-Bruhat function $\Psi$ on $\mathbf A^{n\times d}$, every
$u\in U_X^{(v)}$ and every $g\in U_X$.
\end{lemma}
\leanref{Development/WeilBound/ThetaLifts/EigenvalueEndgame/OrthogonalDecomposition/HodgeTateFromSandwich.lean}{PositivityTheory.ThKEnd.Orth.PairLoc}
\leanref{Development/WeilBound/ThetaLifts/EigenvalueEndgame/OrthogonalDecomposition/PairLocOf.lean}{PositivityTheory.ThKEnd.Orth.pairLoc_of}
\leanref{Development/WeilBound/ThetaLifts/EigenvalueEndgame/OrthogonalDecomposition/SigThConstFin.lean}{PositivityTheory.ThKEnd.Orth.splitB}
\begin{proof}
Put $k=F_v$ and $\varpi=\varpi_v$, let $|\cdot|$ be the normalized absolute
value of $k$, and put $\nu=(\chi_W)_\tau$, $\alpha=(\chi_X)_\tau$,
$\nu'=(\chi'_W)_\tau$ and $\alpha'=(\chi_{X'})_\tau$. These characters have
values of absolute value $1$, and $\nu=\alpha/\alpha'$, as
$\chi_W=\chi_X/\chi_{X'}$. Put $I=\{1,\ldots,n\}\times\{1,\ldots,d\}$ and
$I'=\{1,\ldots,n'\}\times\{1,\ldots,d\}$.

\emph{Reduction to one local map.} Write $f=\sum_{i=1}^mT_i(\phi_i)$ with local
maps $T_i$ at $v$ for $\omega'$ and Schwartz-Bruhat functions $\phi_i$ on
$F_v^{I'}$. Each $T_i(\phi_i)$ is rapidly decreasing, so by
Lemma~\ref{asm-theta-agree-end-orth-lift}(2) the function
$x\mapsto T_i(\phi_i)(x)\overline{\theta_\omega(\Psi)(g',x)}$ is integrable on
$\mathcal D$ for every $g'\in U_X$, and
$\theta^*_\omega(f,\Psi)=\sum_i\theta^*_\omega(T_i(\phi_i),\Psi)$. So we may
assume that $f=T(\phi')$ for one local map $T$ at $v$ for $\omega'$ and one
Schwartz-Bruhat function $\phi'$ on $F_v^{I'}$.

\emph{A form.} For Schwartz-Bruhat functions $a$ on $F_v^{I'}$ and $\Psi_1$ on
$\mathbf A^I$ put $P(a,\Psi_1)=\theta^*_\omega(T(a),\Psi_1)(g)$. As $T(a)$ is
rapidly decreasing, Lemma~\ref{asm-theta-agree-end-orth-lift}(2) shows that the
integral defining $P(a,\Psi_1)$ converges absolutely, that $P$ is
conjugate-linear in $\Psi_1$, and, as $T$ is linear, that $P$ is linear in $a$.
\begin{enumerate}
\item[(a)] If $h\in U_W^{(v)}$, the functions $a$ and $a_1$ are Schwartz-Bruhat
  and $\omega'(1,h)$ acts at $v$ by $a\mapsto a_1$, then
  $P(a_1,\omega(1,h)\Psi_1)=P(a,\Psi_1)$ for every Schwartz-Bruhat $\Psi_1$.
\end{enumerate}
Indeed, $T(a_1)(x)=T(a)(xh)$ for $x\in U_W$, as $T$ is a local map, and
$\theta_\omega(\omega(1,h)\Psi_1)(g,x)=\theta_\omega(\Psi_1)(g,xh)$, as $\omega$
is a representation. So $P(a_1,\omega(1,h)\Psi_1)=\int_{\mathcal D}\Lambda(xh)\,d\mu(x)$
and $P(a,\Psi_1)=\int_{\mathcal D}\Lambda\,d\mu$ for
$\Lambda(x)=T(a)(x)\overline{\theta_\omega(\Psi_1)(g,x)}$. The function
$\Lambda$ is invariant under left translation by $U_W(F)$, as $T(a)$ is rapidly
decreasing
and by Lemma~\ref{asm-theta-agree-anisotropic-automorphy}. The matrix of $W$ is
nondegenerate $\epsilon$-Hermitian
(Lemma~\ref{asm-theta-agree-convergence-signature}), so by
Theorem~\ref{asm-theta-agree-definite-volume} the group $U_W(F)$ is countable and
$\mu$ is invariant under right translations. The computation in the proof of
Lemma~\ref{asm-theta-agree-anisotropic-invariance}, for $U_W$ in place of
$U_X$, gives $\int_{\mathcal D}\Lambda(xh)\,d\mu(x)=\int_{\mathcal D}\Lambda\,d\mu$.

\emph{Dictionaries.} By Lemma~\ref{asm-theta-agree-coinvariants-dictionary}, for
the two Weil pairs and the embedding $\tau$, there are linear models $\iota$ and
$\iota'$ of the second kind for $(\tau,1,1)$, for $M$ and for $M'$; every
$h_0\in GL_d(k)$ is $\pi_\tau(h)$ for some $h\in U_W^{(v)}$; and for
$u'\in U_X^{(v)}$ and $h\in U_W^{(v)}$, the operators $\omega(u',1)$ and
$\omega(1,h)$ read as $\omega_{\nu,\alpha}(\pi_\tau(u'),1)$ and
$\omega_{\nu,\alpha}(1,\pi_\tau(h))$ through $\iota$, and $\omega'(1,h)$ reads as
$\omega_{\nu',\alpha'}(1,\pi_\tau(h))$ through $\iota'$
(Definitions~\ref{asm-theta-agree-coinvariants-local}(1)
and~\ref{asm-theta-agree-finite-split-models}(3)). Let $j$ and $j'$ be the
inverses of $\iota$ and $\iota'$ on the Schwartz-Bruhat functions. For a
function $b$ on a space of matrices over $k$ write $b^{h_0}(Y)=b(Yh_0)$ and
${}^{g_0}b(Y)=b(g_0Y)$. Then
\[
\omega_{\nu,\alpha}(1,h_0)b=\alpha(\det h_0)\,|\det h_0|^{n/2}b^{h_0},\qquad
\omega_{\nu,\alpha}(g_0,1)b=\nu(\det g_0)\,|\det g_0|^{d/2}\,{}^{g_0^{\mathsf T}}b,
\]
and $\omega_{\nu',\alpha'}(1,h_0)b=\alpha'(\det h_0)\,|\det h_0|^{n'/2}b^{h_0}$.

\emph{Pure tensors.} The Schwartz-Bruhat functions $\Psi_1$ with
$P(\phi',\omega(u,1)\Psi_1)=P(\phi',\Psi_1)$ form a linear subspace, as $P$ is
conjugate-linear in $\Psi_1$ and $\omega(u,1)$ is linear and preserves the
Schwartz-Bruhat functions (Definition~\ref{asm-theta-agree-models}(6)). Every
Schwartz-Bruhat function on $\mathbf A^I$ is a finite linear combination of
products $f_\infty\otimes\bigotimes_w\phi_w$ with
$f_\infty\in\mathcal S(F_\infty^I)$, every $\phi_w$ Schwartz-Bruhat on $F_w^I$,
and $\phi_w$ the indicator function of $\mathcal O_w^I$ for almost all $w$. So
we may assume that $\Psi$ is such a product. For a function $\psi$ on $F_v^I$
let $\Psi[\psi]$ be $\Psi$ with $\phi_v$ replaced by $\psi$. Then
$\Psi[\psi]=\psi\otimes r$ with $r$ not depending on the $v$-component. So
$\Psi[\psi]$ is Schwartz-Bruhat when $\psi$ is, and an operator that acts at
$v$ by $\psi\mapsto\psi_1$ maps $\Psi[\psi]$ to $\Psi[\psi_1]$
(Definition~\ref{asm-theta-agree-models}(2)).

\emph{A bilinear form.} For Schwartz-Bruhat functions $a$ on $M_{n\times d}(k)$
and $b$ on $M_{n'\times d}(k)$ put $Q(a,b)=P(j'b,\Psi[j\bar a])$. It is
bilinear, as $P$ is conjugate-linear in its second argument. Let
$\xi(x)=\nu(x)|x|^{-d/2}$. Let $h_0\in GL_d(k)$, and choose $h\in U_W^{(v)}$
with $\pi_\tau(h)=h_0$. Then $\omega(1,h)$ acts at $v$ by
$j\bar a\mapsto j\,\omega_{\nu,\alpha}(1,h_0)\bar a$, and $\omega'(1,h)$ by
$j'b\mapsto j'\omega_{\nu',\alpha'}(1,h_0)b$. So by (a) and the formulas above,
\[
Q(a,b)=P\bigl(j'\omega_{\nu',\alpha'}(1,h_0)b,\ \Psi[j\,\omega_{\nu,\alpha}(1,h_0)\bar a]\bigr)
=\alpha'(\det h_0)\,\overline{\alpha(\det h_0)}\,|\det h_0|^{(n+n')/2}\,Q(a^{h_0},b^{h_0}).
\]
As $|\alpha|=1$, $\nu=\alpha/\alpha'$ and $n+n'=d$,
$Q(a^{h_0},b^{h_0})=\xi(\det h_0)\,Q(a,b)$. If $\xi(\varpi)$ were a real number
$t>0$, then $\nu(\varpi)=t|\varpi|^{d/2}$ would be a positive real number of
absolute value $1$, so $\nu(\varpi)=1$, against the hypothesis. By
Lemma~\ref{asm-theta-agree-definite-dist-bilinear},
$Q({}^{g_0}a,b)=\xi(\det g_0)\,Q(a,b)$ for every $g_0\in GL_n(k)$.

\emph{Conclusion.} Put $g_0=\pi_\tau(u)$, $a=\overline{\iota\phi_v}$ and
$b=\iota'\phi'$, so that $\Psi=\Psi[j\bar a]$ and $\phi'=j'b$. As $\omega(u,1)$
acts at $v$ by $\phi_v\mapsto j\,\omega_{\nu,\alpha}(g_0,1)\iota\phi_v$,
\[
\omega(u,1)\Psi=\Psi\bigl[j\,\omega_{\nu,\alpha}(g_0,1)\iota\phi_v\bigr]=\nu(\det g_0)\,|\det g_0|^{d/2}\,\Psi\bigl[j\,\overline{{}^{g_0^{\mathsf T}}a}\bigr],
\]
so, as $\det g_0^{\mathsf T}=\det g_0$,
\[
P(\phi',\omega(u,1)\Psi)=\overline{\nu(\det g_0)}\,|\det g_0|^{d/2}\,Q({}^{g_0^{\mathsf T}}a,b)
=\nu(\det g_0)^{-1}|\det g_0|^{d/2}\,\xi(\det g_0)\,Q(a,b)=P(\phi',\Psi).
\]
As $\omega$ is a representation,
$\theta_\omega(\omega(u,1)\Psi)(g,x)=\theta_\omega(\Psi)(gu,x)$. Hence
\[
\theta^*_\omega(f,\Psi)(gu)=\theta^*_\omega(T(\phi'),\omega(u,1)\Psi)(g)=P(\phi',\omega(u,1)\Psi)=P(\phi',\Psi)=\theta^*_\omega(f,\Psi)(g).
\]
\end{proof}
% Lean: pairLoc_of ThetaBasic EpsOf FinVolU Unimod LocDist splitB (PairLocOf.lean, about 180
% lines). The reduction is th_addf (th_int); the form P is exists_ext for
% (a, Psi) -> th om mu D (T_i a) Psi g, with th_add and th_smul; (a) is th_right (uniW from
% Unimod and EpsOf, L2.integral_right, theta_rat). splitB gives o1, iota, l1 and r1 for the
% first pair and iota', r2 for the second (dictionary_extended, Lemma
% asm-theta-agree-coinvariants-dictionary). The rest is inv_loc2: exists_inv' for j and j',
% schA_le and intersection_update_fsHerm_neg_one for the pure tensors and Psi[psi], beta for
% Q, TC.locBil with LocDist for the bilinear step (xi = nu N^{-1}), and k2 with omB_left,
% omB_right, apow_real and acts_of_reads for the computations.

\begin{lemma}\label{asm-theta-agree-end-orth-constant}
Let $(B,G,\mathcal X,\omega,\omega')$ be a triple of ranks $(a,a')$ for a theta
base $B$ of a CM field $E$ (Definition~\ref{asm-theta-agree-triple}), and let
$f$ be a rapidly decreasing function
(Definition~\ref{asm-theta-agree-definite-products}(1)) in its set of
differences (Definition~\ref{asm-theta-agree-end-notions}(5)). Let $\mu$ be a
Haar measure on $U_W$ and $\mathcal D$ a fundamental domain for $U_W(F)$ with
respect to $\mu$. Then the theta lift $\theta^*_\omega(f,\Psi)$
(Definition~\ref{asm-theta-agree-definite-products}(2)) is constant on $U_V$
for every Schwartz-Bruhat function $\Psi$ on $\mathbf A^{a\times(a+a')}$, and
$\theta^*_{\omega'}(f,\Psi)$ is constant on $U_{V'}$ for every Schwartz-Bruhat
function $\Psi$ on $\mathbf A^{a'\times(a+a')}$.
\end{lemma}
\leanref{Development/WeilBound/ThetaLifts/EigenvalueEndgame/OrthogonalDecomposition/SigThConst.lean}{PositivityTheory.ThKEnd.Orth.SigThConst}
\leanref{Development/WeilBound/ThetaLifts/EigenvalueEndgame/OrthogonalDecomposition/SigThConstFin.lean}{PositivityTheory.ThKEnd.Orth.sigThConst_fin}
\leanref{Development/WeilBound/ThetaLifts/EigenvalueEndgame/OrthogonalDecomposition/SigThConstOf.lean}{PositivityTheory.ThKEnd.Orth.sigThConst_of}
\leanref{Development/WeilBound/ThetaLifts/EigenvalueEndgame/OrthogonalDecomposition/SigThConstOf.lean}{PositivityTheory.ThKEnd.Orth.chrW_inv_ne}
\begin{proof}
Fix a uniformizer $\varpi_v$ of $F_v$ for every finite place $v$. By
Definition~\ref{asm-theta-agree-triple}(4), $(V,W,M,\beta,\beta/\alpha,\omega)$
and $(V',W,M',\alpha,\alpha/\beta,\omega')$ are Weil pairs for $B$, of
signatures $(a-p,p;a+a'-p-q,p+q)$ and $(q,a'-q;a+a'-p-q,p+q)$. As $G$ is
proper, $a>a'>2(p+q)+1\geq1$, so $a,a'\geq2$. Let $T_0=S_0\cup Q\cup\{v_1\}$,
and let $v\notin T_0$ be a finite place such that the rational prime below $v$
splits completely in $E$. Then $v$ is clean for $\mathcal X$ and $v\neq v_1$.
By Lemma~\ref{asm-theta-agree-definite-inputs}(7), whose statement and proof
concern only the triple, there is an $F$-embedding $\tau\colon E\to F_v$, and
$\nu_\tau(\varpi_v)\neq1$; so also
$(\beta/\alpha)_\tau(\varpi_v)=\nu_\tau(\varpi_v)^{-1}\neq1$. By
Lemma~\ref{asm-theta-agree-end-orth-local}, $f$ comes from $v$ for $\omega$ and
for $\omega'$.

Let $\Psi$ be a Schwartz-Bruhat function on $\mathbf A^{a\times(a+a')}$ and put
$\varphi=\theta^*_\omega(f,\Psi)$. By
Lemma~\ref{asm-theta-agree-end-orth-lift}(3) for the first Weil pair, $\varphi$
is continuous and invariant under left translation by $U_V(F)$. For $v$ as
above, $u\in U_V^{(v)}$ and $g\in U_V$,
Lemma~\ref{asm-theta-agree-end-orth-local-pair}, for the first Weil pair, the
second, whose space $V'$ has rank $a'$ with $a+a'$ the rank of $W$, the
character $\beta/\alpha$ of $W$ in the first, which is the quotient of the
characters $\beta$ of $V$ and $\alpha$ of $V'$, and for $v$, $\tau$ and $f$,
which comes from $v$ for $\omega'$, gives $\varphi(gu)=\varphi(g)$. By
Lemma~\ref{asm-theta-agree-definite-inputs}(4) for $V$, of rank $a\geq2$, and
the set $T_0$, $\varphi$ is constant.

The second claim follows in the same way with the two Weil pairs exchanged: the
character $\alpha/\beta$ of $W$ in the second is the quotient of $\alpha$ and
$\beta$, and $(\alpha/\beta)_\tau(\varpi_v)=\nu_\tau(\varpi_v)\neq1$; $f$ comes
from $v$ for $\omega$; and Lemma~\ref{asm-theta-agree-definite-inputs}(4)
applies to $V'$, of rank $a'\geq2$.
\end{proof}
% Lean: sigThConst_fin = sigThConst_of TriplePair GoodNe ConstOfGood LocDecomp PairLoc ThAut
% (SigThConstOf.lean), with LocDecomp = locDecomp_of, PairLoc = pairLoc_of (LocDist, splitB)
% and ThAut = thAut_of. The set T_0 is insert v1 X.T; chrW_inv_ne gives the value of
% beta/alpha.

\begin{lemma}\label{asm-theta-agree-end-orth-iso}
Let $(B,G,\mathcal X,\omega,\omega')$ be a triple for a theta base $B$ of a CM
field $E$ (Definition~\ref{asm-theta-agree-triple}), with models $M$ and $M'$
and theta integrals $I$ and $I'$, and let $f$ be a cuspidal function
(Definition~\ref{asm-theta-agree-end-notions}(4)) in its set of differences
(Definition~\ref{asm-theta-agree-end-notions}(5)). If $p\geq1$,
then $I(\Phi)$ is orthogonal to $f$ for every $\Phi\in\operatorname{PD}(M)$. If
$q\geq1$, then $I'(\Phi')$ is orthogonal to $f$ for every
$\Phi'\in\operatorname{PD}(M')\cap\mathcal S^\circ(M')$.
\end{lemma}
\leanref{Development/WeilBound/ThetaLifts/EigenvalueEndgame/PlaneRemovalConfiguration.lean}{PositivityTheory.ThKEnd.OrthIso}
\leanref{Development/WeilBound/ThetaLifts/EigenvalueEndgame/OrthogonalDecomposition/SigThConstFin.lean}{PositivityTheory.ThKEnd.Orth.orthIso_fin}
\begin{proof}
Let $G$ have ranks $(a,a')$, and let $\mu_V$, $\mathcal D_V$ and $\mu_{V'}$,
$\mathcal D_{V'}$ be the measures and domains of $M$ and $M'$. By
Definition~\ref{asm-theta-agree-triple}(4), $(V,W,M,\beta,\beta/\alpha,\omega)$
is a Weil pair of signature $(a-p,p;a+a'-p-q,p+q)$ and
$(V',W,M',\alpha,\alpha/\beta,\omega')$ one of signature
$(q,a'-q;a+a'-p-q,p+q)$. As $G$ is proper, $a>a'>2(p+q)+1$ and $4p<a-a'$, so
$a+a'+2\min(a-p,p)<2a$, as in the proof of
Lemma~\ref{asm-theta-agree-decay}. By
Lemma~\ref{asm-theta-agree-end-orth-rapid}, $f$ is rapidly decreasing. So by
Lemma~\ref{asm-theta-agree-end-orth-constant}, for every Haar measure $\mu$ on
$U_W$ and every fundamental domain $\mathcal D$ for $U_W(F)$ with respect to
$\mu$, the theta lifts $\theta^*_\omega(f,\Psi)$ and $\theta^*_{\omega'}(f,\Psi')$
are constant on $U_V$ and $U_{V'}$ for all Schwartz-Bruhat $\Psi$ and $\Psi'$.
Fix a uniformizer $\varpi_{v_1}$ of $F_{v_1}$. As $v_1$ is clean for
$\mathcal X$, Lemma~\ref{asm-theta-agree-definite-inputs}(7), whose statement
and proof concern only the triple, gives $\nu_\tau(\varpi_{v_1})\neq1$ for every
$F$-embedding $\tau\colon E\to F_{v_1}$. So $\nu_\tau\neq1$ and
$(\beta/\alpha)_\tau=\nu_\tau^{-1}\neq1$ for every such $\tau$.

Let $p\geq1$ and $\Phi\in\operatorname{PD}(M)$. By
Lemma~\ref{asm-theta-agree-nonvanishing-schwartz}, $\Phi$ is Schwartz-Bruhat,
and by Lemma~\ref{asm-theta-agree-convergence}(1) for the first Weil pair,
$\Phi$ converges and
$\int_{\mathcal D_V}|\theta_\omega(\Phi)(g,h)|\,d\mu_V(g)\leq C\operatorname{ht}(h)^k$
for some $C$ and $k$. Let $\mu$ and $\mathcal D$ be as above.
Lemma~\ref{asm-theta-agree-end-orth-zero} applies to the first Weil pair: its
signature has $a-p\geq1$ and $p\geq1$, its ranks are $a\geq3$ and
$a+a'\geq1$, and $(\beta/\alpha)_\tau\neq1$ for every $\tau$. So
$\theta^*_\omega(f,\Phi)=0$. By Lemma~\ref{asm-theta-agree-end-orth-perp},
$I(\Phi)$ is orthogonal to $f$.

Let $q\geq1$ and $\Phi'\in\operatorname{PD}(M')\cap\mathcal S^\circ(M')$. By
Lemma~\ref{asm-theta-agree-nonvanishing-schwartz}, $\Phi'$ is Schwartz-Bruhat.
By Lemma~\ref{asm-theta-agree-continuity-pieces}(1),
$\Phi'=\sum_if_i\otimes\phi_i$ with full-rank finite parts $\phi_i$ for $M'$;
Lemma~\ref{asm-theta-agree-continuity-domination} for the second Weil pair
gives the bounds required in Lemma~\ref{asm-theta-agree-continuity-integrals}
for the $\phi_i$; so by Lemma~\ref{asm-theta-agree-continuity-integrals}(1),
$\Phi'$ converges and
$\int_{\mathcal D_{V'}}|\theta_{\omega'}(\Phi')(g,h)|\,d\mu_{V'}(g)\leq C\operatorname{ht}(h)^k$
for some $C$ and $k$. Lemma~\ref{asm-theta-agree-end-orth-zero} applies to the
second Weil pair: its signature has $q\geq1$ and $a'-q\geq1$, its ranks are
$a'\geq3$ and $a+a'\geq1$, and $(\alpha/\beta)_\tau=\nu_\tau\neq1$ for every
$\tau$. So $\theta^*_{\omega'}(f,\Phi')=0$ for every $\mu$ and $\mathcal D$,
and $I'(\Phi')$ is orthogonal to $f$ by
Lemma~\ref{asm-theta-agree-end-orth-perp}.
\end{proof}
% Lean: Orth.orthIso_fin = orthIso_of ThetaBasic TriplePair WeilConv DomSoU SoFin PolSch
% GoodNe SigRapid SigThConst ThZero ThPerp (SigThConstOf.lean), with ThZero = thZero_of and
% ThPerp = thPerp_of. The first paragraph is p, p' and w (TriplePair), hr, hk and hn (GoodNe
% at v1, with chrW_inv_ne); the second is PolSch, WeilConv (Or.inl w), ThZero and ThPerp; the
% third is PolSch, SoFin, dom_sum with dom_of_domSoU, ThZero and ThPerp. The proofs of
% Lemmas asm-theta-agree-end-orth-smooth, -orth-local and -orth-local-pair are written out.

\begin{lemma}\label{asm-theta-agree-end-orth-cusp-lift-unfolding}
Let $H$ be a locally compact, second countable Hausdorff group, $\mu$ a Haar
measure on $H$, $\Gamma\subset H$ a countable subgroup and $\mathcal D$ a
fundamental domain for $\Gamma$ in $H$ with respect to $\mu$.
\begin{enumerate}
\item Let $\Gamma$ act on a countable set $Q$, and let $T_q$, $q\in Q$, be
  measurable functions on $H$ with $T_{\delta q}(\delta y)=T_q(y)$ for all
  $\delta\in\Gamma$, $q\in Q$ and $y\in H$, and with
  $\sum_q\int_{\mathcal D}|T_q|\,d\mu<\infty$. Suppose that
  $\int_{\mathcal G}T_q\,d\mu=0$ whenever $q\in Q$ and $\mathcal G$ is a
  fundamental domain, with respect to $\mu$, for the stabilizer $\Gamma_q$ of
  $q$ in $\Gamma$ on which $T_q$ is integrable. Then
  $\sum_q\int_{\mathcal D}T_q\,d\mu=0$.
\item Let $N\subset H$ be a closed subgroup with $\delta m\delta^{-1}\in N$ for
  all $\delta\in\Gamma$ and $m\in N$, let $\kappa$ be a Haar measure on $N$, and
  let $\mathcal M$ be a fundamental domain for $\Gamma\cap N$ in $N$ with
  respect to $\kappa$ with $\kappa(\mathcal M)<\infty$. Let $K$ and $\varphi$
  be measurable functions on $H$ with $K(my)=K(y)$ for $m\in N$,
  $K(\delta y)=K(y)$ and $\varphi(\delta y)=\varphi(y)$ for $\delta\in\Gamma$,
  and $\int_{\mathcal M}\varphi(my)\,d\kappa(m)=0$ for every $y\in H$. If
  $K\varphi$ is integrable on $\mathcal D$, then
  $\int_{\mathcal D}K\varphi\,d\mu=0$.
\end{enumerate}
\end{lemma}
\leanref{Development/AbelianVarietyTheory/Proofs/TwistedOperators/TsumSetIntegralFundamentalDomainEqZeroByOrbits.lean}{GlobalTheory.OrbitUnfold.tsum_setIntegral_eq_zero}
\leanref{Development/AbelianVarietyTheory/Proofs/TwistedOperators/UnfoldingZeroMeanNormalizedSubgroup.lean}{GlobalTheory.UnfoldNormal.unfold_zero_mean}
\begin{proof}
(1) Let $q_0\in Q$, and for each $q$ in its orbit choose $\delta_q\in\Gamma$
with $\delta_qq_0=q$. Then $T_q(y)=T_{q_0}(\delta_q^{-1}y)$, and the
$\delta_q^{-1}$ are representatives of the cosets $\Gamma_{q_0}\delta$,
$\delta\in\Gamma$, as $\Gamma_{q_0}\delta\mapsto\delta^{-1}q_0$ is a bijection onto
the orbit. So $\mathcal G=\bigcup_q\delta_q^{-1}\mathcal D$ is a fundamental
domain for $\Gamma_{q_0}$, and the union is disjoint up to null sets. By the
invariance of $\mu$,
$\int_{\mathcal D}T_q\,d\mu=\int_{\delta_q^{-1}\mathcal D}T_{q_0}\,d\mu$, and
likewise for $|T_q|$. Hence $T_{q_0}$ is integrable on $\mathcal G$, and the sum
of the $\int_{\mathcal D}T_q\,d\mu$ over the orbit of $q_0$ is
$\int_{\mathcal G}T_{q_0}\,d\mu=0$. As the sum over $Q$ converges absolutely, it
is the sum over the orbits of these partial sums, which is $0$.

(2) The group $N$ is locally compact and second countable. The measure
$\kappa$ is invariant under right translations: for $m_0\in N$,
$\mathcal Mm_0$ is again a fundamental domain for $\Gamma\cap N$, and two
fundamental domains of the countable group $\Gamma\cap N$ have the same measure,
while $\kappa(\mathcal Mm_0)$ is $\kappa(\mathcal M)$ times the modulus of
$m_0$; and $0<\kappa(\mathcal M)<\infty$, as countably many translates of
$\mathcal M$ cover $N$ up to a null set. So $\kappa$ is also invariant under
$m\mapsto m^{-1}$. In the same way, for $\delta\in\Gamma$ the automorphism
$m\mapsto\delta m\delta^{-1}$ of $N$ maps $\kappa$ to a multiple of itself and
$\mathcal M$ to a fundamental domain for $\Gamma\cap N$, so it preserves
$\kappa$.

Let $\mu'$ be the image of $\kappa\times\mu$ under
$\mathrm{sh}(m,y)=(my,y)$, a measure on $H\times H$, and let $G$ be the group of
the $(\delta_1,\delta_2)\in\Gamma\times\Gamma$ with
$\delta_1\delta_2^{-1}\in N$, acting on $H\times H$ by left multiplication. As
$(\delta_1,\delta_2)\,\mathrm{sh}(m,y)=\mathrm{sh}(\delta_1\delta_2^{-1}\cdot\delta_2m\delta_2^{-1},\delta_2y)$,
the measure $\mu'$ is invariant under $G$, and $S=\mathrm{sh}(\mathcal M\times\mathcal D)$
is a fundamental domain for $G$ with respect to $\mu'$: given $\mathrm{sh}(m,y)$,
choose $\delta_2$ with $\delta_2y\in\mathcal D$ and then $\lambda\in\Gamma\cap N$
with $\lambda\delta_2m\delta_2^{-1}\in\mathcal M$, and take
$\delta_1=\lambda\delta_2$. The swap $\tau(x,y)=(y,x)$ preserves $\mu'$, as
$\tau(\mathrm{sh}(m,y))=\mathrm{sh}(m^{-1},my)$ and $(m,y)\mapsto(m^{-1},my)$
preserves $\kappa\times\mu$; and $\tau$ maps $G$-orbits to $G$-orbits, so
$\tau(S)$ is also a fundamental domain for $G$. The function
$\Phi(x,y)=K(x)\varphi(x)$ is invariant under $G$, so its integrals over the two
fundamental domains agree:
$\int_S\Phi\,d\mu'=\int_{\tau(S)}\Phi\,d\mu'=\int_S\Phi\circ\tau\,d\mu'$. Here
$\int_S|\Phi\circ\tau|\,d\mu'=\kappa(\mathcal M)\int_{\mathcal D}|K\varphi|\,d\mu<\infty$,
which justifies the use of Fubini's theorem below. Now, as $K(my)=K(y)$,
\[
\int_S\Phi\,d\mu'=\int_{\mathcal D}K(y)\int_{\mathcal M}\varphi(my)\,d\kappa(m)\,d\mu(y)=0,
\qquad
\int_S\Phi\circ\tau\,d\mu'=\kappa(\mathcal M)\int_{\mathcal D}K\varphi\,d\mu .
\]
As $\kappa(\mathcal M)\neq0$, $\int_{\mathcal D}K\varphi\,d\mu=0$.
\end{proof}
% Lean: (1) is OrbitUnfold.tsum_setIntegral_eq_zero (isFundamentalDomain_iUnion_pull,
% setIntegral_pull, with the decomposition of Y into orbits); (2) is
% UnfoldNormal.unfold_zero_mean, with haar_map_mul_right_eq_of_finite_fundamentalDomain and
% haar_map_conj_normalizing_lattice_eq_self for the invariances of kappa, pairGroup for G,
% shear for sh, smulInvariantMeasure_sheared_product_pairGroup,
% ae_exists_pairGroup_smul_mem_sheared_domain, aedisjoint_pairGroup_smul_sheared_domain and
% map_swap_sheared_product_measure_unimodular. About 275 lines with their helper files.

\begin{lemma}\label{asm-theta-agree-end-orth-cusp-lift-inputs}
Assume the meromorphic continuation of Siegel Eisenstein series in the form of
Theorem~\ref{asm-theta-agree-eisenstein} and the Siegel-Weil formula in the form
of Theorem~\ref{asm-theta-agree-siegel-weil}. Let
$(X,W,M,\chi_X,\chi_W,\omega)$ be a Weil pair of signature $(r,s;r',s')$ for a
theta base $B$ of a CM field $E$ (Definition~\ref{asm-theta-agree-triple}(1)),
with $X$ and $W$ of ranks $n,d\geq1$, with $\min(r,s)=0$, and with measure
$\mu_X$ and domain $\mathcal D_X$ of $M$. Put $s_0=(n-d)/2$, and $\kappa=1$ if
$d<n$ and $\kappa=2$ if $d\geq n$. Let $G_W$ be the doubled group of $W$, with
Siegel parabolic subgroup $P_W$, the matrices $t(p)$ for $p\in P_W$ and the
height $\operatorname{ht}$ (Definition~\ref{asm-theta-agree-doubling}(2)); put
$\iota(h,h')=\operatorname{diag}(h,h')\in G_W$ for $h,h'\in U_W$ and
$h_W(u,u')=u^{\mathsf T}J_W\,c(u')$; and let $Q$ be the set of the cosets
$(P_W\cap G_W(F))\gamma$, $\gamma\in G_W(F)$, each with a fixed representative
$\gamma_q\in G_W(F)$.
\begin{enumerate}
\item Let $\gamma\in G_W(F)$ be such that $\gamma\,\iota(1,\delta)^{-1}\notin P_W$
  for every $\delta\in U_W(F)$. Then there is a nonzero $E$-subspace
  $V\subset E^d$ with $h_W(v,v')=0$ for all $v,v'\in V$ such that, for the
  group $N_V$ of Definition~\ref{asm-theta-agree-end-notions}(4): every
  $m\in N_V$ has $\det m=1$, $\gamma\iota(1,m)\gamma^{-1}\in P_W$ and
  $\det t(\gamma\iota(1,m)\gamma^{-1})=1$; and $\delta m\delta^{-1}\in N_V$ for
  every $m\in N_V$ and every $\delta\in U_W(F)$ with
  $\gamma\iota(1,\delta)\gamma^{-1}\in P_W$.
\item There are a compact subgroup $K\subset G_W$, a continuous height function
  $a>0$ for $K$ (Definition~\ref{asm-theta-agree-doubling}(4)) and a real number
  $t_2$ with the following properties. For every Haar measure $\mu$ on $U_W$
  there is $t_0$ such that $h\mapsto a(\iota(1,h))^t$ is $\mu$-integrable on
  $U_W$ for every $t\geq t_0$. For every real $\sigma>t_2$ there are $C$ and
  $N$ such that for every $x\in G_W$ the sum
  $\sum_{q\in Q}a(\gamma_qx)^{\sigma+d/2}$ converges and is at most
  $C\operatorname{ht}(x)^N$.
\item Let $K$ and $a$ be as in (2), and let $\Phi_1,\Phi_2\in\operatorname{PD}(M)$.
  Then there are a continuous function $f$ on $G_W$ with
  $f(px)=\chi_X(\det t(p))\,|\det t(p)|^{s_0+d/2}f(x)$ for all $p\in P_W$ and
  $x\in G_W$, a connected open set $U\subset\mathbf C$ containing $s_0$, a real
  number $t_1$ and a function $\hat E$ on $\mathbf C\times G_W$ such that, with
  $f_s(x)=a(x)^{s-s_0}f(x)$:
  \begin{enumerate}
  \item[(a)] every $s$ with $\operatorname{Re}s>t_1$ lies in $U$, and the sum
    $\sum_{q\in Q}f_s(\gamma_qx)$ converges unconditionally to $\hat E(s,x)$
    for every $x\in G_W$;
  \item[(b)] $s\mapsto\hat E(s,x)$ is holomorphic on $U$ for every $x$, and
    $\hat E$ is continuous on $U\times G_W$;
  \item[(c)] every $s\in U$ has a neighbourhood $B$ such that, for some $C$ and
    $N$, $|\hat E(s',x)|\leq C\operatorname{ht}(x)^N$ for all $s'\in B$ and
    $x\in G_W$;
  \item[(d)] for all $h,h'\in U_W$,
    \[
    \chi_X(\det h')^{-1}\hat E(s_0,\iota(h,h'))=\frac{\kappa}{\mu_X(\mathcal D_X)}\int_{\mathcal D_X}\theta_\omega(\Phi_1)(g,h)\,\overline{\theta_\omega(\Phi_2)(g,h')}\,d\mu_X(g).
    \]
  \end{enumerate}
\end{enumerate}
\end{lemma}
\leanref{Development/WeilBound/ThetaLifts/LatticeBasics/OffOpenAlt.lean}{GlobalTheory.ThK.Dbl.OffOpen}
\leanref{Development/WeilBound/ThetaLifts/LatticeBasics/OffOpen.lean}{GlobalTheory.ThK.Dbl.offOpen}
\leanref{Development/WeilBound/ThetaLifts/LatticeBasics/PairsStable.lean}{GlobalTheory.ThK.Dbl.HgtInt}
\leanref{Development/WeilBound/ThetaLifts/LatticeBasics/ContinuousFlat.lean}{GlobalTheory.ThK.Dbl.Stab2Cpt}
\leanref{Development/WeilBound/ThetaLifts/LatticeBasics/PairsStable.lean}{GlobalTheory.ThK.Dbl.DblW}
\leanref{Development/WeilBound/ThetaLifts/LatticeBasics/PairsStable.lean}{GlobalTheory.ThK.Dbl.BoxSec}
\leanref{Development/WeilBound/ThetaLifts/LatticeBasics/Hypotheses2.lean}{GlobalTheory.ThK.Dbl.EisContinues}
\leanref{Development/WeilBound/ThetaLifts/LatticeBasics/Hypotheses2.lean}{GlobalTheory.ThK.Dbl.IsEisCont}
\begin{proof}
Let $J$ be the matrix of $W$, so that $h_W(u,u')=u^{\mathsf T}Jc(u')$.

(1) Write
$\gamma^{-1}\Delta=\left(\begin{smallmatrix}A\\B\end{smallmatrix}\right)$ with
$A,B\in M_d(E)$. As $\gamma^{-1}$ preserves $J^\square$ and $\Delta$ is
isotropic for it, $A^{\mathsf T}Jc(A)=B^{\mathsf T}Jc(B)$. If $Z$ is a matrix
over $\mathbf A_E$ with $AZ=0$ and $BZ=0$, then $\gamma^{-1}\Delta Z=0$, so
$Z=0$. Let $S\in M_d(E)$ with $ASA=A$, put $Q=1-SA$, so that $AQ=0$, and
$QZ=Z$ whenever $AZ=0$; in particular $Q^2=Q$. Put $C=BQ$, and let $V$ be the
image of $C$ in $E^d$.

\emph{$V$ is isotropic, and orthogonal to the columns of $B$.} Indeed
$B^{\mathsf T}Jc(C)=B^{\mathsf T}Jc(B)c(Q)=A^{\mathsf T}Jc(AQ)=0$, and so
$C^{\mathsf T}Jc(C)=Q^{\mathsf T}B^{\mathsf T}Jc(C)=0$.

\emph{$V\neq0$.} Suppose $C=0$. Then $AQ=0$ and $BQ=0$, so $Q=0$, that is,
$SA=1$, and $A$ is invertible; taking determinants in
$A^{\mathsf T}Jc(A)=B^{\mathsf T}Jc(B)$ shows that $B$ is invertible, as
$\det J\neq0$. Then $\delta=BA^{-1}$ has
$\delta^{\mathsf T}Jc(\delta)=A^{-\mathsf T}A^{\mathsf T}Jc(A)c(A)^{-1}=J$, so
$\delta\in U_W(F)$, and
$\iota(1,\delta^{-1})\gamma^{-1}\Delta=\left(\begin{smallmatrix}A\\\delta^{-1}B\end{smallmatrix}\right)=\Delta A$.
So $\iota(1,\delta^{-1})\gamma^{-1}\in P_W$, and so is its inverse
$\gamma\,\iota(1,\delta^{-1})^{-1}$, against the hypothesis.

\emph{A criterion.} For $m\in U_W$ and $t\in GL_d(\mathbf A_E)$,
$\gamma\iota(1,m)\gamma^{-1}\in P_W$ with $t(\gamma\iota(1,m)\gamma^{-1})=t$ if
and only if $At=A$ and $Bt=mB$: both say that
$\iota(1,m)\gamma^{-1}\Delta=\gamma^{-1}\Delta t$.

\emph{Generators.} Let $m\in U_W$ fix $V$ pointwise, with $mz-z$ in the
$\mathbf A_E$-span of $V$ for every $z\in E^d$ orthogonal to $V$. As the
columns of $B$ are orthogonal to $V$, $mB=B+CX$ for some matrix $X$ over
$\mathbf A_E$. Put $t=1+QX$. Then $At=A$ and $Bt=B+CX=mB$. Further $QXQ=0$, as
$A(QXQ)=0$ and $B(QXQ)=(mB-B)Q=mC-C=0$; so
$\det(1+QX)=\det(1+Q\cdot QX)=\det(1+QX\cdot Q)=1$, and $t$ is invertible with
$\det t=1$. So $\gamma\iota(1,m)\gamma^{-1}\in P_W$ with $\det t=1$. Next,
$\det m=1$. Let $R=c(C)^{\mathsf T}J^{\mathsf T}$, let $S'$ with $RS'R=R$, and
$P'=1-S'R$. Then $RP'=0$, so the columns of $P'$ are orthogonal to $V$ and
$mP'=P'+CY$ for some $Y$. Transposing $m^{\mathsf T}Jc(m)=J$ and using $mC=C$
gives $Rm=R$, so $P'(m-1)=m-1$; and $RC=(C^{\mathsf T}Jc(C))^{\mathsf T}=0$, so
$P'C=C$. Hence
$\det m=\det(1+P'(m-1))=\det(1+(m-1)P')=\det(1+CY)$. Let $T'$ with
$CT'C=C$; then $CT'\cdot CY=CY$ and $CY\cdot CT'=(m-1)P'CT'=(m-1)CT'=0$, so
$\det(1+CY)=\det(1+CT'CY)=\det(1+CYCT')=1$. The elements $m$ of $U_W$ with
$\det m=1$ and $\gamma\iota(1,m)\gamma^{-1}\in P_W$ with $\det t=1$ form a
group, as $t$ is multiplicative; so they include $N_V$.

\emph{Normalization.} Let $\delta\in U_W(F)$ with
$\gamma\iota(1,\delta)\gamma^{-1}\in P_W$. This element lies in $G_W(F)$, so its
$t$ is a matrix $T$ over $E$, and by the criterion $AT=A$ and $BT=\delta B$.
Then $A(TQ)=0$, so $Q(TQ)=TQ$, and $\delta C=BTQ=BQTQ=C(TQ)$: $\delta V\subset V$.
The same holds for $\delta^{-1}$, as
$\gamma\iota(1,\delta^{-1})\gamma^{-1}$ is the inverse, so $\delta V=V$. Let $m$
be a generator of $N_V$. Then $\delta m\delta^{-1}$ fixes $V$ pointwise, and
for $z$ orthogonal to $V$, $\delta^{-1}z$ is orthogonal to $\delta^{-1}V=V$, as
$\delta$ preserves $h_W$; so
$\delta m\delta^{-1}z-z=\delta(m\delta^{-1}z-\delta^{-1}z)$ lies in the span
of $\delta V=V$. So $\delta m\delta^{-1}$ is a generator, and
$\delta N_V\delta^{-1}\subset N_V$.

(2) Let $K$ be the lattice stabilizer of the doubled space of $W$, compact by
Lemma~\ref{asm-theta-agree-doubling-stabilizer}(1), and $a$ its height
function, continuous and positive by
Lemma~\ref{asm-theta-agree-doubling-stabilizer}(3). For a Haar measure $\mu$
on $U_W$, Lemma~\ref{asm-theta-agree-doubling-integrable} gives $t_0$. The
function $a^{d/2}$ is a section of degree $0$ for $K$ and the trivial character,
which is $\eta^0$ on $\mathbf A^\times$: it is continuous,
$a(px)^{d/2}=|\det t(p)|^{d/2}a(x)^{d/2}$ for $p\in P_W$, and it is invariant
under right translation by $K$. Its flat family is $a^{s+d/2}$. Let $U_1$,
$\hat E_1$ and $t_2$ be given by
Lemma~\ref{asm-theta-agree-doubling-continuation}(1) to (4) for this section.
Let $\sigma>t_2$ be real. Then $\sigma\in U_1$, and for every $x$ the sum
$\sum_{q\in Q}a(\gamma_qx)^{\sigma+d/2}$, whose terms do not depend on the
representatives as $a(px)=a(x)$ for $p\in P_W\cap G_W(F)$, converges to
$\hat E_1(\sigma,x)$; by part (4) of that lemma there are $C$ and $N$ with
$\hat E_1(\sigma,x)\leq C\operatorname{ht}(x)^N$ for all $x$.

(3) Let $(I\sqcup I,e^\square,\omega_2,S,l)$ be the doubling datum of
Lemma~\ref{asm-theta-agree-doubling-datum}(1), and
$f(x)=l(\omega_2(x)(\Phi_1\boxtimes\bar\Phi_2))$. As $\min(r,s)=0$,
Lemma~\ref{asm-theta-agree-doubling-section}(1) shows that $f$ is a section of
degree $s_0$ for $(K,\chi_X)$: it is continuous and has the law in (3). Let $U$
and $\hat E$ be given by Lemma~\ref{asm-theta-agree-doubling-value} for the
situation of Lemma~\ref{asm-theta-agree-doubling-siegel-weil}(1). Then $U$ is
open and connected and contains $s_0$, and (a), (b) and (c) are properties (2),
(3) and (4) of Lemma~\ref{asm-theta-agree-doubling-continuation}, with $t_1$
the $t$ of (2). For (d), let $h,h'\in U_W$ and $x=\iota(h,h')$. The value of
$\hat E$ at $s_0$ and the seesaw identity of
Lemma~\ref{asm-theta-agree-doubling-datum}(1) give
\[
\chi_X(\det h')^{-1}\hat E(s_0,x)=\frac{\kappa}{\mu_X(\mathcal D_X)}\int_{\mathcal D_X}\chi_X(\det h')^{-1}\theta\bigl(\omega_2(x)S(g)(\Phi_1\boxtimes\bar\Phi_2)\bigr)\,d\mu_X(g)=\frac{\kappa}{\mu_X(\mathcal D_X)}\int_{\mathcal D_X}\theta_\omega(\Phi_1)(g,h)\,\overline{\theta_\omega(\Phi_2)(g,h')}\,d\mu_X(g).\qedhere
\]
% Lean: these are the parts of the record OrthIn that cusp_lift and ct_zero use, as the
% closure EndCore.lean provides them. (1) is OffOpen = offOpen (OffOpen.lean, about 256
% lines). (2) is HgtEx = hgtEx_of stab2Cpt (K, a, continuity and positivity, and the
% embedding e = iota with IsEmb), Stab2Cpt = stab2Cpt (K compact), HgtInt = hgtInt_of stab2Cpt
% (t_0), and EisContinues = eisContinues_of EisBL FrameConj applied, as in ct_zero, to the
% trivial character, b = 0 and the section a^{d/2} of degree 0 (t_2 from IsEisCont.half and
% the bound from IsEisCont.grow). (3) is BoxSec = boxSec_of FinLevel stab2Cpt (the section
% f(x) = l(om2(x) boxFun Phi_1 (conj Phi_2)) for (K, chi_X)), SWFull = swFull_of SWAniso
% EisContinues BoxSec (U, E-hat, IsEisCont and the value kappa times the average of the theta
% sum over D_X), and DblW = dblW (the doubling datum (dd, om2, S, l) and the seesaw identity
% detInv chi h' * thetaSum (om2 x (S g box)) = theta Phi_1 g h * conj (theta Phi_2 g h') for
% IsEmb h h' x), which turns the value into (d). SWAniso = siegelWeilAniso_of SiegelWeilStd
% FrameConj FinLevel Stab2Cpt. The derivation of SWFull, BoxSec, HgtEx, HgtInt and Stab2Cpt
% is in Lemmas asm-theta-agree-doubling-*, shared with Lemmas
% asm-theta-agree-nonvanishing-zeta and asm-theta-agree-definite-inputs(10), including Lemma
% asm-theta-agree-doubling-level (FinLevel, Gs.gau, Gs.same).
% In (1), offOpen: frame gives A and B (with iso), exists_gi gives S, K = 1 - S A is Q, the
% isotropy is bC and cC, V /= 0 is the first refine goal, the criterion is Off.par_iff with
% Off.inj, the generators are Gen.cols, det_one and Gen.det, and the normalization is norm
% (par_rat) with Gen.conj.
\end{proof}

\begin{lemma}\label{asm-theta-agree-end-orth-cusp-lift}
Assume the meromorphic continuation of Siegel Eisenstein series in the form of
Theorem~\ref{asm-theta-agree-eisenstein} and the Siegel-Weil formula in the form
of Theorem~\ref{asm-theta-agree-siegel-weil}. Let
$(X,W,M,\chi_X,\chi_W,\omega)$ be a Weil pair of signature $(r,s;r',s')$ for a
theta base $B$ of a CM field $E$ (Definition~\ref{asm-theta-agree-triple}(1)),
with $X$ and $W$ of ranks $n,d\geq1$ and $\min(r,s)=0$. Let $\phi$ be a
rapidly decreasing (Definition~\ref{asm-theta-agree-definite-products}(1)) and
cuspidal (Definition~\ref{asm-theta-agree-end-notions}(4)) function on $U_W$,
let $\mu$ be a Haar measure on $U_W$ and $\mathcal D$ a fundamental domain for
$U_W(F)$ with respect to $\mu$, and let $\Phi_2\in\operatorname{PD}(M)$ be such
that the theta lift $\theta^*_\omega(\phi,\Phi_2)$
(Definition~\ref{asm-theta-agree-definite-products}(2)) is constant on $U_X$ and
$\int_{\mathcal D}I(\Phi_2)(h)\overline{\phi(h)}\,d\mu(h)\neq0$. Then
$I(\Phi_1)$ is cuspidal for every $\Phi_1\in\operatorname{PD}(M)$.
\end{lemma}
\leanref{Development/WeilBound/ThetaLifts/LatticeBasics/OrthAn.lean}{GlobalTheory.ThK.Dbl.OrthIn}
\leanref{Development/WeilBound/ThetaLifts/LatticeBasics/OrthAn.lean}{GlobalTheory.ThK.Dbl.cusp_lift}
\leanref{Development/WeilBound/ThetaLifts/LatticeBasics/CuspLift.lean}{GlobalTheory.ThK.Dbl.ct_zero}
\leanref{Development/WeilBound/ThetaLifts/LatticeBasics/UnfoldC.lean}{GlobalTheory.ThK.Dbl.unfoldC}
\begin{proof}
Put $\chi=\chi_X$, and let $\mu_X$ and $\mathcal D_X$ be the measure and the
domain of $M$. Let $s_0$, $\kappa$, $G_W$, $P_W$, $t(p)$, $\iota$, $Q$ and the
$\gamma_q$ be as in Lemma~\ref{asm-theta-agree-end-orth-cusp-lift-inputs}, let
$K$, $a$ and $t_2$ be given by part (2) of that lemma and $t_0$ by (2) for
$\mu$, and let $f$, $U$, $t_1$ and $\hat E$ be given by (3) for $\Phi_1$ and
$\Phi_2$, with $f_s=a^{s-s_0}f$.

\emph{Preliminaries.}
\begin{enumerate}
\item[(P1)] By Lemma~\ref{asm-theta-agree-convergence-signature} the matrices of
  $X$ and $W$ are nondegenerate $\epsilon$-Hermitian. So by
  Theorem~\ref{asm-theta-agree-convergence-siegel}(1) $U_W$ is locally compact,
  second countable and Hausdorff, $U_W(F)$ is countable and the height is
  bounded on compact sets; $\mu(\mathcal D)<\infty$ by
  Theorem~\ref{asm-theta-agree-definite-volume}; and as $\min(r,s)=0$,
  Lemma~\ref{asm-theta-agree-definite-compact} shows that $U_X$ is second
  countable, $\mu_X(\mathcal D_X)<\infty$, and there is a compact
  $\Omega_X\subset U_X$ such that every $g\in U_X$ has $\gamma g\in\Omega_X$
  for some $\gamma\in U_X(F)$.
\item[(P2)] Heights are at least $1$, and
  $\operatorname{ht}(hh')\leq\operatorname{ht}(h)\operatorname{ht}(h')$ and
  $\operatorname{ht}(\iota(h,h'))\leq\operatorname{ht}(h)\operatorname{ht}(h')$
  for $h,h'\in U_W$: at each place the norm of a product is at most the product
  of the norms, and the norm of a block diagonal matrix is the larger of the
  norms of its blocks.
\item[(P3)] If $p\in P_W\cap G_W(F)$, then $t(p)$ has entries in $E$, so
  $\chi(\det t(p))=1$, and $|\det t(p)|=1$ by the product formula. As
  $a(pz)=|\det t(p)|\,a(z)$ for every $p\in P_W$ (write $z=p'k$ with $k\in K$;
  then $t(pp')=t(p)t(p')$), this gives $f_s(pz)=f_s(z)$ and $a(pz)=a(z)$. Also $\chi(\det\delta)=1$ for $\delta\in U_W(F)$, and
  $\iota(\gamma,\gamma')\in G_W(F)$ for $\gamma,\gamma'\in U_W(F)$. For
  $h\in U_W$, $\iota(h,h)\in P_W$ with $t(\iota(h,h))=h$, and $|\det h|=1$ as
  $\det h\cdot c(\det h)=1$; so $f_s(\iota(h,h)z)=\chi(\det h)f_s(z)$.
\item[(P4)] For every Schwartz-Bruhat $\Phi$ on $\mathbf A^{n\times d}$ there are
  $C\geq0$ and $k$ with $|\theta_\omega(\Phi)(g,h)|\leq C\operatorname{ht}(h)^k$
  for all $g$ and $h$. Indeed, Lemma~\ref{asm-theta-agree-definite-inputs}(2)
  gives $C'$ and $k$ with
  $|\theta_\omega(\Phi)(g,h)|\leq C'(\operatorname{ht}(g)\operatorname{ht}(h))^k$;
  for $g\in U_X$ choose $\gamma\in U_X(F)$ with $\gamma g\in\Omega_X$; then
  $\theta_\omega(\Phi)(g,h)=\theta_\omega(\Phi)(\gamma g,h)$
  (Lemma~\ref{asm-theta-agree-anisotropic-automorphy}), and the height is
  bounded on $\Omega_X$ by some $C_0$
  (Theorem~\ref{asm-theta-agree-convergence-siegel}(1)), so $C=|C'|C_0^k$
  works.
\item[(P5)] By Lemma~\ref{asm-theta-agree-end-orth-lift}(1), for every $k\geq0$
  there is $M_k\geq0$ with $|T(h)|\,|\phi(h)|\leq AM_k$ for every $A\geq0$, every
  function $T$ on $U_W$ invariant under left translation by $U_W(F)$ with
  $|T|\leq A\operatorname{ht}^k$, and every $h\in U_W$.
\end{enumerate}

\emph{Step 1: invariance.} For $s\in U$, $\gamma\in G_W(F)$ and $x\in G_W$,
$\hat E(s,\gamma x)=\hat E(s,x)$. Indeed, for $\operatorname{Re}s>t_1$, right
multiplication by $\gamma$ permutes $Q$, and $f_s(\gamma_q\gamma x)$ depends
only on the coset of $\gamma_q\gamma$, by (P3); so the two unconditionally
convergent sums of (3)(a) agree. Both sides are holomorphic in $s$ on the
connected open set $U$ and agree on its nonempty open subset
$\{\operatorname{Re}s>t_1\}$, so they agree on $U$. In the same way, the sums of
part (2) of Lemma~\ref{asm-theta-agree-end-orth-cusp-lift-inputs} at $\gamma x$
and at $x$ agree.

\emph{Step 2: a cuspidal kernel integral.} For $s\in U$ and $h\in U_W$ put
\[
k_s(h)=\int_{\mathcal D}\hat E(s,\iota(h,h'))\,\chi(\det h')^{-1}\phi(h')\,d\mu(h').
\]
We show that $k_s$ is cuspidal. Let $V\subset E^d$ be a nonzero subspace with
$h_W(v,v')=0$ for $v,v'\in V$, $\nu$ a Haar measure on $N_V$, $\mathcal M$ a
fundamental domain for $N_V\cap U_W(F)$ in $N_V$ with respect to $\nu$, and
$h\in U_W$. By Lemma~\ref{asm-theta-agree-end-kill-unipotent}(2), $N_V$ is
closed in $U_W$, locally compact and second countable, and $N_V\cap U_W(F)$ is a
lattice in $N_V$: there is a compact $\Omega_V\subset N_V$ such that every
$u\in N_V$ has $\lambda u\in\Omega_V$ for some $\lambda\in N_V\cap U_W(F)$. And
$\nu(\mathcal M)<\infty$, by the proof of
Lemma~\ref{asm-theta-agree-end-kill-averages}(1). Let $b$ bound the height on
the compact set $\Omega_Vh$.

(i) \emph{A local bound.} Let $s\in U$, and let $B$, $C$ and $N$ be given by
(3)(c). Let $s'\in B\cap U$, $u\in N_V$ and $h'\in U_W$, and choose $\lambda$
as above. As $\iota(\lambda,1)\in G_W(F)$, Step 1 gives
$\hat E(s',\iota(uh,h'))=\hat E(s',\iota(\lambda uh,h'))$. The function
$T(y)=\hat E(s',\iota(\lambda uh,y))$ is invariant under left translation by
$U_W(F)$, by Step 1 for the $\iota(1,\delta)$, and
$|T(y)|\leq|C|\operatorname{ht}(\iota(\lambda uh,y))^N\leq|C|b^N\operatorname{ht}(y)^N$
by (P2). So by (P5), and as $|\chi|=1$,
\[
\bigl|\hat E(s',\iota(uh,h'))\,\chi(\det h')^{-1}\phi(h')\bigr|\leq|C|b^NM_N .
\]

(ii) \emph{Holomorphy.} For $s\in U$ the function
$\Xi_s(u,h')=\hat E(s,\iota(uh,h'))\chi(\det h')^{-1}\phi(h')$ on $N_V\times U_W$
is continuous, by (3)(b), and bounded, by (i). As $\nu(\mathcal M)$ and
$\mu(\mathcal D)$ are finite, it is integrable on $\mathcal M\times\mathcal D$,
and by Fubini's theorem
$Z(s)=\int_{\mathcal M\times\mathcal D}\Xi_s=\int_{\mathcal M}k_s(uh)\,d\nu(u)$.
By (i), each point of $U$ has a disc around it on which the functions
$s\mapsto\Xi_s(u,h')$, holomorphic by (3)(b), are bounded uniformly in $(u,h')$.
By Cauchy's estimates their derivatives are bounded on a smaller disc, and
differentiation under the integral sign shows that $Z$ is holomorphic on $U$.
As $U$ is connected and contains $\{\operatorname{Re}s>t_1\}$, by the identity
theorem it suffices to show $Z(s)=0$ for $\operatorname{Re}s>t$, where
$t=\max(t_1,t_2,t_0-d/2)$.

(iii) \emph{Bounds.} Fix $s$ with $\operatorname{Re}s>t$, and put
$\sigma=\operatorname{Re}s$. As $a(pz)=|\det t(p)|\,a(z)$,
$f_s(pz)=\chi(\det t(p))\,|\det t(p)|^{s+d/2}f_s(z)$ for $p\in P_W$. Writing
$z=pk$ with $p\in P_W$ and $k\in K$, so that $a(z)=|\det t(p)|$,
$|f_s(z)|=a(z)^{\sigma+d/2}|f_s(k)|\leq C_sa(z)^{\sigma+d/2}$, with $C_s$ the
maximum of $|f_s|$ on the compact group $K$. As $\sigma+d/2\geq t_0$, part (2) of
Lemma~\ref{asm-theta-agree-end-orth-cusp-lift-inputs} shows that
$g\mapsto f_s(\iota(1,g))$ is $\mu$-integrable on $U_W$. Further, for
$x\in U_W$,
\[
\sum_{q\in Q}\int_{\mathcal D}\bigl|f_s(\gamma_q\iota(x,y))\,\phi(y)\bigr|\,d\mu(y)<\infty .\tag{$*$}
\]
Indeed, let $C'$ and $N'$ be given by (2) for $\sigma>t_2$, with $C'\geq0$. The
function $T(y)=\sum_qa(\gamma_q\iota(x,y))^{\sigma+d/2}$ is invariant under
left translation by $U_W(F)$, by Step 1, and at most
$C'\operatorname{ht}(x)^{N'}\operatorname{ht}(y)^{N'}$ by (P2). By (P5) and
Tonelli's theorem, the left side of ($*$) is at most
$C_sC'\operatorname{ht}(x)^{N'}M_{N'}\,\mu(\mathcal D)$.

(iv) \emph{Unfolding.} For $x\in U_W$,
\[
\int_{\mathcal D}\hat E(s,\iota(x,y))\,\chi(\det y)^{-1}\phi(y)\,d\mu(y)=\int_{U_W}f_s(\iota(1,g))\,\chi(\det g)^{-1}\phi(xg)\,d\mu(g).\tag{$**$}
\]
Put $T_q(y)=f_s(\gamma_q\iota(x,y))\chi(\det y)^{-1}\phi(y)$ for $q\in Q$. By
(3)(a) and ($*$), the left side of ($**$) is $\sum_q\int_{\mathcal D}T_q\,d\mu$. The
group $U_W(F)$ acts on $Q$ by $\delta\cdot q=(P_W\cap G_W(F))\gamma_q\iota(1,\delta)^{-1}$,
and $T_{\delta\cdot q}(\delta y)=T_q(y)$: as
$\gamma_{\delta\cdot q}\iota(x,\delta y)=p\gamma_q\iota(x,y)$ with
$p\in P_W\cap G_W(F)$, this follows from (P3), $\chi(\det\delta)=1$ and
$\phi(\delta y)=\phi(y)$. Call $q$ \emph{open} if
$\gamma_q\iota(1,\delta)^{-1}\in P_W$ for some $\delta\in U_W(F)$; the action
preserves the open cosets.

The cosets that are not open contribute $0$. We apply
Lemma~\ref{asm-theta-agree-end-orth-cusp-lift-unfolding}(1) to $H=U_W$,
$\Gamma=U_W(F)$ and the functions equal to $T_q$ for $q$ not open and to $0$ for
$q$ open, the sum of the integrals of whose absolute values is finite by ($*$).
Let $q$ be not open. Its stabilizer $\Gamma_q$ is the group of the
$\delta\in U_W(F)$ with $\gamma_q\iota(1,\delta)\gamma_q^{-1}\in P_W$; let
$\mathcal G$ be a fundamental domain for $\Gamma_q$ on which $T_q$ is
integrable. Let $V'$ be given by
Lemma~\ref{asm-theta-agree-end-orth-cusp-lift-inputs}(1) for $\gamma=\gamma_q$.
Then $\Gamma_q$ normalizes $N_{V'}$, and $\Gamma_q\cap N_{V'}=U_W(F)\cap N_{V'}$,
as $\gamma_q\iota(1,m)\gamma_q^{-1}\in P_W$ for $m\in N_{V'}$. By
Lemma~\ref{asm-theta-agree-end-kill-unipotent}(2), $N_{V'}$ is closed and some
Haar measure $\kappa'$ on $N_{V'}$ has a fundamental domain $\mathcal M'$ for
$U_W(F)\cap N_{V'}$ of measure $1$. Put
$L_q(y)=f_s(\gamma_q\iota(x,y))\chi(\det y)^{-1}$. For $m\in N_{V'}$,
$\gamma_q\iota(x,my)=(\gamma_q\iota(1,m)\gamma_q^{-1})\gamma_q\iota(x,y)$, where
$\gamma_q\iota(1,m)\gamma_q^{-1}\in P_W$ has $\det t=1$, and $\det m=1$; so
$L_q(my)=L_q(y)$. For $\delta\in\Gamma_q$,
$\gamma_q\iota(1,\delta)\gamma_q^{-1}\in P_W\cap G_W(F)$, so $L_q(\delta y)=L_q(y)$
by (P3). Also $\phi(\delta y)=\phi(y)$, and
$\int_{\mathcal M'}\phi(my)\,d\kappa'(m)=0$ for every $y$, as $\phi$ is cuspidal.
By Lemma~\ref{asm-theta-agree-end-orth-cusp-lift-unfolding}(2), for $\Gamma_q$
and $N_{V'}$, with $K=L_q$ and $\varphi=\phi$,
$\int_{\mathcal G}T_q\,d\mu=\int_{\mathcal G}L_q\phi\,d\mu=0$. So
the cosets that are not open contribute $0$.

The open cosets. If $\iota(1,\delta)\iota(1,\delta')^{-1}\in P_W$, then
$\operatorname{diag}(1,\delta\delta'^{-1})\Delta=\Delta t$ for some $t$, so
$\delta\delta'^{-1}=t=1$. Hence $\delta\mapsto(P_W\cap G_W(F))\iota(1,\delta)$ is
a bijection from $U_W(F)$ onto the open cosets. Put
$k(g)=f_s(\iota(1,g))\chi(\det g)^{-1}\phi(xg)$; it is $\mu$-integrable on $U_W$
by (iii), as $\phi$ is bounded. For the open coset of $\iota(1,\delta)$, by (P3)
with $\iota(x,\delta y)=\iota(x,x)\iota(1,x^{-1}\delta y)$,
\[
T_q(y)=f_s(\iota(x,\delta y))\chi(\det y)^{-1}\phi(y)=\chi(\det x)f_s(\iota(1,x^{-1}\delta y))\chi(\det y)^{-1}\phi(\delta y)=k(x^{-1}\delta y),
\]
using $\chi(\det\delta)=1$. So the open cosets contribute
$\sum_{\delta\in U_W(F)}\int_{\mathcal D}k(x^{-1}\delta y)\,d\mu(y)=\int_{U_W}k(x^{-1}y)\,d\mu(y)=\int_{U_W}k\,d\mu$,
as $\mathcal D$ is a fundamental domain and $\mu$ is invariant under left
translations. This proves ($**$).

(v) \emph{Conclusion of Step 2.} By ($**$) for $x=uh$ and Fubini's theorem,
\[
Z(s)=\int_{\mathcal M}\int_{U_W}f_s(\iota(1,g))\chi(\det g)^{-1}\phi(uhg)\,d\mu(g)\,d\nu(u)
=\int_{U_W}f_s(\iota(1,g))\chi(\det g)^{-1}\int_{\mathcal M}\phi(uhg)\,d\nu(u)\,d\mu(g)=0,
\]
the integrand being integrable on $\mathcal M\times U_W$ by (iii), as $\phi$ is
bounded and $\nu(\mathcal M)<\infty$, and the inner integral being $0$ as $\phi$
is cuspidal. So $k_s$ is cuspidal for every $s\in U$.

\emph{Step 3: the theta integral.} By
Lemma~\ref{asm-theta-agree-nonvanishing-schwartz}, $\Phi_1$ and $\Phi_2$ are
Schwartz-Bruhat; let $C_1,k_1$ and $C_2,k_2$ be given by (P4) for them. The
function $h'\mapsto|\theta_\omega(\Phi_2)(g,h')|$ is invariant under left
translation by $U_W(F)$ (Lemma~\ref{asm-theta-agree-anisotropic-automorphy}), so
by (P5) $|\theta_\omega(\Phi_2)(g,h')|\,|\phi(h')|\leq C_2M$ for all $g$ and
$h'$, with $M=M_{k_2}$. The functions integrated below are continuous on
$U_X\times U_W$ (Lemma~\ref{asm-theta-agree-convergence-theta}(2)) and bounded,
and $\mu_X(\mathcal D_X)$ and $\mu(\mathcal D)$ are finite, so Fubini's theorem
applies. As $\theta^*_\omega(\phi,\Phi_2)(g)=\theta^*_\omega(\phi,\Phi_2)(1)$
for every $g$,
\[
\int_{\mathcal D}I(\Phi_2)(h)\overline{\phi(h)}\,d\mu(h)=\int_{\mathcal D_X}\int_{\mathcal D}\theta_\omega(\Phi_2)(g,h)\overline{\phi(h)}\,d\mu(h)\,d\mu_X(g)=\mu_X(\mathcal D_X)\,\overline{\theta^*_\omega(\phi,\Phi_2)(1)} .
\]
By hypothesis this is not $0$, so $\mu_X(\mathcal D_X)\neq0$ and
$\theta^*_\omega(\phi,\Phi_2)(1)\neq0$. Let $h\in U_W$. By (3)(d), with the bound
$C_1\operatorname{ht}(h)^{k_1}C_2M$ for the integrand,
\[
k_{s_0}(h)=\frac{\kappa}{\mu_X(\mathcal D_X)}\int_{\mathcal D_X}\theta_\omega(\Phi_1)(g,h)\int_{\mathcal D}\phi(h')\overline{\theta_\omega(\Phi_2)(g,h')}\,d\mu(h')\,d\mu_X(g)=c\,I(\Phi_1)(h),
\]
with $c=\kappa\,\mu_X(\mathcal D_X)^{-1}\theta^*_\omega(\phi,\Phi_2)(1)\neq0$. As
$s_0\in U$, $k_{s_0}=c\,I(\Phi_1)$ is cuspidal by Step 2. For $U'$, $\nu'$,
$\mathcal F$ and $h''$ as in Definition~\ref{asm-theta-agree-end-notions}(4),
$c\int_{\mathcal F}I(\Phi_1)(uh'')\,d\nu'(u)=\int_{\mathcal F}k_{s_0}(uh'')\,d\nu'(u)=0$,
so $I(\Phi_1)$ is cuspidal.
\end{proof}
% Lean: cusp_lift OrthIn (OrthAn.lean, about 150 lines), with ct_zero (CuspLift.lean),
% unfoldC (UnfoldC.lean) and Lemma asm-theta-agree-end-orth-cusp-lift-inputs for the parts
% of OrthIn. (P1) is top, all (CptAll, Lemma asm-theta-agree-definite-compact), vol, hcpt.fin;
% (P2) is HtEmb = htEmb (heightGL_mul_le, Ht.heightGL_le, one_le_heightGL); (P3) is Law.rat,
% detInv_rat, Law.diag' and absE_det'; (P4) is
% exists_poly_height_bound_theta_of_min_eq_zero (ThetaBasic and Cpt.cov); (P5) is
% IsRapid.bound. Step 1 is iv (An.ident with Law.eis_rat). Step 2 is ct_zero: (i) bk with
% eis_bound and g.cov (UnipUGood); (ii) ck, ik, hol (An.diffOn_integral), An.ident with t =
% max (max t1 t2) (t0 - n/2); (iii) Law.flat, Law.bound (hK), i1 (HgtInt), hfin (with the
% Eisenstein series of a^{n/2} for the trivial character from EisContinues); (iv) is unfoldC:
% cosAct, Neg for the cosets that are not open, neg (OrbitUnfold.tsum_setIntegral_eq_zero,
% OffOpen, mem_stab, UnfoldNormal.unfold_zero_mean), ratE_inj, op and integral_eq_tsum'';
% (v) is integral_integral_swap and the cuspidality hc. Step 3 is hne', hv_th, hmain and hkappa
% with integral_integral_swap_of_bounded_continuous. In the Lean n is the rank of W in
% ct_zero and unfoldC.

\begin{lemma}\label{asm-theta-agree-end-orth-rechoice}
Let $(X,W,M,\chi_X,\chi_W,\omega)$ be a Weil pair of signature $(r,s;r',s')$ for
a theta base $B$ of a CM field $E$ (Definition~\ref{asm-theta-agree-triple}(1)),
with $\min(r,s)=0$. Let $X^\circ$ and $W^\circ$ be lattice spaces of the ranks
of $X$ and $W$; let $X[X^\circ]$ and $W[W^\circ]$ be the lattice spaces with the
matrices and lattices of $X$ and $W$ and the matrices $\theta_w$ and $b_v$ of
$X^\circ$ and $W^\circ$, and $M^\circ$ the model for $(X[X^\circ],W[W^\circ])$
with the isomorphism, vacua, measure and domain of $M$; and suppose that
$(X[X^\circ],W[W^\circ],M^\circ,\chi_X,\chi_W,\omega)$ is a Weil pair of
signature $(r,s;r',s')$. If $I(\Phi)$ is cuspidal
(Definition~\ref{asm-theta-agree-end-notions}(4)) for every
$\Phi\in\operatorname{PD}(M)$, then $I(\Phi)$ is cuspidal for every
$\Phi\in\operatorname{PD}(M^\circ)$.
\end{lemma}
\leanref{Development/WeilBound/ThetaLifts/LatticeBasics/OrthAn.lean}{GlobalTheory.ThK.Dbl.cusp_re}
\begin{proof}
The theta integral of the second Weil pair is $I$, as $M^\circ$ has the measure
and domain of $M$. By Lemma~\ref{asm-theta-agree-rechoice-cartan} for $W$ and
$W[W^\circ]$, which are skew-Hermitian of the same signature $(r',s')$ with the
same matrix, there is $h\in U_W$ at $\infty$ that moves $W[W^\circ]$ to $W$
(Definition~\ref{asm-theta-agree-rechoice-terms}(4)). By
Lemma~\ref{asm-theta-agree-rechoice-move} there is $g\in U_X$ that moves
$X[X^\circ]$ to $X$ such that $T(\operatorname{Pol}(M^\circ))=\operatorname{Pol}(M)$
for every archimedean operator $T$ of $\omega$ at $(g,h)$, and by
Lemma~\ref{asm-theta-agree-rechoice-operators}(1) there is such an operator
$T$.

Let $\Psi\in\operatorname{PD}(M^\circ)$. Then
$\omega(g,h)\Psi\in\operatorname{PD}(M)$: the space $\operatorname{PD}(M^\circ)$
is spanned by the $f\otimes\bigotimes_v\phi_v$ with $f\in\operatorname{Pol}(M^\circ)$,
all $\phi_v$ Schwartz-Bruhat and $\phi_v=\mathrm{vac}_v$ for almost all $v$;
for these
$\omega(g,h)\bigl(f\otimes\bigotimes_v\phi_v\bigr)=T(f)\otimes\bigotimes_v\phi_v$
(Definition~\ref{asm-theta-agree-rechoice-terms}(1)) with
$T(f)\in\operatorname{Pol}(M)$; and $\omega(g,h)$ is linear. Next, $\Psi$ is
Schwartz-Bruhat by Lemma~\ref{asm-theta-agree-nonvanishing-schwartz}, so
$I(\omega(g,1)\Psi)=I(\Psi)$ by
Lemma~\ref{asm-theta-agree-anisotropic-invariance}. As $\omega$ is a
representation, $\omega(g,h)=\omega(1,h)\omega(g,1)$ and
$\theta_\omega(\omega(1,h)\Psi')(x,h')=\theta_\omega(\Psi')(x,h'h)$; so
$I(\omega(g,h)\Psi)(h')=I(\omega(g,1)\Psi)(h'h)=I(\Psi)(h'h)$ for $h'\in U_W$.
By hypothesis, $h'\mapsto I(\Psi)(h'h)$ is cuspidal.

Let $U$, $\nu$, $\mathcal F$ and $h''$ be as in
Definition~\ref{asm-theta-agree-end-notions}(4). Then
$\int_{\mathcal F}I(\Psi)(uh'')\,d\nu(u)=\int_{\mathcal F}I(\Psi)\bigl((uh''h^{-1})h\bigr)\,d\nu(u)=0$.
So $I(\Psi)$ is cuspidal.
\end{proof}
% Lean: cusp_re OrthIn (OrthAn.lean): CartanConj (Lemma asm-theta-agree-rechoice-cartan),
% CartanMove (-rechoice-move) and ArchOp (-rechoice-operators(1)); hgen is the span induction
% on PolData with the first clause of IsArchOp; hI uses I_act (CptAll; the definite case of
% Lemma asm-theta-agree-anisotropic-invariance); key and the last step translate by h.

\begin{lemma}\label{asm-theta-agree-end-orth-constant-term}
Let $B$ be a theta base for a CM field $E$
(Definition~\ref{asm-theta-agree-models}(1)), let $W$ and $W_1$ be
skew-Hermitian lattice spaces of some signatures
(Definition~\ref{asm-theta-agree-doubling}(1)), let $H'$ be an adapted plane in
$W$ over $W_1$ (Definition~\ref{asm-theta-agree-planes}(1)), and let $f$ be a
function on $U_W$.
\begin{enumerate}
\item If $c$ and $c'$ are constant terms of $f$ along $H'$
  (Definition~\ref{asm-theta-agree-planes}(2)), then $c=c'$.
\item If $f$ is cuspidal (Definition~\ref{asm-theta-agree-end-notions}(4)),
  then $0$ is the constant term of $f$ along $H'$.
\end{enumerate}
\end{lemma}
\leanref{Development/WeilBound/ThetaLifts/LatticeBasics/OrthAn.lean}{GlobalTheory.ThK.Dbl.exists_trace_zero_generator_of_isCMExt}
\leanref{Development/WeilBound/ThetaLifts/LatticeBasics/OrthAn.lean}{GlobalTheory.ThK.Dbl.isCT_of_cusp}
\leanref{Development/WeilBound/ThetaLifts/EigenvalueEndgame/UnipotentGroups/UnipUGood.lean}{PositivityTheory.ThKEnd.Grp.UnipLine}
\begin{proof}
Write $H'=(x,y,\iota,\mathfrak b)$ and $h_W(u,u')=u^{\mathsf T}J_W\,c(u')$. By
Lemma~\ref{asm-theta-agree-convergence-signature}, the matrix of $W_1$ has
nonzero determinant, so $(x,y,\iota)$ is a plane in $W$ over $W_1$
(Definition~\ref{asm-theta-agree-end-notions}(1)), and $N_{H'}=N_{Ex}$ by
Lemma~\ref{asm-theta-agree-end-kill-unipotent}. As $h_W(x,y)=1$, $x\neq0$; and
$h_W(u,u')=0$ for $u,u'\in Ex$, as $h_W(x,x)=0$.

(1) By Lemma~\ref{asm-theta-agree-end-kill-unipotent}(2) for $U=Ex$, some Haar
measure $\nu$ on $N_{Ex}=N_{H'}$ has a fundamental domain $\mathcal F$ for
$N_{H'}\cap U_W(F)$ of measure $1$. For every $h\in U_W$,
$c(h)=\int_{\mathcal F}f(uh)\,d\nu(u)=c'(h)$.

(2) As $Ex$ is a nonzero $E$-subspace on which $h_W$ vanishes and $f$ is
cuspidal, $\int_{\mathcal F}f(uh)\,d\nu(u)=0$ for every Haar measure $\nu$ on
$N_{Ex}=N_{H'}$, every fundamental domain $\mathcal F$ for $N_{H'}\cap U_W(F)$
and every $h\in U_W$. So $0$ is the constant term of $f$ along $H'$.
\end{proof}

\begin{lemma}\label{asm-theta-agree-end-orth-an}
Assume the meromorphic continuation of Siegel Eisenstein series in the form of
Theorem~\ref{asm-theta-agree-eisenstein} and the Siegel-Weil formula in the form
of Theorem~\ref{asm-theta-agree-siegel-weil}. Let
$(B,G,\mathcal X,\omega,\omega')$ be a triple for a theta base $B$ of a CM field
$E$ (Definition~\ref{asm-theta-agree-triple}), with models $M$ and $M'$ and
theta integrals $I$ and $I'$, and let $f$ be a cuspidal function
(Definition~\ref{asm-theta-agree-end-notions}(4)) in its set of differences
(Definition~\ref{asm-theta-agree-end-notions}(5)). If $p=0$ and
$q\geq1$, then $I(\Phi)$ is orthogonal to $f$ for every
$\Phi\in\operatorname{PD}(M)$. If $q=0$ and $p\geq1$, then $I'(\Phi')$ is
orthogonal to $f$ for every
$\Phi'\in\operatorname{PD}(M')\cap\mathcal S^\circ(M')$.
\end{lemma}
\leanref{Development/WeilBound/ThetaLifts/EigenvalueEndgame/PlaneRemovalConfiguration.lean}{PositivityTheory.ThKEnd.OrthAn}
\leanref{Development/WeilBound/ThetaLifts/LatticeBasics/OrthAn.lean}{GlobalTheory.ThK.Dbl.orthAn_of}
\begin{proof}
Let $G$ have ranks $(a,a')$. By Definition~\ref{asm-theta-agree-triple}(4),
$(V,W,M,\beta,\beta/\alpha,\omega)$ is a Weil pair of signature
$(a-p,p;a+a'-p-q,p+q)$ and $(V',W,M',\alpha,\alpha/\beta,\omega')$ one of
signature $(q,a'-q;a+a'-p-q,p+q)$. As $G$ is proper, $a>a'>2(p+q)+1\geq1$. By
Lemma~\ref{asm-theta-agree-end-orth-rapid}, $f$ is rapidly decreasing.

\emph{The case $p=0$ and $q\geq1$.} Let $\Phi\in\operatorname{PD}(M)$, let
$\mu$ be a Haar measure on $U_W$ and $\mathcal D$ a fundamental domain for
$U_W(F)$ with respect to $\mu$, and suppose that
$\int_{\mathcal D}I(\Phi)(h)\overline{f(h)}\,d\mu(h)\neq0$. By
Lemma~\ref{asm-theta-agree-nonvanishing-schwartz}, $\Phi$ is Schwartz-Bruhat, so
$\theta^*_\omega(f,\Phi)$ is constant on $U_V$ by
Lemma~\ref{asm-theta-agree-end-orth-constant}. The first Weil pair has
$\min(a-p,p)=0$, and $a,a+a'\geq1$. By
Lemma~\ref{asm-theta-agree-end-orth-cusp-lift} with $\phi=f$ and $\Phi_2=\Phi$,
$I(\Phi_1)$ is cuspidal for every $\Phi_1\in\operatorname{PD}(M)$.

Write $a'=b+2$. By Lemma~\ref{asm-theta-agree-tower-beta} there are lattice
spaces $V^\circ$, $V'^\circ$ and $W^\circ$ such that
$(B,G^\circ,\mathcal X,\omega,\omega')$ is a triple for
$G^\circ=G[V^\circ,V'^\circ,W^\circ]$, a triple datum $G_1$ of ranks $(a,b)$,
adapted planes $H$ and $H'$ such that $(G^\circ,G_1,H,H')$ is a
$\beta$-configuration, and $\omega_1$ and $\omega'_1$ such that
$(B,G_1,\mathcal X,\omega_1,\omega'_1)$ is a triple. Let $M^\circ$ and $M_1$ be
the first models of $G^\circ$ and $G_1$, and $I_1$ the first theta integral of
$G_1$; the first theta integral of $G^\circ$ is $I$
(Definition~\ref{asm-theta-agree-triple}(2)). The first Weil pair of $G^\circ$
has the signature of that of $G$, so by
Lemma~\ref{asm-theta-agree-end-orth-rechoice}, $I(\Psi)$ is cuspidal for every
$\Psi\in\operatorname{PD}(M^\circ)$. By
Definition~\ref{asm-theta-agree-planes}(4), $H'$ is an adapted plane in the
third space of $G^\circ$ over the third space $W_1$ of $G_1$, and the first
spaces of $G^\circ$ and $G_1$ are equal. By
Lemma~\ref{asm-theta-agree-evaluation} for the first Weil pair of $G^\circ$,
$H'$ and $M_1$, there is an evaluation map $e$ for
$(M^\circ,M_1,1,\iota_{H'},\mathcal J_0(H'))$, and
$e(\operatorname{PD}(M^\circ))=\operatorname{PD}(M_1)$. Let $C\neq0$ be given by
Lemma~\ref{asm-theta-agree-ct-beta-zero} for $(G^\circ,G_1,H,H')$ and $e$. As
$p_1=p=0$, Lemma~\ref{asm-theta-agree-nonvanishing} for the triple of $G_1$
gives $\Phi_1\in\operatorname{PD}(M_1)$ with $I_1(\Phi_1)\neq0$; write
$\Phi_1=e(\Psi)$ with $\Psi\in\operatorname{PD}(M^\circ)$.

Let $h_1\in U_{W_1}$. The matrix of $W_1$ has nonzero determinant
(Lemma~\ref{asm-theta-agree-convergence-signature}), so $H'$ is a plane in the
sense of Definition~\ref{asm-theta-agree-end-notions}(1), and by
Lemma~\ref{asm-theta-agree-end-unipotent}(2) there is a Levi element $m$ for
$(1,h_1)$ along $H'$. By Lemma~\ref{asm-theta-agree-end-orth-constant-term}(2),
$0$ is the constant term of $I(\Psi)$ along $H'$, so $(0,0)$ splits $I(\Psi)$
along $H'$ (Definition~\ref{asm-theta-agree-planes}(3)), as $0$ has every law.
Lemma~\ref{asm-theta-agree-ct-beta-zero}, for the triples of $G^\circ$ and
$G_1$, the split $(0,0)$, $h_1$, $m$ and $h=1$, gives
$0=C\,I_1(e(\Psi))(h_1)$. So $I_1(\Phi_1)(h_1)=0$ for every $h_1$, against the
choice of $\Phi_1$. Hence
$\int_{\mathcal D}I(\Phi)(h)\overline{f(h)}\,d\mu(h)=0$ for all $\mu$ and
$\mathcal D$, and $I(\Phi)$ is orthogonal to $f$.

\emph{The case $q=0$ and $p\geq1$.} Let
$\Phi'\in\operatorname{PD}(M')\cap\mathcal S^\circ(M')$, let $\mu$ and
$\mathcal D$ be as above, and suppose that
$\int_{\mathcal D}I'(\Phi')(h)\overline{f(h)}\,d\mu(h)\neq0$. By
Lemma~\ref{asm-theta-agree-nonvanishing-schwartz}, $\Phi'$ is Schwartz-Bruhat,
so $\theta^*_{\omega'}(f,\Phi')$ is constant on $U_{V'}$ by
Lemma~\ref{asm-theta-agree-end-orth-constant}. The second Weil pair has
$\min(q,a'-q)=0$, and $a',a+a'\geq1$. By
Lemma~\ref{asm-theta-agree-end-orth-cusp-lift} with $\phi=f$ and $\Phi_2=\Phi'$,
$I'(\Phi_1)$ is cuspidal for every $\Phi_1\in\operatorname{PD}(M')$.

Write $a=b+2$. By Lemma~\ref{asm-theta-agree-tower-alpha} there are lattice
spaces $V^\circ$, $V'^\circ$ and $W^\circ$ such that
$(B,G^\circ,\mathcal X,\omega,\omega')$ is a triple for
$G^\circ=G[V^\circ,V'^\circ,W^\circ]$, a triple datum $G_1$ of ranks $(b,a')$,
adapted planes $H$ and $H'$ such that $(G^\circ,G_1,H,H')$ is an
$\alpha$-configuration, and $\omega_1$ and $\omega'_1$ such that
$(B,G_1,\mathcal X,\omega_1,\omega'_1)$ is a triple. Let $M'^\circ$ and $M'_1$
be the second models of $G^\circ$ and $G_1$, and $I'_1$ the second theta
integral of $G_1$; the second theta integral of $G^\circ$ is $I'$. By
Lemma~\ref{asm-theta-agree-end-orth-rechoice} for the second Weil pairs of $G$
and $G^\circ$, $I'(\Psi)$ is cuspidal for every
$\Psi\in\operatorname{PD}(M'^\circ)$. The second spaces of $G^\circ$ and $G_1$
are equal (Definition~\ref{asm-theta-agree-planes}(4)), so
Lemma~\ref{asm-theta-agree-evaluation} gives an evaluation map $e'$ for
$(M'^\circ,M'_1,1,\iota_{H'},\mathcal J_0(H'))$ with
$e'(\operatorname{PD}(M'^\circ))=\operatorname{PD}(M'_1)$. Let $C\neq0$ be given
by Lemma~\ref{asm-theta-agree-ct-alpha-zero} for $(G^\circ,G_1,H,H')$ and $e'$.
As $q_1=q=0$, Lemma~\ref{asm-theta-agree-nonvanishing} for the triple of $G_1$
gives $\Phi_1\in\operatorname{PD}(M'_1)\cap\mathcal S^\circ(M'_1)$ with
$I'_1(\Phi_1)\neq0$; write $\Phi_1=e'(\Psi)$ with
$\Psi\in\operatorname{PD}(M'^\circ)$.

Let $h_1\in U_{W_1}$, and let $m$ be a Levi element for $(1,h_1)$ along $H'$, as
above. By Lemma~\ref{asm-theta-agree-ct-alpha-zero} there is $f'_\alpha$ such
that $(f'_\alpha,0)$ splits $I'(\Psi)$ along $H'$ and
$f'_\alpha(mh)=C\,I'_1\bigl(e'(\omega'(1,h)\Psi)\bigr)(h_1)$ for every
$h\in U_W$. Then $f'_\alpha$ is the constant term of $I'(\Psi)$ along $H'$, and
so is $0$, by Lemma~\ref{asm-theta-agree-end-orth-constant-term}(2); by part (1)
of that lemma, $f'_\alpha=0$. With $h=1$ this gives
$0=C\,I'_1(e'(\Psi))(h_1)$. So $I'_1(\Phi_1)=0$, against the choice of
$\Phi_1$. Hence $I'(\Phi')$ is orthogonal to $f$.
\end{proof}
% Lean: Dbl.orthAn_of OrthIn SigRapid SigThConst NonZeroAn TowerA TowerB EvZero CTBigZero
% CTSmallZero (LatticeBasics/OrthAn.lean). The first paragraph is hpair
% (isPair_and_isPair_of_isTriple) and hr; the second is hth and hcusp (cusp_lift); the third
% is TowerB, hcR (cusp_re), EvZero, CTBigZero and NonZeroAn.1; the fourth is LeviEx, hct0
% (isCT_of_cusp), hsplit and h5. The case q = 0 is the same with TowerA, CTSmallZero,
% NonZeroAn.2 and hfa (exists_trace_zero_generator_of_isCMExt). EisBL and SiegelWeilStd
% enter through OrthIn (SWFull and EisContinues, used by cusp_lift) and NonZeroAn. What
% remains open here is in Lemma asm-theta-agree-end-orth-cusp-lift-inputs.

\subsection{Vanishing of differences}\label{subsection-asm-vanishing-differences}

\begin{lemma}\label{asm-theta-agree-end-square-exponent}
Let $n,d,q,q',Y\geq0$ be integers with $q\leq Y$, $q'\leq Y$, $2Y+1<n$ and
$2n\leq d$, with $\mathcal C_q$, $\delta_n$ and $P_{n,d}$ as in
Definition~\ref{asm-theta-agree-convergence-tori}(3) and (5) and the excess
$\Delta$ of Definition~\ref{asm-theta-agree-continuity-sums}(2). Then there is
$m_0\in\mathbf R$ such that
\[
\tfrac12P_{n,d}(s,t)-\delta_n(s)-m_0\,\Delta(s,t)\leq\tfrac12\delta_d(t)-\tfrac14\sum_{j=1}^{q'}t_j-\frac1{4(q+1)}\sum_{i=1}^qs_i
\]
for all $s\in\mathcal C_q$ and $t\in\mathcal C_{q'}$.
\end{lemma}
\leanref{Development/WeilBound/ThetaLifts/FrameEstimates/SquareIntegrableWeilRange.lean}{PositivityTheory.ThKEst.ExpSo}
\leanref{Development/WeilBound/ThetaLifts/FrameEstimates/InitialPairs/SquareIntegrableIntegralHolds.lean}{PositivityTheory.ThKEst.Sq.exponent_bound_orthogonal_case}
\leanref{Development/WeilBound/ThetaLifts/FrameEstimates/InitialPairs/SquareIntegrableIntegralHolds.lean}{PositivityTheory.ThKEst.Sq.major_slack}
\begin{proof}
We take $m_0=d+q+1$. Let $s\in\mathcal C_q$ and $t\in\mathcal C_{q'}$, and put
$\Delta=\Delta(s,t)$, which is at least $0$, the term for $i=0$ being $0$. Let
$u\in\mathcal C_q$ have $u_j=t_j$ for $j\leq\min(q,q')$ and $u_j=0$ for
$j>q'$. Then $\sigma_i(u)=\sigma_i(t)$ for $i\leq q$, so
$\sigma_i(s)\leq\sigma_i(u)+\Delta$ for $0\leq i\leq q$.

(i) $\sum_{i,j\leq q}\min(s_i,s_j)-\sum_{i,j\leq q}\min(s_i,u_j)\leq\sum_{j=1}^q(q-j)u_j+q\Delta$.
For each pair, $\min(s_i,s_j)-\min(s_i,u_j)$ is at most $s_i-u_j$ if
$u_j<s_i$, and at most $0$ otherwise. Fix $j$. As $s$ is nonincreasing, the $i$
with $u_j<s_i$ are those with $i\leq k$, for some $k\leq q$, and
\[
\sum_{i\leq k}(s_i-u_j)=\sigma_k(s)-ku_j\leq\sum_{i\leq k}(u_i-u_j)+\Delta\leq\sum_{i<j}u_i+\Delta,
\]
since $u_i-u_j\leq0$ for $i\geq j$ and $u_i-u_j\leq u_i$ for $i<j$, the $u_i$
being nonnegative. Summing over $j$, and as
$\sum_{j\leq q}\sum_{i<j}u_i=\sum_i(q-i)u_i$, gives (i).

(ii) $\sum_{i\leq q,j\leq q}\min(s_i,u_j)\leq\sum_{i\leq q,j\leq q'}\min(s_i,t_j)$:
the terms with $j\leq\min(q,q')$ agree, those with $q'<j\leq q$ vanish on the
left, and those with $q<j\leq q'$ are nonnegative on the right.

(iii) $\sum_is_i\leq\sum_ju_j+\Delta$, the case $i=q$.

(iv) Put $c=\frac d2-n+\frac1{4(q+1)}$; as $2n\leq d$, $0\leq c\leq d+1$. Then
\[
\sum_{j=1}^q(c+q-j)u_j\leq\sum_{j=1}^{q'}\Bigl(\frac{d-2j+1}2-\frac n2-\frac14\Bigr)t_j=\tfrac12\delta_d(t)-\Bigl(\frac n2+\frac14\Bigr)\sum_jt_j .
\]
For $j\leq\min(q,q')$ the coefficient on the left is at most the one on the
right, as this amounts to $q+\frac1{4(q+1)}-\frac14\leq\frac n2$, and
$n\geq2Y+2\geq2q+2$. For $q<j\leq q'$ the coefficient on the right is at least
$\frac n2-j+\frac14>0$, as $d\geq2n$ and $j\leq Y<\frac n2$. The terms with
$j>q'$ on the left vanish.

Now $\sum_{i,j}\min(s_i,s_j)=\sum_i(2i-1)s_i$, as $s$ is nonincreasing, so
$\delta_n(s)=n\sum_is_i-\sum_{i,j}\min(s_i,s_j)$; and
$P_{n,d}(s,t)=d\sum_is_i+n\sum_jt_j-2\sum_{i,j}\min(s_i,t_j)$, as in the proof
of Lemma~\ref{asm-theta-agree-convergence-cone}(1). Hence
\[
\tfrac12P_{n,d}(s,t)-\delta_n(s)=\Bigl(\frac d2-n\Bigr)\sum_is_i+\frac n2\sum_jt_j+\sum_{i,j}\min(s_i,s_j)-\sum_{i,j}\min(s_i,t_j).
\]
By (ii) and (i) the difference of the last two sums is at most
$\sum_j(q-j)u_j+q\Delta$, and by (iii)
$(\frac d2-n)\sum_is_i=c\sum_is_i-\frac1{4(q+1)}\sum_is_i\leq c\sum_ju_j+c\Delta-\frac1{4(q+1)}\sum_is_i$.
With (iv),
\[
\tfrac12P_{n,d}(s,t)-\delta_n(s)\leq\tfrac12\delta_d(t)-\tfrac14\sum_jt_j-\frac1{4(q+1)}\sum_is_i+(c+q)\Delta,
\]
and $(c+q)\Delta\leq m_0\Delta$.
\end{proof}
% Lean: exponent_bound_orthogonal_case (SquareIntegrableIntegralHolds.lean), with m = d + q + 1;
% u is ext t, (i) is major_slack (with exists_initial), (ii) is h2, (iv) is h4, and
% sum_min_self and posWt_eq give the two identities.

\begin{lemma}\label{asm-theta-agree-end-square-torus}
Let $(X,W,M,\chi_X,\chi_W,\omega)$ be a Weil pair for a theta base $B$ of a CM
field $E$ (Definition~\ref{asm-theta-agree-triple}(1)), with $X$ and $W$ of
ranks $n$ and $d$, put $D=[E:\mathbf Q]$, and let $a$ and $a'$ be Witt tori of
$J_X$ and $J_W$ of sizes $q$ and $q'$ with property (b) of
Definition~\ref{asm-theta-agree-convergence-tori}(3). Let $U$ and $V$ be
measurable functions on $\mathbf R^q$ and $\mathbf R^{q'}$ with
\[
\int_{\mathcal C_q}e^{U(s)-D\delta_n(s)}\,ds<\infty\quad\text{and}\quad
\int_{\mathcal C_{q'}}e^{2V(t)-D\delta_d(t)}\,dt<\infty,
\]
and let $\Phi$ be a Schwartz-Bruhat function on $\mathbf A^{n\times d}$ such
that for all compact sets $\Omega\subset U_X$ and $\Omega'\subset U_W$ there is
$A\geq0$ with
\[
\bigl|\theta_\omega(\Phi)\bigl(a(s)\kappa,a'(t)\kappa'\bigr)\bigr|\leq Ae^{U(s)+V(t)}
\]
for all $s\in\mathcal C_q$, $t\in\mathcal C_{q'}$, $\kappa\in\Omega$ and
$\kappa'\in\Omega'$. Then $|I(\Phi)|^2$ is integrable on $\mathcal D'$ for
every Haar measure $\mu'$ on $U_W$ and every fundamental domain $\mathcal D'$
for $U_W(F)$ in $U_W$ with respect to $\mu'$.
\end{lemma}
\leanref{Development/WeilBound/ThetaLifts/FrameEstimates/InitialPairs/SquareIntegrableOfLatticeTorus.lean}{PositivityTheory.ThKEst.Sq.square_integrable_of_torus_bound}
\begin{proof}
Let $\mu$ and $\mathcal D$ be the measure and the domain of $M$. Property (b)
of $a$ for $\mu$ and $\mathcal D$ gives $K<\infty$ and a compact set
$\Omega\subset U_X$, and property (b) of $a'$ for $\mu'$ and $\mathcal D'$
gives $K'<\infty$ and a compact set $\Omega'\subset U_W$. Let $A$ be as in the
hypothesis for $\Omega$ and $\Omega'$, and put
$c_U=\int_{\mathcal C_q}e^{U(s)-D\delta_n(s)}\,ds$.

\emph{Step 1.} For $t\in\mathcal C_{q'}$ and $\kappa'\in\Omega'$,
$|I(\Phi)(a'(t)\kappa')|\leq KAc_Ue^{V(t)}$. Put $h=a'(t)\kappa'$ and
$f(g)=|\theta_\omega(\Phi)(g,h)|$. The function $f$ is continuous, by
Lemma~\ref{asm-theta-agree-convergence-theta}(2), invariant under left
translation by $U_X(F)$, by Lemma~\ref{asm-theta-agree-anisotropic-automorphy},
and $f(a(s)\kappa)\leq Ae^{V(t)}e^{U(s)}$ for $s\in\mathcal C_q$ and
$\kappa\in\Omega$. Property (b) of $a$, with
$\beta(s)=Ae^{V(t)}e^{U(s)}$, gives
$\int_{\mathcal D}f\,d\mu\leq KAe^{V(t)}c_U$. And
$|I(\Phi)(h)|\leq\int_{\mathcal D}f\,d\mu$, also when $f$ is not integrable,
as $I(\Phi)(h)$ is then $0$.

\emph{Step 2.} The function $I(\Phi)$ is measurable, by Fubini's theorem, as
$\theta_\omega(\Phi)$ is continuous on $U_X\times U_W$ and $U_X$ is second
countable (Theorem~\ref{asm-theta-agree-convergence-siegel}(1)). It is
invariant under left translation by $U_W(F)$, by
Lemma~\ref{asm-theta-agree-anisotropic-automorphy}. By Step 1, property (b) of
$a'$ applies to $|I(\Phi)|^2$ and $\beta(t)=(KAc_U)^2e^{2V(t)}$, and gives
\[
\int_{\mathcal D'}|I(\Phi)|^2\,d\mu'\leq K'(KAc_U)^2\int_{\mathcal C_{q'}}e^{2V(t)-D\delta_d(t)}\,dt<\infty .
\]
\end{proof}

\begin{lemma}\label{asm-theta-agree-end-square}
Let $(B,G,\mathcal X,\omega,\omega')$ be a triple for a theta base $B$ of a CM
field $E$ (Definition~\ref{asm-theta-agree-triple}), with models $M$ and $M'$
and theta integrals $I$ and $I'$, let $\mu$ be a Haar measure on $U_W$, and let
$\mathcal D$ be a fundamental domain for $U_W(F)$ in $U_W$ with respect to
$\mu$. Then $|I(\Phi)|^2$ is integrable on $\mathcal D$ for every
$\Phi\in\operatorname{PD}(M)$, and $|I'(\Phi')|^2$ is integrable on $\mathcal D$
for every $\Phi'\in\operatorname{PD}(M')\cap\mathcal S^\circ(M')$.
\end{lemma}
\leanref{Development/AbelianVarietyTheory/Definitions/Divisors/CTSmallCone.lean}{GlobalTheory.ThK.SqInt}
\leanref{Development/WeilBound/ThetaLifts/FrameEstimates/InitialPairs/TriplePairHolds.lean}{PositivityTheory.ThKEst.Sq.sqInt_of_leaves}
\begin{proof}
Let $G=(p,q,V,V',W,M,M')$, of ranks $(a,a')$, and put $D=[E:\mathbf Q]$. By
Definition~\ref{asm-theta-agree-triple}(4), $(V,W,M,\beta,\beta/\alpha,\omega)$
is a Weil pair of signature $(a-p,p;a+a'-p-q,p+q)$ and
$(V',W,M',\alpha,\alpha/\beta,\omega')$ one of signature
$(q,a'-q;a+a'-p-q,p+q)$. By Lemma~\ref{asm-theta-agree-convergence-signature}
and Theorem~\ref{asm-theta-agree-convergence-siegel}(2), the matrices of $V$,
$V'$ and $W$ have Siegel Witt tori $\lambda_V$, $\lambda_{V'}$ and $\lambda_W$
of sizes $k_V\leq\min(a-p,p)$, $k_{V'}\leq\min(q,a'-q)$ and
$k_W\leq\min(a+a'-p-q,p+q)$. As $G$ is proper,
$2(p+q)+1<\lfloor(a-a')/2\rfloor$ and $2(p+q)+1<a'$; so $4p<a-a'$ and
$k_W\leq p+q<a'$.

For real numbers $c_1,\ldots,c_r$ whose partial sums $c_1+\cdots+c_k$,
$1\leq k\leq r$, are all negative,
$\int_{\mathcal C_r}\exp(\sum_kc_kt_k)\,dt=\prod_k|c_1+\cdots+c_k|^{-1}<\infty$,
by the substitution in the proof of
Lemma~\ref{asm-theta-agree-convergence-cone}(2). Write
$|t|_1=t_1+\cdots+t_r$ for $t\in\mathbf R^r$.

\emph{The first space.} Let $\Phi\in\operatorname{PD}(M)$. It is a
Schwartz-Bruhat function, as in the proof of
Lemma~\ref{asm-theta-agree-end-automorphic}, so $\Phi=\sum_if_i\otimes\phi_i$
with $f_i\in\mathcal S(F_\infty^{a\times(a+a')})$ and Schwartz-Bruhat functions
$\phi_i$ on $\mathbf A_f^{a\times(a+a')}$. Let $N$ be the integer $N_0$ of
Lemma~\ref{asm-theta-agree-convergence-lattice}(5) for $\lambda_V$ and
$\lambda_W$, and for each $i$ let $p_i$ and $C_i$ be given by
Lemma~\ref{asm-theta-agree-convergence-square} for $\phi_i$ and $N$. Let
$\Omega\subset U_V$ and $\Omega'\subset U_W$ be compact.
Lemma~\ref{asm-theta-agree-convergence-lattice}(3) gives $C'_i$ and $K_i$ with
$L_{C_i,N}(g\kappa,h\kappa')\leq K_iL_{C'_i,N}(g,h)$ for $\kappa\in\Omega$ and
$\kappa'\in\Omega'$, and Lemma~\ref{asm-theta-agree-convergence-lattice}(5) for
$C'_i$ gives $K'_i$ with
$L_{C'_i,N}(\lambda_V(s),\lambda_W(t))\leq K'_ie^{DP_{a,a+a'}(s,t)}$. Put
$U(s)=\frac{D(a+a')}2|s|_1$ and $V(t)=\frac{Da}2|t|_1$; by
Lemma~\ref{asm-theta-agree-convergence-cone}(1),
$DP_{a,a+a'}(s,t)\leq2(U(s)+V(t))$. So
\[
\bigl|\theta_\omega(f_i\otimes\phi_i)\bigl(\lambda_V(s)\kappa,\lambda_W(t)\kappa'\bigr)\bigr|^2\leq p_i(f_i)^2K_iK'_ie^{2(U(s)+V(t))},
\]
and, the theta function being additive
(Lemma~\ref{asm-theta-agree-convergence-theta}(2)), $\Phi$ satisfies the bound
of Lemma~\ref{asm-theta-agree-end-square-torus} with
$A=\sum_ip_i(f_i)(K_iK'_i)^{1/2}$. The integral of
$e^{U(s)-D\delta_a(s)}$ over $\mathcal C_{k_V}$ is finite by
Lemma~\ref{asm-theta-agree-convergence-cone}(2), with $n=a$ and $d=a+a'$, as
$k_V\leq p$ and $a+a'+2p<2a$. And
$2V(t)-D\delta_{a+a'}(t)=\sum_jD(2j-1-a')t_j$, whose partial sums
$Dk(k-a')$ are negative for $1\leq k\leq k_W<a'$; so its exponential has a
finite integral over $\mathcal C_{k_W}$. By
Lemma~\ref{asm-theta-agree-end-square-torus} for $\lambda_V$ and $\lambda_W$,
$|I(\Phi)|^2$ is integrable on $\mathcal D$.

\emph{The second space.} Let
$\Phi'\in\operatorname{PD}(M')\cap\mathcal S^\circ(M')$. With $Y=p+q$,
$k_{V'}\leq Y$, $k_W\leq Y$, $2Y+1<a'$ and $2a'\leq a+a'$; let $m_0$ be given
by Lemma~\ref{asm-theta-agree-end-square-exponent} for
$(a',a+a',k_{V'},k_W,Y)$, let $\ell\geq0$ be an integer with $\ell\geq2Dm_0$,
and let $N$ be
given by Lemma~\ref{asm-theta-agree-continuity-lattice} for $\lambda_{V'}$,
$\lambda_W$ and $\ell$. By Lemma~\ref{asm-theta-agree-continuity-pieces}(1),
$\Phi'=\sum_if_i\otimes\phi_i$ with full-rank finite parts $\phi_i$ for $M'$;
for each $i$ let $p_i$ and $C_i$ be given by
Lemma~\ref{asm-theta-agree-continuity-square} for $\phi_i$ and $N$. Let
$\Omega\subset U_{V'}$ and $\Omega'\subset U_W$ be compact. As in the first
case, Lemma~\ref{asm-theta-agree-continuity-full-sums}(2) and
Lemma~\ref{asm-theta-agree-continuity-lattice} give $K_i$ and $K'_i$ with
\[
\bigl|\theta_{\omega'}(f_i\otimes\phi_i)\bigl(\lambda_{V'}(s)\kappa,\lambda_W(t)\kappa'\bigr)\bigr|^2\leq p_i(f_i)^2K_iK'_i\exp\bigl(DP_{a',a+a'}(s,t)-\ell\Delta(s,t)\bigr).
\]
Put $U(s)=D\delta_{a'}(s)-\frac D{4(k_{V'}+1)}|s|_1$ and
$V(t)=\frac D2\delta_{a+a'}(t)-\frac D4|t|_1$. As $\Delta(s,t)\geq0$ and
$\ell\geq2Dm_0$, Lemma~\ref{asm-theta-agree-end-square-exponent}, multiplied by
$2D$, gives $DP_{a',a+a'}(s,t)-\ell\Delta(s,t)\leq2(U(s)+V(t))$. So $\Phi'$
satisfies the bound of Lemma~\ref{asm-theta-agree-end-square-torus}, as in the
first case. The exponents $U(s)-D\delta_{a'}(s)=-\frac D{4(k_{V'}+1)}|s|_1$
and $2V(t)-D\delta_{a+a'}(t)=-\frac D2|t|_1$ have negative partial sums, so
both integrals are finite. By Lemma~\ref{asm-theta-agree-end-square-torus} for
$\lambda_{V'}$ and $\lambda_W$, $|I'(\Phi')|^2$ is integrable on $\mathcal D$.
\end{proof}
% Lean: Sq.sqInt_of_leaves = sqInt_of triplePair_holds (square_integrable_of_lattice_torus ...)
% (square_integrable_orthogonal_case ...) (SquareIntegrableOfLatticeTorus.lean). triplePair_holds
% gives the two Weil pairs; the first space is square_integrable_of_lattice_torus (SiegelAll,
% EpsOf, SqAll, LatCpt, LatTor, square_integrable_integral_holds), the second
% square_integrable_orthogonal_case (SqSo, SoFin, LatFull, exponent_bound_orthogonal_case,
% exp_cancel); both bound each pure tensor by torBd_pure, add them by torBd_sum and conclude by
% square_integrable_of_torus_bound. The cone integrals are cone_exponent_bound_holds and coneAff.

\begin{lemma}\label{asm-theta-agree-end-perp-zero}
Let $(B,G,\mathcal X,\omega,\omega')$ be a triple for a theta base $B$ of a CM
field $E$ (Definition~\ref{asm-theta-agree-triple}), with models $M$ and $M'$
and theta integrals $I$ and $I'$, and let $f$ be in its set of differences
(Definition~\ref{asm-theta-agree-end-notions}(5)). If $I(\Phi)$ is orthogonal
to $f$ for every $\Phi\in\operatorname{PD}(M)$ and $I'(\Phi')$ is orthogonal to
$f$ for every $\Phi'\in\operatorname{PD}(M')\cap\mathcal S^\circ(M')$
(Definition~\ref{asm-theta-agree-end-notions}(4)), then $f=0$.
\end{lemma}
\leanref{Development/WeilBound/ThetaLifts/EigenvalueEndgame/PlaneRemovalConfiguration.lean}{PositivityTheory.ThKEnd.PerpZero}
\leanref{Development/WeilBound/ThetaLifts/EigenvalueEndgame/UnipotentGroups/PerpZeroOf.lean}{PositivityTheory.ThKEnd.Grp.perpZero_of}
\leanref{Development/WeilBound/ThetaLifts/CartanFactorisation/IntegrableOnRight.lean}{PositivityTheory.ThKF.L2.eq_zero}
\begin{proof}
Write $f=I(\Phi)-I'(\Phi')$ with $\Phi\in\operatorname{PD}(M)$ and
$\Phi'\in\operatorname{PD}(M')\cap\mathcal S^\circ(M')$. By
Lemma~\ref{asm-theta-agree-end-automorphic}, $I(\Phi)$, $I'(\Phi')$ and $f$ are
continuous and invariant under left translation by $U_W(F)$. The matrix of $W$
is nondegenerate $\epsilon$-Hermitian
(Lemma~\ref{asm-theta-agree-convergence-signature}), so by
Theorem~\ref{asm-theta-agree-definite-volume} there are a Haar measure $\mu$ on
$U_W$ and a fundamental domain $\mathcal D$ for $U_W(F)$ with respect to $\mu$.

By Lemma~\ref{asm-theta-agree-end-square}, $|I(\Phi)|^2$ and $|I'(\Phi')|^2$ are
integrable on $\mathcal D$. As $|f|^2\leq2|I(\Phi)|^2+2|I'(\Phi')|^2$ and
$|\bar f\,g|\leq|f|^2+|g|^2$, the functions $|f|^2$, $\bar fI(\Phi)$ and
$\bar fI'(\Phi')$ are integrable on $\mathcal D$. By the hypotheses,
$\int_{\mathcal D}I(\Phi)\bar f\,d\mu=0$ and
$\int_{\mathcal D}I'(\Phi')\bar f\,d\mu=0$, so
\[
\int_{\mathcal D}|f|^2\,d\mu=\int_{\mathcal D}\bar f\,\bigl(I(\Phi)-I'(\Phi')\bigr)\,d\mu=0.
\]
So the set $S=\{h:f(h)\neq0\}$ meets $\mathcal D$ in a null set. It is
invariant under $U_W(F)$, and $\mu$ is invariant under left translations, so
$\mu(S)\leq\sum_\gamma\mu(S\cap\gamma\mathcal D)=\sum_\gamma\mu(\gamma(S\cap\mathcal D))=0$,
the sums over $\gamma\in U_W(F)$, which is countable by
Theorem~\ref{asm-theta-agree-definite-volume}. As $f$ is continuous, $S$ is
open, and a nonempty open set has positive Haar measure. So $S$ is empty and
$f=0$.
\end{proof}
% Lean: perpZero_of SqInt IAut: unimodularity_holds gives mu and D, L2.integrableOn_mul the
% integrability of the products, and L2.eq_zero the last paragraph.

\begin{lemma}\label{asm-theta-agree-end-differences}
Let $(B,G,\mathcal X,\omega,\omega',R,R_f)$ be a matched quadruple for a theta
base $B$ of a CM field $E$ (Definition~\ref{asm-theta-agree-matched}), let $H'$
be an adapted plane in $W$ over a lattice space $W_1$, and let $c\in\mathbf C$.
Let $\Phi'\in\mathcal S^\circ(M')$ and let $\Phi$ correspond to $\Phi'$. Then
$f=I(\Phi)-c\,I'(\Phi')$ lies in the set of differences of the triple
(Definition~\ref{asm-theta-agree-end-notions}(5)), and there are pairs
$(f_\alpha,f_\beta)$ and $(f'_\alpha,f'_\beta)$ that split $I(\Phi)$ and
$I'(\Phi')$ along $H'$ (Definition~\ref{asm-theta-agree-planes}(3)) such that
$(f_\alpha-c\,f'_\alpha,f_\beta-c\,f'_\beta)$ splits $f$ along $H'$.
\end{lemma}
\leanref{Development/WeilBound/ThetaLifts/EigenvalueEndgame/TowerCodimensionTwo.lean}{PositivityTheory.ThKEnd.diff_data}
\begin{proof}
By Lemma~\ref{asm-theta-agree-partners}(3), $\Phi\in\operatorname{PD}(M)$ and
$\Phi'\in\operatorname{PD}(M')$. The function $c\Phi'$ lies in the subspace
$\operatorname{PD}(M')\cap\mathcal S^\circ(M')$, and $I'(c\Phi')=c\,I'(\Phi')$,
since $\theta_{\omega'}(c\Phi')=c\,\theta_{\omega'}(\Phi')$. So
$f=I(\Phi)-I'(c\Phi')$ lies in the set of differences.

By Lemma~\ref{asm-theta-agree-split} there are pairs $(f_\alpha,f_\beta)$ and
$(f'_\alpha,f'_\beta)$ that split $I(\Phi)$ and $I'(\Phi')$ along $H'$. By
Lemma~\ref{asm-theta-agree-end-automorphic}, $I(\Phi)$ and $I'(\Phi')$ are
automorphic. The space $W$ is skew-Hermitian of signature $(a+a'-p-q,p+q)$, so
$H'$ is a plane by Lemma~\ref{asm-theta-agree-end-iwasawa}. By
Lemma~\ref{asm-theta-agree-end-constant-terms}(1), the constant term of $f$
along $H'$ is
$(f_\alpha+f_\beta)-c(f'_\alpha+f'_\beta)=(f_\alpha-c\,f'_\alpha)+(f_\beta-c\,f'_\beta)$.
If $T$ and $T'$ have the law $(\chi,z)$, so does $T-c\,T'$. So
$f_\alpha-c\,f'_\alpha$ has the law $(\alpha,a'/2)$ and
$f_\beta-c\,f'_\beta$ the law $(\beta,a/2)$, and the pair splits $f$.
\end{proof}
% Lean: diff_data uses MatchEx (third part), SplitEx, IAut, Iwa, ctBasic.1, hasLaw_sub and
% I_smul.

\begin{lemma}\label{asm-theta-agree-end-compare}
Let $E$ be a CM field with a theta base $B$
(Definition~\ref{asm-theta-agree-models}(1)), let $W$ be a skew-Hermitian
lattice space of some signature, let $H'$ be an adapted plane in $W$ over a
lattice space $W_1$ at some positions, let $\chi$ be a homomorphism from
$\mathbf A_E^\times$ to the complex numbers of absolute value $1$, let
$z\in\mathbf R$, and let $T$ be a function on $U_W$ with the law $(\chi,z)$
(Definition~\ref{asm-theta-agree-planes}(2)). Suppose that $T(uh)=T(h)$ for all
$u\in N_{H'}$ and $h\in U_W$, and that $T(lk)=0$ for every $h_1\in U_{W_1}$,
every Levi element $l$ for $(1,h_1)$ and every $k\in K^\theta_W$
(Definition~\ref{asm-theta-agree-end-notions}(2)). Then $T=0$.
\end{lemma}
\leanref{Development/WeilBound/ThetaLifts/EigenvalueEndgame/TowerCodimensionTwo.lean}{PositivityTheory.ThKEnd.comp_zero}
\begin{proof}
Let $g\in U_W$. By Lemma~\ref{asm-theta-agree-end-iwasawa}, $g=umlk$ with
$u\in N_{H'}$, $m$ a Levi element for $(t,1)$, $l$ a Levi element for $(1,h_1)$
and $k\in K^\theta_W$. By the invariance, the law and the vanishing,
$T(g)=T(mlk)=\chi(t)|t|^zT(lk)=0$.
\end{proof}

\begin{lemma}\label{asm-theta-agree-end-zero}
Assume the meromorphic continuation of Siegel Eisenstein series in the form of
Theorem~\ref{asm-theta-agree-eisenstein} and the Siegel-Weil formula in the form
of Theorem~\ref{asm-theta-agree-siegel-weil}. Let
$(B,G,\mathcal X,\omega,\omega')$ be a triple for a theta base $B$ of a CM field
$E$ (Definition~\ref{asm-theta-agree-triple}) with $p\geq1$ or $q\geq1$, and let
$H'$ be a plane in $W$ over a lattice space $W_1$
(Definition~\ref{asm-theta-agree-end-notions}(1)). If $f$ lies in the set of
differences of the triple and $0$ is the constant term of $f$ along $H'$, then
$f=0$.
\end{lemma}
\leanref{Development/WeilBound/ThetaLifts/EigenvalueEndgame/TowerCodimensionTwo.lean}{PositivityTheory.ThKEnd.zero_of_line}
\begin{proof}
By Lemma~\ref{asm-theta-agree-end-cusp}, $f$ is cuspidal. Let
$\Phi\in\operatorname{PD}(M)$. If $p\geq1$, then $I(\Phi)$ is orthogonal to $f$
by Lemma~\ref{asm-theta-agree-end-orth-iso}; if $p=0$, then $q\geq1$, and
$I(\Phi)$ is orthogonal to $f$ by Lemma~\ref{asm-theta-agree-end-orth-an}. Let
$\Phi'\in\operatorname{PD}(M')\cap\mathcal S^\circ(M')$. If $q\geq1$, then
$I'(\Phi')$ is orthogonal to $f$ by Lemma~\ref{asm-theta-agree-end-orth-iso};
if $q=0$, then $p\geq1$, and $I'(\Phi')$ is orthogonal to $f$ by
Lemma~\ref{asm-theta-agree-end-orth-an}. By
Lemma~\ref{asm-theta-agree-end-perp-zero}, $f=0$.
\end{proof}

\begin{lemma}\label{asm-theta-agree-end-pure}
Assume the meromorphic continuation of Siegel Eisenstein series in the form of
Theorem~\ref{asm-theta-agree-eisenstein} and the Siegel-Weil formula in the form
of Theorem~\ref{asm-theta-agree-siegel-weil}. Let
$(B,G,\mathcal X,\omega,\omega')$ be a triple of ranks $(a,a')$ for a theta base
$B$ of a CM field $E$ (Definition~\ref{asm-theta-agree-triple}), with
$\mathcal X=(\alpha,\beta,S_0,Q)$ and $q\geq1$. If $f$ lies in the set of
differences of the triple (Definition~\ref{asm-theta-agree-end-notions}(5)) and
is pure of law $(\alpha,a'/2)$ (Definition~\ref{asm-theta-agree-end-notions}(3)),
then $f=0$.
\end{lemma}
\leanref{Development/WeilBound/ThetaLifts/EigenvalueEndgame/TowerCodimensionTwo.lean}{PositivityTheory.ThKEnd.sumB}
\leanref{Development/WeilBound/ThetaLifts/EigenvalueEndgame/TowerCodimensionTwo.lean}{PositivityTheory.ThKEnd.baseB}
\leanref{Development/WeilBound/ThetaLifts/EigenvalueEndgame/TowerCodimensionTwo.lean}{PositivityTheory.ThKEnd.stepB}
\leanref{Development/WeilBound/ThetaLifts/EigenvalueEndgame/TowerCodimensionTwo.lean}{PositivityTheory.ThKEnd.split_of_pure}
\begin{proof}
\emph{Translations and re-chosen data.} (a) If a function $F_0$ on $U_W$ is
pure of law $(\chi,z)$, through a plane $H''$ and a constant term $c$, and
$h_0\in U_W$, then $h\mapsto F_0(hh_0)$ is pure of law $(\chi,z)$ through $H''$
and the constant term $h\mapsto c(hh_0)$: the integrals defining the constant
term and the identities defining the law are carried by right translation.
(b) For a re-chosen datum $G^\circ=G[V^\circ,V'^\circ,W^\circ]$, the space of
$G^\circ$ of rank $a+a'$ has the matrix of $W$, so it has the same unitary group
and the same planes as $W$; a function on $U_W$ is pure for one if and only if
it is pure for the other.

The proof is by induction on $p$, for all triples with the theta base $B$, the
character datum $\mathcal X$ and $q\geq1$.

\emph{The case $p=0$.} As $G$ is proper, $2(p+q)+1<a'$, so $a'\geq2$; write
$a'=b+2$. By Lemma~\ref{asm-theta-agree-tower-beta} there are lattice spaces
$V^\circ$, $V'^\circ$, $W^\circ$ such that, with
$G^\circ=G[V^\circ,V'^\circ,W^\circ]$, $(B,G^\circ,\mathcal X,\omega,\omega')$ is
a triple, a triple datum $G_1$ of ranks $(a,b)$, adapted planes $H$ and $H'$
such that $(G^\circ,G_1,H,H')$ is a $\beta$-configuration, and representations
$\omega_1$, $\omega'_1$ such that $(B,G_1,\mathcal X,\omega_1,\omega'_1)$ is a
triple. By Lemma~\ref{asm-theta-agree-end-rechoose} there is $h_0\in U_W$ such
that $f^\circ(h)=f(hh_0)$ lies in the set of differences of
$(B,G^\circ,\mathcal X,\omega,\omega')$; by (a) and (b), $f^\circ$ is pure of
law $(\alpha,a'/2)$. Let $\tilde W$ be the space of $G^\circ$ of rank $a+a'$. The
plane $H'$ is adapted in $\tilde W$ over the space $W_1$ of $G_1$ of rank $a+b$,
which is
skew-Hermitian of some signature, as $(B,G_1,\mathcal X,\omega_1,\omega'_1)$ is
a triple. By Lemma~\ref{asm-theta-agree-end-iwasawa}, $H'$ is a plane; by
Lemma~\ref{asm-theta-agree-end-automorphic}, $f^\circ$ is automorphic; and by
Lemma~\ref{asm-theta-agree-end-move} there is a function $c$ that is the
constant term of $f^\circ$ along $H'$ and has the law $(\alpha,a'/2)$. As $0$
has the law $(\beta,a/2)$, the pair $(c,0)$ splits $f^\circ$ along $H'$.
Lemma~\ref{asm-theta-agree-end-alpha-zero}, for the triple
$(B,G^\circ,\mathcal X,\omega,\omega')$ with $p=0$ and $q\geq1$, gives $c=0$.
So $0$ is the constant term of $f^\circ$ along $H'$, and
Lemma~\ref{asm-theta-agree-end-zero} gives $f^\circ=0$. Hence $f=0$.

\emph{The step to $p\geq1$.} As $G$ is proper, $a>a'>2(p+q)+1\geq1$, so
$a\geq2$; write $a=b+2$. By Lemma~\ref{asm-theta-agree-end-tower} there are
lattice spaces $V^\circ$, $V'^\circ$, $W^\circ$ such that, with
$G^\circ=G[V^\circ,V'^\circ,W^\circ]$, $(B,G^\circ,\mathcal X,\omega,\omega')$ is
a triple, a triple datum $G_1$ of ranks $(b,a')$, adapted planes $H$ and $H'$
such that $(G^\circ,G_1,H,H')$ is a plane-removal configuration, and
representations $\omega_1$, $\omega'_1$ such that
$(B,G_1,\mathcal X,\omega_1,\omega'_1)$ is a triple; then $p_1=p-1$ and
$q_1=q\geq1$. Let $\Sigma_1$ be the set of differences of the smaller triple
and $\tilde W$ the space of $G^\circ$ of rank $a+a'$. As in the case $p=0$,
Lemma~\ref{asm-theta-agree-end-rechoose} gives $h_0$ such that
$f^\circ(h)=f(hh_0)$ lies in the set of differences of
$(B,G^\circ,\mathcal X,\omega,\omega')$ and is pure of law $(\alpha,a'/2)$, and
Lemmas~\ref{asm-theta-agree-end-iwasawa}, \ref{asm-theta-agree-end-automorphic}
and~\ref{asm-theta-agree-end-move} give a function $c$, the constant term of
$f^\circ$ along $H'$, with the law $(\alpha,a'/2)$, so that $(c,0)$ splits
$f^\circ$ along $H'$.

Then $c=0$, by Lemma~\ref{asm-theta-agree-end-compare} for $\tilde W$ and $H'$.
Its law is $(\alpha,a'/2)$. As $a'<a$, the exponents $a'/2$ and $a/2$ differ,
and Lemma~\ref{asm-theta-agree-end-constant-terms}(2), applied to $f^\circ$ and
the parts $c$ and $0$ of its constant term, shows that $c$ is invariant under
$N_{H'}$. Let $k\in K^\theta_{\tilde W}$, $h_1\in U_{W_1}$, and $l$ a Levi
element for $(1,h_1)$. By Lemma~\ref{asm-theta-agree-end-alpha-low}, for every
$k'\in K^\theta_{\tilde W}$ some element of $\Sigma_1$ is the lower part of $c$
at $k'$; let $f_1\in\Sigma_1$ be the lower part of $c$ at $k$. By
Lemma~\ref{asm-theta-agree-end-two-step}, $f_1$ is pure of law
$(\alpha,a'/2)$. By the induction hypothesis for
$(B,G_1,\mathcal X,\omega_1,\omega'_1)$, which has $p_1=p-1$ and $q_1\geq1$,
$f_1=0$. So $c(lk)=f_1(h_1)=0$, and Lemma~\ref{asm-theta-agree-end-compare}
gives $c=0$.

So $0$ is the constant term of $f^\circ$ along $H'$, and
Lemma~\ref{asm-theta-agree-end-zero}, with $q\geq1$, gives $f^\circ=0$. Hence
$f=0$.
\end{proof}
% Lean: sumB (induction on p); the case p = 0 is TowerB, SigRe and baseB (split_of_pure, which
% applies LineMove, then AlphaZero and zero_of_line); the step is TowerCodimensionTwo, SigRe and
% stepB (split_of_pure, AlphaLow, TwoStep, the induction hypothesis and comp_zero with
% ctBasic.2, then zero_of_line). pure_right and pure_reS are (a) and (b); zero_of_right is the
% last step from f^circ to f.

\begin{lemma}\label{asm-theta-agree-end-alpha}
Assume the meromorphic continuation of Siegel Eisenstein series in the form of
Theorem~\ref{asm-theta-agree-eisenstein} and the Siegel-Weil formula in the
form of Theorem~\ref{asm-theta-agree-siegel-weil}. Let
$(B,G,\mathcal X,\omega,\omega',R,R_f)$ be a matched quadruple for a theta base
$B$ of a CM field $E$ (Definition~\ref{asm-theta-agree-matched}) with $q=0$ and
$p\geq1$, let $H'$ be an adapted plane in $W$ over a lattice space $W_1$, and
let $c\in\mathbf C$. Suppose that $f_\alpha(m)=c\,f'_\alpha(m)$ whenever
$\Phi'\in\mathcal S^\circ(M')$, $\Phi$ and $\Phi'$ correspond,
$(f_\alpha,f_\beta)$ and $(f'_\alpha,f'_\beta)$ split $I(\Phi)$ and $I'(\Phi')$
along $H'$, and $m$ is a Levi element for $(1,h_1)$ for some $h_1\in U_{W_1}$.
Then $c$ is a polynomial constant of the quadruple.
\end{lemma}
\leanref{Development/WeilBound/ThetaLifts/ConeHypotheses/PolDense.lean}{PositivityTheory.ThKG.EndA'}
\leanref{Development/AbelianVarietyTheory/Definitions/Divisors/CTSmallCone.lean}{GlobalTheory.ThK.EndA}
\begin{proof}
Let $G$ have ranks $(a,a')$ and $\mathcal X=(\alpha,\beta,S_0,Q)$, and let
$\Phi'\in\mathcal S^\circ(M')$ and $\Phi$ correspond. By
Lemma~\ref{asm-theta-agree-end-differences}, $f=I(\Phi)-c\,I'(\Phi')$ lies in
the set of differences of the triple, and there are pairs $(f_\alpha,f_\beta)$
and $(f'_\alpha,f'_\beta)$ that split $I(\Phi)$ and $I'(\Phi')$ along $H'$ such
that $(f_\alpha-c\,f'_\alpha,f_\beta-c\,f'_\beta)$ splits $f$. By
Lemma~\ref{asm-theta-agree-end-iwasawa}, $H'$ is a plane in $W$ over $W_1$, and
by Lemma~\ref{asm-theta-agree-end-automorphic}, $f$ is automorphic, as a
difference of automorphic functions.

\emph{Step 1: $f_\alpha=c\,f'_\alpha$.} Put $T=f_\alpha-c\,f'_\alpha$; it has the
law $(\alpha,a'/2)$. As $a'<a$, the exponents $a'/2$ and $a/2$ differ, and
Lemma~\ref{asm-theta-agree-end-constant-terms}(2), applied to $f$ and the parts
$T$ and $f_\beta-c\,f'_\beta$ of its constant term, gives $T(uh)=T(h)$ for
$u\in N_{H'}$ and $h\in U_W$. Let $h_1\in U_{W_1}$, let $l$ be a Levi element
for $(1,h_1)$, and let $k\in K^\theta_W$. By
Lemma~\ref{asm-theta-agree-end-stable}, $\omega'(1,k)\Phi'\in\mathcal S^\circ(M')$
and it corresponds to $\omega(1,k)\Phi$. As $\omega$ is a representation,
$\theta_\omega(\omega(1,k)\Phi)(g,h)=\theta_\omega(\Phi)(g,hk)$, so
$I(\omega(1,k)\Phi)(h)=I(\Phi)(hk)$; and right translation by $k$ carries a
constant term along $H'$ to the constant term of the translated function and
preserves laws. So $h\mapsto(f_\alpha(hk),f_\beta(hk))$ splits
$I(\omega(1,k)\Phi)$ along $H'$, and likewise
$h\mapsto(f'_\alpha(hk),f'_\beta(hk))$ splits $I'(\omega'(1,k)\Phi')$. The
hypothesis on $c$, for these data and the Levi element $l$, gives
$f_\alpha(lk)=c\,f'_\alpha(lk)$, that is, $T(lk)=0$. By
Lemma~\ref{asm-theta-agree-end-compare}, $T=0$.

\emph{Step 2: $f_\beta=c\,f'_\beta$.} As $q=0$ and $p\geq1$,
Lemma~\ref{asm-theta-agree-end-beta-zero}, applied to $f$ and its splitting
$(f_\alpha-c\,f'_\alpha,f_\beta-c\,f'_\beta)$, gives $f_\beta-c\,f'_\beta=0$.

\emph{Step 3: conclusion.} By Steps 1 and 2, $0$ is the constant term of $f$
along $H'$. As $p\geq1$, Lemma~\ref{asm-theta-agree-end-zero} gives $f=0$, that
is, $I(\Phi)=c\,I'(\Phi')$. So $c$ is a polynomial constant.
\end{proof}
% Lean: endAB_sph proves EndA as ThKEnd.endA_of Leaves SplitEx MatchEx (TowerCodimensionTwo.lean),
% the fifteen fields of Leaves being Lemmas asm-theta-agree-end-iwasawa, -constant-terms,
% -move, -automorphic, -stable, -rechoose, -tower, -alpha-zero, -beta-zero, -alpha-low,
% -two-step, -cusp, -orth-iso, -orth-an and -perp-zero, SplitEx Lemma asm-theta-agree-split
% and MatchEx Lemma asm-theta-agree-partners; ThKG.endA'_of drops Below. The setup is
% diff_data; Step 1 is h0 (comp_zero with ctBasic.2, sig_isAut, half_ne, matchStab,
% split_trans and I_trans), Step 2 is h0' (betaZero), Step 3 is zero_of_line. The open steps
% left are in the leaves: about 16,000 lines in all beyond the lemmas stated elsewhere.

\begin{lemma}\label{asm-theta-agree-end-beta}
Assume the meromorphic continuation of Siegel Eisenstein series in the form of
Theorem~\ref{asm-theta-agree-eisenstein} and the Siegel-Weil formula in the
form of Theorem~\ref{asm-theta-agree-siegel-weil}. Let
$(B,G,\mathcal X,\omega,\omega',R,R_f)$ be a matched quadruple for a theta base
$B$ of a CM field $E$ (Definition~\ref{asm-theta-agree-matched}) with $q\geq1$,
let $H'$ be an adapted plane in $W$ over a lattice space $W_1$, and let
$c\in\mathbf C$. Suppose that $f_\beta(m)=c\,f'_\beta(m)$ whenever
$\Phi'\in\mathcal S^\circ(M')$, $\Phi$ and $\Phi'$ correspond,
$(f_\alpha,f_\beta)$ and $(f'_\alpha,f'_\beta)$ split $I(\Phi)$ and $I'(\Phi')$
along $H'$, and $m$ is a Levi element for $(1,h_1)$ for some $h_1\in U_{W_1}$.
Then $c$ is a polynomial constant of the quadruple.
\end{lemma}
\leanref{Development/WeilBound/ThetaLifts/ConeHypotheses/PolDense.lean}{PositivityTheory.ThKG.EndB'}
\leanref{Development/AbelianVarietyTheory/Definitions/Divisors/CTSmallCone.lean}{GlobalTheory.ThK.EndB}
\begin{proof}
Let $G$ have ranks $(a,a')$ and $\mathcal X=(\alpha,\beta,S_0,Q)$, and let
$\Phi'\in\mathcal S^\circ(M')$ and $\Phi$ correspond. As in the proof of
Lemma~\ref{asm-theta-agree-end-alpha},
Lemmas~\ref{asm-theta-agree-end-differences},
\ref{asm-theta-agree-end-iwasawa} and~\ref{asm-theta-agree-end-automorphic}
show that $f=I(\Phi)-c\,I'(\Phi')$ is an automorphic function in the set of
differences of the triple, that $H'$ is a plane, and that there are pairs
$(f_\alpha,f_\beta)$ and $(f'_\alpha,f'_\beta)$ that split $I(\Phi)$ and
$I'(\Phi')$ along $H'$ such that $(f_\alpha-c\,f'_\alpha,f_\beta-c\,f'_\beta)$
splits $f$.

\emph{Step 1: $f_\beta=c\,f'_\beta$.} Put $T=f_\beta-c\,f'_\beta$; it has the
law $(\beta,a/2)$. Lemma~\ref{asm-theta-agree-end-constant-terms}(2), applied to
$f$ and the parts $f_\alpha-c\,f'_\alpha$ and $T$ of its constant term, with
the exponents $a'/2\neq a/2$, gives $T(uh)=T(h)$ for $u\in N_{H'}$ and
$h\in U_W$. Let $h_1\in U_{W_1}$, let $l$ be a Levi element for $(1,h_1)$, and
let $k\in K^\theta_W$. As in Step 1 of the proof of
Lemma~\ref{asm-theta-agree-end-alpha}, Lemma~\ref{asm-theta-agree-end-stable}
gives corresponding data $\omega(1,k)\Phi$ and
$\omega'(1,k)\Phi'\in\mathcal S^\circ(M')$, whose theta integrals are split along
$H'$ by the right translates by $k$ of the two pairs. The hypothesis on $c$
gives $f_\beta(lk)=c\,f'_\beta(lk)$, that is, $T(lk)=0$. By
Lemma~\ref{asm-theta-agree-end-compare}, $T=0$.

\emph{Step 2: conclusion.} By Step 1, the constant term of $f$ along $H'$ is
$f_\alpha-c\,f'_\alpha$, which has the law $(\alpha,a'/2)$. So $f$ is pure of
law $(\alpha,a'/2)$ (Definition~\ref{asm-theta-agree-end-notions}(3)). As
$q\geq1$, Lemma~\ref{asm-theta-agree-end-pure} gives $f=0$, that is,
$I(\Phi)=c\,I'(\Phi')$. So $c$ is a polynomial constant.
\end{proof}
% Lean: endAB_sph proves EndB as ThKEnd.endB_of Leaves TowerB SplitEx MatchEx
% (TowerCodimensionTwo.lean), with TowerB Lemma asm-theta-agree-tower-beta; ThKG.endB'_of drops
% Below. The setup is diff_data; Step 1 is h0 (comp_zero, ctBasic.2, matchStab, split_trans);
% Step 2 is sumB (Lemma asm-theta-agree-end-pure) applied to f, pure through H' (hpl).

\subsection{The theta hypotheses and all characters}\label{subsection-asm-theta-hypotheses}

\begin{lemma}\label{asm-theta-agree-constants}
Let $(B,G,\mathcal X,\omega,\omega',R,R_f)$ be a matched quadruple for a theta
base $B$ of a CM field $E$ (Definition~\ref{asm-theta-agree-matched}). Then a
complex number is a constant of it if and only if it is a polynomial constant
of it.
\end{lemma}
\leanref{Development/WeilBound/ThetaLifts/ConeHypotheses/CtSmallCone.lean}{PositivityTheory.ThKG.Leaves.polConst_of}
\leanref{Development/WeilBound/ThetaLifts/ConeHypotheses/CtSmallCone.lean}{PositivityTheory.ThKG.Leaves.isConst_of}
\begin{proof}
Let $c$ be a constant, and let $\Phi'\in\mathcal S^\circ(M')$ and $\Phi$
correspond. By Lemma~\ref{asm-theta-agree-partners}(3) they match, so
$I(\Phi)=c\,I'(\Phi')$.

Let $c$ be a polynomial constant, and let $\Phi$ and $\Phi'$ match, through
$\Phi=\sum_if_i\otimes g_i$, $\Phi'=\sum_jf'_j\otimes g'_j$ and sequences
$s_{k,i}\to f_i$ and $s'_{k,j}\to f'_j$. Put $\Psi_k=\sum_is_{k,i}\otimes g_i$
and $\Psi'_k=\sum_js'_{k,j}\otimes g'_j$, and let $h\in U_W$. For each $k$,
$\Psi'_k\in\mathcal S^\circ(M')$ corresponds to $\Psi_k$, so
$I(\Psi_k)(h)=c\,I'(\Psi'_k)(h)$. Lemma~\ref{asm-theta-agree-continuity}, for
the Weil pair of $M'$, gives $I'(\Psi'_k)(h)\to I'(\Phi')(h)$, since
$\Phi'\in\mathcal S^\circ(M')$. The Weil pair of $M$ has signature
$(a-p,p;a+a'-p-q,p+q)$, and $a+a'+2p<2a$ because
$2(p+q)+1<\lfloor(a-a')/2\rfloor$. So by
Lemma~\ref{asm-theta-agree-convergence} every Schwartz-Bruhat function
converges for $(M,\omega)$, which makes $I$ additive on these functions, and
$f\mapsto I(f\otimes g_i)(h)$ is continuous. Hence
$I(\Psi_k)(h)=\sum_iI(s_{k,i}\otimes g_i)(h)\to\sum_iI(f_i\otimes g_i)(h)=I(\Phi)(h)$.
The limits agree: $I(\Phi)(h)=c\,I'(\Phi')(h)$.
\end{proof}

\begin{lemma}\label{asm-theta-agree-transfer}
Let $(B,G,\mathcal X,\omega,\omega')$ be a triple for a theta base $B$ of a CM
field $E$ (Definition~\ref{asm-theta-agree-triple}), and let $V^\circ$,
$V'^\circ$ and $W^\circ$ be lattice spaces such that, with
$G^\circ=G[V^\circ,V'^\circ,W^\circ]$, $(B,G^\circ,\mathcal X,\omega,\omega')$ is a
triple. Let $M$ and $M^\circ$ be the first models of $G$ and $G^\circ$. If
$I(\Phi)\neq0$ for some $\Phi\in\operatorname{PD}(M^\circ)$, then
$I(\Phi)\neq0$ for some $\Phi\in\operatorname{PD}(M)$.
\end{lemma}
\leanref{Development/WeilBound/ThetaLifts/ConeHypotheses/CtSmallCone.lean}{PositivityTheory.ThKG.Leaves.nz_move}
\begin{proof}
The theta integrals of $G$ and $G^\circ$ are the same. Suppose that $I$
vanishes on $\operatorname{PD}(M)$; we show that it vanishes on
$\operatorname{PD}(M^\circ)$. As in the proof of
Lemma~\ref{asm-theta-agree-constants}, Lemma~\ref{asm-theta-agree-convergence}
applies to the Weil pair of $M$, so $I$ is linear on the Schwartz-Bruhat
functions. It suffices to treat $\Phi=f\otimes\bigotimes_v\phi_v$ with
$f\in\operatorname{Pol}(M^\circ)$, the $\phi_v$ Schwartz-Bruhat and
$\phi_v=\mathrm{vac}_v$ for almost all $v$. Since $M$ is regular,
$g=\bigotimes_v\phi_v$ is a Schwartz-Bruhat function on the finite adeles, and
$\Phi=f\otimes g$. For every $f'\in\operatorname{Pol}(M)$ we have
$f'\otimes g\in\operatorname{PD}(M)$, so $I(f'\otimes g)(h)=0$ for every $h$.
By Lemma~\ref{asm-theta-agree-convergence}(2) the map $f'\mapsto I(f'\otimes g)(h)$
is continuous, and by Lemma~\ref{asm-theta-agree-density}
$\operatorname{Pol}(M)$ is dense in the Schwartz space. So
$I(f'\otimes g)(h)=0$ for every Schwartz function $f'$, in particular for
$f'=f$.
\end{proof}

\begin{lemma}\label{asm-theta-agree-key-beta}
There are signs $e_1,e_2,s\in\{\pm1\}$ such that the following holds. Let $B$ be
a theta base of a CM field $E$, $(G,G_1,H,H')$ a $\beta$-configuration
(Definition~\ref{asm-theta-agree-planes}(4)) with $G$ of ranks $(a,b+2)$, and
$R$ and $R_1$ real maps for $G$ and $G_1$. Then there is $C\neq0$ such that for
all matched quadruples $(B,G,\mathcal X,\omega,\omega',R,R_f)$ and
$(B,G_1,\mathcal X,\omega_1,\omega'_1,R_1,R_{f,1})$ such that nondegeneracy
holds below $p+q$ (Definition~\ref{asm-theta-agree-matched}(6)), with
$\lambda=C\sigma_{e_1,e_2,s}(\mathfrak b_H,\mathfrak b_{H'};-k_\beta)L^{S_0}(\nu,k_\beta+1)$:
\begin{enumerate}
\item if $c_1$ is a polynomial constant of the second quadruple,
  $\Phi'\in\mathcal S^\circ(M')$ and $\Phi$ correspond, $(f_\alpha,f_\beta)$ and
  $(f'_\alpha,f'_\beta)$ split $I(\Phi)$ and $I'(\Phi')$ along $H'$, and $m$ is a
  Levi element for $(1,h_1)$ with $h_1\in U_{W_1}$, then
  $\lambda f_\beta(m)=c_1f'_\beta(m)$;
\item there is $\Phi\in\operatorname{PD}(M)$ such that for every pair
  $(f_\alpha,f_\beta)$ that splits $I(\Phi)$ along $H'$ there is a Levi element
  $m$ for some $(1,h_1)$ with $f_\beta(m)\neq0$.
\end{enumerate}
\end{lemma}
\leanref{Development/WeilBound/ThetaLifts/ConeHypotheses/KeyAOf.lean}{PositivityTheory.ThKG.KeyB}
\leanref{Development/WeilBound/ThetaLifts/ConeHypotheses/KeyAOf.lean}{PositivityTheory.ThKG.keyB_of}
\begin{proof}
Take $e_1,e_2,s$ from Lemma~\ref{asm-theta-agree-local-beta}, and let $B$,
$(G,G_1,H,H')$, $R$ and $R_1$ be given. If there are no triples
$(B,G,\mathcal X,\omega,\omega')$ and $(B,G_1,\mathcal X,\omega_1,\omega'_1)$
with a common $\mathcal X$, take $C=1$. Otherwise fix two such triples. They
give Weil pairs for $M$ and $M'$ (Definition~\ref{asm-theta-agree-triple}(4)),
and $M_1$ and $M'_1$ are regular. Lemma~\ref{asm-theta-agree-evaluation}, for
the Weil pair of $M$, the plane $H'$ and $M_1$, gives an evaluation map $e$
for $(M,M_1,1,\iota_{H'},\mathcal J_0(H'))$ that carries $\operatorname{PD}(M)$
onto $\operatorname{PD}(M_1)$. Lemma~\ref{asm-theta-agree-evaluation-cone}, for
the Weil pair of $M'$, the planes $H$ and $H'$ and $M'_1$, gives an evaluation
map $e'$ for $(M',M'_1,\iota_H,\iota_{H'},\mathcal J_\beta(H,H'))$. By
Lemma~\ref{asm-theta-agree-stabilizer} the lattice stabilizer
$K\subset U_{V'}$ is compact, and there are Haar measures $\mu_K$ and $\mu$.
Lemmas~\ref{asm-theta-agree-ct-beta-zero}, \ref{asm-theta-agree-ct-beta-cone}
and~\ref{asm-theta-agree-local-beta} now give constants $C_0$, $C_1$ and $C_2$,
all nonzero. Put $C=C_2C_1/C_0$.

Let the matched quadruples be as in the statement, and put
$\lambda_2=C_2\sigma_{e_1,e_2,s}(\mathfrak b_H,\mathfrak b_{H'};-k_\beta)L^{S_0}(\nu,k_\beta+1)$,
so that $\lambda=\lambda_2C_1/C_0$.

(1) Since $\Phi$ and $\Phi'$ correspond, $\Phi\in\operatorname{PD}(M)$ and
$\Phi'\in\operatorname{PD}(M')\cap\mathcal S^\circ(M')$.
Lemma~\ref{asm-theta-agree-ct-beta-zero} with $h=1$ gives
$f_\beta(m)=C_0I_1(e(\Phi))(h_1)$, and
Lemma~\ref{asm-theta-agree-ct-beta-cone} with $h=1$ gives
$f'_\beta(m)=C_1I'_1(\operatorname{cone}_\beta(\Phi'))(h_1)$. By
Lemma~\ref{asm-theta-agree-local-beta},
$\operatorname{cone}_\beta(\Phi')\in\mathcal S^\circ(M'_1)$ and corresponds to
$\lambda_2e(\Phi)$. As $c_1$ is a polynomial constant of the second
quadruple and $I_1$ is linear,
$\lambda_2I_1(e(\Phi))(h_1)=c_1I'_1(\operatorname{cone}_\beta(\Phi'))(h_1)$.
Hence
$\lambda f_\beta(m)=C_1\lambda_2I_1(e(\Phi))(h_1)=c_1C_1I'_1(\operatorname{cone}_\beta(\Phi'))(h_1)=c_1f'_\beta(m)$.

(2) As $p_1=p$ and $q_1+1=q$, the second quadruple is nondegenerate, so
$I_1(\Phi_1)\neq0$ for some $\Phi_1\in\operatorname{PD}(M_1)$. As $e$ carries
$\operatorname{PD}(M)$ onto $\operatorname{PD}(M_1)$, $\Phi_1=e(\Phi)$ with
$\Phi\in\operatorname{PD}(M)$. Choose $h_1$ with $I_1(\Phi_1)(h_1)\neq0$. Let
$(f_\alpha,f_\beta)$ split $I(\Phi)$ along $H'$. By
Lemma~\ref{asm-theta-agree-levi} there is a Levi element $m$ for $(1,h_1)$, and
Lemma~\ref{asm-theta-agree-ct-beta-zero} gives
$f_\beta(m)=C_0I_1(\Phi_1)(h_1)\neq0$.
\end{proof}

\begin{lemma}\label{asm-theta-agree-key-alpha}
There are signs $e_1,e_2,s\in\{\pm1\}$ such that the following holds. Let $B$ be
a theta base of a CM field $E$, $(G,G_1,H,H')$ an $\alpha$-configuration
(Definition~\ref{asm-theta-agree-planes}(4)) with $G$ of ranks $(b+2,a')$, and
$R$ and $R_1$ real maps for $G$ and $G_1$. Then there is $C\neq0$ such that for
all matched quadruples $(B,G,\mathcal X,\omega,\omega',R,R_f)$ and
$(B,G_1,\mathcal X,\omega_1,\omega'_1,R_1,R_{f,1})$ such that nondegeneracy
holds below $p+q$ (Definition~\ref{asm-theta-agree-matched}(6)), every
polynomial constant $c_1$ of the second quadruple, and
$\kappa=C\sigma_{e_1,e_2,s}(\mathfrak b_H,\mathfrak b_{H'};1-k_\alpha)L^{S_0}(\nu,k_\alpha)c_1$:
\begin{enumerate}
\item for all corresponding $\Phi$ and $\Phi'$ there is $f'_\alpha$ such that
  $(f'_\alpha,0)$ splits $I'(\Phi')$ along $H'$ and $f_\alpha(m)=\kappa f'_\alpha(m)$
  whenever $(f_\alpha,f_\beta)$ splits $I(\Phi)$ along $H'$ and $m$ is a Levi
  element for $(1,h_1)$ with $h_1\in U_{W_1}$;
\item there are corresponding $\Phi$ and $\Phi'$, an $f'_\alpha$ as in (1), and
  a Levi element $m$ for some $(1,h_1)$ with $f'_\alpha(m)\neq0$.
\end{enumerate}
\end{lemma}
\leanref{Development/WeilBound/ThetaLifts/ConeHypotheses/KeyAOf.lean}{PositivityTheory.ThKG.KeyA}
\leanref{Development/WeilBound/ThetaLifts/ConeHypotheses/KeyAOf.lean}{PositivityTheory.ThKG.RelA}
\leanref{Development/WeilBound/ThetaLifts/ConeHypotheses/KeyAOf.lean}{PositivityTheory.ThKG.keyA_of}
\begin{proof}
Take $e_1,e_2,s$ from Lemma~\ref{asm-theta-agree-local-alpha}, and let $B$,
$(G,G_1,H,H')$, $R$ and $R_1$ be given. If there are no triples
$(B,G,\mathcal X,\omega,\omega')$ and $(B,G_1,\mathcal X,\omega_1,\omega'_1)$
with a common $\mathcal X$, take $C=1$. Otherwise fix two such triples.
Lemma~\ref{asm-theta-agree-evaluation}, for the Weil pair of $M'$, the plane
$H'$ and $M'_1$, gives an evaluation map $e'$ for
$(M',M'_1,1,\iota_{H'},\mathcal J_0(H'))$ that carries $\operatorname{PD}(M')$
onto $\operatorname{PD}(M'_1)$. Lemma~\ref{asm-theta-agree-evaluation-cone}, for
the Weil pair of $M$, the planes $H$ and $H'$ and $M_1$, gives an evaluation
map $e$ for $(M,M_1,\iota_H,\iota_{H'},\mathcal J_\alpha(H,H'))$. By
Lemma~\ref{asm-theta-agree-stabilizer} the lattice stabilizer $K\subset U_V$ is
compact, and there are Haar measures $\mu_K$ and $\mu$.
Lemmas~\ref{asm-theta-agree-ct-alpha-zero}, \ref{asm-theta-agree-ct-alpha-cone}
and~\ref{asm-theta-agree-local-alpha} now give constants $C_0$, $C_1$ and $C_2$,
all nonzero. Put $C=C_2C_1/C_0$.

Let the matched quadruples and $c_1$ be as in the statement, and put
$\lambda_2=C_2\sigma_{e_1,e_2,s}(\mathfrak b_H,\mathfrak b_{H'};1-k_\alpha)L^{S_0}(\nu,k_\alpha)$,
so that $\kappa=\lambda_2c_1C_1/C_0$. We first show: for corresponding $\Phi$
and $\Phi'$ there is $f'_\alpha$ as in (1) with, in addition,
$f'_\alpha(m)=C_0I'_1(e'(\Phi'))(h_1)$ for every Levi element $m$ for $(1,h_1)$.
As $\Phi'\in\operatorname{PD}(M')$, Lemma~\ref{asm-theta-agree-ct-alpha-zero}
with $h=1$ gives $f'_\alpha$ such that $(f'_\alpha,0)$ splits $I'(\Phi')$ and
$f'_\alpha(m)=C_0I'_1(e'(\Phi'))(h_1)$. Let $(f_\alpha,f_\beta)$ split $I(\Phi)$
and let $m$ be a Levi element for $(1,h_1)$.
Lemma~\ref{asm-theta-agree-ct-alpha-cone} with $h=1$ gives
$f_\alpha(m)=C_1I_1(\operatorname{cone}_\alpha(\Phi))(h_1)$. By
Lemma~\ref{asm-theta-agree-local-alpha}, $\operatorname{cone}_\alpha(\Phi)$ and
$\lambda_2e'(\Phi')$ correspond for the second quadruple. As $q_1=0$,
Lemma~\ref{asm-theta-agree-anisotropic} applies to it and $c_1$:
$I_1(\operatorname{cone}_\alpha(\Phi))(h_1)=c_1\lambda_2I'_1(e'(\Phi'))(h_1)$.
Hence $f_\alpha(m)=C_1c_1\lambda_2f'_\alpha(m)/C_0=\kappa f'_\alpha(m)$. This
proves (1).

(2) As $p_1+1=p$ and $q=q_1=0$, the second quadruple is nondegenerate, so
$I'_1(\Phi'_1)\neq0$ for some
$\Phi'_1\in\operatorname{PD}(M'_1)\cap\mathcal S^\circ(M'_1)$. As $e'$ carries
$\operatorname{PD}(M')$ onto $\operatorname{PD}(M'_1)$, $\Phi'_1=e'(\Phi')$ with
$\Phi'\in\operatorname{PD}(M')$, and by
Lemma~\ref{asm-theta-agree-partners}(2) some $\Phi$ corresponds to $\Phi'$.
Let $f'_\alpha$ be as just shown. Choose $h_1$ with $I'_1(\Phi'_1)(h_1)\neq0$
and, by Lemma~\ref{asm-theta-agree-levi}, a Levi element $m$ for $(1,h_1)$.
Then $f'_\alpha(m)=C_0I'_1(\Phi'_1)(h_1)\neq0$.
\end{proof}

\begin{lemma}\label{asm-theta-agree-claim}
Assume the meromorphic continuation of Siegel Eisenstein series in the form of
Theorem~\ref{asm-theta-agree-eisenstein} and the Siegel-Weil formula in the
form of Theorem~\ref{asm-theta-agree-siegel-weil}. Then every matched quadruple
$(B,G,\mathcal X,\omega,\omega',R,R_f)$, for every theta base $B$ of every CM
field $E$, is nondegenerate (Definition~\ref{asm-theta-agree-matched}(6)).
\end{lemma}
\leanref{Development/WeilBound/ThetaLifts/ConeHypotheses/CtSmallCone.lean}{PositivityTheory.ThKG.Leaves.const_of}
\leanref{Development/WeilBound/ThetaLifts/ConeHypotheses/CtSmallCone.lean}{PositivityTheory.ThKG.Leaves.nzV}
\leanref{Development/WeilBound/ThetaLifts/ConeHypotheses/CtSmallCone.lean}{PositivityTheory.ThKG.Leaves.claim_of}
\leanref{Development/WeilBound/ThetaLifts/ConeHypotheses/CtSmallCone.lean}{PositivityTheory.ThKG.Leaves.below}
\leanref{Development/WeilBound/ThetaLifts/ConeHypotheses/CtSmallCone.lean}{PositivityTheory.ThKG.Leaves.claim}
\begin{proof}
It suffices to show that a matched quadruple
$Q=(B,G,\mathcal X,\omega,\omega',R,R_f)$ is nondegenerate whenever
nondegeneracy holds below $p+q$. Then, by induction on $n$, nondegeneracy holds
below every $n$, the case $n=0$ being empty, and every $Q$ is nondegenerate.
So let $Q$ be such a quadruple. By Lemmas~\ref{asm-theta-agree-real}
and~\ref{asm-theta-agree-finite}, every triple has matched real and finite
maps.

\emph{Step 1: $Q$ has a nonzero polynomial constant.} If $p=q=0$, this is
Lemma~\ref{asm-theta-agree-definite}.

Let $q=0$ and $p\geq1$. Properness gives $2(p+q)+1<\lfloor(a-a')/2\rfloor$, so
$a\geq2$; write $a=b+2$. Lemma~\ref{asm-theta-agree-tower-alpha} gives a
re-chosen datum $G^\circ$ with a triple $(B,G^\circ,\mathcal X,\omega,\omega')$,
an $\alpha$-configuration $(G^\circ,G_1,H,H')$ and a triple
$(B,G_1,\mathcal X,\omega_1,\omega'_1)$. Choose matched maps for them, giving
matched quadruples $Q^\circ$ and $Q_1$ with real maps $R_2$ and $R_1$. By
Lemma~\ref{asm-theta-agree-rechoice} it suffices to find a nonzero polynomial
constant of $Q^\circ$. Lemma~\ref{asm-theta-agree-key-alpha}, for
$(G^\circ,G_1,H,H')$, $R_2$ and $R_1$, gives signs $e_1,e_2,s$ and $C\neq0$. As
$p_1+1=p$ and $q_1=0$, $Q_1$ is nondegenerate, with a polynomial constant
$c_1\neq0$; and $Q^\circ$ has the same $p$ and $q$ as $Q$. Put
$c=C\sigma_{e_1,e_2,s}(\mathfrak b_H,\mathfrak b_{H'};1-k_\alpha)L^{S_0}(\nu,k_\alpha)c_1$.
By Lemma~\ref{asm-theta-agree-factors}, $c\neq0$, since
$k_\alpha=\lfloor(a-a')/2\rfloor>2(p+q)+1\geq3$. We show that $c$ is a
polynomial constant of $Q^\circ$, by Lemma~\ref{asm-theta-agree-end-alpha} for
$Q^\circ$ and $H'$. Let $\Phi'\in\mathcal S^\circ(M')$ and $\Phi$ correspond for
$Q^\circ$, let $(f_\alpha,f_\beta)$ and $(f'_\alpha,f'_\beta)$ split $I(\Phi)$
and $I'(\Phi')$ along $H'$, and let $m$ be a Levi element for $(1,h_1)$.
Lemma~\ref{asm-theta-agree-key-alpha}(1) gives $f''_\alpha$ such that
$(f''_\alpha,0)$ splits $I'(\Phi')$ and $f_\alpha(m)=c\,f''_\alpha(m)$. As
$\Phi'\in\operatorname{PD}(M')\cap\mathcal S^\circ(M')$,
Lemma~\ref{asm-theta-agree-split} for the triple
$(B,G^\circ,\mathcal X,\omega,\omega')$ shows that $I'(\Phi')$ has only one
splitting along $H'$, so $f'_\alpha=f''_\alpha$ and
$f_\alpha(m)=c\,f'_\alpha(m)$.

Let $q\geq1$. Properness gives $a'>2(p+q)+1\geq3$; write $a'=b+2$.
Lemma~\ref{asm-theta-agree-tower-beta} gives $G^\circ$, a $\beta$-configuration
$(G^\circ,G_1,H,H')$ and triples as before, and matched quadruples $Q^\circ$
and $Q_1$. By Lemma~\ref{asm-theta-agree-rechoice} it suffices to treat
$Q^\circ$. Lemma~\ref{asm-theta-agree-key-beta} gives signs and $C\neq0$. As
$p_1=p$ and $q_1+1=q$, $Q_1$ has a polynomial constant $c_1\neq0$. Put
$\lambda=C\sigma_{e_1,e_2,s}(\mathfrak b_H,\mathfrak b_{H'};-k_\beta)L^{S_0}(\nu,k_\beta+1)$,
which is nonzero by Lemma~\ref{asm-theta-agree-factors}, since
$k_\beta=\lfloor(a-a')/2\rfloor\geq1$, and $c=c_1/\lambda\neq0$. Whenever
$\Phi'\in\mathcal S^\circ(M')$ and $\Phi$ correspond for $Q^\circ$,
$(f_\alpha,f_\beta)$ and $(f'_\alpha,f'_\beta)$ split $I(\Phi)$ and
$I'(\Phi')$ along $H'$, and $m$ is a Levi element for $(1,h_1)$,
Lemma~\ref{asm-theta-agree-key-beta}(1) gives
$\lambda f_\beta(m)=c_1f'_\beta(m)$, that is, $f_\beta(m)=c\,f'_\beta(m)$. By
Lemma~\ref{asm-theta-agree-end-beta}, $c$ is a polynomial constant of
$Q^\circ$.

\emph{Step 2: if $q\geq1$, then $I(\Phi)\neq0$ for some
$\Phi\in\operatorname{PD}(M)$.} With $G^\circ$, $Q^\circ$ and $Q_1$ as in the
last case of Step 1, Lemma~\ref{asm-theta-agree-key-beta}(2) gives
$\Phi\in\operatorname{PD}(M^\circ)$ such that every pair $(f_\alpha,f_\beta)$
that splits $I(\Phi)$ along $H'$ has $f_\beta(m)\neq0$ for some Levi element
$m$. As $(0,0)$ splits $0$, $I(\Phi)\neq0$. Lemma~\ref{asm-theta-agree-transfer}
gives $\Phi\in\operatorname{PD}(M)$ with $I(\Phi)\neq0$.

\emph{Step 3: conclusion.} Let $c\neq0$ be a polynomial constant of $Q$, by
Step 1. If $q=0$, Lemma~\ref{asm-theta-agree-nonvanishing} gives
$\Phi'\in\operatorname{PD}(M')\cap\mathcal S^\circ(M')$ and $h$ with
$I'(\Phi')(h)\neq0$, and Lemma~\ref{asm-theta-agree-partners}(2) gives $\Phi$
corresponding to $\Phi'$; then $\Phi\in\operatorname{PD}(M)$ and
$I(\Phi)(h)=c\,I'(\Phi')(h)\neq0$. If $q\geq1$, Step 2 gives
$\Phi\in\operatorname{PD}(M)$ with $I(\Phi)\neq0$, and
Lemma~\ref{asm-theta-agree-partners}(1) gives $\Phi'\in\mathcal S^\circ(M')$
corresponding to $\Phi$; then $\Phi'\in\operatorname{PD}(M')$, and
$I'(\Phi')\neq0$ because $I(\Phi)=c\,I'(\Phi')$. So $Q$ is nondegenerate.
\end{proof}

\begin{proposition}\label{asm-theta-agree-derivation}
Assume the meromorphic continuation of Siegel Eisenstein series in the form
of Theorem~\ref{asm-theta-agree-eisenstein} and the Siegel-Weil formula in
the form of Theorem~\ref{asm-theta-agree-siegel-weil}. Then the agreement of
global and local theta integrals and the local matching hold
(Definition~\ref{th-theta-hypotheses}).
\end{proposition}
\leanref{Development/WeilBound/ThetaLifts/Reductions/FiniteClosure/TheoremKOfMod.lean}{PositivityTheory.ThKR.theoremK_head2}
\leanref{Development/WeilBound/ThetaLifts/Reductions/FiniteClosure/TheoremKOfMod.lean}{PositivityTheory.ThKR.theoremK_head2_of_mod}
\leanref{Development/WeilBound/ThetaLifts/SphericalLattices/SphericalClosure/CloseSph.lean}{PositivityTheory.Sph.Tr.closeSph}
\leanref{Development/WeilBound/ThetaLifts/AgreeOfLeaves.lean}{PositivityTheory.ThKG.agree_of_leaves}
\leanref{Development/AbelianVarietyTheory/Definitions/Divisors/MatchingOfParts.lean}{GlobalTheory.ThK.agree_of_parts}
\leanref{Development/AbelianVarietyTheory/Definitions/Divisors/MatchingOfParts.lean}{GlobalTheory.ThK.matching_of_parts}
\leanref{Development/WeilBound/ThetaLifts/ConeHypotheses/CtSmallCone.lean}{PositivityTheory.ThKG.Leaves}
\leanref{Development/WeilBound/ThetaLifts/ConeHypotheses/CtSmallCone.lean}{PositivityTheory.ThKG.Leaves.kConst}
\leanref{Development/WeilBound/ThetaLifts/ConeHypotheses/CtSmallCone.lean}{PositivityTheory.ThKG.Leaves.kBeta}
\leanref{Development/WeilBound/ThetaLifts/ConeHypotheses/CtSmallCone.lean}{PositivityTheory.ThKG.Leaves.kAlpha}
\begin{proof}
\emph{Step 1: the parts.} In the terms of
Definitions~\ref{asm-theta-agree-models} to~\ref{asm-theta-agree-parts}, the
agreement of global and local theta integrals and the local matching
(Definition~\ref{th-theta-hypotheses}) are the following statements, over all
CM fields $E$ and theta bases $B$ for $E$:
\begin{enumerate}
\item[(a)] every triple has decay;
\item[(b)] every matched quadruple has one constant;
\item[(c)] every matched quadruple describes the coinvariants;
\item[(d)] the $\beta$-step holds for some signs $e_1,e_2,s$, and the
  $\alpha$-step holds for some signs;
\item[(e)] every triple has real maps matched to it, and
\item[(f)] every triple has finite maps matched to it.
\end{enumerate}
Parts (a) to (d) make up the agreement and parts (e) and (f) the local
matching.
% Lean: agree_of_leaves L hd hc is agree_of_parts hd L.kConst hc L.kBeta L.kAlpha and
% matching_of_parts L.kReal L.kFin; the parts are KDecay (a), KConst (b), KCoinv (c),
% KBeta and KAlpha (d), KReal (e) and KFin (f), of MatchingOfParts.lean. closeSph builds
% the structure Leaves of twenty-six properties, which are
% Lemmas asm-theta-agree-real to asm-theta-agree-end-beta other than -decay and
% -coinvariants; their order in Leaves is evZero, evCone, stabEx, towerA, towerB, facNe,
% matchEx, leviEx, polDense, weilConv, soCont, splitEx, ctBigZero, ctSmallZero, ctBigCone,
% ctSmallCone, anAll, polRe, nonZeroAn, polZero, locA, locB, endA, endB, kReal, kFin.

\emph{Step 2: parts (a), (c), (e) and (f).} These are
Lemmas~\ref{asm-theta-agree-decay}, \ref{asm-theta-agree-coinvariants},
\ref{asm-theta-agree-real} and~\ref{asm-theta-agree-finite}.
% Lean: hd = ThKEst.kDecay hSiegelSet thetaBasic and hc = Close.kCoinv DictSt are the
% hypotheses of agree_of_leaves; kReal and kFin are fields of Leaves.

\emph{Step 3: part (b).} Let $Q$ be a matched quadruple. By
Lemma~\ref{asm-theta-agree-claim} it is nondegenerate: it has a polynomial
constant $c\neq0$, and $I'(\Phi')(h)\neq0$ for some
$\Phi'\in\operatorname{PD}(M')\cap\mathcal S^\circ(M')$ and some $h$. By
Lemma~\ref{asm-theta-agree-partners}(2) some $\Phi$ corresponds to $\Phi'$. By
Lemma~\ref{asm-theta-agree-constants}, $c$ is a constant of $Q$, and every
constant $c'$ of $Q$ is a polynomial constant, so
$c'I'(\Phi')(h)=I(\Phi)(h)=c\,I'(\Phi')(h)$ and $c'=c$. So $c$ is the only
constant, it is nonzero, and $I'(\Phi')\neq0$ with
$\Phi'\in\mathcal S^\circ(M')$.

\emph{Step 4: part (d).} Let $e_1,e_2,s$ be the signs of
Lemma~\ref{asm-theta-agree-key-beta}. Let $(G,G_1,H,H')$ be a
$\beta$-configuration for a theta base $B$, $R$ and $R_1$ real maps, and $C$
the constant of Lemma~\ref{asm-theta-agree-key-beta}. Let
$(B,G,\mathcal X,\omega,\omega',R,R_f)$ and
$(B,G_1,\mathcal X,\omega_1,\omega'_1,R_1,R_{f,1})$ be matched quadruples with
constants $c$ and $c_1$, and put
$\sigma=\sigma_{e_1,e_2,s}(\mathfrak b_H,\mathfrak b_{H'};-k_\beta)$ and
$L=L^{S_0}(\nu,k_\beta+1)$. By Lemma~\ref{asm-theta-agree-claim},
nondegeneracy holds below $p+q$. Let $\Phi\in\operatorname{PD}(M)$ be as in
Lemma~\ref{asm-theta-agree-key-beta}(2). By
Lemma~\ref{asm-theta-agree-partners}(1) some $\Phi'\in\mathcal S^\circ(M')$
corresponds to $\Phi$, and by Lemma~\ref{asm-theta-agree-split} a pair
$(f'_\alpha,f'_\beta)$ splits $I'(\Phi')$ along $H'$. By
Lemma~\ref{asm-theta-agree-constants}, $c$ is a polynomial constant, so
$I(\Phi)=c\,I'(\Phi')$, and $(cf'_\alpha,cf'_\beta)$ splits $I(\Phi)$. By the
choice of $\Phi$ there is a Levi element $m$ with $cf'_\beta(m)\neq0$. By
Lemma~\ref{asm-theta-agree-constants}, $c_1$ is a polynomial constant of the
second quadruple, and Lemma~\ref{asm-theta-agree-key-beta}(1) gives
$C\sigma L\,cf'_\beta(m)=c_1f'_\beta(m)$. As $f'_\beta(m)\neq0$,
$c_1=C\sigma Lc$. By Lemma~\ref{asm-theta-agree-factors}, $C\sigma\neq0$, so
$cL=(C\sigma)^{-1}c_1$. This is the $\beta$-step.

Let $e_1,e_2,s$ be the signs of Lemma~\ref{asm-theta-agree-key-alpha}, let
$(G,G_1,H,H')$ be an $\alpha$-configuration with real maps $R$ and $R_1$, and
$C$ the constant of Lemma~\ref{asm-theta-agree-key-alpha}. Let matched
quadruples with constants $c$ and $c_1$ be given. By
Lemma~\ref{asm-theta-agree-constants}, $c$ and $c_1$ are polynomial constants;
put
$\kappa=C\sigma_{e_1,e_2,s}(\mathfrak b_H,\mathfrak b_{H'};1-k_\alpha)L^{S_0}(\nu,k_\alpha)c_1$.
By Lemma~\ref{asm-theta-agree-claim}, nondegeneracy holds below $p+q$, and
Lemma~\ref{asm-theta-agree-key-alpha}(2) gives corresponding $\Phi$ and
$\Phi'$, an $f'_\alpha$ such that $(f'_\alpha,0)$ splits $I'(\Phi')$ and
$f_\alpha(m)=\kappa f'_\alpha(m)$ for every pair $(f_\alpha,f_\beta)$ splitting
$I(\Phi)$ and every Levi element $m$, and a Levi element $m$ with
$f'_\alpha(m)\neq0$. As $q=0$, Lemma~\ref{asm-theta-agree-anisotropic} gives
$I(\Phi)=c\,I'(\Phi')$, so $(cf'_\alpha,0)$ splits $I(\Phi)$. Hence
$cf'_\alpha(m)=\kappa f'_\alpha(m)$, and $c=\kappa$. This is the
$\alpha$-step.
% Lean: Leaves.kConst, Leaves.kBeta and Leaves.kAlpha, through Leaves.claim (the
% induction Leaves.below) and Leaves.polConst_of, isConst_of.
\end{proof}

\begin{theorem}\label{asm-every-character-literature}
Assume that Galois representations exist relative to the maps satisfying the
local Langlands laws (Definitions~\ref{gal-galois-representations}
and~\ref{gal-llc-laws}); that for every number field and every finite place
such a map exists; and the weakening of Vandiver's conjecture
(Definition~\ref{sha-vandiver-weakening}). Then for every Dirichlet character
$\chi$ modulo $N\ge1$ and every $s$ with $\operatorname{Re}s>1/2$,
$L(s,\chi)\ne0$ unless $\chi$ is trivial and $s=1$.
\end{theorem}
\begin{proof}
We check the hypotheses of Theorem~\ref{asm-every-character-theta}. The
published results on unitary groups hold by Theorem~\ref{lit-bundle-unitary},
whose proof uses the first two assumptions (Remark~\ref{lit-open-inputs}): the
Galois representations enter twice, once through Caraiani's local-global
compatibility and once as the Galois representations attached to polarized
cuspidal representations. The published local inputs of the theta
correspondence hold by Theorem~\ref{lit-bundle-theta}. The agreement of theta
integrals and the local matching (Definition~\ref{th-theta-hypotheses}) hold
by Proposition~\ref{asm-theta-agree-derivation}, from the meromorphic
continuation of Siegel Eisenstein series
(Theorem~\ref{asm-theta-agree-eisenstein}) and the Siegel-Weil formula
(Theorem~\ref{asm-theta-agree-siegel-weil}).
% Lean: theta_integrals_agree_and_local_matching_of_literature O =
% PositivityTheory.ThKR.theoremK_head2 (eis_bl_of_literature O)
% (siegel_weil_std_of_literature O); the two inputs are the theorems
% PositivityTheory.RgEis.Hd.eisBL and SiegelWeilFormula.Fin.siegelWeilStd, not fields of O.
The covolume bound holds by
Theorem~\ref{cv-covolume-bound-from-literature}, all of whose hypotheses are
proved results. The local Langlands maps exist and the weakening of Vandiver's
conjecture holds by assumption.
\end{proof}
\leanref{Development/AutomorphicArithmetic/ThetaCorrespondence/ThetaPairLifts/ThirdPlaceSeries/EveryCharacterOfLiteratureOfVWin.lean}{UpperTheory.LiteratureTheorems.every_character_of_literature_of_vWin}

\begin{theorem}\label{asm-every-character}
Suppose that Galois representations exist relative to the maps satisfying the
local Langlands laws (Definitions~\ref{gal-galois-representations}
and~\ref{gal-llc-laws}). Then for every Dirichlet character $\chi$ modulo
$N\ge1$ and every $s$ with $\operatorname{Re}s>1/2$, $L(s,\chi)\ne0$ unless
$\chi$ is trivial and $s=1$.
\end{theorem}
\begin{proof}
Maps satisfying the local Langlands laws exist for every number field and
every finite place, by Theorem~\ref{gal-llc-exists}. The Galois
representations give the weakening of Vandiver's conjecture, by
Theorem~\ref{sha-weakening-from-galois-representations}.
Theorem~\ref{asm-every-character-literature} applies, with the assumed Galois
representations in both of their roles.
\end{proof}
\leanref{Development/AutomorphicArithmetic/CharacterBoxes/MainConclusion.lean}{UpperTheory.FinalTheorems.everyCharacter_of_row}

The hypothesis of Theorem~\ref{asm-every-character} is
Theorem~\ref{gal-galois-representations-any-rc}. So the half-plane statement,
Theorem~\ref{mt-half-plane}, holds, and with it the main theorems of
Chapter~\ref{chapter-main-theorems}.

\section{Parameters and the flow of bounds}\label{section-asm-parameters}

\begin{remark}\label{asm-flow}
Fix a member $D$ of a tied family at a scale $x$, with base point $u_0$, and a
word $u$. The comparison inequality of Chapter~\ref{chapter-inert-prime-balance}
bounds the inert-prime sum $\mathrm{ASK}(u_0)$ by
$\mathrm{Main}(u_0)+\mathrm{Err}(u_0)$. Its hypotheses come from the following
bounds, each holding eventually along the designed families.
\begin{enumerate}
\item The gain link (Chapter~\ref{chapter-period-divisibility}) bounds, at
  every place over a prime $p\le X''(n)$ inert in the host field, the class
  sums $\mathrm{VM}_{j,t}(p)$ by the divisibility gains $\mathrm{gain}_j(w)$.
  Summed over $p$, with the chosen constants and the small shares of the
  weight vectors, this gives the loss bound
  $\mathrm{ASK}-G_a\le C|T|mbx^{e_L\varepsilon}$ with $e_L=0$.
\item The gains raise the slope of the division lattice:
  $\mu(N_u)\ge\mu(M_u)+G_a+F_o-E_u-\mathrm{Lo}_C(u)$, and the slope bound
  $\mu(M_u)\ge-C|T|aL_a(1+m/k)$ holds (Definition~\ref{asm-member-bounds}
  and Theorem~\ref{asm-global-link}).
\item The frame inequality of Chapter~\ref{chapter-hermitian-slopes} passes
  from the slopes of the lattices $N_u$ to a lower bound for the slope
  $\mu(CC)$ of the dual of the glued lattice. Its almost tight frame comes from
  the orthogonality link, with errors at most $e^{-x^{\varepsilon/8}}$, and
  the irreducibility link, with row deficit at most
  $e^{-x^{\varepsilon_z/7}}$ (Chapter~\ref{chapter-integral-structure}).
\item The division link (Chapter~\ref{chapter-integral-structure}), with the
  covolumes of Chapter~\ref{chapter-covolumes}, gives the upper bound
  $\mu(CC)+h_u-F_o\le\mathrm{MAIN}(u)+C|T|aL_a+Cm(B_0+N_c)x^{2\varepsilon}$.
\item The division link needs the division bound at twisted places, which
  comes from the Tate-Shafarevich bound of
  Chapter~\ref{chapter-tate-shafarevich} through
  Lemma~\ref{asm-division-from-sha}. This is where $e^{x^{\varepsilon}/2}$
  enters: it is at most $B_0=fme^{x^{\varepsilon}}x^{\varepsilon}$.
\end{enumerate}
The lower bound (3) and the upper bound (4) for $\mu(CC)$ give
$G_a\le\mathrm{MAIN}+\text{error}$, and with (1),
$\mathrm{ASK}\le\mathrm{Main}+\mathrm{Err}$. The balance identity of
Chapter~\ref{chapter-inert-prime-balance} rewrites $\mathrm{ASK}-\mathrm{MAIN}$
as a combination of the sums $G_\psi(a/w)$, and the host conditions give the
window bound
\[
 G_{K_\psi}(x)\le x^{1/2+(e_D+e_L+7)\varepsilon}=x^{1/2+9\varepsilon}.
\]
\end{remark}

\begin{remark}\label{asm-parameters}
The parameters are chosen in the following order. Let $\psi$ be primitive and
$\operatorname{Re}s>1/2$.
\begin{enumerate}
\item Choose $\varepsilon\in(0,1/52)$ below the threshold $\varepsilon_1$ of
  Definition~\ref{asm-tied}, with $9\varepsilon$ at most half of
  $\min(\operatorname{Re}s-1/2,1/2)$. The exponents are $e_D=2$ and $e_L=0$,
  from Theorem~\ref{asm-global-link}.
\item The scales $x$ fill windows $[Y,Y^{3/2}]$ for arbitrarily large $Y$
  (Chapter~\ref{chapter-host-fields}). The Tate-Shafarevich bound of
  Chapter~\ref{chapter-tate-shafarevich} is available on windows $[Y,Y^2]$
  with $Y\ge X^{1/J}$, for twists $1\le k<x^{401}$.
\item At a scale $x$: $L=x^{\varepsilon}$, $b=e^{x^{\varepsilon}}$,
  $a=xe^{x^{\varepsilon}}$, the thin scale $X'=\exp(x^{\varepsilon/3})$ and
  $X''(n)=xnX'$ (Definition~\ref{pd-size-data}).
\item The ranks are adapted to the scale (Chapter~\ref{chapter-host-fields}):
  $x^{2+8\varepsilon}\le n\le2x^{2+8\varepsilon}$ for the $2$-power part of
  the Galois group, $f=2n$, $\sqrt{x^{1-\varepsilon/8}}\le m\le
  2\sqrt{x^{1-\varepsilon/8}}$ lifts, and $x^{9/2+10\varepsilon}\le k\le
  2x^{9/2+10\varepsilon}$ for the first Hermitian space.
\item The auxiliary primes satisfy $\log q\approx Y^{\varepsilon/5}$ and
  $\log d\approx Y^{\varepsilon/5}\log^2Y$; the top classes satisfy
  $\delta_1n/2\le|T|\le2\delta_1n$ with $\delta_1=x^{-\varepsilon/8}$, and
  $|T_2|\ge x^{-\varepsilon}n$.
\item The exponents satisfy $a\le a_t\le a+a/(kx^{13})$ and $b_t\ge b'$;
  the auxiliary primes of the words satisfy
  $b^{m/2}\le T_u\le e^{C_{T_u}mL}$; and the stabilizer denominators divide
  $M(2fN_m)$ (Definition~\ref{asm-designed}).
\item The chosen constants are $C_0=3$, so that
  $1/3\le\dim M_u\cdot\Gamma^2/D^{*}\le3$ with $\Gamma=\exp(C_\Gamma m^5fL)$; the
  thin-scale bound $|g_\chi(y)|\le y\exp(-c_6x^{\varepsilon/8})$ for
  $y\ge X'/3$; and the bounds $|g_\chi(x)|\le x^{1/2+\varepsilon}$ and
  $|\frac{L'}{L}(1,\chi)|\le C_5\log x$ for the characters twisted by $q$ and
  $d$.
\item The middle and last constants of the family, among them the height
  constant $C_H$, are taken large for the threshold of the links
  (Definition~\ref{pd-family}). The tie then holds for all $x\ge x_0$ in
  the set of scales.
\end{enumerate}
\end{remark}
